\documentclass{book}
%\usepackage[letterpaper, portrait, margin=1cm]{geometry}
%\usepackage[letterpaper, bindingoffset=0.2in, left=1in,right=1in,top=.5in,bottom=.5in,footskip=.25in,marginparwidth=5em]{geometry}
\usepackage[letterpaper, left=1in,right=1in,top=.5in,bottom=.5in,footskip=.25in,marginparwidth=1cm]{geometry}
% ---------------------
% mini-table-of-content
% ---------------------
\usepackage{minitoc}
\setcounter{minitocdepth}{1}
\setlength{\mtcindent}{24pt}
\setcounter{secnumdepth}{-2}
%\renewcommand{\mtcfont}{\small\rm}
%\renewcommand{\mtcSfont}{\small\bf}
%\usepackage{setspace}
%\usepackage{tocloft}
%\setlength\cftparskip{-1.2pt}
%\setlength\cftbeforesecskip{1.3pt}
%\setlength\cftaftertoctitleskip{2pt}
%\renewcommand{\cftsecafterpnum}{\hspace*{02.0em}}
%\renewcommand{\cftsubsecafterpnum}{\hspace*{02.0em}}

% ---------------------------
% Chinese Characters Packages
% ---------------------------
\usepackage{fontspec} 
\usepackage{xeCJK}
\setmainfont{Times}
\setCJKmainfont{BiauKai}
\newfontfamily\sblgoodhebrew{SBL BibLit}[Script=Hebrew,Contextuals=Alternate]
\newfontfamily\sblgoodgreek{SBL BibLit}[Script=Greek,Contextuals=Alternate]

\usepackage{ifpdf,cite,algorithmic,url,tikz}
\usepackage[cmex10]{amsmath}

% ---------------------------
% Hebrew Characters Packages
% ---------------------------
\usepackage{polyglossia}
\setmainfont{Times New Roman}

% -------
% General
% -------
\usepackage{multicol}
\usepackage{multirow}
\usepackage{color,colortbl}
\usepackage{xparse}
\usepackage{pbox}
\usepackage{stackengine}
\usepackage{titlesec}% http://ctan.org/pkg/titlesec
\usepackage{tabularx}
\usepackage{xltabular}
\usepackage{titlesec}
\usepackage{makecell}
\newcommand{\sectionbreak}{\clearpage}

\author{
  Editor, Michael Chan\\
  \texttt{michaelchan\_wahyan@yahoo.com.hk}
}
\usepackage{tocloft}

\usepackage{hyperref}
\hypersetup{
    colorlinks=true, % set true if you want colored links
    linktoc   =all , % set to all if you want both sections and subsections linked
    linkcolor =blue, % choose some color if you want links to stand out
}

% ----------
% Afterword
% ----------
\usepackage{marginnote}
\usepackage{sectsty}
\usepackage{ragged2e}
\usepackage{lineno}
\usepackage{xcolor}
\usepackage{paracol}

\begin{document}

\clearpage
%% temporary titles
% command to provide stretchy vertical space in proportion
\newcommand\nbvspace[1][3]{\vspace*{\stretch{#1}}}
% allow some slack to avoid under/overfull boxes
\newcommand\nbstretchyspace{\spaceskip0.5em plus 0.25em minus 0.25em}
% To improve spacing on titlepages
\newcommand{\nbtitlestretch}{\spaceskip0.6em}
\pagestyle{empty}
\begin{center}
\bfseries
\nbvspace[1]
\Huge
{%\nbtitlestretch
\Large
\textbf{旺角浸信會 粵語講道逐字稿 2020-present \\
       Youtube Channel: MongKok Baptist Church
       }}

\nbvspace[1]

{\large
Editor: Michael\\
\texttt{michaelchan\_wahyan@yahoo.com.hk}
}

\nbvspace[1]

{\large
Revision: \texttt{v1.3}\\
Last Update: \today
}


\vfill
\begin{tikzpicture}
    %remove comment for OT cover%\node (0,0) [opacity=0.03]{\includegraphics[width=15cm]{../bible_out/ot_frontcover.png}} ;
    %remove comment for NT cover%\node (0,0) [opacity=0.03]{\includegraphics[width=15cm]{../bible_out/christ_on_cross.png}} ;
    %remove comment for Bible cover%\node (0,0) [xshift=0.8cm, yshift=+2cm, opacity=0.03]{\includegraphics[width=10cm]{./christ_on_cross.png}} ;
    %remove comment for Bible cover%\node (0,0) [              yshift=-2cm, opacity=0.03]{\includegraphics[width=14cm]{./ot_frontcover.png}} ;
\end{tikzpicture}
\vfill

\end{center}

\newpage

\setcounter{tocdepth}{0}
\dominitoc
\begin{multicols}{3}
\addtocontents{toc}{\protect\hypertarget{toc}{}}
\tableofcontents
\end{multicols}

\large
%\twocolumn

% the color definition syntax is as follow:
% \definecolor{name}{system}{definition}
% example: a mono-channel color can be defined as
%          \definecolor{Gray}{gray}{0.9}
% example: an rgb-3-channel color can be defined as
%          \definecolor{LightCyan}{rgb}{0.88,1,1}
%          \definecolor{pink}{rgb}{0.68,0,0.68}

\definecolor{CUV1LightRed}{rgb}{1,0.75,0.75}     % for CUV1
\definecolor{LZZVLightGray}{rgb}{0.9,0.9,0.9}    % for LZZ
\definecolor{KJVVLightGreen}{rgb}{0.75,1,0.85}   % for KJV
\definecolor{CUV2LightYellow}{rgb}{1,1,0.75}     % for CUV2
\definecolor{CNVVLightBrown}{rgb}{1,0.85,0.7}    % for CNV
\definecolor{NRSVLightBlue}{rgb}{0.75,1,1}       % for NRSV
\definecolor{WENLLightPurple}{rgb}{0.95,0.85,0.9}% for WENL
\definecolor{TCV19PaleGreen}{rgb}{0.85,1,0.95}   % for TCV19
\definecolor{MSGVLightWhite}{rgb}{0.98,0.98,0.98}% for MSGV
\definecolor{NETSLightRed}{rgb}{1,0.75,0.75}     % for NETS
\definecolor{JPS1917LightYellow}{rgb}{1,1,0.75}  % for JPS1917
\definecolor{SBLGNTPaleRed}{rgb}{1,0.85,0.80}    % for SBLGNT

\section{目錄\small{(順時)}}
\label{sec:index_chronic}
{ \scriptsize


\begin{xltabular}{\textwidth}{|p{0.15\textwidth} p{0.6\textwidth}|p{0.07\textwidth} p{0.1\textwidth}|}
\hline
約書亞記 5:1-12 & \hyperref[sec:oVGpn0PwX_E]{2020年2月9日 崇拜直播｜約書亞記5\_1-12} & 2020-02-18 & \href{https://youtube.com/watch?v=oVGpn0PwX_E}{\texttt{oVGpn0PwX\_E}} \\
約書亞記 6:1-21 & \hyperref[sec:w8cFfcUefjQ]{2020年2月23日 崇拜直播｜約書亞記6\_1-21} & 2020-02-23 & \href{https://youtube.com/watch?v=w8cFfcUefjQ}{\texttt{w8cFfcUefjQ}} \\
約書亞記 9:1-21 & \hyperref[sec:ruBv0LtC8pg]{2020年3月15日 崇拜直播｜約書亞記9\_1-21} & 2020-03-15 & \href{https://youtube.com/watch?v=ruBv0LtC8pg}{\texttt{ruBv0LtC8pg}} \\
約書亞記 10:1-15 & \hyperref[sec:q99RT1cz_zM]{2020年3月22日 崇拜直播｜約書亞記10\_1-15} & 2020-03-23 & \href{https://youtube.com/watch?v=q99RT1cz_zM}{\texttt{q99RT1cz\_zM}} \\
約書亞記 11:1-15 & \hyperref[sec:pz6DLtcCEWM]{2020年3月29日 崇拜直播｜約書亞記11\_1-15} & 2020-03-30 & \href{https://youtube.com/watch?v=pz6DLtcCEWM}{\texttt{pz6DLtcCEWM}} \\
約書亞記 12:1-24 & \hyperref[sec:1PuKX5mdEMQ]{2020年4月19日 崇拜直播｜約書亞記12\_1-24} & 2020-04-21 & \href{https://youtube.com/watch?v=1PuKX5mdEMQ}{\texttt{1PuKX5mdEMQ}} \\
希伯來書 4:1-11 & \hyperref[sec:0sy2FpLq_lk]{2020年5月3日 崇拜直播｜希伯來書4\_1-11} & 2020-05-05 & \href{https://youtube.com/watch?v=0sy2FpLq-lk}{\texttt{0sy2FpLq-lk}} \\
希伯來書出埃及記 2:1-10 & \hyperref[sec:0WamuVI2aAk]{2020年5月10日 崇拜直播｜出埃及記2\_1-10;希伯來書11\_23} & 2020-05-11 & \href{https://youtube.com/watch?v=0WamuVI2aAk}{\texttt{0WamuVI2aAk}} \\
約書亞記 11:16-23 & \hyperref[sec:y8QuGSs8cQg]{2020年4月5日 崇拜直播｜約書亞記11\_16-23} & 2020-05-17 & \href{https://youtube.com/watch?v=y8QuGSs8cQg}{\texttt{y8QuGSs8cQg}} \\
馬太福音 20:1-16 & \hyperref[sec:XMPoIkBynFQ]{2020年5月31日 崇拜直播｜馬太福音20\_1-16} & 2020-06-01 & \href{https://youtube.com/watch?v=XMPoIkBynFQ}{\texttt{XMPoIkBynFQ}} \\
以弗所書 1:15-23 & \hyperref[sec:PBWvRjXaPlA]{2020年7月5日 崇拜直播(主餐)｜以弗所書1\_15-23} & 2020-07-06 & \href{https://youtube.com/watch?v=PBWvRjXaPlA}{\texttt{PBWvRjXaPlA}} \\
以弗所書 2:4-10 & \hyperref[sec:iBJ_D2WmvXo]{2020年7月12日 崇拜直播｜以弗所書2\_4-10} & 2020-07-13 & \href{https://youtube.com/watch?v=iBJ_D2WmvXo}{\texttt{iBJ\_D2WmvXo}} \\
以弗所書 4:1-16 & \hyperref[sec:hIi89f_Cp4Q]{2020年7月26日 崇拜直播｜以弗所書4\_1-16} & 2020-07-27 & \href{https://youtube.com/watch?v=hIi89f_Cp4Q}{\texttt{hIi89f\_Cp4Q}} \\
以弗所書 6:5-9 & \hyperref[sec:4TNIVRU8_Gc]{2020年8月16日 崇拜直播｜以弗所書6\_5-9} & 2020-08-18 & \href{https://youtube.com/watch?v=4TNIVRU8-Gc}{\texttt{4TNIVRU8-Gc}} \\
以弗所書 6:10-20 & \hyperref[sec:TUMQvpeyy3M]{2020年8月30日 崇拜直播｜以弗所書6\_10-20} & 2020-09-01 & \href{https://youtube.com/watch?v=TUMQvpeyy3M}{\texttt{TUMQvpeyy3M}} \\
但以理書 1:1-12 & \hyperref[sec:Ik8b7Ffs2qU]{2020年9月6日 崇拜直播｜但以理書1\_1-12} & 2020-09-11 & \href{https://youtube.com/watch?v=Ik8b7Ffs2qU}{\texttt{Ik8b7Ffs2qU}} \\
但以理書 2:1-30 & \hyperref[sec:oCGsFJy1z54]{2020年9月13日 崇拜直播｜但以理書2\_1-30} & 2020-09-17 & \href{https://youtube.com/watch?v=oCGsFJy1z54}{\texttt{oCGsFJy1z54}} \\
  & \hyperref[sec:PTdjTBD7nsw]{2020年9月18日「看不見的祝福」佈道會} & 2020-09-19 & \href{https://youtube.com/watch?v=PTdjTBD7nsw}{\texttt{PTdjTBD7nsw}} \\
但以理書 3:16-30 & \hyperref[sec:_Ymdcnu_z5c]{2020年9月20日 崇拜直播｜但以理書3\_16-30} & 2020-09-21 & \href{https://youtube.com/watch?v=_Ymdcnu-z5c}{\texttt{\_Ymdcnu-z5c}} \\
哥林多後書 5:1-10 & \hyperref[sec:jdoVqa_V4mw]{2020年9月27日 崇拜直播｜哥林多後書5\_1-10} & 2020-09-29 & \href{https://youtube.com/watch?v=jdoVqa_V4mw}{\texttt{jdoVqa\_V4mw}} \\
但以理書 4:4-18 & \hyperref[sec:Jfw_Vv2ojxk]{2020年10月11日 崇拜直播(主餐)｜但以理書4\_4-18} & 2020-10-12 & \href{https://youtube.com/watch?v=Jfw_Vv2ojxk}{\texttt{Jfw\_Vv2ojxk}} \\
但以理書 5:1-16 & \hyperref[sec:1ZS5UTvrjn8]{2020年10月18日 崇拜直播｜但以理書5\_1-16} & 2020-10-18 & \href{https://youtube.com/watch?v=1ZS5UTvrjn8}{\texttt{1ZS5UTvrjn8}} \\
但以理書 6:1-19 & \hyperref[sec:YHKUQm4ko_M]{2020年10月25日 崇拜直播(青少主日)｜但以理書6\_1-19} & 2020-10-25 & \href{https://youtube.com/watch?v=YHKUQm4ko_M}{\texttt{YHKUQm4ko\_M}} \\
馬太福音 12:15-21 & \hyperref[sec:MFeunbAAmR4]{2020年11月01日 崇拜直播｜馬太福音12\_15-21} & 2020-11-02 & \href{https://youtube.com/watch?v=MFeunbAAmR4}{\texttt{MFeunbAAmR4}} \\
約翰福音 8:1-11 & \hyperref[sec:PcbwkOrPy_0]{2020年11月15日 崇拜直播｜約翰福音8\_1-11} & 2020-11-16 & \href{https://youtube.com/watch?v=PcbwkOrPy_0}{\texttt{PcbwkOrPy\_0}} \\
詩篇利未記箴言 19:32 & \hyperref[sec:d4oKfBQrMaI]{2020年11月22日 崇拜直播(長者主日)｜利未記19\_32;詩篇71\_9,18;箴言20\_29} & 2020-11-23 & \href{https://youtube.com/watch?v=d4oKfBQrMaI}{\texttt{d4oKfBQrMaI}} \\
路加福音 2:22-38 & \hyperref[sec:CjzmC7ulsYw]{2020年12月6日 崇拜直播｜路加福音2\_22-38} & 2020-12-06 & \href{https://youtube.com/watch?v=CjzmC7ulsYw}{\texttt{CjzmC7ulsYw}} \\
  & \hyperref[sec:MIn0I7TUOfE]{2020年12月11日擁抱幸福分享晚會｜將幸福帶回家 ｜羅乃萱師母} & 2020-12-12 & \href{https://youtube.com/watch?v=MIn0I7TUOfE}{\texttt{MIn0I7TUOfE}} \\
哥林多後書 6:11-18 & \hyperref[sec:m_Hc7J2x7vY]{2020年12月13日 崇拜直播｜哥林多後書6\_11-18} & 2020-12-15 & \href{https://youtube.com/watch?v=m-Hc7J2x7vY}{\texttt{m-Hc7J2x7vY}} \\
路加福音 2:8-20 & \hyperref[sec:TPSrTFQecuk]{2020年12月20日 崇拜直播｜路加福音2\_8-20} & 2020-12-21 & \href{https://youtube.com/watch?v=TPSrTFQecuk}{\texttt{TPSrTFQecuk}} \\
  & \hyperref[sec:lmnMx2HA46Q]{2020年12月24日｜擁抱幸福分享晚會:平安夜晚會 ｜旺角浸信會聖樂團隊} & 2020-12-25 & \href{https://youtube.com/watch?v=lmnMx2HA46Q}{\texttt{lmnMx2HA46Q}} \\
馬太福音哥林多前書 11:28-30 & \hyperref[sec:DokqOqhZ3o0]{2020年12月27日崇拜直播｜羅志強執事｜福音與愛的實踐｜馬太福音11\_28-30;哥林多前書13\_4-5} & 2021-01-03 & \href{https://youtube.com/watch?v=DokqOqhZ3o0}{\texttt{DokqOqhZ3o0}} \\
馬太福音 5:1-10 & \hyperref[sec:bLhELYX_Uw0]{2021年1月3日崇拜直播｜陳明泉牧師｜幸福由虛心開始｜馬太福音5\_1-10} & 2021-01-03 & \href{https://youtube.com/watch?v=bLhELYX_Uw0}{\texttt{bLhELYX\_Uw0}} \\
馬太福音 5:4 & \hyperref[sec:POKxDUqFlWw]{2021年1月10日崇拜直播｜陳冠成牧師｜幸福由哀慟開始｜馬太福音5\_4} & 2021-01-12 & \href{https://youtube.com/watch?v=POKxDUqFlWw}{\texttt{POKxDUqFlWw}} \\
馬太福音 5:5 & \hyperref[sec:YMwn1_GRZYg]{2021年1月24日崇拜直播｜吳真真牧師｜幸福由溫柔開始｜馬太福音5\_5} & 2021-01-24 & \href{https://youtube.com/watch?v=YMwn1_GRZYg}{\texttt{YMwn1\_GRZYg}} \\
馬太福音 5:6 & \hyperref[sec:s_Njf9cQc18]{2021年1月31日崇拜直播｜陳啟興牧師｜幸福由飢渴慕義開始｜馬太福音5\_6} & 2021-02-03 & \href{https://youtube.com/watch?v=s_Njf9cQc18}{\texttt{s\_Njf9cQc18}} \\
哥林多後書 8:1-9 & \hyperref[sec:fM11JJVt3SY]{2021年2月28日崇拜直播｜梁家麟牧師｜慈惠的精神｜哥林多後書8\_1-9} & 2021-02-28 & \href{https://youtube.com/watch?v=fM11JJVt3SY}{\texttt{fM11JJVt3SY}} \\
馬太福音 5:10 & \hyperref[sec:rPutWeJ6_Xg]{2021年3月7日崇拜直播｜陳明泉牧師｜幸福由為義受逼迫開始｜馬太福音5\_10} & 2021-03-07 & \href{https://youtube.com/watch?v=rPutWeJ6_Xg}{\texttt{rPutWeJ6\_Xg}} \\
哥林多後書 8:9-24 & \hyperref[sec:zila_5Bj4Bo]{2021年3月14日崇拜直播｜梁家麟牧師｜均平的原則｜哥林多後書8\_9-24} & 2021-03-14 & \href{https://youtube.com/watch?v=zila_5Bj4Bo}{\texttt{zila\_5Bj4Bo}} \\
馬可福音 10:17-22 & \hyperref[sec:qJa4QRG5SV0]{2021年3月21日崇拜直播｜陳冠成牧師｜感動還需敢動｜馬可福音10\_17-22} & 2021-03-22 & \href{https://youtube.com/watch?v=qJa4QRG5SV0}{\texttt{qJa4QRG5SV0}} \\
馬可福音 12:18-27 & \hyperref[sec:pK4j32Egyx0]{2021年3月28日崇拜直播｜陳啟興牧師｜復活的大能｜馬可福音12\_18-27} & 2021-04-01 & \href{https://youtube.com/watch?v=pK4j32Egyx0}{\texttt{pK4j32Egyx0}} \\
馬可福音 15:21-32 & \hyperref[sec:NcnNHWwVDiw]{2021年4月2日受苦節崇拜直播｜吳真真牧師｜十字架的愛｜馬可福音15\_21-32} & 2021-04-02 & \href{https://youtube.com/watch?v=NcnNHWwVDiw}{\texttt{NcnNHWwVDiw}} \\
約翰福音 20:24-29 & \hyperref[sec:R7EPR0TUgHE]{2021年4月4日復活節主餐崇拜直播｜陳明泉牧師｜疑與信之間｜約翰福音20\_24-29} & 2021-04-05 & \href{https://youtube.com/watch?v=R7EPR0TUgHE}{\texttt{R7EPR0TUgHE}} \\
馬可福音 8:27-9:37 & \hyperref[sec:oKGXX4r5yvQ]{2021年4月18日崇拜直播｜陳妙濃傳道｜就當捨己｜馬可福音8\_27-9\_37} & 2021-04-18 & \href{https://youtube.com/watch?v=oKGXX4r5yvQ}{\texttt{oKGXX4r5yvQ}} \\
哥林多後書 9:1-15 & \hyperref[sec:opD_mMvcV8I]{2021年5月2日主餐崇拜直播｜梁家麟牧師｜甘心作美事｜哥林多後書9\_1-15} & 2021-05-04 & \href{https://youtube.com/watch?v=opD_mMvcV8I}{\texttt{opD\_mMvcV8I}} \\
哈巴谷書 3:16-19 & \hyperref[sec:d1qGfZ_AH3c]{2021年5月16日崇拜直播｜陳冠成牧師｜從聖經看喜怒哀愁~因救我的神喜樂｜哈巴谷書3\_16-19} & 2021-05-16 & \href{https://youtube.com/watch?v=d1qGfZ_AH3c}{\texttt{d1qGfZ\_AH3c}} \\
哥林多後書 10:1-18 & \hyperref[sec:RH7uL12ghP4]{2021年5月23日崇拜直播｜梁家麟牧師｜保羅的勇敢和怯懦｜哥林多後書10\_1-18} & 2021-05-25 & \href{https://youtube.com/watch?v=RH7uL12ghP4}{\texttt{RH7uL12ghP4}} \\
以弗所書 4:26-32 & \hyperref[sec:Czt7dmGniVs]{2021年5月30日崇拜直播｜吳真真牧師｜從聖經看喜怒哀愁 - 怒亦有道｜以弗所書4\_26-32} & 2021-06-05 & \href{https://youtube.com/watch?v=Czt7dmGniVs}{\texttt{Czt7dmGniVs}} \\
詩篇 13:1-6 & \hyperref[sec:9JRB4RUnJVQ]{2021年6月6日主餐崇拜直播｜陳明泉牧師｜從聖經看喜怒哀愁 - 哀嘆中仰望｜詩篇13\_1-6} & 2021-06-19 & \href{https://youtube.com/watch?v=9JRB4RUnJVQ}{\texttt{9JRB4RUnJVQ}} \\
哥林多後書 10:12-18 & \hyperref[sec:bhMQ0ofdb_Q]{2021年7月4日崇拜直播｜梁家麟牧師｜越界問題｜哥林多後書10\_12-18} & 2021-07-10 & \href{https://youtube.com/watch?v=bhMQ0ofdb-Q}{\texttt{bhMQ0ofdb-Q}} \\
雅各書 1:19-27 & \hyperref[sec:VDCXuQp5CNk]{2021年7月11日崇拜直播｜陳冠成牧師｜生命見真章系列 - 聽道為了行道｜雅各書1\_19-27} & 2021-07-19 & \href{https://youtube.com/watch?v=VDCXuQp5CNk}{\texttt{VDCXuQp5CNk}} \\
哥林多後書 11:1-15 & \hyperref[sec:rb0qE8Nt22M]{2021年7月18日崇拜直播｜梁家麟牧師｜保羅為個人身份的自辯｜哥林多後書11\_1-15} & 2021-07-19 & \href{https://youtube.com/watch?v=rb0qE8Nt22M}{\texttt{rb0qE8Nt22M}} \\
雅各書 2:1-7 & \hyperref[sec:n5CNurE7uu0]{2021年7月25日崇拜直播｜陳啟興牧師｜生命見真章系列 - 不以貌取人｜雅各書2\_1-7} & 2021-08-01 & \href{https://youtube.com/watch?v=n5CNurE7uu0}{\texttt{n5CNurE7uu0}} \\
雅各書 2:14-26 & \hyperref[sec:a56U_20AydY]{2021年8月1日主餐崇拜直播｜陳冠成牧師｜生命見真章系列 - 顯明你信什麼｜雅各書2\_14-26} & 2021-08-01 & \href{https://youtube.com/watch?v=a56U-20AydY}{\texttt{a56U-20AydY}} \\
雅各書 3:1-18 & \hyperref[sec:z9XCLH0XmKE]{2021年8月8日崇拜直播｜陳明泉牧師｜生命見真章系列 - 鑄人以言｜雅各書3\_1-18} & 2021-08-08 & \href{https://youtube.com/watch?v=z9XCLH0XmKE}{\texttt{z9XCLH0XmKE}} \\
哥林多後書 11:16-12:10 & \hyperref[sec:OhUT6e4ffZM]{2021年8月15日崇拜直播｜梁家麟牧師｜愚拙的保羅｜哥林多後書11\_16-12\_10} & 2021-08-16 & \href{https://youtube.com/watch?v=OhUT6e4ffZM}{\texttt{OhUT6e4ffZM}} \\
雅各書 4:1-10 & \hyperref[sec:0pujbt8sHO8]{2021年8月22日崇拜直播｜吳真真牧師｜生命見真章系列 - 與神為友,與世俗為敵｜雅各書4\_1-10} & 2021-08-22 & \href{https://youtube.com/watch?v=0pujbt8sHO8}{\texttt{0pujbt8sHO8}} \\
哥林多前書 13:1-13 & \hyperref[sec:bJlzZqbRugo]{2021年8月29日崇拜直播｜何羅乃萱師母｜疫中有愛｜哥林多前書13\_1-13} & 2021-08-29 & \href{https://youtube.com/watch?v=bJlzZqbRugo}{\texttt{bJlzZqbRugo}} \\
哥林多後書 12:1-10 & \hyperref[sec:Xf1UScB62N8]{2021年9月19日崇拜直播｜梁家麟牧師｜保羅不刷存在感｜哥林多後書12\_1-10} & 2021-10-15 & \href{https://youtube.com/watch?v=Xf1UScB62N8}{\texttt{Xf1UScB62N8}} \\
約翰壹書 3:21-24 & \hyperref[sec:aSTV2iY5HmY]{2021年10月10日崇拜直播｜羅志雄牧師｜坦然無懼｜約翰壹書3\_21-24} & 2021-10-17 & \href{https://youtube.com/watch?v=aSTV2iY5HmY}{\texttt{aSTV2iY5HmY}} \\
雅各書 5:13-20 & \hyperref[sec:rfN73DzwrE4]{2021年10月3日主餐崇拜直播｜陳明泉牧師｜生命見真章系列 - 與生活同步的禱告｜雅各書5\_13-20} & 2021-10-19 & \href{https://youtube.com/watch?v=rfN73DzwrE4}{\texttt{rfN73DzwrE4}} \\
馬太福音 14:22-23 & \hyperref[sec:irpyKMaPiJI]{2021年10月17日崇拜直播｜陳明泉牧師｜放心!不要怕｜馬太福音14\_22-23} & 2021-10-21 & \href{https://youtube.com/watch?v=irpyKMaPiJI}{\texttt{irpyKMaPiJI}} \\
雅各書 5:1-6 & \hyperref[sec:LQ_YiQLwO5g]{2021年9月26日崇拜直播｜陳冠成牧師｜生命見真章系列 - 追求真正富足｜雅各書5\_1-6} & 2021-10-21 & \href{https://youtube.com/watch?v=LQ_YiQLwO5g}{\texttt{LQ\_YiQLwO5g}} \\
列王記下 5:1-14 & \hyperref[sec:SlF159ccdOo]{2021年10月31日崇拜直播｜陳冠成牧師｜這樣就能得醫治｜列王記下5\_1-14} & 2021-11-09 & \href{https://youtube.com/watch?v=SlF159ccdOo}{\texttt{SlF159ccdOo}} \\
創世記 39:1-23 & \hyperref[sec:wlHRHi438v8]{2021年10月24日青少主日崇拜直播｜梁煒傳道｜當不幸敲門時｜創世記39\_1-23} & 2021-11-21 & \href{https://youtube.com/watch?v=wlHRHi438v8}{\texttt{wlHRHi438v8}} \\
約書亞記 14:6-15 & \hyperref[sec:ec0Lp80ZOxQ]{2021年11月14日長者主日崇拜直播｜葉筱春牧師｜我還是強壯｜約書亞記14\_6-15} & 2021-11-21 & \href{https://youtube.com/watch?v=ec0Lp80ZOxQ}{\texttt{ec0Lp80ZOxQ}} \\
約翰福音 2:13-22 & \hyperref[sec:N45H8jB07Gc]{2021年11月28日崇拜直播｜陳明泉牧師｜基督帶來的革新:新的聖殿｜約翰福音2\_13-22} & 2021-11-29 & \href{https://youtube.com/watch?v=N45H8jB07Gc}{\texttt{N45H8jB07Gc}} \\
約翰福音 1:14-18 & \hyperref[sec:Msw_I2sXTQc]{2021年12月5日崇拜直播｜陳明泉牧師｜基督帶來的革新:新的恩典｜約翰福音1\_14-18} & 2021-12-08 & \href{https://youtube.com/watch?v=Msw-I2sXTQc}{\texttt{Msw-I2sXTQc}} \\
哥林多後書 12:19-13:4 & \hyperref[sec:Sd44cjxmurg]{2021年12月26日崇拜直播｜梁家麟牧師｜最後的提醒｜哥林多後書12\_19-13\_4} & 2021-12-29 & \href{https://youtube.com/watch?v=Sd44cjxmurg}{\texttt{Sd44cjxmurg}} \\
馬太福音 28:8-20 & \hyperref[sec:oDel5UYyQt0]{2022年1月2日崇拜直播｜陳明泉牧師｜我們的6G - 門訓不斷｜馬太福音28\_8-20} & 2022-01-03 & \href{https://youtube.com/watch?v=oDel5UYyQt0}{\texttt{oDel5UYyQt0}} \\
申命記 6:1-9 & \hyperref[sec:Qke9jW9Cx6E]{2022年1月16日早堂崇拜直播｜陳冠成牧師｜我們的6G - 聖經引導｜申命記6\_1-9} & 2022-01-17 & \href{https://youtube.com/watch?v=Qke9jW9Cx6E}{\texttt{Qke9jW9Cx6E}} \\
哥林多後書 13:1-14 & \hyperref[sec:9NLZUx6GLHU]{2022年1月23日崇拜直播｜梁家麟牧師｜成全聖徒的召命｜哥林多後書13\_1-14} & 2022-01-24 & \href{https://youtube.com/watch?v=9NLZUx6GLHU}{\texttt{9NLZUx6GLHU}} \\
馬太福音 9:35-38 & \hyperref[sec:tgXI8MUm32I]{2022年1月30日崇拜直播｜梁煒傳道｜我們的6G - 胸懷普世｜馬太福音9\_35-38} & 2022-01-30 & \href{https://youtube.com/watch?v=tgXI8MUm32I}{\texttt{tgXI8MUm32I}} \\
腓立比書 1:1-11 & \hyperref[sec:XQBczBTx6j0]{2022年2月6日崇拜直播｜陳明泉牧師｜我們的6G - 凝聚動力 GET CLOSER｜腓立比書1\_1-11} & 2022-02-07 & \href{https://youtube.com/watch?v=XQBczBTx6j0}{\texttt{XQBczBTx6j0}} \\
箴言 27:5-17 & \hyperref[sec:Q_GbM3l6t1s]{2022年2月13日崇拜直播｜陳冠成牧師｜我們的6G - 真誠砥礪 Genuine Relationship｜箴言27\_5-17} & 2022-02-14 & \href{https://youtube.com/watch?v=Q_GbM3l6t1s}{\texttt{Q\_GbM3l6t1s}} \\
約翰福音 1:1-18 & \hyperref[sec:yigtcB2OiSE]{2022年2月20日崇拜直播｜梁家麟牧師｜太初有道｜約翰福音1\_1-18} & 2022-02-21 & \href{https://youtube.com/watch?v=yigtcB2OiSE}{\texttt{yigtcB2OiSE}} \\
希伯來書 1:1-2:4 & \hyperref[sec:cYMHT7F9Rko]{2022年3月6日崇拜直播｜陳明泉牧師｜超越的基督:神親自曉喻｜希伯來書1\_1-2\_4} & 2022-03-07 & \href{https://youtube.com/watch?v=cYMHT7F9Rko}{\texttt{cYMHT7F9Rko}} \\
希伯來書 4:14-5:10 & \hyperref[sec:bMqVNP_dN2s]{2022年3月20日崇拜直播｜陳冠成牧師｜超越的基督:完全的大祭司｜希伯來書4\_14-5\_10} & 2022-03-20 & \href{https://youtube.com/watch?v=bMqVNP-dN2s}{\texttt{bMqVNP-dN2s}} \\
約翰福音 1:1-18 & \hyperref[sec:XlsH8Zitzyg]{2022年3月27日崇拜直播｜梁家麟牧師｜施浸約翰的見證(一)｜約翰福音1\_1-18} & 2022-03-28 & \href{https://youtube.com/watch?v=XlsH8Zitzyg}{\texttt{XlsH8Zitzyg}} \\
希伯來書 7:1-28 & \hyperref[sec:Or14nIn97R8]{2022年4月3日崇拜直播｜陳明泉牧師｜超越的基督:按麥基洗德等次為祭司｜希伯來書7\_1-28} & 2022-04-04 & \href{https://youtube.com/watch?v=Or14nIn97R8}{\texttt{Or14nIn97R8}} \\
約翰福音 1:19-38 & \hyperref[sec:V98EonzJR3Q]{2022年4月10日崇拜直播｜梁家麟牧師｜施洗約翰的見證(二)｜約翰福音1\_19-38} & 2022-04-11 & \href{https://youtube.com/watch?v=V98EonzJR3Q}{\texttt{V98EonzJR3Q}} \\
  & \hyperref[sec:knwfu3yRKQ8]{2022年4月15日受苦節晚會直播｜十架上的你與他/她} & 2022-04-16 & \href{https://youtube.com/watch?v=knwfu3yRKQ8}{\texttt{knwfu3yRKQ8}} \\
希伯來書 8:1-13 & \hyperref[sec:U9MWFdIg8Fc]{2022年4月17日復活崇拜直播｜陳明泉牧師｜超越的基督:新約的中保｜希伯來書8\_1-13} & 2022-04-17 & \href{https://youtube.com/watch?v=U9MWFdIg8Fc}{\texttt{U9MWFdIg8Fc}} \\
希伯來書 9:1-14 & \hyperref[sec:rKFXkQ539_w]{2022年4月24日崇拜直播｜陳冠成牧師｜超越的基督:神的帳幕在人間｜希伯來書9\_1-14} & 2022-04-25 & \href{https://youtube.com/watch?v=rKFXkQ539_w}{\texttt{rKFXkQ539\_w}} \\
希伯來書 11:11-16 & \hyperref[sec:vxnWocuBbhg]{2022年5月8日母親節崇拜直播｜吳真真牧師｜超越的基督:憑著信作母親｜希伯來書11\_11-16} & 2022-05-09 & \href{https://youtube.com/watch?v=vxnWocuBbhg}{\texttt{vxnWocuBbhg}} \\
希伯來書 12:1-13 & \hyperref[sec:qTE37kvGc24]{2022年5月15日崇拜直播｜陳啟興牧師｜超越的基督:人生馬拉松｜希伯來書12\_1-13} & 2022-05-16 & \href{https://youtube.com/watch?v=qTE37kvGc24}{\texttt{qTE37kvGc24}} \\
箴言 31:1-31 & \hyperref[sec:uzCPC0kQbAA]{2022年5月29日崇拜直播｜何羅乃萱師母｜疫情下才德婦人的啟示｜箴言31\_1-31} & 2022-05-30 & \href{https://youtube.com/watch?v=uzCPC0kQbAA}{\texttt{uzCPC0kQbAA}} \\
創世記 39:1-23 & \hyperref[sec:e5I3VklFIWg]{2022年6月5日主餐崇拜直播｜陳明泉牧師｜並肩勇闖系列:耶和華的同在｜創世記39\_1-23} & 2022-06-06 & \href{https://youtube.com/watch?v=e5I3VklFIWg}{\texttt{e5I3VklFIWg}} \\
列王記下 4:1-7 & \hyperref[sec:_Anjyu4UPkQ]{2022年6月12日崇拜直播｜陳冠成牧師｜並肩勇闖系列:神不撇下我們｜列王記下4\_1-7} & 2022-06-13 & \href{https://youtube.com/watch?v=-Anjyu4UPkQ}{\texttt{-Anjyu4UPkQ}} \\
路得記 1:6-18 & \hyperref[sec:8sXdI3Lte_M]{2022年6月26日崇拜直播｜陳妙濃傳道｜並肩勇闖系列:以恩慈相待｜路得記1\_6-18} & 2022-06-27 & \href{https://youtube.com/watch?v=8sXdI3Lte_M}{\texttt{8sXdI3Lte\_M}} \\
列王記下 6:1-7 & \hyperref[sec:c6nbM3DecyY]{2022年7月3日崇拜直播｜陳冠成牧師｜並肩勇闖系列:邀請上帝同行｜列王記下6\_1-7} & 2022-07-04 & \href{https://youtube.com/watch?v=c6nbM3DecyY}{\texttt{c6nbM3DecyY}} \\
路加福音 24:13-35 & \hyperref[sec:sq3x1g07xfo]{2022年7月10日崇拜直播｜羅潔盈博士｜我們的心豈不火熱?｜路加福音24\_13-35} & 2022-07-11 & \href{https://youtube.com/watch?v=sq3x1g07xfo}{\texttt{sq3x1g07xfo}} \\
尼希米記 6:1-16 & \hyperref[sec:SQA6e1_jseU]{2022年8月7日崇拜直播｜吳真真牧師｜並肩勇闖系列:戰勝懼怕｜尼希米記6\_1-16} & 2022-08-07 & \href{https://youtube.com/watch?v=SQA6e1_jseU}{\texttt{SQA6e1\_jseU}} \\
約翰福音 1:29-37 & \hyperref[sec:IiBVTMF9Vak]{2022年8月14日崇拜直播｜梁家麟牧師｜施洗約翰的見證(三)｜約翰福音1\_29-37} & 2022-08-16 & \href{https://youtube.com/watch?v=IiBVTMF9Vak}{\texttt{IiBVTMF9Vak}} \\
以斯拉記 1:1-11 & \hyperref[sec:ZKaifzpEjak]{2022年8月21日崇拜直播｜陳明泉牧師｜並肩勇闖系列:願神與這人同在｜以斯拉記1\_1-11} & 2022-08-22 & \href{https://youtube.com/watch?v=ZKaifzpEjak}{\texttt{ZKaifzpEjak}} \\
羅馬書 15:13 & \hyperref[sec:hASHSP9n8fQ]{2022年8月28日崇拜直播｜區祥江博士｜大有盼望的秘訣｜羅馬書15\_13} & 2022-08-30 & \href{https://youtube.com/watch?v=hASHSP9n8fQ}{\texttt{hASHSP9n8fQ}} \\
民數記 20:2-13 & \hyperref[sec:0H976XcFnC4]{2022年9月11日崇拜直播｜陳冠成牧師｜以石為記:擊打出水的磐石｜民數記20\_2-13} & 2022-09-11 & \href{https://youtube.com/watch?v=0H976XcFnC4}{\texttt{0H976XcFnC4}} \\
約翰福音 1:38-43 & \hyperref[sec:uD1kVchCRPU]{2022年9月18日崇拜直播｜梁家麟牧師｜第一批追隨者｜約翰福音1\_38-43} & 2022-09-19 & \href{https://youtube.com/watch?v=uD1kVchCRPU}{\texttt{uD1kVchCRPU}} \\
撒母耳記上 17:26-40 & \hyperref[sec:PSHqqITZkIo]{2022年10月2日崇拜直播｜陳明泉牧師｜以石為記:擊打巨人的五塊小石｜撒母耳記上17\_26-40} & 2022-10-02 & \href{https://youtube.com/watch?v=PSHqqITZkIo}{\texttt{PSHqqITZkIo}} \\
以賽亞書 28:14-22 & \hyperref[sec:x9IT18Te0VA]{2022年10月9日崇拜直播｜陳冠成牧師｜以石為記:試驗過的房角石｜以賽亞書28\_14-22} & 2022-10-10 & \href{https://youtube.com/watch?v=x9IT18Te0VA}{\texttt{x9IT18Te0VA}} \\
馬太福音 4:1-11 & \hyperref[sec:862eUMKnVnc]{2022年10月16日崇拜直播｜梁煒傳道｜以石為記:試探之石｜馬太福音4\_1-11} & 2022-10-18 & \href{https://youtube.com/watch?v=862eUMKnVnc}{\texttt{862eUMKnVnc}} \\
約翰福音 8:1-11 & \hyperref[sec:YIG0i1d0dV8]{2022年10月30日崇拜直播｜吳真真牧師｜以石為記:誰敢拿起審判之石｜約翰福音8\_1-11} & 2022-10-31 & \href{https://youtube.com/watch?v=YIG0i1d0dV8}{\texttt{YIG0i1d0dV8}} \\
羅馬書 9:30-33 & \hyperref[sec:ucA_hp7yY_0]{2022年11月6日崇拜直播｜陳明泉牧師｜以石為記:跌人的磐石｜羅馬書9\_30-33} & 2022-11-08 & \href{https://youtube.com/watch?v=ucA_hp7yY_0}{\texttt{ucA\_hp7yY\_0}} \\
馬太福音 16:13-19 & \hyperref[sec:y7hPaMojRiM]{2022年11月13日崇拜直播｜陳冠成牧師｜以石為記:建造教會的磐石｜馬太福音16\_13-19} & 2022-11-15 & \href{https://youtube.com/watch?v=y7hPaMojRiM}{\texttt{y7hPaMojRiM}} \\
詩篇 71:5-18 & \hyperref[sec:esPM0k_dZBY]{2022年11月20日長者主日崇拜直播｜羅志雄牧師｜恩典的年歲｜詩篇71\_5-18} & 2022-11-22 & \href{https://youtube.com/watch?v=esPM0k-dZBY}{\texttt{esPM0k-dZBY}} \\
提摩太後書 1:1-7 & \hyperref[sec:4hDfY5nlUxk]{2022年12月4日主餐暨堂慶主日崇拜直播｜梁家麟牧師｜遺產.使命.恩典｜提摩太後書1\_1-7} & 2022-12-08 & \href{https://youtube.com/watch?v=4hDfY5nlUxk}{\texttt{4hDfY5nlUxk}} \\
馬太福音 2:1-12 & \hyperref[sec:6BkHUdPWvFs]{2022年12月11日崇拜直播｜陳明泉牧師｜看見星星的你｜馬太福音2\_1-12} & 2022-12-13 & \href{https://youtube.com/watch?v=6BkHUdPWvFs}{\texttt{6BkHUdPWvFs}} \\
路加福音 2:1-7 & \hyperref[sec:1SOWC0lwIU0]{2022年12月25日崇拜直播｜陳冠成牧師｜展現脆弱的一面｜路加福音2\_1-7} & 2022-12-29 & \href{https://youtube.com/watch?v=1SOWC0lwIU0}{\texttt{1SOWC0lwIU0}} \\
約翰福音 1:43-51 & \hyperref[sec:HUos2vHhz9s]{2023年1月1日崇拜直播｜梁家麟牧師｜掌握召命｜約翰福音1\_43-51} & 2023-01-05 & \href{https://youtube.com/watch?v=HUos2vHhz9s}{\texttt{HUos2vHhz9s}} \\
馬太福音 4:18-25 & \hyperref[sec:WHVDHKq51Y4]{2023年1月8日崇拜直播｜陳明泉牧師｜給當代門徒的15個忠告:真正的追隨｜馬太福音4\_18-25} & 2023-01-09 & \href{https://youtube.com/watch?v=WHVDHKq51Y4}{\texttt{WHVDHKq51Y4}} \\
馬太福音 5:13-16 & \hyperref[sec:Z3AMKkvW_mc]{2023年1月15日崇拜直播｜陳冠成牧師｜給當代門徒的15個忠告:成為世上的鹽和光｜馬太福音5\_13-16} & 2023-01-17 & \href{https://youtube.com/watch?v=Z3AMKkvW_mc}{\texttt{Z3AMKkvW\_mc}} \\
馬太福音 5:17-20 & \hyperref[sec:AMjA5Kj3ujM]{2023年1月29日崇拜直播｜吳真真牧師｜給當代門徒的15個忠告:遵行耶穌所成全的律法｜馬太福音5\_17-20} & 2023-02-04 & \href{https://youtube.com/watch?v=AMjA5Kj3ujM}{\texttt{AMjA5Kj3ujM}} \\
馬太福音 5:21-26 & \hyperref[sec:iHCyyrWPaj4]{2023年2月5日崇拜直播｜陳啟興牧師｜給當代門徒的15個忠告:處理怒氣｜馬太福音5\_21-26} & 2023-02-07 & \href{https://youtube.com/watch?v=iHCyyrWPaj4}{\texttt{iHCyyrWPaj4}} \\
馬太福音 5:27-32 & \hyperref[sec:N8TX9PAocpk]{2023年2月12日崇拜直播｜陳冠成牧師｜給當代門徒的15個忠告:持守聖潔｜馬太福音5\_27-32} & 2023-02-13 & \href{https://youtube.com/watch?v=N8TX9PAocpk}{\texttt{N8TX9PAocpk}} \\
路加福音 10:25-37 & \hyperref[sec:VsBsdArDKAI]{2023年2月26日崇拜直播｜李盛林牧師｜誰是你的同行者｜路加福音10\_25-37} & 2023-03-02 & \href{https://youtube.com/watch?v=VsBsdArDKAI}{\texttt{VsBsdArDKAI}} \\
馬太福音 5:33-37 & \hyperref[sec:O_RnYw4oPiY]{2023年3月5日崇拜直播｜陳明泉牧師｜給當代門徒的15個忠告:成為誠信可靠的人｜馬太福音5\_33-37} & 2023-03-12 & \href{https://youtube.com/watch?v=O-RnYw4oPiY}{\texttt{O-RnYw4oPiY}} \\
馬太福音 5:38-48 & \hyperref[sec:DeJAKAGjSug]{2023年3月19日崇拜直播｜陳冠成牧師｜給當代門徒的15個忠告:像天父完全一樣｜馬太福音5\_38-48} & 2023-03-19 & \href{https://youtube.com/watch?v=DeJAKAGjSug}{\texttt{DeJAKAGjSug}} \\
馬太福音 6:1-4 & \hyperref[sec:XpqKURe4134]{2023年3月26日崇拜直播｜余嘉敏傳道｜給當代門徒的15個忠告:擁有一顆單純服侍主的心｜馬太福音6\_1-4} & 2023-03-27 & \href{https://youtube.com/watch?v=XpqKURe4134}{\texttt{XpqKURe4134}} \\
馬太福音 6:5-15 & \hyperref[sec:F3PyAVrDKHQ]{2023年4月2日崇拜直播｜陳明泉牧師｜給當代門徒的15個忠告:禱告無添加｜馬太福音6\_5-15} & 2023-04-02 & \href{https://youtube.com/watch?v=F3PyAVrDKHQ}{\texttt{F3PyAVrDKHQ}} \\
約翰福音 21:1-17 & \hyperref[sec:HNqZXsZMfMU]{2023年4月9日復活節崇拜直播｜陳冠成牧師｜復活主的顯現｜約翰福音21\_1-17} & 2023-04-09 & \href{https://youtube.com/watch?v=HNqZXsZMfMU}{\texttt{HNqZXsZMfMU}} \\
馬太福音 6:16-23 & \hyperref[sec:Dcz9krAW4uI]{2023年4月16日崇拜直播｜羅志雄牧師｜給當代門徒的15個忠告:禁食為何｜馬太福音6\_16-23} & 2023-04-17 & \href{https://youtube.com/watch?v=Dcz9krAW4uI}{\texttt{Dcz9krAW4uI}} \\
馬太福音 7:1-6 & \hyperref[sec:zKeckdMvwPY]{2023年4月23日崇拜直播｜陳明泉牧師｜給當代門徒的15個忠告:別錯放辨別力｜馬太福音7\_1-6} & 2023-04-24 & \href{https://youtube.com/watch?v=zKeckdMvwPY}{\texttt{zKeckdMvwPY}} \\
詩篇 73:1-28 & \hyperref[sec:PtF_7l7B_qs]{2023年4月30日崇拜直播｜區祥江博士｜如何面對嫉妒的困擾｜詩篇73\_1-28} & 2023-05-01 & \href{https://youtube.com/watch?v=PtF_7l7B_qs}{\texttt{PtF\_7l7B\_qs}} \\
馬太福音 6:25-34 & \hyperref[sec:38XteBMXhDk]{2023年5月7日青少主日崇拜直播｜梁煒傳道｜給當代門徒的15個忠告:不要憂慮｜馬太福音6\_25-34} & 2023-05-08 & \href{https://youtube.com/watch?v=38XteBMXhDk}{\texttt{38XteBMXhDk}} \\
撒母耳記上 1:4-28 & \hyperref[sec:iCj_rgLUlC4]{2023年5月14日母親節主日崇拜直播｜余嘉敏傳道｜苦媽的最佳選擇｜撒母耳記上1\_4-28} & 2023-05-15 & \href{https://youtube.com/watch?v=iCj-rgLUlC4}{\texttt{iCj-rgLUlC4}} \\
馬太福音 13:1-9 & \hyperref[sec:qRN_c71zY24]{2023年5月28日崇拜直播｜張培良傳道｜成為不合理的撒種人｜馬太福音13\_1-9} & 2023-05-29 & \href{https://youtube.com/watch?v=qRN-c71zY24}{\texttt{qRN-c71zY24}} \\
馬太福音 7:13-23 & \hyperref[sec:bwknUBKFdWs]{2023年6月4日主餐崇拜直播｜吳真真牧師｜給當代門徒的15個忠告:走窄路｜馬太福音7\_13-23} & 2023-06-05 & \href{https://youtube.com/watch?v=bwknUBKFdWs}{\texttt{bwknUBKFdWs}} \\
哥林多前書彼得前書 11:3-12 & \hyperref[sec:UnHN8_G394I]{2023年6月18日父親節崇拜直播｜梁家麟牧師｜作頭的男人｜哥林多前書11\_3-12;彼得前書3\_7} & 2023-06-20 & \href{https://youtube.com/watch?v=UnHN8_G394I}{\texttt{UnHN8\_G394I}} \\
約翰福音 2:13-25 & \hyperref[sec:rlyrCSM_7W8]{2023年7月9日崇拜直播｜梁家麟牧師｜潔淨聖殿｜約翰福音2\_13-25} & 2023-07-12 & \href{https://youtube.com/watch?v=rlyrCSM-7W8}{\texttt{rlyrCSM-7W8}} \\
士師記 2:6-3:6 & \hyperref[sec:CAkR8_FGcmQ]{2023年7月16日崇拜直播｜吳真真牧師｜「家醜說出口」系列:信仰傳承失效｜士師記2\_6-3\_6} & 2023-07-17 & \href{https://youtube.com/watch?v=CAkR8-FGcmQ}{\texttt{CAkR8-FGcmQ}} \\
士師記 8:22-35 & \hyperref[sec:iklPHMTALSA]{2023年8月6日崇拜直播｜陳明泉牧師｜「家醜說出口」系列:竟為下代設下網羅｜士師記8\_22-35} & 2023-08-07 & \href{https://youtube.com/watch?v=iklPHMTALSA}{\texttt{iklPHMTALSA}} \\
士師記 11:1-40 & \hyperref[sec:Fr77s7J2G5I]{2023年8月13日崇拜直播｜陳冠成牧師｜「家醜說出口」系列:傷痛不代傳｜士師記11\_1-40} & 2023-08-14 & \href{https://youtube.com/watch?v=Fr77s7J2G5I}{\texttt{Fr77s7J2G5I}} \\
士師記 9:1-6 & \hyperref[sec:bCwRlVK9wDE]{2023年8月20日崇拜直播｜陳明泉牧師｜「家醜說出口」系列:寵出個皇兒｜士師記9\_1-6} & 2023-08-22 & \href{https://youtube.com/watch?v=bCwRlVK9wDE}{\texttt{bCwRlVK9wDE}} \\
士師記 13:8-23 & \hyperref[sec:sDc2iE_sQ_g]{2023年9月3日崇拜直播｜陳明泉牧師｜「家醜說出口」系列:錯焦的敬虔｜士師記13\_8-23} & 2023-09-04 & \href{https://youtube.com/watch?v=sDc2iE_sQ-g}{\texttt{sDc2iE\_sQ-g}} \\
士師記 14:1-9 & \hyperref[sec:Up7U5HlOZfo]{2023年9月10日崇拜直播｜陳冠成牧師｜「家醜說出口」系列:欲望阻不了｜士師記14\_1-9} & 2023-09-10 & \href{https://youtube.com/watch?v=Up7U5HlOZfo}{\texttt{Up7U5HlOZfo}} \\
士師記 16:1-31 & \hyperref[sec:rllwPQsYaK0]{2023年9月17日崇拜直播｜余嘉敏姑娘｜「家醜說出口」系列:過度自信者的求告｜士師記16\_1-31} & 2023-09-18 & \href{https://youtube.com/watch?v=rllwPQsYaK0}{\texttt{rllwPQsYaK0}} \\
士師記 17:1-13 & \hyperref[sec:jXXnWoB7EBk]{2023年9月24日崇拜直播｜陳明泉牧師｜「家醜說出口」系列:私有化的祭司｜士師記17\_1-13} & 2023-09-26 & \href{https://youtube.com/watch?v=jXXnWoB7EBk}{\texttt{jXXnWoB7EBk}} \\
士師記 19:1-30 & \hyperref[sec:_4zbsWU6e2A]{2023年10月1日崇拜直播｜陳明泉牧師｜「家醜說出口」系列:沒有上帝的麻木｜士師記19\_1-30} & 2023-10-03 & \href{https://youtube.com/watch?v=-4zbsWU6e2A}{\texttt{-4zbsWU6e2A}} \\
約翰福音 3:1-15 & \hyperref[sec:e5U96ffu0L0]{2023年10月15日崇拜直播｜梁家麟牧師｜耶穌談重生｜約翰福音3\_1-15} & 2023-10-17 & \href{https://youtube.com/watch?v=e5U96ffu0L0}{\texttt{e5U96ffu0L0}} \\
士師記 21:1-15 & \hyperref[sec:jwnc_J9h4aM]{2023年10月29日崇拜直播｜陳冠成牧師｜「家醜說出口」系列:問題解決了嗎?｜士師記21\_1-15} & 2023-10-30 & \href{https://youtube.com/watch?v=jwnc-J9h4aM}{\texttt{jwnc-J9h4aM}} \\
撒母耳記下 6:16-23 & \hyperref[sec:t62ZJvrM_ng]{2023年11月5日崇拜直播｜陳明泉牧師｜唯獨你是王 - 牧者君王｜撒母耳記下6\_16-23} & 2023-11-06 & \href{https://youtube.com/watch?v=t62ZJvrM-ng}{\texttt{t62ZJvrM-ng}} \\
以賽亞書 7:1-17 & \hyperref[sec:KCjIplGBEcc]{2023年11月12日崇拜直播｜陳冠成牧師｜唯獨你是王 - 以馬內利的兆頭｜以賽亞書7\_1-17} & 2023-11-14 & \href{https://youtube.com/watch?v=KCjIplGBEcc}{\texttt{KCjIplGBEcc}} \\
約書亞記 14:6-14 & \hyperref[sec:SCWtMBw63Gk]{2023年11月19日長者主日崇拜直播｜羅志雄牧師｜老而彌堅的迦勒｜約書亞記14\_6-14} & 2023-11-20 & \href{https://youtube.com/watch?v=SCWtMBw63Gk}{\texttt{SCWtMBw63Gk}} \\
約翰福音 3:1-15 & \hyperref[sec:_oDlWIvKGkE]{2023年11月26日崇拜直播｜梁家麟牧師｜你們要重生｜約翰福音3\_1-15} & 2023-11-28 & \href{https://youtube.com/watch?v=-oDlWIvKGkE}{\texttt{-oDlWIvKGkE}} \\
以賽亞書 36:1-12 & \hyperref[sec:4NVCY6ZFpaA]{2023年12月3日崇拜直播｜陳冠成牧師｜唯獨你是王 - 神必拯救我們｜以賽亞書36\_1-12} & 2023-12-04 & \href{https://youtube.com/watch?v=4NVCY6ZFpaA}{\texttt{4NVCY6ZFpaA}} \\
路加福音 9:51-56 & \hyperref[sec:LudzCxftFrQ]{2023年12月17日崇拜直播｜梁家麟牧師｜耶穌與撒馬利亞人｜路加福音9\_51-56} & 2023-12-19 & \href{https://youtube.com/watch?v=LudzCxftFrQ}{\texttt{LudzCxftFrQ}} \\
馬太福音 2:1-12 & \hyperref[sec:QUWqIbTXwuo]{2023年12月24日聖誕節主日崇拜直播｜陳冠成牧師｜唯獨你是王 - 耶穌帶來的驚與喜｜馬太福音2\_1-12} & 2023-12-25 & \href{https://youtube.com/watch?v=QUWqIbTXwuo}{\texttt{QUWqIbTXwuo}} \\
詩篇 65:1-13 & \hyperref[sec:iN3iqiIAXcs]{2023年12月31日崇拜直播｜陳明泉牧師｜你以恩典作年歲的冠冕｜詩篇65\_1-13} & 2024-01-01 & \href{https://youtube.com/watch?v=iN3iqiIAXcs}{\texttt{iN3iqiIAXcs}} \\
創世記 18:16-33 & \hyperref[sec:2Yp8rtZS_PM]{2024年1月7日崇拜直播｜陳明泉牧師｜十個深化關係的禱告-為別人懇求｜創世記18\_16-33} & 2024-01-08 & \href{https://youtube.com/watch?v=2Yp8rtZS-PM}{\texttt{2Yp8rtZS-PM}} \\
尼希米記 1:1-11 & \hyperref[sec:H22XGgjIH6M]{2024年1月28日崇拜直播｜陳冠成牧師｜十個深化關係的禱告-以神國優先｜尼希米記1\_1-11} & 2024-01-28 & \href{https://youtube.com/watch?v=H22XGgjIH6M}{\texttt{H22XGgjIH6M}} \\
歷代志上 4:9-10 & \hyperref[sec:_RD54XIpk84]{2024年2月4日崇拜直播｜陳明泉牧師｜十個深化關係的禱告-逆轉痛苦的禱告｜歷代志上4\_9-10} & 2024-02-06 & \href{https://youtube.com/watch?v=-RD54XIpk84}{\texttt{-RD54XIpk84}} \\
撒母耳記上 1:4-20 & \hyperref[sec:xitIkhwfBC8]{2024年2月18日崇拜直播｜吳真真牧師｜十個深化關係的禱告-愁苦中的禱告｜撒母耳記上1\_4-20} & 2024-02-18 & \href{https://youtube.com/watch?v=xitIkhwfBC8}{\texttt{xitIkhwfBC8}} \\
約書亞記 4:1-24 & \hyperref[sec:1GcKc4lpPUQ]{2022年9月25日崇拜直播｜陳啟興牧師｜以石為記:這些石頭是什麼意思?｜約書亞記4\_1-24} & 2024-03-13 & \href{https://youtube.com/watch?v=1GcKc4lpPUQ}{\texttt{1GcKc4lpPUQ}} \\
箴言 30:1-14 & \hyperref[sec:8EXqdqFDCIc]{2024年3月17日崇拜直播｜陳明泉牧師｜十個深化關係的禱告-智者的禱告｜箴言30\_1-14} & 2024-03-18 & \href{https://youtube.com/watch?v=8EXqdqFDCIc}{\texttt{8EXqdqFDCIc}} \\
希伯來書 10:19-11:3 & \hyperref[sec:3GKCAsbLUbA]{2022年5月1日崇拜直播｜陳明泉牧師｜超越的基督:親近上帝之路｜希伯來書10\_19-11\_3} & 2024-03-19 & \href{https://youtube.com/watch?v=3GKCAsbLUbA}{\texttt{3GKCAsbLUbA}} \\
路加福音 18:9-14 & \hyperref[sec:l8_t1yVcPHE]{2024年3月24日崇拜直播｜余嘉敏姑娘｜十個深化關係的禱告-稅吏的禱告｜路加福音18\_9-14} & 2024-03-24 & \href{https://youtube.com/watch?v=l8-t1yVcPHE}{\texttt{l8-t1yVcPHE}} \\
約翰一書 3:13-23 & \hyperref[sec:1_gQDpS_bXo]{2021年5月9日母親節崇拜直播｜陳嘉璐醫生｜愛在家中｜約翰一書3\_13-23} & 2024-03-27 & \href{https://youtube.com/watch?v=1-gQDpS_bXo}{\texttt{1-gQDpS\_bXo}} \\
使徒行傳 2:37-47 & \hyperref[sec:7e2_k14U9kk]{2023年1月22日崇拜直播｜陳明泉牧師｜旺丁旺才的家｜使徒行傳2\_37-47} & 2024-03-27 & \href{https://youtube.com/watch?v=7e2-k14U9kk}{\texttt{7e2-k14U9kk}} \\
  & \hyperref[sec:DVdI0qIQgiA]{2024年3月29日 受苦節主餐崇拜直播｜仰望聖潔被殺羔羊} & 2024-03-30 & \href{https://youtube.com/watch?v=DVdI0qIQgiA}{\texttt{DVdI0qIQgiA}} \\
馬太福音 27:62-28:15 & \hyperref[sec:g_Q6zBFvts0]{2024年3月31日復活節崇拜直播｜陳明泉牧師｜無法封住的真相｜馬太福音27\_62-28\_15} & 2024-03-31 & \href{https://youtube.com/watch?v=g_Q6zBFvts0}{\texttt{g\_Q6zBFvts0}} \\
路加福音 24:13-35 & \hyperref[sec:JykGqUs7644]{2024年4月14日崇拜直播｜陳明泉牧師｜十個深化關係的禱告-拿起餅祝謝的禱告｜路加福音24\_13-35} & 2024-04-15 & \href{https://youtube.com/watch?v=JykGqUs7644}{\texttt{JykGqUs7644}} \\
列王記下 6:14-17 & \hyperref[sec:IOUGpe3_PlM]{2024年4月21日青少主日崇拜直播｜馬嘉言姑娘｜讓青少年親眼看見神｜列王記下6\_14-17} & 2024-04-21 & \href{https://youtube.com/watch?v=IOUGpe3-PlM}{\texttt{IOUGpe3-PlM}} \\
使徒行傳 7:51-60 & \hyperref[sec:K_X10pNASkM]{2024年5月5日崇拜直播｜陳明泉牧師｜十個深化關係的禱告-學禱告學做人｜使徒行傳7\_51-60} & 2024-05-07 & \href{https://youtube.com/watch?v=K-X10pNASkM}{\texttt{K-X10pNASkM}} \\
詩篇申命記以賽亞書 22:9-10 & \hyperref[sec:CNrc6_M4HZY]{2024年5月12日崇拜直播｜區祥江博士｜神像母親養育我們｜詩篇22\_9-10;以賽亞書49\_15,66\_12-13;申命記32\_11} & 2024-05-13 & \href{https://youtube.com/watch?v=CNrc6-M4HZY}{\texttt{CNrc6-M4HZY}} \\
創世記 28:10-22 & \hyperref[sec:OGOXjFxluCg]{2022年9月4日崇拜直播｜陳明泉牧師｜以石為記:伯特利的枕首石｜創世記28\_10-22} & 2024-05-26 & \href{https://youtube.com/watch?v=OGOXjFxluCg}{\texttt{OGOXjFxluCg}} \\
腓立比書 1:1-11 & \hyperref[sec:1HAPF8xhKJg]{2024年5月26日崇拜直播｜陳冠成牧師｜靈性·關係·成長－健康的屬靈關係｜腓立比書1\_1-11} & 2024-05-27 & \href{https://youtube.com/watch?v=1HAPF8xhKJg}{\texttt{1HAPF8xhKJg}} \\
希伯來書 12:1-3 & \hyperref[sec:OamF1FcoO7k]{2020年4月10日 受苦節崇拜直播｜希伯來書12\_1-3} & 2024-06-06 & \href{https://youtube.com/watch?v=OamF1FcoO7k}{\texttt{OamF1FcoO7k}} \\
腓立比書 1:27-2:11 & \hyperref[sec:9ogZwUemHN4]{2024年6月16日父親節崇拜直播｜陳冠成牧師｜靈性·關係·成長－與身份相稱的生活｜腓立比書1\_27-2\_11} & 2024-06-16 & \href{https://youtube.com/watch?v=9ogZwUemHN4}{\texttt{9ogZwUemHN4}} \\
約翰福音 4:5-15 & \hyperref[sec:0Rt7bY88Mpk]{2024年6月9日崇拜直播｜梁家麟牧師｜為有源頭活水來｜約翰福音4\_5-15} & 2024-06-16 & \href{https://youtube.com/watch?v=0Rt7bY88Mpk}{\texttt{0Rt7bY88Mpk}} \\
腓立比書 2:12-18 & \hyperref[sec:_sNrPBKmvwc]{2024年6月23日崇拜直播｜羅志雄牧師｜靈性·關係·成長－活潑的生命｜腓立比書2\_12-18} & 2024-06-24 & \href{https://youtube.com/watch?v=_sNrPBKmvwc}{\texttt{\_sNrPBKmvwc}} \\
腓立比書 2:19-30 & \hyperref[sec:l1Zxe_Q_Nro]{2024年6月30日崇拜直播｜陳啟興牧師｜靈性·關係·成長－工作不離三兄弟｜腓立比書2\_19-30} & 2024-06-30 & \href{https://youtube.com/watch?v=l1Zxe_Q-Nro}{\texttt{l1Zxe\_Q-Nro}} \\
腓立比書 3:1-16 & \hyperref[sec:8t3Xeh3MCK8]{2024年7月7日主餐崇拜直播｜陳明泉牧師｜靈性·關係·成長－甚麼偷走了我的快樂｜腓立比書3\_1-16} & 2024-07-10 & \href{https://youtube.com/watch?v=8t3Xeh3MCK8}{\texttt{8t3Xeh3MCK8}} \\
約翰福音 5:18-47 & \hyperref[sec:BxDFSB8mRZw]{2024年7月21日崇拜直播｜梁家麟牧師｜與父上帝平等的耶穌｜約翰福音5\_18-47} & 2024-07-23 & \href{https://youtube.com/watch?v=BxDFSB8mRZw}{\texttt{BxDFSB8mRZw}} \\
腓立比書 4:1-9 & \hyperref[sec:rzRD6n7AW0w]{2024年7月28日崇拜直播｜陳冠成牧師｜靈性·關係·成長－靈性安穩與探索｜腓立比書4\_1-9} & 2024-07-29 & \href{https://youtube.com/watch?v=rzRD6n7AW0w}{\texttt{rzRD6n7AW0w}} \\
耶利米書 20:7-18 & \hyperref[sec:OxtvYBdmNrI]{2021年6月13日崇拜直播｜陳冠成牧師｜從聖經看喜怒哀愁 - 告別憂愁｜耶利米書20\_7-18} & 2024-08-06 & \href{https://youtube.com/watch?v=OxtvYBdmNrI}{\texttt{OxtvYBdmNrI}} \\
民數記 12:1-15 & \hyperref[sec:Hufazqfe3oE]{2024年8月4日崇拜直播｜陳冠成牧師｜正視你的情緒－嫉妒惹的禍｜民數記12\_1-15} & 2024-08-06 & \href{https://youtube.com/watch?v=Hufazqfe3oE}{\texttt{Hufazqfe3oE}} \\
約翰福音 6:1-14 & \hyperref[sec:0SAEU456WxE]{2024年8月11日崇拜直播｜梁家麟牧師｜吃飽了不餓｜約翰福音6\_1-14} & 2024-08-19 & \href{https://youtube.com/watch?v=0SAEU456WxE}{\texttt{0SAEU456WxE}} \\
撒母耳記上 10:9-13 & \hyperref[sec:Ef4j6FO5fXg]{2024年8月18日崇拜直播｜陳明泉牧師｜正視你的情緒－焦慮的糾纏｜撒母耳記上10\_9-13;19\_18-24} & 2024-08-19 & \href{https://youtube.com/watch?v=Ef4j6FO5fXg}{\texttt{Ef4j6FO5fXg}} \\
詩篇 119:105-112 & \hyperref[sec:cerRyXhHqiE]{2024年8月25日崇拜直播｜張培良傳道｜活在上帝的光中｜詩篇119\_105-112} & 2024-08-26 & \href{https://youtube.com/watch?v=cerRyXhHqiE}{\texttt{cerRyXhHqiE}} \\
創世記 16:1-16 & \hyperref[sec:lCXbCQNDTgA]{2024年9月1日崇拜直播｜吳真真牧師｜正視你的情緒－受屈者的出路｜創世記16\_1-16} & 2024-09-02 & \href{https://youtube.com/watch?v=lCXbCQNDTgA}{\texttt{lCXbCQNDTgA}} \\
腓立比書 4:1-9 & \hyperref[sec:uj5FB161Gvc]{2024年2月11日新春崇拜直播｜梁家麟牧師｜一個叮嚀三個祝願｜腓立比書4\_1-9} & 2024-09-16 & \href{https://youtube.com/watch?v=uj5FB161Gvc}{\texttt{uj5FB161Gvc}} \\
約書亞記 1:1-9 & \hyperref[sec:QTbAXdDuZOw]{2024年9月15日崇拜直播｜陳冠成牧師｜正視你的情緒－從神而來的勇氣｜約書亞記1\_1-9} & 2024-09-16 & \href{https://youtube.com/watch?v=QTbAXdDuZOw}{\texttt{QTbAXdDuZOw}} \\
希伯來書 4:1-11 & \hyperref[sec:QHvBwcg9HEw]{2022年3月13日崇拜直播｜羅志雄牧師｜超越的基督:進入安息｜希伯來書4\_1-11} & 2024-09-19 & \href{https://youtube.com/watch?v=QHvBwcg9HEw}{\texttt{QHvBwcg9HEw}} \\
撒母耳記下 12:15-23 & \hyperref[sec:U7nTwzqoXV8]{2024年9月22日崇拜直播｜陳明泉牧師｜正視你的情緒－為悲傷劃界線｜撒母耳記下12\_15-23} & 2024-09-23 & \href{https://youtube.com/watch?v=U7nTwzqoXV8}{\texttt{U7nTwzqoXV8}} \\
約翰福音 6:24 & \hyperref[sec:JCVpdjhfj4w]{2024年10月6日崇拜直播｜梁家麟牧師｜生命的糧｜約翰福音6\_24 - 59} & 2024-10-11 & \href{https://youtube.com/watch?v=JCVpdjhfj4w}{\texttt{JCVpdjhfj4w}} \\
耶利米書 15:15-21 & \hyperref[sec:hd2ElWc7xho]{2024年10月27日崇拜直播｜陳明泉牧師｜正視你的情緒－當你感到自憐時......｜耶利米書15\_15-21} & 2024-10-27 & \href{https://youtube.com/watch?v=hd2ElWc7xho}{\texttt{hd2ElWc7xho}} \\
詩篇使徒行傳 2:1-9 & \hyperref[sec:limrMw42910]{2024年11月3日崇拜直播｜余嘉敏傳道｜詩篇裡的基督－受膏僕人｜詩篇2\_1-9;使徒行傳4\_23-31} & 2024-11-04 & \href{https://youtube.com/watch?v=limrMw42910}{\texttt{limrMw42910}} \\
約翰福音 6:60-7:13 & \hyperref[sec:uf_k9SIAd7s]{2024年11月17日崇拜直播｜梁家麟牧師｜被人厭棄的耶穌｜約翰福音6\_60-7\_13} & 2024-11-20 & \href{https://youtube.com/watch?v=uf_k9SIAd7s}{\texttt{uf\_k9SIAd7s}} \\
希伯來書詩篇馬太福音 95:6-11 & \hyperref[sec:230hZbWKaH0]{2024年11月24日崇拜直播｜陳冠成牧師｜詩篇裡的基督－安息之主｜詩篇95\_6-11;希伯來書4\_1-2;馬太福音11\_28-30} & 2024-11-24 & \href{https://youtube.com/watch?v=230hZbWKaH0}{\texttt{230hZbWKaH0}} \\
希伯來書詩篇 102:3-28 & \hyperref[sec:a62m7d20RCU]{2024年12月1日主餐崇拜直播｜陳明泉牧師｜詩篇裡的基督－永活神子｜詩篇102\_3-28;希伯來書1\_10-12} & 2024-12-02 & \href{https://youtube.com/watch?v=a62m7d20RCU}{\texttt{a62m7d20RCU}} \\
詩篇馬太福音 118:19-26 & \hyperref[sec:_MkLayWrDNI]{2024年12月8日崇拜直播｜陳冠成牧師｜詩篇裡的基督－房角石頭｜詩篇118\_19-26;馬太福音20\_33-46} & 2024-12-09 & \href{https://youtube.com/watch?v=_MkLayWrDNI}{\texttt{\_MkLayWrDNI}} \\
約翰福音 2:1-12 & \hyperref[sec:jH3ZJLIEHsY]{2023年2月19日崇拜直播｜梁家麟牧師｜以水變酒的神蹟｜約翰福音2\_1-12} & 2024-12-12 & \href{https://youtube.com/watch?v=jH3ZJLIEHsY}{\texttt{jH3ZJLIEHsY}} \\
詩篇使徒行傳 110:1-7 & \hyperref[sec:KLXa_sv_cYI]{2024年12月15日崇拜直播｜楊珊珊傳道｜詩篇裡的基督－永遠為祭司的君王｜詩篇110\_1-7;使徒行傳2\_29-40} & 2024-12-15 & \href{https://youtube.com/watch?v=KLXa_sv_cYI}{\texttt{KLXa\_sv\_cYI}} \\
路加福音彌迦書 5:2-5 & \hyperref[sec:75p0crrJvGI]{2024年12月22日聖誕節主日暨嬰孩奉獻崇拜直播｜陳明泉牧師｜萬世之君誕於小城｜彌迦書5\_2-5;路加福音2\_1-20} & 2024-12-22 & \href{https://youtube.com/watch?v=75p0crrJvGI}{\texttt{75p0crrJvGI}} \\
約翰福音 7:14-31 & \hyperref[sec:NqwhnMUzemI]{2024年12月29日崇拜直播｜梁家麟牧師｜令人不安的耶穌｜約翰福音7\_14-31} & 2024-12-30 & \href{https://youtube.com/watch?v=NqwhnMUzemI}{\texttt{NqwhnMUzemI}} \\
申命記 1:19-46 & \hyperref[sec:H2Yyyhsu_Qc]{2025年1月5日主餐崇拜直播｜陳明泉牧師｜朝向應許之地－要懂得選擇｜申命記1\_19-46} & 2025-01-06 & \href{https://youtube.com/watch?v=H2Yyyhsu_Qc}{\texttt{H2Yyyhsu\_Qc}} \\
約翰福音 7:32-52 & \hyperref[sec:YgvzlIV1PPE]{2025年1月12日崇拜直播｜梁家麟牧師｜令人不安的耶穌(二)｜約翰福音7\_32-52} & 2025-01-12 & \href{https://youtube.com/watch?v=YgvzlIV1PPE}{\texttt{YgvzlIV1PPE}} \\
約拿書 4:1-11 & \hyperref[sec:uiuY8tue_2Q]{2020年11月29日 崇拜直播｜約拿書4\_1-11} & 2025-01-17 & \href{https://youtube.com/watch?v=uiuY8tue-2Q}{\texttt{uiuY8tue-2Q}} \\
申命記 3:18-28 & \hyperref[sec:KYvJDDo3N1U]{2025年1月26日崇拜直播｜鄭國邦傳道｜朝向應許之地－寬容與堅拒｜申命記3\_18-28} & 2025-01-26 & \href{https://youtube.com/watch?v=KYvJDDo3N1U}{\texttt{KYvJDDo3N1U}} \\
申命記 4:1-14 & \hyperref[sec:cuIe7N8imyc]{2025年2月2日主餐崇拜直播｜陳明泉牧師｜朝向應許之地－安居樂業的願景｜申命記4\_1-14} & 2025-02-02 & \href{https://youtube.com/watch?v=cuIe7N8imyc}{\texttt{cuIe7N8imyc}} \\
約翰福音  & \hyperref[sec:uDlv_6lnHKk]{2025年2月9日崇拜直播｜梁家麟牧師｜我也不定你的罪｜約翰福音 八 1-11} & 2025-02-09 & \href{https://youtube.com/watch?v=uDlv-6lnHKk}{\texttt{uDlv-6lnHKk}} \\
申命記 4:15-31 & \hyperref[sec:oex63GVBJWc]{2025年2月16日崇拜直播｜陳冠成牧師｜朝向應許之地－禁戒雕刻偶像｜申命記4\_15-31} & 2025-02-16 & \href{https://youtube.com/watch?v=oex63GVBJWc}{\texttt{oex63GVBJWc}} \\
申命記 5:1-6:9 & \hyperref[sec:bBtvGdEATA0]{2025年2月23日青少主日崇拜直播｜楊珊珊傳道｜朝向應許之地－十誡與大誡命｜申命記5\_1-6\_9} & 2025-02-23 & \href{https://youtube.com/watch?v=bBtvGdEATA0}{\texttt{bBtvGdEATA0}} \\
申命記 7:1-11 & \hyperref[sec:C0mf0rNYQN4]{2025年3月2日主餐崇拜直播｜陳明泉牧師｜朝向應許之地－專一｜申命記7\_1-11} & 2025-03-03 & \href{https://youtube.com/watch?v=C0mf0rNYQN4}{\texttt{C0mf0rNYQN4}} \\
  & \hyperref[sec:mfCFFhs9u_c]{2025年3月6日預苦期首日塗灰禮晚禱崇拜直播} & 2025-03-06 & \href{https://youtube.com/watch?v=mfCFFhs9u-c}{\texttt{mfCFFhs9u-c}} \\
申命記 8:1-10 & \hyperref[sec:crDuuHRkfCw]{2025年3月9日崇拜直播｜陳冠成牧師｜朝向應許之地－謹記遵守盟約｜申命記8\_1-10} & 2025-03-09 & \href{https://youtube.com/watch?v=crDuuHRkfCw}{\texttt{crDuuHRkfCw}} \\
申命記  & \hyperref[sec:P0a0tJAFKd4]{2025年3月16日崇拜直播｜羅志雄牧師｜朝向應許之地－恩典中的悔改與更新｜申命記918-29} & 2025-03-16 & \href{https://youtube.com/watch?v=P0a0tJAFKd4}{\texttt{P0a0tJAFKd4}} \\
約翰福音 21:20-23 & \hyperref[sec:LKA6QvIWId8]{2021年4月25日崇拜直播｜陳明泉牧師｜與你何干?｜約翰福音21\_20-23} & 2025-03-18 & \href{https://youtube.com/watch?v=LKA6QvIWId8}{\texttt{LKA6QvIWId8}} \\
申命記 14:1-29 & \hyperref[sec:aQhmLPtYC40]{2025年3月23日崇拜直播｜陳明泉牧師｜從生活細節,區分特殊身份｜申命記14\_1-29} & 2025-03-23 & \href{https://youtube.com/watch?v=aQhmLPtYC40}{\texttt{aQhmLPtYC40}} \\
馬太福音 4:24-5:3 & \hyperref[sec:P1VP0bkSNZc]{2025年3月30日崇拜直播｜張永信牧師｜虛心蒙福｜馬太福音4\_24-5\_3} & 2025-03-30 & \href{https://youtube.com/watch?v=P1VP0bkSNZc}{\texttt{P1VP0bkSNZc}} \\
申命記 15:1-18 & \hyperref[sec:I13jDOQMkLQ]{2025年4月6日主餐崇拜直播｜陳明泉牧師｜朝向應許之地－顧念別人｜申命記15\_1-18} & 2025-04-07 & \href{https://youtube.com/watch?v=I13jDOQMkLQ}{\texttt{I13jDOQMkLQ}} \\
申命記 16:1-17 & \hyperref[sec:CrZ2VP_EO18]{2025年4月13日崇拜直播｜陳冠成牧師｜朝向應許之地－循環生活中記念上帝｜申命記16\_1-17} & 2025-04-13 & \href{https://youtube.com/watch?v=CrZ2VP_EO18}{\texttt{CrZ2VP\_EO18}} \\
  & \hyperref[sec:96uE5WjMIVg]{2025年4月18日受苦節崇拜直播} & 2025-04-19 & \href{https://youtube.com/watch?v=96uE5WjMIVg}{\texttt{96uE5WjMIVg}} \\
羅馬書 8:31 & \hyperref[sec:9gw1hxqzLeo]{2025年4月20日復活節崇拜直播｜陳明泉牧師｜勝利的體會｜羅馬書 8\_31 -39} & 2025-04-20 & \href{https://youtube.com/watch?v=9gw1hxqzLeo}{\texttt{9gw1hxqzLeo}} \\
創世記 22:1-14 & \hyperref[sec:eeqdwBmS648]{2025年4月27日崇拜直播｜林嘉鴻牧師｜生命的最重要｜創世記22\_1-14} & 2025-04-27 & \href{https://youtube.com/watch?v=eeqdwBmS648}{\texttt{eeqdwBmS648}} \\
  & \hyperref[sec:u7h0Nna94Mk]{2025年4月27日主日講道｜胡康偉牧師:天國的筵席與和平之子} & 2025-04-28 & \href{https://youtube.com/watch?v=u7h0Nna94Mk}{\texttt{u7h0Nna94Mk}} \\
申命記 24:1-22 & \hyperref[sec:Wk72Zd8NFhw]{2025年5月4日主餐崇拜直播｜陳明泉牧師｜朝向應許之地－情理兼備的國度｜申命記 24\_1-22} & 2025-05-04 & \href{https://youtube.com/watch?v=Wk72Zd8NFhw}{\texttt{Wk72Zd8NFhw}} \\
約翰三書使徒行傳 16:31-34 & \hyperref[sec:eFUsT6910pQ]{2025年5月11日母親節主日崇拜直播｜吳真真牧師｜喜樂家庭｜使徒行傳16\_31-34;約翰三書1\_4} & 2025-05-11 & \href{https://youtube.com/watch?v=eFUsT6910pQ}{\texttt{eFUsT6910pQ}} \\
使徒行傳 11:1-18 & \hyperref[sec:6xido1UHEiY]{2025年5月18日崇拜直播｜鄭家恩牧師｜摻埋我玩｜使徒行傳11\_1-18} & 2025-05-18 & \href{https://youtube.com/watch?v=6xido1UHEiY}{\texttt{6xido1UHEiY}} \\
申命記 26:16-19 & \hyperref[sec:wyPYrchtH6g]{2025年5月25日崇拜直播｜鄭國邦傳道｜朝向應許之地－自由與順服｜申命記26\_16-19} & 2025-05-25 & \href{https://youtube.com/watch?v=wyPYrchtH6g}{\texttt{wyPYrchtH6g}} \\
申命記 30:11-20 & \hyperref[sec:dxgI8DcWOBk]{2025年6月1日主餐崇拜直播｜陳冠成牧師｜朝向應許之地－務要揀選生命｜申命記30\_11-20} & 2025-06-01 & \href{https://youtube.com/watch?v=dxgI8DcWOBk}{\texttt{dxgI8DcWOBk}} \\
申命記 31:1-8 & \hyperref[sec:Buo5z3C6Kgc]{2025年6月8日崇拜直播｜陳明泉牧師｜朝向應許之地－你們不孤單｜申命記31\_1-8} & 2025-06-08 & \href{https://youtube.com/watch?v=Buo5z3C6Kgc}{\texttt{Buo5z3C6Kgc}} \\
約書亞記 1:1-9 & \hyperref[sec:f__YU0k5QKA]{2025年6月15日父親節主日崇拜直播｜劉佩婷博士｜50歲的剛強壯膽｜約書亞記1\_1-9} & 2025-06-15 & \href{https://youtube.com/watch?v=f-_YU0k5QKA}{\texttt{f-\_YU0k5QKA}} \\
約翰福音 8:12-30 & \hyperref[sec:EfFf_ZeLMAM]{2025年6月22日崇拜直播｜梁家麟牧師｜耶穌與天父的關係｜約翰福音8\_12-30} & 2025-06-22 & \href{https://youtube.com/watch?v=EfFf_ZeLMAM}{\texttt{EfFf\_ZeLMAM}} \\
申命記 32:45-52 & \hyperref[sec:12V0_gKaK_A]{2025年6月29日崇拜直播｜楊珊珊傳道｜朝向應許之地－話別與離別｜申命記32\_45-52;34\_1-12} & 2025-06-29 & \href{https://youtube.com/watch?v=12V0_gKaK_A}{\texttt{12V0\_gKaK\_A}} \\
馬可福音 6:1-13 & \hyperref[sec:cD5y395bElI]{2025年7月6日崇拜直播｜陳明泉牧師｜耶穌給門徒的十課－學懂被拒絕｜馬可福音6\_1-13} & 2025-07-06 & \href{https://youtube.com/watch?v=cD5y395bElI}{\texttt{cD5y395bElI}} \\
馬可福音 6:14-29 & \hyperref[sec:OaFcCMZ21q4]{2025年7月13日崇拜直播｜鄭國邦傳道｜耶穌給門徒的十課－付代價｜馬可福音6\_14-29} & 2025-07-13 & \href{https://youtube.com/watch?v=OaFcCMZ21q4}{\texttt{OaFcCMZ21q4}} \\
馬可福音 6:30-44 & \hyperref[sec:SbKqVXHzrXE]{2025年7月20日崇拜直播｜陳冠成牧師｜耶穌給門徒的十課－給眾人吃飽｜馬可福音6\_30-44} & 2025-07-21 & \href{https://youtube.com/watch?v=SbKqVXHzrXE}{\texttt{SbKqVXHzrXE}} \\
約翰福音 8:31-59 & \hyperref[sec:gyB32xPImr0]{2025年7月27日崇拜直播｜梁家麟牧師｜放棄耶穌的兩大理由｜約翰福音8\_31-59} & 2025-07-27 & \href{https://youtube.com/watch?v=gyB32xPImr0}{\texttt{gyB32xPImr0}} \\
馬可福音 6:45-56 & \hyperref[sec:woNoWOUM52k]{2025年8月3日崇拜直播｜吳真真牧師｜耶穌給門徒的十課－祂是誰?｜馬可福音6\_45-56} & 2025-08-04 & \href{https://youtube.com/watch?v=woNoWOUM52k}{\texttt{woNoWOUM52k}} \\
約翰福音 9:1-41 & \hyperref[sec:qauwj9FGRZc]{2025年8月10日崇拜直播｜梁家麟牧師｜彰顯上帝的作為｜約翰福音9\_1-41} & 2025-08-11 & \href{https://youtube.com/watch?v=qauwj9FGRZc}{\texttt{qauwj9FGRZc}} \\
馬可福音 7:24-37 & \hyperref[sec:iXL9M6i3TXE]{2025年8月17日崇拜直播｜鄭國邦傳道｜耶穌給門徒的十課－離開群眾,真認識主｜馬可福音7\_24-37} & 2025-08-18 & \href{https://youtube.com/watch?v=iXL9M6i3TXE}{\texttt{iXL9M6i3TXE}} \\
馬可福音 7:1-23 & \hyperref[sec:0QX__zW2yAY]{2025年8月24日崇拜直播｜陳冠成牧師｜耶穌給門徒的十課－傳道vs誡命｜馬可福音7\_1-23} & 2025-08-30 & \href{https://youtube.com/watch?v=0QX__zW2yAY}{\texttt{0QX\_\_zW2yAY}} \\
馬可福音 8:1-21 & \hyperref[sec:EMIvj5LfGAU]{2025年8月31日崇拜直播｜楊珊珊傳道｜耶穌給門徒的十課－你們還不明白嗎?｜馬可福音8\_1-21} & 2025-08-31 & \href{https://youtube.com/watch?v=EMIvj5LfGAU}{\texttt{EMIvj5LfGAU}} \\
馬可福音 8:22-9:1 & \hyperref[sec:xAjegbjPwhU]{2025年9月7日崇拜直播｜陳冠成牧師｜耶穌給門徒的十課－基督的本質•門徒的職份｜馬可福音8\_22-9\_1} & 2025-09-07 & \href{https://youtube.com/watch?v=xAjegbjPwhU}{\texttt{xAjegbjPwhU}} \\
馬可福音 9:14-29 & \hyperref[sec:HtBaZ5KPIHs]{2025年9月14日崇拜直播｜陳明泉牧師｜耶穌給門徒的十課－在人不能時......｜馬可福音9\_14-29} & 2025-09-14 & \href{https://youtube.com/watch?v=HtBaZ5KPIHs}{\texttt{HtBaZ5KPIHs}} \\
約翰福音 9:8-41 & \hyperref[sec:LeA_EIgcwcw]{2025年9月21日崇拜直播｜梁家麟牧師｜瞎子信耶穌的故事｜約翰福音9\_8-41} & 2025-09-23 & \href{https://youtube.com/watch?v=LeA_EIgcwcw}{\texttt{LeA\_EIgcwcw}} \\
馬可福音 10:32-45 & \hyperref[sec:JJ124JmNAMw]{2025年9月28日崇拜直播｜吳真真牧師｜耶穌給門徒的十課－走僕人之路｜馬可福音 10\_32-45} & 2025-09-28 & \href{https://youtube.com/watch?v=JJ124JmNAMw}{\texttt{JJ124JmNAMw}} \\
啟示錄 2:1-7 & \hyperref[sec:Jhd45blKy7w]{2025年10月5日崇拜直播｜陳明泉牧師｜耶穌給教會的七封信:起初的愛｜啟示錄2\_1-7} & 2025-10-05 & \href{https://youtube.com/watch?v=Jhd45blKy7w}{\texttt{Jhd45blKy7w}} \\
啟示錄 2:8-11 & \hyperref[sec:mfLhrRj9Qu0]{2025年10月12日崇拜直播｜陳冠成牧師｜耶穌給教會的七封信:至死忠心｜啟示錄2\_8-11} & 2025-10-12 & \href{https://youtube.com/watch?v=mfLhrRj9Qu0}{\texttt{mfLhrRj9Qu0}} \\
約翰福音 10:1-18 & \hyperref[sec:980zAvvXkiM]{2025年10月19日崇拜直播｜梁家麟牧師｜耶穌是好牧人｜約翰福音10\_1-18} & 2025-10-19 & \href{https://youtube.com/watch?v=980zAvvXkiM}{\texttt{980zAvvXkiM}} \\
啟示錄 2:12-17 & \hyperref[sec:SlVA2_DGfXk]{2025年10月26日崇拜直播｜羅志雄牧師｜耶穌給教會的七封信:當悔改｜啟示錄2\_12-17} & 2025-11-03 & \href{https://youtube.com/watch?v=SlVA2_DGfXk}{\texttt{SlVA2\_DGfXk}} \\
路加福音 17:11-19 & \hyperref[sec:LTkMkPBreS4]{2025年11月2日崇拜直播｜羅潔盈博士｜救恩的距離｜路加福音17\_11-19} & 2025-11-03 & \href{https://youtube.com/watch?v=LTkMkPBreS4}{\texttt{LTkMkPBreS4}} \\
詩篇 92:12-15 & \hyperref[sec:5pSwQwIJ1AI]{2025年11月9日長者主日崇拜直播｜區祥江博士｜老年仍發旺-栽在聖殿中的長青樹｜詩篇92\_12-15} & 2025-11-09 & \href{https://youtube.com/watch?v=5pSwQwIJ1AI}{\texttt{5pSwQwIJ1AI}} \\
啟示錄 2:18-29 & \hyperref[sec:H6vb7uu46Wk]{2025年11月16日崇拜直播｜陳明泉牧師｜耶穌給教會的七封信:提防假先知｜啟示錄2\_18-29} & 2025-11-17 & \href{https://youtube.com/watch?v=H6vb7uu46Wk}{\texttt{H6vb7uu46Wk}} \\
約翰福音 10:19-42 & \hyperref[sec:Ndl2CyA1nlc]{2025年11月23日崇拜直播｜梁家麟牧師｜信耶穌到甚麼地步?｜約翰福音10\_19-42} & 2025-11-29 & \href{https://youtube.com/watch?v=Ndl2CyA1nlc}{\texttt{Ndl2CyA1nlc}} \\
啟示錄 3:1-6 & \hyperref[sec:XMVPsTUvcaY]{2025年11月30日崇拜直播｜陳冠成牧師｜耶穌給教會的七封信:務要儆醒｜啟示錄3\_1-6} & 2025-12-01 & \href{https://youtube.com/watch?v=XMVPsTUvcaY}{\texttt{XMVPsTUvcaY}} \\
啟示錄 3:7-13 & \hyperref[sec:XJJ1wNBmbSY]{2025年12月7日主餐崇拜直播｜陳明泉牧師｜耶穌給教會的七封信:謹守與保守｜啟示錄3\_7-13} & 2025-12-07 & \href{https://youtube.com/watch?v=XJJ1wNBmbSY}{\texttt{XJJ1wNBmbSY}} \\
啟示錄 3:14-22 & \hyperref[sec:1zYtfRDiK4g]{2025年12月14日崇拜直播｜陳冠成牧師｜耶穌給教會的七封信:重拾愛主的心｜啟示錄3\_14-22} & 2025-12-14 & \href{https://youtube.com/watch?v=1zYtfRDiK4g}{\texttt{1zYtfRDiK4g}} \\
箴言 8:12-31 & \hyperref[sec:wYKTi6uXSJI]{2025年12月21日聖誕節主日崇拜直播｜高銘謙牧師｜愛智慧｜箴言8\_12-31} & 2025-12-21 & \href{https://youtube.com/watch?v=wYKTi6uXSJI}{\texttt{wYKTi6uXSJI}} \\
約翰福音 11:1-27 & \hyperref[sec:yqcP_iDkMnA]{2025年12月28日崇拜直播｜梁家麟牧師｜復活和生命在我｜約翰福音11\_1-27} & 2025-12-28 & \href{https://youtube.com/watch?v=yqcP_iDkMnA}{\texttt{yqcP\_iDkMnA}} \\
歌羅西書 1:1-14 & \hyperref[sec:DESvXfy__1I]{2026年1月4日主餐崇拜直播｜陳明泉牧師｜門徒DNA:福音介定的身份｜歌羅西書1\_1-14} & 2026-01-04 & \href{https://youtube.com/watch?v=DESvXfy__1I}{\texttt{DESvXfy\_\_1I}} \\
約翰福音 11:28-57 & \hyperref[sec:eZQl9Arirgg]{2026年1月11日崇拜直播｜梁家麟牧師｜相信復活的主｜約翰福音11\_28-57} & 2026-01-14 & \href{https://youtube.com/watch?v=eZQl9Arirgg}{\texttt{eZQl9Arirgg}} \\
歌羅西書 1:15-23 & \hyperref[sec:RDxRU6iPPzM]{2026年1月18日崇拜直播｜陳冠成牧師｜門徒DNA:連於元首基督｜歌羅西書1\_15-23} & 2026-01-18 & \href{https://youtube.com/watch?v=RDxRU6iPPzM}{\texttt{RDxRU6iPPzM}} \\
歌羅西書 1:24-2:5 & \hyperref[sec:sjw5HKR2iIk]{2026年1月25日崇拜直播｜吳真真牧師｜門徒DNA:在基督裡完全｜歌羅西書1\_24-2\_5} & 2026-01-25 & \href{https://youtube.com/watch?v=sjw5HKR2iIk}{\texttt{sjw5HKR2iIk}} \\
歌羅西書 2:6-23 & \hyperref[sec:_GdsDmJHJDI]{2026年2月1日主餐崇拜直播｜陳明泉牧師｜門徒DNA:在基督裡的智慧｜歌羅西書2\_6-23} & 2026-02-02 & \href{https://youtube.com/watch?v=_GdsDmJHJDI}{\texttt{\_GdsDmJHJDI}} \\
歌羅西書 3:1-15 & \hyperref[sec:397tpwFaRoA]{2026年2月8日崇拜直播｜陳冠成牧師｜門徒DNA:棄舊更新｜歌羅西書3\_1-15} & 2026-02-09 & \href{https://youtube.com/watch?v=397tpwFaRoA}{\texttt{397tpwFaRoA}} \\
歌羅西書 3:18-4:5 & \hyperref[sec:jemYOQSwajw]{2026年2月15青少主日崇拜直播｜楊珊珊傳道｜門徒DNA:我們這一家｜歌羅西書3\_18-4\_5} & 2026-02-16 & \href{https://youtube.com/watch?v=jemYOQSwajw}{\texttt{jemYOQSwajw}} \\
約翰福音 12:1-8 & \hyperref[sec:NhZRAevmHb4]{2026年2月22日崇拜直播｜梁家麟牧師｜腳的事奉｜約翰福音12\_1-8} & 2026-02-22 & \href{https://youtube.com/watch?v=NhZRAevmHb4}{\texttt{NhZRAevmHb4}} \\
詩篇 22:1-31 & \hyperref[sec:8qvQyuuLcyA]{2026年3月1日主餐崇拜直播｜陳明泉牧師｜基督之詩篇:我的神!為什麼離棄我?｜詩篇22\_1-31} & 2026-03-02 & \href{https://youtube.com/watch?v=8qvQyuuLcyA}{\texttt{8qvQyuuLcyA}} \\
馬可福音 4:1-20 & \hyperref[sec:yBE6W2fCqSM]{2026年3月8日崇拜直播｜黃偉立牧師｜門徒必備的那種態度｜馬可福音4\_1-20} & 2026-03-09 & \href{https://youtube.com/watch?v=yBE6W2fCqSM}{\texttt{yBE6W2fCqSM}} \\
詩篇 69:1-36 & \hyperref[sec:PGMHgAUtHoA]{2026年3月15日崇拜直播｜陳冠成牧師｜基督之詩篇:急難中的禱告｜詩篇69\_1-36} & 2026-03-16 & \href{https://youtube.com/watch?v=PGMHgAUtHoA}{\texttt{PGMHgAUtHoA}} \\
約書亞記 1:1-9 & \hyperref[sec:p8dj6FAIn5Y]{2026年3月22日崇拜直播｜蔡康怡牧師｜超越功能價值 立足聖約人生｜約書亞記1\_1-9} & 2026-03-24 & \href{https://youtube.com/watch?v=p8dj6FAIn5Y}{\texttt{p8dj6FAIn5Y}} \\
詩篇 16:1-11 & \hyperref[sec:3zRaVYVNQY0]{2026年3月29日崇拜直播｜鄭國邦傳道｜基督之詩篇:不讓你的聖者遇見朽壞｜詩篇16\_1-11} & 2026-03-30 & \href{https://youtube.com/watch?v=3zRaVYVNQY0}{\texttt{3zRaVYVNQY0}} \\
馬太福音 28:16-20 & \hyperref[sec:ltS5lDDLFjU]{2026年4月5日復活節暨主餐崇拜直播｜陳明泉牧師｜告別浮淺,門徒進深:疑惑裡的差遣｜馬太福音28\_16-20} & 2026-04-05 & \href{https://youtube.com/watch?v=ltS5lDDLFjU}{\texttt{ltS5lDDLFjU}} \\
馬可福音 1:35-39 & \hyperref[sec:rsvExPNsP2I]{2026年4月12日崇拜直播｜何梓峰傳道｜告別浮淺,門徒進深:退到曠野為要看到你｜馬可福音1\_35-39} & 2026-04-12 & \href{https://youtube.com/watch?v=rsvExPNsP2I}{\texttt{rsvExPNsP2I}} \\
馬太福音 8:27-35 & \hyperref[sec:Xc4fVvDZWRE]{2026年4月19日崇拜直播｜陳冠成牧師｜告別浮淺,門徒進深:拒絕成功主義｜馬太福音8\_27-35} & 2026-04-19 & \href{https://youtube.com/watch?v=Xc4fVvDZWRE}{\texttt{Xc4fVvDZWRE}} \\
歷代志上 17:15-27 & \hyperref[sec:3bVTrDHjdBs]{2026年4月26日崇拜直播｜鍾家威傳道｜告別浮淺,門徒進深:擁抱限制中的祝福｜歷代志上17\_15-27} & 2026-04-28 & \href{https://youtube.com/watch?v=3bVTrDHjdBs}{\texttt{3bVTrDHjdBs}} \\
哥林多後書 4:7-18 & \hyperref[sec:U29QDtiVI5U]{2026年5月3日主餐崇拜直播｜陳明泉牧師｜告別浮淺,門徒進深:發掘在失意裡的珍寶｜哥林多後書4\_7-18} & 2026-05-07 & \href{https://youtube.com/watch?v=U29QDtiVI5U}{\texttt{U29QDtiVI5U}} \\
路得記 1:19-22 & \hyperref[sec:wdtTzXf4txM]{2026年5月10日崇拜直播｜吳真真牧師｜告別浮淺,門徒進深:在痛苦中找到珍寶母親｜路得記1\_19-22} & 2026-05-10 & \href{https://youtube.com/watch?v=wdtTzXf4txM}{\texttt{wdtTzXf4txM}} \\
約翰福音 13:1-35 & \hyperref[sec:9iKJWzoiFEk]{2026年5月17日崇拜直播｜陳冠成牧師｜告別浮淺,門徒進深:靈命成熟,以愛為尺｜約翰福音13\_1-35} & 2026-05-18 & \href{https://youtube.com/watch?v=9iKJWzoiFEk}{\texttt{9iKJWzoiFEk}} \\
約翰福音 19:1-22 & \hyperref[sec:0dlSUYhlo5s]{2026年5月24日崇拜直播｜盧國偉先生｜在黑暗中靜待黎明｜約翰福音19\_1-22} & 2026-05-24 & \href{https://youtube.com/watch?v=0dlSUYhlo5s}{\texttt{0dlSUYhlo5s}} \\
創世記  & \hyperref[sec:dKG666yV4KA]{2026年5月31日崇拜直播｜楊珊珊傳道｜告別浮淺,門徒進深:衝破過去的藩籬｜創世記 四十一 46-52,四十五 1-8} & 2026-05-31 & \href{https://youtube.com/watch?v=dKG666yV4KA}{\texttt{dKG666yV4KA}} \\
腓立比書 4:1-14 & \hyperref[sec:NSYx8y4kLPg]{2026年6月7日主餐崇拜直播｜陳明泉牧師｜告別浮淺,門徒進深:喜樂的負傷者｜腓立比書4\_1-14} & 2026-06-07 & \href{https://youtube.com/watch?v=NSYx8y4kLPg}{\texttt{NSYx8y4kLPg}} \\
約翰福音 12:9-19 & \hyperref[sec:wKhk31WV7Os]{2026年6月14日崇拜直播｜梁家麟牧師｜尊祂為王｜約翰福音12\_9-19} & 2026-06-15 & \href{https://youtube.com/watch?v=wKhk31WV7Os}{\texttt{wKhk31WV7Os}} \\
詩篇 50:10-17 & \hyperref[sec:c5rLyKd8AcQ]{2026年6月21日父親節主日崇拜直播｜鄭國邦傳道｜告別浮淺,門徒進深:如釋重「父」｜詩篇50\_10-17} & 2026-06-22 & \href{https://youtube.com/watch?v=c5rLyKd8AcQ}{\texttt{c5rLyKd8AcQ}} \\
列王記上 19:1-18 & \hyperref[sec:i6Mr5bzJFRQ]{2026年6月28日崇拜直播｜鍾家威傳道｜告別浮淺,門徒進深:在低谷中重新站起來｜列王記上19\_1-18} & 2026-06-28 & \href{https://youtube.com/watch?v=i6Mr5bzJFRQ}{\texttt{i6Mr5bzJFRQ}} \\
\end{xltabular}
}
\newpage



\section{約書亞記 5:1-12}
\label{sec:oVGpn0PwX_E}
\textbf{2020年2月9日 崇拜直播｜約書亞記5\_1-12}
\newline
\newline
連結: \href{https://youtube.com/watch?v=oVGpn0PwX_E}{\texttt{https://youtube.com/watch?v=oVGpn0PwX\_E}} ~~~~ 語音日期: 2020-02-18
\newline
\newline
\hyperref[sec:code]{\small{< < < PREV SERMON < < <}}
~
\hyperref[sec:index_chronic]{\small{[返順時目]}}
~
\hyperref[sec:w8cFfcUefjQ]{\small{> > > NEXT SERMON > > >}}
\newline
\newline
約書亞記 5:1-12
\newline
\begin{longtable}{cl}
\hline
\hline
章節 & 經文 (和合本修訂版)\\
\hline
5:1 & \begin{tabularx}{0.7\textwidth}{X} 約旦河西亞摩利人的眾王和靠海迦南人的眾王，聽見耶和華在以色列人前面使約旦河的水乾了，直到他們過了河，眾王因以色列人的緣故都膽戰心驚，勇氣全失。 \end{tabularx} \\ \\ \relax
5:2 & \begin{tabularx}{0.7\textwidth}{X} 那時，耶和華對約書亞說：「你要造火石刀，第二次為以色列人行割禮。」 \end{tabularx} \\ \\ \relax
5:3 & \begin{tabularx}{0.7\textwidth}{X} 約書亞就造了火石刀，在哈爾拉勒山為以色列人行割禮。 \end{tabularx} \\ \\ \relax
5:4 & \begin{tabularx}{0.7\textwidth}{X} 約書亞行割禮的原因是這樣：從埃及出來的眾百姓，所有能打仗的男丁，出了埃及以後，都死在曠野的路上。 \end{tabularx} \\ \\ \relax
5:5 & \begin{tabularx}{0.7\textwidth}{X} 這些從埃及出來的眾百姓都受過割禮；但是那些出埃及以後，在曠野的路上所生的眾百姓卻沒有受過割禮。 \end{tabularx} \\ \\ \relax
5:6 & \begin{tabularx}{0.7\textwidth}{X} 以色列人在曠野走了四十年，直到那從埃及出來，全國能打仗的人都消滅了，因為他們沒有聽從耶和華的話。耶和華曾向他們起誓，必不容許他們看見耶和華向他們列祖起誓要給我們的地，就是流奶與蜜之地。 \end{tabularx} \\ \\ \relax
5:7 & \begin{tabularx}{0.7\textwidth}{X} 他們的子孫，就是耶和華興起接續他們的，都沒有受過割禮；因為在路上他們沒有受割禮，約書亞就為他們行割禮。 \end{tabularx} \\ \\ \relax
5:8 & \begin{tabularx}{0.7\textwidth}{X} 全國的人都受了割禮，留在營中自己的地方，直到痊癒。 \end{tabularx} \\ \\ \relax
5:9 & \begin{tabularx}{0.7\textwidth}{X} 耶和華對約書亞說：「我今日將埃及的羞辱從你們身上除掉了。」因此，那地方名叫吉甲，直到今日。 \end{tabularx} \\ \\ \relax
5:10 & \begin{tabularx}{0.7\textwidth}{X} 以色列人在吉甲安營。正月十四日晚上，他們在耶利哥的平原守逾越節。 \end{tabularx} \\ \\ \relax
5:11 & \begin{tabularx}{0.7\textwidth}{X} 逾越節的第二日，他們吃了當地的出產，就在那一天，吃了無酵餅和烘過的穀物。 \end{tabularx} \\ \\ \relax
5:12 & \begin{tabularx}{0.7\textwidth}{X} 他們吃了當地出產的第二日，嗎哪就停止了。以色列人不再有嗎哪了。那一年，他們就吃迦南地的出產。 \end{tabularx} \\ \\
[1ex]
\hline
\hline
\end{longtable}
$^{1}$接著下來是經份.
今天的經文是在.
舊約聖經的267頁.
約書20記五章一到十二節.
舊約聖經267頁.
翻開的請聽我讀.
第五章第一節.
約旦河西阿摩利人的諸王.
和靠海迦南人的諸王.
聽見耶和華在以色列人.
前面使約旦河的水乾了.
等到我們過去.
他們的心因以色列人的緣故.
就消化了.
不再有膽氣.
那時.
耶和華吩咐約書亞說.
你製造火石刀.
第二次給以色列人行國禮.
約書亞就製造了火石刀.
再除皮生那裡給以色列人行國禮.
約書亞行了國禮的緣故.
是因為從埃及出來的眾民.
就是一切能打仗的藍丁.
出了埃及以後.
都死在曠野的路上.
第五節.
因為出來的眾民都受過國禮.
唯獨出埃及以後.
在曠野的路上所生的眾民.
都沒有受過國禮.
以色列人在曠野走了四十年.
等到國民.
就是出埃及的兵丁.
都消滅了.
因為他們沒有聽從耶和華的話.
耶和華曾向他們起誓.
必不容他們看見耶和華向他們列祖起誓.
應許賜給他們的地.
就是流奶與蜜之地.

$^{41}$他們的子孫.
就是耶和華所興起來接續他們的.
都沒有受過國禮.
因為在路上沒有給他們行國禮.
約書亞這才給他們行了.
第八節.
國民都受完了國禮.
就住在營中自己的地方.
等到全裂了.
耶和華對約書亞說.
我今日將埃及的羞辱從你們身上滾去了.
因此那地方名叫吉甲.
就是滾的意思.
直到今日.
以色列人在吉甲安營.
正月十四日晚上.
在耶利哥的平原守逾越節.
逾越節的次日.
他們就吃了那地的出產.
正當那日吃無烤餅和烘的菊.
他們吃了那地的出產.
第二日嗎哪就止住了.
以色列人也不再有嗎哪了.
那一年他們卻吃迦南地的出產.
弟兄姊妹平安.
不論大家在家裡面.
或在教會當中.
今天我們能夠一齊來到敬拜.
一切都是上帝的恩典.
在這段日子.
疫情有很多的改變.
所以教會的行政安排裡面.
很快地要做一些決定.
但是我們很希望盡力做到.
就是希望每一個弟兄姊妹回來.
那個風險可以減到最低.
同時間如果有一些可能在家裡面.
或者生活當中有點不方便的.
都能夠暫時在這段時間.
在家裡面一齊敬拜上帝.

$^{81}$我們很盼望這段日子.
可以快快地過去.
再一次弟兄姊妹一齊聚集在教會當中.
同心一齊敬拜上帝.
而且可以放下我們的口罩.
可以面露多一些的笑容.
今天和大家繼續.
來看邁向新世代的這個系列.
繼續我們去研讀約書亞記.
在約書亞記裡面上一次講到.
整個以色列文約書亞帶領他們.
徒步經過了約旦河.
在這件事經過之後.
又未去到目的地.
第五章第一至第十二節.
就是講到去到迦南目的地.
這個應許之地之前.
和過了約旦河.
這個當中裡面有一個很重要.
界定他們這班以色列人身份的舉動.
就是行割禮和守逾越節.
在這個舉動裡面.
怎樣上帝透過這些重要身份的舉動.
來確立這班人.
讓他們能夠帶著信心來前行.
有時候我們都是一樣的.
很多弟兄姊妹我們人生裡面.
很想有一些目標.
有一些理想.
或者我們人生裡面.
有一些方向想去前行.
但是如果我們不知道自己的位置.
不知道自己的身份.
或者不知道自己的照明.
可能我們就橫衝直撞.
我們覺得想去就去到哪裡.
今日這段經文.
似乎上帝要透過約書亞記裡面提醒我們.
亦都提醒當代的人.
這班在曠野漂流了的第二代.

$^{121}$他們將要進到一個應許之地.
他們的生命是怎樣.
上帝要怎樣繼續使用他們.
希望今日我們從這段經文裡面.
我們找到一些的亮光.
亦都幫助我們繼續生命裡面一齊前行.
好 我們一起看經文裡面所講.
在第五章第一節裡面.
就特別來講到.
當這班以色列民過了約旦河之後.
在約旦河西的阿摩利人.
還有朱王.
還有迦南的朱王聽見.
以色列人這樣徒步的.
經過了約旦河.
約旦河的河水乾了這樣過去.
這班人害怕到不得了.
他們心裡面失去了那種勇氣.
和合本逆膽氣.
其實比較好一點的翻譯是逆勇氣.
我們看看我的powerpoint.
他們失去了這份勇氣的時候.
約書亞就對著他們講.
他們之前未受過割禮的.
這班上一代的那些人.
現在他們要透過這個割禮的舉動.
來除去他們的那個羞辱.
在經文裡面特別記載.
那些從埃及出來的百姓.
舊時第一代在曠野裡面是受過割禮的.
在第二代在曠野當中.
那些大概四十歲或者打猴的那班新世代.
他們全部都未受過割禮.
所以上帝說在進入這個迦南美地之前.
他們必須要接受這個割禮.
以至到他們繼續來承受這份的應許.
我們看看為什麼要做這個舉動.
其實在.
我們看下一章.
其實在創世記第十七章.

$^{161}$第一至第十四節.
其實這個割禮是來自上帝對一班以色列人的一個應許.
這個應許是一個某程度是一個立約的記號.
這個約本身是怎樣呢.
創世記裡面講到上帝和亞伯拉罕來到立約.
本來一個無兒無女的一個老人家.
上帝就將福氣臨到在亞伯拉罕身上.
上帝和這個人來到立約.
告訴他將來有很多很多的子孫.
還有他有成為一個屬於他們自己幫國的一個地方.
他更加要成為多國的爸爸.
多國之父.
所以他們亞伯拉罕和他的後裔.
就要承受上帝傾下的福氣.
成為上帝賜福的一個源頭.
在這個立約裡面上帝特別強調一樣東西.
以色列民他的後代.
特別是亞伯拉罕的後代.
世世代代都要做一個立約的記號.
就是要做這個割禮.
割禮是什麼意思呢.
可能很多弟兄姊妹都知道.
不過不厭其煩再解釋一下.
其實就是在男性的生殖器裡面.
將包皮割了.
做這些禮儀其實是初生孩子.
出世第八天就做這件事.
嬰兒沒有那麼痛的.
大家可以想像一下.
一個嬰孩出世八天就會做這個舉動.
不過歷史一直亞伯拉罕的後裔.
一直言遠流長.
他們未必繼續來守這個割禮.
更加重要就是去到埃及地.
甚至乎去到曠野當中.
他們這班人已經遺忘了立約的身份.
可能在曠野裡面.
第一代都有做過這個舉動.
這個割禮的舉動.
不過在曠野那四十年之後.

$^{201}$那些上一代和那些父母.
都沒有幫他們自己的兒女.
來做這個立約的記號.
所以在上帝的眼中裡面.
或者上帝所講的話裡面.
他強調整個在曠野漂流的新一代.
四十歲打下的.
全部都沒有受過割禮.
上帝就要強調他們必須要再一次來行這個割禮.
值得注意是這個割禮.
在什麼時候來做呢?.
這個是將要進到約旦河.
將要進到一個目的地.
這個舉動是某程度上是一個打仗.
在打仗之前做一個外科手術.
更加重要的是那班人不是小孩子.
是一班成年壯丁.
所有壯丁那些打得仗的男丁.
全部要做一個這樣的舉動.
其實這個舉動是一個很大的冒險.
如果那些敵軍或者迦南地的王之獨.
在以色列做這樣的事.
馬上乘虛而入.
當時可能甚至乎他們身體未復原.
他們就可以令到他們全軍覆沒.
不過這班以色列民不是傻的.
約書亞不是傻的.
大概他們都知道將要去面對一個這樣的情勢.
他們都一定要去進攻.
不過上帝的御令比一切更加重要.
所以這一班人就很聽話.
遵從上帝的吩咐.
在上帝這樣吩咐之下.
約書亞整個所有的男丁全國.
一起去接受這個國禮.
另一個注意的地方就是.
有一個字眼在五章第七節那裡.
在經文裡面特別的來到講到.
他們的子孫就是耶和華所興起來接觸他們的.
都沒有受過國禮.

$^{241}$在這個有個特別的字叫做興起這個字.
興起這個字其實.
剛才講到創世記.
亞伯拉罕和上帝一起去立約.
立約的納字的動詞.
後來還有約書亞記過約旦河納錫維記的納字.
和現在這個興起這個字是同一個字.
當你讀約書亞記你要看回這段經文的時候.
你就會知道.
上帝要透過這件事.
透過這個興起立約的這個舉動.
來提醒他們新的一代.
這班人不是純粹一班軍事家.
或者憑著自己的野心.
他要進入去迦南地.
這一切的舉動.
是上帝自己作主動.
是上帝將這個福氣.
答應了亞伯拉罕.
上帝就將這個承諾兌現.
而且現在就在這一代裡面.
正式的來到去兌現.
整個過程是上帝作為一個主動.
是上帝自己一手要去促成的一件事.
過去裡面.
為什麼這班人領受不到福氣呢.
大概他們可能生命裡面發出怨言.
或者他們活逆上帝.
又或者他們的罪.
以致到他們沒有辦法.
來去承受上帝給他們的這個福氣.
如今不是說真的.
這一代就是比他們更加敬虔.
只不過是上帝將這個福氣.
轉移到他們的下一代裡面.
但同樣上帝都提醒他們.
這一班人他們不是有什麼過人之處.
只不過是上帝藉著他自己的恩典.
他的承諾如今就要兌現出來.
另外一個字眼就是那個地方.

$^{281}$吉甲的字本身就叫做滾出去.
即是今天廣東話說.
爛開的意思.
滾出去.
滾出去是在說什麼滾出去呢.
不是在說這班以色列人滾出去.
他是在說這班整個以色列.
整個群體的集體的那種羞辱.
由昔日上帝答應了這一班以色列人.
很多又很多代.
好幾百年的日子.
但是他們這個承諾.
他們因著他們自己的軟弱.
沒有辦法來實踐開它.
他們的生命.
因為沒有辦法承受上帝的福氣.
外面的人就會嘲笑他.
包括在埃及地的那些人民.
嘲笑他.
說上帝應許.
但是現在還是繼續成為一個圍爐之人.
甚至去到曠野裡面.
成為列國嘲笑的對象.
但上帝去到現在.
他要提醒他們.
因為上帝一手承諾這件事他要完成了.
所以他要將整個以色列民族的那種羞辱.
驅除去.
藉著這個立約的記號.
藉著這個國禮.
重新奠定提振回整個民族.
和上帝立約的那種關係.
告訴他們知道.
今天做了這個國禮之後.
他們認定自己的身份.
是一個將要打仗的一個很重要的關頭裡面.
做一個這麼危險的舉動.
就是要認定.
這班人是上帝所閱立.
是上帝賜福的一群.

$^{321}$上帝要將整個以色列民族的羞恥.
趕走它.
讓他們回復一個榮耀尊貴的身份.
在這裡整個關於國禮的記載.
是一個很不容易的過程.
這班以色列民.
他們專心遵行上帝的吩咐.
他們按著上帝的吩咐而行.
藉著一個很大膽.
很冒險.
也是一個很標誌性的舉動.
來確認他們的身份.
這班人是屬於上帝的.
在經文裡面再繼續說.
五章十至十二節.
在經文裡面.
以色列人在吉甲安營正月十四的晚上.
他們在耶利哥的平原守逾越節.
逾越節的第二日.
他們吃了當地的出產.
就在那一天.
吃了無餃餅和烘過的穀物.
他們吃了當地出產的第二日.
嗎哪就停止了.
以色列人不再有嗎哪了.
那一年他們就吃迦南地的出產.
在這段經文裡面.
短短的幾節.
三次出現一個重複的動詞.
就是吃.
吃些什麼呢.
吃些什麼原來是很重要的.
在這幾節裡面.
過往裡面.
可能在曠野裡面.
他們都守過逾越節.
又或者在出埃及的時候.
他們守過第一次的逾越節.
他們吃的東西.
是埃及的東西.

$^{361}$或者曠野裡面的.
他們擁有的一些物產.
更加重要的是整個以色列民.
他們四十年來的那些主糧.
吃飽肚子的東西.
是天降的嗎哪.
他們吃這些東西.
不過呢.
去到將要進入迦南地這個地方.
這次他們所吃的東西就很不同.
他們這次所吃的東西.
就是吃迦南地當地的出產.
再三的來到去講.
他們要吃迦南地的出產.
而當他們吃的第二天開始.
天降的嗎哪就停止了.
在這裡面.
在聖經裡面很特別刻意地記載.
他們所吃的東西.
吃了的東西不同.
他們守的逾越節的食品.
吃的種類都是一樣的.
不過來源是不同了.
即日由天上降下來.
又或者是埃及地的.
所吃的這些的土產.
如今他們就在迦南地裡面.
將要進入這個應許之地.
他們就吃那裡的東西.
當他們吃了這些東西的時候.
然後天上的那個嗎哪.
那個水喉就關了.
那樣東西就不再供應.
嗎哪的終止.
代表著一個新的時代開始.
這班以色列民已經某程度.
像破釜沉舟一樣.
他們就沒有再有退路.
他們不會再回到約旦河那邊.
過河就會有嗎哪降下來.

$^{401}$大致上帝的心意.
都不是想他們這樣繼續.
來過他們未來的生活.
有一些的生活.
或者在他們未來所走的路.
是不能回頭的了.
上帝一方面.
要這班以色列人.
確認自己的身份.
藉著這個逾越節的禮儀.
不單止是說.
他們的身份被更加肯定.
這個應許又要兌現.
他更加要推動這班以色列人.
整個的歷史.
整個面前要走的路.
就沒有退路.
他們就要繼續這樣走.
他們必須要義無反顧.
帶著信心和帶著盼望.
來信靠上帝.
來讓他們繼續走.
面前的每一步.
在過去四十年裡面.
這個嗎哪.
無可否認.
是幫助整個以色列民.
令到他們衣服沒有穿破.
腳沒有腫.
上帝的公認.
四十年是一個很偉大的神蹟.
但是如今嗎哪停止.
他要令到約書亞所帶領的.
這班以色列民.
不要再看舊事了.
不要再緬懷舊時.
上帝給他們.
不用做天掉下來的.
心靈裡面的滿足.
未必是最富裕的.

$^{441}$但是好像不用憂.
上帝自動降下來的.
那種公認的生活.
如今他的生活方式.
他要改變.
不過這個應許.
是比起之前嗎哪更加好的.
比起昔日在曠野的日子.
會更加好的.
上帝為他們預備一個新的.
美好的將來.
當我們看完這兩個很重要的標誌.
割禮和守這個逾越節.
大概他們過去四十年.
都不是特別去做的.
割禮更加完全沒有做過.
守逾越節可能零星這樣去做.
有這些教導.
但是他們不是提起勁這樣去實踐.
上帝藉著這兩個很重要的禮儀.
去喚醒整個民族的身份.
知道他們自己是哪一個.
他們就繼續前行他們面前的路.
在這裡我們稍為有一些思想.
今天我們每一個在座裡面.
或者在家裡面聽到的一班弟兄姊妹.
我們很宣認耶穌基督是我們的救主.
我們都會認定耶穌是我們生命的主.
他要帶領我們.
我們成為一個基督徒.
不過有時候我們基督徒這個身份.
是否真的成為我們生命裡面決定了.
是有一個決定性來左右我們行事為人.
哪些做哪些不做一個很重要的準則呢.
當我們生命裡面拿捏不準一個很重要的身份的時候.
這個世界想說的東西.
覺得好的東西我們就會去做.
有些人所說的話我們就會跟隨.
或者我們按著我們自己的喜好來做每一個的抉擇.
有時候我們都嗤之以鼻.

$^{481}$我們都不想聽見有些人在批評.
基督徒都做這些事.
基督徒都是這樣的.
我們聽見一來就覺得很羞愧.
作為基督徒都有一些感覺上自己覺得好像很尷尬.
但同時間我們就很憤恨.
有沒有搞錯呀.
即是累街坊的那種感覺.
你做了些壞事人家還要自己去掃.
不過我們自己想.
我們今天我們自己行事為人.
跟我們蒙召的因是不是一定來到去相稱呢.
今日這段經文提醒我們有幾方面.
第一方面在聖經裡面.
上帝很強調他的應許要兌現.
而這個應許是要為了他揀選的子民來到去實踐出來.
當我們每一個被召的一班弟兄姊妹.
我們生命裡面我們都是帶著上帝這一份的應許.
在舊約裡面就是亞伯拉罕的應許.
去到今天我們每一個就是屬於上帝.
每一個都是屬於耶穌基督.
在基督裡的一班基督徒.
我們這班人同受上帝寶貴的應許.
上帝將這個福音帶給我們.
祂期望我們要認定這個很尊貴很美好的身份.
上帝更加要告訴我們.
這個身份不是我們眾多身份裡面.
就好像我們信用卡那裡.
或者不同的會籍那些卡.
其中一張卡.
而是一個最重要的身份.
這個是界定了我們整個人生的.
未來人生路向的一個很重要的身份.
早前有一個機會看一本書.
是關於一些北美的一班中國華裔的基督徒.
他們去到美國這個地方.
他們經歷他們自己的信仰生活.
但是去到中國人去到一個北美的地方.
他要經歷很多文化適應和身份認同的問題.
去到那裡他們會問自己的問題.

$^{521}$究竟他是中國人?.
他是香港人?.
他是美國人?.
還是什麼人呢?.
其實在外國教會裡面.
這本書特別刻畫了內心裡面的掙扎.
後來裡面寫了一個很重要的見證.
他講到有一個第二代的美籍華人.
他的語言能力很強.
然後他講他自己的見證.
他說他的身份是很多元的.
在北美當中.
如果在美國人當中.
他就會視自己為美國人.
在外國西方世界裡面.
他會視自己為一個西方人.
在中國人的群體裡面.
他就會視自己為一個中國人.
跟廣東話和中文的人交往.
然後他說他們的思維全部都能夠答得通.
我們都覺得很厲害.
學貫中西或者每一樣東西都能夠溝通.
他更加講.
不過有時候在外國人裡面有些漏雜.
他們就會用中國人的眼光來批判.
例如對於他們對親情的那些看法.
即是外國可能那樣東西.
那種效度比較淡薄.
但他用中國人的看法.
他就會覺得他尊重別人的文化.
但他心裡面我都是流著中國人的血.
我是重視這些的效度.
同時間他去到中國人裡面.
他會覺得中國人只是重情不重禮.
不講道理的.
有些東西怕肩膀搞掂.
不需要尋真的那種精神.
他心裡面就會用外國人的心態來看.
如果所有東西都是講情不講理.
很容易會有信心.

$^{561}$所以他在這裡面暗地裡就做一些批判.
所以既是認同.
但同時間又有批判.
這個見證最後他寫.
為甚麼他有這樣的能力呢.
就是因為他有一個更加超越了文化的身份.
來奠定他自己最重要的生命.
就是一個基督徒的身份.
基督徒的身份跨越了他眼前不同文化.
不同地域.
不同習性.
是因為這個基督徒的身份.
令到他能夠有一種分辨的能力.
這個是一個決定性的身份.
我想說甚麼呢?弟兄姊妹.
今日我們可能自己都有不同的背景.
我們有不同的專業.
我們有不同的能力.
但我們不可以忘記.
基督徒的身份不是我們魂魂裡面.
其中一個身份裡面平等的一個身份.
我很想我們一齊來努力.
我們活出一個.
或者我們以基督徒.
以成為上帝指明的生命.
這種身份.
來凌駕於我們眼前不同.
我們自己覺得自豪的.
或者我們過去所活著的.
那些不同的身份.
唯有一個更加重要.
有一個超然的基督徒生命.
這個身份才可以左右我們人生.
作不同的決定.
特別在不同的關口裡面.
我們才有.
能夠有明辨是非的能力.
以至到我們可以繼續的前行.
在這段經文裡面.
亦都提醒我們第二點.

$^{601}$我們身份界定我們自己.
這個過程不是容易的.
甚至乎有時我們活出我們的身份.
或者我們強調自己的身份.
我們有時要付上代價的.
有時我們會做一些事被人吃虧的.
我們可能不單止說我們做些甚麼.
我們更加說到有時我們不會做些甚麼.
有所為與有所不為.
在不同的時機裡面.
我們更加要看清自己.
在這個世代裡面.
我們要成為一個怎樣的人.
在這個過程裡面.
可能因著這些身份.
因著這些背景.
令到我們可能與其他人.
與世界的待遇有些不同.
不過我們甘心.
因為我們成為一個基督徒.
我很感動的.
我聽到教會裡面有些弟兄姊妹.
這段時間大家都是.
你封50元利是給我都不是很開心.
你封個口罩給我當利是.
大家很開心的.
我很感動的.
我知道有些弟兄姊妹.
特意下街為了社區裡面.
有些很基層又不能上街的弟兄姊妹.
他們下街幫他們去封口罩.
排隊排很長時間.
有些下去買不到米.
我們有些童工自己玩著袋米.
就怕面對面.
就放在袋米門口按個鐘.
然後就馬上走.
但他們做一些這樣的舉動.
為甚麼他們會這樣做呢.
除了他們心裡面有一種.

$^{641}$來自上帝的愛之外.
他更加知道在這個世代裡面.
基督徒要發光.
在一個很困難的處境裡面.
我們就要努力地為身邊的人.
做一點小小的祝福.
身份從來不是純粹.
一張表格或者一個聚會去界定我們.
身份往往建基這些.
剛才所講之外.
還要加上我們自己的行動.
和我們怎樣去做決定.
在特別困難的日子裡面.
這個身份更加左右我們怎樣行事為人.
在聖經裡面.
他鼓勵這班以色列民將要打仗當中.
他們要行割禮.
他們就要順從上帝的吩咐.
繼續來對勇往直前.
上帝叫我們得到這個身份.
不是令到我們感覺上是一個羞愧的一種生命.
在彼得前書第二章第六節裡面.
他講到看我把我所揀選的.
所寶貴的房角石.
安放在石安.
信靠他的人必不至於羞愧.
在新約裡面在舊約裡面同樣強調一樣東西.
我們每一個弟兄姊妹我們的身份.
雖然我們面對著這個世界.
或者我們有時候會為著我們所擁有的這個基督徒身份.
要付上一些代價.
但這些的際遇或者這些的代價.
不是叫我們垂頭喪氣.
低頭低頭地過我們的生活.
彼得他很想告訴每一個被揀選的弟兄姊妹.
我們的身份有時候我們可能會吃虧一點.
甚至乎我們可能會面對一些危難困難.
但這個身份卻不是叫我們是一個羞愧的身份.
反而是一個榮耀的身份.
上帝叫我們活出一種榮耀的生命.

$^{681}$可能未來的日子.
我們香港或者面對著一些經濟的困難.
或者面對著社會裡面有些疫情.
可能我們都有機會會病.
坊間很容易就會用際遇來定義成為一個人.
例如有時候會用失業的叫失業人士.
失業成為我的身份.
或者我是一個病人.
病成為我終身的一個身份.
不過上帝說.
在上帝的國度裡面.
從來不是用這些的際遇來定義一個人.
他的身份他的價值是怎樣的.
上帝是用他的應許.
是用人的信心來定義一個人的生命.
可能我的經濟變得沒那麼富裕.
可能我生命裡面都面對著一些疾病困難.
我面對這些的際遇.
不等於我永遠是一個很羞愧很痛苦的際遇.
我面對這些的經歷.
但我有上帝在我生命當中.
上帝的應許淋到在我生命裡面.
叫我們可以昂然闊手挺起胸膛.
繼續來承受他的應許.
弟兄姊妹.
今天我們不知道我們今天處於什麼的位置.
無論你是相對比較順境.
很順境.
面對逆境或者和平時差不多.
今天界定我們的身份.
不是這些的際遇.
而是上帝的應許.
今天我們一齊的來讀這段經文.
上帝這些的吩咐就是要讓新一代的以色列民.
找回自己的生命.
找回自己的身份.
以致到他們才可以繼續的來前行.
同樣我們今天.
我們都要找回我們自己基督徒這個寶貴榮耀的身份.
我們一齊的來努力.

$^{721}$可能在當中裡面會有些的吃虧.
但我們不需要懼怕.
也不需要憂慮.
我們倚靠上帝.
我們一齊的來過一個上帝喜悅的生活.
願主繼續來對私恩.
與我們每一個同在.
我們同心祈禱.
阿天父願你幫助我們.
在你話語當中.
讓我們生命有亮光有激勵.
讓我們繼續前行的日子.
讓我們不需要擔憂.
我們邁向一個新的時代的時候.
給我們有足夠的信心.
在疫情當中.
幫助我們活出信望愛的生命.
求你繼續的來對恩待我們眾人.
賜福給眾靈姐妹.
我們獻上禱告.
奉基督耶穌得勝的聖名祈求.
Amen.
\newpage



\section{約書亞記 6:1-21}
\label{sec:w8cFfcUefjQ}
\textbf{2020年2月23日 崇拜直播｜約書亞記6\_1-21}
\newline
\newline
連結: \href{https://youtube.com/watch?v=w8cFfcUefjQ}{\texttt{https://youtube.com/watch?v=w8cFfcUefjQ}} ~~~~ 語音日期: 2020-02-23
\newline
\newline
\hyperref[sec:oVGpn0PwX_E]{\small{< < < PREV SERMON < < <}}
~
\hyperref[sec:index_chronic]{\small{[返順時目]}}
~
\hyperref[sec:ruBv0LtC8pg]{\small{> > > NEXT SERMON > > >}}
\newline
\newline
約書亞記 6:1-21
\newline
\begin{longtable}{cl}
\hline
\hline
章節 & 經文 (和合本修訂版)\\
\hline
6:1 & \begin{tabularx}{0.7\textwidth}{X} 耶利哥的城門因以色列人的緣故，關得嚴緊，無人出入。 \end{tabularx} \\ \\ \relax
6:2 & \begin{tabularx}{0.7\textwidth}{X} 耶和華對約書亞說：「看，我已經把耶利哥城和耶利哥王，以及大能的勇士，都交在你手中。 \end{tabularx} \\ \\ \relax
6:3 & \begin{tabularx}{0.7\textwidth}{X} 你們要圍繞這城，所有的士兵繞城一次，六日你都要這樣做。 \end{tabularx} \\ \\ \relax
6:4 & \begin{tabularx}{0.7\textwidth}{X} 七個祭司要拿七個羊角走在約櫃前。到了第七日，你們要圍繞這城七次，祭司也要吹角。 \end{tabularx} \\ \\ \relax
6:5 & \begin{tabularx}{0.7\textwidth}{X} 羊角聲拖長的時候，你們一聽見角聲，眾百姓要大聲呼喊，城牆就必倒塌，各人要往前直上。」 \end{tabularx} \\ \\ \relax
6:6 & \begin{tabularx}{0.7\textwidth}{X} 嫩的兒子約書亞召了祭司來，對他們說：「你們抬起約櫃來，要有七個祭司拿七個羊角在耶和華的約櫃前。」 \end{tabularx} \\ \\ \relax
6:7 & \begin{tabularx}{0.7\textwidth}{X} 他又對百姓說：「你們向前去圍繞那城，帶兵器的要在耶和華的約櫃前過去。」 \end{tabularx} \\ \\ \relax
6:8 & \begin{tabularx}{0.7\textwidth}{X} 按照約書亞對百姓所說的，七個祭司拿了七個羊角在耶和華面前過去，他們吹著角，耶和華的約櫃在他們後面跟著。 \end{tabularx} \\ \\ \relax
6:9 & \begin{tabularx}{0.7\textwidth}{X} 帶兵器的走在吹角的祭司前面，後隊跟著約櫃走，號角繼續在吹。 \end{tabularx} \\ \\ \relax
6:10 & \begin{tabularx}{0.7\textwidth}{X} 約書亞吩咐百姓說：「你們不可呼喊，不可讓人聽見你們的聲音，連一句話也不可出你們的口，直到我對你們說『呼喊』的那日，你們才呼喊。」 \end{tabularx} \\ \\ \relax
6:11 & \begin{tabularx}{0.7\textwidth}{X} 這樣，約書亞使耶和華的約櫃圍繞那城，把城繞了一次。然後，眾人回到營裡，就在營裡住宿。 \end{tabularx} \\ \\ \relax
6:12 & \begin{tabularx}{0.7\textwidth}{X} 約書亞清早起來，祭司又抬起耶和華的約櫃。 \end{tabularx} \\ \\ \relax
6:13 & \begin{tabularx}{0.7\textwidth}{X} 七個祭司拿七個羊角，走在耶和華的約櫃前，他們吹著角；帶兵器的走在他們前面，後隊跟著耶和華的約櫃走，號角繼續在吹。 \end{tabularx} \\ \\ \relax
6:14 & \begin{tabularx}{0.7\textwidth}{X} 第二日，他們再把城圍繞一次，就回營裡去。六日都是這樣做。 \end{tabularx} \\ \\ \relax
6:15 & \begin{tabularx}{0.7\textwidth}{X} 第七日清早黎明時，他們起來，以同樣的方式圍繞城七次；惟獨這一日他們圍繞城七次。 \end{tabularx} \\ \\ \relax
6:16 & \begin{tabularx}{0.7\textwidth}{X} 到了第七次，祭司吹角的時候，約書亞對百姓說：「呼喊吧，因為耶和華已經把城交給你們了！ \end{tabularx} \\ \\ \relax
6:17 & \begin{tabularx}{0.7\textwidth}{X} 這城和其中所有的都要永獻給耶和華作當毀滅的，只有妓女喇合與她家中所有的可以存活，因為她隱藏了我們所派的使者。 \end{tabularx} \\ \\ \relax
6:18 & \begin{tabularx}{0.7\textwidth}{X} 但你們務必謹慎，不可取那當滅的物，免得你們受詛咒，取了那當滅的物，使以色列全營成為詛咒而遭受災禍。 \end{tabularx} \\ \\ \relax
6:19 & \begin{tabularx}{0.7\textwidth}{X} 只有金子、銀子和銅鐵的器皿都要歸耶和華為聖，放入耶和華的庫房中。」 \end{tabularx} \\ \\ \relax
6:20 & \begin{tabularx}{0.7\textwidth}{X} 於是百姓呼喊，祭司吹角。百姓一聽見角聲就大聲呼喊，城牆隨著倒塌。百姓上去進城，各人往前直上，把城奪取。 \end{tabularx} \\ \\ \relax
6:21 & \begin{tabularx}{0.7\textwidth}{X} 他們把城中所有的，無論男女老少，牛羊和驢，都用刀殺盡。 \end{tabularx} \\ \\
[1ex]
\hline
\hline
\end{longtable}
$^{1}$我們今天要聽神的話是記載在.
約書亞記第六章一至二十一節.
念到的是記載在舊約聖經第二百六十八頁.
第六章約書亞記第一節.
耶利哥的城門因以色列人.
就關得嚴緊.
無人出入.
耶和華曉遇約書亞說.
看啊.
我已經把耶利哥和耶利哥的王.
並大能的勇士都交在你手中.
你們一切兵兵要圍繞這城.
一日圍繞一次.
六日都要這樣行.
七個祭司要拿七個楊角走在約櫃前.
到第七日.
你們要繞成七次.
祭司也要吹角.
他們吹的角聲延長.
拖長.
你們聽見角聲.
眾百姓要大聲呼喊.
城牆就必塌陷.
各人都要往前直上.
亂的兒子約書亞召了祭司來.
吩咐他們說.
你們抬起約櫃來.
要有七個祭司拿七個楊角.
在耶和華的約櫃前.
有對百姓說.
你們前去繞城.
大兵氣的要走在耶和華的約櫃前.
第八節.
約書亞對百姓說了話.
說完了話七個祭司拿七個楊角.
走在耶和華面前吹角.
耶和華的約櫃在他們後面跟隨.
大兵氣的走在吹角的祭司前面.
後隊隨在約櫃行.
祭司一面走一面吹.

$^{41}$約書亞吩咐百姓說.
你們不可呼喊.
不可出聲.
言一句話也不可出你們的口.
等到我吩咐你們呼喊的日子.
那時才可以呼喊.
這樣他使耶和華的約櫃擾城.
把城擾了一次.
眾人回到營裡.
就在營裡住宿.
第十二節.
約書亞清早起來.
祭司又抬起耶和華的約櫃.
七個祭司拿七個楊角.
在耶和華的約櫃前.
時常行走吹角.
大兵氣的.
在他們前面走.
後隊隨著耶和華的約櫃行.
祭司一面走一面吹.
第二天.
眾人把城擾了一次.
就回營裡去.
六日都是這樣行.
第七日清早.
黎明的時候.
他們起來.
照樣擾城七次.
唯獨.
這一天把城擾了七次.
到了第七次.
祭司吹角的時候.
耶和華吩咐百姓說.
呼喊吧.
因為耶和華已經把城搞給你們了.
這城和其中所有的.
都要在耶和華面前毀滅.
只有妓女拉合.
與她家中所有的可以存活.
因為她隱藏了我們所打發的使者.

$^{81}$至於你們.
務要謹慎.
不可取那當滅的物.
恐怕你們取了那當滅的物.
就連累以色列的傳迎.
之傳迎.
受就在.
唯有金子銀子和銅鐵的器皿.
都要歸耶和華為聖.
必入耶和華的腹中.
於是百姓呼喊.
祭司也吹角.
百姓聽見各聲變大聲呼喊.
城牆就塌陷.
百姓便上去進城.
各人往前直上.
將城奪取.
又將城中所有的.
匹軀男女老少.
牛羊和奴.
都用刀殺盡.
聖經讀筆.
各位主內平安.
平安實在在這段日子是很需要.
亦不是我們自己可以去祈求.
是上帝來賜給我們.
在開始之前.
我們一起有一個禱告.
天父上帝我們向你禱告.
向你祈求.
除了我們現在很多的地方.
甚至世界很多的角落.
都面對著這個疫症的問題.
除了我們鄰近中國.
我們香港這個地方.
我們都有很多.
很多這樣的情況出現.
還有我懇求你憐憫我們世人.
知道我們真的有很多的軟弱.
有很多的罪.

$^{121}$還有我們懇求你憐憫.
我們去祈求你的時候.
你就賜下這一份.
是沒有人能夠可以奪去的平安.
天父願意你親自來眷顧.
今日我們仍然可以在禮堂來敬拜.
甚至有弟兄姊妹.
他們在家中都和我們一起敬拜.
實在是上帝你給我們的恩典和福氣.
我們真是願意.
專心一意的來敬拜你.
在這段艱難的日子裡.
我們真的只有仰望你.
因為我們的生命是屬於你.
我們的生命是帶著上帝.
來給我們的照明.
我們怎樣可以在這個世代.
更艱難的日子當中.
去活出上帝你給我們.
生命裡面那種的身份.
求你與我們同行.
與我們同在.
亦都願意今日的訊息.
可以成為我們眾人的安慰.
提醒我們.
讓我們怎樣能夠在這段日子.
去見證你.
誰聽我們在你面前的禱告.
奉主耶穌基督的名.
阿們.
首先在這裡來到.
有一個powerpoint預備.
他們很努力.
因為我預備的時候.
可能和今日.
平時播powerpoint.
是有些分別的.
可能他們又要做一番功夫.
另外就是在程序表裡面的大綱.
有些經文是會有少少的調動.

$^{161}$弟兄姊妹.
剛才Jeff和我們讀了.
這一段的經文.
第六章的一至二十一節.
我們就這樣讀.
因為有很多重複.
好像有些不清不楚.
但如果去找一些聖經裡面的動畫的時候.
我發覺看下去.
真是很精彩.
如果大家在教會裡面.
一直成長的時候.
你都會看到.
這個其實是一個很戲劇性的.
一個聖經裡面的記載.
我們都知道耶利哥城是一個.
城牆很厚.
但它的城市不是很大的地方.
這一場的戰爭.
我們來讀.
但其實嚴格來說.
它不是一個戰爭.
其實神確立了.
約書亞來到承接摩西.
他一直去到約旦河的時候.
他就準備帶領百姓來過約旦河.
來去耶和華賜給他這個迦南地.
約書亞的一生.
在這個時刻就不再一樣.
因為他承擔著帶百姓進迦南的使命.
而在他往後的一生.
都是這樣來事奉神.
經文亦都告訴我們.
其實如果從約書亞來看的時候.
其實這一班是新生代的.
一班的以色列人.
他們要重新進入這個應許地.
其實就是要實踐神所交給他的使命.
究竟約書亞的使命在哪裡來呢?.
如果大家很熟悉民數記的時候.

$^{201}$大兄姐妹都會留意.
約書亞和迦勒當時.
和其他十個探子.
走進迦南地的地方.
來視察地形.
究竟裡面的堅城堡壘是怎樣.
他們回來的匯報確實有這樣的情形出現.
但他們再匯報的時候是怎樣呢?.
原來是物產豐富.
但他們有一個問題.
就是那裡的居民很強大.
所以在那一刻他們建議不要進去.
但我們都知道約書亞和迦勒.
他們是力排眾議.
在這樣的情況之下.
縱然知道是上帝的心意.
但他們仍然拒絕進入去.
結果就是這一班的以色列人.
他們在曠野兜兜轉轉四十年.
但在那個時候約書亞.
已經是上帝賜給他.
只有他和迦勒可以進去迦南地.
所以他的身份和使命.
就在那個時候開始形成.
約書亞記來到這個時候.
第五章就是進入迦南地的準備.
我相信陳牧師.
兩位陳牧師在之前的講道.
都有講過這一部分.
過了約旦河的時候.
其實就是重新與神重整這個關係.
去確認那個的身份.
那個的使命.
所以就為新生代的以色列人.
來進行割禮.
來表明他們這種立約的身份.
然後就有御戰.
當一切就緒.
約書亞就會帶領著.
一班新生代的以色列人.

$^{241}$就要踏上這個征途.
所以去到第六章至第十一章.
其實就是記載著這一段的事跡.
而攻打耶利哥城.
我相信這個是最重要的戰役.
因為這個是第一場的戰爭.
究竟耶利哥城的位置是怎樣呢?.
大家如果看到PowerPoint.
你會看到在右邊的地圖.
近處的死海.
在這裡的翻譯.
它就叫做阿拉巴海.
就是最南端那個.
耶利哥城就在對上一點.
其實這個城市.
它大概有三點多公里.
但是那裡的氣候非常之好.
泉水又多.
是適合種植的.
所以風修是非常之好.
另外也有叫做中樹城之稱.
雖然它比其他的城市略為小.
但是因為它是一個交通的要道.
如果要進入去迦南的中部.
或者是那個福地.
然後向南向北.
這個耶利哥城是一個必經之路.
當時迦南地裡面.
有不同的王.
有他們的勢力.
其實是危機四伏的.
如果要攻打耶利哥城.
我想整個以色列人.
在約書亞帶領的時候.
這個是關鍵的時刻.
其實是不容有失.
所以今天我們就透過.
攻打耶利哥城.
這個記載.
讓我們思考.

$^{281}$使命身份的態度.
如何去實踐我們使命人生.
我們的態度.
我們的心態是怎樣呢.
我相信第一點就是要放下.
我們自以為是的執著.
我們總要去信靠主.
去信服他.
經文就是在第五章十三至十五節.
和六章三節那裡.
在六章第一節的時候.
剛才所讀的.
耶利哥城因以色列人關得嚴緊.
無人出入.
我想這個就是約書亞見到的現實.
當然較早之前.
拉合都來通風報信.
給探子聽說.
迦南耶利哥城的氣氛是怎樣.
說到以色列人浩浩蕩蕩.
過了紅海.
殺到約旦河東的時候.
更加將兩個王都殺死了.
所以在這個時候.
耶利哥城的氣氛.
可以說是膽戰心驚.
丙姐妹.
如果我們在困難的裡面.
我們怎樣做呢.
約書亞在這個重要的時刻.
其實要攻打.
其實一定是有難度.
當我們面對這些危機的時候.
面對這些挑戰.
不知道丙姐妹怎樣應付呢.
我們是否會一如既往.
很堅持尋求神的心意.
等候上帝給你的旨意.
還是以自己的經驗.
經歷人生的閱歷來判斷.

$^{321}$去找一個較容易.
傷害少.
損失少的方案來決定去做呢.
約書亞作為一個經驗豐富的軍事領袖.
可以說是一個軍事家.
他必定會跟你一樣.
計算清楚.
但其實約書亞亦是尋求神的心意.
當時他尋求神的心意的時候.
我想大部份都會捉壇來禱告.
所以禱告的時候.
他會不會這樣說呢.
他說上帝.
你給我一個很完美的軍事部署.
我評估之後.
我能夠去攻打耶利哥城.
因為這是第一將.
不容有失.
如果大獲全勝.
上帝我就順服你了.
我不知道約書亞會不會這樣祈禱.
但有一件事我會知道.
其實約書亞在第一章裡面都提到.
約書亞是一個切實敬畏神.
敬拜上帝的人.
所以他這種虔誠的生命.
能夠讓他在以色列人中成為一個敬虔的領袖.
而他的百姓.
他所帶領的百姓.
都是屬於約書亞來帶領.
我們去看剛剛所看到的經文.
第五章十三至十五節的時候.
簡單來說這段經文.
我用三個H來講解這段經文.
第一個H我們會看到約書亞.
他是順服地敬拜.
Humble worship.
從這段經文十三節至十五節.
他看到誰出現了呢?.
他看到原來在他前面有一個人.

$^{361}$這個人當他問完之後.
他說他是耶和華軍隊的元帥.
當他知道的時候.
聖經裡面告訴我們.
他就馬上來俯伏敬拜.
其實約書亞作為當時的一個領袖.
其實當對方這樣講的時候.
他沒有覺得他的位置被譴責.
反而他知道原來是上帝來幫助他.
讓他知道在他一直想的方法裡.
究竟敵人是怎樣.
我自己的能力是怎樣.
發覺原來不是.
他真的有這種順服謙卑來敬拜.
這位前面的元帥.
接著他有一個聖潔的行動.
一個Holy walk.
因為當他知道這是軍隊的元帥的時候.
他就想請求.
想要一個指示.
我真的要怎樣攻入耶利哥城呢?.
但是答案是告訴他.
你現在所站的地方是聖地.
所以你要將鞋脫去.
這個行動更加的展示.
他真的願意降服在這個元帥.
就算是成為他的副手.
約書亞都願意.
當我們看回第六章第三節至第五節的時候.
其實這個就是第三個H.
這個Heavenly war.
是一個神聖的戰爭.
當約書亞他經歷了.
這個順服敬拜.
這個聖潔行動的時候.
他知道前面這位元帥的將軍.
一定會說出那個戰爭的策略是怎樣.
究竟那個戰爭的策略是怎樣呢?.
就在第三節至到第五節.
上帝將攻城的策略.

$^{401}$就交了給約書亞.
現在只是約書亞.
會不會按著這個軍令.
這個戰術的方案來執行.
弟兄姊妹.
假如你接到這個命令.
可能今天我們容易.
用Fax E-mail去到前線的時候.
不過如果看清楚.
這張紙的時候.
在最頂一行寫著的是什麼呢?.
「國漢,我已經把耶利哥.
和耶利哥的王並大能的勇士.
都交在你手中」.
我想問.
在標題.
所謂的heading.
他先打了這一段的時候.
然後才說他的軍事行動.
大家覺得怎樣呢?.
可能這個耶和華軍隊的元帥.
他沒有說出.
他沒有告訴約書亞.
其實我已經察看了.
看完這個形勢.
其實就是要用這個戰術.
第二節.
這個元帥所說的話.
其實會不會暗示了.
如果不執行.
你用什麼的戰略.
用什麼的方式.
都一定會打輸呢?.
大家一路看.
亦都聽的時候.
嚴格來說.
這個戰術的內容是怎樣的?.
他沒有軍事行動的攻擊.
當中沒有提及任何的武器.
連當時用的.

$^{441}$刀,矛,槍.
其實都是什麼都沒有的.
我想問.
怎樣上陣呢?.
大家玩過戰爭遊戲.
都要有什麼?.
有子彈,有槍.
現在玩戰爭遊戲.
你就是這樣去.
總之避到子彈.
你不死.
走出來限時.
你就贏.
不是這樣.
我看過一齣電影.
《鋼腳靚》.
大家有沒有看過這套電影?.
這套電影原來是第二次世界大戰的時候.
是一個真人真事的改編的電影.
當中的主角.
這個軍人.
他有沒有拿槍上戰場?.
他沒有.
但是.
在他最危急.
戰場打得最激烈的時候.
他居然可以救了百多個同袍.
結果就是.
戰爭完結之後.
他就接受了國家榮譽的勳章.
我相信這個都是史無前例.
但是對我來說.
如果我接到這個命令.
沒有任何武器.
我就一定違抗軍令.
死就死吧.
又或者我靈機一觸.
我身手都很敏捷.
沒有武器.
沒有任何的東西.

$^{481}$我走得快.
只要.
如果在城牆上.
突然間敵人勇猛的時候.
我就左閃右避.
他扔石頭下來的時候也好.
或者飛矛.
飛槍也好.
我就可以走了.
結果就可以.
保了一條命.
但是如果我真的上陣的話.
就算是全軍覆沒.
其實.
於罪何有呢?.
所以對約書亞來說.
如果我們看第六節的時候.
其實約書亞不是做.
我剛才所說的那種決定.
原來他真的在執行軍事的任務.
剛才我提過馴服.
但我也提過屈服.
如果對約書亞來說.
我相信這是一個很有趣的問題.
但我們都知道他是馴服.
上帝的安排.
他完成了這個任務.
其實都值得我們去想一想.
在今時今日.
我們生活在這樣的狀況.
可能相比於約書亞的時代.
真的很不同.
如果我們說馴服的時候.
我們會不會有些難度呢?.
馴服這個題目.
其實我們都說了很多.
會不會說.
是不是有些框框.
框著我們.
令我們做很多的決定.

$^{521}$都不能夠.
很從心所欲.
特別是我人生很重要的前途.
婚姻,家庭,子女.
他們的學業.
要選哪間學校.
這些其實是很大件的事.
如果自己決定.
有空或有時間.
去祈禱讀聖經.
都可以,重大的決定.
當然由自己決定.
這個是不是叫馴服呢?.
又或者我出一點點的犧牲.
上帝,我就是這樣.
來馴服你.
等我過了這個關.
我下一次就真真正正.
來馴服你.
其實馴服和屈服.
我相信兩者是有些分別的.
馴服其實是一種主動的態度.
是放下我們一些的執著.
而是按照聖經.
單單的來去.
依靠馴服上帝.
我想約書亞就給了我們一個.
很好的榜樣.
但如果屈服是怎樣呢?.
屈服我想是.
被迫的感覺.
這好像是一種壓迫感.
如果我不做,心裡面不平安.
甚至.
我又成為我自己一種的挑戰.
因為上帝給的原來和我的不同.
所以實踐使命.
這個身份.
要過得勝的人生.
其實.

$^{561}$我們要面對的.
很多時候都會有不同的危機和挑戰.
可能我們會憂慮.
會有擔心.
其實這個是人之常情.
其實耶穌都勸勉我們.
不要憂慮,不要擔心.
明天自有明天的憂慮.
耶穌其實同時又應許我們.
其實我們每一天都有恩典.
我相信恩典是怎樣呢?.
恩典日日有,日日新鮮.
這個也是聖經裡提醒我們.
對約書亞來說.
他自己的使命的身份.
我相信他是很著緊和在意.
在這麼重要一場戰役.
他親自來到耶利哥城視察.
他遇見了這位軍隊的元帥.
他降服在這個元帥面前.
然後他就知道這個戰事.
原來不是他自己打的仗.
原來這場仗是上帝的仗.
他和上帝打這場仗.
其實他就是要學習信服.
所以他調兵遣將的權力.
軍事的策略.
縱然他是療癒之將.
但他甘願作為副手.
他脫鞋,把鞋脫下.
其實就是表示完全降服.
在這個元帥面前.
弟兄姊妹.
其實我們的人生.
很多時候都說.
人生如戰場.
在這裡要問一問.
究竟我們人生的戰場.
是誰在做元帥?.
誰在做上司?.

$^{601}$還是耶穌是副指揮?.
還是我們將元帥的位置.
交給我們這一位主呢?.
我只是副手.
只是一個副指揮.
在執行上帝給我們的戰術.
所以我們要放下一切.
或者是自以為是.
甚至是我話事的執著.
完全來信靠,信服.
重新在神的面前確立.
上帝給我們的使命身份.
將得勝的生命主權.
讓神來掌管.
因為耶穌才是我們生命的主宰.
第二點就是.
以神為中心.
與神同工.
剛才我們都看過.
三至五節的戰術.
這裡.
我做的powerpoint應該有四項.
因為要遷就IT的問題.
所以就將它分開.
但這個不要緊.
我會再有一些比較簡單的.
比較容易看到.
當他去執行耶和華的戰術的時候.
其實作為一個副指揮.
要作出一些真的要做的軍令.
其實按照第六至第七節的時候.
我們會看到.
他首先會找祭司.
然後就吩咐他們.
按照經文.
我去詳細看的時候.
其實祭司有兩種責任.
他要抬藥櫃.
有祭司抬藥櫃.
有七個祭司.

$^{641}$拿著羊角.
走在藥櫃前.
另外就是吩咐百姓.
第二點應該是吩咐百姓.
百姓也做兩件事.
他做什麼呢.
他負責帶兵器.
另外一些百姓負責繞城.
約書亞有很仔細的軍事命令.
除了調配人手.
調兵遣將之外.
他執行的細節都很清晰.
他的軍令是怎麼樣呢.
祭司要吹羊角.
要發出聲響.
但是同時他發出了一些禁令.
就是不可以呼喊.
不可以出聲.
連一句話都不可以出你們的口.
就是說所有人.
在這個隊伍列陣的時候.
只可以一是腳步聲.
第二就是羊角的吹號聲.
另外的軍令是怎麼樣呢.
他說當我吩咐你大喊吶喊的時候.
你才可以吶喊.
他的部署的隊形.
即是出戰的隊形是怎麼樣呢.
原來是這樣的.
如果按照經文.
我已經將他排列了.
你看看那幅圖.
其實最左手邊應該是.
漏了一些拿兵器的人.
如果那幅圖有拿兵器的人.
我相信這幅是最好的.
按照這裡記載的圖畫.
他有日程.
說完這些隊形.
他怎樣來行動呢.

$^{681}$其實都是按照原稅所吩咐.
給他的做法.
這個就在第8至11節.
15至16節和20節.
這裡是重重複複的.
當我將他來到簡短.
行軍的日程1至6天.
每一天繞成一次.
之後回到營裡休息.
但是很奇怪.
由第2天至第6天.
比第一天早起.
原來打仗可能有時段.
但是時段不同了.
但去到第7天更早.
是要黎明的時候來繞成.
這個是在原稅所講的指令沒有提及.
有什麼行動呢.
行動就是在第7天.
繞成第7次的時候.
軍令說.
剛才我們提及過.
吩咐呼喊的時候.
你們就要大聲呼喊.
但原來軍令裡面加了一句.
因為耶和華已經把城交給你們了.
所以這班軍隊.
他們的行動一呼百應.
加上楊國聲的時候.
整場就倒塌.
百姓就跟著一擁而上.
將城奪取.
我相信這個.
弟兄姊妹.
原來若書雅明白.
這一場的戰役.
雖然說攻打耶利哥城.
但原來實際上.
是一場屬靈的戰爭.
若書雅就是以神為中心的態度.

$^{721}$很嚴謹地執行軍事任務.
當我們剛才看到那個列隊的時候.
其實這個陣營.
就是當他們在西乃山.
按照上帝所吩咐做各樣的物件.
有這個藥櫃.
又有怎樣去制師.
各樣的身份.
他們所做的事.
完全列出來的時候.
其實若書雅知道上帝的心意.
藥櫃的擺列.
原來是以神為中心.
但藥櫃前後有什麼呢?.
是以士兵拿兵器來保護.
其實藥櫃所代表的就是.
上帝與人同在.
是上帝與人相會的地方.
也是上帝與人說話的地方.
如果是這樣的時候.
若書雅為什麼要找一些軍隊.
一些拿兵器的人前後來保護呢?.
我相信若書雅是一個很認真.
她要向百姓表明.
藥櫃是神所吩咐建造的.
她有臨在的標記.
同時也讓在耶利哥城裡面的敵人.
來表明.
宣示這場戰爭是屬於神的.
神會必定為以色列人.
爭戰必定的來取勝.
所以也向周圍的百姓.
向軍隊展示.
其實神的藥櫃.
我們作為人.
是屬於祂的子民.
是要好好保護祂.
不可以讓藥櫃有任何的閃失,差池.
今日弟兄姊妹.
我們每一個屬上帝的人.

$^{761}$我們叫基督徒.
我們因為相信耶穌基督.
藥櫃相信對我們來說.
沒有一個實質的意義.
但它是一個象徵意義的存在.
因為藥櫃裡面有法板.
其實我們內心裡面的空間.
會不會成為心靈的藥櫃.
有上帝的律法,典章.
神的話常常讓我們思想呢?.
同時我們思想神的話的時候.
我們能夠活出神的教導.
來見證上帝給我們的永生之道.
以致不要枉費了神給我們的使命的身份.
如果我們看過列隊的時候.
我們都看到還有些羊角.
祭司拿著羊角吹的角聲.
和以色列人的呼喊聲.
其實角聲是代表什麼呢?.
角聲是代表神的同在.
如果大家還記得在出埃及記的時候.
當摩西在西乃山上.
準備領受神所頒佈的誡命律例的時候.
當時就有羊角聲四起.
其實震動了整個西乃山.
甚至山下的百姓.
原來羊角聲是代表了威嚴的神.
有威榮的神就臨在.
所以羊角聲的吹響.
就是表示神和以色列人一起去爭戰.
當然羊角聲也有一個.
為了作戰勝利而喜慶的含義.
所以如果是這樣的話.
就會明白為什麼這個軍令之前有一句.
這個是耶和華的戰爭一定要戰勝.
另外一方面.
這種羊角聲也有震懾敵人的心理作用.
所謂的心理戰.
那些百姓的吶喊聲其實是怎樣的呢?.
當然你會知道是第七天.

$^{801}$繞成第七次的時候.
才會出聲.
其實就是有助威提升士氣的形勢.
更加造成敵人的心理壓力.
所以耶利哥城.
一定會倒塌.
一定會.
傳言的毀滅其實是神的主動.
出擊.
一定勝利.
人 軍隊.
只是參與在神的工作裡面.
我們要順靠這位.
出軍令的這位上帝.
我們要謙卑配合神的策略.
行神所吩咐 與神同行.
弟兄姊妹.
其實不出聲容易嗎?.
當約書亞提到的時候.
他說不可呼喊 不可出聲.
連一句話都不可出聲的時候.
我想問有什麼作用呢?.
其實如果每個人都出聲.
在行軍的時候會怎樣呢?.
這是不是戰術呢?.
好像在遊行.
每天繞一圈就可以了.
可能會有很多議題.
不如我們的戰術應該改良一下.
約書亞不要這樣打.
很多的批評.
很多的不信任.
其實會影響士氣.
所以祭司軍隊.
他們要執行這個軍令的時候.
其實他們是不需要出聲.
當然實踐的時候一定會很困難.
因為如果不合情合理.
不是經過人生經歷.
去分析的時候.

$^{841}$要我們這樣行其實是很難的.
其實在教會很多時候都會這樣.
執行教會所定的計劃事工.
或者小組 團契.
當這樣執行一定有不同的意見.
是不是說這些意見不好呢?.
是不是一個適當的時候來表達呢?.
其實約書亞的吩咐有一個屬靈的提醒.
聖經裡面提醒我們.
他說靜默有時言語有時.
其實這一句的經文我想大家都很熟悉.
有沒有留意那個次序呢?.
其實原來我們靜默才是首要.
與神同工.
其實是要配合神的戰術.
當要部署執行每一步的時候.
祂要信服.
結果就帶來了勝利的結果.
求主幫助我們今天在這個聲音很多的世代.
學習安靜在神的面前.
領受神的心意.
甚至將我們的內心.
個人的內心裡面來轉化.
讓神檢視我們的生命.
我們的使命方向目標.
教正了我們就能夠.
做任何的事都是以神為中心.
我們多一些少看自己.
以致只看著神.
看著神的時候.
我們就能夠與神同工.
最後第三點.
我們如何能夠活出使命人生呢?.
其實就是戒善要分明.
要堅持活出聖潔.
經文是在第六章的十七至十九節.
很想大家跟我一起讀出這段經文.
預備起.
「這承和其中所有的.
都要在耶和華面前毀滅.

$^{881}$只有妓女拉合和與她家中所有的可以存活.
因為她隱藏了我們所打發的使者.
至於你們務要謹慎.
不可取那當滅的物.
恐怕你們取了那當滅的物.
就連累以色列的全營.
使全營受咒坐」.
原來我們放下自己的身段執著.
謙卑我們自己參與神奇妙獨特的行事計劃.
與神同工.
以神為中心.
其實是一個很寶貴的機遇.
神這樣邀請我們.
我們就要活出神的聖潔.
來過我們得勝的人生.
所以我們生活上要界線分明.
剛才這段經文提醒我們.
「弟兄姊妹」攻打耶利哥的戰役.
大家都覺得很戲劇性.
特別是今天的打仗.
尤其是今天的科技戰爭.
根本沒有可能沒有武器.
可能武器還先進過約書亞的時代.
所以說這場戰役是後無來者.
我們很難想像這樣的戰術能夠打贏仗.
只有摧割人生的吶喊.
就這樣繞成十三次.
結果就是這樣勝利.
但是勝利之後.
我們如何處理所得到的戰利品呢?.
很多時候我們都說勝利會充分頭腦.
充分頭腦的時候其實甚麼都不記得.
所以在這裡.
當約書亞要執行元帥的軍令的時候.
他都有部署.
有很多的步驟.
他說戰利品不可以拿.
為甚麼呢?.
原來耶利哥城還有一個意思.
是月城.

$^{921}$月城的意思其實按照考古學和很多聖經的學者說.
原來城內的人.
應該所有人都是拜月亮.
所以有月城這個稱呼.
原來神要百姓將戰利品據為己有.
是不可以這樣做的時候.
其實和當時古代的國家行軍打勝仗的做法是不同的.
古代行軍幫國打勝仗之後.
戰勝國是怎樣呢?.
戰勝國原來是會用霸主自居的身份.
將手下敗將變成俘虜.
這班是他們的臣民.
所以戰勝國做任何事情.
這班人都要服從.
如果男的強壯的會怎樣呢?.
做奴隸.
沒用的窮的會怎樣呢?.
可能會成為他們的婦女.
可能就會殺死他們.
而女的會怎樣呢?.
可能會成為別人的妻妾.
甚至可能在戰勝國的神廟裡.
做廟祭供奉給神明.
諸如此類的做法.
所以約書亞吩咐軍隊百姓.
他們要將戰利品全部毀滅.
無論是人還是物.
原因是甚麼呢?.
從耶利哥城看到的拜月神.
可想而知整個迦南地區會是怎樣.
是流行不同偶像的崇拜.
在他們的宗教上.
這種崇拜可能包括了一些淫亂.
甚至會將一些嬰孩來祭祀.
去教誨.
所以道德上充滿了很多敗壞.
暴力.
如果不消滅他們.
不全部成為當滅之物的時候.
以色列人如果繼續和他們一起.

$^{961}$用他們的物件.
和他們的人有交往的時候.
他們都會受到感染.
像瘟疫一樣感染他們.
隨波逐流.
因此毀滅耶利哥的人和物.
正正就是代表了.
耶利哥城裡面那種罪惡充斥.
大家還記得.
「所多瑪俄摩拉」嗎?.
毀滅是需要.
毀滅審判是彰顯神的公義.
至於約書亞所吩咐的.
「拉合妓女」是怎樣呢?.
其實她本應是要滅亡的.
但因為她選擇了.
跟從敬拜以色列的神.
同時又暗中幫助探子.
結果以色列人手約.
成家得救.
攻打耶利哥的地方.
是約書亞的第一戰.
所以要很清晰.
迦南人和以色列人.
他們要有一條很清晰的界線.
約書亞正在進行.
上帝給予摩西的命令.
「因為你歸耶和華里神為聖傑的公民」.
「耶和華里神從地上的萬民中」.
「揀選你,特作自己的民」.
所以在這個情況下.
知道以色列人和神.
有一種很獨特的關係.
是神所揀選.
所以神的選民一定要聖傑.
其實聖傑是怎樣呢?.
是要強烈地.
想和其他不潔的東西分別出來.
不可以和他混為一談.
所以以色列人進入迦南地.

$^{1001}$要將當地的文化和偶像.
全部清除.
然後過一個聖傑的生活.
其實在座很多父母都知道.
有時候自己所愛的子女.
會有不好的習慣.
不良的風氣影響.
其實在所難免.
我們都會用盡一些方法去阻止.
我相信上帝都是這樣的心意.
迦南人敬拜偶像,假神.
牽涉了很多這樣的事情.
所以以色列人進入迦南地.
他不只是奪取上帝給他應許之地.
從此快快樂樂地生活.
而是同時要求他的百姓.
立約的子民.
要聖傑分別出來.
不要來到一個敗壞的生活.
所以在第十八節.
約書也很嚴厲地警告.
提醒不要謹慎.
因為如果不是這樣.
拿了一些當滅的物件.
自己就成為當滅的.
是被上帝毀滅.
讓我們今天是神所揀選.
是神的子民.
活在新約的時代.
其實上帝,耶穌,聖靈.
都很重視我們聖傑的生活.
所以我們要在清晰的界線上生活.
所以這條界線不是一個框框.
框住我們.
其實上帝立下的界線.
是要保護我們.
以至我們不要落入犯罪.
跌入撒旦的隱憂.
過界是怎樣?.
過界很平常.

$^{1041}$你和我都會過界.
有時過馬路.
有時真的會過界.
過界就下次不過.
如果踩界又會怎樣?.
大家有沒有踩界?.
在地鐵,火車那些.
Stand behind the yellow line.
坐在界線上.
現在的服務員都說.
站後一點.
我們會聽到甚麼?.
我踩著界線.
我不是在界線前面.
起碼車經過都不會碰到.
會是這樣.
弟兄姊妹.
如果踩界不是按照神的界線來生活.
其實會怎樣?.
會模糊了我們.
正如我們今天.
在這個社會裡.
其實我們沒有辦法避免.
扭曲了的道德倫理.
或者言論.
或者物質.
我們要去享受.
其實我們沒有辦法不受影響.
但我們有沒有清晰的界線.
來給我們很好的生活?.
其實整本聖經.
就有很清楚這種界線的標準.
當中沒有一些陷糊的地方.
我相信.
弟兄姊妹在你們的工作場所.
我自己的工作.
過去的工作也是一樣.
特別是商界.
都會受到很多的衝擊.
都會受到很多的挑戰.

$^{1081}$但我們更需要去倚靠神.
去到最後.
其實我們的人生是帶著使命.
其實我們的心態是怎樣呢?.
其實我們要放下一些執著.
信靠信服主.
以神為中心.
與人與神來同工.
以及我們生活上要有界線.
以及按照上帝教導我們.
要過聖潔的生活.
我們能夠處變北京.
在這個充滿了挑戰.
以及不穩定的世界.
特別是我們現在這幾個月.
開始受到一些疫情.
疫症的威脅之下.
可能都影響了我們各樣的生活.
但是會不會這樣就影響了.
我們使命的身份呢?.
求神幫助我們.
能夠以生命來見證耶穌.
活出不一樣的生命.
我們一起來禱告.
讓我們生命裡面經歷你.
求你教導我們.
好像約書亞那樣.
他來信服信靠你.
他以上帝為中心.
與神同工.
也有清晰的界線.
有聖潔生活的標準.
最後幫助我們.
以致我們在我們現今的世代裡面.
學校,約書亞.
我們能夠去活出.
我們有使命的人生.
我們獻上禱告.
奉主耶穌基督的名.
阿們.

\newpage



\section{約書亞記 9:1-21}
\label{sec:ruBv0LtC8pg}
\textbf{2020年3月15日 崇拜直播｜約書亞記9\_1-21}
\newline
\newline
連結: \href{https://youtube.com/watch?v=ruBv0LtC8pg}{\texttt{https://youtube.com/watch?v=ruBv0LtC8pg}} ~~~~ 語音日期: 2020-03-15
\newline
\newline
\hyperref[sec:w8cFfcUefjQ]{\small{< < < PREV SERMON < < <}}
~
\hyperref[sec:index_chronic]{\small{[返順時目]}}
~
\hyperref[sec:q99RT1cz_zM]{\small{> > > NEXT SERMON > > >}}
\newline
\newline
約書亞記 9:1-21
\newline
\begin{longtable}{cl}
\hline
\hline
章節 & 經文 (和合本修訂版)\\
\hline
9:1 & \begin{tabularx}{0.7\textwidth}{X} 約旦河西，住山區、低地和沿大海一帶直到黎巴嫩的諸王，就是赫人、亞摩利人、迦南人、比利洗人、希未人、耶布斯人的諸王，聽見這事， \end{tabularx} \\ \\ \relax
9:2 & \begin{tabularx}{0.7\textwidth}{X} 就都聚集，同心合意要與約書亞和以色列人作戰。 \end{tabularx} \\ \\ \relax
9:3 & \begin{tabularx}{0.7\textwidth}{X} 基遍的居民聽見約書亞向耶利哥和艾城所做的事， \end{tabularx} \\ \\ \relax
9:4 & \begin{tabularx}{0.7\textwidth}{X} 就設詭計，假扮使者出去。他們拿舊布袋和破裂補過的舊皮酒袋馱在驢上， \end{tabularx} \\ \\ \relax
9:5 & \begin{tabularx}{0.7\textwidth}{X} 將補過的舊鞋穿在腳上，把舊衣服穿在身上，作食物的餅都又乾又長了霉。 \end{tabularx} \\ \\ \relax
9:6 & \begin{tabularx}{0.7\textwidth}{X} 他們到吉甲營中約書亞那裡，對他和以色列人說：「我們是從遠地來的，現在求你與我們立約。」 \end{tabularx} \\ \\ \relax
9:7 & \begin{tabularx}{0.7\textwidth}{X} 以色列人對希未人說：「或許你是住在我附近的。若是這樣，我怎能和你立約呢？」 \end{tabularx} \\ \\ \relax
9:8 & \begin{tabularx}{0.7\textwidth}{X} 他們對約書亞說：「我們是你的僕人。」約書亞對他們說：「你們是甚麼人？是從哪裡來的？」 \end{tabularx} \\ \\ \relax
9:9 & \begin{tabularx}{0.7\textwidth}{X} 他們對他說：「你的僕人是因耶和華－你神的名從極遠之地來的。我們聽見他的名聲，他在埃及所做的一切， \end{tabularx} \\ \\ \relax
9:10 & \begin{tabularx}{0.7\textwidth}{X} 以及他向約旦河東的兩個亞摩利王，希實本王西宏和在亞斯她錄的巴珊王噩所做的一切。 \end{tabularx} \\ \\ \relax
9:11 & \begin{tabularx}{0.7\textwidth}{X} 我們的長老和我們當地所有的居民對我們說：『你們手裡要帶著路上用的乾糧去迎接以色列人，對他們說：我們是你們的僕人。現在求你們與我們立約。』 \end{tabularx} \\ \\ \relax
9:12 & \begin{tabularx}{0.7\textwidth}{X} 我們出來要往你們這裡來的那日，這從我們家裡帶出來的餅是熱的；看哪，現在這餅又乾又長了霉。 \end{tabularx} \\ \\ \relax
9:13 & \begin{tabularx}{0.7\textwidth}{X} 這些皮酒袋，我們盛酒的時候還是新的；看哪，現在已經破裂了。我們這些衣服和鞋，因為路途非常遙遠，也都穿舊了。」 \end{tabularx} \\ \\ \relax
9:14 & \begin{tabularx}{0.7\textwidth}{X} 以色列人收下他們的一些食物，但是沒有求問耶和華的指示。 \end{tabularx} \\ \\ \relax
9:15 & \begin{tabularx}{0.7\textwidth}{X} 於是約書亞與他們建立和好關係，與他們立約，讓他們存活；會眾的領袖也向他們起誓。 \end{tabularx} \\ \\ \relax
9:16 & \begin{tabularx}{0.7\textwidth}{X} 以色列人與他們立約之後，過了三天才聽說他們是近鄰，住在附近。 \end{tabularx} \\ \\ \relax
9:17 & \begin{tabularx}{0.7\textwidth}{X} 以色列人起行，第三天就到了他們的城鎮，他們的城鎮是基遍、基非拉、比錄和基列‧耶琳。 \end{tabularx} \\ \\ \relax
9:18 & \begin{tabularx}{0.7\textwidth}{X} 因為會眾的領袖已經指著耶和華－以色列的神向他們起誓，所以以色列人不擊殺他們。全會眾就向領袖發怨言。 \end{tabularx} \\ \\ \relax
9:19 & \begin{tabularx}{0.7\textwidth}{X} 眾領袖對全會眾說：「我們已經指著耶和華－以色列的神向他們起誓，現在我們不能碰他們。 \end{tabularx} \\ \\ \relax
9:20 & \begin{tabularx}{0.7\textwidth}{X} 我們要這樣對待他們，讓他們存活，免得因我們向他們所起的誓而憤怒臨到我們。」 \end{tabularx} \\ \\ \relax
9:21 & \begin{tabularx}{0.7\textwidth}{X} 領袖對會眾說：「讓他們活著吧。」於是他們照領袖所說的，為全會眾作劈柴挑水的人。 \end{tabularx} \\ \\
[1ex]
\hline
\hline
\end{longtable}
$^{1}$今日的經訓是記載在.
約書亞記九章一至二十一節.
聖經舊約第二百七十二至二百七十三頁.
約書亞記九章第一節.
約旦河西駐山地高原.
並對著黎巴嫩山.
元大海一帶的諸王.
就是黑人,阿摩尼人,迦南人,比利西人,.
希美人,耶布斯人的諸王.
聽見這事就都聚集.
同心合二的要與約書亞和以色列人爭戰.
欺騙的居民聽見約書亞向耶利哥和艾聖所行的事.
就設詭計假充是者.
拿舊口袋和破裂鳳寶的舊皮酒袋.
拖在爐上,將倒過的舊鞋穿在腳上,把舊衣服穿在身上.
他們所戴的餅都是乾的,長了沒了.
他們到吉格營中見約書亞.
對他和以色列人說.
我們是從遠方來的,請求你與我們立約.
以色列人對這些希美人說.
只怕你們是住在我們中間的.
若是這樣,怎能和你們立約呢?.
他們對約書亞說.
我們是你的僕人.
約書亞問他們說.
你們是什麼人?是從哪裡來的?.
他們回答說.
僕人從極遠之地而來.
因聽見耶和華理神的名聲和他在埃及所行的一切事.
並他向約旦河東的兩個阿摩利王.
就是希實本王西王和在阿斯他祿的巴山王岳.
一切所行的事.
我們的長老和我們那地的一切居民對我們說.
你們手裡要帶著路上用的食物.
去迎接以色列人.
對他們說.
我們是你們的僕人.
現在求你們與我們立約.
我們出來要往你們這裡來的日子.
從家裡帶出來的這餅還是熱的.

$^{41}$看吶,現在都乾了,長了,沒了.
這啤酒袋我們姓酒的時候還是新的.
看吶,現在已經破裂.
我們這衣服和鞋.
因為道路甚遠也都穿舊了.
以色列人收了他們些食物.
並沒有求問耶和華.
於是約書亞與他們講和.
與他們立約.
容他們活著.
會眾的首領也向他們起誓.
以色列人與他們立約之後.
過了三天才聽見他們是近鄰.
住在以色列人中間的.
以色列人起行.
第三天到了他們的城邑.
就是基片基費拉比祿基列耶林.
因為會眾的首領.
已經指著耶和華以色列的神.
向他們起誓.
所以以色列人不擊殺他們.
全會眾就向首領發願言.
眾首領對全會眾說.
我們已經指著耶和華以色列的神.
向他們起誓.
現在我們不能害他們.
我們要如此待他們.
容他們活著.
免得有憤怒.
因我們所起的誓.
臨到我們身上.
首領又對會眾說.
要容他們活著.
於是他們為全會眾.
作了劈柴挑水的人.
正如首領對他們所說的話.
六個王決定聯手對抗以色列人.
而機騙人則是.
如經文所說.
設詭計騙以色列人.

$^{81}$與他們立約講和.
事實上他們本來是.
上帝要滅絕審判的對象.
但為何最後機騙人逃過厄運.
我們說成為繼幾里拉合一家之後.
另一個被拯救的迦南族群.
我們今天會和大家看看.
到底是什麼原因.
當然事件中亦記載.
以色列人受騙的經歷.
其實亦很值得我們去借鑒.
所以在這裡.
先和大家看幾幅圖畫.
麻煩音響室.
我想透過這幾幅圖畫.
來測試一下大家的眼睛.
先看第一幅.
你覺得裡面的家具擺設.
是否正常.
正不正常.
正常吧.
好,播放來看看.
原來你發覺用了距離感.
來做到好像.
你看的東西.
好像真實擺設的比例一樣正常.
再看第二幅.
這四條軌道.
是否全部向上斜.
看上去好像是.
但實際呢.
你發現其實只是用了另一個角度看.
其實它們本來是向下斜.
第三幅.
你看到你的.
左手邊.
或者右手邊.
即是你的.
是否有架飛機在這裡.
真還是假的呢.

$^{121}$我們又看看.
在電腦特技效果裡.
真的不知道是真是假的.
最後一張.
是否看到有四個人坐在梳化裡.
看真一點.
是否真的有四個人.
原來有三個是紙版人.
只有一個真人.
還有一張單人梳化.
Ben 姐妹你看到.
製作這幾幅圖片的人.
其實運用了我們說的視覺效果.
令到你看上去.
產生一種錯覺.
以為你眼見的.
就是事實的真相.
其實以色列人當時.
也面對類似的情況.
我們看第24節.
麻煩你.
欺騙人首先在往後.
交代他們為何要用詭計的原因.
他們是因為.
害怕被滅絕.
所以為了活命.
他們就想和以色列人.
立約講和.
但是又害怕.
如果他們真的講真話.
說自己是迦南地的居民.
後果會怎樣.
他們很擔心會被殺害.
於是.
他們就在第六節說.
我們是從遠方來的.
就好像.
不難理解.
可能之前也聽過有市民.
召喚救護車服務.

$^{161}$但其實他之前回過內地.
但如果怕被人問到.
他要講真話的時候.
就很擔心.
別人可能不幫他.
又或者擔心.
如果真的有事就要強制隔離.
為了保障自己.
所以他隱瞞這個真相.
但當然最後一定查到.
而你看到.
欺騙人他要為了.
增加自己說話的可信性.
於是.
他就將那些舊的口袋.
將那些縫補.
破爛了那些舊的皮袋.
放在爐上.
又將那些補過的鞋.
和那些破舊的衣服.
穿在身上.
還帶著那些.
本來是乾又長了毛.
實際原意可能是又乾又碎.
毛其實可以解作.
碎了的意思.
乾的東西自然都碎了.
他試圖透過這些我們說的精心預設的道具.
來誤導以色列人.
相信他們真的從遠方來.
而當以色列人和約書亞他們起疑心的時候.
欺騙人也似乎想好了對策.
他們立刻就來一招.
轉移視線.
如果我們細心看.
第八節開始.
似乎他們都想過.
當人家真的起疑心的時候.
他們要怎樣去回應.
你會看到欺騙人他沒有直接去回答.

$^{201}$約書亞或者說.
以色列人的提問.
他只是一再很謙恭地.
叫自己做什麼.
服人.
叫了自己三次.
我是服人來的.
對以色列人來說這句說話.
合不合聽.
合不合耳.
很合聽的一句說話.
為什麼這麼說.
你想想過往幾百年以色列人在埃及做什麼.
做奴隸.
現在.
竟然反過來有人.
肯做他們的奴隸.
自己一下子就提升.
做了什麼.
做了主人.
肯定感覺飄飄然.
另一方面.
欺騙人亦都說.
他們是因為聽到耶路亞話神的名聲.
聽見耶路亞話的大能.
他怎樣在埃及.
怎樣在約旦河東.
所行的一切事.
所以山長水遠.
專登來到去拜會以色列人.
留意欺騙人.
其實本來他們知道.
近日耶利哥城.
和艾城所發生的事.
不過專登隱瞞不說.
因為他們怕被人識穿.
其實他們本來就是住在附近的.
欺騙人.
試想像下.
你是以色列人.

$^{241}$竟然有一班人.
素未謀面但肯派隨萬難.
從遠方冒名.
來找你.
還表示對你的信仰.
因為他們相信神非常景仰.
還主動來投靠你.
說要做你的僕人.
以色列人.
以前有試過如此威風嗎.
沒有試過.
這樣被人抬舉.
其實自然.
他的戒心會怎樣.
會大大降低.
就在這個時候.
欺騙人再怎樣.
拿預先設計好的道具出來.
林信.
所帶來的東西.
我們出發時在家裡拿來的餅.
還是熱的.
現在已經乾到碎了.
你看看我們這個皮袋.
原本出門口.
放酒進去的時候還是新的.
現在已經破爛了.
你看看我們穿的衣服.
穿的鞋.
因為路程實在很遠.
所以都破舊了.
沒有花沒有架.
你相信我們.
真的從遠方來.
我們真的很想和你立約的嗎.
欺騙人.
一早預設好.
要怎樣做.
所以結果你會看到.
第十四到十五節.

$^{281}$經文這裡說.
以色列人受了他們的食物.
並沒有求文.
耶和華.
不是一定吃了那些.
又乾又碎的餅.
意思是.
以色列人只是單憑.
那些又乾又碎的餅.
來做證據.
單憑他們眼睛.
所見到的事物.
就接納了欺騙人所說的.
是一個事實的真相.
但相比之下.
他們明明.
是要以上帝的名.
來去起誓.
一個這麼重大的事件.
他們竟然怎樣.
不問上帝.
竟然跳過了上帝.
其實上帝曾經.
祝福過他們.
當要做一些重大的判斷.
特別是遇到困難.
遇到疑惑的時候.
可以先透過大祭司身上的.
污淩和土命.
來求文.
過程其實不複雜不難.
但為什麼他們寧願.
倉促自己去做一個判斷.
都不事先去.
求文一下上帝的指引.
這一方面很值得我們去想.
如果你有留意其實.
以色列人剛剛經歷完的.
是一個逆境.
變回順境.

$^{321}$如果你還記得.
處理完阿干犯罪的事件.
之後他們再一次怎樣.
打勝仗.
征服了艾城.
再加上現在.
有人怎樣.
特地來做你的主人.
感覺肯定好不好.
非常良好的感覺.
感覺良好的時候.
自然.
對自己辨別事物的能力.
就會怎樣.
更加有自信更加有把握.
覺得自己可以有能力.
處理到眼前很多的事情.
人在這個.
情況的裡面.
很容易就會忘記了.
自己本來是僕人的身份.
是吧.
當我們被人叫做主人.
被人抬轎多的時候.
就會很容易忘記了.
自己其實是上帝僕人的身份.
一點都不出奇.
你想想順境的時候.
加上感覺良好.
自信心爆棚的時候.
我們就會很容易.
慢慢忘記了.
其實我們每一日都要尋求上帝的指引和帶領.
這個有時都會是.
我們生活一些的寫照.
所以其實.
我想我們沒有人喜歡.
面對逆境.
但是上帝卻是喜歡使用逆境.
上帝卻是喜歡通過逆境.

$^{361}$將迷失了的人.
帶回到他身邊.
迫使我們.
來尋求他.
當然我想.
尋求上帝不會馬上有個答案.
不會馬上解決所有問題.
但是肯尋求上帝的人.
他和上帝的關係.
一定會不斷提升.
同時.
也可以幫自己建立到一種生活的態度.
就是每一日.
先將自己.
先將你要處理的事情.
交託了給上帝.
你要知道.
離開了他.
其實你什麼都不能夠做.
這樣的時候你就會.
懂得尋求上帝的帶領.
等候他的指引.
保守自己.
周瑜最近的情況.
我想大家都類似.
少了出街.
少了很多.
和外面接觸聯繫的機會.
多了很多時間.
特別在家.
但是有時.
又不知道怎麼打.
或者不知道怎麼好好去運用.
特別如果你是做家長.
小朋友讀小學讀幼稚園.
就更加頭痕.
沒有網上功課的時候.
就擔心什麼.
擔心他懶散.
擔心他戶外活動不夠.

$^{401}$要安排很多東西.
到開始有網上學習.
網上上學.
網上功課的時候.
你又煩什麼.
下班之後.
你還要做多另一份工作.
來跟進一下.
他能不能搞定.
有沒有真的依時依後.
按照學校的規定來完成他的學習.
如是者.
這樣都幾多個月.
都過幾兩個月.
繼續這樣好不好.
都挺難受的.
真的.
不是說笑.
沒有病都會變怎樣.
都有病.
所以弟兄姊妹.
我想或許這個都是一個好時機.
我們回到上帝的身邊.
尋求祂的心意.
怎樣在那個.
不屈悶的日子.
怎樣在擔憂的裡面.
求神幫助你.
重新激活過來.
重新去調教.
我們生活的步伐.
和一些模式.
來面對眼前的局面.
早幾天收到.
關信輝執事一個WhatsApp.
可能大家都有.
看過那條片在網上.
關於關信輝執事.
最近發了一條YouTube的影評.
問了他容許我今天在這裡.

$^{441}$分享少少他的見證.
話說其實關執事他.
有一日很悶.
在家裡很屈悶.
呆呆望著灰灰的天空.
很有意境.
正在下沉之際.
他就突然想起.
那些曾經令他很振作.
振奮的電影情節.
很觸動他心靈的對白.
於是決定重新振作.
不再讓自己燒沉下去.
他想想可以做些什麼.
於是他最終想到.
將那部已經鋪滿灰塵的相機.
拿出來.
他要突破他自己.
作出一些新嘗試.
在一個小小的工作間裡.
他在網上.
分享過往的電影.
裡面的情節.
如何鼓勵他.
亦都希望在過程裡.
如何幫助別人.
幫他們帶來一些共鳴.
帶來一些力量.
我看到其實上帝.
透過關道演在網上的分享.
這個新嘗試.
我們說好像在疫情下.
一個救兵.
不單單是首先.
將他自己救起來.
幫他重拾對上帝.
昔日那種熱情.
那種負擔.
他在這裡.
也要將他.

$^{481}$同時也要將其他.
正在下沉的人.
救起來.
弟兄姊妹.
上帝在疫情裡.
他不單單.
要動員關信輝.
執事.
上帝亦都要動員.
我們每一個.
事實上你會留意在這兩個月裡.
上帝已經動員了.
很多弟兄姊妹.
來作出不同的新嘗試.
讓我們可以在疫情裡.
繼續保持和上帝.
和人來聯繫.
你會看到.
其實我們今天.
能夠有網上的崇拜.
不單單只是成年的崇拜.
兒童的崇拜.
少年的崇拜.
其實是因為幕後有一班弟兄姊妹.
台上台下.
每個星期很用心去預備.
又或者有些小組.
亦都有一些新突破.
透過網上的視像會議.
來和弟兄姊妹.
保持一個相交和守望.
亦都有弟兄姊妹.
上帝給他們有感動.
他們自發捐贈一些物資.
來轉送給有需要的人.
今天.
剛好我們一班的傳道同工.
亦都完成了.
40天的崇拜.
我們就在這裡.

$^{521}$亦都完成了40天的.
牧者家書.
當然你會看到.
間中同工還會粉墨登場.
跳舞.
都是一個突破.
弟兄姊妹我們見到其實有很多人.
願意在疫情的限制裡面去突破自己.
來挑戰自己.
為了上帝的緣故.
為了幫助身邊的人.
親近上帝的緣故.
多走一步.
並且這一步.
亦都勒著了.
眺望了身邊的人.
願意起來回應上帝.
在這個世代所發出的呼召.
為的是要.
讓上帝贏得.
更多人的心.
這是經文.
第一個給我們去思考.
的一個課題.
而繼續看.
餘下的經文.
讓我們也透過.
以色列人的經歷.
來反省一些功課.
第十六節.
幾天之後.
以色列人就發覺自己.
中伏了.
中計了.
其實是欺騙人.
是附近的居民.
是當滅的迦南人.
既然是這樣.
處理方法是怎樣.
很簡單.

$^{561}$就是執行上帝的吩咐.
把欺騙人擊殺.
這似乎是一個.
最直接簡單的處理方法.
但是.
領袖們就說了.
不行,不可以這樣做.
因為我們已經納了約.
要容許.
欺騙人存活.
所以問題就來了.
以色列人.
現在處於一個.
兩難的境況裡面.
正所謂.
殺不殺.
欺騙人.
都得罪上帝.
為何這樣說.
因為如果他們選擇.
不消滅欺騙人.
就是違背.
上帝以前的吩咐.
但是如果選擇.
消滅欺騙人又怎樣.
就是違背了.
以上帝的名字.
所起的誓約.
所以.
現在已經有很多大頭佛.
怎樣收科呢.
背誦於是.
就向領袖們發願意.
19到21節.
領袖們.
就站出來解釋.
為何不能夠.
擊殺欺騙人的原因.
以及如果我們真的.
擊殺了欺騙人.

$^{601}$到底會有什麼後果.
他們就說.
我們已經指著上帝的名.
來起誓.
是一件非常之.
嚴肅.
不可以違背的事.
如果我們這一刻.
違約.
推翻了先前的盟誓.
來加害欺騙人.
上帝.
一定會追討.
他的憤怒一定會淋到.
我們的身上.
所以我們要遵守.
在上帝.
面前所立的誓約.
要容許.
欺騙人活著.
在我們中間.
來挑柴.
來劈柴挑水.
領子妹確實.
我們看到.
以色列人之前確實.
是很務實.
是很倉促.
所以做了一個錯誤的決定.
但有一樣.
很值得我們去借鏡.
和欣賞的就是.
他們沒有選擇一錯.
再錯.
沒有選擇用一個新的錯誤.
來處理.
舊的錯誤.
令到事情.
變得更加的爛.
以色列人.

$^{641}$卻是選擇承擔.
犯錯了的後果.
即使是.
受騙吃虧.
仍然要堅守.
以上帝名義.
所起的誓約.
仍然要優先.
尊重維護.
上帝的聲譽.
不單止這樣.
你會看到.
他們還是設法.
來補救之前的過失.
他們很盡力地.
去將那個後果.
轉化為一件.
未事.
為什麼這樣說?我們留意一下.
剛才沒有讀的一段經文.
22到27節.
你看到約書亞在善後.
的過程裡面.
顯出上帝給他一種智慧.
簡單的說就是.
約書亞就召了他們來.
對他們說.
為什麼欺控我們?.
你們離我們.
我們離你們甚遠呢?.
其實你們是住在我們中間.
現在你們是被救助的.
你們中間的人.
必斷不了作牢捕.
為我神的殿.
作劈柴挑水的人.
他們.
即是欺騙人.
回答約書亞說.
我們為你們的緣故.

$^{681}$甚怕喪命就行了這事.
現在我們在你們手中.
你以怎樣.
待我們為善為正.
就這樣做吧.
於是約書亞這樣待他們.
救他們脫離以色列人的手.
以色列人就.
沒有殺他們.
當日約書亞使他們.
在耶和華所要選擇的地方.
為會眾和耶和華的.
壇作劈柴挑水的人.
直到今天.
表面上我們看來.
約書亞確實沒有很嚴格.
遵守上帝.
要滅絕.
迦南人的吩咐.
不過我們停一停.
我們先去.
問一問一個問題.
上帝.
為何要滅絕迦南人.
上帝為何要這樣做.
是否純粹.
一種種族的屠殺.
肯定不是.
因為上帝拯救了.
同樣是迦南人的拉合一家.
同樣我們知道.
上帝拯救了.
本來當滅的欺騙人.
上帝之所以.
要滅絕迦南人.
背後的意思其實是.
要除去.
所有對抗.
上帝的力量.
要除去那些使得以色列人.

$^{721}$離棄上帝的力量.
而現在.
欺騙人他願意怎樣.
他坦白承認自己.
從前說謊的過犯.
他沒有用一個新的謊言.
來掩蓋.
他的過犯.
並且他都願意.
交在約書亞的手裡.
承擔後果.
順從約書亞的吩咐.
永遠做上帝的奠.
和祭壇.
那個劈柴挑水的僕人.
那意味著什麼呢?.
弟兄姊妹.
欺騙人他不單單.
只歸降了.
以色列人.
某程度上.
約書亞更加要他們歸入.
以色列人的信仰.
約書亞現在要將.
欺騙人.
一直留在誰的身邊.
留在上帝的身邊.
一直要活在.
上帝的監管下.
如是者.
他們就很難去.
引誘以色列人.
離棄上帝.
因為這樣.
約書亞最後.
亦都好似經文所講.
將本來當滅的.
欺騙人拯救了出來.
弟兄姊妹我想整件事.
我們學了一些很重要.

$^{761}$處理事情的功課.
首先他提醒我們.
只埋怨.
真的解決不了問題.
與其.
積埋整套的冤屈氣.
一直折磨自己.
折磨身邊的人.
不如坦白接受那些.
發生了的事實.
坦然地去接受.
我們沒有辦法.
改變過去.
第二方面.
趁著我們仍然有機會.
其實我們可以努力.
做很多修補的工作.
做很多善後的工作.
將那些.
後果.
惡果減到最低.
甚至轉化為一些.
未善的旨意.
我想.
忘羊補牢未為晚也.
大家都明白這個道理.
另一方面.
這件事讓我們看到.
再深思細密的人.
再優秀的人.
都會有自己.
失落的時候.
都會有自己的缺點.
自己的短板的時候.
所以.
沒有必要.
所有事情都追求百分百完美.
要追求的是.
我們盡力.
做好每一件事.

$^{801}$做到問心無愧.
其實這樣.
已經很難得.
最後.
這段經文也提醒我們.
你要相信.
上帝對軟弱的我們.
總是有.
恩典念文.
上帝總是給我們機會.
重新做好一件事.
等子們所以.
無論是.
欺騙人.
以色列人.
甚至乎是我們.
這段經文讓我們看到.
每個人總有軟弱.
犯錯的時候.
總有得罪人.
得罪上帝的時候.
但同時我們也看到.
上帝是很肯給我們機會.
來正視.
我們一些過失.
上帝很肯給我們機會.
來改正.
來補救.
因為上帝很願意.
每一個軟弱的人.
回到他身邊.
上帝很願意.
把每一個軟弱的罪人.
留在他身邊.
就好像欺騙人.
所以.
等子們其實來到最後.
你有沒有發現.
上帝.
整件事裡面.

$^{841}$你看看上帝其實.
他是不是特別容易被欺騙.
上帝是不是這麼容易.
被人欺騙.
又或者說上帝是不是這麼.
好欺負被人欺騙了.
也不去反抗.
你知道其實欺騙人.
無論再厲害其實欺騙不到上帝.
為什麼上帝在整件事裡面.
好像沒有出聲.
為什麼上帝.
不第一時間拆穿.
欺騙人的詭計.
為什麼.
上帝不去阻止.
以色列人和他們立約.
其實老實說.
欺騙人的說謊.
肯定是不對的.
以色列人的務實立約有沒有錯.
都有錯.
所以如果要真的追討罪惡.
的話應該怎樣.
每人打五十.
但是為什麼上帝.
沒有追討雙方的罪.
為什麼最後.
上帝要接納.
約書亞的善後方案.
容許.
本來是當滅之物.
的欺騙人.
可以進到.
上帝的殿.
可以親近上帝的祭壇.
這麼接近那位.
聖潔的上帝.
最後一張powerpoint.
其實很值得我們去想.

$^{881}$如果你有留意.
這段經文的結尾.
作者刻意將.
欺騙人的.
由遠離上帝到親近上帝.
的描述.
本來是怎樣的.
23節的上半節.
是被咒詛.
斷不了作路.
23節的下半節.
在上帝的.
恩典和憐憫下.
他變成了.
在上帝的殿.
作劈柴挑水的人.
到27節.
不單止在上帝的殿.
更加在殿裡.
一個很緊要的地方.
上帝的祭壇.
一個.
服侍上帝的地方.
上帝容許.
一群本來該滅的人.
變成了親近上帝.
可以服侍上帝的人.
奇不奇妙.
聽姐妹.
其實上帝就是這樣對待我們.
上帝就是這樣對待.
你和我.
正如保羅所說.
我們本來是可怒之子.
我們本來.
和耶穌基督沒有份.
我們在以色列國民以外.
在所應許的約上.
是局外人.
並且活在世上.

$^{921}$死亡沒有神的境況裡.
就好像.
今日經文裡的「欺騙人」.
但使徒保羅.
亦都說.
可視從前遠離神的人.
如今已經.
得親近了.
並且其實.
你會看到,弟兄姊妹.
直到現在.
我們都會有軟弱.
我們都會有甩鹿的時候.
但上帝仍然堅持.
對我們守約.
思慈愛.
所以弟兄姊妹.
我們要懂得感恩.
感恩的是.
我們仍然有這位.
滿有恩典憐憫.
對我們不離不棄的神.
祂帶領我們.
不只是我們的救主,更是我們生命的主.
在感恩的同時,我們要懂得去回應.
我們要去致力.
去堅守.
和上帝所納的約.
要做一個.
稱職的門徒.
無論在順境,逆境.
有逆境,無逆境.
底下.
無論遇到什麼情況.
我們仍然要心繫.
屬靈的家,心繫.
天國.
就好像今早.
看到一個新聞.
如果大家記得.

$^{961}$311事件.
日本的311事件.
福島的海嘯.
地震.
加上核輻射洩漏.
今年已經是九週年.
固然令到整個福島縣.
重創.
原來也影響到.
當地一個很知名的.
一個行業.
釀酒業.
福島的清酒.
是很出名的.
不過因為311事故.
不單受到波及.
幾乎毀於一旦.
你不難想像.
在重災之後.
要重建家園已經很不容易.
而那些酒商要贏回.
客戶對福島清酒.
這個品牌的信心.
就更加難上加難.
因為你擔心有輻射.
但有一個日本人.
他叫做鈴木大介.
他是福島縣.
原本一個酒廠的老闆.
在311事件後.
他沒有辦法.
要遵從政府的指令.
疏散.
撤離福島縣.
他去到另一邊.
應該是本州那邊.
在山形縣那裡.
買了一間釀酒廠的股份.
繼續釀酒.
但他始終.

$^{1001}$心是他的家園.
心是福島.
所以自從兩年前.
福島漸漸解封之後.
他就經常往返兩地.
即往返山形縣和福島縣.
為什麼要這樣做.
因為他要從福島縣裡面.
去提取那些很好的稻米.
和很好的地下水源.
你知道釀酒這兩樣東西.
是很重要的.
所以他就.
這兩年就不斷來回.
山形縣和福島縣.
將福島縣那些好的稻米和地下水.
帶回山形縣.
釀酒.
簡單嗎.
不簡單,要解決很多問題,很多挑戰.
很多突破,不單止這樣.
他為了要確保.
這個食品的安全.
一點都不馬虎.
福島縣拿回來的地下水.
每小時都要檢測.
那個輻射.
的情況,有沒有被污染.
釀酒的米.
都必須通過比以前.
那個標準更.
嚴格的檢驗.
夠不夠呢,還未夠.
到這些清酒釀好.
要發貨之前.
還要做最後一次的輻射檢測.
來確保.
是安全的.
鈴木他.
將這一批的清酒.

$^{1041}$改了一個很有意思的名字.
叫做.
里程碑.
象徵著.
這一批的酒.
成為了復興家園的里程碑.
並且他計劃.
在明年.
即是311事件的10週年.
他要回到福島縣.
來釀酒.
丁子妙議員.
我想這一件的新聞.
其實提醒我們.
假如當我們.
真的心繫著淑玲的家.
真的心繫著.
天國的時候.
我想無論面對.
再困難的境況.
再多的限制.
你都會怎樣.
你都會回到.
你都會很想去.
親近,服侍上帝.
難阻不到你.
你總是會想方法.
總是會盡你的能力.
來復興這個家.
所以.
確實.
疫情真的.
未完全掌握.
未完全受控.
但我們不能夠忘記一件事.
就是所有的東西.
仍然在上帝的手裡.
仍然在上帝的掌控裡.
福音的工作.
不會因為疫情.

$^{1081}$停止.
對上帝的敬拜.
對上帝的尊崇.
不會因為疫情而改變.
上帝期望我們.
用不同的方法.
將福音.
將上帝的福氣.
繼續延展下去.
在這段時間.
你可以將它看為.
一個困難.
一個逆境.
但同時.
亦都可以是一個復興的里程碑.
因為上帝已經為我們.
預備了.
很多可以服侍他.
服侍別人的機會.
等我們去回應.
等你和我去回應.
願意我們真的放膽.
去求問上帝.
看看上帝.
在未來的日子.
會怎樣用你.
會怎樣突破你的生命.
我們同心圖告.
祈求上帝讓我們.
無論在什麼境況.
我們總是心繫著你.
總是心繫著.
你的國.
總是心繫著.
你給我們的使命.
你所交託給我們的人.
上帝讓我們.
在困難的裡.
讓我們不是單單看到.
看到黑暗.

$^{1121}$讓我們看到.
讓我們不是單單看到.
看到限制.
不是看到軟弱.
讓我們看到你仍然掌管一切.
讓我們看到.
你仍然為我們很多機會.
來服侍你.
親近你.
讓身邊的人帶到你面前.
上帝為此.
讓我們真的.
能夠首先仰望你.
對準你.
學習過一個.
不一樣的生活.
重新調教我們的模式.
我們的步伐.
來貼近上帝的安排.
配合你的安排.
讓我們懂得每天.
先將自己.
將上帝你所託付給我們的事情.
我們交託你.
尋求你的旨意.
讓我們在過往.
我們有軟弱做得不好的地方.
我們相信上帝.
你仍然給我們機會.
去修正 去補救.
以致我們真的可以.
重新起步.
來衝破那些困難.
挑戰.
為你 為自己.
為身邊的人.
締造更多.
美善的圖畫.
讓人在裡面.
尊崇你 認識你.

$^{1161}$讓人在裡面.
經歷到不一樣的福氣.
上帝願你.
因待我們每一個.
或許在當下我們有不同的掣肘.
有不同的難處.
甚至我們可能會有下沉.
發霉的心.
但上帝我們知道.
在你裡面一定有出路.
一定會找到幫助.
願你帶領我們.
願你的靈指引我們.
讓我們真的能夠.
藉著今日經文的提醒.
重新的起步.
靠著上帝你的恩典和憐憫.
我們可以.
跨過.
面前的困難.
繼續在當下.
靠著你站穩 好好的生活.
我們多謝你.
願你聽我們在你面前的禱告.
奉基督耶穌你得勝的名字.
來到祈求 阿們.
\newpage



\section{約書亞記 10:1-15}
\label{sec:q99RT1cz_zM}
\textbf{2020年3月22日 崇拜直播｜約書亞記10\_1-15}
\newline
\newline
連結: \href{https://youtube.com/watch?v=q99RT1cz_zM}{\texttt{https://youtube.com/watch?v=q99RT1cz\_zM}} ~~~~ 語音日期: 2020-03-23
\newline
\newline
\hyperref[sec:ruBv0LtC8pg]{\small{< < < PREV SERMON < < <}}
~
\hyperref[sec:index_chronic]{\small{[返順時目]}}
~
\hyperref[sec:pz6DLtcCEWM]{\small{> > > NEXT SERMON > > >}}
\newline
\newline
約書亞記 10:1-15
\newline
\begin{longtable}{cl}
\hline
\hline
章節 & 經文 (和合本修訂版)\\
\hline
10:1 & \begin{tabularx}{0.7\textwidth}{X} 耶路撒冷王亞多尼‧洗德聽見約書亞奪了艾城，徹底毀滅，處置艾城和艾城的王像處置耶利哥和耶利哥的王一樣，又聽見基遍的居民與以色列人立了和約，住在他們中間， \end{tabularx} \\ \\ \relax
10:2 & \begin{tabularx}{0.7\textwidth}{X} 耶路撒冷人就很懼怕，因為基遍是一座大城，如京城一樣，比艾城更大，並且城內的人都是勇士。 \end{tabularx} \\ \\ \relax
10:3 & \begin{tabularx}{0.7\textwidth}{X} 耶路撒冷王亞多尼‧洗德派人去見希伯崙王何咸、耶末王毗蘭、拉吉王雅非亞和伊磯倫王底璧，說： \end{tabularx} \\ \\ \relax
10:4 & \begin{tabularx}{0.7\textwidth}{X} 「求你們上來幫助我，我們好攻打基遍，因為它與約書亞和以色列人立了和約。」 \end{tabularx} \\ \\ \relax
10:5 & \begin{tabularx}{0.7\textwidth}{X} 於是五個亞摩利王，就是耶路撒冷王、希伯崙王、耶末王、拉吉王和伊磯倫王，聯合上去，率領他們所有的軍隊，對著基遍安營，要攻打基遍。 \end{tabularx} \\ \\ \relax
10:6 & \begin{tabularx}{0.7\textwidth}{X} 基遍人就派人到吉甲的營中約書亞那裡，說：「不要袖手不顧你的僕人，求你趕快上來拯救我們，幫助我們，因為住山區亞摩利人的諸王已經聯合來攻擊我們。」 \end{tabularx} \\ \\ \relax
10:7 & \begin{tabularx}{0.7\textwidth}{X} 於是約書亞和所有跟他一起作戰的士兵，以及大能的勇士，從吉甲上去。 \end{tabularx} \\ \\ \relax
10:8 & \begin{tabularx}{0.7\textwidth}{X} 耶和華對約書亞說：「不要怕他們， 因為我已將他們交在你手裡，他們沒有一人能在你面前站立得住。」 \end{tabularx} \\ \\ \relax
10:9 & \begin{tabularx}{0.7\textwidth}{X} 約書亞就連夜從吉甲上去，猛然襲擊他們。 \end{tabularx} \\ \\ \relax
10:10 & \begin{tabularx}{0.7\textwidth}{X} 耶和華使他們在以色列人面前潰亂。約書亞在基遍大大擊殺他們，在伯‧和崙的上坡路上追趕他們，擊殺他們，直到亞西加和瑪基大。 \end{tabularx} \\ \\ \relax
10:11 & \begin{tabularx}{0.7\textwidth}{X} 他們在以色列人面前逃跑。正在伯‧和崙下坡的時候，耶和華從天上降下大冰雹在他們身上，直降到亞西加，打死他們。被冰雹打死的，比以色列人用刀殺死的還多。 \end{tabularx} \\ \\ \relax
10:12 & \begin{tabularx}{0.7\textwidth}{X} 當耶和華將亞摩利人交給以色列人的那一日，約書亞向耶和華說話，在以色列人眼前說：「太陽啊，停在基遍；月亮啊，停在亞雅崙谷。」 \end{tabularx} \\ \\ \relax
10:13 & \begin{tabularx}{0.7\textwidth}{X} 太陽就停住，月亮就止住，直到國家向敵人報仇。這事豈不是寫在《雅煞珥書》上嗎？太陽停在天空當中，沒有急速落下，約有一整天。 \end{tabularx} \\ \\ \relax
10:14 & \begin{tabularx}{0.7\textwidth}{X} 在這日以前，這日以後，耶和華聽人的聲音，沒有像這日的，這是因為耶和華為以色列作戰。 \end{tabularx} \\ \\ \relax
10:15 & \begin{tabularx}{0.7\textwidth}{X} 約書亞和跟他一起的以色列眾人回到吉甲的營中。 \end{tabularx} \\ \\
[1ex]
\hline
\hline
\end{longtable}
$^{1}$今日的經訓是記載在《約書亞記》.
十章一至十五節.
程序表已經刊載了.
或者可以在我們的屏幕看到.
《約書亞記》十章一至十五節.
耶路撒冷王亞多尼西德.
聽見約書亞奪了艾城.
進行毀滅.
怎樣待耶利哥和耶利哥的王.
也照樣待艾城和艾城的王.
又聽見姬遍的居民.
與以色列人納了和約.
住在他們中間.
就甚懼怕.
因為姬遍是一座大城.
如都城一般比艾城更大.
並且城內的人都是勇士.
所以耶路撒冷王亞多尼西德.
打發人去見希伯倫王可咸.
耶末王彼蘭.
拉吉王瓦非亞.
和伊基倫王底碧雪.
求你們上來幫助我.
我們好功打姬遍.
因為他們與約書亞和以色列人納了和約.
於是五個阿摩尼王.
就是耶路撒冷王.
希伯倫王.
耶末王.
拉吉王.
伊基倫王.
大家聚集.
率領他們的眾軍上去.
對著姬遍安營.
攻打姬遍.
姬遍人就打發往吉格的營中.
去見約書亞.
說:你不要就手不顧你的僕人.
求你速速上來拯救我們.
幫助我們.

$^{41}$因為住在山地阿摩尼人的諸王.
都聚集攻擊我們.
於是約書亞和他一切兵兵.
並帶能的勇士都從吉格上去.
耶和華對約書亞說:不要怕他們.
因為我已將他們交在你手裡.
他們無一人能在你面前站立得住.
約書亞就中夜從吉格上去.
猛然臨到他們那裡.
耶和華使他們在以色列人面前潰亂.
約書亞在姬遍大大地殺敗他們.
追趕他們在伯和倫的上坡路.
擊殺他們直到亞西加和馬基大.
他們在以色列人面前逃跑.
正在伯和倫下坡的時候.
耶和華從天上降下冰雹在他們身上.
直降到亞西加打死他們.
被冰雹打死的比以色列人用刀殺的還多.
當耶和華將阿摩利人交付以色列人的日子.
約書亞就禱告耶和華.
在以色列人眼前說:.
「日頭啊,你要停在姬遍.
月亮啊,你要止在亞倫谷」.
於是日頭停留,月亮止住.
直等國民向敵人報仇.
這事豈不是寫在《瓦薩爾書》上嗎?.
日頭在天當中停住.
不急速下落,若有一日之久.
在這日以前,這日以後.
耶和華聽人的禱告沒有像這日的.
是因耶和華為以色列爭戰.
約書亞和以色列眾人回到吉甲的營中.
弟兄姊妹平安.
平安.
原本這段經文是由梁偉先生宣講.
不過因為他近日腸胃不適.
所以早兩日決定由我負責今天的靜悼.
讓梁先生可以多些休息.
之前珍珍姑娘也為我祈禱.
我開玩笑說.

$^{81}$求上帝像這段經文一樣.
將日頭停住,月亮止住.
讓我可以有多些時間去預備這篇講章.
當然上帝沒有答應這英文.
沒有叫日頭和月亮止住.
但上帝的恩典總是夠用.
也讓我們藉著禱告.
求上帝的恩典繼續臨到我們身上.
我們同心禱告.
上帝我願你醫治梁偉他.
讓他能夠做一個康復.
也帶領我們.
讓我們在宣講聆聽的過程中.
我們可以將榮耀歸給你.
亦可以彼此造就.
我們這樣去禱告.
奉耶穌你的名字來祈求.
其實承接上次我們說到.
羈披人歸降了.
以色列人和他們結盟.
這個消息也傳到另一個部族當中.
就是我們今天一開始看到耶路撒冷王.
再加上他知道近日約書亞.
勢如破竹.
先後滅絕了耶利哥城和艾城.
令到他不能夠不懼怕起來.
好像經文所說.
因為之前很大都還有機片.
英文形容這是一座大的城.
還有裡面有很多勇士.
原本這是一個很好的靠山.
這是一個很好的幫手.
可以幫助耶路撒冷王擋一擋住以色列人.
但是現在沒有了.
羈披人反過來和以色列人聯手.
亦會在將來直接進攻耶路撒冷.
所以在懼怕的裡面.
耶路撒冷王選擇先下手為強.
機片懂得聯盟.
他也懂得.

$^{121}$並且他更厲害.
機片人只和以色列人聯盟.
但是耶路撒冷王決定和四個阿摩尼王結盟.
並且揮軍直上.
來到現在安營在機片前.
準備攻打他們.
這樣做一來可以消滅和機片.
我們說的同外人聯盟的敵人.
同時削弱了以色列人的力量.
另一方面也可以像借鑒一樣.
讓其他的迦南族群看到.
誰敢投靠以色列的時候.
我們會打誰?對付誰?.
看看還有誰敢和約書亞聯盟.
第三方面也可以透過.
這個手段.
嘗試逼機片人背棄和以色列人的盟約.
逼他們轉tie .
改投五王陣營.
不過我們看到第六節.
雖然現在形勢很危急.
但機片人沒有像上一章那樣.
為了保存自己的性命.
行詭詐.
他們沒有這樣做.
他們沒有做牆頭草.
選擇背棄耶和華來保護自己.
我們看到既然機片人已經投靠以色列上帝.
他們就決定一心一意投靠到底.
於是他們就打發人往吉甲去.
去到以色列人的陣營.
希望約書亞可以顧念他們.
速速派兵來營救他們.
千萬不要袖手旁觀.
確實機片人有理由擔心約書亞會不理他們.
因為始終是自己鬼咋在先.
自己騙人在先.
確實應該受到懲罰.
事實上這似乎是一個好機會.
亦被約書亞大可以推搪.

$^{161}$今次機片人受襲.
很明顯是上帝報應在他們身上.
所以我們應該任由上帝.
透過這五個阿摩利王來懲罰和消滅他們.
我們不應該派兵來幫助他們.
約書亞大可以這樣想.
但我們看到約書亞選擇堅守和機片人所立的約.
沒有丟下他們.
並且你留意.
約書亞是否只派一隊人去幫助他們.
不是.
你會看到經文說他是動員所有士兵.
大能的勇士.
總動員前去營救機片.
為什麼這樣?.
因為這不單單是人和人之間的約定.
是雙方都曾經在上帝面前.
以上帝的命來做的約定.
所以即使現在機片人為自己帶來麻煩.
甚至乎帶來危機.
但約書亞仍然選擇堅守.
他在上帝面前所立的約.
來到這裡讓我們看到.
機片人和約書亞一種情況.
就是一個人決定了一心去投靠跟隨上帝.
但卻不等於他之後一帆風順.
他仍然會在生活中遇到風險.
遇到考驗.
但經文亦讓我們看到.
那些專心投靠上帝.
那些謹守他吩咐的人.
上帝又不會叫他們羞愧.
不會叫他們失望.
所以我們最終會看到.
上帝不單單解救了機片的危機.
你會發現上帝更加在危機中使用.
他會推動約書亞走前很多步.
並且一鼓作氣.
將迦南南邊所有城市一一收復過來.
為什麼這樣說?.

$^{201}$我們看第十章的結尾.
剛才我們沒有讀到經文.
40至42節.
這裡說.
「這樣約書亞擊殺全地的人.
就是山地,南地,高原,山坡的人.
和那些地的諸王沒有留下一個.
將凡有氣息的進行滅絕.
正如耶和華以色列的神所吩咐的.
約書亞從迦底斯巴尼亞攻擊到加薩.
又攻擊哥山全地直到機片.
約書亞一時殺敗了這些王.
並奪了他們的地.
因為耶和華以色列的神為以色列爭戰」.
什麼意思呢?.
我們看完以下這幅圖.
我們會更加看到上帝如何帶領和推動約書亞.
在右邊的上方.
看到吉甲和機片嗎?.
看到了嗎?.
其實這是第十章一開始.
約書亞最初處於一個被動的情況.
不是他主動要和這五個阿摩尼王開戰.
他只是因為機片受到攻打.
所以他派兵來幫助機片.
但你會看到.
上帝卻藉機片這個危機.
推動約書亞不斷突破前進.
帶領以色列人.
你跟著機片的箭咀指下去.
帶領他收復所有南邊迦南的城市.
這是約書亞沒有想過的.
他沒有想過要一口氣攻打迦南的南邊.
但上帝很奇妙.
透過逆境和解救機片的危機.
同時也帶領他乘勝追擊.
順勢一口氣擊敗迦南南邊的王.
奪回他們的城市.
一次過完成了上帝的計劃.
另一方面,當我們再看一看.

$^{241}$你會發現這班敵人原本住在高原和山地.
我們說通常的城市在高原和山地.
容易攻破嗎?.
我們說容易守難攻.
守很容易,攻上去是難的.
只要他按兵不動.
緊閉著城門.
就要花更多時間和力氣.
還要逐一攻破.
但現在如何?.
在這個危機中.
上帝反過來成為一個契機.
因為這五個王組成聯軍攻打機片.
但在上帝眼中變成了什麼?.
給了約書亞一個一次過一網打盡的機會.
所以我們看到上帝的作為很奇妙.
他的安排往往超過你和我的想像.
確實當危機擺在面前的時候.
人很自然的反應是怎樣?.
會不會很大膽地走前去和他拼命?.
不會的.
人很多時候的防衛機制就是.
縮後一步,先看好一點.
先保留實力.
但上帝的想法有時候未必是這樣.
上帝會藉著一些環境.
甚至一些逆境.
推我們走前.
或者逼我們不可以不走前.
來開展他的工作.
但上帝不只是叫我們做.
上帝會在沿途給予很多肯定和鼓勵.
提供很多幫助給我們.
使我們由原本很被動的.
變成越來越主動.
使我們原本膽怯.
變得越來越剛強.
越來越積極投入上帝的工作.
來為他爭佔.
直至完成上帝的一些計劃.

$^{281}$正如經文說.
當約書亞決定走出那一步.
他決定放膽去履行.
他和欺騙人在上帝面前所納的約的時候.
奇妙的事情就由那一刻開始.
按著上帝的帶領.
逐一的發生.
當然我們相信.
約書亞決定出兵的時候.
心裡總是會有些忐忑不安.
為什麼這樣說?.
你留意之前耶利哥城和艾城的戰役.
都是怎樣?一對一.
但這次是一個對五個.
以色列民要對著阿摩利的五王.
所以上帝也即時對約書亞.
這個決定作出鼓勵和肯定.
第八節上帝對約書亞說.
不要怕他們.
因為我要將他們交在你手裡.
他們無一人能在你面前站立得住.
上帝立即引證約書亞出兵救欺騙的決定.
是合乎他心意.
並且上帝表明.
他會為以色列人去爭戰.
而約書亞有了上帝派給他的定心丸.
他自然更加可以放膽去做戰略部署.
你會看到這次所用的策略.
與之前兩次的戰役很不同.
還記得耶利哥城的策略是怎樣?.
圍著城13次.
然後呼喊,城牆就倒塌.
艾城呢?.
艾城就不同.
艾城是以退為進的方法.
再加上復兵去打贏.
但今次第九至十一節.
你會看到約書亞用了另一個策略.
就是兵貴神速.
約書亞今次選擇漏夜由吉格起行.

$^{321}$趕去欺騙那裡.
他要漏夜突襲.
殺敵軍一個措手不及.
而你又看到上帝幫助約書亞.
令敵人在以色列人面前.
經文說,賄亂.
給約書亞他們大大的失敗.
而當敵軍落荒逃跑的時候.
上帝又再一次幫助他們.
透過一個神蹟降下兵搏.
事實上這個確實是神蹟.
因為如果純粹說是降兵搏.
有可能是一個什麼?.
這麼巧合,這麼巧合一個自然現象.
但是你會發覺兵搏只降在哪裡?.
只降在敵人身上.
兵搏有沒有眼?.
兵搏也沒有導航.
而只是準確地打落在敵人身上.
所以約書亞一看就知道.
這不是偶然.
是上帝正正在替以色列人爭戰.
他很明白這件事.
結果就像經文所說.
被上帝用兵搏打死的敵人.
比以色列人用刀殺死的還要多.
經文提醒我們.
上帝確實在當中為以色列人爭戰.
使五個王的聯軍都要兵敗如山倒.
毫無抵抗的資歷.
等節目再深入看七至十一節.
麻煩音響室.
繼續看經文.
你會發現作者刻意地.
做出交錯的記錄.
紅色的字是約書亞.
藍色的字是上帝.
大家看到他們的互動嗎?.
當約書亞決定派兵去營救基地時.
上帝就立即肯定他.

$^{361}$當約書亞來到前面時.
他就出現在敵軍前面.
又說要使敵軍賄亂.
當約書亞在機片大大殺敵.
令敵方落荒逃跑時.
上帝協助他.
用兵搏幫助他們殺死敵人.
等節目你會看到這段記載.
顯示了上帝和人的互動.
同工的緊密關係.
約書亞很明白.
即使上帝一開始.
已經說明了這場仗一定會怎樣.
一定會打贏.
一定會贏定.
但並不代表約書亞可以翹起雙手.
不代表他可以坐著甚麼都不用做.
等上帝處理.
約書亞卻要在一個爭戰的過程中.
學習如何與上帝建立關係.
他要在一個很靠近的情況中.
認識到上帝是怎樣的神.
藉著一個不容易的環境.
他一路行軍打仗.
一路要指揮部署.
在一個緊張危急的氣氛中.
他要操練自己.
很敏銳上帝如何參與.
他要操練自己很敏銳.
上帝如何帶著他打這場仗.
亦在裡面累積更多的經驗.
如何配合上帝的計劃.
丁子梅你看到.
約書亞就是這樣.
越來越敏銳上帝的帶領.
有一個故事是這樣說.
有一個台灣原住民.
他長年住在台東的山地上.
住在高山上.
有一天他來到台北找他的朋友.

$^{401}$當他和這位朋友在馬路上行走時.
正準備要過馬路時.
這位原住民突然跟他的朋友說.
你聽聽是色恤的叫聲.
這位朋友很詫異.
我只聽到汽車聲.
我聽不到色恤聲.
這位原住民繼續說.
真的有色恤叫.
我很清楚聽到他們在唱歌.
這位朋友開始笑.
你別玩了.
你看看.
整條路都只聽到車和人.
這麼吵怎麼可能聽到色恤的叫聲.
於是這位原住民的朋友.
二話不說.
立即跑到馬路中央的安全島.
掀開旁邊一塊草叢.
叫他的朋友來看.
你看看這裡不是有色恤嗎.
這位朋友大驚.
他問為什麼你的聽力.
可以這麼敏銳這麼厲害.
在這麼吵的環境.
還可以聽到色恤的叫聲.
原住民回答.
你也可以.
其實我們每個人都可以.
朋友不相信.
他就說哪有可能.
原住民於是就說.
不如我們來做一個實驗.
隨即他從他的褲袋裡.
拿了幾個硬幣出來.
然後向天一拋.
自然就掉到地上.
附近的途人被這些硬幣聲吸引到.
就立即調轉頭幫他撿回硬幣.
於是原住民就說.

$^{441}$你看看.
其實大家對聲音的敏銳程度.
差不多.
只不過.
有些人對錢的聲音.
比較敏銳.
有些人就比較敏銳.
大自然的聲音.
確實了弟兄姊妹.
我們有沒有習慣.
和上帝互動交往.
有沒有習慣敏銳.
他的帶領和他的作為.
正正是這段經文要提醒我們.
你會看到上帝其實很希望.
我們不單單只是運用.
才幹智慧來服侍他.
不只是一個做字.
上帝也期望我們很敏銳.
在整個事奉的過程裡.
他是怎樣為我們增賤.
以致我們可以更加近距離.
認識上帝到底是怎樣的一位上帝.
以致我們更加懂得怎樣配合.
上帝的計劃.
隨時緊貼著他一些部分.
就好像經文裡所要我們注意的一樣.
就好像約書亞.
他在一個不簡單的場景裡.
他一路事奉.
一路學習怎樣和上帝互動.
這亦成為我們的寫照.
最後我們看餘下的經文.
十二到十四節.
這裡說.
原來當上帝.
應許要將阿摩利人交付在.
以色列人手中的日子的時候.
約書亞就在以色列人的眼前.
吩咐日頭和月亮要停留止住.

$^{481}$吩咐月亮和日頭來協助他.
攻打阿摩利人.
而上帝也應允了.
約書亞這個禱告.
確實好像經文所說.
讓日頭在天當中停住.
不急速下落.
若有一日之久.
直到以色列人完全殺敗敵人為止.
等姐妹記載這件事件.
道不道德?.
很道德吧?.
不過重點不是神蹟本身.
因為上帝要行一個這樣的神蹟.
有沒有困難?.
一點困難都沒有.
日月都要伏於上帝的掌管裡.
所以這件事的重點.
是沒有人會想到.
上帝竟然會答應約書亞.
一個這樣的祈禱.
竟然為約書亞.
願意改變日月行常運作的法則.
這是一個重點.
而不是神蹟怎麼可能發生.
怎樣發生.
你會留意十四節.
作者怕我們誤會了.
他補充了一句.
在這日之前.
在這日之後.
耶和華聽人的禱告沒有像這日.
是因為耶和華為以色列爭戰.
作者補充了一句.
這個上帝這樣聽約書亞的祈禱.
是什麼?.
限定的.
是這一日,是這一次.
明天下次約書亞再做類似的祈禱.
上帝會不會應運?.

$^{521}$基本上都未必會.
你和我做這樣的祈禱.
應該都未必會.
所以弟兄姊妹.
經文要我們注意的.
不是神蹟本身.
亦不是我們怎樣可以祈求上帝.
行一個這樣的神蹟.
經文要我們注意的就是.
上帝和約書亞的關係.
你細心想.
其實無論約書亞他做不做這個祈禱.
他最終會不會打贏阿摩利人?.
會不會?.
會,為什麼?.
因為上帝早已答應了.
上帝說我會為你爭戰.
我會搞定阿摩利人.
約書亞不做這個祈禱.
都一樣會贏,一樣可以殺敗敵人.
但是你會見到.
上帝卻又很喜悅.
人在事奉的過程裡.
放膽為上帝的事情.
來做一個祈求.
我相信約書亞當時.
在做這個祈禱之前.
他會不會用科學分析?.
他會不會分析天文.
到底日頭月亮會不會止住?.
然後確認了.
上帝應該會聽我這個祈禱.
叫日月停止.
他都有把握的.
他才會在眾人面前.
做這個祈禱.
大概不會.
因為真的有機會.
上帝不答應.
到時如何交代?.

$^{561}$在眾人面前.
很難下台.
約書亞這個祈禱.
我覺得就好像.
一個爸爸和他的小孩子.
一個很純真的互動關係.
什麼意思呢?.
就好像小孩子對爸爸說.
他說:爸爸.
天快黑了.
但是你交給我做的工作.
還未做完.
不如你嘗試叫太陽伯伯.
遲些下班.
你又嘗試叫月亮姨姨.
遲些開工.
這樣我就有充足的時間.
為你完成工作.
弟兄姊妹.
聽到這個.
這麼純真.
這麼忽發奇想的想法.
你做爸爸的.
你會怎樣回應?.
你會不會跟他說.
很理性地說.
別傻了.
快做吧.
還說.
你會跟他理性分析嗎?.
還是你會把握這個機會.
跟他建立一個更親密的關係?.
你會不會跟他說.
兒子,你這個主意.
你這個想法很有趣.
很有創意.
但更重要的是.
我看到兒子你很有心.
你很愛我.
是因為你愛我.

$^{601}$你重視我交給你做的工作.
所以你才會有這樣的忽發奇想.
好了.
我們一起來求天父.
我們就嘗試叫太陽伯伯.
遲些下班.
叫月亮嬸嬸遲些下班.
丁姐妹你會發覺這個祈禱.
最後無論成不成功.
不知道.
但亦不重要.
因為重點是透過這個祈禱.
建立了一段更親密的父子互動關係.
是嗎?.
同樣正如約書亞.
一路殺敵的時候.
發覺天快黑了.
但敵人還未消滅.
於是很單純地跟上帝說.
上帝,可否吩咐日頭月亮嬸嬸.
讓我可以有更多時間為你爭戰.
讓我完成你交給我的工作.
這個祈禱確實很大膽.
確實很不合常理.
甚至乎有些瘋狂.
大概你和我都不夠膽去做這樣的祈禱.
但很奇妙.
上帝看中了約書亞這個單純的赤子心.
上帝看中了他這份熱誠.
對上帝的愛.
並且有一個很重要的大原則.
經文也說了.
是因為這個想法,這個祈求.
符合了上帝為以色列人爭戰的原則.
所以上帝說.
好,我今次就成全你這個想法.
應你這個祈求.
並且上帝要藉著約書亞這個祈求.
進一步在以色列人面前.
確立他領袖的身份.

$^{641}$讓他之後可以更加有信心.
帶領著百姓繼續前行為上帝爭戰.
《上帝英姊妹》這段經文讓我們看到.
上帝是正正期望我們.
為了他的國度的緣故.
為了福音的緣故.
我們真的不妨可以放膽.
向上帝祈求更大的事情.
為上帝做更多的事情.
而這句話.
亦是出自近代的宣教之父.
克里威廉所說的.
當年克里威廉.
他首次在同宮會議裡面.
他提出了要向海外的彝族.
宣教這個儀仗.
不過卻是遭到當時的主席.
大力的反對.
就這樣說.
年輕人,如果上帝要拯救異教徒.
上帝會自己出手.
上帝會自己做事.
上帝不需要你和我的幫助.
克里威廉並沒有因為這次的反對.
就被打沉.
他仍然保持著.
那份宣教的熱誠.
那個夢想.
到了1792年.
他就和十三位當地的牧師.
建立了第一間宣教差會.
為了教會將福音遍傳.
這個歷史揭開了新的一頁.
而向上帝求大事.
為上帝成大事.
亦都成為了克里威廉的座右銘.
同樣地,兄弟姐妹.
盼望這段經文.
盼望我們藉著.
約書亞和上帝的互動.

$^{681}$這段更緊密的關係.
給予我們一個力量.
給予我們一個方向.
感動我們在這個時代.
特別是在困難的時刻.
不妨去放膽.
尋求上帝.
不妨向上帝求更大的事.
求上帝用我們做更多的事.
好嗎?.
讓我們同心禱告.
回應上帝對我們的提醒.
我們同心禱告.
父上,我求你.
藉著這段經文.
開闊我們的視野.
拓展我們的眼界.
讓我們看得見你的期望.
看得到你喜悅我們.
怎樣去回應你.
特別是在艱難的日子.
讓我們知道上帝你沒有離開.
你仍然要幫助我們.
來成就你福音的使命.
所以上帝為此求你.
加添我們對你的信心.
幫助我們在生活裡面.
我們敏銳.
上帝你真的與我們互動同工.
帶領我們.
知道怎樣配合你的計劃.
來緊隨你的心意.
亦幫助我們重新調教.
生活的優先次序.
不要因為沒有時間.
以此為藉口.
忘記與上帝相交.
在忙碌的日子.
我們更應該愉快地.
花時間與上帝你同工.

$^{721}$上帝讓我們在做每一個決定之前.
在我們做每一件事之前.
求你讓我們有智慧.
讓我們懂得分配時間.
讓我們優先回應你所喜悅的事情.
並且讓我們放膽.
為上帝你求更大的事情.
為上帝你做更多的事情.
以致我們能夠將你自己的福氣.
在困難的裡面一路播送開去.
讓更多的人在逆境的裡面.
經歷到你的帶領.
經歷到你的保守.
我們這樣的禱告,教託.
奉耶穌你的名字來祈求.
阿們.
(影片完結).
\newpage



\section{約書亞記 11:1-15}
\label{sec:pz6DLtcCEWM}
\textbf{2020年3月29日 崇拜直播｜約書亞記11\_1-15}
\newline
\newline
連結: \href{https://youtube.com/watch?v=pz6DLtcCEWM}{\texttt{https://youtube.com/watch?v=pz6DLtcCEWM}} ~~~~ 語音日期: 2020-03-30
\newline
\newline
\hyperref[sec:q99RT1cz_zM]{\small{< < < PREV SERMON < < <}}
~
\hyperref[sec:index_chronic]{\small{[返順時目]}}
~
\hyperref[sec:1PuKX5mdEMQ]{\small{> > > NEXT SERMON > > >}}
\newline
\newline
約書亞記 11:1-15
\newline
\begin{longtable}{cl}
\hline
\hline
章節 & 經文 (和合本修訂版)\\
\hline
11:1 & \begin{tabularx}{0.7\textwidth}{X} 夏瑣王耶賓聽見了，就派人到瑪頓王約巴、伸崙王、押煞王， \end{tabularx} \\ \\ \relax
11:2 & \begin{tabularx}{0.7\textwidth}{X} 和北方山區、基尼烈南邊的亞拉巴、低地、西邊多珥山岡的諸王， \end{tabularx} \\ \\ \relax
11:3 & \begin{tabularx}{0.7\textwidth}{X} 以及東方和西方的迦南人、山區的亞摩利人、赫人、比利洗人、耶布斯人，和黑門山下米斯巴地的希未人那裡。 \end{tabularx} \\ \\ \relax
11:4 & \begin{tabularx}{0.7\textwidth}{X} 他們和他們的眾軍都出來，一大隊人馬，多如海邊的沙，並有極多的戰車戰馬。 \end{tabularx} \\ \\ \relax
11:5 & \begin{tabularx}{0.7\textwidth}{X} 眾王組成聯軍，來到米倫水邊一同安營，要與以色列作戰。 \end{tabularx} \\ \\ \relax
11:6 & \begin{tabularx}{0.7\textwidth}{X} 耶和華對約書亞說：「你不要怕他們。明日這時，我必把他們全部交給以色列人殺滅。你要砍斷他們馬的蹄筋，用火焚燒他們的戰車。」 \end{tabularx} \\ \\ \relax
11:7 & \begin{tabularx}{0.7\textwidth}{X} 於是約書亞和所有跟他一起作戰的士兵，來到米倫水邊，突然攻擊他們。 \end{tabularx} \\ \\ \relax
11:8 & \begin{tabularx}{0.7\textwidth}{X} 耶和華將他們交在以色列人手裡，以色列人就擊殺他們，追趕他們到西頓大城，到米斯利弗‧瑪音，直到東邊米斯巴的山谷。以色列人擊殺他們，沒有留下一個倖存者。 \end{tabularx} \\ \\ \relax
11:9 & \begin{tabularx}{0.7\textwidth}{X} 約書亞照耶和華所吩咐他的去做，砍斷他們馬的蹄筋，用火焚燒他們的戰車。 \end{tabularx} \\ \\ \relax
11:10 & \begin{tabularx}{0.7\textwidth}{X} 那時，約書亞轉回，奪了夏瑣，用刀殺了夏瑣王。先前夏瑣在這些王國中是為首的。 \end{tabularx} \\ \\ \relax
11:11 & \begin{tabularx}{0.7\textwidth}{X} 以色列人用刀擊殺城中所有的人，把他們完全滅盡；凡有氣息的，沒有留下一個。約書亞又用火焚燒夏瑣。 \end{tabularx} \\ \\ \relax
11:12 & \begin{tabularx}{0.7\textwidth}{X} 約書亞奪了這些王的一切城鎮，擒獲了這些王，用刀殺了他們，把他們完全滅盡，正如耶和華的僕人摩西所吩咐的。 \end{tabularx} \\ \\ \relax
11:13 & \begin{tabularx}{0.7\textwidth}{X} 至於造在山岡上的城鎮，除了夏瑣以外，以色列人都沒有焚燒。約書亞只焚燒了夏瑣。 \end{tabularx} \\ \\ \relax
11:14 & \begin{tabularx}{0.7\textwidth}{X} 從那些城鎮所奪的財物和牲畜，以色列人都取為自己的掠物。至於所有的人，他們都用刀殺了，直到滅盡；凡有氣息的，沒有留下一個。 \end{tabularx} \\ \\ \relax
11:15 & \begin{tabularx}{0.7\textwidth}{X} 耶和華怎樣吩咐他的僕人摩西，摩西就這樣吩咐約書亞，約書亞也照樣做了。凡耶和華所吩咐摩西的，約書亞沒有一件偏離不做的。 \end{tabularx} \\ \\
[1ex]
\hline
\hline
\end{longtable}
$^{1}$今日我們要重讀的經文.
是記載在《約書亞記》第十一章一至十五節.
「下所王耶彬聽見這事.
就打發人去見馬頓王,約巴,申崑王,阿薩王.
與北方山地基尼列南邊的阿拉巴高原.
並西邊多爾山江的諸王.
又去見東方和西方的迦南人.
與山地的阿摩利人,赫人,比利駛人,耶布斯人.
並黑門山跟米斯巴地的希美人.
這些王和他們的重軍都出來.
人數多如海邊的沙.
並有許多馬匹車輛.
這諸王會合來到米倫水邊.
一同安營.
要與以色列人爭戰」.
第六節.
「耶和華對約書亞說.
你不要因他們懼怕.
明日這時.
我必將他們交付以色列人傳言殺了.
你要砍斷他們馬的蹄筋.
用火焚燒他們的車輛.
於是約書亞率領一切兵力.
在米倫水邊突然向前攻打他們.
耶和華將他們交在以色列人手裡.
以色列人就擊殺他們.
追趕他們到西頓大城.
到米斯利弗馬欽.
直到東邊米斯巴的平原.
將他們擊殺.
沒有留下一個.
約書亞就照耶和華所吩咐他的去行.
砍斷他們馬的蹄筋.
用火焚燒他們的車輛」.
第十節.
「當時.
約書亞轉回奪了下所.
用刀擊殺下所王.
數來下所在這諸國中是為首的.
以色列人用刀擊殺城中的人口.

$^{41}$將他們進行殺滅.
凡有氣息的沒有留下一個.
約書亞又用火焚燒下所.
約書亞奪了這些王的一切成衣.
擒獲其中的諸王.
用刀擊殺他們.
將他們進行殺滅.
正如耶和華僕人摩西所吩咐的.
至於造在山崗上的城.
除了下所以外.
以色列人都沒有焚燒.
約書亞只將下所焚燒了.
那些成衣所有的財物和生畜.
以色列人都取為自己的略物.
唯有一切人口都用刀擊殺.
直到殺盡.
凡有氣息的沒有留下一個.
耶和華怎樣吩咐他僕人摩西.
摩西就照樣吩咐約書亞.
約書亞也照樣行.
凡耶和華所吩咐摩西的.
約書亞沒有一件懈抬不行的.
各位弟兄姊妹.
主女們平安.
今日望著這個禮堂.
令我聯想到.
在1月26日年初二那天.
聯合崇拜的時候.
我們這個禮堂是坐滿了.
我們還一同去聽梁牧師.
為我們不同的新肖的人.
去預測一下我們今年的運程.
眨眼兩個月了.
今日旺朕.
因為疫情嚴峻的緣故.
就決定取消實體崇拜.
又因為星期五政府.
出了一些星期六刊憲.
禁止四個人以上在公眾集會的條例.
於是我們就變陣.

$^{81}$即刻將我們的崇拜事奉人員減到最低.
今日在這個禮堂裡面我們只有三個人.
在這個日子我越來越感受到.
我們今年的主題.
邁向新世代.
真是越來越迫切.
去回應這個新世代的轉變.
現在全世界都在疫情之中.
新冠狀病毒肺炎.
由武漢,湖北.
傳遍中國.
散播到亞洲.
再去歐洲,美洲.
回過頭再回來我們這裡.
早前曾經出現過.
口罩,廁紙,米糧,抗疫物資的爭奪戰.
現在又再捲土重來.
有時真是想.
這場戰役真是很漫長.
我們估不到.
究竟何時才會完.
何時才會是盡頭呢?.
其實我們每一個人的人生.
都面對大大小小.
有形無形的爭戰.
而事實上對於信徒來說.
是否遵從神的命令和吩咐.
就是我們生命中成敗得失的關鍵.
曾經聽過我一個傳道的好朋友.
她認識一個姐妹.
她和她丈夫自小青梅竹馬.
於是在年輕的時候就結婚.
後來這位姐妹經過幾年事業的搏殺之後.
在事業上得到一定的成績.
然後她就慢慢被這個花花世界所吸引.
她開始覺得她的丈夫和她距離越來越遠.
越來越覺得這個丈夫配不起她.
但反而在她的工作上.
她就認識了一個很投契的弟兄.
教會知道了這兩夫婦面前的危機.

$^{121}$於是就有導師去找她.
關心她們.
並且用聖經的真理去勸導她們.
鼓勵這個姐妹.
要和那位在工作認識的弟兄保持距離.
但姐妹沒有聽這個勸告.
終於就和這個弟兄發展了婚外情.
並且她決意要和她的丈夫離婚.
結果這個姐妹就在她的工作上.
和弟兄分別離開了教會.
姐妹和父母也因此而決裂.
她的前夫也因此而患上抑鬱.
從前很有目標,很有盼望的生命.
變得一塌糊塗.
姐妹因為不能面對令她滿足不到的婚姻關係.
抵禦不了這個世界的引誘.
而不聽神的吩咐和命令.
結果人生變得混亂.
並且令到身邊的人受到傷害.
同樣我們每一個人.
都會面對人生的不同挑戰.
但怎樣我們才可以在這個持續不斷的爭戰裡.
過一個得勝的人生呢?.
今天我們會從約書亞記.
去看看約書亞怎樣過一個持續得勝的人生.
讓我們一同低頭祈禱.
親愛的主,求你的恩光照亮我們.
開啟我們的心靈.
讓我們能夠聆聽你對我們說的話.
讓我們心意更新.
讓我們貼近你的心意.
讓我們行在你的恩光之中.
禱告教託,奉主耶穌基督得勝明智祈求,阿們.
我們今天要分享的經文有一個背景.
就是一個持續不斷的爭戰.
我們知道約書亞帶領著以色列人.
從約旦河東過河,去到約旦河西.
在這片土地上.
約書亞帶領著以色列人.
打了三場很重要的仗.

$^{161}$都是在迦南地的南部.
相信你們仍然記憶猶新.
耶利哥,艾成到最後出兵營救機遍.
每一場都非常精彩.
每一場都引人入勝.
隨後第十章二十八節起.
約書亞帶領著以色列人在迦南南部.
攻佔一座又一座的城.
可以說是征戰連場.
最後到第四十三節.
聖經描述以色列人回到吉格安營.
當南方一個一個的城被以色列人攻佔後.
北方的各民各族自然會想到.
他們會受到威脅而有所行動.
今日我們所說的經文就是約書亞記第十一章一至十五節.
就是說以色列人和迦南北方的聯軍如何征戰.
我們在地圖上看到北方的所在.
我們細心看經文.
第一節「夏爽王耶賓聽見這事」.
這是一個非常關鍵的人物.
這個關鍵的人物在第十節說到.
他是這場戰爭的牽頭人.
是北方聯軍的首領.
而我們聽到夏爽王耶賓.
「耶賓」並不是他的名字.
原來就像埃及王法老一樣.
我們一說「法老」就知道他代表了法老王.
代表了埃及王.
是埃及的皇帝.
如果我們一說「耶賓」.
我們就知道這個是夏爽王.
代表了夏爽的皇帝.
「夏爽王耶賓聽見」.
「聽見了」.
「聽見」這個動詞.
其實分別在五章,九章,十章,十一章的開首.
我們都見過.
在這個部分我們見到紅色的字.
就是「聽見」.
在五章至九章.

$^{201}$我們見到藍色的字.
就是哪些人聽到呢?.
就是迦南地的人.
不同的國.
不同的民族.
他們聽見.
然後綠色的字.
就是他們聽見後有甚麼反應呢?.
我們看完五章,九章.
我們就看看十章和十一章.
同樣的.
都是藍色字代表了哪些人.
紅色字代表了聽見.
綠色字代表了他們有甚麼反應.
你們可以從這些圖表看到.
迦南地各個王,各族,各民.
對這支從約旦河東西廠的以色列人入侵的反應.
我們從他們所知道的.
往往他們聽見的內容.
都是關乎以色列民的事.
夏爽王聽見的.
必然就是關於南部的戰役.
包括了在姬遍的戰役.
神如何從天降下冰雹.
又如何叫日頭止住一整天.
他又會聽見南方不同的城邑.
如何大敗.
被侵略,侵佔.
姬遍的阿摩利王.
如何被殺.
相對於其他王.
這裡沒有形容夏爽的懼怕.
但我相信他應該從這些消息.
他知道他受到極大的威脅.
他知道強敵當前.
他知道這一隊是不能輕視的軍隊.
我們可以從他接下來的行動.
可見一斑.
第九章一至二節提到一件事.
就是一旦河西駐山區.

$^{241}$低地和沿大海一帶.
直到黎巴嫩的諸王.
原來本來已經聚集在一起.
同心合意地要和約書亞.
和以色列人爭戰.
即是說迦南地無論南和北.
本來在第九章的時候.
他們已經聽到以色列人.
在耶利哥艾城作為的時候.
他們已經決定要爭戰.
但這個原先的計劃.
想和以色列人拼命的計劃.
最終因為南方的大城姬遍王.
不願意和南方各國結盟.
反而使用詭計.
向以色列人騙得免死的保障.
而落空了.
結果迦南諸王的如月算盤打不響.
南方的盟軍反而攜不成軍.
而夏所這個地方.
從地理環境裡是極有優勢的地方.
從考古學的發現.
我們見到夏所是一個很大的城都.
位於基里列海附近.
基里列海是甚麼呢?.
其實原來是後來的加利利海.
而這裡是一個主要的貿易路線經過的地方.
是一個繁華盛世舉足輕重的地方.
而夏所王在考古文獻記載.
原來迦南的那麼多皇帝之中.
只有唯一一個夏所王.
與法老的通信是自以王來相稱.
他自稱為王.
其他迦南的所有王都沒有人這樣做.
其實夏所已經是一個實力不小的城市.
但聽到以色列人過河之後.
戰役節節勝利.
大敵當前.
他不像欺騙人那樣騙以色列人去保命.
反而他主動出擊.

$^{281}$聯絡南面那邊重要的三個王.
就是馬頓王,新倫王和阿薩王.
並且我們在第一,二,三節.
我們知道他從北方的南到北.
東到西的各城.
聚集了一隊很大的軍隊.
以夏所為首.
集結在米倫的水邊安營.
準備與以色列人開戰.
我們從第四節看到.
「這些王和他們的重軍都出來.
人數多如海邊的沙.
並有許多馬匹車輛」.
北方的聯軍在這裡擁有了兩個強大的力量.
第一就是人數多如海邊的沙.
我們從第一,二節.
說到各個不同邦國的王和居民.
但第三節其實是形容同一班人.
只是表達他們來自不同的民族.
但也足以令他們的軍隊如海邊的沙般多.
還有他第二樣強大的武器.
就是有許多馬匹車輛.
馬匹和車輛在當時來說.
是非常厲害的武器.
是行軍極為重要的裝備.
我們看到這場戰爭.
關乎兩個敵對陣營.
夏所王知道以色列人來勢洶洶.
他作出了一個他認為可以致勝的部署.
另一邊廂.
以色列人又何嘗不是要面對一個非常強大的敵人呢?.
他們所面對的軍隊.
是他們前所未有的那麼多.
他們遇見敵方的軍備.
更加史無前例地先進和厲害.
規模比起他之前征戰年長的戰役.
實在是蚊比和牛比.
以色列人面對的是夏所為首.
戰士如海沙般多.
擁有厲害武器的北方聯軍.

$^{321}$是極其強大的敵人.
弟兄姊妹.
不知道對於你來說.
甚麼是你的強大敵人呢?.
甚麼是你的強敵呢?.
在你人生裡面.
甚麼會叫你如臨大敵一樣呢?.
是患了重病?.
是婚姻失敗?.
是患了新冠狀病毒?.
或是擔心面前會被裁員?.
破產?.
或親人離世?.
又或是你的強敵是.
在試煉中跌倒?.
到底當我們面對人生強敵的時候.
我們可以有甚麼出路呢?.
讓我們一起來看看第六至第九節吧.
這一場是一場戰役.
米崙水邊之戰.
究竟這場戰是怎樣打的呢?.
讓我們一起去看看PowerPoint吧.
這裡就是夏娑和北部聯軍的戰役.
首先我們會看到紫色的箭咀告訴我們.
迦南北部的聯軍從四面八方來到米崙的水邊.
然後以色列人在哪裡呢?.
雖然上一章最後一節告訴我們.
以色列人回到吉格安營.
但由於在經文裡.
我們知道神和約書亞保證了明日這時.
他們已經取得勝利.
在一天之間他們已經取得勝利.
但吉格距離米崙足足有170公里那麼遠.
所以我們估計第六節所形容的以色列人.
應該已經去到米崙水邊的附近.
所以紫色的箭咀就是以色列人行軍去到米崙水邊的附近.
作為開始.
已經在北方盟軍的安營之地.
大家會留意到一件事.
比起之前的戰役.

$^{361}$很詳細地去記錄每一個戰役的經過.
但今次只是用了三節來形容這場戰役.
第七節告訴我們.
約書亞帶領著軍隊作出突擊.
第八節.
以色列人就擊殺他們.
我們留意紫色的虛線箭咀.
以色列人追趕北方聯軍到西頓大城.
到米斯利弗馬欽.
直到東邊米斯巴平原.
將他們擊殺.
一個都不留.
然後第九節告訴我們.
當他們這班盟軍潰敗的時候.
他們就砍斷他們馬的蹄筋.
用火焚燒他們的車輛.
米崙水邊之戰的開始.
經過和結束.
只是用了短短三句去表達.
似乎略為平淡了一點.
我在預備講章的時候.
我在想有什麼好說的.
但當我仔細細味的時候.
我就會發覺.
他們強敵當前.
而並擁有最先進的武器.
人多到打你不死.
都好像蟻一樣,氣死你.
還要對付你們這班征戰連場.
打到沒有停手.
甚至是疲累不堪的軍隊.
其實論排面.
北方聯軍是沒得輸的.
但是第六節第九節.
聖經的作者經過細心的鋪排.
我們就會看到這場戰爭.
以色列人是怎樣打.
經文的第六節告訴我們.
神明白到以色列人所面對的征戰.
是多麼的可怕.

$^{401}$祂知道約書亞面對強敵的擔心和憂慮.
祂再次鼓勵他不要怕.
我們看到PowerPoint裡面.
約書亞記不同的地方.
例如第一章第九節.
當約書亞要接續摩西領袖岡維.
帶領以色列人攻打迦南地的時候.
神叫他不要懼怕.
第八章第一節.
當他預備要攻打艾城的時候.
又或者第十章第八節.
他預備去迎戰南方的聯軍的時候.
上帝都是堅定不斷的安慰,鼓勵.
賦予他信心.
而且一而再再而深地向他作出應許.
而來到這裡第十一章.
他所作出的應許.
就是明天這個時候.
我必將他們交予以色列人.
傳言殺了.
當人面對強敵.
憂慮擔心是人之常情.
但是我們看到這段經文.
我們看到神對我們的愛.
是何等堅定不移.
他一而再再而深地去安慰,鼓勵.
加我們力量.
並且他給我們應許.
是他最完美的計劃.
大家要留意.
這個應許是很特別.
上帝對以色列人的應許.
和訓練是非常非常到位.
大家還記不記得第十章.
上星期Raymond跟我們分享.
耶穌對付南方聯軍的時候.
南方聯軍的規模比較小.
敵軍沒有那麼先進的武器.
但是你記不記得那場戰役.
除了以色列人親手用刀殺的之外.

$^{441}$神還用大自然的威力.
用冰雹幫他殺死他的敵人.
但神這次的應許是甚麼?.
他要將所有的敵人.
全交給以色列人殺滅.
他在應許裡面.
沒有說在他的手上.
而是以色列人全面殺滅.
這次沒有答應有冰雹.
怎能殺光海沙那麼多的軍隊?.
第二件事.
上帝說明日這時.
明日這個時候.
大家記不記得.
「機片之戰」.
約書亞為了要殺敵.
他向上帝禱告甚麼?.
他向上帝禱告.
太陽,月亮要停住.
要給多些時間去殺敵.
這次他面對的北方聯軍.
比「機片」的人數更加多.
神應許他只需要一天的時間.
就可以打敗如海沙那麼多的軍隊?.
你信不信?.
這個應許是一件看來沒可能的事.
但弟兄姊妹你看看約書亞.
他是毫不猶豫.
因為他相信神的應許.
一定不會落空.
他反而會看.
這是上帝給他的指引.
第七節告訴我們.
於是他和以色列軍二話不說.
到米倫水邊的大軍進行突襲.
他要攻其不備.
有學者告訴我們.
這是神很巧妙的指引和安排.
因為米倫水邊位於山區.
北方聯軍只用這個地方作為集合地點.

$^{481}$集合後.
就會用有利的戰車和戰馬.
到平原開戰.
耶和華的應許是一個清晰的指示.
引導以色列人速戰速決.
為的是要削弱敵軍.
恃有戰車和戰馬的優勢.
結果以色列人將他們殺個措手不及.
而敵方軍隊樂方而逃.
第十一節.
第九節說到.
約書亞也完全遵從神的要求.
砍斷北方聯軍馬的蹄筋.
用火焚燒他們的戰車.
砍斷馬的蹄筋.
目的是令他們的屈機腱受傷.
使他們不能再在戰爭中發揮功效.
耶和華這樣做有什麼目的?.
耶和華這樣做的目的是要摧毀我們以為可以依賴的東西.
祂提醒我們.
我們的征戰究竟需要誰?.
祂要我們認清.
耶和華自己才是我們力量的來源.
強如夏所帶領的北方聯軍.
有龐大的軍隊.
有戰車,戰馬.
耶和華都將他們交給祂所愛的子民.
願意依靠祂的子民的手中.
所以第七節提醒我們很對.
有人靠車,有人靠馬.
但我們要提到耶和華,我們神的命.
弟兄姊妹.
在你的人生旅途中.
你怕什麼?.
你怕失去?.
你怕孤單?.
你怕死?.
你怕窮?.
你怕痛苦?.
你怕不怕上帝?.

$^{521}$其實我們的人生正正就是要怕上帝.
因為你怕什麼呢?.
你就會成為那件事的奴隸.
如果你敬畏上帝.
就成為上帝的奴隸.
成為上帝的僕人.
那你就可以過得勝的人生.
弟兄姊妹.
我們可以選擇.
我挑戰你們.
選擇,怕誰.
最後聖經在第十至第十四節.
記載了如何對付下所和其餘北方城鎮的.
人民和城鎮本身.
對這場北方的戰役做一個總結.
第十至第十一節就提到下所.
第十二至第十四節就提到其他城鎮.
經文提到米倫之戰之後.
約書亞就回到攻下所.
殺掉王和城中一切的人口.
而其他城都是一樣.
但有所不同的是.
在其中有一個城鎮.
但有所不同的是.
下所城和其中一切盡被火燒毀.
而他們的城就得以保留.
其中的財物生畜被以色列人取為掠物.
在這裡處理城鎮的方法.
因為在其他經文也有提及.
所以我不在這裡重複了.
不過我很想提到一件事.
就是在這段經文裡有一個重點.
在第十二和第十五節.
他一再強調.
約書亞在城邦要滅盡人口這件事上.
是要如耶和華,獨人摩西所吩咐的.
而第十五節他再具體.
而且又從正反兩邊去論述.
耶和華如何吩咐他獨人摩西.
摩西就照樣吩咐約書亞.

$^{561}$約書亞又照樣行.
凡耶和華所吩咐摩西的.
約書亞無一件下台不行的.
其實在《申命記》第七章第一至二節.
那裡清楚地說明摩西吩咐他的百姓.
要在進入迦南地德爾維葉的時候.
趕走裡面的國民.
耶和華利神將他們交給你擊殺.
你要把他們滅絕淨盡.
不可與他們立約.
也不可憐恤他們.
摩西很清楚地說明上帝的心意.
進入迦南地.
耶和華利神為了愛他的子民的緣故.
不想他們被這些外族拜偶像的惡行.
影響到以色列人.
所以他要求他們除滅一切的人.
約書亞也一一遵從.
一件都不下台.
今天分享的經文是.
全體以色列人攻佔迦南地的最後一役.
以後的經文會交代各個支派的分地.
事實上分地之後.
他們不是就這樣得享太平.
他們還要按照各支派所被劃分的地區.
在裡面攻佔不同的城邑.
征戰仍然是不斷.
但是在過去幾場戰役裡.
我們看到上帝對以色列人的操練.
並以領袖約書亞為榜樣.
示範如何義無比的信心.
遵行上帝的吩咐.
以信服的心命完成每一場的征戰.
遵行神的話是人生得勝的秘訣.
即使上帝的話不是我杯茶.
不是我所願.
我們仍然要像約書亞一樣.
毫不猶豫.
立即跟從.
絕不下台.

$^{601}$我認識一位對神很有信心的人.
他非常遵從神的道.
在他的人生面對強敵.
他仍然樂觀積極地活出得勝的人生.
他是一位很值得我尊敬的前輩.
他叫做浪曼華博士.
他曾經在中國神學研究院.
跨文化宣教科擔任講師.
他人生走遍五大洲三十多個國家.
一百零一個城市.
他協助過無數的差傳.
或者關愛機構創立或者發展.
我自己對於傳福音.
差傳的領袖深受他影響.
但他給我最大最大的感覺就是.
他凡事尋求上帝.
他曾經分享過.
他自己如何在婚姻和獨身的事上.
去求問上帝.
因為他經常要去到創啟的地區服侍.
他對每一個事奉.
如何開展如何發展.
找甚麼人.
他都時時求問上帝.
他說過他絕對不會走前過上帝.
他一定是緊緊跟在上帝後面.
但2006年他被確診患上肺癌.
最初醫生估計他只剩下兩三年的命.
但最後他在2018年安息主懷.
在這12年裡.
我親眼見到他如何去面對這個強敵.
他沒有被嚇怕.
他熱切地尋問上帝如何使用他.
神吩咐他做的事.
他忠心地靠著上帝給他的力量.
一件一件地完成.
直至他相信神給他的使命已經完結.
他便安心求主接他回天家.
我們所有認識他的人都相信.
他已經打了美好的掌.

$^{641}$跑完當跑的路.
守著當守所信的道.
現在有公義顧面為他存留.
我最近看到這篇文章.
在PowerPoint裡面大家可以看到.
是老益思C.S. Lewis在1942年所寫的.
當時的背景是第二次世界大戰的倫敦.
學校停課.
經濟衰退.
聚會的地方要關門.
跟現在很相似.
文中中文是這樣說的.
撒旦說我要引起焦慮.
恐慌.
恐懼.
我要關閉企業.
關閉學校.
關閉敬拜的場所.
我要叫體育的賽事停止.
我要引起經濟動盪.
這是撒旦的作為.
這是撒旦的詭計.
但耶穌如何面對.
他的回應是.
我將要將鄰居聚在一起.
我要去恢復家庭關係.
我要將晚餐帶到家裡的廚房桌子.
我將幫助那些人放慢生活.
欣賞真正重要的事情.
我將要教我的孩子們.
依靠我而不是依靠世界.
我將要教我的孩子相信我.
而不是他們的金錢和物質的資源.
弟兄姊妹.
我們的人生會不斷遇到強敵.
對你來說.
你人生的強敵是甚麼呢?.
到底你怕甚麼呢?.
你要記住.
你怕甚麼.

$^{681}$你就會成為那件事的奴隸.
如果讓你選擇.
你要怕這些強敵.
還是要怕或敬畏你的上帝呢?.
如果你要過一個持續得勝的人生.
秘訣就是遵從神的吩咐.
遵主之命.
對祂的吩咐.
你一件事都不會太多.
這樣你就會得到一個得勝的人生.
讓我們一起低頭祈禱.
撒美耶和華.
你是我們生命的元帥.
人生裡會遇上有形無形不同的爭戰.
但我們知道我們不需要害怕.
因為你會和我們一起.
你叫我們要細心思想你的話.
並且謹守進行.
一見不懈怠.
上帝我就祈求你幫助我們眾弟兄姊妹.
心裡剛強膽壯.
因著你與我們同行的緣故.
面對生命不同的難關挑戰的時候.
我們仍然緊緊跟著你的吩咐去行.
過一個得勝有力的人生.
我們這樣禱告交託.
是奉主耶穌基督得勝明智祈求.
阿們.
\newpage



\section{約書亞記 12:1-24}
\label{sec:1PuKX5mdEMQ}
\textbf{2020年4月19日 崇拜直播｜約書亞記12\_1-24}
\newline
\newline
連結: \href{https://youtube.com/watch?v=1PuKX5mdEMQ}{\texttt{https://youtube.com/watch?v=1PuKX5mdEMQ}} ~~~~ 語音日期: 2020-04-21
\newline
\newline
\hyperref[sec:pz6DLtcCEWM]{\small{< < < PREV SERMON < < <}}
~
\hyperref[sec:index_chronic]{\small{[返順時目]}}
~
\hyperref[sec:0sy2FpLq_lk]{\small{> > > NEXT SERMON > > >}}
\newline
\newline
約書亞記 12:1-24
\newline
\begin{longtable}{cl}
\hline
\hline
章節 & 經文 (和合本修訂版)\\
\hline
12:1 & \begin{tabularx}{0.7\textwidth}{X} 這些是以色列人在約旦河東，向日出的方向，從亞嫩谷直到黑門山，以及東邊亞拉巴的整個地區所擊殺的王和所得的地： \end{tabularx} \\ \\ \relax
12:2 & \begin{tabularx}{0.7\textwidth}{X} 有住希實本的亞摩利王西宏，他統治的地從亞嫩谷邊的亞羅珥起，包括谷中之城和基列的一半，直到亞捫人邊界的雅博河， \end{tabularx} \\ \\ \relax
12:3 & \begin{tabularx}{0.7\textwidth}{X} 以及從東邊的亞拉巴，直到基尼烈海，又向東通過伯‧耶施末的路，直到亞拉巴的海，就是鹽海，再往南直到毗斯迦山斜坡的山腳。 \end{tabularx} \\ \\ \relax
12:4 & \begin{tabularx}{0.7\textwidth}{X} 又有巴珊王噩，他是利乏音人所剩下的，住在亞斯她錄和以得來。 \end{tabularx} \\ \\ \relax
12:5 & \begin{tabularx}{0.7\textwidth}{X} 他統治的地是黑門山、撒迦、巴珊全地，直到基述人和瑪迦人的邊界，以及基列的一半，直到希實本王西宏的邊界。 \end{tabularx} \\ \\ \relax
12:6 & \begin{tabularx}{0.7\textwidth}{X} 這兩個王是耶和華的僕人摩西和以色列人所擊殺的。耶和華的僕人摩西把他們的地賜給呂便人、迦得人和瑪拿西半支派的人為業。 \end{tabularx} \\ \\ \relax
12:7 & \begin{tabularx}{0.7\textwidth}{X} 這些是約書亞和以色列人在約旦河西所擊殺的諸王，他們的地從黎巴嫩平原的巴力‧迦得，直上到西珥的哈拉山。約書亞按著以色列支派所得的份把這地分給他們為業， \end{tabularx} \\ \\ \relax
12:8 & \begin{tabularx}{0.7\textwidth}{X} 就是赫人、亞摩利人、迦南人、比利洗人、希未人、耶布斯人的地，包括山區、低地、亞拉巴、山坡、曠野和尼革夫。 \end{tabularx} \\ \\ \relax
12:9 & \begin{tabularx}{0.7\textwidth}{X} 這些王是：耶利哥王一人，靠近伯特利的艾城王一人， \end{tabularx} \\ \\ \relax
12:10 & \begin{tabularx}{0.7\textwidth}{X} 耶路撒冷王一人，希伯崙王一人， \end{tabularx} \\ \\ \relax
12:11 & \begin{tabularx}{0.7\textwidth}{X} 耶末王一人，拉吉王一人， \end{tabularx} \\ \\ \relax
12:12 & \begin{tabularx}{0.7\textwidth}{X} 伊磯倫王一人，基色王一人， \end{tabularx} \\ \\ \relax
12:13 & \begin{tabularx}{0.7\textwidth}{X} 底璧王一人，基德王一人， \end{tabularx} \\ \\ \relax
12:14 & \begin{tabularx}{0.7\textwidth}{X} 何珥瑪王一人，亞拉得王一人， \end{tabularx} \\ \\ \relax
12:15 & \begin{tabularx}{0.7\textwidth}{X} 立拿王一人，亞杜蘭王一人， \end{tabularx} \\ \\ \relax
12:16 & \begin{tabularx}{0.7\textwidth}{X} 瑪基大王一人，伯特利王一人， \end{tabularx} \\ \\ \relax
12:17 & \begin{tabularx}{0.7\textwidth}{X} 他普亞王一人，希弗王一人， \end{tabularx} \\ \\ \relax
12:18 & \begin{tabularx}{0.7\textwidth}{X} 亞弗王一人，拉沙崙王一人， \end{tabularx} \\ \\ \relax
12:19 & \begin{tabularx}{0.7\textwidth}{X} 瑪頓王一人，夏瑣王一人， \end{tabularx} \\ \\ \relax
12:20 & \begin{tabularx}{0.7\textwidth}{X} 伸崙‧米崙王一人，押煞王一人， \end{tabularx} \\ \\ \relax
12:21 & \begin{tabularx}{0.7\textwidth}{X} 他納王一人，米吉多王一人， \end{tabularx} \\ \\ \relax
12:22 & \begin{tabularx}{0.7\textwidth}{X} 基低斯王一人，靠近迦密的約念王一人， \end{tabularx} \\ \\ \relax
12:23 & \begin{tabularx}{0.7\textwidth}{X} 多珥山岡的多珥王一人，吉甲的戈印王一人， \end{tabularx} \\ \\ \relax
12:24 & \begin{tabularx}{0.7\textwidth}{X} 得撒王一人，共三十一個王。 \end{tabularx} \\ \\
[1ex]
\hline
\hline
\end{longtable}
$^{1}$接著我們透過經文敬拜主.
我們請看約書亞記十二章一至二十四節.
是記載在舊約二百七十七至二百七十八頁.
經文裡面這樣說.
以色列人在約旦河外向日出之地擊殺二王.
突他們的地,就是從亞亂谷直到黑門山.
並東邊的全阿拉巴之地.
這二王有著希實本阿摩利人的王西雲.
他所管之地是從亞亂谷邊的亞羅爾和谷中的城.
並激烈一半直到阿摩利人的境界亞博河.
與約旦河東邊的阿拉巴直到基利烈海.
又到阿拉巴的海,就是嚴海.
通伯耶西末的路以及南方直到比斯加的山根.
又有巴沙王岳.
他是利弗欽人所盛夏的.
住在亞斯他祿和爾德萊.
他所管理的地是黑門山.
剎那巴山全地直到基述人黃馬加人的境界.
並激烈一半直到希實本王西雲的境界.
這二王是耶和華伯人摩西和以色列人所擊殺的.
耶和華伯人摩西將他們的地.
賜給劉邊人,加德人和嗎哪西半支派的人為業.
第七節.
「約書亞和以色列人在約旦河西擊殺了諸王.
他們的地是從黎巴嫩平原的巴力加德.
直到上西伊的哈拉山.
約書亞就將那地安著以色列支派的宗族.
分給他們為業.
就是黑人,阿摩利人,迦南人,比利西人.
希美人,耶布斯人的山地.
高原阿拉巴,山坡,曠野和南地.
他們的王.
一個是耶利哥王.
一個是靠近伯特利的艾城王.
一個是耶路撒冷王.
一個是希伯倫王.
一個是耶美王.
一個是拉吉王.
一個是伊基倫王.
一個是基瑟王.

$^{41}$一個是底碧王.
一個是基德王.
一個是荷爾瑪王.
一個是阿拉德王.
第十五節.
一個是納拿王.
一個是亞道蘭王.
一個是馬基大王.
一個是伯特利王.
一個是托普亞王.
一個是希伯王.
一個是亞弗王.
一個是拉沙倫王.
一個是馬頓王.
一個是夏索王.
一個是新倫米倫王.
一個是阿薩王.
一個是他納王.
一個是米吉多王.
一個是基底斯王.
一個是靠近加密的約廉王.
一個是多爾山江的多爾王.
一個是吉甲的哥印王.
一個是德薩王.
共計三十一個王.
今日的講題是.
「邁向新世代,真成全」.
各位弟兄姊妹.
能夠和大家一起敬拜上帝.
一起能夠聽上帝的話語.
讚美上帝的恩典.
讓我們一起先禱告.
「慈愛恩主,你是我們的盾牌」.
「無論仇敵,環繞,攻擊我們」.
「無論在千般駭浪當中」.
「或者我們在不能控制的環境下」.
「你是王,你坐著為王」.
「你的同在叫我們出黑暗」.
「你的同在叫我們不怕遭害」.
「能夠乘風破浪」.

$^{81}$「祈求你激進我們的心思意念」.
「祈求你打開我們的心」.
「讓聖靈引導,感動我們」.
「一同聆聽你的話語」.
「並且在這個世代裡潛行出來」.
我們同心禱告.
是奉耶穌基督,聖命祈禱.
阿們.
最近不知道大家有沒有在網上看到這個圖片.
這是一個實習的醫生.
他手拿著的是一個護目鏡.
不過這個護目鏡原來很有價值.
這個護目鏡就是他爸爸在2003年.
當我們香港面對SARS的時候.
他爸爸自願參加前線抗疫團隊.
即是我們所說的Dirty Team.
17年後.
女兒她同樣自願加入.
去照顧一些確診的病人.
爸爸昔日所作的.
對她有一定的影響.
各位弟兄姊妹.
不知道哪一個人曾經影響了你.
不知道又有哪一個生命曾經感動了你.
剛才大家聽我一口氣讀《約書亞記》十二章.
我們大概可以分為兩部分.
一至六節.
它就說到以色列人在約旦河東所擊殺的.
原文是這樣說.
「這些就是那地的王.
是以色列人所擊殺也佔領他們的地」.
作者希望我們的焦點就是看當地的君王.
剛才我讀的是「得他們的地」.
這個「得」字在原文裡面也有「佔領」的意思.
我們來看看地圖.
「莫世帶領以色列人得到河東之地」.
藍的就是亞聯谷.
北邊就是黑門山.
阿拉巴在哪裡?.
阿拉巴就是約旦河.

$^{121}$沿著約旦河.
一個狹長的山谷.
所以經文裡面說「東邊的阿拉巴之地」.
就是指阿拉巴東面所有的地區.
由於作者希望我們專注在一些王.
所以我們接著發覺幾個經文.
他描述到阿摩利人王西弘和巴山王鄂.
他們所管治的地區.
阿摩利人是迦南民族之一.
在摩西時代他們管轄了河東大部分的地圖.
特別我們看民數記二十一章二十七至三十節.
他們吞滅了摩亞人在亞聯河以北的地圖.
奪取了希什本城.
究竟西弘得到的領土有多大?.
南端,剛才我提及了.
就是阿聯谷.
阿聯谷,我們發覺往東就是阿羅爾這個地方.
阿羅爾向北就是廣闊的基烈這個地方.
亞博河,我們很容易認得基烈地.
我們找到亞博河.
亞博河就將基烈地分為兩部分.
而西弘王得到一半.
他的地圖不止這麼少.
他還由阿拉巴的東側直上到基利烈海.
即是我們熟悉的加利利海.
然後他通過帕耶西末直到比斯加的山腳.
剛剛我們說完西弘.
我們就說另一個統治者.
巴山王鄂.
經文裡提到他是彌弗音人.
之前我們的牧師也跟我們分享過.
彌弗音人就是巨人.
在《新面記》三章裡.
很特別,他記載巴山王鄂的床有多大?.
原來是九爪乘四爪.
我們會問,九爪乘四爪即是多少?.
我們可以想像.
大概是四米乘米八,兩道門那麼高.
他過往在阿斯他祿和以德萊.
他得到這兩座城市.

$^{161}$他統治河東北部的地方.
而且他也包括剛才那裡的黑門山.
薩加,巴山全地.
直到基述人和瑪嘉仁的境界.
剛才我不是說過亞伯河將基烈分成一半?.
沒錯,另一半就是巴山王鄂所得.
接下來我們看第六節.
原本的文字裡面,他先這樣說.
「耶和華的僕人摩西和以色列人擊殺了他們」.
亦提到「耶和華的僕人摩西.
將他們的地賜給留便人加德人.
和馬來西亞半支派的人為業」.
我們想到這個爭戰.
我們就想到在《民數記》21章有這樣的記載.
當時以色列人在曠野已經四十年了.
摩西就正正帶他們進入迦南.
是從東邊進入迦南.
他們就求西雲,這個阿摩利人的王西雲.
准許他們穿過他們的地圖.
但是西雲拒絕這個社交的外交的要求.
他們還襲擊以色列人.
結果上帝就讓以色列人打敗西雲王.
接下來《民數記》21章就記載.
以色列人就上瓦遮.
接著他們就去到巴山.
巴山王知道之後.
他就領全軍出戰.
在他們面前有這麼多巨人的軍隊.
但是上帝同樣使以色列人得勝.
這件事的事跡在書篇135篇有記載.
為什麼會記載呢?.
原來這是進應許地一個很重要的階段.
它讓以色列人看到上帝的能力.
上帝的信實.
《民數記》32章說到劉變加德.
他們這兩個支派求摩西能夠得到產業.
摩西就提出要求.
希望他們能夠參與之後的約旦河西的征戰.
然後嗎哪西支派同樣加入.
就回應我們約書亞的第一章.

$^{201}$而《民數記》32章39節.
它特別提到嗎哪西半支派的人.
他們將激烈地阿摩利人全部趕走.
接著我們回到約書亞的12章.
我們看下半部7至24節.
請聽我讀第7節.
約書亞和以色列人在約旦河西擊殺了諸王.
他們的地是從黎巴嫩平原的巴力加德.
直到上西爾的哈拉山.
約書亞就將那地安著以色列人支派的宗族.
分給他們為業.
是的,原文同樣都是先說.
這些就是那地的王.
他想我們的焦點放回這些王身上.
剛才我有讀到巴力加德這個地方.
我們要留意的是.
巴力我們不陌生.
是迦南人所拜的神.
迦南人,迦南地和鄰近的地區.
他們很多時候都以巴力為名.
所以後來15章60節.
我們發現猶大支派得地的時候.
其中一個地方就是基列巴力.
當地有指這些地方有巴力的寺廟.
這正正解釋為何以色列人.
當他們不趕走這些迦南人.
他們不去消滅這些寺廟的時候.
就會影響他們和上帝的關係.
剛才大家聽我一口氣讀這麼多王.
第九節所呈現.
到第24節所呈現的地名和統治者.
原來正正反映在約書亞記6到11章.
戰爭的記載.
我們看看圖.
第九節.
他呼應約書亞記6章.
講到耶利哥城和八張艾城.
事實上我們的確欣賞他們的戰略.
因為他們先打中部.
他們讓迦南地一分為二.

$^{241}$將南北分割.
然後個別擊破.
第十到第十六節.
原來正正呼應約書亞記第十章.
講到由五王所組成的南方聯盟.
第十九至二十四節.
呼應約書亞記十一章.
就是以夏所王為首的北方聯盟.
剛才大家聽我讀的時候.
不知道會否替我緊張.
要讀這麼多王的名字.
不過我相信大家除了替我緊張之外.
我們一定會感受到.
用一個形容.
怎樣形容?.
是大獲全勝.
我們看到以色列人.
他們打這麼多個仗都打贏.
似乎我們越數的時候.
就像數以色列人的戰績.
他們的威水史.
不過.
經文是否想我們這樣呢?.
原來這些數點.
是讓我們在重重的戰爭裡.
我們更加數到上帝的恩典.
更加看到以色列人.
他們在打仗期間對上帝的馴服.
剛才.
鎮壇執事的分享.
他說到這個疫情.
我們很難掌握到.
無論是疫情現在.
或者疫情之後.
我們好像仍然在打不同的仗.
每日電視機的數字.
不多不少.
都影響了我們的心情.
剛才.
經文裡面勉勵我們.

$^{281}$在多少場仗.
無論是得,失.
我們都要數算上帝的恩典.
我不知道今天.
大家在這個疫情裡.
我們數不數到上帝的恩典.
我看到很多很美麗的人性光輝.
我又看到很多肢體.
他們去發掘自己的一些興趣.
我又看到大家很珍惜.
很想見的人.
其中一個恩典.
會否就是我們透過今次的疫情.
是一個好機會.
讓我們反省我們和世界.
和人和上帝的關係.
韋爾斯·Samuel.
在他的一篇哀訴的圖文裡面.
他勉勵我們想三件事.
第一.
究竟是世界因為我們而失常.
還是我們因世界而失常.
第二點.
求上帝指教我們.
活出做人的樣式.
活出群體的樣式.
第三點.
這位牧師他勉勵我們.
最重要的是在疫情當中.
我們求上帝復興我們.
讓我們重新依靠上帝.
或者在這裡我補充一下.
剛才我讀這三十一位的王當中.
有十六位其實是記載在.
剛才我講的.
《約書亞記》六至十一章.
這是一個概要的記載.
例如.
為何我們這樣說呢?.
例如在二十一章二十一節.

$^{321}$就講到以色列人得到事件這個地方.
不過剛才大家應該沒有聽到我讀這個名字.
反過來.
十二節講的基式.
這個地方.
原來不是當時以色列人能夠完全.
盡情地得到這個地方.
是後來到了撒母耳記下五章.
國王時代才能夠得到佔領.
接下來.
我們透過一些經文.
我們看看結構.
我們先看河東這一篇.
一至六節.
我們發現首先講到清單.
然後就講摩西的角色.
接著就講摩西和以色列人.
他們擊殺君王.
然後摩西就將.
將這些支派分給兩個支派為業.
接著我們看另一幅圖.
在七至二十四節裡.
他就講河西.
首先他講約書亞的角色.
然後就講約書亞和以色列人.
擊殺這些君王.
接著約書亞就將這些地.
分給以色列為業.
最後就是這些君王的清單.
我們留意到.
留意到在七至八節裡.
約書亞做了兩件事.
第一.
他就是將他和以色列人.
他們擊殺這些王.
第二.
他們就將這些地.
分給以色列的支派為業.
這就是正正對應摩西所作的.
我們發現到無論是摩西和約書亞.

$^{361}$他們同樣得到上帝的幫助.
我們看到上帝的應許.
得以應驗.
他是多麼信實.
我們又看到.
無論河東,河西.
他們所領受的土地.
都是上帝所賜的.
神的子民無分彼此.
他們展示合一.
亦都讓我們看到合一的寶貴.
或者我們又拿兩個的戰爭來看一看.
在《申命記》裡面.
他就記載到河東西宏之戰.
而早前我們牧者都分享了艾城之戰.
我們發現.
我們看一看的時候.
無論是耶和華的宣告.
這些王他們攻擊以色列.
到後來耶和華將他們交給以色列的手中.
擊敗他們.
奪去城鎮等等.
裡面的編排,結構,元素.
都是很相近.
我想說到這裡.
大家都知道我想說甚麼.
原來貞貞告訴我們.
約書亞就是摩西的承繼者.
究竟承繼者是怎樣出現的呢?.
在《申命記》一章三十八節.
他這樣說.
「事後你戀的兒子約書亞.
他必得進入那地.
你要大免利他.
因為他要使以色列人承受那地遺孽」.
承繼者是怎樣形成的呢?.
原來是上帝的旨意.
裡面描述到約書亞.
就是上帝所揀選.
接著摩西這一棒.

$^{401}$承繼者是怎樣形成的呢?.
我們發現摩西在他死之前.
他已經為百姓預備了一個.
這麼好接他這一棒的領袖.
就是約書亞.
我們發現約書亞常常伴隨著摩西.
他看到摩西當百姓犯罪拜牛毒的時候.
摩西怎樣去處理.
而在《民數記》二十七章.
又提到摩西在會眾面前.
他親自按手祝福約書亞.
摩西好好培育和儲備這個接棒人.
承繼者是怎樣形成的呢?.
有說貓是老虎的師父.
所以貓就把十八般的武藝給了老虎.
但是老虎學會了之後.
就看不起貓了.
他還想一口就吃了老虎的肚子.
不過貓也不愚蠢.
他留下一個招數沒有教老虎.
是什麼招數呢?.
原來就是爬上樹.
所以有傳當老虎一想吃回自己的師父的時候.
貓就怎樣做呢?.
貓就跑到樹上.
老虎就只能夠望貓輕嘆.
我想問問大家.
摩西在培育約書亞的時候.
他有沒有留一手?.
摩西即使知道自己不能夠進入迦南美地.
他仍然無私地帶領約書亞.
繼續邁向上帝所賜給他們的目標.
進迦南美地.
那麼約書亞又怎樣呢?.
作為接這一棒.
約書亞謙謙卑卑地受教.
他同樣得到清晰的意象.
別人一定將他和摩西比較.
但是他不作比較.
他不是要贏過自己的師父.

$^{441}$更加重要的是.
他看到摩西怎樣帶領以色列人.
打這些河東之戰.
在這麼大的戰爭的時候.
他看到自己的師父在何等信靠上帝的時候.
都成為約書亞一個榜樣.
我們再回頭看約書亞的十二章一節.
大家再聽我讀.
「以色列人在約旦河外.
向日出之地擊殺二王.
得他們的地.
就是從雅聯谷直到黑門山.
並東邊的傳阿拉巴之地」.
一開始我講信息的時候.
我有提到「得」字在原文裡面的意思.
是「佔領」.
但是我們發現.
原來在經文裡面.
在第十二章就沒有記載到.
約書亞佔領這些地圖.
而剛才我們也特別提到.
嗎哪西派半支派.
他們將激烈的亞摩列人趕走.
但是我們在約書亞的十五章,十六章,十七章裡面.
我們都發現到.
經文特別提到以色列人.
他們未能趕出一些迦南的民族.
說到這裡.
我們發現摩西在死之前.
他已經將領導權轉交給約書亞.
不過我們看聖經的時候.
我們找不到約書亞死之前.
他預備一個領導人.
結果接下來就出現時思局面.
約書亞盡心盡力去承繼摩西所吩咐他的.
摩西得到上帝的吩咐.
然後約書亞又聽摩西的吩咐.
他照樣行.
約書亞是一個很好的承繼者.
不過.

$^{481}$他沒有傳遞下去.
他沒有傳遞下去下一個.
消極來說.
我們看到這個是事實.
不過這是不是一個結局呢?.
有危就有機.
讓我們積極思想兩個問題.
第一.
為什麼佔領地圖趕出迦南人是這麼重要?.
第二.
究竟哪一個是繼承者?.
哪一個是傳遞者?.
土地在約書亞記是一個很重要的課題.
以色列人的約旦行就是要得地為業.
有學者指出.
我們表面看到的就是以色列人和地的關係.
其實核心最重要的.
就是以色列人和上帝的關係.
以色列人他們到了迦南美地.
他們得安息,得平安.
他們有很豐富的物資.
但是背後是誰?.
是誰保護他們?.
是誰與他們同在?.
是上帝自己.
原來真正的得地.
最寶貴的不是得地.
是得到上帝自己.
所以為什麼經文裡.
一而再再而三提醒我們.
要除去這些迦南人.
除去這些誘惑我們遠離和上帝的關係.
一些人,一些事,一些罪行.
今天我和在家中每一個弟兄姊妹.
我們都得到地了.
我們得到上帝了.
我們就要問.
我們可以做些什麼呢?.
會不會就是透過我們的生命.
我們的生活見證.

$^{521}$我們將耶穌基督.
並且祂的十字架.
我們傳揚出去.
第二樣我們要想.
哪一個是承繼人呢?.
哪一個是傳遞人?.
我自己照一照快鏡.
看一看自己.
我的確不是一個什麼師父.
是一個平凡人.
又好像上個星期.
牧師這樣說.
我們是一個小人物.
我們沒有什麼豐功偉績.
也沒有什麼戰績存留.
但是原來在上帝眼中.
我們每一個都可以承繼.
也都可以傳遞.
或者容許我透過一個電話.
一個影音,一個鏡頭.
一頓飯.
跟大家說一說.
有關黃浸一些生命的故事.
我不知道大家有沒有想念.
我們多加團.
或者女底亞團.
或者等等一些長者.
我想跟大家說.
你們放心.
他們更加在這個疫情裡面.
他們發揮互相守望的精神.
當我們一班牧者打電話給他們的時候.
他們很多時候都是反胃.
問候我們,擔心我們.
我自己其中一個服侍.
很開心就是女傳道會.
女傳道會每次聚會.
有一個很特別的地方.
令我有些回憶.
就是會點名.

$^{561}$他們就會說.
誰就到.
我從一個多加團的肢體知道.
原來有一個姐妹.
在九十年代.
她就是負責聯絡這個崗位的.
她其實本身自己身體不是很好.
但是她因為愛.
因為愛上帝.
她謹守崗位.
她常常掛念這個,掛念那個.
如果肢體一剔到.
沒有出席.
她就立即打給他們.
又或者當剔到的時候.
讀到某些名字沒有出席.
她就很快告訴其他人.
為什麼他們沒有出席.
那是不是她是先知呢?.
當然不是.
是因為她很用心去打給每一個.
不過我們要留意.
他們用的是家居電話.
如果有一些肢體出了街.
那這位姐妹就要打幾次.
打到真的找到這些姐妹為止.
真的很不容易.
後來跟我分享這件事的多加團姐妹.
她就是承接這一棒.
她自己本身身體都是一般般.
她說有時真的不想打電話.
因為比較累.
但是她一想到.
一想到這位肢體.
那麼盡力去慰問其他人的時候.
這個姐妹她反而說.
她真的不召喚這些肢體.
她還不舒服.
而現在她也將這個傳遞給下一位的姐妹.
由受苦到復活.

$^{601}$這個靈修系列的熱誠.
我自己有份參與.
很開心,很感動.
特別後來看回每一段片的成果.
以及大家的代討.
以及大家的使用.
我們有同工.
他藉著這個機會.
他邀請一些天恩堂的.
或者我們往生的少年人.
他去培育他們.
叫他們試試剪片.
試試幕後的工作.
大家的目標.
大家的目標就是對準上帝.
希望透過這件事將福音傳開去.
有人給機會.
有人很謙卑地學習.
時間過得很快.
我還記得上年.
當我上下樓梯.
看到少年部宣傳營地的時候.
那張海報.
也令我很為之一振.
因為很對應這個潮流.
也很能夠捉到年輕人想參加的心.
當時自己已經有一個.
全心全力籌辦.
他們親自去玩.
跟這些少年人玩了幾天.
有些甚至是未見過.
我想問大家.
容易嗎?.
為什麼?.
為什麼這班少年部的導師.
他們願意這樣做呢?.
原來.
我問其中一個.
原來他回想.
以往他的導師.

$^{641}$就是希特勒.
他的導師就是犧牲時間.
跟他們去露營.
他們的導師就用盡心機.
去預備不同的聚會.
這就成為現在這位導師的意象.
我們的少年人.
很多個都有團伙在當中.
到他們接第三棒的時候.
他們又見到他們的導師.
是沒有價值.
願意花這麼多時間跟他們玩.
跟他們去相處.
就成為他們自己的目標.
之前.
我們星期日早上.
我們有一些肢體會回來.
可能是幫忙一些少年主一學的教學等等.
我見到有一個情景.
我見到有一個中學團契的導師.
有一個姐妹.
她買早餐.
她搭訕九點.
她見到店舖開.
她就買.
我都很喜歡吃的早餐.
就給祖月.
然後我就問她.
為什麼你這樣做?.
她就回想昔日她的導師.
就是這樣.
請她吃飯.
陪伴她.
關心她的需要.
我不知道這個祖月.
日後她成為導師的時候.
她會不會同樣受影響?.
她同樣這麼貼心地去關懷她的祖月.
一開始.
我提到這個實習醫生.

$^{681}$我提到他們守護我們市民的健康.
而我們作為上帝的兒女.
我們守護的是人的心靈.
是關於生命的問題.
我們沒有什麼護目鏡.
不過我希望有一天.
有一天.
有人拿起.
拿起是由我們傳遞出來.
他們從我們生命上面.
他見到一些美事.
見到一些見證.
然後他就承接我們這一棒.
重要的是.
傳遞出去.
光環不是屬於我們.
榮耀是歸給上帝.
最後.
我問大家一個問題.
你的生命.
影響了哪一個?.
你的生命.
感動了哪一個?.
有人承繼.
就有人傳遞.
這個就是真誠傳的意思.
祈求我們靠著上帝.
我們每一個的生命.
我們看似為很小的事.
卻能夠影響.
影響現在黑暗.
遇著很多風浪.
他們要面對很多事情的人.
求主幫助我們.
謝謝大家.
\newpage



\section{希伯來書 4:1-11}
\label{sec:0sy2FpLq_lk}
\textbf{2020年5月3日 崇拜直播｜希伯來書4\_1-11}
\newline
\newline
連結: \href{https://youtube.com/watch?v=0sy2FpLq-lk}{\texttt{https://youtube.com/watch?v=0sy2FpLq-lk}} ~~~~ 語音日期: 2020-05-05
\newline
\newline
\hyperref[sec:1PuKX5mdEMQ]{\small{< < < PREV SERMON < < <}}
~
\hyperref[sec:index_chronic]{\small{[返順時目]}}
~
\hyperref[sec:0WamuVI2aAk]{\small{> > > NEXT SERMON > > >}}
\newline
\newline
希伯來書 4:1-11
\newline
\begin{longtable}{cl}
\hline
\hline
章節 & 經文 (和合本修訂版)\\
\hline
4:1 & \begin{tabularx}{0.7\textwidth}{X} 所以，既然進入他安息的應許依舊存在，我們就該存畏懼的心，免得我們 中間有人似乎沒有得到安息。 \end{tabularx} \\ \\ \relax
4:2 & \begin{tabularx}{0.7\textwidth}{X} 因為的確有福音傳給我們像傳給他們一樣；只是所聽見的道對他們無益，因為他們沒有以信心與所聽見的道配合。 \end{tabularx} \\ \\ \relax
4:3 & \begin{tabularx}{0.7\textwidth}{X} 但我們已經信的人進入安息，正如神所說：「我在怒中起誓：他們斷不可進入我的安息！」其實造物之工，從創世以來已經完成了。 \end{tabularx} \\ \\ \relax
4:4 & \begin{tabularx}{0.7\textwidth}{X} 論到第七日，有一處說：「到第七日，神就歇了他一切工作。」 \end{tabularx} \\ \\ \relax
4:5 & \begin{tabularx}{0.7\textwidth}{X} 又有一處說：「他們斷不可進入我的安息！」 \end{tabularx} \\ \\ \relax
4:6 & \begin{tabularx}{0.7\textwidth}{X} 既有這安息保留著讓一些人進入，那些先前聽見福音的人，因不信從而不得進去， \end{tabularx} \\ \\ \relax
4:7 & \begin{tabularx}{0.7\textwidth}{X} 所以神多年後藉著大衛的書，又定了一天—「今日」，如以上所引的說：「今日，你們若聽他的話，就不可硬著心。」 \end{tabularx} \\ \\ \relax
4:8 & \begin{tabularx}{0.7\textwidth}{X} 若是約書亞已使他們享了安息，後來神就不會再提別的日子了。 \end{tabularx} \\ \\ \relax
4:9 & \begin{tabularx}{0.7\textwidth}{X} 這樣看來，另有一安息日的安息為神的子民保留著。 \end{tabularx} \\ \\ \relax
4:10 & \begin{tabularx}{0.7\textwidth}{X} 因為那些進入安息的，也是歇了自己的工作，正如神歇了他的工作一樣。 \end{tabularx} \\ \\ \relax
4:11 & \begin{tabularx}{0.7\textwidth}{X} 所以，我們務必竭力進入那安息，免得有人學了不順從而跌倒了。 \end{tabularx} \\ \\
[1ex]
\hline
\hline
\end{longtable}
$^{1}$等姐妹今日終於來到邁向新世代.
這個系列最後一講.
不過今日我們不是宣講約書亞記.
而是希伯來書第四章一至十一節.
經文裡是這樣說.
我們既蒙留下有進入他安息的應許.
就當畏懼.
免得我們中間或有人似乎是趕不上了.
因為有福音傳給我們.
像傳給他們一樣.
只是所聽見的道與他們無益.
因為他們沒有信心與所聽見的道調和.
但我們已經相信的人得以進入那安息.
正如神所說.
我在路中起誓說.
他們斷不可進入我的安息.
其實造物之功從創世以來已經成全了.
論到第七天.
有一處說.
到第七天神就揭了他一切的功.
又有一處說.
他們斷不可進入我的安息.
既有必進安息的人.
那先前聽見福音的.
因為不信從不得進去.
所以過了多年.
就在大衛的書上又限定一日.
如以上所引的說.
你們今日若聽他的話.
就不可硬著心.
若是約書也已叫他們享了安息.
後來神就不再提別的日子了.
這樣看來.
必另有一安息日的安息.
為神的子民傳流.
因為那進入安息的.
乃是揭了自己的功.
正如神揭了他的功一樣.
所以我們務必竭力進入那安息.
免得有人.

$^{41}$學那不信從的樣子跌倒了.
弟兄姊妹我們留意到.
剛才所讀的這段經文.
他很多次這樣重複了.
兩個對立的字句.
一個就是進入神的安息.
另一組就是硬心不信從.
簡單來說.
作者提醒我們.
心硬不信從的人.
就不能進入上帝的安息.
並且作者說.
其實有前人試過是這樣.
剛才我們特別聽到.
作者舉了一個例子.
第八節說.
若是約書也叫他們享了安息.
後來神就不再提別的日子了.
咦?.
我們之前一直說約書也記.
我們還記得.
約書也不是已經帶領了以色列人.
進行上帝的吩咐.
打贏了一場又一場的聖仗.
開始在迦南地得地為業.
為何希伯來書的作者會說.
他們仍然享受不到安息.
希伯來書的作者這句說話.
又對作為新約的信徒的我們.
有甚麼提醒.
或者說得白一點.
我們會否也像昔日的以色列人.
失去了上帝所賜的安息的應許.
正是今天這段經文.
他要我們來正視.
或者我們首先去理解一下.
在希伯來書所說.
上帝的安息到底是甚麼一回事.
作者說.
神的安息是要竭力進入的.

$^{81}$很明顯他不是說.
一個信徒離世.
息怒歸主的那一種安息.
作者特別引述了《創世記》.
第三章的下半節至第四節.
他這裡這樣說.
其實造物的工從創世以來已經成全了.
輪到第七天有一處說.
到第七天.
神就揭了他一切的工.
我們知道這段經文是引述.
《創世記》第二章一至二節.
說到上帝用了六天.
做齊了天地的萬物.
第七天就進入安息.
到底是否因為連續工作了六天.
覺得很疲累.
所以第七天.
上帝就決定甚麼都不做.
放日假休息一下.
充充電.
大概這可能是.
我們連續工作了五六天.
之後所有的想法.
其實這裡說上帝的安息.
不是說上帝完全靜止.
甚麼都不做.
你還記得耶穌基督說過一句話.
「我負作事直到如今,我也作事」.
提醒我們其實上帝沒有一刻停止過.
對萬物包括你和我的掌管和看顧.
在這裡說揭了工或安息.
其實他們原文是同一個意思.
裡面有欣賞,滿足,喜樂的意思.
其實上帝用了六天.
就已經將萬物做齊了.
不只是做了那些東西出來.
是包括上帝對萬物的期待.
萬物如何運作.
才能滿足上帝的旨意.

$^{121}$這一切上帝說已經做完了.
沒有東西需要加減.
可以說是畫上了一個完美的句號.
於是上帝在第七天就開始.
欣賞自己所完成的一切工作.
享受受造之物如何按照上帝的心意.
一一運作,一一發揮他們的所長.
這就是上帝的安息的意思.
作者表示上帝的安息.
自從創世以來就已經存在.
到底受造之物如何才能享受上帝的安息.
當然不是說我們雙手合十.
甚麼都不做.
而是說甚麼時候萬物按照上帝的旨意運作.
甚麼時候我們聽從上帝的吩咐.
遵從祂的話來生活的時候.
我們也會進入上帝的安息.
享受到與神同行.
與上帝一起工作的滿足和豐盛.
相反,人如果不願意聽.
不願意相信上帝的帶領和吩咐.
即使他多努力尋求生活的安穩.
他仍然無法進入聖經所說的真正的安息.
為了要說明這個道理.
作者引用了一個例子.
他用了第一代出埃及的以色列人.
作為一個借鏡.
事實上整段的渲染由第三章第七節開始.
不過我們直接看第三章的十五至十九節.
他提到第一代出埃及的以色列人.
為何進入不了上帝的安息.
經文這樣說.
經上說:你們今日若聽他的話.
就不可硬著心像惹他發怒的日子一樣.
那時聽見他話惹他發怒的是誰呢?.
豈不是跟著摩西從埃及出來的眾人嗎?.
甚四十年之久又厭煩誰呢?.
豈不是那些犯罪.
廝守島在曠野的人嗎?.
又向誰起誓不容他們進入他的安息呢?.

$^{161}$豈不是向那些不信從的人嗎?.
這樣看來他們不能進入安息.
是因為不信的緣故了.
作者在這裡說第一代出埃及的以色列人.
之所以不能夠進入上帝的安息.
原來並不是因為他們沒有聽過.
不知道上帝的應許和大令.
事實上他們連約都跟上帝納了.
要做上帝的子民.
只不過我們說好像廣東話所說.
「雖心硬,不肯信從上帝的吩咐」.
特別你會見到在摩西五經裡.
每當他們遇到困難逆境.
譬如說沒有水喝.
又或者沒有肉吃的時候.
很多時候就會怎樣呢?.
就會用一個抱怨的態度.
用一個懷疑的態度來回應上帝.
甚至拒絕上帝的大令.
想回到埃及.
我們說探子事件.
就是一個最極致的例子.
某程度上他們會覺得.
摩西你叫我們上去迦南.
去面對那些強大巨大的敵人.
不就是叫我們上去送死嗎?.
他們不單單那一刻不想上去迦南.
甚至乎聖經記載他們想拿石頭.
打死摩西和亞倫.
結果上帝當時.
向那一群心硬不信從的以色列人發怒.
要他們漂流甚至死在曠野裡.
不容許他們進入安息之地.
進入上帝的應許.
進入上帝的安息.
作者就是這樣引用這個例子.
作為一個借鏡.
第一代出埃及的以色列人不行.
第二代呢?.
約書亞所帶領的以色列人.

$^{201}$又能不能夠進入上帝的安息裡?.
我們一直以來.
過往幾個月看約書亞基里會發現.
除了阿干事件之外.
我們似乎看不到他們有.
心硬或者不信從的情況出現.
特別在約書亞記第23章裡.
是這樣記載的.
第23章第一節.
這裡記載.
「耶和華使以色列人安靜.
不與四圍的一切仇敵爭戰」.
已經多日了.
約書亞年紀老邁.
這裡記載上帝使那一群.
已經在迦南地的以色列人.
確實享受了很多年安穩的生活.
不過,到了約書亞的最後一章.
第24章的31節.
作者原來埋下了一個伏筆.
這裡這樣說.
「約書亞在世和約書亞死後.
那些知道耶和華為以色列人所行諸事的長老還在的時候.
以色列人事奉耶和華」.
弟兄姊妹,暗示著什麼?.
就是說,原來到了士司紀的時代.
當約書亞和他那個世代的長老都相繼離世之後.
以色列又開始學他們上一代那樣.
重蹈覆轍.
開始心硬.
不再順從上帝的吩咐.
舉一點點例子,在士司紀裡.
簡單看兩段經文.
你會明白他們當時的境況.
在士司紀一章第28節.
這裡這樣說.
「及至以色列強盛了.
就使迦南人做苦工.
沒有把他們傳言的講出」.
第二章說.

$^{241}$「以色列人行耶和華眼中看為惡的事.
去事奉朱巴力.
利器了令他們出埃及地的耶和華他們列祖的神.
去敬拜別神.
就使四位列國的神.
惹耶和華發怒.
耶和華的怒氣向以色列人發作.
就把他們交在搶奪他們的人手中.
甚至他們在受敵面前再不能站立得住」.
讓姐妹們看到.
新一代在迦南的百姓.
他們這次不是因為遇到逆境.
不是困難所以離開上帝.
相反他們順境.
他們強盛起來.
他們生活改善了.
甚至想提升他們的生活質素.
所以離棄上帝.
簡單說.
當時他們受到迦南人的文化影響.
人人都把戰俘當做奴隸.
人人在迦南地想有好的農作收成.
就一定是拜巴力.
並不止拜巴力.
多拜幾個神明來保自己平安.
所有人都這樣做.
如果我們不這樣做.
那不是很吃虧嗎?.
「執輸」.
於是慢慢上帝的吩咐.
再不像昔日那樣.
是百姓生活的首選.
已經淪為了他們眾多信仰裡.
一個選項.
有需要,有好處,有時間.
我才跟隨.
慢慢變成這樣.
弟兄姊妹我們看到.
以色列人雖然已經安穩在應許之地.
但卻享受不到上帝賜給他們安息的應許.

$^{281}$可不可惜?.
確實真的很可惜.
就等同於我們今天見到.
有人信了耶穌.
卻享受不到耶穌所應許的安息.
一樣的可惜.
原因可能是.
人雖然信了耶穌.
甚至有參與教會的活動.
有了事奉.
但骨子裡.
卻仍然靠著自己.
順著自己的心意來生活.
不太願意相信跟從上帝的吩咐.
不肯讓耶穌來掌管.
改變他們的人生.
就像以色列人在迦南地生活一樣.
你看到他們很努力地打拼.
為了生計.
用盡方法來打造一個安穩的生活給自己.
但因為他們服侍的對象不是上帝.
服侍了其他東西.
錯了對象.
自然就享受不到真正的安息.
所以弟兄姊妹.
作者透過這些例子.
來勸勉今天的我們.
第四章七節下半節說.
「你們今日若聽他的話.
就不可硬著心」.
我們說重要的話會說多少次?.
三次.
這裡也說三次.
作者在整段的勸勉裡說了三次.
「你們今日若聽他的話.
就不可硬著心」.
硬心是怎樣來的呢?.
通常不是突然間爆出來.
硬心我們說是慢慢孕育出來.
就像我們今天聽到上帝的話.

$^{321}$我們說遲一下再跟.
拖一拖再跟.
隔兩天又聽到了.
不過我們又想著.
不要緊的 拖多兩天.
遲一點再做.
隔了一頭半個月.
隔了一年兩年.
上帝又再提醒我們的時候.
我們還是說拖一拖.
遲一點再做的時候.
心就是這樣硬回來.
我還記得上一年.
19年.
19年年初.
我很想買一張電腦台.
來將房間裡的電腦搬出客廳.
來設置好.
左隔隔又隔隔.
結果在一間家品店.
終於找到大減價.
二百多元一張電腦台.
很便宜的 買了回來.
但這張電腦台是要自己組裝的.
買了回來之後.
想著遲一點裝.
等假期才裝.
等沒那麼累才裝.
經過了新年假期.
經過了上年復活節假期.
暑假假期.
聖誕假期.
總是有不同的理由.
藉口讓自己一拖再拖再拖.
拖到什麼時候呢?.
拖到今年三月.
我發覺不可以再這樣拖下去.
再這樣拖下去.
應該拖得煮得再久.
我都不會把台子蓋好.

$^{361}$於是我覺得做一個決定.
今天就蓋好它.
用一整個下午就把電腦台裝好.
但我同樣地.
我們面對上帝的吩咐.
面對祂的說話.
我們有時都會用拖字訣.
我們會想 上帝遲一下吧.
你交託給我的任務.
等我處理好手上的事.
等我有空.
等我沒那麼累.
我就會去回應.
不知道大家還記不記得.
一首歌叫做《明日歌》.
明日復明日.
下一句是什麼?.
明日何其多.
我們再看兩句.
我生待明日萬事成蹉跎.
說到這句很有意思.
西人若被明日淚.
或者有些版本是.
西人苦被明日淚.
春去秋來老將至.
其實很提醒.
明日又有一個明日.
明日是多麼的多.
我一生都在等待明日.
原來什麼事情都沒有進展.
蹉跎了歲月.
如果西人還是這樣.
很辛苦地被明日來拖累.
一天一天地過.
很快就會老去.
同樣地 姐妹.
有些事情.
聖經提醒我們.
不要拖.
面對上帝的話語.

$^{401}$上帝的帶領.
今天聽了就要把握機會去回應.
就要把握機會.
進入上帝安息的應許裡.
來與神同工 與祂同行.
所以作者也因為這樣.
提醒我們.
如何幫助自己.
進入上帝的安息.
他提了兩件事.
第一方面.
在第四章的一到第二節.
他說.
我們既蒙留下.
有進入他安息的應許.
就當畏懼.
免得我們中間或有人.
似乎是趕不上了.
因為有福音傳給我們.
像傳給他們一樣.
只是所聽見的道與他們無益.
因為他們沒有信心.
與所聽見的道調和.
弟兄姊妹原來.
作者很擔心每一代讀者.
他很擔心我們會重蹈.
當年以色列人所犯的錯.
就是聽了上帝的吩咐.
但卻是不肯去順從.
不肯在生活裡.
把上帝的話語實踐出來.
結果雖然有上帝安息的應許.
但卻是進入不到上帝的安息.
所以作者首先要針對.
我們一種態度.
他要求我們畏懼.
你沒聽錯是畏懼.
當然這裡不是說畏首畏尾.
什麼都不敢去試.
什麼都不敢去做.

$^{441}$這裡說畏懼是說.
你要小心謹慎.
去回應上帝給你的應許.
千萬不要覺得是可有可無的事.
因為失去了上帝的安息.
後果作者說是嚴重的.
就正如我想.
在這段時間這幾個月裡.
我們是不是做人都分外.
格外小心謹慎.
你有沒有發覺這幾個月.
你的畏懼心是大了.
你怕什麼.
其實我們很怕感染.
因為你知道感染後.
後果會怎樣.
可以是很嚴重.
甚至有性命危險.
亦會影響你身邊的人.
所以你不會鬆懈.
你不會心存一種僥倖.
覺得沒事的.
我不會這麼幸運中招的.
你不會因為心存僥倖.
所以這段時間.
你不戴口罩就衝出街.
你一定會做足防疫設施.
不單單戴口罩.
還會洗手和人保持距離.
我們為了要減低感染風險.
我們會竭力做好一切.
我們能做的防疫措施.
對不對?弟兄姊妹.
同樣在這裡說畏懼或敬畏.
其實會成為一種推動力.
推動你更加努力.
去遵行上帝的話語.
使我們不敢再怠慢.
輕慢了上帝的道.
亦不敢辜負上帝給我們.

$^{481}$這個安息的應許.
畏懼成為我們一種靠近上帝.
不想失去安息的推動力.
這是第一方面.
作者要我們注意的態度.
另一方面.
作者針對我們一種決心和意志.
他在十一節說.
我們務必竭力進入神的安息.
就是我們要想盡辦法去攻克己心.
想盡辦法去配合.
調教自己的心思意念.
配合上帝的心思意念.
我們知道我們得救.
固然是出於上帝的恩典.
不過上帝要求每一個信徒.
在他成聖的過程中.
真的要下苦功.
不可以偷懶.
不可以孖孖呼呼.
簡單來說就是.
我們信了主之後.
不可以懶惰地過我們信仰的生活.
經文提醒我們.
特別是經文進一步說.
千萬不要想著.
不覺的,沒有人知道的.
沒有人知道我心底裡到底有沒有盡力的.
千萬不要這樣想.
因為在第十二節裡.
經文說我們內心所有的心思意念.
上帝知道一清二楚.
十二節裡這樣說.
我用了新漢語譯本.
因為神的道是活潑的.
是有功效的.
比所有兩刃的劍更鋒利.
甚至可以刺入,剖開.
魂與靈,骨節與骨髓.
也可以分辨心裡的思想和意圖.

$^{521}$弟兄姊妹.
這裡說上帝的道是活潑的.
永遠不會過時.
不要以為勾了.
上帝的道永遠與你我的存活息息相關.
這裡也說上帝的道是有功效的.
他用的字是古代農夫.
用一個篩來篩穀物的動作.
你明白為什麼了.
就是將好的穀物留在篩子裡.
不好的就篩走.
意思是上帝的道能夠將人內心思意念完全分辨出來.
什麼是好什麼是壞.
上帝知道一清二楚.
這裡也說.
上帝的道比兩刃的劍更快.
更利.
就算是骨頭那麼硬.
就算是關節骨髓那麼複雜的構造.
甚至是魂和靈.
這些好像掌握不了的部份.
上帝全都可以掌握到.
上帝知道一清二楚.
並且上帝能夠一一解開.
所以弟兄姊妹.
上帝的道就是這樣壓倒性的.
使我們面對上帝的道的時候.
我們不能夠不謙卑下來.
不能夠不苦心自問.
其實我們有沒有盡力.
我們有沒有竭力去攻克己心.
我們有沒有信服所聽到的道.
還是我們好像昔日的以色列人一樣.
聽是聽了.
不過慢慢地培育了一種硬心.
不肯順從的態度.
事實上這樣是嚴重的.
因為在第十三節裡.
作者更加提醒我們一件事.
這句說話很重要.

$^{561}$「受造之物在他」.
即是在上帝的面前.
「沒有一樣是隱藏的」.
「萬物在他眼前都是赤露闖開的」.
「而我們都要向他交仗」.
這句說話很重要的提醒就是.
無論我們如何收藏掩飾.
自己的心思意念.
原來在上帝的面前是無所遁形.
並且上帝一定會向我們每一個來問責.
我們到底在他心目中.
是否一個忠心的好管家.
有沒有按著他賜給我們個人的恩賜財幹.
來好好地服侍上帝.
進行他的吩咐.
有沒有竭力完成上帝所託付給我們的事情.
所交託給我們的福音使命.
事實上作者之所以說這句說話.
是因為他期望我們.
趁現在就想好將來面對上帝的那一刻.
將來每個人都有這一幕.
當你想到這一幕的時候.
現在就要開始把握機會.
今天你聽了上帝的道.
就要起身回應.
就要努力活出上帝的吩咐.
來竭力與上帝同工.
事實上這段日子.
真的有很多可以與上帝同工的機會.
我們都會發覺.
有很多人在疫症期間.
他們都用不同的方法.
來填補自己心靈裡的洞.
不知道你覺不覺得.
這段時間我想每一個人.
不只是香港人.
世界各地不同的人.
心靈裡都出現了一個洞.
想盡方法去填補.
譬如有人上網去煲劇.

$^{601}$又或者有人學習在家裡煮食.
又或者年輕人可能打機.
打到某隻手提遊戲機斷市.
又或者最近這幾天.
黃金周.
你見到很多人出去踩單車.
出去行山.
又或者迫爆了海灘.
對不對?.
早兩天看到一段新聞.
他說最近美國.
出現了一個前所未有的現象.
他說疫症底下民眾.
因為受困在家裡.
不可以輕易出街.
所以又引發了大量領養.
帶領寵物的需要.
導致各地方的動物收容所.
出現了一個清空的現象.
所有動物領養的.
都被人領走了.
無論是貓,狗,天竺鼠.
甚至乎原來母雞.
都找到新的主人.
有一位從事動物救援機構的創辦人.
他表示.
他說他從事動物救援工作25年.
他未見過這個現象.
讓我們看到.
在這段時間.
確實很多人的心靈空虛.
沒有平安.
所以不同的人都想.
用不同的方法.
去塞回自己心靈裡的洞.
來渴望生活裡.
找回真正的安息.
所以弟兄姊妹.
其實在你和我身邊.
真的有很多.

$^{641}$可以和上帝同工的機會.
我們可以趁著這些機遇.
來竭力幫助你和你身邊的人.
進入上帝的安息裡.
話說我認識一位朋友.
早前他很擔心.
為了他的女兒.
為什麼這樣說呢?.
因為他的女兒之前在美國.
所以到底要決定.
回香港還是留下呢?.
機委又很緊張.
又不知道能否找到機會回來.
又很怕上機的時候.
成全機會感染.
但是留在當地.
也有一個很大的風險.
於是他出現了很多的憂慮.
令到我這位朋友.
精神很緊張.
又睡不安.
又吃不下.
整個人受了一個麻.
我除了和他分析形勢.
到底回來還是不回來.
幫他想想.
回來後如何預備家居隔離.
甚至乎預備什麼物資.
我覺得更重要的是為他祈禱.
讓他知道不單止有人和他同行.
來支援他.
更重要是讓他知道.
其實上帝很愛他.
上帝一定會看顧他和他的家人.
所以我不單止為他祈禱.
我鼓勵仍然未信耶穌的他.
你大膽一點開口去求耶穌.
這位上帝是真實的.
他亦願意看顧你的需要.
他亦會垂聽你的祈禱.

$^{681}$於是乎他亦願意.
自己來開口向耶穌祈禱.
求上帝去憐憫他.
結果上帝亦聽禱告.
保守他的女兒順利回到香港.
亦完成了家居的檢疫.
聖經師讓他經歷到.
上帝給他一種奇妙的平安.
當然我亦順勢將教會.
頌拜的YouTube連結發送給他.
說不定他現在也在看.
我希望他能夠進一步.
認順這位愛他的主.
真的,弟兄姊妹.
上帝預備了很多服侍的機會.
他很期望你和我今天聽了他的禱告.
就去回應.
做一個能夠承擔.
對自己有要求的信徒.
為他真正享受.
聖經所說的豐盛的生命.
耶穌亦是這樣說.
在馬太福音11章28至30節.
我們都很熟悉這段經文.
「凡奴婦擔重擔的人.
可以到我這裡來.
我就使你們得享安息.
我心裡柔和謙卑.
你們當乎我的額.
學我的樣式.
這樣你們心裡就必得享平安.
因為我的額是容易的.
我的擔子是輕生力」.
「額」是說什麼呢?.
在古代是說.
放在牛頸裡.
箍住它.
管轄著它一塊木頭.
讓它可以順服主人.
為主人效力.

$^{721}$耶穌說我們要附他的額.
要順服上帝的吩咐.
為上帝效力.
耶穌說這樣才能享受真正的安息.
確實對現代人來說不容易.
好像沒有什麼賣點.
不太中聽.
每次都被人管住.
什麼都說要順服.
沒有自由.
怎麼會安息?.
怎麼會舒服?.
那就要看我們被誰管束.
事實上這裡說耶穌的額.
就好像我們走行人天橋上.
兩邊的圍欄.
我們走行人天橋.
都知道兩邊會有圍欄.
你有沒有走過天橋沒有頂?.
只有兩邊.
有沒有走過?.
感覺沒那麼安全.
實際上是安全一點.
有沒有走過天橋.
連兩邊都沒有欄?.
只有平板.
你覺得有圍欄好一點.
還是沒有圍欄好一點?.
兄弟姐妹.
其實你有沒有想過.
如果一條天橋沒有兩邊的圍欄.
你覺得走得闊了.
走得窄了.
還是走得自由了?.
兄弟姐妹.
其實如果沒有兩邊的圍欄.
你走天橋的時候.
你會走得窄一點.
你不敢走那麼窄.
為什麼?.

$^{761}$因為你怕失足.
跌死.
相反如果兩邊有圍欄.
你就可以怎樣?.
你走到最盡最寬.
靠著邊走都可以.
因為你知道有東西保護你.
你知道什麼是你的安全網.
你知道不會跌死.
同樣其實.
耶穌的額.
就好像天橋兩邊的圍欄.
保護著你和我.
其實不是要管束.
合制我們的自由.
相反耶穌的額.
是令你和我活得更自由.
你在耶穌的保護下.
你可以自由地.
按著上帝的心意來生活.
不會怕跌倒.
當你信服耶穌的吩咐的時候.
我們才能夠真正享受到.
什麼叫做自由.
什麼叫做安息.
耶穌是這樣說的.
並且耶穌的應許.
當我們願意付他的額的時候.
他會和我們同行.
最後一節經文.
希伯來書四章十四至第十六節.
這裡說.
我們既然有一位.
已經升入高天尊名的大祭司.
就是神的兒子耶穌.
便當持定所承認的道.
因我們的大祭司.
並非不能體恤我們的軟弱.
他也曾凡事受過試探.
與我們一樣.

$^{801}$只是他沒有犯罪.
所以我們只管坦然無懼地.
來到他施恩的寶座前.
為要得連恤蒙恩為作.
隨時的幫助.
我們還記得.
耶穌在下西馬利園的時候.
其實他也像我們一樣.
經歷過人生的掙扎.
也出現過一些.
不想信服上帝的想法.
不過我們也知道.
他最後勝過這一切的試探.
所以耶穌這位.
現在是高天的大祭司.
他不單只能夠體恤.
你和我在信服上帝過程裡.
那些軟弱和掙扎.
耶穌更加可以幫我們.
去勝過這一切.
因為耶穌就是那位創始成中.
比摩西,比約書亞更加超越那一位.
他一定能夠帶領我們.
完成整個成聖的旅程.
而我們要做的就是.
一路竭力信從上帝的吩咐.
也一路竭力仰望.
耶穌的連署和幫助.
讓我們可以在地上.
開始進入上帝所允許的.
這個安息的旅程.
所以弟兄姊妹.
作為這個港道系列的總結.
我很希望我們可以彼此的勉勵.
就是說我們在得救之後.
我們也要繼續過一個得勝的人生.
我們要去恭取.
上帝已經為我們預備好.
每一吋的美德之地.
就是說我們要恭喊在生命裡.

$^{841}$那些仍然不太信得過上帝帶領.
那些仍然不太信服上帝的地方.
聖經說.
你要靠著耶穌.
幫助你逐一的取回來.
同時也要靠著耶穌.
將上帝託付給你的美德之物.
從撒旦的手中.
一一的搶回來.
使得他們和我們都一樣.
回歸上帝.
一起進入上帝的安息應許裡.
就讓我們用一首詩歌.
我們用一首詩歌來回應.
我們今天所聽到的道.
我們同心來宣告.
我們要靠著耶穌過一個得勝的人生.
\newpage



\section{希伯來書出埃及記 2:1-10}
\label{sec:0WamuVI2aAk}
\textbf{2020年5月10日 崇拜直播｜出埃及記2\_1-10;希伯來書11\_23}
\newline
\newline
連結: \href{https://youtube.com/watch?v=0WamuVI2aAk}{\texttt{https://youtube.com/watch?v=0WamuVI2aAk}} ~~~~ 語音日期: 2020-05-11
\newline
\newline
\hyperref[sec:0sy2FpLq_lk]{\small{< < < PREV SERMON < < <}}
~
\hyperref[sec:index_chronic]{\small{[返順時目]}}
~
\hyperref[sec:y8QuGSs8cQg]{\small{> > > NEXT SERMON > > >}}
\newline
\newline
$^{1}$今日我們所重讀的經訓是.
《出埃及記》二章一至十節.
和《希伯來書》十一章二十三節.
《出埃及記》二章是這樣記載的.
「有一個利未家的人.
娶了一個利未女子為妻.
那女人懷孕生一個兒子.
見他盡美就藏了他三個月.
後來不能再藏就取了一個蒲草箱.
抹上石漆和石油.
將孩子放在裡頭.
把箱子擱在河邊的路壁中.
孩子的子子遠遠站著.
要知道他究竟怎麼樣.
發怒的女兒來到河邊洗澡.
他的侍女們在河邊行走.
就看見箱子在路壁中.
就打發一個婢女拿來.
他打開箱子.
看見孩子哭了.
他就可憐她說.
這是希伯來人的一個孩子.
孩子的子子對發怒的女兒說.
「我去在希伯來婦人中叫一個奶媽來.
為你奶這孩子,可以不可以?」.
發怒的女兒說.
「可以」.
童女就去了叫了孩子的母親來.
發怒的女兒對他說.
「你把這孩子抱去,為我奶他.
我必給你公家」.
婦人就抱了孩子去奶他.
孩子漸長.
婦人把他帶到發怒的女兒那裡.
就作了他的兒子.
他給孩子起名叫摩西.
意思說因我把他從水裡拉出來.
希伯來書11章23節.
摩西生下來.
他的父母見他是個盡美的孩子.

$^{41}$就因著信把他藏了三個月.
並不怕皇命.
今天我們會分享一個港島的題目.
是「動盪中的非凡母親」.
今天是母親節.
很開心今天可以和其他兩位母親一起服侍.
剛才我們用服侍來慶祝母親節.
在這裡祝願每一位母親.
母親節快樂.
願上帝的恩惠厚厚地充滿你們.
圍繞你們.
我想每一個人會否認一件事.
就是在家庭裡母親角色的重要.
我想起我自己.
我有四兄弟姐妹.
我們四個化骨龍.
一級一級地下.
當我四歲時,我的妹妹就出生了.
回想起小時候,媽媽真的很辛勞.
不但白天要上班,要去遮擋衣服.
晚上回來,很多事情要做.
洗衣服.
大家記住,那時候沒有洗衣機.
還要家裡一家七口.
晾衣服,煮飯,掃地,拖地.
我還記得當我們生病的時候.
媽媽面乾水濕地照顧我們.
真的睡不著.
不過我還記得一幕.
有一年,臨近新年.
媽媽說家裡的牆身開始剝落殘舊.
媽媽說不如把它塗了.
但家裡其他人完全沒有反應.
結果,媽媽一個人塗了整間屋.
雖然那時的屋不大.
但難處不是塗油.
難處是她一個人搬開雜物.
滾床,櫃,木椅等等.
你們別看我這麼大份.
以為我媽媽也這麼大份.

$^{81}$我媽媽多年來都是一位.
一百磅多的纖纖女子.
在我心目中.
我媽媽真的幫我們遮風擋雨.
在家裡困難的日子.
無論怎樣艱難.
她都會努力賺錢給我們.
我們家裡沒有交租.
她都四處張羅,遮風擋雨.
所以在我心目中,我媽媽是一個超人.
在這裡,不知道她有沒有在看.
但我祝她母親節快樂.
媽媽的愛是很偉大.
不單單我媽媽.
我相信每一位我們的兄姊妹.
我們都經歷過媽媽的愛.
媽媽的愛可以帶來很大的動力.
在1999年.
在土耳其發生了地震.
當年有18000人死亡.
4萬多人受傷.
在地震災後第三天.
工作人員在瓦礫裡去挖.
在發掘的時候.
挖掘了一些瓦礫.
就看到一個媽媽.
和一個七歲大的女兒在當中.
那個媽媽的姿態是怎樣的?.
大家知不知道?.
她是用雙手像掌上壓一樣.
撐起自己的身體.
為的是擋住一些瓦礫或雜物.
可以傷害到她的女兒.
她就這樣苦撐了這麼多天.
但當救援人員把女兒救出來的時候.
女兒其實已經過世了.
她已經沒有氣了.
當再去救媽媽的時候.
媽媽問女兒怎樣?.
因為救援人員不忍心.

$^{121}$怕她沒有求生意志.
所以就瞞著她.
告訴她女兒還沒死.
當把媽媽救出來的時候.
他們送到醫院.
幫她檢查傷臂的時候.
發覺她的傷臂大部分.
大幅度地骨裂.
大家有沒有試過骨裂?.
你能否想像到那種痛楚?.
她還要在大部分地方骨裂.
當時土耳其報紙的頭條.
寫了這五個字.
「這就是母愛」.
其實今天我們也要在聖經裡.
分享一個媽媽.
她在一個非常動盪的時代裡.
憑著她的愛.
憑著對上帝的信.
她把她的孩子拯救出來.
這個媽媽就是約基別.
她不單是剛才第十節說到摩西的媽媽.
我們知道她還是以色列的一位女先知米利暗的媽媽.
也是以色列第一位大祭司亞倫的媽媽.
讓我們一同去禱告.
求主開放我們的心.
讓我們從這位媽媽的身上.
學習一些屬靈的質素.
我們一起祈禱.
你賜給我們有媽媽.
讓我們去經驗.
在人世間裡.
一些很緊密的關係.
一份很真摯的愛.
求你使用我們,幫助我們.
教導我們如何孝敬我們的父母.
又幫助我們.
可以在我們媽媽聖經裡的角色.
一個重要的人物.
約基別的身上.

$^{161}$我們去學習.
生命裡有什麼屬靈的質素.
以至我們不單單自己.
是一個謹憐於你的門徒.
我們又可以.
牧養到我們屬靈的下一代.
我們就恭敬祂將這個時間交給你.
求主,你與我們同在.
開放我們的心.
去聆聽你的話.
奉主耶穌基督得勝名自祈求.
阿們.
聖經在《出埃及記》第一章裡.
告訴我們.
其實約基別.
他們身處在一個什麼時代.
自從雅各.
因為大饑荒.
帶著他11個兒子到埃及投靠約瑟之後.
本來的70人.
經過了400年之後.
他們以色列在埃及的人口.
已經多到200萬這麼多.
在埃及地.
是屬於一個龐大的力量.
但是好景不常.
在《出埃及記》一章八至十節裡說.
他提到有一個新興的王.
這個新興的王為什麼會這樣呢?.
原來是有一群在亞述的閃米特人.
他們後來統治了埃及的大部分地方.
成為了埃及新的統治者.
所以以前那群埃及人.
經歷過約瑟.
幫助他們渡過了一個大饑荒的約瑟.
他們以前因著約瑟的緣故.
對以色列人有幾分的尊重.
但是新王已經不再認識他們.
反而以色列人的人口膨脹到.
令這個新的法老覺得受威脅.

$^{201}$於是他想了三個計謀去對付他們.
第一就是逼他們去做苦工.
奴役他們.
令他們的生活極其的勞苦.
第二他命令修生婆.
要殺死一切以色列人的男嬰.
第三就是《出埃及記》一章二十二節.
他那裡說到以色列人所生的男孩.
要吊在河裡面.
只是留女孩的活命.
以色列的人在一個動盪的時代.
工作被奴役.
甚至要面對一個被滅族的危機.
法老的苛政一天比一天嚴厲殘忍.
當時的父母要肩負起.
去養育保護兒女的事.
是一個非常非常艱難的事情.
我想到的就是螳臂擋車.
完全毫無招架的力量.
不過一到第二章的開頭.
我們就發覺他將焦點轉到一個家庭裡面.
就是我們今天要分享的主角約基別.
究竟我們從約基別三個的歷程裡面.
去看他如何去面對這麼大的危機.
我們會從他一開始生摩西.
然後他育養摩西.
然後他對摩西的影響.
我們一路去看下去.
首先我們看看.
當約基別生摩西的時候.
他究竟生命裡面有甚麼屬靈的質素呢?.
在《出埃及記》二章一至四節.
我們會發覺.
聖經裡面對這個媽媽是沒有名字的.
我們只知道她是一個利未的女子.
直至《出埃及記》第六章二十節.
我們才知道原來摩西的媽媽叫約基別.
而他的爸爸叫做暗蘭.
其實約基別在聖經裡面.
這個名字出現的次數只是兩次.

$^{241}$但他在這裡卻起了一個非常重要的角色.
摩西的父母都是利未人.
約基別的希伯來文的意思.
是解「耶和華是榮耀的」.
猶太人的名字是很有意思的.
從他們的名字我們可以看到.
他們這個家庭.
是要以耶和華的榮耀為重.
他們要彰顯耶和華的榮耀.
代表了他們對上帝的崇敬.
我相信當時一定人心惶惶.
官兵一定是逐家逐戶地去找希伯來的男孩子.
在一個這麼大的威脅下.
約基別竟然違背法老的命令.
心驚膽跳地將這個男嬰收藏三個月.
為甚麼呢?.
我想不難想像媽媽一定是出於愛.
無論如何都竭力.
希望能夠傳流孩子的性命.
但另一件事我更加相信的就是.
聖經裡提到約基別時.
提到他的信心.
我們看到聖經裡說.
因為看到這個孩子「俊美」.
「俊美」這個字其實是同創世記.
一開始神創造天地的時候.
他連續說了很多次.
他說「神看著是好的」.
「俊美」就是和「美好」這個字是一樣的.
我們知道原來摩西是好的.
是美好的,是神眼中最美好的.
我們可以從《使徒行傳》第七章第二十節說.
史提反在申訴時形容摩西.
他是這樣說的.
「摩西生了下來,神看為俊美」.
我相信當約基別.
或者後來聖經也提到.
他的父母包括暗蘭.
看到這個兒子是「俊美」的時候.
是美好的時候.

$^{281}$他們知道在神的眼中.
在這個孩子身上有獨特的安排.
我們去看希伯來書十一章二十三節.
這裡提到摩西生下來.
因為他看到他是「俊美」的孩子.
因著信把他藏了三個月.
就是因為父母對神的信心.
他們竟然不怕皇命,冒死.
要將來可能成為神所器重的摩西保存下來.
我們知道他們的父母.
很知道這是神特別美好的旨意.
相信神會保存他的性命.
去完成神要在他身上的計劃.
正正是因為這個信念.
他們就有勇氣去違背法老的命令.
其實在肥勒比書一章二十八節說.
「凡事不怕敵人的驚嚇.
這就說明他們的沉淪.
你們的救是出於神」.
對神有信心的表現.
就是他們對於敵人不怕.
他們相信神的旨意是最美好.
違抗敵人的命令.
脅於敵人,不是對上帝有信心的人的表現.
我不知道大家在最初疫情開始的時候.
有沒有看《牧者》的家書.
我記得我自己曾經分享過.
因為我兩個兒子正值青春期.
他們的變化很大.
不像以前在小學時期那麼聽話.
我還發覺他們上到中學的時候.
朋輩的影響或周圍文化對他們的影響很大.
大家知道疫情期間.
我們都有很多機會在家工作.
工作多了,會怎樣呢?.
就會磨擦得多.
當我每天工作已經很花心心神.
還要處理家務,煮飯.
還要因為上網學校的功課不太能跟進.
已經身心疲乏.

$^{321}$又沒有耐性,衝突多了.
關係變得很緊張.
我看到他們時下青年有些壞習慣.
不是每個青年都有.
可能是罵人的口不收.
心裡就覺得很沮喪.
我跟丈夫說,我真的要捱到什麼時候.
我跟上帝說,神啊,求你把這個苦杯拿開.
說真的,我真的有點怕了他們.
於是,我最初就用了一個辦法.
就是幫自己拿開苦杯.
怎樣幫自己拿開苦杯呢?.
就是上班,上得就上.
不用見到他們,不用面對他們.
不用有衝突,就不用那麼苦.
不過,正如我在《牧者家書》提到.
當我安靜在神的面前.
神的話語再一次提醒我.
每一個孩子的身上.
神都有祂獨特的計劃.
我們要好好在這個孩子身上.
做好做母親教導的工作.
開始我就轉念了.
我開始不是很怕跟他們吵架.
而是我反而很珍惜.
有得跟他們在青少年期相處的機會.
可以透過這些機會.
我能夠好好地將神的真理跟他們分享.
我不需要借用外面的力量.
我不需要用工作去逃避.
我每一天早上起床禱告.
倚靠上帝求主教導我.
怎樣去教導我的兒女.
兒子,母女.
我相信上帝在他們身上.
必有祂的美意.
我看的不是今天的祂.
我看的是神塑造將來的祂.
所以憑著對上帝有的信心.
摩西的母親約基別.

$^{361}$一直藏著這個兒子三個月.
她的哭聲越來越大.
收不起來.
於是她就想.
怎樣用方法去保存摩西的性命.
聖經說她用蒲草葉柄造了一個盒.
蒲草可以造成堅韌的紙.
其實是一個紙盒.
因為紙盒不防水.
她就塗了一些石漆和石油.
其實是一些從植物提煉出來的油脂.
目的是有防水和隔熱的作用.
然後她將它放在箱子裡.
然後拿去河邊.
放在爐底旁邊.
將它卡住.
有兩樣目的.
爐底可以防熱.
卡住的時候.
就不會那麼容易被水沖走.
在一個再不能隱藏的時候.
如果被冰冰捉到.
摩西的性命會凍過水.
所以約基別想了一個.
非不得已的方法.
就是將它放在河邊.
在河水上.
其實我們覺得.
為甚麼可以這樣呢?.
但你想一下.
不難想像.
如果我們像以前香港.
有些父母.
無法再撫養孩子.
又或者基於種種原因.
他已經不能再承擔.
撫養孩子的責任.
會怎樣呢?.
可能將它放在孤兒院.
或者醫院門口.

$^{401}$他希望一些有心人.
一些善心的人.
可以接他回去.
有一個更適合的環境.
有一個他可以生存的機會.
去給自己的孩子.
約基別在這個非不得已的安排下.
做了一個這樣的決定.
我們看到一個淑齡的母親.
她對神的信心.
她不怕自己失去生命.
她只想孩子的命可以保存下來.
她知道.
聖經告訴我們.
她倚靠神.
神去帶領她,幫助她.
除了信心之外.
我們看到約基別很有智慧.
在《出拉喀記》二章五至十節.
我們看到.
法老的女兒.
來到河邊洗澡.
各位弟兄姊妹.
你有沒有想過.
為什麼法老的女兒會去河邊洗澡?.
為什麼她不在宮殿裡?.
宮殿應該非常豪華,輝煌.
設備應該很好.
為什麼她去河邊洗澡?.
其實我們估計.
是因為宗教禮儀的緣故.
因為在七章十五節.
和八章二十節.
都有兩次記載到.
法老早晨會去河邊洗澡.
所以有些聖經學者都會認為.
約基別設計.
幫助摩西到河邊洗澡.
放置她.
在那個時候是有原因的.

$^{441}$可能是法老的女兒.
會在那個時候出現.
當法老的女兒看到.
這個蒲草箱的時候.
她就奮赴疲累.
去拿這個箱子.
一看.
是一個希伯來人的嬰兒.
那你怎麼知道她是希伯來人呢?.
我相信可能從她的外貌.
她的衣服.
她身上的裝飾.
又或者.
其實可能嬰兒.
有一樣很重要的東西.
就是她受了國禮.
她就知道這是希伯來人的嬰兒.
但是你會想.
法老的女兒.
她的爸爸不是頒佈了皇令嗎?.
不是凡是有希伯來人的嬰兒.
都要將她丟在河裡嗎?.
所以現在這是一個很關鍵的時候.
這個時候摩西哭了.
摩西哭了.
令法老的公主憐憫她的心腸.
另一個關鍵就是姐姐的出現.
摩西的姐姐米利暗很機警.
她駕駛到便利商店.
一見到公主的心意.
她就不慌不忙地自動請纓.
告訴公主.
我幫你去找一個奶媽來照顧小摩西.
大家去留意二章第七節.
孩子的子子對法老的女兒說.
我去在希伯來婦人中.
叫一個奶媽來為你奶這孩子.
我們見到米利暗非常機警.
她有三個很重要的特點在這裡出現了.
第一是她很把握時機.

$^{481}$一駕駛到便利商店.
見到公主的樣子.
她就馬上出去.
第二是她為了找一個希伯來婦女來為她奶這孩子.
她的心腸就是為了這位法老的女兒.
第三就是她找了自己的生母.
找了自己的媽媽來做這孩子的奶媽.
大家想想.
米利暗當時不是一個很大的女孩.
我們相信她不夠十歲.
可能八,九歲左右.
然後又察覺到公主那份憐憫的心腸.
然後又有這樣的建議.
其實我們相信.
媽媽的機別在這件事上.
都應該對她有提醒和建議.
於是米利暗就把機別帶到這位埃及公主那裡.
公主就請了她.
還跟她說請你代我餵養這位嬰孩.
即是這孩子已經成為法老女兒的兒子.
結果被遺棄.
只差一步就要進入死亡的摩西.
再次被她的親生母親抱起.
還要被親生母親撫養她長大.
機別不用偷偷摸摸.
她可以公然養.
還可以賺到錢去照顧她的親生孩子.
我們去到第十節會看到.
法老的女兒幫她起名叫摩西.
摩西這個名字應該是埃及名字.
埃及的名字的意思就是生孩子.
但這本聖經的作者希伯來文.
在翻譯這文章時.
把摩西翻譯成希伯來文.
意思是拉出來.
正正就是講述摩西人生裡一個很重要的責任.
就是把以色列人從埃及這個遺牢之地.
拉出來拯救他們.
我們看到米利暗和約基別的智慧.
他們按照上帝的心意.

$^{521}$精心的策劃.
目的就是幫這個孩子摩西.
成為拯救以色列人出埃及的重要領袖.
成為民族的英雄.
我們真的要好好向這個屬靈的母親學習.
在1920年時.
史太林下令將前蘇聯一個城市.
叫斯塔羅波的城市的基督徒.
和所有聖經全部消除.
所有基督徒捉去勞改.
如果不肯放棄聖經的人.
更加要被即時處決.
在很多年後.
蘇聯解體.
有一個美國教會的報道團去到這個地方.
但他們去到之後.
好像要在莫斯科運來的聖經.
都還未運到.
於是他們就問人怎樣可以找到聖經.
一位年老的長者.
他就告訴他們.
以前史太林就是將聖經收押在倉庫裡.
於是他們就祈禱.
政府也開路.
讓他們在那裡取得勝經.
第二天他們租了幾輛卡車.
還找了當地的大學生和地方專員.
陪他們到這個地方取得勝經.
大部份的人都很忙碌地去搬運聖經的時候.
突然有一個年青人坐在那裡.
一筆成聖.
大家都覺得很驚訝.
這個年青人坐在那裡拿著聖經.
不停地飲入.
誰知一打開.
令他很吃驚的就是.
原來這本聖經的第一頁.
就是他祖母的簽名.
課倉裡有千千萬萬本聖經.
但他竟然在隨手的裡面.

$^{561}$拿起了他祖母親手寫給他簽名的聖經.
當他再翻閱的時候.
他發覺聖經裡面夾了一些很重要的祈禱文.
這個祈禱文就是他祖母為了他去懇切禱告.
希望他的孫兒有一天能夠順處被神使用.
難怪他哭得這麼厲害.
後來這個年青人不單止信了耶穌.
他還將這個美好的見證向全世界的人分享.
屬靈的母親對上帝的話是堅持的.
她對上帝的真理是堅持的.
而且他們會在艱難的裡面.
仍然不住地為他的兒女去禱告.
有時我都會去反省.
我都會去想一想.
其實兒女並不是我們的擁有物.
我們到底有沒有用屬靈的智慧去養育我們的下一代呢?.
究竟我們的子女是屬於誰呢?.
是上帝還是我們自己呢?.
如果我們當我們的子女是擁有物.
我們就會為他鋪搭一條我們認為最好的道路.
但這樣會弄錯.
唯有我們將我們的兒女清楚地知道.
他是屬於上帝的.
我們要鼓勵他走在上帝的道路上.
這才是真真正正屬靈的智慧.
我們要記住.
我們是兒女的托管者.
不是擁有者.
兒女是屬於上帝的.
接下來我們會看看約畸別的影響.
它對於摩西有甚麼影響呢?.
我們從希伯來書11章23至28節可以看到.
大家有沒有想過.
摩西在他頭四十年裡.
他一直在埃及的皇宮生活和成長.
他受的是埃及的教育.
他經驗的是埃及的文化.
他可能會不認識耶和華.
他更加不認識以色列的文化和同胞.
但希伯來書11章24至28節.

$^{601}$那裡很清楚地說.
摩西因著信長大了.
就不肯稱為法老女兒之子.
他不要埃及皇子這個身份.
他寧可和臣的百姓同受苦害.
也不願暫時享受最終之樂.
他看為基督受的靈育比埃及的財物更寶貴.
因他賞望所要得的賞賜.
他因著信就離開埃及.
不怕黃路.
因為他恆心忍耐.
如同看見那不能看見的主.
他因著信就守逾越節.
行灑血的禮.
免得那滅長子的臨近以色列人.
希伯來書對於摩西的信心有很詳細的描述.
為何摩西可以如此有信心?.
為何摩西可以倚靠這位以色列的神?.
其實我們很深信是約基別.
或者他的家庭對他的影響.
他的教導和榜樣影響了他一生.
我們知道約基別在河裡救了摩西.
被救出來的時候.
他就開始做他的奶媽.
奶媽其實不單是餵奶這麼簡單.
在《創世記》24章59節.
當尼伯加出嫁的時候.
他的奶媽底波拉是隨著他一起出嫁.
你要想想出嫁的女孩子不是嬰兒.
不是孩子.
是一個出嫁的少女.
他仍然有奶媽去教導他.
所以奶媽除了餵奶.
做保母的身份之外.
其實也有家庭教師的意味.
將猶太傳統對上帝的敬畏去教導.
新教研究他的兒女.
約基別很有可能在他未入皇宮之前去教導他.
當他入了皇宮之後.
也有機會在那段日子裡.

$^{641}$慢慢去建立他對於民族和信仰的觀念.
他很清楚地告訴他.
你不是埃及人.
你拜的不是埃及的神.
你是希伯來人.
你的上帝是耶和華.
所以當他四十歲見到他的同胞.
他的同胞不是埃及人.
他的同胞是希伯來人.
被埃及人欺壓的時候.
他就將埃及人創手打死了.
所以我們很清楚地看到.
屬靈的母親還有一個很重要的質素.
就是她是一個屬靈的榜樣.
唐崇明牧師在2007年分享了一個見證.
他說他們全家人自小就信主.
三歲的時候爸爸離世.
他就由一位很敬虔的媽媽撫養他成人.
唐牧師的記憶很深刻.
他說七,八歲的時候.
每天起床最先聽到的是媽媽的祈禱.
再大一點.
媽媽每天捉住每一個小朋友.
和他們每天祈禱一個小時.
然後才送他們上學.
這位那麼愛主的媽媽.
就成為了唐崇明很好的榜樣.
成為了他屬靈的啟蒙老師.
所以弟兄姊妹.
若基別是一個非凡的母親.
她三個兒女在屬靈方面都有非凡的成就.
我們從她的身上看到.
要有對上帝充足的信心.
不依靠其他東西.
只信靠上帝.
她有從神而來屬靈的智慧.
去行上帝的旨意.
而且她自己的生命都有美好的榜樣.
我們看到若基別的背後.
有一位非凡的上帝.

$^{681}$神很奇妙地去保守以色列人.
在出埃及往後的發展裡.
我們會看到.
摩西之所以可以成為一代那麼重要的領袖.
就是這位非凡的上帝奇妙的作為.
我記得我曾經和大家分享過一個媽媽的故事.
是前年母親節時說的.
我還未說完.
我今天很想和你分享.
那個媽媽因為她的丈夫.
在工作裡面很勞碌.
半夜突然離世.
她一毛錢都沒有留給這個太太.
她留下的是九個兒女.
所有身邊的人都勸她.
把這九個兒女送給其他人養.
她很掙扎.
她在晚上等孩子熟睡的時候.
她跪下哭了很久.
向神祈禱了很久.
她打開了聖經.
她看到耶利米書49章第11節.
她說:你撇下孤兒.
我必保存他們的性命.
你的寡婦可以倚靠我.
她於是就和上帝說:上帝啊.
我知道了.
是你將這些孩子給我的.
若我作為母親的責任.
你就成為他們的父親.
而這位含辛茹苦堅持去養九個孩子的母親.
就是一位偉大的報道家.
無敵的媽媽.
神的愛使媽媽很有力量去養活孩子.
故事來到這裡還未說完.
因為神不單單使用媽媽.
在1855年的時候.
有一位叫愛德華金寶的先生.
他在波士頓一間教會教主日學.
他經常為上主日學的那群頑皮的孩子.

$^{721}$感到非常頭痛和煩惱.
他曾經一度很想放棄這群孩子.
但他最終決定為他們祈禱.
並且求神幫他一個一個帶領他們信主.
在班上有一個賣鞋的孩子.
好像不太明白救恩的道理.
於是他特意去賣鞋的地方.
在倉庫裡跟這個孩子說上帝奇妙的救恩.
並且他必須自己和上帝建立個人的關係.
不是藉著媽媽.
是要有個人的關係.
帶領他們認罪,悔改,祈禱.
這個人就是慕迪.
接著慕迪在芝加哥的一間大學裡.
栽培了一個大學生.
他叫做查普曼.
查普曼後來跟隨著慕迪.
成為了一個暴徒,暴道家.
而查普曼在一次他舉行的暴徒會裡.
他帶領了一個職業的棒球員.
叫做Billy Sunday去信主.
他毅然放棄了棒球職業生涯.
加入了這個暴徒團和他一起暴徒.
後來查普曼受聘於一間教會.
他沒有再做暴徒團.
Sunday就承接他的棒球.
繼續去暴徒.
他暴徒之後又帶了一個男孩.
叫做Ham信主.
這個男孩信了主之後.
又是很有活力,很有暴徒的熱心.
他竟然租了一輛棺材車.
開著棺材車到處巡遊,報道.
他去到一個北卡羅來納州的地方.
他在一個高中對面街的屋裡舉行暴徒會.
有一個男生他知道了之後.
立刻趕快排完蛋就走去參加這個暴徒會.
不是他對信仰有興趣.
而是他想去賭蛋.
他就是想去搞亂.

$^{761}$但當他去聽這個暴徒會的時候.
Ham的暴徒會的時候.
他被上帝的真理震懾了.
他沒有賭蛋.
到第二天他繼續去參與這個暴徒會.
就在那個暴徒會裡決志了.
這個男孩叫做Billy.
你知不知道他是誰?.
他就是葛培里.
這個近代偉大的報道家葛培里.
他成為了舉世知名的報道家.
但你知不知道.
如果我們看回這一連串的事件.
我們會看到.
上帝的道一個一個.
忠誠地去輔託給他的僕人.
去傳遞這個訊息.
源頭是來自一個卑微到不得了的寡婦.
無敵的媽媽.
而且還是一位堅持不懈的主日學老師.
所以各位親愛的弟兄姊妹.
無論你是媽媽.
或者你在教會裡.
你都承擔起教導的功夫.
是人家屬靈的導師.
我都勸勉大家.
要有充足的信心.
要有屬靈的智慧.
要成為好的榜樣.
去將上帝的恩典一直傳下去.
讓我們一起祈禱.
你是何等的愛我們.
你讓我們得著這份誘因.
主啊我求你幫助我們成為一個屬靈人.
在我們的生命當中.
對你有絕對的信靠.
在我們日常生活裡.
不靠其他.
而是靠著你.
因為我們深信在你裡面.

$^{801}$有大能大力.
我們深信你是親自去成就救恩.
我們深信每一個人的生命.
都有你完美的計劃.
求你使用我們.
賜福給我們.
對你繼續堅持有信心.
也將你的愛傳開.
我們這樣禱告交託.
奉主耶穌基督的聖名致祈求.
阿們.
\newpage



\section{約書亞記 11:16-23}
\label{sec:y8QuGSs8cQg}
\textbf{2020年4月5日 崇拜直播｜約書亞記11\_16-23}
\newline
\newline
連結: \href{https://youtube.com/watch?v=y8QuGSs8cQg}{\texttt{https://youtube.com/watch?v=y8QuGSs8cQg}} ~~~~ 語音日期: 2020-05-17
\newline
\newline
\hyperref[sec:0WamuVI2aAk]{\small{< < < PREV SERMON < < <}}
~
\hyperref[sec:index_chronic]{\small{[返順時目]}}
~
\hyperref[sec:XMPoIkBynFQ]{\small{> > > NEXT SERMON > > >}}
\newline
\newline
約書亞記 11:16-23
\newline
\begin{longtable}{cl}
\hline
\hline
章節 & 經文 (和合本修訂版)\\
\hline
11:16 & \begin{tabularx}{0.7\textwidth}{X} 約書亞奪了那全地，就是山區、整個尼革夫、歌珊全地、低地、亞拉巴、以色列的山區和山下的低地， \end{tabularx} \\ \\ \relax
11:17 & \begin{tabularx}{0.7\textwidth}{X} 從上西珥的哈拉山，直到黑門山下面黎巴嫩平原的巴力‧迦得。他擒獲了那裡的眾王，把他們殺死。 \end{tabularx} \\ \\ \relax
11:18 & \begin{tabularx}{0.7\textwidth}{X} 約書亞和這些王作戰了很長的一段日子。 \end{tabularx} \\ \\ \relax
11:19 & \begin{tabularx}{0.7\textwidth}{X} 除了希未人基遍的居民之外，沒有一城與以色列人講和，都是以色列人作戰奪來的。 \end{tabularx} \\ \\ \relax
11:20 & \begin{tabularx}{0.7\textwidth}{X} 因為耶和華的意思是要使他們的心剛硬，來與以色列人作戰，好使他們全被殺滅，不蒙憐憫，反被除滅，正如耶和華所吩咐摩西的。 \end{tabularx} \\ \\ \relax
11:21 & \begin{tabularx}{0.7\textwidth}{X} 那時約書亞來到，剪除了住山區、希伯崙、底璧、亞拿伯、整個猶大山區和以色列山區的亞衲族人。約書亞把他們和他們的城鎮盡都毀滅。 \end{tabularx} \\ \\ \relax
11:22 & \begin{tabularx}{0.7\textwidth}{X} 以色列人的地中沒有留下一個亞衲族人，只有一些還留在迦薩、迦特和亞實突。 \end{tabularx} \\ \\ \relax
11:23 & \begin{tabularx}{0.7\textwidth}{X} 這樣，約書亞照著耶和華所吩咐摩西的一切話奪了那全地，就按著以色列支派所得的份把地分給他們為業。於是國中太平，沒有戰爭了。 \end{tabularx} \\ \\
[1ex]
\hline
\hline
\end{longtable}
$^{1}$在這個時間我們一起來聆聽一段經文.
這段經文是今天講到的主題.
邁向新世代 征戰不再 挑戰仍在.
經文是記載在約書亞記第十一章.
第十六節至到第二十三節.
請聽我讀出以下的經文.
約書亞奪了那泉地 就是山地.
一大藍地 戈山泉地 高原 阿拉巴 以色列的山地.
和山下的高原.
從上西爾的哈拉山.
直到黑門山下尼巴嫩平原的巴力加德.
並且擒獲那些地的諸王 將他們殺死.
約書亞和這諸王征戰了許多年日.
除了欺騙希美人之外.
沒有一成與以色列人講和的.
都是以色列人征戰奪來的.
因為耶和華的意思是要叫他們心裡剛硬.
來與以色列人征戰.
好叫他們準備殺滅 不蒙憐憫.
正如耶和華所吩咐莫世的.
當時約書亞來到.
將住山地 希伯倫 底碧 阿拉巴 猶大山地.
以色列山地所有的亞納族人剪除了.
約書亞將他們和他們的城邑盡都毀滅.
在以色列人的地沒有留下一個亞納族人的.
只在加薩 加德和阿薩特有留下的.
這樣約書亞照著耶和華所吩咐莫西的一切話.
奪了那全地.
就按照以色列支派的宗族.
將地分給他們為業.
於是國中太平沒有征戰了.
今日的經文是特別記載著約書亞帶領一班弟兄姊妹.
帶領整個以色列的軍隊.
一齊進駐迦南地有好幾年的征戰.
現在讓我們先一齊祈禱.
祈求上帝的靈與我們同在.
讓我們明白聖經的話語.
我們同心祈禱.
天上阿爸父願你對我們每一個講說話.
藉著今天約書亞.

$^{41}$怎樣帶領以色列整隊軍隊進入迦南地.
這幾年的日子上帝你怎樣去因領他們.
求你用同樣的話語來勉勵我們.
讓我們今天學習約書亞整個軍隊一樣.
讓我們帶著謙卑信服的心.
帶著恭昂的鬥志.
一齊邁向新的時代.
願你在此刻當中.
你對我們每一個講說話.
不論在家中聆聽的.
或者在教會當中宣講的.
求你的聖靈與我們同在.
幫助我們眾人.
我們將此時此刻交託給你.
願你賜福.
欣賞我們祈禱奉基督耶穌得勝的聖名祈求.
阿們.
今天我們一齊來看.
我的講章.
我們今天一齊來看這一段經文.
首先講到.
如果你過去連續幾次來聆聽.
我們知道.
約書亞一直有經過不同的戰役.
一齊來到去爭戰.
去到第十一章和第十二章的尾.
其實某程度上那些個別最重要的戰役.
已經是記載完畢.
在第十一章開始.
就是為了整個的戰事.
這幾年的戰事作一些總結.
在第十一章第十六節.
到第二十三節.
如果在地域上.
就是由南至北.
約書亞的軍隊怎樣去打仗.
而在第十二章.
第一至到第二十四節.
就會記載著.
由河東河西.

$^{81}$很容易記的.
這一段是講南去北.
下一次的章節.
就是由河東至到河西.
在地域上.
今天特別的來到去記載.
講到兩個很重要的敵人.
在經年累月怎樣去打勝他們.
首先在第十一章.
第十六節.
到第二十節就講到.
約書亞和整個軍隊.
怎樣去打勝.
迦南地裡面.
南至北的這些諸王.
在經文裡面特別記載著.
在這一班的王.
他們是聯合成為一個聯軍.
來攻打他們.
不過約書亞一次又一次.
這樣戰勝了他們.
那裡講.
以色列奪了那全地山地.
一帶南地.
高原.
阿拉巴.
以色列的山地和山下的平原.
從上西爾的哈拉山.
直到黑門山.
下的黎巴嫩平原.
的巴力加特.
在這裡面我們見到.
地域其實由南至北.
幅員是很廣的.
更加重要的.
他講的不只是地方.
來到去打勝了.
他更加要講到.
很多諸王.
所有的王.

$^{121}$他們就在約書亞的軍隊.
在上帝的帶領之下.
一一的來到去戰勝.
值得留意的就是.
在經文裡面特別記載.
約書亞.
或者整個耶和華的軍隊.
與諸王爭戰了許多的年日.
爭戰了許多年日.
有多少年日呢.
其實在經文裡面沒有特別講的.
不過有些聖經學者.
就去做一些計算.
大概在.
爭戰的年日.
大概有五至七年.
有這五至七年裡面.
綿綿不絕的.
打仗.
一場又一場.
一場比一場更加難打.
就是一個很長的.
韌力戰.
以致他們就可以打贏.
這些聯軍的軍隊.
值得注意.
除了是講那個時間之外.
在五至七年裡面.
其實其他外邦 其他諸王.
都會見到耶和華軍隊的作為.
特別是講到.
很多的.
打仗 行軍裡面.
致勝的秘訣不是在講.
以色列的軍隊有多精良的武器.
或者他們純粹.
在戰術上的運用.
在他們當中裡面.
他們知道.
有一個神叫做耶和華上帝.

$^{161}$與這一隊的軍隊.
與約書亞同在.
以致到他們戰無不勝.
這些諸王見到.
整個以色列軍隊.
打贏仗之後.
他們有什麼反應呢.
經文裡面講.
沒有一城.
與以色列人講和.
只有欺騙的希美人.
跟他們講和.
其他的城 其他的王.
見到他們打贏仗.
軍隊的名來到打仗的時候.
他們都不會.
來到為了這些戰勝.
而有一些.
改變自己的立場.
在這裡面 在近東當時.
軍隊 政治.
和宗教.
三者是一起捆綁在一起.
講和的意思.
不是純粹講政治性 你不要打我了.
算了 講和的意思.
不單止是講.
他們在軍事上放棄.
不再進攻.
更加重要的就是講和.
歸順敵方.
所敬拜的神.
他們在政治上臣服.
在宗教上.
他們都是臣服在.
打贏了的那一方.
或者耶和華軍隊.
帶著耶和華的名號.
這班人似乎.
繼續來拒絕耶和華上帝.

$^{201}$只有欺騙的希美人.
他們願意歸順.
所以在經文裡面.
我們見到之前他們不單止是講和的.
他們還要敬拜上帝.
在殿裡面.
來施跡 原因就是他們.
連自己的信仰.
都改換.
不過在迦南諸王裡面.
他們沒有人願意.
成為一個講和的.
團隊.
反而他們繼續來群起攻之.
他們希望藉著.
他們自己所保佑他們的.
那個神明.
希望幫他們打勝這一場仗.
不過經過五至七年的日子.
其實很足夠時間驗證.
他們的神.
是沒有辦法幫助他們.
只有耶和華上帝可以保護.
和帶領以色列的軍隊.
一齊來打贏.
一次又一次的勝仗.
在經文裡面.
特別在第二十二節.
特別的.
來到.
第二十節裡面.
講到耶和華的意思.
是要使那些.
諸王的心.
剛硬.
來與以色列人作戰.
好使他們全被殺滅.
不蒙憐憫 反被除滅.
正如耶和華所吩咐摩西.
有一個字眼.

$^{241}$我們讀聖經很少會看到這個字.
就是不蒙憐憫.
我們看到這個字.
我們會覺得上帝好像.
毫無懼怕.
覺得這隊軍隊如果不歸向耶和華.
他們就好像面對著.
殺身之禍.
不過和合本採取的翻譯.
就講到上帝是一個很硬的上帝.
但是如果你看仔細.
去看經文的上文.
剛才我們所講的.
願意講和.
其實有五至七年就給那些諸王.
看到上帝的作為.
其實上帝製造了一些空間.
讓其他的諸王可以臣服.
另外一個字眼需要注意.
就是不蒙憐憫.
有些譯本.
或更多的聖經學者.
覺得這個字眼.
和合本未必是最精準的翻譯.
更加好一點.
或更多的翻譯.
就是講不向耶和華懇求.
因為憐憫這個字.
在聖經裡面.
大部份使用都是懇求的意思.
如果你用不向耶和華懇求.
你就會明白.
其實在這五至七年裡面的戰事當中.
這些諸王的心是很硬的.
他們堅決地拒絕耶和華上帝.
他們不向大能的主懇求.
求主施恩.
在他們的國民當中.
祈求神臨在他們.
將禍患改變.

$^{281}$這些諸王他們不是這樣想的.
這些諸王他們的心繼續的去僵硬.
甚至他們不會繼續的去尋求耶和華上帝.
在這樣的情況下.
這就好像一個惡性循環.
這班人的心越硬.
上帝的作為就越強力.
這就好像一個惡性循環.
他們的心越硬的時候.
上帝所給的工作.
上帝所帶領的戰事.
就會越來越轟烈.
這些諸王到了最後.
就要面對著被滅整個城的危機.
還有另外一個字眼.
耶和華使他們的心剛硬.
剛硬的字.
Kazak.
這個字其實是說一種心態.
其實在出埃及記裡面.
摩西直著.
上帝透過摩西來做這個十災.
在埃及發怒的時候.
發怒的心就像這個字.
心是硬的.
拒絕上帝.
不容讓耶和華上帝成為他的主.
或者在他生命當中.
懇求這個主來施恩.
這種剛硬.
使他們面對著這些諸王.
面對一個很大的禍患.
剛硬這個字其實不單止用在諸王身上.
剛硬這個字也是用在約書亞身上.
不過用在約書亞身上.
他就不是用剛硬這個字.
即是那個譯法.
而是在一章第六至第九節裡.
講到剛強壯膽.
壯膽那個字.

$^{321}$就是這個剛硬那個字.
面對著一個很剛硬的諸王.
上帝興起他的僕人.
他就要堅心地.
面向著一個剛硬的世代.
面對剛硬的那些人.
帶著無比的鬥志.
同樣約書亞也帶著上帝給他的那種硬朕.
來打贏這一場又一場的勝仗.
在這裡我們稍微思想一下.
約書亞這五至七年裡.
其中一個致勝的秘訣是甚麼呢?.
除了上帝與他同在之外.
在人的作為當中.
在人的整體裡.
更加重要的就是有一種堅決.
抵擋心硬的人的那種態度.
不過很多時候在我們信仰裡面的實踐.
有時候我們會失去了這一份鬥志.
甚至有時候我們面對著一些挫敗.
或者艱難.
我們就很容易為自己.
或者為身邊的東西作出一個很負面的結論.
甚至有時候講一些晦氣說話.
以致失去了自己那種鬥志.
來將福音或者將上帝的祝福.
藉著我們的生命來展現出去.
有一次在我家裡面有一個這樣的經歷.
家裡有人病患.
都很危急的.
早幾個星期之前.
很危急的時候.
家裡人就安排我們的長輩進了醫院.
家裡的長輩其實都在那裡面有生命的危險.
當他送了進醫院之後.
其他家裡的親友就打給我.
希望我一起祈禱.
一起祈禱的時候.
祈完禱分擔完之後.
那位親友就特別說.

$^{361}$開始就說很擔心.
很多困擾.
那個長輩又未信耶穌.
他自己又好像有很多困難.
數的時候就將所有千年過去的重擔.
一字過數算出來.
當然我在那裡聽.
聽到差不多的時候.
數了一個小時.
都是講一些跟我長輩患病的事情.
沒有關係的事情.
但是講了一些很大的.
他自己很多的困擾.
或者他自己心裡面的鬱結.
給他一個小時聊完之後.
我就跟他說.
親友說說這些說話夠了.
然後我提醒他.
我說不過現在我們家裡的這個長輩.
他正在打仗.
他面對著這個病.
我們一家人要陪他一起打這個仗.
要打贏這場仗.
如果我們太多緬懷自己的過去.
或者覺得自感形傷.
覺得自己很可憐.
很沒有力量的時候.
身邊的人會受到影響.
更加重要的是我們要鼓勵這個長輩.
打贏這場勝仗.
治好這個病.
現在是打仗的時間.
而不是繼續去緬懷.
或者那種感覺很傷感的時間.
雖然我講的那些說話很不客氣.
但是我的親友其實都是我的長輩.
我跟他說完這番說話之後.
他就說我明白了.
我知道了.
我們就繼續這樣來祈禱.

$^{401}$我也用說話來勉勵.
我在病床裡面的長輩.
隔了兩三天之後.
病就慢慢好起來.
醫生也說很奇怪.
為什麼會這麼高士氣.
婆婆.
為什麼你這麼快就好像.
漸入佳境.
然後有一些家人就說.
因為有上帝.
有人為我們祈禱.
我們要打贏這場的仗.
弟兄姊妹.
今天我們面對生活所有不同的艱難.
面對疫情也好.
經濟下滑也好.
或者生活裡面.
林林種種不同的艱難也好.
這些東西是很兇猛的.
撲向我們的生命當中.
我們只不過是不要過份的.
來將我們有一種負面的感受放大.
以致我們失去了那種戰鬥的決心.
當有困難來的時候.
我們就要帶著上帝給我們的那種的綱領.
那種的堅強.
來打贏這場的仗.
如果我們沒有辦法來到這樣的盡力的時候.
很多時候我們還沒有打仗.
已經被外面的環境嚇怕了.
或者令到我們自己好像破了氣一樣.
我們沒有辦法繼續前行.
在生活的艱難也是這樣.
在傳福音的工作也是這樣.
傳福音一點都不是一件很容易的事.
是不是因為很多的艱難.
我們就放棄.
就是因為傳福音是這麼艱難.
所以成功.

$^{441}$大人信主的那種喜樂是很大的.
如果所有東西都是手到拿來.
很容易去實現的時候.
我們那種喜悅不會真的很大.
但是正正在逆流而上.
充滿艱難的時候.
我們靠著上帝的恩典.
也都會堅心地去實踐的時候.
上帝給我們的那種喜樂是很大的.
我們更加分享到那種.
真是屬天的那份喜悅.
約書亞就是帶著那種剛強壯膽.
就是給上帝給他的那種硬朗來打仗.
相反來說.
如果這種硬朗是放在上帝敵黨裡面.
帶來的就是生命裡面.
拒絕上帝那種恩典的處境.
這個是第一點.
約書亞軍隊帶著那種無比的堅毅.
打贏這五至七年的仗.
接下來第21至第23節.
當時約書亞來到將駐山地希伯倫抵迫.
阿拉伯,猶大山地.
以及山地所有的阿納族人剪除了.
約書亞將他們和他們的城邑盡都毀滅.
在以色列人地沒有留下一個阿納族人的.
在阿薩加特和阿薩特樓下的.
這樣約書亞照著耶和華所吩咐摩西的一切.
說奪了那全地.
就按著以色列支派的宗族.
將地分給他們為業.
於是國中太平.
沒有征戰了.
在這裡特別說一個民族.
一個族群叫阿納族.
特別在經文裡面說到.
要迎戰阿納族整體的族人.
阿納本身的意思就是頸項.
為什麼要用頸來形容民族呢?.
其實已經失傳了.

$^{481}$不過很多色經學家估計.
因為這些巨人的頸項很粗大.
如果我們看過拳擊就知道.
頸越粗代表越強壯.
阿納族可能是一群很強悍的.
頸項很粗的.
很有力量的一群民族.
更加在《申命記》裡面第二章十至十一節.
特別在摩西時代已經記載著.
這個阿納族人已經和以色列族曾經面對面.
經文裡面描述的阿納族人叫利梵音人.
利梵音其實是一個譯音.
利梵音的意思就是巨人.
意義就是巨人.
譯音就是利梵音人.
這個巨人其實一直在聖經裡面.
由《創世記》已經出現了.
這些巨人的頸項很粗.
充滿力量的宿敵或敵人.
他們帶著很強悍的戰鬥能力.
特別注意就是這一群阿納族人也好.
或者利梵音人也好.
通常聖經記載不是單獨一個單丁.
全部一定是用眾數來描述利梵音人.
所以你會看到.
以色列面對的軍隊是一群什麼人呢?.
就是一群強而有力.
每個頸項都很粗的大隻佬.
而且這群巨人是整個群族這樣來攻打的.
實力實在太強.
還有一些考古發現.
就是利梵音人.
有些考古的骸骨發現.
利梵音人阿納族人有多高呢?.
大概是四米多的高度.
四米多.
我才一米八.
兩個我再加一些.
加珍珍打橫躺在那裡.
就等於利梵音人那麼巨大.

$^{521}$你想想.
他躺的床也四米長.
在一個那麼巨大的群族裡面.
是一個很大的宿敵.
所以難怪在民數記十三章.
二十二至三十三節這一段.
曾經說過摩西曾經叫十二個探子去迦南地來窺探.
看一下迦南地的虛實的時候.
這十二個探子一看就看到.
果然是流瀨縟物之地.
不過他發現那裡的人是什麼人呢?.
就是這群利梵音人.
也是這群阿納族人.
一群強而有力的一群很大的敵人.
在當時以色列十二個探子裡面有十個.
都帶著一個壞消息去到他們當中.
我們的身形見絕.
我們在這群巨人面前.
好像一個蚱蜢一樣.
死定的.
怎麼打?.
面對著這群人.
就是沒得打的.
十個探子都是報一個休訊.
一個沒得打的訊息.
給摩西知道.
信弟知道這群人到最後.
我們知道整件歷史.
就是他們因為沒信心和發怨言.
他們就要留在曠野四十年.
阿納族人都沒用打.
他們只是看外表.
都令到以色列人在當時第一代已經輸了.
不過在經文裡面.
他要再一次來記載.
這個好像宿敵一樣.
這個曾經治以色列人失敗的族人.
去到約書亞.
去到今時今日.
他打贏他們.

$^{561}$怎樣贏呢?.
所有山地.
所有地方.
裡面那些巨人.
全部都剪除了.
一個都殺得清光.
不容讓在這些地方裡面.
繼續出現.
就是剛才說的那幾個地方.
希伯倫,底碧,阿拉伯,猶太山地.
以色列山地.
所有的阿納族人.
一個都殺得清光.
可能我們看到這些字.
有時候我們會有個題外話.
我們有時候會覺得.
上帝好殘忍.
要殺光這些族人.
這麼殘忍.
不過當你剛才聽到我這樣說的時候.
這些阿納族人從來不是說.
上帝好像感覺上很固執的一班這樣的人.
如果我們用愛情小說的角度.
來看這段經文.
我們有很多東西是不解的.
但如果我們用看荷里活電影.
我們怎樣去打壞人.
打終極的惡人.
那個大佬.
我們要快快將他剷除.
當我們看到那個大佬.
那個惡人被剷除的時候.
我們的心會被提振.
我們看復仇者聯盟那個.
打贏了他.
我們的心就有一種.
一定要贏.
我們今天.
不要用愛情小說的方式來看約書亞這個戰記.
那個是敵人.

$^{601}$這個敵人曾經將以色列整個軍隊羞辱.
甚至乎還沒有打.
已經令整個以色列流落曠野四十年.
他們的習俗.
他們利欺臣的生活.
令到耶和華軍隊寸步難行.
唯有靠著上帝的恩典.
唯有帶著無比的決心.
才可以打贏這一場的仗.
在聖經的作者裡面.
他要將這個宿敵.
或者某程度上是一個大佬.
打贏了他.
來讓整個以色列民族.
記下一個很輝煌的一段歷史.
打贏這場仗不是靠他們.
而是靠耶和華上帝.
不過這一隊的軍隊.
他們自然帶著上帝給他們的那種決心.
來打贏這一場的仗.
特別在十一章第23節裡面說到.
當這個亞納族人被剪除的時候.
除滅巨人.
除滅了這些人之後.
國中太平.
沒有征戰了.
之前有一次又一次.
說的是五至七年的征戰.
終於在打贏了這個亞納族人之後.
一切都歸回太平.
迦南帝就從此就止息了一些很重要的.
或者不停的那些的征戰.
打贏了這個宿敵之後.
讓整個民族就重新好像可以安居樂業.
接下來的篇幅就是說一些分地的那些安排.
不過打贏了這個亞納族之後.
是不是所有問題就解決呢?.
聖經作者其實很仔細.
大家去看的時候一定要留意一點的地方.
在第22節那裡.

$^{641}$雖然就是說在以色列中沒有留下一個亞納族人.
但是只在加薩加特和阿薩特樓下.
有一些這樣的亞納族的人.
雖然約書亞軍隊將這些巨人.
大部分99\%清除.
但是剩下的那個百分比.
原來都是散居在加薩加特和阿薩特這些不同的地方.
聖經的作者說到其實他們沒有徹底的將這個敵人消滅.
不過已經完成了很大的任務.
所以對於約書亞對於整個以色列民族來說.
你們先不要開心得這麼快.
這個隱藏的巨大的敵人依然存在在整個以色列.
或者在迦南地當中.
只不過他們可以某程度控制了整體的局面.
其實在聖經裡面剛才我也有說過.
其實對於那個仇敵.
在舊約開始來描述.
去到一路的先知書也有說.
那個仇敵其實是一個很有型有體.
讓他看得到爭戰的一個對象.
早在其實在創世記第六章四至五節裡面.
就說到有墮落的天使與人間的女子交合.
就生了那些上古英偉之人.
不是說那些人很英偉很帥英偉不凡的意思.
而是說生了那些巨人身形很龐大的那些巨人.
而這些巨人在地上裡面就一味的行惡.
任意而行.
這個就好像一個敵人的初影.
在創世記裡面這樣記載.
跟著在聖經裡面一路一路的延伸.
特別是去到出埃及記或約書亞記.
在迦南地裡面.
這些巨人就集結成為一個很強大的軍隊.
後來被約書亞殲滅之後.
其實在第十三章第三節後一點點的篇幅.
這些亞納族的巨人又再出現.
甚至乎去到薩姆爾記.
大衛用那個基源石擊打的那個哥利亞巨人.
其實都是亞納族的人.
那個哥利亞來自甚麼地方呢.

$^{681}$就是剛才散居了在迦特的那個地方.
哥利亞來自這些亞納族散居了在迦特的一個巨人.
雖然時間隔了很久.
不過這些苗頭這些巨人的敵人的苗頭.
繼續的來到去生存在這個世界裡面.
對於整個以色列軍隊來講.
他們不可以掉以輕心.
不要以為自己打贏了這一次.
或者這幾年的戰役之後.
就等於徹底將所有問題解決.
未來的挑戰仍然很大.
我們中國有句詩.
「志未就志未酬 問君知志幾時酬 志亦無盡量 人亦無盡時」.
如果我們只是在講我們的使命達到了.
做了某些事就完成了.
那個其實不是一個奮發向上的心.
我們人生這個志向完成了.
這個戰役完成了之後.
我們又會有一場新的戰役.
新的戰役再完成了之後.
還有更新的戰役一路一路等著我們.
人生就是不斷地向前.
我們要帶著這一份的戰意.
什麼時候完結?牧師.
等到耶穌基督再回來的時候.
基督耶穌最後將新天新地.
來作一個結算的時候才完成.
在基督未回來的這一段末世的日子.
我們每一個都是面對著一場戰役.
由最初他們面對著一個巨人.
一路一路延展到我們今天新約.
這個敵人不再用一個巨人的方式來呈現.
而是用一個空中掌權者.
用一個無形無體你看不到他.
但是存在迷惑人心.
扭曲人性的魔鬼撒旦.
來到我們心靈裡面成為我們的仇敵.
聽聽保羅在《爾忽所書》裡面他怎樣說.
我還有沒了的話.
你們要靠著主.

$^{721}$依賴他的大能大力.
作剛強的人.
要穿戴神所賜的全副軍裝.
就能抵擋魔鬼的詭計.
因為我們並不是與屬血氣的爭戰.
乃是與那些執政的掌權的管轄者.
幽暗世界的.
以及在天空屬靈氣的惡魔爭戰.
保羅在這裡面.
他講到今天基督徒都要繼續的來爭戰.
不過爭戰的對象不是人.
人不是我們的敵人.
不是我們的仇敵.
我們爭戰的對象是那個我們看不見.
好像籠罩著偏地橫行的空中掌權者.
那個在人心靈裡面敗壞人的生命.
扭曲人性的魔鬼撒旦.
那個是我們面對的仇敵.
仇敵的形式不同了.
但是仇敵依然在我們生命當中.
我們要帶著的就是那種來自神的剛強壯膽.
保羅都是這樣來勉勵基督徒.
我們要剛強穿起整副的軍裝.
那個軍裝其實就是在講耶穌基督的福音.
用上帝的話語.
用上帝的愛.
用上帝的真理.
來到一個充滿敵黨的世界.
一個將人性扭曲.
仇敵的掌控的世界裡面.
我們將人挽救過來.
今日我們都要帶著爭戰的鬥志來傳福音.
不要因為一兩次的失敗.
被人覺得你很無聊.
或者說一些話言不及義.
或者面對著一些人可能敵黨你.
我們就索性將自己放棄了.
繼續傳福音的這個照明.
其實我們都是不斷由失敗裡面累積經驗.
在失敗當中我們不灰心.

$^{761}$不單止一個人不灰心.
而是整個群體我們都繼續不灰心.
我們成為一個有力的團隊.
在面對著這個剛硬的世界裡面.
我們一齊打這場美好的仗.
弟兄姊妹,今日這段經文我們大致看到這裡.
如果用一個字,一句說話來總結今天的經文.
我會說我們要剛強.
面對著一個很深硬的世代.
我們要剛強.
面對著這個昔日戰我們於失敗的巨人.
我們要剛強.
我們的生命不要被過去的失敗成為我們人生的結論.
我們的人生不是被這些挫敗經歷打垮我們.
反而藉著上帝的恩典.
我們一齊努力.
在這個新的世代當中.
讓我們一齊起來.
我們總動員地邁向這個世代.
在這個世代裡面.
我們用盡不同的方法將福音傳.
求主繼續來對待我們.
激勵我們帶著鬥心打面前美好的仗.
我們同心祈禱.
天上下不負我們,多謝你.
求你在我們當中激勵我們的鬥心.
讓我們知道我們今天仍然面對末世裡面有仇敵.
困擾我們,阻攔我們.
讓我們不能將真理傳開.
幫助我們.
讓我們今天看到約書亞他們整隊軍隊怎樣去打仗的時候.
求你讓這些經歷也激勵我們的心.
讓我們能夠一齊為你打拼.
我們知道我們個人的能力有限.
但在你裡面沒有難成的事.
我們就仰望你.
我們也將我們的心靈一一地交託.
讓我們盡力盡心.
來將你給我們的福音使命實踐.
幫助我們在我們生活當中可能都遇到不同的艱難.

$^{801}$讓我們同樣帶著這一份的確信.
戰勝眼前這些一切的挫折.
讓我們挺起胸膛做人.
一齊為你打美好的仗.
願你施恩,願你大靈.
多謝主,垂聽祈禱.
奉基督耶穌得胜的聖明祈求.
Amen.
\newpage



\section{馬太福音 20:1-16}
\label{sec:XMPoIkBynFQ}
\textbf{2020年5月31日 崇拜直播｜馬太福音20\_1-16}
\newline
\newline
連結: \href{https://youtube.com/watch?v=XMPoIkBynFQ}{\texttt{https://youtube.com/watch?v=XMPoIkBynFQ}} ~~~~ 語音日期: 2020-06-01
\newline
\newline
\hyperref[sec:y8QuGSs8cQg]{\small{< < < PREV SERMON < < <}}
~
\hyperref[sec:index_chronic]{\small{[返順時目]}}
~
\hyperref[sec:PBWvRjXaPlA]{\small{> > > NEXT SERMON > > >}}
\newline
\newline
馬太福音 20:1-16
\newline
\begin{longtable}{cl}
\hline
\hline
章節 & 經文 (和合本修訂版)\\
\hline
20:1 & \begin{tabularx}{0.7\textwidth}{X} 「因為天國好比一家的主人清早去雇人進他的葡萄園做工。 \end{tabularx} \\ \\ \relax
20:2 & \begin{tabularx}{0.7\textwidth}{X} 他和工人講定一天一個銀幣，就打發他們進葡萄園去。 \end{tabularx} \\ \\ \relax
20:3 & \begin{tabularx}{0.7\textwidth}{X} 約在上午九點鐘出去，看見市場上還有閒站的人， \end{tabularx} \\ \\ \relax
20:4 & \begin{tabularx}{0.7\textwidth}{X} 就對那些人說：『你們也進葡萄園去，我會給你們合理的工錢。』 \end{tabularx} \\ \\ \relax
20:5 & \begin{tabularx}{0.7\textwidth}{X} 他們也進去了。約在正午和下午三點鐘又出去，他也是這麼做。 \end{tabularx} \\ \\ \relax
20:6 & \begin{tabularx}{0.7\textwidth}{X} 約在下午五點鐘出去，他看見還有人站在那裡，就問他們：『你們為甚麼整天在這裡閒站呢？』 \end{tabularx} \\ \\ \relax
20:7 & \begin{tabularx}{0.7\textwidth}{X} 他們說：『因為沒有人雇我們。』他說：『你們也進葡萄園去。』 \end{tabularx} \\ \\ \relax
20:8 & \begin{tabularx}{0.7\textwidth}{X} 到了晚上，園主對工頭說：『叫工人都來，給他們工錢，從後來的起，到先來的為止。』 \end{tabularx} \\ \\ \relax
20:9 & \begin{tabularx}{0.7\textwidth}{X} 約在下午五點鐘雇的人來了，各人領了一個銀幣。 \end{tabularx} \\ \\ \relax
20:10 & \begin{tabularx}{0.7\textwidth}{X} 那些最先雇的來了，以為可以多領，誰知也是各領一個銀幣。 \end{tabularx} \\ \\ \relax
20:11 & \begin{tabularx}{0.7\textwidth}{X} 他們領了工錢，就埋怨那家的主人說： \end{tabularx} \\ \\ \relax
20:12 & \begin{tabularx}{0.7\textwidth}{X} 『我們整天勞苦受熱，那些後來的只做了一小時，你竟待他們和我們一樣嗎？』 \end{tabularx} \\ \\ \relax
20:13 & \begin{tabularx}{0.7\textwidth}{X} 主人回答其中的一人說：『朋友，我沒虧待你，你與我講定的不是一個銀幣嗎？ \end{tabularx} \\ \\ \relax
20:14 & \begin{tabularx}{0.7\textwidth}{X} 拿你的錢走吧！我樂意給那後來的和給你的一樣， \end{tabularx} \\ \\ \relax
20:15 & \begin{tabularx}{0.7\textwidth}{X} 難道我的東西不可隨我的意思用嗎？因為我作好人，你就眼紅了嗎？』 \end{tabularx} \\ \\ \relax
20:16 & \begin{tabularx}{0.7\textwidth}{X} 這樣，那在後的，將要在前；在前的，將要在後了。」 \end{tabularx} \\ \\
[1ex]
\hline
\hline
\end{longtable}
$^{1}$經訓記載馬太福音20章1至16節.
馬太福音20章1至16節.
是新約的29頁.
因為天國好像家主清早去僱人.
進他的葡萄園做工.
和工人講定一天一錢銀子.
就打發他們進葡萄園去.
約在至初出去.
看見市上還有閒暫的人.
就對他們說.
你們也進葡萄園去.
所當給的我必給你們.
他們也進去了.
約在午正和新初又出去.
也是這樣行.
約在有初出去.
看見還有人站在那裡.
就問他們說.
你們為什麼整天在這裡行站呢.
他們說因為沒有人顧我們.
他說你們也進葡萄園去.
到了晚上.
園主對管事的說.
叫工人都來給他們工錢.
從後來的起到先來的為止.
約在有初顧的人來了.
各人得了一錢銀子.
及至那先顧的來了.
他們以為必要多得.
誰知也是各得一錢.
他們得了就埋怨家主說.
我們整天勞苦受熱.
那後來的只做了一小時.
你竟叫他們和我們一樣嗎.
家主回答其中的一人說.
朋友我不辜負你.
你與我說定的不是一錢銀子嗎.
那你的酒吧.
我給那後來的和給你的一樣.
這是我願意的.

$^{41}$我的東西難道不可隨我的意思用嗎.
因為我作好人.
你就紅了眼嗎.
這樣那在後的將要在前.
在前的將要在後.
請姐妹平安.
有多少人聽過剛才這個比喻.
有沒有?原來很多.
你覺得葡萄園的園主.
他做事公不公道.
有人說公道.
有些有保留.
有些沒有表態.
其實如果按照今天的文化來說.
我們說這些叫什麼.
同酬不同恭.
大家都是出這麼多工資.
但有些做多些時間.
有些做很短時間.
我們會覺得這些不是很合理的待遇.
我曾經看過一篇文章.
都是和這個比喻有關.
他最後問一個問題.
他說如果你是當時的工人.
你會怎樣去寫這個葡萄園的故事.
於是有些人說.
我當然是非常生氣.
所以我會去葡萄園拉橫額.
去抗議園主的不公道.
有另一些人說.
怎麼可以這麼虧蝕.
我們應該用今天的情況.
去哪裡申訴?評機會.
我們一定要拿回公道.
拿回合理的賠償.
當然有些人比較有創意.
他會提議.
我覺得這個園主要去讀一個學位.
是一個工商管理MBA的學位.
為什麼?他好像不太懂管人.

$^{81}$好像人事策劃管理方面.
差一點要去進修.
也有人說.
明天這個園主再出來找工人工作.
我會怎樣?.
我一定閃.
一定避開他.
避到下午四五點左右.
我就出現給他找到.
這樣我就可以怎樣?.
做最短時間的工作.
但是怎樣?.
拿足整天的工資.
大家應該會計算.
好像沒有反應.
不過他繼續這樣說.
但是我發現這個園主.
原來只是上午出來找工人.
下午他就沒有出現過.
結果我整天都吃白果.
沒有飯吃.
並且家人都要和我一起捱肚餓.
丁姐妹.
我們今天聽了這個比喻之後.
我們又會怎樣去.
繼續寫這個葡萄園的故事呢?.
耶穌說這個比喻和天國有關.
天國就好像葡萄園的主人.
出去僱用工人進葡萄園工作.
首先我們要了解什麼是故宮.
有點像深水Po .
有條街有很多布匹.
有沒有印象?.
以前.
現在我不知道.
以前一清早.
在那裡聚集了很多基層人士.
他們來到要碰碰運氣.
看看那一天有沒有人請他們去工作.
搬運布匹.

$^{121}$正所謂手停就會怎樣?.
就是口停.
故宮要有工開.
才能賺到生活費.
就像這段經文所說.
清早的時候.
其實大概早上六點.
園主就開始請第一批工人.
跟他們說好.
一天就一錢銀子.
一天是什麼意思呢?.
一個工作天.
在當時的社會.
就是朝六晚六.
要工作十二個小時.
工作十二個小時.
就會獲得一錢銀子.
一錢銀子大概足夠.
當時一家四口生活一天.
我們先有了這個概念.
園主在六點鐘就請了一班人.
但到了三個小時後.
其實這裡大概是上午九點鐘.
這個園主又再出去找人.
他看到街上仍然有人行站住.
於是又請他們去葡萄園工作.
留意這次園主有沒有說實.
出多少錢工資?.
沒有,他只是承諾.
會給工人應得的工資.
什麼意思呢?.
假如你今天去建局找工作.
老闆說請你.
你自然就會問什麼?.
一個很切身的問題.
每個月出多少工資?.
老闆就說你放心.
會給你應得的工資.
你先說明.
你願意嗎?.

$^{161}$沒有說實多少錢工資.
只說我會給你應得的工資.
你會猶豫,因為不知道多少錢.
但當時的僱工其實沒有那麼多選擇.
因為已經等了三個小時.
再等下去會不會有人請其實不知道.
所以既然有人請.
無論如何都先做.
有工作總是好的.
於是又到了中午十二時.
午正和新初.
即下午三時.
這個園主很特別.
他又繼續出去找工人.
也請了一批工人.
和之前九時的那批一樣.
到最後去到有初.
其實就是下午五時.
園主.
雖然你剛才記得.
是早六晚六.
其實距離下班還有一個小時.
這個園主仍然繼續出去找工人.
他看到街上仍然有人問.
為什麼你們經常在這裡行站?.
他們怎樣回答?.
他們說我們不是不想開工.
不過我們等了一整天.
都沒有人請我們.
於是園主就說.
你們也進來葡萄園工作吧.
留意.
這位園主這次.
連人工這兩個字都沒有提到.
很特別的是.
這班最遲的工人也沒有問.
為什麼?.
心裡有數.
只可以工作多少小時?.
一個小時.

$^{201}$你預計一定賺的錢是少的.
但總比那天好.
零收入.
為什麼?.
所謂小的不吃.
大的也要吃.
賺到多少也好.
買到一個包給兒子和女兒吃.
總比他們要和大人一起捱餓好.
這是當時他們的狀況.
一個小時後.
到了黃昏傍晚的六點.
園主對管事說.
叫工人出來給他們發工資.
最遲來的一班.
下午五點來的一班.
我們先發工資給他們.
於是胡經文說.
五點來工作的那班.
他們每人可以得到一錢銀紙.
雖然工作了一個小時.
但可以拿到全日工資.
正不正?.
正啊.
早上六點的那班.
看到有什麼想法?.
你們有什麼想法?.
那批工作了一個小時的.
竟然可以拿到全日工資.
這次我們就發大財.
小學生的數學題.
明仔在葡萄園工作一小時.
賺取了十元.
強仔在葡萄園工作十二小時.
賺了多少錢?.
也是十元?.
不會的.
小學是教你怎樣應用.
不會是十元.
一定是十二倍.

$^{241}$一百二十元.
很簡單的數學題.
不複雜.
我們工作了十二小時.
肯定不只出一日工資.
這次沒有四倍也有三倍吧.
怎樣也要相公.
這個期望也很合理.
不過最後他們發現.
原來我們也只有一樣.
只有一錢銀紙.
我不知道比作弟兄姊妹.
你有什麼反應?.
你的期望落空了.
你不但生氣.
還覺得很不值.
可能我們第一個反應.
就是有沒有搞錯?.
同酬不同功.
好像我們以前所說.
做又三十六.
不做又怎樣?.
都好像是三十六.
我們最不喜歡這種制度.
很反感.
於是你見到經文裡說.
六點最早開工那批.
就覺得不公平.
很不服氣.
他們埋怨那位主人.
跟他們說.
五點那批只做了一個小時.
我們經文說得好.
受熱勞苦.
曬著來摘葡萄為你工作.
你竟然給大家一樣待遇?.
這樣很不公道吧?.
你擺明欺負我們.
但這個原主.
於是他就回答了.

$^{281}$回答其中一個.
他說朋友.
我有沒有少給你薪水?.
沒有.
我沒有少給你薪水.
我有沒有叫你加班超時工作?.
沒有.
你真的只做了十二小時.
我們不是說好了.
做一天有一池銀紙.
所以我沒有虧待你們.
我只不過是額外.
多給你薪水.
給最遲那批.
難道我的錢.
不可以照我自己的意思來花?.
還是因為我仁慈.
在你們面前做好人.
你們就眼紅.
不值得?.
這個比喻.
耶穌就說到這裡.
弟兄姊妹.
先問問大家.
舉手讓我知道.
如果有得你選擇的話.
你想做哪一批工人?.
想做早上六時的那批.
這班一定是習慣早起.
可能要四五點就起床.
非常好.
有沒有想九點起床的?.
睡眠三小時.
都有想睡眠一點.
太早了六點.
有沒有想中午十二點.
三點的那批?.
沒有?.
這批其實可以很好.
做完家務.

$^{321}$送小朋友放學.
吃飯.
吃下午茶才開工.
三小時之後.
六小時或三小時之後就發工了.
沒有?.
有沒有想做下午五點的那批?.
有的,不過不敢舉手.
其實不好意思.
其他不舉手的呢?.
還有其他想法嗎?.
最好當然是怎樣?.
聖經沒有教不勞而獲.
最好當然是怎樣?.
做半小時.
十五分鐘.
然後就出全日的工資.
這樣就正了.
再問得深入一點.
聽著.
你想上帝這位葡萄園的主人.
他怎樣對待你?.
想不想上帝對待六點的工人?.
即是說.
跟足合約.
依足規矩.
做足十二小時.
就出一日工資.
換句話來說.
只持續一分鐘也好.
都當你沒有履行合約.
一個崩都不給你.
想不想做這一批?.
剛才舉手的立即想一想.
立即沒人敢舉手.
又或者想不想上帝.
早上九點至下午三點那批對待我們?.
即是怎樣?.
會給你應得的工資.
有沒有想過應得的人工.

$^{361}$這幾個字可圈可點?.
有沒有想過?.
如果講得溫柔一點.
我會給你應得的人工.
如果講得嚴厲一點.
我會給你應得的人工.
即是甚麼意思?.
即是你出多少錢的工資.
跟甚麼掛勾?.
跟你的工作小時.
跟你的工作表現掛勾.
主人覺得你值.
那就給你.
覺得你不值.
就不給那麼多.
甚至不給.
這樣好不好?.
想起來又不是很好.
所以弟兄姊妹.
其實不是我們想做哪一批工人.
實情是我們很需要上帝.
好像對待五點那一批.
即是不是跟合約.
你做齊上帝所有要求.
我就給你永生.
亦不是跟工作時薪,工作表現.
我看你做得怎樣.
做了多久,做了多少.
你只是做了五成的工作表現.
我就怎樣?.
我給你五成的救恩.
那糟了.
上面那一部分還是下面那一部分.
上到天價呢?.
自己選擇.
你做到九成的.
我給你九成的救恩.
又是你自己選擇.
想頭還是腳留下來.
其他的上去.

$^{401}$弟兄姊妹,可不可以這樣?.
不可以的.
我們每一個都需要上帝.
好像對待下午五點那批工人.
用恩慈,恩典來對待我們.
否則死定了.
所以我們看到這位玄主.
他沒有按照合約規定.
或是一個公平的原則.
來對待最後那批工人.
而是按照這批工人.
他們真實,真正的需要.
來作出恩待.
否則又如何?.
如果你們做到一個小時.
我給你們十二分之一的工資.
他們很可能全家就要繼續.
活在飢餓和惶恐的裡面.
所以耶穌用這個比喻.
向我們說明了一個道理.
就是上帝的恩典,恩慈.
是在規定,約定之外再加上去的.
上帝的恩典,恩慈.
是在所謂公平之上顯明出來的.
因為我們能夠得救進天國.
不是基於我們本身擁有的條件.
有多好的能力,多好的際遇.
又或者你的表現有多卓越.
為上帝立下了多少功勞.
而是完全出於上帝的恩慈和憐憫.
這個比喻其實是這樣的意思.
換句話來說,耶穌要強調一件事.
或者逆轉當時每一個人的觀念.
就是天國的道理,天國的觀念.
和社會給我們的價值觀是不同的.
社會的法則通常是怎樣?.
如果你的條件越優厚.
你的能力越強.
表現越卓越.
又或者你的際遇越好.

$^{441}$通常你在社會就會怎樣?.
你是名列前茅那一班.
你永遠是最頂尖的那一班.
但是如果你的條件比較差.
你的能力又不夠.
或者你的際遇比別人沒那麼好.
又或者你的表現差一點.
那你就永遠怎樣?.
你可能是跑輸戰底那一班.
所以如果按照社會的法則.
往往就會怎樣?.
在前的就繼續在前.
在後的呢?.
很難上到去.
繼續去在後.
但是耶穌說天國的數目不是這樣計算的.
他說上帝是願意所有的人.
無論你本身的條件是怎樣.
你的能力如何.
你的際遇到底是好壞.
又或者無論你是在前那一班.
最早被請回來的工人.
又或者你是最後最遲受聘的.
即是說無論你在葡萄園裡面.
你的工作量.
你的功勞是多還是少.
因為上帝的恩典,憐憫.
每一個人都可以有機會得救.
都可以有機會進入天國.
分享到上帝的眷顧和豐富.
耶穌是這樣說的.
亦都要扭轉當時每一個人.
或者說今天每一個人的觀念.
是不難理解的.
這個天國的道理.
還記不記得耶穌釘十字架的時候.
他兩旁有哪些人?.
還有兩個囚犯.
其中一個是怎樣?.
就求耶穌在你得國的時候.

$^{481}$紀念我.
於是耶穌二話不說.
就說:好啊.
你今天就跟我一起進到樂園裡面.
嘩!這麼好?.
就是這麼好.
在當時的社會確實.
這個囚犯是在後的.
沒有人看好他.
甚至會看死他.
他亦都不可能像彼得其他門徒.
可以為耶穌付出過什麼.
沒有這個機會.
但因為上帝的恩典和憐憫.
他可以進入天國.
他原本是在後.
但耶穌說你可以在前.
所以弟兄姊妹.
關鍵是人願不願意.
好像聖經所說.
認罪悔改.
接受這位天國主人的邀請.
進到他的葡萄園.
從今以後.
不再為自己而是為上帝而活.
為上帝而打拼.
不過我們要留意一件事.
在當時來說.
耶穌這個比喻亦是暗示.
這個邀請有沒有時限?.
有沒有時限?.
有的.
因為就好像收割葡萄園.
其實都有時限的.
當收割完畢.
主人就不會再繼續出來.
找人進去工作.
所以無論在座當中.
又或者在屏幕前.
如果仍然有未認識耶穌的朋友.

$^{521}$是真的.
這段經文是邀請我們.
要把握一個機會.
接受耶穌基督的邀請.
做一些進入上帝的原子.
做一些去經歷到.
真正的平安和盼望.
就好像我以前認識一位弟兄.
他一出來工作.
賺到第一份薪金.
他就做了一樣的舉動.
就是立即跟他的父母.
買了一份醫療保險.
留意他不是做保險.
所以第一個就找父母開單.
他不是做保險.
而是他覺得父母年紀大.
真的有需要.
給予一份保障.
同樣的弟兄姊妹.
相信了主的我們.
會不會關心.
你身邊的人有沒有那種.
永恆屬天的保障.
有沒有留意這件事.
當然你和我不可以.
代替他信耶穌.
沒有什麼可以代替.
但卻可以介紹耶穌.
給他們認識.
所以趁現在又賣廣告.
記得.
好像大家都認識.
記得訂閱.
記得讚好和分享.
旺朕的YouTube Channel.
記住了.
介紹給你身邊的人.
其實不知道你有沒有留意.
其實不只是教會弟兄姊妹.

$^{561}$有很多不是教會的.
可能是弟兄姊妹.
甚至微信的朋友.
都看這個Channel.
他們都有機會認識耶穌.
認識這個信仰.
所以我邀請你們.
繼續將這個網站推出.
讓更多人可以認識到.
耶穌的好處.
好嗎?.
另一方面.
不知道你有沒有留意到.
其實這個圓主.
他可以怎樣?.
他叫六點.
最早的工人先來發工資.
然後發工資後.
快點回家.
然後到最後.
才恩代五點那一批.
那就什麼事都沒有.
怎會有任何矛盾衝突呢?.
所以你留意耶穌.
是刻意設計這個比喻.
專登激你一下.
來幫我們檢視一下.
心靈到底是否健康.
什麼意思呢?.
為什麼有時我們做基督徒.
甚至熱心事奉.
但總是活得不太開心.
不太喜樂.
或許會否像這個比喻所說.
是因為人生中.
夾雜著一些比較和計較.
就像最早六點開工的工人.
總覺得上帝對我不公平.
總覺得上帝對其他人.
比對我好.

$^{601}$你想想.
其實如果六點的工人.
他又看不到.
還不知道原來五點那一班.
是出同一工資的話.
他會否這麼憤怒呢?.
不會了.
所謂眼不見為乾淨.
看不到就不會眼紅.
又或者相反.
他也看到五點那班工人先發工資.
不過他原來看到.
原主只給了十二分之一工資.
這樣不公平嗎?.
完全公平.
公平嘛.
也不會因為這樣.
埋怨疾到主人處事手法.
對不對?.
不知道大家有沒有這種經歷.
回想小學讀中學.
考完試.
派卷取得成績單的時候.
我有一個經歷.
那次派了一份卷回來.
只有七十分.
開心嗎?.
一般般吧.
七十分即是沒有A.
但又八卦.
又和周圍的人問.
你取了多少分?.
你取了六十分.
你呢?.
你只取了五十分嗎?.
還問.
三十五分?.
最厲害的也是七十五分.
於是瞬間.
開心回來了.

$^{641}$為什麼?.
也是七十分.
沒有變.
沒有加分.
但因為什麼?.
雖然是七十分.
但拉曲線可以取得什麼?.
取得A.
馬上笑了.
當然不敢那麼囂張.
但也試過相反.
有一個經驗.
那次我取了八十八分.
厲害吧?.
取A吧.
於是又八卦.
想炫耀一下.
想著你七十分.
又是八十八分.
但誰知每一個都說.
我九十分.
我九十五分.
你一百分?.
最差的也九十分以上.
於是便馬上.
失落.
很不開心.
非常不開心.
所以弟兄姊妹.
有時做人不開心.
不喜樂.
與你收穫無關.
與你得失多少.
也未必有關係.
與什麼有關?.
與比較.
計較有關係.
所以耶穌提醒我們.
其實原來進天國.
得到上帝的賞賜.

$^{681}$是出於什麼?.
恩典和憐憫.
剛才詩歌也有唱.
作為天國的子民.
其實一定要學習一個功課.
無論學多久.
都要學習一個功課.
就是一定要在人生裡.
少計算自己個人的得失.
多以上帝的榮耀為優先.
你想想.
耶穌有沒有計算過我們?.
有沒有想過.
如果耶穌計算自己得失.
那你和我都大件事.
什麼意思呢?.
假如耶穌對父神說.
不公道啊.
有什麼理由.
我這個獨生子.
要這麼吃虧.
為罪人討價還價?.
不是說明了.
罪人的公價是什麼?.
死.
這本來是世人.
應該得到的人工.
結果是這樣的.
應該和父神永遠隔絕.
所以你想深一層.
最被不公平對待的是誰?.
是耶穌.
因為無罪的主.
要替有罪的我們贖罪.
代替我們還那一筆.
我們這世永世都還不清的債.
讓我們有機會成為上帝的兒女.
深一層為什麼這麼不公平.
這麼不公平.
耶穌也願意?.

$^{721}$因為祂的愛.
因為祂不只是單單顧著自己的事情.
祂也有一個更超越的大局觀.
祂看到上帝傳盤的計劃.
祂看到上帝的需要.
所以總是以上帝的益處.
上帝是否得到榮耀.
作為耶穌基督行事為人的優先考慮.
我記得曾經聽過一位基督徒朋友分享.
他說在辦公室裡.
有其他同事和他的職位相等.
不過工作量卻比他少.
覺得有些不公平,不公道.
於是他和他的部門主管反映.
因為他的經驗能力比其他人優勝.
所以他相對比其他人多一點工作量.
他聽到之後仍然覺得很不值.
不過又不知如何用合宜的方法.
來表達心裡的不舒服.
在靈修的過程裡.
他也不斷去問.
耶穌我可以怎樣做?.
我可以怎樣回應?.
經過一段時間的沉澱.
偶然聖靈提醒他一件事.
你試試在這個情況裡學習感恩.
有什麼好感恩的?.
有好的工作經驗和能力.
這是真實的.
所以才讓他經歷這種試煉.
上帝不會虧待人.
讓你經歷.
即是他覺得你可以承受.
他覺得你可以在裡面有所更新成長.
第二就是要學會.
要預先去感恩.
為什麼要預先去感恩?.
因為上帝既然讓你遇到這個試煉.
他也一定會給你足夠的供應.
足夠的恩典和能力.

$^{761}$你去跨過這個試煉.
上帝已經允許了.
而當他不再用一個計較抱怨的目光.
去看他目前的遭遇的時候.
很特別他開始體會到上帝的心意.
就是要用他在辦公室裡.
幫助祝福更多的人.
於是他除了每天做好自己的工作.
他也花額外的時間走多步.
將他的經驗和同事分享.
令到整個團隊的能力都提升了.
他的工作量也不知不覺輕生了.
整件事他有沒有加薪呢?.
沒有.
有沒有付出多了?.
有.
但是賺了些什麼呢?.
賺了他和同事的關係.
他祝福了更多的人.
更重要的是.
上帝在他的辦公室得到榮耀.
親愛的弟兄姊妹.
有時人生裡太多比較,計較.
不想吃虧.
其實很容易錯過了.
上帝在你身上的一些恩典和計劃.
所以多一些感恩.
才能夠增潤到我們和上帝和身邊的人的關係.
你看看那些工人的比喻.
你想一想他們有什麼可以感恩呢?.
這幾批工人有什麼可以感恩呢?.
先說六點那批.
他們沒有錯.
他們工作的時間是最長.
但待遇是一樣的.
但是他們有什麼可以感恩呢?.
他們是最早被聘請的那批.
換句話說.
他們不需要站在街上擔心.
他們是最早得到生活保障和盼望的人.

$^{801}$感不感恩呢?.
很感恩的事.
並且他們可以替最晚那一批去感恩.
最晚那一批很淒涼的.
最晚那一批不只是為了自己.
很多時候是背負一家的責任.
來到街上.
一直站著.
一個小時都沒有人聘請.
到了十二點,三點都沒有人聘請.
你想想他們怕不怕.
他們如何向家人交代呢?.
很沒有盼望和安全感.
六點那一批知道的時候.
有沒有能力幫助他們呢?.
其實沒有.
但是感恩有一位沒有恩典憐憫的主.
願意出手眷顧他們.
可以為別人感恩.
再說中間那些.
九點至三點那些有什麼可以感恩呢?.
有什麼可以感恩呢?.
你細心想一想.
剛才我們說主人古古怪怪.
六點拯救就行了.
為什麼九點,十二點,三點,五點又要出去請人呢?.
你想想他是不是真的很缺人手?.
還是他不懂規劃人手.
所以經常要出去請人?.
你想想上帝這位天國主人.
他的計劃不會這麼差.
你有沒有留意.
經文裡形容.
由九點至五點那一批人.
其實是什麼人呢?.
經文用兩個字來形容他們.
閒暫.
閒暫不是說他們不想開工.
他們都很想有工作.
很想有一番作為.

$^{841}$只不過限於自己的條件.
自己的能力和際遇.
沒有身邊的人這麼好.
所以只能夠一直站著等.
等著人請他.
但一直都沒有人要他們.
但這位主人.
他是願意一次又一次出去.
其實不是僱用他們.
是幫助他們.
讓他們得到保障和溫飽.
不知道大家小時候.
或現在有沒有打乒乓球?.
有沒有打過乒乓球?.
小學有吧?特別特別.
男生會多一點.
女生可能是跳橡筋繩或其他.
小時候很多朋友都打過乒乓球.
你知道以前資源不太好.
很多時候學校只有兩三張球桌.
一大班人要玩.
通常是怎樣玩?.
有沒有試過猜隊?.
例如有十四個同學.
想打乒乓球.
猜兩隊.
最強的兩個一定是師傅.
可能一人猜包剪鎚.
分成兩隊每隊七個人.
又或者大家鬥球技.
每人一球可以選一個人.
通常怎樣?.
我贏了,我一定選什麼?.
選最強的.
選最厲害的.
不會輸的人先入隊.
選了選了.
最後那幾個一定是怎樣?.
最差的.
所以你記不記得有個.

$^{881}$打乒乓球的術語.
在小學或中學都會出現.
就是什麼?.
有沒有聽過「猜垃圾」?.
沒有聽過?.
奇怪,早堂的男生是喝牛頭的.
沒有聽過「猜垃圾」?.
知不知道什麼是打乒乓球猜垃圾?.
猜到最後剩下的沒什麼人想要.
但又要被迫要那些.
能力差一點.
沒那麼好的條件.
一定是輸球那些.
感覺很不好.
沒有人喜歡成為垃圾.
沒有人喜歡永遠.
沒有人要淪為最後那個.
他們很想玩.
但未必有這樣的能力.
未必有這樣的際遇.
事實上在比喻裡邊.
閒暫的人往往是這種情況.
但很感恩.
有一位沒有恩典憐憫的主人.
他沒有忘記他們.
並且看好他們.
他願意主動出去.
我請你們進去.
進我的園子.
你們也會像六點那群人一樣.
得到生命的保障.
得到平安.
所以弟兄姊妹其實.
這個比喻最後呼籲我們.
其實在你身邊有很多閒暫的人.
特別在今天這個世代.
有很多人仍然活在焦慮惶恐的裡面.
甚至可能我們自己.
也會對前景失去了信心和盼望.
但我們怎樣讓他們知道.

$^{921}$多亂也好.
有一個好消息.
就是天國的主人.
從來沒有忘記過他們和我們.
並且願意尋找.
拯救我們每一個人.
就好像這段經文.
他要我們思考的就是.
仍然繼續在上帝的葡萄園工作的我們.
會不會也跟主人說.
主啊.
其實外面仍然有很多人.
閒暫住.
過著一些沒有保障.
平安盼望的日子.
你可不可以.
差派我們出去代表你.
去請他們盡力的園子工作.
讓他們也可以像我們一樣.
得到上帝的眷顧.
我記得有一本書.
校園出版的.
叫做「活出群體的美好」.
裡面提到一個很有意思的例子.
假如你對著一群加拿大的印第安小朋友.
來到去說.
最先舉手回答問題的.
可以得到一份獎品.
你猜他們這群小朋友會怎樣做?.
原來他們會全部聚在一起.
一起想答案.
然後一起去舉手回答.
很有趣 為什麼會這樣?.
因為他們不能夠接受一個情況.
就是只有一個人.
贏得那份禮物.
其他人都要成為輸的人.
同樣的 兄弟姐妹.
上帝不只是想我們坐在這裡.
等人來.

$^{961}$上帝想我們像身後這塊新的展板.
你要出去的.
你要主動一點.
出去教會以外.
好像那位園主.
或者代表那位園主.
去尋找那些.
未有著落的人.
未認識福音的人.
你去代表耶穌.
邀請他們.
介紹耶穌給他們認識.
弟姐妹求上帝容我們每一個.
讓我們可以一起去回應.
上帝給我們的提醒.
和給我們的託付.
藉著我們來經歷.
讓人去經歷上帝的福氣.
我們同心討過.
將這個比喻.
告訴我們每一個.
讓我們知道.
我們今天每一個在你面前.
其實是出於你的恩典.
其實不是因為我們早了來.
不是因為我們條件比較好.
能力比較強.
際遇被人優勝.
又或者我們為上帝.
你立了很多功勞.
其實是出於上帝你.
無條件的揀選.
和你的愛.
你的恩典.
我們才可以得救進到天國.
所以上帝既然這樣.
幫助我們在人生裡.
我們不要太計較.
自己個人的得失.
不要因為這樣.

$^{1001}$看不到上帝對我們很好.
你給了很多恩典.
很多的供應我們.
更不要因為比較計較.
忘記看不到上帝.
你給我們的旨意.
你給我們的託付.
就是你要用我們.
去祝福身邊的人.
特別是未信的人.
所以求主你幫助我們.
可以去回應你.
讓我們知道.
在弦外仍然有很多人.
閒暫住.
過著顛沛流離的日子.
過著沒有平安.
沒有盼望.
沒有保障的生活.
神話我們都是一樣.
藉著這段經文幫助我們.
能夠重新安穩在裡面.
重新出發.
不只是在教會裡面.
而是主動出去.
就好像那位的原主.
我們代表耶穌.
來尋找失喪的人.
給我們有智慧.
給我們有耐性.
給我們有一個願意捨己的心.
來為上帝擺上.
讓人真的可以.
來經歷福音的好處.
經歷到耶穌基督.
你自己大能的拯救.
求聖靈你幫助我們.
讓我們可以坐言起行.
回應你的教導.
我們這樣的禱告.

$^{1041}$交託仰望.
耶穌基督你得勝的名字.
來去呼求.
阿們.
\newpage



\section{以弗所書 1:15-23}
\label{sec:PBWvRjXaPlA}
\textbf{2020年7月5日 崇拜直播(主餐)｜以弗所書1\_15-23}
\newline
\newline
連結: \href{https://youtube.com/watch?v=PBWvRjXaPlA}{\texttt{https://youtube.com/watch?v=PBWvRjXaPlA}} ~~~~ 語音日期: 2020-07-06
\newline
\newline
\hyperref[sec:XMPoIkBynFQ]{\small{< < < PREV SERMON < < <}}
~
\hyperref[sec:index_chronic]{\small{[返順時目]}}
~
\hyperref[sec:iBJ_D2WmvXo]{\small{> > > NEXT SERMON > > >}}
\newline
\newline
以弗所書 1:15-23
\newline
\begin{longtable}{cl}
\hline
\hline
章節 & 經文 (和合本修訂版)\\
\hline
1:15 & \begin{tabularx}{0.7\textwidth}{X} 因此，我既然聽見你們對主耶穌有信心，對眾聖徒有愛心， \end{tabularx} \\ \\ \relax
1:16 & \begin{tabularx}{0.7\textwidth}{X} 就不住地為你們感謝神，禱告的時候常常提到你們， \end{tabularx} \\ \\ \relax
1:17 & \begin{tabularx}{0.7\textwidth}{X} 求我們主耶穌基督的神，榮耀的父，把那賜人智慧和啟示的靈賜給你們，使你們真正認識他， \end{tabularx} \\ \\ \relax
1:18 & \begin{tabularx}{0.7\textwidth}{X} 照亮你們心中的眼睛，使你們知道他呼召你們來得的指望是甚麼，他在聖徒中所得榮耀的基業是何等豐盛， \end{tabularx} \\ \\ \relax
1:19 & \begin{tabularx}{0.7\textwidth}{X} 並知道他向我們這些信的人所顯的能力是何等浩大，這是照他的大能大力運行的。 \end{tabularx} \\ \\ \relax
1:20 & \begin{tabularx}{0.7\textwidth}{X} 這大能曾運行在基督身上，使他從死人中復活，又使他在天上坐在自己的右邊， \end{tabularx} \\ \\ \relax
1:21 & \begin{tabularx}{0.7\textwidth}{X} 遠超越一切執政的、掌權的、有權能的、統治的和一切有名號的；不但是今世的，連來世的也都超越了。 \end{tabularx} \\ \\ \relax
1:22 & \begin{tabularx}{0.7\textwidth}{X} 神使萬有服在他的腳下，又使他為了教會作萬有之首； \end{tabularx} \\ \\ \relax
1:23 & \begin{tabularx}{0.7\textwidth}{X} 教會是他的身體，是那充滿萬有者所充滿的。 \end{tabularx} \\ \\
[1ex]
\hline
\hline
\end{longtable}
$^{1}$今日的經文是記載在二弗所書一章十五至二十三節.
大家可以翻開新約二百六十五至二百六十六頁.
或者在程序表的內頁也有.
因此,我既聽見你們信從主耶穌.
禱告的時候,想提到你們.
求我們主耶穌基督的神.
榮耀的父.
將那次人智慧和啟示的靈.
上給你們,使你們真知道祂.
並且照明你們心中的眼睛.
使你們知道祂的恩兆有何等指望.
祂在聖徒中得的基業.
有何等豐盛的榮耀.
並知道祂向我們這信的人所顯的能力.
是何等浩大.
就是照祂在基督身上所運行的大能大力.
使祂從死裡復活.
叫祂在天上坐在自己的右邊.
遠超過一切執政的.
掌權的,有能的,主治的和一切有名的.
不但是今世的.
連來世的也都超過了.
又將萬有服在祂的腳下.
使祂為教會作萬有之首.
教會是祂的身體.
是那充滿萬有者所充滿的.
這是上主的話.
弟兄姊妹平安.
一齊來敬拜一齊服侍祂.
接下來七月八月我們會有一個新的講道系列.
這個系列叫做「發放群體的光芒」.
這個題目很熱血的.
其實這個題目是我們一班同工.
我們去思考.
在這段日子裡.
在香港.
或者在我們教會生活裡.
我們應該要有什麼信息呢?.
我們一路去禱告的時候.
我們知道.

$^{41}$今天這個世界很需要教會.
可能這個世界他們都不知道.
他們自己需要教會.
但是我們作為基督徒我們知道.
因為我們面對著這個世代.
不停地改變.
人的心靈深處裡.
充滿了很多的絕望.
困惱.
甚至很多的罪惡.
有時候很多人.
按照自己的能力他們招架不住.
唯有主耶穌基督.
我們藉著.
祂所呼召的聖教會.
我們成為一個亮麗的群體.
藉著我們的生命.
我們將很好的信息帶給這個世界.
所以接下來我們會用一連幾次.
查完已忽所書的六章經文.
內容裡面可能有一些會有重複的.
前前後後.
不過我希望我們讀完已忽所書之後.
我們心靈裡面有更多的確據.
在我們看聖經前.
問大家或者你想一想.
你自己的信主前後有沒有改變.
或者有些人生的經歷.
在我信主之前.
我家人有機會和他聊聊天.
家人經常都說我.
其實生得牛高馬大.
其實是一隻母痴貓.
從小到大都是.
見到有些事情發生.
我很怕事.
小時候都被人欺凌過.
但我不覺得那些是欺凌.
但覺得有些矛盾.
有些衝突不想面對.

$^{81}$或者面對一些失敗挫折的時候.
我最喜歡的就是睡覺.
曾經有一次.
自己覺得自己很失敗.
讀書很差.
又被人取笑.
大被蓋過頭.
一睡就睡了18個小時.
其實不是真的睡著了18個小時.
只不過被蓋過頭.
睡了覺.
覺得好像不用面對一個現實.
覺得這個世界很殘酷.
或者覺得自己沒有能力去面對.
用鴕鳥政策來逃避眼前的責任.
或者我眼前自己的一些失敗.
用睡覺的方式來逃避.
不過信了主之後.
我家人說看著我一個人改變.
他用一種字眼.
說我有一種傻乎乎的盼望.
什麼傻乎乎的盼望呢.
信了主之後.
開始一個人面對困難的時候.
願意面對多一些.
過往逃避責任.
逃避一些困難.
信了主之後.
慢慢懂得面對這些困難的地方.
總之有不同的挑戰.
不是逃避了它.
不是睡覺睡了它.
而是正面來面對.
我也不知道自己是怎樣改變的.
不過家人就說.
耶穌真的幫了我.
由一個很怕事的人.
不會完全變成另一個人.
但至少不像以前那樣.
完全怕事.

$^{121}$躲起來鴕鳥政策.
而是夠膽挺起胸膛.
面對生命裡面有不同的責任.
甚至有這些失敗.
弟兄姊妹.
我不知道你信了主前後.
你自己有沒有這樣的改變.
或者你自己有什麼更新.
但可以肯定的一件事就是.
信仰不單只是說一套腦裡面的知識.
信仰或者基督耶穌帶來的福音.
是扭轉生命.
改變很多生命.
很多困難的一個很重要的福音.
因為知道上帝愛我們.
以至我們有改變的動力.
今天這一段經文.
《以弗所》書裡面第一章.
是一個很重要的引言.
保羅席爵一個很壯闊的描述.
要告訴所有基督徒.
我們的生命不單止可以改變.
我們的生命的改變.
更加成為這個世界很重要的盼望.
當我們認真去讀的時候.
我們會知道我們這個教會.
我們一群基督徒是不得了的一個群體.
我希望我們今天仔細去讀的這段經文.
讓我們生命有更多的亮光和啟迪.
讓我們更加有力量.
迎接眼前不同的挑戰.
我們一起祈禱好不好?.
求上帝的靈幫助我們.
讓我們明白祂的話語和心祈禱.
天父我們恭敬在你根前.
我們恭聽你對我們要說的話.
我們求你.
你的話語成為我們生命的幫助.
在此時此刻.
我們求你.

$^{161}$你的話語進到我們的內心.
幫助孫講的講得清楚.
幫助每位聆聽的弟兄姊妹聽得明白.
求你繼續去說領我們.
恭敬將來的時間交託.
願你賜福.
牽上祈禱奉基督耶穌.
得勝的聖名祈求.
Amen.
以弗所書是一個什麼樣的書卷呢?.
很多弟兄姊妹讀以弗所書.
會覺得裡面的神學好像很博大精深.
甚至乎有很多神學上的論述.
這個是真實的.
事實上當你看以弗所書.
和保羅其他的書信.
他的寫作方式有些不同.
如果你有機會看過不同的保羅書信.
很多時候保羅的書信.
他會有一個很明確的問題.
或者教會面對不同的挑戰.
他會從問題入手.
藉著這些處理問題.
將一個上帝的話語.
將神學藉著這些處境裡面呈現出來.
不過以弗所書的寫作方式就有些不同.
你看到他讀的時候.
他沒有特別去講教會某一個很具體的問題出現.
他很籠統地只是在講一個很大衛的處境.
不過他用得最多著墨的.
他是在講耶穌基督是一個怎樣的神.
有很多悉經學者.
或者今天也有很多學者都一致認同.
其實以弗所書和歌羅西書.
就好像是一個並行的書卷一樣.
歌羅西書的技術和以弗所書是很相似的.
不過在這麼相似的技術裡面.
重點會有些不同.
歌羅西書是保羅寫的前傳.
講基督的前傳.

$^{201}$將基督耶穌的救贖和宇宙為首的中心.
將萬有伏在基督耶穌的腳下.
他將耶穌基督描繪得很淋漓盡致.
在歌羅西書裡面的基督.
是一個掌管宇宙萬物的中心.
他掌管天地萬物.
人信靠他的時候就是倚靠這個主.
所以基督耶穌是超越萬有.
將萬有的豐盛都藏在基督的身上.
這個是前傳.
後傳就是既然基督是一個這麼榮耀的身份.
他創造天地萬物萬有之首.
在他轄下的教會是一個怎樣的教會呢?.
以弗所書就是用一個這麼榮耀的基督的神學論述.
來發揮下去.
就講到教會同樣得著這份榮耀.
教會也要得著基督耶穌所享有的尊榮.
和將萬有的豐富.
不單止耶穌自己擁有.
也將這種豐富的教會彰顯在這個世界當中.
今天我們讀的經文一章15至23節.
是一個祈禱回應著一章第3至14節裡面.
一個很豐富上帝榮耀的.
一個很重要的表述.
一個神學裡面.
給我們一定要知道的上帝的豐饒.
如果你有機會.
你回去仔細再看.
我快速的給你們知道.
一章第6節.
一章12節.
和一章14節.
同樣出現一個的片語.
這個片語是什麼呢?.
就是使他的榮耀恩典得著稱讚.
一章第6節裡面.
他就這樣說.
就如神從創立世界以前.
在基督裡揀選了我們.
使我們在他面前成為聖傑.

$^{241}$無有瑕疵.
又恩愛我們.
就按著自己的意旨所喜悅的.
預定我們藉著耶穌基督得而止的名分.
使他的榮耀恩典得著稱讚.
這恩典是在他愛子女所賜給我們的.
這是一章4至6節.
同樣出現榮耀得著稱讚.
是一章11至12節.
我們在他裡面得了基業.
這原是那位隨己而行.
作萬事的.
照著他旨意所預定的.
叫他的榮耀.
從我們這首先在基督裡有盼望的人.
可以得著稱讚.
榮耀得著稱讚.
在第二次裡面重覆.
第三次重覆.
一章14節.
這聖靈是我們得基業的憑據.
直等到神的民被贖.
使他的榮耀得著稱讚.
在這裡面三次提及.
上帝的榮耀得著稱讚.
為何要這樣講論呢?.
其實原來這裡是在說.
用一個三一神的框架.
用三一神的神學.
來將上帝的恩典來展述出來.
一章4至6節講到.
無論你是甚麼背景的人.
不論你的國籍.
你的膚色.
你的背景是怎樣.
你領受了耶穌基督.
你其實是被上帝揀選的一班子民.
上帝在創世之前.
已經將這個兒子的身份.
賜給每一個屬他的子民.

$^{281}$你和我坐在這裡.
不單止是在說.
旺角浸信會這裡有一個席位.
有一個位置給你.
在天國.
在上帝的聖靈冊裡面.
上帝預備了一個席位給你.
這個席位是創造萬物的主位.
我們所預備的.
我們的身份是尊貴.
教會就是一班聚集.
神兒女的一個群體.
這班人就是要將上帝的榮耀.
藉著這班人來表彰出來.
保羅再繼續說.
這個席位帶來我們生命有甚麼好處.
或者甚麼呢.
一章第十一節就說到.
我們在裡面得到基業.
和合本是譯得到基業.
是拿到一些產業的.
不過更加好的翻譯.
不是說得到基業.
很被動式.
上帝給你一些東西.
在這裡面是用一個字眼.
有些譯本.
我們在他裡面成為基業.
我們在這裡不斷不斷地.
成為上帝的產業.
上帝在萬有為首.
耶穌基督是萬有之主.
但是視為他的產業.
視為他的基業.
去尊重的是一班基督徒.
是教會.
我們這一班人.
承受了上帝很重視的身份.
上帝看我們是寶貴.
上帝看我們是一班.

$^{321}$他所喜悅的子民.
我們能夠承受這個產業.
在這個產業裡面.
因為我們承受了.
所以我們生命裡面.
有一種很穩妥的盼望.
我們將這份盼望.
在基督裡一一.
透過生命實現出來.
然後你讀上去.
牧師很深奧.
又說產業又說過賬.
是不是真的.
是不是騙案.
本來他都怕你不明白.
然後他用你的生活體驗.
再一次去確據給你知道.
一章十四節.
用聖靈的工作告訴你.
這個產業是真的.
這個聖靈與你同在.
當你信主.
聽道.
信道.
然後領受聖靈.
聖靈就像一個罩一樣.
吸入你的生命當中.
不能脫離.
但聖靈時時刻刻與你同在.
聖靈的督責提醒.
聖靈將上帝的豐盛.
呈現在你的心靈當中的時候.
你就會經歷到信仰.
那種豐盛.
這個是上期.
保羅說.
得贖的日子.
再體會的那個更加大.
但在現世裡面.
聖靈與每一個基督徒.

$^{361}$與教會同在.
他就給你一個基業的憑據.
拿著這個基業.
拿著這個東西.
你才經歷到聖靈的工作.
以致到.
你就可以將榮耀.
這個群體.
將上帝的榮耀.
藉著我們的生命.
呈現出來.
所以在第一章裡面.
第三至第十四節.
他講的是一個這樣的神學.
就是上帝將所有的豐富.
他藉著上帝的群體.
藉著教會.
展現出來.
將我們生命.
我們脫胎換骨.
或者將上帝豐富流露出來的恩典.
藉著一群破破爛爛.
一群未必讀書很多.
甚至乎一群自己都覺得自己不是很像樣的一群人.
上帝竟然將這個這麼豐厚的恩典.
用三一神的框架.
來呈現給每一個讀二忽所書的弟兄姊妹.
當我們讀的時候.
我們會覺得上帝很榮耀.
不過我再要繼續問.
今天我們有沒有經歷到這麼豐盛的恩典呢?.
我們今天的教會生活.
我們今天的教會體會.
是一個怎樣的生活呢?.
因為近來疫情關係.
除了看我們自己的網站.
自己的YouTube.
我自己也會回去.
每個星期日都會看自己的.
要看看有什麼改進.

$^{401}$有時候要看看其他同工.
有什麼地方要改良.
看完自己之後又看看其他教會.
又看看別人怎樣領師.
學習一下.
不過我發現有一種很特別的現象.
原來很多領師也好.
或者講道也好.
很有趣的.
原來我們今天久而久之.
有一種很特別的習慣.
原來我們很多時候.
那個講道或者領師的起點.
我們不是由高位出發.
我們很多時候的起點.
是由低位.
由低壓槽的方式來做起點.
你覺得我很奇怪.
不知道我在講什麼.
你留意一下.
很多時候我們去講我們的救恩.
我們第一件事.
我們不是在講上帝的榮耀.
我們通常是在講我們生活很慘.
不是批評.
不過有一種觀察.
很多時候一見到一些頂梓梅領師的時候.
那種形態很悲壯.
我們會覺得這個世界已經氣壞了.
我們覺得我們自己是一群受苦被迫迫的受害者.
我們經常覺得自己很慘.
一群很慘的人聚集在一起唱慘歌.
唱情歌.
但一邊唱的時候流下幾滴眼淚.
心裡面就會覺得一邊唱的時候.
雖然唱的詩歌是很有盼望.
但帶著絕望的心.
來唱這些充滿盼望的詩歌.
甚至一邊唱的時候.
我們覺得自己好像被這個世界所壓迫.

$^{441}$我們是一個沒有招架能力的一群人.
我們心靈裡面.
雖然堂堂上口我們要有盼望.
但心靈深處裡面.
對我們自己.
對教會.
對我們身處的地方.
充滿了絕望.
我們覺得我們毫無招架之力.
我們覺得我們自己只是一群受害者.
不知道你會不會有這樣的體會.
我們從低位出發.
怎樣高都好.
我們都繼續是低位.
唱完之後.
我們心裡面的悲壯感加劇.
但對世界改造的盼望.
或者那種生命裡面流露出來的那種決心.
我們沒有辦法.
有力量來支撐.
如果我們久而久之.
繼續強化這種心態.
我相信我們會集體抑鬱.
我們陷入一種絕望的境地.
以弗所書.
他要扭轉顛覆我們思考的方法.
在接下來的一章十五節到第二十三節.
接下來的幾點.
是值得我們.
藉著以弗所書重新去塑造我們自己.
一個榮耀盼望一個基督徒的身份.
藉著去仔細讀聖經.
提振我們發酸的腿.
垂下來的手.
或者覺得面前無招架之力的大環境之下.
我們還可以挺起胸膛.
有毅力有盼望的過我們的生活.
在第一章第十五至第十八節裡面.
保羅就開始.
按著一個這麼榮耀的三一頌.

$^{481}$上帝的恩典這麼偉大的神學.
他就化為他對一班弟兄姊妹的禱告.
一章十五至十八節是這樣說.
因此我既聽見你們信從主耶穌基督.
親愛的眾聖徒.
禱告的時候.
上提到你們.
求我們主耶穌基督的神.
榮耀的父.
將那賜人智慧和啟示的靈.
賞給你們.
使你們真知道祂.
並且照明你們心中的眼睛.
使你們知道祂的恩兆有何等指望.
祂在聖徒中得的基業.
有何等豐盛的榮耀.
在這裡裡面.
保羅當他在說一個榮耀上帝.
基督成為宇宙的中心核心.
將這個榮耀加給教會的時候.
他為普世的信徒.
聖徒來祈禱.
他求上帝.
這個耶穌基督的神.
榮耀的父.
將賜人智慧和啟示的靈.
賞賜給我們每一個.
這個靈有些人翻譯.
會不會是一種特質呢.
那個人很有智慧.
智慧的靈.
在舊約很多時候都是這樣翻譯的.
不過如果按著上文下理.
保羅去求的就是.
求聖靈將智慧和啟示賜給我們每一個.
那個靈是在說聖靈.
聖靈將智慧和啟示.
將上帝自己揭示自己是一個怎樣的神.
告訴我們每一個.
當我們有這種智慧和啟示的能力的時候.

$^{521}$我們就能夠辨別得出.
究竟我們生命是一種怎樣的狀態.
更加重要的就是.
保羅在說.
聖靈這種的開啟能力.
就好像將我們本來的心地很夢迴.
心裡面好像是蒙了之又.
看不到東西的那個心懷念念.
現在變得明亮.
開啟我的心靈開竅.
以致我們看到上帝的恩兆有多麼的展望.
看到他的基業有多麼的豐盛.
在這裡裡面.
保羅的禱告.
或者我們第一件事.
我們要不停地問自己.
今天我們是否看得到上帝這麼偉大的工作.
還是我們只是注視到這個世界.
都是看到很多困難.
很多慘情的事.
覺得自己很可憐很悲慘.
悲悲慘慘地過著自己.
覺得按著這條路過我們的信仰生活.
保羅他要提醒我們.
基督徒的生命不是悲慘的生命.
基督徒的生命是一個充滿展望.
他沒有應許過我們所有事情是很順利.
不過他應許在人生不同的路程裡面.
不同的挑戰當中.
上帝都會為我們開一條新的出路.
上帝會在我們生命覺得困難裡面.
他會為我們開啟一條新的出路.
他在《箴言三章六節》說到.
我們凡事仰望他的時候.
他就會指引我們.
開為我們開一條新的出路.
我們對於任何的困境.
我們都不要帶著絕望的心態.
去看我們眼前不同的情況.
在預備這段經文的時候.

$^{561}$這一個多兩個星期.
剛好是我打柴的日子.
剛好好一點.
什麼打柴呢?很粗俗.
我覺得自己的狀態是很不濟的.
如果大家有在祈禱會裡面看過.
知道我媽媽入院.
那種心裡面就像坐過山車一樣.
父親節吃完飯後.
第二天就進了醫院.
突然間發高燒.
發高燒初時以為是做了一些化療之後.
一些副作用.
怎料晚上就馬上說.
我媽媽有肺炎.
馬上要入ICU.
入ICU都很擔心.
剛好我探完她之後回家.
醫院打電話來.
說入ICU又要馬上半夜趕進去.
那種感覺很擔心.
不覺得自己累.
只是覺得很擔心.
陪了媽媽一段時間.
那幾天就見到她很虛弱.
說話沒力量,眼睛垂下來.
身體很發性,人很疲乏.
見到她心裡很不安.
很不容易,慢慢燒就退了.
見到她雖然累.
但力量都稍為好一點.
醫生說她有好轉.
肺炎都慢慢可以清除.
就可以回到普通病房.
我打算慢慢會恢復.
一去到普通病房.
第一天馬上,沒多久.
又再發燒.
然後又很緊張.
那種感覺好一點.

$^{601}$然後又再下去.
但醫生說都不知道發生什麼事.
醫生都不知道怎樣.
有些白綢膜剪.
不知道發生什麼事.
令她又再發燒.
那時候心裡很擔心.
我捉住我媽媽的手來祈禱.
我媽媽都沒什麼力量來回應我.
我只說.
媽媽我們繼續.
黃浸有很多弟兄姊妹.
禱告紀念你.
你要加油.
不用擔心.
有病就醫病.
其實最擔心的是我自己.
回到家見到媽媽很虛弱.
又怕家人知道.
回到家受不了.
臨洗澡之前.
跪在地上大聲哭.
我和我的子女.
他們都不知道發生什麼事.
知道的,不過沒想過爸爸會這樣.
在禱告裡面.
我在問上帝.
究竟怎樣.
可以怎樣走.
上帝提醒我.
總有出路.
在禱告裡面有個領受.
上帝愛我媽媽比我愛媽媽更深.
那是祂的淚言.
儘管安心地交托.
日縮歲又酷日.
早辰定別歡呼.
聽了這句說話.
自己舒服些,但心裡也忐忑不安.
但我沒想過第二天.

$^{641}$醫院打電話給我.
很害怕的.
醫生告訴我.
你媽媽一晚之後.
第二天好多了.
燒退了.
她有精神.
我也不敢相信.
沒多久.
家人探望她.
我阿姨探望她.
我媽媽在電話裡面.
和我對話.
聽她的聲音.
很久沒聽她和我對話.
每一天都很虛弱.
整句說話很清晰.
我沒事了,好多了.
在教會那天早上.
關上房門.
自己又在爆哭.
腳軟了.
好像鬆了一口氣.
好像放鬆了一樣.
日縮歲又酷日.
早辰定別歡呼.
在我們覺得很大的困境.
上帝竟然會為我們.
開一條新的出路.
保羅在這裡.
他不是說病患這麼簡單.
他說人生不同的處境.
都是一樣.
不要認定了那個就是困境.
以後就是處於一種階段.
上帝會在我們生命裡面.
作一些我們都意想不到.
或者上帝的奇妙的工作.
他在我們生命裡面.
注入不同的盼望.

$^{681}$在我們面對不同的困難裡面.
上帝總要帶領我們走過.
面前的每一步.
所以在這裡.
說到基督徒.
如果你真的信靠上帝.
你的心靈.
理應是不會有絕望的.
我們信靠主.
所以我們不會絕望.
我知道上帝愛我們.
總會為我們開一條新的出路.
所以在這裡.
保羅特別去說.
他的恩召有何等的指望.
上帝召喚我們.
他不會召喚我們.
就任由我們去做.
他不會召喚我們.
就任由我們自生自滅.
他會用他的大能大力.
托住我們的人生.
幫助我們.
即使面對不同的挑戰.
我們都深信我們生命裡面.
有指望,有盼望,有出路.
我們可以.
藉著上帝的恩典.
走過一些我們覺得很難走的.
不同的路.
在2008年.
倫敦時報.
有一個作者.
他說自己是一個無神論者.
叫做Mathieu Paris.
他在一篇文章裡面.
他特別去說.
對一個無神論者的角度來說.
非洲需要神.
他用一個這樣的主題.

$^{721}$來寫一篇文章.
他講到非洲有很多困難.
有很多不同.
面對著很多的艱難.
一個很貧窮的國家.
去到今天還是面對很赤貧.
很多死亡的事件.
的一個地方.
他說這個地方.
他們需要基督徒.
講這番說話的無神論者.
他說為何會這樣講呢.
他說基督徒這班人是很獨特的.
這班人的信心.
總是為人類.
帶來釋放和舒緩.
在人很多的困難裡面.
這班人.
會帶來釋放和.
舒緩了很多的困難.
他們的生命力.
他們的好奇心.
他們對人的投入感.
能夠改變.
身處的環境.
這班基督徒.
是不是真的很厲害呢.
不是.
只不過是.
心靈裡面我們有上帝.
我們這個群體裡面.
有上帝的力量.
以致我們傻乎乎的.
本著.
耶穌怎樣愛這些人我們就去愛他們.
本著上帝給我們.
有那種生命力.
我們就在困難裡面.
活出那種生命力.
本著那種.

$^{761}$人總是複雜.
很多樣化的.
我們就去了解他.
為人帶來更加好的祝福.
這班人.
就是一班這麼獨特的人.
在世界裡面很多人都很有心.
但能夠展現.
這種生命力.
這個無神論者.
就是基督徒.
他畫了個框.
說是非洲需要.
不過這個世界.
我們深信不是非洲需要.
整個世界都需要.
上帝將這份恩召.
放在一班.
屬天的兒女身上.
放在教會裡面.
我們這班.
信主的人聚集.
我們就是為這個世界.
帶來更加多的盼望.
藉著我們生命的活力.
就將上帝.
那種何等大的指望.
呈現給這個絕望的世界.
知道.
不單止說一個何等大的指望.
保羅再繼續說.
他要將在.
耶穌基督在聖徒裡面的基業.
怎樣封成那種榮耀.
透過我們來展現出來.
我們這一班基督徒.
沒有人敢說自己很有錢.
不過我們那種用錢的方式.
就像一個敗家子一樣.
我們真的.

$^{801}$完完全全.
我們知道.
我們所擁有的東西.
其實不是屬於我們的.
是上帝給我們託管的.
我們就將我們所擁有的東西.
和身邊的人來分享.
我們的基業.
為何可以彰顯上帝的榮耀.
就因為我們不吝惜.
我們不是小家.
計算那一元幾毫.
用了之後就沒有了.
我們將所擁有的.
我們願意和身邊的人分享.
我們的心胸是廣闊的.
我們知道一切都是由神而來.
我們一切都是按著上帝的心意.
來去.
免自己的能力.
來去服侍上帝.
我們是基督徒.
來去服侍身邊的人.
這就是將基業.
何等豐盛.
藉著基督徒 藉著教會.
來去呈現出來.
早兩三個星期有機會和一個.
資深的牧者聊天.
我見到他開了很多時工.
很闊 做的事.
很有意思.
而且做的事全部都是大事.
然後我很好奇問他.
牧師啊.
其實你做這麼多的時工.
你哪裡找資源來呢?.
你怎樣搞呢?.
每件事都這麼大談.
你做這些工作.

$^{841}$你沒有資源怎樣做呢?.
然後牧師就笑了一下.
他就說.
你知不知道我爸爸是誰?.
我不知道.
不太懂.
我爸爸你也懂.
天父啊.
我沒有啊.
我爸爸有很多啊.
他真的開玩笑的說.
任何開展不同的工作.
除了資源之外.
更加要有廣闊的胸襟.
你更加要知道.
叫我們做的不是人叫我們去做.
是天父叫我們去做.
只要我們認定.
是上帝叫我們去做的工作.
的時候.
上帝就會打開天上的窗戶.
供應我們一切所需.
我們就只管.
帶著這份信心去回應上帝.
我們的爸爸是誰?.
天上的父.
當我們認定.
是天父叫我們去做的.
供應我們的時候.
我們就不再需要.
跟著計較.
不再像芝麻綠豆一樣計算到盡.
或者自己的心裡面.
帶著這種狹義的心.
去看今天的福音工作.
眼前如何去用錢.
而是帶著一個更加闊的天國觀念.
來實踐我們的照明.
為世界注入更加多的盼望.
為世界裡面彰顯神的豐富.

$^{881}$保羅不單止在說.
指望和基業的榮耀.
19至21節.
他繼續的來講.
能力.
19節.
他繼續的來講.
能力.
19節.
他繼續的來講.
能力.
19節.
他繼續的來講.
能力.
19節.
他繼續的來講.
能力.
19節.
他繼續的來講.
能力.
19節.
他繼續的來講.
能力.
19節.
他繼續的來講.
能力.
19節.
他繼續的來講.
能力.
19節.
他繼續的來講.
能力.
19節.
他繼續的來講.
能力.
19節.
他繼續的來講.
能力.
19節.
他繼續的來講.

$^{921}$能力.
19節.
他繼續的來講.
能力.
19節.
他繼續的來講.
能力.
19節.
他繼續的來講.
能力.
19節.
他繼續的來講.
能力.
19節.
他繼續的來講.
能力.
19節.
他繼續的來講.
能力.
19節.
他繼續的來講.
能力.
19節.
他繼續的來講.
能力.
19節.
他繼續的來講.
能力.
19節.
他繼續的來講.
能力.
19節.
他繼續的來講.
能力.
19節.
他繼續的來講.
在接下來的篇幅.
或者一直發展下去.
在接下來的篇幅.
這份所書裡的三大最重要的能力.

$^{961}$第一就是.
道德的能力.
罪惡滿盈的社會.
空中掌權者.
或者有很多罪惡觀.
纏繞的.
世界裡.
我們這一班人可以.
力排眾議.
成為社會上的清流.
可以過道德的生活.
成為社會上的清流.
聖靈給我們有道德的能力.
聖靈也給我們有第二種能力.
這種復活能力就彰顯在.
合一的生命當中.
分化.
崩解的人際關係.
互相拉扯的那種.
不同的關係當中.
我們所說的要破碎這些牆.
基督徒不論.
不同的境地.
我們都可以成為一個合一的群體.
我們.
合一的能力.
展示出.
上帝的榮美.
這兩個能力.
之後我們的同工會繼續.
告訴大家在講道裡.
不過第三個能力我們今天再繼續.
去講的就是盼望的能力.
在很多的挑戰.
在很多不同的困難裡.
很多的.
我們覺得不可能的當中.
剛才我說過.
我們還可以.
帶著從神而來的盼望.

$^{1001}$好好地.
活出一個.
有盼望人生的能力.
弟兄姊妹.
今天你信了主.
你稱了耶穌基督為主.
你有沒有一份.
總會深信.
眼前不同的處境.
總會有改變的可能.
我更加要去講.
我們的命是屬於上帝的.
沒有上帝的准許.
任何這個世界裡.
所有的罪惡.
負面的能力.
都不能夠在我們身上.
阻礙我們.
騷擾我們.
我們的命是屬於上帝的.
我們要有這份.
這樣的確據.
當然有些.
悉經學家說這種能力會不會抗衡.
當時以弗所.
在小亞細亞裡那些邪靈精靈.
邪靈的.
能力都會有的.
不過不單止.
是說這些東西會侵害我們.
世界上的罪惡.
或者不同的價值觀念.
可能都侵害著我們.
基督教會的那種生命.
教會的那種生命.
不過我們知道.
我們的命是屬於上帝的.
上帝總會在我們.
不同的難關裡.
開啟一條新的出路.

$^{1041}$我們只要.
有信心地.
堅定地去仰望他.
行上帝要我們.
要走的每一步.
不講晦氣的話.
在挑戰當中.
不做令人洩氣的事.
不針對.
身邊的人.
不再狗咬狗骨.
不再在別人的.
痛苦裡再要.
撩動別人一夜.
反而在別人的痛苦裡.
注入更多從神而來的.
盼望.
我們就成為一個.
有力量.
祝福別人.
不僅不讓外面的晦氣.
負面侵害我們.
我們更加心靈裡知道.
我們認定主一定會.
帶領我們走過人生.
不同的處境,不同的階段.
問題總會在此時此刻.
挑戰總會有一天會過去.
但上帝的恩典.
卻永遠長存.
上帝對我們的恩兆.
上帝給我們的能力.
永遠長存.
最後一個段落.
一章22至23節.
保羅再用這一句.
來作一個很重要的總結.
「有將萬有伏在他的腳下.
使他為教會作萬有之首.
教會是他的身體.

$^{1081}$善那充滿萬有者所充滿的」.
在這裡說到.
剛才他承接著.
在歌羅西書裡面所說的.
基督是萬有之首.
這裡再說一次.
耶穌基督是萬有之首.
宇宙萬物都是伏在基督耶穌的腳下.
這個我們明白了.
我們知道了.
不過保羅不是純粹沒有意義地重複.
這裡他作為萬有之首.
萬有之首和教會有什麼關係.
這個才是他的重點.
既然萬有是屬於他.
被他管轄.
教會也是屬於萬有的其中一份子.
不過在云云的萬有被創造.
或者創造的生命有機體當中.
上帝選召.
和上帝和耶穌有一個特別的連線關係.
就是教會.
所以說到教會是耶穌基督的身體.
他不是單純說一個描述性.
他是在說云云的眾生裡面.
這個宇宙萬物裡面.
上帝選召教會群體.
成為他的身體.
基督是教會的頭.
教會是成為他的身體.
藉著這個身體來顯出上帝的榮美.
更重要的是.
這個身體除了是一個比喻之外.
更加說到上帝將萬有的豐盛.
灌載在這個群體當中.
所以充滿萬有者所充滿的.
教會是一個很有前景的群體.
因為我們所信的主.
是一個這樣的主.
他將萬有的豐盛.

$^{1121}$竟然集中在我們這個群體裡面.
保羅說這個教會的觀念是不得了.
我們過往不會這樣看教會的.
我們覺得是信徒俱樂部.
講一個牧師講到悶悶悶.
那些培肚好不好.
我們用一個機構來看教會.
不過保羅不是這樣看的.
他在說教會是一個榮耀的身體.
是一個被充滿萬有者所充滿的.
一種生命力的群體.
在這裡特別提醒我們.
今天不要輕看我們身處這個群體.
我們極力好好維繫這個基督徒群體的生活.
你跟我一起隔離那個.
雖然有時奇人異事很多.
性格巨星特別多.
妄想特別多.
不過那個是上帝所看好的.
那個是上帝所愛的.
那個是上帝擺在我旁邊.
跟我一起過這個信仰生活的一個人.
上帝珍惜我們每一個.
他叫我們在這個群體裡面.
一步一步地來成長.
我們知道在座可能有些新朋友.
或者有些弟兄姊妹.
可能你來到旺鎮想聽一下道.
但你沒有想加入這個教會群體.
你不是很容讓自己的身份和這個群體有些結連.
但他們都是想和教會保持一個安全距離.
如果你是這樣我告訴你你走了一步.
因為最豐富的就是群體的生活.
聽道不是塑造教會最豐盛的生命.
聽道和一班人一起行道.
那個就是教會精彩的地方.
在我們人生有不同的處境裡面.
有人抵禮我們.
有人搓我們.
我們有一些盲點在群體裡面顯露出來.

$^{1161}$這些都是我們成長很重要的契機.
我也知道可能有些弟兄姊妹.
可能說牧師你有所不知.
我也曾經試過.
可能在旺鎮也好或者在其他教會也好.
我們經歷過受傷害.
我只是想找一間最好的教會安身立命.
我也明白的.
不過我再要說.
如果你想找一間最好的教會.
可能我們很難找.
等到耶穌回來的時候才能等待.
今天活在這個地上的教會.
每個人都是破破爛爛的人.
旺鎮的牧師也破破爛爛.
我們生命也是千瘡百孔.
有很多失去信心的地方.
有時候會有些罪.
有些救我都依然纏繞著我們.
不過我們這群人雖然是這樣.
上帝竟然沒有嫌棄我們.
上帝竟然用自己的生命.
用他的靈.
灌注在我們這個.
一個破破爛爛的群體裡.
上帝沒有嫌棄我們.
我希望大家不要嫌棄你身處的這個群體.
真的.
我們生命裡每一個教會.
總有不同教會的問題.
我也不是一味唱好.
說旺鎮厲害.
旺鎮是全世界最好的教會.
我不敢說這件事.
但我可以真誠的去說.
當我們每一個真誠在這個群體裡.
打開我們的心扉.
讓上帝動一下自己的生命.
真誠的建立這種弟兄姊妹的關係.
勉勵自己以如掩霞.

$^{1201}$活出我們生命裡最大的祝福.
那個教會生活是最美好的.
不需要找到一間最好的教會.
跟著像工廠一樣製造一個一式一樣的信徒.
我們就帶著我們今天這個現況.
去到一個今天這個現況的教會.
被聖靈持續的塗造,更新.
我們的生命就會慢慢趨向成熟.
我們就是在一個這樣的同步之下.
被塑造成為上帝喜悅的子民.
成為上帝喜悅的一個群體.
藉著一個破破爛爛的群體.
被上帝的靈充滿.
顯出上帝的榮耀.
今天這篇信息.
其實當你解釋清楚.
你會覺得不是很複雜.
我希望眾弟兄姊妹今天開這個會堂的時候.
你謹記一句說話.
教會是這個世界的希望.
我們這個群體不是因為做到最好.
而為世界帶來希望.
我們成為世界的希望是因為.
上帝在我們當中.
上帝藉著祂所設立的教會.
彰顯祂的榮耀.
我更加寄望.
旺角浸信會是我們身邊淪捨的希望.
藉著你的愛心.
藉著你的投入.
藉著你在我們當中.
一步一腳印證上帝的恩典.
我們破破爛爛的生命.
在這些破口當中.
就滲露出上帝的光芒.
你的光可以照在人群.
發光地照亮這個.
不是很顯眼.
或者很多不同困難的世界.
求主繼續幫助我們.

$^{1241}$激勵我們注入更多的希望.
為我們自己.
為這個世界.
我們同心祈禱.
天父上帝我們多謝你.
多謝你選照我們成為你的兒女.
多謝你愛我們.
你將榮耀的恩典賜給我們眾人.
我們要打開我們的心扉.
迎接你在我們生命中的塑造.
聖靈求你在我們生命中提醒我們.
有時我們迷失,絕望.
或者灰暗的時候.
求你的光光照我們.
以致到我們生命中.
同樣得著你給我們的盼望.
又將這份盼望帶給我們身邊的人.
我們恭敬將我們每一個交託給你.
耶穌求你掌管.
求你引領.
求你充滿我們眾人.
多謝主上聽祈禱.
奉基督耶穌得胜的聖命祈求.
Amen.
\newpage



\section{以弗所書 2:4-10}
\label{sec:iBJ_D2WmvXo}
\textbf{2020年7月12日 崇拜直播｜以弗所書2\_4-10}
\newline
\newline
連結: \href{https://youtube.com/watch?v=iBJ_D2WmvXo}{\texttt{https://youtube.com/watch?v=iBJ\_D2WmvXo}} ~~~~ 語音日期: 2020-07-13
\newline
\newline
\hyperref[sec:PBWvRjXaPlA]{\small{< < < PREV SERMON < < <}}
~
\hyperref[sec:index_chronic]{\small{[返順時目]}}
~
\hyperref[sec:hIi89f_Cp4Q]{\small{> > > NEXT SERMON > > >}}
\newline
\newline
以弗所書 2:4-10
\newline
\begin{longtable}{cl}
\hline
\hline
章節 & 經文 (和合本修訂版)\\
\hline
2:4 & \begin{tabularx}{0.7\textwidth}{X} 然而，神有豐富的憐憫，因著他愛我們的大愛， \end{tabularx} \\ \\ \relax
2:5 & \begin{tabularx}{0.7\textwidth}{X} 竟在我們因過犯而死了的時候，使我們與基督一同活過來—可見你們得救是本乎恩— \end{tabularx} \\ \\ \relax
2:6 & \begin{tabularx}{0.7\textwidth}{X} 他又使我們在基督耶穌裡與他一同復活，一同坐在天上， \end{tabularx} \\ \\ \relax
2:7 & \begin{tabularx}{0.7\textwidth}{X} 為要把他極豐富的恩典，就是他在基督耶穌裡向我們所施的恩慈，顯明給後來的世代。 \end{tabularx} \\ \\ \relax
2:8 & \begin{tabularx}{0.7\textwidth}{X} 你們得救是本乎恩，也因著信；這並不是出於自己，而是神所賜的； \end{tabularx} \\ \\ \relax
2:9 & \begin{tabularx}{0.7\textwidth}{X} 也不是出於行為，免得有人自誇。 \end{tabularx} \\ \\ \relax
2:10 & \begin{tabularx}{0.7\textwidth}{X} 我們是他所造之物，在基督耶穌裡創造的，為要使我們行善，就是神早已預備好要我們做的。 \end{tabularx} \\ \\
[1ex]
\hline
\hline
\end{longtable}
$^{1}$今日要和大家誦讀的經文是在新約的已忽所書第二章第八到第十節.
新約的已忽所書第二章第八到第十節.
你們得救是本乎因,也因著信,這並不是出於自己,乃是神所賜的,也不是出於行為,免得有人自誇..
我們原是他的工作,在基督耶穌裡造成的,為要叫我們行善,就是神所預備叫我們行的..
願神的話語常在我們心中..
弟兄姊妹平安,今日和大家分享一個主題,就是我們每一個都是上帝的傑作,你們認同嗎?.
我們每一個都是上帝的傑作,說起傑作,話說有一天我回到家,我兒子當我一進門口時,就對我說,.
爸爸,你看看我手上這件Lego,這是他親手造的機械人,你看看這隻手,你看看這武器,這關節,是不是很厲害,我自己造的,是不是很厲害?.
我看到他一卷嘴,說得眉飛色舞,於是當他停下來時,我就對他說,我說,這件真是你的傑作,.
他就說,是的,沒錯,於是他就繼續發表怎樣設計,到底這樣做有什麼目的,有什麼原因,說不停地跟我分享這件傑作..
於是我就說,爸爸說到,不如借你這件作品來分享,好嗎?.
他說,好啊,讓我再改裝一下,修改一下,於是他就埋頭苦幹地繼續砌他的作品,所以現在你看到手上這件和Powerpoint那件已經不同了,又進化了..
《鼎之回國實》,我想在很多人眼中,這件傑作很平凡,不過在設計者的心目中,他真的經過精心的計劃,他真的花時間去打造它,可以說是他一件愛不釋手的傑作..
同樣地,我想透過今天這段經文來提醒,真的,你和我,不折不扣,是上帝親手打造,精心設計的一件作品..
確實我們可能都很平凡,我們都未必有很多優越的本領或者成就,甚至乎可能在生活裡面我們會有很多艱難挫折,我們可能都會被一些身處的環境或者際遇來打低,甚至乎去踐踏..
我們有時都會覺得自己活得多麼無助,多麼淒涼,又或者我們總是有些缺憾,缺點我們改不了,我們生命裡面總會多多少少有失敗過,有軟弱過,有時覺得做人都幾不濟,幾糟糕..
不過,聖經經常提醒我們一件事,就是這一切其實都無法否定上帝在創造的時候,他已經將一個無比的價值賦予給我們..
無論你的生活是怎樣,你現在的情況是怎樣,上帝在創造你的時候所賦予給你的價值都不能夠被奪去,被打倒..
什麼意思呢?舉個例子,我手上有張什麼?500元,價值500元的紙幣,我將它拿一拿,插在一起,它現在值多少錢?有沒有少到?沒有少到,都是500元..
撕爛它,不要這麼快去到這麼盡,我將它扔在地上,踩兩腳,再拿起來,那它價值多少錢?都是500元..
好了,待會我去街市買菜,我去買海鮮,不小心掉在地上,你知道海鮮那層通常是怎樣的嗎?很濕的,很腥的,又濕又腥,又髒,撿不撿?撿,當然撿,沒理由不撿..
你拿去海鮮檔買一條東星斑,給它500元,它收不收?收,一定收,因為它還是一張什麼?500元,沒有變,它本身的價值不會因為它外在的骯髒而折損了,就是這個意思..
確實人生會有很多不完美的地方,但是在上帝眼中,你仍然是貴重,是他親手打造,是他重價買回來..
今天的經文二弗所書第二章第十節說,保羅說我們完事他的工作在耶穌基督裡造成的.工作這個意思可以說是一件精心創造的作品,又或者是一個親手做的藝術品..
保羅說,我們每一個都是上帝匠心獨運地創造的傑作,並且上帝要透過耶穌基督來將我們達至一個完善的地步,保羅是這樣說的..
事實上,如果你有留意整個二弗所書第二章,其實保羅在描述上帝在我們的生命裡,他是怎樣打造你的,並且他為什麼要打造你..
一至三節我們看,看回前一點,保羅在這裡說,其實很簡單地描述了我們信主前的境況,第一句已經說什麼?我們原本是死在過犯裡面,死在過犯之中,我們活在一個糟糕的光景裡面..
一方面我們被社會的風俗,我們被社會的價值觀牽著鼻子走,第二方面就是我們會信服於空中掌權者,或者有人在封音裡面所說,我們會順從著這個世界之王的勢力之下,我們做了很多活逆上帝的事情..
第三方面,他講到我們自己是活在血氣,一個放縱情慾的情況裡面,隨著你的肉體或者心中的喜好來行事.所以,保羅在這裡很清楚地說,和其他未信耶穌的人一樣,我們在上帝的眼中可不可愛?不可愛的,可憎可怒就有,是可怒之子..
就是說,上帝的審判,那個嚴厲的審判,最終必定臨到你和我的身上.本來我們簡單說就是死定的,死定的..
不過,在第四節這裡,保羅有個轉捩點給我們看到,然而,你見到然而,即是可以轉彎的.保羅在這裡說,我想大家一起讀,第四節到第七節,預備起..
好,謝謝.我們很熟悉的一段經文,這裡說,因為上帝他豐富的憐憫和他對我們的大愛,透過耶穌基督在十字架上寫生的拯救,使我們有一個翻身的機會..
不過,保羅也很清楚地說,上帝將救恩賜給我們,他不是單純想幫我們撿回一條命,讓我們將來死後可以上天堂..
保羅,你有沒有留意,他一連三次說了,神叫我們與基督一同活過來,一同復活,一同坐在天上,並且這三個動詞是用過去式..
意思就是,不要搞錯,我們不是要等到將來真的死後上天堂,才真真正正和耶穌聯合..
保羅說,其實從我們信主的一刻開始,我們就要和耶穌一起過一個出死入生的生活模式..
什麼意思呢?其實剛才一至三節,保羅說了我們的慘況,但是當你信了耶穌,和耶穌聯合之後,就不同了..
從前,你是被世界牽著鼻子走的,現在你是跟隨耶穌的帶領,不再被世界的局勢或價值觀帶領你走.這是第一方面..
第二方面,從前我們是信服著空中掌權者,但是現在我們已經和耶穌一同活過來,他讓你擁有勝過撒旦,勝過一切勢力的力量..
換句話來說,他可以讓我們擺脫魔鬼撒旦的權勢,可以勝過罪惡..

$^{41}$第三方面,從前我們是照著自己的肉體或想法行事,現在我們卻可以效法耶穌基督的一個寫記信服的心腸,就好像他在克西瑪利園的禱告一樣,不再照自己的意思,乃是照上帝的意思..
就是說,我們在耶穌的內部要成為一個新造的人,意思是我們藉著耶穌基督的拯救,藉著我們和耶穌基督的聯合,上帝好像,通常如果你電腦開機,發覺不能啟動,輕了又會怎樣?.
就會有一個重設的功能,將你還原到原廠的設定,上帝也要透過耶穌基督將我們重設,恢復你和我的原廠設定,使我們的生活現行,符合上帝原本創造你和我的心意..
不單止如此,背後的原因是上帝要透過我們吸引其他人認識耶穌,上帝要透過我們吸引其他人信靠耶穌,上帝要重新打造我們,成為一件可以擺在人面前的藝術品,透過我們的言行吸引經過的人注意到,然後他停下..
在欣賞這件作品,並不是純粹讚作品本身有多麼美善,更重要的是他對創造者發出一份讚嘆,說這件作品這麼厲害,像大師級的傑作,背後創造的人一定不簡單,我想進一步認識是誰造的..
本來是這樣的意思,所以弟兄姊妹,經文說上帝關注的不是我們死後能否上天堂,而是我們能否在地上成為一件他心目中的作品..
到底我們在人面前是一件怎樣的作品,是上帝關注的,到底是吸引人走近上帝,還是遠離上帝,到底是幫耶穌省招牌,還是幫耶穌拆招牌,我們有沒有按照上帝的吩咐來行善,.
我們有沒有將上帝豐富的恩典,將耶穌基督的恩惠向別人顯現出來,保羅在說這件事..
當然,保羅不想人誤會我們是靠行為得救,所以在八至九節他先包底,他強調得救完全是出於上帝所賜的恩典,.
沒有一個人可以靠好行為來換取,賺到救恩,每一個人都只需要憑著信心,認耶穌為主,他就可以領受上帝所賜的救恩..
保羅很清楚說明這件事,我們所說的恩訊清儀,不過保羅亦表示,行善亦是每一個蒙恩得救的人,他應該要有的表現..
簡單說,行善就是你和我蒙恩得到拯救,一個很實際的證據,如果別人在我們身上看不到有好行為的話,我們確實很難在人前稱得上蒙恩得救的人..
所以第十節,保羅強調,特別我選了一個新普及的譯本,他譯得更加容易明白,他這裡說,我們是上帝的傑作,他在基督耶穌裡重新創造了我們,讓我們能夠做他早已安排給我們做的善事..
他講得很清楚,耶穌基督的拯救就像上帝重新創造你,將你重新設計,來進行上帝早已安排你做的未善事,上帝早已定了..
所以這裡說,行善並不是我們一行一善,也不是我們按自己的想法做好人好事那麼簡單,行善是在說,當你重新實踐上帝創造你的時候,其實他早已在你的生命裡計劃好,安排好一些事情,他很想你去回應..
保羅在說這個意思,就像有一位信徒,他叫史懷哲,不知道大家有沒有聽過他的名字,德國諾貝爾和平獎,也有另外一個叫非洲之父的美譽史懷哲,他的見證是這樣的,他從小就精通兩種的語言,一種是德文,一種是法文,他是一個德國人..
並且他擁有音樂的天賦,他是一個教會的風琴師,不說笑的,他23歲的時候就已經收了哲學和神學的學位,24歲就成為牧師,這位我們說是音樂家,哲學家,神學家,牧師一身的他..
30出頭正是他聲望日昂,正是他虔誠感受的時候,他突然間就拋開這一切的成就,轉而去學醫術,直到他38歲,當他取得醫生的資格之後,他就前往非洲去宣教,到底什麼令他有這麼大的轉變?.
他自己就說,在他29歲的時候,當他讀到一篇關於非洲地區一些民生疾苦的報導的時候,他就說,他知道在非洲大陸的叢林裡面,其實生活了很多不信耶穌的土人,沒有傳教士來幫助他們,也都在他們生病的時候,沒有人醫治他們,給藥物他們..
真的不是,他身邊的親朋戚友都紛紛給他負評,紛紛責備他,有些人為了他在學術上或藝術上的成就而可惜,有些人覺得他是不是傻傻的,是不是神經有問題,但是最令人覺得困惑的時候,是為什麼他要去非洲裡面行醫..
其實以他本身的條件,他已經可以直接去宣教,他又是神學家,又是哲學家,還要是牧師,他直接去非洲裡面傳道,可能更快更實際,不過他自己分享就是他覺得用醫生這個身份來服侍,可以直接幫到他所服侍的對象,因為醫生是直接看病,直接醫病..
不用說這麼多,在日常的生活裡面,可以很直接,身體力行地去宣揚神的愛,耶穌基督的恩典,他是這麼說的..
所以弟兄姊妹,我們看到,假如我們看回30歲前的史懷哲,其實他算不算是上帝手中一件很漂亮的作品,算不算?已經很好了,就算放回現代,都已經很傻了,真的見得人了..
不過原來,在永恆的角度裡面,上帝說你可以再好一點,你可以再被上帝雕得漂亮一點,成為一件更加美善的藝術作品,上帝是這樣給他一個意象..
上帝很想他更加服侍不同的人,向身邊的人更多的顯現上帝豐盛的恩典.同樣弟兄姊妹,其實上帝也要讓我們成為一件更美善的作品,真的..
無論你覺得這一刻,可能你已經是一件不錯的藝術品,可能你已經覺得自己沒有進步的空間,但不是的,在上帝手中,他仍然可以將你雕得更加美善,他仍然有用得著我們的地方..
而事實上,當我們越對上帝保持開放的特化,如我們越是順應上帝的帶領,我們越是放手讓上帝雕得你的生命,我們就越能夠成為上帝心目中那件喜悅的作品..
就好像大家有沒有聽過米高安哲奴這個人,有吧?你可能不認識他,但你也認識他所雕刻的大衛像,有沒有印象?給大家看一看,因為他是一個全身裸體的雕像..
大衛像可以說是舉世知名的傑作,沒有人會反對,但原來大衛像在米高安哲奴開始動工雕刻之前,其實原本是一塊被人廢棄的大理石,為什麼呢?因為上手有個雕刻師不經意失手,在這塊大理石中間鑿穿了一個洞,破壞了大理石的中間..
破壞了大理石的完整度,換句話說,理論上這塊大理石是不能夠再用的,用不著..
不過當米高安哲奴知道這件事之後,他就嘗試去遊說當時的首長,把這塊有缺陷的大理石交給他,交給他創作一個大衛像,他要挑戰一個似乎所有人都覺得不可能的任務..
於是在1501年,當時26歲的他就開始動工,把一些很高的圍板圍住大理石,不讓人看到他如何創作..
然後他就開始動手,人們經過可能每天每夜都只聽到他在創作石頭的聲音,一直在創作.沒有人知道這件作品到底是什麼樣子的..
到了1504年,四周的圍板就打開,一座前所未見的大衛像就聳立在人們的眼前.為什麼說前所未見呢?因為以往雕刻大衛像的藝術家,你猜他們通常都會做一個什麼樣的造型呢?.
我們記得大衛戰勝過哥利亞,原來以前的雕刻家通常會雕刻大衛手拿著巨人哥利亞的人頭,在說一個勝利的場景..
不過米高安哲奴這個大衛像就很不一樣,他不單單是一個全身赤裸的大衛像,他還要是一個手拿著頭石器,看到他都眉頭深鎖的樣子,雙眼仍視著前面的一個年輕人..
有學者分析,這個大衛像面容緊張,但他的姿態卻是放鬆,換句話說,他呈現了一種強烈的對比,就像一個準備好要和哥利亞戰鬥,準備要在瞬間爆發出強大力量的大衛..
所以這件大衛像簡直是動與靜結合,力與美都很平衡,一件大師的傑作..
之後有人問米高安哲奴,為什麼他可以這麼厲害,雕刻出這麼偉大的作品?你猜猜他怎麼回答?他回答,其實我沒有雕刻過他,我只是將不需要的部分削走了..
將他禁錮在石頭裡的生命釋放出來.厲害吧?大師就是大師..

$^{81}$當有人問他怎麼去設計,他回答,我沒有設計過,其實他本來就在裡面,不過我只是將他外面的多餘,不需要的部分拿走,他裡面的生命自然就出來了..
同樣的,弟兄姊妹,其實上帝已經將他美善的旨意一直放在你和我裡面,放在了.創造你的時候已經定好了,只不過有什麼遮住了他?.
你和我的罪,你和我的自我,讓他沒有辦法釋放出來.所以,你和我都需要在人生裡面持續的接受上帝雕琢你的生命,將外面多餘的石頭,就是那些難阻住我們跟隨,信服上帝的那些可能是你的想法..
是你的性格,是你的自我,你要讓上帝慢慢割走它,讓耶穌基督的光芒可以透過你向別人重新展現出來..
真的,弟兄姊妹,其實我們坐在這裡,代表我們仍然一色尚存,上帝仍然讓我們在這裡的時候,亦代表什麼?他在我們身上的雕刻工程完了嗎?還沒完..
所以,我們要對上帝有那份耐性,你也要信得過上帝.意思就是,有沒有帶小朋友去剪頭髮?或者可能你自己小時候都經歷過家人幫你剪頭髮?.
當你坐定定,不亂動的時候,就會怎樣?快,靚,正,就會剪完,對不對?.
但假如剪頭髮的過程中,你經常坐不定,又動不動身,又怎樣?有兩個情況,要麼剪很久,又或者剪錯了,剪衰了,剪得不靚仔,我們要承擔這個後果..
但我們可以放心的是,上帝比你和我更有耐性,上帝會慢慢等,等你的頭髮長了,等你下次坐定定,不再亂動的時候,上帝會繼續幫你修剪..
上帝總有辦法將我們修剪得靚仔靚女,所以可以放心,放心讓上帝在你的生命裡繼續進行雕琢修剪..
最後,保羅再一次說道,我們原本是上帝的傑作,留意傑作這個字是單數,意思是上帝不是要每一個信徒分散,各自成為一件精彩的傑作,這麼簡單..
上帝是要照著自己的計劃,將我們每一個拼合成為一件真正的藝術傑作,所以留意,已忽所書不是寫給個別信徒,是寫給教會全體..
第十九至二十二節,新漢語譯本說道:.
「在這看來,你們不再是愛人和僑民,而是信徒的同胞,是神家裡的人.你們被建立在使徒和先知的根基上.耶穌基督自己就是基石,整座建築在他裡面連接在一起,不斷增長,成為主裡的聖殿.你們在他裡面也一同被建立,成為神藉著聖靈居住的地方.」.
保羅很清楚地說,被上帝建立從來不單單是個人的事,而是整個群體的事..
第十九至二十二節,保羅形容教會是一所上帝的殿,是聖靈的居所..
有耶穌基督成為教會最重要的基石,也有使徒和先知作為教會的基礎.然後,上帝也將外邦人,包括今天你和我,一個又一個從罪惡裡藉著耶穌基督拯救出來,.
將我們雕琢成一件適合上帝用,適合他的尺寸,適合他的心意的石頭,讓我們成為建立教會一塊的好材料,使教會能夠不斷地延續增長下去..
保羅向我們每一個人呈現這幅圖畫,所以弟兄姊妹,保羅是在提醒我們,你和我仍然是教會的一份子,所有人都需要投入參與教會的建立,沒有一個例外..
每一個人都要成為上帝心目中那塊石頭,那塊材料,被建造上去,讓我們可以成為上帝所喜悅的居所..
或許我們覺得我們很小,不太顯眼,微小,我們不是一塊什麼的好材料,但重點是什麼?重點是上帝說你適合我用,你適合上帝用,上帝要用你,不在於你怎樣看待自己,而在於上帝賦予你的價值..
所以他透過耶穌基督來拯救你和我,並且將我們放在一起,來為主發熱發光..
就好像我最近看一本書叫做《讓群體發光》,裡面的作者提到一位牧者叫做Thomas Smith,他在南非創立了一間塗糊教會,即是用陶瓷製造的糊..
他在建立的過程中,他和一群基督徒一起去禱告,並且用幾個星期來持續查考上帝的話語,最後他們停留在哥林多後書四章七節,其實我們很熟悉的這句話:.
「我們有者寶貝放在雅氣裡,要顯明者莫大,那些能力是出於神,不是出於自己.」.
他們被這段經文裡的雅氣和寶貝吸引著,於是他們很想做一個行動來回應上帝給他們的意象,他們很想在坊間找一個適合的塗糊,來作為創立這間教會的可見標記..
最後他們找到一個很貴重的塗糊,他們覺得很適合的塗糊是一個沒有人要的,被人丟到一旁,表面上還是很骯髒,有很多泥,還有幾條裂痕在裡面,他們找到這個塗糊,他們覺得很適合..
不單止這樣,他們又約在一起,可能拿一個袋子,帶著塗糊,然後弄碎它,弄到它碎了,然後每人帶一塊碎片回家,然後他們把自己的禱告寫在那塊碎片,然後又約在一起,將每一塊碎片重新組合,嘗試還原成為那個塗糊..
在黏合的過程,在重新組合的過程,會不會變回完美?不會的,一定比以前可能更加多的裂痕,對不對?.
不過,當他們將一支蠟燭放在裡面的時候,就能夠從這個不完美的器皿裡面,照耀出燦爛的光芒,很有意思,對不對?.
就正如地上每一間教會,每一個教會的群體,包括黃浸,都是不完美的,都有裂痕,有破碎,不過,只要有耶穌基督這個寶貝,放在這個不完美的瓦器裡面,住在我們中間的時候,.
祂的光芒就可以透過你和我的不完美,發放出來.所以,今天的經文一方面提醒我們,你和我不折不扣是上帝的傑作,上帝不單單只拯救了我們,不單單只要塑造我們,雕琢我們,上帝更加要像我們今年的主題一樣,總動員我們..
上帝要將我們所有人,拼合成一個合他用的器皿,來向人前展示出他的微善.真的,弟兄姊妹,我們絕對不是這麼巧合,因為交通方便,因為你住附近,所以在一起.而是上帝在背後精心設計過..
上帝認為我們在一起是合他用的,上帝要用我們,才將我們放在一起.他要透過往生來展現他的微善,所以上帝才讓我們來到這裡,留在這裡繼續成長..
所以弟兄姊妹,很盼望透過今天的經文,邀請大家在未來的日子,做三件事,不難的大家都做得到,甚至正在做.我們每天需要在上帝面前充電,可能你透過祈禱,透過你的讀經,或者其他你恆常的靈修生活,你要在裡面讓上帝雕琢你的生命,這是真的..
我們每一個都需要被上帝這樣更新改造.與此同時,我們也要思考,既然上帝呼召我們一起建立教會的時候,上帝一定有些東西,有些恩賜財幹給了你..
你一定有些東西在這裡是上帝用得著,適合用,可以榮耀祂,可以祝福身邊的人.你禱告,想一想上帝給了你什麼恩賜,並且不要停留在頭腦的階段.你再進一步想一想,你可以怎樣具體擺上,你可以怎樣具體擺上,而不是一千兩的那個埋在土地..
你可以怎樣將上帝給你的恩賜財幹擺上,來建立教會.這是第二方面.第三方面,我們有時忽略了,你要邀請人守望你..
我們不只是做,我們需要建立那個屬靈的關係,不單單是和上帝,也要和你身邊的弟兄姊妹.有時我們很少在團契或小組分享,我們怎樣毀身上帝,並且要弟兄姊妹守望我.我們很少有時..
但原來這件事不可以缺少,不可以缺少.特別如果你仍然未有小組生活,未有團契的話,如果崇拜完,你可以和我聊聊天,你也可以找其他傳道同工,我們一起來幫你,幫你尋找到你的小組相交生活,因為這件事是需要..

$^{121}$你要邀請你身邊的弟兄姊妹來定期為你代禱,來關心守望彼此之間的屬靈生命成長.這是第三方面.重複一次,你需要每一天來到上帝面前,支取力量,來充電..
可以是讀經,可以是靈修,可以是禱告,無必要讓上帝繼續刁鑽你的生命.第二方面,你要同時思考上帝已經給了你什麼恩賜財幹,讓你可以和其他人一起配搭建立教會,並且你怎樣可以具體來落實..
第三方面,你邀請身邊的弟兄姊妹為你去守望彼此的守望,讓我們可以互相去毀身,一起成為上帝合容的器皿..
就讓我們一起祈禱,求上帝幫助我們回應祂今天給我們的提醒,我們同心禱告..
並且不單止給我們永恆的生命,更加要我們和耶穌基督連結,為上帝作見證,成為上帝你心目中所喜悅的一個傑作..
上帝我求你藉著今天的經文幫助我們,讓我們首先能夠認定我們每一個在上帝裡面都是貴重無比,因為有耶穌基督用重價將我們每一個贖回,並且上帝在創造我們的時候已經將你永恆美善的心意放在我們的生命裡面..
上帝我求你幫助我們,讓我們願意保持接受,讓上帝你去雕琢我們的生命,將那些不合你心意,難阻你心意的石頭鑿走,讓我們每一個都可以成為上帝你適用的石頭,讓我們都可以願意放手,讓上帝你將我們每一個拼合在一起..
在教會裡面我們彼此去配搭,互相守望,一起順應耶穌你的帶領,來在人前成為燈台,成為炎,成為光,讓上帝你自己的美善可以藉著我們的行為,藉著我們的善行來彰顯,最終讓身邊的人也都像我們一樣,一同經歷上帝你的拯救,建造,讓教會繼續來增長..
求主你幫助我們,讓我們真的能夠願意落實去回應你今天對我們的提醒,也都一樣我們在過程裡面,我們和弟兄姊妹彼此去守望,互相的代禱,讓我們可以在當中能夠為上帝你去發熱,發光,我們這樣去禱告,交託,仰望,奉基督耶穌你得勝的名字來祈求.阿們..
\newpage



\section{以弗所書 4:1-16}
\label{sec:hIi89f_Cp4Q}
\textbf{2020年7月26日 崇拜直播｜以弗所書4\_1-16}
\newline
\newline
連結: \href{https://youtube.com/watch?v=hIi89f_Cp4Q}{\texttt{https://youtube.com/watch?v=hIi89f\_Cp4Q}} ~~~~ 語音日期: 2020-07-27
\newline
\newline
\hyperref[sec:iBJ_D2WmvXo]{\small{< < < PREV SERMON < < <}}
~
\hyperref[sec:index_chronic]{\small{[返順時目]}}
~
\hyperref[sec:4TNIVRU8_Gc]{\small{> > > NEXT SERMON > > >}}
\newline
\newline
以弗所書 4:1-16
\newline
\begin{longtable}{cl}
\hline
\hline
章節 & 經文 (和合本修訂版)\\
\hline
4:1 & \begin{tabularx}{0.7\textwidth}{X} 我為主作囚徒的勸你們，既然蒙召，行事為人就要與你們所蒙的呼召相稱。 \end{tabularx} \\ \\ \relax
4:2 & \begin{tabularx}{0.7\textwidth}{X} 凡事要謙虛、溫柔、忍耐，用愛心互相寬容， \end{tabularx} \\ \\ \relax
4:3 & \begin{tabularx}{0.7\textwidth}{X} 以和平彼此聯繫，竭力保持聖靈所賜的合一。 \end{tabularx} \\ \\ \relax
4:4 & \begin{tabularx}{0.7\textwidth}{X} 身體只有一個，聖靈只有一位，正如你們蒙召，是為同有一個指望而蒙召， \end{tabularx} \\ \\ \relax
4:5 & \begin{tabularx}{0.7\textwidth}{X} 一主，一信，一洗， \end{tabularx} \\ \\ \relax
4:6 & \begin{tabularx}{0.7\textwidth}{X} 一神－就是萬人之父，超越萬有之上，貫通萬有，在萬有之中。 \end{tabularx} \\ \\ \relax
4:7 & \begin{tabularx}{0.7\textwidth}{X} 我們每個人蒙恩都是照基督所量給每個人的恩賜。 \end{tabularx} \\ \\ \relax
4:8 & \begin{tabularx}{0.7\textwidth}{X} 所以有話說：「他升上高天的時候，擄掠了俘虜，將各樣的恩賜賞給人。」 \end{tabularx} \\ \\ \relax
4:9 & \begin{tabularx}{0.7\textwidth}{X} 既說「他升上」，豈不是指他曾降到地底下嗎？ \end{tabularx} \\ \\ \relax
4:10 & \begin{tabularx}{0.7\textwidth}{X} 那降下的，就是高升遠超越諸天之上的，為要充滿萬有。 \end{tabularx} \\ \\ \relax
4:11 & \begin{tabularx}{0.7\textwidth}{X} 他所賜的有使徒，有先知，有傳福音的，有牧者和教師， \end{tabularx} \\ \\ \relax
4:12 & \begin{tabularx}{0.7\textwidth}{X} 為要裝備聖徒，做事奉的工作，建立基督的身體， \end{tabularx} \\ \\ \relax
4:13 & \begin{tabularx}{0.7\textwidth}{X} 直等到我們眾人在信仰上同歸於一，認識神的兒子，得以長大成人，達到基督完全長成的身量。 \end{tabularx} \\ \\ \relax
4:14 & \begin{tabularx}{0.7\textwidth}{X} 這樣，我們不再作小孩子，中了人的詭計和欺騙的法術，被一切邪說之風搖動，飄來飄去。 \end{tabularx} \\ \\ \relax
4:15 & \begin{tabularx}{0.7\textwidth}{X} 我們反而要用愛心說誠實話，各方面向著基督長進，連於元首基督， \end{tabularx} \\ \\ \relax
4:16 & \begin{tabularx}{0.7\textwidth}{X} 靠著他全身都連接得緊湊，百節各按各職，照著各體的功用彼此相助，使身體漸漸增長，在愛中建立自己。 \end{tabularx} \\ \\
[1ex]
\hline
\hline
\end{longtable}
$^{1}$讓我們一起去聆聽一段經文.
已忽所書第四章一至十六節.
使徒保羅這樣去勉勵我們每一個.
我為主被囚的勸你們.
既然夢召行事為人.
就當與夢召的恩相稱.
凡事謙虛 溫柔 忍耐.
用愛心互相寬容.
用和平彼此聯絡.
竭力保守聖靈所賜合而為一的心.
身體只有一個.
聖靈只有一個.
正如你們夢召同有一個指望.
一主一信一浸一神.
就是眾人的父.
超乎眾人之常.
貫乎眾人之中.
也住在眾人之內.
我們各人夢因.
都是照基督所量給各人的恩賜.
所以經上說.
他升上高天的時候.
努力了受的.
將各樣的恩賜賞給人.
記說升上.
豈不是先降在地下嗎.
那降下的.
就是遠升諸天之上.
要充滿萬有的.
他所賜的有使徒.
有先知.
有傳福音的.
有牧師和教師.
為要承傳聖徒.
各盡其職.
建立基督的身體.
直等到我們眾人在真道上同歸於一.
認識神的兒子.
得以長大成人.
滿有基督長成的新娘.

$^{41}$使我們不再作小孩子.
忠了人的詭計和欺騙的法術.
被一切異教之風搖動.
飄來飄去.
周隨從各樣的異端.
唯用愛心說誠實話.
凡事長進.
憐於元首基督.
全身都靠祂聯絡得合適.
百折各按各職.
朝著各體的功用彼此相助.
便叫身體漸漸增長.
在愛中建立自己.
今天來到發放群體光芒的第四講.
有陳妙龍姑娘在我們當中分享.
剛剛所讀的這段經文.
講題是「合而為一」.
各位導演姐妹平安.
剛才為我們領師的姐妹.
不斷提醒我們唯有耶穌.
我記得這個月.
當我們主任穆斯第一講的時候.
他提到我們敬拜的時候.
雖然環境不是最好.
但是我們仍然存著一個喜樂.
我們不是悠悠愁愁去敬拜主.
因為我們所信的上帝.
祂的確能夠將風浪平靜.
我們一起祈禱吧.
慈愛天父.
萬國萬民都要來讚美你.
耶穌.
你明白我們每一個都有害怕的時候.
你賜下平安.
你又醫治我們.
釋放我們.
讚美主.
你能夠叫風浪平靜.
慈愛恩主.
我們憑信心.

$^{81}$在這個疫情之下.
我們同心合一.
倚靠你.
敬拜你.
因為我們知道.
我們活著就是要讚美你.
求你開我們的心靈.
讓我們聆聽你的話語.
並且踐行出來.
也都求你繼續保守.
我們所住的地方和世界各地.
讓我們能夠在不容易的情況下.
在裡面得著平安.
我們同心的禱告.
是奉耶穌基督.
聖命祈禱.
阿們.
今天和大家分享的時候.
我都希望在座和在家的弟兄姊妹.
都可以翻譯剛才的經文.
《已忽所書》四章.
我在網上看到一個.
尋找上帝的故事.
它說到有四個人.
四個好朋友.
都很想見到上帝.
他們各自對上帝都有自己的理解.
他們無法說服對方.
我們可以放在頭部影片.
他們無法說服對方.
後來他們被困在一個.
很幽暗的房間裡面.
牆是很高的.
只有一隻窗.
他們四個都深信.
他們能夠透過這隻窗.
能夠見到上帝.
這個時候.
他們就想了一個方法.
就好像陳陳諜騎駁馬一樣.

$^{121}$就可以碰到窗.
不過問題來了.
誰可以站在最上面呢?.
由於每個人都很想見到上帝.
所以大家都在等.
等誰開口說.
「我做最低的那個」.
想問大家.
他們怎樣能夠達到共識.
而不會令到大家的關係破裂呢?.
雖然是一個故事.
不過我不知道大家同不同意.
人與人的相處.
人際關係是一個大學問.
有人說.
人生根本就離不開關係.
而基督教基本就是表示關係的一個宗教.
而今天保羅就在《已忽所書》四章.
提到合一.
經文一開始.
其實是有原文.
是有「所以」的.
《新頁本》是這樣說.
「所以我這一個為主作囚徒的勸你們.
行事為人要與你們所蒙的呼召相稱」.
這個「所以」.
就很奇妙地將《已忽所書》頭三章.
保羅說到上帝的恩典和救恩.
正如過往三個星期.
我們三位很好的牧師.
為我們分享.
我們教會是世界的光.
是上帝的傑作.
是福音的使者.
然後連接到第四章.
當我們蒙恩之後.
我們怎樣生活呢?.
有牧者這樣說.
信了主就要向人承認.
而且在生活當中.

$^{161}$毀身出來.
剛才.
無論是牧師.
或者我剛剛讀的《新漢語》.
都提到.
保羅說到自己的身份.
是為主被囚.
如果我們翻到前面的經文.
《史徒行傳》八章三節.
保羅拉誰入監.
是呀.
保羅當時形容叫「掃羅」.
他殘害教會.
將男女拉入坐監.
現在反而是自己被囚.
保羅有什麼感受呢?.
在三章十三節裡.
保羅完全不覺得羞恥.
他覺得他為主被囚.
是一個榮耀.
正正是因為保羅為主被囚.
所以他講以下這番勸勉.
對當時信徒來說.
就很有說服力.
同時保羅都有一個提示.
就是提示當時教會的弟兄姊妹.
當他們持守信仰的時候.
都有機會被囚.
經文裡面提到「呼召」.
聽到「呼召」.
我們不知道會不會連線.
傳世間事奉.
或者是一些傳道人.
不過我們看回.
《以弗所書》二章十三至二十二節.
原來這個「呼召」大家都看到了.
就是我們呼召出來.
成為上帝的兒女.
而且成為一個整體.
經文裡面保羅鼓勵我們.

$^{201}$行兆為人就要和這個「呼召」相稱.
有點像一個天秤.
這個就是主呼召我們成為一個整體.
這裡就是我們的「行兆為人」.
「行兆為人」這個字原文.
原來是「行」.
Walk.
甚麼意思呢?.
出來走.
是否出來走呢?.
原來這個字在跟著《以弗所書》.
這幾章都出現了很多次.
無論是跟著第十七節.
五章的第二節.
十一,十五.
為甚麼它會不斷說呢?.
原來這個「行走」.
「行兆為人」.
就是說我們的日常生活.
我們這個天秤怎麼做好呢?.
既然主用付這麼大的代價.
呼召我們成為一個整體.
我們「行兆為人」是怎樣呢?.
原來保羅說得一句話就說完了.
就是甚麼?.
就是「竭力保守聖靈所賜.
合二為一的心」.
「聖靈所賜,合二為一的心」.
原文就是「聖靈的合一」.
就是指聖靈所賜的.
不用我們找的.
和聖靈帶來的.
合一的意思.
就是我們聯合一起成為一個整體.
竭力是甚麼?.
竭力就是很積極.
用盡一切的方法和努力.
希望能夠達到一個美善的目標.
是要付代價的.
今早我不知道大家有沒有聽到收音機.

$^{241}$以下這幅圖.
就是在網上看到的.
就是巴西的老人院.
發生甚麼事呢?.
今早大家都聽到了.
巴西的確診數字的確比較厲害.
我們都要為他們祈禱.
他們同樣因為疫情原因.
不容易探望長者.
這個新聞就是前一段日子的.
家屬他們很想見到自己所愛的長者.
於是老人院負責人.
他就想一些方法.
他用盡一些想法.
他在網上看一些片等等.
然後就想到有一個叫做雙湧隧道.
當然我看報道的時候.
他們做足防疫措施.
為甚麼要這樣做?.
他們的目的是很美善.
他希望能夠解家屬和長者相思之苦.
他用盡一些方法去想.
竭力就有這個意思.
我們就要問.
牧者,我們合二為一的心.
要保持,要用盡方法嗎?.
是這麼難的嗎?.
我們就要問一個問題.
大家覺得教會合一.
是否一件容易的事?.
教會在我面前.
都有好幾位的肢體和牧者.
教會是由不同的獨立個體組成.
教會人數越多就變得越多元.
多元就意味著我們能夠發現.
發展每一個個體的獨特性.
個個都是寶來的.
個個都是寶藏.
是不能夠替代.
多元,多元同時意味著.

$^{281}$我們個個就有自己獨立的思考.
我們立場,價值觀不一樣.
我們的種族,國籍,文化也不一樣.
家庭背景,生活方式也不一樣.
這就表示,在這麼多不一樣的時候.
我們比較難去找到共識.
因而就會產生分歧.
甚至分裂.
所以合一不是一件容易的事.
其中一樣就是自我.
史托德在一本書裡面.
勉勵我們做信徒.
不要被世界的主義所同化.
其中一樣他特別說的.
就是自戀主義.
原來自戀主義同時就叫做納西斯主義.
納西斯是誰呢?.
納西斯就是希臘神話裡面.
一個很英俊,偏偏風度的美少年.
他在詞中看到自己的倒影.
他深深愛上了自己.
他極度愛戀自己.
這是以極度自我為中心.
我想問大家.
當我們這麼自我的時候.
我們容不容易容納其他人呢?.
容不容易成為一個合一?.
說到這裡.
所以我們也明白了.
為何保羅要我們竭力地做.
保羅呢.
上帝也明白我們的需要.
所以保羅就提出一些新的素質.
大家螢光幕也看到了.
保羅提出第二節和第三節.
或者大家跟我一起讀好不好?.
好,一二三.
凡事謙虛,溫柔,忍耐.
用愛心互相寬容.
用和平彼此聯絡.

$^{321}$竭力保守聖靈所賜.
合二為一的心.
剛才我們說第一樣是謙虛.
謙虛我們參考回.
《臣妾記》二章三節.
在經文裡面他提到.
這些人他們看別人比自己強.
謙虛的人.
就是承認自己在上帝的面前.
沒有什麼可誇.
一切都是出於上帝的恩典.
是很恰當的自我評估.
這個自然就不是自戀主義.
希羅文化裡面.
謙虛就被視為自卑.
所以當時只有奴隸才有這個自卑的態度.
各位弟兄姊妹.
我們是謙卑.
我們不是自卑.
如果沒有謙卑.
我們很容易驕傲.
驕傲就是我們看不上其他人.
特別是我們跟我們立場.
想的很多不同的人.
而驕傲.
大家同不同意.
往往就是成為紛爭.
結黨.
分裂的其中一個很重要的原因.
溫柔是什麼呢.
溫柔我們可以看到.
提多書三章一至二節.
有些補充.
就是不毀謗.
不跟別人爭的.
忍耐.
大家都看到了.
是一個寬大的心.
可能大家即時就想到.
哥林多前書十三章裡面有補充.

$^{361}$原來忍耐就是.
忍無可忍.
都要忍.
而且在忍耐當中.
我們是沒有怨言.
不報復.
用愛心互相寬容.
其實就是講到忍耐.
一個實質的行動.
接著保羅就說.
用和平彼此聯絡.
我們想到聯絡.
不知道會不會想到.
好像打電話一樣.
大家看看我手頭上有一條繩.
原來聯絡這個字.
是綁住的意思.
和平即是和睦.
保羅其實是很.
生動化地提醒我們.
我們這條繩叫做和平.
就將我們大家.
緊緊地連在一起.
我們竭力保守這個.
合而為一的心.
弟兄姊妹.
我們要問的.
剛才保羅說的這些生命的素質.
我們多不多呢?.
很感恩.
和大家相處裡面.
我看到大家很多很美的特質.
不過因為我們都是軟弱.
我們都是罪人.
我們難免都會經歷到紛爭.
我們的關係可能變得僵化.
有機會在這兩個星期.
我和幾位肢體聊過.
在旺鎮合一的事情.
有一位肢體很坦誠地和我說.

$^{401}$原來早段日子.
他和肢體相處的時候.
其實不是什麼大事.
只不過大家可能是觀點,角度.
處事方法不一樣.
就發生了一些不開心的事.
小事.
小事就變成大家沒有理睬對方.
就變成大事.
這位肢體告訴我.
他當時覺得很不開心.
甚至哭了.
因為他覺得不被尊重.
因為他很體會到.
是自己對.
他在等一件事.
他在等對方認錯.
謙卑.
當時這位肢體.
他謙卑.
他去到主耶穌面前的時候.
他發現了一個情況.
原來自己也是一個罪人.
為什麼覺得一定要對方先和他道歉.
就是這個轉捩點.
就是這個謙卑.
能夠靠著主耶穌慢慢有改變.
剛才我也提過.
我和一些肢體聊天的時候.
我也問他們.
為什麼現在大家和好.
原來有一樣東西真的很重要.
就是你很想和對方和好.
未必是牧者勸你.
而是你的心.
你捨不得這個人.
你很想和他和好.
當我們知道要合一之後.
有時我們也會問.
為什麼要堅持這個合一.

$^{441}$有人說信念這個力量很大.
影響我們的行為.
不知道家庭的弟兄姊妹.
或者在場的弟兄姊妹.
有沒有聽過練習生這件事.
有些綜藝節目.
它說到一些藝人或者偶像.
他們出道之前成為練習生.
有些可能是訓練很多很多年.
那些人就訪問他們.
為什麼你們仍然堅持.
不止一個都是這樣說.
因為他們有一個很熱愛.
這個舞台的信念.
那我們要問.
是什麼信念令我們堅持.
要竭力保持這個合一.
是什麼基礎呢?.
剛才牧師也為我們讀了四至六節.
原來因為我們看到.
保羅說到有七個的一.
是我們的認順.
第一個是指到第四節.
一個身體.
我們看回《以弗所書》一章二十二至二十三節.
正正就是說.
教會就是基督的身體.
既然身體就是一個.
那不可能分開.
一聖靈.
我們可以參考二章十八節.
無論是猶太信徒和外邦信徒.
都靠著同一位聖靈.
來到父神面前.
而盼望是指什麼呢?.
我們看回第一章十八節.
裡面說到的盼望.
就是一個永遠的救恩.
是一個復活的盼望.
所以第一組這幾個一.

$^{481}$我們看到.
無論是外邦人和猶太人.
我們都是靠著同一個聖靈所感.
然後歸入基督的身體.
有著同一個盼望.
第二組它說到一主一信一浸.
一主就是主耶穌基督.
一信就是我們有共同的信仰.
和核心對象.
就是對著主耶穌基督.
而浸禮只有一個.
咦?那是不是表示浸禮的形式呢?.
原來不是.
是透過浸禮.
我們弟兄姊妹.
我們公開認信.
就好像正如《羅馬書》六章.
三至六節裡面所說.
我們公開透過受浸.
歸入主基督.
和祂聯合.
所以第二組的一就說到.
我們每一個.
都在信同一個福音.
我們都要靠著同一個.
主耶穌基督的血.
得蒙救贖.
然後也受同一個浸.
能夠歸入主耶穌基督.
最後第六節.
保羅說到父神.
父神是眾人之父.
是一個一神的觀念.
對於猶太人來說.
這個觀念絕對不陌生.
在《申命記》六章四節裡面.
他這樣說.
他說:以色列啊.
你要聽.
耶和華我們神.

$^{521}$是獨一的主.
舊約的聖經.
經常提醒他們.
千萬不要拜其他的偶像.
保羅在這裡提醒我們.
父神是大家的.
是眾人之父.
他用同一個方法.
去拯救我們每一個.
我們都是無分彼此.
來到祂的面前.
所以.
保羅就用這幾節的經文.
去提醒我們.
教會為什麼我們要堅守合一.
不是表面一個功夫.
也不是我們盲目地做.
是因為我們有一個.
同一認順的基礎.
或者我再重複一次.
其實合一.
聖靈已經賜給了我們.
我們教會.
我們弟兄姊妹.
我們要做的.
我們的責任.
就是我們很重視祂.
然後我就想到三寶.
就是保守.
保持.
保護祂.
剛才分享的經文.
就是四章一至六節.
我們發覺很有意思.
保羅明明用了七個一.
來凸顯合一的基礎.
但是在接下來的經文裡.
保羅就用「因此」.
將我們的角度.
由一到多.

$^{561}$為什麼我這樣說呢?.
因為第七節.
原文是有「但是」的.
在這裡告訴我們.
保羅不單只是注視合一.
他也看到多元.
也看到不同的出現.
在這裡保羅展示了.
一和多的關係.
合一是否容不下.
每個獨立的個體呢?.
是否教會真的只有一把聲音呢?.
是否社會只能有一個立場呢?.
家裡是否只有一個人作主呢?.
合一不是單一.
也不是統一.
合一裡有多元.
而多元又成為合一.
構成合一.
是一個很微妙的關係.
大家聽我讀第七節.
「我們各人蒙恩.
都是照基督所量給各人的恩賜」.
我補充一下.
剛才也說了.
有「但是」.
在鏡頭轉一轉.
我預備這個經文的時候.
上帝提了一句說話.
原來我的傳道人曾經問過我.
「喵喵,你有什麼恩賜呢?」.
當時的氣氛是很靜的.
我一時之間想不到.
我不知道弟兄姊妹有沒有試過.
想不到自己有什麼恩賜.
我們會否誤會了「恩賜」的意思呢?.
上一週.
Paul Mok 跟我們分享了.
原來「恩賜」是兩回事.
第一.

$^{601}$我們是要為別人的好處.
上帝給我們這些恩賜.
是為別人的好處.
第二.
就是正正建立教會.
經文裡第七節.
他所說的「恩」.
所說的「恩賜」.
原來是一些美好.
白白給我們的.
也指的是用以事奉的恩典.
在這裡我想勉勵大家兩件事.
都是提醒我自己.
無論我們看自己是怎樣一個人.
很厲害.
或者是低低沉沉.
最近靈性又不是很好.
又很懶惰.
但原來經文裡.
上帝透過保羅已經清清楚楚告訴我們.
我們每一個國人.
你和我都有份.
或者我們覺得小弟兄,小姊妹.
都有份的.
有份做什麼?.
有份加入事奉的行列.
第二.
上帝會賜一些屬靈的恩賜給我們.
叫我們每一個.
回到上帝的家.
都有祂的價值.
接著第八節.
保羅用到詩篇六十八篇十八節.
大家看完綱目.
可以聽我讀.
「你已經升上高天.
努力仇敵.
你在人間就是在叛逆的人間.
受了貢獻.
叫耶和華神可以與他們同住」.

$^{641}$經文六十八篇.
大衛讚美上帝.
讚美上帝.
叫他們得勝仇敵.
讚美上帝保守他的百姓.
經文裡我們看到.
耶和華是用一個君王的身份.
去接這個禮物.
但我們看回第八節.
保羅說.
把各樣的恩賜賞給我們.
就好像這個轉變.
就好像耶穌的恩賜.
是一個禮物.
把禮物賜給我們.
似乎有點不同.
所以.
《色經學者》有不同的見解.
其中有學者指出.
當我們教會.
當我們每一個.
不嫌自己的恩賜大或少.
多或少.
只要我們運用這些恩賜.
去服侍的時候.
原來就是成為禮物.
獻給上帝.
十一節.
它提及到有幾種的恩賜.
裡面提及到有使徒.
先知.
傳福音.
牧師和教師.
大家數數.
然後熟悉聖經的大家.
會否想起.
哥林多前書十二章.
好像有提到.
例如醫病.
辨別諸靈等等.

$^{681}$為何這裡不提及呢?.
我們相信這裡.
主要是保羅提出.
幾個恩賜的代表.
使徒.
在當時來說.
就是授權.
是上帝授權的代理人.
他們用來解明.
傳揚福音.
使徒行傳十一章和二十一章.
我們看到有些先知.
真的會預言將來的事.
不過通常.
當我們說先知的時候.
原來他們主要的職責.
就是向人解明福音.
和上帝的心意.
傳福音.
我們現在聽得多了.
原來傳福音.
這個寫法.
這個職份.
我們在經文裡看得不多.
不過我們參考一下.
原來它是形容.
肥膩和提摩太.
我們想一想.
肥膩和提摩太的對象.
主要是一些非信徒.
重點就是.
向人說好消息.
牧師.
在我面前也有一些牧師.
原來牧師.
是一個牧羊人的意思.
就是帶領和看顧著.
我們這群經常揹的羊.
所以總括來說.
我們看到保羅選擇這幾個恩賜.

$^{721}$這幾個不同的恩賜.
對內來說.
就是一個牧羊.
一個教導.
但對外來說.
就是傳揚福音.
我一輪嘴說了這麼多.
弟兄姊妹.
如果現在我再問你.
你有什麼恩賜.
你想到.
你的恩賜.
是對外多.
還是對內呢.
你是牧羊.
教導.
傳福音.
還是關懷呢.
另一個問題.
我也有一點.
挑戰一下大家.
不知道大家有沒有羨慕.
其他人的恩賜呢.
我剛才看到弟兄姊妹.
牧者和姊妹.
練習的時候.
我聽到一些高音.
我不但完全聽不到.
我音就音了.
別人高音.
我就走音.
我又看到.
我自己的同工.
他們在關懷.
在傳福音.
或者在教導上面.
都比我強.
可能我們就會想.
為什麼.
我們有些恩賜.

$^{761}$是給這個.
或者給這個呢.
耶穌給我們有不同的恩賜.
我就想起.
我們當了一個牧者的名言.
不是用來比較和計較.
目的是什麼呢.
第十二節.
不如大家在家裡.
和面前的弟兄姊妹.
和我一起讀.
第十二節.
一二三.
為要成全聖徒.
各盡其職.
建立基督的身體.
原本的經文裡面.
成全呢.
可能我們.
我們牧師叫名全.
成全我們又聽得多了.
但好像有一點點抽象.
原來成全呢.
說得簡單一點.
就是裝備了.
而各盡其職.
就是我們每一個信徒.
我們都從事.
聖公.
我們是事奉主.
原來目的只有一個.
為什麼這個唱歌能夠高音.
為什麼這個報道.
有音賜.
為什麼這個講信息.
能夠讓我們.
很容易聽得明白呢.
原來目的不是要他們威風.
不是要我們去.
捧他們做偶像.

$^{801}$或者.
我不知道大家有沒有試過.
其實我自己都會.
有一些講完.
我們會追的.
追看他們的信息.
其實沒有問題.
他們的努力我們看不看到呢.
不過問題就是.
如果我們追到.
我們看不到背後.
祂有上帝的恩典.
因為剛才我都講了.
恩賜是神伯伯給我們的.
如果我們看不到.
其實上帝都給我們恩典的時候.
我們就.
真是.
沒有拿到這一份這麼美麗的禮物.
所以弟兄姊妹.
原來.
神給我們每一個都不同.
原因就是.
建立基督的身體.
在這裡我們覺得.
保羅真是很奇妙.
他簡簡單單用幾節經文.
就將多元和合一連結起來.
多元.
不就是構成合一嗎?.
建立是需要時間.
亦都是一個漫長的過程.
目標是什麼呢?.
目標就是.
我不知道弟兄姊妹同不同意.
我們做一段時間.
不就可以實體回來的.
因為真是.
我們.
各個牧者都有自己崗位.

$^{841}$我主要在五樓.
看一看兒童.
其實大家做得很好.
看一看.
嘩.
我發現.
高了.
即是高了.
有些就瘦了.
有些又.
可能大隻了.
但是.
原來經文裡面很簡單.
教會.
我們有目的之後.
我們的目標就很簡單.
長大成人.
長大成人達至什麼?.
達至我們在真道上.
同歸於一.
我們聽上去好像有少少恐怖.
原來它講到.
就是我們在主裡面.
我們有同一個信仰和認識.
而這個認識.
是我們每一個.
都要一起.
都要和祂建立關係.
正如我一開頭所講.
是要走出來的.
是要毀身.
我們怎知道.
這間教會長大成人呢?.
是不是建築物越來越漂亮.
或者是怎樣呢?.
或者人數多.
或者他們質素怎樣呢?.
原來以下的經文.
很簡單.
講了兩點.

$^{881}$第一.
長大成人的教會.
就能夠分辨是非黑白.
長大成人就和小孩子不一樣.
當然.
我們現在的小孩子.
都不容易受騙的.
不過保羅在這裡的確是想講.
小孩子是容易受騙.
他講得很生動的.
好像被風擺下擺下.
搖擺不定.
其實他想講到.
如果一間教會.
不長大成人的話.
他就會對基督的認識.
缺乏真正的認識.
那就很容易中計.
很容易中了一些傳異教的人的計謀.
十五節經文裡面.
原文亦有提到.
是用反而開始的.
即是說.
保羅不打算提我們.
一些負面的事.
他提一間教會長大成人.
原來不中計不止.
每一個弟兄姊妹.
都是用愛心說誠實話.
我以前每次聽這一句話.
都會想.
甚麼叫用愛心說誠實話呢?.
你的誠實和我的誠實.
可能不一樣.
可能有人說.
甚至坊間有人說.
甚麼善意的謊言.
或者包裝了的謊言.
是為對方好.
所以我們參考.

$^{921}$《迦拉太書》四章十六節.
《迦拉太書》四章十六節裡面說.
原來說誠實話.
是講真理.
不是生你的氣的真理.
是真理.
意思是.
我們在生活上.
我們用愛心說誠實話.
就是我們在生活上.
我們走出來的.
不單是在教會.
在任何地方.
甚至對自己都是.
我們有愛.
但都要有真理.
那怎知道真理是甚麼呢?.
當然我們就一起努力.
去認識聖經.
亦都求上帝去幫助我們.
教會能夠全新全面發展.
原來是因為我們連在同一個頭.
是誰呢?.
不是把任何頭貼上去.
是否貼上我們主要目標的頭.
或者貼上你們一直.
牧養你們的道師的頭.
或者過往不同年代的聖徒.
或者一些我們很欣賞偉大的信徒.
是否把他們的頭放上去呢?.
原來不是.
原來經文裡面第十五,十六節.
就講得很清楚.
就是連於主耶穌.
因為主耶穌是滿有能力.
主耶穌是滿有權柄.
祂明白教會真正需要的是甚麼能力.
去面對這個世界.
主耶穌亦都有權柄.
祂叫我們真是好好地去面對.

$^{961}$我們之間的紛爭.
有時當我們不合一的時候.
我們回想.
回想原來上帝就是有這個.
祂有這個命令.
祂有這個提醒.
就是叫我們合一.
但另一方面.
保羅很奇妙.
他說原來這個身體長大成人.
要被建立.
大家看看自己.
是你和我都有份.
是我們各按各職.
彼此相助.
一個都不能少.
我自己想一想.
這麼多的牧者.
我暫時第一個.
都是限聚令.
四個人.
闊正道的.
其實我有時想.
如果我講這篇信息.
看到大家每一個的時候.
我就更加體會到.
甚麼叫做一個呢?.
都不能少.
所以各位弟兄姊妹.
合而為一.
原來能夠有多元的.
原來就是我們不同恩賜的弟兄姊妹.
我們恩著主耶穌基督.
我們憐於祂.
我們共同承擔一個責任.
這個責任就是建立基督的身體.
接下來有少許照片.
給大家看看.
這張是甚麼呢?.
這張就是.

$^{1001}$我記得第一次我們.
幼年部拍片.
其實大家看片.
有時可能都覺得.
好像可以.
有時覺得一般般.
不過不知道大家知不知道.
其實背後.
是有很多人付出.
由寫信息.
錄音.
打燈.
剪片.
甚至我們有一些片裡.
有不同家庭的參與.
我聽到有些家長.
為了大家一起建立基督的身體.
他都要哄兒女拍片.
當然不是要我們威風.
我們不是用來報學校.
但是我們都明白.
有時我們有些內斂.
即是我對著鏡頭其實不容易.
大家可能以為我很喜歡對著鏡頭.
其實不是.
我都有我的內斂.
不過當我想到.
我們一起.
在這個特別的時刻.
我們繼續去傳揚主.
去牧養弟兄姊妹.
我們繼續連繫的時候.
這些少少的害羞或是欺負.
我們就沒有所謂.
這個朱古力是甚麼呢?.
就是不同的家庭.
你看他們多有心.
當我叫他們拍片的時候.
他們真的買齊了.
而且我相信廚房應該一向都很乾淨.

$^{1041}$又有齊所有材料等等.
好了.
接著另一個鏡頭.
早兩天有支體給我訊息.
說他正在看網絡祈禱會.
還有他剛剛拍的那集.
就是我現在給大家看的這一集.
我想他們是否好的東西就回味了?.
應該是的.
但是我想說甚麼呢?.
平時大家見到我們幾個人.
但看旁邊的照片.
原來是這麼大隊團隊.
當你們在網上.
不是很聽到有雜聲的時候.
原來連關冷氣,關風扇.
都要大家合作.
有些可能對聲音特別敏感的.
他們就去負責.
有些可能是在鏡頭上拍得比較好.
我都記得.
以我今天穿的顏色衣服為例.
也是因為專業人士提醒我.
不要經常穿白色衣服.
那我要聽話的.
因為大家都是為了建立一件好的事情.
不過問題來了.
究竟因此.
有沒有分時間性或潮流性呢?.
會不會某些因此被其他因此所代替?.
我自己預備這個經文的時候.
我發現是沒有的.
不過就在某一個時段裡.
可能某些的伺候是比較合適.
但是我想問大家.
如果我們.
先數一數.
我們七個牧者和.
Belly也是小牧者.
或者加上傑成等等.

$^{1081}$我們跳火圈.
我們做什麼都好.
但是如果沒有大家背後一起祈禱.
一起守望.
祈禱,守望都是因此.
那我們拍這麼多片.
做什麼?.
和綜藝節目有什麼不同?.
我又試問.
如果我們只是拍片.
但是我們很多弟兄姊妹.
不是只靠一個畫像.
她可能心靈有重擔.
她正正就是需要你這個.
這麼有關懷因此的人.
給她一個電話.
所以說弟兄姊妹.
真的.
因此我絕對不覺得是怎麼分.
分什麼呢?.
沒有什麼好分.
不過就是要用.
是不是?.
要用.
還有在這裡我希望大家不要害羞.
原來你在構思很多方面.
你是很厲害的.
你播劇,播道.
原來你發覺.
喂.
不要用喂.
不好意思.
表表姑娘.
原來我.
哎.
崇拜片.
我給你一個想法.
給你一個建議.
你這樣拍.
我真的很歡迎.

$^{1121}$因為神給我們的恩賜.
就是我們覺盡奇蹟.
令整個教會很美.
不過教會美不重要.
重要的是什麼呢?.
重要的就是一開始我都提過.
教會是基督的身體.
當教會我們能夠合二為一的時候.
我們就能不折不扣.
我們以基督的身體.
這個姿態出現.
我們站在這個分裂.
大家不容易合一的社會裡面.
我們去宣告這個福音.
我們讓別人能夠相信主.
上帝的榮耀.
我們就正正回應這個口音.
或者用我最後大膽的說.
成也教會.
敗也教會.
當教會分裂的時候.
正正濟嚴重.
衰弱了教會的使命.
但是當我們實踐合一的時候.
教會正正濟彰顯了.
上帝是這個世上.
令人和睦的力量.
一開始的故事.
相信大家都不忘記.
好像現在看這幅圖一樣.
騎駁馬.
無分大小.
我們一起去建立教會的身體.
不過我講回這個故事.
過了一會兒.
阿甲就說.
不如我做最低那個.
接著月,丙,幸運兒就是丁.
於是他們就合手合腳.
大家騎駁馬.

$^{1161}$結果丁真的碰到一把槍.
他就向外望.
但是因為大家不知道.
是否留在家裡留得多.
每個人都變得大了.
所以很快就倒下了.
這個人柱.
我想問大家.
甲,月,丙會問丁一個什麼問題呢?.
是的.
你們說得對.
他們就是問丁.
你能不能看到上帝?.
丁說.
我看到很多東西.
我在短短的一段時間裡.
我透過槍看到很多東西.
但是我看不到上帝.
當時的氣氛很沉重.
突然間.
甲說了一件事.
我醒了.
當我在最低的時候.
你們三個.
在壓著我的時候.
當我盡力支撐著大家的時候.
我看大家是很重要的時候.
我看到上帝.
接著月也有同樣的意思.
到丙了.
丙這樣說.
其實我剛才差點就碰到一把槍了.
不過我回想.
我為了希望丁.
你能夠真的看到上帝.
我自己就甘願.
「呃?」推你上去.
因為你對我很重要.
而在這個時候.
我看到上帝.

$^{1201}$最後丁這樣說.
原來我們一直以為要在窗外.
看到上帝.
但是當我想到.
就是你們每一個.
將機會給我.
建立我的時候.
支撐著我.
我就在你們身上.
看到上帝.
弟兄姊妹.
我們合二為一.
我們讓別人看到上帝.
讓別人看到這種力量.
今天我們找一個時間.
主席會為我們唱回應詩.
就是「同心合二」.
最後我想說一個小經歷.
當我選這首詩歌的時候.
我是選普通話版的.
但是我親愛的同工.
就提一提我有廣東話版.
接著我跟大家說.
我完全沒有任何考慮.
我就想.
好啊.
就唱廣東話版.
為什麼我沒有考慮呢?.
因為不用比.
同工給我在聖樂的知識這麼豐富.
我一定覺得沒問題.
雖然我前一晚.
其實我兩邊的歌詞.
我都有看到.
自己有一個選擇.
到後來我打給今天的領師.
我們談電話.
我們談談天.
大家猜不猜得到.
領師的姐妹.

$^{1241}$她想唱廣東話版.
還是普通話版呢?.
你們猜對了.
她也是想普通話版.
但是她願意.
願意因為整件事好.
她問我.
我怎麼想.
我自己祈禱了之後.
我結果還是選了廣東話版.
這位姐妹.
她絕對有資格.
跟我繼續爭辯.
就是這麼小事上.
大家可以爭辯.
但是她沒有.
因為正如詩歌所說.
當我們放下自己的時候.
神就在我們當中.
別人就看到上帝.
然後我們一起祈禱.
「慈愛天父 讚美你」.
「讚美你 你沒有嫌棄我們每一個」.
「讚美你 你竟然將這個彰顯你和睦的力量」.
這個職份交給我們.
求你幫助我們.
我們難免都會有紛爭.
我們難免都有傷害過別人.
所以主耶穌.
你赦免我們.
我們難免都會有分裂的時候.
但是主多謝你.
你透過經文再一次提醒我們.
我們憐於你.
我們有一些生命的素質.
當我們知道我們的目標.
就是去見證你.
建立你的身體的時候.
求聖靈加我們的力量.
改變我們.

$^{1281}$寬恕我們.
讓我們重新.
真是同心合意去見證你.
讚美主你自己.
我們同心去禱告.
去奉耶穌基督的聖命祈禱.
\newpage



\section{以弗所書 6:5-9}
\label{sec:4TNIVRU8_Gc}
\textbf{2020年8月16日 崇拜直播｜以弗所書6\_5-9}
\newline
\newline
連結: \href{https://youtube.com/watch?v=4TNIVRU8-Gc}{\texttt{https://youtube.com/watch?v=4TNIVRU8-Gc}} ~~~~ 語音日期: 2020-08-18
\newline
\newline
\hyperref[sec:hIi89f_Cp4Q]{\small{< < < PREV SERMON < < <}}
~
\hyperref[sec:index_chronic]{\small{[返順時目]}}
~
\hyperref[sec:TUMQvpeyy3M]{\small{> > > NEXT SERMON > > >}}
\newline
\newline
以弗所書 6:5-9
\newline
\begin{longtable}{cl}
\hline
\hline
章節 & 經文 (和合本修訂版)\\
\hline
6:5 & \begin{tabularx}{0.7\textwidth}{X} 作僕人的，你們要懼怕戰兢，用誠實的心聽從你們肉身的主人，好像聽從基督一般； \end{tabularx} \\ \\ \relax
6:6 & \begin{tabularx}{0.7\textwidth}{X} 不要只在人的眼前這樣做，像僅是討人的喜歡，而是作基督的僕人，從心裡遵行神的旨意， \end{tabularx} \\ \\ \relax
6:7 & \begin{tabularx}{0.7\textwidth}{X} 甘心服侍，好像服侍主，不像服侍人， \end{tabularx} \\ \\ \relax
6:8 & \begin{tabularx}{0.7\textwidth}{X} 因為知道每個人所做的善事，不論是為奴的或是自主的，都必按所做的從主得到賞賜。 \end{tabularx} \\ \\ \relax
6:9 & \begin{tabularx}{0.7\textwidth}{X} 作主人的，你們待僕人也是一樣，不要威嚇他們，因為知道他們和你們在天上同有一位主，他並不偏待人。 \end{tabularx} \\ \\
[1ex]
\hline
\hline
\end{longtable}
$^{1}$弟兄姊妹仰我們繼續直注聆聽上帝的話語.
來敬拜祂.
今日所誦讀的經文是.
《耳窟所書》第六章五至第九節.
史圖保羅這樣說.
「你們作僕人的.
要懼怕戰經.
用誠實的心.
聽從你們肉身的主人.
好像聽從基督一般.
不要只在眼前事奉.
像是討人喜悅的.
要像基督的僕人.
從心裡遵行神的旨意.
甘心事奉.
好像服侍主.
不像服侍人.
因為曉得.
各人所行的善事.
不論是為奴的.
是自主的.
到別案所行的.
得主的賞賜.
你們作主人的.
待僕人也是一理.
不要威嚇他們.
因為知道他們和你們.
同有一位主在天上.
他並不偏待人」.
今天早上我們會繼續.
宣講《耳窟所書》.
如果大家還記得.
其實早幾天.
剛剛剛過去.
星期三晚的祈禱會.
三位的姐妹.
一個叫阿珍.
一個叫阿眉.
一個叫阿妙.
珍 眉 妙這三位姐妹.

$^{41}$其實和我們分享了.
《耳窟所書》第五章裡.
彼此信服的一個吩咐.
並且如何實踐在夫婦.
和親子的關係上.
所以今天我們會繼續看.
保羅如何將彼此信服的教導.
延伸到主人和僕人的關係裡.
當然我們首先要問.
為何使徒保羅要求.
耳窟所教會的信徒.
要在主僕 夫婦.
並且親子關係上.
學習彼此信服呢?.
首先我們要明白.
當時一些社會背景.
我們說在希臘羅馬的文化裡.
其實一向由男士男性主導.
在家裡的丈夫.
那位父親.
並且是家主.
他們是話事.
是屬於強勢的一方.
至於其他的家庭成員.
就好像他的妻子.
他的兒女.
和他家中的奴隸 僕婢.
其實是處於一個被管轄.
被轄制.
需要絕對服從的一個弱勢.
不過.
隨著福音的擴展.
這種局面慢慢開始被打破.
因為信了主的人.
無論他是男還是女.
無論他是猶太人還是外邦人.
無論他是自主的.
或者是遺奴的.
現在都變成上帝的兒女.
大家在神的面前.

$^{81}$都是平等的.
只不過因為這樣.
問題就開始來了.
信了主的丈夫.
那位的父親.
那位的主人.
是否還可以像以前一樣.
來核制他的太太.
他的兒女.
和他家中的奴隸呢?.
他們已經是主內的弟兄姊妹了.
又或者反過來.
已經信了主.
在耶穌基督裡面.
得到釋放自由的那一位太太.
那一位的兒女.
或者說家中的奴隸.
是否仍然要像以前一樣.
信服他的主人.
他的丈夫.
或者說他的父親呢?.
這些其實都是當時.
很多基督徒的家庭.
又或者教會.
很實際會面對到的難題.
所以使徒保羅.
有需要教導他們.
如何彼此相待.
確實在屬靈上.
主,目,雙方.
真是對等的.
是平等的.
但是保羅沒有因為這樣.
就否定了.
他們已經有的職份.
和他們的責任.
剛才所讀的經文.
第五節說.
「你們作僕人的.
要懼怕戰驚.

$^{121}$用誠實的心.
聽從你們肉身的主人.
好像聽從基督一般」.
這裡說恐懼戰驚.
不是說到.
怕得顫顫顫的那種狀況.
而是說一份很恭敬.
很小心謹慎.
恐怕有負所託的態度.
而誠實的心是說.
要帶著一份純正純潔的動機.
全心全意為別人謀求好處.
這當然是所有信徒.
對耶穌基督所要有的態度.
但是使徒保羅也勉勵.
作僕人的.
也要以這種態度.
來聽從他的家主.
這樣做.
和他的家主的品格.
人品.
沒有直接必然的關係.
不是說他敬我一尺.
我就敬他一丈.
老闆對我好得無可奈何.
所以我自然要甘心.
命底地服侍他.
又或者反過來.
如果家主經常苛刻.
處事不公道.
又經常罵人.
還派我做一些厭惡的工作.
這麼難受.
我當然是理他也傻了.
其實保羅在另一本書中.
提摩太前書的六章一至這裡這樣說.
他說.
「凡負厄作勞僕的.
應當看自己主人是配受十分敬重.
免得神的名和道理被人褻瀆」.

$^{161}$其實是什麼意思呢?.
我們要留意.
在當時多神的社會中.
基督信仰不單只是一個新興的宗教.
更不算是一個討好人的宗教.
假如信了主的勞僕.
若不再敬重他的家主.
就很容易被人說病.
你看看我的勞僕.
他信了主之後.
沒有像以前般尊重我.
偶爾還會不聽我的話.
甚至會駁我嘴.
你說信耶穌有什麼好呢?.
所以.
史圖保羅勉勵信了主的僕人.
仍然要繼續敬重他的家主.
特別是他家主未信主的時候.
他要為神的福音的緣故.
作美好的見證.
如果反過來.
他的家主信了耶穌.
又怎麼做好呢?.
提摩太前書六章二至這裡.
保羅繼續這樣說.
他說如果.
勞僕有信主的主人.
不可因為他們是弟兄.
就輕看他們.
倒要加以服侍他.
因為這些受到服侍的.
益處的.
是信主蒙愛的人.
保羅在這裡提醒.
信了主的僕人.
不可以因為他的家主.
現在也是主內的弟兄.
與他同一格.
同一個級別.
就輕看他.

$^{201}$反而反過來.
要更加殷勤地.
來服侍他的主人.
理由一點也不複雜.
既然家主現在也信了主.
是跟隨耶穌.
蒙耶穌所愛的人.
那麼僕人又有什麼理由.
不去愛耶穌所愛的呢?.
簡單說.
愛屋及烏四個字.
那麼僕人到底怎樣.
向他的家主展現.
基督的愛呢?.
其實最直接的方法.
就是在他現有的崗位上.
用另一個身份.
用基督僕人的身份.
全心全意地服侍.
聽從他的家主.
為他的家主.
謀求最大的益處.
當然我們要留意的是.
保羅沒有提及.
家主的吩咐.
就等同於基督的吩咐.
僕人聽從服侍耶穌基督.
永遠是優先過他的家主.
他不能夠盲目聽從.
或者盲目.
支持他的家主一切的要求.
在耶穌基督事主.
這個前提下.
僕人需要拒絕替家主做.
任何上帝不喜悅的事情.
第六節.
史勞保羅在這裡這樣說.
他說.
不要只是在眼前事奉.
像是討人喜悅的.

$^{241}$要像基督的僕人.
從心裡遵行神的旨意.
討人喜悅.
這裡其實.
原本是說.
為犧牲原則.
來取悅別人的那些人.
保羅在這裡提醒.
不可以為了要取悅.
你眼前看得到的主人.
而犧牲了聖經信仰的原則.
放棄遵守上帝的吩咐.
因為耶穌基督才是僕人.
真正要效忠.
要服侍要取悅的對象.
而另一方面.
不要只在眼前的事奉.
亦是提醒.
不要只是表面上.
對家主必恭必敬.
當家主在場的時候.
就憤愛的賣力.
來博取他的歡心.
一旦主人走轉杯的時候.
就陽奉陰違.
可能將他手上的東西.
敷衍來做.
甚至乎偷懶,吞撲.
甚至說他主人的壞話.
雖然這一切.
他肉身的主人未必看得到.
不過.
保羅提醒.
他說在天上的主.
卻是知得一清二楚.
事實上,弟兄姊妹.
有時陽奉陰違的態度.
不單單出現在古代.
其實在現代打工仔身上.
都會發生.

$^{281}$相信我們都很明白.
我還記得在.
十多年前.
我還在公司工作的時候.
在職場工作的時候.
我在一個很大的工作環境裡.
你知道有很多張寫字台.
很多人在工作.
我通常位置比較坐後排.
我會發現前排的同事很特別.
通常他們有幾個.
會在他們的電腦螢光幕裡.
安裝一塊鏡.
我開始就好奇.
看看發生什麼事.
為什麼幾個人安裝一塊鏡.
是不是風水呢?.
我看一看.
那塊是廣角鏡.
很有趣.
不算很大.
大約這麼大.
剛剛安裝在電腦螢幕的角落.
後來我就明白.
有什麼作用.
原來方便他.
因為他坐在前排.
很多人會在他身邊從後經過.
為了知道誰站在他後面.
有沒有人在看他做什麼.
他就安裝了這塊倒後鏡.
以致他可以遠遠看到.
有沒有人盯著他工作.
特別是他的上司.
有沒有看到他工作.
丁姐妹,我想.
其實在公司裡.
我們真正要取悅的是哪一個歡心呢?.
雖然眼看不到.
但我們是否相信.

$^{321}$上帝一直在我們身邊.
看著我們做什麼呢?.
特別是.
現在因為疫情的關係.
我們有時可能要在家工作.
其實這也是挑戰.
當你的老闆,上司.
甚至同事都不再在你視線範圍.
沒有人看到你在家做什麼的時候.
你還可以自律地.
緊守你的工作崗位.
討著上帝的喜悅嗎?.
經文提醒我們.
無論你在辦公室裡.
或者在家工作.
我們都要做好.
上帝託付給我們的工作.
正如在第八節裡.
史圖保羅繼續說.
「神曉得各人所行善事.
不論是為奴的,是自主的.
都必按所行的,得主的賞賜」.
到底能否在上帝的面前.
獲得稱讚,獲得賞賜.
保羅說.
與我們的角色身份沒有必然關係.
無論你是做僕人的.
你是做主人的也好.
上帝看重的是.
我們到底是否在用.
祂喜悅的手段和態度.
來演繹好.
我們在職場上的角色.
因為原來.
我們每一個都需要向.
這位天上的老闆.
問責,交仗.
所以你看到.
史圖保羅不單止.
勉勵僕人要善待主人.

$^{361}$來到第九節.
他也同時勉勵做主人的.
你要用相同的原則.
來對待你的僕人.
因為主僕雙方.
其實都是在服侍同一位.
在天上不偏待人的主.
大家都是耶穌基督的路僕.
到底做主人的.
他又要怎樣在管理僕人的角色上.
好好地服侍耶穌呢?.
其實有兩段經文在福音書裡.
耶穌所說的.
都是我們熟悉的.
可以值得給我們一個參考.
第一段是《馬太福音》第十章四十四節.
這裡大家都熟悉了.
「你們中間誰願為首.
就必作眾人的僕人」.
耶穌對門徒這樣說.
另一段經文就是.
《馬太福音》二十五章第四十節.
耶穌說.
「這些是你們既做在我.
這弟兄中一個最小的身上.
就是做在我身上了」.
其實兩段經文的意思.
不是很難理解.
套用在做家主的身上的話.
雖然他的地位確實比他的奴隸高.
但是作為家主.
他卻要效法.
耶穌基督那種僕人領袖的心腸.
放下他的身段.
來去善待在他家裡.
最最最弱勢的一個.
就是他家中的奴僕.
當他這樣做的時候.
耶穌基督說.
其實就是做在他的身上.

$^{401}$所以保羅勉勵他做主人.
不能夠再像以前那樣.
任意用威嚇的手法.
來去核制他的僕人.
雖然在當時的奴隸社會來說.
這樣做是正常合法的.
亦是很有效地.
令奴隸屈服在你的面前.
譬如我們以前也聽過.
如果你不聽話.
我就不會給你飯吃.
甚至你做不好.
我交託給你的工作.
我就要你沒命.
你說拿著頸就命.
做奴隸的怎會不就範.
但是家主現在不可以再這樣.
因為他家裡的奴僕.
不再是屬於家主的工具.
不再是一件物件.
而是屬於耶穌基督的.
奴僕是耶穌所愛的.
所以家主又怎能夠用威嚇的手段.
來惡待他的奴僕.
得罪耶穌基督.
保羅是說這個意思.
事實上當我們看回.
今天一些的職場文化.
雖然我們不會像昔日奴隸社會那樣.
用一些極端的手法.
來威脅別人去做事.
但是我們會不會.
將別人不知不覺地.
當作一件工具.
當作一個死物呢?.
是很值得我們思想的.
當對方有利用價值的時候.
我們會不會對他好一點.
反過來如果沒有利用價值的時候.
我們就開始看不起他.

$^{441}$甚至對他無禮.
來隨意呼喝.
甚至責罵對方.
頂智慧這段經文提醒我們.
要善待你身邊的同事.
你身邊的下屬.
不可以因為在職權上的優勢.
就來欺負甚至威嚇對方.
因為他們都有上帝的形象.
都是耶穌所愛要拯救的對象.
特別當你的下屬.
未能夠達到你的要求的時候.
甚至有時犯錯令你很失望.
要你背鍋.
有時可能頂你兩句.
不聽你的話.
氣得你死呀死呀的時候.
你就很想怎樣?.
有時將他炒之而後快.
但這段經文提醒我們.
會不會有另一個想法.
就是為著上帝榮耀的緣故.
為著福音的緣故.
你選擇以基督耶穌獨人的心腸.
與眼中不太可愛的同事.
多走一里兩里的路.
並且求上帝動他.
我們說不是把他動走.
是動動他的心.
你求上帝改變他.
亦同時求上帝改變你的心.
幫你有更多的愛心.
有更多的忍耐.
來提醒他.
勸他.
給他機會成長.
一起在辦公室裡.
經歷上帝的恩典.
這是經文給我們的考驗和挑戰.
確實不容易的弟兄姊妹.

$^{481}$但保羅說.
耶穌基督喜悅我們這樣做.
事實上.
耶穌也是這樣對我們.
當我們不太可愛的時候.
當我們不聽耶穌基督的話的時候.
甚至是搏擊他的時候.
耶穌也沒有用威嚇手段來對我們.
你想想.
耶穌也給了我們很多機會.
幫我們成長.
提醒我們.
勸我們.
甚至我們忘記了耶穌的時候.
聖靈也會親自來呼喚我們.
弟兄姊妹.
盼望.
這段經文給到你.
無論你是作為上司的.
作為下屬的.
我們也能夠以基督徒的身份.
來服侍善待身邊的人.
讓耶穌基督.
祂的名字.
讓福音.
能夠.
獲得別人的稱許.
而另一方面.
其實這段經文也給了我們一些其他方面的提醒.
首先我們要明白.
在當時的奴隸.
他們不像今天的打工仔.
他們不可以說.
我辭職了.
我不做了.
我轉工了.
東家不打就打西家.
現在這個主人這麼壞.
他們不可能這樣做.
他們也不會像今天這樣.

$^{521}$有工會.
跟他們爭取一個較好的工作環境.
做奴隸的更加沒有權要求他的主人.
你對我好一點.
客氣一點.
不能這樣說.
所以你看到在當時做奴隸的.
似乎沒有什麼辦法可以改善.
他的工作環境.
不過你留意很特別的就是.
上帝並沒有呼召使徒保羅.
你即時推翻當時一個不公平.
不合理的奴隸制度.
上帝沒有這樣給使徒保羅一個這樣的意象.
反而上帝是呼召.
保羅.
你去傳福音.
給做主人和做僕人.
你讓耶穌基督去改變他們的心.
使他們在現有的崗位裡面.
活出一個更高層次的衛生.
就是以耶穌基督奴僕的身份.
這個新的身份.
透過信服耶穌基督的教導.
來改變他們各自在職場上.
那種待人處事的態度和手段.
願意效法耶穌基督那一份.
捨己犧牲的心志.
願意以耶穌基督的愛.
來善待對方.
這是耶穌基督提供的方法和出路.
正如有一位歷史學家.
他叫做威爾·杜蘭特.
Will Durant.
他觀察得很好.
他說在當時.
羅馬帝王凱撒是透過改變一個制度.
來改變人心.
但是耶穌基督卻是透過改變人心.
來改變當時的制度.

$^{561}$結果你會看到.
福音那種大能.
為人類帶來了那種內在的革新.
導致在很多年之後.
產生了我們說.
廢除奴隸主義.
不單只是一個結論.
更加是一個結果.
應用在我們今天.
職場的生活中.
可能在我們工作環境.
在我們身處工作的場所中.
有很多事情都未完善.
人不是很完善.
公司也有很多制度.
可能仍然有待改革.
但是.
這段聖經提醒我們.
當這一切都未達至完備之前.
請不要停止.
在你的職場崗位裡.
實踐福音真理.
讓耶穌繼續去改變各人的心.
包括你自己.
因為一個制度組織.
一個組織.
無論它的制度.
無論它的架構.
即使變得多麼完善也好.
但是如果在這個組織裡的人的心.
沒有被改變的話.
其實始終.
人與人之間很難相處.
關係始終很難改善.
無論辦公室也是這樣.
在你家裡.
在你教會裡.
其實也是這樣.
改變人心.
勝過改變制度.

$^{601}$又或者弟兄姊妹.
可能在你現在工作的環境裡.
你覺得不是很愉快.
你覺得你的上司.
可能你的同事.
有時候也很難忍受.
很難相處.
你很想去轉工.
找一個新的環境.
找到出路.
當然.
確實.
上帝有可能透過你.
轉換工作的環境.
找到一個更理想的工作.
更適合你.
這個固然也是.
上帝賜給你的福分.
但如果上帝不開路.
他需要你留下來繼續等待.
在你找到下一份好的工作之前.
經文提醒我們.
先去調教工作的心態.
先在你現有的崗位裡.
你發掘在工作中.
上帝給了什麼樂趣.
給了什麼好處.
一定有的.
因為上帝不會煽你.
不會作弄你.
當你找到這些好處.
這些樂趣的時候.
其實就能夠提升.
你工作的動機和動力.
讓你更有力量地.
去面對那些你解決不了的逆境.
那些複雜的人士.
並且可以讓你有更多的愛心.
和你身邊的同事.
好好相處.

$^{641}$為上帝做好見證.
等到將來.
我們來到上帝面前的時候.
交帳的時候.
我們可以贏得他的讚賞.
贏得他的獎賞.
事實上,真正的弟兄姊妹.
聖經沒有反對我們轉工.
但不可以缺少的就是.
無論你去到哪裡工作.
你需要按照聖經的教導.
來校正你的工作觀.
你要找到.
你在職場上的真正身份.
你要知道你其實去到哪裡.
你是在服侍耶穌.
這樣你才能找到工作的原動力.
這樣你才能在工作的時候.
你沒有什麼氣氛的時候.
沒有什麼感覺的時候.
你可以知道向誰.
來去支取力量.
千萬不要純粹為了升職.
加薪.
或者說出名.
純粹因為那裡的上司同事對你好.
我們說這些某程度上.
雖然是工作的動力.
但是屬於短暫.
和會變的.
為什麼這樣說?.
很簡單的一些例子就是.
你剛剛這個月加薪.
這份好處可以幫你延續多久?.
你的滿足度可能在月頭的時候滿足.
可能到了一個星期.
兩個星期之後.
已經不覺得是什麼一回事.
因為不持久.
公司也不會每天都加薪.

$^{681}$來保持你的工作熱誠.
又或者反過來.
你找到一個.
那裡同事很好的工作地方.
所有人都很好.
但假如上帝讓那些人.
因為不同的緣故.
離開了公司.
又或者換了一些.
你覺得很難相處的新同事.
甚至乎是上司下來的時候.
你又可能會覺得.
為什麼會這樣呢?.
所以弟兄姊妹.
不要為了純粹一些.
會短暫或者變幻的因素.
來作為你工作的原動力.
而是要以基督是主.
這個真理.
作為你的首要工作動機.
使你無論在哪裡工作.
遇到什麼人也好.
你都可以帶著.
耶穌基督僕人這個身份.
來討上帝的喜悅.
並且在這個前提下.
你努力靠著耶穌.
來與人彼此的配搭,合作.
我還記得在以前的公司.
一做就做了十二年.
即是說從畢業.
到去蒙召讀神學.
這十二年裡.
我服侍過七個老闆.
七個上司.
當中有韓國人.
第一個老闆是韓國人.
之後有馬來西亞人.
之後有英國人.
之後又變回韓國人.

$^{721}$之後又是一個.
美國回流回來的香港人.
到最後那個就是一個地道.
本地香港人.
這七個人有不同的性格.
有不同的要求.
我也試過被外行人領導.
被一個可能在工作範疇裡.
沒有你那麼熟悉.
沒有你那麼資深的人.
不過他是你的上司.
他可以有權去.
吩咐你怎樣做.
我又試過這樣.
試過有很多不合理.
不合理的要求.
我也試過很想轉工.
很想離開.
不過我會問自己一些問題.
難道上帝真的錯配了.
人在公司裡面的職位?.
難道上帝真的差到.
他不知道應該安排誰做我上司.
才是最好?.
我相信.
發生在我身上的事不是偶然.
一定會有上帝的心意.
所以最後我仍然選擇.
留在那間公司裡面.
當然有不同的就是.
我要學習去調教自己的心態.
譬如遇到我覺得.
很難受的人士的時候.
我會先望下上帝.
禱告求他幫助我.
不單止給力量.
也給我懂得去反省.
到底我是否只能夠去.
信服那些我佩服的人.
到底我是否只能夠去.

$^{761}$聽從那些我覺得厲害的人.
又或者.
到底我是否只能夠.
和那些我覺得友善的人.
來和睦共事.
如果對方我覺得.
我不喜歡他.
我不喜歡他的.
我就會刻意不和他合作.
總是用理由來推搪他的要求.
當我多點這樣想的時候.
其實會發覺上帝.
不知不覺在裡面改變我.
變得比以前謙柔了.
沒那麼急躁.
慢慢地你會開始.
和不同的上司建立一些默契.
建立一些互信.
我會發覺其實上帝.
在不同階段會透過不同的人.
來操練我一些愛心.
我的忍耐力.
和給一些工作的智慧我.
然後上帝又很奇妙的就是.
到了上帝覺得.
時候到了的時候.
上帝又會安排.
我那些上司.
一個一個離開那間公司.
當然到最後上帝連我也帶走.
就是我去夢召全職.
來接受神學裝備.
所以很奇妙的,弟兄姊妹.
當你認定你在工作裡面.
你所做的一切是本乎於基督.
是為了耶穌的時候.
你的工作就會比以前.
顯得更加有意義.
你做事上會格外認真.
你不敢得過且過.

$^{801}$因為你知道上帝看著你.
你對人的態度.
也很肯定會有所改變.
即使你要面對一些很難受的同事.
你不敢會再將他當作一件工具.
一個死物.
甚至將他當成一個問題.
來看待.
因為耶穌也不是這樣對待他們.
你會願意代表耶穌.
來學習祝福對方.
而當我們有這種醒覺的時候.
聖經裡面說.
我們的行事為人.
就會越來越像耶穌.
使得我們可以很不經意地.
在人的面前.
流露出耶穌基督的香氣.
剛才說了我工作的一些背景.
剛才說我第一個老闆是韓國人.
你知道韓國人其中一個飲食習慣就是.
他們每一餐.
可能無論早午晚三餐.
都會吃生的蒜頭.
他們有這個習慣.
我那個韓國人的上司當然不例外.
他每一餐都有蒜頭吃.
所以我們很相信.
我們一班同事很相信.
他的蒜頭味.
已經進入了他的生理系統.
已經進入了他的五臟六腑.
進入了他的血液裡面.
為什麼這樣說呢?.
因為他不需要開口.
他只要經過你身邊.
他就會散發出.
那一陣蒜頭的氣味.
可能你也有這樣的經歷.
譬如你身邊家人吃了蒜頭.

$^{841}$你走近他.
他一開口你就發現.
你今天吃了蒜頭.
但我老闆不用.
他不需要開口.
他身上已經有蒜頭味.
很厲害.
同樣地,姐妹們.
我想耶穌基督的香氣.
會不會在我們的肺腑裡面.
在我們每個毛孔裡面.
以致到就算.
你沒有機會開口跟別人說福音.
但別人也能夠透過你.
工作的態度和手法.
去看到福音.
就正如.
你想起一幅圖畫.
在當日保羅的時代.
一個奴隸和他的主人.
竟然可以彼此不單清卿道弟.
還可以大家和睦相處.
你有沒有想過.
本身這是一個很震撼.
很有力的見證.
他們不需要開口說福音.
看到的人已經很驚訝.
為什麼你的主人.
會對你這麼好.
我的主人對我不好.
又刻薄又尖酸.
又或者反過來.
為什麼你的奴隸.
會對你這麼忠心聽話.
我的奴隸又懶.
又吞撲又打虎頭.
你想想當時這個見證的畫面.
是多麼的震撼.
事實上有些數字顯示.
其實當時羅馬帝國裡面.

$^{881}$有很多這些奴僕關係.
只是奴隸可能也有幾千萬人.
你想想如果就算.
有一成的人信了主.
有幾百萬人這樣來到.
在他的宮葬裡面.
和他的僱主.
有一個這麼和睦的見證的時候.
你想想怎會不翻天覆地.
看到的人他們會很想知道.
發生了什麼事.
為什麼你們的主僕關係.
和以前完全不同.
而這些好的關係.
在我們身上是沒有發生.
在這個時候.
作為信了主的主人和僕人.
他們就有機會開口.
來將他們心中盼望的緣由.
和對方分享.
來和對方開口講福音.
所以弟兄姊妹.
這段經文真的提醒我們.
其實職場.
是我們其中一條很重要的福音戰線.
你想想你的人生每一天.
花了多少時間在辦公室裡面.
你就知道.
這是一個對信徒來說.
很重要的福音陣地.
有人曾經這樣說過.
他說.
其實你很可能.
就是對某個人來說.
唯一並且將會讀到的聖經.
意思是什麼呢?.
你身邊的同事.
他們大概不會自己.
突然拿一本聖經來讀.
甚至他們根本沒有聖經.

$^{921}$他們只能夠通過你的好行為.
看到聖經是怎麼說的.
只能夠通過你的好行為.
看到耶穌是有多好.
我還記得以前公司.
有一個這樣的例子.
在同一層樓有不同的部門.
我認識有另一個部門的主管.
坐在我附近.
有時他的下屬會上來.
和他談一些工作的事情.
他會吩咐他的下屬.
怎樣怎樣做.
這個上司也是基督徒.
不過有時每次.
當他吩咐完下屬之後.
上司就離開了.
他的下屬有時會走過來.
和我說:壞人來的.
即是說他的上司.
壞人來的.
我說:為什麼這樣?.
經常都迫他走.
經常都迫害他.
他說:都不知道為什麼.
信耶穌信到這樣.
弟兄姊妹,我想很多時候.
真的很實際.
有時說多無謂.
別人往往從我們的行為.
直接感受到我們信耶穌.
到底是信成怎樣.
確實.
上帝會給我們.
有一些人.
擁有較多.
說話的能力和恩賜.
可以很自然地.
流暢地.
開口和人分享信仰.

$^{961}$但亦有一些人.
無論你多努力練習.
要你對著別人說福音.
總是說得很急.
總是會有些膽怯.
不過.
重要的不是這件事.
重要的是.
無論上帝給你的恩賜是什麼.
你都有能力為耶穌做見證.
因為耶穌吩咐了.
他們也都應許.
你一定有足夠的能力.
為耶穌做見證.
無論你是用口說.
還是透過你的好行為.
是每一個基督徒都能夠做到的事.
如果你還記得.
耶穌基督在福音書裡.
吩咐門徒傳福音的時候.
要他們怎樣?.
靈巧像什麼?.
像蛇.
我們都記得吧?.
有人問到底什麼是蛇的智慧呢?.
怎樣靈巧像蛇呢?.
除了我們知道蛇是怎樣.
經常會鑽孔鑽縫.
很多小技巧.
但有一位靈修學大師.
叫做魏諾德.
叫做Dallas Wheeler.
他說得很好.
你有沒有見過一條蛇.
經常追著獵物來咬人?.
有沒有見過一條蛇.
四處追人.
狗條街似的來咬人?.
很少.
我們都很少見到.

$^{1001}$通常蛇是怎樣等待獵物的?.
靜靜地.
不動.
假裝不在.
等待牠的目標.
不知不覺地走近牠.
等到機會的時候會怎樣?.
一口撲上去.
當然不是咬牠.
如果用全福音來說.
就是怎樣?.
一撲上去.
就用純良像.
甲子的心.
來溫柔地.
和牠分享信仰.
所以弟兄姊妹.
我想.
今日這段經文提醒我們.
無論你在公司裡面.
甚至你在家裡.
或在教會裡面.
我們都要一起成為.
福音的瀑翼.
我們都要藉著.
我們的好行為.
我們的見證.
吸引你身邊的人.
來走近.
停下.
然後求上帝給我們機會.
可以用口.
來和他們分享見證.
分享福音.
為耶穌.
結更多的果子.
很盼望我們今天透過這段經文.
我們可以.
彼此的激勵.
我們一起為上帝發光.

$^{1041}$一起成為.
上帝心目中.
能夠討.
祂喜悅的子民.
就讓我們在這一刻.
我們一起用一首詩歌來回應.
我們夢照為神子民.
這一首歌.
盼望.
這首詩歌能夠提醒.
我們的身份.
我們是屬神的子民.
我們是耶穌基督的奴僕.
我們.
存活在世上.
是為了.
要服侍耶穌.
為人.
作美好的見證.
讓更多的人.
能夠認識耶穌.
成為.
耶穌的子民.
讓我們同唱這首歌.
\newpage



\section{以弗所書 6:10-20}
\label{sec:TUMQvpeyy3M}
\textbf{2020年8月30日 崇拜直播｜以弗所書6\_10-20}
\newline
\newline
連結: \href{https://youtube.com/watch?v=TUMQvpeyy3M}{\texttt{https://youtube.com/watch?v=TUMQvpeyy3M}} ~~~~ 語音日期: 2020-09-01
\newline
\newline
\hyperref[sec:4TNIVRU8_Gc]{\small{< < < PREV SERMON < < <}}
~
\hyperref[sec:index_chronic]{\small{[返順時目]}}
~
\hyperref[sec:Ik8b7Ffs2qU]{\small{> > > NEXT SERMON > > >}}
\newline
\newline
以弗所書 6:10-20
\newline
\begin{longtable}{cl}
\hline
\hline
章節 & 經文 (和合本修訂版)\\
\hline
6:10 & \begin{tabularx}{0.7\textwidth}{X} 最後，你們要靠著主，依賴他的大能大力作剛強的人。 \end{tabularx} \\ \\ \relax
6:11 & \begin{tabularx}{0.7\textwidth}{X} 要穿戴神所賜的全副軍裝，好抵擋魔鬼的詭計。 \end{tabularx} \\ \\ \relax
6:12 & \begin{tabularx}{0.7\textwidth}{X} 因為我們的爭戰並不是對抗有血有肉的人，而是對抗那些執政的、掌權的、管轄這幽暗世界的，以及天空靈界的惡魔。 \end{tabularx} \\ \\ \relax
6:13 & \begin{tabularx}{0.7\textwidth}{X} 所以，要拿起神所賜的全副軍裝，好在邪惡的日子能抵擋仇敵，並且完成了一切後還能站立得住。 \end{tabularx} \\ \\ \relax
6:14 & \begin{tabularx}{0.7\textwidth}{X} 所以，要站穩了，用真理當作帶子束腰，用公義當作護心鏡遮胸， \end{tabularx} \\ \\ \relax
6:15 & \begin{tabularx}{0.7\textwidth}{X} 又用和平的福音當作預備走路的鞋穿在腳上。 \end{tabularx} \\ \\ \relax
6:16 & \begin{tabularx}{0.7\textwidth}{X} 此外，要拿信德當作盾牌，用來撲滅那惡者一切燒著的箭。 \end{tabularx} \\ \\ \relax
6:17 & \begin{tabularx}{0.7\textwidth}{X} 要戴上救恩的頭盔，拿著聖靈的寶劍—就是神的道。 \end{tabularx} \\ \\ \relax
6:18 & \begin{tabularx}{0.7\textwidth}{X} 要靠著聖靈，隨時多方禱告祈求，並要為此警醒不倦，為眾聖徒祈求。 \end{tabularx} \\ \\ \relax
6:19 & \begin{tabularx}{0.7\textwidth}{X} 也要為我祈求，讓我有口才，能放膽開口講明福音的奧祕， \end{tabularx} \\ \\ \relax
6:20 & \begin{tabularx}{0.7\textwidth}{X} 我為這福音的奧祕作了帶鐵鏈的使者，讓我能照著當盡的本分放膽宣講。 \end{tabularx} \\ \\
[1ex]
\hline
\hline
\end{longtable}
$^{1}$之後是經訓.
今天的經文是在二弗所書.
六章十至二十節.
新約聖經二百七十頁.
(見課).
我還有沒了的話.
你們要靠著主.
倚靠他的大能大力.
作剛強的人.
要穿戴神所設的全副軍裝.
就能抵擋魔鬼的詭計.
因我們並不是遇屬血氣的真賤.
乃是遇那些執政的.
掌權的.
管轄這幽暗世界的.
以及天空屬靈氣的惡魔真賤.
所以要拿起神所設的全副軍裝.
好在磨爛的日子抵擋受敵.
並且成就了一切.
還能站立得住.
所以要站穩了.
用真理當作帶子束腰.
用公義當作護心的借控.
又用平安的福音.
當作預備走路的鞋穿在腳上.
此外.
又拿著信德當作藤牌.
可以滅盡那惡者一切的火箭.
並戴上救恩的頭盔.
拿著聖靈寶劍.
就是神的道.
靠著聖靈.
隨時當方禱告祈求.
並要在此警醒不倦.
為眾聖徒祈求.
也為我祈求.
使我得著口才.
能以放膽開口.
講明神福音的奧秘.
我為這福音的奧秘.

$^{41}$作了帶所憐的使者.
並使我照著當盡的本分.
放膽講論.
今日的講題是.
「發放群體的光芒.
作剛強的人」.
弟兄姊妹平安.
香港的疫情.
越來越漸趨穩定.
期待我們能夠.
恢復實體的崇拜.
我們能夠在實體裡面.
一起敬拜神.
一起互相分享.
繼續去紀念這個世界裡面.
受著疫情影響的不同地方.
求神止息這個病毒.
帶來的人命的傷亡.
面對著這個病毒.
我們都會嚴陣以待.
每天出門口.
我們都會確保自己.
是否已經戴上口罩.
口袋裡面.
衣袋裡面是否有戴上潔手液.
更加有些人嚴謹一點.
會戴上護目鏡.
戴上手套.
可能有些人攻降.
乘飛機會全身穿著保護衣.
保持社交距離.
出門後回到家裡.
我們立刻第一件事情.
徹底清潔我們的身體.
洗澡,換衣服.
這個新型冠狀病毒.
是現在我們每一個人類.
全球的人頭號的敵人.
我們每一天都很認真地看待.
雖然很多時候.

$^{81}$我們現在慢慢漸漸出現抗疫的疲勞.
我們很想呼吸新鮮的空氣.
我們很想出街.
我們很想見人.
但是我們知道.
如果我們有任何的不小心.
有任何的鬆懈.
我們就很容易確診.
這個確診還會感染我們的家人,朋友.
所以我們要繼續堅忍著.
守住這一關.
有專家說這場的真戰.
可能是一個長久的戰.
不會這麼快結束.
可能會維持一兩年的時間.
有些專家更加去說到.
這個病毒可能會成為雙峰感冒.
繼續存在這個世界裡面.
這場真戰的敵人是這個病毒.
是看不見的敵人.
但是它真實存在在這個世界裡面.
並且為人類帶來傷害和威脅.
我們生活在這個世界裡面.
不單止面對這個病毒的真戰.
同樣我們基督徒也都面對著不同的真戰.
在我們日常生活裡面.
遇上大大小小的困難,挑戰.
生老病死.
日常不同的大大小小的困難.
我們都努力去面對著.
在屬靈真戰裡面有不同的戰爭的場地.
第一個場地是我們內心和情慾裡面的交戰.
在羅馬書裡面說到.
因為情慾與聖靈相爭.
而聖靈又與情慾相爭.
這兩者是彼此相敵.
第二場真戰是聖俗的相戰.
我們和這個世界的神.
即魔鬼撒旦.
和這個世界裡面的價值觀.

$^{121}$不斷去爭戰.
我們基督徒持守著真理.
基督徒持守著上帝給我們的價值觀.
第三場的真戰是群體的真戰.
無論在教會,在家庭,在朋友.
每一個地方有人的地方.
不同的人際的關係會有磨擦.
會有誤會.
會被魔鬼撒旦留有地步.
成為真戰的場所.
而第四個真戰的場所.
就是一個屬靈的真戰.
我們的敵人只有一個.
就是魔鬼撒旦.
每一場的真戰都是持久.
從我們信主.
甚至會從我們出生開始.
到我們再見主的面.
每一場真戰都是要持久.
繼續去戰爭.
當我們面對著這個看不見敵人.
這個敵人不斷用不同的方法.
誘惑我們離開上帝.
令我們活在罪的裡面.
今天這段經文裡面.
上帝教導我們.
如何去打這場屬靈的真戰.
如何能夠在屬靈的真戰裡面.
站立得穩.
經文一開始.
他告訴我們要做剛強的人.
渴望我們每一個能夠藉著這段經文.
我們每一個一起成為一個剛強的人.
從我們的心.
我們的靈.
不斷被上帝去挑釁.
這是《已忽所書》裡面最後一個部分.
補充一下總結了上帝救恩的計劃.
和信徒的生活.
有很多的字眼.

$^{161}$屬靈的軍裝.
其實都是前面提及過.
現在再去重複.
不過藉著這段經文.
教導我們如何在現今這個世代裡面.
站立得穩.
抵擋魔鬼撒旦的攻擊.
經文分開三個部分.
每一個部分環環緊扣.
每一個部分都是很重要.
第一個部分.
抵擋魔鬼撒旦.
第十到第十三節.
經文裡面我再去讀出.
第十節.
「我還有沒有了的話.
你們要靠著主.
以教他的大能大力.
作剛強的人.
要穿戴上帝所賜的全副軍裝.
就能抵擋魔鬼的詭計.
因我們並不是與屬血氣的真善.
乃是與那些執政的掌權的.
古特這幽暗世界的.
以及天空屬靈氣的惡魔真善.
所以要拿起上帝所賜的全副軍裝.
好在這磨爛的日子抵擋受敵.
並且承受了一切.
還能站立得住」.
在第一個部分裡面.
保羅他去講到魔鬼這個身份.
他提到我們要認清.
我們的敵人是魔鬼撒旦.
我們要抵擋他.
魔鬼撒旦.
他經文裡面.
第十二節裡面他講到.
「因我們並不是與屬血氣的真善.
乃是與那些執政的掌權的.
古特這幽暗世界的.

$^{201}$以及天空屬靈氣的惡魔真善」.
經文裡面用了四個描述.
多角度去形容這個敵人.
不是人間有血有肉的人.
是有位格的.
是有能力的邪靈權勢.
是超自然.
他也是狡猾.
是聰明.
有智慧.
有謀略.
這幾個稱呼.
不是說撒旦.
有獨立於上帝以外的權力.
與上帝能夠勢均力敵地去抗衡.
更加不是說撒旦真的管轄了天空.
更加不是說他管轄了地上這個世界.
管轄了地下幽暗的世界.
這些不同的領域.
這四個形容詞.
僅僅是強調撒旦與光明對抗.
對抗黑暗的勢力.
其實魔鬼在聖經裡也有不同的稱號.
龍,古蛇,撒旦.
《創世記》裡形容他為狡猾的蛇.
《但與理書》裡形容他為魔君.
《預案福音》稱他為說方者的父.
《多倫多後書》稱他為世界的神.
假裝光明的天使.
《以法所書》稱他為空中掌權者的首領.
《貼上來家後書》稱他為那大罪人.
就是沉淪之子.
《彼得前書》形容他為咆哮的獅子.
等待機會.
《啟示錄》更加形容他為大紅龍.
魔鬼撒旦有不同的稱呼.
有不同的身份.
有不同的工作.
他的目的是什麼呢?.
《預案福音》十章第十節.

$^{241}$我們都很熟悉.
耶穌說.
「刀刃內無非是要偷竊殺害毀壞.
我來了要叫人得生命.
並且得得更豐盛」.
魔鬼撒旦的目的.
就是要來偷取我們的信心.
愛心.
信任.
盼望.
用謊言毀壞我們的自我形象.
扭曲人的價值觀.
顛倒是非黑白.
破壞倫理.
挑撥人與人之間的關係.
以威嚇.
壓迫.
武力殘害踐踏人的生命.
令人走上滅亡毀滅的道路.
魔鬼撒旦用的方法.
他的方式.
他的能力.
經文中提到.
他是一個用詭計.
魔鬼撒旦是控訴者.
他是說謊者.
他會詐騙.
騙人.
他會殺人.
用欺騙.
用詭詐.
來達到他的目的.
魔鬼無處不在.
卻是無孔不入.
在這個世界裡面.
無論是社會.
政治.
法治.
經濟.
任何一個體系.

$^{281}$都被魔鬼利用.
被他的權勢影響.
為了達到他的一個目的.
就是要去破壞人.
根據魔鬼的工作.
除了利用金錢.
一些好處.
情慾.
權利.
去吸引人去犯罪之外.
更加在我們每一個的心裡面.
《真言六章》十六到第十九節.
裡面是這樣說.
「而話所恨惡的有七樣.
連他心所憎惡的共有七樣.
就是高傲的眼.
殺荒的舌.
留無辜人血的手.
圖謀惡計的心.
飛跑恨惡的腳.
吐謊言的假見證.
並抵輕中報殺紛爭的人」.
實在他的手段.
是要去散播絕望.
散播仇恨.
令我們對將來失去了盼望.
他渴望用仇恨.
用報仇蒙蔽我們的內心.
令我們失去愛心.
他用他的謊話充滿著這個世界.
令我們失去信心.
魔鬼撒旦更加會假裝成天使.
魔鬼用的方法騙人的道理.
不一定是完全的錯誤.
有時候是似是而非.
半是半非.
可能表面上以高舉愛.
公義,良心,和平,穩定,自由.
為一個包裝.
似乎很吸引人.

$^{321}$很值得人羨慕.
亦很合情合理.
但這個包裝的裡面.
卻是叫人離開上帝.
傷害人的生命.
帶來群體裡面的破壞.
這是魔鬼撒旦的誘惑.
偽裝的陷阱.
記得2012年的時候.
當時東方閃電在香港教會裡面.
滲入非常嚴重.
他們宣講雷基督.
再來到這個世界裡面.
他們假裝著是舞蹈者.
去親近福音派的弟兄姊妹.
以愛心去關心人.
慢慢去取得人的信任.
帶領弟兄姊妹去查經.
去關心他們.
其中一個深刻的新聞就是.
有一個牧師相信東方閃電的教義.
離開了這個信仰.
對教徒.
對很多未信的人.
產生了很多的影響.
魔鬼撒旦同樣會以患難攻擊信徒.
用苦難令人的心很憂愁.
動搖人的信心.
不再相信祈禱.
不再願意祈禱.
不再去倚靠上帝.
甚至放棄這個信仰.
離開上帝.
魔鬼更加進入人與人之間的關係裡面.
在人與人之間的摩擦裡面.
產生誤會.
產生仇恨.
魔鬼的工作更加令我們的生活陷入盲目的裡面.
特別是當我們有了電話之後.
我們很大部分的生活都被電話所佔據.

$^{361}$心靈裡面的空間缺乏了很多.
對上帝失去興趣.
完全沒有空間.
沒有任何的精神.
去幻想上帝的話.
莫說是要履行上帝的誡命和祂的使命.
網絡世界裡面大量的虛假信息.
而且網絡世界裡面高舉一些錯誤的信息.
高舉金錢,情慾,名利.
網絡世界裡面也不需要為自己說的話付上代價.
大量的網絡欺凌,批鬥,舉報,公審的事情.
當我們打開新聞.
面對著這個荒誕的世界.
充滿謊言,顛倒,黑白是非,指鹿為馬.
當一個人無恥.
他什麼都能夠做得出.
這些信息不斷不斷衝擊我們的良知.
他試圖以威嚇利誘,恐嚇,傷害我們每一個.
這是魔鬼的工作,他的能力.
他有能力,但他卻不是全能.
屬靈的勢力龐大,但他卻不是為所欲為.
在啟示錄裡說到魔鬼撒旦的結局.
啟示錄二十章第十節.
那維說他們的魔鬼被困在流亡的火湖裡.
就是獸和假仙子所在的地方.
他們必晝夜痛苦直到永遠.
這個屬靈戰爭的結果.
從耶穌基督死裡復活.
一早已經獲取勝利.
基督已經勝過黑暗.
勝過一切死亡的權柄.
我們靠著耶穌基督的時候.
我們不需要懼怕.
不過這個魔鬼的勢力仍然殘留著.
仍然在這個世界裡活躍著.
他仍然有一定程度的影響力.
在天空的領域.
在人這個世界的領域.
經文裡形容磨難的日子.
其實是指黑暗,邪惡的日子.

$^{401}$就是我們今天面對的每一天.
基督已經勝過了這個權柄.
我們基督徒當我們願意依靠耶穌基督的時候.
我們不需要去懼怕他.
不過這場勝利還未完全實現.
我們基督徒仍然生活在這個衝突裡面.
需要上帝的裝備來抵擋它.
這是一場宿命的真戰.
是光明和黑暗的對立.
我們每一個都是被上帝所揀選的群體.
站立在光明的一方.
勇敢和黑暗的勢力對抗.
黑暗不斷衝擊我們.
滲透我們.
蠶蝕我們的生命.
甚至黑暗企圖假裝天使.
可惜黑暗最終的結局.
他必定會露出他的真面目.
在日光之下.
上帝會將隱藏的事情去顯明.
黑暗終必不能勝過光.
魔鬼撒旦和他所有有能的大能.
在上帝的裡面都會變得無能.
當我們願意穿戴在上帝所設的全副軍裝.
就能夠成為抵擋魔鬼撒旦的剛強的人.
能夠站立在這.
跟這個勢力能夠抵擋.
當我們知道這個屬靈的世界.
魔鬼的真正存在的時候.
我們就不會過分針對人間任何一個人.
也不會視每一個人為終極的對手.
更加不會視他們為我們的敵人.
將他們妖魔化.
我們真正的敵人.
是背後邪惡的勢力.
不是任何的人.
或許有些人會不斷攻擊我們.
傷害我們.
不過我想說的是.
他們不是魔鬼.

$^{441}$我們真正的敵人.
是魔鬼撒旦.
是這個屬靈氣的惡魔.
不是人間裡面每一個人.
不是我們的家人.
不是我們的朋友.
不是社會裡面任何一個人.
更加不是任何一些組織.
如果教會只是著重於比較次要的目的.
可能會錯過了主要的敵人.
我們可以明白和體諒世人的無知.
他們不知不覺地成為撒旦的工具.
但是我們要做的是.
認清魔鬼撒旦的詭計.
不要讓仇恨,絕望,黑暗.
支配我們的心.
我們用肉體的方法.
不能勝過撒旦.
因此我們需要倚靠上帝.
第一個部分.
上帝已經給了一個全副軍裝給我們每一個.
當我們每一個願意穿上的時候.
上帝就與我們同在.
經文裡面說到.
所以要站立.
所以要站穩了.
用真理當作戴紙束腰.
用公義當作護心鏡遮空.
要用平安的福音當作預備走路的鞋.
穿在腳上.
此外又拿著信德當作盾牌.
可以滅盡那惡者一切的火箭.
並戴上救恩的頭盔.
拿著聖靈的寶劍.
就是神的道.
我們要剛強.
需要穿戴在神所賜的全副軍裝.
剛強不是因為我們能夠堅強.
而是因為由上帝從上而來的保護.
上帝是一切的源頭.

$^{481}$一切的祝福.
一切的保障.
一切的恩典.
都是源自我們的上帝.
我們的主耶穌基督.
當信徒願意與耶穌基督的聯合.
我們能夠靠著祂剛強起來.
耶穌基督會供應我們在屬靈裡面所需要用的一切.
當我們依靠祂大能大力的時候.
上帝能夠叫耶穌基督死裡面復活的大能.
超過這一切執政的掌權的惡魔.
一切的勢力.
我們能夠在上帝裡面.
支取力量.
抵擋魔鬼撒旦的攻擊和詭計.
保羅在這裡用以賽亞書裡面的一些預言.
來描寫耶和華和彌賽亞的軍裝的一些景象.
描寫下來以法所書的全副軍裝.
特別是以賽亞書十一章.
五十九章.
四十九章.
五十二章.
裡面不同的片段.
都是在說彌賽亞穿戴起一個全副軍裝.
以賽亞書在描述這個萬軍之主.
是一個戰士.
他為他的子民伸冤.
上帝的子民需要穿戴上帝所設的全副軍裝.
和這位彌賽亞一起並肩作戰.
上帝已經得勝.
他並且勝過死亡.
他將這個權柄.
他將這個使命交託給我們每一個子民.
當我們願意穿戴上帝所設的全副軍裝的時候.
我們就能夠與屬靈屬世的勢力一起去對抗.
信道正在進行一個生死攸關的屬靈征戰.
如果不小心就會中了魔鬼的殺蛋的詭計.
因此我們需要穿起全副的軍裝.
這個全副軍裝在防守.
在進攻的裡面.

$^{521}$能夠保護我們身體裡面的每一個角落.
信道所穿戴的全副軍裝.
有真理,有公義,有救恩.
表明我們穿上上帝自己和他的屬性.
穿戴全副的軍裝目的是讓我們能夠站立得穩.
抵擋詭計.
每一個上帝的屬性都貫穿在已忽所書.
並不是獨立地存在.
第一真理的大子.
真理的大子在已忽所書裡面.
第一章四章的裡面.
不同的節數的裡面都記載著.
對於羅馬士兵來說.
這些的大子可能是懸掛在軍裝下面.
保護大腿的一個皮革的圍裙.
而不是佩戴著刀劍的一個大子.
或者保護腰的一個腰帶.
腰部穩妥能夠站立得穩.
束上腰的時候.
其實是一個動作.
是代表他們已經準備好真戰.
隨時候命進入戰場.
真理是上帝的屬性.
是教義.
是福音裡面的.
在聖經裡面的真理.
教徒的生活原則.
讓福音八章三十二節.
「你們必曉得真理.
真理必叫你們得以自由」.
今天魔鬼撒旦的攻擊.
他就是一個沒有真理的魔鬼.
他高舉的自由.
其實是叫人放縱.
真理是一個有規範.
叫我們生活裡面有原則.
魔鬼撒旦會利用一些謊言.
詭計.
一些好處.
欺騙.

$^{561}$假裝成一個陷阱.
叫我們陷入裡面.
當我們有真理的腰帶的時候.
我們能夠分享上帝的聖情.
用真誠的品格.
成為一個原則.
抵擋魔鬼撒旦一切的謊言.
第二.
是公義的護心鏡.
這是保守我們每一個清義的心.
對於羅馬士兵來說.
護心鏡是一個遮空的裝備.
有前也有後.
避免他受到毆打或戰的傷害.
在以上書59章第17節.
一般學者都會認為.
這個公義是指上帝稱人為義.
使人在上帝面前有一個合宜的立場.
與上帝保持一個良好的關係.
不過同樣.
公義其實都在倫理道德的事情裡面.
公義的至根與理所當然與合宜其實是一樣的.
只要我們要做一些理所當然的事情.
魔鬼用的方法就是扭曲人的價值觀.
顛倒黑白是非.
使人的心被迷惑被污染.
讓人失去良知.
教徒拒絕在這個世界裡面.
被魔鬼在世界裡面去同化.
竭力讓不公不義的事情.
圍著這些事情去發聲.
追求公義追求真理追求真相.
對良知去守護著.
魔鬼撒旦也會進行控訴質疑.
否定我們的工作.
令我們陷入否定.
令我們陷入自憐自卑的當中.
公義的護心鏡告訴我們.
耶穌基督已經勝過這一切.
祂赦免了我們一切的罪.

$^{601}$要保守我們的心勝過保守一切.
第三件事是平安福音的鞋.
分享平安福音的好訊息.
羅馬士兵的鞋.
他們是叛徒的鞋.
不是一個武器.
是一個行軍打仗.
做一個長遠的一個真戰裡面.
所穿上的.
鞋底裡面會有釘.
釘在戰靴裡面.
目的是讓軍兵.
就是士兵穿上去的時候.
遇到攻擊能夠站立得穩.
不會輕易地被打倒.
走路的鞋其實在原文裡面沒有出現.
新譯本裡面是翻譯.
把和平的福音預備好.
穿在腳上.
其實是在說信徒要積極地去傳揚這個福音.
這個其實是一個畫面.
一個報信者.
他將他在街尾的腳中走遍了這個世界裡面不同的角落.
好像武陸一樣飛快地去到不同的地方去報好訊息.
傳到耶路撒冷.
當這個好訊息被耶路撒冷裡面的居民聽到之後.
耶路撒冷的居民會回應.
平安好訊息.
救恩你的上帝.
我們穿戴起這個鞋.
是為了去傳揚更多好訊息到不同的角落裡面.
第四樣信德當作盾牌.
羅馬士兵其實戴著一個很大的盾牌.
這個盾牌大約是四呎乘兩呎半.
能夠將整個人包圍在這個盾牌裡面.
由兩塊木板和麻布裡面包住.
然後再加上皮和鐵片.
在作戰的時候.
將這個盾牌完全浸在水裡面.
能夠抵擋敵人的火箭和拿著箭頭.

$^{641}$或者招誘所點成的火.
火箭可能使得飛訊徒一些炎熱的攻擊.
帶來強大的殺傷力.
信德就是上帝信實的一個屬性.
無論魔鬼的控訴有多厲害.
環境有多惡劣.
屬靈攻擊有多猛烈.
上帝仍然信實不變.
祂是我們最大的保障.
祂每一個的應許從來沒有任何落空.
第五樣是救恩的頭盔.
永生的確具.
羅馬士兵所戴的頭盔是銅製的.
覆蓋包圍著整個頭部.
上帝已經從死裡復活.
拯救了他的子民.
將每一個相信他的人.
從罪惡裡帶到上帝的身體國度裡.
和上帝一同復活.
勝過了惡者.
上帝的子民有著這個救恩.
還有永生的確具.
耶穌說.
「那殺身體不能殺靈魂的.
不要怕他們.
唯有把身體和靈魂都滅在地獄裡的.
正要怕他」.
當我們穿戴救恩的頭盔的時候.
我們無需去懼怕死亡.
分享一個事蹟.
在1956年的時候.
以艾納奧為首的五個年輕宣教士.
他們去到南美洲厄瓜多爾宣教.
他們去到奧加族.
其實這群人是一群術人族.
他們向他們傳講福音的好信息.
但最後卻被奧加族族人殺死.
他們死在這個地方裡.
他們被殺的時候.
雖然他們佩戴著槍.

$^{681}$但他們卻沒有向他們開槍.
因為他們知道奧加族族人.
還沒有準備好上天堂.
而他們每一個已經準備好.
艾納奧死前用奧加族語.
跟這個族人說.
「你是我的朋友」.
「真實的朋友」.
這五個宣教士.
他們也有他們的老婆.
也有他們的妻兒.
幾歲大的小朋友.
後來他們的妻子.
帶著他們的小朋友.
踏上了仙山殉道的旅程.
告訴當地的奧加族人.
他們上帝仍然愛他們.
之後整個奧加族人.
受洗歸入主的名下.
殺死這群宣教士年輕的.
一個奧加族人.
最後成為了牧師.
他說道.
「全道就是拯救絕對不能失去的生命」.
這五個生命激勵起過千的大專生.
參與在宣教裡面的工作.
最後一項宿命的軍裝.
就是神的刀.
上帝的寶劍.
聖靈和上帝的刀成為武器.
這個武器其實是一把短劍.
大約六十厘米長.
是近身埋伏的時候.
一個重要的防具.
聖靈能夠使這個寶劍有效.
因為它能夠讓這個寶劍鋒利.
真理的聖靈引導我們進入真理.
神的刀就是祂自己的說話.
本有能力.
在聖經裡面.

$^{721}$在希伯來書裡面說到.
「上帝的刀是活潑的,是有效的,比一切兩刃的劍更快.
甚至將靈與魂,節與骨髓都能刺入,剖開.
連心中的思念和主義都能辨明」.
當我們擁有聖經.
就能夠分辨正邪.
透過上帝的真理.
透過聖靈.
我們明白上帝的說話.
不過我們今天怎麼去看聖經呢?.
我們有沒有去看重呢?.
當我們沒有看重聖經的時候.
我們很可能只會跟著大衛的聲音.
只會跟著眾人以為美的事情.
我們跟著去做隨波逐流.
「全副軍裝是穿在身上.
讓我們身體從上到下.
每一個角落,每一個部分.
從外到內都被上帝的恩典,上帝的屬性所包圍住」.
保羅用了四次去重複.
第十一節,十三節,十四節裡面說到.
要站立得穩.
鼓勵信徒要堅守這個原則.
帶著全副軍裝.
其實最重要的.
問我們每一個.
其實我們生命裡面.
從頭到腳.
從外到內.
有沒有不自覺在參與魔鬼撒旦的勢力?.
有沒有不自覺在參與它的活動軌跡裡面呢?.
詩篇的第一篇裡面說到.
「不從惡人的計謀.
不站惡罪人的道路.
不坐洩萬人的座位」.
你有沒有跟從惡人的計謀?.
你有沒有站在罪人的道路?.
你有沒有坐洩萬人的座位?.
求神光照我們.
讓我們從上到下.

$^{761}$從外到內.
每一個部分.
我們都能夠通透.
被上帝的恩典去光照.
努力抵擋魔鬼撒旦.
跟邪惡的勢力劃清界線.
最後一個部分.
也是最重要的部分.
以禱告作真正的武器.
經文裡面繼續說到.
第十八節.
「靠著聖靈隨時多方禱告祈求.
並要在此警醒不倦.
為眾聖徒祈求.
也為我祈求使我得著口才.
能以放膽開口.
說明福音的奧秘.
我為這福音的奧秘.
作了底所念的使者.
並使我照著當進的本分.
放膽講論」.
面對著屬靈的真賤.
我們一定要用屬靈的方法.
用屬靈的兵器.
魔鬼撒旦最怕最怕.
我們去找上帝.
最怕我們去呼求上帝.
經文裡面形容到.
這個真正的兵器禱告.
有五方面.
第一方面.
「靠著聖靈」.
聖靈是啟示我們每一個.
引導我們每一個.
感動我們每一個去祈禱.
《羅馬書八章二十六節》裡面說到.
「況且我們的軟弱有聖靈的幫助.
我們原不曉得怎樣禱告.
只是聖靈親自用說不出來的彈式.
為我們禱告」.

$^{801}$接著經文裡面.
用到四個全字去形容祈禱.
「隨時」是說我們全部的時間.
第二.
第三方面.
「多方」是說我們用盡所有的方法方式.
不同的角度去祈禱.
第四方面.
「祈禱」是要警醒不倦.
是說我們全部的精神.
集中我們的力量.
最後一樣.
「為眾聖徒祈禱」.
是說我們要為的對象.
是全部的人.
想跟大家分享一個數字.
就是「七二四」.
「七二四」其實是一個祈禱的網絡.
「七」其實是指七天.
「二十四」是說二十四小時.
這個祈禱網絡的發起人.
是非洲烏干達的一個牧師.
在1962年烏干達獨立之後.
到1986年.
其實烏干達正在經歷十幾個總統.
烏干達裡面戰火連連.
一個比一個的總統更加殘暴.
尤其是阿敏這個總統執政的時候.
更加將回教成為國教.
殺害了無辜數以百萬計的人.
並且他走出這些邪術.
在當時裡面去橫行無忌.
大量的逼迫基督徒.
當阿敏下台之後.
烏干達接著面對著愛滋病的侵擾.
每三個的市民.
就有一個患上愛滋病.
這是一個不治之症.
教會要封閉.
信徒要躲起來去聚會.

$^{841}$這個國家束手無策.
聯合國衛生組織.
甚至一度去公佈.
這個全球愛滋病裡面.
影響最嚴重的國家.
可以怎樣做?.
面對著腐敗.
面對著黑暗的政權.
面對著病患.
人可以怎樣做呢?.
當時烏干達的牧師.
毛約翰牧師.
他見到一個異象.
他的異象其實是.
上帝給他一個異象就是.
就是要見到一個網.
這個網就要去聯繫在世界裡面.
就是國家裡面不同的地方.
這個網是在說.
上帝要將國家從困境裡面拯救出來.
於是乎他去建立了724的祈禱網絡.
在不同的地方裡面建立祈禱塔.
目的是叫人更多的祈禱.
祈禱能夠成為真善的兵器.
當烏干達整個的人民.
都願意去祈禱的時候.
這個國家有翻天覆地的改變.
祈禱不單只是.
我們將我們沒辦法面對的事情.
去交託給神.
祈禱更加是我們和神一起工作.
我們去面對著無力的事情.
和上帝一起去改變.
當人願意去依靠上帝.
願意去祈禱的時候.
我們就能夠一起和上帝戰勝魔鬼的詭計.
烏干達過了二十多年.
三十多年之後.
整個政權都會改變.
今天這個烏干達.

$^{881}$超過了九成的人是基督徒.
連不治之症,愛滋病.
都從34\%一直下跌.
到3-4\%.
祈禱帶來改變.
祈禱帶來復興.
弟兄姊妹.
你今天祈禱的生活是怎樣的?.
你相信祈禱嗎?.
今天這個世界面對著各種患難.
各種疾病.
人不同的問題.
你相信祈禱嗎?.
你花了多少時間在祈禱裡面?.
本來最後要以不所識的教會.
為他去祈禱.
當保羅正在下在監獄的時候.
他鼓勵這群信徒祈禱的目的.
並不是為了他能夠得以獲釋.
脫離其他的逼迫.
他從頭到尾關心的目的.
是上帝的國.
希望他能夠在被囚的環境之下.
仍然能夠放膽開口.
全講上帝福音的奧秘.
保羅以上帝福音的事情.
為他生命裡面最優先.
無論他是處於富貴.
處於貧賤.
處於患難.
處於信境.
處於任何一個光景裡面.
保羅都以全福音為他最大的照明.
當他叫以不所教會為他祈禱的時候.
他渴望他講出這個福音的奧秘的時候.
能夠正確.
能夠講這個聖靈裡面的大能.
能夠讓人的生命自由.
得著釋放.
希望能夠清楚.

$^{921}$更加在一些君王的面前.
大膽宣認耶穌基督是全人類裡面的救主.
和大家分享這個國家裡面的狀況.
伊朗.
伊朗是全球迫迫基督徒十大最嚴峻的國家之一.
基督徒每天都在經歷各種生命威脅的迫迫.
但同時伊朗其實是全世界教會.
基督徒增長得最快的國家.
很奇怪.
越多的迫迫.
越多的患難.
竟然越多人信耶穌.
國際邊境聯盟的發起人.
就是湯瑪斯他講到.
目前的伊朗的教會.
在這個世界裡面成長得最快.
但是他沒有任何的建築物.
也不是有人的組織.
他也沒有持有任何的物業.
沒有任何的財產.
也沒有心中的領導.
更加沒有任何的宗派分別.
但是人的心卻是十分之火熱.
神在這裡的工作實在令人很震撼.
不過福音卻是大離迫.
湯瑪斯為了伊朗這些基督徒的轉向.
更多的祈禱.
不過他的祈禱.
和一般人所祈的祈禱是不同的.
基本上我們會為了這些人.
脫離迫迫去祈禱.
但是湯瑪斯說到.
不要啊.
你們不要這樣做.
因為迫迫會令到教會增長.
當迫迫停止的時候.
教會的增長就停止.
我們所希望的是福音能夠擴張.
能夠讓伊朗每一個地方的人.
信耶穌.

$^{961}$保羅鼓勵信徒要抵擋魔鬼撒旦.
不單只是他們自己的生命能夠得到保障.
得到保障過後.
剛強起來.
是要幫助其他人.
幫助其他人脫離魔鬼撒旦的核制.
將人從魔鬼撒旦的型裡拯救回來.
轉向基督.
過不多後書第十章四至五節裡面說到.
因為我們雖然在血氣中或行事.
卻不憑著血氣增善.
我們增善的兵器本不是屬血氣的.
乃是在神面前有能力.
可以攻破這個堅固的型裡.
將各樣的計謀.
各樣難阻人認識神那至高之事.
一概攻破.
又將人所有的心意奪回.
使他都信服基督.
這是一個肅靈的增善.
這個增善目的是要拯救人的生命.
透過傳福音.
透過謝伍祈禱.
是一個拯救人靈魂的增善.
教會的設立是要為著傳揚真理.
教會的設立是要抵擋魔鬼撒旦的謊言.
教會的設立是要為著拯救靈魂.
今天我們落在任何光景裡.
你如何去看上帝給你的召命.
上帝給你的使命.
重複今天所說的經文.
三個重點.
第一.
弟兄姊妹.
我們的敵人只有魔鬼撒旦.
今天我們很忙碌.
我們的心思意念.
所忙著要去對付的敵人.
是誰?.
是你的家人?.

$^{1001}$是你的朋友?.
還是地上任何一個人?.
你有沒有試過將他們變成魔鬼.
將他們妖魔化.
更加陷入自相殘殺裡面.
虛耗了我們很多的心力.
第二.
你有沒有逼經意地.
中了魔鬼撒旦的詭計.
令你對上帝失去了盼望.
對人失去了愛心.
對這個世界失去了任何信任.
第三.
亦是最後一點.
我們如何去看待祈禱.
我們有沒有真正用祈禱.
關心患難,痛苦裡面的人.
當我們今天面對著病毒的時候.
我們的心所掛念的是自己.
還是掛念福音.
掛念人靈魂裡面的需要.
求神幫助我們.
求神鞭策我們.
如果我們有任何做得不好的地方.
求神叫我們為罪,為義,為審判.
自己責備自己.
讓我們一起祈禱.
親愛的上帝.
我們求你憐憫.
雖然今天我們是基督徒.
雖然今天我們會去教會.
會看聖經.
我們會口對邊.
口裡說很多道理.
頭腦裡充滿著各人的知識.
主要原來大多我們行事為人.
當我們在上帝你的光照的時候.
主要我們所有心思,意念都很邪惡.
我們被這個世界裡面的價值觀.
被情慾,被患難所勝過.

$^{1041}$主要我們祈求你的憐憫.
求你幫助我們.
拯救我們.
用你的恩典去覆蓋我們.
保護我們.
我們不想有任何一刻.
跟魔鬼撒旦有任何關係.
我們不想有任何一刻.
口裡說我們是光明之子.
但我們卻是在黑暗的裡面.
所以我們相信你很心痛.
所以我們很相信.
你想我們從這些罪惡裡面拯救出來.
所以我們祈求你.
幫助我們今天.
就去站立.
就去堅強.
就去勇敢在這個世界裡面繼續去抵抗.
因為這個本是你拯救我們的目的.
願你親自你的聖靈去挑釁我們的心.
挑釁我們傳福音.
挑釁我們祈禱.
挑釁我們去愛人靈魂的心.
讓我們不再冷淡.
不再去擔顧自己的事.
讓我們為其他人多走一步.
願我們在這個黑暗的世代.
成為明燈.
在這個末世裡面成為燈台.
願意和上帝你.
打這一場勝仗.
我們懇求你幫助我們.
聽我們在面前祈禱.
奉耶穌基督得勝之名祈求.
謝謝大家收看 再見.
\newpage



\section{但以理書 1:1-12}
\label{sec:Ik8b7Ffs2qU}
\textbf{2020年9月6日 崇拜直播｜但以理書1\_1-12}
\newline
\newline
連結: \href{https://youtube.com/watch?v=Ik8b7Ffs2qU}{\texttt{https://youtube.com/watch?v=Ik8b7Ffs2qU}} ~~~~ 語音日期: 2020-09-11
\newline
\newline
\hyperref[sec:TUMQvpeyy3M]{\small{< < < PREV SERMON < < <}}
~
\hyperref[sec:index_chronic]{\small{[返順時目]}}
~
\hyperref[sec:oCGsFJy1z54]{\small{> > > NEXT SERMON > > >}}
\newline
\newline
但以理書 1:1-12
\newline
\begin{longtable}{cl}
\hline
\hline
章節 & 經文 (和合本修訂版)\\
\hline
1:1 & \begin{tabularx}{0.7\textwidth}{X} 猶大王約雅敬在位第三年，巴比倫王尼布甲尼撒來到耶路撒冷，將城圍困。 \end{tabularx} \\ \\ \relax
1:2 & \begin{tabularx}{0.7\textwidth}{X} 主將猶大王約雅敬和神殿中的一些器皿交在他的手中。他就把他們帶到示拿地他神明的廟裡，將器皿收入他神明的庫房中。 \end{tabularx} \\ \\ \relax
1:3 & \begin{tabularx}{0.7\textwidth}{X} 王吩咐太監長亞施毗拿，從以色列人的王室後裔和貴族中帶進幾個人來， \end{tabularx} \\ \\ \relax
1:4 & \begin{tabularx}{0.7\textwidth}{X} 就是沒有殘疾、相貌俊美、通達各樣學問、知識聰明俱備、足能在王宮侍立的少年，要教他們迦勒底的文字和語言。 \end{tabularx} \\ \\ \relax
1:5 & \begin{tabularx}{0.7\textwidth}{X} 王從自己所用的膳和所飲的酒中，派給他們每日的分量，養育他們三年，好叫他們期滿以後侍立在王面前。 \end{tabularx} \\ \\ \relax
1:6 & \begin{tabularx}{0.7\textwidth}{X} 他們中間有猶大人但以理、哈拿尼雅、米沙利和亞撒利雅。 \end{tabularx} \\ \\ \relax
1:7 & \begin{tabularx}{0.7\textwidth}{X} 太監長給他們另外起名，稱但以理為伯提沙撒，稱哈拿尼雅為沙得拉，稱米沙利為米煞，稱亞撒利雅為亞伯尼。 \end{tabularx} \\ \\ \relax
1:8 & \begin{tabularx}{0.7\textwidth}{X} 但以理卻立志，不以王的膳和王所飲的酒玷污自己，於是懇求太監長容他不使自己玷污。 \end{tabularx} \\ \\ \relax
1:9 & \begin{tabularx}{0.7\textwidth}{X} 神使但以理在太監長眼前蒙恩，得憐憫。 \end{tabularx} \\ \\ \relax
1:10 & \begin{tabularx}{0.7\textwidth}{X} 太監長對但以理說：「我懼怕我主我王，他已經派給你們的飲食，何必讓他見你們的面貌比你們同年齡的少年憔悴呢？這樣，你們就使我的頭在王那裡不保了。」 \end{tabularx} \\ \\ \relax
1:11 & \begin{tabularx}{0.7\textwidth}{X} 但以理對太監長所派監管但以理、哈拿尼雅、米沙利、亞撒利雅的管理者說： \end{tabularx} \\ \\ \relax
1:12 & \begin{tabularx}{0.7\textwidth}{X} 「請你考驗僕人們十天，給我們素菜吃，清水喝， \end{tabularx} \\ \\
[1ex]
\hline
\hline
\end{longtable}
$^{1}$弟兄姊妹讓我們一起去聆聽一段經文.
在《但義利書》1-16節中.
經文是這樣說.
猶大王約瓦經在位第三年.
巴比倫王尼布·格尼薩來到耶路撒冷.
將成為群.
主張猶大王約瓦經並神殿中器命的幾分.
交付他手.
他就把這器命帶到士拿地.
收入他神的廟裡.
放在他神的腹中.
王勳父太監長阿斯皮拿.
從以色列人的宗族和貴族中.
帶進幾個人來.
就是年少沒有殘缺.
相貌俊美通達各樣學問.
知識聰明具備.
足能侍納在王宮裡的.
要教他們迦勒底的文字言語.
王派定將自己所用的善和所喝的酒.
每日賜他們一份養他們三年.
滿了三年就叫他們在王面前侍納.
他們中間有猶大族的人.
但以理 哈拿尼亞 米沙利 阿薩尼亞.
太監長給他們起名.
稱但以理為白提沙薩.
稱哈拿尼亞為沙德拉.
稱米沙利為米薩.
稱阿薩尼亞為阿伯尼戈.
但以理卻立志.
不以王的善和王所喝的酒.
玷污自己.
所以求太監長容他不玷污自己.
臣使但以理在太監長眼前蒙恩惠.
受憐憫.
太監長對但以理說.
我懼怕我主我王.
他已經派定你們的飲食.
倘若他見你們的面貌.
被你們同歲的少年人羈授.

$^{41}$怎麼好呢?.
這樣你們就使我的頭在黃那裡難保.
但以理對太監長所派管理但以理.
哈拿尼亞 米沙利 阿薩尼亞的委辦說.
求你斯斯薄人十天.
給我們素菜吃 白水喝.
然後看看我們的面貌.
和用王善那少年人的面貌.
就照你所看的待獄人罷.
委辦便允許他們這件事.
視看他們十天.
過了十天.
見他們的面貌比用王善的一切少年人.
更加俊美肥胖.
於是委辦撤去派他們用的善.
飲的酒 給他們素菜吃.
踏入九月份.
我們會開展另一個港島系列.
叫做「夾縫中的中瀑」.
我們會宣講但以理書.
透過但以理和那三個朋友.
被俘虜到外邦期間.
所發生的幾個宮廷故事.
來學習如何當我們面對一些限制.
一些挑戰 甚至逆境 逼不臨到的時候.
我們可以如何自力堅守信仰的原則.
特別在一些身不由己的環境.
仍然可以像廣東話說的「捐窿捐瀨」.
來設法在生活的一些夾縫中.
見證上帝的帶領和掌管.
事實上我們會看到.
在這六個宮廷的故事裡.
其實帶出了一個真理.
真正主宰歷史和世界局勢發展的.
並不是人世間的霸主君王.
而是上帝他自己.
就像我們剛才所讀出.
但以理書的第一章.
記載了當時我們說在近東世界裡.
一個最有影響力的霸主.

$^{81}$巴比倫王尼布格尼撒.
在猶大王約瓦勁在位的第三年.
他就揮軍圍困耶路撒冷.
強迫猶大國去臣服巴比倫.
尼布格尼撒更將聖殿部分器皿拿走.
帶回巴比倫.
就像戰利品一樣.
放在巴比倫的神廟裡.
這樣做一方面為了要去答謝巴比倫的神明.
來保佑巴比倫的大軍打敗了以色列.
同一時間尼布格尼撒亦都相信.
巴比倫的神比以色列的耶和華更加有能力.
更加的厲害.
巴比倫之所以能夠打敗猶大.
是因為巴比倫的神打敗了耶和華.
至於當時的以色列人.
特別是被俘虜到巴比倫的上帝的子民.
其實他們也在問相類似的問題.
為何猶大國會臣服巴比倫的下層.
為何上帝的聖殿器皿會輕易被敵人拿走.
為何上帝的子民明明是蒙福的.
卻被俘虜到外邦.
甚至淪為人質.
上帝難道真的不及巴比倫的神.
我們一直所信賴的耶和華.
難道已經沒有能力.
又或者他不再保守我們.
這些可能都是當時的以色列人.
他們會疑惑的事情.
不過你看到.
作者在但以理書一開首.
特別在第二節裡.
他清楚記載著.
不是耶和華不及巴比倫的神.
而是上帝故意讓約瓦敬和聖殿的器皿.
交給尼布甲尼撒王.
在第二節裡.
主張猶大王約瓦敬並神殿中器皿的幾分.
交付他(尼布甲尼撒)的手中.
換句話說.

$^{121}$作者開宗明義裡說了一個重點.
就是耶和華仍然是列邦列國的主宰.
在他的眼中.
尼布甲尼撒只是一個他可以使用的僕人.
只是他用來管教以色列的一個工具.
上帝由始至終都掌管著人世間的歷史和世界局勢的發展.
所以弟兄姊妹.
當我們看但以理書的時候.
作者提醒我們.
你要先戴上這個淑妮眼鏡.
你要從這個角度去了解.
到底這個世界是誰作主.
這樣我們才能夠在不同的境況裡.
可以好像故事裡.
但以理和他幾個朋友這樣站立得穩.
事實上你仔細看第一章.
你會發現其實在經文.
在第一章的一開頭和結尾.
其實他記載了兩個的時間坐標.
第一節裡面說的是猶大王若瓦競第三年.
說的是主前的605年.
到結尾了.
第一章的結尾.
第21節.
這裡說的是到了波斯古烈王的元年.
即是說主前的539年.
其實相差了超過60年.
正正就是但以理他在外邦事奉上帝的日子.
在巴比倫的皇宮或波斯的宮廷生活.
其實就成為了但以理見證上帝掌管一個的舞台.
他知道真正主場話事的其實是上帝.
而不是尼泊爾加利沙王或古烈王.
我們說這些如果用今天的類比.
就好像可能是我們的國家主席.
又或者美國總統.
甚至乎俄羅斯的普京.
這些歷史上的大國領袖或當權者.
上帝說不是他們.
是上帝自己掌管著.
所以你看到其實整個第一章裡面.

$^{161}$作者黑葉提到我們.
譬如說剛才說到第二節.
是上帝將若瓦競等人成為了人質.
交給尼泊爾加利沙.
第九節.
其實這裡也說.
神使到但以理和三個朋友.
在大鑑長面前蒙恩惠.
受連述.
是上帝出手幫助.
可以下下一張powerpoint.
並且第十七節.
這裡也說.
上帝在各樣的文字學問上.
賜給他們聰明知識.
是上帝幫助他們.
所以你看到其實短短的第一章經文.
裡面說到無論是國家的興衰.
無論是職場的人際關係.
甚至是個人前途的升跌.
全都在上帝掌握手裡.
經文某程度上給我們有色切的提醒.
相信我們不會有一個疑惑.
香港正正捲入兩個大國角力的漩渦裡.
無論你是什麼政見.
我們都面對這些情況.
外在環境不斷轉變.
你每天都會從不同的渠道裡面.
收到不同的資訊.
來轟炸你的思想.
有很多人都會擔心香港的前景.
擔心自己和家人的前途.
所以對基督徒來說.
這些困惑其實就像一個信仰測試劑.
他問我們一個問題.
到底我們相信誰掌管香港的局勢.
到底是誰保守帶領我們的前途.
假如我們的答案和但義利書一樣.
是上帝的話.
我們在人生做決定的時候.

$^{201}$我們就要好像故事裡面那幾位主角一樣.
我們要用聖經的教導.
用信仰原則來作為行事為人的依歸.
來致力在一些生活的納悶裡面.
你要幫自己體驗到.
上帝仍然掌管.
他仍然帶領著.
就好像第三至七節.
這裡他繼續說.
但義利和那三位朋友.
為什麼被選上去巴比倫的皇宮.
是因為他們本身有一些良好的條件.
你看到他們有優秀的外貌.
有學識.
有強的學習能力.
簡單來說.
用今天的術語.
這幾個就是男神.
所以他們被挑選入皇宮.
並且進行重點的培育.
讓他們有機會一方面.
接觸到當時的超級大國的文化.
可以學習到巴比倫的語文.
律法.
還有歷史.
並且可以學習到.
巴比倫人視為最高等的學問.
包括如何觀星像.
占星.
觀兆.
又或者法術命理.
解夢等等的技能.
並且不單單這樣.
你看到.
他們不單單為這四個人.
免費提供三年的宮廷訓練.
還包什麼?.
包伙食.
不是學生吃的餐蛋飯.
腿蛋飯和凍飲的頹飯.

$^{241}$這裡說什麼?.
每天都可以享用.
頂級的皇室料理.
當時被視為最上等.
最有營養的食物.
簡單來說.
皇帝吃什麼你吃什麼.
比訓練香港的政務官AO還厲害.
務求將這四個人打造到.
高大衛猛.
又容剛煥發.
又有智慧學問.
還學會禮儀.
能夠有足夠的台型.
可以被選上將來.
侍立在皇帝面前.
並且你看到這四個人.
連名字都被改了.
但你改了名字叫做「白提沙薩」.
另外那三個朋友.
被改了名字叫做「沙德拉米薩阿伯尼戈」.
很巧.
不過我曾經教過主人學.
陳玉峰弟兄說.
很容易記住這三個名字.
沙德拉米薩阿伯尼戈.
沙德拉米薩.
即是不要殺.
阿伯來的.
很容易記住.
沙德拉米薩阿伯尼戈.
沙德拉米薩.
阿伯來的.
很容易記住這三個名字.
這三個新的名字其實是說.
原來和巴比倫的偶像神明有關係.
所以你看到整個訓練配套.
加上改名.
很明顯是巴比倫一個管治政策.
要從飲食,教育,名字.

$^{281}$甚至從宗教各方面入手.
來進行同化的工程.
使這四個人逐漸忘記了.
自己原本有的國籍,名字.
身份和信仰.
來幫他們徹底融入.
巴比倫的生活文化.
並且將來可以替皇帝效力.
這是他們的政策.
而面對巴比倫這樣的管治.
其實你不難想像.
但你和這三個朋友.
其實很大掙扎.
首先能夠在皇室工作.
也就是意味著.
他們有機會享受到.
優質的生活.
試想像他們現在是什麼身份.
他們是俘虜,人質.
將來很大機會.
會變成奴隸或階下囚.
但現在帝國給他們.
一個飛黃騰達的機會.
其實有誰會不心動.
不過如果接受了.
皇帝提供的訓練套餐的話.
其實就代表著.
這四個人會不斷被.
外邦的文化和宗教.
入侵他們的生命.
上帝子民的身份.
也會越來越模糊.
到底是選擇接受皇帝的厚待.
還是拒絕呢?.
有沒有選擇?.
頂梓梅坦白說.
其實沒有什麼可以選擇.
因為如果拒絕反抗.
應該就是死路一條.
沒有選擇.

$^{321}$不過有趣的是.
你看到最後這四個人.
他們沒有選擇一直反抗到底.
但他們也沒有選擇.
百分之百全盤接受.
皇帝提供的套餐.
你留意到這四個人.
他們願意學習巴比倫的知識文化.
他們也不介意接受一個.
和巴比倫偶像有關的名字.
不過他們卻在飲食方面.
有一個不同的選擇.
第八節說.
「但以你立志不以皇的膳.
和皇所喝的酒玷污自己」.
當然我們很快就會理解到.
但以你認為皇的食物和酒.
一定是先用來祭祀異教偶像.
又或者但以你擔心.
皇的膳食一定會抵觸猶太的律法.
可能裡面含有猶太律法所說的不潔淨食物.
所以但以你拒絕進食.
確實有這個可能.
不過我們多看一點經文.
第十二至十三節.
但以你求管理他們的委辦.
就是這樣說.
「求你施施僕人十天.
給我們素菜吃白水喝.
然後看看我們的面貌.
和用皇膳那少年人的面貌.
就照你所看的待僕人罷」.
即是說他們只試多久?.
十天而已.
十天後如果行不通.
即是如果吃素菜.
喝白開水.
令自己皮黃骨瘦.
跟那些吃皇膳的少年人比下去的時候.
但以你他們會怎樣做?.

$^{361}$是否立即自盡殉道?.
不是.
他們只是說.
「我們吃回吧,吃回皇的膳食」.
將一切還原.
他們不會連累到委辦.
甚至太過尷尬.
換句話來說.
但以你這裡立志堅持.
在三年的皇宮訓練裡.
只靠素菜和白開水.
似乎不是擔心皇的飲食.
會使自己在禮儀上變成不潔淨.
而是這四個人.
想跟其他同期被選入宮的少年人.
劃一條界線.
將自己分辨出來.
什麼意思呢?.
再看多一段經文.
我們會明白多一點.
第十七至二十節.
剛才我們沒有讀到的經文.
這裡說.
「這四個少年人.
臣在各樣的文字學問上.
賜給他們聰明知識.
但以你有明白各樣的意象和夢想.
尼泊爾加里沙皇預定.
帶進少年人來的日期滿了.
太監長就把他們帶到皇面前.
皇與他們談論.
見少年人中無一人能比但以你.
哈拿利亞,米沙利,亞撒利亞.
所以留他們在皇面前侍立.
皇考問他們一切事.
就見他們的智慧聰明.
比通國的術士和用法術的勝過十倍」.
弟兄姊妹.
我們想一想.
這四個人.

$^{401}$但以你和他們三個朋友.
他們和其他同期選秀的少年人.
有什麼分別?.
你只看下去.
不知道內情的話.
他們真的沒什麼分別.
大家一起上課.
學習相同的巴比倫文化知識技能.
論外表.
大家都是怎樣?.
都是高大衛猛.
都是帥.
都是有學識.
懂王公禮儀.
所有東西都是一樣.
唯一不同的是什麼?.
唯一不同的是一群人.
在三年來.
喝王的酒.
吃王的膳.
但以你和他們三個朋友.
在三年來.
只吃素菜.
喝清水為生.
這就是他們的不同.
所以作者很想強調.
這四個人之所以能夠在人群裡.
突圍高升.
不是靠王的豐富公認.
而是在十七節說.
他們是靠上帝看顧他們.
所以你看到.
不是用我們今天的常識.
當時的人認為.
素菜白開水.
是沒有營養價值.
但以你他們四個.
都會這樣認知.
吃素菜白開水.
是會令自己疲惘骨瘦.

$^{441}$但假如他們這三年來.
竟然可以和那些.
吃王膳的少年人.
比他們強.
相比之下一點都不輸蝕.
甚至比他們帥.
這是什麼原因?.
唯一的解釋.
就是上帝出手.
施恩幫助這四個人脫穎而出.
所以你看到這個就是.
但以你他立志.
不用王的膳和王的酒.
其中一個很重要的原因.
三年以來他們每天都是.
吃素菜喝白開水.
來不斷勉勵提醒自己.
真正保守帶領我們的.
並不是巴比倫王提供的訓練.
提供的優厚待遇.
提供的豐富食物.
而是上帝自己.
是上帝主宰著我們的前途.
全靠上帝和進行祂的話語.
我們才可以傳媒.
才能夠高升.
試想一下弟兄姊妹.
假如這四個人同樣.
都是三年來吃王的膳食的話.
三年來不斷被世界的王.
來到後代填飽自己的生命的話.
這樣他們四個人的心靈.
確實又真的很難不被玷污.
因為日子越久.
他們又越難分別什麼?.
到底是誰在幫助我?.
到底是巴比倫王在幫助我.
還是上帝在幫助我.
養肥我.
使我可以高升?.

$^{481}$如果他們和其他人一樣.
都是吃王的食物和喝王的酒.
就連這個分別點都沒有了.
信仰就真的很模糊.
他們的生命就會被玷污.
同樣弟兄姊妹.
這段經文給我們提醒和挑戰.
我們會否意識到.
即使在世上.
看起來是最美善的事情.
都有可能玷污你和我的生命.
就如豐厚的飲食.
三年的優質教育.
其實都是世上美好的事物.
可以是上帝很大的恩典和祝福.
對不對?.
不過你們知道魔鬼撒旦.
很多詭計.
他也有辦法.
使這些看似美善的事物.
來將你和我叛道.
來引誘我們.
逐漸凡事只懂得仰望.
世界和世上的王的公認.
開始不仰望上帝的保守和帶領.
現實其實往往就是這樣.
我們不難明白.
基督徒被世界的價值觀填飽和蠶食.
通常不是突然.
是一個很緩慢.
往往是一個很緩慢.
不知不覺的過程.
世界的模式就會不經意地.
入侵主導你的生命.
直到有一天.
好像聖經所說.
我們漸漸愛自己.
愛世界.
多於愛上帝.
所以這段經文提醒我們.

$^{521}$千萬不要被世界的聲音.
掩沒了上帝的話語.
也不要為了迎合世界的模式.
而犧牲了上帝的吩咐.
我們剛才也說過.
最近香港的局勢.
引發很多人.
包括基督徒也不例外.
開始想一個話題.
移民與否?.
將不將我的資產.
轉移到其他地方?.
基督徒群體也會討論這些事情.
聽說有些教會.
因為回眾有需要.
所以索性在教會裡.
舉辦一些移民講座.
坦白說.
無論移民或不移民.
沒有對與錯.
因為上帝給予我們.
不同的領袖.
不同的經歷.
重點是無論我們選擇去或留.
也不能忘記.
我們的身份是什麼.
不能忘記我們去到哪裡.
都是上帝的子民.
屬神的兒女.
無論你住在哪個角落.
你也要像但義利和三個朋友.
立定你的心志.
靠主而活.
並且為主而活.
你要將信仰實踐出來.
這是經文提醒我們.
所以我們看到.
但義利和他三個朋友.
對他們來說.
信仰不是很虛浮的概念.

$^{561}$不是一個生活點綴.
特別當他們失去了.
昔日那些安定穩定的生活的時候.
甚至乎現在淪落到.
老年人身不由己的情況.
首先支撐著他們的生命.
讓他們能夠有正面積極的力量和態度.
繼續活下去的是什麼.
不是世界給予他們的待遇和保守.
是信仰.
是上帝給予他們.
繼續正面積極活下去的力量和理由.
而不是世界給予他們的後代和保障.
即使身處外邦.
身不由己.
信仰仍然能夠幫助他們.
提醒他們.
神仍然掌管一切.
神是超越一切地域限制.
上帝的子民仍然可以放心.
投靠他們.
信得過他們一切的安排.
這是經文第一章給予我們的重點.
不過.
你看到這四個人.
他們會覺得只有這個想法不夠.
不可以停留在腦裡.
他們覺得在這個極端的環境裡.
他們要用實際行動幫助自己.
深化和強化這個信念.
變成一個很堅實的信心.
不是光說光想.
他們有行動.
所以你看到這四個人.
他們在能力範圍內.
可以嘗試到的選擇.
他們嘗試在逆境裡.
把握了他們手上的資源和機會.
首先為上帝保存著聖名.
然後當自己站穩陣腳.

$^{601}$就找一些縫隙.
鑽孔鑽縫.
設法將他們的信仰實踐出來.
不單止這樣.
我們看到他們內心等候上帝給予的時機和大靈.
特別是稍後在第二節錄像.
你會看到上帝為他們預備更多的機遇.
讓他們可以為上帝做更大的見證.
所以弟兄姊妹都要經文去勉勵我們.
要將耶穌是主這個信仰宣告.
假如在你心裡你真的覺得.
耶穌是你的主宰的時候.
你就要想辦法貫徹落實在生命裡.
特別是面對今天的社會境況.
即使你和我只是一個小市民.
我們沒有能力去左右世界大局的發展.
但我們可以像這四個人一樣.
透過持續遵循上帝的吩咐.
從而體驗到上帝是你和我生命的主權.
你在當中設法去貫徹你的信仰.
透過一些未必會起眼的事情.
讓自己體驗到上帝是如何帶領保守緊利.
所以不一定要敲鑼敲鼓.
可以從一些很微小.
別人未必發現的事情做起.
可以不經意地.
不一定是口傳.
未必有社交距離的時候.
你可以用網上資訊.
發一個短訊WhatsApp.
或將教會一些信仰資訊.
送給你身邊未認識的朋友.
或你覺得他很有需要的朋友.
又或者你去關顧弱勢的需要.
又或者你珍惜.
上帝當下託付給你每一段的關係.
在裡面好好學習去愛神愛人.
亦都可以成為一個和平之子.
來為身邊可能有很多懼怕擔憂.
疑惑的人送上安慰和同行.

$^{641}$你亦都可以像聖經所說.
成為一個代求者.
守望身邊的人.
為上位的人來代求.
而當我們願意這樣去.
遵行上帝的吩咐的時候.
我們願意在多麼艱難困惑的日子.
我們都堅持遵行上帝的吩咐.
討祂喜悅的時候.
上帝就應許.
祂一定會賜福,施恩,憐憫幫助我們.
並且祂會逐漸透過你不同的事情.
來建立,擴大你對祂的信心.
就好像但義理和他三個朋友一樣.
真的,弟兄姊妹.
我們再看看.
到底但義理如何成功爭取到.
不用吃黃的膳和喝黃的酒.
我們看回第八節.
首先但義理最初游說的是誰?.
是太監長,對嗎?.
不過第十節.
太監長拒絕了但義理的請求.
原因是這是皇命.
假如太監長准許他們不吃黃的膳的時候.
第二天皇發現.
他們被其他同期入宮的少年人.
又瘦又黃的時候.
當皇知道真相的時候.
太監長肯定會人頭落地.
好像他自己所說.
不過,你有沒有留意一個很特別的作者.
刻意將第九節插進去?.
你覺不覺得有點奇怪?.
「神駛到但義理在太監長面前.
眼前蒙恩惠受憐憫.
明明對方已經拒絕了自己的要求.
為何還說我是蒙恩惠受憐憫?」.
其實你試試這樣讀.
讀完第八節,讀完第十節.

$^{681}$其實會比較通順.
但義理請求太監長這樣這樣.
太監長卻拒絕.
這樣讀下去很順暢.
為何作者刻意加第九節.
明明被拒絕了.
還是說他在上帝面前.
神駛到他在太監長面前蒙恩惠受憐恤呢?.
你試試想一想.
其實太監長有沒有需要跟他解釋這麼多?.
太監長的職位很高.
太監長一句「不行」,「不准」.
硬梆梆一句就拒絕了他.
不用跟他說那麼多.
但從他們的對話中.
看來太監長和但義理的關係都很好.
太監長願意和但義理交心.
說自己的軟弱和難處給但義理知道.
說他為何不願意和但義理冒這個風險.
弟兄姊妹.
這就是上帝很微妙的介入.
上帝在一些很微細,不起眼的.
奇妙介入他們短短的對話中.
向但義理施恩憐憫.
雖然他被太監長拒絕.
但他得到一個很重要的資訊.
他能夠摸到太監長的底線是甚麼.
他明白太監長最關注的不是.
他們是否吃黃的膳.
而是他們會否因為不吃.
而被其他人銷售.
是否收的問題,不是吃的問題.
這個資訊很重要.
換句話說,但義理在第十一節開始.
要調整整個策略.
他知道黃的政策是不能改變的.
但甚麼可以調教.
如何執行這個政策.
其實有很多空間和彈性.
正所謂我們都知道八個字.

$^{721}$上有政策,下有對策.
於是經文記載.
他轉移去遊說太監長的手下.
就是管理他所謂監護人的委辦.
求委辦讓他們試一試十天.
讓他們試試在這十天內.
只喝白開水和素菜.
然後,你看到他很聰明.
委辦是你自己決定.
你來判斷我們的容貌.
是否比其他吃黃膳的人差了.
你就按你的決定.
看十天後我們是繼續喝白開水和素菜.
還是吃黃膳.
他將決定權交給委辦手裡.
雖然他有這個建議.
所以萬一試完十天,真的不行.
可能在他們的理解裡.
只不過是上帝沒有預納我們的嘗試.
上帝沒有為我們這一方面開路.
那就唯有先吃飯.
再想其他的策略和方法.
看如何繼續為上帝作見證.
看是否行得通.
你看到四人不會因為信心超級膨脹.
而超越上帝的主權.
你看到他們一方面嘗試.
但一方面又敏銳.
上帝開不開路呢?.
上帝的指意是怎樣呢?.
而對委辦來說,他有沒有吃虧?.
首先你要明白.
他需不需要直接面對尼布甲尼薩王?.
不用的.
他只是面對太監長.
壓力小很多.
假如這十天裡這四人真的瘦了.
他最需要怎樣?.
將場地還原就行了.
逼迫他們吃飯.

$^{761}$你想想,有三年時間去追.
無論如何都要令這四人變胖.
但假如這四人的建議真的行得通.
那這三年王的膳食和王飲用的酒就歸給誰?.
就歸給委辦.
你想想,這麼好.
委辦自然就讓他們去試.
結果第十五節說.
過了十天之後.
他們的面貌確實比那些用黃膳的少年更加進味肥胖.
於是委辦就替他們撤去王的飲食.
容許他們這三年都以素菜清水為生.
想問問大家.
聽了這麼久.
問一個問題.
這三年來有誰知道故事的真相?.
有多少人?.
有.
但義理.
三個朋友.
還有呢.
委辦.
當太監長都知道了.
五至六個,加上上帝.
就這麼多.
但義理就是在沒有人察覺.
只有少許空間作主的情況下.
透過這件很小的事.
他如何經歷上帝在他生命裡的私恩憐憫.
也成為了他們一個信心的起步點.
但義理並不是第一天.
第一章被擄到巴比倫.
上帝就用最強的獅子坑或火窯來鍛鍊他們的信心.
上帝其實是用了很多時間.
透過不同的人生經歷.
讓他們有足夠的機會慢慢累積.
增加對上帝的信靠.
或者簡單說.
如果沒有第一章小事牛刀的信心立志.
其實很難可以累積到幾十年後.

$^{801}$新磡獅子坑那種非凡的信心表現.
所以弟兄姊妹.
經文給我們看到.
信心不只是我們被動等待上帝出手.
不是我們盼望上帝有一天突然用神蹟.
改變所有艱難的局面.
解決我們人生所有的問題.
更多的時候.
上帝是透過我們和他同工.
透過我們所手所作的來成就他的美意.
關鍵是我們能否對上帝保持開放的態度.
我們願不願意在不同的處境裡.
隨著事情不同的變化發展.
我們去不斷深入更新.
調教自己固有的想法或步伐.
來竭盡所能配合上帝的安排和計劃.
這是作者他勉勵我們.
就是要操練這種我們說.
經得起時代考驗的信心.
就好像二十一節怎麼說.
到了古烈王的元年.
但你還在.
即使經歷了改朝換代.
即使經歷了時代的洗禮.
但你還在.
因為上帝幫助他.
上帝看顧著他的前程和前路.
所以弟兄姊妹拋夢.
我們藉著這段經文.
我們一起勉勵和提醒.
我們去重建對上帝的信心和信靠.
不妨像學習他們四人一樣.
由一小步開始.
可能透過你的意念開始落實.
如何去遵行上帝的吩咐.
並且你沿途要敏銳.
上帝微小的大靈和恩典.
當你遇到阻力的時候.
不要輕言放棄.
你要耐心等待上帝的機遇.

$^{841}$上帝不同時間為你打開的路.
又或者你可以停下來想一想.
我還有什麼可以替上帝所做的呢?.
我有什麼可以為上帝工作.
討他喜悅?.
當然最重要的是.
你也要有屬靈的夥伴.
但以你不是一個人.
他有三個朋友.
他們會一起同行.
一起彼此守望.
一起去圍爐取暖.
面對艱難.
盼望上帝藉著這段經文.
去建立我們的信心.
認定他是掌管大局.
掌握我們前途的那一位.
也幫助我們.
能夠堅持走上帝的道路.
在困難的日子.
仍然堅持依靠他.
不依靠這個世界和這個世界的禍.
弟兄姊妹.
願意上帝藉著他的話語.
賜福給我們每一個.
也盼望我們用一首詩歌.
「我知誰掌管前途」.
來回應今天經文對我們的提醒和勉勵.
\newpage



\section{但以理書 2:1-30}
\label{sec:oCGsFJy1z54}
\textbf{2020年9月13日 崇拜直播｜但以理書2\_1-30}
\newline
\newline
連結: \href{https://youtube.com/watch?v=oCGsFJy1z54}{\texttt{https://youtube.com/watch?v=oCGsFJy1z54}} ~~~~ 語音日期: 2020-09-17
\newline
\newline
\hyperref[sec:Ik8b7Ffs2qU]{\small{< < < PREV SERMON < < <}}
~
\hyperref[sec:index_chronic]{\small{[返順時目]}}
~
\hyperref[sec:PTdjTBD7nsw]{\small{> > > NEXT SERMON > > >}}
\newline
\newline
但以理書 2:1-30
\newline
\begin{longtable}{cl}
\hline
\hline
章節 & 經文 (和合本修訂版)\\
\hline
2:1 & \begin{tabularx}{0.7\textwidth}{X} 尼布甲尼撒在位第二年，他做了很多夢，心裡煩亂，不能睡覺。 \end{tabularx} \\ \\ \relax
2:2 & \begin{tabularx}{0.7\textwidth}{X} 王吩咐人將術士、巫師、行邪術的和迦勒底人召來，要他們把王的夢告訴王；他們就來，站在王面前。 \end{tabularx} \\ \\ \relax
2:3 & \begin{tabularx}{0.7\textwidth}{X} 王對他們說：「我做了一個夢，心裡煩亂，想要知道這是甚麼夢。」 \end{tabularx} \\ \\ \relax
2:4 & \begin{tabularx}{0.7\textwidth}{X} 迦勒底人用亞蘭話對王說：「願王萬歲！請將夢告訴僕人，我們就可以講解。」 \end{tabularx} \\ \\ \relax
2:5 & \begin{tabularx}{0.7\textwidth}{X} 王回答迦勒底人說：「這事我已決定，你們若不把夢和夢的解釋告訴我，就必被凌遲，你們的房屋必成糞堆； \end{tabularx} \\ \\ \relax
2:6 & \begin{tabularx}{0.7\textwidth}{X} 但你們若能說出這個夢和夢的解釋，就必從我得到禮物、賞賜和殊榮。現在，你們要把夢和夢的解釋告訴我。」 \end{tabularx} \\ \\ \relax
2:7 & \begin{tabularx}{0.7\textwidth}{X} 他們再一次回答說：「請王將夢告訴僕人，我們就可以講解。」 \end{tabularx} \\ \\ \relax
2:8 & \begin{tabularx}{0.7\textwidth}{X} 王回答說：「我確實知道你們是故意拖延，因為你們知道這事我已決定。 \end{tabularx} \\ \\ \relax
2:9 & \begin{tabularx}{0.7\textwidth}{X} 你們若不將夢告訴我，只有一個辦法對待你們；因為你們彼此串通，向我胡言亂語，要等候情勢改變。現在，你們要將夢告訴我，讓我知道你們真能為我解夢。」 \end{tabularx} \\ \\ \relax
2:10 & \begin{tabularx}{0.7\textwidth}{X} 迦勒底人回答王說：「世上沒有人能解釋王的事情；從來沒有君王、大臣、掌權者向術士、巫師，或迦勒底人問過這樣的事。 \end{tabularx} \\ \\ \relax
2:11 & \begin{tabularx}{0.7\textwidth}{X} 王所問的事很難，除了不與血肉之軀同住的神，沒有人能在王面前解釋。」 \end{tabularx} \\ \\ \relax
2:12 & \begin{tabularx}{0.7\textwidth}{X} 王因這事生氣，大大震怒，吩咐滅絕巴比倫所有的智慧人。 \end{tabularx} \\ \\ \relax
2:13 & \begin{tabularx}{0.7\textwidth}{X} 命令發出，智慧人將要被殺，人就尋找但以理和他的同伴，要殺他們。 \end{tabularx} \\ \\ \relax
2:14 & \begin{tabularx}{0.7\textwidth}{X} 王的護衛長亞略奉命去殺巴比倫的智慧人，但以理用婉言和智慧回應， \end{tabularx} \\ \\ \relax
2:15 & \begin{tabularx}{0.7\textwidth}{X} 向王的大臣亞略說：「王的命令為何這樣緊急呢？」亞略就把事情告訴但以理。 \end{tabularx} \\ \\ \relax
2:16 & \begin{tabularx}{0.7\textwidth}{X} 於是但以理進去求王寬限，好為王解夢。 \end{tabularx} \\ \\ \relax
2:17 & \begin{tabularx}{0.7\textwidth}{X} 但以理回到他的居所，把這事告訴他的同伴哈拿尼雅、米沙利、亞撒利雅， \end{tabularx} \\ \\ \relax
2:18 & \begin{tabularx}{0.7\textwidth}{X} 要他們祈求天上的神施憐憫，將這奧祕指明，免得但以理和他的同伴與巴比倫其餘的智慧人一同滅亡。 \end{tabularx} \\ \\ \relax
2:19 & \begin{tabularx}{0.7\textwidth}{X} 這奧祕就在夜間異象中顯明給但以理，但以理就稱頌天上的神。 \end{tabularx} \\ \\ \relax
2:20 & \begin{tabularx}{0.7\textwidth}{X} 但以理說：「神的名是應當稱頌的，從亙古直到永遠！因為智慧和能力都屬乎他。 \end{tabularx} \\ \\ \relax
2:21 & \begin{tabularx}{0.7\textwidth}{X} 他改變時間、季節，他廢王，立王；將智慧賜給智慧人，將知識賜給聰明人。 \end{tabularx} \\ \\ \relax
2:22 & \begin{tabularx}{0.7\textwidth}{X} 他顯明深奧隱祕的事，洞悉幽暗中的一切，光明也與他同住。 \end{tabularx} \\ \\ \relax
2:23 & \begin{tabularx}{0.7\textwidth}{X} 我列祖的神啊，我感謝你，讚美你，因你將智慧才能賜給我，我們所求問的現在你已指明給我，把王的事給我們指明。」 \end{tabularx} \\ \\ \relax
2:24 & \begin{tabularx}{0.7\textwidth}{X} 於是，但以理進到王所派滅絕巴比倫智慧人的亞略那裡去，對他這樣說：「不要滅絕巴比倫的智慧人，求你領我到王面前，我可以為王解夢。」 \end{tabularx} \\ \\ \relax
2:25 & \begin{tabularx}{0.7\textwidth}{X} 亞略就急忙領但以理到王面前，對王這樣說：「我在被擄的猶大人中找到一人，能將夢的解釋告訴王。」 \end{tabularx} \\ \\ \relax
2:26 & \begin{tabularx}{0.7\textwidth}{X} 王對那稱為伯提沙撒的但以理說：「你能將我所做的夢和夢的解釋告訴我嗎？」 \end{tabularx} \\ \\ \relax
2:27 & \begin{tabularx}{0.7\textwidth}{X} 但以理回答王說：「王所問的那奧祕，智慧人、巫師、術士、觀兆的都不能告訴王， \end{tabularx} \\ \\ \relax
2:28 & \begin{tabularx}{0.7\textwidth}{X} 只有那在天上的神能顯明奧祕。他已把日後將要發生的事指示尼布甲尼撒王。你在床上做的夢和你腦中的異象是這樣： \end{tabularx} \\ \\ \relax
2:29 & \begin{tabularx}{0.7\textwidth}{X} 你，王啊，你在床上所思想的是關乎日後的事，那顯明奧祕的主已把將來要發生的事指示你。 \end{tabularx} \\ \\ \relax
2:30 & \begin{tabularx}{0.7\textwidth}{X} 至於我，那奧祕顯明給我，並非因我智慧勝過一切活著的人，而是為了讓王知道夢的解釋，知道你心裡的意念。 \end{tabularx} \\ \\
[1ex]
\hline
\hline
\end{longtable}
$^{1}$等子妹讓我們一同去聆聽一段上帝的話語.
記載在《但義理書》第二章一至三十節.
經文是這樣說.
「尼泊爾·加利薩在位第二年.
他作了夢,心裡煩亂,不能睡覺.
王吩咐人將術士用法術的,行邪術的和迦勒底人召來.
要他們將王的夢告訴王.
他們就來站在王前.
王對他們說:我作了一夢,心裡煩亂.
要知道這是什麼夢.
迦勒底人用亞蘭的言語對王說:.
願王萬歲,請將那夢告訴僕人.
僕人就可以講解.
王回答迦勒底人說:夢我已忘了.
你們若不將夢和夢的解說告訴我.
必被凌遲.
你們的房屋就必成為糞堆.
你們若將夢和夢的講解告訴我.
就必從我這裡得贈品和賞賜,並大榮耀.
現在你們要將夢和夢的講解告訴我.
他們第二次對王說:.
請王將夢告訴僕人.
僕人就可以講解.
王回答說:我準知道你們是故意遲延.
因為你們知道那夢我已經忘了.
你們若不將夢告訴我.
只有一法待你們.
因為你們預備了謊言亂語向我說.
要等候時勢改變.
現在你們要將夢告訴我.
因我知道你們能將夢的講解告訴我.
迦勒底人在王面前回答說:.
世上沒有人能將王所問的事說出來.
因為沒有君王大臣掌權的.
向術士或潤法術的或迦勒底人問過這樣的事.
王所問的是甚難.
除了不與世人同居的神明.
沒有人在王面前能說出來.
因此王氣憤憤地大發烈怒.
吩咐滅絕巴比倫所有的哲士.

$^{41}$於是命令發出:哲士將要見殺.
人就尋找但以理和他的同伴.
要殺他們.
王的護衛長亞略出來.
要殺巴比倫的哲士.
但以理就用軟言回答他.
向王的護衛長亞略說:.
王的命令為何這樣急速?.
亞略就將情節告訴但以理.
但以理誰進去求王歡善.
就可以將夢的講解告訴王.
但以理回到他的居所.
將這事告訴他的同伴哈拿尼亞.
米沙利亞撒尼亞.
要他們祈求天上的神斯漣敏.
將這奧秘的事指明.
免得但以理和他的同伴.
與巴比倫其餘的哲士一同滅亡.
這奧秘的事就在夜間異象中.
給但以理顯明.
但以理便稱頌天上的神.
但以理說.
神的名是應當稱頌的.
從耿古直到永遠.
因為智慧能力都屬乎他.
他改變時間.
日其輝煌立王.
將智慧賜予智慧人.
將知識賜予聰明人.
他顯明深奧隱秘的事.
知道暗中所有的.
光明也與他同居.
我列祖的神啊.
讚美你.
因你將智慧財能賜給我.
允准我們所求的.
把王的事給我們指明.
於是但以理進去見亞略.
就是王所派滅絕巴比倫哲士的.
對他說.

$^{81}$不要滅絕巴比倫的哲士.
求你領我到王面前.
我要將夢的講解告訴王.
亞略就急忙將但以理.
領到王面前對王說.
我在被擄的猶大人中遇見一人.
他能將夢的講解告訴王.
王問稱為伯提沙薩的但以理說.
你能將我所作的夢和夢的講解.
告訴我嗎?.
但以理在王面前回答說.
王所問的那奧秘事.
哲士用法術術士觀笑的.
都不能告訴王.
只有一位在天上的神.
能顯明奧秘的事.
他已將日後必有的事.
指示彌布甲尼撒王.
你的夢和你在床上.
腦中的異象是這樣.
王啊.
你在床上想到後來的事.
那顯明奧秘事的主.
把將來必有的事指示你.
至於那奧秘的事顯明給我.
並非因我的智慧勝過一切活人.
乃為使王知道夢的講解.
和心裡的思念.
今天早上有陳珍珍姑娘.
在我們當中繼續宣講.
剛才那一段經文.
講題是.
「隔逢中的衷瀑轉危為機」.
請珍珍姑娘.
各位弟兄姊妹平安.
很想和大家分享一個故事.
這個故事是兩方對陣.
而且這個對陣是一個.
強弱懸殊的對陣.
一方是一個九尺多高.

$^{121}$戴著銅柱的頭盔.
穿上盔甲.
有護膝.
有重型的矛和盾和槍的人.
他是一個具有.
很豐富戰士經驗的戰士.
而且在他身邊.
有一大群士兵.
為他叫囂,吶喊,坐鎮.
另一方的對手是誰呢?.
他的對手.
首先他的身形矮了一大截.
而他身上連一些.
很普通的盔甲都穿不起.
手上只拿著一個.
丫叉或叫做脫石的機銀.
和幾塊石頭.
他的職業並不是一個戰士.
而是一個看守羊群的牧羊人.
我想問問你們.
在你們的眼中.
誰會贏呢?.
聰明的你們一聽就知道.
這個故事是在說.
哥利亞和大衛的對陣.
不過大家都知道.
最終勝利的一方是大衛.
因為哥利亞所依靠的.
是他身上的裝備.
他自己的能力.
但是大衛所依靠的.
就是那位與他同在的耶和華.
我們今天說到.
人世間有很多危機.
但是如何可以轉危為機呢?.
關鍵就在於我們依靠誰.
我們由今個月開始.
連續兩個月.
我們都會在港島裡.
分享《但義理書》.

$^{161}$上回港島但義理和他三個朋友.
這四個猶太人成為亡國奴.
進入了巴比倫的宮廷受訓.
他們定義選擇.
不要被這個世界填滿他們的人生.
他們定義要專心養賴神.
不去玷污自己.
以至不去吃黃的膳食和黃飲的酒.
選擇以清水素菜作為他們的膳食.
這是一個靈性的操練.
是一個自己拿來的考驗.
我記得上次Raymond跟我們說.
我們需要在這些操練裡.
耐心等候機遇.
為神做見證.
今天講的第二章.
告訴我們這些考驗已經來到了.
這個考驗是一個極大的威脅.
甚至是一個關乎生死的危機.
不過究竟這個危機.
是一個危機還是一個機遇呢?.
那就要看被考驗的人.
究竟如何面對.
剛才我們讀到.
但義理書第二章一至十三節.
故事發生在尼布甲尼薩在位的第二年.
即是但義理他們被擄的第三年.
因為在巴比倫的曆法來說.
猶太人被擄的第一年是登基年.
被擄的第二年是在位的第一年.
順理成章巴比倫的曆法.
說到尼布甲尼薩在位第二年.
就即是猶太人被擄的第三年.
按照第一章的記載.
但義理他們在巴比倫的宮廷受訓三年.
即是說現在這個時候.
但義理他們已經成為在宮廷裡的.
哲士的一份子.
故事一說到尼布甲尼薩.
面對一個難題.

$^{201}$這個難題是什麼呢?.
他做夢.
但他做完這個夢之後.
他心裡很焦慮,很煩擾.
睡不著覺.
歸為一國之君.
他可以說是要風得風,要雨得雨.
他剛剛才打贏周圍的敵人.
建立了強大的王國.
但突然的這個夢.
就令他輾轉反側.
不能入眠.
極之的焦慮.
我相信你作為都市人都很明白.
如果能夠有一個好睡覺.
其實是一個福氣.
偏偏這個王就遇到這個大能題.
尼布甲尼薩是一個有權有位的王.
於是他召集了一群.
聖經說「術士,用法術的,行邪術的」.
和迦勒底人.
這群是甚麼人呢?.
這群其實是占卜星象.
和會和鬼神通靈的人.
他們也懂得行邪術.
而迦勒底人的意思.
其實是等同占卜的專家.
他們這群人.
就是統稱為「哲士」的一群人.
他們在宮廷裡服侍.
為甚麼他們會被稱為「哲士」呢?.
就是因為外邦人認為.
能夠去測天象,玩法術的.
這些人必定有高人一等的智慧.
所以王認為他們可以幫他們解決.
或者舒緩他們心裡的不安.
於是就要他們來到他們面前.
要求他們將他所發的夢說出來.
然後解夢給他們聽.
為甚麼尼泊爾要這樣做呢?.

$^{241}$是否就像第五節所說的.
「夢我已經忘了」.
他們是否真的忘記了夢呢?.
其實在和合本裡有個「細」字.
「細」字寫著「我已命定」.
在「我已經忘了」這個字後面.
寫著「我已命定」.
其實一般我們認為.
他們不是真的忘記了夢.
「我已命定」的意思是說.
我已經決定了.
我要頒布以下的命令.
就是這個意思.
如果他們真的忘記了夢.
他們之後又怎能判斷.
這份節時所說的夢是否真實呢?.
為甚麼這個夢會令他們如此心煩意亂呢?.
其實他們是知道這個夢的.
但為甚麼他們要這樣做呢?.
是因為他們認為夢所要啟示的東西.
是非常非常重要.
令他們心水不寧.
他們要確保這班人解這個夢.
的答案是真實和可信的.
所以他們要測試.
要這班人說出這個夢的真相.
而他們定義的是甚麼呢?.
他們定義的是兩方面.
如果你們不能解釋這個夢的話.
就要遭受凌遲之刑.
甚麼是凌遲呢?.
凌遲是指要將他們四肢斬斷的刑罰.
是很殘忍的刑罰.
但與此同時他們亦說.
如果你們能夠說出這個夢.
成功地解釋這個夢的時候.
必定會賞賜禮物和享有尊榮.
聖經告訴我們.
到後來這班人.
完全是無能為力的時候.

$^{281}$尼泊爾撒旺勃然大怒.
他的刑罰不單止是凌遲這麼簡單.
而且是要將他們全部殺死.
要滅絕他們.
當面對艱難的時候.
心情很焦慮很不安.
自然我們就會用盡一切.
我們可用的資源去解決.
甚至用一些壓迫的方法.
去令其他人就範.
去協助他解開他心底裡的難題.
和心裡的焦慮.
但結果徒勞無功.
而且還衍生更多更多的問題.
我們看到尼泊爾撒遇到問題.
他靠的是什麼呢?.
我們看到他靠的是自己的王權.
靠的是他身邊的資源.
但結果並不能解決他當前的問題.
我們看完尼泊爾撒旺.
我們看看這班被召入宮的摺屎.
他們這班摺屎除了面對一個難題之外.
還面對一個極大極大的危機.
他們面對的是什麼難題呢?.
最初這班摺屎想像平時那樣.
入宮為王解夢應該沒有問題.
因為他們都是受過訓練.
他們知道如何解夢的一班人.
如果像平時那樣.
王將他所發的夢告訴他們.
按照他們所學的占卜.
觀星象的方法就可以解決.
但這次入宮.
萬萬想不到.
王要求他們將他的夢說出來.
正如第十節.
第十節說到加勒底人說.
根本你給我們這個要求.
世上沒有人能像王那樣問出來.
這是什麼意思呢?.

$^{321}$世上沒有一個位高權重的人.
問過這麼難的問題.
然後第十一節他們說得很對.
除了不同世人.
世人的意思.
另一個說法就是屬血氣的人.
除了不同屬血氣的人同居的神明之外.
沒有人能夠在王面前說出來.
巴比倫是一個多神的國度.
這班節時所信靠的神明.
就是在地上跟他們同居的.
一班所謂的神明.
是人手所雕刻出來的偶像.
他們根本沒可能幫助到這班人.
去解答王的問題.
去完成這個不可能的任務.
於是這個難題接著.
就是關乎一個生命危機的問題.
他們將要面對酷刑.
他們將要面對生命受威脅.
他們的生命和財產都會喪失的一個危機.
如果是你,你會怎樣?.
如果你面對危機,你會怎樣?.
好像這班術士.
他們面對這個危機.
是到了一個極之無助的關口.
這個王跟平時的做法完全不一樣.
根本不按常規.
根本不按平常理.
而且他表明了一些.
用一些殘酷的刑罰去對待他們.
要求他們要說出所發的夢.
王看穿了他們.
他們根本無計可施.
在故意的拖延時間.
一再問王,究竟你的夢是什麼?.
王認為他們根本既然不知道夢是什麼的時候.
其實也不會懂得解那個夢.
所以他認定他們是說謊言的.
對他們完全失去信任.

$^{361}$推斷他們心裡的期望.
就是如果將來時移勢易.
政權被顛覆.
或者因為王的喜好.
而將這件事遺忘了.
忘記了對他們進行私刑的計劃的時候.
他們就有彎轉.
其實他們是無助的.
他們是等運到的.
他們是等這個環境能夠轉變.
我們看到這班外邦的哲士.
他們是靠著什麼方法去面對他們的危機?.
他們其實就是靠著.
他們所信任這班外邦的神明.
但到頭來他們所得到的.
也是無助,無出路,失敗的.
其實在今天我們很多人.
都為著我們未知的未來.
面對一些危機的時候.
我們會去尋求其他神明的協助.
甚至在基督徒當中.
我們都會有人去占卜.
有人去玩塔羅牌.
去提升,借助一些水晶的力量.
或者去尋求一些.
姓鍾的師父,姓蘇的師父.
去了解一下自己的運程.
但我們要知道.
原來倚靠這些力量.
最終都是失敗的.
我們來到第十三節.
我們看到這個危機.
不單是站在皇面前的這班哲士面對著.
我們看到第十三節.
皇要去將命令發出.
所有的哲士都受到刑罰的威脅.
當中包括了我們今天要去學校的主角.
但以理和他們的同伴.
人人都面對同一個危機.
大家的反應都有所不同.

$^{401}$但以理的反應又是怎樣呢?.
我們看到第十四至第十九節.
當但以理面對這個危機的時候.
他的態度是從容不迫.
我們可以看到幾個表現.
第一是他鎮定地回應.
他並沒有逃避這個危機.
當皇要被殺的命令來到的時候.
其實但以理完全不知道發生什麼事.
他只知道現在命令一出.
死就是他現在要去接受的刑罰.
你想想至高無上的皇.
下了這道命令.
還有反抗的餘地嗎?.
我們說這兩個月的主題是.
身處甲縫的中瀑.
我在想他現在在這個甲縫裡.
好像找不到一條縫隙.
究竟怎麼辦呢?.
聖經提到第十四節.
當護衛長來到的時候.
這個阿略來到的時候.
但以理的回應是用軟言回答他.
什麼是軟言呢?.
軟言就是謹慎和機警地回應.
這是必須在危急的關頭裡.
保持鎮定才能做到的行動.
相對於那些哲士在皇面前的表現.
但以理在這一刻的從容不迫.
很值得我們留意和學習.
的確在危急的關頭裡.
能夠保持鎮定.
才能找到那條縫隙來捐的出路.
他問阿略.
他說皇的命令為什麼這麼緊急?.
其實緊急這個字指示出來的.
指示出來的就是說.
為什麼這個命令這麼殘忍和殘酷?.
於是阿略就把事情告訴他.
但以理知道皇所做的夢不能解決.

$^{441}$以至動了殺機.
他知道皇不是無緣無故.
拉起一根筋帶開殺戒.
他明白到問題的關鍵所在.
之後他不是想著保自己的命.
他不是想著如何脫身,如何逃脫.
而是他決定要勇敢地去見皇.
因為之後我們說到.
是解決他和其他哲士.
性命威逼的一個方法.
其實這是一個很冒險的行動.
你想想,如果但以理去到皇的前面.
皇的怒氣未散,正在發難的時候.
他去求皇的歡善.
你不可以肯定皇會不會像他認為的那些哲士一樣.
覺得但以理其實是無能的.
只是用拖字訣來解決方法.
如果皇癲起來,可以即時去問斬但以理.
但但以理並無畏縮.
他相信在他背後的上帝必然會保護他.
於是我相信他去到皇的面前.
一定會用剛才所說的軟言.
去向皇求歡善.
結果皇真的給他一條生路.
他的鎮定和勇敢這個特質.
可能在不同的人身上都可以展現出來.
但聖經告訴我們.
但以理之不同的原因.
就是因為他有一個很強大的後著.
大家看到嗎?.
當他得到歡善的時候.
他第一件事是做什麼?.
是回家.
他回家做什麼?.
他回家是邀請那些和他一起住.
三位出生入死的好兄弟.
哈拿利亞,米沙利和亞撒利亞.
和他一起去禱告.
邀請那位在天上的神.
去介入他們在地上的這件事.

$^{481}$介入他們當中.
他的禱告很簡單.
他知道皇的命令.
皇的夢的意象是從神而來.
上帝才是智慧的源頭.
所以他要找的就是上帝.
他的禱告很簡單.
是很直接的呼求.
直接將他自己所想的告訴神.
祈求天上的神斯連文.
將這個奧秘揭示給他知道.
因為他不想死.
他不想他三個同伴死.
他也不想巴比倫其餘的哲士一起滅亡.
但以你在千軍一髮之際.
仍然能夠保持這個高度的鎮定.
並且夠機警.
謹慎地去面對護衛長.
並且勇敢去見皇.
去取得歡閒之期.
在沒有什麼罅隙的地方找到罅隙.
找到出路.
結果第十九節告訴你.
神就將這個奧秘.
在夜間向但義利顯現了出來.
我們看到但義利和之前的.
尼泊爾賈尼撒王和哲士很不同.
他這一切的行動所表明的.
是他不是靠什麼.
他唯一靠的就是上帝.
我記得十多年前.
我剛回到武會實習的時候.
認識了一個大專團契的姐妹.
和她的媽媽.
這位姐妹就讀在名校裡面.
她的會考成績很優異.
所以在中五會考放榜之後.
她就已經拿到中大的暫取生資.
只要完成中六的課程.
她就可以進入中大讀書.

$^{521}$但就在她讀中六那一年.
她發生了一宗很嚴重的交通意外.
她被車撞昏迷.
腦部受到嚴重的創傷入院.
生命危在旦夕.
她的媽媽是一個很愛主的信徒.
當她媽媽回憶.
在人面前見證這件事的時候.
她說:我也不知道.
為什麼當時的我這麼鎮定.
她說當時想起.
雖然她不是回教會很久.
當時她只是回教會短短一兩個月.
但她已經立即打電話給教會.
懇請大家為她代禱.
她的女兒剛回到我們的大專團契一個星期.
我們只認識她一次.
她也要求教會通知團契導師.
發散團友為她女兒禱告.
她說當她在手術室門口.
她向上帝祈禱.
她的禱告是這樣的.
她說:上帝啊.
我知道你愛我的女兒比我愛得更多.
我也深信你滿有大能.
如果你說你一定要醫好我的孩子.
對你來說是輕而易舉的事.
不過我信服你.
我深信你要他所經歷的.
一定是為他最好的.
我將他交給你.
無論你要他是生是死.
我都知道你必然是對他最好.
她說她的心裡是出奇地鎮定和平安.
然後向上帝做了這樣的禱告.
很奇妙.
結果這個女兒不單救得了.
而且她在這一年後順利入讀醫科.
現在是一個真心服務病人的內科醫生.
這是何等大的奇蹟.

$^{561}$但我更加記得她媽媽當時的鎮定.
當時對上帝的信靠令我很動容.
因為我看到她完全相信.
對上帝的信靠是她唯一的出路.
接著我們看20至30節.
我們相信當人面對生命危險和威脅的時候.
其實你所表現的態度和回應方法.
是源於你的信念.
信主的人我們會問.
其實你信靠的上帝是一個怎樣的上帝?.
讓我們一起看看.
但以你所信靠的上帝是一個怎樣的上帝?.
當但以你在清醒的時候.
去領受到這個王所發出的夢的意象的時候.
他第一件事不是衝去皇宮.
他第一件事是獻上了一個頌讚的讚美詩.
第20節他說.
「神的名是應當稱重的.
從亙古直到永遠.
因為智慧能力都屬乎他.
他改變時候日期.
輝煌立王.
將智慧賜予智慧人.
將知識賜予聰明人.
他顯明深奧隱蔽的事.
知道暗中所有的光明也與他同居.
我列祖的神啊.
讚美你.
因你將智慧才能賜給我.
允准我們祈求的.
在皇的事給我們指明」.
這篇詩非常非常重要.
是整個故事發展的轉捩點.
因為我們在這裡看到.
真正面對這個危機.
這個必死無疑.
無彎可轉的危機.
扭轉的是這位真正有智慧和能力的上帝.
這首詩告訴我們.
他所信教的上帝是一個怎樣的上帝?.

$^{601}$是有智慧和有能力的上帝.
首先他的能力是什麼?.
他說他能夠改變時候日期.
他是一個掌管大自然的上帝.
他輝煌立王.
他是掌管歷史的上帝.
世世歷代.
他的興衰變更.
他都掌管著.
他勝過地上一切的君王.
包括尼布格尼撒.
但你知道.
他無需要對地上的王畏懼.
因為他倚靠的那個王.
萬王之王的上帝.
才是大有能力智慧的上帝.
而這個上帝是滿有智慧.
他的智慧是超脫的.
因為他掌管歷史.
他洞悉一切.
知道過去未來.
並且他的智慧.
是會賜智慧給那些聰明智慧人.
勝過一切的占卜家.
在這個故事裡.
我們可以看到.
神將這個深奧隱蔽的事.
沒有人能夠知道的事情.
指明給但以理知道.
但姐妹我們是上帝的子民.
他都會將他的智慧能力賜給我們.
他才是智慧能力的來源.
他可以將智慧能力賜給我們每一個.
有些解經學家告訴我們.
智慧其實可以分兩種.
一種是我們日常生活裡行事為人.
一些人生的哲理.
是經驗的累積.
這是我們在《詩歌智慧書》.
常見的一些智慧的名言.

$^{641}$但另一種智慧.
就是上帝揭示奧秘的智慧.
是從上帝的啟示而來的智慧.
或者你和我都很羨慕.
很希望有但以理擁有.
這種從神而來.
知道知識奧秘的能力.
但是你知不知道.
其實我們都是有智慧的人.
因為其實在新約裡.
就正正揭示了一個救贖的奧秘.
就是我們能夠.
因著基督耶穌的緣故.
我們得蒙救贖.
和神建立一個和好的關係.
我們的生命是被揀選的一群.
基督教所說的智慧的根本.
其實就是一份關係.
鄒賢告訴我們.
認識耶和華就是智慧的開端.
認識智性者就是聰明.
但以理的經驗也告訴我們.
他有這些智慧的能力的來源.
是源於他和神的關係.
他和三位出生入死的至親好友.
一起向神呼求的一個關係.
然後神就按著他的憐憫.
去作出回應.
即是我們要提醒自己.
在我們日常的生活裡.
我們常常要提醒自己.
其實有一個力量比我們更大.
是在我們自身以外.
那就是我們和基督耶穌的關係.
因為我們和耶穌的關係.
讓我們帶著基督的眼光去看這個世界.
去看各樣的事情.
去教我們如何過世上的生活.
所以在我們的生活裡.
其實對神的信靠.

$^{681}$就自然能夠活出這份智慧.
去到第24至30節.
但以理知道這個夢的時候.
他就立刻聯絡王的護衛長阿略.
請他不要殺其他哲士.
引薦他見王解夢.
當但以理去到王的面前表達的時候.
他跟王說:.
「王啊,你之前在床上做的夢.
是一位天上的神.
要告訴你未來的事.
他現在要藉著這個僕人向你講解.
接下來的31至45節.
但以理就講出這個夢.
我們今天不會提到這一段經文.
但我們必須要留意的就是.
大家有沒有留意.
但以理在24至30節裡.
他一再強調.
並不是他自己有什麼個人之處.
不是他自己比其他哲士聰明.
是因為要你發這個夢的神.
要藉著我向你講解未來的事情.
這也提醒到我們.
我們要得著智慧的動機是什麼.
但以理為了保存多人的性命.
他每一次都很清楚地表達.
他之所以可以做到.
並不是他有什麼個人之處.
他清清楚楚地去高舉上帝.
說明是上帝的作為.
而且他知道這是神的心意.
他只要跟著做.
上帝的旨意就能夠成就.
所以一個智慧的人.
一個能夠領受神智慧的人.
他必須是一個謙卑.
忠心服侍主的人.
他就能夠展現真正的智慧.
最後我們會讀46至49節.

$^{721}$這一節的經文.
剛才我們讀經時沒有讀過.
但我們要去看看一個危機.
最後是怎樣變成一個轉機.
當但以理完成任務後.
當但以理向著這位.
極之地上至高無上.
殘酷不仁.
亦極之不安的尼布格尼撒王.
說完上帝的傲弊事的時候.
尼布格尼撒王立即俯伏在地上.
他面貼於地.
但最吊詭的是.
他向一個亡國奴但以理下拜.
並且奉上公物和香.
他高歌說.
「你既能夠講明這個傲弊.
你們的神成然是萬神之神.
萬王之王.
是傲弊的啟示者」.
這就是這個故事的最高峰.
由這個王講出來.
上帝的宣認是一個很諷刺的位置.
這個掌權的王.
向他手下拜將一群亡國奴跪拜.
是因為這群亡國奴後面的上帝.
才是掌管萬有.
包括世上神靈偶像.
和各國各民各族領袖的上帝.
我們看到當但以理謙卑行的時候.
這位高舉上帝的人.
就是能夠在他的生命裡.
彰顯上帝的作為.
叫人真真正正認識上帝是誰.
於是最後危就化為機.
王使道但以理高聲.
賞賜他極多的禮物.
派他管理巴比倫全省.
又納他做總理.
掌管巴比倫所有的節食.

$^{761}$但以理求王.
他沒有忘記他的朋友.
求王將他三個朋友.
沙特拉,米薩和阿伯尼戈.
一起去管理巴比倫省的事務.
而但以理繼續在這個朝裡面.
去侍納服侍這個皇帝.
其實一個飛離的還和.
在完全沒有預兆的情況下.
本來是要死的.
現在不單止不用死.
但以理還有機會在王的面前侍納.
繼續將神高舉去見證上帝.
危機就化為轉機.
其實人生很長.
現在是一個轉機.
但我們知道.
繼續看著但以理書下去.
我們就知道這個轉機.
將來亦是重重危機的舞台.
危機會不斷地出現.
但我們今天看到的是.
危機能夠轉為轉機的最重要關鍵.
就是專心一意.
信靠上帝.
將上帝彰顯出來.
所以弟兄姊妹.
我勉勵大家.
當我們面對生命的危機或難題的時候.
我們千萬不要像尼泊爾格薩那樣.
去倚仗自己有的權力.
亦不要像那些外邦的哲士.
去求一些神明的幫助.
我們必須像但以理那樣.
單單依靠上帝.
他很清楚他的上帝是智慧和能力的源頭.
而且只要懷著謙卑憐憫.
上帝必會憐憫.
賜他智慧能力.
在面對艱難的危機裡面去見證神.

$^{801}$我先生Paul曾經在港交所工作.
在他工作的生涯裡.
曾經經歷過一個很大的危機.
在2002年裡面.
他當時負責的一個部門.
會公佈一些上市公司的股價.
有一次他的下屬.
將一個當時的大象股的股價的少數點.
而退了一個位.
結果因為那個是大象股的原股.
令到股市的波幅產生了一些波幅.
因為他的下屬出錯.
他的上司問責.
Paul認為雖然是他的下屬直接出錯.
但他覺得他應該要承擔責任.
於是他就向老闆說.
其實我需要承擔這個責任.
當日有不同的報章打電話來.
他們公司的傳訊部.
去查詢這件事.
當時心裡面真的有些事情擔心.
但他說因為他相信.
他做的是上帝要他做的事情.
所以他就將他的前途交託給上帝.
如果因為這件事他要被問責.
或者要承受傳媒窮追不捨的壓力.
甚至要離職的話.
他仍然堅信這是上帝為他預備最好的.
必然會帶領他走更好的路.
我還很記得.
因為這個緣故.
我第二天買了很多份報紙.
因為我擔心他會因為這件事.
要承受極大極大的壓力.
最後在這麼多的報紙當中.
我只找到一份.
裡面有一個小小的篇幅.
將這件事和他的名字登出來.
其實我們覺得很奇怪.
因為其實都是一個頗大的事情.

$^{841}$但上帝真的很奇妙.
因為當日有一個財經界更大的事發生了.
大家未必熟悉.
但就是因為當時有一個先古事件的引發.
以至財經版大幅報道.
因而他這件事就被淡忘.
這件事就被淡化.
亦都不需要被追查.
或者去負甚麼責任.
而且他都一直在工作崗位裡.
服侍到最後.
直至到他讀神學.
所以弟兄姊妹.
其實在危機裡.
我們絕對可以化為轉機.
關鍵是我們能夠依靠我們所信的上帝.
讓我們一起去祈禱.
親愛的主,我們獻上感恩.
我們知道,我們亦都親身經歷.
你是掌管萬有的主.
世上的君王,領袖都在你掌管之中.
上帝,我們更加知道.
你愛你的子民.
你會賜我們智慧和能力去面對.
上帝,我們生活在這個世代.
的而且確有很多危機湧過來.
但我們卻深信.
只要我們謙卑,全心去依靠.
必從你而來得著智慧去面對.
並且能夠在不同的環境中去見證你.
我們這樣仰望,交託.
是奉主耶穌基督得勝,明智祈求.
阿們.
面對很多事情.
我們都會面對很多擔心,憂慮.
但我們今天回應詩歌.
「天父必看顧你」.
讓我們看到我們不需要害怕.
我們在失望的中.
仍然可以找到盼望.

$^{881}$因為天父必看顧你.
\newpage



\section{ }
\label{sec:PTdjTBD7nsw}
\textbf{2020年9月18日「看不見的祝福」佈道會}
\newline
\newline
連結: \href{https://youtube.com/watch?v=PTdjTBD7nsw}{\texttt{https://youtube.com/watch?v=PTdjTBD7nsw}} ~~~~ 語音日期: 2020-09-19
\newline
\newline
\hyperref[sec:oCGsFJy1z54]{\small{< < < PREV SERMON < < <}}
~
\hyperref[sec:index_chronic]{\small{[返順時目]}}
~
\hyperref[sec:_Ymdcnu_z5c]{\small{> > > NEXT SERMON > > >}}
\newline
\newline
$^{1}$各位朋友大家好.
很高興今晚能夠藉著網絡.
藉著電腦或者手機平台跟大家聊聊天.
剛才我們聽了劉牧師講他的見證.
一個高低起跌.
很戲劇性的故事.
其實是很豐富.
很精彩的.
或者我們沒有經過大的考驗磨難.
可能我們的故事沒有那麼大起大跌.
沒有那麼戲劇性.
不過其實我們的故事都同樣精彩.
不過我們比較少為自己過去的日子.
做一些這樣的創作.
包在一起串連一個大的故事.
做一個療痕.
今天因為劉牧師的見證.
所以我會順著大會的主題.
漢不見日祝福.
來跟大家分享.
對我來說這個題目其實有點特別.
我試著跟這個題目去聊.
有點像寫作文的感覺.
我們先看看漢見這個字.
這個世代其實是一個超級相信眼睛.
依賴眼睛的世界.
是資訊爆炸的時代.
什麼資訊什麼主題.
你都可以在網絡找到.
你想炒碟麵.
有十條youtube教你.
有什麼秘訣可以炒得好.
什麼主題你都找到.
並且在任何角落.
你會發現都有一些植入式的廣告.
有人刻意要傳遞訊息給你.
我們接收了太多不同的資訊.
不同的言論.
言人人樹真真假假.
所以我們慢慢學乖了.

$^{41}$我們拒絕輕易相信別人的意見.
我們要求自己親眼看.
親自查證.
聽和讀其實不可靠的.
看得見的才是真的.
這種相信自己眼睛的認定.
很多時候是對的.
現在流行買海外物業.
如果你沒時間親自去當地看.
你緊緊透過香港的代理.
或者看video.
做投資決定.
你被騙的機會是七成以上.
如果你想藉著網上的社交平台去找朋友.
你認識了一個三十多歲.
又高大又帥又溫柔的單身醫生.
你希望和他結成伴侶.
你被騙的機會是多少?.
100\%.
沒什麼僥倖的.
這個世界怎會還有損盤.
投資界的巨人巴菲特.
大家認識他吧.
你認識他多過認識我.
有一句名言.
我一直相信我自己的眼睛.
遠勝於其他一切.
I always believe my own eyes.
always far better than the other.
我相信自己的眼睛.
媽媽從小就教我們.
大眼識人.
我們的眼睛是否不會騙我們自己呢?.
我們是否每件事都會看得正確呢?.
也不是吧.
我們時時看錯了.
也時時做錯了決定.
我們叫做 fooling ourselves.
自己騙自己.
不說其他.

$^{81}$你去服裝店.
自己挑選回來的衣服.
很多時候買回家後都沒有拆袋.
連穿的動機都沒有.
最後要收拾衣櫃.
就連同包裝一起丟掉.
我們自己選伴侶.
不是盲婚丫架吧.
不是父母幫我們做決定吧.
不過這一代的離婚率.
是遠遠高於上一代.
我們不常常做對的決定.
我們貴買便宜賣.
入市出貨的時機.
總是拿捏得不對.
我們跟隨我們的感覺走.
一時衝動.
總是盲目的.
我這種人不適宜做投資.
因為我是盲目的.
不過雖然是這樣.
我們總是覺得.
我們自己看.
自己做選擇.
就算選錯了也無所謂.
這是一個強調自我的世代.
我走我的路.
我為我一生負責.
只要我自己做決定.
就算是錯的也是對的.
如果由別人幫我做決定.
對也不好玩.
不要.
這叫做自我決定.
我們很多時候.
有機會跟年輕一代聊天.
勸一個年輕人.
最有效.
最受落的說法是什麼.
你要對自己好.

$^{121}$你要忠於自己.
這些說話他很喜歡聽的.
如果他認同.
你的勸導是他本來應該做的決定.
他就會甘心聽你的.
否則你訴諸什麼傳統.
權威.
或者其他東西.
浪費時間.
根本不理你.
由自我出發.
每個人憑自己的眼.
認識外面的世界.
去理解真相.
去判斷價值.
然後採取行動.
這是我們理解自己的行事模式.
不過真相往往是.
選擇理解.
選擇一些價值.
然後選擇我們自己的行動.
這是一個選擇.
不一定是正確的.
執著自己的看見.
問題會出在哪裡.
當然就是看不見的東西.
每個人的眼睛都有限.
其實你們特別看到.
我的眼睛是超級小的.
這是我一件很遺憾的事.
所以.
我的眼睛很小.
我看到的東西其實不是很多.
並且是非常選擇性.
非常具選擇性.
這個選擇性.
很多時候不是故意的.
而是不經意.
受制於我們自己的喜惡愛憎.
也受制於我們本來有的觀點.

$^{161}$先說喜惡愛憎這一點.
我們喜歡什麼.
我們不喜歡什麼.
就影響我們的視線.
例如.
男和女對事物的看法是很不一樣的.
關注很不一樣.
我和太太在街上走.
她撞了一下.
然後問我.
你有沒有留意.
剛才走過的那個胖女人.
穿什麼衣服.
顏色配搭得這麼醜.
我總是搖頭.
我沒有留意到.
當然我沒有跟太太說.
不漂亮的胖女人.
我怎會留意.
我只留意漂亮的.
又例如.
每次和太太行山.
我太太是超級怕毛毛蟲.
特別是那些咬著線.
從樹上掉下來做笨豬跳.
她一看見就尖叫.
我告訴你.
每一次都是她看見毛毛蟲.
所以久不久就大叫一聲.
我因為完全不怕.
所以我根本看不到.
這個故事教訓我們.
我們喜歡什麼.
或者不喜歡什麼.
都會左右我們的看法.
我們的喜惡愛憎.
會影響我們的視線.
我其中最討厭的一件事.
就是那些不負責任的狗主.
拖狗在街上大小便.

$^{201}$又不清理.
我看見就很生氣.
覺得很礙眼.
很破壞心情.
不過.
就因為我不喜歡這件事.
我告訴你.
我經常都看見.
這裡有一篤.
那裡有一篤.
真的每次出街都可以破壞心情.
又舉個例.
有些年長的人.
不認同年輕一代的.
某些想法和做法.
他總是看見.
年輕人.
睜眼睜鼻在他面前.
做一些他不喜歡的事情.
所以這些老人家.
總是凡事都看不慣.
整天耳激的.
慨嘆一代不如一代.
是不是這樣.
喜惡愛憎會影響我們的看法.
我們的觀點.
也會影響我們的視線.
我們既有的想法.
也有的立場.
那些所謂pre-understanding.
前理解.
前理解是影響我們的理解.
核子摸象.
先入為主.
這個所謂perspectivism.
觀點主義.
我喜歡一個人.
我總是可以找到很多.
讓我喜歡他的理由.
情人眼裡出西施.

$^{241}$我不喜歡一個人.
我總是可以找到很多.
讓我討厭他的理由.
所以.
身處在同一個社會.
面對每一天發生的.
相同的事情.
你會看見.
不同政治立場的人.
總是找到.
支持他們立場的鐵證.
真是鐵證如山.
任何事情發生.
都不會影響他們既有的立場.
不會改的.
只會不停加強他們既有的想法.
你很喜歡現任美國總統的.
你會繼續喜歡他.
不會因為疫病爆發.
不會因為社會騷亂.
不會因為山火蔓延.
而你不喜歡他.
如果你不喜歡他.
你本來就不喜歡了.
也不需要新的理由.
這個世界分成不同的陣營.
沒有因為我們面對一個相同的世界.
相同的處境.
相同的危機.
相同的問題.
而稍為縮窄我們的分歧.
對立的繼續對立.
撕裂的繼續撕裂.
香港人.
我猜我們對這方面是.
最感同身受的,是不是?.
美國.
有線新聞網絡CNN.
的口號是.
所謂Facts first(事實第一).

$^{281}$事實優先.
不過.
我的感受是.
再多的事實.
都不會減少我們.
大家彼此之間的分歧.
當我們從自己的立場.
我們自身的感受.
我們的喜惡愛憎.
我們的利益.
一定佔主導的地位.
這個就很容易會.
讓我們忽略其他人的感受.
甚至忽略其他人的存在.
我們的視野變得很狹窄.
我們所享用的一切.
別人在我們身上所做的事.
我們會認為是理所當然的.
我們關心的是.
如何繼續保有我們的東西.
擴充我們本來有的權利.
實現我們自己的目標.
並且我們會很容易看.
任何和我們爭奪利益.
威脅我們福祉的人是敵人.
甚至我們會看.
所有人都是我們的潛在對手.
據說.
我只是聽據說而已.
某個年輕人.
因為家裡不肯幫他付房租.
租外面的房子.
他就在網上跟他的好朋友.
討論如何控告他的父母.
這些很前衛的想法.
對我這個頭髮都掉了的上一輩.
是覺得很匪夷所思的.
不過.
我不想美化我們這一代.
我必須說真的.

$^{321}$我在年輕的時候.
也是很自我中心的.
我出生在一個很貧窮的家裡.
我八歲就出來打工.
供自己讀書.
在很艱難的環境中成長.
我很早就認定這個世界是不公平的.
沒有人可憐你.
只有你自己能夠幫助自己.
所以我咬緊牙關.
努力讀書.
考入大學.
改變自己的命運.
我覺得我的意志力是挺強的.
我的目光也很專注.
對準要達到的目標.
其他的事情我是完全不理的.
我常常忽略我身邊的人的感受.
真正是目光如豆.
直到我十八歲那一年.
我成為了基督徒.
生命有了很大很大的改變.
我才對我自己以外的事情.
有更多的關注.
從而我告訴你.
我進入了一個完全不一樣的世界.
成為了基督徒之後.
我當然認識了上帝.
祂是我生命的主宰.
祂掌握了我的一生.
祂默默無聲地.
帶領我.
引導我.
當我知道上帝的存在.
我整個世界立即從2D變成3D.
除此之外.
我對身邊的人.
也多了一份覺察和認識.
我知道我不是一個自給自足的人.
上帝為我的生命.

$^{361}$佈置了很多的天使.
幫助我.
祝福我.
我媽媽.
我的家人.
我的老師.
很多對我有期望的人.
沒有他們.
我根本活不到今天.
沒有他們.
我也失去繼續奮鬥下去的動力.
舉一個生活的小例子.
這是我幾年前的一個覺悟.
我是在長洲上班的.
我每天早上五點起床.
我坐六點十五分的小巴.
然後坐七點鐘的船進長洲.
我發覺小巴司機和渡輪的水手.
都比我更早起床.
沒有他們的盡忠職守.
我進不了長洲上班.
我在我的學校是做院長的.
算是最高領導人.
不過我好知道.
如果沒有身邊的人對我的幫助.
我這個院長是做不成的.
沒有followers.
何來leader?.
為此我對我身邊的人.
是充滿感激之情.
習慣聽我說話的人.
都知道我經常說這個爛笑話.
不過這是我油蔥的說話.
我和我太太結婚三十六年.
我到如今仍然覺得.
我是執導.
能夠有她.
將她的幸福壓住在我一生.
這是我很幸運的事.
蘇蕊台灣的歌手.

$^{401}$有一首《走乾躺埋霧》這首歌.
有幾句歌詞.
「沒有天乃有地.
沒有地乃有家.
沒有家乃有你.
沒有你乃有我」.
我告訴你.
我有一次.
跟別人分享這首歌的幾句歌詞.
我兒子聽了就質問我.
他說:老套啊你們說這些.
我好知道我和我兒子.
兩代之間是有不同的觀點.
這是代溝的問題.
不過我必須說.
這是我生命裡的核心信念.
飲水思源.
我確認我自己.
是在上帝的恩典.
和身邊很多人的恩情的包圍之下.
我才可以存活到今天.
所以我的心總是充滿著感激之情.
我的子女都可以做見證.
和清潔工人.
我在很多好人的包圍之下生活.
我覺得自己超級幸福.
有些事情我看得見.
有些事情我看不見.
但就在我所活的小世界裡面.
我知道也不僅是.
我看到的那一點點.
我看不到的東西遠遠比我看到的多.
看不到的也遠遠比我看得到的更重要.
例如你坐飛機去外地.
我上到飛機就閉目養神.
我看不到駕駛室裡面的機師怎樣操作.
我看不到他和地面指揮中心怎樣聯繫.
我也看不見整個航空交通的操作.
不過這些看不到的東西.
就是我能夠去到目的地的關鍵.

$^{441}$我們必須要承認.
我們自己的眼界不夠廣闊.
視野範圍很有限.
看事情看不完.
並且上上都是短視的.
我又跟大家說一些老歌.
兩個月之前.
許冠傑在尖沙咀開了一場同舟共濟演唱會.
我不知道你有沒有在電視上.
或者在手機裡面看到.
他在演唱會裡面唱了一首古老的名曲.
《沉默是金》.
寫給張國榮的.
其中的歌詞是這樣說的.
冥冥中都早注定你富或貧.
是錯永不對.
真永是真.
任你怎說堅守我本分.
始終相信沉默是金.
我對這兩句歌詞是非常共鳴的.
跟大家分享一下.
這兩句歌詞其實是反映兩個核心的信念.
第一.
有很多事情其實不由我們控制.
剛才劉牧師說人生沒有什麼把握.
有很多事情我們是沒有選擇權的.
你在哪裡出生時空.
你的種族你的父母.
他們是聰明的是笨的.
高矮肥瘦的基因.
你成長的經歷一生很多的際遇.
我不管你信不信宗教.
你都要承認有很多事情你是控制不了的.
例如說這兩三年畢業的大學生.
剛剛碰到環球經濟衰退.
就真的不是很幸運.
有些行業不請人.
你就入不了行.
可能一生都入不了行.
或者你在今年年初開了一間飲食店.

$^{481}$企圖大展鴻圖.
不過都不用我說.
後果一定堪疑.
限聚令禁堂食.
基本上你一極都沒有用.
聖經的說法是.
很多事情是在乎當時的機會.
時機.
中國人說一命二運三風水.
當然我們都不喜歡命運這個說法.
我們都很想反抗命運.
我們不甘心受命運波動.
做唐吉訶德都比命運波動好.
這樣的勇氣我是很佩服的.
不過當我們反抗失敗.
我們仍然忍不住就會說.
時也命也.
天亡我也.
這就是命運.
非戰之罪.
各位朋友.
我不知道你是不是相信上帝.
這個世界很少人是完全的無神論者.
多數都不知道上帝是誰.
也沒有動機去認識真正的上帝.
我希望你們承認.
當你碰到生老病死.
天災人禍.
人的能力真的很有限.
一場疫病.
全球三千萬人確診.
死亡人數差不多一百萬.
誇不誇張.
美國西岸一場山火.
只是中間的俄羅斯.
一周已經有一半人口.
超過五十萬人要被迫離開家園.
很多人奮鬥半生.
什麼都沒有了.
在風和日麗的日子.

$^{521}$我們很難有這種無助無能.
有限的自覺.
我們總是也文也武.
覺得自己樣樣行.
是不是?.
唯一是遇上一些critical moment(危險時刻).
我們才要正視這個事實.
學習謙卑下來.
人是人.
人不是上帝.
人需要上帝.
可能有人包括你.
就會這樣問.
上帝或命運.
是否人在遭遇生命盡頭的時候的一個解說.
很為自己的失敗.
去尋求開脫的藉口.
責任.
是不是?.
不用什麼.
是不是這樣呢?.
有人更加認為.
只有失敗者才需要宗教.
相信宗教的.
總是一些宿命論者.
是不是?.
我不同意這個說法.
完全不同意.
這關涉到我要說的第二點.
也是沉默至今.
這兩句歌詞要傳遞的第二個信念.
我們不僅僅是在.
面對無法抗拒命運的時候.
承認自身的有限.
冥冥中都早注定你苦或貧.
我要說.
我們在繼續奮鬥.
反抗命運的過程裡.
我們都需要.
抓緊一些能夠讓他們超越現實的信念.

$^{561}$歌詞接著說.
是錯永不對.
真永是真.
任你怎說堅守我本分.
始終相信沉默至今.
我怎麼知道真理存在.
正義存在.
我怎麼知道.
強權不是真理.
我可以為堅持真理的緣故去反抗強權.
我為什麼不承認.
正義是由那些強者說了算.
我為什麼可以確定.
堅持真理是必須.
而是能夠的.
那些暴君不會信上帝.
因為他們宣告自己就是上帝.
可以為所欲為.
我們這班無權無勢的小老百姓.
如果我們不信上帝.
我們除了趴在那些人間的暴君面前之外.
我們還可以做什麼.
三十年前.
面對一場人禍的悲劇.
香港人在極度的悲傷和絕望裡面.
我教會有一個年輕人.
很激憤的.
跟我大喊大叫.
他說牧師.
你不要再跟我說上帝.
這個世界根本沒有正義.
沒有真理.
沒有上帝.
我當然明白他激動的心情.
但我沒有安慰他.
我只跟他說.
兄弟.
你要知道你在說什麼.
如果這個世界真的沒有真理.
沒有上帝.

$^{601}$那你還激動來做什麼.
你叫什麼.
是的.
這個世界是弱弱強弛的世界.
誰掌權誰有軍隊.
誰就代表真理.
誰就是上帝.
世界就是這樣.
你不相信嗎.
生氣嗎.
你不怕嗎.
對不對.
為什麼我們可以不甘心.
不需要接受現實.
為什麼我們拒絕承認強權就是真理.
因為我們真的相信真理存在.
為什麼我們拒絕向現實低頭.
因為我們相信正義會最終獲勝.
如果我們在眼見的現實之外.
再不存在任何價值任何理想.
那憑什麼我們可以拒絕現實的壟斷性.
我們仍然敢和現實格格不入.
他今天的現實很遙遠.
很遙遠的理想.
你為什麼還敢追求他.
各位朋友.
你信不信天地有正氣.
很多大智學家其實都說過.
如果沒有上帝的話.
就做什麼都可以.
我們中國人也曾經經歷過.
很多次無法無天的歲月.
真的.
如果沒有上帝.
超越現實的真理和公義.
如何確立.
不要不說.
這個世界是殺人放火金腰帶的.
損害個人利益的愛和憐憫.
如何堅持.

$^{641}$愛人.
幫人.
人不為己.
天誅地滅.
距離現實遙不可及的理想.
如何堅持.
空中留角而已.
不會成功的.
你放棄吧.
人做官我做官.
貪污就月盡千萬.
不貪就注定清貧.
如果沒有上帝.
只有此生沒有來世.
為什麼我不搏一搏.
就算被人抓了.
打靶.
就認命吧.
劉曉波曾經說過.
中人的悲劇.
是沒有上帝的悲劇.
我.
是複修哲學的.
我看過很多哲學書.
我看到不少學者.
拒絕承認.
沒有宗教論理道德是沒有基礎的.
他們努力尋求無神論的道德基礎.
我要很抱歉地說.
我今天也看不到一個.
很具說服力的理論.
上帝存不存在?.
聖經說.
上帝是個靈.
上帝是我們憑肉眼看不見.
看不見.
但憑著信心.
我們可以清清楚楚地.
經歷祂的存在.
宇宙的浩瀚.

$^{681}$人心中的良知.
都告訴我們上帝存在.
仰望星空.
聆聽我們的內心.
你會發現.
上帝就在你身邊.
呼喚著你.
認識上帝.
認識生命的根源.
認識自己的由來.
認識自己所哀憤的.
沒有天那有地.
沒有地那有家.
沒有家那有你.
沒有你那有我.
我是上帝的創造.
我是上帝愛的創造.
因為我信上帝.
相信我永遠不會絕望.
認識我的人都知道.
我是超級樂觀的人.
我總是認定眼前的困難.
不是壓倒性的.
眼前的惡人.
不是所向無敵的.
我更加認定.
天無絕人之路.
這是聖經的教導.
聖經這樣說的.
在哥林多前書十章十三節.
你們所遇見的試探.
無非是人所能夠承受的.
上帝是信實的.
必不叫你們的試探.
過於能夠承受的.
在受試探的時候.
總要為你們開一條出路.
叫你們能夠忍受得住.
你不熟悉聖經.
你不知道我在說什麼.

$^{721}$認識聖經的弟兄姊妹.
會知道.
這是聖經很重要的應許.
我是一個不死的小強.
並且我是一個好樣的小強.
我可以病都病個好樣出來.
死都死個好樣出來.
通常讓我們陷在絕望的原因有兩個.
都是和你的體見有關.
第一.
就是我們體見的東西是選擇性的.
以偏概全.
你經常看到眼前的困難.
看不到身邊已經有的幸福.
是不是.
一個年輕的女孩.
留下因為男朋友撇她.
就決定自殺.
留下為她的死呼天搶地的父母.
留下很多愛她的人心中的遺憾.
我總是會問.
只有一個人不珍惜你.
值得為她死嗎.
不過.
人就是這樣.
我們總是看到自己缺乏的.
沒有的.
看不到我們已經有的.
一直有的.
當我們的眼總是看著自己沒有的東西的時候.
你永遠不會滿足.
你覺得自己是超級貧窮的人.
當你的眼睛看到你有的東西.
你會發覺自己原來是很豐富.
我的太太.
我剛才說過.
她將她一生的幸福壓在我身上.
我看到她.
我就知道.
我是一個很富有的人.

$^{761}$這是第一個.
因為我們看得很片面.
第二個原因.
我們除了看得片面之外.
我們常常也太短視.
太短視.
我們只看到眼前的幾步.
眼前出現困難.
我們就要生要死.
我們看不到這個危機.
或者是未來突破更新.
這是一個很重要的轉折.
天將降大任於我們.
投資界不是常常說的.
大紅市才是做大投資.
賺大錢的機會.
中國人常常說的老話.
福是禍之所倚.
禍是福之所伏.
禍福無常.
有很多東西要看遠一點.
要沉得住氣.
不過除了眼光和膽量之外.
由於我們必須承認.
事實上我們看不到路的盡頭.
堅持下去.
有賴我們的.
除了眼光和膽量之外.
其實是我們的信心和信念.
我們的信念.
你是一個怎樣的人.
你相信什麼.
就影響很大.
我相信我這一生.
是在愛我保護我的上帝手中.
所以我敢繼續向前行.
心存盼望.
因為我信上帝.
所以我不會放棄我固有的價值.
包括自己做人處事的原則.

$^{801}$對人的善念.
信念.
對人有good will.
永遠對人有善意.
我拒絕被現實牽著鼻子走.
在現實中沉淪.
和光同塵.
為兩餐什麼都肯做.
人貪污我貪污.
人說謊我說謊.
我不會向現實或群眾妥協.
人人這樣做.
我就必須要這樣做.
我確認這個世界是上帝掌管的.
他是正義的上帝.
真理的上帝.
所以我認信.
善有善報.
惡有惡報.
正義必會伸張.
真理不會被掩埋.
我只管做對的事.
讓上帝負責一個對的後果.
因為我相信上帝.
所以我不會太心急.
我不會急於在今天.
或今生取得一切的公平.
聖經裡上帝這樣說.
「身冤在我,我必報應」.
沉默至今.
也有這樣的信念說.
是錯就永不對.
錯就是錯,錯不會變成對.
真永是真.
真不會變成假.
謊言說上一萬遍仍然是謊言.
沒有人可以作虛弄假到永遠.
沒有人可以掩蓋真相到永遠.
所以我們不急於為自己做解釋.
婆婆媽媽說來說去.

$^{841}$我們可以不需要解釋.
不竊解釋.
沉默至今.
聖經是清楚告訴我們.
信靠上帝的人必不著急.
今天我們的主題是.
看不見的祝福.
我確認.
我這一生是在上帝的恩典.
和身邊的人的恩情包圍之下.
才能夠活到今天.
我們常常看不到上帝.
也沒有留意到身邊的人.
但我知道我永遠不會孤單.
常被寵愛.
我確認我眼前有很多困難.
有很多我看不透的東西.
不過上帝仍然掌管一切.
他不僅僅掌管我一生.
也掌管著宇宙人生的發展.
世界看起來無疑是惡人當道.
但真正的.
是在上帝的手裡.
面對眼前的香港和世界.
你是不是覺得充滿危機和挑戰.
我們慣常的生活秩序被擾亂.
有很多我們從前視為理所當然的東西.
今天一下子失去了.
面前我們確實會面對很多的失去.
我們的工作,我們的財產,我們的自信.
甚至我們的關係.
都可能會失去.
面前我們都需要做很多重大的決定.
轉不轉工?.
執不執手眼前這盤生意?.
送不送子女出國讀書?.
移不移民?.
這些都是很重大的決定.
都要由你和我去做.
面對改變,面對失去.

$^{881}$面對重大的決定.
你會發現你是一個怎樣的人.
你的性格,你的人格,你的執著,你的信念.
你的價值,你的信仰.
統統是關鍵性.
各位朋友.
如果你決定留在香港.
我希望你能夠大聲唱許冠文的《貼塔靈魂》.
彈到灘外點點輪光.
豈能及魚登在香港.
河水多見浮多求.
且唱一曲歸途上.
認定香港是我家.
是我永遠會押注在這裡的地方.
如果你決定移民.
我希望你能夠堅決地去唱林子祥的《抉擇》.
要將我根和苗再種新土樣.
就算受挫折也當平常.
發揮抉擇力量.
再起我新門牆勝過舊家鄉.
我再說.
面對改變,面對失去,面對重大的決定.
你會發現你是一個怎樣的人.
你的性格,你的人格,你的執著,你的信念.
你的價值,你的信仰.
統統是關鍵性.
你相信什麼呢?.
你對面前有什麼盼望呢?.
最重要的是.
我希望無論你身處何方.
無論你面對什麼處境.
都有耶穌基督在你身邊.
成為你的安慰和保護.
也成為你的信念和價值最強力的支撐.
各位親愛的朋友.
今晚時間不多.
我沒辦法和大家完備地將基督信仰說出來.
不過我認識很多還在信仰抉擇中的人.
其實他們或多或少都聽過基督教很多道理.
如果你不是身邊有一個基督徒朋友或家人.

$^{921}$帶著你去看這個節目.
你是自己不小心迷倒進來的呢?.
我誠懇地邀請你.
稍後填寫回應表.
最重要的是找你身邊的基督徒.
或者你填寫回應表後.
我們會有人直接聯絡你.
帶你回教會.
我希望你在教會可以繼續認識這個信仰.
並且知道這是你人生最重要的倚靠.
我真的很誠懇向你介紹.
這個改變我一生.
改變林牧師一生的基督信仰.
我希望在今天邀請每一個對信仰有一點點的興趣.
又或者你剛才被觸動到的.
你可不可以和我一起閉上眼睛.
我們去做一個祈禱.
如果你身邊有基督徒的朋友.
請你的親友一起祈禱.
祈禱是向上帝講話.
你聽我說的不是完全真的.
我希望你在祈禱裡面.
你接觸你經歷這位上帝.
他親自觸摸你.
向你說話.
如果你覺得你可以的話.
在祈禱的時候.
我講一句.
你又講一句.
你認同的話.
或者你想嘗試的話.
我誠懇的希望你和我一起祈禱.
好嗎?.
我們一起祈禱吧.
親愛的上帝.
親愛的耶穌基督.
今天我聽到你是宇宙間唯一的神.
你是愛我.
為我捨己犧牲的上帝.
如果這是真的話.

$^{961}$求你向我顯現.
進入我的生命裡面.
讓我能夠認識你.
親愛的上帝.
過去我一直走自己的路.
我以為我自己很行.
今天我承認我的生命其實很脆弱.
我心裡也常常有黑暗的力量.
我沒有足夠的勇氣.
信念去面對明天.
所以我求你成為我的上帝.
幫助我.
引領我.
讓我能夠活出你所應許的豐盛的生命.
你和我同在.
我希望在未來的日子.
我可以更加認識你.
各位親愛的朋友.
如果你覺得你準備好了的話.
我更加期望你和我說以下的話.
親愛的耶穌.
我知道你愛我.
為我捨己.
我也知道我自己的軟弱.
我的無能.
所以今天我願意承認.
你是我生命的救主.
你是我的上帝.
我願意將我的生命交託給你.
跟隨你.
認識你更多.
求你進入我的生命.
拯救我.
改變我.
我宣告耶穌基督是我的救主.
是我生命的主.
奉耶穌基督的聖名求.
阿們.
再次祝福每一位.
\newpage



\section{但以理書 3:16-30}
\label{sec:_Ymdcnu_z5c}
\textbf{2020年9月20日 崇拜直播｜但以理書3\_16-30}
\newline
\newline
連結: \href{https://youtube.com/watch?v=_Ymdcnu-z5c}{\texttt{https://youtube.com/watch?v=\_Ymdcnu-z5c}} ~~~~ 語音日期: 2020-09-21
\newline
\newline
\hyperref[sec:PTdjTBD7nsw]{\small{< < < PREV SERMON < < <}}
~
\hyperref[sec:index_chronic]{\small{[返順時目]}}
~
\hyperref[sec:jdoVqa_V4mw]{\small{> > > NEXT SERMON > > >}}
\newline
\newline
但以理書 3:16-30
\newline
\begin{longtable}{cl}
\hline
\hline
章節 & 經文 (和合本修訂版)\\
\hline
3:16 & \begin{tabularx}{0.7\textwidth}{X} 沙得拉、米煞、亞伯尼歌對王說：「尼布甲尼撒啊，這件事我們不必回答你， \end{tabularx} \\ \\ \relax
3:17 & \begin{tabularx}{0.7\textwidth}{X} 即便如此，我們所事奉的神能將我們從烈火的窯中救出來。王啊，他必救我們脫離你的手； \end{tabularx} \\ \\ \relax
3:18 & \begin{tabularx}{0.7\textwidth}{X} 即或不然，王啊，你當知道，我們絕不事奉你的神明，也不拜你所立的金像。」 \end{tabularx} \\ \\ \relax
3:19 & \begin{tabularx}{0.7\textwidth}{X} 當時，尼布甲尼撒怒氣填胸，向沙得拉、米煞、亞伯尼歌變了臉色，命令把窯燒熱，比平常熱七倍； \end{tabularx} \\ \\ \relax
3:20 & \begin{tabularx}{0.7\textwidth}{X} 又命令他軍中的幾個壯士，把沙得拉、米煞、亞伯尼歌捆起來，扔在烈火的窯中。 \end{tabularx} \\ \\ \relax
3:21 & \begin{tabularx}{0.7\textwidth}{X} 這三人穿著內袍、外衣、頭巾和其他的衣服，被捆起來扔在烈火的窯中。 \end{tabularx} \\ \\ \relax
3:22 & \begin{tabularx}{0.7\textwidth}{X} 因為王的命令緊急，窯又非常熱，那抬沙得拉、米煞、亞伯尼歌的人都被火焰燒死。 \end{tabularx} \\ \\ \relax
3:23 & \begin{tabularx}{0.7\textwidth}{X} 但是這三個人，沙得拉、米煞、亞伯尼歌被捆綁著，掉進烈火的窯中。 \end{tabularx} \\ \\ \relax
3:24 & \begin{tabularx}{0.7\textwidth}{X} 那時，尼布甲尼撒王驚奇，急忙站起來，對謀士說：「我們捆起來扔在火裡的不是三個人嗎？」他們回答王說：「王啊，是的。」 \end{tabularx} \\ \\ \relax
3:25 & \begin{tabularx}{0.7\textwidth}{X} 王說：「看哪，我看見有四個人，並沒有捆綁，在火中行走，也沒有受傷；那第四個的相貌好像神明的兒子。」 \end{tabularx} \\ \\ \relax
3:26 & \begin{tabularx}{0.7\textwidth}{X} 於是尼布甲尼撒靠近烈火窯門，說：「至高神的僕人沙得拉、米煞、亞伯尼歌，出來，來吧！」沙得拉、米煞、亞伯尼歌就從火中出來。 \end{tabularx} \\ \\ \relax
3:27 & \begin{tabularx}{0.7\textwidth}{X} 那些總督、欽差、省長和王的謀士一同聚集來看這三個人，見火不能傷他們的身體，頭髮沒有燒焦，衣裳也沒有變色，都沒有火燒過的氣味。 \end{tabularx} \\ \\ \relax
3:28 & \begin{tabularx}{0.7\textwidth}{X} 尼布甲尼撒說：「沙得拉、米煞、亞伯尼歌的神是應當稱頌的！他差遣使者救護倚靠他的僕人，他們不遵王的命令，甚至捨身，在他們神以外不肯事奉敬拜別神。 \end{tabularx} \\ \\ \relax
3:29 & \begin{tabularx}{0.7\textwidth}{X} 現在我降旨，無論何方、何國、何族，凡有人毀謗沙得拉、米煞、亞伯尼歌的神，他必被凌遲，他的房屋必成糞堆，因為沒有別神能像這樣施行拯救。」 \end{tabularx} \\ \\ \relax
3:30 & \begin{tabularx}{0.7\textwidth}{X} 那時王在巴比倫省使沙得拉、米煞、亞伯尼歌高升。 \end{tabularx} \\ \\
[1ex]
\hline
\hline
\end{longtable}
$^{1}$讓我們一同以聆聽聖道去敬拜我們的上帝.
今日要重讀的經文是.
但以理書三章十六至三十節.
撒德拉,米薩,阿伯尼哥對王說.
「利布格尼薩,這件事我們不必回答你.
即便如此,我們所事奉的神.
能將我們從烈火的窯中救出來.
王啊,祂也必救我們脫離你的手」.
即或不言.
「王啊,你當知道我們缺不事奉你的神.
也不敬拜你所立的金像」.
當時利布格尼薩怒氣填空.
向撒德拉,米薩,阿伯尼哥變了面色.
吩咐人把窯燒熱.
比尋常更加七倍.
又吩咐他軍中的幾個壯士.
將撒德拉,米薩,阿伯尼哥捆起來.
瓮在烈火的窯中.
這三人穿著褲子,內袍,外衣和別的衣服.
被捆起來,瓮在烈火的窯中.
因為王命緊急,窯又滲熱.
那台撒德拉,米薩,阿伯尼哥的人都被火焰燒死.
撒德拉,米薩,阿伯尼哥這三個人.
都被捆著,落在烈火的窯中.
那時利布格尼撒王驚奇.
急忙起來對謀士說.
「我們捆起來瓮在火裡的那不是三個人嗎?」.
他們回答王說「王啊,是」.
王說「看啊,我見有四個人.
並沒有捆綁在火中遊行,也沒有受傷.
那第四個的相貌好像臣子」.
於是利布格尼撒就近烈火窯門說.
「至高臣的僕人撒德拉,米薩,阿伯尼哥出來.
賞這女來吧!」.
撒德拉,米薩,阿伯尼哥就從火中出來了.
那些總督,陰差,巡撫和王子.
和那些不同的謀士.
一同聚集看著三個人.
見火無力傷他們的身體.
頭髮也沒有燒焦.

$^{41}$衣裳也沒有變色.
並沒有火燎的氣味.
利布格尼撒說.
撒德拉,米薩,阿伯尼哥的神是應當稱重的.
神差顯示者救護已靠他的僕人.
他們不遵皇命.
捨去己身.
在他們神以外不肯事奉敬拜別神.
現在我幹旨.
無論何方,何國,何族的人.
幫讀撒德拉,米薩,阿伯尼哥之神的.
必被凌遲.
他的房屋必成糞堆.
因為沒有別神能這樣施行拯救.
那時.
王在巴比倫省高聲鳥撒德拉,米薩,阿伯尼哥.
今日有我們的主音牧師陳明柱牧師為我們靜悼.
今日宣講的題目是.
「隔逢中的中瀑,織或不現,我們仍然」.
弟兄姊妹平安.
今日我們能夠繼續看《但義理書》.
這個系列特別講到.
今日的世界或我們身處的環境當中.
有時對我們基督教信仰或對我們每一個所認信的.
我們的生命,我們的身份.
有時外界可能有些質疑.
我們希望透過這個系列.
重新幫助弟兄姊妹.
堅定我們的心智在這個世代裡.
忠心成為上帝的見證.
有一個叫做存在主義的心理學家叫做Victor Franco.
他說到一個人在一個極端的環境之下.
在一個對自己的生命面對重重威脅的一種狀態之下.
什麼人可以生存下來呢?.
有什麼人可以繼續在一些這麼困難的處境裡面.
可以面對這些艱難繼續走在面前的路.
Victor Franco說有三種人.
看看這三種人是不是你其中一種.
第一種就是有一個叫做創造性的價值.
就是說當人在一個極端環境裡面.

$^{81}$他找到自己能夠在那個環境裡面做一些貢獻.
可能是教他自己所長的教材.
又或者是他能夠為那個環境裡面.
添上一些藝術的氛圍.
或者自己盡一點的綿力去改變更新.
或者創造一些新的條件.
如果能夠有一個人有一種創造的能力.
這些人在一些極端環境之下.
他都依然可以繼續的來生存的.
這是第一種.
第二種就是一個人如果他很知道他身邊的朋友.
他的親情.
他知道自己是好被深愛.
他有一種很強力的關係的能力.
即使面對一個多麼困難的處境也好.
因為他知道有人愛他.
或者他知道在他生命當中.
有些人值得他繼續生存下去.
用他的愛來愛身邊他所愛的人的時候.
在這種關係性的價值底下.
這個人可以繼續的有力量的生存.
去到第三種了.
可能他又沒有創造的能力.
甚至乎他處於一種隔絕的狀態.
他沒有辦法經歷到關係的融洽的時候.
第三種也是最高層次的.
那個叫做信念的價值.
在他生命當中他單單凝信.
一個他生命裡面所抱持的信念.
憑著這個信念就可以幫他支撐走過.
生命裡面有很多很困難的處境.
今天這個世界裡面很多人都很特別.
吹捧第三種的那種信念的價值.
有些人會覺得這個信念價值.
如果將人的意志成為一個超人.
有些哲學家尼采就在說.
他有一個強力的意志.
當人擁抱這種強力的意志的時候.
人就可以生存了.
不過在歷史裡面他驗證了給我們知道.

$^{121}$這種信念的價值不是來自於人.
單單很強力的那種意志能力.
而是將我們的意志放在一個永恆的主身上.
說這番話的存在主義者.
Friedrich Frenkel是一個有神論的猶太學者.
他深信一個人如果能夠.
認信自己在天上的主的時候.
當我們生命裡面經歷多少困難也好.
我們都可以憑著這個認信.
我們可以走過生命裡面很多很不容易走過的路.
今天我們要看的《但義理書》第三章.
其實是說一種但義理三個的朋友.
他們如何憑藉他們的信念.
認信這位天上的主.
在一個很困難很極端的環境之下.
他不單止生存下來.
而且在這個處境當中好像一台戲一樣.
讓巴比倫的尼泊爾·格尼薩和眾神子.
看到上帝的大能.
希望今天我們看這段經文的時候.
可以提振我們.
讓我們能夠更加明白上帝的心意.
也讓我們更加有勇氣.
有信心走面前的路.
在我們只是一題看聖經之前.
我們一起祈禱.
求上帝的靈與我們每一個同在.
我們從心祈禱.
親愛主願你幫助我們.
讓我們能夠在你話語當中.
我們看到你給我們每一個生命的亮光.
驅除我們心靈裡面的黑暗.
幫助我們在你話語當中.
我們找到你要對我們所說的訊息.
我們恭敬將來每一個弟兄姊妹交託給你.
幫助在家裡面看著螢光幕的.
你的聖靈與我們每一個同在.
讓我們不受地域所限.
讓我們能夠領受你的話語.
有幫助宣講的講得清楚.

$^{161}$讓你的話語中心被宣講出來.
我們恭敬求你與我們每一個同在.
多謝主垂聽祈禱.
奉基督耶穌得勝的聖名.
Amen.
之前我們在上兩個星期.
我們一起看《但以理書》.
裡面的內容我不重複了.
如果你們想看我們的YouTube頻道.
可以看上兩次的講道.
不過這裡我想講兩個伏線.
在上個星期第二章裡面講到.
尼泊爾加尼沙發了一個怪夢.
在夢裡面.
即是彈以理就幫他解夢.
尼泊爾加尼沙發了這個夢之後.
他怎樣回應這個夢的訊息呢?.
其實在夢裡面.
彈以理說.
王啊你是諸王的王.
天上的神已經將國度權並能力尊榮都賜給你.
凡世人所住的地和走獸.
並天空的飛鳥.
他都交付在你的手.
你就是那金頭.
在這裡.
彈以理.
在講這個夢.
他講到尼泊爾加尼沙.
他夢中所見的一個很奇怪的像.
裡面那個金的.
金柱造的那個頭.
尼泊爾加尼沙聽完之後.
在第二章沒有特別回應.
不過在第三章裡面.
就講到尼泊爾加尼沙就開始.
鑄造一個.
六十爪高.
爪闊的一個金像.
似乎在第三章裡面.

$^{201}$回應著這條伏線.
尼泊爾加尼沙知道自己.
從夢中裡面指示出來.
他是那個金頭.
這個像就好像.
可能是他自己.
又或者是他代借一個新的偶像.
來幫助他統治.
當時整個福原很大的巴比倫王國.
透過這種.
這樣的表達.
讓整個福原廣大的巴比倫的子民.
很多的種族.
都要臣服在.
他的管治之下.
第二條伏線就是.
我們一直都看到但義理書.
雖然我們知道但義理是其中一個.
很重要的人物.
不過聖經裡面一而再再而三.
他都提及.
但義理不是孤身作戰的.
他還有三個好朋友.
米薩.
另外.
阿伯尼哥.
還有沙特拉.
應該是沙特拉米薩.
阿伯尼哥會比較好讀.
這三個朋友.
在第二章裡面就提到但義理怎麼做.
但是在第49節.
二章49節裡面.
他提到但義理求王.
王就派沙特拉米薩阿伯尼哥.
管理巴比倫省的事務.
只是但義理常在朝中.
事立.
大概.
但義理高升.

$^{241}$他都帶了幾個兄弟.
幫他在巴比倫裡面.
有一官半職.
不過經文裡面就提到.
這三個友人.
在第三章裡面就要面對.
一個很重大的考驗.
就是.
尼泊爾夾尼薩駐造金像.
這三個友人.
他怎麼去面對呢.
直到第三章第四至六節裡面.
他就提到.
尼泊爾夾尼薩.
特別的.
在經文裡面交代.
他就駐造了一個這樣的金像.
剛才也有提到.
這個金像.
基本出土文物是怎麼樣呢.
已經無從考究沒有人知道.
聖經學者也沒有畫一.
的結論究竟這個金像.
是尼泊爾夾尼薩.
王自己.
是他製造一個神出來.
或者將他以前巴比倫敬拜的神.
成為一個偶像.
來通國.
來敬拜.
不過姑且不論這個.
神明是一個怎樣的神.
但是他的手段.
是一個很強硬的手段.
當這個像.
駐好了之後又要做一個.
開光儀式.
他就要用很.
壯闊的音樂.
琴,鈴,琵琶,蝨.

$^{281}$和其他樂器.
當聽到.
聽到這些樂器一吹奏的時候.
所有人.
就要面對著.
這個金像.
俯伏下拜.
如果不下拜這個金像.
的時候.
當時王就下了一道很嚴格的命令.
就是要將那些人.
推進去火窯當中.
來被火鑑.
燒死.
這是一個整個通行的.
一個預令.
沒有人夠膽.
違背這個命令.
更重要的是.
跪拜在這個金像的.
身體動作.
當每一個人都趴在那裡.
俯伏下去去拜這個金像的時候.
你不拜.
你站在那裡.
就會像鶴立雞群.
全部人都會看得到.
在過去的第一章.
第二章我們看到.
但以理和他三個朋友.
在一個對他信仰.
很不友善的巴比倫帝國裡面.
他們都忍辱負重.
他們捐農捐拿.
希望找到一個空間.
不單止自己生存.
更加重要的就是在這些.
的夾縫當中.
他希望繼續用他們自己猶太人.
被擄的子民.

$^{321}$雖然他要面對上帝的刑罰.
但他仍然要繼續承擔上帝.
給他們的召命.
過去的日子.
頭兩章講到他們.
真的很低調.
很忍辱負重.
不過去到第三章.
那種對立面就變得白熱化.
因為.
當每個人都俯伏敬拜的時候.
你不去敬拜.
或者你藉機蹲在那裡.
對他們的信仰.
都是一個很重要的違背.
這時候就會有很多人看到.
他面對的.
對神.
不表忠的一種做法.
所以.
很短時間之內.
撒德拉米薩.
阿伯尼戈.
這三個友人.
在不敬拜這個像之下.
就開始.
面對一個很重要的衝突.
再加上.
其實這幾個.
猶太的年輕人.
在朝廷裡面.
很快就冒升起來.
當中有一些宮廷的鬥爭.
有一些迦勒底人.
看到他們不敬拜.
乘機就借這個位置.
來到皇上面前就告余壯.
所以在皇上面前.
這些.
迦勒底人就跟皇上說.

$^{361}$跟尼泊爾夾尼撒王說.
這幾個人又不聽.
你的命令.
不跪拜你所鑄造的.
金像.
這件事是真的這樣.
不過他將這些事告訴王上.
目的就好像製造了.
一些宮廷裡面的鬥爭.
當王上.
聽到這番說話的時候.
在聖經裡面就說到尼泊爾夾尼撒.
就勃然大怒.
很氣沖沖.
就要召見這三個友人.
在第十四至第十七節.
可以這樣說.
或者第十八節.
在這段經文裡面.
不是特別交代細節.
而是在說王和這三個友人的對話.
這三個友人說的說話.
不是很多.
不過他的表達.
是非常精彩.
甚至乎是一個很強烈.
表達他們對耶和華上帝認信的一種表現.
第十四至第十七節.
尼泊爾夾尼撒問他們說.
撒德拉,米薩亞伯尼哥.
你們不事奉我的神.
也不敬拜我所立的金像.
是故意的嗎?.
你們再聽見.
國,滴,琵琶,琴,塞,聲.
和各樣的樂器聲音.
若俯伏敬拜.
我所做的像.
卻還可以.
若不敬拜.

$^{401}$必立時,瀟灑烈火的搖鐘.
有何神能救你脫離我手呢?.
撒德拉,米薩亞伯尼哥對王說.
尼泊爾夾尼撒.
這件事我們不必答你.
即便如此.
我們所事奉的神.
能將我從烈火的搖鐘救出來.
王啊!.
他也必救我們脫離你的手.
在這裡.
首先說到.
王氣沖沖.
叫這三位友人.
去到自己的庚前.
某程度上是一種對質.
王好像給了他最後通牒.
你不敬拜我所鑄造出來的像.
是否故意的?.
是這樣的.
第二天你再聽到.
這些音樂這麼壯闊.
這麼偉大的音樂的時候.
你再去.
去拜我所鑄造出來的像.
那你就可以不用死了.
不過.
如果你違背了這個皇令.
就要將你.
扔進烈火的搖鐘裡.
就要將你扔進烈火的搖鐘裡.
來燒死.
當你真的這樣違抗的時候.
有什麼神可以幫到你?.
尼泊爾說這番話.
其實他講到.
他自己的強權比神更加大力.
他用他自己的方法.
他自己的方法.
用他自己的強權.

$^{441}$來逼使所有人.
包括這幾個.
已經剛剛冒升上來的.
這三個少年的官.
他給最後通牒給他.
不過也給他一個很重要的警告.
當他們不去下拜這個像的時候.
他就面對著這個酷刑.
在尼泊爾的眼中.
其實他的權力大到.
已經是目空一切.
甚至他當自己是神一樣.
所以他講到.
如果你不下拜的時候.
神也幫不到你.
在他的口裡.
大概他對於這一群.
覺得很迂腐他們的信仰的.
這群猶太人.
覺得他們在發瘋.
覺得他們拜一下.
實際一點就搞定.
何必這麼硬頸.
還要面對著火的刑罰.
他用這種強權.
某程度上.
就要屈服這一群.
外邦或猶太人.
一臣信念的.
一臣信念的一群子民.
這三個有人面對著尼泊爾.
這種軟硬的提問.
他們有什麼反應呢?.
他們三個回答得很精彩.
尼泊爾.
這件事.
我們不必回答你.
這件什麼事呢?.
大概是在說.
哪個神可以救你?.

$^{481}$又或者你不去.
下拜金像的時候.
有什麼結果?.
這三個有人.
不再需要在他面前.
在王面前做任何的答辯和解釋.
在最白熱化的階段.
他說到.
我們是信一個唯一天地萬物.
創造天地萬物的主.
我們是敬拜這個神.
所以在這裡.
特別來說.
如果你要將我們.
從烈火中的窯裡.
燒了他們.
即便如此.
我們所事奉的神.
能將我們從烈火的窯中.
救出來.
即便如此.
這個字很有意思.
大概是在說.
即使你要將我.
丟到火裡.
即使如此.
我們都相信.
我們所事奉的神.
能將我們從烈火中.
拯救出來.
王啊.
他也必救我們.
脫離你的手.
在這三個.
有人當中他所看的.
那個偶像.
是一個強權的象徵.
但是在這三個.
有人.
撒德拉米薩阿伯尼.

$^{521}$眼中所看到的神.
不是一個強權.
屈服人在他面前.
下拜的強權.
而是一個拯救的主.
他們心裡確信.
他們所信的神.
是一個拯救的神.
即使面對.
王的這種.
這麼嚴格的禦令.
他們都相信.
他們將他們的生命.
交付在他們所事奉的神手裡.
而不是交在尼泊爾格尼薩的手裡.
即便如此.
是其中一個.
很重要.
我們應信的一種表現.
我們稍為在這裡.
繼續思考.
我們今天.
對於我們的神的觀念.
我不知道大家.
如何去理解.
曾經有些時候.
我們有機會到醫院探病.
或者面對著.
探訪弟兄姊妹.
面對生命裡很多的艱難.
我曾經有一次這樣的經驗.
那時候剛好初出茅廬.
做傳道人.
跟著長輩的牧者一起來探訪.
面對著一些.
很艱難處境的人.
我們如何祈禱呢?.
心裡有時候都打定輸數.
心裡.
永遠都帶著絕望的心.

$^{561}$但是在我的言辭.
不想給別人絕望.
有時候會說一些.
祈禱的圖文無伶兩可.
例如醫病.
心裡覺得這個病人.
可能都治不了.
心裡是絕望的.
但為了不要讓弟兄姊妹.
感覺不良好.
就說上帝求你醫治他.
但還要包底.
不過你不醫治他也要給他平安.
或者給他家人.
可以舒緩.
大包圍式的祈禱.
什麼都祈禱.
表面上就好像包得很緊.
滴水不漏.
在完了探訪之後.
那個很資深的牧者.
他跟我一路乘車回教會.
一路說.
Andis,剛才你的祈禱.
很大包圍.
不過你的心裡.
你信不信上帝會醫治他呢?.
我就很坦白說.
其實我不是很相信.
不過我不知道上帝是怎樣做事的.
所以就大包圍.
這個牧者.
如果你不相信.
說得浪浪上口.
說得好像.
很.
提到很大的那種.
決心.
同時你又為自己說.
那個內容就好像做了一些.

$^{601}$大包圍那樣.
其實你的祈禱是一個什麼的祈禱.
這句話.
我一路心裡.
一路沉澱了很久.
去到今天我都在沉澱.
今天我們.
我們陪一些弟兄姊妹同行.
或者我們為別人祈禱.
我們信不信究竟.
我們的神可以幫到他.
或者我們覺得.
我們做大包圍.
怎樣都中的.
總之就說了所有的.
可能性.
表面上就好像一個很典雅的祈禱.
實際上可能我們對上帝的應訊.
我們只不過是.
好像包裝著一些言辭.
讓別人舒服一點.
不是將別人帶到.
在上帝面前.
更加重要的是.
甚至乎有時候.
在我們的信仰運作裡面.
我們會笑一些弟兄姊妹.
笑某些宗派.
我們覺得他們神神化化地這樣想.
覺得他們很強調某一種神蹟.
或者有一種想法.
他們一腔熱誠.
單純地來求上帝.
很熱切,很多眼淚.
我們心裡面想.
用不用這樣啊?.
是不是真的這樣?上帝是不是一定這樣做?.
我們很多時候會冷眼旁觀.
不過.
你再上深一層.

$^{641}$今天.
我們有沒有.
一份確切的應訊.
即使我們生命處於.
一種這樣的階段.
我們要深信我們的神是有能力.
甚至乎是很願意.
將我們從這些困境裡面.
拯救出來.
當我們笑別人神神化化的時候.
其實可能我們自己都陷入.
一種不信的狀態.
我們不相信我們的神會為我們.
開啟生命的出路.
我們只是為我們自己今天的.
信條或者為我們的生活.
建立一個很穩定的系統.
讓人走向這個系統裡面.
我們就覺得將福音.
傳了給別人.
但以你這三個朋友.
他不是這樣看上帝的工作.
當他們面對一個困境的時候.
他們深信不疑.
他們知道.
即便我們將要面對著.
在搖搖搖.
裡面.
活生生地燒.
他知道他所信的主是有能力.
而且會拯救這幾個人.
在這種險境裡面.
救不出來.
他們所信的神是一個.
真身的神.
救不出來.
他們所信的神是一個.
真身實實是一個拯救的主.
即便如此提醒我們.
今天我們即使面對.

$^{681}$這種困境.
我們信不信上帝.
在這種處境當中.
他要救助我們.
當人沒有辦法去改變的時候.
我們都依然.
確認.
確信我們所事奉的.
耶和華上帝.
必然將我們從這種困境.
裡面救不出來.
這個是.
撒德拉米薩亞伯尼哥.
他第一次.
在這裡他所答辯的.
他對尼泊爾薩王.
那種表達.
可能你會問.
讀到這個時候.
牧師.
在歷史也好.
在我們眼前也好.
有時候未必.
事情如我們所願.
我們祈禱.
牙血都流出.
未必如我們心裡面所願的.
信仰或祈求.
會出現的狀態.
這三個友人.
他繼續說.
即使面對上帝.
不是這樣.
來拯救他.
他們怎樣去.
面對他們呢.
在第十八節裡面.
即或不言.
王啊.
你當知道.

$^{721}$我們缺不事奉你的神.
也不敬拜.
你所立的.
金像.
在經文裡面.
這三個友人.
繼續將他們的信心.
再推高一個層次.
即使上帝.
不在這個堯裡面.
即或不言的意思.
即不是這樣做.
我們的信.
告訴我們.
我們的忠誠.
告訴我們.
也告訴我們.
我們不會事奉你的神.
即使面對死也好.
我都不會下拜.
你所鑄造的金像.
我不會事奉你的神.
也不敬拜你的金像.
即或不言.
這個字眼.
更加表達出一種.
堅決的.
忠心和信心.
面對著很多逆境.
可能上帝真的不是這樣.
來救助我們.
但是我們會不會.
繼續持守這個信仰呢.
可能有些靈姐妹.
面對信仰.
人生裡面總會高高低低.
有時候我們會.
可能面對著這些.
高高低低.
但有時候面對低谷的時候.

$^{761}$我們都會產生那種.
信仰上的質疑.
那種.
為何上帝一直要保守我.
為何要讓我面對著這樣的痛苦呢.
又或者面對著.
這些難關的時候.
為何上帝不將這個苦杯挪走.
卻要我面對.
這些艱難呢.
甚至有些靈姐妹.
面對這些艱難的時候.
他就全盤否認了他所信的.
上帝.
所信的福音.
甚至他過去裡面.
一直在教會裡面的生活.
因為一個.
上帝不是用祈求的方式.
來做的時候.
他就全盤推倒.
以致寧願做一個.
沒有神的人.
過往的認信.
因著這些挫敗.
他一一的來推倒.
如果那個信仰只不過是.
為我們帶來信境.
卻不能夠承托我們的逆境.
大概那個.
都是之前很多哲學家所說.
是一個鴉片.
是一個麻醉我們的腦神經.
讓我們面對著.
這個世界很困難的時候.
有些東西麻醉我們的心靈.
到慰藉.
然後就讓我們舒舒服服的過我們的人生.
基督教的信仰從來不是說.
製造一個精神的鴉片.

$^{801}$給我們.
麻醉我們的中樞神經.
讓我們純粹在好景裡面.
可以解釋一下.
逆境的時候.
感覺上.
讓自己舒服一點.
聖經裡面.
特別說到.
大部分上帝所用的人.
上帝大部分.
真的經歷信仰的人.
他們都會經歷過.
有時候上帝不如.
他的心願來.
聽他的心願祈禱的方式.
來生活.
甚至乎有時候.
我們會覺得.
那個信仰.
或者和我們的生活.
帶來一種很荒謬的差距.
我們很多時候.
心裡面有很多美好的意願.
但是.
現實卻不是.
用我心中的意願.
來如願以償.
來出現.
不過.
這三個.
年輕的官.
他們知道.
他們看到的信仰.
不是純粹眼前的信仰.
他更加知道.
他和神的關係.
純粹為他帶來信徑的.
一個神明.
這個神.

$^{841}$是他值得.
用他自己整個人生中心.
來跟隨的主.
有一個老牧者和我說這段經文.
我覺得很有意思.
他說基督徒要實際一點.
Andas, 基督徒要實際一點.
我聽到他怎樣實際一點.
他說當面對一些很大的難處的時候.
最實際的人.
要跟隨一個最有權力的對象.
是不是?.
我說是呀.
要跟隨最有權力.
如果在尼泊爾格尼撒.
他只是地上的君王.
他掌管你今生的王.
夠大權力了.
不過有一個比他更大.
超越他的能力的主.
就是掌管天地萬物.
掌管歷史的主.
如果.
大家去比對一下.
你跟地上的霸主.
還是跟天上拯救的主呢?.
如果實際一點.
你一定要跟天上.
拯救的主.
因為他可以.
殺你的身體.
但地上的主殺不到你的靈魂.
但拯救的主.
卻可以救拔你的一生.
他用一個很實際的角度.
來說,如果你的生命裡.
你真的要跟隨他的.
你應該要跟天上的主.
用你生命.
整盤的.

$^{881}$忠心投放在他身上的時候.
這個天上的主.
永遠不會待薄你.
你就安然地走過.
你人生的路.
即使眼前.
你祈願的事.
不是如你所願地出現.
但天上的.
救贖主總會有.
最好的安排幫助你.
渡過生命裡面.
高高低低.
堅心地跟隨他.
這三個友人.
向尼布格尼撒.
用這種拜金像.
希望他身邊的神民.
來表達他的忠誠.
但這三個.
年輕的官卻是.
在這些關口當中.
向天上的耶和華上帝.
表達最大的忠誠.
即使性命都不要也好.
他都要堅持.
這種信仰.
決不事奉.
尼布格尼撒所建造的金像.
在這個對話之後.
接著就出現了一個.
很重要的狀況.
接著的第.
二十四節至到第三十節.
就講到整件事的結局.
在第.
二十四節裡面.
接著就講到尼布格尼撒王.
聽見這幾個人.
有這樣的.

$^{921}$這麼大的反抗.
於是他就氣沖沖地將這三個人.
扔進火窯裡面.
來進行擊營.
而且怕他燒不死.
還用七倍的溫度.
來加熱火窯.
一個很高的溫度.
甚至乎台.
這三個友人的.
那個僕人因為熱力太強.
台的人都被燒死.
但將這三個人扔進.
窯中的這三個友人.
他們絲毫無損.
在尼布格尼撒的眼中裡面.
明明是將這三個人.
放進去.
在這個窯裡面的.
但經文裡面就記載著.
他見到的是四個人.
而且並沒有捆綁.
在火中遊行.
不是舉牌那些遊行.
即是走來走去的意思.
在這個火中裡面.
繼續的來來去去.
沒有任何的傷害.
第四個人.
就是剛才那三個友人以外的第四個人.
樣貌好像神子一樣.
就是這樣來.
保護了他們.
尼布格尼撒.
見到這個整個的情景.
他自己就很譁然.
他自己就覺得.
為何會有這樣的事來出現.
於是乎他就將.
撒德拉,米薩,阿伯尼哥.

$^{961}$叫他們出來.
從這個窯中.
拉他們出來.
就對他們說.
撒德拉,米薩,阿伯尼哥.
你的神是應當稱頌的.
他猜顯使者救護.
倚靠他的僕人.
他們不遵從王的命令.
捨去己身.
在他們神以外.
不肯事奉敬拜別神.
現在我要降旨.
無論何方何國.
何族的人.
旁獨撒德拉,米薩,阿伯尼哥.
之神.
必被凌遲.
他的房屋必成為糞堆.
因為沒有別神.
能這樣施行拯救.
那時王在巴比倫省.
高聲了撒德拉,米薩,阿伯尼哥.
這裡是28至30節.
就是這件事的一個結局.
當這幾個年輕人堅持下去的時候.
上帝的神蹟就在這裡.
在火搖中出現.
而尼波加里撒原本是一個.
某程度上用強權.
威逼這幾個人.
去臣服的這個人.
他自己看到這樣的結果.
他自己都在這三個人面前.
去服服.
甚至要稱頌撒德拉,米薩,阿伯尼哥.
背後的那位神.
是他自己親口做的一個見證.
甚至乎.
他也見證著這三個忠心的人.

$^{1001}$寧願捨棄自己的身體.
都不遵從尼波加里撒的皇命.
不肯事奉敬拜別的神.
他們的堅持.
一方面是上帝的保守.
令到整件事好像整個結局扭轉.
但同時間這三個年輕人的堅持.
卻成為一個強而有力的見證.
在一個不順.
在一個敵黨.
甚至乎不友善的處境當中.
成為最鏗鏘,最有力的事過.
成為生命裡面令到人們讚嘆不已的一個很重要的見證.
丙姐妹.
今日我們的信仰.
是一個怎樣的信仰.
會否是面對一些逆境,逆流.
或者一些外在的環境.
你面對一些困難的時候.
你就覺得自己好像很委屈.
甚至乎有時候很容易.
就將你的情緒推到極致.
很容易就感懷身世.
將很多不幸的事件串連在一起.
然後就定義自己都是一個不幸的人.
我們用一個受害者的角度.
來看我們自己的人生.
看我們自己的基督教信仰.
不過整個基督教信仰.
或者十字架的真理.
是講到.
我們今日所信的主.
不是一個受害者.
耶穌基督不是一個受害者.
我們也不是要帶著一種.
可憐他覺得他很慘的角度.
來信這個.
令到我們流眼淚的福音.
來信這個.
我們今日所信的是一個得勝的主.

$^{1041}$有時候我們面對的處境.
未必是如我們所願.
事情未必按我們的期望來發展.
我們人生可能都會面對不同的際遇.
不過在聖經裡面.
特別是今日這段經文提醒我們.
今日即使我們面對.
這些困難,這些逆境.
我們已經沒有辦法再捐窿捐瓦.
我們今日夠不夠膽站起來.
宣告那個主會拯救我們的生命.
無論我們落在什麼狀況.
這個神是有能力來拯救我們.
我們永遠都帶著一份盼望.
將我們的祈願.
向祂面前表達.
更加重要的是.
即或不然,上帝未必按我們的心意而行.
我們都會否將自己的生命.
交付在這個永恆的主裡面.
不存盤否定我們自己.
所信的主.
反而是在這些逆流當中.
我們用我們的信念來承托我們.
以致我們有力量走過.
這些眼前很磨難的日子.
我們這些表達.
不是枉然的.
是為了自己而做的.
我們生命的表現.
更加可以成為這個世界的一場戲.
當這個世界很多人說.
哪裡有利益就去哪裡.
哪裡順就哪裡.
我們徹徹底底做一個實用主義者的時候.
我們的信仰可以變得很失效.
不過今日的經文提醒我們.
當我們知道.
我們大是用這個方式來生活.
但我們都依然保持我們的信念.

$^{1081}$在不同的逆境當中.
我們今日一齊的.
堅心寫去.
即使寫去多深.
我們都信靠這位天上的父.
甚至乎.
在我們生命裡.
有些東西比我們自己的生命.
更加值得我們去堅持.
我們有沒有一個這樣的決心.
來去跟隨這位神.
今日的經文不是很複雜.
當我們讀的時候.
我們讀但義理書.
我們見到他環環相扣.
在一個很不友善.
對信仰很疲乏.
很困難的處境當中.
依然都可以有忠心的信徒.
興起來成為一個美好的見證.
我們今日.
我們願不願意成為一個這樣的人.
一說繼續成為一個自怨自艾.
經常都是為了自己生命不好的制約.
而感觸放棄自己的人.
眼前有兩條路.
你會選擇哪一條呢?.
我會選擇.
即或不然.
我們都依然跟隨我們的主.
我相信我們的主.
一定會帶領我們走過面前的路.
我希望你都和我一樣.
選擇一個這樣的決定.
\newpage



\section{哥林多後書 5:1-10}
\label{sec:jdoVqa_V4mw}
\textbf{2020年9月27日 崇拜直播｜哥林多後書5\_1-10}
\newline
\newline
連結: \href{https://youtube.com/watch?v=jdoVqa_V4mw}{\texttt{https://youtube.com/watch?v=jdoVqa\_V4mw}} ~~~~ 語音日期: 2020-09-29
\newline
\newline
\hyperref[sec:_Ymdcnu_z5c]{\small{< < < PREV SERMON < < <}}
~
\hyperref[sec:index_chronic]{\small{[返順時目]}}
~
\hyperref[sec:Jfw_Vv2ojxk]{\small{> > > NEXT SERMON > > >}}
\newline
\newline
哥林多後書 5:1-10
\newline
\begin{longtable}{cl}
\hline
\hline
章節 & 經文 (和合本修訂版)\\
\hline
5:1 & \begin{tabularx}{0.7\textwidth}{X} 因為我們知道，我們這地上的帳篷若拆毀了，我們將有神所造的居所，不是人手所造的，而是在天上永存的。 \end{tabularx} \\ \\ \relax
5:2 & \begin{tabularx}{0.7\textwidth}{X} 我們在這帳篷裡嘆息，渴望得到那從天上來的居所，好像穿上衣服； \end{tabularx} \\ \\ \relax
5:3 & \begin{tabularx}{0.7\textwidth}{X} 倘若脫下也不至於赤身了。 \end{tabularx} \\ \\ \relax
5:4 & \begin{tabularx}{0.7\textwidth}{X} 其實，我們在這帳篷裡的人勞苦嘆息，並不是願意脫下地上的帳篷，而是願意穿上天上的居所，好使這必死的被生命吞滅了。 \end{tabularx} \\ \\ \relax
5:5 & \begin{tabularx}{0.7\textwidth}{X} 那為我們安排這事的是神，他賜給我們聖靈作憑據。 \end{tabularx} \\ \\ \relax
5:6 & \begin{tabularx}{0.7\textwidth}{X} 所以，我們總是勇敢的，並且知道，只要我們住在這身體內就是離開了主。 \end{tabularx} \\ \\ \relax
5:7 & \begin{tabularx}{0.7\textwidth}{X} 因為我們行事為人是憑著信心，不是憑著眼見。 \end{tabularx} \\ \\ \relax
5:8 & \begin{tabularx}{0.7\textwidth}{X} 我們勇敢，更情願離開身體，與主同住。 \end{tabularx} \\ \\ \relax
5:9 & \begin{tabularx}{0.7\textwidth}{X} 所以，無論是住在身內或住在身外，我們都立了志向要得主的喜悅。 \end{tabularx} \\ \\ \relax
5:10 & \begin{tabularx}{0.7\textwidth}{X} 因為我們眾人必須站在基督審判臺前受審，為使各人按著本身所行的，或善或惡受報。 \end{tabularx} \\ \\
[1ex]
\hline
\hline
\end{longtable}
$^{1}$讓我們專心聆聽一段上帝的話語.
記載在哥林多後書第五章一至十節.
經文這樣說.
我們原知道.
我們這地上的帳棚若拆毀了.
必得神所造.
不是人手所造在天上永存的房屋.
我們在這帳棚裡嘆息.
心想得那從天上來的房屋.
好像穿上衣服.
倘若穿上.
被遇見的時候就不至於赤身了.
我們在這帳棚裡嘆息勞苦.
並非願意脫下這個.
乃是願意穿上那個.
好叫這必死的.
被生命吞滅了.
為此.
培植我們的就是神.
他又賜給我們聖靈作憑據.
所以.
我們時常坦然無懼.
並且曉得.
我們住在身內便與主相離.
因我們行事為人.
是憑著信心.
不是憑著眼見.
我們坦然無懼.
是更願意離開身體.
與主同住.
所以.
無論是住在身內.
離開身外.
我們納了志向.
要得主的喜悅.
因為我們眾人必要在基督台前顯露出來.
叫各人按著本身所行的.
或善或惡受報.
我們恭請梁家騮牧師來到我們當中宣講.
講題是《至詐人生的盼望》.

$^{41}$各位弟兄姊妹.
在這個星期和下個星期.
接連兩個星期.
我們都會繼續看哥林多後書.
那麼.
我就找回一些紀錄.
其實這一篇和下一篇.
我在好幾年前.
曾經作為獨立篇章和大家講過.
最少我們講過這段經文.
所以今天.
你會聽到一些信息.
好像很熟悉似的.
不過.
因為我們現在是.
一卷書順著來講.
整卷哥林多後書來看.
我想我們連繫上下文.
可能我們都會看到不同的東西.
並且.
我對大家不是有很大的期望.
幾年前聽過的東西.
你們大概都不記得了.
都是當新的來聽的.
每早都是新的.
《至詐人生的盼望》.
哥林多後書.
一個中心的信息是.
因著天上屬靈的事實.
保羅興翰人間的遭遇.
哥林多頂子妹挑戰保羅.
講他的賣相.
他的現實.
和他所傳講的信仰不相稱.
講他人間的際遇.
和他所講的福音不相稱.
保羅於是說.
有些事情不是今天我們就可以兌現.
我們是懷著盼望.
去指向將來才兌現.

$^{81}$那就是說.
今天沒有了.
全部都是將來的東西.
又不是.
今天我們是.
兩個世界同時夾纏在我們中間.
信仰不只是指向將來的.
只是未來的東西.
那就今天浪費了.
等運到這樣.
不是.
今天已經有真實的一面在我們的生活裡面.
所以保羅就強調.
他是憑信心.
不憑眼見.
他說那些看不到的東西.
比那些看到的東西更加重要.
我們是腳踏兩個世界.
你記得他前段講到.
寶貝和雅氣.
也講過所謂外體和內心.
就是說.
外在的人和內在的人.
那個 outer being(外在).
和 inner being(內在).
講了這些異象.
我們怕外體是一個扭壞的外體.
但我們裡面的生命.
卻是一個活潑完美的生命.
我們整個樣子.
好像一個爛鋼牙一樣.
不過我們裡面.
卻是有一個寶貝藏在其中.
你看到的只是一個鋼牙.
你看到的只是.
就是那個掉頭.
就是掉骨的那個外體.
你看不到我裡面.
那個真實.
普羅強調外在內在.

$^{121}$順著這樣的思路.
普羅在這裡就將這個 imagery(畫面).
進一步的發揮.
他指出我們今天活在人間.
彷彿好像是我們自己.
這個內在的人.
套在一個外體.
那個肉體或者塵世裡面.
就是說我們好像穿上了一件塵世衣.
穿上了一件肉體衣服.
那件真皮.
是不是?.
普羅就用了舊約.
象棚的意象.
指出我們的肉體.
就好像一個地上的象棚.
講象棚.
是在講什麼呢?.
實際上整段聖經.
普羅都是在講一個.
temporal(時間性)的問題.
就是說時間性的問題.
他強調的不是.
我們的肉體臭皮囊.
就是很敗壞.
很污穢.
他強調不是.
他特別在這裡講的是短暫.
暫時的.
是臨時的.
不是永久的.
他關心的.
主要的就是在講時間的長短問題.
象棚的一個特性就是臨時性.
住象棚的人未必是真的.
招行猛賊這麼誇張的.
不過象棚最少的就是.
很容易建造.
又很容易拆除.
對不對?.

$^{161}$方便在旅途中使用.
遊牧民族住在象棚裡面.
為什麼?.
因為他們是逐水草而居.
沒辦法長期停留一個地方的.
牛羊吃光了附近的草.
他沒理由早上帶牠們去五個小時外的地方吃.
然後晚上又帶牠們五個小時回來.
是不是?.
乾脆就搬過去.
對不對?.
所以.
遊牧民族是不會長久定居一個地方的.
因為不長久定居一個地方.
你建造一些堅固的.
用木頭石頭搭建的房子就不值了.
對不對?.
你花這麼多錢.
這麼多時間建造一間房子.
立刻就拆了它.
不知誰會這麼做.
對不對?.
以色列人最典型住在象棚的時期.
是在什麼時候?.
是他們離開埃及.
在曠野漂流的那四十年.
是不是?.
當他們進入了迦南之後.
他們就學做農耕的生活.
過農耕的生活.
所以.
在曠野漂流的那四十年.
就是住象棚的日子.
因為他們沒有固定的居所.
今天的我們應該沒有多少住象棚的經驗.
我們住的房子現在也很堅固.
所以現在香港打十號球.
理論上倒樓的可能性幾乎是零的.
對不對?.
以前你們還沒出生.

$^{201}$什麼六零年代.
什麼溫代那些.
一吹來之後很多房子都倒塌了.
因為那時候很多住木屋.
木屋是一點點像象棚的.
對不對?.
就是兩下手勢就全部吹倒了.
但是現在香港除了打風還發燙.
就是走到海邊觀浪.
然後就被棚架吹下來.
被打傷.
那些是沒得救的.
是不是?.
其實那些應該自己打屁的.
你還要搞一些救護人員去救你.
真的很不值得.
就是.
今天我們住的房子很安全.
所以我們今天沒有什麼住象棚的經驗.
我這一生唯一有住的象棚經驗就是.
小時候去玩 wild camp(野餐).
是不是?.
就是去海灘或者在山邊紮營.
那些是快樂的經驗.
就是住 wild camp 沒痛苦經驗.
所以住象棚我沒有痛苦經驗過.
是不是?.
我能夠 recall(回想).
就是跟這段聖經相關.
大家可以聊得到.
對得到嘴型的.
我想到的就是.
去旅行.
住酒店.
或者住朋友家.
當然你覺得是 happy(開心)的經驗.
你覺得不是痛苦的經驗.
看看吧.
當然有些事情跟性格有關的.
但是多數人都是.

$^{241}$你短暫去玩.
其實是愉快的.
但是如果你說.
一個月兩個月都住在酒店裡.
就算是五星酒店.
都不見得一定是很愉快的.
除非我們的居所太過狹窄.
又或者是我們跟家人的關係太過惡劣.
否則例如.
你知道這幾個月我們出不去旅行.
是不是?.
香港搞很多 staycation(宅度假).
很多酒店.
就是一個 package(套餐).
又吃又住.
你問我.
其實坦白說.
住.
就算是住文華東方.
住半島.
住 Grand Hyatt(大海灘).
其實那個吸引力都不是很大.
我也住過一次.
就是.
算數的就是.
到底那些飲食的 offer(提供).
夠不夠吸引.
就是.
籠床不如狗窩.
如果在香港.
我為什麼不在家裡睡.
我家裡的房間連廳也比他那間房間大.
是不是?.
雖然有病池也不能開.
只是困在房間裡也不知道該做什麼.
所以其實我覺得我不是很吸引.
住酒店.
如果你一天都要開行李箱去拿東西.
其實不是很愉快的經驗.
住在朋友家裡就更加痛苦.

$^{281}$就算幾個朋友的朋友.
你想一想.
每一天早上就要走出去客廳.
然後跟他聊個多小時才有好意思出門.
晚上回來又要聊個多小時才能交數.
要不然住在家裡.
別人幫你打掃等等.
你完全不跟別人做一些娛樂事業也不行.
是不是?.
其實.
總有一場你就覺得.
不是很過癮.
我也跟你分享過.
我跟太太在還沒退休之前.
我們已經有幾年時間去預備.
我們很喜歡去旅行.
我平時是一個很多出門的人.
我一年有四個多月不在香港的.
那麼.
太太就說.
不如大家揹個大背包.
出去旅行.
我們去玩半年.
逐個地方慢慢去.
我立刻反對.
我是很喜歡旅行的.
不過我也說我不喜歡.
通常去旅行.
超過三個星期.
就開始開到肚子裡.
就算是很好的朋友.
不要說兩夫妻.
也有時候會有些張力.
因為累.
是不是?.
我寧願去一段.
回來享受一下.
再去一段.
比你熬到第四個星期.
第二個月.

$^{321}$其實可能是痛苦.
你覺得自己是漂流.
多於覺得自己是旅行.
是不是?.
所以.
大概你可以想象一下.
那種長時間去旅行.
那種有些漂泊.
漂蕩的感覺.
這個就有一點類似.
住帳篷的經驗.
無論如何.
沒有人會在他短暫居庭的地方建造房屋.
因為建造的時間.
建造的費用太貴.
是不是?.
並且建造了又拿不走.
同樣的.
走在我們短暫人生的旅途上.
如果真的活出一個客旅人生的我們.
我們也要視此生的種種.
是一座帳篷.
而不是任何的.
用日本的說法.
叫做不動產.
不是不動產.
五章一節說.
我們原知道我們這地上的帳篷若拆毀了.
未必得上帝所做不是人所所做.
在天上永存的房屋.
保羅強調.
塵世的生命在死後就換上永恆的生命.
基督徒不需要為地上的帳篷被拆毀.
而太過的惋惜.
因為這是讓我們獲得天上永存的房屋的時候.
舊的不去.
新的不來.
擁有永生盼望的人.
不會太過懼怕人間的苦難和死亡.
你要注意的是.

$^{361}$保羅雖然關懷來生.
但他絕對不是一個厭世者.
地上的帳篷雖然是短暫使用.
卻是不可或缺的.
因為在人間是沒有永存的房屋的.
永存的房屋不在人間.
在天上.
一天我們仍然在人間流傳.
我們就要暫住在地上的帳篷裡面.
所以.
保羅只是提醒我們.
不要太過依戀這個塵世的生命.
但保羅其實沒有去污衊.
沒有否定地上的帳篷的價值.
沒有踩到地底泥.
沒有價值.
不是的.
我們這個身子.
雖然不會是永久的.
不過.
基督教從來都強調尊重生命.
不主張傷害生命.
是不是?.
從來沒有人恨惡自己的身子.
總是保養固息.
這個也是保羅所說的.
《以弗所書》五章三十節.
是不是?.
就算對著一個垂死的病人.
羅先生是以前做護士的.
雖然是精神科.
不過做一個護士.
你面對一個垂死的病人.
你明知道他很快死的了.
你都要盡力減低他的痛楚.
是不是?.
讓他活得有尊嚴.
是不是這樣?.
我們不會說.
慈又死,災又死.

$^{401}$隨便吧.
不會的.
是短暫.
但我們仍然是珍惜.
是不是?.
聖經告訴我們.
七位帳篷的上帝不是我們自己.
所以你不要自己自盡.
有一天.
賜生命的主會召我們離開人間.
屆時他會給我們一個新的身體.
就是保羅在《哥林多前書》十五章.
所說的.
我們不是很明白的.
所謂靈性的身體.
是不是?.
其實我也不知道是什麼意思.
Spiritual body(靈性身體).
這個就是永存的房屋了.
保羅說.
我一直期盼獲得那個靈性的身體的那一天.
五章二到四節.
我們在這帳篷裡嘆息心上.
德亞從天上來的房屋.
好像穿上衣服.
倘若穿上被遇見的時候就不至於赤身了.
我們在這帳篷裡嘆息牢符.
並非願意脫下這一個.
乃是願意穿上那個.
好叫這必死的被生命吞滅了.
我先說.
我很喜歡保羅.
說死亡的.
他說死亡的那種.
就是.
超級焦跡.
超級驕傲.
超級豪氣的那個地步.
他怎麼說的?.
好叫這必死的被生命吞滅了.

$^{441}$我一上上以為.
離開世界.
是生命被死亡吞滅.
死亡.
扎了我們的命.
但是保羅的說法剛剛好相反.
離世.
是死亡被生命吞滅.
是不是很過癮?.
所以.
基督徒的死其實是生.
人間的死是得著永生.
是永生的開始.
是不是?.
所以.
我們這個必死的生命.
會被一個.
不死的生命吞滅.
就好像那些.
吃怪獸那樣.
吃吃吃.
吞掉它.
是不是?.
我們這個必死的身體會被.
永遠的生命所取代.
在這裡保羅轉用了.
穿衣服的意象.
我們死後.
就脫下地上的帳篷.
穿上天上的房屋.
好像換衣服那樣.
舊衣服不脫.
新衣服就穿不上去.
是不是?.
並且.
就算穿上去.
你也不會.
你跑完步.
渾身.
就是抽汗.

$^{481}$你會不會就這樣套一件.
乾淨的衣服上去?.
脫了再穿吧.
對不對?.
就是你乾淨的衣服.
套一件髒的衣服.
就是找來搞的.
是不是?.
所以.
我們不脫掉舊衣服.
就不會穿回一件新衣服.
你一旦脫掉舊衣服.
你就要立刻穿回一件新衣服.
不然你就打大赤肋.
這裡說赤身.
你就跟人肉白相見.
是不是?.
在五章.
二到三節保羅說.
我們在這個帳篷裡面嘆息.
心想得了.
從天上來的房屋.
很想穿上衣服.
倘若穿上.
被異見的時候.
就不至於赤身了.
這裡保羅強調.
嘆息.
後面他又再說.
我們在這帳篷裡.
嘆息勞苦.
同一段經文.
兩次提到.
在塵世裡面的嘆息.
我們又會再問.
保羅是不是真的有一個厭世的傾向.
空泛又廣泛的.
將人生理解為.
一聲嘆息.
那麼.

$^{521}$沒有什麼意義.
浪費力氣.
不值得活下去.
不是.
保羅.
當講到嘆息的時候.
他不是.
in general(一般來說).
就這麼說.
生命就是一聲嘆息.
他真的具體指著.
在過去在現在.
遭遇的各種負面的經歷.
很多讓他嘆息勞苦的事情.
再快樂的人.
再活得滿足的人.
只要你不是兩無人事.
一無關係二無責任.
如果你不是沒有關係沒有責任.
就會有很多讓你嘆息勞苦的地方.
不會沒有的.
不可能沒有的.
並且.
我覺得.
奇怪一點.
你的嘆息和快樂是不相衝突的.
人生最快樂的是什麼?.
不是來自吃喝玩樂.
如果不是的話你的層次很低.
讓人得著最大快樂的是兩件事.
第一是關係.
第二是成就.
第一是關係.
第二是成就.
吃喝玩樂.
都需要有 partner(伴侶).
有伴.
自己一個人去吃.
真的吃得很快樂嗎?.
自己一個人去玩.

$^{561}$玩得很快樂嗎?.
最好就是有人跟你一起 share(分享).
所以關係其實很重要.
我們需要別人.
亦都需要有人需要我.
不要說這些.
不是的.
being wanted(想要).
這樣很重要的.
你回到家.
孩子撲在你身上.
你覺得很過癮嗎?.
他們需要你.
有人需要你.
那種快樂是很真實的.
只是獨自一人.
你是很難享受到那種關係上的快樂.
有人在乎你.
有人介意你.
你亦都介意一些人.
你在乎一些人.
做父母的是不是很變態?.
看到子女吃.
是不是比自己吃開心?.
待會你見到他吃.
你自己會覺得很快樂.
因為真的能夠幫到人.
給予人.
你就會覺得很快樂.
不過.
這就是關係.
成就.
成就不一定是指名成利就.
而是你打拼過.
有一個奮鬥的目標.
這個目標完成了.
你覺得自己沒有白活.
沒有枉過.
是不是?.
每天渾渾噩噩.

$^{601}$不知道做些什麼.
你很難快樂的.
每天都好像放假那樣的.
其實真的一點都不快樂的.
你也有些做一下.
有些細藝.
有些事情.
並且通常都是.
越是難做成的事情.
越是需要付代價的事情.
做完之後.
你會越快樂的.
對不對?.
就是要Ray 牧帶一首詩歌.
基本上這個easy job(簡單的工作).
對他來說.
就是又沒有難度.
亦都產生的快樂很有限的.
你要我帶一首歌.
糟了.
我自己走音走到瘋了.
是不是?.
很艱難.
可能就是避很久很久.
是不是?.
亦都被人罵很久.
又走音了.
這個又不好了.
這個又錯了.
但是我打拼了三個月之後.
然後都做出了挺好的效果.
是不是快樂很多?.
那個是成就.
一個 achievement(成就).
所以永遠都是.
你花三年時間寫好這篇論文.
那種寫好的感覺.
那種解脫的感覺.
都不知道怎麼形容.
是不是?.

$^{641}$你說一晚就鬥起的那些作文.
你不會很快樂的.
是不是?.
我想說的就是.
關係和成就是真的.
為人生帶來最大快樂的兩件事.
多過吃喝玩樂.
真的.
而你發覺這兩件事.
這兩件事.
都一定有嘆息.
有勞苦.
你愛一個人.
你抱一個人上身.
你就有很多的掛慮.
是不是?.
很多的思念.
你對子女不會只有快樂.
沒有憂愁.
有很多的擔心.
是不是?.
就算他長大了.
你還是有很多擔心.
是不是?.
真是.
就是常慧.
百載憂.
是不是?.
就是對著子女.
你為什麼沒有憂慮?.
是不是?.
你.
愛惜這間教會.
你毀身這間教會.
你就會好像保羅說的.
教會的事天天壓在我的心上.
是不是?.
所以.
永遠都是關係帶來快樂.
但通常有夾前很多的.

$^{681}$嘆息.
勞苦.
正如剛才說成就.
如果可以帶給我們最大的快樂.
但成就同樣.
是夾前很多的嘆息勞苦.
是不是?.
保羅說事奉.
常常都強調事奉的勞苦.
你們在主裡面的勞苦不是圖鑒的.
你們記不記得他說的.
勞苦勞苦.
如果純粹是容易的工作的話.
那個事奉是有限的.
永遠都是毀身.
奮進.
然後你勞苦.
所以.
保羅強調.
我們有嘆息.
不是一個厭世的想法.
不是全盤否定生命的價值.
不過.
確實在人間.
不完美的世界.
不完美的他人.
也有不完美的自己.
我們常常會在很多很多的.
嘆息勞苦之中.
事奉完全沒有眼淚是很罕有的.
是不是?.
不完美的自己加上不完美的人.
我們總是會有很多碰撞.
有很多要適應的地方.
是不是?.
這個是勞苦.
保羅補充說.
在五章四節.
我不再只將盤去嘆息勞苦.
並非願意脫下這個.

$^{721}$乃是願意穿上那個.
好叫者必死的被生命吞滅了.
在這裡.
保羅.
要突顯要強調的.
他不是有一個厭世的思想.
我們期盼永恆的生命.
不是因為我們現在活得不耐煩.
因為現在太過難受.
現在太過不好.
而是因為將來.
更好.
我們希望以最好的.
去取代如今的好.
所以不是厭世的.
我不是活得不耐煩.
你問我是不是現在很想捏死自己.
去見上帝.
我不至於.
現在每一天的生活.
我都仍然覺得很過癮.
是不是?.
覺得很快樂.
不過.
我知道將來.
有更好的.
是不是?.
保羅指出.
當我們仍然擁有地上的帳棚的時候.
我們還未與主完全聯合.
雖然很多靈修神學.
常常講如何追求與上帝.
有一個spiritual union(靈性聯合).
有一個屬靈的神秘契合.
不過.
保羅在這裡告訴我們.
只要你活在人間.
你就不用恣意與主完全契合.
沒有這件事的.
沒有這件事.

$^{761}$五章六節.
住在身內便與主相利.
五章八節.
唯有離開這個塵世的軀體.
我們才可以與主同住.
所以.
這裡保羅是將身體和主做一個對照.
你住在身內就與主相利.
你離開身體就與主同住.
你不可能在人間就與主.
可以緊密聯合.
沒有這個情況.
我們羨慕.
但不可能的.
當然對保羅來說.
這不是一個兩難的選擇.
到底你的選擇.
住在身內與主相利.
還是離開身體與主同住.
保羅說我不是要做這個選擇.
因為.
離世和基督同住是好得無比的事.
但是.
活在人間.
也不是太痛苦的事.
並且.
坦白講.
這也不是我們真正的選擇.
你有選擇嗎?.
上帝還沒召你走.
你可以選擇離開身體與主同住嗎?.
你很大隻的.
我告訴你.
你自己戳死自己.
與主同住嗎?.
你狠一點吧.
是不是?.
所以.
其實我們是沒有選擇的.
就算保羅怎樣叫.

$^{801}$誰能夠我脫離這取死的身體.
他都仍然不可以.
用刀片割手的.
對不對?.
所以.
我們只能夠耐心等候.
希望身體得贖的一天.
來臨.
你不能.
不能急.
無論如何.
保羅是立志以得著永恆的生命.
作為他人生的追求.
即使他活在地上的帳篷裡面.
他仍然是朝得著天上的房屋而奮鬥.
他的行事為人.
順應那個屬靈的原則.
五章七節是憑著信心不憑著眼見.
他始至過聖潔和誠實的生活.
五章十節.
準備將來接受基督的審判.
而最最重要的是.
在五章六節和八節.
保羅都強調.
人間所有的經歷.
都不會教他畏懼.
No fear.
Fear not.
不再怕.
一個畏懼基督審判的人.
五章十一節.
是不會畏懼人間任何的力量的.
基督徒的人生.
是一個預知終局的人生.
我們知遇不知近.
確實.
今天下午會發生什麼事我不知道.
明天會發生什麼事.
我不知道.
香港的後日會怎麼樣我不知道.

$^{841}$但是.
世界的終局我知道.
我的終局我知道.
是不是.
我知道.
我死後一定會得著永恆的生命.
因為耶穌基督在離開世界之前.
是對門徒這樣應許的.
祂去.
就是為我們預備地方去.
將來祂要再回來.
接我們去祂那裡.
是不是.
這個《約翻音》十四章二到三節的應許.
我們的主.
已經為我們預備了復活的生命.
弟兄姊妹.
這番話是我以前常常都說的.
為什麼今天很多基督徒不再熱心傳福音.
為什麼很多教會不熱心傳福音.
是因為我們失落了一個末世的信仰.
我們對人的永生永死.
不再關切.
我們的生活.
也不再是末世性.
說到應容這部分.
你會聽我說過.
你會記得我以前說過的話.
今天是一個樂生的年代.
就是 Enjoy life(享受生活).
Enjoy life(享受生活)的年代.
我們追求生活各種的享受.
也積極享受生活.
我不是覺得這是錯的.
我曾經說過題外話.
我兒子有時候說話很刻薄.
不過也很到位.
他經常批評我們這些做傳道人做牧師的.
不太接地氣.
他就說這個人沒有生活.

$^{881}$他也不知道生活為何物.
這樣的人其實是很乏味的.
所以我這裡不是說.
這是一定不好的.
Enjoy life(享受生活)是不好的事.
不過.
凡事不可以太過分.
今天我們樂生到一個很盡的地步.
我們期待生活裡面的享受越多越好.
享受的時間越長越好.
今天的基督徒多數都不渴望主在內.
我們最多是說死後上天堂.
你不用來.
你不用這麼快來搞亂我們的當.
將來我們上去.
並且上去之後.
越遲越好.
對不對?.
講主在內.
我們都知道我們無法確定具體發生的日期.
所以你也不知道日期是不能拖的.
但是如果說死後上天堂.
我們就認為我們有一點點能力.
可以將這個日期延到越後.
越後.
越後越好.
哈哈.
雖然耶穌說過.
我們沒有一個人可以藉著我們的思慮.
使受數多加一刻.
不過我們不是很信邪.
我們總是相信.
控制飲食.
多做運動.
對於延年式益壽.
是不是很有用?.
你記得我這個例子說過.
我看過一個統計的報告.
很好玩的.
很好笑的.

$^{921}$他問一個問題.
那些每天都做運動.
定期運動的人.
和那些沒有做運動的人.
誰比較長命呢?.
那個研究的結果就是.
做運動的人比較長命.
不用這麼悲觀.
做運動的人真的會長命一點.
不過做運動的人大概長多少命呢?.
他們的說法就是.
就是長了他做運動浪費了的時間.
如果我.
每天.
我現在是真的.
我這半年.
用兩個小時行山.
就是我用了十二分之一的時間行山.
我大概比那些不做運動的人.
我會長十二分之一左右.
長他三兩年命.
不過三兩年命就用來行山了.
就是.
中國人有句老話說.
吃多少做多少是注定的.
是不是?.
我太太經常用這句話來說我.
因為我很喜歡吃自助餐.
我太太就說.
一個人.
你一生吃自助餐的數目是注定的.
例如吃一百餐.
就是你一生的 quota(配額).
如果你每個星期去吃.
你就很快過了 quota(配額).
就會死的.
拖錢點吃.
反正你也是吃一百餐.
這些說法很 mystical(神秘)的.
你信就行,不信就行.

$^{961}$我也不是想推廣這些.
現代人.
我們有兩大任務.
這個一定不記得我說過.
第一,現場受命.
第二,現場清出.
這是兩大任務最神聖.
最重要.
並且是畢生的任務.
所以跟這兩個任務相關的生意.
在現在全世界都是最厲害的.
補品啊.
簽體啊.
健體啊.
化妝品啊.
你知不知道每年.
這些生意當然是想笑.
不是一家,是笑.
很可怕的.
養活全世界的境外人口.
還有餘有症的.
很可怕的.
現代社會的一個特徵是什麼呢.
香港就是這樣.
很多藥房.
不過這些藥房不是賣治療性的藥房.
我們西洋菜街.
你看一下就知道了.
那些藥房.
就是這兩樣東西.
現場受命.
現場清出.
做這兩樣的服務.
是不是.
有時候你都覺得挺有趣的.
我們有常常.
在生活裡面很多新淫.
很多慨嘆.
又覺得自己活得不開心.
但又很在意自己的命不夠長.

$^{1001}$你覺不覺得很矛盾.
你的親人離開了.
你就有很多遺憾.
很多抱怨.
今天你和他們吃午飯.
你有沒有雀躍的感覺.
你又沒有.
你覺得很該死.
是不是.
你活得這麼不耐煩.
為什麼又要這麼在意.
長不長呢.
現代人很怕年老.
我真的可以作證.
我過了60歲.
不過總之自己不照鏡就不覺得自己很老.
是不是.
這個例子我講過你聽過.
有一次我和太太在廣州.
逛過一個服裝店.
寫著中老年服裝店.
我太太笑到亂了肚.
誰會進去這間服裝店買衣服.
你說話.
是不是.
沒有的.
怎麼會去呢.
傻的嗎.
我去中老年服裝店.
我就不去了.
是不是.
建道曾經辦過一個登錄講座.
參加的人不多.
可能是80歲以上的人.
辦下事去參加一個登錄講座.
將自己減低20年.
我是60歲我就不會去這樣的講座.
所以這些題目其實沒有什麼好結果.
是不是.
跟你講過這個笑話.

$^{1041}$真實笑話.
我有一次在街上.
就在遠處見到一個比我高五屆的師姐.
在中大的師姐.
我遙遠的叫她.
師姐.
她.
轉頭望一望.
見到我.
立刻飛奔來我面前.
跟我說.
你這個死小子.
甩頭老翁.
你當街叫我師姐.
你想死啊.
暴露她的真實年齡.
想死啊.
人家修飾這麼久.
隱瞞真實年齡.
被我揭穿真相.
從此我學了功課.
我以後見任何異性.
我都只會叫師妹.
絕對不會再有師姐這個名詞.
是不是.
落身的心態影響我們的信仰.
我們常常為自己和身邊的人.
疾病得不到醫治.
壽命不夠長.
而埋怨上帝.
甚至懷疑信仰的有效性.
是不是.
都不靈的.
上帝都沒有保守的.
我們期望用信仰去幫助我們延年益壽.
信仰成為了一個祈福的手段.
為什麼上帝不保守他的兒女長命一點.
這是我經常聽到的問題.
雖然在舊約記載了一些人.
蒙上帝保守又富有又長壽.

$^{1081}$但在新約聖經.
只有一處.
輕輕提及和長壽有關的是什麼.
就是孝敬父母.
沒有了.
沒有說吃補品那些.
沒有說做運動.
做什麼可以使你長壽呢.
除了孝敬父母之外.
我想不到.
如果你找到的話.
告訴我.
可能你那本聖經和我有些不同.
其實沒有.
沒有能為力的.
所以我經常跟別人說.
如果你真的想長壽一點.
其實你信道教可能比信基督教更有效.
他們的連丹,氧氣都是教你長壽的.
基督教其實不是刻意教你長壽的.
我不是說基督徒要追求死亡.
嚮往死亡.
但最少公道一點說.
不要過分恐懼死亡.
逃避死亡.
因為我們比任何人都更加知道生命的真相.
死亡的真相.
並且我們知道死後的去處.
而這個去處.
不要說你很期待.
最少你沒有理由抗拒.
有沒有搞錯.
跟耶穌同在.
你抗拒什麼.
是不是.
你得著永遠的生命.
有什麼好抗拒的.
不會的.
是不是.
所以有幾個信仰的人.

$^{1121}$總不能夠比其他人更加抗拒死亡.
並且為自己和家人的死亡去埋怨上帝.
否則你的永生盼望在哪裡.
希望你不要介意我這樣說.
基督徒如果太過受困於生死.
你的信仰層次很低.
你真的很低B.
我不反對基督徒享受生命.
樂生.
但不要太過分.
不要過分到因為樂生的緣故而抗拒永生.
過去教會為了讓信徒追求天上的事.
思念天上的事.
總是致力教導信徒.
要輕視現世的生命.
並且根據這樣的信念.
訂定很多生活的要求.
基督徒不要參加人間這麼多的娛樂活動.
不要看電影.
不要聽流行歌.
女生不要化妝.
是不是.
我看一個師母的傳記.
李飛吳師母的自傳.
她說到她小時候.
她是針信會的.
她說教會是反對女性化妝的.
反對得很激烈.
女人一定不可以用虛假的東西遮住自己的真面目.
不過.
這個不止我沒有跟你說過.
見道仇人院以前是嚴禁女士化妝的.
到了八十年代.
我已經在見道教書.
那個時候有人提出.
如果姐妹不化妝.
在畢業禮登台拿證書.
就不太好看.
個個都白皙.
很難受.

$^{1161}$於是就一下子准許女同學上台拿證書的時候化妝.
結果效果是怎樣.
那些姐妹以前從來沒有化過妝.
個個化到好像自取威虎山那些.
嚇得你死.
於是我們看見馬上不行.
於是學校就馬上做了一個議決.
請一個外面的專人來教我們同學化妝.
由禁止化妝到第二個禮拜教他們化妝.
你覺得很古怪嗎.
是不是很古怪.
腳圖追求延長壽命.
延長青春.
享受生活.
今天我們不再輕易說追求永恆.
整個社會文化都是在說樂生.
你想家琪去哪裡玩.
待會去哪裡吃飯.
這是我們最主要的話題.
有時我都在想.
腳圖和飛躍圖.
我們生活上有甚麼差異.
人玩我玩 人做我做.
最多是我們腳圖.
我希望是.
我們生活比較簡樸.
我們對物質的追求不要太過強烈.
不過我都不相信我們每個人都做得到.
怎樣識別出我們有永生的盼望.
這都是我在預備這篇訊息時.
我自己想了很久的事.
我又不可以說些抽空的話.
或者我這樣說.
很多時候那個差異只能在我們參與事奉裡區分.
怎樣證明你有永生的盼望.
怎樣證明你思念天上的事.
是因為我們在敬著人間名利的同時.
我們都花時間花資源參與天國的事奉.
投身教會的服侍.
福音的使命裡面.

$^{1201}$我們不是完全做到.
但我們大致仍然是跟著耶穌的教導.
我們同時都有求他的國和他的義.
我們不只是追求吃喝.
我們不是思念吃甚麼喝甚麼的意義.
我們都有思念天上的事.
弟兄姊妹.
學習更多關心上帝家裡面的事.
關心福音的侍貢.
關心那些為福音使命擺上自己的.
我們中間的一些支持的宣教士.
在禱告,在金錢事奉裡面參與他們的工作.
我們自己沒有去.
但我們參與他們去的行動.
藉此表達我們不只是關注眼前的需要.
學習更多的付出.
為教會,為他人的需要.
奉獻自己的時間和金錢.
耶穌的說法永遠是一個重要的指標.
你的財寶在哪裡.
你的心在哪裡.
你是不是珍惜那東西.
就看你的錢壓在哪裡.
這件事我覺得是一個永恆的學習.
我沒有任何要展示.
聖經是禁止我們展示的.
不過像我們這樣的年紀的人.
弟兄姊妹我們彼此挑戰.
不只是我挑戰你.
我們互相挑戰.
我真的認識不少跟我這樣年紀的弟兄.
我們現在是學十三,十四奉獻.
因為真的沒有其他需要.
基本責任盡了之後.
記錄下來.
給自己一些挑戰.
看看是不是可以.
因為十一已經是by default(須要).
你好像真的隨便.
試試給自己一些挑戰.

$^{1241}$逼自己學習.
去關懷在付出裡面獲得最大的快樂.
對現.
聖經說的.
好像很抽空.
好像很高調.
施比受更為有福的這個應許.
是不是.
這個是幫我們去稍稍擺脫現世的追求.
渴想永生.
我覺得是一個很好的切入點.
多些渴想上帝的婉語.
是不是.
多於人間的房屋.
最後.
簡單重溫保羅在這段聖經給我們兩個還有信息.
很快就完結了.
第一.
坦然無懼.
保羅在五章六節和七節.
先後兩次提到.
基督徒不要怕死亡.
不要因為留戀人間而抗拒死亡的悼念.
耶穌說過.
殺身體不可以殺命運的.
不要怕他.
第二.
預備好迎見主.
在五章九到十節.
保羅說.
無論是生是死.
都要致力土主的喜悅.
並且預備好.
將來在審判台前.
向上帝交仗.
我講過我不喜歡剪報.
或者在網上找一些故事出來.
我所講的這些例子.
全部都是生活中親身經歷的例子.
所以例子不多.

$^{1281}$以下這個故事.
我可能跟你講過.
不過這個故事.
是我真的很感動的一個故事.
我以前在一間教會.
有一個弟兄.
一個好弟兄.
跟我差不多年紀.
他告訴我.
他最近和太太慶祝結婚三十週年.
和一班朋友一起慶祝.
那些朋友就起哄.
跟太太講求愛宣言.
當眾講.
於是這個老弟兄.
就跟他的老伴這樣講.
他說.
假如還有來生的話.
我都希望可以再娶你.
感動到他太太眼淚連連.
這個弟兄.
向我複述這番話.
不忘我是一個牧師.
於是馬上.
講完浪漫的話.
馬上很不浪漫的.
就做一些反省.
批判.
自我批判.
牧師.
我剛才講這番話.
是神學上不對的.
我說是的.
真的不對的.
在天堂是沒有嫁娶的.
所以.
第一你沒有來生.
第二你沒有來生.
再娶她可能.
你只能夠做一世夫妻.

$^{1321}$是不是很不浪漫.
講這些話.
於是.
那位弟兄就說.
如果下次有三十五或者四十週年的時候.
我要怎樣講.
講一些神學上對.
但是又夠浪漫的話.
於是我們兩個老男人.
就一起去量度.
真的很好玩.
我們量度了很久.
最後我們夾出來的台詞就是.
太太.
我為過去三十年.
跟你在一起的日子.
為此.
我特別期待.
將你交在耶穌的手中.
你是值得由耶穌親自照顧.
雖然在人間我們獲得過很多快樂.
但我相信.
能夠讓你在耶穌的手裡.
由祂親自照顧.
你會更加開心.
耶穌比我更愛你.
祂比你更大的快樂和滿足.
我希望我們母的弟兄姊妹.
我們珍惜.
但我們不苟前.
我們享受現在.
但我們更加期盼將來.
上帝恩代我們每一個.
\newpage



\section{但以理書 4:4-18}
\label{sec:Jfw_Vv2ojxk}
\textbf{2020年10月11日 崇拜直播(主餐)｜但以理書4\_4-18}
\newline
\newline
連結: \href{https://youtube.com/watch?v=Jfw_Vv2ojxk}{\texttt{https://youtube.com/watch?v=Jfw\_Vv2ojxk}} ~~~~ 語音日期: 2020-10-12
\newline
\newline
\hyperref[sec:jdoVqa_V4mw]{\small{< < < PREV SERMON < < <}}
~
\hyperref[sec:index_chronic]{\small{[返順時目]}}
~
\hyperref[sec:1ZS5UTvrjn8]{\small{> > > NEXT SERMON > > >}}
\newline
\newline
但以理書 4:4-18
\newline
\begin{longtable}{cl}
\hline
\hline
章節 & 經文 (和合本修訂版)\\
\hline
4:4 & \begin{tabularx}{0.7\textwidth}{X} 「我－尼布甲尼撒安居在家中，在宮裡享受榮華。 \end{tabularx} \\ \\ \relax
4:5 & \begin{tabularx}{0.7\textwidth}{X} 我做了一個夢，使我懼怕。我在床上的意念和腦中的異象，使我驚惶。 \end{tabularx} \\ \\ \relax
4:6 & \begin{tabularx}{0.7\textwidth}{X} 因此我降旨召巴比倫的智慧人全都到我面前，要他們將夢的解釋告訴我。 \end{tabularx} \\ \\ \relax
4:7 & \begin{tabularx}{0.7\textwidth}{X} 於是那些術士、巫師、迦勒底人、觀兆的都進來，我將那夢告訴他們，他們卻不能把夢的解釋告訴我。 \end{tabularx} \\ \\ \relax
4:8 & \begin{tabularx}{0.7\textwidth}{X} 最後，但以理，就是按照我神明的名字稱為伯提沙撒的，來到我面前，他裡頭有神聖神明的靈，我將夢告訴他： \end{tabularx} \\ \\ \relax
4:9 & \begin{tabularx}{0.7\textwidth}{X} 『術士的領袖伯提沙撒啊，我知道你裡頭有神聖神明的靈，甚麼奧祕都不能為難你。現在你要把我夢中所見的異象和夢的解釋告訴我。』 \end{tabularx} \\ \\ \relax
4:10 & \begin{tabularx}{0.7\textwidth}{X} 「我在床上腦中的異象是這樣：我觀看，看哪，大地中間有一棵樹，極其高大。 \end{tabularx} \\ \\ \relax
4:11 & \begin{tabularx}{0.7\textwidth}{X} 那樹漸長，而且茁壯，高得頂天，從地極都能看見， \end{tabularx} \\ \\ \relax
4:12 & \begin{tabularx}{0.7\textwidth}{X} 葉子華美，果子甚多，可作所有動物的食物；野地的走獸臥在蔭下，天空的飛鳥宿在枝上，凡有血肉的都從這樹得食物。 \end{tabularx} \\ \\ \relax
4:13 & \begin{tabularx}{0.7\textwidth}{X} 「我觀看，我在床上腦中的異象是這樣，看哪，有守望者，就是神聖的一位，從天而降， \end{tabularx} \\ \\ \relax
4:14 & \begin{tabularx}{0.7\textwidth}{X} 大聲呼叫說：『砍倒這樹！砍下枝子！拔掉葉子！拋散果子！使走獸逃離樹下，飛鳥躲開樹枝。 \end{tabularx} \\ \\ \relax
4:15 & \begin{tabularx}{0.7\textwidth}{X} 樹的殘幹卻要留在地裡，在田野的青草中用鐵圈和銅圈套住。任他讓天上的露水滴濕，和地上的走獸一同吃草， \end{tabularx} \\ \\ \relax
4:16 & \begin{tabularx}{0.7\textwidth}{X} 使他的心改變，不再是人的心，而給他一個獸心，使他經過七個時期。 \end{tabularx} \\ \\ \relax
4:17 & \begin{tabularx}{0.7\textwidth}{X} 這是眾守望者所發的命令，是眾聖者所作的決定，好叫世人知道至高者在人的國中掌權，要將國賜給誰就賜給誰，並且立極卑微的人執掌國權。』 \end{tabularx} \\ \\ \relax
4:18 & \begin{tabularx}{0.7\textwidth}{X} 「這是我－尼布甲尼撒王所做的夢。伯提沙撒啊，你要說明這夢的解釋；我國中所有的智慧人都不能把夢的解釋告訴我，惟獨你能，因你裡頭有神聖神明的靈。」 \end{tabularx} \\ \\
[1ex]
\hline
\hline
\end{longtable}
$^{1}$我們聆聽聖導.
繼續進入去敬拜.
經訓今天是選至.
舊約的《但義理書》第四章.
四至十八節.
大家教會手上的聖經是.
1198頁.
舊約聖經1198頁.
找到之後請聽我讀出.
《但義理書》第四章四至十八節.
我尼布格尼撒安居在宮中.
平順在殿內.
我作了一夢使我懼怕.
我在床上的思念.
並腦中的異象使我驚惶.
所以我降旨.
召巴比倫的一切哲士到我面前.
叫他們把夢的講解告訴我.
於是那些術士.
用法術的.
迦勒底人.
官少的都進來.
我將那夢告訴了他們.
他們卻不能把夢的講解告訴我.
沒後那照我神的名.
稱為白提沙剎的但義理.
來到我面前.
他裡頭有聖神的靈.
我將夢告訴他說.
術士的領袖白提沙剎.
我知道你裡頭有聖神的靈.
甚麼奧秘的事.
都不能使你為難.
現在要把我夢中所見的異象.
和夢的講解告訴我.
我在床上腦中的異象是這樣.
我看見當中有一棵樹.
極其高大.
那樹漸長而且堅固.
高得頂天.

$^{41}$從地極都能看見.
葉子華美.
果子甚多.
可作眾生的食物.
田野的酒瘦.
臥在蔭下.
天空的飛鳥.
縮在枝上.
凡有血氣的.
都從這樹上得食.
我在床上夢中.
腦中的異象.
見有一位守望的聖者.
從天而降.
大聲呼叫說.
佛島這樹.
堪下枝子.
搖掉葉子.
拋散果子.
使酒瘦離開樹下.
飛鳥躲開樹枝.
樹等卻要留在地內.
用鐵圈和銅圈箍住.
在田野的青草中.
讓天路滴濕.
使它與地上的瘦一同吃草.
使它的心改變.
不如人心.
給它一個壽心.
使它經過七期.
這是守望者所發的命.
聖者所出的令.
好叫世人知道.
至高者在人的國中掌權.
要將國賜與誰.
就賜與誰.
或立極卑微的人.
執掌國權.
這是我尼布格尼薩王所作的夢.
百提沙剎.

$^{81}$你要說明這夢的講解.
因為我國中的一切之事.
都不能將夢的講解告訴我.
唯獨你能.
因你裡頭有聖神的靈.
以下是宣講.
有本會的主任牧師陳寗全牧師.
來到宣講.
「隔逢中的中部面對驕傲的人」.
陳寗全牧師.
弟兄姊妹平安.
一起領受上帝的話語.
我不知道你一回來有什麼感覺.
曾經我們在六七月的時候.
我們都有實體.
然後沒多久又要再停的時候.
今天再回來.
我相信大家心裡面有很多很深刻的感受.
事實上在這半年有多的日子.
我們都沒有去過.
大家隔著螢幕隔著電腦看.
相信我們都盡力做好的.
但是和在現場參與崇拜的體會很不同.
希望今天我們能夠在這個時間裡面.
我們親身面對面的經歷.
從神而來的話語.
以及對我們一些的勸勉.
今天我們看的這段經文.
是繼續我們之前看的整個系列.
「隔逢中的中部」.
事實上但以來.
和他三個友人.
一直在經歷著尼布甲尼沙.
或者被魯巴比倫之後.
有幾個朝代.
而在第四章裡面.
大概是尼布甲尼沙皇.
最後一次出現他的名字.
然後之後就某程度上.
就改換了另外一朝.

$^{121}$但是但以來和他三個朋友.
他一直在這些不同的君主之下.
他繼續忠心地去做.
實踐上帝給他們的召命.
一個很不容易的處境.
一個被擄.
好像二等公民的身份.
卻是在一個很重要的位置裡面.
在關鍵的時刻.
他顯露出上帝與他們同在.
而且這一班人的生命.
能夠在一些很極端.
甚至乎有時候.
起路不形於色.
或者起路無常.
的君主之下.
他們不單止能夠生存.
也能夠將上帝的真理展現出來.
今天我們看的第四章.
我們剛才讀的經文.
我們知道了.
尼布甲尼沙皇又再一次發夢.
在第二章裡面之前也是這樣.
到了第四章又再重複一個這樣的現象.
不過這次發這個.
是一個噩夢.
這個夢特別針對的.
就是尼布甲尼沙.
不單止是說他見證了上帝的真理.
其實在尼泊爾格尼撒王秉政的時候.
新巴比倫王朝秉政了大概五十多年的時間.
而第四章裡面出現的歷史事蹟.
大概是他主政的二十多年.
第三章到第四章其實不是緊接著的.
在歷史裡面其實那個時間已經走了二十多年.
聖經學者特別講到.
當新巴比倫王朝秉政到二十多年的時間裡面.
整個國家.
整個巴比倫的幅員.
大到一個地步.

$^{161}$是一個超級大帝國.
他的幅員由波斯灣到地中海.
南邊由埃及至到伊朗.
你想一下中東地區所有.
包括以前很強大的埃及.
在巴比倫帝國之下.
而一個這麼寬這麼幅員這麼廣大.
而且有這麼多民族的人.
來到這個帝國裡面的時候.
他的稅收又很厲害.
於是乎他有足夠的稅收的時候.
他就興建了很多不同很偉大的大型的建築.
其中一個就是.
舊的世界七大奇蹟.
那個奇蹟叫做空中花園.
不是你家住的那個.
不是在52樓有一個隔火層.
變成了花園的空中花園.
在巴比倫王朝裡面.
曾經尼泊爾王爺為了他的愛妃.
已經被我們2000多年前.
已經想到了.
去做一個很大的庭園.
而且一級級十級以上.
一個很高大.
一個很偉大的建築物.
上面再滿花草.
有思鄉病的皇后.
每一次去到這個花園.
在這裡走走散步的時候.
都能夠想起她的家鄉.
建築這麼高的原因就是.
讓她可以遙遙遠遠的.
看到自己的家鄉.
旁邊的鳥語花香.
是一個賞心樂事.
在這裡說到尼泊爾格尼薩王.
安居在宮中.
宮中就是空中花園.
窮盡當時大量的稅收.

$^{201}$而且有很多建築的智慧在當中.
所起的.
今天這個空中花園還有歷史遺跡.
不過只剩下地底.
上面空中的地震和打仗都沒有了.
尼泊爾格尼薩王在一個.
很值得自豪的宮殿裡.
他就覺得自己一切.
都很平順.
很平安很滿足.
過著他自己主政的日子.
大概在這二十多年的經歷裡.
他一切都很滿足.
不過在他最滿足的時間裡.
然後第十節裡就開始記載著.
所有東西都是令他.
心滿意足的王朝.
這個王就發了一個很奇怪的夢.
在第十節到第十七節裡.
特別是第十三節.
十三節裡講到.
我在床上腦裡有個異象.
見有一位守望的聖者.
從天而降.
大聲呼喊說.
佛島這棵樹.
撼下枝子搖掉的葉子.
拋散那些果子.
使酒獸離開樹下.
飛鳥離開樹枝.
樹等卻要留在地內.
用鐵圈和銅圈箍住.
在田野的青草中.
讓天路滴濕.
使他和地上的獸.
一同吃草.
使他的心改變.
不如人心.
給他一個獸心.
使他經歷七期.

$^{241}$這是守望者所發的命.
聖者所出的命.
好叫世人知道.
至高者在人的國中掌權.
要將國賜與誰就賜與誰.
或立卑微的人.
執掌國權.
在這個夢裡面.
他見到一棵很大的樹.
突然之間有天而降的聖者.
將這棵很大的樹.
斬了它.
撼下枝子搖掉葉子.
拋散那些果子.
然後斬了它.
剩下樹等.
然後.
令到本來一棵參天大樹.
裡面寄居了.
或者棲息了很多動物.
那些葉子.
被人搶走了.
棲息了很多動物.
那些野獸動物飛鳥.
全部趕走它.
這棵樹的意象就是這樣.
然後.
第二個意象就是.
要將這個心.
這棵樹的心.
或者將人的心.
改變更換過來.
變成一個野獸的心.
經過七期.
一個完全的周期.
完全的日期.
日期滿了.
才回復給它.
國中來掌權.
第十七節就是這個夢裡面.

$^{281}$一個很重要的.
聖者所說的話.
就是要將國賜予誰.
就賜予誰.
或立極卑微的人.
執掌王權.
這個夢是很奇怪的.
當一個如日方中的尼布甲尼撒王.
看到整個.
他自己的國家.
進入一個高峰的時候.
突然間看到一個這樣的夢.
某程度上是一個很大的惡夢.
但是他不懂得解釋.
於是就找回.
之前他所訓練的一班迦勒底.
迦勒底其實就是巴比倫的.
那些族人.
找那裡通過裡面的術士.
沒有一個人能夠將.
他這個夢可以完完全全.
地解釋出來.
於是尼布甲尼撒王.
再一次.
要找丹爾利.
來將他這個夢.
解釋出來.
我沒有印給大家.
不過在第四章22-27節.
丹爾利就將.
這個夢.
解釋給他知道.
我讀一讀那段經文.
給大家聽.
王啊!這漸漲又堅固的樹.
就是你.
剛才說的夢裡面的樹.
就是尼布甲尼撒王自己.
你的威勢漸漲及天.
你的權柄管道地極.

$^{321}$王記看見一個守望的聖者.
從天而降.
說將這樹砍伐毀壞.
樹等卻要留在地內.
用鐵圈和銅圈.
箍住.
在田野的青草中.
讓天怒滴濕.
使他和地上的獸一同吃草.
直到經過七期.
王啊!講解的就是這樣.
臨到我主我王的時.
是出於至高者的命.
你必被趕出.
離開世人.
與野地的獸同居.
吃草如牛.
被天怒滴濕.
且要經過七期.
讓你知道至高者在人的國中掌權.
要將國賜與誰就賜與誰.
守望者既吩咐.
全留樹等.
等到你知道諸天掌權.
以後你的國必定歸你.
王啊!求你閱立我的諫言.
以施行公義.
斷絕罪過.
以憐憫窮人.
除掉罪孽.
或許你的平安可以延長.
其實這個噩夢.
不是很難解的.
只是夠不夠膽有人解.
剛才講到.
尼泊爾王發夢.
看到一棵大樹被斬.
還有一個新的心.
一個野獸的心.
其實那些事都是很明顯的.

$^{361}$不過.
發夢的人不敢想像.
也不敢推測.
那是自己.
是關於自己的夢.
在朝中的迦勒底人.
如果.
我們不懂得解夢.
我們直接解.
我們都知道是關於尼泊爾王的.
不過.
大概沒有什麼人夠膽.
將夢的真實情況.
告訴尼泊爾王.
更何況.
面對起路無常的.
尼泊爾王.
說一些不中聽的.
隨時朝野殺身之禍.
但去到但以理.
但以理就將夢的現象.
很直接地對他說.
在經文裡.
但以理當他聽到夢的時候.
他自己也有驚訝.
當他要解釋夢的時候.
他自己也有面有難色.
不過尼泊爾王.
在尼泊爾里面有難色.
他就說你儘管將夢的解釋告訴我.
給他有信心.
給他可以放膽.
將夢的信息展現出來.
但以理說的時候.
這個夢的內容.
很清晰.
這棵在夢裡.
參天的大樹.
就是尼泊爾王自己.
他將要被砍伐下來.

$^{401}$那個砍伐的意思.
就是他要離開他的王位.
他的國位.
將要.
某程度要斷絕.
甚至斷絕.
一個這樣的程度.
就是他要離開世人.
和野地的野獸一起同居.
好像那些野獸一樣.
吃草.
經常被下雨.
不會有瓦遮頭.
下雨的時候就濕透了身體.
毛髮會濕透.
經過一個完全的日子.
七期.
完完全全的一個日子.
才回到國中.
來掌權.
在這個夢裡.
特別的針對.
尼泊爾王.
他自己.
一個昔日的霸主.
一個很顯赫.
整個中東地區裡.
一個軍事強人.
去到這一刻.
說白一點.
是一個很折墮的一個境地.
沒理由一個霸主會面對一個.
這麼折墮的一個.
這樣的狀態.
但以你講.
這些的經歷.
不是人為做給他的.
而是天上的主.
來自至高者的命令.
在說.

$^{441}$耶和華上帝.
他要將這個命令.
真的很強力地實施.
在一個當時來說.
沒有人敢挑戰他的.
一個霸王面前.
去到一個.
至高者的面前.
地上任何有能力.
有權力的.
甚至乎每一個人都.
懼怕的尼泊爾夾尼沙王.
但是在上主面前.
他都要屈服.
在上帝的工作.
底下.
這個夢裡要讓他知道.
究竟這個世界裡.
真正掌權的是誰.
是尼泊爾夾尼沙.
還是天上.
至高者的命令.
不過說完這番說話之後.
尼泊爾夾尼沙.
有什麼反應呢.
在經文說.
沒有反應.
經過一年之後.
這個夢才實現出去.
由聽完.
但以你所說的.
這番說話.
到夢的實踐.
在聖經裡特別交代.
有足足一年的時間.
在這一年裡發生什麼事呢.
大概不是說.
尼泊爾夾尼沙王悔改.
然後安著.
但以你的勸諫.

$^{481}$以公義斷絕罪過.
以憐憫窮人除掉罪孽.
來換取更長的平安.
他沒有這樣做.
反而是繼續的來顧我.
在裡面繼續的.
來過他自己奢華滿足的生活.
大家看到.
在第四章裡面.
第28節.
這時都臨到尼泊爾夾尼沙王.
過了十二個月.
他遊行在巴比倫的皇宮裡.
遊行不是舉牌那些.
我再一次說.
散步在一個.
令他很滿足的巴比倫皇宮裡.
30節.
他說這個大巴比倫.
不是用我大能大力.
見為京都.
要顯我威嚴的榮耀嗎.
這裡你再看到.
經歷過這麼大惡夢.
十二個月之後.
尼泊爾夾尼沙王繼續的.
都是帶著囂張不惑的心.
在他的空中花園.
在他的行宮裡.
一路走的時候心裡很滿足.
這個這麼美的皇宮.
不是用我自己的雙手.
用我自己的大能大力.
建基一個新的巴比倫帝國.
令到這個國家強盛.
有一個這麼美好的皇宮.
都是靠我這一雙手.
這雙手就是.
令我要顯出我威嚴的榮耀.
是不是很囂張很威風.

$^{521}$站起來.
全部都是矮人.
自己覺得自己是全世界.
I am the king of the world.
就是全世界的霸主.
不過在聖經裡很緊湊的.
31節.
這話在王的口中尚未講完.
這句話都未講完.
榮耀可能都未講完.
整句完整的句子.
有聲音從天降下說.
尼布甲尼撒有話對你說.
你的國位離開你了.
你必被趕出離開世人.
與野地的獸同居吃草如牛.
且要經過七期.
等你知道至高者在人的國中掌權.
要將國賜與誰就賜與誰.
當時這話就應驗在尼布甲尼撒的身上.
他被趕出離開世人.
吃草如牛.
身被天怒滴濕.
頭髮長長.
口像鸚鵡.
指甲長長.
如同鷹爪.
這裡你看到那個訊息.
那個實現.
第28節開始.
到第32節就講到.
這個夢就實踐出來了.
一個最顛峰的狀態.
最囂張不惑的那一刻.
最狂妄的那句話還沒講完的時候.
天上的使者就將一句審判的說話.
淋到在尼布甲尼撒王的身上.
然後聖經沒有描述他的心靈狀態.
只是在講他的外貌.
他被趕出去世人.

$^{561}$然後就開始像野獸一樣.
手指甲長了.
身上的毛髮長了.
就像野獸一樣.
講回剛才夢裡的同樣說話.
再一次重複在夢裡的實現當中.
今天我們比較掌握多一點.
大概尼布甲尼撒王.
那一刻聽到這句說話.
這句審判之後.
他就患了一種精神病.
今天我們用這些角度來看.
當時不會有精神病.
可能羅莫斯知道.
以前在醫院會見過.
那叫變獸症.
變獸症是一種心理很嚴重的病.
近來今天還有出現這樣的症狀.
就是一個人的心智突然改變了.
腦裡突然不覺得自己是人.
他覺得自己是一隻野獸.
我特別去看過變獸病的外國研究.
有些人真的是這樣的.
我看過一個這樣的文章.
或者相片.
突然有一天他覺得自己是一隻豹.
然後他就模仿豹來到地上走.
然後他就去紋身.
把自己的臉紋成豹的樣子.
然後就吃生的肉.
覺得自己是一隻豹.
不是人.
他就繼續這樣過生活.
就是一個豹人混合體.
在這裡.
大概尼泊爾格尼薩王.
都是一個這樣的精神狀態.
聽完天上的使者審判的說話之後.
他就變成了自己.
成為一隻野獸.

$^{601}$是一隻怎樣的野獸我不知道.
不過他的身體外貌.
就是毛髮變長了.
他的手指甲像鷹爪一樣長.
而且他脫離了自己的王位.
你一隻野獸你是白手指王嗎?.
你不會做到平時.
他主政的時候.
那些人可以做到的事.
他就去吃的東西.
都像牛一樣來吃草.
就過著野獸的生活.
一個囂張跋扈.
一個生命裡面去到頂峰的王.
一句話.
突然之間就扭轉.
成為一個我們覺得很.
很折墮的狀態.
他做不到一個王.
他只能做到一隻.
像坊間所說的野獸一樣.
這個夢就實現了在尼泊爾格尼薩王的身上.
這個不是尼泊爾格尼薩王的結局.
經文裡面再繼續的說.
這種變了獸的.
覺得自己是野獸的時期.
但這個時期是多長呢?.
沒有人知道.
有些聖經學家在這裡計算.
大概可能七至十年的時間.
七至十年.
人生當自己是一隻野獸.
是一個很痛苦的經歷.
在七至十年之後.
四章三十六至三十七節.
這個王回復了自己的神智.
他就這樣說.
那時我的聰明復歸於我.
為我的國榮耀威嚴和光耀.
都復歸於我.

$^{641}$並且我的謀士和大臣都來朝見我.
我尤得建立在國位上.
之大的權柄加增於我.
現在我尼布甲尼撒讚美尊崇.
恭敬天上的王.
因為他所做的全都誠實.
他所行的也都公平.
那行驕傲的他能降為卑.
這裡面說到日子滿了.
尼布甲尼撒王回復自己的神智.
做回一個王.
雖然他還是繼續覺得.
國位榮耀還是回復於我身上.
不過這個囂張跋扈的王.
和昔日的表現有點不同.
同樣都是說現在我尼布甲尼撒.
不過這次說現在我尼布甲尼撒.
在第三十七節裡說到讚美尊崇.
恭敬天上的王.
因為他所做的全都誠實.
他所行的也都公平.
那行驕傲的他能降為卑.
一個囂張跋扈的人.
在這一刻學會入門謙卑.
體會到天上還有一個掌權者.
而且這個掌權者不是要糟蹋他.
而是說天上的掌權者是公平正直誠實.
他甘心樂意回望自己過去這段經歷的時候.
他將最大的尊榮給天上的王.
天上的王這個字在舊約只出現一次.
是一個最大最勁的最高級數.
讚美上帝的字眼.
在舊約只出現這一次.
不過尼布甲尼撒王將一個最大的榮耀.
他要敬拜天上的王.
更重要的是說到那行動驕傲的他能降為卑.
這大概是他為自己這幾十年秉政.
甚至這幾十年在人生高峰再降為野獸之後.
他為自己生命作出一個小小的結論.
天上的王能夠將一個驕傲的人變回卑微.

$^{681}$整件事我們大概理解了經文.
那個夢是怎樣.
不複雜的.
在這裡我想就驕傲這個主題.
我們有幾個層面重新思考一下.
第一.
尼布甲尼撒特別說到.
在他自己一個這麼顯赫的經歷裡.
沒有人可以抵擋他.
但只有天上的一個神可以抵擋他.
甚至天上的神可以將一個這麼大權力的.
可以將他降卑.
降到最卑微的一種狀態.
上帝在我們生命裡.
他有完全的掌權.
可能我今天或我們今天生命裡處於順或逆.
但如果上帝所愛的 上帝所看顧的那個人.
或者上帝要施行的公義.
沒有一件事可以難阻他.
在一個囂張跋扈的人當中.
上帝會用更加極端的方法.
來對付這些囂張跋扈的.
那種生命性情的那些人.
當我們去讀 去看尼布甲尼撒王的故事的時候.
我們心裡可能有一種痛快.
之前那麼壞 糟蹋猶太 糟蹋但以理.
糟蹋他們三個朋友.
現在知道了.
我們心裡讀上去有一種痛快的感覺.
我覺得這些驕傲人活該.
壞格.
我們看電影喜歡看這些.
大奸國面對著最老倒的時候.
我們心裡會透心涼的那種感覺.
不過雖然上帝會將報應放在驕傲的人身上.
或者面對著一些囂張跋扈的人.
上帝不會掩眼不看 坐視不理.
只是什麼時候去實踐上帝的那種能力.
不過在一個上帝主權底下.
一個那麼囂張跋扈的人.

$^{721}$我們身邊其他這些人我們怎麼看這件事呢.
我們是不是樂於.
覺得看你什麼時候知道.
看你怎麼死.
我們心裡有時候可能會有這樣的聲音.
不過當我們去看但以理的時候.
你會發現當但以理為這個王解夢.
或者他知道上帝.
如果他認定上帝的主權能力在這個王身上的時候.
但以理表現出來的說話.
和那種勸勉是表現得很厚度的.
他一來可能在強權之下.
他都避免自己朝野殺身之禍.
但當他知道自己的命是屬於上帝.
面對上帝有一個審判能力在這個人身上的時候.
但以理不是冷冰冰地說.
你死定了 你在做壽了.
而是他製造一個空間.
讓這個人可以有回轉的空間.
他提醒他.
如果你想延長你的平安.
你就要施行公義斷絕罪過.
憐憫窮人 除掉罪孽.
在一個囂張不惑的.
在上帝的審判將要臨到的人身上.
但以理帶著一種厚度的心.
面對一個將要面對上帝審判的霸主.
他為他的生命提供一些出路.
在這裡我們值得他再想.
我們今天都咬牙切齒.
我們看到一些囂張不惑.
有些我們可能工作場景也好.
或者在世界不同的地方也好.
我們看到有些人是這樣.
心裡面甚至是咒詛他.
不過這裡聖經似乎讓我們看到.
在他面對上帝的審判以來.
他的心或者他所呈現出來的那種態度.
不是一種義德.
看到一個囂張不惑的人.

$^{761}$被置於死地.
反而是他希望那個人還有回轉的空間.
心靈裡面能夠再闊一些.
承載一些不太好看的.
我們不太認同的一些人.
我們都希望天父是管教他.
幫助他回到正軌.
而不是看他落在一個無底坑的地獄裡受盡折磨.
今天我們基督徒.
我們今天身處這個世界社會裡.
我想我們都要學習這種思考的方式.
沒錯,看到有些人做得不好.
我們可能心裡面有一些負面的咒詛的話.
但我們祈求神幫助我們.
讓我們的心胸可以有多一些的厚度.
面對過往做得不好的人.
我們不是將他置諸於死地.
或者我們幸災樂禍的心.
來看這些人怎樣折墮.
祈求主幫助我們.
讓我們的心不要成為另外一個野獸.
不要讓我們成為另外一個尼布甲尼薩.
看到那些藥或者那些人被攻打之後.
我們的心就紅起來.
或者覺得自己永遠站在上帝那邊.
可能當我們的心裡面這樣想的時候.
我們不知不覺可能我們都成為上帝下一個要管教的對象.
這是第一點.
在經文裡面說到上帝是會施行公義.
行動驕傲的他要降為卑.
但同時間在我們面對這些人的時候.
我們要學習厚度一點.
為他的生命回轉來祈求.
而不是要咒詛他.
第二點是《箴言》第十六章第十八節.
驕傲在敗壞已先.
狂心在跌倒之前.
我們從小讀到大.
驕兵必敗.
我們都知道敗壞失敗的人.

$^{801}$是在驕傲未失敗就先驕傲.
我們明白這一點.
到最後我們的態度決定了我們做事行為的結果.
可能今天我們會做得好.
但走得不遠.
當我們帶著一種驕傲狂心的時候.
到最後就是為自己埋下了一個陷阱.
讓自己泥足深陷地陷下去.
驕傲其實在《屬靈生命》裡面說的是什麼呢?.
不只是那些渾唇的人才驕傲.
驕傲的本質就是高估自己的能力.
在上帝的面前.
或者在真誠面對面地看自己的時候.
我永遠只覺得自己比別人高人一等.
我永遠覺得自己是鶴立雞群.
覺得自己是全世界最厲害的一個人.
有些人可能帶著這樣的心態成功.
有些人可能繼續帶著這樣的心態失敗.
但成功失敗不是驕傲的條件.
我怎樣看自己才是真正構成了我驕不驕傲.
一個人不能夠正確地看自己的生命.
我們到最後都會製造一種失敗的經歷.
而這種失敗不只是做一件事做得不好.
而是更加在敗壞,在罪裡面面對的失敗.
在試探裡面一次又一次地跌倒.
我曾經說過很多例子.
我對著年輕人也好.
對著一些成年,單身的弟兄姊妹.
我經常都說一個這樣的例子.
很多的年輕弟兄姊妹現在流行的.
單獨一男一女兩個還沒結婚拍拖.
就單獨去旅行.
這個世界很平常.
每個人都是這樣的.
我每一次不厭其煩都會說.
我們作為基督徒.
或者中國人社會裡面.
我們都不要單獨一男一女兩個去旅行.
你進去旅行你自己住在一個房間裡.
你會不會一男一女分房睡.

$^{841}$你告訴我實際上你都不會這樣做.
更加重要的是.
人面對試探.
男與女都好.
在這樣的旅行.
在這樣的狀態之下.
你面對的試探是很大.
你會隨時分前行為.
隨時你都會為自己帶來一些很悔舊的生活.
很多弟兄姊妹.
一些年輕的弟兄姊妹質問我.
牧師你這麼老套.
你又不相信我.
我不是不相信你.
我是不相信全世界的男人和女人.
我特別是不相信弟兄.
我告訴他.
如果弟兄晚上你跟他在房間裡.
你跟他說兩個孤男寡女在房間裡查經.
我真的不太相信.
女孩子又查護法數.
我經常說護法數一發.
頭髮的味道.
男孩子自動.
賀爾蒙就上了.
與其說我不相信你.
倒不如說全世界的男人和女孩子.
我都不相信.
女孩子你不要想著就好了.
女孩子見到男孩子.
有時候手臂很乖.
明明沒有什麼.
老鼠仔都覺得很大隻等等.
情人眼裡出西施.
你不吸引就不會走在一起.
隨時面對一種很大的試探.
很多弟兄姊妹.
經常覺得我老套.
不過我說的是.
我們面對罪的時候.

$^{881}$不要高估自己可以抵抗.
這些人有.
或者這些陷阱的能力.
事實上.
不要說黃浸.
我認識很多年輕的弟兄姊妹.
當我和他這樣勸勉.
他覺得我很懦弱的時候.
沒多久三個月回來.
兩個人去三個人回來.
大家明白我的意思吧.
但是三個人回來.
那個都不是最重點.
而是為什麼不相信人是軟弱的.
為什麼這麼高估自己抵抗.
面對罪的那種能力.
覺得自己免疫的對於這個罪.
當我們不看清楚自己能夠.
不能夠抵抗罪的那種人有.
我們就要好好地.
逃避這些少年人的私慾.
要好好地逃避眼前面對的.
這些人有.
我們才能夠一路一路.
走在上帝的旨意當中.
《箴言》的作者.
或者到今天我們讀的這篇經文.
聖經裡面大量的篇幅提醒我們.
不要太過高估自己.
在地上裡面.
你自己抵抗罪惡的能力.
不要覺得自己在罪裡面永遠免疫.
我們都是謹慎地.
令自己走在上帝的旨意裡面.
我們的屬靈操練.
是幫助我們更加謹慎地過我們的生活.
而不是滿足了自己.
那種純粹感覺好像很親密.
跟別人較勁的那種態度.
所以第二點提醒我們.

$^{921}$不要自視過高.
驕傲會帶來失敗.
我們生命裡面.
不斷不斷地提醒自己.
我們是否陷在這種狀態當中.
第三個提醒是雅各書.
去到新約裡面.
神阻擋驕傲的人.
賜恩給謙卑的人.
在這裡特別講到.
上帝的恩典.
是淋到在那些.
真誠謙卑的人身上.
驕傲的人是上帝阻隔他.
難了他.
為什麼這樣做呢?.
因為上帝自己都謙卑.
上帝的救恩都是用最謙卑.
最卑微的方式.
十字架上面的耶穌就是在講.
這個世界裡面最大的患難.
直著神的愛子耶穌.
用最謙卑.
最受害的.
最痛苦的方式來呈現出來.
當我們知道神的樣式是這樣的時候.
我們沒理由會想.
我們這個信仰.
神就受苦我們就享福.
不是這樣的信仰形態.
上帝反而一路又一路呼籲我們.
學習基督耶穌的謙卑樣式.
當然我們未必去到十字架.
耶穌基督所承受的款額.
但至少我們所過的生活.
我們的心思意念.
絕對不會和耶穌基督的救恩背道而行.
反而是我們仰望著這個.
在十字架上面那位謙卑的耶穌.
勉勵自己一路一路來走這條信仰的路.

$^{961}$更加重要的是.
謙卑和驕傲.
謙卑永遠看著神的樣式.
而驕傲就是看著撒旦的形象.
撒旦做大量的工作.
都是令人的自我越來越膨脹.
覺得自己越來越滿有力量.
在耶穌面對撒旦的試探的時候.
耶穌也是被撒旦叫去殿頂跳下來.
滿足一下.
有神的使者托著你.
為何不去做呢?.
不過耶穌叫撒旦退後面去吧.
我們生命裡.
我們究竟是看著神的樣式來做人.
還是看著這個坊間.
或者這個世界.
甚至撒旦的形象來過我們的生活呢?.
當我們心靈裡充滿了驕傲的心.
真理就不能夠進入我們的心靈裡.
我們會覺得自己的東西都懂得了.
明白了.
不過實際上.
原來我們心靈裡貧乏得可憐.
我們依然只是頭腦上.
知道神的名字.
心裡沒辦法謙卑地跟隨他.
亞國書提醒我們.
真正跟隨主的人是謙卑的人.
不要容讓自己的自我感覺良好.
左右我們走信仰的路.
當我們越謙卑的時候.
我們越發現上帝的恩典.
當我們越驕傲的時候.
我們會發現.
越來越多的上帝的管教臨到我們身上.
到那時候.
如果我們再不去回轉的時候.
我們生命可能會一步一步離棄神.
進入一種上帝不喜悅的狀態.

$^{1001}$今天這段經文.
希望大家掌握到.
不是好心的道理.
但是在生命裡.
驕傲卻是潛藏在每一個人的心靈裡.
就讓我們今天讀好這段經文.
進入一種謙卑的心靈.
希望我們不是成為下一個尼布甲尼撒.
而是成為下一個的但爾尼.
厚道,謙卑.
在一個不明朗.
又或者充滿敵意的世代裡.
與上帝同行.
求主繼續來恩待我們.
\newpage



\section{但以理書 5:1-16}
\label{sec:1ZS5UTvrjn8}
\textbf{2020年10月18日 崇拜直播｜但以理書5\_1-16}
\newline
\newline
連結: \href{https://youtube.com/watch?v=1ZS5UTvrjn8}{\texttt{https://youtube.com/watch?v=1ZS5UTvrjn8}} ~~~~ 語音日期: 2020-10-18
\newline
\newline
\hyperref[sec:Jfw_Vv2ojxk]{\small{< < < PREV SERMON < < <}}
~
\hyperref[sec:index_chronic]{\small{[返順時目]}}
~
\hyperref[sec:YHKUQm4ko_M]{\small{> > > NEXT SERMON > > >}}
\newline
\newline
但以理書 5:1-16
\newline
\begin{longtable}{cl}
\hline
\hline
章節 & 經文 (和合本修訂版)\\
\hline
5:1 & \begin{tabularx}{0.7\textwidth}{X} 伯沙撒王為他的一千大臣擺設盛筵，與這一千人飲酒。 \end{tabularx} \\ \\ \relax
5:2 & \begin{tabularx}{0.7\textwidth}{X} 伯沙撒在歡飲之間，吩咐人將他父尼布甲尼撒從耶路撒冷聖殿所擄掠的金銀器皿拿來，好使王與大臣、王后、妃嬪用這器皿飲酒。 \end{tabularx} \\ \\ \relax
5:3 & \begin{tabularx}{0.7\textwidth}{X} 於是他們把聖殿，就是耶路撒冷神殿中所擄掠的金器皿拿來，王和大臣、王后、妃嬪就用這器皿飲酒。 \end{tabularx} \\ \\ \relax
5:4 & \begin{tabularx}{0.7\textwidth}{X} 他們飲酒，讚美金、銀、銅、鐵、木、石造的神明。 \end{tabularx} \\ \\ \relax
5:5 & \begin{tabularx}{0.7\textwidth}{X} 當時，忽然有人的指頭出現，在燈臺對面王宮粉刷的牆上寫字。王看見寫字的指頭， \end{tabularx} \\ \\ \relax
5:6 & \begin{tabularx}{0.7\textwidth}{X} 就變了臉色，心意驚惶，腰骨好像脫節，雙膝彼此相碰， \end{tabularx} \\ \\ \relax
5:7 & \begin{tabularx}{0.7\textwidth}{X} 大聲吩咐將巫師、迦勒底人和觀兆的領進來。王對巴比倫的智慧人說：「誰能讀這文字，並且向我講解它的意思，他必身穿紫袍，項帶金鏈，在我國中位列第三。」 \end{tabularx} \\ \\ \relax
5:8 & \begin{tabularx}{0.7\textwidth}{X} 於是王所有的智慧人都進前來，他們卻不能讀那文字，也不能為王講解它的意思。 \end{tabularx} \\ \\ \relax
5:9 & \begin{tabularx}{0.7\textwidth}{X} 伯沙撒王就甚驚惶，臉色改變，他的大臣也都困惑。 \end{tabularx} \\ \\ \relax
5:10 & \begin{tabularx}{0.7\textwidth}{X} 太后因王和他大臣所說的話，就進入宴會廳，說：「願王萬歲！你的心不要驚惶，臉不要變色。 \end{tabularx} \\ \\ \relax
5:11 & \begin{tabularx}{0.7\textwidth}{X} 在你國中有一人，他裡頭有神聖神明的靈，你父在世的日子，這人心中光明，又有聰明智慧，好像神明的智慧。你父尼布甲尼撒王，就是王的父，曾立他為術士、巫師、迦勒底人和觀兆者的領袖， \end{tabularx} \\ \\ \relax
5:12 & \begin{tabularx}{0.7\textwidth}{X} 都因為他有美好的靈性，又有知識聰明，能解夢，釋謎語，解疑惑。這人名叫但以理，尼布甲尼撒王又稱他為伯提沙撒，現在可以召他來，他必解明這意思。」 \end{tabularx} \\ \\ \relax
5:13 & \begin{tabularx}{0.7\textwidth}{X} 於是但以理被領到王面前。王問但以理說：「你就是我父王從猶大帶來、被擄的猶大人但以理嗎？ \end{tabularx} \\ \\ \relax
5:14 & \begin{tabularx}{0.7\textwidth}{X} 我聽說你裡頭有神明的靈，心中有光，又有聰明和高超的智慧。 \end{tabularx} \\ \\ \relax
5:15 & \begin{tabularx}{0.7\textwidth}{X} 現在智慧人和巫師都被帶到我面前，要叫他們讀這文字，為我講解它的意思；無奈他們都不能講解它的意思。 \end{tabularx} \\ \\ \relax
5:16 & \begin{tabularx}{0.7\textwidth}{X} 我聽說你能講解，能解疑惑；現在你若能讀這文字，為我講解它的意思，就必身穿紫袍，項戴金鏈，在我國中位列第三。」 \end{tabularx} \\ \\
[1ex]
\hline
\hline
\end{longtable}
$^{1}$今日經訓是舊約聖經的.
《但義理書》第五章一至十六節.
舊約聖經《但義理書》第五章.
它是這樣說.
「伯沙薩王為他的一千大臣設擺筵席.
與這一千人對面飲酒.
伯沙薩歡飲之間.
吩咐人將他父或役祖彌布加尼沙.
從耶路撒冷中所掠的金銀器皿拿來.
王與大臣王后妃嬪好用這器皿飲酒.
於是他們把耶路撒冷神殿庫房中.
所掠的金器皿拿來.
王和大臣王后妃嬪就用這器皿飲酒.
他們飲酒讚美金銀銅鐵木石所造的神.
當時忽有人的紙頭顯出.
在王宮與燈台面相面的墳場上寫字.
王看見寫字的紙頭.
便便了臉色 心意驚慌.
腰骨好像脫節 相失彼此相碰.
大聲吩咐將用法術的和迦勒底人.
並官少的領進來.
對巴比倫的哲士說.
誰能讀這文字.
把講解告訴我.
他必身穿紙袍.
恆大今年在我國中位列第三.
於是王的一切哲士都進來.
卻不能讀那文字.
也不能把講解告訴王.
怕沙剎王就神驚惶.
面色改變.
他的大臣也都驚奇.
太后或亦王后.
因王和他大臣所說的話.
就進入彥公說.
願王萬歲.
你心意不要驚惶.
不要變色.
在你國中有一人.
他裡頭有聖神的靈.

$^{41}$你父在逝的日子.
這人心中光明.
又有聰明智慧.
好像神的智慧.
你父尼布格尼薩王.
就是王的父.
納他為術士.
用法術的和加勒底人.
並官少的領袖.
在他裡頭.
有美好的靈性.
又有知識 聰明.
能圓夢 識謎語.
解疑惑.
這人名叫.
但爾里.
尼布格尼薩王.
又稱他為.
帕提沙剎.
現在可以召他來.
他必解明這意思.
第十三節.
但爾里就被領到王前.
王問但爾里說.
你是被擄之猶大人.
中的但爾里嗎?.
就是我王.
從猶大老來的嗎?.
我聽說你裡頭有神的靈.
心中光明.
又有聰明和美好的智慧.
現在之士和用法術的.
都領到我面前.
為叫他們讀者文字.
把講解告訴我.
無奈他們都不能把.
講解說出來.
我聽說你善於講解.
能解疑惑.
現在你若能讀者文字.

$^{81}$把講解告訴我.
就必身穿紙袍.
恆大今年.
在我國中位列第三.
(梁:姐妹平安).
(梁:我們不經不覺已經來到).
(但爾里書第五章).
(剛才你聽到出現一個新的名字).
(叫做白沙剎).
(讀者到底是誰呢?).
(正確來說).
(他是尼泊爾格尼薩王的外孫).
(因為在亞蘭文).
(基本上沒有孫這個字).
(無論是父子).
(或者是祖父和孫).
(或者外祖父和孫).
(都可以用兒子來相稱).
(事實上).
(尼泊爾格尼薩王).
(在公元562年離世).
(之後他的王位輾轉).
(落到其中一位女婿的手上).
(叫做拿波利道).
(白沙剎就是拿波利道的兒子).
(這兩父子在公元前的550年).
(一起執政).
(拿波利道長住在巴比倫的第二首都).
(一個叫提瑪城的地方).
(而今日經文所說的白沙剎王).
(主要在巴比倫的首都).
(巴比倫城內來攝政).
(所以剛才我們聽到).
(位列第三這幾個字).
(意思是誰能夠講解到).
(牆上文字的意思).
(就能成為巴比倫國的第三把交椅).
(地位僅次於拿波利道和白沙剎).
(經文開首提到).
(白沙剎王擺設了一場千人的盛宴).

$^{121}$(結尾我們未讀到).
(不過結尾提到他那一晚被殺).
(其實).
(當時我們說有另一個大國崛起).
(叫馬代波斯).
(經文的結尾提到白沙剎王被殺).
(而巴比倫國亦被馬代大利烏取代).
(誰是大利烏呢?).
(其實你很熟悉的另一個名字).
(就是波斯王古烈).
(被他取代).
(事實上有文獻記載).
(在巴比倫亡國前幾天).
(波斯的大軍).
(已經在距離巴比倫城不遠的地方).
(擊敗了白沙剎的父親).
(拿破利道的軍隊).
(下一步準備進攻巴比倫城).
(攻打白沙剎).
(這是故事當時的歷史背景).
(不過我們不知道到底是什麼原因).
(白沙剎要舉行這一場千人的盛宴).
(可能是為了壯膽).
(又或者為了穩定人的心).
(又或者他對巴比倫城的城牆).
(覺得很堅固很穩固).
(而且有護城河來保護他).
(所以即使他在國難當前).
(即使敵人已經向他劍拔弩張).
(他竟然還有興致).
(和他的群臣去暢飲).
(不過我們常說).
(酒能夠聯性).
(你看到白沙剎喝著喝著).
(應該喝大了).
(他吩咐人將尼泊爾加利沙皇).
(從前在耶路撒冷聖殿).
(努力取得的金銀器皿).
(拿出來用來和大臣和皇后妃嬪).
(用來喝酒).

$^{161}$(不單止這樣).
(還一邊喝一邊讚美).
(人手所做的金銀和鐵石木).
(所做的偶像).
(即是怎樣?).
(當以色列的上帝無道).
(要公然地褻瀆他).
(就好像你手頭上的這幅名畫).
(這幅是荷蘭作家林布蘭).
(他在1635年的一幅作品).
(名為「白沙剎的盛宴」).
(白沙剎看到).
(有手的指頭在牆上出現).
(對著登台對面的牆上寫字).
(白沙剎當場嚇得).
(英文裡形容面青).
(又心驚又腳震).
(即是他知道這次大事不妙了).
(得罪了以色列的上帝).
(他不知道會有什麼後果).
(於是他召集了).
(所有法術或官少的巴比倫哲士).
(到底誰能夠講解到).
(牆上文字的意思).
(他想他紙袍今年).
(成為巴比倫國的第三把交椅).
(當然他以為重賞之下).
(必會有勇夫).
(不過很可惜).
(經文記載這班哲士).
(都不明白那些字的意思).
(有聖經學者估計).
(為何不明白呢?).
(很可能與那些字的寫法有關).
(因為亞蘭文).
(與英文不同).
(亞蘭文首先是).
(由右邊寫到左邊).
(英文不同).
(英文是左至右).

$^{201}$(是橫向寫).
(亞蘭文是右至左).
(聖經學者估計).
(他們不是用正常).
(右至左的橫向寫法).
(而是橫向寫).
(由上而下寫).
(那些文字變成了).
(一個謎語).
(一個密碼).
(所以他們看不明白).
(無論如何).
(這件事驚動了太后).
(也就是白沙薩的母親).
(尼布甲尼薩的女兒).
(她知道這件事).
(她就進宮安慰白沙薩).
(並向他舉薦一個人).
(她說這個人一定能幫助你).
(因為他才是巴比倫國).
(真正解夢解謎的專家和高手).
(她還說到).
(你外公尼布甲尼薩在位的時候).
(覺得這個人心中有光明).
(擁有像神一樣的智慧).
(無論多麼隱密的夢想).
(那些異象).
(他都有本事來拆解).
(所以她的外公).
(昔日很重用她).
(並且派她管理).
(一切巴比倫的哲士).
(這個人就是但以理).
(白沙薩王大可以召他入宮).
(他一定能幫你解答).
(牆上文字的意思).
(於是王就把但以理召來).
(但你有否留意).
(他對但以理的態度?).
(從他的對話).

$^{241}$(你會看到他擺出一副).
(頗傲慢).
(看不起但以理的態度).
(他就說).
(你就是我外公).
(從猶大努力回來的猶太人).
(但以理嗎?).
(很不客氣).
(我聽說).
(聽說).
(我聽說).
(你的女兒有神的靈).
(看事物很通透).
(你又聰明又有智慧).
(現在我一班專家都幫不了).
(解答不了牆上文字的意思).
(所以唯有找你來試試).
(我聽說你很善於解夢).
(很擅長解決人的難題).
(現在跟你做一個交易).
(你如果能夠).
(解讀到牆上的意思).
(我就想你).
(紙袍今年位列第三).
(你留意).
(白沙殺沒有提到什麼).
(他沒有提到).
(什麼原因).
(為什麼牆上會有字).
(他沒有交代到).
(其實是自己喝大了).
(其實是自己失敗了).
(闖了這個大禍).
(得罪了上帝).
(他隻字不提).
(你亦看到).
(但以你在第十七節).
(他亦不賣帳).
(他沒有對這位).
(心高氣傲的白沙殺賣帳).

$^{281}$(他說什麼).
(我們一起看那幾節經文).
(簡單說第十七節).
(黃,你的禮物可以怎樣).
(可以省下).
(你可以獎賞給別人).
(留給其他人).
(不過他一定會替白沙殺).
(來解答這些文字的意思).
(我想大家一起讀).
(第十八至二十四節).
(好嗎?).
(預備起).
(好,謝謝).
(有沒有印象).
(裡面提到一些經歷).
(其實我們上星期已經看了).
(其實上面提到第四章).
(尼布甲尼撒的一些經歷).
(但你叫白沙殺).
(你回想你外公以前發生過什麼事).
(上帝把國位,權力).
(威榮賜給尼布甲尼撒).
(但他卻心高氣傲,妄自尊大).
(於是上帝將他趕離了皇位).
(甚至令他失去了常性).
(好像野獸般過活).
(直至管教的日子滿了).
(上帝恢復了尼布甲尼撒的人性).
(也恢復了國位).
(從此尼布甲尼撒學會了).
(什麼叫謙卑).
(不再敢在上帝面前狂傲自誇).
(因為他知道真正在人間掌權的).
(不是皇帝,而是上帝).
(所以他借著這個歷史).
(借著白沙殺外公的經歷).
(來指證白沙殺是明知故犯).
(自然是罪加一等).
(白沙殺既然知道這件事).

$^{321}$(很應該吸取教訓).
(但他不但不懂得謙卑).
(他比他的外公更加狂妄).
(連他外公都不敢做的事).
(他就敢做).
(就是拿上帝在聖殿的器皿).
(出來喝酒,不但助興).
(還要專門讚美其他偶像).
(不把榮耀歸給那位).
(掌握著生命氣息).
(那位其實是掌握著前途的上帝).
(簡直是公然的來敵黨).
(挑戰這位真神).
(所以上帝也沒有客氣).
(直接在牆上寫下刑罰).
(讓白沙殺或讀到這段經文的人).
(都知道上帝是興萬不得).
(無論你信不信祂).
(認不認識祂).
(祂也是那位真神).
(至高的人).
(最終必定會被這位至高的上帝).
(降為卑).
(甚至降為最卑微的一格).
(第25至31節).
(但義理開始解讀牆上文字的意思).
(所寫的文字是).
(利利,尼尼,提克勒).
(烏法爾辛).
(講解是這樣的).
(利利就是神已經數算).
(你各的年日到此為止).
(提克勒就是你被稱在天平上).
(顯出你的虧欠).
(皮勒斯就是你的國分裂).
(歸與馬代人和波斯人).
(白沙殺下令人就把紙袍).
(給但義理穿上).
(把金鏈給他戴在頸項上).
(又傳令使他在各中位列).

$^{361}$(第三,當夜加勒底王白沙殺被殺).
(馬代人戴利烏年六十二歲).
(取了加勒底國).
首先和大家一起了解一下.
那幾個文字是甚麼.
甚麼是利利,尼尼,提克勒,烏法爾辛.
下一張,謝謝.
其實很簡單.
烏法爾辛有個「烏」字.
其實是亞蘭文的連接詞.
你當作中文的「和」字.
或者英文的「and」.
意思就是利利,尼尼和法爾辛.
而這三個字.
你可以理解為三個重量單位或貨幣的名詞.
但義理可以將它們看作動詞.
將利利解釋為數算.
提克勒解釋為稱一稱.
至於烏法爾辛即是和,分裂.
所以你見到是數算,數算,稱一稱和分裂.
這個意思.
為何會出現「皮勒斯」這個字呢?.
因為「皮勒斯」這個字是法爾辛的單數寫法.
並且這兩個字很容易令人聯想起波斯的串法.
我用最簡單的方法.
嘗試將一些很複雜奇怪.
其實我都不明白.
是但義理明白的東西告訴大家.
所以最終長尚文字但義理解為.
神已經數算白沙薩你執掌王權的日子.
到此為止.
你亦被上帝稱在天平上.
顯出你對上帝的虧欠.
就如他剛才拿聖殿器皿來喝酒.
和當著上帝面前敬拜其他偶像.
這個虧欠.
所以你的國也要被分裂.
歸給馬代人和波斯人.
這個審判那晚就應驗了.
結果白沙薩被殺.

$^{401}$馬代人大利烏也就是波斯王古烈.
取締了巴比倫國.
整個故事講到這裡.
你看到白沙薩的下場.
讓我們看到一個人無論他身處的地位有多高.
無論他的權力有多大.
甚至他管治一個多宏大的國家或機構.
都敵不過上帝.
人能夠因為擁有自己的權力自由.
就以為可以目空一切.
任意妄為.
人必須要明白一件事情.
原來一切的權力和自由.
不是自己賺回來的.
亦不是其他人給他的.
而是來自上帝的賞賜.
所以人就要苦心自問.
到底我有沒有用正確的方法.
來運用上帝賦予我們的權利和自由呢?.
否則可能就會好像這個故事那樣.
上帝既然有權將自由和權力賜給人使用.
他同樣也有權收回.
不讓你享用.
你看到白沙薩的行為.
某程度上反映了人兩個的生活盲點.
我們有時會自由運用我們的權利.
但做了些錯的事.
你和我可能都有很多這些經歷.
自由運用選擇的權利.
不過做了些錯的抉擇.
這是第一方面.
第二就是人有時有個盲點.
就是認為我喜歡做的事.
就一定是些對的,正確的事.
這不單是白沙薩的盲點.
我想亦是現代人經常踩地雷的盲點.
我喜歡做的事.
我覺得就應該是些對的,正確的事.
但最終往往會讓人走上一條自毀.
甚至滅亡的路.

$^{441}$白沙薩成為了我們的借鏡.
當然我們又未去到白沙薩那一格.
那麼狂妄.
我相信我們搞外賢也好.
我們不夠杯喝東西也好.
我們都不會用主餐的杯來用.
對不對?.
你不會那麼夠膽.
我們不會這樣.
但我想我們都試過.
甚麼叫做「明知不過固犯」.
我想我們都試過.
不用上帝給我們的權利或自由.
來做一些損害自己,損害別人.
又或者讓上帝受虧損的事.
多多少少我們都試過.
特別在講求.
我們說今天很看重那個自由.
很看重那個個人主義.
對不對?.
你和我都很崇尚.
只要不妨礙人,不傷害人.
每一個人都有權利選擇自己喜歡做的事.
對不對?.
聽起來其實很受落.
我有權去選擇我的錢怎樣用.
我有權去選擇.
我和誰拍拖,結婚.
我有權去選擇上網看甚麼.
我有權去選擇去哪間教會.
我有權去選擇我吃甚麼,喝甚麼.
坦白說,弟兄姊妹.
問心,爽不爽?.
可以選擇.
大家靜下來.
我一問這個問題,大家靜下來.
可以選擇,好嗎?.
你反過來想.
我的錢我沒權選擇怎樣用.
我沒權選擇和誰拍拖,結婚.

$^{481}$很慘的.
可以選擇,真的挺爽的.
但你喝深一層.
在我們運用這些選擇自由的時候.
如果把上帝也抽走會怎樣?.
其實我們會製造.
不知不覺製造一個又一個白沙煞出來.
白沙煞也是這樣製造出來.
我抽走上帝.
運用我的自由來隨意使用我的權利.
慢慢就變成白沙煞.
最終我亦失去了.
我原本享有的權利和自由.
就譬如一些很簡單的例子.
我二十歲,三十歲的時候.
我飲食不需要怎樣節制.
你不要看我現在這麼橫.
我最初結婚的時候.
我瘦二十磅以上,比現在.
瀟灑很多.
喜歡吃甚麼都可以.
炸雞腿沒問題.
梅菜扣肉沒問題.
多邪惡的事都可以處理.
但當我繼續放縱我的自由.
喜歡每餐吃甚麼就吃甚麼的時候.
我可能會怎樣?.
很快就甚麼都吃不到.
我會損失很多吃不同的自由.
甚至可能會早點回天家.
對不對?.
你和我也有這樣的經歷.
所以人到中年.
就會很注意自己的健康飲食.
我們需要選擇限制自己的自由.
可能少吃鹹的食物.
高糖的食物.
高脂的食物.
少看電視.
多做運動.

$^{521}$用這些限制來換取更大的自由.
就是我可以留一個健康長壽的身軀.
來給上帝用.
來給你家人.
對不對?.
你選擇限制自己的自由.
來換取更大的自由.
聖經也說了這件事.
耶穌選擇限制自己.
選擇成為人.
來換取更大的自由.
所以弟兄姊妹.
聖經所說的自由.
從來不是做甚麼都可以.
而是讓自己接受適當的規限.
來換取更大的益處.
就是在上帝的規範和管教裡面.
正確地運用.
祂所賦予我們的自由和權柄.
來做永生益人的事.
意思是甚麼呢?.
我曾經說過一個比喻.
我們也有走過行人天橋.
行人天橋一定會有圍欄.
你有沒有試過貼近圍欄來走?.
有沒有?.
沒有?.
安全一點吧?.
不敢吧?.
或者反過來問.
你覺得天橋有圍欄好一點.
還是沒有圍欄好一點?.
覺得有圍欄好一點的舉手.
大部分都舉手.
無所謂吧?.
兩樣都可以.
坦白說.
有圍欄你會有膽子走到哪邊?.
你知道靠著也不會掉下去.
但如果你想想.

$^{561}$天橋沒有兩邊圍欄.
沒有限制怎麼辦?.
你敢不敢走到最邊?.
因為你明知走到最邊會很危險.
你有機會掉下去.
所以圍欄是否有限制?.
但亦是保護.
機械限制更是保護.
以至你可以放膽安心地.
走到最邊也不怕.
因為它保護你.
但相反沒有兩邊圍欄的限制和保護.
你走的地方就窄了.
你的自由其實是少了.
所以弟兄姊妹.
其實上帝的規範.
上帝的管教.
有時好像覺得很不中聽.
但上帝不是要欺壓我們.
不是要甚麼都管著.
拿走我們的好處.
相反上帝是要保護你和我.
幫我們拿到真正的自由.
更大的益處.
因為這樣.
上帝在適當的時候.
會出重手管教尼布格尼撒.
和今日我們看到的白沙撒.
你試想想.
上帝如果要消滅白沙撒.
用不著這麼煩惱嗎?.
上帝可以怎樣?.
上帝如果生氣得震的時候.
直接擊殺就可以了.
用不著搞那麼多.
又或者好一點.
不出聲.
等你今晚也過不了.
等你玩完最後那幾個小時.
然後就找人要你的命.

$^{601}$很簡單.
上帝要消滅白沙撒.
不用搞那麼多.
為何又需要在牆上露一手.
然後又要太后的介入.
然後又要宣召但以理入宮.
然後又要講述.
公尼布格尼撒的經歷.
最後才講宣判.
為何要搞那麼多?.
頂展維在我們想的過程中.
你不妨多想兩段經文.
其實你想想.
聖經裡不止單單是但以理書第五章記載.
上帝用手指頭寫東西.
有沒有印象?.
聖經裡有三次記載.
上帝用他的指頭來寫字.
第一次在哪裡?.
第一次在《昭雅及記》.
上帝用他的指頭在石板裡寫下.
他的法律,律例,誡命.
寫下他對人的要求.
上帝第一次用手指寫字.
是跟人納藥.
第二次.
就是今天你看到的這段經文.
上帝用他的手指在牆上寫下.
對得罪他的人的宣判.
他寫下人對他的虧欠.
第三次呢?.
第三次難想一點吧.
在新約.
在哪裡?.
在《約翰福音》第八章.
耶穌和行淫婦人的故事.
有沒有印象?.
有群人特意抓行淫婦人.
來到耶穌面前.
要求耶穌可以扔石頭嗎?.

$^{641}$解決她.
耶穌怎樣?.
彎下腰.
在地上寫字.
但沒有說寫什麼.
但最後你知道耶穌說了一句.
「你們誰沒有犯罪.
誰出來才扔石頭呢?」.
於是人們一個一個.
由老到少.
由老到少離開.
剩下的婦人.
我們不知道耶穌寫了什麼.
但我們知道最後耶穌.
寬恕了行淫婦人.
上帝成為肉身.
第三次在地上用指頭寫字.
顯示他對罪人願意寬恕,接納.
再問大家剛才的問題.
牆上的字.
太后的舉薦.
外公的借鑒.
最後但以利的宣判.
這一切是純粹出於偶然.
還是上帝仍然不斷給予機會白沙煞?.
仍然等他把握最後的機會.
好像他外公一樣.
懂得回頭?.
弟兄姊妹,很值得我們去想.
如果白沙煞願意回轉悔改.
整件事的結果會否有得拔?.
有得改變?.
有無得改變?.
可能即使悔改.
白沙煞也難逃一死.
因為他真的要承擔.
罪所帶來的後果.
既然如此,可能有人會問.
悔不悔改也死了.
橫死念死,不如怎樣?.

$^{681}$不如享受那幾個小時吧.
享受最後的權利和自由.
是不是這樣?.
但又不是,弟兄姊妹.
你要想一件事.
沒錯.
他願意向上帝悔改.
未必能改寫他當晚要被殺的結果.
但他一定可以改變.
他和上帝的關係.
不要只想那一刻.
想將來很遠的事.
他能夠確保將來可以盡到永恆.
這才重要.
正如你記得.
耶穌在釘十字架的時候.
旁邊有兩個囚犯.
就是最好的例子.
一個橫死念死,無所謂.
另一個呢?.
另一個懂得把握最後機會.
他最後要死嗎?.
他求耶穌紀念他.
他也要被釘死在十字架.
不過耶穌說.
「你今晚跟我去樂園」.
如果白沙剎能夠把握最後機會.
將來在永恆的裡面.
整個結局會改寫.
聖經裡說.
上帝不願一人尋倫.
乃願人人悔改.
不過可惜的是.
白沙剎願意出重賞.
來換取真相.
希望有人告訴他.
牆上寫了什麼.
解釋了什麼.
事實上他們都認同.
但以你的解讀.

$^{721}$否則他不會獎賞他.
但很可惜的是.
他知道了真相.
但他錯過了上帝.
給他最後迴轉的機會.
諷刺嗎?.
其實現實也會這樣.
我們很多事情都知道了.
不過不去跟進.
不去改變.
明知故犯.
所以,弟兄姊妹.
其實你細心看.
第四和第五章.
兩個王的故事.
其實很好對比.
兩個故事的題材很類似.
都有一個驕傲的王.
有上帝迎佛管教臨到.
又有但以理向他講解.
不過不同的是.
尼布甲尼撒王能獲救.
白沙剎卻是滅亡.
關鍵是什麼?.
關鍵是.
迴轉悔改.
中國人喜歡說失敗是什麼?.
成功之母.
是不是真的?.
你想清楚.
失敗是否成功之母?.
失敗和成功兩者.
有沒有必然直接的關係?.
其實沒有.
你覺不覺得一個人.
失敗五次,十次.
儲夠了之後.
就會自動升級成成功?.
有沒有?.
沒有.

$^{761}$不像某信用卡.
消費滿多少.
就會自動由某卡變成某卡.
失敗不會自動變成成功.
我們有很多例子.
都見識過.
什麼叫不斷犯錯.
下一句是什麼?.
有沒有聽過?.
壞人不斷犯錯.
但是死不改過.
有聽過吧?.
我們經歷過很多例子.
不斷犯錯.
下次又來.
再下次又來.
所以你看到.
再多的失敗不會自動變成成功.
裡面需要有一個關鍵.
就是人要從失敗裡.
懂得悔和改.
悔改才能夠生出成功.
才能夠改變.
如果只是不斷的失敗.
無論再多次.
大概只是原地踏步.
甚至越踩越深.
所以有人說笑.
失敗不是成功之母.
最多只是外祖母.
失敗最多只是成功的外祖母.
要先有悔改.
中間要有悔改.
才能夠生出成功.
真正的弟兄姊妹.
無論我們信不信靠耶穌.
最終我們每一個都要像白沙薩一樣.
來到上帝的天秤裡.
讓上帝澄一澄.
看我們夠不夠清.

$^{801}$有沒有額清.
上帝只得一清二楚.
我們到底有沒有虧欠祂.
祂很明白.
但是有一個好消息.
有一個好消息是.
耶穌基督已經代替了我們.
被澄了在上帝的天秤上.
祂承擔了我們那些不夠澄.
對上帝虧欠的罪.
讓我們每一個信靠耶穌的人.
讓我們每一個認耶穌基督為主的人.
你可以得到上帝的接納.
會出不再一樣的生命.
有新的機會.
所以假如.
無論在座當中.
或者在螢幕當中.
有未認識耶穌的朋友.
我邀請你把握這個機會.
來去認順這一位耶穌.
成為你人生的主.
又或者可能在過往的日子.
我們真的誤用.
濫用了上帝給我們的自由.
或者一些權利.
做了一些可能令自己.
令別人失去的.
甚至令上帝受虧損的事.
我也很想鼓勵你.
藉著今天經文的提醒.
來到上帝的面前.
真誠的認罪悔改.
來把握上帝給我們.
回轉的機會.
上帝說過.
憂傷痛悔的心.
祂必不輕寒.
我們也很清楚.
人若認自己的罪.

$^{841}$上帝是信實公義的.
一定要赦免我們的罪.
洗淨我們一切的不義.
真的.
耶穌一直在尋找.
拯救每一個失喪的罪人.
包括你和我.
耶穌要帶我們.
回到祂的正路當中.
最後.
最後我們不能夠.
不看一看但以利.
我們剛才說過.
尼泊爾和尼沙皇.
很賞識但以利.
在他在位的時候.
納但以利成為.
巴比倫哲士們的領袖.
對不對?.
但現在呢?.
現在是誰作主?.
現在是白沙剎執政.
已經風光不再.
正所謂一朝天子一朝臣.
其實白沙剎執政了多久呢?.
剛才我們一開始說過.
公元前550年執政.
到他死那天.
即是這個故事.
今天這段經文的結局是.
539年.
你數一數最少多少年?.
50減39.
11或者12.
算上頭尾.
這11年的裡面.
但以利其實很大機會被人怎樣?.
被人投限自散.
為什麼這樣說?.
因為那個盛宴沒有他份.

$^{881}$如果你還記得.
以前他是管理哲士的.
但來到另一個人話事的時候.
沒有他份.
沒有他.
另外一方面就是.
當白沙剎需要找人解夢解謎的時候.
當然應該找誰?.
當然是找但以利.
有經驗有往績.
他才是巴比倫的第一解夢高手.
但不如他.
應該正確地說.
忘記了他.
忘記了這位前朝的重臣.
但以利被人遺忘了.
被雪藏了.
被人丟了一邊.
但以利沒有因為這樣.
本來就過活.
他沒有因為沒有人看好自己.
就索性躲起來隱居.
不問世事.
終日可能懷念昔日多麼風光.
為自己現在鬱鬱不得志.
每天都在滲水.
你看到他沒有這樣.
他沒有因為被人丟了一邊.
又或者我已經年紀大了.
七八十歲了.
他就開始懷疑.
否定自己的人生.
甚至乎放棄了自己.
他仍然保留著一個有用之身.
等候上帝用他.
有沒有留意這一點?.
為什麼這樣說呢?.
其實我們剛才解釋過這幾個字.
當你懂得看了.
懂得打動看.

$^{921}$原來利利利利.
看到提克勒烏化爾森.
看到他這樣寫.
但這樣是否足夠解釋.
原來是說巴比倫要被取締.
原來是說上帝已經將他拿去稱一稱.
其實很難.
比爵你能解釋嗎?.
很可能是.
字我就懂得看了.
這裡是解數算數算.
稱一稱和分裂.
我懂得看了.
不過這幾個字和白沙薩王有什麼關係.
我就解釋不了.
就這麼多了.
為什麼他可以解釋得這麼詳細呢?.
你不難想像.
但他仍然很懂得當時國內外的局勢.
用今天的說法就是.
他有看時事和新聞.
他還有看時局分析.
他可以知道當時世界發生什麼事.
他腦海裡甚至有幅圖畫.
就是他知道白沙薩執政不行.
因為白沙薩和他之前服侍的老闆.
尼布甲尼沙差幾皮.
差幾條街.
再加上馬代波斯現在這麼強勢.
步步進逼.
他不難推算到.
其實白沙薩遲早完蛋了.
巴比倫國遲早沒有.
另一方面更重要的是.
從他過往的事奉經歷.
他很認識上帝是怎樣做事的.
他知道就算強到像尼布甲尼沙那樣.
上帝都可以即時拉落馬.
要拿走就拿走.
更何況是白沙薩.

$^{961}$上帝要一個這麼狂妄自大的人.
隨時下架.
是很容易的事.
但以他知道有一樣東西.
主權在上帝那裡.
怒氣也沒有用.
他要等.
等上帝的時機.
好了,現在等到了.
但以你知道皇宮牆上的字.
太后的舉薦不是偶然.
是上帝計劃的一部分.
是上帝的安排.
基於他過往對局勢的分析.
對上帝的熟悉.
加上現在牆上的字.
數算,呈一呈,分裂.
這幾個字的引證.
他知道為上帝候命這麼多年.
現在上帝終於要用他.
透過他宣佈.
上帝要出手改朝換代了.
弟兄姊妹.
我想但以你這種態度讓我們看到.
在現實裡.
其實時代真的會改變.
皇朝也會掩沒.
但那些擁有美好靈性的人.
好像但以你這裡所說.
有美好的靈性.
別人看到.
擁有屬靈質素的人.
總是會在這些轉變裡.
被上帝使用.
會留下美好的榜樣.
給你和我.
所以無論今天我們的人生.
到底是怎樣也好.
無論有什麼挫折也好.
即使沒有人看好我們.

$^{1001}$甚至被人遺忘了.
會覺得我們可能沒有作為.
但這段經文提醒我們.
勉勵我們.
上帝仍然在人的國裡掌管.
作為他的兒女.
我們要為上帝保留住.
保持住那一份生命的熱誠.
不要輕易被你現在身處的環境.
打擊你對上帝的信靠.
也不要輕易被你現在.
面對的逆境際遇.
淋熄了你對上帝的那團火.
你要放心.
好像但義理那樣.
上帝仍然掌管著一切.
他不會忘記你.
他一直都看好你.
看顧你.
他也看好你.
他一定會用你.
帶領你.
來完成他的旨意.
但我們要像但義理那樣.
為上帝爭口氣.
不要自己放棄自己.
上帝一定會在我們.
高高低低的人生夾縫裡.
他會用我們.
而我們需要做的.
就是專心等候他的時機.
等候他的帶領.
來堅持成為他忠心的僕人.
因為這樣的緣故.
你要為你的靈命保持住.
來保溫.
來增值.
《丁姐妹盼望》這段經文.
雖然很遠古.
但它所說的很現實.

$^{1041}$很靠近.
其實是你和我的寫照.
願意我們因為上帝的掌管.
我們能夠為上帝.
保持住那份生命的熱誠.
來隨時候命.
給上帝使用嗎?.
\newpage



\section{但以理書 6:1-19}
\label{sec:YHKUQm4ko_M}
\textbf{2020年10月25日 崇拜直播(青少主日)｜但以理書6\_1-19}
\newline
\newline
連結: \href{https://youtube.com/watch?v=YHKUQm4ko_M}{\texttt{https://youtube.com/watch?v=YHKUQm4ko\_M}} ~~~~ 語音日期: 2020-10-25
\newline
\newline
\hyperref[sec:1ZS5UTvrjn8]{\small{< < < PREV SERMON < < <}}
~
\hyperref[sec:index_chronic]{\small{[返順時目]}}
~
\hyperref[sec:MFeunbAAmR4]{\small{> > > NEXT SERMON > > >}}
\newline
\newline
但以理書 6:1-19
\newline
\begin{longtable}{cl}
\hline
\hline
章節 & 經文 (和合本修訂版)\\
\hline
6:1 & \begin{tabularx}{0.7\textwidth}{X} 大流士隨心所願，立了一百二十個總督，治理全國， \end{tabularx} \\ \\ \relax
6:2 & \begin{tabularx}{0.7\textwidth}{X} 又在他們以上立總長三人，但以理也在其中；使總督在他們三人面前呈報，免得王受虧損。 \end{tabularx} \\ \\ \relax
6:3 & \begin{tabularx}{0.7\textwidth}{X} 這但以理因有卓越的靈性，超乎其餘的總長和總督，王想立他治理全國。 \end{tabularx} \\ \\ \relax
6:4 & \begin{tabularx}{0.7\textwidth}{X} 那時，總長和總督在治國的事務上尋找但以理的把柄，為要控告他；只是找不到任何的把柄和過失，因他忠心辦事，毫無錯誤過失。 \end{tabularx} \\ \\ \relax
6:5 & \begin{tabularx}{0.7\textwidth}{X} 那些人就說：「我們要找但以理的把柄，若不從他神的律法中下手，就尋不著。」 \end{tabularx} \\ \\ \relax
6:6 & \begin{tabularx}{0.7\textwidth}{X} 於是，總長和總督紛紛聚集來見王，說：「大流士王萬歲！ \end{tabularx} \\ \\ \relax
6:7 & \begin{tabularx}{0.7\textwidth}{X} 國中的總長、欽差、總督、謀士和省長彼此商議，求王下旨，立一條禁令，三十天之內，不拘何人，若在王以外，或向神明或向人求甚麼，就必扔在獅子坑中。 \end{tabularx} \\ \\ \relax
6:8 & \begin{tabularx}{0.7\textwidth}{X} 王啊，現在求你立這禁令，在這文件上簽署，使它不能更改；照瑪代人和波斯人的例，絕不更動。」 \end{tabularx} \\ \\ \relax
6:9 & \begin{tabularx}{0.7\textwidth}{X} 於是大流士王在這禁令的文件上簽署。 \end{tabularx} \\ \\ \relax
6:10 & \begin{tabularx}{0.7\textwidth}{X} 但以理知道這文件已經簽署，就進自己的家，他家樓上的窗戶開向耶路撒冷。他一天三次，雙膝跪著，在他的神面前禱告感謝，像平常一樣。 \end{tabularx} \\ \\ \relax
6:11 & \begin{tabularx}{0.7\textwidth}{X} 於是，那些人紛紛聚集，發現但以理在他神面前祈禱懇求。 \end{tabularx} \\ \\ \relax
6:12 & \begin{tabularx}{0.7\textwidth}{X} 他們就進到王面前，向王提及禁令，說：「三十天之內不拘何人，若在王以外，或向神明或向人求甚麼，必被扔在獅子坑中，王不是在這禁令上簽署了嗎？」王回答說：「確有這事，照瑪代人和波斯人的例是不可更改的。」 \end{tabularx} \\ \\ \relax
6:13 & \begin{tabularx}{0.7\textwidth}{X} 他們對王說：「王啊，那被擄的猶大人但以理不理會你，也不遵守你簽署的禁令，竟一天三次祈禱。」 \end{tabularx} \\ \\ \relax
6:14 & \begin{tabularx}{0.7\textwidth}{X} 王聽見這話，就甚愁煩，一心要救但以理，直到日落的時候，他還在籌劃解救他。 \end{tabularx} \\ \\ \relax
6:15 & \begin{tabularx}{0.7\textwidth}{X} 那些人就紛紛聚集到王那裡，對王說：「王啊，當知道瑪代人和波斯人有例，凡王所立的禁令和律例都不可更改。」 \end{tabularx} \\ \\ \relax
6:16 & \begin{tabularx}{0.7\textwidth}{X} 於是王下令，人就把但以理帶來，扔在獅子坑中。王對但以理說：「你經常事奉的神，他必拯救你。」 \end{tabularx} \\ \\ \relax
6:17 & \begin{tabularx}{0.7\textwidth}{X} 有人搬來一塊石頭放在坑口，王用自己的璽和大臣的印，封閉那坑，使懲辦但以理的事絕不更改。 \end{tabularx} \\ \\ \relax
6:18 & \begin{tabularx}{0.7\textwidth}{X} 王回到宮裡，終夜禁食，不讓人帶樂器到他面前，他也失眠了。 \end{tabularx} \\ \\ \relax
6:19 & \begin{tabularx}{0.7\textwidth}{X} 次日黎明，王起來，急忙往獅子坑那裡去， \end{tabularx} \\ \\
[1ex]
\hline
\hline
\end{longtable}
$^{1}$《大利文化(大禮務)條例草案》 - 全體委員會審議 - 第6章1至10節.
「大利烏隨心所願,立一百二十個總督,治理通閣,又在他們以上立總長三人,但以理在其中.使基督在他們三人面前回覆事務,免得王受虧損.因這但以理有美好的靈性,所以顯然超乎其餘的總長和總督,王又想立他治理通閣..
那時總長和總督尋找但以理誤覺的把柄,為要參他,只是找不著他的錯誤過失.因他忠心辦事,毫無錯誤過失.那些人便說:我們要找參者但以理的把柄,除非在他神的律法中就尋不著.」.
「於是總長和總督紛紛聚集來見王,說:願大利烏王萬歲,閣中的總長,欽差,總督,謀事和從府彼此相宜,要立一條堅定的禁令,或亦求王下旨要立一條,云云.三十日內,不驅害人,若在王以外,或向神或向人求甚麼,則不得不求王下旨.」.
「王現在求你立這禁令,加蓋玉璽,使禁令決不更改.」.
「照馬代米底亞和波斯人的律例,是不可更改的.」「於是大利烏王立這禁令,加蓋玉璽,但以理知道這禁令蓋了玉璽,就到自己家裡.」.
第十九節:「次日黎明,王就起來,急忙往獅子坑那裡去.」.
「臨近坑邊,哀聲呼叫但以理,對但以理說:「永生神的僕人,但以理啊,你所常事奉的神,能救你脫離獅子嗎?」.
「但以理對王說:願王萬歲!我的神差遣使者,封住獅子的口,叫獅子不傷我,因我在神面前無辜,我在王面前也沒有行過虧損的事.」.
第十六節:「王就甚起來,吩咐人將但以理從坑裡繫上來,於是但以理從坑裡被繫上來,身上毫無傷損,因為信靠他的神,王下令人就把那些控告但以理的人,連他們的妻子兒女都帶來,.
永在獅子坑中,他們還沒有到坑底,獅子就抓住他們,咬碎他們的骨頭,那時大利烏王傳旨:.
「曉遇住在全地各方,各國各族的人說:願你們大享平安!現在我幹旨曉遇我所統合的全國人民,要在但以理的神面前戰兢好鋸.」.
「因為他是永遠長存的活神,因他的國永不敗壞,他的權柄永存無極,他護庇人,答救人,在天上地下施行神蹟其事,救了但以理脫離獅子的口.」.
「如此,這但以理,當大利烏王在位的時候,和波斯王古列在位的時候,大享亨同.」.
(字幕由梁美芬提供).
現在的年青人不容易,面對著不同的挑戰,希望我們的教會能夠成為一間能夠肯定他們,建立他們,支持他們的地方,讓更多年青人能夠認識上帝,事奉上帝..
我的太太和我的小女兒都平安,當找照片的時候,才發現原來我家四口沒有一張合照,可能找個公仔BB代替我吧!.
介紹一下我的家庭情況,太太上個月就生了小孩,今天是小孩,即是小女兒,40日大.原來要照顧一家幾口子的確不容易,老套一點說一句,養兒方知父母恩..
小女兒算是一個有難度要照顧的BB,特別是因為她喝奶不多,而且容易嘔奶,經常需要長時間掃風,特別是在晚上,昨晚也是這樣,她會莫名的去嚎哭,然後不停的進入掃風換尿片,不停的喝奶的循環裡面,難以入睡..
正如今天的主角一樣,但以理,他也經歷了上帝的保守.經文的背景是公元前的五世紀,當時的猶大國已經被巴比倫所列,但以理和其他猶太人被擄到巴比倫..
第一章裡面,但以理和他三個朋友被王選中,要栽培成為哲士.第二章,但以理成為首席的哲士.第三章,但以理更加成為巴比倫的總理.第四章,他更加成為王的私人顧問.第五章,他更加成為國中裡面第三個統治者..
第五章的尾部,描述馬代人,大利烏,或其他日本大劉氏,殺死了白沙薩,成功佔領了巴比倫城市,改朝換代..
在馬代和波斯的統治下,但以理被大利烏王沿用,成為當朝的大臣,幫助改變朝代更換中減少的困難和適應.但以理當時已經接近八十多歲的高齡..
第六章,但以理書,一至九節,這是一個惡人當道的世界.馬代波斯的國,版圖很大,包括波斯,巴比倫,巴勒斯坦,小亞,聖亞和安吉的地方..
大利烏王為了方便管治,在這個國家裡面設立了120個總督,節省了很多行政和瑣碎的事務.他們這些總督的工作,包括為王徵收稅項,而且還會防止區內人民造反..
在這120個總督裡面,他再設立三個總廠,負責監督這些總督的工作.每個總廠管治70個總督..
大利烏王害怕各區的總長因為有貪污的情況,拿了稅項來貪污,據為己有.於是他想成立但以理,成為總長之首..
但以理雖然年紀老邁,在政治生涯中橫跨了巴比倫和波斯兩個王朝,他的工作也超越了120個總督和兩個總長之上,有上帝的靈在他們當中..
於是大利烏王想晉升他成為一個較高的位置,讓他自己管理全國.大利烏王想將但以理升官事件洩露出去,引起其他總長和總督的質度.他們想要找但以理的把柄阻止他能夠升職這件事..
但以理為何會遭到敵視?以下有四個原因.因為但以理被擄到巴比倫身為地下亡國的奴隸,王為何不用自己的同胞而要用外國人?.
但以理年紀老邁為何還要擔任這麼高的職位,應該給其他機會其他人?但以理是巴比倫的重臣,他究竟是否還會效忠馬代波斯的國王?.
但以理青年公正,靈性也好,在他手中,其他人不能以權謀私,更不能貪贓枉法.於是他們想剷除但以理..
但以理行事忠心辦事,毫無缺失,他們知道在但以理的生活中沒有任何過錯,於是他們想在他們的宗教,想在他們的行事為人的守則中找一些把柄,於是設計陷害他,很想讓他下台..
這些總長和總督對王說謊,他們聲稱自己已經建立了一個方案,已經全國人民一致通過,包括但以理也同意,但以理卻不知情..
他提議王利用王虛榮的心,命令30天內,王是唯一敬拜的神明,30天內,可能因為波斯國內也有其他神明要拜,所以這個法例不能維持很久,不能長期實踐,所以如果任何人犯了這條律例,就要被扔到獅子坑內..
王不知總長們的底蘊,以為這班大臣是為了著想,能夠營造一個聖人,明君的形象,建立國民和其他人的權柄威信..
在這個世代內,也有極權的國家,他們更會將自己的領袖神化,放在他們的神像,圖像,在不同的地方叫人去膜拜,他們會更加將耶穌的故事套入自己的生平內,以賺取人民的民心..
從歷史來看,宗教,教育,傳媒,經常在掌權者控制的地方,因為他們覺得這些東西對他們來說是威脅,必定要去控制他們的思想和他們的生活..
這是一個惡人當道的世界,政治角力,權力鬥爭,充滿在世界每一個角落,不論是大國,小國,國家與國家之間都有不同的競爭,特別是當我們看新聞,我們都會看到一些超級大國舉行一些貿易戰爭,拉攏其他小國去支持,以鞏固自己的影響力,成為更強的勢力..
不單是國家與國家之間,國家之內的明顯紛爭更加明顯,總統的選舉,內閣的選舉,黨派之間的選任,地區性,功能性,原來全部都有政治的角力,權力的鬥爭在裡面..

$^{41}$當我們看到這些新聞的時候,我相信不陌生,不同的群體都為了自己利益的關係,你如我乍,互相詆毀,務求得到更多人的支持..
這些惡人的特質,他們是不務正業,能力也有限,於是乎專門在其他人的身上找錯處,令自己升官發財.他們搬弄是非,拉攏權貴,分黨派,作小圈子,互相指責,互揭醜聞,將人污名化..
當然在某些地方裡面,如果背後有人支持,更多的醜聞都能夠繼續連任.他們的詭計為了得到更多的好處,他們以愛心,和平,正義為一個旗號,目的是包裝他們裡面的詭計,無論是高傲的眼睛,煞方的舌頭,流無辜人的血和惡計的心..
他們以半真半假的謊言來包裝,試圖混淆視線,能夠達到目的.不單只是大義利的年代,今天的年代同樣在這個世界裡面發生..
這些事情不單只是在公共的組織,更加在不同的商業的公司,辦公室,學校,甚至乎有人的地方,就有這些黑暗的事情..
因為最根本的問題是人心裡面的罪,因為罪的惡果,令人性的醜惡,自私,貪婪,這一切一切都表露無遺..
特別在過往的一年,相信大家都經歷過很多大大小小,無論是社會性,家庭裡面,都有很多分裂和對立,加上疫情,令到整個社會都陷入了一個負面的氣氛裡面..
今天是青少主日,我們的青年人又在面對什麼樣的困難和挑戰呢?什麼是惡人?什麼是籠罩著他們的生活裡面呢?.
前一陣子有一個聲中的家長說過,不同的家長都在分享他們最困難的地方,就是不知道怎樣去控制,就是不知道怎樣去處理他們子女用手機的問題..
大部分的青年人花了大部分的時間在電話裡面,如果叫停他們,又好像影響關係,但如果用得太長,又怕影響精神,影響學業,影響他們的健康..
有些中日打機,有些會看不同的影片,或者是劇集,有些用不同的程式,瀏覽不同的平台,收集資料,發放資訊..
網絡裡面有大量的消息,正確地使用,去幫助我們學習新的知識和技能..
不過網絡裡面的世界也眼花繚亂,資訊爆棚,裡面的信息也沒有經過過濾和篩選,大量的信息可能都是充滿著網絡裡面的欺凌,批鬥,舉報,公審的事情,更加嚴重的是犯法,成為一個色情的溫床..
早前有一段時間,網上有一個新聞,這段片段是記述一班科技人員的成員,他們想研究透過網絡的兒童拐帶問題的情況,於是其中一個媽媽大約37歲,假扮成一個15歲的少女,在社交平台裡面建立一個虛擬的帳戶,.
在上面上載一些自己改圖的照片,建立一些虛擬的對話和背景,當這個帳戶在這個世界裡面面世的頭一個小時,已經超過七個成年人主動找她去聊天,而頭一個星期已經超過九十個陌生人去找這個女人去聊天,對話裡面的內容是暗示要求可能拿照片,或者試圖威脅這個少女..
之後這個媽媽把這個社交平台的設定,教到她11歲,結果不夠兩分鐘已經有人去上門,不斷發放一些短信去騷擾他們,他們的頭像更加可能是一些性器官,信息的內容也都以色情為主,要求視頻,要求見面..
這個是網絡裡面的情況,加害的人無需走進一個房間,也都無需走進一個地區,同一個國家都能夠進行這些行動,往往年輕人可能受到受傷,他們會覺得羞恥,驚恐,更加將問題藏在自己的心裡面獨自去面對..
這個組織最後將這些部分的犯人能夠捉拿,不過他們去鼓勵,這些捉拿的行動遠遠低於家長給青年人的鼓勵和支持..
今天的年輕人面對著這個無力的世界,有一份無力感,無論是學業還是社會的影響,更加不幸的是原本可以支持他們的家庭,卻是成為另外一個令他們受傷,令他們一個傷心的地方,令他們更加走進這個網絡的世界裡面,希望更加得到人的愛,肯定..
這代年輕人每一個心裡面原來有一顆寂寞的心,更加有一顆受傷的心,面對這個情況我們可以怎樣做呢?.
第二部分,上帝的子民的回應,但以理書六章十至十八節,但以理聽到王的禁令,禁止除了王以外,任何人都不可以向王以外的人去祈禱,面對著斯子坎死亡的威脅,但以理毫不畏懼..
但以理的反應很簡單,他上到這個樓上打開窗戶,向著耶路撒冷的方向去祈禱,故事裡面沒有描述他的言論,或者沒有描述他內心的掙扎,可能為了給讀者一個印象,就是他完全不怕,完全毫不動搖,信服上帝..
如老不幸的事情,有人會怨天尤人,會歸咎自己的不幸,有人會終日享樂,沉醉於一些行為,逃避現實,也有人會覺得保持中立,不問世事,遠離世間的紛爭,就可以解決問題..
更加有人會選擇以惡報惡,最終讓自己的仇恨和苦毒所充滿,令自己的傷害更加深.有人決定以武,有人決定以勇氣去面對,迫適任何代價,反抗到底..
不過這一切的方法,都不是蓋爾利的方法,蓋爾利安守本份,選擇將生命的主權交給上帝,因為他知道上帝有最終的安排,上帝必定會有最公平的安排..
表面上這個禁令只是考驗他們有沒有公開的祈禱,但是對蓋爾利來說,這個考驗卻是問到他最終的一個問題,蓋爾利最終效法的對象是誰呢?.
當地上的法律和天上的法律相抵觸的時候,蓋爾利又如何選擇呢?他陷入這個兩難的困局中,究竟選擇天上的身份,還是選擇地上的馬代邦斯總督的身份,只是拜大利烏王呢?.
如果選擇去敬拜上帝,其實意味著他要放棄過往幾十年的努力,放棄安恕的生活,放棄家庭,放棄一切的權力,更加重要的可能是放棄自己的生命..
有人說,人在江湖,身不由己,留得青山在,哪怕舞柴燒,有時風平浪靜,退一步,海闊天空..
但是蓋爾利大可以改變他祈禱的習慣,就是在暗中去祈禱就可以了,不讓人發現,又或者他忍耐這三十天不去祈禱都可以的,不值得為這三十天而去公開祈禱,以致朝至殺身之禍..
我們也會很體諒人的軟弱,很體諒蓋爾利的情況,我們有時候不需要面對生死的抉擇,但有時候他是因為很多的原因,將信仰,將上帝放在一二邊..
遇到讀書,考試,工作的繁忙,家庭裡面的繁重,有兒女要照顧,身體裡面很累,我們會決定,這個星期不如休息一個星期,不去回教會..
原本有事奉,但忽然有人約我去行山,去玩,不如給自己充充電,說不定放完假回來更有力去事奉,結果放完假回來,通宵去玩,令自己更加累,下次更加不想回教會..
拍拖戀愛之前,遇到一個未進者,覺得自己有信心,不會離開信仰,結果一闖入這個關係裡面,就被另一半拉著離開上帝,更加在教會裡面越來越少見到..
但以你更加重視他國民的身份,多於珍惜他的生命,他不忘記自己是上帝的子民,沒有因為榮華富貴而忘記神,沒有因為誘惑而背棄信仰,更加沒有因為貪圖生死而離開上帝..
他不向大理烏王屈膝,而是向耶路撒冷屈膝祈禱,但以你的祈禱對象不是大理烏王,他相信大理烏王不是他的宗寶,而是耶穌基督耶和華才是他的祈禱對象..
但以你的信心不是天生的,也不是突然一時的衝動,而是每天見機在他和上帝的祈禱,他和上帝的關係,當他親近神,認識神,對上帝慢慢培養信心..
但以你的祈禱有以下幾個特質,第一,他在固定的地方祈禱,每天的習慣就是打開窗去祈禱,於是其他人都看到,他打開窗去向耶路撒冷祈禱,因為耶路撒冷有聖殿,聖殿是上帝的居所,意思是代表向上帝祈求..
不過,當聖殿被巴比倫國滅了猶太國的時候,聖殿已經不在了,而第九章也在繼續,但以你為耶路撒冷的荒涼去祈禱和重建去禱告,他心裡想著他家鄉裡的人,他為了他的同胞,為了他們的安全,求神保守,求神復興..
第三點,他一天三次行切地去祈禱,雖然他工作繁多,但祈禱仍然是他每一天的堅持,他跪在地上的祈禱,這是一個謙卑的態度,他學會尊重神,尊上帝為大,學習要去信服上帝的指示..

$^{81}$但義利一天三次的祈禱,他是有規律,不是故意在上帝的面前,不是故意在王的面前去做反叛,他只是做回他平日的工作..
但面對敵人精心的詭計,但他污昂束手無策,想盡辦法去救大義利,但卻被自己的玉璽所困,不能更改,徒勞無功,在其他神僕的催逼下,最終皇下的命令,要將大義利扔進獅子坑裡面..
或許我們今天不會經歷大義利的考驗,因為我們都在信仰裡面,一個很自由,很安全的環境,可能我們過往經歷比較大的掙扎,會經歷家人的反對,在公司裡面有人要求你去裝香,辦公室裡面有人要求不可以擺十字架,等等等等,這些對於我們來說都是很大的掙扎和挑戰..
但這個世界裡面仍有很多的基督徒,仍有很多的人,他們每天因為信仰的緣故,受著各種各樣的逼迫,受著各樣的挑戰和患難..
每個星期我們教會都會刊登他們的代禱,會刊登為這個世界裡面50個受逼迫不同程度的國家的人的信徒去祈求上帝,求神保守他們,拯救他們,使用他們的見證..
聖經的預言說到這個末後的日子,教徒會受到更嚴重更多的逼迫,如果當我們相信主耶穌基督再來的日子越來越近,這樣我們將來一定會受到更加多的逼迫來到我們的生命,你和我有沒有預備好呢?.
從我們生活的優先開始做起,今天我們很容易著重在神的慈愛,強調著上帝的恩典和救贖,不過我們忽略了上帝的公義,忽略了上帝是興萬不得,忽略了上帝原來都有管教,有懲罰的一面..
我們太習慣安恕的日子,很想尋求一個安穩安定的日子的情況,我們在信仰裡面不可以去下苦功,渴望用世界裡面的價值觀套下來信仰,不過信仰是有要求的,信仰是要付代價,是要我們每一個付出不單只是時間精神金錢,甚至乎要付出我們的生命..
印度的獨立之父甘地,他是一個提倡非暴力和平的主義勇士,這個人非常敬重和喜愛耶穌基督這個偉人,不過為什麼他最後沒有相信耶穌基督呢?.
甘地在他的自傳裡面說道,我很喜歡你們所相信的耶穌基督,不過我不喜歡你們這班基督徒,你們這個基督教不像你們那位的基督一樣..
如果你們這班基督徒能夠真正按照聖經裡面所說的生活原則去做的時候,我相信今天整個印度都能夠成為基督教的地方,很多人會相信耶穌..
有時候基督徒最大的問題不是他們太像基督徒,而是他們不像基督,所以不要怪責其他不相信上帝的無神論者,有時候不是他們不願意相信,而是他們很多時候只是在我們的生活,在我們的行事為人裡面,他們找不到基督..
除了甘地之外,年輕人其實同樣把基督徒的行為看在眼裡,心理學家艾力克森提出了一個成長的階梯,說到成長裡面有八個階段,原來每個階段都有東西要學習和成長..
青少年階段大約是十二到十八歲這個年紀,這個年紀的青年人很想成為成年人,他們很想模仿成年人的行為,這個年紀的人最容易吸收,最容易受到成年人的行為,價值觀,各樣東西所影響,於是他們很想尋求人生的意義,價值,追求肯定感..
若在這個階段裡面尋索失敗,他們很容易採取了社會不被接納的角色,以致表達他們不想要成為的方法,他們難以對人建立長久親密的關係,究竟可以在哪裡學習呢?.
當我們沒有做好教導的工作,沒有做好榜樣,也沒有做好傳遞正面訊息的功能的時候,年青人就在網絡裡面,就在這個世界裡面,就在花花世界裡面,在明星裡面,去找到一些榜樣,讓他們去吸引,讓他們去學習,因為他們覺得這些東西都是很有型,他們很想不計代價,繼續去模仿..
《羅馬書》十二章九到二十節,其實這裡記載了教徒恆士榮的準則,同時也記載了教徒的一些生命特質,愛,公義,誠信,真理,喜樂,和平,仍然是這個世界裡面年青人最真實的渴望..
因為這個疫情,網絡的世界成為了信徒和牧者維持大部分信仰生活的一個平台,而近日這幾個月裡面,社交平台,Instagram,IG,也有不少關於信仰的貼文,也在他們下面的一些hashtag記載著教徒基督教文化運動,Delta Movement,一些標籤..
他們會透過寫作,插畫,以及一些色經等等的不同方法,去讓人認識這個信仰..
年青的信徒在IG裡面發動Delta Movement,希望能夠凝聚不同專業的人在IG,在年青人的世代裡面,推行基督教的風氣,基督教的文化,更多的人去實踐這個信仰,甚至一改以前基督教老土的形象,讓更多年青人可以有份建立教會..
在早兩個月之前,我們教會的年青人也嘗試自己寫一些文章,這些文章都是關於愛情,生活,這一季最近也會寫一下萬聖節,是否應該去拜鬼神,其實都是這些老生常談,不過當再包裝,再進入年青人的世界裡面,他們能夠更容易去接受..
我們很渴望透過這個媒體傳遞更多正面的訊息..
今年的疫情令到很多封印工作可能都會停頓,不過年青人也沒有浪費時間,他們也把握每一個學習聖經的機會,繼續透過主日學,小組團契,崇拜維持他們教會的生活..
今年的裡面,我們也有機會和浸會大學基督徒團體合作,舉辦一些Alpha Course啟發課程,有十幾個大學生有興趣加入這個Alpha Course,開了三組的福音小組,他們大部分都是微信者,對信仰裡面有很多疑問,很想去認識神,每一次開組其實都很開心,.
因為他們心裡面很單純,很想認識上帝是一個怎樣的神,祂是否真實..
我更加聽到中大或科大學園傳道會,他們更加經歷一浪接一浪的復興,每一間學校裡面也有超過一百多人願意認識信仰,開了十幾組的福音小組..
除此之外,我們教會去年去到越南裡面短宣,其實今年原定也想去兩次越南裡面的短宣,不過因為疫情的緣故就暫停,原本以為一切的福音事公都要停頓,但忽然間由八月開始宣教社邀請我們透過Zoom去接觸一些去年英語中心的大學生..
現在我們逢星期四的晚上已經維持了兩個月的時間,跟他們有一個英語班的時間,透過英語跟他們建立關係,目的就是想將這個福音傳給他們,這是一個很興奮的過程,特別我們做嘉賓,當地的機構會給一些車船車馬津貼給我們,.
我們就會將這些錢奉獻給當地的教會,成為他們的事公,這是一個很興奮的過程,很看似是一個逆境,看似什麼都不能做,不過當我們願意信靠上帝,上帝是會開一條新的道路..
最後一個部分,我想告訴大家,這個世界不是黑暗掌權,不是惡人當道的世界,這個世界仍然是天父掌管的世界,仍然是上帝作王的地方..
在《大義利書》第六章十九到二十八節,此日黎明,大利烏王急急腳腳去到獅子坑裡面,看看大義利的命運,大利烏王去到這個坑的旁邊,哀求大義利,發現大義利最後仍然生還,證明他是無辜..
大義利見證夜間神察見他的使者,封住了獅子的口,因為大義利在神的面前是無辜,他在神裡面沒有做任何虧損的事情,大義利從獅子坑裡面走出來,身上沒有任何受傷的痕跡,即使在獅子坑過了一日一夜,都沒有獅子能夠傷害他..
最後,王將控告大義利的人和他的家人扔下獅子坑,他們的後果就是死亡..
這個故事精彩之處在於故事的大反轉,忍耐到底必定得勝,善惡眾頭終有報,原本加害的奸國死在自己設定的獅子坑裡面,原本自認為神和神同等的大利烏王,卻發現自己的無能,能夠稱頌我們的神..
上帝是人生命的主,是歷史的主宰,他獨行其事,進行他拯救的工作..
在整個大義利書第六章裡面,當我們看經文的結構裡面,我們會發現一個平衡的結構,大義利統治全國,一個ABCDE,總之是這個三文治結構,整個漢堡包結構裡面,最重要的是看中間那裡..
王最重要的就是這個故事裡面最重要的角色,這個王在中心裡面出現了兩次,作者很想去肯定王口裡面所說的話,大義利事奉的神是至高至尊,全知全能,公義和公正的神..
大義利烏王命令全國人民崇拜大義利所事奉的上帝,因為他親眼看到這個神蹟,心裡面心存的態度是敬畏大義利所認識事奉的神..
原本禁止人去敬拜神,卻轉變到全國的人能夠一起去敬拜上帝..
原本是一條邪惡的法例,但是因為上帝使用大義利的生命,他的僕人完全能夠扭轉改變..

$^{121}$上帝直著他的僕人去宣認他在外邦世界裡面的工作,宣認上帝裡面的偉大和他的奇事..
神是永活的神,他不是人手所做出來的偶像,他是永活的神,不像其他神一樣會有生有死,他是真正活著的,他的國度永遠長存,永遠不滅..
他是慈悲仁愛的神,樂意拯救人脫離困難,他是獨行其事的神,擁有超自然的能力,隨心所欲,施行神蹟其事..
他是賜生命的神,保佑人的生命,拯救大義利,他也是掌管自然界一切萬物的神,獅子都要聽命於他..
大義利信靠上帝,面對死亡的威脅仍然忠於上帝,黑暗雖然猖狂,但是當我們願意倚靠上帝的時候,忍耐到底,必然得勝..
大義利為什麼要堅持祈禱呢?我相信就像耶穌基督在赫西德瑪利元的禱告一樣,.
我的父啊,若果可以,求你不要讓我喝這苦杯,可是不要照我的意思,只要照你的旨意..
大義利成為上帝貴重的器皿,在一個無神論的地方見證了上帝的偉大,連整個國家,整個人民都要去敬拜,認識上帝..
這是大義利祈禱的堅持,因為他忠於上帝,單單很想榮耀上帝..
在這一季高中大專的主一學,一起上到中國教會的歷史,當中我們都很被上帝的工作,歷世歷代的人物所震撼..
其中我想分享一個人物,成為今天的總結,這個人物是我其中最敬重的一個人物,這個人的中文名譯名叫做李愛銳..
李愛銳的事跡翻譯成電影,有一本書叫做烈火戰車..
1924年的奧運會,李愛銳為了貫徹蘇格蘭長老的教導,遵守安息日聖日的教導,.
他為了榮耀神而放棄了為自己國家奔跑最擅長的100米奧運的項目..
這樣的消息驚動了全世界,當時英國的傳媒嚴厲譴責他的行為,等同於叛國..
當他改跑他最不擅長的400米的時候,一個男的按摩師送了他一句金句:.
「尊重我,我必尊重他」,這個「我」當然是在說上帝,.
帶著神的榮耀,他的目的就是得到上帝的應許.他以短跑的方式完成了400米的賽事,.
最後還破了世界紀錄,拿下奧運金牌冠軍..
他堅持榮耀神的原則,在神的恩典下,讓全世界的人看到神的榮耀..
神果然厚,似乎一切敬畏他的人.英國的傳媒也對他改變了態度,視他為民族裡的英雄..
但是在世俗榮耀最高最高的高峰的時候,他卻選擇履行在神面前的承諾,.
去遠方中國奉獻他的一生.23歲的時候,他在1925年踏上他立志實踐的神裡面的諾言,.
去到當時軍閥割據,戰火燎天的中國,默默地在中心帶著喜樂,在天津一所學校裡服侍年輕人..
1933年,當時有加拿大的記者訪問李愛銳,問道,你真的很高興將生命全面奉上你現在的事奉嗎?.
難道你不留戀在昔日閃耀的鎂光燈,蜂擁的人潮,激動的觀眾,大聲的歡呼,珍惜每一次的慶功宴嗎?.
李愛銳這樣回答道,當然,人有時候很自然地想到這一切,但是我更加喜歡現在的事奉,.
因為現在裡面的工作遠遠超乎你所說的一切.你知道嗎?這一份榮耀是一個永遠不會腦壞的冠冕..
在抗日戰爭中,李愛銳也很努力支援中國抗日軍隊,有一次李愛銳聽到六國的戰士被日本人抓了並要行刑,.
他立刻說到當時的現場,一個老人家告訴李愛銳,原來有一個士兵奄奄一息,尚有氣息,.
李愛銳立刻將這個士兵從屠宰場送到武記院的醫院裡面,李愛銳在這個病房裡面守了三天三夜,直到這個士兵逃離生命的危險..
1941年,當時中國的戰局很不穩定,英國政府勸喻所有的公民,所有的人民都要撤退,離開中國..
5月的時候,李愛銳安排當時華人的太太和妻兒離開中國到加拿大去暫避,.
不過自己卻留在河北省南部繼續參加救治傷民和接濟難民的軍隊服務..
1943年,當他被日軍軍閥佔據了山東省集中營的時候,李愛銳成為人質..
在最黑暗的時候,他仍然發揮人性裡面的光芒,藉著上帝的愛和喜樂,在黑暗的地方,這個集中營帶給人平安,喜樂和盼望..
有一個困在集中營裡的年輕人,蘭格登古海,在他所寫的山東集中營的一本書裡面,他描述了李愛銳的事蹟..
他說,當時李愛銳其實是一個充滿幽默感和熱愛生命的人,他的熱誠和魅力,令大家在那段艱苦的日子裡面很容易去度過..
李愛銳是我所認識的人當中最接近聖人的一位.1945年2月21日,李愛銳死在集中營裡面,主懷安息..
李愛銳雖然沒有像大耳裡一樣逃避了死亡,上帝也沒有施行他的手去拯救他,不過他的生命卻是告訴我們,不單止李愛銳,歷世歷代的信徒,他們願意付上他們的生命,將這個福音傳遞給我們..

$^{161}$認識上帝,服侍上帝,是一個榮耀的事蹟,上帝是值得我們一生去追求,值得我們付上一切,能夠事奉上帝是一個最榮耀的事件..
李愛銳說,你知道嗎?這份榮耀是一個永遠不會腦壞的冠冕,讓我們一起祈禱..
\newpage



\section{馬太福音 12:15-21}
\label{sec:MFeunbAAmR4}
\textbf{2020年11月01日 崇拜直播｜馬太福音12\_15-21}
\newline
\newline
連結: \href{https://youtube.com/watch?v=MFeunbAAmR4}{\texttt{https://youtube.com/watch?v=MFeunbAAmR4}} ~~~~ 語音日期: 2020-11-02
\newline
\newline
\hyperref[sec:YHKUQm4ko_M]{\small{< < < PREV SERMON < < <}}
~
\hyperref[sec:index_chronic]{\small{[返順時目]}}
~
\hyperref[sec:PcbwkOrPy_0]{\small{> > > NEXT SERMON > > >}}
\newline
\newline
馬太福音 12:15-21
\newline
\begin{longtable}{cl}
\hline
\hline
章節 & 經文 (和合本修訂版)\\
\hline
12:15 & \begin{tabularx}{0.7\textwidth}{X} 耶穌知道了，就離開那裡，有一大群人跟著他。他把所有的病人都治好了， \end{tabularx} \\ \\ \relax
12:16 & \begin{tabularx}{0.7\textwidth}{X} 又囑咐他們不要把他宣揚出去。 \end{tabularx} \\ \\ \relax
12:17 & \begin{tabularx}{0.7\textwidth}{X} 這是要應驗以賽亞先知所說的話： \end{tabularx} \\ \\ \relax
12:18 & \begin{tabularx}{0.7\textwidth}{X} 「看哪，我所揀選的僕人，我所親愛，心所喜悅的；我要將我的靈賜給他，他必將公理傳給外邦。 \end{tabularx} \\ \\ \relax
12:19 & \begin{tabularx}{0.7\textwidth}{X} 他不爭吵，不喧嚷，街上也沒有人聽見他的聲音。 \end{tabularx} \\ \\ \relax
12:20 & \begin{tabularx}{0.7\textwidth}{X} 壓傷的蘆葦，他不折斷，將殘的燈火，他不吹滅，直到他使公理得勝。 \end{tabularx} \\ \\ \relax
12:21 & \begin{tabularx}{0.7\textwidth}{X} 外邦人都要仰望他的名。」 \end{tabularx} \\ \\
[1ex]
\hline
\hline
\end{longtable}
$^{1}$經訓在《馬太福音》12章15至21節.
《新約聖經》《馬太福音》12章15至21節.
在新約的17頁.
耶穌知道了就離開那裡.
有許多人跟著他.
他把其中有病的人都治好了.
又祝福他們不要給他全名.
「這是要應驗先知以賽亞的話說.
看啊我的僕人.
我所揀選.
所親愛.
心裡所喜悅的.
我要將我的靈賜給他.
他必將功利傳給外邦」.
「他不憎勁.
不圈養.
街上有沒有人聽見他的聲音.
壓傷的爐位.
他不截斷.
將殘的燈火.
他不催滅.
等他施行功利.
叫功利得勝.
外邦人都要仰望他的名」.
丙姐妹平安.
平安.
很快2020年踏入最後兩個月了.
今年在疫情之下.
時間原本好像過得很慢.
但回頭一看.
時間過得很快.
在今年最後兩個月.
我們希望進入一個新的系列.
在這個系列裡.
我們好像一個插序式.
一起來思想一下.
耶穌這位救主.
這位上帝所差派降生在這個世界裡.
為我們寫在十字架上的救主.
是一個怎樣的救主.

$^{41}$他為我們生命帶來什麼不同.
在這個世界裡.
他加入了歷史之後.
為一個信教的人得著了什麼.
接下來這幾個系列.
這幾次我們整個系列.
都會講到我得著了什麼.
因為耶穌基督的降生.
今天我們會講盼望這個主題.
今天如果講盼望.
很多時候特別是香港社會.
今天不論在疫情之下.
或者在之前的環境裡.
都籠罩著有一種大家好像很絕望的氣氛.
我不知道大家會不會同意.
很坦白講在我的同代裡.
在我自己成長裡.
有些同學或者有些弟兄姊妹.
很多都會想移民.
我問他們為什麼會想移民.
其實目的就是覺得自己在香港這個地方.
已經覺得絕望了.
覺得很想轉換環境.
他將他自己未來的前途.
不再投放在這裡.
他將他自己人生所有的積蓄.
或者他的將來放在別的地方.
他希望在不同的地方裡.
能夠帶來一些新的可能.
當然我不會反對弟兄姊妹移民.
每一個人都有不同的領袖.
都有不同的看法.
不過如果我們只是帶著絕望.
這樣看未來的人生.
無論去到哪裡都好.
如果我們都是帶著這種絕望的心情.
以及對未來的那種體會.
可能我們去到什麼地方.
我們心裡都會有一樣東西.
好像籠罩著我們生命當中.

$^{81}$我們沒有辦法走出去.
我好想今天藉著這一段經文.
去看回究竟我們生命裡.
我們將我們的人生交付給什麼呢.
究竟是給外在的經濟環境.
政治環境,社會環境.
還是將我們的人生交付給一個.
真的可以帶領我們前行的一位上帝.
希望今天我們看這一段馬太福音.
在第十二章裡說到耶穌是誰.
我們的生命是被一個什麼的救主扶托.
從他身上裡.
看看可不可以為我們生命帶來多一點的盼望感.
讓我們有力量走在面前的路.
在我們看聖經前我們一起祈禱.
求上帝幫助我們.
我們同心祈禱.
前愛的主我們恭敬.
將我們每一個的心思意念交託給你.
幫助我們能夠在你華語當中我們得著亮光.
跟著基督耶穌你的降生.
我們生命得著盼望.
求你讓我們更加認清你的旨意.
你是一個怎樣的神.
求你將我們每一個來講明清楚.
以至我們生命裡有方向.
而且我們活在有盼望的生命當中.
恭敬將每一位頂子會交託給你.
無論聆聽的或是宣講的.
給你的聖靈與我們每一個同在.
獻上我們祈禱.
奉基督耶穌得勝的聖名.
Amen.
在馬太福音裡如果大家認識馬太福音的結構.
其實馬太福音最主要整個馬太福音的鋪排.
來自五篇大的講章.
如果比較熟悉的第五至第七章裡.
叫登山寶訓.
耶穌上山的時候教群眾.
然後在第十章裡耶穌藉著差遣十二個門徒.

$^{121}$在差遣的時候對門徒有些吩咐.
這就是門徒差遣的一個講章.
第十章整章就是講耶穌門徒要怎樣行事為人.
到了第十三章那裡.
耶穌又開一個新的段落一個講論.
這個就是關於天國比喻的講論.
這個天國比喻的講論.
目的就是要為門徒或者為聽道的人知道.
因為天國的來臨所以他們生命裡充滿期盼.
他們要等待天國的來臨.
他們的生命當中.
他們不是單單屬於這個世界.
而是天國進入他們生命當中的時候.
產生這個群體裡的一些變化.
到了第二十六章那裡.
就是講到耶穌將要分離的一個講章.
你看到整個鋪排裡面.
林林總總馬太福音用耶穌的講論.
用五篇的講章來構成整個的結構.
今天我們所看的第十二章.
就介乎在第十章的門徒訓練.
吩咐門徒的這個講章.
和第十三章天國比喻裡面夾住了三文治的火腿.
在天國的比喻和門徒訓練當中.
耶穌就將自己是一個怎樣的救主.
在馬太福音裡面.
在他的筆之下來描述.
耶穌是一個怎樣的神.
我們首先看前一點.
第十一章第二節.
一句很有意思的一個問題.
在第十二章第二節裡面.
施浸約翰當時就下了監.
他被捉拿在監牢裡面.
在監裡面他就猜險他的徒弟.
兩個門徒去問耶穌一條問題.
在第十一章第二節.
約翰在監裡聽見耶穌所作的事.
就打發兩個門徒去問他說.
那將來要來的是你嗎?.

$^{161}$還是我們等候別人呢?.
耶穌回答說你們去.
把所聽見所看見的事告訴約翰.
就是乞子看見,騎子行路,長大麻風的得潔淨,.
聾子聽見,死人復活,窮人有福音傳給他們,.
凡不因我跌倒的就有福了..
施浸約翰問這條問題.
其實你回答下去會覺得很有趣.
其實如果你知道施浸約翰其實是.
很早的時候已經成為耶穌傳道的先鋒.
如果要知道耶穌是基督的身份.
施浸約翰大概是最清楚的.
他是第一個知道的.
他是為耶穌來做開路.
不過當一路成長.
一路面對著耶穌傳道的經歷的時候.
施浸約翰也有些質疑.
究竟我等有沒有等錯人呢?.
我等的是否等錯了人.
還是等別人呢?.
我等的不是將要來的那個人嗎?.
為什麼施浸約翰要問這些問題呢?.
在經文裡面提到.
因為約翰在監獄聽見耶穌所作的事.
耶穌基督做了什麼呢?.
其實耶穌基督當然不會犯罪.
不過耶穌基督所做的路線.
或者他所作的工作.
跟施浸約翰心裡面所想的有很大的差異.
當時約翰特別要講悔改的信息.
他特別要將人清晰地知道.
他自己生命裡面犯了很多罪.
離棄了上帝.
叫他們來悔改.
這是施浸約翰很重要的工作.
不過耶穌基督大部分時間做了什麼呢?.
耶穌基督當然有講天國來了.
不過耶穌大部分時間.
跟當時所謂坊間裡面的罪人走在一起.
妓女或者稅吏.

$^{201}$跟他們又喝又吃.
施浸約翰以前是吃蝗蟲野蜜.
施浸約翰覺得跟他所做的.
跟他的想法.
跟他一直覺得將天國使命實踐出來的差距這麼大.
耶穌在這裡對施浸約翰的回答.
他就說你將我所做的事情告訴他.
就是看到盲的人看得見.
不跛的可以走.
但麻風這個世紀絕症可以被潔淨.
聽不到東西的人可以重新聽見.
甚至死人復活.
窮人有福音聽.
將社會上面面對著最大艱難.
甚至被邊緣化的這群人.
當時社會看不起的這群人.
耶穌就走到他們當中.
讓這群人聽見有福音.
聽見上帝直著他.
將這些一切上帝天國的旨意.
直著耶穌基督實踐出來.
耶穌就將這件事回答了施浸約翰.
經文裡面沒有繼續描述施浸約翰對於這個答案滿不滿足.
不過就著這個問題.
除了耶穌對施浸約翰的回答之外.
馬太福音再繼續回答我們看的時候帶著的問題.
為什麼傳道的工作和施浸約翰不同.
對象路線圖.
甚至耶穌走到外邦人.
不是以色列人當中將福音傳開.
是不是等錯人.
在馬太福音裡面馬太繼續描寫.
是呀,他們沒有等錯人.
不過耶穌所做的和這群人所期望的.
他們所想的有一個很巨大的落差.
神做事永遠不是人可以完全想像的.
神的工作往往是超越了人的想像.
當上帝的工作是這樣做的時候.
要調教的不是要調教上帝的工作.
是要調教每一個在地上聽見耶穌所言所行的每一個人.

$^{241}$所以馬太裡面繼續說一些極致的.
某程度上是衝著律法.
但其實耶穌承傳律法的一些舉動.
就是關於安息日的教導.
仔細的可能大家回去再看.
不過我想讓大家知道.
耶穌去到第十二章的時候說到安息日治病.
令到當時固守律法的法利賽人.
都已經萌生殺機.
見到耶穌在安息日裡面.
醫好一個枯乾手的一個人的時候.
這群法利賽人覺得耶穌是要破壞律法.
但在耶穌說.
人只是安息日的主.
安息日是為人所設立.
而不是為了宗教條例來設立.
耶穌就要顛覆了這些想法.
顛覆了很多人對於他作為利賽亞.
這個期望這個舉動的想法.
以致到他要將他天國的道理.
在人意想不到的情況之下傳開.
你看到這個上文下理你會知道.
說的是來的是對的.
耶穌基督就是將要來的那位利賽亞.
那個落差在於人有一個框框.
來看神的工作.
但耶穌不會被這些框框所框住.
耶穌卻是用另類的方法.
用他最嶄新的教導.
但這個是一脈相傳律法的教導.
來將上帝天國的道理來去展明出來.
要調教的就是我們地上的每一個.
在第十二章第十五至十七節.
就是剛才我們所讀的經文.
當耶穌知道那些法利賽人要陷害他.
或者要將耶穌陷在控告的狀態.
甚至要除滅耶穌的時候.
耶穌知道了就離開那個地方.
不和那些法利賽人一齊去做一些爭拗.
有很多人跟著耶穌.

$^{281}$耶穌就將這些病患者全部治好了.
在治好了這個病之後.
耶穌就特別作出一個吩咐.
第十六節.
又祝福他們不要給他全名.
不要將耶穌治好他的這件事.
將耶穌的名字傳揚出去.
在這裡如果很多時候查經的弟兄姊妹.
我們很快就會套答案下去.
為什麼耶穌叫他們不要為他們傳開這個名字.
我們會憑我們的直覺.
或者我們其他篇幅.
其他段落裡面看到的答案.
我們就套進去.
我們第一件事就會說.
因為耶穌傳遞.
如果有很多法利賽人在爭拗.
可能為他做事會帶來很多不方便.
而且耶穌又不想正面衝突.
甚至說到那個時候還沒到.
所以耶穌就不會公開自己的名字.
經文沒有特別將這些答案.
成為馬太福音在這一節裡面.
這個段落裡面的答案.
馬太福音反而不是從這個角度.
來看耶穌基督.
反而是引用舊約爾賽亞書第42章第1至第4節.
這一段的經文.
來解釋耶穌不傳開他的名字.
作為一個應許的來講.
爾賽亞書第42章1至4節是這樣說的.
看吶我的僕人我所扶持所揀選的.
心裡所喜悅的.
我已將我的靈賜給他.
他必將功利傳給外邦.
他不宣揚不揚聲.
也不使街上聽見他聲音.
壓傷的爐位他不截斷.
將殘的燈火他不摧滅.
他憑真實將功利傳開.

$^{321}$他不悔心也不喪膽.
直到他在地上設立功利.
海島都等候他的訓誨.
這裡只是一個節錄.
整個爾賽亞書講的僕人之歌.
篇幅更加長.
不過在這裡.
馬太福音引用這幾節經文.
說耶穌叫他們不要傳開他的名字.
在舊約裡面所說的.
馬太福音說什麼呢.
其實他引用這段經文說.
耶穌是一個很低調的人.
就是這麼簡單.
耶穌很低調.
不過他是一個很低調的僕人.
是神的靈所充滿的.
他的低調和他一直以來.
很多人視為先知的身份.
是一個很大的落差.
如果過往裡面你知道先知.
怎樣令人聽他說話.
他在城門口最多人的地方.
如果我們要傳福音.
我們沒理由躲在角落底.
沒理由躲在櫃底傳福音.
我們一定要在最寬闊的地方.
最多人聽的.
在旺角裡面最多人的.
可能銀行中心門口.
我們就拿著咪高峰.
就在那裡說罪人快點悔改.
我們一般想的先知.
或者想傳揚福音的人.
我們很直覺的舉動就是這樣.
不過馬太引用.
以賽亞先知裡面所說.
他不是說一個很偉大.
什麼什麼的英雄.
他說到耶穌是一個怎樣的先知呢.

$^{361}$或者耶穌是一個怎樣的尼賽雅.
怎樣的僕人呢.
是一個低調的僕人.
是一個不會特別在十字路口裡面說話的一個僕人.
甚至乎有人知道.
體會到他的醫治的能力.
天國的福音.
臨到他身上的時候.
耶穌叫他某程度上做一些反宣傳的舉動.
你不要到處說.
你不要說我出來.
這個是說僕人的那種心智.
說上帝的實言.
上帝的實言就是透過一個這麼低調.
低調到一個地步.
就好像那種反宣傳效果的方式.
來將耶穌的身份表達出來.
與其說耶穌是一個很內斂的人.
就不如說上帝的救贖計劃.
是用那種最低調.
最不起眼.
最謙卑.
最柔和的方式.
來彰顯出來.
他要藉著這一段的僕人之歌.
他要將.
為什麼耶穌叫這些人不要傳揚他呢?.
就是因為舊約應許上帝.
要用耶穌一種最低調的方式.
最謙柔的方式.
來將上帝美好的旨意.
透過一種這麼低調.
最沒有霸氣的方式來傳揚出來.
在第十九至第二十節裡面.
可以繼續的來講.
耶穌的低調.
他講到他不憎敬.
不圈養.
街上也沒有人聽見他的聲音.
壓傷的爐位他不截斷.

$^{401}$窗殘的燈火.
他不催滅.
剛才有講過了.
耶穌他不是一個.
在城門口裡面.
吵鬧的先知.
耶穌反而是用.
用他的行為.
用他所作的功.
來將上帝天國的道理.
來展現出來.
真理是用真正的行為.
用實幹.
不是口舌之爭.
來表達出來.
他是一個不憎敬.
不喊養的先知.
藉著這麼低調的方式.
來傳揚上帝的福音.
在這裡我們稍微思想一下.
今天我們不知道.
今天最吸引我們的眼睛.
最引起我們的注目.
通常是什麼事呢?.
讓我們知道生命裡面可信的.
那件事是可以幫我們人生.
可以依附的.
或者我們生命可以靠賴的.
通常是什麼呢?.
如果.
我以前讀市場學的.
最令人有信心的是什麼呢?.
就是將一件事重複.
例如重複又說.
還要高調的市場學.
就是要宣傳.
令到那件事製造信心的時候.
聽多幾次就會有信心.
所以一些產品出來的時候.
不停不停不停說.

$^{441}$你在電視裡面不停入腦.
你就會自動相信.
就會自動買那件事.
就會相信那個產品.
有時候甚至用一些最嘩眾取寵的方式.
令到人激烈了人家的思考.
改換了他的行為模式.
然後就將這些人.
本來買A牌子的東西.
轉移到B牌子的東西.
這些就是坊間裡面宣傳要做的事.
不過.
在我們的信仰裡面.
或者基督耶穌所呈現的面貌.
某程度上和我們所說的市場學.
和我們所說的宣傳效果.
是完全南轅北轍的.
耶穌基督所做的事.
不是叫我們聽到最大聲的聲音.
甚至走出來最有霸氣的講論.
我們就要多一點信他.
耶穌說.
你看看我低調的那種作為.
有時候.
讓我們的生命可以依附著.
我們可以信靠的.
未必是靠著那些嘩眾取寵的.
大聲疾呼.
或者一些selling point(賣點).
那些銷售的策略.
有一個賣點.
令到人相信.
福音從來不是針對一個人的賣點去做.
耶穌基督他自己的舉動.
卻是將他自己的生命呈現出來.
用最低調的方式.
在這裡再繼續來說.
我們才知道很多時候為了真理會有爭辯.
會有爭辯.
很多弟兄姊妹會上傳福音課程.

$^{481}$以前我們都是互教學.
或者是變道學.
第一件事弟兄姊妹就會問.
牧師如果遇到這些人.
我怎樣爭辯贏他.
我們經常都是這樣.
我們想著叫人信耶穌.
爭辯贏了人.
那個人就會自動信.
不過在經驗裡面.
或者你看到耶穌基督所做.
從來不是一個辯論的遊戲.
信仰從來不是用一種咬頸的方式.
咬贏了.
咬輸了的就說好我信了.
從來不是這樣的.
反而是大部分是在觀察.
信耶穌的人的生命是怎樣的人.
他的言行舉止.
他所做的德行.
他生命裡面對事物的看法.
和他怎樣去回應.
這反而是成為一個最有力.
告訴別人我們所信的主是誰.
而不是說一些很強的咬頸的道理.
當然我們認識真理.
我們能夠明白辯清.
這是重要的.
不過我再繼續說.
耶穌從來不是用一個辯論家的身份.
來討論咬贏了那些法利賽人.
或者當時的普羅市民.
來叫他相信耶穌.
他是用他生命的榜樣.
他不咬頸.
他不憎勁.
他在街上沒有他的聲音.
但是用他自己的生命一個一個.
帶到在上帝的面前.
經文再繼續說.

$^{521}$壓傷的爐位.
他不截斷.
張殘的燈火.
他不催滅.
當我們一讀這兩句的時候.
我們很多弟兄姊妹.
就會把我們的情.
突然間讀入這兩句.
為什麼呢?.
因為一讀這兩句的時候.
就會將自己比喻為.
我真是被壓傷的爐位.
或者覺得自己真是張殘的燈火.
自己很殘.
樣子照鏡子.
然後就說自己是張殘的燈火.
然後就會將我們的可憐.
讀入經文裡面.
覺得耶穌不會壓迫我們.
這樣讀又不是完全錯的.
不過如果看回上文.
講到耶穌的說話.
或者他行事為人的低調.
你知道壓傷爐位.
和張殘的燈火.
就不是描述那個人很慘.
不是純粹說那個人.
落在一種很困苦的境地.
Eugene Peterson在他的信息本聖經裡面.
他這樣去看的.
他在說壓傷的爐位.
他不折斷代表著.
耶穌不會說刻薄的話.
來傷害我們的情感.
張殘的燈火.
他不會摧滅.
他的意義就是在說.
耶穌甚至他的說話.
他的言行是不會逼你埋牆.
讓你覺得他在羞辱你.

$^{561}$所以那個壓傷的爐位.
和張殘的燈火.
那種處理是講到.
耶穌基督那種柔和謙卑的低調.
對於可能有些人都不會咬緊.
對於有些人可能生命裡面.
都還在徘徊著.
究竟真理是怎樣的.
又或者我們轉數沒那麼高的那些人.
耶穌從來不是用話語.
來咄咄逼人.
讓這些人悔改.
罵他.
讓他可以回到上帝那裡.
所以耶穌用一種最柔和的姿態.
面對著不知道怎樣回應的時候.
面對著一些不善於辭令.
甚至生命裡面.
都有很多罪的.
很羞愧的那些人.
耶穌都是用最溫柔的方式.
來呈現他自己是一個怎樣的神.
當你用這個角度來看.
其實滿聖經就能看到.
面對一些遂利.
甚至當時的那些妓女.
那些不起眼的人.
甚至乎是黑人憎恨的那些人.
耶穌呈現的姿態從來不是說.
你看看你多壞.
你看看你多差.
耶穌從來不是用這種方式.
他來到這些人當中.
用最柔和的姿態.
來將天國的道理.
謙謙柔柔地呈現出來.
唯有用一種這樣的態度.
我們生命才能得到真正的釋放.
真的弟兄姊妹.
其實很多時候我們自己知道自己錯了.

$^{601}$聖靈都有提醒我們.
我們是不是真的需要多一個人.
來罵我們呢?.
是不是真的純粹需要用責備的方式.
就可以為我們生命帶來新的改變呢?.
可能我們有時候覺得.
有些特別大聲的聲音.
會吸了我們的耳朵.
但我們信服這些聲音.
是因為我們怕這些聲音.
或者覺得很煩.
算了,我改變吧.
或者我們真心聽見這些呼聲.
我們去改變.
耶穌對我們的呼籲.
他永遠都是用最溫柔的方式.
他的立場很堅定.
但他卻是將天國的道理.
用一種不會有霸氣的方式.
作出一個邀請.
讓人去到他的面前.
所以前一段篇幅他特別說.
「凡勞苦擔重擔的可以到我這裡來」.
「我就使你得享安息」.
這是一個很偉大的邀請.
亦是耶穌基督.
傳揚天國道理的時候.
所呈現出來的那種表達.
十二章十九至二十一節.
耶穌基督.
他用一個這麼謙柔的方式.
到最後他帶來什麼後果呢?.
十九至二十一節.
他就這樣提醒我們.
「厄殉的奴為他不截斷」.
「將塵的燈火他不摧滅」.
「等他施行公理」.
「叫公理得勝」.
「外邦人都要仰望他的名」.
在這裡他提到.

$^{641}$一個這麼謙柔的方式.
沒有霸氣的宣揚真理.
竟然到了最後.
這個公理.
公理這個字很奇怪.
我們平時說公義.
大家會明白多一點.
正義,公義.
藉著一個這麼溫柔謙卑的耶穌.
竟然最後可以得勝.
得勝那些很喧鬧.
得勝那些很暴力.
甚至得勝那些.
最譁眾取寵的聲音.
用最謙柔的方式.
來將這個真理呈現.
而最後這個真理會得勝.
會勝得過.
今天我們的盼望在於什麼呢?.
我們的盼望不是純粹說.
我想有層樓.
沒有,然後上帝你給我吧.
那個為之盼望.
如果我們基督徒的盼望.
只是希望純粹經濟條件改變.
那個層次大概都沒有去到.
耶穌基督所說的.
一直說天國道理的層次.
我不是說那樣東西不對.
不過那個層次永遠都是在這裡.
如果我們是說我們信靠福音.
我們知道耶穌基督.
我們很渴望祂的時候.
我們盼望上帝的旨意臨到地上.
如同行在天上.
上帝的心懷.
上帝的美好.
我們所一直固守的.
那些價值.
在有一天.

$^{681}$在基督耶穌將要來到的時候.
這些一切都要成就.
那種成就方式.
或者一直我們正在努力的.
不是用一些氣勢逼人.
壓著別人那種態度.
卻是用一種最溫柔的方式.
來表達出來.
這一種是一個實現公義.
是一個耶穌基督所用的手段.
甚至我相信是唯一的手段.
今天讓姐妹們想一想.
我們今天會遇到一些很不合理的事情.
我們會怎樣反應呢?.
罵爆她.
用我們的憤怒來表達.
甚至有一個姐妹.
初信主的時候.
她很喜歡罵人.
初信主的.
不是旺丞的,所以大家不用擔心.
多說一點就好.
她做了一些因似調查.
我相信信主無奈.
因似調查到最後.
這個姐妹的分數.
有一個先知的因似.
她很開心地跟我們一班大學同學分享.
她說自己有先知的因似.
我剛剛相信主,不知道什麼叫先知的因似.
我問她是什麼.
然後她拿著一張紙.
很驚訝地讀.
先知的因似就是說.
用聖經的話語.
在這個時代裡面.
作出色情的責備.
和分灰.
說出來很酷.
然後我就不太明白.

$^{721}$我就說,換句話來說.
先知的因似就是用聖經來罵人.
然後姐妹就不太出聲.
我說,如果你說的先知的因似.
就是拿著一本聖經.
到處罵人.
那跟耶穌時代的法利賽人有什麼分別呢.
耶穌所做的和.
先知的因似的弟兄姊妹做的.
有什麼不同呢.
為什麼你這麼渴望.
覺得這麼酷.
我真的帶著很好奇的心來問這個問題.
不過她覺得我當然會質問她.
現在回想起來.
真的挺牙刷.
挺囂張的.
不過這個問題同樣在問我們.
今天我們會不會.
都是自稱為先知的.
我代表神.
我去說上帝的話語.
我就好像戰高幾級.
然後我是神聖.
祂是俗的.
我就要將聖經就打下去.
然後就.
罵到人家可以悔改.
或者.
我們用最溫柔謙卑的方式.
來將基督耶穌的真理.
謙謙柔柔地.
不灰心.
不捨棄.
我們堅信.
上帝的真理將要實現.
但犯不著我們要這樣來動腦.
犯不著要用一些.
破壞了我們生命.
言行舉止的方式.

$^{761}$來將這個真理.
傳揚出來.
今天這段經文不是很複雜.
不過.
大家希望你記住這段經文.
所說的以賽亞書.
裡面的.
應許實現.
那個實現不是純粹說.
那件事.
實現.
在以賽亞書裡面.
所說的應許實現.
就是上帝竟然猜顯一個最謙卑的.
尼賽亞.
來到這個世界裡面.
他不是一個咬頸的先知.
他來到這個世界裡面.
竟然用一種.
最可親的方式.
最令人.
心靈.
最可以貼近上帝的.
方式來領受.
天國的福音.
當我們今天去讀.
這段經文的時候.
當我們看到我們的主是一個這樣的主的時候.
我們生命的盼望.
就在於我們生命.
上帝不是.
不是很喜歡.
挑出我們的壞事.
去邊一邊我們.
上帝也不是喜歡.
要挖苦我們生命裡面的.
破破爛爛之處.
上帝應許過我們.
當我們生命.
有一些不濟的地方.

$^{801}$他就在我們.
當中.
耶穌基督永遠張開雙手.
接納我們的生命.
我們只管將我們的生命.
安然交付在他的手裡.
當我們讀這段經文.
也有第二個帶來我們生命的盼望.
是.
盼望從來不是.
用言語.
或者用說話來講贏一個道理.
盼望卻是.
透過基督耶穌.
他自己的生命來實踐出來.
他的降生.
他的教導.
他的訓練.
他的受苦.
他的犧牲.
基督事件.
基督耶穌自己的生命.
為我們這一群.
活在世界裡面.
不是很理想的世界裡面.
一群人的盼望.
他與我們一起同行.
他明白我們的處境.
當我們知道.
我們是張禪的燈火.
或者是壓傷的蘆葦的時候.
他依然與我們同行.
不會逼我們去牆角那裡.
上帝用最溫柔的方式.
呈現他的救恩.
求主繼續來幫助我們.
今日這個世代.
雖然可能很多人.
咬牙切齒.
可能很多人用很極致的方式.

$^{841}$來面對.
但希望今天.
我們看到上帝.
他依然給機會這個世界.
他依然為我們生命.
帶來一些新的改變.
他亦給機會給我們每一個.
希望我們明白基督耶穌的心腸.
明白祂的旨意.
我們同心祈禱.
基督耶穌願你幫助我們.
讓我們能夠知道.
你是一個怎樣的救主.
在喧鬧的聲音當中.
願你的安靜.
願你的同在.
成為我們生命.
深處的安慰.
驅除我們生命.
裡面很多的暴戾.
幫助我們檢視我們的口舌.
會不會成為.
最強烈的武器.
傷害身邊的人.
檢視我們的心懷念念.
我們有沒有地方.
不理想.
求你催近我們.
貼近我們.
讓我們的心.
被你謙柔的心.
改變.
願你繼續來對先恩與眾弟兄姊妹同在.
我們多謝主的垂聽祈禱.
奉基督耶穌的聖名祈求.
Amen.
\newpage



\section{約翰福音 8:1-11}
\label{sec:PcbwkOrPy_0}
\textbf{2020年11月15日 崇拜直播｜約翰福音8\_1-11}
\newline
\newline
連結: \href{https://youtube.com/watch?v=PcbwkOrPy_0}{\texttt{https://youtube.com/watch?v=PcbwkOrPy\_0}} ~~~~ 語音日期: 2020-11-16
\newline
\newline
\hyperref[sec:MFeunbAAmR4]{\small{< < < PREV SERMON < < <}}
~
\hyperref[sec:index_chronic]{\small{[返順時目]}}
~
\hyperref[sec:d4oKfBQrMaI]{\small{> > > NEXT SERMON > > >}}
\newline
\newline
約翰福音 8:1-11
\newline
\begin{longtable}{cl}
\hline
\hline
章節 & 經文 (和合本修訂版)\\
\hline
8:1 & \begin{tabularx}{0.7\textwidth}{X} 耶穌卻到橄欖山去。 \end{tabularx} \\ \\ \relax
8:2 & \begin{tabularx}{0.7\textwidth}{X} 清早，他又回到聖殿裡。眾百姓都到他那裡去，他就坐下，教導他們。 \end{tabularx} \\ \\ \relax
8:3 & \begin{tabularx}{0.7\textwidth}{X} 文士和法利賽人帶著一個犯姦淫時被捉的女人來，叫她站在當中， \end{tabularx} \\ \\ \relax
8:4 & \begin{tabularx}{0.7\textwidth}{X} 然後對耶穌說：「老師，這女人是正在犯姦淫的時候被捉到的。 \end{tabularx} \\ \\ \relax
8:5 & \begin{tabularx}{0.7\textwidth}{X} 摩西在律法書上命令我們把這樣的女人用石頭打死。那麼，你怎麼說呢？」 \end{tabularx} \\ \\ \relax
8:6 & \begin{tabularx}{0.7\textwidth}{X} 他們說這話是要試探耶穌，要抓到控告他的把柄。耶穌卻彎下腰，用指頭在地上寫字。 \end{tabularx} \\ \\ \relax
8:7 & \begin{tabularx}{0.7\textwidth}{X} 他們還是不住地問他，耶穌就直起腰來，對他們說：「你們中間誰沒有罪，誰就先拿石頭打她！」 \end{tabularx} \\ \\ \relax
8:8 & \begin{tabularx}{0.7\textwidth}{X} 於是他又彎著腰，用指頭在地上寫字。 \end{tabularx} \\ \\ \relax
8:9 & \begin{tabularx}{0.7\textwidth}{X} 他們聽見這話，從老的開始，一個一個都走開了，只剩下耶穌一人和那仍然站在中間的女人。 \end{tabularx} \\ \\ \relax
8:10 & \begin{tabularx}{0.7\textwidth}{X} 耶穌就直起腰來，對她說：「婦人，那些人在哪裡呢？沒有任何人定你的罪嗎？」 \end{tabularx} \\ \\ \relax
8:11 & \begin{tabularx}{0.7\textwidth}{X} 她說：「主啊，沒有。」耶穌說：「我也不定你的罪。去吧！從今以後不要再犯罪了。」〕 \end{tabularx} \\ \\
[1ex]
\hline
\hline
\end{longtable}
$^{1}$今日經訓在約翰福音第八章一至十一節.
在新約聖經第137至138頁.
約翰福音第八章一節.
「於是,各人都回家去了,耶穌卻往橄欖山去,.
清早又回到殿裡,眾百姓都到他那裡去,.
他就坐下教訓他們,.
問是和法利賽人帶著行淫時被拿的婦人來,.
叫她站在當中,就對耶穌說:.
夫子,這婦人是正行淫之時被拿的,.
摩西在律法上吩咐我們把這樣的婦人用石頭打死,.
你說該把她怎麼樣呢?.
他們說者說乃試探耶穌,要得著告她的把柄,.
耶穌卻彎著腰,用指頭在地上畫字,.
他們還是不住的問她,.
耶穌就直起腰來對他們說:.
你們中間誰是沒有罪的,誰就可以先拿石頭打她?.
於是又彎著腰,用指頭在地上畫字,.
他們聽見者說,就從老到少,一個一個的都出去了,.
只剩下耶穌一個一人,.
還有那婦人仍然站在當中,.
耶穌就直起腰來對她說:.
婦人,那些人在哪裡呢?沒有人定你的罪嗎?.
她說:主啊,沒有..
耶穌說:我也不定你的罪,去吧,從此不要再犯罪了..
請問各位姐妹平安,.
還記得上一次講到《但耶律書》第五章的時候,.
我提過一件事情,就是上帝在聖經的記載裡,.
上帝曾經用指頭來到他三次寫字,有沒有印象?.
第一次就是在《出埃及記》,.
上帝就用指頭在兩塊石板上寫上他的誡命律例,.
彰顯著上帝與人立約的記號..
第二次就是在《但耶律書》第五章,.
上帝在粉牆上寫上白沙薩的罪行和刑罰,.
顯示了人對上帝的虧欠..
第三次就是剛剛所讀的經文,《約翰福音》第八章,.
這次我們不知道耶穌在地上寫了什麼,.
不過我們確實知道,.
他最後放過了那個行淫的婦人,.
沒有定他的罪,.
顯示了上帝已經降世為人,.

$^{41}$彰顯他的恩典,並且宣告他對罪人的寬恕和接納..
正如今天我們看到的這段經文,.
耶穌在聖殿裡與百姓教訓他們的時候,.
文士和法利賽人將一個行淫的婦人帶到人群當中,.
並且要求耶穌根據摩西的律法作出裁決..
很明顯,可能你也聽過這段經文,.
很多聖經學者也知道這是一個陷阱,.
某程度上,哪有這麼巧合的法利賽人和文士,.
會抓到一個只是行淫的婦人呢?.
通常這些都是要麼預先收到情報,.
要麼就預先設局,.
他們如何得知有這樣的事情發生?.
第二方面就是,.
為什麼只抓到一個女的,抓不到一個男的呢?.
應該也抓到一個男的吧?.
第三方面就是,其實他們根本知道怎麼做的,.
律法已經說明了可以怎樣處理整件事,.
或者犯事的人,.
其實他們沒有必要專門來諮詢耶穌,.
很明顯,就像經文裡所說,.
他們不是要彰顯公義,.
而是要試探耶穌,.
他們有自己的目的,基本上..
假如耶穌不准許群眾用石頭打死這個婦人的話,.
很自然會有一個陷阱,.
就是被人指控他公然違反摩西的律法,.
那麼,法利出行和民事就會找到把柄,.
將耶穌從猶太的社群裡面排擠出去,.
又或者假如耶穌不准群眾,.
啊,不是,對不起,.
他准許群眾打死這個婦人的話,.
那麼又會掉落到另一個陷阱,.
就是被人控告他濫用私刑,.
觸犯了羅馬人的律法,.
因為當時只有羅馬巡撫能夠有這個權柄,.
可以定別人的死罪,.
所以你看到,.
似乎耶穌答又死,不答又怎樣?.
不答也死,進退兩難,.
但是你看到耶穌也挺淡定的,.

$^{81}$他沒有即時去回答,.
而是在地上寫著什麼?.
寫著字,.
我們不知道內容是什麼,.
但是我們卻看到耶穌一種處之泰然的態度,.
面對著敵人的詭計,.
群眾的催逼,.
還有現場環境你不難想像,.
很吵擾,.
但是你看到耶穌一點都不害怕,.
他沒有心裡面說,.
糟了,這次死了,.
答又不是,不答又不是,.
他也沒有怨天尤人,.
這次被婦人害死,.
無往之災,.
他也沒有急於為自己脫身,.
來迴避這個難題,.
他沒有想過,.
不如壞人一點,.
其實他可以怎樣做?.
你們將婦人交給羅馬巡撫,.
搞定她吧,.
他可以這樣袖手旁觀,.
他沒有這樣迴避,.
他更沒有和傀儡錯人和民事死過,.
反正都是,.
不如踢爆他們冤枉我的陰謀,.
大家最多攬炒,.
他沒有這樣做,.
你看到耶穌做了一個很好的示範,.
無論旁邊的人怎樣催逼他,.
來回應一些難題,.
做裁決,.
耶穌也沒有被人打亂他的節奏,.
相反耶穌反過來,.
反客為主,.
他透過彎腰在地上寫字的動作,.
迫身邊的人先等一會,.
你有理吵,.

$^{121}$但先等一會,.
並且用一個問題扭轉了所有人的目光,.
經文記載,.
忠於耶穌,.
他直起腰泥對眾人說,.
誰若無罪,.
就可以率先拿起石頭,.
扔行淫婦人,.
很奇怪,.
耶穌說完後,.
有沒有催逼他們回答?.
沒有,.
耶穌又彎腰在地上寫字,.
雖然眾人催逼他,.
但他沒有催逼眾人回答,.
他將球交給群眾,.
然後耐心等他們自行判斷這件事,.
這是耶穌處理的方法,.
至於群眾,.
當他們聽到耶穌這個提問時,.
整個氣氛不同了,.
你不難想像,.
由很嘈吵的環境,.
突然全場靜了,.
就像現在這樣,.
靜了,.
並且經文仔細記載,.
群眾從老到少,.
逐一離開,.
最後只剩下耶穌和婦人,.
於是耶穌再次直起腰泥,.
他表示他不去定婦人的罪,.
並且打發她離開,.
祝福她今後不要再犯罪,.
故事說到這裡很簡單,.
不過有幾件事值得我們細心看,.
首先你留意耶穌沒有否定摩西律法在群眾中的權威性,.
不過他透過一個我們說法人心醒的問題,.
讓群眾看到婦人問題的同時,.
也看到自己的本相是什麼,.

$^{161}$本質上耶穌這個問題是提醒大家,.
其實我們都不符合上帝的完全聖潔標準,.
或者群眾確實沒有犯奸淫的問題,.
但是苦心自問,.
有誰夠膽說自己已經100\%.
全部遵守上帝的吩咐和律法,.
對上帝對人一點虧欠都沒有呢?.
耶穌說如果真的有人可以這樣自誇的話,.
不妨第一個走出來向這個婦人掟石頭,.
耶穌是在問一條這樣的問題,.
由原本眾人只是聚焦在.
我們要怎樣解決眼前這個犯婦人,.
怎樣和她割蓆或清界線,.
耶穌用一個問題來去叮醒當時在場所有人,.
大家看清楚一點,.
其實在上帝的眼中,.
在上帝的標準裡面,.
我們全部都是在同一格,.
同一個水平,.
如果我們覺得這個婦人是罪有應得,.
應該向她掟石,.
其實亦都意味著我們應該受上帝相同的刑罰,.
所以大家不能夠寬己嚴人,.
事實上如果上帝真的拿正來做的話,.
按照律法的標準去量度個人的話,.
大家最終都和婦人同一條路,.
就是死路,.
沒有分別,.
結果在場所有的人,.
包括那些原告,.
即是那些人,.
包括那些目擊證人,.
即是民事和法理財人,.
連他們都離開,.
因為他們知道自己,.
無論道德的情操可能相對比別人優越,.
但是自己同樣都有軟弱的時候,.
都有達不到律法標準的時候,.
甚至乎都試過沒有遵守上帝的吩咐,.
大家都意識到其實自己在人生裡面,.

$^{201}$多多少少都得罪過人,.
虧欠過上帝,.
特別這裡你和我都會很有共鳴,.
從年長到年少,.
逐個離開,.
年紀越大的人,.
越有這份虧欠的體會,.
每一個人都經歷過,.
所以你看到,.
民事和法理財人就算他多麼不滿耶穌也好,.
多麼想消滅他也好,.
他都不能否認世人都犯了罪,.
規缺了上帝榮耀這個事實,.
沒有人能夠否定,.
我們其實都很明白,.
作為新約的信徒,.
我們知道舊約的律法,.
其實它有一個功用,.
就是叫人知罪,.
但是律法本身並不能夠將人從罪中拯救出來,.
唯有靠誰?.
你和我今天很明白,.
如果有個新朋友問你,.
你一定答得到,.
靠誰?.
靠耶穌基督的恩典,.
靠祂的救恩,.
人才能夠徹徹底底,.
從罪惡中被釋放出來,.
可以和上帝和好,.
經文其實正正將這個真理,.
很徹底精彩地活現在,.
不單單是婦人身上,.
其實也活現在在場所有人身上,.
為什麼這樣說?.
你留意細心看這段經文,.
耶穌其實不單單沒有定婦人的罪,.
他有沒有定民事和化理錯人的罪?.
沒有,.
如果他拿上上帝的標準去看,.

$^{241}$其實他可以定民事和化理錯人,.
那個鬼渣偽善的罪,.
對不對?.
耶穌有沒有追究其他人的罪?.
你說其他人有什麼罪?.
有什麼罪?.
吃花生的罪,.
用今天的角度看,.
有人需要幫忙在街上跌倒,.
或者有人吵架,.
大家可能見到其他人會怎樣?.
不是第一時間幫忙勸架,.
而是拿出手機拍下,.
上載,.
然後分享,.
很奇怪的現象,.
對不對?.
很多時候我們會袖手旁觀,.
但是耶穌沒有去定這群袖手旁觀的群眾的罪,.
你想像一下,.
多多少少他們也受過耶穌的恩惠,.
耶穌有事有沒有出來說兩句話解圍?.
也沒有,.
又或者這個婦人多麼壞,.
他們也沒有來代祂求情,.
希望人們可以重慶快樂,.
他們很冷漠,.
沒有憐憫之心,.
但是耶穌也沒有定他們的罪,.
不單單這樣,.
耶穌疼愛到他們一個程度,.
就是怎樣?.
幫他們找下台階,.
是不是?.
耶穌沒有當面揭穿這一切罪行,.
反而只是問一個問題,.
讓他們自己審查自己,.
然後離開,.
找下台階給他們,.
你看到耶穌沒有因為眾人的過犯,.

$^{281}$甚至乎在當中人們對他的漠視,.
對他的頂撞,.
而停止他對罪人的接納和寬恕,.
耶穌沒有因為這樣停止對罪人來施恩憐憫,.
不單單這個婦人,.
包括在場所有人,.
包括你和我,.
耶穌對人這份恩慈,.
這份憐憫,.
其實某程度上也提醒,.
挑戰我們,.
我們有沒有帶著上帝這份恩慈,.
來看待那些犯錯的人?.
當別人犯錯的時候,.
你看到,.
甚至乎得罪了你,.
激怒你的時候,.
你在心裡面,.
你會怎樣回應?.
有沒有留意,.
其實你在心裡面有兩個按鈕可以按,.
第一個按鈕就是定罪,.
擲石的按鈕,.
一按著它,.
就好像你開車踩油門,.
一去零到一百,.
砰一聲就會發怒,.
另一個按鈕叫恩慈憐憫,.
你按著它的時候,.
你的接納程度會大很多,.
當面對別人犯錯,.
得罪你,.
甚至激怒你的時候,.
你會首先按哪一個按鈕?.
我記得上星期或前星期,.
我兒子打電話給我,.
我正在工作,.
他很有趣地輕描淡寫,.
「爸爸,家裡發生一件很大的事,.
好像報道來的,.

$^{321}$這是報道,.
報道發生什麼事件,.
我知道其實是他出事,.
肯定是,.
不過他好像一個局外人,.
只是描述這件事件的發生,.
我說什麼事?.
我又很輕鬆地回應,.
「家裡我打爛了一瓶豉油,.
整地都是,.
不過你放心,.
我已經搞定了.」.
很開心,.
報道了這件事件,.
於是我就說,.
「我是,.
那你有沒有受傷?.
沒有受傷吧?.
最重要,.
不要緊,.
等我回來再說.」.
於是兒子就掛了線,.
但掛了線之後,.
我會想,.
其實這個對答都很奇妙,.
都很感恩,.
因為不是正常的我,.
你想想,.
如果我身處現場,.
目擊這個案發的事件,.
我還要很忙,.
做其他事的時候,.
你猜我會先按著哪個按鈕?.
我一定立刻零到一百,.
踩油門就會起火,.
不是吧?.
我已經很忙,.
還添煩添亂,.
這麼小的事,.
都搞到一塌糊塗,.

$^{361}$然後就走開,.
等我收拾,.
怎會這麼好?.
「哎呀,有沒有事?.
有沒有受傷?」.
不會的,.
我只會為了十多元的東西,.
基本上會和他擲石,.
和他割席,.
就是那一刻,.
到我發洩完之後,.
我才會懂得按恩慈的按鈕,.
「下次不要了,.
要求主憐憫.」.
擲錢會其實,.
如果不是聖靈的保守帶領,.
我想我不會關心我兒子的狀況和需要,.
我可能只會關注,.
我怎樣收拾這件事,.
和收拾犯錯的人,.
不知道你有沒有這些經驗和體驗,.
我想唯有聖靈的幫助,.
我們才能夠用上帝的恩慈,.
來看待別人的過失,.
寬恕接納別人的過犯,.
就正如耶穌,.
他沒有因為這個婦人的犯錯,.
第一時間就推走她,.
第一時間就和她割席,.
甚至要她消失在眼前,.
這是問世和法理錯人處理罪人的方法,.
將骯髒的人從社群中剔除,.
完全的撤除,.
那又怎樣?.
自己就可以維持心目中的聖潔和乾淨,.
他們是這樣想,.
但是耶穌卻選擇和這個罪人連結在一起,.
主動向他施恩憐憫,.
並且逐步幫他脫離罪惡,.
走向聖潔,.

$^{401}$很不同吧?.
法理錯人的聖潔,.
所說的聖潔是排外,.
將罪人排走,.
耶穌所說的聖潔是將罪人拉回來,.
幫他越來越近上帝的聖潔,.
兩個不同的系統,.
是因為上帝猜他的兒子降生,.
不是要定人的罪,.
而是要世人因為他得生得救,.
所以這段經文提醒我們,.
在生活中,.
我們不單要意識到,.
以上帝的恩慈,.
來對待那些有過失,.
軟弱的人,.
更加要學習怎樣操練,.
簡單說,.
我們在看見別人的過犯,.
很想扯火,.
很想問責,.
扔石的同時,.
能否看到上帝的恩典和憐憫,.
我們怎樣操練時間距離,.
將它越縮越短,.
曾經聽過一個牧師的故事,.
在他辦公室的桌上,.
放了一個用石做的紙陣,.
不單紮著文章和紙張,.
在石做的紙陣上,.
寫了一個很大的「先」字,.
先後的先,.
你明白了,.
提醒他一個很好的操練,.
每次看見就提醒自己,.
想起耶穌那句,.
你們誰沒有試過犯錯,.
就可以先拿起石頭來扔人,.
一來可以幫助自己抑制犯罪的意念,.
二來可以克制自己,.

$^{441}$很想批評論斷人的衝動,.
透過這個小小的日常操練,.
他發現,.
他遙說接納別人的程度,.
慢慢會有所提升,.
這個紙陣也有另一個描用的方法,.
當有弟兄姊妹走進他的房間,.
怒氣沖沖,.
很想批評別人的時候,.
這個牧師聽完他說之後,.
就會找一個適當的時機,.
將那塊石做的紙陣放在他的手上,.
然後很耐心等聖上發露扔石的人,.
來細心看這塊石頭,.
當他發現那裡寫了個「先」字的時候,.
牧師就會邀請他說,.
其實過鐘是說什麼意義?.
讓他想起昔日耶穌這個教導,.
就是當我們看到別人行差踏錯,.
很想問罪輕施的時候,.
想扔石的時候,.
盼望我們都能夠想到,.
自己其實都是蒙恩的罪人,.
能夠想到耶穌期望我們,.
在領受恩典的同時,.
也都要學會對別人施恩憐憫,.
從而寬恕那些和我們一樣,.
有軟弱過繁的人,.
給一個改過的機會對方,.
正如耶穌給一個改過的機會,.
這位婦人一樣,.
並且如果你細心看這段經文,.
你有沒有留意,.
耶穌是不是只給一次機會她?.
你猜一下,.
問得很奇怪,.
我很喜歡這樣去問聖經裡面的問題,.
雖然經文沒有直接說,.
耶穌給了多少次機會她,.
但是你會留意到,.

$^{481}$耶穌會不斷幫助她,.
為什麼這樣說呢?.
耶穌最後怎樣跟她說?.
耶穌說:去吧,.
從此不要再犯罪了,.
那她會不會再犯罪?.
耶穌有沒有說:去吧,.
從此你不會再犯罪了?.
沒有,.
耶穌說:從此你不要再犯罪了,.
這句話一點法力都沒有,.
沒有法力,.
是信任,.
是耶穌看好她,.
是耶穌會幫助她,.
婦人今後仍有機會面對試探,.
特別可能是性方面的試探,.
仍有機會犯罪,.
甚至再犯姦淫的可能,.
於是我們開始會問,.
為什麼耶穌會這麼兒戲,.
這樣就放過這個人?.
起碼應該怎樣做?.
大家一起想想,.
比將利今天可以怎樣確保她不會再犯罪?.
大家立即靜下來,.
很容易就這樣,.
起碼都寫封什麼?.
小學生?.
是呀,寫封悔過書,.
徹底的懺悔,.
或者罰抄,.
我以後不犯姦淫,.
抄一千次,.
一萬次,.
抄到你都受不了,.
不難的,.
我們一定要有憑據,.
才可以相信那個人真的會改過,.
但你看到耶穌沒有,.

$^{521}$沒有叫他寫悔過書,.
然後要求他,.
你先守兩年行為,.
這兩年我會不斷監視你,.
看你有沒有重犯姦淫,.
沒有我才真正放過你,.
不再追討你的罪,.
耶穌沒有這樣做,.
又或者耶穌沒有仿效中國古代,.
那些犯婦人,.
或者淫婦在面前做什麼?.
刺她清,.
大大隻字,.
淫婦,.
沒有了,.
斷正了,.
你自己都不想出街,.
還有什麼機會犯姦淫?.
但可以在她面前,.
落入一個不可磨滅的印記,.
讓她羞愧,.
這一世都不用出街,.
或者人們見到她都避開,.
耶穌亦沒有這樣做,.
又或者再高明一點,.
好像武林高手,.
耶穌可以怎樣做?.
我廢你武功,.
好像以前那些人對付.
很厲害的,.
懂武功的壞蛋,.
我廢了你武功,.
一掌打在天靈蓋上,.
耶穌亦有辦法求上帝,.
廢她七情六慾的根,.
不難,.
真的確保了,.
這一世她都不會再犯罪,.
不會再犯姦淫的罪,.
但耶穌都沒有這樣做,.

$^{561}$為什麼耶穌可以這麼便宜她?.
好值得思考,.
為什麼耶穌有時這麼便宜我們,.
放過我們一些過犯,.
不用這個婦人付錢,.
聽節目其實你和我都明白,.
是因為誰付了錢?.
是耶穌親自幫她付了那筆錢,.
耶穌親自為她付上了,.
付清了那個罪債,.
在福音書後段我們會看到.
耶穌親自和這個婦人,.
並且和所有的罪人交換了位置,.
代替他們釘死在十字架上,.
承擔了罪所帶來的羞辱和死亡,.
使我們明白上帝的恩典恩詞,.
由始至終在律法之上再加上去,.
是在我們的行為以外,.
額外顯明出來,.
使我們明白人之所以能夠得救進天國,.
不是因為我們的行為表現得多好,.
不是因為我們能討上帝有多多的喜悅,.
而是完完全全出於上帝的恩典,.
所以耶穌沒有等這個婦人改好,.
「我原諒你,你有好的行為,我才接納你.」.
耶穌是首先拯救了這個婦人,.
然後逐步幫助她改變生命,.
所以一定不是一次的寬恕,.
一次的拯救和幫助,.
而是每一次這個婦人軟弱跌倒,.
再犯罪,得罪上帝的時候,.
只要她願意來到上帝的面前悔改,.
耶穌都願意拯救,救到底,.
同樣,神亦不會因為我們的過犯,.
而選擇離棄我們,.
上帝卻是不離不棄地愛我們,.
並且祂看好你和我,.
祂知道靠著祂的恩典和幫助,.
我們終有一日,.
我們的生命會成長並且成熟,.

$^{601}$能夠符合滿足上帝的期望..
所以即使明知這個婦人在改變的過程中,.
她有可能再次遠離上帝,.
有可能再次回到罪中,.
將自己捆綁回到罪的當中,.
但耶穌亦都會願意向她施恩幫助,.
一次又一次地替她解鎖,.
讓她可以從罪惡中重新出來,.
慢慢一步一步走回近上帝..
耶穌不會因為這個婦人糟糕的生命,.
就將她看死,看低..
曾經聽過一個故事是這樣說的,.
有位大學教授,.
有一次他來到鄉村的地方,.
在當地對一群小朋友的表現,.
做了一份評估報告,.
他發現這群小朋友的成長條件一般般,.
幾乎如果用常理來分析,.
他們不太可能有出人頭地的一天..
二十五年後,.
這份報告被另一個年輕學者無意間發現,.
他很好奇,.
想知道這群小朋友現在會變成怎樣,.
於是他就回到當時的地方,.
再做一次追蹤和調查..
他發現這群當初沒有什麼人看好的孩子,.
長大之後有一半都出人頭地,.
有些做了醫生,.
有些做了律師,.
有些做了商人,.
有些做了牧者..
他覺得很驚訝,.
不明白為什麼,.
於是他繼續追查下去,.
他發現這群受訪者他們有一個共同答案,.
就是他們遇到一個很好的老師..
於是這位學者就從他們的口中找回當年這個老師,.
發現她只是一個很平凡的女士,.
沒有什麼特別高的學歷,.
也沒有什麼顯赫的權勢,.

$^{641}$更加沒有多餘的資金來補送這群小朋友去接受高等教育,.
甚至幫他們創業..
這位學者便請教這位女士,.
你是怎樣栽培這群小朋友,.
你的秘訣是什麼?.
她笑著說,.
其實沒什麼,.
你和我都做得到,.
只不過我懂得不斷地愛他們,.
不斷地鼓勵他們,.
他們做錯的時候我包容他們,.
慢慢給他們時間,.
給他們很多機會來改變成材,.
就是這麼簡單..
同樣地,你看看這位婦人,.
相信她不是不知道行淫是不好的,.
她也知道這樣會得罪上帝,.
上帝不喜歡,.
不過改不了,.
你和我都很明白,.
我們都知道很多事情上帝不喜歡,.
不過都做了,.
總是改不了,.
靠自己真的改不了,.
靠自己永遠都像這位婦人那樣,.
擺脫不了一些試探,.
她需要一位無條件接納她,.
愛她,.
關心她,.
並且有能力的上帝來幫她,.
回到那條正路裡面,.
而這位神願意於磚降跪,.
從天上降臨到人間,.
來顯明在婦人面前,.
同樣也顯明在你和我的面前,.
正如耶穌不是不知道我們有些事情不想,.
我們都很希望做一個上帝喜悅,.
百分百完美無瑕的兒女,.
不過做不到,.
或者有時可以,有時不可以,.

$^{681}$但是上帝不要因為這樣而將我們看低,.
祂仍然看好我們,.
祂只是知道我們未做到,.
並且祂願意愛我們,.
在我們的生命裡面不斷挽回,.
鼓勵,幫助我們,.
為我們帶來改變的勇氣,.
力量和盼望,.
所以假如今天我們仍然被某些罪,.
纏繞,多番掙扎,.
還是不能脫身的話,.
我深信上帝對當日婦人的宣告,.
也是對我們的宣告,.
耶穌說:.
「我也不定我們的罪,去罷,.
從此不要再犯罪了.」.
耶穌也是對我們說,.
對我們一份相信,.
一種歡恕,接納和肯定,.
真的,弟兄姊妹,.
耶穌很愛你和我,.
祂不惜用重價將我們買回來,.
祂不惜犧牲祂的兒子,.
來將我們買熟,.
並且祂會看好我們,.
會幫助我們改變我們的生命,.
耶穌已經將這份恩典賜給婦人,.
祂也要將這份禮物賜給你和我,.
你和我已經袋了,.
不過袋不袋得實呢?.
袋不袋得穩呢?.
可否隨時在你軟弱的時候,.
拿出來,.
來幫你站在上帝那邊,.
這是關鍵,.
需要去操練,.
你願不願意相信,.
耶穌真的愛你愛到這麼徹底,.
以至你每次倚靠祂的時候,.
你都可以得勝罪惡,.

$^{721}$你都可以重新開展你的人生,.
經文這樣去勉勵我們,.
最後再說一個故事,.
話說有個畫家的兒女,.
她畫畫很厲害的,.
有一天畫了一幅很漂亮的圖畫,.
畫裡面有很多盛放的蘭花,.
非常雅致,.
她看到之後很開心,.
於是準備收拾她的畫具,.
不過卻不小心將那支黑色顏料的筆,.
甩了一頁,.
於是那些很無情的黑色顏料,.
就像灑了幾滴雨一樣,.
落在那幅很漂亮的,.
盛放蘭花的圖畫上,.
她覺得很難過,.
因為補救不了,.
不知道怎麼辦,.
於是她就流下眼淚,.
她的爸爸聽見就從房間裡跑出來,.
了解了原因之後,.
她就看著這幅已經面目全非的畫,.
想了一會兒,.
於是突然她笑了一笑,.
她就馬上拿回那支,.
沾滿了黑色顏料的和筆,.
拿出來,.
瞬間在她女兒的畫上,.
撩了幾筆,.
奇妙的事就發生了,.
那些原本在畫裡的敗筆,.
突然在她父親的描筆上,.
加工了幾下之後,.
就變成一隻隻快樂飛舞的蜜蜂,.
很奇妙,.
將整幅以花園為主題的圖畫,.
再一次增添了出乎意料的生命力,.
丁姐妹,其實我們人生都是這樣,.
有沒有想過?.

$^{761}$我們平時都像女孩子的角色,.
總是一個不留神,.
就將你的生命,.
怎樣?.
弄花了,.
弄到面目全非,.
差點補救不了,.
但只要你肯來到耶穌的面前,.
讓祂介入,.
改變你的生命,.
那個幾近乎一拜一塔糊塗的生命,.
在耶穌的描筆裡,.
可以有生機,.
可以再次散發出人意外的生命力,.
你相信嗎?.
你相信耶穌可以這樣嗎?.
真的,各位朋友,.
丁姐妹,.
很盼望這段經文能夠幫助我們更加了解,.
體會到耶穌對我們的愛是多麼的長闊高深,.
可以經歷到耶穌的恩典,.
並且能夠用上帝的恩賜,.
來對善待,.
祝福身邊更多的人,.
好不好?.
讓我們一起有個祈禱吧!.
親愛的恩主,.
看到我們原本是多麼的罪孽的心眾,.
重到一個地步,.
耶穌你要親自到城獄身,.
來幫我們還債,.
不可以草草了事,.
同樣你的降生,.
你的捨命也讓我們看到,.
你有多麼的愛我們,.
愛到一個的緣故,.
你不會假手別人來承擔我們的重擔和罪債,.
你願意親自以你的生命,.
來換回我們的生命,.
上帝,我們真是很切實看到你這份的心恩厚愛,.

$^{801}$特別藉著這段經文讓我們重新提醒,.
其實我們每一個都是罪人的本相,.
沒有上帝你的恩典,.
我們根本不能夠存活,.
更何談討你的喜悅?.
上帝幫助我們藉著這段經文,.
我們看到自己,.
看到自己的軟弱,.
看到自己需要耶穌你的救恩,.
並且縱然我們可能覺得被罪困鎖,.
擺脫不了,.
甚至可能人生已經一塌糊塗,.
但我們知道上帝仍然看好我們,.
沒有看壞我們,.
上帝你要我們知道,.
靠著你的恩典,.
我們可以翻身,.
我們可以像這個婦人一樣,.
過新的生活,.
或許我們仍然會跌倒,.
面對試探,.
但我們知道靠上帝你,.
必定最終能夠走到終點,.
因為不是我們自己有什麼本事,.
而是上帝你總會親自挽回我們,.
幫助鼓勵我們,.
上帝讓我們存著這份盼望,.
以致我們盡我們的努力,.
好像聖經所說,.
盡心盡性盡義,.
遵行你的吩咐,.
讓我們努力來討你的喜悅,.
亦相信上帝你會幫我們走到你的懷裡,.
並且藉著這段經文,.
讓我們去善待身邊同樣有軟弱過犯的人,.
讓我們在生活中操練,.
有更多的恩慈來彼此相待,.
讓我們同樣對人心存盼望,.
我們可能會有時被人激怒,.
甚至有時會對人失望,.

$^{841}$盼望上帝你在當中介入,.
讓我們重新找到站立起來的力量,.
讓我們不灰心,.
因為我們知道上帝你不灰心,.
我們這樣禱告交託,.
奉基督耶穌你得勝的名字來祈求,.
阿們..
\newpage



\section{詩篇利未記箴言 19:32}
\label{sec:d4oKfBQrMaI}
\textbf{2020年11月22日 崇拜直播(長者主日)｜利未記19\_32;詩篇71\_9,18;箴言20\_29}
\newline
\newline
連結: \href{https://youtube.com/watch?v=d4oKfBQrMaI}{\texttt{https://youtube.com/watch?v=d4oKfBQrMaI}} ~~~~ 語音日期: 2020-11-23
\newline
\newline
\hyperref[sec:PcbwkOrPy_0]{\small{< < < PREV SERMON < < <}}
~
\hyperref[sec:index_chronic]{\small{[返順時目]}}
~
\hyperref[sec:CjzmC7ulsYw]{\small{> > > NEXT SERMON > > >}}
\newline
\newline
$^{1}$經文化及文化局局長.
利未記十九章三十二節.
在白髮的人面前.
你要站起來.
也要尊敬老人.
又要敬畏你的神.
我是耶和華.
詩篇七十一篇九節.
「我年老的時候.
求你不要丟棄我.
我力氣衰弱的時候.
求你不要離棄我」.
詩篇七十一篇十八節.
「神啊!我到年老發白的時候.
求你不要離棄我.
讓我將你的能力指示下代.
將你的大能指示後世的人」.
真言二十章二十九節.
「強壯乃少年人的榮耀.
白髮為老年人的尊榮」.
(十二節).
等子妹子女平安.
剛才長者詩班和我們一同唱的一首詩歌.
不知道大家覺得如何.
很輕快.
勾起了我以前在教會團契.
很喜歡唱的一首詩歌.
我想問在座當中.
有沒有長者以下的.
請你舉手.
長者以下.
這麼少的嗎?.
每個都長者嗎?.
長者以下的請你站起來.
請你站起來.
在白髮人的面前.
你要站起來.
我們要向在座當中的長者鼓掌.
好.
請坐.

$^{41}$這是聖經教導.
剛才李美記跟我們說的.
但很多時候在現實生活中.
未必是這樣.
特別是在繁忙的香港.
我們每天上班.
當大家在紅綠燈頒馬獻的時候.
很心急要過馬路.
趕著上班的時候.
綠燈一轉的時候.
大家一起急步走.
誰知前面有一個.
有一個婆婆.
伯伯.
他動著拐杖左右搖擺的時候.
你的心會怎樣?.
會否說.
「阻礙著」.
如果旁邊的年青人.
大家鼓掌.
當你看到這樣的情景的時候.
你的心裡是怎樣呢?.
會否有些不屑.
為何對老人家這麼不尊重呢?.
你走慢一點.
又或者在排隊.
今早真的在排隊買早餐.
前面有一個上行者.
他不是很掌的長者.
在那裡等很久.
左一間 右一間.
旁邊有一個.
其實他也是差不多的長者.
不過他又說.
「阻礙著了」.
「你看後面排了十多二十人」.
「其實後面只有六七人」.
他也是這樣說.
我們會否嫌老人家步伐慢呢?.
他們想事情要想很久.

$^{81}$當我預備這篇信息的時候.
我腦海裡也浮現了這些情景.
有時候我會去反省一下.
我會犯了這些毛病.
這些錯誤.
實在我們很需要上帝來憐憫我們.
在長者主日裡.
讓我們更新對長者的態度.
其實我自己也步入長者的.
上行者的年齡裡.
雖然未至於很緩慢的步伐.
但上帝給我一個恩典.
可以在教會裡.
一同來服侍大家.
特別是上行者牧區.
和長者部.
一同來服侍.
我們又回過頭來說.
不知道大家對「老」有什麼定義.
或者怎樣才算「老」呢?.
很多時候.
我們會從社會學的角度來看.
人的生理變化.
可能身體會有多病.
身體的機能會退化.
甚至我們說回功能.
很多事情.
因為年齡的時候.
社會的角色要我們轉變.
可能還可以繼續做.
多做幾年就做.
但結果.
社會的功能.
可能要退下來.
另外醫學.
我想大家都知道.
發展得很快.
甚至有生物老齡醫學.
去探索關於生物衰老的過程.
去研究一下.

$^{121}$人怎樣可以延長壽命.
甚至會在衰老上減慢.
這個我想大家都很喜歡.
去追求長壽.
但人就想長壽.
不想老.
大家想不想?.
我想這是現代人.
在我們現代社會裡面.
對於在老齡化裡面.
其中一個很矛盾.
和很有張力的.
這種情況.
在很久之前.
在峇里島裡面.
有一個很古老的傳說.
他說在那裡有一個很遙遠的山區.
村裡面住了很多村民.
整條村在那裡.
但他們有個習俗.
他們常常把老人.
一去到那個階段.
就把他們獻祭給神明.
久而久之.
結果那條村裡面.
一個老人家都沒有.
文化傳統.
那些智慧都漸漸消失.
後來這班人.
他們很想做一些事.
要團結一下.
想起一間很大的屋.
但發覺他們把木材都砍了下來.
結果他們不知道怎樣分木材.
頭和尾.
分出來去建屋.
因為亂建.
可能結構不穩固.
會不會漏水.
這些都有問題.

$^{161}$他們想辦法怎樣解決.
因為他們知道.
應該是老人家很有經驗.
但他們沒有.
怎樣呢.
有一個年輕人跟他們說.
如果你們不把老人家.
獻祭給神明.
我就會幫你想辦法解決這個問題.
你猜他怎樣解決問題呢.
他終於在家裡.
把老人家帶出來.
給父老和村民看.
原來當這些事發生的時候.
他就把祖父收藏得很隱蔽.
結果這位老人家出來之後.
就告訴他們怎樣分辨木材.
頭和尾.
於是就建了一間很穩固的大屋.
甚至長者都把他們的傳統智慧.
人生歷練.
分享給一班村民.
結果那班村民.
維持了大家彼此良好關係.
互助互愛.
對長者有一種新的體會.
對他們尊敬有加.
有研究推算到2036年.
每四個人就有一個長者.
在座當中的剛才站起來的.
十多年後可能輪到你了.
趁著我們還沒有加入長者行列之前.
其實我們要學習一下.
認識老年那種狀況.
從而擁抱老年.
繼而進入敬老.
愛老 護老.
甚至可以樹立一個榜樣.
幫助一些一直成長階段的人.
他們有所學習.

$^{201}$因為不久將來.
他們也會體會到.
練老練長.
所以其實就要在這個階段學習.
敬老 愛老 護老.
剛才姊妹跟我們讀這段經文.
利未記十九章三十二節.
在白髮人面前站起來.
也要尊敬老人.
又要敬畏你的神.
我是耶和華.
其實這裡是講到.
要擁抱老年的理由.
留意十九章第二節.
第二節是這樣說.
你們要聖潔.
因為我耶和華你們的神是聖潔的.
這就是傳捲利未記的鑰節.
這一節所指向的對象.
其實不單單是利未人.
或者是祭司.
甚至是全以色列會眾.
其實在這裡的經文就是要他們聖潔.
耶和華說他們的神是聖潔.
這是利未記重要的神學框架.
很多人認為利未記只是記載.
利未人如何做祭司.
這一本手冊.
可能其實有些誤解.
當中當然是耶和華神.
藉著摩西頒佈祭司的即時.
和憲制的條例.
以及如何區分潔淨和不潔淨的物件.
以及講到關於聚德屬.
這種祭祀的律例.
其實同時也是利未記當中提到.
關於以色列全會眾一個道德.
律例,律法,祭祀和百姓精神生活的規範.
經文中同時也要求以色列全民要神聖.
因為神說你們要聖潔.

$^{241}$其實在十九章所列明的條例.
這是道德的條例.
神的聖潔領域要延伸到每一個以色列的全會眾.
要讓以色列活出敬虔和聖潔的生活.
因為他們是神國道的身份.
所以他們同時既有這種身份的時候.
就可以分享到神聖潔的屬性.
整篇十九章有三十七節的經文.
在利未記裡面十九章.
我們仔細去看的時候.
其實你會看到很鬆散.
或者好像很雜亂.
但我們仔細去看的時候.
它分了十六小段.
每一段都是以「我是耶和華」作為結束.
而這個內容其實就是.
十戒在日常生活上的形容.
孝敬父母是十戒中的第五戒.
和敬拜神手安息日是十戒中的第四戒.
同時出現在第三節那裡.
表明了父母是神在地上的代表.
應該受到充分的尊重.
在新約裡面亦不時的重述了這個命令.
無論我們的父母出生.
或者社會地位是怎樣.
又或者他們的身體精神狀況是怎樣.
其實我們作為兒女都需要敬愛他們.
按著聖經的教導.
這是人最重要的責任.
三十二節將尊敬老年人,白髮人.
和敬畏神的命令並列.
其實這是呼籲了.
十九章第三節敬畏父母.
或者是孝敬父母的命令.
是要求以色列人要在白髮人.
老年人面前站起來.
我們會尊榮我們的長者.
其實是表示我們對神的敬畏.
因為他們是按著上帝的形象去做.
不要因為他們已經是年紀老邁.

$^{281}$我們就小看他們.
甚至忽略他們白髮為老年人的尊榮.
這是真言裡面最有智慧的教導.
我們要愛神所愛的.
要尊敬神所尊敬的.
有新聞是這樣報導.
一個年老的婆婆傷心欲絕.
她在自己住戶的單位門外打地鋪.
這件新聞雖然不是很轟動.
但也有記者採訪.
這個婆婆可以上樓.
因為她的家庭.
她和兒子,媳婦和兩個很小的兒孫女.
用一個家有長者優先配屋計劃.
來申請上樓.
所以整個家庭一家幾口.
很快就上樓了.
但這個婆婆因為年紀大.
體弱多病.
先後都進出醫院.
到最後兒子覺得要照顧家庭.
小孩還未長大.
要照顧媽媽.
比較吃力.
所以和他的姐弟商量.
之後就決定.
把老人家婆婆送進老人院.
他們都會承擔一些費用.
經常去探望她.
但原來因為這班年青人.
子女出現了很多分歧.
特別是財務上的問題.
有沒有交老人院費.
福利署介入都很困難.
婆婆被迫在單位住戶外.
去住宿.
有家歸不得.
這種被遺棄.
這種難過 孤單 傷感.
心情極度不安.

$^{321}$莫說老來從子.
可能已經不再想這件事.
如何解決面前的處境.
當我看到這些情景時.
一方面是同情她.
另一方面為了這個婆婆有點難過.
現實生活中.
很多長者除了身體衰弱之外.
更加害怕什麼呢.
害怕被家人甚至社會離棄.
浸信會外群社會服務處.
做了一個對長者的調查.
在1617年.
調查關於長者身心健康.
情緒健康的調查.
得出來的結果是怎樣呢.
其中有和家人關係不滿的長者.
一定會有負面情緒.
心情也會不好.
但這個情況其實很嚴重.
有百分之五十受訪.
一提到家庭不滿.
長者的情緒負面原來有這麼多.
另外提到參加活動的長者.
如果不參加活動的長者.
負面情緒比參加團體活動的長者多10\%.
即是說回教會的長者.
負面情緒指數應該很低.
其實還有其他高危因素.
還有經濟壓力.
跟老朋友老友記的關係不好.
甚至有時會覺得自己記性不好.
會帶來很多麻煩.
不便.
你忘記東西.
一出門就看看有沒有關掉爐子.
各樣事情.
又感覺到自己沒有希望.
當中這個調查.
其實沒有包括信仰作為調查項目.

$^{361}$但聖經裡告訴我們.
年長者受到身體衰老.
病患.
情緒有時會表達出驚惶恐懼的描述.
長者在這種狀態裡.
他們如何面對老化呢?.
長者原來可以向上帝.
向神去呼求.
在詩篇71篇第9節18節裡說.
我年老的時候求你不要丟棄我.
我力氣衰弱的時候求你不要離棄我.
神啊我到年老發白的時候.
求你不要離棄我.
詩篇71篇是一篇老人的祈禱詩.
當中.
當你剛才讀了這兩字經文的時候.
其實就反映了.
我們步入老年那種被遺棄的恐懼.
不單單可能是被家人.
被朋友遺棄.
原來他最害怕的是.
覺得被神離棄.
我年老的時候求神你不要離棄我.
我體弱的時候求你不要丟棄我.
求你不要離棄我.
速速幫助我.
遇到身體病患衰弱.
我們的能力一定沒有以前那麼好.
詩篇裡面提到仇敵.
當我們一路看經文的時候.
是會看到仇敵這兩個字.
年老的詩人其實沒有明確地指明.
是不是過去一直與他為敵的敵人.
抑或是他現在.
現時年老.
多病 體弱.
一無是處.
甚至失去了年青的活力.
自覺沒有用這種苦況.
甚至沒有辦法服侍神.

$^{401}$覺得好像被其他人圍瞞.
就算在信仰群體裡好像被丟棄一樣.
但詩人接受現實這種人生.
他的生命老化.
頭髮都白了.
力氣都衰弱的事實.
但他有一個信念.
他因為有信仰.
他一生所倚靠的神.
就是由他在無福之中.
神都一直保護他.
還有神是與他一起面對大大小小.
他多多的患難 危機.
在這些過程裡他都能夠經過.
經歷了各種人生的階段.
他都很平安順利地渡過.
因此詩人仍然相信神是他的保障.
雖然到了這個年紀.
所以這位詩人向神呼求.
求神不要丟棄他.
不要離棄他.
我相信這是詩人過去的經驗告訴他.
他必須抓住神.
要好好地抓住上帝.
還有他要懷著盼望.
到了第十七節的時候.
詩篇七十一篇.
詩人就這樣說.
「直到如今,我傳揚你奇妙的作為」.
接著十八節的時候就說.
「神啊,我到年老發白的時候.
求你不要離棄我.
讓我將你的能力指示下代.
將你的大能指示後世的人」.
這位詩人就是將神的能力.
傳揚到下一代.
向下一代宣告.
這表示什麼呢?.
就是這位年長的詩人.
他要用他活生生的見證.

$^{441}$親自來為神做見證.
因為他知道.
雖然體弱,年老.
不能夠限制他.
年老只是限制了他身體的活動.
但不能阻止他和別人宣講.
在他人生經歷裡.
看到神作為這些事情.
又或者我們從另一個角度看.
詩篇七十一篇.
其實是導向我們來看看.
我們整個人生成長的過程.
因為經文裡有提到.
「從母腹生出」.
就是說在他不同年齡的階段裡.
詩人也知道.
一路成長.
當他成功跨越的時候.
這是自然的進程和規律.
特別說到年老發白的時候.
那種生命歷程是怎樣呢?.
我相信這就是他對生命老化.
和擁抱老年的呼聲.
因為經文一路提到.
詩篇告訴我們.
這位年老的詩人.
以積極面向成為他的動力.
人生必須經過老化的進程.
他能夠見證神的榮美.
和去頌讚神在人身上的作為.
這是上帝給他的恩典.
最終.
老人家這位年老的詩人.
能夠為他自己在神的靈魂有所歸宿.
而來歡呼.
這也是年老長者最所盼望的.
馬克亞瑟基金協會.
做過一個關於成功老化的醫學報告.
報告可能會很長.
但他認定了一句.

$^{481}$聖經很肯定地說.
人的習慣和他所選擇的.
對他們如何度過老年生活.
是有很大的影響.
成功老化是一個很大.
很廣闊.
很複雜的課題.
亦涉及很多專業.
如果是醫學.
甚至分到老齡醫學.
又或者很多細分.
對長者老人家的老化是怎樣.
最早提出關於老化的概念.
在公元前44年.
一位西塞羅.
羅馬帝國時代的哲學家已經提出.
他當時寫了一篇論文.
講到得而的自然老化.
之後很多關於老年課題.
都有很多講述.
寫了很多不同的文章.
他們從不同的角度.
去看老年老化的問題.
今日我們要講成功老化.
其實是一個簡易的定義.
對老化的適應良好.
生理保持最佳狀態.
然後在老年好好享受生活.
這就是順利老化.
又或者積極老化.
病痛是很自然.
即使我們年輕人成長.
都有很多病痛.
老化是指不要容易得病.
又或者失去身體功能.
所以我們要有良好的心智.
身體的功能.
以致我們可以有豐盛的生活.
這就是三個內涵的條件.
所以在生理,心理和社會方面.

$^{521}$這裡也有提示.
生理上我們也知道.
進入長者的情況.
我們也會注重健康.
有規律去飲食,休息.
以致我們做這些事時.
可以預防疾病.
心理上是怎樣呢?.
其實最主要是幸福感.
我們要對老化有正向的態度.
我們要樂觀豁達.
排除生活上的壓力.
更重要的是在心理方面.
我們有宗教信仰時.
令我們的精神有超脫的感覺.
亦維持我們的興趣.
社會方面是怎樣呢?.
我們也知道.
我們會想想退休的生活.
我們要有很好的計劃.
甚至有更多社會支持的網絡.
我們參加社區的服務.
甚至我們做義工.
學習一些新的知識.
幫助我們面對這種生活.
甚至我們做義工.
更重要的是.
有很多調查發現.
在一些年長者裡.
他們的遇難為足.
他們有正面的看法.
如何在老化過程中.
能夠很坦然.
很安然地面對老化過程後.
要去到那裡.
國立台灣師範大學.
也就成功老化的定義.
在論文裡.
說出生命之首.
大家也有手掌.

$^{561}$原來透過手掌.
他告訴我們.
積極順利老化的要素.
我們的手有甚麼?.
由大拇指一直到尾指.
大拇指代表甚麼呢?.
代表我們的心智.
其實我們要有終身學習.
我們要在技能上.
在情感上.
都要不斷地學習.
食指呢.
其實是講到.
我們的宗教信仰.
宗教信仰可以幫助我們.
正面地去理解.
老人生命的意義和目標.
以致我們身心靈全人的發展.
可以達到一個坦然.
甚至成為養生之道.
中指代表我們寶健.
我們在這個年齡層裡.
很多時候都要吃補品.
人家說甚麼都好.
我們又去做.
但原來要維持我們身體健康.
這也是養生之道.
無明指.
即是甚麼都沒有.
沒有甚麼最重要呢?.
原來是沒有負面情緒的干擾.
因為我們一有負面情緒干擾.
就會影響我們.
怎樣去看自己.
怎樣看自己在老化過程裡.
會產生甚麼狀況.
或者不明朗的情況.
美指是怎樣呢?.
美指就是休閒.
其實要忘記我們年齡.

$^{601}$和身份地位.
樂於參與各種活動.
去過的生活.
大英姐妹在這裡.
分享一個長者的見證.
其實這位長者.
今年應該是88歲.
兩三年前.
她參加了教會的長者報道會決志.
她是一位弟兄的爸爸.
和這位爸爸回來的.
之後那時候身體還可以的時候.
她都參加了家樂團.
或者和我們一起參加長者旅行.
但因為她住得太遠.
在新界的村屋住.
所以每次見到這位弟兄.
我都來問候她爸爸的狀況是怎樣.
原來在這兩年裡.
她亦都參加了政府的優先配屋計劃.
她和兩個孫女.
很快就配了一個單位.
搬到天水圍附近的屋村住.
我都去探望過她.
我相信她是一位成功.
或者叫順利老化的長者.
每次去探望她.
她都很樂觀和豁達.
她很健談.
記性很清晰.
我問她平時做甚麼.
她說做運動.
到街上就去打個圈.
但最近都跌倒.
因為年紀大.
她都覺得跌倒.
然後入了幾次醫院.
但她仍然抱著樂觀的心情.
去過她每一天.
因為家裡的孫女對她非常好.

$^{641}$她可以進入一個積極老年的生活.
當我和她討論信仰時.
她說「是的,我已經相信了」.
後來她的兒子也和我說.
「我和爸爸都繼續討論過關於信仰的」.
「很想她再走一步,更加圓滿」.
「不如參加浸禮,不就好了」.
當然她都知道浸禮不是一個得救的條件.
當我去探望她時.
她很幽默地對我說.
「水禮是睡在病床上」.
「最後牧師替她灑水」.
「就是這個,到時再說」.
我說「不是」.
「其實趁著你這麼積極的時候」.
「可以對耶穌,信仰有更多的認識」.
於是她答應.
之後我們真的安排和她在家中.
上浸禮房.
弟兄姊妹成功老化.
其實包括了我們信仰的部分.
我們要積極順利進入老化的弟兄姊妹.
在我們當中有很多.
他們都積極把握這些機會.
來邀請他們所認識的長者.
特別是認識了很久的朋友.
退休的一些同事.
很想他們認識信仰.
分享他們自己信耶穌的豐盛生命是怎樣.
所以要把握機會.
把握事奉的機會.
聖經又提到老年年邁的人.
神格外顯出他的恩典和慈愛.
令他們享受豐富的晚年生活.
聖經常常提到年紀老邁日子滿足而死.
大家有沒有留意.
特別是在舊約的時候.
很常出現的這幾節聖經.
大家怎樣去詮釋呢?.
因為不是很多解釋.

$^{681}$釋經家或聖經學者去解釋.
某某年紀老邁日子滿足而死.
但當我去服侍長者的時候.
其實我很體會到這幾個字.
年紀老邁日子滿足而死.
當這群年長者.
他們是怎樣呢?.
特別是看到現在這群.
剛才和我們獻唱的長者.
其實當中有些接近90歲.
他們仍然用他們的生命來事奉上帝.
所以說年紀老邁日子滿足.
其實死我們每一個人都要面對.
但在他們的老化積極順利的時候.
我相信這是神的恩典.
神的憐憫和神的慈愛.
以撒年紀老邁日子滿足而死.
大衛一樣.
或許我們少認識的.
耶和耶大大祭司.
在猶大國後期的一位.
很重要的先知.
他也是年紀老邁日子滿足而死.
約伯我們更加清楚.
我們很了解他整個人生是怎樣.
到了後期的時候.
聖經裡面描述.
上帝給了他很多財產.
也有很多的兒女.
聖經告訴我們.
約伯年紀老邁日子滿足而死.
讓姐妹長者能夠安享晚年日子滿足.
我們要尊敬他們.
我們可以怎樣做呢?.
保羅在新約書信裡.
勸勉提摩太前書五章的一至二節.
不可嚴責老年人.
只要勸他如父親.
勸少年人如弟兄.
勸老年婦人如同母親.

$^{721}$勸少年婦女如同姐妹.
總要清清潔潔.
我們要耐心對待長者.
孝敬年老相親.
我們回想一下.
曾幾何時我們是嬰孩的時候.
他們為我們做了很多事情.
甚至他們為我們做的事情.
有很多其實我們不知道.
現在他們年紀大.
他們成為長者.
我們怎樣擁抱他們呢?.
網上有一個故事.
父親問兒子.
那些是什麼來的?.
兩父子在屋外的花園遊走.
年長的爸爸看到小動物在前面動來動去.
甚至在花叢上撲來撲去.
他就問那些是什麼來的?.
兒子就說麻雀.
父親又再問那些是什麼來的?.
麻雀呀!麻雀呀!.
爸爸又問那些是什麼來的?.
這個時候兒子又怎樣?.
都說麻雀呀!麻雀呀!.
你做什麼呀?.
你不知道這是什麼意思嗎?.
當兒子繼續問.
爸爸又再問那些是什麼來的?.
兒子就更加按捺不住了.
我都告訴你了,爸爸.
那些是麻雀來的,你知不知道?.
爸爸於是起來離開.
兒子還大聲叫.
你跑去哪裡了?.
看著下一句是什麼?你跌倒了!.
然後爸爸就沒有理他了.
走了進去,沒多久就從家裡走出來.
他手裡拿著一本日記.
他就打開其中一頁.

$^{761}$就叫兒子來讀.
他說你大聲一點讀.
今天我剛剛滿三歲的兒子.
跟我一起在公園裡坐.
有隻麻雀停在我的面前.
我的兒子問我二十一次.
那些是什麼來的?.
我就回答他二十一次.
那些是麻雀.
他一直不停地問我同樣的問題.
而我還是熱情地擁抱著他回答.
一次又一次.
沒有發怒.
持愛地回答他純真活潑的問題.
我們有沒有好好地反省.
在今天的經文提醒我們.
我們應該怎麼做呢?.
我們可以怎麼跟長者相處,溝通.
顯出我們對他的關懷呢?.
我們要擁抱長者.
我們要敬愛,敬老,愛老,護老.
有一位資深的社工.
他在老人服務服務裡面有很多經驗.
分享給我們聽三從四德.
哪三從呢?.
從不放棄.
人是神所創造,所愛的.
無論哪一個年紀.
神都沒有放棄.
所以我們面對長者的時候.
我們不要有任何前設跟他們相處.
跟他們溝通.
不要放棄.
第二就是從不抱怨.
我們要甘心服侍他們.
要按聖經的真理來教導他們.
我們不要懷怒,不要容易地發怒.
第三就是從不變心.
面對長者我們要付出愛心.
倚靠著神的恩典我們盡力去做.

$^{801}$不變形,不走樣.
四德是哪四德呢?.
第一就是要忍得.
跟長者相處,溝通.
讓他們盡情抒發.
可能他會將你成為對罵,發洩的對象.
你就要忍.
之後你給他倒杯水.
讓他一邊喝一邊平伏心情.
我們就處之泰然.
第二就是捱得.
吩咐我們做的事,辦的事.
對長者來說是重要的.
我們不要遲.
我們要為他盡快地完成.
第三就是陪得.
陪長者喝,陪長者吃,陪長者玩,陪長者坐.
但是不說話的長者.
或者他比較說話有困難的.
我們就陪他坐.
他坐多久我們就陪他坐多久.
第四就是傾得.
天南地北無所不談.
我們要聽長者說.
說他那些陳年的舊事.
其實當他分享這些舊事的時候.
當中其實有很多道理.
如果他們還是很想信耶穌的時候.
趁著這個機會跟他們談的時候.
將信仰告訴他們.
今天我們按著聖經的教導.
我們積極擁抱老年.
讓長者能夠積極地.
正面地以聖經為基礎.
在神的能力,在神的恩典.
在耶穌基督豐富的慈愛裡面.
成功地老化.
過一個豐盛得尊榮的生活.
我們作為年輕人.
我們就要放下身段.

$^{841}$以神的愛去擁抱老年.
因為神愛我們每一個.
我們一起有個禱告.
讓我們能夠經歷.
不同階段生命的成長.
主在過去的日子.
你怎樣帶領一班長者.
在他們年輕的時候.
他們怎樣能夠滿有上帝的祝福.
有能力,有氣魄.
來養活自己.
不單單是這樣.
也養活家庭.
讓今天他們所曾經養活的兒女.
都成家立業.
他們無私的奉獻.
他們那種付出.
是上帝你所體重.
他們現在進入老年的時候.
上帝你就是這樣擁抱他們.
讓我們下一代.
怎樣能夠好好的去敬奉.
去敬老,愛老,護老.
願耶穌基督你的愛.
常與我們同在.
幫助我們能夠去尊敬長者.
去擁抱老年.
誰聽我們的禱告.
奉主耶穌基督的名.
阿們.
\newpage



\section{路加福音 2:22-38}
\label{sec:CjzmC7ulsYw}
\textbf{2020年12月6日 崇拜直播｜路加福音2\_22-38}
\newline
\newline
連結: \href{https://youtube.com/watch?v=CjzmC7ulsYw}{\texttt{https://youtube.com/watch?v=CjzmC7ulsYw}} ~~~~ 語音日期: 2020-12-06
\newline
\newline
\hyperref[sec:d4oKfBQrMaI]{\small{< < < PREV SERMON < < <}}
~
\hyperref[sec:index_chronic]{\small{[返順時目]}}
~
\hyperref[sec:MIn0I7TUOfE]{\small{> > > NEXT SERMON > > >}}
\newline
\newline
路加福音 2:22-38
\newline
\begin{longtable}{cl}
\hline
\hline
章節 & 經文 (和合本修訂版)\\
\hline
2:22 & \begin{tabularx}{0.7\textwidth}{X} 按摩西律法滿了潔淨的日子，他們就帶著孩子上耶路撒冷去，要把他獻給主。 \end{tabularx} \\ \\ \relax
2:23 & \begin{tabularx}{0.7\textwidth}{X} 正如主的律法上所記：「凡頭生的男子必歸主為聖」； \end{tabularx} \\ \\ \relax
2:24 & \begin{tabularx}{0.7\textwidth}{X} 又要照主的律法上所說，用一對斑鳩，或用兩隻雛鴿獻祭。 \end{tabularx} \\ \\ \relax
2:25 & \begin{tabularx}{0.7\textwidth}{X} 那時，在耶路撒冷有一個人，名叫西面；這人又公義又虔誠，素常盼望以色列的安慰者來到，又有聖靈在他身上。 \end{tabularx} \\ \\ \relax
2:26 & \begin{tabularx}{0.7\textwidth}{X} 他得了聖靈的啟示，知道自己未死以前必看見主所立的基督。 \end{tabularx} \\ \\ \relax
2:27 & \begin{tabularx}{0.7\textwidth}{X} 他受了聖靈的感動，進入聖殿，正遇見耶穌的父母抱著孩子進來，要照律法的規矩而行。 \end{tabularx} \\ \\ \relax
2:28 & \begin{tabularx}{0.7\textwidth}{X} 西面就把他抱過來，稱頌神說： \end{tabularx} \\ \\ \relax
2:29 & \begin{tabularx}{0.7\textwidth}{X} 「主啊，如今可以照你的話，容你的僕人安然去世； \end{tabularx} \\ \\ \relax
2:30 & \begin{tabularx}{0.7\textwidth}{X} 因為我的眼睛已經看見你的救恩， \end{tabularx} \\ \\ \relax
2:31 & \begin{tabularx}{0.7\textwidth}{X} 就是你在萬民面前所預備的： \end{tabularx} \\ \\ \relax
2:32 & \begin{tabularx}{0.7\textwidth}{X} 是啟示外邦人的光，是你民以色列的榮耀。」 \end{tabularx} \\ \\ \relax
2:33 & \begin{tabularx}{0.7\textwidth}{X} 孩子的父母因論耶穌的這些話就驚訝。 \end{tabularx} \\ \\ \relax
2:34 & \begin{tabularx}{0.7\textwidth}{X} 西面給他們祝福，又對孩子的母親馬利亞說：「這孩子被立，是要叫以色列中許多人跌倒，許多人興起；又要成為毀謗的對象， \end{tabularx} \\ \\ \relax
2:35 & \begin{tabularx}{0.7\textwidth}{X} 叫許多人心裡的意念顯露出來；你自己的心也要被劍刺透。」 \end{tabularx} \\ \\ \relax
2:36 & \begin{tabularx}{0.7\textwidth}{X} 又有位女先知，名叫亞拿，是亞設支派法內力的女兒，年紀已經老邁，從童女出嫁，同丈夫住了七年， \end{tabularx} \\ \\ \relax
2:37 & \begin{tabularx}{0.7\textwidth}{X} 就寡居了，現在已經八十四歲。她不離開聖殿，禁食祈求，晝夜事奉神。 \end{tabularx} \\ \\ \relax
2:38 & \begin{tabularx}{0.7\textwidth}{X} 正當那時，她進前來感謝神，對一切盼望耶路撒冷得救贖的人講論這孩子的事。 \end{tabularx} \\ \\
[1ex]
\hline
\hline
\end{longtable}
$^{1}$由我們同心以聆聽聖道.
去敬拜我們的上帝.
今天的經訓是記載在.
《路加福音》2:22-38.
「按摩西律法滿了的潔淨的日子.
他們帶著孩子上耶路撒冷去.
要把他獻於主.
正如主的律法上所記.
凡投生的男子必稱聖貴主.
又要照主的律法上所說.
或用一對斑溝或用兩隻初甲獻祭.
在耶路撒冷有一個人名叫西面.
這人又公義又虔誠.
素上盼望以色列的安慰者來到.
又有聖靈在他身上.
他得了聖靈的啟示.
知道自己未死以前.
必看見主所立的基督.
他受了聖靈的感動.
進了聖殿.
正遇見耶穌的父母抱著孩子進來.
要照律法的規矩辦理.
西面就用手接過他來.
稱頌神說.
主啊!如今可以照你的話.
釋放僕人安然去世.
因為我的眼睛已經看見你的救恩.
就是你在萬民面前所預備的.
是照亮愛邦人的光.
又是領民以色列的榮耀.
孩子的父母因這論耶穌的話就稀奇.
西面給他們祝福.
又對孩子的母親瑪利亞說.
這孩子被勒.
是要叫以色列中許多人跌倒.
許多人興起.
又要作偽邦的話柄.
叫許多人心裡的意念顯露出來.
你自己的心也要被刀刺透.
又有女先知名叫阿娜.

$^{41}$是阿切之派法內力的女兒.
年紀已經老邁.
重作同女出嫁的時候.
同丈夫住了七年就寡居了.
現在已經八十四歲.
並不離開聖殿.
禁食祈求.
晝夜事奉神.
正當那時.
他進前來稱謝神.
將孩子的事.
對一切盼望耶路撒冷得救贖的人講說.
今天陳明荃牧師會向我們宣講上帝的話.
題目是「因他降生,我們得安慰」.
各位姐妹平安.
今天能夠和大家一起分享上帝的話語.
雖然可能從這個星期開始.
我們又要隔著螢幕.
但我們深信在上帝的話語當中.
無論在不同的地域當中.
上帝都會繼續對我們說話.
今天早上繼續思考.
耶穌來到這個世界的時候.
為我們帶來什麼不同呢?.
我們之前就像倒敘法一樣.
由耶穌的壯年一直講到現在.
今天這段經文就是耶穌出生的第八天.
到聖殿獻祭的一個情景.
其實在路加福音中.
如果留意記載著耶穌的出生.
他是用不同的角度去講.
有不同階層或不同的人.
看到耶穌出生後的反應.
來突顯耶穌的身份是誰.
今天這段我特別集中.
想和大家一起看看西面的老人家.
他的講話和祈禱.
如何將耶穌是一個救主呈現出來.
特別是在十二月.
我們這段時間開始慶祝聖誕節.

$^{81}$我們知道耶穌是誰的時候.
我們慶祝的時候.
我們帶著什麼心情去慶祝.
雖然坊間的市道不太好.
看到旺角區已經沒有過去幾年那麼多的裝飾.
或者有很豐富的佈置.
但我想今年我們就用一個比較簡約的心情.
我們帶著聖經裡面最原初的記載.
耶穌是誰.
我們一步一步的去紀念主.
盼望主的話語.
主的福音在我們生命裡面.
很豐富的呈現出來.
在我們今天看這段聖經之前.
我和大家一起祈禱.
求上主幫助我們.
我們同心祈禱.
天父上帝願你幫助我們.
藉著你的話語重新提醒我們.
我們所稱的救主耶穌是一個怎樣的救主.
幫助我們不單止知道你是誰.
幫助我們能夠按著你的心意過我們的生活.
我們恭敬將來領受你話語的時間交給你.
幫助宣講的講得清楚.
幫助每一位在家中看直播的弟兄姊妹聽得明白.
祈求聖靈與我們每一個同在.
獻上我們祈禱.
奉靠基督耶穌得勝的性命祈求.
Amen.
剛才有說過.
在路加福音裡面.
他講到有不同的角度去描述.
耶穌是誰.
在這裡就在第二章第二十二節開始.
就立刻講到按照摩西的律法.
結證的日子到了.
耶穌基督的父母約瑟和瑪利亞.
將Baby耶穌帶到耶路撒冷獻祭.
獻給上帝.
獻的祭物就是一對班戈或兩隻初甲.

$^{121}$作為祭身來獻在聖殿.
就在這個時間裡.
有一個老人家.
在耶路撒冷裡有一個長者.
叫西冕.
西冕的意思其實是敬虔人.
路加福音在筆下裡再特別的說.
這個人不單是叫敬虔.
他的人生整個人生都是一個.
又公義又虔誠.
素常盼望以色列的安慰者來到.
又有聖靈在他身上.
他得到聖靈的啟示.
知道Baby耶穌將要去到聖殿.
父母來做一些獻祭和感恩.
於是他們就在聖殿裡相遇.
當他們相遇的時候.
在聖經裡就說到這個長者很突兀.
見到Baby耶穌.
父母抱著之後.
這個長者就帶著很豐富的情感.
就走到這個孩子耶穌面前.
就按手在他身上來祈禱.
就說出一番話.
如今你可以釋放你的僕人安然去世.
因為僕人西面已經看見上帝的救恩.
就是照亮外邦人的光.
又是以色列的榮耀.
在這個情景裡我們看見.
這個長者去到聖殿當中.
受到聖靈的感動.
之前沒有相遇過.
或者在耶穌出生的時候.
天使沒有顯現給這個長者知道.
大概這個長者心裡一直在等待.
一個彌賽亞的救主.
等了大半生.
等到年紀雙笨發白.
終於等到了.
在聖靈的引導之下.

$^{161}$他就見到這個孩子耶穌.
值得注意一個很重要的字眼.
在第25節那裡.
這人又公義又虔誠.
素常盼望以色列的安慰者來到.
和合本就加了一個者字.
如果在家裡看聖經.
那個者字下面是有點點點的.
點點點是什麼意思呢.
就是說原文是沒有那個者字的.
所以如果比較貼近原文的.
就是說這個公義敬虔的長者.
素常盼望以色列的安慰來到.
大概可能是和合本的翻譯.
或者不同的中文譯本.
他就很想強調耶穌是那個安慰者.
但是一說到安慰者.
我就會說沒事沒事.
我們心裡面一定會有這種感覺.
就覺得好像一個磁場老人一樣.
有個小孩跌倒了.
沒事沒事.
這樣開一下.
我們心目中的安慰者.
是一個這樣的形象.
不過在聖經裡面所說的.
以色列的安慰.
和我們心目中的安慰者.
有一個很大的距離.
以色列的安慰主要出現在.
什麼舊約書卷裡面.
尼加書有說的.
不過特別清楚說到.
以色列的安慰.
就在以塞瓦書裡面說的.
在以塞瓦書第40章.
第一至第二節裡面.
經文裡面這樣說.
你們的神說你們要安慰.
安慰我的百姓.

$^{201}$要對耶路撒冷說安慰的話.
又向他宣告說.
他征戰的日子已滿了.
他的罪孽赦免了.
他為自己一切的罪.
從耶和華手中加倍受罰.
在這裡上帝藉著以塞瓦先知說到.
以色列整個民族的安慰到了.
你要安慰以色列的眾百姓.
為什麼要安慰.
因為這班百姓.
征戰就是受苦.
受苦的時間到了.
他的罪孽已經被赦免了.
他為自己一切的罪.
從耶和華手中已經受到加倍的刑罰.
這裡說到.
上帝說時間已經足了.
他所面對的.
整個以色列民族面對的苦難已經足夠了.
所以你要接受安慰.
因為上帝的拯救來到.
簡單來說是什麼意思呢.
在以塞瓦書裡面.
其實你可以分為三部曲.
上部份就說到上帝的審判.
上帝的審判就基於以色列離開上帝.
離棄了神的旨意.
多次又一次的離棄耶和華.
去敬拜偶像.
不聽上帝所說的話.
以致他們就要面對上帝的審判.
當上帝的審判來到的時候.
就是足夠了.
上帝的目的不是要糟蹋那些以色列人.
他不是要將他們整個民族置於死地.
而是藉著這些審判.
或者藉著這些苦難.
就好像教訓一樣.
教導了整個民族之後.

$^{241}$然後就進行上帝的拯救.
而在拯救和在審判中間.
就是一個安慰的說話.
三部曲是審判,安慰,然後拯救.
整個以色列民族都是在一個這樣的狀態之下.
一直在等待上帝為他們預備的拯救.
所以在西面裡面說.
我看到你的救恩.
為什麼看到他的救恩呢.
就是知道上帝懲罰以色列人.
或者以色列受害的日子已經足夠了.
如果你這樣說.
也不明白我再用一個例子.
這是我自己小時候的例子.
話說我爸爸在我們家裝修.
然後家裡就美輪美奐.
做了一套組合櫃.
然後組合櫃中間位置.
就買了一套最漂亮的Hi-Fi唱機.
我小時候是唱黑膠唱片的.
如果大家知道.
黑膠唱片上面有一支唱針.
唱盤是順時針轉動.
我爸爸就千叮萬囑.
他說家裡所有東西都很危險.
但這個唱機真的不可以搞.
很貴的.
而且唱片一動.
唱片會花.
唱針會斷.
所以就嚴禁去搞.
有一天.
趁著我爸爸星期日早上還在睡覺的時候.
我一早起床.
他越叫我不要搞的東西.
就越好奇.
我想每個小朋友都是這樣.
於是我就搬一張椅子.
就爬在那裡.
看一看為什麼這個東西這麼有趣.

$^{281}$會轉動呢.
於是就把他扭大了.
開了那部機.
然後在那裡第一件事做的.
就是搓碟.
搓碟搓碟之後.
為什麼要順時針唱才可以呢.
逆時針行不行呢.
於是就把碟反方向轉動.
逆時針就斷了.
一聲噠噠的.
好像爆炸聲一樣.
爆炸聲之後.
我爸爸在房間裡聽到.
以為真的石油氣爆炸.
就衝了出來.
一衝出來去客廳的時候.
看到自己的兒子.
舉起手在唱機上.
然後就看到唱碟.
整部機就好像面對著一個很大的.
就是整部機都要報銷的狀態.
於是我爸爸就馬上喝下來.
拿回來.
就罵了.
然後被爸爸喝的時候.
當然馬上哭了.
每個小朋友都很本能的.
我想起自己為什麼要哭呢.
本能來的.
不過心裡面想.
我夠慘情.
我夠仇富.
我爸爸可能會輕手一點.
但是我爸爸一定不是這樣做.
看到這樣.
他剛剛買了沒多久的Hi-Fi.
這樣搞法.
第一件事.
關了機之後.

$^{321}$就拿一把硬石出來.
就是那些木的鐵尺這麼長.
伸一隻手出來.
問我打多少下.
我說打一下就好了.
一下就弄壞了我的唱機.
於是就忘記了.
打了十下或二十下.
每次都很大力的.
啪啪聲的打下去.
就哭得更厲害.
一邊哭一邊痛.
家人不敢求情.
因為實在是.
這一次實在太離譜了.
痛完之後.
縮在一個角落.
在沙發邊繼續喝.
一直哭.
然後哭了一段時間之後.
怎樣去修科呢.
誰做主動呢.
最後還是爸爸走過來.
主動說痛不痛.
知道痛了嗎.
伸一隻手來看看.
然後塗上一點點藥膏.
然後就安慰.
抱著自己.
這樣教兒子的.
做爸爸的.
錯了的事.
錯了就錯了.
懲罰就懲罰.
懲罰完之後哭完了.
才安慰他.
哭他.
如果我們用這個角度來看.
上帝對以色列人整個民族都是這樣.
尼西.

$^{361}$上帝.
拜偶像.
偏行幾路.
整個民族都是這樣.
離開上帝的時候.
上帝就將刑罰淋到他們當中.
懲罰完了.
足夠了.
痛苦完了.
日子夠的時候.
上帝做主動的.
來跟他們說.
是不是痛了.
現在就是療傷的時間.
你的受苦的日子夠了.
你的罪孽已經赦免了.
你.
加倍的刑罰.
就是長子就加倍刑罰.
受夠了.
以致到.
如今就可以接受上帝的拯救.
所以在這裡.
在說的以色列的安慰背後.
就是在說上帝的.
審判.
上帝的赦免.
和上帝的拯救.
西緬這個老人家.
世世代代.
由以塞瓦書到他在他那一刻.
他一直.
一直看以塞瓦書.
或者一直看聖經.
他知道.
上帝的刑罰是.
是真實的.
他真的要這樣.
來去.
懲治.

$^{401}$改變.
整個活躍的群體.
不過在一個這樣的狀況下.
上帝卻是將一個受苦的群體.
重新挽回過來.
安慰他們.
以致到他們不會.
再繼續這樣來行差踏錯.
所以在這裡所說的那個字眼.
就不是純粹說一個安慰者.
就是說耶穌是一個.
沒事的.
舒服了.
而是說到上帝.
懲罰以色列人的日子夠了.
他要藉著這個孩子.
將來所做的工作.
將這個救恩.
將上帝的拯救.
帶到整個以色列的民族當中.
讓上帝的救恩.
呈現在這個民族裡面.
當你讀這段聖經.
或者你會看路加福音.
你會說.
這是上帝閉門教兒子.
教以色列.
以色列的救恩.
關我們.
我們這些中國人.
這些外邦人什麼事呢.
相對於以色列人.
其實這件事可以不關我們事.
不過路加福音裡面.
他的不足.
他不是純粹停在這裡.
說到那個救恩.
以色列的盼望.
見到之後就停止了.
經文裡面再繼續的去展述.

$^{441}$特別是第29至32節.
這是其中西面一句很重要的土文.
很多的詩歌就叫做Nundimittis.
拉丁文.
就是西面的祈禱.
Nundimittis就是在說.
讓僕人去世.
就是釋放僕人安然去世.
整句話是怎樣說呢.
主阿如今可以照你的話.
釋放僕人安然去世.
因為我的眼睛已經看見你的救恩.
就是你在萬民面前所預備的.
是照亮外邦人的光.
又是你民以色列的榮耀.
在第29至32節裡面.
他講到西面看到上帝的救恩.
臨到在以色列的當中.
不過這個救恩.
不是只給以色列這個民族.
他看的聖經.
或者在路加福音引用西面所說的.
那個祈禱文.
後面就說你在萬民面前預備的.
是照亮外邦人的光.
又是你民以色列的榮耀.
這一整句話.
是來自以塞瓦書後的篇幅.
第49章第6節.
以塞瓦書.
那句經文是這樣說的.
現在他說.
你作我的僕人.
使雅各眾之派復興.
使以色列中得保存的歸回尚為小事.
我還要使你作外邦人的光.
叫你施行我的救恩.
直到地極.
以塞瓦書說的是什麼訊息呢.
去到第49章.

$^{481}$他講到以色列的安慰.
上帝的刑罰.
到了日子到了.
一路一路發展.
去到第49章的時候.
他講到以色列民族被拯救.
都是小事.
不是最重要的重點.
最重要的重點是講到.
以色列整個民族.
或者上帝所選派的耶穌基督.
要成為上帝的僕人.
不單止要將復興放在以色列當中.
而且以色列的復興.
要成為這個世界所有人的光.
這個世界外邦人所光照的上帝的光.
這個救恩不單止要淋在以色列這個民族裡.
更加要淋在整個以色列的.
以外的外邦人的群體裡.
所以這個外邦人的光.
這個救恩淋到全地.
其實是講到西面讀通了以塞瓦書.
他知道耶穌基督來到這個世界.
成為以色列的盼望.
亦都是成為全世界的盼望.
他要將這個救恩.
不單止淋到一個國家裡.
藉著這個以色列的這個小群體.
一個小國家.
重新將上帝的美好.
淋到世界不同各族.
各邦不同民族的人身上.
西面他的滿足.
不單止是純粹滿足在自己的民族得救.
他的滿足卻是知道.
上帝藉著這個英鞋.
將救恩去到更加遠.
帶到更加廣闊的一個境地.
弟兄姊妹,去到讀到這裡.
我們今天思想一下.

$^{521}$今天我們心靈裡面的滿足.
我們的信仰為我們帶來的是一種怎樣的滿足呢?.
很多時候我們會覺得.
我們自己經營著自己的生活多美好.
我們搞得定自己.
或者我們覺得我們做了某一些的行為也好.
或者我們做了一些宗教上的禮儀.
我們覺得已經做足了,很好了.
我們自給自足.
經營我們的小天地.
我覺得我們這個是很棒,很滿足.
不過在聖經裡面.
特別當你仔細去讀路加福音.
其實聖經在催促我們.
我們的世界不能去到這麼窄.
我們不是純粹在信仰裡面滿足了我們自己.
經營我們自己的小天地.
信仰某程度上在催迫我們的世界要廣闊一些.
讓我們的心胸看到這個普世上福音的需要.
我們的滿足不是在說我們滿足在自己關上門的宗教生活.
或者我們上YouTube看黃浸的製作就很滿足了.
我們更加要把這個福音廣傳開去.
我們的心胸越廣闊.
我們越貼近上帝的心懷的時候.
我們才能夠明白基督耶穌的心腸.
在早一段日子我看到一個作家叫Peter Hsu.
他寫了一句說話很有意思.
他說人的心是一對翅膀.
你的心有多大,你的世界就有多大.
但是如果你不打破心的禁錮.
即使我給你整個天空.
你都找不到自由的感覺.
今天我們會不會我們信了主之後.
我們有意無意之間令到我們的心胸好像更加被禁錮一樣.
我們只是滿足一班弟兄姊妹那種很溫馨的感覺.
我們只是純粹滿足了在我們自己生活的形態繼續自給自足.
卻看不到這個世界上福音上的需要.
曾經有些弟兄姊妹跟我說.
牧師其實我們的媒體工作不用做得這麼廣闊.
其實不用做得這麼好.

$^{561}$其實我們現有會有幾百人去看.
那個製作足夠了.
你做了這麼多東西會不會真的容易浪費了資源.
或者有時候會不會搞到捱更晚.
當然對於弟兄姊妹.
他是善意的,怕我們同工過路.
不過我認為我們今天的視野不是純粹直著媒體服務我們眼前.
我們牧養開的弟兄姊妹.
我們所看的胸襟,我們所看的世界.
我們看到福音上的需要可能超越了我們旺角浸信會這個群體.
我們希望我們所做的事情可以祝福更加多未知未信的人.
我們所說的訊息已經不是單單是說一小撮人.
來滿足他們心靈上的需要.
當我們看到這個世界還有很多人.
特別是面對著這個疫情.
這個世界甚至乎在眼前的未來都還是很迷茫的時候.
我們相信福音要進到在這些.
未在旺角浸信會這個群體的這一班人身上.
我們要將我們的翅膀,我們的眼光放遠一點.
以至讓更加多人得著基督耶穌的救恩.
我們的滿足就正如西面一樣.
不是純粹滿足在以色列受到安慰.
我們的滿足要看到外邦人,看到一些未信的人.
看到一些生命裡仍然是仍然是逆逆,未知方向的人.
他們得著福音,看到他們信主.
我們內心才有真正而來的滿足感.
所以弟兄姊妹,在這裡我特別呼籲.
當你看完這段經文,你看到西面的這個祈禱.
我希望都成為我們的祈禱.
讓我們不單止滿足在我們的小群體裡的那種很溫馨的感覺.
嘗試將你自己生命裡的群體拉闊一開.
將你的視野擴闊一點,看到這個世界上不同的需要.
又或者我們看小組的時候,小組我們很多時候的關係很溫馨.
區寒民暖,我們經常都會說我們的生活分享,這是好的.
不過如果我們幾十年都是說這幾個人,都是繼續這樣圍爐取暖.
大概我們就變得很內在,我們的心胸就會變得很窄.
我們的心反而成為籌固了我們享受不了自由的天空.
但當我們的生命將我們的眼光擴闊了的時候.
我們將我們的愛跟別人分享的時候.
我們小組不是純粹互相講安慰的話.

$^{601}$而是在講我們那個群體可以成為更加多人帶來從神而來的安慰.
這個是第二點,西面那個西面縱裡面給我們看到的那種滿足感.
去到第二章第34至35節.
西面繼續對瑪利亞說,這孩子彼納是要叫以色列中許多人跌倒.
許多人興起,又要作偽謗的華病,叫許多人心裡的意念顯露出來.
你自己的心也要被刀刺頭.
在這裡西面就好像一個發出先知的預言一樣.
說這個孩子將來會是怎樣的,幾天大的孩子.
大概他不是按他的行為來說這番話.
西面是按著聖靈的感動,將上帝的旨意.
好像告訴瑪利亞的媽媽知道將來這個孩子會成為怎樣的人.
這個孩子將來是要叫以色列人很多被跌倒.
講到整個以色列盼望竟然令到以色列想法顛覆了他.
甚至好像半腳石那樣把這些以色列人卡住.
他們一直帶著錯誤的誤解.
對於彌賽亞以為是有一個君王,一個政治君王來統管他們.
重新來復國.
耶穌的出現,耶穌的將來不是一個這樣的形態.
他是以色列的王,也是世界的王.
不過他不是一個軍事霸主.
他不是以色列心目中揭竿起義的革命家.
上帝是要施行拯救,在十字架上將福音呈現出來.
所以在這裡,耶穌的將來是要令到很多以色列人大跌眼鏡.
是會令到他們跌倒的.
因為他所講的東西和這群以色列人心目中所想的彌賽亞完全不相符.
所以這些人就會被跌倒.
但同時間他也會興起一些新的人.
叫很多人被興起.
因為這些是真正明白上帝.
真正信靠耶穌成為他們生命的主.
這些人不是勢很強.
只不過是這群人知道上帝的心意.
他們行在上帝救贖計劃當中.
所以有人跌倒,有人興起.
不是世代純粹的更替.
而是上帝要將拯救計劃淋到很多偏頗看法的人身上.
上帝要將人帶到他們的根前.
過程可能會顛覆一些人的想法.
甚至令到人像半腳石般跌倒.
但上帝的目的就是要用耶穌基督的應許.

$^{641}$一種受苦的形態.
來將世界的罪歸到他自己身上.
在經文裡面繼續的去展述.
他說耶穌要成為毀謗的說病.
說病這個字是記號的意思.
一見到耶穌,人們就會毀謗他.
說一些冤枉他,辱罵他的說話.
其實耶穌不是做那些事.
不過冤枉他或是將不好的話套在耶穌身上.
他就更加反映了那些人內心的那種人性.
人的罪在耶穌出現的時候.
就呈現了這些一切.
耶穌的降生就像一面鏡一樣.
將人心裡面的邪惡.
將人心裡面的那些人性顯露無遺.
耶穌沒有做任何錯的事.
不過人心裡的惡在耶穌的事件當中.
他就顯然而見地將人心裡的意念顯露出來.
當耶穌做一個角色的時候.
做媽媽的瑪利亞一定是很痛苦的.
西面說你自己的心也要被刀刺頭.
最終是十字架上的受苦.
不過未上十字架之前.
見到所有人都毀謗耶穌.
有些人要敵視他.
做媽媽的瑪利亞絕對是一個很痛苦的經歷.
不過在聖經裡面說到.
耶穌將來所行的路就是這樣的路.
藉著真理將每個人生命深處的污穢呈現出來.
耶穌就是一面鏡.
將人生命裡面的隱情揭露出來.
說這件事為何這麼重要呢.
跟安慰是有關的.
剛才記不記得我們所說的.
我們的安慰就是要讓以色列人知道自己知罪.
他們的罪被赦免.
然後上帝的安慰才能臨到這班人身上.
所以如果要說上帝的安慰來到的時候.
這個安慰是真實的.
首先要見基於這班人知道自己走錯了路.

$^{681}$要知道罪 知道錯.
如果耶穌的來到是繼續強化我們自己在罪裡面的生活.
大概是一個虛假的安慰.
大概不是引導我們進入上帝真理的安慰.
只不過是一個藏著我們.
是一個糖衣的毒藥.
令我們感覺良好.
不過上帝的安慰從來不是純粹令我們感覺良好.
上帝的安慰很多時候首先要藉著基督耶穌的真理.
呈現我們生命的真相.
我們看不到我們自己是一個怎樣的人.
我們不停去祈求上帝的安慰.
我們是求不到的.
因為我們只是抽絲剝繭地去選擇我們喜歡聽的經文.
繼續去我行我素.
卻不能真正面對著我們生命面對著一些最重要.
我們忽略了的盲點或難題.
但是當基督耶穌的真理來到.
打開了我們生命看到自己.
我們自己的問題的時候.
我們坦然在上帝面前去承認自己的過錯.
將我們自己的罪呈現在上帝面前.
上帝的安慰就會隨著他的恩典來臨到我們每一個身上.
耶穌來到這個世界.
他的角色就像一面鏡子.
這個真理的鏡子就照顯了我們生命裡面的幽暗.
如果我們沒有這個覺醒.
大概那個安慰都是虛假的.
大概我們都是活在自我感覺良好的信仰觀念裡面.
今天這段經文不是很複雜的.
從西面的祈禱裡面.
不過我們知道從這段經文裡面我們看到.
上帝去信守承諾.
在審判和拯救之間.
他要對以色列人作出安慰的話.
但這個安慰的話不是純粹滿足了我們.
是撫平了我們不安和不舒服的感覺.
上帝的話卻是要照顯我們生命裡面的忍情.
上帝所做的工作不單止是針對以色列人.
以色列人成為一個先鋒部隊.

$^{721}$他們生命的改變.
上帝看以色列人就像一個長子一樣.
他要藉著這班人首先歸回.
上帝臨到他們身上.
接著藉著這些經歷成為外邦人.
同樣得到上帝拯救的一個救恩.
整個世界都要得到上帝的安慰.
整個世界都要藉著基督耶穌的來臨.
照顯我們生命的不足.
以致我們才真真正正需要神的安慰.
最後有個例子.
早幾天有機會見到一些小孩子.
近來很喜歡見到小孩子.
和父母怎樣教導子女.
看到一個孩子是這樣的.
見到大人很害怕.
他就經營自己的小天地.
他喜歡躲在椅底玩.
其實沒所謂的.
小孩子總會有自己覺得舒服的地方.
但在玩的過程中.
我故意作弄他.
我就坐在他的椅子上.
坐在他的椅子旁邊.
因為他說椅子是大貨櫃車.
就這樣駛來駛去.
我看到他撞來撞去.
這樣不行.
我就坐在他開的車上.
坐在椅子上.
然後就發脾氣.
小孩子就發狂.
又哭又怎樣.
我最喜歡看到小孩子哭.
我不是變態.
而是看他有什麼反應.
看父母有什麼反應.
當我看到小孩子哭的時候.
我就冷眼旁觀.
我第一時間看到的就是.

$^{761}$他的父母衝進去抱著孩子.
然後就哄他.
哄他的時候就冷靜下來.
平復他的情緒.
抱著他很疼他.
然後就哄他.
叫他跟我說對不起.
其實他沒有對不起我.
不過我看到的就是.
我看到父母有一點負面.
第一件事就是怕他哭.
怕他不開心.
所以第一件事就是安撫他.
令他舒服一點.
然後平復了他之後.
才說有什麼錯.
不過在過程中.
小孩子其實只著重父母.
是否哄他開心.
後來有沒有說自己做錯什麼.
或者對不對.
他不再理會了.
他繼續.
一會兒平復了之後.
就繼續走去玩.
在這件事裡帶來我的一些看法.
如果我們將那個關係逆轉.
我們其實很多時候.
今天的信徒都是這樣.
我們做錯了.
或者我們有些情緒.
第一件事我希望上帝抱著我.
你平復我.
令我感覺良好.
但我們沒有問自己.
究竟我們這些痛苦.
會不會是我們生命的盲點造成.
我們不是很願意接受.
上帝給我們的懲罰.
我們不覺得我們自己.

$^{801}$是一個生命裡有錯誤的人.
我們只是求上帝平復我們.
讓我們舒服一點.
如果我們繼續這樣.
去過我們的信仰.
大概我和這個小孩子沒什麼分別.
我們都是將上帝視為一個慈父.
哄我快樂.
平息我自己心裡的不安.
卻不能夠指引我們正路的一位神.
不過耶穌基督.
不是做這些事.
或者上帝的計劃.
是要讓我們正視.
我們自己今天的問題.
當我們看到.
我們生命的隱情的時候.
上帝的安慰才臨到.
如果逆轉的話.
大概我們將來會成為一個.
生骨大頭菜.
大概我們都是繼續成為一個.
活逆的人.
求主繼續幫助我們.
上帝的安慰.
臨在首先認罪的人身上.
臨在看到自己生命真正真相的人身上.
願主繼續給我們這對慈福.
我們用人.
\newpage



\section{ }
\label{sec:MIn0I7TUOfE}
\textbf{2020年12月11日擁抱幸福分享晚會｜將幸福帶回家 ｜羅乃萱師母}
\newline
\newline
連結: \href{https://youtube.com/watch?v=MIn0I7TUOfE}{\texttt{https://youtube.com/watch?v=MIn0I7TUOfE}} ~~~~ 語音日期: 2020-12-12
\newline
\newline
\hyperref[sec:CjzmC7ulsYw]{\small{< < < PREV SERMON < < <}}
~
\hyperref[sec:index_chronic]{\small{[返順時目]}}
~
\hyperref[sec:m_Hc7J2x7vY]{\small{> > > NEXT SERMON > > >}}
\newline
\newline
$^{1}$大家平安.
在這個疫情的日子.
其實很渴望親眼跟大家面對面聊天.
不過現在都可以的.
我知道我的Facebook專頁.
都應該有人在看.
如果你在看的話.
發訊息給我聽好嗎?.
讓我知道你在看.
是的.
今天給我的題目叫.
「將幸福帶回家」.
這個都是我童年的想像.
不過對我來說.
其實我是活在一個很多陰影的家庭裡面.
但是怎樣將幸福帶回家呢?.
就請聽我慢慢道來.
先說說這個時代好嗎?.
有人叫做「疫時代」.
是的.
疫是疫境的疫.
或者疫情的疫.
我怎樣形容這個時代呢?.
給大家看看.
我用四個「不」去形容.
下一張就是「不健康」.
甚麼叫不健康呢?.
疫情橫行.
健康難保.
我就買了很多顏色口罩.
不過前幾天聽到說.
有顏色的口罩不安全.
現在都想著對著很漂亮的口罩.
不知怎樣辦.
不知你是否像我這樣.
貪漂亮的.
第二.
沒有了安全感.
有人說失去常規就失去保障.
這也是真的.

$^{41}$我們每天出來有時間表.
早上做甚麼?.
送子女上學.
現在呢?.
沒得送.
然後去哪?.
去到餐廳.
原來呢?.
沒東西賣.
或者停業.
很多事情你覺得.
失去了一個日常.
當然有人說新常態.
去適應.
新常態即是沒有常態.
第三.
不知道.
明天如何還不知道.
我現在都有計劃.
一月做甚麼.
二月不敢想.
三月更加不敢想.
不確定.
沒有肯定.
只有不清楚.
和兩手準備.
是嗎?.
實體Zoom.
我們現在Zoom.
和大家.
應該Facebook Live.
還是YouTube.
如何將幸福帶回家.
甚麼是幸福快樂的家庭.
有幾點想和大家討論.
因為我自己做親子教育.
我覺得.
有些事情很想和大家討論.
第一.
健康和樂的家庭氣氛很重要.

$^{81}$一個快樂家庭.
這裡.
有八個元素.
很想和大家分享.
甚麼是健康和樂呢?.
溫暖.
回到家.
不是開著暖氣覺得溫暖.
好像剛才牧師所說.
不是有間屋覺得溫暖.
你要感受到你是受歡迎的.
有些小朋友.
因為我教寫作.
小朋友告訴我.
回到家很害怕.
因為一開門.
媽媽就跟他說.
做了功課沒有?.
溫習了沒有?.
他已經看到媽媽的苦紋在臉上.
覺得冷冰冰.
溫暖.
表示你擁抱你的孩子.
他們想回家.
親密.
大家能夠彼此相愛.
親密不是甚麼都要說.
親密是代表彼此之間.
我經常覺得身體語言很重要.
我爸爸在我們年幼的時候.
訂了一個家庭禮儀.
是甚麼呢?.
每天晚上睡覺前.
一定要跟父母說.
早上好.
親他們兩個的臉.
很親切.
直到孫兒出生.
夫妻協調一致.
支持了解.

$^{121}$開放溝通.
支持了解大家都知道.
溝通就好像打乒乓球.
你打過乒乓球就知道.
我們很多時候會說「削」.
或者叫「下旋」.
那句「削」的說話就是.
你做了功課就削他.
甚麼是開放溝通?.
今天開心嗎?.
有甚麼有趣的事情?.
尊重合作.
尊重孩子.
不要用他的缺點來做他的花名.
所以我不贊成人們叫他肥仔或肥妹.
我小時候.
我現在都黑黑的.
我發覺我的臉都頗黑的.
我的花名叫黑妹.
周圍的人都叫我黑妹.
我喜不喜歡?.
當然不喜歡.
幽默風趣.
大家都知道.
一個家是充滿活力的.
這八點.
你看看.
另外一張圖畫.
這張圖畫是我參加繪畫比賽.
拿到第一名的作品.
請你上去那張圖畫.
那張圖畫是一家三口.
有一片白雲.
白雲之下有一棵巨大而不會凋謝的花.
頒獎禮那天.
看到這一家三口.
我問小妹妹.
因為她畫的.
小學三,四年級.
「甚麼叫一盤不會凋謝的花?」.

$^{161}$「是麵粉花嗎?是膠花嗎?」.
她說「不是」.
「那是甚麼來的?」.
她說「就像父母對我愛」.
「日日新鮮 不會凋謝」.
他們就是要這些.
好 再下去.
Appreciation.
彼此欣賞,鼓勵.
右邊的圖是我最近聽一個訪談.
他說中大教育研究團隊.
說影響學生升讀專上教育.
最影響他的人物是甚麼?.
等等等等.
媽媽.
27.8\%.
爸爸也不差.
12.2\%.
當然是朋友.
所以.
媽媽不要小看自己.
我們可以做的很多.
譬如.
不要嘮叨他.
不要經常告訴他.
「你的數學不行」.
「你多做點數學」.
你可以鼓勵他.
投入的過程.
思考的技巧.
找到潛能.
我最近找了一個.
十句最討厭父母說的話.
跟大家溫習一下.
第一名.
「玩笑些電腦手機」.
第二「你學人家那麼勤力」.
第三「今天做完功課了嗎?」.
第四「再不讀好書」.
「將來沒有前途了」.

$^{201}$「你的東西那麼亂」.
「拜託你收拾一下」.
再看下去.
六到十.
我很不喜歡第六句.
「經常躺在家裡不工作」.
「像個廢人」.
父母可能覺得.
我隨口說說.
但孩子記得一輩子.
我教過一個小朋友.
他父母叫他做廢柴.
我叫他做組長.
他不肯做.
他說「我是個廢柴來的,羅老師」.
我說「你不是」.
「你是一支火柴來的」.
「少少事不懂做」.
「沒有用」.
「孩子一輩子都記得」.
所以我覺得.
我們要小心說.
第十句都不是太好.
「你再這樣」.
「我以後不理你」.
我們現在叫情緒勒索.
可能你會問.
「那說些什麼好呢?」.
根據我的接觸,引導,實踐.
我覺得這十句話.
跟大家分享一下.
我們很想他溫書.
你想溫書後才填.
打機?.
踩單車?.
還是踩完單車才溫書.
讓他有選擇.
但其實沒有選擇.
第二.
今天在學校有什麼有趣發現呢?.

$^{241}$第三.
是我爸爸經常問我的.
在我童年時.
一打電話給他.
「爸爸,我放學了」.
他說「有什麼好消息?」.
我說「想想吧,一定有的」.
變成.
我每天都要想正面積極的事.
第四.
看到新聞.
你問他.
如果你是故事裡的主角.
你會怎樣呢?.
或者你告訴他.
有難題我們一起解決.
我們來一次事後檢討.
考試後我經常跟女兒說.
她考完試.
有時成績不太好.
我們說「賽後檢討」.
等於打完乒乓球一樣.
看看有什麼做得好.
什麼做得不好.
這個很有趣的事.
是我女兒小時候經常說的.
最近.
剛才牧師說我「P牌婆婆」.
我聽到女兒跟孫兒說這句話.
我很開心.
她跟孫兒說.
「兒子,你來看看,很有趣」.
我看著心裡很甜.
因為我覺得.
這就是傳承.
你說的話.
她聽進腦裡.
她繼續做下去.
事情怎樣發生.
可能你孩子很緊張.

$^{281}$你就問他.
發生什麼事?.
誰在場?何時?.
試試等他複述出來.
等他自己說的時候.
他也明白.
第九.
你做得很好.
做得好是什麼?.
懂得體諒別人.
懂得幫助別人.
很多人說我們是「I世代」.
我未必同意.
當我們繼續.
用這些思維.
去影響我們的孩子的時候.
其實就是「V世代」.
看得出你很努力.
很欣賞你.
也要說.
我另外一個工作.
我說過了.
很喜歡教小朋友寫作.
這裡.
有兩個小朋友.
其中一個就說.
我說你怎樣鼓勵自己.
他說我要有.
屬於自己的生命花園.
要去栽種.
另外一個就說.
當我生氣的時候.
我另外一部分.
當然一部分變怪獸.
我不說.
他說我另外一部分.
就會變成天使.
「詩歌止息了憤怒」.
「榮光闖開了善良的眼」.
「聖靈安慰了氣炸的我」.

$^{321}$是不是寫得很好?.
然後他說.
要掃那些.
什麼經文冷靜了我.
第一句就是「不可含怒到日落」.
第二句就是「愛是不輕而發怒」.
「一個個氣沖沖的靈魂」.
「傳避這善意的話感化」.
不要覺得小朋友不懂事.
很多時候.
他們一句兩句.
是我們最大的鼓勵.
第三.
「Play」.
「越玩越親密」.
這個時間.
很多時候.
我們在家裡叫「困獸鬥」.
其實不是的.
是新常態.
有些弟兄姊妹或朋友跟我說.
「嘩!可以跟子女天天玩」.
「玩桌上遊戲」.
「玩不同的東西」.
最近有個爸爸.
我前天跟他聊天.
他說.
他跟女兒去看日落.
中一的女兒.
開車跟他坐在那裡.
我說很浪漫.
他說「是的」.
當他跟女兒看日落的時候.
女兒就會將自己心裡面.
很多的心事.
向爸爸傾訴.
我說「你需要問他嗎?」.
他說「不用」.
跟他很輕鬆的環境.
不是對著家中的鐘.

$^{361}$就已經可以了.
大家參考一下.
另外再下去.
「Place and Presence」.
在疫情當中.
需要將幸福帶回家.
我們要為自己.
營造一個空間.
我們有很多壓力.
是的.
失業.
家中有人生病.
或者自己遭遇一些.
很難以告訴別人的苦衷.
不要說甚麼了.
我們需要一個空間.
我在疫情期間.
讀了一個課程.
叫做「Darkness of Night」.
靈魂的黑夜.
其中有一個章節說.
我們要為自己營造一個.
心靈空間.
其實就是一個盤.
你將古代的圖畫.
公仔放進去.
每天看著它.
我其實每天都是.
心裡面就祈禱.
「神啊,你對我說話」.
「你告訴我有甚麼人」.
「我需要關心」.
「有甚麼人我需要」.
「繼續跟他連結」.
所以我很多電話.
就是早上六時起床.
零收.
打的.
寫的WhatsApp.
真的覺得很奇妙.

$^{401}$另外.
這個空間是包括.
我們都要找朋友的.
當然現在就說.
不要出去喝那麼多咖啡.
可以.
現在我知道有Zoom Lunch.
Zoom Coffee.
Zoom Dinner.
你有沒有聽過.
大家都拿著盤東西在那裡吃.
你開個Zoom.
一樣可以聊天.
沒有東西可以難阻到.
當然.
Be present(存在).
活在當下很重要.
Be present(存在)很簡單.
好像我現在這樣.
沒有手機在手.
你就Present(存在).
你有手機在手.
你這個人就不知飛到哪裡去.
人在心在的時候.
你的孩子.
你的家人.
是感覺到的.
家其實是甚麼呢?.
就是心之所在.
Home is where your heart is.
說起來好像很好聽.
又好像很容易.
但是.
家又是令我們心碎.
心傷.
和心心受傷害的地方.
為甚麼呢?.
最近經常聽這幾個字.
原生家庭.
我們自己的家.

$^{441}$帶給我們愛.
但同樣帶給傷害.
今天.
我的Facebook就寫傷痛.
都沒有想到這麼多人回應.
告訴我到現在都很痛.
被最親的人傷害.
很難度過.
原諒不了他.
怎樣可以修復呢?.
我們可能真的會這樣想.
第五.
第五,我是You and Yesu.
原生家庭.
或者命運.
剛才說的失業.
突然而來的.
遭遇帶來很多困惑.
令我們覺得第一.
不能選擇.
沒得選擇.
父母沒得選擇.
老闆要解僱我.
沒得選擇.
沒得抗拒.
沒得揮.
我們都搞不定這個冠狀病毒.
希望疫苗可以.
都不知道.
第三,不能逆轉.
是沒有彎轉.
沒有出路.
那怎麼辦呢?.
很坦白說.
各位親愛的朋友.
如果真的沒彎轉.
沒得選擇.
我今天羅乃軒不會站在這裡.
在過去.
我信了耶穌差不多四十多年.

$^{481}$我真的經歷了.
遇見神.
覺得人生是有彎轉.
有得選擇.
有選擇.
有得選擇是什麼呢?.
你可以選擇生命.
選擇光明.
《申命記》三十章.
十九到二十節.
它說你要選擇生命.
好讓你和你的後裔.
都可以活著.
愛耶和華你的上帝.
聽從他的話.
不單止這個選擇.
影響我們下一代.
詩篇三十六篇.
第九節也很好.
因為你在那裡.
有生命的源頭.
在你的光中.
我們必得見光.
我們有得選擇.
因為我們要選擇光明.
我們要離開黑暗.
我們要選擇生命.
離開咒詛.
我記得.
我信耶穌的時候.
真的很單純.
牧師對我說.
祈求就得著.
尋找就尋見.
人何時要選擇光明呢?.
是你最黑暗最軟弱的時間.
幾年前.
我經歷一場大病.
知道要動一個很大的手術.
當時.

$^{521}$害怕.
忐忑.
那天是星期二.
我要做一個決定.
我要選兩個醫生.
A醫生還是B醫生.
各有各好的.
我星期一的時候心裡很忐忑.
於是我很記得那天早上.
我看聖經看到機電.
濕陽毛和乾陽毛.
求印證.
我就和上帝祈禱.
耶穌啊.
我很害怕.
只要你給我一個印證.
讓我知道你與我同在.
我就不害怕.
聽起來可能覺得很玄妙.
不過在上帝裡面真的有這樣的東西.
當時.
心裡有個聲音跟我說.
你今天十二點半去這間餐廳.
你會見到你的.
以前有一個朋友.
Dr. Chan.
他會告訴你看哪個醫生.
很離奇.
我連我先生何牧師都不敢說.
因為我覺得太離奇.
但我有沒有做?.
當然有了.
我跟老公說.
去那間餐廳吃東西.
我去.
去到.
這個Dr. Chan.
坐在那裡一個人吃午餐.
我見到他.
打完招呼.

$^{561}$我說你經常來這裡吃嗎?.
他說不是.
我四年沒有來過.
今天剛剛路過.
我來這裡吃一個小時午餐.
我就在那裡遇到他.
然後我就哭了.
我說Dr. Chan.
你真是我的恩人.
我遇見你真是很好.
他叫我來吃飯而已.
你用不著這麼激動.
我說因為.
對我來說.
我問耶穌.
你是否與我同在.
上帝回應我.
然後祂告訴我.
你看哪個醫生.
他年輕一點.
結果看完醫生.
做了大手術.
休息了四個月.
如今.
你看我都健健康康的.
因為我相信.
在上帝的光裡.
我們一定.
一定見到光.
第二有得揮.
有得揮.
聖經有很多經文告訴我.
沒有人能夠讓我們.
與上帝的愛隔絕.
叫患難嗎?.
冠狀病毒嗎?.
失業嗎?.
困苦嗎?.
逼白嗎?.
飢餓嗎?.

$^{601}$赤身老體危險嗎?.
然而靠著我們的主.
在這一切的事上.
已經得勝有餘了.
有沒有經歷過?.
各位在鏡頭前的.
親愛的弟兄姊妹.
你有沒有經歷過?.
甚麼叫得勝有餘呢?.
我最近有一個很特別的經歷.
在臉書的朋友可能知道.
我幾個星期前.
在山頂.
跌了一跤.
在山頂的餐廳.
不小心.
我又沒有看路.
踩了凸起的磚頭.
向前衝.
跌了到地上.
兩隻門牙其中一隻崩了.
當然這裡也跌得很厲害.
我記得跌倒的那一刻.
我都不知怎樣.
整個人有點暈暈的.
我一個星期後.
就要在自己的舞會講道.
講不講呢?.
接不接呢?.
我到底傷成怎樣呢?.
一連串的問題.
當然當時是冠狀病毒.
我都不知道地上是否乾淨.
但我老公很好.
叫我去看何牧師.
他說你先起來.
坐下.
坐下後.
他為我按手做了一個祈禱.
他說「傑索」.

$^{641}$求你給司母.
在五天之內你醫治她.
讓她康復.
難阻撒旦一切的難阻.
她的工作.
讓她可以在同福港島誤傳.
因為我們知道她在裡面是得勝有餘的.
結果.
三天之後.
傷痕好了.
這裡的痛.
開始消退.
而在那星期裡.
最奇妙的是.
我帶了一對夫妻去信耶穌.
我當時是.
肩膀是痛的.
但是.
我不知道你明不明白.
有一種喜樂.
你忘記了那些痛.
那些痛是超越了.
那種喜樂是超越了我剛才說的痛.
在.
愛我們的主神.
在一切的事.
得勝有餘.
Do you believe it?.
另外.
「凡從神生的就勝過世界」.
讓我們勝過世界.
就是我們的信心.
勝過世界就是那個信耶穌.
是神的兒子的人.
另外.
我最近聽了一個討論.
我們未來的孩子需要的.
Skill set.
能力.
我將它跟聖靈的果子.

$^{681}$當我們信了耶穌的人.
有聖靈的果子.
做個對比.
發覺很多東西都像.
我們要有忍耐.
有恩慈.
有良善.
有信實.
我們是可以選擇.
可以戰鬥.
還有最後.
有彎轉.
甚麼叫有彎轉呢?.
新做的人.
新開始.
若有人在基督裡.
他就是新做的人.
舊事已過.
都變成新的了.
我不知道在鏡頭前.
你有沒有這樣想過.
過去好像很拖累我們.
我們一直拖著他走.
走得很辛苦.
覺得前面沒有去路.
就好像我最近見到一對夫妻.
太太有重病.
他們跟我說.
師母.
怎樣開始?.
要做手術.
那怎樣走?.
我說你信靠耶穌.
舊事已過.
祂會讓你的心靈有新的開始.
你有新的力量.
祂會跟你同在去面對.
我最近很喜歡說一句英文.
It doesn't matter.
Let's start all over again.

$^{721}$就好像電腦一樣.
你想不想重新開始?.
想吧.
聽聽我的故事.
我的故事叫做《I am nobody》.
在前一星期淑媚姐訪問.
我幫她改的.
《I am nobody》.
我成長在一個.
上有一個很出名的彈琴姐姐.
然後下有一個.
我爸爸等了很久的兒子的家庭.
我是中間那個類.
我出生只有三磅.
媽媽就給天主教修女幫我受洗.
所以我有個姓名.
叫做Shirley Marie Therese.
是很有型的.
我到現在都不肯改.
我很記得.
當我後來奉獻做傳道.
我媽媽就說.
將我奉獻給主.
是我一出生就做的.
另外就是我弟弟的滿月酒.
我弟弟出生有滿月酒.
我對自己沒有信心.
給大家看看下一張.
是一個低著頭的我.
不敢望鏡頭.
覺得自己很不堪.
覺得自己在家裡多餘出來的.
I am nobody.
我記得以前我爸爸的朋友.
一見到我爸爸就說.
「這個是你女兒羅生」.
「是不是很會彈琴的那個?」.
我爸爸說「不是」.
「沒有人有興趣問我叫什麼名字」.
所以我覺得自己是個nobody.

$^{761}$差不多我讀中學的時候.
我爸爸就幻想了很嚴重的.
應該叫抑鬱症也好.
精神衰弱也好.
重度的精神病也好.
我想都很重.
因為他要用shock treatment.
要用電擊.
他怎樣重呢?.
他經常自殘.
傷害自己.
在家裡無端端會發脾氣.
他和媽媽都是股票經紀.
我們一句話不知道.
爸爸可以將整台電台倒寫.
又或者他問我一些問題.
我回答得不對.
他就拿一把黑色的剪刀.
我們以前那些.
插自己的胸.
或者拿玻璃瓶打自己的頭.
由於我是中間的女兒.
我姐姐已經離開去別處讀書.
我弟弟,爸爸,媽媽不讓他知道.
他們倆吵架.
所以我是他們倆的中間人.
但一個很痛苦的中間人.
因為我當時很憎恨我爸爸.
我覺得我爸爸有病.
我知道他有病.
但他這樣對媽媽.
我覺得很難饒恕.
所以我在中五那年就去了美國讀書.
誰知中六那年回到香港.
剛剛股票大跌.
爸爸媽媽就變成兩隻大閘蟹在家.
有一天回到家.
看到爸爸在我的房間裡.
攤開了我的日記.
我的日記寫什麼呢?.

$^{801}$我說我很憎恨我爸爸.
我很想他死.
我爸爸看完.
就跟我說:女兒.
我這麼疼你.
我送你去外國讀書.
你今天居然說你想我死.
我還記得我跪下.
我說:爸爸.
我不是有心寫的.
那時我生氣.
我爸爸說:不是的.
白紙黑字不能賴.
從今天開始.
你不要叫我做爸爸.
我不認你做女兒.
當時我十七,八歲.
我一想到就進廁所拿著一張滴露.
我想喝.
不過在要喝的時候.
有個聲音跟我說.
你不如找耶穌.
我說:不行.
我在基督教學校讀書.
我當時對基督徒覺得.
他們很軟弱才相信.
又取笑基督徒.
怎麼好意思呢?.
聲音說:你現在也想喝滴露.
為什麼不試試呢?.
結果.
我去到多倫多.
被同學勸我回教會.
我就回了教會.
那天教會是祈禱會.
祈禱會結束後有個執事跟我說.
你是新來的朋友嗎?.
我說:是的.
你叫羅乃軒嗎?.
他說:你有聽過耶穌嗎?.

$^{841}$我說:我都聽過.
那你信了耶穌沒有?.
我說:沒有.
那你想不想信?.
我說:是什麼來的?.
他當時說了一句話.
他說:你知不知道.
在天上有個爸爸.
他一直在找你.
當我在地上的爸爸不要我.
要跟我斷絕關係的時候.
我聽到天上的爸爸在找我.
我說:是的.
他在找我.
他說:是的.
祈求就得著.
尋找就尋見.
他問我.
你覺不覺得自己有罪?.
我以前不覺得.
我有沒有殺人放火.
我哪有罪?.
但那天我覺得自己有.
我說我.
明明有個爸爸.
我應該愛他.
我愛不到.
我很憎恨他.
這是我的罪.
他說:很簡單.
你只要願意認你的罪.
求耶穌基督的寶血洗淨你的罪.
祂就會帶領你的一生.
當信主耶穌.
你和你一家都別得救.
我就是在那天做了這個祈禱.
結果幾個月後.
爸爸告訴我.
他開始原諒我.
他開始回教會.

$^{881}$那張照片是我1981年.
很年輕.
去台灣.
和何牧師一起做宣教士.
爸爸媽媽來探望我.
開開心心.
幸幸福福的全家福.
我看著爸爸的病.
逐漸逐漸好.
看著他和媽媽開始.
恩愛牽著手.
過著很幸福的家庭生活的那一刻.
1994年.
媽媽突然之間中風.
走了.
整個家庭又陷入幽谷裡面.
我還記得當時.
爸爸問我為什麼.
我說我答不了為什麼媽媽會走.
因為我和媽媽關係很好.
我都接受不到她會走.
但我知道.
我要做的就是好好地照顧爸爸.
爸爸在那十年.
很奇妙地.
他真的慢慢慢慢.
在那個情緒低谷.
走出來.
他做什麼呢?.
他每天.
用四個小時.
你不需要看到他什麼字.
我都看不到.
他用他的孫子的練習簿.
抄聖經.
抄四個小時的聖經.
他留給我的遺產.
就是這些練習簿.
而他一直.
熟讀上帝的話的時候.

$^{921}$他的生命開始改變.
他每天.
和他的股票朋友吃飯.
就勸他們.
你們快點信耶穌吧.
所以.
我在爸爸走了之後.
我聽到他有股票朋友信耶穌.
我堂哥信耶穌.
我相信我爸爸是很開心的.
接著看到這張照片.
就是我爸爸九十歲那一年.
他身體開始越來越差.
心臟病,糖尿病.
高血壓令他走路.
是站不穩的.
但當時是我和何牧師結婚二十五週年.
我就說爸爸.
我結婚的時候.
你沒有來參加.
可能你不知道我嫁了一個牧師.
你沒有來.
但我們結婚二十五週年.
很想多走一次婚禮.
你可不可以帶我走出來.
他說好.
結果.
他就.
顫抖的腳.
一步一步帶著我.
讓我多走一次婚禮.
但當時其實他很害怕.
因為他夢見媽媽.
說在天國車站叫他上車.
於是有一天.
他的飛龍叫我.
「你回來吧!你爸爸要死了」.
我說「要死?」.
我沒有理會他.
可能你們有些人也會這樣.

$^{961}$沒有理會.
然後.
飛龍又再叫我.
「回來吧!你爸爸要死了」.
沒辦法.
扔下所有工作.
飛車回到爸爸家.
進到門口我就問飛龍.
「你爸爸要死了嗎?」.
他說「要死」.
門口我就問飛龍姐姐.
「他要死了嗎?」.
他說「不」.
「他只想見你」.
即是爸爸想見你.
當時我不開心.
我有些生氣.
「爸爸!我其實在工作」.
「你為什麼叫我回來?」.
但我知道那個心態不對.
於是我就向上帝祈禱.
「主啊!求你改變我的心」.
讓我好好的和爸爸聊天.
接著我的心開始平靜.
進到房間.
爸爸捉住我的手.
他說「女兒」.
「我很害怕死」.
「你可不可以幫我做個祈禱」.
我說「做什麼祈禱?」.
他說「祈禱就是」.
「我走到好像你媽媽一樣」.
「一模一樣」.
「中風」.
「然後幾天就走了」.
我說「好啊」.
應該是2014年.
我媽媽是2004年.
8月.
其中一個星期六的下午.

$^{1001}$爸爸和我吃了一頓很好的下午.
他整頓就在那裡讚.
上帝給他一個好老婆.
老婆怎樣好.
上帝給他一個好大女兒.
大女兒怎樣好.
上帝給他一個好兒女.
兒女怎樣好.
三兒子怎樣好.
說了兩個小時.
說到我的心都甜了.
我還記得送他上車的時候.
我要衝上前.
我說「爸爸」.
摸了他的額頭一下.
又摸了他的臉一下.
「爸爸」.
「拜拜」.
「拜拜啊,爸爸」.
結果.
就是最後的再見.
最後的再見.
爸爸那晚中風.
昏迷.
幾天之後.
真是如他所願.
回到天教.
我的故事很真實.
很多人問我.
有一個這樣的爸爸.
怎樣走過來.
是慢慢一步一步走過來.
但我很想告訴大家.
如果沒有耶穌.
今天我不會站在這裡.
如果沒有了耶穌.
我看不到很多家庭.
因為耶穌的愛而復和.
我們有很多技巧.
我都懂得.

$^{1041}$但最好的.
所謂方法.
就是順靠主耶穌.
到最後.
很想跟大家.
分享一個.
開心一點的畫面.
除了我們機構的畫面.
歡迎大家收看.
「雙劍合璧」.
或者.
我們其實都有很多故事.
短片在我Facebook專頁.
和家庭發展基金.
不過想和大家分享一個.
很開心的片段.
就是我有一個六姨丈.
這個六姨丈.
是一生篤信耶穌.
他在八十歲那年太太過世.
自己主動入老人院.
因為上帝要他安慰老人家.
九十歲學水彩畫.
畫很多水彩畫送給老人家.
然後一百零一歲的時候.
我有機會去探望他.
我表哥帶我去探望他.
入到老人院原來很有趣.
一個人的個性是怎樣呢?.
你年紀大了就會顯露.
很多人在埋怨,罵人.
但六姨丈很安詳地坐在那裡.
我就問他.
「六姨丈,你怎樣?」.
他說「有耶穌當然好」.
表哥說.
如果上帝有恩典.
你會聽到他背誦的那首.
人生影響他很多的詩.
所以在現在這個時間.

$^{1081}$很想讓你聽聽.
我六姨丈親口背誦的這首詩.
是完全沒有經過剪接.
是他由心底裡相信而活出來的.
四代基督徒到現在.
好,請大家欣賞看看.
(對,那首詩是怎樣的?).
(在套書裡).
(在套書裡?).
(套書).
(講一下那首詩吧).
(他說).
(靈自加靈).
(亦仍是靈).
(是).
(無論加減或隨性).
(前來先行仍依舊).
(多出親民也不能).
(若在靈田加蓋一).
(今時就不圖往日).
(後便靈田越加多).
(數目越來越大不一).
(請問前面一是誰).
(原來就是主耶穌).
(哇,厲害).
(虛空有主便充滿).
(謙卑隨時便有餘).
(記得一字擺在前).
(千祈不可在後面).
(基督若非居首位).
(靈靈依舊不值先).
(哇).
(這首詩是不是很好?).
(虛空有主便充滿).
(謙卑隨時便有餘).
剛才經文說.
若有人在基督裡.
他就是新造的人.
舊事已過.
都變成新的了.

$^{1121}$今日我記得有人在Facebook問我.
如何擺脫爸爸的陰影.
原生家庭的陰影.
我不知道你有沒有在電視機旁邊.
在電腦旁邊看.
我也不知道你信了耶穌沒有.
今天很特別.
我有機會和我的網友.
或者在旺角浸信會的會友.
或者你們的朋友一起來分享.
在這個時間.
很想大家讓你的心對你來說話.
你可以閉上眼睛.
回想剛才所聽到的.
有沒有什麼是真的對你說話呢?.
我深信聖靈會對我們每一個人來說話.
如果聖靈真的對你這樣說.
你很願意在今天.
你選擇在今天.
讓一切有新的開始.
舊事已過.
過去的事抹掉.
有新的開始.
靠著耶穌有光明的未來.
有盼望的未來.
如果你願意這樣.
我請你在鏡頭前.
你用你的右手.
放在你的心上.
向上帝祈禱.
我講一句.
你跟著我講一句.
因為親愛的主耶穌.
我承認我的罪.
我知道我曾經得罪你.
求你十字架的寶血.
洗淨我的罪.
帶領我和我一家.
讓我可以靠著你的恩典.
一生跟隨你.

$^{1161}$奉主耶穌基督的名字求.
阿們.
在四十多年前的某一個晚上.
你問我是否很像基督教信仰.
我認不認識聖經.
我不認識.
但當時我心裡有一個觸動.
耶穌.
我願意認識你.
我願意跟隨你.
當時我十八歲.
直到今天.
祂改變了我一家.
改變了我.
我親眼看到我身邊很多人的家庭被改變.
所以我希望你能夠認識祂.
你能夠遇見祂.
我的分享到此為止.
將時間交給牧師.
祝福.
\newpage



\section{哥林多後書 6:11-18}
\label{sec:m_Hc7J2x7vY}
\textbf{2020年12月13日 崇拜直播｜哥林多後書6\_11-18}
\newline
\newline
連結: \href{https://youtube.com/watch?v=m-Hc7J2x7vY}{\texttt{https://youtube.com/watch?v=m-Hc7J2x7vY}} ~~~~ 語音日期: 2020-12-15
\newline
\newline
\hyperref[sec:MIn0I7TUOfE]{\small{< < < PREV SERMON < < <}}
~
\hyperref[sec:index_chronic]{\small{[返順時目]}}
~
\hyperref[sec:TPSrTFQecuk]{\small{> > > NEXT SERMON > > >}}
\newline
\newline
哥林多後書 6:11-18
\newline
\begin{longtable}{cl}
\hline
\hline
章節 & 經文 (和合本修訂版)\\
\hline
6:11 & \begin{tabularx}{0.7\textwidth}{X} 哥林多人哪，我們對你們，口是誠實的，心是寬宏的。 \end{tabularx} \\ \\ \relax
6:12 & \begin{tabularx}{0.7\textwidth}{X} 你們的狹窄不是由於我們，而是由於你們自己的心腸狹窄。 \end{tabularx} \\ \\ \relax
6:13 & \begin{tabularx}{0.7\textwidth}{X} 你們也要照樣用寬宏的心報答我；我這話正像對自己的孩子說的。 \end{tabularx} \\ \\ \relax
6:14 & \begin{tabularx}{0.7\textwidth}{X} 你們不要和不信的人同負一軛。義和不義有甚麼相關？光明和黑暗有甚麼相連？ \end{tabularx} \\ \\ \relax
6:15 & \begin{tabularx}{0.7\textwidth}{X} 基督和彼列有甚麼相和？信主的和不信主的有甚麼相干？ \end{tabularx} \\ \\ \relax
6:16 & \begin{tabularx}{0.7\textwidth}{X} 神的殿和偶像有甚麼相同？因為我們是永生神的殿，就如神曾說：「我要在他們中間居住來往；我要作他們的神，他們要作我的子民。」 \end{tabularx} \\ \\ \relax
6:17 & \begin{tabularx}{0.7\textwidth}{X} 所以主說：「你們務要從他們中間出來，跟他們分別；不要沾不潔淨的東西，我就收納你們。 \end{tabularx} \\ \\ \relax
6:18 & \begin{tabularx}{0.7\textwidth}{X} 我要作你們的父，你們要作我的兒女。這是全能的主說的。」 \end{tabularx} \\ \\
[1ex]
\hline
\hline
\end{longtable}
$^{1}$讓我們繼續聆聽上帝的話語.
記載在哥林多侯書第六章.
十一至十八節.
史圖普羅這樣說.
「哥林多人啊.
我們向你們口是張開的.
心是寬宏的.
你們狹窄.
願不在乎我們.
是在乎自己的心腸狹窄.
你們也要照樣用寬宏的心報答我.
我這話正像對自己的孩子說的.
你們和不信的.
願不相配.
不要同父一額.
義和不義有什麼相交呢?.
光明和黑暗有什麼相通呢?.
基督和比列有什麼相和呢?.
信主的和不信主的有什麼相干呢?.
神的殿和偶像有什麼相同呢?.
因為我們是永生神的殿.
就如神曾說.
我要在他們中間居住.
在他們中間來往.
我要作他們的神.
他們要作我的子民.
又說.
你們務要從他們中間出來.
與他們分別.
不要碰不見證的物.
我就收納你們.
我要作你們的父.
你們要作我的兒女.
這是全能的主說的.
今天早上有梁家倫牧師.
在我們當中繼續宣講.
這一段經文.
講題是開放與封閉.
將時間交給梁牧師.
旺角浸信會在神學傳統上屬於福音派.

$^{41}$在教義立場上繼承基督主義的傳統.
不需要諱言.
今日基督主義甚至是福音派.
常常給人一個不開放的印象.
因為我們在很多教義和倫理的問題上.
總是採取一個嚴謹保守的態度.
簡單一句說.
我們不是像自由主義那樣.
我們拒絕向這個世界妥協.
例如我們堅持聖經無謬誤的主張.
我們相信耶穌基督是人類唯一的拯救.
我們也反對同性戀,婚外情,性濫交等等的行為.
我們這個對真理執著的態度.
常常招來很多人的批評.
我們必須要承認.
這個態度在這個世界.
在今時今日是非常不受歡迎的.
我常常被我事奉教會的大學生在我面前訴苦.
教會持另一個立場.
和他們在大學校園面對的處境.
基本上是水火不用.
大學是要高舉包容多元的地方.
包容甚至成為了最高價值.
所以如果你敢在大學.
你敢在課堂上說我反對同性戀.
你猜後果會是怎樣.
必死.
你一定會被大批的同學孤立排斥.
不向你掟石也奇怪.
所以他們很多同學說.
他們最多只能夠保持沉默.
不說話.
別人大談同性婚姻的時候.
他們就不說話.
退到一邊.
甚至有些壞人就做兩面人.
見人講人話.
見鬼講鬼話.
在大學校園講一套.
在教會講另一套.

$^{81}$你覺得他們虛偽嗎.
我很多大學生的年輕朋友跟我說.
他問為何我們的心胸要這麼狹窄.
為何我們有排他主義的傾向.
很多人問我們為何不能開放一點.
他們建議我們應該相信真理有很多個.
不同人有不同的真理.
或者真理雖然只有一個.
不過有很多個表達真理的方法.
再退一步.
即使只有一個真理.
有一個表達真理的方法.
但可否有很多個不同對真理的理解和領受呢.
領受不同可以嗎.
我聽到弟兄姊妹的困惑.
不過我總是很無奈地說.
我們豈是因為驕傲的緣故.
然後成為排他主義者呢.
我們豈是只為事排斥別人呢.
最少我事奉了這四十年.
我可以用我自己的見證去說.
如果有認識我的人.
知道我過去不是像我以下所說的話.
請篤爆我.
這麼多年來.
我對婚姻出了危機.
甚至無法彌奉.
無法填補裂痕.
不論是不是基督徒.
我總是存著一個同情和諒解的態度.
只希望伸出手去給一些幫助.
我很少在他們面對困境的時候.
我嚴加責備.
我自己也有幾位.
我現在數到的有三位.
同性戀者的朋友.
我從來沒有以二人自居.
在心態或行動上排斥他們.
拒絕跟他們來往.
我從來不敢忘記耶穌基督的教訓.

$^{121}$沒病的人用不著醫生.
有病的人才用得著.
我來不是召二人悔改.
而是召罪人悔改.
蒙上帝拯救的我.
是一個不折不扣的罪人.
而教會要收納要親近的是罪人.
不是二人.
所以至少在心態上.
在行動上.
我們不是排他主義者.
問題只是接納罪人和接納他所犯的罪.
是截然不同的兩回事.
我們要慎防自明清高的法利賽人心態.
要拒絕以二人的角度去審判一眾的罪人.
但我們不可以無視聖經在教義.
在倫理問題上的清楚教導.
我們不可以篡改上帝的心意.
將罪不管是自己的還是其他人的.
看成不是罪.
就算我們個人願意無條件接納一個犯罪者.
我們也不可以擅自代上帝發言.
說就因為我們個人的喜惡.
我們個人願意接納他.
上帝就應該.
上帝就理所當然無條件接納他.
我們知道.
不是上帝以我人的喜惡來成為他的喜惡.
而是人要以上帝的喜惡.
作為他自己的喜惡.
是不是?.
我們要學習與上帝同步思想.
我們不敢憑自己的意思去剪裁真理.
因為真理不是由我們自己定定的.
與自由主義不同的是.
我們信奉的是一個有衛家,全權,自主的上帝.
而不是某一些單純的倫理的概念.
公義,愛.
基督教是相信一個愛的上帝,公義的上帝.
而不是基督教相信愛,公義本身就夠.

$^{161}$我所信的是一個.
擁有至終真理權威的聖經.
上帝藉著聖經.
向我們去曉喻他不變的心意.
所以聖經是inspired word.
而不僅僅是一本inspiring.
充滿啟發性的一本古典作品.
可以給我們任意的詮釋.
信仰理解的不同.
我們的行事為人就有很大的差異.
我們當上帝是上帝.
當聖經是聖經.
所以我們不敢在大是大非的問題上妄自敬作.
將自己的感受,觀點篤落聖經那裡.
或者自封做上帝的發言人.
宣告他已經成為一個虛無主義.
或者相對主義的上帝.
上帝如果在某一個問題上做了他的啟示.
即使那個啟示不合我的胃口.
不讓我的感受,我的理性所認同.
我們仍然需要恭謹的去領受.
信服,遵行.
讓我們深知.
人在親近上帝的時候.
必須要持絕對開放的態度.
不應該帶著各種宗教或者非宗教的成見.
我們過濾上帝的說話.
規限他的作為.
否則我們就將宇宙的主.
將他馴化.
將他domesticate(家庭化).
變成我們的家臣.
用我們自己的形象去塑造上帝.
上帝以他的形象塑造人.
但很多時候人是用自己的形象去塑造上帝.
我們願意對這個世界開放.
但我們更加願意向上帝開放.
正因為我們願意向上帝開放.
向對真理開放.
很抱歉.

$^{201}$我們就不可以同時期.
向這個世界全面開放.
保羅在哥林多後書六章十一到十三節.
做了一個這樣的對照.
他說.
哥林多教會的弟兄姊妹.
總是以為他是一個心胸狹窄的人.
在神學觀點上難容異己.
在倫理抉擇上.
也不喜歡給人太多的自由.
例如他強迫哥林多的弟兄姊妹.
清理門戶.
將犯奸淫的人開除出教.
又彈劾那些在神學立場上跟他不同的人.
不容許人按著心中的靈感.
自由地去創造新的真理.
但是保羅卻辯護說.
他對哥林多的弟兄姊妹是開放的.
他接納他們.
愛護他們.
向他們傾注所有的感情.
他說.
哥林多人啊.
我們向你們口是張開的.
心是寬宏的.
反倒他說.
是哥林多的弟兄姊妹對他.
對保羅心胸狹窄.
不願意接納他.
不願意聽從他的意見.
他們為什麼不願意接納保羅呢?.
是不是因為他們在思想上.
已經充塞了各種自以為是的開放觀點.
而再沒有空間去留給保羅呢?.
如果不是保羅排斥哥林多人.
而是哥林多人排斥保羅.
那到底是誰心胸狹窄呢?.
誰排斥誰就是誰心胸狹窄了.
是不是?.
是的.

$^{241}$誰是真正心胸狹窄的人呢?.
這是一個觀點和角度的問題.
哥林多人總是以為他們自己是心胸寬宏的人.
因為他們願意向世界開放.
接納各種和基督教信仰相衝突的想法和做法.
凡事不計較不執著.
但是對保羅來說.
向上帝和真理持寬宏開放的態度.
才是最重要的.
他必須要接納上帝.
不按人的個人喜好和想法.
來啟示他的心意.
指示人當走的路.
唯有這樣做.
一個基督徒才是真正心胸寬宏的人.
一個向罪寬宏的人.
必然同時向真理狹窄.
一個向世界寬宏的人.
必然會向上帝狹窄.
相反的.
一個向真理寬宏的人.
就無法不對那些和真理相敵的罪.
持排斥的態度.
他要對上帝開放.
容許上帝在自己生命裡自由的行使主權.
就不能不對自己原本的想法堅決拒絕.
《解體書》五章十七節說.
因為情慾與聖靈相爭.
問題不在於.
誰寬宏誰狹窄.
而在於向誰寬宏向誰狹窄.
這個世界沒有左右抱抱.
左右逢源.
both end 這回事.
沒有一腳踏兩船.
一個人不可能向所有東西開放.
如果他自誇向所有東西開放.
他就根本從來都沒有委身在任何一件事上.
是不是.
他對真理或者信仰的看法.

$^{281}$僅僅是相對主義.
甚至是虛無主義.
他的人生態度也是遊戲人間式.
凡事都無所謂.
無所堅持無所執著.
是不是.
同樣的.
一個人不可以同時愛所有的人.
如果他說他愛所有的人.
他等於他沒有愛過任何人.
他未曾試過為心中所愛的對象.
茶飯不思朝思暮想.
是不是.
愛總是排他的.
我們在婚禮裡.
和身邊的配偶說yes.
我們就同時對其他所有的異性say no.
是不是.
耶穌不是說過類似的話嗎.
馬達福音六章十二十四節.
一個人不可以同時事奉兩個主.
不是污這個愛那個.
就是重這個輕那個.
你不可以既事奉上帝.
又事奉馬門.
我們不可以既稱基督為主.
又同時凡事自作主張.
不可以既跟隨基督.
又對世界倦倦不捨.
沒有拆毀就沒有建立.
沒有捨棄就不會得著.
兼容並包.
這個是那些巴開大同教.
我們中國傳統的民間宗教.
可以炒在一起.
什麼偶像都放在同一個祭壇裡.
這個不是基督教信仰.
基督徒應該向上帝心胸寬宏.
向真理心胸寬宏.
卻不要對罪寬宏.

$^{321}$向敵黨真理的觀念和行為.
闖開我們的心門.
雅各書四章四節說.
凡想要與世俗為友的.
就是與上帝為敵.
接著.
保羅向他的讀者解釋.
怎樣才不向罪寬宏.
他用了五個不相作為說明.
他們的意思其實是相仿的.
信和不信的不相配.
義和不義的不相交.
光明和黑暗的不相通.
基督和比列就是撒旦不相和.
上帝的殿和偶像不相同.
在這五個不相當中.
我們最熟悉的當然是第一個.
信與不信 緣不相配.
不要同父一額.
我們恆常將它應用在婚姻上.
宣稱基督徒不應該和非基督徒通婚.
這個應用我不反對.
不過這個不是經文原來要carry的意思.
因為保羅在這裡說的不是婚姻的問題.
第二 保羅甚至在這段經文裡.
說的不是基督徒和這個世界和人際關係的問題.
他首先說的是基督徒和對自身的問題.
他不是說基督徒不應該和非基督徒同坐一架巴士.
同坐一架地鐵 同父一額.
如果是的話你沒有車錯.
不是說你不可以和非基督徒打同一份工.
同父一額.
如果是的話你就差不多沒有工打了.
好像我這樣在教育機構事奉會好一點.
否則的話你基本上沒得做了.
不是 主要不是說這件事.
因為他後面說什麼.
在六章十六節裡說.
我們的生命是永生上帝的電.
所以其實保羅是針對基督徒的個人內在生命.

$^{361}$整段經文主要是提醒你不要讓罪在你的生命裡.
即是佔一席位.
信與不信不能相判.
基督與比列不能混在一起.
是在說這裡.
在說你裡面.
裡面理論上被基督霸佔了.
當你被基督霸佔了之後.
理論上你沒有位置放個撒旦在這裡.
你這裡要乾淨.
這裡要純水.
這裡不可以混雜.
我們活在一個混雜的世界.
你說基督徒不可以跟非基督徒相交.
這是開玩笑.
連福音使命都不用做了.
我們自己自說自話吧.
關上門.
弄個基督教俱樂部就行了.
不是 他說這裡.
整段經文首先說的是這裡.
這裡要乾乾淨淨.
要純純水水.
這裡不可以信與不信相交.
不可以光明與黑暗相同.
不可以異同不異相交.
這裡不可以.
基督徒活在世上.
既不能夠也不應該.
拒絕與未信的人來往.
相處甚至深交.
是不是可以結婚是另一個問題.
但至少你說我不可以跟非基督徒深交.
只能夠泛泛之交.
那你要全福音就很難了.
怎樣做友誼報道.
所有友誼都是為了過橋.
哄他入局.
就好像那些做推銷的.
找你做朋友其實是想你買他的東西.

$^{401}$這樣卑不卑鄙一點.
我們基督徒就算多支持全福音都不應該這樣做.
友誼不是純粹一個手段.
是不是.
有嫌可能是手段.
你怎麼可能拒絕與非基督徒來往.
這是觀念上在實踐上在使命上完全不通的.
完全不通的說法.
但我們要拒絕.
讓罪在個人生命裡佔有位置.
這才是這段經文教訓的重點.
保羅精印了以賽亞書和以西結書.
四句舊約的經文.
說明上帝願意收納我們做他的兒女.
他要在我們中間來往居住.
因此我們必須要敬虔自首.
使自己與世人分別出來.
不要碰不潔淨的東西.
好使我們身子可以被上帝停留.
上帝願意停住在我們生命裡.
帖撒尼加前書四章七節都說.
上帝召我們.
本不是要我們滇染污穢.
而是要我們成為聖潔.
一句話.
基督徒不可以與罪相交.
不是說不可以與罪人相交.
是不可以與罪相交.
但你必須要與罪人相交.
對不起.
剛剛好相反.
一定要與人相交.
如果我們每個人都不與人相交.
我們這間旺鎮就不用傳福音了.
倒閉都可以了.
是不是.
不可能.
弟兄姊妹.
我們不可以對罪做任何妥協.
這是一個很簡單又寶貴的真理.

$^{441}$不過實踐起來是很不容易的.
身處一個罪惡滔天的世界.
活在一個和教教的道德標準.
自然不同的人群裡.
要完全避免罪的引誘.
剷除所有可能令我們跌倒的因素.
幾乎是不可能的事.
走在鬧市.
方言望去.
你看到的都是那些淫穢的東西.
聽到的也很多是那些下流邪惡的說話.
除非我們與世隔絕.
不與世人並居.
否則根本無可圖遁.
防不勝防.
你所能夠做的.
是盡量在罪惡.
在你自己生命裡萌生的當下.
立刻禁息它.
不讓它發芽生長.
這就是我們中國人說的.
責無尊重.
如果我們的視線.
偶然接觸到一些淫穢不潔的東西.
聽到一些彼俗褻瀆的話題.
盡可能抽身而出.
或者最少關閉我們感官的接收器.
遠離這些令我們心懸意馬的試探.
不要讓自己在罪惡裡耽擱.
如果我們心裡突然冒出一個惡念.
例如對人的疾度仇恨.
淫念或其他污穢的思想.
就要即時的.
狠狠的將它鎮壓.
強力將它除去.
不要讓它在自己腦海裡有進一步滋長的機會.
我們可以甩一下頭.
罵自己兩句.
想這些壞事.
或者逼自己中斷思想.

$^{481}$將注意力轉去一些建設性的事上.
甚至暴力一點.
摑自己一巴.
扔自己胸口一下都可以.
當然你不需要做到天主教那些打虎鞭的做法.
總而言之不要讓魔鬼留地步.
不要含匿在罪惡裡.
我常常說的一句套話.
我們沒辦法防範.
讓自己看到第一眼的罪惡.
但我們必須要努力.
不讓第二眼發生.
我們不需要為第一個惡念負責.
我們要為這個惡念繼續盤旋而負責.
第二才算數.
要完全沒有第一.
除非你是一個超凡入勝的人.
否則很難.
尤其要緊的是.
我們不要讓惡念兌現成為真正的惡行.
弟兄姊妹.
你有沒有誤會我在這裡宣揚行為得救的觀念.
認為基督徒一定要藉著家賢善行來討上帝喜悅.
如果信主之後繼續犯罪.
就會激怒上帝.
就會失落舊恩.
我不是說接納天主教的那些律法主義主張.
認為基督徒一生一定要做到收支平衡.
功過相抵.
有多少merit有多少demerit.
用merit去掩蓋demerit.
如果你離世之前.
還有罪債未曾償清的話.
你就沒辦法直接上天堂.
我沒有這樣的想法.
我關心的其實不是犯一條罪.
甚至兩條罪.
三條罪.
四條罪.
會帶來多少刑罰的問題.

$^{521}$我講這些話講到幾日日日.
我希望你捉到我想講什麼.
基督徒犯一條罪兩條罪三條罪四條罪.
其實不是最主要的問題.
你犯一條罪會有什麼刑罰.
兩條罪有什麼刑罰.
三條罪有什麼刑罰.
上帝是一個完美主義者.
你總之有多幾條罪.
根深就會一個雷劈死你.
我不是這樣想的.
我是講的是.
當你容讓罪在你生命裡面.
會對你的生命帶來多少破壞的作用.
或者孤立的一條罪兩條罪.
未必有很大的影響.
好像保羅在羅馬書五章二十字那裡講.
罪在那裡顯多.
恩典就顯得更多.
罪顯示恩典.
但是犯罪很少是孤立事件.
它總是在人的生命裡面帶來連鎖性的影響.
而這些連鎖性的影響.
我要告訴你才是真正致命的.
當我有意或者不慎犯了一個罪之後.
最好的處理方法是怎樣.
即刻誠心認罪悔改.
不要兜,不要掩飾,不要遮蓋.
是不是.
不過人的反應通常不是這樣的.
能夠即時有錯就認.
這個人是自律性很高的人.
他繼續犯罪的機會都不會太多.
我要講.
因為你要維持這個真誠認罪的態度.
比保守自己不犯罪難度更高.
法理賽人成日都戰戰兢兢保守自己不犯罪.
某程度上不是做不到的.
他們做到的.
不過真誠認罪是難過也重要過保守自己不犯罪.

$^{561}$你記不記得耶穌讚那兩個在聖殿祈禱的人.
讚誰.
是不是.
我們最正常的反應.
是為自己犯的罪尋求各種開脫的藉口.
如果我不可以徹底將我所犯的罪辯解到不成罪.
也要減低他的嚴重程度.
而且要凸顯犯罪的正常性與不可避免性.
人誰無過呢.
很正常.
當然正常.
我看見婦女動淫念當然不正確.
不過美的婦女例外.
對不對.
不美的當然不會動淫念.
美的人之常情.
我也是男人.
男人一個沒辦法.
對不對.
我們很容易去強調正常性.
不可避免性.
藉此去消去因為犯罪而帶來良心的不安.
人總是會為自己的錯誤辯護.
亦都.
我告訴你.
人總是能夠成功的為自己辯護.
我沒見過一個辯護不了的人.
除非他真的蠢到+01.
你知道我們自我辯護.
我們是自說自話.
只有被告沒有原告.
我自己要搞到自己沒罪有什麼難.
都沒有檢控官來指控我.
我都壓死了良心.
性命都沒了位置.
所以我要解釋到自己沒罪太容易了.
如果我犯了某一個罪行.
我一定可以找到一個使自己開脫的藉口.
最一般性的.
例如人誰無過.

$^{601}$性人都會失手.
小癡不染大魚.
是不是.
或者我的道德水準已經是一般水平.
Above average(超過平均).
對不對.
連我兒子在十歲的時候.
他每一次犯錯都找到開脫的理由.
我沒見過他會搞到持窮履屈.
沒試過.
我沒教他做辯護.
他沒讀法律的.
不過總之他就口囊過油.
一定找到理由.
你信不信.
一旦我們為自己的罪找到開脫罪責的理由.
我們之後就很難再理直氣壯的拒絕相同的罪.
因為我們已經為那個罪的存在.
賦予了若干程度的合法性和合理性.
解除了對他的理性的防禦武裝.
在缺乏理性支撐之下.
我們的意志很難作出一個堅強的守護.
結果很容易就會重蹈覆轍.
一件事兩件事.
反正都做了.
反正都覺得合理.
合理不怕再做一次.
再做第三次.
太容易了.
吃第一口毒品.
通常不會讓人染上毒癮的.
所有染上毒癮的都是因為第一口開始的.
是不是.
第一口開始的.
吃一口就上癮.
沒那麼誇張的.
沒那麼激的.
好像雪球一樣越滾越大.
我們對某一個罪採取寬宏的態度.
就很自然的對其他罪都開放起來.

$^{641}$很少人會只是慣性的犯一種罪.
但是對其他的罪嚴加防範的.
我沒見過.
在多數情況之下.
我們或者是對所有罪都責罵.
警醒提防.
或者是都採取放任的態度.
畢竟.
當我自知私生活.
某一個片段不夠檢點.
人格不完整.
表裡不一致.
我們就整體性會失去道德勇氣.
很難再嚴於律己.
一點一滑.
要求自己兌現信仰的要求.
反正真是.
就像剛才所說的.
一件污穢.
兩件一樣污穢.
有什麼分別.
是這樣.
最後.
我就自行解除了對所有罪的防範.
不再要求自己變好.
不再相信自己能夠變好.
我們或者對道德絕望.
或者對自己的道德能力絕望.
我們在追求誠信的路上.
就會自暴自棄.
如果我們仍然留在教會裡.
我們很容易就會成為人格分裂的偽君子.
耶穌批評的那些假無畏善的人.
口講一套.
實際做的另一套.
我們不相信自己所說的理想.
我們甚至會變得憤怒.
我們對人人所說的理想.
都不相信.
教會啊.

$^{681}$牧師啊.
掌執啊.
滿口仁義道德.
做的都是陰間的事.
我怎麼知道.
我看自己就知道.
他怎麼比我更了得.
我是這樣.
人人都是這樣.
我曾經有一次不小心.
被一塊很小的木刺.
插到手指.
我以為我已經把刺都拔出來.
誰知道還有一條刺留在皮膚裡.
結果外面的傷口癒合了.
我看不出有什麼不同.
但裡面的肌肉.
就因為剩下的那條刺.
就潰爛了.
所以不是一條罪兩條罪.
會不會破壞我們的救恩問題.
而是犯罪的行動本身.
會腐蝕我們的屬靈生命.
令我們失去了對罪的防範.
放棄對信仰一致性.
人格完整的堅持.
結果我們整個屬靈生命就會破壞.
所以基督徒要嚴肅的守護自己.
隨從聖靈的引導.
不要讓罪惡的念頭滋生.
不要讓惡念變成惡行.
耶穌曾經很嚴厲的教訓門徒說.
如果你的一隻手或一隻腳跌倒.
就撞下來丟掉.
你只有一隻手或一隻腳進入永生.
強如有兩手兩腳被雕在永火裡面.
如果你的一隻眼跌倒.
就把它彎出來丟掉.
你只有一隻眼進入永生.
強如有兩隻眼被雕在地獄的火雷.

$^{721}$《馬特福音》18:8-9.
你覺不覺得很恐怖.
耶穌說這個密的說話.
但你看清楚.
如果你相信聖經的平白的技術.
在《馬特福音》耶穌說這番說話.
是說過兩次的.
《馬特福音》五章29-30節說了一次.
18章8-9節又說了一次.
即是說這番說話.
如果馬太不是自己不小心重覆了記載.
就是耶穌說了不止一次.
這些很噁心的說話.
這些很恐怖很三級的說話.
我們主不止說了一次.
無論如何.
弟兄姊妹.
今日的訊息其實很簡單.
對自己狠一點.
決絕一點.
但對人間對世界卻不是這樣.
我們要堅持真理.
但我們一定要確定.
甚麼是我們非堅持不可的真理.
盡量在次要的事上.
我們包容一點.
不要在心態上搞全面的對立.
特別不要在心中懷仇恨.
或者有被迫害的感受.
包括如果你是一個未信主的家庭.
你的家人迫你很多事上參與.
或者反對你做很多事.
你都需要持這個態度.
一些最關鍵的問題上.
堅持但不要事事都搞對立.
盡可能與人和善.
不要每件事都搞對立.
當然如何畫線.
都需要上帝給我們智慧.
但無論如何.

$^{761}$千萬不要在家中搞對立.
因為我是基督徒.
我看不起你們這班人.
多些關愛.
多些憐憫.
我們總是記住.
不入虎穴 焉得虎子.
我們現在是張臨期.
我們總是記住.
耶穌基督到誠肉身.
祂來到人間.
就說明我們基督徒絕對不可以獨善其身.
否則我們的主為什麼不在聖潔的天庭就算.
祂為我們示範.
如何來到一個拒絕祂的人間.
我們守住信仰.
說句不做elaboration的說話.
就算面對將來不可知的政治局面.
不要想什麼防禦策略.
教會其實沒有什麼策略可做.
當然你逃跑我不反對.
沒有策略.
守住我們的信仰.
求聖靈幫助我們繼續守住我們的信仰.
讓我們坦然開放去面對即將來臨的各種際遇.
讓我們進入狼群.
不過,make sure,讓繼續是讓.
我們是耶穌基督的讓.
謝謝大家收看 再見.
\newpage



\section{路加福音 2:8-20}
\label{sec:TPSrTFQecuk}
\textbf{2020年12月20日 崇拜直播｜路加福音2\_8-20}
\newline
\newline
連結: \href{https://youtube.com/watch?v=TPSrTFQecuk}{\texttt{https://youtube.com/watch?v=TPSrTFQecuk}} ~~~~ 語音日期: 2020-12-21
\newline
\newline
\hyperref[sec:m_Hc7J2x7vY]{\small{< < < PREV SERMON < < <}}
~
\hyperref[sec:index_chronic]{\small{[返順時目]}}
~
\hyperref[sec:lmnMx2HA46Q]{\small{> > > NEXT SERMON > > >}}
\newline
\newline
路加福音 2:8-20
\newline
\begin{longtable}{cl}
\hline
\hline
章節 & 經文 (和合本修訂版)\\
\hline
2:8 & \begin{tabularx}{0.7\textwidth}{X} 在伯利恆的野外有牧羊人，夜間值班看守羊群。 \end{tabularx} \\ \\ \relax
2:9 & \begin{tabularx}{0.7\textwidth}{X} 有主的一個使者站在他們旁邊，主的榮光四面照著他們，牧羊人就很懼怕。 \end{tabularx} \\ \\ \relax
2:10 & \begin{tabularx}{0.7\textwidth}{X} 那天使對他們說：「不要懼怕！看哪！因為我報給你們大喜的信息，是關乎萬民的： \end{tabularx} \\ \\ \relax
2:11 & \begin{tabularx}{0.7\textwidth}{X} 因今天在大衛的城裡，為你們生了救主，就是主基督。 \end{tabularx} \\ \\ \relax
2:12 & \begin{tabularx}{0.7\textwidth}{X} 你們要看見一個嬰孩，包著布，臥在馬槽裡，那就是給你們的記號。」 \end{tabularx} \\ \\ \relax
2:13 & \begin{tabularx}{0.7\textwidth}{X} 忽然，有一大隊天兵同那天使讚美神說： \end{tabularx} \\ \\ \relax
2:14 & \begin{tabularx}{0.7\textwidth}{X} 「在至高之處榮耀歸與神！在地上平安歸與他所喜悅的人！」 \end{tabularx} \\ \\ \relax
2:15 & \begin{tabularx}{0.7\textwidth}{X} 眾天使離開他們，升天去了。牧羊人彼此說：「我們往伯利恆去，看看所成的事，就是主所告訴我們的。」 \end{tabularx} \\ \\ \relax
2:16 & \begin{tabularx}{0.7\textwidth}{X} 他們急忙去了，找到馬利亞和約瑟，還有那嬰孩臥在馬槽裡。 \end{tabularx} \\ \\ \relax
2:17 & \begin{tabularx}{0.7\textwidth}{X} 他們看見，就把天使論這孩子的話傳開了。 \end{tabularx} \\ \\ \relax
2:18 & \begin{tabularx}{0.7\textwidth}{X} 聽見的人都詫異牧羊人對他們所說的話。 \end{tabularx} \\ \\ \relax
2:19 & \begin{tabularx}{0.7\textwidth}{X} 馬利亞卻把這一切的事存在心裡，反覆思考。 \end{tabularx} \\ \\ \relax
2:20 & \begin{tabularx}{0.7\textwidth}{X} 牧羊人回去了，因所聽見所看見的一切事，正如天使向他們所說的，就歸榮耀於神，讚美他。 \end{tabularx} \\ \\
[1ex]
\hline
\hline
\end{longtable}
$^{1}$第一節會平安.
你看到我們都戴了很有特色的口罩.
就知道我們預備好來迎接耶穌基督的降生.
讓我們都帶著一份興奮期盼來進入今天的崇拜好嗎?.
讓我們一起起立吧.
我們一起祈禱吧.
祈禱.
請不吝點贊 訂閱 轉發 打賞支持明鏡與點點欄目.
主持人:萬民啊!救主耶穌已降生人間.
全地都當歡騰.
併發起大聲迎接他.
向他歡呼歌頌.
地的四極都看見我們神的救恩.
救主耶穌是全地的拯救和盼望.
即使在限聚的聖誕節.
我們仍然能夠在不同的角落.
慶祝救主降生.
還唱普世歡騰.
讓地的四極不單止看見.
更加可以聽見神完備的救恩.
普世歡騰救主降生.
全地全都為我.
萬修羅來祝禰免疫醫科.
宇宙萬物中藥.
國際公民.
誓言共同打仗.
天地共享.
平民與神明.
慶祝禰免疫醫科.
禰人命必頒.
禰也可慶禰了了不定.
祭祀禰受全國祭祀.
告禰作出敬意.
祂主來來是福中人.
主人心口無言.
當主祭祀禰贊美來.
讓道長安分輝.
祭祀禰共意可感全力.
自愛可感全力.
請姐妹讓我們同心獻上祂的禱告.

$^{41}$因為昔日你夕你自己的愛子.
你的獨生兒子.
耶穌基督親自降臨到人世間.
為的是要向世人顯明.
上帝啊,你對我們的愛.
特別當我們還未認識耶穌你的時候.
幻作罪人的時候.
上帝啊,你已經為我們.
預備這份奇妙的救恩.
寶貴的禮物.
上帝啊,我們實在要獻上.
因為你成為人生.
住在我們的中間.
讓我們真的能夠.
看得到你.
如何去明白.
我們的心路歷程.
你也曾經成為人.
有個人的喜怒哀樂悲歡離合.
所以你能夠明白.
我們在生活裡面.
所遇到不同的境況.
並且你能夠成為我們的安慰.
盼望和拯救.
以致我們能夠靠著你.
我們能夠恩著.
認順你為我們的救主.
生命的主的時候.
我們就可以和父神.
恢復那個美好的關係.
成為上帝你的兒女.
成為你所喜悅的子民.
我們今天早上聚集在你的面前.
為的是要尊崇高舉你的名.
盼望聖靈啊,你幫助我們每一個.
讓我們真的帶著一份喜樂的心.
帶著一份期盼的心.
來敬拜你,尊崇你.
並且讓我們將你這份大好信息.
一直向人播送開去.

$^{81}$讓人真真正正體會到.
上帝是你各樣的福氣.
我們這樣的禱告,交托,仰望.
是奉基督耶穌你得勝的名字來祈求.
阿們.
各位請坐.
主耶穌視神榮耀所發的光輝.
視神本體的真相.
聖誕節不但紀念.
道成育身耶穌基督的誕生.
也是宣告神榮耀光輝的降臨.
當天使向牧羊人宣講基督誕生的時候.
天君和鐵鉛歌唱贊美.
在至高之處榮耀歸於神.
在地上平安歸於他所喜悅的人.
天使歌唱在高天.
微妙歌聲綿綿圓.
是週三的花會笑.
和年輕媳婦老娘.
威 英 賀 耶 基督 誕 生.
牧人有何大事.
快樂歌聲快樂甜.
有何氣聲甘心甜.
世傳歌聲美難言.
威 英 賀 耶 基督 誕 生.
天公通光光明燈.
天使所唱是天神.
前來跪拜高聲唱.
威 英 賀 耶 基督 誕 生.
主基督出生的記號.
就是一個嬰孩包著布.
我在馬槽里.
牧羊人就是按著這個記號.
去尋覓嬰孩耶穌.
難以想象.
榮耀救主的誕生.
竟然如此簡樸和卑微.
不過正因為這樣.
一個本來冷冰冰的馬槽.
就頓然充滿安寧.

$^{121}$愛 溫暖和亮光.
甚至成為一個實現.
神拯救人類應許的神聖之所.
我主耶穌.
多安寧神聖嬰孩.
滿窗安安出遊宇.
看東風天外曠野.
聖誕精神充衍.
笑聲應若地得收.
方言我讓你開闊.
天使歌唱過詩諾.
掀滿月光的夜空.
請指梅.
他帶來化燈恩賜.
聖詔天賦心所願.
笑聲哭 恩怨歇息.
上帝應許該赦免.
笑聲應 耶聖賞賜.
燦燒我滿腹怨苦.
不笑我難得忍辱.
卻無存我沒點活.
他願意 永要得願.
盡周多 所所滿足.
上帝愛 使我再難.
做了只有何做.
笑聲應 耶你聲應.
包圍倦世俗慶豐.
音樂取倦大亂河.
你是主基督耶穌.
多歡迎 神聖天下.
拜託安安所有人.
感動方圓 歡歡迎.
聖誕清晨終宴宴.
因我們神憐憫的心腸.
他按自己鎖定的日期.
使他兒子耶穌基督.
為我們降生 受死 復活.
叫清晨的日光從天臨到我們.
要照亮坐在黑暗中.
死忍淚的人.

$^{161}$把我們的腳引到和平的路上.
在這個充滿憂鬱的時代.
願上主的真光.
溫暖我們眾信徒的心.
謹記我們今天所看見.
所聽見的一切.
原諒盼望.
繼續以歡欣同敬拜基督君王.
齊來上主聖祠.
拜請大家歡呼.
讓人來祝福.
來到伴侶們.
慶迎我出宗.
聖愛天使光明.
齊來歡欣同敬拜.
齊來歡欣同敬拜主的燈.
天使一同歡慨.
恭喜到場中.
我聖羊羊.
以煙煙香贊佩禰.
天生有永遠.
唯於天上神神.
齊來歡欣同敬拜.
齊來歡欣同敬拜主的燈.
齊也,感知耶穌.
各聖德外九屬.
傳耶穌來人.
至永守之所.
旗在主的燈.
祝各聖愛便能.
齊來歡欣同敬拜.
齊來歡欣同敬拜主的燈.
讓我們繼續聆聽一段.
關於耶穌基督降生的經文.
記載在《路加福音》第二章.
八至二十節.
經文裡面是這樣說的.
在伯利恆之野地裡.
有牧羊的人.
夜間安著奸刺看守羊群.

$^{201}$有主的使者站在他們旁邊.
主的榮光四面照著他們.
牧羊的人就甚懼怕.
那天使對他們說.
不要懼怕.
我報給你們大喜的信息.
是關乎萬民的.
因今天在大位的城裡.
為你們生了救主.
就是主基督.
你們要看見一個英雄.
包著布.
我在馬槽里.
那就是基督了.
忽然有一大隊天兵.
和那天使贊美神說.
在至高之處.
榮耀歸於神.
在地上平安.
歸於他所喜悅的人.
眾天使離開他們.
升天去了.
牧羊的人彼此說.
我們往伯利行去.
看看所成的事.
就是主所指示我們的.
他們急忙去了.
就尋見瑪利亞和藥室.
又有那英孩窩在馬槽里.
既然看見.
就把天使論這英孩的話傳開了.
凡聽見的.
就差以牧羊之人對他們所說的話.
瑪利亞卻把這一切的事.
存在心裡反復思想.
牧羊的人回去了.
因所聽見.
所看見的一切事.
正如天使向他們所說的.
就歸榮耀與神贊美他.

$^{241}$接下來有陳明荃牧師為我們宣講.
剛才所讀的那段經文.
講題是.
因他降生.
我們得平安.
請姐妹在這裡預祝大家聖誕快樂.
因為耶穌降生.
我們就得享屬天而來的平安.
在這個世界裡面.
只要我們深信主耶穌基督.
成為我們生命的救主.
這份平安就很實在地.
充滿在每一個基督徒的心靈裡面.
今年雖然因為疫情的緣故.
我們沒辦法真正熱鬧地.
坐在教會裡面一起慶祝.
不過我們心裡知道.
真正令我們快樂的.
不是一群人純粹熱鬧地坐在教會裡面.
而是每一個人.
因為生命裡面.
因著基督耶穌而帶來的轉變.
我們得享永恆的福樂.
永生就在我們信主的那一刻開始.
我們的喜樂就由我們信主的那一刻開始.
所以我們的心靈裡面.
永流不息的.
可以一起來慶祝耶穌的誕生.
因為我們知道.
我們的生命是屬於上帝的.
這個信息其實我想大家都不會陌生的.
不過在聖經裡面.
有時候他怎麼說這個大喜的信息.
或者這個平安的福音.
上帝用什麼人來宣講的時候.
其實是很值得我們再三回味.
究竟上帝用什麼方法來幫助我們.
明白上帝救贖的旨意.
如果大家有看電影.
我很喜歡看周星馳的電影.

$^{281}$其中有一部電影叫喜劇之王.
周星馳在那部電影裡面.
那部電影叫喜劇之王.
但整部電影圍繞著喜劇之王.
成為這部電影裡面的主角.
由一個喜劇之王在他的電影生涯裡面.
沒有什麼人重視的.
甚至被人嫌棄的.
但在這部電影裡面.
卻是每一個觀眾都看著這個喜劇之王.
他的心路歷程.
以及將整部電影的焦點放在他的身上.
大家不要誤會.
我不是說電影.
不過同樣的手法.
其實在《路加福音》裡面.
說到耶穌基督的誕生.
某程度上都是用一群Cataphra來去展明出來.
今天我們其中看一群.
連名字都沒有記載在聖經裡面的.
一群牧羊人.
某程度上他們是在上帝計劃.
或者在很多當代社會裡面.
是一群Cataphra.
但這群Cataphra卻是承載著.
或者上帝藉著這群不起眼的小人物.
將上帝一群很重要的信息.
告訴全世界知道.
今天我們希望從這群人身上.
看到榮耀的福音.
那個平安福音是一個怎樣的福音.
在我們仔細看聖經前.
我們一起祈禱.
求上帝幫助我們.
我們同心祈禱.
天父啊你的愛真是美好.
你的救贖計劃真是超乎我們所想象.
我們今天我們等姐妹.
我們一起聚集.
我們一起聆聽.

$^{321}$不單只是一個耳熟能詳的福音.
更加是一個對我們生命裡面.
有很深刻提醒的福音.
求你的愛激勵我們.
幫助我們.
讓我們翻開聖經的時候.
我們明白你的心意.
在這個世代裡面.
你讓我們能夠按著你的心意.
來行事為人.
幫助此時此刻我們聆聽的時候.
讓孫講的講得清楚.
讓聆聽的每一個等姐妹聽得明白.
我們恭敬將來的時間交代給你.
願你賜福.
牽上祈禱奉靠基督耶穌.
得勝的聖明祈求.
Amen.
正如我剛才所講.
在奴家福音裡面.
其實他一連串記載.
特別是去到牧羊人的段落.
他在講上帝席著一群.
在野地或者夜晚裡面.
按著羊群的牧羊人來顯現.
天使向他們出一個大喜的信息.
值得留意的就是牧羊人在耶穌時代.
是一群什麼人呢.
我們如果看聖經.
我們經常覺得牧羊人就是.
在舊約裡面好像在講上帝和羊之間.
是一種很親密的關係.
一個很正面的身份.
是一個很好的職業.
不過很多聖經學者告訴我們.
其實在耶穌時代的牧羊人.
其實不是一份很高尚的職業.
至少不會每個人都爭著.
我一定要做牧羊人.
成為我的志願職業.

$^{361}$當時牧羊是一些很粗重的功夫.
甚至當人家睡覺的時候.
有時候要晚上來看守羊群.
為什麼晚上要看守羊群呢.
在伯尼衡的野地.
或者在這裡的野地裡面.
為什麼要無端端有一個.
這麼奇特的出現呢.
其實在很多的文獻.
我們知道了.
其實這群夜更裡面.
看守羊群的牧羊人.
其實看守的羊群.
就是在耶路撒冷每一年.
或者每一年不停獻祭的那群.
為著將要獻上聖殿.
獻祭的羊羔.
要看守著那些羊群.
因為晚上裡面可能會有其他的野獸出現.
進入羊籃裡面可能會傷害.
甚至會殺害那些.
吃了那些羊群.
所以牧羊人晚上要看守著.
將要在聖殿裡面獻祭的羊群.
讓這些羊群可以安全的.
準備來獻祭.
當你知道這樣的時候.
這群牧羊人.
那群夜班.
其實某程度上好像一群漢奸一樣.
他看著那些羊群.
不會受害.
不會被侵害.
一個這樣的角色.
其實在社會上.
不是一個很標秉.
或者一群很受人尊重的角色.
甚至有些文獻.
講到有一些牧羊人因為夜班.
不是一個很受尊重的行業.

$^{401}$在當代社會.
那些人看不起的.
甚至有些人覺得.
那些沒有什麼技能.
身份裡面沒有工作.
他肯付出努力就能做到的一份工作.
甚至有些人說.
這些三五成群的人.
可能是很品流複雜的等等.
在聖經裡面沒有特別.
講參牧羊人的職業.
不過在路加時代.
當他寫牧羊人.
去夜班裡面看守這些羊群的時候.
大概在當代的讀者裡面.
就知道.
這些是努苦大眾.
這些不是一些很出色的.
很有說服力的.
甚至令到社會上.
每個人都景仰的一個職業或者身份.
如果我們試想一想.
如果我們有一個很重要的信息.
如果一個關乎全世界所有萬族萬民的信息.
如果在你印象中裡面.
如果你收到這樣的情報.
你第一個要通知的是什麼人呢?.
就這個問題我問過我女兒.
我女兒問她.
如果你收到這條.
有一個關於全世界生與死的一個信息.
你第一個告訴什麼人?.
我女兒跟我說.
我第一個告訴國家元首領袖.
因為他們有權力.
所以他們知道這件事的時候.
他們就可以將國家或者所有東西調動.
希望幫助全世界人.
渡過一個這樣的難關.
不過在聖經裡面.

$^{441}$似乎他要逆轉了.
他顛覆了我們想事情的思維.
上帝將一個關乎萬民.
甚至生與死.
一個拯救萬民的信息.
他將這個信息第一個顯現.
或者通報的.
天使要用很大的場面.
要講的那群人.
不是一群KOL.
或者一群很有影響力的人.
反而是在社會上.
是沒有什麼影響力.
不是很顯眼的一群牧人人.
上帝親自選擇這群人.
成為他接收信息的.
第一個去領受耶穌降生的消息的.
這一群奴婦大眾.
一群不起眼的小人物.
上帝就要用他們.
將一個大氣的信息.
告訴他們知道.
第二件事值得注意的.
聖經裡面他特別講到.
在大魏的城裡面.
為你們生了救主.
就是主基督.
在和合本裡面.
我們看到字眼沒有譯錯.
不過他這樣讀上去好像很順.
為你們生了救主.
就是主基督.
我們好像很順地讀了.
看不到那個特性.
其實這句話裡面.
是一個很重要的表達.
一個很重要的人物降世.
就是救主.
主基督.
原文裡面.

$^{481}$其實某程度上要有少少停頓.
你才能讀到原來.
裡面是關乎一個全世界的信息.
一個拯救萬人.
全世界人的一個拯救者.
一個好像王一樣.
掌管萬有.
滿有權力的主.
降到這個世界裡面.
更加重要的是基督.
就是以色列整個民族.
引頸以待.
那個受高者.
他民族裡面.
整個民族在等待著的一個人.
所以說一個救主.
主和那個受高者.
今天在伯利恆.
在大衛的城裡面出生了.
所以在這個說法裡面.
天使在說一個重要的人物出現.
這個身份這麼顯黑的.
扭轉世界整個命運的耶穌基督.
他的降生是一個很重要的信息.
繼續天使在說.
他在哪裡出生呢.
你要看到那些記號.
記號是什麼.
看到那個小孩.
包著布.
我在馬槽裡面.
這個就是記號了.
看到那個小孩.
包著布.
其實每個初生嬰兒都是這樣的.
路加福音裡面多次反而.
一而再再而三地強調.
馬槽這個字.
一個很重要的孩子.
包著布.

$^{521}$而且是一個嬰孩.
但是卻是在那個馬槽裡面降生.
如果你這樣看.
你可能到今天都不知道那個張力.
如果我說.
這個世界裡面將會有一個很重要的人物出生.
你當是一個全世界的.
或者一個國家裡面的領袖.
這個領袖將要出生了.
然後那些人問.
他在哪裡出生呢.
哪家私家醫院.
我們第一件事就會問.
哪家醫院.
走到那裡第一時間報道他.
天使說不是.
是在旺角某個天橋底裡面.
找個紅A膠箱.
栽著那個小孩.
就是未來的那個重要人物.
你聽下去.
這樣都可以.
一般小孩都不會找個紅A膠箱.
在天橋底裡面出生.
為什麼會是一個這麼重要的人物.
用一個這樣的方式.
說白一點.
就好像一個氣嬰那樣.
在那裡出生呢.
當然耶穌的降生.
不是我們心目中所想的.
被遺棄的嬰孩一樣.
只是客店沒有位置.
當時要報名上冊.
在阿古士都.
要將人數來到數點的時候.
因為客店沒有位置的時候.
而瑪利亞就臨盤了.
在一個這麼裂缺的時間裡面.
唯有就在馬槽裡面.

$^{561}$擔下這個主耶穌.
在一個很不容易的空間.
在一個很緊迫的時間裡面.
耶穌就是用一種最卑微的方式.
用一種難以想象.
不可思議的方式.
降生在這個世界裡面.
《路加福音》.
他記載著耶穌的降生.
用一種超乎我們想象.
顛覆我們對於一些很偉大的人的.
降生的一些看法.
反而是在《路加》的記載.
或者聖經的記載裡面.
特別強調.
反而一個這麼重要的基督耶穌.
用最卑微.
最惹怒.
甚至是一個最不起眼的方式.
來降生在這個世界裡面.
更加直著一群社會上.
最不起眼的一群牧羊人.
一群身份不是很顯黑的那些人.
上帝將這個信息告訴他們.
整個過程就充滿了張力.
整個過程.
大概對於讀聖經的人來說.
如果我們真的很認真去看的時候.
上帝的救贖計劃竟然不是.
用強人姿態.
用一種很輝煌.
很誇張.
很有場面的那種形式.
來呈現出來.
甚至是比當代社會低一點點.
甚至是一個很卑微的形式.
將這個救贖的計劃.
來直著這一群人.
直著基督在馬槽的降生.
來呈現出來.

$^{601}$在這裡我們稍微思想一下.
我們今天想的基督教信仰.
是一個怎樣的基督教信仰呢.
會不會有時候我們覺得.
要很大的場面.
要很有榮耀的感覺.
甚至要很多東西很華貴.
才能夠配得上這個信仰.
我想那些東西都不是錯的.
那些東西不是差的.
不過有時候在聖經裡面的信息.
和我們心中所想的東西.
可能會製造一個很大的距離.
甚至有時候我們會覺得.
因為我們呈現的面貌.
可能和我們最初領受.
聖經裡面所記載的福音.
可能背道而馳.
我記得我小時候.
我小時候就開始回到禮拜堂裡面.
特別是聖誕節的時候會聚餐.
甚至會參與聚會.
我小時候我跟我姑婆回教會.
但是小時候家裡不算很有錢.
也不算很窮.
但是我一回到教會.
我有一種感覺.
就是覺得自己好像比下去.
有一種不是很屬於自己的感覺.
為什麼呢.
即使我穿得最漂亮.
其實我穿得最漂亮的都是那些衣服.
家裡沒有那麼殘.
以前沒有特別漂亮的去街衣服.
只有一件衣服沒有那麼殘.
去街就穿那些.
但是去到教會裡面.
看到自己同年紀的那些.
身光頸靚的.
遮住頭.

$^{641}$另外穿的衣服顏色很漂亮.
是名牌子.
不知道你們知不知道.
我小時候那些美美百貨買的.
不過你們可能不知道.
可能在在的老兵知道.
那些衣服在美美百貨買.
我的那些就是.
以前在兜兜裡面.
找一件自己適合的.
那些是在屋村衣服.
一回到教會有種比下去的感覺.
甚至聽那些和自己年紀差不多的那些.
說話就溫文有禮.
走到鋼琴裡面又能彈一首.
自己就什麼都不懂.
感覺上就是覺得.
我不是屬於那個群體.
我覺得我們這些低下階層.
或者我們不是在教會長大的一些人.
根本上和教會好像不關自己事.
所以每年回到教會.
就好像帶著有些見識.
去看一看.
或者學習的形式.
來參與教會的聚會.
甚至有一個很尷尬的例子就是.
去到聖誕節最後的晚會.
到最後牧師會派那些聖誕禮物.
每個小朋友拿一份.
自己貪心那時候.
想多拿一份給自己的妹妹.
還要說謊騙牧師.
我說我沒拿的.
被牧師在聖誕節的場合裡面.
聚會當中大聲罵.
你說謊.
明明看到你拿著聖誕禮物.
那次感覺上自己穿得不光彩.
還要好像貪心那樣.

$^{681}$拿多一份禮物.
雖然我最後拿著兩份禮物回家.
但我心裡面感覺上.
教會是不是屬於我的呢.
耶穌是不是屬於出身較好.
較有錢的人.
或者那個福音是為他們而設的.
對於我們這些傳統的.
或者不是很顯眼的中國人.
這個信仰不是很關自己事.
不過在路加福音裡面.
他說到不是這樣的.
上帝用那些社會上不顯眼的人.
第一個就是要通報給他們.
天使要向他們顯現.
將基督耶穌大喜的信息.
直著他們來宣講.
甚至耶穌自己也是用最卑微的方式.
不是最華麗.
不是最顯眼.
很令人羨慕的方式來降生在這個世界裡面.
卻是用一個最卑微.
最平易近人的方式.
正悄悄地將他救贖的計劃.
帶到這個世界裡面.
上主要用最卑微的方式.
來成就這個最偉大的救贖計劃.
他使用最卑微的人.
來傳講.
甚至領受上帝而來.
一個美好的信息.
經文再繼續來講.
在第二章十四節.
除了天使將信息講完之後.
一大群的天君天使.
在天上很大的場面.
地上可能很冷清.
但天上裡面有榮光.
有一群的天使就在那裡頌唱.
在至高之處.

$^{721}$榮耀歸於神.
在地上平安.
歸於他所喜悅的人.
這是天啟超自然的顯現.
牧羊人看到一個很華貴很誇張的天上景象.
他們不得不信.
不過這句歌頌裡面.
後半句是很有意思的.
我一路看的時候.
我覺得為什麼這麼特別呢.
我們平時看在地上平安歸於他所喜悅的人.
我們就覺得讀上去很順利.
很開心.
不過當你仔細看的時候.
平安是給什麼人呢.
是給上帝所喜悅的人.
有些和合本或者有些版本.
會有括弧寫著有些細字.
譯法就是在地上平安.
喜悅歸於人.
沒有特別說將平安歸於喜悅的人.
喜悅的事成為一個接受詞.
或者是一個主詞.
不過歷代讀聖經和合本的翻譯是準確的.
或者是更加經得起神學的考驗.
在這裡特別說到.
平安不是順手捏來的平安.
不是一個耶穌出世然後全世界都平安.
不是這個意思.
這個平安是有具體的對象.
這個平安是要特別給一些上帝所喜悅的人.
他才得著這份平安.
如果你看到這樣的意思的時候.
你就要繼續問.
哪些人才是上帝所喜悅的人呢.
哪些人呢.
剛才說了.
是不是身份很好.
或者讀很多書.
甚至是教會子弟.

$^{761}$那些就是上帝所喜悅的人呢.
聖經裡面出現喜悅這個字不是很多次.
在路加福音裡面另外再出現喜悅這個字.
在第十章第21至第22節.
那段經文是說到耶穌差派一些門徒.
七十個門徒出去.
兩個兩個來傳福音.
看到平安的福音傳給不同的人的時候.
看到有魔鬼在天上墜落.
是這樣的背景.
然後耶穌在那段時間.
看到門徒的果效.
耶穌被聖靈感動而快樂.
他就這麼說.
因為你將這些事向聰明通達人藏起來.
向英雄就顯出來.
父啊是的因為你的美意本是如此.
一切所有的都是我父交付我的.
除了父沒有人知道子是誰.
除了子和子所願意子是的.
沒有人知道父是誰.
在這段經文裡面看不到喜悅的字.
因為他用另外一個字眼譯的就是美意.
上帝的美意就是這樣.
上帝的美意或者上帝的喜悅是什麼呢.
Eudokias這個字.
他說到上帝按著他的心意.
他特別說一些以為自己很厲害的人.
上帝要將他的旨意.
將他的喜悅隱藏.
對於那些像是英雄的心那麼單純的人.
對上帝的心敞開的那些人.
上帝就喜悅那些人.
向他們顯明上帝.
聖靈而來的那種能力和聖靈而來.
引導人知道耶穌的身份是誰.
所以在這裡面.
講到那個喜悅或者上帝的美意.
其實是按著上帝的心意的.
我們大概沒有辦法知道上帝的心意是誰.

$^{801}$你說牧師你沒回答過吧.
不是.
講上帝喜悅的.
第一件事我可以排除的就是.
不是人的努力.
不是人的功德.
界定了我自己是否被上帝所喜悅.
甚至我沒有罪.
我的生活.
我過去是怎樣.
都不界定上帝喜悅我與否.
上帝的喜悅是按著上帝的心意.
上帝要將恩典臨到這個人身上的時候.
上帝就要將他自己要作成的功.
無論怎樣都好.
他都要將一個不可抗拒的恩典.
放在一個可能自己都看不起的人身上.
上帝要將他的旨意.
定義要哪些人得救.
他喜悅哪些人得著這份平安.
上帝就用他自己大能大力.
按著他的主權來選擇不同的人來得著福音.
這個說法對於以色列人來說是一個很大的顛覆.
為什麼呢.
因為對於以色列人或者對於整個民族來說.
他覺得自己整個民族.
以色列民族是the chosen one.
在世界這麼多邦國裡面.
他們是凸顯出來的.
耶和華上帝所選擇的民族.
其他民族沒有被選招.
是他們自己民族才被選招的一個特別的群體.
所以他們某程度上都有一種屬靈的驕傲.
他覺得自己是被選擇的一群人.
不過去到奴家福音.
隨著基督耶穌的降生.
他說的那個墳墓.
不只是說以色列群體裡面的這一群人.
他的墳墓.
就是那個計數的算式的墳墓.

$^{841}$是全世界所有的萬國萬民.
都是上帝喜悅的範圍.
不過真正得著平安.
真正生命裡面得著基督耶穌的救恩.
是上帝在這一群人裡面選擇萬國萬邦.
甚至乎跨越了國界民族.
他都要在不同的民族萬邦裡面.
選擇一些人得著基督耶穌.
帶來生命的平安.
繼續說平安.
是說上帝喜悅的人得著.
那些平安是一種什麼樣的平安.
我們今天說平安這個字是值得解釋的.
因為我們今天是白答的.
我們說的那些平安.
有時候覺得心裡不舒服.
就是在說平安.
平安這個字.
如果你用聖經新約裡面的上文下理.
如果在《路加福音》第一章77-78節.
早幾十幾節的經文.
撒迦利亞的祈禱.
他在說的平安.
是在說慰助施浸約翰的出世祈禱.
他說這個兒子要帶來平安的就是.
讓人不會在黑暗當中.
過他繼續而來的黑暗生活.
反而是就光進入一種離開黑暗進入光明.
走在平安路上裡面的那種心靈的那種狀態.
簡單來說這個平安第一個的詮釋.
或者第一個說的具體內涵是道德性的平安.
上帝直著他的使者直著基督耶穌.
指明這個世界的人應該要走一條怎樣的正路.
走在正路裡面的時候.
我們生命就不會擔驚受怕.
不需要擔心我們自己怎麼走錯.
按著聖經來做人的時候.
無論路怎麼難走.
我們都可以挺起胸膛.
我們都可以松容不逼.

$^{881}$有一份安全的踏實感來走面前的路.
因為我們知道這條路是上帝告訴我們應該走的正路.
所以那個平安第一個具體的內涵是說.
走在正路裡面夜半敲門也不驚.
那種心空磊落做人光明磊落挺起胸膛.
道德對得住上帝.
知道上帝怎麼教我們的那條道德的路.
那是第一個平安.
第二種平安的內涵除了道德之外.
就是一種關係性的平安.
《爾弗所書》第二章13至14節裡面.
他說到耶穌基督.
祂自己成為我們的和平和和本益.
其實同樣是平安這個字.
耶穌自己是我們的平安.
使雙方合而為一.
拆毀了中間隔絕的牆.
而且以自己的身體終止了冤仇.
那個平安是說關係上不再對立.
不再衝突.
犯罪的人本身和上帝是一種對立的關係.
每一個人都是可怒之子.
上帝要傾倒憤怒在我們每一個全世界所有的罪人當中.
如果按照我們自己背後的記錄.
我們的檔案.
午夜夢回你回想你自己人生的片段.
大部分你都是要搖頭嘆息.
不堪回望.
我們生命都是滿布很多的罪.
我們都是一個可怒之子.
我們和上帝是對立的衝突的.
我們沒有辦法和上帝建立一種和諧的關係.
不過基督耶穌的降生.
祂要讓每一個被選照.
心靈裡面領受福音的人.
祂重新稱基督耶穌為我們生命的主.
我們的罪就歸到祂的身上.
而因著基督耶穌的救贖.
我們再和神之間不再是一種敵對.
或者一種可怒之子的關係.

$^{921}$我們成為可喜悅的兒子.
而塞爾書62章1至5節說到.
好像新郎哥喜悅新娘.
好像上帝喜悅子女一樣.
本來由黑人憎恨上帝憎恨的人.
變成一個討上帝喜悅的一群人.
這份平安就是說我們不需要擔驚受怕.
不需要覺得我們做少了事上帝不喜歡我們.
好像其他宗教那樣.
怕做少了功德.
做少了事上帝隨時會有和緩臨到自己身上.
不是的.
上帝說借著基督耶穌.
借著這個救贖.
我們可以和祂有一個和諧美好的關係.
祂喜悅我們.
祂愛我們.
祂愛我們到底.
耶穌來到這個世界.
就要將人與神之間的那一幢牆拆毀.
借著這一幢牆拆毀.
我們人與人之間也可以拆毀很多的冤仇.
很多的仇恨我們都可以拆毀.
所以這個平安是一個關係性的平安.
是一個神與人之間的關係.
是一個人與人之間的關係的平安.
基督耶穌為我們帶來這樣的轉機.
所以天使他唱這一句話.
不是純粹是在說鏗鏡有力很大場面.
他唱出來的那個信息.
在天使天君裡面要歌頌或者要講的.
是講到耶穌基督.
他就將這個世界那種痛苦.
將世界那種不和諧扭轉.
用最卑微的方式.
用最不起眼的人.
來宣講一個最偉大平安的信息.
上帝喜悅人.
上帝要將這一份行在正路的平安.
將一份和神和諧的平安.

$^{961}$重新賜給你和我.
可能你會問.
那我是不是上帝所選擇的呢?.
我是不是上帝所喜悅的呢?.
在很多的神學家.
我們都沒有辦法給一個終極的答案.
不過在宗教改革或者是.
加爾文的神學裡面講到.
這個選照的預定論.
是要告訴每一個.
當我們心裡面有聖靈感動.
稱基督耶穌為主.
我們決意跟從他的這一群人.
當我們生命裡面領受過天恩的滋味.
這一群基督徒.
他要告訴我們一個確確.
我們所信的神是真實的.
上帝不會那麼輕易.
上帝不是喜歡要扔下你放棄你.
當上帝要將一個救恩.
臨在你的身上的時候.
你知道上帝是選照你.
上帝喜悅你.
上帝在你人生無論什麼境遇都好.
都要將那種不可抗拒的恩典.
無條件的愛.
臨到你的身上.
當我們人生處於可能不同的制約.
患病.
制約不好.
失業.
經濟不好.
甚至犯罪.
這些一切都不會難阻了上帝對我們的愛.
只要我們堅信耶穌對我們的愛.
我們稱他為我們生命唯一的主.
我們不放棄自己.
我們不放棄這個救恩.
上帝是不會放棄我們的.
所以在這裡上帝的愛.

$^{1001}$就是一個超越我們所想所求.
當坊間說你乖我就愛你.
你做得好我就對你好的時候.
上帝說無論你怎樣都好.
我不撇下你成為孤兒.
我喜悅你.
我用永遠的慈愛愛你.
借著基督耶穌的拯救.
你永遠得到這份永恆的福氣.
牧羊人聽到這個消息.
聽到天使這個歌聲.
他們有什麼回應呢?.
2:15-20.
眾天使離開他們升天去了.
牧羊人彼此說.
我們往伯利行去.
看看所成的事.
就是主所指示我們的.
他們急忙去了.
就尋見瑪利亞和藥室.
又有嬰孩餓在馬槽里.
既然看見就把天使論者.
孩子的話傳開了.
凡聽見的就詫異牧羊之人.
對他們所說的話.
瑪利亞卻把這一切的事.
傳在心裡反復思想.
牧羊的人回去了.
因聽見所聽見的一切事.
正如天使向他們所說的.
就歸榮耀與神贊美他.
在牧羊人聽見.
一個如此重要.
關乎萬民.
如此美好的信息的時候.
天使升天去了.
他們的回應是急忙地.
上去伯利行.
聖經特別要講到.
那個速度讓你知道.

$^{1041}$他們不是聽完之後.
都想一想.
反復思量不是牧羊人.
是瑪利亞.
牧羊人聽到這個信息的時候.
他們心裡就知道.
這是上帝所指示他們的.
所以他們就急急腳地.
甚至羊群都被扔下了.
他就急忙地走去.
找瑪利亞和藥室.
和這個英孩的耶穌.
將天使傳給他們的話.
再一次告訴瑪利亞藥室.
他們這個家庭知道.
將這個印證.
重新告訴.
似乎很哀怨.
生耶穌的時候.
某程度上很淒涼.
一個這樣的情景.
這群牧羊人又好像.
在兼顧家庭.
是一個很卑微降生的家庭.
有一群陌生人.
從野地裡面趕過來.
既是慶祝.
同時也印證著上帝.
一個很重要很重要的任務.
交付在孩子身上.
聖經裡面特別說到.
他們三個同作這群牧羊人.
他們衝出去.
去看看所成就的事.
他們不單止將這件事.
告訴藥室瑪利亞.
他們再將他們所經歷的.
所見所聞.
告訴他們身邊的人知道.
他將這一切的事傳揚出去.

$^{1081}$最後他們將榮耀歸給神.
整個的回應是一個很單純.
很直接簡單的回應.
福音可以說得很複雜.
甚至可以說得.
有很多養分和智慧.
不過回應是很簡單的.
將你所聽到的.
將你所領受到的.
將你所經歷到的.
傳揚出去.
立即行動.
不要拖.
不要純粹想.
很多人以為.
想福音.
或者在回答.
福音就會傳開.
我們當然不是傳一個.
我們的福音.
不過任何思考.
如果帶來不了行動.
改變不了這個世界.
真的.
如果我們只是.
永遠都是坐著.
永遠都是想.
等.
福音或者這個世界.
根本就不會.
有任何的改變.
在《奴家福音》裡面說到.
當人領受了.
聽見了的時候.
他們就.
就好像社會上最沒有條件的人.
但他們卻是最.
直接回應上帝的一班人.
在這裡.
我特別對黃澤民的弟兄姊妹說.

$^{1121}$我們很喜歡聽的.
我們很喜歡聽一些很厲害的論調.
我們用了很多心血來預備.
希望將一個好的信息告訴你們.
不過聽了這些信息的時候.
你不要純粹放在自己的口袋里.
滿足自己.
黃澤民的弟兄姊妹很多時候.
未必是黃澤民的香港人.
都是這樣的.
我們喜歡聽.
聽一些厲害人所說的厲害的東西.
聽完之後就覺得.
滿足了自己.
不過最重要的就是.
信息傳給你之後.
其實要借著你的生命來展開.
如果我收藏在心底里.
我們沒有辦法.
繼續進一步傳揚的時候.
福音就會變成一個.
拿著酒杯底的福音.
就是一班人.
閒時拿出來.
回答一下.
感覺良好.
好像談哲學那樣.
說一下人生哲理.
放完酒杯底之後.
然後又再回到你的生活當中.
不過福音從來不是一個.
拿著酒杯底的福音.
福音是當你領受了的時候.
上帝都在催促我們去回應.
傳揚出去.
不是純粹滿足自己那種智慧.
或者感覺良好的一種信息.
這麼簡單.
而是要將這個信息.
經歷歷代.

$^{1161}$一代又一代傳揚出去.
唯有這樣的世界.
才一直借著我們的生命.
來傳揚出去.
聖經裡面繼續說.
特別是說.
可能在世界裡面.
有很多人覺得自己不起眼.
覺得在世界裡面.
其實有時候我都覺得自己是.
Catafalque(貓頭鷹).
不過.
上帝的救贖計劃.
他特別要用那些Catafalque(貓頭鷹).
那些不起眼的人.
反而是覺得自己很有能力.
很標秉的那些人.
如果他的智慧是難阻了他.
來到他有回應的行動.
大概那個福音的棒.
不是交給那些人的.
在哥林多前書第一章26至27節裡面.
他這麼說.
弟兄們.
可見你們蒙照的.
按著肉體有智慧的不多.
有能力的不多.
有尊貴的不多.
神卻揀選了世上愚拙的.
叫有智慧的那些羞愧.
又揀選世上軟弱的.
叫那強壯的羞愧.
聖經直著一群軟弱的.
叫一群卑微的.
叫一群某程度愚拙的一群人.
竟然去傳揚一個.
沒有上帝救贖智慧.
那麼偉大的救贖福音.
牧羊人就是一個最佳的示範.
當上帝在昔日裡面.

$^{1201}$他都用牧羊人.
沒理由不用你.
可能今天我們將自己給下去.
覺得自己不是在這個世界舞台裡面.
或者生活裡面.
成為最重要的焦點所在.
不過上帝說.
我要用你.
我喜悅你.
我要直著你將這個平安的福音.
直著一些不起眼的人.
傳揚開去.
早前有機會和陳兆徹校長.
有一個直話聊天.
他跟我說了一番話很有意思.
他說今天做基督徒.
千萬不要帶著打機的思維.
來過我的人生.
我說什麼叫打機的思維.
他說打機的人永遠.
都是用一個主角的角度.
來看那個遊戲.
所有的挑戰困難.
就好像打機的人是主角.
所以失敗了.
就重新來過.
他說這個打機的思維累死人.
因為現實不是這樣的.
第一不是說失敗了.
遊戲結束了就重新來過.
第二件事.
這個世界根本就沒有那麼多主角.
反而聖經裡面.
拉回我們來看.
這個世界真正的主角是基督耶穌.
我們只不過在.
基督耶穌身邊裡面.
那些綠葉.
每個人的氣氛有不同.
有些人會重一點.

$^{1241}$有些人會輕一點.
有些人可能做這些.
有些人可能做過一些.
但是.
無論我們處於什麼位置.
我們好好地.
演繹我們的角色.
上帝不要介意我們的卑微.
我們也不要介意自己.
我們出生是怎樣.
在這個聖誕節裡面.
運用你的身份.
演繹好你的角色.
上帝喜悅你.
上帝要直著你.
將福音傳開.
求主繼續給你幫助.
\newpage



\section{ }
\label{sec:lmnMx2HA46Q}
\textbf{2020年12月24日｜擁抱幸福分享晚會:平安夜晚會 ｜旺角浸信會聖樂團隊}
\newline
\newline
連結: \href{https://youtube.com/watch?v=lmnMx2HA46Q}{\texttt{https://youtube.com/watch?v=lmnMx2HA46Q}} ~~~~ 語音日期: 2020-12-25
\newline
\newline
\hyperref[sec:TPSrTFQecuk]{\small{< < < PREV SERMON < < <}}
~
\hyperref[sec:index_chronic]{\small{[返順時目]}}
~
\hyperref[sec:DokqOqhZ3o0]{\small{> > > NEXT SERMON > > >}}
\newline
\newline
$^{1}$剛才那首詩歌很有意思.
當詩歌講到.
其實是在約翰福音14章第27節.
經文裡面這樣說.
「我留下平安給你們.
我把我的平安賜給你們.
我所賜給你們的.
不像世人所賜的.
你們心裡不要憂愁.
也不要膽怯」.
這一句說話在約翰福音記載.
其實是耶穌所說的一句說話.
祂告訴我們這個世界的平安有兩種.
一種就是我們所說的.
是世界給你的平安.
第二種就是耶穌所賜給你的平安.
什麼是世界所賜的平安呢?.
世界的平安就是告訴我們.
會因著環境的改變.
而對我們的心態有影響.
環境會帶給你的.
是一個平安還是不平安的訊息.
平安是因著環境而影響.
例如什麼呢?.
我覺得最常見的就是.
在我們的日常生活裡面.
數字會令到我們的平安有起有跌.
例如股市.
股市恆生指數升或跌.
或者你買的股票升或跌.
會影響你的平安.
影響你心裡的安穩.
又或者健康.
不知道大家會不會定期檢查自己的身體呢?.
你的膽固醇指數.
血糖,血壓.
這些指數都會影響我們的心情.
又或者近來我們會去驗身冠肺炎.
究竟我們是呈陰性還是陽性.
這些環境都會影響我們的心態.

$^{41}$或者我們做爸爸媽媽.
我們會很緊張.
子女的分數.
她的考試成績.
她的排行榜排第幾.
這些都會影響我們的心.
我們的心會不會因著外圍環境的數字.
而令到我們有起有跌.
反而這些外在的環境.
就會令到我們的心.
正如聖經裡面所說.
會令到我們很擔心.
令到我們很擔憂.
令到我們很膽怯.
但在這裡說耶穌的平安是甚麼平安呢?.
其實耶穌的平安就是讓我們內心有一種很安全.
很穩妥的感覺.
這些感覺不會因為外在環境的挑戰而影響.
而帶來長期的困擾.
的確我們會有情緒.
我們遇到困難.
遇到一些外在環境的影響.
我們的確會有一些情緒出現.
但真正的平安給予我們的.
就是我們可以去應付面對這些不安的情緒.
而跨過去.
從心底裡有一種安穩.
有一種安恕.
更加是一種可以使人和睦的力量.
並且是一個在人生裡面.
去預見人生一個正面和積極的態度.
其實有耶穌就有平安.
令我想起一個76年前的故事.
在76年前.
就是第二次世界大戰.
1944年的時候.
當時德軍就面對著盟軍的攻佔.
的打仗.
就發動了一個叫做布爾格之戰.
這個戰爭裡面.

$^{81}$有一段最初是不為人所知.
但後來在90年代.
是被傳頌的一個很動人的故事.
這個故事還在2002年的時候.
拍成為一部電影.
那部電影的英文名叫做《Silent Night》.
《寂靜的夜》.
這個故事是在說什麼呢?.
這個故事是在說一個德國的婦人.
她的名字叫做維肯太太.
因為她逃避戰亂.
她就帶著一個12歲的兒子.
叫做弗洛斯.
他就離開了他身處的小鎮.
去到比利時邊境的一個地方.
在那裡住了進一間用來狩獵的小木屋裡.
維肯太太就等她的丈夫.
有一段時間就將一些食物送給他們母子.
就去逃避這個戰役.
就是在1944年的平安夜的夜晚.
那一晚維肯太太打算去宰了家裡的肥雞.
然後打算敲起那隻雞.
就預備和她的丈夫度過這個平安夜的夜晚.
但是那一晚.
有些聲音在她的屋外傳過來.
弗洛斯很興奮.
以為爸爸來了.
於是她就走出屋外打算迎接爸爸.
誰知她發現的是三個美國的美軍.
美國的士兵.
維肯太太他們是不懂英文的.
不過當她見到其中一個士兵.
他是被槍傷.
然後他們是在一個困難困乏的境況裡面.
她就冒著去包庇或者去窩藏敵軍的一個危險.
然後就收容了他們.
她安置了他們.
為他們包紮了傷口.
然後就預備了去宰那隻雞.
還叫了她的兒子.

$^{121}$預備了六個馬鈴薯.
為他們敲一個聖誕的大餐.
過了不久.
外面有人敲門.
弗洛斯以為爸爸來了.
誰知一走出去.
發現的是四個德軍.
在這個時候.
維肯太太去到門口.
見到這些德軍荷槍實彈.
她很冷靜地跟他們說.
今晚是平安夜.
希望你們絕對不要在我們這個地方開槍.
然後各個軍人就把槍放在屋外.
你可以想像得到.
有七個士兵.
在一個狹小的木屋裡面.
氣氛是很緊張的.
但是很奇妙.
維肯太太依然端上了.
她預備了的烤雞和美食.
然後熱情地款待這兩個不同國家的敵人.
然後家裡的氣氛就開始慢慢地鬆弛了.
維肯太太帶著這麼多人.
去作了一個謝飯的祈禱.
她求上帝去止息戰爭.
亦都去保護這些士兵.
當時這七個軍人感到非常非常之感動.
不久後.
德國的士兵就拿出了他們的一瓶酒和麵包.
而美國的軍人亦都拿了他們身上的咖啡.
大家就一起在飯菜上享用.
甚至有一個懂得英文.
又懂得一些醫療訓練的德軍.
就幫那位士兵去拿出他的彈頭.
第二天早上.
德軍用他們的軍用地圖.
指示了他們怎樣走回到美軍的基地當中.
維肯太太對著他們每一個說.
我希望有一天你們平安地回到你們自己的家鄉.

$^{161}$願主祝福保護你們.
賜你們平安.
接著七個士兵就擁抱道別.
分道揚鑣.
每當維肯太太去益述這件事的時候.
她說在飯桌上主耶穌和我們在一起.
在一個戰火的裡面.
當人與人之間面對自己的仇敵.
環境的確非常緊張.
人的心有很多不安.
有很多思緒.
裡面可能有仇恨.
有驚慌.
有絕望.
有擔心.
但是當主耶穌在他們當中的時候.
他們的內心就得到真正的平安.
彼此之間人與人之間有愛.
打破了裡面的隔膜.
有力量去面對一個緊張的惡劣的環境.
就是這份平安.
讓他們能夠經歷耶穌的愛.
事實上耶穌在那段經文裡很清楚地告訴我們.
其實他所賜的平安不像世人所賜的.
他賜給我們的平安才是真正的平安.
但是很可惜.
世上的人很多很多都是用世上的平安.
來滿足自己裡面的不安.
裡面往往都是白費心機.
徒勞無功.
回想起來.
我自己的家就是一個很好的例子.
我還記得我小時候家裡比較窮.
一家七口住在一個狹小的公屋裡面.
爸爸媽媽要謀生.
媽媽身體也有很多病痛.
再加上其實爸爸很多時候.
都會和鄉里的同事打麻將.
或者當鄉里來到香港的時候.
他都會去關心他們.

$^{201}$接濟他們.
家裡的經濟其實都很緊絀.
經常都在想我們怎樣可以賺多點錢.
令我們開心一點.
後來到了中學.
我還記得我們一家人很努力地去賺錢.
放工.
放暑假.
有假期.
我們全部人都一起回到工廠裡面.
一起去工作.
慢慢家裡的生活就得到大大的改善.
我們由一間小型的公司.
搬到一個西貢的三層高的別墅.
我們的工廠由一間小廠開了兩間大廠.
然後又在國內繼續發展業務.
不單止這樣.
我還記得我們小時候.
三四個人擠在一張床上.
但是我們可以改善到一個地步.
就是我們每一個人有一個一百多呎的房間.
我們以前早上是坐巴士上學的.
但是後來我們可以有三個座駕供我們出入.
生活很好.
不過環境對我們來說是一個很大的衝擊.
1998年的一場金融風暴.
將我們的家庭由一個富裕的經濟.
變成一個很大的危難.
當時因為家庭投資失利.
內地的工廠也不能經營.
在香港的工廠也因為有人欠債.
我們要將西貢的房子變成銀主盤.
給銀行去抵押債務.
我們一直追求的安全和幸福.
一夜之間都沒有了.
不過有很多人說.
不要緊,就算窮也好.
一家人齊齊整整就可以了.
但是當我們家庭慢慢富裕起來的時候.
媽媽就發現了爸爸有婚外情.

$^{241}$並且他長期都留在內地.
媽媽當時心灰意冷.
平時多麼堅強的她.
一下子就倒下.
她曾經想過要了斷自己的生命.
所謂一個在感情裡依靠幸福的家庭.
一時之間沒有了.
幾年之後爸爸回到香港.
他回來香港是因為他患了重病.
他染了一個嚴重的抗藥性肺癆.
以及末期的癌症.
我還記得當時他還沒進醫院之前.
每天都要去健康院吃肺癆藥.
他乘車回家去到月華街的街口.
他完全感覺到喘氣和無力.
他坐在公園很久.
結果他還是要截的士.
回去只是距離150公尺的家庭.
爸爸的經歷是什麼都沒有.
錢沒有了.
家庭沒有了.
健康也沒有了.
那是不是就拿不到真正的平安呢?.
很奇妙.
就在這個患病的過程當中.
我爸爸得到人生最寶貴的一件事.
就是認識主耶穌.
信靠主耶穌.
他信了主耶穌之後.
我見證著他去倚靠耶穌.
他向上帝祈禱.
身邊有很多人為他祈禱.
讓他不去懼怕他患了疫病.
甚至他對於將來要離開這個世界.
回到天家.
他心裡也沒有盼望.
但是我看著媽媽.
我知道他心裡很矛盾,很忐忑.
爸爸要離開這個世界.
但是媽媽呢?.

$^{281}$當爸爸要離開這個世界的時候.
如果媽媽當時.
仍然心裡懷著一個未能解開的結.
這樣帶著中老會怎樣呢?.
我心裡一直向神祈求.
希望媽媽都可以得到這份真正的平安.
終於有一天.
我們接到電話.
醫院打來說爸爸不行了.
叫我們去見他最後一面.
當時我們每個人都在不同的地方.
趕到醫院在爸爸的床邊.
我們圍著床看到一個軟弱無力的爸爸.
但是我看到他雙眼一直看著媽媽.
他心裡好像有很多話想跟媽媽說.
我終於聽到他鼓起最大的勇氣.
呼出最後一口氣.
去跟媽媽說對不起.
對不起.
當時媽媽愕然.
她完全不懂得反應.
我在身邊告訴媽媽.
媽媽,爸爸對你說對不起.
你原諒他吧.
媽媽在這個時候已經哭到無法收聲.
她只是不斷不斷不斷地點頭.
然後慢慢她開始說話.
她說我原諒.
我要跟你信耶穌好不好.
爸爸點頭.
接著爸爸就離開了這個世界.
我爸爸在環境裡真的什麼都沒有了.
但是他有耶穌.
他能夠得到真真正正的平安.
並且我媽媽相信了耶穌之後.
她的人生不同了.
她以前是很憂愁.
以前是很不開心.
很不喜歡見人.
但是因為她相信了耶穌.

$^{321}$她在教會裡面有很好的生活.
她還很熱心地去服侍人.
把耶穌這份福氣跟人分享.
我家裡就經歷了一個.
耶穌的平安如何改變我們.
但是人究竟是如何得到.
從耶穌而來的平安呢.
剛才陳牧師有說過.
以色列人所說的平安.
其實是有一個圓滿的意思.
其實當神去做我們的時候.
祂就是盼望我們有一個圓滿.
祂想我們有完整.
祂想我們人生得到真正的平安.
但是為什麼人會沒有平安呢.
因為你和我都有罪.
有罪的意思就是.
我們都自我中心.
我們都目中無神.
我們都不認上帝.
有一位牧師曾經說過.
他說過一個很好的比喻.
對於我們是一個很好的提醒.
他叫我們想像一下.
如果你是一個很本事的兒子.
他讀書很厲害.
做事很厲害.
賺了很多錢.
又買樓收租.
又娶了一個很好的老婆.
生了一些很厲害很孝順的子女.
不過.
唯一有一件事.
他沒有的.
就是他不認你做爸爸.
不認你做媽媽.
你還可不可以說.
這個是乖兒子.
他擁有一切.
但是他缺缺了一件很重要的事.

$^{361}$我們都是一樣.
我們必須去認我們的上帝.
我們的上帝為我們安排最好的路.
讓我們去走.
不過我們拒絕信他.
我們用自己的方法.
去順著自己的慾念.
這樣去行.
就好像我們人生喝了一些鹽水.
喝多了都是頸渴.
我們用自己的方法去解決問題.
結果.
我們的心.
永遠都得不到滿足.
永遠都得不到.
上帝已經為我們預備了那份圓滿.
而且我們還在很多人生的路途上.
還因為要滿足自己的時候.
可能製造了很多人際關係的不和.
傷害自己.
傷害到別人.
不過我們要知道.
其實平安的希伯來文.
是很有意思的.
當你將這個字的前面.
加一個介詞.
即是英文的for.
為了圓滿.
當這個字去組合成功的時候.
其實它的意思就代表了付帳單.
代表了還清債務.
代表了可以恢復這個圓滿.
代表了我們可以同神.
重新建立一個平安和和好的關係.
不過這張帳單由誰去付呢?.
是耶穌.
耶穌很愛我們.
祂在十字架為我們犧牲.
擔當了我們的罪.
並且祂在第三天復活.

$^{401}$祂賜給我們一個豐盛的生命.
又賜給我們一個永恆的盼望.
只要我們能夠相信這位拯救我們大能的主.
信靠祂.
並且能夠將我們的生命交託給祂.
讓祂能夠主宰我們的生命的時候.
就可以得到真正的平安.
真正的和好.
回到真正的圓滿裡面.
這份珍貴的平安.
是上帝為我們預備的.
就好像剛才淑華和Epi告訴我們.
當我們拿到這份平安.
和神和好之後.
上帝寬恕我們.
我們能夠和祂有一份好的關係.
以致我們和自己能夠和好.
也能夠和別人和好.
我們的罪債得到還清.
我們得到上帝完全的寬恕.
得到人生真正的平安.
以賽亞書告訴我們.
聖誕節紀念耶穌這位英雄降生來到這個世界.
他是這麼說的.
「有一英孩為我們而生」.
「有一子賜給我們」.
「政權必擔在他的肩頭上」.
「他名稱為奇妙賊士」.
「全能的神永在的父和平的君」.
就是主耶穌.
祂給了我們真正的平安.
因為主耶穌來到這個世界.
祂經歷過我們人世間的種種關係.
經歷過人生的種種疾苦.
祂明理多過世界上任何一個人.
祂愛你深過世界上任何一個人.
所以祂要給你的.
一定是最好的.
耶穌是永遠都不會離開我們.
可以給我們真正平安的人.

$^{441}$親愛的朋友.
我很小就認識了這位上帝.
認識了耶穌的拯救.
祂賜給我有很大的平安.
讓我在人世間.
其實我都經歷了不同的困難和挑戰.
家庭,經濟,人際關係,錢等等.
很多事情發生在我身上.
但是神真的給了我平安.
我還很記得.
當我的小兒子在五歲多的時候.
發現有天生的心路症的時候.
醫生最初說可以用一個微創的手術.
幫他做好手術之後就可以恢復.
但是結果我們在手術室門外.
等了超過一倍的時間.
他出來醫生說對不起.
做不到.
後來他要做一個開心手術.
開心不是很開心.
是要將肋骨和心臟開刀.
去做一個大手術.
我們人就會想.
我們知道有最好的醫生.
因為我們認識.
知道有最好的醫生安排了為他做.
誰知臨前一晚.
我們發現原來換了一個我們不認識的醫生.
我們心裡很擔憂.
我們心裡很掛心.
但是當我們去祈禱.
去尋求上帝的時候.
我們很清楚知道.
上帝比我們更愛我們的兒子.
我們放心交託.
以致當他在手術室裡面.
我們在門外的時候.
心裡真的有莫大的平安.
親愛的朋友.
這個世界的確可以有兩種平安.

$^{481}$一種平安是世界給你的.
是隨著環境會變動的.
但是一份從耶穌而來的平安.
是永遠不變的.
可以在你的心裡.
有一份超越環境而有的真正平安.
林肯總統說得很對.
他說一個人要選擇有多快樂.
他就可以有多快樂.
其實平安是可以選擇的.
我在這裡誠意邀請你.
去選擇一個從耶穌而來的平安.
一個真正的平安.
讓我們稍稍有一些安靜的時間.
去思考在這個平安夜的晚上.
上帝跟你分享這個平安的訊息.
你要怎樣回應呢?.
\newpage



\section{馬太福音哥林多前書 11:28-30}
\label{sec:DokqOqhZ3o0}
\textbf{2020年12月27日崇拜直播｜羅志強執事｜福音與愛的實踐｜馬太福音11\_28-30;哥林多前書13\_4-5}
\newline
\newline
連結: \href{https://youtube.com/watch?v=DokqOqhZ3o0}{\texttt{https://youtube.com/watch?v=DokqOqhZ3o0}} ~~~~ 語音日期: 2021-01-03
\newline
\newline
\hyperref[sec:lmnMx2HA46Q]{\small{< < < PREV SERMON < < <}}
~
\hyperref[sec:index_chronic]{\small{[返順時目]}}
~
\hyperref[sec:bLhELYX_Uw0]{\small{> > > NEXT SERMON > > >}}
\newline
\newline
$^{1}$兄弟姐妹今天我們會一起去重讀兩段經文.
兩處的經文是來自馬太福音第十一章二十八節到三十節.
以及哥林多前書第十三章四至五節.
我先讀馬太福音.
凡奴婦擔重擔的人可以到我這裡來.
我就使你們得安息.
我心裡柔和謙卑.
你們當負我的厄 學我的洋式.
這樣你們心裡就必得享安息.
因為我的厄是容易的 我的擔子是輕散的.
另外一處經文哥林多前書十三章四至五節.
愛是恆久忍耐又有因此.
愛是不疾度 愛是不自誇 不張狂 不作害羞的是.
不求自己的益處 不輕易發怒 不計算人的惡.
今天為我們正道的是我們親愛的羅志強執事阿叔.
今天宣講的主題是福音如愛的實踐.
弟兄姊妹你們好.
聖誕節過了.
所以我預祝大家新年快樂.
在這裡我很感恩.
自從去年大約八月的時候.
我確診肺炎的時候.
到現在差不多有一個年半的時間.
已經有差不多兩個年的時間.
沒有機會和大家分享神的話語.
所以很感恩神今天在零年給我有機會和大家分享.
在這裡我不單止為自己感恩.
我也要為教會感恩.
作為執事會主席.
我在這裡我願意代表教會.
去多謝我們每一位傳道同工.
和其他幹事的配合.
因為疫情的緣故.
我們同工沒有因為這樣的緣故而放下他們的工作.
反之他們運用各方面的聰明和智慧.
媒體組的幫助.
使我們在媒體上做得相當出色.
不論是網上的崇拜.
不論是網上的祈禱會.
不論是四十天的牧者家書.

$^{41}$不論是刺客心聲.
新世代書俗.
甚至最近的奇怪的禮物.
我們看到同工的努力不懈.
不單止這樣.
因為疫情的緣故.
我們更加多了和不同的福音機構合作.
不論是無家者協會.
新年心.
風雨同路人.
銀行團契.
家庭發展基金等等.
我們很感恩.
我們雖然沒有實質崇拜一段時間.
但我們看到在網上的參與崇拜人數.
比我們以前實質崇拜的時候還要多.
這是很感恩的事.
當然我們也要多謝各部的合作.
不同的部長城市有性的財務.
都是很努力地.
用不同的方法繼續去運作.
牧羊教會幫助弟子們.
而且在疫情當中.
經濟困難的底下.
我們看到.
奉獻不單止沒有下跌.
反而增加了.
我們都是為此感恩.
多謝弟兄姊妹的努力.
我希望從長到青年到少年.
我們每一位弟兄姊妹.
都能夠在接下來的奉獻上.
事奉上.
服侍上.
更加努力.
明年是修國的一年.
我們的主題是福音覆蓋跨平台.
以年輪謝傳遞愛.
我希望當疫情慢慢過去的時候.
我們有實質崇拜的時候.

$^{81}$我們能夠有修國的時間.
看到實質人數的增長.
今日我的題目是.
福音與愛的實踐.
用了兩段很簡單.
大家會背的經文.
福音是什麼?.
其實不用說都懂.
福音就是good news.
好訊息.
什麼是好訊息?.
福音是主耶穌基督.
道成六生來到世間的時候.
為了拯救世人.
最先聽到福音的.
好像上兩星期陳牧師都提過.
是野地裡的牧羊人.
在白地行的野地裡的牧羊人.
我報給你們大喜的訊息.
是關乎萬民的.
因今天在大會的城裡.
為你們伸起救主就是主基督.
牧羊人是第一個聽到好訊息的人.
在主耶穌工作的三年半裡.
他走遍猶太各城各鄉.
傳遍天國的福音.
最後為人類釘身在十字架上.
做了人類的贖造祭.
完成了救贖的工作.
而福音的中心.
在《約翰福音》第三章第十六節.
講得更加清楚.
神愛世人.
甚至將他人的獨生子賜給他們.
叫一切信他的不至滅亡.
反得永生.
人類學家曾經這樣說.
任何時代任何地方的人.
都期盼死亡並非生命的終結.
他們尋求一個人永遠的生命.

$^{121}$福音似乎給了你這個答案.
是能得到永生.
很多時候我們集中的問題就在這裡.
得永生.
但耶穌曾經說過.
當我們信耶穌的時候.
「每日在今世不得八倍.
在來世不得永生的」.
也就是說不單單只是來世.
今世我們都能得到好的福樂.
很多時候我們說.
基督徒的標誌是什麼?.
是愛.
但我常常說.
不單止愛.
還有一樣東西.
是喜樂.
你想想.
盡你多大的愛心.
我看到一個基督徒.
經常愁眉苦臉的.
你叫我信耶穌就難了.
所以我覺得除了愛之外.
喜樂.
活得開心快樂.
是基督徒的第二個標誌.
但問題是.
為什麼我們活得不快樂?.
我用這段經文.
《瑪達福音》所說.
凡勞苦擔重擔的人.
可以到我這裡來.
我便使你們得安息.
我心裡柔和謙卑.
你們當負我的額.
學我的樣式.
這樣你們心裡就必得著安息.
因為我的額是容易的.
我的擔子是輕散的.
耶穌邀請我們.

$^{161}$卸下我們原先背負的重擔.
那些額.
轉而背負耶穌的額.
不是不用你負額.
而是轉負一個較輕散的.
或者更加合適你的額.
我的額是容易的.
容易也可以逆作合適.
在巴勒斯坦.
那些牛.
通常那些額.
那些牛的額.
都是木製的.
很多時候農夫會將牛.
拖到木匠那裡.
先做一個底樣.
一個帽子.
差不多做好了.
量好尺寸.
差不多做好了.
又要拖一隻牛來.
試一試它是否合適.
這個和我們舊時做西裝差不多.
現在很多人做西裝.
是完成的.
舊時我們不是的.
做西裝的時候.
就選了布料.
量身體.
做一個帽子.
然後到兩個星期.
多少時間之後.
又回去.
再量一量.
所以最後出來的西裝.
最合適你了.
當時原來巴勒斯坦都是這樣的.
那些牛.
那些額.
最後是最適合那隻牛的.

$^{201}$有些假人家.
甚至相傳.
當然不知是真是假.
說耶穌當時是巴勒斯坦.
可能是一個最出名的做額的木匠.
如果是這樣的話.
就更顯得這句說話有意思.
耶穌說.
我的擔子是輕散的.
這個不是說擔子不用扶.
而是說.
你扶一個耶穌基督給你的.
輕散而合適的額.
而且愛呢.
能夠將你額顯得更加輕散.
我看過一個遇到的故事.
提到一個人.
看到一個男孩.
十三四歲.
他背上一個七八歲的男孩.
當然是有強腿的.
這個人看見就覺得.
你背上的擔子真是重啊.
那個男孩怎麼回答呢.
他說這個不是重擔.
這個是我弟弟.
愛就能使額變得更加輕散.
但法理賽人呢.
卻將這個輕散的額變得沉重.
被殺種的時候.
不能用混雜的種子殺在田間.
收割的時候.
亦不能夠與別下一群的時候.
不可以再回去取.
當安息的時候.
可行的路程只可以二千集.
林林種種.
很多.
多到我相信他們自己都不記得清楚.
但我們不要少這些法理賽人.

$^{241}$其實現在我們都是一樣.
我們都是現代的法理賽人.
我們都有很多這方面的專家.
我們將信仰無限的演繹.
我們比自己.
而比別人.
更多更多更壓.
有時我覺得有些人.
會喜歡環保.
有些人會追求公義.
或者政見上的改革.
教育制度上的改革.
扶貧.
甚至我聽過一些扶貧的弟姐妹這樣說.
她說最好不要養狗.
養狗就將錢給周濟窮人.
我心想.
吃著餐去旅行都不好.
很難搞.
我甚至很多年前.
我試過一次.
與弟姐妹唱卡拉OK.
有人勸我做什麼呢.
他說你不要唱成全中聲.
我說唱大戲就不行了.
死定了.
所有大戲都是無端神佛.
有時每個人.
神給你有不同的使命和呼召.
這是神給你的輕散而適合你的壓.
你做是沒有問題的.
但不要將這些壓.
同樣加在別人身上.
甚至無限的加在自己身上.
當然有些根本不是壓.
根本與信仰甚至沒有太大關係.
又成為一個壓.
配在身上.
教會都是這樣.
有時我們覺得教會應該扶貧.

$^{281}$或者單身的弟姐妹.
或者單親的弟姐妹.
林林總總.
但教會始終有不同的priority.
在每一年的裡面.
我們何苦將這麼多壓.
放在自己身上呢.
有時我們活得不快樂.
因為我們太多壓在身上.
我希望弟姐妹都能夠明白.
神給你的壓是輕散而適合你的.
我們能夠背負神所呼召我們的.
所給我們的壓.
你就能夠在當中成長.
而活得開心快樂.
但活得開心快樂都不夠.
其實開心快樂是自我的.
給我們自己的見證.
吸引人們.
這些人為何活得這麼開心.
但都不夠.
傳福音.
「熱年輪謝傳遞愛」.
要有愛.
今日很大膽地用了林左前書.
為何我說大膽呢.
因為大家都懂得回.
有甚麼好說呢.
所以我發覺原來有時回口枉.
實際這些經文說甚麼.
我們可能一知半解.
所以我今日大膽用這些經文.
給予一支勉勵.
所以時間關係.
就算我選了兩節都不是全段.
可能我只能夠說這兩節的一部份.
雖然我已預備了.
但時間關係我不能每一句都說.
我只說幾句.
其他我比較簡單的帶過.

$^{321}$第一句.
愛是恆久忍耐又有恩慈.
忍耐甚麼呢.
是否忍耐約了女朋友她遲到呢.
或是忍耐我現在回到教會.
為何這麼久沒有冷氣.
忍耐甚麼呢.
是否忍耐公司遲遲不升職呢.
你就忍耐吧.
其實忍耐在新疾中.
這個慈這個希臘字的意思是甚麼.
是指對人的忍耐.
不是對環境的.
甚至是指一個人受到傷害.
他有足夠的能力去報復.
有足夠機會去報復.
而他不報復.
這種是忍耐.
正如上帝如何忍耐我們一樣.
忍耐並非弱者的行為.
相反是剛強的表達.
有一個很多年前的例子.
很久以前.
我經常都記得.
是我讀書的時候.
在加拿大讀書的時候.
自問的時候耶穌都很乖.
我們四個男生.
在宿舍.
一個兩張床的公寓.
怎樣乖呢.
每晚都去圖書館讀書.
讀完書圖書館關門.
九點多鐘就關門了.
我們回家吃一點東西.
過個祈禱.
都是標準的約學生.
樓上有四個中國人.
他們打樂隊.
我不是打錯.

$^{361}$他們晚上十二點鐘才打樂隊.
在樓上很熱鬧.
如果要報復.
我便去找社監去打他們.
但我們沒有.
因為都是中國人.
有時我會想.
可能我心裡很空虛.
晚上沒有東西想.
唯有打樂隊.
但為什麼說呢.
有一晚社監來我們那裡.
關門.
我們開門.
他說為什麼你這麼吵.
我們沒有吵.
原來我們樓下.
例如我們住七樓.
八樓打樂隊.
六樓投訴七樓.
他以為我們吵.
多慘呢.
是不是可以踩多他兩腳.
但沒有.
最後反而我和樓上的.
五個弟妹朋友.
建立了友誼.
很多時候.
報復並不帶來對你任何的補償.
反之能夠建立仇恨.
看我小時候就知道.
通常報復復仇那些.
復仇後都變成空虛.
一片茫然.
所以為什麼要報復呢.
其實報復對你沒有什麼好處.
但現在你反而不報復的時候.
反而有機會建立友誼.
忍耐的價值在乎恆久.
愛是我們有持久的忍耐的能力.

$^{401}$不單止這樣.
不單止要忍耐.
不單止要持久.
還要加上恩慈.
恩慈按字面解很簡單.
恩慈按字面解就是.
恩典加慈愛結合.
是一種充滿連率.
能夠體恤別人的一種生活狀態.
一種表現.
是指那些得罪了我們的人.
那些對我們不義的人.
我們激勾了我們的人.
我們都要恩慈去待他.
我不知道你有沒有別人.
傷害過你.
煽惑過你等等.
但當你能否這樣對待他們.
當你有機會報復的時候.
很多時候我們覺得.
基督徒應該可以.
能夠恩慈待人有什麼難.
基督徒講愛心.
但事實上當你反觀這個社會的時候.
似乎未必一樣.
未必一定.
反觀歷史都是一樣.
在中古時代我們看到宗教迫害的時候.
其實那種慘烈難以想像.
現代人有時都會.
有時我們因為不同的宗教信仰.
不同的理念.
甚至在教會不同的方向的看法.
我們對其他人絕對沒有恩慈.
這些基督徒應該有的表現嗎?.
沒有一個有愛心的人.
會去愛待他人.
其實作為基督徒.
我覺得我們應該懂得去恩慈待人.
不單不記仇.

$^{441}$不單不報復.
有機會都不報復.
而且去恩慈待他人.
這是愛.
但這樣都不夠.
很多年前我看過一篇文章.
一句說話.
引用它.
講施捨.
施捨是什麼呢?.
施捨是.
借而不忘其還.
這樣叫施捨.
這樣很好.
施捨是什麼呢?.
我施捨給你.
好像我高高在上.
借不同.
同等.
雖然我不期望他還給我.
這是更加好的施捨.
同一道理.
我再說下一個說話.
不計算人的惡.
什麼是不計算人的惡呢?.
計算這個字本身.
出自一個會計的紀錄.
有關.
意思是我們心裡有一本賬簿.
將別人對我們不好的.
記下.
在賬簿裡.
很多時候我們間中看到.
一對很恩愛的夫婦.
已三十年了.
很恩愛的夫婦.
一對弟兄姊妹.
一對兩位弟兄.
三十年的老友記.
因為很小的事情.

$^{481}$罵了.
或者甚至離婚.
哎呀!你說為什麼這樣?.
這麼小事.
去歐洲旅行.
選擇英國還是法國.
這樣也要離婚.
這麼小事.
不是.
那本賬簿寫滿了.
剛剛到那個位置.
就是breaking point.
臨界點.
其實.
我們為什麼要這樣?.
為什麼要記住一些.
別人對我們不好的事?.
有一句格言.
人生格言也不錯.
人情當晨初相見.
道老終無怨恨心.
人情當晨初相見.
道老終無怨恨心.
意思是什麼?.
人與人之間的相處.
如果能夠永遠.
保持在第一次見面的情況.
看到美好的一面的時候.
多好.
美好能夠維持到.
終老的時候.
就能夠維持良好的關係.
人與人之間的相處.
就是這樣.
愛是不懷怨恨的心.
其實.
剛才我說.
施捨是什麼?.
借而不忘記還.
而當別人羨慕你.

$^{521}$對你不好的時候.
我們不單單只有機會.
有能力而不報仇.
而且能夠持久.
因此代他之餘.
更加高的境界是什麼?.
是忘記.
我看過一段短文.
寫到忘與記.
不知道大家有沒有分享過.
年幼的時期.
我以為記得.
勞是真本事.
過目不忘的大腦.
是真天才.
中年以後.
我逐漸領悟到.
忘掉才是真幸福.
忘不了別人的閒言閒語.
人生會披上一層灰色的陰影.
忘不了傷心的往事.
人家會逐漸扭曲.
壯年以後.
我開始向神求見忘之恩.
忘掉過去的輝煌.
這個叫謙卑.
忘掉以往的失敗.
這個叫勇氣.
忘掉從前的創傷.
這個是饒恕.
忘掉昔日的罪過.
這個是感恩.
忘掉朋友的不周.
這些是大方.
忘掉受敵的攻擊.
對你的損害.
給你的傷痛.
這些是愛.
既是聰明.
忘是智慧和修養.

$^{561}$是愛.
愛是恆久忍耐而有恩慈.
愛是不嫉妒.
雖然隔著年光幕.
如果我問.
「鄧小姐,有覺得自己嫉妒他人的請舉手」.
我猜也沒有人舉手.
沒有人說自己喜歡嫉妒人.
但為什麼我不嫉妒呢?.
那羨慕又如何呢?.
你有沒有羨慕他人呢?.
他真的多帥.
他真的多美女.
他想事這麼厲害.
想事這麼快.
他賺這麼多錢.
他的身體這麼健壯.
我們會不會羨慕他人呢?.
羨慕他人這種心態.
我相信很難避免.
或多或少也有.
但我們羨慕他人的情況.
覺得垂涎他人所有.
而希望自己也得到.
甚至覺得別人擁有的東西.
有些不配.
再發展下去.
這就是嫉妒.
再發展下去.
這就是嫉恨.
其實嫉妒很有趣.
被嫉妒的人很冤枉.
為什麼呢?.
他可能不知道你嫉妒他人.
受苦的是嫉妒別人.
這樣是不是很實在呢?.
為什麼你嫉妒他人呢?.
嫉妒心很多時候不會傷害別人.
卻會傷害你自己.
又何苦呢?.

$^{601}$保羅有一段說話.
記載在羅馬斯.
他說:以喜樂的人要同樂.
以哀哭的人要同哭.
說話容易.
實際行走也有些困難.
在教會有一種文化.
如果弟兄姊妹在小組.
有傷病甚至經濟困難.
我們都樂意幫助他們.
以哀哭的人同哭.
我看到很多這樣的例子.
弟兄姊妹對傷病的弟兄姊妹照顧.
甚至我最近在這段年紀.
我感受到.
接受了很多弟兄姊妹的幫助.
這我們做得不錯.
甚至看到有些弟兄姊妹有經濟困難的時候.
我們有很多弟兄姊妹經濟能力較好的.
都樂意幫忙.
無條件的.
但很少以喜樂的人同樂.
我們教會代討事項.
十居其九都是傷病.
有沒有見過弟兄升為處長.
或者他升為總經理.
他達福加薪.
他甚至他的學業有很好的成就.
我們為他感恩.
這很少見到.
我們能與哀哭的人同哭.
但我們不懂得.
要喜樂的人同樂.
甚至我們久而久之形成這種文化.
可能覺得這叫謙卑.
我們覺得如果我們分享這些.
是不是很驕傲.
甚至這種心態影響到我們事奉的心態.
所以我們不追求.
在事奉的崗位上做領導.

$^{641}$我們追求名利.
利就不用說了.
哪有利?.
名都沒有.
為何在教會追求名利?.
在外面世界追求名利很多得士.
我不是說謙卑不對.
先說明.
但謙卑有時不一定是這樣的.
舉個例子.
很多年前.
我曾經負責一個州會.
叫偽似的謙卑.
甚麼意思呢?.
我不知這個字眼是否正確.
當時是這樣寫.
舉個例子.
如果你是一個中學會考理科的狀元.
有個中三仔來跟你說.
哥哥,可否幫幫我?.
我一條數不懂.
你告訴他.
我不知道懂不懂.
人家會怎樣想?.
理科的狀元對中三的數學題目.
不如不要說這句.
我看看怎樣.
這樣就簡單.
有時我們過份了.
過份了.
甚至這些影響我們事奉的心態.
不要嫉妒他人.
甚至不要羨慕.
要怎樣呢?.
去欣賞.
欣賞他人的好處.
不單以哀樂的人同樂.
能夠以喜樂的人同樂.
這是愛.
愛是恆久忍耐又恩慈.

$^{681}$愛是不疾度不自誇不張狂.
不自誇不張狂我不想說.
簡單來說.
不自誇不張狂是甚麼呢?.
不自誇就是不聰明.
不張狂就是吹噓自己.
你回去想想.
為何不聰明不吹噓自己就是愛呢?.
不作害羞的事.
我們不說時間關係.
不作害羞的事不如想想.
甚麼是不作害羞的事呢?.
是害羞的事不做嗎?.
這句意思不是害羞.
不要搞錯.
不作害羞的事的意思.
大概是.
一個不合禮儀的事.
可以作不知禮.
不作害羞的事.
是指不做任何沒有教養粗魯行為.
是這個意思.
當然最後我想說.
不求自己的益處.
最後一句不說.
最後一句說不輕而發怒.
當然不是說不能生氣.
偶爾生氣也可以.
我剛才說耶穌也會入法二道.
但問題是.
生氣也是一樣.
不要胡亂生氣.
生氣會傷身.
不可含怒對人惡.
生氣也是一樣.
生氣是用別人的錯來懲罰自己.
有時你真的很蠢.
我們不要這樣做.
最後我說的原因是甚麼?.
不求自己的益處.

$^{721}$這個說話.
細心想一想.
不太合理.
為甚麼?.
人都是自私一點.
怎可以不求自己的益處?.
求別人的益處.
真的難一點.
有人說世上有兩種人.
一種是堅持自己的權利.
一種是常常記得他人.
記得他自己責任的人.
另一種是常常想別人欠他.
一種是永遠不忘記他欠人.
這個世界第一種.
現在很多人都是這樣.
這個社會,這個世界,這個政府.
欠了他.
我們實在是這樣.
但是.
人都幾自私.
基督徒可能好一點.
完全不求自己的益處.
是否這麼容易?.
有個俗語說.
人不為己.
天誅地滅.
我想問你知不知道上一句是甚麼?.
我十個月讀了九個九都不知道.
這句說話不是這樣.
整句說話是.
人身為己.
天經地義.
人不為己.
天誅地滅.
意思是甚麼呢?.
其實是出自佛家.
意思是甚麼呢?.
佛家的意思是.
不殺身,不偷渡,不邪淫,不妄語,不兩善等等.

$^{761}$其實這句說話很正面.
人身為己,就要你這樣為己.
其實是為人.
為他人.
不偷渡,不邪淫,不妄語,不殺身.
都是為他人的好處.
如果你不這樣做.
你就天誅地滅.
但現在我們不是這樣.
可能人都自私.
這個解釋好一點.
俗語說是過癮一點.
適合自己.
人不為己,天誅地滅.
這是一個講求自我的世代.
我們講求自我的權利.
忽略別人的權利.
我們強調我們自由做任何事.
我們忘記別人也有應有的自由.
因為我們太過自我中心.
自我中心這句說話好一點.
其實自私是更加合適.
但人是否就是這樣呢?.
人是否就是這樣呢?.
就是不能愛自己嗎?.
這未必定.
當你看周圍的弟兄姊妹甚至朋友.
有時都經歷或接觸到.
不同感人的真實故事.
例如父母因為子女的緣故.
願意將部份器官切割出來.
是為了子女的健康.
現在少一點.
以前看兒女殘片最多.
兄弟姐妹人多.
讀書沒錢.
哥哥姐姐最大那兩個.
不讀書了.
出來為了弟妹的緣故.
去謀生.

$^{801}$使他們能有好的教育.
現在我們在國內扶貧都見過類似例子.
當然為了國家,民族或大眾利益.
我們見過不少人.
願意犧牲自己都有.
這就是愛.
愛使這些變得不一樣.
愛使這些以為自己的事.
都能夠為他人.
愛能夠世人戰勝自我的自私自利.
而甘心去愛他人.
而付出一切.
無名的傳道者.
編雲波.
我引用他幾句說話.
他說自己的手甘心放下世上的享受.
似自己的腳甘心到苦難的道路上奔走.
選中這條不自由的道路.
並非無奈.
相反地.
卻是他正是大膽地用了自己的自由.
人因為愛的緣故.
其實真的可以放下自我的利益.
是為他人著想.
今日的分享大概到這裡.
我想跟大家說的是.
其實要呼音的傳人出去.
有兩點.
一是我們要有好的見證.
活得開心快樂.
為什麼活得開心快樂呢?.
不要將無謂的騙子.
放在自己身上.
或者放在他人身上.
神給你的騙是剛剛好的.
剛剛合適你自己的.
不要說對錯.
但你不要將你認為自己的.
硬加於他人身上.
你做好你自己的本分.

$^{841}$就可以了.
你就會活得輕散.
活得快樂.
但這樣都不夠.
我們還有愛心的實踐.
愛不是念口黃.
愛是能夠在實踐裡面走出來.
當別人對你傷害過的時候.
你不但有機會有能力去報復.
而不報復.
而且能夠恩慈地待他.
而且能夠最後將他對你不好的東西.
從你腦袋中抹掉.
愛是不嫉妒.
我們不要去嫉妒他人.
甚至不羨慕.
他是欣賞.
與喜樂的人同樂.
與哭泣的人同哭.
而且我們是能夠因為愛的緣故.
放下自己.
不求自己的積極處.
他是為他人的好處.
明年的主題是.
其中後面那句.
「結連能謝傳遞愛」.
願這個愛.
能夠借助我們的見證.
傳遞開去.
願神祝福大家.
\newpage



\section{馬太福音 5:1-10}
\label{sec:bLhELYX_Uw0}
\textbf{2021年1月3日崇拜直播｜陳明泉牧師｜幸福由虛心開始｜馬太福音5\_1-10}
\newline
\newline
連結: \href{https://youtube.com/watch?v=bLhELYX_Uw0}{\texttt{https://youtube.com/watch?v=bLhELYX\_Uw0}} ~~~~ 語音日期: 2021-01-03
\newline
\newline
\hyperref[sec:DokqOqhZ3o0]{\small{< < < PREV SERMON < < <}}
~
\hyperref[sec:index_chronic]{\small{[返順時目]}}
~
\hyperref[sec:POKxDUqFlWw]{\small{> > > NEXT SERMON > > >}}
\newline
\newline
馬太福音 5:1-10
\newline
\begin{longtable}{cl}
\hline
\hline
章節 & 經文 (和合本修訂版)\\
\hline
5:1 & \begin{tabularx}{0.7\textwidth}{X} 耶穌看見這一群人，就上了山，坐下後，門徒到他跟前來， \end{tabularx} \\ \\ \relax
5:2 & \begin{tabularx}{0.7\textwidth}{X} 他開口教導他們說： \end{tabularx} \\ \\ \relax
5:3 & \begin{tabularx}{0.7\textwidth}{X} 「心靈貧窮的人有福了！因為天國是他們的。 \end{tabularx} \\ \\ \relax
5:4 & \begin{tabularx}{0.7\textwidth}{X} 哀慟的人有福了！因為他們必得安慰。 \end{tabularx} \\ \\ \relax
5:5 & \begin{tabularx}{0.7\textwidth}{X} 謙和的人有福了！因為他們必承受土地。 \end{tabularx} \\ \\ \relax
5:6 & \begin{tabularx}{0.7\textwidth}{X} 飢渴慕義的人有福了！因為他們必得飽足。 \end{tabularx} \\ \\ \relax
5:7 & \begin{tabularx}{0.7\textwidth}{X} 憐憫人的人有福了！因為他們必蒙憐憫。 \end{tabularx} \\ \\ \relax
5:8 & \begin{tabularx}{0.7\textwidth}{X} 清心的人有福了！因為他們必得見神。 \end{tabularx} \\ \\ \relax
5:9 & \begin{tabularx}{0.7\textwidth}{X} 締造和平的人有福了！因為他們必稱為神的兒子。 \end{tabularx} \\ \\ \relax
5:10 & \begin{tabularx}{0.7\textwidth}{X} 為義受迫害的人有福了！因為天國是他們的。 \end{tabularx} \\ \\
[1ex]
\hline
\hline
\end{longtable}
$^{1}$讓我們一同聆聽上帝一段話語.
記載在《馬太福音》第五章一至十節.
經文是這樣說.
耶穌看見許多的人就上了山.
既已坐下.
門徒到他跟前來.
他就開口教訓他們說.
虛心的人有福了.
因為天國是他們的.
愛動的人有福了.
因為他們必得安慰.
溫柔的人有福了.
因為他們必承受地土.
忌渴無意的人有福了.
因為他們必得飽足.
連戌人的人有福了.
因為他們必蒙連戌.
清心的人有福了.
因為他們必得見神.
使人和睦的人有福了.
因為他們必稱為神的兒子.
為而受迫迫的人有福了.
因為天國是他們的.
今天早上陳明樹牧師會率先.
透過剛才所讀的八福經文.
來講第一個福氣.
就是虛心的人有福了.
因為天國是他們的.
講題是「幸福由虛心開始」.
弟兄姊妹平安.
在新的一年裡.
盼望大家繼續經歷從神而來的幫助.
無論我們生命是怎樣也好.
我們知道我們的命是屬於上帝的.
我們生命裡面.
只要倚靠上帝的時候.
我們會發現上帝有美好的恩典在我們生命當中.
剛才Raymond已經跟我們說過.
接下來這個系列就是一個關於八福宣講的系列.
我們不同的同工就會就著.

$^{41}$耶穌基督在這裡所講的登山寶訓.
八節的經文.
來講述一下究竟真正的福氣.
真正的幸福.
是一回怎樣的事.
在這段經文裡面.
好像很簡單幾節的經文.
但是當你真的細心去讀的時候.
你會發現當中裡面有很深的智慧.
以及對我們生命有非常大的幫助.
在我們落聖經前我跟大家分享一個片段.
有一個中年的男子.
大概跟我差不多的年紀.
有一次跟他一起吃下午茶.
吃一些很簡單的餅乾.
喝一杯茶.
其實那些餅乾就是.
獨立包裝的檸檬夾心餅.
不知道大家有沒有吃過.
大家吃的時候.
不外如是.
普通夾心餅.
我看到這個中年人就撕開袋子.
然後拿那塊夾心餅出來吃.
但是吃的時候吃得很滋味.
怎樣吃呢.
他就撕開了它.
然後.
中年人跟牧師聊天.
他就撕開它之後.
就這樣舔.
一邊吃的時候就吃得津津有味.
舔一下左邊又舔一下右邊.
然後舔舔的時候.
有些檸檬醬夾在手上.
手指也被吮了.
然後最後一口就把餅乾放進口裡.
我看到他就很好笑.
很奇特.
這些檸檬夾心餅怎麼吃得這麼滋味.

$^{81}$其實什麼味道這麼出奇呢.
為什麼吃得這麼開心呢.
這位中年人士就說.
其實我不是吃那個檸檬夾心餅.
我是吃童年.
他說其實在家裡小時候.
沒什麼錢的.
最奢侈的就是.
下午的時候家人能夠有機會.
開這些餅乾.
多快呢.
在一個很貧困的家庭裡面.
能夠有這塊餅乾.
每一口都是很滋味.
所以連手指也要舔一下.
因為每一個味覺.
都充滿了一個很豐富.
童年裡面很缺乏.
但是一塊餅乾裡面恩典的回憶.
這個人其實.
中年人大概以他的身家買了很多塊.
整家廠都可以買到.
不過.
一塊餅乾.
他帶來生命裡面一種很奇特的滿足.
後來這個中年人提醒我.
你要吃到這塊餅乾的味道.
你要回到兒童以前貧窮的那種經歷.
你才能吃到這塊餅乾的真味.
這句話對我很深提醒.
可能人生裡面我們經歷不到人生的味道.
或者經歷不到人生的幸福.
可能我們某程度.
我們未必回到以前的處境.
當我們退半格.
吞口半步的時候.
可能我們發現這個世界.
或者我們生命裡面的恩典處處.
今天我想和大家看的.
虛心的人有福了.

$^{121}$因為天國是他們的.
大概我們用這個角度來看.
究竟耶穌基督今天帶了一個怎樣的福音給我們.
在我們仔細看聖經前我們一起祈禱.
求上主的靈幫助我們.
同心祈禱.
天父我求你去幫助我們.
我們知道我們生命裡面.
因為認識到你.
我們是一個最有福的兒女.
幫助我們.
在我們聆聽.
研讀你的話語當中.
你的話語成為我們生命的幫助.
在新的一年開始.
我們祈求你打開我們的心扉.
幫助我們.
讓你的話語.
充充滿滿在我們心靈當中繼續的來到動工.
我們恭敬將來領受你話語的時間交給你.
幫助宣講的講得清楚.
幫助聆聽的弟兄姊妹聽得明白.
奉基督耶穌得勝的聖名.
Amen.
在《馬太福音》裡面有一個很重要的結構.
要在講這個系列之前要給大家知道的.
《馬太福音》其實是來自五篇很重要的講章.
以及五段耶穌的言行.
或者祂所做的神蹟奇事.
祂所過的那種生活來構成的.
簡單來說.
你當是一個十個段落的分段.
在第五章第一節到第二節.
剛才說到.
耶穌見很多人上了山.
就寄以坐下.
門徒到他跟前來.
他就開口教訓他們說.
值得注意的是.
那個分段就是.

$^{161}$耶穌在第五章二節那裡.
耶穌開口教訓他們說.
然後在第七章二十八節.
我沒有印給大家.
我讀給大家聽.
七章二十八節裡面.
他特別說到.
耶穌說完了這話.
眾人就稀奇他的教訓.
因為他教訓他們.
正像有權柄的人士.
不像他們的民事.
其實在五段的講章裡面.
全部都是會以耶穌開口教訓他們.
以及耶穌所說的話就說完了.
這個就是一個講章.
同樣的結構在十章第一節.
以及第五節.
耶穌也是這樣開口講論.
直到十一章第一節.
耶穌就結束他的講話.
十三章第一節.
也是這樣開始.
十三章第五十三節.
就是耶穌結束他的講話.
十八章第一節.
耶穌又再一次親口的跟他們講論.
十九章第一節就結束耶穌的講話.
最後一個段落就是二十四章第一節.
直到第二十六章第一節.
在這五個耶穌講論當中裡面.
就夾雜著耶穌所行的神蹟.
將耶穌所作的.
回應他所說的天國的道理.
所以馬太福音如果簡單來分段.
就是耶穌所言和耶穌所行.
他就卡拉卡拉著.
來將真理呈現開去.
所以在馬太福音裡面.
那個結構其實很工整的.

$^{201}$你用一個這樣的角度來看的時候.
耶穌不單止是說那些文士那樣.
就是純粹說的.
耶穌說完他的天國的道理之後.
他就在生活當中.
在他的全道生活裡面.
用神蹟.
用經歷.
用故事來呈現.
耶穌怎樣將他所說的話.
實踐出來.
用一個這樣的結構你會看到.
就是一路一路說的時候.
馬太福音將耶穌所言所行.
貫穿成為整個馬太福音裡面.
講述一個怎樣的福音.
耶穌是一個怎樣的救主.
當你知道一個這樣的結構之後.
去到八福.
耶穌所說的登山部分.
某程度是耶穌第一篇的講章.
在五章第一和第二節.
剛才我所說過.
耶穌上到山上.
就開口教訓他們.
來將天國的真理教訓給門徒.
以及那些坐植在山上的回眾.
同樣的經文.
如果你留意.
同樣的篇幅的內容.
在路加福音裡面也有記載.
不過路加福音在六章十七節的記載.
不是講耶穌上山那裡教導.
是耶穌下山.
在平原裡面.
教導那些門徒和群眾.
不是聖經裡面抄錯.
很多人說會不會是.
瑪太寫錯東西.
所以路加要糾正他.

$^{241}$應該在山下做的教導.
在山上是一個錯誤的地點.
你可以這樣理解.
不過如果我們純粹是在講.
聖經裡面講的錯誤的資訊.
是沒有什麼營養價值的.
我們看聖經.
有些靈修學家說不是瑪太寫錯.
瑪太所描述的耶穌的講論.
在這裡不是要強調一個正確的位置.
他在講上山就像一個靈修之旅一樣.
上到山上.
一群群眾特別要跟耶穌上山.
某程度是一群很渴望天國真理的群眾.
除了渴求耶穌為他們治病外.
更重要的是這群人.
他們很想聽到天國的真理是怎樣的真理.
他們上到山上.
跟隨耶穌和這一群門徒.
他們要聽聽天國真理是怎樣的事.
所以靈修學家說.
如果你用靈修的角度去看登山寶訓.
內容的詮釋就會變得很不同.
當我們今天去看這段經文的時候.
我們都是用一個.
耶穌帶我們上山.
耶穌幫我們重新去檢視一下.
脫離一下我們現有的形形異異的生活.
從一個好像生活忙忙碌碌的形態裡面.
上到山上領受一個新的教導.
或者改變我們思維的一些想法.
所以在這裡面.
如果你用這個角度.
你會發覺很有意思.
登山寶訓就是在說.
對我們生命的寧情裡面重新製造了一些新的體會.
你從這個角度去看.
再繼續看第三節.
虛心的人有福了.
因為天國是他們的.

$^{281}$在這裡面.
耶穌一直不斷地鋪陳說有福了.
有福的人.
一些不同的取向.
或者生命裡面作不同的操練的人.
這些人的生命是有福的.
原文裡面用有福這個字就是Makario.
希臘文.
你不懂希臘文不要緊.
你大概知道這些人是有福的.
不過在說那些有福的.
前面說那些什麼人有福呢.
不是說那些很有錢的人有福.
又或者不是說那些人很快樂就很有福.
如果你留意八福裡面說.
大部分所說的東西.
在當代社會或者我們今天這個時代.
我們都不覺得那些是什麼福氣.
怎麼會一個人虛心.
直譯其實是心靈貧窮.
覺得貧窮的人怎麼會開心.
又或者那些人不開心哀動的人.
怎麼會是一個有福的人.
在這裡面.
耶穌就用上山的一個靈修的角度.
他要顛覆一下我們究竟對福氣的一些看法.
可能我們今天社會上我們覺得累積有足夠的錢財.
又或者我們生活的際遇很好.
我們覺得那些是福氣.
我想我們都不需要否定這件事.
我們新年或者我們每一天.
我們都希望我們生活一切的際遇都是美好的.
不過耶穌他要給我們看見的那種福氣.
不單止是表面的際遇很平順.
他要給我們看到生命裡面更加隱含著更加大的一種的祝福.
那種祝福不是世人或者一般人看得見的.
是一種在一些我們覺得不是太好的經歷或者處境裡面.
我們發現上帝賜福的那種經歷.
有一個神學家叫做Martin Buber.
他解釋這段經文的時候他說.

$^{321}$生命的福氣其實往往很多時候是用隱藏的方式.
一種隱藏的幸福能夠幫助我們.
可以有力量來抵抗生命所有的不幸.
特別是在當代的社會裡面.
個個都是勞苦大眾.
個個都是面對著很多生命裡面的困境.
甚至耶穌說的這班群眾不是一班有錢人.
不是一班吃飽飯在摸酒杯底.
研究人生哲學的人.
不是.
耶穌對著的是一班勞苦大眾.
不過耶穌藉著說有福的講論.
說的是一個很偉大的講章.
他要將這班人的思維重新調教一下.
幫助他們看見生命裡面有些不是太理想的處境當中.
原來有一種隱藏的福氣.
是當我們經歷到的時候.
我們會發現我們生命裡面有力量抵禦生命所有一切不幸的遭遇.
這種的幸福就是真實的幸福.
所以我在大家的講章或是在螢光幕裡面看到一個字眼.
就是Macharismus.
就是另外一個希臘文的字眼.
和Macharioi這個字是近似的.
不過他所說的意思是說這種福氣不是一種抽象的福氣.
雖然我現在說得好像很抽象.
不過這種福氣是說你生命可以經歷的.
稍為你將你的思維改變一下.
你會經歷得到這一份的祝福.
所以幸福不是一個很遙遠.
不是一個人生終結的時候你才經歷得到的.
在今世裡面當你的思維重新去調教.
你明白上帝心意的時候.
你就能夠活在一種有福的處境當中.
耶穌就用有福了作為他登山的一部分.
也作為接下來他所講論的一個引言.
他很偉大的去宣告.
真的世界雖然有很多困苦有很多不幸.
但是當我們找到上帝.
找到那種隱含的生命裡面的祝福的時候.
當我們找得到的時候.

$^{361}$我們生命會經歷到這一份的幸福.
就是Mercuriousness.
就是可以經歷得到生命裡面那種圓滿.
和上帝所賜下的福氣.
另外一個字眼.
剛才也有輕輕講過.
就是虛心這個字.
虛心這個字是和合本所翻譯的.
當我們一讀虛心的時候.
我們會覺得謙虛.
虛心受教.
或者我們會覺得好像虛懷若谷.
這個譯法不是錯的.
他譯了某一程度的意思.
但本身來講.
那個字面的意思就是心靈貧窮.
心靈裡面覺得貧窮.
我們經常都是這樣的.
我大概如果在螢光幕面前問你.
你覺得哪個人家裡很有錢.
多有錢的人也好.
都覺得自己沒有的.
包括我在內.
你問我有沒有錢.
耶穌講的心靈貧窮.
那個虛心的意思不是在講.
你認自己身家有多少錢.
這裡是在講一種心靈的覺醒.
虛心就是在講.
我們知道.
我們自己.
我們真正生命裡面.
究竟處於一種什麼樣的境況.
稍後我會再仔細的解釋.
不過在這裡裡面.
特別的要讓我們知道.
這個字是顛覆了我們永遠一直想的.
多就等於福的意思.
反而在這裡裡面.
耶穌在講心靈自覺.

$^{401}$貴乏.
缺少.
這是一個福氣.
在字面裡面.
虛心這個字不只是在講一種謙卑.
或者是一種學東西的態度.
在這裡裡面.
他的意思更加要擴闊到.
一種心靈上的覺醒.
覺得自己有多少.
體會一下自己生命裡面.
活在一種什麼樣的狀態.
第三個字眼.
天國就是他們的.
天國在猶太人.
特別在.
馬太福音裡面.
如果寫給猶太人的信徒看.
天國這個字眼對於當代的人來說.
就會覺得.
耶穌掌權.
耶穌是做政治的君王.
以色列人復國.
講的是上帝的國道.
上帝的公義臨到.
有一個國界.
有一個.
國力的延伸.
這個感覺上.
讀法上.
如果天國用猶太人的讀法.
就是在講一個政治上的天國.
不過在這裡裡面.
耶穌要糾正那些.
法利賽人或者很多以色列人的想法.
因為曾經有些法利賽人.
問耶穌.
究竟有些人說天國在這裡.
有些人說天國在那裡.
究竟天國在那裡呢.

$^{441}$耶穌說.
天國不是在這裡或那裡.
而是在你的心裡面.
不是純粹在心裡面.
感覺上很良好就是天國.
耶穌再一步的.
來講.
天國是在講上帝掌權.
在哪裡.
容讓上帝作主.
那個處境.
就是天國.
一方面當然.
上帝將來會臨到.
這個世界.
將要臨到的一個天國.
就是上帝全盤的統治.
但同時間天國.
也在這一刻開始來發生.
就是當我們生命裡面.
稱基督為主.
上帝在我們生命裡面.
來掌管的時候.
天國就在我的生命裡面.
發揮果效.
當我們這一班人以耶穌基督.
的教導.
以上帝的心意.
以上帝作為我們生命的主權的時候.
這一班人.
就是成為天國.
所以在這裡面講的天國.
我們稍微要想一想.
不是在講一個國界.
而是在講一種上帝主權的彰顯.
你定義了這三個字眼.
有福.
虛心.
天國.
我們就要看看究竟這三者.

$^{481}$究竟在這裡面.
為我們帶來什麼福氣.
心靈貧窮.
在這裡面所講的心靈貧窮.
剛才我有講過.
是自覺自己.
一種覺醒.
覺得自己.
沒有,不足夠.
相對於盧家福音同樣記載著.
這一段經文.
十節.
耶穌舉目看著門徒說.
你們貧窮的人有福了.
因為神的國是你們的.
盧家福音就強調.
人的物質貧窮.
他特別強調.
耶穌是一個顧念窮人的主.
紀念那些生活很傑具的.
沒有錢的那些人.
那些人就有福了.
上帝要將福氣.
淋到他們身上.
盧家用這個框架來解釋.
耶穌所講的講論.
不過在馬太福音裡面.
耶穌的講論.
他特別將貧窮.
加上心靈兩個字.
如果當你這樣看的時候.
馬太就不是純粹.
在講物質上.
多與少的問題.
而是心靈裡面.
你體會自己究竟.
是多或是少.
是一種心靈裡面的一種自覺.
中國人.
或者保羅書信有時候也會講.

$^{521}$知足常樂.
我萬事都覺得滿足.
在這裡面.
心靈貧窮好像和知足兩個字.
好像相反.
不過其實在這裡面.
他要強調的.
在知足之前.
你要問問自己.
你究竟你的生命處於什麼狀況.
是一種什麼的狀況.
靈修學家古倫神父.
在這裡處理這段經文.
處理得非常好.
他說心靈貧窮.
有三個層面.
來看.
第一個層面就是.
你的人生一無所求.
你的人生的貧窮.
你心靈裡面的貧窮.
因為你生命裡面.
你沒有辦法.
和上帝再去討價還價.
當你求的時候.
你沒有條件和上帝交換.
你根本在你生命裡面.
認清主宰.
上帝的時候.
你不可以為你自己做一些功德.
來換取上帝.
有更加多的祝福在你生命當中.
人活在這個世界上.
其實一切.
我們和上帝.
是沒有議價的能力.
有時候有些弟兄姊妹.
在年頭的時候立志.
曾經我聽過.
有些年輕的弟兄姊妹.

$^{561}$我問他們有什麼立志.
他們說.
我希望上帝在這一年.
讓我的成績好一些.
讀書的成績.
各方面都有進步.
我說都不錯.
然後他們就交給他了.
如果他讓我成績進步.
我就會勤力讀聖經.
我說這麼奇怪.
為什麼要成績進步才讀多些聖經.
當然了.
這樣才能回報上帝.
讀聖經不是為上帝而讀.
是為你而讀.
上帝將話語賜給你.
其實你能夠有機會讀.
是一個福氣.
為什麼你會用一個這樣的角度.
和上帝好像條件交換.
然後這個年輕人.
再想.
不如我乖一點.
希望換取上帝.
給我多一些的祝福.
這是你自己的責任.
是你人生的本份.
你活得好.
到最後對你自己生命有益處.
這個你可以成為.
和上帝交換的條件嗎.
一邊數下去.
一邊數.
或者賺錢.
到最後一邊數的時候.
這個年輕人發覺.
牧師給我一條生路.
好像我什麼都不行.
我說不是.

$^{601}$不是我不讓你走生路.
而是看清楚我們自己生命的本相.
當我們去祈求的時候.
我們其實哪有.
和上帝議價的空間.
所有東西.
都是上帝給我們的.
我們不會用眼前.
我們覺得好的東西.
或者我們犧牲一點點的東西.
來換取更多的福氣.
我們在上帝面前.
根本就沒有.
議價還價的能力.
我們只是在上帝面前.
祈求恩典.
所以.
不要想著用你的功德.
用你的好的東西.
來換取上帝的祝福.
在地上的每一個人.
我們根本就沒有.
這種能力.
古倫神父提醒我們.
不要用一種討價還價.
條件交換.
你認清這個時間.
或者認清你自己生命的處境.
你就會發覺.
你其實就是.
所想的貧窮就是這個意思.
你心靈裡面有這份覺醒.
第二個層面.
一無所有.
承接著剛才所說的.
我們沒有辦法用功德.
來換取上帝更多的恩典.
甚至在我生命裡面.
所擁有的一切.
那些其實都不是.

$^{641}$我們所真正擁有.
我們在.
活在地上裡面.
生命所有一切的物質.
生命所有的一切.
擁有我們覺得很好的東西.
其實上帝只不過是.
託付給我們.
去管著.
那些東西不是屬於我們.
真正的擁有主權.
是來自上帝.
上帝不介意.
也很慷慨.
在我們這一生裡.
他將這一切.
付託給我們.
讓我們好像一個管家.
很好地去管理.
包括你的房子.
包括你的車.
包括旺角浸信會這個地方.
包括我們生命所有.
覺得我們覺得好的東西.
不好的東西.
這些一切.
其實我們無最終.
訴諸於那種終極的擁有權.
上帝只是.
託付這些東西.
給我們來管理.
他將這些.
我們看為好的事物.
他將這些物質.
或者我們覺得珍惜的東西.
他交付給我們.
目的不是要我們純粹.
自己滿足自己的喜樂.
他叫我們好好地管理的時候.
在過程可能我們會.

$^{681}$享受到某一些狀況.
帶來一些滿足.
更加重要的是.
這些一切要祝福身邊更多的人.
更加重要的是.
既然這東西.
主權不是屬於我的.
有一天可能上帝有主權.
要收回.
賜予的是耶和華上帝.
收取的也是耶和華上帝.
我們赤身而來.
我們來的時候沒有帶任何東西來.
我們有一天.
都會離開這個世界.
我們都不會帶任何東西.
回到天府.
在短促的這幾十年裡.
雖然好像累積了很多東西.
不過.
一切都是上帝.
託付給我們作為一個管家.
好好地來管理他.
當我們認識了.
我們醒覺.
我們生命只不過是一個管家.
我們對於物質.
我們對於很著緊的東西.
我們可以釋懷.
其實我們本來就是一個很貧窮的人.
我們赤身而來.
我們什麼都沒有.
只不過是上帝將.
不同的東西.
加給我們.
豐富我們.
給我們有機會來使用.
也讓我們在過程當中.
經歷.
使用這些東西的時候的那種喜樂.

$^{721}$不過有一天.
他會主權.
收回他.
所給我們的一切.
我們所擁有的東西.
當我們明白這件事的時候.
我們生命裡面可以.
減低那種執著.
那件事不是我的.
那個人不是虧負我的.
那個人如果真的.
老弱了.
或者令我感覺上有種不忿.
其實老弱了是上帝的事.
當然.
我們今天生命裡面.
我們也要好好地保存.
我們生命一切的擁有.
但是最終的擁有權.
我們知道.
不是屬於我們.
我們一無所有.
古倫神父再繼續第三講.
我們其實生命.
一無所知.
我們對於.
明天是怎樣.
我們也不知道.
耶穌說.
我們生命.
未來的日子明天如何.
我們也不能夠把握.
野地的花.
可能第二天就被摧毀.
就會消殘.
但是上帝也會將.
美麗的顏色.
賦予給這些野草和野花.
我們生命.
不是比這些更加寶貴嗎.

$^{761}$不過.
我們寶貴歸寶貴.
但是明天或者日後是怎樣.
我們其實真的想.
相信有算命.
或者預知未來嗎.
沒有的.
沒有一個人能夠知道未來是怎樣.
我們對於明天如何.
我們一無所知.
甚至對於這個世界所有萬物.
今天我們覺得對的東西.
明天可能就被推翻.
以前的人說.
吃雞蛋不要吃蛋黃.
吃蛋白就好了.
然後沒多久就有一個新的研究.
說不,吃蛋黃好.
不要吃蛋白.
有些人說整隻吃好一點.
有些人說吃半隻.
不停在說.
我們對於現象的世界.
都一無所知.
更何況是創造天地萬物的主.
我們其實對他的認識.
以我們所有知的.
人生的智慧.
我們沒有辦法.
知道上帝整全的面貌.
我們對於生命.
一無所知.
不過上帝說.
你放心.
將我們的生命交給他.
既然我們一無所知的時候.
我們不要以我們智慧.
引以為傲.
因為今天的智慧.
可能成為明天的愚蠢.

$^{801}$我們生命裡面.
沒有一樣東西.
是可以完全來吹捧的.
一無所知.
亦都類似和學本剛才所說.
我們變得虛心.
我們今天.
我們帶著.
我們有機會認知錯誤.
有機會我所知的東西.
可能是偏頗的.
這個心懷.
以致我們帶著.
一種這樣的虛心來接受真理.
當我們能夠有這樣的.
醒覺的時候.
我們生命.
就真的知道.
什麼叫做心靈上的貧窮.
耶穌基督說.
當我們體會到.
這種心靈的貧窮的時候.
我們這班人.
是有福的.
這個福氣不是將來.
而是當下.
因為當我們自覺.
原來我們今天所擁有的一切.
原來都是恩典來的.
當我們自覺.
我們根本就沒有辦法用功德.
來換取上帝的祝福.
當我們自覺.
我們只是想像中的這麼厲害.
當我們生命裡面.
看回我們自己的本相.
將我們打回原形的時候.
我們會發現.
生命其實有什麼.
值得我們去依戀呢.

$^{841}$有什麼值得我們.
很執著.
來到去.
就是經常都要抓住.
不放過這些人.
不放過這些事呢.
心靈的貧窮.
會讓我們重新將.
我們的心靈打回原形.
真的.
今天這個世界裡面很多人不開心.
很多人覺得.
很痛苦.
原因就是我們沒有辦法.
將我們的內心打回原形.
去回最原始的狀況.
當你想清楚.
其實生命裡面.
一切都是恩典來的.
這一切都不是屬於你的.
甚至你不是想像中這麼有智慧的時候.
你會發覺.
我還可以和上帝.
怎樣去強辯呢.
我生命裡面還有什麼.
和上帝可以討價還價呢.
心靈貧窮.
幫我們重新.
人生好像歸零一樣.
讓我們重新體會到.
自己生命的本相.
在這些本相的時候.
當我們經歷到這些.
體會.
發現了自己是這樣的時候.
我們就可以知足.
我們就可以知道.
生命擁有的一切.
原來真的.
超乎我所想所求.

$^{881}$我們生命裡面擁有.
在人生這幾十年的日子.
原來我們不是白過.
而是上帝用不同的恩典.
在我們這個本來.
貴乏的人生裡面.
添上大量的色彩.
大量的美好.
大量的恩典.
是由我自己努力換回來的.
一切都是來自上帝.
心靈貧窮.
讓我們心靈可以歸零.
讓我們真的看到.
我們生命的本相.
耶穌說不單止這班人.
將他們的心靈.
歸零的時候.
可以經歷到福氣.
耶穌繼續的對他們說.
他說這一班人.
天國就是屬於他們.
天國的意思我剛才也說過.
這個就是耶穌基督.
掌權的一種狀態.
不是純粹說地域.
或是幫國.
或是政治上的國家.
國土的概念.
而是說我們心靈裡面.
我們可以被上帝來掌管.
掌管.
今天很多人.
我們口裡可能在說.
主啊主啊.
但我們心裡其實是在說.
主啊.
祈求後面我們覺得很重要的東西.
我們不想失去.
所以我們就用力.

$^{921}$猛力的來求.
求天父成就我們.
我們懇切去求.
其實我們求天父給我們心目中想要的東西.
多過我們求天國.
求上主在我們心靈裡面.
來掌權.
但當我們能夠放下.
我們生命裡面.
我們覺得很重要的東西.
我們一切都是.
恩典的時候.
我們會漸漸發現.
原來生命裡面.
有一個主比我們.
生命裡覺得好的東西更加重要.
唯有賜予.
生命賜予.
厚此白物的主.
比起我們想要的.
那個才是.
我們生命裡面最重要的那個位置.
我們不要將我們.
心靈的位置.
被其他這個世界的物質.
或者人與事.
來佔據我們心靈的空間.
如果我們的心靈.
是被這些人與事.
佔領了的時候.
我們就好像衣亂了他們.
到最後我們就好像成為一個.
扮偶像的人.
我們心靈沒辦法經歷到.
從神而來的快樂.
相反來說.
發現自己的貴佛.
承認基督耶穌的主權.
我們的生命.
我們就能夠經歷到.

$^{961}$從神而來的那種臨在.
值得注意的是.
那個字眼.
耶穌說天國是他們的.
天國是他們的.
那個「是」在原文是在說.
不是未來式.
是現在式.
當你發現自己的貴佛.
發現自己.
生命裡面一切不屬於你.
你找回上帝的主權.
天國.
現在.
就成為在你生命裡面.
天國就屬於你.
你就進入在天國當中.
那個是一個.
現在式.
現在當下這一刻.
你就能夠經歷到天國.
不過很可惜.
在《馬達福音》裡面也有繼續說.
19章21至23節.
有些人和耶穌面對面.
會遇耶穌欣賞的.
有一個少年的官.
他問夫子我做些什麼.
可以承受永生.
耶穌說.
撇下你的一切周濟窮人.
你將來會有.
財寶在天上.
而且你就要撇下一切跟隨我.
和耶穌會遇.
而且耶穌是欣賞這個.
年輕有為的年輕人.
也問耶穌一條很有意思的問題.
但當耶穌發出這個邀請的時候.
這個年輕人.

$^{1001}$聽到耶穌這樣的呼籲.
他就悠悠愁愁地離開了.
為什麼呢.
《馬達福音》繼續說.
因為這個年輕人.
擁有的財寶實在太多了.
擁有地上的東西.
擁有太多了.
擁有所有人羨慕的東西.
不過卻是失去了.
即使和耶穌面對面的時候.
耶穌作出一個邀請.
他卻是在天國裡面.
他卻是失落的.
所以聖經裡面說到這個年輕人.
悠悠愁愁地離開.
不是耶穌.
要令他變得憂愁.
而是他所擁有的東西.
令他成為他生命的半腳石.
他不能夠輕省地.
跟隨這位主.
他不能夠經驗到天國.
相反對於貧苦大眾.
可能他的機會成本.
opportunity cost很低.
放低就放低吧.
他反而很輕省地.
他就可以進入.
基督耶穌的邀請的國度裡面.
財主進入天國是困難的.
財主未必是在說.
家財萬貫.
財主可能是在說.
我們固執.
我們覺得我們很重要的東西.
我們就拿著.
執著我們很想擁有的東西.
那些東西不能夠失去.
我們可能.

$^{1041}$帶著一份那麼強烈的.
那種心.
即使我稱呼我主.
背後我其實就是在說.
主啊,你不要叫我放棄.
這些東西.
生命如果我們固執.
沒有辦法經歷.
天國而來的.
那種滿足.
那種隱藏.
但其實是一個很.
歸零的那種幸福.
相反來說.
願意放手.
承認自己生命.
一切都是恩典.
容讓上帝帶領自己人生.
去走的時候.
我們可以豁然開朗.
稱基督耶穌為主.
不單止口說.
祂救贖我.
將我從罪裡拯救出來.
稱基督為主.
更加是我生命的主.
我願意放棄.
我覺得生命裡面很重要的東西.
一路一路放棄.
將位置放下.
我將基督放在我生命裡.
最重要的位置.
當你一路一路輕散.
你會發現.
當你一路一路輕散.
放下這些生命裡面.
纏繞著你的好東西.
好的重要的東西.
放下這些.
依戀和牢牢.

$^{1081}$慢慢你以基督的心意.
以耶穌基督的心懷.
成為我生命裡面.
最重要的關注的時候.
我們一路一路.
就會覺得.
生命是很輕散.
其實生命可以過得很簡單.
我們不需要為著.
我們所強求的東西.
再繼續的.
來執著.
心靈貧窮的人有福了.
因為天國.
是他們的.
這句說話就是這個意思.
我早前有機會.
看了一個很好的故事.
台灣作家叫許榮哲.
寫了一個很精彩的故事.
我節錄一部份.
因為故事未完.
話說有一個人.
叫阿甲.
阿甲要上路.
他帶著四塊餅.
一路趕路.
一路去到一個目的地.
在他趕路.
走這條道路的過程當中.
他遇見.
和他一起走同一個方向的阿悅.
阿甲帶了四塊餅.
給他作為.
在這條路程裡面.
來充飢的餅.
很平凡的一個餅.
不值錢的.
但帶著這四塊餅一路走.
在路上遇到阿悅的時候.

$^{1121}$阿悅帶了兩個餅.
大家因為.
同一條路徑和一個方向.
去某個村莊.
阿甲和阿悅就認識了.
大家很投契.
一路聊的時候很開心.
到一路入夜的時候他們就坐下.
大家打開自己的行李.
拿自己的餅出來吃.
阿悅只有兩塊餅.
阿甲就說.
不如這樣吧,大家同路,大家這麼投契.
一起分餅吃.
於是大家就.
把六塊餅兩個人.
打開一起的來吃.
當他吃開的時候.
然後又見到一個人.
氣沖沖的又趕路.
因為入黑了,他又.
入村莊又路又不夠.
而且那個餅就說我沒有帶餅.
阿甲和阿悅在那裡吃.
於是阿餅就走過來.
他就說可不可以和他一起.
因為幾塊餅乾而已,算不得什麼.
所以阿甲和阿悅就.
一起在路邊.
享受這六塊餅.
吃完之後.
大家過了一晚夜,很開心.
經歷過一個.
很快樂的晚上之後,第二天.
繼續來趕路.
當他趕路的時候.
阿悅就說其實我走到這裡.
我都要和大家分離.
我要走東邊.
你們繼續走西邊的方向.

$^{1161}$於是和他們道別.
在道別的時候,阿餅就說.
雖然我沒有帶我的餅乾.
但我手頭上有六塊金幣.
作為多謝你們.
兩個接待我.
昨晚一起.
可以吃到這麼豐富的晚餐.
這麼美好的.
能夠和你們一起.
我就覺得很開心,於是將這六塊金幣.
送了給.
甲和悅.
金幣和餅的價值相差很遠.
當.
拿到這六塊金幣的時候.
甲和悅一路一路.
繼續趕路.
走的時候.
因為金幣的價值實在太豐富.
甲和悅就.
去研究究竟這六塊金幣.
要怎樣分呢.
究竟是.
每人三個這樣分.
還是按著有多少個餅.
甲就有四個餅.
就拿四個金幣.
悅就拿兩個金幣.
大家去爭論拿多少個金幣的時候.
走著走著.
就吵架了.
差不多進到村莊口的時候.
大家甚至乎.
咬到臉紅耳赤.
差不多要大打出手了.
悅在這個時候就覺得很委屈.
為什麼.
一定要分少一點呢.
甲又會覺得為什麼這麼貪心.

$^{1201}$分多一點呢.
於是就氣沖沖的.
想打甲.
但是在混亂的過程當中.
甲和悅就有人來勸他們.
然後就找一個智者.
來到他們的頻道.
悅就氣沖沖的.
一個人就走去找那個智者.
找那個智者就.
告訴他剛才那件事.
整件事應該怎樣分.
數字應該怎樣算.
我應該分三個金幣.
甲又說分兩個.
特意這樣.
剋扣我一個金幣.
然後這個智者聽完整件事之後.
然後他就問.
悅.
如果六個餅.
你們三個人一起吃.
每個人平均應該吃多少個餅.
那.
悅就說按著計數.
每人就吃兩個.
智者就說嗯.
對了 那你本來帶多少個餅.
來趕路的.
那悅就說.
我就帶兩個.
那智者就跟悅說.
喂 你真是貪心.
你吃你自己那兩件餅.
你還想拿別人的金幣.
當悅聽到他講這番說話的時候.
突然間醒一醒.
對了.
我吃我自己那兩件餅.
如果沒有這六個金幣的出現.

$^{1241}$我和甲已經是本來好朋友.
但為什麼六個金幣之後.
我和他反而會大打出手.
會成為敵人呢.
這個智者就將悅的想法歸零.
你吃你自己那件餅.
為什麼要.
何苦這樣.
來糾纏.
要去分這些金幣呢.
這個故事還沒完.
有機會我借好本書給大家看.
但我告訴你.
這個故事.
我借了一本書給大家看.
不過在這裡.
這個智者講的.
那個很重要的.
讓我們知道.
我們生命裡面歸零.
看回其實.
我們生命擁有什麼.
我們苦苦糾纏.
我們.
經常覺得很痛苦的東西.
我們覺得失去.
其實是不是真的失去呢.
還是.
我們很執著.
那些我們覺得很重要的東西.
我們沒想過.
有一天.
那個主人會拿走.
可能.
我們人生裡面覺得很多痛苦.
我們覺得很.
很折磨.
可能就是被這些我們覺得很重要的東西.
牽緊我們的內心.
以致到我們的心靈不平安.

$^{1281}$但是.
今天這段經文提醒我們.
看清楚我們自己生命的本相.
我們一無所有.
我們一無所求.
我們更加是一無所知.
我們一無所求.
我們更加是一無所知.
我們沒有能力用我們的功德和上帝討價還價.
我們只不過是受託的管家.
我們更加對明天.
對我們人生一無所知.
所有的東西.
都是基於上帝.
找回生命最重要的源頭.
找回生命的主.
我們生命就可以輕省.
幸福就由這個.
心靈貧窮開始.
求主繼續幫助我們.
在新的一年裡.
讓我們生命可以歸零.
找回自己的方向.
找回自己生命裡.
最重要的東西.
幸福 隱藏的幸福.
由這一種心態開始.
願主繼續給予我們祈福 眾人.
\newpage



\section{馬太福音 5:4}
\label{sec:POKxDUqFlWw}
\textbf{2021年1月10日崇拜直播｜陳冠成牧師｜幸福由哀慟開始｜馬太福音5\_4}
\newline
\newline
連結: \href{https://youtube.com/watch?v=POKxDUqFlWw}{\texttt{https://youtube.com/watch?v=POKxDUqFlWw}} ~~~~ 語音日期: 2021-01-12
\newline
\newline
\hyperref[sec:bLhELYX_Uw0]{\small{< < < PREV SERMON < < <}}
~
\hyperref[sec:index_chronic]{\small{[返順時目]}}
~
\hyperref[sec:YMwn1_GRZYg]{\small{> > > NEXT SERMON > > >}}
\newline
\newline
馬太福音 5:4
\newline
\begin{longtable}{cl}
\hline
\hline
章節 & 經文 (和合本修訂版)\\
\hline
5:4 & \begin{tabularx}{0.7\textwidth}{X} 哀慟的人有福了！因為他們必得安慰。 \end{tabularx} \\ \\
[1ex]
\hline
\hline
\end{longtable}
$^{1}$讓我們一同以聆聽聖道去敬拜我們的主.
今天我們會重讀的經訓在.
《馬太福音》五章一至十節.
(試試書).
耶穌看見這許多的人就上了山.
記耳坐下門徒到他跟前來.
他就開口教訓他們說.
虛心的人有福了.
因為天國是他們的.
哀動的人有福了.
因為他們必得安慰.
溫柔的人有福了.
因為他們必承受地土.
饑渴無意的人有福了.
因為他們必得飽足.
連戌人的人有福了.
因為他們必蒙連戌.
清心的人有福了.
因為他們必得見神.
使人和睦的人有福了.
因為他們必稱為神的兒子.
為易受逼迫的人有福了.
因為天國是他們的.
今天陳冠成牧師宣講的經文.
是在《馬太福音》五章第四節.
哀動的人有福了.
因為他們必得安慰.
今天所宣講的主題是.
「幸福由哀動開始」.
請陳冠成牧師.
弟子們平安.
剛才珍珍牧師讀出了.
八福這段經文.
耶穌祂很期望天國的子民.
可以在生活中實踐.
例如謙卑,溫柔,締造和平,連戌人等.
這八個美善的品格.
並且耶穌說.
你這樣做是蒙福的.
不過當我們看今天福氣時.

$^{41}$耶穌說哀動的人有福了.
聽起來會否覺得有點奇怪.
一方面哀動和使徒保羅所強調.
要上上起落.
這個教導好像格格不入.
到底一個哀動的人.
哀動的人如何經歷福氣呢?.
又或者一個人明明哀動.
為何我們說是一種幸福呢?.
到底是不是上帝很喜悅.
祂的子民終日都很悲傷.
很憂愁地過生活呢?.
耶穌說哀動的人有福了.
到底是什麼意思呢?.
在這裡我們不妨先看兩段.
和哀動有關的經文.
第一段記載在詩篇的三十八篇裡.
可以在螢幕播放這段經文.
謝謝.
詩人說:我幾乎跌倒.
我的痛苦常在我面前.
我要承認我的罪孽.
我要因我的罪而憂愁.
耶路說:求你不要撇棄我的神.
求你不要遠離我.
拯救我的主.
求你快快幫助我.
詩人因為自己罪孽的緣故.
哀求上帝來憐憫.
來拯救他.
而另一段在以賽亞書的第四十章裡.
說:要對耶路撒冷說安慰的話.
又向他宣告說.
他征戰的日子已滿了.
他的罪孽赦免了.
他為自己一切的罪.
從耶和華手中加倍受罰.
先知以賽亞某程度上.
對耶路撒冷來說安慰的話.
原因是他們過往的罪孽.

$^{81}$已經被上帝赦免.
他們因為犯罪而被上帝管教.
甚至要被捕.
這個日子將來也會結束.
上帝要重新挽回他們.
所以這兩段經文讓我們看到.
哀悼是和罪惡有關.
哀悼並不是因為犯罪.
而招致自己的利益受損.
這麼簡單.
不是因為這方面而感到哀悼.
譬如說.
有一個學生假如他考試作弊時.
被發現.
學校要記他一個小過.
他因為要受罰.
他很擔心可能在同學面前會失去面子.
又或者被家人知道會罵他一頓.
所以他流淚痛哭.
他很希望因為這個緣故可以重慶發樂.
但卻沒有為自己的過犯痛悔.
雖然他流淚痛哭.
但這不等於他哀悼.
又或者我們熟悉的士師記.
他記載以色列人上上.
因為外敵的踐踏.
嚴重影響他們的生計.
打擾他們的生活.
他們感到很懊惱,很困擠.
哀求耶和華幫助他們擊退外敵.
但事過境遷後.
他們又再姑太復萌.
繼續拜偶像.
不聽上帝的吩咐.
沒有正視或根治他們罪的問題.
沒有根治他們一再踐踏上帝吩咐的惡行.
這種我們說因為自身利益受到損害.
而向上帝發出的哀求.
亦不是耶穌所說的那一種哀動.
耶穌所說哀動的意思是.

$^{121}$一個人為自己或世人的過犯而感到痛悔.
並且知道人犯罪最受傷害的.
不是自己或別人的損失.
而是上帝他自己.
他很明白人犯罪是虧缺了上帝的榮耀.
辜負了上帝對人的期望.
人犯罪傷透了上帝的心.
所以哀動的人.
他不會再希望上帝繼續傷心下去.
他會很希望自己或別人能夠回轉.
來得到上帝的寬恕和接納.
就好像我們很熟悉大衛的一篇詩篇.
《詩篇51》大家也有印象吧.
《詩篇51》是大衛在犯奸淫後.
被先知拿丹斥責而所寫成的一篇詩篇.
當中我們看到他如何表達他為罪而哀動.
這裡說.
「神啊,求你按你的慈愛連恤我.
求你豐盛的慈愛塗抹我的過犯.
求你將我的罪孽洗除淨盡.
並潔淨我的罪.
因為我知道我的過犯.
我的罪常在我面前.
我向你犯罪唯獨得罪了你.
在你眼前行了這惡.
以至你責備我的時候顯為公正.
判斷我的時候顯為清正」.
大衛知道他不單傷害了人.
不單傷害了拔士巴.
或拔士巴的丈夫烏尼亞.
他知道他最得罪的是上帝.
是他親手破壞了和上帝美好的關係.
他知道上帝有權審判自己的不義.
將他消滅.
但當我們繼續看這篇詩篇時.
我們會看到大衛很明白一件事.
他認識到上帝有恩典有憐憫.
並且他在這裡說.
神願意看重那些為憂傷痛悔的人.
為罪憂傷痛悔的人.

$^{161}$所以大衛認罪悔改.
他求神塗抹赦免他的過犯.
求上帝不要丟下他.
並且求上帝幫助他做正直的人.
結果我們看到.
大衛的哀動蒙上帝他越立.
上帝赦免他.
給他一個機會改過自新.
大衛正正體驗耶穌基督所說.
哀動的人有福了.
因為他們必得上帝所給他們的安慰.
事實上「安慰」這個字.
不是拍拍對方肩膀.
給對方一個擁抱.
然後說一些很呵護的話.
「沒事了,不要哭了」.
不是這樣說.
「安慰」其實有意思是說.
得力量.
並且「安慰」這個字.
在希臘文裡和「保衛施」是同一個字根.
給我們看到凡願意為罪而哀動的人.
上帝一定會透過聖靈保衛施.
賜下應許.
賜下安慰和力量.
首先就是.
當人為罪哀動的時候.
他會得到上帝赦罪的保證和應許.
我們很熟悉約翰一書的一章九節.
這裡說.
「我們若認自己的罪.
甚是信實的.
是公義的.
必要赦免我們的罪.
洗淨我們一切的不義」.
我們可以和上帝重新和好.
罪亦不再在我們和上帝中間.
重新作梗的話.
這個應許本身已經是極大的安慰.
極大的福氣.

$^{201}$而另一方面.
即或我們有一些自身的軟弱問題.
仍然未能徹底解決.
又或者你看到.
世上仍然有很多.
不法不義的事情.
每天都發生.
但是聖靈會安慰我們.
提醒我們有幾件事.
第一件事就是.
上帝沒有走開.
上帝仍然看顧我們.
掌管這個世界.
上帝亦不會任由這些惡人惡事.
橫行猖狂.
雖然有很多狀況.
是不合上帝的心意.
是上帝不肯.
是上帝不起穴.
不過.
這個世界仍然在上帝的掌管中.
這個仍然是天賦的世界.
萬事萬物最終都要向上帝的計劃發展.
而在上帝這個計劃中.
他可以瞬間消滅罪惡.
他可以阻止惡人對我們的加害.
當然有時上帝會容許.
一些禍患出現在我們中間.
沒有即時消滅他們.
但是上帝卻會在當中.
加添你和我的力量.
來面對這些艱難.
面對這些惡人惡事.
上帝有時亦會運用他奇妙的大能.
來使用那些逆境.
那些人的罪惡.
反過來為上帝效力.
成就他們一些超越的旨意.
事實上.
聖靈提醒我們.

$^{241}$上帝原本的心意.
是要祝福我們每一個.
雖然罪惡,苦難.
仍然在我們身邊出現.
但是上帝加給我們的安慰和力量.
足夠勝過這一切.
這是祂給我們的保證和安慰.
所以弟兄姊妹.
當耶穌在今天一個.
我們說充滿罪惡,不信的世代裡.
祂跟我們說.
哀動的人有福了.
因為他們必得安慰.
你應該知道.
耶穌不是在說一些離地的話.
祂是在指明一條真正的夢福的道路.
祂要我們相信.
當我們真的起來.
為身處罪惡世界哀動的時候.
是真的會經歷福氣.
所以.
有一件事我們要留意的就是.
不哀動.
不為罪而哀動.
無論是你的罪.
或世人的罪而哀動的話.
不是說你少了福氣這麼簡單.
為什麼這樣說呢?.
你試試反轉這句話讀一讀.
你會明白.
耶穌說哀動的人有福了.
因為他們必得安慰.
換句話來說.
不哀動的人會有什麼?.
有禍了.
為什麼呢?.
因為他們得不到上帝的安慰.
你說得大不大?.
得不到上帝的安慰.
其實是禍.

$^{281}$所以弟兄姊妹.
我們首先需要為自己的罪哀動.
為過往我們和上帝的疏離.
可能是不聽祂的吩咐.
可能是明知故犯.
來痛悔和改變.
不單止承認我們得罪了上帝.
我們需要有新的行為.
我們需要改變自己.
證明我們的痛悔是真心的.
事實上改變包含了兩個層面.
首先.
我們需要決心離開惡道.
不再亂棧在過往罪的生活當中.
不再明知故犯.
做一些上帝不喜悅的事情.
我最近聽到一首詩歌.
不知道你有沒有聽過.
叫做《早知道》.
是一首普通話的詩歌.
頗有意思的.
一開始是唱副歌.
副歌是這樣說的.
早知道.
其實我也早知道.
你話語曾經教導.
早知道.
其實我也早知道.
你屢次提醒勸告.
其實寫出了.
世人那種明知故犯.
那種無奈軟弱的光景.
接著這首歌會唱出幾句很責心的歌詞.
早知道驕傲會讓我跌倒.
我就會向謙卑倚靠.
早知道謊言會讓人離我遠去.
我就立志誠實道路.
早知道荒蕩會讓我被痛苦纏繞.
我就不敢再次向他擁抱.
早知道叛逆會讓我無處可逃.

$^{321}$我就從開始回轉就好.
來到結尾的時候.
這首歌會拖慢節奏地唱出最後一句.
只是我到現在才真正知道.
有種唏噓 有種酸的感覺.
但我想卻唱出了很多人的心聲和現況.
我們早就知道.
我們明知道.
譬如歌詞所說.
驕傲 謊言 荒蕩 叛逆.
並不是好事.
聖經裡也強調.
叮囑我們要遠離這些事情.
但人有時真的很奇怪.
人是喜歡偏行己路.
明知不可為而為之.
當我聽到這首歌詞時.
回想起一句說話.
我會用來警惕自己.
或勉勵我的兒子.
就是在上帝未出手管教你之前.
你自己先好好約束自己.
在上帝未出手管教你之前.
你先好好約束你自己.
帕默這句說話成為我們彼此的勉勵.
真的弟兄姊妹.
上帝想我們及早對付生命中的陋習.
那些私慾 軟弱 罪惡.
上帝很想我們靠他決心離開惡道.
免得到時他出重手管教的時候.
我們真的會後悔莫及.
事實上耶穌已經為我們開闢了救恩的道路.
你和我都知道.
但耶穌不能替你和我悔改.
那是我們的責任.
那是我們當盡的份.
至於改變除了離開惡道之外.
其實還有另一個更加積極的層面.
就是你要有好的行為.
就是你要進行聖經的真理.

$^{361}$就是你要靠著上帝的話語.
為上帝結出果子來.
詩篇第一篇我們很熟悉.
為喜愛耶和華的律法.
就耶思想 者人便為有福.
他要像一棵樹栽在溪水旁.
安時後結果子.
葉子也不枯幹.
凡他所做的盡都順利.
我們很記得這句話.
如果你更加記得的就是.
耶穌在曠野面對撒旦時.
祂怎樣做?.
祂用上帝的話語.
來擊退撒旦.
升哥撒旦的試探.
同樣作為耶穌的跟隨者.
上帝的話語對我們很重要.
事實上詩篇這裡其實.
間接提醒我們.
最有效遠離罪惡的方法是什麼?.
不是你千方百計避開.
不做這樣 不做那樣.
你做一樣東西便可以.
你喜愛上帝的話語.
你讓上帝的話語充滿.
你自然就遠離罪惡的試探.
情況就好像.
不知道有沒有兄弟姐妹.
你喜歡看劇集.
或者喜歡玩遊戲機.
你喜歡看劇集 喜歡玩遊戲機.
喜歡到一個地步.
就夜去思想.
怎樣就夜思想?.
你上班.
你吃飯.
你逛街.
甚至你做夢的時候.
你還在想.

$^{401}$怎樣打機.
怎樣煲劇.
腦裡充斥著畫面和想法.
你還有沒有時間去想其他東西?.
大概可能沒有.
你將存在的時間精力.
放在你喜愛的事情上.
你其實沒有空間.
去再想再做其他東西.
同樣這段經文教導我們.
當我們願意不斷被上帝的話語.
充滿的時候.
無論是什麼境況.
無論你身處什麼生活崗位.
你做什麼都好.
你都會自然地想一想.
我怎樣可以在裡面.
實踐聖經真理呢?.
我怎樣可以為上帝多走一步.
來討祂的喜悅呢?.
這樣的話.
你和我還有什麼多餘的時間和精力.
去想那些不太好的.
那些會得罪上帝的事情.
是嗎?.
你想想.
所以弟兄姊妹.
我們可以檢視一下.
在這段日子.
我們的讀經.
我們靈修.
我們親近上帝的質素.
到底是怎樣呢?.
特別在疫情下.
我想我們少了很多消遣活動.
我們多了很多時間留在家裡.
我們怎樣去善用呢?.
我們有沒有好好親近上帝.
甚至認真地.
去研讀祂的話語.

$^{441}$讓聖經的真理塑造我們的生命.
來結出聖靈的果子.
事實上我知道在這段日子.
很多弟兄姊妹.
他們回到聖經裡面.
他們從聖經裡面.
找到生活的亮光和出路.
譬如說有些人.
他在聖經裡面找回他的照明.
他願意去回應上帝給他的感動.
為他所看見的.
那些他覺得不公道.
不公平的事情.
來發聲.
他求上帝來念文.
幫助在當中受傷害的人.
他也求人在當中能夠回轉悔改.
雖然在過程裡面.
他很體會到.
別人對他的忠言逆耳.
那一種失望.
他很體會到.
他不斷地呼求.
但是環境一直不去改變.
那份無力感.
那份無助.
但是他沒有放棄上帝給他的照明.
因為他知道上帝沒有放棄.
上帝沒有放棄人.
沒有放棄這個世界.
上帝仍然掌管著大局.
所以他仍然願意.
繼續堅守他的崗位.
願意去順服.
仰望上帝一切的安排.
有些弟兄姊妹在上帝的話語裡面.
他得到醫治.
就好像詩篇裡面的詩人.
他會將自己的苦情.
將自己的困局.

$^{481}$一一來向上帝傾訴.
他知道困難逆境.
總有一天會過去.
上帝會將這一切都撥亂反正.
所以他不急於來求上帝.
盡快去解決他自己的困局.
他不急於來求上帝.
消滅眼前一切的罪惡困難.
他反而靜候下來.
祈求上帝給他一個平靜安穩的心.
讓他可以有力量.
來生活.
和逆境共存.
亦有人在這段時間裡.
透過上帝的話語.
他重新和人聯繫.
他比以前更加珍惜.
和家人和朋友相處的時刻.
他會去關顧他的朋友.
他的家人.
特別是在這段日子裡.
他們身心靈的需要.
他亦知道在這些日子裡.
他可以為上帝多走一步.
就是留意身邊的鄰舍.
需要什麼.
並且他願意身體力行地.
代表耶穌來關顧幫助他們.
因為他知道.
在這些小生活裡.
也有上帝命定的福氣在當中.
等我們每一個人去領受.
和傳揚出去.
所以「弟兄姊妹攀望」.
在新的一年裡.
我們真的要用更多的時間和心力.
來好好研讀.
好好沉澱思考.
並且想一想如何實踐上帝的話語.
從而我們真的可以得到.

$^{521}$幾方面的能力.
得到勝過試探的能力.
得到好好管理我們的心思.
我們的情緒的能力.
亦得到行善結果子的能力.
使我們可以在一些動盪不安的日子裡.
不單止我們自己首先站穩.
並且我們可以去祝福身邊的人.
讓上帝在我們的生命裡得著榮耀.
最後.
不單止我們要為自己的罪哀動.
更加要為世人.
為世界的罪哀動.
我知道有時不容易.
特別當.
可能我們也有過這些情況.
當你看到世界的罪案發生.
當你看到身邊的人犯罪.
你可能會出現兩個即時的反應.
第一個就是你很憤怒.
特別是當你的利益受損的時候.
你很想馬上去還擊.
你很想指罵對方的惡行.
甚至你去咒詛他.
又或者另一個好像鐘擺一樣.
你會去到另一個極端.
就是你很冷漠.
你已經沒有反應.
你當作看不到.
繼續生活.
反正也與我無關.
我又沒有能力去做什麼.
但是我們不妨細心想.
這是否合乎耶穌對門徒的期望?.
話說有一個故事是這樣說的.
有一個已故的拳王.
拳擊高手.
叫做喬路易斯.
有一次他與他的朋友開車出外的時候.
在路上不小心.

$^{561}$與別人的車發生了碰撞.
對方下車.
不理會三七二十一.
狠狠地罵了他們一頓.
羞辱了他們一番.
路易斯的朋友非常憤怒.
當對方離開之後.
他就氣沖沖地.
去質問路易斯.
他說:喂!.
你是拳王.
你的拳術這麼好.
為什麼你剛才不打對方一身?.
為什麼你剛才不還擊?.
對方做得不對.
路易斯的回答不單風趣幽默.
並且很有涵養.
他說:如果有人侮辱.
歌王卡羅素.
卡羅素是誰呢?.
是當時很著名的意大利男高音.
他說:如果有人侮辱歌王卡羅素.
你猜卡羅素會不會為侮辱他的人獻上一曲?.
會不會為他唱一首歌?.
這是一個很有智慧的回應.
聽姐妹.
下次.
當我們看到別人的過犯.
當我們遇到不公義的對待.
當我們很生氣.
很想第一時間還擊.
很想痛罵對方的時候.
盼望這個故事.
也讓我們想起一些事情.
提醒我們.
你的生命是用來做什麼?.
你的生命是用來榮耀上帝.
上帝給你各樣的恩賜財幹.
是要詛咒別人.
而不是用來報復.

$^{601}$用來攻擊.
或者計算別人的惡.
謹記上帝給我們各樣的才華.
是要用在討上帝喜悅的事情上.
千萬不要因為別人的犯錯.
因為你看到一些惡人惡事.
就輕易挑動到你的情緒.
或者說.
我們也懂得一句說話.
不要為了別人的過犯.
而懲罰你自己.
不要讓自己的怒氣.
牽著你的鼻子走.
你要讓聖靈帶領你.
你要讓上帝的話語掌管你.
更不要以惡還惡.
讓自己不知不覺變成惡人.
失去了上帝的福氣.
得罪了上帝.
來到最後.
耶穌不單要求門徒.
要為世人的罪哀悼.
耶穌有沒有親身示範?.
耶穌有.
你看福音書.
耶穌怎樣親身示範.
為世人的罪哀悼.
首先你看到耶穌.
祂沒有置罪人而不顧.
祂主動接觸他們.
祂和他們交往.
為他們禱告.
甚至身體力行到一個地步.
是怎樣呢?.
祂願意為他們寫明.
祂願意為那些得罪他們的人.
來求上帝寬恕,赦免.
因為祂知道他們不曉得.
他們不知道其實是得罪了上帝.
他們不知道自己正在走一條滅亡的路.

$^{641}$所以上帝會憐憫罪人.
耶穌會為世人.
包括你和我.
預備救恩.
而現在.
耶穌吩咐我們.
你要去將這個好消息傳出去.
使更多的罪人.
好像適應你和我一樣.
回頭.
悔改回轉成為天國的子民.
經歷上帝的福氣.
提醒我們.
不要再心硬.
不要再採取視而不見.
或者冷眼旁觀的態度.
來看待你身邊的人的過犯.
或者世人的罪.
經文提醒我們要學習耶穌.
為罪人來禱告.
最初可能真的不太習慣.
看到別人的過犯.
看到世人的罪.
要我們去學習.
為他們禱告.
很難的.
有時你心裡有另一個律就會出來.
特別當那些人很討厭的時候.
你最想看到他們怎樣.
我們很想看到惡人.
看到罪人跌倒.
我們很想看到那些壞人怎樣.
回不了身.
這是我們另一個律.
但當我們看到聖經裡說.
上帝不願意一人沉淪.
乃願人人悔改的時候.
就挑戰我們.
要有另一個律.
來管理我們的心.

$^{681}$就是表示我們需要和耶穌同行.
來效法他的心腸.
渴望上帝的旨意.
在地上可以彰顯.
看到那些很討厭的人.
都有回轉悔改的機會.
和我們一樣.
最終可以活在救恩的裡面.
所以假如.
如果你覺得為罪人禱告很困難的話.
可以先為你自己禱告.
可以先求上帝來安慰你.
求祂加力量給你.
讓你可以效法耶穌.
求祂幫助我們.
當我們真的願意為罪人.
為世人禱告的時候.
為他們的罪而哀動的時候.
我們真的可以經歷到.
經文這裡所說.
從上帝而來的那份依治和安慰.
而除了為罪人禱告之外.
有一件事我們不能忘記.
無論任何境況.
你都需要實踐耶穌頒佈的呼音使命.
既然世人最大的罪.
是不認識真神.
不願意相信跟隨他的話.
最直接的回應方式是什麼?.
就是你和他傳呼音.
就是你透過呼音工作.
改變他的心.
改變世人的心.
帶領他們信主歸向神.
事實上.
一個願意為世人哀動的基督徒.
他最終一定會邁向一條道路.
就是他會慢慢樂意參與呼音工作.
他會喜悅看到罪人悔改這一幕.
正如耶穌喜悅看見這一幕一樣.

$^{721}$你很熟悉的失淺失揚的比喻.
你很明白.
當有罪人悔改.
誰最開心?.
天上的父最開心.
你能不能夠喜悅耶穌所喜悅的呢?.
關心耶穌所關心的呢?.
是經文給我們的挑戰和考驗.
最後很想分享昨天的一個經歷.
不知道大家有否看過明珠台.
有一套《2012世界末日》.
這套片看過N次.
經常播放.
播到我都熟悉.
當然你知道2012世界末日沒有發生.
但我今次看的時候.
有一段說話很觸動我.
如果你記得.
那裡有個演員扮美國總統.
他原本打算乘飛機到救生船逃難.
他可以在世界末日生還.
但他最後選擇不走.
他最後選擇把世界末日的訊息說出去.
但其他國家的元首早已說不要說.
因為一說出去就恐慌.
因為救不了所有人.
但美國總統說.
確實我不能救全世界的人.
但我仍然覺得應該跟他們說出真相.
讓那些無法乘船的人.
仍然有時間跟家人道別.
仍然有時間擁抱家人.
仍然有時間彼此要恕.
這段說話很觸動我.
同樣我們無法解決很多事情.
我們自己救不了多少人.
但我們可以把真相,真理告訴他們.
我們無能力撥亂反正的事情.
誰有能力?.
誰可以徹底把他們完完全全拯救出來?.

$^{761}$是你和我相信的上帝.
我們的主耶穌基督.
所以弟兄姊妹盼望.
讓今天上帝的話語.
重新來教訓我們的天國觀.
讓我們真的能夠緊緊跟隨耶穌.
這位天國君王的腳蹤.
來活出祂所喜悅的八福美善品格.
使你和身邊的人.
真的可以經歷上帝同在的福氣.
好嗎?.
讓我們同心禱告.
主耶穌說幫助我們.
幫助我們每一個.
說要跟隨你.
相信你的門徒.
我們真的擁有你的心腸.
我們真的願意為罪而哀動.
我們願意為自己的過犯.
來到上帝的面前.
真誠的悔改.
因為我們真的得罪了上帝.
傷透了上帝的心.
上帝我們知道.
你有恩典有憐憫.
你願意赦免我們的過犯.
不單如此.
你更加會幫助我們改變.
讓我們真的可以離開惡道.
讓我們願意實踐遵行.
你寶貴的話語.
使我們自己和身邊的人都能夠蒙福.
上帝又幫助我們.
讓我們看見別人的罪.
看見世界的罪的時候.
我們不要冷漠.
我們真的願意有耶穌基督的心腸.
來為這一切禱告.
來為這一切求上帝.
你幫我們懂得怎樣去回應.

$^{801}$讓我們真的有耶穌基督的視野.
讓我們看到上帝.
你對罪人的心意.
正如你對我們的心意一樣.
你不希望一人沉淪.
哪怨人人都悔改.
主也以至我們真的能夠.
懷著憐憫的心腸.
來為罪人禱告.
來渴望他們能夠接受福音.
能夠迎信你.
回轉悔改.
並且看見.
上帝在當中的喜悅.
天父上帝為此讓我們.
不要忘記福音的使命.
讓我們在任何時候.
我們都能夠.
無論得時不得時.
我們都去實踐.
以致我們真的可以經歷到.
耶穌為我們的同行.
以致我們真的可以.
切切實實地.
將福音的真理.
傳揚開去.
讓人都有機會.
經歷到上帝的拯救.
和各樣的福氣.
多謝上帝.
聽我們在你面前的禱告.
奉耶穌基督.
您得勝的名字來祈求.
阿們.
坦白說.
在我們的生命裡.
我們每一天都能夠.
有神與我們一起過.
我們每一天.
我們一首回應詩歌.

$^{841}$「有你同行」.
(人海裡最終找到是你).
(有悲有喜都竟得起).
(我深信這一切愛我願比).
(勇敢去走至人世便有你).
(回望人生裡你與我一起).
(何懼冷風帶給的痛悲).
(若是兩樣無力的跌倒了).
(是你扶持我在未來不知).
(神啊 你依然向我賀登歌歸).
(祈求賀登眾多).
(我約束的被害殺更多).
(我睡醒的時候陪我你同坐).
(他永遠都不會撇下我).
(愛動了的人他便安慰我心).
(他永遠都不會撇下我).
(引領我做人就出黑暗).
(他永遠都不會撇下我).
(愛動了的人他便安慰我心).
(他永遠都不會撇下我).
(引領我做人就出黑暗).
(他永遠都不會撇下我).
主席.
全能的上帝 一切美善.
傳遍恩惠的賜予者.
我們感激你的公認.
父啊 求你指教我們.
回應你 獻上自己.
並且一切從你而來的豐盛.
好叫我們不單用嘴唇去讚美你.
更加以生命的全部向你獻情.
主啊 無論我們在工作 家庭 事奉.
我們的心思 意念.
我們的待人接物.
面對生活挑戰的種種態度.
我們都願意討你的喜悅.
求主你幫助我們成為一個哀動的人.
好使我們能夠竭力去明白你的話語.
明白你的心意.
過討你喜悅的生命.

$^{881}$成為活祭獻情給你.
我們這樣好告交託.
奉主耶穌基督的聖名致祈求.
阿們.
\newpage



\section{馬太福音 5:5}
\label{sec:YMwn1_GRZYg}
\textbf{2021年1月24日崇拜直播｜吳真真牧師｜幸福由溫柔開始｜馬太福音5\_5}
\newline
\newline
連結: \href{https://youtube.com/watch?v=YMwn1_GRZYg}{\texttt{https://youtube.com/watch?v=YMwn1\_GRZYg}} ~~~~ 語音日期: 2021-01-24
\newline
\newline
\hyperref[sec:POKxDUqFlWw]{\small{< < < PREV SERMON < < <}}
~
\hyperref[sec:index_chronic]{\small{[返順時目]}}
~
\hyperref[sec:s_Njf9cQc18]{\small{> > > NEXT SERMON > > >}}
\newline
\newline
馬太福音 5:5
\newline
\begin{longtable}{cl}
\hline
\hline
章節 & 經文 (和合本修訂版)\\
\hline
5:5 & \begin{tabularx}{0.7\textwidth}{X} 謙和的人有福了！因為他們必承受土地。 \end{tabularx} \\ \\
[1ex]
\hline
\hline
\end{longtable}
$^{1}$我們以聆聽神的話來繼續敬拜.
今日的經訓是.
《馬太福音》第五章一至十節.
耶穌看見許多人上了山.
寄以坐下.
門徒到他跟前.
他就開口教訓他們說.
虛心的人有福了.
因為天國是他們的.
哀動的人有福了.
因為他們必得安慰.
溫柔的人有福了.
因為他們必承受地土.
饑渴無意的人有福了.
因為他們必得飽足.
連戌人的人有福了.
因為他們必蒙連戌.
清心的人有福了.
因為他們必得見神.
使人和睦的人有福了.
因為他們必稱為神的兒子.
為易受逼迫的人有福了.
因為天國是他們的.
今日仍然有幸福的宣講系列.
今日的講題是「幸福由溫柔開始」.
溫柔的人有福了.
因為他們必承受地土.
我們恭敬的請吳珍珍牧師來宣講.
我們聆聽上帝的話語.
今日「百福」來到第三講.
溫柔的人有福了.
因為他們必承受地土.
大家不要以為今天講溫柔.
是特意找我上來講.
因為很多人不會從我身上聯想到溫柔.
平時說話粗粗魯魯.
走路又大大咋咋.
不過我想說在我人生中.
和溫柔有幾年的時間是很接近的.
只有很近的距離.

$^{41}$是什麼呢?.
不知道大家有沒有聽過.
以下這一句話.
是關於幾間在廣州來的學校的學生.
就是「真光珠」「嶺南牛」「培正」「馬嬲頭」.
「培道女子」「溫柔柔」.
這裡有很多培正柔和馬嬲頭.
我不是培正柔,而是培道女子.
所以我和溫柔柔有很近的距離.
不過這純粹是題外話.
究竟我自己在培道中.
最感恩,得到最多的.
就是我在培道中學裡信了主.
在裡面,我真的學習從主的道裡.
學習如何做一個溫柔的人.
這是今天我們八福要說的題目.
「讓我們思一切祈禱」.
因為我們在你的國度裡成為你的子民.
我們可以有一個很超然的身份.
我們可以在你裡面享受無限的福樂.
上帝大我求你.
開我們的眼睛指教我們.
讓我們預備好我們的心.
去聽你自己的教導和說話.
我們恭敬仰望交托.
奉主耶穌基督得勝,明治其球.
阿們.
在這個世代裡.
溫柔到底是怎樣的呢?.
溫柔是否有價值的呢?.
我很想和大家看一看.
1987年的一個運動廣告.
大家都看到這個廣告.
是一位籃球運動員.
他的名字叫查理斯·巴赫利.
他當年是NBA裡面一個好巨.
被公認為是一個好巨統治力的大前鋒之一.
然後這間運動公司就邀請他拍一個廣告.
我被這個廣告的題目.
這個標語深深地吸引了我的眼睛.

$^{81}$這個廣告的標語是什麼呢?.
The meek may inherit the earth.
but they won't get the ball.
溫柔的人是可以承受地土.
不過他不可以讓你搶到球.
似乎在溫柔競爭社會的角度.
溫柔好像沒有什麼價值.
沒有什麼市場的價值.
我還記得很多年前.
我在長洲觀賞了一場球賽.
那場球賽是由建道的同學迎戰.
如果我沒有記錯.
好像是浸會神學院的同學.
如果我沒有記錯.
因為每年都會有不同神學院集合在一起.
會有一個友誼賽.
我還很記得當年有一個同學.
在我們當中都是出名很謙柔.
當日他擔任後裔.
後防的司職.
但我覺得他完全把溫柔表露無遺.
第一,我看到他.
你記住他是後防.
所以他攔截對方的前鋒.
去帶球前進是他的責任.
但我看到他很適當地保持距離.
當他靠近對方碰到球時.
他真的會說聲抱歉.
我經常在場邊想.
「大哥,打球,踢球」.
「你這樣怎可以?你一定會輸」.
是的,的確我們要問.
溫柔到底在這個競爭激烈的社會中.
有沒有它的價值呢?.
我們身處在一個競爭不斷的世界.
如果你不夠強大,力量.
你是沒有立足之地.
你不會得到在世人眼中有利的東西.
包括實質的名利,金錢.
得到人的認同.

$^{121}$人與人之間很自然地成為競爭者.
溫柔在這個世界中變得很吃虧.
而且當我們越要擁有的時候.
我們以為擁有就可以滿足的時候.
原來誰不知會令我們更加沒有安全感.
更加不安.
所以其實如果我們在這個世界中.
去爭取心中的幸福.
其實世界告訴我們.
要用強大的力量.
但聖經卻告訴我們.
真正得到內心的滿足.
能夠在這個地上爭取立足之地.
可以爭取一日長短相對.
告訴我們甚麼.
原來是我們在上帝的角度裡.
用幸福去操練我們的幸福.
以至我們可以承受地土.
究竟承受地土是一回甚麼事呢?.
詩篇37篇第11節.
它說道.
「但謙卑人必承受地土.
以豐盛的平安為樂」.
詩篇37篇裡.
總共有四節經文提到.
神要讓他的子民承受地土.
這個應許.
祂要祂的子民.
蒙恩恩賜繼承祂所賜予的土地.
對於猶太人來說.
承受土地是非常有重要的意義.
當耶和華神和以色列的先祖.
亞伯拉罕立約的時候.
祂應許他賜迦南地為業.
這就是承受地土.
對神的子民來說.
能夠進入迦南地.
得享安息,平安.
而詩篇37篇就說.
以豐盛的平安為樂.

$^{161}$去到新約.
彼得前書又說.
地土代表上帝賜給我們.
各式各樣豐富的屬靈產業.
這些產業是不扭壞的.
是不癲烏,不衰殘.
上帝賜給你這個基業.
祂應許我們.
我們可以在神面前站得住.
可以領受祂的恩典.
豐足的恩惠.
得到祂的眷有.
得到祂所賜豐盛的一切.
當上帝的子民.
可以承受神的產業的時候.
我們要明白到.
其實地土這個字本來的意思是甚麼.
你想像一下.
地土本來就是一班小農民.
他們擁有自己耕種的地土.
我們想像一幅圖畫.
天國的這塊地圖.
和地上世界的地圖是分開來的.
是完全不一樣的地圖.
在這個世界的地圖.
有它的法則,有它的制度.
而這些法則我們很努力地去開墾.
很努力地去耕種這個世界的地圖的時候.
反而令我們不安.
令我們很不安寧.
但是上帝這個天國的地圖.
卻帶給我們這份平安.
帶給我們上帝的恩典,看顧和保守.
但是在這個天國的地圖裡面.
我們是用甚麼法則,用甚麼方法.
去耕種,開墾這塊土地呢?.
《馬太福音》五章五節告訴我們.
溫柔的人就可以了.
到底甚麼是溫柔呢?.
很多人以為溫柔就是體貼.

$^{201}$溫柔,體貼.
很多人以為溫柔很體貼.
溫柔是善解人意.
有很多時候會覺得溫柔是沒有力量的.
不是很夠力量.
是很軟弱,很柔弱的.
但是溫柔這個字.
我們在英文裡面.
剛才我們也說到.
「meek」.
在《牛津字典》我查一下這個字是怎樣翻譯.
它就會說是「溫順謙恭」.
而中文譯本的不同聖經.
都有不同的翻譯.
《和合本》,《新譯本》和《新漢語譯本》.
都是譯作「溫柔」的.
而《和合本修訂版》就譯作「謙和」.
而《呂鎮忠譯本》就譯作「柔和」.
其實原文希臘文這個字.
溫柔這個字.
很多學者.
很多聖經學者.
甚至用不同語言.
包括德文,英文,中文等等的學者.
都認為每一個簡單的字.
或者簡單的詞.
是可以代表希臘文的意思.
例如德文.
德文翻譯希臘文的時候.
他會譯作「不用暴力」.
會譯作「溫柔敦厚」.
「溫和」等等.
其實有很多的.
蘊含了很豐富的意思.
我嘗試從三個向度.
就是溫柔.
究竟對於自己.
我們怎樣對自己溫柔.
我們怎樣對人溫柔.
我們怎樣對上帝溫柔.

$^{241}$這三個向度去演繹一下.
明白一下溫柔的意義在哪裡.
首先第一個.
我們會去了解一下.
我們怎樣溫柔地對自己.
其實人的內心是很脆弱的.
我們很想去肯定自己.
藉著擁有一些東西.
讓自己有一個存在的感覺.
覺得自己很有價值.
在這個競爭激烈的社會裡.
就令我們變得很有競爭性.
很有侵略性.
就算你自己本身不是很想去爭.
不是很想去爭的人.
但因為你身處這個社會的時候.
當攻擊來的時候.
你為了捍衛的時候.
你都會反擊的.
其實真是一個很可悲的現實.
人不單只對別人.
其實很多時候.
我們的內心都很不平靜.
我們的內心都因著.
對自己很多的不滿.
我們心裡是得不到那份平靜.
在耶穌的時代.
一班宗教領袖說.
要求那班上帝的子民.
做很多很多的行為.
訂出很多很多的規條.
這樣反而誤導了百姓.
令人覺得自己內心很有罪疚感.
虧欠感.
很枯乾.
找不到出路.
得不到真正的平安.
古倫神父在一本書.
叫做《幸福人生的八條道路》.
《八福》那本書裡面.

$^{281}$形容當今的世代.
都是一個暴力的世代.
既是社會氣氛的影響.
又或者是我們自己成長的時候的經歷.
令到我們的內心.
帶來一種很大的不平安.
不安全的感覺.
甚至古倫神父說.
我們一班身為上帝的子民.
有時候對自己都很暴力.
很攻擊性.
我們有時都會很嚴苛地去批判自己.
去踩低自己.
這樣令到我們內心的罪惡感.
苦無出路.
但是我們又要在外表上.
去扮演一個很好的基督徒的角色.
所以有時候我們很分裂.
我們外表上好像做一些很好的行為.
但是內心就受一些很嚴苛的批評去推動.
但是這樣反而令到我們的生命.
覺得更加軟弱無力.
使我們對自己的批判更大.
不單止這樣.
對自己的批判之餘.
我們對身邊的人的批判都變得很嚴厲.
令到自己更加苦毒.
更加多怨恨.
「溫柔」是「八福」的第三幅.
我們之前兩位牧者跟我們說到.
當我們遇見基督的時候.
我們的意念就發覺.
我們的心靈本來就是一無所有.
在靈裡面我們是貧窮無助的.
我們會為自己的罪感到哀動.
明白到自己的破碎,軟弱,有限,缺乏和需要.
這就是我們的本相.
耶穌告訴我們我們要「溫柔」.
就是因為祂要提供一個更好的方式.
幫助我們去對待自己.

$^{321}$這個方式就是我們勇敢地去接納自己的本相.
以致我們可以向上帝求更多的恩典.
上帝從來都沒有要求我們.
可以靠自己的能力去完美.
祂沒有要求我們完美才愛我們.
祂很清楚我們每一個.
沒有人能夠靠自己的力量去做一個更好的人.
上帝更加沒有要用你事奉的效果.
你做出來的果效去稱你,去評核你.
祂計算的是你的忠心.
甚至我們要提醒自己.
若果我們用自己能夠做到的好成績.
而去感到滿足的話.
其實我們反而是驕傲的.
而不是溫柔.
我記得我小兒子K3的時候.
有一個很深刻的片段.
我小兒子的性格是有一點完美主義的.
我記得一個K3的小孩.
他晚上十一點多還在做功課.
其實回來的時候.
只有兩頁中文字讓他去抄寫.
但是他做到十一點多.
為什麼呢?.
因為他覺得自己寫的字不漂亮.
只要彎度不夠彎.
角度不夠90度.
他就去抄.
他一個字可以抄十幾次.
我看到他這樣做.
我真的很生氣.
我陪你做功課,你做到十一點多.
還要看著你不斷地抄.
我不斷地跟他說.
很漂亮的,不要抄了.
但結果變本加厲.
我在外面一直去否定他內心那種.
覺得不漂亮的感覺.
對他來說,其實是多了一份迫害.
但是當我去反思.

$^{361}$當我去想怎樣幫他上課.
當然這個媽媽不斷上課去明白的時候.
我發覺原來我要體諒他內心.
自己是真的欣著他那份.
覺得不漂亮而不安的感覺.
我接納他,我明白他.
原來這樣是可以幫到他.
於是我就坐在他身邊.
他一邊做功課,一邊抄寫的時候.
我就說:是嗎?覺得不漂亮嗎?.
是呀,你覺得不漂亮不舒服吧?.
我一直是這樣.
我發覺他慢慢地越抄越少.
由抄十次,變成抄七次,變成抄兩次.
對,讓他知道,接納到.
其實自己內心有這種不完美.
以致他找到出路.
所以上帝很願意我們靠著溫柔.
對自己不要那麼苛刻.
以致和自己和諧.
讓我們能夠用一個更好的方式去對待自己.
這就是溫柔.
另外上帝說到溫柔的第二個向度.
就是對別人溫柔.
我形容為對別人謙和.
當你善待自己的時候.
你不再對自己過份批判.
不再慨嘆自己不及隔離的好.
自己不會無時無刻心想.
如何和別人比較的時候.
自然地你就不會成為一個.
常常要和別人在心裡競爭的人.
不需要用一個很攻擊的方法去對待自己.
你的內心就可以承受地圖.
承受上帝豐富的恩惠.
上次上帝所賞賜的不朽的產業.
因為這份滿足.
我們就不需要再強加自己的意念.
而那些意念是對自己有利的意念.
去加在別人身上.

$^{401}$使別人覺得你很好競爭.
很侵略性.
剛才提到新漢語譯本.
對於「溫柔」這個詞的解釋.
它的翻譯是「溫柔」.
但它在這裡做了一個備註.
它的備註是.
希伯來文當中.
「貧窮」和「溫柔」是同義詞.
「貧窮」是強調對上帝傳言的倚靠態度.
而「溫柔」是強調對人溫和的態度.
我們先講解「溫柔」對人的態度.
究竟對人謙和是甚麼特質?.
在聖經中.
在希臘哲學中.
「溫柔」是一種品格.
原來「溫柔」代表甚麼?.
就是完全不會發怒.
和經常發怒的中間.
是指自己的發怒.
為自己的發怒.
即是自私的發怒.
和為不義發怒.
即是不自私的發怒.
它表達的是.
溫柔的人不是每件事都無所謂.
每件事都謙讓.
他們都有勇敢去面對衝突時刻的能力.
而重點是他們不是從自身維護角度出發.
我這樣想起了舊約聖經的一個人物.
他就是摩西.
在民數記十二章第三節說.
「摩西為人極其謙和,勝過世上的眾人」.
「謙和」這個字.
就是指「溫柔」這個字.
這句經文其實有背景.
背景是甚麼?.
背景是摩西的姐姐和哥哥米利暗和阿倫.
他們對摩西做了一些誣衊.
並且妒忌他作先知的職份.

$^{441}$然後就加了這句說話.
「摩西為人極其謙和,勝過世上的眾人」.
當這件事發生的時候.
摩西沒有主動做出任何的還擊.
反而是耶和華召了他們三人見他.
耶和華很清楚地說.
摩西這個僕人是忠心,是盡忠.
是上帝的揀選.
於是耶和華向米利暗和阿倫發怒.
指責他們毀謗.
還刺下大麻風在米利暗身上.
摩西受毀謗.
但結果他哀求上帝去醫治米利暗.
對於自己被毀謗誣衊.
摩西反而求主去醫治對他作出傷害的人.
他沒有忌恨.
他沒有去報復.
就像亞歷士多德所形容.
溫柔的人就是沒有憎恨,沒有苦毒.
沒有那些憤怒要去報復.
其實謙和的人.
去承受地土.
去得到上帝的任避.
在這件事上.
上帝甚至為摩西出頭.
為他討回公道.
但我們看到摩西對人的謙和.
但摩西也會有發怒的時候.
當摩西帶領以色列人出埃及的時候.
他面對很多民眾對他發怨言.
他也是謙和地面對.
但有一次我們記得.
當他在西乃山領受十誡.
當他下來的時候.
看到一群以色列民在拜金牛毒的時候.
他看到他們放縱.
於是他露髓法板.
然後把法板打碎了.
放火燒了金牛毒.
然後把金粉灑在水面上.

$^{481}$要以色列民喝.
甚至殺死了三千個以色列民.
摩西這樣的舉動讓我們想到.
他和溫和是風馬牛不相及.
但上帝在《食經經》裡說他謙和.
是因為他畏著上帝的公義.
以色列民的叛逆.
他發怒.
這個發怒是不自私的發怒.
是因為神聖不得到滿足而發怒.
不過你可知道.
當摩西做完一連串的舉動的時候.
他就上山為以色列民代禱.
他甚至和神說.
「神啊,如果可以的話.
求你把我的名字從生命冊上除去」.
他為了救他的同胞.
他為了愛他的子民.
他可以放下自己.
這就是謙和.
弟兄姊妹.
摩西的謙和不是一朝一夕的.
我們從他的成長就會發現到.
其實他控制不了自己的法路.
以至到他殺死了欺負自己同胞的人.
到後來他成為神的僕人.
經歷了無數的事情.
他任上帝如何容他.
帶領以色列人出埃及.
每一次當人發怨言的時候.
他都來到上帝的面前.
求上帝的指引.
甚至到最後上帝.
因為摩西他沒有完全相信神.
而自行發出命令的時候.
上帝懲罰他不可以進入迦南地.
但摩西到了申命記最後.
他仍然堅守他的崗位.
呼籲他的百姓一定要謹守.
遵循上帝的吩咐.

$^{521}$摩西的溫和.
摩西的謙和.
是經歷了人生無數次的經歷和鍛鍊.
上帝給他的每一個鍛鍊.
都能夠塑造出謙和的質素.
而這個謙和的質素.
我們發覺是一個完全信服的質素.
正如剛才我們講到.
《新漢語》曰:「希伯來文」的.
「溫柔」和「貧窮」是同義.
而「貧窮」就是對於上帝完全的態度.
我們就知道溫柔的人.
一定有一個向度.
就是對上帝信服.
其實這個字除了在登山寶訓之外.
「溫柔」這個字除了在登山寶訓之外.
在《馬太福音》還有兩處經文出過.
耶穌說:「我心裡柔和謙卑.
你們當乎我的厄學我的樣式」.
「柔和」這個字和「溫柔」是同一個字.
到二十一章五節.
「要對錫安的居民說:.
看啊!你們王內到你者裡.
是溫柔的.
又騎著驢,就是騎著驢驅子」.
「王」是「溫柔的」.
「溫柔」亦是《白福》裡面.
第三個質素「溫柔」的同義詞.
同一個字.
其實希臘文的「溫柔」.
很多時候可以形容為.
一隻野馬被馴服.
受控的一個樣子.
這就是馴服.
耶穌說:「我的厄要當乎我的厄」.
「厄」就是農夫用來耕田的工具.
「富厄」就是馴服的意思.
就是服役的意思.
耶穌祂自己就是謙和的表表者.
溫柔的表表者.

$^{561}$祂的樣式很值得我們學習.
祂是上帝的兒子.
祂是至尊的君王.
祂是高高在天.
但馬太福音說祂是溫柔的王.
菲勒比書二章告訴我們.
祂本有上帝的形象.
但祂不堅持與神同等.
反倒虛己.
取了奴僕的形象.
成為人的樣式.
既有人的樣子.
就謙卑自己.
全心馴服.
以至於死.
且死在十字架上.
耶穌正正展現了.
祂是一個沒有權柄的人.
謙卑自己.
絕對信服在上帝主權的裡面.
以至能夠成就上帝的計劃.
在十九世紀.
有位著名的浸信會牧師.
他叫施布珍.
他提到溫柔的時候.
提到一個這樣的例子.
他說溫柔就是完全信服上帝的意思.
他舉了一個例子.
就是有一位在城中很有學識的醫生.
他去到索爾茲伯利的大草原.
與一個牧羊人聊天.
他問牧羊人.
「你想今天的天氣怎樣呢?」.
「適合牧羊呢?」.
牧羊人說.
「其實只要是上帝安排的天氣」.
「就是好天氣」.
醫生說.
「那即是任何天氣都是?」.
「不是有些天氣」.

$^{601}$「會令你做得辛苦一點嗎?」.
牧羊人說.
「不是,當然不是」.
「上帝無論是什麼天氣」.
「上帝的憐憫都圍繞著我」.
醫生說.
「其實做牧羊人都很辛苦」.
牧羊人說.
「嗯,也是的」.
「的確是花了一些氣力」.
「要勞苦多一點」.
「不過總比我在這裡好食懶飛好」.
醫生說.
「那你做牧羊人的時候」.
「是否遇到很多困難呢?」.
牧羊人說.
「是的,都挺多困難的」.
「不過這樣就可以令我」.
「不用去到五光十色的城鎮裡」.
「被很多東西吸引」.
「不用面對很多誘惑」.
「我還可以在這個大草原裡」.
「用很多時間去默想上帝」.
「令我心裡很滿足」.
「我相信上帝放我在什麼位置」.
「就是最好的位置」.
我覺得這個牧羊人很值得我們學習.
他對上帝安排給他的任何境遇.
他沒有埋怨.
他很安心地.
活在上帝擺放他的位置裡.
他不會以自己為中心出發.
認為上帝不同的處境在阻礙他.
就算上帝要去督責他.
教訓他的時候.
他仍然對上帝有信服.
會向上帝求恩典.
他深信雖然牧羊人是低微的工作.
但他會將他升高.
享受從上帝而來.

$^{641}$百般的屬靈福氣.
原來溫柔可以讓我們有一個很強大的力量.
進入上帝的旨意裡.
在一個不自私的堅持當中.
去忍耐地走上上帝想我們達成的目的.
所以一個溫柔的人.
他是完全信服上帝一切的安排.
認為上帝的安排是最美好.
一位神學家Michael Green.
他說過:「德性的往往不是那些很聰明,很強壯的人.
而是在神的面前極之渺小的人.
因為他們低微到一個地步.
就是上帝很容易就能舉起他們.
不需要擔心他們會變得不可一世.
神可以放心地將他的產業交託給他們.
承受無論是在地上或天上的一切產業.
而不會令到這個溫柔的人造成虧損」.
所以弟兄姊妹.
天國的子民是蒙福的.
因為我們可以承受地土.
我們可以領受一切豐富的恩惠.
拿到不扭壞的福氣.
有主的眷由保護.
為我們出頭.
使我們得到真真正正的平安.
我們要操練成為一個溫柔的人.
我們要對自己溫柔一點.
不要對自己作出嚴苛的批評.
我們要對別人謙讓.
但我們要為公義去發牢.
我們心裡不報仇,不怨恨.
能夠信服在上帝的旨意下.
這就是溫柔.
我聽過一個真人真事溫柔的故事.
素斯的丈夫被黎巴嫩的恐怖份子綁架.
於是他繼續留在黎巴嫩裡面.
一直尋求得到釋放.
但他的姿態從來沒有作出敵對.
指責或查問.
綁架丈夫祖榮的人.

$^{681}$他反而顯出一個人看來很軟弱.
很溫柔的表現.
因為他留在當地的緣故.
他更加認識黎巴嫩的地方.
他關心那個國家的兒童.
他還將他的所長.
就是他很熟悉的音樂.
帶到為黎巴嫩兒童而建設的文化中心裡.
他很愛這班小朋友.
當時他完全不知道.
這份溫柔的力量是那麼大的.
當他很用心地照顧這個文化中心的小朋友時.
他的丈夫祖榮在被綁架的地方.
他從來只吃碎麵包.
細細片片的乾乳酪.
後來換了做一些暖的食物.
還有朱古力吃.
還有水果.
而綁架他的人更加為他額外加了一些毛毯.
一些短的襪.
甚至問他:你想要甚麼聖誕禮物?.
結果他得到一本聖經作為他的聖誕禮物.
後來綁架的人被祖榮逃離被綁架的地方.
可以和他的太太素詩重聚.
他們才知道原來素詩一旦被綁架.
才知道原來素詩一直對這個地圖有需要的人的愛.
原來是一個溫柔的力量帶來那麼大的改變.
所以溫柔並不是世人看得那麼軟弱.
反而溫柔卻是一個很強大的力量.
弟兄姊妹,我們都願意彼此勸勉.
做天國的溫柔人.
不只是陪女子才是溫柔人.
我們一起做天國公民溫柔人.
\newpage



\section{馬太福音 5:6}
\label{sec:s_Njf9cQc18}
\textbf{2021年1月31日崇拜直播｜陳啟興牧師｜幸福由飢渴慕義開始｜馬太福音5\_6}
\newline
\newline
連結: \href{https://youtube.com/watch?v=s_Njf9cQc18}{\texttt{https://youtube.com/watch?v=s\_Njf9cQc18}} ~~~~ 語音日期: 2021-02-03
\newline
\newline
\hyperref[sec:YMwn1_GRZYg]{\small{< < < PREV SERMON < < <}}
~
\hyperref[sec:index_chronic]{\small{[返順時目]}}
~
\hyperref[sec:fM11JJVt3SY]{\small{> > > NEXT SERMON > > >}}
\newline
\newline
馬太福音 5:6
\newline
\begin{longtable}{cl}
\hline
\hline
章節 & 經文 (和合本修訂版)\\
\hline
5:6 & \begin{tabularx}{0.7\textwidth}{X} 飢渴慕義的人有福了！因為他們必得飽足。 \end{tabularx} \\ \\
[1ex]
\hline
\hline
\end{longtable}
$^{1}$領子會讓我們這段時間.
恭聽一段上帝的話語.
馬太福音第五章第一節至第十節.
.
既以坐下,門徒到他跟前來,耶穌就開口教訓他們說:.
「開心的人有福了,因為天國是他們的;.
哀動的人有福了,因為他們必得安慰;.
溫柔的人有福了,因為他們必承受地土;.
饑渴無意的人有福了,因為他們必得飽足;.
連戌人的人有福了,因為他們必蒙連戌;.
清心的人有福了,因為他們必得見神;.
使人和睦的人有福了,因為他們必稱為神的兒子;.
為義受逼迫的人有福了,因為天國是他們的.」.
今天正道的經文是在五章第六節:「饑渴無意的人有福了,因為他們必得飽足.」.
今天為我們宣講的是陳啟興牧師.
今天宣講的題目是「幸福由饑渴無意開始」.
現在恭請陳啟興牧師.
.
阿弟姐妹早晨.
很感恩有機會再一次來旺朕分享神的話語.
我們同心一起禱告.
我們將講的,聽的交託給上帝.
親愛的主,我們去仰望你.
如果我們的生命的焦點離開了你.
我們甚麼都不是,也甚麼都做不到.
任何有意義的生命都不能夠活出來.
主,幫助我們能夠去找到你的心意.
進行你的召命去過我們的一生.
讓今天你的靈去引導我們進入真理.
讓我們聽得明白.
讓我們的心很想去跟從.
讓我們的手有力去活出你託付我們的使命.
同心禱告仰望奉耶穌基督的聖名.
阿們.
在開始的時候.
很想讓大家去看一套電影的片段.
一幅相.
有沒有人看過這套電影?.
是否在座?.
在電腦背後的我聽不到大家說話.

$^{41}$在座有沒有人看過這套電影?.
大家可能在座看不到圖片.
這套電影是《鹿鼎記》第二集《唇紅膠》.
話說劉松仁飾演的陳近南.
跟徒弟周星馳飾演的韋小寶.
有一段對話.
韋小寶問師傅.
自從康熙登位之後.
四海聲平國泰民安.
為什麼還是要推翻康熙呢?.
接著陳近南回答.
其實這句話.
也讓我在這套電影中間.
除了搞笑之外.
有很多的反省.
陳近南回答.
因為這是我畢生的志願.
韋小寶若有所思就自言自語.
為了你的畢生志願.
就要打仗.
要人民受苦.
韋小寶不明白.
要達成你的畢生志願.
其實可能更加多人受苦.
所以我們有沒有想過.
究竟我們的畢生志願是什麼呢?.
很多人都有畢生志願.
畢生志願就是你人生最渴望的東西.
最影響你的人生.
你會調動.
你會去調整一切的資源.
會向著這個方向走.
去達成你的畢生志願.
你的畢生志願是什麼呢?.
如果你的畢生志願是錯的.
你整個人生都會走歪了.
韋小寶不明白.
他不明白他的師父為了你的畢生志願.
就要打仗.
就要人民受苦.

$^{81}$在人生的不同階段.
可以有不同的志願.
如果你現在在這個階段.
你的志願是錯的.
你現在的人生就走歪了.
這裡的錯不是道德的錯.
不是社會倫理上的錯.
是指在神眼中是錯的.
弟兄姊妹.
如果你現在的心願和志願.
不是神所賦予你的志願.
可能我們會走一條錯的路.
我自己小時候有一個志願.
不是老師叫我作品的志願.
我小時候很喜歡踢球.
我的志願就是做一個足球評述員.
因為我小時候已經很心近事.
踢不到球.
很想做一個足球評述員.
我經常去聽叔叔.
不是你們的羅執士.
是林尚義,蔡文健,何鑑剛,何靖剛.
我經常聽他們講球.
我很想做一個足球評述員.
是我一生當時的夢想.
上帝很有恩典.
由很想講球變成了講道.
上帝的恩典帶領.
這不是你的人生.
我帶領你不是去講球.
弟兄姊妹.
如果我們的畢生志願.
我們很想人生追求的東西.
離開了上帝的心意.
恐怕我們的人生真的會失落.
今天在第八篇.
我們講到第四篇.
是講我們最渴望的是什麼?.
我們最渴望的是神嗎?.
好像在舊約.

$^{121}$詩人說.
神啊!我的心切無理如鹿.
切無溪水.
我的心渴想神.
就是永生神.
我們最渴望的是神的義嗎?.
好像主耶穌說.
我們要先尋求神的國和神的義.
我們人生不是尋求自己的畢生志願.
我們是尋求神賜給我們的畢生志願.
我想在講第四篇之前.
我想先宏觀去看看.
究竟八篇是怎麼樣.
我們借助了過渡神學院院長.
郭文慈牧師的八幅圖解.
他的解釋也不錯.
可以幫我們整全概念.
關於八幅.
所以如果你想在座看.
就要看電視機了.
他分了三個層次.
先是最基礎的虛心.
虛心就是靈裡的貧窮.
我們要靈裡變貧窮.
才能夠融入上帝的真理.
第二個層次就是態度.
對自己要哀動.
對別人要溫柔.
對基督要飢渴無異.
然後帶出四個行動.
清心,連率,使人和睦.
和為義受逼迫.
我覺得也不錯.
可以幫我們整全理解.
我們每一個特質.
可能也有關係.
有互動的關係.
八幅就是天國子民的特質.
有這些特質的人就有福了.
我們華人最重視的是福.

$^{161}$接下來的新年.
我們很多的揮春.
都是跟福有關.
不過似福不同彼福.
我們華人所說的福.
或者我們人所說的福.
重點就是擁有.
有錢,有權,有健康,有長壽.
有就平安.
沒有就很卑鄙.
這就是我們人所說的福.
而八福所說的福.
肯定不是這個福.
另外在聖經有一個福字.
用希臘文就是 "Eulogiel".
這個字一般是想說.
神的祝福,神的賦予,神的賜予.
是由位分大的去恩代位分小的.
但是八福也不是說這一種.
因為八福的福字是 "Machairios".
不是這個字.
不是說這一種的似福.
而八福的福字.
其實意思就是指一種很喜樂.
很滿足的生命狀態.
是內在的.
不是外面的東西能夠給我們.
錢,權,健康,長壽.
有我們很感恩.
沒有也很平安.
而八福就是說這一種的生命狀態.
所以當我在說第四福的時候.
我首先要交代清楚.
我們所說的福.
是這一種的福.
上個星期講完就講第三福.
就是講到溫柔.
上個星期講完就自誇.
"培度女子溫柔柔".
我就很認識這位講員.

$^{201}$我可以很印證.
他真的很溫柔.
溫柔的人是什麼呢?.
是對神信服.
你們記得上個星期講完講到三個層次.
對神信服.
今天第四福是講基漢武二.
其實基漢武二的重點就是.
我們看慕上帝.
我們看慕神的旨意成就.
能夠基漢武二的關鍵是對神溫柔.
能夠信服上帝.
我們對基漢不會陌生.
兄弟姐妹可能現在十點多.
如果你還沒吃早餐.
你也相當饑渴.
饑渴是指我們非常切勿.
很想得到.
會用盡任何方法.
要花費多少資源也好.
都想要得到的東西.
這就是我們非常饑渴.
非常渴慕.
我想先講講我們處身香港.
我覺得我們香港人.
有四件事是十分饑渴.
第一件事.
不用說大家都知道.
就是錢.
同意嗎?.
我們對錢是非常饑渴.
錢不是萬能.
其實每個人都相信.
你信不信耶穌.
每個人都相信.
錢不是萬能.
但是沒有錢就萬萬不能.
這句話連基督徒也相信.
沒有錢萬萬不能.
但是當我們認真去看我們的信仰.

$^{241}$我們知道我們沒有上帝.
就萬萬不能.
沒有錢我們依然可以很平安.
很喜樂地生活.
反而如果我們對錢很看惡.
很饑渴.
我們就會養成一個貪念.
很貪.
很想有錢.
才會安心一點.
我們的安全感放在錢上.
所以保羅說貪財是萬惡至近.
所以我們對很多事情有饑渴.
對錢.
其實某程度上也有饑渴.
我們要很小心.
如何將這件事情放下.
由上帝去取代.
有錢也不夠.
有些人對權力很饑渴.
我想大家今天看香港的社會.
如果你是為官的.
或者在議會裡面.
權力的爭奪.
權力的鬥爭很明顯.
但是大家想想職場也不會難以明白.
你做公.
你在職場裡面你發覺很多的權力鬥爭.
為了權的位置的頭銜.
可能有很多很腐敗的.
一些不見光的.
黑暗的事情都會做出.
這種權力鬥爭不僅在職場.
在家庭也有.
家庭也有權鬥.
最傳統就是兩夫婦有權力鬥爭.
如果家裡有四大長老.
可能就更加複雜.
權力鬥爭.
如果你勾結外部勢力.

$^{281}$當然我是在說來自菲律賓.
來自印尼的外部勢力.
就更加複雜.
大家在爭權.
教會應該是上帝的地方.
會不會有權鬥呢?.
弟兄姊妹是有的.
我不是在說旺朕.
在我認識的一些教會裡面.
長輩和教睦在爭權.
不是爭錢.
因為教會一般不是跟我們在商界的錢不太一樣.
但是權力真的會讓人腐敗.
錢除了權我們教會很饑渴之外.
另外就是色.
色包括什麼呢?.
包括性的關係.
兩性的關係.
我們的慾望.
淫念.
很多有錢人.
例如我聽到很多明星.
很多球星.
有錢有權有名.
他們最渴望的.
就是性的關係.
性的滿足.
所以美國神學家.
庫斯德他所說.
Money, Sex, Power.
就是錢,權,色這三樣東西.
最勞弄控制人心.
他們經常很渴望得到這些東西.
因為我在香港觀察多了.
特別是這十幾年.
我多加了一樣.
就是什麼呢?.
就是樣貌.
原來我們很渴望.
就是我們外面被人看到的東西.

$^{321}$姐妹們你有沒有想過.
花多少時間,資源,金錢.
可能是處理你的髮色.
你的外表.
有些不滿意自己的樣子.
去南韓整整他.
整我們的身形,身材.
甚至我們穿什麼衣服.
什麼牌子的衣服.
戴什麼牌子的錶.
開什麼車.
別人看到我們的.
我們的外觀.
我們花多少時間,資源,心機去搞呢?.
你想一想.
每天我們去處理我們外表的時間多嗎?.
還是去讀聖經多呢?.
姐妹們.
這些全部其實都不是邪惡的東西.
不是邪惡的.
彼得提醒我們.
不要以偏頭髮,大金色,穿美衣為裝飾.
當然他那裡是鼓勵女性.
但我想你們每一個都應該這麼想.
我們要耐心,長久的溫柔.
這個才是重要的.
弟兄姊妹.
其實我們有很多東西我們會渴慕.
所以我經常都有一句提醒我自己.
我說「潛權式樣」.
「自見為長」.
如果我們人生渴慕「潛權式樣」這幾樣東西.
我們好像圍了幾面牆在我們的生命裡面.
我們困在這裡.
我們跳不出來.
除了上帝的恩典救拔我們.
好像我自己當年.
一個很爛賭的人.
上帝救拔我.
從一個錢的追求裡面.

$^{361}$轉去看上帝.
弟兄姊妹.
如果我們擁有這些「潛權式樣」.
是虛浮的福.
因為你一定會沒有.
一定會離你而去.
不是真正的福.
我們饑渴的對象.
不會是這些.
我們要很小心.
我們饑渴的對象是什麼呢?.
《聖經》說.
「饑渴無義」.
我們的對象是「義」.
什麼是「義」呢?.
我們不要局限在公義,正義.
正義就是justice.
「饑渴無義」這個義字.
英文就是righteous.
是指神認為對的東西.
就是指神的旨意.
神的心意.
我們渴慕的應該是神的旨意.
我們的畢生志願是什麼呢?.
應該就是神給我們的照明.
這就是我們人生應該去渴慕的東西.
不是去做我自己想做的事情.
我是做神想我做的事情.
很多地方都會問.
神給我什麼照明呢?.
其實我覺得照明.
你不需要完全.
從頭到尾完全明白.
看透了.
知道人生怎麼走.
面對照明.
我覺得最重要.
就是渴慕和信服.
我很渴慕知道.
上帝在我生命裡面的心意是怎樣.

$^{401}$我知道的.
我會跟.
我也不知道怎麼跟的.
走得多遠.
做得多多.
做到什麼果效.
我不知道.
也不需要知道.
但是我會忠心跟進去.
我舉一個例子.
在舊約.
摩西經常給我一個很大的提醒.
他四十歲前就在皇宮好好吃好住.
他完全不知道.
原來上帝會呼召他.
帶領他的子民.
出拉琴.
到了他四十歲後.
他打死人了.
接著受訓練.
這個訓練也有四十年.
到了八十歲才真的踏上這個照明.
我想摩西真的很久很久都不知道原來我的照明.
就是要帶領同胞出拉琴.
但是當他帶領同胞出拉琴的時候.
我想他也不知道.
自己是進不了迦南地的.
你想一下處境.
如果你是打工.
上司的公司跟你說.
你有試用期的.
你有 probation(試用期).
但是四十年.
你願意嗎?.
打死也不願意.
四十年之後才工作.
接著跟你說.
你怎麼做都好.
你不能升職.
你進不了迦南地.

$^{441}$弟兄姊妹摩西不需要完全知道.
他只要跟著上帝的帶領.
一步一步走就可以了.
剛剛新約了另一個很激勵我的.
就是使徒保羅.
當他在大馬式路上決志.
他跟從上帝.
他走這條路.
艱難到不得了.
如果你問我.
他信主前.
在一個公會裡面.
做一個尊貴的議事.
對基督徒有權又抓又鎖.
信主之後被人又抓又鎖.
又坐牢.
又被人打.
被人出賣.
又遇到海難沉船.
弟兄姊妹.
保羅這條路他一直都是信服.
一次宣教旅程.
兩次宣教旅程.
三次宣教旅程.
所以對於趙明.
我覺得最重要.
就是我們很渴望.
我很想知道.
上帝你想我怎樣.
我不是走一條我想怎樣走的路.
當我知道的時候.
我就走.
憑信心憑勇氣.
走出這一步.
信服可能很艱難.
但如果這是上帝給你的召命.
一步一步走下去.
我真的有時候都回望.
如果上帝呼召我獨身.
我就不會追親親.

$^{481}$我就不會結婚.
如果上帝呼召我獨身學.
踏上全職事奉的路.
我就不會留戀在金融界.
弟兄姊妹有時候選擇不容易.
選擇.
我們需要很細心去尋索.
去求問.
但當你知道的時候.
我們就踏上.
但是可能今天.
我們的問題就是.
我們未必很想花費很多精力.
很多時間.
真的跪在上帝面前.
去求問.
你想我怎樣.
坦白說在過往這兩年.
香港有很多很多的困難.
在這兩年.
就是2019年有社會運動.
去年有疫症.
我最多的禱告是什麼呢?.
就是問.
上帝你想我怎樣.
在社會運動裡面.
上帝你想我怎樣.
在疫症裡面.
上帝你想我怎樣.
你的召命是什麼.
當我們知道的時候.
我們就勇敢地走上去.
當我們很渴望.
上帝的旨意.
在我們的生命中成就的時候.
我們就必得飽足.
雖然必得飽足的希臘文.
是一個將來式.
重點不是時間性.
而是應許.

$^{521}$一定會應驗.
所以中文就翻譯成必.
其實必得飽足是一個字.
我們一定是得到飽足.
你不需要去為神憂慮.
祂有沒有能力.
祂會不會做.
飽足不是指我們餓.
這個字可以指餓.
真的吃飽.
飽足在這裡是指內心.
真真正正的滿足.
我們記得奧古斯丁的名句.
人的內心有一個空間.
或者一個空洞.
唯有神才能夠填滿.
當主耶穌再來的時候.
在新天新地.
在啟示錄第七章第十六節.
我們不再飢不再渴.
當主來的時候.
我們完全在上帝的心意裡面.
再不再存在的時候.
我們就有完滿的飽足.
我們完全滿足.
行在神旨意中的人生.
我常常覺得是最滿足最滿足的人生.
所以我也是常常會這麼想.
人生真的不是做我想做的事情.
人生是做神想我做的事情.
行在主的旨意中.
即使有苦難有逼迫.
你也是滿足的.
這個我不詳細說了.
因為在下面第八幅.
為義受逼迫的人.
也是有福的.
在那裡可能各位講完.
會再詳細一點去分析.
所以道理很簡單.

$^{561}$我們要看無神的旨意.
在我們生命裡面去成就.
這是最滿足最滿足的人生.
但是為什麼我們經常做不到呢?.
不知道你有沒有去想這個問題.
我有時候也會想這個問題.
為什麼我不渴望上帝的旨意.
我很渴望自己的想法是達成的呢?.
我服侍職場裡面.
有時候聽到很多的人際關係.
但多數你聽到的是什麼呢?.
這個世界上有一種動物是最討人厭的.
這種動物的名字叫做上司.
是很討人厭的.
每個下屬都覺得上司是很有問題的.
當一個很差的上司.
在下屬眼中一個很差的上司.
他們的禱告是怎麼樣的呢?.
通常禱告就是.
我們股票的名字.
就是中移動.
終於要移動的上司.
祈禱就是上司快點走.
上帝你拿走他吧.
這是我苦難的來源.
弟兄姊妹.
我們要飢渴無疑.
我們要尋求上帝的心意.
在這個關係裡面.
究竟是怎麼樣的呢?.
是不是真的他的問題.
還是我也有問題呢?.
所以當我常常去想.
為什麼不能夠飢渴無意的時候呢?.
我想起兩件事情.
第一件事情就是我們的焦點錯了.
我們的焦點放錯了.
我們的焦點在哪裡呢?.
是不是在錢權式樣子那裡呢?.
我們的焦點是不是就是我們的工作.

$^{601}$我們的家庭.
我們的事奉.
我們所擁有的東西.
是我們常常去思考去尋求的焦點.
當我們的焦點放錯了.
我們整個人生一定會放錯.
因為我們渴望的是那些.
不是上帝.
不是上帝的心意.
我不是說錢權式樣子是邪惡.
我不是說你的工作,家庭,事奉.
你自己所擁有的東西.
那些不需要你.
不是這樣的意思.
我們的焦點永遠都是上帝和他的心意.
我們明白了上帝的心意.
從上帝的眼光去看我們的工作.
看我們的家庭.
看我們擁有的財富.
看我們的人際關係.
看我們教會的事奉.
坦白說.
事奉從來都是做上帝想我做的事情.
事奉不是做我自己很喜歡做的事奉.
弟兄姊妹.
如果我們的焦點放錯了.
我們永遠或者很難去明白上帝的心意.
很難明白上帝賦予我們的照明.
當我們人生離開了上帝的照明.
我很肯定.
這不是一個夢幅.
你不會有長久的平安,喜樂.
不會有長久的滿足.
可能有很短暫,一時之間的刺激快樂.
但你不會有一個在內心所說的飽足.
第二個問題.
就是我們憑眼見而不是憑信心.
保羅很清楚地說.
我們行事為人就憑信心.
不是憑眼見.

$^{641}$Live by faith.
Not by sight.
但我們剛剛反過來.
很多時候我自己也是這樣.
我有時候回想去年.
我當年有機會去安靜.
回想去年疫情的時候.
我有些觀察.
在疫症一開始第一波的時候.
大家最緊張的是什麼呢?.
有沒有口罩?.
有沒有物資?.
有沒有洗手液?.
有沒有紙巾?.
有時候看到很多人在家裡.
紙巾上到天花板.
不斷地去囤積.
我在想.
如果我憑眼見.
家裡沒有口罩.
那就慌張一點.
家裡沒有這些物資.
就慌張一點.
但如果我們憑信心.
我們知道上帝掌管.
我們不是去囤積.
我們是去施羽.
弟兄姊妹.
我們很多時候.
我們不能夠饑渴無意.
原因就是我們眼見的東西.
已經令我們的心.
我們的焦點錯了.
離開了上帝的心意.
在前年社會運動.
到疫情.
到今年很多弟兄姊妹都開始想.
我被弟兄姊妹問了不止一次這個問題.
究竟我要不要移民.
去還是留?.

$^{681}$弟兄姊妹.
憑眼見.
你可能覺得香港不是九流之地.
但我們要憑信心.
其實去還是留是一個屬靈的決定.
上帝要你留.
不要因為這裡很艱難你離開.
上帝要你離開.
不要因為這裡你留戀很多東西.
我不走.
是一個屬靈的決定.
去還是留?.
憑信心.
信德國上帝.
不是憑眼見.
而是看四邊的狀況是怎樣.
弟兄姊妹不是叫我們閉上雙眼.
完全不理周圍發生的事.
我們很觀察我們周圍發生的事.
然後憑信心交回給上帝.
你想我怎樣?.
你的召命是甚麼?.
當我們的焦點擺錯了.
當我們常常憑眼見而不憑信心.
我們很難去看顧.
上帝的心意.
我們很難知道上帝的心意.
很難走在上帝的旨意當中.
我們要怎樣?.
我們要怎樣去操練機悍無意?.
我沒有什麼秘方.
我自己來來去去.
近十幾二十年都是這兩招.
但是這兩招很見效.
真的很見效.
希望大家可以去嘗試一下.
第一件事就是親近神.
很簡單.
親近神.
很多弟兄姊妹經常問.

$^{721}$陳牧師究竟神的召命是甚麼?.
神的旨意是甚麼?.
好像很神秘.
神不會在晚上報夢給我知道.
也不會有一個金甲神人.
晚上下來告訴我他的心意.
也不會像舊約那樣.
面對面告訴我他想我怎樣.
我怎麼知道呢?.
我的答案也很簡單.
上帝的心意.
上帝的普遍召命.
聖經說了.
我們就做聖經說的那些.
這是普遍召命每一個信耶穌的人.
都要去進行.
不用問那麼多.
不用問應不應該做.
是問我怎樣去做.
但上帝也給我們有獨特的召命.
是只給你在今時今日.
在此時此地.
祂給你的召命.
這個你就要認真去尋索.
但我的意見就是.
我一直在尋索我的獨特召命的時候.
那些繼續去尋索.
你會做你已經知道的那些.
那個普遍召命.
上帝在聖經很清楚吩咐.
我們要使萬民做門徒.
我們要傳福音.
我們要愛人如己.
我們要行公義好憐憫.
我們一定要做.
頂多我們怎麼知道呢?.
我們那個獨特召命.
親近神.
所以不要太輕忽你每天的靈修.
看聖經,祈禱.

$^{761}$很重要.
你親近神.
我經常用一個比喻就是兩夫妻.
兩夫妻的關係好.
很親近.
你實名對方想怎樣.
關係不好.
你明他想怎樣都不想做.
所以弟兄姊妹.
當你和上帝親近.
你是自自然然知道上帝想你怎樣.
就好像一個屬靈人.
你自自然然會做一些屬靈的決定.
會做一些屬靈的合神心意的事情.
所以不要輕忽你每個星期去崇拜.
不要覺得這些每個星期都有的.
這個星期錯過了就下個星期.
敬拜上帝和祂見面.
不過很敬畏祂的關係.
是非常非常重要.
很能夠幫助你明白上帝的心意.
你自己查聖經.
在團契裡面.
是不是都是隨隨便便的呢?.
是不是去到才知道.
今天查些什麼經文.
為什麼我不好好預備.
用心去研讀聖經.
明白上帝的普遍召命.
普遍召命和上帝給你的獨特召命.
是有些關係.
不會完全是偏離.
弟兄姊妹.
所以第一個.
你說對方也好.
你說不說你也知道也好.
就是親近神.
你想想.
在2021年.
自己在親近神這方面.

$^{801}$比起2020年.
有沒有進步呢?.
我們不會求一個完美的狀態.
也沒有一個完美的.
最後誰給分數是耶穌基督.
我們不需要代耶穌基督給分數.
我們只是求進步.
2021年的我.
在親近神方面.
比2020年的我.
我要進步.
比2019年的我要再進步.
弟兄姊妹.
很重要.
第二件事.
就是經歷神.
你真的夠親近.
你是知道祂想你怎樣.
可能你有迷茫.
有掙扎.
你大體知道上帝想你怎樣.
上帝想你走一條什麼路.
我鼓勵你.
勇敢的.
有信心的.
踏上這條路.
去經歷祂.
有一點像走獨木橋.
好像左右點一樣.
可能有些危難.
可能有些迷茫.
可能不完全知道.
走神的路的特色就是.
我們不會完全知道整條路怎麼走.
唯有上帝掌管一切.
所以我自己呢.
有一句我自己的名言.
我自己經常說的.
就是「行步見步」.
不是「見步行步」.

$^{841}$是「行步見步」.
走出上帝所帶的第一步.
你會看到下一步.
弟兄姊妹.
經歷神是很重要的.
當你越經歷神.
你越明白這位上帝是怎樣.
你會越想親近跟從祂.
你越親近神.
你會越明白上帝對你的照明.
你就會踏上這條路.
憑著上帝給你的信心.
勇氣.
膽量.
踏上這條路.
所以說了兩個操練.
機悍無疑的秘訣.
可能你們心裡也覺得.
你不說我也知道.
的確是.
我不說你也會知道.
但是千里之行.
始於足下.
希望大家可以在年頭.
開始立志.
我操練.
親近神.
經歷祂.
以致我們一生都可以.
很渴望上帝的心意去成就.
弟兄姊妹.
當你很有質素地去親近神.
有信心和勇氣地去經歷神.
我很肯定.
這是上帝的應許.
你必定得到飽足.
能夠過一個喜樂滿足的人生.
我們同心圖告.
親愛的主.
因為你是一個很樂意去啟示給我們.

$^{881}$告訴我們.
你的心意的上帝.
我們很愚鈍.
我們很多事情都摸不透.
但是你向我們顯明.
主求你幫助我們.
讓我們很上上學武.
你自己的旨意.
在我們的生命裡成就.
以致我們能夠經歷喜樂.
滿足的人生狀態.
主幫我們校正焦點.
離開那些真的能伴我們.
去側觀我們的東西.
讓我們單單去仰望你.
我們單單對準你.
讓你的使命成為我們人生的畢生志願.
成為我們的目標.
我們的方向.
以致我們見你面的時候.
能夠討你的喜悅.
求你幫助我們.
奉耶穌基督的聖名.
阿們.
\newpage



\section{哥林多後書 8:1-9}
\label{sec:fM11JJVt3SY}
\textbf{2021年2月28日崇拜直播｜梁家麟牧師｜慈惠的精神｜哥林多後書8\_1-9}
\newline
\newline
連結: \href{https://youtube.com/watch?v=fM11JJVt3SY}{\texttt{https://youtube.com/watch?v=fM11JJVt3SY}} ~~~~ 語音日期: 2021-02-28
\newline
\newline
\hyperref[sec:s_Njf9cQc18]{\small{< < < PREV SERMON < < <}}
~
\hyperref[sec:index_chronic]{\small{[返順時目]}}
~
\hyperref[sec:rPutWeJ6_Xg]{\small{> > > NEXT SERMON > > >}}
\newline
\newline
哥林多後書 8:1-9
\newline
\begin{longtable}{cl}
\hline
\hline
章節 & 經文 (和合本修訂版)\\
\hline
8:1 & \begin{tabularx}{0.7\textwidth}{X} 弟兄們，我們要把神賜給馬其頓眾教會的恩惠告訴你們： \end{tabularx} \\ \\ \relax
8:2 & \begin{tabularx}{0.7\textwidth}{X} 他們在患難中受大考驗的時候，仍然滿有喜樂，在極度貧窮中還格外顯出他們樂捐的慷慨。 \end{tabularx} \\ \\ \relax
8:3 & \begin{tabularx}{0.7\textwidth}{X} 我可以證明，他們是按著能力，而且超過了能力來捐助，主動 \end{tabularx} \\ \\ \relax
8:4 & \begin{tabularx}{0.7\textwidth}{X} 再三懇求我們，准他們在這供給聖徒的善事上有份； \end{tabularx} \\ \\ \relax
8:5 & \begin{tabularx}{0.7\textwidth}{X} 並且他們所做的，不但照我們所期望的，更照神的旨意先把自己獻給主，又給了我們。 \end{tabularx} \\ \\ \relax
8:6 & \begin{tabularx}{0.7\textwidth}{X} 因此，我們勸提多，既然在你們中間開始這慈善的事，就當把它辦成。 \end{tabularx} \\ \\ \relax
8:7 & \begin{tabularx}{0.7\textwidth}{X} 既然你們在信心、口才、知識、萬分的熱忱，以及我們對你們的愛心上，都勝人一等，那麼，當在這慈善的事上也要勝人一等。 \end{tabularx} \\ \\ \relax
8:8 & \begin{tabularx}{0.7\textwidth}{X} 我說這話，並不是命令你們，而是藉著別人的熱忱來考驗你們愛心的真誠。 \end{tabularx} \\ \\ \relax
8:9 & \begin{tabularx}{0.7\textwidth}{X} 你們知道我們主耶穌基督的恩典：他本是富足，卻為你們成了貧窮，好使你們因他的貧窮而成為富足。 \end{tabularx} \\ \\
[1ex]
\hline
\hline
\end{longtable}
$^{1}$無錯 我們深信上帝是扶持我們向前行的上帝.
我們今天繼續帶著信心.
帶著感恩的心去聆聽上帝的話語.
今日的經訓是在哥林多後書第八章第一節至第九節.
弟兄們 我把神賜給馬其頓宗教會的恩告訴你們.
就是他們在患難中受大試煉的時候 仍有滿足的快樂.
在極窮之間 還格外顯出他們樂捐的口恩.
我可以證明他們是按著力量 而且也過了力量.
自己甘心樂意的捐助.
在三的求我們 准他們在這供給聖徒的恩情上有份.
並且他們所作的 不但照我們所賞望的.
更照神的旨意 先把自己獻給主 又歸附了我們.
因此我勸提多 既然在你們中間開辦這慈慰的事.
就當辦成了 你們既然在信心 口才 知識 熱心.
和待我們的愛心上 都格外顯出滿足來.
就當在這慈慰的事上 也格外顯出滿足來.
我說者說 不是吩咐你們 乃是藉著別人的熱心.
試驗你們愛心的實在.
你們知道我們的主耶穌基督的恩典 祂本來富足.
卻為你們成了貧窮 叫你們因祂的貧窮可以成為富足.
今天早上 哥林多後書八章一至九節的宣講.
是由我們親愛的梁家倫牧師為我們宣講 主題是慈慰的精神.
現在有請梁牧師.
各位親愛的電影姐妹平安.
我們來到第十二次講哥林多後書.
我們看第八章 第八章到第九章.
是哥林多後書的第二個大段落.
在這兩章的聖經裡面.
保羅呼籲哥林多教會的電影姐妹協助他完成為耶路撒冷教會募款的任務.
我們用三次的時間來看這兩章的聖經.
保羅為耶路撒冷教會收捐款.
這是他事奉生涯其中一件比較重要的事.
這件事可以考據的應該是在主後55年左右的時間.
由於他不只是向一間教會 並且都應該綿延了多久.
所以應該整個籌款的行動是維持了相當長的一段時間.
保羅和在耶路撒冷的濕爾斯圖和信徒群體其實沒有多少關連.
他不是他中間的一員.
他不是由使徒親自帶領信主的.
保羅在耶路撒冷復活之後被耶穌呼召.
他的使徒資格是由耶穌親自給他的.

$^{41}$保羅的成長背景和信主前的經歷.
使他和其他絕大多數的出身加利利的使徒其實格格不入.
保羅立志向外邦人傳福音.
並且在傳福音的過程裡建構出一套福音神學.
極大程度的擺脫了猶太教的基本.
使得多數使徒覺得不安.
他們不只是說錯 但總是有點不安.
所以保羅和耶路撒冷教會一直存在著幾微妙的緊張關係.
保羅從來沒有想過在真理 在事公策略上的讓步.
全部不一樣.
不過他總是希望做一些事情.
去表達他對耶路撒冷教會的關愛 善意 舒緩雙方的關係.
耶路撒冷有很多信徒經濟陷入困境裡.
這是長期的情況.
不是偶然 突然發生天災人禍造成的.
為什麼呢?.
因為保羅和耶路撒冷教會裡多數信徒都是出身貧困.
有密切的關係 並且他們很多是來自外地的.
十二使徒都沒有多少個是耶路撒冷人 是不是?.
他們來自外地 他們寄居在耶路撒冷.
他們很難在耶路撒冷找到工作.
找到謀生之志.
打魚的在耶路撒冷怎麼打魚呢?.
對不對?.
種植的也沒有東西可以種.
所以他們找不到工作.
因為耶路撒冷雖然不是嚴格意義的基督教全球總部.
我們可以懷疑有沒有一個這樣的總部.
不過耶路撒冷的教會確實也承擔著關顧各地教會的責任.
使徒常常要奔波各地.
那怎麼也要浪費生活費的 是不是?.
支出自然就大增.
還有 使徒在耶路撒冷.
他們首當其衝的面對猶太教的攻擊.
被捕坐牢的信徒.
以及他們家屬的生活也需要人照顧.
還有一些在家庭 在社區被迫害.
走到教會尋求庇護 尋求供養的信徒.
多數是婦女 也為數不少.
凡此種種原因.

$^{81}$都使耶路撒冷教會長期處於缺乏的狀態之中.
相對地 位處小阿西亞和希臘半島的教會.
雖然信徒仍然多數出身基層.
但是在商業繁盛的地區.
生活的條件總是比耶路撒冷要好 是不是?.
保羅在傳教的同時.
就立志呼籲各地的信徒為耶路撒冷貧乏的信徒捐獻.
實際上你會看到有很多處聖經記載.
保羅為耶路撒冷勸捐或收捐的故事.
包括《少林行傳》24章17節.
在哥林多前書16章14節有提及.
《迦太書》2章10節有提及.
有趣的是 連保羅沒有親自去過的羅馬教會.
他寫羅馬書的時候.
也同樣勸教會捐助耶路撒冷.
羅馬書15章25到32節.
所以哥林多後書8到9章只是其中一個技術.
不過因為這是用了兩章的篇幅來講.
你可以說是最長的一個技術.
要指出的是.
位處小阿西亞和希臘半島的教會.
其實多數都不是很富有.
地區好一些.
信徒多數來自社會的低層.
所以那些地方的教會.
只是比耶路撒冷情況要好一些.
相對地比較充裕一些.
他們可以說是窮鬼中的附庸.
就好像你跟孫明梅去探訪緬甸或孟加拉的平民.
你平日在香港經常喊窮.
你在當地也會覺得自己挺有錢.
你突然間覺得自己變了富有.
這是相對的,不是絕對意義的富有.
保羅誇讚馬其頓教會的信徒.
在8章2到3節.
他們在患難中受大試煉的時候.
仍有滿足的快樂.
在極窮之間還格外顯出他們樂捐的厚恩.
可以證明他們是安著力量.
而且有過了力量.

$^{121}$自己甘心樂意的捐助.
這兩節的經文你會看到.
在患難中受大試煉.
在極窮之間過了力量.
這些騙語都說明馬其頓教會不是很多錢的.
他們不是在有餘的情況下做捐獻.
而是在自己不足的情況下去捐獻.
我們沒有太多資料探索.
馬其頓教會到底遭遇了什麼大試煉.
怎麼會落在極窮之間.
極窮之間我們理解為赤貧.
為什麼會去到這個地步.
不過不知道並不重要.
因為這根本不是這段經文的重點.
保羅論述的對象是哥林多教會.
不是馬其頓教會.
就好像媽媽跟兒子說.
你知不知道隔壁B仔有多厲害.
B仔是不是真的很厲害其實不重要.
你不需要尋根究底.
媽媽只是期望她自己的兒子跟著她說的話去做.
所以馬其頓教會為什麼會有這樣的情況.
根本不是保羅要講的重點.
他是勸哥林多教會跟著做.
保羅之所以提到馬其頓教會的榜樣.
是希望藉此激勵哥林多教會.
維持希臘半島的哥林多教會.
應該比馬其頓教會還要富裕一點.
所以馬其頓教會能夠做到.
他們就應該能夠做到.
如果我們跳去看八章六節.
那裡講到因此我勸提多.
既然在你們中間開辦者辭位的事就當辦成了.
這句說明保羅最初猜險提多.
前往哥林多教會處理他們糾紛的時候.
就已經順便要求提多在哥林多教會為耶路撒冷教會勸捐.
勸捐不是提多去哥林多的重點.
是搭單做的.
不過保羅當然是真的希望哥林多教會參與這一項的辭位工作.
我們不確定提多是不是完成了任務.

$^{161}$我也不知道哥林多信徒對為耶路撒冷奉獻這件事有什麼反應.
或者因為提多在處理教會糾紛的時候.
出現了一些磨擦火花.
可能因此影響了哥林多信徒對捐獻的參與也說不定.
由於可能 我說可能而已.
哥林多教會的反應不理想.
所以保羅在這一封信這裡再落密.
我們再加兩錢肉緊.
希望鼓勵哥林多教會有更大的參與.
在這一段經文裡.
保羅藉著誇讚馬其頓信徒的榜樣.
然後向哥林多信徒募捐.
在他講這番話的同時.
其實保羅是為奉獻或者捐獻提出了三個基本的原則.
這段聖經不是在講為什麼聖光奉獻.
是講為人的生活去籌款.
所以奉獻和捐獻兩個字.
有時候聽到我會互融.
因為我奉獻給上帝.
但我不會奉獻給陳明泉牧師.
我們講捐獻多一點 特別照顧他的生活.
所以我有時候會將兩個字混淆來講.
請大家不要介意.
三個基本原則.
第一 捐獻是任何時候任何情況都可以做的行動.
金錢捐獻是全天候任何情況都可以參與的行動.
前面講到小馬聖和希臘半島的教會.
大多是不富裕的.
馬其頓教會更是處於患難中受大撕裂和極窮之間.
因此他們所做的就是過了力量的捐獻.
他們貌似沒有條件和能力去捐獻.
但他們仍然可以捐獻.
什麼人都需要捐獻.
人人都需要捐獻.
凡是領受了上帝恩典之國.
領受了上帝恩典.
我沒有多餘的說話.
領受了上帝恩典是講事實的.
但你不一定是意識到的.
你又自覺自己領受了上帝恩典.

$^{201}$這樣的人不論男女老幼.
不論富貴貧窮都需要捐獻.
有青少年說我還沒出來工作.
我只是拿家裡的零用錢.
我有什麼資格捐獻呢.
有初職的青年說.
我剛剛畢業出來打工.
每月賺了一萬多兩萬元.
很可恥的.
扣除給政府的關東還款.
我自己食用都不夠.
怎麼捐呢.
沒得捐.
有成年信徒跟我說.
我要供樓.
要存錢做生意周轉.
支付子女在外國的讀書費用.
哪裡有餘力捐呢.
有老年信徒埋怨.
我現在只靠微薄的退休金.
政府的生果金生活.
今天不知明天有什麼能力捐獻呢.
是的.
除非我們富有像李嘉誠.
然後我們就可以搞慈善基金會.
除非我富有像Bill Gates.
我就會捐出絕大部分財產來做慈善公益.
如今窮鬼如我.
根本沒資格談捐獻.
王震.
捐獻任何人.
王震,弟兄姊妹.
相對於李嘉誠.
你是不是也是窮鬼.
你不認不得的.
我們個個都是窮鬼.
我們都沒資格.
沒能力捐獻.
弟兄姊妹.
你知道我在說反話.

$^{241}$真相是.
聖經的教導是.
無論我們落在一個怎樣的經濟環境.
在富裕在貧窮在豐年在荒年.
自覺寬鬆自覺緊絀.
我們都可以捐獻.
亦都需要捐獻.
原因呢.
就這裡有兩個原因.
第一就是前面說過的.
富有和貧窮是相對的觀念.
相對於李嘉誠.
我們個個都是窮人.
但身邊仍然總有一些比我們貧窮.
需要我們跟他們分享.
我們自己所有的人.
是不是.
就好像馬其頓教會.
有能力捐獻給耶路撒冷教會一樣.
第二個原因是.
捐獻的本質.
是上帝恩典的表達.
捐獻反映了我們確實領受了上帝的恩典.
這些恩典是足夠我們使用之外.
還有餘剩.
可以跟其他人分享.
並且藉著捐獻.
我們將上帝的恩典.
都去覆蓋其他人.
讓他們都可以分尖.
甚至是經歷上帝的恩典.
在這兩張聖經的起手.
保羅提到馬其頓教會的榜樣的時候.
他特別強調.
第一節.
弟兄們.
我把上帝賜給馬其頓教會的恩告訴你們.
他不是誇讚馬其頓教會的慷慨.
或對人的愛心.
或對上帝的信心.

$^{281}$而是指出.
這是上帝給他們的恩典.
馬其頓教會得到上帝所給的恩典.
如何看得出來?.
就是第二節.
他們在患難中受到大試煉的時候.
仍有滿足的快樂.
在極窮之間.
還格外顯出他們樂捐的口恩.
上帝給馬其頓信徒的恩典.
帶到一個地步.
即使他們在最不利的患難情況之中.
即使他們被各種大試煉.
一邊折騰一邊消耗他們所有的東西.
他們仍然流露出滿足的快樂.
有一個答堂的樣子.
並且仍然樂意作出捐獻.
這個老人家.
我剛剛好.
以下的道理我以前講過很多次.
如何證明我們是蒙恩的人?.
藉著我們在人面前大講自己如何蒙恩的見證故事?.
講見證是好的.
但作用是有限的.
多講無謂 行動最實際 是不是?.
最重要的證明是.
我們成為有恩慈的人.
領受恩典 Grace.
使得我們成為恩慈的人.
Gracious.
耶穌在《馬達福音》18:33節.
責問那個領受了恩典的僕人的說話.
你不應當連恤你的同伴.
像我連恤你嗎?.
這句說話表明.
變得恩慈是領受了恩典的自然後果.
也是領受了恩典後的信仰和道德責任.
有責任變得Gracious.
同樣.
如何證明我們有豐盛的生命?.

$^{321}$如何證明我們活得豐富 每樣都有?.
我拿個錢包出來給你看.
裡面是塞滿錢的.
是沒有意思的.
拿我的銀行存摺出來養.
看一看那個圈.
也很無聊.
最好的證明是什麼?.
我們慷慨地跟別人分享.
藉著分享證明我們有餘有靜.
好像耶穌所說.
一棵好樹是由它所結的果子證明.
恩慈是聖靈所結的果子.
藉著恩慈證明我們是領受了.
上帝豐富的恩典的人.
中國一句老話.
對於那些孤寒的人.
叫做刻薄寡恩.
從信仰的角度來說.
寡恩不僅是指他對人不太願意給恩典.
也指著他自覺領受的恩典稀少.
弟兄姊妹.
你們是不是蒙恩的人?.
你們是不是得著上帝豐盛的恩典?.
如果是就將這個事實表現出來.
這是第一點.
第一個原則.
任何情況下我們都可以捐獻.
都可以奉獻.
只要你自覺豐富.
自覺領受上帝的恩典.
你就可以.
第二個原則.
捐獻是參與和權利.
金錢捐獻是參與和權利.
捐獻使我們成為侍工團隊的一份子.
當我們用金錢捐獻給一個侍工的時候.
我們就在那個侍工上有份.
八章四節說.
再三求我們.

$^{361}$准他們在這供給聖徒的恩情上有份.
求我們准他們有份.
有份的意思是團契相交.
就是說能夠埋堆.
能夠入伙.
這句話翻成廣東話.
就是收買他們玩.
馬其頓教會再三求保羅.
在支持耶路撒冷教會的事上.
收買他們玩.
他可以埋堆.
我以下的說話可能顯得有點老病.
今天是一個富裕的社會.
讀書期間很少需要去打工.
很多我認識的大學生.
在大學時代都沒有去兼職.
完全不用.
家裡已經供應足夠.
我成長的年代完全不一樣.
我小學階段就已經要做童工兼職.
不僅僅是賺自己的零用.
更加要幫補家計.
不過正因為很早.
就要分擔家裡的開銷費用.
我對家裡就有很強的承擔感.
我是家裡的骨幹成員.
我對維持這頭家.
有推卸不到的責任.
我從來沒有怪過我媽媽.
沒有給足夠的零用錢我.
她根本沒有給過我.
或者沒有幫我買最新款的手機.
球鞋.
也沒有嫌家裡的飯菜不好.
讀書環境不好.
原因呢.
我是有份養家的人.
我在教會也有相同的認定.
我信主後.
參加一間剛剛成立的屋村教會.

$^{401}$不是很多的人.
裡面大部分都是小孩子.
或者是老人家.
捐獻能力低是不在話下的.
事奉人手都很緊張.
所以我受洗不到一個月.
我就開始教主的學了.
接著就陸續承擔.
兩個團契的導師工作.
忙到邊都直了.
教會成立沒幾年.
母會就決定讓我們在五年內滯養.
每年cut五分之一的補貼.
財政壓力非常沉重.
弟兄姊妹咬緊牙關.
勉勵承擔.
我把自己的關人窿都要拿出來.
不過幾年後.
我們不僅僅自立了.
我們還勇敢地開了一間分堂.
在另一個屋村.
承辦一個幼兒中心.
期教會.
不過很慘.
剛剛開始的時候.
那個幼兒中心.
營運得不算很理想.
收生很不夠.
但是你請了一個隊人在那裡.
要給人工的.
每個月虧損很嚴重.
教會又陷入另一場財政危機.
那個時候我擔任義務堂主任.
就和一班的執事.
我們真的大家協定.
準備抵押自己的屋.
來幫教會周轉.
因為去到最後.
我們都勉強過關.
未需要走到這個地步.

$^{441}$我和一班立志承擔教會.
財政責任的弟兄姊妹.
從來不會覺得自己是教會的局外人.
繞著手指著教會指著點點.
是的 教會.
各樣情況都不算理想.
也有很多問題.
但我們就是教會的一份子.
所有問題都和我們自己有關.
是我們做得不好.
不夠好.
我們願意為教會負責.
我們願意為教會找數.
捐獻是埋堆.
捐獻是成為其中一份子.
你只是說出席.
在場.
現在很講這件事.
不會使我們成為局內人.
享受各種服務.
不會使我們成為局內人.
觀察批評.
都不會使我們進入局中.
到最後你批評得越多.
你越會抽離.
抽身.
唯有在行動上付出和承擔.
才使我們成為真正的成員.
金錢捐獻.
不是唯一的承擔方式.
卻是非常具標誌性.
耶穌說過.
財寶在哪裡.
心就在哪裡.
連錢都不肯拿出來的人.
很難有真正的關切.
真正的擁抱 是不是.
現在我們來到第三點.
第三個原則.
捐獻使人的道德和生命變得完全.

$^{481}$金錢捐獻是屬於生命的完成.
捐獻使人變得完全.
在八章七到八節.
保羅在講述了馬其頓的見證故事之後.
就直接教導哥林多信徒說.
你們既然在信心,口才,知識,熱心.
和待我們的愛心上都格外顯出滿足來.
就當在這詞彙的事上.
也格外顯出滿足來.
我說者說.
不是吩咐你們.
乃是藉著別人的熱心.
試驗你們愛心的實在.
哥林多是一間驕傲的教會.
身處知識水平高的希臘半島.
信徒自覺各方面的條件.
都比其他人強.
他們有智慧,有能力.
他們甚至因此看不起保羅.
保羅沒有否定他們這樣的驕傲自覺.
自我感覺良好的情況.
卻是順勢地說.
你們真的覺得自己樣樣都第一嗎?.
那麼拜託.
在金錢的捐獻上也造成第一名吧.
格外顯出滿足來.
指的是充盈,滿溢,達到完全.
保羅的意思是.
唯有加上慷慨的金錢奉獻.
然後才證明哥林多信徒的愛心充足完備.
然後才證明他們在各方面都豐盛圓滿.
很有趣.
這些說法跟我們日常是很不同的.
很多人都說.
講錢失感情.
保羅的說法是剛剛相反的.
講錢不僅沒有失感情.
反而是證明我們對人的感情.
講錢才有感情.
講錢才得感情.

$^{521}$試驗你們愛心的實在.
意思就是.
講金就是講心.
講心就是要講金.
中國人是很實際的民族.
非常明白.
試驗你們愛心的實在這句說話.
用金錢去贈禮.
我們通常的叫法是做人情.
這是人情.
給錢是做人情.
子女孝順父母.
我不知道你們是不是.
定期的供奉是必須的.
過年過節我見到老人家.
你不可以只虛寒問暖兩句.
最好是怎樣.
就是給錢給他.
對他來說.
最感覺到你給他的溫暖.
給他的愛心.
是不是.
除了講錢失感情之外.
很多人常常說.
金錢僅僅是身外物.
金錢是觸物.
是次要的事.
給錢證明不到什麼.
其實我們教會的人.
常常都會順著這個想法.
例如我這個牧者.
明明在港台上.
是想教導信徒.
有關金錢奉獻的道理.
不過我總是婉轉的這樣說.
新約聖經其實也沒有.
所謂奉獻的嚴格要求.
是不是.
只是要求我們所有奉獻.
我們奉獻全心全人.

$^{561}$才是最重要的.
金錢不是很重要的部分.
或者一個執事.
主日崇拜收捐之後.
帶領謝土的時候也這樣說.
我們剛才奉獻了一點點的金錢.
不過這不是最重要的.
我們更加願意奉獻我們的生命.
求上帝閱立.
你剛才只是放了20元進去.
那真的很少量的金錢.
少量的金錢是真的.
不過那20元就顯示你願意奉獻生命所有.
是不是.
很怪.
無論如何.
我們不喜歡說錢.
金錢奉獻重要不重要.
一個關鍵的問題是.
端在於我們自己對金錢的看法.
我們需要誠實的自問.
金錢在我們心中的位置有多高.
如果我們自己很看重那個dollar sign(美元標誌).
你就不可以說奉獻金錢不重要.
你送禮.
如果送給你重視的人.
你珍惜的人.
例如你的女朋友.
老婆是不是最珍惜我不知道.
女朋友一定是最珍惜的.
你總是要搞盡心思.
想最好的禮物.
最貴重的禮物送給她.
以前追求我太太的時候.
都試過將我所有存摺的錢一次過都洋出來.
買了一隻很便宜的手鏈給她.
這個才是送禮.
是不是.
就不會像你現在.
將人家送給你的賀年禮物.

$^{601}$轉手送給別人.
我還有幾盒金沙等著送.
有什麼機會.
如果我去她家我就送.
那些有什麼情意是沒有的.
你將一些你不要的.
你不想見到的東西拿走.
是不是.
你奉獻給上帝的東西.
總是你心中覺得最珍貴的.
是不是.
是吧.
我們珍惜金錢.
我們珍惜我們的命.
我們才將我們心中認為最好的東西送給上帝.
是不是.
所以金錢重要不重要.
看你怎麼說.
你自己想清楚.
是不是.
你真的覺得那些是糞土.
麻煩給我一些糞土.
是不是.
金錢無疑是新的愛物.
但是金錢對信仰的意義其實都很關鍵.
對基督徒來說.
生命中最大的偶像.
就是錢.
和上帝爭奪在生命中的中心位置的都是錢.
所以耶穌沒有提醒我們.
你們不可以既事奉上帝又事奉觀音.
因為你多數都不會的.
但是他在說.
一個僕人不可以事奉兩個主.
不是惡者過.
愛那個.
就是重者過.
輕那個.
你們不能又事奉上帝.
又事奉馬門.

$^{641}$耶穌將上帝和馬門相提並論.
視為競爭的兩造.
是不是.
可見.
在耶穌眼中.
金錢真的很重要.
是不是.
金錢真的很次要的事.
不是的.
捐獻金錢是信仰一個很重要的表達.
可以無段的說.
連錢都不肯捐獻的人.
根本談不上其他的奉獻.
不要說什麼全人奉獻.
我和身邊的人看到的經歷是.
金錢捐獻常常是我們學習全人奉獻的第一步.
你剛剛信主.
教會就教你要學奉獻.
最少每個星期主一崇拜.
預備一些錢.
可能一兩元.
放進去.
然後學習做月捐.
金錢捐獻是我們學習全人奉獻的第一步.
但金錢捐獻.
常常是我們學習全人奉獻的最後一步.
常常有老年的弟兄姊妹和我說.
他們年紀大了.
做不到什麼事奉.
只能夠奉獻一些錢.
讓年輕的弟兄姊妹有氣有力.
你們做多一點.
這十年八載.
我最少很深刻印象的是.
我曾經幫兩個肢體.
處理他們最後的財務安排.
一個是我所敬重的回音牧師.
年近九十的他.
有一次約我喝茶.
他告訴我.

$^{681}$他把他住的那間屋都賣了.
分成三部分奉獻.
其中一部分捐給神學院.
他跟我說.
他請我做見證.
他是剛脫脫來到人間.
又剛脫脫的離去.
他沒有霸佔任何東西.
另一個是一個加拿大的姊妹.
是一個醫生的太太.
那個醫生早幾年已經死了.
這個姊妹.
把她自己多年珍藏的所有首飾都賣了.
在離世之前.
她要做妥所有的財務安排.
什麼都奉獻了出來.
對我來說.
這是全人奉獻.
這是全所有奉獻的最完美的呈現.
納至以他們作為我的榜樣.
老人家最喜歡派錢給晚輩.
我媽媽最喜歡給孫子們收錢.
因為他們剩下派錢的能力.
老得只剩下錢.
窮得只剩下錢有負面的含義.
但老得只剩下錢就沒有什麼不對.
這是自然法則.
有一天我們都會這樣.
跟大家說個小故事.
在新年假期之後.
約了一個年長的弟兄喝咖啡.
我們去了一間他不熟悉的咖啡店.
因為那個地點是我提議的.
他還找了一些時間去找那間咖啡店.
強調他不熟悉的原因是.
當他走進這間咖啡店之後.
他竟然立刻就拿出一疊利是.
每個人都派.
每個店員派.
又走去那個吧桌又派.

$^{721}$那天是年初六.
我是見過很多人.
平日去的酒樓食肆.
向員工講派利是.
這是很普遍的做法.
一方面新春期間.
哄人開心皆大歡喜.
另一方面打賞建立感情.
鞏固關係.
希望來年獲得更佳服務.
當然有些人也會很享受.
發一點錢給一大堆人圍著.
梁先生新年快樂.
陳先生馬道恭成.
講盡一切恭維祝福的話.
很有面子.
我身邊很多人都喜歡這樣做.
不過好像前面說的那個弟兄.
他去到一間剛剛初步到的咖啡店.
講派利是就有點奇怪.
這間店離他家辦公室也很遠.
他應該不會想在這裡建立飲食新據點.
所以他這樣派發.
投資回報是沒有的.
跟他坐下來聊天.
聊了一個多小時之後.
弟兄就講到他這十幾年心態的轉變.
他說他老了.
很喜歡現在對人做一些小愧憎.
例如去餐廳.
那些比較高級的.
通常洗手間有人在那裡駐守.
他就會做一個比較豪爽的打賞.
一出手就一張紅底.
他說對他來說.
這是一個小錢.
不過收到的人就得到一筆意外之財.
付出的沒有痛苦.
收的就很快樂.
所以收資其實不是相稱.

$^{761}$付出一點點收穫大大.
所以他說他很樂意將這個變成他的生活習慣.
他說他甚至不只是想慈惠工作.
他只是想讓身邊的人開心一下.
做一點點美善的事.
是善事不是慈事.
送贈的快樂.
我相信不少弟兄姊妹都有過的經驗.
甚至成為了你們的生活信念和習慣.
對人慷慨.
對人不計較.
樂於和人分享自己所有的.
確實是基督徒的良好習慣.
希望我們每個人都已經養成.
施比受更為有福.
聖經這一句寶貴的教導.
不是說施贈者.
藉著今天的施贈來換取明日的福氣.
施比受更為有福氣.
今日施贈明日有福.
不是將施贈變成一個投資獲利的手段.
而是施贈者在施贈的當下.
就已經覺得幸福快樂.
讓人快樂.
看到人快樂.
覺得自己能夠幫到人快樂.
自己就很快樂.
施贈者.
自己正在經歷快樂.
施贈本身就帶來福氣.
所以不需要講積積功德.
將來在天堂上有更多回報.
其實不是這樣想的.
今日我做的時候都覺得很快樂.
我都覺得很快樂.
我希望大家都是這麼快樂.
奉獻給教會是很快樂的事.
能夠有能力去幫一個人是很快樂的事.
今日的訊息就講到這裡.
剛才那位弟兄.

$^{801}$亦都傳給我一個訊息.
他說這些付出是他能力範圍內的事.
區區小錢 芝麻小事.
他豪得起.
我和大家賣個小小的關子.
豪得起.
這是一個很有趣的說法.
下一次的講道.
就會集中講豪得起的問題.
我們下次節目再見 拜拜.
\newpage



\section{馬太福音 5:10}
\label{sec:rPutWeJ6_Xg}
\textbf{2021年3月7日崇拜直播｜陳明泉牧師｜幸福由為義受逼迫開始｜馬太福音5\_10}
\newline
\newline
連結: \href{https://youtube.com/watch?v=rPutWeJ6_Xg}{\texttt{https://youtube.com/watch?v=rPutWeJ6\_Xg}} ~~~~ 語音日期: 2021-03-07
\newline
\newline
\hyperref[sec:fM11JJVt3SY]{\small{< < < PREV SERMON < < <}}
~
\hyperref[sec:index_chronic]{\small{[返順時目]}}
~
\hyperref[sec:zila_5Bj4Bo]{\small{> > > NEXT SERMON > > >}}
\newline
\newline
馬太福音 5:10
\newline
\begin{longtable}{cl}
\hline
\hline
章節 & 經文 (和合本修訂版)\\
\hline
5:10 & \begin{tabularx}{0.7\textwidth}{X} 為義受迫害的人有福了！因為天國是他們的。 \end{tabularx} \\ \\
[1ex]
\hline
\hline
\end{longtable}
$^{1}$弟兄姊妹仍然以一份聆聽神的話.
這份謙卑,恭敬的心.
來聆聽今天宣讀的聖經.
經文是在《馬太福音》第五章一至十節.
《馬太福音》五章一至十節.
耶穌看見許多的人就上了山.
既已坐下,門徒到他跟前來.
他就開口教訓他們說.
「虛心的人有福了,因為天國是他們的.
愛動的人有福了,因為他們必得安慰.
溫柔的人有福了,因為他們必承受地土.
饑渴無意的人有福了,因為他們必得飽足.
憐恤人的人有福了,因為他們必蒙憐恤.
清心的人有福了,因為他們必得見神.
使人和睦的人有福了,因為他們必稱為神的兒子.
為而受逼迫的人有福了,因為天國是他們的.
這是今天講題的經文.
在《馬太福音》第五章第十節.
「幸福由為而受逼迫開始」.
這是今天陳明全牧師的講題.
我們現在專心聆聽這段宣講.
我們請陳牧師.
各位弟兄姊妹平安.
今天能夠和大家一齊來領受上帝的話語.
慘不得眼下,這個百福的系列.
今天就去到最後一講.
為而受逼迫的人有福了,因為天國是他們的.
其實之前在我們篇幅裡面也曾經講過.
「饑渴無意的人有福了,因為他們必得飽足」.
在那一次的講道,Tormund也講過.
大概關於渴慕神的義.
今天這個跟那個不是對立,或者是講另一件事.
反而是承接著之前的主題.
上一次是講我們渴慕取得生命裡面.
上帝的義,或者上帝的旨意.
在今天裡面我們為上帝的旨意.
我們願意付上代價.
講起義,其實在坊間裡面有不同的定義.
對於義的看法,講公義.
有些是從經濟角度,怎樣去分配.

$^{41}$以致到分配是公義的.
有些是從發展機會.
不同階層的人都可以同樣有.
一起發展自己的一些能力.
有些講政治等等,不同坊間.
講公義其實是一個大題目.
是一個很闊的題目.
不過很多時候在教會裡面.
講上帝的義,或者講義.
很多時候我們會夾雜著坊間裡面不同的定義.
來講的,所以講著講著.
就好像大家自說自話,大家各自表述.
然後大家就像牛頭不對馬嘴.
究竟講的東西好像不是很對.
今天我想和大家拉回聖經裡面.
怎樣看聖經的義.
或者特別是針對在馬太福音裡面登山寶訓.
對於義的看法.
對於義的時候耶穌怎樣去講述.
一個天國的子民.
怎樣去活現一個上帝義的一種生命.
如果大家記得我講過的.
在八福裡面的八句說話.
不是純粹是一個標題.
或者是純粹是一句說話.
就可以濃縮所有要講的東西.
其實是登山裡面的一個引言.
所以這八句說話.
其實後面的講論.
它會互相呼應的.
所以今天我們講的這一段經文.
怎樣去為義受迫迫的人有福了.
因為天國是他們的.
我會用多一些後面登山寶訓的教導.
來展現究竟耶穌所講的是一回甚麼事.
我們今天怎樣來實踐.
在這個時間我想邀請大家和我一起祈禱.
我們需要聖靈幫助我們才能夠明白這段話語.
我們同心祈禱.
阿爸天父願你在我們當中帶領我們.

$^{81}$讓我們能夠明白你的心意.
在這個世代當中.
讓我們帶著你的話語來行事為人.
幫助我們在此時此刻我們聆聽你的話語的時候.
願你親自對我們講說話.
幫助每位聆聽的頂子們聽得清楚.
幫助孫講的讓人聽得明白.
恭敬將來的時間交託給你.
願聖靈你引導我們眾人.
獻上我們祈禱.
奉靠基督耶穌得勝的聖命祈求.
Amen.
剛才講過了.
義的定義.
義如果用在聖經裡面比較概括一點的字眼.
或者比較容易理解的字眼.
就是上帝的旨意.
這個旨意不是說上帝要逼我們去做某一些我們不喜悅或者我們辛苦的事.
這個旨意是上帝美善的旨意.
上帝創造這個世界萬物.
或者上帝在這個世界掌權當中.
上帝祂有祂的心意.
上帝很想祂的子民.
安著祂美善的心意.
來過他們的生活.
不過在這裡講美善的旨意.
還要繼續仔細理解.
究竟美善的旨意包括什麼呢.
包括三方面.
第一方面就是道德性的旨意.
上帝在這個創世的世界裡面.
祂有祂的主權.
上帝在人與人之間的關係.
人與大自然.
或者人與神之間.
有一個上帝希望我們.
隔守的一個生命準則.
這個是道德性的旨意.
在聖經裡面有大量的篇幅教導我們.
怎樣行事為人.

$^{121}$應該怎樣做人.
怎樣去對待人.
接下來的登山寶訓.
都有很多篇幅講這件事.
這個是道德性的旨意.
所以道德是本身上帝美善的彰顯.
祂期望所有的弟兄姊妹.
所有天國的子民.
都覺得上帝認為好的事我們都去遵守.
這是第一個道德性的旨意.
第二個旨意就是救贖性的旨意.
救贖性的旨意是講什麼呢.
世人都犯了罪.
如果純粹按人的能力來實踐道德性的旨意.
人其實都是活在罪裡面.
因為人沒有辦法透過自己的能力.
來活現一個完美道德的生活.
所以我們需要唯一的救主.
基督耶穌.
耶穌將我們從罪裡面救出來.
羅馬書經常講.
我們被稱為義.
不是因為我們做了很多好的道德.
或者做了很多好的事.
而是因為基督耶穌.
基督耶穌將我們的罪換在十字架上.
用祂的生命來換取我們的生命.
所以這是一個救贖性的心意.
耶穌繼續在福音書裡面.
或者登山寶法裡面都繼續講.
不是純粹講做些什麼.
而是我們每一個人能不能夠在上帝的名下.
在上帝的任命之下成為神國的子民.
這是救贖.
上帝救贖的旨意就是要臨到我們每一個人當中.
這是第二部分.
第三個反而是我們最關心的.
就是弟兄姊妹的個別主.
生命裡面神的旨意.
這個女孩是不是我的老婆呢.

$^{161}$我做這份工作好.
還是住在這個區份好呢等等.
這是每一個人很獨特自己的生命故事.
上帝在我們生命當中.
究竟怎樣帶領我走未來的路呢.
我自己個人的前路是怎樣呢.
這是個別性的旨意.
當然在這個篇幅裡面.
耶穌基督特別針對的是什麼呢.
就是救贖的旨意.
和道德的旨意.
所以在接下來的篇幅裡面.
耶穌不會特別說將來的擇偶條件.
他不會說很多的.
因為這是個別性的.
但是他對於一般道德和上帝的救贖.
耶穌就用很多的篇幅.
甚至他要極力的讓每一個聆聽的子民.
他真的要收得到這個訊息.
不過當我們明白了.
這是行道德的旨意.
是上帝好的計劃.
和上帝救贖的恩典之後.
走這條路是什麼容易的路呢.
耶穌繼續說這是一條窄路.
這條路是很難走的.
我們以前很多時候.
有意無意都會進入一種想法.
我們經常覺得.
做上帝行上帝的義的時候.
我們做每一樣事情都會諸事大吉.
我們覺得我們行義.
我們從那件事順不順的角度來判斷那件事.
是不是行在上帝的旨意當中.
所以那件事順頭順路.
我們就覺得這是上帝的旨意.
不過在耶穌的登山寶訓裡面.
或者特別是在說道德救贖的上帝的計劃.
上帝的旨意當中.
走這條異議的路的時候.

$^{201}$一點都不容易.
荊棘滿途.
很多人都面對很多的困難.
雖然在舊約裡面說到.
我們遵從耶和華的吩咐.
我們所做的盡都會順利.
一個因果神學.
不過這個因果神學不是說.
所有子民在生活裡面.
做一件好事我們就有好的報應.
這個因果神學是說整個以色列.
整個國度整個群體.
整個國運.
是一個整體而言.
不是一個個別我們每一個人走自己的路.
每一個人走在上帝的路裡面.
都是荊棘滿途.
都是面對著很多不同的困難.
所以我們要謹記一件事.
走上帝的旨意.
未必一定順利.
甚至耶穌預告.
走上帝的旨意.
追求神的旨意的時候.
反而會帶來很多的荊棘.
很多的困難.
如果在這裡面特別說.
為了義我們甚至要面對逼迫.
生命當中我們持守上帝的道德.
持守上帝的救恩.
在生活裡面.
我們可能會面對很多的困難.
耶穌繼續說.
那個逼迫這個字.
這個字很吸引我們的眼球.
一說到逼迫.
第一件事會想起什麼呢.
我們會覺得自己很慘.
我們會覺得自己有很多的困難.
我們活在一個充滿鬥爭的世代.

$^{241}$我想這個我們都不需要否定.
不過我自己平心而論.
我這一代.
我們信耶穌的香港這一代.
相對於我上一代.
或者再早幾十年前.
或者一百年前的我們的先賢.
今天真的舒服很多.
今天我們教會的冷氣.
冷到要穿.
夏天的時候我們都要穿外套的.
如果舊時幾十年前沒有冷氣.
你說要秘密聚會.
在一個暗室裡面.
還要逼到不得了.
沒有社交距離的限制.
那時候大家只是渴望.
在一個很極端的環境當中.
可以爭取到聽到上帝的話語.
甚至很多人因為信耶穌.
被人笑耶穌仔.
被人作弄.
上一代.
或者我們之前的那些.
信的先賢.
他們面對的東西.
和我們今天真的很難完全比較.
所以我今天我不會說.
我們真的完全活在一個迫不及待的世代.
說到我們自己真的好像.
一百幾十年前.
我們先賢面對的那種困難.
不過.
為了我們願不願意付上代價.
這個對我們來說就很有營養價值.
今天可能我們沒有了一些外在的困難.
所以普遍來說.
信仰對我們變得.
就是首到難來.
信耶穌成為一個很容易得到的恩典.

$^{281}$甚至在我們生活當中.
我們很少為了信耶穌要付上代價.
我們可能過去沒有想過要付上代價.
不過去到這段經文裡面.
他重新問我們.
究竟這個信仰是否開著冷氣.
拿著酒杯底.
大家很舒服地.
來到去快快樂樂地過我餘生.
還是我們生命裡面.
耶穌期望我們.
去為著這個福音.
為著祂的道德旨意.
我們自己要付上一點代價.
我們願意為這個信仰.
我們可能犧牲我們生命裡面的一部份.
我們覺得順理成章的東西.
我再繼續來說.
在這段經文裡面.
其實神學家潘博華.
我們經常讀的.
他講到今天.
很小的基督徒很多時候面對一個問題.
我們領受一個cheap grace.
廉價的恩典.
我們覺得上帝只不過是.
在我們覺得不舒服的時候祈求.
或者是一個.
生命裡面很安逸的環境當中.
有時候有些無病身淫之下.
呼求的那個上帝.
那個合逢中的上帝.
潘博華說.
如果我們純粹的信仰壓縮到一個這樣的.
那一區中的信仰的時候.
上帝在我們生命裡面.
完全沒有任何的份量.
反而是我們.
當領受了這個寶貴的恩典的時候.
我們願意為祂來擺上.

$^{321}$那個生命才能夠真正的經歷上帝.
今天很多基督徒或者很多弟兄姊妹說.
我想經歷上帝啊.
我們想著那個經歷上帝.
就是純粹有些東西很順頭順腦.
或者天上有一些.
很奇特的神蹟騎士.
淋在我身上為之.
那些叫做經歷上帝.
不過基督耶穌特別在登山裡面說.
如果真是經歷上帝的人.
第一件事去問他願不願意為這個信仰擺上代價.
他願不願意為著福音的緣故.
他願不願意為著上帝的公義.
上帝視為重要的道德價值.
在他生命裡面他持守.
為了有可能持守這些立場的時候.
他面對生命裡面有些困難.
但是他都在所不辭.
當他願意這樣做的時候.
他才能夠經歷上帝的豐富.
經歷上帝的生命價值.
用這個百福來說.
他才能夠經歷到天國.
如果生命裡面不願意為上帝付上任何代價.
大概我們都沒有辦法經歷到天國.
在我們生命裡面的主權.
就是那種來自神的主權.
在我們生命裡面的彰顯.
在這段經文裡面.
和其他八福不同.
耶穌講完第五章第十節之後.
耶穌繼續的即刻來做一些補充.
他就為著這個畏異仇迫伯再繼續的來對他說.
人若因我辱罵你們.
迫伯你們.
捏造各樣的壞話誹謗你們.
你們就有福了.
應當歡喜快樂.
因為你們在天上的賞賜是大的.

$^{361}$在你們以前的先知.
人也是這樣迫伯他們.
五章十一至十二節.
在這裡面.
耶穌再繼續的來講.
為何人追求神的旨意.
追求神的義要面對迫伯呢.
他第一件事是他立刻告訴你.
因為人因為他而辱罵這一班跟隨者.
因為耶穌.
單單是因為耶穌這個名字.
因為我跟隨他.
不需要任何罪名.
不需要看我的行為.
這個世界就是帶著一種這樣的偏頗的價值.
來判斷一個人.
扣了一個人帽子.
這個人就是一個耶穌仔.
所以就是這樣面對著生命裡面承受的迫伯.
因為耶穌的緣故.
我記得有一點點噁心.
以前我認識一些基督徒.
早一代他們在外面吃飯.
被同事作弄.
就是一起吃飯的時候.
因為要借飯祈禱.
其實說得上是甚麼見證.
當他們低頭祈禱.
因為他是基督徒.
身邊的人還不是這樣.
就作弄他.
當他祈禱的時候.
就在他口中吃的那一口飯.
就吐口水.
混入他自己正在祈禱的弟兄那碟飯.
他祈禱之後吃.
那碟飯也挺好吃的.
其實他不知道很多人的口水.
已經混入了.
那些人就在笑.

$^{401}$每一次那個弟兄問.
為甚麼每次吃飯之後.
身邊的人都在笑.
他又不覺得有甚麼.
後來他有一次.
祈禱之後特意嘲笑.
閉上眼睛.
但又不祈禱.
那些人在做甚麼呢.
就在做剛才那些舉動.
就知道被人作弄了.
這些都是小作弄.
更加大的.
其實不是說幾十年前的事.
是今天的事.
今天在.
不是在香港.
在很多稱為基督教的國家.
外國的社會裡面.
當你稱自己為基督徒的時候.
很多時候就會引起那些仇恨言論.
在你的社交媒體.
你自己說基督徒的價值.
按照聖經怎樣教.
將這些說話分享出來的時候.
隨時會受到一些叫做haters.
那些仇恨言論攻擊你.
你想說一句對的話.
會大量蜂擁而至.
你越出名.
就越大量的.
這些haters.
這些仇恨者.
就在網上攻擊.
更加不要說在大學社會.
是一個最開放自由的.
當你強調你自己基督徒生命的價值.
當強調你自己是一個基督徒的時候.
單單是為了這種身份.
都已經大量的人.

$^{441}$質你.
嘲弄你.
這個也包括我自己在內.
不是我被人嘲弄.
而是我嘲弄別人.
在我讀大學的時候.
其中一個最喜歡作弄的.
就是我.
在我未信主的時候.
有些基督徒會走來和我談道.
但是我喜歡令他們很難堪.
說一些話.
令他們回答不到的時候.
覺得他糟蹋他們.
那個感覺真的很過癮.
不過當我信了主的時候.
特別是想起之前學完傳道會.
有幾個弟兄走來找我傳福音.
令他們面灰灰離開.
但當我信了主的時候.
我發覺將來我也要學他們一樣.
我不是覺得那些人很可惜.
而是覺得那些人.
那些才是真正的生命.
那些才是真正的信仰英雄.
罕罕而談.
那個並不是說那些真正的信仰英雄.
會在裡面灰灰地.
承受著很多人的調侃.
或者是有些人來嘲弄.
但依然逆流而上.
將福音帶給人那些.
那些才是真正的信仰英雄.
基督徒面對逼迫原因.
因為世人拒絕了基督.
單單是因為基督耶穌的緣故.
就逼迫他們.
第二個可能性.
提摩泰後書第四章三至五節.
因為時候要到.

$^{481}$人必厭煩純正的道理.
以躲發揚.
就隨從自己的情慾增添好些斯苦.
並且掩耳不聽真道.
偏向荒謬的言語.
你卻要凡事謹慎.
忍受苦難.
做全道的功夫.
盡你的職份.
保羅提醒提摩泰.
在一個末世的世代裡.
時候到了.
大概是耶穌升天之後.
又再回來之前的這段日子.
有什麼事發生呢?.
人們對於純正的道理.
對的事.
越來越覺得不中聽.
覺得那些歪的.
譁眾取寵的.
那些挑換你的情慾的.
反而要拜師學藝.
人會將自己所追求的東西.
將自己的人生.
寧願交給控制人的情慾的師父.
都不去跟一些真正傳揚.
上帝道理的師父.
那些不中聽的.
通常都是用在耶穌的道理當中.
聖經的話語已經等同於沉悶.
或者等同於說教等等.
人對於純正的道理.
以躲法尋.
當然我們也要反省.
我們有時候說上帝的道.
會不會說到他悶了.
可能也是.
不過我要問的核心是.
很多時候我們覺得對我們生命重要的.
或者這個來自上帝的救恩.

$^{521}$來自基督耶穌的拯救的恩典.
聖經裡面的教導.
我們經常覺得我們聽過了.
我小學時已經讀基督教小學.
我已經明白了.
不要再跟我說這些說教的說話了.
對於純正的道理.
以躲法尋.
所以在壞事裡面.
這個世代的人.
寧願在壞事情慾的事裡面.
為自己添更多的師父.
早前有機會跟一個YouTuber.
網紅聊天.
這個網紅是正路的網紅.
也不是很紅.
但他說什麼呢.
他說今天在網上的世界.
在YouTube.
在媒體的世界裡面.
怎樣令到你的點擊率增加呢.
他說其實離不開兩樣東西.
娛樂性和話題性.
他又說自己沒有娛樂性.
就說話題性.
但說話題通常.
什麼叫做最多人談論的話題呢.
通常不是純正的道理.
你按照聖經的角度.
或者按照聖經的故事.
聖經的敘事.
聖經的價值觀.
去說你自己純正的道理.
你預計了你只會自己為內.
自己家人看而已.
因為大部份的人.
今天去尋訪的.
也不是尋訪這些訊息.
反而你說爭拗.
說一些很熱議.

$^{561}$大家為之罵的題目.
在搜尋網裡面.
你一看這些字.
全部都是熱議搜尋的.
用這些字來做你的節目.
你就多人看了.
不過到最後是什麼呢.
就是做一場大龍鳳.
就罵來罵去.
到最後就好像.
抒發了心裡的鬱結.
罵了一部份.
但對生命有沒有什麼幫助呢.
這個youtuber和我說的慨嘆.
他說今天要做真正的內容.
對生命有營養價值的內容.
沒什麼人看.
說的那些很熱議.
很有話題性的其他題目.
大家就會一直按著來看.
這個現象可能都反映在我們的生命裡面.
各位弟兄姊妹.
今天對我們生命重要的東西.
我們願意放多少力量去追尋呢.
又或者在你生命裡面.
你講這些話的時候.
會不會有些人覺得你很悶.
悶蛋.
所以就拒絕你.
甚至會整蠱你.
玩你.
覺得你是這個世界不入流的人.
耶穌基督說.
當這個世界是拒絕基督的世代.
是一個拒絕真理的世代.
真的跟隨耶穌的人難免要面對這些逼迫.
甚至要為走這條窄路的時候.
擺上很多生命的代價.
不過呢.
耶穌不是強調這個世界有多恐怖.

$^{601}$人有多逼迫.
耶穌的蝙蝠其實不是說很多被逼迫.
那種痛苦那些東西.
耶穌不是想鋪陳一個這樣的世界觀給我們.
接下來你會覺得有一點變態.
如果那些人逼迫你們.
捏造各樣的壞話.
毀謗你們.
你們有福了.
應當歡喜快樂.
因為你們在天上的賞賜是大的.
在你們以前的先知.
人都是這樣來逼迫他們.
耶穌要轉換我們的角度.
我們明白這個世界會這樣來逼迫基督徒.
或者我們為了信仰要付上代價.
不過當你付上代價的時候.
你要為這件事來歡喜快樂.
如果別人毀謗你.
你更加知道你的生命是有福的.
為什麼呢.
在逆流而上裡面.
你才知道原來你才是真正的中流砥柱.
如果你只是隨波逐流.
信耶穌信得順風順水.
沒有什麼意義上的挑戰和危機.
大概可能你不是很認真地跟隨這個信仰.
耶穌說反而面對這個世代裡的逆流.
你知道你的生命是有福的.
這個福氣除了在今生裡知道.
你真的站穩在聖經的教導裡.
耶穌再更加遠一點.
不是說不實體.
他說將來這些人在天國裡得到極大的賞賜.
可能在今生裡他要承受代價.
但是在永恆裡.
每一個站得住腳的人.
在天上裡都得到大的賞賜.
這種賞賜不是純粹說.
讓你感覺上奢華一點.

$^{641}$而是在天國裡.
我們慧烈和昔日一起被迫害.
和我們看的上帝屬靈偉人的先知.
我們的生命和他們一樣.
他們也是被迫.
他們也是逆流而上.
如果我們今天繼續逆流而上的話.
我們的生命在上帝眼中和他們是對等的.
如果上帝要賞賜他們.
同樣都要賞賜今天.
你願意持守信仰的每一個弟兄姊妹.
你都會得到天上極大的賞賜.
今天我們在地上營營役役.
我們今天在生活當中.
我們做很多事情.
我們希望做一些事情有價值.
買一層樓我們都想升值.
我們做每一件事都是希望.
為了我們今世去做一些儲蓄和投資.
我們希望今天所擺出的將來有更大的回報.
如果我們在地上日子我們都這樣想我們的生活.
我們更加要去問.
在我們今世裡面我們所努力預備的.
在我們的永生將來在天國裡面.
我們更加要努力地為自己去積存.
今天如果我們純粹看地上的富足生活.
看永恆的角度.
我們這個眼光實在太短淺.
弟兄姊妹要鼓勵你.
在我們今天這個世代裡面.
我們需要不單止是為我們今世的人生來努力.
我們更加要為我們的永恆.
我們做一些有永恆價值的努力.
為著上帝的道德旨意.
為著上帝的救恩.
站立得穩.
我們才能夠站立起來.
我們才能夠和先知並列.
一起得到天上的賞賜.
接著你會說.

$^{681}$牧師你今天解釋的東西好像很理論性.
我們可以怎樣實踐.
耶穌也不是一個純粹的理論家.
然後那個蝙蝠.
其實在登山部分裡面.
他講完這個引言之後.
後面就講一些更加具體的.
怎樣去實踐面對著一個逼迫世代裡面.
或者特別是針對當時的法利賽人.
怎樣在一個這樣的社會裡面.
來去活現一個上帝喜悅.
得到賞賜的生命.
為什麼耶穌會用法利賽人.
你留意一下.
為什麼在馬太福音也好.
特別是馬太福音很多的.
特別是他有一個假想的對象.
來去拆解的.
講很多道德的生活.
他特別用針對法利賽人來說的.
為什麼呢.
他不是和法利賽人有仇的.
不過法利賽人就是當時.
猶太社會裡面的道德代言人.
不過在耶穌看到這群道德代言人當中.
背後或者是這群人的雜性.
在他們承傳律法.
然後再加上口傳律法.
到他們那種迂腐.
甚至是嫉妒真理的心.
耶穌一一看在眼內.
如果這群道德代言人成為接下來.
神國子民要學校的無街的時候.
耶穌說這些粉色的墳墓.
只是將人帶領去死亡.
將難擔的擔子.
割在普羅大眾身上.
所以耶穌某程度上.
他就要解構那些法利賽人所做的.
那種表面的道德行為.

$^{721}$再進入一個深層道德.
究竟上帝的旨意在這群人生中.
究竟是一回什麼事呢.
我怎樣去實踐呢.
耶穌說了幾件事.
第一件事.
他第一個給的錦囊就是要超越.
五章二十節.
我告訴你們.
你們的義若不勝於民事和法利賽人的義.
斷不能進天國.
剛才說的.
為義受逼迫的天國就是他們.
另一面就是說.
如果我們的義不勝過民事和法利賽人的義.
我們不能夠進入天國.
耶穌在這裡不是否定了.
那些法利賽人所做的行為.
他沒有說那些人是錯的.
不過他所做的那些只不過是.
移民救養.
是一些表面的事情.
如果我們跟隨神的人.
跟著這些移民救養.
跟著他們的形式來過我們的信仰生活.
是不可以的.
進入不了天國的.
我們今天每一個要超越.
這群法利賽人所生命所行過的義.
他怎樣說超越呢.
其實接下來五章開始.
十七節開始.
一路打下去整個登山寶訓.
他也是在說這件事.
給幾個例子你們知道.
耶穌說我不是要廢掉先知或律法.
我們不是要廢掉的.
我們是要承傳的.
律法裡面一點一滑都不能夠少的.
所以這些事情一定要教訓人來進行.

$^{761}$不過面對著那些仇恨.
或者你憎恨的弟兄.
令你向弟兄動怒的.
難免你要受審判.
所以罵不是打.
罵弟兄是拉街的就是說他垃圾.
難免受到工會的審判.
說弟兄是摸你的.
摸你的其實很多時候我們都不知道.
簡單用廣東話.
謀你的你就知道了.
意思就是這樣.
說弟兄是謀你的.
難免他會受到地獄之火.
這裡他不是在說跟他打架結怨.
而是說你要超越法利賽人的看法.
你怎樣跟人相處.
如果你憲制的時候想起有個弟兄跟你懷怒.
跟你有怨.
你就要怎樣憲制.
放下那些禮物在壇前.
跟你的弟兄和好了.
然後才跟上帝來去憲制.
你要超越.
他再繼續說.
他說犯奸淫.
以為有一些不聞不慮的婚前性行為等等.
不是的耶穌說.
見到婦女動人念.
你已經犯了奸淫.
耶穌一直說.
他在說我們生命的標準.
要比法利賽人比民事所說的.
他們做了很多規條.
不過你生命的義你對上帝的追求.
是要比他們更進一步更加超越.
這個不是要求別人去做的.
這個超越其實是反求諸己的.
法利賽人除了說口述傳統.
他要製造一個外在的框框.

$^{801}$一個道德社會.
一個移民的社會之外.
其實他們整個社會就在面對一個困局.
那些人做錯了事.
法利賽人或者民事就拿著律法.
你犯了這條.
是一個比訓導處更嚴厲的一種.
懲罰性的社會.
不過耶穌說.
律法的要求.
不是用這些律法去質人.
說人.
用律法要求自己生命得到提升.
明白上帝的美善.
透過上帝的教導.
我們重新發現上帝的美善旨意是怎樣.
而針對的對象不是隔離的那個.
或者我家裡的老公老婆有多差.
往往是針對自己.
反求諸己.
在五章四十四節裡面.
耶穌特別的來講.
不要愛你們覺得很可愛的人.
為迫迫你們的人來禱告.
反而你們要完全像天父一樣.
同時間你要明白.
在你的生命裡面.
你不是一個所謂道德完美的人.
天父只是用他天上的慈愛來接納我們這些不完美的人.
不要看見弟兄裡面有些錯誤的行為.
看見你們那些小小的一條刺.
就看不到自己眼中的良莫.
反求諸己.
每一樣東西.
用律法來操練自己的生命.
這是耶穌給我們第二個.
面對著這個對信仰不太友善的世界裡面.
反而我們要成為一個更加被他們超越的.
一班行在上帝旨意的一班人.
最後一點.

$^{841}$我們追求神的義.
不是要將這個世界鬧到稀巴爛.
這樣才算是痛快.
大家知道.
我曾經在突破裡面做過市場營銷.
其中一句說話不是市場營銷.
而是我內心很重要的一句說話.
與其就坐黑暗.
不如燃燒自己.
蘇銀沛在突破創立的時候.
說了一句這樣的說話.
這句說話成為了整個機構.
甚至成為了我今天.
我行事為人.
我們自己怎樣做人的指標.
這句說話說到這個世界有些黑暗.
你鬧多它.
或者你面對著這些不理想的世界.
你可以繼續咒坐它.
你可以罵它.
你每次祈禱就咒坐它.
不過我這樣做有什麼價值.
成為一個粉身碎骨的人.
成為一個你做初一.
我沒辦法做十五.
但我在心裡恨你.
對這個世界有什麼幫助.
既然如此.
為什麼不轉念.
在我生命當中.
我燃燒自己.
發出一點點的光芒.
在一個不理想的世界當中.
我和它們有些不同.
在它逼迫我也好.
在它質我也好.
甚至我要付上代價也好.
我喜歡我喜愛上帝的義.
所以我願意放上代價.
可以牙刷一點的表達.

$^{881}$當別人覺得我們沒價值的時候.
我們說我喜歡我喜愛上帝的義.
這就是我們生命的態度.
在《箴言》第24章19-20節.
不要為作惡的心懷不平.
也不要嫉妒惡人.
因為惡人終不得善報.
惡人的燈也必熄滅.
這句話說什麼呢.
今天面對這個世界.
可能有很多事情很荒謬.
你咬牙切齒.
但我不會因為這樣.
覺得那些人真是不義.
但心裡在罵什麼呢.
因為這些不義而享受到那些成果.
我心裡是在嫉妒它.
如果沒有其他規範.
可能我都像它一樣.
在聖經裡說的嫉妒惡人就是這個意思.
我們說的公義或義.
為什麼呢.
我們心心不忿.
就是因為它得到多一點.
而不是因為義本身有重要的價值.
今天我們看見這個世界.
不要咬牙切齒.
憤憤不平.
或者覺得別人的東西虧欠了你.
覺得自己虧了很多.
今天的虧本可能是表面的.
但我生命裡.
上帝的義在我生命中沒有虧蝕.
我追求上帝的義.
是一個好的美善的旨意.
我甚至乎.
我願意在生命裡擺上代價.
我都要追求一些好的.
上帝美善的旨意在我生命中彰顯.
弟兄姊妹.

$^{921}$今天這一段的經文大致說到這裡.
我問大家幾個問題.
神的義對你的份量有多少.
你同不同意.
或者你覺不覺得這個義是好東西.
追求這份有長遠永恆價值的好東西.
你願意擺上多少代價.
你願不願意在今天.
似乎走這條路很不容易的過程當中.
依然笑著感恩的心.
去實踐出來.
因為知道你和這個世界不同.
在我們今天最後的問題.
今天聽完這篇道之後.
你怎樣願意為上帝的義.
來付上代價.
不要說逼迫了.
付上代價.
你願意今天開始怎樣去為上帝的義付上代價.
求主繼續幫助我們.
激勵我們.
\newpage



\section{哥林多後書 8:9-24}
\label{sec:zila_5Bj4Bo}
\textbf{2021年3月14日崇拜直播｜梁家麟牧師｜均平的原則｜哥林多後書8\_9-24}
\newline
\newline
連結: \href{https://youtube.com/watch?v=zila_5Bj4Bo}{\texttt{https://youtube.com/watch?v=zila\_5Bj4Bo}} ~~~~ 語音日期: 2021-03-14
\newline
\newline
\hyperref[sec:rPutWeJ6_Xg]{\small{< < < PREV SERMON < < <}}
~
\hyperref[sec:index_chronic]{\small{[返順時目]}}
~
\hyperref[sec:qJa4QRG5SV0]{\small{> > > NEXT SERMON > > >}}
\newline
\newline
哥林多後書 8:9-24
\newline
\begin{longtable}{cl}
\hline
\hline
章節 & 經文 (和合本修訂版)\\
\hline
8:9 & \begin{tabularx}{0.7\textwidth}{X} 你們知道我們主耶穌基督的恩典：他本是富足，卻為你們成了貧窮，好使你們因他的貧窮而成為富足。 \end{tabularx} \\ \\ \relax
8:10 & \begin{tabularx}{0.7\textwidth}{X} 我在這事上把我的意見告訴你們，是對你們有益，因為你們開始辦這事，而且起此心意已經有一年了。 \end{tabularx} \\ \\ \relax
8:11 & \begin{tabularx}{0.7\textwidth}{X} 如今就當辦成這事，既然有願做的心，也當照你們所有的去辦成。 \end{tabularx} \\ \\ \relax
8:12 & \begin{tabularx}{0.7\textwidth}{X} 因為人只要有願做的心，必照他所有的蒙悅納，並不是照他所沒有的。 \end{tabularx} \\ \\ \relax
8:13 & \begin{tabularx}{0.7\textwidth}{X} 我不是要別人輕鬆，你們受累，而是要均勻： \end{tabularx} \\ \\ \relax
8:14 & \begin{tabularx}{0.7\textwidth}{X} 就是要你們現在的富餘補他們的不足，使他們的富餘將來也可以補你們的不足，這就均勻了。 \end{tabularx} \\ \\ \relax
8:15 & \begin{tabularx}{0.7\textwidth}{X} 如經上所記：多收的沒有餘，少收的也沒有缺。 \end{tabularx} \\ \\ \relax
8:16 & \begin{tabularx}{0.7\textwidth}{X} 感謝神，把我對你們的熱忱同樣放在提多心裡。 \end{tabularx} \\ \\ \relax
8:17 & \begin{tabularx}{0.7\textwidth}{X} 他固然聽了我的勸告，但自己更加熱心，自願往你們那裡去。 \end{tabularx} \\ \\ \relax
8:18 & \begin{tabularx}{0.7\textwidth}{X} 我們還差遣一位弟兄和他同去，這人在傳福音的事上得了眾教會的稱讚； \end{tabularx} \\ \\ \relax
8:19 & \begin{tabularx}{0.7\textwidth}{X} 不但這樣，他也被眾教會選派跟我們同行，把所交託我們的這捐款送到了，為的是榮耀主，也表明我們的好意。 \end{tabularx} \\ \\ \relax
8:20 & \begin{tabularx}{0.7\textwidth}{X} 我們這樣做，免得有人因我們收的捐款多而挑剔我們。 \end{tabularx} \\ \\ \relax
8:21 & \begin{tabularx}{0.7\textwidth}{X} 我們留心做好事，不但在主面前，就是在人面前也是這樣。 \end{tabularx} \\ \\ \relax
8:22 & \begin{tabularx}{0.7\textwidth}{X} 我們又差遣一位弟兄同去。這人的熱忱，我們在許多事上屢次考驗過，現在他因為深深信任你們，就更加熱心了。 \end{tabularx} \\ \\ \relax
8:23 & \begin{tabularx}{0.7\textwidth}{X} 至於提多，他是我的夥伴，為服事你們作我的同工。至於那兩位弟兄，他們是眾教會的使者，是基督的榮耀。 \end{tabularx} \\ \\ \relax
8:24 & \begin{tabularx}{0.7\textwidth}{X} 所以，你們務要在眾教會面前向他們顯明你們的愛心和我所誇獎你們的憑據。 \end{tabularx} \\ \\
[1ex]
\hline
\hline
\end{longtable}
$^{1}$讓我們繼續聆聽上帝的話語.
記載在哥林多後書第八章九至二十四節.
斯圖普羅這樣勉勵提醒.
在哥林多教會的弟兄姊妹.
你們知道我們主耶穌基督的恩典.
祂本來富足.
卻為你們成了貧窮.
叫你們因祂的貧窮可以成為富足.
我在這事上把我的意見告訴你們.
是願你們有益.
因為你們下手辦這事.
並且起此心意已經有一年了.
如今就當辦成這事.
既有願作的心.
也當照你們所有的去辦成.
因為人若有願作的心.
必蒙月立.
乃是照他所有的.
並不是照他所無的.
我原不是要別人輕省你們受累.
乃要均平.
就是要你們的富裕.
現在可以補他們的不足.
使他們的富裕.
將來也可以補你們的不足.
這就均平了.
如經上所記.
多收的也沒有餘.
少收的也沒有缺.
多謝神.
感動提多的心.
叫祂待你們殷勤.
像我一樣.
祂固然是聽了我的勸.
但自己更是熱心.
情願往你們那裡去.
我們還打發一位弟兄和他同去.
這人在福音上得了宗教會的稱讚.
不但這樣.
他也被宗教會挑選.

$^{41}$和我們同行把所託予我們的這捐資送到了.
可以榮耀神.
又表明我們樂意的心.
這就免得有人因我們收的捐錢很多.
就挑我們的不是.
我們留心行光明的事.
不但在主面前.
也在人面前.
就在人面前也是這樣.
我們又打發一位弟兄同去.
這人的熱心.
我們在許多事上屢次試驗過.
現在他因為深信你們.
就更加熱心了.
論到提多.
他是我的同伴.
一同為你們勞碌的.
論到那兩位弟兄弟.
他們是宗教會的使者.
是基督的榮耀.
所以你們務要在宗教會面前.
顯明你們愛心的憑據.
並我所誇獎你們的憑據.
今天早上有梁家倫牧師在我們當中.
繼續宣講剛才那段經文.
講題是「均平的原則」.
恭請梁家倫牧師.
弟兄姊妹早晨.
我們繼續看哥林多後書八到九章.
我會用三次的時間去講.
保羅談到金錢捐獻.
講到詞彙的工作.
我們這裡是第二篇.
不是我想囉囉嗦嗦講這個課題.
講多一點.
去賺錢.
不是.
因為聖經在這裡有很多很豐富的教導.
值得我們慢慢去思考.
保羅繼續勸導哥林多信徒.

$^{81}$為耶路撒冷信徒的生活需要去捐款.
一般來說.
當我們呼籲人為有需要的人捐款的時候.
我們總是會強調那些有需要的人的悲慘情況.
他們多麼可憐.
多麼迫切需要別人的援助.
藉以激起那些被勸捐的人那些憐憫的心腸.
然後慷慨解囊.
是不是.
就像那些慈善機構.
呼籲人去捐助非洲的災民.
就會刊登一些瘦骨淪衰.
眼神無助的那些非洲人的照片.
小孩兒童是最容易激發人同情心的對象.
所以慈善機構總是特別報導.
那些代救援國家的兒童的悲慘故事.
他們缺衣缺食.
他們沒有受教育的機會等等等等.
這裡我沒有反對慈善機構慣用的宣傳策略.
有效用的就是好策略.
沒有什麼好批評的.
我只是想說的是.
當保羅多次在不同的聖經地方.
為耶路撒冷.
捐的時候.
他完全沒有用過這樣的策略.
他沒有提到耶路撒冷的信徒.
如今落在一個多麼悲慘的境況.
沒有提到如果其他教會袖手旁觀不幫的話.
他們會有多少人餓死.
保羅根本沒有將論述的重點.
放在那些代救援的對象之上.
沒有幾句話提到耶路撒冷教會的情況.
保羅僅僅是集中在那些要募捐的對象.
指出他們需要作出捐助.
捐助對他們的信仰有多大的重要性.
換言之保羅關心的不是.
別人需要我們捐助.
而是我們需要對人捐助.
這是一個很值得注意的現象.

$^{121}$我在上次講完道.
2月28日那一堂道.
那一堂叫做慈慰的精神.
就是第一篇.
講完之後.
我坐下來.
一起唱回應詩.
然後我才有這個發現.
那天的回應詩是.
讓你的心破碎.
第一節的歌詞就是.
讓世界的需要破碎你的心.
餓者得飽足.
傷者得安慰.
送上你的乾糧.
施出一杯水.
服侍人如耶穌.
祂有使用你.
唱的時候.
我突然發覺.
我剛才講的訊息.
跟這段歌詞完全不配合.
我沒有提及過一句.
那些代救援者的悲慘.
沒有提及到要對.
其他人施以慈悲憐憫.
為什麼呢.
講慈慰.
為什麼不提及這個世界有多少人.
住在飢餓狀態裡面呢.
講慈慰怎麼可以不提及.
香港有多少人.
住在tong 房.
或者露宿街頭呢.
怎麼可以不提及這些呢.
然後我.
立刻很仔細地看了自己的講章.
再看一看《斷聖經》.
我確定我是依書直說的.
你可以查看.

$^{161}$我是如實地.
跟著保羅的講述內容.
思路去講的.
我沒有遺漏什麼.
如果是這樣的話.
沒有強調受助者的那個悲慘情況.
有什麼迫切需要.
這個不僅僅是我.
是保羅.
保羅為什麼這麼粗心大意呢.
為什麼會遺漏這個.
這麼重要的重點呢.
因為.
保羅.
在前面那段聖經.
保羅提出.
慈慰是任何人.
在任何時間都可以參與的行動.
只要你給那些有需要的人富有.
只要你自覺得著上帝豐富的恩典.
你就應該參與的了.
慈慰是一個權利來的.
金錢捐獻使我們跟人.
結連團契.
成為其中一份子.
金錢捐助使我們的信仰生命得以完全.
實踐.
以及信仰成熟的重要部分.
這三點.
都是強調.
慈慰金錢捐獻對捐助者的意義.
沒有任何一點是關乎.
受助者有什麼價值.
是不是.
所以.
簡單的說法是.
我們不是因為看到人有需要.
然後才突然想起要捐獻.
而我們本來就有捐獻的需要.
我們要跟其他人分享我們自己所有的.

$^{201}$我相信.
保羅不會是原則上.
反對提出別人的需要的.
然後喚起其他人關注和同情心的.
我相信他之前.
就是跟.
哥林多信徒一定有說過.
就是.
耶路撒冷教會的情況.
是不是.
如果完全不知道.
就是耶路撒冷弟兄姊妹有多大的需要.
為什麼無端端哥林多信徒要捐錢呢.
對不對.
不會完全沒有提過的.
你不可以說他們盲從保羅的話.
保羅要他捐什麼就捐什麼.
就是完全不用你捐去做什麼.
不會這樣的.
是不是.
不過.
我相信是保羅以前曾經提過.
他也假設哥林多信徒已經知道捐款的目的.
這裡他就不再重提.
保羅在這裡最關心的.
是哥林多信徒的信仰光景.
而他們是不是願意參與捐獻.
是他們信仰很壞的其中一個寒暑表.
金錢捐獻.
幫助別人.
是基督徒的屬靈生命的重要實踐.
也是他們信仰.
是不是達到完全的一個試金石.
所以保羅要強調的是.
哥林多信徒需要捐獻.
而不是耶路撒冷信徒有待他們捐獻.
正如保羅在羅馬書.
15章27節.
說到.
小亞細亞的信徒要捐助耶路撒冷的時候.

$^{241}$他怎麼說呢.
「這固然是他們樂意的.
其實也算是所欠的債.
因為外邦人既然在他們屬靈的好處上有份.
就當把養生之物供給他們」.
所強調的都是捐獻者需要施予.
而不是受助者需要幫助.
外邦人的教會欠了猶太人教會的恩情債.
如今要藉著做慈惠來償還.
做慈惠不過是還債.
欠債還錢是天經地義的.
難道你要答謝還錢的人嗎?.
是不是?.
這是第一點我們在這兩章聖經看到的東西.
然後我們來看第二個方向.
保羅強調.
叫我們甘心做我們能力所及的善事.
在八章十一節.
保羅說.
「因為你們下手辦這事.
而且起此心意已經有一年了.
如今就當辦成這事.
叫願作的心也當照你們所有的去辦成」.
保羅去年已經委託提多在哥林多教會推動募捐.
哥林多信徒也曾經表示願意參與.
保羅如今寫信催他們.
好頭好尾將這件事盡快辦成.
我們不少人也有這樣的經驗.
我自己試過很多次.
曾經有感動要為一個人或一個機構.
一個事工去奉獻.
不過很壞沒有立刻行動.
沒有立刻寫支票.
結果過了兩三天感動就沒有了.
又覺得自己先吧.
是不是?.
所以如果我們相信.
為善的動機是出於聖靈的感動的話.
就要小心不要將那個感動消滅了.
曾經立願起此心意.

$^{281}$保羅提醒他.
你就要努力還願付諸實踐.
信心沒有行為是死的.
保羅特別強調.
哥林多教會是願意做這件事的.
起此心意有願作的心.
這是他們自願的.
不是出於被迫勉強.
下次我們會特別提到這一點.
甘心樂意的捐獻是第一要緊的原則.
辭位如果出於勉強的話.
就一切價值都蕩然無存了.
八章十二節保羅說.
因為人若有願作的心必蒙月立.
乃是照他所有的.
並不是照他所務的.
這句話重提奉獻的兩個條件.
很簡單的.
第一是有心.
第二是有力.
有願作的心就是有心.
照他所有的就是有力.
雖然前面保羅提到.
馬其頓教會是過了力量的去捐獻.
不過他沒有用這個例子作為奉獻的標準.
多數人其實是不需要過了力量去奉獻的.
不需要像窮寡婦那樣.
把養生的,買菜的錢,吃飯的錢都捐了.
只要按著我們的感動和能力.
跟別人分享我們有餘的就很好了.
我們偶然會歌頌一些超凡的見證故事.
平日我們只需要做凡人的信仰實踐.
日日叫背十字架為主寫明.
開口埋口搬霍華.
遇上國安法的頒布.
就第一時間快點走.
本福華第一句說話就是.
如果叫基督呼召我們,就是叫我們去死.
如果我們不是真的這麼信,又這麼做.
我們說這些話,很多時候只會造反見證.

$^{321}$既叛倒信徒的心,又讓非信徒.
得到嘲諷的把柄.
寧可踏實一點.
盡量拒絕喊那些激情的口號.
專注做我們能力所及的事.
上帝是合情合理的上帝.
他對人的要求總是照他所有的.
並不是照他所無的.
上帝不會無理取鬧.
耶穌沒有計較窮寡婦所奉獻的兩個小時是多麼寒酸.
遠遠比不上財主奉獻的數額.
他是讚嘆他所做的已經是他能力所及.
新聞記33章25節的原則.
日子如何,力量如何.
保羅在哥林多前書10章13節的原則.
我們所受的試探.
都是我們所能夠承受的.
我們做的事,上帝對我們的要求.
永遠是在我們的潛力之中.
我們自己的可能.
上帝將來不會怪責我為何不豪捐一億兩億.
他會責問我.
為何不向你地上困苦窮乏的弟兄鬆開手.
說完照他所有的.
不是照他所無的原則之外.
之後保羅再說到另一個原則.
就是君平.
8章13到15節.
我原不是要別人輕省你們受累,乃要君平.
就是要你們的婦孺現在可以補他們的不足.
使他們的婦孺將來也可以補你們的不足.
這就是君平了.
如經上所記.
多收的也沒有餘,多少收的也沒有缺.
保羅強調他無意將哥林多信徒的財富轉移到耶路撒冷信徒身上.
使哥林多信徒由富有變貧窮.
使耶路撒冷信徒由貧窮變富有.
他當然是希望哥林多信徒在滿足了自己生活的需求之後.
有些婦孺.
就拿那些婦孺來跟其他有需要的人分享.

$^{361}$補他們的不足.
換言之,哥林多信徒只需要用他們餘下的資源跟其他人分享就夠了.
辭位工作不會太影響他們的生活水平.
當然,如果哥林多信徒能夠像馬其頓信徒那樣.
在極窮之間還格外顯出他們落捐的口音.
他們按著力量又過了力量自己甘心落意的捐助.
這就很難能可貴.
這很比耶穌稱讚奉獻兩個小錢的窮寡婦.
不僅是奉獻有餘,連養生也拿了出來.
不過聖經仍然只鼓勵我們做能力之內的奉獻.
沒有將過於能力的奉獻變成一個宗教的規定.
保羅其實是用一個很小心的態度去談捐獻的問題.
在八章十節,他很鄭重地說.
我在這事上把我的意見告訴你們,是與你們有益.
在捐獻的事上,保羅沒有從上帝來的絕對真理.
他只是憑他個人豐富的信仰和事奉經驗得出來的意見.
意見是可以作為參考,但不需要奉為金科玉律.
保羅相信他的意見對哥林多信徒是有益的.
在哥林多前書十章二十三節,保羅曾經說過一個倫理原則.
凡事都可行,但不都有益處.
用益處來考慮一件事,要做不要做.
這是一個實用主義或功利主義的考慮.
不是絕對意義的真假對錯.
有沒有好處?不是對錯.
所以有人問我,是否必須做十一奉獻?.
如果不做有甚麼效果?會否影響救恩?.
我的答案是,不知道.
那些度宿,孤寒財主將來會否在陰間受苦.
遙看窮人將來享受福樂.
如《路加福音》十六章十九至三十一節.
是否真的會這樣?.
做度宿將來一定會在地獄哀嚎設置?我不知道.
我們拒絕周濟窮人會否不小心苦待耶穌.
每日會被棄絕並接受永刑.
你記不記得《馬利福音》二十五章.
耶穌說過的山羊比喻.
我又說我不知道.
我只能說,我們彼此提醒,小心一點.
因為耶穌只是用一個預渡故事的形式去說這些.
聖經沒有很硬朗地將它變成一個清楚白紙黑字的真理.

$^{401}$不周濟窮人就落地獄.
至少我的聖經沒有,我不知道你的有沒有.
沒有這樣說.
所以我不可以找一兩個故事將它推到盡.
我覺得要小心,我不可以去到這麼盡.
我們沒辦法做一個一刀切的論斷.
總而言之,正面的說法要多一點說.
盡力幫助有需要的人.
例如《雅各書》二章十五章,聖經有很多.
求約還多,周濟窮人是多到不得了的教導.
認定上帝對我們的要求是.
行公義,好憐憫,全謙卑的心和上帝同行.
這個大家要懂得背,我們要做.
這個聖經對我們清楚的要求.
正面說這件事就夠了.
我們來看看均平這個原則.
均平其實不是這麼複雜,只是公平的意思.
Fairness,公平.
在普羅眼裡.
有些人擁有用不完的金錢或物質.
另外有些人就樣樣缺乏.
這個就是不公平的.
如果富有的人能夠將他們剩餘的分給貧窮的人.
富有的人不會因為剩餘而浪費.
貧窮的人也能得到基本的飽足.
就像十四節說的.
要你們的富餘,現在可以補他們的不足.
使他們的富餘將來也可以補他們的不足.
這就是均平了.
普羅再引述七十四日本的《昌克局記》十六章十八節.
多收的也沒有餘,少收的也沒有缺.
來證明他所說的道理.
不少學者將普羅的說話連繫到舊約的相關教導.
藉以發展出一套又一套很嚴密的聖經經濟學.
或者叫猶太人的經濟學.
例如宣稱普羅是信奉一個叫經濟平等主義.
Egalitarianism.
我覺得這些說法是過頭了.
這些解釋或許是添足的.
必須要承認普羅根本不是經濟學家.

$^{441}$舊約包括摩西和大衛在內.
沒有多少人懂得現代經濟理論.
誰懂得呢.
所以從普羅或者聖經整體的教導.
發展出一套有系統的.
值得今天參考的,不要說全盤進行.
參考的經濟法則,我都覺得是緣木求魚.
做不到的.
但是聖經是禁止放債取利的.
聖經有宣揚知足常樂.
這是在貴佛時代的經濟觀念.
我要說這些跟我們現代人的經濟觀念是格格不入的.
是不是.
今天政府不斷鼓吹消費.
你會不會覺得這是違反聖經的教導.
他派五千元你拿不拿.
如果你硬著頭皮說不拿,我會尊敬你多些.
鼓吹消費,是鼓吹消費.
救經濟就要消費.
這是現代經濟觀念.
是不是跟聖經的觀念很格格不入.
所以保羅不是一個經濟學家.
保羅的公平或是均平的觀念是常識性的.
訴諸的是人的直觀和當下的感受.
跟任何嚴謹的學術辯證是完全不相干的.
就好像我丟書包.
我們中國的詩人杜甫所描述的情況.
朱門走肉臭,路有凍死骨.
杜甫同樣對現代經濟理論一竅不通.
當他見到有人吃不完的走肉.
有人呆死街頭的時候.
他就覺得不對,不應該發生這樣的事.
這個貧富懸殊的現象是不公平的.
杜甫沒有興趣,沒有能力.
他極主要是沒有能力去研究.
富人為何自負,窮人為何餓死.
政府要用什麼方法扭轉這個分配不公平的情況.
是不是要實行共產主義制度.
杜甫什麼都不懂.
他只是傻乎乎的.

$^{481}$從一個不理想的現實.
直接跳到一個不能夠實現的理想.
安得廣廈千萬間,大被天下寒時區歡顏.
你記得他說過這句話.
這是他的憧憬,他的夢想.
你問杜甫如何做到這樣的情況.
土地從何而來,金錢從何而來.
我保證杜甫是啞口無言的.
你同意嗎?.
別吃蛋.
不過,常識和直觀的想法是不是很幼稚.
很反智,很沒有價值呢?.
又不是.
常識和直觀,intuition.
通常是和我們的良知連在一起.
緊扣著我們的人性和基本信念.
我甚至可以mystical(神秘)地說.
和我們的DNA有關.
或者和上帝做人的時候.
埋在人生命裡的natural law(自然法)有關.
所以你可以說.
常識,直觀,common sense(普通意識).
intuition(intuition).
很原始,很粗糙.
不過是赤字之情.
我認定.
人類社會能夠維持發展到今天.
所最寶貴的想法和信念.
我們同情那些缺乏的人.
我們希望有福同享.
就是這樣.
不要過分高舉那些專業精英.
professionalism(專業).
elitism(偏見).
你們個個都是精英.
不要照鏡子看過尖尖自喜.
讓人類社會文化延續到今天.
上上不是一小撮的經濟學家.
而是很多的愚夫愚婦.
當我在街頭.

$^{521}$掏錢給一個黑衣的時候.
你不要跟我說.
我這樣的做法是助長了好吃懶飛的風氣.
破壞了太弱留強的經濟法則.
我不懂經濟理論.
我出於自然簡單的憐憫的心.
就好像孟子說的.
惻隱之心.
是不是.
我想幫小忙.
做.
信心做一下.
我爸爸媽媽.
省吃儉用.
將最多的錢留給子女.
連孫子將來公樓首期都預備了.
你罵他吧.
你鞏固著社會的相對貧窮的問題.
造成最大的不公平.
你有沒有聽過.
不少學者都說.
遺產是造成相對貧窮問題.
根深蒂固的禍首.
你這樣罵我爸爸媽媽.
他傻傻的.
他一句都聽不明白.
他只是按著他的本性和良知去行事.
正如《史篇》103篇說.
父親如何連戌他的子女.
做爸爸媽媽就是想給最好的子女.
做錯了嗎.
破壞社會嗎.
不知道.
常識直觀.
保羅的想法其實很簡單.
很蠢蠢的.
太簡單.
很多時候蠢蠢的.
有人太豐富.
有人在缺乏裡面.

$^{561}$這樣是不對的.
這是不公平的.
你有三隻雞腿.
他連一隻都沒有.
這樣是不對的.
不公平的.
你不要爭辯.
你三隻雞腿都是你自己辛苦賺回來的.
獲取的過程沒有任何不公平.
問題根本不是在於.
你獲取這三隻雞腿的手段.
是不是公平.
而是當你手拿三隻雞腿.
而身邊的人手忙忙.
一隻都沒有的時候.
這樣就不公平了.
太多和沒有的並存現象就是不公平的.
這樣不公平.
是不是根源著一個不公平的經濟制度.
分配制度.
是不是和地緣經濟.
Geo-economics.
耶路撒冷貧窮.
希臘半島富裕相關.
保羅是.
全部都不明白.
他也只知道一個簡單的解決方法.
鼓勵豐富的人將他們有餘的.
和那些貴佛的分享.
你最少將其中一隻雞腿給他.
大家都可以吃.
這樣就公平了.
問題看得很簡單.
解決方法也很簡單.
保羅真是too simple and too naive.
真的.
小時候.
媽媽幫我和我妹妹.
各人買了一支雪條.
我妹妹拆雪條袋的時候.

$^{601}$小心一掉.
雪條掉到地上.
沒得吃.
於是大哭.
我媽媽沒有罵我妹妹.
為什麼弄掉了雪條.
指出妹妹自己害自己沒得吃.
這是活該.
自作孽.
與人無尤.
竟然吩咐我將雪條給妹妹吃兩口.
我怎麼需要為我妹妹弄掉雪條負責任.
對不對.
但我媽媽看來.
如果我吃雪條的時候.
妹妹沒得吃.
就是不對的.
我要和她分享.
本來是我的東西.
這樣就對了.
這樣就叫公平了.
各位弟兄姊妹.
各位觀眾.
你明不明白我媽媽的思路.
邏輯.
就是這樣.
弟兄姊妹.
我是一個簡單的人.
我很喜歡保羅的公平觀念.
我也覺得這是歷代教會.
今日教會所普遍奉行的慈惠制度.
我們鼓勵弟兄姊妹按照自己的能力.
將有餘的捐獻出來.
幫助有需要的人.
有錢出錢.
有力出力.
幫得多少就幫多少.
我們過去是不是一直這樣做.
我相信這樣的做法.
不可以改變整個社會的結構.

$^{641}$但總是或多或少能促進社會的公平.
是不是.
當然.
我好知道.
今日不斷有思想前衛.
行動激進的信徒告訴我們.
慈惠工作不等於社會關懷.
更加不等於社會公義.
教會只是做慈惠工作.
不追求社會公義.
就是在社會關懷問題上缺席.
區區的慈惠工作.
根本無助解決社會上的貧窮問題.
貧窮問題.
poverty(貧窮).
這個問題比個別的窮人poor(貧窮).
更加根本.
如果我們沒辦法消除.
社會上的貧窮問題.
沒辦法消除貧富不均這個經濟制度.
根本沒辦法.
根本上杜絕窮人的產生.
我們所做的慈惠工作是毫無價值的.
只是讓自己良心好過一點.
也減低社會不公平.
沒有改變.
減低它的殘酷性.
所以這個叫偽善.
這個是粉飾星評.
從某個角度說.
慈惠工作或者福利事業.
是鞏固了不公平的分配制度的幫兇.
至少讓窮人可以久言殘喘一點.
造成鴉片的作用.
團會起來一起奮起.
去推翻這個不公義的制度.
無論如何.
在錢惠,鳩濤的眼中.
慈惠,扶貧解決不了相對貧窮的問題.
公屋制度是改變不了堅尼系數的.

$^{681}$相對貧窮的問題才是社會不公平的根本.
教會如果要參與社會關懷.
派飯是沒有意義的.
唯一有意義的是去政府總部示威抗議.
要求全面改革香港的稅制.
打倒地產霸權,打倒官商勾結.
當然如果我們在示威之外.
能夠參與更激烈的行動就更好了.
如果能夠串連全球各地的機構.
一起來促進經濟公義.
這個社會運動就更好了.
是啊,你給錢在養一個柬埔寨的女童.
使她得到教育的機會.
能夠解決柬埔寨的貧窮問題嗎?.
能夠解決亞洲的貧窮問題嗎?.
肯定不可以.
基於all or nothing的思路.
不可以全盤解決問題的做法.
就等於沒用.
我養兒童是沒用的.
我再說,我是一個思想很簡單的人.
我又不是讀經濟的.
所以我完全不明白剛才說的那些深的理論.
我只是在感情上很憎恨別人批評說.
你這個做法是治標不治本.
你這個做法是頭痛醫頭,腳痛醫腳.
基於我對人的罪性的認定.
基於我對基督教末世觀的信念.
這個世界是會終極敗壞的.
人間不可以建成天國.
所以我不信有治本的做法.
我們能不能夠發明一個完美的教育制度.
醫療制度,政治選舉制度,經濟分配制度呢?.
即使我們推行更激烈,更徹底的革命運動.
將現存的一切都弄壞,推倒.
我們就能夠建立一個最理想的社會制度嗎?.
對不起,在我的信仰系統裡是沒有這個可能性的.
現實問題總是一計死一計命.
今天解決了問題的答案,明天可以成為問題.
今天的答案變成明天的問題.

$^{721}$所以只能夠沉著氣,逐一去應對.
今天我幫助一個人,他明天真的會死的.
是的,我知道他會死的.
所以所有救助都只有短暫的價值.
不過我認定短暫的價值還是有價值的,對不對?.
要說《屬靈八股》,我會很清楚地說.
除了帶領人信耶穌之外.
任何的救濟,改良或者革命都是治標,不治本的.
是不是?傳福音是治本唯一有永恆價值的東西.
同樣的,所有醫生都是頭痛醫頭,腳痛醫腳的.
所以醫生要分科,針對病徵,尋出病因.
盡力應對解決,是不是?.
那些宣稱能夠根本上解決人的健康問題的.
就我的經驗看都是神醫,神醫就是神棍.
毫無例外.
你們可能太年輕了.
我記得我年輕的時候.
不是聽過保濟苑的廣告是能醫百病的嗎?.
能醫百病是什麼藥呢?.
如果王若仙再生,可能我會相信他的.
他給我吃一顆藥丸吧.
所以聽姐妹,讓我們繼續高舉保羅君平或者公平的原則.
以有制無,樂於與人分享自己所有的.
亦樂於接受別人分享自己所沒有的.
前者,就是你分享給別人的時候不要驕傲.
接受別人的分享也不用自卑.
快樂施與,快樂接受.
盡力行善,從近處做起.
謹記耶穌的話,馬利福音25:40.
我實在告訴你們.
這些事你們既作在我這弟兄中一個最小的身上.
就是作在我身上了.
我相信我們會讓這個社會變得更加公平.
亦會更加有人情味.
不要懷疑自己所能給予的那兩個小錢的價值.
只要上帝喜悅我們這樣做就已經有最大的價值.
至於上帝會否將五餅二魚餵飽五千人.
會否施行以水變酒的神蹟.
那是上帝的事.
None of our business 與我們無關.

$^{761}$我們只管跟隨上帝.
我以前在你們中間經常說的老套說話.
一個lover,一個愛者是不可以太聰明的.
我們沒有太多的智慧.
所以我們被主呼召傻乎乎的跟隨他.
傻乎乎的愛上帝.
愛身邊的人.
就是這樣.
我已經用了所有時間去講八章十三到十五節.
所以十六到二十三節的內容只能夠草草收場.
這八節經文我跟大家很快講一下.
基本上傳遞三個訊息.
第一,保羅大大誇讚提多.
他讓保羅差派前來處理哥林多教會的問題.
保羅說提多的熱情和主動性不比他低.
十六節.
即使他好像是被保羅差派過來的.
其實是他自己想來的.
他很關心哥林多信徒的需要.
十七節.
第二點,保羅不僅派了提多來.
還有另外一個弟兄跟他同行.
保羅沒有講到那個弟兄的名字.
只是提到這個人在福音上得了宗教會的稱讚.
十八節.
宗教會也舉手贊成讓他同行處理運送捐款的問題.
十九節.
保羅這樣做是要保證運送捐款的過程不會出現差錯.
客觀的差錯都不是問題.
主觀的誤會保羅都不想.
就算沒有錯,但有人覺得錯.
保羅都不想.
保羅要讓辭位的行動保持光明磊落,清潔順傳.
就這一點,聖經給我們一個很重要的榜樣.
我們教會一直奉行不渝的.
繼續堅持.
教會凡是牽涉到錢.
一定不可以交給任何人單獨處理.
就算是廣被信任,得高望重的領袖都不可以.
另外,參與處理捐款的人.

$^{801}$必須本來就是得高望重的人.
他們不可以由領袖或任何人自行挑選.
是需要由信徒公開選舉.
不要給人思想受授的印象.
無論如何,教會的錢銀問題一定要建立互相監察的制度.
有清楚帳目,能夠跟眾人交代.
不要招來謠言誤會,破壞教會合一.
本來21節的說話.
是非常重要的一句說話.
我們留心行光明的事.
不但在主面前,就算在人面前.
也是這樣.
覺不覺得擲地有聲?.
這是我們教會的金科玉律.
錢銀問題一定要清清楚楚,乾乾淨淨.
第三點.
寫這封信的時候,保羅一方面催促.
哥林多信徒盡快完成他們的捐款承諾.
另外又增派一位弟兄過來.
這位無名者,同樣是熱心的信徒.
自設服侍哥林多教會的22節.
連同之前派來的保羅和某弟兄.
如今在哥林多教會的有三個人.
保羅對這三個人一起稱讚.
在23節,論到提多,他是我的同伴.
為你們勞碌的.
論到那兩位弟兄,他們是宗教會的使者.
是基督的榮耀.
除了提多是保羅親密的助手之外.
其他兩個弟兄.
保羅都稱他們是宗教會的使者.
指的是他們被保羅差派往返各地教會.
處理保羅所交付的任務.
他們是保羅的使者.
也是服侍宗教會的使者.
好,整段經文的訊息.
主要就是這三點.
其他資訊的內容我們就講到這裡.
我們下次會講第九章.
看金錢捐獻的第三篇.

\newpage



\section{馬可福音 10:17-22}
\label{sec:qJa4QRG5SV0}
\textbf{2021年3月21日崇拜直播｜陳冠成牧師｜感動還需敢動｜馬可福音10\_17-22}
\newline
\newline
連結: \href{https://youtube.com/watch?v=qJa4QRG5SV0}{\texttt{https://youtube.com/watch?v=qJa4QRG5SV0}} ~~~~ 語音日期: 2021-03-22
\newline
\newline
\hyperref[sec:zila_5Bj4Bo]{\small{< < < PREV SERMON < < <}}
~
\hyperref[sec:index_chronic]{\small{[返順時目]}}
~
\hyperref[sec:pK4j32Egyx0]{\small{> > > NEXT SERMON > > >}}
\newline
\newline
馬可福音 10:17-22
\newline
\begin{longtable}{cl}
\hline
\hline
章節 & 經文 (和合本修訂版)\\
\hline
10:17 & \begin{tabularx}{0.7\textwidth}{X} 耶穌剛上路的時候，有一個人跑來，跪在他面前，問他：「善良的老師，我該做甚麼事才能承受永生？」 \end{tabularx} \\ \\ \relax
10:18 & \begin{tabularx}{0.7\textwidth}{X} 耶穌對他說：「你為甚麼稱我是善良的？除了神一位之外，再沒有善良的。 \end{tabularx} \\ \\ \relax
10:19 & \begin{tabularx}{0.7\textwidth}{X} 誡命你是知道的：『不可殺人；不可姦淫；不可偷盜；不可作假見證；不可虧負人；當孝敬父母。』」 \end{tabularx} \\ \\ \relax
10:20 & \begin{tabularx}{0.7\textwidth}{X} 他對耶穌說：「老師，這一切我從小都遵守了。」 \end{tabularx} \\ \\ \relax
10:21 & \begin{tabularx}{0.7\textwidth}{X} 耶穌看著他，就愛他，對他說：「你還缺少一件：去變賣你所有的，分給窮人，就必有財寶在天上；然後來跟從我。」 \end{tabularx} \\ \\ \relax
10:22 & \begin{tabularx}{0.7\textwidth}{X} 他聽見這話，臉就變了色，憂憂愁愁地走了，因為他的產業很多。 \end{tabularx} \\ \\
[1ex]
\hline
\hline
\end{longtable}
$^{1}$接著我們繼續以聆聽聖道去敬拜我們的主.
今日的經訓是在《馬可福音》10章17至22節.
耶穌出來行路的時候.
有一個人跑來.
跪在他面前.
問他說.
良善的夫子,我當作什麼事才可以承受永生?.
耶穌對他說.
你為什麼稱我是良善的?.
除了神一位之外,再沒有良善的.
誡命你是曉得的.
不可殺人.
不可姦淫.
不可偷盜.
不可作假見證.
不可窺婦人.
當孝敬父母.
他對耶穌說.
夫子,這一切我從小都遵守了.
耶穌看著他.
就愛他.
對他說.
你還缺少一件.
去變賣你所有的分級窮人.
就必有財寶在天上.
你還要來跟從我?.
他聽見這話.
念想就變了色.
悠悠愁愁的走了.
因為他的產業很多.
今天,陳冠成牧師會為我們宣講上帝的訊息.
題目是「感動還需感動?」.
弟兄姊妹平安.
還有差不多兩個星期就到復活節.
所以今天很想和大家看一段關於耶穌基督.
他在踏上耶路撒冷受苦之旅其中一個經歷.
其實這段經文我們都很熟悉.
不單止馬可科音.
馬太科音和路加科音都有記載這段經歷.
經文裡說到有一個人來找耶穌.

$^{41}$馬太科音補充了一些資料.
說他其實是一個年輕人.
路加科音說他是做官的.
而馬可科音說故事的結尾.
交代了他有很多產業.
是一個富有的人.
他擁有我們說令人羨慕的生活條件.
不過他覺得自己欠缺一件事.
沒有永生的把握.
於是他專程來向耶穌求教.
我們看到整個圖畫都很特別.
他很主動來找耶穌.
但來到最後.
他又很憂愁地離開了耶穌.
當你看到這個情景.
讓你想起什麼呢?.
話說在一個聚會裡.
有一群社會賢達.
有一些學者.
甚至有一些教會的牧者.
他們聚首一堂.
雖然他們的職業都不同.
但他們都有一個共通點.
就是他們經常會去到不同的地方.
來跟別人分享.
如何讓自己的人生更加精彩.
其中有一個人在當中半開玩笑地說.
我們每一年在台上都對著不少人.
甚至乎有千上萬的人來說話.
但大家覺得.
有多少人聽完我們說話之後.
真的願意跟著去做呢?.
大家就靜了一靜.
之後陸續有人開口.
有人說我猜有兩成左右.
有人說我猜只有一成.
也有人說我猜百分之一左右.
事實上他們都很認同.
就是很多人聽到他們的分享之後.
真的深受感動.

$^{81}$但亦都慨嘆可能只有少數的人.
生命有所改變.
我想正如今日.
無數的人都深被耶穌的教導.
被祂的言行為之感動.
但又有多少人願意遵行耶穌的吩咐呢?.
當然我們沒有統計過有多少人.
但我盼望.
我們都是其中的一個.
就如今日這個講題所說.
我們不單對耶穌的吩咐.
對祂的教導.
不只是感動.
更加勇敢地用行動來實踐.
以行動來回應所聽到的道.
所以為此我們一起先有個祈禱.
我們交託給上帝幫助我們.
聖靈啊,願你與我們同在.
讓我們真是柔和謙卑.
讓我們有悟性能夠明白.
聖經的真理.
耶穌你對我們的教導和吩咐.
也讓我們不單單只心裡有感動.
更加讓我們願意用行動來回應你.
切切實實地跟隨你.
為此我們同心獻上我們的祈禱.
靠上帝幫助我們.
奉基督耶穌.
你的冥之來祈求.
阿們.
看回這段經文.
少年才主不單主動跑來找耶穌.
是嗎?.
經文很仔細地形容.
他更加不介意放下他的身份和面子.
在眾目睽睽之下.
向耶穌來到下跪.
你看到他真的很謙虛.
還有追求永生的誠意.
看得出.

$^{121}$更加看到他真的很介意自己有沒有永生.
或許他真的想將現在所擁有的一切幸福.
延續下去.
甚至到永遠.
我們不知道他是刻意抬舉耶穌.
還是他真的很仰慕耶穌.
他用一個猶太人用來稱呼上帝的稱號.
「良善」.
來稱呼眼前這位拿撒勒人的夫子.
我們看到耶穌不接受這個稱呼.
耶穌又說.
為什麼稱我為良善呢?.
除了神以外沒有良善的.
我們今天當然知道耶穌是良善的.
但不是因為他身為人.
在地上做了很多好人好事.
行了很多神蹟奇事.
不是因為他的好行為.
所以被稱為良善.
而是因為我們今天知道他是上帝的兒子.
和父神同等.
所以是良善的.
但明顯少年才主.
他未知道耶穌的真正身份是誰.
就用上帝的尊稱.
「良善」來稱呼一個猶太人.
耶穌不能認同這一點.
其實從少年人這方面的說話.
我們不難看到他整個思路是怎樣.
首先.
大概他不會去求教一個被自己稱為「渣」的人.
對不對?.
他一定是認為在他心目中.
耶穌是比他優越很多的一個猶太老師.
否則不需要去求教他.
他亦相信耶穌的好行為.
耶穌的善行一定比自己更加好.
更符合他心目中的良善標準.
所以才這樣稱呼耶穌.
而他之所以問.

$^{161}$我當作甚麼才可以承受永生.
是因為他認為.
我都可以透過善行.
達到上帝的標準.
從而換取到永生.
不過我只是未找到竅門.
我只是未知道要做對甚麼善行.
所以我唯有找專家.
向耶穌尋求一個更專業的意見.
大概這也是少年才主當時的想法.
所以,丁姐妹.
假如今天同樣有人問你.
我應該做甚麼事情.
做甚麼善行.
才可以承受永生.
相信我們今天不難回答.
我們今天可以用已忽所書.
代表耶穌回答他.
大家都懂得背吧?.
你們得救是本乎因.
也因著信.
並不是出於自己.
乃是神所賜的.
也不是出於行為.
免得有人自誇.
我們一定會和對方強調.
人得救只能憑信心.
只能信靠上帝藉著耶穌基督.
在十字架上已經成就的救恩.
這樣才能進入天國.
得享永生.
人無法透過自己的好行為.
透過善行.
來作為交換的條件.
我們都很明白.
今天我們信仰的重要教導.
但你看到耶穌沒有這樣回答少年才主.
既然少年才主很在意做甚麼.
耶穌就順著他的思路.
來跟他說做甚麼.

$^{201}$要做甚麼來入手.
你看到耶穌在問少年才主.
猶太人所熟悉的十誡.
特別是關於人與人之間關係的誡命.
那些律法.
測試少年才主和上帝的關係是怎樣.
耶穌問.
誡命你曉得的.
不可殺人.
不可姦淫.
不可偷渡.
不可作假見證.
不可窺負人.
當孝敬父母.
我們可能會奇怪.
為何耶穌不問頭四誡.
關於愛上帝的吩咐.
大概可能耶穌也認為.
人是否真的愛上帝.
對上帝忠心.
往往能夠從他.
有沒有實踐上帝給他的愛人的吩咐.
這樣顯示出來.
其實這也是新舊約一致的教導.
就連耶穌自己也一樣說.
這些事如果做在最少的弟兄上.
又如何?.
就是做在耶穌身上.
你看到少年才主的回答是.
老師拉比.
這一切上帝的律法吩咐.
我至少就做了.
我至少就遵行.
當然潛台詞是.
但我還沒有永生的把握.
我是否有甚麼善行做漏了.
還是有甚麼我未知.
他很心急想了解答案.
你從耶穌的反應.
就是耶穌看見他.

$^{241}$就愛他.
大概少年才主.
也不是在自言自語.
他真的有真材實料.
他確實從小就遵行上帝的律法.
是一個敬虔的猶太人.
事實上他通過了第一關.
耶穌對他第一重的考問.
他不單止合格.
還怎樣?.
滿分.
所以你看到耶穌很愛他.
認為這個人是一個可造之材.
於是怎樣?.
如果你還記得.
可能很多年前.
我們會考之後.
就要高考A級.
會考合格.
還要科4A.
上帝就直接考A級.
高考.
耶穌直接給他一個.
更高級別的考驗.
耶穌就說.
你還缺少一件.
近了.
差一步.
到底是甚麼?.
耶穌.
你快點告訴我.
讓我完成它.
我就可以承受永生.
耶穌就說.
你去變賣你所有的分給窮人.
就必有財寶在天上.
你還要來跟從我.
等一等.
耶穌.
你又說是一件事.

$^{281}$又說差一件事.
為甚麼我數來數去.
有多少件事?.
不對數.
不止一件.
你留意其實原本.
經文的原文.
其實不止一個命令.
一個吩咐.
裡面包括去.
去賣.
變賣.
分給.
來.
跟從.
原來有幾個命令.
為甚麼這麼多命令.
其實耶穌都說.
只有一件事.
好像有點屈機.
確實.
我們其實不難理解.
變賣所有分給窮人.
是否極大的善行?.
確實是.
我們今天都不容易做到.
但亦不能因為這麼大的善行.
人就可以進到永生.
交換到永生.
我們又很明白.
只不過變賣所有分給窮人.
有甚麼效用?.
你跟我都不難理解.
兩個人掉進水裡.
你先救誰?.
先救誰?.
耶穌問一條二選一的問題.
耶穌要求他變賣所有分給窮人.
其實是迫使這個少年財主.
很嚴肅地面對自己.

$^{321}$如果你只能夠二選一.
你選愛上帝.
還是愛你的產業?.
你選吧.
不可以無伶兩可.
不可以左右逢源.
不可以又渴望永生.
又想保留你在地上的財富.
不放手直到永遠.
所以.
耶穌向少年財主發出這麼多命令.
確實又只向一見一.
用舊約的說法是甚麼?.
就是第一條的誡命.
除了我之外.
你不可有別的神.
用新約的說法是甚麼?.
你跟我都很明白.
耶穌經常強調.
撇下所有專一的跟隨主.
所以無論這個少年人.
他表示自己多麼渴望親近上帝.
多麼渴望追求永生.
我們經常說.
甚麼是最誠實?.
身體真是最誠實.
在這條二選一題裡.
迫使這個少年人知道.
他自己當下真正最關心的是甚麼.
其實是甚麼?.
所以你看到耶穌.
這條題目是度身訂做給他.
耶穌度身訂做這條變賣所有.
分級窮人的考驗.
給這位少年財主.
原因是馬可交代了.
他有許多的產業.
或許是他人生.
或許是他當下.
他覺得最重要.

$^{361}$他覺得最需要倚靠的事情.
最難放手的事情.
但對於其他人.
耶穌或上帝要他撇下的事情.
可能不一樣.
不一定是錢財.
不一定是產業.
而是其他事情.
譬如說彼得.
使徒彼得.
耶穌要他撇下的是甚麼?.
是他的工作.
是他賴以為生的職業.
以至他能夠專一.
跟隨主.
由一個捕魚的漁夫.
變成一個得人的漁夫.
這是他要他撇下的.
或許對於舊約.
亞伯拉罕.
他夠有錢.
他有許多的產業.
上帝有否吩咐他撇下所有跟隨神?.
我們看見沒有.
但要他撇下甚麼?.
那件事很重要.
將你的命根.
將你的心肝寶貝以撒.
給上帝.
來顯明到底.
你是愛上帝.
還是愛上帝給你的好處.
二選一.
或者耶穌也有些事情要撇下.
父上帝要他撇下甚麼?.
撇下你的尊榮.
身為神兒子的尊榮.
撇下你的生命.
來顯明.
你體重父神的意思.

$^{401}$過於自己的意思.
所以其實耶穌.
稍後你會看見他通過二選一的考驗.
就是他捨棄他的生命.
讓我們得到生命.
事實上撇下所有的意思.
他的核心意義不是說.
你扔下世上所有一切.
斷掉你的六親.
然後遠離塵世.
個個全職入神學院事奉.
然後全時間服侍主.
這可能是其中一條路.
但不是必然.
撇下所有的焦點.
其實是說除了上帝之外.
任何的事情.
可能是金錢.
可能是你的權利.
可能是你的名譽.
可能是你的安屍生活.
或者某段關係.
任何這些事情.
如果上帝開口.
你願不願意放手.
這才是撇下所有的核心意義.
如果上帝開口.
要你放下某件事.
你願不願意放下.
但並不是那件事.
本身沒有價值,沒有意義.
而是相比起耶穌.
相比起永生.
這一切都變得次要.
所以弟兄姊妹.
其實今天耶穌這條題.
說祂已經在考問你.
又或者某年某月某日.
終會臨到你的面前.
因為人生總有些際遇情況.

$^{441}$你只能夠二選一.
對不對?.
要麼選愛上帝.
還是選愛自己.
正如耶穌迫少年人看清楚.
永生是什麼一回事.
要麼你100\%全人進入永生.
要麼就沒有永生.
不會你的頭.
你的手進入永生.
你的腳留下.
不能講價.
不能中間落密.
所以假如有一天.
耶穌要我們撇下.
我們生命裡一樣.
最不捨得.
最重要很難放手的事情.
以至我們能夠更加專心地.
愛祂,跟隨祂.
我們會否都像少年財主一樣.
真的很苦惱.
真的很掙扎.
所以.
平時的操練就很重要.
不會一考試的時候.
就抱頭痛哭.
平日的操練相當重要.
免得當我們真的要做二選一.
選愛耶穌或選愛自己的時候.
不會卡住了.
又或者不會像這個故事結尾一樣.
悠悠愁愁地離開了耶穌.
有兩方面.
首先.
不要限定了上帝祝福你的方式.
不要限定了上帝要怎樣賜福你.
這其實就是少年財主.
他最糾結的信仰的位置.
你想想.

$^{481}$如果我們想清楚少年財主.
他的掙扎可能是.
耶穌你要我多守幾條律法.
這樣來贊助他永生.
有沒有問題?.
沒有問題.
多做一點沒有問題.
反正他也認定上帝的福氣是怎樣.
靠自己多勞多得.
正如我們也會這樣.
為上帝多做一點事奉沒有問題.
但為上帝放棄一些東西.
就萬萬不能.
他限制了上帝對他祝福的方式.
又或者說.
耶穌你要我多收窮人.
來換取永生.
我也接受.
你要我捐七成八成身家.
我覺得也值得.
如果因為這樣得到永生.
起碼還有什麼?.
還有兩三成.
還有你不要忘記.
他做聖行的.
他做官的.
他年輕嗎?.
年輕.
即是說什麼?.
就算我捐了七八成甚至九成.
我還有很多時間.
青春可以找回來.
但現在耶穌要他怎樣?.
不只是要他周濟窮人.
而是要他變做窮人.
變做窮人不止要怎樣?.
還要跟隨耶穌.
跟隨耶穌對財主來說.
意味著什麼?.
我不單止這一刻變成窮人.

$^{521}$我還要窮到永遠.
我有生之年我都要窮到永遠.
我還要窮到永遠.
很奇怪.
我們兄弟姐妹.
同一個場景.
兩個人看到不同的東西.
耶穌看到的是.
這樣你在天上將會變得富有.
四年才主看到的畫面.
原來是另一幅.
是甚麼?.
這樣我在地上立即變成窮人.
為何會這樣?.
是因為他限定了.
上帝祝福給他的方式.
在他的信仰生命裡太多限定詞.
譬如說.
我做得越多越好.
上帝就會越祝福我.
這是他的限定詞.
又譬如說.
上帝的祝福只能是一條加數.
只可以再給我更多好處.
上帝的祝福不能是一條減數.
不能拿走我已經有的好處.
這是他的限定詞.
弟兄姊妹.
我們會否在信仰裡.
有很多這些限定詞.
限定了上帝給我們的祝福呢?.
譬如說.
我們會否不自覺地.
在我們的禱告生活裡.
我們限定了上帝.
上帝我希望你何時何地.
用甚麼方法來幫我解決.
我現在眼前的困難.
最好這樣就好了.
最好不要那麼辛苦.

$^{561}$我們有時會很狹窄地理解上帝.
我們只是期望上帝.
用我們所想,所認知的方法.
來幫助,帶領我們.
這樣做可能是人之常情.
但這樣做有一個盲點.
就是我們慢慢會忘記.
上帝可以用不同的方法.
上帝可以用你和我.
都想像不到的方式.
甚至乎可能是一些.
你不想碰到的際遇.
來建立我們的生命.
上帝可能在某年某月某日.
在你人生裡.
不跟你解釋.
沒有徵詢過你的同意.
那又怎樣?.
他就拿走你的東西.
又或者他容許一些禍患疾病.
突然間臨到你的身上.
會令到你很困惑,很辛苦.
甚至乎可能會埋怨上帝.
你為何這樣作弄我.
但你細心想.
保羅說過甚麼?.
神是那位叫萬事都互相效力的上帝.
叫那些愛上帝的人.
總是能夠得到益處.
這些沒有人想遇到的遭遇.
背後其實會不會都是上帝的祝福?.
會不會都有上帝的美意.
只不過化了妝.
你還未看得清.
所以弟兄姊妹.
真的.
假如我們真的想經歷.
上帝各樣的福氣.
這是我們心所想的,是嗎?.
我們就不要難了上帝各樣祝福你的方式.

$^{601}$可能不是你喜歡.
但上帝說對你來說剛剛好.
有甚麼操練可以幫我們呢?.
譬如說.
我們人生總有不同的習慣.
我們總有一些行之有效.
做人處事的方法.
我們一直都在用.
會不會試試改變一下.
增加自己的彈性?.
會不會用一些新的方法.
去面對相同的事情.
來拉闊自己的彈性.
從而拉闊你對上帝的看法?.
譬如說你每天上班.
都是走那條路.
試試走第二條路.
你會體會到其實都去到目的地.
時間不是差很多.
但你會看多了很多東西.
試過甚麼?.
又或者可能你每天早餐都是甚麼?.
早餐A,麻煩你.
試試早餐BCD.
試試自己煮.
增加自己的彈性.
拉闊自己的目光.
你看到可能上帝的祝福都在裡面.
又或者.
回到公司.
你習慣第一件事是做甚麼?.
先查一下今天有甚麼行程.
有甚麼要做.
試試先祈禱.
看看上帝的臉色.
凡事交託給上帝.
然後才開工.
整件事完全不同.
輕省很多.
又或者用新的模式.

$^{641}$來跟別人互動.
譬如說.
先不要太快開口.
先說一套.
不要太快回覆對方的電郵.
或者WhatsApp.
特別是那些令你很混亂的.
等一等.
跟上帝聊天.
上帝為你預備的時機.
一定比你自己所想好.
如果你相信他是上帝的話.
特別我知道在疫情下.
我們有很多不習慣的事.
我們不習慣Zoom.
對不對?.
我們不習慣用視像的方式來開小組.
但今天這段經文提醒我們.
不要限制了上帝賜福給你的方式.
先試試.
試試用Zoom.
試試Zoom裡有不同的周會模式.
看看上帝的祝福在哪裡.
又或者在疫情下.
你發覺原本有些弟兄姊妹你認識.
他們沒有加入小組.
對不對?.
但現在因為疫情的緣故.
連平時每個主日.
崇拜見面的機會都沒有了.
你試試用不同的方式.
正好你用Zoom的方式.
重新聯繫他們.
邀請他們入組.
關心他們.
我相信他們都很渴望見到大家.
很希望在疫情下.
仍然有相交的機會.
即使是視像進行.
弟兄姊妹,真的.

$^{681}$不要限制了上帝賜福給我們的方式.
不妨操練一下.
用新的形式來體驗上帝的福氣.
你會體驗到.
有意想不到的收穫和驚喜.
另一方面.
耶穌其實提醒我們.
要有小孩子的信心.
或許.
我們在信仰上真的會有疑惑.
有些不明白的地方是真的.
但同時耶穌亦要求我們.
要認定.
上帝不會騙你的.
為何這樣說?.
雖然我們今天沒有機會詳細看.
但如果你看回上文.
馬可在少年才主的故事中.
上文是交代.
有班人帶小孩子來.
希望耶穌祝福,按手.
門徒覺得他們都搞定了.
但耶穌是斥責他們.
不可以這樣.
因為耶穌對門徒說.
「我實在告訴你們.
凡要承受神國的.
若不像小孩子.
斷不能進去」.
正好是回應下文少年才主.
他很想承受永生的故事.
所以.
從另一個角度看.
其實耶穌對少年才主的吩咐.
是否他能做到?.
其實是做得到.
關鍵只是是否能做到.
不是他能否做到.
而是能否做到.
是他是否百分百相信.

$^{721}$耶穌所說的是真理.
剛才我們提及亞伯拉罕.
提及亞伯拉罕獻耳殺的事情.
你試試代入他的角色.
是否掙扎,糾結?.
其實很難做決定.
他等了多少年才有耳殺?.
聖經記載.
他七十五歲.
上帝就開始跟他說話.
將要成為大國.
很多的英雄.
到他多少歲耳殺才出世?.
一百歲.
等了二十五年.
到他現在養大了耳殺.
可能都是年輕人.
背得起憲法制的柴.
等了那麼久.
上帝才將耳殺賜給我.
現在我養到那麼大.
你又說甚麼?.
你要拿回?.
很難的.
正路可能我們都會.
跟身邊的人商量.
跟太太商量.
跟撒拉商量.
一起祈禱.
一起求問.
上帝是否真的想我們這樣.
看看大家是否感動.
然後才做決定.
對不對?.
一定會.
但你猜.
如果撒拉知道這件事.
會怎樣?.
如果你.
特別是姐妹們.

$^{761}$你是作為母親.
你有甚麼反應?.
知道亞伯拉罕見耳殺.
你是撒拉的話.
你會不會說.
好啊.
那你帶兒子去吧.
早點回來.
我煮飯等你.
會不會?.
不會的.
讓撒拉知道.
亞伯拉罕見耳殺.
一定怎樣?.
我不會跟你拼命.
肯定跟你拼命.
所以雖然.
創世記沒有交代.
不過我想大概.
亞伯拉罕都是瞞著撒拉.
去見耳殺.
不容易的決定.
但是這麼多年.
亞伯拉罕和上帝的互動.
提醒他一件事.
就是這麼久以來.
上帝有沒有規附過我.
他可能有規附上帝.
但是上帝是信實的.
從來沒有規附過亞伯拉罕.
這次亦不會例外.
他相信.
所以即使上帝現在的吩咐.
是多麼難做決定.
多麼不合理.
多麼瘋狂.
他都願意.
多難都好.
他仍然堅持相信一件事.
就是上帝命定要賜福給我.

$^{801}$不是要為難我.
不是要作弄我.
所以.
他寧願在這一刻.
在這個二選一題裡.
他寧願選上帝.
選上帝的吩咐.
而不選.
我死都要保留上帝賜給我的福氣.
他是這樣選的.
所以我們看到.
亞伯拉罕不是單單信心之父.
這麼簡單.
你知道他那時超過了一百歲.
一個這麼年長的人.
他仍然可以保持.
對上帝保持著一個孩子的心.
仍然沒有規限上帝賜福給他的可能性.
他相信上帝不會騙我.
我記得有一個.
很深刻我和我兒子的互動.
幾年前.
我不記得他幾歲.
應該是讀一二年班.
大概.
有一次他放暑假.
我帶他去太空館.
打算帶他去看天象廳.
可能是關於太空科幻的戲.
你知道未開場.
到處去逛逛.
你記得太空館有個精品廊買手信.
裡面有各式各樣的玩具.
包括顯微鏡.
他拿著顯微鏡.
他知道是顯微鏡.
他很想擁有.
因為沒有這個玩具.
他亦知道顯微鏡用來做什麼.
平時我一定會這樣回應他.

$^{841}$很自然地說.
你有很多玩具.
不要買了.
放下.
然後轉頭就走了.
為什麼這麼深刻.
因為那次我不是這樣回應他.
大概可能因為.
我以前的工作.
用很多精密的顯微鏡.
甚至電子顯微鏡.
我都會用.
所以我一看他拿上手的顯微鏡.
我就知道是什麼.
他是不是玩具.
其實我不肯定.
可能真的有顯微鏡的功能.
但肯定是假的.
於是我就跟他說.
你這個顯微鏡的放大功能很差.
你不是想浪費錢.
買一個很差的顯微鏡吧.
他聽到之後.
馬上說不想.
然後從他的面部表情.
我知道他開始嫌棄.
他手上那個被我批評很差的.
顯微鏡.
一點掙扎都沒有.
就放下.
跟我一起走.
在這件事我體會到.
什麼叫做小孩子的心.
就是當他發現原來.
這個世界上有更好的東西.
擺在他眼前的時候.
等他的時候.
他不會死撇手上那個.
次好的東西不放.
這個是小朋友的心.

$^{881}$第二件事就是.
小朋友會完全相信爸爸所說.
他不會質疑爸爸.
有一把你是不是點我.
你是不是說真的.
你真的懂.
他一句都沒有問過.
就很簡單.
聽話.
照做.
我從來沒有說你放手.
但他自己明白.
他放手.
所以弟兄姊妹.
無論你的年歲是多少.
無論你信了主多久.
盼望我們都常常對天父.
耶穌常傳那顆孩子的心.
你要認定.
耶穌和祂的吩咐.
是你的首選.
而不是你的後備方案.
你的backup plan.
什麼叫做backup plan.
後備方案.
就是你先試自己的方法.
衝完一輪.
撞完一輪你發覺不行.
又怎樣.
就求上帝.
求上帝.
又或者什麼為之.
當上帝是backup.
就是你問上帝.
上帝給了你的心意.
你覺得我不喜歡這樣.
所以ignore了他.
delete了他.
丟下了他.
當上帝是後備方案.

$^{921}$真的弟兄姊妹.
其實你想想.
如果我們凡事真的學習.
讓上帝居首位.
其實人生很多問題.
可以輕省面對.
因為你知道上帝會幫你.
上帝會同行.
很多事情其實可以迎刃而解.
你心裡的那些勞苦受煩.
可以一掃而空.
如果你相信上帝不會騙你的.
所以當我們遇到困難.
甚至考驗臨到.
又或者像少年才主那樣.
上帝逼你二選一的時候.
你不要這麼快.
就想想自己有什麼吃虧了.
不要太快前途將會受虧損.
假如我遵行上帝吩咐的話.
你嘗試改換你的思想模式.
當困難考驗.
或者上帝給你的挑戰.
臨到的時候.
第一件事不要想避開.
不要想放棄.
你想想上帝.
相反上帝不會無緣無故這樣做.
一定是有原因.
你不明白的意思.
提醒自己停下來.
真的要慢慢去想清楚.
並且提醒自己.
上帝由始至終是賜福給你的.
祂不是要作弄你.
不是要為難你.
可能有些事情你會覺得不舒服.
但你要相信一句說話.
就是上帝不會讓我們經歷.
我們受不了的試煉或者試探.

$^{961}$所以相反地想.
當這些試煉試探臨到的時候.
其實是說上帝看得起你.
上帝認為你像少年才主.
你已經通過了某個階段的考驗.
你不差,很好.
上帝愛你.
所以給你更高級別的挑戰.
為的是要你更加走近祂.
如果當我們懂得反過來想的時候.
其實我們就成長了.
我們就更加願意親近上帝.
而不是像少年才主那一刻.
悠悠愁愁地離開.
所以最後,弟兄姊妹.
回到我們今天的講題.
我們對耶穌今天的教導.
不單要有感動.
心裡不單要有感動.
更加要勇敢地去行動.
你要夠膽走出那一步.
我記得以前認識一位.
弟兄.
其實他的上年生命很好.
他不單在教會裡面.
有很好的生活.
有信仰生活,有事奉.
他還主動去神學院.
接受夜間的信徒訓練.
有一天他約了.
他就說他很想更加走一步.
全時間去夢照,事奉上帝.
但他卡住一些東西.
因為他說上帝沒有很明顯.
很清晰地告訴我.
Yes or No.
他糾結的是.
假如我真的選錯.
那我現在那份退休金.
沒有了一部分.

$^{1001}$很難抉擇我.
怎樣幫他去看這件事.
我沒有特別跟他分析.
我認識他,他是一個很好的基督徒.
我也不知道上帝是否真的呼召他.
去全職事奉他.
我知道很簡單.
先遞表.
你先走一步.
你要相信上帝不會.
假裝你做錯了一個決定.
上帝就降和給你.
或者不祝福你.
即使你所謂人生的選錯.
上帝都有辦法撥亂反正.
賜福你.
但你要先遞表.
你不遞表,上帝怎會告訴你是否真的.
你不遞表,然後回到.
徵詢你牧者的.
問他是否同感一靈.
你怎知道上帝的回應是怎樣.
上帝一定有辦法.
難阻你的,如果不是他的心意.
神學院不會收你.
但你要走出那一步.
同樣地,兄姐妹.
今天我們很多時候.
都會被人提醒我們.
當你收到懷疑短訊的時候.
訊息的時候,就要怎樣.
就要怎樣.
核實,就要Fact Check.
確認了是真的.
才去回應,才去做決定.
這個確實是需要的.
但對於追求真理.
我們不能單靠百分之一百.
核實,Fact Check.
有些事上帝是隱密.

$^{1041}$不會告訴你的.
有些事是你要先有行動.
走出那一步.
運用信心來跨過.
我們的信仰.
不可能凡事都百分之一百.
核實了,才去做.
所以弟兄姊妹.
都要經文提醒我們.
很多事情,特別是信仰的事情.
你只有先踏出一步.
你只有先勇敢,有行動.
上帝才會繼續告訴你.
祂要怎樣帶領你.
那幅圖畫才會越來越清晰.
所以,在這裡.
很想去呼籲大家.
確實,有時最難就是走那第一步.
走出去,反而輕散.
當你勇敢起動的時候.
上帝就會告訴你.
下一步要怎樣做.
很想呼籲大家.
假如你很有感動.
很想在這個時段.
或者在人生某個範疇.
你真的需要為上帝作出一些改變.
上帝給了你感動.
我很想為你祈禱.
求上帝給你勇氣.
給你決心,攬出去.
就好像彼得踏出船身那一下.
走出去.
這是第一個呼籲.
又或者假如.
在螢光幕的你.
你未認識耶穌.
你很想進一步了解.
我們也邀請你憑信心走出那一步.
你可以打電話給王浸.

$^{1081}$來了解更多信仰.
來找我們.
來認識這位.
真的很愛你的耶穌.
來認識.
祂愛你愛到一個地步.
祂願意為你犧牲,捨命.
祂願意拯救你.
讓你真的經歷豐盛的永生.
又或者假如你已經是基督徒.
我鼓勵你.
不要只滿足於得到救恩.
得到天堂的入場券.
上帝更加期望我們.
專心一意跟隨祂.
順從祂的吩咐.
實踐科林的使命.
這才是耶穌對我們最大的期望.
耶穌想我們得到的.
不只是祂的好處.
是得到祂自己.
所以很想為這三個的呼籲.
我們一起同心禱告.
請我們低頭.
我們一起去祈禱.
主耶穌說我們.
多謝你的拯救.
你告訴我們.
你多麼愛我們.
愛到一個地步.
你願意犧牲捨命.
你願意藉著十字架上.
那個捨身流血.
來顯示上帝對我們的愛.
也顯明你對上帝的愛和忠誠.
你願意撇下.
以致我們能夠得著.
你會讓我們效法.
不單我們心裡有那份感動跟隨你.
我們真的用行動來實踐.

$^{1121}$特別當我們在生命裡.
有某些範疇.
我們真的很想突破成長.
我們很想為上帝而更新.
求主不單感動我們.
求你幫我們.
求你用你不同的方式.
來幫我們走出這一步.
以致我們更加走近你.
更加經歷到上帝.
你奇妙的帶領和福氣.
讓我們知道.
無論人生是什麼境況.
無論什麼環境.
不能夠難阻到.
上帝你給我們的愛和賜福.
因為上帝你定義.
要用我們去成就你的工作.
你要在當中讓我們成長.
更加明白你.
更加體會什麼是豐盛的生命.
事實上這樣的事情.
不是要等到我們死後在天間.
而是在地上我們就已經可以經歷到.
上帝我又願你去幫助.
每一位未認識你的.
願你賜福這些新朋友.
讓他們真的可以憑信心.
行前一步.
讓他們願意在當中.
放開懷抱.
放開一些擔心.
一些懷疑或.
讓他們相信.
上帝是很愛他們的.
愛到一個地步.
你親自來找他們.
並且不止一次.
而是多次.
因為上帝你定義要得著他們.

$^{1161}$也期望他們真的能夠把握機會.
得著耶穌.
來做他們的救主.
做他們人生的主宰.
為此我求上帝幫助我們每一位.
讓我們真的能夠成為你的兒女.
上帝我又願你幫助每一個基督徒.
讓我們不滿足於現狀.
因為上帝你要我們成長.
就好像耶穌你希望.
這個少年才主行多一步一樣.
讓我們真的能夠更加專注跟隨你.
讓我們真的能夠為上帝你的緣故.
我們學習如何放下次要的事情.
放下那些難阻我們.
更加走近上帝你的事情.
上帝我們知道不容易.
但這確實是你對我們的期望和心意.
並且你一定會幫我們成就.
因為你最終希望我們得著的.
不是純粹上帝你給我們的好處.
給我們的福氣.
上帝你希望我們得著的是耶穌祂本人.
是主耶穌祂自己.
所以主啊.
幫助我們讓我們藉著今天的訊息.
我們真的聽到,行到.
有感動亦都敢於去行動.
回應你的愛,回應你的呼籲.
我們這樣的禱告,交託.
是奉基督耶穌你得勝的名字來去呼求.
阿們.
\newpage



\section{馬可福音 12:18-27}
\label{sec:pK4j32Egyx0}
\textbf{2021年3月28日崇拜直播｜陳啟興牧師｜復活的大能｜馬可福音12\_18-27}
\newline
\newline
連結: \href{https://youtube.com/watch?v=pK4j32Egyx0}{\texttt{https://youtube.com/watch?v=pK4j32Egyx0}} ~~~~ 語音日期: 2021-04-01
\newline
\newline
\hyperref[sec:qJa4QRG5SV0]{\small{< < < PREV SERMON < < <}}
~
\hyperref[sec:index_chronic]{\small{[返順時目]}}
~
\hyperref[sec:NcnNHWwVDiw]{\small{> > > NEXT SERMON > > >}}
\newline
\newline
馬可福音 12:18-27
\newline
\begin{longtable}{cl}
\hline
\hline
章節 & 經文 (和合本修訂版)\\
\hline
12:18 & \begin{tabularx}{0.7\textwidth}{X} 撒都該人來見耶穌。他們說沒有復活這回事，於是問耶穌： \end{tabularx} \\ \\ \relax
12:19 & \begin{tabularx}{0.7\textwidth}{X} 「老師，摩西為我們寫下這話：『某人的哥哥若死了，撇下妻子，沒有孩子，他該娶哥哥的妻子，為哥哥生子立後。』 \end{tabularx} \\ \\ \relax
12:20 & \begin{tabularx}{0.7\textwidth}{X} 那麼，有兄弟七人，第一個娶了妻，死了，沒有留下孩子。 \end{tabularx} \\ \\ \relax
12:21 & \begin{tabularx}{0.7\textwidth}{X} 第二個娶了她，也死了，沒有留下孩子。第三個也是這樣。 \end{tabularx} \\ \\ \relax
12:22 & \begin{tabularx}{0.7\textwidth}{X} 那七個人都沒有留下孩子。最後，那婦人也死了。 \end{tabularx} \\ \\ \relax
12:23 & \begin{tabularx}{0.7\textwidth}{X} 在復活的時候，她是哪一個的妻子呢？因為他們七個人都娶過她。」 \end{tabularx} \\ \\ \relax
12:24 & \begin{tabularx}{0.7\textwidth}{X} 耶穌說：「你們錯了，不正是因為不明白聖經，也不知道神的大能嗎？ \end{tabularx} \\ \\ \relax
12:25 & \begin{tabularx}{0.7\textwidth}{X} 當人從死人中復活後，也不娶也不嫁，而是像天上的天使一樣。 \end{tabularx} \\ \\ \relax
12:26 & \begin{tabularx}{0.7\textwidth}{X} 論到死人復活，你們沒有念過摩西書中《荊棘篇》上所記載的嗎？神對摩西說：『我是亞伯拉罕的神，以撒的神，雅各的神。』 \end{tabularx} \\ \\ \relax
12:27 & \begin{tabularx}{0.7\textwidth}{X} 神不是死人的神，而是活人的神。你們是大錯了。」 \end{tabularx} \\ \\
[1ex]
\hline
\hline
\end{longtable}
$^{1}$現在讓我們聆聽一段經文.
經文是記載在《馬可福音》第十二章.
第十八節至到第二十七節.
請聽我讀出這段經文.
「撒都該人常說沒有復活的事.
他們來問耶穌說.
夫子,莫使為我們寫著說.
人若死了,撇下妻子,沒有孩子.
他兄弟當娶他的妻,為哥哥生子立後.
有弟兄七人.
第一個娶了妻,死了,沒有留下孩子.
第二個娶了他,也死了,沒有留下孩子.
第三個也是這樣.
那七個人都沒有留下兒子.
沒了那婦人也死了.
當復活的時候,她是哪一個的妻子呢?.
因為他們七個人都娶過她.
耶穌說:你們所以錯了.
豈不是因為不明白聖經.
不曉得神的大能嗎?.
人從死裡復活,也不娶也不嫁.
乃像天上的使者一樣.
輪到死人復活.
你們沒有念過摩西的書經的篇上所載的嗎?.
神對摩西說:我是亞伯拉罕的神.
以撒的神,雅各的神.
神不是死人的神,乃是活人的神.
你們是大錯了.
今天為我們宣講這段經文的.
是銀行團契的總幹事陳啟興牧師.
幾天之後就不再叫銀行團契了.
就叫做金融業福音團契.
無論怎樣也好.
他們都希望將他們的.
傳福音的領域一直擴展.
昔日由銀行的團隊.
接下來就會擴充到整個金融業.
我不阻礙時間了.
現在將時間交給陳啟興牧師.
為我們宣講今天的主題.

$^{41}$復活的大能.
(主持人:各位在現場的弟兄姊妹早安).
(很感恩再一次有機會在旺朕).
(分享神的話語).
(其實我在主日講道).
(很少教會找同工招呼講完).
(這個很少).
(我知道在這個星期).
(旺朕的同工會已決).
(找一個同工招呼講完).
(這個很窩心的安排).
(雖然最後還是要我請).
(好我們重新禱告).
(我們怎麼都想不到).
(你會到繩肉身來到這個世界).
(我們更加驚訝).
(無法想像你為我們死).
(讓我們有生命).
(讓我們有復活的盼望).
(我們有永生).
(主你這個復活的大能).
(深深觸動和影響我們).
(讓我們今天的生命).
(能夠反映和回應我們的永生).
(主求你用你的話語去幫助我們).
(我們同心禱告仰望).
(奉耶穌基督的聖名).
(阿們).
(在有一年的復活節期間).
(有一對香港夫婦).
(跟了一個旅行團去到聖地).
(他們不是基督徒).
(他們純粹去旅行觀光).
(誰知道在耶路撒冷發生了一個意外).
(我太太意外身亡).
(我導遊就幫忙打點).
(這位旅團友的身後事).
(這位導遊就跟當地的).
(一些殯儀代理聯絡).
(他說如果是葬在當地).

$^{81}$(大約是五百元美元).
(整個喪禮安排).
(但如果要將遺體運回香港去安葬).
(就要一萬元美元).
(這位丈夫就想都不用想).
(很堅持要將太太運回香港).
(他說多錢都可以給).
(所以導遊和團友就盛讚這位丈夫).
(真的愛妻好).
(非常愛惜太太).
(不想她葬在異鄉).
(但這位太太就說其實不是).
(我跟太太的丈夫說).
(我跟太太的關係很差).
(為什麼我一定要運回香港呢).
(因為我聽過一個故事).
(聽說二千多年前).
(有一個人死在這裡).
(他會復活的).
(所以我一定要運走太太).
(當然這是笑話).
(人會不會復活呢?).
(人死後有沒有生命呢?).
(在基督教的喪禮).
(我想大家都多數去過).
(我們不會期望死者).
(在棺木裡面翻生).
(我們沒有這個期望).
(但是我們會相信).
(我們會期望這位基督徒).
(這位離世的基督徒).
(有新的生命).
(這是我們相信的盼望).
(根據幾年前的統計).
(信耶穌的人).
(包括基督教,天主教,東正教).
(在世上大約有23億).
(穆斯林大約有18億).
(印度教徒有11億).
(佛教徒有5億).

$^{121}$(加起來已經超過).
(全世界的四分之三的人口).
(這些人不相信人死如燈滅).
(他們相信死後仍然會有生命).
(如果計算民間的信仰).
(我相信全世界有八成以上的人).
(都相信人死後還有東西).
(還有一個生命).
(如果我問你們或是問基督徒).
(你們是否相信人死後會復活?).
(其實每個人都相信).
(我們行善的會復活得生).
(作惡的會復活定罪).
(行善的代表相信耶穌的人).
(作惡的代表不相信耶穌的人).
(所以我們相信有復活).
(不過有一批人不相信有復活).
(這一批就是殺都該人).
今天這段經文是殺都該人盤問耶穌.
不是盤問.
是挑戰耶穌.
在第十八節.
殺都該人嘗試沒有復活時.
就來問耶穌.
他們來見耶穌.
如果留意殺都該人.
是去到馬可第十二章.
就是第一次.
也是唯一一次出現在馬可福音.
殺都該人是什麼人呢?.
他是猶太人的其中一派.
其他有法利賽人.
有愛色尼人.
有粉瑞黨人.
所以殺都該是其中一派.
他的名字可能是源於.
大衛和所羅門時代的祭司薩都.
有這個可能性.
所以他們和祭司.
和猶太公會的關係很深.

$^{161}$根據猶太歷史學家約瑟夫的記載.
殺都該人是社會較少數的上流人士.
他們是有錢的人.
也是一些有權力的權貴.
不過殺都該人的神學立場很保守.
他們只相信摩西五經.
摩西五經沒有提到復活.
所以他們不會接受.
不會相信復活.
即使你看回舊約.
以塞爾蘇.
以西結書.
以但以理書.
都有提到復活.
但是摩西五經沒有.
他們不接受.
殺都該人也不相信有天使和鬼魂.
在公元七十年後.
聖殿被毀之後.
這群人銷聲匿跡.
再沒有他們很多的記載.
有學者研究認為殺都該人.
最重視的是妥拉和權力.
所以摩西五經和猶太公會.
正正就是殺都該人所倚仗的工具.
他們很重視五經.
以及很重視權力.
殺都該人來找耶穌做什麼呢?.
如果留意在上一段.
希律黨人和花里塞人.
用一個政治問題.
去陷害耶穌.
問他什麼呢?.
應不應該納粹比凱撒.
是一個政治想去陷害耶穌.
在這一段.
殺都該人用一個神學問題.
人死後有沒有復活呢?.
他們的立場是沒有的.
所以用這個問題去挑戰耶穌的教導.

$^{201}$去打擊耶穌的權柄.
所以殺都該人這一次問耶穌.
同樣也是不懷好意.
在新約絕大部分提及殺都該人的場景.
都是在聖殿.
所以我們相信這段對話.
可能有機會在聖殿中發生.
殺都該人只是信摩西五經.
所以他提出一個.
在《新命記》第25章.
5-6節的一個律例.
這個律例怎麼說呢?.
就是說兄弟住在一起.
如果其中一個死了沒有兒子.
死者的妻子就不可以出去嫁給陌生人.
其丈夫的兄弟應該盡兄弟的本分.
娶她為妻.
與她同房.
婦人生的長子要歸在已故兄弟的名下.
免得他的名在以色列中逃去了.
這就是這個律例.
這個規矩我們叫做弟娶兄相.
或者如果你用英文的音譯.
叫做利美拉特的婚姻.
就是兄弟會娶了已故兄弟的太太.
為他的兄弟立後.
在《創世記》我們看到.
其實有弟娶兄相的例子.
所以你看到.
其實這些殺都該人.
看摩西五經.
都是看一點不看一點的.
其實明顯有一個例子.
我們知道猶大的長子叫做爾.
他的息父叫做他馬.
當爾死了之後.
猶大就叫爾的弟弟.
你應該娶哥哥的太太.
拉瑪為妻.
為他立後.

$^{241}$這個很明顯.
有弟娶兄相的規矩.
按照這個規矩.
殺都該人就想了一個很假設性的問題.
兄弟七人.
娶了老婆之後逐一死去.
弟娶兄相.
一個一個嫁給一個兄弟.
如果將來復活.
他是哪一個人的太太呢?.
因為他們七個人都娶過他.
所以結果就會一塌糊塗.
他們想用這個假設性的問題.
去證明沒有復活.
摩西也反對復活的.
摩西也這麼說的.
如果大家留意《新命記》.
第25章這一段.
原意是什麼呢?.
是為死者留後.
原意根本不是想討論復活.
這段經文跟復活是沒有關係的.
殺都該人就是用來設計一個假設性的情況.
用來攻擊耶穌.
所以你這麼看.
其實殺都該人看聖經.
他沒有留意到暢世紀有這個例子.
他們看一點不看一點.
他們不是很明白聖經.
不是很完全明白聖經的真正意思.
我想用殺都該人這個段落.
去提出來.
其實我們去反省.
我們會不會也是.
其實有時候不是很明白聖經.
甚至是曲解聖經.
誤用聖經真理呢?.
我跟很多弟兄姊妹去談的時候.
大家都很用心去.
進行十一奉獻.

$^{281}$十一奉獻最早.
建於暢世紀.
亞伯拉罕去獻給默基西德.
一直到舊約最後的馬拉基書也有.
我們將十分之一獻出.
看看神會不會闖開天上的窗戶.
傾福給我們.
甚至無處可用.
所以弟兄姊妹很進行十一奉獻.
少一個崩也不行.
因為出了糧一定要十一奉獻.
但反過來多一個崩也不行.
不要多一個十一奉獻.
很規矩地.
剩下十分之九.
是誰的呢?.
當然是自己用.
我已經給了十分之一旺振.
剩下的我自己用.
究竟聖經是否這樣解釋呢?.
今天我不是想解釋十一奉獻.
我們很多時候會誤解了.
我奉獻了十分之一.
剩下的百分之九十.
就是我自己用.
我們對聖經.
其實有時候可能是.
不是很完全明白.
好像殺刀該人一樣.
其實不明白也不要緊.
其實我對聖經很多都不懂.
但我不滿意自己不懂.
我嘗試努力學習.
馬丁路德說.
他說聖經不是舊的.
也不是新的.
而是永遠.
牛頓這位科學家說.
聖經的真理比世界上任何歷史.
更加可靠.

$^{321}$德國的大文豪戈德.
他說如果我坐牢.
只能帶一本書進去.
我就會選擇聖經.
美國的總統林肯.
他說聖經是神賜給世人最美好的禮物.
包含所有能夠得到幸福的真理.
連這個非常好戰的拿破侖.
他說聖經不是一本普通的書.
是一個生命的機體.
能夠征服所有反對聖經的.
你信聖經也好.
不信聖經也好.
很多人去接觸過聖經.
如果你真的認真去理解聖經.
你會覺得.
不是一本書.
是承載你的生命.
去支撐你的生命.
是上帝的真理.
所以.
在這個段落.
就是第一個段落.
殺刀該人這個提問.
去反映殺刀該人對聖經是一般般的.
不是很認識.
我們對聖經是怎麼樣的呢?.
我們對聖經會不會也是.
不是很認識.
但又不是很肯竭力去明白.
去做一個無愧的工人.
按著正義去分解真理的毒.
所以我想停一停.
在這裡給大家想一個形容.
這個形容.
我在舞會也提出過.
也推行過.
我們叫做一小時預備查經.
在教會裡面.
我們有團契.

$^{361}$我們有小組.
我們有成長班.
有很多不同的群體.
我們多數都查經.
我們教會很少有群體不查聖經的.
總有一些州會去查聖經.
我很建議就是.
每當有查聖經機會的時候.
我就去準備.
有準備和沒有準備的查經是差很遠的.
如果你有準備查經.
經過一年,兩年,三年.
你對聖經是很不同的.
不用複雜.
一小時就可以了.
當時我就鼓勵弟兄姊妹.
每有查經機會.
就是團聚查經.
可能一個星期一次.
可能兩個星期一次.
我們就用一個小時去準備.
也很足夠了.
如果你日積月累.
你對聖經的理解.
會有相當不同的感受和體驗.
你會經歷聖經真的活潑的道.
在你的生命裡面會產生功效.
我們不要像殺刀戈人那樣.
知道一點不知道一點.
誤解聖經.
亂用聖經.
鼓勵旺朕的弟兄姊妹.
去嘗試.
嘗試一小時預備查經.
我相信.
不用很久.
你的生命可能會被聖經改變.
真的明顯.
我們看看耶穌怎麼回應.
第二個段落.

$^{401}$就是24-27節.
耶穌怎麼回應呢?.
耶穌說.
你們錯了.
所以錯了.
豈不是因為不明白聖經.
不曉得神的大能嗎?.
很有趣的.
殺刀戈人倚重的兩大武器.
摩西五經和權力.
耶穌的回應就是.
你們錯了.
你們不明白聖經.
你又不知道神的權能.
好像左右開弓.
就打了他們兩巴掌.
人從死人中復活之後.
也不嫁.
也不娶.
好像天使一樣.
當然不是人復活之後.
我們的本質變成了天使.
耶穌不是這個意思.
好像天使一樣.
是指天使沒有嫁娶.
天使也不會再死亡.
所以我們復活之後.
我們不會有嫁娶.
也不會再有死亡.
如果你這樣看.
或者你這樣理解聖經.
復活的生命.
不是地上生命的延續.
我們死了之後.
有一個延續在天上的生命.
是一個截然不同的生命.
是一個完全跟你地上生命不同的生命.
不是你和我可以想像到的生命.
這個生命唯有神的大能.
才能夠做得到.

$^{441}$復活這個字在原文絕大部分都是被動式.
在聖經上我們叫做神性的被動式.
Define passive.
即使我們.
什麼呢?.
我們只能夠讓神的大能去復活.
不是我們能夠去復活.
也是神的大能.
唯有神的大能.
讓我們有這麼不同的生命.
復活之後是怎樣的呢?.
很多人都問這個問題.
我以前初信也好.
信了一段時間.
我也問這個問題.
究竟復活之後的我.
是十八歲朝氣勃勃的我.
還是八十歲氣力衰殘的我.
皺皮的我呢?.
復活之後.
能不能夠認得回來呢?.
我能不能夠認得太太呢?.
既然沒有嫁娶.
為什麼喪禮經常提到.
你有重罪的盼望.
能夠見到配偶呢?.
其實很多這些問題.
我自己當時也去思考.
嘗試去看聖經有沒有解答.
當然最詳細的解答就是在.
哥林多前書第十五章.
那裡.
說到復活之後.
是一個不會扭壞.
一個很榮耀的軀體.
這個軀體.
一定比任何歲數的你好.
你不用擔心.
如果你很帥.
像劉德華一樣.

$^{481}$復活之後會不會變成像陳牧師一樣.
這個陳牧師.
不是那位.
那位是相當帥的.
不用擔心.
復活之後的你是一個非常榮耀.
不會扭壞.
比任何生命階段.
都要美好的一個生命.
你想想.
在天家.
就不會再有死亡了.
所以你不需要結婚生子.
沒有這個問題.
也沒有罪惡的.
人與人的關係比夫妻更加甜蜜.
在今生.
在人生.
我們最親密的關係.
就是夫妻的結合.
但是夫妻的結連.
其實也是有缺陷的.
因為兩夫妻都有罪性.
我們也有問題的.
但是將來在天家就是千千萬萬的信徒.
在天家裡面結連.
這個結連是沒有罪.
在當中難阻的.
所以在天家我們見到配偶.
當然我們會歡喜快樂.
但是我們完全了.
在一個我們想像不到這麼美好的完美世界.
與神永遠同在.
這是一個無法想像的美好.
假如我們知道.
我們相信神的大能.
我們就會相信復活.
也不會以地上的生命.
我們今天經歷的.
去妄斷天上的生命是怎樣怎樣.

$^{521}$好像這群殺到該人一樣.
殺到該人用摩西五經.
耶穌都引用在《初拉及記》.
應該是第三章.
耶和華在荊棘的火焰中.
向摩西顯現這一段.
神跟摩西說.
我是亞伯拉罕的神.
是以撒的神.
是雅各的神.
是什麼意思呢?.
當我們去細讀《創世記》.
我們知道在第十七章第七節.
耶和華向亞伯拉罕說.
我要與你以及你世世代代的後裔.
建立我的約.
成為永遠的約.
是要作你和你後裔的神.
就是以撒,雅各.
你所有後裔的上帝.
耶和華說什麼呢?.
是永遠的約.
不會因人的死亡而終止.
這是永遠的約.
人死後這個約仍然有效.
為什麼呢?.
因為人會復活.
所以耶穌補充.
神不是死人的神.
而是活人的神.
在事實上耶穌教導有復活.
很清楚駁斥了殺到該人.
他自己親身去證實有復活.
他就是從死亡中.
被天父父神復活的那一位.
他是初熟的果子.
我們相信他.
我們會跟著他復活得永生.
殺到該人去否定復活.
是錯的.

$^{561}$是大錯特錯.
在上世紀有一個美籍的懷著科學家.
他叫徐道國.
這個徐道國.
雖然不是很出名.
不過他有哺乳動物細胞遺傳學之父的稱譽.
他在美國德州的研究中心.
收藏了超過300種瀕臨絕種動物的基因.
他盼望後世的科學家能夠利用這些基因.
能夠重造這些將來會絕種的動物.
這個橋段在侏羅紀公園應該也見過.
徐道國說什麼呢?.
他說如果透過我的工作.
在200年之後.
小朋友能夠看到這些已經絕種的生物.
我想我已經留下了一份相當有價值的遺產給這個世界.
但我想一想.
當然這是很大的貢獻.
但即使這些動物復活.
它也會死的.
這些動物可能能夠生存一個時間.
但有一天仍然有機會會再度絕種.
我們不需要什麼基因重建.
其實我們在主裡面的人一定會復活.
耶穌留下這個遺產.
比任何人留下的遺產.
更加有價值多不知多少千萬倍.
我們不需要搞盡腦汁去想一些事情.
我們只要相信耶穌.
就能夠有復活的盼望.
蘇格拉底說.
人會不會復活呢?.
他說我希望會.
但沒有人知道.
所以我們今天.
如果從一個信仰角度.
我們在這方面比蘇格拉底厲害.
我們知道.
我們有復活的把握.
坦白說這段經文很淺白.

$^{601}$對我們基督徒有什麼用呢?.
如果我問大家.
你們信不信有復活.
我想大家全部一致都會答信.
你信耶穌的人.
都會信復活.
但我們想深一層.
我們會不會頭腦相信.
我們今天的生命.
是不是反映將來復活這個事實呢?.
這個值得我們多想一點.
我們今天所投身的.
是不是真的有永恆的價值呢?.
復活的確具.
至少帶給我們兩件事情.
我自己常常去思想.
去想想.
究竟在我生命裡面.
有沒有這兩件事情呢?.
第一就是.
復活帶給我們有盼望.
這個復活的盼望.
帶給我們真真正正的平安.
無論在疫情裡面.
有很多人死亡.
今天已經有.
全球超過200多萬.
接近300萬人死亡.
很多人在很多哀傷裡面.
但是復活的盼望.
給我們有真真正正的平安.
讓我們的眼光超越在地上生命的終點.
能夠看得遠一點.
在基督教的喪禮裡面.
大家應該很多都經歷過.
安息主懷的人.
去到一個更美好的家鄉.
這是一個很大的安慰.
這個復活的盼望帶來的平安.
我們每個人都會欣然接受.

$^{641}$相信大家都會接受.
大家都會經歷.
但是第二件事情.
就是復活的生命.
如何決定你今天的生命.
這個更加值得我們去思考.
這個復活的生命.
有沒有影響到你今天的生命呢?.
未來決定現在.
你同不同意嗎?.
你如何為你未來復活的生命.
去預備你的今天呢?.
其實2021年好像剛剛到.
但是已經想了想.
已經過了差不多四分之一.
人生匆匆.
你如何去投資永恆.
很值得我們認真去思考.
在很多年前.
葛培里在世的時候.
就被邀請到華盛頓.
他接受美國政府頒發的國會金章.
當時的總統和國會的一些議員.
就在場頒發國會金章.
他們請葛培里牧師說幾句話.
葛培里這麼說.
他說親愛的總統.
國會議員.
以及各位政府官員.
請問在這個危鵝紅衛的大廳裡面.
你們看到什麼呢?.
那些人就四處看.
就說歷代總統的雕像.
歷代總統的畫.
葛培里牧師就問.
你們看到他們有沒有什麼共同點呢?.
他們的衣著.
他們的勳章.
他說都不是.
他說這些歷史人物有一個共同點.

$^{681}$就是他們全部都死了.
林肯死了.
傑弗遜死了.
羅斯福死了.
所有這些偉大的人物.
都成為歷史.
只留下雕像和畫.
他們現在在哪裡呢?.
我希望你們.
就是這一群國家的領袖.
能夠思考這個嚴肅的問題.
將來你要在哪裡過你永恆的生命呢?.
十九世紀英國的藝術家John Ruskin.
他也是一個詩人.
以及一個哲學家.
他說了一段讓人去三思的話.
他說每一個活在地上的人.
都應該思考三個問題.
我從哪裡來.
我要去哪裡.
在中間我要怎麼去度過呢?.
很值得我們思考.
我們知道我們從哪裡來.
從神的創造而來.
我們知道我們會去哪裡.
我們會永恆地與愛我們的主一起.
但是我們中間怎麼過呢?.
我們中間過的生命.
是不是能夠反映.
我們將來與主同在的永恆生命呢?.
所以應用就是.
今天你真的投身在什麼事情上呢?.
你奮身,毀身.
你花很多時間精力在什麼事情上.
在復活的生命中.
這些有沒有價值呢?.
如果沒有價值.
你怎麼要調整你的生命.
將你的生命放在有永恆意義.
有永恆價值的事情上呢?.

$^{721}$所以我想到最後總結一下.
這段經文.
很簡單的兩件事.
第一個段落就是從薩都蓋人一個很愚蠢的提問.
我們要反省.
我們對聖經的理解.
其實是否足夠呢?.
或者是否好呢?.
我鼓勵的應用就是.
很簡單.
一小時預備查經.
從有機會我們一起查經.
我們用心去學習.
到不用團契查經.
不用小組查經.
你自己會查聖經.
你自己有問題.
你會在聖經裡面.
找到答案.
找到智慧.
我們從耶穌一個極有智慧的回應.
我們應該深深相信.
我們有復活的生命.
我們要反省的就是.
今天.
我們是否正在預備.
或者是預備好.
去迎接復活的生命呢?.
未來決定現在.
我今天要如何去調整我的生命.
預備進入永恆的生命裡呢?.
很坦白說.
在2021年這幾個月.
很多人離開.
很多名人.
很多身邊的人離開.
其實生命不是很穩定.
生命可以隨時離開這個世界.
即使我們長命到百歲.
我想我們也可能只有幾十年.

$^{761}$轉眼就過了.
很多人已經過了一半.
到了人生的下半場.
我們是否應該去預備永恆.
多於在人生的幾十年裡.
去搞盡腦汁.
很多的籌算呢?.
這個很值得我們去深思.
我今天.
其實我的角色是一個.
就是說復活的信息.
像什麼呢?.
像是思洗約寒.
我主要是預備下個星期.
在復活節.
那一位去說一個真的復活的信息.
希望帶給你有更大的幫助.
我們同心同告.
請教主在臨近這個復活節.
讓我們真的去深深思想.
為什麼你要復活?.
你的復活帶給我們有什麼衝擊?.
我們終於在聖經裡明白.
原來你復活也令我們可以復活.
我們有永恆的生命.
主幫助我們.
我們今天所作的.
我們能夠投資於永恆.
在永恆裡面有價值.
將那些沒有意義沒有價值的.
幫助我們.
讓我們有力量去放下.
去捨棄.
主我們要靠著你.
因為唯有你可以令我們復活.
也唯有你可以令我們今天.
能夠過一個指向復活的生命.
願你祝福皇上每一個弟兄姊妹.
我們有力的.
在今天為你而活.

$^{801}$預備好未來.
我們同心同告仰望.
奉耶穌基督的聖名.
\newpage



\section{馬可福音 15:21-32}
\label{sec:NcnNHWwVDiw}
\textbf{2021年4月2日受苦節崇拜直播｜吳真真牧師｜十字架的愛｜馬可福音15\_21-32}
\newline
\newline
連結: \href{https://youtube.com/watch?v=NcnNHWwVDiw}{\texttt{https://youtube.com/watch?v=NcnNHWwVDiw}} ~~~~ 語音日期: 2021-04-02
\newline
\newline
\hyperref[sec:pK4j32Egyx0]{\small{< < < PREV SERMON < < <}}
~
\hyperref[sec:index_chronic]{\small{[返順時目]}}
~
\hyperref[sec:R7EPR0TUgHE]{\small{> > > NEXT SERMON > > >}}
\newline
\newline
馬可福音 15:21-32
\newline
\begin{longtable}{cl}
\hline
\hline
章節 & 經文 (和合本修訂版)\\
\hline
15:21 & \begin{tabularx}{0.7\textwidth}{X} 有一個古利奈人西門，就是亞歷山大和魯孚的父親，從鄉下來，經過那地方，他們就強迫他同去，好背耶穌的十字架。 \end{tabularx} \\ \\ \relax
15:22 & \begin{tabularx}{0.7\textwidth}{X} 他們帶耶穌到了一個地方叫各各他（翻出來就是「髑髏地」）， \end{tabularx} \\ \\ \relax
15:23 & \begin{tabularx}{0.7\textwidth}{X} 拿沒藥調和的酒給耶穌，他卻不受。 \end{tabularx} \\ \\ \relax
15:24 & \begin{tabularx}{0.7\textwidth}{X} 於是他們把他釘在十字架上，抽籤分他的衣服，看誰得甚麼。 \end{tabularx} \\ \\ \relax
15:25 & \begin{tabularx}{0.7\textwidth}{X} 他們把他釘十字架的時候是上午九點鐘。 \end{tabularx} \\ \\ \relax
15:26 & \begin{tabularx}{0.7\textwidth}{X} 罪狀牌上寫的是：「猶太人的王。」 \end{tabularx} \\ \\ \relax
15:27 & \begin{tabularx}{0.7\textwidth}{X} 他們又把兩個強盜和他同釘十字架，一個在右邊，一個在左邊。 \end{tabularx} \\ \\ \relax
15:28 & \begin{tabularx}{0.7\textwidth}{X} / \end{tabularx} \\ \\ \relax
15:29 & \begin{tabularx}{0.7\textwidth}{X} 從那裡經過的人譏笑他，搖著頭，說：「哼！你這拆毀殿、三日又建造起來的， \end{tabularx} \\ \\ \relax
15:30 & \begin{tabularx}{0.7\textwidth}{X} 救救你自己，從十字架上下來呀！」 \end{tabularx} \\ \\ \relax
15:31 & \begin{tabularx}{0.7\textwidth}{X} 眾祭司長和文士也這樣嘲笑他，彼此說：「他救了別人，不能救自己。 \end{tabularx} \\ \\ \relax
15:32 & \begin{tabularx}{0.7\textwidth}{X} 以色列的王基督，現在從十字架上下來，好讓我們看見就信了呀！」那和他同釘的人也譏諷他。 \end{tabularx} \\ \\
[1ex]
\hline
\hline
\end{longtable}
$^{1}$我們繼續聆聽經文來敬拜主.
經文是記載在《馬可福音》15章21至32節.
有一個古利來人西門.
就是亞歷山大和魯夫的父親.
從鄉下來經過那地方.
他們就勉強他同去.
好背著耶穌的十字架.
他們帶耶穌到了各個他地方.
各個他翻出來就是獨留地.
那沒藥調和的酒給耶穌.
他卻不受.
於是將他釘在十字架上.
念叩分他的衣服.
看是誰得什麼.
釘他在十字架上是自初的時候.
在上面有他的罪狀.
寫的是猶太人的王.
他們又把兩個強盜和他同釘十字架.
一個在右邊一個在左邊.
從那裡經過的人辱罵他.
搖著頭說:.
「唉!這七位聖殿三日又建造起來的.
可以救自己從十字架上下來吧」.
祭司長王文氏也是這樣戲弄他.
彼此說:.
「他救了別人不能救自己」.
以色列的王基督.
現在可以從十字架上下來.
叫我們看見就信了.
罵和他同釘的人也譏笑他.
接著吳錚錚牧師透過經文勉勵我們.
題目是「十字架的愛」.
一首在十九世紀寫成的聖詩.
靠近十架.
歌詞是由一位在美國失明的女詩人所寫的.
Fanny Grosby一出生六個星期的時候.
她對眼睛發炎.
後來醫生治不好她.
沒想過就由那天開始.
她的雙眼就瞎了.

$^{41}$後來在她的傳記中.
她提到有一位弟兄.
很可憐她.
認為她自六個星期大開始已經瞎了眼.
一定是很自卑 很自憐.
很多怨恨.
但是她告訴弟兄.
「我巴不得一生下來.
我就是一個「核」字」.
弟兄非常不明白.
「為什麼?」.
Fanny的回答是.
「如果是這樣.
當我有一天去到主那裡的時候.
第一眼 我的眼睛所看見的.
就是我親愛的主耶穌」.
為什麼Fanny幾乎一天生就失明.
她仍然可以這麼愛主耶穌.
這麼願意靠近十字架.
十字架為什麼會這麼奇妙呢?.
英國歷史學家湯恩比在他的著作中.
曾經提到有一個中國女孩子.
來到英國深造.
在她假期的時候.
她去到一個基督徒的家裡.
做孩子的保姆.
當她去到那間屋的時候.
她看見大廳的中間.
懸掛著一幅耶穌釘十字架的油畫.
她覺得很不自在.
於是她問這位主人.
「為什麼你要在這個廳裡.
掛一幅這麼殘酷的圖畫?」.
主人聽到之後.
他不以為然地說.
「這幅不是一幅殘酷的圖畫.
這幅是一幅愛的圖畫」.
其實十字架往往給人的感覺.
都會有兩種.
一種就好像中國女孩子這樣想.

$^{81}$感覺到很殘酷.
因為它是一個形具.
另一個感覺就是愛.
因為耶穌在十字架裡.
為我們犧牲.
很多人.
尤其是未信主的人.
都搞不清楚.
很不明白十字架上的耶穌.
是一個號刑下的受害者.
還是一個愛人的救世主呢?.
就好像清朝的高官李鴻章.
他曾經有一個疑問.
他說一向崇尚武力和尊敬的西洋人.
為什麼會崇拜一個釘在十字架上的失敗者?.
由今天星期一開始到現在.
教會的靈修系列的目擊者.
分別都說了有很多不同的人.
親眼目擊耶穌釘十字架.
他們心裡都存著這份疑惑.
彼拉多,巴拉巴拉,.
巴拉爆,冰釘都存在.
的確.
今天讓我們好好反省.
究竟十字架對我們的意思是什麼.
讓我們同心一起祈禱.
主啊.
無論我們是何等樣的人.
我們處身在何等樣的境況.
你在十字架裡的寶血.
都是為我們而流.
多謝你.
由你引導我們.
讓我們若果眼睛是黑眼的.
你今晚就去跟我們說話.
幫助我們去明白十字架的意義.
在現在的境況的我.
應該怎樣去回應.
求你引導.
多謝你聽祈禱.

$^{121}$奉主耶穌基督的聖名自祈求.
阿們.
今晚我們會從《馬學福音》十五章.
來思想耶穌在十字架裡面.
給我們的一些提醒.
在第21到22節.
耶穌經歷了受審.
他被判釘十字架.
到他被釘釘棄輪.
現在他準備去刑場.
行刑的地方叫「葛國他」.
意思是一個骷髏骨頭.
是因為他的地形的形狀.
好像一個骷髏骨頭.
又或者是因為那裡是一個經常行刑的地方.
葛國他的位置.
其實在耶路撒冷西北邊.
和耶穌被審的那個地方.
其實相距只是幾百公尺.
不過.
因為羅馬人為了要殺一儆百.
通常會將釘十字架的囚犯遊街示眾.
根據第一世紀羅馬一位著名的教育家.
昆體良的記載.
他說釘十字架之前.
會帶著犯人.
專程行經最多最多人的街道.
讓最多最多的人看到.
令人們能夠害怕.
不要犯罪.
所以其實這條通往刑場的道路.
一點也不短.
十字架打直的那一條木.
早已經在刑場了.
而耶穌所背負的十字架的橫樑.
就是要走去葛國他.
因為這條路確實不是太短.
而耶穌之前也受盡癡刑.
他被鞭打得很虛弱.
根本沒有力氣去背起刑具.

$^{161}$羅馬的癡刑其實很殘忍.
癡刑所用的工具就是一個短節鞭.
是一條皮製成的鞭.
但上面是鑲了一些骨碎片和圓塊.
有時會有一些小勾.
當鞭打到犯人身上的時候.
就會將他的皮撕開.
皮開肉爛.
犯人往往會打得痛不欲生.
奄奄一息.
耶穌很虛弱.
他走不動.
羅馬的士兵強迫一個叫西門的人.
來背負刑具.
因為羅馬的官員.
是可以有權隨時徵用任何人幫他們做事的.
而這個西門其實是北非人.
古利內人.
北非古利內就是現在的利比亞.
一些色經家認為.
其實是一個外國人.
他住在外地.
他來到耶路撒冷是為了過如意節.
而在這個時候他就被迫去代耶穌背負刑具.
在《使徒行傳》十三章第一節裡.
記載有一個叫做「離結的西面」.
西面其實是西門的另一個說法.
「離結」的意思是指皮膚黑的人.
所以北非的西門.
可能就是《使徒行傳》十三章第一節所說的西面.
馬可在這裡更加提到兩個人.
就是西門的兩個兒子.
亞歷山大和魯夫.
對於馬可的福音讀者來說.
他們這班羅馬信徒.
可能就算不認識西門也好.
他們也會認識亞歷山大和魯夫.
聖經裡記載亞歷山大的事.
我們不太清楚.
但魯夫可能就是羅馬書十六章十三節.

$^{201}$提到那位羅馬的蒙主選擇的信徒.
如果西門是《使徒行傳》十三章一節的西面.
他在安提厄教會裡.
他是受聖靈感動去按手差派保羅和巴拿巴.
作第一次傳道旅程的領袖之一.
而魯夫是保羅所親愛的弟兄.
保羅形容他視魯夫的母親.
就好像自己的母親一樣.
所以我們可想而知.
西門和魯夫他們都是早期教會裡面的領袖.
或者在靈命裡面比較突出的人.
西門本來是居於外邦的.
他們兩個兒子的名字.
其實都是用外邦的名字來起名.
為何他們一家會成為教會的領袖.
他們的靈命會突飛猛進.
因為西門被迫代耶穌背負十字架.
本來是一百萬個不情願.
不過我們相信當他去到葛國他.
他見到耶穌死的時候的景象.
他聽到耶穌在十字架上的震撼說話.
「父啊,赦免他們.
因為他們所做的他們不知道」.
我相信西門改變了.
也帶來他們兩個兒子的改變.
耶穌教導他們的門徒.
要背起十字架去跟從他.
但當耶穌被捉拿的時候.
原本跟從他的那群門徒四散.
但當耶穌在富瀛牆的路上幾乎陷落的時候.
西門就背起十字架.
跟主一起走這條苦路.
到了營牆.
他見到耶穌臨死的景象.
甚至在15章39節有提到.
十字架下那位白夫長評價耶穌.
這個真是神的兒子.
於是西門的人生被折服.
作了一個很大的改變.
弟兄姊妹.

$^{241}$有時我們會很不情願.
甚至被迫去面對一些逆境.
但你有沒有留意.
有時會帶來一些正面的改變.
甚至人生都會得到逆轉.
就好像西門他們一家一樣.
關鍵在哪裡?.
第一.
不要怕逆境.
因為神叫萬事都互相效力.
叫愛神的人得益處.
上帝必然有美好的心意在背後.
第二.
就是要背負耶穌的十字架.
為何要背負耶穌的十字架?.
就是你接受十字架的救恩.
擁抱十字架的大愛.
信靠十字架的大能.
你就有能力去面對逆境.
逆境亦會帶來正面和積極的影響.
曾經有一位姐妹.
她在讀大學的時候.
她立志她的人生要奉獻給上帝.
要做一個傳道人.
不過大學畢業之後.
她發覺神沒有開路.
於是她就打算找工作.
暫時在工作上工作.
但當她工作的時候.
她接二連三地遇到一些.
不好脾氣的上司.
又因為她比其他人.
都更加盡責和奮心工作.
所以很多艱難的工作.
都要她去完成.
有一晚.
經過了這麼久以來.
高壓的生活.
她終於承受不住.
她對著上帝抱怨.

$^{281}$「主啊!為何是我?.
這麼難相處的人.
這麼麻煩的工作.
為何偏偏要讓我遇到呢?」.
突然間她心裡有一把微小的聲音.
跟她說.
「親愛的孩子.
為何不可以是你?」.
「親愛的孩子.
為何不可以是你?」.
她擦乾了眼淚.
她開始思考這條問題.
「為何不可以是我?」.
當她再度安靜在神面前的時候.
她對神說.
「神啊!多謝你.
你看得起我.
多謝你願意教導我.
我願意成為在職場裡的宣教士.
我甘心接受.
你對我作出任何的挑戰」.
於是乎.
以後無論遇到甚麼困難.
她不再跟神說.
「為何是我?」.
她反而說.
「神啊!多謝你.
因為你愛我.
所以才會是我」.
所以弟兄姊妹.
我們不要怕逆境.
我們要擁抱十字架的救恩和愛.
因為其實現在的我們.
是可以帶來愛.
因為其實逆境是可以帶來成長.
可以帶來祝福.
可以帶來重生.
來到第23節.
耶穌已經到了葛國他.
「使兵拿沒藥調和的酒給他喝.

$^{321}$沒藥調和的酒是有麻醉作用的.
因為耶穌當時.
真的受盡肉體極度的痛苦.
但使兵未必是出於好意.
因為麻醉的作用.
可以延長犯人的生命.
延長他的痛苦」.
耶穌拒絕去喝這個酒.
我想.
耶穌拒絕.
是因為他不需要這些短暫的麻醉.
為了拯救世人.
他願意清醒地去喝這個苦杯.
因著他這一份清醒.
他後來在十字架上.
所說的十加七言.
表現了驚人的威力.
正如剛才所說.
不單止折服了西門.
折服了十字架上另一個囚犯.
也折服了十字架下那位白夫長.
還有.
耶穌受盡十字架的痛.
他清醒地背負了我們的罪.
和罪帶來的後果.
才可以和我們每一位.
身處痛苦的人說.
「我明白你一切的痛苦.
你所經歷的.
我曾經經歷過」.
所以耶穌絕對絕對去明白.
我們的痛苦.
弟兄姊妹.
其實面對生命的困難和逆境.
我們都會選擇逃避.
我們甚至會選擇.
用一些方式去麻醉自己.
但是這裡告訴我們.
耶穌鼓勵我們.
不要這樣.

$^{361}$不要活在麻醉自己的境況.
不要活在假象裡面.
因為這樣是沒有幫助的.
我們始終要面對.
因為.
神恩典救容.
世上有很多人.
都會做很多事去幫自己麻醉.
有些人說吸煙.
可以幫他忘記煩惱.
有些人沉醉在往來的世界裡面.
我有些朋友長期.
不眠不休地去煲劇.
說他們的感覺是暫時可以脫離一下.
那個困境和痛苦.
但是.
其實.
我們始終都要面對.
要去經歷這個痛.
唯有靠著十字架的恩典.
我們可以做得到.
我們可以去面對這個痛苦.
對於猶太人來說.
絕對不能夠接受.
他們的彌賽亞是這樣死的.
對於外邦人來說.
神也不會蠢到這樣去死.
所以保羅說得對.
我們全被釘十字架的基督.
這對猶太人是絆腳石.
對外邦人是愚蠢.
因為十字架的道理.
在那滅亡的人是愚蠢.
在我們得救的人.
卻是神的大能.
神的智慧高過我們的智慧.
我們猜不透祂的心意.
但是藉著祂自己.
降生受苦.
道成肉身的啟示.

$^{401}$聖經真理的啟示.
我們今天很清楚知道.
十字架是救恩.
是神捨己犧牲的愛.
是猶太人沒有想過的神蹟.
也是希臘人沒有想過的智慧.
更加可能是在在.
當今這一代的人.
沒有想過的人生答案.
如果我們的人生.
沒有十字架會是怎樣?.
如果我們的人生.
沒有十字架會是怎樣?.
可能我們見到很多人.
為名為利為權為色.
為滿足自己的慾望.
而付出一生.
結果到最後.
這些事情全部都會過去.
因為這些事情.
全部都沒有永恆價值.
不單止未信神的人.
就算是信徒.
我們有時也值得.
好好反省.
因為我們有時.
在這條天路上.
我們也會失去焦點.
失去方向.
好像星期三的祈禱會.
我們一起走進.
耶穌門徒彼得的心.
他是耶穌門徒當中的門徒.
他用自己的方法.
找尋他生命的出路.
他一直以為自己.
被其他人出眾.
更加值得.
做眾人的領袖.
不過.

$^{441}$愛他的老師.
主耶穌.
不容許他活在這個假象裡.
不容許他再在這個麻醉裡.
耶穌當著門徒的面前.
去揭穿他.
今晚.
雞叫兩片二線.
你要三次不認我.
撕開了彼得.
他外表那種有為熱血.
但其實.
內裡的自己.
都是滿足自己的生命.
讓他體無完膚.
清清醒醒地.
去面對自己內心那種驕傲.
看到他內心.
原來仍然很自我.
原來他看到他的內心.
仍然死不透.
唯有.
他能夠認真去面對這個痛處.
這個醜陋的自己.
他才可以.
全然去倚靠這位.
被釘十字架的主.
他才能夠.
成為一個真正.
謙卑的義僕.
成為合神心意.
在初期教會的領袖.
其實.
追求滿足自己自我中心的人生.
其實的確是很苦的.
但十字架的愛.
告訴我們.
認識耶穌.
效法他捨己的愛.
這才是人生最好的出路.

$^{481}$我們愛.
因為神先愛我們.
讓我們學習放下自己.
去愛身邊的人.
讓我們真是親身去經歷到.
施比受更為有福.
人生活出豐盛.
才是我們永恆的盼望.
去到第29節到32節.
我們看到耶穌除了肉體上.
經歷極大的痛苦和挑戰外.
我們看到.
耶穌都經歷了一個.
屬靈上極大的挑戰.
29到32節.
耶穌受了三批人的譏笑.
第一批.
就是剛剛經過的人.
第二批.
就是祭司長和文士.
第三批.
就是和他一起釘十架的囚犯.
他們怎樣挑戰.
耶穌的屬靈權柄.
路過的人說.
你這個三日建造聖殿的.
救救你自己吧.
祭司長和文士說.
以色列的黃基督.
從十字架上下來.
我們就相信你.
我們看到.
馬太福音27:40.
更加說.
如果你是神的兒子.
你就救,你就跳下來吧.
各位.
耶穌真是神的兒子.
他真是以色列的王.
他真是基督.

$^{521}$他真是能夠.
三日建造聖殿.
如果他少了一分的定力.
如果他不是.
決意要進行天父的使命.
他真的可以從十字架下來.
不過耶穌有屬靈的權柄.
但是他卻自限.
不用.
這就是他面對屬靈挑戰的地方.
其實耶穌為何不從十字架下來.
讓人們可以相信他.
耶穌其實是要找.
真正信心的人.
因為.
看見才相信.
不是真正的信心.
信心不是.
看到神蹟的結果.
不過信心是.
經歷神蹟的條件.
而且.
耶穌不從十字架下來.
為了的是要.
彰顯完全的愛.
救世君的創立人.
卜威廉他曾經說過.
正因為耶穌甘願.
受十字架的苦.
他沒有從十字架下來.
我們才相信他.
耶穌為人類而死.
他來到世上的目的.
就是要彰顯.
上帝完全的愛.
倘若他拒絕這個十字架.
真的從十字架下來.
那就代表了.
上帝的愛是有限的.
原來有些事.

$^{561}$耶穌是不願意替人做的.
原來有一個界限.
上帝的愛是超越不了的.
但是耶穌.
彰顯了完全的愛.
他為人走了.
一整段的道路.
甚至死在十字架上.
這個做法.
就是要告訴我們.
上帝的愛是無止境的.
世上沒有一件苦差.
是他不願意為我們做的.
世上沒有任何一件艱鉅的任務.
甚至死在十字架上.
他不肯為我們擔當.
你有沒有試過.
被人質疑.
質疑你所信的耶穌.
你有沒有試過被人挑戰.
你的屬靈權兵.
有沒有人曾經批評過你.
你究竟會不會講道.
你是組長.
你是大茶經.
你是我們的領袖.
耶穌榜樣告訴我們.
當我們受到挑戰的時候.
我們不是第一時間去抗辯.
去反擊.
要常常念既.
成就天父的旨意.
就是要愛對方到底.
以建立對方對上帝的信心.
為最重要的任務.
耶穌被釘十架.
他的身心靈受著極大的痛苦.
不過我相信.
耶穌不是要我們很感性地.
覺得他很慘.

$^{601}$很可憐.
耶穌在十字架.
所受身心靈的痛苦.
是要表達.
他捨己的愛.
他為了救我們.
他愛我們到底.
即使要付上.
多大的犧牲.
印度有一位報道家叫孫大信.
他致力向西藏人傳福音.
有一次孫大信.
和一個西藏人要去高山.
傳道.
但不幸遇到一場大風雪.
他們見到山邊有一個人奄奄一息.
孫大信堅持要救這個人.
但他的同伴認為.
這樣會令到他們三個人.
都一起葬身雪海.
結果那個同伴選擇了自己離開.
孫大信唯有自己.
揹著這個傷者下山.
但奇妙的是.
因為他們身體的接觸.
以至他們能夠獲得溫暖.
最終沒有凍死.
但當他們下到山.
到村口的時候.
很可惜他發現了.
他同伴的屍體.
孫大信他體重別人的生命.
結果他救了人.
也救了自己.
他的同伴.
體重自己的生命.
結果他喪掉生命.
捨己的愛就是以別人的益處優先.
十字架所彰顯的.
就是完全捨己的愛.

$^{641}$彭建輝牧師說過.
沒有十字架的福音.
不是福音.
沒有十字架的基督是假基督.
弟兄姊妹.
十字架的愛.
不單單是感性的愛.
也需要我們用理智.
和意志去回應.
既然耶穌基督願意.
救贖我們.
我們又如何回應這份愛呢?.
你又願意如何毀身基督.
追求成長.
做一個真真正正愛神愛人的人呢?.
愛人的人呢?.
\newpage



\section{約翰福音 20:24-29}
\label{sec:R7EPR0TUgHE}
\textbf{2021年4月4日復活節主餐崇拜直播｜陳明泉牧師｜疑與信之間｜約翰福音20\_24-29}
\newline
\newline
連結: \href{https://youtube.com/watch?v=R7EPR0TUgHE}{\texttt{https://youtube.com/watch?v=R7EPR0TUgHE}} ~~~~ 語音日期: 2021-04-05
\newline
\newline
\hyperref[sec:NcnNHWwVDiw]{\small{< < < PREV SERMON < < <}}
~
\hyperref[sec:index_chronic]{\small{[返順時目]}}
~
\hyperref[sec:oKGXX4r5yvQ]{\small{> > > NEXT SERMON > > >}}
\newline
\newline
約翰福音 20:24-29
\newline
\begin{longtable}{cl}
\hline
\hline
章節 & 經文 (和合本修訂版)\\
\hline
20:24 & \begin{tabularx}{0.7\textwidth}{X} 那十二使徒中，有個叫低土馬的多馬，耶穌來的時候，他沒有和他們在一起。 \end{tabularx} \\ \\ \relax
20:25 & \begin{tabularx}{0.7\textwidth}{X} 其他的門徒就對他說：「我們已經看見主了。」多馬卻對他們說：「除非我看見他手上的釘痕，用我的指頭探入那釘痕，用我的手探入他的肋旁，我絕不信。」 \end{tabularx} \\ \\ \relax
20:26 & \begin{tabularx}{0.7\textwidth}{X} 過了八日，門徒又在屋裡，多馬也和他們在一起。門都關了，耶穌來，站在當中，說：「願你們平安！」 \end{tabularx} \\ \\ \relax
20:27 & \begin{tabularx}{0.7\textwidth}{X} 然後他對多馬說：「把你的指頭伸到這裡來，看看我的手；把你的手伸過來，探入我的肋旁。不要疑惑，總要信！」 \end{tabularx} \\ \\ \relax
20:28 & \begin{tabularx}{0.7\textwidth}{X} 多馬回答，對他說：「我的主！我的神！」 \end{tabularx} \\ \\ \relax
20:29 & \begin{tabularx}{0.7\textwidth}{X} 耶穌對他說：「你因為看見了我才信嗎？那沒有看見卻信的有福了。」 \end{tabularx} \\ \\
[1ex]
\hline
\hline
\end{longtable}
$^{1}$讓我們繼續以聆聽主的道去敬拜我們的神.
今日所要宣講的經文在約翰福音20章24至29節.
那十二個門徒中有稱為底土馬的多馬.
耶穌來的時候 他沒有和他們同在.
那些門徒就對他說:我們已經看見主了.
多馬卻說:我非看見他手上的釘痕.
用指頭探入那釘痕 又用手探入他的肋旁.
我總不信.
過了八天 門徒又在屋裡.
多馬也和他們同在 門都關了.
耶穌來站在當中說:願你們平安.
就對多馬說:伸過你的指頭來 摸我的手.
伸出你的手來 探入我的肋旁.
不要疑惑 總要信.
多馬說:我的主 我的神.
耶穌對他說:你因看見了 我才信.
那沒有看見 就信的有福了.
今天由陳明全牧師為我們正道.
宣講的講題是「疑與信之間」.
各位弟兄姊妹平安.
三分一人的歌聲都是很震撼的事.
和你過去幾個月在家裡敬拜.
看著螢幕也不會唱.
我們只是聽 今天能夠在現場.
一起將歌聲獻給上帝.
我相信上帝是閱立我們獻上這麼美好的歌聲.
但除了現場的感覺好之外.
我們更加知道我們唱的是一個很重要的訊息.
就是關於我們生命復活的訊息.
基督耶穌在二千年前復活了.
他的生命扭轉整個世界.
扭轉我們每一個人生.
因著他的復活 我們從死刃的幽谷.
從一個被罪捆綁 被罪籠罩的生命.
徹底地改喚.
當我們稱耶穌基督為我們生命救主的時候.
我們就成為一個復活的人.
永生就在我們信主的時間就開始了.
所以我們的歌聲是要歌頌那位復活主.
歌頌那位讓我們生命得救的主.

$^{41}$不過在我們唱頌一個復活的.
一個很重要的訊息的時候.
我們必須要問 復活究竟是一個什麼訊息呢?.
其實當你仔細看聖經的時候.
特別是在新約裡面.
復活這個教義其實是貫穿整個新約的.
在保羅所寫的書卷裡面.
如果沒有復活的訊息.
正如保羅自己在哥林多前書十五章裡面說.
我們所做的一切都是枉然的.
保羅用他的生命 用這麼多的力量.
用他整個人生 他要宣講的.
不是一個道基督教化.
而是一個很偉大的人物耶穌來到這個世界裡面.
單單為我們死.
保羅更加要用盡他的力量來講.
這個這麼偉大的夫子.
這個神的兒子為我們死在十架上面.
更加為我們每一個人復活.
他戰勝死亡.
他將黑暗的權命打低.
每一個跟隨他的人.
我們生命都可以得著永恆的生命.
復活是聖經裡面一個不可或缺的.
一個很重要的教義.
復活是我們生命裡面.
不可或缺的一個很重要的訊息.
當沒有了復活這個教義的時候.
或者沒有了復活這個訊息的時候.
我們會怎樣呢.
保羅說我們生命就好像.
一個不能夠跨越生死的悲劇.
我們的人生可能比那些未信主的人.
更加沒有展望.
保羅更加說.
如果基督沒有復活.
耶穌只是一個悲劇英雄.
他做不到我們的救主.
他也帶不到我們走出黑暗的權勢.
但是基督的復活.

$^{81}$如果我們確切的來到去經歷.
確切的相信.
這個訊息是會改變我們的.
講這一堆的話我們都一定認同的.
否則你都不會今天來到教會當中.
一起聆聽這一篇道.
對於復活的教義我們建基在一個.
什麼的根基上呢.
曾經有一些批評基督教.
或者有一些哲學家.
他說基督徒是一班什麼人呢.
其實只是滿有一班宗教情感的一班人.
當他不停強調復活.
可能就是將一個很熱血的訊息.
不斷又不斷的來到去講.
以致到將人那種防衛機制拆低.
所以大家就帶著一種強烈的情感.
去篤信復活的這件事.
在他們來說.
他們當然不深信有復活.
不會相信有耶穌基督.
死裡復活的這個現實.
不過在聖經裡面.
他卻提醒我們.
我們對於復活這個很重要的訊息.
這個教義.
不是建基於一種宗教純粹的那種感覺.
或者一種情感.
也不是一種環境形象出來的結果.
保羅在宣講的時候.
特別強調用一個客觀的現實.
有五百人見到耶穌的復活.
顯然被門徒看到.
顯然被婦女們知道.
在歷史裡面有記載著這件事實.
保羅想將我們從宗教的情感拉下來.
去到我們的理性上.
我們所信的是一個.
經得起現實考驗的信仰.
而約翰福音或者約翰.

$^{121}$他就用另一個角度告訴我們.
我們今天所信的是一個摸得到.
經歷得到的上帝.
今天這一段關於多瑪的記載.
是一段很寶貴的經文.
傳統以來我們有時候很批判多瑪.
今天我們不單止為多瑪說一句話.
更加從多瑪經歷裡面.
拉我們看回基督耶穌確切地復活.
去將祂自己顯現給他的門徒.
以及讓我們知道.
今天我們仔細讀這段聖經.
希望除了我們心裡面.
綻放出那種宗教的情操之外.
更加讓我們的信心有根有機.
我們一起祈禱好嗎.
求上帝的靈幫助我們.
讓我們明白這段華語.
同心祈禱.
基督耶穌願你在我們當中裡面.
幫助我們.
讓我們能夠明白你的華語.
讓我們能夠聽到祂的讀音經文.
重新讓我們知道.
基督耶穌是一個怎樣的神.
幫助我們眾弟兄姊妹.
讓我們能夠聆聽的時候.
我們聽得清楚.
幫助孫講的.
讓經文講得清楚明白.
讓人領受的時候.
我們有信心走面前的路.
我們恭敬將來的時間交託給你.
願主你賜福.
牽手我們祈禱.
奉靠基督耶穌得勝的聖名祈求.
在約翰福音第20章24-25節.
在這裡特別強調門徒.
耶穌剩下的只有11個門徒.
其中一個叫低土馬.

$^{161}$低土馬的意思其實是花名.
孖生仔的意思.
沒有什麼特別內容.
不過這個多馬和其他門徒的經歷有些不同.
當耶穌基督復活了.
顯現給門徒知道的時候.
多馬就不在他們當中.
在聖經裡面提到.
一群門徒見到耶穌的復活.
顯現給他們知道.
他們就興高采烈.
他們的心裡面就一定帶著極大的開心.
迎接救主耶穌.
不過多馬不在他們當中.
聖經裡面提到.
他沒有與他們同在.
然後那些門徒即使怎樣說.
已經看見主了.
多馬就帶著一種很冷靜.
或者跟他們有一點距離感.
多馬就說.
我非看見他手上的釘痕.
用指頭探入那釘痕.
又用手探入他的肋旁.
我總不信.
在這裡面多馬所說的話.
似乎對於那些興高采烈的門徒.
好像淋了一個大冷水一樣.
他們經歷過十字架的痛哭.
那種極大的悲傷.
看到耶穌在十字架上受死的掙扎.
然後到了第三天.
經歷耶穌的復活.
內心心情就像過山車一樣.
不過當他們坐過山車的時候.
好像沒有參與.
耶穌復活的多馬.
說的這句話就是.
我不信的.
除非我看見耶穌手上的釘痕.

$^{201}$用指頭探入那釘痕.
用手探入他的肋旁.
這三件事.
他的釘痕,肋旁.
用指頭探入身體的傷患.
這是受苦的印記.
釘十字架裡面一個很重要的記號.
如果多馬說他看不見這些.
他是不會相信的.
你會覺得那種冷水.
有時候覺得這種反論調.
都很磨人的.
有時候我們甚至覺得.
對於一個真理.
或者一個我們覺得無所謂.
我們已經順理成章的一個真理.
有人這樣面對挑戰.
是很納悶,不善未疑的.
不過多馬真的很不客氣地說.
他要看見,他要用手指碰過.
他要探入去,他才相信.
在他的邏輯裡面.
不是隨便一件事出現就等於他相信.
是要A,B,C三件事.
都要同時間出現.
看見耶穌基督受傷的記號的時候.
他才相信.
否則的話,他不相信.
這是一個多馬.
某程度上是一個很冷靜的分析.
一個邏輯的思維.
多馬在這一刻裡面都還未相信.
不過多馬的不相信.
和之前的蝙蝠說到.
耶穌面對著一班法利賽人的不相信.
有什麼不同呢?.
法利賽人如果我們看聖經就知道.
從頭到尾,他們都是一班拒絕相信的人.
在聖經裡面,馬可福音特別有一個記載.
耶穌用七個餅和幾條魚.

$^{241}$在五餅二魚之後,再行多個神蹟.
就是七個餅和幾條魚,再餵飽四千人.
一班法利賽人看著整個神蹟.
耶穌基督所做的這件事的進行.
然後就餵飽了這幾千人.
其實是一個很偉大的神蹟.
在講完這個神蹟之後.
或者在耶穌做完這個神蹟之後.
一班法利賽人再去挑戰耶穌.
走去耶穌面前說.
夫子,如果你是神的話,請你做一個神蹟出來.
在聖經裡面特別加一個主腳.
這班法利賽人去問耶穌這個問題.
不是想叫耶穌證明他是神的兒子.
而是想找一些把柄出來挑戰耶穌.
在雞蛋裡面挑骨頭,挑出一些錯誤.
證明耶穌不是神的兒子,不是來自神的拉比.
法利賽人就是用這樣的心態去做.
然後耶穌怎樣回應這班法利賽人呢?.
他說,在一個不信的世代裡面.
做多少神蹟都沒有用.
所以,在聖經裡面,我們可以看到.
在一個不信的世代裡面.
做多少神蹟都沒有用.
你不信就不信.
即使看完有幾個餅,幾條魚,餵飽四千人的神蹟都不相信.
再去做更多的事,他們都是拒絕相信.
這班人就是一班完全拒絕耶穌.
不論耶穌做多少神蹟奇事.
做多少記號在人沒辦法實踐的作為當中.
這班人都是帶著他們自己的偏頗.
帶著他們心靈的那種硬心.
拒絕基督耶穌.
就是生命裡面來自神的拉比,來自神的兒子.
和那班法利賽人在多瑪的角度是有點不同的.
多瑪就是一直跟著耶穌的其中一個門徒.
他很清楚,他很凝信耶穌基督所行.
他所言所行,他生命裡面所作的教導.
全部一切都是來自上帝的.
不過在復活這件事裡面,怎樣令到他去確信呢.

$^{281}$他心裡面,如果從他的情感或是從他的角度去想.
剛剛才經歷完,要接受耶穌的死也是一個很不容易的階段.
就是任何你所愛的人,他經歷死亡.
在你的心目中,你都要有一段時間來消化.
你都要接受這個現實.
在這個現實還未接受得夠的時候.
第三天,就要接受一個死了的人.
在十字架裡面接受酷刑的一個基督耶穌.
竟然重新復活起來.
實在太過不可思議了.
更何況是那班門徒自己說的.
雖然那班是他的兄弟,是他的老友.
是一起經歷過的.
不過整件事實在匪夷所思.
哪有這個世界裡面曾經出現過這樣的事.
從來都未發生過.
所以是一個難以置信的一件事.
所以在多瑪來說.
他的不信不是建基於他好像法利賽人那種的心眼.
而是一個史無前例.
在過去裡面都沒有可以參考的一件事.
竟然淋到他第一身上.
所以他就作出一個好像條件式的表達.
如果他見不到,他不會驗證過.
他不會相信基督耶穌復活的事實.
傳統以來去讀這段經文.
甚至乎有些,我們有時查經也好.
多瑪寫著不信,多瑪總不信這個字.
我們很快就將他對號入座.
甚至乎帶著一種批判的角度.
來看多瑪這種表達.
我們甚至乎會覺得.
懷疑本身是一個信仰的毒素.
我們覺得如果有些問題.
或者按著聖經這樣說.
如果提出一種質疑,是不是這樣說的.
他未必是基於一種挑戰的角度.
基於一種尋求真理,尋真的精神.
但是有些聖經學家.
甚至乎在我們華人社會裡面.

$^{321}$對於這些說法我們會覺得很納悶.
甚至乎覺得不止淋冷水.
還會覺得有時搞事.
我不知道你們會不會是這樣.
在我自己年輕一點的時候.
在我未來旺朕,在我自己的母會查經.
暗信主,所有事情都是很新奇.
有些查經的問題,導師這樣說.
是不是真的這樣.
我是很坦誠地問,是不是真的這樣.
有時會覺得,是不是Anders.
每件事都遮住了.
或者覺得我好像是搞事分子.
但我心裡不是要挑戰權威.
或者純粹覺得那件事是不可信.
而是想帶著一份真誠,尋真的精神.
究竟那段經文是說什麼.
滿足於一個純粹別人說出來的二手答案.
事實上在聖經裡面對於多瑪的懷疑.
稍後我會再說.
不過在這裡特別強調.
約翰福音或者其他聖經裡面.
不是用負面的角度來看多瑪這種驗證.
有時可能這種懷疑,這種驗證.
會令身邊的人覺得很冷冷水,掃興.
或者覺得好像有些搞事.
將一些我們已經認信的真理.
從頭反過來再去問.
我們會覺得有點不耐煩.
不過在聖經裡面.
從來都不是將多瑪成為一個批判的對象.
反而在這裡上帝樂意將一個真理.
給多瑪來知道.
其實弟兄姊妹你們知不知道.
今天在信仰裡面.
最危險的不是一個所有東西照單傳修的信仰.
或者那些人問問題.
我們覺得不照單傳修的時候.
那些人是危險人物.
在團契裡面要找一個強勁的導師來應付他.

$^{361}$那些人不是最危險的人物.
最危險的人物反而是我們不求甚解.
我們有大量的熱誠.
但我們其實不知道自己信什麼.
我們不知道真正的基督耶穌.
他所啟示給我們的是一個什麼信息.
我們只是帶著一腔的熱誠.
純粹來過我們的宗教生活.
但我們其實心底裡面.
對於耶穌復活是一回什麼事.
不是很問的.
不是很想知道.
我們知道大家都是這樣做.
我們就跟著去吧.
曾經我把這段經文和我女兒分享.
我問她女兒你讀這段經文有沒有什麼問題.
她就不出聲.
她就悶悶不樂.
她說不要緊你有問題就問吧.
我女兒有種感覺說.
爹地我怕問問題你覺得我不孝.
覺得我挑戰她的權威.
不是挑戰我的權威.
她說很沒禮貌.
不懂得問問題有多沒禮貌.
她說其實有時真的.
對於聖經的事.
這樣說就要這樣信.
我說你可以問問題.
她說問多一句不信.
就是犯罪.
她去到一個這樣的極致.
我跟她說不是.
我說看聖經如果不問問題.
那個才是最危險.
我們才落入了一種.
只有宗教情感的生活.
但信仰不清楚的時候.
我們內裡空空如也.
外面有很多活力.

$^{401}$那個信仰才是最危險.
甚至很多異端都是這樣.
很多異端就是斷章取義.
不明白聖經上下文所說的話.
然後就用他的熱誠來經營.
他們自己認為對的信仰.
有時頂指會容許我說.
有時我聽到查經.
我們會寫一些資料.
查經的時候有些組員問問題.
問到最後組長答不到.
通常會出一張黃牌.
知不知道會出什麼黃牌.
不是聖經這樣說.
他出一張紅心D.
那張紅心D是什麼呢.
陳牧師是這樣說的.
每次都是這樣.
解釋不到的時候.
牧師是這樣說的.
但有時聽到.
其實我不是這樣說的.
我都不是這樣解釋的.
不過大家想出一張紅心D.
不出一張大D就是聖經這樣說.
就想著陳牧師說的.
就搞定了.
有時我都聽到.
我都沒說過這些.
但奉為金科玉律.
今天.
教會最大的問題.
不是我們失去了宗教的熱誠.
我們失去了那份.
對信仰那份執著.
我們真的要知道尋真的那種精神.
陳牧師這樣說不一定是真的.
當然我不是拆自己的招牌.
我盡力告訴你.
一個真實的聖經.

$^{441}$不過真正導引我們.
為我們生命提供答案的.
不是陳牧師怎樣說.
而是聖經怎樣說.
而是真正有權威的.
放在那本.
有66卷聖經.
啟示而成.
成典的聖經裡面怎樣說.
我們今天.
有時帶著健康的懷疑.
有時帶著虛心.
帶著真誠來發問.
那個反而是對我們的生命信仰.
帶來一個真正的養分.
令我們的信仰.
可以變得有根有機.
不要不求甚解.
多瑪在這裡.
為門徒帶來一些納悶.
淋冷水或麻煩.
不過多瑪在這裡.
他比其他門徒有些不同.
他要親身經歷.
他要問問題.
他要用一個驗證的方式.
來看見復活的基督耶穌.
在第26至27節.
既然他這樣問.
在八天之後.
耶穌就回覆他.
過了八天門徒又在屋裡.
多瑪也和他們同在.
門都關了.
耶穌來站在他們當中.
願你們平安.
就對多瑪說.
伸過你的指頭來.
摸我的手.
伸出你的手來.

$^{481}$探入我的肋旁.
不要疑惑.
總要信.
在第26至27節裡.
耶穌回答了多瑪的條件式的疑問.
到了第八天.
特別是耶穌這次的顯現.
在聖經裡面.
特別是朝著多瑪的問題而去.
所以耶穌特別向多瑪作出一個邀請.
來吧,你看看.
用你的手指來.
摸摸我釘痕的手.
伸出你的手來.
摸摸他的肋旁.
有一個名畫.
影著多瑪.
樣子好像很疑惑.
伸手入去耶穌的肋旁.
摸摸那個洞洞.
摸到是一個真實的.
復活的耶穌.
摸到那個釘痕.
摸到那個.
曾經被釘刺透的身軀.
一個在十字架上受苦的印記.
本來是一個痛苦的經歷.
但在多瑪眼前.
就呈現了一個救贖的記號.
在多瑪伸出手去驗證的時候.
或者耶穌作出一個這樣的邀請的時候.
耶穌要解決他的疑惑.
讓他真誠地相信.
基督耶穌就是復活的救主.
在這裡.
一個很不明顯.
不過是一個很重要的現實.
人每一件事都帶著問題.
去到上帝面前.
如果剛才我所說的.

$^{521}$我鼓勵大家去聖經.
或者在信仰裡面尋求真的時候.
我們的尋真精神.
其實不足以純粹讓我們知道上帝是誰.
或者復活的現實.
在這裡.
是上帝主動.
藉著他的愛子耶穌.
將他自己向人啟示自己.
在神學裡面說的啟示的意思.
不是尋人啟示的事情.
而是由上而下表達的意思.
是在說上帝從天而落地.
將他自己是誰.
他的計劃是怎樣.
上帝採取一種主動的方式.
將他的救贖計劃.
告訴普世所有人知道.
是上帝自己作主動.
讓人知道上帝是誰.
上帝的救贖計劃是怎樣.
所以即使多瑪經這麼多的疑問.
如果耶穌不解答.
他都會永遠帶著這個疑問.
繼續過他的人生.
不過上帝不是這樣的神.
不是喜歡跟人捉迷藏的神.
上帝是將他自己告訴世界知道.
告訴這個徒弟多瑪.
讓他看到這個就是復活的主.
解開了疑惑之後.
他就要信靠.
不要疑惑 總要信.
這個信仰是上帝首先.
將他自己啟示出來.
將一個真實 將一個真理.
活到一個地上裡面.
讓人知道耶穌基督真真正正活過在歷史當中.
所以你看約翰福音.
由第一章到最後.

$^{561}$甚至約翰一書.
約翰都是用這個角度去寫.
道成了肉身住在我們中間.
充充滿滿地有恩典有真理.
我們也見過他的榮光.
正是父獨生子的榮光.
約翰一書一章而至.
論道者起初原有生命的道.
就是我們所聽見 所看見.
親眼看過 親手摸過的.
這生命已經顯現出來.
我們也看見過 現在又作見證.
將原與附同在.
且顯現與我們那永遠的生命.
傳給你們.
那麼急急急急說什麼呢.
約翰福音說.
耶穌不是一個道讚出來的神話故事.
真理不是一個冰冰冷冷的概念.
一個想法.
一個大家覺得製造一個英雄人物出來.
真理活然在這個世界裡面.
真理成為肉身.
這個肉身與人相近.
這個肉身還要與這班門徒生活過.
這班門徒摸過.
實在的 有溫度的.
那個釘痕的手 那個隔了底的洞.
是真實的.
這不是一個悲劇的英雄.
一班人因為很悲傷之下.
幻想出來 復活的這件事.
約翰告訴我們.
我們所信的主是在這個世代裡面.
在他的時空裡面.
介入在他們當中.
與他們一起生活過.
更加重要的是 神竟然死了.
一個永活的神竟然經歷死亡.
一個荒誕絕倫的說法.

$^{601}$但就是一個正正的方式.
在人看為愚蠢.
但在上帝眼中.
就是要叫所有自以為是的人可以跌倒.
用一個真實的信仰.
來呈現甚麼叫做真實和虛假的世界.
今天我們可能很多時活在很多大話當中.
甚至我們有時信以為真.
有一段時間我去到大學進修.
讀輔導.
除了讀輔導之外.
其中有幾科是講後現代的思潮.
其中一個講法叫做social constructionism.
社會建構主義.
可能現在大學生都會讀過這件事.
這個講論的激戰是講甚麼呢.
就是這個世界沒有真理.
所謂人所講的真理是甚麼呢.
是一班人在一個社交圈子裡.
大家信以為真的事.
大家建構出來的.
大家講著講著就成真.
當這件事成真的時候.
大家就確信.
就像滾雪球一樣.
將這件事成為一個永恆的真理.
這班人為了令自己的心靈.
有那種穩妥感.
所以會極力用更多的事.
圍繞著這個倒傘出來的真相.
social constructionism 他說.
這個世界哪有真相.
沒有的 所以你不用求真.
尋求真實感就行了.
人們感覺到真實就行了.
任何一個神話.
講到大家聲淚俱下.
大家有種情感投入.
有真實感就行了.
不過我們的信仰拒絕這種說法.

$^{641}$要有人提醒我們.
我們今天所信的.
甚至曾經自己門外都質疑的事.
他記錄下來.
我們今天不是信任一個倒傘出來的神話故事.
上帝活在這個歷史空間裡.
上帝倒乘肉身.
藉著基督耶穌將這個真理呈現出來.
我們所信靠的是一個進入歷史的真理.
我們所經歷的不是一個謊言.
更加重要的是.
這個謊言為何我們可以這麼確切呢.
因為基督耶穌將自己啟示出來.
如果他不經歷過這些神蹟.
我們沒有人可以相信.
可能大家早一代會看過一本書.
一本書叫做鐵證代判.
不知道大家有沒有看過.
鐵證代判這本書是甚麼呢.
今天不看了.
這麼厚的枕頭書.
這本書的名字最適合我們戰鬥.
很多人都以為是枕頭書.
但這本鐵證代判的書是有意思的.
作者原本是想做甚麼呢.
就是想砌基督教.
這個世界哪有這樣的信仰.
他覺得那些基督徒是很愚蠢的一班人.
所以他就想做一些理據.
來砌基督徒的論據.
覺得不成立的,不合邏輯的等等.
他就想用這些歷史真相來砌基督教的信仰.
不過砌砌砌的時候.
怎知道又有新的發現.
又有不同的理據出現.
砌砌砌,這本鐵證代判.
本來由砌基督徒的一本書.
變成一本護教的書籍.
砌砌砌,連這個人自己都.
最後的結論就是.

$^{681}$原本我是要攻擊基督教.
但到了今天,當我找到這麼多的證據的時候.
我不得不相信基督耶穌就是神的兒子.
這個神的兒子在歷史裡面進行拯救.
這個基督教信仰.
不單止是經得起考驗.
而且是我們生命裡唯一可以相信的真理.
今天我們除了我們在復活.
在一個令我們很澎湃的信息之外.
我們更加要知道.
我們所信的不是製造一個真實感.
很澎湃感受的信仰.
而是一個真真實實在歷史裡面出現的信仰.
20章第28至29節.
在這個段落裡面最後一個小節.
當馬雪探完耶穌的肋旁和手的時候.
他就說,我的主,我的神.
耶穌對他說,你因看見我才信.
那沒有看見就信的就有福了.
在這裡面,雖然有幾個字.
我的主,我的神.
我們平時自己祈禱都會這樣說.
不過在聖經裡面.
特別在門徒當中.
最近傍講這番說話的是彼得.
你是神的兒子,你是基督.
你是神的兒子.
最近傍的了,都沒有這句這麼高的規格.
直稱耶穌為神為主.
大概在新約裡面是多瑪這裡的表達.
這句話帶著的不單止是一種驗證.
證明你是甚麼,冷冰冰.
這句話的語氣帶著強烈的情感.
在經歷過尋求真相之後.
探過那個洞.
看過那個釘痕之後,摸過了.
見到眼前這個耶穌.
這個是我的主,我的神.
這一種感嘆,這一種的形訊.
帶著強烈的那種應訊的情感.

$^{721}$多瑪在這裡面.
將他最重要的應訊表達.
在這裡面記錄下來.
從此之後.
多瑪就成為一個真正驗證基督耶穌.
去到各城各鄉.
繼續延續著門徒心智的.
一個很重要的門徒.
他的應訊代表著他找到真相.
他生命裡面找到那位真正復活的主.
耶穌怎樣去評論多瑪整個過程呢?.
這句不是特別對多瑪說的.
其實是對我們說的.
耶穌對他說.
我才信,那沒有看見就信,就有福了.
在當代裡面,多瑪親身的驗證.
得出的就是一個我的主,我的神.
耶穌基督就是我的主,我的神的結論.
不過今天我們沒有辦法.
再親手摸摸耶穌.
我們未必會再.
耶穌顯現給我們看釘痕.
如果你看到這件事你都會害怕.
不過今天多瑪.
他的經驗,他的驗證.
成為了我們今天相信信息的一部份.
在多瑪當代沒有聖經這樣記載.
他就第一身.
但是到了今天我們就直接讀這段聖經.
我們就從這段經文裡面.
我們就知道基督耶穌是一個真正的主.
哥林多後書五章十六節.
保羅都是這樣告訴我們.
所以我們從今以後.
雖然憑著外貌認過基督.
如今卻不再這樣認他了.
我們今天認順基督.
我們認到他.
就不是純粹憑肉體外貌認基督.
是憑著聖靈在我們內心裡面幫助我們.

$^{761}$讓我們經歷到聖靈的同罪.
神的經歷.
以致我們明白這個信仰.
今天當我們認順的時候.
我們是否經歷上帝呢?.
我們是否純粹回到.
二千年前多瑪的模式.
還是我們知道我們今天已經按著聖經.
憑著聖靈的感動.
我們便認到基督耶穌是一個怎樣的神.
有些牧者更加極致.
有這樣的經歷.
在前年的時候.
如果大家記得我們曾經有一班.
執事童工去過台灣福氣教會.
去學習幸福小組.
我記得很深刻.
那個牧師跟我說了一件事.
他說每一個基督徒.
因為在台灣有很多民間信仰.
他特別說了一件事.
怎樣幫助弟兄姊妹經歷到信仰呢?.
他說每一次都會帶剛開始信主.
去到那些最迷信.
最多別異宗派.
或者被鬼附的那些人.
帶那些人去經歷一下.
怎樣經歷呢?.
帶他們去趕鬼.
我聽到有些無骨悚然.
為什麼要做這些事?.
你不怕嚇到他們嗎?.
真正信主的人怕什麼?.
有罪的人才會害怕.
接著我說為什麼要做這件事?.
那個牧師說.
我們今天有很多基督徒.
朗朗上口.
奉基督耶穌的名.
講一下耶穌的名.

$^{801}$說在他嘴邊.
但基督耶穌的名.
帶著權能他們沒有經歷過.
他說真正讓他們經歷到這件事.
他就帶他們去做那些趕鬼.
奉基督耶穌的名.
將那些污鬼趕走.
那些人匿匿在目.
看著整件事發生.
不是那些人很厲害.
不是那些人很聰明.
只是單單憑著認信.
憑著基督耶穌的名.
將這些污鬼趕走.
讓他們經歷到信仰的真實.
他們就是這樣經歷信仰.
兄弟姐妹.
當然我們未必一定要用這種方式.
不過基督耶穌的復活的名.
基督耶穌的名.
今天不要成為我們純粹.
掛在口邊嘴巴舔舔.
隨口說出來的東西.
有一份認信.
祂是我們的主.
我們的神.
我們生命裡面.
要真實地經歷這個信仰.
真實體會這個信仰.
是一個怎樣的復活.
大能的信仰.
在你生命當中.
你可以經歷得到的.
只要你放開你的內心.
不要被純粹單憑一個約定俗成的答案.
滿足了它.
你就放棄了繼續尋真的精神.
今日這段經文提醒我們.
當晚的尋真.
他尋求一個真理.

$^{841}$上帝會向他來對啟示.
是上帝採取一個主動.
向人告訴他.
他是一個怎樣的神.
他更加告訴我們.
今日我們不是信一個虛擬.
或者道讚出來的一個神話故事.
而是一個親手摸過.
親身經歷過.
一起生活過的基督耶穌.
在一班門徒當中.
記錄下他這些生命的事跡.
讓每一個看聖經的人.
可以因此而信基督耶穌.
就是我們生命的主.
我們今天.
除了我們帶著激昂.
來歌頌我們的上帝.
我們更加要.
繼續尋求真相.
尋求那位真實為我們.
釘身十架的主耶穌.
尋求那位帶領我們走過生命.
一切憂谷的上帝.
當我們尋求他.
他會向我們顯現.
多瑪色日講.
我的主 我的神.
今天我們都一起凝信.
耶穌基督.
你是我的主.
我的神.
\newpage



\section{馬可福音 8:27-9:37}
\label{sec:oKGXX4r5yvQ}
\textbf{2021年4月18日崇拜直播｜陳妙濃傳道｜就當捨己｜馬可福音8\_27-9\_37}
\newline
\newline
連結: \href{https://youtube.com/watch?v=oKGXX4r5yvQ}{\texttt{https://youtube.com/watch?v=oKGXX4r5yvQ}} ~~~~ 語音日期: 2021-04-18
\newline
\newline
\hyperref[sec:R7EPR0TUgHE]{\small{< < < PREV SERMON < < <}}
~
\hyperref[sec:index_chronic]{\small{[返順時目]}}
~
\hyperref[sec:opD_mMvcV8I]{\small{> > > NEXT SERMON > > >}}
\newline
\newline
馬可福音 8:27-9:37
\newline
\begin{longtable}{cl}
\hline
\hline
章節 & 經文 (和合本修訂版)\\
\hline
8:27 & \begin{tabularx}{0.7\textwidth}{X} 耶穌和門徒出去，往凱撒利亞‧腓立比附近的村莊去。在路上，他問門徒：「人們說我是誰？」 \end{tabularx} \\ \\ \relax
8:28 & \begin{tabularx}{0.7\textwidth}{X} 他們對他說：「是施洗的約翰；有人說是以利亞；又有人說是先知中的一位。」 \end{tabularx} \\ \\ \relax
8:29 & \begin{tabularx}{0.7\textwidth}{X} 他又問他們：「你們說我是誰？」彼得回答他：「你是基督。」 \end{tabularx} \\ \\ \relax
8:30 & \begin{tabularx}{0.7\textwidth}{X} 於是耶穌切切地囑咐他們不可對任何人說起他。 \end{tabularx} \\ \\ \relax
8:31 & \begin{tabularx}{0.7\textwidth}{X} 從此，他教導他們說：「人子必須受許多的苦，被長老、祭司長和文士棄絕，並且被殺，三天後復活。」 \end{tabularx} \\ \\ \relax
8:32 & \begin{tabularx}{0.7\textwidth}{X} 耶穌明白地說了這話，彼得就拉著他，責備他。 \end{tabularx} \\ \\ \relax
8:33 & \begin{tabularx}{0.7\textwidth}{X} 耶穌轉過來看著門徒，斥責彼得說：「撒但，退到我後邊去！因為你不體會神的心意，而是體會人的意思。」 \end{tabularx} \\ \\ \relax
8:34 & \begin{tabularx}{0.7\textwidth}{X} 於是他叫眾人和門徒來，對他們說：「若有人要跟從我，就當捨己，背起自己的十字架來跟從我。 \end{tabularx} \\ \\ \relax
8:35 & \begin{tabularx}{0.7\textwidth}{X} 因為凡要救自己生命的，必喪失生命；凡為我和福音喪失生命的，必救自己的生命。 \end{tabularx} \\ \\ \relax
8:36 & \begin{tabularx}{0.7\textwidth}{X} 人就是賺得全世界，賠上自己的生命，有甚麼益處呢？ \end{tabularx} \\ \\ \relax
8:37 & \begin{tabularx}{0.7\textwidth}{X} 人還能拿甚麼換生命呢？ \end{tabularx} \\ \\ \relax
8:38 & \begin{tabularx}{0.7\textwidth}{X} 凡在這淫亂罪惡的世代，把我和我的道當作可恥的，人子在他父的榮耀裡與聖天使一同來臨的時候，也要把那人當作可恥的。」 \end{tabularx} \\ \\ \relax
9:1 & \begin{tabularx}{0.7\textwidth}{X} 耶穌又對他們說：「我實在告訴你們，站在這裡的，有人在沒經歷死亡以前，必定看見神的國帶著能力臨到。」 \end{tabularx} \\ \\ \relax
9:2 & \begin{tabularx}{0.7\textwidth}{X} 過了六天，耶穌帶著彼得、雅各、約翰，領他們悄悄地上了高山。他在他們面前變了形像， \end{tabularx} \\ \\ \relax
9:3 & \begin{tabularx}{0.7\textwidth}{X} 衣服放光，極其潔白，地上漂布的人沒有一個能漂得那樣白。 \end{tabularx} \\ \\ \relax
9:4 & \begin{tabularx}{0.7\textwidth}{X} 有以利亞和摩西向他們顯現，並且與耶穌說話。 \end{tabularx} \\ \\ \relax
9:5 & \begin{tabularx}{0.7\textwidth}{X} 彼得對耶穌說：「拉比，我們在這裡真好！我們來搭三座棚，一座為你，一座為摩西，一座為以利亞。」 \end{tabularx} \\ \\ \relax
9:6 & \begin{tabularx}{0.7\textwidth}{X} 彼得不知道說甚麼才好，因為他們很害怕。 \end{tabularx} \\ \\ \relax
9:7 & \begin{tabularx}{0.7\textwidth}{X} 有一朵雲彩來遮蓋他們，又有聲音從雲彩裡出來，說：「這是我的愛子，你們要聽從他！」 \end{tabularx} \\ \\ \relax
9:8 & \begin{tabularx}{0.7\textwidth}{X} 門徒連忙向周圍觀看，不再看見任何人，只見耶穌同他們在一起。 \end{tabularx} \\ \\ \relax
9:9 & \begin{tabularx}{0.7\textwidth}{X} 下山的時候，耶穌囑咐他們說：「人子還沒有從死人中復活，你們不要把所看到的告訴人。」 \end{tabularx} \\ \\ \relax
9:10 & \begin{tabularx}{0.7\textwidth}{X} 門徒將這話存記在心，彼此議論「從死人中復活」是甚麼意思。 \end{tabularx} \\ \\ \relax
9:11 & \begin{tabularx}{0.7\textwidth}{X} 他們就問耶穌：「文士為甚麼說以利亞必須先來？」 \end{tabularx} \\ \\ \relax
9:12 & \begin{tabularx}{0.7\textwidth}{X} 耶穌說：「以利亞的確先來復興萬事。經上不是指著人子說，他要受許多的苦和被人輕慢嗎？ \end{tabularx} \\ \\ \relax
9:13 & \begin{tabularx}{0.7\textwidth}{X} 我告訴你們，以利亞已經來了，他們任意待他，正如經上指著他說的。」 \end{tabularx} \\ \\ \relax
9:14 & \begin{tabularx}{0.7\textwidth}{X} 他們到了門徒那裡，看見有一大群人圍著他們，又有文士和他們辯論。 \end{tabularx} \\ \\ \relax
9:15 & \begin{tabularx}{0.7\textwidth}{X} 眾人一見耶穌，都很驚奇，就跑上去向他問安。 \end{tabularx} \\ \\ \relax
9:16 & \begin{tabularx}{0.7\textwidth}{X} 耶穌問他們：「你們和他們辯論甚麼？」 \end{tabularx} \\ \\ \relax
9:17 & \begin{tabularx}{0.7\textwidth}{X} 眾人中的一個回答：「老師，我帶了我的兒子到你這裡來，他被啞巴的靈附著。 \end{tabularx} \\ \\ \relax
9:18 & \begin{tabularx}{0.7\textwidth}{X} 無論在哪裡，那靈拿住他，把他摔倒，他就口吐白沫，牙關緊鎖，身體僵硬。我請過你的門徒把那靈趕出去，他們卻不能。」 \end{tabularx} \\ \\ \relax
9:19 & \begin{tabularx}{0.7\textwidth}{X} 耶穌回答：「唉！這不信的世代啊，我和你們在一起要到幾時呢？我忍耐你們要到幾時呢？把他帶到我這裡！」 \end{tabularx} \\ \\ \relax
9:20 & \begin{tabularx}{0.7\textwidth}{X} 他們就帶了他來。那靈一見耶穌，就使他重重地抽風，倒在地上，翻來覆去，口吐白沫。 \end{tabularx} \\ \\ \relax
9:21 & \begin{tabularx}{0.7\textwidth}{X} 耶穌問他父親：「他得這病有多久了呢？」父親說：「從小的時候。 \end{tabularx} \\ \\ \relax
9:22 & \begin{tabularx}{0.7\textwidth}{X} 那靈屢次把他扔在火裡、水裡，要治死他。你若能做甚麼，求你憐憫我們，幫助我們。」 \end{tabularx} \\ \\ \relax
9:23 & \begin{tabularx}{0.7\textwidth}{X} 耶穌對他說：「『你若能』，在信的人，凡事都能。」 \end{tabularx} \\ \\ \relax
9:24 & \begin{tabularx}{0.7\textwidth}{X} 孩子的父親立刻喊著說：「我信；求你幫助我的不信！」 \end{tabularx} \\ \\ \relax
9:25 & \begin{tabularx}{0.7\textwidth}{X} 耶穌看見眾人都跑上來，就斥責那污靈說：「你這聾啞的靈，我命令你從他裡頭出來，再不要進去！」 \end{tabularx} \\ \\ \relax
9:26 & \begin{tabularx}{0.7\textwidth}{X} 那靈大喊一聲，使孩子猛烈地抽了一陣風，就出來了。孩子好像死了一般，以致眾人多半說：「他死了。」 \end{tabularx} \\ \\ \relax
9:27 & \begin{tabularx}{0.7\textwidth}{X} 但耶穌拉著他的手，扶他起來，他就站起來了。 \end{tabularx} \\ \\ \relax
9:28 & \begin{tabularx}{0.7\textwidth}{X} 耶穌進了屋子，門徒就私下問他：「我們為甚麼不能趕出那靈呢？」 \end{tabularx} \\ \\ \relax
9:29 & \begin{tabularx}{0.7\textwidth}{X} 耶穌對他們說：「非用禱告，這一類的邪靈總趕不出來。」 \end{tabularx} \\ \\ \relax
9:30 & \begin{tabularx}{0.7\textwidth}{X} 他們離開那地方，經過加利利；耶穌不願意人知道， \end{tabularx} \\ \\ \relax
9:31 & \begin{tabularx}{0.7\textwidth}{X} 因為他正教導門徒說：「人子將要被交在人手裡，他們要殺害他；被殺以後，三天後他要復活。」 \end{tabularx} \\ \\ \relax
9:32 & \begin{tabularx}{0.7\textwidth}{X} 門徒卻不明白這話，又不敢問他。 \end{tabularx} \\ \\ \relax
9:33 & \begin{tabularx}{0.7\textwidth}{X} 他們來到迦百農。耶穌在屋裡問門徒說：「你們在路上議論的是甚麼？」 \end{tabularx} \\ \\ \relax
9:34 & \begin{tabularx}{0.7\textwidth}{X} 門徒不作聲，因為他們在路上彼此爭論誰最大。 \end{tabularx} \\ \\ \relax
9:35 & \begin{tabularx}{0.7\textwidth}{X} 耶穌坐下，叫十二個使徒來，說：「若有人願意為首，他要作眾人之後，作眾人的用人。」 \end{tabularx} \\ \\ \relax
9:36 & \begin{tabularx}{0.7\textwidth}{X} 於是耶穌領一個小孩過來，讓他站在門徒當中，又抱起他來，對他們說： \end{tabularx} \\ \\ \relax
9:37 & \begin{tabularx}{0.7\textwidth}{X} 「凡為我的名接納一個像這小孩子的，就是接納我；凡接納我的，不是接納我，而是接納那差我來的。」 \end{tabularx} \\ \\
[1ex]
\hline
\hline
\end{longtable}
$^{1}$讓我們一起聆聽一段上帝的話語.
記載在《馬可福音》第八章二十七至三十八節.
及第九章三十至三十七節.
在大家手上的新約聖經第五十九及六十一月.
聖經這裡是這樣說.
耶穌和門徒出去往凱撒利亞肥勒比的村莊去.
在路上問門徒說:人說我是誰?.
他們說:有人說是施浸的約翰.
有人說是以利亞.
又有人說是先知裡的一位.
又問他們說:你們說我是誰?.
彼得回答說:你是基督.
耶穌就禁戒他們不要告訴人.
從此他教訓他們說.
人子必須受許多的苦.
被長老濟斯長和問世氣絕.
並且被殺.
過三天復活.
耶穌明明說者說.
彼得就拉著他勸他.
耶穌轉過來看著門徒就責備彼得.
說:撒旦退我後邊去吧.
因為你不體貼神的意思.
只體貼人的意思.
於是叫眾人和門徒來對他們說.
若有人要跟隨我.
就當捨己背起他的十字架來跟從我.
因為凡要救自己生命的.
必喪掉生命.
凡為我和福音喪掉生命的.
必救了生命.
人若賺得全世界賠上自己的生命.
有甚麼益處呢?.
人還能拿甚麼生命呢?.
凡在這淫亂罪惡的世代.
把我和我的道當作可恥的.
人只在他父的榮耀裡.
和聖天使降臨的時候.
也要把那人當作可恥的.
第九章三十節.

$^{41}$他們離開那地方經過加利利.
耶穌不願意人知道.
於是教訓門徒說.
人子將要被交在人手裡.
他們要殺害他.
被殺以後.
過三天他要復活.
門徒卻不明白這話.
又不敢問他.
他們來到加伯隆.
耶穌在屋裡問門徒說.
你們在路上議論的是甚麼?.
門徒不作聲.
因為他們在路上彼此爭論誰為大.
耶穌坐下.
叫十二個門徒來說.
若有人願意作首先的.
他必作眾人末後的.
作眾人的用人.
於是領過一個小孩子來.
叫他站在門徒中間.
又抱起他來.
對他們說.
凡為我名接待一個將者小孩子的.
則是接待我.
凡接待我的.
不是接待我.
乃是接待那差我來的.
聖經獨筆.
今天早上有陳美龍姑娘.
在我們當中分享剛才這段話語.
講題是「就當寫記」.
(字幕).
各位弟兄姊妹早安.
內心覺得很開心.
因為我嫂嫂原來我有年紀.
沒有這樣跟大家分享訊息.
原來上年4月19日分享訊息.
當然不計其他.
剛剛一年了.

$^{81}$我們一起禱告.
榮耀的救主.
你配得讚美.
你配得萬失跪拜.
萬口承認你是主.
你是王.
我們同心呼求.
求你賜我們力量.
離開我們的罪恨.
我們一心一意.
尋求你的心意.
願你的憐憫恩典.
充滿這裡.
充滿我城.
我們都祈求你開我們的耳朵.
開我們心靈的眼睛.
讓我們聆聽你的話語.
不單是聆聽.
並且淺行出來.
我們同心禱告.
是奉耶穌基督.
聖名祈禱.
阿們.
小野二郎.
小野二郎是誰呢?.
小野二郎被稱為.
日本壽司第一人.
或有人稱他為壽司之神.
他現在已經九十五歲.
1950年.
他開始做壽司的師傅.
而且他也是.
最大年紀米芝蓮的三星廚師.
他的生命都頗傳奇.
不少人都驚訝.
原來小野先生.
他除了做壽司的時候.
他沒有戴手套.
他無時無刻.
他都戴著手套.

$^{121}$例如他外出.
一年四季.
包括將來的夏天.
這麼熱.
他都戴著手套.
又或者他連睡覺.
他都戴著.
為什麼他要這樣做呢?.
原來他就是為了.
保護他這一雙手.
有人問他.
為什麼你要保護你雙手?.
原來小野先生他這樣回答.
原來同樣是壽司.
如果雙手沒有皺紋.
做出來的壽司是更加美味.
小野二郎他為了壽司.
他為了傳統的料理.
他為了客人.
他全力以赴.
奮身付出.
祂因為愛我們.
祂付上祂的生命.
然後託付我們.
甘願作祂的門徒.
我們延續福音這個大使命.
今天我們能夠為主耶穌.
能夠為福音.
能夠為未信的人.
我們付上什麼代價呢?.
《馬可福音》八章二十二節.
至十章五十二節.
這是一個很緊湊的段落.
也是《馬可福音》的核心.
耶穌就是願意立志跟隨祂的門徒.
就說做門徒的真正意義.
今天我們看其中兩段.
剛才主席為我們讀的經文裡.
我們聽到「凱撒利亞·菲勒比」.
這個被冠上凱撒稱號的地方.

$^{161}$原來是奉承,討好羅馬君王.
所以這不是地名那麼簡單.
它是代表「成」,「拜」,「榮」,「辱」的象徵.
耶穌就在這個地方.
祂問門徒「我是誰?」.
這個問是有一個考察的意思.
「比德就達,你是基督」.
基督是彌賽亞受高者.
在猶太信仰裡.
君王,祭司是受高者.
「比德」這個認信.
主耶穌在《馬太福音》十六章裡.
稱讚「比德」是有福的.
不過耶穌就說.
你們不要把這個身份告訴別人.
為什麼呢?.
原來在耶穌時代.
彌賽亞同樣聯想為一個首領,一個王.
他是一個帶領以色列人推翻羅馬政府.
一個王,一個民族救星.
這就和耶穌的使命大大不同.
所以耶穌不想其他人誤解祂的身份,祂的工作.
接著,剛才的經文31至38節.
耶穌就和「比德」門徒對話.
我們發現原來和耶穌這麼親近.
跟隨耶穌的門徒.
他們都誤解了耶穌的身份.
31節,當耶穌親口第一次說.
原來祂是受苦的彌賽亞的時候.
「比德」有什麼反應?.
「比德」拉著祂,阻礙祂.
原來這裡有一個意思.
「比德」說耶穌是滿口謬論.
大膽嗎?.
大膽吧.
原來「比德」不能接受耶穌受死,受苦.
「比德」攔阻了上帝在耶穌裡面的旨意.
「比德」抵擋了上帝的工作.
結果這次,耶穌沒有再稱讚祂有福.
耶穌責備祂.

$^{201}$耶穌和祂劃清界線.
不過我們看回經文.
耶穌的對象不單止「比德」.
因為33節裡面,祂同樣說「門徒」.
接著,耶穌就和「比德」和「門徒」一起說.
一定要透過受苦的彌賽亞.
去明白耶穌的使命.
千萬不要帶著錯誤的理由或盼望.
David Arthur De Silva 德籍爾雅.
他說,耶穌這個受苦的彌賽亞的揭示.
一方面喚起我們人一個真正的信心.
另一方面,祂興起盡心盡力.
遵行主道路的門徒.
耶穌表明自己的身份之後.
祂就鼓勵我們34節要「捨己」.
「捨己」是耶穌給門徒一個很重要的託付.
「捨己」,「捨去」.
原來這個字,我們看看經文.
同樣出現在《馬可福音》.
不過就在《馬可福音》的14章30節.
當時耶穌是預言「比德」,不認耶穌.
原來這個「捨去」的意思就是斷絕關係.
和這個人,我不認得他.
究竟今天弟兄姊妹和我一起作門徒.
我們不認得哪一個呢?.
我們和哪一個斷絕關係呢?.
原來是和「救」我.
我們有一個問題要問.
「救」是一定不好的嗎?.
我記得上兩個星期.
有機會去探望一個老姐妹.
她和她的家人都和我分享.
她說:喵喵,其實以前的人好像比現在更加捱得.
我自己立即點頭.
似乎「救」不是不好.
那為什麼我們要和「救我」斷絕關係.
我們說不認得這個「救我」呢?.
那「救我」的定義是什麼呢?.
我們看看剛才和大家分享的經文.
我們看看上文下理.

$^{241}$原來我們從耶穌責備「比德」這件事.
耶穌說「比德」,你體貼.
體貼在新漢語譯本,譯作「思想」.
在呂鎮忠譯本,譯作「意念」.
原來耶穌說「比德」,你體貼.
你全心全意將你的心思,意念,思想.
放在一些人的事.
不過這些人的事和上帝的事是對立的.
原來這個「救我」,我們要斷絕的「救我」.
就是一個肉體情慾的「救我」.
是一個自我中心.
是一個目中無神的「救我」.
當我預備的時候,我都問我自己.
我和這個「救我」一刀兩斷了沒有?.
畢哲思在《聖經讓你想得不一樣》這本書裡.
他有這樣的描述.
他說保羅在《迦拉太書》五章.
他列舉了不同的罪.
例如姦淫,拜偶像,邪術.
似乎我們都會對付這些罪.
不過作者問我們會否容易接納.
一些我們接納了的罪呢?.
咦?我們有罪我們會接納的嗎?.
原來他說的是一些我們思想.
我們心裡想的.
我現在看到大家.
我其實真的未必猜到大家在想什麼.
當然希望大家是同心我們仰望.
主和我們說話.
大家看到我可能都會想.
「喵喵你是不是有點緊張?」.
或者「你好想全心全力和我們分享」.
對呀!大家猜到.
不過有時人的內心我們猜不到.
人的思想我們猜不到.
原來有些罪別人看不到.
我自己怎樣?我自己都看不到.
然後我自己欺騙自己.
例如什麼呢?.
例如我們論斷.

$^{281}$我們很快就會看到.
或者說到別人的錯.
我們不知道事實的真相.
我們就批評別人.
作者說最可惜,最可怕的是.
連我們都不知道這些原來是罪.
而我們將這些念頭.
就放在一些疑惑裡面.
我們掩蓋著它.
我們以為對.
我們有沒有出路呢?.
在同一本書裡面.
作者勉勵我們.
我們可以嘗試.
當我們發現我們有這些意念.
我們看到別人的恃.
看不到自己的良目.
我們論斷別人我們說的是非.
我們覺得這樣不對,那樣不對的時候.
我們嘗試謙卑.
我們嘗試誠實地和主說.
我們靠神的恩典.
我們將這些不美的雜草.
在我們心裡面.
這些不美的雜草我們拔出來.
《路加福音》九章同樣有記載到寫己.
而九章五十七至六十二節.
它對寫己有另一個層次的詮釋.
或者我想大家打開你的大綱.
我想大家跟我一起請讀.
你們大綱裡面的經文.
「他們走路的時候」.
好,謝謝弟兄姊妹.
「他們走路的時候」.
「有一人對耶穌說」.
「你無論往哪裡去」.
「我要跟從你」.
「耶穌說:你有動」.
「你同你妓女有和」.
「至善人自然會有尋求的地方」.

$^{321}$多謝大家.
這段經文裡面.
我們對寫己有另一個看法.
剛才的寫己.
就是和一些同神對立的舊我.
我們和他一刀兩斷.
這次的寫己更加有挑戰性.
因為這段經文裡面.
有人主動願意跟從耶穌.
甚至到天涯海角都願意跟從主.
耶穌和他說.
狐狸,鳥仔都有他們的安樂窩.
都有他們的安書區.
不過耶穌跟從耶穌的沒有.
耶穌要挑戰願意主動跟從祂的人.
祂會問:你願不願意放下你的安書區.
這個安書區可能就是一個環境.
一個很安全的環境.
亦可能是大家厲害的東西.
大家的才幹,大家的恩典.
這就奇怪了.
不好的東西我們不要.
為什麼這麼好的東西我們都不要呢?.
為什麼呢?.
原來耶穌要我們學習信心的功課.
我們一起冒險.
我們一起用更大的勇氣.
我們去寫己,服侍主.
然後就換來更大的祝福.
我想大家看看身邊的弟兄姊妹.
很久沒有見到這麼多弟兄姊妹.
其實剛才我唱詩歌.
我真的有感動.
因為我們同心呼求阿爸父.
我們黃浸弟兄姊妹.
我們同心年輕人有年輕人.
弟兄姊妹有弟兄姊妹.
我們都很關心我們本地的福音工作.
不過同時我知道大家同步.
關心不同地方的福音工作.

$^{361}$無論是越南,國內事工.
或者是不同的地方.
特別早前,例如緬甸有事的時候.
我知道不少弟兄姊妹.
都會很盡心盡力去禱告,去呼求阿爸父.
如果我們回到大陸,回到國內事奉.
我想問問大家,有什麼是必殺技?.
大家有不少不同的答案.
我自己想,這個東西我沒有.
但是姓古的藝人也沒有.
就是流利的普通話.
因為我們的普通話真的很普通.
但是原來我們去國內事奉.
流利的普通話就是必殺技.
我們當中有一個肢體他分享到.
他很想透過這個恩慈回國內事奉.
正所謂大展拳腳.
不過他的心也是一條心為了主.
不過正所謂英雄無,勇武之地.
雞同鴨講.
原來國內不少人是講地方方言的.
我有我流利的普通話.
他有他精湛的國內方言.
原來大家是講不到的.
在這個時候,這位肢體他唯有怎樣?.
他要別人聽福音,他要關心別人.
他可不可以堅持說.
我一定要繼續操我流利的普通話.
不行了,他要靠耳朵,他要靠心.
他要更細心聽別人講.
而且用一個我們全部人都很厲害的功能.
就是國際語言,身體語言,我們去溝通.
結果這個肢體他學到一課.
他學到要放下自己一些恩慈,厲害的東西.
接著我就問他.
那你有什麼體會?.
你明明可以大展拳腳.
明明是一些好的東西.
他就說.
如果我厲害,我未必那麼靠著耶穌.

$^{401}$如果我厲害,我未必那麼體會到.
耶穌那麼愛人的心腸.
我們試問耶穌.
祂降卑得那麼厲害去親近我們.
其實我們應該聽不明白祂說什麼.
因為我們是太低層次.
但是耶穌祂卻願意降卑去明白我們.
所以姐妹覺得很能夠體會到耶穌的心腸.
祂更細微地服侍的意義.
各位弟兄姊妹.
原來社稷有時是一個近乎挑戰的社稷.
是一些耶穌要我們放下我們一些厲害的東西.
一些我們的經驗.
一些我們很有把握.
或者我們的安書區.
或者這一年我們很明白什麼是安書區.
我們有很多限制.
但偏偏因著這些限制.
耶穌就大大使用我們.
我們透過不同的方法.
我們仍然去愛人,愛神.
大家喜不更了不起?.
剛才一直跟大家說.
和這個一刀兩斷.
這樣你孤獨精了.
你什麼都沒有了.
所以經文裡很特別.
耶穌叫我們拿起一些東西.
祂叫我們背起一些東西.
祂說有些東西.
各位弟兄姊妹.
我們不可以不拿著.
《馬可福音》八章三十四節.
背起他的十字架.
耶穌給我們的十字架.
是耶穌給的.
耶穌為我們每一個度身訂做.
這件衣服,這十字架.
是大家能拿起的.
是剛剛好的.

$^{441}$是適合的.
由於是剛剛好,適合.
大家是獨一無二.
無人可以代替大家.
耶穌使用我們每一個.
在不同的場合去見證祂.
不過弟兄姊妹.
我們看清楚背起的.
原來不是一件西裝外套.
不是一件恤衫.
是一個十字架.
我們每人的十字架不同.
但十字架的象徵是一模一樣.
十字架的象徵就是全力擺上.
是你的心思意念.
甚至連你的生命都擺上.
上心一層一層.
我們有多少個會為主殉道呢?.
耶穌要求我們的是.
為祂盡心擺上的心.
這個心反映了耶穌在我們裡面有多重要.
一個愛我們的主.
在我們心裡面有多重要.
當祂對我們的愛不保留的時候.
我們又如何回應呢?.
其實我現在在五樓服事的時候.
我想後面有些導師都跟我一樣.
有時特別之前我們可以吃東西.
有時我見到弟兄姊妹.
會叫小朋友給我食物.
或者把他們喜歡的玩具給我.
大家猜他們通常會給我嗎?.
我見大家點頭.
我想告訴大家.
他們要衡量過是否喜歡.
如果他們很喜歡.
他們會有少許掙扎.
如果他們有很多包.
他們不太介意跟我分享.
不過我想告訴弟兄姊妹.

$^{481}$我應該沒有吃過很多檸檬.
包括我們的導師.
但我想說的是.
我們很心甜.
因為我們是否真的要他們的食物.
是否真的要他們心愛的玩具.
老實說有些玩具.
心愛到我要了都不會玩.
我坦白說.
但就是他們的心.
捨己與佩喜.
我不知道耶穌接下來.
要我們走的是什麼路.
但我們的心願意為耶穌付出多少.
我們說了捨己對神.
如何去跟從神.
我記得上次穆者在受苦節崇拜裡.
他也有說到.
捨己要作在人的身上.
捨己對人又如何呢?.
剛才主席為我們讀了九章三十至三十七節.
大家可以翻開.
《馬可福音》九章三十至三十七節.
耶穌提出捨己對人是如何的.
簡單來說.
捨己就是作眾人的僕人.
耶穌再次提出.
祂自己的身份是一個受苦的彌賽亞.
不過門徒的反應是怎樣的?.
我們覺得很惋惜.
是爭誰的大.
他們要威風,要有權,要有利.
這種爭大就是人的事.
就是體貼人的事.
是自我中心的事.
是與捨己完完全全相反的事.
聖經裡邊怎記載耶穌的反應呢?.
耶穌每次都循循善誘地愛我們.
愛門徒.
耶穌對他們說.

$^{521}$如果你跟從我,你捨己的.
你就作眾人的僕人.
眾人我們很容易理解.
是一個範圍.
代表所有人.
那為什麼耶穌接著會說小孩子呢?.
因為在希羅的時代.
無論是猶太文化.
或者是希臘文化,羅馬文化等等.
小孩子是被邊緣化.
視為脆弱,無利可圖.
視為麻煩的人.
所以耶穌提醒我們.
我們服侍所有人之餘.
我們服侍的深度.
就是一些被忽略.
別人覺得麻煩.
或者我們不能從他身上得到什麼的人.
我們就去服侍他.
上帝透過我們接待他.
我們歡迎接納這個人.
我們接待他.
讓他同樣經歷.
上帝與他同在.
上帝給他的祝福.
接著跟大家想一想.
究竟誰是小孩子呢?.
上星期我坐地鐵.
我看到有一個婦人.
她很熱心.
她看到地鐵門.
有一個懷孕的媽媽.
拿著東西.
她害怕別人撞到她.
她很熱心.
其實她自己也站在門口.
她很大聲說.
有沒有人讓座?.
我就讓不了座.
為什麼呢?.

$^{561}$因為我只是站在那裡.
我沒有位子坐.
我看到她說了兩三次.
我心想.
別人都沒有反應.
你還說.
但是她說了第三次之後.
結果真的有人讓位給懷孕的婦人.
對於這個女士來說.
在她眼中.
這個孕婦.
就是有需要的人.
是她的小孩子.
我們也會有愛心.
我們也會看到.
不少人都是小孩子.
而且也是實況.
別說別人.
我們自己也是小孩子.
我們也需要別人接待我們.
不過會不會有時.
我們焦點放錯了.
我們會質問.
我們會埋怨.
為什麼我看到有這個人有需要.
你看不到.
為什麼你看到.
你懶得去接待他.
你隨便去服侍.
而我是傾巢而出.
去服侍他呢?.
想不到連社稷.
我們也有機會去比較.
我們用了我們的角度來看社稷.
本來明明是上帝的事.
現在變成人的事.
勞雲在《向下的移動》.
基督的社稷之路裡.
他講到謙卑.
他說如果有一天.

$^{601}$當我們想我們比其他輪捨更謙卑.
原來我們目的只有一個.
我們希望得到周圍的人稱讚我們.
這不是耶穌的期望.
在《哥唔多前書》13章3節裡.
這樣記載.
「我亦將所有周濟窮人.
又社稷身叫人焚燒.
卻沒有愛.
仍然與我無益」.
一個不欣在基督的愛.
沒有愛動機的社稷.
動機錯了.
沒有人得益.
耶穌不喜悅.
如果有一天.
我還是跟別人比較.
我覺得我比別人做得好.
勞雲給了我們出路.
我們嘗試將這些掙扎.
這些負面的事.
我們跟主說.
而且他也鼓勵我們找一些.
你信任的同路人.
屬靈的導師.
談一談這件事.
反過來.
我不知道大家看新聞.
每天有大量的需要.
我們很想關心這些事.
我們很想關心世界的痛苦.
不過如果事與願違.
我不知道大家久而久之.
會不會產生無力感.
由於有這份無力感.
有人分析.
其中一個反應.
就是我們關閉我們的心.
我們不聽.
我們不理.

$^{641}$本來適當的抽身.
是不錯的方法.
不過會不會當我們處理得不好的時候.
我們會產生一種譴責.
最後最慘的是.
明明我們有一份基督徒的.
愛人的積極熱情.
都沒有了.
在《祈禱的探險》這本書裡.
作者提出一個方法.
我最近也在學習.
他說:你嘗試縮小世界.
什麼意思呢?.
既然世界有這麼多需要.
令我們有無力感.
當我們看到一些新聞.
一些需要.
一些身邊的人需要的時候.
你嘗試找一個人.
或者找一個家庭.
你當他就像你的朋友一樣.
你想一想他這個環境.
你就為他祈禱.
所以我自己早兩天.
我也為一個村的火燭.
那個家庭我禱告.
我嘗試代入.
如果現在這個家人.
他是什麼心情呢?.
我們一起試試好不好?.
潘博華他這樣說.
他說:社稷絕對不是一連串自我折磨.
或者禁慾主義的孤立行動.
社稷我們就是看著主耶穌.
然後我們就緊緊跟隨祂.
社稷是一個旅程.
大家都想旅行.
我們這個旅程.
我們有耶穌親自帶領我們.
所以我們就很安心了.

$^{681}$旅行的精彩.
是因為有很多意想不到的收穫.
擴闊我們的眼睛.
不過旅程也都會有意外.
有時我們真的不知道怎樣應變.
就好像當我們說社稷跟從主.
我們作門徒.
我們事奉的時候.
我們難免有時會消盡.
我們會疲乏.
我們會灰心.
我們會有壓力.
我們會失望.
感恩的是.
社稷不單只是一個旅程.
也是一個過程.
過程的意思.
就是我們可以學.
我們可以失手.
我們可以失誤.
不過主一定會繼續教我們.
而且是我們一個成長的機會.
最後我想和大家一起看.
《馬可福音》八章三十一節.
我們一起讀.
一,二,三.
從此.
他教訓他們說.
人子必須受許多的苦.
被長老,祭司長和民事氣絕.
並且被殺.
過三天復活.
耶穌基督這個受苦的彌賽亞.
被釘不是故事的結局.
結局是剛才大家讀最後的說話.
復活.
復活是甚麼呢?.
復活是德性.
復活是甚麼呢?.
復活是好氣在後頭.

$^{721}$耶穌很明白.
我們作門徒,寫記.
這條路是辛苦的.
耶穌知道.
既然復活是德性.
復活是好氣在後頭.
復活是主耶穌先擺上的時候.
我們的擺上就是值得.
願神賜福給大家.
\newpage



\section{哥林多後書 9:1-15}
\label{sec:opD_mMvcV8I}
\textbf{2021年5月2日主餐崇拜直播｜梁家麟牧師｜甘心作美事｜哥林多後書9\_1-15}
\newline
\newline
連結: \href{https://youtube.com/watch?v=opD_mMvcV8I}{\texttt{https://youtube.com/watch?v=opD\_mMvcV8I}} ~~~~ 語音日期: 2021-05-04
\newline
\newline
\hyperref[sec:oKGXX4r5yvQ]{\small{< < < PREV SERMON < < <}}
~
\hyperref[sec:index_chronic]{\small{[返順時目]}}
~
\hyperref[sec:d1qGfZ_AH3c]{\small{> > > NEXT SERMON > > >}}
\newline
\newline
哥林多後書 9:1-15
\newline
\begin{longtable}{cl}
\hline
\hline
章節 & 經文 (和合本修訂版)\\
\hline
9:1 & \begin{tabularx}{0.7\textwidth}{X} 關於供給聖徒的事，我本來不必寫信給你們； \end{tabularx} \\ \\ \relax
9:2 & \begin{tabularx}{0.7\textwidth}{X} 因為我知道你們的好意，常對馬其頓人誇獎你們，說亞該亞人預備好已經有一年了。你們的熱心感動了許多人。 \end{tabularx} \\ \\ \relax
9:3 & \begin{tabularx}{0.7\textwidth}{X} 但我差遣那幾位弟兄去，要使你們照我的話預備妥當，免得我們在這事上誇獎你們的話落了空。 \end{tabularx} \\ \\ \relax
9:4 & \begin{tabularx}{0.7\textwidth}{X} 萬一有馬其頓人與我同去，見你們沒有預備好，就使我們所確信的反成了羞愧；你們的羞愧更不用說了。 \end{tabularx} \\ \\ \relax
9:5 & \begin{tabularx}{0.7\textwidth}{X} 因此，我想必須鼓勵那幾位弟兄先到你們那裡去，把從前所應許的捐款預備妥當，好顯出你們所捐的是出於樂意，不是出於勉強。 \end{tabularx} \\ \\ \relax
9:6 & \begin{tabularx}{0.7\textwidth}{X} 還有一點：「少種的少收；多種的多收。」 \end{tabularx} \\ \\ \relax
9:7 & \begin{tabularx}{0.7\textwidth}{X} 各人要隨心所願，不要為難，不要勉強，因為神愛樂捐的人。 \end{tabularx} \\ \\ \relax
9:8 & \begin{tabularx}{0.7\textwidth}{X} 神能將各樣的恩惠多多加給你們，使你們凡事常常充足，能多做各樣善事。 \end{tabularx} \\ \\ \relax
9:9 & \begin{tabularx}{0.7\textwidth}{X} 如經上所記：「他施捨，賙濟貧窮；他的義行存到永遠。」 \end{tabularx} \\ \\ \relax
9:10 & \begin{tabularx}{0.7\textwidth}{X} 那賜種子給撒種的，賜糧食給人吃的，必多多加給你們種地的種子，又增添你們仁義的果子。 \end{tabularx} \\ \\ \relax
9:11 & \begin{tabularx}{0.7\textwidth}{X} 你們必凡事富足，能多多施捨，使人藉著我們而生感謝神的心。 \end{tabularx} \\ \\ \relax
9:12 & \begin{tabularx}{0.7\textwidth}{X} 因為辦這供給的事，不但補聖徒的缺乏，而且使許多人對神充滿更多的感謝。 \end{tabularx} \\ \\ \relax
9:13 & \begin{tabularx}{0.7\textwidth}{X} 他們從這供給的事上得了憑據，知道你們宣認基督，順服他的福音，慷慨捐助給他們和眾人，把榮耀歸給神。 \end{tabularx} \\ \\ \relax
9:14 & \begin{tabularx}{0.7\textwidth}{X} 他們也因神極大的恩賜顯在你們身上而切切想念你們，為你們祈禱。 \end{tabularx} \\ \\ \relax
9:15 & \begin{tabularx}{0.7\textwidth}{X} 感謝神，因他有說不盡的恩賜！ \end{tabularx} \\ \\
[1ex]
\hline
\hline
\end{longtable}
$^{1}$今日經訓是記載在哥林多後書9章1至15節.
新約聖經252頁.
哥林多後書9章1至15節.
「論道攻及聖徒的事.
我不必寫信給你們.
因為我知道你們樂意的心.
常對馬其頓人誇獎你們.
說:『阿該亞人預備好了,已經有一年了.
並且你們的熱心激動了許多人.
但我打發那幾位弟兄去.
要叫你們照我的話預備妥當.
免得我們在這事上誇獎你們的話落了空.
萬一有馬其頓人與我同去.
見你們沒有預備.
就叫我們所確信的反誠了羞愧.
你們羞愧更不用說了.
因此,我想不得不求那幾位弟兄先到你們那裡去.
把從前所應許的捐資預備妥當.
就顯出你們所捐的是出於樂意.
不是出於勉強.
少眾的少收,多眾的多收.
這話是真的.
各人要隨本心所著定的.
不要作難,不要勉強.
因為捐得樂意的人是神所喜悅的.
神能將各樣的恩惠多多的加給你們.
使你們凡事上上充足.
能多行各樣善事如經上所記.
他施捨錢財,周濟貧窮.
他的仁義傳道永遠.
那似眾給撒種的,似糧給人吃的.
必多多加給你們種地的種子.
又增添你們仁義的果子.
叫你們凡事富足,可以多多施捨.
因為辦這功級的事.
不但補聖徒的缺乏.
他們從這功級的事上得了憑據.
知道你們承認基督信服他的福音.
多多捐錢給他們和眾人.
便將榮耀歸於神.

$^{41}$他們也因神極大的恩賜.
顯在你們心裡.
就切切地想念你們.
為你們祈禱.
.
各位親愛的弟兄姊妹.
先講一點家事.
我們親愛的志強弟兄.
在幾天前離開了我們.
雖然也有一點心理準備.
但大家應該有很多很複雜.
難寫的感情.
我仍然記得那兩個字.
一個是Legacy.
我猜志強弟兄在我們中間幾十年.
有很多很多的故事存留在我們中間.
那些故事我們可以懷緬,可以回憶.
同時他也留下很多他事工的痕跡.
摸過我們中間很多的人.
那些會一代一代成為了.
仍然活著的Legacy.
成為我們的傳統.
成為我們的遺產.
今天很開心.
我見到幾位年輕的弟兄姊妹大師.
我覺得這是一個對志強弟兄的離去.
一個最好的應答.
我們基督徒不只是往後看的人.
往後看讓我們有很多很多的恩情.
讓我們可以記住.
但最重要的是我們向前看.
第一我們有一個終極的盼望.
我們知道耶穌基督會回來.
我們大家最終會在那裡.
永遠相聚.
而近一點的.
我們看到的是.
上帝在不同的世代.
繼續興起更多更多的人.
毀身給他.

$^{81}$我很喜歡聽到年輕人求異象.
我求主真的將心的激動.
放在你們每一個人的心裡面.
讓你們看到當前可以做的事.
沒有人需要重覆上一代的事.
你們不會是第二代的梁家倫.
不會是第二代的羅志強.
你們會做比我們更多更大的事.
望朕的盼望.
是在於我們有上帝.
我們也有一代接一代的弟兄姊妹.
好我們回來講今天的訊息.
我是第三次連續三次講有關金錢捐獻的問題.
如果你不記得.
之前我講過什麼的話或沒聽過的話.
你可以上網也看得到.
2月28日.
3月14日.
兩次.
為什麼會講三次呢?.
是因為跟著聖經.
那麼多後書.
8-9章連續兩章都是講金錢奉獻.
保羅囉嗦不是我囉嗦.
今天我是一口氣講完第九章.
將整章講完.
所以要講得很快.
保羅吩咐哥林多信徒履行他們的承諾.
完成為耶路撒冷信徒捐款的計劃.
第九章是比較麻煩的一章.
因為一開始收信人好像不是哥林多人.
而是亞該亞人.
你看九章二節就知道了.
並且內容上也跟前面有些重覆.
所以有不少人主張.
這一段其實是屬於另一封信.
後來編輯將它們合在一起.
在當時期.
保羅同時間照顧,牧養好幾間教會.
覆蓋小亞細亞和整個希臘半島.

$^{121}$他為耶路撒冷籌款.
他籌款的對象不只是一間教會.
不只是哥林多.
也包括了小亞細亞其他的教會.
所以當他委託他的助手.
替他將向某地的信徒籌款的信息傳遞的時候.
他可能同時也請助手向其他教會帶信函.
一個使者一個行程.
完成幾個任務.
好像今天那些包裹速遞.
不是只做一宗生意的.
通常一車就做很多宗生意的.
可能是信道這樣.
也有人相信.
這是同一封信.
保羅寫信給教會.
常常是用一個公開信的形式寫的.
在一個地方的信徒中間讀完.
然後又會去另一個地方再讀一次.
在旺鎮讀完去大角咀又讀一次.
就是這樣.
所以書信有第一的收信對象.
也可能同時搭單有第二對象,第三對象的.
不要浪費嘛.
反正當時沒有什麼很強烈的私隱觀念.
私人信函是可以公開的.
當時什麼都可以公開的.
就算你犯罪.
教會也會公開處理.
公開打屁股的.
不會怕冒犯了你的信譽.
不會這樣的.
所以保羅這封信主要是寫給哥林多教會的.
但是第九章前面一段是例外的.
是特別跟阿該亞的信徒說的.
保羅之所以沒有特別為阿該亞的信徒寫一封信.
而是搭單在這裡這樣做.
是因為他覺得沒有太多特別的信息要跟他們說.
在九章一到二節他這麼說.
他說:論到供給信徒的事,我不必寫信給你們.

$^{161}$因為我知道你們樂意的心.
常對馬其頓人誇獎你們.
說阿該亞人預備好了,已經有一年了.
並且你們的熱心激動了許多人.
保羅從前已經跟阿該亞信徒詳細說過有關捐獻的事.
他們不僅明白,並且已經切實進行.
有很好的表現.
這個博得保羅很開心.
保羅就在馬其頓教會的弟兄姊妹面前.
大力誇讚阿該亞很棒.
阿該亞很值得你們效法.
將阿該亞的帖堂變成榜樣.
所以他不覺得阿該亞還需要一些很特別的信息.
但為什麼他這次又特意寫一段信息給他呢?.
原因可能是保羅這個老人家.
就算說自己信心爆棚也好.
心裡還是有很多牽掛的.
老人家總是會有很多話要說的.
他對阿該亞信徒有信心.
但仍然要做一些提醒.
不怕一萬,最怕萬一.
第四節就提到萬一了.
第三,第四節他這麼說.
"但我打發那幾位弟兄去,要叫你們照我的話預備妥當".
"免得我們在這事上誇獎你們的話落了空".
"萬一有馬其頓人與我同去,見你們沒有預備".
"就叫我們所確信的,反正叫羞愧,你們羞愧,更不用說了".
很好笑的.
近東地區的人跟中國人一樣.
都是恃感文化,說shame的.
很怕丟臉.
沒面子.
保羅也要用羞愧一詞.
就是丟臉.
保羅首先在馬其頓人面前讚賞了阿該亞信徒.
但他又怕.
"慘了,我讚賞你們,你們天上有地下無".
"現在我派人來捐的時候,如果你們沒有預備,豈不是很大擊".
"那我又沒面子,你們又沒面子,對不對?".
"那我日後又如何對馬其頓弟兄姊妹說話呢?".

$^{201}$為了怕丟臉,出洋相.
保羅就要做些手腳.
九章五節.
"因此,我想不得不求勒幾位弟兄".
"先到你們那裡去,把從前所應許的捐資預備妥當".
"顯出你們所捐的出於樂意,不是出於勉強".
保羅於是就事先派人去向阿該亞信徒通水吹吹風.
"預備好,我派人來收錢的,你記得要預備好".
收錢人來到的時候.
又可能真的有馬其頓的弟兄姊妹來觀察.
看看誰綁架,看保羅吹得這麼大.
這班人是不是真的很正.
如果來到,沒錢了,那不就很大擊.
那不就很尷尬.
我們中國人,我們香港人,很會做這些事.
裝飾門面,是不是.
有視學官,或者觀摩的實習老師.
來課室觀課.
老師早就跟學生說好了,對不對.
我現在問問題.
懂得回答舉右手,不懂得回答舉左手,對不對.
不要舉錯,舉錯就糟糕了,對不對.
於是一問問題,竟然全班人都舉起手.
厲害吧,真是全國十強學校.
清華大學附中,附屬小學.
或者重慶巴屬小學,這些十強小學.
我這些什麼街坊福利會.
什麼紀念小學,竟然都有這樣的表現.
我們中國人很流行這些.
裝飾門面.
就這樣,所以這一段就是.
保羅跟阿居亞信徒通水.
我遲些要你收錢,麻煩你.
預備好,不然很容易上的.
就這樣.
給阿居亞信徒的訊息就只於第五節.
之後保羅所說的.
就是對金錢捐獻的一般性的教導.
又回到主題.
在這裡我要做一個補充.

$^{241}$你們多數都不是很聰明的.
如果你有仔細聽我之前兩次對金錢.
捐獻的講道.
你會發覺我有一節的經文.
是我故意跳了不講的.
我想故意去講.
這節經文很重要.
就是八章的第九節.
八章第九節說.
你們知道我們主耶穌基督的恩典.
祂本來富足.
祂為你們成了貧窮.
叫你們因祂的貧窮可以成為富足.
這節經文為什麼重要.
因為這節經文是講耶穌的榜樣.
講耶穌的榜樣也不重要.
開玩笑的.
這節經文講很簡單的東西.
凡事有代價.
救助也是一個價值的交換.
耶穌用祂的富足.
來換取我們的富足.
因此祂也替代了我們成為貧窮.
你知道富足和貧窮.
不僅僅是經濟上所有意義的.
保羅說.
耶穌由富足變貧窮.
是指著他甚至.
被上帝遺棄.
失去一切的東西.
在新學上.
這個叫做替代性的贖罪.
Substitutionary Atonement.
耶穌基督替代了我們.
有關耶穌基督的替罪犧牲.
我不在教義上做解釋.
先跟大家講一個不太對稱的例子.
今年年初因為疫情的緣故.
其實有很多很棒的東西.
包括香港藝術節.

$^{281}$竟然全部免費.
在網上任看.
我不知道有多少人看過.
我看了很多齣歌劇.
Orchestra.
當然都是錄製的.
但很愉快.
其中我看了一齣線上節目.
一個捷克的歌劇.
歌劇很有名.
叫做耶路發.
那裡講到.
在捷克農村.
有一個女孩叫做耶路發.
她未婚懷孕.
生了一個兒子.
她媽媽叫葛坦尼卡.
為了保護女兒.
把她和.
女兒和嬰兒.
都收起來.
然後跪求.
經手的男人.
娶她的女兒.
耶路發.
但那個男人拒絕了.
負情漢.
葛坦尼卡轉去.
求另一個.
一直喜歡她的女兒.
的男人.
娶耶路發.
那個男人答應了.
但那個男人有點擔心.
我要接受.
情敵的兒子.
將他變成我的兒子.
葛坦尼卡.
看到這個男人.
有些猶豫掙扎.

$^{321}$她把心一橫.
於是將那個.
嬰兒帶去結了冰的河.
將他.
淹死了.
幸好她女兒.
得到釋放.
可以自由開展新的生活.
並且她跟著.
回去和她女兒耶路發說.
她的兒子不幸逝世.
捷克人.
都是信天主教的.
所以當葛坦尼卡.
決定殺死她的孫子的時候.
她心裡想的是.
從此我要放棄.
永生的盼望.
成為終極被咒詛的人.
結果春天河水融化.
那個嬰兒的屍體被發現.
葛坦尼卡殺嬰的行為.
是一槍事發.
她要接受法律的制裁.
她的女兒耶路發.
也不肯原諒她的母親.
所以葛坦尼卡.
做母親的.
在被捕的時候她這樣說.
她說我是一個被咒詛的.
惡毒婦人.
從此我女兒也不會在我身上.
珍惜任何遺傳.
女兒的生活.
會和我徹底割蓆.
換言之.
我母親獨自承擔所有罪孽.
接受所有人對她的憎恨.
包括.
她一心想抱想救的女兒.

$^{361}$也憎恨她.
我們也不會贊同.
殺害無辜的.
嬰兒的做法.
我們也可以批評.
葛坦尼卡很蠢.
應該還有其他方法救自己的女兒.
不過在當時.
一個很保守的天主教的農村社會.
也不真是.
有什麼方法.
可以好好想.
幫助女兒保護.
名譽.
特別是被男人遺棄的恥辱.
葛坦尼卡選擇了.
她認為唯一可行的路.
然後承擔所有罪孽和恥辱.
承擔所有人對她的憎恨.
也甘心放棄.
人間的親情.
來生的盼望.
這真是名副其實的傳言蛇氣.
我知道不應該將.
母親的抉擇和耶穌基督.
相提並論.
因為葛坦尼卡只是一個平凡無能的婦人.
她所做的都是殺人犯罪.
和耶穌基督擁有.
上帝的身份和能力並且完全沒有犯罪.
是完全不同的.
我只是從這個歌劇.
我聯想到.
耶穌基督讓人類.
能夠擺脫和上帝永遠.
分離的命運.
確實是付出了.
沉重到匪夷所思的代價.
他捨棄不僅僅是生命.
更加是他所有的尊榮和權利.

$^{401}$創造主.
竟然願意為受造物.
捨命.
這會是怎樣的一件事呢?.
如果我們對葛坦尼卡.
值不值得為他的女兒.
做這樣的犧牲和疑問.
那麼上帝為人捨棄一切.
就更加是匪夷所思了.
值不值得?.
最近.
這個問題.
我一個以前的學生.
一個傳道人.
給了我一個很好的啟迪.
王信成傳道.
他在北角帶領一間基層的教會.
參與很多.
服務基層人士的工作.
信成上上在他的WhatsApp群組.
分享他事奉的反省.
今年.
農曆新年.
他和我們分享.
他默想.
有關揮霍的觀念.
有一個街坊.
從教會手中獲得了一小筆的金錢補助.
眾人以為他當然是.
拿去買罐頭.
那個街坊不是.
那個是媽媽.
他竟然帶他的兒子.
去吃了一餐麥當勞.
信成就說.
他一餐麥當勞.
對一個基層家庭的開支來說.
會不會太奢侈呢?.
信成他.
自問自答.

$^{441}$真是很奢侈.
然後他又說.
對一個媽媽來說.
他花上最大的努力.
掙扎生活.
只是想兒子開心.
過活.
信成再這樣說.
他這樣想.
真是很不恰當的奢侈.
但是.
上主竟然將他的兒子.
白白賜給本該受罰的人.
難道不是在見證.
上主揮霍的愛嗎?.
上帝的恩典.
從來都不是僅僅足夠.
而是慷慨奢侈的賜予我們.
這就是上主揮霍的愛.
反之是我們.
這群人跟隨著教導.
信成說.
他想到這些家庭.
可以在農曆年間.
可以毫無顧忌地吃一餐.
小朋友喜歡的麥當勞餐.
這個不僅僅是.
放肆.
更是讓他們看到.
上帝揮霍的恩典.
今年.
21日的訊息.
我是一個漢堡愛好者.
我很喜歡吃漢堡.
很喜歡吃漢堡包.
不過我不吃麥當勞.
不用解釋了.
沒想到對一個基層家庭來說.
去麥當勞.
都可以用揮霍一個詞.

$^{481}$你覺得有些.
真的冒冷汗嗎?.
我更加沒有想過.
可以用揮霍一詞.
去描述上帝對我們的愛.
是的.
揮霍這個詞真的用得很正.
很正.
上帝拯救你.
或者比較有智慧的長遠投資.
本小利大.
百倍回報.
上帝拯救我肯定就是一個揮霍.
嚴重浪費.
浪費了.
我是一個無能無恥.
也無聊的人.
不僅是罪人.
更是庸人廢人.
四不足惜.
在我身上壓重著.
可以預計.
這首歌怎能如此.
像我這樣罪人.
也能蒙主寶血救贖.
更加荒謬的是.
上帝不僅一次性.
這樣壓著在我身上.
不停加碼.
到如今他還死不悔改.
不肯止舌離場.
這首歌.
雖然我曾不是一度.
是百度叫你傷懷.
但是你愛卻又將我沉沒.
你說上帝揮霍不揮霍.
一個昔日又濫交又吸毒的愛滋病患者.
值不值得你救助.
一個嫖賭飲吹.
樣樣來搞到妻離子散.

$^{521}$一貧如洗的百夜公.
值不值得你照顧他中老.
是不是浪費社會資源.
今天你還可以問.
那些又自大又卑鄙的印度.
值不值得我們去救助.
你可以不停問這些問題.
值不值得.
你一問值不值得.
就不值得.
人不為己天誅地滅.
全世界只有我才值得你全程付出.
其他人全部浪費資源.
對不對.
我只愛我自己.
任何人都不值得我幫助.
金錢捐獻是不值得.
付出給人會不會浪費資源.
社稷救人是否揮霍無度.
真過你想想.
我們這一班第一世界大都會的專業人士.
去救援那些第三世界農村的貧民.
這樣的價值交換.
值不值得.
我只愛我自己.
任何人都不值得.
我只愛我自己.
任何人都不值得.
我只愛我自己.
值不值得.
當然不值得.
在人間.
糾纏愛是否可能.
愛是否值得的問題.
是永遠沒有出路的.
唯有當你.
看著十字架的耶穌.
然後才知道愛.
既可能.
又值得.

$^{561}$保羅講述了耶穌基督的榜樣.
放棄自身的富足.
好讓我們成為富足.
用這個榜樣來鼓勵我們.
和身邊的人.
分享我們的富足.
然後我們看.
第九章一般性的教導.
保羅提出第一個捐獻的原則是.
甘心樂意.
九章五節下到七節.
顯出你們所捐獻的.
是出於樂意而非勉強.
然後他再說.
要隨本心所決定的.
不要作爛不要勉強.
因為捐獻樂意的人是上帝所喜愛.
捐獻是自願的.
不是被迫不是攤牌.
不是一個政治或社會的任務.
不是為了討好權貴.
討好什麼人.
如果不是出自內心的自願.
上帝就不會.
越立你捐獻的東西.
前面保羅.
在介紹馬其頓教會的建設時.
也說過.
可以證明他們是按著自己的力量.
過了自己的力量.
自己甘心樂意的去捐助.
保羅說.
他們的捐助是甘心樂意.
並且不僅是量力而為.
是不自量力.
過了他們自己的能力.
這更加證明.
那是出於自願.
不是出於勉強.
我連續三次提及.

$^{601}$金錢捐獻的問題.
有人提及我.
我自己也提及自己.
在今天的處境.
可能有些犯忌.
我不知道你們是否.
也經常上網.
網上有不少網紅.
KOL.
是肆意攻擊教會.
其中一宗罪是教會斂財.
有人善意.
傳了一個陳儀.
的youtuber.
一段短片給我.
他在那段短片中.
大肆攻擊元朗某教會的.
牧者.
用威逼利誘的手段.
要求信徒做十一奉獻.
陳儀又嘲弄.
那個牧者.
本來是耕田人.
如今住大屋.
肯定在教會是掠水斂財.
他又說.
他有很多錢.
但他不斷斂財.
我發覺.
這個youtuber.
曾經發過不止一段短片.
責罵教會.
對教會非常hostile.
他大概是一個.
ex-church member.
離開教會的人.
唯有自覺自己在教會受過傷害的人.
才會和教會.
有不共戴天之仇.
不停地攻擊.

$^{641}$我在此沒有意思回應.
坊間對教會的攻擊.
我只是承認.
真的我們是.
處於一個對教會非常不友善的環境.
反對教會的聲音.
是來自四方八面.
藍營批評教會.
是社會運動的幕後黑手.
中央政府也相信.
教會和和平演變.
中國的西方力量.
是勾結的.
但與此同時.
黃營也指責教會.
守護社會建制.
保守勢力.
攔阻年輕信徒參與激烈的社會改革行動.
違背了聖經教導.
我們是兩頭不是人.
都受攻擊.
要答辯是沒有什麼需要.
只能夠做好我們自己.
謹守崗位.
實踐使命.
無法迴避別人的攻擊.
但我們不會因為害怕.
批評.
而改變我們的信息和使命.
我們在說那麼多後書.
我們在說第八第九章.
第八第九章是.
有關金錢捐獻的問題.
我們不可以避忌.
我不以聖經為恥.
聖經這樣說我們就這樣說.
當然我們要謹慎.
按著聖經的講論去說.
不要加鹽加醋.
不要用不恰當的理由.

$^{681}$或者威嚇.
要求你們大力捐獻.
然後大力渲染.
捐得多就上天堂.
捐得多就在上面.
可以住洋樓養番狗.
甚至說這是人間最聰明的投資.
擅自代上帝.
給奉獻者.
在人間有百倍回報.
或者恐嚇那些不願意奉獻的信徒.
你們偷了上帝的財寶.
將來要受地獄的刑期.
你們要被人家的財寶.
捐了上帝的財寶.
將來要受地獄之火.
你們可以做見證.
我這幾堂課都沒有說過這些.
我過去從來都沒有說過這些.
聖經是講壇的至高規範.
聖經沒有說的.
我們不會做任何隱身添加.
聖經說了的.
我們都不敢漏說.
照說.
再說.
在這裡聖經的主題.
鼓勵信徒幫助那些有需要的肢體.
捐獻.
做慈惠捐獻.
所以說捐給印度就對了.
捐給那些山區就對了.
這裡不是說.
提高牧者薪酬福利.
叫你捐錢.
不是.
所以.
用今天教會財務的分類.
這裡是說慈惠奉獻.
說慈惠捐獻.

$^{721}$不涉及我這個講員的私人利益.
我也不需要慈惠捐獻.
保羅在九章五節說.
不是出於勉強.
勉強這個詞有貪婪的意思.
保羅是提防.
有人批評他出於貪婪的心.
才扼殺信徒的錢.
所以他在後面.
也兩次強調.
在十二章的十七到十八節.
他和他的同工.
沒有佔過哥林多信徒的便宜.
他呼籲他們捐錢.
不是出於自利.
慈惠不涉及個人利益.
無論如何.
金錢捐獻一定要出於自願.
不可以夾雜任何的勉強.
這個不是保羅第一次說.
是說了不止一次的原則.
接著我們要說.
另一個原則.
除了不可以勉強之外.
另外要說的是快樂的捐獻.
我上次說到.
施比授更為有福.
這個聖經的原則.
我講到施贈者.
本身就已經讓人覺得.
施贈本身.
捐錢本身.
已經讓你覺得很幸福.
很快樂.
你不需要藉著行善積福.
將來有百倍回報.
然後才覺得快樂.
我這樣說完之後.
有人批評我.
如果行善.

$^{761}$是讓自己快樂.
那不就像施治尼.
你只是為了.
加值給自己.
加添自己的好處.
那又有什麼道德價值可言呢.
他還引用耶穌在登山寶訓中.
對祈禱和禁食的教導說.
如果人間有任何好處.
那天上就不會有好處了.
人間拿了回報.
還會有嗎?沒有了.
我看到這個回應就覺得很好笑.
今天說什麼都會有急著來.
都會招來批評.
我的回應很簡單.
首先我真的希望.
你和我.
如果我們真的願意.
捐錢去幫助有需要的人.
我真的不是為了.
將來得回報.
你不是出於想回報.
你不是因為將來有大獎品.
不是的.
你喜歡事奉.
事奉是我開心的事.
確實我不是為了回報.
而做這些事.
我沒有積欽德.
為來生,為子孫.
種善因德善火的觀念.
所以.
如果我們在行善的過程中.
自己快樂.
將來在天上就不會有賞賜.
那我就會說沒有賞賜就沒有賞賜.
對不對.
有什麼關係呢.
我們就求今日快樂.

$^{801}$我服侍是為了今日我開心.
我的金錢奉獻是為了今日快樂.
我幫助人服侍人.
是為了今日快樂.
不是為了長遠投資.
不是.
我求人得人.
我不求將來回報.
那行不行.
那就不用爭論.
好像我前面提過一次.
那個弟兄.
他很喜歡送錢給人.
並且他說.
給出去的微不足道.
收到的人覺得是意外之財.
付出的覺得少出.
收穫的覺得收得多.
收支很不平衡.
所以很抵.
這個想法套在經濟學上.
就是將.
事物的用途.
予以最大化.
最優化.
optimization of utility.
power of power.
是不是.
多餘的錢放在.
有錢人身上.
那一百元不值錢.
帶給你的快樂很少.
你將它給一個.
窮寡婦.
那一百元對她很大.
你不是給了很多出去.
但她收了很多.
是不是很抵.
所以錢放在你身上.
是剩下的.

$^{841}$是浪費的.
放在有需要的人身上.
賺了便是有好處.
抵了.
所以.
我們的錢.
在我們付足的人身上.
所能製造的快樂.
遠遠不及放在窮人口袋裡.
是不是.
這是law of diminishing marginal return.
是不是.
所以.
將多餘的轉贈別人.
讓別人快樂.
省下位置.
衣服穿多了.
穿不完.
你又多個藉口賣衣服.
是不是很快樂.
弟兄姊妹.
我講這麼多.
我只是想簡單講.
我希望我們基督徒.
我們喜歡我們自己的生活方式.
今天我們參與服事.
今天的金錢奉獻.
在身邊的人.
我未信主的媽媽.
很浪費錢.
覺得自己很無聊.
做出錢出力這麼愚蠢的事.
是不是很蠢.
但我自己.
我希望你們也一樣.
覺得自己很幸福快樂.
我覺得.
這將我們的生命.
做了一個optimization.
最優化的使用.

$^{881}$我不事奉.
我每天去打麻將.
或者傻乎乎.
做個cushion.
放在沙發上對著電視.
對不對.
一夜打麻將.
你肯定更醜.
你覺得你參與了.
事奉後的樣子好了嗎.
真的變得帥.
真的變得好.
我覺得你們每一個.
神清氣爽.
所以.
奉獻事奉.
不是為將來得快樂.
是今天就已經快樂.
我希望你們同意.
我的說法.
這是我們大家共同的經驗.
共同的信念.
我以基督徒為榮.
以我過基督徒的生活方式為榮.
星期日.
上午是不是.
排隊找位喝茶的機會.
我不覺得.
我覺得回來教會很快樂.
是不是.
燒著來施贈.
施贈後.
開懷歡笑.
比你哭著來施贈好.
施贈後.
心如刀割.
施贈者感到幸福.
被感到倒霉.
是不是要強很多.
這也是聖經的教導.

$^{921}$行善者要.
自願甘心.
不可以有絲毫勉強.
所以只能燒著做.
不可以哭著做.
為有需要的人.
有慈悲憐憫的警方.
哭著做.
為捐出的錢.
哭著做就無謂了.
你還會將好事變成壞事.
對不對.
九章七字說.
不要作爛不要勉強.
捐得樂而.
捐得喜樂.
頂子妹.
我希望.
這是我們每一個人的事奉態度.
我們的奉獻.
我們的付出.
全部都是既是.
甘心也都是喜樂.
當然保羅也不是完全沒有提到.
金錢捐獻為人帶來的好處.
接著.
八節到十五節.
這裡有幾個教導.
我很快地跟大家講.
第一.
上帝會額外給我們保護.
和上報.
讓我們不只因為金錢捐獻而變得貴乏.
八章九節.
上帝能將各樣的恩惠多多.
加給你們.
凡事常常充足.
能多行各樣的善事.
在這裡保羅沒有說.
會給百倍的補償.

$^{961}$只是說.
你捐了出去.
也不會缺乏.
變成你繼續有能力捐獻.
是這樣的.
你不要說我現在投資一塊錢.
然後賺十塊.
然後到街上.
一踢就踢了一張百塊紙.
你省一點不要想這些.
不是這樣的.
我們仍然沒有缺乏.
這是我們大家可以有的見證.
我們不會因為捐一次.
就變成窮光蛋.
就再沒有能力捐第二次.
我們繼續有能力.
還可以常常捐出去.
第二.
上帝允許讓施贈者的德行建立.
人格成全.
聖經.
111.
第九節.
72.
詩.
112.
他捨棄.
施捨錢財.
周濟貧窮.
他的人義傳道永遠.
奉獻給人.
和人分享.
是建立我們道德人格的一個途徑.
這是.
生命的殺種.
少眾的少收.
多眾的多收.
這才是真理.
我們做一些.

$^{1001}$好像給了人的事.
我們吃虧了.
其實不是.
我們是在鍛造我們自己的人格.
我們的氣質.
一個有愛心.
一個對人有同情共感的人.
你覺得真的不一樣嗎.
我們在鍛造自己的氣質.
金錢捐獻除了幫助捐獻者建立.
崇明生命和道德之外.
也會造就其他人的生命.
受助者得到物質的幫補.
也會視這些幫補是來自上帝的恩典.
就將榮耀歸給上帝.
我們永遠.
藉著我們對人的愛心.
成為了我們福音的見證.
最後.
保羅用一句讚美文做結束.
因祂有說不盡的恩賜.
一切好處都來自上帝.
我們之所以能向人捐獻.
是因為上帝給我們極大的豐富.
我們之所以向人有恩賜.
是因為我們領受了上帝的恩典.
愛心捐獻.
能夠使人將榮耀歸給上帝.
保羅自己示範.
他所說和他所做的.
他說:我都希望將榮耀歸給主.
子強弟兄用他一生.
去見證了怎樣將榮耀歸給主.
我們很多弟兄姊妹.
我們都繼續踐行這個精神.
我們參與.
我們服事.
我們付出.
因為我們想將榮耀歸給我們所愛的主.
\newpage



\section{哈巴谷書 3:16-19}
\label{sec:d1qGfZ_AH3c}
\textbf{2021年5月16日崇拜直播｜陳冠成牧師｜從聖經看喜怒哀愁~因救我的神喜樂｜哈巴谷書3\_16-19}
\newline
\newline
連結: \href{https://youtube.com/watch?v=d1qGfZ_AH3c}{\texttt{https://youtube.com/watch?v=d1qGfZ\_AH3c}} ~~~~ 語音日期: 2021-05-16
\newline
\newline
\hyperref[sec:opD_mMvcV8I]{\small{< < < PREV SERMON < < <}}
~
\hyperref[sec:index_chronic]{\small{[返順時目]}}
~
\hyperref[sec:RH7uL12ghP4]{\small{> > > NEXT SERMON > > >}}
\newline
\newline
哈巴谷書 3:16-19
\newline
\begin{longtable}{cl}
\hline
\hline
章節 & 經文 (和合本修訂版)\\
\hline
3:16 & \begin{tabularx}{0.7\textwidth}{X} 我聽見這聲音，身體戰兢，嘴唇發顫，骨中朽爛，在所立之處戰兢；但我安靜等候災難之日臨到那上來侵犯我們的民。 \end{tabularx} \\ \\ \relax
3:17 & \begin{tabularx}{0.7\textwidth}{X} 雖然無花果樹不發旺，葡萄樹不結果，橄欖樹也不收成，田地不出糧食，圈中絕了羊，棚內也沒有牛； \end{tabularx} \\ \\ \relax
3:18 & \begin{tabularx}{0.7\textwidth}{X} 然而，我要因耶和華歡欣，因救我的神喜樂。 \end{tabularx} \\ \\ \relax
3:19 & \begin{tabularx}{0.7\textwidth}{X} 主耶和華是我的力量，他使我的腳快如母鹿，又使我穩行在高處。這歌交給聖詠團長，用絲弦的樂器。 \end{tabularx} \\ \\
[1ex]
\hline
\hline
\end{longtable}
$^{1}$今日我們所選讀的經訓.
是在舊約聖經.
哈巴谷書.
第三章十六至十九節.
是教會的聖經舊約.
一千二百八十三頁.
但都列印在大家手上的程序表的內頁.
哈巴谷書三章十六節.
「我聽見耶和華的聲音.
身體賤敬 嘴唇發賤 骨中扭爛.
我在所立之次 鑽進.
我只可安靜等候災難之日臨到.
犯警之民上來.
雖然無花果樹不發旺.
葡萄樹不結果.
橄欖樹也不效力.
田地不出糧食.
圈中絕了羊.
棚內也沒有牛」.
第十八節想和大家一起讀出.
「然而我要因耶和華歡恩.
因救我的神喜樂.
主耶和華是我的力量.
他使我的腳跨於無陸的堤.
又使我穩行在高處.
這歌交予靈長.
用詩言的樂器」.
願神祝福大家.
頂智慕 平安.
(同學)平安.
在五,六月.
幾位同工會透過四次宣講.
分享一個主題.
我們會從聖經裡.
從聖經裡看.
四個我們會有的情感反應.
分別是喜,怒,哀和愁.
今天我們首先會說「喜樂」.
我們說「喜樂」.
基本上在基督徒也很熟悉的詞語.

$^{41}$「喜樂」.
無論在新約或舊約.
都經常提到這兩個字.
對不對?.
不過在分享之前.
我想先和大家退後一步.
其實我們不妨想一想.
在你心目中.
甚麼是「喜樂」?.
又或者怎樣才可以所謂「活得喜樂」?.
特別當真的有人問你.
可能是你身邊的一個弟兄姊妹.
突然間這樣問你.
或者你家人這樣問你.
或者一個很想認識信仰的人問你.
基督徒經常說「喜樂」.
是甚麼來的呢?.
你又會怎樣回答他呢?.
讓我們先有個祈禱好嗎?.
我們求上帝開我們的心竅.
阿爸天父願你的靈願我們同在.
讓我們有智慧有悟性.
能夠明白你自己的教導.
讓我們有柔和謙卑的心.
能夠實踐出來.
讓我們在生命裡面.
能夠榮耀你.
能夠造就身邊的人.
為此我們同心獻上禱告.
奉基督耶穌.
你得勝的名字來祈求.
阿們.
說起來其實.
不知道有沒有聽過.
古代人其實很強調.
人生有四大快樂.
不知道有沒有聽過呢?.
第一個.
我先試著說出來.
看看你記不記得之後的三個是甚麼.

$^{81}$第一個很大的快樂就是.
金榜提名事.
又點頭了.
古代科舉制度.
你高中狀元.
代表著甚麼呢?.
出了名.
不只是出名.
你的前途.
甚至你會得到很多的好處.
利益.
名勝利就.
簡單說四個字.
第二個很大的快樂是甚麼呢?.
第一個就是金榜提名事.
第二個呢?.
厲害厲害.
洞房花燭夜.
很實際.
你和所愛的人.
結為一體.
成家立室.
洞房花燭夜.
第三個呢?.
久旱逢甘露.
想不想起這個畫面.
其實照字面理解.
就是你口渴了很久.
突然間即使是一些露水.
你都覺得真的很開心.
很滋潤.
又或者在農業社會裡面.
你最希望見到這件事.
其實今天也不難理解.
久旱逢甘露.
有沒有突然間喜歡吃自助餐呢?.
應該是吧.
過去兩年難一點.
很多時候限聚令.
但一放寬一點.

$^{121}$你見到自助餐又如何呢?.
訂滿了.
久旱逢甘露.
很久都沒有吃.
現在可以吃了.
又或者可能這兩年.
有沒有聽過一個新興的名詞.
叫做「偽出國」.
偽即是真偽.
即是假出國.
意思是甚麼呢?.
不只是發生在香港.
可能不同地方的人.
他們在疫情下.
都不能去外地旅遊.
那如何止止咳嗽呢?.
有些旅行公司很懂得做生意.
我們都爭取到一些航班.
讓你上飛機.
讓你出發.
去到當地國家的上空繞一圈.
然後回來.
有些貼身到.
譬如你去日本.
可能由台灣出發去日本.
在日本上空繞一圈.
然後回到台灣.
還會給你一份日本的手信.
幫你買了.
讓你有甚麼呢?.
有去了旅行的感覺.
讓你止咳.
雖然真的去不到.
不能到當地.
但也讓你有坐飛機的感覺.
久旱逢甘露.
最後一個人生的快樂是甚麼呢?.
他鄉如故知.
都不難理解.
當你在情感上和人有聯繫.

$^{161}$特別是你所愛,你所熟悉的人.
見到面.
譬如疫情的時候.
我們很多時候都在Zoom見面.
但當有實體崇拜.
你就會體會到那種珍貴.
大家見到面的時候.
好像他鄉如故知的感覺.
令你很開心.
但不知道大家有沒有留意.
剛才說的這四種快樂.
其實某程度上.
是和外來事物有關的.
有一個直接關係.
是因為在我們人生中.
遇到一些好事.
很值得開心的.
所以我們心中有份喜悅感.
有它的意義.
也很有它的價值.
但我們說的這不是聖經所說的喜樂.
聖經裡說的喜樂.
其實不是純粹說一種情感.
很多時候它更加指向一種.
我們說的積極的生活態度.
就是無論我們落入一個什麼環境.
總是對上帝.
心存一份信心和盼望.
因為我們認定.
上帝已經掌管一切.
祂也會帶領我們.
祂也會看顧不會虧待我們.
並且我們能夠以.
上帝能不能夠得到榮耀.
上帝的旨意能不能夠得到彰顯.
作為我們歡喜快樂的原因.
這是聖經裡所說的.
喜樂的真正焦點.
換句話來說.
聖經裡說的喜樂和外在環境.

$^{201}$到底是好還是壞.
到底有沒有好的事情.
臨到我們身上.
又或者我們個人的願望和需要.
是否一定達成了.
原來沒有一個必然的關係.
話說.
有一個猶太小朋友.
有一天他對自己的母親說.
「媽媽,如果你買一輛單車給我」.
「我就會很喜樂」.
於是他母親就很有智慧地.
翻開聖經來對他說.
「兒子」.
「喜樂其實不是由外在而來」.
「不是由某個人給你的」.
「喜樂是你自己的責任」.
這句話很有意思.
因為聖經裡也強調.
「喜樂是你自己的責任」.
去培育.
去實踐的一種屬靈品格.
你還記得史奴保羅所說.
聖靈所結的果子.
其中一種特質就是.
「喜樂」.
可以說是每一個基督徒.
一種責任.
特別你經常聽到.
保羅有一句金句.
「要靠主常常喜樂」.
說一次都不夠.
還要說多一次.
「我再說,你們要喜樂」.
不是說在一個好的環境.
是在一個不容易的環境裡.
我們要靠主喜樂.
好像一個命令.
背後也代表著.
是我們辦得到的.

$^{241}$因為上帝會幫助我們.
事實上.
我們也明白.
一個喜樂的人生.
對我們身心也很有益.
我們也很熟悉.
那句金句.
「喜樂的心」是什麼?.
乃是良藥.
相反「憂傷的靈」.
是「骨枯肝」.
其實以上這些聖經真理.
我們也聽過.
我們也明白知道.
不過實踐起來.
大家覺得容易嗎?.
有時又好像「喜樂」不來.
現實生活很多場景.
有時又很難叫我們「喜樂」.
對不對?.
不知道大家有沒有這種體會.
特別可能我們有時也聽到.
身邊會有人說.
等我找到一份理想的工作.
我就會「歡喜快樂」.
等我找到我所愛的另一半.
我就會「喜樂」.
等我賺夠.
到時我就「喜樂」.
等我捱大子女出生.
到他們事業有成.
那時一定會「歡喜快樂」.
到我退休.
不用工作的時候.
就真真正正「喜樂」.
不知道上述一些想法.
你有沒有聽過.
甚至出現在你的腦海裡.
但我們也明白.
最後即使這一切目標都達成.

$^{281}$也不見得我們會像聖經所說.
「歡喜快樂」.
因為人生總有其他難題.
是你和我會遇上.
但又解決不到.
是否我們不努力呢?.
也不是.
是否我們對成功的執著不夠強烈.
亦似乎不是.
很可能問題出在次序上.
我們將「喜樂」的次序.
有時會調轉.
我們會將「喜樂」當作一個「副產品」.
意思是我先要達成.
某個人生的目標和心願.
然後我就會「喜樂」.
「喜樂」是次要的副產品.
變了.
或許你經常聽到人說.
「我今天很不開心」.
因為這樣.
因為我想得到的東西得不到.
「我今天很不喜樂」.
因為我不想失去的東西失去了.
「我今天很不喜樂」.
因為有誰這樣對我.
我們將「喜樂」當作副產品.
但聖經是要改變我們這種想法.
它要求我們.
你先學懂甚麼是「喜樂」.
先學懂「喜樂」.
原因是因為.
上帝透過耶穌已經救了我們.
並且已經恢復了.
這位造物主和我們的真正關係.
我們不再是地上的奴僕.
我們是上帝的僕人.
是上帝的兒女.
或許這個世界有很多事情.
我們都應付不來.

$^{321}$很混亂.
但我們確知.
最終上帝會將一切撥亂反正.
所以我們仍然可以「喜樂」.
並且我們能夠帶著這份「喜樂」的力量.
向著上帝給我們的目標和人生方向.
前行.
聖經是這樣處理.
先學懂「喜樂」.
然後透過這份「喜樂」的力量.
去面對你生活不同的場景.
今天所分享的經文.
《哈巴谷先知》.
它就是一個最好的例子.
當然在這裡要交代一點.
這封書信的背景.
這封先知書的背景.
《哈巴谷先知》.
大概在猶太國亡國前幾十年.
他擔任先知.
他曾經服侍的有.
約西亞王到約瓦競王.
約西亞王我們最熟悉的.
就是他當時在宗教.
或者在政治上推動的一些改革.
令到國家當時有一段短暫的復興的時候.
不過後來很不幸.
他突然間戰死沙場.
於是皇位就輾轉到約瓦競的身上.
聖經記載約瓦競是一個惡的王.
不遵行上帝的吩咐.
所以國家也開始越來越敗壞.
百姓也越來越遠離上帝的吩咐.
而在這個時候.
哈巴谷充滿了很多的焦慮.
很多的困惑.
他有很多的不明白.
於是乎他就在《哈巴谷書》裡面.
向上帝提出一些申訴.
他問上帝.

$^{361}$他說:上帝.
猶太國現在越來越敗壞.
為什麼你不出手.
你要袖手旁觀?.
為什麼你容許.
有一些不公義.
甚至乎暴力的事情出現?.
你不出手去管教你的百姓.
那些不聽話的人.
又或者甚至乎出手去懲治惡人?.
於是乎上帝就回答了.
他會差派巴比倫.
作為審判猶太的工具.
並且時後近了.
猶太國也將要被巴比倫所滅亡.
上帝不回答猶自可.
他這個回答令到.
哈巴谷才更加的困惑.
原因是.
上帝你是公義的.
你有什麼理由用殘暴不仁的巴比倫.
一個這麼邪惡的國家.
來管教懲治你的百姓?.
沒錯上帝.
我們是做得不好.
我們是得罪你.
但你沒有理由用一群比我們更差的人.
來教訓我們的.
你沒有理由用這麼奇怪的方式.
來讓我們變回好的.
才會有一個更加的疑惑.
於是乎上帝再一次回答他.
「我旨意已經決定了」.
「你將它寫下」.
意思是不會改變.
審判一定會臨到.
你就安心等候這一天的來到.
上帝亦表明.
公義的人必定會因為對上帝的忠信而存活.
至於巴比倫.

$^{401}$亦要因為他自己的不義.
來承擔代價.
臣最後亦會懲罰他.
讓他消滅在地上.
上帝說得很清楚.
弟兄姊妹,停在這裡.
如果你是哈巴谷先知.
試試代入他的心情.
聽到上帝這麼清楚的回答.
你滿意嗎?.
收不收貨?.
會否因為明白了上帝的旨意.
就馬上變得心情愉快喜樂?.
會否呢?.
我看大家都靜下來.
大概都很難消化.
那時不容易.
就算我們放回今天.
我們明白了上帝的旨意.
我們知道上帝是這樣決定.
不等於我們會喜樂地去接受.
對不對?.
我們可以接受.
不過用另一個態度.
譬如是甚麼?.
我們可以.
沒辦法.
上帝你都決定了.
我不可以推翻.
我不鋸.
但我都要接受.
我接受,但我無奈.
我可以消極.
未必一定是喜樂地去接受上帝的心意.
對不對?.
確實.
又真的不容易去喜樂.
在座當中.
大部分都是上班一族.
我們都有自己的上司.

$^{441}$對不對?.
如果讓你選擇.
你想選擇一個.
厲害的上司.
好的上司.
還是選擇一個差的上司.
壞的上司?.
其實都不用問.
可以讓我們選擇.
還是選擇一個讓我們優秀.
讓我們良善的人.
來帶領管理我們.
但現在.
之所以很難喜樂.
就是我要接受.
我讓一些比自己更差.
更壞的人來教訓.
來管我們.
怎麼能夠?.
甚至他們會破壞我們的生活.
破壞我們的安寧.
確實真的很難去喜樂.
另一方面.
先知很不容易去喜樂的原因.
是因為.
假如先知自己都是壞人.
他自己都是帶頭得罪上帝.
所以落得如此下場.
那又如何?.
四個字.
心甘命瀝.
做了錯事.
受罰.
很正常.
但是.
先知雖然不是完美.
但應該都沒有做過.
什麼傷天害理的壞事.
他都沒有立心不良去害過誰.
是哪些人不好?.

$^{481}$是我身邊的人不好.
所以我無辜被牽連.
我都要面對各破家亡.
雖然其實不太關我事.
所以不容易去喜樂.
另一個更加難喜樂的原因是.
我無能為力改變這個困局.
因為上帝旨意已定.
就算我有天大的本領.
我都改變不了.
我只能夠眼白白看著.
十七節描寫的情況出現.
是什麼?.
無花果樹不發旺.
葡萄樹不結果.
橄欖樹不效力.
田地不出糧食.
圈中絕了羊.
棚內沒有牛.
我眼白白看著這些事情出現.
不用同時間出現.
就算每個月出現一樣.
都夠難受.
意味著什麼?.
無花果樹.
你明白嗎?.
就是以色列人其中一個.
很主要的食糧.
特別曬乾了可以儲存很久.
隨時用來沖洗.
葡萄樹呢?.
不是葡萄.
是葡萄樹所有的酒.
包括自己享用.
也可以用來做生計.
橄欖樹.
油.
不單是煮食.
是獻祭的油.
牛羊呢?.

$^{521}$牛羊不單是有沒有牛肉吃.
影響到他們的宗教生活.
獻祭受影響.
五穀也不用說了.
你會看到.
眼見自己身處的地方.
無論是自己的生活.
經濟.
生計.
甚至是宗教生活.
都一直每放月下的時候.
整個社會慢慢陷入蕭條.
甚至崩壞.
何來喜樂?.
但奇怪的是.
第十八節才說什麼?.
我仍然可以歡喜快樂.
很瘋狂吧?.
怎麼可能仍然可以歡喜快樂呢?.
你看到他沒有說.
等我見到無花果樹再次發旺.
見到葡萄樹再結果.
橄欖樹為我再效力.
田地有糧食.
圈中滿羊.
棚裡有牛.
一切雨過天青.
我就會喜樂.
他沒有用這個邏輯來思考.
他卻說了什麼?.
然而.
這兩個字可圈可點.
代表什麼?.
就算這一切沒有變好.
就算這一切繼續崩壞.
就算所有身邊的人都認為.
看到這個場景.
都開心不下.
喜樂不下的時候.
我仍然可以喜樂.

$^{561}$原因是他自己交代了.
我仍然可以因上帝歡欣.
因救我的神喜樂.
哈伯曲很明顯.
他看到喜樂不是外在加給他.
是來自上帝.
不是外面的事物影響他有沒有喜樂.
是因為他結連上帝.
所以他能夠喜樂.
並且他知道上帝會給他力量.
去和那些禍患相處.
所以才知道他用他的親身經歷.
告訴我們秘訣是什麼.
接下來.
首先你看到.
第十六節先知這裡說什麼?.
「我聽見耶和華的聲音.
身體顫抖,嘴唇發顫.
身體顫驚,嘴唇發顫.
骨中扭壞.
我在所站之處顫驚.
我只可安靜等候.
災難之日臨到.
犯警之民上來」.
這裡給我們第一個提醒.
其實人很多時候.
喜樂不了是什麼原因?.
可能是我們不願意接受.
上帝在我們生命裡的一種安排.
那些不舒服.
那些會讓人痛苦的安排.
我們不太接受.
即使是出於上帝的旨意.
但我們怎樣也不想淋到我們身上.
於是我們自然就會啟動一些.
我們所謂的防衛機器.
可能是我們會設法迴避現實的出現.
我們會自我安撫說.
「沒事的.
就算上帝是這樣決定.

$^{601}$未必會發生的」.
自我安慰一下.
又或者我們會拼命擺脫現實.
我會想盡辦法.
來試圖解決眼前的困境.
解決不了我就離場.
讓自己安心一點.
又或者我會將問題.
最壞的都是它.
不是它我弄成這樣.
我會試圖迴避.
其實自己也有些要承擔的責任.
甚至第四種防衛機制就是.
我會將情緒向我身邊無辜的人發洩.
掩飾我自己.
內心其實很焦慮.
很恐懼.
可能以上這些防衛機制.
我們多多少少也試過.
有時可能也會有短暫的效用.
但一定不會給我們真正的喜樂.
原因是.
這些方法都有一個共通點.
它沒有處理到我們和上帝的關係.
你有沒有想過?.
它沒有將我們和上帝結連在一起.
所以不會有喜樂.
你看到哈巴菊才知道.
它之所以這樣也能夠有喜樂.
在於它很坦然接受.
兩件事發生在它身上.
第一.
它接受上帝有權.
不按它的期望去做事.
上帝有這個主權.
我稱得它為上帝.
認得它為主.
就代表它說了算.
不是我說了算.
它有權不按我的旨意出牌.

$^{641}$這是第一.
第二就是它很坦然接受.
上帝容許人的生命中.
有苦難和痛苦出現.
它接受這件事.
但它不代表我玩完了.
它接受有苦難和痛苦.
但不代表我們已經玩完了.
因為上帝仍然掌管.
上帝沒有玩完.
一切仍然在它的掌握裡.
就好像我聽過一個例子.
我們也熟悉宗教改革加瑪丁路德.
我們也認識他.
有一次他也曾經面臨人生很大的困局.
他極其的受煩.
他太太看到.
故意在他面前.
穿著一件裝服.
大概全身都黑黑的.
瑪丁路德看到就很驚訝.
他就問:誰去世了?.
他太太故意跟他說.
上帝死了.
瑪丁路德很怒氣.
馬上回應:胡說.
上帝哪有死?.
上帝是永活的.
怎會死呢?.
他太太就知道.
瑪丁路德上吊了.
於是就笑著跟他說.
是的.
上帝是永活的.
你自己也懂得說.
既然如此.
他一定會幫你渡過難關.
那你現在又何必這麼愁呢?.
瑪丁路德他恍然大悟.
重拾從上帝而來的信心.

$^{681}$和喜樂.
弟兄姊妹提醒我們.
喜樂的關鍵其實不在於.
環境有沒有改善.
聖經裡強調.
喜樂的關鍵是.
重新跟上帝結連.
你重新連接他.
你就會有喜樂.
正如耶穌也說.
在世上我們會有.
會有什麼呢?.
苦難.
你能不能說阿們呢?.
基督徒應該要說阿們.
當耶穌說.
在世上你們是有苦難的.
你應該要說阿們.
因為耶穌沒有違反.
《商品說明條例》.
說得很清楚.
無論你是非基督徒還是基督徒.
都要面對苦難.
分別就是.
基督徒明白下一句是什麼.
不過你們可以放心.
因為.
誰贏了世界呢?.
不是我們.
耶穌說他搞定了.
所以我們可以放心.
所以弟兄姊妹.
其實真的每一個人.
在世上.
你會碰到各自的難關.
每個人都有他自己.
要遇見的重擔擔子.
或者逆境.
但是喜樂的人總是會找到方法.
來跟上帝結連.

$^{721}$正如如果你有留意哈巴谷先知.
第三章.
剛才我們也有讀.
最後其實是一首什麼呢?.
是一首歌.
一首詩歌.
如果你看回三章的開頭.
其實是一個禱告.
既是禱告.
也是詩歌.
它提醒我們其實.
我們每個星期.
或者甚至每天.
向上帝祈求.
禱告.
唱詩.
敬拜.
其實是彼此提醒.
就算我們身處的環境變得多惡劣.
上帝仍然是掌管著.
就算我們經歷多麼嚴峻的局面.
上帝仍然愛我們.
祂會賜下平安.
詩歌很多時候是這樣用的.
你的禱告是這樣用的.
事實上你會留意到.
聖經裡經常強調.
平安喜樂就像一對.
屬靈的孿生兄弟一樣.
平安是我們知道上帝掌管.
所以即使在風浪尖裡.
你仍然可以安穩.
而喜樂就是進一步.
既然如此.
你就享受在風浪中.
上帝與你的同在.
然後怎樣?.
繼續事奉祂.
是這樣的意思.
不是拼在一起.

$^{761}$平安喜樂加在我們身上.
是要讓我們在困難裡.
繼續為上帝作見證.
事實上我聽過神學家尼布爾.
他有一段很著名的圖文.
他就說.
神啊.
請賜給我安寧.
去接受那些.
我無法改變的事情.
賜給我勇氣.
去改變那些我能夠改變的事情.
也讓我智慧.
能夠懂得分辨.
這兩者的不同.
這些寺門事實上.
都可以成為我們一個提醒.
有些事情可能.
上帝是想你去改變.
你放膽去.
有些事情上帝定了.
祂想你安靜接受.
但是怎樣去分辨.
求聖靈幫助帶領我們.
讓我們懂得怎樣去回應.
事實上你會看到.
哈巴谷先知既然知道.
上帝旨意已決.
他就不再試圖.
花自己的心思意念.
去改變他將要面對的困境.
他亦沒有嘗試去迴避.
苦難的臨到.
又或者將責任推卸給別人.
他選擇什麼呢?.
經文裡清晰說.
他安靜等候.
巴比倫的臨到.
他坦然接受.
上帝審判行動即將出現.

$^{801}$不過.
我們要搞清楚.
安靜等候這四個字.
並不代表我們什麼都不做.
翹起雙手.
繼續過自己的生活.
只等.
你看到哈巴谷先知.
他會有些堅持.
特別有兩方面的堅持.
第一個就是.
戰驚.
他不是純粹說.
哎呀!糟了.
即將有禍患臨到.
國破家亡.
所以怕得震驚.
他這裡說戰驚.
不是純粹說因為災難本身.
戰驚是因為自己.
都要和其他人一樣.
他都要經歷上帝的管教.
和審判.
無論是巴比倫.
無論是猶大.
甚至乎先知本人.
在那個情況裡.
在當時的裡邊.
都要面對上帝同一個法則.
是什麼?.
如果你是行惡惡人.
最終會滅亡.
如果你是義人.
因為對上帝的忠信.
最終會得救.
先知都要面對這件事.
換句話來說.
戰驚.
警醒就能夠幫助先知.
在惡人災難當道的日子.

$^{841}$他怎樣?.
他不會亂來.
他不會自己因為怒氣.
而自己去以惡還惡.
找那些惡人來報仇.
找他們去悔氣.
他不會因為這樣.
而令自己也變成惡人.
離棄了上帝.
而是他將出手的主權.
管教的主權.
交回給上帝.
等上帝出手.
這是第一樣.
確保自己警醒.
不會在一個亂局裡.
自己不知不覺都變成惡人.
第二樣是確保在災難臨到的日子.
自己可以穩站在.
上帝公義的道路上.
換句話來說.
他進一步.
不單止自己不行惡.
更進一步要行上帝的義.
他要堅守.
上帝已經給予他.
交託給他的使命.
第十九節說.
主耶和華是我的力量.
他使我的腳快如舞碌的蹄.
又使我穩行在高處.
原本的意思是.
直譯是.
上帝使他.
好像舞碌的腳.
可以讓他穩行在不同的高處.
其實是指焦點不是站穩.
不是,對不起.
是指焦點是站穩.
而不是快速.

$^{881}$重點是站穩.
不是快速.
不是說上帝一定會保守我.
在困難裡.
所有人都出事.
我不出事.
上帝不是保守我.
在困難裡.
我被人逃離現場.
快一步.
趨吉避凶.
不是這個意思.
如果你有時看.
譬如BBC也好.
或者國家地理雜誌.
有時拍攝動物.
鹿.
牠們會在懸崖峭壁裡.
牠們會怎樣呢?.
牠們可以站穩.
一點點距離牠們都可以站穩.
是嗎?.
意思就是這樣.
就是說.
才知道他很明白.
上帝在困難裡.
神是會以力量促邀的.
使牠的道路仍然可以順傳.
好像舞碌一樣.
可以穩行在那些高處.
那些山崖峭壁裡.
為的是什麼?.
繼續完成上帝的吩咐.
站穩是代表要繼續.
完成上帝的吩咐.
不是自己人生的安全.
不是那回事.
換句話來說.
站穩在高處的焦點.
不是說先知一定得到.

$^{921}$上帝各樣的保護.
賜福不是純粹這麼簡單.
而是我在亂局裡.
我可以穩守上帝託付給我的使命.
所以你看到.
先知為什麼會起落.
是因為他在一個困難的環境裡.
他全力投靠上帝.
全力依靠上帝.
包括進行上帝的吩咐.
更加謹守上帝的使命.
所以.
一個所有人都看為.
最凶險,最危險的地方.
他說是最穩妥的地方.
他說是最喜樂的地方.
沒錯.
雖然我沒有能力去搞定眼前的局面.
但是.
上帝有能力.
並且上帝允許了.
最終他會將一切搞定.
所以.
先知就將他的目光.
轉移回在上帝已經託付給他的使命當中.
他要拔喻這些事情.
他說這個才是正路.
這樣.
就會經歷到.
從上帝而來的喜樂.
所以頂展回到我們所身處的環境.
活在一個罪惡世界裡.
其實無論你是一個基督徒.
或者你是非基督徒.
我們每天都一定會面對一些不合理.
甚至乎不公平的事情.
但是這段經文今天挑戰我們.
我們將目光放在哪裡?.
純粹放在一些惡人,惡事身上.
被他們卡住了.

$^{961}$以致我們沒有辦法會出上帝.
喜樂的旨意.
還是像哈巴菊先知勉勵我們一樣.
將你的目光.
由惡人,惡事.
轉移回到上帝身上.
他邀請我們看看.
上帝的大能是什麼.
看回上帝.
在歷史的作為是怎樣.
他怎樣一一撥亂反正.
看回上帝在聖經裡.
給了我們什麼應許.
看回上帝是怎樣愛我們.
他是怎樣已經拯救我們.
而這一切再不是你頭腦裡.
純粹一個觀念.
而是變回.
當你遇到困難,危難的時候.
那個是你的支撐位.
以致即使你見到局面崩壞到.
就算像哈巴菊先知那樣.
但你仍然可以對上帝有信心來依靠.
因為你開始會找回.
上帝做你的目的是什麼.
上帝做你的意義是為了什麼.
這裡看到.
雖然無花果樹,葡萄樹,橄欖樹.
田地牛羊,牛圈羊棚都不效力.
但會不會動搖.
上帝和你的關係.
會不會動搖你和上帝的關係.
其實哈巴菊正在問這條問題.
當你的經濟.
當你的生計,你的前景.
都受到阻滯的時候.
會不會動搖你和上帝的關係.
應該是不會的.
並且上帝仍然會抓緊我們.
因為我們不是為無花果樹,葡萄樹,橄欖樹,牛羊,田地.

$^{1001}$不是為生計,經濟,前途而活的.
我們被造由始至終是為誰而活.
為上帝而活.
換句話來說.
我們是以上帝的榮耀.
祂的旨意有沒有成就.
上帝在我們身上是否獲得滿足.
以這一個為我們人生的喜樂.
這是哈巴菊才知道.
她持守著一個信仰真理.
亦盼望我們可以在困難的裡面效法祂.
所以就算耶路撒冷.
真的崩壞.
甚至要被擄到外邦生活.
她仍然可以繼續去服侍神.
仍然可以繼續為上帝做見證.
活出上帝的吩咐.
經歷到上帝和她同在的那種喜樂.
沒有人能夠奪去的喜樂.
所以頂子妹這裡經文提醒.
我們喜樂由始至終不是說.
我們解決一切眼前的問題.
如果是這樣我們應該永遠都找不到.
因為每天都會有一些事情.
你是掌握不到.
你是無能為力.
但是好消息是.
上帝有能力.
上帝已經拯救了你.
上帝並且說一切最終會撥亂反正.
所以這幅音是提醒我們.
不只是得救.
是找回我們所有的東西.
是找回上帝做我們的目的和意義.
就是認定我們由始至終.
不是為了經濟生計前途而活.
是為了上帝而活.
事實上.
真正喜樂的人.
他會很靈巧在這方面.

$^{1041}$當他面對逆境的時候.
他會盡快把自己的心思意念.
回到上帝的旨意身上.
特別當一切生活的情況.
不像自己預期的時候.
特別當此路不通的時候.
他懂得像哈巴谷先知一樣.
換個角度去想.
我如何繼續事奉上帝.
為上帝而活.
就正如他可以在困難的裡面.
他仍然找到他有使命.
他仍然可以透過他的宣講.
透過哈巴谷書.
去安慰服侍他當代的同胞.
甚至到現在.
上帝也在使用他.
在安慰我們.
來完成上帝給他的歷史任務.
同樣看回我們今天.
其實我們也可以善用.
上帝仍然給我們的資源.
給我們的機會.
去活出那個喜樂的生命.
如果當我們把眼目.
放回上帝身上的時候.
特別在疫情裡面.
我想很多事情.
我們都會變得不順暢.
很多的掣肘.
但是我們能否同時看到.
上帝已經預備了.
一些科幻機會.
我們能否看到?.
特別如果你有留意.
其實崇拜每次.
可能這一兩年.
不過不過每次都有新朋友.
今早又有四個.
到底怎樣上帝為我們預備?.

$^{1081}$我們能否看到?.
在自己很多困難的處境裡面.
能否將人生的焦點.
撥回去.
放回上帝的身上.
看見上帝.
想我們怎樣去回應他?.
並且將其他.
你應付不來的事情.
安心交託給上帝.
特別你會看到.
很多人因為這一兩年.
對環境,對前途.
失去了昔日那種掌握的時候.
正正是他們需要福音的時候.
正正是他們的心.
變得柔和,柔軟的時候.
亦是教會,福音工作.
一個殺種,灌溉.
甚至收割的日子.
所以結束的時候.
我想用耶穌基督一段比喻.
來分享.
大家還記得.
《路加福音》第十五章.
耶穌不是一口氣說了三個比喻.
記不記得?.
第一個是.
失錢.
失去錢.
賺回.
第二個呢.
不是第一個.
失陽.
賺回.
第二個呢?.
失錢.
第三個呢?.
浪子回頭.
一連說了三個比喻.

$^{1121}$而這三個比喻.
我發覺都是以一個什麼狀況結尾.
賺回.
失而復得.
喜樂.
看到耶穌對我們的心意嗎?.
原來這個世界上真的有些喜樂.
不受外在環境影響.
而是真真正正.
和我們有沒有履行.
耶穌的吩咐有關.
就是領人歸主.
上帝應許了.
當我們參與在裡面的時候.
當我們有份回應.
福音使命的時候.
我們必然和上帝一同歡呼喜樂.
也唯有我們真的去嘗試.
我們真的踏出這一步.
我們就會體驗領受到.
上帝這個真真實實的應許.
弟兄姊妹你願意嗎?.
讓我們同心禱告.
阿爸天父我們多謝你.
因為你已經拯救我們.
讓我們出死入生.
讓我們今天真的能夠再去思考.
到底你所說.
你的吩咐喜樂是怎麼一回事.
或許過往.
我們有些人.
調轉了關係.
我們總是覺得.
我們要完成自己人生一些需要.
目標滿足.
然後我們才有那種喜樂.
以致我們真的不能夠.
經歷到上帝你所說的.
那種喜樂的生命.
但上帝多謝你.

$^{1161}$藉著今天的經文提醒我們.
我們要先學習和你連結.
從你支取那份平安喜樂的時候.
我們才更有力量.
去和逆境共舞.
也真的去對準.
上帝你做我們的使命和意義.
上帝幫助我們.
縱然可能我們身邊有很多事物.
未必會為我們效力.
未必順著我們的意願來發生.
但你應許過這一切.
都不能夠動搖你和我們的關係.
不能夠動搖你給我們的喜樂.
那份屬天的喜樂.
沒有人能夠奪去.
沒有任何環境能夠影響.
是因為當我們真的找回.
你創造我們的意義的時候.
當我們真的去回應.
耶穌基督你大使命的吩咐的時候.
我們也要同時去經歷.
上帝你的應許.
就是那份領人歸居的喜樂.
上帝幫助我們.
讓我們真的在.
同義的日子裡.
我們更加的謹靠你.
我們更加去把握機會.
來回應你的呼召.
來榮耀你.
來造就身邊的人.
來和上帝你一起歡喜快樂.
求你幫助我們.
預立我們在你面前的禱告.
奉耶穌說你得勝的命之.
來祈求.
阿們.
\newpage



\section{哥林多後書 10:1-18}
\label{sec:RH7uL12ghP4}
\textbf{2021年5月23日崇拜直播｜梁家麟牧師｜保羅的勇敢和怯懦｜哥林多後書10\_1-18}
\newline
\newline
連結: \href{https://youtube.com/watch?v=RH7uL12ghP4}{\texttt{https://youtube.com/watch?v=RH7uL12ghP4}} ~~~~ 語音日期: 2021-05-25
\newline
\newline
\hyperref[sec:d1qGfZ_AH3c]{\small{< < < PREV SERMON < < <}}
~
\hyperref[sec:index_chronic]{\small{[返順時目]}}
~
\hyperref[sec:Czt7dmGniVs]{\small{> > > NEXT SERMON > > >}}
\newline
\newline
哥林多後書 10:1-18
\newline
\begin{longtable}{cl}
\hline
\hline
章節 & 經文 (和合本修訂版)\\
\hline
10:1 & \begin{tabularx}{0.7\textwidth}{X} 我－保羅與你們見面的時候是溫和的，不在你們那裡的時候向你們是勇敢的，如今親自藉著基督的溫柔和慈祥勸你們。 \end{tabularx} \\ \\ \relax
10:2 & \begin{tabularx}{0.7\textwidth}{X} 有人認為我們是憑著血氣行事，我認為必須敢於對付這等人；我但求在那裡的時候，不必這樣勇敢。 \end{tabularx} \\ \\ \relax
10:3 & \begin{tabularx}{0.7\textwidth}{X} 我們雖然在血氣中行事，卻不憑著血氣爭戰。 \end{tabularx} \\ \\ \relax
10:4 & \begin{tabularx}{0.7\textwidth}{X} 因為我們爭戰的兵器本不是屬血氣的，而是憑著神的能力，能夠攻破堅固的營壘。我們攻破各樣的計謀， \end{tabularx} \\ \\ \relax
10:5 & \begin{tabularx}{0.7\textwidth}{X} 和各樣攔阻人認識神的高壘，又奪回人心來順服基督。 \end{tabularx} \\ \\ \relax
10:6 & \begin{tabularx}{0.7\textwidth}{X} 我已經預備好了，等你們完全順服的時候來懲罰所有不順服的人。 \end{tabularx} \\ \\ \relax
10:7 & \begin{tabularx}{0.7\textwidth}{X} 你們只看事情的外表。倘若有人自信是屬基督的，他要再想想，他屬基督，我們也屬基督。 \end{tabularx} \\ \\ \relax
10:8 & \begin{tabularx}{0.7\textwidth}{X} 主賜給我們權柄，是要造就你們，並不是要拆毀你們；我就是為這權柄稍微誇口也不覺得慚愧。 \end{tabularx} \\ \\ \relax
10:9 & \begin{tabularx}{0.7\textwidth}{X} 我說這話，免得你們以為我寫信是要恐嚇你們。 \end{tabularx} \\ \\ \relax
10:10 & \begin{tabularx}{0.7\textwidth}{X} 因為有人說：「他信上的語氣既嚴厲又強硬，他本人卻軟弱無能，言語粗俗。」 \end{tabularx} \\ \\ \relax
10:11 & \begin{tabularx}{0.7\textwidth}{X} 這等人當明白，我們不在那裡時信上怎麼說，見面時也必怎麼做。 \end{tabularx} \\ \\ \relax
10:12 & \begin{tabularx}{0.7\textwidth}{X} 因為我們不敢將自己和某些自我推薦的人並列相比；他們用自己度量自己，用自己比較自己，是不明智的。 \end{tabularx} \\ \\ \relax
10:13 & \begin{tabularx}{0.7\textwidth}{X} 我們不願意過分誇口，但是我們只在神劃定的界限內誇口。這界限甚至擴展到你們那裡。 \end{tabularx} \\ \\ \relax
10:14 & \begin{tabularx}{0.7\textwidth}{X} 我們擴展到你們那裡時並沒有越過了自己的界限，其實我們是首先到你們那裡傳基督福音的。 \end{tabularx} \\ \\ \relax
10:15 & \begin{tabularx}{0.7\textwidth}{X} 我們不靠別人所勞碌的過分誇口；我們只希望你們信心增長的時候，所劃定給我們的範圍也能夠因著你們更加擴展， \end{tabularx} \\ \\ \relax
10:16 & \begin{tabularx}{0.7\textwidth}{X} 使福音得以傳到你們以外的地方，而不在別人的範圍之內，以別人所成就的事誇口。 \end{tabularx} \\ \\ \relax
10:17 & \begin{tabularx}{0.7\textwidth}{X} 但「要誇耀的，該誇耀主」。 \end{tabularx} \\ \\ \relax
10:18 & \begin{tabularx}{0.7\textwidth}{X} 因為蒙悅納的，不是自我稱許的，而是主所稱許的。 \end{tabularx} \\ \\
[1ex]
\hline
\hline
\end{longtable}
$^{1}$接下來有林永祥大兄.
為我們宣讀一段經文.
經文亦印在電視機的內部.
也是呼應今天的講道主題.
「保羅的勇敢和怯懦」.
由梁家騮牧師考慮當中來正讀.
先把時間交給林永祥大兄.
(教育局局長).
今天要重讀的經文.
是哥林多後書十章一至十八節.
請大家打開聖經新著部分.
第252頁.
哥林多後書一至十八節.
「我保羅就是與你們見面的時候是謙卑的.
不在你們那裡的時候.
向你們是勇敢的.
如今親自藉著基督的溫柔和平勸你們.
有人以為我憑著血氣行事.
我也以為必須用勇敢待這等人.
求你們不要叫我在你們那裡的時候.
有這樣的勇敢.
因為我們雖然在血氣中行事.
卻不憑著血氣征戰.
我們征戰的兵器本不是屬血氣的.
乃是在神面前有能力.
可以攻破堅固的形裡.
將各樣的計謀.
各樣難阻人認識神.
的那些至高之事一概攻破了.
又將人所有的心意奪回.
使他都順服基督.
並且我已經預備好了.
等你們十分順服的時候.
要責罰那一切不順服的人.
你們是看眼前的嗎?.
倘若有人自信是屬基督的.
他要再想想他如何屬基督.
我們也是如何屬基督的.
主賜給我們權柄.
是要造就你們.

$^{41}$並不是要敗壞你們.
我就是為這種權柄稍微誇口.
也不至於慚愧.
我說者說.
免得你們以為我寫信要威嚇你們.
因為有人說.
他的信又沉重又厲害.
及至見面卻是氣貌不揚.
言語躁俗的.
這等人當想.
我們不在那裡的時候.
信上的言語如何.
見面的時候行事也必如何.
因為我們不敢將自己.
和那自薦的人同列相比.
他們用自己度量自己.
用自己比較自己.
乃是不通達的.
我們不願意分外誇口.
只要照神所量給我們的界限.
他扣到你們那裡.
我們並非過了自己的界限.
好像扣不到你們那裡.
因為我們早到你們那裡.
傳了基督的福音.
我們不像著別人所牢固的.
分外誇口.
但指望你們信心增長的時候.
所量給我們的界限.
就可以因著你們更加開展.
得以將福音傳到你們以外的地方.
並不是在別人界限之內.
即著他現成的是誇口.
但誇口的當指著主誇口.
因為忘越立的.
不是自己稱許的.
乃是主所稱許的.
聖經獨筆.
主持人:鄭若驊.
Margaret Elizabeth Barber 她的一句名言.

$^{81}$為己無所求,為主求一切.
I want nothing for myself, I want everything for the Lord.
對於一些針對她個人的批評.
例如她的樣子,氣質,做事手法.
我放在心裡,連自辯的動機也懶得考慮.
反正她對自己的賣相,公共形象也沒有什麼期望.
那有什麼所謂?.
別人喜歡與否不會為她增加或減少任何東西.
但是對她所傳揚的福音真理.
對信徒是否切實地接受和進行她的教訓.
我非常在乎,堅持到底,決不退讓.
保祿的勇敢和怯懦.
充分反映了教會事公上一個很重要的原則.
對事不對人.
我們在應用上就說對事不對人的原則.
我把它變成兩小點來說.
第一,我們秉持對事不對人的原則.
在所有人事行政的處理上.
在事公策略的抉擇上.
我們只關注福音事公的推展.
弟兄姊妹的益處.
教會未來的發展.
我們執事會,各部會只是關心這幾點.
弟兄姊妹有什麼益處?.
對教會發展有沒有幫助?.
對福音事公有沒有推展?.
我們只考慮這一點.
不要夾前太多個人的得失成敗的考量.
尤其切忌不要重視太多個人的面子和感受.
你知道教會出現人事糾紛.
甚至出現教會政治,church politics.
甚至鬧分裂.
很少因著聖經真理或福音策略路向.
這些偉大的大是大非的問題.
多數是夾雜了個人的義氣之爭.
誰和誰有積怨.
誰對誰有成見.
對方不尊重我.
我的面子受損.
或者我們總是堅持糾纏.

$^{121}$事情當然要由我決定,我拍板.
為什麼他代我做拍板?.
誰決定比決定什麼更重要?.
再糟糕的.
我們學不學.
學了外面的政治遊戲,政治策略.
專事破壞.
計算自己堅持的立場不會獲得通過.
專注讓對手的提案都不獲得通過.
一拍兩散.
寧願癱瘓我組織.
我們常常為了固存自己的面子.
而不固存教會的大局.
在那些孜孜孜孜的問題上.
堅持到底,絕不妥協.
即教會事公因此受損.
亦在所不惜.
固存面子不固存大局.
面子比大局更重要.
教會叫了很多無聊的交.
浪費資源,浪費福音機會.
造成很壞的見證.
是不是很可惜的事?.
我們往朕.
就我所知,是不存在這些情況的.
不過,你不要告訴我.
在香港教會這樣的情況不出現.
這是值得我們彼此互相勉勵提醒.
我常常說,現在我講到你們耳朵都起繭.
基督徒的自我,不可以太強.
不要太有自覺,不要太有自覺性.
太有自覺的人.
上上就會在兩極兜來兜去.
一是自大,一是自卑,自憐.
Self-Gory, Self-Pity.
很糟糕.
保羅堅持對事不對人的原則.
不理會,不care(在乎)自己的得失.
只關心教會的福祉.
我們大家勉勵.

$^{161}$我們要學這個原則.
同樣,第二方面.
對事不對人的原則也提醒我們.
別人針對我們,或者討論事情的時候.
對自己不要太個人.
他是不是在插我,他是不是在壓我.
不要想太多這些.
但是對人都不要太個人.
我和同工商討教會的事.
就算各有不同的意見.
我認定我們大家都是為教會的好處.
所以不要夾前個人的喜惡愛憎.
不要特別憎那些和我的意見不同的同工.
珍珍搏斥過我.
所以我就憎了她.
傾事不是後面的恩怨.
不要將對手看為敵人.
不是的,你沒有對手.
同工來的.
除非對方已經偏離了真理.
說了一些嚴重傷害教會.
做了一些嚴重傷害信徒的事.
否則,我始終認定她是我的弟兄,我的姊妹.
不要輕易產生芥蒂,建立成見.
甚至反目成仇.
我們在旺鎮服侍.
你們有很多樣化的選擇.
不喜歡梁家倫的人.
一定可以喜歡陳詠全.
不喜歡梁家倫沒有頭髮的人.
可以喜歡Ray 沐,很多頭髮.
所以你不喜歡我是沒有問題的.
只要你不要因此不喜歡旺鎮就行了.
你繼續喜歡旺鎮.
你不喜歡一個人也無所謂.
喜歡另一個.
大家很快樂的.
保羅對提摩太太.
在提摩太太後書三章十四到十五節的提醒.
我覺得很有意思,我不會詳細解釋.

$^{201}$不過我讀一讀.
提摩太太後書三章十四到十五節.
她說:若有人不聽從我們者信上的話,要記下他.
不和他交往,叫他自學羞愧.
但不要以他為仇人,要勸他為弟兄.
這真是對事不對人.
在真理上絕不妥協.
但在身份,在關係上卻不切斷.
千萬不要做不成同工就連兄弟也做不到.
這是最糟糕的事.
本來大家是很老死的.
進了執事會之後.
然後就反目成仇.
以後都不會再談話.
大不大濟.
這真是很糟糕.
對旺鎮的弟兄姊妹.
剛剛被主接去的志強弟兄.
其實是為我們留下一個很美麗的事奉榜樣.
我進了旺鎮很遲.
我能夠再記錄多一個.
有人告訴我還有更多.
但我不知道,我沒有見過.
馬美清姐妹.
你覺不覺得她們都是相同性格的人.
對主效忠.
但對人很溫柔.
我覺得旺鎮真是很幸福.
為己無所求.
為主求一切.
我們有一些很好的淑寧前輩.
為我們留下皆未腳踵.
在座當中都有這樣的弟兄姊妹.
很多的.
不過我們歌頌死人不要歌頌生人.
歌頌生人有幾好意思.
我不是說你不好.
不過我不說你.
你安息你拜祭我才說.
不說不說.

$^{241}$仍然在生的.
大家努力.
我們還要繼續努力.
我們不說的.
所以我們有很好的淑寧傳統.
求主幫助我們繼續這個淑寧傳統.
旺鎮會繼續走在上帝的旨意裡.
上帝賜福你們每一個.
\newpage



\section{以弗所書 4:26-32}
\label{sec:Czt7dmGniVs}
\textbf{2021年5月30日崇拜直播｜吳真真牧師｜從聖經看喜怒哀愁 - 怒亦有道｜以弗所書4\_26-32}
\newline
\newline
連結: \href{https://youtube.com/watch?v=Czt7dmGniVs}{\texttt{https://youtube.com/watch?v=Czt7dmGniVs}} ~~~~ 語音日期: 2021-06-05
\newline
\newline
\hyperref[sec:RH7uL12ghP4]{\small{< < < PREV SERMON < < <}}
~
\hyperref[sec:index_chronic]{\small{[返順時目]}}
~
\hyperref[sec:9JRB4RUnJVQ]{\small{> > > NEXT SERMON > > >}}
\newline
\newline
以弗所書 4:26-32
\newline
\begin{longtable}{cl}
\hline
\hline
章節 & 經文 (和合本修訂版)\\
\hline
4:26 & \begin{tabularx}{0.7\textwidth}{X} 即使生氣也不要犯罪；不可含怒到日落， \end{tabularx} \\ \\ \relax
4:27 & \begin{tabularx}{0.7\textwidth}{X} 不可給魔鬼留地步。 \end{tabularx} \\ \\ \relax
4:28 & \begin{tabularx}{0.7\textwidth}{X} 偷竊的，不要再偷；總要勤勞，親手做正當的事，這樣才可以把自己有的，分給有缺乏的人。 \end{tabularx} \\ \\ \relax
4:29 & \begin{tabularx}{0.7\textwidth}{X} 一句壞話也不可出口，只要隨著需要說造就人的好話，讓聽見的人得益處。 \end{tabularx} \\ \\ \relax
4:30 & \begin{tabularx}{0.7\textwidth}{X} 不要使神的聖靈擔憂，你們原是受了他的印記，等候得救贖的日子來到。 \end{tabularx} \\ \\ \relax
4:31 & \begin{tabularx}{0.7\textwidth}{X} 一切苦毒、憤怒、惱恨、嚷鬧、毀謗，和一切的惡毒都要從你們中間除掉。 \end{tabularx} \\ \\ \relax
4:32 & \begin{tabularx}{0.7\textwidth}{X} 要仁慈相待，存憐憫的心，彼此饒恕，正如神在基督裡饒恕了你們一樣。 \end{tabularx} \\ \\
[1ex]
\hline
\hline
\end{longtable}
$^{1}$經訓記載在二福所書四章二十六至二十七節.
和三十一至三十二節.
在新約聖經二百六十八頁.
程序表內頁已經有記載.
生氣卻不要犯罪.
不可含怒道日落.
也不可給魔鬼留地步.
三十一節.
一切苦毒.
怒恨,憤怒,讓罵,毀謗.
並一切的惡毒.
都當從你們中間除掉.
並要以恩慈相待.
全面吻的心彼此饒恕.
正如神在基督裡饒恕了你們一樣.
(自動翻譯).
各位弟兄姊妹平安.
很久沒有在主日崇拜和大家分享上帝的道.
大家有否看到我今天有什麼新形象?.
有呀,今天有什麼新形象呢?.
束起了頭髮,是嗎?.
大家在禮堂裡知道.
這裡的熱氣比較上升.
所以要束起頭髮.
讓自己清爽一點.
但也有另一個原因.
因為很多時候.
我之前在崇拜講道.
也會把我們的崇拜放在YouTube.
我崇拜完之後.
回家就會收到一些人把YouTube的連結發給我.
通常那裡會有我的大頭像.
大頭像就會說.
「某某在講道」.
那「某某」是誰呢?.
「某某」是現在很紅的節目監製.
和一對很受歡迎的男子組合的經理人.
我收過不止一次.
我今天刻意改變了形象.
但為什麼我會這樣呢?.

$^{41}$有些人問我「你生氣嗎?」.
正好,今天我們就要從聖經裡去看.
究竟生氣有什麼教導.
有些人說.
情緒是這個世界的共通語言.
我們人的面部有44條肌肉.
如果這些肌肉產生不同的變化.
和協調,運作.
就會製造出不同的表情.
很奇妙的.
你做同一個表情.
來自世界不同文化,不同地方.
講不同語言的人.
都會說「你是開心」.
「你是很厭惡」.
「你很憤怒」.
是不是很特別?.
是不是很奇妙?.
原來人與生俱來有情緒.
而有一些研究告訴我們.
情緒與生俱來有六種.
大家都很熟悉.
喜,怒,哀,驚.
還有厭惡.
厭惡包括內疚.
第六樣是驚訝.
有六種表情.
而這些情緒都有其功能.
會因應不同情緒而產生.
激發不同的反應.
而這些反應.
對於人與人之間互動的關係.
有著很深遠和重要的影響.
舉個例子.
例如你很開心.
原來這種開心會激發我們的反應.
就是讓我們很想跟別人接近.
讓我們很想跟別人連結.
但當我們害怕的時候.
可能我們跟別人之間的影響.

$^{81}$就是我們會呆了.
或者可能馬上逃離現場.
不過亦有可能我們會反擊.
而憤怒的情緒.
憤怒的情緒會讓我們傾向.
有一些自我保護.
但另一方面亦會是一些行動.
會讓我們採取行動.
甚至令我們勃起了.
振奮了,要去行動.
而情緒有分正面和負面.
我想問大家.
你喜歡正面還是負面情緒?.
正面吧.
人人都喜歡正面情緒.
人傾向逃避負面情緒.
而且當有負面情緒的時候.
就需要用很大的力氣.
花很多心機.
才能夠去應付和處理.
大家有沒有留意到.
我們在五月和六月.
有四個講題.
講題的系列是.
從聖經的角度看起.
怒哀愁.
今天我會講負面情緒.
憤怒.
今天我會用.
《耳窟所書》四章26至27節.
和31至32節.
去學習了解.
上帝對於我們有憤怒情緒的時候.
有什麼教導.
不知道大家每一次聽到.
第26節這樣說.
生氣卻不要犯罪.
不要含怒到日落.
你會想到什麼呢?.
我告訴你.

$^{121}$我每次都想起這兩個弟兄.
哪兩個弟兄呢?.
有一個叫陳弟兄.
有一個叫羅弟兄.
陳弟兄和羅弟兄說.
我們查聖經.
得出一個結論.
我們兩個人協議了.
按照聖經的教導.
不可含怒到日落.
我們就答應了彼此.
以後兩夫妻如果有什麼爭執.
就一定要在睡前.
要傾清清楚.
如果萬一.
睡前談不攏.
就先不睡.
直至談妥為止.
羅弟兄很欣賞.
不是羅牧師.
其實是羅弟兄.
不好意思.
我沒有拍攝.
羅弟兄就說.
哇!這麼好.
你們這麼用心應用聖經.
好!我想問你們實踐得如何.
陳弟兄就說.
嗯!實踐得很好.
我們做得非常好.
他就問.
可否舉些例子來聽聽.
例如三天前.
我和太太有爭執.
因為我們有這個協議.
所以我們在三天內沒有睡覺.
這是笑話一則.
但其實是真的.
當我們讀聖經的時候.
我們實踐聖經的教導.

$^{161}$是否即是說.
我們一發怒.
就要在睡前處理.
而沒有細究.
如何處理和面對.
所以我們今天也很想.
從已忽所書這段聖經.
學習一下.
上帝對於憤怒有何教導.
讓我們一起祈禱.
因為我們不是機械人.
我們是有血有肉的人.
你讓我們每一個人.
都有自己很豐富的情感.
讓我們從這些情感中.
有互動和結連.
亦可以認識自己的限制.
和作出回應的互動和需要.
上帝大,我們知道.
當我們成為你的跟從者.
成為基督徒的時候.
我們就過著新的樣式.
新的生活.
我們很需要學習.
究竟在你的真理和真道裡.
我們如何面對憤怒.
如何在人生中.
活出一個真是討你喜悅的.
行出你的心意的生命.
我們就恭敬將這段時間.
交託在你的手中.
求天父你親自帶領.
多謝你聽祈禱.
奉主耶穌基督得勝.
明知其求.
阿們.
已忽所書第四章.
一開始就說到.
一些從前未信主的人.
如今是跟從主的信徒.

$^{201}$他說你們的行事為人.
應當與蒙召的恩相稱.
我們信了主之後.
我們不是自己了.
我們再不是單打獨鬥.
孤軍作戰.
我們是基督的身體.
我們有一位主.
祂要讓我們在生命裡.
與我們同行.
也願意跟從祂的人.
能夠在生命裡結連和合一.
所以我們應該.
必須除去舊人的生活模式和態度.
不再去順從迷惑人的私慾.
我們現在應該是穿上新人.
過一個新的生活.
而這個新的生活.
就是從真理而來的人意.
和聖潔的生活.
現在保羅就是告訴信徒.
各位弟兄姊妹.
我們有新的身份.
我們有新的生活.
到底我們的生活應該是怎樣呢?.
我以前會憤怒.
我現在可不可以憤怒?.
我以前的憤怒.
是會任憑自己的情緒.
心意,想法去反應.
現在呢?.
現在我應該怎樣處理?.
以我自己為例.
我想跟大家說.
其實我小時候經常被人說小器.
小器.
怎樣小器呢?.
我經常覺得.
人們說一點點話.
都會令自己覺得很冒犯.

$^{241}$然後心裡就很生氣.
我記得有一次.
很深刻的.
我小時候不吃魚.
有一頓飯.
我媽媽就擺了五碟菜出來.
這五碟菜裡面有三碟是魚.
結果我整頓飯就很生氣.
然後我整頓飯都是焦口焦面.
還有一件菜都不合適.
只是白飯.
然後回到房間裡.
還是很生氣.
氣難一下.
於是我就寫日記.
我就將我媽媽的惡行記下.
然後就去以洩我心頭的不忿.
但現在我不會了.
但信了耶穌之後.
我又不知道為什麼.
現在我經常都不敢生氣.
我不敢生氣.
每件事都是有愛.
有包容.
這些只是無心之失.
就算他放了我一百次飛機.
也好.
也有他的原因.
我原諒他.
所謂既不敢怒也不敢言.
說到憤怒.
好像不知道該怎麼說.
古希臘哲學家亞歷斯多德這樣說.
任何人都會憤怒.
但要對正確的人.
以正確的程度.
在正確的時間.
懷著正確的動機.
以正確的方法去生氣.
是不是很難.

$^{281}$我看到有些弟兄姊妹在笑.
我心想.
你生氣這件事也要這麼優雅.
為什麼要這麼優雅.
但不要緊.
聖經裡面有教導我們.
有告訴我們.
究竟憤怒是一回什麼事.
我們應該怎樣去回應.
我們來看看第26節.
第26節說到生氣卻不要犯罪.
這就像一個箴言式的諺語.
有兩個動詞.
一方面是肯定的動詞.
生氣.
另一方面是否定的動詞.
不犯罪.
但這兩個動詞都是命令式.
肯定的命令是你要生氣.
否定的命令是不犯罪.
而生氣這個字.
在新約裡面出現過八次.
除了在啟示錄裡說到.
龍和愛邦的憤怒之外.
到這裡就是說人.
其他都是說人的憤怒.
其實這裡說生氣是一個被動詞.
是指外界的一些刺激元.
激使我們有情緒.
激使我們的憤怒有些情緒.
是無可避免的.
除非我們不和別人交往.
可能身邊沒有東西激怒你.
於是有人就說.
既然是這樣.
我們置身於無可避免的環境裡.
而這個又是命令語氣的生氣.
即是說如果你不生氣.
你就不正常了.
是不是這個意思呢?.

$^{321}$我自己就不太傾向這個解釋.
因為我很難接受到.
如果一個被聖靈充滿的信徒.
他常常都因為外在的環境.
刺激他的時候.
他就常常生氣.
我覺得這個不是真正的解釋.
這不是命令我們生氣.
雖然它是用命令的語氣去說.
我們不需要證明我們正在生存.
或者我們是人所以我們要生氣.
不是這樣的.
所以我是比較認同在文法上.
和合本修訂版的解釋.
它的譯法是.
你即使生氣都不要犯罪.
有些人就會探討.
這個生氣這個生氣.
是不是在說發疑怒呢?.
因為神都會發疑怒的.
在舊約裡面.
七十事易本.
這個生氣這個字出現了79次.
有39次是說上帝會發疑怒.
上帝因為人的犯罪.
上帝因為人的錯誤.
他憤怒.
於是有人就說.
如果這個生氣是否在說.
有學者就認為.
會不會是這裡說的生氣.
是在說發疑怒呢?.
甚至他們會延伸到.
會不會是在說教會裡.
因為要處理一些紀律的問題.
而去談論這件事呢?.
不過因為在《已忽所書》.
沒有其他的論證去證明這件事的時候.
我就會比較傾向.
這個不是《已忽所書》所關心的觀點.

$^{361}$我個人的看法是.
其實生氣是會影響人和人.
和人和神之間的關係.
我們從《已忽所書》四章的上下文來看.
那個背景是說.
如果我們不好好的控制怒氣.
我們就是過著一個舊人的生活模式.
可以帶來極大的影響.
會破壞教會的合一.
破壞人的連結.
所以無論你的生氣是否合理.
是否因為異怒.
又或者是因為有人冒犯了你.
對你不公平.
你對人的期望被那個人藐視.
或者質疑你的價值.
或者你的信念被僭越.
聖經是告訴我們.
生氣是不可以不受控的.
生氣是不可以被駕馭的.
我們個人.
因為怒氣容易令我們犯罪.
我在想難道你生氣合理.
你就可以犯罪嗎?.
這個絕對不對.
所以我們要知道.
接下來的時間.
我們要知道的就是.
憤怒其實可以破壞關係.
我們就要為這件事設限.
接下來的經文就用了三個不.
來告訴我們.
控制怒氣的規範是甚麼.
在說這件事之前.
我想分享一個猶太拉比.
對於憤怒有三個階段的描述.
我們在研究這段經文時.
會明白得更深.
這位猶太拉比叫阿巴拉罕特韋爾斯基.
他是一位精神科醫生.

$^{401}$他告訴我們憤怒有三個階段.
第一個階段叫憤怒.
憤怒是憤怒的感受.
即是說我因為一些刺激.
有些刺激源令我有些情緒出來.
這就是憤怒的感受.
第二步.
當有憤怒的情緒出現時.
我有甚麼回應呢.
他用的字是暴怒.
憤怒的第二階段就是暴怒.
不是穿衣暴怒.
是很暴力,很大暴發雷霆的憤怒.
我們的回應是可以繼續憤怒下去.
是可以極端地憤怒下去.
我們可以罵,打一頓.
第三個階段就是問一個問題.
就是你的怒氣要維持多久呢?.
當我們問這個問題時.
就是在說不能釋懷的怨恨.
這是第三階段.
就是藏在心裡,忿忿不平.
心裡有一種怨恨感.
這就是第三個階段.
先有這個觀念.
我們就進入第一個步.
不犯罪.
如果有人令到我們.
剛才說的第一階段.
有些事刺激了我們.
使我們有憤怒的情緒.
而這個感受.
或者我們可能是一個很本能.
或者從小到大培養出來.
我們可能都會有一個感受.
就是我生氣了.
然後我們第二階段.
那個反應就非常重要了.
聖經告訴我們.
這個反應有一個規範.

$^{441}$就是不犯罪.
不犯罪這個字在希臘原文.
是指當你頭紮一個錨.
或者一支錶.
向著一個方向的時候.
你頭紮不到.
你miss the mark.
這樣就叫犯罪.
如果你不能夠中這個錨.
那你就是犯罪了.
意思就是說.
你不能夠達成上帝給我們的目標.
我們走錯了路.
所以其實犯罪的人.
這裡很想強調一件事.
這個犯罪不是無心之失.
不是說我憤怒.
我不經意的.
人人都會生氣的.
不是.
他很想告訴你.
這個犯罪是有意圖走錯路.
有意識地去錯.
不是無心之失.
我們和神的生氣去比較.
我們知道剛才說神會有生氣.
但是神的生氣是可以受控的.
一個神學作家叫J.I.Packer.
他這樣說.
他說神的生氣永遠不會反覆無常.
不會自我縱容.
不會暴躁.
不會卑鄙.
神的生氣永遠都是受控.
不會過火.
對於不義的惡行.
他會作出正確而必須的反應.
但是人很容易.
很容易被怒氣所控制.
憤怒有時很容易.

$^{481}$令我們失去判斷的理智.
以致我們控制不了自己的思想.
控制不了舌頭.
控制不了我們的行為.
我們去得罪上帝.
得罪人.
我們的憤怒.
不是只是很生氣.
然後我打你一身.
罵你一頓.
這個憤怒才叫犯罪.
我們的憤怒.
可以是當情緒爆發的時候.
我們內心想的事情.
我們之後背後對人的重傷.
或者心裡去想.
我真的翹起雙手.
放長雙眼.
看看你能好多久.
這些都是憤怒的展現.
這些都是犯罪.
《已忽所書》四章二節說.
當我們性靈充滿的時候.
我們可以控制怒氣.
我們可以適時發怒.
我們可以不被情緒操控的情況下.
去作出對對方的對質或責備.
有解經學家提醒我們.
可以用一把尺去衡量.
究竟我們的生氣是否犯罪.
他說如果你的生氣是犯罪的話.
你會發覺你越罵越開心.
越罵越興奮.
你會令自己更加自大.
更加自義.
如果沒有犯罪的生氣.
往往會作出合神心意的回應.
有時你會越罵.
你心裡會有哀傷.
你心裡會有為他而擔心.

$^{521}$有很多不同的哀傷情緒夾雜在其中.
我們知道.
聖經提醒了我們.
第一個規範就是我們不要犯罪.
第二就是我們不要累積怒氣.
我們去到不可含怒到日落.
含怒是一種憤怒的狀態.
直譯是積壓或極化的怒氣.
積壓或極化的怒氣.
是一種很濃很深的憤怒.
就像剛才拉比問的第三個階段.
究竟你的怒氣要憤怒多久.
要持續多久.
然後他就用了日落作為一個指標.
日落不是代表剛才陳弟兄說.
不睡覺代表日落.
日落在猶太人來說有一個很重要的意思.
日落代表一天的完結.
裡面說的是一件很重要的事.
猶太人的日落是說憐憫的公認.
以至讓人避免艱難.
日落原來有這麼美好的意思.
當一天完結的時候.
你身邊的人的缺乏.
他的艱難.
有沒有得到公認.
有沒有人憐憫去公認他.
以至他能夠避免這個艱難呢.
猶太人有個傳統.
他說如果因為彼此的憤怒.
彼此的指控.
而不能夠包容和復和的話.
太陽是不應該下山的.
猶太人是這樣說的.
所以他說不可含怒到日落.
我覺得很有意思.
如果我們未能夠好好控制我們的怒氣.
對於你,對於對方.
其實都是一個艱難.
實在是需要那份憐憫去跨越這份艱難.

$^{561}$不過我經常說.
我發覺其實.
因為剛才提過怒氣是一個負面情緒.
人往往會避開.
我發覺有時不是人故意不去處理怒氣.
而是往往.
我覺得基督徒都不知道自己生氣.
或者不讓自己生氣.
他對於憤怒是採取一個否定或抑壓的心態.
你聽我講就知道.
有一次我約了一班年青人.
我以前來旺鎮做青年牧羊的工作.
我約了一班年青人.
幾個女生說.
「珍珍姐」.
他們叫我珍珍姐.
「珍珍姐,我們接下來考完試一起去唱K」.
大家知道我是一個很正經.
我覺得爭取時間去唱K.
浪費時間的事不是太好.
但沒關係,牧羊年青人.
就陪他們去唱.
於是就畫了一天.
那天打算跟他們去玩.
去癲.
誰知真的氣死你.
一個小時分別接到不同型的電話.
「我不來了」.
只剩下一個.
我們兩個去哪呢?.
晚上回到家.
我對著Paul.
不斷罵.
「這班年青人有甚麼搞錯?」.
「我已經抽了時間」.
這樣說.
說了很久.
他就跟我說.
「你不要生氣」.
「生氣來幹甚麼?」.

$^{601}$「小女孩就是這樣的」.
我很大聲地告訴他.
「我沒有生氣」.
「你這樣也不是生氣?」.
「那是甚麼?」.
我們很怕.
不知為何.
牧羊做人的榜樣.
生氣好像不是一件好事.
我們會否定自己生氣.
但如果我們不肯去面對這份氣.
我們沒有辦法處理.
沒有辦法不累積.
我們不理會.
其實就累積了.
所以.
拉比有個提議.
他說.
晚上請你記下.
如果那天有人惹你生氣.
你心裡有憤怒.
你寫下你做了甚麼回應.
因為通常人家惹我們生氣.
我們很容易做一個本能的反應.
那個本能的反應.
可能都很不理智.
頭腦不清晰.
一些很情感的回應.
但當夜晚.
當那股氣.
頭腦清醒一點.
冷靜下來的時候.
我們就可以運用.
我們本能的知性腦袋的部分.
去反問一下自己.
「我是否生氣?」.
「我這樣的回應方式」.
「處理的方式是否最好?」.
「有沒有其他方式更好呢?」.
以致到.

$^{641}$當我們下一次再有被惹怒的狀況.
我們是否可以用一個更好的方法去回應呢?.
而事實上.
當我有這樣的操練的時候.
在我長大一點.
我又再目仰另一班年青人.
有一次這班年青人.
開會.
臨開會的時候.
就一個個傳送一些訊息出來.
就說.
「我來不了」.
「我今天有點不舒服」.
「OK」.
「糟了,我忘記了」.
「我沒有記下來」.
「OK」.
然後.
「哎呀」.
「我的最好朋友今晚生日」.
「我要跟他慶祝」.
「嘩」.
結果.
一個個申請缺席.
一半人不夠出席.
「你生氣嗎?」.
「咦,你們沒有反應?」.
「不敢生氣嗎?」.
「不敢怒亦不敢言」.
我可以告訴大家.
我生氣的.
我真的生氣.
但我發覺這個生氣是複雜了很多.
這個生氣是代表了我的失望.
這個生氣是代表了我的心痛.
這個生氣是有一點恨鐵不成鋼.
你們是一班職員.
團契是你們帶來的.
你們都可以在生活的優先次序裡.
放得這麼低.

$^{681}$這是什麼事奉態度呢?.
我們會生氣.
我們真的會生氣.
但裡面有很多情緒.
很多情感.
當我冷靜下來的時候.
我就用一個.
沒有這麼生氣的.
一篇電郵.
告訴他們.
我的看法,我的感受.
我相信對他們來說.
是一個造就.
下次開會就沒有人敢知道.
沒有人敢不來.
是真的.
其實我們真的需要.
去處理這份情緒.
並且這個怒氣的情緒.
是夾雜了很多其他東西.
我們需要冷靜下來.
讓自己消化一下.
將自己真正的感受.
跟大家分享.
如果.
我們不處理的話.
就好像.
《已忽所書》四章三十一節說.
一切的苦毒.
怒恨.
讓罵,毀謗.
並一切的惡毒.
都從你們中間除掉.
這裡有.
五個形容詞.
苦毒是代表了.
未能夠放開一直以來.
對一些事情不滿的情緒.
怒恨.
是一些外露的憤怒.

$^{721}$是一些內藏的憤怒.
讓罵就是大聲的.
在那裡吵架.
而不講就是.
去重傷別人,傷害別人的.
名聲.
所以原來.
我們真的.
千萬不要累積.
這些憤怒的情緒.
是會帶來.
像剛才那樣的五個形容詞.
都是圍繞著.
指著當我們沒有好好處理.
情緒的時候,那份憤怒可以怎樣.
彰顯出來.
最後第三個不就是.
魔鬼留地步.
吸這個字.
是命令式.
它說的是.
要讓這個.
持續的情況,即是被魔鬼留地步.
的這個情況要立刻停止.
地步其實是位置的意思.
是一點點.
罅隙都不可以留給魔鬼.
為什麼?.
因為魔鬼的名字就是.
重傷.
魔鬼的名字就是誹謗.
魔鬼的名字就是欺控.
魔鬼的名字就是說謊.
在魔鬼裡面是沒有真理.
它的目的就是.
叫人分化.
它的目的就是叫人的關係不和.
我們積怨了這些憤怒.
就會讓我們.
分歧越來越大,甚至我們.

$^{761}$心裡常常會想怎樣報復.
有時候我都會遇到.
一些不處理自己的怒氣的人.
他們否認.
或者抑壓.
放在心裡.
結果苦毒.
老恨就一直隱藏在裡面.
慢慢.
魔鬼就會利用.
這些積壓.
令到放在我們心裡的東西.
放大了出來.
會形成一樣我覺得很可怕的事情.
就是令到自己.
成為受害者.
令到自己.
覺得自己正確.
將自己的憤怒合理化.
接著就可能.
出現了.
趁著一些機會.
積壓成憤怒.
接著就是讓罵和誹謗.
這些受害人的冤屈.
向外宣洩.
無論是背後受傷.
或者公然的對質.
或者言語上的衝突.
越吵越烈.
無論是對於自己.
或者對於對方.
都是很深的傷害.
如果在家裡.
只會令到家人之間.
的關係越來越差.
家就會散.
孩子會離異.
子女經歷這麼多的憤怒.
到他們一可以.

$^{801}$有一點點能力.
他們就遠走高飛.
如果在教會.
人與人之間.
就會分黨分派.
四分五裂.
所以弟兄姊妹.
的那麼多的罵詭.
都不可以容忍.
殺彈有機可乘.
我們知道.
當第一階段.
有憤怒的情緒.
我們去回應的時候.
我們的確有些事可以做.
有些人教我們.
要行動之前數十下.
我們經常教媽媽.
逃離現場.
不要讓更恐怖的事發生.
這些是絕對是.
我們有些方法去幫助我們.
不要讓第二階段的回應.
造成更大的傷害.
但是我想說的.
就是第三點.
我們真正的出路.
是第32節.
並要以恩慈相待.
全憐憫的心.
彼此饒恕.
正如神在基督裡.
饒恕了你們一樣.
恩慈相待.
這個字原文.
與基督這個字很相似.
有色影學者說.
當我們恩慈的時候.
我們就是仗基督.
是有一個好的心腸.

$^{841}$是一個願意關心別人的感受.
諒解他的難處.
的一個心腸.
能夠回應憤怒的行動.
其實是要由內心帶出來的.
內心有這幾個元素.
《箴言十九章》11至12節說.
人有見識.
就不輕易發怒.
寬恕人的過失.
就是自己的榮耀.
王的憤怒.
好像獅子嘲笑.
他的恩典卻如草上的甘露.
這裡說得很清楚.
其實憤怒.
最終的出口.
就是我們能夠.
蒙基督的饒恕.
以至到.
領受這個十字架的饒恕.
全因此憐憫.
和饒恕的心.
有人說.
憤怒的那一天.
應該是復和的一天.
這個世代.
告訴我們.
有仇不報.
但在《羅馬書》12章.
告訴我們.
你不要被惡所成.
要以善勝惡.
因為.
不要為自己身冤.
寧可給主的憤怒留地步.
因為經上記著說.
主啊!身冤在我.
我必報應.
你的仇敵.

$^{881}$若果我呢.
就給他吃.
若果學呢.
就給他學.
以善勝惡.
在當中必須要有體諒.
有憐憫.
也有寬恕.
其實回想.
我當初有幾年.
當家庭主婦.
我發覺脾氣暴躁.
我經常.
劈頭罵兒子.
罵到他散落.
我記得有些時候.
看到廚房亂七八糟.
心裡很不舒服.
於是就大罵一頓.
但其實.
兒子其實是嘗試.
自己做一些.
柚子蜜給我喝.
我只是弄到亂七八糟.
又或者.
我會因為兒子寫字這麼難看.
很生氣.
你都不認真.
你有沒有搞錯.
這麼不小心.
但其實.
我後來知道.
原來要寫好一手字.
是很不容易.
小肌肉的配合.
原來他寫字會累.
他已經很努力.
但他很累.
但我罵了他一頓.
其實在我們日常生活裡.

$^{921}$我們很容易生氣.
但我們有沒有想過.
其實我們真的需要.
有恩慈,有憐憫.
有寬恕.
一位德州的牧師.
去探望一位男士.
他去到他的辦公室.
看到一張照片.
他隨口說.
很漂亮的照片.
但當他說這句的時候.
那位男士的眼睛就開始紅.
然後他就看到.
有雷珠在眼眶裡.
滾動著.
然後他就問.
先生,你為什麼哭?.
那位男士說,這位是我太太.
當年我婚姻出軌.
對家人造成.
極大極大的傷害.
他決定.
我太太決定要帶著孩子走.
我看到這樣的情況.
我突然醒了.
我知道自己錯了.
我覺得很愧疚.
於是我立即和第三者斷絕來往.
然後,我知道我沒有資格.
我實在很愧疚.
但我求我的太太.
寬恕我.
結果,我太太真的寬恕了我.
這張照片.
是她寬恕了我後.
不久才拍的.
當我每一次看著這張照片的時候.
我就想起.
太太對我的寬恕.

$^{961}$她願意.
繼續站在我的旁邊.
和我走我的人生.
其實當你說這張照片很漂亮的時候.
我心裡很知道.
這張照片最漂亮的地方.
是我太太救了我的性命.
我的生命.
好像重生一樣.
可以回到一個軌跡.
另外,第二件事.
我覺得可以幫助我們.
在憤怒的過程中.
有一個出路.
就是我們要減低刺激源.
什麼叫減低刺激源呢?.
就像我剛才所說.
家裡吃飯,五碟菜有三碟魚.
我就會生氣.
這是我小時候.
但當我們成為.
基督的跟隨者.
我們的舊人會脫去.
我們的新人會穿上.
第四章21-24節.
說,如果你們聽過.
他的道,領了他的教.
學了他的真理.
就要脫去你們從前行為上的舊人.
這個舊人.
是因為私慾的迷惑漸漸變壞.
要將你們的心智改換一生.
心智改換一生.
並且穿上新人.
這個新人是照著神的形象做的.
有真理的仁義.
和聖潔.
心智轉換這個字.
其實改換這個字.
是一個不斷重複的動作.

$^{1001}$其實我們做了基督徒之後.
我們的生命.
是要不斷不斷的.
心智改換.
我們不會一信了主.
就變成了另一個人.
在我的生命裡.
我要常常去.
審查自己的內心.
究竟.
是否被上帝的真理.
和仁義在我心裡.
是嗎?.
剛才的《箴言十九章》說.
人有見識就不容易發怒.
我覺得說得很對.
如果當我們的見識.
是上帝的真理.
我們照著上帝的形象去做.
充滿了他的真理.
就不會有那麼多東西.
容易激怒我.
當我們每一次.
憤怒的時候.
我們有自己的審查.
了解我是否過著一個舊人的生活.
以致到.
我.
究竟我的人生.
我的生命.
有沒有得罪神的地方.
有沒有得罪人的地方.
其實自然地.
我的生命正在改變.
我的生命亦不需要那麼防衛.
我的生命.
在旁邊的刺激元.
亦越來越少.
令我生氣.
我看一本書叫.

$^{1041}$Emotionally Healthy Church.
中文譯名叫.
靈情步步高.
教會的不知.
我覺得很特別.
他說你回教會的時候.
你會發覺有些人.
你是特別.
怕了他.
或者在走廊見到他,你會想掉頭走.
或者是.
這個人好像不太好相處.
有些問題.
但這本書很有趣.
他說你一定覺得別人有問題.
不過他說.
你有沒有想過.
你對別人有什麼不喜歡.
我很想和大家分享.
當我預備這個講章的時候.
上帝的靈都提醒我.
我都覺得.
很想和大家分享.
做了幾個圖那麼久.
我們會否認自己的憤怒.
我們會逃避自己的憤怒.
甚至.
有些時候.
我覺得,這也是我要.
認罪的地方.
我要承認自己的愚昧.
因為當我讀聖經的真理的時候.
我用了這個真理.
去稱別人.
而不是自己.
這個人活該我生氣.
這個人.
真的不行.
真的.
我覺得很慚愧.

$^{1081}$我們的生命.
不應該只停留在.
我知道.
當一些刺激源來到的時候.
我沒有把握這個機會.
去問上帝.
我應該怎樣反應.
反而很多時候.
就成為自己的武器.
成為自己是受害者.
成為.
憤恨籠罩的基礎.
所以.
真理沒有帶來改變.
真理成為我殺人的武器.
我覺得很慚愧.
所以.
我很想提醒.
我知道上帝提醒我.
我也相信.
上帝會提醒大家.
讓.
這些怒氣.
在神的內裡好好處理.
如果我們.
常常懷著憤怒.
常常向人討罪.
心裡常常都嫌這個嫌那個.
我們真的要檢視一下.
我是靠著自己生活.
還是我真的蒙上帝的恩典.
赦免去過生活.
我是否常常要求.
主的憐憫幫助.
以致到.
我才懂得怎樣去.
處理怒氣.
與人相處.
用憐憫.
用搖書的心.

$^{1121}$與我們的弟兄姊妹相交.
與我們的家人相交.
剛才說過.
憤怒會帶來行動.
有些憤怒會叫我們退縮.
譬如.
有些情緒會叫我們退縮.
譬如羞愧.
覺得悲哀.
我們可能會收起自己.
但有些情緒會叫我們向前.
譬如喜怒.
喜和怒.
怒會叫我們自我保護.
亦會叫我們帶來行動.
當我這樣說的時候.
想起一套卡通片.
不過可能.
年輕人未必知道.
我在說什麼.
大家有沒有看過《大力水手啪牌》?.
你們額頭和領頭.
就在表現你們的年紀.
《大力水手啪牌》.
就是《The Papaya Sailor Man》.
大家看到都會心微笑.
我知道你們跟我差不多.
《大力水手啪牌》.
有一個女朋友叫奧利芙.
她有一個死對頭叫布魯托.
叫布魯托.
回望回看.
這套戲.
我覺得很悶.
因為每次都是布魯托去搶奧利芙.
然後.
大力水手很憤怒.
因為他是男生.
然後他會怎樣?.
吃大力菜.

$^{1161}$他不敢行動.
亦不知如何回應.
亦不知如何贏布魯托.
他就吃大力菜.
他一吃大力菜就會變成.
《The Papaya Sailor Man》.
然後打暈他.
搶奧利芙.
我自己覺得.
有時憤怒可以帶來正向的行動.
我們可以贏撒旦.
我們可以和這個世界不同.
我們可以勝利.
是什麼呢?.
那罐大力菜.
記住不是有仇不報.
那罐大力菜就是我剛才說的.
恩慈憐憫.
彼此饒恕.
讓我們一起祈禱.
天父我們多謝你.
我們實在獻上感恩.
我們多謝你.
因為我們已經不是舊人.
我們穿上新人.
我們能夠靠主你的恩典.
好讓我們的生命變得不一樣.
我們不會被情緒所困.
所駕馭.
我們不會被情緒牽引我們走.
反而因著這些情緒.
我們可以轉化.
我們可以帶來正向的行動.
我們可以勝過撒旦的氣洪.
勝過撒旦要叫我們.
在人與人之間的怨恨割裂.
反而是要我們.
因為愛,基督的饒恕的關係.
我們可以走在一起.
主願意我們每一個人.

$^{1201}$手裡都有大力菜.
就是憐憫,恩慈.
彼此饒恕.
讓我們能夠勝過憤怒.
讓我們能夠靠著你得著的.
與人和睦,和平之子的生命.
我這樣禱告,交託,仰望.
是奉主耶穌基督得勝的明智祈求.
阿們.
足夠法定人數.
主席.
\newpage



\section{詩篇 13:1-6}
\label{sec:9JRB4RUnJVQ}
\textbf{2021年6月6日主餐崇拜直播｜陳明泉牧師｜從聖經看喜怒哀愁 - 哀嘆中仰望｜詩篇13\_1-6}
\newline
\newline
連結: \href{https://youtube.com/watch?v=9JRB4RUnJVQ}{\texttt{https://youtube.com/watch?v=9JRB4RUnJVQ}} ~~~~ 語音日期: 2021-06-19
\newline
\newline
\hyperref[sec:Czt7dmGniVs]{\small{< < < PREV SERMON < < <}}
~
\hyperref[sec:index_chronic]{\small{[返順時目]}}
~
\hyperref[sec:bhMQ0ofdb_Q]{\small{> > > NEXT SERMON > > >}}
\newline
\newline
詩篇 13:1-6
\newline
\begin{longtable}{cl}
\hline
\hline
章節 & 經文 (和合本修訂版)\\
\hline
13:1 & \begin{tabularx}{0.7\textwidth}{X} 耶和華啊，你忘記我要到幾時呢？要到永遠嗎？你轉臉不顧我要到幾時呢？ \end{tabularx} \\ \\ \relax
13:2 & \begin{tabularx}{0.7\textwidth}{X} 我心裡籌算，終日愁苦，要到幾時呢？我的仇敵升高壓制我，要到幾時呢？ \end{tabularx} \\ \\ \relax
13:3 & \begin{tabularx}{0.7\textwidth}{X} 耶和華－我的神啊，求你看顧我，應允我！求你使我眼目明亮，免得我沉睡至死； \end{tabularx} \\ \\ \relax
13:4 & \begin{tabularx}{0.7\textwidth}{X} 免得我的仇敵說「我勝了他」；免得我的敵人在我動搖的時候喜樂。 \end{tabularx} \\ \\ \relax
13:5 & \begin{tabularx}{0.7\textwidth}{X} 但我倚靠你的慈愛，我的心因你的救恩快樂。 \end{tabularx} \\ \\ \relax
13:6 & \begin{tabularx}{0.7\textwidth}{X} 我要向耶和華歌唱，因他厚厚地恩待我。 \end{tabularx} \\ \\
[1ex]
\hline
\hline
\end{longtable}
$^{1}$願以我們以聆聽聖徒繼續我們的敬拜.
今日我們經訓有選了一段經文.
是在《舊約聖經》.
《舊約聖經》的詩篇第十三篇.
是一首由大衛作的詩.
交予當時的領長作為歌頌神.
詩篇第十三篇 請聽我讀出.
耶和華啊 你忘記我要到幾時呢.
要到永遠嗎.
你掩面不顧我 要到幾時呢.
我心裡愁酸 終日受苦 要到幾時呢.
我的仇敵 升高壓制我 要到幾時呢.
耶和華 我的神啊 求你看顧我 應運我.
使我眼目光明 免得我沉睡至死.
免得我的仇敵說 我勝了他.
免得我的敵人在我搖動的時候起落.
但我倚靠你的慈愛.
我的心因你的救恩快樂.
我要向耶和華歌唱 因祂用厚恩待我.
弟兄姊妹平安.
一齊來到敬拜上主.
剛才讀的一篇詩篇.
平時我們比較少讀的.
因為我們對於一些比較陽光.
比較亮麗的經文.
我們會容易吸引我們的注意.
如果剛才你認真去聽.
或者你去看的時候你會發覺.
在剛才的詩篇裡有一種陰沉在裡面.
雖然詩人最後都會有一種認順.
但在詩篇裡面特別呈現了.
詩人心裡面有一種哀愁.
有一種難以.
純粹用口說出來的那種沉鬱.
今天我們就用這段詩篇一齊思想.
今天這個系列裡面的第三個題目.
就是從聖經裡面看喜怒哀愁.
哀嘆中的仰望.
剛才讀的這篇詩篇.
可以在聖經歸類的文體裡.

$^{41}$我們叫做哀歌.
如果你有心.
或者仔細去看聖經.
其實在詩篇裡面.
有些學者粗略估計.
哀歌佔了詩篇裡面三分之一的篇幅.
大概在百五篇的詩篇裡面.
有五十篇.
是類似一種這樣的形態來表達的.
當然哀歌除了是在說.
詩人裡面的仇婦之外.
其實裡面也有一些出路.
人生裡總不會所有時間都是平順的.
在低潮的日子.
或者有時候在生命裡面.
有一些失落哀愁的日子裡面.
藉著這些哀歌.
可以給我們生命裡面有更加深入的體會.
當我們去明白了這個時候.
我們生命的承托能力會大很多.
希望今天我們一起看這個詩篇.
可能你今天生活很開心.
不過打到底而已.
讓我們能夠從這個詩篇裡面找到亮光.
讓我們人生裡面有方向.
有力量走在面前的路.
我們一起祈禱好嗎.
求上主幫助我們.
同心祈禱.
阿爸天父願你在我們當中裡面.
你對我們說話.
幫助我們眾弟兄姊妹.
今天我們一起來去研讀你話語的時候.
你親自對我們將你的心意展明.
幫助我們在你面前.
讓我們能夠明白你的心意.
又讓我們在我們生活當中實踐你的旨意.
恭敬將我們每一個弟兄姊妹交代給你.
幫助他們聆聽的時候聽得明白.
又讓孫講的將你話語解釋清楚.

$^{81}$帶來亮光.
以致我們生命有力量走在面前的路.
多謝主誰聽祈禱.
奉基督耶穌得勝的聖名祈求.
Amen.
在詩篇裡面.
剛才講了第十三篇.
第十三篇其實不是很複雜.
六節的經文可以分為三個段落.
其實很工整的.
第一二節是第一個段落.
三四是第二.
五六就是第三個段落.
用重複的字眼幫助我們.
為詩篇做一個簡單的分段.
在第一個段落裡面.
耶和華說你忘記我要到幾時呢?.
要到永遠嗎?.
你轉念不顧我要到幾時呢?.
我心裡愁善終日愁苦要到幾時呢?.
我的仇敵升高壓制我要到幾時呢?.
同樣出現一次又一次的重複字眼.
就是要到幾時呢?.
這是詩人在上帝面前的哀求.
講到眼前有一些.
他覺得很磨難.
很痛苦的日子.
這首詩篇是大衛所寫的.
究竟大衛在什麼狀態下去寫呢?.
其實我們在詩篇裡面看不到的.
不過我們可以想像.
一個上帝的僕人.
上帝所揀選的君王.
人生裡面都有經歷低潮.
在他低潮當中他寫下這首哀歌.
大概他面對的磨難日子.
他都好像看不到盡頭一樣.
要到幾時呢?.
也是哀歌裡面.
其實很多時候出現的這種字眼.

$^{121}$我們一看到這個字眼.
大概我們就會明白.
這是一個哀歌的文體.
也是詩人面對著一種痛苦的狀態下而寫的.
在這裡大衛或者詩人.
你會看到他們用的字眼是很直率的.
甚至你會覺得有一點離經半道.
因為我們不會這樣向上帝祈禱的.
他講到上帝要忘記他要到幾時.
是不是永遠不看他.
掩面不看他.
是不是永遠呢?.
他心裡面很煩悶.
要到幾時呢?.
仇敵壓著我.
要到幾時呢?.
這個是帶著沒什麼力量.
但心裡面有說不出的那個嘰嘰咕咕.
那種心態向上帝祈禱.
他帶著一份最大的真誠.
但是這些我們平時不會.
我們自己祈禱或者我們教人祈禱.
我們不會鼓勵頂子妹向上帝投訴的.
不過這裡大概詩人真的向上帝投訴.
他昔日經歷過很平順的日子.
經歷上帝的同在.
上帝幫助他或者戰勝仇敵.
在生命裡面有很多很彪炳的經驗.
不過去到這一刻似乎他經驗不到.
他找不到上帝.
他面對生命這些困難的時候.
他向上帝發出哀求.
發出一種投訴的聲音.
我們不用這麼快去做一個判斷.
甚至乎我們都不應該.
有時候我們很簡單用一個道德.
或者我們今天所說的價值觀念.
套進去就說.
這些就是不愛主.
愛主的人應該去到什麼時候都很緊靠主.

$^{161}$現在這些時間做投訴.
就是說他的靈性不夠好了.
我們不用這麼快去做這樣的判斷.
事實上在聖經裡面.
如果剛才你聽到的篇幅.
詩篇三分一的篇幅.
都是用一種這樣的聲音.
用一種這樣的心態.
向上帝表達.
大概在聖經裡面.
似乎都想告訴我們一個很重要的道理.
其實我們心裡面有一些.
很困難的日子裡面.
其實我們心裡面.
可能有時候有一種對上帝的投訴.
其實我們可以向上帝表明出來.
不過今天我們很容易.
會落入一種道德.
或者有一種樣板式的信仰.
樣板式的信仰就是說.
無論怎樣都要講正面的.
樂觀的積極的.
用一些標準答案.
來解釋我們整個人生.
但不過我們心裡有時候都不是這樣想的.
如果我們純粹用一種標準答案的方式.
來面對我們的主的時候.
我們的心可能想這件事.
但我們的口就說另一件事.
就是口是心非.
變成我們永遠的信仰.
都不能夠真的納入我們的根基裡面.
進入我們的心靈當中.
來將我們真實的面貌向上帝表明.
在預備這個講章的時候.
我自己都問自己.
我有沒有經歷過這樣的階段呢.
我不敢說我自己真的很困難.
我不覺得自己的生命.
我的際遇很差.

$^{201}$我不覺得.
不過我生命裡面都會有一些.
很不容易處理的一些難關.
早幾個月前家人相繼.
兩個離開.
我都很想用一個標準答案來提升自己.
覺得自己作為一個牧者.
我一定要用這個面貌來呈現上帝.
但我發覺我不行.
如果我真誠去禱告的時候.
我心裡面真的在問上帝.
為什麼我要經歷這樣的事.
為什麼我要家人接二連三這麼短時間離開.
我這麼努力去愛惜他們.
為什麼他們都不能夠長時間在我身邊.
我心裡面都是咕咕叫叫.
有很多這些.
發這些冷咋向上帝面前.
不過當我真誠地.
將我心裡面的說話向上帝講的時候.
我發現我自己的心靈裡面開始有一些地方.
有一些改變.
我心靈裡面開始覺得.
有些地方開始上帝令到我.
開始明白人生的道理.
特別是我們沒有辦法.
人生裡面總會經歷生老病死.
我們沒有辦法憑著人的智慧和能力.
可以改變到這個現實.
更重要的是.
離開了的我的家人.
其實他們每一個都已經很真心很堅心地.
在上帝的懷裡面.
他們篤信基督.
所以我們在永恆裡面.
在將來的天間我們可以相遇.
不過我今天.
如果我不將這種我的心態.
向上帝表明的時候.
我就繼續在一個樣板.

$^{241}$剛才說的那些可以是一些標準答案.
但如果你不發過冷咋去體會這些標準答案的時候.
我們是消化不了的.
再說一些入獄的.
我的同工知道的.
我的同工知道我家人這樣的時候.
在WhatsApp裡面安慰我.
我發他們.
不單發上帝冷咋.
我發他們脾氣.
我的同工可能都不覺得.
我說你不要給我一個標準答案.
我不接受.
你們不要安慰我.
在WhatsApp裡面是這樣的.
過了同工之後.
每個人都面都青了.
看到我這樣回答.
有一刻我真的變成一個.
發冷咋的人.
直到我真的在上帝面前這樣.
將我最真誠的面貌向他呈現.
我發現.
那些再不是一些標準答案.
不是一些標準答案.
那些是我可以拿實上帝的應許.
但如果我不經歷過一個情緒.
這樣真誠的流露的時候.
那些一切的應許.
都進不了我們的內心裡面.
但他私人還是帶著一份這樣的情懷.
去到上帝的庚前.
他藉著自己真誠的表達.
他不用一種樣板式的信仰.
他不再給一個標準答案.
來壓低自己的情緒.
而是用他自己真誠的面貌向上帝表明.
以至這個信仰.
真的可以有根有機在生命裡面落實.
說得盡一點.

$^{281}$其實基督耶穌在十字架上.
他都是發出哀歌的.
基督耶穌在十字架上說.
我的神我的神.
為什麼離棄我.
面對著一個重大的苦難.
他那一刻的心靈裡面的那種痛苦.
他就經歷到好像上帝離開了他一樣.
耶穌基督都將他心靈裡面最真實的話.
藉著詩篇的內容.
他再一次向神祈禱.
上帝為什麼你離開我.
私人就是這樣.
將他自己走到上帝的庚前.
除了用真誠的面貌之外.
或者不用標準答案之外.
另外還有一種心態.
有時候我們都要避免的.
就是一種叫做犬儒主義.
我不知道大家有沒有聽過犬儒主義.
犬儒就是古希臘其中一種哲學.
很多時候或者在古希臘裡面.
覺得犬儒是一種美德.
犬儒代表著什麼呢.
就是說他否定了世間裡面一切的熱情.
他覺得人生不需要放太多熱情進去.
對於那些說得很厲害的話.
那些人是不相信的.
更加不相信一些義正詞嚴的話.
不相信正義的呼喚.
不相信生命可以有改變.
他們生命用一種比較抽離冷漠的方式.
來面對著眼前這些很難的疾苦.
面對著生命裡面這些難處.
他們會將自己變成一個超理性的人.
一種冷眼旁觀的人.
對生命不再有盼望.
不再有盼望的時候就不會有失望.
犬儒主義大概是一個這樣的.
不過去到極致的犬儒主義.

$^{321}$不單止他自己獨善其身不再相信.
他也嘲笑那些人依然向上帝.
向生命裡面的主來發出哀求.
他覺得那些人是傻的.
可能我之前有講過一個例子.
我認識一個律師.
他打官司見慣了很多罪惡.
有一次聊天的時候.
這個律師就跟我說.
你們是牧者嗎?.
他用一種嗤之以鼻的態度.
他當我是朋友.
但他覺得我們基督徒是白費唇舌.
我問他為什麼要這樣想.
他說他有時候去過教會.
但他說因為他做律師.
他見過很多的案件.
他覺得殺人放火金腰帶.
修橋足露寶屍鞋.
他見過很多不公義的事.
因為那些人有錢打官司.
屈枉正直.
他見到太多荒謬的事.
所以他覺得這個信仰不真實.
他甚至對於基督徒追求真正的.
你所說的盼望.
你的真誠.
你的愛來追求你生命眼前的目標.
他覺得這些一切是荒謬.
當我一直跟他說.
我打算跟他傳福音的時候.
他用一種輕視的態度跟我說話.
到最後他說.
既然這個世界這麼荒謬.
我就要用更荒謬的方式去過這個世界.
Anders 你不要跟我傳福音.
你所說的話我不會相信.
犬儒大概就是覺得.
自己活在一個荒謬的世界裡面.
他不再覺得生命裡面有一個主宰可以帶領他.

$^{361}$他否定神.
否定生命裡面一切美善的價值.
他用更荒謬的方式來過他的生活.
不過哀歌的作者.
或者我們作為一班的基督徒.
雖然我見到這個世界這麼荒謬.
見到有些事是很不理想.
或者我們生命裡面有些難處是很難受的.
無論有多少困難.
但我們認定一件事.
我們生命裡面還有一位主.
我們不會因為眼前的困難和困境.
而否定我們生命裡面還有一個比我們辛酸的對象.
哀歌為何成為一個重要的問題.
因為人經歷一些很難的處境當中.
我們生命裡面還有一個神可以讓我們化下難渣.
他會知道我們的苦情.
他會在我們最困苦的時間裡面.
當我們祈求的時候.
他可能會為我們出手.
扭轉我們眼前這些苦況.
我們生命不會被所有這些困難排山倒海的東西.
擠壓著我們以至否定我們的信仰.
犬儒主義表面上好像活在這個世界.
但犬儒主義者某程度上是行屍走肉.
因為生命裡面他不再有盼望.
他過比這個世界更加荒謬的生活.
哀歌的詩人卻是在荒謬的世界裡面.
選擇向這個掌管天地萬物的主傾心吐怡.
將他心靈裡面的痛苦向這位神表達.
所以,弟兄姊妹,當我們看哀歌的時候.
我們第一件事是要認定的.
我們生命裡面的難處.
有一個主會聆聽我們.
甚至有些說話離經半道.
不是很教導式.
但在你的思圖當中,在你的個人祈禱裡面.
這個真實的面貌是要呈現出來.
不要太快被一些model answer.
即是成為你的內容.

$^{401}$反而你向神真心吐怡.
這些就不再成為model answer.
而是成為上帝對你的應許.
以至到你有一些情感擋住.
你沒有辦法接受這些說法.
但當你的情感表達了的時候.
你就會慢慢答落.
你就會知道這個應許實在太真實,太寶貴了.
這是第一個段落.
給我們一些很重要的提醒.
我們再看第二個段落.
十三章第三至第四節.
耶和華我的神,求你看顧我,應允我.
求你使我眼目明亮.
免得我沉睡至死.
免得我的仇敵說我姓了他.
免得我的敵人在我動搖的時候起落.
在十三章第三至第四節裡面.
語氣上改變了.
還有重複的字眼.
祈求,求你看顧我,應允我.
求你使我眼目明亮.
這個祈求是私人改變了他的語氣.
另外一個字眼就是免得我沉睡至死.
免得我的仇敵說我姓了他.
免得我的敵人在我搖動的時候起落.
你看到嗎?.
祈求和免得是第三至第四節裡面.
私人轉念的禱告句式.
你知道這是另外一個段落.
在這裡面,私人剛才由一種哀求,哀嘆.
轉化成為他向上帝祈求.
在人無力改變的情況下.
或者在一些很痛苦的經歷當中.
這個私人又再向上帝發難.
轉化成為祈求的語氣.
看顧我,應允我,就對應上帝掩面不看.
忘記他開頭詩篇起首的兩句說話.
祈求上帝在他痛苦當中.
聽他祈禱,應承他.

$^{441}$看著我,不要忘記我.
他心靈裡知道上帝繼續與他同在.
他知道上帝如果要出手的時候.
上帝一定可以幫助他.
值得注意的是,在這個祈禱裡面.
還有兩個很重要的兩句說話.
求你使我眼目明亮,免得我沉睡至死.
在這裡面特別提到眼目明亮這個字.
哀痛的人為何要尋求上帝.
要祈求主令他眼目明亮呢?.
眼目明亮是一個象徵的手法.
不是眼睛老花馬上消失.
不是這個意思,大家明白的.
如果你深入去看的時候.
眼目明亮是指他有聰明智慧.
他想得通,他的思想想得透.
這個思想想得透帶來一些重要的結果.
就是免得自己沉睡至死.
沉睡就是指與明亮是相對的字眼.
沉睡就是指自己失去了盼望.
失去了生命裡的堅持.
一個活著的人好像睡了一樣.
這個就是沉睡,完全不開竅.
完全好像自己不知道自己在做甚麼一樣.
帶著一種絕望的態度過自己的生活.
帶著絕望的態度直至死亡.
在私人的祈禱裡面.
他祈求上帝除了他的際遇改變之外.
他更加祈求上帝幫助他自己.
他有一些的聰明智慧.
令自己不要落入絕望至死的狀態.
這種聰明是一種怎樣的聰明呢?.
私人沒有直接講的.
不過我相信這種聰明.
就是對自己內在的心有更加多的體會.
明白人生,明白人生的痛苦.
也明白自己的哀傷去到甚麼位置.
今天容許我直接講.
其實我們有時在網上看.
我們其實滿佈這個社會裡面.

$^{481}$有時有些資訊是塑造我們的想法方法.
當我們在看的時候.
其實悲情絕望的說話是滿佈在你的facebook裡面.
在你的生活當中.
有時我們未必自己這麼負面這麼悲觀.
不過當你自己覺得有少少情緒的時候.
再看這些東西.
再看這些內容的時候.
你的情感就像搭升降機一樣會提升.
會一直提升你,一直煲你.
令你覺得這個世界是絕望的.
甚至今天的訊息.
Google會幫你.
令你變成這樣的情緒.
當你不開心的時候.
你打如何令自己開心一點.
其實它會繼續打負面的資訊在你的Google裡面.
我們被很多的言論.
很多的說話來擠壓我們.
我們自己的情緒被身邊很多的話語.
或者被很多渲染出來的一些矯情的內容.
令我們看這個世界變得灰暗.
或者加劇了我們自己那種悲哀的情緒.
又或者有時我們面對自己.
都可能有一些低潮.
但是我們心裡沒有一種聰明智慧.
來為著這種負面的想法去勒一勒或是劃一條界線.
我們久而久之.
好像成為一個情緒失控的人.
一有負面情緒就好像開快車一樣.
一踩油門就去.
令自己一夜就去到極致.
有一個這樣的例子.
十多年前我外公去世了.
我外公去世做安息禮拜.
做安息禮拜的時候.
有一些親友.
因為老人家.
認識了幾十年的老街坊.
都來殯儀館表達一點心意.

$^{521}$我見到一個情景很奇怪.
我認識一個老街坊.
我也認識的.
就是外公那一代的那些.
平輩,街坊.
三十多年的街坊.
但不等於三十多年的交情.
大家明白這句說話嗎?.
認識了三十多年.
但不等於很緊密.
只認識一個名字.
或是有一些接觸.
不過是一個女街坊.
我外公去世了.
是一個女街坊.
到了喪禮的占用遺容儀式裡.
這個女街坊占用了我外公的遺容.
在過程中哭得呼天搶地.
抱著棺木來哭.
我和我家人臉色一沉.
為何會這樣呢?.
我外婆也沒有到達.
其實我外婆也沒有到達.
哭得呼天搶地.
抱著棺木哭得很厲害.
我們也不知道發生甚麼事.
哭了.
結束後問她.
我們是家人.
還要安慰她.
節哀順變.
整件事很奇怪.
我們問她為何會有這麼大反應.
她說不開心時當然會踩紅褲子.
所以她看到這樣.
其實我交情未到此地.
但她覺得不開心時就踩紅褲子.
於是在那一刻有極致的表現.
我外婆當然是咀咕咕.
她說我還未到此地.

$^{561}$如果別人不知道.
以為我外公和她有甚麼關係.
其實是沒有關係.
不過只是將自己的情感推到極致.
詩人在此說的這句話.
求你使我眼目明亮.
免得我沉睡至死.
如果你用這個角度去想.
每個人都有經歷.
我都可以真誠地去表達.
但去到某些位置.
我們要求上帝幫助我們.
看我們的心.
我們的悲哀或苦情.
要懂得分辨去到甚麼位置.
不要每一天都踩紅褲子.
將你心裡的一切都說出來.
如果我們這樣.
我們不會有自我察覺和覺醒.
可能我們就永遠終日.
活在一種沉溺在自己的哀傷當中.
可能我們整生人.
繼續在這個哀傷裡.
帶著絕望直至死亡.
詩人這句說話.
如果我們用今天的現況.
是很值得我們繼續深思.
另外一個字眼.
在這裡特別提到仇敵.
免得我的仇敵說我姓了他.
免得我的敵人在我動搖的時候起落.
敵人是誰呢?.
如果你說大衛的詩篇.
可以將大衛曾經的敵人數點出來.
可能是巨人哥利亞.
又或者是蘇羅等等.
或者是他自己的兒子亞沙隆.
這些可以將他視為大衛的仇敵.
不過如果我們只是將仇敵或敵人.
純粹套入某個人身上.

$^{601}$我們未必能夠用到這個詩篇.
事實上我們生命裡未必有這麼多的敵人.
在我們的日常生活裡.
其實不是四處為敵來生活.
我們未必活在一種戰爭的狀態裡.
不過敵人這個字.
其實可以形象化.
將你生命裡面的一些難關.
將它成為敵人的面貌.
呈現出來.
以至你不是保羅所說的.
九章二十六節哥林多前書說.
我奔跑不是無定向.
我的鬥權不是打空氣.
我將我生命裡面的一些困難.
我將它成為具體的形象.
我針對它來處理.
而不是打空氣.
我這樣說很複雜.
我講個例子你就明白.
這是近代裡面敘事治療.
其中一個很重要用的方法.
如果有一個抑鬱症的姐妹.
她走來問我們求助.
特別找我來求助的時候.
除了我說要處理眼前身體等等之外.
我就會告訴她.
抑鬱症如果給你比喻.
是一個怎樣的敵人.
是一個怎樣的形態.
你可不可以描述給我聽.
我講個例子.
姐妹說抑鬱症就像一隻獅子一樣.
一隻強悍的獅子.
這隻獅子經常攻擊她.
令她有種很大的痛苦.
我就會問這隻獅子何時沒那麼兇.
這隻獅子有沒有一段時間沉睡了.
你會好一點.
她說可能她夠睡的時間.

$^{641}$或者她和朋友在一起的時候.
這隻獅子沒那麼兇.
然後我們就想辦法對付這隻獅子.
我用的這個方法.
或者我們所想的敵人.
其實就是將我們生命裡面的一些難關.
我們將它具體化.
你不要套用在一個人身上.
你不是以人為敵.
我的目的不是叫你這樣.
只不過是將一些困難.
一些處境.
打倒你的東西.
你將它形象化.
你不是像一天.
做一拳又一拳.
你都不知道在哪裡出拳.
你自己覺得很被動.
很辛苦.
但當你總結了某一些.
現況有些困難的時候.
你將它.
看到一個比較清晰的敵人的輪廓.
你會幫到自己.
好像有方法去針對這些問題.
問題不會打垮了你.
而是你可以和這個問題.
可以有策略去面對它.
在詩人裡面說.
不想他的敵人要打贏他.
不要在他自己站立不住.
動搖的時間來喜樂.
在詩人心目中.
他很清晰有一個敵人的形象.
是一個怎樣的敵人的形象.
在他憂傷.
在他低落的時候.
他很清楚他要對付的是一個怎樣的對象.
不要讓這個仇敵在生命裡面跨性.
同樣我們今天.

$^{681}$在哀傷的日子.
在沉淪.
或者覺得痛苦.
不知道怎樣走的日子裡面.
將你的問題.
將它形象化了.
你可以比喻為一種動物.
比喻為一種你覺得很討厭的.
一種電影形象等等.
具體地將這個問題現形了出來.
你就有力量去對付它.
你就不會處於一種捱打的狀態.
我們再看第三個段落.
詩篇十三章第五至六節.
但我倚靠你的慈愛.
我的心因你的救恩快樂.
我要向耶和華歌唱.
因他厚厚地恩待我.
詩人他在第三個段落裡面.
當你見到一些字.
就會轉折的.
就是但.
我要怎樣.
這個也是語氣上的改變.
開頭是祈求的.
最後一個段落是對神的一種倚靠.
他拿著上帝.
圍著上帝的應許.
上帝給他的恩典.
他站在這些恩典當中.
然後給他有力量去面對眼前的挑戰.
詩人他特別講到.
他倚靠上帝的慈愛.
這個慈愛絕對不是一個.
純粹口說出來的.
一個堆砌出來的字眼.
詩人真實在他生命當中.
經歷過上帝的慈愛.
經歷過上帝怎樣愛他.
幫助他.

$^{721}$讓他生命裡面經歷到上帝的恩典.
新約所講我們嘗過天恩的滋味.
就是經歷過信仰.
經歷過上帝對我的那種口袋.
詩人站在一個感恩的位置裡面.
有時候我們.
很容易會落入一種很奇怪的謊言當中.
當我們在痛苦裡面.
我們很多時候會有一種體會.
就是將一些美好過去的回憶.
否定了它.
又或者講那些舊的東西.
有什麼用呢?.
現在很慘.
將舊的東西.
成為今天的對照.
令到今天覺得自己更加痛苦.
詩人的思路是相反的.
今日既然這麼痛苦.
你就想起你過去的.
上帝在你生命裡面的慈愛.
救恩.
以及怎樣去恩待自己.
用上帝這個主.
用上帝舊時儲蓄給我們生命裡面的美好.
這些恩典我們儲起了.
這些恩典儲起的時候.
到我們真的面對艱難的時候.
我們可以慢慢撕掉.
我們有本錢.
來跟眼前.
令到打垮我們.
或者令到我們很困擾的處境.
來抗衡.
跟這些處境我們有力量去面對.
如果連這些儲蓄都沒有.
大概我們虛怯到一個地步.
我們什麼都沒有本錢去面對.
不過詩人告訴我們.
其實在她生命當中.

$^{761}$她大大經歷上帝的恩典.
弟兄姊妹.
甚至我們每一個.
生命裡面.
再難走也好.
都會有足夠的恩典在我們生命當中.
這些一切的恩典.
就好像上帝為我們生命儲蓄一樣.
儲夠蓄的時候.
我們就有力量去打這些眼前.
這些很多不容易的處境.
我昨天有機會看回我的舊同事.
他可能也落入了一種困難當中.
他在Facebook上寫了一句很有意思的話.
他說思念可能很多人會覺得遺憾.
但思念也可能是力量的傳源.
思念過去美好的日子.
是會給我們有力量.
讓我們知道.
我們人生不是上帝所撇棄.
我們帶著一些思念.
可能是家裡的人.
有些生命裡面上帝相遇的片段.
這些一切都會幫助我們.
走眼前不容易走的路.
更加重要的是.
在生命當中.
這種思念.
在這些困難裡面.
重新來思考.
其實是幫助我們的生命.
我們更加有承載的能力.
如果我們一路快樂的日子.
人生裡面很平順.
所有事都很快樂.
我們很羨慕.
但將來天國裡面.
大家講自己的見證的時候.
一生人都很平順.
沒有甚麼故事可以講.

$^{801}$不過.
我們生命裡面有些眼淚.
有些不容易的處境.
但又有恩典在我們生命當中交織而來.
我們生命就精彩很多.
甚至將自己信仰的經歷.
告訴別人的時候.
你告訴別人的不是單單是.
一個很順頭順路.
上帝幫助你.
處理好所有事.
幫你難關.
處理好.
那種信仰觀.
這都是好的.
不過更加重要的是.
我們所信的基督耶穌.
這個信仰如何承載我們生命的苦難.
既然這位主都是生命的苦難.
在哀傷裡面.
都是交織著祂的即時的時候.
我們生命都是跟隨耶穌基督的足跡.
一路走面前的路.
我們的生命故事.
和我們所信的主.
我們大概都是一致的.
我們和坊間的人不同.
我們不是純粹追求追吉避凶.
而是我們追求一個.
可以承載我們生命.
不同的喜怒哀樂的.
一個很重要的信仰.
如果生命是一個良器.
如果我們只是單一盛葬快樂.
我們這個容器實在太淺.
承載能力實在太低.
但我們生命裡面.
如果這個容器裡面.
有一些傷痛.
有一些眼淚.

$^{841}$有一些苦澀.
可能在過程裡面被作.
或者勉強去撐開它.
但當我們一路再回首一看的時候.
哇!這個生命的容器.
確實可以成為很多人的祝福.
這個容器有深度地.
來活出一個基督教的信仰.
丁姐妹.
今天不知道你落在一個什麼處境當中.
可能你很平順.
我恭喜你.
今天這篇詩篇也適合你聽.
因為我們沒辦法.
既定地想著我們人生的快樂.
可以直到前面.
今天這篇詩篇提醒我們.
我們生命要有承載的能力.
同樣你可能正在經歷一些哀傷.
正在經歷一些不容易的難關.
我深信這篇詩篇同樣幫助我們.
這篇詩篇提醒我們.
我們在上帝面前表達我們的愛情.
不是不屬靈的.
不是靈性不好.
反而是當我們真誠向上帝的時候.
上帝才在我們心靈裡面來對我們工作.
這篇詩篇更加提醒我們.
私人祈求祂眼睛明亮.
我們時刻要察驗我們的心靈.
不要讓我們的心絕望至死.
這篇詩篇更加提醒我們.
生命裡面我們不是處於捱打狀態.
當我們覺得很痛苦的時候.
嘗試將你令你痛苦的敵人現形出來.
你知道他現形的時候.
你會懂得怎樣去對付他,面對他.
甚至你會懂得在不同的時間裡.
透過上帝的恩典陪你走過眼前的路.
不要忘記上帝一直沒有離開我們.

$^{881}$上帝的救恩.
上帝昔日在我們生命裡面的過去.
這些都是一些很重要的資產和寶藏.
我們可以拿著這些恩典走面前的路.
更加重要的是.
今天你所流過的眼淚.
他日成為勝在別人眼淚的皮袋.
你經歷過的痛苦.
你會成為別人的祝福.
願主繼續幫助我們.
讓詩篇第十三篇.
繼續對我們每一個人說話.
我們同心祈禱.
上主願你幫助我們.
我們恭敬在你面前.
領受你懷孕.
改變我們憂傷的心.
讓我們在覺得痛苦裡面.
我們仰望你.
以致我們不會失腳.
幫助眾弟兄姊妹.
願你的話語成為每一個的幫助.
以致我們生命裡面.
有人力,有恩典.
有那種從你而來的聖靈提醒的能力.
以致我們人生不會活到失去盼望.
反而無論怎樣都好.
我們緊緊依靠你.
祝福我們眾人.
讓我們今天繼續實踐你給我們的教導.
帶領我們走面前的路.
多謝主上聖聽祈禱.
奉基督耶穌得勝名祈求.
Amen.
\newpage



\section{哥林多後書 10:12-18}
\label{sec:bhMQ0ofdb_Q}
\textbf{2021年7月4日崇拜直播｜梁家麟牧師｜越界問題｜哥林多後書10\_12-18}
\newline
\newline
連結: \href{https://youtube.com/watch?v=bhMQ0ofdb-Q}{\texttt{https://youtube.com/watch?v=bhMQ0ofdb-Q}} ~~~~ 語音日期: 2021-07-10
\newline
\newline
\hyperref[sec:9JRB4RUnJVQ]{\small{< < < PREV SERMON < < <}}
~
\hyperref[sec:index_chronic]{\small{[返順時目]}}
~
\hyperref[sec:VDCXuQp5CNk]{\small{> > > NEXT SERMON > > >}}
\newline
\newline
哥林多後書 10:12-18
\newline
\begin{longtable}{cl}
\hline
\hline
章節 & 經文 (和合本修訂版)\\
\hline
10:12 & \begin{tabularx}{0.7\textwidth}{X} 因為我們不敢將自己和某些自我推薦的人並列相比；他們用自己度量自己，用自己比較自己，是不明智的。 \end{tabularx} \\ \\ \relax
10:13 & \begin{tabularx}{0.7\textwidth}{X} 我們不願意過分誇口，但是我們只在神劃定的界限內誇口。這界限甚至擴展到你們那裡。 \end{tabularx} \\ \\ \relax
10:14 & \begin{tabularx}{0.7\textwidth}{X} 我們擴展到你們那裡時並沒有越過了自己的界限，其實我們是首先到你們那裡傳基督福音的。 \end{tabularx} \\ \\ \relax
10:15 & \begin{tabularx}{0.7\textwidth}{X} 我們不靠別人所勞碌的過分誇口；我們只希望你們信心增長的時候，所劃定給我們的範圍也能夠因著你們更加擴展， \end{tabularx} \\ \\ \relax
10:16 & \begin{tabularx}{0.7\textwidth}{X} 使福音得以傳到你們以外的地方，而不在別人的範圍之內，以別人所成就的事誇口。 \end{tabularx} \\ \\ \relax
10:17 & \begin{tabularx}{0.7\textwidth}{X} 但「要誇耀的，該誇耀主」。 \end{tabularx} \\ \\ \relax
10:18 & \begin{tabularx}{0.7\textwidth}{X} 因為蒙悅納的，不是自我稱許的，而是主所稱許的。 \end{tabularx} \\ \\
[1ex]
\hline
\hline
\end{longtable}
$^{1}$我們今天要聽神的說話.
是記載在哥林多後書第十章.
十二至十八節.
哥林多後書十章十二至十八節.
新約聖經的第二百五十三頁.
找到的請聽弟兄讀出.
哥林多後書第十章十二節.
因為我們不敢將自己和那自薦的.
人同列相比.
他們用自己度量自己.
用自己比較自己.
乃是不通達的.
我們不願意分外誇厚.
只要照神所量給我們的界限.
救到你們那裡.
我們並非過了自己的界限.
好像救不到你的那裡.
因為我們早到你們那裡.
傳了基督的福音.
我們不將著別人所勞碌的分外誇厚.
但指望你們信心增長的時候.
所量給我們的界限.
就可以因著你們更加開展.
得以將福音傳到你們以外的地方.
並不是在別人界限之內.
藉著他現成的是誇厚.
但誇厚的.
當自著主誇厚.
因為蒙越立的不是自己清去的.
乃是主所清去的.
聖經獨筆.
各位親愛的兄弟姐妹.
我們來到第十六次講哥林多後書.
越界的問題.
保羅在這裡直接回應.
那些將他和那些超級仕途做比較.
並且說他不如那些超級仕途的評論.
他首先回應的對象.
當然是那些提出這樣的說法的弟兄姊妹.
哥林多的弟兄姊妹.

$^{41}$不過不可避免的.
因為他要答辯.
他不是不如人.
所以很自然的.
他也會提出一些對超級仕途的負面評價.
保羅首先用一個諷刺性的口吻說.
十二節.
因為我們不敢將自己和那自薦的人同列相比.
他們用自己度量自己.
用自己比較自己乃是不通達的.
不敢這裡指的不是沒有膽量.
而是我不會狼到放肆到.
口硬無恥到這樣的程度的意思.
就是說我做不到這樣的地步.
英文是fondary(誇張).
或者是audacity(膽怯).
總之是我受不了.
那是過了我的界線.
保羅敢去哥林多懲處那些不服從的信徒.
但他不敢和那些哥林多的信徒.
所推崇的超級仕途去做比較.
為什麼呢?.
因為這個超出了他作為基督仕途的道德底線.
和做人的原則.
那些人的所作所為.
是作為基督仕途的他不能夠做的.
他們一個很重要的特徵是什麼呢?.
是自我推薦.
self-commendation(自我推薦).
自我宣傳.
自我誇耀.
這種自己訂定一個標準去量度自己.
自己為自己打分.
自己宣告自己滿分.
這樣的做法.
保羅說我怎樣都做不出來.
正如他在前面說.
凡是誇口的只能夠指著主誇口.
是不是?.
他不可以用任何形式去舉薦他自己.

$^{81}$三章一節.
他也不可以高舉人的智慧和能力.
要誇的只能夠誇自己的軟弱.
十一章一節到十二章十三節.
這個是作為基督所猜派的傳道者.
就是他的身份.
以及他要傳講的道.
就是信息所限定了的.
他的身份和信息.
都限定了他不可以做這些事情.
他們玩的遊戲他玩不到.
他們高舉的做法.
他自己認為是說不通的.
這裡說的是不通達的意思.
大家都在不同的平台上.
怎麼可以互相比較呢?.
各位姐妹.
我們知道這個真的是一個推銷和自我推銷.
self-advertise(自我宣傳)的年代.
每個人都是推銷員.
不是推銷你賣的產品.
就是推銷你自己.
其實推銷自己比推銷你的產品更重要.
因為你自己就是你所賣的產品的最重要的品牌.
是不是?.
很多店的老闆和生產者.
如今都不喜歡退居幕後.
而跑到鏡頭面前.
親自介紹自己所做.
自己所出產的產品.
是不是?.
所以呢.
特別是現在你經常看YouTube.
你真的認識了很多咖啡師.
製餅師.
廚師.
設計師.
是不是?.
從前.
我只知道街口那一家餅店的蛋撻好吃.

$^{121}$我不知道那個製餅師傅是誰.
今天呢.
那個製餅師傅就會自己粉墨登場.
穿著全新的白色工作服.
然後解說他一絲不苟地用十八個工序.
你去精心烘焙出一個蛋撻來.
是不是?.
一個涼茶店的老闆.
好熱衷.
擔任他店的廣告片的主角.
至於那些在巴士車身擺成特務尖士邦.
這樣的姿勢的那些補習社的名師的廣告.
我猜我們都見怪不怪了.
我要說.
這樣子自我宣傳的風氣.
也會影響教會的.
我答應.
某個地區的教會主令陪寧聚會.
他首先就要求.
說完.
才寫一段自我介紹的文字.
不是說你的個人簡歷.
而是要你講一下你自己.
有些什麼要讓人必聽的特點.
這個真的很受不了.
很奇怪的事情.
明明是你們請我.
現在接著要說.
我有些什麼值得你請的理由.
怎麼說呢?.
例如我高大英俊.
我的聲音洪亮.
我的說話幽默.
我的思想獨到.
我的信息接地氣.
對不起.
剛才說的話我通通都沒有.
接著在聚會前不久.
大會又要求.
說完自拍一條短片.

$^{161}$為陪寧會做宣傳.
自己宣傳自己.
然後將那個宣傳片放在不同的堂會裡面播放.
我是一個過時落伍的人.
我很討厭.
我簡直要用討厭這個字.
就是自我宣傳.
多次拒絕.
配合他們做.
於是被人罵得飛起來.
如今學乖了.
先小人後君子.
在答應邀約之前.
我首先講明.
我只是負責預備講章.
我不做宣傳的.
否則我就會臨時退縮不說的了.
這樣做會減少一些直接衝突的.
不過在背後被人指點點.
說我不合時宜.
這個就不用說了.
大家都知道的了.
是不是?.
今天很多教會.
很多機構.
當要作聲一些主管級的同工.
也會沿用外面商業世界的一貫做法.
要求同工自己申請.
並且要跟別人競爭的.
要積極爭取上位.
我真的要很慶幸.
我是屬於上一輩的人.
我工作的時候.
這樣的風氣還沒有流行的.
我的時辰生涯.
所有任命或者升遷.
都是被人冤枉出來的.
自己從來沒有爭取過.
如果今天.
在任命我做神樂院的院長之前.

$^{201}$我要在院董會裡面.
回答院董的提問.
你覺得你有什麼優點.
值得我們委任你做院長呢?.
我會立刻說.
沒有.
什麼都沒有.
不用考慮我.
我說真的.
如果幾年前.
羅志強執事是這樣問我.
梁牧師.
請你說明一下.
你有什麼資歷.
為什麼值得王震找你做顧問牧師呢?.
我保證今天我不會在大家中間.
因為我一個理由都找不到.
保祿很厭惡.
這些自我宣傳的做法.
他拒絕自己.
舉薦自己.
並且.
在其他人.
要他跟那些超級使徒做比較的時候.
他首先就棄權.
我認輸.
我承認我不及他們.
因為我們不敢將自己和那自薦的人.
同列相比.
他不做這些比較.
不過.
保祿不是純粹因為謙虛.
所以拒絕自薦.
他拒絕自薦是有兩個理由的.
一方面.
好像前面所說的.
這是他的傳道身份和他所傳的道.
所加給他的約束.
他沒有辦法自我誇耀.
不可以吹噓他自己的能.

$^{241}$這是第一個原因.
但還有第二個原因.
第二個原因.
保祿說.
他在光陰多教會裡面.
他擁有一個特別的條件.
或者叫優勢.
是包括那些超級使徒在內的所有人.
都比不上他的.
光陰多教會是他一手傳福音.
栽培建立起來的.
信徒就是他的戰書.
所以保祿說.
我根本不需要另寫戰書.
第一個原因是他不能夠.
第二個原因是他不需要.
當保祿說到他和光陰多信徒的特殊關係.
他在光陰多教會的不可替代的角色的時候.
保祿就提出了一個界限的觀念.
我們不願意憤害誇口.
只要照上帝所量級我們的界限.
救到你們那裡.
我們並非過了自己的界限.
好想救不到你們那裡.
因為我們早到你們那裡.
傳了基督的福音.
我們不將著別人所勞碌的憤外誇口.
但指望你們信心增長的時候.
所量級我們的界限.
就可以因著你們更加開展.
得以將福音傳到你們以外的地方.
並不是在別人的界限之內.
藉著他現成的事誇口.
13到16節.
保祿說.
在上帝畫給他的地界那裡.
光陰多教會是屬於其中.
這裡和學本用了一個很罕見的舊字.
即是接觸或拉上關係的意思.
然後保祿說.

$^{281}$我們早到你們那裡.
傳了基督的福音.
指著保祿是第一個去光陰多傳福音的人.
在那裡開拓建立教會.
保祿說.
他不是一廂情願.
甚或是自作多情的.
和光陰多信徒攀附關係.
他和他們當然有關係.
或者更加準確的說.
是他們和保祿有割不斷的關係.
沒有保祿的傳福音.
保祿的奴婦.
就沒有第一批信徒信主.
沒有教會的建立.
當然我們說.
沒有保祿.
上帝也可以興起另一個人.
是不是?.
不是時必要保祿的.
好.
你說上帝這個元素.
那我就要說.
但是現在事實上.
是保祿.
是上帝派遣了保祿在那裡建立教會.
所以保祿在那裡建立教會.
保祿是這教會的屬靈生父.
這個是上帝的旨意.
你要說上帝.
你要說事實上.
這就是上帝的旨意.
對不對?.
所以保祿說.
那麼多教會是上帝量給他們的界限.
就好像說.
上帝畫給他的地盤.
就像昔日.
以色列人進入迦南之後.
各個支派分一個地.

$^{321}$保祿獲賓派到哥林多建立教會.
由於哥林多是保祿的界限.
所以當保祿介入處理哥林多教會出現的人事問題.
信仰問題的時候.
當保祿將哥林多信徒的信仰.
表現成為自己的成績單的時候.
保祿是沒有過界的.
沒有愚忿的.
沒有攀雞思神.
沒有著囚溝佔.
沒有褻奪人的位置.
沒有侵佔人的權益.
保祿說我們不是過了自己的界限.
好像救不到你們那裡.
因為我們祖導你們那裡傳來基督的福音.
保祿說的話全部都是恰如其分的.
自己種的樹.
收你枝葉.
吃樹上的果子.
是不是理直氣壯?.
這是舊約的說法.
看守無花果樹的.
必吃樹上的果子.
心中的真言二十七章十八節.
又說自己.
因為我想不到更好的例子.
我看到神學院服侍了三十七年.
我事奉生危.
總之我出道.
有95\%都是在神學院裡面.
當然我不是說神學院是我一手建立的.
是我建立的.
我們有很多同工很多義工.
他們為神學院盡心擺上.
他們做的比我更多.
不過.
如果這三十七年裡面.
多數的事工.
我說和我都有關係.
我又推卸不了的責任.

$^{361}$我必須要分擔.
或者分享神學院的所有榮同慾.
我就會當仁不讓.
特別是做院長那十三年.
不可能有什麼事工和我完全無關.
我可以洗手撇清關係.
是不是?.
我知道將來.
在審判台前和耶穌基督交仗的時候.
見到一定是佔我很大很大的比例.
不過.
我在旺鎮只是幾年的時間.
並且參與的程度非常低.
多數只是維持一個客卿的身份.
所以.
我只是見證.
各位牧者各位執事在這裡的打拼.
見證上帝怎麼使用你們.
我不會視這裡是我.
呈交給耶穌基督的成績單.
是不是?.
所以.
用保羅的說法.
見到是我的界限.
旺鎮不算.
為什麼保羅要特別強調.
哥林多是上帝畫給他的地界呢?.
非常可能的原因是.
當保羅透過書信.
處理哥林多教會裡面的一些錯謬的事件.
一些紛爭的時候.
曾經有人就指出.
保羅你既然已經不在我們中間服侍.
你就不應該再插手過問我們的事情.
保羅你在哥林多教會再沒有管治權.
沒有 jurisdiction(管治權).
保羅對此作出清楚的回應.
他有權過問哥林多教會的事務.
他不僅僅是在他們中間.
服侍過一段時間的一個過客.

$^{401}$他有兩個特殊的身份.
第一.
他是他們的屬靈父親.
這間教會是他一手建立的.
就好像我媽媽那樣.
她經常都這樣責罵我.
你再厲害.
我都是你媽媽.
我不責罵你.
我不罵你.
誰罵你呢?.
你聽到你媽媽說這些話.
你怎麼做呢?.
當然是收口了.
鞠躬.
媽媽會有界限的嗎?.
會做好工作.
沒有責任嗎?.
沒有的.
一天是媽媽.
她一天還沒死.
她一天都可以餓你.
就是這樣.
保羅是哥林多教會的屬靈生父.
這是一個無論如何都擺脫不了的責任關係.
第二.
保羅是仕途.
不是一般的傳道者.
這一點.
後面保羅在哥林多學會再說的.
他是仕途的身份.
接著保羅再說了一句.
我們不仗著別人所勞碌的憤外誇口.
這句話可能是映射暗示其他人.
特別是被吹捧的超級仕途.
有搶佔別人工作成果的嫌疑.
雖然他們曾經在哥林多教會服侍.
並且他們對教會其實都有相當的貢獻.
至少在哥林多前.
舒保羅都承認.

$^{441}$我栽中了.
阿波羅囂貫了.
他沒有否認阿波羅有貢獻.
是不是?.
阿波羅是有貢獻的.
不過.
他們所佔的份額一定不及保羅.
所以不應該將教會歸到他們自己名下.
變成他們的成績單.
否則他們就是仗著別人所勞碌的憤外誇口.
我們知道.
保羅其實是一個頗另類的傳道人.
他以一個很另類的情況成為耶穌基督的仕途.
耶路撒冷的教會群體總是有些格格不入的.
保羅和那些耶路撒冷的嫡系門徒.
其實不是很能混在一起的.
他獨來特往.
他沒有受耶路撒冷教會指派.
也沒有得到他們怎樣的資助或者支援.
他領受的是向外邦人傳福音的意象.
致力做開方的工作.
向美德之地挺進.
他曾經做過兩個宣告.
在《加拉提書》一章十一到十二節他講.
弟兄們我告訴你們.
我數來所傳的福音不是出於人的意思.
因為我不是從人領受的.
也不是人教導我的.
而是從耶穌基督啟示你的.
這是他所講的福音.
不是受教於耶路撒冷的仕途.
然後羅馬書十五章二十節他也講.
我立了志向.
不在基督的名被稱過的地方傳福音.
免得建造在別人的根基之上.
這裡是講他的事奉使命.
所以你會看到的是.
保羅指有人說他在別人的根基之上建造.
仗著別人所勞碌的誇口.
這些批評保羅是特別皮膚過敏的.

$^{481}$總之是非常敏感.
最討厭聽到這些.
我沒有佔人便宜.
我沒有將別人的功勞納在我自己的帳下.
不過保羅也不是對.
仗著別人所勞碌的誇口.
這句話徹底拒絕.
什麼都要自己做自食其力.
他說有一種情況是他可以接受的.
就是哥林多信徒熱心傳福音.
將教會一直擴展開去.
所以他們結束的果子.
就可以視為是保羅的間接成就.
保羅的界限也因此往外擴充.
15到16節上半節.
我們不仗著別人所勞碌的憤外誇口.
但指望你們信心增長的時候.
所以量級我們的界限.
就可以因著你們更加開展.
得以將福音傳到你們以外的地方.
這樣的說法我猜我們也不會很陌生.
做保險或者營銷的弟兄姊妹.
你們也知道這個故事.
那些高階主管或做師傅的.
可以將他屬下的僱員或徒弟的營業額分成.
我自己做也有用.
我的徒弟或下線的員工.
他做出來的份額我也可以分成.
做武術教頭的.
徒弟自立門戶.
所收的徒弟也被視為自己的名下.
關德興,壽桃,梁家麟.
我開過武館.
將我師父也放在一起.
這叫開枝散葉.
我教出來的學生.
他日如果他們事奉有成.
我也會在其中沾光.
他的成就也會是我的成就的一部份.
無論如何.

$^{521}$保羅相信哥林多教會是上帝劃定給他的事奉範圍.
他在這裡勞苦.
他的事奉成果也在這裡被評定.
所以他將哥林多教會視為自己的戰訓.
他是心安理得的.
那些超級使徒就不行.
他們的成績單和哥林多教會的關係有限.
大部分的成果也不是他們勞碌得來的.
如果藉著哥林多教會為他們自己誇口.
他們就等於褻奪他人的成果.
在這裡我們可以對界限做一些思考.
Metron這個字.
界限這個字.
可以譯作地界.
確實有地域的意思.
好像黑社會說地盤一樣.
不過主要不是指地域.
不是指地域.
而是指著一個事工.
一個群體.
或者是上帝指定給我們.
要我們去承擔的對象.
你可以說是一個指定.
是用這個指定來評定我們的成績標準.
就好像耶穌的比喻裡面說.
主人交託五千,兩千和一千給不同的僕人.
這個托付就成了評定他們將來是不是良善忠心的程度的標準.
每個人所派定的界限.
即是評估標準都不一樣.
有人五千,有人兩千,有人一千.
按材受績.
有人多些,有人少些.
關鍵的是我們在我們所領受的界限裡面是不是盡忠.
每個人的地界是不是一早就確定了,固定了.
不可以增減呢.
我們先不說預定論問題.
根據耶穌的比喻所說.
我們如果在一個小事上忠心.
上帝就會派定我們承擔更大的事.
這個《路加福音》十九章十二到二十七節所說的.

$^{561}$所以這些派織是會按時,按需要,按表現而有所增減.
保羅在這裡更加提出一個新的角度的想法.
他說他的事奉界限是可以多次的,可以倍增的.
如果哥林多信徒都能夠信心增加.
信仰生命,成長,事奉而心不減.
他們就成為了保羅的延長的外延.
他們能夠去保羅去不到的地方.
做保羅做不到的事.
哥林多信徒延續傳福音的使命.
帶領更多人成為主的門徒.
信心一直增長.
保羅就自覺他的事奉邊界隨之擴大.
並且是母子境向外伸展.
弟兄姊妹.
這個道理教訓我們一件很簡單但很重要的事.
在教會裡面做人的工作.
永遠是最根本性,最長遠性.
耶穌基督公開事奉生涯只有三年半.
哪怕他再努力走遍各城各鄉傳揚天國的福音.
他能夠走的地方仍然有限.
傳講的對象仍然有限.
但他呼召並且訓練了十二門徒.
讓他們繼續他的使命.
結果門徒和門徒的門徒.
和門徒的門徒的門徒.
歷代的門徒就將福音傳遍了天下.
我們在某一個事奉崗位裡面有優秀的表現.
這個當然是好事.
但如果二十年來只有我夠資格在這個崗位之上.
那就不是太好的事了.
盡力栽培下一梯隊的人才.
讓自己可以被替代.
這個才是健康的情況.
弟兄姊妹我們在教會裡面多做人的工作.
少做事,多做人少做事.
在任何時候任何崗位.
都要栽培他人成長做目標.
是這樣的旺盛的界限才能夠開展.
我不知道的.
所以不是我故意的.

$^{601}$今天我們領唱的小組是一個新的配搭.
初登場.
你覺得是好事嗎?.
我們教會繼續有新的人被發掘出來.
擺上不同的事奉的祭.
這是非常非常精彩的事.
無論如何.
保羅強調他沒有藉著事奉去建立個人王國的企圖.
他不是要藉著事奉招收更多徒子徒孫.
壯大自己的門派.
建立他的江湖霸業.
保羅強調一切都是屬於上帝.
如果他誇口.
只能夠指著主去誇口.
17到18節.
但誇口的當子就主誇口.
因為蒙月立的不是自己稱許的.
而是主所稱許的.
這個教訓.
保羅曾經與哥林多信徒講過.
在哥林多前書一章31節.
這個教訓也不是保羅自己發明的.
是轉引自耶利米書九章23到24節.
保羅強調他沒有自誇.
哥林多信徒確實是他勞碌得來的成果.
這教會可以證明他的身份和成就.
這是理直氣壯的事.
但保羅同時知道.
教會雖然由他一手建立.
但起關鍵作用的不是他.
而是上帝.
唯有上帝使生命成長.
所以他只配搭了上帝的工作.
哥林多教會是上帝的工作.
最大的榮耀要歸給上帝.
人如果將上帝佔主要份額的工作.
據為己有.
聲稱這是他自己的勞碌的成績.
就是褻奪上帝的榮耀.
當然人在參與上帝的工作中.

$^{641}$可以分享其中的成果.
不過佔的份額多少.
表現多少.
成就多少.
就不是看工作本身.
因為你無法分別.
你不可以猜我做的事奉.
上帝佔七成功勞.
我佔三成.
你怎樣分辨出來.
你無法分辨.
你唯一可以做的.
我根本不計算有什麼可以分享.
我也不看在工作中.
我有多重要.
我只看一件事.
就是上帝開金口讚我.
上帝怎樣讚我.
我就得著多大的成績.
你明白我說什麼嗎.
你自己不可以計算.
黃浸有多大成績.
無法計算.
一點都計算不到.
是沒有.
你只能看上帝怎樣讚我.
他說我有就有.
他說我多好就多好.
只能看上帝開金口.
這是我們.
想確定我們成就.
唯一的方法.
上帝說的才算數.
他說有就有.
多就多.
不要自我計算.
不要自我評定.
絕對不要評定.
我們不要說生人.
還是說死者.

$^{681}$志強弟兄.
今天在天父懷裡.
他要接受耶穌基督的評斷.
評斷的標準不是.
他拿.
死之前.
我們放在他的棺材裡.
燒了他幾年旺贈的成績單.
然後他拿成績單.
在上帝面前養.
省一點.
每樣都不關他的事.
教會什麼時候扣了堂.
多少人接受水禮.
這些成績全部屬於上帝.
一樣都不可以.
說自己不可以.
他只能夠聽上帝.
耶穌基督對他的評斷.
耶穌基督會不會說.
你就是又良善.
又忠心的僕人.
這個才是他.
唯一能夠獲得的成績單.
得著主的稱讚.
只有他的稱讚.
是算數的.
並且只有他的稱讚.
才能帶來終極性的果效.
但誇口的.
當之主誇口.
因為蒙月立的不是自己稱許的.
而是主.
所稱許的.
這是我們每一個人的事奉.
才能得到終極的指標.
耶穌基督對我們的評斷.
稱讚永遠是最關鍵.
無論如何.
我們今天學了最少兩個功課.

$^{721}$第一是地界的觀念.
我們必須在上帝面前.
特別是在基督計劃裡面.
發現上帝派定我們.
參與和承擔的位置.
用耶福所書的說法.
這是我們所承受的基業.
我們知道自己奮鬥的目標.
或存活的理由.
我們有長期毀身服侍的對象.
亦有將來.
賴以在基督審判台前.
交帳的成績單.
弟兄姊妹.
我不能代你們去確認地界.
這需要你們自行.
從上帝那裡尋求和認定.
例如說.
有人認定他的工作.
他的專業.
他的生意.
他的事業.
是他榮耀上帝的地方.
他認真做好一個物理治療師.
一個社會工作者.
一個保險從業員的角色和責任.
盡力在每一個場景.
榮神益人.
這是上帝畫定給他的地界.
亦有人將家裡.
特別是子女的成長.
年老父母的照顧.
患病家人的長期承擔.
視為上帝畫給他的地界.
當然亦有人將實踐.
某一個政治和社會的理想.
完成某一個人民.
或者是藝術的境界.
視為他終身的志業.
視為上帝畫給他的地界.

$^{761}$都是可以的.
不一而足.
總的來說.
我不覺得我們需要將地界.
完全畫定在.
刻意的福音使命.
或者教會的事奉之內.
是嗎?.
我們每個人都有一條命.
上帝派定給我們的責任.
那就是我們的地界.
上帝給我們的託付.
那就是我們的地界.
這個地界的觀念.
我覺得其實是很正.
很正.
聖經其實.
講到上帝給我們的恩典的時候.
亦都講到一個portion.
一個份.
是我杯中的份的說法.
上帝給我的.
我認定.
這是我值得.
是我當佔的份.
上帝派定給我的工作.
給我的任務.
我亦都視.
這是上帝給我的份.
我的家庭是上帝給我的恩典.
但我的家庭亦都是上帝給我的duty.
是嗎?.
我疼我的老婆.
不是你們疼我的老婆.
對不對?.
每個人疼自己的老婆.
這是上帝派給我的任務.
我努力做好我的任務.
這個地界.
界限的觀念.

$^{801}$我覺得是一個非常好的提醒.
我們每個人bear in mind.
知道我們自己.
活著一個很重要的責任.
承擔好上帝給我的份.
我們不一定需要有光光烈烈.
很清楚的.
大意象.
很多弟兄姊妹.
像我這樣都是小人物.
我們簡單的.
認定我的生活是上帝.
給我的恩典.
給我的責任.
又認定旺鎮是我屬靈的家.
我全心全意參與在其中.
我做好主日學.
團契每一個角色.
行山的每一個角色.
都是很重要.
為其他人的信主成長守望.
像保羅說到恩賜的時候.
按我們所得的恩賜各有不同.
說義言就當照著信心的程度說義言.
執事就當專一執事.
教導就當專一教導.
勸化就當專一勸化.
施捨就當誠實.
治理就當殷勤.
憐憫人就當甘心.
這些好像很瑣碎很微小的事奉.
是很重要的.
弟兄姊妹.
我真的要代教會多謝你們每一個.
沒有你們認定.
你們毀身承擔的每一個崗位.
旺朕是沒辦法發展到今天的.
所以我們每個人.
都有或大或小的地界.
你認定毀身的.

$^{841}$就是你的地界.
第二.
除了地界的觀念之外.
我們也認識了.
栽培別人成長.
而擴展地界的觀念.
我們努力傳福音.
努力事奉.
但我們也要盡力去發掘.
栽培人才.
讓更多人參與傳福音和各種的事奉.
教會才會不斷地擴展.
旺朕的地界才可以與日俱增.
在外面社會.
所謂企業承傳的問題.
我們教會不說企業.
但教會也一樣說傳承.
我們不僅僅是自己做.
我們是栽培更多人.
可以做更多的事.
我們自己.
我們教會的地界.
然後才可以不斷地擴展.
這個課題其實有很多值得.
深思的地方.
我們日後才會再說.
我們自己的地界.
也讓我們看到我們地界擴展的可能性.
\newpage



\section{雅各書 1:19-27}
\label{sec:VDCXuQp5CNk}
\textbf{2021年7月11日崇拜直播｜陳冠成牧師｜生命見真章系列 - 聽道為了行道｜雅各書1\_19-27}
\newline
\newline
連結: \href{https://youtube.com/watch?v=VDCXuQp5CNk}{\texttt{https://youtube.com/watch?v=VDCXuQp5CNk}} ~~~~ 語音日期: 2021-07-19
\newline
\newline
\hyperref[sec:bhMQ0ofdb_Q]{\small{< < < PREV SERMON < < <}}
~
\hyperref[sec:index_chronic]{\small{[返順時目]}}
~
\hyperref[sec:rb0qE8Nt22M]{\small{> > > NEXT SERMON > > >}}
\newline
\newline
雅各書 1:19-27
\newline
\begin{longtable}{cl}
\hline
\hline
章節 & 經文 (和合本修訂版)\\
\hline
1:19 & \begin{tabularx}{0.7\textwidth}{X} 我親愛的弟兄們，你們要明白：你們每一個人要快快地聽，慢慢地說，慢慢地動怒， \end{tabularx} \\ \\ \relax
1:20 & \begin{tabularx}{0.7\textwidth}{X} 因為人的怒氣並不能實現神的義。 \end{tabularx} \\ \\ \relax
1:21 & \begin{tabularx}{0.7\textwidth}{X} 所以，你們要除去一切的污穢和累積的惡毒，要存溫柔的心領受所栽種的道，就是能救你們靈魂的道。 \end{tabularx} \\ \\ \relax
1:22 & \begin{tabularx}{0.7\textwidth}{X} 但是，你們要作行道的人，不要只作聽道的人，自己欺騙自己。 \end{tabularx} \\ \\ \relax
1:23 & \begin{tabularx}{0.7\textwidth}{X} 因為只聽道而不行道的，就像人對著鏡子觀看自己本來的面目， \end{tabularx} \\ \\ \relax
1:24 & \begin{tabularx}{0.7\textwidth}{X} 注視後，就離開，立刻忘了自己的相貌如何。 \end{tabularx} \\ \\ \relax
1:25 & \begin{tabularx}{0.7\textwidth}{X} 惟有查看那完美、使人自由的律法，並且時常遵守的，他不是聽了就忘，而是切實行出來，這樣的人在所行的事上必然蒙福。 \end{tabularx} \\ \\ \relax
1:26 & \begin{tabularx}{0.7\textwidth}{X} 若有人自以為虔誠，卻不勒住自己的舌頭，反欺騙自己的心，這人的虔誠是徒然的。 \end{tabularx} \\ \\ \relax
1:27 & \begin{tabularx}{0.7\textwidth}{X} 在神—我們的父面前，清潔沒有玷污的虔誠就是看顧在患難中的孤兒寡婦，並且保守自己不沾染世俗。 \end{tabularx} \\ \\
[1ex]
\hline
\hline
\end{longtable}
$^{1}$這次真是讀經的時候了.
《雅各書》第一章十九至二十七節.
新約聖經三百二十一頁.
接著請聽我讀.
聽道和行道.
我親愛的弟兄們.
這是你們所知道的.
但你們各人要快快的聽.
慢慢的說.
慢慢的動腦.
因為人的怒氣並不成就神的義.
所以你們要脫去一切的污穢和盈餘的邪惡.
全溫柔的心領受那所栽種的道.
就是能救你們靈魂的道.
第二十二節.
只是你們要行道.
不要單單聽道.
自己欺控自己.
因為聽道而不行道的.
就像人對著鏡子看自己本來的面目.
看見走後.
隨即忘了他的相貌如何.
唯有詳細察看那全被使人自由之律法的.
並且事上如此.
這人既不是聽了就忘.
乃是實在行出來.
就在他所行的事上必然得福.
二十六節.
若有人自以為虔誠.
卻不勒住他的舌頭.
反欺控自己的心.
這人的虔誠是虛的.
在神我們的父面前.
那清潔沒有玷污的虔誠.
就是看顧在患難中的孤兒寡婦.
並且保守自己不沾染世俗.
.
今天我們繼續第二次宣講《雅各書》.
其實剛才讀的經文.
我想大家都很熟.

$^{41}$聽道行道.
《雅各》提醒我們要快快聽.
慢慢說.
慢慢動腦.
其實是說如何聽上帝的道.
如果你有留意.
可以打開程序表的廣東大綱.
十八節這裡也說了.
剛才我們沒有讀的.
即是上一節.
他說「他」即是上帝.
「上帝按自己的旨意.
用真道生了我們.
叫我們在他所做的萬物中.
好像初熟的果子」.
然後.
剛才我們一開始十九節所聽到的經文就是.
《雅各》繼續說.
這是讀者.
即是每一個信徒.
已經知道的事情.
就是上帝要透過他的真理建立我們.
並且他繼續勉勵讀者.
當他們去領受真理的時候.
要快快的聽.
慢慢的說.
慢慢的動腦.
不過先停在這裡.
有少少問題我們要搞清楚.
理論上每一個基督徒都渴慕聽不到上帝的真理.
對不對?.
會不會有些信徒或基督徒.
因為聽見真理.
就會發怒?.
有沒有?.
有沒有試過?.
聽到上帝的真理會動了怒氣.
會發怒.
可能會有.
否則《雅各》就不用叫我們.

$^{81}$在領受真理的時候.
要慢慢的動腦.
對不對?.
事實上我們有這些例子.
你看聖經其實有沒有留意.
最典型的.
我一說你就知道.
若拿.
若拿是否一個.
慢慢的動腦的人.
當他聽到上帝的吩咐.
不是的.
他很快就動腦了.
如果你記得約拿書說到.
上帝猜險若拿才知他來尼美城.
來宣告一個審判的訊息.
是吧?.
但一開始若拿是左閃右避的.
死都不肯去.
為什麼?.
他是不是害怕?.
他是不是害怕去到.
亞述敵人的國家.
去到敵人的陣地.
還要說一些不吉利的話.
害怕自己生命有危險?.
他一點都不害怕.
他害怕的是什麼?.
為什麼不去?.
他害怕的是.
假如當他們聽到上帝的真理.
萬一他們真的接受悔改.
那又如何?.
那上帝就會赦免他們.
他最怕,最不想見到這件事.
但當然事實往往是怎樣?.
你最害怕見到的東西.
通常真的會出現發生.
於是乎.
尼里未成當聽到若拿的宣講之後.

$^{121}$全城上下的人都真的披麻蒙灰.
來向上帝認罪悔改.
結果上帝也赦免他們.
沒有像原本的想法那樣.
降災來毀滅他們.
若拿開心嗎?.
有人得救.
很奇怪吧?.
有人得救,有人接受真理.
反而一個先知會不開心.
因為他認為.
這一班人這麼壞.
應該怎樣?.
該死的.
不應該同情,可憐,原諒他們.
但上帝竟然放過這班壞人.
若拿就覺得.
上帝的道理很不合理.
宗教不是這樣計算的.
很不公道.
若拿覺得自己先有道理.
於是乎是怎樣?.
他向上帝發脾氣.
大大的發怒.
你看到人在聽見真道.
遇見上帝的安排時.
他不一定是開心的.
可以是動怒氣.
另一些實際的例子.
我們也看過了.
新約主耶穌.
在四福音中.
當耶穌來尋找甚麼人.
尋找的他自己也說.
我不是來找義人.
我是來找罪人.
於是乎你常常見到.
耶穌在公開場合.
會和稅吏,罪人交往.
甚至吃喝.

$^{161}$哪些人不喜歡?.
明明很好.
真理得到彰顯.
福音得到講傳.
但哪些人不開心?.
當時的宗教領袖不滿意.
為甚麼?.
覺得耶穌在破壞規矩.
挑戰我們所信奉的傳統.
所以他們不但不會快快地聽.
簡直是甚麼?.
不聽.
簡直是掩耳.
對耶穌的真理完全聽不入耳.
甚至他們發怒.
發怒到一個地步.
要怎樣?.
要讓耶穌消失.
所以聖經有很多例子.
都說到.
或者對比今日.
當人在社會生活得久.
打滾得越久.
很多時就會累積了.
一些很牢固的想法.
或者價值觀念.
當中有些是好.
有些亦都不是太好.
但這樣就越容易.
跌落一個試探.
就是他會不自覺地.
自己也有尺子.
他會用自己的尺子.
不單量度他身邊的人.
還會量度甚麼?.
量度上帝說的話是否合聽.
他會用他的主觀意願.
來到我們常說的選擇性收聽.
甚至進行.
那些自己認為合理的聖經教導.

$^{201}$那些覺得合聽的上帝身輔.
很多時如果當上帝.
不按照他的主觀意願.
來行事安排的時候.
又或者上帝一些真理教導.
真的衝擊著他一些想法.
一些行為的時候.
就會怎樣?.
就會馬上啟動防衛機制.
我們都很明白.
一啟動防衛機制.
一是選擇對抗真理.
不跟.
又或者甚至乎.
好像剛才我們所說的例子.
他會發怒.
用怒氣,憤怒的方式.
來回應上帝.
正如《雅各書》.
如果你還記得上次.
我們說第二節的時候.
雅各就勸.
我的弟兄們.
你們落在百般的試驗中.
都要怎樣?.
都要以為大喜樂.
你覺得當你聽到這句說話的時候.
或者你想起當時的情景.
是否所有的讀者.
每個信徒聽到這段說話.
落在百般的試驗的時候.
都要大喜樂.
你會不會馬上說.
你覺得會不會每個人都是這樣?.
大概不會.
似乎未必每個人聽到這番真理.
都會馬上欣然去接受.
特別當假如真是當事人.
他真正經歷很大的試驗.
真是落入一個苦難的裡面.

$^{241}$他可能不單單不會多謝雅各.
反而會怎樣?.
用一些我們今天有時會說的術語.
吐槽.
向雅各吐槽.
即是說.
覺得你唱高調而已.
受苦的又不是你.
你當然是風涼.
你教的東西.
有點違反人情常理.
不太接地氣.
於是就會怎樣?.
反彈.
甚至乎憤怒.
所以其實雅各在這裡提醒我們每一個.
面對上帝的真理的時候.
他希望我們懂得快快的聽.
慢慢的說.
慢慢的動腦.
在這裡快和慢.
不是說速度的意思.
是說我們接收上帝的真理.
接收聖經教導的時候.
的一種態度應該是怎樣.
快其實背後的意思是說.
你要敏銳地聽.
細心地聽.
慢其實是說.
你要謹慎.
你不要急著急.
這樣的意思.
雅各提醒我們.
第一方面要留心聽上帝的道.
先聽完.
先聽清楚.
如果不明白.
可以再聽.
再問.
或者你讀過一些輔導課程.

$^{281}$你也會知道.
當聆聽的時候要怎樣?.
要複述.
對方說完之後.
你要複述是否這個意思.
以致你不會聽錯,聽漏.
上帝的道一樣.
你反覆探索一下.
看看自己有沒有聽錯,聽漏.
誤解了上帝的真理.
這是第一樣雅各提醒我們的.
第二方面就是.
要謹慎回應聖經的教導.
意思是你不要急於批評.
甚至論斷.
那些你認為已經有點不合時宜.
好像不太合時宜.
或者很離地的聖經教導.
你覺得好像不切實際.
但是慢慢來.
不要急於下定論.
甚至你和上帝爭拗.
甚至談判.
我經常覺得有一句提醒很重要.
我們經常說上帝是永活的上帝.
對不對?.
應該沒有人會反對.
什麼叫做永活的上帝?.
永活是否他永遠存在活著?.
當然是.
好像欺騙大家一樣.
大家不出聲.
怕我安裝彈弓.
永活的上帝固然是.
他所說的永遠的活著.
昔在今在永在.
是不是只是活著這麼簡單?.
其實當我們說上帝是永活的時候.
包括他的話語.
他的教導都是永活的.

$^{321}$意思是什麼?.
以前有用.
現在有用.
將來永遠都有用.
上帝的道都是永活的.
不會過時.
不會勾的.
是不是?.
有可能我們覺得不中聽.
不舒服.
是上帝搞錯.
還是我們只不過這一刻.
未能完全領悟和明白呢?.
這就是雅各提醒我們.
所以要謹慎回應.
你知道上帝不會搞錯.
只不過可能我們在這一刻.
未能完全理解背後的意思.
所以.
雅各提醒我們.
不要急於對聖經的教導.
太快就有定案.
不要急於下結論.
你認為上帝這樣吩咐我.
是不可能的.
不合理的.
我怎麼做呢?.
你要愛那些迫害我的人.
不合理的.
那些人這麼壞.
你還要我愛他,幫他.
不合理的.
我做不到.
不要這麼快下定論.
這是第二方面.
第三方面.
雅各提醒我們.
慢慢的動路.
意思是好好控制你的情緒.
控制你的脾氣.

$^{361}$特別當我們一些人性的陰暗面.
被真理光照的時候.
當上帝吉中你一些死穴的時候.
吉得你很痛,很不舒服的時候.
你很想發怒.
很想去抗衡真理的時候.
甚至很想和上帝對著幹.
好像若拿一樣.
上帝叫你去東.
你就去西的時候.
實際上有這個方法.
有沒有聽過六秒原則?.
就是你嘗試深呼吸.
忍住一口氣六秒.
深呼吸不是這樣.
不是像跳跳虎那樣.
不是那種.
這樣代表什麼?.
這樣代表準備作戰.
如果你用胸腔來呼吸.
你試做一次就知道.
即是準備去again作戰.
相反你是複式呼吸.
用肚子向下呼吸.
唱歌的人會明白.
向下吸深一點.
自然就會怎樣?.
還要忍六秒.
其實我們人體的機制.
如果你有聽過一顆叫杏仁核.
真的像杏仁那麼大顆.
在你的腦裡面.
那顆東西是管理你的情緒.
原來你深呼吸.
數六秒就會讓你這顆杏仁核.
沒那麼活躍.
沒那麼容易生氣.
會減低你一些壓力的荷爾蒙.
你就會舒緩到.
沒那麼快.

$^{401}$沒那麼容易被自己有怒氣.
但不會爆出來.
讓自己冷靜一點的時候.
就提醒自己.
沒錯我真的對上帝的教導.
對聖經真理感到懷疑.
感到不公道甚至激氣.
我不明白上帝為什麼要我這樣做.
我不明白為什麼上帝要這樣.
安排一些我不喜歡的人.
或者事情發生在我身上.
但是我知道上帝沒有搞錯.
一定有好的東西在.
一定有你的心意等著我去領受.
這就是雅各所說.
萬萬動腦的意思.
不是說不可以有情緒.
是怎樣控制你的情緒.
以至到像雅各所說.
先處理好你的情緒.
你的脾氣.
這樣我們才有空間去成就上帝的義.
來進行祂的吩咐.
活出祂對我們的期望.
所以雅各這裡勉勵我們.
你要好好預備自己.
他用一個當時都明白的比喻.
你要像耕田一樣.
你就是那塊田.
你要好好預備你那塊心靈的田.
成為一個好的土壤.
就是對上帝存著一個.
柔和謙卑的態度.
你樂意接受他的管教.
因為你知道.
可能會不中聽不舒服.
但一定對你好.
並且好像二十一節所說.
不單止要柔和謙卑.
還要除草.

$^{441}$就是將你.
裡面可能有些上帝不起悅的品行.
一些慾望.
好像拔雜草一樣.
你有決心拿走它.
事常被上帝的話語.
來光照你的生命.
讓上帝的道.
真的可以在你身上.
暢通無阻.
這樣才能夠在我們身上紮根成長.
才可以發生改變.
好像雅各這裡說.
拯救我們的身心靈.
當然.
這一段經文.
快快的聽 慢慢的說 慢慢的動腦.
很自然也可以進一步應用在.
一些人際關係上.
我們都會明白之所以.
很容易和人起爭執.
起衝突.
其實很多時是什麼問題?.
是接收的問題.
是聽的功課.
做得未夠完善.
我們很多時會選擇去接收.
對方說什麼.
我們最想聽到我們喜歡的說話.
又或者我們反過來.
我們已經有既定的假設.
來聽對方說話.
我們覺得他一定是這樣.
我們帶著這份執著.
來按自己主觀的偏見.
來解讀對方的意思.
明明都不是這樣.
不過可能你已經有個前設.
所以無論對方怎樣說.
你都覺得他一定是一般般.

$^{481}$甚至乎一旦.
對方的言詞真的擊中你.
觸碰到你的情緒按鈕時.
你一啟動就怎樣?.
你就開始不想和對方理性地傾談.
甚至乎你腦內會有很多.
憤怒尖銳的詞彙拿出來.
你準備還擊.
於是不用我再說.
情況一發不可收拾.
不知道上面這些畫面.
有沒有出現在我們身上.
我相信多多少少都會有.
因為我們都是塵土血氣.
試回想一下.
上一次你和人爭執.
甚至引起衝突.
是什麼時候?.
記不記得?.
是什麼地方?.
在教會?.
還要是在事奉的時候?.
在辦公室?.
甚至在家裡?.
和誰一起發生衝突?.
和你的同事?.
還是姐妹?.
你的子女?.
你的父母?.
還是你的配偶?.
當時有沒有說了一些.
其實不應該說的話?.
有沒有因為這樣.
不經意破壞了雙方的關係.
多了嫌隙?.
不單止沒辦法成就上帝的義.
甚至以前你為上帝所做的美好見證.
就是因為那次衝突.
現在都怎樣?.
一筆勾銷了.

$^{521}$會不會因為這樣.
感到後悔呢?.
早知當日.
就不那麼衝動了.
早知那天.
就不說那句話.
真的弟兄姊妹.
我想我們都明白.
特別對未認識耶穌的朋友來說.
可能偶然一次炒一次大鑊.
就會怎樣?.
就要搞很久了.
你一次發脾氣.
就足以令他懷疑.
重新懷疑.
信耶穌是否真的這麼好?.
我聽過一個很好的故事.
話說有個企業家.
他數來都被人說行事很穩陣.
做事很穩陣.
很少爆大鑊.
他經營的公司也有很好的業績.
在他準備退休的時候.
一定有個接任人.
接任人就問他.
你幾十年這麼成功.
有什麼秘訣?.
過兩招來好啊.
他就笑笑口說.
其實真的沒有什麼特別秘訣.
不過我學會一件事.
就是.
謹記在你很生氣的時候.
不要說那麼多話.
在你很生氣的時候.
不要做一些重大的決定.
所以很感恩.
我在任內.
都沒有什麼重大的過失.
出現過.

$^{561}$都沒有和身邊的人.
吵過大鑊.
其實短短一句話.
足以讓他的繼任人.
受用一生.
同樣雅閣這裡短短.
一句說話.
快快的聽.
慢慢的說.
慢慢的動腦.
細心的聆聽.
謹慎的回應.
控制你的情緒.
都同樣讓我們可以受用一生.
事實上有很多.
研究情緒的專家都說.
他說我們一個人.
無論你學了多少東西.
你有多少學歷.
你多麼聰明.
有多少學識,經驗.
原來當你處於一個.
極之憤怒的狀態的時候.
你猜他的智商剩下多少歲?.
你猜呢?.
剩下五歲.
人一旦處於極度的憤怒.
他以往擁有的理性.
他以往擁有的思考.
意脅.
表達能力.
他的儀態全部都會怎樣?.
剩下五歲.
以至到他所說的東西.
他所做的決定.
你不理性嗎?.
不理性.
很情緒化.
一點都不成熟.
明明他真的想做好一件事.

$^{601}$但通常出來後會怎樣?.
壞了一件事.
我們都經歷過.
所以我很喜歡聽一個故事.
話說猶太人有一個這樣的故事.
就是說.
說到亞伯拉罕.
他有一次請客.
請一個外邦的人.
到他的帳篷裡吃頓飯.
但他沒有想過.
一邊吃飯.
對方是外邦人.
當然是捧自己所信的信仰.
當然是踩低亞伯拉罕所信的上帝.
一邊批評上帝.
是怎樣怎樣.
亞伯拉罕聽到後.
當然是越來越生氣.
結果就惱怒罵罵.
不吃了.
你走吧.
逐客令趕走他.
到了晚上.
上帝又向亞伯拉罕在夢中顯現.
跟亞伯拉罕說.
亞伯拉罕.
其實剛才那個人.
我忍了他幾十年.
我一句怨言都沒有發.
我內心在等他回頭.
你只不過和他相處了幾十分鐘.
你就生氣得受不了.
失去儀態.
怎樣繼續成為上帝兒女?.
怎樣為我作見證?.
弟兄姊妹是真的.
下次假如當.
有人氣得你跳來跳去的時候.
不妨提醒自己.

$^{641}$上帝可能都忍了這個人很久.
但上帝都怎樣?.
繼續忍.
上帝都在等他.
接受福音.
回轉.
回頭.
我又何必太在意.
這一刻.
我自己的得失是什麼呢?.
我們要謹記.
衝突最少要多少人才成事?.
多少個?.
兩個人.
除非你懂得自己和自己衝突.
那些是另一個主題.
今天不說了.
最少兩個人才成事.
即是說沒有你幫忙是吵不成架.
是嗎?.
想想是嗎?.
你看看身邊.
如果是你的配偶.
沒有你幫忙是吵不成架.
你可以跟他說.
所以下次你不妨試下.
當對方越急的時候.
你會怎樣?.
你就越豬油膏越慢.
是嗎?.
當對方越大聲.
你就越細聲.
當對方越多話.
就越少話.
這樣通常.
吵不成架.
因為有趣的人吵架.
就好像照著鏡子.
他看到對方越生氣.
他怎樣?.

$^{681}$他唯有升級.
他就越生氣.
你越多話.
他就越多話.
但他看到你不出聲.
越少話.
很平和.
他會忍嗎?.
支持不了很久.
吵不成架.
所以真的,弟兄姊妹.
盼望雅各這個.
我們很熟悉的教導.
快快的聽.
慢慢的說.
慢慢的勞.
真的能夠幫我們.
締造和睦.
成就上帝的義.
我們繼續看餘下的經文.
雅各在這裡表示.
到底上帝的道.
到底真理有沒有在我們的生命裡.
產生效用呢?.
有個指標的.
看看你有沒有行動.
就知道一清二楚了.
雅各說.
第二二至到.
二十四節這裡.
和本的翻譯.
焦點在聽道和行道那裡.
是兩個動作行為的比較.
但原本是什麼意思呢?.
我印了一段經文.
在中間那段.
新譯本也好一點.
我想請大家讀一讀.
二十二至二十四節新譯本.
預備起.

$^{721}$(上帝的語言).
好,謝謝.
其實你留意了.
特意括了出來的.
其實他的焦點不是純粹說.
聽道,行道.
然後叫我們離身.
好像不太關自己事.
你覺得哪樣好一點.
他在這裡說的是.
聽道的人和行道的人做一個比較.
然後他.
雅各.
要我們選擇.
你想做一個什麼樣的人.
他要強調的是.
一方面信徒要傳柔和謙卑的心.
來接收聆聽上帝的道.
但同時.
不可以只聽.
他也要我們成為一個行道的人.
有行動的.
將聽見的真理實踐出來.
否則他會說.
你騙自己.
你是自欺.
但是.
人真的有時會有錯覺.
就是你頭腦上認同某個道理.
有時你會以為.
你會成為這樣的人.
覺不覺得?.
很有趣的.
最常見的例子就是.
我們絕大部份的人都很想.
都認同做運動.
保持健康.
保持體態的健康健美.
好不好?.
好!沒理由不好.

$^{761}$在在當中都有人在進行.
所以你會去收集一些資料.
坊間哪些方法會好一點.
你是想自己.
變得健康一點.
變得健美一點.
沒有人不想的.
但他不會因為我們只聽.
只看.
只想著.
你就會自動變瘦變健美.
對不對?.
如果是就好了.
你見我就知道了.
是不是?.
你覺得我有沒有定期提醒自己.
要做運動要保持健康.
絕對有.
你覺得我有沒有經常想著.
瘦20磅就好了.
好像剛剛入職黃針的時候.
有沒有想著?你覺得我.
是不是?.
晝夜思想.
但你看到我的身形.
你就知道我的行動很有限.
我用最低的成本.
吃少一點.
是不是?.
即是幻想.
這樣應該都可以.
是吧?.
少少少少這樣撕走脂肪.
你看到假如我只是.
一路說一路想.
一路想著沒有行動.
沒有改變我的生活習慣.
我只是怎樣?.
自欺.
就等於我上網站.

$^{801}$塗一副叫什麼?.
瘦身鏡.
知道什麼叫瘦身鏡嗎?.
你去時裝店那塊就是瘦身鏡.
有沒有發覺你在那裡照.
總是漂亮點,瘦點.
一回到家.
咦?.
為什麼還了?.
是不是我眼花呢?.
是那塊鏡幫你瘦了.
有的.
你打瘦身鏡這三個字.
有大把可以選擇.
如果我只是買一塊.
我很想瘦身.
我很想保持健康.
但我只是不想付出代價.
我用最低的成本.
買一塊瘦身鏡回來日照夜照.
頂多都只是怎樣?.
哄自己開心.
但對身體沒有實際幫助.
所以雅各用照鏡的比喻.
他用得很好.
他說.
信徒若只是聽到.
不行到.
其實就好像他對著鏡子.
看到自己被罪污染.
那個不完善的本相.
但又怎樣?.
什麼都不做就走了.
很快就忘記了.
剛才自己看到的是什麼.
忘記自己實際是什麼狀況.
請問.
他用照鏡來做什麼呢?.
他用鏡子來做什麼呢?.
雅各正在問這件事.

$^{841}$周禎如.
我想今早大家有沒有照鏡?.
大部分都有.
你每天睡醒.
出門口之前.
上崇拜之前.
你會照鏡子.
還要不止照一次.
照幾次.
照鏡子是為了什麼?.
看看自己的儀容.
妥不妥當.
看看眼屎洗乾淨了沒有.
牙膏漬有沒有擦乾淨.
頭髮有沒有亂.
女士們可能化妝好不好.
看看自己妥不妥當.
妥不妥當就撿起來.
男士可能半身鏡就夠了.
女士又怎樣?.
一定要有一塊全身鏡.
看看自己的衣著打扮怎樣.
覺得起碼過到自己.
這樣才出門口見人.
否則.
我們也會問自己.
那我要一塊鏡來做什麼?.
如果我照完之後什麼都不做.
當看不見.
那我照來做什麼?.
所以雅各這裡說得很好.
什麼是我們那塊鏡?.
什麼是我們人生生命那塊鏡?.
他說就是上帝那些.
傳遞使人自由的律法.
上帝的真理就是那塊鏡.
就是照清我們本相那塊鏡.
上帝的道能夠監察我們的生命.
將我們那些隱而未見的幽暗.
都通通照出來.

$^{881}$所以雅各他勉勵我們.
在聽上帝的道的時候.
不要水過鴨背.
不要聽完之後.
又忘記了.
而是要像他所說.
你要去.
花時間花功夫.
你要去詳細察看.
你要願意花時間.
可能有些弟兄姊妹.
你會馬上做筆記.
然後回去思想.
或者你有本聖經.
你抄下來回去再論.
你要花時間.
將聖經的教導反覆思考.
你要讓它沉澱.
你要讓它消化.
好像吃東西一樣.
要點時間.
你要咀嚼.
要消化.
慢慢地將上帝的道吸收.
真正成為你生命的養分.
如果不是.
就像雅各這裡所說.
照完鏡.
看到自己一些盲點問題.
沒事了,走吧.
那照來做什麼?.
並且雅各這裡說.
我們要時尚.
不單將上帝的道吸收.
消化.
還要身體力行.
實踐上帝的道.
簡單說就是有入.
也有出.
這樣雅各說.

$^{921}$這樣才是蒙福.
否則會怎樣?.
頂多你都是自肥.
你只是吸收上帝的道.
但不去實踐.
你只是自肥.
也不算是蒙福.
已故的丹麥神學家.
奇克果.
他在說雅各書這段教導時.
有一個很精闢的提醒.
他說.
你不要只看著鏡子.
而忘記看鏡中的你自己.
明白嗎?.
你不要只看著鏡子.
「這鏡子很美」.
「這鏡子照得很清楚」.
但不看鏡中的你.
他提到我們有一個盲點.
很多時候我們花很多精力.
去研究這鏡子是怎樣的.
這鏡子用什麼來做?.
為何它照得這麼清晰?.
但忘記了在照鏡中.
看到上帝想你怎樣改變.
看到你自己的本相.
所以頂子妹這段很大的提醒.
當我們研讀上帝的話語.
當我們查經的時候.
不是純粹研究一下.
這鏡子有多美.
聖經有多好.
有多博大精深的道理.
是說我看聖經,查經.
聽到我要生命改變.
我要生命改變.
我要成長.
我要配合上帝.
這才是我們照鏡.

$^{961}$才是我們領受真理.
真正的目標.
事實上.
其實你和我很明白.
你家裡不會因為有鏡子.
你自動會變美.
是嗎?.
對不對?.
我們有不少鏡子.
我們買多些回來.
會否自己變美了?.
不會的.
你買多幾本聖經.
屬靈參考書.
你不去消化,實踐.
不會特意屬靈了.
又或者說.
我們不會不斷照著鏡子.
照日照夜照.
就照美了自己.
不會的.
那塊是瘦身鏡.
我們一定要一路照.
一路執整自己.
才會越來越妥善.
越來越祈禮.
正如我們要一路查考聖經.
一路反覆思想.
一路改變自己.
實踐真理.
我們才會越來越有上帝的形象.
所以.
雅各提醒我們.
我們查考聖經.
不是純粹為了增長知識.
更重要的是.
我們要去反省.
實踐我們的生命.
雅各說.
這樣才真真正正體會到.

$^{1001}$什麼是上帝的福氣.
雖然在過程中.
你可能會不舒服.
你會覺得辛苦.
但是上帝說.
這樣是有益的.
所以.
弟兄姊妹.
雅各將兩個.
將兩個選項.
將兩條路.
現在擺在我們眼前.
到底是選擇.
認真地聽道行道.
成為一個蒙福的人.
還是選擇.
欺騙自己.
只聽.
只看.
但是不實踐呢?.
相信我們.
懂得怎麼選擇.
相對我們.
知道要怎麼選擇.
對不對?.
好了.
看完最後兩節.
在雅各這裡說.
原來有人可能會跌入盲點.
他會出現一些錯誤的假設.
他會以為.
我覺得自己敬虔就是敬虔.
我覺得自己虔誠.
就是虔誠.
但是他沒有實質虔誠的內涵.
雅各這裡說.
這樣同樣等於.
欺騙自己沒有分別.
是不是?.
因為.

$^{1041}$一個人到底在上帝面前.
是否虔誠.
是誰決定的?.
是誰定的標準?.
是否人定的標準?.
很多時候.
未信的人就會覺得.
一個基督徒什麼是虔誠?.
他每個星期都上崇拜.
他定期上教會有很多生活.
他一定是一個很虔誠的人.
這個也可能是.
但是.
真正的標準取決於上帝.
在上帝眼中.
他的敬虔標準是怎樣.
這個才重要.
所以承接著之前的教導.
雅各就引入一些實際的生活例子.
來講明.
到底上帝眼中.
虔誠標準是什麼.
首先雅各說到.
其中一個很實際的表現.
就是控制自己的舌頭.
那些可能沒什麼益.
幫不了人.
造就不了人的話.
傷害人的話.
就不要說了.
並且要清楚自己.
說話背後的動機.
你是否真的想榮耀上帝.
祝福身邊的人.
以致我們可以謹慎自己說話的內容.
不單止內容.
好的內容.
還要加上好的語氣.
好的速度.
是吧.

$^{1081}$我們都明白.
內容對.
有時也會炒大鑊的.
要有適當的語氣和速度.
另一方面.
慎言控制舌頭.
更加包括在某些情況下要怎樣.
保持克制.
閉口不言.
有時不說話才是智慧.
我有很多這些體會.
有時對著小朋友也會很生氣.
於是乎我聽過一個教導.
挺好用的.
我想對所有人都適用.
用你的眼神.
是吧.
用你的眼神.
平時我對著小朋友.
和顏如色的多數會.
依起兩個面珠燈.
又說又笑.
突然有一刻.
你這樣看著他.
小朋友很容易觀察到.
突然你平時笑嘻嘻.
經常開心.
突然很嚴肅.
而不夠嚴肅.
可以拿走眼鏡.
增加你的殺氣.
是吧.
然後他就會知道.
咦?不妥.
變了.
你不說話更害怕.
為什麼?.
什麼事呢?.
他開始會用這方面去想.
糟了,是否我做錯事.

$^{1121}$惹事惹爸爸生氣.
他開始有所謂的警醒.
他知道你正在生氣.
但不知道生氣什麼.
他就會開始定了.
你不用很大發怒.
你知道你表達了你的不滿.
但怎樣?.
好像聖經所說.
生氣,但不犯罪.
他又收到你的不滿.
於是他吃飯也會吃快一點.
馬上去收拾玩具.
事實上,弟兄姊妹.
我們很明白.
中國人都說.
病從口入.
就什麼?.
禍從口出.
或者說.
《箴言》也說.
多言多語難免有過.
禁止咀唇.
是由自慰.
人有時說話.
往往反映你是怎樣的人.
從你的說話就知道.
你是否一個敬畏上帝的人.
所以弟兄姊妹求上帝.
求祂的靈幫我們.
好像傳道書所說.
說話有時.
但靜默也有時.
有智慧地.
在適當的時候說適當的話.
讓我們的言語.
真的能夠成就上帝的意.
另一方面.
在上帝眼中的敬虔表現.
是說我們有沒有關懷.

$^{1161}$社會弱勢的需要.
特別是阿葛這裡說到.
孤兒寡婦.
就是當時那些沒有能力.
沒有足夠的能力照顧自己的群體.
在當時的社會沒有所謂.
福利基金,社會保障.
我們也明白.
孤兒寡婦必須依靠別人幫助.
否則她難以為生.
阿葛這裡提醒我們.
做基督徒.
不是一個隱居自肥的群體.
不是不問世事的.
他是要活在人間.
他要和他的靈社結緣.
阿葛提醒我們.
對社會弱勢的關懷.
是每一個領受福音的人.
一種納己,納人的表現.
一定會有.
阿葛說這樣.
也是一種敬虔的表達.
當然.
阿葛也提醒我們.
在關懷社會需要的時候.
我們同時也要保持警醒.
不要讓自己沾染了世俗.
不要跟這個世界完全打成一片.
不要被世俗裡的.
似是而非的歪理影響.
阿葛提醒我們謹記.
基督徒除了關心社群外.
他更加有一個很重要.
很特殊的身份和使命.
就是我們是福音的大使.
代表耶穌將基督傳遍開去.
我們做一切背後的目的.
也是為了這樣.
所以弟兄姊妹今天.

$^{1201}$不是純粹說孤兒寡婦.
今天要面對的可能是.
很多單親家庭的出現.
你知道越來越嚴重.
你見到不少.
可能是爸爸或媽媽.
要自己獨力照顧小朋友.
要面對不同的生活挑戰困難.
教會如何跟他們同行.
或如何跟身邊的福音機構.
善用社會資源的配合.
來輔助他們呢?.
這既是上帝給我們今天.
給教會一個挑戰.
但反過來看.
也是一個很好的福音機遇.
所以你會留意到.
這幾年上帝一直帶領我們這樣走.
一直開展了很多對基層的服事.
當中你看到其實有成效.
有不少的論社透過教會.
可能不單止獲得物資的供應.
他們的心靈得到關顧.
他們知道旺角有基督徒很關心他們.
他們知道信耶穌的人很愛他們.
甚至耶穌也很愛他們.
更重要的是.
上帝通過這些愛心服事.
讓得救的人數加添給我們.
所以.
弟兄姊妹盼望上帝在我們中間.
興起更多的人.
去回應祂的道.
回應祂的呼召.
謹記並不是上帝不夠人用.
所以要多找些人出來服事祂.
實情是怎樣?.
就好像今天這段經文所說.
上帝是很想把福氣給我們.
不單止給我們的論社.

$^{1241}$是給我們.
給教會,給你自己.
上帝很希望我們透過聽到祂的道之後.
然後實踐出來.
透過聽到,行到.
真真正正領受到.
從上而來的福氣.
讓我們同心祈禱.
上帝我們多謝你.
再一次藉著你自己的教導.
甚至我們以熟能詳的教導.
你喚醒我們.
讓我們重新對著上帝.
你真理的鏡子.
來審查自己的生命.
上帝幫助我們每一個人.
來如雅各所說.
快快的聽,慢慢的說,慢慢的動腦.
讓我們有一個柔和謙卑的心.
來受教.
讓我們敏銳聆聽上帝的呼聲.
讓我們多花時間沉澱,消化.
不要太急於回應.
甚至批評,挑戰上帝的真理.
更加讓我們好好控制.
我們的情緒,我們的脾氣.
因為雅各所說.
人的怒氣不能夠成就上帝的義.
所以上帝求你幫助我們.
我們真的能夠去.
柔和謙卑來到你的面前.
我們知道我們會有軟弱.
並且我們知道上帝是完全的.
或許有些安排,有些教導.
我們真的一時三刻.
我們覺得不明白.
我們覺得不合理.
我們覺得做不到.
但是求上帝幫助我們.
讓我們真的能夠慢慢的明白.

$^{1281}$慢慢的一路一路走近你.
順服進行你.
上帝我又求你保守我們同人的關係.
在建立關係的過程中.
同樣學懂敏銳去聆聽.
學懂控制自己的言語.
保守自己.
可能會生氣.
但是卻不犯罪.
上帝又求你讓我們可以實踐.
將今天我們所領受的.
將上帝給予我們一些意象.
一些你所託付給我們的人.
我們懂得去回應.
我們起來去回應.
並且我們深信.
上帝既然讓我們看到這些圖畫的時候.
你都已經有供應.
我們只要靠著你.
我們只要順著你的帶領的時候.
我們就能夠榮耀上帝你.
也能夠祝福身邊的人.
並且讓我們自己都成為無福的人.
為此我們將每一個弟兄姊妹交托.
願你興起我們.
來到去聽道行道.
願你聽我們的禱告.
奉基督耶穌說.
你得勝的名字來到祈求.
Amen.
\newpage



\section{哥林多後書 11:1-15}
\label{sec:rb0qE8Nt22M}
\textbf{2021年7月18日崇拜直播｜梁家麟牧師｜保羅為個人身份的自辯｜哥林多後書11\_1-15}
\newline
\newline
連結: \href{https://youtube.com/watch?v=rb0qE8Nt22M}{\texttt{https://youtube.com/watch?v=rb0qE8Nt22M}} ~~~~ 語音日期: 2021-07-19
\newline
\newline
\hyperref[sec:VDCXuQp5CNk]{\small{< < < PREV SERMON < < <}}
~
\hyperref[sec:index_chronic]{\small{[返順時目]}}
~
\hyperref[sec:n5CNurE7uu0]{\small{> > > NEXT SERMON > > >}}
\newline
\newline
哥林多後書 11:1-15
\newline
\begin{longtable}{cl}
\hline
\hline
章節 & 經文 (和合本修訂版)\\
\hline
11:1 & \begin{tabularx}{0.7\textwidth}{X} 但願你們容忍我小小的愚蠢；請你們務必容忍我。 \end{tabularx} \\ \\ \relax
11:2 & \begin{tabularx}{0.7\textwidth}{X} 我以神嫉妒的愛來愛你們，因為我曾把你們許配給一個丈夫，要把你們如同貞潔的童女獻給基督。 \end{tabularx} \\ \\ \relax
11:3 & \begin{tabularx}{0.7\textwidth}{X} 我只怕你們的心偏邪了，失去那向基督所獻誠懇貞潔的心，就像蛇用詭詐誘惑了夏娃一樣。 \end{tabularx} \\ \\ \relax
11:4 & \begin{tabularx}{0.7\textwidth}{X} 假如有人來，傳另一個耶穌，不是我們所傳過的；或者你們另受一個靈，不是你們所受過的聖靈；或者接納另一個福音，不是你們所接納過的；你們居然容忍了！ \end{tabularx} \\ \\ \relax
11:5 & \begin{tabularx}{0.7\textwidth}{X} 但我想，我一點也不在那些超級使徒以下。 \end{tabularx} \\ \\ \relax
11:6 & \begin{tabularx}{0.7\textwidth}{X} 雖然我不擅長說話，我的知識卻不如此。這點我們已經在每一方面各樣事上向你們表明了。 \end{tabularx} \\ \\ \relax
11:7 & \begin{tabularx}{0.7\textwidth}{X} 我貶低自己，為了使你們高升，因為我白白地傳神的福音給你們，難道這算是我犯了錯嗎？ \end{tabularx} \\ \\ \relax
11:8 & \begin{tabularx}{0.7\textwidth}{X} 我剝奪了別的教會，向他們取了報酬來效勞你們。 \end{tabularx} \\ \\ \relax
11:9 & \begin{tabularx}{0.7\textwidth}{X} 我在你們那裡有缺乏的時候，並沒有連累你們一個人，因為我所缺乏的，那些從馬其頓來的弟兄都補足了。我向來凡事謹慎，將來也必謹慎，總不要連累你們。 \end{tabularx} \\ \\ \relax
11:10 & \begin{tabularx}{0.7\textwidth}{X} 既有基督的真誠在我裡面，在亞該亞一帶地方就沒有人能阻止我這樣自誇。 \end{tabularx} \\ \\ \relax
11:11 & \begin{tabularx}{0.7\textwidth}{X} 為甚麼呢？是因我不愛你們嗎？神知道，我愛你們！ \end{tabularx} \\ \\ \relax
11:12 & \begin{tabularx}{0.7\textwidth}{X} 我現在所做的，將來還要做，為要斷絕那些尋機會之人的機會，不讓他們在所誇耀的事上被人認為與我們一樣。 \end{tabularx} \\ \\ \relax
11:13 & \begin{tabularx}{0.7\textwidth}{X} 那樣的人是假使徒，行事詭詐，裝作基督的使徒。 \end{tabularx} \\ \\ \relax
11:14 & \begin{tabularx}{0.7\textwidth}{X} 這也不足為奇，因為連撒但也裝作光明的天使。 \end{tabularx} \\ \\ \relax
11:15 & \begin{tabularx}{0.7\textwidth}{X} 所以，他的差役若裝作公義的差役也沒有甚麼大不了。他們的結局必然跟他們的行為相符。 \end{tabularx} \\ \\
[1ex]
\hline
\hline
\end{longtable}
$^{1}$我要為大家讀出的經文是在聖經新約的哥林多後書第十一章一至十五節.
是大家手上的聖經新約二百五十三到二百五十四頁.
新約哥林多後書十一章一至十五節.
聖經這樣說:但願你們寬容我這一點慾望.
其實你們原是寬容我的.
我為你們起的憤恨 原是神那樣的憤恨.
因為我曾把你們許配一個丈夫.
要把你們如同貞潔的童女獻給基督.
我只怕你們的心或偏於邪.
失去那向基督所傳純一清潔的心.
就像蛇用鬼咋又挖了下蛙一樣.
假如有人來另傳一個耶穌.
不是我們所傳過的.
或者你們另受一個靈.
不是你們所受過的.
或者另得一個福音.
不是你們所得過的.
你們容讓他也就罷了.
但我想我一點不在那些最大的使徒以下.
我的言語雖然粗俗.
我的知識卻不粗俗.
這是我們在凡事上向你們眾人顯明出來的.
我因為白白傳神的福音給你們.
就自居卑微 叫你們高聲.
這算是我犯罪嗎?.
我規附了別的教會.
向他們取了功架來給你們效力.
我在你們那裡缺乏的時候.
並沒有累著你們一個人.
因我所缺乏的.
那從馬其頓來的弟兄們都補足了.
我向來凡事謹守.
後來也必謹守.
總不至於累著你們.
既有基督的誠實在我裡面.
就無人能在雅歌亞一帶地方阻擋我者自誇.
為什麼呢?.
是因我不愛你們嗎?.
這有神之道.
我現在所作的 後來還要作.

$^{41}$為要斷絕那些尋機會人的機會.
使他們在所誇的事上也不過與我們一樣.
那等人是假使徒 行事詭詐.
裝作基督使徒的模樣.
這也不足為怪.
因為連撒旦也裝作光明的天使.
所以他的差役.
若裝作仁義的差役.
也不算稀奇.
他們的結局.
必然照著他們的行為.
.
各位聽眾平安.
真的很希望平安在我們每一個人的生活.
在我們每一個人的心裡面.
在遙遠向在泰國曼谷.
沙敦堂的弟兄姊妹打個招呼.
今早看新聞.
也看到泰國的疫情.
心中也有些忐忑.
為你們祈禱了很短的一段時間.
希望你們在那邊一切平安.
我們繼續看哥林多後書餘下的篇章.
在上面一講.
保羅已經開始為他的仕途身份做辯護.
然後來到這裡.
從一章一節到十二章的十三節.
保羅做了一件令人詫異的事.
就是他用了非常長的篇幅.
為他的愚拙誇口.
所以不少人稱這半章的聖經.
為愚人宣言.
一個愚蠢的笨人的自白.
這些背景的事情其實以前已經講過了.
保羅對哥林多教會心情很複雜.
一方面他非常愛他們.
這是他一手建立的教會.
他不可以不愛惜.
教會再有問題他也不可以置之不理.
人家可以劈炮他不可以的.

$^{81}$他怎樣也要跟他們糾纏下去.
另一方面.
這間教會確實也為他帶來太多的傷害.
愛心就傷害同樣心.
他建立的教會.
這班信徒竟然背叛他.
轉投另外的傳道者.
並且搬其他人的權威去打壓他.
因此保羅在自辯的時候.
忍不住滋生一些負面的情緒.
覺得自己實在很蠢.
死蠢死蠢死蠢.
在十一章他提自己蠢.
提了至少四次.
為什麼我墮落到這樣的田地.
要在子女面前為自己的身份和使命去做辯護.
我們以前也說過這些.
在那麼多信徒乃至其他人的眼中.
保羅的信仰保羅的生活確實是非常奇葩的.
絕對不令人羨慕.
他決意接受所傳的道.
所信奉的耶穌基督帶來的規範.
要他放棄人間的智慧和能力.
以愚厥者的姿態生活和事奉.
很彰顯上帝的大能.
他不可以流露自己的書.
他不可以顯示他聰明.
他要用愚厥的面貌來彰顯上帝的大能.
你覺不覺得這個難度很高.
特別對我們這些忍不住也要露一手.
也要顯示自己風光的人.
要強行禁止自己是不是很難.
保羅說福音規定了他不可以凸顯自己.
再加上保羅實際上客觀.
他整個事奉生涯都確實是倒霉加倒霉.
苦難不絕.
表面上看真是一個被咒詛的人生.
而他們還好意思到處說一個使人蒙福的福音.
真是奇怪.
一個以苦為甜.

$^{121}$以咒詛為祝福的思想和實踐.
會不會真的被人覺得一個瘋了的人才有的反應.
保羅雖然他自己不斷為自己的愚厥去做辯護.
聲稱為基督的緣故成為愚厥.
講了那麼多前書四章十節.
不過保羅也不只是自說自話.
他心底裡是雪亮的.
他真的知道其他人的確對他最渺渺.
在前書四章十三節他也說.
直到如今人還把門看作世上的污穢萬物中的渣子.
他真的不知道別人看他是什麼樣子.
不會那麼奇怪吧.
我自我感覺良好的同時.
我真的知道其他人捏著鼻子.
根本不想坐在我的位置.
你坐地鐵那麼擠迫.
你坐下去之後.
左邊右邊各兩個位置都沒有人坐.
你都知道什麼事了.
人們當然是當你是麻瘋病.
他沒理由不知道.
最痛苦的當然是.
他所愛的兒女竟然嫌棄他.
被嫌棄的感覺實在傷害太大.
舉個例子.
你們有沒有試過.
心底裡流露出來.
嫌棄你媽媽很蠢.
有沒有試過.
千萬不要舉手.
如果真的被你媽媽知道.
你會有一點這樣的感覺.
她一定大受傷害.
你想想你明白我說什麼嗎.
就是這樣.
是嗎?.
是呀,我是很蠢.
不過我養大你.
你讀到博士又如何.
對不對?.

$^{161}$是我供你讀到博士的.
對不對?.
保祿的感覺就是這樣.
兒女嫌他死蠢.
那種受傷害的感覺是很大.
既然保祿是一個死蠢的人.
所以他這次就來一次.
愚拙大盤點.
將自己種種愚拙的行徑.
表列成為一張愚拙清單.
充分以一個愚拙者自居.
以愚者為榮.
這就是愚者之言.
由十一章一節一直到十二章的十三節.
我們分三次去看完這一大段聖經.
我們先看十一章的一到十五節.
背景這段聖經不需要多說.
保祿藉著第二封信.
一個措辭嚴厲的信.
已經迫使哥林多教會.
懲治了那些犯罪的人.
基本上將最嚴重的問題解決了.
不過解決了眼前的大問題.
不等於解決了他和哥林多教會的關係.
因為哥林多教會仍然有弟子不喜歡保祿.
不喜歡繼續不喜歡.
特別保祿你竟然用殺手簡迫我們去懲治人.
我們不可以不做.
但心裡不順不開心就一定有.
老實說我們都知道.
有人喜歡誰有人不喜歡誰.
你是沒辦法迫的.
如果真的在一個團契裡有些人不喜歡你.
那你有什麼辦法.
有人不喜歡我就不喜歡.
我不可以迫他喜歡我.
這是沒辦法做的.
不過這種對立猜忌的氣氛.
就傷透了保祿的心.
反對保祿的人.

$^{201}$由於不可以再挑起任何.
教義上的事端.
於是就只能夠搞一些關係上的離間.
他們就吹捧一些曾經在他們中間服侍過的.
其他的傳道人.
那些超級使徒.
誰比誰好.
做了很多這樣的格價比較.
藉此想貶低保祿.
當然他們捧出來的那些超級使徒.
不一定是真的想和保祿對著幹的.
只不過是被他們夾了出來.
但是總之就是發生這樣的事.
王震有些弟兄姊妹不喜歡陳明全牧師.
和執事會做的決定.
於是他們就自己出去.
去徵詢其他大目的意見.
有人從吳忠文牧師聽到一個看法.
有人從陳欣明牧師聽到另一個看法.
總之就和王震執事會做的決定.
剛剛好不一樣.
他們就搬那些不同的說法來.
去砌.
你看看人家的大目名目都不是這樣看的.
你們為什麼做這麼蠢的事.
為什麼做這麼蠢的決定.
更糟糕的是.
那些反對我們教會的領導層的人.
還專門邀請那些超級使徒.
重訪這間教會.
他們曾經在這裡服侍過.
現在叫他們回來.
那些超級使徒見到有人邀請他們.
當然就來.
他們不知道背後的內容.
他們不知道他們來到.
就等於為反對派撐牆.
壯大了反對派的聲勢.
吳小學者就十一章第四節.
他們說假如有人來.

$^{241}$傳另一個耶穌不是我們所傳過的.
或是你們受另一個靈.
另受一個靈不是你們所受過的.
或是另得一個福音不是你們所得過的.
這個說如果或是.
其實也可能是說.
真的有實際的那些反對保羅的人.
因為保羅根本不是國民教會.
他們專門找回阿波羅.
找回某些大使徒來.
他們中間坐鎮說話.
更加振振有詞.
你看看保羅很差的.
我們現在找回來的大目.
看法一不一樣.
這樣就造成保羅被迫要和那些超級使徒正面對決.
陳明泉牧師大戰陳欣凌牧師.
我小時候.
晚上我下班回家.
他們走來跟我說.
爸爸我們星期六不如去海洋公園.
不如我們去吃自助餐.
我第一個反應不是.
為什麼要去.
我第一個反應是.
你們問過媽媽沒有.
媽媽怎麼看.
你明白我說什麼嗎.
如果媽媽說了話.
我還有什麼好意見.
你告訴我.
你想家裡世界大戰嗎.
不會的.
就是這樣.
老老實實.
總之媽媽說了話.
我的意見就是她的意見.
我不會還有其他意見.
我沒理由在家裡製造兩種聲音.
兩個意見.

$^{281}$這是公道道.
同樣很多時候都會有其他教會的弟兄姊妹.
走來跟我說教會的問題.
我總是要很謹慎的.
先聽你們教會怎麼決定.
為什麼這樣決定.
不要多多事幹.
不要自覺自有智慧.
不斷指手畫腳.
說說怎麼做.
被人指控就無謂.
反正我也不覺得自己比其他人有智慧.
這個我們要明白.
如果不是牽涉到絕對真理.
我必須要說.
我們所做的所有決定.
都有相確的餘地.
永遠有的.
要不要多買一層.
要不要開分堂.
我們永遠無法有百分百同意.
或者只覺得這是唯一對的.
但你搬弄不同的權威.
來互相抗衡.
結果是怎樣.
結果除了分裂之外.
是沒有其他可能的.
所以保羅面對的是這樣的局面.
所以保羅要奮起為自己的仕途身份.
他自己的Postulate去辯護.
他要為自己的資格.
自己的條件誇口.
甚至要直接批評.
那些跟他意見不一致的超級仕途.
前面保羅才剛剛說到.
他不可以自誇.
只能夠以主誇口.
十章十七節.
如今他又以自己誇口.
他自己也知道.

$^{321}$他前言不對後語.
自相矛盾.
所以他覺得很不好意思.
所以他首先就道歉.
但願你們寬容我這一點的愚拙.
其實你們原是寬容我的.
十章一節.
寬容我這一點的愚拙.
應該就翻譯成為.
寬容我這一個愚拙的表現.
其實你們原來是寬容我的.
應該就翻譯成為.
希望你們多多包容.
不是說知道他們已經寬容.
整句意思是.
我這麼愚蠢的表現.
希望你們原諒.
我但願你們能夠原諒我.
保羅有什麼愚拙的表現呢?.
需要哥林多教會原諒他呢?.
就是他要將自己和那些假師傅.
或者超級仕途去做比較.
證明他更加值得哥林多信徒信賴和跟從.
這個彷彿就好像兩個男人.
為了爭奪一個女人.
使盡手段爭風吃醋.
保羅接著在十一章二節說.
我為你們起的憤恨.
原是上帝那樣的憤恨.
因為我曾把你們許配一個丈夫.
要把你們如同貞潔的童女獻給基督.
保羅承認他的表現.
其實憤恨應該翻譯成是嫉妒.
我是有些嫉妒.
這個嫉妒保羅馬上就解釋.
其實和上帝對以色列民的嫉妒一樣.
在舊約聖經上帝明講他是一個忌邪的上帝.
忌邪其實是jealous.
真是上帝是一個嫉妒的上帝.
上帝是不喜歡他的子民.

$^{361}$三心兩意見異思遷一腳踏兩船.
你效忠他就只能效忠他.
你不可以只心思思有件事是不行的.
保羅說我的嫉妒和上帝的嫉妒其實一樣.
上帝是忌邪的上帝.
嫉妒的上帝七一及二十章五節.
嫉妒的他不容忍百姓既拜上帝又拜巴力.
要求他們專一事奉耶和華.
許配一個丈夫.
保羅說他的爭女心態一如上帝.
他不可以容忍那麼多信徒.
跟那些全講錯誤道理的人相交.
是因為他要為他們預備.
將他們成為一個貞潔的新娘子.
他是父親要將女兒許配給人.
新郎是耶穌基督.
舊約多次用新郎和新娘去比喻耶和華和以色列的關係.
通常一做新郎新娘的比喻.
這個比喻的引伸教訓.
就是以色列民意要保持貞潔.
基本上就是在講貞潔的問題.
一講結婚.
因為以前還是嚴守一夫一妻制.
結婚就不可能有外遇.
所以荷西亞書1到3章.
以色列書16章.
以色列書50章.
講到婚姻基本上都是在講貞潔.
新約保羅在《依法所書》第5章.
講到教會和基督的關係.
比輔為丈夫和妻子.
你記不記得他還要強調什麼.
要使妻子聖潔無玷污.
要強調這一點.
其實一講到婚姻.
就是在講貞潔.
啟示錄用高揚的新婦.
去比喻那些最終得勝的基督徒.
都是在講這樣的話.
愛情總是排他主義.

$^{401}$沒有嫉妒就沒有真愛.
連上帝都表達對他的子民.
移情別戀的嫉妒.
何況保羅.
在3到4節.
保羅說.
我只怕你們的心或偏於邪.
失去那向基督所傳順一清潔的心.
就像蛇用鬼渣誘惑了夏娃一樣.
假如有人來傳另一個耶穌.
不是我們所傳過的.
或你們另受一個靈.
不是你們所受過的.
或是另得一個福音.
不是你們所得過的.
你們容讓他也就罷了.
先做一個註解.
最後那句話.
你們容讓他也就罷了.
也是錯譯的.
我不想講舊約.
不想講和學本的很多錯譯.
但這段聖經是.
因為保羅整段經文.
是一個他從來未試過的風格.
就是很多諷刺.
有骨的.
遞遞下的說話.
所以這句說話.
你聽到也覺得沒可能.
你們容讓他也就罷了.
其實真正的意思是.
這樣你們都忍得住?.
這樣你們還能忍得住?.
其實不是鼓勵他們容忍.
是質疑你們為什麼能容忍.
如果有人跟你們講另一個福音.
傳另一個耶穌.
你這樣都忍得住?.
不是吧?.

$^{441}$是不是?.
保羅說的是.
我怕你們的心.
偏斜失去對基督的純權清潔的心.
就好像蛇用了鬼柵誘惑了夏娃一樣.
保羅擔心.
那些信徒經受不到試探.
於是受欺騙.
失去聖潔誠實的心.
這說到那些超級使徒所傳的另一個耶穌.
另一個福音.
這是很激進的說法.
是不是真的他們在傳另一個福音呢?.
是什麼福音呢?.
很多學者就將初期教會所有能夠找到的異端都搬出來.
那些是不是洛斯底主義的呢?.
是不是幻影說的呢?.
是不是什麼以便利派的呢?.
是不是那些猶太人說的政治彌賽亞呢?.
前面說的幾樣東西如果你不明白都無所謂的.
因為根本上都是異端.
異端不用明白.
總之是不是這個是不是那個.
另一個福音.
如果你記得.
保羅在《皆泰書》其實都曾經做過同樣的指控.
說彼得或者其他一些人在傳另一個福音.
不過《皆泰書》說的叫別的福音.
不是說另一.
其實是一樣的意思.
別的福音.
《皆泰書》我們知道的問題是.
當時期有一些猶太教化的基督教.
那些主張一個基督徒要先做猶太人.
先接受各禮.
先守所有的禮儀.
然後才可以成為基督徒.
這個可以叫做巴勒斯坦化的基督教.
《皆泰書》有這樣的問題.
但是哥林多教會應該就沒有這個問題.

$^{481}$因為《皆泰》在小亞細亞就比較猶太化.
但是哥林多其實是在希臘半島.
在希臘半島.
外邦人為主.
所以我就不相信哥林多.
有很多人是要主張要人先守各禮.
先做猶太人.
然後才做基督徒.
應該這個不是在哥林多教會出現的問題.
我相信保羅在這裡所說.
那些說人是另一個耶穌.
另一個福音.
不應該是說那些超級使徒.
是在傳說徹底的異端的基督論.
異端的救贖論.
另一個耶穌就是異端的基督論.
另一個福音就是異端的救贖論.
不應該是.
我相信真正的問題.
是對教義的.
對信仰的.
生活應用的問題.
就不是教義本身有差異.
否則保羅在哥林多書信.
應該更嚴詞的去批判這班超級使徒.
不會像這裡說的那麼婉轉.
用我們今天生活的例子.
我們是福音信仰者.
我們一直反對那些財富和健康的福音.
那些Wealth and Health的Gospel.
我們反對那些Prosperity Theology.
那些成功神學.
嚴格來說.
成功神學的人所傳的耶穌和福音.
和我們都沒有分別.
他們都相信耶穌基督的神人異性.
道成肉身.
他們沒有反對耶穌基督的替死待贖.
所以在基督論和救贖論其實沒有差異.
但成功神學只是宣稱.

$^{521}$由於在人間活得貧窮和謙卑的耶穌.
如今已經升到高天.
享受最高的尊榮.
所以今天我們要跟隨的是.
已經上了高天尊榮的耶穌.
我們不開最少的Benz.
怎樣能夠配得上尊榮的耶穌呢?.
對不對?.
我們不穿一件名牌西裝.
怎樣能夠配得上.
已經上了高天.
擁有一切榮華的耶穌呢?.
他在應用上.
做了一些我們不能夠接受的觀點.
所以成功神學的錯謬處.
不是基督論和救贖論的教義.
而是今天我們信奉的耶穌的形象的描繪.
和接受福音所要帶來的人間後果.
如果我們說成功神學是傳另一個耶穌.
另一個福音.
我們反對的不是基督論和救贖論的教義.
而是教義的應用.
我相信很多教會裡面的保羅反對派.
他們所推崇的超級使徒.
也可能有成功神學相類似的主張.
例如他們認為成功的傳道者.
需要用成功的面貌.
用顯揚的手法去事奉.
彰顯他們的智慧和能力.
好讓人心悅誠服.
他們有些人也認為.
基督徒應該從舊約的律法釋放出來.
所以不再受舊約的.
或者眾多宗教的道德規條所限制.
可以隨心所欲做事.
所以哥林多前書也說到.
他們有些人做了一些雙方敗得的事.
他們認為無所謂的.
走肉穿腸過餓.
佛祖心中留.

$^{561}$做多少壞事都不影響我們的命命.
他們有這樣的想法.
保羅就和他們格格不入.
保羅傳講一個愚拙的十字架的道理.
保羅雖然不主張要嚴守舊約的律法.
但是他對道德的部分是很執著的.
要求基督徒活出高水準的道德行為.
所以很可能保羅對福音的應用.
和這班超級使徒不一樣.
保羅跟著和超級使徒做比較.
第五節.
但我想我一點不在那些最大的使徒之下.
最大的使徒其實是一個諷刺性的說法.
不一定是地位上最大的.
所以你不用想到是西門彼得.
而且保羅用了眾數.
我們的聖經也譯為那些.
所以不是說一個.
最大應該是一個.
比較大的可以有很多個.
最大應該是一個.
所以修定本和合本也不再譯為最大的使徒.
而是超級使徒.
那些可能不是指一個人.
而是指好幾個人.
保羅說我和任何人比.
我都不覺得我彼得差過他們.
跟著他說.
第六節.
我的言語雖然粗俗.
但我的知識卻不粗俗.
這是我們在凡事上向眾人顯明出來的.
保羅承認他的言語粗俗.
言語粗俗除了說他講使者之言之外.
主要是指著他的口才不及人.
因為以前的人.
古人和我們不同的.
今天你們聽我說到.
要求我接地氣.
要求我講一些人話.

$^{601}$但你不知道上一輩不是這樣.
上一輩想著上台.
應該講一套特別的語言.
這樣才顯示出率領.
各位親愛的弟兄姊妹.
我們今天來到上帝的台前.
連語調也不同了.
你不知道嗎.
你沒聽過嗎.
一定要用特別的語調.
講那些四六文.
平時你不會講的東西.
這樣才顯示出率領.
以前的人演說是講很特別的言語.
所以其實講保羅粗俗.
不一定是講粗口.
而是講他不用那些.
懂得演說的表達方法.
保羅可能沒受過演說訓練.
可能更重要.
保羅確實是口才不太好的人.
保羅是嘰吉的.
賣相不好.
講話不流暢.
這是保羅先天的限制.
可能是因為這個緣故.
保羅有很多吃虧的地方.
還有就是.
這裡再提到一點.
保羅沒接受江浦多教會給他的錢.
這也成為了江浦多教會.
批評保羅的其中一點.
7-11節.
我以為白白傳上帝的福音給你們.
就自居卑微叫你們高升.
這算是我犯罪嗎.
我歸附了別的教會.
向他們取了功架來給你們效力.
我在你們那裡缺乏的時候.
並沒有累著你們一個人.

$^{641}$因我所缺乏的.
那從馬其頓來的弟兄們都補足了.
我向來凡事謹守.
後來也必謹守.
總不至於累著你們.
既有基督的誠實在我裡面.
沒有人能在雅歌亞一帶地方阻擋我者自誇.
為什麼呢.
只因我不愛你們.
只有上帝知道.
這裡提到了兩個地方.
一個叫馬其頓.
一個叫雅歌亞.
我們記得這是希臘半島兩個恆省.
哥林多是雅歌亞省的首府.
而菲勒比就在馬其頓.
有些哥林多的信徒認為.
保羅之所以不接受他們的金錢.
而選擇接受馬其頓的教會.
顯示保羅是大世超.
不愛哥林多.
只愛馬其頓的教會.
保羅為此赴願.
保羅說我非常愛哥林多教會.
我不想接受你們的錢.
是因為我不想構成對你們的負擔.
不想引起任何的誤會.
我因為白白傳上帝福音給你們.
總不至於累著你們.
這段經文提到累著.
我們相信.
這裡主要指的意思不是財務上的包袱.
因為哥林多教會挺有錢的.
哥林多教會有挺多中產信徒.
所以其實要哥林多教會拿錢支持保羅.
也不會使得他們經濟有困難.
累著主要不是指經濟上.
而是觀念上.
保羅擔心收他們的錢.
就會產生紛爭.

$^{681}$惹來誤會.
叛倒他們.
很多學者相信.
保羅不接受哥林多教會的錢.
卻是接受馬其頓教會長期的援助.
這裡指的馬其頓教會是指菲律賓教會.
你看腓立比書一章五字就知道了.
保羅說他們長期支持他.
為什麼保羅要菲律賓教會的錢.
不要其他教會.
包括哥林多教會呢.
是因為保羅是一個很另類的傳道人.
他自設向不同地方的人傳福音.
他沒辦法長期留在一個地方.
而唯有就他的經驗裡.
只有菲律賓教會有宣教的心胸.
不計較他們奉獻支持的傳道人.
沒有在他們中間工作.
這個看腓立比書四章十五到十七我們可以看到的.
其他教會通常都是.
我出錢請你.
我當然希望你留下事奉.
你到處走.
不知像什麼樣子.
所以那些教會不接受保羅這樣的情況.
保羅當時的情況.
今天其實都很普遍.
講一個我們的生活經驗.
我們神學院要訓練宣教士.
其實面對很多困難.
我們通常要求一個宣教系的同學.
在畢業之後.
先在本地教會事奉兩到三年.
然後才接受警察會.
警察派出去宣教.
為什麼呢.
因為宣教都是做報道直堂.
都是做牧羊.
問題是你去到宣教工場.
通常都不會有資深的牧師督導你.

$^{721}$你自己要獨當一面.
你只是在神學院讀完哪會認識.
所以最好就是在這間教會.
例如陳明泉牧師那班.
一路督導他.
讓他可以學會.
然後才出去.
所以我們通常都要求一個畢業生.
先找一間教會.
事奉兩三年.
做儲備宣教士.
不過我告訴你.
我們過去幾十年的經驗.
是很難找到的.
90\%以上的教會都不肯.
我請一個畢業的傳道人.
我希望他坐定幫我們教會.
做兩三年.
剛剛熟絡了的人就走.
對我們教會有什麼好處.
傳道人走馬燈.
對教會只有壞處沒有好處.
我所認識的宣教系學生.
碰壁處處.
十間有八間不肯.
除非回自己的母會.
基本上都找不到的.
沒有教會的肯定.
就是這個心態.
你說他錯我也不覺得錯.
總之就是這樣.
以我為主.
我們自己的考慮.
這個是很難.
所以頂尖會.
我們肯不肯做蝕本生意.
宣教就是超級蝕本.
將最好的人推出去.
又出錢又出力.
你覺得很蠢嗎.

$^{761}$真的很蠢.
保羅接受菲律賓教會的援助.
不接受哥林多教會一分錢.
不過這樣的想法.
這樣的做法.
就被哥林多教會指責.
他們和哥林多教會的關係很生疏.
保羅說我向來凡事謹守.
後來也不必謹守.
保羅說我在錢銀上面.
我很有原則.
我拿著我自己的原則去做.
我不是一個三心兩意的人.
最後12到15節保羅批評.
他說那班超級仕途.
有些好像裝扮成光明的天使的撒旦.
然後保羅說.
我現在做的事.
我將來也要做.
為了斷絕那些尋機會的人的機會.
簡單說最後這句說話的意思.
保羅說的是.
我現在在書信上罵他們.
我的立場不是一時一樣.
我保羅預告過他很快會再去哥林多.
當我去到哥林多的時候.
保羅說我不客氣的說.
我會去清理門戶.
我會去把那班超級仕途趕走.
我不會讓他們繼續在教會裡面影響你們.
好.
這是整段聖經.
做一點點應用.
其實這一段聖經.
是我看保羅的書信.
特別是看哥林多後書.
唯一一段我找不到怎樣應用的方法.
完全沒有應用.
因為保羅的表現實在太奇葩了.
他得到一個根本給不到恰當的評價.

$^{801}$更加不要說找他做我們的榜樣.
整段聖經保羅都是一個自白.
坦誠表達他對哥林多信徒的負面感受.
有嫉妒有委屈有怨懟甚至憤怒.
保羅是一個性情中人.
他一方面討論一些很抽象的神學話題.
但另一方面他又直率赤裸裸的表達他對上帝對人的感情.
理性中有感性.
理性不掩蓋他的感性.
所以我真的想說.
我讀保羅的書信.
我真的不明白這是什麼人.
這個偉大的仕途.
他那個不知道他怎樣想.
他思想複雜.
但他心靈單純.
你要他講那些深奧的神學理論.
他講到天花亂墜.
要消化他講什麼都很難.
消化不良.
但是當講到宗教感情的時候.
卻又是最簡單最直接的.
在加拉提書保羅和彼得死過.
講恩信稱義的道理.
紮馬完全不賣人的帳.
但在殺氣騰騰的竊戰中間.
保羅竟然講了加拉提書二章二十節的一段說話.
我已經與基督同釘十字架.
現在活著的不再是我.
那是基督在我裡面活著.
並且我如今在肉身活著.
是身信上帝的兒子而活.
他是愛我.
為我捨己.
這八個字我趴在這裡.
保羅表達他心中對基督最深的熱愛.
他願意為這個救他.
愛他的基督奉獻一切.
他是愛我.
為我捨己.

$^{841}$保羅執實的不是耶穌基督是前人類的救主.
這些這麼宏大的說話.
他只是說耶穌愛我.
他愛我.
他是我的救主.
他是愛我的主.
他捉住的請求不是偉大的屬靈定律.
若不流血就罪不得赦.
而是基督愛他.
基督愛他.
天方地老海枯石爛.
任憑全世界只剩下他一個人.
基督都會為他犧牲捨己.
所以最刻骨銘心的不是上帝愛世人.
而是他是愛我.
他是愛我.
弟兄姊妹.
這個認定是不是你的認定.
你的信仰最深層.
最內核是不是這句說話.
什麼是我的基督信仰.
他是愛我.
為我捨己.
他是愛我.
為我捨己.
這是不是你切實的經驗.
他是愛我.
為我捨己.
是不是你的神學知識的第一句話.
是不是你屬靈生命的第一經歷.
我要說如果是的話.
你和我都會活得不一樣.
是不同樣的.
同樣的.
保羅在書信裡面.
亦毫不掩飾表達他對弟兄姊妹的關愛.
他對肥立彼的弟兄姊妹怎麼說.
我所親愛所想念的弟兄們.
你們就是我的喜樂.
就是我的冠冕.

$^{881}$你是我的月亮.
你是我的太陽.
這好像是跟情人說的話.
在替我看前書.
保羅都說我們基督者愛你們.
不但願意將上帝的福音給你們.
連自己的性命都願意給你們.
因為你們是我們所疼愛的.
保羅從來不是用一個故宮.
一個專業人士的角色.
心態去事奉.
按張工作半斤八兩.
精思分明.
上班就售賣我工作時間.
我的精力 我的專業.
回家就回復我私人世界.
做回真實的自己.
對保羅來說.
事奉是他生命的全部.
他全心全意壓在他服侍的對象身上.
毫無保留傾出一際.
連命都不要.
這個大情大性的保羅.
在哥林多後書直率的質問信徒.
在你們移情別戀.
效忠其他超級使徒.
他反覆問我做錯了什麼.
我在哪裡對不起你們.
為什麼你們這樣對我.
十一章十一節.
為什麼呢 是因為我不愛你們嗎.
只有上帝知道.
他連上帝都搬了出來.
天地共鑒 我不愛你.
不是吧.
我要說.
保羅在整段聖經裡面.
修昔對表達他的嫉妒.
他的吃醋.
我們看到一段醋意濃郁的聖經.

$^{921}$今日的職場文化.
和保羅的事奉表現.
是南轅北轍的.
我兒子在一間公營機構工作.
他告訴我.
整間公司.
那種安全距離比疫情更嚴重.
同時之間.
保持距離 維持公事關係.
午飯都不一起吃.
平日有限度交流.
公事之外.
只是一些沒有營養的社交客套說話.
絕對不講私事.
總之有交情.
沒有友情.
我.
看過不少職場書.
都在說.
絕對不要在工作裡面.
夾帶任何感情.
在政府部門尤其如此.
弟兄姊妹.
你覺不覺得保羅的事奉表現.
完全違反了.
今日我們公營的職場倫理和工作操守.
你再熱愛工作.
你都不應該將你自己.
所有時間精力感情.
甚至條命都壓在你的工作上.
有沒有誇張一點.
你不應該和你.
服侍的對象發生感情.
要保持安全距離.
你在工作的時候.
不應該有太多喜惡愛憎.
更加不應該和同工爭訪合躁.
你不可以要求.
同工信徒愛你多一點.
我要告訴你.

$^{961}$如果今日旺朕有同工.
像保羅這樣的表現.
例如兒童團契.
有個導師很疼他的學生.
希望那班學生喜歡他多於喜歡其他導師.
你怎樣寫.
如果你是管理兒童部.
你會怎樣.
很煩.
用不用這樣.
不行.
所以你問我.
怎樣應用這段聖經.
我都不懂.
難道要你學嗎.
我都學不到.
很難學.
難道我叫你們.
大家合躁嗎.
我只能說.
割斷上下文.
我苦心自問.
我講我自己.
我對耶穌基督的愛.
是否太過正常.
我對弟兄姊妹的愛是否太過節制.
我知道我做不到.
像保羅這個.
拼盡 癲狂.
愚蠢的樣子.
的樣子.
所以我真的.
只能是一個關懷者.
一個關懷者.
我不可以是一個愛的.
我不是一個愛者.
我知道我真的.
怎樣做都做不到.
保羅這種癲癲得得的地步.
我心裡是否完全.

$^{1001}$是否完全.
是否完全是這種癲癲得得的地步.
我都不敢說.
你會不會都想.
變成這樣癲癲得得的人.
很難想.
不過.
或者我們最後.
講一點正面說話.
我們可不可以要求自己.
稍微癲一點.
不用太過正常.
太過理性.
你猜我們做不做到.
求主.
感動我們.
\newpage



\section{雅各書 2:1-7}
\label{sec:n5CNurE7uu0}
\textbf{2021年7月25日崇拜直播｜陳啟興牧師｜生命見真章系列 - 不以貌取人｜雅各書2\_1-7}
\newline
\newline
連結: \href{https://youtube.com/watch?v=n5CNurE7uu0}{\texttt{https://youtube.com/watch?v=n5CNurE7uu0}} ~~~~ 語音日期: 2021-08-01
\newline
\newline
\hyperref[sec:rb0qE8Nt22M]{\small{< < < PREV SERMON < < <}}
~
\hyperref[sec:index_chronic]{\small{[返順時目]}}
~
\hyperref[sec:a56U_20AydY]{\small{> > > NEXT SERMON > > >}}
\newline
\newline
雅各書 2:1-7
\newline
\begin{longtable}{cl}
\hline
\hline
章節 & 經文 (和合本修訂版)\\
\hline
2:1 & \begin{tabularx}{0.7\textwidth}{X} 我的弟兄們，你們信奉我們榮耀的主耶穌基督，就不可按著外貌待人。 \end{tabularx} \\ \\ \relax
2:2 & \begin{tabularx}{0.7\textwidth}{X} 若有一個人戴著金戒指，穿著華麗的衣服，進入你們的會堂，又有一個窮人穿著骯髒的衣服也進去， \end{tabularx} \\ \\ \relax
2:3 & \begin{tabularx}{0.7\textwidth}{X} 而你們只看重那穿華麗衣服的人，說：「請坐在這裡」，又對那窮人說：「你站在那裡」，或「坐在我腳凳旁」； \end{tabularx} \\ \\ \relax
2:4 & \begin{tabularx}{0.7\textwidth}{X} 這豈不是你們偏心待人，用惡意評斷人嗎？ \end{tabularx} \\ \\ \relax
2:5 & \begin{tabularx}{0.7\textwidth}{X} 我親愛的弟兄們，請聽，神豈不是揀選了世上的貧窮人，使他們在信心上富足，並承受他所應許給那些愛他之人的國嗎？ \end{tabularx} \\ \\ \relax
2:6 & \begin{tabularx}{0.7\textwidth}{X} 你們卻羞辱貧窮的人。欺壓你們，拉你們到公堂去的，不就是這些富有的人嗎？ \end{tabularx} \\ \\ \relax
2:7 & \begin{tabularx}{0.7\textwidth}{X} 毀謗為你們求告時所奉的尊名的，不就是他們嗎？ \end{tabularx} \\ \\
[1ex]
\hline
\hline
\end{longtable}
$^{1}$本主日的經訓是.
《雅各書》二章一至七節.
由陳玉峰弟兄為我們宣讀.
之後有.
金庸業 灰暗團契總幹事.
陳啟興牧師在我們當中靜悟.
他的題目是.
生命見真章系列的不二貌取人.
《雅各書》第二章一至七節.
今日的經訓是記載在新約的《雅各書》.
第322頁.
教會的聖經322頁.
請大家翻開.
或者看程序表.
我聽外讀出.
第二章一節.
我的弟兄們.
你們信奉我們榮耀的主耶穌基督.
便不可按著外貌待人.
若有一個人戴著金戒指.
穿著華美衣服.
進你們的會堂去.
又有一個窮人.
穿著骯髒衣服也進去.
你們就重看那穿華美衣服的人.
說:請坐在這好位上.
又對那窮人說:.
你站在那裡?.
或坐在我腳椅下?.
這豈不是你們偏心待人.
用惡意斷定人嗎?.
我親愛的弟兄們.
請聽.
神豈不是揀選了世上的貧窮人.
叫他們在這信上付足.
並承受他所應許.
及那些愛他之人的國嗎?.
你們反倒羞辱貧窮人.
那付足人豈不是欺壓你們.
拉你們到公堂上去嗎?.

$^{41}$他們不是褻瀆你們所敬奉的尊名嗎?.
敬文獨筆.
願神賜福常讀他們的話語並且遵行的人.
弟兄姊妹早安.
早安.
很感恩陳牧師容許我不用穿西裝.
這個很感恩.
我相信大家不會以貌取港.
不穿西裝沒問題.
順便西裝也不穿.
衣袖也摺起來.
為甚麼呢?.
因為這段經文一點也不難明白.
你聽一次也明白.
很難做.
很難應用.
鼓勵大家對於雅各書.
很多經文都是.
「挾起衣袖去實踐」.
這才是關鍵.
但我們要靠聖靈.
心心禱告.
請主我們求你的靈.
引導我們進入真理.
讓真理挑戰我們.
轉化我們.
求主給我們一顆謙卑柔軟的心靈.
我們願意去聽.
我們願意去改.
我們願意我們的生命更加土地起源.
不再以貌取人.
更加不願意以外表的東西.
去衡量自己.
求主幫助.
奉耶穌基督的聖名.
阿們.
在很多年前.
我有一個經歷.
我想是二,三十年前.
有一晚我下街跑步.

$^{81}$因為我有跑步的習慣.
那時候還住在公屋.
我穿得很難搭的.
我記得是一件背心.
一條短褲.
還穿著一雙白飯衣去跑步.
晚上滿面鬚根.
樣子不是很好看.
跑完步大概十一點多.
回到家.
在樓下等電梯.
我身邊有一個中年婦人和她一起等電梯.
只有我們兩個.
電梯門開了.
我進電梯了.
要回樓.
怎知道她不進.
從她看著我的眼神.
過了二,三十年我還記得.
我猜她看見我的樣子.
我猜她猜我是壞人.
所以就被指責.
隔壁不好.
進電梯不好.
跟著這個人.
被人以貌取人不是很好受.
我到今天還記得.
當年一個很短很小的生命片段.
小姐妹.
香港人很重視外觀.
我們每年花在外觀的消費.
是相當驚人的.
我曾經看過一個市場調查.
他訪問了香港1200位.
15至50歲之間的女性.
每年用於護膚的消費.
平均每人大約一萬元.
好像不是很多.
原來隨著年紀增長.
女士的花費在護膚上.

$^{121}$就會一直有急速的增長.
在30至39歲.
平均是一萬一千元.
40歲以上的組別.
平均每年差不多一萬四千元.
一個相當大的增長.
翻查當年.
做這個調查.
當年的人口普查.
香港15至50歲之間的女性.
有208萬人.
你再算一算.
原來每年花在護膚的這群人.
是200億.
200億.
是相當不小的金額.
但你想想這只是護膚.
還沒算髮型.
瘦身.
修甲.
去南韓整整樣子.
這些還沒算.
還沒算15至50歲以外的女性.
還沒算男性.
所以如果你加起來.
這個金額是天文數字.
人是很重視外表的.
其實憑外表去看人.
以貌取人是一個非常普遍的現象.
一個心態.
我們以外貌也不單單是樣貌.
包括我們穿什麼衣服.
我們可能的學歷是什麼.
你告訴別人你是博士.
可能又比起你還沒畢業.
可能又不同.
你自己的職銜等等.
所以美國神學家.
傅斯德.
Richard Forster.

$^{161}$他有一本很出名的書.
就是《屬靈操練禮讚》.
在這本書的第一句.
他說什麼呢?.
他說膚淺是這個世代的就坐.
只是看外表.
以貌取人.
看表面的.
是我們這個世代的就坐.
聖經如何看待以貌取人呢?.
剛才的經文一開始就說.
不可以按著外貌待人.
我們今天看亞國.
他如何教導我們.
不要以貌取人.
在第二章的第一節.
他說我的弟兄們.
亞國是對信耶穌的人說.
這一番話.
你們信奉我們榮耀的主耶穌基督.
便不可按著外貌待人.
按外貌待人.
希臘文是一個複合字.
是由兩個希臘文組成.
第一個就是prosopon.
這個就是解釋你的樣子.
或者你的面貌.
或者面具的意思.
另一個就是lumbano.
這個字就是解釋接受.
或者拿著.
就是receive.
或者拿著的意思.
所以如果你純粹按字面翻譯.
這個字可以解釋.
拿著面具.
是什麼意思呢?.
就是你看到什麼人.
你戴什麼面具.
就是你看著那個人的衣著.

$^{201}$他的容貌.
他的職業.
他的社會地位.
他的知名度.
他有沒有錢.
等等.
你會用不同的面貌.
或者面具.
去對待他.
在原文這句結構很有趣.
原文的意思就是.
你們如果有對主耶穌基督的信心.
就不可以有以貌取人這件事.
這個"有"字是一個命令式的動詞.
就是說如果你是信耶穌.
就一定不可以以貌取人.
如果你以貌取人.
你就不是信耶穌.
雅各是很斬釘截鐵.
以貌取人是一個違背信仰.
跟你的信仰背道而馳.
兄弟姐妹.
以貌取人不是一些微小的小問題.
在信仰上是大是大非的問題.
信耶穌的人是不會以貌取人.
在《申命記》第十章十七到十八節.
耶和華是萬神之神萬主之主.
是偉大強有力可畏的神.
不看人的情面也不受賄賂.
神的性情是不以貌取人.
祂期望信祂的人也是這樣.
以貌取人是違反神的性情.
是嚴重的.
我希望在這一節能告訴我們.
是一個很嚴重的生命問題.
不是一些小問題.
我們留意《雅各》如何形容耶穌.
在這一節.
「榮耀的主耶穌基督」.
如果你記得.

$^{241}$應該是上幾次你去看《雅各書》.
在第一章第一節.
《雅各》沒有用「榮耀的」來形容主耶穌基督.
為什麼這裡刻意要加上「榮耀的」呢?.
有什麼用意呢?.
我推斷《雅各》用「榮耀的主耶穌基督」.
是想表達耶穌在世的時候.
也被人憑外貌去對待.
這個木匠的兒子.
在馬槽出世.
無階型美容.
無財無勢.
以致別人看不到他的榮耀.
《雅各》似乎要暗示.
我們曾經這樣對待主耶穌基督.
我們以貌取神.
我們不要再這樣對待別人.
也不要再這樣錯下去.
孔子也說.
「以貌取人,失之止於」.
而聖經說的更加嚴重.
「以貌取人,是有違信仰的」.
我們有沒有輕視他的嚴重性.
和他的破壞力呢?.
弟兄姊妹.
希望大家看完這一節.
你應該很小心.
我們有沒有這樣去看待別人呢?.
《雅各》接著舉出一個實例.
去印證「以貌取人,是不可取」.
在原文第二到四節.
其實是一句長句子.
它提到一個衣著華麗的人.
和一個衣著很骯髒的窮人.
進入會堂.
他們受到極之不同的對待.
極之不公平的待遇.
「眾父輕貧」是一個以貌取人的.
很明顯的例子.
也是我們偏心待人的一個例子.

$^{281}$所以在這一段經文.
你偏心待人.
以貌取人.
眾父輕貧.
都是指同一個問題.
我們犯了同一個信仰的問題.
在第二節一開始.
有一個「若」字.
「若」有一個人.
這個「若」字在文法上.
我們叫做第三類條件句.
就是說所描述的東西.
是有可能是真的.
是有可能發生的.
「不可按外貌待人」.
這個詞語的文法結構.
我們叫做一個否定詞.
加一個命令式的動詞.
是表示要停止某一項.
你進行中的一個行動.
所以較精準的譯法是.
「不可再按外貌待人」.
如果我們這樣解釋.
修信人是曾經犯過這個毛病.
因此這個例子.
應該是一個真實的情況.
而且用真實的情況.
去支持雅各的論點.
也是較有說服力.
現在讓我們看看.
這是一個什麼場景.
他說.
有一個衣著華麗.
一個衣著骯髒的人.
進去會堂.
會堂是猶太人最會敬拜的地方.
所以這裡可以指一個崇拜的場景.
但有沒有其他可能性呢?.
有幾個因素值得我們去參考.
第一.

$^{321}$其實這段經文的大段落.
就是二章一到十三節.
這個大段落裡面.
如果你看原文.
有審判官.
公堂.
並違犯法.
受審判.
審判.
等等.
這些法律.
或者法庭的用語.
是重複又重複出現.
暗示這個會堂.
可能是一個宗教法庭的場景.
第二個因素.
坐在這裡.
站在那裡.
坐在我腳下.
這些字句就較適合去形容.
一個法庭的場景.
多於一個崇拜的場景.
一個崇拜.
你來到.
你不會坐在某某腳下.
所以第三個因素就是.
第三節.
我的腳下.
這個可能是指法官的腳下.
一個人坐在法官的腳下.
這個就是將他視為一個敗訴者.
所以更加凸顯.
這是一場很偏私的審訊.
你穿得這麼差.
你一來到.
你坐在那裡.
你一定輸.
不用審.
你一定輸.
按照這些理解.

$^{361}$這裡可能是指一個宗教法庭的場景.
可能更加合理.
我想起我自己的一個經歷.
我的外父以前在內地做一些小生意.
他被人騙.
也被人告.
其實告他那群人是想侵吞他的工廠.
我出錢替我外父打官司.
但是每次都輸.
沒有一次是贏的.
後來知道原來他的對頭是認識法官.
這當然是相當若干年前的事情.
我認為是打不了的.
這些官司是打不了的.
是一些不公平的官司.
於是我勸我外父結束了內地的生意.
面對一些不公平的審訊.
不公義的審訊.
我不知道你最氣憤.
你最生氣的是誰.
是那些以貌取人.
偏心待人的人.
那些不公義的人.
這個有錢就多幫一些.
這個很窮.
這個沒有財沒勢的.
他一定輸.
可能是法官.
可能是其他人.
雅各稱呼這些重負輕貧的人.
是偏心待人.
用惡意去評斷人.
偏心待人這個希臘文.
都是由兩個字組成.
兩部分.
一部分是分割.
一個是審判.
就是說你很割裂地去判斷.
你是很側重一方的.
當然側重的是有錢的那方.

$^{401}$這樣做就是心懷惡意的審判官.
其實用惡意評斷人.
在第四節.
這個字的原文.
它的執法就是心懷惡意的審判官.
以貌取人的人是坐了神的位置.
我自己就是審判官.
而且是一個心懷惡意的審判官.
在美國加州.
很多年前有一對老夫婦.
我不記得我在黃浸有沒有講過這個故事.
如果有的話大家可以再聽一次.
有一天這對夫婦.
穿著一些很便宜很舊的衣著.
去拜會哈佛的校長.
沒有事先約好的.
直接去拜會哈佛的校長.
校長的秘書一看這兩個老人家.
就很肯定你跟哈佛都沒有什麼業務上的往來.
就跟這對老夫婦說.
校長很忙的.
不會見你們的.
那對夫婦就說.
不要緊.
我們有時間.
我們可以等.
過了幾個小時.
秘書沒有理會他們.
希望這對夫婦知難而退.
但是這對夫婦一直在等.
這個秘書也沒有辦法.
最後就通知哈佛的校長.
或者你跟這對老夫婦.
說兩句打發他們離開.
校長很不耐煩.
但是也同意.
希望可以很快地打發這對夫婦離開.
這個校長很有價值.
也很不甘心.
不情願去見這對夫婦.

$^{441}$坐下來.
這個老太太跟校長說.
我們的兒子去年就死了.
但是他很喜歡哈佛.
我們很想在哈佛.
可以在校園裡面.
做一些事情去紀念他.
這個校長不僅沒有被感動.
還覺得很好笑.
他說我們不可以為每一個很喜歡哈佛而死的人.
去建立一個雕像去紀念他.
如果這樣做.
我們整個哈佛校園.
就好像一個墓園一樣.
這個老太太很快地說.
不是的.
我們不是這個意思.
我們不是要建立什麼雕像去紀念兒子.
我們是想捐一棟大樓給哈佛.
這個校長就看著這對老夫婦的衣著.
這麼破舊.
這麼殘舊.
你們知不知道.
要建一棟大樓要多少錢?.
當時他說要超過750萬美元.
當時是說十九世紀末.
要750萬美元.
這個時候那對老夫婦就沒有出聲.
這個校長很高興.
他覺得你們聽完之後就會走了.
這個老太太就轉頭跟她老公說.
750萬美元可以建一棟大樓.
我們為什麼不建一棟大學去紀念兒子呢?.
這個老師就同意.
這個校長覺得很困擾.
為什麼會這樣呢?.
後來這對夫婦.
就是利蘭史丹佛先生和他太太.
就離開了哈佛.
在1891年在加州就建立了.

$^{481}$很著名的史丹佛大學.
這個是流傳故事.
有些是真的.
但肯定有些不是真的.
如果你們上網知道有些不是真的.
但這個故事也很有教育意義.
這位哈佛校長就以貌取人.
很用惡意去評斷一個人.
成為笑柄.
說我們知道不應該以貌取人.
很清楚我們不應該這樣.
按照這個實例我們不應該這樣.
但原因在哪裡呢?.
雅各在第五到七節就告訴我們有兩個原因.
在第五節.
「我親愛的弟兄們請聽.
神豈不是揀選了世上的貧窮人.
叫他們在信上付足.
並承受他所應許給那些愛他之人的國慕.
神揀選了世上的貧窮人.
叫他們在信上付足.
不以貌取人.
我們不重負輕貧.
不偏心待人」的第一個原因.
這就是神的心意.
神的心意是這樣的.
「神揀選了世上的貧窮人」.
就是說從世界的標準看來是貧窮的.
就是物質上的貧窮.
神讓他們在信心上付足.
就是從信仰的角度看來.
他們是付足的.
所以很多窮人.
露宿者.
基層都可以信耶穌.
都可以平安喜樂.
覺得很滿足.
因為上帝揀選了他們.
但是你一想.
神是不是也有些問題呢?.

$^{521}$祂不是重負輕貧.
祂是重貧輕負.
是不是?.
祂是揀選了世上的貧窮人.
如果神是不是又以貌取人呢?.
我們不要誤會.
當然不是.
所以雅各很清楚地說.
能夠承受神的國.
能夠得救的條件.
不是貧窮.
而是信心.
是愛神的人.
「愛」這個字的文法.
是表達一種持續進行的愛.
你唯有信神.
你不斷愛神.
你才能夠承受神的國.
貧窮本身.
是不可以讓你進入神的國.
所以神不是重貧輕負.
更加不是重負輕貧.
雅各在這裡第五節.
是在說你們重負輕貧.
輕視貧窮人.
但神仍然看顧這群貧窮人.
就是說神的心意是不應該重負輕貧.
也不應該以貌取人.
到了第六節.
其實一開始有一個「但」字.
而且將「你們」這個詞語.
放在這一句的開始.
是表示一個強調.
神不會重負輕貧.
但是你們就去羞辱貧窮人.
在禎延的17章5節.
「譏笑窮乏人的是藐視造他的主」.
如果我們去羞辱貧窮人.
其實就是羞辱創造這些貧窮人.
各位上帝.

$^{561}$而且羞辱這個動詞的詞態.
是表示他們實際曾經做過.
這些羞辱貧窮人的行為.
他們曾經侮辱過神.
這一群所謂相信神的人.
不僅如此.
這一群富足人還有三宗罪.
在最後一節.
他們欺壓你們.
將你們拉去公堂.
接受一些不公平的審訊.
又褻瀆你們求告時所奉的尊名.
也就是耶穌基督.
那些富足人的真面目是這麼壞的.
你們還要看重這些富有人家.
你們簡直不將神放在眼內.
雅各的意思就是這樣.
結論就是.
你憑什麼去做你判斷的準則.
如果你是憑外表.
你看看他是不是有財有勢.
看看他是不是有什麼效力.
有什麼社會地位.
那就會更為負興奮.
如果你是看人的真面目.
看人的內心.
你就不會重負興奮.
所以聖經說.
「夜禍說看人」.
就不像人看人.
人是看什麼的?.
外表.
神是看什麼?.
內心.
在撒慕爾基.
所以我們不重負興奮的第一個原因是神的心意.
第二個原因就是你判斷的準則.
應該不是看人的外在貧富.
他的外表.
是看人的內心.

$^{601}$看人的生命.
在泰國的歷史上.
有第一位的女總理.
大家應該都知道是誰.
英祿.
就是他信的妹妹.
英祿有美女總理的稱譽.
她應該是美女.
但是現在因為濫權獨職.
她就被通緝.
浪望海外.
曾幾何時.
海峽兩岸有兩個帥哥.
在政壇有兩個帥哥.
不知道大家有沒有印象.
一個就是當時商務部部長薄熙來.
另一個就是當時的台灣市長馬英九.
稱為政壇商壁.
亞洲周刊曾經找這兩位帥哥做封面.
後來薄熙來貪污受賄.
濫用職權.
現在坐牢.
判無期徒刑.
馬英九本人非常厭惡這種以貌取人的評論.
台灣有一位官員.
每一次當記者稱讚馬英九是帥哥.
他就很氣憤.
去糾正.
這是對馬英九的侮辱.
難道他沒有其他的才華.
可以說你只是說樣貌嗎?.
這個世代有時候真的好像Richard Forster.
膚淺到一個地步.
我們連判斷這些政治人物的準則.
都是看樣子.
在美國曾經做過一個統計.
原來帥哥或帥女候選人.
當選的機會明顯高過像我這樣的人.
是明顯的.
這個是帥一點的.

$^{641}$你的印象分也加一點.
當我們想起我們的信仰.
我們在第一節說.
這位榮耀的主耶穌基督.
是我們所尊崇所信奉.
他最榮耀的時候是怎樣的?.
是頭戴荊棘冠棉.
血流披面.
滿身鞭傷.
衣不閉體.
可能還沒有穿衣服.
掛在十字架上.
說「成了」的時候.
這就是他最榮耀的時候.
所以我們不要以貌取神.
我們所相信的上帝最榮耀.
不是看樣子.
所以如果你問基督徒.
以貌取人對不對?.
十個基督徒.
十個都說不對.
但是我們明知道不對.
我們常常都犯這個毛病.
人靠衣裝.
佛靠金裝.
其他信仰.
所供奉的那些.
是做得很漂亮的.
讓你去供奉.
這是我們常犯的毛病.
在台南有一個真人真事.
有一位姓張的男子.
是很無辜的.
他天生的樣子.
就是一個壞人.
一個賊的樣子.
他無論去超市買東西.
去下班回家.
通常都會被警察截下.
要身分證.

$^{681}$要看的.
他試過一個月.
50次被警察截下.
他破了.
我不知道這是不是見歷史記錄.
他最高紀錄.
就是一天有7次被警察截下身分證.
看看你是不是賊.
其實他就是身家清白的.
他完全沒有前科.
一個良好市民.
他被警察查到破了紀錄.
這位張先生就說.
沒辦法.
因為我的樣子長得很像賊.
張先生就自嘲他可能長得有問題.
但是警方就很鄭重去否認.
我不是看人的樣子去查的.
可能張先生出入的地方比較複雜.
所以我們先去查.
這是警方的辯解.
不過張先生就很無奈地說.
他出入的地方.
一就是公司.
一就是家.
一就是超市.
他說難道這樣也算複雜?.
所以其實我們知道不應該以貌取人.
我們又知道是違背信仰.
不是神的心意.
是用了錯判斷的準則去看人.
但是我們經常都是這樣.
我們真的經常都是這樣.
我們學到什麼?.
有什麼應用呢?.
最後我有三個反省.
這三個反省都很幫助我自己.
第一個就是.
如果我以貌取人.
我也會這樣去看自己.

$^{721}$我只會注重自己的外在,外表.
外貌,學位,職涵,身家.
你擁有多少東西.
你開什麼牌子的車.
你住的屋子有多大.
我忽略了我內在的建立和成長.
我的自我形象.
我只會建立在外在會失去的東西上.
其實很危險的.
如果你這樣看自己.
這樣去建立自己的成長.
很危險的.
所以.
第一.
自我形象不要建立在外在會失去的東西上.
應該建立在與神的關係.
這個不會失去.
永遠都是長存.
上帝永遠與你一起.
所以弟兄姊妹.
不需要花很多時間精力去搞我們的髮型.
買什麼衣服.
我們外表被人看到的東西.
不需要.
我們端莊就可以了.
我們需要建立的是裡面和神的關係.
有寶貝放在瓦器裡面.
我們的寶貝是我們有正確的自我形象.
是因為我們有神.
做我們生命的主宰.
不是我們外表有多少東西.
那些是一定一定會失去.
有一個美軍在越戰打仗毀容.
回到美國.
接受當時很著名的電台主持人.
Barbara Waters 的訪問.
當記者訪問她的時候.
她就拉下面罩.
對著鏡頭告訴別人.
這個就是我的樣子.

$^{761}$一個毀容的樣子.
她說這個就是神的形象.
弟兄姊妹.
我覺得這位美軍很有勇氣.
以及她的價值觀是很值得我們去效法.
我們的自我形象是建立在哪裡呢?.
我很鼓勵大家的應用是.
嘗試減少時間精力金錢.
在建立在那些外在.
你終會失去的東西上.
第二個反省.
我們深信神愛每一個人.
即使是一個你一看就很討厭的人.
可能是你公司的同事,上司.
可能是你的家人,某個親友等等.
如果我們相信神都愛他.
他也是神所很寶貴的人.
你就會少一些以貌取人.
多一些去接納他.
你這樣看他.
他就順眼得多.
我真的試過.
一些人可能是很.
甚至當年我在教會,牧會的時候.
有一些會友.
我很坦白說.
雖然錄影我也要說.
有一些會友真的很討厭.
我們教了一個禮堂.
四邊有走廊.
我看到走廊.
我一暈路走.
我不想看到他.
我不想跟他聊天.
很煩,很討厭.
很尖酸,黑薄,很不屬靈.
但後來我學到一件事.
這也是神所愛的.
這也是神所救的.
耶穌在十字架上也為他流血.

$^{801}$為他捨身.
我有什麼權.
我不愛他呢?.
所以.
你的身邊,你的家人,同學,同事,朋友.
有沒有一些人你眼中很討厭的.
你是否很相信神都愛他呢?.
你越相信.
你就越少以貌取人.
去看待他.
如何幫助你對一些不太可愛的人看得順眼一點.
不要以貌取人呢?.
我的做法就是多一點發掘這個人.
其他方面的優點.
一個人可能有99\%都很討厭你.
不過.
你有沒有努力去發掘他1\%的優點呢?.
一定有的.
沒有一個人是完全的邪惡.
也沒有一個人是完全的正直.
所以他再多不好也好.
他總有一些優點.
你可以去發掘.
我們有沒有努力去做這件事呢?.
還是我們一看到他已經.
以片蓋全.
他整體性是非常之討厭你呢?.
如果我們努力去做.
我們不要針對一些不可愛的地方.
我們會發現原來他也是很可愛的.
有可愛的一面.
我們就不會以貌取他.
對自己.
對別人.
我們不要去以貌取人.
這是上帝非常非常不喜願的.
第三個反省.
可能是最重要的反省.
如何克服以貌取人呢?.
其實說是很容易的.

$^{841}$就是一個金科玉律.
就是愛人如己.
就是這段經文的下一段經文.
第二章的八到十三節.
我相信下個星期會說這一段.
其實這是一個大段落.
一到十三節是一個大段落.
今天點出一個問題.
應對雅各建議的對策.
就是守住這個金科玉律.
這個至尊的律法.
愛人如己.
你做到你就不會以貌取人.
這個我不詳細說.
因為下個星期Raymond應該會跟大家分享.
我很希望這三個反省.
這三個反省幫了我很多年的.
如何看自己.
如何看別人.
努力去學習愛人如己.
如果我們一直去操練.
我們越來越得神的喜悅.
我們越來越能夠得神的稱讚.
我們有一天都會見主面.
每個人都會.
見到主面的時候.
我們想耶穌稱讚我們跟著祂.
忠心良善的一撥人.
還是耶穌說.
我所說的話為何你們不做呢?.
還是姐妹願意我們一起努力?.
這樣我相信.
你真的這麼做.
很慌張.
很溫暖.
是一個很可愛的地方.
不僅熱還很火熱.
如果你做的話.
你的家庭會很溫暖.
如果你這麼做的話.

$^{881}$你的職場會很溫暖.
我們同心吐告.
你自己沒有以貌取人.
我們相當的活躍.
相當的不可愛.
真的很討厭你.
但你仍然去愛我們.
愛到一個地步.
在十字架為我們犧牲.
你的愛在這裡顯明.
主幫助我們.
同樣的以你的眼光去看自己.
建立我們的自我形象.
用你的眼光去看別人.
更加幫助我們.
有你的愛.
你一份捨己犧牲的愛.
去愛身邊的人.
求主祝福旺鎮.
祝福我們的家庭.
祝福我們所身處的職場.
我們的公司學校.
以致我們在.
我們不同的崗位裡面.
能夠活出你的真理.
能夠見證榮耀你自己.
我們同心吐告.
奉耶穌基督的聖名.
阿們.
\newpage



\section{雅各書 2:14-26}
\label{sec:a56U_20AydY}
\textbf{2021年8月1日主餐崇拜直播｜陳冠成牧師｜生命見真章系列 - 顯明你信什麼｜雅各書2\_14-26}
\newline
\newline
連結: \href{https://youtube.com/watch?v=a56U-20AydY}{\texttt{https://youtube.com/watch?v=a56U-20AydY}} ~~~~ 語音日期: 2021-08-01
\newline
\newline
\hyperref[sec:n5CNurE7uu0]{\small{< < < PREV SERMON < < <}}
~
\hyperref[sec:index_chronic]{\small{[返順時目]}}
~
\hyperref[sec:z9XCLH0XmKE]{\small{> > > NEXT SERMON > > >}}
\newline
\newline
雅各書 2:14-26
\newline
\begin{longtable}{cl}
\hline
\hline
章節 & 經文 (和合本修訂版)\\
\hline
2:14 & \begin{tabularx}{0.7\textwidth}{X} 我的弟兄們，若有人說自己有信心，卻沒有行為，有甚麼益處呢？這信心能救他嗎？ \end{tabularx} \\ \\ \relax
2:15 & \begin{tabularx}{0.7\textwidth}{X} 若是弟兄或是姊妹沒有衣服穿，又缺少日用的飲食； \end{tabularx} \\ \\ \relax
2:16 & \begin{tabularx}{0.7\textwidth}{X} 你們中間有人對他們說：「平平安安地去吧！願你們穿得暖，吃得飽」，卻不給他們身體所需要的，這有甚麼益處呢？ \end{tabularx} \\ \\ \relax
2:17 & \begin{tabularx}{0.7\textwidth}{X} 信心也是這樣，若沒有行為是死的。 \end{tabularx} \\ \\ \relax
2:18 & \begin{tabularx}{0.7\textwidth}{X} 但是有人會說：「你有信心，我有行為。」把你沒有行為的信心給我看，我就藉著我的行為把我的信心給你看。 \end{tabularx} \\ \\ \relax
2:19 & \begin{tabularx}{0.7\textwidth}{X} 你信神只有一位，你信得很好；連鬼魔也信，且怕得發抖。 \end{tabularx} \\ \\ \relax
2:20 & \begin{tabularx}{0.7\textwidth}{X} 你這虛浮的人哪，你願意知道沒有行為的信心是沒有用的嗎？ \end{tabularx} \\ \\ \relax
2:21 & \begin{tabularx}{0.7\textwidth}{X} 我們的祖宗亞伯拉罕把他兒子以撒獻在壇上，豈不是因行為得稱義嗎？ \end{tabularx} \\ \\ \relax
2:22 & \begin{tabularx}{0.7\textwidth}{X} 可見信心是與他的行為相輔並行，而且信心是因著行為才得以成全的。 \end{tabularx} \\ \\ \relax
2:23 & \begin{tabularx}{0.7\textwidth}{X} 這正應驗了經上所說：「亞伯拉罕信了神，這就算他為義」；他又得稱為神的朋友。 \end{tabularx} \\ \\ \relax
2:24 & \begin{tabularx}{0.7\textwidth}{X} 這樣看來，人稱義是因著行為，不是單因著信。 \end{tabularx} \\ \\ \relax
2:25 & \begin{tabularx}{0.7\textwidth}{X} 同樣，妓女喇合接待使者，又放他們從另一條路出去，不也是因行為稱義嗎？ \end{tabularx} \\ \\ \relax
2:26 & \begin{tabularx}{0.7\textwidth}{X} 所以，就如身體沒有靈魂是死的，信心沒有行為也是死的。 \end{tabularx} \\ \\
[1ex]
\hline
\hline
\end{longtable}
$^{1}$今日經分係記載響雅各書二章十四至二十六節新約聖經三百二十二頁.
雅各書二章十四至二十六節新約聖經三百二十二頁.
我的弟兄們若有人說自己有信心卻沒有行為有什麼益處呢.
這信心能救他們嗎若是弟兄或是姐妹赤身露體有缺了日用的飲食.
你們中間有人對他們說平平安安的去吧願你們穿得暖吃得飽.
卻不給他們身體所需用的這有什麼益處呢.
這樣信心若沒有行為就是死的必有人說你有信心我有行為.
你將你沒有行為的信心指給我看.
我便藉著我的行為將我的信心指給你看.
你信神只有一位你信的不錯魔鬼也信卻是賤精.
虛浮的人你願意知道沒有行為的信心是死的嗎.
我們的祖宗亞伯拉罕把他兒子以殺獻在壇上.
豈不是因為行為輕易嗎.
可見信心是與他的行為並行.
而且信心因著行為才得成全.
這就應驗經上所說.
亞伯拉罕信神這就稱為他的兒.
他又得稱為神的朋友.
這樣看來人稱已是因著行為不是單因著信.
妓女拉合接待使者又放他們從別的路上出去.
不也是一樣因行為輕易嗎.
身體沒有靈魂是死的.
信心沒有行為也是死的.
請姐妹平安.
剛才所讀的亦是另一段我們都很熟悉的經文.
信心與行為.
你看到雅各在這段聖經裡.
他不再與讀者繞圈轉彎抹角.
他直接與讀者信仰見真章.
來檢視到底我們所說的信心是否真的.
你看到雅各很清楚地說.
真正的信心是一定會產生相應的行動.
我們每天都在運用這個信心.
你今天有沒有帶雨傘?.
有多少人有?.
即是你相信甚麼?.
不帶那些呢?.
打算借人家的那把?.
立即買一把?.
給自己機會買一把?.

$^{41}$很簡單的.
你出門前相信今天一定會下雨.
你就會帶雨傘.
相反,我相信今天全日都陽光普照.
所以我今早就特意洗好衣服.
還要特意晾在窗台.
讓它曬乾.
我放完崇拜回家,剛剛好乾了.
雅各強調信心與行為是分不開的.
但他從來沒有否定信心是很重要的.
他表明有信心.
但他表明有信心不代表甚麼都不用做.
他一再指出.
信心要加上行為才完整.
才得以稱為義.
並且他跟我們說一件很重要的事.
不只說一次.
你看到整段經文已經說了四次.
他說如果只停留在言語上,口頭上.
但沒有付諸行動的信心是怎樣的?.
死的,虛的.
對自己的信仰生命沒有甚麼幫助.
對別人也沒有甚麼益處.
他講得很清楚,你看看那幾段經文.
他舉了一個例子.
假如看到身邊有弟兄姊妹陷入生活困境.
去到一個甚麼地步?.
吃不飽,缺乏,穿不夠,不夠暖.
去到這個地步時.
如果只跟他們說「你平平安安地去吧」.
「願上帝供應你日用所需」.
「我相信上帝會保守你」.
雅各說如果只停留在口頭上的信心.
如果去到這個情況仍然是這樣的話.
對那些正在捱餓.
可能特別是冬天.
不夠衣服穿,在抵寒冷的人來說.
有沒有幫助?.
其實沒有甚麼實際養份.
當然,雅各這裡不是說.

$^{81}$言語上的關心和代求一點價值都沒有.
不是這樣說.
而是提醒我們.
當我們真的誠心誠意祝福對方.
想對方好的時候.
我們祈求上帝供應他所需要的時候.
我們自己會否同時被感動?.
好像聖經經常說.
「動了慈心」.
雅各在問這件事.
我們願不願意為我們所祈求.
我們所相信的事情.
既然我們期望它發生.
我們會否自己付諸實際行動?.
我們會否主動去想一想.
我為了上帝的緣故.
上帝讓我遇到這件事的緣故.
我也為了對方.
怎樣可以多走一步?.
我自己有甚麼資源可以提供?.
又或者我可以善用周邊的資源.
幫對方搬到甚麼物資.
來更具體地同行祝福他?.
來守望扶持他.
過渡他在這一刻.
人生最艱難的一段路程.
跟他走一段?.
雅各在說這件事.
確實我們沒有能力解決.
這個世界所有人的難題.
這是上帝的工作.
亦是他才有這樣的能力.
但我們在寄望上帝出手幫助別人的同時.
我們不可以忘記一件事.
真的不可以忘記.
上帝通常是怎樣出手?.
上帝是怎樣出手幫助你?.
你想想,很多時是透過人.
上帝很多時的出手方法.
就是感動差派不同的人.

$^{121}$透過他們出手來替上帝工作.
這是經歷,對吧?.
上帝會使用不同人點點滴滴的付出擺上.
來成就他的計劃.
不單止幫助到一些有需要的人.
其實亦是建立我們自己.
當你每次付出愛心行動的時候.
其實背後代表甚麼?.
你每次付出,每次為別人擺上.
其實是每次顯明你對上帝的信心.
點顯明法.
其實你真的相信耶穌所說的.
愛倫社,愛人如己的吩咐.
你如果不相信,你不會這樣做,對吧?.
你因為相信愛人如己,愛倫社的吩咐.
這是真理,這是美善.
這是對人,對自己都有益處.
你這樣相信,所以你也這樣實踐.
來顯明你的信心是活的,不是死的.
亞哥透過這個例子提醒我們.
我記得十多年前,那時神學院還未畢業.
收到一筆獎學金,很開心.
說多不多,說少不少.
但也足夠畢業後去旅行和太太.
但我同時收到兩個消息.
其中一個神學院同學在班土會內代土.
他說生計開始結局.
太太突然失業.
又開始帶著肚子.
自己作為丈夫也很無奈.
因為神學院的實習費是有限的.
很拮据,又幫不上忙.
很艱難.
他這樣分享,當然為她代土.
但上帝給我一個感動.
其實兩件事是否偶然發生的呢?.
上帝突然給我一筆獎學金.
是否想讓我自己用呢?可能是.
但也可以回應這位同學眼前的急需.
另一位同學是他的師弟.

$^{161}$剛入神學院第一年.
但教會對他的資助不足夠.
他擔心能否捱過這幾年.
我又有感動,能否用這筆獎學金幫助他.
結果用這筆錢分成兩筆.
轉贈給他們.
在過程中,有幾個體會,學到幾個功課.
首先,上帝才是他們的救主.
但上帝真的會使用人在某個時刻.
與別人同行.
上帝會使用我某個時刻去輔助對方一把.
這是真實的.
另一方面,我這樣做.
固然可以解兩個同學一個燕眉之急.
確實能使他們暫時不需要為學費或生活擔憂.
但這個行動有一個更深入的意義.
其實是兼顧他們的信心.
讓他們體會到上帝真的聽禱告.
上帝如何聽禱告?.
就是透過某個人,不知為何.
會將一些金錢轉贈給他們.
上帝真的看顧他們.
上帝會在適當的時候.
奇妙地賜下足夠的恩典.
我想他們明白這一點.
沒錯,錢真的會有一天花完.
但如果他們對上帝有這份信靠.
這份信心可以長久.
這份信心可以在日後他們遇到困難時.
繼續承載他們.
可以像今次一樣.
隨時仰望上帝幫助和供應他們.
我想這個心意比金錢更大.
事實上這件事不只是幫助對方.
其實是建立我對上帝的信心.
確實那筆錢不是我的.
是上帝的.
我們說我們只是一個管家.
那時候就要體驗什麼是管家.
用行動體驗,你才是管家.

$^{201}$你要學習妥善.
如何和別人分享上帝給你的資源.
你確信你這樣做.
上帝雖然不會說出口.
但祂會起源你.
另一方面,我想整件事.
真的體會到讓我確信.
什麼叫上帝使萬事互相效力.
叫愛上帝的人得到益處.
很多時候是透過你的行動體驗出來.
所以,反過來,我有時會這樣想.
將來如果我在任何事上.
不一定是金錢.
可能我事奉到沒力的時候.
缺乏的時候.
其實我都可以一樣放膽.
隨時向上帝傾訴祈求.
上帝亦可能會像那時候那樣.
透過不同人來鼓勵我,幫助我.
不知道大家最近有沒有看奧運?.
有吧?.
有沒有看舉重?.
舉重比較冷門,大家有看嗎?.
有時不經意按了過去都看了.
我在裡面舉重這個項目.
學了一個成語叫.
叫舉重躍興.
有沒有聽過這個成語?.
舉重躍興.
字面解釋就是.
舉起很重的物件.
但給人感覺很輕,很容易.
你看到,譬如那個是50公斤級的賽別.
那個選手其實只有五十多公斤.
他怎樣舉起一個超過200公斤的那個叫什麼?.
又不是啞鈴,總之那幾塊鐵頂.
一個五十多公斤的人.
怎樣舉起200多公斤的東西?.
最直接的方法就是.
拿走鐵頂就行.

$^{241}$只舉起那支槓.
就立即舉重躍興.
我們很快會想.
怎樣設法將那件物件減重,令它變輕.
這樣就很容易舉起.
但其實我想舉重躍興的意思不是這樣.
它是說物件的重量不變.
但你怎樣調教自己的心態.
你不要經常想它有多重.
反過來想.
你要怎樣將它輕輕的舉起.
舉得多輕,舉得多輕散.
很有哲學的意義,舉重躍興.
你怎樣通過正確的心態,正確的方法.
反覆的操練.
以至你可以從容不迫地.
完成一件很大的事情.
如果從信仰的角度看.
舉重躍興的秘訣就是什麼呢?.
就是不要刻意避開.
逃避上帝真的給了你那些託付.
祂可能將某些人.
他的需要,他的缺乏,真的託付給你.
不要逃避上帝給你的這個重任.
而是先調教自己.
上帝為何給你?.
為何?.
是否上帝想整你?.
是否上帝不愛你?.
是否上帝難為你?.
這些其實全部都是減磅的方法.
你想將重擔拿走.
不是這樣,而是怎樣?.
是因為上帝看得起你.
是因為上帝喜悅你.
是因為上帝知道你靠著祂舉得起.
所以,祂才將這個福事交給你.
其實是這樣的意思.
先調教了你的心境.
心境改變,事奉自然怎樣?.

$^{281}$輕省.
而當我們繼續願意尋求上帝的幫助.
用上帝提供的方法,資源.
不再靠自己的努力亂來.
亂事奉的時候.
你的事奉會變得更加輕省.
因為你用上帝給你的方法和資源.
久而久之,你慢慢會發現.
其實是誰擔起這個擔子?.
是上帝和你.
是上帝默默地在幫助我們.
真正想辦法的是祂.
真正出錢出力的也是祂.
我們只是,像保羅所說.
是器皿.
你自結就行了.
你應上帝的機,上帝就要用你.
祂會幫你.
漸漸你也會開始明白.
原來再重的擔子.
耶穌也說過.
祂也應許過.
你可以靠著祂,而揹得輕省.
所以.
其實雅各提醒我們.
信心和行為好像很沉重.
但其實祂叫我們放心.
祂從來沒有說過.
我們要解決世界裡所有的重擔.
上帝也沒有這樣說過.
但上帝同時也沒有叫我們.
見到別人有缺乏的時候.
你就守旁觀,當看不到.
你留意雅各這裡舉的例子是什麼?.
若是看見弟兄或姐妹的缺乏.
某個的缺乏.
吃不飽,穿不亂.
意思是.
祂沒有說過我們要.
無止境地為全世界付出一切.

$^{321}$你看見某個弟兄或姐妹的缺乏.
其實就是指你身邊的人.
上帝讓你看見.
上帝安排你去接待.
你去回應某個弟兄姊妹.
或某個可接觸的鄰舍親友.
在他們某個需要援手的時刻.
你跟他們分享上帝.
和上帝給我們的資源.
很實際了,今天幾號?.
你收到什麼?.
還沒有收到?真的?還沒有登記?.
用回今天的經文.
你怎樣善用那些消費券?.
你會不會想.
自肥之餘.
亦都服侍身邊有需要,有缺乏的人.
來建立他們對上帝的信靠.
亦都同一時間顯明.
你真的對上帝愛人如己.
這個吩咐的相信和信服.
很實際的一個挑戰.
我們繼續看亞覺怎樣說.
他似乎已經預計到.
當他講信心和行為的時候.
一定會有人反駁他的教導.
不是講好了恩信,清義嗎?.
基督徒得救我們都懂得背.
本符是什麼?.
本符是恩典,亦都恩著.
信嘛,我們都懂得背.
唯獨信心.
總之我知道自己信上帝.
我對上帝很有信心.
這樣就可以了.
你會不會這樣想?.
這樣又不是錯的.
不要覺得罪大惡極.
亞覺對於這樣的回應.
他就打蛇除棍上.

$^{361}$把你口中所講的.
看不到行為的信心.
拿出來看看.
把你口中所講的.
沒有行為的信心.
展示給我看看.
讓我看看是什麼一回事.
然後我亦都透過我的行為.
把我所講的信心.
具體地給你看.
亞覺的潛台詞是什麼?.
對方可不可以拿出來給你看?.
拿不拿到?.
拿不到.
因為沒有人能夠.
把沒有行動的信心拿出來.
去存在這種信心.
相反.
亞覺說他的信心是有根有據的.
可以透過一些行動.
向人展現.
來驗證這是真的信心.
等於展會亞覺進一步提醒我們.
當我們講信心的時候.
我們都要告訴自己.
或者向別人作一個見證.
信心不會隱形的.
信心不是隱形戰機.
你看不到我.
或者用廣東人的術語.
口講無憑.
行動才是最實際.
透過我們一些行為表現.
就可以向自己或別人顯明.
我其實在信什麼.
我到底信得有多深入.
但你留意.
亞覺不是要和保羅唱對台戲.
他不是要挑戰唯獨信心的教導.
他只是要比較兩個信心.

$^{401}$第一個信心是有行動的信心.
另一個就是沒有行動的信心.
從而突顯了人的狀況.
雖然口稱信上帝.
但身體才是最誠實.
他的行為才會反映.
他是在信還是未信上帝.
他又舉了一個例子.
對一個猶太人.
特別是虔誠的猶太人來說.
上帝是獨一真神.
是他們從小就已經教導.
認信的事情.
上帝是獨一真神.
亞覺表示這個認信.
完全正確.
絕對無誤.
但如果只是停留在腦袋上的認信.
其實不一定真是信靠上帝.
因為亞覺說的.
和上帝一直對著幹的魔鬼撒旦.
他也是這樣信的.
他也相信上帝的存在.
他也相信上帝是獨一真神.
他更加深信不疑.
但他這個相信.
沒有為他更親近上帝.
沒有讓他更信服上帝.
相反是更加遠離.
更加懼怕搏逆上帝.
簡單來說.
魔鬼對上帝的相信.
對他一點正面的好處也沒有.
聽姐妹.
假如亞覺同樣問我們這個問題.
我們是否相信上帝是獨一真神?.
你的答案是什麼呢?.
假如我們答「我信」.
我相信亞覺會再問深入一點.
問第二個問題.

$^{441}$你這個認信背後對你有什麼意義?.
為你帶來了什麼實際的改變?.
想一想.
是更加親近上帝.
更加倚靠信服上帝.
還是你仍然和上帝保持距離.
若即若離?.
亞覺說.
我們的行動.
我們的行為.
就是你最誠實的答案.
他在說這件事.
如果我們信了耶穌一段年.
我們就能夠信上帝.
如果我們信了耶穌一段年.
但沒有人看得到.
我們自己也不察覺.
我們有什麼行為改變的話.
而我和上帝的關係.
總是他有他的工作.
我有我的忙碌的話.
我有困難.
我就懂得第一時間向他祈求.
但當他來敲我的門的時候.
我總是跟他說.
我很忙,遲些再說吧.
我遲些回應你.
甚至忘記了回應他.
久而久之,我也忘記了.
開始忘記了.
上帝對我有什麼要求.
有什麼期望.
這樣.
亞覺就會問我們一次.
喂.
雖然我們說自己信耶穌,信上帝.
但有沒有想清楚.
你其實在信上帝什麼.
你在信祂什麼.
會不會只是在信.

$^{481}$上帝會給我一些好處.
上帝會保佑我.
他將來一定會給我.
天堂的入場券.
但是.
我又不想他干預過問.
我現在是怎樣生活.
我怎樣做人.
亞覺說.
這樣我們很難說自己.
真的相信上帝是獨一真神.
為什麼這樣說?弟兄姊妹.
因為當我們說.
我們真心相信上帝是獨一真神.
其實背後的意思是什麼.
我確認我不會跟其他人.
我確認我只是跟上帝.
效忠上帝一個.
不會跟其他事物.
因為唯有上帝說的是真的.
是對的.
其他的都是怎樣.
甚至我有時的想法.
我有時的行徑都是假的.
這就是你宣認上帝是獨一.
背後真正的意義.
只有上帝說的話可以信.
只有他的安排帶領是美善.
是有益.
這就是我們認上帝是獨一.
背後一個真正的意義.
我怎樣證明.
我不是口說.
我是真的相信.
用雅各的說話.
你走出來就是了.
你致力在日常生活裡.
將上帝的教導.
將聖經的話語活出來.
實踐出來.

$^{521}$過程裡可能我們會有軟弱.
可能會有活逆的時候.
可能會衰左倒退.
但你繼續照著上帝的真理.
你認罪悔改.
然後求耶穌的寶血.
潔淨你,赦免你.
然後你再繼續重新上路.
進行上帝的話語.
直至你行完上帝.
給你在世上這條信仰的旅程.
去到終點.
見耶穌那一面為止.
你這樣做的話.
你就真的相信.
上帝是獨一真神.
真的,弟兄姊妹.
如果我們真的相信.
上帝是獨一的.
其實我們真的要用行動.
來跟自己,跟別人證明.
上帝的話語是可靠的.
是有效的.
我們也會透過行動.
來顯明.
凡是委身信服上帝的人.
他真的會望上帝保守.
望上帝賜福.
就好像雅各.
再舉兩個見證.
兩個見證人.
亞伯拉罕和拉合.
你會看到.
亞伯拉罕獻耳撒這件事.
為什麼說和信心有莫大的關係.
是因為.
很簡單.
上帝賜給亞伯拉罕一切福氣和應許.
都要透過誰來承傳和延伸下去.
透過誰?.

$^{561}$耳撒.
所以才讓他和撒拉生了耳撒.
如果沒有了耳撒.
那代表什麼?.
就是上帝先前所說.
一切的祝福.
一切的應許.
就會完全沒有了.
化為烏有.
所以為什麼我們說.
這是一個相當糾結.
相當掙扎的信心考驗.
上帝說我等了多少年?.
二十多年了.
足足二十五年了.
你才真的給我耳撒.
現在你為什麼要我還給你?.
我不明白你到底想我怎樣.
你真的想祝福我.
還是想作弄我?.
但是.
亞伯拉罕厲害的就是.
即使如此.
搞不清整件事.
他都決定什麼?.
用今天的術語就是做死忠.
Die hard fans.
即使如此.
我不明白.
但我都甘願將耳撒獻上.
來顯明對上帝的信靠.
因為我相信.
上帝才是我的主宰.
他有權賞賜.
也有權收回.
他很認定這件事.
第二件事.
他很相信有些安排.
有些吩咐.
我真的不明白上帝.

$^{601}$為什麼要這樣.
搞到我這麼辛苦.
但他也明白.
上帝不一定要向他解釋清楚一切.
但他仍然堅持相信.
上帝有原因.
有些好東西在後面等著我.
他相信上帝確實是愛他.
要賜福給他.
他相信上帝永遠不會騙他.
不會虧待他.
也因為他這份信心.
他做了一個行動.
來決定.
如果人生有些事情.
只能夠二選一的話.
他寧願選擇什麼?.
他寧願選擇上帝自己.
他寧願相信上帝的安排.
也不相信自己的判斷.
如果人生只有二選一的話.
他寧願選擇.
我忠於上帝的吩咐.
而不保留上帝給了我的好處.
他的認信是這樣.
丁子梅·亞伯拉罕.
通過對上帝這種至死忠心的表現.
來確認.
確信.
上帝是永遠正確的.
而這份態度也在挑戰.
我們對上帝的信心.
假如.
上帝真的突然收回.
他先前賞賜給你的東西.
可以是你的財富.
可以是你的健康.
可以是你的家人.
可以是一些關係.
你如何取捨?.

$^{641}$你的反應是怎樣?.
你捨不捨得放手?.
你會不會埋怨上帝.
明明已經要將束縛賜給我.
已經放在口邊.
已經吃了.
現在你要我套出來?.
你要親手拿走.
你之前給了我的福氣.
你會不會因為這樣.
不相信他.
所以死撐不放?.
還是你仍然相信.
上帝的安排.
總是對的.
總是好的.
你要捨棄.
才會得到新的東西.
捨得.
你願不願意相信.
上帝拿走你的東西.
其實是要給你更好的東西.
以至你繼續相信.
他的安排和帶領.
又或者.
當上帝要我們為他.
犧牲擺上的時候.
你願不願意?.
有人說過一個很實際的測試.
我們是一發薪水就奉獻.
還是薪水後才奉獻?.
你是一發薪水就奉獻.
還是薪水後才奉獻?.
你說有什麼不同.
都是那筆錢.
都是那個數字.
有沒有不同?.
有不同.
很不同.
一發薪水就甘心樂意奉獻給上帝.

$^{681}$是表明.
我擁有的不是我賺回來的.
是出於上帝的恩賜和供應.
我確信上帝恩典是夠用.
不會因為我奉獻了而有缺乏.
上帝總有辦法讓我能夠豐盛.
所以我不會月尾有剩的.
才給上帝.
而是我會憑信心.
月頭就奉獻.
來表達我對上帝的尊重和絕對信任.
這正是一個小小的操練.
以致將來你遇到更大的挑戰時.
你可以運用不斷建立的信心.
最後再看另一個例子.
看拉合如何接待探子.
這段記載是在約書亞記.
大家都熟悉的第二章.
拉合知道當時以色列人的上帝很厲害.
厲害到連當時的埃及法老.
或兩國的阿摩利王都被上帝打敗.
而現在以色列人已經臨近耶利哥.
拉合相信上帝已經將耶利哥城交給上帝.
交給以色列人手.
他亦相信上帝一定會毀滅耶利哥城.
自己所身處的城市.
所以他決定冒險.
協助兩個以色列探子回到自己的陣地.
拉合相信和和平平接待上帝的使者.
其實是在表達甚麼?.
就是和和平平地接待上帝本人.
接受上帝的旨意和安排.
其實是這個意思.
我相信上帝亦會和和平平地.
到兵荒馬亂的時候.
他會接待我.
而事實上拉合因著對上帝的信靠.
他所產生的行動.
不只是救了自己.
最後救了那些人?.

$^{721}$救了他全家.
他很大膽地和兩個探子立約.
我現在恩待你們.
亦都指著上帝來向我起誓.
亦都要恩待我的父家.
救活他們.
你看到一個.
你相信他應該學識不多.
出於基層生活的婦女.
他不惜冒我們說今天是一個甚麼風險?.
叛國的風險.
他憑信心加上行為.
將自己全家的前途壓在上帝身上.
亦都最終為他贏取了上帝的稱許.
藉著他的信心和行為.
在聖經上說得清為宜.
弟兄姊妹,我想他的信心.
令我們想起很多年前.
1924年巴黎奧運會.
有一個男子400米的冠軍.
Eric Liddell.
叫做艾克里.
中文譯音的見證.
有否聽過?.
可能你看過《烈火戰車》這部戲.
不是港產片那一套,是外語片.
艾克里是當時英國男子.
100米短跑的.
基本上他一落場就贏了.
因為當時他已經是紀錄保持者.
基本上他是一個奪金的希望.
所以全國上下都看好他的預期.
而比賽的賽程表公佈了之後.
艾克里發現原來他比賽那天.
100米那天.
糟了,是主日,是星期日.
他本身是一個虔誠的基督徒.
他認定持守主日崇拜.
比參與奧運比賽更加重要.
他認定上帝才是他生命的第一位.

$^{761}$他也堅持自己所相信.
而做了一個所有人都覺得是瘋狂的決定.
就是為了崇拜.
而放棄競逐奧運100米金牌的機會.
這個消息傳到英國的時候.
全國震驚.
上至皇室,下至平民百姓都不接受.
甚至有輿論已經認定他是叛國者.
因為他不以國家利益為優先.
事態當然陷入一個僵局.
在那個時候,另一位同樣是英國的運動員.
叫做蔡林.
把自己400米的參賽資格讓給了艾利克.
希望他仍然有機會代表英國參賽.
如果真的拿到獎牌,事情可以慢慢平息.
但是艾利克是100米的選手.
100米的選手擅不擅長跑400米?.
很少,很少這樣的人.
沒辦法,都要頂硬上.
唯有短期接受400米的針對訓練.
然後毅然上去比賽.
很有趣的是.
他自己說,頭200米奮力為上帝跑.
400米的賽事,如果只計算200米.
他真的可以憑時間拿到金牌.
但你知道他全力跑了200米.
後來計算,後面的200米通常是怎樣?.
通常是沒有力量的,因為前面不留力.
但他說,前面的200米,我全力為上帝而跑.
後面的200米,其實是上帝加力讓我像飛車一樣.
結果,他就像電影一樣.
像一個烈火戰車一樣在賽場上飛奔.
原本沒想過可以贏的.
結果上帝很奇妙地讓他贏了那一屆的400米金牌.
最奇妙的是什麼呢?.
他不只贏了400米金牌,他還贏了200米銅牌.
原本想失去的100米金牌.
就由另一個英國選手拿了.
所以,因為這樣,原本反對艾利克的輿論都慢慢消失.
上帝奇不奇妙?.

$^{801}$上帝說,你尊重他,他也看重你.
事實上不止這樣.
當他攀登上運動界高峰的時候.
他毅然回應昔日對上帝的立戰.
他去了中國,他的出生地在那裡做了宣教士.
直到他回到天家那一天.
可能我們覺得,艾利克的故事太匪夷所思.
哪會發生在我身上?.
竟然為了一番崇拜而放棄世人認為難得的機會.
但是他這份信心和行動的堅持.
真的讓上帝得到榮耀.
他的故事也都剛才說過了.
之後被拍成電影.
到今天仍然成為很多人的見證.
《天之門》的《雅各》提醒我們.
我們要有一個真正的信心.
《天之門》的《雅各》提醒我們.
人生本來就是一場賽跑的考驗.
一個信心考驗.
沿途上,上帝會設不同的挑戰困難給你.
目的是要試煉你和我對上帝的信心.
是否真的堅持將你人生壓在上帝身上.
壓在他的吩咐上.
使你的生命真的可以像聖經所應許.
邁向完全和豐盛.
並且揚得上上帝的稱讚.
弟兄姊妹,你相不相信上帝這份心意?.
你相不相信?.
相信的話,你在未來的日子.
又會如何用實際行動.
來顯明你今天所相信的呢?.
讓我們同心同告.
祈求上帝幫助我們.
讓我們不單單只口中認信.
你是那位獨一的真神.
讓我們在生活中如實反映.
你才是唯一.
你才是絕對.
其他的事物,包括我們自己.
都是次要.

$^{841}$以致我們真的用行動來顯明.
我們願意遵從你的話語.
願意實踐你在我們身上的使命.
來回應我們所信的是真真實實.
我們所信的是有根有機.
正如今日經文所說.
上帝幫助我們.
讓我們更放膽去接待身邊的人.
讓我們不要看為是一個重擔.
讓我們看為上帝.
讓我們一次愛你,服侍你的機會.
讓我們知道上帝是你和我們一起去承載.
並且其實是你去承載.
讓我們不要怕.
讓我們真的能夠放膽.
將我們全人押注在你的身上.
敢於向你擺上.
敢於將我們的人生前途擺上.
特別在這個動盪的世代.
讓我們真的憑信心投靠你.
讓我們憑信心緊緊跟隨你.
直到完成在地上的旅程.
見你面的日子.
我們這樣的禱告交託.
是奉基督耶穌你得勝的名字來祈求.
Amen.
\newpage



\section{雅各書 3:1-18}
\label{sec:z9XCLH0XmKE}
\textbf{2021年8月8日崇拜直播｜陳明泉牧師｜生命見真章系列 - 鑄人以言｜雅各書3\_1-18}
\newline
\newline
連結: \href{https://youtube.com/watch?v=z9XCLH0XmKE}{\texttt{https://youtube.com/watch?v=z9XCLH0XmKE}} ~~~~ 語音日期: 2021-08-08
\newline
\newline
\hyperref[sec:a56U_20AydY]{\small{< < < PREV SERMON < < <}}
~
\hyperref[sec:index_chronic]{\small{[返順時目]}}
~
\hyperref[sec:OhUT6e4ffZM]{\small{> > > NEXT SERMON > > >}}
\newline
\newline
雅各書 3:1-18
\newline
\begin{longtable}{cl}
\hline
\hline
章節 & 經文 (和合本修訂版)\\
\hline
3:1 & \begin{tabularx}{0.7\textwidth}{X} 我的弟兄們，不要許多人做教師，因為你們知道，我們做教師的要接受更嚴厲的審判。 \end{tabularx} \\ \\ \relax
3:2 & \begin{tabularx}{0.7\textwidth}{X} 原來我們在許多事上都有過失；若有人在言語上沒有過失，他就是完全的人，也能勒住自己的全身。 \end{tabularx} \\ \\ \relax
3:3 & \begin{tabularx}{0.7\textwidth}{X} 我們若把嚼環放在馬嘴裡使牠們馴服，就能控制牠們的全身。 \end{tabularx} \\ \\ \relax
3:4 & \begin{tabularx}{0.7\textwidth}{X} 再看船隻，雖然甚大，又被強風猛吹，只用小小的舵就隨著掌舵的意思轉動。 \end{tabularx} \\ \\ \relax
3:5 & \begin{tabularx}{0.7\textwidth}{X} 同樣，舌頭是小肢體，卻能說大話。看哪，最小的火能點燃最大的樹林。 \end{tabularx} \\ \\ \relax
3:6 & \begin{tabularx}{0.7\textwidth}{X} 舌頭就是火。在我們百體中，舌頭是個不義的世界，能玷污全身，也能燒燬生命的輪子，而且是被地獄的火點燃的。 \end{tabularx} \\ \\ \relax
3:7 & \begin{tabularx}{0.7\textwidth}{X} 各類的走獸、飛禽、爬蟲、水族，本來都可以制伏，也已經被人制伏了； \end{tabularx} \\ \\ \relax
3:8 & \begin{tabularx}{0.7\textwidth}{X} 惟獨舌頭沒有人能制伏，是永不靜止的邪惡，充滿了害死人的毒氣。 \end{tabularx} \\ \\ \relax
3:9 & \begin{tabularx}{0.7\textwidth}{X} 我們用舌頭頌讚我們的主—我們的天父，又用舌頭詛咒照著神形像被造的人。 \end{tabularx} \\ \\ \relax
3:10 & \begin{tabularx}{0.7\textwidth}{X} 頌讚和詛咒從同一個口出來。我的弟兄們，這是不應該的。 \end{tabularx} \\ \\ \relax
3:11 & \begin{tabularx}{0.7\textwidth}{X} 泉源能從一個出口發出甜苦兩樣的水嗎？ \end{tabularx} \\ \\ \relax
3:12 & \begin{tabularx}{0.7\textwidth}{X} 我的弟兄們，無花果樹能生橄欖嗎？葡萄樹能結無花果嗎？鹹水也不能流出甜水來。 \end{tabularx} \\ \\ \relax
3:13 & \begin{tabularx}{0.7\textwidth}{X} 你們中間誰是有智慧有見識的呢？他就當在智慧的溫柔上顯出他的善行來。 \end{tabularx} \\ \\ \relax
3:14 & \begin{tabularx}{0.7\textwidth}{X} 你們心裡若懷著惡毒的嫉妒和自私，就不可自誇，不可說謊話抵擋真理。 \end{tabularx} \\ \\ \relax
3:15 & \begin{tabularx}{0.7\textwidth}{X} 這樣的智慧不是從上頭下來的，而是屬地上的，屬情慾的，屬鬼魔的。 \end{tabularx} \\ \\ \relax
3:16 & \begin{tabularx}{0.7\textwidth}{X} 在何處有嫉妒、自私，在何處就有動亂和各樣的壞事。 \end{tabularx} \\ \\ \relax
3:17 & \begin{tabularx}{0.7\textwidth}{X} 惟獨從上頭來的智慧，先是清潔，後是和平、溫良、柔順，滿有憐憫和美善的果子，沒有偏私，沒有虛偽。 \end{tabularx} \\ \\ \relax
3:18 & \begin{tabularx}{0.7\textwidth}{X} 正義的果實是為促進和平的人用和平栽種出來的。 \end{tabularx} \\ \\
[1ex]
\hline
\hline
\end{longtable}
$^{1}$今日經訓是新約的《雅各書》三章一到十八節.
大家可以打開新約的三百二十三頁.
請聽我讀出.
我的弟兄們,不要多人作師父,因為曉得我們要受更重的判斷,原來我們在許多事上都有過失.若有人在話語上沒有過失,他就是完全人,也難立住自己的全身..
我們若把雀環放在馬嘴裡,叫它馴服,就能調動它的全身..
看!旋脊雖然甚大,又被大風吹逼,只用少少的陀,就隨著陀的意思轉動.這樣,舌頭在白體裡也是最少的,卻能說謊..
看!最少的火能點著最大的樹林,舌頭就是火.在我們白體中,舌頭是個最惡的世界,能污穢全身,也能把生命的輪子點起來,並且是從地獄裡點著的..
看!國內的走獸,飛禽,昆蟲,水族本來都可以制服,也已經被人制服了..
唯獨舌頭沒有人能制服,是不知識的惡物,滿了害死人的毒氣.我們用舌頭頌讚那為主,為父的,又用舌頭就助那照著神形象被造的人,頌讚和就助從一個口出來..
我的弟兄們,這是不應當的.全員從一個眼裡能發出甜苦兩樣的水嗎?.
我的弟兄們,無花果樹能生橄欖嗎?葡萄樹能結無花果嗎?鹹水裡也不能發出甜水來..
你們中間誰是有智慧有見識的呢?他就當在智慧的溫柔上顯出他的善行來..
你們心裡若懷著苦毒的嫉妒和紛爭,就不可自誇,也不可說謊話抵擋真理..
這樣的智慧不是從上頭來的,乃是屬地的,屬情慾的,屬鬼魔的,在何處有嫉妒,紛爭,就在何處有紛亂和各樣的壞事..
唯獨從上頭來的智慧,先是清潔,後是和平,溫柔柔順,滿有憐憫,多結善果,沒有偏見,沒有假冒,並且使人和平的是用和平所栽種的異果..
弟兄姊妹平安!.
在生活中我們可以特別反省多一點,實踐聖經給我們的教導..
講一個很特別的實驗,在2018年,有一間跨國的傢俬公司做了一個很特別的實驗,.
在一個地方,所有客觀環境一致,放了兩個盆栽,兩個盆栽都是用膠箱進去,裡面播放一些聲音,.
右邊的盆栽播放著不停的咒罵人奴的聲音,在盆栽裡面播放,左邊的盆栽播放著一些讚美,善言的聲音,.
困住了,類似港台那麼寬,所以客觀環境大致不變,這個實驗做了30天,兩個盆栽最初的情況是一樣的,.
30天之後,播放著很多罵人或者很多噪音,人裡面很多咒罵聲音的盆栽,.
竟然枯萎了,葉子開始枯黃,甚至有凋謝的跡象,很多讚美聲音的盆栽,慢慢的,一路茁壯地成長,.
整個實驗放上網上,其實這個實驗是想講一個校園欺凌,反霸凌事件,人的說話對植物都很有影響,何況對人呢?.
很多人不相信這個實驗,但是真的檢查過,每天來看,那些公司的職員都每天淋水,在這樣的狀態下,聽壞話的盆栽都凋謝,.
聽好話多的盆栽就有一個很好的環境來成長,這個不是屬於最嚴格的,在實驗室裡面的一個科學的實證,.
不過在這件事裡面,似乎給我們帶出一個新的看法,原來人的說話,聲音是帶著一些對生命裡面有改變的,可能你會覺得那盆盆栽都不知道是不是真的,.
不過對人的生命來說,怎樣去栽種一個人,怎樣去打造一個人,我怎樣接收不同的聲音,對我們的生命有很大的幫助,.
如果在座的老師知道,在學生學校裡面整天說髒話,見一見家長,大概見到家長都知道發生了什麼事,.
我們把什麼聲音放進我們的耳朵,決定了我們整個人生,繼續面對人生的路怎樣走,都是朝著那些聲音來塑造我們,.
今天我選取的這段經文亞國書第三章一至十八節,正正就是亞國要提醒我們今天口裡所說的話,某程度上他比較嚴厲或者警惕性的說話多一些,.
他講到今天的舌頭,我們說的話怎樣對人有影響,我們怎樣去運用我們的舌頭,.
真的,一句好的說話,或者人很多時候因為很正面的說話,可能從一個失敗的境地可以振作起來,.
但也有一些很成功的人士因為講錯了一句話,令自己身敗名裂,在聖經裡面,亞國書也好,在箴言也好,甚至在耶穌基督的教導也好,.
他要示範究竟怎樣運用我們的語言,來栽種建立一個人,而不是一個拆毀我們生命的對象,.
希望今天我們讀這段經文的時候,給我們多一些反省,也不要覺得很難實踐,我們就由我們生活細節,一點一步一步開始,希望這段經文成為我們的幫助,.
在我們知曬看聖經之前,我們一起祈禱,求上主幫助我們,同心祈禱,真愛主我們祈求你在我們當中幫助我們,.
讓我們能夠在你的話語裡面,我們知道怎樣行事為人,藉著亞國對我們一番的勸勉,更新我們,.
特別是今天我們要檢視我們平時所說的每一句話,求你的話語更新調教我們,讓我們不會走在你不起眼的道路當中,.
求你恩代頌弟兄姊妹,讓弟兄姊妹聽的時候不單聽得明白,也能夠有力量在生活實踐,幫助宣講的,能夠清晰地將你的話語,.

$^{41}$來講述宣講,讓弟兄姊妹聽得明白,恭敬將來的時間交托,願耶穌基督你帶領,獻上祈禱奉基督耶穌得勝的聖名祈求,Amen..
我們一起看,亞國書第三章一至五節,我的弟兄們不要多人作師傅,因為曉得我們要受更重的審判,.
原來我們在許多事上都有過失,若果有人在話語上沒有過失,他就是完全人,也能勒住自己的全身,.
我們若把爪環放在馬嘴裡叫他馴服,就能調動他全身,看,船隻雖然甚大,又被大風吹逼,只用少少的舵,.
就隨著掌舵的意思轉動,這樣洩透在白體裡也是最少的,卻能說大話..
在這幾節的經文裡面,首先他用師傅這個即時,作為一個起始點來講,為什麼會用師傅呢?.
不難理解的,因為當時的師傅,或是今天的師傅也好,最主要的即時是一個嘴唇的即時,直著我的嘴巴,直著我說的話,.
可能是一些很微細的,一句很簡單的話,但卻有一個很重要,很巨大的效果..
在這裡,雅各特別要講到,對於師傅或老師這個即時,你們不要輕忽,因為這個即時,雖然是簡單的幾句話,.
但卻會受到上帝照著我們所說的話,來有更重的審判..
在教會裡面,做教師的,做團長的,導師的,這個即時是不簡單的即時..
不是說話很有趣,很懂得搞氣氛,那個就是可以做導師,他要講的就是來自聖經的真理,.
按著上帝的心意來做教導,這個不是用強權或是用手臂瓜來左右你人生的路,只是用幾句話,.
或者每個星期以魚目染,用文化或是用他所說的內容來感染大家,令大家走在正途當中..
這個在今天的後現代社會,今天講的這種影響力,這種教師,或者是師傅,輔導員,.
甚至乎坊間所說的KOL,意見領袖,其實都是類似這樣的角色,.
他只講幾句話,但能夠左右人不同的方向和人生做的不同決定..
在後現代或者近代,有一些研究,他將說話成為一種soft power,柔性的權力,.
這種權力不是在制度上或是在人的威權下實踐,純粹是透過說話..
有一個作家叫做Joseph Nye,他特別說這個柔性權力是在國際政治中使用的,.
伊拉克戰爭,他用這樣的戰爭背景,幾句說話可以將整件事風回路轉地扭轉過來,.
他不需要用一兵一卒,但卻是藉著幾句說話將人生命改變..
這個概念柔性的權力後來被應用到在所有的助人行業當中,.
特別是在座的老師,可能是公司的上司,家中的母親或父親,.
每天都會說話,有時行使權力的方式不是用制度或強權,.
有時只講一句,或每天都講那句,這就成為了一種塑造人的權力..
如果我們小心謹慎使用的時候,可以幫助聆聽的人或聽到的人走在一條對的道路上,.
相反來說,如果這個柔性的權力或話語的權柄你不用了,.
你就會將人推向墮落的邊緣,甚至會令到一個人走向滅亡的生命..
雅各特別的去叮囑在當時教會裡面的一班師傅,做教導的那些,.
你不要想著簡單地說幾句說話,你就代表著神的代言人..
師傅這個字其實是一個很高級別的字,大概是在耶穌傳道的時候,.
那些門徒稱呼耶穌就是師傅這個字,所以和合本有時譯夫子,有時譯老師,.
其實這個字就是在說一個身份很崇高,帶領一個人走在一個對的道路,.
跟隨那個師傅做些什麼,他就示範他的生命是一個怎樣的生命,.
那些門徒和學生就跟著他走.雅各在這裡面說,.
既然師傅這個身份是一個這麼有權力,這麼尊貴的身份,.
在教會裡面就不要輕易讓人成為師傅這個身份..
這樣說這件事,好像跟我們這個世代說的很勵志的事有些差別,.
我們今天鼓勵多些人做保羅,帶多些提摩太,我們坊間有很多這些說法,.

$^{81}$甚至有更多的導師上來,不過當你用心去看聖經的時候,.
你會發現雅各對於師傅,對於生命師傅,或者用華語即是的人,.
他不是放得這麼寬,或者不是要令到教會每個人都成為老師,.
那間教會就變得很輕和,反而是他要叫信徒自己謹慎,.
你知不知道你今天所說的話,你要接受雙倍的,或者更大的審判,.
所以這個即是是一個很嚴謹的即是,我記得我們教會的執事,.
羅志強執事阿叔,他以前經常都說一句話,我們一起談,.
他說導師的身份很高尚的,陳牧師,剛剛我入職做青少年傳道的時候,.
因為剛剛要重新建立那個團隊,沒什麼人,.
然後我就想說會不會找這個幫忙做導師,阿叔就說,.
你謹慎一點,陳牧師,那時候是陳先生,謹慎一點,師傅或者生命導師這個職事是很崇高的,.
不要輕易將這個權柄交給一些生命未配合到的人,.
如果你輕易將這個職事交給不成熟的人,對那班年輕人,對你身邊的牧羊,.
不會帶來卑憶的,我記住他這一句話,對於我初出茅廬來說,.
我以前經常覺得多些人服侍,多些人在話語裡面幫幫忙就好了,.
不過經過這十幾年的驗證,我也看著我們的執事他們自己怎樣的榜樣,.
我知道這句話正正就是印證亞國書所說的這番話,.
在教會裡面不要輕易去想或者輕易地讓人站在一個話語職事裡面做師傅這個身份,.
除了不要輕易地在教會裡面去做,好像很負面之外,其實他正面的來看,.
他應該要找一些真是忠心良善,有教導恩賜的人,.
讓他擔當在教會裡面可以教導的責任,如果他能夠好好地教導,對教會是有幫助,.
而不是很濫用盡所有不同的方法,讓所有人都成為導師,這樣就令到教會有一個更新改變,.
這個是第一個很重要讓我們知道的,第二樣東西就是,.
Eugene Peterson在譯這句的時候也是很有意思的,他的英文譯本就是Don't be in any rush to become a teacher,.
不要急於或者倉促地做別人的師傅,除了教會這個處境不要輕易地讓人成為老師之外,.
對於個別信徒來說,你也不要急於自己很快很想成為別人的意見領袖,.
或者用你的話語成為別人的師傅,這是對我們個別人士,對我們個別信仰裡面,.
我們看我們的說話變得很謹慎的角度,來詮釋這一句聖經,.
事實上我們今天很想用我們的話語,憑著我們三寸不爛之舌來影響別人,.
沒有人喜歡你說的話,沒有人聽,我們也希望我們所說的話,.
我們所作的一些建議,或者我們為別人好,給出的一些方法,我們也希望別人聽,.
不過如果我們很倉促,或者手捏牙奈幾個都未驗證到的一些說話,.
就成為我們想影響別人的心態,在雅各裡面叫我們要謹慎,.
我們口說的話,隨時會成為上帝對我們的審判,.
今天很多年輕人有機會接觸一些中學生,他們的志願是做什麼呢?.
以前是做飛機師,或者不同,今天的志願是做網上的KOL,.
他們會想著在網上說幾句話,影響別人就等於那些意見領袖,.
每個人都想做KOL,不過通常聽到這些東西,除了笑笑口之外,.
我都會挑戰一下他,你憑有什麼努力令別人會聽你說話呢?.
是不是將你的說話放上YouTube,別人就會自動聽呢?.

$^{121}$為什麼全世界這麼多人,點擊率這麼高,那些人會聽的,是基於什麼原因呢?.
如果你說的話是正路的,怎樣吸引別人去追求,而且令到那個人會跟隨呢?.
這些就是坊間裡面的KOL作出的努力,這些努力很低調,你看不到的,.
不過如果你純粹這樣想,想著將你的影片放上YouTube,.
你一定會被人看的,事實上有時候我都害怕,我們祈禱會放上去,.
有沒有人看呢?有些弟兄姊妹說,我都有看的,給我一點安慰,.
不過你不要想著放上YouTube,放上一個平台,就會自動蜂擁而至,.
背後有很多的努力,甚至乎有時候你說的話,你用盡心機預備也好,.
未必有人覺得是好的,在我們的話語裡面,亞國提醒我們,要承受更重的審判,.
每一句話,我們所說的都要謹慎,免得我們落入上帝的憤怒當中,.
亞國會繼續的說,原來我們很多事都有過失的,特別過失的是什麼呢?.
在話語裡面,如果我們說話,如果在話語上沒有過失,這個人就是完全的人了,.
在亞國的觀點裡面,其實很多人的錯失,是在話語裡面說錯話,.
說錯話不單只是說一時口快,令到別人感受不好,有時候可能說錯一句話,.
會摧毀一個人,一句說話,令到自己過往努力,可能付諸一旦,.
我記得有一個老師,有一個負面一點的例子,在上課的時候,.
口快,這樣也不會的,這麼垃圾的,這麼簡單說一句話,我們今天聽下去也很不滿意,.
但對那個學生來說,差點面臨著被說的學生,差點面臨著要自殺的風險,.
原來人性是一個很脆弱的人性,同時間我們也要留意,面對著身邊很脆弱的人,.
我們更加要謹慎我們自己說的每一句話,如果我們不是能夠造就人,.
我們的說話盡量不說,如果我們在話語裡面減少缺失,.
我們就走向一個完全成熟的,慢慢成為一個完全人,在這裡說到靈修的角度,.
走向一個靈修完美的,完全的,上帝喜悅的人的生命典範,.
雅各再繼續用比喻,用了連續幾個比喻來說,你看看舌頭在你人生裡面是一個小小的器官,.
就好像一個馬嘴裡面的爪環,一隻這麼巨大的馬,怎樣令它轉方向和奔跑?.
一個小小的爪環就能夠控制它全身,一隻很大的巨大的船,.
怎樣令這隻船可以面對著狂風巨浪,面對狂風巨浪它都會搖搖擺擺,.
但怎樣去逃避這些狂風巨浪?只要一塊小小的舵,將它改變了方向,.
整隻大船就能夠離開一個風浪的位置,或者令到它,這麼大的船就可以駛離開一個危險之地,.
這個小小的舵只是被掌舵的人轉動,同樣你的舌頭都是一樣,.
你的舌頭是一個小小的器官,但你說的就有日積月累的效果,.
每天不斷地說的話,不單影響身邊的人,你也在影響你自己,.
一年365天,每天都說那十句話,一年就3650句,那句負面的話,.
我們人生就是朝著某一個方向,被我們的話語塑造,.
在近代心理學裡面,有一個字眼叫做語境,Verbal Context,.
這個Verbal Context就是人活在一個什麼樣的說話背景,.
有些說話可能是別人說給你聽的,但更多的說話是自己說給自己聽的,.
有時候可能我們面對某一些處境,我們不停地說某一些負面的判斷,.
或者某一些我們可能失去盼望的一些話語,和聖經的教導是相違背,.
我們久而久之,即使我讀多少聖經也好,我也及不上我們自己說的這些說話的影響的塑造力,.

$^{161}$到最後我們被我們自己口中的負面話來調教,甚至一路一路偏離聖經,.
這個就是Verbal Context語境裡面帶來的果效,今天我們不可以小看我們說話的威力,.
我們反而要不停地檢視我們的生命,可以怎樣重新打造我們的舌頭,.
慢慢地將我們成為一個完全的人,有句說話是這樣說的,言語發自心聲,.
詞令寄於學問,知不知道是從哪裡出來的?是唐狄生作的一句歌詞,.
這句不是聖經,是帝女花說的,不過這句跟聖經裡面接下來的兩點是很有關係的,.
你的言語發自你的內心,你的心在想什麼,跟你口中說的話是相符的,.
第二,你的詞令寄於學問,你有沒有智慧,你的詞令反映了你的學問,你的知識,.
我們從這兩點繼續看接下來的兩段經文,三章第五至十二節,.
(文)看!最小的火能夠點著最大的樹林,舌頭就是火,在我們百體當中,舌頭是個罪惡的世界,.
能污穢全生,也能把生命的輪子點起來,並且從地獄裡點著的各類的走獸,飛禽,昆蟲,水族本來都可以制服,.
也已經被人制服了,唯獨舌頭沒有人能制服,是不知識的惡物,滿了害死人的毒氣..
我們用舌頭頌讚那為主為父的,又用舌頭咒詛那朝著神形象被造的人,.
頌讚和咒詛從一個口裡出來,我的弟兄們,這是不應當的,.
全源從一個眼裡能發出甜苦兩樣的水嗎?我的弟兄們,無花果樹能生橄欖嗎?.
葡萄樹能結無花果嗎?鹹水裡也不能發出甜水來..
這個小段,他講到我們剛才說的一個小小的舌頭,帶來的果效,帶來的後果是很大的,.
星星之火可以燎原,舌頭就好像一點火把,一個火源一樣,小小的火而已,.
但是可以把整個森林燒毀,也可以把生命的輪子點起來,並且從地獄點著..
他這裡所講的就是舌頭帶來的那種負面,你看到全部都是用負面例子,舌頭講的那句惡話,.
令人充滿怒火,生命的輪子如何點起來呢?這是當代的典故,可能今天也有些失傳,.
不過我們看回後面從地獄裡面點著的,我們看到大概是可能令到生命好像火輪一樣,衝向一個滅亡,.
就是用一個這樣的意象,但是這個意象究竟有多少成真呢?這個其實已經失傳了,.
不過從他所講的我們知道,就是這個星星之火,小小的火把,.
可以把一種擴散的果後,講的那句小小的東西,卻會樹成一個大錯..
另一個字眼要特別留意的,多次出現的就是那個制服這個字,.
他說大自然世界,昆蟲走獸,所有東西,人類都可以制服,但是制服不了人口中這個小小的器官,.
就是你的舌頭,沒有人能夠制服它,這個舌頭可以看到,像雅各裡面說的,是一個惡物,害死人的,.
對雅各來說,舌頭是一個要攻克己身,叫身服我一個很重要的粗煉的一個身體,很重要的器官,.
如果今天我們很多弟兄姊妹要斷食減肥,要令自己的肚腩消滅,變回那六球腹肌,.
如果我們有這樣的意志,我們更加要更加強力的意志,來訓練我們的舌頭,.
成為一個更加容易沉淪的一個身體的器官,因為它帶來的果後實在太大,.
對於我們來說,我們很多時候不為意的,我們會覺得有意無意之間,我們會覺得自己說話都幾好,.
我們幾有口才,我們不會覺得自己很口臭,我不知道你同不同意,我們經常都覺得自己說話很謹慎,.
不過原來我們想像中的自己和真實的自己是很不同的,有一次偶爾聽一些錄音,.
手機很多時候會自己錄音,不知道為什麼,有時候Google真的會無端端錄音,.
錄音後會聽到自己說的話,說的東西都幾討厭,我經常覺得自己悔改了,.
但原來不是完全悔改,當你聽自己錄音的時候,你會發覺自己說的話不是你想像中那麼得體,那麼典雅,.
我相信天上的阿伯父對我地上的牧師所說的每句話都記錄在案,可能我得失過很多人,.
可能讓很多弟兄姊妹處於不能夠目仰他的,因為我的洩頭,而我拆毀了她的生命,.

$^{201}$不要高估自己所說的話,反而是如實地檢視自己,有些方法很值得我們用的,.
有時候可以憑心情去問問你的家人,問問你家人,其實我說話是不是都令你很辛苦呢?.
你的家人經常跟你相處,他們會聽到你最真實說話的一面,.
你問教會弟兄姊妹,他們都會客客氣氣的,通常都是負面的事正面說,.
不過在家裡的家人最真實的一面,他們可以告訴你,甚至我們更加要刻意問問我們的長輩,.
我們信了主,我們對自己的洩頭有沒有塑造的果效,.
或是我們在理念上信了主,但我們的行為,特別是我們說話,我們依舊顧我,.
所以這裡第一點要提醒我們,我們在日常生活的說話,問問你的家人,問問你的長輩,.
問問跟你經歷過很近距離接觸的,可能你身邊一些很好的同事,問問他們,.
你說話是一個怎樣的人,如果你有空間去接納別人,.
你講實話,其實是反映了你的屬靈光景,你的心在想什麼,你的口就會說什麼,.
耶穌基督在馬福音第七章十五節,同樣跟門徒說,在外面進入的不能夠玷污人,.
唯有從裡面出來的才能玷污人,你的心裡面的惡念才會污穢人,.
你吃了一些污糟的東西,你自然會肚瀉,不會污穢到你,.
但你的說話,你的心裡面帶著這些不好的惡念,你會污穢身邊的人,.
同樣的經文,一銀兩面,在二賽亞書裡面,當二賽亞先知見到耶和華上帝,.
他見到上帝的榮光,第一件事,他反省的是什麼呢?.
他說,禍災我滅亡了,因為我是個嘴唇不潔的人,.
住在嘴唇不潔的民眾,又因我親眼看見大君王萬君之耶和華,.
見到上帝讀聖經的時候,即時要反對自己,我們的口是一把怎樣的口?.
在亞國當中,你說多少頌讚的話,說多少Hallelujah,讚美主,.
但轉過頭來說很多得罪,咒詛,謾罵,是非,說這麼多的東西,.
對亞國來說,那些讚美是沒有意思的,因為那些讚美只不過是熟齡套語,.
可以是技術性字眼,不信主的人也可以說到,.
不過,弟兄姊妹,當我們認真讀聖經的時候,.
我們有沒有想過我們內心的話語充滿了什麼呢?.
我們的內心有什麼惡念,就會在嘴唇裡面表達出來,.
不能賴其他人,每個人都這樣說,我繼續說,.
如果我們深信上帝對我們的生命有塑造,如果我們催向上帝,.
我們不會輕易被這些外面世界的話語繼續容讓塑造我們自己的生命,.
求主幫助我們,我們的話語反映了我們自己的靈性,.
反映了我們的屬靈光景,更加重要的是,那些神聖的話說得多,.
未必會令到你變得神聖,真相是,如果說一些聖經的話語,.
要容讓它檢視我們自己,看到我們自己的話,看到我們自己的內心,.
那些話語才有塑造的功能,那些讚美才是真誠的讚美,.
雅各說,一個泉源又出甜水又出苦水嗎?不可能的,.
同一個嘴巴,怎麼可以兩樣東西都一起有呢?.
基本上,苦就可以蓋過甜,將壞的東西蓋過好的東西,.
你說多少句神聖的話語,及不上那種對你負面說話帶來的壞影響,.
雅各提醒我們,我們不要覺得自詡說很多屬靈套語就變成一個神聖的人,.

$^{241}$禁戒自己的嘴巴,不要說那麼多負面的話,檢視我們的內心,.
問身邊和我們一起相處的人,我的嘴巴長成怎樣,來重新改變我們自己..
詩篇十九章十四節,詩人都是這樣操練自己,.
也說我的盆石是我的救贖主,願我口中的言語,心中的意念在你面前蒙月立,.
不是純粹說讚美的話蒙月立,我心中的意念,我說的那句是否真誠的讚美,.
我口中有沒有惡言的檢視,成為我生命中最重要的讚美神的一種方式..
有些弟兄姊妹或牧師,如果我真的很有心,知道我很想操練自己,.
我改變我自己說話的方式,我可以怎樣做呢?.
其中一個在靈修裡面有一個禁言的操練,不說話,.
這個不是我作的,而是很多外面的修會,或者基督教的團體,.
有些靜修地方都會訓練這一個禁言的操練,.
禁言的意思,第一,可能你一個星期撥一段時間,安靜,.
安靜的意思不是純粹staycation,今天覺得安靜是staycation,.
在酒店住一下,ok的,不是錯的,不過在酒店裡面,如果真的想操練,.
那天不要說太多話,不說話,就是閱讀,看聖經,.
容讓自己的口真的不說任何一句的聲音,.
有些天主教修院,甚至他的飯堂寫到明,.
吃飯的時候不要談話,那是完全安靜的修院,.
有一次我還說那個湯很鹹,被人刮著我,覺得我沒禮貌,.
說那個湯很鹹,其實是那裡不准說話,那裡是靜悄悄的,.
這個操練對於特別喜歡說話的弟兄姊妹來說,是一個很大的挑戰,.
我昨晚跟我女兒說,我女兒其實也很多口,我問她,這麼嚴的操練,.
她說,爹地,我死定了,如果不說話,就是正正這樣,你才要操練,.
喜歡說話的人,試一下操練自己不要說話,想清楚才說,.
這個是特別是一些修院式的操練,第二個方法,在日常生活裡面操練,.
讓你的知識置於你的說話之前,這個是廖炳棠博士,.
在我建築神學院的靈修老師,他在他的屬靈修神學裡面有教的,.
他在說怎樣操練舌頭,操練舌頭就是操練安靜,.
不要讓舌頭跑在我們知識的前面,你的口不要比你的腦快,.
不要比你對一件事物的掌握和知識快,所以對談論的事,.
如果你知道多一點,你的知識多一點,你知多一點,你可以多說一點,.
你少知識,你對那件事不掌握,你知道多少你就說多少,.
少知識就少說話,沒有知識就不要說話,這個就是在我們平時生活裡面的操練,.
如果你能夠將你的知識置於你的舌頭之前,.
你講話就會慢慢有一種操練和一種意識,你不會想講甚麼就講甚麼,.
言論自由不是在講你講甚麼就講甚麼,言論自由是在講你能不能講一些真誠的話,.
是否講得體的話,不是喜歡講甚麼就講甚麼,.
這個不是聖經裡面給我們講的言論自由,而是對生命裡面,.
我們今天重視的是怎樣去操練我們,讓我們的心不會被惡念,被惡語來充斥我們的嘴巴,.
這是第五至十二節的第一個觀點,第十三至十八節的第二個觀點,.

$^{281}$這裡重複出現的就是智慧這個字,很多的識經學家將這個段落不放在舌頭上,.
因為覺得智慧好像沒有講話,好像是另一個新的主題,.
不過當你細心去看的時候,你看到第十三節,你們中間誰是有智慧有見識呢?.
這句話就回應之前講的三章第一節裡面,那些不要多人作師傅這個主題,.
師傅有見識有智慧才能做人家的師傅,他就用這個角度繼續引伸,.
你怎樣去將你的見識智慧表達出來,他這裡面有正反的智慧,.
來自上帝的智慧就是在你當在智慧的溫柔上顯出你的善行,.
你的智慧就在溫柔和善行當中表達出來,這些就是屬神的智慧,.
你們心裡若懷著苦毒,嫉妒和紛爭,就不可自誇,也不可說謊話抵擋真道,.
這樣的智慧不是從上頭來的,乃是屬地的,屬情慾的,屬鬼魔的,.
在何處有嫉妒紛爭,就在何處有擾亂和各樣的壞事,.
唯獨從上頭來的智慧先是清潔,後是和平,溫涼柔順,.
沒有憐憫,多潔善果,沒有偏見,沒有假冒,並且使人和平的,.
是用和平所栽種的異果..
在這裡面講的智慧,額角在講,有屬地的智慧,屬鬼魔的智慧,.
也有屬神的智慧,這些智慧表達出來,如果屬神的智慧,.
就會將上帝的善行和溫柔,後面也有個清單,.
清潔,和平,柔順,憐憫,善果,沒有偏見,沒有假冒,使人和平,.
這些善果或者美德就彰顯了出來,.
不過如果屬地的智慧,同樣也有這些果實出來,.
這些果實是什麼?就是自誇,質度紛爭,抵擋真道,.
這些負面的人際關係,就會在地上的智慧裡面呈現出來,.
這裡很清晰,從你的果子辨別一下,.
在你的生命裡面,你流露著一種什麼智慧?.
曾經有個弟兄姊妹問我,她很多次轉工,.
她很慨嘆,時不與我,經常沒有一個好的工作環境,.
最初我也安慰她,也是的,還年輕,找一份新的工作,.
不要放棄,鼓勵她,一直做的時候,沒多久又辭工,.
沒多久又再辭工,沒多久又再辭工,一年裡面可能辭兩三次工,.
我覺得未必是純粹工作環境的關係,.
然後我問了一條比較尖銳的問題,.
我就問弟兄姊妹有沒有想過,.
為什麼你經常覺得不想跟隨你的這些,.
或者轉工,或者工作環境很不理想的那些狀況,.
為什麼要跟著你走,好像天氣先生那塊烏雲一樣跟著你呢?.
究竟是工作環境的問題,還是你的出現會令這些放大了呢?.
這件事在弟兄姊妹裡面從來沒有想過,.
她一直想都是自己出來打工,.
她沒有想過原來她的出現可能會製造更多的紛爭,.
製造更多人際關係的矛盾,.

$^{321}$說得白一點,可能是弟兄姊妹在她的生命裡面出了一些問題,.
所以每一個狀況都出現問題,.
我這樣說好像很負面,.
也有一些弟兄姊妹在她的工作裡面,.
無論她轉哪一份工,.
那個公司都會帶來很多的笑聲,很多的快樂,.
很多很美好的回憶,很多很好的關係,.
她的出現又為那個地方帶來一些很正面的果效,.
我都會問,你的出現會放大了負面的東西多一些,.
還是放大了正面的東西多一些呢?.
同樣的弟兄姊妹,在這裡面,亞國都提醒我們,.
我們今天在我們身處的狀態當中,.
當你出現的地方會不會多了一些質度紛爭,.
甚至是是非非,在你公司的茶水間,.
在你公司的某個暗黑的角落,.
會不會成為你談論,或者糟蹋某個人的一個地方,.
或者是歷數人家不是的一個地方,.
我們心裡面可能會覺得很不公義,.
但反過來,我們都在另類裡面,.
我們用我們自己的嘴巴,就坐,說是道非,.
令到那個人永不翻身,我講這番說話都是這樣講,.
我不是用一個道德高地來講這些說話,.
而是我很明白,我作為一個很喜歡說話的人,.
原來我自己很多時候都會是這樣,.
我們只是互相面對面來,不要容讓自己落入在舌頭亂講話,.
不造就的人身上的這些試探當中,.
我們經常以為我們講一些評論別人,.
我們就顯出我們自己很有智慧,.
我們比別人高,我們比別人叻,.
但實際上是我們用熟地的智慧,.
來互相咬相吞,踩著別人,或者我們用這個舌頭來插在身邊很多人,.
求主幫助我們,如果我們不好好制服我們的舌頭,.
我們做多少努力,多少善功都好,.
我們一句話可以摧毀很多很多我們覺得很好的果效,.
我們傳多少福音都好,當我說一句是非,.
將一個人踩到地底來糟蹋一個人,.
正如耶穌基督說的,就好像將石頭生在小子身上,.
將他掉到水裡一樣,讓他無法浮起來,.
我們將一個人變得沉淪,我們所說的話,.
如果我們不謹慎下去,我們會成為一個拆毀這個世界的人,.

$^{361}$近來我讀這段聖經的時候,我在想,.
究竟我們平時要說的造就的話,我們什麼叫造就呢?.
是不是什麼都要稱讚呢?盲讚一個人就等於造就呢?.
又或者看到一些人做得不理想的事情,.
將我覺得理想的想法套在別人身上,這些叫造就呢?.
我看一本書,有一個作者我很喜歡的,.
之前也說過,叫做《老樣的貓頭鷹》,.
他有一本書的名字叫做《迷人的雲彈贏得尊重》,.
《窩囊的好人忘求認同》,其中一個chapter,.
其中一個他在說一些人的說話,.
他說有些人說話會怎樣說呢?.
他說生活中總會有一些閒人,.
他寫的很直率,.
習慣了對別人生活橫加指責,指指點點,.
他們有意無意嘲笑別人的努力,驚擾著別人的幸福,.
他們的邏輯是,我這樣是為你好,.
所以我做的是對的,你是錯的,.
我這樣是關心你的,所以你一定要聽我的,.
這樣代表著我很在乎你,你一定要改,.
你得回復我喜歡的狀態,我是為你好,.
所以你不能不領情,否則你就是自私,無情無義..
即是那些人喜歡指點江山的人,.
然後他再繼續說,他說真正活得好的人,.
其實忙於什麼呢?忙於熱愛自己的生活,.
只有那些過得不好的人,.
才喜歡用自己的無聊,庸俗淺薄,.
去驚擾別人的幸福,指點別人的生活..
這句話真是有一點點尖酸,.
不過我想,今天我們說的詛咒的話,.
我們覺得為別人好,.
我覺得我們為了身邊的人,.
想要帶來好處,但其實我是否真心去背負這個人,.
我真的為著我所說的話,來對他負責任地幫助他,.
或是那個人的生活不合我口味,.
我喜歡指點一下他江山,.
嘲弄他生命付出的努力,.
在別人做得不太好的裡面,.
來增高自己,獻自己有智慧..
有時候可能我們走錯路了,.
真正造就的話,可能是我們看到別人,.

$^{401}$雖然在很多艱難的處境裡面,.
他們都在找一些努力,去克服眼前這些困難..
有些人可能他生活跟我們不一樣,.
但他都很忠誠地活在上帝的話語當中..
有些人可能感覺上傻傻笨笨,.
不過可能他正正是實踐聖經的真理,.
跟這個世界背道而馳,.
反而我們立心不良,.
不需要負責任地批評別人,.
只不過是否定別人的努力..
讀完這套聖經之後,.
我立志,我希望我的說話,.
是看到別人的努力,幫助別人,.
找到別人努力實踐信仰的那種單純的心..
我更加希望每一次,.
等姐妹跟我聊完之後,.
她的心靈會有力量,.
不會覺得自己無能為力,.
不會覺得自己是一個窩囊,.
反而是有更加決心,.
來實踐聖經給我們每一個的教導..
如果我們每一個這樣實踐,.
可能我們說不上是一個師父,.
但是在我們的話語裡,.
我們慢慢操練自己,.
成為一個完全的人..
求主繼續幫助我們,.
藉著今天這一段經文提醒我們,.
讓我們好好善用我們的嘴巴,.
記住,一個小小的器官,.
卻是帶來生命很大的果效,.
至於是美好的果效,.
還是負面的果效,.
看你怎樣說下一句話..
求主繼續幫助我們..
\newpage



\section{哥林多後書 11:16-12:10}
\label{sec:OhUT6e4ffZM}
\textbf{2021年8月15日崇拜直播｜梁家麟牧師｜愚拙的保羅｜哥林多後書11\_16-12\_10}
\newline
\newline
連結: \href{https://youtube.com/watch?v=OhUT6e4ffZM}{\texttt{https://youtube.com/watch?v=OhUT6e4ffZM}} ~~~~ 語音日期: 2021-08-16
\newline
\newline
\hyperref[sec:z9XCLH0XmKE]{\small{< < < PREV SERMON < < <}}
~
\hyperref[sec:index_chronic]{\small{[返順時目]}}
~
\hyperref[sec:0pujbt8sHO8]{\small{> > > NEXT SERMON > > >}}
\newline
\newline
哥林多後書 11:16-12:10
\newline
\begin{longtable}{cl}
\hline
\hline
章節 & 經文 (和合本修訂版)\\
\hline
11:16 & \begin{tabularx}{0.7\textwidth}{X} 我再說，誰都不可把我看作愚蠢的；即使你們把我當作愚蠢人，那麼，也讓我稍微誇誇口吧。 \end{tabularx} \\ \\ \relax
11:17 & \begin{tabularx}{0.7\textwidth}{X} 我說的話不是奉主的權柄說的，而是像愚蠢人具有自信地放膽誇口。 \end{tabularx} \\ \\ \relax
11:18 & \begin{tabularx}{0.7\textwidth}{X} 既然有好些人憑著血氣在誇口，我也要誇口了。 \end{tabularx} \\ \\ \relax
11:19 & \begin{tabularx}{0.7\textwidth}{X} 你們是聰明人，竟能甘心容忍愚蠢人！ \end{tabularx} \\ \\ \relax
11:20 & \begin{tabularx}{0.7\textwidth}{X} 假若有人奴役你們，或侵吞你們，或壓榨你們，或侮辱你們，或打你們的臉，你們居然都能容忍。 \end{tabularx} \\ \\ \relax
11:21 & \begin{tabularx}{0.7\textwidth}{X} 說來慚愧，在這方面好像我們是太軟弱了。然而，我說句蠢話，人在甚麼事上敢誇口，我也敢誇口。 \end{tabularx} \\ \\ \relax
11:22 & \begin{tabularx}{0.7\textwidth}{X} 他們是希伯來人嗎？我也是。他們是以色列人嗎？我也是。他們是亞伯拉罕的後裔嗎？我也是。 \end{tabularx} \\ \\ \relax
11:23 & \begin{tabularx}{0.7\textwidth}{X} 他們是基督的用人嗎？我說句狂話，我更是。我比他們忍受更多勞苦，坐過更多次監牢，受過無數次的鞭打，常常冒死。 \end{tabularx} \\ \\ \relax
11:24 & \begin{tabularx}{0.7\textwidth}{X} 我被猶太人鞭打五次，每次四十減去一下； \end{tabularx} \\ \\ \relax
11:25 & \begin{tabularx}{0.7\textwidth}{X} 被棍打了三次，被石頭打了一次，遭海難三次，一晝一夜在深海裡掙扎。 \end{tabularx} \\ \\ \relax
11:26 & \begin{tabularx}{0.7\textwidth}{X} 我又屢次行遠路，遭江河的危險，盜賊的危險，同族人的危險，外族人的危險，城裡的危險，曠野的危險，海中的危險，假弟兄的危險。 \end{tabularx} \\ \\ \relax
11:27 & \begin{tabularx}{0.7\textwidth}{X} 我勞碌困苦，常常失眠，又飢又渴，忍飢耐寒，赤身露體。 \end{tabularx} \\ \\ \relax
11:28 & \begin{tabularx}{0.7\textwidth}{X} 除了這些外表的事以外，我還有為眾教會操心的事天天壓在我身上。 \end{tabularx} \\ \\ \relax
11:29 & \begin{tabularx}{0.7\textwidth}{X} 有誰軟弱，我不軟弱呢？有誰跌倒，我不焦急呢？ \end{tabularx} \\ \\ \relax
11:30 & \begin{tabularx}{0.7\textwidth}{X} 我若必須誇口，就誇我軟弱的事好了。 \end{tabularx} \\ \\ \relax
11:31 & \begin{tabularx}{0.7\textwidth}{X} 那永遠可稱頌之主耶穌的父神知道我不說謊。 \end{tabularx} \\ \\ \relax
11:32 & \begin{tabularx}{0.7\textwidth}{X} 在大馬士革的亞哩達王手下的提督把守大馬士革城，要捉拿我， \end{tabularx} \\ \\ \relax
11:33 & \begin{tabularx}{0.7\textwidth}{X} 我被人用筐子從城牆上的窗口縋下，逃脫了他的手。 \end{tabularx} \\ \\ \relax
12:1 & \begin{tabularx}{0.7\textwidth}{X} 雖然自誇無益，我還是不得不誇。我現在要提到主的異象和啟示。 \end{tabularx} \\ \\ \relax
12:2 & \begin{tabularx}{0.7\textwidth}{X} 我認識一個在基督裡的人，他在十四年前被提到第三層天上去；或在身內，我不知道，或在身外，我也不知道，只有神知道。 \end{tabularx} \\ \\ \relax
12:3 & \begin{tabularx}{0.7\textwidth}{X} 我認識的這樣的一個人—或在身內，或在身外，我都不知道，只有神知道— \end{tabularx} \\ \\ \relax
12:4 & \begin{tabularx}{0.7\textwidth}{X} 他被提到樂園裡，聽見隱祕的言語，是人不可說的。 \end{tabularx} \\ \\ \relax
12:5 & \begin{tabularx}{0.7\textwidth}{X} 為這人，我要誇口；但是為我自己，除了我的軟弱以外，我並不誇口。 \end{tabularx} \\ \\ \relax
12:6 & \begin{tabularx}{0.7\textwidth}{X} 就是我願意誇口也不算狂，因為我會說實話；只是我絕口不談，恐怕有人把我看得太高了，過於他在我身上所看見所聽見的； \end{tabularx} \\ \\ \relax
12:7 & \begin{tabularx}{0.7\textwidth}{X} 又恐怕我因所得的啟示太高深，就過於高抬自己，所以有一根刺加在我身上，就是撒但的差役來折磨我，免得我過於高抬自己。 \end{tabularx} \\ \\ \relax
12:8 & \begin{tabularx}{0.7\textwidth}{X} 為了這事，我曾三次求主使這根刺離開我。 \end{tabularx} \\ \\ \relax
12:9 & \begin{tabularx}{0.7\textwidth}{X} 他對我說：「我的恩典是夠你用的，因為我的能力是在人的軟弱上顯得完全。」所以，我更喜歡誇耀自己的軟弱，好使基督的能力覆庇我。 \end{tabularx} \\ \\ \relax
12:10 & \begin{tabularx}{0.7\textwidth}{X} 為基督的緣故，我以軟弱、凌辱、艱難、迫害、困苦為可喜樂的事；因為我甚麼時候軟弱，甚麼時候就剛強了。 \end{tabularx} \\ \\
[1ex]
\hline
\hline
\end{longtable}
$^{1}$今日的經訓是記載在.
哥林多後書第十一章十六至三十三節.
在聖經新約第二百五十四至二百五十五頁.
這裡記載了保羅作仕途所受的患難.
請聽保羅的說話.
我在說人不可把我看作愚忘的.
縱然如此也要把我當作愚忘人接納.
叫我可以略略自誇.
我說的話不是奉主命說的.
乃是讓愚忘人放膽自誇.
既有好些人憑著血氣自誇.
我也要自誇了.
你們既是精明人.
就能甘心忍耐愚忘人.
假若有人強你們作奴僕.
或侵吞你們.
或擄掠你們.
或舞蔓你們.
或打你們的臉.
你們都能忍耐他.
我說話是羞辱自己.
好像我們從前是軟弱的.
然而人在何事上勇敢.
我說句愚忘話.
我也勇敢.
他們是希伯來人嗎.
我也是.
他們是以色列人嗎.
我也是.
他們是亞伯拉罕的後裔嗎.
我也是.
他們是基督的僕人嗎.
我說句狂話.
我更是.
我比他們多受勞苦.
多下監牢.
受鞭打是過重的.
誣死是屢次有的.
被猶太人鞭打五次.
每次四十減去一下.

$^{41}$被棍打了三次.
被石頭打了一次.
遇著船壞三次.
一晝一夜在深海裡.
又屢次行軟路.
遭江河的危險.
盜賊的危險.
同族的危險.
外邦人的危險.
城裡的危險.
曠野的危險.
海中的危險.
假弟兄的危險.
受勞碌.
受困苦.
多次不得睡.
又饑又渴.
多處不得食.
受寒冷.
赤身怒體.
除了這外面的事.
還有為宗教會掛心的事.
天天壓在我身上.
有誰軟弱我不軟弱呢.
有誰跌倒我不焦急呢.
我若必須自誇.
就誇那關乎我軟弱的事便了.
那永遠可稱頌之主耶穌的父神.
知道我不說謊.
在大馬式的阿里達王手下的提督.
把守大馬式城.
要捉拿我.
我就從窗戶中在康子裡.
從城牆上被人墜下去.
脫離了他的手.
我們繼續看保羅的自白.
前面一段批評超級仕徒的說話之後.
保羅又回到段落的主題.
愚拙人 蠢人.
在前面保羅已經呼籲過那麼多信徒.

$^{81}$要包容他.
在此他再做一次請你們包容我的呼籲.
十一章十六到二十節.
我再說人把我看作愚妄的.
縱然如此又把我當作愚妄人接納.
叫我可以略略自誇.
我說的話不是奉主名說的.
是奉主名說的.
既有好些人憑著血氣自誇.
我又要自誇了.
你不記事精明人.
就能甘心忍耐愚妄人.
有人強你作勞瀑.
或侵吞你們或擄掠你們.
或舞慢你們或打你們的臉.
你們都能忍耐他.
縱然如此.
退一步說.
保羅說你不要當我傻的.
就算你當我傻.
也請你包容我的傻.
就算我現在蠢蠢的.
為何我蠢蠢的呢.
因為保羅說我開始要自誇了.
我要說我自己的事.
保羅一直不喜歡說自己的事.
他不喜歡說自己的強.
他連自己的弱都不喜歡說.
你知我們很多時候都會碰到一些人.
他很喜歡說自己的問題.
自己的困難自己的需要自己的危機.
明明是別人說自己的困難.
他都會說自己有類似的經驗.
然後說一大堆自己的事.
很多時候這是一個.
關注人的做法.
保羅不喜歡.
保羅喜歡成為一個.
關心其他人的人.
不喜歡說自己.

$^{121}$所以保羅強不說.
弱都不說.
覺得他被迫要說自己.
所以他覺得這是一個蠢人才做的事.
所以保羅認定他做蠢人.
他就預先聲明.
這裡所說的話不是來自基督的訓示.
而是隨人而,憑血氣做的.
並且我是仿效那些.
哥林多教會裡自誇的愚蠢人.
意思是甚麼呢.
你要跟我打爛仔交.
做那些挖眼療陰.
扯頭髮吐口水.
好,我奉陪你一次.
不過先此聲明.
這些爛仔交不是黃飛鴻師父教我的.
我自小做街童的時候.
自己無私自通的.
所以我現在打的招式.
跟黃師父無關.
千萬不要破壞他的清譽.
保羅說不是從主領受的.
而是我自己做出來的.
保羅再用一個諷刺的口吻說.
你們既是精明人.
就能夠甘心忍耐愚妄人.
假如有人強你們作奴僕.
或者侵吞你們.
或者擄掠你們.
或者打你們的面.
你們都能忍耐他.
這句話翻成今天的說話就是.
我是蠢人.
而你們自己覺得自己很聰明.
你們有多聰明呢.
你們聰明到一個地步.
竟然可以被人去核制你們.
剝削你們佔你們便宜.
侮辱你們虐待你們.

$^{161}$他們這樣待你們.
你們還能忍耐.
真是逆來順受.
你們真是聰明到絕頂.
這裡保羅說了一大堆諷刺性的話.
說了五個動詞.
剝削佔便宜.
侮辱和虐待.
我們問問知道保羅所說的情況.
是不是真的說那些超級仕途或者假師傅.
在哥林多信徒身上做了那麼多傷害.
核制之餘什麼.
剝削之餘什麼.
其實我們都不太知道.
不過當然保羅整個主題都是說.
那些假師傅或者超級仕途.
會使得哥林多信徒離開真道.
失卻失落正確的信仰.
對保羅來說.
沒有什麼比剝去一個人的正確信仰更嚴重.
所以他認為這是一個最大最大的傷害.
所以那些核制,剝削,佔便宜,侮辱.
到底是什麼都已經變得不重要.
保羅說.
你在那些假師傅面前.
顯出你們的軟弱.
到一會兒在我這個仕途面前.
在我這個淑寧生父面前.
你們就顯得也文也武.
好像很堅強.
好像我人善就被人欺負.
我扮個善良的樣子就被你們欺負.
保羅接著說.
我從前對你們太過溫柔.
可能是做錯了.
我真是對不起你們.
保羅接著說.
人在何時相勇敢.
我說句如網話.
我也勇敢.

$^{201}$人是anyone.
指的是那些高舉自己的人.
保羅說.
其實他們可以,我也可以.
What they are, I am also.
你說他們可以,我也可以.
其實我也有條件做剛強人.
這句說話翻成廣東話.
你別當我是Lulu.
我也可以用強者面貌示人.
所以保羅就像開水喉一樣.
說了一大堆他能夠剛強的經歷.
最初他和那些超級仕途做比較.
他們有的,保羅也有.
然後再說他獨有的經歷.
保羅有的,他們沒有.
先說共同有的經歷.
Me too,他有,我也有.
他們是希伯來人嗎?我也是.
他們是以色列人嗎?我也是.
他們是亞伯拉罕的後裔嗎?我也是.
他們是基督的僕人嗎?我說句狂話.
我更是.
從種族,從血緣出身.
乃至從信仰身份.
保羅說我沒有一樣東西不及其他人.
他們是我通通都是.
希伯來人強調的是血緣的純正.
保羅說我和他們一樣正統希伯來人.
以色列人強調的是被揀選的上帝子民的一員.
亞伯拉罕的後裔強調他是有份繼承上帝給亞伯拉罕的應許.
所以保羅用以色列人的身份和超級仕徒比較.
證明他們有的,我通通都有.
不過保羅其實要強調的是最後一句.
他們是基督的僕人嗎?我說句狂話.
我更加是.
因為前面說到希伯來人,以色列人,亞伯拉罕的後裔.
最多是說到人有我有,不會有比較性.
不會說你是香港人,我更是香港人,什麼意思?.
大家都是香港出生,就是香港人,對不對?.

$^{241}$沒有所謂更是,沒有比較的,是就是,不是就不是.
除非說我是混種的,不過這也不是很重要.
香港人和混種沒有關係.
就算是歐洲血緣的,你在香港出生也是香港人.
巴基斯坦血緣的,你也是香港人,香港出生.
所以沒有比較可說的,但是基督的僕人不是天生的.
所以就有得比較了,所以他說是不是基督的僕人呢?.
這是最重要的身份,保羅說我更是.
我更是,怎樣更是呢?.
接著他就說出他和那些假仕途,或者是那些超級仕途的不同之處.
比較了同,接著就說不同,哦,不同處就很多了.
由23節,下半節一直到26節.
我比他們多受勞苦,多下監牢,被受劫鞭打是過重的,誣寫是屢次有的.
被猶太人鞭打五次,每次四十,搞起一下被棍打了三次.
被小孩打了一次,遇著損壞三次,一走一夜在深海裡.
又屢次行遠路,遭江河的危險,盜賊的危險,同族的危險,外邦人的危險,城裡的危險,曠野的危險,海中的危險,賈帝卿的危險.
保羅細數他為這個仕途身份所付出的各種代價.
勞苦,我比他們更勞苦,被囚,我比他們坐監更多.
受鞭打,我比他們打得更多,遭生命的危險,誣寫是屍體.
吃生菜這麼多,是是有的.
受勞苦主要指著各樣的困苦,主要是身體上的痛苦.
冒死,即是在鬼門關死亡的邊緣.
保羅強調在所有仕途中間,他拿了四個第一名.
他的受勞苦率排行第一,他坐監次數排行第一,他被打的次數排行第一.
他和死神擦過的次數排行第一,四甲超級狀元,四個五星星.
這真是光輝英雄的紀錄,不過這樣的紀錄令人厭倦嗎?.
你欣不欣賞,你期不期待成為保羅這樣的樣子呢?.
你可以看到保羅的經歷,和我們常說的成功神學所宣講的.
信耶穌會得到上帝特別的保守,健康財富陸續有利,是完全佔不上邊的.
只可以說保羅全說的是失敗神學,倒楣神學.
想勞苦?信耶穌,想坐監?信耶穌,五行欠打?信耶穌,要命長?信耶穌.
保羅再為這四種代價做了量化的統計.
他先說後面那兩種,受鞭打,他說被猶太人鞭打了五次.
每次四十減去一下被棍打了三次,被石頭打了一次.
為什麼這麼清楚,每次四十減去一下這麼奇怪呢?.
因為猶太人根據新命紀二十五章二到三節的規定.
對犯事者的鞭打不可以超過四十,否則明明是打就變成殺.
所以每次只能夠四十,他記得被猶太人打過五次,一定是四十下.
不過他不想誇大,有一次不知為什麼那個人手軟,或者是不記得數錯.
打少了一次,他就說明,我被人打了五次,每次四十.

$^{281}$就差了一下,就三十九,就是這樣,本來是不是很心碎清?.
接著說到被棍打,被棍打沒有明講,不過這個應該是指著被羅馬地方政府的官員打的.
這個是羅馬政府給他的刑罰,雖然他是羅馬公民理論上是不應該受體罰的.
不過很多地方政府根本沒有執行這個規定的.
至於被石頭打的經驗,在《少行傳》有記載的.
那個記載在《少行傳》十四章五節,那個應該是發生在二哥念的.
被石頭打應該不是來自官方當局的,因為如果來自官方當局.
石頭打是很嚴重的,最嚴重的,因為叫stone to death,是被石頭扔死的.
所以這個理論上是來自部門,民間做的,不是官方做的.
《少行傳》和《保羅書信》其實沒有記載保羅的生平.
所以我們沒辦法就保羅數的經歷的次數,夾到他的經歷裡面.
哪一次在哪裡發生,我們數不到的.
接著保羅講到冒死,閻名祥和死神察覺.
他講到遇著船壞三次,一咒一謝,在深海裡,屢次行遠路.
遭江河危險,盜賊危險,同仇危險,外邦人危險,城裡危險.
曠野危險,海中危險,假兄弟危險.
船壞三次,保羅經常要坐船,往來地中海各地.
遇到船壞一點都不奇怪.
屢次行遠路,渡海可以坐船,陸路就要靠雙腳.
江河危險,大概是指著河水暴漲的時候,渡河危險.
盜賊同仇,外邦人,指著那些加害他的不同人.
城裡,曠野,海中,指著那些發生危險的地點.
最後保羅特別提到假弟兄,指著他在教會服侍的時候.
和那些傳義端的人對抗,所受到的迫害.
再來是為他所說的勞苦這件事,做具體的說明.
單單說勞苦,可能很籠統,即是我陪老婆去街買東西.
我也覺得很勞苦,超級勞苦.
或者你們聽我說半英頭道,也覺得很勞苦,很難受.
所以一定要具體說,勞苦是甚麼勞苦.
保羅說他的勞苦包括了內在兩部分.
外面的事,受勞碌受困苦,多次不得睡又飢又渴.
多次不得食受寒冷赤身勞體,這些就是外面的事.
我們沒辦法具體知道是甚麼事,不過猜也猜到.
多次不得睡可能指著要行夜路,或者要警醒禱告.
或者因為擔心很多事而睡不著.
多次不得食或者沒錢買食物,或者途中供應不上.
或者因為宗教原因禁食,赤身勞體指著沒錢買衣服.
或者被盜賊打劫,剝肛豬一無所有.
這些都可能,我們沒辦法具體知道是甚麼事件.
保羅最想說的是他連最基本的生活需要.

$^{321}$吃喝穿衣睡覺都想像到,他過不到一個人的起碼生活.
然後說到內在的勞苦,他說為宗教會掛心的事.
天天壓在他心上,這是我很感動的說話.
保羅將外在的勞苦,遭遇生命危險的.
和他掛念教會,擔心教會的需要.
而有的那種心理壓力,是相提並論.
保羅說具體說,有誰軟弱我不軟弱呢?.
有誰跌倒我不焦急呢?.
他將弟兄姊妹的需要完全包在身上.
他們難過我就難過,他們跌倒我就困苦.
如果你將身邊所有人的需要都包在身上.
真是先天下之憂而憂,你一定是常懷憂患.
一定有強烈的憂患意識.
保羅的誇口清單大部分都是他受患難的經歷.
雖然在這些患難中上帝都答應了他.
但整張清單仍然壓倒性是負面的事.
保羅承認這一點,不過他強調這是他刻意的佈局.
因為他說誇口,如果人要自吹自擂.
只能夠指著主誇口,不能誇口自己.
只能高舉主的大能,主的拯救.
但經歷主大能的前提是人自覺軟弱.
經歷主的拯救的先決背景是人遭遇患難.
沒有人的軟弱,如何經歷主的大能?.
沒有患難,如何需要拯救?.
所以保羅在上一章三十節說.
我若必須自誇,就誇那關乎我軟弱的事便了.
在講述患難清單時,保羅突然加插了一段小故事作為清單結束.
三十一至三十三節,我永遠可稱頌之主耶穌的父上帝知道我不說謊.
他要很鄭重,用誓願的方法說,我以下說的一點都沒有誇口.
在上帝可以知道我沒有誇張.
什麼事呢?.
在大馬色,阿彌陀佛王下手下提督.
把守大馬城要捉拿我.
我從窗戶中從康子裡從城牆被人追下去,脫離了他的手.
這對保羅來說是驚心動魄的經歷.
被人追殺,午夜逃亡,並且從城牆上掉下去.
到底是什麼事呢?我們無法知道.
總之是保羅記憶猶新,想起都會打冷顫,流冷汗的故事.
我那次被人追殺,很嚴重.
就是這件事,作為他苦難清單的總結.

$^{361}$我們讀保羅的苦難經歷到這裡就夠了.
保羅為什麼會罕罕而談他的受苦經歷呢?.
他心理不正常,以苦為樂,以苦為榮.
他受希臘哲學那種否定肉體,否定人的慾望的思想影響.
當然不是.
這段內容好像集中在受苦經驗那裡.
但其實保羅要傳遞的訊息根本不是受苦.
而是他要強調忠於使命,使命必達.
為了完成基督教父的福音使命.
保羅不用聽李靖的提議,他不看鬼片也膽生毛.
龍潭虎穴,刀山火海,他也敢闖進去.
天災人禍,被打坐監受辱,一切代價.
他眼眉不皺就付出,甚至他連性命都可以不顧.
保羅列舉他受過的種種苦難和凌辱.
都是為了說明他的無比決心,他的堅毅忠誠.
保羅說了一句豪氣乾雲的話.
他們是基督的僕人嗎?我說一句狂話,我更是.
他的苦難清單目的就是要說明這句話.
他們是基督的僕人嗎?我更是.
基督的僕人是保羅的一張名片.
是他一生前進的起點,也是他要奮鬥的目標.
又是Alpha,又是Omega,我事奉的起點,因為我是基督的僕人.
我事奉的終點,要證明我是基督的僕人.
我一生要證明這件事.
這個讓我聯想起最近看過一個新聞的故事.
不是刻意找的,是看到之後很有感動.
7月河南鄭州豪雨成災,最嚴重的事故地點是5號地鐵線.
發生水浸,造成十幾人死亡.
在看新聞報道的時候,又說了一個小故事.
說到一個被車廂關在裡面的女乘客.
最後僥倖獲救,不過獲救的時候,因為缺氧的緣故,都昏迷了一段時間.
當她醒過來的時候,聽到身邊有人大聲喊,有沒有醫護人員.
她立刻彈起來,她就回應,她說我是醫生.
然後她和另外兩個起來說自己是醫護人員的人,三個人.
他們在災場奮鬥了幾個小時,為那些被水浸的地道救出來昏迷的傷者進行心肺復甦.
她後來憶述,她一共為十個人進行了急救.
不過她說她不知道那些人最終有沒有獲救.
這個女士就是鄭州大學第一附屬醫院的中學科的醫生.
她的姓氏是秦,秦始皇的秦,英雄豪傑的傑,林是樹林的林.
今年32歲,你在網上很容易找到她的資料.

$^{401}$我是醫生,我覺得這四個字千棍尋蟲.
這四個字蓋過了我都是病人.
甚至更加簡單的,我都是人,內心的呼喚.
不管自己的身體狀況如何,不管自己的心有多麼懼怕.
希望可以盡快離開地鐵月台,上地面,對不對.
回家休息,是不是?.
我是醫生,這四個字就寫到秦,傑,林義無反顧的撲向病人,遇難者的身邊.
做一個醫生要做的事.
同樣道理,我是基督的僕人.
這個告白也寫到保羅,忘卻生命的危險,肉體的痛楚,生理心理的自然渴求.
弟兄姊妹,我知道我講什麼都會被人罵.
這段日子應該啞,不應該講話.
不過被人罵就繼續講下去.
我真的越來越厭聽到一些講法.
牧師都是人,基督徒都是人,基督徒都有七情六慾的辯解.
在這幾年,我們聽到的類似說話實在太多了.
因為牧師都是人,並且人的身份比牧師更大.
所以當遇到危險或危機,牧師當然是自然於不顧.
保住自己的命更重要,對不對?.
畢竟你們這麼多人加起來都不夠我老婆和我那對子女貴重,對不對?.
更加重要的是,我沒有你們這批羊,我還有很多羊可以讓我去牧養.
為什麼一定要抱著一批不放呢?對不對?.
天涯何處無羊牧,對不對?.
更加有牧者,我要說的是厚顏無恥,我要用這些字眼去說.
有人會立即上綱上對號,我要說厚顏無恥宣稱.
我梁家倫才是全香港教會,全環教會最需要保護的人才,你明不明白?.
所以棄羊自救,棄你們這班又便宜又賤的羊去救我自己.
這是最恰當,最有智慧的抉擇.
我才是教會唯一值得保護的人才,對不對?.
你要知道這些說法是得到很多教會有識之士股掌贊同的.
按照這樣的邏輯,一個接受了這麼多年訓練,並且IQ一定很高.
並且一定是讀書很厲害的一個專業醫生.
怎麼值得冒生命危險去營救一個學歷,資歷,身酬,社會地位都比他低的新冠肺炎患者呢?你說,值不值得?.
找個醫生去救一個這樣的病人,又不可以弄死那個醫生,對不對?.
那些便宜的就送出門吧.
我們聽太多牧師也侍人的說法.
正因為人的身份比牧師更大,苗條淑女,君子好求,食色聖也.
所有婚外情都是人之常情,情有可原的,對不對?.
無論如何,人的身份最大,我是人,我先是人,我最重要是人.
牧師只是一個職業,一份工作,只是在特定時空,扮演特定角色.

$^{441}$所以你見到,我進來黃浸才打一條領帶,對不對?.
我在街上也不會打,對不對?.
牧師這個身份在我生命裡不佔主導性,不佔主宰性.
你知道我在牧師之外,我還有一個丈夫,一個父親,一個中大校友.
我是香港人,我是橋牌愛好者,我是亞仙奴勇等,對不對?.
牧師只是我眾多身份的其中之一,和其他身份是互不干涉的.
分庭降禮,沒有誰最大,誰主宰誰的說法,沒有.
我不是否定牧師也是人,基督徒有七成六欲的說話事實性,我承認的.
我也不是站在道德高地去批評達不到高要求的人.
我對人的軟弱跌倒,對理想和現實存在差距,我是趴在這裡,我充分明白,充分接受.
我要說我只覺困擾的是,很多剛才說的話的人,竟然是理直氣壯,視為理所當然的.
一點怯懦都沒有,我有份愛情被你捉到,捉到就捉到,對不對?.
不行嗎?牧師是人,對不對?你會跌倒,我會跌倒,人人都可能有三拉四拉,有什麼奇怪?.
那種理直氣壯令我受不了,將實言變成應言令我受不了.
並且為了擁抱實言,我們要踐踏我們所說的理想,那些說什麼將基督徒,將牧師,將他神化的說法.
其實是低國主義式的,是文化霸權,於是我們可以將一切理想摧毀.
牧師就由我們的事實去定義,我們是怎樣就是怎樣,什麼叫牧師?就是看我的樣子,我剛剛在戳魚蛋,牧師就戳魚蛋.
就是這樣,不是由聖經所說的理想去定義,耶穌沒有發言權,我們事實上有多爛,牧師就是怎樣去定義.
我只提牧師在這裡,因為我自己是牧師,自己批評自己沒那麼尷尬.
我的說話其實適用於醫生,適用於教師,適用於任何要承擔責任的崗位.
身份蘊含使命,使命必達,要壓上自己作為代價,當然故宮不會這樣想,不會這樣做.
如果你說找牧師來舉例不好,牧師只是一份工作,一個職業,不值得賣掉你的青春,你的家庭和你的生命,那我要問,那基督徒呢?.
或者基督的僕人呢?是不是只是你眾多身份之一?是不是同樣不具備宰制性,主導性呢?.
我是基督的僕人,我更是,在你和老婆吵架,吵到臉紅耳綠的時候,我是基督的僕人,這個身份,這個宣告有沒有發言權?.
在你和身邊的人一起去爭奪,一起去競賽,其他人不擇手段,你覺得你都應該和光同尋,不擇手段的時候,我是基督的僕人.
這個身份有沒有啟示?有沒有王牌,特朗普的作用?有沒有?有沒有?.
保羅為了區區身份的認定,就壓著他自己一生,他的愚昧,他的執迷,他的癲狂,真是到了一個無可救藥的地步.
今日的經文其實是一張癲狂病歷單,弟兄姊妹,你會不會同樣自稱我是基督的僕人?你怎樣理解這個自稱?.
你又準備為這個自稱展示什麼質素?願意為他押上什麼代價?我們彼此互勉.
\newpage



\section{雅各書 4:1-10}
\label{sec:0pujbt8sHO8}
\textbf{2021年8月22日崇拜直播｜吳真真牧師｜生命見真章系列 - 與神為友,與世俗為敵｜雅各書4\_1-10}
\newline
\newline
連結: \href{https://youtube.com/watch?v=0pujbt8sHO8}{\texttt{https://youtube.com/watch?v=0pujbt8sHO8}} ~~~~ 語音日期: 2021-08-22
\newline
\newline
\hyperref[sec:OhUT6e4ffZM]{\small{< < < PREV SERMON < < <}}
~
\hyperref[sec:index_chronic]{\small{[返順時目]}}
~
\hyperref[sec:bJlzZqbRugo]{\small{> > > NEXT SERMON > > >}}
\newline
\newline
雅各書 4:1-10
\newline
\begin{longtable}{cl}
\hline
\hline
章節 & 經文 (和合本修訂版)\\
\hline
4:1 & \begin{tabularx}{0.7\textwidth}{X} 你們中間的衝突是哪裡來的？爭執是哪裡來的？不是從你們肢體中交戰著的私慾來的嗎？ \end{tabularx} \\ \\ \relax
4:2 & \begin{tabularx}{0.7\textwidth}{X} 你們貪戀，得不著就殺人；你們嫉妒，不能得手就起爭執和衝突；你們得不著，是因為你們不求。 \end{tabularx} \\ \\ \relax
4:3 & \begin{tabularx}{0.7\textwidth}{X} 你們求也得不著，是因為你們妄求，為了要浪費在你們的宴樂中。 \end{tabularx} \\ \\ \relax
4:4 & \begin{tabularx}{0.7\textwidth}{X} 你們這些淫亂的人哪，豈不知道與世俗為友就是與神為敵嗎？所以，凡想要與世俗為友的，就是與神為敵了。 \end{tabularx} \\ \\ \relax
4:5 & \begin{tabularx}{0.7\textwidth}{X} 經上說：「神愛安置在我們裡面的靈，愛到嫉妒的地步。」你們以為這話是徒然的嗎？ \end{tabularx} \\ \\ \relax
4:6 & \begin{tabularx}{0.7\textwidth}{X} 但是他賜更多的恩典，正如經上說：「神抵擋驕傲的人，但賜恩給謙卑的人。」 \end{tabularx} \\ \\ \relax
4:7 & \begin{tabularx}{0.7\textwidth}{X} 所以，要順服神。要抵擋魔鬼，魔鬼就必逃避你們； \end{tabularx} \\ \\ \relax
4:8 & \begin{tabularx}{0.7\textwidth}{X} 要親近神，神就必親近你們。有罪的人哪，要潔淨你們的手！心懷二意的人哪，要清潔你們的心！ \end{tabularx} \\ \\ \relax
4:9 & \begin{tabularx}{0.7\textwidth}{X} 你們要愁苦，悲哀，哭泣；要將歡笑變為悲哀，歡樂變為愁悶。 \end{tabularx} \\ \\ \relax
4:10 & \begin{tabularx}{0.7\textwidth}{X} 要在主面前謙卑，他就使你們高升。 \end{tabularx} \\ \\
[1ex]
\hline
\hline
\end{longtable}
$^{1}$《新約聖經》第4章1至10節.
第4章1至10節.
你們中間的真戰鬥毆是從哪裡來的呢?.
不是從你們八體中戰鬥之私欲來的嗎?.
你們貪戀還是得不著?.
你們殺害嫉妒又鬥毆真戰也不能得.
你們得不著是因為你們不求.
你們求也得不著是因為你們妄求.
要浪費在你們的燕樂中.
你們這些淫亂的人.
豈不知與世俗為友就是與神為敵嗎?.
所以凡想要與世俗為友的.
就是與神為敵了.
你們想經上所說是圖言的嗎?.
神所賜住在我們裡面的靈.
是戀愛至於嫉妒嗎?.
但祂賜更多的恩典.
所以經上說.
神阻擋驕傲的人賜恩給謙卑的人.
故此你們要崇服神.
勿要抵擋魔鬼.
魔鬼就必離開你們逃跑了.
你們親近神神就必親近你們.
有罪的人要潔淨你們的手.
深懷異議的人要清潔你們的心.
你們要愁苦悲哀哭泣.
將喜笑變作悲哀.
歡樂變作愁悶.
勿要在主面前自卑.
主就必要你們聖告.
聖經獨筆.
各位弟兄姊妹平安.
我在這裡和曼谷沙吞堂的弟兄姊妹說聲平安.
記得兩年前我們一家四口去到曼谷旅行.
都有機會到沙吞堂和一眾弟兄姊妹一同去敬拜主.
我先生陳牧師也有機會在當中服事.
多謝你們的接待.
我們能夠在跨地域裡.
透過媒體我們可以一起去敬拜上帝.
這是一個很大的福氣.

$^{41}$讓我們一同去學習上帝的話.
今天來到我們一個專題系列.
生命見真章的第六講.
大家還記不記得之前.
三位陳牧師.
都和我們分享了雅各書的教導.
他在當中提醒我們這群信主的人.
我們不要單單只是在頭腦上去認信有這位上帝.
我們必須在生命裡經得起考驗.
提醒我們要在聽到之後.
我們必須要將生命的道去走出來.
無論在我們日常生活的態度.
待人接物.
不要以貌取人.
要將神的話記在心.
並且走出來的時候.
我們要留意我們的說話.
因為我們所行的就是反映我們的所信.
所以生命才見到真章.
今天來到第六講.
我們的題目是「與神為友,與世俗為敵」.
經文是繼續在第四章第一至十節.
其實在第三章十三至十八節.
就提到基督徒可以向神求智慧.
當我們求到智慧之後.
我們的生命就會滿有憐憫.
結滿美好的果子.
並且在當中會由和平的人.
用和平栽種正義的果子.
不過在第一節一開始.
就帶出了一個現象.
就和和平背道而馳.
反而教會就充斥著爭戰鬥毆.
紛爭不斷.
這段經文就以此帶出了什麼是私慾.
指出了教會的信徒在靈性上如何淫亂.
如何與世俗做朋友.
如何與神為敵.
我們稍後會一起去研讀這段經文.
在讀這段經文之前.

$^{81}$讓我們一起將時間.
將我們的心獻上給上帝一起禱告.
主啊!求你閱立我們的敬拜.
是由心發出的敬拜.
願以我們全天秘的心.
受你的話語教導,督責.
使我們歸正.
提醒我們如何讓我們的生命.
與我們所信的相情.
主啊!幫助孩子.
讓我去說出你的話.
能夠讓弟兄姊妹明白.
讓弟兄姊妹心靈裡被提醒.
也願以聖靈引導我們.
得著我們的心.
讓我們能夠回應你對我們說的話.
我們這樣禱告,交託.
是奉主耶穌基督得勝的明智祈求.
阿們.
我們說到今天這段經文.
一至六節就提醒了我們.
告訴我們.
其實我們有些時候.
我們這些基督徒.
都要被提醒.
我們竟然會與世俗為友.
與神為敵.
第一節至第三節.
你們中間的真戰鬥毆.
是從哪裡來的呢?.
不是從你們八般八體中.
戰鬥的私欲來的嗎?.
你們貪戀還是得不著.
你們殺害嫉妒.
又鬥毆,真戰也不能得.
你們得不著是因為你們不求.
你們求也得不著.
是因為你們妄求.
要浪費在你們的言惡中.
這裡一開始就帶出了.

$^{121}$一個很重要的問題.
就是在教會裡有一個現象.
原來你們中間的真戰鬥毆.
是從哪裡來的呢?.
這裡就是說.
原來在群體裡產生了紛爭.
從這個現象說出來.
就是告訴我們.
其實這些就是與世俗為友的表現.
和合本翻譯成為真戰鬥毆.
好像很大件事.
是不是在教會裡有打架呢?.
打架,打仗.
其實原文這個字的確有戰爭的意味.
有戰爭的意味.
也有打架的意味.
但是我們知道.
它是想彰顯一件事.
就是裡面有一些很大張力的敵對關係.
用這個字的原意就是想突出一件事.
就是當中有紛爭,有衝突.
和一些口頭上的爭拗.
所以當和合本修訂版再去翻譯這段經文的時候.
它所翻譯的字眼就是衝突和爭執.
不過我相信雅各用這個字是有它的意思.
因為紛爭的本質的確帶有戰爭的意味.
就是為了自己所信,所堅持,所捍衛的東西.
要用盡一切的方法去打敗對方,去贏.
我們看到在這個紛爭裡.
當我們與世俗為友.
我們會有一些問題出現.
這個問題就是原來我們人與人之間的人際關係.
出現了一些的嚴規和問題.
當中會有紛爭.
對於這個問題.
雅各是用另一條問題去回應.
他的問題就是說.
「不是從你們八體中戰鬥的私欲來的嗎?」.
這個是一個問題.
不過這是一個反問句.

$^{161}$那個意思是肯定的.
就是說這些紛爭是從哪裡來的?.
這些紛爭就是我們肢體之間的私欲.
是我們肢體之間的私欲.
私欲這個字的希臘原文.
其實是解釋我們英文用的一個字.
叫做享樂主義.
原來是說我們為了享樂而有的慾望.
在雅各書這裡的私欲.
多一些是指一些罪惡的慾望.
有學者去分析.
他說其實這個欲望,私欲的本質.
就是自我中心,自誇,貪戀,虛榮,憎勁,惡毒.
以至會用一些掌控權力.
而帶來一些紛爭問題.
他說我們是八體中爭戰而來.
八體的意思是可以有兩個解法.
一個解法就是在我們內心裡各樣的爭戰.
又或者在教會群體裡我們的肢體.
肢體與之間的爭戰.
其實兩種解法都有人說.
不過就是顯示了一個很激烈的.
你爭我奪的局面.
我們看到與世俗為友.
除了人際關係會出現問題之外.
我們看到與世俗為友一個很重要的特色.
就是自我中心.
是用自己的慾望.
原來自己的慾望就是宇宙的中心.
就是宇宙的一切.
會為了宇宙而不擇手段.
然後我們一直讀到第二,第三節.
我們就會看到原來.
私欲是有一個盡情的.
是描繪著一個盡情的.
第二,第三節說.
你們貪戀還是得不著.
你們殺害嫉妒.
又鬥毆爭戰也不能得著.
我們先說這裡.

$^{201}$第二,第三節其實這段經文裡.
出現了三次的得不著.
一次的不能得.
如果讀完二,三節的話.
這段經文告訴我們一件很恐怖的事.
就是當我們發現.
有一些從自我出發的私欲.
引誘一開始的時候.
原來就會越走越邪惡.
這裡,剛才我讀的第二節.
上半節,你可以分兩段.
兩個並排的句子就是.
你們貪戀得不著,就殺人.
你們嫉妒不能得手,就起爭執和衝突.
原來貪戀就是你自己沒有的東西.
你很想用一些圖謀方法去奪取它.
甚至說殺害.
然後,他說嫉妒不能得手.
嫉妒就是我們看到別人比我們好.
別人擁有的東西比我們好的東西.
我們又無法得到的時候.
接著就會越演越烈.
得不著並沒有帶來我們停一停.
想一想,究竟我們是否需要呢?.
反而是令我們越來越進逼.
越來越用不同的方法和手段.
以致去想爭取這東西.
甚至是殺害,引起爭戰.
目的就是滿足我們自己的慾望.
大家覺得殺害這個字很誇張.
但其實耶穌在《馬太福音》和《十誡》說到.
不可殺人的言論時.
當我們向弟兄動怒,或罵弟兄時.
就好像殺害人一樣.
其實這個殺害可以作為比喻.
為了我們要佔上風.
為了我們要得到我們的東西.
而詆毀別人的人格.
這也可以是一個殺害的比喻.
所以你會發覺與世俗為友.

$^{241}$還有一個很特別的重點.
就是我們會存在一個貪婪的心.
是不停貪婪的心.
「你們得不著,是因為你們不求.
你們求也得不著,是因為你們妄求.
要浪費你們的言諾中」.
其實這些貪戀和嫉妒的人.
是很少祈禱的.
因為他們自以為是.
因為他們滿心驕傲.
他們平時遇到問題時.
並不會求問神.
他們會依靠自己.
沒有意識去專注依靠上帝.
所以他們很少去祈禱.
所以雅各說「你們得不著,是因為你們不求.
你們根本沒有當上帝去幫你解決問題.
帶領你去解決問題的方法是一個方法」.
「就算你們求也得不到.
何時會求?就是發覺靠自己就解決不了.
那就為了上帝吧」.
危機來到時,覺得很擔心,很憂心.
才來找上帝「上帝,你幫我吧!求你引導我吧!」.
所以他們的祈禱都是妄求.
妄求的意思是「妄」字.
其實是含著罪惡的目的.
所以當他們求的時候.
可能都是因為懷著一些自己的私慾.
或者動機而被神拒絕.
這裡我們看到,其實與世俗為友的人.
他們和神的關係原來都被破壞了.
他們將神錯置在他們的生命裡.
神成為達到他們的目的的工具.
神並不是他們生命終極的目標.
所以在《成卷雅各書》裡.
我們之前一直看.
我們看到其實雅各在教會裡的紛爭,結黨的現象.
就是一群帶著私慾,貪戀,妒忌的心去出發.
以致他們要爭奪做師父.
沒有好好勒住他們的舌頭.

$^{281}$牢籠他們的會眾,去分化他們.
其實在紛爭的本質就是捍衛自己的信念.
其實有時我覺得戰爭.
我們有時也會為正義而戰.
但是雅各提醒我們.
這些紛爭來到的時候.
我們往往就是想怎樣去贏對手.
敵人有什麼不對.
對方有什麼錯.
以致我可以攻擊他,我可以贏.
但是雅各在提醒我們.
我們要留意一下.
這些紛爭的背後.
看的不是對方,看的是自己.
我裡面懷著什麼動機.
有沒有私慾去爭戰.
所以在一些爭拗,衝突,爭執出現的時候.
我們檢視的應該就是自己.
檢視自己的心思意念.
檢視自己究竟誰是我的盟友.
是世俗還是我們的上帝呢?.
然後我們去到第四召第六節.
這裡很清楚地說到.
其實原來與世俗為友的人.
就是與神為敵.
與神敵對.
你們這些淫亂的人.
豈不知與世俗為友.
就是與神為敵嗎?.
所以凡想要與世俗為友的.
就是與神為敵了.
你們想經上所說的是圖言的嗎?.
神所賜住在我們裡面的靈.
是戀愛至於嫉妒嗎?.
但他賜更多的恩典.
所以經上說.
神阻擋驕傲的人.
次因及謙卑的人.
雅各,耶穌的兄弟.
我想他作為教會的領袖.

$^{321}$面對紛爭,爭執,衝突不少.
但這位教會的領袖.
始終都是為主殉道.
為了繼續傳揚主的福音.
他被鞭打.
甚至被人在聖殿上摔下來.
其實一個完完全全不怕紛爭.
讓上帝做他的盟友.
站在上帝旁邊的人.
他說了一句話.
「你們這些淫亂的人」.
淫亂的意思不是在聖上的淫亂.
淫亂的意思是.
在舊約裡的不忠的以色列民.
被喻為淫婦.
當以色列國敗偶像.
就是對神不忠.
雅各很嚴厲地指出.
當這些基督徒.
宣稱自己和上帝納了這份盟約.
向人宣稱他們是屬於神的人.
卻卻是與世俗為友.
在羅文化裡.
朋友是一件嚴肅的事.
當你作為朋友的時候.
代表了他們的觀點一致.
而且對對方效忠.
我們要留意到.
其中經文裡兩次重覆提到.
與世俗為友就是與神為敵.
而在下半節裡說.
所以凡想要與世俗為友.
「想要」這個字.
並不是不想與世俗為友.
而是經過思考.
經過考慮而作出的選擇.
一群信徒.
因為在他們的生命裡.
對上帝的信越來越薄弱.
越來越不信.

$^{361}$越來越信自己.
和世界的價值觀認同.
於是就與世俗成為朋友.
當信徒以享樂成為他所貪戀的.
以世界的價值觀成為他自己的偶像.
他們就是與世俗為友.
在他們的生命裡.
在他們的行為裡.
就會表現出.
好像亞國書裡說.
嚴貧,愛婦,毀謗人.
有野心,嫉妒,享樂,紛爭.
甚至重傷人.
這些就是與世俗為友的表現.
這些就是對上帝不忠.
與神為敵的表現.
我們不難理解.
與世界做朋友.
就是與神為敵.
但第五節這段經文.
我們要稍為留意一下.
因為是一段比較難解的經文.
「你們想經上所說是圖然的嗎?.
神所賜住在我們裡面的靈.
是戀愛至於嫉妒魔」.
「神所賜住在我們裡面的靈.
是戀愛至於嫉妒魔」.
其實有不同的解法.
經文的翻譯.
在不同的版本都有不同的譯法.
這裡說的主詞.
就是在和合本這段經文裡.
所說的主題.
「嫉妒的人」是誰?.
就是住在我們裡面的靈.
但其實和修版有另一個譯法.
它也就是說.
「神愛安置在我們裡面的靈.
愛到嫉妒的地步」.
「神愛安置我們裡面的靈.

$^{401}$愛到嫉妒的地步」.
我比教會取向這個版本.
首先,人是按照神的形象做的.
所以我們裡面有靈.
這個靈是神所賜的.
而神是很愛我們.
祂愛我們愛到一個地步.
為我們死.
祂說祂嫉妒.
為什麼嫉妒?.
因為這個靈和世界做朋友.
神會嫉妒嗎?.
其實《川普記》20章第5節說.
「不可跪拜那些像.
也不可事奉他.
因為我耶和華你的神.
是忌邪的神」.
「忌邪」這個字就是嫉妒.
耶和華是忌邪的神.
在舊約和新約.
也有不同的經文表達.
上帝和子民的關係.
是新郎和新婦的關係.
可以想像當一段婚姻.
出現了第三者的時候.
難道被拋棄的那個.
他可以不嫉妒嗎?.
這個嫉妒其實含有兩種意思.
第一個含義就是.
神的身苦因為轉移跟隨世俗.
就與神為敵.
招致神的憤怒.
第二個含義就是.
不過神是好丈夫.
祂有著慈愛和憐憫.
祂期望祂的子民悔改.
所以祂會寬恕他.
這個嫉妒的含義.
到了下一節,第六節.
如果單看第四,第五節.

$^{441}$我們會覺得很絕望.
因為當我們在生活中.
選擇了有時跟隨世俗的價值的時候.
我們就是與神為敵.
但是第六節這樣說.
但他賜更多的恩典.
所以經上說.
神阻擋驕傲的人.
賜欽給謙卑的人.
經文第六節第一開始就是「但」.
這裡說了一件很重要的事.
就是神賜下恩典.
不過這裡很清楚地告訴我們.
他賜下恩典的對象是誰.
他賜下恩典的對象.
就是那些謙卑的人.
但他亦很清楚地說.
他不要賜恩典.
他要抵擋的是誰.
就是驕傲的人.
驕傲的人就是這群.
與世俗為友的人.
驕傲的人是不認同神的價值觀.
他可能聽很多道.
他可以懂得很多神學知識.
他甚至可以是教師.
不過他的生活中.
不表現到他是跟從上帝.
有些學者說.
驕傲不單是罪的清單上其中一項.
反而驕傲是罪的原由.
是令罪更加加有.
最慘的是驕傲的人.
不承認自己有問題.
一位諾貝爾文學獎的得主.
加布列賈西亞·馬奎斯.
在1982年得到諾貝爾文學獎.
他寫了一本書叫做.
《愛在瘟疫蔓延時》.
我們在祈禱會也曾經提及這本書.

$^{481}$他提及到.
有一天因為一塊肥皂.
而將一段婚姻瓦解.
究竟是怎樣的呢?.
原來太太本來負責.
家中浴室提供毛巾,肥皂,衛生紙等等.
有一天她忘記了更換肥皂.
她丈夫很誇張地說.
「我整個星期都不能用肥皂」.
於是她很生氣.
極力否認.
雖然她明知這是事實.
但她不肯承認.
因為丈夫用誇張的責備形式對待她.
於是接下來的七個月.
他們就分房睡.
吃飯時也沒話說.
馬奎斯說.
一塊肥皂為何能毀掉一段婚姻呢?.
因為雙方都不肯停下來.
想一想.
我們宣告停戰.
我們宣告做錯了.
我們宣告跟對方說對不起.
所以驕傲的人.
就是不承認自己問題的人.
但上帝要賜恩的.
就是給那些謙卑的人.
就是靠著他美善的恩賜.
讓人可以勝過罪.
蒙寬恕與神和好.
其實說到底.
與神為友的人.
就是一個謙卑的人.
有一顆謙卑的心.
如何是謙卑的人呢?.
如何可以與神為友呢?.
我們就要看第七至十節.
第七至十節說.
「故此你們要信服神.

$^{521}$不要抵擋魔鬼.
魔鬼就必離開你們逃跑了.
你們親近神.
神就必親近你們.
有罪的人要潔淨你們的手.
心懷異議的人要清潔你們的心.
你們要受苦,悲哀,哭泣.
將喜笑變為悲哀.
歡樂變作愁悶.
不要在主面前自卑.
主就必叫你們升高」.
這裡.
雅各用了一個日文字的遊戲.
第六節說「神阻擋驕傲人」.
「阻擋」這個字.
在希臘文中是Anti-Tesla.
是說甚麼意思呢?.
就是對立.
是在一個神服和制服的對立面.
即是不肯神服上帝.
是一個對立面的Anti-.
第七節說「你們要信服神」.
「信服」這個字.
這個字是指「之下」.
一個位置是對立.
一個位置是在制服,神服之下.
其實我們與神為友.
就是不站在神對立面.
而是降服在上帝的下面.
這個就是信服.
接著在七節到第十節.
總共用了十個命令式動詞.
雅各將耶穌對人的那個質度.
他告訴我們是因為他很愛我們.
他告訴我們我們不是沒有出路的.
當我們與世俗為友的時候.
當我們忘記了上帝的恩情的時候.
他告訴我們.
我們最重要做這十個命令式動詞.
第一就是要信服神,抵擋魔鬼.

$^{561}$第一個字.
即是那十個命令式詞語.
在第七節就彰顯了.
就是信服和抵擋.
信服神.
就是甘於讓自己居於上帝的權柄下.
服從神.
遵行祂的命令.
對祂表示尊重和服從.
信服就是這個意思.
然後我們要抵擋.
抵擋的對象是誰?.
是魔鬼.
我們知道魔鬼不容易對付.
不過不要緊.
聖經很清楚告訴我們.
魔鬼與上帝對立.
但是上帝會給力量我們勝過魔鬼.
好像馬太福音的主統文一樣.
「不叫我們遇見試探」.
意思是叫我們不要在試探中屈服.
不要在試探中屈服.
是給力量我們抵擋惡者的試探.
這裡告訴我們.
如果我們面對這些試探.
唯一的方法就是依靠主給我們的力量.
當我們依靠主的力量的時候.
魔鬼也要逃跑.
我們要注意魔鬼是欺控者.
他不能戰勝他.
所以魔鬼只要我們靠著神.
他就可以逃離.
很明顯.
當我們先信服神的時候.
我們才可以抵擋魔鬼.
大家記不記得迪士尼有一套卡通片.
叫做《獅子王》.
有一場戲講到主角鮮巴.
他面對著一大群土狼.
土狼看到他.

$^{601}$「嘩!這隻小獅子周圍沒有人支援」.
看著他的時候流口水.
鮮巴在後面.
用他很輕微的力量去咆哮.
引來的就是那群土狼的恥笑.
他們越來越進逼.
預備去吃這個大餐的時候.
面上戴著笑容.
面上在恥笑他.
但突然間鏡頭看到他們全班靜下來.
停下來.
然後眼睛突然間瞳孔放大.
然後就走了.
為甚麼?.
因為鮮巴的爸爸在他後面.
這位又大又壯.
不可小覷的獅子王就在他後面.
其實我們對抗撒旦.
抵抗他的就是這樣.
我們知道勝過撒旦的是主自己.
是祂釘十架為我們而死.
從死裡復活.
征服了罪和死亡.
以至到我們能夠在罪裡面.
在死亡裡面勝過.
有復活的生命.
撒旦是被主打倒的.
所以我們只要靠這位征服撒旦的主.
和他一起的時候.
我們就可以對抗撒旦.
聖經裡面說:親近神.
我們要親近神.
神就必親近你們.
有罪的人要潔淨你們的手.
心懷疑的人要清潔你們的心.
這裡給了我們一個很大的提醒.
亦給了我們很大的信心.
就是神告訴我們.
你放心吧.
你來親近我.

$^{641}$祂是很樂意很樂意親近我們.
就好像浪子在外面走完一轉.
回到祠父身邊一樣.
我們和神的關係.
上帝期望毫無保留地去結連.
上帝很想在我們身邊的每一樣東西.
祂都可以有份參與.
如果當有一些事情.
你不讓神去管.
要用自己的問題去解決的時候.
你就要問自己究竟為什麼.
我還記得我年輕的時候.
在談戀愛的事上.
我是特別不讓神去管.
結果就碰到焦頭爛額.
因為我很想靠自己的方法去處理.
去解決.
但結果就讓我和神的關係.
越來越遠離.
我知道有些弟兄姊妹.
會因為做錯事而不敢面對神.
沒有面見祂.
沒有面見人.
是的,我們有這個心是正常的.
以賽亞書59章1至2節說.
「耶和華的傍臂並非縮短,不能拯救.
耳朵並非發沉,不能聽見.
但你們的罪孽使你們與神隔絕.
你們的罪惡使他掩戀不聽你們」.
原來是罪孽使我們與神隔絕.
不過我們看到第8節下面說.
我們仍然可以親近上帝.
就是罪人要潔淨你們的手.
心懷異議的人要清潔你們的心.
下半節告訴你們.
如果我們有憂傷痛悔的心.
如果我們能夠清醒知道.
我們在心思,在行為上要跟從神的時候.
上帝絕對願意和我們打破這個阻隔.
這裡說罪人就是不遵從主命令的人.

$^{681}$他叫我們潔淨自己的手.
即是請你在行為上回到正軌.
請你聽到律法的時候.
你就跟著去行.
心懷異議的人.
就是對主不忠又愛世界又愛神的人.
他說你要清潔你的心.
你要留意不要被這個俗世沾污你.
生命才能夠有機會釋放和自由.
生命才能夠與神做復和的工程.
這裡說的是我們真的要加緊.
要好好整理自己的生命.
有時信主的人都會這樣.
我們會固執,我們會不知罪.
我們既想擁有世界.
也想有神的祝福.
在雅各看來.
這些都是我們典型罪人的本質.
但這裡告訴我們.
提醒我們.
我們一定要謹慎.
一定要聽上帝的話.
一定要去親近祂.
一定要俯伏祂.
在祂面前謙卑自己.
在1992年10月4日的《Our Daily Bread》.
有一篇靈修文章.
題目是《我們只是一》.
We were just one.
文章說到二次世界大戰的時候.
希特勒命令所有宗教團體.
都要聯合一起.
以便讓他們能夠容易控制.
當中有弟兄會的弟兄姊妹.
一半決定去服從.
加入這個團體.
但一半拒絕.
那些服從的人.
在往後的日子覺得輕鬆一點.
但那些拒絕的人.

$^{721}$在往後的日子就面對著.
嚴厲的迫害.
幾乎他們每一個家庭.
都有人死在集中營裡.
戰爭結束之後.
這些人中間充滿了怨恨和苦毒.
最後他們決定了.
不行,我們一定要去治癒這個情況.
於是這些群體裡少了一個領袖.
就在一個安靜的地方進行退休.
幾天以來.
他們很努力地親近上帝.
每個人都願意花時間去祈禱.
根據基督的命令.
他們去檢查自己的內心.
然後他們走在一起.
報導這件事的人.
叫做佛朗西斯謝弗.
他問他:你到最後發生了什麼事?.
然後那位被訪問的朋友這樣回答.
We were just one,我們只是一.
他們回答.
當他們向上帝承認自己內心的苦毒.
承認自己對對方的敵意.
並且向上帝宣告他們願意信服.
屈服在上帝的命令下.
聖靈叫他們合一的事就在他們中間發生.
愛就充滿了他們的心.
化解了他們的仇恨.
當愛在信徒當中運行的時候.
這些強烈的分歧.
就會因著這份認愛.
去展現了耶穌基督真正跟隨者的標誌.
當我們願意親近上帝.
當我們願意立志整理好生命.
行神的吩咐的時候.
我們是謙卑的.
接著第九節說.
你們要「仇苦,悲哀,哭泣.
將喜笑變作悲哀,歡樂變作愁悶」.

$^{761}$這裡有四個命令詞語.
「仇苦,悲哀,哭泣,變作」.
這幾個字其實好像別人在喪禮中所形容的表情一樣.
不過這裡很想表達的是.
面對災難,面對自己所犯的罪.
帶來的不幸,不幸的那種悲痛.
那種的仇苦,那種的哭泣.
這是一個表明自己悔改的心智.
是表明了我們願意轉換對罪的看法.
用神的眼光看待.
這才是真正悔改.
所以信徒啊.
我們真的要有超脫的眼光.
不讓世界裡一時的享樂,慾望.
那些喜樂就已經充滿滿足我們的心.
而是我們要知道.
原來我們能夠畏罪,悲哀.
才能獲得真真正正的自由.
最後第十節這樣說.
「勿要在主面前自卑.
主就叫你們升高」.
其實第十節和第六節是一個相呼應.
我們說到其實這個謙卑的背後.
全是恩典.
其實由我們信主開始.
豈不是已經是恩典嗎?.
全是靠著我們有什麼能力.
而令我們得著這份救恩嗎?.
我還記得年輕時有位阿姨跟我說.
「你記住找老公就要找個愛你」.
「多過你愛他的人」.
你們會不會這樣勸你們的子女?.
「找個他愛你,多過你愛他的人」.
因為你愛他多些,你會很辛苦.
但弟兄姊妹,其實我們很幸福.
因為我們有一個超級超級愛我們的上帝.
他的愛是無窮無盡.
我們很幸福.
因為我們得救是本福恩.
我們並不是因為我們的身份.

$^{801}$有什麼豐功偉績.
要求上帝用性命去救我們.
而是他實在憑著一個恩典去愛我們.
所以當我們不用驕傲去抬高自己.
當我們謙卑自己主動去信服上帝的時候.
神就將我們升高.
世界教我們要升高.
要得榮譽,地位,金錢,權勢.
才是真正的滿足.
要用自己的方法去追求,去圓滿人生.
不過,上帝這裡很清楚地告訴我們.
真正的升高是在救恩裡面有份.
是有上帝同行的人生.
是有上帝恩典不絕傾注的人生.
這個方法就是全然倚靠上帝的謙卑.
最後弟兄姊妹.
《雅各書2:22-23》這樣說.
可見信心是與他的行為並列.
而且信心因著行為才得成全.
這就應驗經上所說.
亞伯拉罕信神,這就算為他的意.
他又得稱為神的朋友.
在《雅各書2:23》告訴我們.
亞伯拉罕是神的朋友.
他並不是一開始就這麼信心.
他曾經用自己的方法去解決問題.
他也當上帝的說話是無道.
不過,因為他因著他的謙卑.
信服,悔改,行主的旨意.
上帝就將他升高.
弟兄姊妹,我們都很願意大家一起.
努力去做上帝的朋友.
他要求我們的是我們謙卑下來.
信服他,抵擋魔鬼.
能夠有一個悔罪的心.
願意去親近他.
他就必定親近我們.
讓我們一起祈禱.
因為在你的裡面.
我們去經歷與你為友的福氣.

$^{841}$上帝啊,你很愛我們.
你願意在我們的生命裡升高我們.
讓我們一生在你恩典的傾注下去同行.
大主,我們不得不承認.
在你面前,有時我們的信心的確是軟弱.
我們會選擇這個世界的價值.
以致我們不能夠全心全意地去跟從你.
天父啊,求你寬恕我們.
求你憐憫我們.
讓我們能夠在你面前謙卑下來.
有悔罪的心.
以致我們的人生裡得蒙你自己的帶領.
過一個精彩又行在你心意的人生.
我們恭敬仰望教託主.
奉主耶穌基督的聖名自祈求.
阿們.
\newpage



\section{哥林多前書 13:1-13}
\label{sec:bJlzZqbRugo}
\textbf{2021年8月29日崇拜直播｜何羅乃萱師母｜疫中有愛｜哥林多前書13\_1-13}
\newline
\newline
連結: \href{https://youtube.com/watch?v=bJlzZqbRugo}{\texttt{https://youtube.com/watch?v=bJlzZqbRugo}} ~~~~ 語音日期: 2021-08-29
\newline
\newline
\hyperref[sec:0pujbt8sHO8]{\small{< < < PREV SERMON < < <}}
~
\hyperref[sec:index_chronic]{\small{[返順時目]}}
~
\hyperref[sec:Xf1UScB62N8]{\small{> > > NEXT SERMON > > >}}
\newline
\newline
哥林多前書 13:1-13
\newline
\begin{longtable}{cl}
\hline
\hline
章節 & 經文 (和合本修訂版)\\
\hline
13:1 & \begin{tabularx}{0.7\textwidth}{X} 我若能說人間的方言，甚至天使的語言，卻沒有愛，我就成為鳴的鑼、響的鈸一般。 \end{tabularx} \\ \\ \relax
13:2 & \begin{tabularx}{0.7\textwidth}{X} 我若有先知講道的能力，也明白各樣的奧祕，各樣的知識，而且有齊備的信心，使我能夠移山，卻沒有愛，我就算不了甚麼。 \end{tabularx} \\ \\ \relax
13:3 & \begin{tabularx}{0.7\textwidth}{X} 我若將所有的財產救濟窮人，又犧牲自己的身體讓人誇讚，卻沒有愛，仍然對我無益。 \end{tabularx} \\ \\ \relax
13:4 & \begin{tabularx}{0.7\textwidth}{X} 愛是恆久忍耐；又有恩慈；愛是不嫉妒；愛是不自誇，不張狂， \end{tabularx} \\ \\ \relax
13:5 & \begin{tabularx}{0.7\textwidth}{X} 不做害羞的事，不求自己的益處，不輕易發怒，不計算人的惡， \end{tabularx} \\ \\ \relax
13:6 & \begin{tabularx}{0.7\textwidth}{X} 不喜歡不義，只喜歡真理； \end{tabularx} \\ \\ \relax
13:7 & \begin{tabularx}{0.7\textwidth}{X} 凡事包容，凡事相信，凡事盼望，凡事忍耐。 \end{tabularx} \\ \\ \relax
13:8 & \begin{tabularx}{0.7\textwidth}{X} 愛是永不止息。先知講道之能終必歸於無有；說方言之能終必停止；知識也終必歸於無有。 \end{tabularx} \\ \\ \relax
13:9 & \begin{tabularx}{0.7\textwidth}{X} 我們現在所知道的有限，先知所講的也有限， \end{tabularx} \\ \\ \relax
13:10 & \begin{tabularx}{0.7\textwidth}{X} 等那完全的來到，這有限的必消逝。 \end{tabularx} \\ \\ \relax
13:11 & \begin{tabularx}{0.7\textwidth}{X} 我作孩子的時候，說話像孩子，心思像孩子，意念像孩子；既長大成人，就把孩子的事丟棄了。 \end{tabularx} \\ \\ \relax
13:12 & \begin{tabularx}{0.7\textwidth}{X} 我們現在是對著鏡子觀看，模糊不清；到那時，就要面對面了。我如今所認識的有限，到那時就全認識，如同主認識我一樣。 \end{tabularx} \\ \\ \relax
13:13 & \begin{tabularx}{0.7\textwidth}{X} 如今常存的有信，有望，有愛這三樣，其中最大的是愛。 \end{tabularx} \\ \\
[1ex]
\hline
\hline
\end{longtable}
$^{1}$今日選讀的經訓是記在哥林多前書13章1至13節.
在新約聖經237頁.
一首愛的頌讚.
我若能說萬人的謊言.
並天使的話語.
卻沒有愛.
我就成了名的羅.
響的拔一般.
我若有先知講道之能.
也明白各樣的傲備.
各樣的知識.
而且有傳備的信.
叫我能夠誕.
卻沒有愛.
我就算不得什麼.
我若將所有的州際窮人.
又寫幾身叫人焚燒.
卻沒有愛.
仍然與我無益.
愛是恆久忍耐.
又有恩慈.
愛是不嫉妒.
愛是不自誇不張狂.
不做害羞的事.
不求自己的益處.
不輕易發牢.
不計算人的惡.
不喜歡不義.
只喜歡真理.
凡是包容.
凡是相信.
凡是盼望.
凡是忍耐.
愛是永不止息.
先知講道之能.
終必歸於無尤.
說方言之能.
終必停止.
知識也終必歸於無尤.
我們現在所知道的有限.

$^{41}$先知所講的也有限.
等那完全的來到.
這有限的必歸於無尤了.
我作孩子的時候.
話語像孩子.
心思像孩子.
意念像孩子.
既成了人.
就把孩子的事丟棄了.
模糊不清.
到那時就要面對面了.
我如今所知道的有限.
到那時就全知道.
如同主知道我一樣.
如今常傳的.
有信,有望,有愛.
這三樣其中最大的是愛.
今天我們很高興邀請到.
何娜軒師母.
在我們當中分享訊息.
何師母是家庭發展基金總幹事.
今天我們.
她在我們當中分享的訊息是.
「逆中有愛」.
是一個我們.
真是一個很貼切的題目.
在我們當中分享.
我們將時間交給師母.
師母:大英姐妹.
平安.
這次來到感受很不同.
不理她了.
我們是這樣的.
逆中甚麼都會發生.
我記得上次是否說報道會.
「張幸福大會」.
如果我沒有記錯.
上次應該叔叔還在.
我如何形容你們教會和我們.
是我的好鄰舍.

$^{81}$亦是我的同工教會.
不知道大家知不知道陳師母.
現在是我們的同工.
如果你看到我的Facebook.
變得漂亮了這麼多.
全部都是她的功勞.
我很感恩.
因為都是.
好像那天.
如果我沒有記錯.
是陳牧師跟她分享的時候.
她問我.
師母,你有甚麼需要?.
我說有啊.
我找到一個同工.
我的需要是這樣這樣.
她說不如.
想想我太太.
我嚇了一跳.
就是一餐飯.
我覺得上帝很真實.
今天跟大家分享的題目叫.
「逆中有愛」.
哥林多前書十三章.
一到十三節.
今天還有一樣特別的.
我也要介紹.
何牧師在這裡.
我的另一半何志迪牧師.
很難得.
多數是我跟他都有說.
這麼巧他今次有空.
「逆中有愛」.
疫情當中.
我們看到的周圍環境是怎樣.
給大家看一下一些PowerPoint.
不知道大家有沒有聽過.
家庭幸福指數.
香港排第七十八位.
還未出.

$^{121}$慢慢來.
我說就行了.
另外.
一個幸福指數.
家庭幸福指數只有六分.
不多.
不要緊,你聽我說.
不過最近有個叫開心指數.
有沒有聽過.
剛剛出的.
上個星期.
就說開心指數.
有六點四七分.
而開心指數最高的是甚麼呢?.
應該是五十到七十歲的人士.
他們最開心是甚麼?.
兩元坐車.
何牧師經常都說.
然後我們再看下去.
失業的壓力增加.
在最近一個調查就說.
家裡最常出現的三個字.
叫做困獸鬥.
有沒有聽過.
最可怕的是甚麼呢?.
小學三年級的學生.
會戒守.
自殘.
有人問我的調查是否準確.
我說是準確的.
因為你看調查.
是出處理自己.
另外.
下個星期有甚麼大事發生在家裡.
知不知道.
海學.
家長對子女升小一壓力.
很大.
不知道這裡有沒有升小一的家長.
大家關心一下他們.

$^{161}$在這個情況當中.
在這麼多的變化裡.
我經常都問.
我們有甚麼渣拿呢?.
我在《疫情》裡.
其實說了三篇.
一篇叫《疫中有信》.
一篇叫《疫中有亡》.
這一篇叫做《疫中有愛》.
因為我相信.
信亡愛.
其中最大的.
是甚麼?.
愛.
好像跟你這麼巧.
難道前書十三章.
剛才都讀過了.
我們可以.
飛了那三章.
說一些前言.
這裡的前言.
說有甚麼卻沒有愛.
這裡說的是萬人的謊言.
並天使的華語.
這裡可能說到很多.
屬靈的恩賜.
你又懂得方言.
你可能又懂得先知講道.
不過如果在家庭來說.
很多時候我們很希望.
我們的子女有天分.
有口才.
最近我才.
因為要照顧孫子上學.
剛才牧師問我最近忙甚麼.
照顧孫子上學.
我發覺每個人都在學很多東西.
起碼是一樣.
幼稚園而已.
四樣.

$^{201}$起碼四樣.
他說不學不行了.
你的時代過時了.
我的時代不用學.
我們即使有這些.
先知講道才能明白.
各樣的學秘.
各樣的知識.
我想對現在的家庭來說.
我們看很多書.
懂很多東西.
讀很多學位.
子女就要進神校.
大家知不知道甚麼是神校?.
那間呢.
你非入不可那間呢.
以前說有100間.
現在不知有多少間.
一定要進那些學校.
甚至說我約將所有的周濟窮人.
又寫幾身叫人焚燒.
我們很少有人會這樣.
寫幾身叫人焚燒.
不過如果用現代的說話.
來形容香港的家庭.
就是我用盡我的錢.
我犧牲.
原來香港很多家長.
為了子女搬屋.
我最近才聽.
我兒女本來在九龍.
進了一間學校.
要搬去南區了.
全家搬的.
真的會這樣.
都是為了孩子好.
我記得一個見證是這樣.
一個大學的教授.
他是普林斯頓大學的教授.
後來就進了一間神學院教書.

$^{241}$我就問他.
教授教授.
為何你會信耶穌.
你在一間這麼有名的大學.
還要進來神學院教書.
他說你有所不知.
他說我在國內出生.
從小我媽媽就跟我說四個字.
你要好好讀書.
我就把你送到香港.
他很好讀書.
然後媽媽就說.
我就把你送到香港那一家.
南校.
你猜猜.
能不能進去?.
可以.
然後媽媽每天跟他說.
你要好好讀書.
我把你送到美國普林斯頓大學.
結果.
進了.
拿到雙大學.
拿到院士榮譽.
生了兩個兒子.
將媽媽在香港接過來住.
有一天.
他放學.
他就見到他媽媽指著兩個乖孫.
跟他們說什麼?.
你要什麼?.
好好讀書.
像你爸爸一樣.
進普林斯頓大學.
這個教授聽到這句.
他整個人有點崩潰.
我說為什麼?.
他說我問我自己.
我的人生是不是只有一個目的.
就是怎樣?.

$^{281}$好好讀書.
進普林斯頓大學.
那怎麼辦?.
我的人生是不是就這樣循環?.
我的孫子.
於是有人叫他去查經班.
他去了查經班.
查考聖經.
後來讀信耶穌.
後來教神學.
我記得他的故事.
我們有了這些東西.
沒有愛也算不得什麼.
接著就說八個不.
不疾度.
原來疾度的意思.
很簡單.
為什麼你有.
我沒有.
為什麼你能進那間學校.
我的子女不能進.
不自誇.
自誇就是光環.
我記得.
大家都知道我做親子教育.
過去十多二十年.
去了很多學校.
有一次我去那間學校.
我就跟家長說.
我喜歡派禮物.
不過今天沒有.
我說給你一分鐘.
寫下你的子女長處.
寫得最快的那個有獎.
於是有一個巨心.
我寫到了.
拿出來.
你猜他怎麼讀.
第一長處是.
鋼琴第三級比賽.

$^{321}$彈巴哈第一名.
第二是羽毛球比賽.
初級組第一名.
第三是中文朗誦比賽.
讀了十五個第一出來.
我不知道你聽完感覺如何.
媽媽很自豪.
這些就是他的優點.
太太.
如果你的女兒.
有一天拿不到第一.
她還有沒有一些東西.
你欣賞她.
她想不到.
我沒有給她禮物.
我想全港唯一一個.
她沒有拿過禮物的.
我們回去想想.
你女兒如果沒有這些光環.
她還有沒有優點.
她的優點是不是光環.
不張狂.
自我中心.
自誇.
覺得自己最了不起.
人家呢.
真的有學生這樣說.
有媽媽走來跟我說.
她說我兒子去到一間新學校.
認識不到朋友.
我說他厲害嗎.
厲害.
為什麼他認識不到朋友.
她說兒子跟我說.
I am smart.
they are stupid.
不配跟他做朋友.
然後她就問我.
那怎麼辦.
有沒有辦法.

$^{361}$我說有.
你跟他說.
你聰明嘛.
不過你聰明層次低一點.
想不想高一點層次.
想.
你聰明.
人家是stupid.
教到人家都想高一點層次.
還有高一點層次.
是不是.
你聰明.
人家是stupid.
人家比你還厲害.
你懂得為他鼓掌.
這個就真的高層次了.
你回去跟兒子說.
你現在層次低一點.
不作害羞的事.
原來另一個解釋是.
不要無禮貌.
這個有少少失傳.
我記得有個校長有一天打給我.
他說羅小姐.
你可不可以來我們學校講禮貌.
我說你家學校應該挺好的.
他說你有所不知.
我說什麼事.
他說我們考收學生那天.
每個人都有禮貌.
媽媽拖著女兒進來.
然後對著我.
校長早安.
誰知道開學那天.
店裡走走店裡過.
當校長沒讀.
校長就說.
你可不可以來幫我說.
起碼叫他們都叫一聲.
校長早安.

$^{401}$不求自己的益處.
不要只是為了自己.
也要求別人的益處.
要顧慮到別人.
不輕易.
法律說的可能是EQ.
大家不知道有沒有覺得.
可能現在的情緒挺容易生氣.
有沒有哪位開車的弟兄.
姐妹都可以.
我不要性別歧視.
有沒有開車的.
沒有嗎?.
坐車也可以.
我現在發覺.
如果綠燈.
我不是馬上踩油門.
就已經被人打.
不止被人打.
還要車子轉到我旁邊.
指著我罵.
我唯有跟他打.
只慢了一秒.
不輕易.
法律不計算人的惡.
是不要儲存記憶.
不要將不好的記憶.
放在腦裡.
人會很不開心.
不喜歡.
不義.
原來有個解釋很好.
是我們不會幸災樂禍.
這八個不.
然後到愛的積極面.
就四個凡.
凡是包容.
原來這個字是有圖像性的.
就像一個屏幕.
遮蔽起來.

$^{441}$我們家應該凡事包容.
我很喜歡那個形容詞.
叫做埋舟.
有沒有玩過竹伊恩?.
有些有.
有些不知道我說什麼.
這個是年齡分界.
竹伊恩.
埋舟.
就是說你去到那裡.
沒有人罵你.
你安全.
You are safe.
家裡就應該是我們孩子的舟.
教會也是我們的舟.
凡事相信.
不是叫我們什麼都很輕易去相信.
但也不需要什麼都懷疑.
要拿到那個平衡.
凡事盼望.
就是在一切最不利的情況下.
都要盼望有好的結果.
我昨天看了一齣戲.
我和何牧師看了一齣戲.
很感動.
還沒上演.
出戲的故事.
是講一對盲人的基督徒夫婦.
他有一個開眼的女兒.
這個女兒每天和這對盲人的夫婦生活在一起.
但他們沒有喜樂.
沒有盼望.
我看完之後覺得很有趣.
他們活在黑暗中.
家裡沒有開燈.
但他們如何彼此相愛.
而這種愛如何感染到他們的孩子.
就算最不利的情況下.
我們依然會看到盼望.
依然有光.

$^{481}$凡是忍耐的另一種易發.
就是經歷一切困難.
那個愛依然堅強.
依然包容.
其實這些都是前言.
現在才講到今天那篇講章.
「我亦終有愛」是怎樣呢?.
有三個提醒.
很想和大家分享.
第一是有限.
圈住了.
第二是把孩子的事丟棄了.
第三是主知道我.
我們先講第一.
所知道的有限.
這裡說我們所知道的有限.
才知所講的也有限.
等那完全的來到.
這有限的.
必歸於無有了.
這裡的原文是說.
我們的恩賜都是暫時的.
有限英文很好.
叫Limited.
我們的公司就叫有限公司.
所以別人叫我開公司.
都一定要加個Limited.
否則就要負債纍纍.
這些都會停的.
當主再來.
那個完美的完全的來到.
這些都會停止的.
所知道的有限.
是我們所看的有限.
我們的眼界有限.
我們的認知有限.
我們的認識有限.
甚至我們對我們的孩子.
都是有限的.
同不同意嗎?.

$^{521}$今天你看到他這樣.
他將來就會變的.
所以我應用在我們的生活上.
在我們的親子關係.
或者對人的關係上.
我覺得有三點很重要.
第一.
不要看死.
不要憑著表面他對你的反應.
你就去定他是一個怎樣的人.
或者是一個怎樣的孩子.
我教小朋友寫作.
教了十幾年.
有很多令我印象深刻的孩子.
其中有一個是一個女孩.
我教四堂課.
我記得第一堂.
她進來課室.
就趴在那裡.
睡覺.
全程睡覺.
我叫她同學.
你叫什麼名字?.
沒有反應.
看都不看我.
第二堂進來.
又是這樣.
趴在那裡.
我有問她.
你知道我的事而不捨.
你叫什麼名字?.
沒有反應.
第三堂.
還問不問她?.
第一次問.
都問了.
只有四堂.
都是沒有反應.
一路.
全程好像睡覺.

$^{561}$最後一堂.
問不問?.
問.
最後一堂.
我叫她自由提.
任她們說話.
突然發覺.
她看著我.
告訴我她叫什麼名字.
我問她問題.
她有些很有創意的答案.
我抱她.
我說.
你講話真是很好.
她看著我.
真的嗎?.
下課.
我就說.
你媽媽會不會來接你?.
不會.
誰來接你?.
我鄰居.
一個媽媽.
我就等.
等到鄰居來.
我就說.
真好.
她今天終於開口.
最後一堂.
真的很開心.
她說.
真的嗎?.
我也很開心.
我問她媽媽呢?.
這位太太沉默了一會.
她說她媽媽.
其實很想女兒.
可以進你的班.
不過在上課.
開課幾天前.

$^{601}$她回到天家.
她媽媽去世了.
所以她不說話.
一直都不說話.
我記得那天.
我抱著她.
我說.
你真的很厲害.
她看著我.
真的嗎?.
這故事.
教會了我什麼呢?.
是我很久之前聽過一句話.
他說.
每一個悲傷的面孔.
背後都有一個故事.
每一個憂愁的面孔.
不會無緣無故的.
一定背後有一個故事.
我也很感激這位鄰舍.
她能夠將別人的孩子.
當作自己的孩子去照顧.
當然不要看得過重.
將孩子成為生命的重心.
特別是子女是青春期的父母.
我不用你舉手.
舉手顯露了你的年齡.
請你有失戀的準備.
特別是有孩子的媽媽.
她長大了.
她可能會拍拖.
請你放心.
我最近在網上寫.
最重要是隻眼開隻眼閉.
人們說.
麻煩你詳述什麼是隻眼開隻眼閉.
我真的放得開.
看得開.
不要看法.
不要覺得是這樣的.

$^{641}$不能改變.
坊間有兩種思維.
一種是成長型.
一種是固定型的思維.
固定型的思維是.
不要搞我.
挑戰不要搞我.
努力是沒有用的.
別人成功就能夠威嚇我.
但成長型的思維是不斷學習的.
我最近訪問了一位老年學的教授.
我覺得我要訪問他.
我問他.
對於坊間叫60歲的長者老人家.
奇英.
你怎麼看?.
我不知道大家喜不喜歡這個名字.
我就一般般.
雖然我已經可以被人叫奇英.
我問他.
這些人都失敗了.
各位哥哥姐姐聽著.
這些人都失敗了.
現在叫什麼?.
四零.
第一年齡.
我們完全要依靠家人幫助.
第一零.
第二零.
我們完全不用獨立自主.
自己養自己.
第三零.
一半靠別人幫忙.
一半自己都可以.
退休的狀態.
第四零.
要靠別人照顧自己.
我還記得那位姐姐.
我跟她差不多.
她問我.

$^{681}$你現在是全職工作嗎?.
我說是.
那即是第幾?.
第二.
全職第二.
我最近都看到.
有兩種智慧的能力.
一種叫流體智力.
即是分析,推理,解決問題.
年輕的時候我們可以.
但年紀越大.
就是晶體智力.
高層次一點.
是我們透過人生的閱歷.
運用已有的知識.
將我們的經驗.
將我們比那些壞.
或是比那些好.
傳遞出去.
那是屬於第三零.
或是第二,第三零的人士.
甚至我覺得第四零都有.
不代表他要別人照顧他.
是沒有的.
給大家看看另外一張.
我也很喜歡.
就是教我們怎樣看待孩子.
這幾句說話.
他說慢熟的孩子.
其實很謹慎,慢熱.
多話說其實很熱情.
固執的孩子很有主見.
衝動的孩子其實是action.
行動.
我很喜歡這句.
甚麼?.
分心的孩子其實是甚麼?.
很有想像力.
我小時候.
我記得小學四年級.

$^{721}$上數學課.
突然被老師叫我起來.
他說羅麗萱.
你讀這條數學題.
那我就讀了.
小明有原子彈.
一千四百二十粒.
小芬有原子彈.
一千二百四十粒.
共有多少粒?.
大家聽得出來嗎?.
問題在哪?.
聽得出來嗎?.
你不要閉上眼睛.
問題在哪?.
原子甚麼?.
彈啊!.
我還記得老師當時拍檯.
羅麗萱.
我就看到你望著窗又吠了.
這個世界怎會有一千二百四十粒原子彈?.
是原子甚麼?.
粒啊!.
我們那個年代是原子粒.
受孕記.
但事實證明.
分心的孩子.
是充滿想像力的.
最重要我們不要限制甚麼?.
不要限制甚麼?.
不要限制我們的上帝.
我很喜歡下面這幾句說話.
他說.
Don't limit God.
God can do whatever He wants.
上帝要做甚麼就做甚麼.
God can use whoever He wants.
上帝要使用誰就使用誰.
God can speak however He wants.
God can go wherever He wants.

$^{761}$God can move wherever He wants.
我很記得.
由於聽了太多關於SCN的故事.
即是特殊學習需要的孩子.
有一次.
我遇到一位基督徒的精神科醫生.
我就問他.
幼稚園又說甚麼過度活躍.
那些又不能救的.
又說自閉.
結果呢?.
他跟我說了他兒子的故事.
他說我是精神科醫生.
我兒子很小的時候.
我跟他做評估.
知道他是自閉的.
我說那你怎麼辦?.
我沒有當他是自閉症的孩子.
我只當他一個身份是甚麼?.
是我的兒子.
是天父所愛的兒子.
他說這個就是我認定的身份.
然後這位爸爸.
每天為兒子祈禱.
他跟我分享的時候.
他說你信不信神蹟?.
他說我兒子到中三那年.
所有的徵象都沒有了.
除了耶穌.
還有誰可以醫治他呢?.
Don't limit God.
不要限制我們的上帝.
我們的上帝是一個又真又活的上帝.
他是隨他幾意.
但他是愛我們的上帝.
第二.
我們要做的事是action.
將孩子的事吊起.
就是將孩子的自我中心.
表現和行為.

$^{801}$要扔掉.
Tinder的註釋說.
不是說時間過了.
孩子的個性.
我們的自我中心會過去.
他是用動詞的.
要決斷一點.
要專心一意.
要扔開這些東西.
要拿定主意.
結束孩子的事.
甚麼是孩子的事呢?.
舉幾個例子.
自我中心.
所以我曾經說過一篇道.
叫做孝敬父母.
我覺得要講回.
第二.
情緒失控.
生氣卻不要犯罪.
不可含怒度日落.
第三.
孩子的事是做著做著容易放棄.
忘記背後努力.
面前要向前飆竿直跑.
第四.
是等不到.
但這裡說.
等候他的.
都是有福的.
我不知道有沒有跟大家說過.
不過我都很喜歡講這個故事.
就是我為甚麼會做家庭發展基金.
為甚麼會走出來講家長教育呢?.
當然很多人說.
因為現在的家長不行.
我都在學.
一邊教一邊學.
其實是因為一個媽媽.
差不多二十年前.

$^{841}$我講一個講座.
叫做中一求生方程式.
那時候已經很難搞了.
即是升中要求生.
是教育局辦的講座.
我是講完.
我去到講.
講完之後有個媽媽走出來.
說:羅女士.
我想請教你.
我兒子小六升中一.
不過經常掛著打機.
我就用兩種方法教他.
我說哪兩種?.
你猜到的.
一是罵.
然後呢?.
打.
他說都沒有用.
可以怎樣?.
我們兒子來的.
我說鼓勵一下他.
甚麼叫鼓勵?.
我說走過去.
擁抱他一下.
拍一下他肩膀.
跟他說:加油啊.
努力啊.
他聽完.
低頭.
想了很久.
看著我.
很坦白告訴你.
我這麼大個女兒.
未被人擁抱過.
抱過.
親過.
鬥過.
我不知道怎樣擁抱我的兒子.
當時我唯一可以做的.

$^{881}$難道跟他說.
你把手伸出來.
像魂渡國 微屈.
我說你過來.
太太.
我擁抱他.
哭了.
因為有人排隊問話.
我記得他看著我.
這樣啊?.
擁抱啊?.
我說是啊.
你記住了.
離我遠去.
我心裡有句話.
很想跟他說.
我想說四個字.
耶穌愛你.
不可以說.
因為那不是教會的聚會.
於是我開始了我的事工.
他既然請我去.
很多人都知道我是師母.
我就說了.
直到有一天.
去到一間學校.
是我二十年前去的那間.
就是剛才說見到的奶奶那間.
我又去說.
我以為沒什麼人出來.
說完這樣的故事.
說完之後.
來了三個媽媽.
第一個.
我就像那個媽媽那樣.
被人擁抱過.
擁抱我行不行?.
行啊 擁抱你.
第二個.
我十八歲還未婚懷孕生子.

$^{921}$我很討厭我的兒子.
因為他奪去了我的將來.
但是.
我很想有人擁抱我.
又擁抱了.
第三個.
沒話說.
只是走過來.
擁抱擁抱.
現在不行了.
但是當時.
我能夠做的就是擁抱他們.
告訴他們.
有盼望的.
在這件事發生.
大概兩個月之後.
我去到一間.
叫做名校.
我就說了這個故事出來.
說了剛才的故事出來.
你知道有什麼發生嗎?.
當中有三個媽媽.
一個在左邊.
一個在中間.
一個在右邊.
一路走.
但是.
沒有人理會她.
不知道大家明白嗎?.
有人看到她哭的.
翹著手.
沒有人理會她.
我是唯一一次.
在講座當中.
我停下了.
我講不下去.
我說對不起.
我要停.
我說我問自己.
我做了十幾年家長教育工作.

$^{961}$所謂何事.
是不是來告訴你.
你的子女.
進了好學校.
然後找到一份好職業.
娶或者嫁到一個好人家.
然後又成立美滿的家庭.
然後又找到好學校.
如果.
我所宣講的信念.
只是自我的話.
我說我是一個.
我當時是用英文.
我是一個total failure.
我是一個完全失敗的.
家庭教育工作者.
我說我希望是什麼?.
愛淪陷如同此劇.
我盼望當我們.
不只是想念我們的事.
我們看到周圍的人的仇苦.
所以各位家長.
可不可以停下你們現在犧牲.
望下哪個家長.
需要給她一個擁抱.
你去擁抱她.
結果.
我看到在講座當中.
他們擁抱你.
他們給錢安慰你.
不再是漠不相關的人.
我們扔掉這些孩子的屎.
我們除了看到我們之外.
看到我們的淪陷.
你又望下周圍.
很難得今天回來.
弟兄姊妹望下.
不用用眼神打招呼.
我們不是漠不相關.
我們在主裡面.

$^{1001}$彼此的連繫.
第三點.
主知道我.
這裡說.
模糊不清.
到那時就要面對面了.
我們讀一次這一節.
一二三.
模糊不清.
到那時就要面對面了.
我如今所知道的有限.
到那時就全知道.
如同主知道我一樣.
最近《鏡》很出名.
我曾經想過這篇《道》.
就一定很多人聽.
不過如果你看回經文.
翻看以前的《鏡子》.
原來是模糊不清的.
這就是那時哥林多的鏡子.
看得清楚嗎?.
模糊的.
看不清楚.
但有一天.
是同鏡.
有一天會面對面見到耶穌.
但我很喜歡這幾個字.
叫做主知道我.
主其實認識我.
主知道我的處境.
主知道我的困難.
主不會丟下我們不利.
繼續下去這一節.
我們有讀.
這一兩節.
一二三.
親愛的弟兄啊.
我們現在是神的兒女.
將來如何還未顯明.
但我們知道主若顯現.

$^{1041}$我們必要像祂.
因為必得見祂的真體.
所以你們要完全.
像你們的天父完全一樣.
我們可以做什麼?.
我們要像耶穌.
我們的愛要愛得像耶穌.
那完全就像天父的完全一樣.
如何學習呢?.
下一章.
這裡說.
我遭遇患難受苦.
你的命令卻是我所喜愛的.
在疫情當中.
上帝的說話.
就是我們最大的安慰.
是我們的燈.
是我們的光.
我經常說.
2011年拿了兩本日記回來.
跟我說不如我們開始抄聖經.
於是.
如是者2011年.
現在是2021年.
以前覺得靈修.
看一段聖經.
很容易水過鴨背.
但當我們願意抄的時候.
真的很不同.
將上帝的話語.
能夠刻在寶上面.
和心板上面.
亦在今年7月.
基督出版社邀請我們.
寫了詩篇靈修.
連字40天.
這本靈修.
還有一些鋼筆字帖.
我鼓勵弟兄姊妹慢慢讀.
慢慢回答上帝的話.

$^{1081}$主知道上帝知道我們的處境.
知道我們有多困難.
多難受.
所以在疫情當中.
我其實有好幾個特別的見證.
可能你聽過我堂哥上半部分.
但現在和你說下半部分.
我堂哥在這裡.
可能你見過這張.
他在幾年前.
我走去和他說.
堂哥你是我羅家唯一.
一個未信耶穌.
亦是爸爸唯一的親人.
爸爸很想你信.
爸爸那時這樣這樣.
信了耶穌.
結果堂哥一聽.
就說他也要信.
就信了.
如果我沒有記錯.
我上次有和你說.
堂嫂怎樣信.
堂嫂很想信耶穌.
但當時她已經說不出話.
我們又不知道怎樣和她傳遞.
在她昏迷離流之際.
堂哥和我說.
堂嫂想信耶穌.
可不可以?.
我想一想.
於是怎樣?.
我和牧師在ICU探望她的時候.
我記得她就這樣祈禱.
主啊.
你知道堂嫂生前說過很想信你.
但她現在昏迷.
清醒指數得二.
我求你.
讓她知道其實她可以信靠你.

$^{1121}$說到那裡.
堂嫂睜開眼睛.
捉住我的手.
是不是很震撼?.
然後我就幫她做了一個決志祈禱.
她沒有跟著我說.
但我相信她的心是跟著我說.
今年五月.
堂哥的身體出現了很大變化.
心臟有事.
送了入ICU.
當時的肝酵數是三千多.
有沒有記錯?.
離流之際.
她的兒子打電話給我.
他說「嬸啊」.
「我爸爸進了醫院」.
「現在又不能探望」.
「那可以怎樣?」.
「醫生已經做盡了」.
「可以怎樣?」.
「可以怎樣?弟兄姊妹」.
「祈禱」.
「是嗎?」.
我也是這樣回答他.
他說「祈禱?」.
「是的,不過也沒有什麼用」.
我便閉嘴.
「主啊,可以怎樣祈禱?」.
我便想起昔日的諾惠娟姐妹.
她告訴我.
她在離流之際.
她見到耶穌.
耶穌給她力量.
我的祈禱就是這樣.
「主啊,如果你連恤堂哥」.
「你讓他在昏迷之中見到你」.
第二天.
我收到電話.
我侄子打來.

$^{1161}$「我爸爸醒了」.
「醒了」.
「是嗎?發生什麼事?」.
「他見到耶穌了」.
「他還見到耶穌」.
「帶著堂嫂來跟他說加油」.
兩個月後.
他換了一個很大的心臟起搏器的手術.
結果手術很成功.
第二天.
他錄了一段說話給我們.
播給大家聽,好嗎?.
OK了.
Peter,Shirley,你們好.
很開心.
多關心我.
多關心我的一家人.
我也是蒙著耶穌的保佑.
能夠很順利.
很多事情諸事順利.
可以這麼說.
當然除了已經開始的New Zealand Project.
對我來說.
最近一生的工作非常順利.
而且.
說後的情況非常好.
恢復得很快.
我已經有這種感覺.
主給我一種新的生命.
在這個關鍵的時刻.
對我和我家人的關心.
希望你們在事業健康方面.
能夠有所成就.
特別是聽到Shirley最近的書展.
有很大的成就.
我是為你高興.
I am proud of you.
從來沒有想過跟堂哥說這句話.
其實我上一次聽.
是我爸爸跟我說的.

$^{1201}$接著通常都是長輩才會說.
如果你問我.
生命裡面真的充滿感恩.
堂哥現在跟她說電話.
我姐姐最近跟她說電話.
她一看到耶穌就哭了.
心中都是感恩的淚.
見到她會跟我們打長途電話.
叫她上海的兄弟姊妹信耶穌.
她說的I am proud of you.
都想在這個時候跟大家分享一點.
這幾個月上帝很大的恩典.
我是一個做夢要做作家的人.
不過家裡一直都不贊成.
覺得做作家沒有出色.
我入行是因為後援雜誌的主編.
願意收我做徒弟.
我開始學寫.
不過整條路都不是很順.
有鼓勵我 有欣賞我.
亦有些叫激勵你.
你明不明.
激勵就是那些.
你都不知道寫什麼.
但是我覺得上帝很好.
在我十八歲那年.
當我告訴別人我很想學寫作.
沒有人肯收我做徒弟的時候.
我就向著天做了一個小女孩的祈禱.
親愛的耶穌.
我很喜歡寫東西.
你可不可以給我機會.
就是這樣.
上帝應允了這個小女孩的祈禱.
讓我在不同的關卡裡面.
有機會學寫作.
直到今年書展居然辦到了.
書展的主題叫做「心靈勵志」.
我從來沒有想過.
十大被選的作家是其中之一.

$^{1241}$其中還有三位都是基督徒.
而這本書是一本有信仰的書.
最後的章節是講印第.
入選了其中一類心靈勵志類的十本好書入圍.
入圍後我已經很開心.
但原來還有一個大獎.
我還記得.
我跟自己說.
出版社跟我說.
你不如等吧.
等看Zoom頒獎.
你有機會拿大獎.
我說不可能的.
上屆那個叫陳美玲.
我說一定不是我的了.
我跟何牧師.
你也知道我們會認識.
去了Staycation.
去了做我人生第一個迪士尼環保袋.
全心全意在製作.
完全沒有看電話.
到電話來.
何牧告訴我.
你拿了獎.
我哭得很厲害.
是感恩的淚.
是覺得順序.
真是聽一個小女孩的祈禱.
接著何牧在此說.
接到何牧師.
又拿了另外一個獎.
那個獎叫傑出校友獎.
我們神學院大學.
應該是九月二十五日.
頒獎.
當然我們還有一個獎.
那張相片沒有拿出來.
就是我們第二個乖孫出世.
所以今天回去要煮飯.
要照顧孫子.

$^{1281}$今天一本得獎的書.
還有一套靈獸.
我不知道價錢是否合適.
總之是特價.
機道書樓在外面.
歡迎大家可以購買.
因為有些弟兄姊妹說去不到書展.
另外再下去.
這個我也要說說.
人生有很多新開始.
知道這兩幅畫是誰畫的嗎?.
何牧師.
我從來不知道他會畫畫.
他覺得人生走到第三靈的階段.
是甚麼都可以學.
甚麼都可以做.
我正在等他開畫展.
好嗎?.
好,再下去.
我的專頁.
我們今晚也有雙劍合璧.
我們今晚會有一個直播節目.
會說挫敗.
有家庭發展基金的網頁.
師母來了也做了很多.
接著回應師.
我分享到這裡.
我們一起有個祈禱.
親愛的主人.
因為你是那個又真又活.
又愛我們的主.
我們不知道今天在甚麼光景裡.
天父大求你.
讓我們知道我們是蒙你所愛.
是你選擇的.
你知道我們內心的孤單無助.
求你對我們每一個人來說話.
你保守我們.
保守我們的家人.
與我們眾人同在.

$^{1321}$禱告,交託,仰望.
奉救主耶穌基督.
得勝的名字.
\newpage



\section{哥林多後書 12:1-10}
\label{sec:Xf1UScB62N8}
\textbf{2021年9月19日崇拜直播｜梁家麟牧師｜保羅不刷存在感｜哥林多後書12\_1-10}
\newline
\newline
連結: \href{https://youtube.com/watch?v=Xf1UScB62N8}{\texttt{https://youtube.com/watch?v=Xf1UScB62N8}} ~~~~ 語音日期: 2021-10-15
\newline
\newline
\hyperref[sec:bJlzZqbRugo]{\small{< < < PREV SERMON < < <}}
~
\hyperref[sec:index_chronic]{\small{[返順時目]}}
~
\hyperref[sec:aSTV2iY5HmY]{\small{> > > NEXT SERMON > > >}}
\newline
\newline
哥林多後書 12:1-10
\newline
\begin{longtable}{cl}
\hline
\hline
章節 & 經文 (和合本修訂版)\\
\hline
12:1 & \begin{tabularx}{0.7\textwidth}{X} 雖然自誇無益，我還是不得不誇。我現在要提到主的異象和啟示。 \end{tabularx} \\ \\ \relax
12:2 & \begin{tabularx}{0.7\textwidth}{X} 我認識一個在基督裡的人，他在十四年前被提到第三層天上去；或在身內，我不知道，或在身外，我也不知道，只有神知道。 \end{tabularx} \\ \\ \relax
12:3 & \begin{tabularx}{0.7\textwidth}{X} 我認識的這樣的一個人—或在身內，或在身外，我都不知道，只有神知道— \end{tabularx} \\ \\ \relax
12:4 & \begin{tabularx}{0.7\textwidth}{X} 他被提到樂園裡，聽見隱祕的言語，是人不可說的。 \end{tabularx} \\ \\ \relax
12:5 & \begin{tabularx}{0.7\textwidth}{X} 為這人，我要誇口；但是為我自己，除了我的軟弱以外，我並不誇口。 \end{tabularx} \\ \\ \relax
12:6 & \begin{tabularx}{0.7\textwidth}{X} 就是我願意誇口也不算狂，因為我會說實話；只是我絕口不談，恐怕有人把我看得太高了，過於他在我身上所看見所聽見的； \end{tabularx} \\ \\ \relax
12:7 & \begin{tabularx}{0.7\textwidth}{X} 又恐怕我因所得的啟示太高深，就過於高抬自己，所以有一根刺加在我身上，就是撒但的差役來折磨我，免得我過於高抬自己。 \end{tabularx} \\ \\ \relax
12:8 & \begin{tabularx}{0.7\textwidth}{X} 為了這事，我曾三次求主使這根刺離開我。 \end{tabularx} \\ \\ \relax
12:9 & \begin{tabularx}{0.7\textwidth}{X} 他對我說：「我的恩典是夠你用的，因為我的能力是在人的軟弱上顯得完全。」所以，我更喜歡誇耀自己的軟弱，好使基督的能力覆庇我。 \end{tabularx} \\ \\ \relax
12:10 & \begin{tabularx}{0.7\textwidth}{X} 為基督的緣故，我以軟弱、凌辱、艱難、迫害、困苦為可喜樂的事；因為我甚麼時候軟弱，甚麼時候就剛強了。 \end{tabularx} \\ \\
[1ex]
\hline
\hline
\end{longtable}
$^{1}$今日的經訓是在哥林多後書第十二章一至十節.
記錄在新約聖經第252頁.
以及在程序表上.
請聽我讀出.
「我自誇固然無益,但我是不得已的,如今我要說道主的顯現和啟示..
我認得一個在基督裡的人,他前十四年被提到第三層天上去,.
或在身內我不知道,或在身外我也不知道,只有神知道..
我認得這人,或在身內或在身外我都不知道,只有神知道..
他被提到樂園裡,聽見隱秘的言語,是人不可說的.為這人我要誇口,.
但是為我自己,除了我的軟弱以外,我並不誇口..
我就是願意誇口也不算狂,因為我必說實話..
只是我禁止不說,恐怕有人把我看高了,過於他在我身上所看見所聽見的,.
又恐怕我因所得的啟示甚大,就過於自高,.
所以有一根刺加在我肉體上,就是撒旦的差役要攻擊我,免得我過於自高..
為這事我三次求過主,叫這次離開我..
他對我說,我的恩惠夠你用的,因為我的能力是在人的軟弱上顯得完全,.
所以我更喜歡誇自己的軟弱,好祈福基督的能力能浮避我..
我為基督的緣故,就以軟弱,凌辱,急難,迫不及待,困苦為可喜樂的..
因我什麼時候軟弱,什麼時候就剛強了..
恭請梁牧師..
鄧小平:各位節目早上好,我們繼續看哥林多候書..
最近我們看了很多場不同的比賽,最多的當然是體育的競賽,.
從東奧到殘奧到現在舉行的全運會,可能有人看了香港小姐決賽,.
還有各種的足球比賽,烹飪比賽,歌唱比賽,.
人總是喜歡做這個那個的比較,證明誰比誰優秀,.
證明誰才是世界第一,史上最強..
所有的比賽都要設定一些評分的標準,要比什麼,要怎麼比..
有些東西是無從比較的,或者是主觀性太強,或者可比性不大..
例如我女兒早上七點多,她一邊做運動一邊看電視,.
有一個外國烹飪比賽,竟然找一個營養師去為美食評分..
你告訴她,我不知道她中間有沒有營養師,營養師能夠說什麼?.
不用猜也知道,脂肪太高,纖維太低,卡路里太高,是不是?.
你覺不覺得這是超級荒誕的事情,追求健康飲食,.
根本不要外出吃飯,更不要說去吃什麼星級美食,對不對?.
有人一直吃著鵝肝魚子醬魚豬片皮鴨,.
然後一面就說膽固醇太高了,太不健康了,那你想怎樣?.
你告訴她,她不僅不健康,她邪惡,是不是?.
很奇怪,很奇怪..
保羅和超級使徒都在比較,在比賽..
前面那段,保羅詳細的說,在事奉歷程裡面,.

$^{41}$他的受苦的清單,說明他為信仰所付的代價遠比那些超級使徒偉大..
我比他們多受勞苦,多下監牢,受鞭打是過重的,誣死是屢次有的..
這個比較的目的,是要證明他在23節所說的話,.
他們是基督的僕人嗎?我講一句狂話,我更危..
保羅要做的說明是,他們有的,我也有,他們沒有的,我也有..
接下來,保羅說的故事,我們今天讀的這一段,才是最叩人心眼的..
保羅首先再作一個聲明,12章一節,.
我自誇固然無益,但我是不得已的,如今我要說道主的顯現和啟示..
第一個動詞是「現在式」,其實是在說「我現在仍在自誇中」,.
他不是自誇一次,他之前已經在自誇,現在再說,我就是自誇的..
不過現在是到了自誇的高峰,保羅說,他知道自誇是沒有價值的,.
不過他覺得他被迫要說這些,目的是要向他的弟兄姊妹說明,.
當有人挑戰保羅的權威,保羅的地位的時候,他們可以有言可答..
值得注意的是,保羅說,我現在說的話是和主的顯現和啟示有關的事情..
而他所用的主的顯現和啟示都是用眾數的,.
即是說,在保羅的記憶中,顯現和啟示都是不止一次的,.
不只是限於他待會要說的三層天的經歷,.
不過由於保羅不好意思自誇過度,所以他只說一件,.
所以當保羅想到上帝在他身上的顯現啟示的時候,.
他不只說一次,他有很多次.第二到第六節,他這樣說,.
我認得一個在基督裡面的人,他在前十四年被提到第三層天上去,.
或在身內我不知道,或在身外我也不知道,只有上帝知道..
我認得這人,或在身內或在身外我也不知道,只有上帝知道..
他被提到樂園裡聽見隱蔽的言語是人不可說的,.
為這人我要誇口,但是為我自己,除了我的軟弱之外,我並不誇口..
我就是願意誇口也不算狂,因為我必須說實話,.
只是我禁止不說,恐怕有人把我看高了,.
過於他在我身上所看見,所聽見的..
這是一個很厲害的經歷,保羅被提到三層天,第三層天的經歷..
先說什麼叫三層天,猶太人一直相信天有很多層,.
舊約常常有,你們應該讀過,天和天上的天,.
你記不記得這個說法,證明天有很多種,很多層,不是一層天,而是很多層天..
在這麼多層的天裡,不同層有不同的活物存在,.
通常比較低層的就是比較低層次的靈界事物的存在,.
什麼執政,掌權,或者撒旦叫做空中屬靈氣的惡魔,.
都是空中的,都是在上面的,不過這些是低層的,.
最高層應該就是上帝所在的地方,對不對?.
三層天理論上應該是最高層的天,.
雖然猶太人對天有多少層,從三層到十層都有人說的,.
最多說是七層天,不過我照這段經文看,.

$^{81}$保羅應該相信天只有三層,第三層就是最高的那層,.
就是看到上帝的那層,沒理由上帝在低層,對不對?.
所以如果他在第三層看到上帝的話,第三層應該就是最高的,.
所以保羅應該相信三層天,保羅沒有直接提自己的名字,.
不過靖主後面他講到一根刺的故事,.
很明顯這個主角一定是他自己,他不提自己的名字的原因,.
或者好像那些拉比傳統,那些Rabbinic Tradition,.
通常都將自己講自己的故事的時候,推到第三身去講法,.
借個人,不講自己,也都很可能是一個謙卑的講法,.
保羅不好意思,如果自誇的話,就用我來講,.
通常就將他推到第三身去講,.
保羅講這個事情發生在14年前,.
就是他事奉生涯的早中期,.
如果他是在主後55到56年間寫成哥林多後書的,.
那件事應該就是數前15年,就是主後42年,.
這個時候保羅應該還在安提阿,.
保羅標明那個具體的時間,.
證明他是第一次將這個故事告訴哥林多教會,.
之前14年他都沒有講過,他沒有提過,.
他不僅沒有跟哥林多教會講過,.
他也沒有跟任何人講過,.
不講的原因,或者是上帝不讓保羅講,.
或者是保羅認為講這件事對信徒是沒有益處的,.
所以盡量不要講,.
保羅的想法其實跟我們今天真的很不同,.
我們今天連吃什麼早餐都馬上貼上社交網站,.
跟人分享,.
你想一想,想像一下,.
如果我昨天真的有被上帝提升到天堂的經歷,.
我會第一時間WhatsApp全世界,.
告訴我的親戚朋友,.
然後這個星期日就算不是我講,.
我都會跟羅莫說交換交換,.
讓我上來先,.
就算我本來是講哥林多學說的,.
不講哥林多學說,沒什麼要講,.
就講我的故事,對不對?.
一定要馬上講出來,.
因為這個經歷太吸睛了,.
太驚心動魄了,太難得了,.

$^{121}$如果你想一下,我今天崇拜我的講題是,.
梁家倫在天堂上去下來,.
或者天堂真相大揭秘,.
我保證我們接下來的視頻會過十萬甚至到百萬的,.
你信不信?.
或者可以想一下,.
給一百塊錢才可以看,.
十萬人看就可以收一千萬,.
多過癮,馬上可以買上一層,.
買多一層,我們不是三層樓,.
只買多一層樓,.
是不是?.
這麼刺激的題目,.
怎麼可以藏起來呢?.
保羅這樣的超凡經歷,.
為什麼不馬上到處講呢?.
不講是很古怪的,.
保羅被提到三層天,.
上面跟神秘者,.
其實那個很可能是上帝,.
這個經歷有什麼特別的價值呢?.
其實我們也不知道有什麼價值,.
我們只知道的就是,.
對保羅自己來說,.
特殊的超自然經歷,.
他不止一次,.
他其實有很多次,.
最有名的當然是信主的經歷,.
在大馬士革路上遇見主,.
然後弄盲了一雙眼睛,.
你記不記得?.
這件事一定驚心動魄,.
另外,他見到馬其頓人的異象,.
而改變他事奉的路向,.
你記不記得?.
那個也是一個很重要的事件,.
記載在少林寺十六章的故事,.
並且保羅在事奉的歷程裡面,.
常常都被主介入,.
或者攔阻,.

$^{161}$或者引渡,.
所以他不是一次經歷主,.
常常都有主的介入,.
不准他去哪個地方,.
要他去哪個地方,.
另外,還有他跟西拉坐牢的時候,.
有天使把整個監獄翻轉,.
你記不記得?.
然後玉卒想自殺,.
然後保羅向他傳福音,.
然後他全家得救,.
你記不記得這個故事?.
這些都是在我們看來,.
都是極難得的,.
很厲害的數人經歷,.
你也沒有這麼多這樣的經歷說,.
保羅一生,.
其實這些故事是不少的,.
如果真的要比較三長天經歷,.
跟剛才我說的那些經歷,.
其實我不覺得在保羅的生命,.
和事奉上,.
三長天經歷有多少的影響,.
大馬士革路上的經歷,.
使他信主,.
改變他一生,.
馬其頓人的意象,.
改變了整個事奉,.
他傳教的路向,.
那麼三長天經歷,.
有什麼大影響呢?.
對他的事奉,.
對他的宣教,.
有方便了嗎?.
不知道有沒有,.
如果有的話,.
保羅應該到處說這件事,.
他都沒有說,.
所以有什麼影響我也不知道,.
不過,.

$^{201}$不說對保羅,.
對事奉有什麼影響,.
從吸睛和呃Like的角度,.
大馬士革的事件,.
一定不及三長天的經歷,.
兩者根本不是同一個檔次,.
大馬士革路上的經歷,.
即是什麼經歷?.
即是耶穌闖入你的生命裡面,.
和你相遇,.
並且改變你一生,.
我們中間大部分弟兄姊妹,.
都有這個經歷,.
你可能沒有那麼戲劇性,.
但耶穌也曾經踩過你的牆,.
動過你的,.
我們這樣信主,.
這個經歷我們都有的,.
三長天的經歷就完全不同,.
不是耶穌進入保羅的生活世界,.
而是保羅進入上帝的世界,.
那個本來仍然有血肉之軀的禁地,.
如今保羅可以進入,.
你說這個多大的殊榮,.
保羅立即變成一個特權階級,.
我覺得保羅這個經歷,.
比聖經你找到的經歷,.
比見到耶穌變相,.
和摩西和以利亞在山上對話,.
或者約翰在拔河島上見到的末日異象,.
都不及保羅這個經歷,.
耶穌登山變相,.
或末日異象都是發生在人間,.
只是時間上是之後的事,.
但三長天不是在人間,.
樂園是人死後才能去到,.
你明白我說什麼嗎?.
你行山時遇到周潤發,.
你和他一起拍照,.
周潤發光顧你的小店,.

$^{241}$吃了一碗牛雜粉,.
又一起合影,.
很厲害,.
我的故事很簡單,.
周潤發約我去他家,.
在他家的露台喝了杯咖啡,.
沒什麼好說的,.
我們說些沒什麼營養的事,.
就是這樣,.
你告訴他,.
和不同檔次,.
你碰到周潤發,.
我去了他家,.
和他家談無聊事,.
就是這樣,.
很明顯,.
就算三長天的經歷沒有影響,.
沒有改變保羅的生命和事奉,.
都大大提高了他的身份和地位,.
保羅成為了上帝特邀的天堂嘉賓,.
還沒死就可以去天堂參觀,.
你數給我聽有誰,.
有誰,.
真的不誇張說,.
保羅是贏了舊約的先知,.
新約的使徒,.
可以先上天堂,.
上帝有話跟他說,.
還沒那麼奇怪,.
如果上帝沒話跟他說,.
都叫他去天堂走一趟,.
你就更加覺得他特殊,.
你明不明白,.
這樣都可以被他當一當,.
這是什麼特權,.
不過保羅是自己認定,.
宗教經驗純粹是私人的事,.
是上帝跟他之間的事,.
不要說這個特殊的三長天經,.
他之前說過,.

$^{281}$他講方言的經驗,.
他都盡量不跟別人說,.
因為這是私人的,.
僅僅是對他適用,.
不會對其他人帶來任何好處,.
後果,.
保羅自己說過,.
他講方言比別人還多,.
不過他不喜歡說,.
他只喜歡跟別人說教導性的話,.
分享個人的宗教經驗,.
如果沒帶來造就別人的後果,.
就是自誇的行徑,.
目的是抬高自己的屬靈地位,.
這是不可取的,.
保羅常常說,.
他在上帝面前可以發瘋,.
可以癲狂,.
但在弟兄姊妹面前,.
他就會約束自己,.
會緊守,.
他不會放肆的,.
不能夠藉特殊的宗教經驗去抬高自己,.
甚至不用特殊經驗,.
去證明自己的仕途身份,.
和事奉的地位,.
保羅只是說,.
他願意用他格外的勞苦,.
付出的代價,.
和事奉產生的客觀果效,.
來證明自己的屬靈地位,.
這個做法,.
確實不是我們現代人能夠理解,.
會這樣做的,.
我的老師,.
Gordon Fee,.
在寫到這一段的時候,.
他也特別說到,.
他說這個真的和我們現在的想法完全不同,.
我們通常是用人的屬靈經驗,.

$^{321}$來成為我們去測驗,.
人的事奉,.
人的靈性的一個重要的證據,.
是不是?.
哇,.
上帝都和我說話,.
證明我的屬靈地位,.
上帝都和我說話,.
證明我所說的話,.
每句都是金句,.
對不對?.
不過,.
保羅不是這樣想的,.
真的不是這樣想的,.
當然,.
我們仍然可以追問的,.
為什麼不能夠分享呢?.
最少分享都可能會產生正面後果的,.
例如,.
激勵弟妹做信仰的追求,.
提供他們的屬靈胃口,.
屬靈品味,.
屬靈眼光,.
這樣行不行?.
當你聽到保羅分享他的三層天經歷的時候,.
我們就立刻更加行切祈禱,.
希望有一天,.
我們都可以攀登相同的境界,.
這樣好不好?.
保羅說不好,.
就是因為可能有這個後果,.
所以保羅才一定不要說,.
一定不要說,.
我講這兩點,.
接著這兩點,.
第一,.
特殊的神秘經驗純粹是上帝的賜予,.
跟人的追求,.
人的努力是完全沒有關係的,.
不管你怎樣重讀重讀哥林多後書,.

$^{361}$或者找回天主教的聖伊納爵,.
大德蘭的靈修作品,.
左看右看,.
或者你每個禮拜都去摩星嶺,.
去長洲,.
那些天主教的退修中心,.
去做神操,.
都不會使得你更容易獲得這些神秘經驗,.
屬靈經驗如果不是人的作為,.
不是人的心理操作,.
而是上帝真實臨在的話,.
根本不是人能夠操控的,.
甚至不可以人預定,.
或者預期的,.
那些什麼說什麼,.
我們先攀登到屬靈的黑夜,.
那個Dark Night of Soul,.
然後等候被接上天,.
這些說法,.
我可以說,.
通通全部都是胡說八道,.
不知所謂,.
在舊的歷史裡面,.
一些所謂靈修大師,.
常常訂定了一些靈魂通往上帝之路,.
The Soul's Journey Unto God,.
就是這樣的說法,.
他們相信他們可以規劃出一條通天之路,.
這些說法全部都是狂妄荒謬的,.
我不是天主教徒,.
我不會將某些人的說法,.
他叫聖人,.
封了他的姓,.
所以他說的話對我沒有權威,.
我只是問有沒有聖經根據,.
我要說是沒有,.
是零,.
如果你找到的話,.
你駁我吧,.
他們只是例如,.

$^{401}$基路柏有三對翅膀,.
等於有六個階梯,.
你告訴我,.
有什麼關係,.
老作跟神棍說的話,.
不是差太多,.
你憑什麼說,.
六個翅膀就等於六個階梯,.
我看來看去都看不到有什麼關係,.
有什麼關係,.
是不是,.
如果不是上帝親自來接引我們,.
沒有人可以腳高一對,.
可以離上帝近一公分,.
沒有人可以定定一條爬升上帝的路,.
沒有,.
是沒有,.
保羅得著三層天的經歷之前,.
他有沒有去過那些退休中心做神修,.
有沒有經過什麼心靈黑夜的階段,.
我不覺得他有,.
他至少沒有這麼說,.
用一個很禪的說法,.
給你的,.
怎樣都給你,.
不給你的,.
再神神化化都沒有用,.
這是第一點我要說的,.
第二點,.
如果保羅的分享吸引了很多信徒,.
做所謂更深度的屬靈追求,.
每個星期都去摩星嶺或長洲,.
這不僅是徒勞無功,.
浪費時間心力,.
更加會將我們整個基督信仰的中心,.
將它錯置,.
將基督的人生目標本末倒置,.
我們不需要,.
不需要,.
更深更高的屬靈經歷,.

$^{441}$單單認識耶穌基督,.
並他釘十字架就已經夠了,.
我們都不應該在人間追求什麼神秘的經驗,.
活出一個怎樣超凡脫俗的人生,.
遵守聖經裡面,.
上帝顯明的旨意,.
努力踐行福音使命,.
這才是最符合上帝心意的人生,.
我們在人間是不需要上三層天,.
連賞望都不需要,.
腳踏實地在地上生活和事奉就足夠,.
我用回耶穌會規的說法,.
任何人力圖攀登某個屬靈的地位,.
是超過上帝命令他應該達到的程度,.
這是一個驕傲,.
上帝造我們成為人,.
我們就老老實實做人,.
不要妄想成為超人,.
神人,.
再說,.
總有一天,.
我們會被提上天,.
在空中和耶穌基督相遇,.
那時候開始,.
我們就有無限的時間,.
可以和上主永遠在天堂一起,.
我們在人間只有數十寒暑,.
在天堂的時間卻是永恆,.
今天為什麼我們不把握好,.
只有一世夫妻的時間,.
好好和你老婆度過每一分每一秒呢?.
純粹開玩笑地說,.
我不敢高攀三層天,.
如果有人告訴我,.
例如陳牧師告訴我,.
可以攀登二層天,.
只需要跟著這個屬靈大師,.
這個Guru,.
操練三百小時,.
我都會計算,.

$^{481}$三百個小時等於,.
如果我和老婆喝茶,.
每次一個小時,.
花了我和老婆三百次喝茶的時間,.
我就寧願選擇吃蝦餃,.
不去操練,.
就正如去旅行,.
我們去三個星期意大利,.
第二個星期就有人說,.
很想吃雲吞麵,.
你不知道搞錯,.
三個星期在意大利,.
你回到香港,.
一天三餐雲吞麵都可以,.
有什麼理由在意大利,.
想吃雲吞麵,.
你有沒有搞錯,.
在意大利專注吃好意大利菜,.
這麼美味,.
回香港才吃雲吞麵,.
所以我們將來在天堂,.
我們自然會在那裡享受永生,.
永生,.
今日真是好好把握珍惜,.
和家人在一起的時間,.
很多姐妹,.
是否會煮飯今晚家人吃,.
兒女回來吃飯,.
媳婦回來吃飯,.
做好今晚的晚餐,.
這個更重要,.
我們看看這個經歷,.
按著字面的理解,.
保羅在經歷裡面,.
被提到三層天,.
然後他又提到樂園這個字,.
在那裡聽到神秘的言語,.
我們很簡單的解釋,.
不和大家糾纏,.
我相信三層天就是樂園,.

$^{521}$樂園就是三層天,.
不是兩個經驗,.
是一個經驗,.
所以是一次的旅程,.
不是兩次的旅程,.
值得注意的是,.
在保羅被提的技術裡面,.
他有兩次提到,.
或在身內,.
或在身外,.
保羅不確定,.
他自己在被提的時候,.
是否連著他的身體,.
還是靈魂出竅,.
靈魂被提,.
他不確定,.
保羅在這裡說兩次,.
他不確定,.
我不知道,.
我覺得一方面是他如實的說,.
他不知道,.
但另一方面,.
都真的顯示他的謙卑,.
通常我們說這些,.
這麼超勁的經歷的時候,.
總是要證明我們有實力,.
證明我們厲害,.
但保羅竟然在這裡唱反調,.
不斷強調他的ignorance,.
我無知,.
我不知道,.
你可以看到他真的不想show off,.
不想演戲,.
另一方面,.
實際上不確定,.
也可能是真的,.
因為這些出神的經驗,.
aesthetic experience,.
根本不是人的理性,.
能夠充分去分析,.

$^{561}$去理解的,.
不過我們也可以從這裡看到,.
保羅並沒有否定肉體的價值,.
宣稱肉體與靈性無關,.
如果保羅真的靈慾二分,.
相信肉體是卑賤的,.
他根本就不需要說,.
他就污毆了,.
我們會把肉體加上去,.
第三層天的可能,.
根本不可能,.
原則上不可能,.
但保羅說我不sure,.
可不可能,.
我不記得,.
我不確定,.
我的身體是否連在一起上去,.
最少他沒有排除這個可能性,.
我們的身體也可以是神性,.
可以在上面,.
保羅在天堂聽到有跟他說話的人,.
不過他沒有說明是誰,.
我們不確定是誰跟他說話,.
有可能是天使,.
但我相信更有可能是上帝,.
因為跟著說的一根刺的經歷,.
就是上帝加給他的,.
保羅只是不喜歡直接提上帝的名字,.
這也是猶太人的傳統,.
覺得上帝太神性,.
所以不要直接說上帝的名字,.
保羅把整件事件理解為顯現和啟示,.
顯現可以明白,.
他見到上帝,.
不過啟示是什麼呢?.
我們其實不是很清楚,.
保羅說這個啟示會使他自告十二章七節,.
不過他沒有說啟示的內容是什麼,.
那個不知名的跟他說話,.
說了什麼他都沒有提到,.

$^{601}$無論如何三層天的經歷很簡單,.
保羅只是想哥林多信徒知道他有這個經歷,.
在面對質疑保羅的人的時候,.
有言可答,.
但保羅不想說更多的東西,.
在五到六節,.
保羅說,.
為這人我要誇口,.
但是為我自己,.
除了我的軟弱之外,.
我並不誇口,.
我就是願意誇口也不算狂,.
因為我必須說實話,.
只是我禁止不說,.
恐怕有人把我看高了,.
過於他在我身上所看見,.
所聽見的,.
這些說法是很惱教的,.
保羅說,.
被提到三層天的經歷,.
確實是獨特超凡的,.
他說他就算用這個來誇口,.
去自誇都絕不過分,.
不過保羅是不喜歡自誇,.
只是喜歡說他的軟弱,.
不過無奈的是,.
保羅知道一旦說這個故事,.
就算他如何低調都沒有用,.
他都會立即顯示,.
不用添油加醋,.
只要實話實說,.
眼球都會凸出來,.
所以保羅只能選擇不說,.
這是他十四年來,.
沒有說過這個故事的原因,.
然後我們由三層天的經歷,.
去到一斤刺的經歷,.
我們其實不確定,.
一斤刺的經歷,.
和三層天是否相同的事件,.

$^{641}$相同的經歷,.
不過按照上下文的編排,.
視為同一個事件是合理些,.
第七到十節,.
我因所得的啟示甚大,.
過於自高,.
所以有一斤刺加在我的肉體上,.
就是撒旦的差役要攻擊我,.
免得我過於自高,.
為這事,.
我三次求主,.
叫這刺離開我,.
他對我說,.
我的恩典夠你用的,.
因為我的能力在人的軟弱上顯得完全,.
所以我更喜歡誇自己的軟弱,.
好讓基督的能力浮避我,.
我為基督的緣故,.
就以軟弱凌弱急難,.
逼百困苦為可喜樂,.
因為我什麼時候軟弱,.
什麼時候就剛強了,.
保羅真的承認,.
他的啟示是甚大,.
這是一個很厲害的肅靈經歷,.
並且這個經歷會使他自高,.
使他自命不凡,.
並且保羅說,.
不是他自己覺得自命不凡這麼簡單,.
是上帝都覺得,.
給了他這個經歷,.
會使他肅毛肅逆,.
耀武揚威,.
上帝都怕他自高,.
所以上帝要壓低他,.
要把一根刺的經歷給他,.
所以這個經歷確實太厲害,.
不過當說到一根刺的經歷時,.
保羅就用第一人稱去描述,.
說好事保羅不喜歡說自己的名字,.

$^{681}$提醒自己,.
說壞事保羅就非常踴躍地說自己,.
不過一根刺的經歷技術其實都是語言不暢,.
甚至和前面三層天的經歷比較更加簡陋,.
他只是說到上帝擔心他,.
因三層天的經歷而自高,.
就把一根刺加在他身上,.
加刺給他的是誰呢?.
很明顯就算提到撒旦的名字都和撒旦無關,.
撒旦怎會怕人自高,.
撒旦會愁死你不自高,.
你越驕傲,.
撒旦就越開心,.
你越驕傲就越離開上帝,.
撒旦怎會喜歡你謙卑,.
撒旦都幾純良,.
當然是上帝怕他自高,.
所以經歷其實是上帝給他的,.
他曾經想過是撒旦的差役整我,.
不過他轉念一想就知道那不關撒旦事,.
是上帝懲治他,.
這個轉念是我們常常有的經驗,.
最初碰到負面事物時,.
我們馬上想到撒旦整我,.
不過當我靜下來,.
我知道我這一生是上帝掌管,.
沒有上帝容許,.
沒有任何屬靈屬世的事可以整我,.
所以一定是上帝有個原因,.
就算是撒旦整我,.
也是上帝容許他整我,.
目的是有一個正面的後果在我身上,.
我要學到這個功課,.
保羅說到將好像是撒旦的經歷,.
將它轉化成為上帝給他的功課,.
給他的教訓,.
這個很重要,.
這個一根刺的經歷是什麼呢?.
是沒有人知道的,.
我們只可以確定的是,.

$^{721}$這個經歷或這個事件,.
發生在保羅身上,.
對他構成不是一次性,.
而是長期的困擾,.
但又未至致命,.
他死了,.
只是長期困擾,.
這是什麼經歷呢?.
很多人只能按字面解釋,.
因為字面上有根刺,.
是加在他身上,.
所以我們只能按字面解釋,.
身上指著身體,.
可能是疾病,.
長期病患,.
有人說保羅可能有耳鳴,.
或是有偏頭痛,.
這些都是在身上發生的事,.
有人說保羅可能有長期疾病,.
並且是很困擾性的疾病,.
如果偏頭痛也很辛苦,.
可能是很辛苦,.
這是有可能的,.
另外一根刺在身上,.
身體,.
你未必能理解是一個身體,.
你可能理解是教會,.
基督的身體,.
可能人們的理解是一根刺在身上,.
可能是上帝安排了一個很麻煩的梁家麟,.
在保羅的團隊,.
常常為保羅的事奉帶來很多麻煩,.
你可以想像,.
一個事奉團隊如果有一個很困擾的人,.
一個永遠說負面說話,.
永遠做拆毀性工作的人,.
是否很困擾?.
如果你相信基督的身體,.
身的意思是基督的身體,.
在保羅的事奉團隊有一個麻煩友,.

$^{761}$保羅常常求上帝,.
請你拯救,拯救,拯救,拯救,.
不過上帝不應運,.
都有可能,.
可以是他身上的毛病,疾病,.
也可以是他團隊的問題,.
還有可能其他的,.
因為我們只是猜測,.
沒有人具體知道,.
到底一根刺在身上是指什麼意思,.
無論如何,.
保羅總是很不喜歡,.
很想拔走,.
並且一直拔不走,.
因為這是一件長期存在,.
長期困擾的事,.
保羅將這個經歷確定了,.
是上帝對他的管教,.
所以十二章九到十節說,.
所以我更喜歡誇自己軟弱,.
很叫基督的能力奉庇我,.
我為基督的緣故,.
以軟弱,凌辱,急難,迫不困苦,.
為可喜樂,.
因為我什麼時候軟弱,.
就什麼時候剛強,.
保羅用正面的態度,.
去看他身上的患難,.
他再列舉一張患難的清單,.
就好像他經歷的軟弱,.
經歷的凌辱,.
經歷的急難,.
經歷的迫不困苦,.
他沒有說內容,.
因為前面的苦難清單都說了,.
總之他再重覆,.
我遭遇這些負面的經歷,.
我都視之為我的榮譽,.
他以他的受苦為榮,.
為什麼?.

$^{801}$因為上帝看得起我,.
讓他繼續在我身上,.
做那個拆毀,毀壞,傾覆,建立災籍的工作,.
我以此為榮,.
並且以此為樂,.
這個真是一個很不得了的保羅,.
我們做一點總結思考,.
很多人都喜歡擦存在感,.
駁上鏡,駁出位,.
爭取曝光,爭取發現的機會,.
在不同的場合找一些自我宣傳的空間,.
娛樂圈的人固然要這樣做,.
女性純粹靠身材,.
靠穿的衣服來出位,.
有時候賣一點肉都要,.
總之吸不到精,.
記者不肯拍你,.
男性只能靠言語,靠行為,.
常常都說寧願被人罵,.
都不要被人忘記,.
被人遺忘是最可怕的事,.
如果沒辦法以正面形象示人,.
負面形象都好過沒形象,.
無論如何都要有出鏡,.
好過沒出鏡,.
公眾人物,政治人物都不例外,.
常常都要沒事找事,.
找個藉口露面,.
在媒體面前罵人,.
或站在道德高地提出要求的批評,.
例如政府為何一次過派5000元消費券這麼奇怪,.
最低限度一個月派一次5000元,.
並且不應該只成年人才可以派,.
BB都可以派,.
甚至媽咪一證實就立即派,.
女性懷孕要進補,.
所以立即派,.
你叫吧,不用代價,.
總之叫下就可以出鏡,.
可以上下位,.

$^{841}$我們生活世界見盡這些事,.
說些無營養價值的話,.
唯一目的是證明自己存在,.
換來別人的注目,.
我看過一篇討論存在感的文章,.
一個內心富足的人從來不會去拆存在感,.
他知道自己的存在,.
亦知道自己活著的意義,.
有時他會寂寞,.
不過當他被聚焦時,.
他不以默起,.
亦不會失態,.
更加不會沒有素質地停不下來,.
停不下來,.
一個內心強大的人,.
他知道自己要甚麼,.
亦清楚自己想表達甚麼,.
並且他知道存在感不是拆出來的,.
不需要拼命去展示才能炫耀自己,.
找個安靜的角落,.
自己aware到自己的存在,.
這是最重要的,.
不要在人群中迷失自己,.
這樣的存在是可怕的,.
你自己也不知道你的存在感是甚麼意思,.
只有弱者才會狼狽地不停去拆存在感,.
保羅是一個非常不喜歡顯露自己的人,.
盡可能低調,.
畢竟他知道自己的身份,.
他知定自己是耶穌基督的奴僕,.
做奴隸是不應該出鏡的,.
你可以做所有事情,.
但在鏡頭前顯露的只是你那個主人,.
怎會把自己的頭露出來,.
他是沒有位置的,.
他知道,.
保羅很少說自己,.
只是在哥林多教會,.
因為他的仕途身份,.
是他的事奉被質疑,.

$^{881}$所以他才說多些有關他的仕途身份,.
他的權威的事情,.
這不是他慣常的做法,.
保羅刻意不喜歡說自己,.
他只想人們見到耶穌基督,.
這樣的風範,.
是我們不容易學到的,.
和大家說些個人事,.
可以展示一些關鍵點,.
這幾年我幾乎每天都和太太行山,.
我們通常主要在博弗林水塘行山,.
沿途見到很多野生動物,.
野豬和蜥蜴最多,.
有很多蜥蜴有他們自己固定的地盤,.
固定的活動時間和地點,.
他們最喜歡伏在山間旁的石頭邊,.
一坐就站著,.
總之就站著,.
兩三個小時以上,.
因為我上山時見到牠,.
下山時也見到牠,.
總之就站著,固定地不動,.
所以我們都當牠們是我們的朋友,.
去到這個位置找一條蜥蜴,.
和牠打個招呼,.
見一見牠,.
然後再行山,.
有幾個固定的位置,.
經常見到,.
幾乎每天見到,.
有一次我太太見到一條蜥蜴時,.
看著溪水,.
他們最喜歡在水邊,.
可能水邊有些蚊蠟,.
有些昆蟲可以給牠吃,.
我不知道是不是,.
動也不動,.
好像在思考人生,.
在思考叢生,.
不是人生,.

$^{921}$我太太說,.
這些蜥蜴大概是最沒有存在感的動物,.
完全沒有存在感的動物,.
牠們的保護色,.
牠們完全隱藏在自然環境裡面,.
你不刻意去留意,.
刻意去看,.
你根本不會發現牠們在那裡,.
而牠們也不喜歡,.
人或其他動物注意到牠們,.
牠們是想隱藏自己,.
所以不被注意對牠們來說是最自在,.
我想說的是,.
我們也未必要找聖經講得那麼盡,.
有時候測一下存在感是沒有問題的,.
不用太過正義,.
因為聖經的做法的確和這個世界的遊戲規則不同,.
和我們的想法的距離太遠,.
不過我真的很希望我們做了基督徒,.
我們有一份因著信仰而來的淡定感,.
就算有一天,.
我這個糟老頭,.
穿著短褲在街上走,.
沒有人看到我,.
我都會抬頭望天,.
我深信我主看著我,.
我在祂那裡被珍惜,.
被看重,.
就好像我經常在山上看到一塊塊的樹葉,.
我上上都很感動,.
這些樹葉不會被人看到,.
沒有人留意它的存在,.
它只是就這樣去展示它一塊葉的綠,.
去訴說上帝的創作,.
不需要被人注意,.
它做的是活出它自己,.
我求主幫助我們,.
可以確認保羅所得的一根刺的啟示,.
我的恩典是夠你用的,.
因為我的能力在軟弱的人身上,.

$^{961}$才顯得完全,.
這一種自足感,.
這種淡定感,.
這種朝夕,.
我覺得是我們大家都求上帝給我們擁有,.
是這個淡定自足的感覺,.
可以讓我們坦然,.
安然,.
寬然去面對我們日後高高低低的日子,.
上帝恩代我們.
\newpage



\section{約翰壹書 3:21-24}
\label{sec:aSTV2iY5HmY}
\textbf{2021年10月10日崇拜直播｜羅志雄牧師｜坦然無懼｜約翰壹書3\_21-24}
\newline
\newline
連結: \href{https://youtube.com/watch?v=aSTV2iY5HmY}{\texttt{https://youtube.com/watch?v=aSTV2iY5HmY}} ~~~~ 語音日期: 2021-10-17
\newline
\newline
\hyperref[sec:Xf1UScB62N8]{\small{< < < PREV SERMON < < <}}
~
\hyperref[sec:index_chronic]{\small{[返順時目]}}
~
\hyperref[sec:rfN73DzwrE4]{\small{> > > NEXT SERMON > > >}}
\newline
\newline
約翰壹書 3:21-24
\newline
\begin{longtable}{cl}
\hline
\hline
章節 & 經文 (和合本修訂版)\\
\hline
3:21 & \begin{tabularx}{0.7\textwidth}{X} 親愛的，我們的心若不責備我們，在神面前就可以坦然無懼了。 \end{tabularx} \\ \\ \relax
3:22 & \begin{tabularx}{0.7\textwidth}{X} 我們一切所求的，就從他得著，因為我們遵守他的命令，行他所喜悅的事。 \end{tabularx} \\ \\ \relax
3:23 & \begin{tabularx}{0.7\textwidth}{X} 神的命令就是：我們要信他兒子耶穌基督的名，並且照他所賜給我們的命令彼此相愛。 \end{tabularx} \\ \\ \relax
3:24 & \begin{tabularx}{0.7\textwidth}{X} 遵守神命令的，住在神裡面，而神也住在他裡面。從這一點，我們知道神住在我們裡面，這是由於他所賜給我們的聖靈。 \end{tabularx} \\ \\
[1ex]
\hline
\hline
\end{longtable}
$^{1}$今日經訓是約翰一書三章二十一至二十四節.
在大家手上的聖經新約三百三十九至三百四十頁.
約翰一書三章二十一至二十四節.
親愛的弟兄.
我們的心若不責備我們.
就可以向神坦然無懼了.
並且我們一切所求的.
就從他得著.
因為我們遵守他的命令.
行他所喜悅的事.
神的命令就是叫我們信他兒子耶穌基督的名.
且照他所賜給我們的命令彼此相愛.
遵守神命令的.
就住在神裡面.
神也住在他裡面.
我們所以知道神住在我們裡面.
是因他所賜給我們的聖靈.
弟兄姊妹平安.
平安.
我們一起來領受神的一道.
在之前我們一起有個禱告.
主耶穌我們向你禱告.
向你呼求.
我們是軟弱.
我們是無能力.
大主你仍然用你的話語.
兼顧我們每一個人的心.
叫我們坦然無懼來到你的面前.
願主你的話語.
就植入在我們每一個人的心坎裡.
由有聖靈來帶領我們.
以致我們能夠在這個世代.
繼續去見證你.
將上帝你的恩典愛與人分享.
願上帝你得著當得的榮耀.
禱告奉主耶穌基督的名.
阿們.
題目是坦然無懼.
我相信應該與膽量有關.
我去看聖經的時候.

$^{41}$其實有個人物.
聖經也說到與他有膽量有關.
這個就是約書亞.
約書亞是摩西的助手.
他協助摩西帶領以色列人離開埃及.
去到神應許他留乃與物之地.
當以色列人去到摩亞平原的時候.
準備去渡過約旦河的時候.
摩西離世.
神選立約書亞成為摩西的接班人.
在這個時候.
約書亞要帶領以色列人走從未走過的路.
他如何帶領一群勃逆以色列民進入去應許之地呢?.
神親自對約書亞應許.
說會與他同在.
神不會撇棄他.
不會丟棄他.
神又多次吩咐.
叫他剛強壯膽.
不懼怕不驚惶.
只要進行神的律例.
謹守誡命.
就能夠順利地帶領以色列人進入去迦南應許之地.
到最後約書亞確實是.
完成了神交付給他艱鉅的使命.
甚至他以自己的榜樣來呼籲百姓.
要事奉神.
要跟隨神.
就必須剛強壯膽.
不懼怕不驚惶.
今天.
執事剛才跟我們讀這段經文.
《約翰一書》.
他的作者.
被稱為耶穌所愛的門徒約翰.
他是貼身地跟隨耶穌.
他和他的兄弟.
曾經請求耶穌.
說:耶穌將來你得國的時候.
我們兩兄弟,兩個人.

$^{81}$在你一左一右.
在你的身邊.
約翰甚至見到.
其他人來奉耶穌的命.
來趕鬼的時候.
但又不跟隨他們.
約翰亦向耶穌投訴.
有一次.
耶穌和門徒.
來到耶路撒冷.
聖經記載他們行經了撒瑪利亞人的城鎮.
但當地人見到耶穌仍然繼續向著耶路撒冷的方向走.
所以沒有接待他們.
約翰又請求耶穌.
吩咐降火來消滅他們.
在這些事件中.
我們看到約翰的性格.
相當火爆.
亦很剛烈.
他跟隨主.
但仍然不明白耶穌基督的心腸.
他甚至缺乏愛心.
又不友善地對待其他人.
但一直看下去.
約翰經歷了相當年日的磨練之後.
他的生命有很大的改變.
他更親自體會到.
耶穌基督犧牲的愛.
那個真正的意義.
約翰是第一世紀的使徒.
當時教會面對極大的艱難和逼迫.
外面有政治壓力.
有異教文化衝擊.
在教會內有很多紛爭.
而當時約翰所牧養的教會.
其中亦有一些假教師.
他們看準了這個機會.
來動搖信徒的信心.
他們有些主張.
就說得救和永生的確具.

$^{121}$和我們道德的生活.
又或者關懷.
相愛.
這些是沒有關係的.
甚至我們日常生活的行為.
可以放縱.
根本犯罪.
和我們的信仰是沒有關係的.
這些假教師甚至宣示自己.
有特別的屬靈經歷.
又有特殊的宗教經驗.
和一些知識.
他比起其他人更有能力.
知道自己如何去到神那裡.
針對這些情況.
約翰就寫了《約翰一書》.
來鼓勵和勸勉信徒.
要堅定持守所領受的道.
過分別為性的生活.
敬虔.
坦然無懼地回到神的面前.
說回我們今天所讀的經文.
在第21至24節.
我以親愛的.
有些日本人稱為「親愛的弟兄」.
我來做這個小段.
因為我最喜歡約翰這樣稱呼.
親愛的.
因為這是很有親切感.
也很有親和力.
親愛的是約翰常用的字.
如果大家去看《約翰一書》.
有六次這麼多.
當約翰說到關於愛的段落時.
他就更常用這個稱呼.
所以.
約翰來稱呼這班門徒時.
這班門徒當時受到良心責備.
又或者受到很多恐懼的攻擊.
當約翰所牧養的信徒.

$^{161}$聽到約翰這樣跟他們說.
親愛的.
我相信對他們來說.
是非常大的安慰.
但又有很多的識經家.
一些學者.
他們就以19至24節為一個段落.
其中一個因素.
就是在第19節說.
從此就知道.
這就是開始.
我們再看回到第24節的時候.
那裡是我們所以知道.
有譯本會加上.
我們所以從此知道.
其實這一段是首尾呼應.
所以我今天會從第19節開始探討.
坦然無懼這個課題.
從此我們知道.
我們是屬真理的.
約翰在這裡強調.
遵行彼此相愛的命令的人.
他們是蒙受恩惠.
是得到最大的心靈安穩.
恩典的一群人.
愛,信服,真理.
在約翰書信裡.
是來回重複地講述的三個主題.
對於約翰來說這是非常重要.
所以他每次講述的時候.
其實他是以新的向度.
更加細緻地深入地討論.
愛,信服,真理.
如果我們可以看回第二章的時候.
約翰首先以光明與黑暗.
來講及彼此相愛這個主題.
而到了第三章.
即是今天我們所講的.
約翰仍然以愛來作主題.
但約翰論述的時候.

$^{201}$就是以神的兒女來作為開始.
來鼓勵信徒彼此相愛.
約翰鼓勵信徒在神的面前坦然無懼.
因為我們是屬於神.
是神的兒女.
經文第十九節至二十一節.
從此就知道我們是屬真理的.
並且我們的心在神面前可以安穩.
我們的心若不責備我們.
神比我們心大.
一切事沒有不知道的.
親愛的弟兄.
我們的心若不責備我們.
在神面前就可以坦然無懼.
約翰鼓勵信徒在神的面前坦然無懼.
我們留意到在第十九節.
從此就知道我們是屬於真理.
從此是甚麼意思呢?.
在呂鎮忠譯本.
從此這個字譯為「在這一點」.
究竟這一點是由哪裡開始呢?.
很多色經學家認為是指第十八節.
第十八節所說的是甚麼呢?.
這裡是說到關於相愛.
我們相愛不要只在言語和舌頭上.
總要在行為和誠實上.
讓姐妹們知道愛是要有行動.
是真實可見的.
以這個行為來表達出來.
第十八節提到行為和誠實.
誠實在原來的文字.
就是在真理上.
讓姐妹們更加明顯.
神是愛.
住在神裡面的就要相愛.
在我們日常的生活裡.
我們以行動.
真誠地實踐相愛.
這就是我們信徒的生活.
這就是在生活中實踐真理.

$^{241}$之後約翰繼續在第二十節說.
我們的心若責備我們.
神比我們的心大.
一切事沒有不知道.
約翰其實是擔心信徒.
被一群假教師.
敵基督者誘惑.
這些人否定約翰所教導.
甚至是否定耶穌基督.
過往他自己親自的教導.
是叫門徒切實地彼此相愛.
約翰在這裡鼓勵.
約翰說我們不要受到心所責備.
究竟心是指什麼呢?.
是指情感,意願,知識,良心和思想.
神是會饒恕屬祂的人.
以恩典盛載屬祂的兒女.
無論我們的心如何責備我們.
神都是以慈悲對待我們.
在二十三節中.
這是一個總結.
這裡約翰說.
神的命令就是叫我們.
信祂而指耶穌基督的名.
召祂賜給我們的命令彼此相愛.
請弟兄姊妹留意這一節的經文.
有一個字是很重要的.
這是一個信字.
一個信耶穌基督的人.
他相信耶穌基督的話.
就會實踐基督吩咐的愛.
弟兄姊妹,你信耶穌基督的.
你在人的面前展現愛心.
是因為你結出了因信心而生的果子.
這就是愛.
所以一個真正對神的心.
會是怎樣的心呢?.
我們再看二十一節.
這裡好像是重覆來講.
我再讀一次.

$^{281}$我們的心若責備我們.
神比我們的心大.
一切事沒有不知道.
親愛的弟兄.
我們的心若不責備我們.
在神面前就可以坦然無懼.
如果我們再看.
約翰不斷講心的時候.
我們會發覺.
二十節和二十一節.
好像有些矛盾.
二十節其實是一個負面的敘述.
我們的心若不責備我們.
但到了二十一節的時候.
約翰就說我們的心若不責備我們.
這是一個正面的陳述.
約翰在這裡重覆提醒我們.
真正屬於神的.
是神的兒女.
在神的面前.
我們的情感.
我們的信靠.
必然會坦然無懼.
很多時候我們未必是這樣.
我們會有良心的責備.
又會有一些遭遇.
來攻擊我們.
但約翰說這裡.
因為神的心比我們大.
在神裡面.
祂有寬恕的心.
有慈悲憐憫的心.
祂能夠勝在我們的良心.
所以當良心責備我們的時候.
我們在神的面前.
已經得到神為我們解決.
我們的心就可以安穩.
弟兄姊妹.
曾幾何時.
大家有沒有感受到.

$^{321}$撒旦的攻擊.
撒旦的厲害.
撒旦的權勢.
和一些邪惡.
他們不斷地來攻擊我們.
我和大家其實都是面對這些.
有時候真的會很軟弱.
甚至會陷入犯罪試探的裡面.
約翰就在這裡來鼓勵我們.
勸勉我們要仰望神.
要信靠耶穌.
當我們這樣做的時候.
我們就會得到確據.
因為我們是神的兒女.
是不會再受到責備.
我們能夠去面對這些權勢.
我們能夠坦然無懼.
這是一個信心.
很想在這裡分享.
關於一位舊團友.
他對神的信心.
是在幾年前的事.
相隔和弟兄失去聯絡.
失聯大概有二十多年.
在幾年前我們重遇他.
當時他得到重病.
在一間領養醫院的舒緩病房.
接受舒緩治療.
他能夠重拾信心.
回到愛他到底的耶穌身邊.
其實有一段.
我覺得是一段奇妙的經歷.
是一位女團友.
她在這間醫院.
是專職醫療人員.
她需要在病房工作.
她會看看病人的情況.
她看到這位團友.
病歷上的名字很熟悉.
跟她聊了幾句.

$^{361}$後來她跟我們說.
原來是我們失去聯絡的一位舊團友.
一直聯絡的時候.
這位團友都願意我們去探望她.
當我們去探望她的時候.
就說回我們以前在團契生活的點滴.
我們都跟她了解.
這麼多年我們失去聯絡.
她去了哪裡.
她有沒有回教會.
她都很坦誠地說.
她沒有回教會.
甚至她沒有將自己的信仰.
生命的經歷.
跟家人分享.
因為她在一個舒緩病房.
我們了解應該不會有很多日子.
在這個地上.
每次去探望她.
我們都會跟她唱詩.
讀經去祈禱.
我們團友當中有幾個都是傳道人.
有些是牧師.
我們除了這樣做之外.
不如跟她守主餐.
去兼顧她對上帝的信心.
這位女團友就在病房裡.
打點安排.
我們有一間比較靜.
沒有受到騷擾的病房.
跟她唱詩,讀經,祈禱.
甚至跟她守主餐.
我們可以重逢這位弟兄.
去兼顧她的信心.
在幾天之後.
這位弟兄回到天家.
原來在早前.
認回這位舊團友的時候.
就已經有其他弟兄.
聯絡了她的前妻和兩個兒女.

$^{401}$所以當出席這位弟兄的.
簡單喪禮上.
她的家人都可以去送別她.
這位舊團友能夠在神和人面前.
確定她這一份信靠.
和她這一份信心.
去仰望上帝.
她能夠坦然無懼地回到天家.
接著我們看回第22節.
到24節.
並且我們一切所求的.
就從她得著.
因為我們遵守她的命令.
行她所喜悅的事.
神的命令就是.
叫我們要信她.
而指耶穌基督的名.
且照她所賜給我們的命令.
彼此相愛.
遵守神命令的.
住在神裡面.
神也住在她裡面.
我們所以知道.
神住在我們裡面.
是因她所賜給我們的聖靈.
信與愛是一起的共存.
這是耶穌基督的吩咐.
是命令我們相信她.
我們就要彼此相愛.
這是23節裡所說.
彼此相愛就成為我們信徒群體.
在主面前坦然無懼的.
一條重要的鎖匙.
約翰在上文.
在第三章裡.
已經舉出兩個例證.
說明這種愛.
經文可以看到.
三章11至17節.
首先約翰以一個負面的例子.

$^{441}$提醒我們每一個信耶穌的人.
應當彼此相愛.
這是從起初所聽見的命令.
他更要警戒.
警告我們信徒.
不可以效法該人.
為什麼呢?.
因為他殺了他的兄弟.
這就是屬於惡者.
為什麼該人要這樣做呢?.
在憲制的事情上.
發生了什麼事呢?.
為什麼在憲制完結後.
該人要動殺機呢?.
原來是心出於疾度.
心出於不高興.
該人惱羞成怒.
弟弟,伯伯.
究竟你在憲制的時候.
你對神做了什麼?.
這就是出於.
該人心中的想法.
一個家庭.
家人不是最寶貴.
是應該要珍惜的嗎?.
家庭關係應該是最重要的.
家庭裡.
我們彼此分享.
彼此交流.
是互相傾訴的地方.
更應該是彼此相愛的地方.
互相扶持的地方.
弟兄姊妹教會是家.
旺角浸信會也是大家的家.
我們應該也是這樣.
不要因為心出了問題.
「凡恨他弟兄就是殺人的.
你們曉得凡殺人的.
沒有真理存在他裡面」.
所以弟兄姊妹心是非常重要.

$^{481}$我們也清楚知道.
我們需要常常留心和警惕.
在這個世界上.
充斥著很多.
以人為主導的.
一些心思,意念,喜好.
求主保守我們的心.
忠心仰望耶穌基督.
將我們的心單單歸給耶穌.
之後,約翰就以一個正面的例子.
一個例證.
就是講到耶穌基督為我們寫命.
在第十六,十七節.
耶穌是犧牲自己.
祂就為我們帶來了愛的恩典.
這個恩典正正就是出於神示愛.
主為我們寫命.
我們從此就知道何為愛.
主為我寫命.
我從此就知道何為愛.
約翰描述關於.
耶穌所展現超然的愛.
犧牲的愛.
祂沒有就這樣在這裡停留.
耶穌為我們寫命.
在這個高峰上.
讓姐妹約翰更加誠實地指出.
我們也當為弟兄寫命.
這是一個平衡句.
約翰這樣寫的時候.
耶穌為我們寫命.
我們也當為弟兄寫命.
約翰就是將人和基督.
耶穌的行動連繫在一起.
耶穌甘願放下自己寫命.
因此約翰就說.
你們也應該為了人的緣故.
同樣放下自己.
放下一切.
放下自己應有的權利.

$^{521}$來捨棄生命.
約翰所說的.
正正就是耶穌基督自己.
曾經這樣表白.
你們要彼此相愛.
像我愛你們一樣.
這就是我的命令.
人為朋友寫命.
人的愛心沒有比這個大的.
我們能夠秉承這種愛的目的.
不是就是不會愛得埋怨.
愛得勉強.
愛得不甘心.
相反我們應該愛得輕鬆.
愛得安心.
當愛是可以被看到.
被聽到.
被觸摸的時候.
人們又看到.
又聽到.
又感受到我們所展現的愛心.
就知道你和我.
是屬於神.
是神的兒女.
耶穌基督犧牲的愛.
為你為我.
寫命成為贖罪制.
聖經裡說.
耶穌本為大祭司.
應該負責獻祭.
他竟然獻上自己.
他這個實際的行動.
是觸得到的愛.
在結束之前.
很想分享幾個例子.
上世紀八十年代.
大家聽過愛滋病.
是當時一個新的病症.
就像今天我們面對新冠疫情一樣.
當時愛滋病.

$^{561}$不知道如何傳染.
又沒有辦法阻止.
他是一個任何藥物.
都沒有藥物.
如何控制受感染的人.
情況就像我們今天所面對.
當時在英國.
愛滋病的情況威脅很大.
所以很多人都在恐慌.
甚至他們的速度.
令他們有恐慌.
因為速度太快.
甚至有很多人因為這個病.
而離開世界.
甚至對這種病患.
愛滋病患者.
有些人會尊敬甚至排斥他們.
當時戴安娜王妃.
去探望醫院的時候.
她沒有戴手套.
冒著感染的風險.
她很從容地表達出.
她對人的愛.
對病患者的愛.
她親自和一個愛滋病患者握手.
這一幕當時被形容為很震撼.
回到2020年.
去年年初.
新冠病毒的時候.
初期爆發得很厲害.
甚至感染死亡的數字.
都上升得很快.
醫療設備都很緊張.
特別是呼吸機.
大家有留意和記憶的時候都知道.
很多地方沒有呼吸機.
因為病人實在太多.
受感染之後.
呼吸很多時候都很困難.
並不是每個病人都有呼吸機.

$^{601}$又有消息傳出.
有些受感染的長者.
自願放棄使用呼吸機.
其實這是表示什麼呢?.
是表示他們自願放棄.
醫療的機會和權利.
他們為什麼這樣做呢?.
相信是出於愛.
這群受感染的長者.
他們很想一群年輕.
仍然能夠為社會做貢獻的時候.
他們能夠得到適切的治療.
康復之後.
能夠繼續貢獻社會.
回到教會.
在過去一段時間.
我們都做了很多的服事.
特別是最近中秋節的時候.
牧師在這裡呼籲大家.
有少許的奉獻.
去籌集幾十份禮物.
結果數字都很多.
有二百多份.
當中有弟兄姊妹.
不單單來奉獻禮物.
甚至願意跟我們一起.
去派送一些禮物.
去家房.
去一些他們不認識的社區朋友.
甚至是一些獨區的長者.
又有一些少年人.
當同工來分享.
他們說有些家庭需要一些大型的家具.
或者一些電器用品的時候.
其實有人已經送了.
不過是需要搬運.
沒有人不包搬運.
徵召教會的年輕人來幫忙.
這些年輕人一收到邀請.
他們很雀躍地說.

$^{641}$可以有服事的機會.
但當完成後.
我去了解一下.
他們做得很好.
同工說不是.
原來不是搭電梯.
是將家具和電器.
原來是逐層樓走上去堂樓.
他說他們真的很累.
但原來不只是一次.
他們需要準備.
要做後備.
又要再搬運的時候.
又有需要的時候.
就要繼續去做.
他說不要緊.
他們說不要緊.
可以來服事.
這些例子其實說明什麼呢?.
可能付出的或者犧牲的.
是很少.
原來不去打機.
不去放棄喜歡做自己的事.
甚至一些娛樂.
以愛心去幫助有需要的人.
其實這就是顯出.
看到的愛.
觸摸得到的愛.
和感受得到的愛.
鄧姐妹有一句很熟悉的金句.
至於我和我家.
我們必定事奉耶和華.
我相信大家在家裡.
如果可以放到這一段金句.
都會放在當晚的客廳.
這一句就是約書亞挑戰.
以色列人信靠神.
敬畏神.
愛神愛人.
的一句宣告的名言.

$^{681}$又明白約書亞所表述的嗎?.
在離世之前.
年老年邁的約書亞.
就是這樣語重心長的叮囑.
在迦南地.
以色列百姓.
他們要謹守遵行耶和華的律法典章.
敬拜祂.
事奉祂.
過聖潔的生活.
約書亞的一生剛強壯膽.
他坦然無懼.
來到神的面前.
去到神的懷中.
約翰這一位.
是稱為耶穌所愛的門徒.
他見到耶穌犧牲的愛.
他又聽到.
又見到耶穌.
祂勸我們彼此相愛的教導.
約翰晚年的時候.
他就是重複這樣.
勸勉信徒.
在艱難的時代.
雖然受到各方的壓力.
在苦難裡.
祂提醒我們.
仍然要坦然無懼.
堅持我們所相信的.
和提醒我們要彼此相愛.
而約翰最終.
他也在神的面前.
坦然無懼.
弟兄姊妹.
你以信心切勿愛主嗎?.
你願意立志以看得到的愛.
聽得到的愛.
回應主耶穌為你獻上的.
捨命之愛嗎?.
你能夠坦然無懼地面對上帝嗎?.

$^{721}$我們一起來禱告吧.
因主耶穌.
我們實在感恩.
主我們知道有很多軟弱.
甚至有很多懼怕.
我們不知道如何面對前面的路.
前面的人生.
但你就藉著約翰的教導提醒.
我們重要的是相信耶穌基督.
我們重要的是在耶穌基督的命令裡.
我們能夠群體彼此相愛.
以致我們能夠一起.
藉著耶穌基督.
能夠走這條信心的路.
過我們能夠為主來作見證.
一個聖潔的生命.
我們來到你的面前.
你坦然無懼.
願你垂聽我們的禱告.
也幫助我們能夠來到.
繼續的靠賴你過每一天.
奉主耶穌基督的名禱告.
\newpage



\section{雅各書 5:13-20}
\label{sec:rfN73DzwrE4}
\textbf{2021年10月3日主餐崇拜直播｜陳明泉牧師｜生命見真章系列 - 與生活同步的禱告｜雅各書5\_13-20}
\newline
\newline
連結: \href{https://youtube.com/watch?v=rfN73DzwrE4}{\texttt{https://youtube.com/watch?v=rfN73DzwrE4}} ~~~~ 語音日期: 2021-10-19
\newline
\newline
\hyperref[sec:aSTV2iY5HmY]{\small{< < < PREV SERMON < < <}}
~
\hyperref[sec:index_chronic]{\small{[返順時目]}}
~
\hyperref[sec:irpyKMaPiJI]{\small{> > > NEXT SERMON > > >}}
\newline
\newline
雅各書 5:13-20
\newline
\begin{longtable}{cl}
\hline
\hline
章節 & 經文 (和合本修訂版)\\
\hline
5:13 & \begin{tabularx}{0.7\textwidth}{X} 你們中間若有人受苦，他該禱告；有人喜樂，他該歌頌。 \end{tabularx} \\ \\ \relax
5:14 & \begin{tabularx}{0.7\textwidth}{X} 你們中間若有人病了，他該請教會的長老們來為他禱告，奉主的名為他抹油。 \end{tabularx} \\ \\ \relax
5:15 & \begin{tabularx}{0.7\textwidth}{X} 出於信心的祈禱必能救那病人，主必叫他起來；他若犯了罪，也必蒙赦免。 \end{tabularx} \\ \\ \relax
5:16 & \begin{tabularx}{0.7\textwidth}{X} 所以，你們要彼此認罪，互相代求，使你們得醫治。義人祈禱所發的力量是大有功效的。 \end{tabularx} \\ \\ \relax
5:17 & \begin{tabularx}{0.7\textwidth}{X} 以利亞與我們是同樣性情的人，他懇切地祈求不要下雨，地上就三年六個月沒有下雨。 \end{tabularx} \\ \\ \relax
5:18 & \begin{tabularx}{0.7\textwidth}{X} 他又禱告，天就降下雨來，地就有了出產。 \end{tabularx} \\ \\ \relax
5:19 & \begin{tabularx}{0.7\textwidth}{X} 我的弟兄們，你們中間若有人迷失了真理而有人使他回轉， \end{tabularx} \\ \\ \relax
5:20 & \begin{tabularx}{0.7\textwidth}{X} 這人該知道，使一個罪人從迷途中回轉，會從死亡中把他的靈魂救回來，而且遮蓋許多的罪。 \end{tabularx} \\ \\
[1ex]
\hline
\hline
\end{longtable}
$^{1}$今日的講道經文是記載在《雅各書》五章十三至第二十節.
在新約聖經三百二十五頁.
你們中間有受苦的,他就該禱告.
有喜樂的,他就該歌頌.
你們中間有病了的,他就該請教會的長老來.
他們可以奉主的名,用油抹他,為他禱告.
出於信心的祈禱,要救那病人,主必叫他起來.
他若犯了罪,也必蒙赦免.
所以你們要彼此認罪,互相代求,使你們可以得醫治.
二人祈禱所發的力量,是大有功效的.
以利亞與我們是一樣聖情的人.
他懇切禱告,求不要下雨,雨就三年零六個月不下在地上.
他又禱告,天就降下雨來,地也生出土產.
我的弟兄們,你們中間若有失迷真道的,有人使他迴轉.
這人該知道,叫一個罪人從迷路上轉回.
便是救一個靈魂不死,並且遮蓋許多的罪.
多謝主席.
剛才的歌聲很好聽,一邊聽的時候,心裡有一種振奮.
靠基督德性不是那麼轟烈,原來可以按住德性.
不過我們生命裡面,如果有像詩歌那種淡定,或是歌詞裡面的確信.
我們生命裡面人生的路會容易走得多.
今天一齊看的雅各書,大概讓我們用一個這樣的角度來看我們的生命.
今天是最後一講,特別是這個段落,講到關於祈禱的一個很重要的訊息.
有時候我們知道祈禱是基督徒生命裡面不可或缺的一個很重要的操練.
不過弟兄姊妹一講到祈禱的時候,曾經有跟我說.
牧師,其實有時候祈完禱之後,人家說重生得力,但是我祈完禱之後還累了,還辛苦了.
我說你怎樣祈禱?你不是聲嘶力竭地祈禱嗎?.
你用什麼方法令到你現在祈禱會變得那麼累呢?.
他回答說,有時候我跟教會的祈禱事項,弟兄姊妹不同,所以我真的代禱,但是寫得很長.
我每一次都要背了它,然後再去看,有時候又不記得,又要找找看.
整個過程好像很用神一樣.
另外他說,牧師傳道你們祈禱,用字很典雅,我用來用去都是很粗俗的字眼.
我經常都不懂得用更加典雅,用得很動聽的說話來講.
有時候說得不夠直接的時候,又說用語用得不準,又不合神學.
整件事真的累到不得了,所以牧師不知道怎樣去祈禱才說有弟兄姊妹所說的重生得力.
大概我們不需要很快去批判弟兄姊妹所說的祈禱,說覺得很吃力的那件事,你很快批判了他.
反而我很欣賞他很真誠地將他那種掙扎和祈禱裡面面對的處境給我們知道.
事實上祈禱如果按照聖經裡面說,不是說到一件很複雜,要令到字眼堆砌得很典雅.
或者要用到很超脫的形式來祈禱,反而祈禱成為我們生命裡面很自然的一部份.
今天亞國書裡面所說的祈禱,特別是回應在第一章裡面亞國開頭的時候說的那句話.

$^{41}$這一種的生命是一個使你們成全完備毫無缺陷的生命.
成全完備毫無缺陷,很多的聖經家說是整個亞國書裡面一個很重要的主題.
簡單來說就是一個基督徒如何在生活裡面變得更加完備,很成熟的呢.
不是透過一些很外在或者很超脫的,忍在深山裡面的操練.
而在生活當中實踐上帝給我們的教導.
而今天祈禱就作為這個亞國書裡面最後一個的篇幅.
也是對我們來說生命裡面一個很重要的篇幅.
因為今天的內容幫到我們,讓我們找到祈禱的真義.
更加讓我們生命裡面踏實地朝著成全完備的方向來努力操練自己.
我們一起祈禱好嗎?求上主的靈幫助我們,讓我們明白話語和心祈禱.
天父上帝願你在我們當中對我們說話.
求你繼續來激勵我們,讓我們在你話語裡面知取養分.
更加明白禱告的真諦,求你親自來引導我們.
讓我們進入真理,聖靈求你對我們說話.
又幫助宣講的說得清楚,幫助聆聽的每位弟兄姊妹聽得明白.
恭敬將來的時間交給你,願你賜福.
多謝主上聽祈禱,奉基督耶穌得勝的聖名祈求,Amen.
我們繼續看,亞國書第五章.
其實這個段落有很多的《釋經學家》為第十二節開始的.
這裡是說一個關於喜事的問題,因為今天時間的問題.
如果有機會我再和大家做一個查經,看看這一節究竟如何處理.
我們今天跳到第十三節,到第二十節,最後一個段落.
在經文裡面,首先第一句十三節裡面.
你們中間有哪一個受苦的,你就要禱告.
有喜樂的,他就應該歌頌.
你們中間有什麼處境,接下來第十四節都是這樣重覆的.
他在說不同的處境,弟兄姊妹生活裡面面對不同的歷程或者際遇的時候.
我們都要將這一切成為我們的祈禱,和神來建立關係.
在聖經裡面說到,我們是哪一個,我們從母腹裡面,上帝已經知道我們是哪一個.
我們的生命是屬於祂的,我們的行事為人,我們生命所有的一切,上帝已經完全掌握.
甚至我們的名字,我們自己的名字都沒有,上帝已經知道我們在暗中受造.
所以在基督教的信仰裡面,或者我們今天按著聖經的神學觀念來說.
上帝不會忘記我們,不會不知道我們是哪一個.
在我們生命裡面經歷過的一切的際遇,上帝都瞭如指掌,上帝都會看在眼內.
甚至我們人生的一切,我們生命的軌跡怎麼走,是上帝為我們所預備的.
我們生命某程度上帝預定了我們生命走一條這樣的軌跡.
所以上帝不是對我們生命一無所知.
如果基於一個這樣的角度來說,我們向上帝去表達我們自己的際遇.
無論是我們的受苦,或者我們的喜樂.
那些一切不是純粹將事情稟告一次給我們天上的阿爸父知道.

$^{81}$如果是這樣的時候,天父什麼都知道,我是誰都知道.
那我們向上帝說些什麼呢?其實是說我們的心.
說我們的內心裡面的體會,我們的感受.
或者面對著這些際遇的時候,我自己是怎麼想的.
我的體會是什麼呢?就好像聊天一樣.
將你心裡面的說話,來向上帝傾心吐熱.
我曾經聽過一個很簡淺,但對我來說很重要的一個比喻.
就是張無愛牧師,他那次在神學院裡面有一次教祈禱.
張牧師說這句話的時候,他就很甜蜜地說.
祈禱是什麼呢?就是一個小孩,你當他是一個小孩回幼稚園.
他回幼稚園的時候,回到學校看到有很多新奇刺激的體會.
回到學校裡面看到老師學到一些新的知識.
在學校裡面有同學可能跟他一起玩或者欺負他.
這些一切的經歷,在他放學的時候,媽媽就牽著他的手.
一路走的時候,小孩就將他上學所有的事情都說出來.
媽媽最想聽的是什麼呢?.
不會是A字怎麼寫,B字怎麼寫,或者老師怎麼教那個內容.
反而媽媽想聽的就是,這個孩子經歷過一個上學生活的時候.
他有什麼體會,他的快樂,或者令他哭,令他不安.
這些情緒,這個媽媽就很想聽到.
聽完之後,媽媽就會抱一抱這個小朋友.
加他力量,鼓勵他繼續走面前的路.
這是一個很簡單的比喻.
不過當我聽的時候,原來我們天上的阿爸父是這樣聽我們祈禱的.
我們不需要將我們所有生活裡面的細節,巨細無遺.
好像資料性的那些,來稟告一次給上帝知道.
因為上帝已經知道了,不過上帝想知道的就是我們的心.
在我們生命裡面一些很重要的體會.
我們人生裡面不會所有際遇都是平順的.
有高有低的時候,我們生命裡面,我們怎麼去經歷上帝.
或者怎麼去面對著生命裡面這些不同的處境.
將你們的心路告訴上帝知道.
值得留意的就是,有喜樂的,他就該歌頌.
有喜樂的,那個喜樂的字是沒有譯錯的意思.
不過我們覺得不夠道.
當你覺得很高興,當你覺得自高氣揚的時候.
我們覺得很興奮,當你心靈裡面很快樂的時候.
你要將這些快樂的事,向上帝歌頌他,表達出來.
所以上帝不是純粹接收我們生命裡面難處的一個上帝.
在我們生命裡面,人生的高與低,所有起起伏伏,甚至乎平淡如水.

$^{121}$我們的心都有很多東西可以體會.
我們就將我們的心呈現在上帝的面前.
邀請上帝介入在我們生命當中.
唯有這樣的時候,我們的生命,我們的禱告才是貼地的.
相反來說,如果我們只是像品臣這樣的方式.
資訊,information(資訊)告訴上帝的時候.
我們覺得上帝好像不掌握我們.
我們用了大量的時間來講我們自己的處境.
其完之後,其實都沒有心機再講你自己內裡的需要.
因為你都已經累了.
亞國書提醒我們,生命裡面,順與逆,或者快樂受苦.
我們都應該要禱告.
當我們將我們心靈裡面的呼聲向神表明的時候.
我們就跟上帝的關係拉近一點.
禱告是邀請上帝介入在我們的心靈裡面.
邀請上帝跟我一起同行生命的一切.
這是第一個亞國書教我們祈禱的重點.
第二個重點,在第十四至第十五節.
你們中間有病了的,他就該請教會的長老來.
他們可以奉主的名用油抹他,為他禱告.
出於信心的祈禱,要救那病人,主必叫他起來.
他若犯了罪,也必蒙赦免.
在第二個重點裡面,很多時候我們很快一讀.
我們就純粹就很…不是很理會字眼上面的使用.
我們就覺得出於信心的祈禱,上帝一定應允.
上帝一定如我所願來報答我們.
我們會拿著這一句純粹的應許,直譯地看.
不過在這裡的用字很小心.
首先就是你們中間有病了的,和合本有疫病了.
其實很多譯本都已經在說你們中間有哪一個是有軟弱的.
後來因為後面用油來抹他,所以是說身體的軟弱為主.
不過如果用這個角度,純粹說病了,然後抹油,就好像下藥一樣.
去理解這一節的經文.
不過原文裡面說的軟弱,除了身體軟弱之外.
更加多一點的就是灰心這個字.
你們當中有沒有人病了灰心的.
有些重病的弟兄姊妹,或者有些長期病患的.
生命裡面因為面對著這種疾病.
所以病到一個處境就是人失去鬥志.
失去了對神的信心,灰心地過他的生活.

$^{161}$所以病了不只是在說肉體上的軟弱.
當一個人面對著身體和心靈的軟弱的時候.
他應該找教會的長老或者今天的教會領袖.
重要的地方不是純粹抹油的動作,而是要用奉主的名.
整個過程是藉著伊治的主,藉著耶穌基督的名字.
來為這個人做一些抹油,為他來禱告.
抹油是什麼意思呢?.
抹油其實是一個宗教禮儀,早代教會會多一點.
今天就少一點了,其實在當代社會裡面.
在亞國書的時代,你說用藥或者有醫生看病.
沒有我們今天那麼普遍,今天醫務所多過7-11.
香港人病了去看醫生是很容易的一回事.
不過在耶穌基督的時代也好,或者在亞國書的時代也好.
在當時來說,一個人病了,醫學沒有那麼昌明的時候.
一個人要面對著,怎樣去熬過一個病是一件很艱辛的事情.
很多人未熬過都已經離開了這個世界.
所以病怎樣得到醫治呢?除了看到醫生,遍尋很多真的能醫治病的醫生之外.
第二件事就是藉著信仰來幫助他.
而這個抹油就好像有一點藥性的油.
但是對於今天來說,你可能只是用一瓶白花油,用萬金油來塗一下.
你不敢說抹一下這個油可以醫治你的病.
是會舒緩你,舒服一點,但是是不能夠醫治你的病.
不過在當代社會裡面,將這個抹油成為一個宗教儀式.
就是一邊塗油,一邊祈禱,令一個身體很軟弱,心靈很灰心的人.
藉著長老一種很外顯形式的抹油,奉主名祈禱.
讓他可以救拔這個病人.
第十五節裡面說,出於信心的祈禱,就是要救那病人.
主必叫他起來,在雅各裡面很小心用這個字,就是用起來這個字.
起來這個字的意思是闊一點的,可以是病得醫治的.
耶穌叫那個癱子,拿你的玉子起來,行走吧.
這個是醫病,一個跛了的人叫他起來,他的腳就得到醫治.
不過這個起來這個字,不是純粹說病得醫治的.
可以是說他的病依然在,但是他的心失意了,灰心易冷了.
因為這個信心的祈禱,令一個面對病患當中的人.
他不會失去那份信心和盼望.
他依然帶著有力量的心,來走面前的路.
雅各在這裡很小心用這個起來這個字.
因為上帝的主權可以叫一個病患者得醫治.
他不否定一個這樣的可能.
但是他看重的,除了身體軟弱可以得到醫治之外.

$^{201}$他更加看重的是心靈很軟弱的那些人.
有些人病了很長的日子.
但是因為面對著很多病帶來的困難.
有時也失去了生存的鬥志.
這個時候,那個抹油的禱告既令他舒緩一點.
同時也給一些長老的祈禱來激勵他.
令他更加有力量,有信心地面對他生命裡人生要走的每一步.
所以在這裡,這個抹油的祈禱.
或者出於信心的禱告.
那個幫助就是有些人心灰意冷.
特別是一些病患者,有些很沮喪的人.
當有長老用這種這麼凶險的方式來幫助他的時候.
讓他可以重新起來,有力量走面前人生的路.
在這裡,我們做傳道同工也有很多探病的經驗.
有時我們面對著很多的生死.
或者面對著頂展會病患,醫院探病.
我想羅牧師會更加多,曾真牧師.
他們真的跑醫院跑得很多.
我們自己同工也會問,究竟我們跑那麼多醫院.
我們走到去為頂展會祈禱.
其實我們在做什麼角色呢?我們有什麼幫助到呢?.
其實進到醫院的時候,有時我也覺得自己左手左腳.
尤其是有些要檢查的護士為了那個病人要去看病.
你又要站在床邊.
但你站在那裡的時候,護士就好像很厭惡的眼神.
不要妨礙到所有人,有時我們也覺得自己很硬朗.
如果我長得短,小一點,不要那麼長手長腳.
我坐在床頭也不會妨礙到別人工作.
但最壞的就是我每次都是這樣.
還要蹲在那裡,屁股又夾著別人.
我也覺得自己很硬朗.
每一次去醫院探病,我為了什麼呢?.
但當我讀到這篇經文的時候,我終於通了.
我目的不是醫生,我沒有辦法治療他身體上的病.
但當我可以握著他的手,握著頂展會的手.
陪他一起度過生命的一刻的時候.
我知道給他多一點信心,多一點鬥志.
來面對身體的艱難的時候.
他會有力量走面前的路.
兩三個月前,我自己長大後有一個弟兄.

$^{241}$他年輕五十多歲,但因為腦裡長了腫瘤.
後來經歷過很多電療,腦部開始萎縮.
這個弟兄很好,我和太太拍拖的時候.
他可以做電燈泡,因為讀神學又沒錢.
又怕我們捨不得,開著車載我和阿維利去兜風.
我們兩個拍拖之後,他載我們各自回家.
他是一個好人,在我心裡是一個很愛錫不可多得的.
很明白我的弟兄,但當他患了腦瘤的時候.
我記得那次他的家人叫我去探望他.
我上去他家探望他,那時候見到他身體已經大不如前.
很活躍的,到那一刻他已經很軟弱.
他以前開著電單車很英明神武,但那一刻我見到他的手.
因為做了很多療程,他的身體已經開始不太能動.
他的手有捲曲,身體已經消瘦,沒力.
去到在床邊的時候,我都不知道可以為他做什麼.
他是一個牧者,他很愛錫我的一個兄弟.
去到那一刻我對他的病完全沒辦法去做.
但我知道那一刻我可以拿著他捲曲的手.
我望著他的眼,他說話都說得不清楚.
但在那一刻我可以拿著他捲曲的手.
那一刻我們好像把心重新貫通一樣.
我去耳邊細細聲說:放心,上帝願你同在.
雖然見到他的身體,我是很心痛.
但我知道我的祈禱對他來說,為他帶來一些安慰.
就是這種緊握和幫他面對生命這些難處的祈禱.
希望這個弟兄可以安然度過他生命最後一段階段.
探驗完那次之後沒多久他就回到天家.
他就很安然地在家人的陪伴之下回到天父那裡.
面對身體軟弱的人,作為基督徒我們未必是醫生.
甚至是醫生也好,都沒有百分之一百可以醫好肉體軟弱病的把握.
但基督徒我們卻可以在這些病患者裡面.
可以跟這些弟兄姊妹一起,可能病到他自己都失去了生存的勇氣和盼望.
但是我們的禱告就幫助他,他的信心有力量走面前的路.
所以這個起來帶來的就是讓他的心不再灰心而冷.
不再失去盼望.
這個禱告,這個沒油的禱告的重點不是在於油抹在他身上.
而是做一些很外顯的方式幫助一個病患者.
一個失去盼望的人有力量走面前的路.
經文再繼續的,他若犯了罪也必蒙赦免.
患病和犯罪在耶穌的眼中裡面.

$^{281}$這件事比較難,耶穌問那些法利場人.
在耶穌來講兩件事都不難.
同樣在我們基督徒裡面.
我們不是靠我們自己的能力令這個人可以起來.
只是我們奉主的名幫助這個生命裡面軟弱的.
如果他生命裡面還有一些糾纏他的罪.
藉著這個禱告讓他生命可以輕省,可以放心.
可以得到蒙赦免,以至他可以安然走面前人生的路.
可以病得醫治,但更重要的是讓他的灰,失望的心可以重新起來.
所以弟兄姊妹,今天我們如果講禱告的真義.
如果講禱告裡面真是很到位的禱告是什麼呢?.
是一個出於信心的祈禱.
不單只是我們自己的信心.
也是我們信心如何燃起灰心喪志的人的信心.
我們一起緊密的禱告,是會讓人的生命可以重新起來.
讓發酸的腿,垂下來的手,心靈裡面疲乏的人.
可以重新站立起來,以至有力量走面前的路.
出於信心的祈禱,是雅各提醒我們第二個我們需要實踐的重點.
第三個重點,十六節到第二十節.
「所以你們要彼此凝聚,互相代求,使你們可以得醫治.
二人祈禱所發的力量是大有功效的.
以令雅與我們是一樣性情的人.
他懇切祈禱,求天不要下雨.
天就三年零六個月不下在地上.
他又禱告,天就降下雨來,地也生出土產.
我的弟兄們,你們中間若有迷失真道的.
有人使他迴轉,這人就該知道.
叫一個罪人從迷路上轉回,便是救一個靈魂不死.
並且遮蓋許多的罪」.
在這段經文裡面,最吸我們的眼球.
我們看到的,我們覺得很開心.
二人祈禱所發的力量是大有功效的.
我們經常面對一些艱難都會強調.
二人的祈禱是大有果效的,究竟誰是二人呢?.
在雅各書裡面,似乎沒有什麼上文下理去描述.
究竟那個二人是誰?.
當然,如果你按著他的內容,那個二人.
大概是按著雅各所講的倫理道德觀念.
行事為人謹慎,有很好的道德的一群人,那些叫二人.
這未必是錯的.

$^{321}$不過,如果純粹是這樣,似乎你看不透裡面的內容.
我們反而要看多一點,另外一段經文.
這是耶穌怎樣定義什麼人叫二人.
我讀給大家聽,在路加福音第十八章,十一至十二節.
在高段經文裡面是這麼說的.
(念經).
我告訴你們,這人回家去被那人搗亂,被那人打死,被那人殺死.
(念經).
在路加福音裡面,耶穌看到法利賽人和瑞利.
特別是針對法利賽人,他看到他們自意識的祈禱.
這是當代社會,耶穌時代裡面二人的代言人.
一講起二人就想起法利賽人,文士,宗教領袖.
但耶穌就用一個這樣的比喻來講.
(念經).
不過耶穌說面對瑞利,他望天都不敢,隨著自己的心口.
耶穌說這個瑞利比法利賽人倒算為異.
耶穌看這個瑞利的義比法利賽人更高.
因為至高的必降為卑,至卑的必升為高.
在講什麼呢?二人是哪一個呢?.
其實二人也是一個很空泛的定義.
不過如果你覺得你是那個二人,大概進入了一個自以為異的圈套.
你就會反而拿不到這句經文.
反而你覺得自己不是那個二人,你的祈禱在耶穌眼中倒算為異.
這樣說好像頭都爆了,簡單來說.
其實這個世界裡面所講的二人從來都不是蹲在那個款.
或者我們站在道德高地,覺得自己做了很多的操練.
懂很多祈禱裡面的術語,或者我們懂得一些很深奧的道理.
在耶穌的眼中裡面,真正的二人只是赤誠地闖開自己的心懷意念.
看到自己生命裡面的不足,看到自己的貴乏,自己的卑惡.
二人很多時候就是標榜自己做了多少為上帝而做的事情.
不過在耶穌眼中看為異的人往往的是看到自己生命裡面的不足.
達不到上帝的要求,很自卑地去到上帝的面前.
祈求上帝赦免他生命的罪.
在這段經文裡面,如果講到那個二人.
其實就是一種對自己很有意識,很知道自己發生什麼事的人.
我們不會以自己的善行作為墊高自己的一個工具.
來講我祈禱特別靈,我祈禱特別勁.
反而是讓自己放下在一個和其他罪人沒什麼分別的.
看到自己生命不濟的一個人.
在這裡面,既然一個二人不是想像中.

$^{361}$我們經常覺得那種道德高地的一班人的時候.
雅各他勸勉我們,所以你們要彼此認罪.
互相代求,使你們可以得到醫治.
我們不是想像中那麼神聖.
我們的祈禱不是比別人更高超.
我們只不過是在這個世界裡面.
繼續在一個不太理想的世界裡面.
繼續努力,形形異異,努力打拼的一個人.
我們生命裡面都滿有很多的罪.
所以在我們生命當中.
弟兄姊妹之間要是一個互相認罪的.
互相代求的一種關係.
所以弟兄姊妹,你留意一下你自己的小組生活.
你的信仰生活.
當你一直和別人接近.
信仰裡面越來越貼近的時候.
你會不會看到別人的壞處越來越多呢?.
看到別人不理想的地方.
你會不會越來越懵整.
而覺得自己越來越神聖呢?.
如果當你面對一個這樣的處境的時候.
我們要快快悔改.
不要看見別人不理想的地方.
繼續去鑽挖別人.
然後將自己覺得不像他們那麼差.
我就會高尚一點.
反而面對一個不理想的人.
你更加要反省自己的生命.
我會不會像那個我批判得很硬的人一樣.
其實我在上帝眼中也是一個罪人.
這個世代需要彼此認罪代求的一個世代.
而不是互相指責,互相批評.
或者是別人不是的一種關係.
如果互相指責或看見別人的東西.
只是眼中釘.
你經歷不到亞國所說的那種完全的生活.
反而是一個不理想的世代.
看到別人的問題.
也看到自己的悲污.
我們坦誠在上帝面前.

$^{401}$一起來認罪代求的時候.
我們的生命才可以得到醫治.
得到更新和改變.
亞國再繼續說.
我們這樣祈求.
他用一個舊約先知的例子.
以利亞作為一個再做一個的展述.
他說以利亞和我們都是一樣性情的人.
以利亞先知和我們一樣性情.
舊約或猶太人傳統裡面.
他沒有經歷死亡.
直至被上帝提了上天.
這個很厲害的先知.
和我們相提並論.
他說一樣性情,不是一樣能力.
不要讀錯.
這個先知他怎樣令到.
以色列民知道上帝的審判.
三年零六個月不下雨.
然後一祈禱就下雨.
其實他都是將自己的禱告.
懇切地放在上帝面前.
以利亞是一個有真性情的人.
見到這個世界有很多不滿.
見到以色列民很不濟.
不過他都是求上帝.
懇切地祈禱.
亞國書更加引用他說.
其實以利亞都是一個失敗的先知.
當有人追殺的時候.
當他自己面對著很多困難的時候.
他都想自己快快取取了斷.
了斷了自己.
發晦氣.
覺得只有自己做事.
亞國說我們和以利亞的性情一樣.
一方面很厲害.
但同時都很不堪一擊.
一有事來到.
我們都不是想像中那麼強.

$^{441}$但是上帝都聽他的祈禱.
懇切地祈禱的時候.
上帝就聆聽.
這種懇切的態度.
更加要放在我們人與人之間的關係裡.
我的弟兄們.
你們中間若有迷失真道.
有人使他迴轉.
這人就該知道.
叫一個罪人從迷路上轉回.
便是救一個靈魂不死.
並且遮掩許多的罪.
禱告帶來的應用最後是什麼呢.
是人際關係.
是一些你覺得那些人很不堪入目的那些人.
你要挽回他們.
亞國為什麼會這樣寫呢.
其實這是耶穌基督的教導.
如果你有機會回去看福音書.
每一次耶穌講完的禱告教導之後.
接下來就是人際關係.
例如你獻祭.
耶穌說你獻祭的時候.
想起有人得罪你.
你要放下祭物.
立即和弟兄復和.
那才是獻祭.
這個教導到最後的應用.
就是人與人之間.
或者看到那些罪人.
你如何去挽回他.
或者放下自己的成見.
復和.
那就是祈禱裡面最佳的應用.
亞國也是用這個角度去寫亞國書.
在這裡最後講到你祈禱了.
你懇切祈禱了.
到最後你不要放棄身邊那些沉淪的人.
你知不知道.
如果我們在真道裡面拯救一個人.

$^{481}$就叫一個人靈魂不死.
叫這個人不會再落入生命的困難當中.
整個亞國書在這裡就收尾了.
很突兀的.
不過這個也是對我們今天.
很多的弟兄姊妹一個很重要.
好像提醒我們.
我們下來的下一章要怎樣做.
我們今天的禱告.
我們會不會純粹將資訊告訴上帝.
好像品臣一樣.
亞國書提醒我們.
我們與上帝的關係不是品臣.
而是一個爸爸媽媽和子女之間一種親密的關係.
將你的心聲告訴他.
上帝聽的.
上帝很想知道你的心靈空間是怎樣.
亞國書更加提醒我們.
今天我們需要信心的祈禱.
這個信心的祈禱.
主要是幫助那些病患.
或者生命裡面心灰意冷的人.
你信心的祈禱.
是會幫助灰心意冷的人.
重新有力量走面前的路.
叫他生命可以重新起來.
讓他不會落入絕望的那種光景裡面.
第三個重點.
我們的祈禱.
我們不知道誰是義人.
我們也不要純粹自封為義人.
如果從耶穌的角度來說.
真正的義人是敞開我們的心靈.
向上帝承認自己的罪.
自己的問題.
那個反而倒算為義.
帶著這一份謙卑.
帶著一份懇切在上帝面前去求的信心.
將祈禱化為復和的力量.
將人帶到上帝的庚前.

$^{521}$這個就是在《雅各書》裡面.
告訴我們一個祈禱很重要的信息.
不知道在過去這段日子.
你看《雅各書》你生命有什麼得著.
我盼望我們整個系列.
整個直著雅各裡面.
跟弟兄姊妹一起分享.
一起去領受的懷與.
希望我們生命.
真的在這個世代裡面.
看到一班基督徒生命的真章.
不要成為一班若不堪封.
無病呻吟.
或者生命裡面失去方寸的一班人.
即使時代是怎樣.
世界是怎樣也好.
上帝的真理永恆不變.
當我們按著這個真理.
站立得穩的時候.
我們的生命就可以成為別人的祝福.
我們可以信靠的一班人.
求主繼續幫助我們.
讓我們的生命.
繼續看到上帝給我們的真章.
喜樂有力量地走.
我們人生每一步.
我們一起祈禱.
天父上帝願你幫助我們.
在我們的生命裡面.
我們知道我們不可以失去你的救恩.
我們更加求你.
讓我們檢查我們的心靈.
我們會不會有時候.
自以為是.
或者覺得自己已經很飽足.
但實際上可能我們是.
實際上是貴乏.
或者是生命裡面很多的卑污.
求你的寶血潔淨我們.
更加幫助我們在這個世代裡面.

$^{561}$成為一個禱告的人.
讓我們的信心禱告達到在你跟前.
又將很多軟弱的.
或者生命裡面灰心的人.
達到在你跟前.
藉著我們的努力.
讓人生命得著更新和改變.
又求你在我們事奉的過程裡面.
讓我們找到生命的出路.
讓我們生命裡面.
有些人可能曾經.
和我們之間有嫌隙的.
讓我們有勇氣.
和他們來復和.
實踐禱告最重要的.
給我們的應用.
求你繼續來恩代我們眾弟姐妹.
祝福我們眾人.
走面前的路沒有上帝的恩典.
多謝主上聽祈禱.
奉基督耶穌得勝的聖名祈求.
\newpage



\section{馬太福音 14:22-23}
\label{sec:irpyKMaPiJI}
\textbf{2021年10月17日崇拜直播｜陳明泉牧師｜放心!不要怕｜馬太福音14\_22-23}
\newline
\newline
連結: \href{https://youtube.com/watch?v=irpyKMaPiJI}{\texttt{https://youtube.com/watch?v=irpyKMaPiJI}} ~~~~ 語音日期: 2021-10-21
\newline
\newline
\hyperref[sec:rfN73DzwrE4]{\small{< < < PREV SERMON < < <}}
~
\hyperref[sec:index_chronic]{\small{[返順時目]}}
~
\hyperref[sec:LQ_YiQLwO5g]{\small{> > > NEXT SERMON > > >}}
\newline
\newline
馬太福音 14:22-23
\newline
\begin{longtable}{cl}
\hline
\hline
章節 & 經文 (和合本修訂版)\\
\hline
14:22 & \begin{tabularx}{0.7\textwidth}{X} 耶穌隨即催門徒上船，先渡到對岸，等他叫眾人散去。 \end{tabularx} \\ \\ \relax
14:23 & \begin{tabularx}{0.7\textwidth}{X} 疏散了眾人以後，他獨自上山去禱告。到了晚上，只有他一人在那裡。 \end{tabularx} \\ \\
[1ex]
\hline
\hline
\end{longtable}
$^{1}$今日要為大家讀出的經文是在大家手上的聖經新約瑪太福音第十四章二十二到三十三節.
是在大家手上的聖經第二十二頁.
聖經這樣說.
耶穌隨即吹門徒上船.
先到達那邊去.
等他叫眾人散開.
散了眾人以後.
他就獨自上去禱告.
到了晚上只有他一人在那裡.
那時船在海中.
因風不順.
被浪搖晃.
夜裡四更天.
耶穌在海面上走.
往門徒那裡去.
門徒看見他在海面上走.
就驚慌了.
說:是個鬼怪.
害怕喊叫起來.
耶穌連忙對他們說.
你們放心.
是我.
不要怕.
彼得說:主.
如果是你.
請叫我從水面上走到你那裡去.
耶穌說:你來吧.
彼得就從船上下去.
在水面上走.
要到耶穌那裡去.
只因見風勢大.
就害怕.
將要沉下去.
便喊著說:主啊救我.
耶穌趕緊伸手拉著他.
說:你這小信的人.
為什麼疑惑呢?.
他們上了船.
風就拒了.
在船上的人都拜他說.

$^{41}$你真是神的兒子了.
請大家靜一靜.
各位姐妹平安.
一齊來到敬拜上帝.
不知道大家在年輕的時候.
有沒有做過一些傻事.
年輕人的標誌就是做傻事.
在小弟年輕的日子.
曾經和一班同學.
一起去到大嶼山的長沙.
去度假屋.
去宿營.
也不是什麼傻事.
不過晚上覺得悶悶的.
於是在晚上漆黑一片的沙灘.
幾個男孩就說.
不如我們立即下去沙灘游泳.
那時候大概晚上.
可能是十二點多一點鐘的時間.
其實挺危險的.
我媽媽臨回天家之前.
我沒有將這件事告訴她.
很重要.
當我們一起去沙灘游泳的時候.
在遠處游泳的時候.
覺得很爽.
又沒有夠醒.
天氣又很舒服.
但在遠處游泳的時候.
我看到在海平線那裡.
有些白色的東西在飄.
然後就很害怕.
但害怕的時候又不敢出聲.
大家是一班男孩.
害怕的時候會很懦弱.
所以就說.
我有點不舒服.
不過你看那裡好像有些白色的東西.
我先游回岸邊.
然後再跟同學說.

$^{81}$其實害怕.
但又要假裝想回去休息一下.
我幾個同學就說.
你害怕而已.
不用害怕.
然後他們自己就跑過去.
跑去那些白色飄在那裡.
我看到他們的時候.
我真的非常害怕.
還跑過去.
這麼恐怖的東西.
還跑過去游泳.
我跟同學就相反方向.
我就回到岸邊.
他們就一直游出去海中心.
那時候因為水不是太漲.
游了一會之後.
就去到那個白色位置.
我離遠遠的.
手遮著眼睛.
糟了.
他們會不會有事.
腦裡就有很多畫面.
如果有事的話.
我會不會拖一隻船出去救他們.
或者馬上回家等等.
在一直害怕的時候.
那班同學就出去的時候.
就說不用怕.
馬上把手拉起來.
原來是一個白色的膠袋.
被一個木的裝柱卡住了.
有風在吹.
所以就在那裡飄來飄去.
自己嚇自己.
看到他們這樣的時候.
我就鬆了一口氣.
自己嚇自己.
膠袋來了.
然後他們再離遠.

$^{121}$還有事.
還有什麼.
手就馬上拉起來.
我心裡就想.
糟了是不是人頭.
我真的怕得不得了.
拉起來之後.
原來有個浮波.
有個網和海草在那裡.
浮在那裡.
拉起來.
遠看就像人頭.
其實是一個浮波.
多了這幾個同學.
他們沒有懼怕.
幫我解了我心裡.
很多想像恐怖的場景.
如果我不處理.
可能那晚在長山的道家屋.
我睡不著.
這個故事.
說的是一件事.
原來人裡面有很多恐懼.
有些恐懼.
其實是自己嚇自己的.
更多的恐懼.
有些心理學家說.
原來恐懼有時候很厲害.
是來自有些人的想像力特別豐富的人.
特別多恐懼.
因為有些事情不是真實出現.
但在腦子裡面已經繪形繪聲.
有很多事情.
某程度上威嚇你.
或者你的心靈裡面.
有很多舊時積累了的.
有些事情還沒處理.
以致到某時某刻.
將這些恐懼.
或者將這種感覺.

$^{161}$籠罩在你的生命當中.
剛才那些只是人生小片段.
但原來人生裡面.
其實你和我都有不同.
害怕的事情.
恐懼的事情.
如果這些恐懼的事情.
我們處理的時候.
會令到我們的生命裡面.
會慢慢積累.
甚至會將其他的感受.
炒在一起.
有時候我們會將厭惡和恐懼.
混在一起.
最簡單的例子就是蟑螂.
有些女士很害怕蟑螂.
害怕到臉都青了.
其實你一腳就能夠打死牠.
但為什麼你會這麼害怕牠呢?.
牠說很噁心.
其實噁心也不需要害怕.
牠看到我都很噁心.
其實牠不會很害怕我.
其實同樣.
有時候將恐懼混在一起.
很多事情.
就會成為一個.
合制著我們心靈.
很多人生.
很多很奇怪或很荒謬的事情.
就會出現了.
今天我們從這個角度來看看.
這一段經文.
其實在《馬可福音》裡面.
他講到恐懼和信心.
特別是在《馬太福音》第十四章.
至第十七章.
這個篇幅叫做信心.
和恐懼的篇幅.
耶穌用不同的場景.

$^{201}$特別是一班門徒.
去經歷這個全道的訓練.
一個門徒訓練的過程裡面.
他們都需要挑戰.
生命裡面有很多的恐懼.
藉著耶穌的同在.
和耶穌的教導.
一步一步幫助這一班.
曾經都經歷過風浪的一班漁民.
或者一班男士.
幫他們可以更上一層樓.
讓他們的生命可以承擔這個的詔命.
希望今天我們看這一段經文的時候.
給我們有多一點的亮光啟迪.
或者一路深思.
想想我生命裡面有沒有一些恐懼.
其實籠罩著我們.
以致我們今天生活.
好像沒有辦法再踏前.
希望今天這一篇訊息.
可以幫助到我們.
我們一起祈禱.
求聖靈幫助我們.
讓我們明白話語.
同心祈禱.
天父上帝願你在我們當中.
帶領我們眾人.
求你對我們講說話.
幫助每位弟兄姊妹聽得明白.
幫助孫講的講得清楚.
願你的話語將你的心意傳遞出來.
讓我們一同領受.
讓我們在這個世代.
成為你喜悅的子民.
恭敬將來的時間交託.
願耶穌基督你親自帶領.
獻上祈禱.
奉基督耶穌得勝的聖名.
Amen.
好 我們一起看聖經.

$^{241}$在馬達荒第十四章第二十二節.
這裡講到耶穌催門徒上船.
去到岸的另一邊.
為什麼要這樣做呢?.
其實是因為耶穌收到一個消息.
就是斯贈若翰比希律.
已死亡.
耶穌收到這個消息的時候.
其實是想上山.
獨自來祈禱.
安靜一下.
不過在祂做這件事的時候.
有很多的群眾跟著祂.
耶穌看到這一群人孤苦流離.
於是用五個餅.
兩個魚餵飽五千人.
使得這一群人.
某程度上.
肚腹又飽足.
又跟他們講道.
跟他們同在一段日子.
接著做完這個神蹟之後.
耶穌真的要休息了.
耶穌真的要祈禱了.
將那些眾人.
安排他們回到自己的地方.
接著耶穌和門徒.
離開那個地方.
那麼耶穌在過程裡面.
因為要自己獨自來祈禱.
所以安排那十二個門徒.
坐船離開岸邊.
然後他自己一個人.
就上山上面祈禱.
這個就是那個篇幅.
耶穌隨即催促門徒上船.
先渡到那邊去.
等他叫眾人散開.
散了眾人以後.
他就獨自上山禱告.

$^{281}$去到祈禱的時候.
悼念晚上只有他一人在那裡.
那時船在海中.
因風不順.
被浪搖撼.
夜裡四更天.
耶穌在海面上走.
往門徒那裡去.
門徒看見他在海邊走.
就驚慌了說.
是個鬼怪變害霸.
喊叫起來.
這個22-26節.
他就初步講到那個情景.
耶穌要獨自祈禱.
他就剩下門徒.
來休息一下也好.
或者先往另外岸的一邊.
給他們一些歇息的空間.
這個過程某程度都是門徒訓練.
你會見到.
耶穌不是純粹叫他們上船.
在過程中.
他們這群門徒是有經歷的.
他們獨自面對一個大風大浪的海面.
耶穌就一個人自己去祈禱.
如果你知道耶穌一直和門徒在一起.
無論什麼風浪.
什麼處境.
耶穌都和這群門徒同在.
不過特別是到了這個時間.
面對這些風浪.
只剩下那12個門徒.
當這12個門徒去到風浪的時候.
四周都風雨交加.
好像前幾天打風一樣.
差不多天亮.
四日更.
天差不多破曉時分的時候.
耶穌祈完禱.

$^{321}$他就直接從山裡面走到海面上.
行捷徑.
不過耶穌的做法就是.
從海面上走到門徒的地方.
從船的位置過去.
在走過去的過程.
對耶穌來說.
他是神的兒子.
他能夠這麼做.
對我們來說.
我們覺得是很順理成章的.
不過對於門徒來說.
他從來沒有見過耶穌用這樣的方式.
走到海上走.
甚至在他們的心裡面.
剩下這12個門徒的海中心.
又大風大浪.
視野又不清晰.
他們看到海面上有些東西飄來飄去.
他們第一個反應不是想起耶穌.
他們第一個反應是想起.
試過鬼怪.
在和合本譯鬼怪.
其實譯的字就沒有馬太方寫的字那麼小心.
因為這裡說到耶穌像鬼一樣.
其實也有一點.
感覺上是怕神學上製造了一些爭拗.
原本的字就用了幽靈這個字.
Fantasma.
Fantasma可能你都不知道.
如果你有看過一套舞台劇叫做Phantom of the Opera.
Phantom這個字就來自這個字.
幽靈或者是一個靈.
不知名的.
他寫得很含蓄.
他不想用鬼來形容.
這班門徒看到的東西.
我們中國人我們接受.
我們沒有一個類似的字眼.
不過這班門徒看到的是一個幽靈.

$^{361}$在上面飄浮起來.
在那個過程裡面.
其實馬太寫這班門徒的心路.
比較細緻去寫.
不是一看到就馬上害怕.
首先門徒看到耶穌.
或者看到有東西飄浮在海面上.
就驚慌.
這個驚慌是說一種心緒不寧.
馬太同樣用另一個字.
就是耶穌出世.
希律知道有一個新生王出世.
內心很不安.
驚慌的字本身就是很不安.
很憂心很擔心.
大家就有一些心裡面感覺.
很不舒服.
然後就說了.
哎呀,是幽靈.
說完之後才變害怕.
喊叫起來.
所以那個害怕.
其實之前還有些前奏.
第一,驚慌.
第二,大家說.
越說越驚.
不是越說越覺得心裡越說越舒服.
越說越害怕.
因為大家知道的事情都差不多.
然後大家就去投射.
想像.
然後一說是幽靈的時候.
就開始害怕.
害怕到一個地步就大叫.
喊叫起來.
可想而知.
你看看那個情景.
我們今天用讀者角度來看.
我們覺得這群人很好笑.
這樣都會這麼害怕.

$^{401}$大驚小怪.
一身袋油進去.
在一個風雨交加的艇上.
看到這樣的情景.
大家又這樣說.
耶穌又不在身邊.
如果有幽靈真的侵害他們.
他們憑什麼來驅走他們呢?.
這群門徒就有一種這樣的心態.
有一種這樣的情緒在.
當我們看到一個這樣的情況之後.
其實有一點值得去留意.
那個驚慌.
正如剛才所說.
不是因為分享了的時候.
會令到那個驚慌.
或者驚嚇.
害怕感減低了.
反而透過人之間.
互相的互動.
門徒那12個門徒在說.
這一種的恐懼感.
會累積起來.
事實上這種的恐懼感.
累積.
恐懼有時候大家都在想.
憑著自己的想像.
會有這樣的東西.
大家互相在說的時候.
就加強了.
那個想像力的威力.
以致那個想像空間.
越來越多.
大家越來越擔心.
越來越害怕.
越來越恐懼.
不知道我們人生裡會不會有時候.
都會面對這樣的情況呢?.
有一些的恐懼.
我們以為說出來會舒服一點.

$^{441}$因為人有一種傾向.
是想表達出來的.
但表達完之後.
告訴那個人.
那個人比你還要害怕.
然後他再將一個更大的情緒.
蓋在你身上的時候.
你就和他之間就同時害怕.
而且大家的想像空間.
越來越多.
越來越大.
以致將這個恐懼感.
籠罩在你自己的生命裡.
這一種的懼怕.
來自一種.
沒辦法去查證.
沒辦法去面對恐懼的.
那種客觀.
同時間也是透過人與人之間.
那種的溝通.
將這種恐懼感.
來加劇.
做一個例子的.
我不知道大家小時候有沒有聽過這樣的例子.
所有的小學,中學.
都說自己是亂葬崗.
我不知道你有沒有聽過.
你有沒有誰說什麼.
那個什麼.
你同不同意嗎.
每個人說這些都是這樣.
以前我被一個師兄嚇我.
我小學的時候.
我那時候少不更事.
不知道,在跳飛機.
然後他說.
我不認識他,他說這個三號.
不要跳,我說為什麼.
那裡的位置.
死過很多人.

$^{481}$他這樣跟我說.
然後我就跳,永遠都是.
我那時候很害怕.
我懷疑那個師兄.
是在耍我,不過我當真話.
來聽,於是我就將.
這些話,我和我平時.
跳飛機的同學,告訴他.
你不要跳.
三號的位置.
一二號你都不要跳.
因為隔離,很近.
怎麼會這麼小的階磚.
所以一二三號都不要跳.
然後這件事再繼續傳開去.
其他同學說,這麼恐怖.
四五六都不要跳.
基本上那個框住.
全部大家都不要跳.
有一次,小息.
小朋友都知道,打小息鐘.
大家不准動,不准什麼.
很害怕,很不幸.
我就經過那裡.
小息的時候.
小學很傻的.
小學生打鐘.
老師說,靜止所有活動.
但大家變了定格.
有時候是這樣的.
你們以前都是這樣.
我那時候剛好.
走過跳飛機,但腳.
不敢插下去.
所以我是單腳,整個動作.
因為很害怕.
我其實一路.
害怕,害怕到中學,其實我都很害怕.
這件事,後來我.
受不了,我真的和家人說.

$^{521}$我的家人說,這個世界.
有什麼地方沒死過人.
哪間學校不是亂葬崗.
一講這件事.
我真的被人騙了.
很多年,跳飛機.
其實沒任何事,跳了這麼多年.
不過.
自己心裡和同學去談.
越談就將這種.
感受放大了.
在一些.
大家患得患失.
不知道的事上,大家這樣.
去講,就會將這種恐懼感.
加強了,其實.
我們人生裡都是一樣的.
有些恐懼可能是類似這些幽靈.
但有些恐懼是說.
聽旁說.
早幾年,我們以前的沙士.
說鹽是可以怎樣.
大家很害怕,盲搶鹽.
或者不同.
的方式.
我們覺得這個世界實在很不友善.
當很不友善,我們又不知道.
怎樣做的時候,大家講著講著.
就將這些東西成為了.
一個很大的.
將小小的東西擴大.
以致那個恐慌越來越大.
在這裡裡面.
大致這個恐懼感.
變了由十二個門徒.
開始在船裡面.
擴大了出來.
耶穌見到.
這班門徒的時候.
他有什麼反應呢?在經文裡面.

$^{561}$就繼續說.
耶穌連望對他們說.
你們放心,是我不要怕.
那麼.
耶穌知道這班人真的害怕.
一來.
耶穌是一個很超自然的方式.
在水面上走.
都會嚇到那些門徒.
不過另一方面,那些門徒開始想過了.
而且亂想像.
所以耶穌急忙的.
向他們說.
走在水面上的是我.
你們不用怕.
值得留意的是那個字眼.
和合本已經很盡力了.
不過我們沒辦法.
看得到,你們放心.
原文譯就是,我是.
不要怕.
那個我是.
如果大家一起聽斯班唱詩.
我是葡萄樹.
耶穌在約翰福音裡面不停說.
我是.
七次的陳述.
說的是我是.
就是說耶穌在這個世代裡面.
他是神的兒子.
他要將他自己的身份.
表露出來.
這是一個很重要的神學傳述.
一個神學的宣告.
沒有人敢說我是.
在聖經裡面說.
耶穌是我是.
我是自有永有.
在舊約裡面說耶和華上帝.
這是一個身份的宣稱.

$^{601}$所以叫這班門徒.
不用擔心.
除了在澄清.
看到的不是幽靈之外.
耶穌也將他自己.
神聖的一面.
我是,來告訴一些門徒.
不過和合本盡力了.
如果他就這樣.
說你們放心.
我是,不要怕.
你們就不知道為什麼.
不過你剛才大概聽到我這樣說的時候.
是一個神學的.
一個專稱.
耶穌將他自己的身份告訴他們.
希望令到他們可以.
放下他們的疑慮.
可以放下他們的擔心和懼怕.
可以放心下來.
不過彼得和一班門徒不同.
他似乎不是很滿足.
這個答案.
因為凡是零都要驗證.
彼得可能很小心.
彼得就下一個條件句.
跟耶穌說.
主,如果是你的.
請你叫我從水面上走到你那裡.
彼得都挺大膽的.
他和其他門徒不同.
他不是純粹接受了答案.
他說,你說是你的.
如果,都是同樣的字.
如果你是,請叫我.
從水面上走到你那裡.
彼得用一個.
條件的方式.
他很想經歷.
耶穌的我事.

$^{641}$耶穌所說的.
是不是真的是耶穌呢?.
他用自己來教課.
他自己踏入水裡.
用他自己.
在急風大浪的海面上.
好像耶穌一樣走.
那就是夫子.
就是他們.
所熟悉的耶穌基督.
所以他就作出.
這樣的請求.
或者這樣的驗證.
來問耶穌.
一般來說,如果我們看.
我們會覺得大言不慚.
真是沒有大沒小的.
不過耶穌的反應.
就不是這樣的.
耶穌說,你來吧.
Come.
耶穌是接受.
彼得.
所作出的一個請求.
所以他邀請彼得.
去到這個.
海面上.
好像他一樣.
來行走.
耶穌說,你來吧.
當彼得聽到.
耶穌這樣說的時候.
同樣的字眼.
會重覆的.
彼得就從船下去.
在水面上走.
要到耶穌那裡去.
只因見風大浪大.
將要沉下去.
便喊著說,主啊救我.

$^{681}$在這裡裡面.
你見到那個句式.
開頭是沒有重覆的.
不過後半段又再重覆.
字眼.
害怕.
喊叫.
這個和之前門徒在船上.
見到耶穌在水面上行走.
說是一個鬼怪.
然後大家就驚慌.
害怕.
喊著說.
同樣的字眼.
在彼得這樣的驗證下.
再一次出現.
其實某程度.
耶穌已經說了他是誰.
彼得也和.
耶穌對話.
不過彼得就像.
自己拿來吹一樣.
是要親身去驗證.
於是乎.
就經歷了一個這樣的情景.
當他下到水面上.
見到風急浪大.
見風甚大.
就害怕了.
就沉下去.
於是就害怕.
就喊著說主啊救我.
在這裡面.
和之前的.
害怕喊叫不同.
這裡是今次是彼得.
自己主動去驗證.
這個信仰.
驗證眼前這個.
稱為夫子.

$^{721}$的耶穌.
他用自己.
去體會一下.
看看這個信仰.
看看這個耶穌是否真的.
他設下一個條件.
他就用自己.
如果能夠像耶穌.
在水面上行走的時候.
他就用這個條件.
用這個套路.
來去認證眼前這個.
就是耶穌基督.
不過今次的害怕.
是因為他不想.
想像是幽靈.
這一次的害怕.
是因為見到周圍的風高浪急.
見到周圍風雨甚大.
他見到周圍的環境.
實在是太過.
驚濤駭浪.
所以他自己就害怕.
就沉下去.
某程度上來說.
彼得.
他沉下去.
他沒有辦法在水面上行走.
可能能夠走一兩步.
不過馬太他強調的不是.
彼得走了幾步之後.
然後沉下去.
其實馬太不是這樣寫的.
馬太強調是.
彼得直接沉下去.
快要淹死了.
然後才叫.
然後耶穌就趕緊伸手拉住他.
你這個小信的人為什麼而惑呢.
某程度上.

$^{761}$彼得所設下的套路.
那個條件.
在他的經歷裡面.
是不成立的.
明白我的意思了嗎.
他說如果像你這樣能夠走到水面上.
我就知道你就是耶穌了.
不過結果他.
是走不到水面上.
他沉下去.
他後來被耶穌的手拉上來.
然後耶穌說.
他是小信的人.
為什麼而惑.
在耶穌.
應許他.
叫他過去的時候.
彼得其實他想的那個條件.
是成立不了的.
不過在他沉下去的時候.
最諷刺的地方就是.
他竟然.
向著這個.
他想要挑戰.
或者他要驗證的這個人.
主啊救我.
其實他自己已經認定了那個是主.
在他驗證的過程裡面.
其實都不需要那個答案.
在他這個經歷裡面.
他會去到耶穌的真實.
而這個真實更加是.
耶穌在水面上拉他上來.
讓他安然地.
回到那艘船裡.
在這裡裡面.
所以在這段幾節的經文.
不過你可以看到.
其實彼得.
想著用人的智慧.

$^{801}$用人的套路.
用人的方法來找.
找耶穌.
上帝似乎.
不會說你那個是錯的.
所以不讓他.
耶穌的姿態是你.
你去試一下.
不過在驗證的過程裡面.
他的害怕不是來自耶穌自己的本身.
而是來自周圍的環境.
風高浪急.
以致到他自己失去了信心.
他就擔心.
所以他就沉下去.
簡單來說.
令他沉下去的不是耶穌.
神聖或是耶穌真實的身份.
耶穌真實的身份.
不會因為人做的事.
而變得更加堅實.
自有永有.
我是,耶穌說的.
沒有東西可以證明耶穌.
反而是耶穌自己證明自己.
不過在這裡面.
當人用我們的智慧.
來想.
去找上帝的時候.
有時候我們可能面對周圍的環境.
有很多困難.
以致我們失去了.
尋求上帝的心.
結果我們就像彼得一樣.
沉下水.
不過沉下水不代表.
眼前的耶穌不是真的.
反而是在他沉下水的那一刻.
他更加經歷到.
是真實的耶穌.

$^{841}$是真實的神的兒子.
在這樣的親歷其境之下.
他親身體會到.
耶穌的實在.
接著在.
第32至33節.
當這群人.
上船.
彼得上船之後.
風就住了,在船上的人都拜他.
說:你真是神的兒子.
其實.
神的兒子這個字眼.
早就出現了.
這群門徒.
其實早就知道了.
某程度上.
在概念上他們.
明白耶穌是神的兒子.
這裡不是說耶穌說.
我是神的兒子,你看到了.
不是這樣的,是這群門徒.
他最終的結論.
他經歷了這些時刻的時候.
他們再一次.
強化他眼前.
這個神的兒子.
這個很重要的.
一個身份.
以前他們純粹放在腦子裡.
但經歷過這個晚上的時候.
這個神的兒子.
不是純粹一個腦的形式.
而是一個.
深深的經歷.
這段聖經大致.
剛才的脈絡或內容.
我解釋給大家聽.
大家理解了.
在這段經文裡有三個要點.

$^{881}$或許給我們一些啟迪.
我們值得再慢慢看下去.
第一.
人生裡有些恐懼.
這段經文特別處理一些恐懼.
他在說一些門徒的恐懼.
來自他們看到一些.
他們不知名的東西.
他們不想到那是耶穌基督.
不過自己人的想像力.
或者人自己之間的互動.
將一種恐懼.
不斷地擴散出來.
或者一路一路.
容讓這些恐懼.
佔據你的內心.
今天我們作為一個.
跟隨神的人.
其實耶穌基督.
在福音書也好.
或者整個基督教信仰也好.
其實祂要掃除.
人生裡很多不必要的恐懼.
不過有時候我們自己.
活在恐懼裡也不知道.
或者我們有時候.
覺得害怕就避開了.
保護了自己.
不過有些害怕是不必要的.
有些是自己嚇自己的.
又或者有些恐懼.
是來自童年有些陰影.
或者過去有些經歷.
這些東西就炒在一起.
厭惡感.
剛才我說過了.
厭惡感或者一些你不想經歷的事.
和一些形象.
炒在一起.
以致成為你生命裡.

$^{921}$一些很困擾你的恐懼.
如果你說小小的.
就對你人生沒什麼大影響.
不過如果有時候.
某些人的形象.
某一類人.
或者某一種的心態.
某一種處境.
令到你很窒息很害怕.
我們就需要正視這些恐懼.
如果這些恐懼不是正視它.
在你人生裡.
總會去到你低潮的日子.
這些東西就會浮出來.
就會令到你.
生命裡就好像癱瘓了一樣.
以致到沒辦法.
再三地.
去經歷到上帝生命裡的那種風聲.
可以這樣拋一點點書包.
有些弟兄姊妹.
曾經說過.
他有一些驚恐症.
其實某些驚恐症.
可能是來自過去成長裡.
有一些缺口.
又或者有些鮮為人知的秘密.
找東西蓋住.
很怕別人知道.
一種怕揭露自己.
讓人知道某一些秘密.
以致到遮遮掩掩.
後來這些東西就開始.
混雜了有不同的狀況.
就成為心裡一種很大的驚恐.
如果我們不處理.
我們純粹說.
吃個保陰丹就算了.
其實我們沒有解決到.
我們心層裡.

$^{961}$心靈裡的恐懼.
很多時直著禱告.
直著輔導.
是會將心靈裡的恐懼.
現形出來.
知道自己怕什麼.
厭惡什麼.
以致到我們可以對症下藥.
心靈的恐懼.
就能夠處理得到.
當心靈的恐懼處理了的時候.
你的人生就豁然開朗.
你就變得輕散和自由.
我們知不知道.
我們心靈裡有沒有什麼恐懼呢?.
值得用一些時間.
來處理它.
每個人都有的.
所以你不需要擔心.
我覺得自己好像.
有些人怕這些.
有些人怕那些.
每個人都會有的.
在屬靈生命裡.
如果能夠好好處理.
面對這些恐懼.
將這些恐懼能夠現形出來.
你就不會像是被一些.
無形的手控制著你的人生.
令你的生命.
動彈不得.
反而你知道的時候.
你就懂得怎樣去對付它.
這是第一點.
耶穌要直著.
祂自己說的「我是」.
又或者.
我們覺得好像.
彼得很衝動.
但是彼得不甘於純粹.

$^{1001}$是腦裡聽到答案.
他要親身來經歷.
他要面對著這種恐懼感的時候.
他自己親身下水.
來面對.
某程度上來說.
彼得是比我們更加勇敢的.
他這種面對恐懼的心.
到最後可能好像笑話一樣.
浸了下水.
要求耶穌救他.
不過耶穌在他生命裡面所作的.
就成為彼得屬靈生命裡面.
一個很堅實的一次的確據.
以至到他心悅誠服.
知道這個眼前的.
就是神的兒子.
他的生命而且有多一份的踏實感.
而其他那十一個門徒是沒有的.
這是第一點.
第二點就是.
繼續用彼得的角度來說.
去正視這種恐懼.
面對這個過程.
彼得下水.
不是彼得下水的時候.
變得風平浪靜.
風平浪靜是整件事結束了.
耶穌救他之後.
那個風才止住的.
但是彼得驗證的時候.
那個海面都是驚濤駭浪.
如果不是彼得也不會害怕的.
是不是?.
他去驗證耶穌的過程.
他要知道耶穌的我事.
原來那個過程都充滿了很多的挑戰.
這個正視恐懼的過程.
可能都會面對很多的困難.
甚至這些困難.

$^{1041}$可以壓倒性地讓它沉下去.
正視的過程.
可能心裡面.
可能將一些舊時的體會.
或者害怕的東西浮上來.
可能你會更差.
你會更害怕.
你可能更加會喊叫.
不過彼得的害怕和喊叫.
和之前門徒的害怕和喊叫不同.
他知道他眼前這個.
他將他自己的生命.
交付在這個主裡面.
在人生裡面經歷多少困難都好.
當我們去正視這個恐懼的時候.
耶穌一定會將手拉我們上來.
祂要讓我們經歷到.
祂生命裡面那種的真實.
我們只要願意.
好像彼得那樣.
去正視生命裡面.
這些不純粹是一個答案.
或者生命裡面有些恐懼.
他真的自己去體會的時候.
到那個過程可能不順利.
失敗.
但耶穌總是不會放棄我們.
反而祂會在我們最脆弱.
最危急的時候.
祂會拯救我們.
幫助我們.
所以弟兄姊妹.
生命裡面如果我們有些恐懼.
有些處理不到的事.
面對它的時候.
我們只管去到上帝面前.
將我們的心思意念.
將我們的一切交給祂.
上帝不會輕看.
上帝也不會忍心看著我們.

$^{1081}$繼續嚇死了下去.
或者不去面對.
反而是在風高浪急的海面上.
祂拯救彼得.
同樣也會拯救我們每一個.
當我們真的能夠去正視這個過程的時候.
我們的生命.
我們就會上多一課.
我們對神的信心.
就會加添多一層.
第三個層面.
剛才那些門徒說.
在船裡面.
試過幽靈.
不過耶穌真正的身份.
不是那個幽靈.
不是那個靈體.
他們最終的結論是.
神的兒子.
很多的宗教.
坊間裡面有不同的宗教.
除了說道人向善之外.
其實很多的宗教都要做一件事.
就是令人心裡面.
產生一種懼怕感.
你留意一下.
每一種宗教不是純粹.
讓你道人向善的一種美德.
它總會有些位置.
令到你會懼怕的.
當在聖經裡面.
我們反過來看.
聖經裡面說的那個懼怕.
從來不是說怕這個世代.
或怕我們功德做得不好.
我們真正裡面.
真正要敬畏.
聖經裡面用敬畏這個字.
只是神自己一個.
當我們知道我們生命裡面.

$^{1121}$只需要怕一個對象.
其他東西我們就變得.
不需要太過擔心,憂慮和不安.
因為我們找到生命裡面的核心.
再加上我們敬畏的神.
會怎樣對待我們呢?.
祂不是要操控我們.
或者令到我們痛不欲生.
或者在一個懼怕裡面.
繼續經營我們的信仰生活.
而是在敬畏的過程裡面.
上帝說的.
祂以永恆的愛.
愛你和愛每一個.
祂是一個愛的上帝.
這種愛驅除我們心裡面.
那些不必要的懼怕.
這個愛令到我們心靈可以踏實.
基督教和很多其他宗教不同.
基督教用很多的方式.
祂要告訴你.
我們只管將我們的心.
交付給這個很愛惜我們的神.
祂會將平安賜給我們.
祂不會令到我們變得神經質.
不停在很多事情上很害怕.
反而是我們經歷人生.
無論風風浪浪都好.
我們知道我們有神看著我們.
我們就可以踏實安穩.
我們就不需要用一種患得患失.
很害怕的心態.
來過我們的信仰生活.
弟兄姊妹.
今天這段經文不複雜.
我想大家甚至都耳熟能詳.
不過這段經文同樣挑戰我們.
我們今天心靈裡面.
會不會有些恐懼感.
籠罩著我們.

$^{1161}$我們要正視.
我們正視的方式不是跟別人不停說.
說完之後自己嚇自己.
越說越害怕.
而是要回到神的兒子身上.
正視自己生命的恐懼.
好像彼得那樣嘗試踏出水面裡面.
可能我們都會沉下去.
不過耶穌會拉我們上來.
他更加讓我們經歷到這個拯救.
耶穌在我們生命裡面.
無論我們怎樣面對處境.
多麼艱難.
他都要與我們同在.
他要讓我們知道.
心悅誠服.
發自內心裡面.
我們高歌的這位就是神的兒子.
求主繼續幫助我們.
在這個世道有很多東西嚇你.
不過讀完今天這段經文之後.
給我們多些踏實感.
讓我們生命不再變得成黃成紅.
而是有信心有平安.
走面前人生的路.
原主繼續帶領我們.
祝福我們.
\newpage



\section{雅各書 5:1-6}
\label{sec:LQ_YiQLwO5g}
\textbf{2021年9月26日崇拜直播｜陳冠成牧師｜生命見真章系列 - 追求真正富足｜雅各書5\_1-6}
\newline
\newline
連結: \href{https://youtube.com/watch?v=LQ_YiQLwO5g}{\texttt{https://youtube.com/watch?v=LQ\_YiQLwO5g}} ~~~~ 語音日期: 2021-10-21
\newline
\newline
\hyperref[sec:irpyKMaPiJI]{\small{< < < PREV SERMON < < <}}
~
\hyperref[sec:index_chronic]{\small{[返順時目]}}
~
\hyperref[sec:SlF159ccdOo]{\small{> > > NEXT SERMON > > >}}
\newline
\newline
雅各書 5:1-6
\newline
\begin{longtable}{cl}
\hline
\hline
章節 & 經文 (和合本修訂版)\\
\hline
5:1 & \begin{tabularx}{0.7\textwidth}{X} 注意！你們這些富足人哪，要為將要臨到你們身上的災難哭泣、號咷。 \end{tabularx} \\ \\ \relax
5:2 & \begin{tabularx}{0.7\textwidth}{X} 你們的財物腐爛了，你們的衣服被蟲子蛀了。 \end{tabularx} \\ \\ \relax
5:3 & \begin{tabularx}{0.7\textwidth}{X} 你們的金銀都生鏽了；這鏽要證明你們的不是，又要像火一樣吞吃你們的肉。你們在這末世只知道積蓄錢財。 \end{tabularx} \\ \\ \relax
5:4 & \begin{tabularx}{0.7\textwidth}{X} 工人給你們收割莊稼，你們剋扣他們的工錢；這工錢在喊冤，而且收割工人的冤聲已經進入萬軍之主的耳朵了。 \end{tabularx} \\ \\ \relax
5:5 & \begin{tabularx}{0.7\textwidth}{X} 你們在地上享奢華宴樂，把自己養肥了，等候宰殺的日子。 \end{tabularx} \\ \\ \relax
5:6 & \begin{tabularx}{0.7\textwidth}{X} 你們定了義人的罪，把他殺害，他沒有抵抗你們。 \end{tabularx} \\ \\
[1ex]
\hline
\hline
\end{longtable}
$^{1}$今日的經訓是記載在新約《雅各書》五章一至六節.
在你們手上的聖經是在三百二十四節與三百二十五頁.
上帝的話語這樣說.
你們這些富足人啊 應當哭泣,豪圖.
因為將有苦難臨到你們身上.
你們的財物壞了 衣服也被從子咬了.
你們的金銀都長了瘦.
那瘦要證明你們的不是.
又要吃你們的肉 兒童火燒.
你們在這末世只知積攢錢財.
工人給你們收割莊稼.
你們虧欠他們的工錢.
這工錢有聲音呼叫.
並且那收割之人的怨聲.
已經入了萬軍之主的耳了.
你們在世上享美福 豪而樂.
當宰割的日子竟驕讓你們的心.
你們定了二人的罪 把他殺害.
他也不抵擋你們.
謝謝各位評判.
很快亞國書的系列只剩下兩講.
連同今次 下星期明珠牧師會講最後一講.
不過在分享之前.
就剛才我們所聽到的經文.
我必須問大家一個問題.
你覺得你有錢嗎?.
你覺得你富有嗎?.
立即搖頭了.
有誰敢承認自己有錢呢?.
所以我不敢叫大家舉手.
因為應該沒有人敢.
還有你剛才聽的經文是怎樣的?.
好像衝著有錢人而來.
死都不承認.
你會看到亞國書用了一些很嚴厲的字眼.
整個亞國書最嚴厲的都是這六節.
來責備一些富有的人.
表示他們需要哀哭.
哭號.
甚至要面對一些很艱難的結局.

$^{41}$經文裡面提到將要面對被宰殺的日子.
有苦難淋到他們身上.
好像火燒一樣.
不過我們明白.
上帝並不憎恨有錢的人.
我也明白.
上帝也樂意祝福.
將他的福氣賜給富有的人.
譬如我們很熟悉的亞伯拉罕.
或約伯或大衛.
我們看到上帝很願意賜福給富有的人.
不過我們也明白.
新約裡面提到貪財是萬惡之根.
上帝憎惡貪心貪愛錢財的人.
如果你有留意.
特別在科學書裡面.
耶穌對於貪心貪財.
所發出的警告其實不少.
有些英學者甚至說.
所發出的警告遠多於犯奸淫的警告.
提醒我們一件事.
其實我們沒有一個人能夠誇口.
說自己對貪婪的試探有免疫.
否則耶穌也不需要說那麼多次.
不過很有趣的是.
我做了牧者十多年.
我發覺有些姐妹會因為.
憤怒的緣故來找我聊天.
可能會因為婚姻的緣故來找我聊天.
因為驕傲的緣故來找我聊天.
不過我從未遇過.
有些姐妹說有貪婪的問題要處理.
不知道你覺不覺得.
貪念似乎不是那麼容易被人察覺出來.
話說美國有一個很出名的神學家.
叫做Timothy Keller.
中文譯做提摩泰凱勒牧師.
他有次分享他曾經舉辦過.
七場關於罪的講座.
他的太太很快就跟他說.

$^{81}$你相信我,一提到貪婪那一場.
肯定是最少人聽到的.
你相信嗎?.
結果果然.
關於情慾方面的罪.
關於憤怒的罪.
關於驕傲的罪都坐無虛次.
唯獨沒有人會認為自己是貪心.
為何貪婪的罪那麼容易被察覺出來呢?.
提摩泰凱勒牧師就繼續分享.
其實從社會文化角度來看.
金錢偶像這個試探.
很聰明.
他很懂得在社會上不同的經濟階層.
將自己包裝隱藏起來.
事實上你會明白.
無論你有沒有錢.
你身邊一定會有人比你更富有.
比你更奢華.
一定會找到.
你一比較就找到.
所以在比較之下.
每個人總會覺得.
我現在的生活質素.
我現在擁有的財富.
比起某某其實都只不過是甚麼.
小事而已.
都只不過是小無見大無.
於是我們都會自然產生一種想法.
就是其實我不是貪的.
我只不過是想怎樣.
我想多賺一點.
我想儲多一點.
我想享受多一點.
合理嗎?.
很合理.
不過份了.
我又不犯法.
沒錯的.
即使其實他的生活已經相當富裕.

$^{121}$甚至乎是奢華.
他依然會有這個想法.
正如《雅各書》所說.
他正是要向一班貪婪.
貪愛財富的人發出一個嚴厲警告.
很希望他們能夠及時.
回頭,醒悟過來.
你看看經文.
到底他們貪婪貪財.
去到一個甚麼地步呢?.
簡單來說有三件事.
第一方面,第二節都說了.
因為他們是儲的.
即賺他們的財物.
即賺他們的金錢.
第二節說到他們的財物壞了.
衣服也被蟲子咬了.
這裡特別說財物壞了.
其實是甚麼呢?.
很大機會是說他們儲糧食.
因為在當時古代世界.
有糧食充足.
其實是一個財富的象徵.
將糧食儲起來.
去到一個地步.
多到根本吃不完.
最後就會.
可能像米過期.
被穀牛鑄成.
最後變壞.
寧願擺在這裡.
爛掉也不拿來分給有需要的人.
另一方面,他們積存的衣服.
也多到怎樣呢?.
當然是穿不完.
最後被蟲子鑄成,沒用.
這個情況我們今天也有.
你打開衣櫃.
也有不少衣服.
可能你未穿過,未剪牌.

$^{161}$甚至可能穿過一兩次.
我們也積存過不少衣服.
最後可能因為.
折舊了.
或者可能身形不同.
不合身,最後要怎樣呢?.
好一點就轉贈.
又或者要丟掉.
所以其實.
現在製造衣服的製造商.
用環保物料.
就預了怎樣呢?.
預了你浪費.
預了你丟掉.
可能你的衣服.
如果是用環保物料製造.
放進衣櫃.
你會發現.
幾年後,突然有些粉末出現.
已經開始分解.
免得你放進堆填.
它分解不到.
所以物料已經可以自行分解.
你看到其實.
這是第一方面.
另一方面,你又見到.
亞國形容這班貪婪的人.
有錢又多到沒地方用.
好像廣東話說.
他們的金紙銀紙.
長期擺到生鏽.
其實原則上.
金和銀很難生鏽.
這裡原本的意思是什麼呢?.
其實似乎是說.
鋪滿塵.
自然會怎樣呢?.
沒有光澤.
不反光.
好像生鏽了.

$^{201}$亞國說這些鏽.
將來會成為審判他們的罪證.
指證這班人.
雖然很富有.
不過不務正業.
他們的人生只顧積賺財寶.
不過其實實際累積著什麼呢?.
上帝對他們的憤怒.
這是亞國對他們的嚴厲指控.
另一方面.
亞國再說.
第五節.
這班貪婪的人.
經常都怎樣呢?.
開派對.
經常搞派對.
吃最好的東西.
喝最好的紅酒.
享受生真海味.
總之我有權.
我有錢.
我想怎樣享受都可以.
這是他們做人的態度.
亞國形容他們其實是什麼呢?.
驕養自己的心.
意思是什麼呢?.
其實是說他們在上帝面前.
其實不懂得死.
只顧著放縱享樂.
養肥自己.
卻不察覺怎樣呢?.
最終養肥了的自己.
原來是被誰割的呢?.
等著被上帝割.
事實上不難理解.
我想活在今天.
我們很強調個人消費主義.
這個世代.
我們很崇尚的是怎樣呢?.
我們有自由.

$^{241}$我們有權去運用.
我們所賺的錢.
做自己喜歡的事.
爽不爽?.
其實真的挺爽的.
我可以駕馭自己所擁有的.
我的錢喜歡怎樣花怎樣用.
是我的自由.
有它的吸引力.
是真的.
但是你想深一層.
聖經裡面說.
當人在享樂的時候.
將上帝完全抽走的時候.
其實就會出現亞國書.
這裡有一個所說.
貪得無厭.
放縱自己的人.
最終失去的不單止是.
他所擁有的財寶.
更是失去了他的生命.
失去了他和上帝的關係.
第三方面.
亞國他斥責.
有一班貪婪的人.
他們自稱是「貪婪」.
\newpage



\section{列王記下 5:1-14}
\label{sec:SlF159ccdOo}
\textbf{2021年10月31日崇拜直播｜陳冠成牧師｜這樣就能得醫治｜列王記下5\_1-14}
\newline
\newline
連結: \href{https://youtube.com/watch?v=SlF159ccdOo}{\texttt{https://youtube.com/watch?v=SlF159ccdOo}} ~~~~ 語音日期: 2021-11-09
\newline
\newline
\hyperref[sec:LQ_YiQLwO5g]{\small{< < < PREV SERMON < < <}}
~
\hyperref[sec:index_chronic]{\small{[返順時目]}}
~
\hyperref[sec:wlHRHi438v8]{\small{> > > NEXT SERMON > > >}}
\newline
\newline
列王記下 5:1-14
\newline
\begin{longtable}{cl}
\hline
\hline
章節 & 經文 (和合本修訂版)\\
\hline
5:1 & \begin{tabularx}{0.7\textwidth}{X} 亞蘭王的元帥乃縵在他主人面前是一個偉大的人，得王的喜悅，因為耶和華曾藉他使亞蘭人得勝。他雖然是大能的勇士，卻染上了痲瘋。 \end{tabularx} \\ \\ \relax
5:2 & \begin{tabularx}{0.7\textwidth}{X} 亞蘭人成群出征的時候，從以色列地擄了一個小女孩，她就服事乃縵的妻子。 \end{tabularx} \\ \\ \relax
5:3 & \begin{tabularx}{0.7\textwidth}{X} 她對女主人說：「我希望主人去見撒瑪利亞的先知，他必能治好主人的痲瘋。」 \end{tabularx} \\ \\ \relax
5:4 & \begin{tabularx}{0.7\textwidth}{X} 乃縵去告訴他主人說，從以色列地來的女孩如此如此說。 \end{tabularx} \\ \\ \relax
5:5 & \begin{tabularx}{0.7\textwidth}{X} 亞蘭王說：「你可以去，我也會送信給以色列王。」於是乃縵手裡帶十他連得銀子、六千舍客勒金子和十套衣裳去了。 \end{tabularx} \\ \\ \relax
5:6 & \begin{tabularx}{0.7\textwidth}{X} 他帶著這信給以色列王，說：「現在你接到這信，看哪，我派臣僕乃縵到你這裡來，你要治好他的痲瘋。」 \end{tabularx} \\ \\ \relax
5:7 & \begin{tabularx}{0.7\textwidth}{X} 以色列王讀了信就撕裂衣服，說：「我豈是神，能使人死使人活呢？這人竟派人來，叫我治好一個人的痲瘋。你們要知道，看，這人是找機會來跟我吵架的。」 \end{tabularx} \\ \\ \relax
5:8 & \begin{tabularx}{0.7\textwidth}{X} 神人以利沙聽見以色列王撕裂衣服，就派人到王那裡，說：「你為甚麼撕裂衣服呢？讓那人到我這裡來，他會知道以色列中有先知。」 \end{tabularx} \\ \\ \relax
5:9 & \begin{tabularx}{0.7\textwidth}{X} 於是乃縵帶著車馬到了以利沙的家，站在門前。 \end{tabularx} \\ \\ \relax
5:10 & \begin{tabularx}{0.7\textwidth}{X} 以利沙派一個使者，對乃縵說：「去，在約旦河中沐浴七次，你的肉就必復原，你會得潔淨。」 \end{tabularx} \\ \\ \relax
5:11 & \begin{tabularx}{0.7\textwidth}{X} 乃縵卻發怒走了。他說：「看哪，我以為他必定會出來，到我這裡，站著求告耶和華－他神的名，在患處上搖手，治好這痲瘋。 \end{tabularx} \\ \\ \relax
5:12 & \begin{tabularx}{0.7\textwidth}{X} 大馬士革的亞瑪拿河和法珥法河豈不比以色列的一切水更好嗎？我難道不可以在那裡沐浴而得潔淨嗎？」於是他生氣，轉身走了。 \end{tabularx} \\ \\ \relax
5:13 & \begin{tabularx}{0.7\textwidth}{X} 他的僕人近前來，對他說：「我父啊，先知若吩咐你做一件大事，你豈不做嗎？何況是吩咐你去沐浴，得潔淨呢？」 \end{tabularx} \\ \\ \relax
5:14 & \begin{tabularx}{0.7\textwidth}{X} 於是乃縵下去，照著神人的話，在約旦河裡浸了七次。他的肉復原，好像小孩的肉，他就潔淨了。 \end{tabularx} \\ \\
[1ex]
\hline
\hline
\end{longtable}
$^{1}$經文:舊約,四百六十一頁.
亞蘭王的元帥乃曼在他主人面前為尊為大..
印耶和華曾直他使亞蘭人得勝,他又是大能的勇士,此事長了大麻風..
先前亞蘭人成群的出去,從以色列國擄了一個小女子..
這女子就服侍乃曼的妻,她對主母說:.
「巴不得我主人去見撒瑪利亞的先知,必能治好他的大麻風.」.
乃曼進去告訴他主人說:.
「以色列國的女子如此如此說:.
「亞蘭王說:你可以去,我也達信於以色列王.」.
於是乃曼帶銀紙十他年得,金紙六千射克勒,以上十套就去了..
且帶信給以色列王,信上說:我打發臣服乃曼去見你,你接到這信,就要治好他的大麻風..
以色列王看了信,就撕裂衣服,說:.
「我豈是臣能使人死,使人活呢?.
這人竟打發人來,叫我治好他的大麻風..
你們看一看,這人何以沉駒攻擊我呢?」.
神人以利沙聽見以色列王撕裂衣服,就打發人去見王,說:.
「你為什麼撕了衣服呢?」.
那時乃人到我這裡來,他就知道以色列中有先知了..
於是乃曼帶著車馬到了以利沙的家..
站在門前,以利沙打發一個侍者對乃曼說:.
「你去在約旦河中沐浴七回,你的浴就必復原.」.
疑得結井,乃曼卻發怒走了,說:.
「我想他必定出來見我.」.
站著後告耶和華他神的名,在幻處以上搖手,治好這大麻風..
大麻式的河,阿巴拿和法爾法,豈不比以色列的一切水更好嗎?.
我在那裡沐浴,不得結井嗎?.
於是氣紛紛的轉身去了,他的僕人進前來對他說:.
「我父先知若吩咐你作一件大事,你豈不作嗎?何況說你去沐浴,而得結井呢?」.
於是乃曼下去,照著神人的話在約旦河裡沐浴七回,.
他的浴復原,好像小孩子的浴,他就結井了..
乃曼帶著一切跟隨他的人,回到神人那裡,站在他面前的說:.
「如今我知道,除了以色列之外,普天下沒有神,現在求你修點僕人的禮物.」.
聖經讀筆,願神賜福常讀他話語並且遵行的人..
丙姐妹平安!.
剛才玉峰弟兄跟我們讀乃曼的大麻風得到醫治,我想大家都很熟悉吧?.
特別可能有小朋友正在上兒童主義學,都會聽過這個故事..
這個故事基本上由兩組對比鮮明的人物帶出一個拯救訊息..
有哪兩組人呢?第一組是來自社會頂層的強者,包括哪些人呢?.
包括主角,乃曼元帥,經文一開始就說,他曾經替亞蘭人,亞蘭即是今天敘利亞一帶的地方,.
他曾經替亞蘭人打贏過仗,所以深得亞蘭王賞識..

$^{41}$這位我們說有戰績,戰績標榜,他加上地位顯赫,並且聖經裡說他是一個很有能力的勇士,.
不過他卻患了大麻風,這種在當時很難能夠醫好,並且具有很高傳染性,需要隔離的一種惡疾..
並且這個病會將乃曼元帥他擁有的成就漸漸的失去,甚至乎他的性命可能也會因為這個病而逐漸的摧毀瓦解..
另外,兩個來自頂層的強者,分別就是亞蘭王和他鄰國的以色列王,.
當時亞蘭王是比以色列王強勢,你看他的信便知道,所以這兩個王很特別,他們的關係一時會結盟,一時又會打仗,.
但無論他們多麼有權有勢,原來在這件事上他們幫不了乃曼,幫不了乃曼擺脫他的大麻風..
這三個強者就是故事中其中一個組別.而另一個組別則是來自社會基層的僕人,包括哪些人呢?.
首先出場的是乃曼妻子的婢女,另外還有乃曼的僕人和以利沙的僕人..
這三個僕人在當時可以說是無權無勢,但很特別的是他們在乃曼得醫治這件事上產生了一個很關鍵的作用..
我們首先看看乃曼妻子的僕妻,她是一個被亞蘭人擄走的以色列小女孩,她知道主人的病情,於是便跟她所服侍的主母說,.
巴不得主人可以去薩馬利亞,當時以色列國的首都,找到那裡的先知便可以了,當地的先知一定能夠醫好主人的大麻風..
你不難想像,在得到這個消息之前,乃曼應該以她的人脈和權勢,她已經試了很多方法,但是沒有效..
所以當她一聽到有一個這樣的消息,竟然在以色列有一個人可以治好自己,所以經文說她立即和亞蘭王稟告,我們說所謂申請什麼?出外就醫?.
亞蘭王批准,並且寫信向以色列王,你看到這封不是拜託人的信,是一封給壓力人的信,向以色列王施壓,大概意思是你接到這封信後,你給我放聰明點,你給我跟緊點這件事,我現在打發乃曼來找你,你要幫她身上的大麻風完全康復..
於是乃曼就帶著這封信,連同經文說,銀紙有十他年得,金紙有六千四赫勒,還有十套衣裳,就是這樣去找以色列王,這是故事的開頭..
不過先停一停,做一點我們說好像看完電影有些影評,大家不知道有沒有發覺到底整件事是不是有人說漏了些東西,還是有人聽錯了些東西?.
為什麼這樣說呢?你看到從乃曼和阿南王的對話,他們只提到什麼人?提到以色列王,提到以色列國,他們隻字不提要去找先知幫忙,但是一開始明明那位很卑微的婢女已經說了什麼?.
乃曼你只要去薩馬利亞見先知,就可以得到醫治,很簡單很直接.換算不是乃曼,而是其他得到大麻風的平民百姓一聽到這個消息就會怎樣?.
就會很懂得自然向這個婢女追問下去,你快點告訴我這個先知叫什麼名字?他住在哪裡?請告訴我,讓我可以立刻直接去找他幫忙,對不對?.
其實我們今天也不難理解,假如你咳嗽了很久,看著也不斷尾,於是你身邊突然有個朋友說,你這個情況我知道有什麼辦法,我認識一個醫生在旺角,他一定可以處理你的病,因為他已經幫很多人醫斷尾部的咳嗽,你聽到一定會怎樣?.
你就會很自然順便追問,醫生的診所在哪裡?電話多少號?他收不收新診?價錢多少?還有哪天休息?假如你真的不懂得路,你會直接問對方,可不可以告訴我怎樣去?.
對不對?我們不需要花時間去找,我們當然想盡快得到醫治,所以你看到奶媽明明應該去找誰?去找以利沙,但為什麼她會去錯地方,又找錯對象?她去了找以色列王,為什麼整件事你看回頭,明明很簡單的事情變得迂迴複雜了?.
這方面很值得我們去思考,但其實不難理解,我們今天很明白這個道理,就是所謂識人比識字好,能夠找到高層直接幫你出面,這件事一定會快捷很多,辦得妥當很多,.
所以你看到亞蘭王也好,奶媽也好,其實他們不會在乎以色列的先知是誰,住在哪裡都不要緊,因為他們習慣在他們的世界裡,我身為高層,當然是和高層對話,所以我直接去找以色列國的最高主事人,然後怎樣?給他壓力,不夠的話再用錢訂他,.
吩咐他替這件事辦妥,這樣就最有效,你看到亞蘭王和奶媽他們用的方法,其實和今天他們所奉行的遊戲規則都很相似,只要夠才容世大,他就可以將對手,甚至被他虐的人,統統壓低,逼他就範,來滿足自己任何的需要,.
所以你會看到在亞蘭王也好,奶媽也好,在他們人生的成功字典裡,有幾個字是找不到的,找不到請求,軟弱,屈就,謙虛,代表弱者的詞語,很難在他們的成功字典裡找到,.
通常找到的是什麼字呢?我要怎樣施壓,威逼利誘,強制,講排場,這些字就多數在他們的字典裡找到,事實上,十他年得的銀子,其實是一個什麼概念呢?.
在當時可以買到什麼呢?原來根據《聖經》列王記尚,他記載昔日另一個以色列王,阿摁尼,他要用多少錢才能買起整座薩馬利亞山呢?.
以色列首都,以色列首都的金紙,要多少錢才能買起整座薩馬利亞山呢?用兩他年得就行了,現在出多少倍的價錢?出五倍價錢,亞蘭王用以色列首都地價的五倍價錢收買以色列王,如果嫌不夠的話,不要緊,.
我再給你六千四百萬的金紙,即是七十公斤的黃金,你聽到七十公斤的黃金,我不知道多少錢,我馬上要上google裡面查一查,在現在的市值是超過三千萬港幣,現在,你想想,在當時的購買力是一億億的,.
這麼多錢,很難會不心動,很難會不跪下的,作為以色列王,但是怎知道,當以色列王看到他這一封信的時候,他激動到撕爛衣服,他很生氣,很委屈,但是又沒有辦法,.
他講了幾句話,我不是上帝,我怎能夠治好奶萬的大麻瘋,所以他覺得整件事不是有心找他幫忙治病,而是怎樣?他直接說,亞蘭王你是有心留難我的,你只不過是借奶萬這個病來撩交打,找藉口來攻擊我,因為我根本沒有能力幫到奶萬,.
你看到整件事很特別的就是,奶萬一心想著靠他那套成功的哲學,可以帶來醫治,殊不知,這個想法卻是被以色列王一句話就戳爆了,就是原來錢和權力都幫不了你,只有上帝幫到你,只有上帝可以治好你,.
其實你看回整件事,確實他又講得對,真是上帝才可以幫到他,並且上帝不單單有能力幫到奶萬,其實上帝願不願意幫助奶萬?很願意,你看回第一節,誰讓奶萬打勝仗?是上帝,上帝一早已經介入了奶萬的生命裡面,.
只不過奶萬不知道,所以你看到第一節,奶萬能夠替亞蘭人打勝仗是出於上帝的幫助,但當然上帝期望奶萬的不單是得到病,得醫治這個好處,上帝要奶萬得到的是你要認識誰才是真神,.
所以你看到上帝首先要拆除的是奶萬一直奉行的人生價值觀,以為靠著大人物,靠著龐大的財力,勢力就能夠成就任何大事,包括醫病這件事,.
上帝卻說即使是你眼中看為那些無權無勢,很卑微不起眼的小人物,上帝說我一樣有辦法用他來幫到你,來替上帝成就大事,話說我聽過一個很有趣的故事,.
有一個患了嚴重失眠的朋友,他看了醫生也改善不了,所以如果你有長期失眠的習慣,就知道很纏擾,睡也睡不著,於是他沒有辦法,唯有回教會,看看上帝能否幫到他,於是他第一次參加崇拜,.
當日崇拜的詩歌很喜樂,令整個崇拜氣氛高漲,但這位朋友發現坐在他前排的一位很年長的弟兄,他居然像嬰孩一樣睡著了,多嘈的音樂詩歌也無法吵醒他,於是這位患有嚴重失眠的朋友立即說,.
哇,喜出望外,原來信耶穌可以這麼好睡,於是他就這樣信了耶穌,後來他講到這個見證的時候,聽到其他弟兄姊妹就說,其實你不知道,你不是第一個,這個弟兄很厲害,他之前就是這樣釣釣釣,幫耶穌釣了幾個人信主,.

$^{81}$不過我說這個故事不是鼓勵大家在崇拜中釣魚,而是想說,我們很多時候會有想法就是受到社會崇尚一些所謂英雄的主義,哪些人來拯救你,不是那些差的人,一定是那些厲害的人才能救你,.
我們受一些成功主義的影響,我們有時候總會覺得自己太過平凡,又沒有什麼能力,又沒有什麼才華,資源又不足,很難為上帝成就大事的,我們都會這樣看自己,但是弟兄姊妹上帝不是這樣看你和我,.
其實這段經文特別反映,上帝很多時候是樂意透過一些平凡起眼的人和一些方法來成就大事,達成上帝一些的心意,所以當你繼續看下去的時候,雖然乃曼沒辦法,以色列王也沒辦法,但是上帝有辦法,.
上帝透過在第八節裡面,神人以利沙聽到以色列王撕裂衣服,就打發一個僕人去見以色列王,問他為什麼撕裂衣服,不用這麼受,可以使那個人來到我,即是以利沙的家,他就會知道以色列中間有先知,.
其實以利沙的僕人宣告這句話,其實和第三節婢女所說的一樣,只不過重新說,雖然乃曼錯過了婢女提供給他的提示,搞到自己又去錯地方,又去找錯人,錯過了醫治的機會,但是上帝有辦法將他帶回來正路,.
上帝是一位有恩典有憐憫的神,他給第二次機會乃曼,並且透過一個很平凡的僕人,一個不起眼的僕人,將乃曼帶回正路,所以你看到今次第九節,乃曼終於去對地方,找對人了,.
他來到以利沙家門前,他帶著居馬,相信也帶著一大批金銀和衣服,來到以利沙的門口,不過沒有想到他又再遇到另一次挫敗,今次是怎樣呢?.
乃曼一心想著,他已經構思好劇本,我畢竟是一個有頭有臉的人,我帶著一隊車隊,帶了這麼多金銀珠寶準備送給以利沙,我這麼屈就,這麼犧牲自己,以利沙你也應該怎樣?.
你也應該出來迎接我,他起碼覺得要這樣,但你竟然派一個僕人來見我,那就是不給面子,所以乃曼非常憤怒,不過這樣也算了,他預計以利沙會像他自己所說,.
我估計他也會站在我面前求告,然後說他上帝的名,然後在我身上的大麻風患處做些動作,搖搖手,治好我的大麻風,大概他曾經在亞蘭國自己的地方也見過亞蘭國的先知,也是這樣用類似的形式幫人治病,所以他已經有一個前設,以利沙起碼也應該要這樣..
但現在以利沙吩咐僕人叫乃曼做些什麼?這裡文雅一點,沐浴七次,你回去洗澡七次,這樣你就會康復,在乃曼的角度他認為,你只是擺明耍我,敷衍我,如果這麼簡單,洗個澡,洗幾次澡就康復了,.
我為什麼不去那些河水,水質比較好的,位於亞蘭大麻色的河流,我為什麼不去那裡沖,那裡的河水比你們以色列的河水多多聲,我為什麼要專門來到你面前受你的氣,你只是存心留難我,你不治就算了,不幫就算了,我走..
弟兄姊妹,我們又停一停,想一想為什麼乃曼這麼憤怒,甚至憤怒到放棄醫治,憤怒到連命都不要,為什麼這麼憤怒?.
你看到表面上是因為他滿懷希望,卻遭到冷待,他覺得很受委屈,所以氣憤難平,但是我們不妨細心想,之所以這麼憤怒,真正的原因恐怕是來自他內心那種非常不安的感覺..
為什麼這麼說呢?因為乃曼已經不知道自己可以做什麼,他覺得自己已經做不到什麼,能做的都做了,接下來還能做什麼?.
試想想,如果以利沙先知不是吩咐他沐浴七次,而是說你想我治好你,可以,加多一倍診金,他會不會給?沒問題.又或者說你想我治好你,可以,你幫我替以色列國打兩座城池回來,有沒有問題?.
沒問題,為什麼?因為金錢,權力,能力,一向都是他很傻的範疇,他是習慣用金錢,能力,財幹解決他所有遇到的問題,所以你要他做多一點沒問題,絕對沒問題..
所以為什麼沐浴對乃曼來說這麼困難,這麼難接受,這麼煩躁?困難在於什麼?困難在於太容易,太容易,有沒有想過?.
他這麼煩躁是因為你說的太容易,有沒有可能這麼容易就行?你叫我去洗澡洗澡七次這麼容易的事,不需要任何財幹,不需要我付出任何努力,甚至連小孩子,連最弱的人都辦得到的事,你要我去做的話,其實你就等於要我相信,乃曼,其實你的病你什麼都不用做,你站在這裡花期等等,過兩天會好起來..
所以為什麼他會這麼不安,這麼煩躁?不難理解的,弟兄姊妹有沒有試過看鐵打?有沒有試過?可能都有,通常去鐵打有什麼程序?.
你會看那個醫師,可能會開一缸,如果你扭傷腳,可能會浸一缸,開一缸藥酒給你浸,暖暖的,浸完之後,可能會搓一下,甚至會針灸一下,或者敷藥,有一系列的程序,通常會不會痛?.
都會有的,所以如果你很介意痛,你不會看鐵打,有一系列的程序,但假如有一天看鐵打,那個醫生很有名的,你去試一下,但他只是開了一缸藥酒給你浸完,可以了,擦乾淨之後,花期,明天來浸,浸七天你會好起來,你會有什麼感受?你還信不信?.
你可能會好像奶慢,如果這麼簡單,我為什麼要特地來這裡做這麼簡單的事?所以,弟兄姊妹,在奶慢的心中,假如她真的去沐浴,其實她就要承認,我是很無能,我是很無助,我是很脆弱,並且她要接受一個她一直都不習慣的事情,就是原來看病是免費的..
原來得著上帝的醫治是免費的,無能,軟弱,免費,這三個字基本上不會出現在她那本成功的字典裡面,是一些對她來說既陌生又害怕,不懂得怎樣面對的詞語,.
所以為什麼她要發怒,掉頭走,來掩藏她內心的恐懼和不安.真的,弟兄姊妹,你會看到奶慢厲害嗎?有沒有才幹呢?絕對有,並不是上帝不起願這些事情,.
不過在這件事上,她的才幹和成就正正是難了她經歷上帝醫治最主要的障礙.事實上,可能有時候我們都會經歷過,有時候人越厲害就越難接受一些事情,不知道你有沒有發覺..
有時候人越厲害就越難接受為什麼那件事不是按著我的心意發生,為什麼我遇到的際遇不是我想像奶慢那樣,我明明人生那麼美好,怎麼會出現大麻風,有我搞不定的事情,我明明以為以莉莎會出來禮待我,為什麼她好像羞辱我一樣,為什麼事情不是照著我的意思做,為什麼這個世界不是按著我的心意走,很難接受..
又或者很難接受為什麼我傾盡全力都擺平不了我身邊的困難,我很難接受,我動用了那麼龐大的財力物力,還拜託了亞蘭王施壓,原來都治不好我的病,為什麼會這樣?.
我都很難接受做那些不肖一顧的事情,我竟然要聽一個比我低很多級的僕人的吩咐,還要做一件我覺得那麼愚蠢,那麼沒意思的事情,我覺得很討厭,我很難接受..
更加難接受的是,幫我治好我的病的人,他說我一個腺都不收你,很難接受,因為在我的世界裡面,你知道我要付出多少代價才能坐到元帥的位置,我要犧牲多少才能達到今時今日擁有的成就,這個世界上不是凡事都有一個價錢嗎?.
更何況是將我起死回生,為什麼你說一個腺都不收我?所以,弟兄姊妹,你會看到,從乃曼身上我們開始會明白,為什麼有時候越厲害的人越難信耶穌,越厲害的人越難接受福音的真理,因為總是想著上帝不可能白白將救恩賜給我的..
一定是這個救恩太便宜,是一個廉價的恩典,所以免費送給所有人都無所謂,卻是不知道實情,上帝的拯救代價實在太貴,太高昂,貴到一個地步,有沒有人可以付得起?沒有..
有沒有人可以有資格獲得?都沒有.但是,上帝又知道每一個人都需要這個救恩,需要上帝拯救,所以那怎麼辦呢?上帝唯有替我們付,上帝唯有藉著耶穌基督在十字架上釘生受死,透過這個當時所有人認為最卑微的僕人,.
用一個最卑微,沒有人看得起的方法,來替世人付上贖價,為所有人帶來醫治和拯救,雖然很多人不願意接受這個真理,甚至聽完之後,可能好像賴慢一樣,掉頭就走了..
不過,這段經文有一個好消息讓我們看到,人放棄了,上帝有沒有放棄?上帝沒有放棄過,上帝又多猜了另一個僕人,今次到第三個僕人出場,是賴慢自己的僕人..
第十三節,賴慢的僕人上前對主人賴慢說:「我父啊,若先知吩咐你作一件大事,你豈不作?何況說你去沐浴而得潔淨呢?」.
這個僕人很厲害,他一眼就看穿了主人的內心,他知道主人的心態就是大人物只會做大事,所以他很有智慧地用幾句說話將主人引導回去,回轉,我們說所謂..
事實上,你留意他這幾句說話,有時我們用來勸勉別人的時候,心態和技巧是值得我們學習的..
首先,你看到他如何動之以情,他稱呼賴慢做什麼?「我父啊」,即是說什麼?我是用一個父子的關係來關心你,我不是要跟你說教,我是用父子情來關心你,我想你好,我不是說你差,我厲害,我想你好..
這是第一方面,第二方面,你留意到他說的話,我知道先知吩咐你做一件大事,你其實是很願意配合的..

$^{121}$這個僕人很懂得運用一個發問的技巧來喚醒賴慢的初心,他來求什麼?求醫治,其實什麼方法原則上他都會配合..
他很懂得透過發問技巧,將很憤怒,很有情緒的賴慢,冷靜下來,你是來求醫治的,我知道你會配合..
然後他再說,何況是目欲這麼小的事情,背後的意思是什麼?沒有說出來,不過是說,怎麼會難到你呢?你一定做得到..
其實他很厲害的地方是他沒有迫他主人去做決定,他製造了一個空間,主人你自己想,到底轉不轉頭,去不去若但何,主人你知道應該怎麼做..
他是等賴慢你自己心甘情願,你自己調頭,你自己改變主意.這個僕人他很明白賴慢的盲點就是這個世界哪有這麼簡單就能得醫治,所以我怎麼都不試..
於是這個僕人就很厲害,他改了一點字,他反過來邀請賴慢跳出框框想一想,如果這個世界這麼簡單就能得醫治的話,你為什麼不試?賴慢就是我怎麼都不試,這麼簡單..
僕人就說這麼簡單你為什麼不試?轉一轉框,跳出框框想一想,反正做了都沒有壞,為什麼你不換一個角度去看整件事,說不定有新的體驗呢?.
這就是這個僕人幾句說話的厲害之處,當然最厲害的就是上帝可以透過一個平凡的僕人幾句有愛心有智慧的說話,就幫賴慢這個大人物走回近上帝..
所以最後想請大家讀14至15節,印在程序表裡面,最後那段落14至15節,預備起..
謝謝,你看到這次是大團圓結局,賴慢最終選擇依從以利沙的僕人的吩咐去做,結果上帝讓他經歷的,是否只是病得醫治這麼簡單?.
不是,你知道雖然我們沒有患過大麻瘋,但大麻瘋就算醫好,身體上仍然會留下什麼?疤痕,現在呢?現在沒有了,並且他的肌膚比以前好,好像小孩子一樣..
你看到上帝讓這個選擇信服遵行的賴慢,他經歷的不是純粹得到醫治,而是得到什麼?重生.重生的不只是他的肉體,更是他回復到一個小孩的心,一個柔和謙卑的心..
因為原本高傲自負的他,現在竟然願意在以利沙面前,降伏謙卑地承認以色列的上帝才是獨一的真神.不單止他的身體能夠得到復原,賴慢和上帝的關係也復原了..
所以賴慢向我們呈現的不只是一個病得醫治的故事,其實是他在傳遞一個福音好訊息給我們,真正弟兄姊妹.你相不相信有時信耶穌得醫治,耶穌的幫助就是這麼簡單?.
我認識一對基督徒夫婦,其中一位來找我,他說其實我婚姻一直很纏繞,因為結婚十多年,關係不太好,我們也有試過找牧者介入,見過父道,但總是處理不到,沒有什麼改善..
他來到我面前,大家談如何做好,十多年了,醫治不好,其實我也沒有什麼辦法,不過我知道,跟他說既然今天我們仍然能夠會面,代表上帝仍然在,上帝仍然沒有放棄你的婚姻,所以我們今天才能坐在一起談這件事..
換句話來說,上帝覺得還可以拔,還可以處理,我們要有信心和盼望,先處理好我們跟上帝的關係..
另一方面,我們開始談,他很想很快問一些相處技巧,如何夫婦改善溝通,先放下這些,先去依靠那位設立婚姻,賜福守望婚姻的上帝,勝過我們去依靠夫婦相處的技巧,並不是那些不重要,但要先回到最根本的那件事..
首先,你要做的是納至每天親近上帝,原因是你要是一個成熟的人,才會是一個成熟的情人,你要是一個成熟的基督徒,才會成為一個成熟的基督徒,或者成熟的基督徒夫婦,丈夫或者妻子..
他就說,眼睛有些疑惑,就是這麼簡單?真的可以這樣做,就能夠搞好夫婦關係?我就反過來問,為什麼上帝的醫治一定要很複雜?為什麼上帝的醫治一定要很複雜?或者說,如果這麼簡單,真的可以搞好,為什麼你不試?反正試了也沒有壞處..
再進一步,我就說,先不要想怎樣改變對方,只想怎樣改變自己,還是想怎樣依靠上帝改變你的脾氣,改變你眼中的良睦,並且你為對方代徒就行了..
靠上帝也要感動他,預備他的心,讓他可以走近上帝.我們不能明白夫婦和上帝的三角關係,上帝在上面,夫婦在下面,如果夫婦都走近上帝,他們自然怎樣?走近,縮小了,三角形縮小了,他們就近上帝,他們也會自然近上帝,對不對?.
如果相反,如果只我走近上帝,他不走近,那怎麼辦?他會不會越來越遠?會不會?我說不會的,你走近上帝,你住在上帝的裡面,上帝會主動找他,上帝不是那麼被動,永遠都等人來找他,上帝會主動尋找拯救失散的人,你住在上帝的裡面,上帝自然會帶你慢慢走近他,你離不開他..
道理很簡單,信不信?擠不擠?並且能不能堅持下去?坦白說,我不知道事後的發展會是怎樣,但我確信一件事,聖經裡面說的,上帝不會虧待每一個願意謙卑尋求親近他的人,你怎樣都會得著上帝,上帝也會有最好的安排..
於是乎就是這樣,經過幾年,上帝又派不同的牧者,瀑人來介入,開始有成效,開始大家復和..
弟兄姊妹,其實講這個故事就像今天的經文,邀請我們,我們會否都像奶饅那樣,幻想了一些我們很難治好的大麻瘋呢?可能是你的性格,是你的脾氣,一些陋習,甚至乎一些上癮的問題,又或者你和身邊人的關係,我們總是靠自己,搞不定,治不好,一直很影響,很纏繞你和身邊的人..
假如上帝透過今天的經文,邀請你和我,去約旦河沐浴七次,你相不相信?你願不願意?當然,不是真的叫我們坐飛機去聖地沐浴七次..
其實,沐浴七次有兩個很重要的意義,就是七,我們很明白,你要信服上帝是完全的大靈,完全的信服,不要和上帝角力,不要和上帝爭話事權,事實上不容易的,因為我們看病,很多時候都看兩個醫生,有沒有發覺?.
第一個就是看醫生,第二個醫生是誰?是自己,自己會做醫生,有沒有試過這些經驗?看完醫生,醫生叫你戒口,叫你不要吃這些藥,叫你準時吃藥,然後你就會按自己的判斷,加加減減,你會做醫生,所以就不斷尾..
看醫生可能真的可以這樣,但在屬靈的事上,上帝說你要一心一意降服,叫沐浴七次,就沐浴七次,所以在藥彈河沐浴七次,另一個重點就是你要有耐性等候上帝的工作,什麼意思呢?就是,奶漫沐浴一次,兩次,三次都沒有效,那你會怎樣?還沐浴嗎?沐浴到第六次都沒有效,那你就會開始懷疑,第七次行不行?.
繼續嗎?正如我們一直尋求上帝,但都沒有效的時候,我們還繼續嗎?沐浴第七次嗎?這個就是經文給我們的挑戰和提醒,即使你沐浴到第六次都沒有效,都不要放棄,不要抱怨,你要相信上帝是在工作,並且你更加要相信上帝是愛我們,祂不會丟下我們,祂一定有救用的恩典和我們同行..
世上頂指的屬靈生命的意志和成長,正正就是考驗我們,你願不願意放下你的幾件聖件,一些執著,甚至一些能力和尊嚴,有時還要放棄你的理性,來專一尋求上帝的帶領,順服祂的吩咐,你又願不願意接受上帝今天向你發出的邀請,去藥彈河沐浴七次,.
讓上帝改變,醫治你的生命,讓我們得到真正的復原,讓我們同心禱告,祈求上帝,願你幫助我們,不單止聆聽,更加反省實踐經文的吩咐,.
上帝事實上,作為一個罪人,我們總會在人生裡,有一些我們無論如何都搞不定的大麻風在我們身上,我們很想靠自己的能力來醫好他,又或者我們將他放在一旁,不理他,但是仍然一直纏繞著我們的成長,我們和身邊人的關係,上帝盼望你藉著今天這段經文,開我們的心竅,讓我們知道上帝從來沒有放棄過我們,.
並且上帝你總有辦法,透過身邊人的提醒,透過簡單一句話的挽回,將我們走迷了的我們帶回到你的身邊,主也幫助我們,我們知道無論我們身上有多難醫好的惡疾也好,上帝你有辦法,並且你的辦法出乎我們所想,可能是很簡單,關鍵是我們願不願意相信你,信靠你,並且堅持到底..
主啊,求你藉著這段經文,幫助我們的生命能夠成長,幫助我們在生活裡面,並不單單是得著你的福氣,更重要的是得著上帝你自己,總是相信信服上帝你一切所吩咐的話,願你聽我們在你面前的禱告,奉基督耶穌你,得勝的名字,你都可以祈求,阿們..
謝謝大家收看 再見.
\newpage



\section{創世記 39:1-23}
\label{sec:wlHRHi438v8}
\textbf{2021年10月24日青少主日崇拜直播｜梁煒傳道｜當不幸敲門時｜創世記39\_1-23}
\newline
\newline
連結: \href{https://youtube.com/watch?v=wlHRHi438v8}{\texttt{https://youtube.com/watch?v=wlHRHi438v8}} ~~~~ 語音日期: 2021-11-21
\newline
\newline
\hyperref[sec:SlF159ccdOo]{\small{< < < PREV SERMON < < <}}
~
\hyperref[sec:index_chronic]{\small{[返順時目]}}
~
\hyperref[sec:ec0Lp80ZOxQ]{\small{> > > NEXT SERMON > > >}}
\newline
\newline
創世記 39:1-23
\newline
\begin{longtable}{cl}
\hline
\hline
章節 & 經文 (和合本修訂版)\\
\hline
39:1 & \begin{tabularx}{0.7\textwidth}{X} 約瑟被帶下埃及去。有一個埃及人波提乏，是法老的官員，是護衛長，他從那些帶約瑟下來的以實瑪利人手中把約瑟買了去。 \end{tabularx} \\ \\ \relax
39:2 & \begin{tabularx}{0.7\textwidth}{X} 約瑟在他埃及主人的家中，耶和華與他同在，他是一個通達的人。 \end{tabularx} \\ \\ \relax
39:3 & \begin{tabularx}{0.7\textwidth}{X} 他主人見耶和華與他同在，又見耶和華使他手裡所辦的事都順利， \end{tabularx} \\ \\ \relax
39:4 & \begin{tabularx}{0.7\textwidth}{X} 約瑟就在主人眼前蒙恩，伺候他主人，主人派他管理家務，把一切所有的都交在他手裡。 \end{tabularx} \\ \\ \relax
39:5 & \begin{tabularx}{0.7\textwidth}{X} 自從主人派約瑟管理家務和他一切所有的，耶和華就因約瑟的緣故賜福給那埃及人的家；凡家裡和田間一切所有的，都蒙耶和華賜福。 \end{tabularx} \\ \\ \relax
39:6 & \begin{tabularx}{0.7\textwidth}{X} 波提乏把他一切所有的都交在約瑟手中，除了自己所吃的食物，其他的事一概不知。約瑟英俊健美。 \end{tabularx} \\ \\ \relax
39:7 & \begin{tabularx}{0.7\textwidth}{X} 這些事以後，約瑟主人的妻子以目送情給約瑟，說：「你與我同寢吧！」 \end{tabularx} \\ \\ \relax
39:8 & \begin{tabularx}{0.7\textwidth}{X} 約瑟拒絕，對他主人的妻子說：「看哪，一切家務我主人一概不知，他把所有的都交在我手裡。 \end{tabularx} \\ \\ \relax
39:9 & \begin{tabularx}{0.7\textwidth}{X} 在這家裡沒有人比我更大，除你以外，他也沒有留下一樣不交給我，因為你是他的妻子。我怎能行這麼大的惡，得罪神呢？」 \end{tabularx} \\ \\ \relax
39:10 & \begin{tabularx}{0.7\textwidth}{X} 她天天這樣對約瑟說，約瑟卻不聽從她，不與她同寢，也不和她在一起。 \end{tabularx} \\ \\ \relax
39:11 & \begin{tabularx}{0.7\textwidth}{X} 有一天，約瑟進屋裡去辦事，家裡沒有一個人在那屋子裡， \end{tabularx} \\ \\ \relax
39:12 & \begin{tabularx}{0.7\textwidth}{X} 婦人就拉住他的衣服，說：「你與我同寢吧！」約瑟把衣服留在她手裡，逃出外面去了。 \end{tabularx} \\ \\ \relax
39:13 & \begin{tabularx}{0.7\textwidth}{X} 婦人看見約瑟把衣服留在她手裡逃到外面， \end{tabularx} \\ \\ \relax
39:14 & \begin{tabularx}{0.7\textwidth}{X} 就叫了家裡的人來，對他們說：「看，他帶了一個希伯來人到我們這裡戲弄我們。他到我這裡來，要與我同寢，我就大聲喊叫。 \end{tabularx} \\ \\ \relax
39:15 & \begin{tabularx}{0.7\textwidth}{X} 他聽見我放聲大喊，就把他的衣服留在我這裡，逃出外面去了。」 \end{tabularx} \\ \\ \relax
39:16 & \begin{tabularx}{0.7\textwidth}{X} 婦人把約瑟的衣服放在身邊，直到他主人回家， \end{tabularx} \\ \\ \relax
39:17 & \begin{tabularx}{0.7\textwidth}{X} 就用這樣的話對他說：「你帶到我們這裡來的那希伯來僕人進來要調戲我， \end{tabularx} \\ \\ \relax
39:18 & \begin{tabularx}{0.7\textwidth}{X} 我放聲大喊，他就把衣服留在我身邊，逃到外面。」 \end{tabularx} \\ \\ \relax
39:19 & \begin{tabularx}{0.7\textwidth}{X} 主人聽見他妻子對他說的話，說：「你的僕人就是這樣對待我」，就非常生氣。 \end{tabularx} \\ \\ \relax
39:20 & \begin{tabularx}{0.7\textwidth}{X} 約瑟的主人把他抓起來，關在監獄裡，就是王的囚犯被關的地方。於是約瑟在那裡坐牢。 \end{tabularx} \\ \\ \relax
39:21 & \begin{tabularx}{0.7\textwidth}{X} 但耶和華與約瑟同在，向他施恩，使他在監獄長的眼前蒙恩。 \end{tabularx} \\ \\ \relax
39:22 & \begin{tabularx}{0.7\textwidth}{X} 監獄長就把監獄裡所有的囚犯都交在約瑟手下；在那裡的一切事都由他處理。 \end{tabularx} \\ \\ \relax
39:23 & \begin{tabularx}{0.7\textwidth}{X} 任何交在約瑟手中的事，監獄長一概不察，因為耶和華與約瑟同在，耶和華使他所做的都順利。 \end{tabularx} \\ \\
[1ex]
\hline
\hline
\end{longtable}
$^{1}$今日經訓中的經文在《創世記》三十九章一至二十三節.
在舊約聖經四十九至五十頁.
《創世記》三十九章一至二十三節.
請聽我讀出.
若失被大廈埃及去.
有一個埃及人是法老的內臣.
護衛長波提弗從那些大廈內的二十瑪利人手下.
馬了他去.
若失住在他主人埃及人的家中.
耶和華與他同在.
他就百事順利.
他主人見耶和華與他同在.
又見耶和華使他手裡所辦的盡都順利.
若失就在主人眼前蒙恩.
侍候他主人.
並且主人派他管理家務.
擺一切所有的都交在他手裡.
自從主人派若失管理家務.
和他一切所有的.
耶和華就因若失的緣故.
賜福與那埃及人的家.
凡家裡和田間一切所有的.
都蒙耶和華賜福.
波提弗將一切所有的都交在若失手中.
除了自己所吃的飯.
別的事一概不知.
若失原來受雅俊美.
這事以後.
若失主人的妻.
以沐送情給若失.
說:你與我同枕吧.
若失不從對他主人的妻說:.
看那一切家務.
我主人都不知道.
他把所有的都交在我手裡.
在這家裡沒有比我大的.
並且他沒有留下一樣不交給我.
只留下了你.
因為你是他的妻子.
我怎能作這惡大惡得罪神呢?.

$^{41}$後來他天天和若失說:.
若失卻不聽從他.
不與他同枕.
也不和他在一處.
有一天若失送屋裡去辦事.
家中人沒有一個在那屋裡.
婦人就拉著他的衣裳說:.
你與我同枕吧.
若失把衣裳雕在婦人手裡.
跑到外邊去了.
婦人看見若失把衣裳雕在她手裡.
跑出去了.
就叫了家裡的人來.
對他們說:.
你們看.
他帶了一個希伯來人.
進入我們家裡.
要戲弄我們.
他到我這裡來.
要與我同枕.
我就大聲喊叫.
他聽見我放聲喊起來.
就把衣裳雕在我這裡.
跑到外邊去了.
婦人把若失的衣裳.
放在自己那裡.
等著他主人回家.
就對他如此如此說:.
你所帶到我們.
這裡的那希伯來人.
還來要戲弄我.
我放聲喊起來.
他就把衣裳雕在我這裡.
跑出去了.
若失的主人聽見他妻子.
對他所說的話說:.
你的僕人如此如此待我.
他就生氣.
把若失下在監裡.
就是黃的囚犯被囚的地方.

$^{81}$於是若失在那裡坐監.
但耶和華與若失同在.
向他施恩.
使他在施玉的眼前蒙恩.
施玉就把監裡所有的囚犯.
都交在若失的手下.
他們在那裡所辦的事.
都是經他的手.
凡在若失手下的事.
施玉一概不測.
因為耶和華與若失同在.
耶和華使他所作的.
重刀順利.
頂子妹平安.
平安.
多謝一班少年的敬拜隊.
師班的獻送.
求神繼續去興起他們.
期望他們越事奉上帝.
越認識上帝.
越去愛上帝.
每次我們見到年輕人一起去敬拜.
帶領我們一起去.
去敬拜上帝的時候.
我們都覺得很有力.
很有盼望.
很美好.
自己有一段時間沒有在大堂裡面講道.
可能當中有些新朋友未必很認識我.
我叫做梁偉.
是教會青少年的傳道人.
我預備了一張照片去介紹自己.
這張是我自己的全家福.
我太太也是一個傳道人.
我有兩個很可愛的小朋友.
如果幸福指數100分是滿分的話.
我猜我應該有120分.
不過.
就是.
在我們這段時間裡面.

$^{121}$家裡正在經歷一個很大的變化.
八月份的時候.
因為我的小女兒大約一歲.
但是她發展緩慢.
我們就帶她去看醫生.
因為她不懂得走.
不懂得坐.
還有爬也不太懂得.
看醫生的時候.
醫生說她患了一個很罕有的病.
可能將來沒辦法正常地生活.
我們夫婦都來不及消化這個消息.
太太已經要立即辭去教會的工作.
全職陪伴著小女兒經常出入醫院.
在醫院裡面.
我們都看到很多病患的兒童.
很多父母都忙碌照顧著自己的小朋友.
如果大家多留意新聞.
其實新聞最近好像多了很多交通的意外.
一場的交通意外.
破壞了一個原本美好完整的家庭.
失去經濟支柱.
失去很重要的親人.
人生的路上的確不是平坦.
不是一帆風順.
人生不如意事.
十成八九.
一個突如其來的改變.
都會令到我們措手不及.
這兩年的裡面.
我們正在經歷很多的變化.
社會運動,疫情.
對我們的生活有很大的影響.
有不少人離鄉別井.
希望尋找一個更合適的家.
不單止香港,緬甸,阿富汗等國家.
他們都面臨突如其來的改變.
令到他們陷入在戰爭混亂的裡面.
今年的暑假.
我們青少年部辦了一個營會.

$^{161}$其中一個晚會裡面.
我們邀請青年人去默想.
如果他們的生命經歷一些傷痕.
不幸,難過,傷痛.
他們就在他們的手裡面.
用一支紅筆畫一畫.
代表著你們所經歷的傷痕.
然後我們邀請他們去到台前.
用耶穌基督的寶血洗清他們這些傷痕.
很意外.
一個十幾歲的年輕人.
有些手是佈滿了很多紅色的痕.
究竟一個年輕人.
他生命裡面經歷著什麼傷痕呢?.
是來自他的童年,家庭,學校.
還是生活呢?.
今天我們需要去學習的.
可能是如何在逆境.
如何在困難裡面.
仍然繼續走下去.
今天這段經文是記載在.
《創世記》第39章.
約瑟生命裡面一段故事.
也是他生命裡面最重要的一個轉捩點.
當時約瑟只是十七歲.
正在經歷被家人出賣,流落異鄉.
究竟這個年輕人如何面對著不幸呢?.
當他賣到埃及一個人的時候.
失去一切.
他是不是還相信上帝是賜福.
上帝是和他同在.
什麼原因令到他繼續走下去呢?.
當不幸臨到的時候.
你又會如何去面對呢?.
讓我們一起祈禱.
親愛的天父.
我們承認我們的生命是脆弱.
是軟弱,是無力.
唯有依靠你,捉緊你.
和你一起去同行.

$^{201}$我們才能夠有帕望.
求你的聖靈臨到我們當中.
讓我們認識,明白上帝你的說話.
要將你的說話.
採在我們心裡面.
教導我們,指引我們.
讓我們走在你心意裡面.
聽我們在你面前祈禱.
奉耶穌基督得勝之名而求.
阿們.
經文我會把它分成三個部分.
第一個部分是一至到第六節.
不如邀請大家一起去讀第一至第二節的經文.
預備,一二三.
約瑟生長在一個大家庭.
他的父親是雅各.
雅各有四個老婆.
總共生了十二個兒子.
這個就是以色列的十二國之派.
拉傑是雅各第二任的太太.
他所生的就是約瑟和卞雅敏.
但拉傑在生卞雅敏的時候去世.
所以約瑟從小就沒有媽媽的照顧.
而他媽媽只生了兩個小朋友.
他們在家裡是最小的.
第三十七章裡面.
雅各偏愛約瑟.
為他做了一件菜衣.
引起了其他十個兄長的不滿.
加上約瑟經常在父親面前告狀.
將他們的惡行告訴父親.
約瑟又發了兩個夢.
兩個夢告訴他.
意思就是十個兄長和他父親都要向他下拜.
因為種種各樣的原因.
這十個兄長很不喜歡他.
對他懷恨在心.
於是很想設計陷害他,殺害他.
不過最後約瑟就被賣到以色列人裡面.
再帶到埃及的地方.

$^{241}$在第三十七章約瑟當時被兄弟賣到埃及.
他那個年紀是十七歲.
第四十章約瑟成為埃及的宰相.
為王解夢.
當時他的年紀是三十歲.
在這十三年的裡面.
約瑟分別在波提弗的家裡做奴隸.
也在監獄裡成為囚犯.
當約瑟經歷被出賣.
成為奴隸.
其實當時約瑟一定很難過,很傷心.
可能有很多不同的情緒.
傷心,難過,憤怒,不滿,迷惘.
甚至乎很多埋怨在裡面.
但在不幸的裡面.
他失去一切.
他會不會懷疑上帝.
上帝的拯救,上帝的賜福.
好像在那一刻他裡面全部都不見了.
他也不知道他將來會成為一個國家裡面的宰相.
不過我想約瑟唯一能夠知道的是.
他需要好好地活下去.
約瑟用三個不去面對著他的光景.
首先他不活在過去.
原本他的過去萬千寵愛在一身.
有爸爸的寵愛,有彩衣,有工人.
吃的,穿的都是尚好.
過著無憂的童年生活.
不過一夜之間.
經歷十個兄長的出賣.
從人生的高峰跌到低谷.
原本天之驕子卻是成為奴隸.
沒有家人,沒有彩衣,沒有人侍候.
約瑟被帶到埃及.
一個言語不通.
膚色,種族完全不一樣的地方.
外邦人的裡面.
可能約瑟每一天都希望.
這些事情不是真的.
希望自己做夢.

$^{281}$希望回到以前的生活.
有些人面對不幸.
可能都很希望這件事情沒有發生.
如果這件事情發生是多麼好.
如果沒有疫情是多麼好.
我們可以到處去旅行.
如果沒有社會事件多麼好.
有些朋友和某些家人.
繼續暢所欲言.
如果爸爸沒有患病.
如果那天沒有上酒吧.
如果沒有離開教會.
如果我聽媽媽的話.
這些事情種種就不會發生.
不過原來當人繼續活在這些傷痛.
活在過去的時候.
沒有辦法走出去.
走不出這個困局.
既然事情已經發生.
是沒有辦法回頭.
學不,約瑟不怨天尤人.
約瑟的生活生長在一個.
很複雜的家庭裡面.
媽媽自小離開.
得不到母愛.
身邊十個的兄長.
都有爸爸媽媽.
他們都比自己大.
比自己有能力.
童年裡面的約瑟.
他經常將哥哥的惡行.
告訴兄長和爸爸.
用破壞的方式.
表達他裡面的不滿.
他得不到幸福.
也不想其他人得到.
約瑟使用這些方法.
能夠得到父親的注意.
父親的寵愛.
不過當他被賣到埃及的時候.

$^{321}$他也沒有將這些罪咎.
歸咎在其他人的身體裡面.
他反而找到工作.
得到波提弗的賞賜.
有些人面對著不幸.
會將這些責任感.
這些責任和結果.
歸咎在其他人的裡面.
將不滿發洩在其他人身體裡面.
經常怨天尤人.
尋找機會報仇.
當全世界都是他的敵人.
覺得全世界做錯.
和他為敵.
他們的口頭禪就是.
最壞的都是他們有問題.
最壞的都是爸爸小時候經常打我.
最壞的都是家裡沒有錢.
最壞的都是老師教書教得不好.
最壞的都是某某高官的策略.
最壞的都是老闆這樣針對我.
最壞的都是老婆和我的關係.
最壞的都是朋友經常慫恿我.
他們會將自己應付的責任.
完全射落到其他人的裡面.
在他們口中經常喊喊聲.
沒有任何表情.
甚至很憤怒.
很多裡面的不滿.
這些不滿.
只會增強自己不幸的感覺.
第三個不.
弱失沒有自暴自棄.
雖然受到打擊.
受到人生裡面沒有失去了任何一切.
不過弱失沒有自我放棄.
沒有自殘.
沒有放棄他的人生.
繼續積極工作.
做好每一份的工作.

$^{361}$有些人面對著一個不幸的疑慮.
打擊他們的信心.
一蹶不振.
為什麼命運偏偏選中我.
認定自己是不幸者.
深信自己沒有將來.
認定他是人生裡面的輸家.
選擇逃避成為隱性.
終日留在家裡打機度日.
過著放縱自我.
微難的生活.
和損友一起圍爐取暖.
最終作樂.
不願意面對現實.
花大量的時間打機.
在IG Facebook YouTube Netflix.
甚至在不良場所裡面.
其實他們不想面對著.
這個殘酷的現實.
他們可能有些口頭禪就是.
唉!幹嘛?有什麼用?.
用不著這麼認真.
放鬆點,放鬆點吧.
我喜歡這樣,怎樣?.
沒可能的了.
沒用的了.
為什麼要這麼努力?.
沒用的了,不用管我.
你任由我吧.
不斷的放負.
發放負能量.
走近他裡面.
你心情都會覺得一沉.
的確,放負能量很容易.
不過對事情沒有幫助.
我們唯一能夠做的.
是積極的面對著這個困難.
挑戰和不幸.
所以弱失.
他首先要做的是.

$^{401}$他正面的去處理.
第一個要處理的是.
他自己的情緒.
面對著不幸.
傷心難過是一定的.
不過有時候我們基督徒.
有時候都會很快跳過這一步.
認為基督徒一定要喜樂.
一定要笑面迎人.
面對著傷心難過的事情.
而認為傷心流淚都是不好.
其實內心的感受.
一定要抒發.
上帝創造人.
是很特別.
他給我們喜怒哀樂.
這是身體機制.
是在調節我們的生命.
是上帝給我們一份禮物.
有時候我們沒辦法向前.
可能是因為我們有一些感受.
未去處理.
我們需要正視.
自己內心的感受.
有一本書叫做.
培育高EQ的靈命.
原來一個成熟的基督徒.
必須要有一個成熟的情緒的處理.
面對著傷痛.
我們其實需要放聲大喊.
如果可以的話.
找一個你值得信任的人.
放聲大喊.
傳達人牧師.
永遠都是你最好的朋友.
開放他們的門.
去和大家分享.
傳達人的房間裡經常有一個垃圾桶.
裝著很多紙巾.
不是因為冷氣冷.

$^{441}$而是因為這些紙巾不是傳達人.
而是盛載著很多弟兄姊妹的眼淚.
當然教會裡面的弟兄姊妹.
你的家人.
你的朋友.
我們的上帝.
永遠都是隨時樂意.
聽我們一切的困難.
面對著困難.
不要自己面對.
要和別人分享.
你會發現.
原來我們共同面對著這些掙扎.
第二個我們要學習的是承認失敗.
接受現在的光景.
既然事情已經發生.
我們唯一能夠做的是.
防止事情再次發生.
和防止事情變得更加差.
因此我們需要檢視自己的生命.
特別是處理罪的問題.
如果有毒癮.
走癮.
甚至淫亂的問題.
要立刻停止.
處理心裡的苦毒.
學習要書.
正視上帝.
發現恩典.
靠著上帝的恩典.
改掉這些壞習慣.
當弱失被罵去波提佛的家裡的時候.
他都要學習做好份內事.
不去打其他小報告.
爭鳴奪利.
第三.
當人面對著大災難.
大問題的時候.
有時候都會質疑自己.
覺得自己沒用.

$^{481}$對於弱失來說.
從前萬千寵愛在一身.
十指不沾揚春手.
有工人服侍他.
不過今天他成為一個奴隸.
要服侍其他人.
要去日曬雨淋.
在外面做奴隸的工作.
體力勞動.
簡單重複的工作.
不需要技巧.
不過當弱失慢慢掌握這些技巧的時候.
他被提升到屋裡面.
做一些技巧性的工作.
當波提佛又見他在屋裡面做得好的時候.
升他成為一個大內總管.
管理著整個官邸的事務.
在整個過程裡面.
弱失的技能不斷去提升.
從小事.
從簡單事開始.
他發現自己有能力.
有才幹.
自信心慢慢被建立.
得到其他人的信任.
可以做更多的事.
每次的不幸.
其實提醒著我們.
我們過去的一套方法.
未必再適用.
我們需要去調整,微調.
再次去重新出發.
當上帝安排我們的功課和環境裡面.
我們學習到的時候.
上帝會賜福給我們.
結果上帝賜福給弱失.
讓他步步高升.
得上司.
得同事.
得其他人的喜悅.

$^{521}$在剛剛暑假裡面有一齣戲.
這齣戲的名字叫做《媽媽的神奇小子》.
改編自香港的神奇小子.
殘奧運金牌得獎者蘇華偉.
和他媽媽熱血的故事.
蘇華偉在剛剛出生的時候.
就患上了黃蛋病.
導致他腦裡面的受損.
注定一世都會躍聽.
和四肢勁亂的殘障人士.
只能夠靠其他人的照顧.
但這個蘇媽媽沒有放棄.
訓練他蘇華偉學習怎樣去走路.
更發現原來他有跑步的天賦.
蘇華偉自小就加入了殘障的田徑隊.
在1996年的時候.
更加參與殘障奧運會裡面.
為香港帶來歷史性的第一個金牌.
蘇華偉五度出戰殘障人士奧運會.
和各樣的比賽.
獲獎無數.
他更加是殘奧男子100米200米.
的世界紀錄保持者.
媽媽的神奇小子裡面.
有一幾句金句想和大家分享.
輸在起跑線不要緊.
最重要知道終點在哪裡.
沒有人當你是普通人.
那你就做一個不普通的人.
雖然你一出生就是輸.
不過不代表你會輸一世.
第二個部分39章7到第18節.
我又邀請弟兄姊妹幫我讀7到第12節.
可能有少一場.
我邀請弟兄讀7節.
姐妹讀第8節相似.
預備123弟兄.
(粵語).
其實約瑟已經很厲害.
面對這個逆境.

$^{561}$他努力去做好.
不過當他努力的過程裡面.
不代表試探不會去淋到他生命裡面.
這個試探從來沒有離開過他.
聖經形容約瑟瘦瓦俊美.
瘦瓦俊美即是說他的樣子帥.
他的體格強健.
高大衛猛英俊瀟灑.
這個高大衛猛英俊的外表.
是引起波提弗的妻子的吸引.
聖經經文裡記載了這三次的試探.
第七節.
(粵語).
當代聖經的譯本是.
波提弗的妻子看中了約瑟.
向他眉目全情.
要引誘他上床.
看中是貪婪的注視.
是要得到人或事.
認定那樣東西是他的獵物.
不會走得開.
其實舊約聖經裡面.
很少人物像這個婦人那樣.
不知廉恥.
公然向男人表達.
跟他上床.
只要性不要愛.
波提弗的妻子.
沒有用甜言蜜語.
只是下了一道命令.
對於年輕貌美的約瑟來說.
這是一個誘惑.
約瑟十幾二十歲.
他血氣方剛.
自己一個孤身寡人.
心理上其實也有需要.
而且波提弗身居要職.
他的妻子是身份的象徵.
其實一定很漂亮.
加上當時奴隸制度的.

$^{601}$道德標準很低.
性開放是很隨便.
只要他一開口.
約瑟一定要聽從.
信奉波提弗的妻子.
能夠令自己有情慾.
以及在工作裡面也有好處.
若果得罪主母的話.
後果可能不堪設想.
約瑟同樣面對權力的試探.
面對與波提弗一樣的試探.
約瑟在官邸裡.
其實他已經位高權重.
波提弗將一切交給他.
除了他的食飯和妻子.
這個權力的誘惑.
就是與波提弗一樣.
當他得到波提弗的妻子.
就能夠與他一樣.
第二個試探更恐怖.
是天天的試探.
第十節後來.
他天天和約瑟說.
其實真的很恐怖.
這個妻子每天告訴他.
你跟我同寢吧.
每天都說.
如果有十天.
即是他說十次.
這十次裡面.
只要有一次站立不穩.
就會泥足深陷.
沒辦法自拔.
第三個試探也是最危險.
就是獨個兒的試探.
一個人的時候.
第十一節.
有一天約瑟進屋裡辦事.
家中沒有一個人.
在那個屋裡面.

$^{641}$家中沒有一個人.
強調的是沒有人.
有一天.
色經學家去解釋.
可能是波提弗的妻子.
精心設計好的一個局.
屋裡沒有一個人.
妻子用力抓住約瑟的衣裳.
用引誘的方法不行.
唯有強行和他同寢.
沒有人的試探.
是最危險的試探.
因為沒有人知道.
你想做甚麼就做甚麼.
可能你會想.
我都沒有高大衛猛的外表.
我也有穩定的關係.
可能這些試探.
都未必淋到我的身上.
不過其實這些試探.
不需要在特定的場景.
特定的人物.
才能夠發生.
這些試探.
每天都在我們的電話裡面.
發生中.
網絡的世界.
每天在引誘我們.
一個電話.
一個程式.
一個手指.
我們可以不受地區的限制.
很簡單.
很方便.
進入一個淫亂犯罪的當中.
很多交友的社交媒體.
總是容易有一些.
裸露的鏡頭和視頻.
只要輕輕一按.
成千上萬的圖片和視頻.

$^{681}$隨意入下載.
當你瀏覽完.
沒有人知道.
早前有段新聞.
提到TikTok.
即抖音裡面的危害.
特別是對年青一代的危害.
未成年一代的危害.
是很嚴重.
這裡面實在有太大量.
暴力色情不良的一些訊息.
有些人會戒手.
亦有些人會有裸露的鏡頭.
小朋友.
即中學生.
根本沒有辦法去意識到.
陷入這些試探裡面.
當年青人花大量的時間.
在這些軟件裡面的時候.
其實他需要的是.
我們從旁的指引和分辨.
因此面對這些試探.
他失去他的態度.
他的思想.
他的行為.
他能夠堅定.
持守他的界線.
首先弱失做的是.
訂立清楚一條線.
這條線他怎樣都不可以越過.
面對著主母很厲害的試探.
用他的眼睛.
用他的身體.
用語言.
不斷誘惑他.
每天這樣誘惑他.
最後用暴力去誘惑他.
令到他屈服.
每一個試探.
越來越厲害.

$^{721}$一次比一次更加厲害.
他清楚指明.
他不會和他同流合污.
不會和他一起犯罪.
更加教導他.
他為什麼不會這樣去做.
弱失不理他.
不聽他的說話.
盡量避開他.
不過在避無可避的時候.
他唯有三十六計走為上將.
衣服可以扔下.
但他的純潔卻不可以扔下.
他當機立斷.
乾淨利落.
採取最正確的方法.
寧可冒著生命的危險.
亦都不可以讓試探勝過.
弱失清楚知道試探的後果.
什麼都可以做.
但不是每一樣東西都有益處.
有些東西會帶來後果.
這些後果甚至不能承擔.
當知道後果的時候.
我們就要離開.
違反道德的事情.
我們不應該做.
違反聖經的教導的時候.
我們不應該做.
不榮耀上帝對人有損的事情.
我們都不應該做.
這是我們基督徒的一條線.
當我們面對試探.
我們要小心.
我們要謹慎.
我們可以做的是追求聖潔.
和上帝有親密的關係.
背誦聖經的經文.
遠離試探.
令到容易我們跌倒的環境.

$^{761}$和眼睛納若.
若果有些吸引出現.
我們就轉移視線.
避免漫無目的的上網.
建立興趣.
將精力放在運動.
放在你的興趣裡面.
和人分享掙扎.
用禱告去守望你的生命.
第二點.
若失做的是尊重人.
忠於他的主人.
若失向忠於他的主人.
盡忠和盡責.
不再越過主人規範下的規範.
他忠於波提佛.
寧死都不會做出對不起主人的事情.
忘恩負義.
這個妻子是屬於波提佛的.
波提佛和他的妻子是夫婦關係.
一切婚姻以外的性都是犯罪.
每個人做每件事.
不是不會影響人.
是一定會影響人的.
要記住我們今天所擁有的.
不單只是自己的努力.
其實我們今天所擁有的.
更加是背負著其他人的信任和期望.
要對我們信任的人.
期望我們的人負責任.
無論是我們的太太.
我們的丈夫.
我們的兒女.
我們的父母.
我們其他的家人.
朋友.
弟兄姊妹.
你的同事.
你所認識的人.
我們都在這個關係網裡面.

$^{801}$千萬不要令其他人失望.
我們每一個人都很著緊你.
很重視你.
和你有關係.
我們每一個人.
都是被帶著信任和期望.
建立信任可能要做很多事.
要經歷長時間才能夠被建立.
但其實破壞關係.
可能因為一件小事.
很快建立的關係完全被破壞.
沒辦法修補.
最後一點.
就是約瑟去敬畏神.
不去得罪神.
約瑟面對著主母的誘惑.
寧可前途盡毀.
甚至冒著生命的危險.
他都不去得罪上帝.
他不寧願得罪和他同在.
賜福他的神.
所以他很嚴厲去看待罪這件事.
神和他同在.
讓他升高.
一切都是上帝的工作.
不過如果你留意經文的時候.
在這段三十九章一開始.
提及神與約瑟同在.
在最後當他下監.
神亦與他同在.
不過在試探的時候.
經文沒有描述神與他同在.
當約瑟受試探的時候.
其實他最需要神.
不過作者沒有描述.
上帝的同在.
其實可以幫助約瑟逃離這些試探.
不過在整件事情發生的裡面.
上帝好像刻意退後.
讓約瑟面對著這個考驗.

$^{841}$其實新約聖經裡面記載.
耶穌面對著試探.
聖靈帶耶穌去到這個曠野.
面對著魔鬼撒旦試探的時候.
聖靈退去.
當耶穌勝過三個試探的時候.
聖靈又再出現.
為他去施洗.
被他去充滿.
在試探的現場的時候.
好像神不在裡面.
面對著試探.
我們以為沒有人知道.
我們以為神都缺席.
不過原來神一直在看.
人在做 天在看.
我們做得好.
上帝一定會賞賜.
但當我們犯罪 得罪神.
神一定會審判.
我們得罪人可以彼得過.
不過得罪上帝卻是沒有辦法逃離.
上帝一定會審判到底.
惡人必有惡報.
上帝必定會去追究到底.
有一個很有恩賜的傳道人.
在他年輕的時候.
28歲.
他去到不同地方傳媒報道.
和講上帝的訊息.
有一天他住在一個教會會友的家裡.
他休息.
但在深夜的時候.
有一個屋的屋主.
18歲的女孩.
起來告訴傳達人.
傳達人 你很有魅力.
我很喜歡你.
我想和你上床.
這個傳達人很堅決地拒絕.

$^{881}$這個女孩流著淚.
很尷尬地回到自己的房間裡.
七年之後.
這個傳達人收到一個回禮.
這個回禮好像是一個新娘子送給他的回禮卡.
他覺得奇怪.
明明他最近都沒有參加任何婚姻的聚會.
為什麼會有回禮呢.
原來當他打開這封信的時候.
是那個他拒絕的女孩寄信給他.
因為她送了最寶貴的禮物.
就是他的貞操.
約瑟選擇做正確的事情.
忠於自己.
忠於他的主人.
忠於他相信的人.
忠於他的上帝.
不過約瑟的敬虔和正直.
沒有幫他脫離誣告和監禁.
他雖然面對著試探.
他勝過.
不過他卻需要付出沉重的代價.
這個是他不妥協的結果.
最後一個部分.
第39章19到23節.
我邀請弟兄姊妹幫我讀三節經文.
19到21節.
一起去讀吧.
預備1,2,3.
(念經).
上帝同在.
是整個39章裡面的重點.
第2節,第5節,第21節,第23節裡面.
都在重複說這個主題.
上帝同在.
神的同在在約瑟的生命裡面.
有很多不同的工作.
令到他成功.
令到他順利.
令到他在其他人面前.

$^{921}$眼前蒙恩.
令到他的工作順利.
令到他的主人.
能夠將一切的事情交託在他裡面.
雖然他面對著困難,迫不還難.
不過上帝同在.
令他在波提法的面前.
令他在監獄長的面前.
能夠蒙恩.
管理一切所要做的事情.
原來神的同在不是抽象的概念.
也不是看不到的事情.
而是實實際際.
是看得到的印證.
有約瑟的地方就有祝福.
這個祝福及至周圍的人.
耶和華的同在不受環境的限制.
無論約瑟在哪裡.
在迦南地,在監獄.
在波提法的家裡.
上帝都和他同在.
上帝在整個人生裡面.
他掌管著,保護著.
當約瑟被賣下監.
甚至高高低低.
耶和華都和他同在.
上帝沒有應許天色常藍.
不過他卻應許.
一生都不會撇下我們.
在我們人生裡面.
也有高高低低.
每一個階段,每一個困境.
上帝從來都不會缺席.
教徒面對著人生的無常.
生,老,病,死.
無例外.
也都沒有免疫力.
教徒都會中肺炎.
教徒都會患癌症.
教徒都會有意外.

$^{961}$我們都是人.
我們和其他人一樣.
面對著生命的那份脆弱.
有些不幸.
可能因為我們自己做得不好.
有些不幸.
可能是其他人加害我們.
不過其實.
有些不幸.
是生命裡面的無常.
沒有辦法解釋的.
未必一定是有益處.
未必是有一定的原因.
不過上帝的同在.
他是很重要.
上帝在我們開心快樂的時候.
他為我們擊掌.
他為我們跳舞.
上帝在我們難過傷心的時候.
他說他用一個疲勞.
裝著我們每一個的眼淚.
上帝說當我們有憤怒.
當我們不滿的時候.
他願意傾聽我們裡面的感受.
成為我們最佳的觀眾.
當我們失敗自責的時候.
上帝什麼都沒有說.
只是擁抱著我們.
我們唯一能夠做的.
依靠的.
就是這位永恆不變的上帝.
有上帝同在.
無論我們遇到任何的光景.
任何的景況都好.
上帝都會給足夠的恩典我們.
讓我們走得過.
讓我們跨得勝.
上帝永遠都在.
只要我們主目去看神的時候.
神一定有恩典.

$^{1001}$祂是站在我們旁邊.
和我們一起面對.
面對著人生的無常.
除了上帝之外.
其實我們生命都很重要.
我們生命的品格.
我們的素質.
聖靈裡面.
在我們生命裡面工作都很重要.
若想在生命裡面.
其實坐著很厲害的過山車.
整天都高高低低.
人生有很多的起跌和低谷.
他第一個低谷.
就是被家人拋棄.
被人出賣.
第二個低谷.
就是被人誣告.
身敗名裂.
第三個的不幸.
就是有人答應了他.
不過他是忘記了他.
不過這一切的低谷.
都在磨練著.
弱失的生命.
這一切的環境.
都成為磨練弱失生命.
的一個很重要的品格.
他被出賣.
被誣告.
被忘記.
都不重要.
這一切都成為一個牧人行.
經過多番的折磨.
弱失的生命.
脫言不一樣.
可以能夠承擔大事.
有人說一個偉大的夢想.
其實是需要一個偉大的品格去承載.
特別是人生裡面的低谷.

$^{1041}$他願意努力去通過的時候.
他就能夠去完成一個民族裡面的大夢.
可以想像一個人.
如果要坐上一個重要的位置.
宰相的位置.
其實無論是他的能力.
他的品格.
他的要求一定很高.
從他在波提弗裡面.
家裡裡面開始.
到他在監獄.
到他去到法老的倉庫.
弱失每一個地方.
他都是在磨練他人際關係.
管理的能力.
在逆境裡面的能力.
這一切都像一個地獄訓練.
得來不易.
有個偉大的中國報道家宋尚哲.
他有一本書叫《失而復得的日記》.
總結了宋尚哲為甚麼他能夠.
行得過這個令程.
他裡面對苦難有幾樣描述.
他說大患難必定有大的復興.
神所重用的工人.
都是在患難裡面受到被詛咒.
神要藉著患難.
使人到謙卑.
神在患難裡面有恩典.
經歷困苦.
才能夠去憐恤其他人.
安慰其他人.
最重要的.
原來上帝不單止磨練人的品格.
和人一起.
最後一點就是.
上帝是在承受他的救恩.
當弱失他被賣.
在這個井裡面叫苦連天.
當他目睹他的親人.

$^{1081}$親手把他賣給以實瑪利亞人.
當他目睹著.
去到埃及波提法的家裡.
當他努力工作仍然被人誣告.
當他在監獄裡面為人解夢.
卻被人忘記.
種種的不幸.
其實當時弱失.
面對這種不幸.
是不明白的.
他不知道.
不過直至到後來.
他五十多歲的時候.
才能夠看得到.
人生裡面整個的圖畫.
原來他過去種種裡面的不幸.
是上帝完成他救恩的工作.
特別是在第五十章.
他和家人重逢.
他爺爺離開.
他兄弟向他請求饒恕的時候.
他終於明白過來.
他和他的兄弟這樣說.
第五十章十九節.
弱失對他們說.
不要害怕.
我豈能代替神.
從前你們的意思要害我.
但神的意思.
原是好的.
要保存許多人的性命.
成就今天的光景.
現在你們不要害怕.
我必養活你和你們的婦人孩子.
可能我們正在經歷生命裡面.
很多的破碎.
失望.
孤單.
我們不夠膽看將來是怎樣.
不夠膽有任何的幻想和期望.

$^{1121}$不過弱失的生命正在鼓勵我們.
我們的生命.
一定是正在經歷磨練.
只要我們願意對準神.
神和我們同在.
每一個點都能夠成為轉彎.
年輕的弱失.
從來沒有想過自己是做夢.
原來這個夢.
有朝一日會這樣成真.
他有朝一日會成為埃及的宰相.
拯救他的夢.
拯救他的家庭.
但原來上帝的工作.
不單要拯救他的家庭.
更加是一個偉大的任務.
這個任務就是.
完成亞伯拉罕的.
上帝給他的應許.
成為大國.
這個大國是需要國民.
國土和國法.
神讓弱失成為宰相的目的.
是要讓以色列人人數加多.
令以色列人在埃及這個文化.
大國的庇佑之下成長.
人數多到一個地步.
令到埃及王都擔心.
展開出埃及這個故事.
十幾歲的年青弱失.
他的生命正在經歷翻天覆地的改變.
十幾歲的年青人.
生命真的可以有翻天覆地的改變.
因為他們正在經歷生命裡的大喜大樂.
若果他們能夠在生命裡的大喜大樂裡.
他們認識上帝.
他們品格被塑造.
和上帝同行的時候.
有不一樣的目光.
今天我們除了在不幸裡.

$^{1161}$我們正在經歷上帝的時候.
我們又如何去看待.
我們今天所面對的不幸呢?.
其實上帝不單要祝福我們的生命.
更加要使用我們的生命.
去祝福這個世代裡的人.
最後我去說.
當我和太太.
正在經歷孩子裡的難處的時候.
我們去到一個.
我們去醫院.
然後又認識了很多不同的家庭.
我們發現原來這個香港裡面.
真的有很多很破碎.
小朋友有很嚴重的疾病.
罕有病的家庭.
我們在過程裡面去.
去誤會.
主動去傳訊息給他們.
因為有群組去說.
他的小朋友抽筋.
抽筋了半個小時.
又很嚴重.
有很多不同的問題.
到五歲十幾歲都還沒有學會走路.
各樣的事情.
我們就簡單傳訊息給他祈禱.
原來我們在患難的裡面.
都可以和人同行.
今天你面對什麼患難.
困苦和不幸呢?.
神說祂愛你.
祂願意去擁抱你.
更加願意使用你的不幸.
成為其他人的祝福.
最後讓我們一起祈禱吧.
(靜音).
弟兄姊妹.
今天你面對什麼不幸呢?.
是家人?.

$^{1201}$是自己的生命?.
是工作?.
還是人際關係?.
上帝說祂會和你同在.
祂會給你力量.
一定可以走過.
你身邊有沒有一些人.
其實都在面對不幸呢?.
你覺得你可以怎樣去做.
分享上帝的愛呢?.
他們需要你.
有沒有一些人物.
在你腦海裡浮現.
你很想為他們祈禱嗎?.
我們一起祈禱吧.
深愛的上帝.
因為你愛我們到底.
你不會離開我們.
在我們人生高高低低的起跌裡面.
你和我們同在.
你應許你會給我們力量.
我們求你.
這個世代裡面有很多人.
經歷不幸,患難.
特別是年青人一代.
有很多挑戰.
求你教導我們多走一步.
去擁抱他們.
讓他們能夠經歷上帝的愛.
求你真摯,加力量.
要我們同在.
聽我們祈禱.
奉耶穌得勝之名求.
阿們.
\newpage



\section{約書亞記 14:6-15}
\label{sec:ec0Lp80ZOxQ}
\textbf{2021年11月14日長者主日崇拜直播｜葉筱春牧師｜我還是強壯｜約書亞記14\_6-15}
\newline
\newline
連結: \href{https://youtube.com/watch?v=ec0Lp80ZOxQ}{\texttt{https://youtube.com/watch?v=ec0Lp80ZOxQ}} ~~~~ 語音日期: 2021-11-21
\newline
\newline
\hyperref[sec:wlHRHi438v8]{\small{< < < PREV SERMON < < <}}
~
\hyperref[sec:index_chronic]{\small{[返順時目]}}
~
\hyperref[sec:N45H8jB07Gc]{\small{> > > NEXT SERMON > > >}}
\newline
\newline
約書亞記 14:6-15
\newline
\begin{longtable}{cl}
\hline
\hline
章節 & 經文 (和合本修訂版)\\
\hline
14:6 & \begin{tabularx}{0.7\textwidth}{X} 猶大人來到吉甲，約書亞那裡，基尼洗族耶孚尼的兒子迦勒對約書亞說：「耶和華在加低斯‧巴尼亞指著我和你對神人摩西所說的話，你都知道。 \end{tabularx} \\ \\ \relax
14:7 & \begin{tabularx}{0.7\textwidth}{X} 耶和華的僕人摩西從加低斯‧巴尼亞差派我窺探這地的時候，我剛四十歲。我把心裡的話向他報告。 \end{tabularx} \\ \\ \relax
14:8 & \begin{tabularx}{0.7\textwidth}{X} 雖然同我上去的眾弟兄使百姓膽戰心驚，我仍然專心跟從耶和華－我的神。 \end{tabularx} \\ \\ \relax
14:9 & \begin{tabularx}{0.7\textwidth}{X} 那日，摩西起誓說：『你腳所踏之地必要歸你和你的子孫永遠為業，因為你專心跟從耶和華－我的神。』 \end{tabularx} \\ \\ \relax
14:10 & \begin{tabularx}{0.7\textwidth}{X} 現在，看哪，耶和華照他所說的使我活了這四十五年。當以色列人在曠野飄流的時候，耶和華曾對摩西說了這話。現在，看哪，我已經八十五歲了。 \end{tabularx} \\ \\ \relax
14:11 & \begin{tabularx}{0.7\textwidth}{X} 現今我還很健壯，像摩西差派我去的那天一樣；無論是戰爭，是出入，我現在的力量和那時的力量一樣。 \end{tabularx} \\ \\ \relax
14:12 & \begin{tabularx}{0.7\textwidth}{X} 請你將耶和華那日所說的這山區給我。那日你也曾聽說，這裡有亞衲族人，以及寬大堅固的城，或許耶和華會照他所說的與我同在，我就把他們趕出去。」 \end{tabularx} \\ \\ \relax
14:13 & \begin{tabularx}{0.7\textwidth}{X} 於是約書亞為耶孚尼的兒子迦勒祝福，把希伯崙給他為業。 \end{tabularx} \\ \\ \relax
14:14 & \begin{tabularx}{0.7\textwidth}{X} 所以希伯崙成了基尼洗族耶孚尼的兒子迦勒的產業，直到今日，因為他專心跟從耶和華－以色列的神。 \end{tabularx} \\ \\ \relax
14:15 & \begin{tabularx}{0.7\textwidth}{X} 希伯崙從前名叫基列‧亞巴；亞巴是亞衲族最尊貴的人。於是國中太平，沒有戰爭了。 \end{tabularx} \\ \\
[1ex]
\hline
\hline
\end{longtable}
$^{1}$《約書亞記》14章6到15節.
大家可以打開舊約的280頁.
請聽我讀出.
那時猶大人來到吉甲見約書亞.
有基里族耶夫尼的兒子加納對約書亞說.
耶和華在加底斯巴利亞指著我.
與你對神人摩西所說的話.
你都知道了.
耶和華的僕人摩西.
從加底斯巴利亞打發我窺探者地.
那時我正四十歲.
我按著心意回報他.
然而.
和我上去的眾弟兄使百姓的心消化.
但我專心跟從耶和華我的神.
當日摩西起誓說.
你腳踏之地.
定要歸你和你的子孫永遠為業.
因為你專心跟從耶和華我的神.
自從耶和華對摩西說這話的時候.
摩西照他所應許的.
使我存活這四十五年.
期間以色列人在曠野行走.
可那.
我現今我八十五歲了.
我還是強壯.
像摩西打發我的那天一樣.
無論是增節是出入.
我的力量那時如何.
現在還是如何.
求你將耶和華那日.
應許我的這山地給我.
那裡有阿納族人.
並寬大堅固的城.
你也曾聽見了.
或者耶和華照他所應許的.
與我同在.
我就把他們趕出去.
於是約書亞為耶夫尼的兒子加納祝福.
將希伯倫給他為業.

$^{41}$所以希伯倫作了基尼洗族.
耶夫尼的兒子加納為的產業.
直到今日.
因為他專心跟從耶和華以色列的神.
希伯倫從前名叫基列阿爸.
阿爸是阿納族中最尊大的人.
於是國中太平.
沒有戰爭爭前了.
今天很高興有粉嶺浸信會榮休主任牧師.
葉紹春牧師來到我們當中宣講.
他今日的講題是「我還是強壯」.
有請葉牧師.
最終以來往浸港島.
可以用真面目跟弟兄姊妹相聚.
特別可以一齊分享神的話.
我今年十月有機會第一次.
用長者的身份去海洋公園玩.
你拿著一張長者卡.
去到長者的櫃台.
讓他看看你的樣子.
對一對就給你一張免費的即日入場券.
不需要在網上事先登記.
不用排隊.
非常方便.
感覺到留在香港做長者特別開心.
特別好.
我退休幾年.
這幾年都有很多機會善用長者的優惠.
享用了很多餐飲的美食.
我才發覺有些地方.
其實55歲已經開始當長者優惠.
所以我吃少了十多年.
所以我先跟大家分享.
不要浪費了.
55歲已經當長者了.
那些優惠價非常之低.
所以真的很感恩.
下星期日就是長者日.
有更多優惠.
周圍都是優惠.

$^{81}$周圍都是免費的給長者.
所以做長者是很被重視的.
特別是旺鎮.
今天有長者的主日崇拜.
特別很重視.
很欣賞.
很感恩教會有這麼多長者.
其實長者都是小孩子.
越是長者就越是小孩子.
我們的心態是.
我們的精神也是.
我試過有幾次跟年青人去行山.
他們睡著覺.
我還跟他們一起行山.
我還睡覺.
起床就走了.
很累啊.
現在十多歲的小孩子.
其實都沒有耐力.
不夠我們行山.
所以我們感恩.
聖經說人生七十.
若果強壯可以到八十.
現在普遍人都是很健康的.
沒有地方是超過八十歲的.
全部都在這邊.
分區坐的.
八十歲了.
人生七十.
強壯就到八十.
現在普遍人都健康長壽.
特別在過去一年.
香港的長壽數字.
女性八十八歲.
你可以說有段時間捱.
但你可以說有段時間享受.
男性都不差.
八十二歲.
所以香港的長壽數字.
在亞洲是很高的.

$^{121}$我們去看我們今天的大綱.
引言說.
我們的人生是一個.
豐盛的人生.
因為在約翰福音第十章第十字說.
我來了是要叫人得生命.
並且得的更什麼.
我們就填了豐盛.
很快速.
已經寫了答案在這裡.
我們會用三個比喻.
來說一下我們的人生.
第一個比喻就是.
人生就像一條河流.
越流就越.
可以出答案了.
人生就像一條河流.
越流就越深.
越流就越廣.
這就是豐盛的人生的寫照.
什麼叫豐盛呢.
可以說是很開心.
很幸福.
很滿足.
很多元化.
很多姿多彩.
人生幾十年.
會遇到的事情越來越多.
有很多事情都是意料之外發生的.
但還有很多.
就是意料之外發生.
我們的人生的經歷.
是越來越豐富.
我們的生命.
就像河流一樣越來越有深度.
更加能夠掌握人世間的人情世故.
幾十年的人生裡.
有很多經歷.
一直影響我們周邊的人.
特別是我們的兒孫.

$^{161}$也祝福了很多生命.
人生就像貨櫃車一樣.
當沒有東西載著貨物的時候.
在路上很嘈吵.
但如果貨櫃車是裝滿了貨物的話.
在路上就會覺得很安靜.
我們很多人生的生活智慧.
都是這樣累積而成.
智慧是有別於聰明.
智慧是生活的應對.
有很多人很聰明.
讀書很厲害.
但他沒有智慧.
他說一句話也沒有智慧.
所以智慧是應對生活的能力.
面對壓力的動力.
跟聰明不同.
跟學識不同.
所以我們的生活智慧.
是不斷累積的.
我們有沒有發覺.
我們的生活上有很多認知.
在過去幾十年都錯的.
到了幾十歲才明白原來是這樣.
才慢慢調教.
有沒有試過這樣的經歷.
我們的知識增加了.
我們的經驗也豐富了.
也純熟了.
這時候會比我們的生命.
來得更加謙卑.
因為發覺原來.
我們雖然經驗豐富.
但有很多人比我們更豐富.
更資深.
於是我們會變得謙卑.
特別是在事奉的生命上.
我們的生命.
會來得更加謙卑.
更加安靜的事奉.

$^{201}$因為我們有資料.
貨櫃車裝滿貨物.
走來走去都安靜一點.
我們的生命有屬靈的能力.
屬靈的智慧.
屬靈的悟性.
我們裡面有資料有份量.
我們的事奉就會安靜.
如果我們的生命沒有資料.
沒有份量.
我們的事奉就白雜不堪.
好像我初信的時候的生命.
滿有熱誠.
是要讓人看到我的能力.
要人欣賞我.
要人叫我.
請我.
做事.
去顯示我自己的能力.
其實在教會裡面.
最有能力的事奉.
就是能夠隱藏在耶穌基督裡面的事奉.
人生就好像一棵大樹一樣.
樹蔭.
充裕樹根.
也是粗壯的.
有時候我們看到百年老樹.
它的形態突出.
氣勢不凡.
我們會停下來多看幾眼.
我今天來到經過荔枝角.
美孚站的時候.
經過荔枝角公園.
看到幾棵老樹.
真的就好像我所說的.
因為我最近去過荔枝角的花園.
走了一段時間.
幾十年都沒有去過.
退休之後才逐個花園去享受.
老樹的形態突出.

$^{241}$氣勢不凡.
真的會讓人停下來多看幾眼.
它們裡面有數不清的根.
深入地底.
令這棵老樹能夠健康地成長.
我們信仰的生命.
幾十年.
仍然要繼續在神的話語裡扎根成長.
就算你人生遇到多大的風浪.
你都不會懷疑神.
不會離棄神.
不是有些人回來教會.
聽到別人說一句話.
就說「我不事奉了」.
別人無意說了一句話.
你說「我不回教會了」.
因為人一句說話.
你就放棄這麼寶貴的天堂的福祿.
真的不值得.
正如詩篇第一篇第二節說.
我要喜愛耶和華的律法.
就耶思想.
這人變為有什麼?.
有福了 有福氣.
所以心中有神的話語就有福氣了.
我們要將神的話語藏在心裡面.
免得我們得罪神.
我們千萬不要將人的話語放在心裡.
因為說的人已經忘記了.
他就睡到咕嚕咕嚕的.
你氣死了.
整晚想他為什麼這樣說我.
為什麼要這樣說這個話.
有什麼意思呢.
整晚說到胃痛.
不開心.
但說的人都忘記了.
已經睡到咕嚕咕嚕的.
所以我們將神的話語放在心裡.
我們就有福氣.

$^{281}$但如果我們將人的話語藏在心裡.
我們就只有激氣.
用了三個比喻說出人生.
希望看回自己的生命.
看回自己在教會成長的經歷.
今天特別在舊約裡面.
找了一個人物.
叫加勒.
他和約書亞一起是.
這十二個探子的其中一個.
從加勒的生命裡面.
我們要學會怎樣去看待.
我們要學會怎樣依靠神的恩典.
過著一個豐盛的人生.
剛才弟兄已經讀了一段經文.
我們不再重讀這段經文.
我會從三方面去思想.
怎樣靠著神的恩典.
過著一個豐盛的人生.
第一點.
我看我們的大綱.
我們信息第一點.
因為專心跟從神.
加勒可以明白神的心意.
我們一起讀一讀.
約書亞記第十四章第八節.
邀請弟兄姊妹一起讀.
約書亞記第十四章第八節.
看回那段文章.
第八節我們一起讀.
預備開始.
(念念以復我上去歷種歷經史百世的心思和行).
(但我專心跟從耶和華和色心).
四十年在曠野的路.
一定經歷了不少.
也改變了很多.
在這四十年的曠野路裡面.
人都變得成熟了.
每餐吃嗎哪.
每餐吃鵪鶉.

$^{321}$每餐吃礦泉水.
沒什麼味道.
可能有一點甜味.
今年我從美國回來.
在酒店隔離了十四天.
第一個星期都可以的.
到第二個星期.
這間酒店真的好.
因為它很節省.
應該是一大批貨.
訂食物.
所以全部都是那類的食物.
只是不同味道.
有時甜點 有時辣點 有時鹹點.
到第二個星期真的不想吃.
但不吃又不行.
不過我們已經用了一個大皮箱.
大皮箱裝了很多零食.
很多杯麵.
預備去熬過這十四天.
十四天我都熬不了.
何況四十年.
吃同樣的食物.
這些操練.
會令人更加的成熟.
迦勒他記得.
四十五年前.
摩西如何揀選他和約書亞.
和另外十個探子.
成為十二個探子.
從迦底斯巴尼亞這個曠野地方.
猜他們可以去到哪裡.
他們就去到那裡.
從迦底斯巴尼亞這個曠野地方.
猜他們可以去到迦南窺探.
迦勒和約書亞.
是十二個探子裡的唯一兩個.
告訴摩西.
這是一個什麼樣的地方.
約書亞和迦勒.

$^{361}$回報了一個叫做有誠實.
又可靠.
又客觀的報告.
Honest and Accurate Report.
這個在我們今天的世代很重要.
這十二個探子裡.
只有他們兩個.
能夠擺出一個誠實的心.
將他們看到的東西.
向摩西回報.
而是正確的事實.
其他的十個探子.
他們眼見的狀況.
他們心裡自己懼怕.
他們自己不安.
於是根據民數記第十三章.
第32,33節說.
他們就回來講謠言.
他們就回來報惡信.
迦勒.
他知道他為什麼可以.
按著心意去回報摩西.
他將一個又誠實又正確的.
迦南地的現況.
相對那十個探子說.
不要去啊.
迦南地的人會吃人的.
這是民數記的記載.
迦勒他知道.
他可以和別人不同.
敢於承擔.
敢於成為少數.
我小時候回教會.
都是在旺角附近的地方.
所以每一次回家.
都經過旺鎮樓下.
經過樓下.
有一間餐廳叫做美爾廉.
不知道有沒有人有印象.
我們經常打完球.

$^{401}$團契完就在那裡吃.
吃完之後.
走著走著和朋友聊天.
大家都是走那條路回家.
走著走著.
回到我家樓下.
還沒聊完又要再聊.
於是去我家樓下的餐廳.
叫做別不同.
你有沒有吃過?.
即是老旺角老街坊.
再聊一聊.
最喜歡別不同的煎蘿蔔糕.
很香啊.
迦勒他知道他能與別不同.
能夠分別為性.
分別為性不是叫你做聖人.
不笑.
不說話.
不吃東西.
不是這些聖人.
分別為性是代表.
你能夠和這個世界分別出來.
和罪惡分別出來.
和那些不屬神的人分別出來.
這就叫分別為性.
迦勒他知道他只能夠與別不同.
能夠敢於承擔.
成為少數.
是因為他願意.
專心地跟從神.
其實我們只要每天緊貼地跟從神.
我們便會很容易明白神的心意.
迦勒他按著留在心裡的感動.
去回報給摩西.
是一種平安的心日.
吸心意.
當我們知道神的心意之後.
而有祂的信服的話.
心裡便會有平安.

$^{441}$所以你去尋求神的心意.
而有平安的時候.
這就是神對你的回應.
神不單給平安在你的心裡.
神也給平安在你周邊的人和家人.
如果你求神要不要走.
神給你很平安的說我要走.
但你家人個個都不平安.
這不是神的心意.
因為神感動你.
亦會感動你周邊的人.
這十個探子.
心裡有很多掙扎.
很多不安.
他們帶來的訊息.
令聽的人.
接受的人都大大不安.
很恐懼.
弟妹們.
記得.
當你願意去尋求神的心意的時候.
你就會尋見.
因為這是神的應許.
神會給每一個尋求祂的人.
有一份出人意外的平安.
不單給你.
也給你周邊的人.
我們看筆記大綱的第一點.
因為專心跟從神.
迦勒在我後面其實是括著了.
我現在邀請大家.
你要明白神的心意.
你就每天緊貼著神.
如果你願意學習.
好像迦勒這樣.
能夠敢與別人不同.
成為少數.
成為一個分別為勝的人.
你在這個括弧上寫了「我」字.
或簽了你的名字.

$^{481}$去學習.
怎樣能夠每天緊貼神的心意.
我們看第二點.
十四章第九節.
因為專心跟從神.
迦勒可以得著神的賞賜.
祝福和應許.
我們一起讀第十四章第九節.
一起讀第十四章第九節.
一起讀 預備開始.
當日摩西起誓說.
你腳所踏之地.
定要歸你和你的子孫.
永遠為業.
因為你專心跟從耶和華我的神.
我們不要忘記神的應許.
四十五年了弟兄姊妹.
迦勒仍然沒有忘記神的應許.
神的應許是絕對不會落空的.
祂是一個信實.
可靠的神.
神一定會負責.
祂對你一切的承諾.
只要你不忘記.
很多時候神的應許.
就成為你這個時候.
繼續向前衝的動力.
亦是我們的盼望和活力.
弟兄姊妹請你記得.
一定在你生命當中最有信心的時候.
將你的應許記在心裡.
在你最有能力的時候.
你要將你的應許記在心裡.
要在你最生活平順.
無風無浪的日子.
將神的應許記在你的心裡.
因為當你困難的時候.
當你軟弱的時候.
你是拿不起.
拿不動這本聖經.

$^{521}$我在2019年.
曾經因為一次大病.
在威爾斯親王醫院住了28天.
他們的細心照顧.
眼睛真的蒙到完全看不到東西.
手震得像是在震動.
手震得像是無法回電話.
腳無法走動.
想讀聖經真的不行.
我更加深信.
弟兄姊妹.
在你最健康,最有能力,最有信心.
最風平浪靜的時候.
你要將神的應許記在心裡.
當你有困難的時候.
當你的家人面對生命的挑戰的時候.
或者當你的家庭面對這些壓力和苦難的時候.
就成為了你的祝福.
成為了你家人的祝福.
為什麼很多時候在我們生活面對困難的時候.
我們軟弱,更軟弱.
就是因為我們平時沒有做足準備.
沒有將神的應許.
在我們最有信心的時候.
藏在我們心裡.
45年了弟兄姊妹.
迦勒仍然沒有忘記神對他的應許.
當他們進入迦南地的時候.
打了幾場仗.
他們的狀況開始穩定.
開始分地.
在分給十二支派地之前的第一個項目.
就是要將地.
希伯倫地分給迦勒.
一個肥美的希伯倫地.
就成為了迦勒和他的子孫.
一個敬畏神的人會有很多福氣.
而且敬畏神的後代.
也會大大的蒙福.
神的應許.

$^{561}$大部分已經記載在這本聖經裡.
舊約三十九.
新約二十七.
加上六十六卷.
一千一百八十九章.
三萬一千一百零一節.
全部都是神的說話.
是神的應許.
你要將它放在你的心裡.
啟示錄.
擔心我們忘記神的祝福和應許.
說你要寫下.
你要寫下.
要寫下神所有的應許.
弟兄姊妹.
我們可以學習每天.
來抄寫聖經.
抄一節也好.
抄一段也好.
抄一章也好.
你用毛筆也好.
用鋼筆也好.
用鉛筆也好.
用蠟筆也好.
總之你是抄.
抄寫聖經.
讓你能夠有更深刻的印象.
如果自己沒有這種能力抄.
我們就在團契.
有個周會一起去抄聖經.
抄一節也好.
抄一段也好.
抄一章也好.
我們看回我們的大綱第二點.
因為專心跟從神.
迦勒可以得到神的賞節.
祝福和應許.
如果你願意在你的生命上.
得到更多神的祝福和賞節.
在這個窟窿裡面請你寫上「我」字.

$^{601}$或者簽上你的名字.
願意納至你在你有信心的日子.
在你健康的日子.
在你有能力的日子.
在你風平浪靜的日子.
多些將神的應許放在心裡.
在有一天.
你的家庭.
你自己個人面對危機的時候.
就成為了你和你家人的祝福.
我們看第三點.
約書亞記第十四章第十到第十五節.
因為專心跟從神.
心中滿載神的話語.
迦勒因此建立了一個健康的自我形象.
一起跟我讀.
約書亞記第十四章第十四節.
我們一起讀.
預備起.
所以希伯倫作了基尼洗族耶夫尼的兒子迦勒的產業.
直到今日.
因為他專心跟從耶和華以色列的神.
神的話語是有能力的.
祂能夠塑造我們的生命.
好像一塊鏡子.
讓我們更加清楚自己.
知己知彼就能夠百戰百勝.
神的話語可以每日成為我們腳前的燈.
路上的光.
神沒有一次過將你的人生告訴你.
神也不想你用任何方法.
知道你一生要走的路.
因為擔心你知道了之後.
你不夠膽走下去.
每早晨.
這都是新的.
你的盛世極其廣大.
每日有神新的恩典.
其實祂的話語成為你腳前的燈.
路上的光.

$^{641}$引帶你走人生神的路上.
神的話語更加能夠令你的能力勝過試探.
免得跌倒.
當加勒仍然說.
我是強壯的.
我還是強壯的.
他不是亂說的.
他充滿信心地說.
45年之後.
今天他說.
我仍然強壯.
身體狀況可以為神打仗.
得到這塊山地.
他充滿信心地去作戰.
因為他知道.
他深信神的應許.
只要有神的同在.
他一定可以奪取這塊山地.
結果他的自信不是盲目的.
經文說他85歲的時候.
真的將大敵人亞納族趕走.
然後從此希伯倫就國中太平.
耶穌同樣在約翰福音第十六章33節.
跟我們說.
我會給你平安.
不過這個世上是有困難.
信耶穌的人都會病.
信耶穌的人都有困難.
但耶穌說.
我給你有平安.
你放心依靠我.
因為我已經得勝了這個世界.
當加勒華說.
我仍是強壯的時候.
你會想起在《新滅記》裡面的.
三章第25節.
日子如何.
力量也如何.
在我那28天住院的過程當中.
我更加認識了我自己.

$^{681}$更加認識了我的神.
所以很深信.
保羅在哥林多後書4章第16節說.
所以我們不要喪膽.
外體雖然毀壞.
但內心不卻.
一天生事一天.
加勒能夠這樣建立.
一個這麼有自信的自我形象.
是因為他專心謹從神.
他心中充滿的不是人的話語.
而是神的話語.
一直他經歷過神的帶領.
供應和保守.
他深信神的話語.
能夠在他生命當中.
有很多次印證是真的.
是可靠的.
是真實的.
我們看回我們的大綱.
因為專心謹從神.
加勒滿載神的話語.
就建立一個健康的自我形象.
邀請你立志.
因為專心謹從神.
我心中就滿載神的話語.
我因此能夠建立一個健康的自我形象.
如果你願意立志.
請你填上這個「我」字.
或者簽名一起學習.
最後.
祝福我們教會的長者.
能夠開開心心地過著一個慢活的生活.
不要再追車不要再追船了.
我試過追.
真的要兩天才回氣.
不要跌倒了.
不過我跌過很多次.
有一次越南跌倒之後.
半年都走不動.

$^{721}$帶著拐杖.
不過唯一的好處就是.
我跌倒之後的下個星期.
我就要去首爾.
接著由首爾直接飛去美國.
唯一的好處就是.
我坐著輪椅師母貼著我.
過關特別快.
有人載你到門口.
不過不要跌倒.
要小心.
其實沒得小心的.
跌就跌的了.
我以前常說.
我牧羊的時候說.
你小心點不要跌倒.
其實當我跌過四五次之後.
發覺不是說不跌就不跌.
原來是要跌就跌的.
不過這樣都要小心.
保重身體.
好好享受放慢了的生活.
不要每天的活動一個接一個.
呼吸不周.
特別當你幫忙照顧孫子的時候.
你更加不容有失.
跌倒時流血.
你要快點自我報案.
如果沒有下次給你看.
我的孫子在美國.
所以我不可以說謊話.
因為客廳全部有攝影機拍著.
當我女兒出去辦事的時候.
她回來會看攝影機.
為什麼孫子這裡腫了.
瘀了.
所以我們在她還沒回來之前.
就用電話告訴她.
跌倒了 孫子.
現在沒有第二次給你看.

$^{761}$所以我們要很小心.
每次我們吃東西的時候.
要吃慢一點.
多咬幾口.
享受一下那些味道.
不像以前趕上班.
十分鐘就吃完午餐.
多咬十口 多咬二十口.
不要吃到飽就睡覺.
早點吃飯.
不要吃那麼飽.
吃四成五成就夠了.
求主幫助我們過著慢活的生活.
要享受子孫的愛惜.
不要說謊話.
這麼浪費錢.
其實很甜.
買這麼爛的葡萄.
吃葡萄很好吃.
其實很開心.
告訴鄰居.
我的女兒買給我吃的.
多讀經祈禱.
守望下一代.
因為時間關係.
我每天的祈禱都是為我的孫子.
為我的女兒.
求健康平安之外.
我求多一樣東西.
就是求我的女兒和我兩個孫子.
能夠比我更愛主.
更敬虔神.
這是我每天為我的女兒和孫子.
祈禱的主要內容.
擴充社交網絡.
我抱著孫子出街的時候.
會引來一些目光.
有些對話.
不認識別人.
但當你抱著孫子的時候.

$^{801}$就自然有朋友走過來和你聊天.
問你很多問題.
我們也多了和長者接觸.
因為可能參加了很多長者的活動.
無論社區中心或是教會.
表示我們的福音對象.
又闊又大了.
所以鼓勵大家.
當你的孫子或兒子為你擺生日宴的時候.
你可以在你的宴席之前.
先祈禱感恩.
或是讀一段聖經.
因為有很多夥計會看著你.
是一個好的見證.
有很多未信的親人.
會看到你的見證.
所以你要珍惜很多機會.
在不同的網絡裡.
為神得到更多的信息.
我退休了.
在退休之前.
經常出入醫院和殯儀館.
我最近在殯儀館裡.
帶領主理安息禮拜的時候.
見到一個堂倌.
他很感動地告訴我.
今年他患了一個大病.
快要死了.
可以救回.
認識了這堂倌十多年了.
我就問他.
你信不信耶穌.
他說想.
在殯儀館大廳裡.
我們一起祈禱.
我就帶領他信耶穌.
我退休之後.
也擴充了我的社交圈子.
我約了這班堂倌喝茶.
我跟他們說.

$^{841}$我約了這班堂倌喝茶.
有些老前輩已經退休了.
他們要回香港.
我跟他們說一起喝茶.
我要將福音傳給他們.
這個堂倌.
是我去殯儀館這麼久.
第二個帶領信耶穌的堂倌.
堂倌很喜歡站在那裡聽你講道.
聽了幾十年.
他們都有感動.
擴充你的社交圈子.
為神得到更多的生命.
最後預備將來的事.
我這個月有三次.
跟女兒聊我之後的事.
是很開心的.
突然間有話題.
聊著聊著就聊著.
要怎樣處理我的骨灰.
要怎樣處理我的財富.
要怎樣處理.
這個月剛好有三次.
我跟女兒聊這些問題.
我很感恩我的爸爸.
在他要回天護之前的幾個月.
他已經叫我們子女.
回到他的床邊.
講了他的心意.
預備好一切.
唱什麼詩歌.
哪個牧師講道.
哪個牧師主禮.
已經告訴我們.
我們到最後照辦.
在去年.
我們為一個101歲的姐妹.
主持安息禮拜.
在她70歲的時候.
她已經將她要讀的聖經.

$^{881}$和唱的詩歌告訴我們.
我已經寫在電話裡.
在她101歲的安息禮拜裡.
我們就用她所選的詩歌.
和她所讀的聖經.
為她預備這個安息禮拜.
能夠在她70歲的時候.
能夠預備自己的將來.
是一個信心的表現.
也請大家能夠更加.
未來到這個階段的時候.
多認識死亡將要發生的事情.
也都做好準備.
讓你和你的兒女.
都可以有這份福氣.
可以用最好的準備.
去到神的懷抱裡.
無後悔.
最後.
很簡短的.
給予五位非長者提醒.
我們要多陪伴他們.
要聆聽他們的傾訴.
說到重複的.
重複都不要緊.
我們的孫子.
做一個動作.
女兒說.
做多一次給阿公看.
再做多一次.
做幾次都開心.
老人家和你說話.
說過了.
爸爸不要說這麼沉悶.
你的女兒說話.
說幾次就行.
爸爸.
說多幾次你就覺得沉悶.
其實都很不公道.
我們要細心照顧長者的身體.

$^{921}$在家居上有沒有一些安排.
要改動.
令她生活得更加好呢?.
在這個疫情期間.
我和我女兒生活的時候.
我女兒.
在網上為我買了一大堆的衣服.
很開心.
又和我剪了兩次頭髮.
剪成怎樣我也很開心.
而且最感動第二次.
她是帶著肚子和我剪頭髮.
沒有估計的.
大肚子都可以拿剪刀.
我的孫女都很健康.
沒有問題.
照顧長者的需要.
留心他們的需要.
我們做父母.
我們會去學和子女溝通.
我們都要學和長者溝通.
在我母養過程中.
我見到很多女兒.
一見到爸爸就罵他.
喝他.
我聽了都很心痛.
為什麼女兒會在公眾場合.
喝爸爸呢?.
學爸爸.
其實我們都需要學和他們溝通.
最後我們同樣要預備.
和長者暫別的時候.
他們信了耶穌沒有?.
要不要做傳福音?.
不像我作為牧師.
我很多時候要去到深切治療部.
傳福音.
那些姑娘醫生覺得你傻.
已經插了十幾條喉.
你還在傳福音.

$^{961}$你信不信?.
你信就祈禱吧.
因為家人要牧師傳福音.
你不傳.
很壞.
那就傳吧.
在深切治療部.
在傳福音.
沒有反應.
那些姑娘看著你.
你都不知道怎樣做.
照做吧.
做了牧師.
一定要放下面子.
可不可以早一點.
未昏迷前去傳福音?.
求主幫助我們.
我們有準備的心.
和長者暫別.
有很多年輕人.
或者依賴父母.
依賴老人家的人.
很容易失落.
求主幫助我們.
不好意思遲了一點.
求主讓信息.
祝福我們每個長者.
一起祈禱.
獻上讚美.
獻上感恩.
因為你讓我們知道.
我們還是強壯.
只要有神.
應許與我們同在.
所以我們會更加明白.
看到自己的狀況.
求主祝福旺盛.
每一位弟兄姊妹.
無論什麼年紀.
都在你面前.

$^{1001}$無印無褔.
仰望交託.
奉主耶穌基督得勝名求.
阿們.
好嘛.
\newpage



\section{約翰福音 2:13-22}
\label{sec:N45H8jB07Gc}
\textbf{2021年11月28日崇拜直播｜陳明泉牧師｜基督帶來的革新:新的聖殿｜約翰福音2\_13-22}
\newline
\newline
連結: \href{https://youtube.com/watch?v=N45H8jB07Gc}{\texttt{https://youtube.com/watch?v=N45H8jB07Gc}} ~~~~ 語音日期: 2021-11-29
\newline
\newline
\hyperref[sec:ec0Lp80ZOxQ]{\small{< < < PREV SERMON < < <}}
~
\hyperref[sec:index_chronic]{\small{[返順時目]}}
~
\hyperref[sec:Msw_I2sXTQc]{\small{> > > NEXT SERMON > > >}}
\newline
\newline
約翰福音 2:13-22
\newline
\begin{longtable}{cl}
\hline
\hline
章節 & 經文 (和合本修訂版)\\
\hline
2:13 & \begin{tabularx}{0.7\textwidth}{X} 猶太人的逾越節近了，耶穌上耶路撒冷去。 \end{tabularx} \\ \\ \relax
2:14 & \begin{tabularx}{0.7\textwidth}{X} 他看見聖殿裡有賣牛羊和鴿子的，還有兌換銀錢的人坐著， \end{tabularx} \\ \\ \relax
2:15 & \begin{tabularx}{0.7\textwidth}{X} 耶穌就拿繩子做成鞭子，把所有的，包括牛羊都趕出聖殿，倒出兌換銀錢之人的銀錢，推翻他們的桌子， \end{tabularx} \\ \\ \relax
2:16 & \begin{tabularx}{0.7\textwidth}{X} 又對賣鴿子的說：「把這些東西拿走！不要把我父的殿當作買賣的地方。」 \end{tabularx} \\ \\ \relax
2:17 & \begin{tabularx}{0.7\textwidth}{X} 他的門徒就想起經上記著：「我為你的殿心裡焦急，如同火燒。」 \end{tabularx} \\ \\ \relax
2:18 & \begin{tabularx}{0.7\textwidth}{X} 因此猶太人問他：「你能顯甚麼神蹟給我們看，表明你可以做這些事呢？」 \end{tabularx} \\ \\ \relax
2:19 & \begin{tabularx}{0.7\textwidth}{X} 耶穌回答他們說：「你們拆毀這殿，我三日內要把它重建。」 \end{tabularx} \\ \\ \relax
2:20 & \begin{tabularx}{0.7\textwidth}{X} 猶太人說：「這殿造了四十六年，你三日內就能重建嗎？」 \end{tabularx} \\ \\ \relax
2:21 & \begin{tabularx}{0.7\textwidth}{X} 但耶穌所說的殿是指他的身體。 \end{tabularx} \\ \\ \relax
2:22 & \begin{tabularx}{0.7\textwidth}{X} 所以他從死人中復活以後，門徒想起他曾說過這事，就信了聖經和耶穌所說的話。 \end{tabularx} \\ \\
[1ex]
\hline
\hline
\end{longtable}
$^{1}$以下我們聽神的話.
請打開我們的新約聖經.
約翰福音第二章.
新約的一百二十七頁.
我們會中讀約翰福音第二章的十三至二十二節.
約翰福音第二章第十三節.
猶太人的雨接近了.
耶穌就上耶路撒冷去.
看見店裡有賣牛羊甲子的.
並且並有兌換銀錢的人坐在那裡.
耶穌就拿繩子作成鞭子.
把牛羊刀趕出店去.
倒出兌換銀錢的人的銀錢.
推翻他們的桌子.
有對賣甲子的說.
把這些東西拿去.
不要把我付的電當作買賣的地方.
他的門徒就想起經上記著說.
我為你的電心裡焦急如同火燒.
因此猶太人問他說.
你寄作這樣的奢侈.
還顯甚麼神蹟給我們看呢?.
耶穌回答說.
你們拆毀這電.
我三日內要再建立起來.
猶太人便說.
這電是四十六年才叫做成的.
你三日內就再建立起來嗎?.
但耶穌這話是以他的身體為電.
所以到他從死裡復活以後.
門徒就想起他說過這話.
便信了聖經和耶穌所說的.
請你讀筆.
弟兄姊妹平安.
我們一起來繼續.
在十一月份聚焦在基督耶穌的身上.
這幾次的主題我們都說了.
我們會用約翰福音作為一個很重要的藍本.
從約翰福音裡面的寫作方式.
來了解一下耶穌基督是誰.

$^{41}$當然在座裡面.
你信了主的你都知道耶穌就是救主.
這是絕對是對的答案.
不過再深入一點.
或者再有些字眼.
或者在聖經裡面怎樣將耶穌介紹出來.
其實我們值得再了解深入一點.
當我們了解深入一點的時候.
其實會為我們的信仰.
會建立一些更加深厚的根基.
這些根基會影響我們如何敬拜神.
這些根基會決定我們如何傳福音.
今天我們要看的是關於.
耶穌說自己指著身體為殿的這段經文.
一般來說很多人處理這段經文.
通常入首點就是說耶穌很生氣.
耶穌趕走一些人在聖殿裡面.
耶穌好像跟平時傳福音的那種很謙柔.
或者那種謙謙君子的形態很不同.
很多時候會用這個角度來入位.
不過今天我們處理這段經文.
我們不是從耶穌的情緒來入切點.
我們從耶穌所說的話和門徒去明白.
耶穌指著身體為殿為聖殿.
這個角度去重新理解.
究竟按福音介紹的耶穌是一個怎樣的耶穌.
雖然今天好像比較神學一點.
希望大家不會覺得很奸詐.
不要讓大家覺得睡眠不足的時間到了.
希望大家真的可以追著來聽.
讓我們更加深入了解.
究竟福音裡面所說的耶穌是那個殿.
究竟是一回什麼的事.
在這段時間我們很需要聖靈幫助我們一起祈禱.
求聖靈幫助我們開啟我們的心靈和心祈禱.
天父上帝求你在我們當中幫助我們.
讓我們能夠明白你的話語.
讓我們能夠真切準確地理解你是一個怎樣的神.
求你幫助我們眾弟兄姊妹.
讓我們去聆聽你的時候.

$^{81}$你親自對我們每一個人說話.
又幫助宣講的說得清楚.
讓你的經文和你的話語成為我們生命的幫助.
恭敬將來的時間交託.
願主繼續來對我們賜福.
獻上祈禱奉靠基督耶穌得勝的聖名.
Amen.
我們一起看看之前跟大家讀過的.
約翰福音第一章第十四節裡面有一節的經文是這樣說.
大家會很熟悉的.
道成了肉身 住在我們中間.
充充滿滿地有恩典有真理.
我們也見過他的榮光.
正是負獨生子的榮光.
這個約翰福音第一章十四節裡面.
約翰很斬釘截鐵地介紹耶穌是誰.
他講到耶穌是道成肉身的神.
接著值得留意的是兩句短句.
住在我們中間.
充充滿滿地有恩典有真理.
我們也見過他的榮光.
正是負獨生子的榮光.
住在我們中間 見過耶穌的榮光.
究竟約翰在說什麼呢?.
我們經常都念口黃 很熟悉這段經文.
不過你想想 耶穌住在.
在說道成肉身的主 住在他當中.
可能也一起生活 是不是?.
但見過耶穌充滿榮光.
將上帝的榮光顯現出來.
其實耶穌真的發光.
在約翰福音沒有記載過.
頂曬窿都是他的門徒登山變相的時候.
見到耶穌登山變相的時候發光.
其他的時間都是平常人一個.
住在他們當中.
但為什麼約翰要用這句話來形容耶穌呢?.
首先住在我們當中.
住在我們中間的住的意思.
其實舊約或希臘文的字根.

$^{121}$是在說回眸的字.
耶穌好像回眸一樣.
行走在我們人世間裡面.
這個住字是這樣的意思.
榮光就是說到舊約裡面敬拜神.
回眸上帝顯現出來的時候.
那種榮耀發出來.
代表著上帝和一群神的子民同在.
上帝的榮光充滿在回眸裡面.
充滿在敬拜當中.
所以住在他們當中和發出榮光.
這是在說舊約裡面.
上帝藉著回眸和以色列同行.
在出埃及或是在曠野的時候.
上帝和他們同在的一個名證.
上帝住在這群人當中.
除了神跡奇事之外.
更加是上帝將他自己的榮光.
顯現在敬拜生活當中.
接下來當以色列人到了迦南地.
立國了.
大衛成為了君主.
當大衛住在香柏木.
還有一個很漂亮的宮殿的時候.
大衛說我自己住在這麼漂亮的宮殿.
以色列已經立了國.
他就成為了一國之君.
他自己住在這麼漂亮的皇宮.
上帝還住在回眸裡面是不行的.
於是大衛王就說.
不如這樣吧,我要做一個宮殿.
做一個聖殿給上帝.
讓上帝來居住在當中.
同樣的字眼,那個住字和榮光出現.
就在聖殿的建造裡面.
當然上帝和大衛有一番對話.
大衛想做一個聖殿給上帝來居住.
但是上帝收到大衛這個心意.
他一開始在聖經裡面說的不太好.
不是說想做一個宮殿給他.

$^{161}$上帝說好,我住進去.
不是這樣的.
上帝說我是萬有之主.
天地都容納不了我.
你給我一個宮殿吧.
我怎麼會住在人手所做的宮殿裡面呢.
後來一路再談的時候.
大衛的神學開始明白了.
然後他說我知道人手所做的.
是不能夠給你居住的.
你可不可以真的.
閱立我心裡面的這種敬拜.
於是上帝就閱立了這個心意.
不過就說大衛的手流人的血.
不能夠做一個這樣的宮殿.
就讓他的兒子所羅門來建造這個宮殿.
所以那個聖殿其實嚴格來說.
那個idea(想法)是由大衛來想的.
真正建殿的就是所羅門來做.
如果你想一下整件事是怎樣呢.
就好像你小時候的家.
我小時候的兒子也是這樣.
在家裡拼了一個LEGO.
拼了一個房子.
差不多一大塊的,拼得很漂亮.
然後說爸爸,這間房子給你住.
這間士多房給你住.
我兒子也是這樣說.
這間士多房給你住.
那我那間主人房呢.
是媽媽住的.
這個世界就是這樣.
很不公道的.
一間士多房住也好.
這麼大的人.
我怎麼住進去那間LEGO拼的士多房呢.
雖然他有這樣的心.
但是在實際裡面是做不到的.
同樣道理.
一個創造天地萬物的主.

$^{201}$怎麼會做在人手所做的一個建築物裡面呢.
其實是不容許的.
不過上帝閱立了大衛和所羅門的心意.
所以第一個聖殿.
所羅門建造的時候.
上帝真的閱立了這個心意.
在歷代志下.
印了給大家的經文.
六章十八節.
所羅門也是這樣向神祈禱的.
他說:神果真與世人同住在地上嗎?.
看!天上和天上的天.
尚且不足你居住的.
何況我所建的這殿呢?.
在這裡.
所羅門拿捏得準的神學.
整個天都不夠上帝所居住.
天外有天都不夠上帝居住.
一個區區的人物.
在地上建一個建築物.
怎麼夠上帝來住呢?.
上帝你是不是真的會住在這裡呢?.
上帝回覆所羅門的祈禱.
在歷代志下第七章裡面.
當所羅門祈禱完畢的時候.
就有火從天上降下來.
燒盡凡濟和別的濟.
耶和華的榮光充滿了電.
因為耶和華的榮光充滿了耶和華的電.
所以濟事不能進電.
這是歷代志下七章一至二節.
在這裡說什麼呢?.
所羅門說:天都不夠你居住.
你是不是真的住在人的小屋裡呢?.
上帝的回應是.
閱納了一個祭物.
上帝將榮光充滿了耶和華的電.
然後就用這個形式來回應了所羅門的祈禱.
是,他住在他們當中.
榮光充滿.

$^{241}$將聖殿成為一個上帝閱納的地方.
然後七章十二至十六節.
就是一個很重要的一段祈禱.
這個如果你可以的話.
我希望大家在家裡背了它.
耶和華向所羅門顯現.
對他說:我已經聽了你的祈禱.
也選擇這個地方作為祭祀我的電雨.
我若使天閉塞不下雨.
或使蝗蟲吃這地的出產.
或使瘟疫流行在我民當中.
這稱為我民的子民若是自卑.
禱告尋求我的面.
轉來他的惡行.
我必從天上垂聽.
赦免他們的罪.
醫治他們的地.
我必睜眼看.
則而聽在此處所獻的祈禱.
現在我已經選擇這電.
分別為聖.
使我的名永遠在其中.
我的眼我的心也必常在那裡.
這是歷代志下的七章十二至十六節.
當上帝的榮光充滿了在所羅門建造了電雨的時候.
上帝再一次很清晰的.
除了榮光充滿了電之外.
他還說話告訴所羅門.
我閱立這個地方.
我的心我的眼會在這裡.
我會選擇這個地方.
我會以這個地方分別為聖.
我要聆聽神的子民祈禱.
這個祈禱的重要內容是什麼呢?.
如果有禍患在這個國家裡面.
在這群人當中.
這些人去到電裡面重新找回上帝.
自卑.
尋求上帝的面.
在他面前來承認自己的罪.

$^{281}$和轉離他的惡行.
這些真誠的禱告.
這種悔罪的禱告.
上帝必然垂聽.
上帝必然會止住這些災禍.
赦免他們的罪.
去醫治他們的地.
所以這個聖殿.
最重要的內容.
最重要的目的.
其實是聆聽子民的祈禱.
上帝閱立了這個地方.
他要將他的榮光.
好像合法化了這個地方.
是上帝所指定的敬拜場所.
這個地方是上帝所命定的.
上帝OK.
雖然神學上是說不出的.
不過上帝說OK啦.
照住在這裡.
不過最重要的就是.
這個地方是神與人之間.
連結的一個很重要的地方.
在神學上來說.
這個聖殿.
和天上的符是最接近的.
這裡有一條專線.
直接接駁到天符那裡.
聖殿就是一個這樣祈禱的地方.
所以在舊約裡面.
所羅門所建的聖殿.
是一個禱告的殿.
君王,祭司,以色列民.
就藉著這個地方.
來將自己的禱告.
達到上帝的根前.
不過上帝說的話.
不是停在這裡.
不是說他們認罪,聆聽祈禱.
19,20節那裡才是重要的.

$^{321}$歷代志下.
倘若你轉去.
驕氣我指示你律例誡命.
去事奉敬拜別神.
我必將以色列人.
從我所賜給你們的地上.
拔出根來.
並且我為幾名.
所分別為聖的殿.
也必捨棄不管.
使他在萬民中作笑談.
避饑憔悴.
歷代志下第七章.
19-20節.
他講到就是說.
這個殿是閱立去祈禱的.
不過如果這班人繼續.
一意孤行去敬拜別神.
離開上帝.
上帝說我不要這個地方.
我驕氣這個聖殿.
因為這個聖殿.
對上帝來說.
是人自己想去.
敬拜神的一個地方.
上帝認可這個地方.
但如果這班人的心不在上帝那裡.
有這個殿是沒有用的.
沒有意思的.
所以上帝說如果你們.
去驕氣上帝.
敬拜別的神.
這個聖殿上帝都會驕氣.
令到這班以色列民.
作為笑談.
被譏笑被愛俗.
來踐踏.
這個是上帝在聖殿.
見完之後.
所羅門的祈禱.

$^{361}$的回應.
上帝就是一個這樣的教導和回應.
在歷史裡面這個第一聖殿.
在所羅門之後.
有一段輝煌日子.
金碧輝煌.
很多人都無名而至.
以色列民都會敬拜.
不過這個聖殿.
存在短短.
幾百年的時間.
去到主前.
587年.
大約這個年份.
以色列民因為.
一意孤行離開上帝.
所以整個以色列民就.
被擄巴比倫.
在他被擄巴比倫的時候.
你不要想著這班人走了.
巴比倫大軍都將這個聖殿.
遊鄰.
破壞它 擄掠裡面最值錢的東西.
將金.
將燈台.
或者最值錢的東西搬走.
連藥櫃都搬走了.
所以藥櫃放在正殿的藥櫃.
都不知道去到那裡 最重要的東西都沒有了.
這個聖殿就被毀了.
這個就是主前.
587年.
以色列民因為離棄上帝.
所以上帝真的.
按他所說的話.
就從此.
第一個聖殿就這樣.
破壞了.
去到.
雖然說這個聖殿被破壞.

$^{401}$但上帝其實心目的.
不是要令到以色列民永遠都是.
陷於一種痛苦邊緣.
上帝說這班以色列民.
被擄70年之後.
他要回到自己的地方.
重回.
去到耶路撒冷.
回到自己的地方.
上帝真的這樣走了.
去到70年之後.
另外一個建殿的猶太省省長.
就是君王的後裔.
索羅巴伯.
馬代波斯王.
古列.
不是踢球的古列字.
古列 波斯王.
他就發諭旨.
讓索羅巴伯可以回去重建.
這個聖殿.
這個就是第二個聖殿.
在主前的516年.
開始重新.
去修造.
不過你可以想像.
一個被擄的子民.
其實沒有什麼資源.
所以重新修造的這個聖殿.
其實等於原本的聖殿.
所羅門的聖殿.
只有三分之一的規模.
但是也能夠有一個敬拜.
屬於自己的地方.
主前516年.
這個就是他們.
所認定的第二個聖殿.
在這個第二個聖殿裡面.
一直有500多年的日子.
就成為.

$^{441}$以色列民去敬拜神的地方.
不過當以色列民.
有了一個這樣的.
第二聖殿的時候.
他們就繼續跟隨上帝.
在《先知書》裡面說.
不是的.
《馬拉基書》裡面說.
聖殿都沒有人看管.
那些祭司很敷衍.
在《馬拉基書》裡面說.
沒有人看管的.
全部都是做祭祀.
大家不是用心去做的.
所以上帝在《馬拉基書》裡面說.
你把門關上.
你的教會你的聖殿.
你都不是敬拜的.
關上門不要讓人進來.
你自己勾.
上帝說了這麼晦氣的話.
原因就是一班以色列民.
即使被擄有第二聖殿.
他們的心都不是放在上帝那裡.
他們都是表面有一些儀式.
心靈和上帝離開的.
所以聖殿.
第二個聖殿.
都被毀在歷史裡面.
在主前的168年.
被安提亞古四世.
這個敘利亞的王.
就一直.
因為要希臘化了這個.
以色列民.
所以第二個聖殿都被毀.
都是處於一種.
很慘烈的階段.
這一次被毀比上一次更慘的是.
這一次被毀之前.

$^{481}$安提亞古這個王要褻瀆信仰.
要羞辱以色列人.
他把另外一個外邦的神像.
放在第二個聖殿裡面.
來叫人敬拜.
說是一個很大的褻瀆.
但是.
以色列人自己都放棄敬拜神.
所以更大的羞辱.
就在這些外邦人裡面.
教訓了一班的以色列民.
接下來就有很長的日子.
以色列民連聖殿都沒有.
但是你剛才讀經文.
去到耶穌時代.
為什麼會有一個聖殿呢?.
原來去到主前第37年.
耶穌出生37年.
那個大希律.
羅馬帝國的大希律.
為主前第37年.
大希律.
為主前第37年.
討好那些以色列人.
於是就重新給錢.
來建立一個.
以色列民.
或者猶太人的一個宮殿.
讓他們敬拜上帝.
你不要以為.
大希律是一個很敬虔的人.
讓他來敬拜.
其實不是這樣.
很多歷史學家找到原因.
其實做一個這樣的聖殿.
製造就業機會.
讓猶太人有份工作.
來做.
聖殿建好之後可以收聖殿稅.
這是為當時的.

$^{521}$羅馬帝國.
方便徵稅的一個方法.
所以去到耶穌.
那個年代的聖殿.
其實是由.
外邦的大希律.
所興建的.
這個地方.
是一個商業中心.
幫助收取.
猶太人的聖殿稅.
和製造一些.
就業機會.
當然在猶太人角度來說.
這個聖殿.
他不當是大希律.
不是當外邦興建的.
是收緝的.
所以有機會你看回猶太人的點注.
他說這個是將第二聖殿.
所羅巴巴的聖殿重新收緝.
這是2.0.
第二聖殿的2.0.
他不會承認.
一個由外邦人所興建的.
來成為.
以色列人敬拜的中心.
反而是說這個是君王的後裔.
不過是外邦來幫他們.
重新收造這個地方.
說了很長的歷史.
說到很長.
說什麼呢.
聖殿你會看到.
本身是由人手.
所做一個敬拜神的地方.
是人手的產物.
不是上帝命定要做的.
在.
以色列民.

$^{561}$失去了敬拜神的.
那種內涵和精髓的時候.
這個聖殿.
就隨著他們.
失去了那種敬拜神的心.
就被.
外邦來毀滅.
歷史裡面.
可能說是.
敘利亞王.
或者是巴比倫.
破壞了這個聖殿.
不過在聖經裡面說到.
真正破壞聖殿的.
不是那些外邦.
那些外邦人只不過是上帝的奴僕.
真正破壞這個敬拜.
聖殿的.
是以色列人.
是那些猶太人自己.
是他們首先離棄了神.
首先他們的心離開了上帝.
以致到這個地方.
失去了內涵的時候.
那些宗族就被外邦.
所蹂躪.
在經文裡面更加說到.
這個聖殿雖然.
以色列人學了很多教訓.
但是他們都不是很掌握得到.
究竟他們的生命出了什麼問題.
反而這個聖殿.
他的位置越擺越高.
舊時有一個禱告的殿.
後來就成為一個民族的中心.
當這個國家.
以色列國猶太被滅的時候.
這個更加成為.
民族裡面的核心.
成為他文化的中心.

$^{601}$這個地方裡面就成為了.
很多不同人訴求聚集.
一班以色列人.
血統啊,文化啊.
生活,存留啊.
所有東西的中心點.
甚至去到耶穌時代.
這班人製造就業機會的一個地方.
所以是一個經濟的中心.
聖殿就扮演了.
不同的角色.
每一個角色似乎一路一路.
令到聖殿.
敬拜神的本質.
一路一路離開了.
去到剛才我們所讀的約翰福音.
解釋了這個聖殿的歷史之後.
我們再看約翰福音.
你就會知道.
為什麼耶穌會這麼憤怒.
耶穌的憤怒是在說.
這個聖殿本來是如月折近.
是一班人想去.
去守節,紀念上帝.
救贖的恩典.
不過耶穌在聖殿裡面.
看到什麼呢?.
看到賣牛,羊,鴿子和兌換銀錢的人.
坐在那裡.
這些其實是.
這些牛,羊,鴿子.
和兌換銀錢.
其實是方便了那些人來獻祭的.
製造了一個方便.
令到獻祭的人.
可以在聖殿裡面.
搞定這些東西.
一站式服務.
這是一個挺聰明的做法.
還有一些典著說.

$^{641}$這些一站式服務.
確保在.
殿裡面買的祭身.
一定合乎那個QC.
因為由祭司來鑑定.
是OK的.
你可以看到整條產業鏈.
在這裡.
整條產業鏈.
本身是方便到人有什麼問題呢?.
不過這個地方.
成為了做生意的地方.
有一些文獻更加說.
譬如那些甲子.
或者獻祭那些.
其實可以.
以色列民或者.
猶太人可以自己戴上去.
以前因為沒有規定一定要在殿裡面買.
不過那些.
祭司鑑定那些.
看著那些.
祭身合不合格的那些祭司.
他們有一個決定權.
當你自己戴一隻白甲.
去的時候那個祭司說.
有白甲眼的.
不是很合規矩.
這隻不穴立你放了它.
買我們這隻.
然後你調高來賣.
如果那個人上到去獻祭.
他被迫一定要買的.
整件事就變了.
就是.
淘汰了那些自己努力.
自己獻祭的心.
成為了一個只是集中在.
一小撮人裡面.
來賺錢或者建立特權的.

$^{681}$一個敬拜制度.
耶穌在這裡面.
見到這一班人.
見到這些兌換銀錢的人.
所以就大發雷霆.
姑且不論.
一些歷史因素.
或者剛才所說的.
那些例證.
是不是完全是真實呢.
但是可以肯定的.
一件事就是見到.
一個以祭身為中心.
要獻祭要獻什麼.
什麼祭身.
然後在這裡兌換銀錢.
你見到那個聖殿.
已經是亂成了一個買賣.
祭身的一個地方.
雖然本意是要獻祭.
不過那些人.
的著重點就不再是祈禱.
那些人.
那個重點就放在.
那些祭身合不合格.
那個祭身賣多少錢.
在這些這樣的.
交易上.
耶穌見到.
這樣的情況於是就趕走.
這些人為了電.
心急如焚如同火燒.
因為.
這個是本來作為.
禱告的電如今.
成為了一個買賣的地方.
耶穌再.
繼續的.
來到去說.
猶太人問他 哇 耶穌你做這些事.

$^{721}$怎麼獻神蹟啊 你這樣褻瀆.
聖殿 就是說.
某程度是去.
玩穿了聖殿.
的這個派對.
就是守節的時候 你搞到那些人.
雞飛狗走 搞到那些.
祭身都走開了.
做這些事怎麼獻神蹟啊.
你怎麼去代表神.
來作先知 來到.
獻神蹟給這些人看啊.
耶穌說.
你們拆毀這個電.
我三天內要再建立起來.
他講到那個.
這個電被拆毀.
不是真的.
建築物那種拆毀.
而是那個電的精神.
那種祈禱.
是被這些人所拆毀.
這些人.
表面上很重視.
這個敬拜生活.
實際上是很輕視.
這個敬拜生活.
在這一班人裡面就.
完完全全表露了那種.
那種矛盾位在這一班人身上.
很重視.
那種敬拜.
但活出來是一種很輕視.
那種狀態.
耶穌說這一班人是拆毀.
這個電的一班人.
當然這個拆毀這個電也都講到.
耶穌指著他自己身體為電.
就是他.
被這一班人所殺害.

$^{761}$被猶太人釘上十字架.
他講的.
也是一個這樣的.
一個的角度.
講到這個電被破壞.
是一班以色列人.
將他自己的.
敬拜生活.
破壞了.
以致到有諸淵內.
必營諸淵外.
將來這個電都一定被拆毀.
我們稍微聽完這麼多歷史.
或者一些講解之後.
我們一起去思考.
在我們生活裡面.
我們會不會有時候.
我們表面上有些東西可能很重視敬拜神.
但當我們去到某些位置的時候.
我們經常覺得很重視的東西.
到最後過了頭來.
我們很荒誕地.
令到我們對這個信仰變得很輕視.
我們可能很強調的.
某一些很外在的東西.
但最真正的內涵.
我們卻是忽略了.
講一個小小.
負面的小小的例子.
我們教睦同工.
幫很多弟兄姊妹.
做婚前輔導.
我們幫很多弟兄姊妹主持婚禮.
通常婚禮都是安排在禮拜六那裡.
如有雷同實屬巧合.
我們曾經幫很多對年輕的弟兄姊妹.
做安排婚禮.
預備所有東西.
幫他們打理他們的婚姻.
他們很著重在教會.

$^{801}$特別是新娘子會.
他們很著重在教會.
特別是新娘子會.
他們很著重在教會.
特別是新娘子會.
他們很著重在教會.
特別是新娘子會.
喜歡在教會裡行禮.
特別是新娘子會.
喜歡在教會裡行禮.
他們很重視.
當他們在禮拜六婚禮.
在教會裡行完婚禮之後.
在禮拜日崇拜.
很多時我都見不到.
很多弟兄姊妹回家崇拜.
很重視婚禮在教會裡進行.
但禮拜日未必很重視.
他們會繼續回來教會參與崇拜.
我和太太問我老婆.
我怕我罵人時會罵自己.
我問我老婆.
我們結婚後第二天沒有崇拜.
我問我老婆.
我老婆說當然有.
拖著一個快死的身軀.
喝完後都累死了.
但第二天回到自己的墓會崇拜.
他問我可不可以主任牧師.
牧師還在台上點名.
在教會報告中.
在台上點名.
我心想現在回想起來一額汗.
我記得牧師說.
我做這些事不是為了尷尬.
在教會行禮只是人生的一段時間.
真正重要的是一生敬拜主.
我作為牧者.
我希望提醒你一生要敬拜神.
不是為了一刻很大場面.

$^{841}$在教會行禮很莊嚴.
繼續一齊敬拜神.
你作為一個傳道人.
你也要教弟兄姊妹.
教導他們不單在教會行禮很莊嚴.
而是整個人生.
每一天的敬拜.
都要好好實踐敬拜的生活.
弟兄姊妹.
在這裡提醒我們.
今天我們會否只是集中.
在某一點時間.
我們覺得敬拜神很重要.
但在其他生活中我們放棄了.
或是我們覺得很著重禮拜日.
早課時間敬拜神.
但在我的生活中.
我都是過一個以自我為中心的生活.
那不是一個敬拜神的生活.
對於一個基督徒來說.
我們是全天候的.
在我的工作崗位.
在我的家庭當中.
在我的教會生活.
我們人生是整全地來對敬拜神.
不會是一刻時間.
我們集中這一點做好它.
就等於敬拜好上帝.
我們的生命是全然地交託給主.
那個生命才是被上帝所閱立.
而不是強調某一刻.
某一點.
某一點做好它.
那個執著就等於.
敬拜上帝.
在這裡.
猶太人的荒謬.
在耶穌裡面.
就嚴嚴地呈現出來.
耶穌說.

$^{881}$如果我們失去了這種敬拜神的內涵.
我們等同於.
拆毀這個殿無二.
我們和拆毀這個敬拜生活.
差不多.
耶穌再繼續說.
這個殿用四十六年才造成.
你三天就能夠建立起來嗎?.
但耶穌說.
這話是以他身體為殿.
所以.
到他從死裡復活以後.
門徒就想起他說過的話.
便信了聖經和耶穌所說的.
這是約翰福音2:20-22.
耶穌.
在這裡.
他要將一班.
弟兄姊妹.
或者一班神國的子民.
將他們的焦點重新調教.
在這裡.
約翰福音第一次就很清晰地說.
耶穌是指著自己的身體為殿.
這個殿已經不再是.
一個建築物.
在某一些節期裡面.
做好敬拜生活.
如果當耶穌以自己的身體為殿.
即是說.
耶穌以他自己成為人神之間.
連結的橋樑.
耶穌如果走到他們當中.
敬拜就不再拘泥在一個地方.
某時某刻.
而著重的就是我們奉誰的名.
我們以誰作為.
我們中間的橋樑來敬拜神.
在舊約裡面說的聖殿.
就是上帝承認的.

$^{921}$分別為聖的一個地方.
來敬拜神.
去到基督耶穌的降生.
上帝以基督耶穌.
那個不是人手所造的建築物.
而是神自己呈現了自己.
讓他成為人神之間的橋樑.
來敬拜上帝.
接下來的篇幅在約翰福第四章.
21至24節.
耶穌再清楚地跟.
撒瑪利亞婦人說.
耶穌說:婦人,你當信我,時候將到.
你們拜父的不在這個山.
也不在耶路撒冷.
時候將到,如今就是了.
那真正拜父的.
是用心靈和誠實.
來去拜他.
因為父要這樣的人拜他.
神是個靈,所以拜他的.
必須用心靈和誠實來去拜他.
那個誠實的字其實.
也不是很到位的.
現在所有的悉經學家.
都一致認同是心靈和.
真理來去敬拜神.
敬拜上帝的.
是用心靈.
用靈和用真理.
真理是甚麼?.
其實那個真理是指耶穌.
用靈和基督耶穌.
不是.
驅離在甚麼地方,甚麼形式.
甚麼時間來去敬拜.
而是直著靈.
直著真理.
我們每一個人都可以去到上帝面前.
跟上帝傾心吐意.

$^{961}$上帝一齊的來去祈禱.
去到耶穌的降生.
那個舊約裡面.
的聖殿已經被顛覆.
因為那個是人手所做的.
短暫的.
去到新約,去到耶穌的時候.
一個新的敬拜模式就出現了.
因為我們可以.
奉基督耶穌的名.
我們每一個基督徒就可以.
一齊的來向上帝.
來去禱告.
保羅再繼續用這個主題.
來講.
今天的聖殿在哪裡.
聖殿不是在講.
旺角浸信會.
六樓這個禮堂.
這個是一個.
敬拜的場所.
要安心出行才可以進去.
真正敬拜的殿.
是基督徒自己一班.
一班的關係.
教會,教會是基督的身體.
我們這一班人.
一齊聚集.
無論去到哪裡.
總之我們是奉主的名.
安著真理,直著基督耶穌.
我們就可以敬拜神.
我們可以一齊來到香神.
來到去傾心吐意.
一齊的來到去禱告.
哥林多前書三章十六.
至十七節.
或者哥林多後書六章第十六節.
保羅是用這個神學.
來講.

$^{1001}$基督徒是.
基督的身體.
我們這一班人的聚集.
就是奉基督的名.
一齊和上帝來到去連結.
直著基督耶穌.
和神來到去彼此結連.
不過在.
哥林多前書裡面.
雖然講的這個神學很勵志.
我們已經不再驅離.
而且好像很方便.
不過保羅去.
嚴嚴的.
哥林多前書裡面提醒我們.
基督徒這個群體.
你不要輕忽它.
既然這個是神聖的殿.
你不要任意來破壞它.
如果在這裡面有結黨紛爭.
破壞上帝的殿.
保羅很嚴厲的講.
上帝會毀壞那個人.
在這個關係裡面.
分化.
在這個關係裡面.
容讓罪.
在當中裡面.
容讓很多的私慾邪情.
在我們這個群體裡面出現的時候.
這個群體如果被瓦解.
上帝會找.
我們每一個人來算帳.
雖然我們今天好像很方便.
但上帝不是為我們製造.
一個方便的敬拜生活.
上帝是為我們.
我們今天不是.
將那個信仰.
驅離在一個地方.

$^{1041}$在一個時間裡面.
我們知道我們的生命.
全部是屬於上帝.
倘若我們生命裡面有罪.
分化這個群體.
上帝.
不容讓這些事情在我們群體裡面出現.
今日這一段經文.
我希望大家聽下去.
你理解聖殿.
約翰福音.
用聖殿來講耶穌.
不過他講耶穌.
他不是講舊約裡面的電.
如何應許.
不是用這個角度.
是講舊約裡面的聖殿有很多不足.
而且是人手所毀壞.
上帝真正設立新的電.
是耶穌自己.
藉著耶穌.
我們每一個人都可以向神.
直接的來聯繫.
今日教會就成為神的電.
我們這個關係.
是上帝為我們設立的.
我們珍惜這份關係.
我們彼此相愛.
就彰顯了這份.
是上帝的美意在我們當中.
若果我們相咬相吞.
我們就好像破壞了.
這個電裡面的一班人.
上帝隨時會找我們.
來算帳.
我們今日.
選擇走一條甚麼道路.
最後補充一個例子.
給大家知道.
耶穌所敬拜的那個電.

$^{1081}$就是耶穌所清除的.
這個2.0的.
大希律所造的電.
在主後70年.
就被提多羅馬將軍.
再一次來拆毀.
人手所造的電.
有一天都會歸於無有.
不過如果上帝所設立.
以耶穌基督為聖殿.
這個關係.
可以直到永遠.
求主繼續來幫助我們.
讓我們明白上帝的心意.
\newpage



\section{約翰福音 1:14-18}
\label{sec:Msw_I2sXTQc}
\textbf{2021年12月5日崇拜直播｜陳明泉牧師｜基督帶來的革新:新的恩典｜約翰福音1\_14-18}
\newline
\newline
連結: \href{https://youtube.com/watch?v=Msw-I2sXTQc}{\texttt{https://youtube.com/watch?v=Msw-I2sXTQc}} ~~~~ 語音日期: 2021-12-08
\newline
\newline
\hyperref[sec:N45H8jB07Gc]{\small{< < < PREV SERMON < < <}}
~
\hyperref[sec:index_chronic]{\small{[返順時目]}}
~
\hyperref[sec:Sd44cjxmurg]{\small{> > > NEXT SERMON > > >}}
\newline
\newline
約翰福音 1:14-18
\newline
\begin{longtable}{cl}
\hline
\hline
章節 & 經文 (和合本修訂版)\\
\hline
1:14 & \begin{tabularx}{0.7\textwidth}{X} 道成了肉身，住在我們中間，充充滿滿地有恩典有真理，我們也見過他的榮光，正是父獨一兒子的榮光。 \end{tabularx} \\ \\ \relax
1:15 & \begin{tabularx}{0.7\textwidth}{X} 約翰為他作見證，喊著說：「這就是我曾說：『那在我以後來的先於我，因為在我以前，他已經存在。』」 \end{tabularx} \\ \\ \relax
1:16 & \begin{tabularx}{0.7\textwidth}{X} 從他的豐富裡，我們都領受了恩典，而且恩上加恩。 \end{tabularx} \\ \\ \relax
1:17 & \begin{tabularx}{0.7\textwidth}{X} 律法是藉著摩西頒佈的；恩典和真理卻是由耶穌基督來的。 \end{tabularx} \\ \\ \relax
1:18 & \begin{tabularx}{0.7\textwidth}{X} 從來沒有人見過神，只有在父懷裡獨一的兒子將他表明出來。 \end{tabularx} \\ \\
[1ex]
\hline
\hline
\end{longtable}
$^{1}$今天早上要誦讀的經文是.
《約翰福音》一章十四至十八節.
請打開聖經新約部分125頁.
《約翰福音》一章十四至十八節.
道成了肉身 住在我們中間.
充充滿滿的 有恩典 有真理.
我們也見過他的榮光.
正是父獨生子的榮光.
約翰為他作見證 喊著說.
這就是我曾說.
那在我以後來的.
反成了在我以前的.
因他本來在我以前.
從他豐滿的恩典裡.
我們都領受了.
而且恩上加恩.
律法本是藉著摩西傳的.
恩典和真理都是由耶穌基督來的.
從來沒有人看見神.
只有在父懷裡的獨生子.
將祂表明出來.
聖經獨筆.
頂子妹平安.
我們今天繼續以《約翰福音》.
去了解究竟耶穌基督是哪一個.
連續有幾次.
在《約翰福音》一章第十四節.
我們其實重複又重複地.
來看這一句經文.
第一次我們說道成肉身.
以耶穌的真相來呈現在.
這個塵世的每一個人當中.
上個星期我們以聖殿的角度來看.
耶穌住在我們當中.
將祂的榮光來呈現出來.
今天我們只剩下.
那一小句的那句子句.
就是充充滿滿地.
有恩典有真理.
讀上去其實很順暢.

$^{41}$不過其實這句話.
這一句聖經是一句很特別的聖經.
希望今天我們從.
上帝一個新的恩典.
耶穌基督帶來一個新的恩典.
是一個怎樣的恩典呢?.
用這個角度來思想一下.
今天我們是否過著一個.
有恩典有真理的生活.
同樣我們明白耶穌基督的工作.
對我們整個靈性是很有幫助的.
希望今天可以幫我們.
更加深化我們的生命.
我們直著思想耶穌基督.
讓我們生命裡面更加踏實.
更加知道.
啟示出來的聖經.
耶穌基督真是無與倫比.
真是一個生命裡面很重要的救主.
我們一起祈禱.
求聖靈幫助我們.
讓我們明白話語和心祈禱.
阿巴天父求你在我們當中.
你對我們每一個人說話.
幫助我們讓我們能夠直著你的話語.
能夠真是明白基督耶穌的救恩.
和對生命的塑造有多重要.
求父你在我們當中.
帶領我們眾人.
幫助眾弟兄姊妹.
讓我們能夠在聆聽的過程裡面.
我們每一個都聽得明白.
又讓孫講的.
忠心和清晰的將你的話語.
來講論出來.
讓弟兄姊妹一同的領受.
生命同得造就.
我們恭敬將來的時間交託給你.
願基督耶穌你親自來帶領我們眾人.
獻上祈禱奉基督耶穌得勝的聖名.

$^{81}$Amen.
好,我們繼續去看約翰福音.
在講約翰福音一章第十四至十八節.
很多的識經學家.
或者很多的神學家都會同意的.
約翰不是抽空這樣.
純粹想像耶穌是哪一個.
然後寫這個聖經出來的.
大概你要明白.
使徒約翰或者第一代的信徒.
他們看的聖經.
他們看的是什麼內容呢?.
除了經歷耶穌基督親身的教導之外.
他們是看舊約聖經的.
而舊約裡面.
特別是摩西五經.
是成為了當時使徒裡面.
去將對比.
或者將上帝的應許實現.
是一個很重要的素材.
他們藉著默想.
舊約聖經的一些內容.
來幫助他們理解基督耶穌.
我們今天也不例外.
我們也是用這個角度來看.
我們先看一段經文.
《出埃及記》第34章第6節.
在這裡跟那個恩典與真理是有關的.
《出埃及記》第34章第6節.
我讀給大家聽.
耶和華在他面前宣告說.
耶和華,耶和華,是有憐憫有恩典的神.
不輕而發怒,並有豐盛的慈愛和誠實.
這句經文.
無端端為什麼上帝會說他這樣的話呢?.
這句話是上帝自我的宣示.
在摩西的面前.
為什麼會說他這樣的話呢?.
其實在《出埃及記》的上文下.
是在說一群以色列人做了一個金流毒.

$^{121}$上帝很生氣,懲罰了他.
懲罰完之後.
上帝再一次藉著摩西.
要將西乃山這個約頒布給一群以色列人.
上帝要求這群以色列人進入迦南地.
原本經文上文下.
上帝說你們進去吧,我保守你們.
我不和你們進去.
我不知道你們是怎樣的人.
在《出埃及記》裡面這樣說.
不過摩西夫婦這樣去求上帝.
於是上帝就說.
好吧,我就陪你們一起進去.
在頒布將要進入迦南地之前的西乃之約.
上帝就用了這句話作為自我宣稱.
自我宣稱為什麼跟這段經文有關呢?.
特別留意內容.
耶和華說:.
「祂是有憐憫有恩典的神.
不輕而發怒.
並有豐盛的慈愛和誠實」.
豐盛的慈愛和誠實.
這是和合本裡翻譯的那句話.
其實如果你用那句話.
你回顧剛才我們看的約翰福音一章十四字.
「充充滿滿地有恩典有真理」.
其實你看看那個意思.
其實是很一致的.
是那樣東西來的.
是同一個字眼來的.
不過去到和合本.
或者我們自己.
因為有舊約的希伯來文.
和希臘文的語言上的差距.
所以我們覺得好像有一點分別.
所以在這裡說的那個恩典.
有恩典有真理.
其實是說有恩典及真理.
而且有豐盛的恩典和真理.
這是上帝描述自己的一個形容詞.

$^{161}$這個形容詞不是隨便用的.
這個形容詞是上帝用這句字眼.
來說上帝的作為.
上帝納若.
祂都是指著自己.
我一個怎樣的神.
有恩典有真理的神.
現在和你們一起的.
來做這個西乃之約.
約翰是用一個這樣的神學觀念.
來說耶穌基督是一個怎樣的神.
除了祂是一個道成肉身.
真理成為肉身之外.
祂是一個新的聖殿.
住在我們當中.
有榮光光照出來.
祂更加是一個神治我的宣稱.
祂是充滿了恩典和真理.
更加重要的是.
約翰福音裡面寫這句說話寫得很小心.
因為這句說話的恩典字.
我們經常說的.
我們念口旺的.
我們有很多恩典.
我們經常都.
基督徒拐口唇邊的字眼.
對於約翰來說.
這不是一個拐口唇邊的字眼.
是一個很謹慎的字眼.
整個約翰福音.
只有一章十四.
十六和十七節才出現恩典的字.
二十多章的經文.
是在引言的部分才說恩典這個字.
而放在一個引言的部分.
接下來他怎樣去演繹.
但是他都盡量不會將這個恩典的字.
成為一個念口旺.
或者拐口唇邊的一個慣性的用語.
在約翰的處理裡面.

$^{201}$這是一個很謹慎的字眼.
這是在說神自己的宣稱.
是上帝自己的工作.
他對人的那種心意.
更加重要的就是.
那個恩典和真理這一句的子句.
其實是放在第一章十四節裡面.
最後的那一部分.
如果你用真的比較好的翻譯本.
呂鎮忠的譯本是這樣的.
他的次序.
道成了肉身住在我們中間.
我們也見過他的榮光.
正是乎獨生子的榮光.
充充滿滿地有恩典有真理.
如果你這樣看.
你會容易明白很多.
充充滿滿有恩典有真理.
插在這裡又見到榮光.
他將這個充充滿滿有恩典有真理.
放在這個句式最後的部分.
他要告訴大家的是甚麼呢.
這個上帝道成肉身.
有榮光住在我們當中.
這個上帝是有恩典和真理.
在這裡我們稍為知道了.
處理了經文裡面.
其實有這樣的呈現出來.
我們再看多一點.
就是出埃及的.
剛才給大家看到的那一節經文.
講到上帝說自己的宣稱.
他是一個有憐憫有恩典的神.
不輕易發怒.
且有豐盛的慈愛和誠實.
慈愛和誠實.
就對應了恩典與真理.
在誠實這個字.
在和合本裡面譯.
上帝不說謊話.

$^{241}$我們有時候會這樣理解錯誤.
或者我們見到字眼會這樣想.
其實那個誠實.
或者那個真理本身.
就是講到上帝是一個很真實.
信守承諾的上帝.
上帝是一個守約思前愛的神.
上帝講過出來的說話.
他講得出就做得到.
他答應了以色列民.
是一個這樣的英雄.
無論以色列民經歷什麼都好.
因為這是一個立約的關係.
所以這班人即使叛逆也好.
上帝都要將這班人挽回正軌.
得到上帝賜下的祝福.
誠實更加好一點的翻譯.
就是信實.
上帝是一個信實的主.
上帝答應過人的事.
他不會失言.
他不會說了就不算數.
反而是上帝答應了人的.
答應了他的子民的.
他對每一個世上的人作出的應許.
他都要一一地實現出來.
如果你用這個角度去看.
其實不難理解的.
在《若行封印》裡面講到.
耶穌就是舊約裡面自我宣稱的.
有憐愛,有恩典.
也是信實的主.
去到新約基督耶穌.
他就將上帝的恩典和信實.
藉著他的生命.
來呈現在人世間的生命當中.
當然你理解完這裡之後.
對外有什麼關係?.
牧師好像很抽象.
在講神學宣稱.

$^{281}$我們再看多一點.
在這段經文裡面第十四節.
和第十七節.
同樣都會有一些字眼.
很細微不過是很重要的.
十四節那裡講到.
道成了肉身住在我們當中.
充充滿滿有恩典有真理.
我們也見過他的榮光.
見到這個榮光.
第十七節或者第十八節.
從來沒有人看見神.
只有在乎的腐壞裡的獨生子.
將祂表明出來.
看見是一個重要的字眼.
上帝的光在舊約裡面.
不會讓人看見.
一看見就即死.
因為人是有罪的人.
看見上帝的榮光.
所以舊時在回望裡面.
最聖潔的人才能進入至聖所.
或者最內室裡面.
和上帝相遇.
不是面對面看上帝.
摩西和上帝那麼近.
他也要用一塊柏子來遮住.
這個榮光不是所有人都看見.
原則上沒有人可以看見上帝的榮光.
不過去到新約.
去到基督耶穌的道成肉身.
人可以看見他的榮光.
牧師你講的東西也很抽象.
見到耶穌的榮光.
那是甚麼呢.
經文再繼續說.
其實摩西也在問這條問題.
我們再繼續看出埃及的三十三章.
十七節到二十節.
我也印了給大家那段經文.

$^{321}$出埃及的第三十三章.
十七節到二十節.
耶和華對摩西說.
你在這所求的我也要行.
因為你在我眼前無了因.
並且我按你的名認識你.
摩西說求你顯出你的榮耀給我看.
將你的榮光給我們看看吧.
上帝摩西這樣求.
十九節.
耶和華說我要顯出我一切的恩慈.
在你眼前經過.
宣告我的名.
我要恩待誰就恩待誰.
又說你不能看見我的面.
因為人見我的面不能存活.
這裡說的是上帝的榮光.
摩西說給我一點榮耀.
讓我知道你是神.
讓我知道有些確距.
讓我可以經歷大靈神的指紋.
我們繼續前行.
一個很重要的確距.
除了會幕顯現之外.
可不可以讓我看看.
套下你的光.
上帝說我顯出的是我的憐憫.
不會顯出他的榮耀.
又或者用另一個角度說.
上帝的榮耀是以上帝的憐憫呈現出來.
上帝不會無端滿足人的那種求知或好奇心.
來將一個很魔法的狀態.
發光地呈現出來.
讓你看到他的樣子.
上帝說他以他的作為.
以他憐憫的作為.
將他自己給弟兄姊妹知道.
在舊約沒有人看見上帝的榮光.
不過去到基督耶穌的處境.
或者到城肉身的時候.

$^{361}$耶穌就將他自己的憐憫.
他對人是怎樣對呢.
就反映了上帝對人的那種心腸.
藉著上帝的作為.
將榮光呈現出來.
所以在約翰福音裡面.
榮耀這個字其實和神蹟.
其實是記號這個字是相提並論的.
耶穌每做完一個神蹟.
加拿大第一神蹟做完之後.
他講到顯出上帝的榮耀.
神蹟的目的目.
本身是要講上帝的榮耀.
在耶穌的世代上彰顯出來.
在你和我的生命裡面彰顯出來.
我們為什麼要有神蹟.
為什麼有些東西不是人可以做到的.
是神的作為.
這個就是將神的榮耀.
將神的榮光呈現在這個需要憐憫的大地裡面.
我們每一個就是經歷神的恩典.
在約翰福音裡面用這個角度來告訴人.
上帝的憐憫.
藉著基督耶穌.
藉著耶穌所行的神蹟.
告訴這個世界.
我們必須要將他的榮耀.
一一來顯現出來.
耶穌基督.
或者由摩西到耶穌基督.
上帝一直的作為都是這樣來做的.
不過耶穌有什麼地方會更加超越呢?.
約翰福音一章十六節裡面就這樣說.
在他豐滿的恩典裡我們都領受了.
而且恩上加恩.
律法本是藉著摩西傳的.
恩典和真理都是由耶穌基督來的.
從來沒有人看見神.
只有父懷女獨生子將祂表明出來.
耶穌比摩西更加超越的.

$^{401}$首先在這裡裡面.
他講到我們正在過一個恩上加恩的一種生活.
恩上加恩是一個什麼樣的狀態呢?.
後面裡面說.
他說律法是藉著摩西傳的.
恩典和真理是藉著耶穌基督傳的.
同樣都是翻譯上我們不太明白.
原來他說的律法也是一個恩典.
律法是藉著摩西賜予給大地的.
所以在聖經裡面就說到.
一群橫衝直撞,一群亂來的人.
因為上帝的律法.
令他們可以走在上帝的心意.
整個以色列文體律法是一個恩典.
不是一個純粹有些規條逼你去做.
今天如果小學生或中學生拿一本手冊.
看到學生規條.
你不會覺得那是學校想獎給你的禮物.
不過在聖經裡面.
他顛覆我們的想法.
如果我們自己亂來過生活.
我們人生帶來的就是痛苦和失望,悲哀.
但是如果有上帝的指引過我們的生活的時候.
我們的生命就可以過得有秩序.
不會過一個任意而為生命裡面痛苦的一種經歷.
律法對於以色列人來說是一份禮物.
建基於在這個恩典當中.
耶穌再要再多一層.
就是恩典與真理的完美結合.
就在耶穌基督的生命當中.
所以摩西舊時帶來的那個律法.
就是鼓勵以色列人實踐一個信服的生活.
但是到了基督耶穌.
恩典與真理就是過著一個信心的生活.
我們今天不是要廢除律法.
不過我們今天我們知道.
我們的生命不只是只有律法.
就是只有一些規條.
還有一大堆要我們遵守的一些條例.
耶穌就要帶來一個新的恩典.

$^{441}$就是直著他的生命.
我們可以將恩典與真理好好地平衡.
來生活.
約翰在這裡將摩西與耶穌做一些比較.
他不是要貶低摩西舊時律法的生活.
他說如果摩西帶來的律法.
這個恩典生活.
在以色列人當中他們都覺得這麼棒的時候.
他們都這麼重視的時候.
更加超越這個生命就是耶穌自己.
以色列人更加要擁抱這種生活.
我們就是活在一個.
以基督的生命來呈現出來的那種恩典與真理的平衡.
耶穌基督帶來一個更超越.
更大的恩典.
在每一個弟兄姊妹.
在每一個基督徒的生命當中.
說了很多神學論述.
可能你會覺得抽象.
我稍微總結一下這一句話.
恩典與真理.
這句話其實是上帝自己的呈現.
約翰用恩典與真理這一句字眼.
來講到我們所信的耶穌是神自己.
而且這個恩典與真理這個恩典.
比舊約裡面的恩典更加超越更加優勝.
舊約裡面摩西所講的是一個順服律法的生活.
不過去到基督耶穌.
是在講一個恩典與真理平衡的生活.
藉著基督耶穌.
我們不是要靠行為來得救.
不是靠著順服來得到生命的豐盛.
我們單單純靠耶穌.
藉著基督耶穌.
我們就明白上帝的心意.
我們就活出真理與平衡的生活.
講到恩典與真理的平衡.
或者我們對於恩典的觀念.
有時候我們基督徒會搞錯了.
早前有機會認識一些情侶來聊天.

$^{481}$那個女孩發脾氣.
找我們聊天.
找牧者.
發脾氣是什麼呢.
她說我的男朋友送了一束花給我.
送了一束花也發脾氣.
她說男朋友不懂得做.
為什麼呢.
因為忘記了周年紀念.
男孩很用心地預備了一束玫瑰.
送給女朋友.
當他送這束花給女朋友的時候.
女孩收到這束花.
很開心 很漂亮.
因為心花怒放.
但這個女孩心水星.
拿了束花之後.
也數了一下多少枝玫瑰.
她真的數了.
送了多少枝呢.
送了11枝.
如果你懂得送花的話.
很浪漫的.
因為送11枝玫瑰.
第12枝就說那不是玫瑰.
是我自己的.
但女孩似乎不太知道這些.
她就說男朋友.
為什麼你送11枝玫瑰給我.
差一枝你也不買給她.
你這麼孤寒.
男孩就說那枝就是我自己的.
我將我全人送給你.
女孩就板著臉.
說你這麼孤寒.
她心裡想.
這個男孩這麼沒趣.
這麼孤寒.
第二樣就覺得.
這個男孩所說的那句話.

$^{521}$將自己送給她.
她覺得是不是這樣.
這麼噁心 這麼老套.
將自己送給別人.
你孤寒就孤寒.
還要找個藉口.
於是就生了男朋友的氣.
送了東西給他之後.
還要生氣.
還要發脾氣等等.
還要不開心找牧者談.
這件事我們怎樣處理.
我不多說了.
不過這個女孩所呈現的.
其實很多時候是我們基督徒的生命.
上帝給了一個恩典給我們.
我們接納了.
有時候將一些.
不要說次要.
不是終極的東西.
都賜了給我們.
可能我們生活裡有富足.
我們生命裡擁有一切美好恩典.
我們收到的時候.
我們感恩 我們很快樂.
但到最後上帝說.
他將自己的獨生子賜給你.
他將自己送給你.
你就會很實際地說.
哪有人這樣說話的.
這麼老套.
有時候我們會覺得.
我們寧願不要那個賜福的源頭.
我們要上帝給我們.
我們覺得生活裡有好的東西.
如果上帝說將他自己.
送給我們的時候.
我們覺得很奇怪.
甚至我們心裡 口裡.
可能都會說.

$^{561}$口裡說我尊稱耶穌為主.
但我心裡就是想要多些上帝.
來伴隨我們生活裡一切的恩典.
我們要生活裡令我們快樂的東西.
我們寧願有這樣的態度.
我們都不要真正的主.
上帝將他送給我們的時候.
我們拒絕.
在《約翰福音》裡說.
上帝是一個賜福的源頭.
上帝是一份禮物.
送給這個世界每一個弟兄姊妹.
送給每一個上帝.
所預備恩典接收的人.
神愛世人.
甚至將他的獨生子.
賜給我們每一個.
叫一切相信他的不致滅亡.
反得永生.
不過我們的生命.
我們的運作是怎樣呢.
我們要上帝獨生子所賜下.
給我們生活的恩典多些.
我們不是很要耶穌.
我們生命裡面.
很多時拒絕生命恩典.
最重要最大的禮物.
我們拒絕它.
但我們要上帝給我們的獎金和獎品.
不是要上帝給自己.
在《聖經》裡面.
約翰講到耶穌是一個最大的禮物.
送給我們.
不過世人不接受.
世人拒絕.
即使很多人.
生命裡面很需要這個主也好.
他都想走回舊時的道路.
想去摩西.
想去舊時的次好那份禮物.

$^{601}$來繼續過去的生活.
因為這是歷史裡面.
他們所認定的.
他們很穩陣的.
他們不陌生的事.
不過基督耶穌.
他卻要我們抹清眼睛去看清楚.
生命裡面最需要的是什麼.
我們再繼續去思考.
耶穌基督講到.
祂是充充滿滿.
有恩典和真理.
有恩典和真理是耶穌基督祂自己.
是祂自己的生命.
是祂這個福音裡面一個很重要.
上帝宣稱的那種內涵.
當我們接受了這個恩典的時候.
不單只是說一份永生的禮物賜給我們.
在我們自己繼續思想的時候.
這份恩典的觀念.
其實都是塑造我們每一個基督徒的生命.
我有三點.
今天當我們明白.
基督耶穌就是那個恩典與真理的時候.
我們就用這個恩典與真理的福音.
來成為我們生命的養分.
提一提我們今天活著是一個怎樣的人.
第一點.
當我們知道我們領受了這個有恩典與真理的福音.
我們的生活.
我們就不要在我們的人生裡面.
將恩典和真理完全切割.
上帝講的那個恩典.
是一定包含真理的.
上帝所講的那個真理.
一定包含恩典的.
不會兩樣東西分割了.
然後不同時間放給你.
那些是錯誤的神學.
想到上帝這一刻是公義的神.

$^{641}$那一刻是恩典的神.
這樣就很機械化地看上帝對人的心態.
在我們領受的福音是這樣的.
上帝所講的永遠都是有恩典有真理的一個宣講.
在耶穌所講論.
或者見到有些人做錯了事.
耶穌的宣稱永遠都是你的罪已經赦免了.
不過後面交給主角不要再犯.
這個是耶穌裡面不停不停宣講的訊息.
在我們生命裡面.
同樣我們都是要以這個恩典與真理.
成為我們做人處事的那種很重要的價值觀念.
早幾天有機會與一個政府的高層聊天.
了解一下管著這麼多人你怎樣做的呢.
其實他是一直升得很快的.
這個弟兄姊妹.
然後我問他你有什麼令你升得這麼快.
為什麼可以做得這麼好.
管著這麼多人.
很複雜的.
他就告訴我.
其實他作為上司.
他都是將真理和恩典實踐在他的職場生活當中.
總括來說他是一個恩慈的上司.
他對下屬是保住下屬的.
所以令下屬對他很有安全感.
放心地來他工作.
那他怎樣令下屬放心地工作呢.
他說有鑊千萬不要卸給下屬.
有鑊自己揹.
大的鑊他更加要揹.
所以他說有鑊氣的東西他自己會頂的.
下屬就很放心地來他工作.
團隊就會有表現.
我覺得都明白.
講完這件事之後他又加個註腳.
他說這是恩典那一面.
真理那一面是怎樣呢.
他說對於任何下屬要令下屬尊重自己.
他很多事情下屬做錯了.

$^{681}$他一定要清晰地讓下屬清楚自己知道自己做錯了什麼.
他講了一個例子.
有一次有個鑊氣.
他講到下屬揹了大鑊.
口頭上又傻傻的在電話裡面.
又對不起.
然後叫他直接報告.
那幾個上司一齊來道歉.
口頭上報告了就ok.
整件事都講出來了.
這個上司就收貨了.
口頭報告ok了還要寫文字報告.
你剛才所講的東西你寫出來.
我就會幫你想想怎樣砌怎樣做.
然後這個下屬就回去馬上寫報告.
寫完報告之後.
寫的東西和他講的東西完全是兩套東西.
因為寫的東西全部都維護著自己.
然後又講到自己多麼身不由己.
又或者很多環境因素令他犯了這個錯誤.
而這個錯又不是他直接錯.
簡單來說文字報告就是卸鑊給其他人.
當這個上司看到這個報告的時候.
他就召了他進來.
然後在他面前大發脾氣拍桌子來罵他.
拍桌子罵他目的不是要發洩.
是要讓他知道這個過程你要心口如一.
在他面對面講的時候講真相.
寫東西的時候就閃閃避避.
把自己的真相掩飾.
用文字包裝自己是一個受害的無辜者.
這個上司說做錯事我一定幫你頂.
但如果你不真誠地面對自己生命的問題.
這個是品格的問題我是不會幫你頂的.
然後就拍桌子罵他.
然後就拒絕收這個文字報告.
然後我再問你是不是要繼續頂他.
他說不是.
叫他回去寫過的意思是要收一個真實的報告.
目的不是要罵爆他或者拒絕他.

$^{721}$不過他要說作為一個可以令人信服和尊重的上司.
我不可以一方面純粹做好人.
他有鑊我全負責.
而是要讓他們知道.
他們生命的錯誤是什麼.
有真理有恩典.
你才要令到身邊的人來尊重你.
有一個好的平衡.
直到你的生命才真的活現了這種.
信仰給我們的養分.
弟兄姊妹我不知道在你的生命裡面.
你會不會真的有很好的柔合.
你會不會把真理和恩典分工合作.
你做這個好我做這個差的人.
上帝期望的是我們每個人的生命都要活現這個.
有真理有恩典的生命.
很多家庭都說罵兒子就找媽媽罵.
媽媽做醜人.
爸爸就負責給零用錢做好人.
在黑白裡面分工.
這麼蠢的事有些人都會忌諱.
媽媽都不知為什麼忌諱.
永遠喜歡做醜人.
但實際上其實不是這樣來分的.
在我們生命裡面.
我們對人處事.
待人接物.
我們都是有恩典有真理.
整體我們是一個有恩典的人.
不過不要因為講恩典我們就認同所有錯的事.
都姑息了互斷.
來將錯誤的事完全接納.
那個叫做恩典.
聖經裡面不是講一個這樣分割的生命.
而是一個整全的生命.
第二件事.
關於恩典.
我們今天接受這個恩典.
我們除了要平衡之外.
我們今天帶著什麼態度.

$^{761}$來接納上帝這個福音的恩典呢.
有很多朋友或是一些弟兄姊妹.
覺得這個福音一點都不是一個令人快樂的福音.
我曾經問過很多朋友.
我的大學同學.
去鼓勵他信主.
他經常說我心裡是信的.
不過我不會跟你們信的.
這句話有點距離.
他說他心裡是信的.
但他不會跟我們信.
為什麼跟我們信和信有什麼差異呢.
他說跟你們信你們基督徒.
其實是一群沒有趣的人.
他告訴我們.
他說你們人生裡沒有經歷過快樂.
用不著說得這麼誇張吧.
他說是有的.
他說有很多人領受了.
本來他說家人信了主之後.
本來都快樂的.
不過信了一段日子之後.
就經常用聖經來罵家人.
罵過媽媽.
罵家人.
就說她用箴言來說她愚蠢.
又或者用其他的東西來說她迷信.
等等.
他掌握了聖經的知識.
但聖經的知識成為了一個批判世界的工具.
他覺得.
他看到基督徒成為一群沒有趣味的人.
看到一群基督徒.
經常都是.
看到人就是.
肆意的來批評.
這群人不單批評一些迷信者.
也弟兄姊妹之間互相批評.
他說.
看到這個基督徒的生命.

$^{801}$他覺得是一個沒有趣的生命.
是一個不快樂的生命.
當然這個朋友所說的.
都是一些很片面的說話.
未必每個人都是這樣的.
反正不是這樣的.
不過我們值得去留意一下.
當我們領受了這個福音的恩典.
我們會不會很容易成為一個吹毛求疵的人.
我們聖經成為了我們批判世界.
批判別人的一個工具.
而不是幫助自己成為聖傑.
幫助自己更加貼近上帝的一個重要的啟示.
反而我們用聖經的話語.
來批判這個世界之外.
也令到身邊的人有一種窒息的感覺.
這個恩典其實不是叫我們變得憤世嫉俗.
這個恩典其實是要打開我們的心靈.
是要叫我們懂得體會上帝的那種豐盛.
這個恩典是要叫我們的心胸變得更加廣闊.
這個恩典是會帶來喜樂.
遵從上帝旨意的人是快樂.
而不是去強迫身邊的人.
或者用上帝的話語來審判別人.
以自己走高一格來比較自己.
覺得自己是神聖一點.
在這裡獲取快樂.
聖經的恩典是叫我們生命裡.
更加掌握上帝的心腸.
我們領受了這個恩典的時候.
我們快快樂樂地跟別人經歷這個恩典.
而不是去批判別人.
第三方面.
恩典不是一班人獨享的.
一小撮人的權利.
或者當我們領受了這個恩典的時候.
我們更加明白.
本身叫做恩典其實不是我們應得的.
我們應得的罪的功架是死.
如果我們應得的是.

$^{841}$因著罪而滅亡是死亡.
但現在反而因著基督耶穌我們有生命.
我們每一樣東西都不是應得的.
當我們知道我們生命裡沒有一樣東西是應得的時候.
我們都學習.
在我們生命裡.
我們會不會分享多一點恩典給身邊的人.
懂得來分享恩典.
你的生命會貼近上帝多一點.
還記不記得有一個比喻.
耶穌基督說的.
他說那個主人免去了一個僕人一千萬兩銀子的債.
因為他見到他苦苦哀求.
他沒有辦法還的了.
所以這個主人就免去了這個僕人的債.
不過轉過頭來這個僕人原來還有另外一個債的.
欠他十兩銀子.
這個債的父母就苦苦哀求.
用同樣的話來求這個僕人免去他的債.
不過這個僕人卻捉住他的頸就罵他.
欠錢不用還嗎?.
就打他一身.
耶穌用這個比喻說.
這個僕人他不知道他自己是免債的嗎?.
耶穌說你這個奴才.
你要求我我就把所欠的都免你.
你不應當連問你的同伴.
像我連恤你一樣嗎?.
這是《馬太福音》十八章三十二至三十三節.
我們得蒙連恤.
得蒙恩典.
同樣我們都要以恩典來對待身邊的人.
我們生命裡面不單止將福音說出來.
講論的福音是重要的.
不過你生命裡面活出的福音.
在你出現的地方就成為你群體裡面的福音.
那個是真章.
是生命裡面最重要的東西.
如果當我們說決志信主領受了耶穌這個恩典.
但我們的生命是成為別人的和音不是福音.

$^{881}$這個人出現的時候.
一群人都說我有和.
這個人出現的時候很慘.
都是很荒謬.
我們生命我們活出福音的那種內涵.
當我們紀念基督耶穌是一個恩典與真理.
這個神是他自己的時候.
在我們的生命裡面我們學習不要那麼憤憤不平.
學習打開我們的心牆.
當我們覺得很難做的時候.
想一下上帝都是這樣來恩待我.
我們就學習或者逼自己來恩待身邊的人.
這個是違逆了我們自己本身罪惡的本性.
不過當我們信了主的時候.
我們不要被自己罪惡的本性來克制我們.
反而想一下上帝怎樣對待我們.
我們就有多一點的力量.
有些神聖的力量來對待身邊的人.
今天這段經文.
希望解釋的部分.
好像用多一點舊約來講解.
好像是複雜一點.
不過其實聖經裡面一脈相承.
由舊約到新約.
都是指著基督耶穌的生命來講的.
律法是一個上帝給的恩典.
不過終極的是要藉著基督耶穌.
將恩典與真理呈現在這個世界當中.
我們今天每一個領受了這個恩典與真理.
領受了基督耶穌他自己的救恩的時候.
我們生命都用這份恩典與真理.
來塑造我們的生命.
讓我們的生命.
人們看到我們的時候.
都要將榮耀歸給神.
更加重要的是.
當我們能夠活出恩典與真理的生命.
我們會體會到人們對信仰的那份尊重.
也都會珍惜在生命裡面.
能夠認識了我這個人.

$^{921}$求主繼續幫助我們.
帶領我們走面前的路.
\newpage



\section{哥林多後書 12:19-13:4}
\label{sec:Sd44cjxmurg}
\textbf{2021年12月26日崇拜直播｜梁家麟牧師｜最後的提醒｜哥林多後書12\_19-13\_4}
\newline
\newline
連結: \href{https://youtube.com/watch?v=Sd44cjxmurg}{\texttt{https://youtube.com/watch?v=Sd44cjxmurg}} ~~~~ 語音日期: 2021-12-29
\newline
\newline
\hyperref[sec:Msw_I2sXTQc]{\small{< < < PREV SERMON < < <}}
~
\hyperref[sec:index_chronic]{\small{[返順時目]}}
~
\hyperref[sec:oDel5UYyQt0]{\small{> > > NEXT SERMON > > >}}
\newline
\newline
哥林多後書 12:19-13:4
\newline
\begin{longtable}{cl}
\hline
\hline
章節 & 經文 (和合本修訂版)\\
\hline
12:19 & \begin{tabularx}{0.7\textwidth}{X} 你們一直認為我們是在你們面前為自己辯護嗎？其實，我們本是在基督裡當著神面前說話。親愛的，一切的事都是為了造就你們。 \end{tabularx} \\ \\ \relax
12:20 & \begin{tabularx}{0.7\textwidth}{X} 我怕我再來的時候，見你們不合我所期望的，而你們見我也不合你們所期望的。我怕有紛爭、嫉妒、憤怒、自私、毀謗、讒言、狂傲、動亂的事。 \end{tabularx} \\ \\ \relax
12:21 & \begin{tabularx}{0.7\textwidth}{X} 我怕我再來的時候，我的神使我在你們面前蒙羞，並且又因許多人從前犯罪，行污穢、淫亂、放蕩的事，不肯悔改而悲傷。 \end{tabularx} \\ \\ \relax
13:1 & \begin{tabularx}{0.7\textwidth}{X} 這是我第三次要到你們那裡去。「任何指控都要憑兩個或三個證人的口述才能成立。」 \end{tabularx} \\ \\ \relax
13:2 & \begin{tabularx}{0.7\textwidth}{X} 對那些犯了罪的人和其餘所有的人，正如我第二次見你們的時候曾說過，現在不在你們那裡再次說：「我若再來，必不寬容。」 \end{tabularx} \\ \\ \relax
13:3 & \begin{tabularx}{0.7\textwidth}{X} 因為你們想求證基督是否藉著我說話。基督對你們並不是軟弱的，而是在你們裡面大有能力的。 \end{tabularx} \\ \\ \relax
13:4 & \begin{tabularx}{0.7\textwidth}{X} 他因軟弱被釘在十字架上，卻因神的大能仍然活著。我們在他裡面也成為軟弱的，但對你們，我們將因神的大能而與他一同活著。 \end{tabularx} \\ \\
[1ex]
\hline
\hline
\end{longtable}
$^{1}$今日經訓記載於《哥林多後書》十二章十九至十三章四節.
在您手上的《聖經》新約二百五十五至二百五十六頁.
《哥林多後書》十二章十九節.
你們到如今還想我們是向你們分數.
我們本是在基督裡神面前說話.
親愛的弟兄啊!一切的事都是為你為造就你們.
我相信我來的時候是你們不合我所想望的.
你們見我也不合你們所想望的.
又怕有紛爭,嫉妒,擾亂,結黨,毀謗,殘嚴,狂傲混亂的事.
且怕我來的時候我的神叫我在你們面前談鬼.
又因許多人從前犯罪,行污穢,奸淫,邪道的事不肯悔改.
我就憂愁.
十三章一節.
這是我第三次要到你們那裡去.
憑兩三個人的口作見證.
句句都要定準.
我從前說過.
如今不在你們那裡遊說.
正如我第二次見你們的時候所說的一樣.
就是對那犯了罪的和其餘的人說.
我若在內必不寬容.
你們既然尋求基督在我裡面說話的憑據.
我必不寬容.
因為基督在你們身上不是軟弱的.
在你們裡面是有大能的.
他因軟弱被釘在十字架上.
卻因神的大能仍然活著.
我們也是這樣同他軟弱.
但因神向你們所顯的大能也必與他同活.
(十三章一節).
各位姐妹連續兩個星期都是我.
請大家原諒.
繼續說聖誕快樂.
雖然今天是boxing day.
但對我們來說.
其實每天都是慶祝耶穌基督誕生的日子.
上星期的講到哪裡.
都收到一些弟兄姊妹的反應和回響問題.
今天先做一點Q and A.
兩個很有意思的問題.

$^{41}$值得簡單認一認.
第一個問題.
有弟兄問我.
上次我說到宗教或者基督教.
特別是聖誕節.
在非宗教世界就被商業化.
在基督教世界就被宗教化.
他問不知道宗教化是什麼意思.
宗教化是什麼問題.
宗教化簡單說.
有兩個很重要的特徵.
第一個就是將它禮儀化.
這是很簡單的.
第二就是很容易將它抽象化符號化.
舉一個簡單的例子.
今天我們很多基督徒喜歡.
配戴十字架的飾物.
十字架的鑽石.
十字架的項鍊.
對我們來說.
有七萬多個十字架的款式.
做得很可愛.
不過很難想像.
在羅馬帝國時代.
耶穌基督時代.
十字架會給人一個很可愛的感覺.
不可能的.
十字架是一個羞辱痛苦可怕的記號.
我們將它藝術化.
將它變得很可愛.
我不反對.
我不是反對藝術的人.
我們要直面真相.
如果十字架是一個很可愛的事.
耶穌基督釘十字架是一個很好玩的事.
我們的救恩就真的用來搞了.
一定不是.
十字架是一個痛苦的記號.
正因為十字架是一個極其羞辱痛苦的記號.
我們耶穌基督被掛在木頭上.

$^{81}$才有真正的意義.
所以我們要小心.
有時候我們會將一些東西.
宗教化了.
將它藝術化了.
將它美化了.
我們就看不到背後.
真正的歷史的意義.
屬靈的意義.
這是第一點要講的.
第二個有人問我的問題.
他問如何謙卑一點.
就能夠看到耶穌基督是我們的答案.
我們的謙卑和看到耶穌基督是答案有什麼關係.
是的.
耶穌的誕生.
是否好像真正有關於這個世界.
在我日常生活中.
我眾多的問題.
耶穌基督好像都不是一個很有效的答案.
這是我們常常面對的困擾.
信仰不再有效.
信仰和我們的生活無關.
這是我們常常面對的困難.
面對的問題.
是真的.
如果你問現在股市大跌.
是否要撈底的時候.
你想打開聖經來找答案.
那你省一點.
沒有這些答案.
什麼時候入市.
現在是否適宜買樓.
我不知道.
你問珍珍牧師.
省一點.
她很差的.
她不會告訴你.
聖經也說不出什麼準確的答案.
我們日常生活太多的問題.

$^{121}$你發覺聖經是沒有答案的.
不過.
我們日常遇到的是.
當我們大家禮時.
睡在病床.
特別是在死亡床的時候.
你會發覺.
原來你平時在江湖行走.
你關心的那些問題.
又好像沒有什麼意義.
而信仰突然變得很有意義.
信仰突然變得很有意義.
耶穌基督為我死.
使我有復活的生命.
對我突然間就有意義.
所以關鍵的問題在於.
我們平時被很多花花碌碌的問題.
佔據著我們的心.
那些問題是不是真的.
對我們來說最關重要的問題.
我們是不是真的只問這些問題的人.
只不過有些問題我們覺得.
too remote(太遙遠).
太遙遠.
對我們好像不是即時有效.
不是即時有需要.
很多時候我們真的.
回到我們人生最基本的問題.
我是誰.
我有什麼追求.
有什麼需要.
然後你就發覺.
耶穌基督是不是答案.
非常是答案.
耶穌基督準確地去回答.
我生命裡面我最需要的問題.
只不過有時候我覺得自己.
很有東西.
裝模作樣.
我覺得我自己.

$^{161}$不會問這麼低層次的問題.
我問一些高端的問題.
然後你會發覺.
耶穌太低端了.
耶穌回答不了我們.
所以弟兄姊妹.
將自己還原.
直面耶穌.
然後你會覺得.
信仰真是我們的準確答案.
道路.
真理.
生命.
不是假的.
好.
回到今天要說的.
我們又回歸哥林多後書.
還有兩次.
我們不過要跨年.
跨完一年.
明年第一次都會繼續說.
通常臨尾的經文都是瑣碎碎的.
都是瑣碎碎碎的.
我們很快.
因為太多瑣碎的東西.
所以就像上主一學一樣.
有很多筆記在後面.
我不會全部說那些筆記的東西.
如果大家有興趣.
想看得仔細一點.
你自己上網翻看.
訂一個定格.
你自己再看那些powerpoint.
我不會讀那些powerpoint.
每句都讀.
保羅在結束他要為自己做的一大段自辯之後.
他為自己的自辯做了一句總結.
十二章十九節.
你們如今幻想我們是向你們分數.
我們本是在基督裡.

$^{201}$當上帝面前說話.
你們如今幻想.
這是一個完成式.
意思就是.
你們是否一直以為.
這是那麼多信徒對保羅的一個成見.
也是造成他們很多問題的關鍵.
在基督裡.
當上帝面前說話.
在臨後二章十七節都有同樣的說法.
在基督裡.
和耶穌基督聯合是他事奉的來源.
能力的來源.
當上帝面前.
保羅純粹說.
我所做的一切.
都是上帝在上.
天眼悄悄.
他可以監察知道我是怎樣的人.
整句話.
你們如今幻想我們是向你們分數.
是什麼意思呢.
原來是一個問題.
你們是否一直以為.
我是羅列證據來和你們申辯.
希望證明你們過去所說的是錯的.
對我的指控是錯的.
猜測是錯的.
並且.
我是想證明.
你錯我對呢.
保羅說不是.
我們只是在上帝面前.
想說事實真相.
上帝監察一切.
會知道我們每句都是實話.
並無虛言.
人間的爭論.
總是想分出是非對錯.
證明誰對誰錯.

$^{241}$證明自己有理.
《老九福音》十章二十九節.
辯論的結果是一方贏.
另一方輸.
但上帝其實不是想.
知道誰贏誰輸.
上帝關心的是事實真相.
所以我經常責罵我的兒子.
將來在天上.
沒有生意沒有工作.
律師是沒有飯吃的.
律師就是找論點.
為自己的家人.
爭取最大的好處.
事實不是最大的問題.
他們喜歡討論論點的問題.
但全知的上帝.
論點是沒有意思的.
你說來說去.
有沒有這樣做過.
有沒有發生過.
有沒有偷過東西.
有沒有殺過人.
這些是他最關心的問題.
說多無謂.
辯論是沒有意思的.
上帝只關心事實.
上帝不關心論據.
不關心論據.
上帝沒有疑點.
本來我不是想爭取贏你們.
我不是想證明你們錯.
我對.
我只是想在上帝面前.
坦誠地說話.
說真實的事.
這是最關鍵的.
很多時候我們.
以為我們在跟別人辯論.
我們很容易會立即.

$^{281}$觸起各種防禦攻勢.
Defensive mentality.
立即起了槓.
要證明我贏.
你錯.
證明我對.
所以要努力去找對方的.
argument(爭論)的漏洞.
哪怕那些漏洞只是一些枝節的東西.
總之就砌死對方.
轉移話題就夠了.
不要讓他在我的弱點糾纏.
將問題轉移.
很多時候.
我們就算是教會.
我們開會.
為辯論而辯論.
為討論而討論.
有時候討論完之後.
都不知道到底重點在哪裡.
假如我們不是有辯論的心態.
不是你贏我輸.
我贏你輸.
這樣的心態.
我純粹只是想就事論事.
想知道這件事的真相.
了解到底甚麼是對的.
應該做甚麼是不應該做.
你會發覺問題簡單很多.
本來我完全沒有意思.
想在你面前去做辯論.
我沒有的.
我不是想證明自己有智慧有能力.
我只不過想澄清誤會.
重建關係.
本來我真誠表達.
我對信徒的感情.
很想馬上去到他們那裡.
因為他思念他們.
但又怕去到他們那裡之後.

$^{321}$看到一些不理想的情況.
讓他傷心難過.
相見爭如不見.
又想去又怕去.
十二章二十到二十一節.
我怕我再來的時候.
見你們不合我所想望的.
你們見我又不合你們所想望的.
又怕又紛爭.
執導,嬲嬲,結黨,毀謗.
殘嚴,狂傲,混亂的事.
即是我再來的時候.
我上帝叫我在你們面前慚愧.
又因許多人從前犯罪.
行污衊,奸淫,邪辦的事.
我便會憂愁.
本來他怕三件事.
第一.
他怕他和那麼多信徒仍然有嫌隙.
互相不接納對方.
你看我不順眼.
我也看你不順眼.
總是在對方的一言一行.
你只要對方有成見.
你當然看到他有毛病.
就算沒有毛病.
你也聽到毛病.
雞蛋裡挑骨頭.
這樣的討論永遠只是吵架收場.
這個情況大概在保羅第二次訪問哥林多教會的時候就出現了.
很多時候我們經常說要溝通.
溝通確實是解決分歧一個重要的出路.
不過溝通的前提是我們大家能夠心平氣和.
坐下來真的傾聽.
並且不僅僅是說.
是聽對方說他的觀點.
聽對方的論點.
盡量不是用駁斥.
而是用一個同情,諒解的角度.
試試設身處地.

$^{361}$明白對方到底在想什麼.
說什麼.
然後才有機會拉近雙方的距離.
你可以願意修訂自己原來的觀點,看法.
走近對方的一步.
這是很重要的.
如果我們只是各自表述.
互相表達自己的立場.
寸土不讓.
根本不準備妥協.
不準備任何修訂的話.
坦白說.
怎樣溝通都沒有用.
溝通是多餘的.
dialogue是monologue(單語).
基本上只有一個字說.
各自表述有什麼意思.
這是第一個問題.
他擔心.
第二個他擔心的是.
那麼多教會仍然有各種內部的爭端.
整個教會充滿負面的情緒,言行.
確實他們的爭端可能是因為保羅而起.
有擁保羅派,有反保羅派.
兩派互相對立是真的.
但爭端也可以跟保羅無關.
因為很可能的就是.
教會裡有部分信徒常常說負面的東西.
凡事懷疑,凡事不信,凡事唱反調.
如果一旦有這樣的情況.
教會要多爛就可以有多爛.
你不要告訴我.
教會不會這樣.
有永遠的反對派.
永遠唱反調嗎.
你覺得不可能發生.
你覺得這件事不可思議.
因為我們教會不是這樣.
不過我要說一個月前.
有個傳道人跟我分享他的事況難處.

$^{401}$他的教會就剛剛好在這樣的情況裡.
沒得搞的.
兩派互相對立.
總之凡事敵人贊成的.
我們必要反對.
你告訴我這樣的教會怎麼搞下去.
沒得搞的.
第三個保羅害怕的是.
過不多信徒中間仍然有犯罪的人.
他們會因為保羅的來到而感到羞愧.
不過覺得羞愧也不是最可怕的.
最可怕的是.
他們有人可能會不覺得羞愧.
這樣就很大件事了.
我離開我原來那個.
事奉幾十年的中派.
其中一個原因就是.
我不想處理太多教會的紀律問題.
我是很怕.
真的怕得要死.
處理教會的紀律問題.
特別是那些大祭那些.
牧師犯罪.
或者教會要開革一班人.
要我去處理這些事.
我是焦頭爛額的.
很痛苦.
有人拒絕承認犯罪.
不斷為自己的錯誤行為強辯.
甚至質疑教會.
和聖經的權威和有效性.
這當然帶來很多的衝突.
就算坦白認罪.
願意接受處分.
願意接受挽回的安排.
他們的跌倒仍然會傷害很多人的心.
甚至叛倒很多人的信仰.
對犯事者的家人.
特別是太太,父母,子女.
傷害是非常嚴重的.

$^{441}$你想像一下.
假如我范姦蔭.
在教會公開處理.
我太太,我子女如何自處.
很艱難.
犯事者越在教會的高位.
從事教導領導工作.
他們的失德見證.
會動搖很多弟兄姊妹的心.
這不用說了.
保羅說了他這三個大的懼怕.
我們說愛侶沒有懼怕.
《約翰一書》四章十一節.
不過這只要對上帝.
對未來的信心來說.
我們對未來沒有懼怕.
對上帝沒有懼怕.
但在人際關係上.
你說沒有懼怕是假的.
總是帶著很多很多的懼怕.
是的.
一切在上帝的手裡.
但我不知道.
如何在上帝的手裡.
對不對?.
你對你的子女有沒有擔心?.
非常擔心.
非常非常多.
怎會沒有擔心?.
特別是做事我不怕.
做事我可以容易控制.
做人.
人從來都不是我控制的.
你帶著一群年輕人.
你可以控制他們嗎?.
限制一點吧.
除了不斷的惡惡惡惡.
惡到你自己都覺得煩.
不要說對方覺得煩.
之外.

$^{481}$真的為他祈禱.
沒有其他事可以做.
上帝照著.
我們人.
你說沒有懼怕.
沒有憂慮.
不可能.
你會看到.
保羅常常在他的書信裡.
表達他對.
宗徒和宗教會的掛念和憂慮.
在前面的《哥林多後書》.
十一章28-29節.
他講.
他的苦難清單.
除了外面的事.
還有為宗教會掛心的事.
天天壓在他身上.
保羅常常為教會.
有很多很多的焦慮.
在《腓立比書》.
保羅因著.
和公爾巴夫提的患病.
也表達極度的憂慮.
他講憂常加憂.
很擔心.
可能你會說.
做仕途.
做傳道人.
不是應該要信心爆棚嗎?.
為什麼這麼多擔心?.
為什麼這麼多憂慮?.
你不是應該對上帝有信心?.
對未來有信心?.
對人有信心嗎?.
不是的.
憂慮就是憂慮.
當然.
你還可以說.
就算你憂慮.

$^{521}$也不要說出來.
說出來那麼容易受傷.
說出來就使自己.
那個屬靈地位.
屬靈形象受到損害.
你知道最好的領導人.
永遠處變不驚.
對不對?.
人家跟你說那個噩耗.
其他人都嘴巴張開.
你都好像沒聽到.
若無其事.
這樣才夠屬靈.
這樣才夠招凡脫俗.
保羅不是.
你會看到很有趣.
保羅從來都不介意說他自己.
心裡各種憂慮.
各種難過.
各種憤怒.
甚至他那些恩怨情仇.
好像失寵的怨婦.
他不介意說這些.
他不覺得噁心.
因為保羅.
坦白說.
從來都不想在我們面前.
顯示他是一個強者.
他從來不介意.
他是一個軟弱者.
他喜歡將他軟弱的一面.
跟別人分享.
你知道.
如果你看球賽就知道我說什麼.
一個身價一億英鎊的守門員.
如果擺烏龍.
一定會被人插到飛起.
但如果一個.
做了很多年後備.
臨時替補出場的守門員擺烏龍.

$^{561}$觀眾的反應就包容很多.
說長說短.
整個身體都僵硬了.
沒辦法了.
撲不到有什麼奇怪的.
同樣的.
人們會質疑.
為什麼一個強者.
例如陳牧師.
會犯某些低級錯誤.
做出這麼蠢的行為.
弟兄姊妹.
我在這裡鄭重聲明.
我是如假包換的弱者.
我絕對不是強者.
我絕對不是能者.
所以如果我犯錯.
如果我有時候.
不小心露出我的.
失望,沮喪,意志消沉.
請你不要詫異.
也不要有太強烈的反應.
我不犯錯.
你才詫異.
保祿從來都不需要.
以一個全能者的面貌顯現.
是的.
我在你面前.
我從來不扮強者.
弱者就是弱者.
我有很多擔心.
很多掛念.
很多焦慮.
有時候會比較消極.
請你包容.
我是一個很差的人.
我們教會只有一個英雄.
一個強者.
他叫耶穌.
我們沒有人需要以強者自居.

$^{601}$我們再看這段經文.
保祿提到.
有八宗會造成教會危機分裂的罪.
紛爭,嫉妒,嬲怒,結黨,毀謗,.
殘嚴,狂傲,混亂.
我不詳細解釋這八件事.
你們自己看筆記.
我有做過研究.
因為沒有解經.
所以就跳了它.
總之這八件事.
都是關涉到個人的性格,行為,表現.
特別是人際關係的處刑.
如果信徒之間沒有彼此憎恨的心.
即是沒有嫉妒.
如果真的不輕易發怒.
沒有嬲怒.
如果不自誇張狂,沒有狂傲.
不隨意攻擊,殘嚴,散播謠言,毀謗.
不搞對立,紛爭,分裂,結黨.
不破壞秩序,混亂.
教會就可以維持聖靈所賜的合一.
反之有這八樣徵狀表現的教會.
一定戰火連天.
最終造成教會分裂.
保羅說他很擔心.
他在哥林多教會見到這八件事.
我們大家都擔心.
在我們皇鎮見到這八件事.
很擔心在我們的信徒群體中間.
見到這八種行為.
OK.
除了這八種會造成分裂的罪外.
保羅說他也不想見到教會有道德上的犯罪.
例如污衊,奸淫,邪辦.
三種行為不需要分開處理.
因為他是用同一個定貫詞.
包含這三種行為.
所以這三種行為應該是同義詞.
即是說是性道德的犯罪.

$^{641}$哥林多信徒曾經出現過嚴重的性倫理問題.
你看哥林多前書.
一節就說到.
在他們中間有人取自己的繼母.
有這樣的亂龍行為出現.
所以哥林多教會曾經有過這些性倫理問題.
當然我們可以問.
這些問題是否仍然存在呢?.
是否沒有處理過呢?.
是否教會仍然有很多分裂的敗壞.
道德的敗壞.
是否有這麼多呢?.
可能有.
我們不知道.
不過我簡單說.
如果我們看《迦太書》.
五章十九到二十一節.
那裡同樣有列舉一個教會的罪惡清單.
你會發覺教會的罪惡清單.
跟哥林多後書所說的很接近.
所以我們可以這樣推斷.
這是保羅提到教會的罪惡的時候.
慣常想到的罪狀.
所以不一定是真的哥林多教會出現這麼多問題.
保羅一提到教會犯罪的時候.
例如他不小心提到旺朕的罪的時候.
他就會提到.
你們不要出現八宗分裂的罪.
加三宗懲獄的罪.
是否真的我們有呢?.
不一定.
保羅只是列舉而已.
哥林多教會雖然問題很多.
但未必爛到全面敗壞的地步.
我們可以有一點信心.
教會如果敗壞到完全沒救.
保羅寫信給他們都浪費時間.
省口水就好了.
我們又教主一學.
保羅計劃第三次訪問哥林多.

$^{681}$在第一次保羅去哥林多的行程.
他建立了哥林多教會.
第二次行程是做了一個簡短的訪問.
在那裡他和哥林多教會有極大的衝突.
讓各人都感到痛苦.
為了幫助大家去了解.
保羅和哥林多教會的複雜關係.
也明白之前兩次行程的來龍去脈.
我將Rodmartin整理出來的大事記.
跟大家讀一次.
保羅和哥林多教會的關係.
首先他建立了哥林多教會.
你看《少行傳》18章一節開始.
你會看到.
然後保羅離開哥林多教會去伊弗所.
接著保羅寫了第一封信給哥林多教會.
這封信其實已經失去了.
第一封信不是哥林多前書.
是之前還有第一封信.
因為在哥林多前書五章九節有提到.
他寫了之前一封信.
所以哥林多前書是第二封信.
總之哥林多前書之前有一封信.
有人相信臨後六章十四節到七章一節.
是這封已經遺失的信一部分.
這個大家不用理會.
總之保羅寫了第一封信.
然後有弟兄姊妹.
格萊士家裡的人寫信給保羅.
告訴保羅哥林多教會現在鬧分裂.
差不多同一個時期.
保羅收到哥林多教會正式的信.
求問如何處理一些教會重大課題.
崇拜的問題.
和教外人關係的問題.
尋求他的指引.
所以保羅回覆哥林多教會.
應答各種問題.
這就是我們今天的哥林多前書.
這封信是由提多親自送到哥林多.

$^{721}$所以哥林多前書是第二封.
之前有第一封.
現在是第二封.
提摩泰奉派去哥林多.
處理保羅教徒的一些問題.
這在哥林多前書四章十七節有提到.
與此同時.
哥林多教會再爆發嚴重的人事衝突和分裂.
領頭人攻擊保羅.
提摩泰沒有能力招架和處理.
被迫離去.
回到伊弗所.
收到提摩泰的報告之後.
保羅就親自去哥林多.
這就是他第二次去哥林多的行程.
要親自去處理那些反對他的領頭人.
不過保羅在教會裡就被羞辱.
處理不了.
那些人抵制他.
最後他帶著極大的創傷離去.
回到伊弗所.
所以這被稱為憂愁之旅.
或者痛苦之旅.
保羅在第二次行程失敗之後.
他在伊弗所寫了一封極其嚴厲的信.
壓響他的仕途權柄和一切威信.
勒令哥林多教會必須要隔退那個領頭人.
這封信被稱為流淚的信.
或者是嚴厲的信.
Severe Letter.
這封信也遺失了.
或者有人說這是臨後的十到十三章.
提託負責送這封信過去.
之後他去德羅亞和保羅會合.
所以嚴厲的信是第三書.
是第三信.
保羅本來計劃在德羅亞見提託.
但是遇上不到.
結果他去了馬其頓.
然後才找到提託.

$^{761}$提託見到保羅之後.
告訴保羅.
哥林多教會最惡劣的情況已經過去.
信徒表達願意順服保羅的帶領.
這是在哥林多後書七章六到十六節.
我們看到的.
不過很多學者就這樣估計.
提託是誤判的處境.
其實反對保羅的人沒有消聲匿跡.
仍然存在.
提託只看到樂觀那一面.
沒有看到消極那一面.
因為聽到提託的喜訊.
保羅就寫信給哥林多教會.
這就是今天的哥林多後書.
有人說這只是寫了一到九章.
一方面表達感恩.
另一方面希望可以掃除教會殘餘的反對勢力.
去做善後.
這封信是在馬其頓寫的.
由提託和兩位弟兄去送遞.
所以哥林多後書是第四書.
是不是很複雜?.
是第四書.
在送出哥林多後書沒多久.
保羅聽到哥林多教會又出現新的麻煩.
原來反對他的信徒連同一些假師傅.
挑戰保羅的權威.
提出和保羅不一樣的教導.
保羅擔心教會的問題沒有徹底處理.
所以再寫另一封信.
很多人猜測這是臨後十到十三章.
如果你真的相信.
哥林多後書是兩封信的聯合.
一到九是一封.
十到十三是第二封.
十到十三就是第五書.
保羅沒多久再訪問哥林多.
這是他第三次的行程.
這是《少行傳》二十章第二節有記載.

$^{801}$他留了一段比較長時間.
差不多一年.
然後他在哥林多教會寫羅馬書.
是不是很複雜?.
整個大事表就是這樣.
保羅前後去了三次哥林多.
寫了五封信.
如果你不相信哥林多後書是兩封信.
寫了四封信.
四封或者五封信給哥林多.
就是這樣.
有這個表就容易跟得上.
到底發生什麼事.
保羅說.
他會準備第三次去他們中間.
這是我第三次要到你們那裡去.
憑兩三個人的口作見證.
句句都要定準.
我從前說過.
如今不在你們那裡又說.
正如我第二次見你們面的時候.
再說一樣.
對那些犯了罪的或其餘的人說.
我若再來必不寬容.
你們既然尋求基督在我裡面說話的憑據.
我必不寬容.
因為基督在你們身上不是軟弱的.
在你們裡面是有大能的.
保羅宣告他第三次要去哥林多.
行程已經確定.
並且他用新聖記十九章十五節的說法.
我這個是一定兌現的承諾.
有人做見證的.
句句可以定準.
為什麼保羅要誓神疲軟說他一定會去呢?.
因為他曾經因著某些原因.
答應去又沒有去到.
所以他被人以為他是一個沒有口齒的人.
又或者被人以為他是一個沒有擔待的人.
不想面對衝突.

$^{841}$所以說去又不敢去.
保羅在這裡很硬朗地說.
我一定去.
我說了去我就一定去.
至於他去這次行程的目的.
除了探望闊別多時的弟兄姊妹之外.
也要處理教會的根本性問題.
賈師傅的問題.
各種敵對保羅的弟兄姊妹的問題.
希望教會不再長期困擾在人事衝突.
或者義肩糾紛的裡面.
保羅前面表達過他.
擔心他再來這次行程會失敗.
而且怕我來的時候.
我上帝叫我在你們面前慚愧.
十二章二十一節.
什麼叫慚愧呢?.
這方面指著的是.
保羅擔心弟弟妹羞辱他.
但這也不是他最擔心的事.
他最擔心的是他怕在上帝面前慚愧.
意思是什麼?.
他怕他此行又再失敗.
有負上帝所託.
所以在上帝面前慚愧.
比在弟兄姊妹面前蒙羞更加嚴重.
保羅很擔心這次行程不成功.
所以他很強硬地說.
我要來了.
我很快來.
我一定來.
你們不要心存僥倖.
我來到就是做紀律處分的時候.
我若再來.
必不寬容.
我必不寬容.
他說兩次.
不寬容.
意思是他不會再給機會.
至於不寬容是做什麼呢?.

$^{881}$保羅對哥林多教會的犯罪者.
能夠做什麼懲處呢?.
最大的可能是將他們趕出教會.
哥林多前書五章二節有提到這一點.
趕出教會在我們今天是很平常的事.
只不過沒有旺朕的會籍.
香港一千三百多間教會.
此處不留人自有留人處.
不去這間就去另一間.
同一條街也有幾間.
我們斜對面有一間崇恩浸信會.
浸信會也有很多間.
很多間也有得去.
但是在保羅時代.
離開被趕出教會.
並不是這麼簡單的事.
他不僅僅在人間組織意義上失去了教會的會籍.
保羅的眼中.
更加在肅明上宣告被懲處的人.
不再享有舊恩.
所以保羅的說法不是趕出教會剝奪會籍.
而是像馬會將一些人剝奪了會籍.
真是好笑.
而是將他們交給撒旦.
保羅用這個字眼.
將他們交給撒旦.
即是說他們不再屬基督.
這可以說是一個社會放逐.
在保羅時代.
人的社交網絡其實很小.
並且信仰很少是個人的事.
信仰牽涉到人生很多層面.
通常是整個家族信耶穌.
當信主耶穌.
一家都必得救.
所以如果一個人被開除出教.
某程度上意味著.
他要切斷和他的家人.
他所屬的社群的聯繫.
所以這個問題是非常非常嚴峻.

$^{921}$就像傳統的農村社會.
將一個人趕出教.
其實是將他隔離他整個家族.
整個社區.
所以例如在印度.
一些比較傳統宗教的地方.
你要他改信基督教.
對他影響很大.
他真的可以被整個家族趕出.
這個影響是很深遠的.
和香港真的不同.
因為宗教在我們生活中佔的份量很少.
你不覺得嗎?.
其實很少.
不回這間教會就回另一間教會.
你都不用被人趕出教.
你都常常轉換教會.
三朝兩日換一間.
有什麼所謂呢?.
並且有些教會是擺明車馬.
專收人家遺棄的人.
有間教會.
你在另一間教會犯了奸淫.
和太太離婚.
有婚外情另娶.
你可以去那間教會做傳道人.
他講得很明白.
我們教會是講恩典的.
我們講恩典的.
沒有問題的.
所以你不需要停聖職.
這間教會不收我.
我就去另一間.
很容易的事.
所以今天開除出會.
變成一個都幾沒有意義的行為.
坦白講.
在我都未開除他出會之前.
他就已經出會了.
那你怎樣?追殺他?.

$^{961}$你有什麼可以做?.
教會沒有什麼可以做的.
你千萬不要出黃榜.
公告天下.
告到你甩褲子.
你不可以的.
你不可以將一間教會出的問題.
公諸天下.
告訴另一間教會都會被告.
會告死你的.
所以這是無可奈何的.
今天是無可奈何的事.
我經常舉這個例子.
我明知我女朋友要撇我.
於是我立即搶襲.
在她未開口講之前.
我撇你.
不是你撇我.
是我撇你.
保持自己的尊嚴.
所以今天.
我們常常都沒有機會.
去開除人家出教會.
人家先自我開除了.
無緣無故.
保羅出了黃牌.
告訴頂智媒.
他再來的時候.
他將會嚴厲去懲處.
那些繼續犯罪的人.
不過.
當保羅講這樣的話.
我們可以知道一件事.
他不是真的.
很想將這些人開除出教.
甚至保羅一直都說.
他不喜歡教會扮演一個審判者的角色.
教會不是法庭.
審判定罪不是教會的主要工作.
我們要傳講的是耶穌基督.

$^{1001}$對人犯罪的赦免的訊息.
這才是福音.
審判定罪.
完全的定罪.
真正的定罪.
永遠留在耶穌基督再回來之後才做.
耶穌講過.
他再回來的日子叫審判的日子.
保羅都宣告過.
現在是月立的時候.
現在是拯救的時候.
教會不是以懲罰.
作為主要的功能.
傳講懲罰的訊息.
也不是教會的主調.
我們不會說多一個少一個無所謂.
我們卻是認定一個也不能少.
無論如何.
保羅對那麼多教會.
信徒採取任何形式的行動.
他的出發點.
他的動機都是出於愛.
為了造就他們.
為了希望他們能夠悔改.
傳道者可以為信徒的警方憂傷.
可以有挫敗,沮喪的感覺.
不過一定不可以懷恨.
就算面對反對者.
我們都不可以懷恨.
所以保羅在這裡表白.
他不再寬容.
要懲治那些犯罪拒絕悔改的人.
不是一個以恨還恨的報復行動.
這是一個救人不是害人的行動.
就算真的沒辦法救那個被懲治的人.
他都希望可以防範更多人.
因為相同的罪而沉淪.
總結簡單的一句話.
保羅宣告他即將不寬容.
說這樣的話.

$^{1041}$就是表明他現在仍然是一個寬容的時候.
教會現在是寬容的時候.
好像你媽媽呢.
我打的了.
一二二個半二又四分三.
諸如此類.
現在仍然是寬容的時候.
好我們就停在這裡.
十三章一到四節還沒講完.
我們下一次就會講完整篇《戈林德學書》.
\newpage



\section{馬太福音 28:8-20}
\label{sec:oDel5UYyQt0}
\textbf{2022年1月2日崇拜直播｜陳明泉牧師｜我們的6G - 門訓不斷｜馬太福音28\_8-20}
\newline
\newline
連結: \href{https://youtube.com/watch?v=oDel5UYyQt0}{\texttt{https://youtube.com/watch?v=oDel5UYyQt0}} ~~~~ 語音日期: 2022-01-03
\newline
\newline
\hyperref[sec:Sd44cjxmurg]{\small{< < < PREV SERMON < < <}}
~
\hyperref[sec:index_chronic]{\small{[返順時目]}}
~
\hyperref[sec:Qke9jW9Cx6E]{\small{> > > NEXT SERMON > > >}}
\newline
\newline
馬太福音 28:8-20
\newline
\begin{longtable}{cl}
\hline
\hline
章節 & 經文 (和合本修訂版)\\
\hline
28:8 & \begin{tabularx}{0.7\textwidth}{X} 婦女們急忙離開墳墓，又害怕，又大為歡喜，跑去告訴他的門徒。 \end{tabularx} \\ \\ \relax
28:9 & \begin{tabularx}{0.7\textwidth}{X} 忽然，耶穌迎上她們，說：「平安！」她們就上前抱住他的腳拜他。 \end{tabularx} \\ \\ \relax
28:10 & \begin{tabularx}{0.7\textwidth}{X} 耶穌對她們說：「不要害怕！你們去告訴我的弟兄，叫他們往加利利去，在那裡會見到我。」 \end{tabularx} \\ \\ \relax
28:11 & \begin{tabularx}{0.7\textwidth}{X} 她們去的時候，看守的兵有幾個進城去，把所發生的事都報告祭司長。 \end{tabularx} \\ \\ \relax
28:12 & \begin{tabularx}{0.7\textwidth}{X} 祭司長和長老聚集商議，就拿許多銀錢給士兵， \end{tabularx} \\ \\ \relax
28:13 & \begin{tabularx}{0.7\textwidth}{X} 說：「你們要這樣說：『夜間我們睡覺的時候，他的門徒來把他偷去了。』 \end{tabularx} \\ \\ \relax
28:14 & \begin{tabularx}{0.7\textwidth}{X} 若是這話被總督聽見，有我們勸他，保你們無事。」 \end{tabularx} \\ \\ \relax
28:15 & \begin{tabularx}{0.7\textwidth}{X} 士兵收了銀錢，就照所囑咐他們的去做。這話就在猶太人中間流傳，直到今日。 \end{tabularx} \\ \\ \relax
28:16 & \begin{tabularx}{0.7\textwidth}{X} 十一個門徒往加利利去，到了耶穌指定他們去的山上。 \end{tabularx} \\ \\ \relax
28:17 & \begin{tabularx}{0.7\textwidth}{X} 他們見了耶穌就拜他，然而還有人疑惑。 \end{tabularx} \\ \\ \relax
28:18 & \begin{tabularx}{0.7\textwidth}{X} 耶穌進前來，對他們說：「天上地下所有的權柄都賜給我了。 \end{tabularx} \\ \\ \relax
28:19 & \begin{tabularx}{0.7\textwidth}{X} 所以，你們要去，使萬民作我的門徒，奉父、子、聖靈的名給他們施洗， \end{tabularx} \\ \\ \relax
28:20 & \begin{tabularx}{0.7\textwidth}{X} 凡我所吩咐你們的，都教導他們遵守。看哪，我天天與你們同在，直到世代的終結。」 \end{tabularx} \\ \\
[1ex]
\hline
\hline
\end{longtable}
$^{1}$今日的經訓是記載在.
《馬太福音》28章18至20節.
在大家手上的聖經新約第46頁.
請聽上帝的話語.
耶穌進前來對他們說.
天上地下所有的權柄都似給我了.
所以你們要去使萬民作我的門徒.
奉父子聖靈的名給他們施洗.
凡我所吩咐你們的都教訓他們遵守.
我就常與你們同在直到世界的末了.
弟兄姊妹新年快樂.
祝願大家在新的一年裡.
生命更進步更長進.
來到新的一年.
我們開始第一次的港道系列.
我們用6G.
我們的6G作為我們這個系列.
弟兄姊妹問牧師6G是什麼來的?.
外面坊間最新的上網都是去到5G.
教會教新科技.
其實6G不是教會教什麼新科技.
我們是說以6個G字頭的英文字母.
成為教會一個未來長遠發展的核心理念.
第一個G就是Generating Disciple.
教會是成為一個門徒不斷產生門徒的一個群體.
我們希望旺角浸信會在未來的日子裡.
不單止是讓人來到這裡敬拜.
我們實踐的那個在地上的目標.
我們要建立更多的門徒.
所以我們在今年開始的時候.
我們就建立門徒.
門徒不斷作為我們第一次崇拜的港道題目.
不過不知道你想起門徒訓練.
或者你成為門徒.
你會想起什麼?.
坊間裡面有很多關於門徒訓練的書.
我都看不少的.
不過是不是看完這些書.
拿了這些教材之後.
我們就自動成為一個門徒呢?.

$^{41}$其實在聖經裡面.
今天這一段經文.
耶穌基督如果細緻一點去看的時候.
他會有一些比較重要.
但是我們平時忽略了的教導.
希望今天我們在這段經文裡面.
我們找到一些線索.
找到一些方法來栽培我們.
讓我們生命慢慢走向一個門徒的階梯.
不過在仔細看聖經之前.
跟大家講一個經歷.
不知道大家有沒有試過的.
有沒有大家試過背著背囊.
做背包客走去世界不同地方去旅行.
有沒有試過的?.
有些試過.
但大部分都沒有試過.
我聽到一個雷同的例子.
不過將它改換了.
話說有一個背包客去到澳洲.
他就搭火車.
就想著去到他前往的目的地.
這個背包客一路走的時候.
他知道反正都去到澳洲這個地帶密駁的地方.
當然要走走的.
他問過人.
由他搭完火車之後.
去到那個目的地.
大概半個小時的路程.
他心想半個小時.
背著背囊都可以的.
於是乎他就搭火車去到那個地方.
於是就徒步.
就想著用半個小時.
走去他心目中想到的目的地.
誰知走了半個小時之後又走半個小時.
為何走這麼久都還未到呢?.
最後就有機會.
打聽了附近的商舖.
他說半個小時是車程.

$^{81}$他說如果是車程.
在澳洲的車開一百公里.
如果半個小時是有五十公里.
五十公里步行走到天黑都還未走完.
再加上那段時間太陽開始下山.
開始收.
這個背包客就很擔心.
死了都不知道怎樣走多久才去到目的地.
於是乎他就迫於無奈.
就在路上這樣折手.
折那些順風車.
折了很久都搞來搞去都還未找到一個人來接載他.
他很沮喪快要累到沒力氣之際.
有輛車就停下來了.
就問他朋友你想去哪裡?.
就車來.
因為這個背包客實在太累了.
於是上車再說.
於是乎他就馬上上了車.
就馬上睡著了.
司機就覺得讓他休息一會.
隔了一會之後.
司機就轉頭問這個背包客.
朋友你的目的地去哪裡?.
讓我載你過去吧.
我都差不多到了.
這個背包客就說.
你的車實在太舒服了.
坐在椅子上的座椅很舒服.
你開著音樂這麼好聽.
冷氣又足夠.
你又這麼好人.
這裡就是我的目的地.
如果你是那個司機.
你開車過來會怎樣?.
你可能會趕他走.
你會覺得這件事很荒誕.
不過很多時候人生都像坐背包客一樣.
我們坐順風車.
但我們不知道自己原本想去的地方.

$^{121}$坐一些情車覺得很舒服之後就不去.
很多基督徒都是這樣的.
回到教會之後.
這個地方的人又溫馨.
又有這麼好笑容.
四萬港口.
那些傳道同工.
那些弟兄姊妹.
這裡的椅子又舒服.
音樂又好聽.
冷氣雖然正在調整.
但現在都還可以.
這個地方簡直是一個無與倫比的地方.
這個地方就是我人生的目的地.
不過如果我們帶著這樣的信念.
或者這樣的看法.
渡過我們的信仰生活.
似乎就和耶穌基督對我們的召命.
或者對耶穌基督剛才我們所讀的瑪達福音裡面.
就會製造一個很大的距離.
今日這一段經文希望.
當我們仔細讀的時候.
我們找得到.
我們生命不是追求一個安舒.
或者令到我們覺得很舒適的生活.
我們的生命是要朝著上帝為我們預備的目的地來進發.
門徒訓練.
或者我們生命裡面成為一個門徒.
朝著上帝在我們生命裡面不同的計劃.
我們朝著這個方向一齊努力.
希望今日這段經文幫到我們.
讓我們在新的一年裡面.
我們帶著這個召命過我們的人生.
我們一起祈禱.
求上帝幫助.
天父上帝求你在我們當中對我們每一個講說話.
幫助我們中弟兄姊妹.
讓我們能夠在你的話語裡面.
我們知道你要我們所行的路.
讓我們行走人生的路不偏離.

$^{161}$幫助孫講的這段經文講得清楚.
讓弟兄姊妹聽得明白.
在生活裡面實踐你給我們的召命.
又傳記在心裡.
讓我們一生不偏離.
恭敬將來的時間交託.
願你賜福.
獻上祈禱.
奉靠基督耶穌得勝的聖名祈求.
Amen.
孟大福音在這裡裡面其實有一個很簡單的記載.
耶穌基督講這個的大使命.
其實是在他實踐完這個救贖工作.
在他釘身在十字架第三天復活之後.
將要升天之前.
他就對門徒講的這一番說話.
耶穌進前來對他們說.
天上地下所有的權柄都賜給我了.
所以你們要去.
使萬民作我的門徒.
奉父子聖靈的名給他們施浸.
凡我所吩咐你們的都教訓他們遵守.
我就常與你們同在.
直到世界的末了.
在經文裡面.
耶穌基督將這個很重要的使命告訴門徒.
在他將要升天的時候.
他要將一個福音的棒.
將一個使命交付給他剩下來在地上的這十一個門徒.
這個過去裡面有三年多的日子.
跟隨耶穌的這一班人.
耶穌就要將這個棒交給他們.
值得注意的是.
為什麼這班門徒要去實踐這個使萬民作.
耶穌基督門徒的這個使命呢?.
耶穌怎樣講這句話呢?.
你仔細去看.
耶穌進前來對他們說.
天上地下所有的權柄都賜給他了.
耶穌為什麼要這麼強調.

$^{201}$天地的權柄都在他身上呢?.
其實如果你細心去想.
耶穌基督一直都有這個權柄的.
耶穌不是在十字架之後.
上帝無端端後來.
你這麼偉大.
後來才加個權力給他.
或者給一個天地的主權給他.
不是這樣的.
在我們認識的聖經裡面.
在約翰福音裡面都講到.
基督耶穌鍍成了肉身.
他本來就是創造天地萬物的主.
萬有都是藉著基督耶穌而做的.
所以基督耶穌本身就有天地的權柄.
不過因為要作這個拯救的計劃.
所以耶穌暫時將他的權柄收起來.
在菲勒比書第二章六至十一節.
保羅就有一個很清晰的注釋告訴我們.
他本有神的形象.
不以自己與神同等為強奪.
反倒虛己.
取了奴僕的形象成為人的樣式.
既有人的樣子.
就自己卑微.
全心信服以至於死.
且死在十字架上.
所以神將他升為至高.
又賜給他那超乎萬名之上的名.
叫一切天上的地上的和地底下的.
因耶穌基督的名無不屈膝.
無不口稱耶穌基督為主.
使榮耀歸於父神.
保羅為了耶穌基督的權柄.
他要做一些注釋.
本來耶穌就是萬有的主.
他擁有天地之間所有的權柄.
因著拯救你和我.
他放下了自己的那種能力.
他不以與神一樣.

$^{241}$來與上帝爭取.
在菲勒比書裡面這樣說.
不以自己與神同等為強奪.
反而耶穌基督用一種最謙卑的方法.
謙虛自己.
謙卑到一個極致.
就是死在十字架上.
耶穌基督在這裡就像下樓梯一樣.
來描述他.
本來他擁有高天的能力.
但是一步一步將自己放下.
放下 再放下.
甚至乎犧牲在十字架上.
所以耶穌基督本來有權柄.
但是他主動的來放棄.
成為一個最卑微的樣式.
來成就救贖的計劃.
當他死在十字架之後.
第三天 上帝叫他復活.
上帝將原本他擁有的權柄.
重新一次歸回在基督耶穌那裡.
所以保羅說.
天地萬物 天上的 地底下的.
和地上的.
因著基督耶穌的名.
無不屈膝.
無不口稱耶穌基督為主.
使榮耀歸於父神.
所以在剛才耶穌所說的天上地下.
所謂的榮耀是有個階段的.
本來他擁有榮耀 擁有萬有.
後來他自己放下.
現在再重新拿回.
他要宣告地上的一切.
都是伏在耶穌基督的主權之下.
沒有一件事可以走出耶穌基督的授戰那當中.
即使未信主的人.
基督耶穌都是他們的主.
不過他們不知道自己是.
伏箭在基督耶穌的威權之下.

$^{281}$耶穌基督為什麼要頒佈使命的時候.
要說一句這樣的說話呢.
其實耶穌說這番使命.
就不像以前所說的輕描淡寫的一個吩咐.
當他升天的時候.
他要將一個門徒必須要.
實貫切執行的一個命令.
實踐在他的生命當中.
所以他說.
所以你們要去使萬民作我的門徒.
那個因為是什麼呢.
因為耶穌擁有權柄.
因為基督耶穌擁有天上地下所有的權柄.
他用這個大帝來處理.
告訴你他按著天地主的身份.
命令我們每一個人都要承擔住.
這個大使命的召命.
使萬民作基督耶穌的門徒.
教會做門徒訓練的工作.
弟兄姊妹建立門徒.
不是一個選擇性.
或者我們按著我們的喜好.
可做可不做的一個選擇.
當教會搞不清的時候.
很多的教會就會變成一個娛樂機構.
容許我直接說.
很多弟兄姊妹聚集為了感覺良好.
所以我們會做一些我們覺得.
自己哄自己快樂的一些程序或者節目.
不過地上的教會.
目的存在不是要製造更多的娛樂節目.
地上的教會的存在.
是要建立更多的門徒.
我們每一個基督徒.
我們每一個信主的人.
我們都不可以選擇性地.
這些偉大的東西交給梁偉做.
或者交給年輕人做.
我自己免於做這個門徒.
建立門徒的這個職事.

$^{321}$耶穌說的這個是一個普世性.
每一個人都要承擔著的一個命令.
耶穌以天地的主這個身份.
命令我們每一個.
所以我們就不可以選擇.
我們每一個都要承擔著.
耶穌給我們的這個命令.
這個是第一點我們需要留意.
耶穌以他的權柄命令我們.
實踐這個建立門徒的職事.
第二件事.
耶穌用他的權柄.
他都告訴門徒知道.
這個世界是屬於他的.
這個世界天上地上和地底下的.
所有萬物都是伏在耶穌的主權之中.
即使可能有些未知未信的人.
他不覺得自己是耶穌為他的主.
但是耶穌說.
他都是這些人.
這些萬物的主.
他要將這個主權宣示出來.
如果我們要將門徒訓練.
講述出來的時候.
我們不是純粹講一個訓練課程.
令到人活得更加好.
我們最終要宣講的內容.
就是基督耶穌是這個世界的主.
是我們生命裡面最重要.
不可以失去的一個最重要的訊息.
耶穌以自己的權柄.
以天地的主來自稱.
他要告訴我們.
他要命令我們.
要將這個訊息告訴這個世界.
都必須要知道的一個很重要.
福音的內容.
所以你看到.
為什麼我們要做門徒訓練呢?.
耶穌的命令.

$^{361}$和這個也是耶穌基督的心意.
他要將他自己的身份.
將他的能力表達出來.
第二件事就是.
既然我們知道.
原因是耶穌基督有這個權柄.
他給的結果就是.
最重要的命令字眼.
就是使萬民作我的門徒.
一講起門徒這個字.
其實很多人和很多字眼.
混淆了.
我用一點時間解釋給大家聽.
門徒在聖經裡面.
你猜猜出現多少次.
在新約裡面.
多不多呢?.
多少呢?.
261次.
大家在程序表記住.
門徒這個字眼在聖經裡面出現261次.
你經常都說的.
你介紹自己的信仰是什麼?.
我是基督徒.
基督徒這個字眼出現多少次呢?.
三次.
其實我們經常說的.
我們只是出現三次.
有一個多一點點的就是信徒.
還要便宜那種.
便宜信徒.
信徒在聖經裡面出現多少次呢?.
九次.
你看到聖經的技術.
你都知道那個側重面.
它講門徒多一點.
講基督徒和信徒.
這兩個字眼是少一點的.
基督徒其實怎樣理解呢?.
或者這個字眼本身有什麼意義呢?.

$^{401}$其實在早代教會裡面.
基督徒這個字眼.
其實是一個貶義詞.
我之前都和大家分享過.
那些人歸信耶穌基督.
以耶穌為主.
那些外邦人或者那些未信主的人.
看到那些人信耶穌.
他們的道德標準又高.
去過那些神殿祭過的那些祭偶像的肉.
他又不吃.
這班人又很規矩的.
眼前見到這些人被人欺負又不出聲.
所以世俗那些人就笑第一代的.
那些信主的人就叫基督徒.
其實簡單來說就是基督仔.
如果等於同義詞就是傻仔.
信耶穌而變成一個傻奸奸那些人.
這個是一個貶義的字眼.
這個字眼由安提柯教會開始.
大概一班基督徒很謹守自己的生活.
所以當時的世俗那些人.
就執著這些人的身份就笑他了.
基督徒這個是一個這樣的字眼.
不過後來到了第二世紀.
在護教學那些神學家.
去到羅馬君王面前.
他要展明他的信仰的時候.
他就用這個基督徒這個字眼.
重新來說基督徒一班信耶穌的人.
其實是一班很規矩的人.
即使被人嘲笑取笑.
他們都堅持行這個上帝的道.
所以去到第二世紀去到今天.
基督徒就成為我們今天的尊稱.
我們會稱呼自己為基督徒.
由一個貶義慢慢在文化裡面改變.
就成為一班信主的跟隨者.
所以我們稱呼基督徒.
其實就沒有什麼很重要的內涵.

$^{441}$在聖經裡面不是說很多著跡.
是描述性的一個字眼.
第二個字眼就是信徒.
信徒這個字眼在聖經裡面.
就不是相對於世俗.
就相對於教會裡面的管理.
信徒和那個聖職人員就做一個對比.
信徒就是廣泛的弟兄姊妹.
信主在教會裡面的弟兄姊妹.
而聖職人員就是教會的領袖.
信徒和那個領袖某程度上.
在屬靈權兵上好像有些差距.
不過他講的是一個身份.
和人生裡面實踐信仰的場景不同.
教會領袖就是成為上帝的牧羊人.
帶領教會為教會犧牲自己.
那些就叫做聖職人員教會領袖.
而其他那些可能一直在教會裡面生活.
繼續來到生活裡面作見證的.
那些就叫做信徒.
所以你看提摩太前書裡面你會知道.
保羅都是這樣講的.
不要叫人小看你年輕.
總要在言語行為信心愛心清潔上.
都作信徒的榜樣.
一個聖職人員要成為信徒的榜樣.
信徒這個字就對比聖職人員.
有一個這樣的描述性的向度給大家知道.
所以不會有人.
或者你以信徒作為你人生目標.
因為信徒都是在描述你和聖職人員之間的對比.
這個都是一個描述性的字眼.
叫你的生命和這些聖職人員.
和這些律師傳道分別出來.
但是沒有講到褒義或者貶義.
去到今天因為可能有不同的課程.
或者令到大家更加明白.
所以我們在信徒裡面再加一個.
平字.
平信徒.

$^{481}$不過弟兄姊妹.
雖然你是平信徒或者平信徒.
但是你擁有的恩典.
你所接納的是一個重價的恩典.
我們每一個都是在上帝的主權底下.
我們都是被稱呼上帝的子民.
弟兄姊妹.
我們沒有分地位高與低.
我們只是在講我們實踐信仰的那個場景不同.
那你講了基督徒和信徒.
但是對於門徒的那個內涵就豐富很多.
有261個次數來講這個門徒生命.
你會看到那個內容是豐富很多的.
不過我歸納出.
作為門徒有兩個很重要的身份和內涵.
第一個門徒就是學生.
門徒一生是學東西的.
我們每一個我們在學什麼呢?.
我們在學耶穌基督的榜樣.
以耶穌成為我們生命的師父.
按著這個師父的教導和言行舉止.
我們就跟著他.
拉著他的衣服.
觀察耶穌基督.
或者在聖經裡面所記載的耶穌基督的言行.
我們在生活裡面一步一步來實踐出來.
學生不能高過先生的.
我不會說我做的事叫先生跟著我的.
我們只是跟著耶穌基督的足印.
一步一腳印.
來按著他的吩咐.
實踐信仰的生活.
今天對於使徒關係或者學生關係.
我們都開始有一點混淆.
我們覺得學校就是一個教育機構.
老師只不過是課堂裡面很短暫的.
教你東西的一個資訊傳遞者.
不過在聖經裡面所說的使徒關係.
比起我們今天所說的教學.
甚遠得多.

$^{521}$那種是一種身體力行.
一個生命裡面所有東西.
交流很緊密的一種狀態.
我想起我自己.
我家裡是做製衣的.
在我爸爸生前的時候.
他也有收徒弟做製衣.
因為我爸爸也是學師學出身的.
我記得以前在我成長的時候.
我看著我爸爸在那間山寨廠.
收幾個學徒.
那些哥哥每個都十多歲.
現在少年部那些.
就是我爸爸那些學徒.
我差不多一年多了.
十多歲就跟我爸爸學師.
某程度上來講.
我看到最粗重或者最厭惡性那些.
最悶的那些.
我爸爸都會給徒弟做的.
因為他拜師學藝.
這些東西當然要做的.
沒理由師父做完.
徒弟就不滿意.
我記得以前剪線頭.
搬貨.
或者做那些物流.
等等那些.
以前是有些徒弟.
粗重的那些.
全部由他們來做的.
你可以這樣講.
某程度上他就像廉價勞工一樣.
所有東西就跟進做完.
但是我爸爸不單止叫他們做這些東西.
到我爸爸去到一些大廠.
去談生意.
或者去做一些統籌.
例如訂船,訂紙箱那些.
他就會帶徒弟出去.

$^{561}$跟進學習.
別人在做事的時候不要問.
別人在做事的時候問.
會被人割頭.
旁邊就抄筆記.
或者觀察一下怎樣做.
師父做了什麼.
徒弟在旁邊觀察.
學著學著.
學到差不多滿師.
他就會跟我爸爸講.
我記得之後.
我學完了.
我想滿師.
我就出去自己開廠了.
就會有封利是給他.
祝福他.
他就自己出去開廠了.
在使徒關係裡面.
生活緊密到一個地步是什麼呢?.
以前那些哥哥.
我爸爸的徒弟.
跟我們家一起住的.
以前我們住在廠那裡.
我記得有一個情景很有趣的.
我爸爸那時候因為趕課.
就睡在那張財床那裡.
那些徒弟就墊著地上.
有些衣服墊著地上.
睡地上.
有些厚的雪褸墊著地上睡.
就像現在這種天氣.
晚上睡一睡就冷了.
我爸爸就很懶.
但又看到徒弟們冷就縮一下.
扶開一些被子.
於是就在財床上找兩隻腳趾公.
用腳趾公的腳趾.
鉗衣服就蓋著那些徒弟.
讓他們不要冷.

$^{601}$做師父就是這樣做的.
但是蓋一蓋又要扶被子.
蓋一蓋又要扶被子.
我記得第二天早上見到.
我爸爸一起來就看到.
整個地都是衣服.
徒弟們把衣服扣住.
把衣服抹在地上.
整個地都是衣服.
但是師父和徒弟之間.
就是一起來生活.
我講這件事.
其實某程度上耶穌和他的門徒.
是這樣的生活.
不是純粹是一個短暫式.
我教了你東西.
然後一句謝謝就搞定.
是說一起生活.
徒弟學生觀察老師.
怎樣來做.
他就照板煮碗.
或者在旁邊問問題.
到不明白的時候.
在老師或師父做完事之後.
就去發問.
藉此解開他心中的疑團.
如果說門徒訓練.
其實我們應該要這樣跟隨耶穌.
用一種這樣的方式.
在師父面前.
我們要放下社會裡.
引以為傲的身份.
可能在裡面.
你是一個很顯赫的人.
可能你是一個CEO.
但在基督耶穌面前.
你都是一個學徒.
我們生命裡.
我們都是看著耶穌基督.
叫我們做的工作.

$^{641}$可能在你的工作崗位裡.
這些粗重的事情.
或者聯絡這些瑣碎的事情.
你叫個秘書那些做.
不過在主裡面.
當耶穌基督叫你去這樣做的時候.
如果耶穌基督.
呼籲我們要傳福音.
要在街上派傳單和傳福音.
我們覺得很尷尬也好.
我們都要勉為其難.
我們都要勉強自己.
或者我們都要自我鼓勵.
來實踐這些教導.
門徒.
就是要學習師父的.
所言所行.
耶穌他怎樣做的時候.
我們就按著他的樣式.
實踐我們的生活信仰.
這是門徒第一個意思.
門徒第二個意思.
是跟隨者.
跟隨者的意思.
是說到.
這個師父是我跟硬的.
我效忠在.
這個師父那裡.
我不會.
一路跟著他.
心裡一直在說.
這是我的.
我不會跟他.
也不會做一些事背叛他.
跟著師父來背叛他.
這麼荒謬的事你都會做.
不過今天這個世界很多人都是這樣.
稱耶穌基督為主.
但背地裡有時候.
做一些反信仰.

$^{681}$反福音的工作.
不過如果我們貫徹.
做門徒的生活.
我們是一個跟隨者.
看著基督耶穌.
在心靈裡.
我們百分百效忠於他.
我們的命.
我們的人生.
以他來定義我們.
而不是我們喜歡怎樣.
過我們自己.
自由自在的生活.
所以在這裡.
耶穌基督.
所說的命令.
用了這麼多篇幅告訴大家.
他用了一個很強烈的語氣.
叫我們.
要建立.
跟隨基督耶穌的門徒.
不要.
按著我們自己的喜好.
不是按著我們自己.
覺得自由自在.
覺得教會很安恕地過我們的宗教生活.
我們覺得自己成為一個.
平信徒就足夠了.
反而是.
我們成為一個.
緊緊跟隨主的學生.
我們聽主的吩咐.
我們也跟隨他.
效忠於他.
耶穌以這種命令.
來訓練.
那十一個門徒.
如果你抽空一點.
退後一點來看一下.
這十一個人都很誇張.

$^{721}$因為這十一個人.
領受了耶穌基督的吩咐.
這個世界就.
翻天覆地.
來改變.
由十一個人演變成為一個.
普世的.
基督教.
由十一個人.
朝夕相對三年多的日子.
後來他們就成為.
一班勇猛的.
一班人.
將人一個一個帶到在上帝的面前.
藉著.
耶穌對他們的訓練.
耶穌的吩咐.
他們又將這個吩咐來訓練其他的人.
一代又一代.
二千多年.
毫不間斷的.
是今天的教會.
弟兄姊妹.
我們今天都在接受這個棒.
今天旺角浸信會都成為.
一個這樣的教會.
我們不可以成為一個娛樂機構.
我們必須要成為一個.
建立門徒.
實踐門徒訓練的一個教會.
可能我們有些地方.
做得不夠成熟.
但是我們不可以因為做得不成熟.
我們就不去做這件事.
反而是要令到我們自己.
變得更加成熟.
變得更加有毅力.
來刻意地.
讓人成為一個門徒.
那你說.

$^{761}$牧師那怎麼做呢?.
28章19至20節.
上面.
耶穌教了給我們.
所以你們要去.
使萬民作我的門徒.
奉父子聖靈的名.
給他們施浸.
凡我所吩咐你們的.
請遵守.
在這節的.
子句裡面出現.
四個動詞.
最重要的動詞就是使萬民作我的門徒.
這是一個重要的.
命令,這個動詞成為.
耶穌基督呼召我們最重要的.
內涵,剛才我說了.
我不重覆了,不過.
怎樣去實踐呢?其他那三個動詞.
就是.
實踐了怎樣去將.
使萬民作我的門徒.
耶穌就說了.
第一個動詞,去.
我們要出去.
我們要成為一個.
打破自己.
很多的脅惡.
或者很多.
害羞,面子.
那些的規範.
我們要成為一個勇敢的人.
我知道.
我經常都說.
我覺得我自己都OK.
未至於很窄.
現在的人說很窄.
很害羞,我也不覺得自己是.
不過我回看我自己的童工.

$^{801}$我自己都很慚愧.
舉個例子,梁偉那些.
買一棵菜.
跟人家談半小時.
買一棵菜都會講福音.
我那次跟他去台灣.
跟人家談,好像自己的鄉下.
但是.
這些天生的.
是搭訕高手.
所有事情談著談著就會談到福音.
真的很厲害.
曾珍牧師也是很厲害.
樓下幾個,每個禮拜回來.
有一個派單張的姐姐.
派賣眼鏡.
那些,我們平時.
只見到人家派傳單,拒絕人家.
曾珍跟他們做了朋友.
跟他們.
抱頭抱脖子.
我心想,你們用不著這麼誇張.
那些社交能力.
強到不得了.
我跟那些官家叔叔.
打個招呼,我已經超越了.
我的安書區.
我常常覺得自己是這樣的.
不過,我的性格.
是這樣的,或者外向一點.
內向一點,但是我們不可以.
因為我們的性格.
好像我這樣,覺得走出一步.
已經做了很多事情.
我們永遠不可以限制自己.
永遠都是在這個階段.
裡面停留的.
我們未必像梁偉.
曾珍,或者其他牧者.
這樣這麼外向.

$^{841}$但是我們可不可以比過往的自己.
走前一步.
主動多一點呢.
我們的社交圈子.
我們傳福音的那個.
我們的圈子.
那個同心圓,可不可以闊一點.
我們對人接觸的.
可不可以宏大一點.
不要只是接觸.
來來去去,那兩三個.
那幾個摯友.
就停止了我們.
擴展我們社交圈子的.
那種.
那種進取的心.
我們未必成為一個.
很外向的人,但是可以比.
昔日的自己,可以主動一點.
聖經裡面.
告訴我們,去這個字眼.
我們就是要打破了.
我們自己很多,給自己很多的枷鎖.
可能我們覺得.
很想要面子,又很怕別人拒絕.
等等那些東西,令到自己.
寸步難行.
在新的一年裡面.
我們不如放下了這些.
基本我們.
令到我們不夠膽主動的.
一些心態.
反而讓我們可以.
勇敢踏前的,讓我們不只是.
在一個安書區裡面.
來到去純粹.
感覺這個信仰.
更加重要的是.
有時候要讓自己刻意.
令到自己面皮厚一點.

$^{881}$其實別人拒絕你.
其實很平常.
別人不跟你.
聊天,其實沒什麼所謂.
進不了你的肉.
不會說心靈很大受傷.
其實想想別人拒絕你.
你天天都在拒絕別人.
有些電話打給你,說你要不要借錢.
你都拒絕別人.
其實我們很容易拒絕別人.
為什麼被別人拒絕.
會覺得這麼受傷害呢?.
放下自己.
一些不必要.
但很重要的面子.
面皮厚一點.
這個會令到我們.
可以放膽出去.
一個實踐.
基督耶穌給我們門徒訓練.
一個很重要的照明.
我記得有一次.
在我們.
去到一些很敏感的地方要傳福音.
普通話又不太行.
但是那個傳道人說.
Anders,去的了.
那說什麼?.
你不試試,你怎知道自己在說什麼?.
你不說,你也不知道自己在說什麼.
只是說話.
用一些很噁心的普通話.
去說.
但是噁心也好.
你願意誠懇溝通.
你會發覺.
我的信仰又踏前一步.
就是結結巴巴.
但是.

$^{921}$主動去做的時候.
你會經歷到那份滿足感.
但是如果我在這些位置.
我退後,不試.
我的人生又變得蒼白.
我又浪費了訓練自己的機會.
所以弟兄姊妹.
第一個.
門徒訓練怎樣去做?.
第一個,去.
走出你的安恕區.
第二個重要的字眼.
就是施浸.
奉父子聖靈的名.
給他們施浸.
施浸的意思.
我們想著是否帶了他們去信主.
令他們去浸禮.
我們就做了我們要做的功夫呢?.
施浸的意思.
不是純粹說做一個禮儀.
那個禮儀在聖經裡面.
不是那個重點.
成為門徒才是重點.
不過在.
初代教會裡面那個施浸.
其實是一種.
視死如歸的一個表達.
當他將全身.
抹入水中.
他就表達了.
他要和死同埋葬.
同復活.
他要和自己過去.
一刀兩斷.
舊時他拜偶像.
他藉著這個浸禮告訴所有人.
一刀兩斷.
甚至乎他因為信仰的緣故.
他信耶穌.

$^{961}$他可能整個宗族都要.
和他為仇.
他社交斷裂的了.
但是他都要堅心來做這個浸禮.
代表著.
視死如歸.
和自己的過去一刀兩斷.
這個施浸.
與其說一個禮儀.
更加準確一點說的就是.
一個公開表達自己.
對基督耶穌的效忠.
禮儀是一次的.
但是門徒生命是一世的.
他用這個禮儀.
來表達他一生跟隨主.
是一種很極端.
很強力的一個禮儀.
表達自己.
我們會再走生命裡面的回頭路.
領姊妹.
我要在裡面挑戰大家.
在座裡面我相信有很多人已經接受浸禮.
但是我們會不會.
做了浸禮,做了禮儀之後.
在生活當中.
在我們人生裡面.
我們走回頭路.
我們跟自己的過去畫一條清晰的界線.
我們覺得做了禮儀之後.
我們已經滿了詩.
我們覺得我們的生命.
停止在這個終點站.
其實耶穌說的那個不是終點站.
那個是中途站.
我們的生命不是停留在浸禮裡.
也不要代表我做了浸禮之後.
沒有人監察著自己.
或者怎樣.
求主憐憫我們.

$^{1001}$在新的一年裡.
這個浸禮是說我們一生堅實地跟隨著.
不要容讓自己走回頭路.
過回舊時在罪中.
或者沒有信心的那些日子.
不要再眷戀自己在罪或遺憾當中.
繼續打轉的生活.
當我們做了這個浸禮的時候.
我們就挺起胸膛.
曾經有些大英姐妹問我.
牧師.
其實現在很多教會.
一信主就可以浸禮.
為何黃浸的浸禮要搞五位多東西.
又要回來一年.
還要讀浸禮班.
又要寫得救見證.
還要預演一次.
講一次那個見證.
還要考問信德.
過五關斬六將.
為何不簡單.
有些教會一信.
就立即在泳池.
或者在沙灘浸.
我們搞那麼多東西.
會不會難了我們福音的工作.
我通常的答法.
我都想簡單一點.
不過如果這是上帝吩咐.
或者是我們的宗派裡面的強調.
是要門徒才可以浸禮.
我們對於這個浸禮.
就不可以掉以輕心.
在我們浸信會裡面所講的浸禮.
是在講門徒才可以浸禮.
是一個切實跟隨主的人.
才用這個浸禮來表達.
他自己的信仰.
不是一個輕忽的舉動.

$^{1041}$而是一個慎重.
想清想楚.
與自己過去斷絕的一個決心的表現.
基督耶穌叫門徒.
藉著這個浸禮.
來幫助那些將來的.
那些人.
去做他們的夢想.
去做那些將來的.
一些門徒.
讓他們過著一個認信.
決心跟隨主的一種生活.
所以有一些.
譯本剛才在讀的那段經文.
奉父子聖靈的名.
給他們施浸.
有一些譯本就會譯為.
受浸歸入.
父子聖靈的名下.
那個是在講歸入.
是一個認信的表現.
藉著這一份的認信.
走他自己門徒生命裡面.
一個重要的里程.
這是第二個動詞.
第三個動詞.
教訓他們遵守.
很多時候我們看到教訓.
我們看到教訓這個字很快就停在這裡.
要教要講很多東西.
對的沒錯的.
不過教訓的目的.
不是我講了一些東西出來.
是聖經要講.
很多弟兄姊妹經常說.
對生活不滿意.
或者有些信仰卡住.
牧師你會不會在講道裡面.
又講一次出來.
我說OK我會講的.

$^{1081}$不過目的不是要講出來.
遙遠發功.
就等於門徒訓練.
講之外最重要的.
就是要大家一起遵守.
講了不遵守其實浪費力氣.
講沒有行動.
這代表著法利塞人那種說教的生命.
不過基督耶穌說.
教會的目的.
門徒訓練就是.
言教也要新教.
幫助弟兄姊妹遵守.
上帝的話語.
我們願意實行.
我們一起來踐行.
基督耶穌給我們的照明.
早陣子.
我聽過一些教會令我很羨慕.
他們講他們的小組生活.
他們說他們的小組生活.
其實弟兄姊妹分組不是常常分的.
我們旺鎮就大概一段短時間.
一兩年左右我們會做一次重新分組.
他們不分組.
我說為什麼會這樣做呢.
他說一種這麼深度的同行.
其實目的就是要弟兄姊妹.
過一個極之交代的生活.
因為有足夠時間.
來相處.
所以弟兄姊妹生命的盲點.
其他組員是知道的.
在這個小組裡.
不是常常分的.
因為這些罪或生命裡的破口.
不容易跟陌生人說的.
所以在一個很長遠.
關係很密切的小組裡.
這些人就會打開自己的心扉.

$^{1121}$將自己生命的困難說出來.
在查經的時候.
不是純粹說一個對的道理.
當這個對的道理.
當基督耶穌的吩咐說出來的時候.
明白了.
跟我生命的距離有多遠.
在我實踐聖經的教導的時候.
我有什麼掙扎呢.
這些掙扎就在這個組裡.
互相分享.
大家互相一起來同行.
這是一個坦誠的表現.
將自己揭出來.
揭出來之後.
其他人就不是笑他或嘲笑他.
揭出來之後.
其他人就一起來守望.
這個不完美的弟兄.
每一個人都將他自己不完美的一面.
交代出來.
我們就互相砥礪.
在這個小組裡.
我們一起來交代.
一起來跟進不完美的地方.
但是下一年.
或者過了一段時間之後.
我們有什麼地方進步呢.
我們就要這樣互相交代.
丁子會.
你今天你的小組生活是怎樣的呢.
甚至乎.
牧師我沒有想過參加小組.
我只是想聽到就夠了.
原來.
我們今天.
如果我們信仰.
沒有一種深度.
將自己揭出來的那種表達.
那種真誠.

$^{1161}$那種勇氣.
大概我們都是浮游在一個很低層次的層面.
根本我們就只聽到.
那種真誠.
根本我們就只聽到.
是沒有辦法做到門徒訓練.
唯有我們的生命.
揭開.
我們一起來接受教導.
一起來抵禮遵守.
有行為的改變.
才有生命的改變.
求主繼續的來幫助我們.
基督耶穌.
祂叫我們.
門徒訓練.
是要幫助人領受教訓.
和一同的來遵守.
不要浪費了.
教會這個群體.
不要純粹.
要人家的恩典.
不要弟兄姊妹的生命.
對我們一些.
甚至乎帶著愛心的責備.
或者一些恩典同行.
我們的生命.
不是孤軍作戰.
而是一班人走.
這個門徒訓練的路.
耶穌基督說了這幾個動詞.
他說.
耶穌基督說了這幾個動詞.
說了怎樣做之後.
他又給了一個應許給我們.
耶穌基督說.
我就償願你們同在.
直到世界的末了.
如果你要經歷神的同在.
你一定要建立門徒.

$^{1201}$今天很多人想.
神同在是什麼呢?.
很隱蔽的那些祈禱.
看到天裡面.
那些很神秘的經歷.
或者.
好像有一種超自然的能力.
或者好像有一種超自然的能力.
祈禱那種.
很厲害的能力.
或者那種環境氣氛.
令到很有宗教感.
我們覺得那些叫神同在.
不過.
在聖經裡面.
神同在從來和使命走在一起.
你不實踐上帝的召命.
你沒辦法經歷神同在.
可能我們偶爾會有一些的祈禱.
會有一些從神而來的感悟.
不過上帝不是這樣純粹讓你過渡你的宗教生活.
上帝要我們帶著祂的使命.
經歷祂的同在.
有時候神的同在.
未必是一些很超自然的東西.
不過你會看到有些神跡.
在人身上出現.
今年是2022年.
2005年.
17年.
我在混浸17年的時候.
我看到上帝同在很多的見證.
我曾經見過一些年輕人.
學校都放棄了他.
在教會.
說他是天王級的人.
教會的那些.
有些牧者罵他.
趕他出去不讓他崇拜.
放棄他.

$^{1241}$覺得他阻礙地球轉.
很多人都放棄他.
但只要用心.
帶他成為門徒.
告訴他.
他是耶穌所愛的.
這個奇奇怪怪的人.
我看著上帝由這個.
邊緣的年輕人.
慢慢成為.
一個可以委以重任.
為教會福音事工.
努力的一個.
中流砥柱.
在他心靈裡面.
常常想著福音可以怎樣傳開.
可以怎樣去理解聖經.
不是一個.
而是很多個.
當你願意.
當你和上帝一起實踐這個照明的時候.
你一定會見到.
這些生命的改變.
是神同在的神跡.
是世界都放棄的人.
但上帝沒有放棄他.
同樣可能我們自己生命也是.
當我們不放棄自己.
我們認定上帝有心意在我們當中.
我們是上帝的門徒.
我們可以學耶穌.
哪怕我們生命有多破爛.
我們闖開我們的心扉.
被主使用.
被主來動我們.
更新我們.
被主使用我們.
更新我們.
門徒的生命.
你就會經歷到不斷的豐富.

$^{1281}$不斷的插位.
但也是一個不斷進步的生命歷程.
不要介意自己有時候很醜.
將自己的錯漏.
將自己不理想的地方.
向上帝或一些弟兄姊妹.
去呈現出來.
能夠呈現自己的弱處.
那個人生命才是最幸福的.
因為我們找到自己的空間.
讓自己繼續有成長的位置.
但這個世界是在說.
包裝自己.
隱藏自己的弱點.
但在聖經裡面.
在這個門徒裡面.
我們就在這個主.
栽培的主裡面.
呈現我們的弱點.
當我們越容易打開我們的心扉.
被上帝動我們的時候.
我們就越能經歷從神而來的成長.
弟兄姊妹.
今天這段經文不是很複雜.
已熟能詳的一段經文.
可能你已經識背了.
但在我們生命裡面.
我們有沒有貫徹地.
讓自己的生命成為一個門徒呢?.
倘若你已經對信仰很認真.
我盼望你不要純粹徘徊在知識層面裡.
在你生命裡面.
你繼續將人帶到在上帝面前.
栽培更多基督耶穌的門徒.
讓他們的生命可以獨當一面.
讓他們一生走在基督耶穌的旨意裡面.
讓他們一生跟著耶穌的足跡而行.
可能你生命裡面.
有時覺得自己有不少面子.
在基督耶穌面前.

$^{1321}$不要說甚麼面子.
說一些你覺得很重要的事.
在基督耶穌面前.
最重要的就是我們是否確切地跟隨他.
求主繼續來恩待我們.
在新的一年.
我們彼此立志.
我們要成為一個門徒.
我們要建立更加多的門徒.
讓我們挺起胸膛.
成為上帝所使用的人.
上帝恩待你們.
(影片完結).
\newpage



\section{申命記 6:1-9}
\label{sec:Qke9jW9Cx6E}
\textbf{2022年1月16日早堂崇拜直播｜陳冠成牧師｜我們的6G - 聖經引導｜申命記6\_1-9}
\newline
\newline
連結: \href{https://youtube.com/watch?v=Qke9jW9Cx6E}{\texttt{https://youtube.com/watch?v=Qke9jW9Cx6E}} ~~~~ 語音日期: 2022-01-17
\newline
\newline
\hyperref[sec:oDel5UYyQt0]{\small{< < < PREV SERMON < < <}}
~
\hyperref[sec:index_chronic]{\small{[返順時目]}}
~
\hyperref[sec:9NLZUx6GLHU]{\small{> > > NEXT SERMON > > >}}
\newline
\newline
申命記 6:1-9
\newline
\begin{longtable}{cl}
\hline
\hline
章節 & 經文 (和合本修訂版)\\
\hline
6:1 & \begin{tabularx}{0.7\textwidth}{X} 「這是耶和華－你們的神所吩咐要教導你們的誡命、律例、典章，叫你們在所要過去得為業的地上遵行， \end{tabularx} \\ \\ \relax
6:2 & \begin{tabularx}{0.7\textwidth}{X} 好叫你和你的子孫在一生的日子都敬畏耶和華－你的神，謹守他的一切律例、誡命，就是我所吩咐你的，使你的日子得以長久。 \end{tabularx} \\ \\ \relax
6:3 & \begin{tabularx}{0.7\textwidth}{X} 以色列啊，你要聽，要謹守遵行，使你可以在那流奶與蜜之地得福，人數極其增多，正如耶和華－你列祖的神所應許你的。 \end{tabularx} \\ \\ \relax
6:4 & \begin{tabularx}{0.7\textwidth}{X} 「以色列啊，你要聽！耶和華－我們的神是獨一的主。 \end{tabularx} \\ \\ \relax
6:5 & \begin{tabularx}{0.7\textwidth}{X} 你要盡心、盡性、盡力愛耶和華－你的神。 \end{tabularx} \\ \\ \relax
6:6 & \begin{tabularx}{0.7\textwidth}{X} 我今日吩咐你的這些話都要記在心上， \end{tabularx} \\ \\ \relax
6:7 & \begin{tabularx}{0.7\textwidth}{X} 也要殷勤教導你的兒女。無論你坐在家裡，走在路上，躺下，起來，都要吟誦。 \end{tabularx} \\ \\ \relax
6:8 & \begin{tabularx}{0.7\textwidth}{X} 要繫在手上作記號，戴在額上作經匣； \end{tabularx} \\ \\ \relax
6:9 & \begin{tabularx}{0.7\textwidth}{X} 又要寫在你房屋的門框上和你的城門上。 \end{tabularx} \\ \\
[1ex]
\hline
\hline
\end{longtable}
$^{1}$以下我要為大家讀出聖經舊約新命記第六章一至九節.
是教會的聖經舊約二百二十二頁.
聖經這樣說.
「這是耶和華你們神所吩咐教訓你們的誡命,律例,典章.
使你們在所有過去得違約的地上進行.
好叫你和你子子孫孫一生敬畏耶和華你的神.
謹守他的一切律例,誡命.
就是我所吩咐你的.
使你的日子得以長久.
以色列啊,你要停,要謹守遵行.
使你可以在那流那雨蜜之地得以幸福.
人數極其增多.
正如耶穌你列祖的神所應許你的.
以色列啊,你要停.
耶和華我們神是獨一的主.
你要盡心盡性盡力.
愛耶和華你的神.
我今日所吩咐你的話.
都要記在心上.
也要殷勤教訓你的兒女.
無論你坐在家裡.
行在路上.
躺下,起來.
都要談論.
也要繫在手上為記號.
戴在額上為經文.
又要寫在你房屋的門窗上.
並你的城門上」.
弟子們,平安.
來到我們錄音系列的第三講.
聖經的引導.
Guiding Biblically.
強調教會其中一個很重要的核心價值.
就是讓上帝的話語.
來引導我們生活的一言一行.
話說猶太人有一個故事是這樣說的.
他決定在這個世界上.
揀選一個民族作為他的子民.
於是開始到處尋找.
他首先找到希臘人.

$^{41}$上帝對他們說.
如果我成為你們的上帝.
你們做我的百姓.
你們會為我做些甚麼?.
希臘人對他們說.
如果你成為我們的神.
我們做你的子民的話.
我們一定會用最精緻的藝術雕刻.
和最崇高的哲學作品來尊崇你.
來稱頌你.
上帝聽見後微笑.
但沒有接受.
只回答「多謝你的好意」.
然後離開了.
接著上帝又找到羅馬人.
對他們說.
如果我成為你們的上帝.
你們做我的百姓.
你們會為我做些甚麼?.
羅馬人回答.
偉大的宇宙君王.
我們是一個很擅長建築的民族.
假如你成為我們的上帝.
我們做你的子民的話.
我們會以你來命名.
建很多宏偉的建築景點.
並且我們會完善整個交通網絡.
方便所有百姓.
很方便地前來這些偉大的建築物面前.
來敬拜你.
上帝聽到後又微笑.
但他亦沒有接受.
只回答「多謝你的好意」.
然後又離開了.
後來上帝找到一群.
很懂得精打細算.
經常與不同的人做生意的猶太人.
上帝又問他們同一番說話.
如果我成為你們的上帝.
你們做我的百姓.

$^{81}$你們會為我做些甚麼?.
猶太人回答.
上帝,我們這個民族.
雖然不太精於藝術.
哲學或建築.
但我們是一個很懂得傳講故事的民族.
如果你成為我們的上帝.
我們做了你的百姓之後.
我們一定會將你的事情.
傳給我們的後代.
讓更多的人認識你.
明白你的心意.
當然上帝比猶太人更加精打細算.
於是就說「好」.
我們就一言為定.
弟兄姊妹.
我們看到上帝很期望子民.
用他的生命故事.
將上帝自己的話語傳頌出去.
讓更多的人認識上帝.
讓上帝的話語可以聽見.
看得到.
正如剛才我們所聽到的申命記一樣.
事實上剛才那段經文的背景是說.
摩西帶領以色列人到埃及.
經過了四十年漂流抗疫的生活.
現在他們又來到迦南前面.
未進入,但已經很近.
在這個時候摩西就重新.
跟百姓講述他們和上帝所納的約.
提醒他們唯有遵行上帝的律法.
單單敬拜事奉上帝.
才能在迦南蒙上帝的賜福.
包括經文所說.
可以活得長久.
可以經歷上帝的豐富供應.
並且他們的人數可以增加.
相反,如果他們不遵守上帝律法的吩咐.
他們就會失去上帝給他們這一切的應許.
不過回頭我們要想一想.

$^{121}$為什麼摩西要搶先跟百姓提醒呢?.
明明他們都未進入迦南.
原因是當百姓一踏入迦南地就會怎樣?.
有很多人教他們的.
有很多人會教他們怎樣過一個好的生活.
事實上在迦南地一個多神主義的地方.
充滿著不同的文化,信仰,風俗.
這一切都會搶著教他們.
怎樣可以在當地生活繁榮.
怎樣可以安居樂業.
譬如說如果你想你的農作物有好收成.
你就怎樣?.
你就拜巴力.
跟我們一樣敬拜巴力.
他會保佑你.
又或者說如果你想在迦南地生活.
你想得到更多的資源.
更大的影響力和保障.
最快的方法就是跟外邦人結盟.
甚至乎通婚.
你就很自然地可以壯大你本身的民族.
你看到在迦南地其實充斥著各式各樣.
聲稱可以幫助人改善生活.
提升生活質素的人間智慧.
不過摩西在這裡搶先提醒他們.
在他們未進入迦南之前就預先說明.
真正使人夢幻的智慧是什麼?.
是敬畏耶和華.
在迦南地這群百姓是否蒙福取決於他們是否持續敬畏上帝.
這個可以說是作為上帝指紋一個最基本的做人態度.
那怎樣才算是敬畏上帝呢?.
摩西這裡其實很清楚地說.
有一個很具體的方法.
就是謹守進行上帝的一切吩咐和教訓.
簡單地說你敬畏他就聽他的話.
他說什麼你就照做.
這就是敬畏最基本的實際操作.
讓上帝的話語成為生活的準則和指引.
使這群百姓可以持續在上帝面前保持聖潔的狀況.
遠離上帝所憎惡的事情.

$^{161}$注意這裡我們經常說保持聖潔.
其實他不是說上帝要求這群百姓脫離世俗.
去到一個一塵不染的地方來生活.
保持聖潔是說透過進行上帝的話語.
使這群百姓可以在其他民族的中間分別出來.
有分別.
擁有和遵守上帝的律法.
就能使以色列人保持在萬民之中的獨特身份.
意思是說舉個例子.
當他身邊的人.
當那些不認識上帝的人.
他們都認定要發達.
要生活好.
就拜偶像.
就要靠假神.
靠人的智慧.
人的手段.
這樣生活才會長久.
才會蒙福.
人口才會加增.
但是上帝的子民卻認定.
不是這樣的.
真正的生活智慧.
上帝說得很清楚.
是來自對他們的信靠.
這個民族之所以能夠存活.
蒙福.
壯大並且長久.
是因為他們敬畏上帝.
並且謹守進行上帝一切所吩咐的教訓.
這就是我們經常提到.
在聖經裡提到.
甚麼是分別為聖.
甚麼是與萬民有分別的實際表達.
以色列人持續通過上帝的話語.
其實他們是持續認順.
誰是他們的上帝.
持續遵循上帝的話語.
就是認順耶和華確實是我們的神.
我們也確實是他們的子民.

$^{201}$上帝也必然按他的應許.
持續保守賜福我們.
而相反.
百姓若不切實進行上帝的吩咐.
我們經常在聖經裡看到.
就是又敬拜上帝.
但同時又敬拜其他事物.
其他的假神偶像.
其實就等於變相將自己.
排除在上帝的應許裡.
不遵守上帝的律法.
其實變相放棄了.
自己屬上帝身份的獨特性.
事實上新約也是一樣.
我們經常說.
「認耶穌為主」.
其實實際操作是甚麼?.
不是純粹一個口號.
也不是單單說.
我們決志那一刻的事情.
「認耶穌為主」是說.
我們信主後.
我們持續要遵循耶穌的吩咐.
來持續確認.
耶穌你真是我們的主.
來持續確認.
我們真是屬於耶穌的門徒.
所以和以色列人一樣.
其實上帝也要求我們.
恆常遵守聖經真理的教導.
致力過一個敬畏上帝.
與萬民有分別的生活.
所以其實今天這段經文.
首先要檢視我們.
在現實生活裡.
你和我是如何看待上帝的話語.
我們是否確信.
好像詩篇所說.
上帝的話語是我們腳前的燈.
路上的光.

$^{241}$是真是確實使人蒙福的智慧.
還是我們總是欠缺一些.
追求聖經真理的胃口.
屬靈上我們其實實際上.
很多時候都是抹黑過日子.
這是經文首先要我們思考的.
曾經有人這樣半開玩笑形容.
他說現代人生活是忙碌的.
我想我們也很認同.
工作時間也很長.
所以今時今日肯固定參與崇拜.
肯聽牧者講解聖經真理.
這已經是標準的基督徒.
假如他還願意每一天來靈修.
無論看聖經.
或用軟件來讀經靈修的話.
可以說是優質的基督徒.
至於那些肯定時花時間.
研讀,查考聖經.
甚至思考如何實踐上帝的教導.
這類基督徒可以說是稀有的品種.
無論怎樣.
確實現代人生活繁忙.
但我們明白無論多忙都會吃飯.
假如我們真的視聖經真理.
為生命所必需的元素.
就算多忙也會盡力看聖經.
特別當耶穌說人活著.
是靠上帝口裡所說的一切話的時候.
其實是提醒我們每一個.
作為新約的信徒.
我們需要用聖經真理來作為生命糧.
到底我們對上帝的話語的認識程度.
投入程度又是怎樣呢?.
我想我們每一個都擁有聖經至少一本.
但我們有沒有時常透過聖經來親近上帝呢?.
事實上我們很明白.
追求聖經真理.
往往不是我們有沒有時間的問題.
而是我們願不願意抽時間的問題.

$^{281}$我想我們都很認同.
周禎如.
當你拍拖戀愛的時候.
你認真去愛一個人.
你很想追求他和他發展關係的時候.
無論你多忙都會怎樣呢?.
你不會不約他吧?.
你不會不陪他吧?.
你真的愛一個人.
你必定會花時間.
去認識他,了解他.
對不對?.
同樣.
其實我們是否願意花時間.
追求聖經真理.
研讀上帝的話語.
實際是測試我們和上帝的關係.
我們和上帝的距離.
某程度上從我們的讀經生活反映出來.
我們是否真的敬畏他.
全心全意來愛他呢?.
我想請大家一起讀.
第四至五節,好嗎?.
程序表都印了.
預備起.
以色列人,你要聽.
耶和華我們臣是獨一的主.
你要盡心盡聖盡力.
愛耶和華你的神.
好,謝謝.
我們很熟悉這段經文.
特別是以色列人.
世代相傳一個信仰的總綱.
用兩字的經文濃縮了.
他們和上帝的關係.
這裡說「聽」.
實際是在說「聽從」.
聽從上帝的吩咐.
透過這個行動來回應.
上帝對他們的揀選和呼召.

$^{321}$上帝要求他們.
我們很熟悉「盡心盡聖盡力」.
簡單說就是你整個人百分之一百.
是屬於上帝的.
你也需要百分之一百.
全心全意愛你的上帝.
因為上帝是我們獨一的主.
我們不妨探討一下.
獨一的主其實是什麼意思呢?.
換句話來說.
獨一的主其實是說.
世上只有上帝是真神.
其他的都是虛假的偶像.
這是獨一的意思.
也就是說唯有上帝掌管一切.
包括他會很明白.
作為人真正需要的是什麼.
作為人什麼才對他有益.
唯有上帝才能一清二楚.
所以上帝是我們唯一.
要效忠要服侍的主人.
也就是說我們只需要.
愛上帝服侍上帝一個就夠了.
不需要因為懷疑擔心.
上帝的安排,供應,保護.
會不會不夠呢?.
於是乎就怎樣呢?.
我多靠幾個.
我多拜幾個.
我多信幾套世俗的哲學.
獨一的意思就是這樣.
換句話來說.
盡心盡性盡力.
這裡是說上帝的子民不能夠因為.
我想賺多點利益.
賺盡點人生保障.
我就怎樣呢?.
很多時候最大的問題.
不是我們完全否定上帝.
不愛上帝.

$^{361}$是我們會左右逢源.
正如舊約上上強調.
「又愛上帝又敬拜別神」.
或者新約裡面說.
我們又愛耶穌.
不過又愛世界.
事實上弟兄姊妹.
在一個變幻莫測的世界裡.
上帝提醒我們.
無論這個世界變成怎樣.
無論你自身的遭遇如何.
你仍然要相信.
這位獨一的真神掌管一切.
你仍要相信.
信服祂的安排.
信服祂的吩咐.
不要讓環境.
不要讓你的際遇.
甚至你的心情.
影響我們對上帝的專注和忠心.
確實盡心盡性盡力.
愛上帝追求聖經真理.
並不一定會感覺良好.
不一定會很浪漫.
因為聖經很多時候.
要拆悔改變我們的生命.
包括.
有時要我們認罪悔改.
要我們去饒恕.
曾經得罪我們的人.
又或者要將你一直抓緊.
你覺得很安全的事情.
你要學習放手.
所以上帝的話語和提醒.
很多時候會令你坐立不安.
如果沒有了一個愛上帝的心.
其實我們真的很難堅持.
去看聖經.
去堅持實踐上帝的吩咐.
事實上弟兄姊妹.

$^{401}$愛上帝不是一個感覺.
不是一種口號.
而是背後說.
為上帝犧牲付代價.
無論生活是順還是逆.
甚或可能因為遵行上帝的吩咐.
而令自己有吃虧虧損的機會.
但當我們仍然願意在這樣的情況.
堅持去追求真理.
實踐真理.
會出聖經的教導.
我們說這就是我們正在愛上帝.
我們和上帝的關係最緊密.
一個最務實的表現.
所以為了這個緣故.
每一個基督徒基本上都要下定決心.
親近上帝的話語.
事實上如果我們連上帝的話語都不認識.
我們又如何談得上.
如何切實進行.
最基本是看上帝的話語.
而不是很多時候.
我想有時現況確實是.
我們可能寧願有時間選擇看劇集.
看迪士尼,Netflix.
都不看聖經.
又或者我們會落入一個狀況.
我間中有事的時候.
我有需要的時候.
我才會揭聖經來看.
話說有一個平時不太看聖經的基督徒.
他很想轉工尋找一些人生的方向.
於是他想透過聖經求問上帝的心意.
他禱告完之後就開始翻開聖經.
結果他一揭就揭到路加福音裡面.
耶穌對他說.
「變賣你的一切,分給窮人,.
就必有財寶在天上,.
你還要來跟從我?」.
於是他就想.

$^{441}$不可能的.
上帝怎可能要我賣掉所有家當.
甚至去讀神學.
做一個傳道人呢?.
他不相信上帝會這樣說.
於是他又祈禱.
又再揭聖經.
這次他揭到的不是路加福音.
他揭到楊福音.
楊福音這裡說.
「你門若遵行我所吩咐的.
就是我的朋友了」.
他就想.
糟了.
難道耶穌真的要我做傳道人?.
否則他會不友好我?.
不會吧.
他不太相信.
於是他又再祈禱.
再揭聖經.
這次他揭到馬太福音.
他一揭就揭到關於猶大那一段.
這裡就說.
猶大把銀錢吊在電裡.
就出去吊死了.
糟了,他就害怕了.
他就怕了.
沒錯,我真的不想賣光我的家當.
我真的有些亂賺我現在所擁有的東西.
但是上帝也沒理由要我自殺.
他不相信.
他又祈禱.
又揭.
這次又揭到路加福音.
你猜耶穌對他說什麼?.
「你照樣去行」.
弟子們,雖然這故事純屬搞笑虛構.
不過它真的反映了一些現實.
確實反映了一種.
用消費,用功利主義模式.

$^{481}$來讀聖經的心態.
即是平時無事無幹.
就會將上帝的話語放在一邊.
忽然間,我生活很有需要.
我想找到一些指引幫助和帶領的時候.
我就會看聖經.
希望找到幫助.
確實,用這樣的方式.
是有看到聖經的.
甚至乎真的找到一些他覺得適合的資訊.
不過,他沒有愛上帝.
有看到聖經,但是沒有愛上帝.
確實,我們會否也落入這種境況呢?.
曾經,當初信主回教會的時候.
在團契裡,有一群團友一起.
互相聊天.
問大家有沒有看完一本聖經?.
於是不同的團友就表達.
我看了新約.
另一個就說,我還看了詩篇.
有些說,我出到埃及.
最後也去到迦南地.
看完五經和約書亞記.
最後,有一個團友就分享.
我也有看完整本聖經一次.
但是他的臉色有點尷尬.
於是有人問.
這樣不就好了,看完整本聖經一次.
何必那麼難為情呢?.
他就說.
我看完的那本是兒童版聖經.
即是很多圖畫.
沒有太多文字.
大人的那本那麼厚.
怎麼看得完呢?.
弟兄姊妹,會不會.
這也是我們的情況呢?.
整本聖經那麼厚.
我剛才也數過.
1680頁新約加舊約.

$^{521}$以我們的和合本來計算.
可能我們一輩子都沒有讀過一本.
超過一千頁的書.
更惶恐的是聖經呢?.
如果沒有了愛上帝這個心.
作為你的推動力.
我們確實很難細水長流.
每天都堅持點點滴滴.
去研讀上帝的話語.
通常我們的情況.
真的好像以色列人一樣.
我們出到埃及.
不過是怎樣呢?.
通常停留在曠野.
就過不了民數機.
所以弟兄姊妹.
其實讀聖經不能夠沒有了愛.
正如我們愛上帝.
不是出於一時的衝動.
不是一刻的需要.
同樣愛沒有上帝的話語.
也不能夠忽冷忽熱.
你需要用愛心.
細水長流地聽從.
遵守聖經的教導.
就好像經文第六節所說.
設法將上帝的話語.
牢記在心裡.
所謂將上帝的話語.
牢記在心裡.
其實不難理解.
我們有類似的經驗.
你試試回想讀書的時代.
你是怎樣將課本.
上課的知識和理論.
牢記在心裡.
記得嗎?.
首先你要上學.
你要上課.
你亦要自己看書.

$^{561}$你要做練習,做功課.
你亦要定期溫習.
如果容許的話.
你要用時間思考那些知識.
那些理論沉澱一下.
最好就和同學一起分享討論.
更加重要的是.
你要經過大大小小的測驗考試.
透過各種不同的反覆操練.
不知不覺將你從小學.
到中學甚至大學.
一系列的知識和理論.
印在你的心裡.
對不對?.
其實追求聖經真理是一樣的.
你需要看,你需要讀.
你需要去查考,需要去沉澱.
需要去實踐.
通過你生活大大小小的考驗.
來去驗證上帝所說的是真的.
上帝所說的能夠令你蒙福.
唯有這樣.
上帝的話語才可以.
踏實在我們的生命裡.
如果純粹是聽.
很快就會流失.
事實上頂梓妹.
今日我們說追求聖經的指引.
就代表我們需要在日常生活裡.
下苦功來操練.
正如第七節所說.
無論你坐在家裡.
行在路上,躺下,起來.
這裡說的是一個人整天的活動.
上帝提醒我們.
在整天的活動裡.
在你生活不同的範疇裡.
你要讓上帝的話語.
滲透到每一個場景.
這裡特別提到.

$^{601}$父母有責任恩勤教導他的兒女.
意思是用上帝的話語.
來磨練子女的生命.
當然過程是說.
父母和子女.
一起用上帝的話語.
磨練彼此的生命.
讓我們看到.
家庭教育其實是信仰裡.
很重要的一環.
原因是教會往往.
由很多家庭所組成.
家庭成員越是追求聖經真理.
實踐聖經真理.
這間教會就會怎樣?.
自然會越來越健康.
事實上有人曾經說過一句.
很睿智的說話.
他說經常在一起讀經.
祈禱的家庭不會破碎.
提醒我們.
到底今天我們送給子女.
最好的禮物是甚麼呢?.
並不是金錢.
並不是物質.
而是我們真的有沒有.
花時間用時間.
陪他們一起追求.
實踐上帝的話語.
事實上這是上帝.
給每一個基督徒家庭.
一個的福分.
一種樂趣.
當然可能有父母會覺得.
牧者或其他某些弟兄姊妹.
他們比我們更加.
擁有豐富的聖經知識.
和數年經驗.
他們一定會比我們教得好.
所以我們還是留給他們.

$^{641}$教我們的小朋友會好一點.
請你謹記.
教會永遠不能取締.
父母的教導角色.
事實上你會很明白.
小朋友回教會的時間.
遠不及他在家的時間.
他們吸收到多少.
聖經真理的教導.
就有賴平日爸爸媽媽.
在他們面前是怎樣.
追求和實踐真理.
雖然我們確實不知道.
小朋友或教會的下一代.
他們的屬靈生命.
何時會開花結果.
但作為他們的父母.
作為他們的屬靈長輩.
我們確實有必要身體力行.
第八節說.
「繫在手上為記號.
戴在額上為經文」.
當然它實際的應用.
不是將經文鎖在手上.
或戴在額上.
它是說.
展現如何用聖經真理.
來管理你的手.
即你的行為.
和你的頭.
即你的心思意念.
我們如何在下一代面前.
身體力行地帶頭展現.
我們真是用聖經真理.
管理我們的心思行為.
這很值得我們思考和操練.
事實上我們不難想像一幅圖畫.
假如我們讓下一代.
讓我們的子女上主一學認識真理.
假如我們沒有身體力行.

$^{681}$其實我們在給一個甚麼訊息出來.
你覺得小朋友會接收到甚麼訊息.
其實很混亂.
你試試代入他們的世界.
你會想到.
為何大人常常叫我上主一學.
看聖經,學習聖經真理.
為何他們不是這樣?.
為何我看不到他們這樣?.
到底發生了甚麼事?.
小朋友會將這個疑惑.
藏在心裡.
等到他們可能踏入中學階段.
開始有機會自主的時候.
就可能會怎樣?.
就有機會流失.
不再上主一學.
不再看聖經.
為甚麼?.
原因是.
他們身邊的大人不斷用行動.
告訴他們.
信耶穌不用那麼認真的.
信耶穌不一定要看聖經.
上主一學.
所以弟兄姊妹.
其實這段經文今天提醒我們.
為了上帝的教會.
為了我們屬靈下一代的緣故.
我們需要締造一個良好的.
屬靈成長環境給他們.
我們一定要千方百計.
向他們展示一個.
很重要的訊息.
無論你對聖經的知識有多少.
我們都需要展示給他們知道.
一起認識聖經的真理.
一起分享實踐上帝的話語.
這是一件很無比的事情.
這是一件很棒的事情.

$^{721}$有益的事情.
不要將上帝的話語.
藏在一旁.
就像第九節所說的.
最後一段經文.
又要寫在你房屋的門框上.
並你的城門上.
這裡提醒我們.
要讓上帝的真理.
在我們不同的生活場景裡.
彰顯出來.
讓你身邊的人看見.
城門這裡特別是說.
在古代世界裡.
社交最活躍的地方.
有很多買賣.
有很多審判.
社交活動.
當時都是在城門口舉行.
它提醒我們.
過往我們可能會想.
信仰歸信仰.
生活歸生活.
我們會將信仰.
和我們家庭甚至社交.
切割出來.
但是上帝這裡重新提醒我們.
不是這樣的.
而是要將祂的話語.
徹底融入你的生活.
你的生計裡.
因為唯有這樣.
我們才能夠在不同的生活場景.
全方位地體驗到.
上帝給我們的祝福.
全方位地體會到.
上帝給我們的帶領和保守.
唯有這樣.
我們才真是.
好像今天的經文所提醒我們.

$^{761}$可以過一個敬畏上帝.
愛上帝的生活.
盼望今天這段經文.
能夠重新喚醒.
我們對上帝話語的熱誠.
我們同心祈禱.
上帝我們衷心的多謝你.
你不單單只揀選拯救我們.
你更加刺客你寶貴的話語.
來成為我們生命的提醒和指引.
上帝我們知道你對我們的期望.
你願意我們真是透過你的話語.
來過一個敬畏你.
分別為聖的生活.
以致我們真是能夠.
在萬民當中有分別.
並且可以吸引他們歸向上帝.
讓他們明白真正人生的智慧.
來自對上帝你的信靠和敬畏.
來自怎樣去活出.
上帝引領你在聖經裡的真理教導.
上帝我們會有軟弱.
我們會有試探.
在忙碌的生活裡.
求你幫助我們.
讓我們欣著你對我們的愛.
我們都願意以愛來回應.
以愛來回應怎樣研讀你的話語.
讓我們在生活裡.
重新重視聖經裡的教導.
我們不單單研讀.
我們也切實去進行.
我們不單單造就自己.
也致力去造就.
你所託付給我們身邊的每一個.
讓教會擁有更多.
願意進行你話語的門徒.
讓教會擁有更多.
願意實踐你話語的家庭.
讓教會成為一間.

$^{801}$願意忠心謹守.
進行你話語的群體.
為此我們同心獻上祈禱.
奉耶穌基督.
你得勝的明天來祈求.
阿們.
\newpage



\section{哥林多後書 13:1-14}
\label{sec:9NLZUx6GLHU}
\textbf{2022年1月23日崇拜直播｜梁家麟牧師｜成全聖徒的召命｜哥林多後書13\_1-14}
\newline
\newline
連結: \href{https://youtube.com/watch?v=9NLZUx6GLHU}{\texttt{https://youtube.com/watch?v=9NLZUx6GLHU}} ~~~~ 語音日期: 2022-01-24
\newline
\newline
\hyperref[sec:Qke9jW9Cx6E]{\small{< < < PREV SERMON < < <}}
~
\hyperref[sec:index_chronic]{\small{[返順時目]}}
~
\hyperref[sec:tgXI8MUm32I]{\small{> > > NEXT SERMON > > >}}
\newline
\newline
哥林多後書 13:1-14
\newline
\begin{longtable}{cl}
\hline
\hline
章節 & 經文 (和合本修訂版)\\
\hline
13:1 & \begin{tabularx}{0.7\textwidth}{X} 這是我第三次要到你們那裡去。「任何指控都要憑兩個或三個證人的口述才能成立。」 \end{tabularx} \\ \\ \relax
13:2 & \begin{tabularx}{0.7\textwidth}{X} 對那些犯了罪的人和其餘所有的人，正如我第二次見你們的時候曾說過，現在不在你們那裡再次說：「我若再來，必不寬容。」 \end{tabularx} \\ \\ \relax
13:3 & \begin{tabularx}{0.7\textwidth}{X} 因為你們想求證基督是否藉著我說話。基督對你們並不是軟弱的，而是在你們裡面大有能力的。 \end{tabularx} \\ \\ \relax
13:4 & \begin{tabularx}{0.7\textwidth}{X} 他因軟弱被釘在十字架上，卻因神的大能仍然活著。我們在他裡面也成為軟弱的，但對你們，我們將因神的大能而與他一同活著。 \end{tabularx} \\ \\ \relax
13:5 & \begin{tabularx}{0.7\textwidth}{X} 你們總要省察自己是否在信仰中生活；你們要考驗自己。除非你們經不起考驗，你們自己豈不應該知道有耶穌基督在你們裡面嗎？ \end{tabularx} \\ \\ \relax
13:6 & \begin{tabularx}{0.7\textwidth}{X} 我希望你們知道，我們並不是經不起考驗的人。 \end{tabularx} \\ \\ \relax
13:7 & \begin{tabularx}{0.7\textwidth}{X} 我們祈求神使你們不做任何惡事；這不是要顯明我們是經得起考驗的，而是要你們行事端正，即使我們似乎經不起考驗也沒有關係。 \end{tabularx} \\ \\ \relax
13:8 & \begin{tabularx}{0.7\textwidth}{X} 我們不能做任何對抗真理的事，只能維護真理。 \end{tabularx} \\ \\ \relax
13:9 & \begin{tabularx}{0.7\textwidth}{X} 當我們軟弱而你們剛強時，我們也歡喜。我們所祈求的是：你們能成為完全人。 \end{tabularx} \\ \\ \relax
13:10 & \begin{tabularx}{0.7\textwidth}{X} 所以，我不在你們那裡的時候，把這些話寫給你們，好使我見你們的時候不用照主所給我的權柄嚴厲地待你們；這權柄原是為造就人，而不是為摧毀人。 \end{tabularx} \\ \\ \relax
13:11 & \begin{tabularx}{0.7\textwidth}{X} 末了，弟兄們，願你們喜樂。要追求完全；要接受鼓勵；要同心合意；要彼此和睦。如此，慈愛和平的神必與你們同在。 \end{tabularx} \\ \\ \relax
13:12 & \begin{tabularx}{0.7\textwidth}{X} 你們要用聖潔的吻彼此問安。 \end{tabularx} \\ \\ \relax
13:13 & \begin{tabularx}{0.7\textwidth}{X} 眾聖徒都向你們問安。 \end{tabularx} \\ \\ \relax
13:14 & \begin{tabularx}{0.7\textwidth}{X} 願主耶穌基督的恩惠、神的慈愛、聖靈的感動常與你們眾人同在！ \end{tabularx} \\ \\
[1ex]
\hline
\hline
\end{longtable}
$^{1}$今日重讀的經文是在.
哥林多後書十三章一至十四節.
是在新約聖經二百五十六頁.
請聽我讀出十三章.
「這是我第三次要到你們那裡去.
憑兩三個人的口作見證.
句句都要定準.
我從前說過如今不在你們那裡遊說.
正如我第二次見你們的時候所說的一樣.
就是對那犯了罪的和其餘的人說.
我若再來必不寬容.
你們竟然尋求基督在我裡面說話的憑據.
我必不寬容.
因為基督在你們身上不是軟弱的.
在你們裡面是有大能的.
他因軟弱被釘在十字架上.
卻因神的大能仍然活著.
我們也是這樣同他軟弱.
但因神向你們所顯的大能.
也必與他同活.
你們總要自己審查.
有信心沒有.
有要自己試驗.
豈不知你們若不是可棄絕的.
就有耶穌基督在你們心裡嗎?.
我卻盼望你們曉得.
我們不是可棄絕的人.
我們求神叫你們一見我是都不作.
這不是要顯明我們是蒙越立的.
是要你們行事端正.
任憑人看見我們是棄絕的罷了.
我們凡事不能抵擋真理.
只能輔助真理.
即使我軟弱.
你們剛強.
我們也歡喜.
並且我們所求的.
就是你們作完全人.
所以我不在你們那裡的時候.
把這話寫給你們.

$^{41}$好叫我見你們的時候.
不用照主所給.
我的權柄嚴厲地待你們.
這權柄原是為造就人.
並不是為敗壞人.
還有沒了的話.
願弟兄們都喜樂.
要作完全人.
要受安慰.
要同心合意.
要彼此和睦.
如此人愛和平的神.
必常與你們同在.
你們親嘴問安.
彼此無要聖潔.
眾聖徒都問你們安.
願主耶穌基督的恩惠.
神的慈愛.
聖靈的感動.
常與你們眾人同在.
聖經獨讀.
一年半的哥林多後書的歷程.
今天是最後一次了.
這一講是接續上一講.
十三章的一到四節.
上次是講了大半的了.
今天提多一點.
保羅宣告他們三次的哥林多之旅.
會是一個不再寬容的行程.
不寬容不是因為他個人心胸狹窄.
而是來自基督的憤世.
保羅說你們既然尋求.
基督在我裡面說話的憑據.
我必不寬容.
因為基督在你們身上不是軟弱的.
在你們裡面是有大能的.
作為基督耶穌的使徒.
行事為人.
自當遵照耶穌基督的教訓和榜樣.
保羅問.

$^{81}$哥林多信徒想知道.
耶穌跟他說了什麼.
耶穌要求他堅持真理.
抵制罪惡.
嚴懲犯罪.
而拒絕悔改的人.
你們既然尋求基督在我裡面說話的憑據.
有人說.
哥林多信徒和一些人.
是在質疑保羅的使徒權威.
問你有什麼憑據證明你是使徒.
保羅現在就說.
他會作為使徒.
行使使徒的權柄.
處分那些犯罪的人.
保羅從來不介意在信徒面前示弱.
並且以自己的軟弱誇口.
不過哥林多不表示.
他所信奉的耶穌基督都是弱者.
弱者保羅跟隨弱者耶穌.
不是.
剛剛好相反.
保羅的示弱是為了要彰顯基督的能力.
基督的大能要在他這個軟弱的人身上.
充分顯現出來.
所以信徒不要誤會基督是弱者.
永遠是好好先生.
以輔導員的面貌出現.
沒有脾氣.
沒有立場.
也不執著.
我們做什麼他都無所謂.
我們怎麼犯錯.
他都願意執手尾.
包底.
最多說不要緊.
下次不要了.
不是.
因為基督在你們身上不是軟弱的.
在你們裡面是有大能的.

$^{121}$基督在我們裡面是有大能的.
我覺得這句話真的很重要.
弟兄姊妹.
基督在我們身上佔了一個什麼位置.
他只是我的參謀.
我的軍師.
幫我出主意.
但是留待我對所有事情做最終的決定.
他是我的輔導員.
只是想了解我.
同情我.
肯定我.
明白我的情慾.
我的犯罪.
是多麼正義和神聖.
並且為我不斷犯罪去提供藉口.
耶穌基督的說話在我生命裡面有多大的權威性.
宰制性.
不管我認同不認同.
理解不理解.
我都要堅決執行.
抑或.
那個是參考而已.
跟不跟由得我.
簡單說一句.
耶穌基督是不是真是你的主.
在你生命裡面有沒有安放祂的寶藏.
我們說基督在人裡面的大能的時候.
我們常常想到的就是基督賜大能給我們.
例如我們有行神即騎士的能力.
使人歸信祂的能力.
我說一句你們立刻全部低頭來信主.
你說多過癮.
基督有大能在我身上.
但是這裡說的是.
基督在我身上所彰顯的大能.
就是說我是基督大能的那個動的對象.
不是藉著我發功來將大能輸出去.
是不是.
我們首先自己要在基督的大能面前俯首稱臣.

$^{161}$全然信服.
一個期盼在事公上彰顯基督大能的人.
必須要在生命上信服基督的大能.
我們在宣講基督是主的福音之前.
我們一定要活出這個基督真是我的老闆.
基督真是我生命的主的生命見證.
你沒有這個見證你說什麼信息.
這裡我還要再注意的就是.
因為基督在你們身上不是軟弱的.
在你們裡面是有大能的.
這裡用的是眾數的你們.
一方面要說基督是主宰.
這個適用於所有信徒.
另一方面也凸顯.
這句話是針對信徒群體.
就是教會你們是指教會說的.
基督是教會的元首.
教會必須要在所有層面的活動和事公上.
彰顯基督的主權和能力.
在進行真理,剷除罪惡的任務上切實執行.
這裡不只是說個人.
也在說教會.
在說我們旺朕.
基督在旺朕身上是有大能的.
說到基督的能力.
保羅指出基督同時存在軟弱和剛強的弔詭性.
第四節.
他軟弱被釘在十字架上.
他因上帝的大能仍然活著.
我們也這樣和他軟弱.
但因為上帝向你們所顯的大能.
也別與他們同活.
耶穌基督被釘死在十字架上.
似乎凸顯他軟弱無能的一面.
但他從死裡復活.
彰顯了勝過死亡的上帝大能.
基督軟弱的外貌背後.
是上帝的大能.
當我們和耶穌基督同死的時候.
也陪伴他們在軟弱裡.

$^{201}$我們和他們一起復活.
彰顯了上帝的大能.
同樣在我們身上彰顯.
Power and Weakness.
我們是軟弱的人.
我們教會沒有什麼能力.
我們執行什麼紀律.
人們不受我們是沒有辦法的.
他要走掉我們是沒有能力的.
我們是沒有能力的.
不過我們自己認定.
我們所宣講的福音真理.
我們所高舉的聖經.
是有永恆的再世能力.
我們自己好像沒有.
是人們不聽我們說話.
但我們知道我們所說的.
是有終極的能力.
能力不是來自我們.
我和你省一點.
是來自上帝.
我們來到第二段.
個人的罪和教會的問題.
這段我們很快地過一過.
很多點很快地過一過就行了.
前面保羅提到個人的教會.
或者存在著八宗分裂的罪.
和三宗懲獄的罪.
我們不是很能夠確定.
是不是真的教會有這樣的問題.
又或者是不是同一班人犯這麼多的罪.
有多少信徒捲入其中.
我們都不是很知道.
如果我們對教會沒有絕望.
大概我們不會假設.
我們沒有辦法想像.
整間教會全體會眾.
都陷在這樣的罪裡面.
如果這樣的話.
整間教會都不用恨了.

$^{241}$不用恨了.
我們不會這樣想的.
不過對保羅來說.
或者對我們日後的基督徒來說.
如果那個罪是長期在教會裡面盤踞.
或者有些人長期在教會裡面犯罪.
拒絕認罪拒絕悔改.
仍然可以安然無事.
這樣.
先說過那麼多教會.
一定是存在著一些結構性的問題.
對不對.
所以重點不是在於有人犯罪.
千萬不要詫異扮天真.
有基督徒犯罪.
有基督徒跌倒.
你傻的.
你照鏡子.
我們都會犯罪.
我們每個人都會跌倒.
那個不是問題.
我不是說犯罪不是問題.
總之我不會詫異這個情況出現.
這個不是很奇怪的情況.
不過如果有人.
例如在我們中間的鮑爾拉.
她可以每個星期帶無私藥回來.
我們大家就當她沒有存在著這種情況.
人家說不知道就沒有問題.
她知道都不處理.
那個就不是那個人的問題那麼簡單.
而是黃浸的問題.
你明不明白.
一定是黃浸的問題.
如果我們教會對人犯罪.
長期犯罪是視若無睹.
那個才是問題所在.
真正的關鍵.
所以保羅在這裡說那麼多教會的情況.
他關心的不僅僅是個別人犯罪的問題.

$^{281}$而是教會如何處理犯罪的問題.
正如他在比哥林多的第二封信.
哥林多前書也說過這件事.
在哥林多前書的五章一到五節.
說風雲在你中間.
有淫亂的事.
有人收他繼母.
但為什麼你們竟然可以不處理.
為什麼你們還可以驕傲.
自高自大不處理這件事.
這個才是保羅最震動.
最覺得需要大刀闊斧去斬下去的事.
我們知道不久之前.
有一個很知名的音樂家.
因為嫖娼而被中國的公安行政拘留.
中國音樂家協會和很多團體.
都立即褫奪了他各種的名銜和職務.
我們知道這樣的事件.
是有政治操控的成份.
不過我們不會因為政治操控.
而抹殺這件事的正面信息.
公眾人物要承擔社會道德的責任.
這不只是中國如是.
是中外皆然.
美國紐約州的前州長科莫.
曾經因為性醜聞在2021年8月辭職.
你可以說有情婦包二奶.
在華爾街誇張一點.
五個就有三個.
有什麼奇怪的.
為什麼這麼大驚小怪.
不行嗎.
但你是政治人物很抱歉.
因為你是政治人物.
就被人抓了.
公眾人物是有特別道德責任的.
我們教會作為信仰的事.
凡單位也是這樣.
我不是說凡事要嚴懲.
教會是講恩典的地方.

$^{321}$教會是挽回的地方.
我們所有紀律的處分都是為了想挽回人.
所以不是要一棍打死.
不過總的來說.
我們艱難極.
我們都不能迴避去處理.
令到我們軟弱跌倒的情況.
帶著眼淚.
帶著很多的痛苦困惑.
但我們怎樣也要正視.
背後反映的是.
我們對上帝的敬畏.
對真理的敬畏.
是壓倒性大於一切人際關係.
信仰測試第三點.
保羅挑戰那些信徒.
說他們正在接受一個考試.
要考驗他們是否真的信循真理.
是否真的有信仰在他們內.
13:5-7.
你們總要自己審查.
有信心沒有.
要自己試驗.
豈不知你們若不是可棄絕的.
就是耶穌基督在你們心裡嗎?.
我卻盼望你們曉得.
我們不是可棄絕的人.
我們求上帝.
叫你們一件惡事都不作.
這不是要顯明我們是蒙越立的.
而是要你們行事端正.
任憑人看我們是被棄絕的罷.
你們總要自己審查.
有信心沒有.
這句說話很有趣.
保羅要求那些信徒.
攻心自省.
檢視是否真的有信仰.
有信心沒有.
原文是在信仰中.

$^{361}$信心不是一件你可以擁有的東西.
也不是你可以計算的數量.
而是我們身處的境況.
生命的狀態.
我們在信仰中.
我們是信者.
I am believing.
I am a believer.
我是一個信仰者.
有耶穌基督在我們心裡.
在信仰中.
我亦是在基督中.
我在基督裡面.
基督在我裡面.
這是耶穌基督的應許.
我在信仰中.
就是我在基督裡面.
這不僅僅是一個感覺.
有基督充滿同在的感覺.
而是基督的大能和主權.
要在我們裡面彰顯出來.
也要自己試驗.
保羅是常常說試驗.
你記得他常常說過.
有火的試驗.
去考驗我們每一個人的生命工程.
在哥林多前書.
三章十二到十五節的說法.
在哥林多後書的二章九節.
他寫這封信是要試驗你們.
看你們是否凡事順從.
後面說到為耶路撒冷捐書.
他也要說.
這也是一個考驗.
看我們是否能夠.
顯出我們對人的愛心.
保羅其實想說的是.
我們生活裡面的大小際遇.
其實都是我們的信仰考驗.
也就是一個試金石.

$^{401}$檢查我們是否真的有信仰.
我們是否真的在信仰中.
我們是否真的在基督裡面.
所以永遠的就是在我們.
遭遇各種開心不開心的事情.
我們遭遇各種喜怒哀樂.
高低起伏.
香港我們不停在驚嚇中.
一個大的挫折來.
小的挫折又來.
好像一波未平一波又起.
我們是否在信仰中.
常常真的能夠顯出來.
遭遇危機.
遭遇轉變.
遭遇挑戰的時候.
我們笑不笑得出.
我們夠不夠淡定.
我們怎樣處理.
我們怎樣回應.
我們會否影響和身邊的人的關係.
我們的煩躁.
我們的焦慮.
我們對人的無耐心.
那些全部都反映著.
我們的生命質素.
我們的信仰質素.
我們是否真的在信.
是否真的在信.
基督徒要將我們生活中的大小遭遇.
轉化成為一個信仰的實踐.
從負面的角度.
是一個信仰的考驗.
是的.
教會每個星期都很辛苦.
上星期三分一.
這星期就沒有了.
不可以聚會.
很多事情都安排好.
又要推倒它.

$^{441}$是否很煩惱.
不過就是這樣.
我們是否在信仰中.
我們視一切無法由我們外力控制的東西.
都是上帝為我們鋪設的場景.
要顯示我們生命的質素.
信仰的質素.
所以弟兄姊妹.
淡定些來.
我們現在在做一個考試.
你們總要自己審查.
有信心沒有.
這是一個很好的問題.
我們要問自己是否在信仰中.
我們是否真實的基督徒.
真實的信者.
十三章六到九節.
上半節.
我們來到第四點.
剛強和軟弱.
這點我以前講了很多.
我腳攀望你們不曉得.
我們不是可棄絕的人.
我們求上帝.
叫你們一件惡事都不作.
不是要顯明我們是蒙越納.
而是要你們行事端正.
任憑人看我們是被棄絕的.
我們凡事不能抵擋真理.
只能輔助真理.
即使我們軟弱.
你們剛強.
我們也歡喜.
可棄絕.
蒙越納.
這裡講了兩個字.
如果跟前面接著來講.
那個試驗來講的話.
可棄絕即是通過測試.
即是在輸送帶那裡.

$^{481}$被甩出來的次貨.
那些叫可棄絕.
即是蒙越納即是通過測試.
即是那些牛丸有牛肉.
那些就是蒙越納.
牛丸沒有牛肉的那些.
就是可棄絕.
總之是通不通得過測試.
保羅盼望那些信徒.
一方面確定他的使徒身份.
是經得起測試的.
他不是被上帝棄絕的.
不過就算他們仍然看不起保羅.
仍然對保羅的身份和事奉不信任.
保羅說也不要緊.
我知道自己在上帝面前是被越納的.
不是被棄絕的.
保羅與同時又盼望.
並不僅是盼望.
他積極向上帝祈禱.
希望那些信徒都是蒙越納.
而不是被棄絕的.
他們不再犯罪.
不再做惡事.
在這裡保羅重申.
不要當那些信徒是對手.
只能贏或輸.
保羅說他不是想證明.
我對你們錯.
我希望我們大家都蒙越納.
我不需要在你們面前顯現.
我是剛強那一面.
保羅再說.
他寧願以弱者的身份出現.
他不想執行使徒的權柄.
不需要縮小肌肉證明身份.
他不想.
他寧願繼續以弱者的面貌在中間存在.
如果你們不需要我扮強者.
我就不扮.

$^{521}$很多教會都設立人事及紀律小組.
處理一旦出現人事糾紛.
信徒犯罪跌倒後的跟進事宜.
假如你被選入我們教會的小組.
你會否介意一年沒開過一次會.
都沒事要處理.
我猜不介意.
年終交報告.
交給執事會 交給會友大會.
我們今年的事工是零蛋.
沒事.
大家不會質疑你們為何不做事.
能夠不做就好了.
每個部門都要做事.
唯獨這個部門最好不做事.
紀律不用.
我們教會沒有紀律問題.
所以這個小組是形同虛設.
這是最好的事.
我希望在你們中間演我這個仕途的肌肉.
其實如果我需要在你們中間.
動用我這個仕途的權柄.
我這個仕途是失敗的.
不是成功的.
正如一個小學老師.
一個中學老師.
是否需要蹲起教師款.
要威迫你威嚇學生.
然後才讓學生聽話.
如果不用 大家嘻嘻哈哈.
這樣就已經很和諧.
很順利完成教學任務.
是不是最好.
不用掩耳朵.
我越不需要蹲起教師款.
我這個教師是越成功的.
我越需要扎實話.
我這個教師就是失敗的.
保羅也說.
我不需要彰顯我這個仕途的權柄.

$^{561}$我們凡事不敵黨真理.
只能輔助真理.
這句話好像是一個公理.
一個exam(論證).
不用論證.
所有人包括保羅在內.
都只能夠依據真理行事.
敵黨真理就是指你所做的事.
和真理相衝突.
輔助真理就是說.
你做的事是validate(證明).
印證了真理.
我們所做的事.
都正在證明真理是真理.
好.
接著我們來到.
作完全人這一點.
保羅陳述他對哥林多信徒的.
牧養目標13章的9到10節.
並且我們所求的就是你們作完全人.
所以我不再你們那裡的時候.
或者說捨給你們.
好叫我見你們的時候.
用召主所給我的權柄嚴厲待你們.
權柄原是為造就人.
並不是為敗壞人.
作完全人是建立使它變得完全的意思.
依法所書4章12節.
成全聖徒也是這個字.
就是作完全人這個字.
成全聖徒就是使它變得完全.
使聖徒變得完全.
緊接13章11節.
保羅再度提到要作完全人.
還有沒有的話.
願底門都起來.
要作完全人.
連續兩次提到作完全人.
不過第9節說的是名詞.
這裡11節是動詞.

$^{601}$兩次提到作完全人是什麼意思呢.
保羅說的不是一個抽離現實.
絕對意義的完全標準.
他沒有完美主義的想法.
基督徒可以一舉手一投足.
全部適合規矩適合體統.
凡事都不犯錯.
沒有這樣的完全人.
特別他在跟哥林多信徒說.
他們沒有犯罪嗎.
他們只能夠變成完全人嗎.
不可能.
最少這個不是絕對意義的完全.
保羅只是要求他們要作完全人.
要作比完全重要.
我是不是很厲害.
要作完全人.
從你如今的境況開始.
努力改過千線.
完善自己.
我不是期望你一下子達到完全的地步.
你要求我們不斷進步.
朝完全的方向進發.
我們再完全進步.
希望自己一下子達到一個完全的境界.
這是不可能的任務.
你提這些東西只能夠叫人有挫折感.
省一點吧,我做不到的.
但是針對現實的一個缺失.
做完善的工作.
希望自己變得越來越好.
這是永遠有奮鬥的空間.
在菲利比書三章十二到十六節.
這種對作完全人的理解.
可以充分彰顯出來.
我們很熟悉這段聖經.
「這不是說我已經得著了,已經完全了.
我那是竭力追求,或者可以得著基督耶穌.
所以得著我.
但應門,我不是以為自己已經得著了.

$^{641}$我只有一件事,就是忘記背後面前的.
向著標竿直跑,以得上.
在基督耶穌裡,從上面至我那裡得的獎賞.
所以我們中間凡事完全人.
總要存這樣的心」.
後面不讀了.
保羅提到凡事完全人,十五節.
他不是說在信仰或者道德已經達到完全的人.
因為根本不存在著這樣的人.
保羅前面兩次說他不到.
我不是以為我已經得著了,已經完全了.
保羅自己也說自己不行.
他為什麼後面無端端有一句.
我們中間凡事完全人.
凡事完全人的意思就是.
你立志改善自己.
不是已經達到完全.
這也是聖經對我們的要求.
我們努力向聖潔.
向完善的地步進發.
不過上上甚至不是一個雷進.
一個絕對二的雷進.
我是由零,一年後變一.
然後變二,好像升班一樣.
我常說我們基督徒的生命.
可以是行前三步,退後兩步.
可以是上下下下,上下下下.
是不是這樣?上上都是這樣.
不過只要我們今天繼續立志做好.
在我們立志做好的當下.
我們就是要作完全人.
所以弟兄姊妹,這是我們彼此的立志.
接著尾二一段,最後的文案.
十一到十三節.
還有沒了的話.
願在弟兄們都喜樂,要作完全人.
又說要作完全人.
要受安慰,要同心合意,要彼此和睦.
如此仁愛和平的上帝標.
常與你們同在,你們親嘴問安.

$^{681}$彼此務要聖潔.
眾聖徒都問你們安.
在最後的丁玲這一段話.
保羅一口氣說了六個提醒.
要喜樂,要作完全人.
要受安慰,要同心合意,要和睦.
還有要聖潔.
這六個提醒都是針對十到十三章的內容.
或者甚至是將一到九章連在一起的整卷歌頌多後書.
所說過的所有教訓重說一次.
可以說是一個總結性的教導.
這也是他所說的「還有沒了的話」的意思.
「還有沒了的話」不是說一些之前沒有說過的東西.
而是最後總而言之,總括而言.
是一個總結.
十一節的「要喜樂,要作完全人」.
「要受安慰,要同心合意,要和睦」.
都是現在式的命令句.
「要喜樂」,保羅說到他寫這封信.
和他之前所做的所有事情.
都不是為了增加信徒的苦楚.
是希望信徒快樂.
他們的快樂是保羅最大的期待.
「要作完全人」是呼應十三章九節的吩咐.
保羅期待他們改變現在不對的行為.
「要受安慰」,這是和睦用一個被動的語態去翻譯.
「受安慰」,就是被動,別人安慰我.
但其實你可以把它變成一個中性的語態.
意思就是彼此安慰,彼此鼓勵.
不只是被動,也有主動成分.
這跟後面他要說的「同心合意」和「和睦」意思相近.
「同心合意」和「睦」是呼應前面所提到的.
要處理分裂教會的錯誤而做的吩咐.
「如果保羅的吩咐,等於指每次實進行.
保羅應許,如此人愛和平的上帝必尚與你們同在」.
「人愛和平」,用這個去描述上帝.
上帝是愛,上帝也是一個賜和平.
而不是賜憂慮,賜憤怒,賜認罰的上帝.
我們聽從上帝的命令,做上順命的兒女.
上帝就向我們彰顯祂的人愛和和平的一面.

$^{721}$「你們親嘴問安,勿要彼此聖潔」.
親嘴問安,這是猶太人的習慣.
親嘴問安是保羅在書信結束時常有的吩咐.
哥林多前書十六章,羅馬書十六章,.
貼撒雷加前書五章和彼得前書五章都有類似的要求.
值得注意的是,羅馬書跟這裡有完全一樣的要求.
「你們親嘴問安,彼此勿要聖潔」.
羅馬書十六章六十六節.
這個說明是保羅一個很慣性的說法.
「你們親嘴問安,彼此勿要聖潔」.
所以不是特別針對哥林多教會說的特別說話.
眾聖徒都問你們安,這個前面也說了.
但答案是很多人替托保羅問他們安.
那就不說那些承認的人了.
好了,我們來到最後一段.
十三章十四節.
願主耶穌基督的恩惠.
上帝的慈愛,聖靈的感動.
尚與你們眾人同在.
這是一個用三一格式做的祝福語.
我們日後教會基本上多數情況.
如果不是說舊約的祝福語.
多數都是這一段.
願主耶穌基督的恩惠.
上帝的慈愛,聖靈的感動.
尚與我們眾人同在.
基督的恩惠.
赴上帝的慈愛,不需要做太多解說.
恩惠,慈愛,大家都明白.
聖靈的感動是甚麼意思呢?.
《穆拉菲勒比書》二章一節曾提到.
聖靈有甚麼交通.
與此同意相近.
感動或交通.
都是指聖靈的團契,聖靈的臨在.
我們可以與聖靈有交流,有互動.
受祂的啟迪.
這是感動的意思.
三一的上帝帶給我們.
恩惠,慈愛和團契.

$^{761}$這是上帝給我們的福氣.
亦要彰顯在我們的肢體關係裡.
是不是?.
在此,由老生常談.
補充我經常說的話.
不過我所說的話.
有非常堅實的聖經根據.
我要補充指出.
雖然日後教會常常.
基本上最出鏡多的.
就是這個三一格式的祝福禦罕.
不過在十三封保羅的書信.
每封都有祝福語.
在十三封的祝福語裡.
這樣的三一格式只出現了一次.
一次罷了.
沒有了.
壓倒性最多的情況.
保羅的祝福只是提耶穌的名字.
看PowerPoint.
我不讀經文.
七次.
祝福是耶穌基督.
願我主耶穌基督的恩.
常和你們同在.
諸如此類.
七次正提基督.
另外有四次.
沒有提到上帝的名字.
父子姓名都沒有提.
願恩為常和你們眾人同在.
另外.
二不鎖書六章.
一處提到父上帝和耶穌基督.
不提上帝的名字有四次.
父子並提的一次.
父子姓名合提的一次.
正是提耶穌的七次.
正提和合提在內.
提到耶穌基督的名字的祝福.

$^{801}$有九次之多.
這說明什麼呢?.
例如我祝福你的時候.
我提到耶穌基督賜福大家.
耶穌基督的恩惠與你們同在.
這是絕對標準的.
因為聖經最多出現的七次.
十三節裡的七就超過一半.
是這樣的祝福.
我以前經常說.
三一論不是聖經直接記載的教義.
而是後世教會依據聖經的教導而建構的真理.
從歷史上.
尼西亞君士坦丁堡會議和他們的信經.
我們知道.
三一論建構只有一個目的.
一個,一個,沒有二個,沒有二個,只有一個.
就是保衛耶穌基督的完全神聖.
只有這個功能.
三一論只有這個功能.
因為它要將兩個表面上相衝突的真理.
猶太人認定只有一個上帝.
但是第二個我們基督徒要舉高的真理.
耶穌基督是上帝.
你要將這兩個互相矛盾的真理.
加在一起.
只有一個上帝.
不是兩個上帝,不是三個上帝.
但我們又要高舉.
耶穌基督是上帝.
為了要和諧這兩個彼此衝突的教義.
所以我們就說上帝是三一.
又是三,又是一,又是一,又是三.
所以三一論的精神.
是要確保耶穌基督的絕對神聖.
用這個作為我們基督信仰的核心.
所以三一論是基督論.
最少它的原初制定的時候是基督論.
是基督論.
三一論旨在突顯耶穌基督的神聖地位.

$^{841}$它依據的聖經.
當然在新約裡面.
舊約有些輔助.
但是要強調耶穌基督的神聖.
你當然要在新約看.
我要說.
在新約聖經裡面.
為我們提供三一論論證最多的是兩個人.
一個叫保羅,一個叫約翰.
我們先說保羅.
一個叫保羅,一個叫約翰.
我們在這次之後.
我們會開始一個最少三年的旅程.
就是講約翰的呼音.
我們會看約翰.
是保羅和約翰.
離開保羅和約翰.
你無法砌出三一論.
對不起,我要說的.
是嗎?.
無可能砌得到三一論出來.
新約的基礎就是保羅和約翰.
對不對?.
為我們提供三一論的主要證據的是保羅和約翰.
正如十三章十四節.
是一個很堅實的基礎.
但是我要說.
保羅和約翰是兩個為我們找到三一論證據的主要來源.
但是你同時也會看到.
保羅和約翰比我更激烈.
是耶穌基督的鐵粉.
我遲些講約翰的呼音的時候跟你論證.
他是耶穌的鐵粉.
整卷的約翰呼音.
是基督中心都不肯的.
基督中心,基督中心,基督中心.
基督是核心,基督是核心,基督是核心.
基督信仰的核心是耶穌基督.
他是我們拯救和崇拜的對象.
保羅從來沒有興趣做什麼三一平衡.

$^{881}$在提到上帝的時候.
要等量三分之一至三分之一,零分之一.
沒有.
他壓倒性的篇幅.
是講耶穌基督.
他是耶穌基督的僕人.
耶穌基督的使徒.
他畢生宣揚.
是耶穌基督並他釘十字架的呼音.
因為十三章十四節的教導.
我們做三一格式的祝福是一定沒有問題的.
但不是表示一定要三一格式的.
我再說,我們單獨提耶穌基督是完全正義的.
因為有最多聖經支持.
我碰過不少這樣的牧者.
我們教會沒有.
有些牧者甚至連這個唯一的聖經根據都要擅自更改.
他們覺得這是一個錯的說法.
於是他們就改為.
願賦上帝的慈愛,是耶穌基督的恩惠.
聖靈的感動.
父子聖靈的命令才是正確的.
於是就把祝福改了.
這才是符合我們所理解的三一格式.
你把耶穌基督放在賦上帝之前.
就不符合邏輯.
先父後子,又不符合教義.
為什麼子會大過父呢?.
高過父呢?不行.
我常說,這是徹頭徹尾.
將人為的教義凌越於聖經的啟示之上.
我不是聖經主義者.
我不是說我們要跟著聖經的指令.
聖經的說法,有一點偏離,也不行.
我沒有這個意思.
所以你祝福的時候,你找聖靈來排頭.
我也無所謂的.
聖子第二,聖父第三,無所謂.
我不覺得這是問題.
我不是說聖經這樣說就一定要跟足.

$^{921}$但我的意思是.
我們如果高舉一個教義.
好像說聖經的記載也可能是錯的.
那就糟糕了,你明白我說什麼嗎?.
你說什麼我都不覺得是錯的.
但如果你說我要修改聖經的說法.
那對不起,我就一定背你.
對不對?.
保羅寫三一的方程式.
將耶穌先放在賦上.
怎麼可能是錯的呢?.
保羅是錯的人,怎麼可能不是錯呢?.
對不對?.
如果聖經是錯的,那還有什麼好說的呢?.
永遠是聖經優善.
我們知道浸信會.
我們是聖經中心.
聖經是壓倒性最大的.
比所有的教義都大.
我不是說教義沒用.
我不是說三一論沒用.
不過不要後世我們弄很多三一論的新論.
我覺得很多時候是威脅了我們基督中心的信仰.
基督教會,基督徒高舉耶穌基督.
傳揚耶穌基督.
耶穌基督是我們的一切.
耶穌基督,All for Jesus.
He is everything to me.
祂是我們的一切.
好像是老生常談.
我胡說八道.
不過我希望大家確認的是.
保羅,約翰是我們找到三一論主要的聖經根據.
但是他們毫無例外.
都是基督中心的主義者.
我們教會,我們每個人同樣都是.
你問我信什麼呢?.
我信那個為我降世,為我寫命救贖我的耶穌基督.
祂是我的至愛,祂是我的至寶,祂是我的一切.
阿們.

\newpage



\section{馬太福音 9:35-38}
\label{sec:tgXI8MUm32I}
\textbf{2022年1月30日崇拜直播｜梁煒傳道｜我們的6G - 胸懷普世｜馬太福音9\_35-38}
\newline
\newline
連結: \href{https://youtube.com/watch?v=tgXI8MUm32I}{\texttt{https://youtube.com/watch?v=tgXI8MUm32I}} ~~~~ 語音日期: 2022-01-30
\newline
\newline
\hyperref[sec:9NLZUx6GLHU]{\small{< < < PREV SERMON < < <}}
~
\hyperref[sec:index_chronic]{\small{[返順時目]}}
~
\hyperref[sec:XQBczBTx6j0]{\small{> > > NEXT SERMON > > >}}
\newline
\newline
馬太福音 9:35-38
\newline
\begin{longtable}{cl}
\hline
\hline
章節 & 經文 (和合本修訂版)\\
\hline
9:35 & \begin{tabularx}{0.7\textwidth}{X} 耶穌走遍各城各鄉，在他們的會堂裡教導人，宣講天國的福音，又醫治各樣的病症。 \end{tabularx} \\ \\ \relax
9:36 & \begin{tabularx}{0.7\textwidth}{X} 他看見一大群人，就憐憫他們；因為他們困苦無助，如同羊沒有牧人一樣。 \end{tabularx} \\ \\ \relax
9:37 & \begin{tabularx}{0.7\textwidth}{X} 於是他對門徒說：「要收的莊稼多，做工的人少。 \end{tabularx} \\ \\ \relax
9:38 & \begin{tabularx}{0.7\textwidth}{X} 所以，你們要求莊稼的主差遣做工的人出去收他的莊稼。」 \end{tabularx} \\ \\
[1ex]
\hline
\hline
\end{longtable}
$^{1}$主席.
《弟兄姊妹平安》.
繼續是我們的系列廣告.
《我們的六字》.
開始的時候想和大家分享一個故事.
從前有三個木匠.
他們每天都拿著同一樣的工具.
鋸,錘,不同的釘子.
放在木板上.
第一個木匠他覺得他的工作很無聊.
每天重重複複.
無論他建什麼.
他都完全沒有任何概念.
因為他覺得這些都不重要.
他覺得全部東西都是一樣.
都是按著按著按著.
按著要求完成.
能夠賺取生活裡面所需的.
那就已經足夠.
第二個木匠他同樣都是這份工作.
每天重重複複.
不斷地按著按著按著.
不過他知道原來當他將木板按在一起的時候.
能夠興建一個美麗的房間.
而且他也想像到將來有一個很偉大的人物.
住在這個房間裡面.
他一邊工作一邊想著製成品的模樣.
覺得這份工作很有意義.
幫到人.
他希望用心預備.
希望客人喜歡.
第三個木匠他是最認真和最用心.
雖然每天都是重重複複地將釘子按在木板上.
不過他知道他興建的是一座很宏偉的聖殿.
雖然他的工作都很微不足道.
不斷地按著按著.
但是他一想到將來興建的模樣是一個聖殿.
人們能夠去敬拜神.
他就更加用心地去服侍.
同樣一個教徒.

$^{41}$我們不同的人的服侍.
可能都好像其中一個木匠一樣的心態.
其實無論我們的崗位是甚麼.
可能是我們重複地唱歌.
彈琴.
教聖經.
拍一個福音的單張.
甚至乎擺一張椅子.
但是當我們有不一樣的心態.
更加廣闊的眼界的時候.
我們就有不一樣的經歷和眼光.
今天繼續我們的系列講到.
我們的六珠.
第四個G.
Global Perspective.
胸懷普世.
透過耶穌基督的榜樣.
我們一起去學習.
如何擁有一個更加廣闊的Global Perspective.
讓我們一起在普世的大圖畫裡面.
清楚知道我們的崗位.
聖經一開始的時候.
裡面是這樣說.
第35節.
耶穌走遍各城各鄉.
在會堂裡教訓人.
宣講天國的福音.
又醫治各樣的疾病.
馬太福音的第五章到第七章.
是登山的部分.
而第八章和第九章裡面.
記載著耶穌不同的醫治.
他醫治麻瘋病人.
他在加拿大裡面講鬼.
醫治癱瘓.
醫治血漏病的婦人.
讓會堂死去的女兒能夠復活.
而馬太福音第九章.
35到38節.
正好是這兩章聖經裡面的一個總結.

$^{81}$耶穌走遍各城各鄉.
耶穌在地上帶著他的門徒.
最主要服侍的地方是加利利.
猶太的聖經學者約瑟夫說.
在耶穌的時代.
加利利大概有204個.
大大小小裡面的城鎮.
作者刻意地強調.
耶穌走遍全部.
每一個城市.
每一個鄉村.
甚至連在撒瑪利亞人的地方.
他都去傳福音.
這個是有普世的意思.
其實他裡面是想去強調.
每一個城鎮.
每一個城市.
每一個文化.
每一個族群.
都需要耶穌.
耶穌同樣重視.
耶穌的傳道工作涵蓋這一切.
不單止走遍每一個地方.
他也都為每一個人而來.
在他服侍所接觸的人.
他的教導.
傳福音.
醫治.
每一個過程.
他都是充滿著上帝的愛在當中.
他要讓門徒去明白.
福音不單止是給猶太人.
福音是給每一個人.
每一個族群.
每一個人.
他都需要耶穌.
耶穌到城域新來到這個世界.
為了要將這個福音傳遍給萬民.
在約翰福音三章十六節.
神愛世人.

$^{121}$甚至將他的獨生子.
賜給他們一切信他的.
不至滅亡.
凡得永生.
彼得後書三章九節.
主所應許的尚未成就.
有人以為他是擔言.
其實不是擔言.
乃是寬容你們.
不願一人有一人沉淪.
乃願人人都悔改.
這個是耶穌的心意.
耶穌的心意就是希望.
世界上每一個民族的人.
都能夠聽聞這個福音.
當耶穌升天.
和門徒道別.
留下最後的說話.
就是頒佈大使命.
叫門徒去洗萬民作主的門徒.
當耶穌升天之後.
教會再承接著耶穌福音的工作.
在萬國萬人裡面.
宣講上帝的福音.
因此歷世歷代的教會.
上帝興起了一班忠心的信徒.
很多無名的傳道者.
宣教士.
他們越洋過海.
將福音傳到地極裡面.
二百年前.
神使用馬勒信.
將福音帶到中國.
翻譯聖經.
大人順主.
打好美好的聖經基礎.
而一百五十年前.
戴德森更加將福音.
帶到中國內陸的地區.
戴德森至少對中國的福音.

$^{161}$已經有負擔.
他清楚知道當時清朝的晚期.
中國面對著各樣的疾病.
以及苦困.
裡面有很多人民的需要.
他立志將來去到中國宣教.
十九歲的他開始去學醫.
過著艱苦的生活.
他的母親很擔心他.
帶他去回信.
他說每年有一千二百萬的中國人死去.
他們沒有神.
他們沒有盼望.
我們不可以置之不理.
他覺得他一定要去為中國做一些事情.
不能夠就此罷休.
之後幾個月.
他就坐船去到中國裡面.
去做船內的工作.
大約過了十年.
一百六十四年的時候.
戴德森患了病.
他需要回英國裡面.
去治療這個疾病.
當時他在英國的時候.
出版了一個刊物.
中國屬靈的需要和託付.
他不斷呼籲人去到中國.
內地不同的省份裡面.
因為裡面每個月都有過百萬的人.
因為疾病.
因為困苦.
他們死去.
沒有聽過福音.
在一次崇拜的裡面.
他看到一千個基督徒.
很喜樂地去聚會.
敬拜札美神.
不過他想到有千萬個中國人.
他們因為沒有聽福音.

$^{201}$而離開.
他不能夠忍受.
他心裡難過.
於是他去向神祈禱.
跪下求神去參見更多人.
去到中國裡面去傳福音.
他認為英國的人.
能夠領受很多上帝的祝福.
不能夠眼白.
白白地看到很多中國人.
因為不知道福音.
而走向滅亡的道路.
每天有三千三百人死亡.
每個月有一百萬的中國人.
沒有聽過福音.
走向滅亡.
戴德森深愛中國.
為中國人的福音消極.
因著戴德森的福事.
他帶動了世界各地不同的教會.
差不多有成千上萬的宣教士.
來到中國裡面傳福音.
不單止傳福音.
更加學習耶穌.
辦學,醫療,戒毒,孤兒院.
回應中國裡面的需要.
1950年代.
因著內地的政策改變.
宣教士需要離開.
香港就成為了外國宣教士的一個中轉站.
希望香港能夠成為一個猜拳的基地.
當有一天.
這個福音的門再打開的時候.
能夠將福音再次傳遍中國.
因此不同的宗派都派人來到香港.
無數的神父,穆者.
在這裡建立教會.
透過辦學,醫療,很多社會福利機構.
讓香港人有機會聽聞福音.
根據2019年的年報.

$^{241}$香港的基督徒,天主教徒.
大約有100萬.
而恆常聚會的基督徒大約有30萬.
當中其實有很多基督教的機構.
大學的大專院校有3間.
中學有超過180間.
小學有200間.
260間幼稚園.
130間幼兒中心.
20多間的神學院.
有各樣的出版社.
也有老人院,醫療中心,兒童病院.
超過了730個基督教的機構.
原來今天我們能夠很自由地去敬拜.
很自由地去認識上帝.
有豐富的資源.
原來一切都是因為上帝派了很多宣教士.
很多無數如雲彩般的見證人來到香港.
祝福我們.
讓我們有機會聽聞福音.
我們華人領受福音超過了200年.
領受了很多的祝福.
是時候接福音的棒子.
參與在普世宣教的裡面.
香港的信徒在國際上扮演著一個很重要的角色.
特別是連接中國大陸的宣教士到海外.
記得在幾年前.
我去到一個穆斯林國家裡面去短孫.
那時候遇到十幾個中國的年輕人.
他們大約二十多歲.
他們是內地教會的宣教士.
他們分享他們的蒙召經歷.
他們是十八歲蒙召.
那時候一起參與一個營會.
大約有四百人一起蒙召.
經過半年的基礎訓練之後.
其中有一百個人.
參拜到非洲的國家裡面.
接受兩年的宣教訓練.
學習了謀生的技能.

$^{281}$過著信心的生活.
不過過程裡面有很多人捱不住.
完成了之後.
他們按照警察會.
按照教會的要求.
需要派到不同的國家.
而這十幾個人.
就被派到穆斯林的國家裡面.
創啟的地方裡面去傳福音.
他們在當地學習語言.
不過因為他們的語言基礎差.
語言成為他們傳福音很大的障礙.
加上要有合法的身份留在當地.
要有謀生的技能.
當中出現了不同的困難.
他們和不同團隊裡面的合作.
也都有各樣的挑戰.
一個資深的宣教士去提出.
其實香港的人.
可以幫助中國大陸的宣教士.
適應迎向國際這個平台.
今天的香港信徒.
其實無論在經濟能力.
教育的程度.
聖經的知識.
神學裡面的普及.
人際關係.
語言能力.
國際視野.
其實都有得天獨厚的優勢.
面對普世猜全.
我們都是責無旁貸.
我們怎樣可以啟動呢?.
經文裡面繼續去說.
第36節.
「他看見許多的人就憐憫他們.
因為他們困苦流離.
如同羊沒有牧人一般」.
耶穌看見很多人.
有黑眼的.

$^{321}$有瘸腿的.
有患大麻瘋的.
有碎裂的.
有婦人.
有兒子死了的.
也都有愛邦人.
有官長.
有要水的.
有要食物的.
很多人冒名來找耶穌.
當每當耶穌進入一座城裡.
合城的人都帶著各樣的需要.
來找耶穌.
耶穌由天光做到天黑.
第二天早上起床的時候.
耶穌退到抗夜祈禱.
祂的門徒馬上找上門.
因為有很多人清早起來找耶穌.
耶穌積極回應人的需要.
耶穌關注人裡面的需要.
看到他們.
認識他們.
重視他們.
今天的資訊很發達.
甚至我們可以說是資訊氾濫的年代.
可能一部手機.
可以上網搜查到各樣的資訊和情況.
香港人都很善於利用一些資訊的人.
無論是最近的疫情.
什麼時候有多少的數字確診.
哪裡有事.
香港人都會得到最新的消息.
特別是教育局的措施.
家長最關注.
我們都會留意各種的優惠.
有好吃的.
有好玩的.
各樣潮流各樣的興趣.
不過耶穌唯一的關注.
是人的需要.

$^{361}$祂將人的需要放在一個很優先的次序.
因為耶穌祂知道祂要為人而來.
解決人裡面的需要.
今天當我們的眼光.
不再關注自己的時候.
方言世界.
其實我們會發現這個世界裡面不同的情況.
疫情裡面最關注的是世界各地感染的數字.
疫苗接種的情況.
原來香港第一針疫苗都超過了.
八成的人打了針.
第二針就超過七成的人.
但原來非洲裡面.
很多大陸裡面的國家.
疫苗接種率都不夠5\%.
這個是一個很嚴重的問題.
病毒最初由法特的國家.
傳到這些貧窮的國家.
我們都慢慢有不同的保障.
口罩或者其他搓手液.
但原來他們每一天.
連最基本的保障都沒有.
難民的問題.
根據聯合國2020年的數字.
大概有八千多萬.
敘利亞為之最多.
其次是委內瑞拉,阿富汗,南蘇丹,緬甸等等.
因為他們正在經歷長期的戰亂.
大量的人民流離失所.
遭受迫迫衝突暴力的情況.
特別是去年的時候.
阿富汗和緬甸的政局的變遷.
導致更多人流離失所.
當我們看新聞的時候.
我們看到很多很破碎的照片和短片.
糧食的問題.
根據世界聯合國的糧食署的統計.
2017年全球每年超過八億的人.
是正在捱餓.
隨著氣候的極端變化.

$^{401}$很多地方面臨旱災.
農作物失收.
加上新冠疫症問題更加嚴重.
全球的人口販賣.
現在全球大約有2700萬人.
成為了奴隸.
比以前的黑奴時代.
1500萬人是高出兩倍.
這些人被禁錮.
被迫去做一些他們不想做的事情.
大部分都是女性.
而且大部分都是18歲以下.
未成年.
更何況是貧富的縣市.
窮人受剝削.
各地的政府貪污腐敗.
很多很多淋淋腫腫的問題.
當我們方言世界的時候.
原來世界的人都生活在水深火熱.
每日的國際都有很多悲劇.
很多事情發生.
有很多新聞被報導.
但有很多每日發生的事情.
他們不再成為新聞.
亦不再被報導.
可能看這些數字.
我們都會覺得有很大的無力感.
不知道怎樣去入手.
不過就好像那個男孩.
在海灘裡面.
去拯救海星的故事裡一樣.
成千上萬的海星.
他的確救不完.
不過對於被拯救的海星來說.
的確是在改變個別的命運.
耶穌面對著這個龐大的需要.
他看到他們.
他更加動了持心.
憐憫他們.
憐憫這個字的字根.

$^{441}$是在說腸和內臟.
因為在希伯來人的文化裡面.
他認為最深層是他的腸和內臟.
憐憫是在說他內心.
有很強烈的感受和情感.
耶穌是充滿著憐憫.
對著他所服侍的對象.
他能夠心同感受去幫助他們.
在耶穌整個服侍裡面.
他由始至終都帶著憐憫的心.
去接觸每一個人.
這個不單只是一個情感.
更加去幫助他.
和人連結在一起.
分享上帝天國的奧秘.
當他伸手摸死人.
摸病人的時候.
在當中分享上帝的愛的時候.
上帝透過憐憫.
醫治很多人.
很多人的身,靈.
得到醫治.
耶穌深深明白他們的需要.
感同身受.
在大兒子和小兒子的故事比喻裡面.
大兒子和他的父親.
其實同樣都見到小兒子.
因為放蕩的生活.
過著顛沛流離,倒楣的日子.
不過不同的是.
這個父親見到小兒子.
動了持心.
想像和他親嘴.
宰殺一隻牛去慶祝.
不過大兒子因此而很憤怒.
在好撒瑪利亞人的比喻裡面.
祭司和民事.
同樣都見到被強盜打個半死.
有沒有人沒有穿衣服的人.
但他們沒有上前幫忙.

$^{481}$唯有那個撒瑪利亞人.
他見到那個人的需要.
動了持心.
上前去包紮.
能夠帶他去旅館裡面被照顧.
原來最重要的.
不單止是我們要去看到.
更加是我們有沒有動我們的持心.
非洲的馬達加斯加.
是全世界十大貧窮的國家.
裡面的人平均每日.
都沒有兩元的收入.
即是連一杯珍珠奶茶都買不到.
雖然他們有豐富的資源.
但因為貪污腐敗.
人民過著很辛苦艱難的日子.
他們沒有食物,沒有清水,沒有電話.
亦都沒有教育的機會.
我想這些資訊我們都不難去得到.
不過我們會不會有些時候.
聽得多就變得麻木.
沒有感受呢?.
記得幾年前和太太去到馬達加斯加裡面去服侍.
宣教士帶我們去到當地的打石場.
在當地裡面有很多的打石場.
當你站在打石場裡面.
望落去工場裡面.
滴答滴答滴答滴答滴答.
很多的聲音充滿著這個打石場.
原來這些聲音.
是那些錘敲下去的釘的聲音.
在打石場裡面的人.
全部都是用簡單的工具.
一個錘仔,一個釘.
將大塊的石頭打打打.
打到好像小塊的石頭.
將小塊的石頭再打打打.
打到好像碎石.
用來做建築.
可能在香港我們有很大型的機器.

$^{521}$不過在貧窮的國家裡面.
他們沒有大型的機器.
他們的耳朵長期沒有任何的保護.
受到嚴重的破壞.
在烈日當空裡面.
經常工作沒有水喝.
他們的手經常勞損.
眼睛的視力都是受損.
但是他們堅持工作.
我們遇到一家五口.
他們帶著嬰兒帶著小朋友.
一邊生活一邊工作.
因為無家打石場.
都成為他們棲息的地方.
他們彷彿是古埃及的以色列人一樣.
過著為奴的日子.
世世代代都沒有辦法脫離貧窮.
每一次我們帶著短孫隊去打石場.
很多短孫隊的眼淚不斷流下來.
這個情景不單只是在白天裡面很觸動他們.
晚上甚至過了一段時間.
他們都會記掛著打石場的人.
為他們的難過.
為他們的事情傷心流淚.
從前我們都覺得.
我們對這個世界裡貧窮的地方.
的資訊都掌握很多.
不過這些都是頭腦上的認知.
當我們親眼看到一切的時候.
這些眼淚都在表達我們裡面的悲痛.
當然我們沒有辦法去每一個地方.
親身感受到每一個場景的困苦.
但我們面對著世界裡面的問題.
我們不可以麻木.
我們需要仍然有感受.
耶穌在馬太福音24章裡面.
講到末世的一個情景.
他說末世裡面會有天災會有人禍.
教徒都會受迫迫.
然後十二節裡面說.

$^{561}$只因不法的事增多.
許多人的愛心才漸漸冷淡了.
當這些不法的事.
當這些悲慘的事一天一天發生.
原來我們的愛心是會不斷地冷淡.
愛心或者憐憫其實是需要操練的.
去對抗各種的冷淡.
教徒的愛心其實不是透過讀經.
不是透過上教會自然能夠增長.
有時候反而我們讀聖經的時間越多.
越對愛心的理解越多.
錯誤地以為我們都是一班有愛心的人.
香港人的生活節奏很急速.
會收到很多的資訊.
沒有空間去處理.
很快地去看了這些資訊.
然後沒有空間去處理.
久而久之我們對很多的資訊.
很多的數字都變得盲目.
再加上現在的媒體.
令我們的專注力都降低.
現在無論是Facebook.
IG或者抖音TikTok.
15秒一個的視頻.
令我們更加缺乏耐性.
我們需要怎樣去操練呢?.
我相信每一個人都很忙.
但忙的裡面我們可以選擇.
在每一次服事的裡面.
都專心做好每一樣的服事.
全神貫注在那個對象裡面.
耶穌都很忙的.
祂由早到晚都在服事.
但祂每一次的服事.
就是專心全情投入去接觸那個對象.
將自己的需要放下.
想著的是對方的需要.
將整個心完全給了對方.
我們未必會對所有事情都有憐憫.
我們亦未必能夠解決很多的問題.

$^{601}$不過我們可以選擇.
在我們遇到的事情裡面.
我們全心全意地投入.
進入對方的世界.
試下心和感受.
去對人產生興趣.
對人產生疑問.
其次當我們見到很多不幸的事情.
我們有時候很快地先入為主.
下了一個判斷.
當我們看新聞的時候.
下面都會有很多的評論去說.
為甚麼會造成這些問題的原因.
讓我們停一停.
不是站立在道德的高地裡面.
斥責問題的原因.
反而是去為對方的處境.
去多想一些對方的困難.
不只是批評或者是指責.
可能對事情都未必有幫助.
從積極態度去想想怎樣去解決.
第三.
耶穌的憐憫是帶著行動的.
相信我們很多人都未必.
可能很有空間時間.
能夠去到不同的地方去幫人.
不過香港人是一個很被祝福的地方.
特別是在金錢上.
我們可以捐獻.
試試每個月捐一些你覺得很信任.
或者你認識的機構.
每個月幾百元.
可能我們吃少一餐比較豐富的食物.
其實對我們來說.
我們的生活不大影響.
不過對於我們所要去福祉的對象來說.
他們卻有很大的幫助.
當我們願意去憐憫其他人.
我們對人的敏銳度都會增加.
會減少評論.

$^{641}$跟人再次連結.
分享上帝的聖情.
操練我們的生命.
讓我們更加像耶穌.
耶穌憐憫他們.
不單只是看到他們沒有飯吃.
身體上受到各樣疾病的煎熬.
耶穌更加看到的是.
他們屬靈裡面的需要.
耶穌看到很多人困苦流離.
好像羊沒有牧人一樣.
耶穌眼前的對象不是一些難民.
也不是非洲沒有食物的人.
也不是被人口販賣的人.
不過耶穌看到他們困苦流離.
困苦如文是備受迫迫和威脅.
流離是指被遺棄.
要倒下的意思.
說的是這群人生活在困苦.
他們生活一點都不好過.
在耶穌的年代.
猶太的領袖沒有好好保護他們的子民.
在羅馬政權的壓迫下吃盡苦頭.
生活裡面各樣的艱苦.
人民都是無助和迷茫.
馬太形容這群人.
如同羊沒有牧羊人一樣.
羊沒有方向感會迷失.
當羊沒有牧羊人的時候.
沒有任何的保護.
會成為野獸撕裂的對象.
任由宰殺.
在2005年的時候.
土耳其有一個蓋亞集的牧羊人.
他吃早餐的時候忽略了四圍奔跑的羊.
先是一隻羊掉下山崖.
結果成千隻羊跟著掉下山崖.
釀成了四百多隻羊死亡.
這個新聞反映了原來沒有牧羊人的生活.
就好像人沒有了福音和羊主一樣走向滅亡.

$^{681}$當我們看到世界的需要.
不單單是看他們表面的條件好像很好.
他們不需要耶穌.
只有窮人才需要耶穌.
但他們同樣的宿命裡有很多需要.
他們需要救恩.
就好像日本一樣.
日本是一個很先進的國家.
我們都會去旅行.
會被他們的食物,文化,禮貌所吸引.
看他們的生活.
好像覺得他們很好.
不需要耶穌.
但日本的色情文化,動漫文化.
長期工作,家庭的問題.
婚姻的離婚率很高.
人的心靈極度貧乏.
自殺率也很高.
人生活在迷茫和空虛的裡面.
同樣面對著貧窮.
其實他們都是沒有糧食,沒有水.
不過最重要的是耶穌看到他們.
是因著沒有福音的緣故.
過著艱苦的生活.
不單要提供他們生活裡的所需.
更重要的是讓他們認識上帝.
德蘭修女在她服侍的過程裡.
她服侍在這個世界裡被遺棄的人.
她同時留下了很多寶貴的話.
有兩段她說.
這個世界裡充滿著痛苦.
包括身體上,物質上和精神上.
而這些人所經歷的痛苦.
可以歸咎於另外一些人的貪婪.
物質上和肉體上的痛苦.
可能來自飢餓,無家可歸和各樣疾病.
但是最大的痛苦是孤單,寂寞.
沒有人愛,沒有任何一個朋友.
當她越來越感受到的時候.
原來人生旅途裡最大的疾病.

$^{721}$是被人遺棄.
德蘭修女說到.
世界上最嚴重的貧窮不是食物裡的缺乏.
而是愛的缺乏.
有些人不滿足於他所擁有的一切.
他也不知道如何去承受苦難.
他掉進無底的深淵裡.
心靈裡的貧窮.
他很難去消除和戰勝.
今天這個世代裡面面對著這個疫情.
疫情帶來人與人之間很大的隔閡.
特別是一些長者,基層的人.
他們被隔離,長期病患者.
這個疫情裡加深了人內心的需要.
人顯得更加空虛,迷茫和孤單.
我們能否看到這些需要呢?.
最後一個部分.
第37至38節.
於是對門徒說.
你收的莊家多,作供的人少.
所以你們當求莊家的主.
打發工人出去收他們的莊家.
莊家多.
在孫格學裡有個詞叫做十四十之昌.
是形容北緯十度到四十度之間.
包括北非,中東,亞洲這些地區.
而十四十之昌大部分居住的都是穆斯林.
印度教徒和佛教徒.
這些地區超過六十多個國家.
沒有任何宣教機構,教會和宣教士在裡面.
大多數的人都沒有機會聽到這個福音.
有將近四十億的人口在裡面.
而全球90\%最貧窮的人.
也都在這個十四十之昌裡面.
裡面的失業率很高.
超過了一半.
很多世界各地的猜拳基金扶貧都會在裡面.
不過原來佔的百分率只有幾個百分比.
當每天看到新聞的時候.
他們的環境,恐怖主義,不公平的貿易等等.

$^{761}$原來在十四十之昌裡面.
大約有61\%的人口都在裡面.
六千個民族的人都沒有機會聽到耶穌的福音.
在全球首望名單裡面.
今年的名單裡面有一些轉變,更新.
過往的二十年.
北韓一直佔領全球十.
全球迫不得已基督徒最嚴重的國家的榜首.
不過今年阿富汗卻是取代了北韓.
成為了全球基督徒最危險的國家.
裡面全球首望名單裡面記錄了超過三點六億的基督徒.
因為信仰的緣故遭受到各樣程度的迫不得已.
比去年增加了二千萬.
這些數字就是七個基督徒裡面.
有一個基督徒受到迫不得已的迫不得已.
而這些數字都是過往二十九年裡面最嚴重和數字最高.
可能這兩年的疫情告訴我們.
地球只是一條地球村.
一個國家,一個城市發生的問題,一個病毒.
都可以傳染給這個世界.
其實同樣的,基督徒的迫不得已問題都越來越嚴重.
越來越被傳開.
好像聖經預言末世的情景將來有不大迫不得已.
我們要做的,是要更加主動地去傳揚這個福音.
在全球的裝甲好像很多.
有五十多億的人沒有聽過福音.
但原來宣教士只有四十二萬在當中.
而香港猜測的宣教士大約是六百個.
他們平均年紀,即是五十歲.
面對著世界五十多億的人.
每個宣教士要傳福音的人數以千萬計.
單靠宣教士,這個福音何時才能夠傳開呢?.
耶穌基督不單止將這個全福音普世萬民的責任.
給宣教士,更加給我們每一個基督徒.
第38節,所以你們當求裝甲主.
差派工人去收他的裝甲.
宣教是上帝的工作.
所以必須用上帝的方法,用祈禱去進行.
記得在十年前,我參加了一次日本的短訓.
聖誕節時,我們去到日本金澤的地方辦聖誕營會.

$^{801}$短訓隊裡有超過十多個國家的人.
一起來服侍日本人.
其中一個崇拜的.
他看到上帝有很多工作.
他深深感受到上帝沒有忘記日本.
接著他帶領我們為日本去認罪祈禱.
牧師跪在地上,眼淚不斷流下.
特別是為了日本過去的年日.
對東南亞的地區造成很多的侵戰,發動戰爭.
他去認罪悔改.
他受上帝的憐憫,求上帝的醫治.
求上帝差派更多的人去到日本的地方.
幫助日本人認識福音.
幫助日本人得聞好消息.
他的祈禱很震撼.
在場的每一個人都被感動,不斷流眼淚.
希望神差派更多人去到日本的裡面.
我們愛不愛我們所身處的地區.
我們有沒有去為香港有這個祈禱.
這個牧師因為深愛日本,深愛上帝.
所以他經常跪下為日本祈禱.
有一個莫拉維的弟兄會.
他們發動了超過一百年的祈禱運動.
這個莫拉維弟兄會始於十五世紀.
他們都是很愛主的一班人.
在十八世紀裡面,有一個愛主的領袖出現了.
這個人叫做親神多夫.
他被上帝的十字架的愛所感動.
他願意奉獻給上帝.
他不單覺得上帝愛貴族.
更愛世界上每一個人.
每當他見到耶穌被釘在十字架上.
頭戴著荊棘冠冕.
他就去祈禱問上帝.
上帝啊,我可以為你做些什麼.
他被大大感動.
他願意將所有的時間,生命,錢財.
一切奉獻給上帝,被上帝所用.
這個莫拉維弟兄會.
全部都是一些很愛主,很委身的人.

$^{841}$在1927年,大約三百個莫拉維弟兄會.
他們成立了一個叫做主護村的地方.
成為他們的棲息地方.
他們最初為了九個女孩去為神祈禱.
求神去醫治他們.
求神去幫助他們.
深受感動.
直至之後他們發動了二十四小時的祈禱.
就像利未記六章十三節裡面說到.
在壇上必常有燒著的火不可熄滅.
莫拉維弟兄會不斷祈禱.
二十四小時的祈禱,每天的祈禱.
這個祈禱持續了一百年.
是一個祈禱鏈.
在這一百多年的時候.
有超過三千個宣教士.
被差派到世界各地去傳福音.
遍佈世界各地的各落.
祈禱是幫助我們與神連結.
祈禱更加幫助我們與世界裡的人連結.
上帝將宣教的託付交給我們.
讓我們能夠守望這個世界.
面對著世界的需要.
我們怎樣去做呢?.
有些小的實踐.
第一,其實現在資訊平台都很豐富.
有個平台叫做Spotify.
裡面有兩個頻道.
一個叫做「角度換日線」.
一個叫「使命門徒」.
每個星期都會訪問世界各地不同的情況.
無論是華人也好,穆斯林也好.
或者歐洲國家,非洲國家裡面的需要.
他們去看上帝在那個地方裡面的工作.
訪問上帝,訪問當地的人.
上帝在那裡的作為.
開闊了很多的眼界.
第二,你也可以Like很多的資訊.
無論是在Facebook,定期訂閱.
當然網站的頻道也要訂閱.

$^{881}$「躺開的門」.
留意更新一些祈禱的需要.
你們可以組成一些小組.
家庭的祈禱守望.
我們的家庭其實每個星期都會.
有些家庭崇拜.
每次都會介紹不同地方的需要.
告訴小朋友.
記得早兩天,小燕子.
我是大女兒,三歲多.
她和另一個朋友在公園玩.
她朋友就去了小食亭買食物.
買了一顆珍寶珠.
吃進口,覺得很好.
那時候小燕子沒有吃.
因為我們沒有買零食給她.
她回家後就跟我說.
她沒有珍寶珠不要緊.
因為非洲裡面很多小朋友.
其實他們都沒有食物.
原來我們一點點的教導和服侍.
小朋友都看在眼內.
能夠改變他們的眼光.
可能我們因為疫情的緣故.
未必可以去到外地.
不過我們可以由本地的警察團隊開始做起.
全港有五六十萬個的.
少數族類在香港.
包括巴基斯坦尼泊爾.
印尼或其他群體在當中.
去年的時候,我們大專團體.
其實都開始每一季去服侍尼泊爾人.
他們都很友善.
亦都有很多不同的需要.
很願意和人接觸.
雖然我們都身處在家中.
但我們同樣都可以透過網絡傳福音.
過往的疫情,其實我們都透過Zoom.
在東南亞一個重的國家裡面.
和大學生去接觸.

$^{921}$每個星期和他們用英語對話.
和他們開聖經單去傳福音.
上帝在這個世代裡面.
放下了很多的資源.
重點是我們怎樣去回應他呢?.
在新的一年,祝福大家.
生命成長,我們的心,我們的眼.
能夠好像耶穌一樣.
望祂所望的.
為祂所愛的東西去召集.
得人如得魚,成為這個世代裡面的祝福.
讓我們最後一起祈禱.
親愛的主.
你愛我們.
在我們還沒有認識你的時候.
其實你已經派了很多人.
來到中國,來到香港.
裡面將這個福音傳給我們.
這群無名的人獻上他們的時間.
青春血汗.
讓我們能夠得到這個福音.
神啊,今天我們有很多福音的好處.
當我們方言世界的時候.
有很多的需要.
我們不能夠每一個需要都去滿足.
大神啊,我們祈求你.
告訴我們.
我們今天要怎樣去回應某些特定的需要.
不需要全部做完.
但我們最重要的是.
我們要去回應你.
最後我們最懇求的是.
你拆拍更多的工人.
拆拍人去到世界各地.
很多人都很需要你.
他們不單止生活裡面.
過著各樣的疾病.
各樣的問題.
他們在心靈裡面都有很多空虛.
主耶穌求你憐憫他們.

$^{961}$求你幫助他們.
拆拍更多人.
求你讓我們的心.
常常去記掛著這些不同的需要.
我們用祈禱.
我們用金錢.
我們用我們有能力的方法.
去支援他們.
聽祈禱奉耶穌特性之名.
拜拜!.
\newpage



\section{腓立比書 1:1-11}
\label{sec:XQBczBTx6j0}
\textbf{2022年2月6日崇拜直播｜陳明泉牧師｜我們的6G - 凝聚動力 GET CLOSER｜腓立比書1\_1-11}
\newline
\newline
連結: \href{https://youtube.com/watch?v=XQBczBTx6j0}{\texttt{https://youtube.com/watch?v=XQBczBTx6j0}} ~~~~ 語音日期: 2022-02-07
\newline
\newline
\hyperref[sec:tgXI8MUm32I]{\small{< < < PREV SERMON < < <}}
~
\hyperref[sec:index_chronic]{\small{[返順時目]}}
~
\hyperref[sec:Q_GbM3l6t1s]{\small{> > > NEXT SERMON > > >}}
\newline
\newline
腓立比書 1:1-11
\newline
\begin{longtable}{cl}
\hline
\hline
章節 & 經文 (和合本修訂版)\\
\hline
1:1 & \begin{tabularx}{0.7\textwidth}{X} 基督耶穌的僕人保羅和提摩太寫信給住腓立比、在基督耶穌裡的眾聖徒，以及諸位監督和執事。 \end{tabularx} \\ \\ \relax
1:2 & \begin{tabularx}{0.7\textwidth}{X} 願恩惠、平安從我們的父神和主耶穌基督歸給你們！ \end{tabularx} \\ \\ \relax
1:3 & \begin{tabularx}{0.7\textwidth}{X} 我每逢想念你們，就感謝我的神， \end{tabularx} \\ \\ \relax
1:4 & \begin{tabularx}{0.7\textwidth}{X} 每逢為你們眾人祈求的時候，總是歡歡喜喜地祈求， \end{tabularx} \\ \\ \relax
1:5 & \begin{tabularx}{0.7\textwidth}{X} 因為從第一天直到如今，你們都同心合意興旺福音。 \end{tabularx} \\ \\ \relax
1:6 & \begin{tabularx}{0.7\textwidth}{X} 我深信，那在你們心裡動了美好工作的，到了耶穌基督的日子必完成這工作。 \end{tabularx} \\ \\ \relax
1:7 & \begin{tabularx}{0.7\textwidth}{X} 我為你們眾人有這樣的想法原是應當的，因為你們常在我心裡；無論我是在捆鎖中，在辯明並證實福音的時候，你們都與我一同蒙恩。 \end{tabularx} \\ \\ \relax
1:8 & \begin{tabularx}{0.7\textwidth}{X} 我以基督耶穌的心腸切切想念你們眾人，這是神可以為我作證的。 \end{tabularx} \\ \\ \relax
1:9 & \begin{tabularx}{0.7\textwidth}{X} 我所禱告的就是：要你們的愛心，在知識和各樣見識上，不斷增長， \end{tabularx} \\ \\ \relax
1:10 & \begin{tabularx}{0.7\textwidth}{X} 使你們能分辨是非，在基督的日子作真誠無可指責的人， \end{tabularx} \\ \\ \relax
1:11 & \begin{tabularx}{0.7\textwidth}{X} 更靠著耶穌基督結滿仁義的果子，歸榮耀稱讚給神。 \end{tabularx} \\ \\
[1ex]
\hline
\hline
\end{longtable}
$^{1}$今日經訓是《肥臘比書》一章十一節.
新約聖經二百七十三頁.
《肥臘比書》一章一至十一節.
基督耶穌的僕人保羅和提摩泰.
寫信給凡住肥臘比.
在基督耶穌女的眾聖徒和諸位監督.
諸位執事.
願因平安從神我們的父.
令主耶穌基督歸於你們.
每逢為你們眾人祈求的時候.
嘗試歡歡喜喜的祈求.
因為從頭一天直到如今.
你們是同心合意的興旺福音.
我深信那在你們心裡動了善功的.
必成全者功.
直到耶穌基督的日子.
我為你們眾人有這樣的意念.
原是應當的.
因你們常在我心裡.
無論我是在困所之中.
是變名證實福音的時候.
你們都與我一同得恩.
我體會基督耶穌的心腸.
切切地想念你們眾人.
這是神可以給我作見證的.
我所禱告的.
就是要你們的愛心.
在知識和各樣見識上多而有多.
使你們能分辨是非.
作誠實無過的人.
直到基督的日子.
並靠著耶穌基督.
結滿了人兒的果子.
叫榮耀稱讚歸於神.
多謝各位.
主席.
各位弟兄姊妹平安.
在新的一年裡.
祝大家展翅上騰.
身體健康.

$^{41}$很快就經歷一個新的年頭.
希望在這一年裡.
雖然好像疫情為我們展開農曆新年.
不過我們深信上帝會保守我們.
我們帶著上帝所賜給的喜樂和平安.
我們安然地度過這一段日子.
我希望我們很快地.
從回到教會的禮堂裡.
一起來敬拜上帝.
今天我們一起來看第五個G.
我們的六G.
凝聚動力.
Get closer(靠近).
中英文翻譯上有一點出入.
當然我們最想說的內涵.
在旺朕的價值觀.
或者在我們教會裡面.
我們有一個很重要的價值.
就是讓姐妹聚集在一起.
不是純粹是一群等待被餵養.
或者純粹是消費宗教一群的群體.
而是我們希望我們一群人聚集在一起.
我們成為一個改變的動力.
為你的社會.
為你身處的環境.
或者為你的家庭.
為你的職場.
我們這一群弟兄姊妹聚集.
是會帶來一些更新和變化的.
所以我們的聚集.
是有目標和有方向感的.
今天我就很想選取.
肥勒比書的《隱言》或者《開場白》.
作為去展述第五個G.
Get closer(靠近).
凝聚動力的一個很重要的經文.
事實上當你仔細去看這段經文的時候.
其實保羅的語調和肥勒比教會.
特別你看到的哥林多前書.
或者其他書卷.

$^{81}$有一點不同的.
肥勒比書是一份情書.
是保羅將他在上帝面前.
所領受的一個很深的情.
來呈現給弟兄姊妹之一.
希望今天我們去讀這一個經文的時候.
或者我們仔細去看的時候.
我們不單止心裡面.
會更加感受到保羅的心腸.
更加我們真的可以凝聚動力.
讓旺振成為一個更新社會.
在福音運動之下.
我們成為上帝一個重要的接棒者.
希望今天我們一起領受到這個訊息.
在我們仔細講這段經文之前.
我想講講動力或者凝聚動力的.
坊間裡面都有很多不同的著作.
或者心理學也好.
文化學也好.
都有些人都會做這些工作.
有些做法是怎樣做呢?.
怎樣凝聚人的動力呢?.
或者凝聚人呢?.
有幾樣東西的.
第一就是高舉一個叫共同利益.
大家一班人走在一起.
大家一起謀生.
大家有共同的.
希望得到的一些利益或者方向.
將這一班人藉著利益.
來將所有人串連在一起.
這是其中一個在坊間裡面.
說凝聚動力是一個很重要的手法.
第二個就是將一班人.
困在一個密封環境裡面.
你留意一下.
有時候有一些很極端.
甚至乎我們看過一些邪教.
他們將那些弟兄姊妹也好.
或者那些信奉者.

$^{121}$他們就會與這個世界與世隔絕.
讓他們過著一個密封式的生活.
當他們與世界隔絕的時候.
身邊的人就變成了.
他一個很生命裡面很重要的人.
在很多的異端邪教裡面.
會用這些手法.
來令到人凝聚這些動力.
令到他們得到他們心目中想要的東西.
第三種就是敵我分明的做法.
強調這個世界有很多敵人.
這些自己人混在一起.
那些才是真正的格見因.
所以敵我分明的思想.
就令到那些人的動力可以更加凝聚.
最後可能有其他這些因素之後.
再加上一個叫做魅力型的領袖.
有一個打造一個明星級的人.
講的一舉十一頭竹.
或者要抄筆記的.
他講每一樣東西.
令到所有人都是心悅誠服地來跟從.
這個就是坊間所講的凝聚動力的方法.
不過你看肥立比書.
保羅不是用這些方法的.
聖經裡面教我們怎樣去重新.
凝聚這群基督徒群體生命裡面有動力.
今天在第一章一至十一節.
都有很豐富的內容.
我們一起祈禱.
求上帝幫助我們.
讓我們明白這段經文的內容.
同心祈禱.
主願你在我們當中帶領我們.
幫助我們眾弟姐妹.
讓我們不論在家中.
或者我們一起在不同的地域裡面.
讓我們專心一意尋求你的話語.
願你親自對我們講說話.
幫助孫講的講得清楚.

$^{161}$聆聽的每位弟兄姊妹聽得明白.
恭敬將此時此刻交託給你.
願你賜福.
獻上我們的禱告.
奉靠基督耶穌得勝的聖名.
Amen.
我們一起看.
在肥立比書裡面第一章一至第二節.
其實是一個比較普遍的開場白.
基督耶穌獨人保羅和提摩泰寫信.
給凡住在肥立比和在基督裡的眾聖徒.
和諸位監督 諸位執事.
願因為平安從我們的父.
並主耶穌基督歸於你們.
在這裡的祝福語其實跟一般的.
祝福語的格式是近似的.
不過在這裡裡面.
保羅強調他是一個僕人的身份.
這個僕人的身份.
他帶著的是一份人際之間的謙卑.
但同時間他是效忠那位基督耶穌的.
所以他帶著既是上帝差派來的使者.
上帝是他的主人.
他是效忠於基督耶穌.
但同時間在人裡面.
他就是用一種比較平易近人的方式.
他不是用一個大D的形式.
來跟教會講說話的.
他是用一種很平凡.
用一種僕人的心態.
來將這封書卷向領袖姐妹文安.
他亦用一個這樣的身份來自稱.
藉此來建立領袖姐妹.
跟她一同來同步領受書卷的信息.
我想特別注意的是.
第一章三字五節.
這是第一個段落.
「我每逢想念你們.
每逢為你們眾人祈求的時候.
嘗試歡歡喜喜地祈求.

$^{201}$因為從頭一天直到如今.
你們是同心合二地.
慶王復陰」.
在這裡你會看到保羅.
寫這封書卷的時候.
他如何凝聚一群.
菲律賓教會領袖姐妹的動力.
除了他將自己的身份.
放到一個謙卑的位置外.
他第一件事就是訴諸於情.
他用一種很深情的方式.
來表達他對領袖姐妹之間.
那份很緊密的關係.
你留意一下那個字眼.
是情書.
可能你寫給你男女朋友.
老公老婆你會先寫這些字.
「我每逢想念你們」.
一般來說.
你見保羅平時在不同的書卷裡面.
很堅強的.
他走出來就有一種.
仕途的風範.
他要辯明真理.
不過對於菲律賓教會來說.
他用一個很柔情的表達.
「想念你」.
「想念」這個字.
epiphoto.
這個字是希臘文.
我讀到嘰嘰嘰.
不過這個字在原文裡面.
帶含著一種很深的情感.
是一些生死之交.
他寫信去表明那份關係.
就會用「想念」.
在原文裡面用這個字.
保羅說到他和鼎姐妹之間.
有一種很深厚的情.
這份情是由起始到如今.

$^{241}$繼續建立起來.
他帶著這份情.
來為鼎姐妹感恩和祈求.
他的感恩和祈求.
是一種怎樣的祈求呢?.
是歡歡喜喜地.
為他們祈求.
你想想他每一次合十雙手.
閉上雙眼想起這班.
菲臘比的鼎姐妹.
他心裡面是永不息的那種歡樂和快樂.
他想起他們的時候.
心裡面是甜蜜的.
雖然保羅寫這本菲臘比書的時候.
他是坐牢的.
但當他在一些很困難的處境裡面.
他想起這班人.
想起這個群體.
心裡面經常會逢.
這個字.
即是說沒有例外的.
每一次都是這樣.
想起這班人.
心裡面有一種永不息的感恩和快樂.
為何保羅會有一種這麼強烈的感恩.
為何保羅會有一種這樣的心態呢?.
在經文裡面他繼續說.
第五節.
因為從頭一天直到如今.
你們是同心合意地慶王復陰.
從頭一天.
在說哪一天呢?.
《史獨行傳》第十六章十二節就說到.
這個教會是如何建立的.
我讀給大家聽那一節經文.
從那裡來到菲臘比.
就是馬其頓這一方的頭一個城.
也是羅馬的駐防城.
我們在這城住了幾天.
菲臘比教會.

$^{281}$是由《史獨行傳》裡面記載著.
就是馬其頓的呼聲.
保羅領受了馬其頓.
聖靈的帶領之後.
就由福音亞洲傳到歐洲.
而菲臘比就是歐洲的第一站.
在這一站裡面.
保羅經過一個比較駐防的城市.
他帶了一些婦女.
女底亞那些.
尊貴的婦女來信主.
這個教會就由那天開始.
開始一路一路的被建立.
保羅在他們當中裡面來講論.
建立了一班信徒.
這班弟兄姊妹.
很努力地追求上帝.
而當保羅寫這個《菲臘比書》.
由第一天建立教會到寫書的日子.
有些色經學家講.
大概有十一年的時間.
這十一年裡面無間斷.
保羅都是看見這個群體.
不斷不斷地將福音傳開它.
雖然是一個很不容易的處境.
在《菲臘比》裡面都會面對一些逼迫.
將一個外邦的.
對於那些歐洲人來講.
是一個外邦的神.
帶到他們自己的地方裡面.
很多的抗拒.
很多人誤解基督教的信仰.
不過這班人不會因為這樣的緣故.
就放棄福音.
保羅為著這班弟兄姊妹的努力.
為著這班弟兄姊妹所擺上的那種心.
他心裡面經常的.
都是永留不息的那種感恩.
另外還有一個很特別的.
在《菲臘比書》裡面有些很重要的字眼.

$^{321}$很多節.
所以我快點講給大家聽.
一章二十七和四章三節.
那裡講的是齊心努力.
二章十七和十八節.
講到一同喜樂.
然後三章十節.
講到所受的苦.
大家有份.
然後四章十四節.
一同分擔.
在這一堆的動詞裡面.
保羅一定會加了一個字.
就是SYN.
就是英文Synergy.
SYN就是這個字.
這個SYN的意思就是大家一起.
無論努力.
快樂.
痛苦.
和所有的重擔.
都是一起的.
在所有的動詞前面.
保羅刻意就要講到.
他不是孤軍作戰.
這個教會也不是自己.
來經營他們自己的信仰生活.
在不同的地域裡面.
保羅和這一班的弟兄姊妹是同心的.
除了動詞之外.
一章七節講到分享.
二章二節講到心志相同.
二章二十五節到四章三節.
講到一同作功.
二章二十五節講到一起作戰.
三章十七節講到一同效法保羅.
三章二十一節.
和他一起.
與他相似.
這些名詞同樣都是講到.

$^{361}$都是加上SYN.
或者希臘文宣.
講的就是.
這個群體和保羅.
在生命所有的場景當中.
特別在信仰.
在福音使命上.
這一班人一起來彼此努力.
雖然可能保羅接下來.
去到巡迴不同地方傳福音.
但是保羅收到的消息.
或者保羅的觀察知道.
這一班弟兄姊妹.
他不會因為保羅不在他們當中.
他們就停止了福音的動力.
反而在不同的崗位裡面.
大家一起來配搭.
一起來彼此建立.
一起來努力.
為上帝來建立他們生命的榜樣.
在這裡我們稍為停一停.
來思考一下.
今天我們的弟兄姊妹關係.
建基的基礎是什麼呢?.
有很多弟兄姊妹經常都會說.
牧師我們要增加教會裡面.
很多的溫馨感受.
要令教會的木印多一些.
我同意的.
不過如果我們只不過是.
將教會成為一個純粹關懷的機構.
我們其實走到某些位置就會覺得瓶頸.
因為聖經裡面所說的那種弟兄姊妹的關係.
從來不是純粹說.
弟兄姊妹互相關懷就足夠.
那份關懷是要推己及人的.
這份關係是要不停.
藉著福音的使命來拓展開它.
所以如果我們要關係更加密切的時候.
我們更加需要有一份使命感.

$^{401}$一份來自聖經的福音使命感.
有這份使命感的時候.
我們聊天的話題.
我們的眼界才可以切磋.
很人性化地說.
每個星期回來讓你關心.
關心到差不多你都不知道關心什麼.
查家宅什麼都查了.
家中的原生家庭.
幾代原生家庭都查了.
到最後都是有限的.
不過如果我們是一起來參與福音工作.
一起來為上帝打拼的時候.
我們的故事就精彩得多了.
有時候出去吃檸檬.
有時候在福音的過程中遇到一些難關.
或者在我們人生裡面實踐上.
或者在福音的照明裡面.
我們有很多的掙扎.
我們生命裡面因為有外在比我們關係更重要的東西在.
所以當我們彼此集結在一起的時候.
我們生命就會有很多火花.
我們的關係就會變得更加牢固.
我相信很多弟兄姊妹都經歷過.
去短孫.
或者去到不同地方宣教.
最精彩是什麼地方呢?.
除了帶領當地人信主之外.
而是回來的那一程飛機或者那一程車.
經歷過這麼多事情之後.
大家將那個經驗在車裡面一起聊天.
或者怎樣將自己的生命的經歷融入去服務的群體裡面.
大家彼此交流.
那個聊天就到位很多.
深度有很多.
對我們來說生命裡面更加難忘.
保羅在這裡面.
為什麼他和這班弟兄姊妹有這麼強的凝聚力呢?.
福音使命.
因著福音使命.

$^{441}$將這班人拉在一起.
因著福音使命.
讓帶領信主的保羅每逢想起他們.
心裡面就會歡喜快樂.
心裡面念念不忘這班群體.
這班人就一齊走上福音的道路.
另外一個字眼可能需要補充一點.
一章第五節.
因為從頭一天直到如今你們都是同心合意地興旺福音.
本身興旺福音這個字很好聽的.
我們經常讀的.
不過本身興旺這個字就不是很大的重點.
興旺的意思是一齊傳揚福音.
比較好的翻譯是在福音的事上大家一起有份.
他講到的就是有份參與.
所以強調的是同心合意.
多於興旺福音的興旺這個字.
有時走在一起.
可能大家都是集體吃檸檬的.
不過集體吃檸檬也好.
保羅說在大家一起去傳揚福音的過程裡.
大家一起同心一起努力去實踐.
丁姐妹,我不知道你選擇參與在YouTube裡面看崇拜.
你會為了甚麼.
我們很誠意,旺角浸信會很誠意告訴大家.
我盼望我們的關係不是純粹對著攝影機.
純粹聽聽好聽的講道.
或者不同的講員分享聖經的話語.
這是重要的.
不過更重要的是我們這班基督徒.
不要只滿足在家裡看著這部電視,YouTube.
我們需要集合在一起.
我們需要實體的去建立彼此的關係.
我們生命裡面有這個福音使命.
我們才可以為這個世界.
為這個社會帶來一些改變.
純粹被動接收信息.
改變不了生命.
你可能腦裡面聽了很多信息.
但是集結在一起.

$^{481}$成為更新.
成為一個有使命的團隊.
這個信仰生活.
比我們純粹看電視,看YouTube精彩得多.
我希望每一位我們一起彼此努力.
我們一起成為一個集結起來的有使命的群體.
我在這裡借題發揮一些經歷.
早陣子有些同道跟我聊天.
他說旺朕現在有四,五百人.
你牧養四,五百人會不會很吃力?牧師.
跟我聊天.
坊間很多人都會覺得.
一間教會有幾百人,幾千人,幾百人.
甚至上萬人.
是一個很吃力的做法.
有些弟兄姊妹很關心地問.
你會不會掏空自己?.
你要嘔心瀝血地牧養弟兄姊妹.
一個人的需要已經不是少了.
幾百人的需要.
來到牧養就很痛苦.
他用這個角度來看.
面對這些問題.
很多人都問過我.
通常我就笑笑的說.
怎會痛苦?.
Facebook幾百人.
我都覺得嫌少.
更何況.
這班不是我掏空自己.
給她食物的羊.
我不用用我的血肉來餵他們.
他們懂得自己吃.
自己謀生.
我看這班弟兄姊妹.
是我的朋友.
這班.
是我生命裡面.
不可或缺的群體.
因為我有這幾百人.

$^{521}$所以我們事奉更加有力量.
因為這幾百人.
是我生命裡面.
可以成為一生中生的朋友.
我生命裡面.
我覺得我不孤單.
我們.
每一位弟兄姊妹在教會裡面.
不是純粹滿足了.
那個宗教場景.
大家一起聚集的敬拜.
我們之間還有一份.
很深厚的情.
這份情緊緊地維繫著我們.
我們的生命.
因為有這一份情.
是繫上上帝的恩典.
所以我們這班人.
在我們每一次的聚會.
每一次的聚集.
我們不會覺得累.
我們每次見面是充電的.
那些牧者或其他人問我.
聽到我這個答案.
他們都很詫異.
他們有個結論.
Anderson難怪你的樣子這麼年輕.
你每天都好像活在天堂一樣.
不過我是真的.
在旺鎮裡面.
去不到天堂的這個層次.
但我們都不要將它成為地獄.
我們這個地方.
是一個近似天堂的寓言.
我們弟兄姊妹.
不同的人聚集在一起.
不是要掏空我們的心靈.
而是我們彼此豐富.
我們這一班人是帶著使命.
一起來去分享生命.

$^{561}$喜與樂.
不同的擔子.
我們共同來承擔.
本來以使命來建立這一班人.
他亦都訴諸於情.
來講這一班的弟兄姊妹.
是他們一起打仗的一班弟兄姊妹.
所以保羅講第一個的動力.
保羅講的凝聚力.
就是一個有使命感.
有情的一份的關係.
好,我們再看第二點.
在第一章第六至第八節.
我深信那在你們心裡動了善功.
必成全者功.
直到今生的最後.
我心中有信.
必成全者功.
直到基督耶穌的日子.
我為你們眾人有這樣的念念.
是應當的.
因你們常在我心裡.
無論我在困所之中.
是變命福音的時候.
你們都與我一同得恩.
我體會基督耶穌的心腸.
切切地想念你們眾人.
這是神可以給我作見證的.
同樣都是.
豐富情感的一個小段落.
你會見到.
保羅.
圍著這一班弟兄姊妹.
他心裡.
想起的就是一班弟兄姊妹.
那種努力.
他知道.
這一班弟兄姊妹心靈裡.
有上帝與他們同在.
所以他們是自發地.

$^{601}$來實踐上帝要他們.
所作的善功.
是上帝要他們所預備.
他們所踐行的福音使命.
這些工作不是.
保羅要威逼.
來有他們的.
而是上帝在每一個信徒的心靈裡.
喚發起他們.
自發地來實踐.
福音召命.
所以保羅說.
這一班人是上帝在.
這個群體裡所做的工作.
令到這一班人也很有動力地.
自動自覺地.
來實踐福音召命.
更加重要的就是.
這一班人這麼努力.
實踐的時候.
保羅說他去到什麼狀態.
去坐牢.
或者.
至高氣揚.
帶著上帝的福音.
來辨明福音的時候.
他也知道.
身處不同地方.
但是他們的心是在一起.
在將來得贖的日子.
在基督將來的日子.
這一班人和保羅.
將要一同得著恩典.
一同得恩.
保羅帶著一份確信.
這一班人.
和他一起領受.
同一個的恩典.
不單止在今世.
也在將來基督耶穌的日子.

$^{641}$他也會一起來領受.
這麼美好的恩典.
我們今天我們弟兄姊妹.
你們有沒有深信呢?.
可能你以前坐在教會.
或者你的小組的弟兄姊妹.
將來你會和她一起上天.
一同得著上帝的恩典.
當你知道.
這個人會陪你.
一起去到.
上帝所預備美好的.
那種狀態的時候.
我們心裡面就會.
少一些.
有時候人與人之間.
那種性格不合的.
那些矛盾.
當你明白.
人生裡面.
我們覺得此時刻有些矛盾.
但是到了終極.
將來上帝也要救贖她.
也救贖你.
大家在天家裡面.
繼續得著美好的恩典的時候.
你就可以輕看眼前.
有一些矛盾和衝突.
在保羅眼中.
裡面就看見這一班弟兄姊妹.
你猜這班人是不是完美到.
沒有瑕疵呢?.
在聖經裡面.
說不是的.
有些人要氣死保羅.
有兩個姐妹.
有阿諜秦都基.
也經常吵架.
為了服侍的緣故.
這間教會跟我們地上的教會.

$^{681}$其實也很近似.
也會有些意見不合.
爭拗,矛盾.
不過在保羅眼中.
這些一切.
都是此時此刻短暫.
在將來得贖的日子裡.
這班弟兄姊妹.
會跟她一起得著.
永恆的恩典.
永生蓋過了眼前.
一些.
大家不同.
不咬弦的地方.
保羅再繼續說.
我體會基督耶穌的心腸.
切切地想念你們眾人.
這是神可以給我作見證的.
基督的心腸.
在說什麼呢?.
如果你看接下來的.
二章五節.
要以基督耶穌的心為心.
基督的心腸.
是一個整體.
將自己放下.
謙卑極致.
然後要建立的.
一個上帝的工作.
他用一種.
耶穌基督的心腸.
放下了自己.
放下自己覺得很重要的.
一些看法.
反而是他將自己的焦點.
放在一個.
上帝所喜悅的群體身上.
上帝藉著這一群人.
成就了上帝的美事.
藉著一些.

$^{721}$性格可能很多出入.
人生裡面可能有很多.
矛盾,衝突.
但那群人都是上帝所喜悅的.
那群人都是上帝所愛的.
這一群人.
和保羅一樣.
他將來一定得享永恆的恩典.
保羅帶著一份.
有一份胸襟.
以基督耶穌的心腸.
放下謙卑自己.
來從這個角度.
來看待一個群體.
這個群體就變得很可愛.
早前.
看過一個這樣的比喻.
他說人與人之間的相處.
你用顯微鏡看東西.
還是用望遠鏡.
如果你用顯微鏡.
來看東西.
如果你上過實驗室.
你看過了.
顯微鏡下有些細微的細菌.
放得很大.
你會覺得很恐怖,很噁心.
將你的皮膚放大一萬倍.
你看到很多病菌.
骯髒的東西.
你就會覺得.
很難受的.
你用顯微鏡.
看人際關係.
你去注視一個人的.
錯處.
你都會發現那個人是這樣的.
這麼奇怪.
三尖八角.
奇形怪狀.

$^{761}$但其實那個人是這樣看你的.
你用顯微鏡看你.
大家互相用顯微鏡看.
其實大家的人際關係是很痛苦的.
不過相反來說.
你用望遠鏡.
好像望天一樣來看一個人.
你就會看到晚上的星空.
繁星閃閃.
每一個星都會.
發出亮麗的光芒.
人際關係裡面都是一樣的.
如果我們用一個.
較為遠的角度來看.
我們身邊的人.
我們嘗試不要用顯微鏡.
來挑他的錯處.
而是用望遠鏡來看這個人.
生命裡面的美麗特質.
你會發現.
每一個弟兄姊妹生命都可以.
發光.
近看可能是千瘡百孔.
好像月球表面一樣.
但遠看就是一個.
明亮的月光.
發出光芒.
我們生命.
如果我們這樣看自己的群體.
我們弟兄姊妹.
我們能夠離遠一點.
抽遠一點來欣賞他.
發現他生命.
美善的地方.
你會發現你活在.
一個可愛的群體裡面.
相反來說.
你經常用顯微鏡來看.
你會覺得那個群體很難忍受.
很可悲.

$^{801}$你經常都會活得不開心.
你更加很難投入在那個群體裡面.
那個關鍵是.
你用顯微鏡或望遠鏡.
來看你身邊的人.
弟兄姊妹.
我很盼望我們旺朕的.
弟兄姊妹.
我們這個群體.
我們真的用望遠鏡來看一個人的生命.
我們用一個.
前景來看他.
我們嘗試在他生命裡面的光彩的地方.
在他生命的特質裡面.
發現他那種光芒.
人總會有一些雜質的.
不過上帝不會創造一件垃圾出來.
上帝是在人的軟弱當中.
呈現他榮耀的大能.
每一個人都能夠可以發揮上帝的光芒.
我們就帶著這一份心腸.
來找找弟兄姊妹.
生命裡面的美麗的那一面.
你這樣做的時候.
你每一天都會這樣為自己.
帶著這樣的角度來看人.
你會發現這個世界實在太美麗.
你身處的群體.
實在有太多的潛質和可能性.
你身處的群體實在有太多的潛質和可能性.
你做人也會有很多自然變得舒坦的地方.
保羅就是這樣看弟兄姊妹.
看他們每一個粒粒皆星.
雖然有不同的性格.
但上帝在這一群人裡面.
在他們的心靈裡面動了善功.
這一群人會和保羅一起.
來領受一同美好的恩典.
剛才說過.
有很美麗的人生.

$^{841}$有很多弟兄姊妹.
那是不是在菲律賓教會裡面.
或者在菲律賓這個城市裡面.
就無懈可擊呢?.
是不是一個差不多等於天堂的地步呢?.
我們看後一點的篇幅.
一章十五至十八節.
有的傳基督是出於嫉妒紛爭.
也有的是出於好意.
這一等是出於愛心.
知道我是為變名福音設立的.
這一等傳基督是出於結黨.
並不誠實.
意思要加增我.
群所的苦楚.
這又何妨呢?.
或是假意或是真心.
無論怎樣.
基督究竟被傳開了.
我便歡喜.
並且還要歡喜.
這一句經文我是很喜歡的.
保羅寫這句經文寫得很得意.
如果你用憂心忡忡看.
你體會不到比率.
保羅那種得意的心.
他說在教會.
有些人傳福音是基於.
增敬嫉妒.
可能為自己立馬.
這樣的角度.
這是保羅說的.
假意不是真心的.
是結黨的.
不是真誠將福音帶給身邊的人.
他說.
教會有這樣的人.
當然有些人是出於愛心.
去引證保羅所傳的福音.
是真確的.

$^{881}$這一班人和保羅.
遙距的彌代同步.
這一班人和保羅的心態一致.
愛心和假意.
兩者放在一起.
的時候.
就會成為一個黨派之爭.
又或者大家彼此之間.
有一種角力在這裡.
保羅再繼續說.
這一班出於假意的人傳福音.
為了什麼呢?.
就是要加增我.
困所的苦楚.
整句話如果用廣東話.
翻譯你會傳神很多.
這一班人出於結黨.
傳福音都不是禮真的.
意思是什麼呢?.
就是要令到我坐牢的時候.
都要氣死我.
保羅是這樣說的.
然後保羅再說.
那又怎樣呢?.
假意也好,真心也好.
無論怎樣,基督都已經被傳開.
而且我開心都來不及了.
保羅其實用了.
這樣的角度.
來回應那些.
帶著假意傳福音.
他也是在傳福音的.
不過心裡面有一些不是很純正動機的人.
對於這些人.
保羅.
想著這一班人可能會氣死他.
但他們一輩子都想不到.
基督究竟被傳開.
保羅為此.
都開心得來不及了.

$^{921}$在這裡.
裡面說的是什麼呢?.
在一個凝聚的群體裡面.
我們沒有辦法.
看到弟兄姊妹的動機.
即使我們.
看見也好.
我們要有一個更加高的.
層次.
來承載.
可能生命裡面.
信仰裡面的雜質.
在保羅的角度來說.
這個群體.
只要福音被傳開.
他就穩贏了.
贏定了.
他帶著一種贏定的心態.
來看即使.
有些人真心也好.
假意也好.
福音被傳開.
對我們來說.
他的心已經.
開心得來不及了.
今天我們很多時候.
弟兄姊妹有一種.
教會也好.
人與人之間的關係也好.
我們很多時候很著重.
我的看法或感受會不會被照顧.
如果我們.
越強調這一點.
在信仰裡面我們就越不開心.
因為這個世界.
總會有人會衝撞到我的情感.
有人會衝擊我的面子.
甚至有時候.
有些說話可能.
不太合聽.

$^{961}$不過保羅的秘訣是.
只要基督被傳開.
他就為此而快樂.
我們心裡面.
有時候.
人與人之間.
可能他未必有惡意衝撞自己.
又或者有些人.
真的帶著假意要氣死你.
不過.
我們要從更高層次去看.
如果他所做的事.
是將福音傳開.
他所做的.
好像增敬一樣.
但同時也帶著福音的果效.
我們就嘗試.
包攬這些.
有些雜質的.
一些努力.
在我們心靈裡.
也不會為這些雜質.
忿忿不平.
反而讓我們的心胸.
來擴闊自己.
我們心裡面.
不要太過介意.
別人給不給自己面子.
不要太過介意.
自己的意見是否被接納.
不要太過介意.
自己的感受.
快樂是否被別人.
所重視.
如果我們生命裡有一些東西.
比我們剛才所說的感受,面子更重要.
就是福音工作.
就是福音.
基督是否被傳開.
我們生命裡就變得輕散一些.

$^{1001}$在我們的人際關係裡.
我們有一個比自己.
個人更重要的東西.
來超越.
我們就.
沒那麼容易在人際關係中.
受傷.
求主繼續幫助我們.
我們生命裡需要一個.
更加崇高的福音照明.
照管著我們的心靈.
這個福音照明不單只是.
我們要傳的訊息.
也是操練我們內心.
以及幫助自己變得更加.
尋口踏實.
以致我們不看自己太過重要.
我們人生走的路.
就會走得輕散一些.
好,我們最後一個小段落.
一章九至十一節.
我所禱告的就是你們的愛心.
在知識和各樣的見識上.
多而又多.
使你們能分別是非.
作誠實無過的人.
直到基督的日子.
並靠著耶穌基督結滿了人意的果子.
叫榮耀稱讚歸與神.
保羅在這個小段落裡.
他用一個禱告作為.
這個小段落的結束.
他的禱告裡.
有幾個字眼需要特別的.
來讓大家知道.
第一個就是.
愛心在知識上多而又多.
愛心,待會再解釋.
在這裡說的.
知識的字是代表著什麼呢?.

$^{1041}$是對真理.
對神的認識.
在哥林多前書裡說的.
知識會使人漸漸更新.
在保羅書信裡用的.
知識是說真理性的認識.
所以越認識神.
越體會到自己的限制.
也體會到.
上帝那種永恆的心意.
這個知識.
這個知識就是在說.
對於聖經.
對於真理的那種認知.
那種透徹的理解.
保羅在說愛心在.
這個神學.
或者在.
對真理的掌握裡.
一路一路豐富起來.
這是第一個層面.
第二個就是各樣的見識.
這個各樣的見識.
更加好一點的翻譯就是.
對事物的.
分辨能力.
這個是在說判決.
審斷生活的.
那種能力.
有些的聖經學家在說.
更加多的實踐智慧.
這個見識就是在說.
實踐真理的時候.
會審時度勢.
會看看這個社會這個世界發生什麼事.
將真理.
到位的實踐.
所以保羅.
他為.
肥納比教會的祈禱.

$^{1081}$其實是很整全的.
他不是單方面說對神認識.
拿著聖經.
就掌握了整件事.
保羅他的祈禱是.
既是對神的認識加深.
同時間他也有那種.
實踐聖經對社會.
對世界裡的.
那種見識分辨的能力.
兩者加起來.
這就成為了他的愛心.
實踐的基礎.
如果沒有這兩樣東西的時候.
愛心很多時候純粹對你好.
區域混亂等等.
但如果當有.
知識和見識.
混合在一起的愛心.
那個造就性就大很多.
保羅的祈願是.
這一班的弟兄姊妹.
他有一份很深厚的情.
很深的愛.
但他的愛不是純粹.
表面.
或者層次.
要一路一路提升.
既是有實踐的智慧.
更加是有.
豐富來自認識神的內涵.
兩者合在一起.
就成為了一個很重要的幫助.
這兩者再加起來.
還有一個能力.
就是使你們能夠.
分別是非.
分別是非這個字.
更加好一點.
我翻譯一下.

$^{1121}$解釋就是.
認清了是與非.
還有對與錯.
黑與白.
美與善.
好與更好.
當愛心能夠有.
剛才說的知識和見識.
的融合之後.
令到這一班人.
可以有更加豐富的分辨能力.
有些好的東西.
有些愛心.
到位地去實踐.
如果.
有一些做扶貧機構的.
等於姐妹說.
一味送東西給他.
就等於對基層的人.
或者貧窮的人好.
我們要顧及他的尊嚴.
顧及他.
心裡面的感受.
亦都給他有上進的動力.
整全地來考慮一個人的需要.
你用這個角度.
來幫助一個人.
那你的分辨能力就強.
你的愛心就更加有根有機.
保羅就是一個這樣的角度.
來幫助眾弟姐妹.
最後這一班人.
在一個這樣的氛圍之下.
最終會結出美好的人意果子.
生命裡面被上帝稱為義.
他本質上.
當他信耶穌.
就是成為神的兒女.
被稱為義.
但在他生命裡面.

$^{1161}$再不斷地結滿人意的果子.
讓這種義成為真實.
生命裡面.
封實的.
可以觀察得到的.
經歷得到的這些果子.
我盼望我們.
黃浸.
我們要講六個G裡面的價值觀.
我們的愛心.
是真的有果子結出來.
我們的生命.
是真的有果子結出來.
生命.
我們這班人的聚集.
不是純粹講這班人很溫馨.
而是一個溫馨的.
群體裡面.
我們可以幫助到更加多的人.
我們有向度地.
來實踐愛心.
我們這個群體.
既是內聚也是外向.
我們可以讓這個世界知道.
基督徒是一班.
可以生命裡面這麼美麗的一班人.
今天這篇訊息.
不難理解的.
是不是?弟兄姊妹.
在這第五個G.
get closer.
我們怎樣和基督耶穌.
走得更加近和教會走得更加近呢?.
今天保羅教會我們.
首先我們要有.
強力.
有.
很重要的使命感.
放在我們前面.
這份使命感.

$^{1201}$會左右我們繼續怎樣.
前進.
第二件事.
我們這個群體的凝聚力.
挑戰我們每一個弟兄姊妹.
我們看身邊的人.
是不是可愛的呢?.
我們嘗試用望遠鏡.
看看每一個人生命發出的光輝.
讓我們.
懂得去發掘.
欣賞身邊的人.
第三,在我們個人來說.
輕看自己的感受.
我們用一個.
高的層次來看.
我們能夠有一個.
胸襟包容一些.
可能衝正我的.
或者和我們心態有些不同.
不過你看到.
他是有福音果效的人.
我們就帶著這個框架.
來和別人相處.
我們沒有辦法將所有和自己.
一色一樣走在一起的人.
聚集在一起,這個世界沒有可能有.
這樣的事,但我們有一個.
胸襟,有一個框框.
承載著.
這些不同式樣的人.
我們生命就變得豐富.
我們有一份穩贏的心態.
因為基督就已經被傳開.
最後,我們更加有.
成熟的愛心.
我們懂得在知識和.
見識上多而又多.
地被建立起來.
我們這個群體,我相信.

$^{1241}$我們會充滿前景.
我們生命會被.
上帝大大使用.
繼續激勵我們.
記住第五個G.
Get Closer.
\newpage



\section{箴言 27:5-17}
\label{sec:Q_GbM3l6t1s}
\textbf{2022年2月13日崇拜直播｜陳冠成牧師｜我們的6G - 真誠砥礪 Genuine Relationship｜箴言27\_5-17}
\newline
\newline
連結: \href{https://youtube.com/watch?v=Q_GbM3l6t1s}{\texttt{https://youtube.com/watch?v=Q\_GbM3l6t1s}} ~~~~ 語音日期: 2022-02-14
\newline
\newline
\hyperref[sec:XQBczBTx6j0]{\small{< < < PREV SERMON < < <}}
~
\hyperref[sec:index_chronic]{\small{[返順時目]}}
~
\hyperref[sec:yigtcB2OiSE]{\small{> > > NEXT SERMON > > >}}
\newline
\newline
箴言 27:5-17
\newline
\begin{longtable}{cl}
\hline
\hline
章節 & 經文 (和合本修訂版)\\
\hline
27:5 & \begin{tabularx}{0.7\textwidth}{X} 當面的責備勝過隱藏的愛情。 \end{tabularx} \\ \\ \relax
27:6 & \begin{tabularx}{0.7\textwidth}{X} 朋友加的傷痕出於忠誠；敵人的親吻卻是多餘。 \end{tabularx} \\ \\ \relax
27:7 & \begin{tabularx}{0.7\textwidth}{X} 人吃飽了，厭惡蜂房的蜜；人飢餓了，一切苦物都覺甘甜。 \end{tabularx} \\ \\ \relax
27:8 & \begin{tabularx}{0.7\textwidth}{X} 人離故鄉漂泊，就像雀鳥離窩四處飛翔。 \end{tabularx} \\ \\ \relax
27:9 & \begin{tabularx}{0.7\textwidth}{X} 膏油與香料使人心喜悅，朋友誠心的勸勉也是如此甘美。 \end{tabularx} \\ \\ \relax
27:10 & \begin{tabularx}{0.7\textwidth}{X} 你的朋友和父親的朋友，你都不可離棄。你遭難時，不要上兄弟的家去；相近的鄰舍強如遠方的兄弟。 \end{tabularx} \\ \\ \relax
27:11 & \begin{tabularx}{0.7\textwidth}{X} 我兒啊，你要做智慧人，好叫我的心歡喜，使我可以回答那辱罵我的人。 \end{tabularx} \\ \\ \relax
27:12 & \begin{tabularx}{0.7\textwidth}{X} 通達人見禍就藏躲；愚蒙人卻前往受害。 \end{tabularx} \\ \\ \relax
27:13 & \begin{tabularx}{0.7\textwidth}{X} 誰為陌生人擔保，就拿誰的衣服；誰為外邦女子作保，誰就要承當。 \end{tabularx} \\ \\ \relax
27:14 & \begin{tabularx}{0.7\textwidth}{X} 清晨起來大聲給朋友祝福的，就算是詛咒他。 \end{tabularx} \\ \\ \relax
27:15 & \begin{tabularx}{0.7\textwidth}{X} 下雨天連連滴漏，好爭吵的婦人就像這樣； \end{tabularx} \\ \\ \relax
27:16 & \begin{tabularx}{0.7\textwidth}{X} 攔阻她的，就是攔阻風，又像用右手抓油。 \end{tabularx} \\ \\ \relax
27:17 & \begin{tabularx}{0.7\textwidth}{X} 以鐵磨鐵，越磨越利，朋友當面琢磨，也是如此。 \end{tabularx} \\ \\
[1ex]
\hline
\hline
\end{longtable}
$^{1}$讓我們同心以聆聽聖道去敬拜我們的上帝.
今日要讀出的經訓是在.
《箴言》二十七章五至六節.
九節和十七節.
《箴言》二十七章第五節.
多面的責備強如背地的愛情.
朋友家的傷痕出於忠誠.
仇敵連連親嘴卻是多餘.
第九節.
高柔與香料使人心喜悅.
朋友誠實的勸教也是如此甘美.
第十七節.
鐵磨鐵磨出人來.
朋友傷感也是如此.
今日來到我們的6G的第六講.
《真誠抵禮》.
我們請陳冠誠牧師.
謝謝梅惠平安.
正如剛才珍珍牧師所說.
來到我們這個系列最後一講了.
真誠關係.
說的是信徒之間的真誠關係.
到底是如何建立和維繫.
其實真誠關係.
苦心自問.
我們每一個對於真誠關係的課題.
很大機會是學了一輩子的課程.
畢竟當我們越來越與其他信徒.
近身相處,事奉的時候.
其實你會看到不單單是自己的不足.
自己的甩漏.
其實你也會看到別人的弱項.
甚至乎這些的錯處.
有時候我想每個人都曾經試過這種掙扎.
就是你很想與對方說.
將你所見到.
將你所發現對方一些問題說出來.
很希望對方知道並且幫他改善.
但你又會擔心可能會絕口不屑.
擔心造就不到人之餘.

$^{41}$反而得罪了對方.
破壞了彼此的關係.
於是你就會選擇不出聲.
還是不說好了.
又或者另一種情況就是.
有人真心真的想你好.
他很勇敢地說出你的錯處.
說出你的一些弱項.
其實是一件好事.
不過很自然我們有時候又會因為自尊心作祟.
想保護自己不想面對自己的挫敗.
於是乎又會不知不覺地用一些逃避.
或者辯駁的方式.
來回應對方一種善意的提醒.
甚至乎有時候會令到面左左不歡而散.
但之後你又會很後悔.
很憤怒自己.
其實為何明明是一件好事.
為何會搞成這樣.
頂尖的我想.
不知這個是否你和我都會遇到.
在人際關係相處裡面的一些掙扎.
但無論如何.
當上帝在暢世紀裡面說到.
「那人獨居不好」的時候.
其實就已經表明我們不是生活在一個孤島裡面.
我們是一定需要透過.
和其他人其他的基督徒來互動.
讓我們的屬靈生命可以成長.
甚至乎邁向成熟.
使我們作為耶穌基督身體的一部分.
可以聯絡得更加通透.
更加緊密.
當然我們都很明白.
在現實生活裡面.
兩個罪人走在一起.
不會自動一拍即合.
是需要一些磨合的過程.
這樣才能夠在事奉的路上.
可以配搭發揮得更好.

$^{81}$就好像今天這個講題裡面.
其中講到「禮禮」這兩個字.
「禮」是指那塊磨刀石.
幼細的那一面.
至於「禮」是指粗糙的磨刀石.
其實你會看到上帝在我們的生命裡面.
透過和不同弟兄姊妹.
不同性情的人來磨合.
來磨練.
練淨我們每一個人的生命.
簡單就是上帝有時會用一些粗糙的.
好像沙子那樣.
省走我們一些缺點.
又或者用一些很幼細的沙子.
將我們美善優雅的一面呈現出來.
其實我以前在工作時.
都累積了一些打磨的經驗.
就是有時候我要將一件物件.
按不同的步驟技巧來將它磨得仔細.
不能夠隨便的.
譬如最簡單就是.
你要去轉換不同粗幼的沙子.
技巧上來說.
可能你最初由200號轉到400號,600號.
號數越大代表它越幼細.
甚至差不多到尾的時候.
你又要用拋光布來拋到它亮亮的.
你需要有耐性.
又需要有愛心.
這樣你磨出來的東西才能磨得亮麗.
磨得幼細.
另一方面就是.
在你打磨的過程裡面.
有一件事是不能缺少的.
大家知不知道是什麼?.
沒有了這件事.
你很難磨到幼細的東西.
就是水或者打磨劑.
令到你打磨的時候.
那件部件不會吸住.

$^{121}$亦不會過熱.
更加可以幫你磨得越來越幼細.
同樣要打磨一段真誠的關係.
我們都必須要有好的打磨劑.
就是上帝的恩典和他的話語.
正如剛才所讀的《箴言》裡所說.
「當面的責備強於背地的愛情.
朋友家的傷痕出於真誠.
受敵連連親嘴卻是多餘」.
其實裡面的意思.
作者用一個對比的方法.
來表達.
原來在一段真誠的關係裡面.
有些時候我們寧願選擇當面責備.
總比我們保持緘默是有益的.
確實.
說到責備.
我想沒有什麼人喜歡.
無論你是責備人的那個.
又或者被人責備的那個.
沒有什麼人會想見到這兩個字.
相比起責備.
其實愛情就多人受落很多.
即使是虛假的愛情.
都仍然有人喜歡去聽.
但是《箴言》的作者這裡.
很清楚提醒我們.
當面的責備強於背地的愛情.
原因是有時朋友對我們當面責備.
往往是出於.
他對我們的忠誠和對我們的愛.
沒錯.
可能會令我們覺得.
有時很沒人性.
被人罵的時候.
會覺得很傷感情.
有時甚至乎很難堪.
不過《箴言》這裡說.
如果真是你的好朋友.
真是愛你的弟兄姊妹.

$^{161}$他們真的有時寧願冒著.
破壞關係的危險.
直言不諱地.
指出我們生命的一些錯誤.
來給我們一些善意的忠告.
希望幫我們回到上帝的心意裡.
其實我想這個道理不難理解.
特別是作為父母或子女.
你很明白.
就是做父母的時候.
當然不希望盡量減少.
當面責備子女.
因為真是會破壞關係.
有時見到可能子女激氣的時候.
你寧願省一口氣.
都不想來到大家面左左.
大家相安無事.
那就輕鬆自在.
我們確實有時真是會這樣想.
但是有些時候你明知道.
他一錯再錯的時候.
你因為愛的緣故.
你因為緊張子女成長的緣故.
你又真是會冒這個風險.
寧願坦白告訴他.
他錯處是甚麼.
同樣正如箴言這裡所講.
如果朋友.
我們的弟兄姊妹.
看到我們一錯再錯.
但卻是畏首畏尾.
寧願殉情顧著我們的面子.
不想得罪我們的話.
其實箴言這裡講.
和仇敵年年親嘴一樣.
對我們其實沒有甚麼好處.
換句話來說.
在一段真誠的關係裡.
原來不能單單靠著妥協.
來維持.

$^{201}$有時愛一個人.
其實無可避免.
你要勇敢地.
指正對方一些盲點.
因為唯有這樣.
你才有機會糾正對方.
避免他繼續行差踏錯.
其實聖經裡有很多經典的例子.
特別你會最深刻印象就是.
在《創世記》裡.
當蛇引誘夏娃的時候.
我們常常會問一個問題.
就是當時阿當去了哪裡?.
原來你看看《創世記》裡.
作者有暗示.
其實阿當一直都在場.
因為當蛇和夏娃對話的時候.
一直是用「你們」一個眾數.
當時的世界只有兩個人.
夏當和夏娃沒有第二個.
而夏娃的回答亦是「我們」.
而另一方面.
在中文和學本裡的翻譯是.
我們常常會看到一個圖畫.
是夏娃先吃了分辨善惡樹的果子.
然後才遞給阿當吃.
但你留意其實原文的意思其實是.
當夏娃在吃果子的時候.
阿當是在她旁邊.
和她一起.
也就是說《創世記》裡面.
其實一直都在暗示.
阿當由頭到尾都在場.
只不過他選擇不表態.
他沒有責備蛇之餘.
亦沒有阻止夏娃犯錯.
以致我們很明白.
罪就是這樣進入了世界.
所以你首先看到.
上帝首先向誰問責?.

$^{241}$是向阿當,不是夏娃.
因為他選擇了對罪保持妥協和緘默.
他沒有順服上帝交託給他的吩咐.
確實我們明白的.
在生活裡面有時.
沉默是禁.
沉默亦是人生一種智慧.
但假如在一段關係裡面.
總是對對方的過錯保持沉默的話.
不願意去指正篤責.
其實是否真的對大家好呢?.
這正正是箴言這裡所提醒我們.
以我們思考的問題.
我曾經聽過一句很有意思的話.
當你珍惜一件事的時候.
其實就代表你要放下另一件事.
當你珍惜某件事物.
其實意味著你要對另一件事放手.
套入我們今天的內容裡面.
當你真是珍惜同人一段關係的時候.
代表有時你要放下.
你總是做好人.
總是遷就,總是妥協的角色.
畢竟其實在你生命裡面.
上帝是託付了一些人.
讓你去照顧,守望和提醒.
上帝是要你放膽.
來向他愛心說誠實話.
來對他說.
朋友,這樣做真的不對.
真的不合上帝心意.
不要繼續這樣.
願意上帝幫你改變.
上帝一定會.
我們都願意守望你,跟你同行.
當你珍惜一段關係的時候.
上帝要你這樣說.
而反過來就是.
假如我們珍惜一段關係的時候.
亦都意味著我們自己.

$^{281}$都要放下一些面子和尊嚴.
來虛心聽下.
你身邊的朋友,弟兄姊妹.
他們所說的真心話.
確實,始終我們很明白.
做人難免真的會犯錯.
所以很多情況在裡面.
上帝都會透過別人.
來指出我們的弱點.
指出我們一些問題.
讓我們可以有修正,改善的機會.
我最記得就是在我讀神學的時候.
教導舊約的老師,他很好.
因為我們每個同學都要.
交功課,交論文給他.
他其實都會每次在功課裡面.
很希望給一些很.
真誠的評語給同學.
讓他們有所改進.
不過後來他說了一件事.
因為有些同學真的不敢.
面對別人的批評.
別人很正面的評價.
所以他學聰明了.
他給同學一個選擇.
就是如果你想老師去.
不只是給你分數.
你想有老師正面的評語.
你在你的功課的第一頁打顆星號.
如果你不打顆星代表.
你不想要評語.
我自己去想一件事.
確實面對別人赤裸裸地.
說自己的弱項是不舒服的.
但如果總是這樣過於膽怯的時候.
不敢去面對老師真的實話實說.
其實很難有改進.
唯有你放膽真的.
接受老師的愛心評論.
在你的功課第一頁打顆星的時候.

$^{321}$讓他可以誠實指出.
我們這份論文.
這份功課有什麼不足.
有什麼可以改善的地方的時候.
這樣才能學到東西.
所以頂指會真的.
我想假如有一天.
有人願意來到你的面前.
鼓起勇氣.
真誠地指正我們過錯的時候.
可能你真的會覺得難堪.
覺得心裡有點不舒服.
但請你謹記今天這裡的提醒.
朋友家的傷痕出於忠誠.
提醒自己不要太快生氣.
你嘗試換一個角度去看一看.
讓你有一位願意向你說真心話.
真的給一個心出來愛你的朋友.
並且嘗試鼓起勇氣.
聽一聽對方到底想說什麼.
這就是箴言給我們做人的智慧.
繼續看.
第九節這裡又說.
「高柔與香料使人心喜悅.
朋友誠實的勸教也是如此甘美」.
這裡說到高柔和香料.
高柔就是當時用油,香料.
或者乳香調製而成.
敷在痛處或傷口裡.
其實有震痛的作用.
能使皮膚柔軟得到醫治.
香料就很簡單.
就好像古代的化妝品香水.
塗在身上會散發香味.
令人心曠神怡.
箴言這裡說.
「忠誠朋友的勸勉.
就像高柔和香料一樣.
原來是可以安慰人的心靈.
帶來醫治的作用.

$^{361}$亦能使人生命有美化滋潤的功效.
讓人真的.
當他嘗到的時候.
會散發出一個喜悅的感覺」.
提醒我們.
在一段真誠的關係裡.
不單要我們愛心說誠實話.
亦要求我們用合宜的說話.
來規勸我們身邊的人.
就像高柔和香料一樣.
原來需要預先調製好.
調和好.
我們才能說出口.
讓人能夠聽到的時候.
更加受樂和受教.
否則我們都試過撞板.
明明是一番好意想對方好.
但可能因為說話的內容.
說話的語氣或時機.
不適當的時候.
很多時候你會發覺.
事與願違弄巧反拙.
不單止令對方誤會.
甚至會激怒對方.
所以其實.
說到勸說.
我們可以說是一門藝術.
你需要花點心思來設計.
想想怎樣說.
不可以想都不想.
單刀直入.
正如我們今天沒有讀到的另一句真言.
其實我們很熟悉的.
就是一句話說得合而.
就好像甚麼.
「金蘋果在銀網紙裡」.
其實它的意思是甚麼呢?.
就是說一句恰到好處的話.
就好像一個金蘋果.
放在一個銀盤上.

$^{401}$可以怎樣?.
互相輝映.
你想像一下.
如果金蘋果放在金盤.
或者銀蘋果放在銀盤.
有沒有這麼漂亮?.
沒有了.
因為不能夠互相輝映.
沒有這麼漂亮.
但金蘋果放在銀網紙.
放在銀盤裡.
其實是刻意精心設計.
一件高貴的藝術擺設.
讓每一個人看到的時候都會讚嘆.
我們不禁要問.
怎樣才算是說話恰到好處.
讓人去受落,受教呢?.
其實你翻開聖經裡面.
都給了我們幾個原則.
首先最簡單就是.
當你想勸說別人的時候.
你不要大庭廣眾說.
怎樣都要顧及對方的感受和面子.
你可以拉在一旁.
趁沒有人的時候.
找一個適當的機會.
跟對方去說.
讓對方在未聽到你的勸說之前.
他才收到你的善意.
他知道你其實都顧及他的面子感受.
所以才拉在一旁.
這樣你會發覺你的勸說.
通常會有效一點.
對方聽下去會受落一點.
即使他未必認同你的見解.
但先收了你的心意.
而另一方面就是.
說到勸說.
其實真的要精簡.
真的要點到即止.

$^{441}$所以你會發覺.
箴言有很多.
無論是父親對兒女.
或者長輩對後輩的勸說.
在箴言裡面你會看到.
很多是格言.
很精短.
很容易記,很容易明白.
一語中的就夠了.
千萬不要做我們所謂.
一個北風型的勸說者.
即是什麼意思呢?.
我想大家小時候都聽過.
太陽和北風的故事.
太陽和北風要做一個比賽.
看看誰可以有辦法.
讓路人脫去外衣.
北風於是就不斷吹.
打算吹走他.
結果我們都知道.
路人的外衣越抓越緊.
因為大風.
反過來太陽又怎樣?.
用溫暖,用熱力,用熱情.
讓路人覺得熱,溫暖.
所以自然很輕鬆地.
脫下外衣.
這裡說.
不要做北風型的勸說者.
是說什麼呢?.
提醒我們.
說話多不一定好.
我們不斷勸說對方的時候.
不一定有力量.
有人說得很好.
當你勸說人的時候.
就好像製作PowerPoint.
即是怎樣呢?.
Point.
你的重點.

$^{481}$要精簡,要到位.
這樣才有力量,有Power.
所以我曾經聽過一個故事.
話說已故的其中一位美國總統.
朗奴列根.
我想大家都還記得他.
他很喜歡和身邊的幕僚.
分享一個故事.
話說有個主日.
有一位年輕牧師.
去一間小教堂裡講道.
當他去到的時候.
他發覺整間教堂空空的.
只有一個會眾.
來聽道.
於是他就問那個來聽道的人.
只剩下你一個.
那我怎樣做好呢?.
還繼續講嗎?.
那個人就回答.
其實我是養牛的.
我今天.
來聽道.
你這個問題其實頗有意思.
如果我的農場裡.
只有一頭牛的話.
我都會照樣餵草給牠們吃.
這位年輕牧師聽到這番說話.
就覺得被激發起他的使命感.
對,就算一個人.
我都要努力去.
將上帝的話語傳遞給他.
於是他就開始講道.
並且越講越有力,越起勁.
他滔滔不絕地.
將他盡可能講到的東西.
都一次過講出來.
當他講完之後.
他又問這個來聽道的人.
你覺得聽完道之後感覺怎樣?.

$^{521}$那個人有少少的疲累.
有些疲態.
就跟他說.
其實我只不過是養牛.
其實我都不知道怎樣去回應你.
不過我想起在我日常生活裡.
如果我的農場裡.
只有一頭牛的話.
我一定不會將所有的草.
一次過餵給牠們吃.
我很喜歡這個故事.
因為其實很有啟發性.
確實.
如果我們的勸說太過冗長.
對方往往不單單只.
聽不及,吸收不到.
甚至可能會覺得怎樣?.
覺得你煩?.
覺得你囉嗦.
所以索性怎樣?.
生意仔?.
甚至會不停跟你唱反調.
對著幹.
所以其實.
禎賢這裡講得很好.
我們的勸說的內容.
是需要精心設計.
需要簡潔.
當然.
更加要有耐性.
要用溫柔的語氣.
來向對方講.
好像另一段經文.
禎賢第25章15節這裡說.
「恆常忍耐可以勸動君王.
柔和的舌頭能拆斷骨頭」.
其實他讓我們回到古代的一些場景.
我們會明白.
在古代的君王.
多數都是大權在握.

$^{561}$操生殺之權.
所以.
無論你是一個.
官員也好.
或者你是皇帝身邊的皇后.
或者王子也好.
你有時要去勸服一個君王.
真的不容易.
甚至乎真的會冒險.
因為始終你激到龍顏大怒的時候.
他真的有機會降旨.
罰你.
甚至乎.
取去你的性命.
而禎賢這裡提供一個.
當我們去勸勉君王的時候.
勸勉一個這麼有權威的大人物的時候.
怎樣可以讓他聽得入耳.
甚至乎不發怒呢?.
禎賢這裡提醒我們.
你要用柔和的聲調來去勸勉.
並且你要耐心等待.
君王自己去改變.
其中一個很經典的聖經例子就是.
大家記得《以斯帖女王》.
你會發覺.
如果細心看《以斯帖記》.
好幾次當他面對著.
阿哈徐老王要向他游說的時候.
他總是開口怎樣呢?.
「王若以為美,求你」.
「王啊,你覺得如果這樣是好的,求你」.
他總是以這一句做開場白.
他不會給人一個很催逼的感覺.
而是他很溫柔地.
邀請對方配合.
然後亦都很耐心期待.
對方真的會配合,會去改變.
所以你看到在《以斯帖記》裡面.
這招很殺食.

$^{601}$阿哈徐老王每次都真的被以斯帖.
他這樣不覺不覺地勸服了.
所以有人說得很好.
他說當你勸說人的時候.
你要想一幅圖畫.
就好像下雪一樣.
好像飄雪一樣.
萬不經意地.
輕輕地下雪,下雪,飄雪.
就是這樣滲入地下.
融入人的心裡.
來到他潛移默化.
不知不覺地溫柔地勸服了對方.
當然,第四方面少不了的就是.
當我們在進行勸說之前.
你真的要祈禱.
你要祈禱求主為你預備.
不單止預備你.
亦都預備好對方的心.
真的,其實你和我很明白.
靠著我們的天性.
很難真的說話恰到好處.
人總有情緒,脾氣.
人的說話總有甩漏.
所以我們要常常仰望.
祈求我們的主耶穌.
你會明白在聖經裡面.
你見到耶穌.
祂的勸說智慧.
祂的態度,祂的時機.
永遠都最好.
最恰到好處.
所以你祈求祂的時候.
你亦都可以.
好像耶穌一樣.
願耶穌真的給你這份智慧.
這份時機.
好了.
我們看最後一段經文.
這裡提到.

$^{641}$「鐵磨鐵,磨出印來,朋友相感」.
這裡原來和本有括著.
原文是磨朋友的面.
即是在說品格的意思.
「朋友相感,也是如此」.
提一提我們就是.
原來真誠的關係.
是需要互相去磨.
互相去抵禮.
鐵和鐵磨擦.
才會磨出鋒利,亮麗的刀印.
不過我們有時會不禁問.
其實一把鐵不就很好嗎?.
為何一定要常常.
和另一把鐵磨擦,磨練?.
為甚麼?.
是因為鐵有一個特性.
很容易擺放久了就會生鏽.
變得不鋒利.
所以兩把鐵要經常去磨練.
才能保持亮麗,鋒利.
同樣其實.
弟兄姊妹,我想.
你也不喜歡被人磨你的臉.
不舒服.
被人磨你的品格,性格.
不是一件.
或不一定是一件開心喜樂的事.
但卻又是.
在我們的生命中.
一個必需的過程.
一個必需的元素.
信徒之間其實真的需要.
互相去磨,互相去抵禮.
將你那些粗派的脾氣.
那些品格.
在和人相處的裡面.
在和弟妹們事奉的裡面.
磨磨磨,磨滑它,磨走它.
又或者將你本身裡面.

$^{681}$有那些美好的特質.
不過被獸跡遮蓋.
你又和其他人磨磨磨.
將這些美善顯露出來.
唯有這樣.
整個教會群體的生命才可以怎樣?.
保持亮麗,不生鏽.
來到有活力,有更好的應變能力.
就好像我挺喜歡看一套節目.
叫做「叢林的法則」.
大家有沒有看過?.
「叢林的法則」.
是一個韓國製造的綜藝節目.
每一次都會有.
或者每一輯都會有一個任務.
就是一班人要去一個叢林.
或者去一個荒島裡面.
生存幾天.
他們沒有什麼物資.
可能只是帶著一些工具.
他們一班人裡面.
可能有八,九個.
每個人都有不同的性格.
有不同的才華.
或者有不同的經歷.
他們去到叢林或荒野裡面.
就要大家一起合作.
因為什麼物資都沒有.
他們需要分工.
商量和互相配合.
哪些人負責預備睡覺的地方.
哪些人去想辦法生火.
哪些人去想山上.
地面上有什麼可以吃.
又或者懂得游泳的人.
就要出海去打魚.
或者拿海螺來沖基.
他們在過程裡面.
真的要互相配合.
他們要學會放下自己的執著.

$^{721}$要增加自己的包容性.
亦要學習互相欣賞和鼓勵.
所以你會看到.
在這個節目裡面.
給予我們一個道理.
一個人能夠做到的事.
真的很有限.
但一班人一起同心協力的時候.
就可以完成一個又一個.
不可能的挑戰和任務.
所以其實這個節目都頗長壽.
做了十多年.
為什麼長做長有.
其實都是因為.
它反映了在現實生活裡面.
人很嚮往一些生活情景.
就是有一班人.
他們有著不同的性格氣質.
他們有不同的意見.
有不同的經歷.
但原來當他們走在一起.
他們有共同的目標的時候.
他們就可以去互相切磋去磨練.
他們透過他們的合作和補位.
就可以拓闊大家的視野.
可以增加大家包容的胸襟.
亦可以給予自己.
更多做人的彈性.
甚至乎可以.
在一個互相抵禮的情況下.
去逼你產生更多.
應變生活場景的能力.
過程辛苦嗎?.
是辛苦的.
很多汗水.
甚至乎會有眼淚.
有時甚至乎會吵架.
有磨擦.
但最終你會發現.
總是會互相祝福.

$^{761}$總是可以將彼此的關係.
拉得更近更緊密.
就好像我想.
大家都應該有看過.
旺朕的跳舞賀聯片.
你會看到.
甚至乎有些幕後花絮.
一班同工他們一起去練習.
甚至乎有些就.
不怕拋頭露面.
在街上跳.
不怕醜.
我自己說得難聽一點.
五十歲人都要學跳舞.
但我會發現.
如果只有我一個.
我打死都不會這樣做.
想都不會這樣做.
一定不會.
但很有趣的是.
一班人走在一起.
你竟然會這樣做.
因為有同一個目標.
同一個心智.
很想在疫情下.
給大家更多的祝福.
一些活力和正能量.
你會發現在互相帶動下.
沒錯是辛苦的.
但一班人.
又不太察覺.
大家跳著跳著.
嘻嘻哈哈笑著.
事就這樣成了.
所以回頭看.
其實很感恩.
亦都很興奮.
事實上.
弟兄姊妹.
有人說得很好.

$^{801}$就是說.
一個人其實.
真的可能會走得快一點.
因為你不需要.
跟著別人的步伐.
或者被別人拖慢你.
一個人所以走得快一點.
但一起走.
才會走得久一點和遠一點.
這是非洲一個.
很出名的諺語.
可能和他們的生活很有關.
因為交通不太方便.
基本上交通工具就是.
他們兩條腿.
所以很多時候.
特別是婦女會怎樣.
背著很多東西.
在頭上或者肩膀上.
但一班人走.
沿途聊天.
互相打氣.
笑一下.
就會由一個城鎮.
走到另一個城鎮.
不會太累.
開心很多.
同樣真的.
如果當我們肯放下.
自己一些想法.
或者過往一些.
固有的做事模式的時候.
當你肯多些.
和身邊的人來切磋.
來互相刺激.
碰撞的時候.
其實你會發現.
這樣才有新的契機湧現.
這樣才能夠.
突破到一些.

$^{841}$瓶頸位.
我們經常說.
這樣才可以令到你.
走得更遠.
正如科學家研究.
眼睛群是怎樣飛的.
他發覺他們.
原來一起飛的時候.
真的會飛得遠一點.
你發現很多時候.
你看到的紀錄片.
或者相片.
眼睛群總是怎樣.
V字形這樣飛.
有一隻帶頭.
然後一直散開.
好像V字一樣.
科學家研究.
一出眼群理論.
他發覺就是.
當他們這樣飛的時候.
會比一隻眼.
即是一隻孤眼自己飛.
會飛多百分之七十的路程.
原因是什麼呢?.
他們分析過.
當一群眼這樣.
V形飛的時候.
當後面的眼.
看到前面的眼.
拍動雙翼的時候.
原來有一個鼓勵的作用.
即是帶頭的.
都奮力去拍翼.
一直向前飛的時候.
後面的都會被帶動.
而反過來.
當後面的眼在叫的時候.
又會給前面的眼.
一種鼓舞.

$^{881}$好像加油一樣.
讓前面的眼繼續前行.
而當帶頭的眼覺得累的時候.
它又會自動退後.
而後面的眼又會自動補位.
又或者假如.
當有眼生病.
或者受傷的時候.
就會有另外的兩隻眼.
飛下來協助它.
或者保護它.
直到它康復.
又或者直到它離開世世.
然後.
那兩隻留下的眼.
又會嘗試趕上原先的眼群.
又或者它會另外組織新的隊伍.
繼續完成它遷徙的旅程.
真的弟兄姊妹.
今天這段經文提醒我們.
在屬靈生命的成長裡面.
其實就好像一趟.
不斷創造建立關係的旅程.
你要和其他的信徒.
去結伴同行.
在過程裡面.
你要敢於聽.
敢於說.
弟兄姊妹的真心話.
誠實話.
而當你來到.
向弟兄姊妹作出勸勉的時候.
你的說話也需要調和.
需要設計.
也需要禱告.
讓耶穌基督帶領彼此雙方.
來成長.
最後在一段真誠的關係裡面.
我們也不少得.
你需要有事奉.

$^{921}$和弟兄姊妹一起合作相交.
在過程裡面.
互相去琢磨.
互相去砥礪.
唯有這樣.
我們才能夠.
更有力量地.
一起去實踐完成.
耶穌基督託付給我們的使命.
所以在這個時候.
也讓我們唱一首回應詩.
一首在程序表裡面刊登了的.
叫做《伴我成長》.
一首我在團契.
在二十年前唱的一首歌.
讓我們用這首詩歌來回應.
我們今天所領受的聖經教導.
《伴我成長》.
\newpage



\section{約翰福音 1:1-18}
\label{sec:yigtcB2OiSE}
\textbf{2022年2月20日崇拜直播｜梁家麟牧師｜太初有道｜約翰福音1\_1-18}
\newline
\newline
連結: \href{https://youtube.com/watch?v=yigtcB2OiSE}{\texttt{https://youtube.com/watch?v=yigtcB2OiSE}} ~~~~ 語音日期: 2022-02-21
\newline
\newline
\hyperref[sec:Q_GbM3l6t1s]{\small{< < < PREV SERMON < < <}}
~
\hyperref[sec:index_chronic]{\small{[返順時目]}}
~
\hyperref[sec:cYMHT7F9Rko]{\small{> > > NEXT SERMON > > >}}
\newline
\newline
約翰福音 1:1-18
\newline
\begin{longtable}{cl}
\hline
\hline
章節 & 經文 (和合本修訂版)\\
\hline
1:1 & \begin{tabularx}{0.7\textwidth}{X} 太初有道，道與神同在，道就是神。 \end{tabularx} \\ \\ \relax
1:2 & \begin{tabularx}{0.7\textwidth}{X} 這道太初與神同在。 \end{tabularx} \\ \\ \relax
1:3 & \begin{tabularx}{0.7\textwidth}{X} 萬物都是藉著他造的，沒有一樣不是藉著他造的。凡被造的， \end{tabularx} \\ \\ \relax
1:4 & \begin{tabularx}{0.7\textwidth}{X} 在他裡面有生命，這生命就是人的光。 \end{tabularx} \\ \\ \relax
1:5 & \begin{tabularx}{0.7\textwidth}{X} 光照在黑暗裡，黑暗卻沒有勝過光。 \end{tabularx} \\ \\ \relax
1:6 & \begin{tabularx}{0.7\textwidth}{X} 有一個人，是從神那裡差來的，名叫約翰。 \end{tabularx} \\ \\ \relax
1:7 & \begin{tabularx}{0.7\textwidth}{X} 這人來是為了作見證，是為那光作見證，要使眾人藉著他而信。 \end{tabularx} \\ \\ \relax
1:8 & \begin{tabularx}{0.7\textwidth}{X} 他不是那光，而是要為那光作見證。 \end{tabularx} \\ \\ \relax
1:9 & \begin{tabularx}{0.7\textwidth}{X} 那光是真光，來到世上，照亮所有的人。 \end{tabularx} \\ \\ \relax
1:10 & \begin{tabularx}{0.7\textwidth}{X} 他在世界，世界是藉著他造的，世界卻不認識他。 \end{tabularx} \\ \\ \relax
1:11 & \begin{tabularx}{0.7\textwidth}{X} 他來到自己的地方，自己的人並不接納他。 \end{tabularx} \\ \\ \relax
1:12 & \begin{tabularx}{0.7\textwidth}{X} 凡接納他的，就是信他名的人，他就賜他們權柄作神的兒女。 \end{tabularx} \\ \\ \relax
1:13 & \begin{tabularx}{0.7\textwidth}{X} 這些人不是從血生的，不是從情慾生的，也不是從人的意願生的，而是從神生的。 \end{tabularx} \\ \\ \relax
1:14 & \begin{tabularx}{0.7\textwidth}{X} 道成了肉身，住在我們中間，充充滿滿地有恩典有真理，我們也見過他的榮光，正是父獨一兒子的榮光。 \end{tabularx} \\ \\ \relax
1:15 & \begin{tabularx}{0.7\textwidth}{X} 約翰為他作見證，喊著說：「這就是我曾說：『那在我以後來的先於我，因為在我以前，他已經存在。』」 \end{tabularx} \\ \\ \relax
1:16 & \begin{tabularx}{0.7\textwidth}{X} 從他的豐富裡，我們都領受了恩典，而且恩上加恩。 \end{tabularx} \\ \\ \relax
1:17 & \begin{tabularx}{0.7\textwidth}{X} 律法是藉著摩西頒佈的；恩典和真理卻是由耶穌基督來的。 \end{tabularx} \\ \\ \relax
1:18 & \begin{tabularx}{0.7\textwidth}{X} 從來沒有人見過神，只有在父懷裡獨一的兒子將他表明出來。 \end{tabularx} \\ \\
[1ex]
\hline
\hline
\end{longtable}
$^{1}$讓我們繼續聆聽上帝的話語.
其實都是一段關於為耶穌基督作見證的故事.
記載在《約翰福音》第一章一至十八節.
經文裡是這樣說.
「太初有道,道與神同在,道就是神」.
「這道太初與神同在」.
萬物是藉著他造的.
凡被造的沒有一樣不是藉著他造的.
生命在他裡頭,這生命就是人的光.
光照在黑暗裡,黑暗卻不接受光.
有一個人是從神那裡差來的,名叫約翰.
這人來,為要作見證,就是為光作見證.
叫眾人因他可以信.
他不是那光,乃是要為光作見證.
那光是真光,照亮一切生在世上的人.
他在世界,世界也是藉著他造的.
世界卻不認識他.
他到自己的地方來,自己的人倒不接受他.
凡接待他的,就是信他名的人.
他就賜他們權柄,作神的兒女.
這等人不是從血氣生的,不是從情欲生的.
不是從仁義生的,乃是從神生的.
他到成了肉身,住在我們中間.
充充滿滿的有恩典,有真理.
我們也見過他的榮光,正是父獨生子的榮光.
約翰為他作見證,喊著說.
這就是我曾經說,那在我以後來的.
反成了我以前的,因他本來在我以前.
從他豐滿的恩典裡,我們都領受了.
而且因上加恩.
律法本是藉著摩西傳的.
恩典和真理都是由耶穌基督來的.
從來沒有人看見神.
只有在富懷裡的獨生子,將他顯明出來.
今天早上有梁家倫牧師在我們當中靜悼.
講題是「太初有道」.
在疫情期間和大家問安.
在未來大概兩年多三年的日子裡.
我會和大家走過約翰福音之旅.
斯祿約翰受聖靈感動.

$^{41}$要撰寫一部有關耶穌基督的生平和教訓的傳記.
這也可以說是他個人的回憶錄.
甚至一部分是自傳.
因為他是耶穌基督一個密切的同代的見證人.
陪伴耶穌走過最後那三年的時間.
彼此有很多交集.
不過,約翰不僅要為耶穌做傳記.
更要見證這位在人間活動過33年的拿撒勒人耶穌.
就是舊約眾先知所應許要來的彌賽亞.
所以他不僅是講故事人.
他更是見證者.
再進一步.
他要揭示一個重大的秘密.
是舊約眾先知都不知道的.
這位彌賽亞不僅是上帝的代言人.
是上帝自己.
上帝道成肉身降世為人.
親自來啟示他自己並施行拯救.
如此,約翰不僅是講故事者.
不僅是見證人.
更是一個揭秘者.
所以,這本書不僅是歷史記述.
也是神學論述和信仰見證.
如同約翰在20章的結尾所說.
耶穌在門徒面前行了許多神蹟.
會記載在書上.
但記者些時.
要叫你們信耶穌是基督,是上帝的兒子.
並且叫你們因信了祂.
就可以因祂的名得生命.
由於這本書的要旨是神學論述和信仰見證.
所以,它的開首並非耶穌的生平故事.
而是有關耶穌的神聖本質的講論.
在頭十八節裡.
約翰急不及待地要告訴我們.
耶穌的本相是誰.
耶穌是道,是上帝.
作者亦將整本書裡的所有重要的詞彙和概念都搬出來.
光,生命,上帝的獨生子,恩典,真理.
將這些觀念共惹一爐,連為一體.

$^{81}$所以可以說,這十八節已經鋪陳了整本書所講的神學論述的要點.
全都講完了.
不過,我不是自己打自己.
我們要退後一步去想.
頭十八節的經文,也可以說是耶穌的生平故事.
跟《婦女福音》不同.
約翰沒有將耶穌人間歷史的開端.
定在瑪利亞誕生的嬰孩之上.
而是定在一個人肉眼看見,理性亦認知的奧秘之上.
道成肉身.
是的,道成了肉身.
這才是人間耶穌歷史的開端.
我們以為因今天在大衛的城裡為你們生了救主,就是主基督.
這是耶穌故事的開端.
對所有人類來說,父母當然是我們的源頭,對不對?.
但耶穌不是.
約瑟和瑪利亞不是耶穌的真正起源.
另外,耶穌跟世界的關係,也不是從伯利恆的馬槽開始的.
約翰要揭示一個秘密.
當世界被造之初.
這個道已經跟萬有有關係.
約翰告訴我們,道成了肉身.
這是人間耶穌故事的真正開端.
所以,一章一到十八節既是神學論述.
又是生平故事,又是信仰見證.
三者合而為一.
這段偉大的聖經,介紹了道Logos.
大家明白了嗎?.
這個希臘詞彙,跟中文所說的道,意思很接近.
都是多層面的,三個層面.
第一,道是言說,是message,是word.
第二,道是指事物構成的原理,道理,對不對?.
是萬有運作的法則,reason,principle,形而上者為之道.
第三,道亦都連上上帝,天道.
即是說,它是一個實體,不僅僅是道理,也不僅僅是言說,是一個實體.
是一個有意志的天的外顯.
我們先來看第一點,道和上帝.
約翰在第一節劈頭就指出道和上帝的關係.
他指出兩個重要的道理.
第一,道從太初就已經存在,太初有道.

$^{121}$在最始初的時候就已經有道.
始初,即是說沒有之前,因此道是沒有開始的.
這是從側面說明道的時間性,它是先存的,pre-existence,或者說是永恆的.
道太初與上帝同在,從一開始,或者從開始的開始,我講的很急.
道一直與上帝同在,如果上帝無始無終,道就是無始無終,對不對?.
因為道沒有始初,所以它才成為萬物的起源.
耶穌在人間,你可以說祂有一個時間的開端,但是道是沒有開始的.
其實在約翰後面的寄述,耶穌多次漏口風,說明祂與父一樣,其實是沒有開始的.
在8章58節,耶穌說過:我實實在在告訴你們,還沒有阿巴拉罕,就有了我.
在17章5節,祂在祈禱的時候說:父啊,現在求你使我與你共同享榮耀.
就在未有世界以前,我與你所有的榮耀.
耶穌是否漏口風說了祂是永恆?祂是無始的.
太初這個想法,讓我們聯想到創世記的第一句.
「起初,上帝創造天地」,起初和太初都是指著一個絕對意義的開端.
先於一切事物的時間,或者更加準確地說,是先於時間的時間.
創世記告訴我們,上帝在起初就開始了祂的創造工作,這是萬有的起點.
約翰也告訴我們,道在太初就已經存在,就已經與上帝一起創造萬有.
這是萬有的起點,這是第一點我要說的.
道與上帝同在是永恆的.
第二點,道與上帝的關係,他說道的是道既是上帝,又與上帝有分別.
約翰明白指出道是上帝,道是上帝.
有不少人認為,約翰說道上帝的時候,沒有給一個定義的開端,一個「The」在前面.
說明道不是絕對意義的上帝,只是具有神性,有神性.
但是,D.A. Carson,這些學者說得很好,新約有太多的例子.
那些的遺詞,Predicate,是沒有定觀詞的,但都可以是獨特的.
The God,不是God,不明白我說什麼,這段就無所謂了.
總之道是上帝,但是約翰的說法就比較扭曲.
第一節,道與上帝同在,道就是上帝.
如果你單單說道是上帝,The word was God,我們很容易明白.
但同時又說The word was with God,道與上帝同在,這樣就有點扭曲了.
A.如果是B,你又怎會同時說A.B.在呢?.
你明白嗎?我是梁家倫,你明白嗎?.
但是跟著我說,我和梁家倫在一起,那就是說什麼?.
我如果是梁家倫,那要什麼在一起呢?.
唯有我們認識教的三一論,我們才明白,耶穌基督既是上帝,又和副上帝有分別.
他是一章十四節和十八節所說的副獨生子.
所以第一節講論道與上帝的關係,告訴我們兩個道理.
第一道是永恆的,第二道是與副有分別的上帝.
然後我們看到第二大點,道與創造.
約翰指出道是上帝創造萬有的憑藉,萬物是藉著他做的,凡被做的,每一樣不是藉著他做的.

$^{161}$第三節,第十節也說世界也是藉著他做的.
我們其實不是很明白,坦白說我不明白.
我不明白上帝創造萬有的時候,為什麼需要有一個藉著,一個agent,一個媒介.
不是從無中創造萬有嗎?並不是說有就有,命立就立嗎?.
上帝創造不需要任何工具材料,也不需要任何工序,任何過程.
你不知道如何講到上帝創造的時候需要一個媒介.
神學家為了要解讀這一段,於是只好將上帝創造的媒介定位在上帝的言說之上.
上帝講什麼,就有什麼,所以上帝是藉著言語來創造.
就是creation by word,平情而論,我有些七人台的感覺.
藉著言語來創造,這是一個畫蛇添足,空道無物,是一個廢話的說法.
我教書要開口說話,所以我是藉著言語去教書.
有人是藉著言語去管理,有人是藉著言語去醫病,有人是藉著言語去做買賣.
我不僅是藉著言語去做模樣,我還藉著我的笑容去做模樣,藉著我的眼神去做模樣,藉著我的手勢去做模樣.
你告訴我,你告訴我.
神學家搞盡腦汁,將上帝的說話,明明是說話,變成一個實體.
所以就是道,就是聖旨,例如一節常常被人拿來作為illustration的經文.
詩篇33篇第6節,那裡說到朱天,藉耶華的命而造,萬象是藉他口中的氣而成.
很多人將這篇經文變成三一論的論據,命是指聖言,就是指聖旨,氣是指聖靈,所以就三一.
這篇經文就是證明三一上帝參與在創造計劃裡面,父在創造,子和靈都在創造.
不過我是一個簡單的人,我最不明白這些複雜的東西.
我在想,朱天藉耶華的命而造這句話,為什麼我們不將子和靈同時放在耶華這個詞上面呢?.
在舊約裡面,凡是耶華都只是指著父嗎?不是吧?.
那不如簡單地說,父指聖靈,朱天藉著父指聖靈,也就是耶華的命而造,就簡單了.
為什麼要零零碎碎地說,命就只指著子,而不是耶華是指著子和靈呢?.
為什麼要雅謀整餅,子就是指著命,靈就是指著氣,而耶華就只是指著父呢?為什麼要這樣分呢?.
無論如何,聖經說萬有是藉著道而被創造的.
這個除了指著道有份於創造的行動之外,更加有一個很重要的含義.
就是道和萬有的關係.
因為在第十一節,約翰就說出一個很重要的話.
他道自己的地方來.
這裡強調,世界和萬有原來都是屬於道的.
道不是這個世界的外來者,更加不是一個入侵者.
蘭吉版的俗語說,有貧自遠方來,非奸即盜.
所以說,這個不知不覺地闖入這個世界,是不是有些別有用心呢?.
我希望各位弟兄姊妹,當耶穌踩你的牆,干涉你的想法和做法的時候.
你不會說你是誰,你有什麼權干涉我的事.
我父母都不管我,你管我什麼?.
你要記住,你的牆是他的牆,他有權.
他道自己的地方來.
世界是屬於上帝,是屬於耶穌基督的,是不是?.

$^{201}$他有權,返來他自己的地方.
呼上帝藉著道,藉著耶穌基督來創造世界.
這是幾乎整本新約聖經的共證.
約翰這樣說,保羅這樣說,希伯來書和啟示錄都這樣說.
曾經有學者討論,這個觀念到底是誰發明出來的呢?.
誰抄誰呢?.
至少在我們現有的四卷科幻書裡面.
你沒有發現人間耶穌曾經和門徒說過這個道理.
我有份參與創造,你有沒有聽過?.
耶穌有沒有在人間說過?沒有嗎?.
我們很難想像這麼神秘的創造起源.
這麼重要的道理,是由個別的使徒自己想出來.
然後所有的使徒都認,人人都認.
你覺得很古怪嗎?.
所以我們可以比較安心地說.
這個要是聖靈刻意啟示給各使徒.
你可以伸入各卷聖經的作者.
作為一個重要的基督論的真理.
你說聖子就是基督,其實已經足夠了.
但你要強調,他有份參與創造,其實是一個很有趣的說法.
希伯來書說,父上帝藉著他的兒子去創造諸世界.
希伯來書一章二字,就在這個末世藉著他的兒子曉遇安門.
由早已立他為承受萬有者,也曾藉著他創造諸世界.
保羅在《歌羅西書》也指出,愛子是那個不能看見的上帝的像,是守生的.
在一切被做的以先,因為萬有都是靠他做的.
無論是天上的,地上的能看見的,不能看見的.
是有為的,主治的,執政的,掌權的.
一切都是藉著他做的,也是為他做的.
他在萬有之先,萬有都是靠他而立.
我們要注意一句說話,為他而做.
為他做的說話,套用,這些只是雕書包,你可以不理的.
亞歷士多德的四因說,藉著他做是動力因,Efficient Cause.
為他做的是目的因,Final Cause.
為什麼上帝要做世界呢?是為基督而做的世界.
保祿是特別看重目的因的.
多次指出,上帝是藉著基督創造萬有.
目的是要使萬有將來復歸基督,因為世界是為基督而做出來的.
所以萬有最後要復歸基督.
以法所使一章九到十節,都是照他自己所預定的美意.
叫我們知道他指的奧秘,要照所安排的.

$^{241}$在日期滿足的時候,就是天上,地上一切所有的.
都在基督裡面同歸於一.
萬有最終要復歸耶穌,因為世界是為他而做的.
啟示錄的作者,如果我們相信是約翰的話.
也都稱耶穌基督是上帝創造萬物之上為元首.
啟示錄三章十四節.
元首有起始,源頭和統治者的意思.
所以道,即是耶穌基督既是萬有的根源,也是萬有的歸宿.
是萬有的主,擁有者,統治者.
從創世記一章,我們看到有形的宇宙萬物是藉著道而被造的.
用歌羅西書,希伯來書的說法.
我們就知道這個世界是之外任何的世界.
上帝創造諸世界,即不止一個世界,都是藉著道而被造的.
所以無論是天使,各種靈體,都是藉著道而被造的.
所以上帝不只是藉著道創造了一個物理世界.
我們今天看到的世界,其實有很多不同的世界.
都是同樣屬於耶穌基督.
萬有因基督而被造,也為基督而被造.
我們可以看到萬有和道的密切關係.
約翰提這段經文,我們不是完全明白的經文,因為是奧秘來的.
不是要告訴我們道的源頭,道是沒有源頭的.
而是要告訴我們人類的源頭,我們自己的源頭.
道是我們的源頭,飲水思源,道就是我們的源.
我們來到第三大點,生命和光.
約翰引進兩個觀念,生命和光.
一章四到五節,生命在他裡頭,生命就是人的光.
光照在黑暗裡,黑暗卻不接受光.
生命和光是貫穿整卷約翰福音的重要觀念.
並且都和救贖有關連.
如果萬有是藉道而被造.
萬有,包括我們人在內,都是藉道而獲得生命.
不過,約翰在這裡要指出的不是道給我們原初的受造的生命.
那個塵世的生命.
而是要指出道要給人新的生命.
前面說到道在創造的時候的地位和作用.
我是要指出道和萬有本來的關係.
但是道之成為肉身,道之來到人間.
不是僅僅要跟你認親戚,恢復你以前和他本來有的關係.
原來我們是有親戚關係的,不是.
道成肉身目的是要賜人新的生命,和人建立新的關係.

$^{281}$所以,這裡說的生命,是強調道是生命之源.
是強調道要賜給我們新的生命.
正因為約翰強調的是新生命,不是舊生命.
是第二次創造,即是救贖,而不是第一次創造.
所以後面他就說,十二十三節.
凡接待他的人,就是信他的人.
他就賜他權柄作上帝的兒女.
而不是從血氣生的,不是從情慾生的,也不是從仁義生的.
乃是從上帝生的.
新生命不是來自血氣和情慾,而是來自上帝.
這也是我們所說的重生.
你一定會記得,四卷荒書獨特的一個技術.
第三章,耶穌和尼哥底姆,或者我們有個叫尼哥德姆.
談到重生的道理,一個很重要的道理.
就是第三章,所以這是約翰福音的一個重要主題.
重生,賜新的生命.
生命在他裡頭.
這句話的意思是,生命是由道而出.
道是生命的源頭,也可以說道擁有生命.
他有派生命,賜生命的主權.
約翰在後面也指出.
耶穌和父上一樣擁有生命的掌握權.
五章二十六節,因為父怎樣在自己有生命.
就自給他而止,也照樣在自己有生命.
我講這些,是想讓你看見.
約翰在這裡鋪陳的東西,其實後面全部呼應,全都講了.
並且用耶穌的口來講這裡做的神學論述.
所以不是單一任意的說法.
生命就是人的光.
這句話最簡單的意思就是,生命,道就是人的出路.
光就是出路.
將生命和光連在一起說,多數是舊俗有關的.
希臘人將智慧和拯救是和光連繫的illumination,光照你.
黑暗是和死亡同義詞,光明是和舊俗同義的.
從知識層面講,黑暗,意涵無知,光明就是啟迪真理.
人落在黑暗的境地,無出路,為了被光照亮才找到拯救之路.
耶穌基督是光,所以祂是道路,祂是真理,祂是生命.
伊罕既然提到光,就立即講到黑暗.
光照在黑暗裡,黑暗卻不接受光.
我們記得,上帝在創造這個世界之前,一切都是混沌和黑暗的.

$^{321}$你記得創世記一章二節的說法.
上帝首先做的是什麼?是光,讓光驅走黑暗,三節.
光是第一天的創造,你要記住,它是比太陽,比任何其他光體的創造更早.
一章十五十六節才開始講的.
光的源頭是上帝,不是太陽,千萬不要誤會.
上帝就是光,在祂裡面毫無黑暗,這個若然一書,一章五節的說法是什麼?.
當上帝介入,上帝作為在黑暗裡面引進光明.
上帝是光的來源,眾光之父.
所以人離開上帝,就重新陷落在黑暗裡面.
所以當聚進入世界之後,黑暗就是世界和世人的寫照.
因為人落在黑暗裡面,道才要成肉身,將光重新引入人間,道進入人間.
光就照亮人間的黑暗.
可惜的是,落在黑暗境地裡面的人接受道,拒絕接受道,也拒絕接受光.
不接受,有人亦不理解,不認識,或者不能征服.
這個動詞是過去式,說明這是過去一個事實性的情況.
而不是說黑暗永遠都不能接受光,不是這樣的.
曾經拒絕過光.
他在世界,世界也是藉著他做的,世界卻不認識他.
他到自己的地方來,自己的人倒不接待他.
落在黑暗境地裡面的人既不認識光,也不接待光.
生命和光是約翰福音非常重要的主題.
耶穌在後面多次強調他是生命.
例如,十一章二十五節,十四章六節,也多次強調他是光.
八章十二節,九章五節.
生命和光是耶穌基督常用的自稱.
我是光,我是生命.
好,我們來到最後一個大段,道成肉身.
一看,展述了道的神聖和萬有的關係.
目的是要點出一個石破天驚的真理.
道成肉身.
一章十四節,道成了肉身.
住在我們中間,充充滿滿的恩典和真理.
我們也見過他的榮光,正是父,獨,生,子的榮光.
道成了肉身.
也就是說,上帝成為人.
全知全能,無所不在的上帝.
竟然降臨人間,變成和我們一樣的凡人.
這是何等令人詫異,難以令人置信的事.
我曾經很多次說過.
教會有別於其他有神論的宗教.

$^{361}$我們不僅相信宇宙間有一個上帝.
我們是認信,二千年前在加里尼湖邊生活的拿撒勒的耶穌.
祂是上帝.
耶穌是上帝,才是基督教最重要,最核心的信息.
我再說,為什麼我們稱祂為基督教,而不是上帝教.
我們信仰的對象是耶穌基督,是道成肉身.
既是人,又是上帝的耶穌基督.
如果上帝曾經來過人間.
並且祂不是快閃擋,旋風式的即來即走.
也不是神靈附體的「噠噠噠噠」.
上了某一個先知,某一個神人的身.
而是實實在在成為人,誕生在人間.
經歷從嬰孩到少年,到成人的整個人生歷程.
既是上帝,又是人的耶穌.
就是我們能夠認識上帝,明白上帝的心意.
與上帝建立關係的唯一途徑.
我們說唯一.
不是說上帝從前沒有猜過任何先知來做代言人.
有!舊有聖經說述過很多列祖先知曾經做過上帝代言人.
希伯來書一章一字都這樣說.
不過,與耶穌基督的正牌上帝相比.
所有先知,所有上帝代言人都要靠邊站.
有什麼地位?沒有的.
我們今天以為陳牧師不來.
於是就找了梁家倫做代堂主任.
但他突然走出來,那代堂主任有什麼意義?.
難道我還說話嗎?靠邊站,對不對?.
有了正牌,正印,代表還有什麼好做?.
沒有什麼好做,難道他還說.
我知道陳明全牧師的心意,你說什麼?.
陳明全牧師生了一棟在這裡,需要你跟他說他的心意是什麼?.
當耶穌基督作為上帝,親自來說上帝是誰.
上帝對我們有什麼要求的時候.
你告訴我,其他的代理人,代言人有什麼意義?.
沒有了,所以耶穌基督是唯一.
耶穌是唯一,我要告訴你這是很霸道的說法.
耶穌宣稱過,我是道路真理生命,若不積著我,每人能到乎那裡去.
是不是很霸道?.
使徒也和應這個霸道的說法,除他以外,必無拯救.
因為天下人間,每次有別的命,我們可以靠得救.

$^{401}$是不是很exclusivist(排他主義).
是的,不過上帝就是這麼霸道的.
六宮粉內無顏色,倚天一出,誰與爭風?.
從來沒有人看見過上帝,只是在腐壞裡的獨生子將祂表明出來.
歷史上,特別是在舊約的基述裡面,上帝曾經以不同的形式向人顯現過.
不過,上帝顯現過是一件事,只有肉眼的凡人,無論是亞伯拉罕,瓦戈,柯西,大衛乃至以利亞.
嚴格地說,他們沒有看見上帝,沒有的.
所以,約翰是斬釘截鐵地說,從來沒有人看見過上帝.
你沒理由挑戰聖經的說法,他說沒有,你不可以說有.
從來沒有人看見過上帝.
人要看見上帝,只能夠張開眼睛,看著耶穌基督.
耶穌基督是上帝在人間最完美的彰顯,沒有任何比耶穌更完美的說出上帝是誰.
上帝親自來說,有更真實,更準確的可能嗎?沒有.
聖賢成了肉身的道,他的道性,他的神性,很大程度上是被遮蓋的.
這是不得已的事,所以和耶穌同代的法利賽人都認不出耶穌是上帝.
不過,這是不得已的事,未成肉身的道,能夠給只有肉身,只能夠看見肉身的我們看得到嗎?.
我們能夠接觸本來是一個零的上帝嗎?當然不能夠.
所以,耶穌基督是上帝謙卑虛己,成為人的樣式,很讓人能夠接觸到他.
所以,耶穌基督就是人能力所及,最大程度上對上帝的認識.
沒有他比,更加貼近上帝的原樣.
神學上說,耶穌基督是為人類的上帝,God for man.
這個意思其實很簡單,就像我們為兒童編寫一本能夠看得懂的兒童聖經一樣.
耶穌基督就是人類版的上帝,你明白我說什麼嗎?.
就是人最大可能能夠明白的上帝的版本.
所以,我們要認識上帝,看耶穌,要知道上帝的心意,問耶穌.
要和上帝建立關係,找耶穌.
求我們的主,今天和我們每一個人親近.
\newpage



\section{希伯來書 1:1-2:4}
\label{sec:cYMHT7F9Rko}
\textbf{2022年3月6日崇拜直播｜陳明泉牧師｜超越的基督:神親自曉喻｜希伯來書1\_1-2\_4}
\newline
\newline
連結: \href{https://youtube.com/watch?v=cYMHT7F9Rko}{\texttt{https://youtube.com/watch?v=cYMHT7F9Rko}} ~~~~ 語音日期: 2022-03-07
\newline
\newline
\hyperref[sec:yigtcB2OiSE]{\small{< < < PREV SERMON < < <}}
~
\hyperref[sec:index_chronic]{\small{[返順時目]}}
~
\hyperref[sec:bMqVNP_dN2s]{\small{> > > NEXT SERMON > > >}}
\newline
\newline
希伯來書 1:1-2:4
\newline
\begin{longtable}{cl}
\hline
\hline
章節 & 經文 (和合本修訂版)\\
\hline
1:1 & \begin{tabularx}{0.7\textwidth}{X} 古時候，神藉著眾先知多次多方向列祖說話， \end{tabularx} \\ \\ \relax
1:2 & \begin{tabularx}{0.7\textwidth}{X} 末世，藉著他兒子向我們說話，又立他為承受萬有的，也藉著他創造宇宙。 \end{tabularx} \\ \\ \relax
1:3 & \begin{tabularx}{0.7\textwidth}{X} 他是神榮耀的光輝，是神本體的真像，常用他大能的命令托住萬有。他洗淨了人的罪，就坐在高天至大者的右邊。 \end{tabularx} \\ \\ \relax
1:4 & \begin{tabularx}{0.7\textwidth}{X} 他所承受的名比天使的名更尊貴，所以他遠比天使崇高。 \end{tabularx} \\ \\ \relax
1:5 & \begin{tabularx}{0.7\textwidth}{X} 神曾對哪一個天使說過：「你是我的兒子；我今日生了你」？又說過：「我要作他的父；他要作我的子」呢？ \end{tabularx} \\ \\ \relax
1:6 & \begin{tabularx}{0.7\textwidth}{X} 再者，神引領他長子進入世界的時候，說：「神的使者都要拜他。」 \end{tabularx} \\ \\ \relax
1:7 & \begin{tabularx}{0.7\textwidth}{X} 關於使者，他說：「神以風為使者，以火焰為僕役。」 \end{tabularx} \\ \\ \relax
1:8 & \begin{tabularx}{0.7\textwidth}{X} 關於子，他卻說：「神啊，你的寶座是永永遠遠的；你國度的權杖是正直的權杖。 \end{tabularx} \\ \\ \relax
1:9 & \begin{tabularx}{0.7\textwidth}{X} 你喜愛公義，恨惡罪惡；所以神，就是你的神，用喜樂油膏你，勝過膏你的同伴。」 \end{tabularx} \\ \\ \relax
1:10 & \begin{tabularx}{0.7\textwidth}{X} 他又說：「主啊，你起初立了地的根基，天也是你手所造的。 \end{tabularx} \\ \\ \relax
1:11 & \begin{tabularx}{0.7\textwidth}{X} 天地都會消滅，你卻長存；天地都會像衣服漸漸舊了； \end{tabularx} \\ \\ \relax
1:12 & \begin{tabularx}{0.7\textwidth}{X} 你要將天地捲起來，像捲一件外衣，天地像衣服都會改變。你卻永不改變；你的年數沒有窮盡。」 \end{tabularx} \\ \\ \relax
1:13 & \begin{tabularx}{0.7\textwidth}{X} 神曾對哪一個天使說：「你坐在我的右邊，等我使你的仇敵作你的腳凳」？ \end{tabularx} \\ \\ \relax
1:14 & \begin{tabularx}{0.7\textwidth}{X} 眾天使不都是事奉的靈，奉差遣為那將要承受救恩的人服務的嗎？ \end{tabularx} \\ \\ \relax
2:1 & \begin{tabularx}{0.7\textwidth}{X} 所以，我們必須越發注意所聽見的道，免得我們隨流失去。 \end{tabularx} \\ \\ \relax
2:2 & \begin{tabularx}{0.7\textwidth}{X} 既然那藉著天使所傳的話是確定的，凡違背不聽從的，都受了該受的報應； \end{tabularx} \\ \\ \relax
2:3 & \begin{tabularx}{0.7\textwidth}{X} 我們若忽略這麼大的救恩，怎能逃避呢？這拯救起先是主親自講的，後來是聽見的人給我們證實了。 \end{tabularx} \\ \\ \relax
2:4 & \begin{tabularx}{0.7\textwidth}{X} 神又按自己的旨意，更用神蹟奇事、百般的異能，和聖靈所給的恩賜，與他們一同作見證。 \end{tabularx} \\ \\
[1ex]
\hline
\hline
\end{longtable}
$^{1}$今天的經文是希伯來書一章一到四節.
二章一到第四節.
希伯來書一章一到第四節.
「神既在古時.
直著眾先知多次多方地曉遇列祖.
就在這末世直著他而止曉遇我們.
又早已立他為承受萬有的.
也曾直著他創造諸世界.
他是神榮耀所發的光輝.
是神本體的真相.
常用他權能的命令托住萬有.
他洗淨了人的罪.
就坐在高天至大者的右邊.
他所承受的命.
既比天使的命更尊貴.
就遠超過天使」.
第二章一到第四節.
「所以你們當越發鄭重所聽見的道理.
恐怕我們隨流失去.
那直著天使所傳的話.
既是確定的.
凡剛犯叛逆的.
都受了該受的報應.
我們若忽略這魔大的救恩.
怎能逃罪呢?.
這救恩豈先是主親自說的.
後來是聽見的人給我們證實了.
神又按自己的旨意.
用神跡.
其事.
和諸和百般的異能.
並聖靈的恩賜.
同他們作見證」.
今天宣講的題目叫.
是一個系列的講道.
「處於基督神親子曉遇」.
有請陳牧師.
各位弟兄姊妹平安.
今天我們開始一個新的講道系列.
之前講過6G系列.

$^{41}$這6G其實在之後不同的系列裡面.
都會滲入去.
希望大家在聽的過程中.
你會明白教會的核心價值更加多.
這次我們接下來的十多次的講道系列.
用希伯來書來做一個主題.
叫做「超越基督」的系列.
為什麼我們會用希伯來書呢?.
首先其實希伯來書.
在《新約聖經》裡面.
雖然我們不知道誰是作者.
或者讀者.
甚至我們都不是很清楚.
但是你看它的內容.
它特別講到基督耶穌那種超越.
那個基督論其實對於我們今天.
對整個基督教信仰來說.
是一個很重要的養分.
在希伯來書裡面講的那個耶穌.
比起舊約還要知道.
或者不同的身份角色.
不同的人物都更加超越.
這種超越不單只是說耶穌很厲害.
在每一次的勸勉當中.
其實都是針對在希伯來書的讀者群裡面.
對他們的生命.
知道耶穌這麼超越.
對他們的生命是有一些很重要的勸勉和提醒的.
我們深信我們跟工隊一起去研讀書卷的時候.
我們知道.
我們今天在這個世代裡面.
我們都面對不同的挑戰和困難.
如果我們的信仰沒有一種超越性在這裡.
就是著重在眼前的生活.
我們經常講很落地.
但是如果純粹看東西很落地.
但是沒有一個超越的眼光.
我們都是在地上.
營營業業地打平.
但是當我們生命裡面.

$^{81}$我們知道我們所信的神是一個怎樣超越的神.
特別是在希伯來書裡面.
我們找到這些養分的時候.
即使我們在生活裡面很落地地生活.
但是我們的眼睛可以看著天上的阿爸父.
望在高天的基督耶穌.
讓我們生命裡面更加有盼望.
有力量.
我們更加深信.
有一個超越的信仰.
就能夠為我們帶來超越的生命.
所以在接下來的這十多次.
希望大家真的好好地.
一起來思考這一段經文.
我們教會除了在講道裡面用這個系列.
在小組查經我們都會一起來配合.
我們會講多一點在舊約裡面的典故.
在舊約裡面看希伯來書的作者.
新舊約對照是怎樣去解釋的.
我們幫大家補一下課.
知道一下舊約在講什麼.
從講道到查經.
希望讓我們能夠在這幾個月裡面.
我們真的好好地掌握到希伯來書的信息.
講起希伯來書.
我要跟大家做一點補課的.
事實上在希伯來書.
你一直對照和其他書卷去看的時候.
它的寫作方式.
或者它呈現的面貌是不同的.
首先.
最大家關心的就是希伯來書的作者是誰呢.
坊間裡面有不同的.
識經學家都有不同的說法.
不過到最後還是沒有結論的.
因為在聖經的內文當中.
真的沒有清楚講到究竟哪一個.
如果你看保羅書信.
保羅很清晰寫的.
我使徒保羅寫信比教會.

$^{121}$他在內文裡面很清楚講到.
道明自己的身份.
但是希伯來書裡面沒有特別.
提及作者是誰.
所以任何識經學家所講的.
都是一種猜測.
都是按著不同的論據來估計誰.
有些說法是使徒保羅.
不過這個很多都已經推翻了.
有些是在講亞歷山大的阿波羅.
又或者在講當代裡面有不同的教會領袖.
總之.
我們今天去看.
我們沒有辦法立一個很重要的結論.
究竟希伯來書是誰寫的.
除了作者不知道是誰之外.
讀者我們都不知道是誰.
你看希伯來書的書名叫做希伯來書.
有些剛剛信主的人問.
希伯來書不是放在舊約給希伯來人看的嗎.
不是.
希伯來書整個書卷.
原文所寫的是在當代裡面的希臘文所寫的.
希伯來書的讀者.
其實在聖經裡面沒有特別指明.
說是給哪一個背景的基督徒.
他的信仰狀況.
他的祖宗是怎麼樣的.
沒有特別去講的.
不過在隱含裡面我們大概知道.
這一群跟猶太人跟希伯來人有關的.
但是如果寫給希伯來人.
為什麼不用當代的希伯來語.
或者用亞拿曼等那些.
耶穌講的語言來寫呢.
無從稽考.
沒有辦法知道.
大概我們可以估計.
可能這一群的猶太人.
他們未必已經在聚居在巴勒斯坦.

$^{161}$或者以色列的那個位置.
可能他們已經散居在不同的外邦地.
他們帶著猶太人的血統.
他們都繼續承接著猶太人的某一些生活習慣.
他們的血統裡面可能都承載著一些這樣的生命.
不過他們對於語言.
或者那個世界觀已經是很希臘化的了.
在希伯來書的作者.
似乎他不是很介意.
這群人究竟他的血脈.
究竟是怎麼樣去承傳著猶太血統這件事.
反而希伯來書的作者在行文之間.
他就很著重這一群讀者.
他們自己的生命狀況是怎麼樣的.
有些字眼如果你有機會的話你去看.
其實在每一次的講論之後.
他有一些勸勉句的.
特別講到那些勸勉是講那些讀者的那個淑齡光景的.
比較負面一點的.
有些講他們說他們隨流失去.
有些說他們停止聚會.
有些是講他們的手是發酸.
發軟蹄.
信仰是講軟弱無力.
又或者有一些是走回去舊約那些.
走回那些懊悔.
那些不同的禮儀當中.
你會看到其實作者很關心的就是這一群人.
其實一直走回頭路.
他們的信仰變得軟弱無力.
這一群人對於聖經不是不懂的.
不過他們活出來的信仰就是.
面對很多挑戰.
很多艱難.
所以我們總結來說.
我們不知道希伯來書的作者.
沒辦法有一個結論性的.
究竟作者是誰.
讀者我們不肯定是哪一個教會.
一個具體的群體.

$^{201}$不過看他的勸勉.
很針對性他們面對信徒的問題.
大概可能是一群猶太基督徒.
他們已經離開了猶太的地方.
不在巴勒斯坦地了.
可能散居了.
甚至乎在不同的外邦的城市.
他們正在經歷基督教信仰.
他們對信仰有認識.
不過他們的生命依然面對世俗的衝擊.
顯得軟弱無力.
希伯來書就是在這樣的狀態之下.
用優美的希臘文來寫給當代的讀者.
也寫給今個世代裡面.
你跟我一起來看.
希望我們進入希伯來書之旅.
我們生命裡面更加踏實.
我們一起祈禱.
求上帝幫助我們.
讓我們明白話語.
同心祈禱.
上主願你的聖靈與我們同在.
幫助我們讓我們能夠.
在查考你話語的時候.
你的亮光光照我們的生命.
讓我們走在光明當中的時候.
我們能夠更加貼近你.
敏銳你對我們所說的話.
以致我們行事為人.
被你所閱立.
恭敬將來.
我們講解和聆聽的時間交託給你.
幫助我們每人.
讓我們打破不同的時空的合制.
願你的話語親自進入我們的內心.
幫助孫講講得清楚.
幫助聆聽聽得明白.
獻上祈禱.
奉求基督耶穌得勝的聖名.
Amen.

$^{241}$好.
在希伯來書裡面第一章第一節至到第四節.
這是一個引言的部分.
即使是引言的部分.
也已經是相當有力量.
而且很精彩的.
在一章一節至到第二節.
神記載古時直著眾先知.
多次多方曉遇列祖.
就在者末世直著他而止.
曉遇我們.
在第一二節裡面.
希伯來書的作者就鏗鏘有力地在說.
上帝一直在說話.
由古時直著眾先知.
曉遇以色列的列祖.
到今天直著基督耶穌.
曉遇我們.
很簡單的一句說話.
但是你會看到有一些對比在這裡.
首先他用的字眼.
他跟一般的書信不同.
他不是寫著.
我寫信給你.
或者是用文字的方式.
希伯來書的作者很刻意地在說.
是上帝說話.
上帝曉遇.
是上帝親自對人來溝通.
直著先知.
直到今個世代.
直著基督耶穌來對我們說話.
值得留意的是.
先知和耶穌基督.
他用神的遺址.
兩者有什麼分別呢.
在舊約裡面是說多次多方.
多次就是在說那個數量.
上帝不是一次地說給某一個先知聽.
全部都聽同一句話.

$^{281}$在舊約聖經裡面.
不斷地記載.
先知領受上帝的話語.
不停地在不同的時代裡面.
對不同的以色列民來說話.
所以那個次數是非常繁多的.
今天我們去看這句經文.
我們知道.
其實舊約先知怎樣對以色列民所說的話.
全部都記載在舊約的聖經裡面.
所以今天我們所體會的.
或者我們去理解.
先知舊約裡面多次的說話.
就記載在舊約當中.
同樣道理.
多方就是在說不同的方法.
舊約先知你不要只想著.
我要說你要聽.
有這些先知的.
不過有一些先知為了要令到當時.
那些百姓的心很硬.
又或者聽習慣了.
又或者上帝說話的時候.
他聽到已經有點麻木了.
那些先知要用不同的方式.
來引起百姓的注意.
有一些呢.
如果你有機會去看.
有一些舊約先知是行為藝術家.
他在城門口裡面做一些很奇怪的舉動.
又或者做一些當時受到別人非議的.
何西雅先知就是了.
故意要娶一個淫婦做為自己的妻子.
然後就成為話題了.
大家就會說這樣也可以的.
然後就成為話題之後.
先知就藉著這些話題.
這個百姓對他的注意.
他就說上帝的心意和訊息.
先知所用的方法.

$^{321}$是不同的方法.
有時候會用神蹟.
有時候會用他自己的經歷.
有時候會用不同的形態.
形態藝術的表現.
行為藝術等等.
目的就是要什麼呢.
不是要標旗立意.
而是要百姓要聽.
上帝有話要你知道.
在希伯來書的作者裡面.
簡單一句而已.
其實已經濃縮了整個舊約聖經裡面.
上帝一直對以色列民來說話.
不過到了我們今代.
到了當代去看希伯來書的讀者角度來說.
希伯來書的作者告訴他們.
現在不再是那個時代了.
現在是藉著神兒子一個.
來將上帝的話語.
去去曉喻他們.
舊約是眾多.
到了耶穌基督是獨一.
舊約是多次多方.
但是到了耶穌基督.
在他身上完全聚焦會聚.
在一個道成肉身的基督耶穌的身上.
在這個神的兒子身上.
一次性將話語表達出來.
更加重要的是.
舊約裡面的多次多方.
除了是上帝不斷提醒之外.
其實在重複的背後也有一些漸進的.
在神學裡面我們叫這個叫做漸進的啟示.
就是說上帝在這個時代說的信息.
隨著時代的一路進展.
他會在現有之前說的東西會加一點.
會加一點.
一直聚焦到基督耶穌的身上.
所以舊約先知除了說.

$^{361}$上帝對當代直轄先師說的上帝的心意之外.
也在說上帝一個很重要的救贖計劃.
就是將整個歷史.
將整個人類的生命.
匯聚在一個人身上.
這個就是神兒子耶穌基督.
所以那個多次多方.
其實就朝著一個中心點去的.
就去著耶穌基督那裡.
而耶穌道成肉身的時候.
就完完全全終極性地將上帝救贖的啟示.
直著他的生命呈現了出來.
你會看到這個比較.
其實除了說上帝不停說話之外.
他更加說到最氣肉.
最重要的信息.
就是直著基督耶穌.
讓他的生命來表達出來.
除了說剛才不同的內容.
就是那個聚焦性之外.
值得留意一個字眼叫做末世的曉喻.
直著神兒子在末世裡面的那個曉喻.
那個末世那個字眼.
我們今天一看.
你一打末世.
你在Google也好.
什麼也好.
你會看到很多東西的.
什麼印度神童等等那些.
不過聖經裡面不是滿足我們.
好像看荷里活片那種.
對末世的那種看法的.
這個末世其實在神學.
或者在聖經裡面所說的.
就是在說耶穌基督的降生.
和耶穌基督再回來的日子.
這一段的.
驅間這一段的時段.
聖經裡面就叫做末世了.
所以這個末世是在說.

$^{401}$如果去到我們今天.
就是在說耶穌回來之前的.
這一段我們生活的那個時代.
我們就叫做末世.
在經文裡面強調這個末世.
就是說到.
耶穌已經在說一個終末的信息了.
他的生命裡面流露了一個新的時代.
就是一個恩典.
藉著基督耶穌.
打開了一個恩典的大門.
讓外邦人都能夠得救.
讓上帝的計劃藉著基督耶穌來實現出來.
但是同時間.
這一段的時間不是無限量的.
基督耶穌再回來的就是結算的日子了.
如果這個結算的日子.
很多人都仍然在救恩以外.
或者很多人都是繼續信了.
但是就離棄了信仰.
當結算的日子.
他們的生命就會面對著和幻.
但是在這一段結算的日子.
他都依然保守著他的信心.
他的生命就會得到永恆的那種福樂和賞賜.
所以那個末世.
在說的其實是一個有時限性的.
這個曉如會有生效日期的.
到今天如果我們仍然在生效日期裡面.
我們接受活出這個信息.
我們就繼續保持著它.
這樣我們的生命是會得著福氣.
相反來說在生效日期裡面你放棄了.
到真的結算的日子裡面的時候.
我們生命就變得很危急了.
很困難了.
甚至乎已經是沒辦法.
再在救恩的門裡面經歷上帝的救贖.
簡單來說.
你會看到了.

$^{441}$救約先知的說話聚焦在耶穌身上.
耶穌就一次過.
綜合性地將上帝最重要的救恩信息告訴我們.
我們在這裡稍微思考一下.
我們知道上帝一直對我們說話的.
不過很多弟兄姊妹關心的.
上帝說的話很多時候關心什麼呢.
我們將來的老公是怎樣的呢.
我的工作好不好呢.
我買不買房子好呢.
我移不移民好呢.
我做路好.
我們在這些裡面尋求上帝的聲音.
對的.
敬虔的心態.
我們在人生做不同的重要決定.
我們都要邀請上帝.
在我們生命當中來做的.
這個我覺得我們非常之重要.
不過在尋求這些聲音.
在尋求上帝的指引之前.
我們必須要知道.
其實上帝說了很多.
整本聖經的說話我們不可以忽略.
在這個前提之下.
我們才去尋求上帝.
在我們個別每一個弟兄姊妹的心意.
如果我們只是尋求上帝個別的心意.
不去找.
或者不去明白.
原來上帝一早說了的話給我們知道.
我們就好像那種信仰觀念.
就好像和坊間裡面.
或者民間信仰就沒有分別了.
因為我們只是覺得那個神.
是幫我們催吉避凶.
或者幫我們做有智慧的決定.
但是我們不去想一下.
整個信仰的根基.
我們的根基在什麼那裡.

$^{481}$在希伯來書的作者裡面提醒我們.
我們今天.
我們要好好地有深度地.
來明白基督耶穌是一個怎樣的神.
祂用什麼身份.
祂用什麼角度來告訴我們.
這個信仰是多麼重要.
不要被這個世界很多的困難.
或者有很多的事.
來分散了我們的注意力.
我更加留意到.
或者我自己有時候也會這樣.
我們上網.
特別是世界凹凹亂亂的時間.
生命裡面面對很多困局.
香港面對疫情等等.
在這些位置裡面.
如果你留意一下那些什麼字眼搜尋榜.
末世這個字眼也排得挺高的.
大家會看到戰爭,災難這些.
在搜尋榜裡面也是一個挺高的位置.
因為那些人擔心,害怕.
不過很多時候.
我們會害怕,擔心.
我們寧願找外在的方法.
也不去看聖經究竟怎樣去教導我們.
怎樣去過我們的生活.
甚至有時候我們會將頂紙媒傳的WhatsApp.
一些很危急的.
又能幫助你的那些.
我不用考慮的.
很多人都是按個按鈕.
WhatsApp傳出去.
這樣我們就覺得心安理得.
又或者幫到人了.
不過有時候裡面說的話.
有時候大家我也覺得要fact check(事實查核)一下.
因為在過程裡面.
可能傳的那些是一些謬誤的信息.
最好笑你當笑話來看.

$^{521}$早前我又收回一個WhatsApp的.
傳遞回來的那些信息.
我們顧問牧師梁家倫.
用梁家倫牧師的名字來署名寫的.
其實這篇信息兩三年前.
都已經有了.
不要浪費疫情.
兩三年前已經在傳了.
到了第五波又再傳.
但是最好笑的其實這篇東西.
其實不是梁院長寫的.
但是我最初的原本.
我也有記得.
他寫梁家倫的家字也寫錯字的.
我懷疑其實可能寫的真的是家.
家頓的家.
梁家倫.
他最後寫的可能是這個人寫的.
不過人們說這麼有意思.
強行將他的家字改回家庭的家.
然後還要加上他的title(頭銜).
前建築神學院的院長.
然後重新包裝.
就再廣傳開去.
但是整件事情.
根本不是梁院長寫的整篇文.
但是大家瘋傳.
然後就說梁院長說得好.
其實那篇內容.
你認識梁院長之後.
其實不是一貫他說話的方式.
不過.
裡面的東西未必是錯的.
但是大家瘋傳.
或者大家很幫得上忙.
覺得如果用梁家倫這樣寫.
就會覺得真的有地位一點.
如果陳明全.
不要浪費我的RAM.
不要浪費什麼.

$^{561}$我們今天用托名方式.
來傳遞某一些資訊.
有些是有建設性.
有些是錯誤的.
但是如果我們沒有去查過.
就去在這個時代裡面.
瞎扯的把東西傳遞出去.
可能帶來的不是生命的果效.
在聖經裡面提醒我們.
在末世當中.
耶穌也在《馬達福音》裡面說.
面對著末世的日子裡面.
很多假的先知.
有很多偽的信息會傳遞出來.
你們不要被這些東西所迷惑.
反而要回到聖經.
看回上帝的話語.
究竟怎樣去幫我們安身立命.
而不是依靠外面的WhatsApp.
或者很多錦囊那些.
來過你自己的生活.
我不是說那些東西一定是錯的.
不過我們的生命.
如果我們都相信這些信息.
我們更加要相信一些.
對我們生命永恆有幫助的信息.
就是聖經的話語.
「神而子基督耶穌,曉遇我們」.
上帝今天在我們當中繼續的來說話.
這是第一,第二節裡面所說的信息.
我們繼續看第二至第四節.
「又早以納他為承受萬有的.
也曾藉著他創造諸世界.
他是神榮耀所發的光輝.
是神本體的真相.
常用他的權能命令托住萬有.
他洗淨了人的罪.
就坐在高天至大者的右邊.
他所承受的名既比天使的名更尊貴.
就遠超過天使」.

$^{601}$這是一章第二至第四節.
在這裡裡面.
希伯來書的作者.
很簡短,扼要.
但是用了七個小段片語.
來講到基督耶穌.
神子要講話了.
他用什麼身份,用什麼權柄來講呢?.
他用這七個小片語.
來將耶穌基督的身份.
他的兒子超越先知和天使的位置.
來表達出來.
我有個圖給大家.
這是《釋經學家》裡面.
他將那七個片語.
有一個循環來表達.
他在說什麼呢?.
第一就是承受萬有.
在說耶穌基督本身.
他的身份是承受萬有.
是一個世界擁有權的擁有者.
接著就講到創造.
接著就講到神本體的光輝和真相.
在這裡就開始.
由他自己的擁有權.
講一些比較外在性的.
創造啊.
或者他的光輝,榮耀.
接著第四就講到他掌管萬有.
第五就去到好像.
我的圖畫不夠好.
應該是最低點的.
他就是從他自己本身擁有天地萬物的身份.
好像一個很大最遠的距離就是.
道成肉身.
來到這個世界.
十字架裡面承擔.
結證了人的罪.
這個是人世間裡面最低的身份.
接著之後.

$^{641}$當他完成救贖之後.
上帝將他坐在至高者的右邊.
將他坐在高星.
接著之後第七個小段落就是.
承受比天使更尊貴的名.
你會看到這個循環.
一頭一尾的講承受尊貴的名和承受萬有.
承受這個字就講到他自己.
他本身就是擁有者.
這個世界是屬於這個聖子耶穌的.
不過他經歷了一個歷程.
這個歷程如果不講出來.
很多人都未必很清晰來知道.
希伯來書的作者就用一句又一句.
這樣來畫一個這樣的圓圈.
你用這個圓圈來理解就可以了.
整個救贖的計劃.
將基督耶穌是一個怎樣的神的兒子.
呈現出來.
這個救贖計劃只有耶穌才做得到.
也是因著他這樣來做.
令到他所講的話對我們生命是重要的.
因為他來到這個世界.
他就是要洗淨我們的罪.
當我們看了這個循環之後.
在第二章開始.
我今天沒有引用經文.
第五節至第十四節.
其實都會繼續解釋.
為什麼無端端講天使.
在第五節引用的舊約.
一直到第七節.
他比天使暫時微小一點.
視他榮耀尊貴為冠冕.
並將他所做的都派他管理等等.
這些舊約指著基督耶穌的經文.
希伯來書的作者就用這個循環來解釋.
在舊約裡面講到為什麼耶穌會比天使小一點.
有些古卷說他暫時比天使小一點.
本來他是高過他的.

$^{681}$因為他道成肉身.
承受這個世界的罪.
潔淨了所以他暫時比天使低.
在他的能力那邊.
然後上帝再將他高升.
所以他對於天使能力上低一點.
他只是暫時性的一種狀態.
希伯來書的作者講到.
我們所信的基督耶穌.
你不要純粹說他是神就很厲害.
他講到這個神不單止對你說話.
他也講到他的救贖.
他對我們每一個人來到這個世界裡面的使命.
為著你和我的罪.
他釘身在十字架.
他某程度是一種虛己自己.
用自己的生命來換取他所愛的世上所有的子民.
換取他們一個新造的生命.
而這個過程.
他的犧牲不是他的終結.
他的犧牲之後上帝再將他高升.
再一次承受萬有.
他就經歷過一個這樣的歷程.
來幫助我們.
所以當我們去看希伯來書.
作者告訴我們.
今天有一個很為你的主.
來對你說話.
這個主本身是有能力.
你根本拍馬都追不到.
我們是受造之物.
他是創造者.
你沒有辦法和他比擬.
不過這個高天的主.
竟然他藉著他的犧牲來救贖我們.
如果他對我們所說的話.
我們不去聽的時候.
我們就忽略了這個主對我們生命最重要的信息.
舊約裡面很多人說.
但是他們才知道.

$^{721}$他們自己的生命犧牲不到來為我們.
不過唯有基督耶穌.
他的犧牲十價.
他換取了我們一個新的生命.
這個就是希伯來書的作者裡面.
他要告訴我們.
今天告訴你的那個.
不是純粹一個冰冰冷冷的道理.
而是一個用生命來演繹出來的.
一個很重要的信息.
所以基督耶穌說的話.
我們必須要聽.
我們話鋒一轉.
我們跳去第二章一至四節.
在這個小段落.
其實就是某程度是第一章裡面.
有一些簡單的應用.
一些的勸勉.
二章第一節.
所以我們當越發鄭重.
所聽見的道理.
恐怕我們隨流失去.
那藉著天使所傳的話.
既是確定的.
凡干犯活逆的.
都受了該受的報應.
我們若忽略這麼大的救恩.
怎能逃罪呢.
這救恩起先.
是主親自講的.
後來是聽見的人.
給我們證實了.
神又按自己的旨意.
用神蹟,騎士和百般的義能.
並聖經的恩賜.
同他們作見證.
在這裡裡面他講到.
既然你知道講話的份量.
上帝給我們講了這麼多重要的訊息.
我們如何去回應呢.

$^{761}$第一個字眼.
我們當越發鄭重所聽見的道理.
好像有點奇怪.
其實意思是.
各位你們知道講話的那個是.
是這個主.
你們就打醒十二分精神來聽.
帶著豎起耳朵來聽.
那種豎起耳朵是代表著什麼呢.
假設有一個很重大的意外.
忽然間有個廣播告訴你.
那條逃生路是怎麼走.
無論怎樣.
那些人在說話.
你都要豎起耳朵.
不要吵.
我正在聽他們說什麼.
然後你就帶著這一份的專注度.
來聽究竟有什麼可以得生命.
那個越發鄭重是在說這個意思.
打醒十二分精神.
豎起耳朵來聽.
這個對你.
關乎你生命的道理.
然後作者再繼續的來講.
如果你不去打醒十二分精神去聽.
你就好像那些船在大海裡面.
面對著很大風浪的那些大船.
但是沒有下過撈.
下去穩妥船的那個位置.
如果那艘船沒有下撈的時候.
風吹雨打那艘船會怎樣.
會翻沉.
會吹去不知去那裡.
隨流失去就是在說這個意思.
如果你不打醒十二分精神.
來聽上帝的道.
特別是基督耶穌所說的那些重要的道理.
你就好像那艘船.
沒有下撈的船一樣.

$^{801}$隨著水流去到哪裡就去到哪裡.
隨時生命都受到很重大的威脅.
作者他要告訴我們.
今日基督耶穌說.
我們要謹慎的來去聆聽.
他再繼續講.
剛才講到那些讀者.
回到猶太的規條.
或者猶太人法理賽的那種生活.
那些頒佈下來的.
他們覺得除了先知之外.
更加重要是說天使頒佈的律例典章.
是由一些很重要的份量的人來講的.
作者說既然你對天使所說的都這麼緊張.
知道有報應就不聽.
干犯了上帝的道他都有報應.
你沒理由聽這些.
就不聽一些更加重要的信息.
就是基督耶穌所說的信息.
他講到這麼大的救恩.
如果我們忽略了它.
我們怎樣能夠逃避這個罪責.
怎樣逃避了救贖我們的恩典.
讓我們自己陷在罪當中.
在這裡我們知道是一個警告.
一個勸勉也是一個警告.
我再繼續去想.
其實很多頂展媒包括你和我.
其實很多時候.
我們對於福音聽了很多次.
在聖經我們看了很多遍.
又或者有一些道理我們都已經覺得是.
以熟能詳.
覺得我們已經懂了.
不過在我們生命當中.
這些懂了的道理.
在我們生命裡面的果效有多大.
有機會和一些監獄事供的牧者聊天.
他講了一個令我吃驚的觀察.
沒有實質數字的.

$^{841}$不過是觀察.
我那次進去有個男監.
進去監獄裡面做一些報導.
他說其實是沒得做的.
我說為什麼.
他說每個人都相信主了.
監獄信了主.
監獄事供都很厲害.
他說不是的.
很多人進監獄之前已經相信了主.
他講到監獄事供裡面.
他搞笑的不是沒得做.
而是很多進去面對牢獄的.
或者犯了事的人.
進去經歷牢獄.
其實他們很多都已經在年輕的時間.
回過教會.
知道福音.
或者在他讀書的生涯裡面.
他已經缺過字了.
甚至乎有些曾經是做.
教會裡面有些位置的.
不過這個世界太多洪流了.
再加上他所信的和他的生活.
他產生一種脫節感.
他不想將他的信仰.
活在他的生活當中.
反而是選擇在教會聽一套.
在生活裡面活另外一套.
以致到當他面對著一些挫折.
或者什麼.
他就選擇了這個世界錯的方法.
來去過他的生活.
拿錢或者用不同的東西來犯事.
很多.
這個牧者告訴我.
大概有七成的很多.
曾經在監獄裡面的.
曾經缺過自信的主.
這個數字究竟是不是真的我不知道.

$^{881}$不過他跟我說.
我也不覺得在經驗上這個是錯的.
在我們基督徒人生裡面.
我們很多道理我們聽過.
我們都甚至乎.
留下很多很感性的眼淚.
我們為著這個信仰.
為著基督耶穌.
我們都可以拋頭顱灑血.
說出來可以這樣.
但是去到生活可能我們都放下了.
還是用這個世界的方式.
過我們自己的生活.
希伯來書的作者提醒我們.
如果你知道.
講話給你聽的那個是.
是一個這麼有份量的神子基督耶穌.
你不要逃避.
你不要放棄.
你也不要輕視了這個所講的內容.
最近有機會看一個教育家.
叫做Pakir Palmer.
基督徒教育家.
他寫了一本書叫做.
《隱藏的整傳朝向不可分割的生命》.
這本書很有趣的.
很有趣的.
有機會借給大家看.
不過他後面講到.
有很多人在這個世代裡面.
人的生命是支離破碎的.
很片段式.
或者後現代的那種特色.
人的生命是不整傳的.
很強調這一點.
不過Pakir Palmer他說.
生命不是這樣的.
要是整傳的.
今天令到生命這麼破碎的原因.
其中一樣東西就是聆聽出了問題.

$^{921}$有什麼東西阻礙了我們聽到聲音.
但是聽不到內容呢.
他第一樣東西是講到忽略.
這句也是希伯來書的作者所講的.
我們知道那些東西重要.
但是我們刻意迴避了它.
搖頭一動.
當作沒有聽過.
那個叫做忽略.
忽略這個救恩.
代表著你知道.
但是你不將這件事貫徹在我生命當中.
在Pakir的講法裡面.
這個忽略是令到我們聽不到內容.
經歷不到我們所信基督教的那種豐富.
這個令到我們好像屬靈裡面失明一樣.
第二樣東西.
他講到很多時候我們人生命裡面會有些過濾.
我們覺得這些只是口水.
自動耳朵關起來.
可能你每個星期聽YouTube講道理都會這樣.
飛快一點.
Anders講得這麼慢.
飛快一點.
快速地覺得.
或者選擇性地聽一些覺得是新穎的內容.
不過那個忽略.
或者那個過濾.
令到我們聽的東西變得很片面.
片面的信息就製造了片面的人生.
以致我們不能夠活出一個整全.
對信仰完完全全經歷的那種生命.
第三樣東西他講的就是抗辯.
那些人可能跟你本身有牙齒印.
總之他說什麼.
你就覺得一定是錯的.
因人而廢言.
因為你本身有個立場在這裡.
你自己帶了你自己對那個人的看法.
或者對了你自己的想法.

$^{961}$你就當了你說的是真理.
你覺得你自己那個是真理而且是結論性的.
當其他有不同的說法說進來的時候.
你第一樣東西就是撐.
而且你就要拒絕.
帶一個這樣的抗辯.
未必一定要破壞關係.
不過這種抗辯或者心裡面有一種覺得.
我已經所知的結論性的這種態度.
其實是一種很驕傲的態度.
以致我們的耳朵聽不到還有很多.
我們不知道或者我們忽略了的.
一些生命裡面很重要的真理.
又或者可能你會覺得那個牧者.
那個同綱或者那個弟兄姊妹說的那些.
都是你已知的東西.
你很強調你自己今天我Here I stand.
我自己這個位置.
就抗拒了他給的一些真誠的勸勉.
這種抗辯抗拒的態度.
也都令到我們耳朵屬靈裡面.
就屬靈的失明.
柏球他再繼續說.
今天靈魂.
你生命裡面你的靈魂需要些什麼呢.
不是需要聽一些你自己覺得對的.
或者你喜歡的東西.
那些很瑣碎的.
或者你用很多東西填塞你.
你的腦袋的那些資訊.
是一些瑣碎的東西.
靈魂不是需要這些東西的.
靈魂是需要真理.
因為這個真理會救贖你整個的生命.
當你忽略了這些真理的時候.
你的生命就陷起一種屬靈的失衝.
又或者甚至乎你已經偏行幾路.
過你自己想過的生活.
沒有了上帝在我生命裡面.
左右我們所走的路.

$^{1001}$求主幫助我們.
當我們生命裡面如果變得偏聽.
帶著自己的立場來抗拒別人所說的話.
我們就沒有辦法經歷到信仰的那種豐富.
柏卡不是純粹批判的.
他在說有一個很重要的應用.
他說要逃避或者要改變這種聽不到東西的困難.
他說一個人聽是有困難的.
你就要一班人一起聽.
他用一個概念叫信任圈.
這個是寫東西裡面一些重新去定義.
其實所說的信任圈就是我們教會的小組.
不過這個信任圈他在說一些很重要的生命特質.
說和聽大家都要帶著一份的專注和尊重.
在別人所說的時候.
我們不要這麼快去批判他.
給他建議.
或者帶著自己的情緒來聽別人所說的話.
盡力讓自己專注聽別人經歷的信仰.
聽足夠了才給你心底裡一些真誠的建議給他.
今天我們就是建立不到信任圈在教會裡面.
原因就是我們太快為別人作出一些結論.
我們太快覺得我已經掌握了別人的生命.
又或者我們太快為自己所說的話.
我們覺得自己已經是完備.
在帕克的建議裡面.
教會成為一個信任圈.
讓人可以安心說他自己心底的說話.
也都讓這個群體可以容納真理.
聆聽別人的故事.
唯有這樣來建立一個信任圈.
我們一起集體恭敬地聆聽.
我們才可以好好地活出這個信仰.
所以各位弟兄姊妹.
可能你是在這段疫情裡面.
你初次來參與旺朕的這個網上崇拜.
又或者你可能過往裡面.
你自己在經歷信仰生活有些失望.
你想離開教會 離開這個實體教會.
你覺得上網聽就足夠.

$^{1041}$我想說不足夠的.
如果信仰只是我們每個星期打開YouTube.
純粹用這個虛擬的方式.
我們只是捕捉了一部份.
但是生命裡面集體的聆聽.
集體經歷那種信任塑造和改變更新.
你就會錯失了這個.
信仰裡面這麼重要這麼大塊的信息.
更加重要的是.
教會是基督耶穌的身體.
是聖子耶穌直著這個群體.
將真理傳遞開去.
當我們專心聆聽的時候.
不單只是說我們聽到聲音.
在這個群體裡面我們會聽得到內容.
這個小段的結束.
作者說道 神安著自己的旨意.
用神蹟的旗幟和百般的賢能並聖靈的恩賜.
和他們作見證.
這個見證這個字跟剛才第一章說的有點不同.
剛才說到先知多恃多方來講說話.
但是去到第二章第三第四節.
人會應上帝最好的方式.
是見證.
是證實基督耶穌所說的話是真的.
在這個群體裡面.
我們一起來經歷.
一起來見證耶穌基督在這個世代裡面的工作.
也希望我們今天帶著一份恭敬的心.
豎起耳朵聽聽對我們說話的那位基督耶穌.
在這個世代裡面.
在旺角浸信會這個群體.
對我們要說什麼話.
求主幫助我們.
謝謝大家.
\newpage



\section{希伯來書 4:14-5:10}
\label{sec:bMqVNP_dN2s}
\textbf{2022年3月20日崇拜直播｜陳冠成牧師｜超越的基督:完全的大祭司｜希伯來書4\_14-5\_10}
\newline
\newline
連結: \href{https://youtube.com/watch?v=bMqVNP-dN2s}{\texttt{https://youtube.com/watch?v=bMqVNP-dN2s}} ~~~~ 語音日期: 2022-03-20
\newline
\newline
\hyperref[sec:cYMHT7F9Rko]{\small{< < < PREV SERMON < < <}}
~
\hyperref[sec:index_chronic]{\small{[返順時目]}}
~
\hyperref[sec:XlsH8Zitzyg]{\small{> > > NEXT SERMON > > >}}
\newline
\newline
希伯來書 4:14-5:10
\newline
\begin{longtable}{cl}
\hline
\hline
章節 & 經文 (和合本修訂版)\\
\hline
4:14 & \begin{tabularx}{0.7\textwidth}{X} 既然我們有一位偉大、進入高天的大祭司，就是耶穌—神的兒子，我們應當持定所宣認的道。 \end{tabularx} \\ \\ \relax
4:15 & \begin{tabularx}{0.7\textwidth}{X} 因為我們的大祭司並非不能體恤我們的軟弱；他也在各方面受過試探，與我們一樣，只是他沒有犯罪。 \end{tabularx} \\ \\ \relax
4:16 & \begin{tabularx}{0.7\textwidth}{X} 所以，我們只管坦然無懼地來到施恩的寶座前，為要得憐憫，蒙恩惠，作及時的幫助。 \end{tabularx} \\ \\ \relax
5:1 & \begin{tabularx}{0.7\textwidth}{X} 凡從人間挑選的大祭司都是奉派替人辦理屬神的事，要為罪獻上禮物和祭物。 \end{tabularx} \\ \\ \relax
5:2 & \begin{tabularx}{0.7\textwidth}{X} 他能體諒無知和迷失的人，因為他自己也是被軟弱所困， \end{tabularx} \\ \\ \relax
5:3 & \begin{tabularx}{0.7\textwidth}{X} 因此他理當為百姓和自己的罪獻祭。 \end{tabularx} \\ \\ \relax
5:4 & \begin{tabularx}{0.7\textwidth}{X} 沒有人可擅自取得大祭司的尊榮，惟有蒙神所選召的才可以，像亞倫一樣。 \end{tabularx} \\ \\ \relax
5:5 & \begin{tabularx}{0.7\textwidth}{X} 同樣，基督也沒有自取作大祭司的榮耀，而是在乎向他說話的那一位，他說：「你是我的兒子，我今日生了你。」 \end{tabularx} \\ \\ \relax
5:6 & \begin{tabularx}{0.7\textwidth}{X} 就如又有一處說：「你是照著麥基洗德的體系永遠為祭司。」 \end{tabularx} \\ \\ \relax
5:7 & \begin{tabularx}{0.7\textwidth}{X} 基督在他肉身的日子，曾大聲哀哭，流淚禱告，懇求那能救他免死的神，就因他的虔誠蒙了應允。 \end{tabularx} \\ \\ \relax
5:8 & \begin{tabularx}{0.7\textwidth}{X} 他雖然為兒子，還是因所受的苦難學了順從。 \end{tabularx} \\ \\ \relax
5:9 & \begin{tabularx}{0.7\textwidth}{X} 既然他得以完全，就為凡順從他的人成了永遠得救的根源， \end{tabularx} \\ \\ \relax
5:10 & \begin{tabularx}{0.7\textwidth}{X} 並蒙神照著麥基洗德的體系宣稱他為大祭司。 \end{tabularx} \\ \\
[1ex]
\hline
\hline
\end{longtable}
$^{1}$今日重讀的經訓是記載在希伯來書4章14節至5章10節.
希伯來書4章14節.
我們既然有一位已經升入高天尊榮的大祭司.
就是神的兒子耶穌便當持定所承認的道.
因我們的大祭司並非不能體恤我們的軟弱.
他也曾凡事受過試探.
與我們一樣 只是他沒有犯罪.
所以我們只管坦然無懼地來到施恩的寶座前.
為的是要得連率 蒙恩維 作隨時的幫助.
第五章.
凡從人間挑選的大祭司是奉派替人辦理屬神的事.
為要獻上禮物和贖罪祭.
他能體諒濫愚蒙的和失迷的人.
因為他自己也是被軟弱所困.
故此他理當為百姓和自己獻祭贖罪.
這大祭司的尊榮沒有人自取.
惟要蒙神所召 像亞倫一樣.
如此 基督也不是自取榮耀作大祭司.
乃是在乎向他說.
「你是我的兒子 我今日生你的那一位」.
就如經上又有一處說.
「你是照著墨跡使得的等次永遠為祭司」.
基督在肉體的時候.
既大聲哀哭 流淚禱告.
懇求那能救他免死的主.
就因他的虔誠蒙了英雲.
他雖然為兒子.
還是因所受的苦難學了信從.
他既得以完全.
就為凡信從他的人成了永遠得贖的根源.
並蒙神照著墨跡使得的等次.
稱他為大祭司」.
今日為我們宣講上帝的道.
是陳冠誠牧師.
今日的講題是.
「超越的基督 完全的大祭司」.
丁姐妹平安.
很快就已經來到希伯來書的第五章.
如果你發現其實希伯來書有一個特色.
就是作者會將讀者所認識的.

$^{41}$很多舊約的事物.
與耶穌基督一一比較.
之前那幾章其實你已經看到.
已經顯明了耶穌基督的分別.
比起眾善之.
所有的天使 摩西 甚至約書亞.
都更為尊貴和優勝.
現在作者亦都將耶穌基督.
與舊約的大祭司作一個比對.
強調耶穌不單止具備.
大祭司的身份和職份.
並且比他們都超越.
更值得信徒來依靠.
事實上自從出埃及以來.
大祭司在以色列人的心目中.
可以說是在地上最高的宗教權威.
大祭司代表著神和人之間的中間人.
或者我們常說的一個中保的角色.
他既代表百姓向上帝獻屬罪祭.
並且尋求上帝的寬恕.
他同時亦代表了上帝宣告百姓.
獲得上帝的接納和赦罪.
所以為了這個緣故.
經文裡面其實看到.
面對著當那些可能因為無知愚昧.
或者迷失而得罪上帝的百姓.
當他們來到大祭司面前的時候.
大祭司會好好管理自己的情緒.
他不會對百姓過分嚴苛.
總是能以一個寬宏體諒的心.
來接納他們.
畢竟我們很明白.
假如大祭司對每一個犯罪的人.
都充滿憤怒甚至嫌棄的話.
其實真的很難心甘情願.
來替他們獻祭並且代求.
至於大祭司為何能夠體恤百姓的軟弱.
經文裡面亦很清楚交代了.
是因為他自己.
其實亦都像他所代表的百姓一樣.

$^{81}$他會有軟弱所困的時候.
甚至乎會被試探所剩的時候.
我記得我自己在讀小學的時候.
曾經被老師來到去揀選.
維持科堂秩序.
特別在可能轉堂的時候.
以前學校有些規矩就是.
不讓到處走.
不讓聊天.
老師就指派了.
轉堂的時候就負責幫我.
記下那些不守規矩的同學的名字.
到時交回給老師.
讓老師去處理.
起初覺得很飄飄然.
挺威風的.
因為老師看好你.
賞識你.
所以找你幫忙維持秩序.
但後來就開始覺得.
做著做著都挺羞愧.
為甚麼呢?.
因為後來自己都怎樣?.
不守規矩.
自己都被人記下名字.
所以就會覺得.
自己其實很不配做這個角色.
明明負責維持秩序.
但卻又不守秩序.
很體會到其實自己和其他人一樣.
都會有軟弱.
行差踏錯的時候.
同樣經文這裡說.
雖然大祭司他的身份和職份.
是很尊貴的.
他是奉派替人辦理.
要屬上帝的事情.
但本質上.
他和所有的百姓一樣.
都是罪人.

$^{121}$他怎樣為百姓獻上贖罪祭.
他也照樣為自己獻上.
他怎樣為人向上帝求情.
他也需要首先祈求上帝.
赦免自己的過犯.
所以經文這裡說.
他之所以獲得大祭司的榮耀.
不是靠自己可以贏取回來.
而是出於上帝的揀選和呼召.
將尊榮加在他身上.
就正如以色列人的第一任大祭司.
亞倫他也是這樣.
同樣經文繼續說.
看回耶穌基督.
他大祭司的身份和尊榮.
也不是自己爭取回來的.
作者你會看到.
在聖靈的引導之下.
他分別引用了兩段詩篇.
詩篇第二章的七節.
和詩篇一十篇的第四節.
來表示其實耶穌基督.
也是蒙上帝所選召.
以上帝兒子的身份成為大祭司.
並且他很強調耶穌基督.
不是按照利美茲派亞倫的體系.
來承繼大祭司的職份.
而是照著默基使德的等詞.
永遠作大祭司.
什麼意思呢?.
在這裡很簡單地交代默基使德的等詞.
默基使德的背景.
其實在《創世記》十四章裡.
當時你會發現亞伯拉罕還未改名.
還在叫「亞伯蘭」的時候.
換句話來說.
不要說亞倫.
說的是利美茲派都還沒有的時候.
默基使德就已經在撒冷.
即耶路撒冷這個地方作王.

$^{161}$並且他也是至高神的祭司.
我們看到默基使德可以為亞伯蘭祝福.
接受他的十一奉獻.
很明顯你會看到.
誰的位分高一點.
當然是祝福人的位分大一點.
所以你看到默基使德的位分高於亞伯蘭.
更加比起亞倫高了多少班.
換言之.
耶穌基督既然照著默基使德的等詞作大祭司.
也就是說.
雖然沒錯.
他和亞倫等人間的大祭司一樣.
都是做大祭司.
但是論字排輩.
耶穌基督大祭司的等級和等次.
遠遠超過亞倫和他之後的所有任人.
這是作者很想表達的.
耶穌基督的超越.
而另一方面.
耶穌也曾經以人的身份.
去面對世上各種的試探.
並且他沒有犯罪.
沒有離棄上帝.
所以和人間的大祭司很不一樣.
耶穌基督的超越在於.
他既能夠體恤人面對罪惡.
面對試探那種軟弱.
但是同時他又不被試探軟弱所困.
他不單單不需要像人間的大祭司一樣.
替自己獻贖罪祭.
並且我們知道.
他自己更加成為了人永遠的贖罪祭.
永遠成為了神和人之間的鐘堡.
將人徹底在罪惡中拯救出來.
所以亦是因為這個緣故.
作者他勉勵正是.
面對著不同的試煉,試探.
甚至迫不得已的讀者.
你們要懂得選擇.

$^{201}$到底是任由自己被軟弱,罪惡.
或者試探所勝過.
選擇可能太辛苦了.
不如走一條舒服一點的道路.
走一條闊一點的道路.
萬一真的跌倒,衰落.
那就靠人間的祭司,大祭司.
來幫忙獻贖罪祭給上帝.
這樣就可以繼續保持著和上帝的關係.
事實上這個方式.
沿用了很多年.
自從有大祭司的制度.
這個就是他們的挽回機制.
但是作者勉勵他們.
既然有更超越的大祭司已經臨到.
已經成就一切的時候.
我們到底還是選擇那條舊路.
還是堅持認耶穌基督為主.
抓緊祂的應許.
幫助我們過渡人生一些最艱難的時候.
來到竭力遵守主導討上帝的喜悅呢?.
這是作者勉勵每一位讀者.
包括我們要去揀選抓住的道路.
我記得曾經有一位姐妹跟我分享.
就是在很多年前.
她的先生對她不忠.
認識了一個小三.
有外遇了.
結果就想和她分開.
甚至搬出去住.
這位姐妹覺得很掙扎,很苦惱.
她很想放棄.
畢竟不是自己出了問題.
她曾經想過算了.
還是不要繼續.
反正不是我的錯.
為什麼我要這樣忍她.
這樣被她欺負.
為什麼要這麼辛苦.
但是在她最軟弱,最無助的時候.

$^{241}$聖靈卻提醒她.
提醒她當日和丈夫.
在上帝面前所許下的婚姻承諾.
提醒她.
即使她的配偶走差踏錯.
軟弱跌倒.
但她怎樣仍然堅守.
她在上帝面前所許下的承諾.
所以這位姐妹最後做了一個決定.
沒錯,她沒有能力阻止對方離開.
但她也堅持不離婚.
並且她同時也求上帝.
在她的配偶身上動功.
日子不容易過.
但最後她的配偶卻是迷途知返.
回到家中和她一起生活.
確實仍然有很多.
在他們二人之間的婚姻相處.
關係破碎的裂痕要修補,要處理.
但上帝卻幫助她.
過渡了最艱難的時刻.
當她在最困難的時候.
她在捉緊上帝的應許和幫助的時候.
她就可以勝過那些試煉和試探.
同樣提起姐妹珍.
確實當我們在生活中面對一些.
似乎信仰也幫助不了我們解決困難.
甚至我們覺得上帝也不出手幫助我們.
他又不聽我們禱告的時候.
我們心裡很自然就會有些想法.
不如我們也試試跟自己的意思去做.
試試不聽上帝的禱告.
可能會更舒服一點.
我們不如靠人的智慧.
來解決眼前的困難.
這不單止是讀者當時面對的試探.
其實今天我們同樣面對.
不過希伯來書的作者卻提醒我們.
不要再亂砍回到舊路上.
或者好像今年的主題是6G.

$^{281}$假如真的已經有6G面世的時候.
你沒理由回去還用2G和3G.
確實日子越艱難的時候.
聖經勉勵我們.
你越是要捉緊耶穌基督.
這一位完全的大祭司.
這個才是正路,這個才是出路.
特別在第七節裡.
作者引用耶穌基督在.
赫西瑪利云的經歷.
來提醒我們.
為什麼耶穌真的可以幫我們.
勝過試探和持守真道.
第七節開始記載.
那位道成肉身的耶穌基督.
祂曾經大聲痛哭.
流淚禱告懇求那能救他免死的主.
也因為耶穌基督的虔誠.
祂的禱告蒙了英雲.
我們還記得在科目書裡.
記載耶穌基督在赫西瑪利云的三次禱告.
其實都是重覆著一個重點.
我相信大家都可以說出來.
就是父啊,假如可以的話.
求你免去我釘十架這個苦杯.
不過不要照我的意思.
照父你的意思來代行.
其實耶穌基督難道不清楚.
父神的意思嗎?.
他絕對清楚.
他從小聖經記載.
他就以父的事為念.
他也很清楚知道自己來到世上.
是要背上十字架.
替人類來受罪,受死.
但是他也會復活.
並且你留意在科目書裡.
他已經一早跟門徒預告了.
這件事情一定會發生.
並且不只預告一次,是三次.

$^{321}$可以說在理性上,頭腦上.
甚至心態上.
其實耶穌基督似乎都做好了準備.
不過,在赫西瑪利云裡記載.
當那個苦難真的逼近那一刻的時候.
身為人的他和我們原來一樣.
都會出現心靈固然願意.
但是肉體卻是軟弱的時刻.
耶穌基督他靠什麼來勝過自己這種脆弱呢?.
就是很多時候最容易被我們忽視.
覺得好像做了也沒什麼用的方法.
就是祈禱.
你留意耶穌雖然祈求父神.
拿走祂的苦杯.
解決祂眼前的苦難.
但是祂同時也不斷地高舉.
父上在祂生命裡的主權.
並且祂求上帝加添祂的力量.
幫祂走完這條苦路.
結果,祂的禱告蒙上帝應運.
也就是說,在祂等候被出賣,被捉拿的過程裡.
耶穌是一面掙扎,一面來禱告.
一次不夠,就兩次吧.
兩次不夠,三次.
持續的禱告,你留意.
沒有改變到客觀環境.
苦難仍然是越來越近.
但是持續的禱告裡.
上帝卻是幫助耶穌可以保持警醒.
讓祂可以信服.
留下在赫西瑪利園.
祂沒有逃跑.
直到那些出賣,捉拿祂的人來到.
所以你不難想像,假如.
那些出賣,捉拿耶穌的人.
一天還沒到的時候.
耶穌就會怎樣?.
可能不止三次.
可能有第四次,第五次.
甚至更多的禱告.

$^{361}$直到赫西瑪利園的主義在祂身上成就為止.
換句話來說.
弟兄姊妹,赫西瑪利園的祈禱.
不單止展示了耶穌基督.
祂有人性脆弱的一面.
更重要的是耶穌教會我們.
我們可以怎樣面對那些負面情緒.
那些不合上帝心意的想法.
特別當面對困難的時候.
很多時候我們聽過不同的人.
會這樣來鼓勵我們.
你不要想得太過負面.
正面一點.
要保持正向的思維.
但是我們也很明白.
那些負面感覺.
那些負面想法和情緒.
有時候不是說停就停.
想不想就不想.
有時候你越不想想的時候.
它越是會不自覺地出現在你的腦海裡.
特別可能在這段日子.
我們也很體會到.
對某些人來說.
越是叫他不要想疫情.
你不要擔心.
不用怕自己感染.
也不要擔心前途會是怎樣.
其實回想起來.
幫助不一定很大.
那些負面情緒和想法.
仍然會隨著他每天所經歷的事情.
他所接收到的資訊.
繼續不請自禮.
很影響他的身心健康.
甚至會癱瘓他的生活.
在這些情況裡.
我們不禁會問.
其實向上帝祈禱有幫助嗎?.
頂尖會有幫助嗎?.

$^{401}$聖經的答案是肯定的.
一定有.
如果不是.
他也不會讓我們看到.
耶穌在黑西瑪利園的這一幕.
我們看到耶穌.
祂是否靠著禱告.
來面對自己那些.
不合上帝心意的負面想法嗎?.
在整個過程裡.
我們看到耶穌透過禱告.
其實不斷將那些負面.
不合上帝心意的想法.
一路外化出去.
不讓他們再繼續控制.
蠶食自己的心靈.
不讓他們影響耶穌基督.
來進行上帝的旨意.
不單止這樣.
你看到耶穌的禱告.
也是同時不斷將上帝的旨意.
內化在自己的心裡.
一路去按實它.
來進一步削弱.
自己內裡的負面感受和想法的威力.
使自己有更多的力量.
來貫徹實踐.
以負的事為念.
這個人生目標.
所以弟兄姊妹.
從耶穌基督的親身演繹.
我們看到.
禱告是絕對有功效.
關鍵是我們到底向上帝求什麼.
我們純粹一遇到困難.
就求上帝幫我們去解決.
免我們受苦.
丁冠仕.
我們效法耶穌基督.
他在黑西瑪利園的禱告.

$^{441}$一方面向上帝.
傾訴自己的軟弱.
自己的黑暗面.
自己內心的種種掙扎.
但同時我們祈求上帝.
他自己的更超越的力量和旨意.
臨到我們身上.
進到我們的心裡.
掌管我們全人的心思意念.
以致我們真的能夠.
好像耶穌基督那樣.
順服上帝帶領的安排.
直至上帝的旨意.
在我們身上完全的實現.
這是耶穌基督.
他親身對我們的示範和教導.
正如在第八節裡.
你會看到.
耶穌基督雖然是上帝的兒子.
他和神擁有這麼獨特親密的關係.
但是這個關係.
不單止沒有令他在順服上帝這件事上.
得到豁免.
反而是怎樣呢.
他和父上帝這種親密關係.
反而是要求他活出一個比任何人更大的順服.
這裡說.
基督因苦難學了順服.
留意其實他的意思不是說.
耶穌基督起初其實就.
不太懂得什麼叫順服上帝.
或者說不是很.
知道自己要順服上帝.
所以要慢慢學.
好像我們經常俗語說.
經一事長一智.
耶穌基督的順服是怎樣呢?.
是經一苦長一智.
從頭到尾學回來的.
從不懂學到懂.

$^{481}$直至完完全全順服上帝.
並不是這樣的意思.
這裡說基督因苦難學了順服.
是在說.
透過接受十字架這個受苦的經歷.
令到耶穌順服的生命.
達到一個最巔峰,最完全的境界.
讓他真的徹徹底底地.
體會到順服神.
實際上是什麼滋味.
不是只是一個頭腦.
或者一個口號的認知.
而是他很真實地.
通過一個終極的考驗.
來將他的順服達至巔峰.
而他這個順服的目的和果效.
就是為每一個凡順服他的人.
奠定了一個永遠得救的根基.
說的不是純粹我們信耶穌.
想天加這件事.
那個根基是說.
耶穌基督他好像一個先驅一樣.
他自己帶頭.
不單單他自己完全順服上帝.
他可以幫助之後每一個跟隨他的人.
都順服上帝.
每一個納至認耶穌為主.
與父神和好的人.
無論在他的靈魂生命裡面.
又或者在他生活日常的裡面.
甚至在他的工作事奉裡面.
都可以徹底經歷到.
耶穌基督他的拯救大能.
就是有可以不落入罪惡的試探網絡的可能.
並且都有力量來到他持守真道.
完成上帝所託付給他.
在人生不同崗位裡面.
每一個的使命.
這個就是經文所說.
耶穌基督他之所以能夠成為.

$^{521}$永遠的天上大祭司.
是因為每一個跟隨他的人.
耶穌都可以徹底拯救他們.
離開罪惡的網絡.
並且過一個土生起月的生活.
所以弟兄姊妹.
我們是否渴望耶穌基督這一份拯救大能.
這是經文想我們去思考.
假如是的話.
我們就要效法耶穌基督.
不要再藏起自己的軟弱掙扎和黑暗.
不要再掩飾在心裡.
不跟耶穌來傾訴.
你留意.
在第四章十六節裡面.
經文鼓勵我們.
只管坦然無懼來到私人的寶座前.
為要得連續夢因為作隨事的幫助.
這個的應許.
是上帝賜給歷世歷代的信徒.
至今仍然有效.
耶穌這位體恤我們的大祭司.
他自己也親身嘗試過.
他也曾經想逃避在十架的苦杯.
他曾經和我們一樣.
有著人性的軟弱.
但是他最終也勝過.
並且他可以幫助我們.
好像他一樣.
勝過我們人生所面對的種種試探.
不過可能會有人去疑惑.
耶穌基督其實他生活在二千多年前的世代.
他是不是真的明白.
我受的苦.
我所面對的困難和艱辛呢?.
請問我們會不會也有這個疑惑.
覺得耶穌其實也不理解我們的困難呢?.
在福音書裡面曾經提醒過我們.
耶穌基督在他出來事奉.
在他出來傳道之前.

$^{561}$他首先通過的是撒旦的試探.
那裡說的是四十天的試探.
撒旦用盡不同的方法.
來攻擊他.
來引誘他.
你不難想像那四十天的角力.
撒旦一定無所不用其極地.
務求要將耶穌拉下來.
務求要讓耶穌不聽從上帝的吩咐.
用盡所有的方法手段.
來攻擊試探耶穌.
換句話來說其實我們可以放心.
我們所面對的苦難.
我們所面對的困境.
耶穌一定體會過.
他一定可以理解到.
因為他所受的試探和苦難.
可以告訴我們有過之而無不及.
他一定懂得你和我的內心掙扎是什麼.
他會體諒,他也會連述.
我記得有一位牧者.
他曾經分享過.
他和他的小朋友的一個經歷.
就是小朋友小時候.
有一個尿床的習慣.
很多時候睡到半夜.
不自覺地賴濕了.
不單止弄髒了衣服.
還弄髒了床單.
於是牧師很多時候.
會被小朋友吵醒.
因為每次小朋友尿床之後.
就會敲門叫醒牧師.
「尿了」.
牧師就會上下升嵩地起來.
幫他慢慢清潔.
換過衣服和床褥.
第二晚又是這樣.
第三晚又是這樣.
不容易,因為很多時候.

$^{601}$會令自己睡得不夠好.
後來過了幾天之後.
突然有一晚牧師一睡就睡到天亮.
他覺得很好.
「這麼好,今晚沒有尿床」.
於是他就去兒子的房間拍門.
他見到小朋友醒了.
為何不起來?.
躲在被窩裡.
牧師發覺他臉上有顏色.
於是就揭開被子.
突然一陣很特別的味道撲出來.
其實原來小朋友也有尿床.
牧師很奇怪.
問「為何你不叫醒我?」.
「我睡了一整晚」.
兒子覺得很羞愧.
因為他不想讓父親知道.
怕不好意思.
怕又要令爸爸很不方便.
又要幫他換衣服,換床單.
令爸爸睡得不好.
牧師對兒子說.
「傻子,你要跟我說」.
「你這個困難只會短暫」.
「隨著你的成長」.
「一定會改變到,改善到」.
「你不跟我說」.
「整晚黯然無言」.
「雖然我睡得好像很好」.
「但其實我還不開心」.
同樣的,姐妹們.
真的.
千萬不要將你心裡的幽暗.
那些脆弱.
那些覺得你很醜.
不想讓人知道的事情.
掩蓋著,不讓耶穌知道.
我們軟弱,得罪了上帝.
耶穌固然不開心.

$^{641}$但是,如果我們掩蓋著.
不告訴祂.
迴避祂,當沒事發生.
又或者想自己去解決.
其實就會令耶穌更加心痛.
所以,真的.
哪怕是我們生命中.
最脆弱,最黑暗,最羞愧的事情.
耶穌不介意.
祂很想我們告訴祂.
你只管坦然無懼.
來到祂的面前.
不要掩蓋著.
向祂傾訴.
最後,很想用一個故事.
來結束今天的訊息.
話說有位牧者.
他拿著一杯水.
半杯水.
他上到台上準備分享.
台下的觀眾一定以為.
他想說一個他們聽過的訊息.
到底你看到的這杯水.
是沒有了半杯還是還有半杯呢?.
但沒想到這位牧師.
他把這杯水舉起.
他就問.
大家覺得這半杯水有多重呢?.
台下的人紛紛說.
四安士左右,八安士.
有些人說二十五安士.
於是這位牧師就分享.
其實這杯水有多重.
在於我們拿的時間有多長.
如果你只拿一分鐘.
這杯水真的一點也不重.
如果你拿一個小時.
你的手可能會有一點麻痺.
如果你拿著這杯水超過一天的話.
我相信你可能要去看鐵打.

$^{681}$或者問你治療.
因為你的手可能已經硬了.
甚至乎肌肉可能已經受損.
雖然這杯水的重量沒有變過.
但時間拉長了.
這杯水對我們來說.
就會越拿越沉重.
同樣在生活裡面.
有很多壓迫,有很多重擔.
甚至乎有很多負面想法.
就好像這杯水.
你拿一會兒的話.
讓它輕輕佔據你的心靈.
可能沒有什麼問題.
你拿久一點的話.
你的心情就會開始變沉重.
開始有一點煩躁.
開始容易出錯.
影響身邊的人.
如果你拿超過一天的話.
你就會開始感到無力.
甚至乎你的人生開始會有部分癱瘓了.
動不了.
無法正常發揮它應有的功能.
沒有辦法做一些令自己有益的事情.
所以,弟兄姊妹.
我們需要在適當的時候.
學會放下那些沉重的壓力.
擔子,甚至是負面想法.
不要讓它窒礙了我們.
你也可以完成上帝託付給我們.
祭司的工作.
就是竭力.
將我們自己和別人帶到耶穌面前.
這位天上大祭司的面前.
使我們和他們一樣.
都可以在黑暗中得到拯救.
進入光明來和上帝和好.
事實上.
信徒皆祭司也是浸信會一直所強調的.

$^{721}$提醒我們不單單只擁有祭司的身份.
更重要的是我們要竭力去持守.
這個尊貴的身份和職事.
就好像經文裡.
耶穌基督向我們發出的應許.
你只管坦然無懼.
來到施恩助前.
你可以隨時向上帝支取力量.
來抵擋罪惡的試探和引誘.
如果當你真的不幸犯罪跌倒的時候.
你也可以隨時來到耶穌面前.
向祂認罪悔改.
祂已經為你完成了挽回祭.
你只要來到祂面前誠心悔改.
必定可以得到祂的寬恕和接納.
你也可以隨時靠著上帝恩典的幫助.
來持守真理.
完成你作為祭司的身份和使命.
直至我們真的走完上帝給我們的人生路.
就讓我們同心圍著祈禱.
天父幫助我們認清.
我們每一個在新約下.
都是你所揀選的君尊的祭司.
幫助我們能夠每一天都過一個.
出黑暗入光明的生活.
確實我們會有軟弱.
甚至我們會把這些軟弱收起來.
不敢跟別人說.
甚至連神也不敢跟你說.
但主也求你幫助我們.
讓我們真的像耶穌基督那樣.
我們可以隨時來到你的面前赤露闖開.
哪怕是我們覺得最幽暗,最軟弱,最羞愧的事情.
其實耶穌說你也很希望.
我們一一向你傾心圖義.
因為你體恤我們的軟弱.
你也曾經歷過人世間種種的試探.
並且你勝過這一切.
你有能力幫助我們面對這些的幽暗.
並且可以勝過.

$^{761}$來過一個土神喜悅的生活.
所以求上帝你幫助我們.
每當試探臨到的時候.
讓我們透過禱告來擺上.
也平時祈求上帝讓我們有足夠的力量和恩典.
讓我們像耶穌那樣勝過試探.
走一條蒙福,土地喜悅的道路.
上帝幫助我們也讓我們可以持守祭司的職份.
就是竭力使人與上帝和好.
來成為別人和上帝的橋樑.
來帶領人來到耶穌的面前.
認識你,跟隨你.
經歷上帝你的恩典和保護.
為此我們同心獻上祈禱.
奉基督耶穌你得勝的名字.
你也可以祈求.
阿們.
\newpage



\section{約翰福音 1:1-18}
\label{sec:XlsH8Zitzyg}
\textbf{2022年3月27日崇拜直播｜梁家麟牧師｜施浸約翰的見證(一)｜約翰福音1\_1-18}
\newline
\newline
連結: \href{https://youtube.com/watch?v=XlsH8Zitzyg}{\texttt{https://youtube.com/watch?v=XlsH8Zitzyg}} ~~~~ 語音日期: 2022-03-28
\newline
\newline
\hyperref[sec:bMqVNP_dN2s]{\small{< < < PREV SERMON < < <}}
~
\hyperref[sec:index_chronic]{\small{[返順時目]}}
~
\hyperref[sec:Or14nIn97R8]{\small{> > > NEXT SERMON > > >}}
\newline
\newline
約翰福音 1:1-18
\newline
\begin{longtable}{cl}
\hline
\hline
章節 & 經文 (和合本修訂版)\\
\hline
1:1 & \begin{tabularx}{0.7\textwidth}{X} 太初有道，道與神同在，道就是神。 \end{tabularx} \\ \\ \relax
1:2 & \begin{tabularx}{0.7\textwidth}{X} 這道太初與神同在。 \end{tabularx} \\ \\ \relax
1:3 & \begin{tabularx}{0.7\textwidth}{X} 萬物都是藉著他造的，沒有一樣不是藉著他造的。凡被造的， \end{tabularx} \\ \\ \relax
1:4 & \begin{tabularx}{0.7\textwidth}{X} 在他裡面有生命，這生命就是人的光。 \end{tabularx} \\ \\ \relax
1:5 & \begin{tabularx}{0.7\textwidth}{X} 光照在黑暗裡，黑暗卻沒有勝過光。 \end{tabularx} \\ \\ \relax
1:6 & \begin{tabularx}{0.7\textwidth}{X} 有一個人，是從神那裡差來的，名叫約翰。 \end{tabularx} \\ \\ \relax
1:7 & \begin{tabularx}{0.7\textwidth}{X} 這人來是為了作見證，是為那光作見證，要使眾人藉著他而信。 \end{tabularx} \\ \\ \relax
1:8 & \begin{tabularx}{0.7\textwidth}{X} 他不是那光，而是要為那光作見證。 \end{tabularx} \\ \\ \relax
1:9 & \begin{tabularx}{0.7\textwidth}{X} 那光是真光，來到世上，照亮所有的人。 \end{tabularx} \\ \\ \relax
1:10 & \begin{tabularx}{0.7\textwidth}{X} 他在世界，世界是藉著他造的，世界卻不認識他。 \end{tabularx} \\ \\ \relax
1:11 & \begin{tabularx}{0.7\textwidth}{X} 他來到自己的地方，自己的人並不接納他。 \end{tabularx} \\ \\ \relax
1:12 & \begin{tabularx}{0.7\textwidth}{X} 凡接納他的，就是信他名的人，他就賜他們權柄作神的兒女。 \end{tabularx} \\ \\ \relax
1:13 & \begin{tabularx}{0.7\textwidth}{X} 這些人不是從血生的，不是從情慾生的，也不是從人的意願生的，而是從神生的。 \end{tabularx} \\ \\ \relax
1:14 & \begin{tabularx}{0.7\textwidth}{X} 道成了肉身，住在我們中間，充充滿滿地有恩典有真理，我們也見過他的榮光，正是父獨一兒子的榮光。 \end{tabularx} \\ \\ \relax
1:15 & \begin{tabularx}{0.7\textwidth}{X} 約翰為他作見證，喊著說：「這就是我曾說：『那在我以後來的先於我，因為在我以前，他已經存在。』」 \end{tabularx} \\ \\ \relax
1:16 & \begin{tabularx}{0.7\textwidth}{X} 從他的豐富裡，我們都領受了恩典，而且恩上加恩。 \end{tabularx} \\ \\ \relax
1:17 & \begin{tabularx}{0.7\textwidth}{X} 律法是藉著摩西頒佈的；恩典和真理卻是由耶穌基督來的。 \end{tabularx} \\ \\ \relax
1:18 & \begin{tabularx}{0.7\textwidth}{X} 從來沒有人見過神，只有在父懷裡獨一的兒子將他表明出來。 \end{tabularx} \\ \\
[1ex]
\hline
\hline
\end{longtable}
$^{1}$讓我們繼續專注在耶穌基督的話語身上.
讓我們跟祂聆聽.
《楊樓福音》一章一至十八節.
經文裡面是這樣對我們說.
「太初有道,道與神同在,道就是神」.
「這道太初與神同在,萬物是藉著他做的」.
「萬被造的,沒有一樣不是藉著他做的」.
「生命在他裡頭,這生命就是人的光」.
「光照在黑暗裡,黑暗卻不接受光」.
「有一個人是從神那裡差來的,名叫約翰」.
「這人來,為要作見證,就是為光作見證」.
「叫眾人因他可以信」.
「他不是那光,乃是要為光作見證」.
「那光是真光,照亮一切生在世上的人」.
「他在世界,世界也是藉著他做的」.
「世界卻不認識他」.
「他到自己的地方來,自己的人都不接待他」.
「凡接待他的,就是信他名的人」.
「他就賜他們權柄,作神的兒女」.
「這等人不是從血氣生的」.
「不是從情慾生的」.
「也不是從仁義生的,乃是從神生的」.
「道成了肉身,住在我們中間」.
「充充滿滿的有恩典,有真理」.
「我們也見過他的榮光,正是父獨身子的榮光」.
「約翰為他作見證,堪著說」.
「這就是我曾說,那在我以後來的」.
「反成了在我以前的,因他本在我以前」.
「從他豐盛的恩典裡,我們都領受了」.
「而且因上加恩」.
「律法本是藉著摩西傳的」.
「恩典和真理都是由耶穌基督來的」.
「從來沒有人看見神」.
「只有在父懷裡的獨生子,將他顯明出來」.
今天繼續有梁家倫牧師在我們當中.
分享約翰福音的第二講.
講題是「施洗約翰的見證」.
恭請梁家倫牧師.
梁家倫牧師:香港現在是一個很困難的環境.
容許我在這裡先多謝.

$^{41}$每一個在醫療系統已經崩壞的警方之下.
仍然站穩每一個醫療位置的醫護人員.
多謝你們.
我也多謝那些在老人院或其他殘疾院舍.
一個這麼高危的地區.
一個這麼受害的地區.
繼續守護,繼續打拼的每一個工作人員.
多謝你們.
我們繼續看約翰福音.
我們從今天開始.
我們會有一連三次的時間.
我們會去講約翰福音裡面所記載的施洗約翰.
我們會看見約翰福音有三段經文講到施洗約翰.
不過沒有什麼內容是關乎施洗約翰本人.
如果我們想從約翰福音裡面看到施洗約翰的生平故事.
不是看到多少.
因為約翰福音的作者很清楚將約翰定義為一個見證者.
他只有一個功能,一個使命.
就是為耶穌基督做見證.
所以我們將這三次的題目都叫做施洗約翰的見證.
我們今天講的是第一講.
仍然看回約翰福音的序言.
一章一到十八節裡面的施洗約翰的技術.
我們想講一講.
約翰福音一章一到十八節的序言.
既是一個有關耶穌基督的生平故事.
也是約翰所做的一個關於道成肉身的神學論述.
以及一個有關耶穌基督的身份和使命的信仰見證.
生平故事,神學論述,信仰見證.
三者合而為一.
當講到信仰見證的時候.
斯諾約翰坦承他不是耶穌的第一個見證者.
在耶穌的公開事奉生涯那裡.
第一個見證耶穌的特殊身份和使命的人是誰?.
是施洗約翰.
在十八節的序言裡面.
約翰有兩次提到施洗約翰.
我知道我們教會應該叫施浸約翰.
不過因為我不懂得很容易轉台.
所以請大家原諒.

$^{81}$第一個的技術是在一章六到八節.
有一個人是從上帝那裡差來的名叫約翰.
這人來為要作見證.
就是為光作見證.
叫眾人因他可以信.
他不是那光.
那是要為光作見證.
然後一章十五節第二段.
約翰為他作見證喊著說.
就算我曾說那在我以後來的.
反成要在我以前的.
因他本來在我以前.
施洗約翰和使徒約翰同名.
不過兩個人應該沒有什麼關係.
施洗約翰為什麼要這麼稱仙.
施洗約翰要特別凸顯他的地位呢?.
在講到道成肉身的時候.
要兩次插段提到他呢?.
施洗約翰有什麼重要呢?.
我們看罪言的這個段落.
我們知道約翰的開局是.
他首先宣告道成肉身的事實.
然後指出這個成了肉身的道.
就是上帝藉以創造萬有的道.
並且就是上帝自己.
但是約翰其實不僅是想告訴我們.
有道成肉身這件事.
而是要指出那個死了不久.
很多在耶路撒冷的猶太人.
仍然印象猶新的那撒勒耶穌.
他就是成了肉身的道.
他就是上帝.
約翰揭示道成肉身的道理.
目的不是說明上帝有能力變成人.
上帝願意變成人.
上帝曾經變成人.
God has personified.
原來上帝曾經變成人.
你們知不知道?.
約翰要指出.

$^{121}$耶穌基督就是上帝變成人的面貌.
The being of God personified.
道成肉身的論述.
端在於指出.
耶穌就是成了肉身的道.
端在於確立耶穌的神聖.
正如我多次說過.
深一論的教義.
目的不是幫助人去窺探.
或者想像上帝本體的奧秘.
他只有十萬零一個功能.
唯一一個功能.
就是宣告耶穌基督是上帝.
同樣的.
宣告道成肉身.
目的不是在於說明道曾經成了肉身.
而是指出.
那個在人間活過擁有肉身的耶穌.
就是道.
就是上帝.
這個不是抽象的道理.
是已經發生了的事實.
Actually happened.
道成了肉身.
住在我們中間.
充充滿滿的有恩典有真理.
我們有見過他的榮光.
正是父獨生子的榮光.
道成肉身不僅僅是一個可能.
一個設想.
而是已經成為事實.
這個已經成為肉身的道.
曾經住在我們中間.
就是我們見過的他.
我們見過他的榮光.
他.
耶穌基督.
我們見過耶穌基督.
見過從他身上透視上帝的榮光.
我們可以論證一個道理.

$^{161}$但我們不會論證一個事實存在.
正如你可能不知道我在說什麼.
以前在宗教哲學上.
很多批評本體論證.
Ontological arguments(道理論證)的人都會指出.
我們不可以將一件事定義到它存在.
We cannot define God.
We cannot define something into existence(我們無法定義神,我們無法定義某些東西進入存在).
存在不是論證出來的.
你不會定義上帝.
他本身是Nothing greater can be conceived(沒有什麼能夠被想像的).
是最偉大的存有.
於是就等於他存在.
你怎麼描述他都沒有關係.
存在不是論證出來的.
存在是經驗到的.
親眼看到,親耳聽到,親手摸過.
而經驗的就成為見證人.
所以漢福音的序言.
講述道成肉身.
不是一個道理的論證.
而是一個事實的宣告.
耶穌基督製成了肉身的道.
這個宣告需要有目擊者去支持.
住在我們中間.
我們也見過他的榮光.
約翰和其他很多使徒信徒.
都是這個事實的見證人.
而史實約翰更加是最重要的見證人.
平情而論.
在觀念上接納上帝變成人的可能性.
已經挑戰了我們理解的極限.
上帝是自由的靈.
怎麼可以變成有限的塵軀呢.
上帝是尊貴神聖的.
怎麼會輸尊降貴變成凡人呢.
如果這樣都難.
那進一步.
要你接納一個出身不是好東西的拿撒勒.
無階型美容.

$^{201}$且被釘死在十字架上的耶穌基督是上帝.
這個真是匪夷所思.
荒天下之大謬.
不過約翰福音正是要揭示這個事實.
上一講的起手.
我們已經引述了這一次.
這段說話.
史上約翰在20章的結尾說.
但記者些時.
要叫你們信耶穌是基督是上帝的兒子.
並且叫你們信了他.
就可以因他的名得生命.
史上約翰不僅要證明耶穌是救主.
是救若應許要來的彌賽亞.
更要證明的是.
耶穌是上帝的兒子.
耶穌是上帝.
這是一個非常難做的見證.
極難取信於人.
不過這是作者承擔的第一任務.
是所有仕途付上一切代價都要傳講的信息.
是歷代教會包括你和我在內.
要向這個世界所做的見證.
耶穌是上帝.
耶穌是盛了肉身的道.
我們都是為耶穌做見證.
哥林多前書一章六節.
提摩太后書一章八節.
在作者撰寫《楊福音》的時候.
耶穌已經死了.
並且復活了升天了.
他和其他仕途是親眼目睹.
耶穌復活和升天的事實.
他們對耶穌是上帝這個真理.
比較容易接納.
你推前一點.
在耶穌生前和耶穌同時代的人.
幾乎都無法接受耶穌是上帝這個說法.
我明明白白眼前見到的.
一個和我差不多樣子的人.

$^{241}$生育的一個人.
並且我還知道他在很多人間的條件.
例如他出生和外貌和社會地位都比我差.
我怎會信他是創天造地的上帝呢.
他在我面前行過神蹟.
就算他行過神蹟.
你最多只能夠承認他是一個二能者.
一個行二能者.
人間不缺乏行二能的人.
舊約有很多先知都行得.
其他宗教都有.
又或者再退一步.
我就算相信他有特殊的地位.
他有他屬靈身份被人尊貴.
但我都不會把他當成是上帝去拜的.
我認識不少弟兄姊妹.
很熱切追求神蹟奇事.
常常參加很多神蹟大會.
甚至主動飛去新加坡,泰國等其他地方參加.
他們見到那些聖經出油,滿天金粉.
很多病患者宣告得到醫治.
我就很興奮.
大嘆法力無邊.
回到香港他們很興奮地和我分享他們所見所聞.
當然順便傳醫治的福音.
邀請我加入他們.
我總是聽了就算.
我沒有和他們辯論.
他們問我.
這些神醫如果不是神人.
又怎會有上帝的特殊能力.
行這麼大的神蹟呢?.
你解釋吧.
我的回應很簡單.
大無其不由.
我無法解說所有異常的現象.
你們所說的神蹟.
我不知道.
其中可能有裝假或誇大成份.
不過我沒有證據.

$^{281}$所以我不會這樣論斷.
但是.
我同樣無法排除.
其他宗教的大師.
例如從前香港人最熟知的泰國的白龍王.
或是密宗的法師.
他們都可能行神蹟.
如果完全是假的.
為何會有這麼多追隨者呢?.
我無法論斷一定是假的.
就算是真的.
我解釋不了.
但我不會因此.
當他是神明去拜他.
我不會.
耶穌能夠行異能.
我們要相信他是上帝.
兩者沒有必然關連.
A能夠行異能.
不可以推證A是上帝.
所以為何.
福音書裡.
耶穌每次在人面前.
稍微透露他的神聖身份.
將他和天父等量齊觀.
所有人的反應是怎樣?.
極其憎惡.
極其厭惡.
然後認定他在說一些.
褻瀆或是聖經最喜歡說的.
潛望的話.
耶穌在說潛望的話.
例如在約翰福音第十章.
當耶穌.
做了一個很大膽的宣告.
我與父緣為一.
嘩!的時候.
猶太人怎樣做呢?.
拿石頭扔他.
並且據事後解釋.

$^{321}$我們不是為善事.
因為耶穌行善事.
用石頭扔你.
是為你說了一些潛望的話.
又因為你是一個人.
但你竟然將自己.
當成上帝.
十章三十三節.
不僅僅同代人.
無法接納耶穌的神聖身份.
連跟隨他的使徒.
也不見得明白.
或接受耶穌是上帝的道理.
我們看約翰福音十四章.
八到九節就知道了.
我講了一大堆的東西.
我們看到.
施西約翰的.
見證的特殊價值.
他在耶穌還未出道之前.
在眾人尚未顯出.
他的智慧和能力.
就已經見證.
耶穌是舊約所預告要來的那一位.
施西約翰是人間.
最早確認耶穌是永生上帝的兒子的人.
西門彼得的認訊.
八到十六章十三到二十節.
最早要早兩年.
約翰福音的序言.
有兩次提到施西約翰.
必須要講.
這兩次的提到.
都是很噁心.
很奇怪.
第一.
兩個寄述都是.
在神學論述.
例如在第五節.
光照在黑暗裡.

$^{361}$黑暗卻不接受光.
第六節就引入約翰.
有一個人要為光做見證.
在講了三節有關約翰的技術之後.
然後第九節.
又回到對光的神學論述.
光是真光照亮一切生在世上的人.
然後一直到第十五節.
第二次又提到施西約翰.
如期來給人一個.
無端端生硬插入的感覺.
約翰為他作見證.
喊著說.
第一.
兩個提到約翰都是很奇怪.
強行插入.
第二.
序言對施西約翰的技術.
是無頭無尾的.
比較詳盡的講法.
是六到八節.
有一個人是從上帝撈來的.
約翰是來為要作見證.
就是為光作見證.
因他可以信.
他不時那光也是為光作見證.
單看這三節經文.
除了名字我知道他叫約翰之外.
我對約翰本人是一無所知的.
因為沒有片言隻語介紹過他.
我只知道他的來源.
從上帝差來的.
知道他的使命.
為光作見證.
他不是那光.
但來源和使命.
都是外在的.
沒有牽涉到約翰的個人資料.
是不是.
我告訴你.

$^{401}$有一個人是烏克蘭來的.
他要.
向香港求救.
你知不知道那個人是甚麼樣子.
我有沒有說過那個人是甚麼樣子.
沒有說過他任何個人的事.
是不是.
約翰的相貌.
他的家庭背景.
約翰的生平事蹟.
不生成就.
他死了還是在生.
這三節經文都沒有提到.
所以.
除非我們假設.
約翰福音的第一受眾.
是已經認識斯蟻約翰.
對他的生平事蹟.
耳熟能詳.
根本不用作者做任何介紹.
否則你無法解釋到.
約翰為何會做這個無厘頭的插段.
就像我在Facebook.
在IG.
分享.
我昨天行山的時候.
碰到周潤發.
我不用在說這句話之後.
馬上補充.
周潤發是無線第三期.
藝員訓練班的畢業生.
曾經拿過三屆香港電影金像獎的.
最佳男主角.
我不用說.
我可以假設絕大多數的香港人.
包括你們在內.
可能有人不認識.
但我猜多數人認識.
斯蟻約翰.
應該是.

$^{441}$同時代猶太人中間.
家傳戶曉的人物.
這個人竟然要勞駕.
分封王希律安提柏.
親自處死他.
你想想他多重要.
多有知名度.
約翰福音的作者.
找他作為耶穌基督的第一個見證人.
真是十足份量.
請用新旗牌生抽.
連周潤發都說好.
請相信耶穌是上帝的獨生子.
連斯蟻約翰都做見證.
無論如何.
但就約翰福音的序言提供的資訊.
我們不知道斯蟻約翰是誰.
其實知道沒有什麼內容.
第三.
證珠前面所說的.
是約翰福音的一句話.
證珠前面所說.
作者做了一個很噁心.
又沒有頭沒有尾的論述.
並且沒有告訴我們斯蟻約翰是誰的種種事實.
我們有理由因此相信.
作者在這裡根本不是要介紹斯蟻約翰.
而是僅僅要指出.
道是耶穌基督.
這個真理早已經有人做見證.
就是第一個見證人.
我們不需要知道斯蟻約翰是誰.
根本不是重點所在.
我們只要知道耶穌基督是成了肉身的道.
這個事實已經有人站出來做見證.
見證者就是約翰.
講個故事吧.
我一次坐飛機從美國回香港.
我坐的是經濟位.
靠窗那三個座位的走道位.

$^{481}$坐我旁邊的乘客.
是非常麻煩的.
都未起飛.
就已經問乘務員.
拿了兩次飲品.
然後又吵吵鬧鬧批評乘務員服務態度不好.
總之就是喧巴嘈吵.
搞到乘務員很困擾.
我同排靠窗.
坐著一個西人.
他忍不住就跟乘務員說.
你不用理他.
這個人是想殺賴.
想博.
換一個升級去商務位.
我亦都即刻就跟乘務員說.
我去做見證.
你們的服務態度是沒有問題的.
一切都是這個人無理取鬧.
那個麻煩乘客.
不甘心.
還想繼續大吵大鬧.
不過前後左右排的人.
都受不了.
一起大吵大鬧.
於是他就不甘心.
不情願都要安靜下來.
飛機起飛之後.
他大概不好意思跟我們一起坐.
於是就.
假裝去廁所.
然後就找到最後排.
有一個位置.
於是他就自己調位調到他那裡.
我去廁所的時候見到他.
他已經.
安靜地睡覺.
被牽連的乘務員.
很感激我旁邊的西人.
和我仗義執言.

$^{521}$不久之後.
乘務員的領班和副機長.
即是穿著機師制服.
他們更加用文字記述了整件事的始末.
然後要求西人和我.
先看了然後簽名.
然後我們兩個人.
都各自留下我們的名片.
包括我們的聯絡電話.
以防那個麻煩乘客.
事後投訴惡人先告狀.
我們就可以為.
乘務員和航空公司的.
清白去做見證.
我告訴你.
那個西人和我在.
底部要下機的時候.
那個領班.
親自來.
每人送了一支紅酒給我們.
在這次的事件裡.
那個西人和我.
都是見證人.
我們留下姓名和聯絡方法.
不過沒有人查問.
我們的個人故事.
我們是誰.
這個根本不重要.
重要的是我們目睹了整件事的經過.
因此我們可以為他做見證.
西約翰是見證者.
他不是最嚴重的主角.
西約翰為什麼有資格成為耶穌基督的第一個見證人呢.
聖經特別強調.
他是從上帝那裡.
差來的一張綠折.
因此他那個見證人的地位.
是由上帝安排的.
耶穌的身份的秘密.
大概也是上帝透露給他聽的吧.

$^{561}$在一張33到34折.
約翰曾經這樣說.
我先前不認識他.
只是那差我來的.
用水施洗的.
對我說.
你看見聖靈降下來.
住在誰的身上.
誰就是用聖靈施洗的.
我看見了.
就證明這是上帝的兒子.
這段說話說明.
約翰對耶穌的認識是有一個過程的.
不是一開始就確定耶穌的神聖.
耶穌在接受洗禮的時候.
那個榮耀的景象.
天開了.
這就是最能夠兼顧約翰對耶穌的信心的證據.
我們最後說一點.
約翰個人的故事.
有關約翰個人的背景.
我們要跟其他的福音書才能知道.
約翰福音將約翰的地位抬到很高.
但有關他的技術其實很少.
有一點特別值得注意.
由於約翰是一個很普遍的名字.
很普通.
聖經有很多個約翰.
所以我們通常叫他施洗約翰.
施浸約翰.
這也是馬可和馬太兩本福音書的叫法.
以致讓他可以跟其他約翰做區別.
但約翰福音從來沒有叫他施洗約翰.
只叫他做約翰.
並且約翰福音也沒有詳細記述.
他在約旦河的公開服事.
如何教訓人.
如何為人施洗.
特別為耶穌施洗的事蹟.
三卷福音都有記述.

$^{601}$但約翰福音沒有.
只是在約翰福音一章25到34節.
記述約翰和人的一段對話裡.
間接提到他曾經為耶穌施洗.
並且做過一些解說.
但不是記載那個故事.
而是記載約翰的一段說話.
我們知道約翰是耶穌的親戚.
他的父母是撒迦利亞和以利沙伯.
父親是祭司.
所以應該是屬於利美茲派.
他的家鄉在.
安基林.
我們今天如果參觀聖地團.
多數會去這個地方.
這裡有一個施洗約翰的教堂.
屬於方濟國會.
約翰和耶穌有親戚關係.
所以他從小應該認識耶穌.
不過這並不是表示.
他在童年的時候就知道耶穌的身份.
他最多從父母口中.
知道馬尼亞所誕生的兒子.
是不尋常的人.
有不尋常的誕生經歷.
約翰在《約翰福音》裡面曾經兩次強調.
我先前不認識他.
一章31節33節.
指的是他過去還不知道耶穌是聖經所要求的那一位.
撒迦利亞夫婦晚年得子.
所以應該是在約翰童年的時候.
父母就相繼去世.
在父母死後.
約翰就加入了一個黑虎的修道團體.
類似我們很熟悉的昆蘭社團.
在約旦河附近曠野居住.
耶穌後來也有加入這個群體.
並且接受約翰的洗禮.
不過他在約翰遇害之後.
耶穌就離開約旦河邊.

$^{641}$轉去加利利地區.
開始他三年半的全道生涯.
約翰福音對斯一約翰最重要的介紹是.
他是從上帝那裡差來的.
只能來為要作見證.
就是為光作見證.
叫他眾人因他可以信.
聖經作者清楚告訴我們.
斯一約翰為耶穌作見證.
這是他一生最重要的使命.
他的全活就是為耶穌作見證.
跟我們一般人.
我們活了一把年紀.
然後我們信耶穌.
然後我們領受異象.
領受一個召命.
這個領受的異象和召命是後加的.
並且不一定對現.
我年輕的時候.
我身邊很多人領受過很多光光烈烈的異象.
也做過很多光光烈烈的納智.
不過好像後來沒有對現.
沒有對現又有什麼影響呢.
最少暫時看不到.
上帝沒有擊殺他們.
但是斯一約翰不是.
斯一約翰一生是為一個召命而被佈局.
而被誕生的.
你記得在他出世之前.
天使對撒迦利亞的預告.
路加福音一章十三到十七節就可以知道.
所以約翰確實是從上帝差來的.
並且他是為一個使命而被差遣到世上.
他一生就只有一個使命.
他的使命是作為耶穌基督的先驅和見證人.
他所有要做的事情.
包括他的屍洗.
都是為了兌現這個使命.
在第一章三十三節他講.
「茲特拉差我來用水施洗的對我說」.

$^{681}$他很知道他是被差來做屍洗的.
他知道誰差他.
差他要做什麼.
耶穌的上帝獨生子的身份.
不是約翰從一開始就知道.
不過耶穌的獨特性是他知道的.
當他為耶穌施洗的時候.
他見到那個榮耀的景象.
他就心悅誠服地確認.
耶穌基督是上帝的獨生子.
因為這是上帝親自講的話.
「這是我的愛子」.
耶穌一心要做一個預備者的角色.
他認定他是預備他的道.
收藉他的道.
他在不知道那個要來的是誰之前.
他已經知道他是開路的那個.
他是一個開路先鋒.
預備人的心靈.
好迎見他們的主.
他為人施洗一個赦罪的洗禮.
但是約翰竟然承認.
這個洗禮是沒有什麼赦罪的能力.
真正能夠除掉世人罪孽的是耶穌基督.
所以這是他對耶穌基督的身份的第一個描述.
一章二十九節.
看下上帝的羔羊除去世人罪孽.
當除去世人罪孽的耶穌來了.
約翰的職事就再沒有價值.
這就是為什麼他日後宣認.
耶穌必然興旺.
他自己必然睡眠.
施洗約翰的見證者的身份.
讓我想起春秋時期和氏璧的故事.
大家很熟悉吧.
在楚國有一個叫做辯和的人.
有一次偶然在楚山中尋到一個還沒割削的薄玉.
他拿去獻給楚國的國君.
荔王.
荔王叫一個玉匠來去鑒別.

$^{721}$那個玉匠說.
只是碎料 只是普通石頭.
荔王認為辯和是個騙子.
就將辯和的左腳斬下來.
楚荔王死了之後.
武王做了國君.
辯和又拿這塊薄玉獻給武王.
武王又找玉匠鑒定.
玉匠說.
只是普通石頭.
武王同樣認為辯和是騙子.
然後將他的右腳也斬了.
武王又死了.
文王繼位辯和.
抱著那塊薄玉在楚山的山腳痛哭了幾天幾夜.
眼淚都哭乾了.
據說哭得出血.
文王聽到這件事.
就派人去問辯和.
他說天下間被人斬了兩隻腳的人很多.
為什麼你要哭得這麼淒涼.
辯和回答說.
我不是傷心.
我自己的腳被斬了.
我悲痛的是.
保玉竟然被看為是石頭.
忠誠的好人被看為是騙子.
這才是我最傷心的原因.
文王於是就叫玉匠.
認真加工雕琢這塊薄玉.
結果就發現.
這真是一個稀細的保玉.
於是就將它命名為和氏璧.
和氏之璧.
弟兄姊妹.
我們活在一個不信的世界.
很多人對信仰不感興趣.
冷嘲熱諷.
不過我們有幸得見耶穌基督的榮光.
被祂拯救過來.

$^{761}$我們也因此得到一個見證人的身份.
要向世人見證.
耶穌基督是上帝的獨生子.
耶穌基督是人類唯一的拯救和盼望.
我們知道見證者的身份不容易做.
不過從第一個見證人.
施洗藥案開始.
見證人也不是很受歡迎.
不是很容易完成使命.
求主幫助我們.
我們站穩見證人的身份.
努力承擔.
福音使命.
第一要緊的是.
我們要享受.
有幸成為耶穌基督的見證人.
這個身份.
我再說.
這段日子很艱難.
我昨天在荃灣坐巴士.
看到一對年老的夫婦.
男的大概有七十多八十歲.
搖搖欲動.
他竟然一個人.
一隻手拿兩袋八公斤的米.
兩隻手就拿了三十二公斤.
然後就上巴士.
由荃灣坐到葵涌村.
我和太太看到眼都乾了.
拿這麼多米.
吃多久呢.
不過我們想.
對老人家來說.
能夠有三十多公斤米.
大概兩個月的安全就有保證.
小老百姓.
只能夠在有限的資源裡.
有限的能力裡.
為自己去經營生活的保障.
想到這件事.

$^{801}$我和太太突然間有種很有趣的幸福感.
是呀 日子艱難.
但我們仍然能夠.
在一個相對安恕的環境裡.
我們為我們今天.
活在一個這樣的世代.
我們仍然向上帝獻上感恩.
等姐妹.
為我們的基督徒身份.
為我們繼續得到上帝的恩眷.
為我們繼續能夠有限度的.
去承擔祂的使命.
我們獻上感恩.
上帝賜福你們.
謝謝大家!.
\newpage



\section{希伯來書 7:1-28}
\label{sec:Or14nIn97R8}
\textbf{2022年4月3日崇拜直播｜陳明泉牧師｜超越的基督:按麥基洗德等次為祭司｜希伯來書7\_1-28}
\newline
\newline
連結: \href{https://youtube.com/watch?v=Or14nIn97R8}{\texttt{https://youtube.com/watch?v=Or14nIn97R8}} ~~~~ 語音日期: 2022-04-04
\newline
\newline
\hyperref[sec:XlsH8Zitzyg]{\small{< < < PREV SERMON < < <}}
~
\hyperref[sec:index_chronic]{\small{[返順時目]}}
~
\hyperref[sec:V98EonzJR3Q]{\small{> > > NEXT SERMON > > >}}
\newline
\newline
希伯來書 7:1-28
\newline
\begin{longtable}{cl}
\hline
\hline
章節 & 經文 (和合本修訂版)\\
\hline
7:1 & \begin{tabularx}{0.7\textwidth}{X} 這麥基洗德就是撒冷王，是至高神的祭司。他在亞伯拉罕打敗諸王回來的時候迎接他，並給他祝福。 \end{tabularx} \\ \\ \relax
7:2 & \begin{tabularx}{0.7\textwidth}{X} 亞伯拉罕也將自己所得來的一切，取十分之一給他。他頭一個名字翻譯出來是「公義的王」，他又名「撒冷王」，是和平王的意思。 \end{tabularx} \\ \\ \relax
7:3 & \begin{tabularx}{0.7\textwidth}{X} 他無父、無母、無族譜、無生之始、無命之終，是與神的兒子相似，他永遠作祭司。 \end{tabularx} \\ \\ \relax
7:4 & \begin{tabularx}{0.7\textwidth}{X} 你們想一想，這個人多麼偉大啊！連先祖亞伯拉罕都拿戰利品的十分之一給他。 \end{tabularx} \\ \\ \relax
7:5 & \begin{tabularx}{0.7\textwidth}{X} 那得祭司職分的利未子孫，奉命照例向百姓取十分之一，這百姓是自己的弟兄，雖是從亞伯拉罕親身生的，還是照例取十分之一。 \end{tabularx} \\ \\ \relax
7:6 & \begin{tabularx}{0.7\textwidth}{X} 惟獨麥基洗德那不與他們同族譜的，從亞伯拉罕收取了十分之一，並且給蒙應許的亞伯拉罕祝福。 \end{tabularx} \\ \\ \relax
7:7 & \begin{tabularx}{0.7\textwidth}{X} 向來位分大的給位分小的祝福，這是無可爭議的。 \end{tabularx} \\ \\ \relax
7:8 & \begin{tabularx}{0.7\textwidth}{X} 在這事上，一方面，收取十分之一的都是必死的人；另一方面，收取十分之一的卻是那位被證實是活著的。 \end{tabularx} \\ \\ \relax
7:9 & \begin{tabularx}{0.7\textwidth}{X} 我們可以說，那接受十分之一的利未也是藉著亞伯拉罕納了十分之一， \end{tabularx} \\ \\ \relax
7:10 & \begin{tabularx}{0.7\textwidth}{X} 因為麥基洗德迎接亞伯拉罕的時候，利未還在他先祖的身體裡面。 \end{tabularx} \\ \\ \relax
7:11 & \begin{tabularx}{0.7\textwidth}{X} 那麼，如果百姓藉著利未人的祭司職任能達到完全—因為百姓是在這職分下領受律法的—為甚麼還需要按照麥基洗德的體系另外興起一位祭司，而不按照亞倫的體系呢？ \end{tabularx} \\ \\ \relax
7:12 & \begin{tabularx}{0.7\textwidth}{X} 既然祭司的職分已更改，律法也需要更改。 \end{tabularx} \\ \\ \relax
7:13 & \begin{tabularx}{0.7\textwidth}{X} 因為這些話所指的人本屬別的支派，那支派裡從來沒有一人在祭壇前事奉的。 \end{tabularx} \\ \\ \relax
7:14 & \begin{tabularx}{0.7\textwidth}{X} 很明顯地，我們的主是從猶大出來的；但關於這支派，摩西並沒有提到祭司。 \end{tabularx} \\ \\ \relax
7:15 & \begin{tabularx}{0.7\textwidth}{X} 倘若有另一位像麥基洗德的祭司興起來，我的話就更顯而易見了。 \end{tabularx} \\ \\ \relax
7:16 & \begin{tabularx}{0.7\textwidth}{X} 他成為祭司，並不是照屬肉身的條例，而是照無窮生命的大能。 \end{tabularx} \\ \\ \relax
7:17 & \begin{tabularx}{0.7\textwidth}{X} 因為有給他作見證的說：「你是照著麥基洗德的體系永遠為祭司。」 \end{tabularx} \\ \\ \relax
7:18 & \begin{tabularx}{0.7\textwidth}{X} 一方面，先前的誡命因軟弱無能而廢掉了， \end{tabularx} \\ \\ \relax
7:19 & \begin{tabularx}{0.7\textwidth}{X} （律法本來就不能成就甚麼）；另一方面，一個更好的指望被引進來，靠這指望，我們就可以親近神。 \end{tabularx} \\ \\ \relax
7:20 & \begin{tabularx}{0.7\textwidth}{X} 再者，耶穌成為祭司，並不是沒有神的誓言；其他的祭司被指派時並沒有這種誓言， \end{tabularx} \\ \\ \relax
7:21 & \begin{tabularx}{0.7\textwidth}{X} 只有耶穌是起誓立的，因為那位立他的對他說：「主起了誓，絕不改變。你是永遠為祭司。」 \end{tabularx} \\ \\ \relax
7:22 & \begin{tabularx}{0.7\textwidth}{X} 既是起誓立的，耶穌也作了更美之約的中保。 \end{tabularx} \\ \\ \relax
7:23 & \begin{tabularx}{0.7\textwidth}{X} 一方面，那些成為祭司的數目本來多，是因為受死亡限制不能長久留住。 \end{tabularx} \\ \\ \relax
7:24 & \begin{tabularx}{0.7\textwidth}{X} 另一方面，這位既是永遠留住的，他具有不可更換的祭司職任。 \end{tabularx} \\ \\ \relax
7:25 & \begin{tabularx}{0.7\textwidth}{X} 所以，凡靠著他進到神面前的人，他都能拯救到底，因為他長遠活著為他們祈求。 \end{tabularx} \\ \\ \relax
7:26 & \begin{tabularx}{0.7\textwidth}{X} 這樣一位聖潔、無邪惡、無玷污、遠離罪人、高過諸天的大祭司，對我們是最合適的； \end{tabularx} \\ \\ \relax
7:27 & \begin{tabularx}{0.7\textwidth}{X} 他不像那些大祭司，每日必須先為自己的罪，後為百姓的罪獻祭，因為他只一次將自己獻上就把這事成全了。 \end{tabularx} \\ \\ \relax
7:28 & \begin{tabularx}{0.7\textwidth}{X} 律法所立的大祭司本是有弱點的人，但在律法以後，神以起誓的話立了兒子為大祭司，成為完全，直到永遠。 \end{tabularx} \\ \\
[1ex]
\hline
\hline
\end{longtable}
$^{1}$弟兄姊妹我們進入以聆聽聖道.
繼續敬拜我們的主.
我們三一的上帝.
我們繼續希伯來書.
來打開第七章的一至三節.
和第十一節至二十五節.
希伯來書第七章的第一節.
「這墨基西德就是撒冷王.
又是至高神的祭祀.
本是長遠為祭祀的.
他當亞伯拉罕撒拜諸王回來的時候.
就迎接他.
給他祝福.
亞伯拉罕也將自己所得來的.
取十分之一給他.
他頭一個名翻出來就是仁義王.
他又名撒冷王.
就是平安王的意思.
他無父無母無族譜.
無生之始無命之終.
乃是與神的兒子相似」.
你們想一想.
我們來到第十一節.
「從前百姓在利未人祭祀職任以下受律法.
倘若藉這職任能得完全.
又何用另外興起一位祭祀.
照墨基西德的等次.
不照亞倫的等次呢?」.
祭祀的職任既已更改.
律法也必須更改.
「因為這話所指的人.
本屬別的支派.
那支派裡從來沒有一人事後祭壇.
我們的主分明是從猶大出來的.
但這支派摩西並沒有提到祭祀.
倘若照墨基西德的樣式.
另外興起一位祭祀來.
我的話更是顯然易見的了.
他成為祭祀.
並不是照屬肉體的條例.

$^{41}$乃是照無窮之生命的大能.
因為有給他作見證的說.
你是照著墨基西德的等次永遠為祭祀.
先前的條例因軟弱無益.
所以廢掉了.
律法原來一無所成.
就引進更美的指望.
靠這指望我們便可以進到神面前.
再者.
耶穌為祭祀並不是不起誓立的.
至於那些祭祀.
原不是起誓立的.
只有耶穌是起誓立的.
因為那立他的對他說.
主起了誓 決不後悔.
你是永遠為祭祀.
既是起誓立的.
耶穌就作了更美之約的忠報.
那些成為祭祀的.
數目本來多.
是因為有死阻隔.
不能長久.
這位既是永遠常存的.
他祭祀的職任就長久不更換.
第25節.
凡靠著他進到神面前的人.
他就能拯救到底.
因為他是長遠活著.
替他們祈求.
今日的宣講題目就是超越的基督.
按默基使德等次為祭祀.
我們將時間交給陳詠全牧師.
弟兄姊妹平安.
今日我們繼續以希伯來書.
作為今天我們要領受的訊息.
在我們講這個訊息之前.
曾經有一個片段.
有一位弟兄姊妹問我.
牧師 其實在聖經裡面.
講很多關於憲制的經文.

$^{81}$但已經進入到一個現代文明社會.
其實這些關於憲制.
都是講舊約或舊時的封建習俗的形式.
其實去到我們今日的現代文明社會.
憲制這個字眼或這種概念.
原則上是不適用的.
所以其實我們不用講那麼多.
為何經常要講這些內容.
其實又好像有時空差距.
我對於弟兄姊妹這個提問.
我覺得很有意思.
我們今日進入到一個現代文明社會裡面.
很多東西的運作.
和舊時幾千年前的信仰形式已經很不同.
不過不是因為形式不同.
我們就抹殺了在聖經裡面關於憲制的教導.
今日我們看的這段希伯來書.
關於墨基使德的內容.
一方面應答了我們為何不用.
像前人或舊時的信仰實踐.
要拖羊或牛或捉鳥仔來電獻祭.
另一方面在整個神學觀念裡面.
他告訴我們今日我們不需要獻祭.
不是說上帝廢除了獻祭這個形式.
而是那個祭用了另外一個方式.
有人代我們獻祭.
今日希望從這段墨基使德的教導.
或者希伯來書的作者.
和我們用幾細緻的方式來展述.
究竟基督耶穌為我們做了甚麼.
我們從這個角度來理解這段經文.
老實說在希伯來書裡面.
關於墨基使德的講論.
不容易立即理解.
需要一些時間來消化.
希望今日我們由一個概覽.
似乎我們在當中裡面.
我們拿到一些亮光.
然後在我們的查經當中.
我們再深入地思想.

$^{121}$我相信上帝會繼續不斷地對我們講說話.
我們一起祈禱.
求上帝幫助我們.
讓我們同心祈禱.
阿伯的父求你在我們當中帶領我們.
讓我們能夠進到經文的世界.
明白你的心意.
讓我們能夠在今日所看的經文當中.
你親自對我們每一個講說話.
我們恭敬將我們每一個的心交給你.
求你將你的華語刻在我們的心板當中.
讓我們聽後不會忘記.
更加讓我們終生持續地實踐你的華語.
幫助我們銘記基督耶穌的教導.
幫助我們每一個能夠清楚明白.
你在這個世代裡面所作的功.
以及你的心意.
求你就恩待我們眾人.
幫助宣講講得清楚.
聆聽每一個聽得明白.
獻上祈禱奉靠基督耶穌得勝的聖名.
Amen.
在第七章裡面.
希伯來書就特別引進一個人物叫做默基西德.
其實默基西德這個人物.
和亞伯拉罕或者聖經裡面一些很偉大的領袖.
其實在聖經裡面所講默基西德這個名字.
這個人的篇幅不是很多.
大概來自兩處的經文.
第一處就是在《創世記》第十四章.
第二處就是在詩篇一八一十篇第四節.
希伯來書的作者就用這兩段的經文.
其實就某程度上用一個名不經傳.
以前可能看聖經都不為意的.
不過這兩段的經文來引申.
或者我們今天在神學上講的預表.
藉著默基西德.
來透過他的記錄下來的事蹟.
來預表基督耶穌是一個怎樣的神.
是一個怎樣的大祭司.

$^{161}$首先我們看看希伯來書第七章一至三節.
這默基西德就是撒冷王.
又是至高神的祭司.
本是長遠為祭司的.
他當亞伯拉罕擊敗諸王回來的時候.
就迎接他 給他祝福.
亞伯拉罕也將自己所得來的取十分之一給他.
他投一個名 翻出來就是仁義王.
又名撒冷王 就是平安王的意思.
他無父無母無族譜無生之始無命之終.
乃是與神的兒子相似.
在希伯來書的作者裡面.
他介紹默基西德.
默基西德的背景內容是怎樣的呢.
來自《創世記》第十四章.
剛才我講過了 十七節到二十節.
這段經文有幾節我讀給大家聽.
亞伯蘭擊敗基大老馬和他的同盟王回來的時候.
所多瑪王出來.
在沙媚谷迎接他.
沙媚谷就是王谷.
又有撒冷王默基西德帶著餅和酒出來迎接.
他是至高神的祭司.
他為亞伯蘭祝福.
說願天地的主 至高的神賜福予亞伯蘭.
至高的神把敵人交在你手裡是應當稱頌的.
亞伯蘭就把所得拿出來的十分之一給默基西德.
在《創世記》裡面就是這麼簡單的.
那個介紹 就是說亞伯拉罕.
那時候連名字都沒有改叫亞伯蘭.
打贏了仗.
然後就將戰勝.
回到所多瑪和撒冷王面前.
其實是他們兩個出來迎接亞伯蘭.
然後默基西德就為他獻祭.
拿著餅和酒出來迎接.
然後亞伯拉罕就將十分之一的戰勝品.
奉獻給默基西德.
整件事其實很簡單.
不過在希伯來書的演繹裡面.

$^{201}$他有幾樣東西很細緻地告訴我.
第一樣東西 默基西德這個名字.
默基西德這個名字是很可行的.
為什麼呢 其實他的意思是很深遠的.
默基的意思其實是譯音.
他就說到Maki Melaki.
這個希伯來文是說我的王的意思.
王其實他的意思就是一個管治者.
一個統治者.
或者一個類似領袖的身份.
在Sedic 即是洗德的譯音.
就是在說公義的意思.
默基西德其實是在說一個公義 仁義的王.
那個譯音就是默基西德.
譯義就是仁義王.
這個人出現就是一個這樣的身份.
另外一樣東西就是.
希伯來書的作者更加說到.
他是薩蘭這個地方的王.
這個王的名字叫仁義.
但是他管治的就是薩蘭這個地方.
薩蘭是什麼地方呢.
應該很多學者都在說.
是舊約裡面的士劍這個地方.
就是說在士劍有一個王.
叫做默基西德 仁義王.
他是來自這個地方.
那個時候叫做薩蘭.
薩蘭其實就是譯那個意思.
意思就是譯那個音.
就是平安的意思.
所以在這個人那裡.
他在描述他.
就是希伯來書的作者.
就用他的名字和他的根據地.
來描述他的身份.
是一個仁義和平安的王.
這個王除了是一個管治者之外.
他也是一個至高神的祭司.
如果按著時序來說.

$^{241}$在聖經裡面第一個出現至高神的祭司.
就是默基西德.
我們很多時候一想起大祭司或者祭司.
我們第一樣想起的就是亞倫.
在舊約裡面和摩西的哥哥.
不過在《創世記》裡面.
他的描述是說按著歷史.
第一個出現的祭司就是默基西德.
當這個默基西德.
這個祭司出現的時候.
他是慶祝.
為著亞伯蘭打聖像回來的.
他是代表著上帝.
所以在聖經裡面說到.
他和一般的那些祭司.
或者那些民間裡面信仰的那些不同.
這個人某程度上是代表上帝.
來祝福亞伯蘭的.
他帶著上帝的祝福.
用他的身份 用他的位份.
按著這個獻祭的舉動.
聖經的作者就很清楚記載著.
他出來迎接他.
到最後亞伯蘭就將他十分之一.
送給默基西德.
作為一個感恩的一種奉獻.
聖經裡面再繼續的.
加多一點點的描述.
他講到這個默基西德.
如果在《創世記》裡面.
順民夏里完全沒有提及過他.
誰生了他.
為什麼他會做了事件.
或者殺了這個地方的王.
全部都沒有提及.
在一些沒有提及的地方.
希伯來書的作者就做了一些文章.
他就說這個人是無父無母.
無族譜 無生之始 無命之終.
他講到這個默基西德.

$^{281}$他沒有背景.
他就用這個沒有背景的介紹.
來類比這個人沒有這些背景.
其實最終是來自神.
那個沒有的意思不是石頭爆出來.
很神奇的.
而是講到這個人好像來自神的一個使者.
一個神的 至高神的祭司一樣.
所以在希伯來書的作者裡面.
他很小心的.
他講到這個沒有背景.
沒有父母 沒有族譜.
無生之始 無命之終.
這個這樣的狀態.
或者這種在創世記裡面這樣的技術.
就好像基督耶穌一樣.
他是自有永有的主.
他也是這個世界歷史終結的主.
他用這個東西.
他沒有技術他的背景.
來類比基督耶穌.
是一個這樣的神.
是一個這樣的祭司.
所以在這個創世記裡面.
某程度上來說.
你覺得希伯來書的作者有一點矛盾.
因為沒有寫的東西.
他就當他是沒有.
就好像無限上綱.
不過那個重點.
其實默基西德只不過是借助這個人.
來描述基督耶穌.
是一個來自神的.
一個至高神的祭司.
他是一個仁義王.
是一個帶來和平的王.
他就是為著亞伯蘭.
為著眾百姓.
來獻上上帝的祝福.
同時間也代表著地上的人.

$^{321}$來將感恩 將一個祭.
來獻情給上帝.
所以在原則上.
在希伯來書的介紹裡面.
他是借助了創世記的那個技術.
講到默基西德是哪一個呢.
來引入一個大祭司的主題.
當我們這樣解的時候.
我們大概明白了.
就是說默基西德.
他和阿倫的祭司的體系有點不同.
他是歷史上第一個的祭司.
更加有一點不同的地方就是.
默基西德的身份和阿倫是不同的.
和阿倫的體系.
阿倫就是利未之派.
整個宗族就是祭司.
他整個支派.
終身就是這樣排著隊.
來做一個祭司的系統.
但默基西德和他們不同.
你不知道他從哪裡出來的.
不過他就代表神.
同時間他除了懷著.
至高神的祭司這個身份之外.
他也是一個王的身份.
他帶著王和祭司.
帶著一個很尊貴的身份.
來為亞伯蘭獻上祝福.
在這裡我們大概理解到.
默基西德是誰呢.
我們在這裡有一個輪廓.
不過難度就在接下來的解釋.
在七章六到十一節.
獨有默基西德不與他們同袍.
不與他們同袍是說阿倫的同袍.
倒收納亞伯拉罕的十分之一.
為那蒙應許的亞伯拉罕祝福.
從來為份大的給為份少的祝福.
這是博不倒的理.

$^{361}$在這裡收十分之一是必死的人.
但在那裡收十分之一的.
有為他作見證的說.
他是活的並且可說.
那收十分之一的利美.
也是藉著亞伯拉罕納了十分之一.
因為默基西德迎接亞伯拉罕的時候.
利美已經在他先祖的身上.
從前百姓在利美祭司職任以下受律罰.
倘若藉者職任得以完全.
又何用另外興起一位祭司.
照默基西德的等次.
不照阿倫的等次呢.
這是希伯來書第七章六至十一節的內容.
這裡說什麼呢.
其實好像很迂迴.
不是很直接說.
他首先說到.
他說默基西德和坊間裡面所有人.
認定的那個祭司系統.
阿倫的那個系統.
不同的.
因為默基西德.
是接納亞伯拉罕十分之一的奉獻.
但是阿倫的支派.
只是接受亞伯拉罕之後的子孫.
所以大家的級數是不同的.
說到級數不同.
為什麼作者要說這番話呢.
其實他要回應.
在詩篇110篇第四節的經文.
耶和華起了誓.
決不後悔說.
你是照著默基西德的等次.
永遠為祭司.
這一句的經文.
是詩篇大衛寫下去的.
無端端他真的看《創世記》.
他知道有一個默基西德作為永生神的祭司.
比阿倫更早.

$^{401}$而詩篇那裡沒有特別去解釋.
究竟默基西德的等次的等級是什麼.
接下來希伯來書的作者就告訴大家.
默基西德的等次高過阿倫的等次.
那個等級是不同的.
他現在就告訴你.
默基西德的等次是一回什麼事.
七章六至十一節.
某程度上就一路演繹告訴你.
默基西德的等次是高階的一種奉獻.
那這種奉獻是怎麼樣呢.
首先講過了.
他是為亞伯拉罕.
是為以色列人的始祖來獻祭的.
所以他收納了十分之一.
經文裡面繼續的來講.
他說位分高的就給位分少的祝福.
在希伯來書或者在聖經裡面講到.
亞伯拉罕已經在整個以色列民族裡面.
是一個始祖.
列祖來的.
最頂的那個應該是最地位崇高的.
不過希伯來書的作者說.
不是的.
比亞伯拉罕更加崇高.
除了神之外.
就是默基西德了.
為什麼呢.
因為位分高的才可以祝福位分低的.
而默基西德是藉著這個獻祭.
來祝福亞伯拉罕.
而亞伯拉罕又回應十分之一.
來奉獻給他.
所以這是一個更高的身份.
位分高的就祝福位分低的.
這是中國人的社會不難明白的.
我們都明白.
不過現在這樣講好像很複雜.
你想想新年的時候.
誰給你紅包.

$^{441}$你不會讓你兒子給你紅包.
你是父母給兒子紅包.
就是一份祝福的金錢.
即使我們成長了.
我們給錢父母了.
我們都不會說我封封利是給母親.
中國人嚴格上來講.
按照我們的尊卑是不會這樣說話的.
我們只是說給錢他買東西吃.
給錢他來一份心意.
不會說我祝福上一輩的.
沒有的.
不過今天這個世代.
人們好像混亂了.
不會看這些長幼的了.
不過在聖經裡面講到.
那個長幼有序.
位分高的就祝福位分低的.
他藉著這種表達去講到.
默基西德比亞伯拉罕更加地位高尚.
他是用一個祝福者的身份.
來呈現在聖經當中.
而默基西德的等次.
比亞倫更高的就是.
因為默基西德是祝福以色列人的列祖.
當時利未人在哪裡呢.
聖經說.
利美已經在他先祖的身中.
那時候沒有現在的醫學科技.
不知道怎麼說.
其實簡單來說.
利美都還沒有出生.
原文說還在腰那裡.
簡單來說.
當這個憲制舉動實踐的時候.
祝福這個先祖亞伯拉罕的時候.
亞倫利美之派都不知道在哪裡.
都還沒有出生.
都還沒有出現.
所以.

$^{481}$默基西德的等次的班數.
追溯到一個歷史的起始點.
默基西德就是一個這麼超班的.
一個祭司來祝福亞伯拉罕.
在七章的第四節裡面.
其實又再重複又重複多次.
你是照著默基西德的等次.
永遠為祭司.
這個默基西德的等次.
其實就是要類比基督耶穌.
耶穌作為一個祭司的身份.
祂是和默基西德的等級是一樣的.
是一個人間的制度上.
都還沒有的時候.
耶穌就好像默基西德那樣.
來祝福身邊的人.
成為永生神的祭司.
是利至神.
無始無終無父無母.
是在說祂是利至神.
祂也是歷史的終結.
祂是終結者.
也是起先者.
按照默基西德的等次.
來祝福身邊的人.
亞倫祭司的後代和默基西德.
有幾樣東西不同的.
除了剛才說的時間上.
或者是剛才說的列祖的憲制之外.
在聖經裡面更加說到.
亞倫這個祭司制度.
其實就是和律法一起來走的.
有律法.
整個律法頒佈了.
同時間頒佈了一個敬拜的制度.
律法和敬拜.
在亞倫之後的祭司.
是一起的並行來實踐出來的.
因著人的這個制度.
如果你認識舊約.

$^{521}$你會知道.
整個利未知派.
他整個宗族.
所有人輪流做祭司.
而大祭司就是.
祭籤到最後被選照出來.
在這個入場券下.
在利未知派下.
每個人都有機會.
不過只有這個利未知派.
才有可能.
大家就好像輪班一樣.
輪班仔.
但是什麼時候會更替呢?.
就是那個大祭司.
可能任期到.
或者他死的時候.
就轉一轉裝.
然後第二個人就輪值的.
來到他上了.
亞倫或者是.
接下來利未的祭司制度.
是用人的制度.
來輪替交換.
來實踐向上帝獻祭的任務.
更加重要的是.
他講到這些祭司.
他能夠勝任的.
為什麼呢?.
因為這些祭司.
對人生命的軟弱.
是很明白的.
因為他自己本身是一個軟弱的人.
他很體恤人的軟弱.
上一次Raymond講的.
一個祭司.
他因為明白他跟你.
對人的那種生命.
那種對罪的掙扎.
很理解.

$^{561}$所以他從人性的軟弱角度.
來代表人.
向神獻祭.
求上帝贖罪.
同時間也求神.
閱納一些卑微的人.
擺上的一切.
亞倫祭司的制度.
就代表著軟弱的人.
用軟弱的身軀來表達.
但是這個祭司的軟弱.
也代表著他有一個缺點.
就是祭司會死的.
他也是血肉之軀.
他有一天離開世界的時候.
那個制度在輪值.
但是當他死的時候.
他的守望就好像暫時停結了一樣.
以致這個制度.
不可以言言不絕地.
徹底地.
有人代表世上的人.
向上帝祈求.
更加重要的就是.
大祭司在亞倫的系統裡面.
極其量代表的.
只是軟弱的人.
他只是代表著一個血肉之軀.
向上帝獻上禮物.
所以在希伯來書五章一至四節裡面.
希伯來書的作者在五章已經說到了.
凡從人間挑選的大祭司.
是奉派替人辦理贖神的事.
為要獻上禮物和贖罪祭.
他能體恤.
那愚夢和失迷的人.
因為他自己也是被軟弱所困.
故此他理當為百姓和自己的獻祭贖罪.
這大祭司的尊榮沒有人自取.
為要蒙神所召像亞倫一樣.

$^{601}$在聖經裡面.
他講到這個身份.
為人獻祭是一個很尊貴的身份.
是上帝選召的.
亞倫整個的支派都是一個很尊貴的身份.
不過尊貴極.
他都是從人性的軟弱出發.
他都是用一個軟弱者的表達來獻祭.
不過默基使德.
或者默基使德的等次作為祭司就有不同.
七章三節剛才講過.
默基使德.
他是無父無母.
無族譜.
無生之始.
無命之終.
乃是與神的兒子相似.
他成為祭司並不是照著肉體的條例.
乃是照著無窮之生命的大能.
這是七章十六節.
他講到默基使德與亞倫不同的是什麼呢?.
默基使德的選召.
默基使德等次的祭司.
是講到是上帝親自appoint的.
是上帝選召他.
上帝叫這個人成為一個祭司.
他是按照上帝永活的.
不老的生命的大能.
來承擔這個職事.
所以基督耶穌演繹默基使德的等次身份.
代表著他與人間祭司不同.
耶穌基督是永遠存活的.
他有一個永生的生命.
與那些會死的祭司.
那個級數完全不同.
那個永活的祭司.
用永遠的祭.
來幫助每一個軟弱的人.
來將祭物獻給上帝.
同時這個永活的祭司.

$^{641}$因為是上帝親自選派的.
所以他也是代表著神.
一個永生的上帝.
將上帝的祝福賜予地上的每一個弟兄姊妹.
所以這麼迂迴,這麼複雜.
牧師講很多字眼.
有時候都不容易.
其實要講什麼呢?.
簡單來說.
希伯來書的作者說.
原來這個世界裡面.
除了亞倫的憲制.
我們常常都想到之外.
原來比這個本身已經很尊貴.
上帝所命定的.
亞倫的憲制的制度.
原來在這個制度之上.
還有一個更加超越的一種憲制.
而這種憲制就是按照默基使德的等次.
作為祭司.
而這個祭司是誰呢?.
這個祭司不是默基使德.
這個祭司是耶穌基督.
耶穌基督就是扮演著.
或者是在實踐一個大祭司的職份.
用一個永遠,永活的生命.
來代表我們.
來獻祭給上帝.
同時間也代表上帝.
將來自上帝的福氣.
藉著基督耶穌.
來祝福我們地上的每一個.
所以七章十八節到第七章二十三節.
經文是繼續這樣說.
先前的條例因為軟弱無益.
所以廢了了.
律法原來一無所成.
就引進更美的指望.
靠者指望我們便可以進到神面前.
再者,耶穌為祭司並不是不起誓立的.

$^{681}$至於那些祭司原不是起誓立的.
只有耶穌是起誓立的.
因為立他的對他說.
主起了誓,決不後悔.
你是永遠為祭司.
既是起誓立的.
耶穌就作了更美之約的鐘簿.
那些成為祭司的數目本來多.
是因為有死阻隔.
不能長久.
經文裡面很奇妙或巧妙地用了詩篇.
第四節第110篇裡面.
說永遠為祭司的字眼.
為什麼耶穌是永遠為祭司呢?.
第一,因為耶穌是永活的神.
祂不會被死亡隔絕祂.
更重要的是多次又多次提醒讀者.
耶穌是上帝起誓而立的.
不過他寫的方法.
有時希伯來書的作者都頗頑皮.
他用的是雙重否定句.
耶穌為祭司並不是不起誓立的.
雙重否定句,負負得正.
即是說耶穌是起誓而立的.
詩篇裡面就是在說這一點.
上帝用誓言,用指著自己起誓的方式.
上帝定義要立一個以默基使德為等次的祭司.
這個人就是耶穌基督.
耶穌是上帝確實地命定.
要為這個世代裡面建立一個新的憲制制度.
他不是要否定憲制.
不過他要藉著基督耶穌.
要將一個更超越的憲制呈現在這個世界當中.
而那個祭司就是上帝起誓而立的耶穌基督.
藉著他就成為這個新約的祭司.
當然在這裡面再引入一個字眼.
就是更美的約的中保.
下一講我會為大家更加講下.
什麼叫做更美的約的中保.
不過簡單來說.

$^{721}$你知道耶穌的身份成為祭司.
他不是單純代表軟弱地上的每一個人.
他更加是代表神.
兩邊來代表.
來藉著他這個中間人.
來將上帝的祝福臨到這個世界裡面.
也將人的讚美和敬拜.
藉著他達到天父的根前.
經文裡再繼續講.
他說耶穌不受死所限.
他不會死的.
或者這樣說.
他是一個永活的上帝.
他的死只是成就在十字架上的救贖的死亡.
代替我們死.
但是死後第三天復活.
他是一個永活的主.
他復活的生命.
他永活的權能成為了一個有力的.
一個新的憲制制度的一種生命表現.
在一個這樣的狀態之下.
他面對著讀者或者我們每一個弟兄姊妹.
舊的那些軟弱無力的.
舊有憲制制度.
是好的 是尊貴的.
不過有一個更好的放在你的根前.
你沒理由走回舊的路.
你要拿一個新的路.
走一條更新更活更有效.
更徹底的道路.
來將我們帶到上帝的根前.
如果你都覺得很抽象.
用一個很簡單的比喻.
大家現在手機用的是一般來說4G.
有些頂展可能會追得更貼.
用5G.
旺朕用6G.
不過那6G不是在說科技.
不過我們時代不停在進步.
我們新的東西取代舊的東西.

$^{761}$舊的問題.
新的已經完全應答了.
如果你忽然拿那個大哥大水壺出街.
你都會覺得很搞笑.
你當玩具一樣.
同時間你說很好用.
比現在那5G手機好.
那些人說不要搞了.
不要玩了.
人家覺得你惡搞.
甚至你拿大哥大水壺出來都用不到.
因為現在的網絡完全都不配套.
同樣道理.
我們今天的信仰.
為什麼不用憲制.
因為有一個更高的.
耶穌基督已經做了這個更超越的憲制.
代替了我們.
所以我們不需要好像舊約那些.
以色列人,猶太人那樣.
拿一隻羊去到聖殿某個地方.
某個處境來將.
祭僧來獻上給上帝.
去到今天.
因為基督耶穌已經成就了這個.
更超越的那種憲制.
祂直著.
祂自己的生命.
將我們帶到在上帝面前.
甚至乎那隻祭僧.
祂用最好的祭僧來獻上.
就是祂自己.
十字架上的基督耶穌.
就是一個最徹底的祭僧.
祂將祂自己來獻上.
成為了最大的挽回制.
將人與神本身.
罪惡的人和聖潔主的對立.
完全解決了.
所以我們不再需要像舊約.

$^{801}$或者以色列人那樣.
走回一條這樣的舊路.
用這種方式來表達.
和上帝之間的敬拜和仰慕.
新的路取代舊的路.
我們不要走回舊有的敬拜模式.
上帝已經用終末的.
用最完美的方式.
來呈現了人神之間.
有一個新的關係的建立.
為何希伯來書的作者.
會說這番話呢.
其實在《紡紋夏利》你會知道.
在第六章裡面說到.
有一些懊悔死行.
去那些猶太人.
做回守節旗的事情.
或者有一些信徒.
走著走著這個信仰的路.
本來是藉著基督的.
怎知道走著走著.
開始灰心了.
就去回舊有的那些.
移民舊養.
做回某一些律法的規條.
希望從這裡可以得到一些信仰的養分.
去到上帝的根錢.
得到上帝的眷顧.
那條是舊路.
希伯來書的作者說.
不要再走那條路了.
浪費力氣了.
有一條更好更新的路.
你們應該要選擇的.
不要再放棄了你們一直堅守對的路.
而走回錯的舊路.
反而繼續要堅持走一條新的路.
上帝必然會引進你.
帶你去到上帝的面前.
去到今代了.

$^{841}$去到我們今天身處的處境.
如果你知道.
有一些特別的宗派.
或者有一些教會.
和我們今天主流所說的有點不同.
和聖經的教導甚至有些違背.
有一次有這樣的經歷.
我和我太太出去逛街.
在商場那裡.
有一個不知名的教會.
就走去問我老婆.
問關於主餐的問題.
他原本說好像做問卷那樣.
了解一下做主餐的.
然後就問我太太.
然後他特別抓住一句.
他說哥林多前書第十一章那裡.
不按你手煮的杯.
那個不按理的道理是什麼意思呢?.
我老婆就回答.
上文夏利就說.
聖經的人混亂聖經的道.
酒醉的飢餓等等.
然後他說不是的.
你不懂的.
那個是師奶來的.
我老婆說.
你都不懂的.
他說那個不按理是不按逾越節的道理.
就是說我們今天基督徒.
或者普世教會.
新約教會裡面所守的主餐.
是錯的.
不按道理的.
因為不按著舊約的逾越節方式.
來守主餐.
所以就應該被保羅責備.
應該被新約所責備.
然後我老婆就.
嘩 怎麼會有這樣的事情呢?.

$^{881}$其實那段時間.
我進了廁所.
然後在廁所出來.
我看到有兩個女士和我老婆聊天.
我老婆看到朋友.
然後我老婆說.
你問他吧.
他又不說我是牧師.
他說你問他吧.
他可以回答你的.
他就問我剛才的問題.
不按理.
然後就解釋.
你不按逾越節的方式守主餐.
我就有點不客氣.
你又不是猶太人.
你守什麼逾越節呢?.
普世教會都不是這樣做的.
而是你們特別這樣做.
通常強調自己有些不同.
和整個信仰不同.
通常都是異端.
那個人不說話.
然後又好像回答不了我.
然後就找個大佬來.
再來跟我再辯論.
都是說同樣的東西.
不過只不過是口才好一點.
然後我說.
你們都不看希伯來書的.
你們只看《斷章取義》.
過那麼多前書上文.
你都不是說逾越節的問題.
真的沒說話.
然後還要我留電話.
我說不如你留個電話給我.
我告訴你真正的信仰是什麼.
多過聽你們說.
你們自己要回去舊約.
要守不同禮儀的道理.

$^{921}$不過我可以很坦白說.
今天很多類似教會.
他們想回去舊約一些禮儀.
如果你用舊約的鑄棚節那些.
你作為一個教導.
了解舊約是可以的.
不過如果你用那些節旗.
成為我們要守的節.
好像猶太人那樣守.
那個就完全違背了.
希伯來書作者.
他在說我們又要走回舊路.
要做回這些事才得救的謬誤.
我們不要再進入那些誤區.
今天如果有很多人質問我們.
為什麼不做回這些事的時候.
我們可以立即去說起.
因為基督耶穌.
用了一個更超越的憲制.
超越了舊約.
這個移民舊養.
這種的律法規範.
來實踐信仰.
基督耶穌用他的生命.
實踐了一個更加超越.
我們跟著基督耶穌的這條路.
來實踐我們今天的基督教信仰.
所以弟兄姊妹.
不要走回那條舊路.
不要覺得那個才是幫我們.
今天我們信的基督耶穌.
才是我們唯一應該要守的出路.
希伯來書作者就在說.
他好像有一個義憤.
不過他強調不要再回頭.
不要走回舊時這些不可行的道路.
他講到律法是一無所成的.
如果你見到歷史裡面.
如果是可以的.
耶穌是不需要來的.

$^{961}$他就講到舊約逐步逐步.
漸進的啟示就講到.
基督耶穌是一個終極的救贖.
直著他.
我們每一個人去到上帝跟前.
去到七章二十五節.
這一段經文某程度是一個結論的部分.
也都是對我們來說一個很窩心.
很直接的教導.
凡靠他進到神面前的人.
他都能拯救到底.
因著他是長久活著.
替我們祈求.
七章二十五節就講到.
基督耶穌是按著默基使德的等次.
永遠為祭司.
對於我們的意義是什麼呢?.
我們每一個人可以靠著他.
都可以去到上帝面前.
拯救到底是一個很有意思的字眼.
拯救到底那個字.
其實在新約是很稀有的.
在這裡和在《路加福音》出現過.
但是我們按照上文下例.
我們知道那個到底其實有三個層面.
我們值得去注意.
第一層面就是時間的層面.
基督耶穌的拯救.
耶穌基督的獻祭.
是在歷史裡面.
直至世界的末了.
祂的救贖不會因為世代的更替.
或者有任何的改變而停止的.
除非祂再回來的日子.
在祂還沒回來的日子.
世界的歷史還沒終結的時候.
祂的救贖.
在整個人類社會世世代代裡面.
永遠長存.
上帝的愛.

$^{1001}$上帝的救贖.
上帝為我們獻祭的那種祭.
繼續的不停的延展.
所以直至到整個世界終結.
都是繼續的維持下來的.
這個是宏觀的歷史角度.
那個拯救到底也是對我們個人生命的角度.
當上帝是這樣照管這個世界.
同樣的拯救到底也是拯救我們個人的生命.
人生裡面可能都會有些差池.
有時候我們可能會有反叛.
我們會有迷失的時間.
但是如果上帝要救我們.
祂的恩典不會因為我們放棄了.
祂就算了這個人沒得救了.
不理我們.
反而是基督耶穌持久的.
祂為我們每一個祂所選召的人.
來守望.
祂會救我們救到底.
所以我們心裡面有一份確信.
我們知道原來這個上帝不會離開我的.
極其量是我們離開上帝.
但上帝從來不會離開我們.
我們靠賴這份基督耶穌的那個獻祭的恩典.
祂會拯救我們拯救到底.
這是第二個在個人層面來看.
第三個到底的意思除了時間之外.
也就是說那個程度.
不會說有些東西是不能救的.
有些人是不夠的.
有些人是脫離了救贖的計劃以外.
沒有的.
那個救到底是說.
既然基督耶穌在十字架上犧牲.
祂就不是很吝嗇.
要救少少人.
或者某些人我不救的.
祂說到.
祂的心意是願意萬人得救.

$^{1041}$不願一人沉淪.
祂對整個世界呈現了基督耶穌的愛.
或者這個世界沒有一樣東西.
可以難阻我們走到在上帝面前.
或者只有人的罪和人自己.
沒有一個罪是上帝不可以赦免和拯救.
我們靠賴基督耶穌.
我們每一個都可以去到在祂的根前.
聖經裡面再繼續說.
耶穌怎樣去對我們.
除了那個救贖.
那個獻祭不是一個很靜態的造型.
這樣就算了.
祂說到這個大祭司.
是不斷圍著這個世界.
圍著每一個罪人來祈求.
那個祈求是什麼呢?.
我想起耶穌為彼得.
在說一個豪情壯語.
說怎樣都不離開他的時候.
耶穌說西門西門.
撒旦想要得獎你.
好像西墨寺一樣.
但我已經為你祈求.
叫你不至失去信心.
你回頭過來.
你要堅固你的弟兄.
路加福音22章31至32節.
耶穌為這個世界祈求.
這個人世間很多人都可能失迷.
或者誇海口.
我是很堅心地跟隨主.
不過當我們都不知道自己想怎樣的時候.
耶穌已經為我們祈求.
叫我們不要失去信心.
在人生跌跌盪盪的日子.
在迷失犯罪.
在遺憾充斥.
或者充滿了愧疚的生命當中.
爛融融.

$^{1081}$但是上帝叫我們不自動搖.
不會失去僅餘的信心.
我們就憑藉這些信心.
可以因著基督耶穌.
再一次回到上帝的根前.
去到私人寶座面前.
得到上帝的祝福.
如果你有機會看到一些書.
有一個國學大師叫林語堂.
如果大家有機會看到他的書.
他是學貫中西近代的一個.
很出名的作家.
他有另外一個別號叫優墨大師.
他寫的文是用了西方的那種優墨.
但是用中國人的處境來表達.
這個國學大師林語堂.
本身他是一個牧師的兒子.
他是信義代來的.
從小到大他想著要做牧師.
所以在他年輕讀書很厲害的時候.
在他身處的年代.
一九三幾年.
他考入了當時的精英的上海約翰.
聖約翰大學.
高材生這樣畢業.
他原本想著讀神學的.
在裡面被安納做一個牧師.
但是進去的時候.
人家覺得他很反叛.
經常問一些問題.
覺得好像挑戰.
而不是真誠去尋求信仰.
到最後他就放棄了這個信仰.
寧願按著他的學習.
找回中國人的文化.
他就走到西方社會.
林語堂在西方社會裡面很吃得開.
因為將中國文化呈現出來.
在他四十歲的時候.
他自己說.

$^{1121}$他是一個無神論者.
他覺得這個世界沒有永生.
他就否定了他自己是一個基督徒.
他採取了一種人民主義的態度.
放棄了自己.
這個從小到大所持守的信仰.
直到他年老的日子.
特別是他第二個女兒自殺之後.
他再尋索究竟真正的人生意義是什麼呢?.
在他的信仰之旅.
有一本書他講了.
他說.
他三十多年來.
他唯一的宗教就是他自己.
就是人民主義.
他覺得理性的指引.
可以不假外求.
這個世界不斷的進步.
人越來越自信.
但是這個世界沒有變得越來越好.
人越來越聰明.
但是缺少了越來越對天那份虔誠和謙恭.
他覺得.
他好像那三十年走迷了自己一樣.
直到他年老的時候.
他重新看回他二時學的聖經.
他看到耶穌基督的祈禱.
耶穌說接待他就好像接待一個小孩子一樣.
他知道耶穌在十字架上的祈禱.
「父說赦免他們.
因為他們所作的他們不知道」.
他再重新去看回.
舊時二時的那個信仰.
他才發覺.
基督教的信仰.
聖經裡面所啟示的真氣.
才是這個世界.
這個人類.
重新要走的一個出路.
所以他寫的那本書.

$^{1161}$信仰之旅.
從異教徒到基督徒.
這個某程度是他生命裡面的一個見證.
我要說些什麼呢?.
不知道弟兄姊妹你今天.
你自己會不會生命裡面走一走這個信仰的路.
走著走著覺得有些事情都不可行.
然後就可能尋索一些其他的道路.
尋求一些新意.
尋求一些刺激.
想著那些可以為你信仰帶來一些養分.
誰知越尋求越穩的時候.
可能我已經久而久之.
不知不覺地.
我就走了一條偏差的路.
不過今日經文提醒我們.
其實我們一直以基督耶穌為中心的信仰.
是對的路.
堅守著它.
不要放棄.
同時間.
倘若我們生命真的迷失了.
上帝是一個拯救我們到底的上帝.
祂不會放棄你.
唯有直著祂.
我們才可以去到上帝的根前.
求主繼續來幫助我們.
基督耶穌是按著默基世德的等詞.
永遠為祭司.
基督耶穌為我們生命獻上最美好的祭.
\newpage



\section{約翰福音 1:19-38}
\label{sec:V98EonzJR3Q}
\textbf{2022年4月10日崇拜直播｜梁家麟牧師｜施洗約翰的見證(二)｜約翰福音1\_19-38}
\newline
\newline
連結: \href{https://youtube.com/watch?v=V98EonzJR3Q}{\texttt{https://youtube.com/watch?v=V98EonzJR3Q}} ~~~~ 語音日期: 2022-04-11
\newline
\newline
\hyperref[sec:Or14nIn97R8]{\small{< < < PREV SERMON < < <}}
~
\hyperref[sec:index_chronic]{\small{[返順時目]}}
~
\hyperref[sec:knwfu3yRKQ8]{\small{> > > NEXT SERMON > > >}}
\newline
\newline
約翰福音 1:19-38
\newline
\begin{longtable}{cl}
\hline
\hline
章節 & 經文 (和合本修訂版)\\
\hline
1:19 & \begin{tabularx}{0.7\textwidth}{X} 這是約翰的見證：猶太人從耶路撒冷差祭司和利未人到約翰那裡去問他：「你是誰？」 \end{tabularx} \\ \\ \relax
1:20 & \begin{tabularx}{0.7\textwidth}{X} 他就承認，並不隱瞞，承認說：「我不是基督。」 \end{tabularx} \\ \\ \relax
1:21 & \begin{tabularx}{0.7\textwidth}{X} 他們又問他：「那麼，你是誰？是以利亞嗎？」他說：「我不是。」「是那位先知嗎？」他回答：「不是。」 \end{tabularx} \\ \\ \relax
1:22 & \begin{tabularx}{0.7\textwidth}{X} 於是他們對他說：「你到底是誰，好讓我們回覆差我們來的人。你說，你自己是誰？」 \end{tabularx} \\ \\ \relax
1:23 & \begin{tabularx}{0.7\textwidth}{X} 他說：「我就是那在曠野呼喊的聲音：修直主的道。」正如以賽亞先知所說的。 \end{tabularx} \\ \\ \relax
1:24 & \begin{tabularx}{0.7\textwidth}{X} 那些人是法利賽人差來的。 \end{tabularx} \\ \\ \relax
1:25 & \begin{tabularx}{0.7\textwidth}{X} 他們就問他：「你既不是基督，不是以利亞，也不是那位先知，那麼，你為甚麼施洗呢？」 \end{tabularx} \\ \\ \relax
1:26 & \begin{tabularx}{0.7\textwidth}{X} 約翰回答：「我是用水施洗，但有一位站在你們中間，是你們不認識的， \end{tabularx} \\ \\ \relax
1:27 & \begin{tabularx}{0.7\textwidth}{X} 就是那在我以後來的，我給他解鞋帶也不配。」 \end{tabularx} \\ \\ \relax
1:28 & \begin{tabularx}{0.7\textwidth}{X} 這些事發生在約旦河東邊的伯大尼，約翰施洗的地方。 \end{tabularx} \\ \\ \relax
1:29 & \begin{tabularx}{0.7\textwidth}{X} 第二天，約翰看見耶穌來到他那裡，就說：「看哪，神的羔羊，除去世人的罪的！ \end{tabularx} \\ \\ \relax
1:30 & \begin{tabularx}{0.7\textwidth}{X} 這就是我曾說『那在我以後來的先於我，因為在我以前，他已經存在』的那一位。 \end{tabularx} \\ \\ \relax
1:31 & \begin{tabularx}{0.7\textwidth}{X} 我先前不認識他，如今我來用水施洗，為要使他顯明給以色列人。」 \end{tabularx} \\ \\ \relax
1:32 & \begin{tabularx}{0.7\textwidth}{X} 約翰又作見證說：「我曾看見聖靈彷彿鴿子從天降下，停留在他的身上。 \end{tabularx} \\ \\ \relax
1:33 & \begin{tabularx}{0.7\textwidth}{X} 我先前不認識他，可是那差我來用水施洗的對我說：『你看見聖靈降下來，停留在誰的身上，誰就是用聖靈施洗的。』 \end{tabularx} \\ \\ \relax
1:34 & \begin{tabularx}{0.7\textwidth}{X} 我看見了，所以作證：這一位是神的兒子。」 \end{tabularx} \\ \\ \relax
1:35 & \begin{tabularx}{0.7\textwidth}{X} 又過了一天，約翰同兩個門徒站在那裡。 \end{tabularx} \\ \\ \relax
1:36 & \begin{tabularx}{0.7\textwidth}{X} 他見耶穌走過，就說：「看哪，神的羔羊！」 \end{tabularx} \\ \\ \relax
1:37 & \begin{tabularx}{0.7\textwidth}{X} 兩個門徒聽見他的話，就跟從了耶穌。 \end{tabularx} \\ \\ \relax
1:38 & \begin{tabularx}{0.7\textwidth}{X} 耶穌轉過身來，看見他們跟著，就問他們說：「你們要甚麼？」他們對他說：「拉比，你在哪裡住？」（「拉比」翻出來就是老師。） \end{tabularx} \\ \\
[1ex]
\hline
\hline
\end{longtable}
$^{1}$讓我們一同去聆聽一段上帝的話語.
記載在約翰福音第一章十九至二十八節.
經文裡是這樣說.
「約翰所作的見證記載下面.
猶太人從耶路撒冷差祭司和利未人到約翰那裡問他說.
你是誰?.
他就明說並不隱瞞.
明說:我不是基督.
他們又問他說:.
這樣你是誰?.
是以利亞嗎?.
他說:我不是.
是先知嗎?.
他回答說:不是.
於是他們說:.
你到底是誰?.
叫我們好回覆猜我們來的人.
你自己說你是誰?.
他說:.
我就是那在曠野有人聲喊著說.
修籍主的道路.
正如先知以賽亞所說的.
那些人是法理賽人猜來的.
他們就問他說.
你既不是基督.
不是以利亞.
也不是先知.
為什麼施浸呢?.
約翰回答說.
我是用水施浸.
但有一位站在你們中間.
是你們不認識的.
就算那在我以後來的.
我給他解鞋帶也不賠.
這是在約旦河外百大里.
約翰施浸的地方作的見證.
今早有梁家倫牧師.
為我們分享剛才的約翰福音.
講題是「施洗約翰的見證」.
恭請梁家倫牧師.

$^{41}$我們繼續講約翰福音所記述的施浸約翰.
今次是第二講.
我們是經文是約翰福音一章十九到三十七節.
施洗約翰以典型的先知面貌.
在約旦河邊出現.
生活刻苦.
穿著落塘毛造的粗衣.
吃的在曠野能夠找到最疲勞的食物.
蝗蟲,野蜜.
野蜜不是指蜜糖.
是椰棗樹的果.
他離群居住.
講的教訓非常具社會批判性.
嚴詞譴責各種罪行.
宣告上帝的審判.
不信情面不畏權勢.
很多人都相信他是上帝所興起的先知.
施浸約翰在猶太人社會中間產生很大的激盪.
很多人前來聽他的教訓.
有人聽從他的呼籲.
接受了公開的悔罪洗禮.
也有些人跟隨他做了他的門徒.
馬太福音三章五到六節這樣說.
那時耶路撒冷和猶太傳遞並約旦河一帶地方的人.
都出去到約翰那裡承認他們的罪.
在約旦河裡受他的洗.
這段技術意味約翰的影響已經不限於約旦河一帶地區.
而是遍及整個猶太人的社會.
包括耶路撒冷.
約翰的影響力是大到一個地步.
威脅著猶太人的政權.
昏封王希律安提柏因為受希羅底和他的女兒的唆使.
竟然將約翰殺害.
這段經文約翰福音一章十九到三十七節提到.
一開始說有人前來查問耶穌.
十九節猶太人從耶路撒冷.
猜祭司和利未人到約翰那裡.
問他說你是誰.
前來查問他的是祭司和利未人.
他們是從耶路撒冷遠道猜險來到這裡.

$^{81}$猜派他們是誰呢?.
要查證約翰身份的是誰呢?.
很明顯一定是耶路撒冷的公會.
公會主要由大祭司家族控制.
其中不少是祭司和利美派的人.
另外一章二十四節說那些人是法利賽人猜來的.
法利賽人是當時猶太人中一個比較重要的宗教團體.
他們是經乾派主張嚴守律法的一點一.
一世紀的教會歷史家約瑟夫告訴我們.
當時期的法利賽人一共有六千人左右.
他們未必有足夠數量控制公會.
不過法利賽人卻控制了不少會堂的管理權和教導權.
所以他們是很有影響力的.
如果施西約翰已經產生廣泛的影響力.
不少人就猜測他是彌賽亞或者是以利亞.
公會的人一定不會掉以輕心.
坐視不理.
任何有這樣的影響力的宗教人物出現.
都有潛在的政治影響力.
或者他們會鼓動群眾危害宗教和社會秩序呢.
我們講句題外話.
楊寬陰裡面常常提到猶太人.
多數都不是指一般的猶太人.
而是指那些有權有勢的猶太教領袖.
並且多數都是敵視反對耶穌基督的人.
無論如何.
祭司利未人被差遣來問施西約翰.
第一個問題是你是誰.
這不是一般的問聖神明誰.
而是問他的屬靈身份和使命.
簡單的說法就是問他是否彌賽亞.
因為施西約翰的父親是撒迦利亞的祭司.
他們都是利未人.
所以很多祭司和利未人都認識約翰.
他們不會連他是誰都不知道.
公會老遠派人去查問約翰.
第一個關心的不是他叫什麼名字.
而是他是不是彌賽亞.
因為社會上有這樣的謠傳.
約翰很坦白地說我不是基督.

$^{121}$基督就是彌賽亞希臘文的叫法.
確定他不是基督.
查問他的人就放下第一個心頭大石.
污口了這個可能性.
他們接著問第二個問題.
那你是不是以利亞.
為什麼要提到以利亞呢.
事源在猶太人的拉比傳統裡面.
以利亞是擁有一個很特殊的地位和召命.
由於聶王記下二章十一節記載.
以利亞不是正常逝世.
而是帶著肉身升天.
很多人就猜想.
上帝為什麼要讓以利亞保持肉身升天呢.
因為他要保持他的臭皮囊.
將來再用.
所以以利亞的臭皮囊不會腐化.
因為將來他仍然會再來到人間.
被上帝再派來人間執行另外的使命.
所以以利亞不是彌賽亞的等級.
不過以利亞是為彌賽亞做預備.
所以以利亞的尼林.
應該是要先於彌賽亞的尼林.
以利亞是彌賽亞的前驅.
馬拉記書正是這樣預言的.
馬拉記書四章五節.
看吧,大而可畏之日未到以前.
我必猜顯,先知以利亞到你們那裡去.
由於人們猜想.
以利亞要先來,然後上帝的拯救才會發生.
所以不少人猜思西約翰就是要再來的以利亞.
從外觀上看,思西約翰和以利亞確實很相似.
都是穿毛衣,腰束皮帶.
也是嚴詞斥責人的罪惡.
撒迦利亞書曾經有趣地說.
有些假先知故意穿毛衣去騙人.
撒迦利亞書十三章四節.
說明穿毛衣是扮先知.
這是當時的標誌.
所以約翰和以利亞的外貌很相似.

$^{161}$實際上不僅約翰.
連耶穌也被人懷疑過是不是以利亞.
耶穌那次在馬拉福音十六章.
被門徒突擊測驗.
你們說我是誰.
有人猜他是以利亞.
無論如何,猶太人問思西約翰第二個問題.
你是不是以利亞.
以利亞約翰也一口否認.
我不是.
第二塊心頭大石也放下了.
猶太人再問第三個問題.
你是不是先知呢.
是那先知嗎.
這裡只說那先知,沒有提到先知的名字.
多數學者都相信.
這是對應摩西在離開世界之前.
向以色列民所做的應許.
這是記載在新領紀十八章十五到十九節.
即是說上帝會在末世的時代.
興起像摩西一樣的先知.
他要傳講上帝的話,要百姓遵守.
約翰再否認,我也不是那個先知.
他接二臨三,否認了各種對他的身份的臆測之後.
查問者確實已經放下所有的心頭大石了.
不過,他接著就產生一個疑惑.
如果約翰什麼都不是,那到底他是誰呢.
他們是前來查問他的身份的嘛.
總不能帶著所有否定的答案就離開的,是不是.
於是他們再肆意不捨,再問施厚生約翰二十二節.
你到底是誰.
叫我們好回覆,猜我們來的人.
你自己說,你是誰.
好玩吧.
他們不再用各種道聽途說得來的身份的臆測.
來去查問約翰.
而是要約翰自己申報自己的身份.
換句話說,他們不再問是非題了.
現在問的是填充題.
請約翰自行填寫,他是誰.

$^{201}$約翰的回答很巧妙.
其實整段英文都是很好玩的.
約翰沒有說明自己的身份.
他只是引述伊賽亞書四十章三節的一段說話.
說明他的使命和功能.
二十三節.
我就是那在曠野有人聲喊著說.
修籍主的道路正如先知伊賽亞所說的.
他的使命就是來為尼賽亞做預備.
好像在皇帝出巡之前.
總是有人為他先修橋鋪路.
做各種事先的安排,是不是.
約翰就是尼賽亞的開路先鋒.
人家問約翰你的身份是什麼.
約翰就用他的使命去回答.
查問施厄約翰的猶太人代表團中的法利塞人.
再追問下去.
二十五節.
你既不是基督,不是以利亞,也不是先知.
那為什麼你施洗呢.
法利塞人是比較注重禮儀的.
特別是對遵守禮儀的細節很執著.
所以他們就追問施厄約翰的水禮的具體情況.
洗禮是猶太人中間一個比較慣常的禮儀.
特別是用在赦罪方面.
不過一般來說都是尋求赦罪者自己做的.
自己浸自己的.
不需要別人幫我浸的.
如今約翰在河邊替人施浸.
這是一個創新性的做法.
對法利塞人來說.
創新即是違反傳統.
即是沒有舊有的做法作為一個背後.
他就質問你這個做法從何得來.
誰授權你可以這樣做.
你依據的是什麼權柄這樣做.
約翰對施浸權柄的問題.
同樣採取沒有直接回答.
甚至可以說是答非所問的回應.
在26,27節他這樣說.

$^{241}$我用水施洗.
但有一位站在你們中間是你們不認識的.
就是那在我以後來的.
我給得解鞋帶也不配.
人們問他權柄的依據.
約翰回答我什麼權柄也沒有.
特別跟接下來要來的比較.
我們之間的差距是天和地的差別.
我就是為那個要來的解鞋帶.
即是做奴隸做的事還是不配.
即是我做他的奴隸也沒有資格.
約翰呼音記載這句說話可能是不完整的.
施習約翰在這裡.
首先說的是我是用水施洗.
即是說他的媒介.
洗禮的媒介是水.
但就沒有說到下一個要來的要怎麼做.
用什麼媒介.
如果我們看馬福音.
即是約翰的原話是.
我是用水幫你們施洗.
但他就要用聖靈幫你們施洗.
馬太和路加就更詳細一點.
我是用水幫你們施洗.
但有一個能力比我大的要來.
他用聖靈和火幫你們施浸.
水是跟聖靈跟火相對照.
約翰這是一個答非所問的回應.
亦都引含一個意思.
猶太人所查問的三個身份.
彌賽亞,以利亞和先知.
都是屬於那個要來的耶穌.
他自己不是這三個身份.
但將來要來的同時擁有這三個身份.
約翰強調.
他是接著要來的那位開路先鋒.
跟要來的那個比較.
他是不重要的.
所以請不要把焦點放在我身上.
而應該把焦點集中在.

$^{281}$那個要來的耶穌基督那裡.
這段經文的開頭是.
十九節那裡說.
約翰所作的見證記在下面.
最後一句是二十八節.
這是在約旦河外伯大利.
約翰施洗的地方作的見證.
前後都強調這是約翰作的見證.
這是給人一個很官方很正式的.
提供證供的一個印象.
以下就是約翰的證供.
約翰的證供到此為止.
所以是前後兜著.
二十八節說故事發生的地點.
是約旦河西岸一個叫伯大利的地方.
這跟科林書常常說的伯大利不一樣.
科林書常常說的伯大利.
是在耶路撒冷南部附近.
靠近前往耶利哥的路上.
在約旦河的東岸.
這是馬大,瓦利亞和拉薩路的住所.
但約旦河西岸的伯大利到底在哪裡呢?.
到今天學者仍然眾說紛紜.
沒有一個確定的說法.
好,整段聖經都讀完了.
我們做一些思考應用.
我們記得約翰福音的佈局.
作者首先指出.
道成肉身這個偉大的事實.
這個事實發生在耶穌的身上.
然後說我們都是他的見證人.
他首先提出的是第一個見證人.
第六到第八節.
有一個人是從上帝那裡猜來的名叫約翰.
為要作見證,就是為光作見證.
然後在十九節到二十八節.
就說到施洗約翰的生平事蹟.
重申約翰不是光,不是利塞亞.
只是為利塞亞的來到做預備.
然後在二十九到三十四節就敘述.

$^{321}$約翰為耶穌施浸.
從而將主角帶出場.
我猜你看過《射雕英雄傳》.
金庸的名著.
頭幾章先說葛肖天和楊鐵心的故事.
然後就說到丘處機和江南七怪.
郭靖和黃蓉這兩個主角是山山來遲.
不是一開始第一章就說他們.
主角進場之前.
總是需要一些跑龍套的小角色在前面開路.
施洗約翰就是扮演這樣的角色.
他是引領耶穌進場的先驅.
剛才我們看的那段聖經.
正是約翰為耶穌做見證.
約翰在這裡是講的重點.
但不是主角.
他要傳遞的訊息是一連串的不是.
我不是基督.
我不是以利亞.
我不是先知.
我是不重要的.
和那個要來的比較.
我是微不足道的.
我不是那個光.
我只是為光做見證.
聖經故事的橋段是猶太人查問約翰的身份.
但從約翰口中帶出來的訊息.
卻無關乎約翰的身份.
只是道出了基督的身份.
即使基督還沒有事實上出場.
但是作者藉著猶太人對約翰的查問.
從而就已經交托了耶穌基督是誰.
猶太人多次追問約翰你是誰.
一個問了三次.
19節 21節 22節.
約翰都沒有回答.
他只說耶穌是誰.
約翰是誰是不重要的.
耶穌是誰才是重要的.
我們知道約翰的出現和他的特殊使命.

$^{361}$為猶太人社會帶來很多困擾.
所以有很多有關他的身份的猜測.
有很多謠傳.
這些謠言傳播之廣 影響之大.
竟然勞駕到撒冷的工會派人來調查.
我們不知道這些猜測.
這些謠言對約翰本人有什麼影響.
他有什麼感受.
被猶太人的工會調查.
大概不會是很愉快的事吧.
你都不會喜歡被廉署招待你喝咖啡.
你不喜歡被警察查你的身份證.
你不喜歡在過海關的時候.
被選中要做詳細檢查.
你一碰到就很煩.
同樣的我相信.
約翰都不會喜歡被人家去盤問.
不過被人懷疑被人誤會.
雖然一般來說不是很愉快.
不過如果被人誤會自己是名人偉人.
那就不一定很不愉快.
我在街上走.
如果有人認錯.
我猜我是梁朝偉或是譚詠麟.
或是台灣的歌星趙全.
唱「我很醜但我很溫柔」的歌手.
我是很高興的.
如果你誤會我是趙全.
我的樣子都很像他.
有一次我講完道.
有個弟兄走來跟我聊天.
他說他聽過我很多港島的錄音帶.
我是很少有錄音帶的.
上去一聽就知道他認錯了我.
可能認錯其他人.
不過我與其問他認錯李思敬或溫偉耀.
我就很客氣的跟他說.
我的樣子和Vic Warren和華利克都很像.
讓你誤會了真的不是很好意思.
有個好朋友跟我說了一個笑話.

$^{401}$他的女兒今年三十多歲.
已經是一女之母.
最近跟朋友去美國一個地方的賭場.
打算不是賭錢.
而是進去吃那些平價自助餐.
通常在賭場中間有個吃自助餐的位置.
誰知道賭場的門衛.
覺得這個三十多歲的女人.
Baby face 罵人.
就猜她不夠十八歲.
堅決不讓她進場.
她即使拿自己的護照出來.
對方都猜你是借別人的證件.
結果最終都沒有讓她進場.
女兒就打電話跟她媽媽.
即是我好朋友 素婦.
就在那裡嘔吐.
我就相信這些素婦只是說說而已.
心裡偷偷歡喜.
所有姐妹不管年齡.
都不介意別人誤會她不夠十八歲.
這真是美麗的誤會.
這聖經讓我們看到.
施厄約翰帶給人的困惑.
人總是傾向用舊有的知識和經驗.
去詮釋新的事物.
用舊有事物的分類系統.
去識別新的事物.
他們曾經以彌賽亞 以利亞.
和先知去詮釋約翰.
但是被約翰一一否定.
他們就無法為約翰做分類.
這真是非常困擾的事.
約翰為什麼不肯直接說他是誰呢.
或者他認為.
我作為開路先鋒.
根本不重要.
昔日大人物出場.
總是有些人明羅喝渡.
拿著儀仗在前面走.

$^{441}$或者今天總統出場.
當然有幾個穿著黑色西裝的保鏢.
在旁邊紮馬.
你需要知道那個保鏢叫什麼名字嗎.
不需要.
根本不重要.
他們只是明羅喝渡的人.
約翰沒有獨立的身份.
沒有獨立的使命.
他的身份和使命是連著基督的.
所以他不是什麼人.
他只是基督的什麼人.
約翰是先知.
是彌賽亞.
但他沒有用這個來做自我定義.
別人問他身份.
你會發覺很有趣的是.
約翰都是用使命來作回應.
我就是那在曠野.
有人聲喊著說.
修籍主的道路.
正如先知彌賽亞所說的.
約翰拒絕用一個名稱去描述自己.
連開路先鋒這個詞他都沒有用.
他只是用動詞去具體說明.
他實際的工作和工作的目標.
我們在這裡小小默想一下.
這個棄名詞.
而選擇用動詞的回應.
其實很有意思.
從前有位同事.
黃錫賢博士曾經說過一句名言.
我不是教書的.
我是教人的.
這句話蘊含極大的智慧.
教師的職責是教學生.
傳道授業計劃.
達不到這個目標.
師不師.
教師都不是教師.

$^{481}$我是一個傳道人.
但我最重要的職責.
是準確有效的傳講上帝的話.
我沒辦法履行這個任務.
傳道人這個名銜是白叫的.
我是一個牧師.
我最重要的職責是帶引信徒.
屬靈生命有所成長.
盡不到這個責任.
講什麼牧師這個身份.
都是思維素餐.
我們有時候講一個職銜.
一個稱呼.
我們要強調的是他的地位.
他的權力範圍.
包括他會獲得什麼待遇.
物質的回報.
或者社會上的尊敬.
你的職業是什麼.
你的職位是什麼.
你會有什麼的職權.
有什麼的待遇.
有人問我的職業.
我說我是醫生.
我是律師.
我是會計師.
我是精選師.
連面上的笑容都立即收起.
強調職務.
和強調職銜有所不同.
職務凸顯的是個人的責任和使命.
這個不是正態的.
是動態的任務.
並不是已經獲得的位置.
而是有待完成的使命.
Not yet.
聖經有多名銜的.
職銜的.
舊約的君王先知和祭司.
新約的仕途監督長老執事.

$^{521}$不過聖經.
其實是更加重視動態的.
責任和使命.
並且常常以.
是否履行使命來決定.
一個人是否擁有一個身份.
例如保羅.
在《少行傳》20章.
記載他和.
已忽了的長老做的.
臨別贈言.
我做了什麼.
我在你中間做了什麼.
這個比我是什麼人.
更重要.
職銜是名字.
使命是實.
如果你不是實踐某個使命.
你名不符實.
名字是沒有意思的.
例如保羅對.
哥林多信徒弟兄姊妹.
這樣說.
就算其他人懷疑他的仕途身份.
他們也不應該有懷疑.
因為他們是.
福音.
結出的果子.
他們是他的戰書.
保羅用自己的工作果效.
去證明他的身份.
另外.
《二十下書》九章六節.
也很有趣.
將耶穌基督稱為.
奇妙測試全能的上帝.
永在的父和平的君.
這些全部是講使命.
不是職位.
否則耶穌基督.

$^{561}$既是子又叫永在的父.
永在的父不是講職銜.
是講使命.
同樣耶穌基督被稱為.
以馬來利.
同樣是指他要實踐的使命.
上帝同人同在.
修補破損了的神人關係.
當時在約翰.
差派門徒查問.
耶穌是否來塞亞的時候.
耶穌的回應.
一如約翰本人.
耶穌的回應是.
他一連串的自己的作為.
《瑪利福音》十一章四到五節.
你們去把所聽見.
所看見的事.
告訴約翰.
就是乞子看見.
瘸子行走.
長大麻風的得潔正.
聾子聽見.
死人復活.
窮人有福音.
全給他們.
證明他的身份.
有些基督徒.
宣稱.
舊約上帝有很多名.
諸如什麼耶和華以納.
耶和華尼西.
這些對不起我要講.
是對聖經誤讀.
聖經上帝只有一個名.
就是耶和華.
前面講的那些名不是上帝的名.
是上帝在某一個地方.
顯現了.
成就一個偉大的作為.

$^{601}$人為了紀念這件事.
就把那個地方.
用這個名字.
這個地方的名不是上帝的名.
無論如何.
這也說明聖經重視的是實際的作為.
而不是耶和華尼西.
強調的上帝.
做了什麼.
名字.
是有關上帝作為的一個標示.
使命比名銜更重要.
教會的行政架構.
有不同的職位.
不同的職稱.
長老 執事 部長 組長.
都是為了.
時宮運作的需要而設立的位置.
但是無論是.
仙子 傳福音.
還是牧師和教師.
最大的目標.
是成全聖徒.
國盡其職 建立基督的身體.
職位是次要的.
功能才是主要的.
畢竟.
不管你有什麼名銜.
真正的身份仍然是耶穌基督的僕人.
這是做使徒的.
第一身份.
他們首先的自稱.
僕人就是僕人.
沒什麼巴閉.
我認識一間.
教會.
教會分成四個群體.
粵語 英語 普通話.
和青少年.
每個群體都有一個執事會.

$^{641}$每個執事會都有8-10個人.
執事會之上.
有一個長老會.
是統管整間教會的.
成員有7個人.
我數一數 他的教會長老.
執事加起來接近50人.
而教會年大年細.
連同小朋友 才500多人.
換言之 長執.
佔了總人數十分之一以上.
人人頭頂都有一頂帽子.
從正面的角度來說.
這是人人參與事奉.
這是美好見證.
從負面的角度來說.
如果這麼多主事人.
教會耗費在開會的時間.
維持行政體制的資源.
就很多了.
更加不要說.
會有比較多人士.
行政問題出現.
官位越多 行政架構越臃腫.
效率就越低.
這是柏金遜定律.
唯有.
每個信徒不管在教會.
當中居住什麼位置.
都認定他的責任.
是踐行基督的使命 誡命.
唯有每個信徒.
都認定事奉的目標.
是讓基督的名被高舉.
讓福音得以有效傳開.
讓更多人能夠分享福音的好處.
我們前面說的.
行政人士問題.
才有望輸解.
教會只有一個真正的老闆.

$^{681}$就是耶穌基督.
其他所有人都是跑龍套.
都是預備主的道.
修正他的道路.
因此都是不重要的 不需要受注目.
不需要被榮耀.
正如約翰知道.
在人間自稱為先知維尼爾利亞.
可以得到很多榮耀.
人們會擁戴他 支持他.
但他咬定.
他唯一要承擔的職務.
就是預備將來要來的那一位.
而跟要來的那位相比.
他一文不值.
連解鞋帶都不配.
更不要說榮耀.
地位.
目標端正了.
心態調整了.
教會群體才可以.
同心合二 興旺福音.
弟兄姊妹.
讓我們校正我們事奉的焦點.
執著我們要.
完成的使命.
不是執著自己的身份職稱.
不要看重自己在教會的身份地位.
不要計較自己能夠獲得.
什麼權力和待遇.
好像施洗約翰一樣.
認定自己就是耶穌基督的僕人.
是他的開路先鋒.
我們唯一的責任是將耶穌介紹給人.
高舉基督.
彰顯基督 讓基督的名得著.
最大的榮耀.
我們都是不重要的人.
都是微不足道的小配角.
只有基督才是真正的主角.

$^{721}$我記得.
一個天國同行者說過這句話.
Live to be forgotten.
So that Christ can be remembered.
讓我們每個人都被忘記.
好叫.
我們大家都只記得耶穌基督.
這是保羅強調.
Not I.
But Christ.
非我為主的真實含義.
不強調角色.
只強調使命.
這才不會造成身份遺失.
施洗約翰為我們樹立了一個.
很美好的榜樣.
讓我們知道.
如何確定目標.
專注使命.
不至陷落到權力.
名聲的網羅裡面.
最後.
我們一起和應.
約翰所說的話.
我不配.
我連為他解鞋帶.
我都不配.
弟兄姊妹.
這是很違反我們.
今天的現代人的觀念.
我們今天.
常常說恩賜.
常常說召命.
我們總是繪形繪聲地說.
每個人都從上帝那裡.
領受一個獨特的恩賜.
一個獨特的召命.
這個恩賜.
這個召命是他一個人擁有的.
所以非他莫屬.

$^{761}$離開他.
就沒有人可以承擔.
所以每個人都是獨特的.
我要說這些話.
就是今天的人.
很喜歡聽,不過是沒有營養價值的.
我在.
事奉生涯裡面.
從來沒有這個發現.
我是.
上帝特別訂做.
為見到神學院.
做院長的.
如果沒有我出山的話.
見到就倒塌了.
見到就迷失方向.
我沒有這個發現.
我只知道.
屬中無大將.
療法作先鋒.
我只知道上帝在故意.
一個不稱職的人在這個位置.
裡面.
讓人清楚看到.
真正見到的院長.
是上帝.
不是梁家綸.
不需要一個稱職的人.
因為上帝是上帝.
祂在人的軟弱之上.
才顯得祂的能力的完全.
所以,弟兄姊妹.
不要覺得自己.
有多不可或缺.
覺得自己有多厲害.
永遠認定.
上帝.
竟然.
肯讓我.
去參與.

$^{801}$祂偉大的使命.
讓我在其中.
雖然微不足道.
但足以讓我.
擺上我的一生.
奮盡.
我所有的智慧和能力.
我能夠.
這樣參與,我榮幸.
但我知道.
我是可以被取代.
因為我本來.
不是一個夠資格.
不是一個配的人.
我們一班不配的人.
一起去宣認.
耶穌基督.
才是我們的主角.
耶穌基督.
才是這個世界.
所需要的答案.
耶穌基督,願你得到最大的榮耀.
主席.
\newpage



\section{ }
\label{sec:knwfu3yRKQ8}
\textbf{2022年4月15日受苦節晚會直播｜十架上的你與他/她}
\newline
\newline
連結: \href{https://youtube.com/watch?v=knwfu3yRKQ8}{\texttt{https://youtube.com/watch?v=knwfu3yRKQ8}} ~~~~ 語音日期: 2022-04-16
\newline
\newline
\hyperref[sec:V98EonzJR3Q]{\small{< < < PREV SERMON < < <}}
~
\hyperref[sec:index_chronic]{\small{[返順時目]}}
~
\hyperref[sec:U9MWFdIg8Fc]{\small{> > > NEXT SERMON > > >}}
\newline
\newline
$^{1}$我們所愛的主在二千年前為我們的罪被釘死在十字架上.
今晚我們在網上聚集.
我們不是要開一個追思會.
不是要搞到氣氛很悲哀,很憂愁.
不是這樣的.
我們不是只來這裡緬懷一件舊事.
紀念耶穌基督的可歌可泣.
為我們捨棄他的生命.
因為我們知道耶穌基督已經復活了.
他已經藉著復活戰勝了死亡和罪惡.
並為我們每一個人開出一條又新又活的救恩之路.
那與一個消防員殉職死了.
我們一起很哀傷.
因為他死了就是死了.
不同的,耶穌基督已經復活了.
所以他的受苦死亡只是一個過程.
只是一個環節,不是那個終點.
我們在這裡做些什麼呢?.
一方面我們傳說耶穌基督給我們的救恩.
另一方面我們思考十字架這個事件對我們的意義.
我們要記住十字架事件不只是屬於耶穌的.
他是作為我們的先行者,作為我們的示範.
十字架事件是為我們每一個人而發生的.
他不僅僅是屬於主的,也是屬於我們的.
耶穌基督曾經跟我們說過的.
他的苦杯我們都要喝的.
這個十字架是我們每個人都要背負的.
所以我們不只是去見到耶穌的一個事件.
我們要在耶穌的事件裡面去見到我們自己.
去見到我們這裡每一個人.
我們的命運,我們的將來.
所以我希望受苦節是一個讓我們去沉思.
去反省的一個時候.
不一定是讓我們很哀傷,去悼念一個不幸殉難的人.
不是這樣的.
我想將今天的訊息分為兩部分.
第一部分我們說的是.
我們怎樣陪同我們的主走過他一生的路.
經歷他的一生.
我們會從一些聖經的人物.

$^{41}$特別是一些婦女來思考這件事.
耶穌基督被捕之後.
他的追隨者大到逃跑狀.
先是所有的門徒在黑西瑪利園各自奪路而逃.
連最勇敢的彼得曾經敢拔刀.
後來也都在大祭司的院子被人認出來之後.
而趕快逃跑.
馬科音十四章五十節那裡說.
門徒都離開他,就是耶穌基督逃走了.
約翰在後來折返伽略山.
得以見到耶穌基督最後一面.
並且被耶穌受託去照顧他的母親.
其他的門徒有沒有偷偷折返.
不過不管怎樣.
他們都曾經在耶穌受苦和死亡的重要關頭.
他們缺席,看不到約翰的重要片段.
沒有分享到耶穌基督在人生最後的時光.
只是有幾個婦女,竊而不捨地跟隨主.
無論在受刑的時候,被釘的時候.
甚至是復活後的第一時間裡面.
馬科音十五章四十到四十一節那裡說.
馬科音十五章四十七節那裡說.
耶穌安葬的時候,他們負責買香膏去糕耶穌.
十六章一節,他們沒有錯過任何一個片段.
沒有在耶穌遭遇危難的時候離棄過他.
馬科音十六章記載了他們的名字.
包括了以下幾個人.
末大拉的瑪利亞,小亞國和約西的母親瑪利亞,還有撒羅米.
這裡有兩個女性都叫瑪利亞.
末大拉的瑪利亞,小亞國和約西的母親瑪利亞.
瑪利亞是猶太人婦女中一個非常普遍的名字.
就好像我那個年代,叫小英小麗.
或者較後的年代,叫做詠詩,叫做曉彤等等.
都是一個很普通的名字.
誰是小亞國和約西的母親瑪利亞呢?.
最有可能的就是耶穌的母親.
因為馬科音六章三節那裡說.
耶穌基督不是瑪利亞的兒子.
亞國,約西,猶太,西門的長兄嗎?.
如果亞國和約西是耶穌的兄弟.

$^{81}$那麼亞國和約西的母親就是耶穌的母親.
當然我們知道羅馬天主教的聖經學者是反對這個說法的.
因為他們堅持耶穌的母親瑪利亞是永遠的童貞.
所以耶穌的兄弟其實只是他的堂兄弟.
所以這裡亞國和約西的母親瑪利亞.
可能就是另有其人了,在天主教的說法裡面.
就不是耶穌的母親了.
我們當然不需要跟隨天主教的說法.
所以我們最方便的說法就是.
耶穌的母親一直跟隨著耶穌.
至於穆達拉的瑪利亞,我們不會對這個名字陌生的.
特別如果很多年前有一本很邪門的書叫達文西密碼.
有人看過這本書的話,都會相信那個作者的歪論.
就說穆達拉的瑪利亞是耶穌的秘密太太,秘密情人.
今天我們不說這些了.
很多人相信穆達拉的瑪利亞是耶穌曾經拯救過一個妓女.
盧卡夫欽第七章記述.
當耶穌在伯大利去到一個法利賽人的家庭坐直的時候.
有一個女人用香膏勾耶穌的腳,聖經說她是一個妓女.
耶穌沒有嫌棄她,反而說因為她的罪多.
上帝對她的赦免也多,她對耶穌基督的愛也多.
然後耶穌就公開赦免了她的罪.
這個伯大利的婦女是不是穆達拉的瑪利亞呢?.
我們是沒有辦法百分百確定的.
但是盧卡夫欽第八章就提到,穆達拉的瑪利亞曾經被鬼附.
耶穌將附在她身上的七個鬼都趕了出來.
盧卡夫欽八章第二節,這件事就一定是百分百確實的了.
穆達拉的瑪利亞是耶穌基督最忠實的女門徒.
我們將馬福音的記述和其他福音書做比較.
馬福音二十七章五十六節也有提到這些婦女的名字.
包括穆達拉的瑪利亞,亞國和約西的母親瑪利亞.
並西比泰的兩個兒子的母親.
馬太提到西比泰的兒子,門徒亞國和約翰的母親.
約翰福音就說,在耶穌基督釘十字架的時候.
他所愛的門徒約翰都在場.
耶穌親自將照顧母親的責任交託給約翰.
如果約翰的母親都在場,那約翰在場就比較合理.
至於約翰福音十九章二十五節就說.
站在耶穌十字架旁邊的有他的母親和母親的姐妹.
並加羅巴的妻子瑪利亞和穆達拉的瑪利亞.

$^{121}$有學者估計馬福音所說的撒羅米.
就是馬太福音所說的西比泰的兩個兒子亞國和約翰的母親.
而亞國和約翰的母親就是耶穌的母親瑪利亞的姐妹.
如果這個估計是正確的話.
那耶穌和亞國和約翰,即是這兩個門徒是有親戚關係的.
這個為什麼亞國和約翰的母親曾經開口.
要耶穌將來容許亞國和約翰分坐在耶穌的左右邊.
照顧親戚,稱善一點,是不是很合理的要求.
將這三組的經文合併起來.
我們可以估計圍繞在耶穌旁邊的大概有四位婦女.
包括耶穌的母親瑪利亞,他的姐妹西比泰的兩個兒子的母親.
穆達拉的瑪利亞和加羅巴的妻子瑪利亞.
這些婦女和耶穌基督的關係都很深厚.
耶穌的母親和他的姨媽,不知是姨媽還是阿姨,我估計是姨媽.
是不在話下,他們都是看著耶穌長大成人.
在耶穌的傳道生涯一直是相伴.
穆達拉的瑪利亞早在耶穌還在北面的加利利地區事奉的時候.
就已經跟隨他了,馬可福音十五章四十一節.
這些婦女或者他們是為耶穌和使徒預備飯食,照顧耶穌的起居.
他們一直跟隨耶穌,從加利利去到耶路撒冷不離不棄.
如果不是耶穌的使徒團隊裡面是沒有女性.
我們可以說他們應該和彼得雅各約翰並列為耶穌的使徒.
沒錯,他們不屬於耶穌最核心的團隊.
最後晚餐的時候都沒有預留他們的位置.
他們都沒有上到赫西瑪利園,不過也可能有的.
但是沒有記載到,我們相信在耶穌晚年的傳道生涯裡面.
多數場合這幾個婦女都是一直不離不棄地跟著的.
他們為何能夠這麼專一地跟隨主.
一個解釋是他們或者獨身,或者寡居,或者年老.
沒有家庭的婦女可以自由往來,這是客觀的條件.
主觀地他們都和主有密切的私人關係,有特殊的經歷.
例如穆達拉的瑪利亞是曾經被耶穌基督趕鬼醫治過,拯救過.
這個出死入生的經歷使他們非常信從耶穌.
不管耶穌的遭遇是怎樣,他們都死心塌地跟隨到底.
他們未必好像彼得那樣能夠說得出好聽的話.
彼得說過,主啊,你有永生之道,我們還跟隨誰呢.
他們未必懂得說這句話,不過他們用實際行動將這句話兌現出來.
你救過我,我不跟你還跟誰?.
這是他們可能心裡最關鍵的想法.
當然我們還可以說一件事,這班婦女之所以一直不離不棄地跟隨耶穌.

$^{161}$是因為她們都是極少人物,極少人物.
她們是普羅大眾,不用說了,在男女地位懸殊的社會.
婦女更加沒有地位,簡單說她們是小人物.
所以她們跟隨耶穌,在耶穌受審受苦被殺的場合.
她們都跟隨著耶穌,其他門徒都走了,原因很簡單.
不在於這班婦女夠勇敢,願意犧牲一切.
而是在於她們沒甚麼好犧牲,小人物,nothing to lose,沒甚麼可以輸.
她們不構成對政權,對那些當權派任何的威脅.
根本那些當權派連理她們都多餘,根本不理她們.
在上世紀六十年代七十年代,中國教會面臨嚴重的打壓.
在那個時候能夠留下繼續傳福音做見證的都是一班老太太,老伯伯.
她們不是因為勇敢站出來,而是因為她們沒甚麼可以輸.
Nothing to lose,這是給我們一個很好的提醒.
我們越是自覺自己是somebody,越是自覺自己了不起.
越是覺得自己有很多東西,很容易會輸.
我們就越不敢跟隨主,越不敢捨棄.
唯有是你自覺一無所有,一無所事.
你就會好像那班婦女那樣,總之跟就跟吧,死就死吧.
那她們就可以跟隨到底.
婦女的專一跟隨,真的對我們有很多提醒.
我們被人間的事務纏繞得越少,我們就越能夠專一跟隨主.
沒錯,領受上帝的恩典是驅動我們跟隨主的主要動力.
不過,只是領受恩典,不等於我們對施恩的主一定有相應的回應.
你記得路加福音17章那個技術.
耶穌有一次同時醫好了十個長大麻瘋的人.
只有一個回來答謝耶穌,我們領受了恩典.
我們很容易成為了忘恩負義的人.
不是我們本性良薄,而是因為我們世務纏繞,生活忙碌.
我們思慮煩擾,很多成年弟兄姊妹都有類似的經歷.
工作越來越忙,時間都被困擾,心也被困擾.
愛主的心,你想淡薄都不行.
耶穌基督說,人到我這裡來,若不愛我勝過愛我的父母.
妻子兒女弟兄姊妹和自己的性命,就不能作我的門徒.
這句話,我們常常認為我們的主是小器,是霸道.
對我們無理苛求,我們沒有想到,他其實是描述性.
說明我們身邊所有的人和事,都可以構成對我們的跟隨主一個妨礙.
是的,不勝過,就被勝過.
愛主不勝過一切,就被一切勝過愛主.
使徒逃跑了,婦女留下了,他們在耶穌身邊.
伴隨耶穌上十字架,伴隨耶穌入墳墓,也見證了他的復活.

$^{201}$馬福音的長結尾,那裡這樣說.
馬福音十六章九節,在七日的第一清早,耶穌復活了.
就先向穆達拉的瑪利亞顯現,耶穌從她身上趕出七個鬼.
穆達拉的瑪利亞是第一個見證耶穌復活的人,早於一切的仕途.
這班婦女是第一批經歷一生的信徒.
他們見證過耶穌基督的受苦,死亡和復活.
不過他們不是什麼出類拔萃的人.
以穆達拉的瑪利亞為例,在聖經的記述裡.
她沒有說過任何一句有一點點智慧含量的話.
問過一句有智慧含量的問題.
甚至你留意一下,聖經沒有寫到她跟耶穌有對話過.
我相信一定有聊過天,但沒有記下任何對話.
她都沒有在耶穌的工作裡扮演過什麼重要的角色.
不過從路加福音八章提到她的名字之後.
她就一直跟隨耶穌,由加利利到耶路撒冷.
從城裡到加略山,從加略山到空墳墓.
她完完整整經歷了耶穌基督的一生.
保羅這樣說,耶穌的一生是跟我們每一個人有密切的關係.
所有信徒都要跟主一同經歷祂的受死,埋葬和復活.
我們跟耶穌同死同埋葬同復活.
羅馬書六章五節八節這樣說.
我們若在祂死的形狀上與祂聯合.
又在祂復活的形狀上與祂聯合.
我們如果是與基督同死,就信必與祂同活.
腓立比書三章十到十一節都說.
使我認識基督,曉得祂復活的大能.
並知曉得和祂一同受苦,效法祂的死.
或者我也得以從死裡復活.
弟兄姊妹,耶穌基督在二千年前.
在大祭司蓋亞法和秦府彼拉多的院子裡受審的時候.
在築樓地被釘死的時候.
在亞利馬大的藥室所預備的墳墓裡被埋葬的時候.
我們都一定不在場,不在其中的,是不是?.
但是我們今天是不是與主同活.
是不是天天都披戴基督,好像加拉太書三章二十七節所說.
我們是不是將耶穌的生平變成我們的生活.
以至我能夠活出祂的生命.
保羅說,耶穌基督釘十字架已經活化在我們眼前.
加拉太書三章一節,我們已經與基督同釘十字架.
現在活著的不再是我,但是基督在我們裡面活著.

$^{241}$加拉太書二章二十節.
所以不要錯過耶穌基督任何一個片段.
不要有任何時間離開祂.
什麼叫做與基督同活,什麼叫做活出基督.
一方面,這是指著我們遵行主的吩咐,學校的榜樣.
並且刻意設想,怎樣更加有效地跟隨祂,以祂的心為心.
好像一首我很喜歡的詩歌,說:直到我摸到你愛世人的心.
我希望我可以摸得到我的主愛世人的心.
這是活出基督第一方面的意思.
另一方面的意思是,用基督的生平去詮釋我們所有的遭遇.
我們自己在生活裡面去發現基督.
也在基督的生平裡面去發現自己.
以後有機會跟大家詳細說這件事.
怎樣將我們的生活基督化.
我們知道如果一切發生在我們生活裡面的事.
都是為了基督,都是出於上帝.
我們就可以淡定安然地做人.
這是第一部分我想跟大家說的.
我們跟基督同行,經歷祂的一生.
另一方面,第二部分我想跟大家說的是.
把握機會去事奉主.
我們今天是壽苦節,本來我不應該鬥到復活節的那部分.
留待星期日,我們在復活節聚會的時候會有詳細的講論.
不過,講到勃大拉的瑪利亞.
我們都要輕輕提一提她在復活的那一天的反應.
勃大拉的瑪利亞,這班婦女不僅僅忠心耿耿地跟隨主.
她們更加把握每一個機會,每一個時間去服侍主.
在耶穌死後第三天,安息日剛剛過去.
人剛剛停止,quarantine,禁足,准出街.
太陽仍然未出來,天色仍然黑暗.
勃大拉的瑪利亞已經第一時間撲去耶穌的墳墓.
她沒有放棄服侍主的任何機會.
她心裡想,這是她最後可以為主做的事.
就是為即將腐爛,在近冬這麼熱的環境,屍體是很快臭很快爛的.
在耶穌的屍體腐爛到不似人形的時候,先為他的屍體塗上香膏.
這件事有什麼特別的屬靈意義,我們不知道.
有什麼實際的果後,我們也不知道.
反正傳統是這樣做,總之這是習俗.
勃大拉的瑪利亞努力將它做到最好.
為此教會相傳,百大來用香膏,.

$^{281}$膏耶穌的女子就是勃大拉的瑪利亞.
勃大拉的瑪利亞最擅長幫人塗膏,最少幫耶穌也塗過兩次.
勃大拉的瑪利亞把握機會,做她能力範圍裡面最好的事.
她將最寶貴的香膏帶到耶穌面前,把握最後的機會去高抹耶穌.
其他人不欣賞她,說她浪費,她也不懂得怎樣自辯.
反正她也不知道這樣做是不是很有價值很有意義.
我本著心想這樣做,想做些好事給耶穌,就這樣.
耶穌卻為她的作為賦予神聖的理由,說為她的腳去高抹.
是為耶穌的安葬而做,並且強調這個故事要被日後傳頌.
勃大拉的瑪利亞不是首先有一個神聖理由,然後去做一個服事.
她是先做一個服事,然後耶穌為她的服事賦予一個神聖的價值.
這班婦女全心全意服事主,從來沒有期望得到特別的注意.
不要說甚麼上報,一切都是出於自發的感情表達.
她們思念耶穌,愛慕耶穌,做能力範圍所給的事.
就好像我媽媽不知道我每天忙甚麼,只是間中.
她還年輕的時候,現在她已經九十歲當然不行.
她比較年輕的時候,她間中去我家煲煲湯給我喝.
然後每一次打電話的時候就提醒我穿衣服,不要喝那麼多生冷的東西.
不要吃那麼多生冷的東西,這個對我的事奉有甚麼直接幫助.
恐怕我媽媽自己也不確定,不過我告訴你.
我媽媽在我身上所說的話,所做的事,對我means a lot.
幫助之大,只有我自己知道.
我們努力做一件美事在耶穌身上,把握機會,竭盡所能地服事祂.
弟兄姊妹,我們活在一個高度分工的世界裡面.
日常生活所使用的東西,即使是很簡單的東西.
例如一支筆,都是由很多人互相協作生產出來的.
有人負責製造塑膠的原材料,有人將原材料啤成膠模.
也有人製造原子筆,這個高度分工的經濟體系.
不需要任何人整體負責管理,指揮每一個環節生產甚麼東西.
生產多少,如何配合供求,經濟學一個價格理論指出.
社會的供求由價格自由調節,不需要政府,不需要任何人去管理.
一個美國的經濟學家訪問中國,和中國政府的物資部官員聊天.
那個官員正在計劃訪問美國,想清楚美國的物資供應是如何安排.
他問美國的經濟學家,應該和哪個部門的官員接觸.
美國經濟學家告訴他,美國根本沒有這樣的政府部門.
也沒有人管理物資的供應,價格的調整.
同樣每一個基督徒努力做好自己的本份.
就是他相信上帝交付給他的責任.
然後上帝就用他無形的手,將每一個人的努力拼合成為一幅完整的圖畫.
我們都未必知道全局,我們都未必很清楚自己的位置.

$^{321}$我們都未必很清楚我們所做的有多大的意義.
但我們每個人努力做好自己的部份.
上帝這個編劇和導演,他就會合在一起成就其他的事.
好像你看到我在錄製這個壽苦節的訊息.
你沒有看到有兩位弟兄坐在我前面,他們在做很多事.
其中一位弟兄比我更早來到這裡做set up(安排).
做很多很多的事.
我們各自在做我們自己,不是很知道有多重要.
不過肯定是很重要的部份.
上帝合在一起,他就會使我們這個聚會.
就會使我們教會的服事幫了很多很多的人.
自從被耶穌基督拯救之後,.
梅達拉的瑪利亞將他一生掛疊在耶穌身上.
從前他被鬼附的時候,活在魔鬼的捆綁之中.
後來得到耶穌基督拯救他,他人生方向改變了.
他就一心一意跟隨耶穌到底.
當耶穌死了之後,瑪利亞的人生霎時間失去了方向.
雖然耶穌基督曾經多次告訴門徒他會從死裡復活.
不過連仕途都不知道耶穌在說什麼.
就不要說沒有知識的梅達拉,瑪利亞.
他根本沒有想像耶穌復活的事.
所以,梅達拉,瑪利亞在復活之日去到墳墓.
發現耶穌的屍體不見了的時候,他不是想像耶穌復活了.
他心中其實是有強烈的失落感,孤單無助.
之後有兩個門徒來到發現相同的景象.
那兩個門徒以為是有冰釘偷了耶穌的屍體,於是就離開了.
但是,梅達拉,瑪利亞就留在墳墓旁邊,哭了,倦倦不捨.
耶穌向他顯現,問他為什麼哭的時候,他說有人拿走了耶穌.
在他心中,這個死了的冷冰冰的屍體仍然是他唯一的主.
他以為跟他說話的是一個圓釘.
於是他就跟那個人說,你知不知道誰人拿走了耶穌.
你告訴我,他們把耶穌放在哪裡?他想拿回.
梅達拉,瑪利亞只是想跟耶穌在一起,不管耶穌是死是活.
耶穌呼喚他的名字,他彷彿聽到一把熟悉的聲音.
然後他才認出這個是主,他恢復人生的希望,黑暗一掃而空.
所以,弟兄姊妹,我想說什麼呢?.
我想說的是,我們不要認為這些婦女是很勇敢,很出色的人.
他們比平凡還要平凡.
馬火音記載,當天使吩咐他們向門徒宣告.
耶穌已經復活了的訊息的時候.

$^{361}$他們竟然因為害怕的緣故,不敢說.
馬火音16章8節那裡說,他們就出來從墳墓那裡逃跑.
又發抖又驚奇,什麼也不告訴人,因為他們害怕.
是不是很陰春啊,這班婦女,是不是很差勁.
基本上是很差勁的一班人.
你想想,我們也可以將心比心的.
一班小婦女,看到這麼奇異,這麼獨特的事件.
墳墓門打開,屍體失蹤,天使顯現,怕是很正常的.
誰碰到這些超自然的事件,不會很驚惶呢.
所以,這班婦女很害怕,他們沒有執行天使的命令.
他們不是屬靈為人.
馬火音是以婦女的害怕作結,是不是很壓突.
不過,是這樣的,人的害怕,說明他們對上帝的大能敬畏.
人的害怕,凸顯了上帝的大能.
人的害怕,也說明了人的無助和無能.
我們不知道面前會發生什麼.
我們不知道上帝在面前有什麼特別的作為.
就好像今天,我不想將一個簡單的安慰的訊息給大家.
面對排山倒海的疫情,面對香港的經濟乃至政治的困局.
我們有沒有害怕,我們有沒有無助無能的感覺.
有,這個是我們今天的本相.
不過,我期望在這個起點,我們會見到上帝叫我們敬異的作為.
我們不知道前景,我們只知道守住我們的本份.
做我們應該做的事.
無論如何,受苦節不是一個專門懷念一個人的悲慘故事.
令到大家很傷感,不管我們如何為耶穌的受苦和死亡哀傷.
我們都知道這個故事已經有一個快樂的結局.
那個快樂的結局不用等星期天發生.
我們今日之成為基督徒,我們生命得改變.
不就是因為這個快樂的結局,是不是?.
耶穌基督已經復活,戰勝死權,祂已經得勝.
我真的很希望我們看到,耶穌的死亡和復活彰顯上帝的大能.
一個叫我們害怕的大能.
我們不是用哀傷作結,是用詫異,用懼怕作結.
上帝竟然為我們這一班這樣的人死而復活.
我們能夠不懼怕嗎?.
我們能夠在此時此刻的香港去成為基督徒.
在一個什麼前景都看不通的環境下,我們敢去為祂作見證.
宣告我們的主已經受苦死亡,祂即將復活.
我們能夠不懼怕嗎?.

\newpage



\section{希伯來書 8:1-13}
\label{sec:U9MWFdIg8Fc}
\textbf{2022年4月17日復活崇拜直播｜陳明泉牧師｜超越的基督:新約的中保｜希伯來書8\_1-13}
\newline
\newline
連結: \href{https://youtube.com/watch?v=U9MWFdIg8Fc}{\texttt{https://youtube.com/watch?v=U9MWFdIg8Fc}} ~~~~ 語音日期: 2022-04-17
\newline
\newline
\hyperref[sec:knwfu3yRKQ8]{\small{< < < PREV SERMON < < <}}
~
\hyperref[sec:index_chronic]{\small{[返順時目]}}
~
\hyperref[sec:rKFXkQ539_w]{\small{> > > NEXT SERMON > > >}}
\newline
\newline
希伯來書 8:1-13
\newline
\begin{longtable}{cl}
\hline
\hline
章節 & 經文 (和合本修訂版)\\
\hline
8:1 & \begin{tabularx}{0.7\textwidth}{X} 我們所講的事，其中第一要緊的就是：我們有這樣一位大祭司，他已經坐在天上至大者寶座的右邊， \end{tabularx} \\ \\ \relax
8:2 & \begin{tabularx}{0.7\textwidth}{X} 在聖所，就是在真帳幕裡作僕役；這帳幕是主所支搭的，不是人所支搭的。 \end{tabularx} \\ \\ \relax
8:3 & \begin{tabularx}{0.7\textwidth}{X} 凡大祭司都是為獻禮物和祭物設立的，所以這位大祭司也必須有所獻上。 \end{tabularx} \\ \\ \relax
8:4 & \begin{tabularx}{0.7\textwidth}{X} 他若在地上，就不用作祭司，因為已經有照律法獻禮物的祭司了。 \end{tabularx} \\ \\ \relax
8:5 & \begin{tabularx}{0.7\textwidth}{X} 他們所供奉的本是天上之事的樣式和影像，正如摩西將要造帳幕的時候，神警戒他，說：「要謹慎，一切都要照著在山上指示你的樣式去做。」 \end{tabularx} \\ \\ \relax
8:6 & \begin{tabularx}{0.7\textwidth}{X} 如今耶穌已經得了更優越的事奉，正如他作更美之約的中保；這約原是憑更美之應許立的。 \end{tabularx} \\ \\ \relax
8:7 & \begin{tabularx}{0.7\textwidth}{X} 第一個約若沒有瑕疵，就無須尋求第二個約了。 \end{tabularx} \\ \\ \relax
8:8 & \begin{tabularx}{0.7\textwidth}{X} 所以神指責他們說：「主說，看哪，日子將到，我要與以色列家和猶大家另立新的約； \end{tabularx} \\ \\ \relax
8:9 & \begin{tabularx}{0.7\textwidth}{X} 不像我拉著他們祖宗的手領他們出埃及地的時候，與他們所立的約；因為他們不恆心守我的約，所以我也不理他們；這是主說的。 \end{tabularx} \\ \\ \relax
8:10 & \begin{tabularx}{0.7\textwidth}{X} 主又說：那些日子以後，我與以色列家所立的約是這樣：我要將我的律法放在他們的心思裡，寫在他們的心上；我要作他們的神，他們要作我的子民。 \end{tabularx} \\ \\ \relax
8:11 & \begin{tabularx}{0.7\textwidth}{X} 他們各人不用教導自己的鄉親和自己的弟兄，說：你要認識主；因為從最小的到最大的，他們都要認識我。 \end{tabularx} \\ \\ \relax
8:12 & \begin{tabularx}{0.7\textwidth}{X} 我要寬恕他們的不義，絕不再記得他們的罪惡。」 \end{tabularx} \\ \\ \relax
8:13 & \begin{tabularx}{0.7\textwidth}{X} 既然神提到「新的約」，那麼第一個約就成為舊的了；而那漸舊漸衰的必然很快消逝了。 \end{tabularx} \\ \\
[1ex]
\hline
\hline
\end{longtable}
$^{1}$今日的經訓系記載在希伯來書8章1至13節.
「我們所講的事,其中第一要緊的.
就是我們有這樣的大祭事.
已經坐在天上至大者寶座的右邊.
在聖所,就是真帳幕裡作執事.
這帳幕是主所知的.
不是人所知的.
凡大祭司都是為獻禮物和祭物而設立的.
所以這位大祭司也必須有所獻的.
他若在地上,必不得為祭司.
因為已經有照律法獻禮物的祭司.
他們供奉的是本是天上的事的形狀和影像.
正如摩西將要做帳幕的時候.
蒙神警誡他說:.
「你要謹慎.
作各樣的物件都要照著在山上指示你的樣式」.
如今,耶穌所得的職任是更美的.
正如他作更美之約的中寶.
這約願是憑更美之應許立的.
那前約若沒有瑕疵.
就無需尋求後約了.
所以主指責他的百姓說.
日子將到.
我要與以色列家和猶大家另立新約.
不像我拉著他們祖宗的手.
令他們出埃及的時候.
與他們所立的約.
因為他們不幸心守我的約.
我也不理他們.
這是主說的.
主又說.
那些日子以後.
我與以色列家所立的約乃是這樣.
我要將我的律法放在他們裡面.
寫在他們心上.
我要作他們的神.
他們要作我的子民.
他們不用各人教導自己的香輪.
和自己的弟兄說.
「你該認識主.

$^{41}$因為他們從最小的到最大的都別認識我」.
我要寬恕他們的不義.
不再紀念他們的罪險.
既說新約就以前約為舊了.
但那佔舊佔衰的.
就必快歸無有了.
今日我們繼續希伯來書的訊息.
題目是「超越的基督 新約的中堡」.
請陳明荃牧師.
.
我們生命因著基督耶穌的復活.
我們生命的詮釋.
對我們的過去.
對人類整個的歷史.
對每一個人生命深處裡面的遺憾.
因著基督耶穌的救贖.
我們生命可以脫胎換骨地去改變.
因著這位為我們.
死在十字架上的基督耶穌.
並在第三天復活.
我們可以看著他的生命.
我們知道我們將來的日子.
他要將屬天的榮耀.
交給每一個上帝所選照的子女和百姓.
我們今日我們一齊的來得享這份光榮.
得享這份喜樂.
從來不是人的努力可以換取回來的.
一切都是坐在高天的這位大祭司.
將這個生命充滿快樂的恩典賜給我們每一個.
所以今天我們是歡喜快樂.
我們的心裡面帶著一份興奮.
帶著一份新的詮釋來看我們自己的生命.
在希伯來書的作者同樣都是用這個角度.
來幫助一班.
以前可能經歷過很多軟弱.
或者生命裡面好像停滯不前的一班基督徒.
希伯來書的作者.
他看基督耶穌的受死復活.
他不是純粹聚焦在一個事件.
或者一個單元裡面.

$^{81}$希伯來書的作者用一個更闊的鏡頭.
由耶穌降臣肉身.
接著在地上的來到去生活.
接著在十字架上面受苦.
接著再升天.
整個的歷程是一個獻祭的歷程.
他將耶穌基督在地上的那33年的日子.
或者特別是聚焦在他受苦的地上生活的日子.
再升天.
成為一個整體.
他用耶穌基督的經歷.
來到去勉勵我們.
甚至乎希伯來書的作者再要推多一層.
他要說.
基督耶穌所做的事件是一個.
不單止是一個救贖的事件.
而是一個獻祭的.
一個大祭司所做的一個獻祭的一個舉動.
所以你會看到的.
在希伯來書裡面經常出現.
或者去到今天我們所讀的經文是一個某程度.
一個幾核心的經文.
就是講到這個大祭司.
已經升到去高天了.
這個升上高天的大祭司.
整件事怎樣去幫我們.
每一個生命可以有一些的更新和改變.
曾經有些弟兄姊妹都不約而同問.
今天我們還需要不要獻祭呢.
從希伯來書的角度來說.
耶穌已經為我們獻上最大的祭.
我們不用像舊約那些.
猶太人或者以色列人那樣做很多泛民俗節.
不過呢.
希伯來書的作者不是就這樣講了就算了.
他今天這一段經文比較仔細一點.
講的是這個大祭司耶穌怎樣來獻祭.
祂的邏輯,祂的思路是怎樣呢.
希望是今天這一段經文來幫助我們.
更加重要的有一些很重要的字眼就是中保.

$^{121}$中保呢.
我們今天理解的中保就是中間人的意思.
究竟耶穌怎樣成為一個大祭司又成為一個中間人.
來挽回.
或者使神人之間的關係進入一個新的階段呢.
希望今天這一段經文我們有些啟迪.
讓我們能夠明白上帝的心意.
我們一起祈禱嘛.
求上帝幫助我們讓我們明白祂的話語和心祈禱.
因為你的復活讓我們整個全息都可以變得不一樣.
對我們生命的看法.
對我們人生所走的路我們更加有確實的那種確具.
求你今天你藉著這一段經文.
你成為大祭司的工作.
來提醒我們.
究竟你怎樣去成為我們人與神之間的中保.
幫助我們眾人.
讓我們去領受你的話語的時候.
讓我們每一個都聽得明白.
又讓孫講的講得清楚.
忠心地將你的話語.
來去展明給弟兄姊妹知道.
讓我們明白的時候.
我們心裡面有確信.
也都有能力.
恭敬將來領受你話語的時間交託.
願你賜福.
獻上祈禱奉靠基督耶穌得勝的聖名祈求.
Amen.
在第八章一至到第五節裡面.
希伯來書的作者這樣講.
我們所講的事其中第一要緊的.
就是我們有這樣的大祭司.
已經坐在高天.
至大者寶座的右邊.
在聖所就是神真像幕裡.
作執事.
這幕是主所知的.
不是人所知的.
凡大祭司都是為獻禮物和祭物設立的.

$^{161}$所以這位大祭司也必須有所獻的.
他若在地上必不得為祭司.
因為已經有照律法獻禮物的祭司.
他們供奉的事.
本是天上的事的形狀和影像.
正如摩西章造帳幕的時候.
蒙神警戒他說.
你要謹慎.
作各樣的物件都要照著在山上指示你的樣式.
在這裡.
希伯來書的作者說.
整件事的救贖計劃.
他是用憲制的角度去表達的.
那個憲制在希伯來人或者以色列人來說.
憲制是需要指定的地方.
不是到處都可以在街角那些.
那些是犯罪的.
那些休壇就是這樣.
希伯來書的作者說.
憲制是需要有一個聖所.
有一個祭壇的.
是上帝指定的一個地上憲制的地方.
地上的憲制的那個聖所.
是照著天上的那個聖所.
是按著吩咐來設立的.
所以不論是在曠野的那四十年.
用一個回望的方式.
搭建出來的一個敬拜的地方.
還是後來去到所羅門所興建的聖殿.
整件事都是一脈相承的.
整個的設計不是純粹覺得漂亮.
我就自己做一些東西出來.
丁丁查查這樣.
我就做一個地方來敬拜神.
不是這樣的.
整個回望.
整個地上的聖所.
是照著天上的那個聖所.
是上帝作為那個設計師.
命令摩西.

$^{201}$命令後來的人.
來按著這些規格.
來建造回望或者聖所的.
所以地上的聖所.
在某程度就是天上上面那個聖所.
整個的藍圖.
或者按著天上敬拜的場所.
來建立而成的.
將來我們基督耶穌再回來.
我們都會在天上聖所裡面去敬拜的.
我們都會按著這個藍本來做的.
不過在希伯來書的作者裡面.
他講到地上的聖所和天上的聖所.
有什麼分別呢.
那個分別就是層次差很遠.
地上的聖所只不過是.
按著天上的聖所.
在這裡裡面講到第五節.
他們供奉的事是按著天上的事物的形狀和影像.
形狀和影像.
當代的用語已經用到盡了.
他的意思就是說參照上面的那個形態.
還有那個樣本.
在地上裡面設立一個這樣的方式來去敬拜.
那個形狀和影像不是那個實體.
只不過是一個人為.
按著上帝的吩咐而做的.
如果你覺得好像很抽象.
你就可以理解是一個模型.
不過這個模型是人裡面參與當中.
如果你有機會去香港的太空館.
早一段時間應該會重開.
他有一個模型就是銀河裡面八大行星.
現在是八大不是九大.
八大行星的模型地球.
金木水火土星.
做了一個很大的模型.
進去之後很偉大很漂亮.
那個是一個模型.
那個模型就按著科學家或者實證.

$^{241}$來投射銀河系.
很壯闊的空間濃縮成為一個小地方.
讓你用一個肉身可以摸得到的方式.
來接觸金木水火土那八大行星.
在地上的聖所.
類似就是一個這樣的模型.
讓人進到一個地方裡面.
模擬的或者初常經歷天上敬拜的那種形態.
將來天上的聖所比地上的宏偉得多了.
整件事就很不同了.
在希伯來書的作者裡面.
地上的敬拜地上的聖所就是一個模型.
一個有不同的人來值班.
按著不同的律法所定下的祭司的班子.
每一個人就在這裡裡面.
來帶領一班以色列民眾一起來去敬拜.
提到外話.
在敬拜裡面的寺廟不是每個人都可以進去的.
有一些在律法裡面.
今天沒有給大家圖片.
有機會我穿給大家看.
在舊時的大祭司穿的袍子裡面有件藍色的內袍.
在內袍的腳底.
即是袍子的三腳.
綁著一些鈴鐺和石頭.
連在袍子的底部.
做動作的時候就會鈴鈴鈴鈴.
為什麼要做鈴鈴鈴鈴呢?.
不是以前《異天屠龍記》裡面那些很嬌俏的女生.
打著鈴鐺告訴別人我在這裡.
那個鈴鐺其實是告訴那些在旁邊服侍的祭司知道.
如果大祭司獻祭的時候聽不到這些鈴鐺聲.
很長時間都聽不到.
證明著這個大祭司身上有罪.
獻祭的時候被上帝擊殺.
這個鈴鐺就是告訴他.
那個人這麼久沒有聲音他就死了.
他就要叫那些祭司從這個致勝所裡面把他拉出來.
雖然大祭司在聖經裡面說到是一個很尊貴的身份.
但也是一個高危的身份.

$^{281}$因為他代表整個以色列民族獻上這個祭.
如果他自己潔淨不夠足夠.
又或者身上有罪他還沒處理夠.
而代表百姓獻這個贖罪祭的時候.
他隨時就在致勝所被上帝擊殺.
所以他這個身份是尊貴.
也是一個高危的身份.
在整個的建制裡面.
按著嚴格的律法來執行.
所以在這個律法底下裡面.
每一個人的建制都是膽戰心驚.
帶著一份極度的敬虔.
回到這裡.
他說到地上的敬拜就是一個這樣的敬拜.
按著律法的要求來將上帝所喜悅的祭來獻上.
而天上的敬拜或者耶穌基督.
作為大祭司所獻的祭就有些不同.
天上的祭 耶穌的祭壇.
你會看到整個聖經裡面.
沒有說到耶穌進去致勝所裡面做獻祭的祭壇.
或者耶穌從來都沒有做一些祭司裡面殺生那樣東西.
在聖經裡面說.
當然耶穌不會做這些事.
但在經文裡面特別說.
因為耶穌是天上的唯一大祭司.
地上的祭祀儀式.
由地上的祭司來做的.
所以經文裡面特別說.
地上的帳幕的設立.
由祭司按著律法.
已經有律法來獻上禮物.
這是第四節裡面說.
在地上裡面做這個祭司.
不過耶穌都依然要獻祭.
祂要獻一個更加高層次的祭.
真正的祭壇.
或者上帝所命令耶穌所要獻上自己的祭壇在哪裡呢.
十字架.
十字架就是呈現.
將最大的生畜.

$^{321}$耶穌將自己的生命獻上的祭壇.
十字架就是耶穌獻祭的祭壇.
祂將自己一字過的獻上上帝.
所以經文裡面特別說.
大祭司一定有東西要獻上.
說的是耶穌所獻的.
不是純粹口說或者賦予祂的意思.
而是真實的獻上自己.
希伯來書的作者頭頭尾尾用很多次的經文.
除了這裡之外.
在第九章裡面說到.
祂的身份是除掉罪.
九章二十六節.
十三章十二節.
祂說到耶穌用自己的血叫百姓成聖.
獻祭是要殺生的.
殺生是說耶穌在十字架上犧牲流血.
讓百姓成聖.
然後一章第三節.
祂的血是洗淨人的罪的寶血.
所以耶穌流的血是一個獻祭流的血.
不是一個受害者.
是一個悲悲慘慘.
是一個被世人唾棄的.
一個可憐蟲流出來的.
孤孤單單很悲慘流出來的那些血.
耶穌流出來的血.
就好像昔日的祭身獻上一樣.
祂的血是有潔淨的果效.
穿著祂自己擺上的生命.
祂將自己完完全全獻上.
成為最完美的祭.
十字架就是祭壇.
獻完祭之後.
耶穌就升上高天.
在至高者的右邊.
永遠的來對作王.
那個升上去就好像舊時的祭物.
被燒上去的那些祭身.
燒上去的那些氫煙.

$^{361}$一路向天上上騰的那些煙.
這個比喻就是耶穌升上高天的比喻.
兩個都是講到達到上帝的根前.
耶穌的獻祭.
完完全全被上帝所閱納.
成為這個世界上最完美最大的祭.
這個就是耶穌裡面.
或者希伯來書裡面所講.
耶穌成為大祭司裡面所獻的祭.
這個希伯來書的作者.
再要進一步.
他要再講這個獻祭.
除了是講一個更偉大的祭之外.
耶穌基督所獻的祭.
耶穌不是成為天上委派下來.
做一個禮儀師.
做一些禮儀這麼簡單.
我們很多時候.
如果純粹用地上的祭祀來想耶穌.
我們會到到樽頸位.
就是極其量他做一個最勁的禮儀師.
天上降下做一個最勁的禮儀.
希伯來書作者說.
不是.
他不是做一個禮儀師.
他要做實質的事.
他的獻祭背後帶著一個更加具體.
更加重要的訊息.
這個是八章六節裡面所講.
如今耶穌所得的職任是更美的.
正如他作更美之藥的中補.
這藥原是憑更美的應許納藥.
在這裡你就明白了.
耶穌作為祭司的身份.
最重要是做些什麼呢.
他要做上帝和人之間的中補.
昔日祭司都是做這些事.
不過用禮儀的方式.
做這些祭.
每一次性,每年性,節期性的.

$^{401}$來做一些禮儀.
來拉近人與神之間的關係.
但拉來拉去.
人與神之間的關係都是在律法當中的規條.
不過耶穌作這個中補.
就將這個關係拉得更近.
首先解釋一下中補的意思.
中補其實就是.
不單只是說中間人.
站在人神之間的位置.
中補的意思.
是要將對立的兩方面.
拉在一起.
他要將因為犯罪的人.
和聖潔的上帝的鴻溝.
藉著基督耶穌將兩個對立面.
拉在一起.
藉著基督耶穌.
耶穌代表著人.
也代表著神.
這是地產商代理的說法.
他代表兩方面的.
這個中補不是單方面表達那份誠意.
而是代表著地上的人.
也代表著神自己親自挽回了人神之間的距離和鴻溝.
地上的人在希伯來書第五章第七至第十節.
裡面這樣說.
基督在肉體的時候既大聲哀哭流淚禱告.
懇求那救他免死的主.
就因他虔誠蒙了英文.
他雖然為兒子.
還是因所受的苦難學了順從.
他既得以完全就為凡信眾他的人.
永遠得救的根源.
並蒙神照著默基使德的等字稱他為大祭司.
這段五章第七至第十節說的是什麼呢?.
他說到耶穌成為人的樣式.
他有人的那種掙扎.
面對著十字架之前.
他也是大聲痛哭流淚禱告.

$^{441}$耶穌面對著苦難.
他也求上帝去免去這個苦悲.
他也是戴著人性.
戴著人的掙扎面對苦難的時候.
那種痛苦憂愁.
在聖經裡面.
希伯來書作者就用這段經文.
用耶穌在赫西瑪利園的祈禱.
來描述耶穌的人性.
耶穌是百分之一百的人.
他是無罪的神兒子.
他是百分之一百的人.
不過面對著十字架苦難.
他也好像人的軟弱一樣.
他也會有痛苦掙扎.
面對著死亡的威脅.
他也有一種很難完全表達出來的那種沉鬱.
希伯來書作者說.
他在這些沉鬱.
在這些痛苦裡面.
完全將人性表露無遺.
不過在他的生命當中.
他在這些苦難當中裡面.
他學懂或者他最終選擇信服.
用信服的方式來實踐.
他成為人世間所有人的代表.
他是代表著一個最完美的人.
面對著掙扎.
他也依然的來到他走上十字架上.
所以耶穌作為一個中保是代表人的.
他用最完美的人的方式來代表.
上帝的計劃不是找一個最懶的人.
不是找一個人渣來代表自己.
而是找一個最人裡面沒有人可以及得上.
最頂級的一個生命的人.
來代表這個世界所有的人走上十字架.
他的受苦就是將人的所有罪歸到他身上.
所以耶穌是成為人.
他的中保一方面代表人.
同時代表神我們也說很多.

$^{481}$或者他自己就是神.
他就是基督就是彌賽亞.
他是神自己.
但同時他是三位一體裡面.
神兒子第二個位格.
他用這個身份.
用神的生命.
來將上帝和罪人之間的鴻溝.
藉著他的憲制來挽回回來.
早陣子有時候看過一本書.
他說這個世界裡面有很多中保.
他用的字是中間人.
就是說這個世界你不要想著隨著科技發展.
中間人會越來越少.
那本書叫中間人經濟.
這個世界的中間人會越來越多.
不過有不同形態的中間人.
他其中說的一件事我覺得很像耶穌基督的那種表達.
他就是在說隔絕型的中間人.
說什麼呢?.
隔絕型的中間人就是在說.
在很多對立面裡面.
這個人扮演的角色.
他不是要強調自己多威風或多厲害.
而是用他的努力.
來將一些很多對立紛爭的層面.
盡量來將這些人放在一起.
他說這個世界有很多中間人.
用一些不道德的方法.
例如其中一個方法就是將兩方面.
撩起那些矛盾和對立.
撩起那些矛盾和對立之後.
這個中間人感覺上就很有價值.
為他才可以傳話.
不過面對著隔絕型的中間人.
他說不是的.
他的角色是要將兩個對立面.
一路一路拉近.
讓這些對立減到最低.
讓那份和諧的關係可以重新來展現.

$^{521}$耶穌基督就是在隔絕.
人和神之間的隔絕當中成為中堡.
將犯罪的人可以走近在聖潔的上帝當中.
積著他的見濟.
令到人可以走到上帝的根前.
耶穌基督是神與人之間的橋樑.
如果沒有了他.
我們只是像舊約那樣用叮叮叉叉.
用一些見濟.
用禮儀想著靠近上帝一點.
但越做這些事其實可以只有外貌.
但心裡是可以剛硬的.
在聖經裡邊.
接下來的篇幅說到.
猶太人或者以色列人.
你做那麼多事都好.
他們的心都是剛硬.
恆心地不聽從上帝.
有這些禮儀又如何呢.
所有一切都是外在的.
最重要的就是.
心靈與上帝彼此建立.
上帝的心意就是要.
將人與神之間重建一份和諧.
將那份對立.
將人犯罪的光景徹底扭轉.
人可以走到上帝的根前.
所以經文裡再說.
這個中保目的是要做什麼呢.
就是要實踐這個更美的約.
更美的約一講的時候.
我們覺得.
一給字眼看就覺得很美.
或者更加美妙.
中文是這樣翻譯的.
用這個角度來看.
當然是在說神與人之間.
大家走近一步.
越來越有進步的.
最美的約通常都是在說進展.

$^{561}$大家在進步下去.
大家越走越近.
不過在聖經裡邊鋪陳的圖畫不是這樣.
更美的約.
不是因為上帝更加愛人.
人就更加愛神.
朝著一個這樣的交匯點來去.
是上帝一直都是這樣.
越來越愛人.
只不過人就越來越走遠.
人越來越走下坡.
這個更美的約.
基於的就是.
人越來越失敗.
人性越來越卑污.
人性越來越差的那種心態.
所以在八章九節裡面這樣說.
他說.
不像我拉著他們祖宗的手.
令他們出埃及的時候.
與他們所納的約.
因為他們不恆心守我的約.
我也不理他們.
這個更美的約.
本身為什麼會被建立起來呢.
原來是在說神人之間.
那個距離越來越遠.
因著基督耶穌.
要將這班人拉在一起.
那些百姓不恆心守上帝的約.
在保羅的用語是說.
這班百姓是心硬.
心裡剛硬.
在舊約裡面說到這班以色列民.
離棄上帝的那種心態.
就好像一個石頭造的心一樣.
越來越偏行幾路.
上帝對於一些石心.
對於一些偏行幾路.
心裡剛硬的人有什麼反應呢.

$^{601}$上帝不理他們.
保羅說任憑他們.
任由他們犯罪.
在世界裡面你會看到.
在羅馬書裡面說到.
這個世界很多人.
用不同形式建立他們的偶像.
任意而行.
心裡面做自己想做的事.
我就是最大的神.
以至到生命裡面.
處於一種和上帝敵對的關係.
上帝給他的懲罰.
最大的懲罰就是任憑他.
任由他.
上帝就看你.
繼續怎樣走向滅亡之路.
不過神人之間.
這不是上帝創造的心意.
如果面對這個困局.
唯有基督耶穌成為中寶.
將這兩個方面.
將失敗的人.
和將任憑人犯罪的上帝.
重新拉在一起.
用憲制的方式.
將兩下重新合而為一.
更美的約是說到.
上帝願意用更大的恩慈.
藉著耶穌基督神的兒子.
來挽回神人之間的關係.
我們再繼續看.
第1至13節.
引用舊約.
具體如何建立新的約.
如何重新建立新的關係.
那前約若沒有瑕疵.
就沒處尋求後約了.
所以主指責他的百姓說.
日子將到我要與以色列家和猶太家.

$^{641}$另立新約.
不像我拉著他們祖宗的手.
領他們出埃及的時候.
與他們所立的約.
因為他們不恆心守我的約.
也不理他們.
我也不理他們.
這是主說的.
主又說.
那些日子以後.
我與以色列家所立的約乃是這樣.
我要將我的律法放在他們裡面.
寫在他們心上.
我要作他們的神.
他們要作我的子民.
他們不用各人教導自己的香輪.
和自己的弟兄.
說:你該認識主.
他們從小到大都必認識我.
我要寬恕他們的不義.
不再紀念他們的罪險.
既說新約就以前約為舊了.
但那佔舊佔衰的.
就必快歸無有了.
七節至第十三節說的是什麼呢?.
說的是神人之間的舊關係.
舊約的關係.
舊約的關係是一個律法的關係.
上帝頒佈了律法.
叫以色列人遵守.
用割禮成為一個記號.
上帝是藉著律法.
頒佈給一群以色列人.
這群子民透過守律法.
來彰顯與這個世界的不同.
分別為聖.
這是舊約的精神.
不過因為這群人.
漸離漸遠.
離開舊約律法.

$^{681}$離開這個律法.
所以這個律法漸漸就好像失效一樣.
上帝說這個舊約.
有瑕疵的.
不是最完美的那種神人之間的方式.
因為律法的方式就是由外而內.
來叫人.
來與上帝建立關係.
上帝說.
去到耶利米書.
耶利米書裡面.
就將這一段經文.
其實希伯來書的作者.
將耶利米書的整段經文.
31章31至34節.
抄下來.
在舊約聖經裡面講新約.
為什麼呢?就是說這群人都失敗了.
這群人都已經漸離漸遠.
整個以色列國都亡國了.
不過亡國之際.
上帝直著先知就告訴他們.
上帝有一個新的關係要和他們講.
他重新.
立一個新約給他們.
這個新的約就是在說.
上帝要將律法.
不是寫在法版裡面.
而是由內而外.
來建立他們之間的關係.
不過在解釋這個新約之前.
我很想藉這個機會.
釐清一下.
概念上.
新約和新約聖經的分別.
在這裡面講的.
那個約是一個關係的約.
在耶利米書就講到.
上帝要和以色列人.
整猶太人和全世界.

$^{721}$建立一個新的關係.
一種相處模式.
這個就是那個新約.
那個預言.
這個預言.
這個約什麼時候開始實現呢?.
就是在耶穌基督降生開始.
這個約就一直一直實現.
直至十字架成就了.
耶穌升天了.
整個救恩完備了.
這個約就正式生效.
所以這個約的定定.
和真正實現的約.
是有一段時間的區間的.
但是在舊約聖經裡面.
就講到這個新的關係.
就已經講明了.
在耶利米書裡面.
在耶穌還沒有出生的時候.
這個新約.
這個新的關係.
上帝已經藉著先知.
來表達出來.
新約聖經.
就是將耶穌基督的救贖.
耶穌怎樣去實踐這個新約.
很多的見證人和耶穌的事蹟.
甚至乎耶穌之後.
升天之後的教會怎樣建立.
那個記載下來的文本.
和成為一個權威.
就成為我們的新約聖經.
所以這個關係.
是講著.
藉著這個聖經.
來讓我們明白,展明出來.
所以你說兩者是否完全一致呢?.
也不是的.
不過新約聖經就記載著.

$^{761}$一個新約的關係.
因為耶穌基督的實踐.
所以這個新約就有效了.
但是在舊約裡面.
已經在講這個新約的關係.
所以希望弟兄姊妹不要混亂了.
新約聖經和新約.
新約是在講新的關係.
而新約聖經.
就是在講我們今天.
由耶穌出生開始記載.
和之後直到啟示錄.
就是後結的.
這二十七卷的經文.
就成為新約聖經.
我們今天讀聖經.
我們作為一個基督徒.
我們怎樣讀聖經呢?.
我們由新約聖經.
由基督耶穌的復活.
來理解舊約.
我們不像猶太人.
他們用律法來理解上帝的工作.
今天我們不是猶太人.
我們理解信仰.
我們是由基督耶穌的生命.
耶穌的救贖開始的.
甚至乎.
整本聖經的全釋.
我們怎樣看上帝的計劃.
我們都是由新約聖經.
來成為我們全釋的框架.
我們稱自己為基督徒.
我們是跟隨耶穌的人.
我們是用跟隨耶穌的眼光.
來看舊約聖經.
所以在舊約裡面.
我們會看到耶穌基督的影子.
我們會用新約來解釋舊約.
原因就是.

$^{801}$因為我們每一個都是在基督裡面得救的一班人.
我們是新舊約.
這樣來讀聖經.
全釋聖經.
希伯來書的作者.
也想將這件事告訴大家.
我們今天領受了一個新的關係.
從一個基督耶穌獻祭的關係.
我們成為神的子民.
來領受這個新的約.
所以這個新約有三個很大的重點.
這三個重點.
你記住這三個英文字母.
HIM.
HIM代表什麼呢.
幫大家記住.
第一個.
H代表著HEART.
Tablet of HEART.
心板.
Tablet不是你手上的iPad.
那叫做Tablet.
在聖經裡面說的那個法板叫做Tablet.
不過舊時的Tablet就寫在石頭上.
律法的那個法板就寫在石頭那裡.
如今上帝要將他的律法寫在什麼.
那個Tablet就寫在什麼呢.
不是你的iPad.
而是你的內心裡面的Tablet.
你的心板上面.
律法.
寫在你的心靈裡面.
這個新關係就說到.
當你信耶穌基督.
以基督耶穌成為大祭司.
你的心靈裡面就會領受上帝的真理.
很有趣的.
一個剛剛一信主的人.
生命就馬上可以改變的.
為什麼會有這樣的神蹟呢.

$^{841}$我認識有一個弟兄.
其實那個是我外公.
我外公吃了煙.
他以前很頑固的.
在他年老的日子要信主.
他說他病危要信主.
然後他一信主的時候.
他說舊的東西全部丟掉.
包括這些壞習慣.
吃了六十年吸煙.
說戒就戒.
我說用什麼方法戒.
求上帝.
上帝叫他這樣做.
他就可以脫胎換骨.
這樣去改變.
當你信耶穌的時候.
心靈就有一個.
好像法板刻在你的內心裡面.
所以希伯來書的作者.
在說那個新的.
或者耶利米書的引用.
就是我要將我的律法放在他們的裡面.
寫在他們的心上.
我要作他們的神.
他們要作我的子民.
這個和舊約裡面很不同.
舊約是由外而內的.
現在這個新約就是.
由裡面.
由心開始發出.
上帝的華語.
上帝的真理.
上帝和人之間的關係.
在裡面來表達出來.
我們印證我們自己是基督徒.
我們印證自己是否屬上帝的子民.
你檢視一下你的內心.
心裡面有沒有上帝的華語.
心裡面會不會為罪為義為審判.

$^{881}$自己責備自己.
心靈裡面會不會有聲音提醒你.
雖然很多時候說聖靈的聲音.
會為我們帶來平安.
這個是對的.
不過聖靈的聲音.
在你犯罪的時候會為你帶來不平安.
當你犯罪或者離棄神.
你的心靈裡面.
你的敏銳度會變得越來越強.
如果你有這種印證.
我恭喜你.
你是一個真誠的基督徒.
你成為神的子民.
你以神成為你生命裡面的神.
律法刻在心靈裡面.
Heart.
這個是第一個新約裡面的表徵.
第二個.
承接著上一個的.
第十一節.
你們不用各人教導自己的香輪和自己的弟兄說.
你該認識主.
因為他們從小到大都必認識我.
認識這個字很重要.
舊約裡面認識神的.
為什麼呢.
透過中間媒介.
透過大祭司禮儀和教導.
透過先知的責備.
用教導的方式.
來讓人認識神.
不過在新約裡面.
上帝內在人心靈當中.
上帝會教導你.
那個關係不是遙遠的.
用zoom來教導你.
現在是親身跟你見面.
比面對面更近.
因為他進入到你的心裡面.

$^{921}$指教你.
聖靈會內在你的內心裡面.
所以第二個字是inspired(啟發).
by Holy Spirit.
inspire.
上帝會inspire我們.
會illuminate(亮出).
光照我們.
生命裡面有上帝直接跟我們說話.
聖靈會開導我們的心靈.
讓我們能夠親身認識上帝.
更加重要的是.
舊約上帝的開導是選擇性的群體.
是以色列民.
是那些猶太人才有機會學到上帝的華語.
偶爾有一些外邦人想學.
或者想經歷上帝的律法.
他們都可以的.
在聖殿裡面有一個叫外邦人院.
在外邦人跟以色列人親疏有別.
他們就在外邦人的外圍裡面.
來沾染學習.
他們不單是這樣學習.
他們也要進入那個文化.
背後的律法.
有很多外在的東西.
你要跟著做整套律法的package(組合).
那你就學到東西了.
不過去到新約.
他在說上帝內在已經不再驅離.
用那些外在的條文.
文化或者宗族.
已經不是成為一個限制.
上帝在各城各邦.
全世界所有人.
上帝的靈都可以inspire(啟發)我們每一個的內心.
我們每一個人都可以直接領受上帝的恩典.
我們再不需要額外.
還要添加一些其他的agent(代理人).
或者其他的中間.

$^{961}$裡面再多一個中間人來幫我們.
我們直接了當地領受上帝對我們的教導和光照.
所以這個約第二個重點就是inspired by God.
inspired by Holy Spirit.
這個I就是上帝直接教導我們.
第三個八章十二節.
我要寬恕他們的不義.
不再紀念他們的罪險.
八章十二節說到.
上帝和人之間的關係.
再不是一個對立.
不是一個純粹看能不能做到行為.
那個道德水平來和人建立關係.
上帝將他的本性.
將他的本質.
愛的本質呈現給這個世界.
他要寬恕.
他要憐憫地上的罪人.
memorize their sin no more.
不再紀念他們生命裡面的罪.
這個不再紀念不是廉價的.
不再紀念的原因就是因為基督做了最大的獻祭.
做了最大的功夫將人和神拉在一起.
所以我們生命裡面不再受罪所束縛.
我們上帝就不再紀念我們所犯過的罪.
這段經文很簡單.
解釋就很複雜.
但帶出來的訊息很直接了當.
耶穌為我們成為挽回制.
成就了最大的祭司的職份.
耶穌藉著他中保的生命.
具體實踐了這個新約.
這個新約就是記住.
Him心靈裡面heart這個心的律法.
他直接開導我們進入我們的內心對我們說話.
他赦免我們的罪.
memorize our sin no more.
不再紀念我們的罪犬.
讓我們生命裡面重新有力量奔走面前的路.
希伯來書的作者他要說.

$^{1001}$耶穌的中保有什麼重要呢.
其中一個其實是.
他講到一些軟弱無力的生命的弟兄姊妹.
可能受著過去的遺憾.
受著生命裡面有很多難阻.
可能我們會放棄.
我們看這個世界.
看自己的軟弱看得很大.
以致到我失去這個救恩.
希伯來書的作者面臨我們.
不要再看自己這些的過去.
不要再為自己生命裡面那些千瘡百孔的事.
再在這個傷口裡面.
再去緬懷這個傷口有多痛.
又或者覺得自己多不濟.
希伯來書的作者想透過耶穌成為大祭司.
他要叫我們勇敢起來.
讓我們將發酸傷腿垂下來的手.
我們可以重新振作.
挺起你的腰骨.
仰望這個創始成中的基督.
體會這個大祭司的心腸.
將自己繼續實踐在這個信仰裡面.
復活的基督告訴我們.
我們的生命不是一個軟弱無力的生命.
復活或者整個建制.
基督耶穌的大祭司的建制告訴我們.
上帝不惜任何代價.
祂都要救贖我們.
祂體重我們.
雖然我們很不濟.
但是上帝卻是用這個重價挽回救贖我們.
就讓我們從這個復活節開始.
我們挺起腰骨.
我們仰望這個基督.
一起奔走這條上帝為我們所預備的天路歷程.
願主繼續來幫助.
\newpage



\section{希伯來書 9:1-14}
\label{sec:rKFXkQ539_w}
\textbf{2022年4月24日崇拜直播｜陳冠成牧師｜超越的基督:神的帳幕在人間｜希伯來書9\_1-14}
\newline
\newline
連結: \href{https://youtube.com/watch?v=rKFXkQ539_w}{\texttt{https://youtube.com/watch?v=rKFXkQ539\_w}} ~~~~ 語音日期: 2022-04-25
\newline
\newline
\hyperref[sec:U9MWFdIg8Fc]{\small{< < < PREV SERMON < < <}}
~
\hyperref[sec:index_chronic]{\small{[返順時目]}}
~
\hyperref[sec:vxnWocuBbhg]{\small{> > > NEXT SERMON > > >}}
\newline
\newline
希伯來書 9:1-14
\newline
\begin{longtable}{cl}
\hline
\hline
章節 & 經文 (和合本修訂版)\\
\hline
9:1 & \begin{tabularx}{0.7\textwidth}{X} 原來連第一個約都有敬拜的禮儀和屬世界的聖幕。 \end{tabularx} \\ \\ \relax
9:2 & \begin{tabularx}{0.7\textwidth}{X} 因為那預備好了的帳幕，第一層叫聖所，裡面有燈臺、供桌和供餅。 \end{tabularx} \\ \\ \relax
9:3 & \begin{tabularx}{0.7\textwidth}{X} 第二層幔子後又有一層帳幕，叫至聖所， \end{tabularx} \\ \\ \relax
9:4 & \begin{tabularx}{0.7\textwidth}{X} 有金香壇和四周包金的約櫃，櫃裡有盛嗎哪的金罐、亞倫那根發過芽的杖和兩塊約版； \end{tabularx} \\ \\ \relax
9:5 & \begin{tabularx}{0.7\textwidth}{X} 櫃上面有榮耀的基路伯罩著施恩座。有關這一切我現在不能一一細說。 \end{tabularx} \\ \\ \relax
9:6 & \begin{tabularx}{0.7\textwidth}{X} 這些物件既如此預備齊了，眾祭司就不斷地進第一層帳幕行拜神的禮。 \end{tabularx} \\ \\ \relax
9:7 & \begin{tabularx}{0.7\textwidth}{X} 至於第二層帳幕，惟有大祭司一年一次獨自進去，沒有一次不帶著血，為自己獻上，也為百姓無意所犯的過錯獻上。 \end{tabularx} \\ \\ \relax
9:8 & \begin{tabularx}{0.7\textwidth}{X} 聖靈藉此指明，第一層帳幕仍存在的時候，進入至聖所的路還沒有顯示。 \end{tabularx} \\ \\ \relax
9:9 & \begin{tabularx}{0.7\textwidth}{X} 那第一層帳幕是現今時代的一個預表，表示所獻的禮物和祭物都不能使敬拜的人在良心上得以完全。 \end{tabularx} \\ \\ \relax
9:10 & \begin{tabularx}{0.7\textwidth}{X} 這些事只不過是有關飲食和各種潔淨的規矩，是屬肉體的條例，它的功效是直到新次序的時期來到為止。 \end{tabularx} \\ \\ \relax
9:11 & \begin{tabularx}{0.7\textwidth}{X} 但現在基督已經來到，作了已實現的美事的大祭司，經過那更大更全備的帳幕，不是人手所造，也不是屬於這世界的； \end{tabularx} \\ \\ \relax
9:12 & \begin{tabularx}{0.7\textwidth}{X} 他不用山羊和牛犢的血，而是用自己的血，只一次進入至聖所就獲得了永遠的贖罪。 \end{tabularx} \\ \\ \relax
9:13 & \begin{tabularx}{0.7\textwidth}{X} 若山羊和公牛的血，以及母牛犢的灰，灑在不潔的人身上，尚且使人成聖，身體潔淨， \end{tabularx} \\ \\ \relax
9:14 & \begin{tabularx}{0.7\textwidth}{X} 何況基督的血，他藉著永遠的靈把自己無瑕疵地獻給神，更能洗淨我們的良心，除去致死的行為，好事奉那位永生的神。 \end{tabularx} \\ \\
[1ex]
\hline
\hline
\end{longtable}
$^{1}$今日經訓在希伯來書第九章一至十四節.
是在新約聖經希伯來書第九章一至十四節.
請聽我讀出.
原來前曰有禮拜的條例和屬世界的聖幕.
因為有預備的帳幕.
頭一層叫作聖所.
裡面有燈台,桌子和陳設餅.
第二晚之後又有一層帳幕.
叫作至聖所.
有金香爐.
有包金的藥櫃.
櫃裡有聖嗎哪的金冠.
和亞倫法幹娜的像.
並兩塊藥板.
櫃上面有榮耀基督帕的影.
罩著施恩座.
這幾件我現在不能一一細說.
這些物件既如此預備齊了.
眾祭司就常進投一層帳幕.
行拜神的禮.
至於第二層帳幕.
唯有大祭司一年一次獨自進去.
沒有不帶著血為自己和百姓的過錯獻上.
聖靈如此指明.
投一層帳幕.
迎傳的時候.
進入至聖所的路還未顯明.
那投一層帳幕.
作現今的一個表樣.
所獻的禮物和祭物.
就著良心說.
都不能叫禮拜的人得以圓全.
這些是連那飲食.
和珠缽洗作的規矩.
都不過是屬肉體的條例.
命定到振興的時候為止.
但現在基督已經來到.
作了將來美事的寨祭司.
經過那更大更傳遍的帳幕.
不是人手所做.

$^{41}$也不是屬乎這世界的.
並且不用山羊和牛毒的血.
乃用自己的血.
只一次進入聖所.
成了永遠贖罪的事.
若三羊和公牛的血.
並沒有牛毒的灰.
灑在不潔的人身上.
尚且叫人乘勝身體潔淨.
何況基督藉著永遠的靈.
將自己無憾無恥獻給神.
他的血豈不更能洗淨你們的心.
除去你們的死痕.
使你們事奉那永生神嗎?.
(記者拋出墊子).
各位姐妹平安.
(各位姐妹平安).
我們有多久沒有實體面對面.
一齊敬拜上帝了.
已經超過三個月了.
所以我們很感恩今天可以.
回到上帝的家.
回到教會.
一齊去敬拜.
事奉上帝.
一齊去聆聽上帝的話語.
今天我們會繼續分享希伯來書.
正如剛才《樹根》執事.
他和我們所讀的經文.
《不經不覺》已經來到第九章.
你留意到.
上星期名存牧師.
他和我們分享到第八章.
作者說到.
耶穌基督透過.
他在天上聖所的大祭司的服侍.
成為了神和全人類.
更美的宗寶.
耶穌基督要和我們另立新的約.
將上帝的律法寫在人的心裡面.

$^{81}$並且透過聖靈的引導.
使人可以明白真理.
認識真神.
現在我們來到第九章.
其實整個第九章到十章的十八節.
在這段內.
你會留意作者進一步詳細講解.
到底基督是如何成為完美的贖罪祭.
以及耶穌基督所獻的贖罪祭.
能夠比舊約的大祭司.
所獻的更為完美.
作者準備和讀者分享這件事.
不過留意到作者在交代這一切之前.
他又先講一講.
舊約會幕的結構.
以及當中一些相關的條例.
大家可以打開程序表.
或者在廣道大綱裡面已經印了.
你會看到作者在第二至第五節.
如果用一個圖畫來解釋.
就好像印在程序表裡面的會幕結構.
主要分了三個範圍.
有外緣,有聖所和致聖所.
第二至第五節就特別描述了.
會幕裡面的兩間房.
聖所和致聖所的擺設.
有經文裡面提及的金燈台.
陳設餅的桌子.
有金香爐或者所謂的金香壇.
還有藥櫃.
其實這一切擺設.
裡面蘊含了很多屬靈的意義.
譬如金燈台的光.
象徵著上帝的真光,真理.
又或者陳設餅的餅.
象徵著上帝的公認和保守.
不過作者很快就跟我們說.
他沒有意跟我們分享這方面的教導.
他想我們集中看兩間房.
一間叫聖所,一間叫致聖所.

$^{121}$集中看這兩個的空間.
和當中一些相關的敬拜條例.
首先在第六節裡說了.
聖所是一般祭司.
他們事奉上帝活動的範圍.
他們可以進去替自己和替百姓.
執行敬拜的禮儀.
作者說得很清楚.
至於致聖所.
只容許大祭司一個人進去.
並且他不可以隨便進去.
他要按照舊約律法的規定.
他每年有多少次機會進去?.
只有一次.
經文說得很清楚.
他只限於每年一次贖罪日.
帶著祭身的血.
為自己和百姓會中的罪.
獻上贖罪祭.
經文說得很清楚.
作者亦在聖靈的引導下.
開始講解.
其實頭一個帳幕.
他所說的是第一間房.
即是聖所.
其實聖所的存在.
就好像一個障礙物.
他很具體地說.
其實聖所固然有它敬拜的作用.
但你反過來看.
它是一個障礙物.
聖所一日的存在.
即是反映什麼?.
你便看不清楚怎樣走進至聖所.
遮住了.
令到進入至聖所的路.
是未顯明的.
換句話來說.
百姓根本沒有辦法.
直接和上帝雙交.

$^{161}$不過我們其實明白.
真正難阻人走近上帝的.
其實不是聖所這間房.
其實是什麼?.
什麼難阻我們親近上帝?.
我們今天知道答案了.
是人的罪.
人的罪和上帝的聖結為敵.
構成了我們和上帝雙交的阻隔.
換句話來說.
作者他想我們明白.
舊約敬拜系統裡面.
聖所的存在其實就好像.
象徵著百姓裡面的屬靈狀況.
其實根本未完全得到潔淨.
百姓雖然有獻祭.
不過他們所獻的祭身.
並不能夠使他們的良心.
完全潔淨.
作者自己也說得很清楚.
在第十節.
你會看到作者進一步解釋.
他說不單單指這方面.
他說連同舊約各樣的飲食.
洗手,潔淨條例.
有其效用.
不過只能提供禮儀上的潔淨.
不能夠將人和上帝的關係.
徹底恢復正常.
這其實我們也很明白.
這種狀況要等到什麼時候.
才可以徹底更新?.
就是等到耶穌基督來.
等到耶穌基督來到.
他獻上完美的贖罪祭的時候.
這種狀況就可以徹底得到更新.
來到這裡我們不難理解.
耶穌基督做了什麼?.
他在十字架上成就救恩的時候.
聖殿發生了什麼事?.

$^{201}$聖殿至聖所對出的萬子裂開了.
從上而下裂開兩半.
象徵著人和上帝的阻隔已經怎樣?.
透過耶穌基督的救贖來徹底打破了.
藉著十字架上.
耶穌基督的寫生流血.
他不單止滿足了.
律法裡面所說.
若不流血罪就不得赦免.
這個要求.
同時他也表明了.
聽命勝於獻祭的屬靈真理.
你看到耶穌基督對上帝的信服.
他犧牲自己的性命.
他親自成為了救約贖罪日的大祭司.
和祭僧.
成為了罪人和上帝之間的忠報.
徹底取代了救約大祭司.
和祭司的角色和功能.
並且比救約的獻祭更加超越.
怎樣超越?.
首先我們會明白.
牛羊在不知情的情況下.
在他非自願的情況下被宰殺.
在他不知情的情況下.
替罪人來受死.
但基督耶穌.
他很清楚上帝的心意.
他知道在父上帝整個的救贖計劃裡面.
他要為人來預備贖罪祭.
他要成為祭僧.
所以他亦甘願去信服.
為罪人犧牲.
並且我們知道.
耶穌基督為何是完美無瑕的祭僧.
他有沒有得罪過上帝?.
他沒有.
他從未背棄過上帝.
所以他是完美無瑕的祭僧.
很自然沒有罪的主才能真正有資格.

$^{241}$替我們罪人來頂罪.
而另一方面.
耶穌基督為何能超越舊約的大祭司.
是因為他不像亞倫體系的祭司.
他會被生死.
或者罪.
又或者病患.
來影響甚至間斷了他中保的角色.
你知道舊約的大祭司.
他不會廿四小時後命.
沒有可能.
但耶穌基督他已經從死裡復活.
戰勝了死亡.
並且他已經幫我們徹底擺脫了.
罪在我們身上的權勢.
他現在已經升入高天.
坐在父上帝的右邊.
永遠執掌王權.
也意味著這位天上大祭司.
他所獻的贖罪祭.
是永遠有效.
不受時空,不受環境所限制.
我們隨時可以透過耶穌基督.
來到上帝的面前.
並且我們也明白.
這位大祭司.
他也可以體恤人的軟弱.
不單止體恤.
他也自己嘗試過各種試探.
他也勝過並且有能力幫我們勝過試探.
所以為何.
來到這裡希伯來書的作者.
他說到.
耶穌基督為何可以用自己的血.
只需要一次進入至聖所.
就成為了永遠贖罪的祭.
一次過為我們所有的罪.
包括你和我將來所犯的罪.
如仙一筆過付上永遠的贖價.
使每一個.

$^{281}$今天我們很明白.
凡信耶穌的人.
凡認耶穌為主的人.
都能夠徹底和上帝和好.
不再需要走舊約的路.
要反覆.
透過祭司.
透過祭僧.
透過獻祭.
才可以接近上帝.
耶穌基督他親自成為了.
更大更完備的天上的會幕.
我們可以通過祂.
直接來到上帝的面前.
敬拜和事奉祂.
以上就是整個希伯來書第九章.
至第十章裡所交代.
為何耶穌基督能夠成為更美的祭.
你會看到.
上帝他一直很致力地.
提供方法,提供路徑.
讓罪人可以走近上帝.
和祂相交.
我們很明白.
特別在新約裡.
不單止希伯來書.
新約的頭一卷書《馬太福音》.
作者已經講明.
耶穌基督來到是怎樣.
記得我們剛才唱的詩歌.
「以馬來利」嗎?.
記得吧?.
「以馬來利」代表甚麼?.
上帝說祂要與人同在.
耶穌基督第一次來到時.
已經很清楚講明.
祂要與人同住.
要住在人的中間.
並且祂不是單說.
祂用祂的生命.

$^{321}$特別用祂十字架的生命.
清楚表明祂是多麼渴望.
與罪人同住.
不惜犧牲自己的性命.
祂講得很清楚.
又或者去到啟示錄.
新約最後一卷書.
主題是同樣的.
當耶穌基督第二次來的時候.
他亦是這樣說.
印在程序表裡.
啟示錄第21章第三節.
「看!神的障礙在人間.
祂要與人同住.
他們要作祂的子民.
神要親自與他們同在.
作他們的神」.
這份心意在整個新約裡.
都是很一致的.
上帝要住在人的中間.
無論祂第一次來.
或者祂將來再來的時候.
都是這份心意.
不過我們要留意.
舊約是否這樣?.
是否只有新約才有這份心意?.
不是的.
你看回上帝與人同住這份心意.
一早在舊約已經出現.
在《出海及記》裡.
講到當上帝與子民立約之後.
上帝很快就說.
「我要與你們同住」.
而方式就是透過會幕的設立.
來顯明祂要與人建立一份.
很緊密的關係.
我們看另一張圖.
如果我們只看會幕結構.
沒有那麼明白上帝的心意.
你看以色列人在曠野.

$^{361}$如何紮營,如何安營這幅圖.
它就更加具體反映.
上帝如何住在人的中間.
你看到這幅圖.
十二個支派圍著上帝紮營.
他們在曠野.
再內圍一點就是利未人.
上帝揀選分別為聖.
在會幕中服侍他們的人.
圍繞著會幕的四周.
而會幕就在所有十二個支派的營中間.
在正中間.
你現在就明白.
甚麼是住在百姓的中間.
真的很具體.
並且這不是人主動提出的.
如果你有留意.
昌埃及記就是.
人是怕上帝走近.
他們只叫摩西去跟上帝說.
我們怕上帝,我們留在山下.
你會發現不是人主動想走近上帝.
反而是上帝主動說.
我要跟你們親近.
於是他就吩咐摩西.
在立約之後就叫回眾.
為上帝建造會幕.
並且不是按照人的設計.
而是按照上帝的心意建造會幕.
當然來到今天希伯來書.
我們很明白.
在地上的會幕.
其實就是反映在天上的影像.
所以你看到.
上帝很實際就是.
跟一群已經被拯救.
昌埃及和他們立約的子民.
怎樣才可以更加認識上帝呢?.
就是同住.
意思是甚麼呢?.

$^{401}$可能大家對我的認識.
主要在教會,在講壇.
你見到我了.
現在戴了口罩,你可能不認識我了.
特別你是新朋友.
你不知道我實際的模樣是怎樣.
那你要怎樣才進一步認識我呢?.
跟我多聊天.
吃飯.
或者一起事奉.
但都是僅限於.
我在教會.
可能在街外.
你認識不到我在家裡發生甚麼事.
你要怎樣才認識我呢?.
就是說.
不如來我家住幾天吧.
來我家住幾天的時候.
如果你願意.
你會認識到我甚麼呢?.
你起碼知道.
我在家有沒有穿鞋.
有沒有穿拖鞋.
或者我要不要求大家穿拖鞋.
你參觀完我的洗手間.
你就會知道我會不會不小心.
或者說.
不小心路過我的房間.
你就知道我有沒有收拾床.
進入廚房.
你就會看到我平時喜歡煮食.
所以如果.
跟一個人同住的時候.
你就更加近距離.
去認識到他的喜好是甚麼.
他的生活習慣是甚麼.
並且.
他家的規矩是怎樣.
所以你看到.
整個會幕的設計.

$^{441}$不是純粹一個禮儀.
是說一群.
剛剛跟上帝立約的子民.
他們是初哥.
他們透過神蹟認識上帝.
透過紅海.
過紅海認識上帝.
透過十災認識上帝.
透過上帝可能在曠野裡.
的供應認識上帝.
但都不足夠.
跟上帝一起住.
住在上帝的房間裡.
來認識上帝.
當然上帝是靈.
神是靈.
其實祂不需要一間房子.
但為了要他要輔奏人.
讓他們明白.
透過會幕這個很實際的空間.
告訴人兩件事.
一方面祂很想人親近祂.
跟祂同住.
另一方面也透過會幕.
敬拜的條例告訴人.
你進到這間房子.
或者說進房前.
要守甚麼規矩.
裡面反映上帝的喜好.
祂對人的期望.
所以你看到這幅圖.
上帝吩咐了十二個支派.
包括利未人.
要圍繞著會幕安寧.
其實本身已經很實際.
上帝用行動.
用一件實物.
來告訴人.
祂不單單只要與人同住.
並且要住在中間.

$^{481}$祂要成為百姓生活的中心.
而神吩咐祭司和大祭司.
可以代表會眾.
從會幕外走入聖所和至聖所.
其實這個舉動本身反映甚麼?.
神是願意人親近祂.
並且上帝亦提供方法.
就是你怎樣可以親近祂.
當然我們可能會問一個問題.
上帝既然主動說.
要住在百姓中間.
為何又要設定這麼多規矩.
要求.
甚至限制著以色列人.
在會幕的活動範圍?.
你知道不是所有以色列人.
都可以進入聖所至聖所.
一般百姓.
即使想限制也好.
他最多只可以在會幕門口範圍.
他不能進入.
唯有祭司和大祭司才能進入.
所以用我們今天.
現代人的角度.
我們經常會問.
為何不開放整個聖所至聖所.
讓所有百姓都可以進入.
讓他們隨時可以自由親近敬拜上帝.
這就符合我們今天.
自由敬拜的要求.
我們要問這個問題.
是不是上帝好像很矛盾?.
又想人親近祂.
但又限制人.
在聖所裡的活動範圍.
好像很矛盾.
不過正正這個矛盾.
將來就反映上帝聖潔的屬性.
上帝固然希望人親近祂.
但上帝也是聖潔的神.

$^{521}$和人有分別.
和罪人有分別.
上帝對親近祂的人.
其實是很講究的.
不可以隨便.
上帝雖然和我們同住.
但祂從來沒有說過.
我們可以和祂同屋共住.
雖然我們和上帝同住.
上帝樂意和我們同住.
但我們不可以在同一間屋.
和上帝一起做主人話事.
由始至終.
是上帝在主場.
敬拜親近祂的規矩.
也需要由上帝來決定.
換句話來說.
上帝並沒有因為耶穌基督對我們的拯救.
而降低了祂對聖潔的要求.
你還記得在《美記》裡.
上帝對子民說.
祂是聖潔的.
所以祂的子民也要聖潔.
又或者來到新約.
耶穌對每一個跟隨者說.
我們要像天父完全一樣.
要求沒有改變.
沒有因為耶穌基督的來到.
而改變了上帝對聖潔的要求.
聽起來好像很困難.
上帝要我們聖潔.
要我們完全.
我們覺得我們自己做不到.
不過弟兄姊妹.
當上帝要求我們完全.
要求我們聖潔的時候.
祂是認真的.
祂不是說笑.
不過祂有恩典.
祂會幫我們完成這件事.

$^{561}$祂的意思是.
我們必須認真信教祂.
為祂而活.
靠著祂來成聖.
靠著祂來得以完全.
在耶穌基督的徹底拯救之下.
你和我不再是舊約的一般百姓.
我們升了職是什麼?.
在新約之下.
你記得我們有軍專的祭司.
其實這件事很好.
如果用回回望來說.
我們從前不能進去聖所.
現在可以了.
甚至你可以通過耶穌基督.
你可以進入聖所.
其實我們稱之為祭司.
換句話來說.
如果你是祭司.
那你要怎樣?.
你就要每天跟著大祭司.
來工作.
聽大祭司的吩咐.
聽天上的大祭司.
耶穌基督的吩咐.
來事奉敬拜上帝.
包括兩件事.
第一件事.
我們既然由普通百姓.
升職為祭司.
自然對什麼更加敏銳?.
對罪更加敏銳.
你知道祭司進入聖所之前.
他要做很多潔淨的條例.
我們一樣.
要保持聖潔.
才可以親近上帝.
我們要對罪更加敏銳.
不要讓罪靠近我們.
我們要靠著耶穌基督.

$^{601}$他已經勝過試探.
他也可以幫我們勝過世上的試探.
這是第一件事.
第二件事就是成為君尊的祭司.
我們也要靠著耶穌.
來堅持信服上帝的吩咐.
這是上帝對所有敬拜祂的人.
無論是舊約或新約.
都一致的要求.
不過我們看到.
特別在舊約裡.
神的指紋是明白的.
不過跟不跟又是另一回事.
神的指紋很多時候.
特別在舊約裡.
反映他們輕慢了上帝這兩個要求.
甚至他們做了一件事.
就是嘗試改變上帝的要求.
用自己的方法試圖走近上帝.
例如他們雖然有做獻祭的動作.
但他們的行事為人.
卻是跟上帝的教導背道而馳.
他們覺得只要有聖殿.
有獻祭.
我就可以親近上帝.
舊約的指紋是這樣.
問題來了.
我們呢.
我們會在新約裡.
知道所有舊約的問題.
我們又會否試圖改變上帝.
祂要求我們聖潔完全.
這個標準呢?.
當然作為新約下的信徒.
我們不會像希伯來書的讀者.
在面對試探的時候.
走回頭路.
走回舊路.
想著我生活差點不要緊.
我行事為人差點不要緊.

$^{641}$我靠獻祭.
我靠濟生的血.
來拉近上帝的距離.
我們不會這樣做.
因為我們不會像舊約那樣獻祭生.
但也不見得我們一定會靠耶穌.
來勝過我們所面對的罪惡試探.
有時我們可能寧願.
將罪.
將對上帝的不信服.
來包裝一下.
遮一遮它.
來淡化罪的嚴重性.
譬如說.
我不是犯罪.
我只是軟弱.
有甩漏.
我不是故意有心.
不想聽上帝說話.
又或者.
我們會有這個藉口.
我錯是因為對方錯在先.
我只是以牙還牙自保.
不怪得我.
又或者另一種包裝方法就是.
我這樣怎算是罪呢?.
身邊每個人都是這樣.
不是問題.
不過聖經裡很清楚說.
無論我們怎樣包裝.
怎樣淡化.
怎樣遮掩我們的軟弱.
我們的罪.
上帝對我們的聖潔要求.
是不改變的.
是不變的.
上世紀的圖書牧師.
很出名的牧者.
他曾經說.
他說我們的困難.

$^{681}$源自於我們不願意.
按照上帝的本相.
來接受上帝.
並且我們努力.
按照上帝的吩咐.
來改變自己的生命.
來配合上帝的旨意.
相反我們會怎樣呢?.
我們會試圖改變上帝.
我們會試圖.
令上帝更加像我們的形象.
即是將上帝拉低.
不是我們提升.
囚禁上帝.
是將上帝嘗試拉下來.
和我們接近.
如果我們對罪感到安恕.
為何上帝也不可以像我們這樣呢?.
事實上這種想法.
會令很多信徒.
不單止不認真對付罪惡的試探.
反而更會嘗試.
拉低上帝聖潔的標準.
以為這樣我們可以和上帝更近.
這很可能是我們的回頭路.
所以「鼎姐妹經文」這裡.
很值得我們去審查.
你和我今天怎樣看待罪的試探呢?.
我們會怎樣對付罪呢?.
我們會否當罪就像.
只中了COVID那樣.
不一定會死的.
所以對罪不需要設防.
任由罪人伏在我們的門口.
不趕走他之餘.
還要讓他進入我們的屋.
進入我們的生命.
還是我們很相信.
認定聖經所說.
罪的功架就是.

$^{721}$死和上帝隔絕.
只不過因為耶穌牺牲了自己.
和我們結了婚.
替我們負上必死的功架呢?.
事實上我們有沒有想過.
到底我們的罪債有多重呢?.
雖然我們知道.
但我們有沒有想過有多重呢?.
其實重到一個地步.
連那一位滿有慈愛憐憫的上帝.
他不能夠很簡單地說一句.
「算了,我不跟你計了,不要緊」.
他不能夠這樣.
就將我們的罪一筆勾銷.
因為上帝是聖潔的.
他不能夠因為愛我們的緣故.
就對我們的罪隻眼開隻眼閉.
他亦不能夠降低自己聖潔的標準.
但他又不忍心消滅我們.
那怎麼辦呢?.
你和我很明白.
上帝選擇一個方法就是.
他犧牲他自己唯一的愛子.
耶穌基督.
來還清我們不能解決的罪債.
這一份上帝的恩情.
其實我們是應該很明白的.
但重點是.
如何回應呢?.
這才是重點.
明白就明白.
重點是我們如何回應上帝這份恩情.
曾經有一個猶太人的民間故事是這樣說的.
他說.
耶穌基督他升天之後.
有一次他靜悄悄.
微服出巡回到人間.
看看他拯救過的人到底發生什麼事.
他在路邊見到一個行乞的人.
身體正常.

$^{761}$有手有腳.
於是乎他就問這個人.
這個人就回答.
其實我本來是癱了.
但幾年之前.
耶穌他治好我.
但我一時高興.
四處開派對.
享樂.
幾年後.
我的錢都花光了.
現在就要行乞了.
耶穌難過搖搖頭.
他又再走了幾條街.
他看到另一個醉酒佬.
躺在街上.
全身酒味.
不省人事.
耶穌唯有私下問旁邊的人.
他們就說.
這個人本身患了大麻瘋.
幾年前有個叫耶穌的人治好了他.
所以他活到今天.
不過沒想到.
他那麼不愛惜自己.
現在搞成這樣.
耶穌最後很慨嘆地說.
難道世人就是這樣浪費.
我所賜的恩典和拯救?.
弟兄姊妹這個故事雖然虛構.
不過都很寫實.
它很反映到.
可能我們的屬靈狀況.
或許都是這樣.
但同時亦提醒我們.
不可以再活在罪中之樂.
寧願糟蹋上帝的恩典.
寧願繼續讓耶穌基督傷心.
都不肯和罪認真地.
一刀兩斷劃清界線.

$^{801}$事實上弟兄姊妹.
今天我們看這段經文.
第九章十四節裡面說得很清楚.
我想邀請大家一起讀.
在程序表的尾二那段經文.
希伯來書九章第十四節.
預備起.
「我看基督是一對典禮的.
將自己和阿罌知義和仇恨.
在你歧覓及世證的每一生.
除了對你每一世恨.
只你本是我的阿靈真神」.
這段經文說得很清楚.
其實不需要怎樣的釋經.
你和我也明白.
耶穌基督將我們徹底拯救出來的結果是怎樣.
你是可以事奉永生上帝.
意思是甚麼呢?.
我們沒有藉口.
我們沒有藉口說自己做得不夠好.
說自己還有罪污.
未夠資格.
不是時候將自己奉獻給上帝使用.
耶穌基督說我已經完成了.
你是可以事奉永生的上帝.
你相信耶穌那一刻已經可以了.
上帝不再需要我們.
好像舊約的子民那樣.
你要定期獻上贖罪祭.
先處理好和上帝和好的關係.
那你才可以事奉.
因為耶穌基督已經一次過替我們獻上.
使我們可以徹底有能力擺脫罪的核祭.
但我們很明白一件事.
當你和我和上帝和好之後.
是不是就是這樣?.
是不是你有你,我有我?.
不是的.
我們還要繼續去增潤和上帝的這段關係.
正如結婚的你.

$^{841}$可能你會很明白.
你和配偶吵架的時候.
修補了之後.
是不是就這樣算了?.
不是的.
你要繼續去維繫增值這段關係.
這段關係才可以繼續持續發展.
越來越好.
和上帝的關係一樣.
上帝拯救了我們.
這段關係是要保持.
是要增值的.
所以.
我們雖然不需要再像舊約的子民那樣獻贖罪祭.
但我們要獻什麼祭?.
我們要持使徒保羅.
他在羅馬書十二章所說.
要天天將身體當作活祭獻給上帝.
謹記這個提醒.
很重要的是.
天天獻上活祭.
是我們的份內事.
耶穌基督不會和我們代勞.
這是我們的責任.
耶穌基督獻的是贖罪祭,挽回祭.
但活祭是你和我信了耶穌之後.
要每天向上帝獻上.
意思是什麼?.
什麼是活祭呢?.
首先.
我們要明白.
建制本身其實有犧牲的意思.
對動物的製生來說.
它的犧牲是什麼?.
就是流血.
就是死亡.
但對於我們來說.
是什麼?.
為上帝犧牲.
最極致的當然是為主殉道.

$^{881}$但在現實裡.
上帝不是要我們每次都自己去送死.
上帝的心意是什麼?.
是要我們活.
不是要我們死.
所以耶穌基督將我們拯救之後.
祂會繼續幫助我們成為一個活著.
但徹底信服上帝的祭物.
再說一次.
耶穌基督的設計拯救了我們之後.
是要我們活著成為一個.
信服上帝的祭物.
這就是活祭的意思.
試想像一下.
在舊約當百姓將祭物.
特別是將動物的製生.
放在上帝的壇的時候.
其實就代表他們願意.
放棄對祭物的掌控.
放在祭壇就不關你事.
是上帝的了.
歸了給上帝.
所以上帝要怎樣處理.
放在祭壇上的祭物.
就由上帝作主.
再不是由憲制的人作主.
換句話說.
當我們立志將自己獻給上帝.
就無法作主.
因為你變成了祭物的製生.
上帝要怎樣用你.
這是上帝的主權.
再不是我們自己作主.
所以你不會看到.
在舊約聖經記載.
有人獻祭.
突然說我還是不想這樣獻.
不如按我的方法這樣獻.
不會這樣的.
你更加不會看到.

$^{921}$舊約聖經記載.
那隻躺在壇上的製生.
自己彈起來說.
我想這樣獻.
不會的.
無法作主.
無法給意見.
說上帝我想怎樣獻給你.
獻了就是上帝作主.
作為祭物,作為製生.
根本沒有和上帝討價還價的空間.
就連最完美的製生耶穌基督.
祂都說什麼?.
照上帝你的意思.
不要照我的意思.
所以大英姐妹.
你看到就是.
在利未記裡面.
反覆一開始就說.
獻祭的條例.
是不是人不懂得獻祭?.
不是的.
在沒有律法之前.
亞伯拉罕就已經懂得獻祭.
所有人都懂得獻祭.
只差獻給誰.
所以你看到利未記.
重覆說獻祭的條例.
是說獻祭的極致.
不是照人的意思.
是照上帝的規矩.
透過每日反覆遵守獻祭的條例.
每日反覆做一個.
屬靈操練聽上帝說.
不是聽人說.
重覆的獻祭.
其實有這個作用.
就是操練一個.
信服聽命的生活.
事實上大英姐妹.

$^{961}$通過持續地學習.
我們就學會了.
怎樣為之信服聽命上帝.
同樣.
當我們成為活祭的時候.
我們表明這樣的時候.
意思就是我們願意操練.
從前可能是我想怎樣.
我要怎樣.
轉移到上帝你想我怎樣.
我就怎樣.
其實就是這樣的意思.
透過不斷操練.
將自己成為活祭.
意味著我們不斷學習.
將主權交回給上帝.
既然我們在上帝的壇上.
就不能夠限制上帝.
在我們身上拿多少.
怎樣拿.
和何時拿.
一切由上帝作主.
我們只需要信服.
配合上帝的旨意.
這樣就必蒙月立.
所以大英姐妹.
今日的經文.
很好的提醒.
一方面上帝是定義.
要住在我們中間的.
我們亦都要靠祂.
靠著耶穌基督.
來過一個對罪說不.
對上帝說是的聖潔生活.
並且祂必定會這樣幫我們.
第三方面.
我們亦都要成為活的祭物.
天天獻給上帝來使用.
所以.
我們一起祈禱.

$^{1001}$我們求上帝幫我們回應.
祂今日對我們所發出的呼召.
我們立定心志.
來與上帝同行.
我們同心祈禱.
七位恩主.
因為我們之所以可以.
近距離地.
敬拜,事奉上帝.
本來就不是必然.
是因為耶穌基督.
你用重價將我們贖回.
以致我們今日.
才可以走近聖潔的上帝.
讓我們學會珍惜.
每一次來到你面前的機會.
讓我們有所準備.
因為我們知道上帝引領.
要求走近的人.
要聖潔,要完全.
我們做不到.
但知道靠著耶穌你.
我們可以凡事都能做.
上帝藉著今日的經文提醒我們.
我們要天天向你獻上活祭.
特別當我們立定心志.
要獻給上帝的時候.
當我們成為祭物.
擺在上帝你的壇上的時候.
就意味著我們願意學習.
一點一滴將我們的主權.
交回給上帝使用.
無論上帝你怎樣用我們.
我們都甘心樂意.
上帝是因為你是我們的主.
你是我們的神.
你配得我們這樣回應你.
幫助我們.
讓我們真是願意的.
竭力的來與神同行.

$^{1041}$求耶穌你幫助我們.
求上帝你閱立我們這份心意.
奉基督耶穌你特性的名字.
你都去祈求.
阿們.
謝謝大家.
\newpage



\section{希伯來書 11:11-16}
\label{sec:vxnWocuBbhg}
\textbf{2022年5月8日母親節崇拜直播｜吳真真牧師｜超越的基督:憑著信作母親｜希伯來書11\_11-16}
\newline
\newline
連結: \href{https://youtube.com/watch?v=vxnWocuBbhg}{\texttt{https://youtube.com/watch?v=vxnWocuBbhg}} ~~~~ 語音日期: 2022-05-09
\newline
\newline
\hyperref[sec:rKFXkQ539_w]{\small{< < < PREV SERMON < < <}}
~
\hyperref[sec:index_chronic]{\small{[返順時目]}}
~
\hyperref[sec:qTE37kvGc24]{\small{> > > NEXT SERMON > > >}}
\newline
\newline
希伯來書 11:11-16
\newline
\begin{longtable}{cl}
\hline
\hline
章節 & 經文 (和合本修訂版)\\
\hline
11:11 & \begin{tabularx}{0.7\textwidth}{X} 因著信，撒拉自己已過了生育的年齡還能懷孕，因為她認為應許她的那位是可信的； \end{tabularx} \\ \\ \relax
11:12 & \begin{tabularx}{0.7\textwidth}{X} 所以，從一個彷彿已死的人竟生出子孫，如同天上的星那樣眾多，海邊的沙那樣無數。 \end{tabularx} \\ \\ \relax
11:13 & \begin{tabularx}{0.7\textwidth}{X} 這些人都是存著信心死的，並沒有得著所應許的，卻從遠處觀望，且歡喜迎接。他們承認自己在地上是客旅，是寄居的。 \end{tabularx} \\ \\ \relax
11:14 & \begin{tabularx}{0.7\textwidth}{X} 說這樣話的人是表明自己要尋找一個家鄉。 \end{tabularx} \\ \\ \relax
11:15 & \begin{tabularx}{0.7\textwidth}{X} 他們若想念所離開的家鄉，還有回去的機會。 \end{tabularx} \\ \\ \relax
11:16 & \begin{tabularx}{0.7\textwidth}{X} 其實他們所羨慕的是一個更美的，就是在天上的家鄉。所以，神並不因他們稱他為神而覺得羞恥，因為他已經為他們預備了一座城。 \end{tabularx} \\ \\
[1ex]
\hline
\hline
\end{longtable}
$^{1}$今日經訓在新約三百一十七頁.
希伯來書十一章十一節至十六節.
新約聖經三百一十七頁.
希伯來書十一章十一節至十六節.
「因著信,連撒拉自己雖然過了生育的歲數,還能懷孕.
因他以為那應許他的是可信的,.
所以從一個彷彿已死的人就生出子孫,.
如同天上的星那樣眾多,海邊的沙那樣無數..
這些人都是存著信心死的,並沒有得著所應許的,.
卻從遠處望見,且歡喜迎接,.
又承認自己在世上是客裡,是寄居的..
說者用話的人,是表明自己是找一個家鄉,.
他們若想念所離開的家鄉,.
還有可以回去的機會,.
他們卻羨慕一個更美的家鄉,就是在天上的..
所以神被稱為他們的神,並不以為恥,.
因為他已經給他們預備了一座城.」.
各位弟兄姊妹平安.
今日是母親節.
我在此祝願各位母親.
願你們生命充滿著上帝豐富的恩典,.
精彩地去過美好的人生..
在我心目中,我認識的很多母親,.
她們都是super mom,.
你們覺得嗎?.
我覺得在這個複雜,幻變得這麼快的社會裡,.
我們的母親真是一人身兼多職..
舉例,她們除了上班之外,.
她們在家裡又要做膳食部部長,.
又要做營養學家,.
還要做家務專員,.
如果家裡有外傭的話,.
她還要做職業培訓,.
再加上要管理人事,.
所以她是人事部的經理..
小朋友回來後,.
要教功課,所以她們是功課輔導員,.
也是專職的補習老師..
星期六,日要想精彩的節目,.
她是康樂組組長,旅遊大師,.

$^{41}$如果小朋友生病了,.
媽媽要做什麼?.
醫護人員..
如果家裡有人有喜慶事,.
是不是要做喜慶專員?.
又要做財務大臣,.
現在疫情這樣,.
小朋友要在家裡,.
媽媽又要處理IT,.
我覺得現在的媽媽.
真的全部都是超級媽媽,.
贊不贊成?.
不過媽媽應付這麼多事情,.
已經很吃力了,.
但她們除了應付日常生活外,.
她們也需要應付社會對她的期望,.
或者社會文化對她的看法..
舉個例子,.
如果媽媽是工作,.
她們會被稱為「雙職母親」,.
你有沒有聽過「雙職父親」?.
好像不用做的,爸爸在家裡..
這樣代表了在社會裡,.
媽媽打兩份工是理所當然的,.
但爸爸可以拿著一個免死金牌,.
不是,拿著一個准許證,.
可以在他的事業裡打拼..
再加上近年有很多壓力,.
我聽到有些媽媽跟我分享,.
社會運動帶來了年輕人,.
或者新一代跟父母之間的鴻溝..
另外在家裡,.
因為疫情關係,.
又要在家工作,.
很多的衝突,.
我覺得現在的媽媽擔子越來越重,.
挑戰越來越多,.
究竟在這麼巨大的壓力下,.
我們如何做一個好的媽媽的角色?.
或者我們問,.

$^{81}$何時才能看見家鄉?.
我們經常說這句,.
今天我們很想從希伯來書11章,.
11至16節,.
去了解一下,.
其實我們如何憑著信心,.
去做一個基督徒的母親的角色?.
說到這裡,.
可能有些等級姐妹說,.
「那我可以從這一刻開始睡覺了.」.
你也不是媽媽,.
不用聽吧?.
不是,.
憑著信心,.
在這個地上去活好上帝所賦予我們,.
在這個地上的角色,.
去擁抱這個使命,.
是我們每一個人都應該要聽的..
讓我們一起去祈禱吧!.
上帝啊,求你公抹使女的口,.
讓我孫說的是主你的話,.
也願聖靈在眾人的心裡幫助我們,.
讓我們去明白,.
你想在我們這個那麼多挑戰的年代,.
在我們的身份角色裡面,.
我們如何可以憑著信心,.
去活一個精彩的人生?.
主啊,我就在這裡懇切地祈求你,.
引導我們,幫助我們,.
開我們的眼睛,.
開我們屬靈的耳朵,.
讓我們打開我們的心,.
去回應你的道..
我們這樣禱告交託,.
奉主耶穌基督的聖名自祈求,.
阿們..
上星期陳明泉牧師,.
一直在跟我們說希伯來書第十章,.
今天我們會去到第十一章,.
上星期第十章,.

$^{121}$他告訴我們,.
我們這群信徒,.
要繼續堅持著對主的信心,.
三十九節就告訴我們,.
以至到靈魂可以得救,.
我們可以保住我們的命..
接下來第十一章,.
整章都是講信心的,.
牧師在最後跟我們說,.
信心有幾個定義,.
信心是要我們堅持的,.
信心是預見的,.
信心是我們可以去到神面前的途徑,.
今天我們會再進去看看,.
到底我們可以如何憑著信心去行..
十一章有一個很精彩的特色,.
就是整本聖經裡有十八個恩著信的字,.
這十八個恩著信的字,.
列舉了一大群信心為人的榜樣,.
可以說是一個對信心的重篇..
來到這裡,.
我們第十一節,.
其實是描述其中一位列祖,.
他自己的信心經歷,.
但是他坐落在第八節至第十九節,.
是關於亞伯拉罕四個恩著信的功課..
第八至十節,.
他提到兩件事,.
第一是亞伯拉罕的夢照,.
他如何有信心,.
第二是亞伯拉罕如何在一個應許之地,.
去作一個客女..
這兩個都是亞伯拉罕的兩項信心的表現..
來到第十一節,.
「恩著信連撒拉自己,.
雖然過了生育的歲數,.
還能懷孕,.
因他以為那應許他的,.
是可信的,.
所以從一個彷彿已死的人,.

$^{161}$就生出了子孫,.
如同天上的星那麼眾多,.
海邊的沙那麼無數.」.
這裡提到亞伯拉罕的第三個恩著信..
有一些解經家提到,.
撒拉生出兒子二撒,.
究竟是亞伯拉罕有信心,.
還是撒拉有信心呢?.
有人提問,.
因為在一些譯本中,.
有些爭議,.
有些譯本將第十一節的「他」,.
譯成女性的「他」,.
但有些譯本將「他」,.
譯成男性的「他」..
為什麼會有這個爭議呢?.
第一,.
由第八節開始說恩著信,.
就是亞伯拉罕,亞伯拉罕,.
最後的又是亞伯拉罕,.
但為什麼第三個恩著信又是撒拉的呢?.
這是第一,.
會不會是他跳走了呢?.
第二,.
原來是一個文字的解釋,.
第一,.
第十二節說到.
「從一個彷彿已死的人生出子孫」,.
「一個」的「人」和「已死」,.
這兩個字都是陽性字,.
知道什麼是陽性字嗎?.
就是masculine,.
即是紅性,.
藍性..
第十二節說到一個男人.
快要死了,.
出了子孫,.
這就是亞伯拉罕,.
第二,.
第十一節的橫能懷孕的翻譯,.

$^{201}$這個翻譯是比較關鍵的,.
因為在希伯來文裡,.
這個字也是一個陽性字,.
而且這個字也是一個主動的字,.
什麼意思呢?.
意思是男性將精子置於母體之內,.
是男性主動的..
所以這樣說,.
其實應該是在說亞伯拉罕..
所以我自己比較選取的,.
好像新漢語譯本的譯法,.
新漢語譯本的譯法是,.
「因著信,雖然撒拉自己不能生育,.
亞伯拉罕又上了年紀,.
他是男性的他,.
還能夠作父親.」.
我自己這樣看,.
不過雖然我們說,.
第三個因著信主持的主題是亞伯拉罕,.
不是直接說撒拉,.
但是我很想說一件事,.
生孩子的事,.
可不可以沒有女人?.
當然不可以了,.
所以在生孩子這件事上,.
在亞伯拉罕在100歲,.
撒拉90歲,.
可以生孩子這件事上,.
撒拉的確是有份,.
而且我還認為他是有因著信去行的,.
為什麼呢?.
如果大家熟悉《創世記》12章,.
亞伯拉罕的事蹟,.
一直打下去的時候,.
當12章的時候,.
亞伯拉罕和撒拉還沒有改名,.
他們還叫做「亞伯蘭」和「撒萊」的時候,.
神已經分別在12章,13章,15章,.
跟他們說了他的應許,.
他的應許就是,.

$^{241}$他要以迦南為業,.
他要成為大國,.
他的子孫就像凱撒一樣多,.
還要萬國恩去得福,.
早就說了,.
但是你要留意,.
第11章已經說明,.
撒拉是不辱的,.
當神在第12章的時候,.
第一次向亞伯拉罕說他的應許的時候,.
原來撒拉已經有一個既定的事實,.
就是他不辱的,.
在人的眼中,.
但是,.
當亞伯拉罕領受這個應許的時候,.
他一直想他的下一代,.
其實他沒有想過是撒拉生的,.
到這個時候,.
耶和華還沒有跟他說,.
他要從撒拉而出的後裔,.
才是他所揀選的後裔,.
所以在第12章,第13章,第15章,.
亞伯拉罕似乎有一個傾向,.
有些解經學者認為,.
亞伯拉罕心裡想繼承他,.
作他的後裔是羅德,.
作他的後裔,.
第15章是說他的管家以便以謝,.
直到第17章,.
才有一個明確的說法,.
就是原來他要從撒拉而出,.
但是我作為一個女人,.
我是一個人的老婆,.
撒拉作為亞伯拉罕的妻子,.
其實她跟丈夫一樣,.
她知道耶和華應許要賜她丈夫有後裔,.
丈夫也一直在尋找上帝的心意,.
到底是誰?.
但是她苦等了十年,.
當時她已經75歲,.

$^{281}$年紀這麼大,.
於是在創世記16章說她做什麼?.
她找自己的婢女下嫁代替生孩子,.
當時的風俗是,.
其實我的婢女是我的產業,.
她幫我和丈夫生孩子,.
兒子也是我的,.
她這樣想,這是第一,.
第二就是,.
其實撒拉是焦急於上帝的應許,.
何時可以兌現,.
所以她向亞伯拉罕提出建議,.
而亞伯拉罕又認同了,.
然後就有二十瑪利出世了,.
所以在創世記第十九章,.
亞伯拉罕已經99歲,.
撒拉已經89歲了,.
神與亞伯拉罕立約的時候,.
亞伯拉罕以為是二十瑪利,.
是上帝所揀選的,.
但是到那個時候,.
上帝才跟他說,.
撒拉生的,.
才是耶和華所揀的猴子,.
這樣看來,.
撒拉所做的,.
是對上帝應許的確信,.
她對上帝的應許是有期盼的,.
她跟亞伯拉罕一樣,.
她對上帝應許他們有猴子,.
是一樣的期盼,.
不過她熱心有餘,.
她用了自己的方法,.
她走得比上帝快,.
但是亞伯拉罕同意,.
亞伯拉罕可以稱為信心之父,.
其實她也是一個配合丈夫有信心的女人,.
亞伯拉罕無兒無女,.
她能夠成為大國,.
其實對於他們來說是天方夜譚的事,.

$^{321}$他們可以相信是非常難得的事,.
但是他們在整個旅程裡,.
他們不明白,.
他們所以一而再,.
再而三地尋問,.
其實就是他們確信上帝的應許,.
只不過是上帝的心意,.
往往不是一次過向我們展露,.
而是逐一一步一步,.
這樣展現出來,.
所以我們看到,.
亞伯拉罕可以成為大國,.
必須要有後知,.
而且在這個八歲生子的信心考驗上,.
上帝預埋撒拉是有份的,.
在這件事上,.
我們看到一樣東西,.
其實原來信心不是個人的事,.
不是一個人的事,.
是夫婦的事,.
是整個家庭的事,.
如果夫婦的家庭目標是一致的話,.
起碼這個家庭可以向前走,.
如果兩個夫婦是道不同不相為謀的話,.
結果大家各自走的時候,.
有時我會看到孩子被五額分屍,.
很可憐,很淒涼,.
所以我很想趁著母親節,.
我們會提醒各位弟兄姊妹要孝順媽媽,.
不過我現在也特別跟丈夫說,.
其實你的太太可以成為一個好母親,.
丈夫是很關鍵,.
丈夫能否同心是很關鍵,.
如果夫婦都是信主,.
可以一同愛主,一同服侍主,.
一同跟從主,.
整個家庭都會蒙福,.
而且也會藉著你的家庭去祝福別人,.
不過說到底,.
最重要的是目標是什麼?.

$^{361}$另外我們看到,.
信心會突破人的軟弱和限制,.
超越我們的人生經驗,.
因為神是可信,是信實的,.
因為剛才讀到第十二節是這樣說的,.
是嗎?.
第十一節說,.
因他以那應許他的,是可信的,.
第十一節尾是這樣說的,.
很多人說,.
當耶和華在《創世記》第十八章時.
跟阿伯拉肯說,.
撒拉生的才是上帝所說的後製,.
撒拉偷聽到,他做了什麼?.
暗笑,是嗎?.
他在偷笑,.
他這樣記載第十二節說,.
「我以以衰敗,我主也怒賣,豈能有者喜事?」.
這個,真是在懷疑上帝,.
好像一個沒有信心的表現,.
大家不要說撒拉,.
其實阿伯拉肯也是在笑,是嗎?.
其實第十七節,.
阿伯拉肯聽到的時候,.
他也是伏在地上笑,.
他心裡說,.
「一百歲的人還能生孩子嗎?」.
「撒拉已經九十歲,還能生養嗎?」.
其實阿伯拉肯也在笑,是嗎?.
不過我在想,難怪他們起初是不相信的,.
他們兩個都已經過了生育的年紀,.
生理上的角度,.
其實覺得這是絕對不可能的事,.
而且,他這裡說什麼?.
他彷彿已經是一個死了的人,.
即是說,有一隻腳已經踏入棺材了,.
不過雖然現在政府的官員告訴我們,.
六十歲都是青年,.
是嗎?大家是不是青年?.
但是其實他們說六十歲是青年,.

$^{401}$也是在說,如果你保持得好,身體好,.
你還有精力可以繼續工作,是嗎?.
還有精力可以服務社會,是嗎?.
但是有沒有說,.
「你六十歲是青年,可以生孩子?」.
沒有啊,.
因為生孩子,或者生理上可以生孩子,.
其實這是一個違反常理的事情,.
所以其實薩拉和阿伯拉肯聽到這件事之後,.
他們覺得不能盡信,很懷疑,.
是不是真的?.
但是大家記不記得,.
當薩拉暗笑之後,.
耶和華說了一句什麼話?.
記得吧?.
「耶和華豈有難成的事嗎?」.
我相信就是這句話,.
我們就知道,.
縱然人是有軟弱,.
人是有限制,.
不過上帝是信實的,.
是可信的,.
所以阿伯拉肯和薩拉.
可以將那些看不到的東西,.
例如上帝的能力,.
看不到的,.
又或者他們還沒看到的東西,.
例如應許的兌現,.
他們看為這是事實,.
這就是信的本質,.
重點是上帝是信實,.
上帝有大能,.
是耶和華豈有難成的事,.
他說了的話,.
就算多有限制,多困難,.
都可以成真..
去年我母親節的時候,.
大家看不到我,.
我去了港島,.
其中在我港島之前,.

$^{441}$有一位見證的嘉賓,.
他的見證令我很印象,.
因為她是一個小兒麻痺症的媽媽,.
不單止,.
不只是她小兒麻痺症,.
她是小兒麻痺症,.
有些嚴重要坐輪椅的那種,.
重要的是她的丈夫,.
也是一個小兒麻痺症的患者,.
兩個看來極有限制,.
你沒想過,.
或者他們自己都沒想過,.
上帝會賜一個女兒給他們,.
他們怎樣可以照顧女兒呢?.
他說,.
他起初都很沒信心,.
覺得自己應付不來,.
而實際上,.
生活裡面真的有很多困難,.
你想像到,.
坐輪椅,.
要照顧一個BB女兒,.
她的出世,.
另外,.
其實當BB越來越長大,.
越來越多需要照顧的時候,.
他就決定,.
因為想專心一意照顧女兒,.
因為實在太累了,.
他辭去了一份,.
他做了22年的公務員工作,.
因為他很想,.
用他有的精力和能力,.
全心全意照顧女兒,.
所以在他的經濟上,.
都沒以前那麼鬆動,.
沒以前那麼有條件,.
但他決定這樣做,.
他很努力照顧,.
他最深刻的印象是,.

$^{481}$他的女兒很可愛,.
在他小學的時候,.
他和他身邊的同學說,.
你們不要小看我媽媽,.
你們當然以為我會照顧我媽媽,.
他和他的小朋友說,.
他和他身邊的同學說,.
是我媽媽照顧我,.
而且照顧得無微不至,.
沒錯,他是傷殘,.
他有很多限制,.
但這對夫婦對上帝很有信心,.
他們克服了很多困難,.
他分享了很多不同的,.
與女兒之間的環節,.
因為沒有錢,.
所以買東西的時候,.
要跟人計較,.
但卻在當中,.
與女兒談了無數的人生道理,.
他說,神應許他,.
敬畏他的人,.
一樣也不缺,.
他和他的女兒就是這樣分享,.
所以他向上帝求智慧,.
沒有錢,.
上帝仍然對他的家有供應,.
他覺得自己沒有智慧,.
不懂得教,.
他看到上帝的教導和幫助,.
從來沒有離開他,.
他說,我的家,.
一樣也不缺,.
弟兄姊妹,.
我們看到亞伯拉罕和撒拉,.
在生育的事上,.
其實在整個歷程中,.
的確有他們的限制,.
更加在他們回應上帝的呼召,.
以及確信應許的路上,.

$^{521}$他們真的有軟弱,.
真的有懷疑,.
真的有真理的方法,.
他們也試過在等候的過程中,.
失去了耐性,.
用自己的方法去解決,.
但是很感恩,.
或者很值得我們留意的是,.
希伯來書的作者仍然視他們對上帝的信靠,.
是一個值得我們學習的榜樣,.
上帝應許撒拉生兒子,.
上帝應許後子到撒拉生兒子,.
足足有二十五年的時間,.
無論環境多麼不順利,.
無論結果如何,.
他們依然願意相信,.
而且他們在過程中,.
不是停下來,.
他們會問,.
他們會按照上帝的心意去走,.
去跟,.
去試,.
所以弟兄姊妹,.
有信心的信徒,.
必須要堅持去走,.
去明白神的心意,.
去實踐上帝所願,.
來到第十三至第十六節,.
除了我們有信心,.
我們必須要向前行之外,.
更加重要的是,.
我們要有一個超越的眼光,.
在第十三至第十六節,.
這些人都是全著信心死的,.
並沒有得著所應許的,.
但卻從遠處望見,.
且歡喜迎接,.
又承認自己在世上是客裡,.
是寄居的,.
說這樣的話的人,.

$^{561}$是表明自己要找一個家鄉,.
他們若想念所離開的家鄉,.
還有可以回去的機會,.
他們卻是沒有一個更美的家鄉,.
就是在天上的,.
所以神被稱為他們的神,.
並不以為此,.
因為他已經給他們預備了一座城,.
第十三至十六節的段落,.
是講述信心的一個小小的總結,.
這段正正就在整個族長的榜樣裡,.
即是在亞伯拉罕,.
以撒和雅各族長的信心表的榜樣描繪裡,.
正正中間,.
其實主要就是想告訴我們,.
我們在地上生活的時候,.
尋找天上家鄉的末世性,.
即是說我們憑著信心,.
必須要有超越的眼光,.
這有個很重要的字,.
就是「更美的家鄉」,.
不知道在座的媽媽,.
或者你們做父母的,.
或者你們自己的人生裡,.
看見家鄉了嗎?.
媽媽,.
你的家鄉在哪裡?.
他能夠成功大學畢業,.
找到一份好工作,.
成家立室,.
是不是你的家鄉?.
不過,我所知的就是,.
養兒一百歲,.
長一路九十九,.
一世是媽媽,.
終生是媽媽,.
一日是媽媽,.
終生是媽媽,.
是不是?.
所以其實,.

$^{601}$我們的家鄉不是在地上,.
如果我們的眼睛是看著地上的家鄉,.
我們的心永遠不安寧,.
永遠不安寧,.
我們會失去焦點,.
所以究竟,.
什麼才是我們應該追求的呢?.
希伯來書的作者,.
講到一個很殘酷的現實,.
我們看看那群列祖,.
他們究竟追求什麼?.
這裡說他們追求一個近味的家鄉,.
事務一個近味的家鄉,.
但是,.
希伯來書的作者,.
第十三,第十四節,.
一開始就講到,.
其實到他們死的那一天,.
仍然沒有得著應許,.
不過,.
大家不需要絕望,.
因為聖經仍然說,.
他們死的時候仍然是有信心的,.
因為他們遠望家鄉的時候,.
是用信心的眼睛,.
才看到這個應許的兌現,.
而且,.
他形容到他們的心情是歡喜的,.
是擁抱這個應許的,.
而且他們很清楚自己在地上,.
只不過是寄居的異鄉者,.
阿巴拉罕,以撒同亞國,.
到死的那一天,.
都未能實現上帝應許他們,.
而迦南維地成為大國,.
地上萬族因許得福的應許,.
但是,.
他們的榜樣告訴我們一件事,.
他們到死仍然能相信,.
是因為他們很明白,.

$^{641}$在地上的祝福,.
一些應許的實現,.
只不過是我們先祥與主同在的屬靈福氣,.
這些應許真正的兌現,.
是在未來,.
所以上帝為了給我們這份盼望,.
在地上給予一些保證,.
所以你會看到這班列祖一點遺憾都沒有,.
不單止這樣,.
這班列祖在地上一直熱切地等待,.
享望著一個家鄉,.
即是第十六節所說的,.
在天上,近末的家鄉,.
比起他們在地上支搭帳篷的生活,.
一座有根有機的城,.
令他們更加嚮往,.
所以這班列祖的目光,.
就是放在天上,.
天上就是神與我們同在的地方,.
亞伯拉罕這班列祖是多麼渴望有神同在,.
得到神應許,.
和他有親密的屬靈關係的時間,.
在舊約裡面,這個景象是模糊的,.
但當耶穌基督來到這個世界之後,.
新約告訴我們,.
因著基督在十字架裡面的成就,.
我們是可以有著這份救贖的恩典,.
我們在地的時候,.
我們得蒙救贖,.
當我們離開這個世界的時候,.
我們可以去到天上,.
與主同在,.
所以列祖的眼光,.
是一個超越的眼光,.
他們向前望的不是地上的事,.
而是放在近尾,.
在天上的家鄉,.
另外第十五節,.
那裡提到,.
他們若想念所離開的家鄉,.

$^{681}$還有可以回去的機會,.
其實這一節提醒我們要堅定忍耐,.
才是信心的表現,.
這裡提及的第一,.
他們提及的家鄉,.
不是提及亞伯拉罕離開吳爾,.
吳爾是他的家鄉,.
他用了「他們若懷念」,.
其實是暗示他們一點也不懷念,.
因為當他們去到迦南作客旅的時候,.
其實他們有機會回去的,.
但他們並沒有這樣做,.
亞伯拉罕幫爾薩回鄉找老婆的時候,.
都是叮囑他的僕人,.
千萬不要帶兒子回家鄉,.
忍耐邁向一個仍然未兌現的應許,.
忍不住要調轉頭走,.
是一件很尋常的事,.
所以希伯來書的作者,.
千方百計去叮囑這班讀者,.
警告他們,.
在第六章甚至用很多字眼提醒他們,.
千萬不要學那班在以色列人,.
在出埃及,.
在曠野流浪漂流的這班列祖,.
千萬不要學他們,.
以免他們沒有重新懊悔的機會,.
所以他叮囑我們要持守,.
要堅持千萬不要走舊路,.
所以我們看到,.
這班列祖的榜樣,.
他們能夠去到上帝應許的近美家鄉,.
能夠堅持忍耐,.
繼續向前行,.
絕對不回頭,.
是我們的榜樣,.
說到向目標進發,.
是需要勇氣,.
是需要堅持,.
我最近看過一個Youtube,.

$^{721}$有一個年青人,.
他那句說話是甚麼?.
他說為了得到夢寐以求的自製物業,.
你可以去到多盡?.
這個年青人有一句名言,.
自己不吃,層樓都要吃,.
所以由他大學開始的時候,.
他限定自己每日只准用50元,.
他走路上學上班,.
不去任何社交活動,.
在家開著Youtube,.
當場卡拉OK,.
去公共浴室洗澡,.
不拍拖,.
因為拍拖要花錢,.
不過他很乖,.
他照給家用,.
結果他用了七年的時間,.
不,應該這樣說,.
他儲了,.
他整份工資有七成可以儲蓄,.
結果他用了八年時間,.
儲了三百萬,.
然後就置業了,.
弟兄姊妹,.
這個人為了地上的一間樓,.
捍衛他的近尾家鄉,.
可以去到這麼盡,.
可以有這麼多忍耐,.
可以有這麼堅持,.
可以有這麼犧牲,.
我們有沒有想過,.
我們可以為這個未來,.
永恆天上近尾的家鄉,.
可以去到多盡,.
我們要搞清楚,.
如果我們的家長,.
我的家鄉是我兒子大學畢業成家立室,.
買了一套樓,.
我現在就會為了這個預備,.

$^{761}$於是,.
就怎樣?.
供書教學,.
讀書就是一切,.
放棄其他東西,.
因為要達到這個家鄉的目標,.
但是,無論是媽媽,.
無論是我們這裡哪一個,.
我們的家鄉,.
是不是上帝給我們的這份救贖,.
是不是這個天上近尾的家鄉?.
第十六節告訴我們,.
「所以神被稱為他們的神,.
並不以為此.」.
其實就是要公開承認這班列祖,.
上帝認他們,.
「是的,這班是我的子民,.
如果你站在人的當中,.
上帝說,.
是的,你是跟從我的認真的基督徒,.
是何等的榮譽?」.
上帝是信實的上帝,.
祂不出意反意,.
「如果當人持守所信,.
恆心忍耐,.
即使上帝所應許在地上的事.
未能兌現,.
祂都要與他們同在,.
給他們近尾的家鄉.」.
媽媽,.
我們可能像剛才一開始那樣說,.
我們身兼多職,.
但如果我們是向著近尾家鄉進發的時候,.
我們要做一個怎樣的母親呢?.
早幾年我看了一本書叫做《母親的使命》,.
它裡面說到我們作為屬靈的母親,.
我們到底有什麼使命呢?.
它說到原來我們的使命是.
僕人,導師,強大的朋友,靈魂的守護者,.
管家,設計者,牧養者,社稷的愛..

$^{801}$嘩!你聽到之後害怕嗎?.
很厲害!.
我們要分辨,.
有時我們會看我們的孩子在今生有多少成就,.
在地上擁有多少,.
但我們有沒有著緊上帝怎樣看我們呢?.
這本書對我們的提醒是,.
上帝要我們成為一個怎樣的母親..
好像很難,.
但這本書有一個最大的安慰,.
它說母親是耶穌的門徒..
我讀到那一章的時候,.
我心裡覺得很溫暖,.
因為耶穌是我們生命的師父,.
祂明白我們的處境,.
祂會安慰我們,.
祂會教我們怎樣做..
剛剛那星期《星宗志行》收身放榜,.
大家知不知道?.
其實三年前我經歷過一次,.
小孩中一的時候,.
我們滿心有很多計劃,.
以為上帝實食無千雅,.
他和三個同學一起選了一間,.
他心目中很想入的中學,.
就是哥哥讀的那一間..
但當那天派位的時候,.
他那對失望的眼神,.
我現在還記得,.
他三個同學都入了那間學校,.
唯獨他沒有入..
我心裡也覺得很難過,.
他派了我去一間比較低階的學校,.
而且不是那一區的..
我們那天到處去,.
去了很多間學校,.
直到去到一間學校,.
那間學校說,.
「好,你下午做兩個半小時的測驗,.
你先去吃飯吧.」.

$^{841}$我兒子在吃飯的時候,.
他問了一個問題,.
他問我,.
「為什麼天父這麼壞?.
為什麼他這麼安排我?」.
我真的不懂回答,.
我只記得我慌忙之中,.
又教他人生的大道理..
上帝給….
我說過自己的經歷,.
上帝沒有給我,.
其實也是為我好的,.
大致是這樣..
到我兒子進去考試,.
其實我坐在禮堂的時候,.
我也哭了,.
我不期然地流眼淚,.
我覺得很累,.
我覺得很失望,.
我也覺得很沮喪,.
我和我兒子一樣問,.
「上帝啊,為什麼?.
為什麼你這樣對我的兒子?」.
但是當我安靜下來的時候,.
我發覺這個超越的眼光,.
看著近美家鄉的眼光,.
很能幫助我,.
很能幫助我..
我問自己,.
我的人生,.
我作為母親,.
我的孩子,.
他人生追求的是什麼?.
真的要讀一間Bandman的學校?.
真的要有成就?.
做醫生?.
我發覺,.
原來在這個時候,.
我看到我的兒子,.
在他失望,.

$^{881}$在他絕望的那一刻,.
他問誰?.
他問上帝,.
他心裡有上帝,.
他也知道,.
其實上帝是有主的,.
是不是?.
還有,.
之後他慢慢向神呼求,.
「神啊,求你讓我補考,.
讓我進入那間」,.
他慢慢呼求,呼求,呼求..
這不是他人生最美好的功課嗎?.
人生不是會有失敗嗎?.
人生不是會事事如意的,.
是不是?.
但是他可以找回上帝,.
他可以和上帝結連..
所以弟兄姊妹,.
我也很願意鼓勵大家,.
願意我們都認定,.
上帝的應許是會兌現的,.
我們眼望近美的家鄉,.
勇敢向前行,.
勇敢去問,.
上帝的信實會幫我們突破很多障礙,.
而且我們憑著堅忍,堅信,.
就可以去到上帝為我們預備的近美家鄉..
謝謝大家.
\newpage



\section{希伯來書 12:1-13}
\label{sec:qTE37kvGc24}
\textbf{2022年5月15日崇拜直播｜陳啟興牧師｜超越的基督:人生馬拉松｜希伯來書12\_1-13}
\newline
\newline
連結: \href{https://youtube.com/watch?v=qTE37kvGc24}{\texttt{https://youtube.com/watch?v=qTE37kvGc24}} ~~~~ 語音日期: 2022-05-16
\newline
\newline
\hyperref[sec:vxnWocuBbhg]{\small{< < < PREV SERMON < < <}}
~
\hyperref[sec:index_chronic]{\small{[返順時目]}}
~
\hyperref[sec:uzCPC0kQbAA]{\small{> > > NEXT SERMON > > >}}
\newline
\newline
希伯來書 12:1-13
\newline
\begin{longtable}{cl}
\hline
\hline
章節 & 經文 (和合本修訂版)\\
\hline
12:1 & \begin{tabularx}{0.7\textwidth}{X} 所以，既然我們有這許多見證人如同雲彩圍繞著我們，就該卸下各樣重擔和緊緊纏累的罪，以堅忍的心奔那擺在我們前頭的路程， \end{tabularx} \\ \\ \relax
12:2 & \begin{tabularx}{0.7\textwidth}{X} 仰望我們信心的創始成終者耶穌，他因那擺在前面的喜樂，輕看羞辱，忍受了十字架的苦難，如今已坐在神寶座的右邊。 \end{tabularx} \\ \\ \relax
12:3 & \begin{tabularx}{0.7\textwidth}{X} 你們要仔細想想這位忍受了罪人如此頂撞的耶穌，你們就不致心灰意懶了。 \end{tabularx} \\ \\ \relax
12:4 & \begin{tabularx}{0.7\textwidth}{X} 你們與罪惡爭鬥，還沒有抵抗到流血的地步。 \end{tabularx} \\ \\ \relax
12:5 & \begin{tabularx}{0.7\textwidth}{X} 你們又忘了神勸你們如同勸兒女的那些話，說：「我兒啊，不可輕看主的管教，被他責備的時候不可灰心； \end{tabularx} \\ \\ \relax
12:6 & \begin{tabularx}{0.7\textwidth}{X} 因為主所愛的，他必管教，又鞭打他所接納的每一個孩子。」 \end{tabularx} \\ \\ \relax
12:7 & \begin{tabularx}{0.7\textwidth}{X} 為了受管教，你們要忍受。神待你們如同待兒女。哪有兒女不被父親管教的呢？ \end{tabularx} \\ \\ \relax
12:8 & \begin{tabularx}{0.7\textwidth}{X} 管教原是眾兒女共同所領受的；你們若不受管教，就是私生子，不是兒女了。 \end{tabularx} \\ \\ \relax
12:9 & \begin{tabularx}{0.7\textwidth}{X} 再者，我們曾有肉身之父管教我們，我們尚且敬重他，何況靈性之父，我們豈不更當順服他而得生命嗎？ \end{tabularx} \\ \\ \relax
12:10 & \begin{tabularx}{0.7\textwidth}{X} 肉身之父都是短時間隨己意管教我們，惟有靈性之父管教我們是要我們得益處，使我們在他的聖潔上有份。 \end{tabularx} \\ \\ \relax
12:11 & \begin{tabularx}{0.7\textwidth}{X} 凡管教的事，當時不覺得快樂，反覺得痛苦；後來卻為那經過鍛鍊的人結出平安的果子，就是義的果子。 \end{tabularx} \\ \\ \relax
12:12 & \begin{tabularx}{0.7\textwidth}{X} 所以，你們要把下垂的手舉起來，發酸的腿挺直； \end{tabularx} \\ \\ \relax
12:13 & \begin{tabularx}{0.7\textwidth}{X} 要為自己的腳把道路修直了，使瘸了的腿不再脫臼，反而得到痊癒。 \end{tabularx} \\ \\
[1ex]
\hline
\hline
\end{longtable}
$^{1}$願意我們一同用聆聽神的話語.
繼續我們的敬拜.
今日我們會選讀的經文是.
在新約聖經的希伯來書第十二章一至十三節.
新約聖經希伯來書十二章.
他在這裡這樣說.
「我們既有者許多的見證人.
如同雲彩圍著我們.
就當放下各樣的重擔.
脫去容易纏累我們的罪.
全心忍耐.
不那擺在我們前頭的路程.
仰望為我們信心創始成中的耶穌.
他因那擺在前面的喜樂.
就輕看羞辱.
忍受了十字架的苦難.
便坐在神寶座的右邊.
那忍受罪人這樣頂撞的.
你們要思想免得疲倦灰心.
你們與罪惡相爭.
還沒有抵擋到流血的地步.
你們又忙了.
那勸你們如同勸兒子的話說.
我兒你不可輕看主的管教.
被他責備的時候也不可灰心.
因為主所愛的他必管教.
又鞭打凡修立的兒子.
你們所忍受的.
施慎管教你們.
待你們如同待兒子.
焉有兒子不被父親管教的呢.
管教原是眾子所共受的.
你們若不受管教.
就是失子不是兒子了.
再者我們曾有新生的父管教我們.
我們尚且敬重他.
何況萬靈的父.
我們豈不更當信服他的身嗎.
新生的父都暫隨己義管教我們.
唯有萬靈的父管教我們.

$^{41}$是要我們得益處.
使我們在他的聖潔上有份.
凡管教的士.
當時不覺得快樂.
反覺得受苦.
後來卻為那經念過的人.
結出平安的果子.
就是義.
所以你們要把下垂的手.
髮宣的腿挺起來.
也要為自己的腳.
把道路修直了.
使缺子的不自歪腳.
反得全愈.
足夠法定人數.
今天主要我們很開心請到陳牧師為我們正道.
主題是超越的基督人生馬拉松.
有請陳牧師.
丁姐妹早安.
在家裡的丁姐妹早安.
在我們聽和講之前.
我們要謙卑地同心道告.
親愛的主我們懇求你的聖靈幫助我們.
如果不是你同在你幫助.
我們沒有辦法去明白.
也沒有辦法去行出真理.
我們的人生會一塌糊塗.
當我們走這條人生路的時候.
我們很仰望你自己說領我們.
我們更加求你在今天.
你用你自己的話語.
成為我們的幫助.
我們的引導.
我們仰望你奉耶穌基督的聖名.
阿們.
我不知道在在有沒有人走過毅行者.
如果有走過可不可以舉手.
我們很驚訝.
完全沒有人走過毅行者.
OK.

$^{81}$當然毅行者不容易.
我二十多年前就走過一次.
也是我的第一次.
也是我的最後一次.
不容易的.
我和幾個弟兄姊妹就去組隊參加毅行者.
毅行者是走一百公里的麥理浩徑.
事前要有很多準備和練習.
或者你做過支援團隊可能也會理解.
要練體能.
一百公里不是輕散地走.
哪一段要夜行.
通常是第三段再開始.
走到晚上.
晚上走上下山不容易的.
哪段上下斜最多.
要組織支援隊.
support team.
要在哪一個checkpoint去支援.
支援要給什麼物資.
給什麼服務.
給什麼食物.
很多事情要安排.
其實人生是一場長途賽.
就好像走毅行者.
要有很好的準備.
你才能走一個美好的人生.
不要像和尚一樣.
我們要做一日和尚.
敲一日鐘.
我們要好好準備我們的人生.
人生就好像跑馬拉松.
我不知道在座有沒有人跑過馬拉松.
王晝不運動嗎?.
我明白了.
但我看來不是全部人都應該是長者.
可能大家都想成熟一點.
我都沒有跑過.
我最長都是跑渣打十公里.
我的體能也有限.

$^{121}$但原來跑十公里.
也是要事前很多操練和準備.
不是就這樣貿貿然.
換件衣服就去.
無論走毅行者和跑十公里.
去到某些位置.
如果你試過.
你是想放棄.
去到某些位置.
走毅行者.
走完第七段.
我們叫去到懸礦坳這個地方.
體力也去到七七八八.
然後前面是什麼呢?.
就是要跨越香港最高的山.
大帽山.
所以這是第八段.
如果你少點堅持.
想放下.
算了已經這麼累.
還要跨越這麼高的山.
如果跑十公里.
有時候我跑到第七八公里.
開始沒什麼氣.
開始沒之前那麼有氣.
徵狀是什麼呢?.
橫隔膜開始頂著有點痛.
如果沒什麼堅持.
其實就會慢慢放慢.
慢慢從跑變成走.
所以到那些位置是要堅持的.
其實希伯來書的作者都鼓勵我們.
所有讀者.
我們人生是一場長途賽.
你就當是像是毅行者或者馬拉松.
或者十公里一樣.
你中間一定有些困難.
有些逆境.
一定有的.
有些位置要堅忍地走下去.

$^{161}$一定要很多的仰望上帝.
堅忍地走下去.
這段經文有很多運動和比賽的用語.
我們看看作者怎麼教導我們.
去走這一場或者跑這一場人生的馬拉松.
在第一節.
我們既有這許多見證人.
如同雲彩圍著我們.
就當放下各樣的重擔.
脫去纏累我們的罪.
全心忍耐.
不拿擺在我們前頭的路程.
仰望為我們信心.
創始成中的耶穌.
在這一段.
這一句主動詞就是"奔"字.
第一至第二節上半段.
是"奔"字.
奔我們人生前面的路程.
其實路程這個字就是指競賽.
就是指賽跑.
所以這段有很多運動的圖像.
人生就像賽跑一樣.
但是怎麼跑呢?.
作者用了三個分詞.
連於"奔"字.
去表達我們怎麼奔走人生的路程.
這三個分詞.
第一個就是"有"字.
就是第一節.
我們有許多見證人.
第二個就是放下或者脫去.
其實原文只有一個字.
其實原文的寫法就是.
我們脫去各樣的重擔.
以及容易纏累我們的罪.
就是這個字.
放下或者脫去.
第三個就是仰望.
這三個分詞就表達.

$^{201}$怎麼去奔走我們人生的路呢?.
我們看第一個.
我們有許多見證人.
都明顯是指上一章.
就是上個星期講完所講的.
他的信心偉人.
這些如同雲彩圍繞我們的信心偉人.
是我們人生沿途的啦啦隊.
是鼓勵我們什麼呢?.
有信心.
你走人生路.
有信心和沒有信心.
有很大分別.
我們走人生路遇到困難.
最重要是有信心.
我們不要單單.
我信耶穌.
在在應該.
全部都應該信耶穌.
我相信.
不要單單只是信耶穌.
我們要信得過耶穌.
這個和單單信耶穌是有分別的.
要信得過耶穌.
遇到困難.
信耶穌和不信耶穌.
是有分別的.
這是明顯的.
大家同樣是基督徒.
有信心和沒有信心.
分別明顯很大.
同樣的困難.
有信心和沒有信心.
有很大的差別.
我們眼前有很多困局.
但是要憑信心.
不要憑眼見.
我們緊緊跟從這一位.
我們看不到的.
這位無形的領航者.

$^{241}$耶穌基督.
信心是第一件事.
我們要好好把握去抓住.
第二個分詞就是放下或者脫去.
這個詞是過去不定式.
就是表達我們事前要準備的.
這個不是在跑才準備的.
事前要已經預備好放下和脫去.
我們準備好要輕省上路.
就是用輕省來表達.
有沒有人跑步會背著重磅?.
有沒有人跑一百公尺.
你會說背著大背囊?.
不會的.
你會不會向毅行者說.
背著幾本書在背後?.
不會的.
你是會輕省上路的.
我記得我當年走這個毅行者.
真的最後上到大帽山頂.
有支援隊來.
我有點驚訝.
有支援隊過來帶了蛋糕來.
跟一個隊員慶祝生日.
又切蛋糕又喝又吃.
你想走得快也挺難的.
輕省.
我們去這些長途賽.
要輕省才走得快.
結果我們用了四十多個小時.
才能夠完成.
再過幾個小時就會被淘汰.
因為毅行者最長的時間是四十八小時.
你過了就不會當你達標.
你要走得遠.
要走得輕省.
經常都是這句.
Travel far.
你要Travel light.
你走得遠一定要輕省.

$^{281}$輕省是什麼呢?.
是代表要脫去.
最重要就是脫去罪.
什麼都不重要.
最重要是脫去罪.
最前類我們走一個美好人生的就是罪.
聽聽姐妹.
現在什麼是你的重擔呢?.
或者我這樣問.
現在什麼罪是阻礙你跟從耶穌的美好人生呢?.
想清楚.
如果你想很久也沒有呢?.
只有兩個可能性.
一是你是一個聖人.
你跟耶穌一樣真的沒有罪.
二就是我們說謊.
如果我們說自己沒有罪.
真理不在我們裡面.
我們是說謊的.
你試試安靜坐下來半個小時.
想想這個月你有什麼.
有沒有想過一些壞事?.
有沒有說過一些不應該說的話?.
有沒有做過一些不應該做的行動?.
一定有.
有沒有憤恨?.
有沒有冤道?.
有沒有不滿?.
有沒有妒忌?.
一定有.
會阻礙我們走美好的人生.
作者補充了一件事.
就是要忍耐.
忍耐這個字很特別.
在原文是一個複合字.
我們叫做 "hupomone" 這個字.
這個字有兩個部分組成.
第一個 "hupo" 英文就是 "under".
就是在之下.
"manow" 這個字其實是 "remain".

$^{321}$就是維持在.
所以直譯這個字就是 "remain under".
中文就是 "維持在之下".
所以忍耐這個字的意思就是.
保持在某種狀況之下.
就是你能夠活在一個困難逆境下的能耐.
忍耐是一種仰望神.
我們等待上帝出手.
將盼望放在神的心態.
忍耐不是很被動無奈地.
我們屈服在一個逆境.
忍耐是主動地憑信心仰望神的介入.
這個就是忍耐.
當年我們行藝人者.
走到第十段.
就是最後一段.
也是最長的一段.
有十多公里的一段.
有隊員開始產生幻覺.
前面有房子.
不是前面是樹葉.
開始幻覺體力嚴重透支.
但是終點就在前面.
最後一段.
繼續還是放棄呢?.
我們全隊都是信著.
我們忍耐.
堅持下去.
最後雖然是四十多個小時.
也完成了整段的路程.
人生是馬拉松.
是長途賽.
不是說速度.
是說耐力.
信心.
輕散.
以及忍耐.
我認為是人生勝利組的三大要素.
如果你想說一個得勝的人生.
這三樣東西你不可或缺.

$^{361}$缺一不可.
第三個分詞就是仰望.
這個分詞是現在時態.
我們不斷要仰望耶穌.
這個不能停.
還有在原文.
耶穌這個字是放在最後一句.
是強調的意思.
我們不斷仰望那位.
不是其他.
是耶穌.
這個單單是耶穌.
不是錢,不是權力.
你今天的成就.
你美滿的家庭.
是耶穌.
這位耶穌是怎樣的呢?.
哈恩那擺在前面的喜樂.
不是,應該是仰望那位.
為我們信心創始成終的耶穌.
耶穌是我們信心的創始者.
就是我們信心的來源.
和領航者.
耶穌也是我們信心的完成者.
是靠祂.
我們的信心才能夠完全地表現出來.
所以你看看時錄.
耶穌是什麼?.
耶穌是阿拉伯.
是俄媒家.
是首先也是末後.
祂是始也是終.
祂是我們信心的創始成終者.
所以如果你當一場跑步.
耶穌是起點的司令員.
也是終點的裁判.
我們一定要看著祂來走人生的路.
為什麼要看著祂呢?.
祂有什麼特別呢?.
我們繼續看下去.

$^{401}$哈恩擺在前面的喜樂就輕看羞辱.
忍受了十字架的苦難.
便坐在神寶座的右邊.
但因受罪人這樣頂撞的.
你們要思想免得疲倦灰心.
耶穌就是人版.
我剛才說的那幾件事.
耶穌就是人版.
祂對父神有完全的信心.
祂不會比罪前淚.
祂沒有罪.
祂忍受了十字架的苦難.
也忍受了罪人的搏逆頂撞.
因為祂看到前面的喜樂.
可以坐在神寶座的右邊.
同得榮耀.
弟兄姊妹為什麼要仰望耶穌呢?.
祂就是我們的樣板.
祂不只是高歌在寶座上.
你跟隨我吧.
祂來到這個世界告訴我們.
怎樣憑信心.
怎樣脫去罪.
怎樣很輕散的.
也是很有目標.
很有方向的.
很有信心.
很有忍耐的.
走這條路.
其實我們一樣可以.
如果我們看到人生前面有喜樂.
我們有一個永生的福樂.
相對於地上的那些.
永生的福樂是完全重要很多的.
當我們看到永生的福樂的時候.
我們可以跟隨耶穌.
走這條人生的路.
如果我們常常能夠仰望思想耶穌.
就不會疲倦灰心.
在亞歷士多德的文獻裡面.

$^{441}$他用疲倦灰心去形容什麼呢?.
就是賽跑的人.
到了終點體力不支.
倒下.
就跌下了.
就是疲倦灰心所描述.
我們可以不疲倦灰心.
因為我們仰望耶穌.
真的會得力.
我相信大家總有這些經驗.
在你人生很困難的時候.
你仰望耶穌.
你會得力.
有一次有一對夫婦在沙灘漫步.
不是我和珍珍.
這是一個比喻.
看不到很遠.
有棵大樹.
老公指著那棵樹跟老婆說.
你看到那棵大樹.
我們比賽就向那棵樹走.
太太抗議.
你是男孩子運動健將.
體能比我好.
你當然要快一點.
老公說.
不是鬥快慢.
是鬥誰跑得最直到那棵大樹.
太太明白.
就一起開始起步.
丈夫就大踏步向那棵樹走.
太太就小心翼翼.
每踏一步就看後面.
是不是直線.
最重要保持直線.
一步一腳印就去到大樹.
當然是老公先.
太太不驚訝.
老公先去到不驚訝.
不是鬥快.

$^{481}$是鬥直.
她看自己那條線.
太太看.
多直.
她再看老公那條線.
更加直.
她很驚訝.
為什麼會這樣呢?.
老公就解釋.
這棵樹是我的目標.
我只是看著這個目標.
完全沒有其他雜念.
我直去.
我不會看下面走得怎樣.
我直去.
走出最短最直的路.
這是一個很有意思的預言.
我們看著是耶穌.
你就走到最直最好的人生路.
你不要看其他.
不要被其他東西偷看.
你的錢,工作.
不是那些不重要.
用耶穌的眼光去管理你的錢,工作,家庭,人際關係.
你人生的夢想等等.
兄弟姐妹.
我們仰望著耶穌.
耶穌或者祂所給我們的使命.
其實就是我們的目標.
我們朝著這個目標進發.
就能夠走出人生的正路.
好像保羅說什麼呢?.
我只有一件事.
就是放下攝影機背後.
努力面前.
向著標桿直跑.
仰望著耶穌.
你有能力走出人生的正路.
但在人生的路上.
你一定有困難.

$^{521}$因為長途賽.
當然有些地方你是很想放棄.
或者很困難.
很不容易.
所以我們人生同樣有逆境有苦難.
作者跟著四到十一節.
鼓勵我們嘗試用一個角度去看.
將這些看為神的管教.
我們接受神的管教.
人生很不同.
因為神的管教對我們是怎樣的呢?.
一定有益.
神從來不會去戲弄祂的子民.
祂一定對我們有益.
我們看第四到十一節.
"你們與罪惡相爭,還沒有抵擋到流血的地步,你們又忘了那勸你們如同勸兒子的話.".
這裡已經有一點.
他轉了一個圖像.
他從一個賽跑.
轉到一個鬥拳的圖像.
就像保羅在臨前也說.
我奔跑不像無定向的.
我打拳不像打空氣.
從奔跑變成打拳.
作者這裡寫有一點諷刺意味.
如果你看下去.
與耶穌相比.
你們遠遠不及耶穌那麼艱難.
不用抵擋到流血的地步.
然後又說你們又忘記了那些勸你們的話.
忘記這個字在新約.
只是這裡出現了這一次.
就是這個忘記這個字.
不是我們一般用來表達忘記的希臘文.
這個忘記是一個比較強的字.
他表示什麼呢?.
你完全忘記完全都忘記了.
所以說作者好像在諷刺這群讀者.
你們也不是那麼辛苦而已.
說過你們又不記得.

$^{561}$在諷刺.
說過什麼呢?.
他說:"我兒,你不可輕看主的管教,.
被他責備的時候也不可灰心,.
因為主所愛的,他必管教,.
又邊打帆所收納的兒子..
你們所忍受的是神管教你們,.
待你們如同兒子,.
焉有兒子不被父親管教的?.
管教源是眾子所共受的,.
你們若不受管教就是私子,.
就是私生子,不是兒子.".
作者引用了《針言》三章十一,十二節.
神愛我們,先管教我們.
神當我們是親生子女,先管教我們.
所以管教是一個愛的記號.
他愛才管教.
也是一個親生子女的身份.
大家很容易明白.
你愛自己的子女.
和愛主日學的小朋友.
和愛你在街上碰到的小朋友.
有沒有分別?.
很大分別,是不是?.
你愛你的子女和愛街上碰到的小朋友一定不同.
所以如果神當我們是子女.
去愛我們.
我們要很清楚這是一個愛的記號.
我們是祂的親生子女.
頂多是自己的寶貝.
所以你應該很珍惜神的管教.
面對神的管教.
作者說不要輕看和灰心.
輕看我嘗試了一個.
現代的描述.
就是當神沒有到.
神管教你當祂沒有到.
我們想起今天我們的疫情.
一顆我們肉眼看不到的病毒.
搞到我們全世界都翻天覆地.

$^{601}$大家都明白了.
但如果我們將疫情看為管教.
疫情提醒我們什麼?.
不要驕傲.
我們不要以為人類很厲害很有能力.
一顆小小的病毒我們都搞不定.
又封城又吃不了東西又聚不了會.
我們謙卑一點看.
我們的東西很有限.
神正在用這個疫情去提醒我們.
我們很有限.
很有限.
我們不要浪費了疫情的教訓.
更加不要當神沒有到.
所以在疫情期間我常常都在想.
究竟神想教我什麼呢?.
在當中你有什麼領袖學習和應用呢?.
灰心就是對前景沒有盼望.
現在面對兩場大型社運之後.
很多人都很灰心很失望.
移民走了.
神是任何政權之上的掌權者.
給我們盼望的是神.
不是任何的制度和權力.
所以我們不用灰心.
坦白說在任何地上的權柄之上就是上帝.
對於神的管教.
我們不要輕看和灰心.
因為我們是神的寶貝兒子寶貝女.
我們是神所愛的.
如果你持著這個心態.
很認真去對待這個身份.
你人生就會很輕散.
面對這些管教這些逆境苦難.
你會輕散很多.
再者我們曾有生生的父管教我們.
我們尚且敬重他.
何況萬靈的父.
我們豈不更當信服他得生麼?.
作者還有第二個原因.

$^{641}$就是他比較肉身的父和萬靈的父.
萬靈就是有更大權柄.
有完全權柄的神.
生生的父我們尚且敬重.
何況是神.
我們要信服神的管教而得生命.
這裡有「曾有」這個字.
就是表達這個情況要結束了.
我們要轉向接受神的管教.
為什麼要這樣呢?.
為什麼要轉向給神的管教呢?.
因為生生的父都是暫隨己意管教我們.
唯有萬靈的父管教我們.
是要我們得益處.
使我們在他的聖潔上有份.
肉身的父是暫隨己意去管教.
一定有瑕疵的.
但是神是完全的.
他的智慧更加高.
他知道怎樣管教我們是最好的.
你想想這個世界誰最明白你呢?.
不是你的配偶.
不是你的家人.
是神.
祂創造你的.
祂當然明白你.
但是沒有人明白你多過其他人.
是上帝的.
最明白你,最愛你的上帝.
祂給你的管教一定是最好的.
神管教我們不是要我們讀書好一點.
找工作好一點.
賺錢多一點.
做事順利一點.
不是這些.
而是要我們在他的聖潔上有份.
目的就是在他的聖潔上有份.
所以我們今天查這段經文.
下一節第十四節.
他說到我們要追求聖潔.

$^{681}$非聖潔.
沒有人能見主.
其實是在繼續延續這個思路.
弟兄姊妹.
我們既然接受人的管教.
我們更加要接受神的管教.
因為神的管教是聖於人的管教.
從這兩節第九第十節.
我們得出三個原因.
第一.
神是萬靈的父.
祂有更大的權柄.
第二.
神比人有更高的智慧.
第三.
神懷著更崇高的目的.
去管教我們.
當我們理解.
祂的管教這麼好的時候.
我們就很甘心情願.
服祂的管教之下.
接著還沒完的.
凡管教的事.
當時不覺得快樂.
反覺得愁苦.
後來卻為那經過經練的人.
結出平安的果子.
就是義.
管教當時不會覺得很快樂.
這個大家都應該有很多經歷.
或者你的子女有很多經歷.
不覺得很快樂.
被人騰條去打.
當然有辛苦的.
當然覺得愁苦的.
但是經過鍛鍊.
有益處.
有平安.
有義.
義就是對.

$^{721}$就是能夠明辨.
什麼是對.
什麼是錯.
所以在希伯來書第五章十四節.
他們的心竅.
集練得通達.
就能分辨很大.
經過鍛鍊.
操練.
我們知道什麼是對.
什麼是錯.
懂得如何走人生的路.
所以神的管教.
如果你看完這一段.
你會發覺.
是對我們有益.
是有百利而無一害.
所以我們要很小心進入上帝的管教.
在人生路上一定有的.
我想起一個很經典的例子.
就是在神簡宣以色列人.
以色列人有十二支派.
十二支派就是亞國的十二個兒子.
他們有兩個孫子.
升級成為兩個兒子.
我就想起他們的經歷.
如果你看《創世記》三十七章.
到四十四五章.
其實講了經歷是什麼呢?.
就是亞國的十個大兒子.
從劉變大哥一直數下去.
對於約瑟第十一個兒子.
很不滿.
原因是什麼呢?.
因為這個人經常背著.
在父親面前數哥哥的不是.
在哥哥的心目中.
這是臥底奸細.
經常在他身邊戳我們背.
以及父親又很偏心.

$^{761}$所以他們很生氣.
生氣得想殺了他.
有一次他們去北邊放羊.
亞國就找約瑟去探望哥哥.
在北邊.
山高皇帝遠.
父親不在身邊.
這還不是好機會下手.
遠遠看到約瑟來.
這個發夢的來了.
要殺了他.
但是大哥劉變比較好心.
不如我們先把他關在坑裡.
暫時不要殺他.
想救他.
但是關在坑裡.
然後他走開了.
然後猶大就建議.
不如把他賣了在埃及.
剛好有一群二十馬里的商販.
賣了埃及.
就賣了約瑟給商販.
不見了弟弟.
就拿回錢.
約瑟就去到波提弗.
一直發展發展.
他做了埃及的宰相.
大家知道大饑荒.
亞國就派十個大兒子.
去埃及擲糧.
這件事過了二十二年.
亞國約瑟.
從十七歲到三十九歲.
他的哥哥不認得他.
他認得這位哥哥.
如果你看第四十二章.
《創世記》.
你會很深印象.
很有趣的.
這位哥哥看到他.

$^{801}$他冤枉他.
你是奸細.
是窺探這個地的.
為什麼誣衊他是奸細?.
他明知道他們不是奸細.
他想他記得他二十二年前.
出賣兄弟的惡行.
接著就將他收監.
是不是?.
你將他收在坑裡.
現在將他收監.
他們真的記得了.
終於記得二十二年前.
怎樣害這個弟弟.
所以劉變就說.
當時提醒了大家.
不要害他.
現在他的血在追討我們.
接著當他們拿了糧食.
回去的時候.
約瑟偷偷找人.
將錢放回他們的行囊裡.
他們說為什麼我們的行囊有錢呢?.
約瑟再一次.
繼續提醒他們.
當年你失去一個兄弟.
只剩下錢.
這次你又失去一個兄弟.
又只剩下錢.
這次失的是西面.
約瑟將西面扣留在埃及.
你們回去.
帶著你的小弟變亞斌過來.
我就相信你說的是真的.
你不帶過來.
我就確認你們是奸細.
後來帶了變亞斌過來.
約瑟繼續測試他們.
又給變亞斌五倍食物.
當年你這麼討厭拉傑的兒子約瑟.

$^{841}$你會不會討厭拉傑的第二個兒子.
變亞斌呢?.
他們沒有.
最後的測試就是.
將變亞斌留在埃及.
你回去.
整家人帶過來就行了.
不過變亞斌要留在埃及.
不行啊.
變亞斌留在埃及.
爸爸死定了.
這兩個生命相連.
白頭人會悲悲慘慘下陰間.
搞這麼多事情.
其實做什麼呢?.
在這麼多事情上.
當年這十個哥哥.
在出賣兄弟這件事上跌倒了.
然後神就管教他們.
他用約瑟的管教.
一而再再而三.
三番四次去測試他們.
在出賣兄弟這件事上.
你們會不會再跌倒.
他們都經過了.
他們沒有出賣西面.
他們沒有妒忌變亞斌.
最後也沒有出賣變亞斌.
他們重新站起來.
所以神的管教是有益的.
這十兄弟.
當然加上最後變亞斌.
和約瑟兩個兒子.
以法連嗎哪西.
成為十二支派的始祖.
所以如果這樣看.
上帝的管教.
可能我們一開始不覺得.
怎麼管教呢?.
但其實對我們很有益.

$^{881}$很有幫助.
所以我們要忍耐.
繼續有信心.
不要在管教裡面犯罪.
到了這段落的最後一段.
第十二和十三節.
第十二節.
所以就像一個結論.
我們要把下垂的手.
髮宣的腿挺起來.
這裡應該引自.
《以賽亞書》三十五章第三節.
當你看到那裡.
是在說跟從神的人.
怎麼走這條神聖的路.
和這裡的內容很吻合.
這裡又回到一個運動的圖像.
你要挺起手腳.
就是剛強起來.
準備啟動.
準備去比賽.
第十三節.
也要為自己的腳把道路修直了.
使騎子不自歪腳.
反得全愚.
第十三節的道路.
和第一節的奔.
是首尾呼應.
表達我們每一個人.
都要跑這條人生的馬拉松.
要把道路修直.
就是要直跑.
不要搖搖擺擺.
不要三心兩意.
一心一意.
跟從耶穌.
仰望著那位耶穌基督.
你能夠走得好.
可以幫助一些騎子.
就是騎腿的人.

$^{921}$代表一些軟弱的人.
不自歪腳.
歪腳的原文就是脫教脫教.
但意思就是走差路.
走偏路.
當我們剛強地跟從神.
去走這條人生的馬拉松.
同時可以鼓勵.
幫助一些軟弱的人.
不偏離正路.
你會成為別人的榜樣.
幫助別人去走.
所以是自助助人.
利人利己.
人生是一場長途賽.
重點是要達標.
而不是首先達標.
這是很重要的.
我們不是鬥快.
不是看速度.
是看你有沒有信心.
忍耐和輕散.
最重要的是.
你是否仰望著耶穌.
來走這條人生的路呢?.
你如何操練你的信心.
忍耐和輕散.
和仰望耶穌呢?.
沒有什麼捷徑.
我猜我應該也說過了.
其實我去很多不同的地方.
都是說三聖的操練.
我有沒有說過?.
三聖的操練.
三聖.
第一就是聖經.
這是你不能避開的.
好好研讀操練.
跟從聖經.
一個人不是常常去研讀聖經.

$^{961}$是很困難的.
第二就是聖靈.
不要不理會聖靈的提醒.
督責令你的掙扎.
多些安靜.
多些禱告.
勉於聖靈在心裡.
給你的感動.
第三就是聖徒.
不要不回皇城.
不要不回教會.
不要沒有了聖徒相交的生活.
回教會不等於有聖徒相交的生活.
有些人很孤零零.
回來坐下來.
崇拜就走了.
不跟別人相交.
不說話.
不回團契.
如果你在聖經,聖靈,聖徒.
這三聖的操練做得好.
我也挺相信.
你走這條人生路.
你是看著耶穌.
跟著耶穌走.
記得耶穌是起點線的司令員.
又是跟你一起同跑的啦啦隊.
還有在終點的裁判就是他.
這個無所不在的耶穌.
在你的人生路上至關重要.
在跑的時候.
能夠有信心.
有輕省,有忍耐.
很多時候這三聖是很能幫助你的.
而且這三聖不是捷徑.
是你一生都要做的.
不會畢業的.
做到你見主面.
但你會越來越開心.
越來越喜樂.

$^{1001}$越來越平安.
這條人生馬拉松.
你要走的時候.
要跑的時候.
一定有困難.
有逆境.
甚至苦難.
你想不想徒然受苦?.
大家都不想吧?.
經歷苦難.
浪費了他.
浪費了上帝的一番苦心.
你試試轉念吧.
這就是神對我的管教.
多些問.
我現在經常禱告.
多些問這個問題.
神,你要我學什麼功課?.
當我面對困難的時候.
神,你要我學什麼功課?.
No pain, no gain.
通常有困難.
你學得最多.
我們常常說.
我們要贏在起跑線.
這個說法很奇怪.
決定一場比賽的勝負.
在起點還是終點?.
沒理由在起點的.
什麼贏在起跑線呢?.
就是想早點起步.
偷步吧.
決定你人生的成敗得失.
在你人生的終點.
當我們走完人生的路程.
主耶穌會在終點.
評價我們整個人生的成敗得失.
我們要贏在終點線.
你想贏在終點線.
可以好好用心學習這段經文所說的一些原則.

$^{1041}$和真理.
我相信對你的人生.
是有很大的幫助.
我們同心同告.
親愛的主,我們再一次去感激你.
你願意去啟動我們.
你賜予我們信心.
你帶領我們去相信你,認識你.
你與我們同行.
你天天與我們同在.
你又在終點等我們.
我們很想奔向你的懷抱.
我們很想走一條你所喜悅的人生正路.
求主幫助我們.
對你有信心.
幫助我們很堅定地去對付罪.
幫助我們懷著忍耐.
仰望著你去走我們的人生路.
縱然在中間有很多的苦難.
有很多逆境.
但我們相信.
你給我們的考驗.
同樣地,你給我們有出路.
我們再一次去感激你.
求你幫助.
帶領我們成為信心的創始成就者.
我們同心圖告仰望.
奉耶穌基督的聖名.
阿們.
\newpage



\section{箴言 31:1-31}
\label{sec:uzCPC0kQbAA}
\textbf{2022年5月29日崇拜直播｜何羅乃萱師母｜疫情下才德婦人的啟示｜箴言31\_1-31}
\newline
\newline
連結: \href{https://youtube.com/watch?v=uzCPC0kQbAA}{\texttt{https://youtube.com/watch?v=uzCPC0kQbAA}} ~~~~ 語音日期: 2022-05-30
\newline
\newline
\hyperref[sec:qTE37kvGc24]{\small{< < < PREV SERMON < < <}}
~
\hyperref[sec:index_chronic]{\small{[返順時目]}}
~
\hyperref[sec:e5I3VklFIWg]{\small{> > > NEXT SERMON > > >}}
\newline
\newline
箴言 31:1-31
\newline
\begin{longtable}{cl}
\hline
\hline
章節 & 經文 (和合本修訂版)\\
\hline
31:1 & \begin{tabularx}{0.7\textwidth}{X} 瑪撒王利慕伊勒的言語，就是他母親教導他的。 \end{tabularx} \\ \\ \relax
31:2 & \begin{tabularx}{0.7\textwidth}{X} 我兒，怎麼了？我腹中生的兒，怎麼了？我許願而得的兒，怎麼了？ \end{tabularx} \\ \\ \relax
31:3 & \begin{tabularx}{0.7\textwidth}{X} 不要將你的精力給婦女，也不要有敗壞君王的行為。 \end{tabularx} \\ \\ \relax
31:4 & \begin{tabularx}{0.7\textwidth}{X} 利慕伊勒啊，君王不宜，君王不宜喝酒，王子尋找烈酒也不相宜； \end{tabularx} \\ \\ \relax
31:5 & \begin{tabularx}{0.7\textwidth}{X} 恐怕喝了就忘記所頒的法令，顛倒所有困苦人的是非。 \end{tabularx} \\ \\ \relax
31:6 & \begin{tabularx}{0.7\textwidth}{X} 可以把烈酒給將亡的人喝，把酒給心裡愁苦的人喝， \end{tabularx} \\ \\ \relax
31:7 & \begin{tabularx}{0.7\textwidth}{X} 讓他喝了，就忘記他的貧窮，不再記得他的苦楚。 \end{tabularx} \\ \\ \relax
31:8 & \begin{tabularx}{0.7\textwidth}{X} 你當為不能自辯的人開口，為所有孤獨無助者伸冤。 \end{tabularx} \\ \\ \relax
31:9 & \begin{tabularx}{0.7\textwidth}{X} 你當開口按公義判斷，當為困苦和貧窮的人辯護。 \end{tabularx} \\ \\ \relax
31:10 & \begin{tabularx}{0.7\textwidth}{X} 才德的婦人誰能得著呢？她的價值遠勝過寶石。 \end{tabularx} \\ \\ \relax
31:11 & \begin{tabularx}{0.7\textwidth}{X} 她丈夫心裡信賴她，必不缺少利益； \end{tabularx} \\ \\ \relax
31:12 & \begin{tabularx}{0.7\textwidth}{X} 她終其一生，使丈夫有益無損。 \end{tabularx} \\ \\ \relax
31:13 & \begin{tabularx}{0.7\textwidth}{X} 她尋找羊毛和麻，歡喜用手做工。 \end{tabularx} \\ \\ \relax
31:14 & \begin{tabularx}{0.7\textwidth}{X} 她好像商船，從遠方運來糧食， \end{tabularx} \\ \\ \relax
31:15 & \begin{tabularx}{0.7\textwidth}{X} 未到黎明她就起來，把食物分給家中的人，將當做的工分派女僕。 \end{tabularx} \\ \\ \relax
31:16 & \begin{tabularx}{0.7\textwidth}{X} 她想得田地，就去買來，用手中的成果栽葡萄園。 \end{tabularx} \\ \\ \relax
31:17 & \begin{tabularx}{0.7\textwidth}{X} 她以能力束腰，使膀臂有力。 \end{tabularx} \\ \\ \relax
31:18 & \begin{tabularx}{0.7\textwidth}{X} 她覺得自己獲利不錯，她的燈終夜不滅。 \end{tabularx} \\ \\ \relax
31:19 & \begin{tabularx}{0.7\textwidth}{X} 她伸手拿捲線桿，她的手掌把住紡車。 \end{tabularx} \\ \\ \relax
31:20 & \begin{tabularx}{0.7\textwidth}{X} 她張手賙濟困苦人，伸手幫助貧窮人。 \end{tabularx} \\ \\ \relax
31:21 & \begin{tabularx}{0.7\textwidth}{X} 她不因下雪為家裡的人擔心，因為全家都穿上朱紅衣服。 \end{tabularx} \\ \\ \relax
31:22 & \begin{tabularx}{0.7\textwidth}{X} 她為自己製作被單，她的衣服是細麻和紫色布做的。 \end{tabularx} \\ \\ \relax
31:23 & \begin{tabularx}{0.7\textwidth}{X} 她丈夫在城門口與本地的長老同坐，為人所認識。 \end{tabularx} \\ \\ \relax
31:24 & \begin{tabularx}{0.7\textwidth}{X} 她做細麻布衣裳來賣，又將腰帶賣給商家。 \end{tabularx} \\ \\ \relax
31:25 & \begin{tabularx}{0.7\textwidth}{X} 能力和威儀是她的衣服，她想到日後的景況就喜笑。 \end{tabularx} \\ \\ \relax
31:26 & \begin{tabularx}{0.7\textwidth}{X} 她開口就發智慧，她舌上有仁慈的教誨。 \end{tabularx} \\ \\ \relax
31:27 & \begin{tabularx}{0.7\textwidth}{X} 她管理家務，並不吃閒飯。 \end{tabularx} \\ \\ \relax
31:28 & \begin{tabularx}{0.7\textwidth}{X} 她的兒女起來稱她有福，她的丈夫也稱讚她： \end{tabularx} \\ \\ \relax
31:29 & \begin{tabularx}{0.7\textwidth}{X} 「才德的女子很多，惟獨你超過一切。」 \end{tabularx} \\ \\ \relax
31:30 & \begin{tabularx}{0.7\textwidth}{X} 魅力是虛假的，美貌是虛浮的；惟敬畏耶和華的婦女必得稱讚。 \end{tabularx} \\ \\ \relax
31:31 & \begin{tabularx}{0.7\textwidth}{X} 她手中的成果你們要賞給她，願她的工作在城門口榮耀她。 \end{tabularx} \\ \\
[1ex]
\hline
\hline
\end{longtable}
$^{1}$我們今天要聽神的話是記載在.
《箴言》第31章十至到三十一節.
請聆聽神的話.
《箴言》三十一章第十節.
「才德的婦人誰能得著呢?.
她的價值遠勝過珍珠.
她丈夫心裡已靠她.
必不缺少利益.
她一生之丈夫有益無損.
則尋找羊肉和麻甘心用手作工.
她好像相信從遠方運糧來.
未到黎明她就起來.
把食物分給家中的人.
將當初的工分派皮類.
她想得田地就買來.
用手所得之利栽種葡萄園.
她以能力蓄養.
使旁被有力.
她覺得所經營的有利.
她的燈中夜不滅.
得手拿撚扇桿.
手把放扇節斜.
她將守周濟困苦人.
伸手幫助窮困婦人.
她不忍下雪為家裡的人擔心.
因為全家都穿著珠紅衣服.
她為自己製作繡衣,毯子.
她的衣服是細麻和紫色布作的.
她丈夫在城門口與本地的長老同坐.
為眾人所認識.
她作細麻布衣常出賣.
又將腰帶賣於商家.
能力和威而是她的衣服.
她想日後的警方就會起笑.
她開口就發自慰.
她事實上有人事的法則.
但觀察家務並不吃閒話.
她的兒女起來稱她有福.
她的丈夫也稱讚她.
說「才德的女子很多.

$^{41}$唯獨你超過一切」.
「嚴厲是虛假的.
美容是虛浮的」.
為敬畏惹惑的婦女便得稱讚.
願她享受操著所得的.
願她的工作在城門口榮耀她.
成議會續畢.
很高興今天有我們熟悉的.
家庭發展基金總幹事.
何羅乃香之母在我們當中分享.
她今天要宣講的題目是.
疫情下才德婦人的啟示.
我們恭敬她.
節目.
各位早晨.
來到這裡好像回到另外一個家.
在這裡.
大家可能知道你們的師母.
是我的同工.
今天還有另一個同工.
Carolyn 也來了.
也不知道怎樣形容那種心情.
覺得我九龍也有一個家.
今天還是五月.
還是母親節.
所以選擇了「疫情下才德婦人的啟示」.
在過去幾個月或一年.
我說甚麼題目.
都有這三個字「疫情下」.
因為實在還在.
她還沒走.
雖然我們祈禱.
我每天都在祈禱.
讓我擺擺她的尾巴.
但在這裡.
「疫情之下」.
有些壓力來源.
也跟大家看一下.
你還看不看「確診大廈」?.
現在.

$^{81}$有沒有看的?.
沒有吧?.
我以前不單看「確診大廈」.
我還看「確診大廈」.
五百呎那裡.
我家在附近嗎?.
現在.
尾空.
沒時間了.
工作改變.
裁員減薪.
有些失去工作.
大家知道我辦公室.
都是你們隔幾步.
很多店舖關了.
你們隔離有幾間食肆.
我喜歡的都沒有了.
沒錯.
轉變很快.
但另外在家庭.
你知道我做親子教育.
每次去不同的家.
聽到都是.
學校聽到三個字.
「困獸鬥」.
我覺得是這個角度.
你覺得混在一起.
「困獸鬥」.
但同樣有另一個角度.
叫「天倫樂」.
我教小朋友寫作.
我問小朋友.
「你在家裡看到甚麼?」.
他說從來沒有見過爸爸.
原來這麼能玩.
最好爸爸經常陪他.
原來爸爸對他們的影響都大.
當然更深的是.
當我們長期被困.
長期困在一個空間的時候.

$^{121}$觸到我們的裡面.
內在形象.
怎樣看自己呢?.
個性特質ABC.
今天不詳細說.
不過那種個性簡單說.
跟我也有點像.
急,快.
還有.
他說最不好的C個性.
你知道是甚麼嗎?.
Cancer characteristic.
容易有癌症的人.
把事情圍起來.
生氣也不生氣.
每一天都會爆發出來.
期望過高.
但我覺得最後最重要是.
信念和堅持.
為何我會選擇這篇《財德婦人》呢?.
我每年都要說母親節.
可能大家都明白.
師母就要說母親節.
就要說人物.
我從來都不比這篇.
為甚麼呢?.
因為太完美了.
我看過一個牧師.
是台灣牧師.
他寫了一個現代版.
你知道我喜歡作畫.
我寫了一個現代疫情版.
是甚麼呢?.
他老公工作步步高升.
房子越來越大.
他覺得娶了這個太太.
是他一生最正確的選擇.
同意嗎?.
那些老公.
他的興趣是找原物料.

$^{161}$有機黃豆.
這是我女兒正在做的.
定期去屈嘟嘟和慢嘟嘟買藥.
每天六點鐘起床做早餐.
所以不准吃外賣.
叫菲傭放假不要出去.
每天做運動.
生酮飲食法.
知道是甚麼嗎?.
知道嗎?.
你們試過了.
168有試過嗎?.
我試了兩個星期.
不行.
因為不能吃早餐.
不行.
怎樣都不行.
我最厲害的是最後.
家中有多款可更換的滅菌口罩.
和消毒藥水.
何牧師說你的消毒藥水.
多到可以洗澡.
不要再買了.
不過.
你知不知道最卡住我.
最初不敢說的是甚麼呢?.
是第22節.
第22節說甚麼呢?.
他為自己製作繡花毯子.
他的衣服是細毛和紫色布造的.
我最弱就是這環.
話說我和何牧師拍拖的時候.
有一天他的褲子有個洞.
他跟我說.
女朋友你幫我縫製它.
褲子有個洞怎樣縫製呢?.
都要用布去縫製.
我不是.
我將那個洞拿出來.
然後用線繞著.

$^{201}$然後給他.
他回答.
不如我自己縫製.
所以我一直對自己的縫紉.
手工.
是毫無信心的.
我覺得還是不要說了.
不要說這篇了.
很難說.
後來在疫情期間.
經常看這篇.
不知為甚麼.
經常反反覆覆地看.
我們找找背景.
給大家看背景.
前半篇.
就是利慕伊勒王.
據聞是一個外邦的王.
他的名字叫奉獻給神的意思.
媽媽對他勸戒.
我們不說這個.
後半段.
是兒子稱讚媽媽.
是不是很好?.
都挺好.
但是我再找再找.
發覺有個解經家這樣說.
他說其實.
整本《箴言》說甚麼?.
智慧.
而這本最後一章.
就用一個才德的婦人.
作為一個結束.
是拿女人做結束.
各位姐妹有沒有覺得.
與有榮焉.
覺得都.
是不是挺開心?.
一個智慧.
然後落在一個才德婦人身上.

$^{241}$這裡說.
智慧不是空洞的老生常談.
一條條的規矩.
是怡人法.
是一個女士.
智慧完全體現在她身上.
我覺得很開心.
於是開始看.
看的時候.
我又發覺.
嘩.
都很多節.
又.
總共有21節.
我看過不同的港島.
可以說10篇的.
你不用怕.
今天一篇而已.
我呢.
就選取其中一些.
我自己覺得有領受的.
可能解經.
有少許解經.
不過是一些生命的領受.
跟大家分享這四點.
才德婦人.
勝過珍珠.
第一點最重要.
丈夫信任有益無損.
持家有道未遇仇謀.
待人慷慨.
周濟群乎.
第一個.
才德婦人勝過珍珠.
這裡原來有很多解釋的.
一會兒有紅寶石.
一會兒有珊瑚.
什麼都有的.
我們要了解什麼叫做才德.
才德.

$^{281}$在聖經的解釋裡有幾個.
一.
是你的能力.
才能才幹.
效率能幹.
簡單的字.
厲害.
他厲害的.
行的.
可以的.
第二.
就是品格.
原來品格都很重要.
是品格高尚.
of noble character.
我初初以為一二都對.
誰知很多解經家用第三.
說他很勇.
Billion Roman.
為什麼呢?.
因為他工作.
職場.
家庭.
他都得勝.
他都行.
我覺得沒有所謂.
大家想想.
但最重要是什麼呢?.
很難找的.
他的價值.
勝過珊瑚.
紅寶石.
或者叫珍珠.
你喜歡選什麼都可以.
好了.
第二件事.
他愛惜自己.
我經常看這個婦人.
我覺得她應該不愛美.
因為美容是虛浮的.

$^{321}$她應該穿衣服都不理會.
不是的各位.
這裡說她為自己製作繡花毯子.
她的衣服是細麻和紫色布做的.
紫色是舊約聖經高貴.
布是漂亮的.
而她製作這些衣著打扮是什麼呢?.
我看回一個解經家.
他說.
The wife dresses herself in a way.
becoming her station.
讓人看到她是怎樣.
Avoiding the.
這個英文我嘗試去讀.
Sordid simplicity and ostentatious luxury.
意思是什麼呢?.
不是隨便的.
隨便的.
隨便拿件衣服就穿上.
再去見人.
但她又不是.
是一身名牌.
說得簡單一點.
她莊重.
我想中國人有句話很好.
「穿得莊重,恰如其分」.
我記得我剛信耶穌的時候.
我們的牧師很嚴厲.
他說「你們姐妹回來.
Sunday best」.
每次星期日回來都要看.
所以本來.
我換了雙鞋.
我本來打算穿球鞋.
但我見到.
我陳目都是皮鞋.
Sunday best.
我臨出門口回去換了雙鞋.
敬拜上帝.
Sunday best.

$^{361}$第三.
她對未來樂觀.
能力和威儀是她的衣服.
能力剛才說過是才幹.
威儀是美麗尊榮.
最有趣的是《詩篇》93篇第一節.
大家看不看得到?.
都看得到吧?.
我們現在看有點暗.
但還可以.
以威嚴為衣穿上.
以能力蓄要.
跟上帝的形容很相似.
有沒有想過?.
形容她和形容神是很相似的.
看《詩篇》.
然後.
她說.
她對未來樂觀是想到.
日後的景況就笑了.
各位姐妹.
各位弟兄.
想到2023年你現在.
會怎樣?笑不笑得出來?.
這就是那種.
我們對未來的看法.
她為什麼可以這樣?.
問題是.
為什麼她可以.
又看自己.
又看未來這麼樂觀呢?.
我有兩個看法.
第一.
是她知道自己的身份價值.
She knows who she is.
她是誰?.
她這裡說.
唯獨你超過一切.
她的身份,她的價值.
是超過一切的東西.

$^{401}$為什麼我這麼有感觸呢?.
因為做了好幾年.
很多年婦女成長課程.
在那個成長課程.
其中有一課叫「原生家庭」.
我叫姐妹寫小時候.
家裡怎樣說她.
原來很多包括我.
OK?包括我.
很自卑.
我覺得.
我出生.
只是一個過客.
因為爸爸要追兒子.
媽媽五次小產生我.
再兩次小產生我弟弟.
我覺得我是多餘的.
沒有人這樣說.
但是自己的感受.
我四歲.
我弟弟出生.
有一個親戚走過來.
跟我說.
「妹妹」.
「你爸爸追到兒子了」.
「以後沒有人疼你了」.
我不知道為什麼她說這樣.
但是這個東西.
對我整個人的自愛價值.
是一個很大的標誌.
貼在這裡.
我接觸很多姐妹.
甚至不一定是姐妹.
弟兄.
在疫情期間.
可能遭遇很多挫折.
可能你心裡會有句說話.
「我都沒用了」.
「我死定了」.
「我不行了」.

$^{441}$「老闆要解僱我」.
這些我們叫做.
自我批判.
內心的控訴.
有沒有呢?.
有.
她好像一個很頑皮的小孩.
在你心情低落的時候.
她就「咻」.
「鼠」出來.
我覺得.
我們要認定什麼呢?.
我們要認定在創造裡.
我們的身份.
上帝創造我們.
做了男人之後.
祂說什麼?.
「Good」.
記不記得?.
「都是好的」.
做了女人之後.
接著說什麼?.
「甚好」.
英文翻譯得不好.
不是.
中文翻譯得不好.
英文是「Very good」.
即是有男有女.
是「Very good」.
這個就是我們的價值.
不是隨隨便便.
被一些東西去動搖的.
這裡有三句說話.
叫做「I'm good」.
「I'm whole」.
「I'm beautiful」.
可以加上男.
弟兄可以說.
「I'm handsome」.
可以嗎?.

$^{481}$這個應該是.
差不多二十年前.
有人給了我一本書.
叫做《流奶與物之地》.
我看.
講述.
救贖.
自己的身價,價值.
這三句我很喜歡.
「I'm good」.
是嗎?.
無論如何.
上帝創造我是.
「Very good」.
是好的.
不要管.
我父母怎麼看.
甚至.
第二.
「I'm whole」.
我有沒有結婚.
有沒有生兒子.
我都是完整的.
在上帝眼中.
第三.
「I'm beautiful」.
無論我多少歲.
在上帝眼中.
「I'm beautiful」.
「I'm handsome」.
我們試試讀出來.
試試心裡說出來.
因為.
我經常說.
如果說三次.
你會記得.
弟兄就說.
「I'm handsome」.
說得到.
是嗎?.

$^{521}$女的就說.
「I'm beautiful」.
我們來吧.
「I'm good」.
「I'm whole」.
「I'm beautiful」.
我說.
你跟我說兩次.
OK?.
1,2,3.
「I'm good」.
「I'm whole」.
「I'm beautiful」.
應該有「Handsome」的聲音.
OK.
再來一次.
1,2,3.
「I'm good」.
「I'm whole」.
「I'm beautiful」.
記住它.
剛才.
我在台下.
突然間回想起一幕.
1976年.
剛剛信耶穌.
有個營會.
要找詩琴.
我就膽粗粗.
走上台做詩琴.
講者是一個很喜歡罵人的講員.
你不要理他是誰.
也不需要猜.
去到最後一場.
我還記得他的歌.
叫做.
「是假的愛」.
「滿足我心」.
「是假的愛」.
「滿足我心」.

$^{561}$我用C Major彈.
最容易的.
突然間那個講者說.
「我們要E Flat」.
詩琴在不在.
「詩琴可能可以」.
「我不行」.
「E哪裡Flat到」.
整個呆在那裡.
然後.
那個講員就說.
「我們不要詩琴」.
「請詩琴下台」.
當時我.
流著眼淚.
覺得那個服侍不被欣賞.
覺得一個很大的否定.
下台了.
但同時在那天.
我收到一封信.
是我爸爸願意信耶穌.
我跟上帝說.
「主啊,如果我爸爸願意信耶穌」.
「我一生奉獻給你」.
但又被人趕下台了.
到他呼召.
這個牧師一定有呼召.
他最後又呼召.
「上帝要低調的人才」.
「要低調的學者」.
「要低調的詩琴」.
說出來.
當時有七,八百人.
我看到周圍的人.
全部都差不多站起來了.
說奉獻.
「我出不出去?」.
「如果你是我,出不出去?」.
「出不出去?」.
「出」.

$^{601}$「陳牧,你明白我」.
「出」.
「不要緊的」.
「我出去吧」.
我走出去.
剛才數一數.
四十多年了.
1976年.
到現在.
如果當時我不出.
我不回應.
我不知道會怎樣.
但最有趣的是甚麼呢?.
我看到那場港島裡面.
有些人做了很出色的校長.
有些人做了神學家.
接著我就問他.
「你記不記得那年」.
「有個東領會叫做主阿里是誰」.
「有個詩琴被人趕下台」.
他說「不記得了」.
原來是別人不記得的.
記得的是我.
我只是一個人.
我很好.
我好.
我漂亮.
《倚創亞書》43章第四節.
我們讀這節經文.
「因我看你為寶為尊」.
「又因我愛你」.
「所以我使人代替你」.
「使你幫人替換你的生命」.
記住那四個字.
上帝看我們是甚麼.
Precious and Honor.
很寶貴.
「甚至我們放棄自己」.
「他也不會放棄我們」.
你信不信?.

$^{641}$再給前書二章九節.
我們又讀.
一二三.
「惟有你們是被揀選的族類」.
「是有君尊的祭司」.
「是聖潔的國度」.
「是屬神的子民」.
「要叫你們宣揚那照你們出黑暗」.
「入光明者的美德」.
是不是很可愛?.
有沒有覺得自己被揀選的?.
還是不被揀選的?.
「君尊的聖潔的國度」.
「召我們離開黑暗」.
「進入光明」.
很重要的.
但我們要做的是甚麼?.
為甚麼我們能夠這樣?.
有一句說話.
最近我跟何牧師聽一篇.
說得很有得著.
我們覺得勇比我們有身分.
但他反過來說.
「勇氣不給你身分」.
「你的身分給你勇氣」.
你是甚麼身分.
就給你有勇氣去做.
我知道我是誰.
沒有人能動我這個身分.
我憑著我這個身分.
是上帝君尊的祭司.
我就去做.
出黑暗,入光明.
為甚麼他能夠這樣?.
第二樣我覺得很重要.
「敬畏耶和華」.
「恐懼主」.
你有沒有查過這個世界有很多「恐懼主」?.
很奇怪.
我有一天上網.

$^{681}$查甚麼叫「恐懼主」.
你知不知道有甚麼「恐懼主」?.
「氫屬性恐懼主」.
很怕水.
「技術性恐懼主」.
我們每個人都有.
「技術性恐懼主」.
不見了電話很害怕.
還有一個叫「電子恐懼主」.
我初時以為是怕八達通.
原來不是.
是怕八字.
即是八的號碼.
會有人這樣害怕.
你會問.
有沒有「新冠肺炎」?.
有沒有?.
有.
病徵暫時未列出.
但現在正在研究.
一定有.
我們會很害怕.
很多的擔心.
但這裡說的不是這個.
是「恐懼主」.
我記得我每次講道都說這句.
甚麼叫「敬畏耶和華」?.
Oxford Chambers 說過一句說話.
我覺得是我一生的一個Motto.
他說甚麼呢?.
If you fear God, you have nothing to fear.
廣東話翻譯出來是.
如果你敬畏神,你就無憂怕.
同意嗎?.
那就說一句.
跟我一起說一次.
1,2,3.
如果你敬畏神,你就甚麼?.
無憂怕.
跟隔壁說一次.

$^{721}$是嗎?.
如果我們敬畏神,我們怎樣?.
無.
憂.
怕.
Amen.
艷麗虛假,美容虛浮.
為敬畏耶和華的婦女.
得著稱讚.
我剛才走來的時候.
我都說我仍然在學甚麼叫敬畏神.
敬畏神是言讀聖經.
屬讀神的話.
跟上帝親近.
《箴言》第九章第十節.
我們一起讀.
1,2,3.
敬畏耶和華是智慧的開端.
認識至聖者便聰明.
很熟悉的.
看聖經.
靈修.
但我給一個浪費的.
另一邊.
大家想想有甚麼.
阻擋我們認識耶和華的壞習慣.
我發覺.
我生命或我靈性有一個很大的轉捩點.
梓梅你會明白,弟兄不明白.
我戒了講電話.
甚麼叫戒了講電話?.
我以前很喜歡講電話.
我曾經創下.
晚上十二時講到第二天早上九時的紀錄.
在我中學的時候.
不講這筆了.
我覺得.
跟姐妹,好朋友講電話很高興.
一有不開心.
「哎,A」.

$^{761}$「你怎樣?」.
「是嗎?你真的這樣?」.
「是,A不好」.
「打去給B」.
「這樣這樣」.
但在差不多十多年前.
我忘記是十多還是二十年前.
有一天靈修.
突然間有一個意念戒掉了.
戒不戒得到?.
戒得到的.
因為戒掉了.
你猜好處是甚麼?.
你專心仰賴上帝的話.
不再將自己的苦水不停地倒.
不停地講.
所以今天挑戰大家.
除了熟讀聖經之外.
我們求聖靈告訴我們.
有沒有壞習慣.
影響我們跟上帝之間的交往.
而使我們不能熟讀.
不能謹記上帝的話.
第二,敬畏神是凡事以榮耀神為首位.
這裡二創亞書.
「這些百姓」即是我們.
「是述說上帝的美德」.
這裡說耶穌也說.
「在地上榮耀上主」.
你今天選的歌很好.
很喜歡這首歌.
兩首歌我都很喜歡.
我要信服.
忠心.
繼續信服上帝的大靈.
信服上帝的指引.
《羅馬書》十二章二節.
和《馬太福音》六章十節.
我們一起讀好不好.
一二三.

$^{801}$「不要效法這個世界.
只要心意更新而變化.
叫你們察驗何謂神的善良.
純存可喜悅的旨意.
願你的國降臨.
願你的旨意行在地上.
如同行在天上」.
「願你的旨意行在地上」.
基本上是每天.
靈修完禱告最後一句.
主啊.
願你的心意行在我們生命.
行在我們香港的地方.
行在我的家庭.
行在我周圍的人.
所愛的人身上.
右邊很好.
它說「Obedient」.
「信服」是一個選擇.
是我們其實很有限.
我們完全看不明白.
上帝會這樣工作.
你有沒有這樣的感覺.
很多時候上帝介入我們的生命.
我們想都想不到.
我在過去這一年.
或者這幾個月開始.
每天靈修之後.
做了一個祈禱清單.
祈禱清單裡面.
家人所愛的人一定每天祈禱.
每天起碼有五個.
我現在加到十個的人.
為他們家庭.
為他們禱告.
然後到心裡有感動.
覺得他們可能會有事.
我就發訊息給他們.
你發覺到原來他們真的需要.
話說有一天.

$^{841}$我收到一個很緊急求救的.
Facebook message.
師母我一定要見你.
我家請你打電話給我.
但不知為何看完之後.
不見了那張東西.
我找不回.
我童工沒有看過.
因為它落到我個人的.
不認識的人那裡.
我找不回.
我覺得這麼緊急.
它又說有個牧師要找我.
我找不到.
我想隔了幾個月.
有一天又會彈出來.
彈出來後我立刻問它.
你怎樣怎樣.
接著它說.
其實何牧師會認識我老公.
喂?.
它說它是三十年前.
幫我老公收針的.
不過現在我們都沒有回教會.
結果我做了些甚麼.
我和何牧約了他們吃飯.
剛剛吃完.
遲些我們希望他們回教會.
帶他們回教會.
上帝會都很奇妙的帶領.
它說我初初以為上帝不理我.
任由我.
我說不是的.
上帝會盯著你的.
上帝珍惜你.
上帝重視你.
上帝愛你的.
祂不會扔你自利於不顧的.
再下去.
敬畏神就是經歷神.

$^{881}$但在你有赦罪之恩.
叫人敬畏你.
下面馬太福音第六章二十六節.
我們一起看.
讀一讀.
一二三.
你們看天空的飛鳥.
它們既不中也不收.
又不收積在倉裡.
你們的天父尚且養活它們.
你們不比它們貴重.
不知道你們看這句有沒有感覺.
這句話是我最近.
聽到一個見證.
一個弟兄告訴我.
它為何會迴轉.
這個弟兄.
幾十年前信了耶穌.
很熱心.
帶著家裡.
帶著妹妹信.
接著他回大陸工作.
沒有再去教會.
完全沒有.
別人跟他說不理.
他覺得不用.
到後來.
他的人生遭遇了很大變化.
你不要想什麼變化.
你也猜到.
然後.
有一天他回到香港探望妹妹.
在妹妹的房間裡.
無意中看到這節經文.
你們不比它們貴重.
他感受到.
上帝還是很在乎他.
他當時跪在妹妹的露台.
痛哭流涕.
然後跟神說.

$^{921}$我要迴轉.
那天.
他將這個故事告訴我.
師母.
你要跟弟兄姊妹說.
不要覺得上帝不理.
我不理自己.
但是耶穌也關心我.
是不是很奇妙?.
一個敬畏神的人的好處.
詩篇34篇第七節.
耶和華的使者在敬畏神的人的四位.
安迎答救他們.
在疫情最嚴峻的時候.
我就貼了這節經文.
因為我每天都要出去見人.
主人求你.
安迎.
保護我們.
另外.
再下一個詩篇31篇第十九節.
我們讀一讀.
一二三.
敬畏你.
投靠你的人.
你為他們所積存的.
在世人面前所施行的恩惠.
是何等的大呢?.
信不信?.
上帝還有很多很多恩惠給我們.
所以想起未來應該是樂觀的.
再下詩篇103篇第十一節.
一二三.
父親怎樣憐恤他的兒女.
耶和華也怎樣憐恤敬畏他的人.
其實我現在說敬畏說得很少.
如果去找敬畏舊約經文是很多的.
不過我將其中一二跟大家說.
上帝很憐恤我們.
很知道我們最大的軟弱在哪裡.

$^{961}$很明白.
甚至會.
意志.
我曾經是一個暈的專家.
甚麼叫暈的專家呢?.
是我有耳水不平衡.
一年都暈一次到兩次.
一暈起來甚麼都做不到.
我想你有耳水不平衡.
你知道我說甚麼.
困擾我的.
去服侍坐飛機,坐車暈.
有一年去以色列.
我還記得.
我還要帶著防暈手鐲.
去到加利利湖.
有沒有人去過?.
以色列加利利湖.
有沒有人去過?.
上去.
那個人會唱歌.
那個外國人.
我們整艘船在那裡.
我當時覺得加利利湖的湖水很平靜.
有些人在那裡敬拜.
然後他就唱.
How great thou art.
你真偉大.
我當時心裡被上帝的偉大所感動.
亦不理三七二一.
你猜我做了甚麼.
我跪在地上.
我說耶穌.
你知我要服侍你.
但我整天暈.
很難救.
主啊.
求你醫治我的暈.
讓我不用吃藥.
不用戴手鐲.

$^{1001}$讓我服侍你.
起來.
拔手鐲.
船搖動.
我沒有暈.
搭快艇乘風破浪.
我沒有暈.
應該是十年前的事.
到現在我想我的同工都未聽過我.
說我會以誰不平衡.
父親怎樣連戌.
他的兒女.
耶和華.
怎樣連戌.
敬畏他的人.
《得道主的喜愛》詩篇147篇.
十到十一節.
我們讀.
「他不喜悅馬的力大.
不喜愛人的腿快.
耶和華喜愛敬畏他.
和盼望他慈愛的人」.
我們的社會推崇力大.
快.
準.
但耶穌體重的是.
敬畏他和盼望他慈愛的人.
如果我沒有記錯.
上面那句話是.
奧斯威特·戴德森說的.
內地會的創辦人.
「神想你擁有更好的東西.
遠遠比富有和金錢更好.
而那就是.
不需要靠祂的依靠.
上帝要你得到的東西.
不比錢財.
什麼都更加好.
就是完全無助.
對上帝的依賴.

$^{1041}$有沒有經歷過?.
開人心很低谷.
我經歷過.
無法祈禱.
哭了八小時.
哭到無法收聲.
覺得自己沒用.
覺得自己不配再服侍上帝.
何牧師想扶起我.
我說你不要扶我了.
你任由我趴在地上哭.
但在這樣的過程中.
兩個星期後我要講培靈會.
陳牧你應該明白多大壓力.
在那個星期日.
我不用講到遇上,回到教會.
我再問上帝.
我可以不服侍你.
我可以取消所有東西.
因為我不配.
誰知.
有一首歌.
我不知道有沒有播給大家聽.
叫《緊握你的手》.
連曲連詞.
下載了進我的心.
整首歌是安慰我自己.
叫我緊握著天父的手.
天涯海角.
信靠信服.
跟著祂哭.
我估不到就是因為這首歌.
我站起來.
我經歷到上帝的實在.
我看到上帝是恩待的.
講了很久.
講了第一點.
因為第一點很重要.
其他都重要.
但第一點真的很重要.

$^{1081}$第二.
丈夫信任有益無損.
他給丈夫有自信.
他丈夫心裡依靠他.
必不缺少利益.
這裡當然說他很支持丈夫.
不過從另外一個角度.
我想說.
姐妹.
我們要學.
從來不要說丈夫一句壞話.
因為你說著說著.
你就這樣看他.
對丈夫.
她的丈夫是成功的.
你看丈夫在城門口.
與本地長老同坐.
為眾人認識.
有地位有名望.
大家都尊敬他.
當丈夫有一些榮譽的時候.
其實都很開心.
我記得我去年看到何牧師.
拿了他學校的傑出校友獎.
真的很高興.
這裡第二樣東西.
他特別計算什麼呢?.
給你看看.
他開口就發自為他.
舌上有人情的法則.
他不是縱容的.
他也有原則.
而我稱呼他一個親子專家.
因為《箴言三十一章》二十八節說.
他的子女都稱他為父.
他的丈夫也稱讚他.
是不是他很厲害?.
應該是他教導的.
姐妹也好.
弟兄也好.

$^{1121}$所以我要學一件事.
叫做「巴閉」.
什麼叫「巴閉」呢?.
就是嘴巴上筆.
在疫情期間.
我們很多時候有很多埋怨.
我正在看一本書.
叫《與人共贏》.
是一個管理學的大師寫的.
他說我們對人.
經常都會用一個鎚子.
「扑」你.
你這裡做得不好.
「扑」.
冷氣要大一點.
「扑」.
看到的都是別人的不好.
他說當我們真的要跟別人溝通的時候.
請放下鎚子.
眼鏡布.
眼鏡布幫我們怎樣?.
自己看得清楚一點.
又幫人擦一下.
大家又看得清楚一點.
下面這個.
其實整個都是師母幫我做的.
我也訪問了不同的人.
就說在疫情當中.
很多人都說.
疫情當中我們內心很多對話.
你不如列舉.
其實有12點.
不如現在選3點.
我中招了.
遲早問題而已.
不是的.
朋友也中招.
問一下他怎麼辦.
兩種想法.
是不是?.

$^{1161}$第二.
家裡菲傭姐姐.
星期天要出去.
她不聽話.
解僱她.
這是親身經歷.
有誰家裡有菲傭姐姐?.
舉手.
有沒有?.
她放不放假?.
疫情期間.
有些不放.
「你好ye ,我不行」.
我最初就.
動之以驚嚇.
「你出去會有事的」.
「不可以出去」.
「No」.
接著有些說.
「你有了」.
「I give you money」.
「OK」.
「Don't go out」.
「No」.
我就問「為什麼?」.
她說「媽媽」.
「我從星期天到星期日都在家」.
「OK」.
「如果你不讓我出去」.
「我無法呼吸」.
她說「I will die」.
差不多這樣跟我說.
我當時.
告訴大家.
當時.
我的臉是紅的.
劍拔弓張.
為什麼?.
因為女兒告訴我.
她身邊朋友中招.

$^{1201}$全部跟菲傭有關.
「當然是菲傭了」.
「他們讓她出去」.
「她染了」.
「公公婆婆也染了」.
「媽媽」.
「不要了」.
但是.
我當時是差那句.
「你再這樣」.
「我解僱你」.
我就是沒有說那句.
我等自己過了兩天.
發覺.
「是的」.
「如果我是她也很慘」.
「沒有放假」.
「兒子畢業看不到」.
「丈夫進醫院照顧不到」.
「星期日跟幾個朋友」.
「吃飯也不行」.
「你是不是很過份?」.
「是的」.
於是又再跟她聊.
「OK」.
「你可以出去」.
「但你可不可以」.
「早點回來?」.
「是嗎?」.
「早點回來」.
「還有」.
「你可以跟朋友一起吃飯」.
「但你可不可以」.
「分開吃飯?」.
「是嗎?」.
「兩個吃吧」.
「不要八個吃」.
她又回答.
結果.
過去這幾個星期.

$^{1241}$全部都自動做回來了.
我沒有叫她.
原來.
當我們的關係.
不是劍拔弓張的時候.
甚麼都容易華為.
但我們甚麼都用錘子的時候.
就很難搞了.
好,下去.
「持家有道,未雨綢繆」.
她很懂得計劃.
買東西.
但我覺得最重要的.
「叉尋找洋湧和麻」.
最重要的兩個字.
叫「甘心」.
我不知道大家明不明白.
當你甘心樂意做那件事.
你很開心沉在裡面.
你完全不覺得辛苦.
昨天我累了.
我昨天早上.
去了買菜.
因為.
晚上要乖孫吃飯.
買了很多東西.
然後.
去趕了.
差不多一個多兩個小時的拖.
然後回去.
又陪乖孫玩.
玩完之後.
又預備今晚.
「是嗎?今晚」.
「今天不是星期日嗎?」.
「是呀」.
「雙劍合璧」.
真是累得.
十點半就上床睡覺.
但有一句說話.

$^{1281}$令我.
現在才說.
我乖孫在吃飯的時候.
突然間.
「婆婆,可否抱你一下?」.
「可以」.
「無緣無故吃飯,為何要抱我?」.
「我真的很喜歡你」.
「可否看到」.
「你甚麼累都沒有了」.
沒錯.
所以昨晚睡得很好.
睡眠質素.
標告訴我.
94分.
很高.
甘心.
開心.
樂意.
持家有道,未雨綢繆.
他不是自己做完的.
未到黎明他就起來.
將食物分給家中的人.
分給人做.
不是覺得自己.
一定甚麼都可以.
都.
再下去.
他是.
不偷懶.
這裡說.
他覺得所經營的有利.
他燈中夜不滅.
我都查過甚麼叫做「中夜不滅」.
是否真的不睡覺.
有些說應該不是.
何牧說.
他如果.
中夜不滅.
那他夫妻.

$^{1321}$應該都有一些.
大家都有一些溝通.
這裡沒有寫得很多.
總之他很勤力的.
他觀察家務並不吃閒飯.
在疫情期間.
我經常都鼓勵.
很有壓力的.
家人弟兄姊妹.
一定要有甚麼呢.
Me time.
和Family time.
自己安靜的時間.
和家人一起.
天倫樂的時間.
我其中一個天倫樂的時間.
就會和乖孫.
去公園.
有他玩甚麼就玩甚麼.
最初和他撿樹葉.
撿到滿夜一天.
他說婆婆.
撿樹枝.
我說撿樹枝來幹甚麼.
他原來拿樹枝打欄杆.
原來不同的欄杆.
有不同的聲音.
打得很高興.
因為我教他.
從小到大唱那首歌.
是甚麼.
「和遊上高山住木」.
一直打.
樹枝斷了.
然後.
我就想.
撿吧.
撿了十五分鐘.
都找不到一條硬的樹枝.
我跟他說.

$^{1361}$「剃頭」.
「回家吧」.
「不要撿了」.
他突然說了一句.
「婆婆不可以放棄」.
誰跟你說的.
你.
那怎麼辦.
硬著頭皮.
繼續找.
找了十五分鐘.
這枝漂亮嗎.
硬的.
這枝叫做.
「永不放棄樹枝」.
(笑).
跟他一起打.
拿回家.
我給自己鼓勵.
不要放棄.
是嗎?.
不要輕易放棄.
好好下去吧.
「待人慷慨,周濟困苦」.
這裡說.
「才德婦人」.
她是「stretch out and open her arms」.
你過來吧.
我擁抱你.
「伸手行多一步」.
去幫助窮乏的人.
疫情期間有不同的人問我.
要關心可以做些甚麼.
又或者有人在醫院.
我們可以怎樣安慰他.
以下我寫了一些東西.
譬如送小禮物.
有些安慰卡.
送服務.
幫他煮食.

$^{1401}$或者特別確診的送食物給他們.
幫忙照顧.
在右邊.
不知道看不看得到那些複華.
是很激勵我的陸議長.
他很年輕就相信了耶穌.
他一直很正向.
經常想怎樣幫助別人.
80歲他太太過世.
他自己申請入老人院.
他的兒子問他.
「你為甚麼要自己入老人院?」.
他說「我想入去鼓勵老人家」.
90歲學水彩畫.
這幅是他畫的.
漂亮嗎?.
95歲在雪梨開了畫展.
他跟我說.
「千萬不要說自己老」.
「多少歲都可以學東西」.
他就將這些畫送給他身邊.
需要被鼓勵的人.
到最後.
這裡我叫做「Appraisal」.
嚴厲是虛假的.
美容是虛浮的.
這裡說他兒女稱讚他有福.
然後說他唯獨你超過一切.
最後想用一個最近的見證.
來結束.
疫情下.
其實我也艱難.
疫情下.
經常說.
為甚麼有信心.
為甚麼有盼望.
但在兩個月之前.
何牧師無意中做了血液檢查.
醫生給到的消息是說.
有擔心.

$^{1441}$要再做骨髓檢查.
當時.
你問我.
也不知怎樣說.
但依然要說到.
依然要每天帶著正向前行.
在這些日子.
上帝給了很多很特別的安慰給我.
「你看看下一張相」.
「我不知你有沒有去訂那些公仔」.
去那些嘉年華.
我以前不敢訂的.
因為我媽媽說.
「You are a bad luck girl」.
「用你的名字買股票」.
「那個股票就跌了」.
所以我不敢訂任何東西.
不敢抽獎.
抽乒乓球也不抽.
信了耶穌更加大條道理.
「我不貪圖這些東西」.
「你不要騷擾我」.
但何牧師很知道我心底很心結.
在很多年前.
天馬纜還有嘉年華的時候.
他抓了我去.
可能也跟大家說過.
抓了我去玩.
然後用二百元左右.
買了硬幣給我訂.
但訂之前最重要是那個祈禱.
你知道媽媽覺得自己是壞運女孩.
求你的大能將她的標誌.
今天撕掉.
我去訂.
訂了兩個大公仔.
訂彩虹.
中彩虹.
訂樽仔.
中樽仔.

$^{1481}$當時我覺得.
將我這個壞運女孩的標誌撕掉.
也沒有再特別去訂.
但在這一天.
拍照的這一天.
我覺得壓力很大.
何牧師的健康報告.
未知是怎樣.
等完又說要等.
那個星期.
我想師母就知道我有四個託.
完全不同的託.
很多要求.
我覺得壓到我呼吸不了.
於是我就選擇去逛商場.
在商場無意中又見到這些彩虹天地.
那我.
你知道我很喜歡做小女孩的祈禱.
我在心裡做的祈禱.
「耶穌啊,我很喜歡布甸狗」.
「如果你恩待我」.
「你讓我今天拿到一隻布甸狗」.
進去後.
我右手受傷了.
你們現在看不到.
我右手的肌肌仍然有些痛.
我扔,很難扔.
扔極也不中.
好,左手.
我最看不起的就是左手.
一扔,中了.
布甸狗.
我們那天.
兩個人用了二百元.
整個床的禮物都是我們的.
有幾十張貼紙.
有六隻小公仔.
大小公仔.
練習寶.
我還記得當時我們是整個場內.

$^{1521}$年紀最大的那兩個.
那些人看著我們很羨慕.
我又很驕傲地走出來.
但我的眼睛又淚光.
因為我知道.
是上帝的恩待.
我知道是上帝.
連續我們發自心底的禱告.
各位弟兄姊妹.
無論你遭遇什麼情況.
我很想你知道.
我們是上帝看顧的.
祂會維持著你不放.
祂會讓我們的心能夠回轉.
我們一起低頭.
在這個時候.
我心裡有一個感動.
你不用告訴我是什麼.
可能.
在弟兄姊妹當中.
心底有一些標誌.
是原生家庭貼在我們的心上.
又或者在這段日子裡.
我們對自己一些批判.
責備.
但今天.
我深信.
無論是經文還是聖靈.
都很想告訴你.
我們可以將這個標誌撕掉.
如果有弟兄姊妹有感動.
很想將這個標誌撕掉.
請你將你的手放在心上.
我們一起祈禱.
求主幫我們撕掉它.
你可以心裡跟著我祈禱.
親愛的主耶穌.
你知道.
我心裡的標誌是什麼嗎?.
求你幫助我.

$^{1561}$提醒我.
讓我知道我是你所重視的.
珍貴的.
讓我知道.
我是I am good.
I am whole.
I am beautiful.
或者handsome.
求主幫助我們撕掉這些標誌.
由願你的恩惠,慈愛與我們同在.
我們這樣交託.
仰望奉救主耶穌基督得勝的明智求.
願上帝祝福我們每一個.
我們明天見.
\newpage



\section{創世記 39:1-23}
\label{sec:e5I3VklFIWg}
\textbf{2022年6月5日主餐崇拜直播｜陳明泉牧師｜並肩勇闖系列:耶和華的同在｜創世記39\_1-23}
\newline
\newline
連結: \href{https://youtube.com/watch?v=e5I3VklFIWg}{\texttt{https://youtube.com/watch?v=e5I3VklFIWg}} ~~~~ 語音日期: 2022-06-06
\newline
\newline
\hyperref[sec:uzCPC0kQbAA]{\small{< < < PREV SERMON < < <}}
~
\hyperref[sec:index_chronic]{\small{[返順時目]}}
~
\hyperref[sec:_Anjyu4UPkQ]{\small{> > > NEXT SERMON > > >}}
\newline
\newline
創世記 39:1-23
\newline
\begin{longtable}{cl}
\hline
\hline
章節 & 經文 (和合本修訂版)\\
\hline
39:1 & \begin{tabularx}{0.7\textwidth}{X} 約瑟被帶下埃及去。有一個埃及人波提乏，是法老的官員，是護衛長，他從那些帶約瑟下來的以實瑪利人手中把約瑟買了去。 \end{tabularx} \\ \\ \relax
39:2 & \begin{tabularx}{0.7\textwidth}{X} 約瑟在他埃及主人的家中，耶和華與他同在，他是一個通達的人。 \end{tabularx} \\ \\ \relax
39:3 & \begin{tabularx}{0.7\textwidth}{X} 他主人見耶和華與他同在，又見耶和華使他手裡所辦的事都順利， \end{tabularx} \\ \\ \relax
39:4 & \begin{tabularx}{0.7\textwidth}{X} 約瑟就在主人眼前蒙恩，伺候他主人，主人派他管理家務，把一切所有的都交在他手裡。 \end{tabularx} \\ \\ \relax
39:5 & \begin{tabularx}{0.7\textwidth}{X} 自從主人派約瑟管理家務和他一切所有的，耶和華就因約瑟的緣故賜福給那埃及人的家；凡家裡和田間一切所有的，都蒙耶和華賜福。 \end{tabularx} \\ \\ \relax
39:6 & \begin{tabularx}{0.7\textwidth}{X} 波提乏把他一切所有的都交在約瑟手中，除了自己所吃的食物，其他的事一概不知。約瑟英俊健美。 \end{tabularx} \\ \\ \relax
39:7 & \begin{tabularx}{0.7\textwidth}{X} 這些事以後，約瑟主人的妻子以目送情給約瑟，說：「你與我同寢吧！」 \end{tabularx} \\ \\ \relax
39:8 & \begin{tabularx}{0.7\textwidth}{X} 約瑟拒絕，對他主人的妻子說：「看哪，一切家務我主人一概不知，他把所有的都交在我手裡。 \end{tabularx} \\ \\ \relax
39:9 & \begin{tabularx}{0.7\textwidth}{X} 在這家裡沒有人比我更大，除你以外，他也沒有留下一樣不交給我，因為你是他的妻子。我怎能行這麼大的惡，得罪神呢？」 \end{tabularx} \\ \\ \relax
39:10 & \begin{tabularx}{0.7\textwidth}{X} 她天天這樣對約瑟說，約瑟卻不聽從她，不與她同寢，也不和她在一起。 \end{tabularx} \\ \\ \relax
39:11 & \begin{tabularx}{0.7\textwidth}{X} 有一天，約瑟進屋裡去辦事，家裡沒有一個人在那屋子裡， \end{tabularx} \\ \\ \relax
39:12 & \begin{tabularx}{0.7\textwidth}{X} 婦人就拉住他的衣服，說：「你與我同寢吧！」約瑟把衣服留在她手裡，逃出外面去了。 \end{tabularx} \\ \\ \relax
39:13 & \begin{tabularx}{0.7\textwidth}{X} 婦人看見約瑟把衣服留在她手裡逃到外面， \end{tabularx} \\ \\ \relax
39:14 & \begin{tabularx}{0.7\textwidth}{X} 就叫了家裡的人來，對他們說：「看，他帶了一個希伯來人到我們這裡戲弄我們。他到我這裡來，要與我同寢，我就大聲喊叫。 \end{tabularx} \\ \\ \relax
39:15 & \begin{tabularx}{0.7\textwidth}{X} 他聽見我放聲大喊，就把他的衣服留在我這裡，逃出外面去了。」 \end{tabularx} \\ \\ \relax
39:16 & \begin{tabularx}{0.7\textwidth}{X} 婦人把約瑟的衣服放在身邊，直到他主人回家， \end{tabularx} \\ \\ \relax
39:17 & \begin{tabularx}{0.7\textwidth}{X} 就用這樣的話對他說：「你帶到我們這裡來的那希伯來僕人進來要調戲我， \end{tabularx} \\ \\ \relax
39:18 & \begin{tabularx}{0.7\textwidth}{X} 我放聲大喊，他就把衣服留在我身邊，逃到外面。」 \end{tabularx} \\ \\ \relax
39:19 & \begin{tabularx}{0.7\textwidth}{X} 主人聽見他妻子對他說的話，說：「你的僕人就是這樣對待我」，就非常生氣。 \end{tabularx} \\ \\ \relax
39:20 & \begin{tabularx}{0.7\textwidth}{X} 約瑟的主人把他抓起來，關在監獄裡，就是王的囚犯被關的地方。於是約瑟在那裡坐牢。 \end{tabularx} \\ \\ \relax
39:21 & \begin{tabularx}{0.7\textwidth}{X} 但耶和華與約瑟同在，向他施恩，使他在監獄長的眼前蒙恩。 \end{tabularx} \\ \\ \relax
39:22 & \begin{tabularx}{0.7\textwidth}{X} 監獄長就把監獄裡所有的囚犯都交在約瑟手下；在那裡的一切事都由他處理。 \end{tabularx} \\ \\ \relax
39:23 & \begin{tabularx}{0.7\textwidth}{X} 任何交在約瑟手中的事，監獄長一概不察，因為耶和華與約瑟同在，耶和華使他所做的都順利。 \end{tabularx} \\ \\
[1ex]
\hline
\hline
\end{longtable}
$^{1}$接著是眾聽聖道敬拜的時間.
請黃潔明牧師為我們讀聖經創世記.
接著我們一起聆聽宣講耶和華的同在.
今天的經分是記載在創世記第39章.
請聆聽上帝的話.
約瑟被帶下埃及去.
有一個埃及人是法老的內臣.
護衛長波提弗.
從那些帶下他來的二十馬利人手下.
買了他去.
約瑟住在他主人埃及人的家中.
耶和華與他同在,他就百事順利.
他主人見耶和華與他同在.
又見耶和華使他手裡所辦的盡都順利.
約瑟就在主人眼前蒙恩.
侍候他主人.
並且主人派他管理家務.
把一切所有的都交在他手裡.
自從主人派約瑟管理家務.
和他一切所有的.
耶和華就因約瑟的緣故.
賜福於那埃及人的家.
凡家裡莫填間一切所有的.
都蒙耶和華賜福.
波提弗將一切所有的都交在約瑟的手中.
除了自己所吃的飯.
別的事亦愧不至.
約瑟原來羞瓦盡美.
這是以後約瑟主人的妻.
以慕頌情給約瑟說.
你與我同尋伯.
約瑟不從對他主人的妻說.
看那一切事務.
我主人都不知道.
他把所有的都交在我手裡.
在這家裡沒有比我大的.
並且他沒有留下一樣不交給我.
只留下了你.
因為你是他的妻子.
我怎能作這大惡得罪神呢.

$^{41}$後來他天天要和約瑟說.
約瑟卻不聽從他.
不與他同尋.
也不和他在一處.
有一天約瑟進屋裡去辦事.
家中沒有一個在那屋裡.
婦人就拉著他的衣裳說.
你與我同尋吧.
約瑟把衣裳丟在婦人手裡.
跑到外邊去了.
婦人看見約瑟把衣裳丟在她手裡.
跑出去了.
就叫了家裡的人來對他們說.
你們看他帶了一個希伯來人進入我們家裡.
要欺弄我們.
他到我這裡來.
要與我同尋.
我就大聲喊叫.
他聽見我放聲喊起來.
就把衣裳丟在我這裡.
跑到外邊去了.
婦人把約瑟的衣裳放在自己那裡.
等著他主人回家.
就對他如此如此說.
你所帶到我們這裡來的那希伯來人.
真要欺弄我.
我放聲喊起來.
他就把衣裳丟在我這裡.
跑出去了.
約瑟的主人聽見他猜似對他所說的話.
說你們僕人如此如此待我.
他就生氣.
把約瑟下在監裡.
就是黃的囚犯被囚的地方.
於是約瑟在那裡坐監.
但耶和華與約瑟同在.
向他施恩.
使他在施玉的眼前蒙恩.
施玉就把監裡所有的囚犯.
都交在約瑟的手下.

$^{81}$他們在那裡所辦的事.
都是經他的手.
凡在約瑟手下的事.
施玉一概不察.
因為耶和華與約瑟同在.
耶和華使他所作的盡都順利.
經文同學.
.
弟兄姊妹平安.
我們這次開始一個新的講道主題.
並肩用廠主題.
這個主題一來是回應教會.
我們六十週年.
我們用這個作為我們對未來的一些寄望.
同時間也是一個很深遠的神學觀念.
因為我們基督徒的生命不是孤軍作戰的.
既有上帝與我們同在.
也有弟兄姊妹與我們每一個同在.
我們就是大家拍肩膀.
一起走明天的道路.
當然在人生裡面有高高低低.
有時候順境 有時候逆境.
不過當我們認清了我們生命不孤單的時候.
我們就有力量走面前的路.
今天跟大家選的這段經文.
我想作為這個系列的開始.
當我們說人生不孤單的時候.
我們看的不是純粹看我們的友情家庭關係.
我們第一樣看的就是神與我們之間的關係.
坊間裡面有一個心理學理論叫依附理論.
不知道大家有沒有聽過.
這個依附理論其實在說什麼呢.
人在關係 特別是在親情或者在友情感情關係上面.
有些人有四種的形態.
我講給大家聽 你自己看看那種形態.
第一種就叫做安全感型的依附.
就是說你對人是很自在的.
對於陌生人也好 對於自己身邊的人也好.
你可以將你最深層的一面.
按著你的安全感來跟別人分擔和分享.

$^{121}$這種就是安全型的依附.
第二種就開始有一點負面了.
這個就是焦慮型或者抗拒型的依附.
什麼意思呢.
就是說我很擔心別人不愛自己.
所以經常在人際關係裡面.
在親情 親密關係裡面.
經常都會患得患失的.
想做些事情來討好別人.
但當別人可能一個眉頭眼額.
你會覺得很不安 很焦慮不安.
這種就是焦慮型或者抗拒型的一種依附型態.
第三種就叫做逃避型.
什麼意思呢.
就是說那個朋友 家人.
他的存在與否.
你基本上就像關上了心門一樣.
漠不關心 麻木的.
他在不在對你來說沒有影響.
你不要想著這件事很高超.
而心裡面有一種保護自己.
很怕進入一種親密關係.
所以這一種漠不關心 麻木的一種表現.
形成了自己有時候在人際關係中會抽離出來.
第四種 這種更難處理的就是disorganized.
這種叫混亂式的依附關係.
有時候很焦慮.
有時候忽然間很熱情.
但忽然間又很擔心.
將剛才那幾種型態混合在生命裡面.
所以令自己進退兩難在人際關係裡面.
不知道大家在依附型態裡面覺得自己像哪一種.
依附型態的心理學家.
發展出很多不同的心理輔導理論.
大家經常聽到的就是.
童年陰影 或者父母在你年輕的日子裡離棄過你.
以致你會有一種依附型態裡面的扭曲.
所以你就會沒有安全感.
通常都會做的就是.
回到你小時候的舊時 找回那些片段.

$^{161}$然後有一個意識.
然後在那裡做一些心理治療.
坊間五花八門的輔導也好.
很多人都是用一個這樣的方式.
用這個理論來建構一個心理輔導的.
一個明白對方的心態的一種做法.
在聖經裡面沒有特別去講這個理論對與錯.
不過今天我們所看的聖經.
這個主人翁約瑟.
大概可以從一個依附理論.
來成為他生命裡面的一個出路.
約瑟在他年輕的日子 很小的時候.
他的媽媽就因為難產.
生了小孩之後就去世了.
大概這個年輕人在他成長的階段裡面.
失去了媽媽的母愛.
在他自己的人生裡面.
雖然他爸爸偏心他.
但是他有十個哥哥針對著他.
在一個對他很不友善.
又很溺愛的環境之下.
他大可以成為一個很混亂的人.
不過今天在聖經裡面講到.
上帝與這個年輕人同在.
幫他譜寫出一個新的人生.
希望今天我們從約瑟的這一段.
很長的經歷.
不過我們抽取一個片段.
特別強調耶和華語約瑟同在的這一段經文.
我們看看上帝怎樣幫助這個年輕人.
也看看上帝怎樣藉著他的同在.
幫助我們每一個.
我們一起祈禱吧.
求上帝幫助我們.
讓我們明白這段經文和深祈禱.
天父上帝願你在我們當中帶領我們.
讓我們能夠在你的華語當中.
找到生命的亮光.
賜福給我們眾人.
讓我們聆聽的時候.

$^{201}$不單止聽到一些訊息.
更加是聽盡我們內心.
你要對我們所講的每句話.
藏在我們的心靈裡面.
實踐在我們生活當中.
帶領我們走面前的路.
我們求你幫助眾弟兄姊妹.
每一個都聽得明白.
幫助孫講的講得清楚.
忠心的將你的華語.
帶給眾弟兄姊妹.
恭敬將來的時間交託.
願你賜福.
獻上祈禱奉靠祈禱耶穌得勝的聖名.
阿們.
好 那我們一起看.
在《創世記》第39章.
第一節到第六節是第一個段落.
這個段落很明快的講到.
約瑟在他剛才講到.
十個哥哥針對他.
他的爸爸就偏心在他身上.
媽媽在他很年幼的日子已經離開了.
在一個有溺愛又敵對的環境之下.
十個哥哥就嫉妒他.
於是有一天約瑟去找哥哥的時候.
那十個哥哥心裡就很憤怒.
於是就藉著這個機會.
就將他賣給了以實瑪利人.
賣到埃及這個地方.
所以第一節裡面講到.
約瑟被帶到埃及去.
有一個埃及人就是法老的內臣.
護衛長波提弗.
從那些大廈他來的以實瑪利人.
手下買了他去.
約瑟就住在他主人埃及人的家中.
在這裡講到其實是一個很淒涼的狀況.
用今天的言語就是說.
他那十個哥哥拐賣了自己的弟弟.

$^{241}$將他自己的弟弟賣給了拐子.
這個拐子將約瑟賣了下去的時候.
就做別人的奴僕.
被賣到別人家裡做奴僕.
就注定了這一生是沒有的.
本來是家中裡的爸爸的天子嬌子.
家裡被溺愛的一個年輕人.
瞬間之間就成為埃及.
去到外邦裡面的一個家庭裡最低下的傭工.
可以想像的是.
面對一個環境失去了家人的支援.
即使以前有十個哥哥針對他.
也有人傾談一下.
有人針對一下.
在這裡基本上就寄人籬下.
言語不通.
在埃及他講開希伯來話.
去到埃及所有東西都是一個新的環境.
其實他可以自暴自棄.
但可以淒淒涼涼悲悲慘慘地過他的生活.
在伊芙理論裡面說的.
他完全沒有一個社交群體為他作為一個支援.
不過在聖經裡面說到.
他真正的支援不是人際關係.
聖經裡面馬上就告訴你們.
耶和華與弱勢同在.
上帝與這個.
上帝所命定將來生命裡面承擔一個更大的責任.
這個年輕人同在.
上帝與這個年輕人同在.
不過這種同在就不是帶來一種心靈慰藉這麼簡單.
在經文裡面.
創世記的作者很仔細地多次用了那個字.
就是手下.
就是在弱勢手裡面所作的功.
這個字眼和弱勢所作的功.
來表達耶和華與這個年輕人同在.
在第二節我們繼續看.
弱勢住在他主人埃及人的家中.
耶和華與他同在.

$^{281}$他就百事順利.
當主人見耶和華與他同在.
又見耶和華使他手裡所辦的盡都順利.
弱勢就在主人眼前蒙恩.
侍候他的主人.
並且主人派他管理家務.
把一切所有的都交在他的手裡.
自從主人派弱勢管理家務.
和他一切所有的.
耶和華就因弱勢的緣故賜福那埃及人的家.
凡家裡或田間一切所有的都蒙耶和華賜福.
把一切都交在弱勢的手裡.
除了自己所吃的飯.
別的事一概不知.
在這個孤立無援的狀態裡.
上帝與這個年輕人同在.
上帝具體的同在是如何表現出來呢?.
這個只是老僕.
上帝賜他所做的一切都能夠順利.
順利這個字我們中國人很喜歡看的.
我們覺得做的事很順很容易做好.
這是一個意思.
不過聖經裡面所說的順利的字.
比我們字面看更加深遠.
順利的意思在原文裡更加說到是昌盛.
令到那件事有好的結果.
我們今天中國人說的順利.
就是順頭順道很有效率地完成了.
不過順利如果用原文的意思.
不是說那件事很快完成.
而是說那個效果很壯大.
那個效果很顯著.
舉個例子.
英文沒有寫的.
我只是想像.
每個家宅每個老僕都站在田裡淋花.
在田裡種植.
其他人也是這樣淋.
弱勢也是這樣淋.
其他人淋的花也會長.

$^{321}$所以就長了.
不過弱勢淋的花就特別壯大.
結果更大.
根更加深.
所有的東西跟其他人的效果是很不同的.
弱勢所做的每一件事都有顯著的效果.
在他生活當中.
可能大家都是洗碗.
不過弱勢洗碗不單洗好碗.
弄到整個廚房都整整有條.
打理得一乾二淨.
大家明白了嗎?.
順利的意思是說.
弱勢自己努力.
再加上上帝的祝福.
令到身處的環境.
本來每天都是這樣的.
經過他的手的時候.
所有東西的果後變得很多.
很大.
而且整個環境都是被更新改變.
更重要的是.
上帝與的同在.
不是一個正態的同在.
上帝的同在.
是表現出在弱勢所做的功.
因為弱勢的努力.
所以上帝加倍的席著他.
做出來的東西帶出來更大的果效.
在經文裡面再說.
因為弱勢做得這麼好.
如果你是主人.
看到這麼多果效.
你看到他做事做得這麼好.
你會怎樣?.
多給他一些工作.
在經文裡面這樣說.
因為看到他盡都順利.
他在主人眼前蒙恩.
所以主人就派他管理家務.

$^{361}$把一切所有的都交付他的手裡.
可以這樣說.
如果懶的人聽見.
這不是福音.
因為你想想.
如果你做的事情盡都順利.
做得這麼好.
意味著你將要做多一些.
不過在經文裡面就是這樣說.
弱勢他沒有逃避任何生命裡面的責任.
因為上帝與他同在.
做每一樣事都能夠做得好.
於是他承擔更加多的責任.
在家庭裡面.
由一個卑微的奴婢.
一個奴僕.
但一直成為一個家裡面的管家.
經文裡面甚至說到極致.
波提弗在他的家裡面.
他只想他自己吃一餐飯.
他吃的東西.
他想吃些什麼.
其他一概也不知.
因為弱勢幫他打點了整個家庭.
所有的一切.
但在第一個.
我們從這個段落.
我們看到上帝的同在.
上帝的同在在這裡裡面.
更新我們今天對上帝同在的理解.
很多弟兄姊妹覺得.
我要經歷神的同在.
我要走到深山.
我要走到一些很隱蔽的深處.
或者我要不工作.
經歷神的同在.
要休息.
很多時候弟兄姊妹都會說這件事.
不過在《創世記》.
特別是說到弱勢的經歷.

$^{401}$一個奴隸沒有辦法.
他只能夠說.
波提弗啊.
我這個時間.
我靈性不是太好.
所以我需要放假一個星期去退休.
沒有這些東西的.
沒有這支歌唱的.
他只是默默地在他生命裡面.
做好他眼前的工作.
一路工作的時候.
一路祈願上帝與他同在.
事實上他所做的事.
和上帝的同在兩者是沒有衝突的.
不過今天我們很多弟兄姊妹.
我們一想.
經歷上帝的同在.
我們就會不做所有事.
放下眼前的任務.
放下所有事.
我們覺得完全去除了生命所有的責任和工作.
我們才能夠經歷上帝的同在.
在這裡經文提醒我們.
上帝的同在原來往往在人的作工裡面.
具體地去實現出來.
有一本書.
我很喜歡的作者 Parker Palmer.
之前也說過了.
另外一本著作.
叫做《行動靈修學》.
什麼意思呢.
很多人.
這個作者是說.
很多人想著靈修.
或者經歷上帝的同在.
很多時候都是.
那個脈觀就退到山野當中.
上到山上面.
不過這個作者是說.
真正展現基督徒生命的活力.

$^{441}$往往在你的生活場景當中.
你的作工.
你生命裡面具體的實踐.
你在工作裡面的觀察.
其實那個是更加展現了你生命的活力出來的.
你的活力.
或者你真正活著.
不是在那個山林曠野裡面.
你的活著在你生活崗位當中.
上帝不是純粹在深山.
在曠野.
或者我們躲在一個退休營地裡面.
才與我們同在.
在我們形形異異.
在我們覺得逆流而上的工作生活當中.
上帝施恩與我們同在.
所以這個大家我們必須要留意.
若是在一個很不利的環境之下.
藉著.
也就是說與他同在.
在他的工作.
在他的生命當中.
帶來一個新的處境.
帶來一個新的境遇.
這個是第一個.
第二個.
剛才我也有說過.
上帝同在意味著什麼呢?.
承擔更加多的責任.
凡有的會加給他.
沒有的.
連他所有的都拿走.
不是在說錢.
而是在說責任.
在這裡面說到上帝與約瑟同在.
他藉著他的生命.
他為身邊的人帶來好處.
帶來祝福.
他所做的一切.
他因為做得更加理想.

$^{481}$所以會承擔更加多的責任.
而這個更加多的責任.
帶來的就是更大的祝福.
整個的家庭.
整個的位置.
因著他的存在.
就帶來新的一片天.
在這裡面同樣都要問我們自己.
弟兄姊妹.
你覺得在你身處的環境.
當你處於那個工作.
或者家庭.
或者鄰舍你所住的區份當中.
你的出現.
會不會為身邊的人.
帶來更加多的祝福呢?.
還是帶來更加多的麻煩呢?.
今天容許我直接的說.
有時候基督徒很有趣.
我們會有意無意之間.
我們會對這個世界有很多的不滿.
很多時候那種不滿感.
那種憤怒.
有意無意之間.
在我們生活裡面表現了出來.
我們未必祝福了那個環境.
但是我們為那個環境帶來.
有時候有些怨氣.
我們覺得這個不好.
那個不好.
這個不好.
久而久之.
我們就為那個環境帶來了.
有一種很難以理解的壓力.
不過弱失的存在.
它不是為那個環境.
為一個對它不友善的處境.
帶來更加多困難.
反而它的存在.
帶來更加多的祝福.

$^{521}$我心願上帝在我們當中裡面.
每一個都能夠成為弱失的生命.
當我們身處的那個位置.
我的存在.
我的出現.
是令到人會帶來一些快樂.
我所做的功.
會帶來更加大的果後.
這個是第一個段落裡面.
在《創世記》的作者裡面.
講到上帝與弱失同在的時候.
就帶來一些.
波提弗家庭裡面.
一些革命性的改變.
好了,接著我們再看第二個段落.
第二個段落.
由第七節.
至到第十七節.
在這個段落裡面.
有幾大的篇幅.
不過和剛才的段落很不同.
耶和華與弱失同在.
在這裡裡面再沒有出現過.
在這個小節裡面.
接著再出現.
耶和華與弱失同在.
在第21至第23節又再出現.
一頭一尾講到上帝與弱失同在.
不過去到中間這個段落.
特別是在講.
弱失面對著主母對他的勾引.
或者是色誘.
在這裡裡面.
上帝沒有特別講到.
他與弱失同在.
那件事是怎樣的呢?.
在經文裡面講到.
第六節他下了一個注腳.
弱失原來瘦瓦盡美.
弱失很帥.

$^{561}$這麼帥 招來橫禍.
今天這個世界很喜歡帥的.
買豬肉也喜歡帥的.
大家都是有這個取態的.
希望大家回教會不是喜歡帥的.
否則大家會失望.
在經文裡面講到弱失很帥.
他的主母見到他這麼帥.
第七節.
弱失主人的妻.
以沐送情給弱失.
眉來眼去.
暗示.
說很直接.
你與我同寢吧.
在華本寫得很厲害.
但在這裡大家明白的.
講到有一些不倫的關係.
想令弱失和他有一些親密的.
越了軌的關係.
不過在第八節講到.
弱失不從.
接著對主人的妻說.
看一切的家無.
我主人都不知道.
他把所有的都交在我手裡.
在這家裡沒有比我大的.
並且他沒有留下一樣不交給我.
只留下你.
因為你是他的妻子.
我怎能作這大惡得罪神呢.
在這裡弱失鏗鏘有力的.
來去抵擋主母的勾引.
他第一句說話是.
在這家裡我要忠於波提弗.
這個主人.
在這家裡所有的都交付給我.
他甚麼都不知道.
是我管理一切.
不過唯獨你是沒有交給我.

$^{601}$因為你是他的妻子.
很清晰.
整個脈絡.
他講到他不會做一些不忠的事.
不過到最後.
弱失訴諸於的.
是一個甚麼的處境呢.
他訴諸於的就是.
他所做的這件事除了得罪主人之外.
他更加強調的就是.
我不可以作這些大惡得罪神.
他心裡面他想著的.
他是帶著上帝的同在.
在他的工作崗位當中.
面對著這些的試探和引誘.
一方面他要忠於他眼前的.
他在這裡講了他眼前的位置.
他身處的.
他的地界是在做這些事的.
不過更加重要的.
訴諸於那個終極.
他在講他的生命.
不可以得罪神.
在經文裡面再繼續的.
剛才我說過.
上下文沒有再提及上帝與弱失同在.
只是提到弱失因為不聽從主母.
但這個主母天天都繼續.
鍥而不捨的去勾引他.
直至到有一刻.
整間屋孤男寡女.
這個主母要拉著他的衣襟.
弱失連衣服都不要了.
然後他就立刻逃跑.
避免自己在這件事裡面跌倒.
就是因為這件衣服.
所以這個主母就用這件衣服去誣衊.
就說弱失想調戲他.
所以留下這件衣服作為證據.
其實這件衣服是說弱失逃走的證據.

$^{641}$不過卻成為了主母冤枉他的一個物證.
於是他的主人就很生氣.
就將弱失下到監獄裡面去.
在這麼長的段落裡面.
又說沒有與弱失同在.
在這個最大的困難.
在最大的關口裡面.
你可以說上帝為何不保守弱失.
為何在這麼危急的關口裡面.
為何上帝不在天上找個天君天使.
來阻止這個主母或是擊殺這個主母.
這樣就不用去受試探.
整個事件發生不是這樣.
上帝就好像隱藏了.
隱密了他自己一樣.
上帝在很隱藏的狀態之下.
弱失的生命才表現出一種真正的章法.
他生命的真章就在這些處境裡面表現出來.
在弱失的心裡面.
上帝的同在除了代表著順利.
更加代表著他生命裡面有一個道德的律.
這個律會左右他做人生的決定.
在我們人生裡面.
有時可能你覺得很大的困難.
在你的職場也好.
在你的生活裡面也好.
有很多這些挑戰,限制.
甚至有很多的誘惑,引誘.
你可能會求上帝出手.
將這些東西阻止了.
神跡地停止這些東西出現.
不過很多時候上帝對人的工作未必是這樣.
有時上帝隱密的同在.
上帝隱密他自己.
有時也要考驗一下我們每一個跟隨他的人.
我們的心有沒有上帝.
在弱失對著他主母和他所表現當中.
他知道這麼大的惡.
他不可以得罪神.
在上帝好像隱密當中.

$^{681}$但他心裡依然經歷上帝的同在.
上帝的同在體現出.
在他面對試探引誘的時候.
所作的每一個決定.
在這裡,弟兄姊妹.
我們都要繼續去問.
我們作為一個基督徒.
或者我們有時信上帝了.
我們祈求神的同在.
我們會不會求上帝的國.
上帝的義在我們生命裡面與我們同在.
上帝的義未必是上帝在超然.
用神跡騎士來臨到這個世界上.
有時上帝的義.
上帝要我們所實踐的.
往往是由我們發展內心.
我們因為敬畏上帝.
和這個世界裡面有一些做法.
完全的來割裂.
以致到我們走在上帝的心意當中.
似乎上帝在裡面好像很隱藏.
但我們心裡面.
或者在我們的行動當中.
我們表現出我們成為一個屬神的人.
今日我們在我們生活裡面.
我們有沒有這一份決心.
我們願不願意成為一個這樣的人.
來表現和這個世界不同.
不同身邊的誘惑來同流.
不被這些試探引誘所激導.
以致到我們陷在萬劫不復當中.
這是第二點我們需要.
一直去祈求上帝去光照我們.
我們再看第三個段落.
十九節至到第二十三節.
「若失的主人聽見.
他的妻子所說的話」.
即是誣蔑若失的話.
「說你的僕人如此如此待我.
他就生氣把若失下到監獄.

$^{721}$就是黃囚犯被囚的地方.
於是若失在那裡坐監.
但耶和華與若失同在.
恆他示恩.
使他在司獄的眼前蒙恩.
司獄就把監獄所有的囚犯.
都交在若失的手下.
他們在那裡所辦的事.
都是經他的手.
凡若失手下的事.
司獄一概不察.
因為耶和華與若失同在.
耶和華使他所作的盡都順利」.
很面熟的句式.
「上帝與若失同在.
凡他手下所作的盡都順利」.
當若失被誣蔑.
主人把他下到監倉裡面的時候.
他再一次經歷上帝與若失同在.
不過這裡雖然字句是一樣.
上帝與若失同在的來龍去脈差不多.
不過你可以知道處境已經變得很不同.
上帝之前與若失同在.
是在波提佛的家庭裡面.
所以這是一個比較友善的環境.
你做好那些事甚至有稱讚.
你做那些事會有發揮的機會.
不過如今去到監倉裡面.
監倉是一個什麼地方呢?.
是一個充滿了憤恨,充滿了悲哀.
面對生命眼前的絕望.
又或者每一個人都覺得自己是冤枉的.
一個被囚禁的地方.
失去自由.
在人生裡面某程度上是一個.
所有人人生的低谷.
都是在這裡裡面發生的.
是一個這麼烏煙瘴氣.
你不要想著今天現代的監獄.
昔日的喧賓難酬.

$^{761}$那個地方充滿怨氣寒微.
人裡面的那種憤世嫉俗.
在一個這樣密集的環境下.
可想而知最接近地獄的地方就是監獄.
在一個這樣的處境當中.
上帝如藥室同在.
在監獄裡面.
上帝都讓他成為一個很特別的使者.
來扭轉他自己的生命.
亦都扭轉了監倉裡面很多人的看法.
在經文裡面特別的來講.
施育將監獄裡面所有的囚犯都交在藥室的手下.
怎樣交呢?.
大家都坐監獄.
大家都是囚友.
怎樣交呢?.
我在這裡想像.
我相信其實不是在說要管理那些囚犯.
而是施育或者管理監倉的人最怕的是甚麼呢?.
就是那些囚犯作亂.
而且在當中裡面有暴力的行為.
搞到監倉很混亂自相殘殺.
那樣東西是最頭痛最麻煩的.
不過藥室的出現.
就是一班充滿了怨氣.
充滿了絕望.
充滿了憤怒的人當中裡面.
藥室就在他們當中.
施育將那些囚犯交付在藥室手裡面.
大概可能藥室就成為一個和平之子.
在一班這樣不理想的環境裡面.
他就活現出一種新的關係.
可能是調解他們當中裡面很多的憤怒.
後來的篇幅講到他為人解夢.
排難解婚.
訴說心事.
藥室竟然在監倉裡面.
成為一個這樣的人.
在他所做的一切.
不單止令到監倉裡面變得平靜.

$^{801}$更加是好像成為監倉裡面.
這班充滿了絕望的人.
心裡面的一個出路.
一個精神上的一種處境.
在經文裡面對我們今天的體會是甚麼呢.
可能在你的生活當中.
可能你會面對一個頗難的處境.
可能有些弟兄姊妹會說.
我認同有些工作崗位是頗困難的.
未必是上司下屬.
而是工種本身製造了大量的困難在當中.
早前有機會.
我不知道是不是.
我去到那個稅局.
我體會到人生.
因為等一份文件.
我坐了四個小時.
在四個小時裡面我看見人生百態.
那裡在稅局裡面很有趣.
那個地方罵人罵到飛起.
那些粗口.
那個環境很嘈吵.
我看見那裡.
我很佩服裡面工作的那些官員.
或者那些在稅局裡面工作的人.
我看見他們不溫不火地回答.
是這樣的.
填一張表就可以了.
坐一坐等一等.
那個人罵完粗口.
不用那麼嘈吵.
在那裡坐一坐.
填完就搞定了.
不用那麼生氣.
然後有些熟客.
他們已經開了.
有些送文件的.
不用快進.
要急.
類似那些.

$^{841}$經常在那裡大罵.
又是你嗎.
又那麼嘈吵.
在咪裡面是這樣說的.
但我覺得.
捍衛官子.
我覺得那班人.
EQ實在太高了.
然後見到準時的時候.
因為我等那份文件都過時了.
所以我等到.
他們平時五點半就收工.
我等到六點半.
他們都未收工.
還在工作.
離開的時候.
又搞定一天.
又走了.
這樣都可以.
在我心裡面.
如果我每天都面對這樣的處境.
我發瘋了.
我回去跟老婆說.
不過在這裡.
或者可以從這個角度來看.
我們有時在我們工作.
或者在我們身處的環境.
我們未必能夠將.
我的存在.
令到這個地方.
變成一個天堂.
這個實在叫價太高了.
在我們日常生活裡面.
我們未必能夠做到這樣的.
不過當我的出現.
當基督徒有上帝的同在.
我總是可以將一個.
不是很理想的地方.
惡化成為一個地獄.
這個地方.

$^{881}$可以帶來一點點的祝福.
在這個地方.
當大家很憤怒.
很多怨氣的時候.
你的出現.
你可以撫平.
疏導.
甚至可以為這個環境.
惡劣程度減輕.
若是在監倉裡面.
有耶和華與他同在.
在這樣的處境當中.
所有事情都變得順利.
監倉在這裡.
所作的功變得順利.
就不是很昌盛.
有更加多囚犯坐下去.
不過是在他的管理.
在他的同在之下.
所有人都可以變得心平氣和.
在一個覺得近似地獄的監獄當中.
若失的出現.
讓這個地方可能改換了.
未至於變得這麼惡劣.
一個這樣的狀態.
弟兄姊妹.
今天我們看《段段經文》.
為我們帶來有什麼新的看法呢?.
我們很想.
我們在說教會在旺角.
這個區份裡面.
我們怎樣去經歷上帝的同在呢?.
除了將這個地方變得美輪美奐.
我們很有人情味之外.
如果教會.
旺角浸信會在旺角這個社區.
在我們身處的這個地方.
我們可以令到身邊的人.
會經歷到上帝的同在.
會經歷到上帝的賜福.

$^{921}$這就是我們真正體現到.
上帝與我們同在的那種生命力.
在我們個人層面裡面.
我們每一個領受了福音.
在這裡接受建立成為一個門徒.
當我帶著這份福音的時候.
我身處的環境.
可能我會令到所有人工作更加昌盛.
我會令到環境變得更加好.
倘若上帝隱密的.
我們面對著試探隱憂.
好像上帝很隱密.
見不到上帝插手在這些事當中的時候.
我們心裡面依然有上帝.
我們心裡面有上帝同在.
所以我們做的決定.
我們人生所行的堅持的路.
我們知道我們按著上帝的心意而行.
可能做了這些決定會帶來一些代價.
不過我們知道.
我們怎可以作這些大惡來得罪上帝呢?.
我們就帶著這份單純的信念.
活好我們生命.
好好地在上帝隱藏的日子.
我們就發揮一種這樣的生命力.
甚或可能在一個很難的處境裡面.
在你的出現.
你會令到這個地方帶來一點點的溫情.
在一個接近地獄的狀態之下.
你令到它不會進入變成更地獄.
反而是讓它可以在人間當中呈現到.
在一些不理想的環境裡面.
呈現一點點的美與善.
求主幫助我們.
今日當我們去讀《創世記》裡面.
關於耶和華同在的時候.
我們就不要這麼靜態.
去躲在深山找上帝.
當我們讀通這段經文的時候.
我們知道我們更加有力量.

$^{961}$有那份勇氣和膽識.
在我們生活當中尋求耶和華同在.
在我們生命裡面亦都活現一種.
耶和華同在的生命.
今日從這段經文裡面.
求主繼續來對我們講說話.
亦都希望我們的生命成為別人的祝福.
\newpage



\section{列王記下 4:1-7}
\label{sec:_Anjyu4UPkQ}
\textbf{2022年6月12日崇拜直播｜陳冠成牧師｜並肩勇闖系列:神不撇下我們｜列王記下4\_1-7}
\newline
\newline
連結: \href{https://youtube.com/watch?v=-Anjyu4UPkQ}{\texttt{https://youtube.com/watch?v=-Anjyu4UPkQ}} ~~~~ 語音日期: 2022-06-13
\newline
\newline
\hyperref[sec:e5I3VklFIWg]{\small{< < < PREV SERMON < < <}}
~
\hyperref[sec:index_chronic]{\small{[返順時目]}}
~
\hyperref[sec:8sXdI3Lte_M]{\small{> > > NEXT SERMON > > >}}
\newline
\newline
列王記下 4:1-7
\newline
\begin{longtable}{cl}
\hline
\hline
章節 & 經文 (和合本修訂版)\\
\hline
4:1 & \begin{tabularx}{0.7\textwidth}{X} 有個先知門徒的妻子哀求以利沙說：「你的僕人，我丈夫死了，他敬畏耶和華是你所知道的。現在有債主來，要帶走我的兩個孩子給他作奴隸。」 \end{tabularx} \\ \\ \relax
4:2 & \begin{tabularx}{0.7\textwidth}{X} 以利沙對她說：「我可以為你做甚麼呢？告訴我，你家裡有甚麼？」她說：「婢女家中除了一瓶油之外，甚麼也沒有。」 \end{tabularx} \\ \\ \relax
4:3 & \begin{tabularx}{0.7\textwidth}{X} 以利沙說：「你到外面去向所有的鄰舍借器皿，要空的器皿，不要少借。 \end{tabularx} \\ \\ \relax
4:4 & \begin{tabularx}{0.7\textwidth}{X} 然後你回家，關上門，你和你兒子在裡面把油倒在所有的器皿裡，倒滿了就放在一邊。」 \end{tabularx} \\ \\ \relax
4:5 & \begin{tabularx}{0.7\textwidth}{X} 於是婦人離開以利沙去了。她關上門，把自己和兒子關在家裡。他們把器皿拿給她，她就倒油。 \end{tabularx} \\ \\ \relax
4:6 & \begin{tabularx}{0.7\textwidth}{X} 器皿都滿了，她對兒子說：「再給我拿器皿來。」兒子對她說：「沒有器皿了。」油就止住了。 \end{tabularx} \\ \\ \relax
4:7 & \begin{tabularx}{0.7\textwidth}{X} 婦人去告訴神人，神人說：「你去賣了油還債，你和你兩個兒子可以靠著所剩的過活。」 \end{tabularx} \\ \\
[1ex]
\hline
\hline
\end{longtable}
$^{1}$今日重讀的經文是記載在.
「列王記下」四章一至七節.
在舊約聖經的四百五十九頁.
「列王記下」四章一節.
有一個先知門徒的妻.
哀求以利沙說.
「女僕人 我丈夫死了.
他敬愛耶和華是你所知道的.
現在有寨主來.
要娶我兩個兒子作奴僕」.
以利沙問他說.
「我可以為你作甚麼呢?」.
「你告訴我 你家裡有甚麼?」.
他說.
「婢女家中除了一瓶油之外.
沒有甚麼」.
以利沙說.
「你去 向你眾倫社謝空氣命.
不要小謝.
回到家裡 關上門.
你和你兒子在裡面.
將油倒在所有的氣命裡.
倒滿了的放在一邊」.
於是 婦人離開以利沙去了.
關上門.
自己和兒子在裡面.
兒子把氣命拿來.
她就倒油.
氣命都滿了.
她對兒子說.
「再給我拿氣命來」.
兒子說.
「再沒有氣命了.
油就止住了」.
婦人去告訴神人.
神人說.
「你去賣油還債.
所剩的 你和你兒子.
可以靠著度日」.
繼續我們永闖.

$^{41}$丙堅永闖系列的第二講.
有人說其實.
失去這件事是.
我們人生成長的必修科.
大家同不同意嗎?.
很多寶貴的經驗.
你有沒有發覺.
其實都從失去開始.
大家一定經歷過.
你出生的時候失去了什麼?.
失去了在媽媽肚子裡.
很安穩的子宮.
你失去了那東西.
離開了.
但你得到的是什麼?.
得到的是.
爸爸媽媽或者家庭成員.
對你近距離埋身的照顧.
你經歷到人與人之間的溫暖和愛.
就正如我們今天所說的.
列王記這個故事.
你會發現其實都是一個.
失去的故事.
由失去開始.
卻是在失去的過程裡.
他這一家人經歷了.
上帝給他們奇妙的帶領.
其實你有留意.
其實整個列王記下的第四章.
都在說以利沙行的神蹟.
一共有四個.
在這些神蹟裡.
其實都與一些小市民的生活困難.
息息相關.
當事人不知道.
如何面對眼前的困境的時候.
以利沙就總是憑著信心.
將那些遇到困難.
有需要的人.
帶到上帝的面前.

$^{81}$然後上帝又用意想不到的方法.
來彰顯他如何與人同在.
一方面解決了人所面對的危機.
同時也讓人更加認識.
上帝那種真實和可靠.
就好像剛才經文所讀到.
那個四人家庭.
你會發覺裡面有爸爸.
爸爸是一個怎樣的人呢?.
是先知以利沙的門徒.
經文裡說.
他是一個敬畏上帝的人.
但很可惜.
他已經離世.
回到天父的懷裡.
他有兩個兒子.
多大呢?.
經文沒有說.
不過你大概猜到.
可能用今天的類比.
就是小學生.
或者至老初中生的年紀.
因為他們是有能力.
幫媽媽做家務.
幫媽媽搬空氣泥.
來協助媽媽倒油.
但似乎還未到.
可以出來工作.
謀生賺錢的年紀.
否則他們可以怎樣呢?.
他們就可以有自己的工作.
來還債.
所以有兩個兒子.
不是很大.
但也幫得上忙.
至於這個家庭.
這個敬畏上帝的家庭.
為什麼會欠下一筆債項呢?.
大概一定不是因為.
生活的揮霍.

$^{121}$也不是因為有什麼不良行為.
很可能是因為.
主力賺錢的那位爸爸.
患病.
本來是經濟支柱.
現在可能沒有辦法工作.
來賺錢養家.
而且還要花錢去醫病.
一家人就唯有.
借債來生活.
所以.
媽媽這個時候.
就很辛苦了.
她的丈夫離開了.
失去了摯愛的同時.
全家人的擔子.
也落在她的身上.
再加上經文這裡說.
債主臨門.
將要帶走他們兩個的兒子.
來做奴隸抵押.
事實上.
舊約律法裡面是容許欠債者.
來將自己或自己的兒女.
賣去做奴僕.
來償還債務.
當然.
做奴僕是有上限的.
根據律法.
只可以做奴僕六年.
到第七年.
債主就要讓奴僕.
來去重獲自由.
債務也一筆勾銷.
所以.
對當時很多欠債的人來說.
這個可能是一個很公道.
很合適的處理方法.
但是.
媽媽是不是這樣想?.

$^{161}$這個婦人很明顯不是這樣想.
我們看到.
是吧?.
失去丈夫.
慘不慘?.
很慘.
自己又沒有能力還債.
又救不了兩個兒子.
眼巴巴看著.
他們即將被人帶去做奴僕.
你是媽媽.
你捨不捨得?.
不要說媽媽.
爸爸也不捨得.
不要說做六年.
就算只是六個月也捨不得.
因為那種感受實在太難受.
一下子.
好像失去了所有家人.
對這個婦人來說.
確實是很陰公.
很淒涼.
而在她走投無路的情況下.
她想起.
她丈夫的主人.
先知爾麗莎.
於是她向他哀求.
她說.
你的僕人.
我的丈夫死了.
你知道他是敬畏上帝的.
現在債主來了.
要拿走我兩個兒子.
來做僕人.
做奴僕.
大家聽到嗎?.
婦人的哀求.
多多少少夾雜著很複雜的情緒.
譬如有很無助.
又或者對現況.

$^{201}$對生活.
有一些不滿.
一個敬畏上帝的家庭.
難道真的要弄成這樣?.
當然我們會明白.
婦人似乎不是.
希望爾麗莎幫她還債.
身為先知門徒的妻子的她.
大概都會知道.
爾麗莎也不會有很多錢.
婦人只是想爾麗莎理解.
她自己的慘況.
並且協助她尋求上帝的介入.
來面對全家人這一刻的危機.
而我們看到爾麗莎.
她也用兩個問題來回應.
她首先問.
我可以為你做些什麼?.
需要問嗎?.
最好當然是.
幫她還錢.
幫她處理了債務的危機.
讓她兩個兒子不用做奴隸.
其實很清楚的.
不用問的.
事實上.
爾麗莎她提出的.
是一個我們說的.
好像用了修辭技巧的問題.
她不是期望婦人給她一個答案.
她其實是在表達.
她自己對婦人的關懷.
她很願意和她們一家人同行.
我可以為你們做些什麼好呢?.
不如我們一起想想.
於是你會看到.
第二個問題就馬上出現.
告訴我.
你家裡有什麼.
這是爾麗莎第二個問題.

$^{241}$目的是引導這個婦人.
試想想.
家裡還有什麼資源.
可以幫你渡過難關?.
婦人的回答是.
婢女家中除了一瓶油之外.
沒有什麼.
和學本這裡的翻譯.
其實比較客氣.
原本實際.
是在說.
家裡什麼都沒有.
只剩下一瓶油.
你會看到其實.
婦人的回答其實是情緒先行.
她雖然仍然自稱是婢女.
在先知爾麗莎面前.
仍有尊重和尊卑之分.
但是多多少少.
我們會看到.
她的回應隱藏著.
心中那份焦慮.
甚至乎有一些怨憤.
我的丈夫敬畏神.
聽先知的吩咐教導.
但是他死了之後.
現在什麼都沒有.
還要留下一大筆債務.
我兩個兒子都快沒了.
你還問我家裡有什麼.
如果有.
我就不需要來找你幫忙.
沒有了.
家裡沒有什麼可以幫到我.
弟兄姊妹確實.
如果我們可以再來一次.
我們可能不會問這個問題.
先知在這個時候.
問婦人家裡還有什麼.
給人感覺好像不懂得人情世故.

$^{281}$人家已經很慘了.
家途四壁.
走投無路.
以利沙很應該主動提供協助.
特別是一些實際的協助.
譬如可以怎樣做.
以利沙有這麼多徒弟.
這麼多門徒.
可以叫齊他們.
大家一起來籌集資金.
幫這個婦人還一筆債.
解決她眼前的嫌棄之急.
又或者.
以利沙身為先知.
可以出面直接和債主交涉.
不如寬限多一段時間.
等小朋友長大了.
他們可以出來工作.
就可以還錢.
又或者.
他可以.
為婦人一家來祈禱.
求上帝施行神蹟.
直接賜下金錢.
幫這個婦人一家還債.
弟兄姊妹確實.
以上的建議.
都是好的.
不過在這一刻.
卻不是上帝給以利沙的感動.
我們知道上帝的確.
可以從無變有.
他可以直接.
賜下足夠的金錢.
給婦人還債.
解決她眼前的危機.
不過我們明白.
有些時候.
上帝不用這個方式.
上帝反而會引導我們看一看.

$^{321}$即使面對困難.
看看你手上.
還有什麼在手.
看清楚一點.
然後上帝就會運用.
這些資源.
來和我們一起面對困難.
不單止這個故事.
其實新約我們也很熟悉的.
「五餅二魚」.
大家也有印象吧.
其實也是這樣的.
最初門徒一口咬定是什麼?.
沒有啊!耶穌.
我們手上沒有資源.
來餵飽這麼多人.
不過.
耶穌是怎樣回應的?.
你看到耶穌沒有放棄.
耶穌是鼓勵他們.
看看你手上有多少食物.
然後最後.
門徒找到的是.
我們也很熟悉的.
「五個餅」.
「兩條魚」.
耶穌就透過.
這僅僅.
微小的資源.
並且讓門徒一起去參與.
「五餅二魚」這個神蹟.
結果最後我們也很清楚.
不單止餵飽了五千人.
當門徒收拾的時候.
他們發現了多少籃子食物?.
十二籃子的食物.
足夠供應他們日後上路.
來使用.
同樣我們看到.
這個婦人家裡.

$^{361}$雖然只剩下一瓶油.
相信不會很多.
微小的資源.
但卻足夠在上帝手中.
成就她的工作.
「以尼沙」就說.
「何丹,你去吧」.
「向你的鄰舍」.
「中鄰舍」是否俗話本譯?.
其實原文是說.
「所有的鄰舍」.
「向你所有的鄰舍」.
「來去借空的氣命」.
不要少借.
其實意思是.
盡量多借.
正面很多.
然後回到家.
關上門.
你和你的兒子在裡面.
將油倒在所有的氣命裡.
倒滿了就放在一邊.
「以尼沙」的指示很清晰.
不需要解經我們都明白.
不過問題是.
「頂子妹」如果你是「以尼沙」.
你夠不夠膽說出口?.
夠不夠膽?.
你夠不夠膽跟婦人說.
「你有一瓶油」.
想清楚.
「你有一瓶油」.
「你去向王浸所有的頂子妹」.
「借空瓶回來」.
盡量借.
然後你把空瓶拿回家後.
關上門.
然後你自己去倒油.
倒滿了就放在一邊.
其實匪夷所思嗎?.

$^{401}$其實不容易接受.
因為不是你理性邏輯.
慣用的.
想法.
不合常理.
所以我自己.
真誠去問自己.
我是「以尼沙」的話.
其實我都不夠膽說出口.
因為一開口對方會怎樣?.
可能會罵到我狗去淋頭.
這樣都算是一個方法?.
你胡說什麼?.
不幫忙.
也不要耍我.
所以我相信.
「先知以尼沙」這個吩咐.
必定是來自上帝給他的指引.
他是確信上帝會透過一個.
這麼奇特的方式.
來拯救這一家人.
如果不是.
大概不夠膽開口.
真的是空頭支票.
你覺不覺得?.
是上帝保證.
他才夠膽開口.
事實上.
你見到「以尼沙」這個吩咐.
其實暗示了一件事.
借到多少空氣的氣泥.
借得越多的話.
婦人得到的油會怎樣?.
就會越多.
婦人可以從上帝手上.
支取到多少供應.
取決於.
她對先知的吩咐.
信心到底有多大.
所以.

$^{441}$來到這裡我們發現.
「以尼沙」的任務.
完成了沒有?.
基本上已經完成了.
上帝要託付他所說的話.
他已經說完了.
你會見到稍後他沒有參與.
借氣泥和倒油的工作.
事情之後怎樣發展.
神蹟會不會出現.
就要看.
婦人的一家是否完全信靠.
信服「以尼沙」的指示.
之後我們很快便看到.
婦人得到先知的指引之後.
她立即行動.
她沒有半點猶豫.
其實弟兄姊妹這裡反映了一種.
很值得我們學習的屬靈素質.
確實到這一刻.
她不知道油到底實際是怎樣來的.
她亦不知道.
是否真的會有源源不絕的供應.
但是.
她既然求問得先知幫手.
就不應該質疑先知的指引.
問得.
就不要爭辯.
否則.
找先知來做甚麼呢?.
這是婦人反映了一種很好的屬靈素質.
就等於我們求問上帝的帶領的時候.
當上帝指引了的時候.
雖然我們不太明白.
甚至有懷疑.
但是.
如果我們不照跟的話.
那又求問上帝來做甚麼呢?.
所以你看到這個婦人.
她很單純,很簡單.

$^{481}$就是.
「一於照祝先知吩咐」.
先做了.
做了之後.
再看上帝的帶領.
她要做的當下.
就是首先盡她的努力.
向她所有的鄰舍.
借凶的器命回來.
你看到.
她的信服.
她的信心.
可以說是這次神蹟出現一個很重要的關鍵.
坦白說.
婦人失去了丈夫.
又債務纏身.
快要失去兩個兒子.
她有足夠的理由來抱怨她的人生的際遇.
她可以本皮.
可以自怨自艾.
但是她選擇的是.
運用對上帝的信靠.
信服上帝透過先知給她的帶領.
結果.
她這一份對上帝很單純的倚靠.
就讓她經歷到.
神蹟的出現.
確實弟兄姊妹原來經文告訴我們.
要體會上帝的同在和帶領.
和我們信主多久沒有絕對的關係.
在乎於.
當下你信不信上帝的話語.
信不信得過上帝總是有活人之路.
所以你將你的人生.
壓注在上帝的吩咐裡面.
經文是這樣說的.
我還記得認識一位弟兄在以前的教會.
他曾經分享過他很想.
放下他的工作來去蒙召.
入神學院裝備做傳道人為主事奉.

$^{521}$但他一直都.
持不住主意.
他一直求問上帝.
他覺得上帝應該是這樣帶領我.
但他又擔心.
我到底會否收錯上帝的指引.
於是他有一次分享的時候.
我就跟他說.
你放心.
上帝不只是感動你.
會感動教會.
感動你身邊的人.
來告訴你.
到底他是否想你放下工作.
蒙召來服侍他.
你要做的首先第一步是甚麼?.
先遞表.
先報名.
你先照做.
上帝就會有下一步的指引.
確實,丙姐妹.
人生.
有時不如意的事.
真是十常八九.
其實你信不信耶穌都一樣.
有沒有發覺?.
不會信耶穌之後少了病痛.
不會信耶穌之後少了.
哀傷難過的時候.
但如何在這些逆境裡分別出來.
就是對生活.
我們少一點抱怨.
少一點疑惑.
多一點感恩.
確信上帝不會丟下我們.
無論甚麼時候.
祂總有妥善的安排.
從而我們懂得將你的憂慮.
將你一些需要交託給祂.
讓祂去帶領.

$^{561}$然後你去配合.
這正是今日經文裡.
給我們的提醒和挑戰.
至於這裡.
先知為何要求婦人.
和她兩個小朋友.
借了一些空氣命之後.
回家要關上門.
為何要這麼神秘.
不要讓人知道?.
有沒有想過?.
我們不難想像這幅圖畫.
當婦人向她所認識的.
所有街坊借氣命的時候.
無論是出於好奇.
或者關心.
街坊都可能會問.
你借這麼多氣命來做甚麼?.
不如我幫你拿回家.
順便去婦人家裡.
了解一下到底發生了甚麼事.
假如在這個時候.
婦人或者小朋友.
不經意地說出口.
其實是先知以利沙.
吩咐我們將這瓶油.
倒滿所有的氣命.
你猜會不會有人有反應?.
我們說人多口窄.
怎麼可能?.
你這樣也相信?.
如果這樣的話.
婦人一家的信心.
難免會被動搖.
所以關上門.
其實是說.
有時候我們需要.
不要讓外在的聲音.
影響你對神的信靠.
不要讓外在的環境.

$^{601}$打擾了我們和上帝的互動.
我們和上帝的同工.
這個美好的時刻.
正如婦人這一家.
我曾經聽過一位弟兄說過.
他說其實他自己的家境.
不算很富裕.
有一段日子他曾經.
經歷過很嚴重的經濟困難.
但他仍堅持.
他要做十一奉獻.
很多人都勸他.
出於關心.
你錢都不夠用.
不如暫時不做.
你自己的家庭.
不夠用.
不如暫時先不要做奉獻.
等你日後環境好一點.
你才補償也不遲.
他很明白.
別人是一番的好意.
但他卻認定.
奉獻不是取決於.
自己夠不夠用.
而是取決於.
他對上帝的敬畏.
和對上帝的信靠.
關鍵不是他有多少錢.
關鍵是無論他有多少錢.
他確信上帝的供應是足夠的.
所以他仍堅持在艱難的日子.
不停他的十一奉獻.
最終上帝也帶領他.
渡過了難關.
確實真的,弟兄姊妹.
有時我們面對信心的挑戰考驗時.
我們需要像婦人一樣.
把門關上.
不要讓太多人知道.

$^{641}$動搖了我們對上帝的信靠.
到了第六節.
經文很清楚地說.
當所有借回來的器皿.
都裝滿油之後.
油是怎樣的?.
停止了.
神蹟完了.
你是婦人,你會怎樣?.
你是婦人和小朋友.
所有的器皿都裝滿油.
家裡有很多油.
於是我們應該怎樣做?.
我們應該賣油.
這次上帝賜那麼多油給我們.
我們賣掉它.
就有足夠的錢來還債.
婦人是否這樣做?.
經文有一個細節值得留意.
婦人沒有立即歡喜若狂.
賣掉所有油.
還債.
然後開心地和兩個兒子一起生活.
為什麼?.
因為油是誰的?.
油是上帝的.
不是這個婦人.
確實婦人很努力地.
按照先知的指示.
盡力把空的器皿借回來.
但是器皿裡的油.
由始至終是屬於上帝.
所以.
婦人另一個很好的屬靈素質.
不單止她對上帝很有信心.
她也對上帝的供應.
不敢擅自作主.
來運用.
當油裝滿的時候.
油止住了的時候.

$^{681}$她第一時間向先知回報.
求問他下一步怎樣做.
然後你見到先知說.
你將這些油賣了還債.
剩下的.
你和你的兩個兒子.
好好地過活.
所以你看到.
神跡的焦點.
其實不是說那些裝滿的油.
不是那些器皿.
而是婦人的信心.
以及她所擁有的智慧.
就是她雖然貧窮.
雖然在艱難的裡面.
但她仍然保持一份.
敬畏上帝的心.
她知道那些東西不是她.
是上帝.
所以問了上帝.
她才去用.
事實上我們知道.
猶太人聰明.
我們也知道.
很多猶太人都拿到諾貝爾獎.
有人統計過.
猶太人在全世界.
只佔0.3個百分比的人口.
不過有13\%的諾貝爾獎.
都是猶太人拿到的.
很有可能和他們.
從小教導小朋友的法則有關.
他們認定.
智慧比一切都重要.
幾乎所有猶太家庭的小朋友.
在小時候都被人問一條問題.
假如有一天.
你家大火燒著了.
你帶什麼走?.
他們每一個都會被人問這條問題.

$^{721}$假如你家大火要逃難了.
你會帶什麼走?.
如果小朋友回答的是.
金銀珠寶,錢等等.
媽媽就會跟他說.
記住.
你要帶走的不是錢.
也不是珠寶.
而是智慧.
因為智慧.
沒有人可以從你身上拿走.
只要你活著.
智慧就會永遠伴隨我們.
這是猶太人從小教導小朋友的法則.
我們很明白.
什麼是智慧.
聖經裡說的.
可能你也會立即背.
敬畏上帝.
就是智慧的開端.
或者所謂的總綱.
就像這四人家庭.
一開始我們看到.
爸爸是一個敬畏上帝的人.
那媽媽呢?.
媽媽原來也是.
在這件事件裡.
反映了她對上帝的一份敬畏.
與錢無關.
與困難無關.
與平日的生活如何耳濡目染有關.
其實我們還記得.
在這個故事開頭.
媽媽是否覺得很憂愁?.
覺得很絕望.
先知問她家裡還有什麼的時候.
她就回答.
「沒有什麼」.
只剩下一瓶油.
到故事的結束.

$^{761}$「沒有」這兩個字.
有沒有發覺又再出現?.
這次是兒子跟她說的.
這次的「沒有」.
再不是絕望,憂愁.
而是充滿喜樂和盼望.
原因是.
所有的氣味都裝滿了油.
再沒有氣味了.
你看到婦人在患難的裡面.
當她覺得什麼都沒有的時候.
上帝就透過這個神蹟提醒她.
「不是」.
「看清楚一點」.
其實還有很多很寶貴美善的事物.
環繞著你們一家.
首先有什麼?.
這個家裡.
有一個很好的爸爸.
很好的丈夫.
他不單是以利沙的門徒.
他也是一個敬畏上帝的人.
他雖然現在不在了.
但是他留下了一個很好的.
屬靈傳統給他的家人.
為什麼家裡遇到困難的時候.
做媽媽,做妻子的.
懂得尋求以利沙的幫助?.
很大可能是因為.
她平時見丈夫都是這樣做的.
我丈夫有事的時候.
他不會單打獨鬥.
他會找一群門徒.
來一起尋求上帝的幫助.
所以,弟兄姊妹.
特別我們做人丈夫的.
做人爸爸的.
這段經文很有意思.
讓我們提醒.
我們的屬靈誡令很重要.

$^{801}$特別在有困難的時候.
你的屬靈生命.
能夠成為你家人.
一個穩定的力量.
一個不可或缺的力量.
你能夠帶領你的家人.
一起依靠上帝.
一起尋求上帝的幫助.
一起尋求教會.
弟兄姊妹的指引.
不單止有一位好丈夫.
有一位好爸爸.
這個家庭還有什麼?.
還有什麼?.
有一群很好的街坊.
有一群很好的鄰舍.
圍繞著這個家庭.
第一節這裡.
經文沒有翻譯到神髓.
因為翻譯了神髓會很掉漆.
這裡其實是說.
眾先知門徒的妻子裡的.
其中一位妻子.
其實很掉漆.
第一節不是說.
只有一個先知門徒的太太.
而是有很多先知門徒的太太.
在裡面的其中一位.
換句話來說.
婦人其實不孤單.
不過可能因為.
慘況的緣故.
情緒的緣故.
她看不到.
她有屬靈群體.
有一群街坊在她身邊.
支持她.
有很多的好姐妹.
都願意出來幫忙.
所以.

$^{841}$以利沙.
向這個婦人的提醒就是.
你盡量向你的鄰舍.
向你的好姐妹.
來盡量借凶的平子.
只要每個人運用.
手上上帝給他們的資源.
就能夠幫助到.
你一家渡過危機.
就能夠看到神跡的出現.
確實沒有人願意面對.
失去的場景.
失去的功課.
但是上帝卻可以將失去.
化成一種祝福.
讓人與人之間學懂聯繫.
讓人更加珍惜.
彼此的守望和關懷.
所以上帝姐妹.
我相信.
假如今天的經文.
是發生在黃浸的話.
發生在我們這個屬靈群體的話.
我相信大家.
願不願意將你家裡空的器皿借出來.
願不願意?.
有點猶豫.
沒有的話,你會怎樣?.
你會借一個給他.
甚至乎弟兄姊妹好心.
買一個回來說.
「夠的,夠的」.
你怎樣也會借一個給他.
確實真的弟兄姊妹.
上帝往往是透過我們每一個人.
去多走一步.
來祝福身邊有需要的人.
幫他們渡過難關.
第三.
不單止有丈夫.

$^{881}$也有好姐妹.
這個家庭還有什麼?.
還有兩個兒子.
能不能幫忙?.
能幫忙的.
其實他們也有參與神蹟的發生.
兩個兒子.
一個幫忙拿器皿.
可能裝滿之後.
有氣力的那個就放在一旁.
他們也是上帝所賜的產業.
我們俗稱所謂信義代.
他們也有份參與這次神蹟.
親身體驗上帝的同在.
他們得到信心成長.
所以弟兄姊妹.
我們也不妨製造多一些機會.
給我們的小朋友.
給年輕人.
參與上帝的工作.
讓他們也能親身經歷上帝的同在.
真的.
其實人生有很多寶貴的功課.
很多時候都是從小學起.
從小就給他們機會體驗.
而不是純粹說給他們聽.
最後.
這個家庭有什麼呢?.
當然是上帝.
上帝從來沒有離開過這個家庭.
總是顧念他們的需要.
垂聽他們的呼求.
確實在事件的裡面.
婦人失去了丈夫.
兩個小朋友失去了父親.
以利沙失去了好的徒弟.
但是對上帝而言.
卻是將祂忠心的婦人.
接回祂的身邊.
釋去祂在地上的奴婦.

$^{921}$確實是會讓人傷感.
上帝很明白.
上帝亦都連述.
所以.
祂親手從這個家庭收回的.
祂自己親自來填補.
裡面的傷痛.
裡面的空洞.
所以我們見到這個故事.
雖然是以死亡來開始.
但結束的其實是一個活著.
盼望的訊息.
因為神沒有丟下他們.
所以弟兄姊妹.
假如上帝亦都在你身上.
收回了一些很重要的事情.
可能是一些人.
一份工.
一些供應.
可能你現在仍然在憂傷的裡面.
仍然未能釋懷的時候.
願意今日的經文.
上帝都使用.
來為你送上醫治和安慰.
亦都給予我們有生活的力量和盼望.
因為上帝必定用祂自己.
來代替祂在我們身上所收回.
讓我們同心祈禱.
天父我們知道.
你必與我們同在.
確實在世上.
我們會經歷到艱難.
正如耶穌所說.
但同時你會賜下足夠的平安.
你會幫助我們.
上帝我們確實人生會經歷很多失去的功課.
在裡頭有很多不願意.
在裡頭有很多傷心難過.
甚至到這一刻.
我們仍然未痊癒.

$^{961}$上帝我們知道.
求你用我們明白的方式.
來幫助我們.
指引我們.
求你使用我們身邊的人.
來互相守望.
讓我們在患難的裡面.
經歷到和人之間的結連.
我們體會我們原來不是孤單.
我們會有失去.
但在失去的裡面.
我們亦得到很多寶貴的功課.
我們體會到別人對我們的守望.
我們亦可以去守望別人.
更重要是我們明白.
上帝你從來沒有丟下我們.
並且幫助我們.
讓我們在不明白的裡面.
讓我們在難過的裡面.
我們仍然憑信心認定.
上帝是你本為善.
你的慈愛永遠長存.
到最後.
當我們見你面的時候.
上帝我們就知道.
這一切最終.
在你裡面都有完美的答案.
為此幫助我們帶著信心.
我們可能和那些艱難.
和那些困惑來共處.
讓我們仍然在其他方面.
有力的來生活.
享受你賜予給我們的.
亦努力守住.
你所託付給我們的使命.
讓我們可以在信.
在逆境裡面.
繼續與你同在.
繼續與你同行.
為此我們同心獻上祈禱.

$^{1001}$同基督耶穌你得勝的名字.
來祈求.
Amen.
主席.
\newpage



\section{路得記 1:6-18}
\label{sec:8sXdI3Lte_M}
\textbf{2022年6月26日崇拜直播｜陳妙濃傳道｜並肩勇闖系列:以恩慈相待｜路得記1\_6-18}
\newline
\newline
連結: \href{https://youtube.com/watch?v=8sXdI3Lte_M}{\texttt{https://youtube.com/watch?v=8sXdI3Lte\_M}} ~~~~ 語音日期: 2022-06-27
\newline
\newline
\hyperref[sec:_Anjyu4UPkQ]{\small{< < < PREV SERMON < < <}}
~
\hyperref[sec:index_chronic]{\small{[返順時目]}}
~
\hyperref[sec:c6nbM3DecyY]{\small{> > > NEXT SERMON > > >}}
\newline
\newline
路得記 1:6-18
\newline
\begin{longtable}{cl}
\hline
\hline
章節 & 經文 (和合本修訂版)\\
\hline
1:6 & \begin{tabularx}{0.7\textwidth}{X} 拿娥米與兩個媳婦起身，要從摩押地回去，因為她在摩押地聽見耶和華眷顧自己的百姓，賜糧食給他們。 \end{tabularx} \\ \\ \relax
1:7 & \begin{tabularx}{0.7\textwidth}{X} 她和兩個媳婦就起行，離開所住的地方，上路回猶大地去。 \end{tabularx} \\ \\ \relax
1:8 & \begin{tabularx}{0.7\textwidth}{X} 拿娥米對兩個媳婦說：「你們各自回娘家去吧！願耶和華恩待你們，像你們待已故的人和我一樣。 \end{tabularx} \\ \\ \relax
1:9 & \begin{tabularx}{0.7\textwidth}{X} 願耶和華使你們各自在新的丈夫家中得歸宿！」於是拿娥米與她們親吻，她們就放聲大哭， \end{tabularx} \\ \\ \relax
1:10 & \begin{tabularx}{0.7\textwidth}{X} 對她說：「不，我們要與你一同回你的百姓那裡去。」 \end{tabularx} \\ \\ \relax
1:11 & \begin{tabularx}{0.7\textwidth}{X} 拿娥米說：「我的女兒啊，回去吧！為何要跟我去呢？我還能生兒子作你們的丈夫嗎？ \end{tabularx} \\ \\ \relax
1:12 & \begin{tabularx}{0.7\textwidth}{X} 我的女兒啊，回去吧！我年紀老了，不能再有丈夫。就算我還有希望，今夜有丈夫，而且也生了兒子， \end{tabularx} \\ \\ \relax
1:13 & \begin{tabularx}{0.7\textwidth}{X} 你們豈能等著他們長大呢？你們能守住自己不嫁人嗎？我的女兒啊，不要這樣。我比你們更苦，因為耶和華伸手擊打我。」 \end{tabularx} \\ \\ \relax
1:14 & \begin{tabularx}{0.7\textwidth}{X} 兩個媳婦又放聲大哭，俄珥巴與婆婆吻別，但是路得卻緊跟著拿娥米。 \end{tabularx} \\ \\ \relax
1:15 & \begin{tabularx}{0.7\textwidth}{X} 拿娥米說：「看哪，你嫂嫂已經回她的百姓和她的神明那裡去了，你也跟你嫂嫂回去吧！」 \end{tabularx} \\ \\ \relax
1:16 & \begin{tabularx}{0.7\textwidth}{X} 路得說：「不要勸我離開你，轉去不跟隨你。你往哪裡去，我也往哪裡去；你在哪裡住，我也在哪裡住；你的百姓就是我的百姓；你的神就是我的神。 \end{tabularx} \\ \\ \relax
1:17 & \begin{tabularx}{0.7\textwidth}{X} 你死在哪裡，我也死在哪裡，葬在哪裡。只有死能使你我分離；不然，願耶和華重重懲罰我！」 \end{tabularx} \\ \\ \relax
1:18 & \begin{tabularx}{0.7\textwidth}{X} 拿娥米見路得決意要跟自己去，就不再對她說甚麼了。 \end{tabularx} \\ \\
[1ex]
\hline
\hline
\end{longtable}
$^{1}$今日的經文在路得記的第一章6到18節.
舊約聖經的331頁.
331頁路得記第一章6至18節.
翻開的請聽我讀.
他就與兩個兒婦起身.
要從摩押地歸回.
因為他在摩押地聽見耶和華眷顧自己的百姓.
賜糧食與他們.
於是他和兩個兒婦起行離開所住的地方.
要回猶大地去.
拿我米對兩個兒婦說.
你們各人回娘家去吧.
願耶和華恩代你們.
像你們恩代已死的人與我一樣.
第九節.
願耶和華使你們各在新夫家中得平安.
於是拿我米與他們親嘴.
他們就放聲而哭.
說不然我們必與你一同回你本國去.
拿我米說.
我女兒們啊 回去吧.
為何要跟我去呢.
我還能生子作你們的丈夫嗎.
我女兒們啊 回去吧.
我年紀老邁不能再有丈夫.
即或說我還有子望.
今夜有丈夫可以生子.
你們豈能等著他們長大呢.
你們豈能等著他們不嫁別人呢.
我女兒們啊 不要這樣.
我為你們的緣故甚是受苦.
因為耶和華伸手攻擊我.
第十四節.
兩個兒婦又放聲而哭.
我二伯與婆婆親嘴而別.
只是老德捨不得拿我米.
拿我米說.
看啊 你嫂子已經回他本國.
和他所拜的神那裡去了.
你也跟著你嫂子回去吧.

$^{41}$老德說.
不要催我回去不跟隨你.
你往那裡去 我也往那裡去.
你在那裡住宿 我也在那裡住宿.
你的國就是我的國.
你的神就是我的神.
你在那裡死 我也在那裡死.
也葬在那裡.
除非死能使你我相離.
不然冤野和華重重地降伐於我.
拿我米見老德定義要跟隨自己去.
就不再勸他了.
足夠法定人數.
各位弟兄姊妹早晨.
剛才的經文.
我們聽到有三個孤苦無依的婦女.
他們真誠的對話.
為什麼他們家中的男性全部缺席呢?.
我們可以看回上文一章一至五節.
原來當時是事先秉正的時代.
大概就是主前1200年至到1020年間.
一個本來住在伯利恆一個平凡的家庭.
一家之主他叫做爾利米勒.
他為了逃避當時的饑荒.
他就跟太太拿我米和兩個兒子.
一行四人移居去摩訶.
因為當時摩訶是有糧食.
可惜十多年後.
爾利米勒和兩個兒子都相繼去世.
就剩下三個寡婦.
沒有丈夫他們沒有兒子.
加上寡婦在當時來說.
還包括孤單被遺棄無助的意思.
他們是弱勢的群體.
我們不難想像他們面對生命無常.
悲歡離合經濟困境.
所處的社會是思時代.
聖經怎樣形容?.
他說當時各人任意而行.
我們簡單理解沒有往往.

$^{81}$加上外族侵襲所以內憂外患.
這三個寡婦他們艱辛他們無助.
他們不知所措.
今日各位弟兄姊妹.
只要我們留心用心.
我們看看自己或我們身邊的人.
我們發現有不少故事.
甚至不少衝擊.
可能是家人或我們自己.
有不舒服甚至有重病.
又或者我們突然被解僱.
婚姻,兒女問題.
莫名的情緒打擾.
社會的不公不義.
所以路得記不只是記錄一個大時代.
一個家庭,一個小富士.
他更加道出我們的心聲.
我們的掙扎.
讓我們看回經文.
三衣為命的他們.
為何他們打算回到伯利恆呢?.
我們看經文六至十四節.
他就與兩個兒婦起身.
要從摩押地歸回.
因為他在摩押地聽見耶和華眷顧自己的百姓.
賜糧食於他們.
於是他和兩個兒婦起行.
離開所住的地方.
要回猶大地去.
那奴彌對兩個兒婦說.
你們各人回娘家去吧.
願耶和華恩待你們.
將你們恩待已死的人與我一樣.
願耶和華使你們各在新夫家中得平安.
於是那奴彌與他們親嘴.
他們就放聲而哭.
說:不然.
我們必與你一同回你本國去.
十一節.
那奴彌說:我女兒呀.

$^{121}$回去吧.
為何要跟我去呢?.
我還能生子做你的丈夫嗎?.
我女兒呀.
回去吧.
我年紀老邁,不再有丈夫.
即話說我還有子夢.
今夜有丈夫可以生子.
你們豈能等著他們長大呢?.
你們豈能等著他們不嫁別人呢?.
我女兒呀.
不要這樣.
我為你們的緣故甚是受苦.
因為耶和華伸手攻擊我.
兩個兒婦又放聲而哭.
娥姨伯與婆婆親嘴而別.
只是.
奴得捨不得那娥迷.
第六節.
出現了一個字.
出現了眷顧.
眷顧這個字原本就有探望.
和查看的意思.
如果引申到的.
是有眷顧,賜福.
不過也由於試探望和查看.
也有責備,討罪的意思.
相比十多年前.
為了糧食.
那娥迷跟丈夫去到摩崖.
這次大家看到.
那娥迷的眼睛.
他不是再看到糧食.
他是看到賜糧食背後.
以色列的上帝.
他知道真正賜福恩待人.
是這位上帝的時候.
他就選擇回到伯尼恆.
回到上帝那裡.
但我們看第八節.

$^{161}$本來他們是喜行的.
為何那娥迷這麼快轉了心意呢?.
叫兩個媳婦回到娘家.
原本經文用了兩個很重要的動詞.
而且是命令式的.
就是「去」,「回去」.
那娥迷不想他們回到伯尼恆.
為何會用到娘家而不是父家呢?.
我們可以參考《創世記》二十四章.
當時伯尼恆的老僕人.
為了以撒娶妻子.
就遇見了李伯嘉.
李伯嘉清楚知道老僕的心意的時候.
原來《創世記》二十四章.
是說他回去.
回去和他媽媽.
和他媽媽的人談婚事.
所以原來聖經刻意用娘家.
代表那娥迷想他們回到摩押.
回到摩押.
能夠找到新的夫君.
能夠有嫁娶的事.
接著那娥迷祈求耶和華恩待他們.
「恩待」是路得記很重要的一個字.
意思是忠於文約的慈愛和信實.
既然是文約.
就是雙方面.
亦都特別重視大家的關係.
大家的態度和行動.
亦都帶有永不改變的意思.
當這個字用在上帝對人的恩待的時候.
是代表上帝對我們的恩愛,眷顧和上帝守約.
他答應了的.
他會守住.
當這個字形容人對人的恩待的時候.
是著重信任,恩情和毀身.
說到這裡.
我想問問大家.
為何那娥迷不留他們在身邊呢?.
試問三個寡婦.

$^{201}$其實誰最需要恩待呢?.
相比之下.
那娥迷是受到雙重的喪親之痛.
首先,她的丈夫死了.
現在到兩個兒子.
兩個年輕的寡婦還可以再分身子.
但是那娥迷沒有這個機會.
所以原來那娥迷不是善變.
而是她想清想楚.
她才有這個充分的勸告.
所以第八節我們看到第一個理由.
因為那娥迷回想這兩個這麼好的媳婦.
昔日對她的恩情.
接著,想跟大家看十三節.
十三節她這樣說.
「我為你們的緣故甚是受苦」.
因為耶和華伸手攻擊我.
「我為你們的緣故甚是受苦」.
聽起來不太舒服吧?.
這句話.
不過我們再看清楚意思.
是代表那娥迷想說.
「我比你們更苦」.
原來那娥迷親身經歷了多年喪夫之痛.
她心想.
「你們跟著我有什麼好處?」.
所以她不忍心兩個媳婦抱她後塵.
手寡.
加上她也擔心摩訝的媳婦.
在以色列能否立足.
所以她一而再再而三.
勸他們回到摩訝的娘家.
當然,到這一刻.
有人會問.
她會否以退為進呢?.
她會否另有謀算.
不是說真心話呢?.
但是剛才.
無論是讀經.
或者我和大家一起看經文.

$^{241}$我們發現無論第八節.
十一至十三節.
十五節.
那娥迷是四次勸勞你回去.
那娥迷是真心真意.
希望這兩個媳婦.
能夠找到一個好的人家再嫁.
重建自己的家園.
結果.
其中一個媳婦.
「娥爾巴」.
她接受了那娥迷的勸導.
真的回到娘家.
從那娥迷的行動和選擇.
我們看到她對上帝的心.
也看到她對媳婦的恩慈.
第一.
經文怎麼說呢?.
她只聽到上帝眷顧自己的百姓.
她就相信了.
而且在最困難的時候.
她沒有被困難困住.
她反而選擇察覺.
相信上帝仍然在.
於是她就憑信心.
回到上帝那裡.
經歷上帝重整和上帝的關係.
第二.
我們看到那娥迷.
本身都經歷了.
丈夫過世的傷痛.
她將心比己.
她體會兩個媳婦的痛.
的確.
如果按人來說.
有兩個這麼好的媳婦在身邊.
我很有大條道理留他們在身邊.
但是她處處為他們著想.
她寧願自己孤單獨居.
她不單動之以情.

$^{281}$更說之以理.
我們看13,14節.
13,14節其實就是參考.
《申命記》25章5至10節.
原來當時猶太人有個習俗.
如果一個媳婦.
丈夫過世的時候.
丈夫的兄弟需要娶她.
而且為死去的丈夫.
留下一個下一代.
但是娥迷一點都沒有用到.
這件事來威脅他們兩個.
她反而很坦誠表達.
即使她今天結婚.
馬上生育.
都要等二十多年.
才可以有一個兒子給他們.
所以娥迷想來想去.
這也不是出路.
她想那麼多.
她千方百計這樣去說.
是希望媳婦不要掛心.
她處處為別人著想.
再和大家看13節.
13節和大家看的是最後一節.
因為耶和華伸手攻擊我.
很奇怪啊!各位弟兄姊妹.
明明一開始我們一直聽的情況.
都是娥迷看到上帝的祝福.
希望上帝恩待他們.
為什麼突然間到了這一節經文.
她控訴上帝呢?.
有學者分析.
因為她基於她相信上帝.
她為兒婦祝福.
但是事實上.
她的經歷.
她真真實實體會到.
上帝是打擊她.
其實這個並沒有矛盾.

$^{321}$各位弟兄姊妹.
當我們有情緒.
我們有難過.
我們有傷痛的時候.
你是否常常用理性的角度去說服自己.
讓自己分身.
不要專注這些感受.
我們不是應該常常處於正能量的狀態嗎?.
或者我們信主多年.
我們會覺得我們憂傷.
我們痛.
我們軟弱.
就是不屬靈.
我們抑壓自己的情緒.
詩篇34篇18節.
她這樣說.
耶和華靠近傷心的人.
拯救靈性痛悔的人.
各位弟兄姊妹.
在這個時代.
我們要正視.
我們要給自己有抒發情緒的機會和空間.
哪怕你是什麼年紀.
你信主?.
你未信主?.
你在哪個階層?.
進一步看.
懂得放個識字.
大家懂得寫嗎?.
我看見大家用手掌識字.
懂得當個識字.
大家懂得寫嗎?.
懂得當個識字.
台灣作家.
亦同時是精神科治療師.
施以樂.
他在一篇文章裡這樣說.
原來我們不單要釋放我們的情緒.
更加要適當地抒發.
何謂適當抒發呢?.

$^{361}$吃東西我們要均衡.
要適當.
意思是我們的亮度.
我們要適可而止.
還有一樣東西.
我不知道大家同不同意.
你要找適合的人跟你談.
有時我們找一些人.
可能他不斷向我們點頭.
但可能他的建議未必是最全面.
但相反.
有些人可能我們少跟他談天.
但他們可能的閱歷.
他們熟悉的對上帝的了解.
會給我們一些正面意見.
我們不妨去找他們.
所以我們看到南俄米.
他在這裡我們看到.
原來我們要正視自己的情緒.
亦要釋放自己的情緒.
結果南俄米說完這番說話之後.
兩位師傅有沒有表示?.
有沒有?.
有呀!.
兩位都有表示.
其中一個怎樣?.
走了.
回去了.
俄爾巴他聽了南俄米真誠的勸告.
他做了一個較為穩妥的選擇.
他真的回到摩訶的娘家.
不過我們看到他走的時候.
他們大聲痛哭.
我們可見婆媳之情是何等深厚.
面對俄爾巴真的走了.
真的走了.
少了一個多年陪伴自己的好媳婦.
南俄米有沒有說半句怨言?.
有時我們說就這樣說.
心裡真的捨不得.

$^{401}$我想這裡有些媽媽明白吧?.
特別近年我們可能有些家庭離開了香港.
可能我們說就說.
我祝福你.
你去到不同地方.
繼續帶著神的使命去闖.
但是內心呢?.
可能就經常數著他走的日子.
然後內心不舒服.
甚至我們軟弱.
正視自己的情緒的時候.
可能我們會有怨言.
會控訴.
甚至不知道南俄米會不會跟爸爸說幾句話.
我都說了.
你看到了.
但是聖經有沒有說各位弟兄妹妹?.
沒有.
原來真正我們去恩待別人.
就是我們要真心真意去祝福別人.
不過呢.
我們留意.
俄爾巴這個回去.
他就回到摩訶當地的神寺裡.
這是事實.
接下來我們看看路德.
既然剛才我說恩待是建基於雙方的關係.
我們看到南俄米對息婦們的恩待.
路德又怎樣對南俄米呢?.
我想跟大家一起讀一章十四至十八節.
如果大家的聖經或程序表都有.
如果聖經就是331頁.
好,一二三.
俄爾巴與婆婆親嘴而別.
此事路德捨不得俄米.
南俄米說:看啊!你嫂子已經回他國和那拜的神那裡去了.
你也跟著你嫂子回去吧.
路德說:不要催我回去不跟隨你.
你往那裡去,我也往那裡去.
你在那裡住宿,我也在那裡住宿.

$^{441}$你的國就是我的國,你的神就是我的神.
你在那裡住宿,我也在那裡住宿.
也葬在那裡,除非死能使你我相離.
不然,怨耶和華重重地降伐與我.
南俄米見路德定義要跟隨自己去.
就不再勸他了.
大家可以繼續拿著這幾個經文.
因為我想跟大家看看裡面的結構.
第十六節:不要催我回去不跟隨你.
十七節正正回應這裡.
路德說:如果我離去.
怨耶和華重重地降伐與我.
然後十六節.
你往那裡去,我也往那裡去.
你在那裡住宿,我也在那裡住宿.
我們看回十七節.
你在那裡死,我也在那裡死.
也葬在那裡.
剛才和大家看到.
十六,十七節是互相呼應.
不過最重要的是哪一節經文呢?.
各位弟兄姊妹.
最重要哪一節經文呢?.
我聽到了.
最重要的是中間.
你的國就是我的國.
你的神就是我的神.
各位弟兄姊妹.
聖經巧妙地寫出這個結構.
讓我們看到這是整段的核心.
路德這個回覆.
整段的核心也是中心的部分.
就是你的國就是我的國.
就是說以色列人是你的同胞.
也是我的同胞.
你所信的上帝.
婆婆,你所信的上帝.
就是我所信的上帝.
甚至他用咒詛式起誓.
表明自己要跟隨拿我米這個決心.

$^{481}$他這個聲明不是隨口說的.
要付很大的代價.
甚至犧牲.
包括什麼呢?.
包括本來他可以跟另一個媳婦一樣.
他可以留在摩押.
年輕的他可以找到新的夫家.
這是關係.
一個幸福.
一個女士的終身幸福.
這個犧牲是什麼呢?.
他要移居一個陌生的地方.
伯利恆.
作為一個異地的外族人.
他沒有地位,沒有身份.
而且很大的冒險.
各位聽到摩押,聽到什麼呢?.
聽到申命記23章3至6節.
民數記25章1至2節.
摩押從來都是一個負面的形象.
因為摩押人引誘以色列人犯罪.
敬拜別臣.
所以拿我米回去.
所以路德聽拿我米的勸告回去.
路德是很冒險.
他有犧牲.
亦是一個不容易的選擇.
但他卻願意.
所以我們看到第一.
路德發現原來人生真正的保障.
不是丈夫.
不是摩押地.
不是新的夫君.
亦不是摩押地的神.
是拿我米所信的上帝.
於是他就一條心.
決心,憑著信心.
放下自己的信仰.
他別無他選.
他一心一意跟隨婆婆的信仰.

$^{521}$他歸信委身於以色列的上帝.
這成為他人生的轉捩點.
看到台下的弟兄.
不知道大家還記得.
跟你另一半有盟約的時候.
是否說過怎樣也不分開呢?.
無論病苦等等.
十七節:「除非死能使你我相離」.
路德故事動人的地方就是.
作為一個媳婦.
他是徹底對拿我米不離不棄的跟隨.
經文周刊沒有特別記載.
但如果大家有聖經.
可以看回第21節.
我讀給大家聽.
這是拿我米最後一章的說話.
「我滿滿的出去,而我話使我空空的回來.
而我話過沒於我,全能者使我受苦.
既是這樣,你們為何還叫我拿我米呢?」.
拿我米是甜的意思.
剛才的經文是拿我米.
她真的跟路德兩個人回到伯利恆.
但聖經記載她的話.
她說她是空空的回來.
各位弟兄姊妹.
拿我米是否空空的回去呢?.
她身邊有一個對她不離不棄的媳婦路德.
但她能否看到他呢?.
她看不到.
當我們拼命為一個人的時候.
我們在她身邊.
但對方當我們隱形的時候.
各位弟兄姊妹,你心情如何呢?.
你會否一直覺得很不舒服呢?.
不過經文很奇怪.
何謂對人的恩慈呢?.
路德不離不棄地跟隨拿我米.
路德沒有受影響.
所以第二章記載路德把握機會.
為他和自己的婆婆拿我米求生計.

$^{561}$他沒有在家裡坐著.
他主動去農田工作.
這是出自於真摯的愛.
是一個毀身,是一個忠誠.
我們充分看到他對拿我米不離不棄的恩情.
說到恩慈.
希望大家可以閉上眼睛想一想.
其實我們在逆境當中.
我們容不容易看到上帝的恩慈呢?.
剛才我們看到人與人之間的恩慈.
上帝恩慈我們能否看到呢?.
請聽我讀一章廿二節,二章一節.
他們到了伯尼衡.
正是開始收割大物的時候.
二章一節.
拿我米的丈夫爾利米娜的親族中.
有一個人名叫波阿斯.
是個大財主.
各位電影姐妹.
原來在一章最後,二章第一節.
早已埋下伏筆.
我們看到.
從來上帝都是透過不同的人.
不同的事介入我們的生命.
叫我們從悲劇當中走出來.
在這裡.
上帝透過不同的人和事.
介入拿我米和路德的生命當中.
他改寫他們的悲劇.
顯出他們的豐盛恩慈.
一開頭.
路得記一開頭第一章.
曾經出場的男士全部死了.
但是第二章一開始.
就記載波阿斯.
這個不是一個普通的人.
這是誰?.
他是爾利米娜有血緣關係的親人.
大財主.
我們聽下去.

$^{601}$我們容易理解.
他在經濟上幫到他們.
不過我們看深一點.
波阿斯這個名字更重要的是代表能力.
有能力就能解決困難.
由於路得記是大家比較熟悉的聖經.
所以我都省略接下來的內容.
但我們從第二,第四章.
我們發現波阿斯不單止為路德拿我米解決困難.
而且是超越了舊約摩西律法所規定的.
簡單和大家看看第二章.
當時路德去一塊田地撿墨水.
這麼尷尬這個田地是誰的?.
波阿斯,大家很熟悉聖經.
利尾記二十三章二十二節.
有這樣的記載.
「你們在自己的地收割莊稼時.
不可割盡田的角落.
也不可拾取莊稼所吊內的.
要把他們留給窮人和寄居的.
我是耶和華你的上帝」.
在舊約聖經裡,上帝特意有這個規條.
上帝憐憫有需要的貧窮人.
不過,波阿斯怎麼做呢?.
波阿斯在路德的二章十五至十六節裡.
都說阿斯吩咐那些僕人.
你們任由他吧.
如果路德一時之間不知規矩.
他拿了沒有掉下來的,你也讓他拿.
甚至還特意掉了一些在地上.
讓路德拿.
僕人當然聽主人的話.
剛才一開始說拿訛米的時候.
我有提到當丈夫過身的時候.
他的兄弟有需要娶太太.
不過,這個律例的意思是兄弟取為相.
不過,大家如果看清楚路得記.
波阿斯是不是路德先夫的直系弟兄?.
不是.
原來波阿斯可以拒絕路得記第三章的要求.

$^{641}$當路德求波阿斯對他照顧.
讓他成為他的太太的時候.
其實波阿斯是不需要做的.
所以各位弟兄姊妹.
波阿斯和路德的結合.
原來不只是履行律法這麼簡單.
也不是因為無可推搪.
是近親的責任.
其實是波阿斯極之極之愛路德對他的恩情.
波阿斯對路德的恩慈.
令到路德和婆婆拿訛米走出困境.
正正就反映上帝對人的豐盛憐憫和眷顧.
超乎我們所想.
最後剛才提過.
路德是一個魔壓的女子.
是一個外族.
是一個外邦人.
是一個犯罪的地方的外邦人.
但他同樣能夠相信上帝.
我們可見上帝從來沒有棄絕任何的人.
上帝的救恩是臨到每一個弱勢社群.
每一個人當中.
重要的是我們回轉.
所以今天我想問問大家兩件事.
第一,你現在是否落在困境呢?.
你身邊有沒有人對你恩慈相待呢?.
他沒有挖苦你.
他對你不離不棄.
他為你著想.
他帶你去認識上帝.
相反來說.
各位弟兄姊妹,今天是六月最後一個星期.
在過往日子裡,特別是疫情.
在社會事件之後.
弟兄姊妹,你有沒有對你身邊的人有恩慈相待呢?.
世界實在太多不幸.
路德的故事從來都不是一個.
只是發生在聖經的故事.
但上帝從不離場.
他知道亦聽到人間的疾苦.

$^{681}$他希望透過我們去恩待有需要的人.
弟兄姊妹,你有沒有這些人在你心上呢?.
心動不如行動.
在這裡我希望大家在下半年.
如果神讓你有這些人在心裡.
你就恩待他們.
你就將上帝的愛帶給他們.
讓他們跳出這個悲劇.
讓他們的生命重得恩典.
2020年7月29日.
曾經有一段這樣的新聞.
食肆首日全面禁堂食.
多人午飯時間在室外公眾地方用膳.
當時香港疫情嚴重.
全日禁堂食.
記者訪問清潔工人張女士.
你們在哪裡吃東西呢?.
張女士她這樣說.
我們不會回辦公室留食.
我們會在垃圾站旁邊吃東西.
當我看到這段新聞的時候.
我自己禱告.
作為一個,不要說是牧者.
作為一個信徒.
我看到這段訪問的時候.
我怎麼做呢?.
於是我禱告.
但因為我膽小.
熟悉我的人都知道我膽小.
我就想找人陪我一起去買些飲品.
送給一些清潔的工人.
大家明白我完全沒有一個職業的.
看低他們.
但那天早上下著大雨.
陪我的弟兄他說.
大雨啊.
去會不會阻到他們呢?.
但我自己想.
他們下著大雨都幫我們清潔.
又在垃圾房旁邊吃東西.

$^{721}$我為什麼不可以和他們送一點溫暖呢?.
於是我和弟兄趁著沒怎麼下雨.
雨勢小小的.
神聽了我們的禱告.
我們就拿上飲品.
我還記得其中一個清潔姐姐.
她這樣對我說.
她說.
都是叫美女.
美女,帥哥.
我們這些人被人看不起.
為什麼你們會來探我們?.
就由那天開始.
多年來兩年.
我和我的好朋友.
我們定期去探訪清潔工人.
有時我趁著不同的同工吃午飯.
我就去送一些物資.
和帶他們去信主.
但不是我一個人去做.
清潔姐姐跟我說.
有些教會特意買了飯盒.
每個月第一個星期三.
早上十一點多.
大家留意他們十一點多吃東西.
然後和他們分享福音.
各位弟兄姊妹.
心動不如行動.
上帝從來都是透過我們.
將恩慈帶給身邊的人.
求幫助我們.
我們一起禱告.
在這裡我有兩個禱告.
第一個禱告.
各位親愛的弟兄姊妹.
你現在的心靈是否真的很困苦?.
你需要人恩代理.
你現在告訴神.
祂知道.
有時有些事說不出來.

$^{761}$但祂知道.
我願神讓你在黃浸.
讓你在身邊人.
在家人,在鄰舍,在社區找到.
能夠和你同行.
恩代理的人.
第二件事,各位弟兄姊妹.
我們信主這麼多年.
我們說作為一個社區的明燈.
有沒有一些人.
原來一直放在我們心上.
我們要將你的愛帶給他們呢?.
如果有的話,你和上帝說吧.
你行動吧.
不要害羞.
豈不知,波亞斯這樣接待了路德和拿訛米.
竟然成為大衛家族.
整個救恩的一個大家族.
上帝同樣透過我們.
將祂的大愛帶給每一個看重的人.
從來沒有一個人是上帝忽視的.
只是我們人看不到.
求主幫助我們.
願主聽我們同心的禱告.
在奉耶穌基督的聖名祈禱.
阿們.
\newpage



\section{列王記下 6:1-7}
\label{sec:c6nbM3DecyY}
\textbf{2022年7月3日崇拜直播｜陳冠成牧師｜並肩勇闖系列:邀請上帝同行｜列王記下6\_1-7}
\newline
\newline
連結: \href{https://youtube.com/watch?v=c6nbM3DecyY}{\texttt{https://youtube.com/watch?v=c6nbM3DecyY}} ~~~~ 語音日期: 2022-07-04
\newline
\newline
\hyperref[sec:8sXdI3Lte_M]{\small{< < < PREV SERMON < < <}}
~
\hyperref[sec:index_chronic]{\small{[返順時目]}}
~
\hyperref[sec:sq3x1g07xfo]{\small{> > > NEXT SERMON > > >}}
\newline
\newline
列王記下 6:1-7
\newline
\begin{longtable}{cl}
\hline
\hline
章節 & 經文 (和合本修訂版)\\
\hline
6:1 & \begin{tabularx}{0.7\textwidth}{X} 先知的門徒對以利沙說：「看哪，我們在你面前居住的地方，那裡對我們太窄小了。 \end{tabularx} \\ \\ \relax
6:2 & \begin{tabularx}{0.7\textwidth}{X} 讓我們往約旦河去，各人從那裡取一根木料，在那裡為自己建造居住的地方。」他說：「你們去吧！」 \end{tabularx} \\ \\ \relax
6:3 & \begin{tabularx}{0.7\textwidth}{X} 有一人說：「請你與僕人同去。」他說：「我可以去。」 \end{tabularx} \\ \\ \relax
6:4 & \begin{tabularx}{0.7\textwidth}{X} 於是以利沙與他們同去。到了約旦河，他們砍伐樹木。 \end{tabularx} \\ \\ \relax
6:5 & \begin{tabularx}{0.7\textwidth}{X} 有一人砍樹的時候，斧子的頭掉在水裡，他就喊著說：「不好了！我主啊，斧子是借來的。」 \end{tabularx} \\ \\ \relax
6:6 & \begin{tabularx}{0.7\textwidth}{X} 神人說：「掉在哪裡了？」他把那地方指給以利沙看。以利沙砍了一塊木頭，拋在水裡，就使斧子的頭浮上來了。 \end{tabularx} \\ \\ \relax
6:7 & \begin{tabularx}{0.7\textwidth}{X} 以利沙說：「拿起來吧！」那人就伸手拿起來了。 \end{tabularx} \\ \\
[1ex]
\hline
\hline
\end{longtable}
$^{1}$接下來是經訓.
今天的經文是舊約聖經第462頁.
經文是列王記下六章一到七節.
請聽我讀.
先知門徒對以利沙說:看啊,我們和你所住的地方過於窄小,求你容我們往約旦河去,各人從那裡取一斤木料建造房屋居住.他說:你們去吧.有一人說:求你與僕人同去.回答說:我可以去.於是以利沙與他們同去,到了約旦河就砍伐樹木..
有一人砍樹的時候,斧頭掉在水裡.他就呼叫說:哀哉!我的時候,斧頭掉在水裡.他就呼叫說:哀哉!我主啊,這斧子是借的.神人問說:掉在哪裡了?他將那地方指給以利沙看..
以利沙砍了一斤木頭,拋在水裡,斧頭就飄上來了.以利沙說:拿起來吧!那人就伸手拿起來了..
各位姐妹平安!多多少少都有類似剛才經文所說的經歷,你不見了一些你覺得很重要的東西,然後你呼求上帝,上帝又很奇妙地幫你找回.有沒有試過這些經歷?多多少少有..
我記得很久以前也有跟大家分享過一個故事,就是以前我戴隱形眼鏡,還未拍拖的時候,可能是想帥一點.有一次戴著隱形眼鏡打籃球,比賽來的,打到差不多完場前一兩分鐘左右,突然跟別人爭球的時候,戳掉了一隻..
大家有沒有試過見隱形眼鏡去找的經驗?找到的機會大一點還是小一點?不容易,在家裡可能容易找,但你想像在籃球場中場,比賽還未完的時候,你又不能叫暫停,要撿回眼鏡才找,是不可以的..
你只能夠呼求上帝那一刻,就是上帝,這副眼鏡不是借回來的,是買了無奈的,怎麼辦呢?於是一兩分鐘過後,吹完球,突然站在中場,一看見有東西反光,見到有隻隱形眼鏡,有一隻,大概你也不會想,是不是我的呢?先撿起來,撿了回去,當然沒有可能立刻戴,之後洗乾淨,戴回去..
你又覺得這次倒數對,我只能夠當這隻眼鏡是我的,不是其他人的.也能體會到上帝那一刻,不是純粹幫自己找回一副眼鏡,找回隱形眼鏡,而是你更加體會到,當你在急難的時候,呼求上帝,上帝真的會回應,上帝真的會顯現他和你是同在,他是多麼真實..
所以那一次數年經驗,我想價值遠勝於那隻隱形眼鏡,比起那副隱形眼鏡貴重得多.同樣,今天我們看這段經文,其實焦點表面上好像是,善至以利沙,透過神蹟,幫他的僕人找回那副不見了的斧頭,表面上好像是這樣..
但當我們看深一點的時候,你會發現,其實這個神蹟是要顯明,上帝,或者說我們的主,是很願意和他的子民同行,並且很樂意介入他們的生活當中..
當我們看這段經文的時候,你會留意到,他有一個背景,居住的環境不足,於是一班先知門徒就打算自發前往約旦河,來砍伐樹木,建造房屋..
這個自發行動其實是很好的,不過你看到他們都不敢擅作主張,他們首先尊重他們的主人以利沙,問准,請求他批准..
結果你看到,以利沙很爽快地答應他們,你們去吧.而在這個時候,你會看到經文突顯,有一個門徒是與別不同的,他請求以利沙和他們一起去..
以利沙同樣很爽快地答應,於是一班人就好好的前往約旦河,砍伐樹木,建造房屋..
在這裡你會發現,雖然短短幾個經文,在同一件事裡面有兩班門徒,有一班是怎樣的呢?.
絕大部分的門徒只是向以利沙說了自己的計劃和想法,然後期望以利沙准許他們自行前往執行,這是大多數的門徒的想法..
問准主人,沒有想邀請主人和他們一起去..
為什麼會這樣呢?為什麼不需要邀請主人和他們一起去呢?.
經文裡面沒有很清楚的交代,不過我們可以試試跟他們回答..
有什麼可能呢?可能他們忘記了,太匆忙了,忘記了邀請以利沙一起去..
又或者可能都知道以利沙是顯赫的大先知,事務一定是很繁忙的,是吧?.
沒空陪我們,所以還是不要搞擾他了..
又或者其實我們都相信主人應該不是斬柴,不是做巢夫出身的,也不是做建築的..
所以對砍伐樹木建屋這些工作,其實應該叫他去也幫不了什麼忙..
又或者可能更加簡單直接就是,其實斬樹這麼小事,自己都能搞定..
既然自己應付到,就不需要勞煩主人和我們一起去了..
我們看到有可能很多的原因,但是無論如何,主人以利沙沒有硬要跟著去..
既然你們沒有邀請,那就隨便他吧..
他讓他們自己去完成工作,這是大多數門徒和主人的互動..
不過你會看到有一位門徒是經文裡面凸顯的,.
就是他在這件事上,你會發覺他是不是屬於少數?.
他是的,在這群的群體裡面,他是異類..
因為只有他期望不單單主人應運他們的計劃,.
他還會希望主人和他們一起去..
為什麼他有這個需要呢?.
很可能是他沒有斬柴伐木的經驗,又或者他自己知道體力一般般,.

$^{41}$又或者他本身都是鵪鶉底,很害怕的工作,擔心會有危險,會出意外..
所以覺得找主人一起去,自己會安心一點..
覺得這樣凡事會順利穩妥一點..
同樣無論如何,你看到主人以利沙沒有拒絕這個僕人的請求,.
她很樂意和他們一起去,接受這個邀請..
弟兄姊妹,在這裡我們看到兩種僕人和主人的互動方式,.
哪一種較為近似我們和上帝的關係呢?.
其實都頗值得去想..
不知道大家以前信主的時候有沒有這樣的體會,.
可能你事無大小,你都會向上帝說,.
甚至你有本靈修筆記,每天都記下你和上帝的互動,.
什麼你都很想和上帝分享,想他知道,.
你很想上帝帶領你走前面每一步,.
不過可能回教會久了,一年兩年之後,.
我們對教會的聚會越來越熟悉,知道如何運作,.
我們和弟兄姊妹又越來越熟絡,.
事奉的經驗亦累積得越來越多,自信心大了,.
但和上帝的關係會否因此變得世故了呢?.
會否漸漸覺得,這麼小的事自己就能處理,.
不用動都去麻煩上帝,不用動也找他..
確實是的,不難理解,.
可能我們今天都用這個模式和上帝互動,.
譬如你拿著很多重物,其實你很需要人幫忙,.
剛好有弟兄姊妹經過,看到,.
「需要幫忙嗎?我們一起拿吧.」.
於是你很快就會衡量自己能否應付,能否處理,.
能處理的,你就會勸對方,.
「不用了,我能處理,多謝你.」.
可能是出於禮貌,或為他人設想,.
不想小小事都開口,麻煩別人,對不對?.
確實,這是一個好的做人態度,好的處事方式,.
不過,我們和上帝的關係,會否經常都這麼客客氣氣呢?.
這又是另一個我們需要去思考的,.
特別是這裡,經文很想我們看到那一個門徒,.
其實他作出這個邀請,有沒有掙扎呢?有沒有壓力呢?.
你不妨想想,其他身邊所有人都沒有這個情況需要,.
其他人都沒有邀請主人同行,.
為何唯獨我有這個需要?.
是否我有什麼問題?.
我比較驚青有亂?我不夠熟齡?.

$^{81}$還是我比較差勁?.
又或者,我這樣問的時候,.
主人會否嫌我煩,都不想理我呢?.
弟兄姊妹,經文裡很清楚要我們想,.
其實真的不需要和上帝這麼見外,.
不用怕麻煩他,.
上帝其實很約立我們隨時來到他的面前,.
將你的需要,哪怕只是微不足道的一些事情,.
你告訴他,事實上,在你人生,.
在你每一日,無論大事或小事,.
上帝不單單喜悅我們來尋求他的心意,.
好像大多數的門徒一樣,.
上帝亦更期望我們邀請他同行,.
就好像那一個門徒一樣,.
邀請他來參與,.
因為聖經裡面講得很清楚,.
如果我們離開了上帝,.
其實我們什麼都不能做,.
所以弟兄姊妹,其實在我們的生活裡,.
我們其實很需要邀請上帝來參與和介入,.
我們會否有意識地邀請他,.
進到我們生活裡面面對的每一個場景,.
而不是單靠自己的努力來解決或作出判斷,.
譬如說,大家今天出門前,.
有沒有預先向上帝禱告交托?.
沒有,.
有沒有感恩我們在風雨下都可以回來,.
其實不是必然的,.
其實很值得去和上帝說,.
哪怕你只是微小的事,.
原來不是必然,.
我們有沒有將今天的崇拜,.
我們所面對,所接觸的每一件事情,.
每一個弟兄姊妹,.
都預先交托給上帝,.
讓上帝帶領我們行事為人?.
又或者,經歷完崇拜之後,.
下午你吃飯前,.
除了你的謝飯祈禱之外,.
會否為你今天所領受的,.

$^{121}$或者你今天接觸到弟兄姊妹的一些需要,.
亦都去禱告交托,.
願意上帝使用你?.
又或者可能到今晚,.
睡覺前,.
你知道你明天將會面對什麼人,.
和什麼事情,.
你有什麼需要,.
你有什麼擔憂,.
你可不可以,.
亦都在你睡覺前,.
祈禱,告訴上帝?.
弟兄姊妹,確實當我們持之以恆的話,.
其實你一定會和上帝的距離拉得越來越近,.
因為當你這樣做的時候,.
你發覺其實,原來上帝就在你身邊,.
你會活得比從前更有平安,.
更加有力量去面對每一天的挑戰,.
因為你知道上帝真的會看管你,.
你相信上帝會保守你,.
而這是經文第一個讓我們思考的,.
再看下去,第五節,.
這裡說那個僕人,.
他在砍樹的時候怎樣?.
斧頭掉下水,.
我們覺得很奇怪,.
為什麼這麼差勁?.
拿著斧頭一砍,整個斧頭掉下水?.
其實他原文的意思,.
不是整把斧頭掉下水,.
是在說斧頭的頭,.
即是金屬的部分,.
即是鐵的部分,.
你不難想像古代是怎樣,.
斧頭分兩部分,.
上半部分就是金屬用來砍樹很利器的部分,.
下半部分就和木頭綁在一起,.
或者怎樣去結合在一起,.
所以很大機會,.
那個僕人是因為斧頭鬆了,.

$^{161}$又或者他沒有檢查清楚,.
所以一砍下去的時候,.
斧頭的金屬部分就鬆脫,.
然後就掉下水,.
不過無論怎樣也好,.
因為這個斧頭是借回來的,.
不見了自然會有負那位物主,.
甚至乎可能會破壞雙方的關係,.
所以你看到這個僕人很慌張,.
慌張到他立即向以利沙,.
他的主人來呼求,.
慘了主人,這把斧頭是借回來的,.
怎麼辦呢?.
以利沙也似乎體會到,.
理解到這位門徒心裡的擔憂和焦慮,.
不單止沒有責備他,.
並且第一時間出手相救,.
問這個僕人斧頭掉到哪裡?.
僕人就將斧頭掉到水的位置,.
指給以利沙看,.
奇怪的是,.
以利沙沒有打算親自下水,.
和他一起去找斧頭,.
而是要出動一個神跡,.
來幫這個僕人找回斧頭,.
所以有兩個問題很值得我們去想,.
其實找東西而已,.
尋找生物而已,.
為何要用到神跡,.
好像這麼大一坪賣,.
很奇怪吧?.
又或者,.
為何要大費周章,.
又要用一斤小樹枝,.
掉到斧頭掉下來的位置,.
為何不直接,.
反正都是行神跡,.
直接吩咐斧頭在水上浮起就行了,.
為何要這麼麻煩,.
做這麼多細微的動作,.

$^{201}$這兩個問題很值得我們去想,.
原來鐵,.
在當時是一個貴重的金屬,.
我們相信都理解,.
不同今日,鐵是很廉價,.
在古代,.
當時是一個貴重金屬,.
遺失了借回來的斧頭,.
不單止有機會破壞,.
這個僕人和鄰舍的關係,.
更加會招致甚麼危機?.
找不回的時候,.
就會有債務危機,.
你知道做仙芝,.
或者做仙芝的門徒,.
是很貧窮的,.
因為他們要去興建這個不見了的斧頭,.
所以如果真的找不回的時候,.
他很大機會要成為奴隸一段時間,.
來還清這個債務,.
所以我們就會開始明白,.
為何仙芝爾利沙要動用到神跡,.
來介入這次的事件,.
首先第六節說,.
斧頭飄上來,.
原文的意思是,.
拋進去裡面的樹枝,.
令到斧頭浮起來,.
不是斧頭自己浮起來,.
雖然我們看不到,.
樹枝和斧頭,.
到底在水裡面,.
實際是怎樣互動,.
但從這段經文,.
你會看到,.
可以想像到,.
水裡面的畫面是怎樣,.
就是樹枝主動潛到水裡面,.
然後找到水底下沉,.
將它浮起來,.

$^{241}$所以這個神跡,.
不單單是說斧頭浮起來,.
更加是樹枝下沉,.
基於物理定論,.
斧頭重,.
它無法自己救自己,.
浮起來,.
所以它必須要借助外來力量,.
來將自己救起來,.
所以它必須要借助外來力量,.
來將自己救起來,.
所以它必須要借助外來力量,.
來將自己救起來,.
誰是外來力量?.
就是同樣違反物理定論,.
竟然可以.
向下沉那條樹枝,.
用新約的角度,.
你會更加明白,.
就好像耶穌基督,.
透過十字架那一根木頭,.
將無法自救的你和我,.
救起來一樣,.
所以同樣,.
經文這裡給我們看到,.
那個身陷債務危機的僕人,.
其實是無法自救的,.
他需要呼求他的主人,.
救他脫險,.
所以我們看這段經文的時候,.
就不是單單表面看到,.
真是很感恩上帝,.
透過神蹟幫我找回,.
不見了的東西,.
而是要顯明,.
我們的本質,.
其實和這個僕人,.
和那個遺失的斧頭,.
都是一樣,.
你無法自己救自己,.

$^{281}$有一位慈愛的主,.
他願意隨時,.
隨聽我們向他發出的呼求,.
並且,.
他會在適當的時候出手,.
甚至用你意想不到的方式,.
將我們從日常,.
大小的困局裡面,.
很奇妙地拯救出來,.
這就是經文,.
給我們看到那幅異像,.
話說其實,.
我平時每天,.
都是一家人,.
一早就出門,.
我就是負責鎖門的,.
我很有交代,.
很有首尾,.
每次都使命必達,.
一定鎖好門,.
包括星期日,.
星期日我通常會早點出門,.
因為要回來預備,.
所以我很有首尾,.
照樣鎖好門,.
不過我太太和兒子,.
還在裡面,.
有時我會,.
我會反鎖好門,.
令他們不能開門,.
太有責任,.
試過一兩次,.
在崇拜的日子,.
當我醒來的時候,.
原來我已經身處旺角,.
我家住在大埔,.
你知道沒有辦法去拯救他們,.
你很想去拯救他們,.
不過當你發覺的時候,.
你發覺不能回去,.

$^{321}$於是你唯有向主呼求,.
於是老婆又打來,.
你鎖好門了,.
不能回來,.
你便知道有多惡極,.
作為牧師竟然將自己的師母,.
和兒子,.
困在家裡,.
令他們不能崇拜,.
不能參與事奉,.
多麼惡極,.
當你呼求上帝的時候,.
上帝有感恩,.
有幾次上帝找了我岳父,.
因為他住在附近,.
他有我們家的鑰匙,.
在門外拯救我們,.
幸好他沒有出門喝茶,.
兩次都是這樣,.
我呼求的時候,.
上帝透過我岳父,.
將我太太和兒子,.
從家裡拯救出來,.
事實上頂指不會是真的,.
我們遇到困難,.
第一個自然反應,.
我們會用自己的方法,.
用自己的力量去解決,.
不過經文給我們另一個選項,.
就是我們會用自己的方法,.
用自己的力量去解決,.
我們會用自己的力量去解決,.
有些頂梓梅會說,.
我會穩陣些,.
我會先檢查一下斧頭,.
確保它不會再鬆脫,.
我才會繼續用它來砍伐樹木,.
又或者有頂梓梅說,.
其實我還未驚完,.
我很擔心會再次遺失,.

$^{361}$讓我休息一下,.
讓我安心一點,.
讓我考慮一下還繼續不繼續,.
有些人會有個抗議的想法,.
不如和隔壁的交換,.
他覺得自己拿斧頭也不太好,.
斧頭又鬆脫,.
不如找個厲害的交換來用,.
讓自己安心一點,.
不同的想法,.
頂梓梅,如果你是那個門徒,.
你會怎樣做?.
這麼難得上帝先幫我找回斧頭,.
多大都繼續用它來完成上帝的工作,.
還是這麼難得我先找回斧頭,.
多大都好好保管,.
不要再讓它不見,.
不要冒險劈樹,.
免得我再一次受驚懼怕..
頂梓梅確實不容易,.
我們有時很兩難,.
到底是冒險進取,.
還是保守安全些?.
曾經聽聞瑞典人有個傳統,.
是保守的,不是冒險的,.
他們是這樣的民族,.
話說很久之前,.
有個來自瑞典人的家庭,.
當中有個小女孩,.
她很喜歡收集很漂亮的瓷器,.
但由於她的家境不是富裕,.
而當時瓷器是很貴重的物品,.
所以她的父母.
只能夠在她重要的日子,.
例如生日,畢業禮或浸禮,.
只能夠送她一件瓷器,.
而當這個小女孩收到瓷器時,.
她擔心會刮花或打爛,.
所以她不捨得用,.
她會很小心地.

$^{401}$將每一件瓷器包得漂亮,.
然後放入盒子裡,.
再將這些瓷器放在閣樓裡,.
她很希望等到將來.
有些很重要特別的日子,.
她就拿出來用了,.
當然最後我們知道.
所謂特別重要的日子從來沒有出現過,.
直到這個女孩.
她做了外婆,.
甚至乎之後離開了這個世界,.
她一生很珍而重之的禮物,.
一直都沒有打開過來用,.
後來到她的女兒.
發現這批瓷器時,.
她決定.
讓他們重見天日,.
將他們物盡其用,.
雖然可能在過程中.
會刮花.
會撞破或弄破,.
但卻能將這些瓷器.
本身的價值和意義.
盡情地發揮出來..
弟兄姊妹,.
你又如何看待.
上帝給你的禮物?.
你會擔心很有風險,.
怕受傷,.
所以一直將上帝賜給你的禮物.
收起來,原封不動不敢用?.
還是你覺得.
上帝賜給你的禮物.
既然如此珍貴,.
就不要收起來,.
應該打開.
與人一起享用?.
弟兄姊妹,很值得我們去思考,.
什麼才是.
沒有浪費上帝.

$^{441}$所託付給你,.
所賜給你的禮物和恩賜?.
確實沒有人喜歡去冒險,.
沒有人喜歡遇到風險,.
每個人都想一帆風順,平平安安,.
但回想你的人生,.
其實很多重要寶貴的成長功課,.
基本上與冒險有關,.
你試想想,.
你自己在小朋友的時候,.
或你做父母,.
你照顧小朋友的時候,.
他們學行的時候,.
有沒有風險?.
一定有,.
小朋友學踩單車的時候,.
有沒有風險?.
一定有,.
如果你身邊的人,.
沒有適當地放手,.
讓你去冒險的時候,.
有很多重要的技能,.
有很多寶貴的經驗,.
其實我們原來這一世都學不到,.
我們這一世都掌握不到,.
唯有別人肯放手,.
你肯去冒險,.
你才會掌握到,.
同樣有些事情,.
其實上帝是定義放手,.
讓你和我去完成,.
從而我們可以經歷到,.
什麼叫信心的成長,.
上帝雖然放手,.
但祂不會走開,.
祂會在旁邊,.
有需要的時候,.
祂一定會像這句經文,.
出手幫你,.
讓你可以繼續前進,.

$^{481}$完成祂交託給你的事情,.
就好像我想起一個圖畫,.
就是你學駕駛,.
當你出路面的時候,.
你不是很害怕,.
因為你的師傅會幫你,.
駕駛師傅,學車師傅,.
當你斷線的時候,.
其實你很害怕,.
人生第一次斷線,.
你很害怕,出不出路,.
你放心,.
你的師傅會幫你,.
有需要他會幫你扭tie ,.
有需要他會救你,.
不會讓你撞車,.
不會讓你撞,.
同樣上帝會幫我們,.
他會在風險裡面,.
提醒我們真正要捉實的,.
不是失而復得的苦頭,.
而是要捉實,.
上帝願意和我們同行,.
上帝願意垂聽我們的呼求,.
上帝願意隨時私恩拯救,.
這個的應許和恩典,.
你捉實這件事的時候,.
再沒有任何的事情,.
能夠難阻你,.
為上帝工作,.
為上帝冒險,.
真正的弟兄姊妹,.
我們能為上帝做多少事情,.
能在信仰的路,.
事奉的路走多遠,.
其實關鍵不是我們自己有多少實力,.
而是在於上帝加多少恩典給我們,.
如果我們能夠放膽善用上帝所加給我們的,.
我們的生命會有更多的色彩,.
更有更多的可塑性,.

$^{521}$所以讓我們一起去祈禱,.
求神幫助我們,.
生命裡面有些突破,.
我們不再停留在純粹量力而為的信仰階段,.
我們只是祈求上帝賜給我們和我們能力相稱的任務,.
而是學習我們量因而作,.
我們放膽向神祈求,.
與任務相稱的恩典和能力,.
讓我們同心祈禱,.
上帝我們多謝你拯救我們,.
事實上不是我們找到你,.
是你親自來尋找呼召我們,.
讓我們可以成為你寶貴的兒女,.
成為君尊的祭祀,.
上帝盼望我們捉緊你這份救恩,.
不單在天堂永恆天國裡面,.
我們在日常生活裡面,.
我們也在捉緊你的拯救,.
因為在人生大大小小的事情裡面,.
其實我們離開了你,.
我們不能做什麼,.
我們沒有辦法自救,.
唯有我們向上帝你呼求,.
求主你幫助我們,.
讓我們真的能夠時刻與你連結,.
無論是大小的事情,.
我們都願意與你分享,.
邀請你同行,.
在今晚裡面,.
我們首先不是自己去自救死去,.
而是懂得交託給上帝你,.
願意上帝你幫助我們,.
讓我們向你求的是更大的恩典,.
更大的能力,.
以致我們可以為上帝你做超乎想像的事情,.
願意上帝你恩待我們,.
按你的心意來幫助我們,.
求你聽我們的禱告,.
基督耶穌您的特性的命志祈求.
Amen.

\newpage



\section{路加福音 24:13-35}
\label{sec:sq3x1g07xfo}
\textbf{2022年7月10日崇拜直播｜羅潔盈博士｜我們的心豈不火熱?｜路加福音24\_13-35}
\newline
\newline
連結: \href{https://youtube.com/watch?v=sq3x1g07xfo}{\texttt{https://youtube.com/watch?v=sq3x1g07xfo}} ~~~~ 語音日期: 2022-07-11
\newline
\newline
\hyperref[sec:c6nbM3DecyY]{\small{< < < PREV SERMON < < <}}
~
\hyperref[sec:index_chronic]{\small{[返順時目]}}
~
\hyperref[sec:SQA6e1_jseU]{\small{> > > NEXT SERMON > > >}}
\newline
\newline
路加福音 24:13-35
\newline
\begin{longtable}{cl}
\hline
\hline
章節 & 經文 (和合本修訂版)\\
\hline
24:13 & \begin{tabularx}{0.7\textwidth}{X} 同一天，門徒中有兩個人往一個村子去；這村子名叫以馬忤斯，離耶路撒冷約有二十五里。 \end{tabularx} \\ \\ \relax
24:14 & \begin{tabularx}{0.7\textwidth}{X} 他們彼此談論所發生的這一切事。 \end{tabularx} \\ \\ \relax
24:15 & \begin{tabularx}{0.7\textwidth}{X} 正交談議論的時候，耶穌親自走近他們，和他們同行， \end{tabularx} \\ \\ \relax
24:16 & \begin{tabularx}{0.7\textwidth}{X} 可是他們的眼睛模糊了，沒認出他。 \end{tabularx} \\ \\ \relax
24:17 & \begin{tabularx}{0.7\textwidth}{X} 耶穌對他們說：「你們一邊走一邊談，彼此談論的是甚麼事呢？」他們就站住，臉上帶著愁容。 \end{tabularx} \\ \\ \relax
24:18 & \begin{tabularx}{0.7\textwidth}{X} 兩人中有一個名叫革流巴的回答：「你是在耶路撒冷的旅客中，惟一還不知道這幾天在那裡發生了甚麼事的人嗎？」 \end{tabularx} \\ \\ \relax
24:19 & \begin{tabularx}{0.7\textwidth}{X} 耶穌對他們說：「甚麼事呢？」他們對他說：「就是拿撒勒人耶穌的事。他是個先知，在神和眾百姓面前，說話行事都大有能力。 \end{tabularx} \\ \\ \relax
24:20 & \begin{tabularx}{0.7\textwidth}{X} 祭司長們和我們的官長竟把他解去，定了死罪，釘在十字架上。 \end{tabularx} \\ \\ \relax
24:21 & \begin{tabularx}{0.7\textwidth}{X} 但我們素來所盼望要救贖以色列民的就是他。不但如此，這些事發生到現在已經三天了。 \end{tabularx} \\ \\ \relax
24:22 & \begin{tabularx}{0.7\textwidth}{X} 還有，我們中間的幾個婦女使我們驚奇：她們清早去了墳墓， \end{tabularx} \\ \\ \relax
24:23 & \begin{tabularx}{0.7\textwidth}{X} 不見他的身體，就回來告訴我們，說她們看見了天使顯現，說他活了。 \end{tabularx} \\ \\ \relax
24:24 & \begin{tabularx}{0.7\textwidth}{X} 又有我們的幾個人往墳墓那裡去，所發現的正如婦女們所說的，只是沒有看見他。」 \end{tabularx} \\ \\ \relax
24:25 & \begin{tabularx}{0.7\textwidth}{X} 耶穌對他們說：「無知的人哪，先知所說的一切話，你們的心信得太遲鈍了。 \end{tabularx} \\ \\ \relax
24:26 & \begin{tabularx}{0.7\textwidth}{X} 基督不是必須受這些苦難，然後進入他的榮耀嗎？」 \end{tabularx} \\ \\ \relax
24:27 & \begin{tabularx}{0.7\textwidth}{X} 於是，他從摩西和眾先知起，凡經上所指著自己的話都給他們作了解釋。 \end{tabularx} \\ \\ \relax
24:28 & \begin{tabularx}{0.7\textwidth}{X} 他們走近所要去的村子，耶穌好像還要往前走， \end{tabularx} \\ \\ \relax
24:29 & \begin{tabularx}{0.7\textwidth}{X} 他們卻強留他說：「時候晚了，天快黑了，請你同我們住下吧。」耶穌就進去，要同他們住下。 \end{tabularx} \\ \\ \relax
24:30 & \begin{tabularx}{0.7\textwidth}{X} 坐下來和他們用餐的時候，耶穌拿起餅來，祝福了，擘開，遞給他們。 \end{tabularx} \\ \\ \relax
24:31 & \begin{tabularx}{0.7\textwidth}{X} 他們的眼睛開了，這才認出他來。耶穌卻從他們眼前消失了。 \end{tabularx} \\ \\ \relax
24:32 & \begin{tabularx}{0.7\textwidth}{X} 他們彼此說：「在路上他和我們說話，給我們講解聖經的時候，我們的心在我們裡面豈不是火熱的嗎？」 \end{tabularx} \\ \\ \relax
24:33 & \begin{tabularx}{0.7\textwidth}{X} 於是他們立刻起身，回耶路撒冷去，看見十一個使徒和與他們正在一起的人聚集在一處， \end{tabularx} \\ \\ \relax
24:34 & \begin{tabularx}{0.7\textwidth}{X} 說：「主果然復活了，已經顯現給西門看了。」 \end{tabularx} \\ \\ \relax
24:35 & \begin{tabularx}{0.7\textwidth}{X} 於是，兩個人把路上所遇到，和耶穌擘餅的時候怎麼被他們認出來的事，都述說了一遍。 \end{tabularx} \\ \\
[1ex]
\hline
\hline
\end{longtable}
$^{1}$接著我們有田慶仙弟兄為我們讀經.
在《奴家福音》24章13至15節.
接著我們今天很開心.
有見到神學院實用神學助理教授.
羅傑盈博士為我們正道.
主席.
請聽我讀出聖經新約《奴家福音》.
24章13至15節.
正當那日門徒中有兩個人往一個村子去.
這村子名叫以馬五師,來耶路撒冷約有二十五里.
他們彼此談論所遇見的這一切事.
正談論相問的時候.
耶穌親自就近他們和他們同行.
只是他們的眼睛迷糊了,不認識他.
耶穌對他們說:你們走路彼此談論的是什麼事呢?.
他們就站著,面上戴著袖戎.
兩人中有一個名叫格劉巴的回答說:.
你在耶路撒冷作客,還不知道這幾天在哪裡所出的事嗎?.
耶穌說:什麼事呢?.
他們說:就是那撒勒人耶穌的事.
他試過先知,在神和眾百姓面前,說話行事都有大能.
祭司長和我們的官府竟把他解去,定了死罪.
釘在十字架上.
但我們數來所盼望要贖以色列民的就是他.
不但如此,而且這是成就,現在已經三天了.
而且我們中間有幾個婦女使我們驚奇.
她們清早到墳墓那裡,不見她的身體.
就會來告訴我們,說看見了天使顯現,說她活了.
又有我們的幾個人往墳墓那裡去.
預見的正如婦女們所說的,只是沒有看見她.
耶穌對她們說:無知的人啊,先知所說的一切話.
你們的心信得太遲鈍了.
基督這樣受害,又進入她的榮耀,豈不是應當的嗎?.
於是從摩西和眾仙子起,凡經上所指著自己的話.
都給她們講解明白了.
將近她們所去的村子,耶穌號召還要往前行.
她們卻留她說:時候晚了,日頭已經平息了.
請你和我們住下吧.
耶穌就進去,要和她們住下.
到了坐直的時候,耶穌拿起餅來,祝謝了.

$^{41}$張開眼睛,提及她們,她們的眼睛明亮了.
這才認出她來,忽然耶穌不見了.
她們彼此說:在路上她和我們說話.
給我們講解聖經的時候,我們的心豈不是火熱的嗎?.
她們就立時起身,回耶路撒冷去.
正遇見十一個使徒和她們的同人聚集在一處.
說:主果然復活,已經現給西門看了.
兩個人就把路上所遇見和抹餅的時候.
怎麼被她們認出來的事都述說一遍.
願主祝福大家.
(主持禮物).
各位網上的弟兄姊妹,主來平安.
主果然復活了.
在經文34節裡,給我們一個崇拜的傳統中.
一個很重要的經文.
為什麼我們要星期日,或主日,或禮拜日去聚集呢?.
就是因為復活的基督帶來新的生命.
所以我們也和旁邊的弟兄姊妹彼此問安.
願復活主的平安與你同在.
(主持禮物).
主果真復活了,Christ is risen indeed, Hallelujah.
多謝剛才的歌頌.
Hallelujah是什麼意思?.
Hallelujah就是讚美,你們要讚美.
Hallelujah就是你們要讚美耶和華.
邀請大家在座位中思考.
我們現在每人都戴著口罩.
雖然師伯很努力去歌頌.
依然聽到他們的歌聲,戲的運用都很好.
我們會否想起.
我們以前能夠更加無拘無束的.
不用戴著口罩.
或者邀請大家思考一個令你覺得火熱.
或心裡感到異常溫暖的時刻.
因為這段經文提醒我們.
我們的心在我們裡面豈不是火熱的嗎?.
在二十五年前,我是一個中六的學生.
當時一個暑假都去三個營會.
有很多不同的,例如聖樂營.
或者青少年的營.

$^{81}$也有教會的下令營.
很多弟兄姊妹請假去聽神的道.
很努力地寫下筆記.
那時候真的覺得心裡很火熱.
一聽到上帝的話語.
雖然對一個中學生來說是難的.
因為經常坐著聽.
但當一群人一起去聽和看神的話語的時候.
神就給我們一個很不同的視野.
讓我們重新看到我們心裡的夢想.
DSE還有十天就放榜.
不知道你的子女或孫子,孫女.
會否都是這一屆都很慘的DSE同學.
因為他們這幾年基本上沒有什麼營會.
也經常在網上做教學.
我們思考二碼五師之路.
可能你未必是DSE的考生.
現在這麼迷茫等放榜.
但是從這兩個門徒.
我們再思考.
我們現在在2022年7月裡.
會否走在這條路上.
也有時候像這兩個門徒一樣.
走著走著停著.
在不知道前方的方向是怎樣的路上.
「戚」「戚」是行字的「察」字.
代表左腳和右腳.
在第17節看到這兩個門徒站著面帶受容.
復活的基督跟他們說.
「你們在談什麼?」.
從他們的身體語言已經感覺到.
他們很悲傷.
他們被眼前的失望所遮蓋.
他們看不到耶穌.
如果復活的基督穿牆過壁.
在我們之間.
他會對我們說什麼呢?.
會否像《路加福音》24章的故事一樣.
耶穌第一句說的是.
「無知的人啊!Foolish people!」.

$^{121}$「你的信心信得太遲鈍了!」.
讓我們用一兩分鐘.
回想一個令我們感到火熱的時刻.
讓我們思考.
我們能否看到耶穌在我們生命當中呢?.
如果看不到的時候.
是什麼令我們火熱的心變成灰心呢?.
有什麼大事令我們看不到耶穌在我們當中?.
祂和我們同行,一直同行.
但我們不斷說自己的事.
不斷在我們的悲傷和哀悵當中.
這個圖是「金計」.
是日本人用黃金修補的工藝.
當我們覺得很多事物都很碎片化.
或不太清楚人生下一步.
或家人流散.
或疫情不知何時完結.
我們是否相信復活的基督能帶來轉變的生命呢?.
最近我們的氣溫都有明顯的不同.
近105年最冷的勞動節.
當日平均氣溫只有14.5度.
但4月又很熱.
所以在思考這個問題的時候.
令我想起有些同學身水身汗.
我的神學院在長洲.
走到禮堂的時候.
一量氣溫,嘀嘀嘀嘀.
37.5度,快38度了.
同學,你休息一分鐘.
一分鐘後又回到37度.
我們的身體會自動調節.
不會太低,也不會太高.
但在靈性當中.
上主經常責怪我們會否不冷不熱.
或我們如何保持熾熱的心.
亦在靈裡面不枯乾.
在當中看到耶穌的同行.
在大家的周刊裡.
也有大綱和經文.
當我在教不同的聖樂或崇拜的時候.

$^{161}$我們都希望崇拜有更新.
最近亦有不同的宣教營.
在神學院幫一班DSE的同學.
思考前路的時候.
其實兩個的核心都是一樣.
希望我們可以在生命當中.
遇見基督,生命得到轉化.
這個生命的轉化是180度.
本身這兩個門徒很不開心.
哀悼他們失落了的信心.
失落了本身的盼望.
但耶穌的同行讓他們再一次發現.
基督真的復活了.
在福音書裡總共有35個神蹟.
但不要忘記還有兩個神蹟.
第一是耶穌基督被聖父從死裡復活.
第二是耶穌向不同的門徒顯現.
在當中讓我們看到.
原來復活的基督沒有死.
而是他的生命要帶來我們生命當中的影響.
《路加福音》二十四章十三節開始.
同一天門徒中有兩個人要往一個村子去.
這村子名叫以馬五斯.
來耶路撒冷約有二十五里.
這二十五里相對現今大約三,四個小時.
約十一公里的路程.
他們彼此談論所發生的一切事.
如果我們用一個想像.
這兩個門徒走的時候很沉重.
走得很慢.
但他們終於開口說話.
經過了三天.
耶穌死了三天之後.
面對這個哀傷.
他跟隔壁的人說.
究竟發生了這些事,為何會這樣呢?.
在失落當中,他們陷於憂愁.
在失望當中,他們彼此再回顧.
他們其實如果用心理學的角度.
是經歷了一個很重大的創傷.

$^{201}$PTSD.
看著耶穌在十字架上被羞辱.
而且血一滴一滴流乾了.
面對這麼大的創傷.
而且已經三天了.
在猶太人的傳統裡.
三天是甚麼意思呢?.
他們認為人的靈魂和身體.
在三天之後,如果死了三天之後.
靈魂就會離開了一個腐化的身體.
因為當時沒有這麼好的冷藏.
或保存肢體的技術.
在當中他們完全絕望.
因為覺得耶穌真的死了.
面對這個哀傷.
很重要的是他們有同伴.
可以互相說出來.
而他們面對的失落.
不單單是失去信心.
而且他們還眼睛迷糊.
看不到耶穌,失魂落魄.
對於這兩個門徒來說.
人生失去的可能彼得到的多.
因為他們認為.
這位拿撒日蘭耶穌是一個先至.
他要幫助以色列民復國.
但為何他竟然被羞辱在十字架上呢?.
在他們的哀悼當中.
耶穌和他們同行.
解答了他們很多不明白.
他們的不明白有甚麼在當中呢?.
我們從他們的答案來看.
他們認不出耶穌.
因為他們心裡對耶穌已經有既定的想法.
他們認為耶穌應該在政治上.
讓他們帶來以色列復國.
或是權柄.
從猶太教的盼望當中.
他們認為耶穌是一個有能力的先知.
耶穌和他們同行時.

$^{241}$跟他們說:無知的人啊!.
先知所說的一切.
你們的心信得太遲鈍了.
於是基督就解釋.
為何基督要受這些苦難.
以及進入他的榮耀.
二十六節這一段經文.
其實是解釋他們一些.
部分對,但部分錯的信仰的疑團.
我們稍後再詳細地看.
這裡我們先看.
他從摩西和眾先知起.
凡經上所指著自己的話.
都給他們作了解釋.
在這裡我們看到.
為何我們要一齊聽上帝的話語.
而且在這個並肩勇闖的系列當中.
大家都聽到很多舊約.
不同的角色.
或是舊約給我們看到的一些亮光.
包括約瑟.
但和耶穌有什麼關連呢?.
我們也希望復活的主.
幫我們更加明白.
如何從舊約到新約當中.
可以看到耶穌的身份.
第三個段落是.
「騰出空間,邀請他者」.
這是整個故事的轉捩點.
他們快要到達要去的村子的時候.
耶穌好像還要往前走.
他們卻強擠耶穌說.
「時候晚了,天快黑了」.
「請你和我們一起住吧」.
香港人的家裡通常沒有空間.
讓其他人住下來.
但作為汪津.
不知道大家是否住在旺角.
我小時候是住在旺角.
我們能否也騰出空間.

$^{281}$邀請其他人.
在這條路上一起.
對於耶穌來說.
是一個陌生人.
但當他們看到他者的需要.
或是邀請這位陌生人.
進入他的生命當中的時候.
很不同的就是.
他們坐下來一起用餐.
耶穌賴在餅內.
把餅給他們.
在餅中.
他們眼睛明亮.
認得出耶穌.
耶穌卻在他們眼前消失了.
耶穌也消失了.
跟一開始一樣.
他們看不到耶穌.
但他們的心靈很不同.
感恩餅,眼目光明.
在路上.
耶穌跟我們講解聖經的時候.
我們的心在我們裡面.
豈不是火熱的嗎?.
當我們再看的時候.
為何我們要聚集去敬拜呢?.
原來我們需要彼此.
我們需要透過餅.
透過一切的敬拜.
成為蒙恩的渠道.
讓我們看到神的恩典是臨在的.
而且面對沒有改變的.
正如主席所說.
世間可能沒有平安.
但耶穌給予我們平安.
在我們心裡.
讓我們用感恩的眼睛.
去看待所有事情.
從此第五點.
就是他們勇敢的前行.

$^{321}$展開使命.
於是他們立刻起來.
回到耶路撒冷去.
看到十一個門徒.
和他們一起聚集一處的人.
他們見證主果然復活了.
已經顯現給西門看了.
於是兩個人在路上遇到的.
和耶穌抹餅的時候.
如何被他們認出來的事.
都述說了一遍.
所以我認為.
耶穌抹餅的時候.
被認出來.
也是一個很大的神蹟.
因為這個神蹟令到.
本來這兩個門徒.
有眼但不能見.
見不到耶穌,見不到盼望.
有腳,但不知去哪走好.
回到鄉下還是去哪裡好呢?.
耶穌和他們同行.
這一席話改變了他們.
不單是肉身上的意志.
也是心靈裡的意志.
重新火熱起來去事奉他.
所以從這五點裡.
我們再思考.
在我們現今2022年7月裡.
我們面對疫情好像死不斷氣.
又或者是面對教會當中.
可能有不同的青黃不接.
又或者是在很多人際關係上.
本身已經有的不同立場的撕裂當中.
我們怎樣可以重燃對主對人火熱的心呢?.
不知道大家有沒有去過聖地?.
聖地會不會有二馬五師這個地方?.
其實好像一個神話.
基本上有四個可能性.
但是來耶路撒冷.

$^{361}$聖經說的二十五里.
我們不是很能夠找到一個肯定的地點.
但是這個El Kibbebe在耶路撒冷.
西北偏西約13公里.
在伯特利西南約9公里.
距離《路加福音》所描述的有一點相似.
1943年因為考古學的發達.
再發掘的時候發現.
有第一世紀的時候古村的遺物在村裡.
而且還有一座是第三世紀所建的聖所.
和一個拜占庭的教堂.
經過歷史上529年毀於撒馬利亞人手.
在530年又被波斯人重圍.
改為回教寺.
十字軍的時候又建了一座羅馬式的聖堂.
到了19世紀聖方濟買了這個地.
在上面建了一間教堂和修道院.
名字叫做格劉巴堂.
格劉巴就是我們在經文當中.
找到耶穌復活之後.
二馬五思的兩個門徒之一.
格劉巴是甚麼人呢?.
有學者認為他是雅各的父親.
也有東方的學者或學派認為.
這兩個門徒其實就是格劉巴和他的妻子.
因為她強迫耶穌住下來.
所以從這個圖中可以看到.
她是一對夫婦.
和耶穌一起走二馬五思的路.
一起坐寂.
也有近代美國的幾幅藝術品.
可以看到一對夫婦在快餐式的中國餐館聊天.
那誰是耶穌呢?.
如果留意,她沒有畫到耶穌的腳.
就是中間那個和他們一起.
二馬五思可能是一個.
你不知道去哪裡的時候.
先去吃點東西.
然後在旁邊,原來耶穌在那裡.
和他們一起去傾談.

$^{401}$一起去講解.
原來這個信仰,信心是這樣理解的.
所以我們再看這段經文的時候.
甚麼意思是眼睛迷糊,內心遲鈍呢?.
聖經裡面說.
他們因為不明白舊約的聖經和經文.
都是指著耶穌去講的.
摩西五經創出利民生和先知書.
由以賽亞書到馬拉基書.
這些對我們來說都是比較陌生和硬的經文.
但透過耶穌在二馬五思的路上.
他開了這兩個門徒的眼睛.
讓他們本身被蒙蔽.
可能是被自己對於耶穌的理解.
又或者是在當中太哀,太悲傷.
所以看不到耶穌是誰.
但當一起擦餅的時候.
他們的眼睛就明亮了.
我們再看經文門徒甲留巴.
看不明白甚麼呢?.
他說耶穌是一個人.
那撒勒人耶穌是一個先知.
但其實在神學上我們認為.
施洗約翰是最後一個先知.
因為先知的身份,角色就是指明誰是微賽亞.
第三,他希望耶穌來救贖以色列民復國.
他認為中間有幾個婦女使我們驚奇.
意思是那些婦女好像很神奇.
但我真的看不到耶穌.
那些婦女說看到復活的基督.
但我看不到.
所以耶穌解答憂愁的門徒.
基督就是微賽亞.
他要救贖的不只是以色列民.
而是救贖萬民.
耶穌基督的復活要帶我們進入榮耀當中.
從哀悼到看見.
我們看到在崇拜當中.
我們感恩去擦餅的時候.
讓這兩位門徒重拾火熱的心.

$^{441}$這位在丹麥的畫家.
最近在香港也很常用他的畫.
但這幅給我們看到.
耶穌的釘行手.
他擦開這個餅.
在葡萄汁的杯子當中.
看到他的樣子的時候.
周圍的人有不同國籍的難民.
有受傷的人.
我們有武刀這樣騰出空間.
讓我們社區當中有需要的人.
一起來到主的聖桌當中.
去擦餅去感恩.
這個感恩獻上為祭.
是崇拜當中第三個下而上的向導.
因為神呼召聚集我們.
去聽他的話語.
也以感恩為祭去獻上.
再擦餅我們出去.
如果有理科的人.
這個圖是兩個光子撞擊的時候.
用一些很特別的方法.
放大很多倍的時候所發現的.
聖餐的向導不單是下而上.
也像這個圖這麼漂亮.
當我們說要作光.
作光明之子的時候.
光子即是一個光的元素.
除了發光也要發熱.
我們有沒有在當中.
彰顯基督在我們生命當中的熱能呢?.
在日本有一個金計.
就是用黃金修補的工藝.
在2011年311地震當中.
有一些藝術家收集不同的陶瓷碎片.
在災禍當中.
災難當中,苦難當中.
我們不明白.
也不知道什麼是復活.
什麼是復活的奧秘.

$^{481}$或者新的創造.
但是透過藝術品.
藤川真這個藝術家讓我們看到.
原來我們生命當中很多的破碎.
透過基督.
透過用很貴重的黃金.
將我們黏在一起的時候.
這件藝術品更加漂亮.
更加有價值,更加有獨特性.
而且在當中讓我們看到.
即使生命當中有不同的傷痕.
但主是我們的治癒者.
祂能夠讓我們不同破碎的撕裂的關係.
可以得到修復.
所以怎樣去保持熾熱的心呢?.
在宣導會的四重福音當中.
其中有一個「游平」.
象徵聖靈傳遞的作為.
包括醫治我們身心靈所有的創傷.
但願我們相信這位復活的基督.
不單單拯救我們.
也是醫治我們.
使我們成聖.
也是必定會再來.
而且在當中我們能夠喜樂的.
再一同的與先聖先賢一起歌頌.
在疫情期間.
我也有家人患病過世.
在一些突發的情況下.
聽不到家人最後一句遺言就過世.
是很悲傷的一件事.
但當我們相信耶穌是治癒者的時候.
但願我們靠著喜樂的「游」.
醫治的「游」.
在哀悼的當中.
願聖靈以喜樂的「游」.
去告我們.
使我們充滿喜樂.
有活潑的生命.
我們亦彼此的.

$^{521}$向旁邊祝福.
願聖靈以醫治的「游」.
去告你.
使你身體,靈魂得以醫治.
在哀悼的當中.
我們祈求主去填補我們心靈當中的空虛.
在當中.
執著神給我們應許哀悼的人有福.
因為我們要見到神.
見到耶穌和我們同行.
耶穌的醫治.
可能在平常的生活當中.
可能在你從旺角回到家中.
或者帶領你家裡的DSE.
或者是中六的學生.
或者是更小的小朋友.
包括我家裡有一個侄子.
是幼稚園畢業的.
但他好像沒讀過書.
就畢業了.
在平常的生活當中.
但願我們都能夠看見.
耶穌藉著聖靈和我們同在.
透過日常.
在崇拜當中.
可以見到,摸到,領受的上帝的恩典.
體會祂的臨在.
祂的恩典是夠你用的.
我們彼此勉勵主的恩典夠你用.
願主賜福大家.
我們一起有個禱告.
你沒有撇下我們為孤兒.
你賜下聖靈.
讓我們能夠稱耶穌基督為主.
當我們生命當中有不同的破碎.
願主你以你的身體寶血.
以你的話語.
去醫治我們內心的創傷.
透過我們的崇拜.
我們的歌頌.

$^{561}$透過我們的生活當中.
再用感恩的心去看主你.
在我們生命當中.
日常當中.
每個細碎的恩典.
願聖靈你像今季的破碎器皿.
用溫柔的手將我們纏繞.
以你的大能讓我們看到.
在很多不明白當中.
其實有主你的旨意.
在面對疫情.
在面對流散的挑戰當中.
願基督你向我們顯明.
你在我們生命當中.
你的話語有能力.
能夠醫治我們的身心靈.
願你以喜樂的由高我們.
以致我們再唱得救的樂歌的時候.
能夠榮神益人.
奉主耶穌基督得勝名求.
\newpage



\section{尼希米記 6:1-16}
\label{sec:SQA6e1_jseU}
\textbf{2022年8月7日崇拜直播｜吳真真牧師｜並肩勇闖系列:戰勝懼怕｜尼希米記6\_1-16}
\newline
\newline
連結: \href{https://youtube.com/watch?v=SQA6e1_jseU}{\texttt{https://youtube.com/watch?v=SQA6e1\_jseU}} ~~~~ 語音日期: 2022-08-07
\newline
\newline
\hyperref[sec:sq3x1g07xfo]{\small{< < < PREV SERMON < < <}}
~
\hyperref[sec:index_chronic]{\small{[返順時目]}}
~
\hyperref[sec:IiBVTMF9Vak]{\small{> > > NEXT SERMON > > >}}
\newline
\newline
尼希米記 6:1-16
\newline
\begin{longtable}{cl}
\hline
\hline
章節 & 經文 (和合本修訂版)\\
\hline
6:1 & \begin{tabularx}{0.7\textwidth}{X} 參巴拉、多比雅、阿拉伯人基善和我們其餘的仇敵聽見我已經建造了城牆，沒有破裂之處在其中，那時我還沒有在城門安門扇； \end{tabularx} \\ \\ \relax
6:2 & \begin{tabularx}{0.7\textwidth}{X} 參巴拉和基善就派人來見我，說：「請你來，我們在阿挪平原的村莊見面。」其實，他們想要害我。 \end{tabularx} \\ \\ \relax
6:3 & \begin{tabularx}{0.7\textwidth}{X} 於是我派使者到他們那裡，說：「我正在進行大的工程，不能下去。我怎麼能離開，下去見你們，而讓工程停頓呢？」 \end{tabularx} \\ \\ \relax
6:4 & \begin{tabularx}{0.7\textwidth}{X} 他們這樣派人來見我四次，我都用這話回答他們。 \end{tabularx} \\ \\ \relax
6:5 & \begin{tabularx}{0.7\textwidth}{X} 參巴拉第五次同樣派僕人來見我，手裡拿著未封的信， \end{tabularx} \\ \\ \relax
6:6 & \begin{tabularx}{0.7\textwidth}{X} 信上寫著：「列國中有風聲，基善也說，你和猶太人謀反，所以你建造城牆。據說，你要作他們的王， \end{tabularx} \\ \\ \relax
6:7 & \begin{tabularx}{0.7\textwidth}{X} 並且你派先知在耶路撒冷指著你宣講說，『在猶大有王。』如今這些話必傳給王知，現在請你來，我們一起商議。」 \end{tabularx} \\ \\ \relax
6:8 & \begin{tabularx}{0.7\textwidth}{X} 我就派人到他那裡，說：「你所說的這些事，一概沒有，是你心裡捏造的。」 \end{tabularx} \\ \\ \relax
6:9 & \begin{tabularx}{0.7\textwidth}{X} 他們全都要使我們懼怕，說：「他們的手必軟弱，不能工作，以致不能完工。」現在，求你堅固我的手。 \end{tabularx} \\ \\ \relax
6:10 & \begin{tabularx}{0.7\textwidth}{X} 我到了米希大別的孫子，第來雅的兒子示瑪雅家裡；那時，他閉門不出。他說：「我們可以在神的殿裡，就在殿的中間會面，鎖住殿門，因為他們要來殺你，要在夜裡來殺你。」 \end{tabularx} \\ \\ \relax
6:11 & \begin{tabularx}{0.7\textwidth}{X} 我說：「像我這樣的人豈會逃跑呢？像我這樣的人豈能進入殿裡保全生命呢？我不進去！」 \end{tabularx} \\ \\ \relax
6:12 & \begin{tabularx}{0.7\textwidth}{X} 我看清楚了，看哪，神並沒有派他，是他自己說預言攻擊我，是多比雅和參巴拉收買了他； \end{tabularx} \\ \\ \relax
6:13 & \begin{tabularx}{0.7\textwidth}{X} 收買他的目的是要叫我懼怕，依從他犯罪，留下一個壞名聲，好讓他們毀謗我。 \end{tabularx} \\ \\ \relax
6:14 & \begin{tabularx}{0.7\textwidth}{X} 我的神啊，求你記得多比雅、參巴拉、挪亞底女先知和其餘的先知，因他們行這些事，要叫我懼怕。 \end{tabularx} \\ \\ \relax
6:15 & \begin{tabularx}{0.7\textwidth}{X} 以祿月二十五日，城牆修完了，共修了五十二天。 \end{tabularx} \\ \\ \relax
6:16 & \begin{tabularx}{0.7\textwidth}{X} 我們所有的仇敵聽見了，四圍的列國就懼怕，愁眉不展，因為他們知道這工作得以完成，是出於我們的神。 \end{tabularx} \\ \\
[1ex]
\hline
\hline
\end{longtable}
$^{1}$《新教育法例》第十三章 第十一至十六節 舊約聖經596頁.
「森巴拉·多比亞 阿拉伯人基善和我們其餘的仇敵聽見我已經收完了城牆.
其中沒有破裂之處 哪時我還沒有安門善.
森巴拉和基善就打發人來見我 說:請你來 我們在阿羅平原的一個村莊相會.
他們卻想害我 於是我差遣人去見他們 說:我現在辦理大功不能下去 應能停工下去見你們呢.
他們這樣四次打發人來見我 我都如此回答他們.
森巴拉第五次打發僕人來見我 手裡拿著尾風的信 信上寫著說:.
「愛邦人中有風聲 加斯姆也說:你和猶大人謀反 修造城牆 你要作他們的王.
你又派先知在耶路撒冷指著你宣講 說在猶大有王.
現在這話必傳於王之 所以請你來與我們彼此相議」.
我就差遣人去見他說:「你所說的這事 一概沒有 是你心裡列造的 他們都要使我們懼怕.
二司說:他們的手必軟弱 以致工作不能成就 神啊 求你堅固我的手」.
我到了米希大別的孫子 蒂萊雅的兒子 史瑪雅家裡 哪時他閉門不出.
他說:「我們不如在神的殿裡會面 將殿門關鎖 因為他們要來殺你 就是夜裡來殺你」.
我說:「像我這樣的人豈要逃跑呢? 像我這樣的人豈能進入殿裡保存生命呢? 我不進去.
我看明神沒有差遣他 是他自己說話攻擊我 史多比亞和森巴拉賄賣了他.
賄賣他的緣故是叫我懼怕 依從他犯罪 他們很全揚惡言 毀謗我.
我的神啊 多比亞 森巴拉 女仙貞 羅亞底和其餘的仙姿要叫我懼怕 求你紀念他們所行的這些事」.
「已六月二十五日 城牆修園了 共修了五十二天 我們一切仇敵四處的外邦人聽見了便懼怕 愁眉不展 因為見這工作完成是出乎我們的神」.
聖經獨笆.
各位弟兄姊妹平安.
放暑假 有一個家庭 他們的母親建議去旅行 今年很久沒有去旅行 她和家人商量 她說不如去日本 爸爸說不要 會地震.
有人說不如去美國 我怕有喉痘症 不如去歐洲 我怕到時機場有罷工 不如去非洲 不行那裡有藥劑 我很怕.
隨便去東南亞也不算不好 因為回到香港又怕訂不到酒店 做隔離.
去旅行那麼麻煩 留在家裡 兒子打開冰箱 看見有冰淇淋 他拿出來吃的時候 媽媽說「你看看這冰淇淋 我怕有除害劑污染」.
兒子覺得沒趣 拿出一個令人有驚喜滿足願望的朱古力.
爸爸說「我怕有三門氏菌污染 前陣子不是回收的嗎?」 嘩 你覺得這個家庭充滿什麼呢?充滿懼怕 充滿害怕.
人生真的有很多令我們懼怕的事 懼怕的確會令我們裹足不前.
有些應該說的事我們不敢說 有些應該做的事我們不敢做.
上世紀有個英國哲學家叫羅素 他說了一句話.
他說如今對於美好世界的最大障礙就是恐懼 就是懼怕 懼怕對這個世界 對我們人類造成很多破壞.
不過認真看 懼怕未必一定是負面的 因為有時我們的恐懼能夠保護我們避開一些危險.
當你怕被車撞 你就會留心遵守交通規則去保護自己 幫助自己 是不是?.
但是懼怕也有它很破壞性的一面 尤其是我們這一班信徒.
我們要認清楚其實我們的仇敵撒旦 最喜歡用懼怕成為一個很厲害的武器.
去讓我們離開上帝 去讓我們不能夠去作成上帝給我們的使命.
所以 弟兄姊妹 我們真的要小心 我們千萬不要被懼怕拉著我們的人生走.
人生因為害怕而去活 於是有時候我們在工作裡面 我怕吃虧.
所以有時候我會用一些不符合真理的心態 或者一些不符合真理的意思去工作.
我怕沒錢 於是我很怕沒安全感 我怕所以我就拼命用我所有的時間用來賺錢.

$^{41}$用所有的方法用來賺錢 我很怕別人看不起我.
所以我要做一些事情令人看得起我 即使那些事情未必符合神的心意.
弟兄姊妹 我們不要小看懼怕的威力 但是我們今天要看的經文就是講尼希米.
尼希米記第六章一至十六節 我們就從這個經文裡面.
我們去明白究竟我們怎樣可以戰勝懼怕 讓我們一起祈禱.
天父我們祈求你幫助我們 幫助我們不要陷在一個懼怕的人生.
幫助我們能夠在生命裡面去定精向著你 依靠你.
好讓我們的生命覺得活潑 覺得合理心意 覺得豐盛.
求你的靈在我們當中帶領我們 幫助我們去明白你要跟我們講的話.
你對我們生命的提點 我們這樣禱告交託 是奉主耶穌基督的聖名自祈求 阿們.
今天尼希米記第六章一至十六節 我一定要講一些背景.
要講很長 簡單來說 其實猶太人在公元前586年被巴比倫滅亡之後.
他們的一些猶太人就被擄到巴比倫裡面.
當中有一些老弱殘兵 有些比較脆弱的人就留在耶路撒冷.
雖然他們是這樣 他們被擄 他們流散.
但是他們卻對於自己國家的宗教和文化並沒有忘記.
到後來波斯國滅了巴比倫國 古列王下令猶太人可以重新回到耶路撒冷去重建他們的家園.
於是有不少猶太人回到以色列地開始重建耶路撒冷 包括重建第二聖殿.
下星期陳明全牧師跟大家分享以色列記就是講這件事.
而今天我們講的是尼希米記.
記載了隨後尼希米帶領以色列民重建城牆的經過.
尼希米本來在波斯王阿達薩西身旁是一個很重要的官員.
他是酒政來的 但因為他是一個猶太人.
他心裡掛念他的同胞.
所以他向王提出要回到耶路撒冷重建城牆.
本來以為王答應了 開綠燈 一切美好了.
誰知在尼希米記一直記載 從第二章開始.
你已經看到困難重重 有外憂有內患.
第五章講到內患就是當中興建城牆的工人都在投訴.
說那些瓦礫太多 供桌子不夠氣力.
他們當中又有些貴族和官長貪心壓榨百姓.
第六章講到外憂.
經文提到有三個很重要的人物.
就是森巴拉,多比亞和阿拉伯的基善.
他們就是外憂的仇敵.
在這段很清楚地告訴我們.
經文裡又重覆地告訴我們.
仇敵要對付他們的伎倆是甚麼?.
就是使他們懼怕.
從第九節,第十三節,第十四節.

$^{81}$甚至我們今天還沒提及的第十九節.
第六章裡有四個很清晰地告訴我們.
仇敵的伎倆是要使他們懼怕.
尤其是尼希米要懼怕.
但很奇妙的在第十五節至第十六節.
你會看到四圍的列國就懼怕.
因為上帝的工程完成了.
所以我們今天就來看看.
究竟這群仇敵用甚麼招數.
令尼希米那麼懼怕呢?.
我們一起看第一至第四節.
「森巴拉多比亞阿拉伯人基善和我們其餘的仇敵聽見我們已經收完了城牆.
其中沒有破裂之處.
那時我還沒有安門時.
森巴拉和基善就打發人來見我.
說請你來我們在阿羅平原的一個村莊相會.
他們卻想害我.
於是我差一點去見他們.
說我現在辦理大功不能下去.
應能停工下來見你們.
他們這樣四次打發人來見我.
我都如此回答他們」.
其實森巴拉原來就是撒瑪利亞的省長.
位於耶路撒冷的北邊.
多比亞是住在阿滿.
在耶路撒冷的東邊.
另外基善可能在阿拉伯.
位於耶路撒冷的南邊.
第四章第七節亦提到有些叫阿薩特人.
他們在耶路撒冷的西邊.
你想像一下.
耶路撒冷周圍東南西北都佈滿仇敵.
他們之前有份去遊論耶路撒冷.
到城牆快要建好的關頭.
還差十個城門要建好的時候.
他們把握最後一擊.
要將他們的領袖尼希米一拳擊中.
他們的計劃是派人去跟尼希米說.
「我們去阿羅平原的村莊見面」.
他們的語氣是表示.

$^{121}$好像之前是仇敵關係很緊張.
現在我要跟你示好.
不如我們出來見面緩和緊張的氣氛吧」.
他們邀請的地點都很有智慧.
阿羅平原位於耶路撒冷西北面三十多公里.
其實這地方就在耶路撒冷和撒瑪利亞的中間點.
很多解經家說這是一個中立性的區域.
可以大大減低尼希米的戒心.
他們很有策略.
如果他們說「我們過來耶路撒冷見面」.
尼希米就一定有很嚴密的保護.
但如果出到阿羅平原的話.
他們要綁架他,要殺他.
其實就輕易很多.
所以這是一個非常危險的鴻門宴.
不過尼希米洞悉他們的陰謀.
他說「我現在正在進行大工程.
我怎可以離開去見你們呢?.
我怎可以讓工程停頓呢?」.
仇敵四次引誘.
尼希米四次堅定拒絕.
我們的仇敵魔鬼都會用一些計謀.
使我們失去焦點.
企圖阻止我們邁向一個有意義.
倚靠上帝豐盛的生命目標.
四次邀請不行.
於是就出第二招.
他們就第五次打發人去.
不過這次打發人去就很巧妙.
他就送一封信.
但那封信是沒有封口的.
正常這些這麼高規格的信件.
是會封口的.
當時封口是怎樣呢?.
就是用一些葡草紙或皮革.
上面寫了一些意思.
然後就捲起.
然後就用一條繩線.
帶著有印的圖製小球.
圖製小球.

$^{161}$這樣封住捲住.
收信的人收到那封信密函的時候.
就打碎圖製的小球.
才可以讀那封信.
但這群仇敵竟然用一封沒有封口的信.
目的是甚麼?.
目的就是這個人帶著這封信.
去到哪裡.
裡面的內容就傳到哪裡.
確保信內容可以傳開去.
信內容是怎樣說的?.
就是這樣說的.
第六節.
「愛邦人中有風聲,迦斯慕(基善).
你同猶大人謀反,修造城牆.
你要作他們的王」.
這封信的內容原來是一些謠言.
是捏造他們,說他們.
其實建城牆的目的是甚麼?.
不是說我想甘願做一個.
在波斯王國裡的猶大省.
是你們根本就是想密謀造反.
自封為王.
說敵人用這一招造謠.
抹黑,恐嚇.
尼熙米這次也是一樣.
他沒有中計.
第九節說.
「他們都要使我們懼怕.
意思說他們的手必軟弱.
以至工作不能成就.
神啊,求你堅固我的手」.
其實尼熙米認清他們這群人.
在製造謠言,製造恐慌.
目的就是要他們懼怕.
懼怕是真的會令他們害怕.
會真的令他們放軟手腳.
真的可能要停工.
讓他們知道.
他們不是想去圖謀造反.

$^{201}$原來受敵的作為.
是可以叫我們放軟手腳.
叫我們不敢去做我們應該做的事.
第三個計謀來了.
第十至第十四節.
「我到了米希大別的孫子.
蒂萊瓦的兒子,斯瑪瓦加里.
那時他閉門不出.
他說:我們不如在神的殿裡會面.
將殿門關鎖.
因為他們要來殺你.
就是夜裡來殺你.
我說:像我這樣的人豈要逃跑呢?.
像我這樣的人豈能進入殿裡保存生命呢?.
我不進去.
我看明神沒有差遣他.
是他自己說話攻擊我.
是多比亞和森巴拉賄賣了他.
賄賣他的緣故是要我懼怕.
依從他犯罪.
他們很全然惡言,誹謗我.
我的神啊!多比亞,森巴拉,呂先知,羅亞底.
和其餘的先知要我懼怕.
求你紀念他們所行的這些事.
兩次的招數.
叫李希米離開耶路撒冷.
去阿羅平縣跟他們談.
這一招用不著.
於是他們設計了一招.
不如在耶路撒冷裡面見面吧.
這次這一招厲害了.
不只是城外,耶路撒冷以外那幾個國家的人.
去攻擊李希米.
這次是賄賂他們當中.
自己猶太人裡面的人.
這個人是誰?.
這個人是施瑪雅.
從這段經文裡知道.
他是一個講預言的先知.
他是一個可以入聖殿的祭司.

$^{241}$他是一個人所共知的宗教領袖.
但是他竟然被賄賣.
李希米來見他.
他彷彿講了一個神諭.
講了一個預言.
告訴他:仇敵要來殺你.
經文直譯就是今晚就要來殺你.
他還很好.
他彷彿提供了一個出路給他.
那個出路是什麼?.
你跟我在聖殿裡面會面吧.
我們將電門關鎖.
敵人就進不來了.
你就在一個安全的處境.
同樣,李希米也是沒有中計.
他知道上帝沒有猜這個所謂先知,祭司,宗教領袖.
去講上帝的話.
他反而知道他講這番話.
是多比亞牽頭的這班人去指使,設的局.
目的是什麼?.
目的是使李希米懼怕.
是要他自己擔心自己的安危.
還引誘他,迫使他做出一些.
得罪上帝,得罪人的舉動.
為什麼這樣說?.
他是一省之長.
如果他逃跑.
就是一個恐懼,懦弱和自我保護的表現.
他只顧著自己的安危.
他怎樣成為省內居民的榜樣?.
他必定在他的百姓面前無地自容.
第二,其實一般的信徒是不可以入聖殿.
即使他是位高權重的省長.
他都沒有資格入聖殿.
因為這不是上帝賦予他的權利.
這是犯上宗教的罪行.
是得罪上帝.
所以敵人這次的計謀更加厲害.
第一,第二次,如果他中計的話.
以色列人還懂得分清敵,我.

$^{281}$但如果這次他中計.
李希米就會在以色列民面前大失威信.
身敗名裂,成為眾矢之的.
整場興建的計劃就會遭受破壞.
我們很快看完這個故事.
是否已經說完了?當然不是.
我們看到的是,仇敵會用盡方法.
使基督徒的生命失去焦點.
會令我們害怕,膽怯,變得軟弱.
將焦點放在自己身上.
甚至會犯下嚴重的錯.
最終不能完成上帝給我們的使命和託付.
但我們感恩,我們看到的是.
李希米沒有中過任何奸計.
為何他會有生命上的質素?.
我們一起去看看.
他如何成為我們這群信徒的榜樣.
他可以最終戰勝恐懼.
我覺得最重要的是.
他清楚知道,仇敵要使他懼怕,軟弱.
但上帝使我們剛強,堅固.
李希米如何告訴我們.
或他的生命如何展現.
他相信上帝使我們剛強,堅固.
第一,我看到的是.
他很清楚和專注地回應上帝的召命.
首先,我們知道.
神感動李希米回到耶路撒冷建造城牆.
當初他是亞達薩西王身邊的重要官員.
身居要職.
但當他聽到他的同胞.
在耶路撒冷遇到阻撓.
遇到靈性的威脅.
城牆被拆毀,城門被焚燒.
他極其心痛.
他向上帝禁食,祈禱.
等候了足足四個月.
他很清晰這是上帝的呼召.
於是他大膽向王提出.
回到耶路撒冷重建城牆.

$^{321}$他也很清楚.
回到耶路撒冷重建城牆.
他放下自己的安殊環境.
放下自己的高薪厚職.
很切實地回應上帝的呼召.
在他整個建造城牆裡.
沒有一天安寧.
但他接受一次又一次的挑戰.
我們卻在他身上看到的是.
他高度地尊重上帝.
對他作出呼召的使命進發.
我們看到當敵人找他的時候.
他很清楚地告訴他.
「如果你現在叫我離開」.
「我現在做的大事」.
「我人生裡第一要做的事」.
「我就要放下」.
「我不會放下去做一些」.
「我人生最重要的事情」.
「而去照顧其他東西」.
當仇敵恐嚇他要保命的時候.
第十一節他說.
「像我這樣的人豈要逃跑呢?」.
他很清楚地記得.
他是耶路撒冷城的省長.
他很清楚自己的崗位.
他很清楚自己的任務.
當他很清楚自己的目標.
自己的方向的時候.
他不會因為其他旁支而失去分寸.
而不知道自己在哪裡.
會有很多擔心.
專注在自己.
而不專注在上帝的事.
弟兄姊妹.
其實我們呢.
有時候我們很容易很害怕.
我們會很容易失去方寸.
但我很想問大家.
我們知不知道自己的目標是什麼?.

$^{361}$我們知不知道自己的召命是什麼?.
你知不知道上帝想你怎樣?.
上帝對於我們這群信徒的召命.
是很清晰.
祂要我們愛神.
祂要我們愛人.
祂要我們傳福音.
祂要我們去作耶穌基督的.
使人作耶穌基督的門徒.
祂要我們過聖潔的生活.
如果我們因為這個世界.
對我們製造了很多的懼怕.
而放下了這個召命.
沒有去踐行信徒的積分.
我們就是過著一個.
被懼怕拉扯著的人生.
所以弟兄姊妹.
我們很需要學尼熙米.
詩篇第十六篇第八節說.
「我將耶和華常擺在我面前.
因祂在我右邊.
我便不自動搖」.
當我們清晰我們作為基督徒的使命.
召命的時候.
這些挑戰.
這些令我們懼怕困難的事.
不會令我們動搖.
我們這個主題是.
這麼久講到的主題是.
「並肩勇闖」.
旺枕六十年了.
我們要邁向下一個六十年.
是需要無比的勇氣去爭戰.
但我經常都很想問.
我們「並肩勇闖」去哪裡?.
弟兄姊妹.
其實我相信每一個人.
上帝都有擺照明在你們的心裡.
我們會因著我們的感動.
我們的負擔.

$^{401}$我們心裡有些心之所繫.
我們有我們的恩賜.
有我們的能力.
我們有我們的性格.
我們有我們的際遇.
我們在旺枕.
弟兄姊妹.
當我們「並肩勇闖」.
當旺枕在這個時代裡面.
有這麼多的挑戰.
有很多教會因為疫情的關係.
因為有很多會眾移民的關係.
很多教會正在萎縮.
但旺枕「並肩勇闖」.
我們正正是福音的工作.
沒有停過地不斷地擴展.
我們有很多的需要.
好像我服侍的地方.
我服侍有年部.
我們其實少了很多的導師.
我們現在缺乏很多的導師.
在疫情之下.
很多的小朋友發展緩慢了.
因為他們沒有機會出去外面社交.
沒有機會有很多的訓練.
很多的刺激幫他們成長.
我們很缺乏這些人手.
我們亦都照顧了很多基層的群體.
很多的家長.
他們因為經濟的困難.
因為家庭裡面.
疫情裡面常常都有困獸鬥的情況.
家人有很多的需要.
弟兄姊妹.
當我們說我們要去關顧他們.
要將福音帶給他們的時候.
你願不願意和我們一起「並肩勇闖」.
你願不願意按著上帝在你身上.
個人的照明.
在你的個性,你的際遇.

$^{441}$你願不願意去回應他.
第二,我會看到的就是.
李希米是一個隨時隨地向神禱告的人.
《李希米記》是一本.
有一個神學的主題.
就是李希米和上帝有很親密.
很活潑,很深切的關係.
《李希米記》記載了李希米禱告十二次.
在第六章.
記載了其中的兩次.
分別是第九節和第十四節.
第九節和第十四節有一個共通點.
就是當敵人要叫他懼怕的同時.
他就向上帝禱告.
第九節他求上帝堅固他的手.
第十四節他求上帝要詆毀抹黑傷害他的受敵.
當一個人.
當一個人常常與主親近.
常常禱告.
我發覺他很懂得祈禱.
我發覺他很懂得上帝的心意.
我發覺慢慢他的自我越來越放低.
上帝的真理,上帝的心意.
是放在了他的心裡.
當人誹謗他.
說他要叛變的時候.
他沒有為自己伸冤.
他只說了一句話.
「你所說的這事一概沒有.
是你心裡捏造的」.
接著他就祈禱.
求上帝堅固他的手.
我覺得他很有型.
如果被人誹謗,抹黑.
明明自己不是.
你實開記者招待會.
列明很多的實例,證據.
說自己不是.
但他不是.
他沒有為自己辯護過一句.

$^{481}$他只是對對方說「我不是」.
然後他就祈禱.
他將自己的需要.
將自己的聲譽放在一邊.
乾脆交給上帝.
到了第十四節.
當人要謀害他.
誘使他得罪上帝的時候.
他也繼續沒有停工.
集中焦點建造城牆.
但他的禱告是.
求主紀念多比亞,森巴拉.
女先知羅亞底.
和其餘先知所做的惡事.
面對抹黑.
他沒有唱衰他們.
他沒有以惡報惡.
他只求上帝公義的審判.
並且保護他免受敵人的傷害.
所以尼希米.
他一個這麼活潑的靈命.
和神這麼多相交的領袖.
他不慌不忙.
他可以經常保持清晰的頭腦.
有分辨的能力.
一看就知道敵人的意圖.
我很深信這份能力.
就是因為他經常和上帝緊密相交.
他很清楚上帝的心意.
所以我看到的是.
弟兄姊妹.
當我們奮勇向上帝.
給予我們的目標邁進的時候.
與此同時.
我們和上帝的關係.
是必須非常緊密.
以至我們自己人.
慢慢的心思意念想法.
就像我們的上帝.
第三.

$^{521}$第十五至第十六節.
「以六月二十五日成祥修完了.
共修了五十二天.
我們一切仇敵.
四圍的愛邦人聽見了便懼怕.
愁眉不展.
因為見者工作完成.
是出乎我們的神」.
成祥在這麼多艱難挑戰下.
用了五十二天去建成.
是不可思議的事.
第十六節.
你看看那群仇敵.
聖經如何形容他們?.
聖經說.
現在害怕的不是尼希米.
而是猶太人.
現在害怕的是誰?.
現在害怕的是那群仇敵.
他們害怕到.
聖經的原文說.
他們垂下他們的眼睛.
表達的是他們意志消沉.
表達的是.
當初尼希米回去建成祥之前.
這個城是被欺凌.
被欺壓.
是很卑賤,很羞愧.
但今天成祥被建起來的時候.
這個城.
這群子民.
可以昂然抬起頭.
因為上帝的工作在他們當中.
但奇蹟地完成了.
而這群仇敵.
就變得羞愧,羞辱.
我覺得這裡很好的反思.
讓我們一起有很好的反思.
其實我們值得思考.
我們要害怕.

$^{561}$應該害怕誰?.
我們面對無數的懼怕.
如何可以化軟弱.
變成堅固?.
根源就是要敬畏神.
敬畏神這個畏懼.
其實就可以除去其他的畏懼.
害怕神.
就不用害怕其他東西.
我們看到第六章出現了五次懼怕.
四章十四節都出現了兩次.
那兩次的懼怕是甚麼?.
那兩次的懼怕告訴我們.
是不要怕他們.
當記得主是大而可畏的.
敬畏神就不需要懼怕困難.
不需要懼怕仇敵.
我想起哥林多後書四章八至九節.
提到保羅說.
我們四面受敵,卻不被困住.
心裡作難,卻不致失望.
遭迫伯,卻不被掉棄.
打倒了,卻不致死亡.
原因是甚麼?.
是第七節.
我們有者寶貝放在瓦器裡.
要顯明者莫大的能力是出於神.
不是出於我們.
我們只是瓦器.
就算尼希米.
就算保羅是一個多麼偉大的領袖.
他們都要坦然承認.
我們只不過是瓦器.
但我們的能力為何會有這麼大的智慧.
不需要畏懼.
就是因為有寶貝在我們的生命裡.
有耶穌基督的救恩在我們的生命裡.
所以我們應該要思考.
我們不要將焦點放在害怕的事情上.
我們反而要轉念思考.

$^{601}$神啊,你要我們做甚麼?.
有一年,阿Paul是我的老公.
他在港交所裡做了一件事.
就是他們要公佈一些消息給外界.
然後讓人去做一些買賣.
那次很特別的就是.
他們將一個小數點.
不知道是撥前還是撥後.
很大件事.
你可以想像在港交所裡.
你公佈消息撥前和撥後.
是影響了很多.
當時在南華早報頭版的一個角落.
就寫了他的名字.
其實那個出錯當然是他的同事做的.
但作為上司,他是責無旁貸.
不過他心裡很想逃避.
他很害怕,他很害怕承認責任.
雖然不是他做錯.
但他可能最初也會想方法.
明明是我的下屬做錯的.
但是他心裡很害怕失去工作.
很害怕,不知道如何面對老闆.
但他只想了一句.
「神啊,你要我做什麼?」.
「作為你的門徒,你要我做什麼?」.
結果他敲門進去老闆房.
坦然承擔責任.
所以,結果,放心.
他不是因為被解僱而讀神學.
他是繼續好好地做好這份工作.
直到上帝呼召他為止.
我看過一篇文章.
那篇文章說空中飛人在練習的時候.
下面有一個很大的安全網.
雖然是這樣.
當他們剛剛開始練習的時候.
他們覺得很害怕.
就算有安全網,他們也很害怕.
究竟如何克服呢?.

$^{641}$原來,除非他們有掉進安全網的經驗.
當他們掉進去的時候.
他們才實實際際地經歷到.
這個安全網是真的安全.
掉進去是沒事的.
當他們擁有了很多掉進去的經歷之後.
他們就越來越不怕掉進去.
他們對於安全網的信心就增加了.
給了他們勇氣.
其實一樣,當我們做上帝的事的時候.
我們很害怕掉出去安恕區外.
不過當我們要對神有信心.
當然我們要對神有信心.
這樣我們才可以除去懼怕.
我們可以選擇畏懼的人生.
讓懼怕追著我們的人生.
但我們也可以選擇.
上帝在我們前頭的人生.
是一個平安的人生.
唯有敬畏上帝.
才能夠除去恐懼.
化軟弱為堅固.
夠了.
\newpage



\section{約翰福音 1:29-37}
\label{sec:IiBVTMF9Vak}
\textbf{2022年8月14日崇拜直播｜梁家麟牧師｜施洗約翰的見證(三)｜約翰福音1\_29-37}
\newline
\newline
連結: \href{https://youtube.com/watch?v=IiBVTMF9Vak}{\texttt{https://youtube.com/watch?v=IiBVTMF9Vak}} ~~~~ 語音日期: 2022-08-16
\newline
\newline
\hyperref[sec:SQA6e1_jseU]{\small{< < < PREV SERMON < < <}}
~
\hyperref[sec:index_chronic]{\small{[返順時目]}}
~
\hyperref[sec:ZKaifzpEjak]{\small{> > > NEXT SERMON > > >}}
\newline
\newline
約翰福音 1:29-37
\newline
\begin{longtable}{cl}
\hline
\hline
章節 & 經文 (和合本修訂版)\\
\hline
1:29 & \begin{tabularx}{0.7\textwidth}{X} 第二天，約翰看見耶穌來到他那裡，就說：「看哪，神的羔羊，除去世人的罪的！ \end{tabularx} \\ \\ \relax
1:30 & \begin{tabularx}{0.7\textwidth}{X} 這就是我曾說『那在我以後來的先於我，因為在我以前，他已經存在』的那一位。 \end{tabularx} \\ \\ \relax
1:31 & \begin{tabularx}{0.7\textwidth}{X} 我先前不認識他，如今我來用水施洗，為要使他顯明給以色列人。」 \end{tabularx} \\ \\ \relax
1:32 & \begin{tabularx}{0.7\textwidth}{X} 約翰又作見證說：「我曾看見聖靈彷彿鴿子從天降下，停留在他的身上。 \end{tabularx} \\ \\ \relax
1:33 & \begin{tabularx}{0.7\textwidth}{X} 我先前不認識他，可是那差我來用水施洗的對我說：『你看見聖靈降下來，停留在誰的身上，誰就是用聖靈施洗的。』 \end{tabularx} \\ \\ \relax
1:34 & \begin{tabularx}{0.7\textwidth}{X} 我看見了，所以作證：這一位是神的兒子。」 \end{tabularx} \\ \\ \relax
1:35 & \begin{tabularx}{0.7\textwidth}{X} 又過了一天，約翰同兩個門徒站在那裡。 \end{tabularx} \\ \\ \relax
1:36 & \begin{tabularx}{0.7\textwidth}{X} 他見耶穌走過，就說：「看哪，神的羔羊！」 \end{tabularx} \\ \\ \relax
1:37 & \begin{tabularx}{0.7\textwidth}{X} 兩個門徒聽見他的話，就跟從了耶穌。 \end{tabularx} \\ \\
[1ex]
\hline
\hline
\end{longtable}
$^{1}$今日的經訓是約翰福音第一章二十九至三十七節.
在聖經新約第一百二十六頁.
請聽施浸約翰為上帝的兒子耶穌基督神的羔羊作的見證.
此日約翰看見耶穌來到他那裡.
就說:看那神的羔羊除去世人罪孽的.
這就是我曾說:有一位在我耳後來.
反成了在我耳前的,因他本來在我耳前.
我先前不認識他,如今我來用水施浸.
為要叫他顯明給以色列人.
約翰又作見證說:我曾看見聖靈彷彿甲子從天降下.
住在他的身上,我先前不認識他.
只是那差我來用水施浸的.
對我說:你看見聖靈降下來.
住在誰的身上,誰就是用聖靈施浸的.
我看見了,就證明這是神的兒子.
再次日,約翰和兩個門徒站在那裡.
他見耶穌行走,就說:看那,這是神的羔羊.
兩個門徒聽見他的話,就跟從了耶穌.
我們繼續看斯德林約翰對耶穌基督的見證.
斯德林約翰接受了從耶路撒冷公會來的人對他的身份和使命的查問之後.
第二天,他在約旦河邊和耶穌相遇.
為什麼耶穌在那裡呢?.
因為沒多久前,耶穌主動去找約翰接受他的浸禮.
這是《馬太福音》3:13-17的記載.
然後,在接受浸禮之後,他留在約翰所屬的修道群體.
這個群體類似昆蘭社團的修道群體.
耶穌在約旦河邊的曠野住了一段日子.
直到斯德林約翰坐牢之後,他才離開.
回到加利利,然後開始自己的傳道生涯.
耶穌住在這個修道群體中間.
所以接下來的兩天,斯德林約翰也碰到耶穌.
這不奇怪,因為群體很小.
很有趣的是,《約翰福音》沒有直接提及耶穌受浸的重要事件.
只是記述了斯德林約翰的兩段說話.
就是他為耶穌所做的見證.
作為耶穌受浸的情景的一個說明.
這樣的鋪排,用一個轉述的方法.
而不是直接記述那個事件.
跟《三卷福利福音》是完全不同的.
《三卷福音》基本上都是記述整件事件.

$^{41}$但是《約翰福音》沒有提及這件事.
《約翰福音》寧願引述斯德林約翰的說話.
沒有直接去敘述耶穌受浸的故事.
這大概是貫徹他在福音書的序言所說.
斯德林約翰一生其實只有一個使命.
不是在約旦河邊為人施浸教訓人真理.
不是招收門徒,講傳認罪和赦罪的福音.
而是為耶穌做見證.
是的,斯德林約翰一生只有一個使命.
就是為耶穌做見證.
我們重溫《約翰福音》一章六到八節的經文.
有一個人是從上帝那裡猜來的,名叫約翰.
只因來為耀作見證,就是為光做見證.
叫眾人因他可以信,他不是那光,那是要為光作見證.
這是在序言,《約翰福音》一開始就說的.
所以斯德林約翰的說話.
就是他一生要做的工作.
他在見證著.
見證的目的就是叫眾人能夠相信耶穌.
用今天講的方法.
當施浸約翰為耶穌做見證的時候.
他在做事.
他為人施浸,甚至為耶穌施浸.
都不算在做事.
唯有他開口為耶穌做見證.
他才是在做事.
即是實踐他的核心使命.
我努力預備講章.
我認真地講道.
我是不是在做傳道的本業呢?.
也是,也不是.
傳道人當然要承擔講道的職責.
但傳道人的真正本業是傳揚耶穌基督.
見證耶穌基督.
如果我所預備的,我所傳講的.
不是扣定在為耶穌做見證.
這個核心使命,終極使命之上.
我就是失焦迷路,忘記本業.
嘉義文說,基督徒的人生只有一個目標.
就是榮耀上帝,見證基督.

$^{81}$基督徒一生的使命是榮耀上帝,見證基督.
我由衷盼望.
我們每個人,每天都在做事.
各位皇朕的兄弟,做事的話.
我們一起做事,見證我們的主.
這是我們最核心的使命.
我們今天讀的聖經.
其實是記載了斯証約翰三段見證的說話.
第一段說話是在約翰看到耶穌來到約旦河邊的時候所講的.
這是有關耶穌的身份的介紹.
「看!上帝的羔羊除去世人罪孽的.
這就是我曾說有一位在我以後來的.
反成了在我以前的,因他本來在我以前.
我先前不認識他,如今我用水來施浸.
為要叫他顯明給以色列人」.
約翰用了一個舊約的意象.
一個imagery,就是羔羊來介紹耶穌.
「看!上帝的羔羊除去世人罪孽的」.
耶穌好比在御節被殺的羔羊.
那個passover lamb.
說要用血去塗抹去遮蓋世人的罪.
這個稱呼在這段聖經一共出現了兩次.
第二天,當施浸約翰和兩個門徒在約旦河邊再遇耶穌的時候.
也是直接稱呼他「看!這是上帝的羔羊」.
所以約翰用了兩次「上帝的羔羊」去描述耶穌.
耶穌是上帝的羔羊,意思就是上帝所預備的羔羊.
上帝定義犧牲耶穌的性命.
讓耶穌替代人類承擔罪責.
在猶太人的傳統中.
羔羊被殺通常是指死亡的替代.
就好像創世記22章.
上帝為亞伯拉罕預備了一枝羊羔去替代耳殺.
大家記得這個故事吧.
今天,因著聖經的傳統.
所以我們在日常生活中常常會說什麼「代罪羔羊」.
即是找一隻不幸的羊來替我們付罪.
替人找數.
人犯錯,羊代成罪責.
如果施禎約翰認定耶穌是彌賽亞.
那麼彌賽亞要替人贖罪,替死.

$^{121}$這不應該是他或其他猶太人對彌賽亞的一貫理解.
猶太人理解彌賽亞應該是強者而非弱者.
應該是以榮耀和大能的面貌顯現.
打敗所有壓迫猶太人的敵人.
不過,舊約聖經並不是完全沒有彌賽亞是受苦的人的說法.
在《爾賽亞書》53章6-7節也有提及彌賽亞承擔眾人的罪孽.
就像羊羔被牽到待宰殺之地.
所以,施禎約翰如果將彌賽亞的使命理解為除去世人罪孽的上帝羔羊.
可能是源自這段舊約聖經作為基礎.
當然,在彌賽亞還未出現之前.
誰也不知道彌賽亞會做些什麼.
所有關於彌賽亞會做些什麼.
他是強者還是弱者.
他是失滅受敵還是代罪犧牲.
這些說法全部都是推論.
沒有人能夠原則上排除彌賽亞會替人贖罪受死的可能性.
在《楊仁福音》後面的11章49-52節.
當耶穌宣告他要被殺害之後.
很有趣的,那裡記載了一段說話.
大祭司蓋亞法這樣說.
彌賽亞藉著犧牲自己的性命去救贖整個國家是不無可能的事.
即是說,連大祭司蓋亞法都不排除彌賽亞可能不是以強者的面貌.
而是以替死的面貌出現.
所以,如果連蓋亞法都不否認彌賽亞有贖罪替死的使命.
施禎若汗有這個想法就不是不可思議了.
不過,我們也可以將施禎若汗這個形象和他個人的使命連繫起來.
《盧格弗音》記載,當阿拿和蓋亞法當大祭司的時候.
撒迦利亞的兒子若汗在曠野裡.
上帝說留不到他.
他來到約旦河一帶的地方宣講悔改的洗禮,洗罪,得赦.
施禎若汗是從上帝那裡領受悔改的浸禮,洗罪,得赦.
斥責人的罪惡,呼籲人認罪,悔改.
然後接受浸禮這樣的使命.
這是他一生的使命.
但是若汗的心裡知道這個悔改的浸禮的功用很有限.
非常有限公司.
沒有辦法真正解決人罪惡的問題.
唯有是耶穌的替死犧牲才真正能夠除掉世人的罪孽.
他做不到.
他一直在河邊做浸禮其實都做不到.

$^{161}$所以他見到耶穌的時候就看到.
這個才是真真正正除掉世人罪孽的那一位.
換言之施禎若汗是從自己所承擔的使命.
以及對自己的使命有限性的了解.
然後去凝訊宣告耶穌基督上帝的羔羊.
這是除掉世人罪孽的那一位.
我知道我做不到.
然後我看到是他才做得到.
我不能,主能.
施禎若汗沒有成就什麼屬於他自己的功業.
他徹頭徹尾是一個鋪路者.
為耶穌基督這個主角的登場做預備.
預備他的道,修正他的路.
一個很甘心,很出色去扮演回鄉的角色的一位.
然後施禎若汗沒有糾纏耶穌的救贖使命這個主題.
他卻是指出耶穌和他兩個人的關係.
他說耶穌就是他先前多次跟眾人提過的那一位.
三十節他說,這就是我曾說,有一位在我以後來的,反成了在我以前的,因他本來在我以前.
前面一章十五節其實已經引述施禎若汗的這一句說話.
所以這句說話其實是說第二次的.
耶穌的出道比施禎若汗遲,但是後發先至,很快就爬了他的頭.
這不是偶然發生的情況,而是從永恆的角度來說.
先傳永存的耶穌本來就在若汗的前頭.
前頭不僅僅是時序上的先後,地位和重要性.
若汗接著說,我先前不認識他,如今我來用水施浸.
我要叫他顯明給以色列人.
在前後這兩段正詞裡面,施禎若汗都反覆強調,我先前不認識他.
我先前不認識他,三十一節,三十三節.
這不是指著他之前完全不認識耶穌,他跟耶穌是親戚,怎會不認識呢?.
而是他從前是不知道,是未覺悟耶穌是彌賽亞.
直到他跟耶穌受浸的那一刻,他見到上帝為耶穌的神聖撐牆.
上帝做的見證,然後才明白耶穌的屬靈身份.
這是第一段施禎若汗的見證.
然後三十二到三十四節是第二段的見證.
那裡就講到有關耶穌的浸泥,將耶穌受浸的情景詳細描述出來.
三十二到三十四節.
我曾看見聖靈彷彿甲子從天降下,住在他的身上.
我先前不認識他,自剩那差我來用水施浸的.
對我說,你看見聖靈降下來住在誰的身上.
誰就是用聖靈施浸的.

$^{201}$我看見了,就證明這是上帝的兒子.
三卷符律福音都有記述耶穌在約旦河受浸的時候.
有聖靈如甲子從天降下降在他身上.
然後天上有上帝發聲宣告耶穌是上帝的愛子.
是天父所喜悅的那一位.
我不需要引述經文,大家很熟悉這個故事.
施禎若汗將一個客觀的歷史事件變成他個人的敘述.
他看見我曾看見完成式一個perfect tense的事態.
當耶穌受浸的時候.
他看到聖靈彷彿甲子從天降下住在他的身上.
可能你也看過一些宗教的圖畫,也說到這些情景.
耶穌在河水中間站立,上身打赤Lap .
然後雙手合十,然後若汗一邊將水倒在他頭上.
而耶穌的頭頂有一隻甲子,好像象徵聖靈降臨.
而在天上有天父的一團雲,一道光柱照在耶穌身上.
可能你也看過一些這樣的圖畫.
弗雷夫音只是描述歷史事件,沒有說明聖靈如甲子降下有什麼含義和作用.
施禎若汗做事後的陳述,當然在講故事的同時,加上他的神學解釋.
他說到聖靈的降下和留在,或本在住在耶穌身上.
是要作為證據,證明耶穌得著父上帝親自以聖靈施禮.
這要證明耶穌是上帝的兒子.
我們下星期有浸禮,大家也見過很多次浸禮.
在浸禮時,施禎的牧師也會一面說,我奉父,聖子,聖靈的名給你施禎.
你想像一下施禎若汗描述當時的情景是什麼.
父上帝將聖靈傾倒在耶穌身上,然後宣告你是我的愛子,我喜悅你.
天父上帝找聖靈來幫耶穌施禮.
耶穌一方面受水浸,另一方面受聖靈的洗禮.
施禎若汗用水幫耶穌行洗禮,但父上帝用聖靈去幫耶穌行洗禮.
換言之,耶穌同時接受水和聖靈兩重的洗禮.
施禎者分別是若汗和父上帝.
當然若汗用水當然是低級一點,父上帝用聖靈的洗禮是高級一點.
所以他所說的是雙重的洗禮.
舊約通常將耶和華的靈稱為高立,接受的稱為受高者.
賜下上帝的靈,即擁有上帝的身份,權柄和能力.
《少行傳》裡試圖介紹耶穌基督時.
在少行傳十章三十八節都這樣說.
上帝怎樣以聖靈和能力高拿撒勒人耶穌,這都是你們知道的.
理賽亞其實就是受高者的意思.
耶穌接受上帝所施以聖靈的洗,顯明祂就是理賽亞.
所以耶穌的受洗是若汗和父上帝共同完成的.

$^{241}$名副其實是神人合作.
不過施贈若汗沒有因此沾沾自喜.
覺得自己和上帝竟然平起平坐.
在耶穌受浸的事件上平分秋色,是不是這樣.
施贈若汗卻知道他自己所做的其實是微不足道.
上帝的行動才是關鍵性.
耶穌受人的浸禮,這是盡諸般的義.
上帝用聖靈幫他施洗,才證明他是上帝的兒子.
幫助人認識耶穌的神聖身份.
在我們日常的生活和事奉中.
我們幾乎時時都是神人合作.
我們做,上帝也在做.
我們的主學是神人合作.
我們的主學老師用心去預備,用心去教導.
上帝也在裝備我們每一個弟兄姊妹.
我們所做的一切都是神人合作.
不過我們所有事奉者都要像施贈若汗一樣.
看到自己所做的其實沒有什麼巴閉.
上帝的作為才是關鍵性,是不是.
就像保羅經常說的,唯有上帝叫生命成長.
有人栽中了,有人囂觀了.
不過後面那句才是重要的.
唯有上帝叫生命成長.
所以保羅怎麼說呢.
算不得什麼,什麼都算不得什麼.
做主席的算不得什麼,領事的算不得什麼.
唯有上帝能夠啟迪我們,幫助我們信生命成長.
所以每一個事奉者,一方面我們努力預備.
另一方面我們都提醒自己,我們算不得什麼.
上帝的作為才是關鍵性.
我們不跟上帝爭風頭,爭榮耀.
我們將所有榮耀都歸給他.
施禎一漢總結他的見證,34節.
我看見了,就證明這是上帝的兒子.
他兩番指出,先前他沒有耶穌神聖身份的認識.
31,33節,但在耶穌受浸的事件裡.
他就充分認識了.
上帝的兒子,這個不僅僅是身份,也是使命.
耶穌基督是那個The Chosen One.
上帝所揀選的那個.

$^{281}$因為約翰沒有跟上帝爭風頭.
沒有太當自己是一回事.
他就在事奉的過程裡,有新的啟迪,新的發現.
他認識到他之前的不認識.
努力調整和改變自己的想法.
幾個星期前,我看一集電視新聞的記錄.
講到政府為貧窮家庭的小朋友推行的私有計劃.
有一個專業人士擔當導師.
定期帶窮家的小朋友參與一些社會活動.
去海洋公園,開開他的眼界.
那個導師在接受訪問的時候說.
他在幫助小朋友開眼界的過程裡.
他自己其實也開了眼界.
有很多很多的學習.
弟兄姊妹,我們在事奉的過程裡.
是否也對上帝,對人,對自己有很多新的發現呢?.
我們是否覺得自己在進步中?.
只要你不覺得自己很有實力.
你不是真的覺得自己隨手行Default mode都可以做到.
你們唱了十幾年,幾十年.
閉上眼睛都可以唱.
如果我們是這樣,你不會學到什麼新的東西.
認定每一次都是一個新的機會.
然後你會有新的發現.
在事奉裡不停學習,不停發現.
我現在在幫助意大利那邊.
做一個我以為做了37年很熟悉的事.
但我要說其實很多驚嚇,也有很多驚喜.
每一次都是新的學習.
一個打亂了我既有的想法和做法的新學習.
這是很感動的事.
你可以確定,施禎若翰在講這個見證之前.
他早已經替耶穌受了浸.
所以耶穌受浸的時間.
是在29節所講的次日.
甚至是早過施禎若翰和耶路撒冷的調查團見面的時候.
再早一點.
不過,總之約翰福音的作者沒有記述耶穌受浸這個歷史事件.
只是引述施禎若翰的口述見證.
這是為了凸顯施禎若翰最主要的身份和使命.

$^{321}$我們一開始就說.
施禎若翰這個名字是我們給他取的.
不是很準確的.
其實他不是一個施浸者.
他其實是一個見證者.
John the Baptist不是一個最準確對他的描述.
是John the Witness.
這才是真正他的身份.
他不是施浸者.
他是見證者.
他最重要的責任不是幫助耶穌施浸.
幫助其他人施浸.
他最重要的責任是為耶穌基督做見證.
這是他存在的唯一目的.
施禎若翰做了兩個見證之後.
第二天,他和兩個門徒再次遇上耶穌.
他又做了一句話的見證.
35到37節.
再次日,約翰和兩個門徒站在那裡.
他見耶穌行走就說.
看!這是上帝的羔羊.
兩個門徒聽見他的話就跟從了耶穌.
看!上帝的羔羊.
其實這句話是重複他第一個見證.
我們剛才說了的第一個見證的內容.
沒有什麼新意.
不過從效果上說.
這句話就很重要.
因為施禎若翰是和他兩個門徒說的.
這導致他兩個門徒轉頭耶穌門下.
成為耶穌的第一批門徒.
然後我們看約翰福音.
你就會看到約翰福音將故事的重點.
不再放在施禎若翰身上.
轉移到耶穌和他的門徒身上.
這是我們接下來要說的故事.
兩個門徒聽到他們夫子的介紹.
立刻就明目張膽跳船.
離開施禎若翰,轉移跟隨耶穌.
這兩個門徒是誰呢?.

$^{361}$一個是聖經說的安德烈.
另一個是聖經沒有說的.
我們相信是二擇其一.
或者是寫約翰福音的約翰.
或者是肥力.
我們下一次會詳細解釋這個推論.
總之不是約翰就是肥力.
這兩個人如果一直是施禎若翰的門徒.
緊跟隨施禎若翰身邊.
聖經沒有說還有其他門徒跟隨他.
理論上這兩個門徒跟他們師父的關係應該比較親密.
他們為什麼可以跳槽就跳槽背叛師門.
這麼無情無義呢?.
事實上我們知道他們的跳槽是師父吩咐的.
所以不算是忤逆師父的心意.
沒有丟師父的面子.
他們是遵使命離開師父.
另投名主.
施禎若翰為什麼要吩咐這兩個門徒跳槽呢?.
當然是他認定耶穌一定希望自己一定須微.
這是三章三十節所說的.
所以那些門徒跟隨耶穌比跟隨自己更有前途.
他希望幫助兩個門徒去覓得更好的前途.
另外施禎若翰也可能猜到自己命不久矣.
現在吩咐門徒跟隨耶穌.
有些託辜的行動.
再說他之所以特別吩咐安德烈或約翰或肥力轉跟耶穌.
可能是因為這兩個人都比較年輕.
安德烈是很年輕的.
如果是約翰也是很年輕的.
可能這兩個人是施禎若翰眾多門徒最年輕的兩個.
年輕人有比較大的塑造空間.
更改行車道也比較容易.
所以施禎若翰特別指這兩個人跳槽.
換學校,我這間學校要倒閉了,請你早點換.
換另一間好一點的學校.
不過就算師父做了一個離奇的吩咐.
要他們另投明主.
安德烈為什麼他們會這麼順攤就範聽從呢?.
這是背叛師門的行動.

$^{401}$難道他們跟師父沒有培養出感情.
所以真的揮一揮衣袖,不帶著一片雲彩就走.
我們相信施禎若翰除了吩咐兩個門徒投靠耶穌之外.
應該還跟他們說了一些有關耶穌的秘密.
耶穌就是舊約所應許要來的彌賽亞.
然後才促使他們毅然作出轉換師門的決定.
這也解釋了安德烈和其後他跟他的兄長彼得說.
一章41節,我們預見彌賽亞了.
為什麼他知道呢?應該是施禎若翰跟他說的.
他自己沒有發現耶穌是彌賽亞.
耶穌也沒有說過他是彌賽亞.
不過這是施禎若翰跟安德烈說的.
如果施禎若翰告訴門徒耶穌就是彌賽亞.
門徒轉而跟隨耶穌就很合理了.
因為猶太人一直期盼就是彌賽亞的降臨.
他是整個民族的期望.
能夠遇到他,怎麼會錯過呢?.
就算師父很重要,不過彌賽亞比師父更重要.
我們中間假如你不是第二代第三代的基督徒.
你不是信二代,信三代,信四代.
如果你是第一代的基督徒,你跟我們都一樣.
我們都經歷過像安德烈那樣的割捨.
背叛師門的那種經歷.
因著信耶穌,我們離開了昔日慣常的生活方式.
離開了一直交往緊密的群體.
平時一起打麻將,喝酒的朋友.
我們捨棄了從前效忠跟隨的對象.
什麼密宗的大師.
童年的時候上了契的觀音大師等等.
甚至我們違背了我們的家族傳統.
抗拒了父母和家人,例如老公的旨意.
背棄了昔日曾經對偶像,親友做過的承諾.
這些一切都是我們在信耶穌的時候.
要即時繳付的社會代價.
信耶穌總是意味著生命的一個斷層.
我經常說耶穌基督不可能是我們生活的補充.
不是補充,永遠是一個顛覆.
一個subversion,一個顛覆.
不是補充,不是延伸.
所以我們一定是背叛師門,割斷過去.

$^{441}$轉頭耶穌重新開始.
我自己在天主教中學讀書.
我一直以來都對天主教有好感.
從小就叫基督教是世反派.
Protestant就是世反派,是叛徒.
我自小在一個艱苦的環境長大.
我自己對所有權威都很憎恨.
我一直認定靠自己,孤身闖路.
我是一個充滿敵意的世界成長.
仇恨是我驅動自己前行的最重要的動力.
這一大堆的事情,都因為信耶穌的緣故.
馬上就割斷,就放下.
我背叛了身邊的人.
最重要的是我背叛了我以前的自己.
背叛就是一個顛覆.
信耶穌是意味著這樣的一個割斷.
正如安德烈他們跟隨耶穌.
是意味著他們背叛師門.
和昔日劃清界線.
彼得曾經對耶穌說.
主啊,你有永生之道.
我們還歸從誰呢?.
但我們對跟隨耶穌都是一個堅決.
不顧一切,義無反顧,死心塌地.
37節講安德烈這兩個前師尊約翰的門徒跟從耶穌.
不過,這裡講跟從耶穌要比較寬鬆地解釋.
他們主動跟隨耶穌.
耶穌去哪裡,他們就跟到哪裡.
接近耶穌,一路聽耶穌的教訓.
不一定是指耶穌已經正式收納他們做門徒.
因為從這段經文前後看.
我們找不到耶穌在這個時候.
有發出過來跟隨我的呼召.
耶穌設立十二個使徒.
更加是在比較後的時間.
所以有人或許會發現.
約翰福音對耶穌呼召門徒的技術.
跟三卷符類福音在次序上其實很不同.
對於這些經卷之間的差異問題.
我們可以應答的解釋就是.

$^{481}$約翰只是記述各個門徒怎樣跟耶穌相遇.
被他吸引成為他的追隨者.
不一定是說耶穌已經對他們發出做門徒的呼召.
你知道跟隨者和成為門徒是不同的.
就是耶穌有沒有呼召他們.
這是最重要的關鍵.
我們看完整段聖經.
一個很簡單的總結.
我們不僅僅總結.
今天總結之前我們說兩次施浸約翰的見證.
我們用了三次時間去說施浸約翰的見證.
我相信我們對約翰的謙卑和忠誠.
應該有很深刻的印象.
在第一講.
我們從約翰的悔言中看到施浸約翰終身的任務.
就是為光,為耶穌做見證.
在第二講.
面對來自耶路撒冷的調查團查問約翰的身份.
約翰以他什麼都不是作為回應.
約翰不是自我形象低落,不是自覺一無是處.
而是當他跟要來的耶穌比較.
他的所有事都變成不是,不值一提.
在第三講.
今天使浸約翰介紹耶穌基督,直接為耶穌做見證.
他跟成年後的耶穌最親密的關係就是替他施浸.
而就在這個事件裡.
他發現了耶穌的神聖本質.
從而他甘願在這個被他年輕的親戚面前.
做他的開路先鋒,做他的見證者.
施浸約翰是見證者.
見證者不可以只講自己的經歷和感受.
否則他就是見證自己.
很多時候我們在教會聚會聽到的見證.
都是自我吹噓,自我宣傳,多於講耶穌.
要凸顯自己的非凡經歷,凸顯自己的虔誠,自己的識見.
真正的見證者一定專以為耶穌做見證.
以耶穌作為說明的對象.
讓聽的人聽到耶穌,認識耶穌,而不是認識見證者.
我們在約翰幾段的見證裡面.
聽到的訊息都是關於耶穌.

$^{521}$約翰的訊息從來沒有失去焦點.
約翰一生的事報沒有迷失方向.
皇上弟兄姊妹,我們講的都是很簡單的事.
我們每個人都是施浸約翰.
只要我們發現自己是施浸約翰.
我們就很有效可以扮演施浸約翰的角色.
我們什麼都不是,也都不重要.
約翰的講法講得很絕.
連幫耶穌玩馬的資格都沒有.
他把自己壓到這麼低.
我幫耶穌玩馬的資格都沒有.
我們僅僅是開路先鋒,是見證者.
讓所有踏足這間教會的人.
都看到耶穌基督是真光.
也都讓他們的生命得到耶穌的光照.
弟兄姊妹,是真的.
我們一起立志.
只要我們真誠忠心扮演施浸約翰的角色.
我們就會在這裡,在我們的教會裡面.
我們會看到天開了,聖靈將甲子降下.
天父上帝會在這裡親自開聲宣告.
耶穌基督是他的愛子.
耶穌基督會在這裡得到最大的榮耀.
而我們每一個來到我們中間的人都會看到.
耶穌基督是他們生命的出路.
耶穌基督是他們唯一的拯救.
上帝閱立我們在祂面前的立志.
\newpage



\section{以斯拉記 1:1-11}
\label{sec:ZKaifzpEjak}
\textbf{2022年8月21日崇拜直播｜陳明泉牧師｜並肩勇闖系列:願神與這人同在｜以斯拉記1\_1-11}
\newline
\newline
連結: \href{https://youtube.com/watch?v=ZKaifzpEjak}{\texttt{https://youtube.com/watch?v=ZKaifzpEjak}} ~~~~ 語音日期: 2022-08-22
\newline
\newline
\hyperref[sec:IiBVTMF9Vak]{\small{< < < PREV SERMON < < <}}
~
\hyperref[sec:index_chronic]{\small{[返順時目]}}
~
\hyperref[sec:hASHSP9n8fQ]{\small{> > > NEXT SERMON > > >}}
\newline
\newline
以斯拉記 1:1-11
\newline
\begin{longtable}{cl}
\hline
\hline
章節 & 經文 (和合本修訂版)\\
\hline
1:1 & \begin{tabularx}{0.7\textwidth}{X} 波斯王居魯士元年，耶和華為要應驗藉耶利米的口所說的話，就激發波斯王居魯士的心，使他下詔書通告全國，說： \end{tabularx} \\ \\ \relax
1:2 & \begin{tabularx}{0.7\textwidth}{X} 「波斯王居魯士如此說：耶和華天上的神已將地上萬國賜給我，又委派我在猶大的耶路撒冷為他建造殿宇。 \end{tabularx} \\ \\ \relax
1:3 & \begin{tabularx}{0.7\textwidth}{X} 你們中間凡作他子民的，可以上猶大的耶路撒冷去，重建耶和華－以色列神的殿，他是在耶路撒冷的神；願神與這人同在。 \end{tabularx} \\ \\ \relax
1:4 & \begin{tabularx}{0.7\textwidth}{X} 凡存留的人，無論寄居何處，那地的人要用金銀、財物、牲畜幫助他，還要為耶路撒冷神的殿甘心獻上禮物。」 \end{tabularx} \\ \\ \relax
1:5 & \begin{tabularx}{0.7\textwidth}{X} 於是，猶大和便雅憫的族長、祭司、利未人，凡是心被神感動的人都起來，要上耶路撒冷去建造耶和華的殿。 \end{tabularx} \\ \\ \relax
1:6 & \begin{tabularx}{0.7\textwidth}{X} 四圍所有的人都拿銀器、金子、財物、牲畜、珍寶支持他們，此外還有甘心獻的一切禮物。 \end{tabularx} \\ \\ \relax
1:7 & \begin{tabularx}{0.7\textwidth}{X} 居魯士王也把耶和華殿的器皿拿出來，這些器皿是尼布甲尼撒從耶路撒冷掠取，放在自己神明廟中的。 \end{tabularx} \\ \\ \relax
1:8 & \begin{tabularx}{0.7\textwidth}{X} 波斯王居魯士派米提利達司庫把這些器皿拿出來，點交給猶大的領袖設巴薩。 \end{tabularx} \\ \\ \relax
1:9 & \begin{tabularx}{0.7\textwidth}{X} 它們的數目如下：金盤三十個，銀盤一千個，刀二十九把， \end{tabularx} \\ \\ \relax
1:10 & \begin{tabularx}{0.7\textwidth}{X} 金碗三十個，備用銀碗四百一十個，其他器皿一千件。 \end{tabularx} \\ \\ \relax
1:11 & \begin{tabularx}{0.7\textwidth}{X} 金銀器皿共有五千四百件。被擄的人從巴比倫上耶路撒冷的時候，設巴薩把這一切都帶了上來。 \end{tabularx} \\ \\
[1ex]
\hline
\hline
\end{longtable}
$^{1}$今日重讀的經文是記載在.
《以斯拉記》第一章一至十一節.
在舊約聖經的五百七十七頁.
《以斯拉記》一章一節.
波斯王古列元年.
耶和華為要應驗藉耶利米口所說的話.
就激動波斯王古列的心.
使他下詔通告全國說.
波斯王古列如此說.
耶和華天上的神.
已將天下萬國賜給我.
又祝福我在猶大的耶路撒冷.
為他建造殿宇.
在你們中間凡作他子民的.
可以上猶大的耶路撒冷.
在耶路撒冷重建耶和華以色列神的殿.
只有他是神.
願神與這人同在.
凡剩下的人.
無論寄居何處.
那地的人要用金銀財物.
生畜幫助他.
另外也要為耶路撒冷神的殿.
甘心獻上禮物.
於是猶大和辯雅敏的族長.
祭司利未人.
就是一切被神激動他心的人.
都起來要上耶路撒冷去建造耶和華的殿.
他們四處的人就拿銀器.
金紙,財物,生畜,珍寶幫助他們.
另外還有甘心獻的禮物.
古列王也將耶和華殿的器皿拿出來.
這器皿是尼布甲尼撒從耶路撒冷掠來.
放在自己神之廟中的.
波斯王古列派副官米提利達.
將這器皿拿出來.
按數交給猶大的首領切巴薩.
器皿的數目記在下面.
金盤三十個.
銀盤一千個.

$^{41}$刀二十九把.
金碗三十個.
銀碗之次的四百一十個.
別樣的器皿一千件.
金銀器皿共有五千四百件.
被擄的人從巴比倫上耶路撒冷的時候.
切巴薩將這一切都帶上來.
請各位請坐.
弟兄姊妹平安.
今天已經是我們說的並肩用廠系列最後一講.
不知道大家去聆聽並肩用廠這個系列的時候.
你心裡有沒有一些很深刻的體會.
我們由創世記一直講到整個以色列.
到今天講到他被擄歸回的一段很長的舊約的一些技術.
在整個以色列人.
整個以色列民族的歷史裡面.
我們看到上帝怎樣去對待這個民族.
這個民族其實離不開這四個歷史進程.
第一個進程就是講人.
由一個人的一個家庭亞伯拉罕.
一直延展到十二支派成為一個民族.
在舊約聖經裡面就講著整個以色列就是這樣建立起來.
當有了人之後下一步就是要有地方.
所以由本來自己流散的一些遊牧民族.
慢慢一路一路就進入流浪軟墨之地.
後來就成為了一個有自己圍牆自己領土的一個國家.
領土Land.
這個是第二個的歷史進程.
當有了地方之後.
第三樣東西他們就要問.
怎樣去管治.
怎樣去延續整個的民族一直維繫下去.
接著就講King皇帝.
大家覺得皇帝是一個很舊的政治思想.
在當代來說是最流行的.
即是每個外族外邦.
每個都有自己的王管治.
以色列民族求上帝給一個王.
好像外邦一樣來管治他們.
令到他們國家真的進入一個體系來建立起來.

$^{81}$所以第三個歷史進程就是King.
最後一個某程度也是整個以色列民族裡面最大的高峰.
聖殿.
舊時昔日的敬拜由回眸.
以色列人去到哪裡.
那個回眸就去到哪裡.
或者回眸引領他們去到不同的地方來敬拜.
但當建立了一個王國.
建立了自己的疆土.
有自己的王的時候.
大衛就求上帝建立一個聖殿.
所以上帝就應許在所羅門當中.
由所羅門開始建立一個華麗的宮殿.
讓上帝.
他們的神學就是讓上帝住在那裡.
聽他的百姓來祈禱.
不過更加重要的是這個聖殿.
其實不是說上帝住在那裡.
而是代表著這個地方.
這個參照著回眸.
也都參照著天上敬拜的這個場所.
告訴以色列人.
神與他們同在.
這個聖殿是說整個以色列文化.
或者整個以色列敬拜上帝的一個歷史的高峰.
這個是最輝煌的時間.
不過以色列人當他們知道有自己的聖殿.
有自己的地方.
他們卻是跟隨了外邦來去拜別的神靈.
拜偶像或者貪圖一些方便,財富.
甚至乎他們覺得外邦的東西比耶和華上帝更好.
所以就離棄了耶和華上帝.
就拜了外邦的神.
然後上帝就進行管教.
直至到整個以色列.
整個猶大民族亡國.
剛才記不記得那四個進程.
People, Land, King, Temple.
即是人,地方,王和聖殿.
但是隨著被擄之後.

$^{121}$聖殿沒有了.
城牆失去了.
沒有了自己的國土.
王被廢棄.
西底加王被挖去眼睛.
被擄巴比倫.
很慘烈的故事.
剩下的就是這一班以色列民.
這一班人經歷過高高低低.
經歷過生命.
整個民族的高潮.
但是因為他們離棄上帝.
所以被上帝管教.
所以他們只剩下這一班人.
其實在歷史的舞台裡面.
不同文化,不同種族的人.
都有這樣出現過.
不過一個民族可以在一段很短時間.
就淹沒在整個歷史舞台裡面.
但是上帝竟然用這一班.
剩下來的以色列民.
用這一班沒有了王.
沒有了國土.
沒有了敬拜聖殿的一班人.
怎樣去維繫呢?.
怎樣去繼續他們未來的信仰生活呢?.
今天我們讀的.
這個系列裡面最後一段經文.
以斯拉技為我們提供答案.
上帝有沒有離棄這一班人呢?.
還是上帝要在這一班人裡面.
另外有一些心意.
藉著他們的經歷和故事.
勉勵我們今天一班在地上.
信了耶穌的基督徒.
很多人以為舊約還舊約.
跟我們沒有關係.
不過其實上帝是啟示的一本聖經.
藉著以色列人.
也都勉勵我們今天在地上.

$^{161}$信了耶穌的一班基督徒.
歷史是一貫的下來的.
我希望今天我們去看的時候.
作為這個系列的總結.
也都讓我們更加知道.
我們今天基督徒怎樣看昔日.
他們這班以色列人怎樣走.
我們按著上帝的心意.
繼續走我們人生的路.
我們一起祈禱.
求上帝幫助我們同心祈禱.
生下主願你在我們當中.
你帶領我們.
讓我們能夠明白你的話語.
讓我們從以斯拉技裡面.
看到上帝的心腸.
看到上帝的心意.
讓我們能夠讀你的話語的時候.
我們的心被激勵.
今天我們每一個弟兄姊妹.
聆聽的過程.
讓我們不單止明白.
我們更加求你.
你對我們的內心說話.
激勵我們.
讓我們奮勇地走在面前的路.
幫助孫講的講得清楚.
忠心地將你的話語.
向弟兄姊妹展明.
恭敬在我們時間交托.
願你賜福.
獻上我們的禱告.
奉靠基督耶穌.
得勝的聖名.
Amen.
我們一起看.
我做了一個表給大家.
這個表在大家的程序表.
或者你看powerpoint都看到.
其實很多弟兄姊妹.

$^{201}$有時候會亂.
就是說.
以色列人被擄.
那些年份是怎樣呢.
那些年份有一兩年的出入.
不奇怪的.
因為不同的參考資料.
有不同的.
我快點告訴大家.
究竟那件歷史事實是怎樣.
剛才說到以色列人.
就亡國.
在公元586年.
南國.
就是直接是說整個以色列民族.
剩下來的猶大.
被巴比倫帝國所滅.
耶路撒冷淪陷.
這個就歷代至.
三十六章十七至二十節.
沒有了.
就是所有國家就.
整個以色列就.
被一個強大的巴比倫帝國所吞滅.
那皇被擄走了.
西帝加被挖去眼睛.
他想反抗.
因為反抗不了.
於是就被苦害.
精壯的一些年輕的.
有生產價值的.
全部就帶到巴比倫那裡.
來做亡國奴.
當然在巴比倫裡面.
他都會禮遇一些有能力的人.
不過整體來說.
國家沒有了.
身份沒有了.
到了五百三十八年公元前.
巴比倫這個這麼強大的帝國.

$^{241}$都覆滅.
巴比倫昔日的管治.
就是用強權用政治.
來管壓不同.
幅員很大的整個.
巴勒斯坦中東等等的地區.
甚至乎有部分的歐洲地方都.
被巴比倫所吞滅.
不過幅員太大.
如果經常都用強權壓迫.
無以為計的.
所以五百三十八年.
就被另外一個新的.
新一代的帝國.
馬代波斯.
或者叫波斯國.
聖經裡面說的.
就所取代.
第一任.
或者他們新的波斯國.
那個領袖就是古列王.
有一些聖經有.
易賽老士或者居老士.
同一個人來的.
譯音不同的.
古列王.
為了方便我之後都用古列這個字.
古列王或者整個馬代波斯.
取代了巴比倫之後.
他和巴比倫的做法不同.
他要讓各個民族.
來做一些自行管治的功夫.
所以在古列王一登基.
古列王元年.
他就說.
叫猶太人可以回去重建耶路撒冷.
這個是公元五百三十八年.
公元前五百三十八年.
好了.
但是一回去.

$^{281}$是不是馬上扔掉聖殿建好呢.
不是.
其實有一些妖樣的時間.
為什麼呢.
因為他說建.
但是不是等著沒有反對的.
有一些的外邦.
特別是說.
以前說有不同集居.
有以色列.
又有其他民族血統.
在亞述帝國之下做的政策.
撒瑪利亞人.
就大力地來阻止聖殿的興建.
所以公元前五百三十八至五百二十.
這十幾年裡面.
就停滯不前.
沒有辦法.
很大規模地將聖殿復工.
繞一繞十幾二十年.
到真的開始興建.
就是公元前五百二十年.
然後正式修建.
搞定了.
公元前五百一十六年.
你看到在以色列記裡面.
其實按著這些年份.
我寫了相應的內容給大家.
大家知道了.
回去慢慢自己看看.
到了以斯拉.
以斯拉記的作者.
他不是書卷裡面.
最重要很早出現的人.
去到中段.
去到第七章才出現.
因為在說四百五十八年公元前.
那時候才回到耶路撒冷.
以斯拉回去有另外一個任務.
就是興建一個硬件的聖殿.

$^{321}$他要回去重新展明律法的意思.
所以四百五十八年.
就是以斯拉回耶路撒冷.
到了公元前四百五十五年.
到四百四十三年.
就是尼希米.
就是另外一個皇身邊的酒正.
回去重建.
剛才說的是聖殿.
現在這個是耶路撒冷的城牆.
耶路撒冷是一個城.
聖殿是一個建築物.
是復元闊一點.
來收足城牆.
尼希米某程度是最後階段.
來將這個城牆重建.
整個某程度歷史被毀的聖殿.
到重建歷史很長的.
這幾十年.
甚至差不多一百年的.
這一段的歷史.
就是我們今天主要去讀的.
待會我們都會參照去看.
不過有些地方想大家再留意.
剛才說到整個以色列民族.
他被外邦.
被巴比倫吞併.
在歷代次那裡.
歷代次第36章22至23節.
有記載著.
有記載著的歷代志下.
他記載著怎樣呢.
就是說到整個國家已經被滅了.
接著他用的字眼就是.
這個就應驗了.
利未口所說的話.
地要享受安息.
因為地土芳良.
便守安息.
寂滿了七十年.

$^{361}$他這句話說什麼意思呢.
歷代志下那裡說.
整個民族被擄之後.
那個地方.
本身他們居住的那個地方.
就變成了一個荒土.
但這個荒土是守安息的.
那個年期是七十年.
順帝是在說.
懲戒這些的以色列民.
但同時間也好像重新重設.
重整整個以色列民族.
是一個怎樣的民族.
接著在經文裡面.
他後面加了一個註腳.
就是說到寂滿了七十年之後.
波斯王元年.
耶和華為要應驗.
耶利米口所說的話.
就激動波斯王古烈的心.
使他下詔書通告全國說.
波斯王古烈如此說.
耶和華天上的神.
已將天下萬物賜給我.
又祝福我在猶大的耶路撒冷.
為他建造殿宇.
你們中間凡作他的子民.
可以上去.
願耶和華他的神與他同在.
這個是歷代志下三十六章.
最後的一句.
也是以色列的第一句.
剛才大家去聽的經文.
是延續著歷代志所寫的內容.
既是一個章節的結尾.
也都是代表著新的章節的開始.
在聖經的作者裡面.
特別是以色列的講述.
以色列民不是因為被擄.
他的歷史就終結了.

$^{401}$沒有了.
而是一個新的一頁正式開始.
這一頁就是講到.
首先由古烈王開始.
來感動他們.
感動古烈王來激動他.
令他頒下一個詔書.
重建耶路撒冷的聖殿.
在這裡有一點點歷史資料.
很多歷史學家.
或者有些學者問.
為什麼無端端古烈這麼有心.
叫一些以色列民.
回去建造聖殿.
是不是他信了耶穌.
信了上帝 耶和華上帝.
然後他就將波斯帝國的國教.
就將猶太教成為了.
舊約聖經就成為他們整個國教呢.
一定不是的.
因為這個有些出土的文獻說.
就是波斯王出的這個詔書.
原來除了給在以色列的民族之外.
同樣的格式.
原來都有給其他.
不同的民族的幫國.
這個某程度是一個政策.
這個政策就是講到.
那個波斯帝國因為福源很大.
他沒有辦法一個國家.
一個波斯人來管治.
所有不同福源這麼大的民族.
所以就給他們一些空間.
去自己自治.
而自治的方式.
主要運用的就是用宗教的方式.
讓這一班人在宗教的規範之下.
來管理好他們自己的生活.
所以不是古烈王信了耶和華上帝.
然後就將他頒佈成為國教.

$^{441}$不是這樣.
而是因為在一個國家裡面.
有不同的政策.
所以他用這個比較明智一點.
一個懷柔的手段.
來幫助不同的民族.
特別是坐這個順風車.
幫這個猶太重新去建立.
他們自己的聖殿敬拜.
這個你明白了.
這個本身是一個政策.
不過在以斯拉記裡面.
他怎麼寫呢.
以斯拉記裡面.
他寫的內容.
他不是純粹講政治考慮.
他不是從政治角度來評論.
古烈王的這個是不是德政.
在經文裡面劈頭.
以斯拉就是說.
這個是應驗應許.
而你不是在他作為先知的過程.
說70年成為荒土.
現在從70年差不多到了.
他就回去.
而這個歷史的進程.
就好像在一個人為的因素底下.
上帝就將這個事件來到應許出來.
除了在耶利米之外.
其實在以塞瓦書也有這樣講.
以塞瓦書還講了古烈這個名字.
以塞瓦書第45章第一節和13節.
他那裡這樣寫的.
我耶和華所高的古烈.
我參乎他的右手.
鮮烈國狂伏在他面前.
要放鬆烈王的腰帶.
使城門在他面前敞開.
不得關閉.
13節.

$^{481}$我憑公義興起古烈.
又要修直他一切的道路.
他必建造我的城.
釋放我被擄的民.
不是為功架也不是為賞賜.
這是萬君之耶和華說的.
在以色列的先知著作裡面.
其實已經講了.
這件事已經是這樣了.
還要講了古烈這個名字出來.
不過古烈當然不知道.
他不是看舊約聖經的.
不過整件事就是.
一個這樣的進程下.
來到上帝引領新一代的.
馬代波斯的古烈王.
有些人說是梟雄.
不過就是這樣被激動.
以致整個聖殿重建就這樣開始了.
在這裡稍微給我們一些地方值得留意.
古烈的心為什麼會有這個決定呢.
是上帝激動他.
激動這個字稍後會再出現.
是一個很特別的字眼.
今天我們講信制度.
但昔日來說.
如果講王是講那個領袖.
一個怎樣的領袖.
就決定了整個國家的命運.
那個領袖怎樣想呢.
就決定了他做些什麼的政策.
上帝是一個掌管歷史.
掌管不同人的心的上帝.
他不是純粹在我們當中.
一些很細微細眼的生活裡面.
輕輕鬱鬱的那種上帝.
上帝是可以超然去到.
人類整個列邦的歷史裡面.
不同君王都是在上帝的掌控當中.
我們說上帝是一個滿有權柄的上帝.

$^{521}$他的能力範圍.
不是我們在我們一群基督徒.
或者信的人當中.
整個世界還沒信的那些君王.
所有的民族.
都是在上帝的掌管下.
今天我們不妨.
雖然我們有時候.
我們都是這樣凝信的.
不過在我們的生活當中.
我們經常覺得.
上帝只是影響我們.
上帝只是影響著那群基督徒一小撮人.
聖經裡面不是這樣說的.
上帝是影響全世界.
掌管萬有的權力.
當大家都不知道.
為什麼會這樣發生的時候.
上帝正在做事.
激動君王激動領袖的心.
用他的大能大力.
來改變一些人為不可能發生的事情.
竟然上帝就在古列的心裡面.
激動.
所以上帝是激動獵王的心的王.
在萬王當中還有一個更高的王.
那個就是我們所信的上帝.
在經文裡面再繼續說.
古列他的王位.
這個很有趣.
他在說古列的王位是怎麼來的呢?.
不是軍事.
或者他承接著.
所有自己很多優越的條件.
在以斯拉記裡面.
或者歷代志下.
那個王位是上帝高納他的.
在耶和華高納他.
我們常常想著耶和華高納大衛.
高納不同歷代的彌賽亞.

$^{561}$我們覺得那條線才是上帝高納.
不過在馬代波斯這個外邦的王.
上帝都是用高納這些.
那個高納其實是賜給他有權力.
他的權力都是來自上帝.
今天如果我們這樣讀聖經.
當我們這樣明白的時候.
我們不要將這個世界.
有時候有一些未信主的想法.
或者純粹是看到不是基督徒.
那些東西所有都是邪惡.
或者所有東西不是上帝所用.
我們要拉闊一點.
上帝原來在這個世界裡面.
未信主的人.
可能都是成為上帝要使用的器皿.
上帝使用的瀑人.
上帝藉著這些人來幫助我們.
這個落在一個沒那麼轟烈的層面.
一個生活一點的見證.
我認識一個還不是基督徒的弟兄.
你聽不聽到這個意思.
還不是基督徒的弟兄.
他還沒有決志.
還沒有真的回教會.
不過他經常幫我忙.
幫我忙傳福音.
人們不開心的時候.
你要找耶穌.
他做到一件事.
就是為他帶的家人祈禱.
他帶家人來信主.
他會配合我的工作.
但他自己還沒有信.
但他就藉著他的努力.
和大家配搭.
他給我有機會去到他家.
他身邊的人那裡.
來帶那些人來決志.
我也覺得有時很奇怪.

$^{601}$我們覺得基督徒才可以傳到福音.
是的.
基本上是正路的想法.
不過有時.
有些徘徊在信與疑之間的朋友.
一些未信的弟兄.
將來的基督徒.
原來上帝都在用他們.
來為他的國度慢慢去拓展出去.
如果在宣教裡面更加見到.
有一些很難基督徒.
在教會裡面進入到的部落群體.
但有一些族長.
見到一班基督徒做好見證.
他們會鋪橋搭路.
幫助基督徒進入最有需要的人當中.
他可能自己還沒有信.
不過他見到這班人是好.
他願意成為中間的橋樑.
上帝感動他們.
讓他們可以拓展上帝的國度.
所以在這裡.
第一個讓我們知道.
上帝甚至未信者.
我們覺得不是他血緣的那些以色列人.
波斯人上帝都會使用.
更加重要是為何上帝要使用他們.
因為這個其實回應那個70年的應許.
耶利米蘇之前說到.
以色列民被擄70年.
在整個70年的日子裡面.
面對著亡國的痛苦.
一班已經是亡國之後.
已經不是頭等公民.
是很低下的一班公民.
即或他們有些聰明才智.
為他們累積到一些財富.
又或者他們在官場裡面有一官半職.
但是以他們的實力根本沒有辦法.
將整個民族來興起.

$^{641}$但是上帝卻是應許過這班以色列民.
一班曾經犯罪離棄上帝的人.
上帝為他們重新建立一個群體.
藉著外邦的人.
來重新建立耶路撒冷的聖殿.
70年如果你計算一下.
其實在猶太國被滅586年.
到古列王下令猶太人可以重建聖殿.
這個區間其實不足70年.
你減一減數字就知道了.
48年而已.
那70年的應許是不是即刻應驗呢?.
不是,漸進式.
這48年代表著上帝沒有忘記整個以色列民.
不過聖殿真的重新興建.
其實就是加上那段十多年的繞繞繞繞.
那些撒瑪利亞人的反對.
有這些節節節節.
加起來才等於70年.
上帝輕手了.
原本在巴比倫手下可以捱足苦70年.
但是上帝打了一個七折給他.
接著在這些繞繞繞繞當中.
這一班以色列民帶著盼望.
帶著很多逆流而上.
但是他們繼續重新興建上帝所命定的聖殿.
當古列王開了第一步之後.
整個以色列那些叫如下的人.
舊時被擄的.
又或者有些人散居了不同地方的人.
48年了.
兩代了.
按句落葉了.
settle down.
在那些地方裡面.
還要拔一根筋回來.
今天很多弟兄姊妹移民了回流.
都不會想這些.
不過在經文裡面怎樣說呢.
上帝不是感動或者激動古列王的心.

$^{681}$在經文裡面我們再看.
在第四至第五節.
以色列記.
凡剩下的人無論寄居何處.
那地的人要用金銀財物生畜幫助他.
另外也要為耶路撒冷神的殿金心獻上禮物.
我們再看前一點.
在第一章裡面第三節.
在你們中間凡作他子民的.
可以上猶大的耶路撒冷.
在耶路撒冷重建耶和華.
以色列神的殿.
只有他是神.
願神與這人同在.
經文為何要這樣說古列.
要頒佈一個這樣的命令呢.
剛才說過是上帝感動他的心.
同時間在這裡面更加說到.
一個很重要的觀念就是.
古列都希望促成那些百姓去做這件事.
他祈求的在他的皇令當中.
他說到如果誰去活耶路撒冷.
就願上帝與這個人同在.
願上帝的恩典臨到這班人的心上.
然後在第五節.
於是猶大和汴雅敏的族長.
祭司利未人.
就是一切被神激動他心的人.
都起來要上耶路撒冷去建造耶和華的殿.
所以在這裡面你會看到.
這班人為何會上去呢?.
除了激動古列的心.
亦激動了後面的數字統計.
四萬多人的心.
同樣被上帝激動.
有時我們一個人被激動.
心很容易的.
兩個人的難度就高一些.
即是一千人就難一些.
我在旺鎮港島幾百人.

$^{721}$幾百人都被激動很難.
我都不覺得自己可以.
因為不是我不是人的人為.
而是上帝激動了.
在散居兩代的一班子民.
一班雨下的一班人.
重新讓他們的心可以回到耶路撒冷.
在二章六十八節裡面.
說到激動他們的心.
他們怎樣去回應呢?.
有些族長來到耶路撒冷和耶和華殿的地方.
便為神的殿.
甘心獻上禮物.
要重新建造.
平平無奇一句.
二章六十八節.
激動回應就是獻上禮物.
去上到那班領袖.
即是那班開荒牛的.
重新開荒的.
去到耶和華殿獻上禮物.
不過獻上禮物.
這是和合本裡面的翻譯.
我們覺得好像一個動詞.
又不是錯的.
不過其實在很多聖經學者.
或者一直看回來還沒記.
原來那個甘心獻上禮物.
不是純粹說那個心.
或者大家做一些舉動.
其實更加好的翻譯是說獻上甘心祭.
甘心祭.
甘心祭其實在利未記裡面.
其實有記載著.
在利未記裡面說到.
那些以色列人.
興建或者祭祀.
敬拜上帝除了獻上.
沒有殘疾的生畜.
或者用農作物沒有發酵的.

$^{761}$來獻上平安祭之外.
原來有一個.
level低一點的.
獻的禮物.
不是最漂亮的.
可能有一點殘疾.
可能那隻牛有疤痕.
在屁股那裡.
以前如果用煩祭就不能獻了.
但那隻牛整體都能吃.
都挺好的.
用這些來擺去另外一個層面.
就是用這個來獻上甘心祭.
甘心祭是不是上帝完全閱立呢?.
經文不會特別去說的.
不過這個祭是強調的.
就是那些百姓獻上的.
未必是他們能力裡面最好的.
人之常情的.
被擄的一班人.
你有多可以呈現上好的牛羊.
上好的東西給上帝呢?.
他們不是那個處境的.
不過在他們眼中.
他們能夠能力範圍.
他們覺得好的東西.
他們想回應上帝的一些禮物.
他們可以獻上的.
在這裡裡面.
整個聖殿重新重建起來.
是由一班昔日的.
一班回來的.
一班的首領.
一班漁民.
從他們能力當中.
獻上甘心祭來開始.
所以在《以色列記》裡面.
說的那個獻祭.
這個聖殿獻祭.
除了說繁祭那些最禮儀性.

$^{801}$最標誌性.
最重要的那些之外.
多次兩三次這樣重複.
這班人是獻上這些甘心祭的禮物.
在他們的能力盡了.
他們沒有得再背了.
還可以想的.
在一個不是最完美的祭物當中.
他們就將他們的心.
甘心樂意地.
來呈獻給上帝.
重新建立在聖殿之中.
今天這個甘心祭對我們很有意思.
我知道今天.
雖然香港整體還是富裕.
但我們更加整體明白.
其實每個弟兄姊妹.
都有不同的難處.
我們在生活.
在很繃緊的一種壓力之下.
我們未必可以做得到.
好像早幾代.
或者有些很well off.
很好弟兄姊妹.
可以將品質最好的那些東西.
來獻上.
不過我們卻是可以.
有些不是最完美.
但令我們可以.
用我們勉勵自己.
盡自己的力量來獻上.
祭物不是最完美.
但那個心是我們甘心樂意獻上的.
我們就將這一切.
擺上給上帝使用.
今天我們要杜絕一種想法.
各位弟兄姊妹.
今天坊間有時在說.
有一種完美主義的事奉觀念.
什麼意思呢?.

$^{841}$上帝是聖潔.
完美的.
舊時舊約是這樣建制的.
所以大家做的事.
每一樣東西都要.
基本上是要完美的.
來擺上給上帝使用.
如果不完美又如何?.
就踢他下台.
不給他.
我們將那個事奉門檻設到很高.
基本上只有我幾個人.
才可以做得到.
所以那個才為之完美的祭.
不過至少在以斯拉記裡面.
那個重現聖殿的獻上.
都不是用一個極高的門檻.
來叫弟妹動員他們來服侍.
反而是甘心樂意.
不是很理想的.
不要緊.
我們就一起將我們能力.
將我們甘心樂意擺上來獻上.
當我們能夠接納.
或者我們鼓勵弟妹獻上甘心祭.
這種甘心祭有一個效果.
就是快樂.
你有份.
你不是要去到最頂端.
你不是要最高層次.
最美的那些東西.
雞雞地的那些東西.
但是你有份.
容許我說.
施班員是不是每個聲音都美到最嘹亮呢?.
不是.
我有些沙啞.
我有些音唱不到.
但是我用我這個喪志.
我願意擺上.

$^{881}$被上帝來使用.
甘心祭.
我未必像一些大有錢人那樣.
可以奉獻很多的錢.
我只是一個寡婦的窮小.
幾個的小錢.
但那個是我養生.
我很努力省回來.
甘心祭.
我們今天教會需要更加多的甘心祭.
不用和別人來做一些比較.
他做得高.
做得低.
他做得美一點.
或者做得雞雞地.
不要緊.
甘心獻上.
上帝就會收納.
在《春海及記》其實是在說回眸的開始.
上帝都是這樣吩咐摩西的.
他說凡甘心樂意獻上禮物的.
你們可以收下.
給上帝.
歸給上帝用.
上帝不是要我們最美的東西.
反而他想要我們甘心樂意.
我們願意擺上自己.
在我們能力裡面最好的.
來獻上.
這個就是甘心祭.
聖殿重建.
由甘心祭作為一個起點.
今日教會的生活.
我們怎樣去重建呢?.
由甘心祭開始.
由你覺得眼前可以擺上的一些能力.
不是最完美的東西.
但你走出那一步.
不是最理想.
但你願意將你不是最理想的東西.

$^{921}$分給大家.
一齊努力.
這個神的家.
你會經歷.
我有份.
我屬於這裡.
我願意為這個群體.
為這個家擺上努力.
除了說甘心祭的獻上.
這個是以色列的頭半部.
一至第六章.
剛才說的硬件的興起.
接著.
以色列的未完的.
在第七章開始.
說到以色列回到七至十章.
說的是以色列的第二個部分.
這個群體重新興建一個硬件的聖殿之後.
他們繼續要做些什麼努力呢.
在第七章裡面就開始說了.
或者我先看看我們第七章第25節.
我沒有印給大家.
我讀給大家聽.
七章25至26節.
其中一個節錄的預令.
以斯拉亞.
你要照著神.
照著你神賜你的智慧.
將你所明白你神律法的人.
納為事司.
審判官.
治理河西的百姓.
使他們教訓一切不明白神律法的人.
凡不遵行你神律法.
和王命令的人.
就當束束定他的罪.
或致死.
或充軍.
或抄家.
或囚禁.

$^{961}$這裡是王令.
是吧.
就是說.
叫以斯拉.
要教導人們.
來聽王的命令和律法.
但是你留意一下.
頒佈這個命令的是什麼人呢.
是接下來.
波斯的.
下來的君主.
阿達舍西王.
一個外邦的王.
然後就叫以斯拉.
這個猶大血統的.
就是祭司.
文士.
在王身邊的一個重要的人物.
來將律法教給以色列人.
提醒他們.
叫他們好好地.
在律法當中.
來行事為人.
否則就有些惡果.
有致死,充軍,抄家,囚禁等等.
不過在這裡.
以斯拉同樣都用同樣的不足.
或者同樣的角度去講.
聖殿由古列王激動的心而起.
而管理以色列的百姓.
這個都是由外邦的阿哈徐老王開始.
來做一個很大的鼓勵.
叫以斯拉用律法來教導以色列人.
所以在第七章開始.
就是在講這個律法.
怎樣去重新建立.
如果講這個聖殿是一個精神文明.
聖殿的興建一至六章.
是在講那個外在.
裡面那種敬畏上帝的精神.

$^{1001}$就要在接下來表達出來.
因為時間關係我沒辦法仔細去講.
我想抽少少給大家知道.
在第九章第一至第二節那裡.
以斯拉認罪.
為甚麼認罪呢?.
「這事做完了眾首領來見我說.
以色列民和祭司並利未人.
沒有離棄迦南人,黑人,比利洗人.
耶布斯人,阿滿人,摩訶人,埃及人,阿摩利人.
仍效法這些國的民行可憎的事.
因他們為自己和兒子取了這些外邦女子為妻.
以致聖潔的族類和這些國民侵襲,混雜.
而且首領和官長在這事上為罪魁」.
以斯拉在後半段做了甚麼呢?.
那班之前興建聖殿建築物的那班領袖,祭司.
即是這班有身份的人.
他們做了一個硬件.
但是竟然在他們的行為,生活當中犯了可憎的事.
離棄了上帝.
甚麼事呢?.
他的子女或者他自己取那些外邦女子為妻.
今天我們不介意.
不過對於當代以色列一個復元很大的波斯帝國.
一個小小的民族如果在那裡侵襲.
他們就隨時不單止是血緣被侵襲.
而是他們的信仰,恩著和外邦來通婚.
他們都會隨時瓦解.
以斯拉要重建的不單止是一個聖殿.
而是因為有聖殿而下的信仰生活.
所以他要繼續宣讀律法.
勉勵那些教導那些以色列民.
走回上帝深意的軌道當中.
當這班人重新看律法.
真的明白了,真的被教導了.
第十章就講到那個反應.
以斯拉禱告認罪,哭泣,苦伏在神殿面前的時候.
有以色列中的男女孩童聚集到以斯拉那裡.
成了大會.
眾民無不痛哭.

$^{1041}$蜀以蘭的子孫耶傑的兒子士加尼對以斯拉說.
我們在此地娶了外邦的女子為妻.
肝飯了我們的神.
然而以色列人還有指望.
現在當與我們的神立約.
優者一切的妻.
棄絕他們所生的.
照著我主和那恩神命令戰經之人所預定的安律法而行.
你起來,這是你當辦的事.
我們必幫助你,你當奮勉而行.
以斯拉的十章一至四節.
重新頒佈之後就開一個陪靈會.
這個大會就是重新一起祈禱,一起教導.
頂展會或者當時的一些領袖就起來.
我們和上帝要悔改.
甚至適日做錯的事他們要斷言斬斷.
重新來到建立整個民族.
有一個核心來到過他們的信仰生活.
最後那句很有意思.
我們必幫助你,你當奮勉而行.
不是用威嚇的話.
說你們不聽就將你們囚禁懲治.
這個是皇的命令.
不過在以斯拉或者在他的群體裡面.
我們幫助你.
你做得不好,你面對過去的一些困難.
我們幫助你.
我們會給你勇氣奮勉而行.
帶著上帝給他的決心和勇氣.
走回一個聖潔的軌道上.
今日我們看的經文很闊.
所以我沒有辦法完全去仔細.
不過我很想大家一路看以斯拉記.
你見到今日我們基督徒都用這段經文.
自己互相勉勵.
講句真心話.
今日世界很多香港人講.
基督徒要留散,分散在地方才是要走的路.
不過聖經從來,或者至少我們看以斯拉記.
分散不是我們的目標.

$^{1081}$聚集在一個散亡,失望,困難當中.
我們都要勉勵地聚集在一起.
可能我們面對很多挑戰.
但我們的聚合在一起是有意思的.
昔日當然有個聖殿建築物來維繫.
去到今天我們基督徒不是再用一個物質的聖殿.
來維繫我們之間彼此的聚.
而是用基督耶穌以他自己的生命作為殿.
奉基督耶穌的名,我們這一班人聚集.
我盼望我們不是要標榜或者鼓勵弟兄姊妹要散開.
反而在一個大家這麼混亂,紛亂之際.
我們更加要勉勵自己聚合在一起.
坦白說,我不會否定弟兄姊妹有不同的想法要移民等等.
我們不會否定的,是嗎?.
不過在神學裡面,大家面對著很多香港社會所說的議題.
我們未必跟香港社會的想法來過我們基督徒的人生.
反而我們可否看回上帝在我們生命裡面.
他激動我們的心的是甚麼呢?.
今天我們留在這個地方.
我們留下來是帶著失望,悲哀還是留下.
還是我們帶著上帝給我們的激勵.
奮勉地將人帶到上帝的殿當中.
帶到我們這個群體裡面.
今天以色列記提醒我們.
真正令到一班弟兄姊妹,一班神的子民聚合在一起.
不是人的作為.
因為甚至起聖殿的領袖自己都成為罪魁.
自己都有錯的地方.
不過最重要的是安著上帝.
安著上帝的話語.
讓真理讓我們重新結合在一起.
盼望在未來的日子裡面.
我們大家攜手.
今年教會所說的並肩用創.
那個意思就是這樣.
是一個個個都覺得亂走的世代.
我們勇敢走在一起.
我們挺起胸膛.
為未來的日子我們一起來打拼.
願主繼續激動我們的心.

$^{1121}$讓以色列繼續對我們講說話.
\newpage



\section{羅馬書 15:13}
\label{sec:hASHSP9n8fQ}
\textbf{2022年8月28日崇拜直播｜區祥江博士｜大有盼望的秘訣｜羅馬書15\_13}
\newline
\newline
連結: \href{https://youtube.com/watch?v=hASHSP9n8fQ}{\texttt{https://youtube.com/watch?v=hASHSP9n8fQ}} ~~~~ 語音日期: 2022-08-30
\newline
\newline
\hyperref[sec:ZKaifzpEjak]{\small{< < < PREV SERMON < < <}}
~
\hyperref[sec:index_chronic]{\small{[返順時目]}}
~
\hyperref[sec:0H976XcFnC4]{\small{> > > NEXT SERMON > > >}}
\newline
\newline
羅馬書 15:13
\newline
\begin{longtable}{cl}
\hline
\hline
章節 & 經文 (和合本修訂版)\\
\hline
15:13 & \begin{tabularx}{0.7\textwidth}{X} 願賜盼望的神，因你們的信把各樣的喜樂、平安充滿你們的心，使你們藉著聖靈的能力大有盼望！ \end{tabularx} \\ \\
[1ex]
\hline
\hline
\end{longtable}
$^{1}$我們以下要聆聽神給我們的說話.
我們今天要聽神的說話.
記載在羅馬書15章13節.
新約聖經第222頁.
羅馬書15章13節.
但願使人有盼望的神.
因信將諸般的喜樂,平安.
充滿你們的心.
使你們藉著聖靈的能力大有盼望.
羅馬書15章13節 讀筆.
我們今天很榮幸請到歐長廣博士.
歐長廣博士跟我們講道.
歐博士是中國神學研究院的輔導科教授.
他的講題是「大有盼望的秘訣」.
我們將時間交給歐博士.
謝謝.
我想可以將經文打上來就可以了.
不用將我的帥哥相片放在這裡.
好,大家早上好.
早上好.
其實疫情又去到八千多宗了.
有沒有盼望清零?.
沒有吧?.
香港回歸二十五週年.
又有新的特首了.
這位特首有沒有給我們盼望?.
一點點也有吧?.
好像信心指數高了一點.
其實盼望是我們人生裡面很重要的信念.
讓我們在疫情,困難裡面能夠繼續走上去.
所以今天一反我過往講道.
通常是講一段經文多的.
今天就講一節經文.
因為這節經文裡面有很豐富的訊息.
想跟大家分享.
我們先看這節經文.
和合本的翻譯多了一點.
這裡說「但願使人有盼望的神」.
但原本的意思是沒有使人這個字.
「願盼望的神」.

$^{41}$就是將焦點放在我們的神是一個有盼望的神.
而不是祂給我們盼望.
然後就「將諸般的喜樂,平安充滿你們的心」.
「使你們藉著聖靈的能力大有盼望」.
盼望是相當重要的.
如果一個人沒有盼望的話.
最差的情況就會去到自殺,輕生.
就好像對前路,前景再沒有盼望.
所以如果能夠在人生裡面找到盼望.
對我們來說是相當重要的.
我想用一點時間和大家看看盼望的神.
為什麼會這樣來形容我們的神呢?.
God of Hope.
我自己是研究心理學的.
在研究的過程中.
近年有一個叫「正向心理學」.
研究哪些人有盼望.
是容易有盼望的.
我們先看人.
一個人為什麼會有盼望呢?.
有兩個主要的元素.
一個人有盼望.
第一個是「Way Power」.
一個人能夠找到出路.
找到方法解決問題.
例如塞車.
總有方法離開那個場面.
逃出生天.
總之有很多辦法.
很有方法.
我們稱為「很有盼望」.
因為他能夠解決面前的問題.
所以他有盼望.
第二個特徵.
一個有盼望的人是.
「Will Power」.
「Will Power」是什麼意思呢?.
是一個很堅毅的信心.
他有持久力.
不會很快放棄.

$^{81}$知道逆境在.
他繼續堅持下去.
到最後是柳暗花明又一尊.
總會有出路.
如果有這兩方面特徵的人.
我們就稱為「有盼望的人」.
我想問問大家.
你們有沒有這兩個特徵呢?.
「Way Power」.
很多辦法.
很多東西有出路.
有方法處理.
「Will Power」.
有沒有一個堅毅的信心.
持久力.
有沒有?.
有嗎?.
那好啊.
怎麼來的呢?.
不過盼望這東西.
基督徒的心理學家都很有興趣.
所以就加多了一個向度.
是由「W」這個字開頭.
大家猜到嗎?.
就是「Wait Power」.
等候的能力.
為什麼這樣說呢?.
羅馬書八章二十四,二十五節.
我讀給大家聽.
我們得救在於盼望.
可視看得見的盼望.
就不是盼望.
但我們若盼望那看不見的.
我們就耐心等候.
其實我們基督徒.
有別於一般人.
我們有一個等候的能力.
其實這次新官.
真的沒想過這麼多年.
SARS很短時間就走了.

$^{121}$怎麼這麼久都沒過呢?.
等了多久?.
誰想去旅行.
想到發燒了.
這麼久都沒去旅行.
原來等待是很考功夫.
但如果我們是一個有盼望的人.
我們有出路.
我們有堅毅的能力.
我們有等候的力量.
如果我將這三個特質.
放在神那裡.
神有沒有力量?.
我記得有首歌.
大家可能很喜歡唱.
英文是God will make a way.
大家記得.
神會開一條路.
就算前面好像沒路.
祂做工作的能力.
我們眼未必能夠看見.
祂會為我們開一條新的出路.
我們相信神是無所不能的.
所以祂有很多方法.
去為我們解決前面的困境.
神有沒有 will power?.
有沒有堅毅的能力?.
其實有.
祂忍我們人類忍了這麼多年.
然後才派耶穌來.
主在來的日子.
還在等待.
祂的旨意很堅定.
我們的神是一個有 will power 的神.
我最欣賞神就是祂的 wait power.
祂等待浪子回頭.
你想多辛苦,等了多久.
我想我們人很多時候不是太明白.
為什麼還要等這麼久.
但是神的 timing(時機).

$^{161}$往往不是我們能夠了解的.
可能當事情過去之後.
我們再看那些 timing(時機).
就會覺得很奇妙.
原來神是有祂自己的時間去工作.
所以願盼望的神.
給我們因為我們的信.
將各樣的喜樂和平安的心.
充滿你的心.
我想這段經文中間這段.
想挑戰一下大家.
其實這段日子.
你的心充滿什麼多呢?.
充滿什麼?.
憂慮,低沉.
好像沒什麼忙.
你的心充滿什麼呢?.
其實一個人.
曾經說過我們要保守我們的心.
勝過保守一切.
所以你的心被什麼充滿呢?.
其實很影響我整個人的生活.
如果充滿了憂慮.
我想很多媽媽都會憂慮.
過了一晚鐘.
可能又要忘課了.
害不害怕?.
想起都害怕.
困獸鬥在家裡忘課.
心裡面會不會有很多憂慮呢?.
不平安呢?.
或者很不開心.
很困住呢?.
所以這裡.
當神給我們盼望的時候.
有兩個特點.
除了證明你剛才有沒有.
什麼way power.
view power 之外.
你的生活或者你的內心.

$^{201}$充滿什麼呢?.
其實顯明你是有沒有一個盼望的人.
兩個特質.
第一喜樂.
開心嗎?.
有沒有喜樂?.
第二.
有沒有平安?.
先說喜樂.
其實喜樂.
我們基督徒的喜樂和快樂不同.
快樂是無限聚令.
大家可以一起喝茶.
很開心.
這是一個透過社交.
或者在生活裡面.
給我們歡愉的感覺.
但基督徒說的那種喜樂.
是很奇怪的.
保羅坐牢也有喜樂.
不是局限於環境.
就算環境未如預期.
未如理想.
但我們心裡面仍然有那種喜樂.
有沒有呢?.
其實你能夠流露喜樂.
也能夠讓身邊的人看到.
你是一個有盼望的人.
第二個特點就是平安.
心裡面平不平安呢?.
如果你有一個盼望.
你知道將來你的前景.
或者將來會是怎樣.
其實是能夠給你今天的內心平安.
盼望和平安的關係是很緊密的.
如果你有盼望.
我既然信神.
我有盼望.
我今時今日就不應該被思慮煩擾.
為明天憂慮太多.

$^{241}$令到我心裡面沒有平安.
所以這兩個是我們的挑戰.
你問問自己.
不要說很久.
想起昨天.
昨天你的心.
是充滿了甚麼感受.
充滿了甚麼情緒.
是擔心顧慮,憂愁.
還是平安喜樂呢?.
當神給我們盼望的時候.
會將喜樂和平安充滿我們的心.
我巴不得.
鄧小姐.
就算在這些困境.
疫情未過.
或者市道未好.
這些環境外面的事.
不要打沉我們.
我們仍然有一種喜樂的心.
同樣有一個平安去過我們的生活.
因為我們知道盼望的源頭在哪裡.
這裡就說了一個秘訣.
我們怎樣可以得到盼望呢?.
這句金句就說得很清楚.
是因信.
因你們的信.
這個信其實有點像一個小孩.
很信父母.
他很安心地睡在懷抱裡.
最近其實我在這個疫情下.
有幾件事很開心.
第一就是我做了外公.
我女兒生了一個女兒.
我很想念她.
沒有跟她一起住.
經常想著何時去探望她.
今晚有機會了.
她有一段時間住在我家.
我發覺我挺會抱BB.

$^{281}$她哭了來到我手.
都可以令她安詳地睡著.
我懂得抱懂得搖她.
我又滿足她又開心.
所以我想作為一個信徒.
我們的天父其實就是我們的天父爸爸.
我們有沒有一個嬰孩的單純心.
去到他的懷抱裡.
我們信得過.
風大雨大都沒有關係.
就像耶穌在船上的門徒.
耶穌照樣睡.
門徒又怕死.
但其實耶穌可以平靜風浪.
我們的天父很愛我們.
所以我們在信仰生活裡.
有一個很單純.
很願意去到天父的懷裡.
去倚靠他.
去信得過他.
其實就是得到盼望的源頭.
我一開始問大家.
你信不信特首.
或者不要說自己.
不要說香港.
其實起初我想特朗普下台後.
拜登對中國會溫和些吧.
其實還是靠不住.
不是我們心目中所想的.
所以其實我們人很奇怪.
我們很想將盼望放在領袖身上.
當然如果是一個敬虔愛神的領袖.
他會帶領我們.
但很多領袖都是強差人意.
不知道大家見了這麼多特首.
其實都心灰意冷.
不敢抱太大的盼望.
所以這裡其實是說.
我們會不會將我們的信心.
能夠帶到天父.

$^{321}$唯有神最信得過.
唯有神才能給我們真正的盼望.
接著.
這裡說.
使你們藉著聖靈的能力帶有盼望.
我想我們內心很多時候會交戰.
或者很多掙扎,矛盾.
其實我們的內心如何可以.
教得準,對準神.
拿到喜樂和盼望呢?.
其實是靠聖靈在我們裡面的工作.
但聖靈的能力,聖靈的同在.
跟我們的心有沒有對準是很相關的.
如果我們還是用自己的方法.
左撲,右撲.
為很多事情奔跑.
沒有靜下來.
去看到神自己的工作.
所以我想在這些疫情下.
其實不能出街這麼多.
其實我們香港人是不行的.
一有得出街就出.
所以這幾個星期街上都很多人.
其實我們如果.
有時覺得疫情將我們被迫停下來.
可以觸摸一下自己的內心.
多些向神開放.
聖靈的能力,祂的工作.
是要我們配合的.
如果我們的心太過忙亂.
我們對不準神.
就算祂想發功.
有能力出來.
其實我們都沒有去接受.
所以這裡說.
藉著聖靈的能力.
所以我們要開放去聽.
或者去給聖靈工作才可以.
羅馬書八章二十六節這裡這樣說.
「況且我們的軟弱有聖靈幫助」.

$^{361}$「我們本不曉得怎樣禱告」.
「只是聖靈親自用說不出來的嘆息替我們禱告」.
有沒有時候當你很灰心,喪志.
連禱告的力量都沒有.
你在那些時候會否感應到.
其實聖靈是在陪伴你呢?.
聖靈是用祂說不出的嘆息替我們禱告呢?.
當我們肯去禱告的時候.
其實那個轉變就來了.
我不知道大家有沒有一個感覺呢?.
有時候你很抗拒去祈禱.
覺得很艱難,很灰心,失意.
那就不去禱告了.
但如果我們兩三位弟兄姊妹走在一起.
一起祈禱的話.
祈禱是有一個很神奇的作用.
本來我們是很低沉,很喪志.
但我們一開聲祈禱的時候.
不知為何我們會用神的眼光去看事情.
會教訓我們對困難,逆境的方向,態度.
而在這個心意的轉變的時候.
其實我們就得到鼓勵.
繼續走下去.
所以大家不要停止聚會.
不要停止和你的弟兄姊妹一起祈禱.
不要停止你自己的思圖.
因為當你開聲去祈禱的時候.
其實就讓聖靈在你心裡去工作.
藉著聖靈的能力.
當我們一祈禱.
其實我們就用了神的眼光來看事情.
在祈禱之前.
我們還是用人的計算,籌劃.
這是一個很特別的轉換.
所以我們需要多些禱告.
最後.
這裡說.
「使你們藉著聖靈的能力大有盼望」.
大有盼望.
大有盼望的英文翻譯是Overflow with Hope.

$^{401}$Overflow.
就好像你倒水,水杯滿了.
就流出來.
說不好的時候.
我們就說那個人情緒化.
將情緒spill over(沖淡)出來.
就是燙傷其他人.
這是一個負面的形容.
這裡Overflow.
就是說你這個人充滿了盼望.
很積極,很樂觀,很有信心.
這種情操會影響你身邊的人.
幫助到其他人.
所以這裡其實是說.
我們可不可以帶給這個沒有盼望的世界.
一個新的希望.
如果我們對神有足夠的信心.
又經歷到祂給我們的喜樂和平安.
我們應該會散發到.
有盼望的人生在我們身邊的親友,同事.
他們已經很頹廢了.
或者對香港沒有盼望.
所以離開英國,加拿大.
移民逃跑.
留下的人又走不了.
是困在這裡還是怎樣呢?.
還是我們看到可能有更多機會.
有一個新的局面.
你們的主題是什麼?.
「並肩勇闖」.
即是一個新的時代.
我們能否給予.
好像困在一個城裡的人.
還是我們看到有新天新地.
看到有將來的盼望.
以致我們將我們的盼望.
流射出來.
影響我們身邊的人.
他們得到我們的祝福.
我們一起再將這段經文讀一次.

$^{441}$回答這句金句的含義.
整句經文裡面.
寫羅馬書結尾時.
是對信徒的祝福.
心願都是祝福到我們.
你慢慢讀的時候.
將這些元素記在心裡.
然後將來我們的生活.
我們可以.
為何我現在這麼憂愁.
好像這麼不平安.
回過頭來.
會否問自己.
我是否對神不夠信靠.
我信不信神是一個有盼望的神.
行不行?.
將這句金句除了明白了結構.
和裡面的真理含義之外.
都對我們的生活有應用.
去檢視我們是怎樣.
我們有沒有盼望.
我們有沒有喜樂平安.
沒有的時候有什麼原因.
是否需要校正.
是否需要親近神.
在聖靈裡面取得能力去面對.
行不行?.
我們一起讀這句金句兩次.
開始.
但願使人有盼望的神.
因順將諸般的喜樂平安.
充滿你們的心.
使你們藉著聖靈的能力大有盼望.
羅馬書十五章十三字.
再來多一次.
但願使人有盼望的神.
因順將諸般的喜樂平安.
充滿你們的心.
使你們藉著聖靈的能力大有盼望.
羅馬書十五章十三字.

$^{481}$心願這次的經文.
在未來的日子或星期.
成為大家生活的鼓勵.
將時間交給主席.
\newpage



\section{民數記 20:2-13}
\label{sec:0H976XcFnC4}
\textbf{2022年9月11日崇拜直播｜陳冠成牧師｜以石為記:擊打出水的磐石｜民數記20\_2-13}
\newline
\newline
連結: \href{https://youtube.com/watch?v=0H976XcFnC4}{\texttt{https://youtube.com/watch?v=0H976XcFnC4}} ~~~~ 語音日期: 2022-09-11
\newline
\newline
\hyperref[sec:hASHSP9n8fQ]{\small{< < < PREV SERMON < < <}}
~
\hyperref[sec:index_chronic]{\small{[返順時目]}}
~
\hyperref[sec:uD1kVchCRPU]{\small{> > > NEXT SERMON > > >}}
\newline
\newline
民數記 20:2-13
\newline
\begin{longtable}{cl}
\hline
\hline
章節 & 經文 (和合本修訂版)\\
\hline
20:2 & \begin{tabularx}{0.7\textwidth}{X} 會眾沒有水，就聚集反對摩西和亞倫。 \end{tabularx} \\ \\ \relax
20:3 & \begin{tabularx}{0.7\textwidth}{X} 百姓與摩西爭鬧，說：「我們恨不得與我們的弟兄一同死在耶和華面前。 \end{tabularx} \\ \\ \relax
20:4 & \begin{tabularx}{0.7\textwidth}{X} 你們為甚麼領耶和華的會眾到這曠野，使我們和我們的牲畜都死在這裡呢？ \end{tabularx} \\ \\ \relax
20:5 & \begin{tabularx}{0.7\textwidth}{X} 你們為甚麼領我們從埃及上來，把我們帶到這壞的地方呢？這地方不能撒種，沒有無花果樹、葡萄樹、石榴樹，也沒有水喝。」 \end{tabularx} \\ \\ \relax
20:6 & \begin{tabularx}{0.7\textwidth}{X} 摩西、亞倫離開會眾面前，到會幕的門口，臉伏於地；耶和華的榮光向他們顯現。 \end{tabularx} \\ \\ \relax
20:7 & \begin{tabularx}{0.7\textwidth}{X} 耶和華吩咐摩西說： \end{tabularx} \\ \\ \relax
20:8 & \begin{tabularx}{0.7\textwidth}{X} 「你拿著杖去，和你的哥哥亞倫召集會眾，在他們眼前吩咐磐石湧出水來，水就會從磐石流出，給會眾和他們的牲畜喝。」 \end{tabularx} \\ \\ \relax
20:9 & \begin{tabularx}{0.7\textwidth}{X} 於是摩西遵照耶和華所吩咐他的，從耶和華面前拿了杖去。 \end{tabularx} \\ \\ \relax
20:10 & \begin{tabularx}{0.7\textwidth}{X} 摩西和亞倫召集會眾到磐石前。摩西對他們說：「聽著，你們這些悖逆的人！我們要叫這磐石流出水來給你們嗎？」 \end{tabularx} \\ \\ \relax
20:11 & \begin{tabularx}{0.7\textwidth}{X} 摩西舉起手來，用杖擊打磐石兩下，就有許多水流出來，會眾和他們的牲畜都喝了。 \end{tabularx} \\ \\ \relax
20:12 & \begin{tabularx}{0.7\textwidth}{X} 但是耶和華對摩西、亞倫說：「因為你們不信我，沒有在以色列人眼前尊我為聖，所以你們必不能領這會眾進入我所要賜給他們的地去。」 \end{tabularx} \\ \\ \relax
20:13 & \begin{tabularx}{0.7\textwidth}{X} 這就是米利巴水，因以色列人與耶和華爭鬧，耶和華在他們面前顯為聖。 \end{tabularx} \\ \\
[1ex]
\hline
\hline
\end{longtable}
$^{1}$今日的經訓是在舊約文掃記.
第20章1至13節.
在堂會的聖經舊約188頁.
請聽我讀出神的話語.
正月間以色列傳會眾.
到了尋的曠野.
就住在迦底斯.
米利安死在那裡.
就葬在那裡.
會眾沒有水殼.
就聚集攻擊摩西,亞倫.
百姓向摩西爭鬧說.
我們的弟兄曾死在耶和華面前.
我們恨不得與他們同死.
你們為何把耶和華的會眾.
領導者曠野.
使我們和牲畜.
都死在這裡呢.
你們為何逼著我們出埃及.
領導我們到這壞地方呢.
這地方不好殺種.
也沒有無花果樹,葡萄樹,石榴樹.
又沒有水殼.
摩西,亞倫離開會眾.
到會幕門口,俯伏在地.
耶和華的榮光向他們顯現.
耶和華曉以摩西說.
你拿著杖去.
和你的哥哥亞倫.
招聚會眾.
在他們眼前.
吩咐盤石發出水來.
水就從盤石流出.
給會眾和他們的牲畜殼.
於是摩西照耶和華所吩咐的.
從耶和華面前取了杖去.
摩西,亞倫就招聚會眾到盤石前.
摩西說.
你們這些背叛的人聽我說.
我為你們試水.

$^{41}$從這盤石中流出來摸.
摩西舉手.
用杖擊打盤石兩下.
就有許多水流出來.
會眾和他們的牲畜都喝了.
耶和華對摩西,亞倫說.
因為你們不信我.
不在以色列人眼前尊我為聖.
所以你們必不得領者會眾.
盡我所賜給他們的地去.
這水名叫米里巴水.
是因以色列人向耶和華爭鬧.
耶和華就在他們面前顯為聖.
(拍手).
請姐妹平安.
相信剛才所讀的經文.
大家都熟悉.
摩西和亞倫帶領了以色列人.
接近四十年.
為何到最後因為擊打了盤石兩下.
就不能進入應許之地.
我們今天嘗試找出原因.
剛才所讀的經文.
一開始提到他們來到迦底斯.
迦底斯是在四十年前.
當時是民數記十三至十四章.
記載以色列人派了十二個探子.
上去窺探迦南地.
結果回來報信的時候.
十二個探子裡有十個報惡信.
說不應該上迦底斯.
只有約書亞和迦勒說.
這是上帝的應許.
是一定沒有問題的.
事實上細心看這兩次的經歷.
其實有很多相似的地方.
特別是大家可以打開程序表的內頁.
廣路大綱裡其實做了一個對比.
分別是他們兩次在迦底斯.
這個地方所發生的事情.

$^{81}$包括民數記十三至十四章.
和今天我們看的經文民數記第二十章.
你會發覺其實無論是地方.
或者裡面所發生的經歷.
大致上都是一模一樣.
你看到十四章裡說.
百姓向摩西和阿倫發怨言.
為什麼要帶領我們來到埃及.
出埃及來抗疫這個地方.
於是乎上帝顯現.
懲罰那些不信的以色列人.
不准他們進入應許之地.
而今天我們看到二十章.
相隔接近四十年後.
百姓又再向摩西和阿倫爭鬧.
又是在說為什麼要帶領我們出埃及.
來到這個鬼地方.
來抗疫這個鬼地方.
於是乎上帝又再顯現.
不過他這次懲罰的不是以色列人.
而是不信的摩西和阿倫.
上帝不容許他們二人.
帶領百姓進入應許之地.
所以你看到兩段經文.
雖然相隔接近四十年.
但其實它帶有相同的訊息.
特別留意兩次都涉及到.
爭鬧和不信.
原來這兩件事.
會令上帝的子民.
他們錯失上帝的應許.
看回今天的經文.
事件的起因是說.
匯仲因為沒有水喝.
於是就聚集攻擊摩西和阿倫.
你看到他們是給這兩位領袖說話.
看回那些字句.
我們恨不得和那些從前.
死在上帝面前的弟兄.
跟他們一起死.

$^{121}$都比現在沒有水喝更淒涼.
摩西和阿倫你們怎樣帶領我們?.
為什麼你們要逼我們離開埃及.
帶我們來曠野這個壞地方.
來折騰耶和華的匯仲?.
潛台詞其實是.
我們現在不是對上帝不滿.
我們是耶和華的匯仲.
我們只不過是投訴摩西和阿倫.
你們兩位領導無方.
連累我們受苦.
我們不是要投訴上帝.
事實上你會看到.
摩西和阿倫他們不是第一次.
面對以色列人抱怨出埃及的場面.
他們很有經驗.
他們知道誰才可以將事件擺平.
其實只有上帝出手才能解決整件事.
於是乎你看到他們很快離開爭鬧的場景.
他們兩個人來到會幕的門口.
尋求上帝的指引和幫助.
第八節裡說.
上帝也吩咐摩西拿著象徵著上帝權柄的象.
和阿倫一起焦醉匯仲.
然後在匯仲眼前吩咐盤石發出水來.
這樣就會有水給匯仲和他們的生畜飲用.
於是兩人就照上帝的吩咐而行.
你看到到目前為止.
摩西和阿倫有沒有犯錯?.
有沒有?.
其實沒有的.
他們一直都表現得很好.
之後你又看到.
摩西和阿倫焦醉匯仲.
來到盤石面前.
於是摩西就舉手用象擊打了盤石兩下.
結果有沒有水流出來?.
有的.
結果有水流出來.
確實解決了匯仲因為沒有水飲用的紛爭.

$^{161}$對不對?.
不過問題來了.
事件不是解決了嗎?.
但為什麼上帝都要懲罰摩西和阿倫.
到底他們出了什麼問題?.
就是我們今天要去細看的.
沒錯,你會看到.
首先摩西用象擊打了盤石兩次.
這是不是上帝的吩咐?.
不是.
上帝沒有這樣吩咐他.
不過問題不是因為摩西做錯了什麼動作.
或者用錯了方法.
這麼表面.
事實上當我們看另一段經文.
詩篇106篇32至33節.
詩人讓我們窺探一下.
當時摩西的心路歷程.
到底是怎樣.
我們去看看他的內心世界.
我想請大家同讀程序表.
旁邊那頁.
第一段經文詩篇106篇32至33節.
這裡是用和合本修訂本.
請大家一起讀.
預備起.
好,謝謝.
你看到詩篇裡表示.
摩西在米利巴水事件中受到虧損的原因.
是不是因為他擊打了盤石?.
隻字不提.
而是在說他被哪些人惹怒.
被匯眾惹怒.
心裡不舒服.
忍不住說了什麼.
急躁的話.
或者說是霧失的話.
所謂急躁的話.
其實就是我們今天看到第20章第10節裡.
他怎樣和匯眾針鋒相對.

$^{201}$這裡說另一段經文.
和合本修訂本也比較傳神.
他怎樣焦醉了匯眾.
然後和他們針鋒相對呢?.
這裡特別說.
聽著,你們這些活逆的人.
我們要叫這盤石流出水來給你們喝嗎?.
你看到.
百姓強調自己是耶和華的匯眾.
但摩西怎樣反諷他們?.
在摩西心中他們是不是耶和華的匯眾?.
摩西說你們不是.
你們是叛徒.
是背叛上帝的人.
難道還要我去服侍你們.
給你們水喝?.
為什麼會停在這裡?.
為什麼摩西這麼不服氣?.
為什麼他這麼不服氣?.
我們看回很多時候在五經裡.
正常的情況,以往發生的情況.
當百姓埋怨攻擊摩西和亞倫的時候.
上帝通常會顯現.
然後怎樣?.
罰誰?.
會懲罰匯眾,向匯眾發怒.
替摩西和亞倫主持公道.
好像替他們平反.
然後當摩西和亞倫眼見.
百姓受罰受苦哀求的時候.
這兩個領袖就會怎樣?.
就會挺身而出.
在上帝面前替百姓代求.
求上帝寬恕他們.
我們記得很多時候都發生在五經裡.
但這次呢?.
這次的情況相反.
上帝這次沒有責備他們的百姓.
反而是怎樣?.
順從他們飲水的要求.

$^{241}$並且吩咐摩西和亞倫.
你們要在匯眾面前吩咐盤石.
代表我去給他們飲水.
你看到上帝這次選擇的不是懲罰百姓.
而是連恤,示人.
為什麼會令摩西這麼氣定?.
為什麼?.
是因為摩西自己的期望落空了.
他期望上帝會怎樣?.
上帝你這次不是應該支持我和亞倫嗎?.
你不是應該替我主持公道嗎?.
他們經常陪伴你.
又屢勸不改.
為什麼上帝你還要對他們這麼好?.
明明這次是他們不對.
你不但不罵他們.
還要我們回頭服侍他們.
上帝你罵罰罰他們.
然後才叫我吩咐給他們飲水.
這樣我也會比較害怕.
我會有點氣.
現在整件事好像是我和亞倫不對.
好像要我們向他們認低威.
摩西忍不忍得這口氣?.
忍不住.
受夠了.
結果你看到摩西.
他沒有按照上帝的心意.
上帝是叫他用口吩咐盤石出水.
但現在摩西就照自己的意思.
他用口怎樣?.
他和匯眾爭鬧.
並且擊打盤石.
來發洩自己的不滿和憤怒.
事實上我們不難體會摩西這份軟弱.
你想像一個人.
他長期來到去仆心仆命事奉.
但是得不到肯定?.
得不到別人的肯定.
讚賞和鼓勵.

$^{281}$反而遭到別人很多批評.
投訴和攻擊.
事奉的人很容易會怎樣?.
有時真的會失去那份耐性和熱心.
甚至會去到另一個極端.
用一個埋怨的方式.
將自己內心的苦毒無限放大.
扮演了一個受苦者的角色.
一個受害者的角色.
最後他很可能會選擇什麼?.
不出聲.
離場.
不再事奉.
又或者像摩西那樣會怎樣?.
他會繼續事奉,不想停.
但是他帶著很多憤怒.
很多苦澀.
來面對他服侍的對象.
我曾經記得.
有位弟兄姊妹來找我傾談.
分享到他在工作上面對一些試探.
很多時老闆會要求他幫老闆隱瞞客戶.
這位弟兄姊妹覺得這樣真的不太好.
他其實不想繼續.
於是在言談中.
我們去想有什麼可以對策.
我鼓勵他有沒有想過.
不妨考慮一下會否轉換一下環境.
離開試探的地方.
找一份新的工作.
並且要確信既然上帝給你有這個責心.
你要相信上帝會有預備,有安排.
但是你有沒有想到.
這位弟兄姊妹突然間很生氣.
他發怒.
他罵我的建議很離地.
不體貼.
他需要養家的經濟壓力.
於是他開始埋怨.
為什麼上帝要給他這麼多不好的際遇.

$^{321}$為什麼對我這麼不好.
那一刻我不單單感受到對方有情緒.
其實我也感受到.
我也有情緒.
因為我無端端被人省了一輪.
當時我有兩個選擇.
一是像摩西那樣跟他針鋒相對.
又或者我有另一個選擇.
就是求上帝給我有更多的愛心.
更多的耐性去承載他這一刻的需要.
他的軟弱,他的情緒.
讓上帝的旨意可以在我和他之間彰顯.
結果,死亡了,靜默了一段時間之後.
我決定沉著氣,繼續聽他神.
因為我知道人的怒氣不能夠承受上帝的意義.
事實上也曾經和大家分享過.
有很多研究情緒的專家也說過.
一個人無論你有多麼多的經驗.
有多麼高的學歷.
當他處於憤怒狀態的時候.
記不記得他的智商剩下多少歲.
記不記得.
五歲.
想像一下一個五歲的小朋友.
他做的事情有時會不理智.
以致當一個人無論多有經驗多有學識.
當他處於憤怒的情況.
他說的話和行為都變得不成熟.
明明那件事他很有能力可以辦妥.
但最後通常是窩河收場.
你和我大概也有很多這樣的經驗.
明明很多事情如果讓你再做一次.
你一定會好好管理你的情緒.
不會翻我的事.
不過當憤怒控制不了的時候.
很多時候就會好事都變壞事.
就好像摩西那樣.
其實整件事反過來.
如果你有機會讓他再來一次.
我肯定他一定不會這樣做.

$^{361}$事實上從他的經驗從他的能力.
他絕對有可能將整件事辦得更好.
更加榮耀上帝更加彼此造就.
不過在承諾之下他就壞了大事.
或許他想著我和匯仲吵吵架.
發洩一下敲一下盆石.
沒事的.
誰知就出事了.
你看到第十二節.
上帝對摩西和亞倫的責備.
上帝說他們是甚麼?.
不信.
沒有在以色列人的眼前尊上帝為大.
其實這句說話是甚麼意思呢?.
首先我們要明白.
給百姓有水飲解他們口渴的困苦.
是誰的心意?.
是上帝的旨意.
亞倫和摩西的職責只是帶著上帝的像.
代表上帝來吩咐盆石出水.
但摩西現在卻是用事奉來解決甚麼?.
解決自己的私怨.
匯仲現在眼裡看到的.
只是摩西的苦讀.
能否看到上帝大能的恩典?.
未必看到.
摩西在百姓面前高舉的是自己的血氣.
沒有高舉上帝的能力和恩典.
也就是說.
兩位領袖沒有在百姓眼前.
沒有在這件事上尊上帝為聖.
而另一方面.
其實全場誰才有資格斥責匯仲是活逆的人?.
只有一位.
只有上帝才有資格判斷誰才是活逆的人.
但上帝這次沒有這樣做.
上帝選擇寬恕放過他們.
反而是摩西自己受不了.
他就做了審判的角色.
他代替上帝去審判.

$^{401}$其實反映他骨子裡.
他服不服上帝的裁決?.
其實是不服的.
他不願意接受.
相信上帝這次善待匯仲.
不替自己主持公道.
這個安排是好的.
摩西不願意接受.
於是他就自己出手看天行道.
所以摩西不願意信服.
上帝因待百姓的決定.
其實就是不信.
因為他不相信.
接受上帝就是對罪人滿有恩慈的上帝.
所以上帝的判決很正確.
裁定他們不信.
不信上帝的裁決是未善.
於是自己就出手干預.
至於阿倫.
我們要說一下他.
他其實什麼都沒有做.
他有沒有出聲?.
沒有.
他有沒有敲打盆石?.
沒有.
那為什麼要罰他?.
關他什麼事?.
他什麼都沒有做,關他什麼事?.
就是因為他什麼都不做.
就是因為他什麼都沒有做.
所以關他什麼事?.
如果你還記得.
上帝是吩咐阿倫.
你要和摩西一起朝聚會眾.
並且一起吩咐盆石出水.
而阿倫在整件事裡.
他就是需在一起保持緘默.
他既沒有和摩西一起吩咐盆石出水.
也沒有阻止糾正摩西所犯的錯.
換句話來說.

$^{441}$如果摩西的問題是.
他不應該做的,他做了.
那阿倫的問題是什麼?.
應該做的,他不做了.
所以兩個人都要負上犯錯的代價.
就是上帝說.
你們不能夠繼續帶領百姓進入應許之地.
而經文大概的講解就到這裡.
接下來我們有幾個問題要問.
覺不覺得上帝這次罰得太重?.
覺不覺得?.
坦白一點.
不用擔心.
你覺就是覺.
覺不覺得上帝這次罰得很重?.
重的,因為進入迦南.
對兩個人來說.
特別是摩西是他們畢生的心願.
就這樣就沒有了.
因為一次的犯錯.
就沒有了.
上帝是不是對摩西和阿倫太絕情?.
這是我們很多人都會問的問題.
畢竟武功都怎樣?.
武功都有勞.
給一次機會都可以吧.
我們多多少少會替這兩位領袖覺得很不值.
覺得很惋惜.
但其實我們都明白到.
上帝多給誰的就會怎樣?.
多向誰收取.
我們明白這個道理就是.
上帝給人的位分.
給人的權柄越大.
他對人的要求自然越高.
你看到摩西和阿倫是不是一般普通的百姓?.
他們不是普通的百姓.
摩西是上帝的代言人.
換句話來說.
當摩西向百姓說什麼,做什麼.

$^{481}$其實就代表什麼?.
代表上帝向百姓宣告什麼,做什麼.
現在摩西擅自向每眾說和做.
上帝沒有吩咐過的事情.
我們叫這些做什麼?.
僭越.
或者假傳聖旨.
死不死得?.
其實是死得的.
至於阿倫的職責是什麼?.
他是大祭司.
大祭司的職責是什麼?.
憲制.
憲制背後一個更深層意義其實是.
教導百姓遵守上帝一切的吩咐.
現在阿倫怎樣?.
他自己帶頭,第一個就不遵守了.
所以你看到當領袖自己.
都不遵守上帝的吩咐.
或者說得正確一點.
都背叛上帝的時候.
試問上帝又怎樣能夠讓他們.
繼續帶領自己的百姓前行?.
所以你看到.
上帝收回他們帶領的工作.
免得他們再叛倒人.
得罪上帝.
一方面其實是保護兩個領袖.
免得他們再得罪上帝.
免得他們再帶錯路.
同時亦要他們搞清楚.
其實什麼是事奉的本質.
沒錯,從人的角度真的很可惜.
很遺憾.
帶領了幾十年.
最後想入迦南的心願.
就一舖被毀.
但從事奉的角度來看.
經文想我們用另一個角度思考上帝的判決.
到底事奉是成就誰的心意.

$^{521}$是成就上帝的旨意.
還是成全人自己的意願.
所以從這個角度來看.
上帝不容許摩西亞倫再帶領百姓進入應許之地.
除了責罰他們的不信服.
要他們負上代價之外.
其實你試試反過來看整件事.
上帝其實是給你另一個信服的功課.
現在你們犯了錯,我這樣判.
服不服?.
上一次我吩咐你們,你們不服.
所以你們亂來.
現在我再判,你們服不服?.
如果你們願意接受上帝現在這個判決是好的.
沒錯,你們是不能進入迦南.
但你們蒙不蒙上帝月立?.
仍然可以蒙上帝月立.
仍然有上帝的同在.
你們不會一錯再錯,連上帝也沒有.
相反,你們仍然選擇不信服上帝這次安排.
你硬要進入迦南.
我讓你進入.
上帝會否起月?.
都不會起月.
甚至上帝如果不和你們同在.
那就真的大件事.
所以弟兄姊妹,經文給我們看到.
信心的標記是甚麼?.
信心的成長是甚麼一回事?.
當我們信心很多時仍是幼嫩的時候.
你會發覺很多時上帝會透過我們的經歷.
來祝福,把各樣的好處賜給我們.
讓我們從那些好處,福氣裡體驗到.
上帝的美善和真實.
對不對?.
我們都經歷過.
但隨著信主的年日越來越長的時候.
是否只靠這些?.
是否只靠上帝的好處來支撐信仰?.
是否只靠上帝如何祝福我們.

$^{561}$讓我們的心願成就來建立信仰?.
你知道不是這樣.
不再是這樣.
如果是就大件事.
我們既然相信上帝是真神.
不是那些承傳我們人心裡意願的偶像.
我們就應該更加關注自己所得.
漸漸邁向更加關注.
上帝所得到的是甚麼.
並且我們努力尋求.
上帝在我們身上的計劃.
讓上帝的旨意可以在我們身上.
透過我們的信服來承傳.
當然這並不代表我們將會一無所有.
相反.
你會體會到上帝讓我們得到的.
超過你我的想像.
就正如摩西.
你知道他之後還有掠.
在申命記他繼續掠上帝.
讓我進入迦南.
但上帝斬釘截鐵說不行的時候.
你看到他最後都要選擇信服.
他信服上帝這個決定.
有生之年進不去迦南.
但你再看申命記的時候.
你會看到上帝對他仍然有恩典有憐憫.
上帝換另一個方法來厚待他.
特別申命記的結尾記載.
上帝親自帶摩西遠望整個迦南地.
他進不去.
但上帝親自導賞他.
你可以看回你的子孫之後.
就在這個地方承受上帝的福氣.
並且將迦南地的版圖逐一展示給摩西看.
最後也將摩西親手埋葬.
上帝厚不厚待他?厚待他.
沒有其他人有這樣的福分.
上帝是刻意來施恩憐憫給他.
還記不記得摩西死了之後有沒有醒過來?.

$^{601}$還記得在福音書裡耶穌登山變象嗎?.
摩西就好像在睡夢中被叫醒一樣.
出現在耶穌基督面前這一幕.
如果你留意聖經記載.
當時他們處於迦南邊界.
一個叫做凱撒尼亞菲勒比的地方.
附近的高山.
其實讓我們看到.
上帝愛他終生僕人的方式.
不一定要承傳他當下的意願.
上帝是要將更好的將來賜給摩西.
所以其實整段經文讓我們看到.
我們可以互相提醒.
在你和我的人生裡總有失望的時候.
就算事奉都不例外.
但不等於上帝不愛我們.
有些事情我們當下真的很想得到.
死難求上帝給你.
但上帝不應允.
是否他故意留難我們?.
不是.
上帝將會有更好的安排.
所以謹記當你和我在最失意的時候.
上帝從來沒有離開過我們.
上帝從來沒有放棄過我們.
所以我們都要選擇.
持守對他的信服和忠心.
就是這段經文給我們一個很重要的提醒和挑戰.
最後,第十三節.
這一節很重要.
我想請大家一起讀出來.
最後那段經文.
民數記二十章十三節.
預備起.
好,謝謝.
誒?.
會眾不是和摩西亞倫爭鬧嗎?.
是不是?.
是.
不過,經文怎麼看整件事?.

$^{641}$經文表示他們不只是和摩西亞倫爭鬧.
是和上帝爭鬧.
讓我們看看整個爭鬧的本質到底是什麼.
原來屬靈群體之間的爭鬧和紛爭.
那種彼此不信服,不配合.
背後其實和誰角力呢?.
是和上帝.
意思就是.
不滿意上帝安排某個人服侍你.
或者不滿意上帝安排你服侍某個人.
所以,本質仍然是和上帝爭鬧.
不高興上帝在你生命裡.
令你不舒服的安排.
而爭鬧的結果.
就是今天故事所看到.
整個群體都受到虧損.
不能榮耀上帝.
不能尊上帝為聖.
所以不單止摩西亞倫要負責任.
百姓在這件事上也要負責任.
確實他們飲水的訴求是可以理解的.
上帝最終也解決了滿足他們的需要.
但是如果從頭再來可以再做一次的話.
如果用一個更合宜的方式.
來向摩西亞倫表達他們的訴求的話.
大概就不會引致一個針鋒相對.
兩敗俱傷的場面.
導致他們失去了兩位很重要的領袖的帶領.
等於我們知道在屬靈群體裡.
我們真正的敵人.
只有一位是誰?.
是魔鬼撒旦.
其他可能你覺得不滿意的弟兄姊妹.
那些不是你的敵人.
那些都是你的手足.
上帝所愛的人.
但是在現實生活裡.
基督徒總是會花很多精力.
來處理教會人際的問題.
甚至可能會視某一個.

$^{681}$激親你令你不高興的弟兄姊妹為受敵.
或者這就是正正撒旦常常用來分化我們的伎倆.
他很想轉移我們的焦點.
使我們將有限的資源能力.
將你很寶貴的光陰.
耗費在紛爭裡.
再沒有能力來執行上帝所託付給你的.
所以不要中計.
確實我們很理解.
有人的地方總是有爭拗.
避不了.
因為你和我舉手.
是不是壞人?你覺得.
你覺不覺得你是壞人?.
不覺得吧,我也不舉手.
但是是不是罪人?.
沒有人舉手的.
我們真的不是立心不良.
想搞亂當,想傷害弟兄姊妹.
因為我們真的不是壞人.
但我們都是罪人.
難免會有碰撞.
難免會有彼此虧欠的時候.
特別當我們只是著眼實現自己的願望.
當會眾只顧著沒有水喝而投訴.
當摩西只顧著上帝為何不為我平反而出事的時候.
當我們都忽略了上帝期望我們合一彼此去理解的時候.
教會的問題就越來越多.
所以其實這段經文對我們每一個都有一個很好的提醒.
就是我們需要互相守望.
我們不要被紛爭模糊了.
影響了我們對上帝的尊重.
以致我們才有力量去集中.
實踐上帝所託付給我們的使命.
真的,你想像一下.
其實如果教會的弟兄姊妹都能夠齊心一起起來事奉上帝.
少一點去比較自己的得失的時候.
其實很多人與人之間相處的問題.
那些糾紛其實自然會迎刃而解.
不用那麼複雜.

$^{721}$當我們願意對準上帝的時候.
很多時候撒旦的詭計就會不攻自破.
所以求主監察潔淨我們的心思意念.
幫助我們真的無論對人對事.
給我們有一個更廣闊的胸襟.
更有耶穌基督的柔和謙卑的心腸.
以致我們的言行都能夠榮耀上帝.
造就身邊更多的人.
我們同心祈禱.
主啊,幫助我們.
讓我們真的可以對準你.
上帝啊,我們其實每一個都有軟弱.
我們很多時候即使我們事奉上帝你.
但是有時候我們都會心懷異議.
有時候我們會過於體重自己的需要.
而忘記了上帝你.
忽略了身邊的弟兄姊妹.
以致我們有時候很難去彰顯合一.
我們中了撒旦的詭計.
我們聚焦了在很多人事紛爭上.
以致我們不能夠對準你.
上帝求你藉著今天這段經文.
再一次認清我們的本相.
我們不是立心不良.
但我們都有血氣.
我們是罪人.
我們很需要上帝.
在你的寧時刻提醒保守我們.
我們也很需要弟兄姊妹互相扶持幫助.
以致我們真的可以專注在你.
和你的福音工作上.
以致我們真的能夠同心.
以上帝你的益處為優先.
而放下輕看自己的成敗得失.
上帝幫助我們整個群體.
以致我們真的.
當有怒氣的時候.
當很想埋怨去負面的時候.
緊守我們的心.
讓我們可以真的去.

$^{761}$透過禱告專心來到你的面前.
經歷上帝你的潔淨和平安.
以致我們再有力量.
回到教會裡面.
和弟兄姊妹一起去復和.
一起去服侍你.
為此我們獻上祈禱.
奉耶穌基督.
你的名字來祈求.
阿們.
\newpage



\section{約翰福音 1:38-43}
\label{sec:uD1kVchCRPU}
\textbf{2022年9月18日崇拜直播｜梁家麟牧師｜第一批追隨者｜約翰福音1\_38-43}
\newline
\newline
連結: \href{https://youtube.com/watch?v=uD1kVchCRPU}{\texttt{https://youtube.com/watch?v=uD1kVchCRPU}} ~~~~ 語音日期: 2022-09-19
\newline
\newline
\hyperref[sec:0H976XcFnC4]{\small{< < < PREV SERMON < < <}}
~
\hyperref[sec:index_chronic]{\small{[返順時目]}}
~
\hyperref[sec:PSHqqITZkIo]{\small{> > > NEXT SERMON > > >}}
\newline
\newline
約翰福音 1:38-43
\newline
\begin{longtable}{cl}
\hline
\hline
章節 & 經文 (和合本修訂版)\\
\hline
1:38 & \begin{tabularx}{0.7\textwidth}{X} 耶穌轉過身來，看見他們跟著，就問他們說：「你們要甚麼？」他們對他說：「拉比，你在哪裡住？」（「拉比」翻出來就是老師。） \end{tabularx} \\ \\ \relax
1:39 & \begin{tabularx}{0.7\textwidth}{X} 耶穌說：「你們來看。」他們就去看他在哪裡住。這一天他們就跟他同住；那時大約是下午四點鐘。 \end{tabularx} \\ \\ \relax
1:40 & \begin{tabularx}{0.7\textwidth}{X} 聽了約翰的話而跟從耶穌的那兩個人，其中一個是西門‧彼得的弟弟安得烈。 \end{tabularx} \\ \\ \relax
1:41 & \begin{tabularx}{0.7\textwidth}{X} 他先找到自己的哥哥西門，對他說：「我們遇見彌賽亞了。」（「彌賽亞」翻出來就是基督。） \end{tabularx} \\ \\ \relax
1:42 & \begin{tabularx}{0.7\textwidth}{X} 於是安得烈領西門去見耶穌。耶穌看著他，說：「你是約翰的兒子西門，你要稱為磯法。」（「磯法」翻出來就是彼得。） \end{tabularx} \\ \\ \relax
1:43 & \begin{tabularx}{0.7\textwidth}{X} 又過了一天，耶穌想要往加利利去。他找到腓力，就對他說：「來跟從我！」 \end{tabularx} \\ \\
[1ex]
\hline
\hline
\end{longtable}
$^{1}$響鐘.
《聖經》第十三節.
第十三節 第126 言.
《聖經》第十三節 第126 言.
耶穌轉身來,看見他們跟著,就問他們說:你們要什麼?.
他們說:拉比,在那裡住?.
耶穌說:你們來看,他們就去看他在那裡住..
這一天便與他同住,那時若有新政了..
聽見約翰的話,跟隨耶穌的人,的兩個人,.
一個是西門彼得的兄弟安特烈,.
他先找著自己的哥哥西門對他說:.
我們遇見彌賽雅了..
於是跟他去見耶穌..
耶穌看著他說:你是約翰的兒子西門,你要稱為基伐..
又次日,耶穌想要往加里尼去,遇見肥納,.
就對他說:來,跟從我吧..
主持人:各位平安!.
我們第五次講約翰福音,.
在講這一段經文的內容之前,.
我們首先補充一下約翰福音沒有講的東西,.
三卷符類福音都提到耶穌受洗後,.
到曠野禁食,接受魔鬼的試探,.
然後才開始傳道和訓練門徒的生涯..
但約翰福音卻是直接從耶穌受洗,.
一跳就跳到他呼召門徒..
約翰福音的技術,起始於施洗約翰對耶穌的見證,.
然後就跳到耶穌呼召門徒,.
中間其實沒有記載很多東西,.
約翰沒有記載耶穌誕生的故事,.
沒有講到他的成長,.
連受洗都是用施洗約翰的益術方法間接地講出來..
在耶穌接受水禮之後,.
也沒有記述他曠野禁食四十晝夜的故事..
從第一章一節一直到三十七節,.
你可以看到約翰福音的技術基本上沒有什麼故事性,.
是一個神學論述,而不是一個傳記的題材..
就算記錄施洗約翰的說話都是他的獨白,.
不是跟任何人對話..
所以真正的故事是由今天開始,.
一章三十八節之後才開始..

$^{41}$值得一提,講一些常識給大家聽聽..
現在疫病蔓延,我們常用一個字叫quarantine,.
quarantine即是隔離,.
這個字原文是意大利文,.
叫做quarantina,就是四十天..
這個字是源自耶穌在曠野禁食四十晝夜的,.
四十天的自我隔離,遠離人群,.
在上帝面前獨處,.
這是靈魂得潔淨的途徑..
中世紀有疫病發生,.
從疫區來的人,或者曾經接觸過染疫者,.
都要首先被要求隔離三十天,.
三十天的意大利文就是quantina,.
後來就增加到四十天,就是quarantina..
威尼斯人為了防範鼠疫,.
要求從疫區來的船要停泊在岸邊四十天..
Quantina這個字一直沿用,.
所以四十天就是四十天隔離的意思..
後來四十天隔離的要求被廢止,.
但是這個字仍然繼續用..
所以今天我們在香港,.
回來香港要隔離的酒店都叫做quarantine,.
雖然我們沒有四十天,.
我們上個月剛剛減到,.
由七天減到三天,.
據說,今天看報紙說,.
下個月可能是零加七,.
普世寬騰,.
七天你覺得隔離的成本太高了,.
也有很多人覺得,.
困在酒店七天是極大的受罪,.
所以很好笑的,.
那些隔離酒店一定有個大大的字,.
我會把那個薩馬利亞防止自殺會的電話號碼給你,.
誰看不開就打電話去求救..
中世紀隔離四十天,.
你是不是覺得有點匪夷所思?.
約翰福音部有關quarantine的技術,.
我們回去看這一段聖經,.
一章三十八節,.

$^{81}$我們看看耶穌如何建立他的核心門徒團隊,.
耶穌基督第一個的門徒是誰呢?.
這兩本書有不同的技術,.
約翰福音記載安德烈和另一個人同時成為耶穌的第一批追隨者,.
他們如何成為耶穌的追隨者呢?.
我上次也有講過,.
不是耶穌親自呼召他們,.
在宏宏眾生裡面將他們拔出來,.
也不是他們自己心儀耶穌,.
很想主動跟隨耶穌,.
卻是出於施洗約翰的撮合,.
他們兩個本來是施洗約翰的門徒,.
施洗約翰大概自知命不久矣,.
不想繼續帶他們,.
然後將他們送到耶穌那裡,.
換言之,.
耶穌是接收了施洗約翰兩個的門徒,.
這兩個門徒,.
一個是安德烈,.
在他的轉介之下,.
他哥哥西門彼得也跟隨了耶穌,.
至於另一個沒有標示名字的門徒,.
多數人都相信這本書的作者是約翰,.
因為在整卷約翰福音裡面,.
常常提到一個沒有具名的門徒,.
通常可以假設是約翰自己,.
因為他不好意思說自己的名字,.
所以很多地方都是這樣的,.
他沒有說名字,.
你假設是約翰就可以了,.
約翰後來也將他的兄弟雅各,.
帶到耶穌的帳下,.
不過也有人相信,.
這個沒有標示名字的門徒是肥力,.
因為後面有提到肥力,.
肥力就將拿丹葉帶到耶穌的庚前,.
所以很有趣的,.
安德烈就介紹彼得,.
約翰就介紹雅各,.
肥力就介紹拿丹葉,.

$^{121}$一個介紹一個,.
一個帶一個,.
好像傳銷活動一樣,.
介紹的如果不是親兄弟,.
就是近親,.
這個就是我們經常說的.
Friendship Evangelism,.
友誼報道,.
最典型的運作方法,.
接受信仰的人,.
立刻成為傳播信仰的人,.
而他們的親朋好友,.
他們本來有的人脈關係,.
就成為了傳播的第一批對象,.
我們在當中,.
多數人信耶穌,.
都是因著別人的介紹,.
很少人主動上教堂,.
去尋找上帝,.
很少很少,.
很少人主動去研究,.
哪個宗教最可信,.
我不知道我們中間,.
有沒有一些剛回來教會,.
沒多久的日子,.
我相信你們多數都是,.
被你們基督徒的親朋好友,.
誣陷回來的,.
很少是主動回來的,.
好像我媽媽在病床裡面,.
呃頭信耶穌,.
不過我很相信,.
對不起,.
我心裡面,.
即使我少信也好,.
我相信祂賣我醬,.
餓了這麼久,.
於是點頭,.
我覺得Bingo成功了,.
但是老實說,.

$^{161}$祂對信仰有多少了解呢?.
其實我自己都心知肚明,.
是這樣的,.
你覺得很奇妙,.
很化學,.
但是真的很多老人家,.
信耶穌是這樣的,.
兒女餓得多,.
相信吧,.
你就當相信吧,.
你說祂對三位一體,.
對《士路信經》有多少了解呢?.
都不用問了,.
這些事情,.
正因為我們都是在親朋的邀請下,.
才接觸耶穌,.
所以我們除了不認識耶穌是誰之外,.
很多時候,.
特別是起初的時候,.
我們自己都不知道自己為何需要耶穌,.
耶穌在哪方面能夠幫助我,.
一切都是夢察察的,.
如果你們剛剛被邀請回來教會聚會,.
你們不僅僅不知道,.
我這個講完想說什麼,.
你也不知道自己想要什麼,.
又不知道自己想怎樣,.
正因為安德烈是夢察察的,.
那兩個被施拉約翰點去跟耶穌的門徒,.
都是夢察察的,.
所以我們會看到以下一段有趣的對白,.
安德烈兩個人來到耶穌那裡,.
但不知道跟耶穌說什麼,.
施拉約翰叫他們去跟耶穌,.
可能跟他們說了很多話,.
說了什麼耶穌是上帝的兒子,.
上帝的羔羊,.
後來的成了先來的,.
這些這麼抽象的事情,.
他們一句都聽不明白,.

$^{201}$太深奧了,.
不過師父吩咐他們跟耶穌,.
並且指出耶穌是萬人期待的彌賽亞,.
他們就來了,.
他們遠遠跟著耶穌,.
都不敢貼近耶穌,.
也不敢離開,.
他們不知道自己想怎樣,.
或者應該做什麼,.
真的不知道,.
耶穌很早就注意到有兩個人,.
閃閃縮縮跟在他後面,.
他沒有向他們發出呼召,.
卻是向他們問,.
你們要什麼,.
三十八節,.
你們要什麼,.
這個話可以解作你們在尋找什麼,.
或者像廣東話所說的,.
你們想怎樣,.
這是耶穌常常問的問題,.
就像兩個核子前來找他醫病,.
他們的需要其實已經非常明顯,.
但耶穌還是刻意問他們,.
要我為你們做什麼,.
瑪利福音20章32節,.
每個人都要主動去到耶穌面前尋求他的幫助,.
但是你尋求他的幫助之前,.
你一定要先問你自己,.
你想怎樣,.
在我們對耶穌有什麼期待之前,.
我們要先問我們對自己有什麼期待,.
我們的期待是塑造我們的前景,.
你尋求什麼,.
你想怎樣,.
你想實現什麼,.
十八世紀德國大哲學家康德,.
曾經說過,.
人類有三個最基本的問題,.
我能夠知道什麼,.

$^{241}$我要做什麼,.
我可以期待什麼,.
我能夠知道我要做什麼,.
我能夠期待什麼,.
中國人臨急抱佛腳,.
如今正在經歷病患,遭遇經濟困難,.
家庭關係破損的人,.
大概就知道自己需要什麼,.
因為那些是urgency,.
生活上有需要,.
身體方面有需要,.
感情上有需要,.
我們來耶穌那裡,.
耶穌是全能的上帝,.
祂能夠幫助我們解決問題,.
滿足我們的需要,.
除了有要解決的困難,.
即時的需要之外,.
或者我們懷抱著一些長遠的期望,.
希望升讀大學,.
可以找到一份合適的工作,.
找到一個可以同孩伴侶,.
儲夠錢上車供樓,.
這些不一定是最迫切的,.
但卻是我們的重大需要,.
一些important but urgent,.
我不會亂開旗票,.
我不確定上帝是否一定會按訂單,.
供應我們的需要,.
不過我們可以這樣說,.
我們將一切的想望,.
神聖的,世俗的,偉大的,渺小的,.
都拿到耶穌面前,.
這個願望最少構成一個真實的動機,.
讓我們在耶穌面前去掠,.
聖經是這樣應許的,.
應當一無掛慮,.
將你們所需要的告訴上帝,.
上帝所賜出人以外的平安,.
必在基督耶穌裡保守你們的心懷意念,.

$^{281}$聖經沒有說到上帝有求必應,.
我們心想事成,.
但聖經給了我們向上帝祈求的權利,.
所以有就說吧,.
最少這個構成我們信仰的動機,.
有不少基督徒說他無法經驗上帝的作為,.
原因是他們對自己沒有期望,.
對上帝也沒有期望,.
慣性回教會參加活動,聽道,.
也有一點點的事奉,.
不過一切都是default mode,.
習慣使然,.
我們不期待在生活裡遭遇上帝,.
聽上帝跟我們說話,.
我們不期待上帝動我們,.
改變我們的生活,.
沒有要求,沒有期待,.
我要說,你有恩典也看不見,.
有神蹟顯現,.
你也會視若無睹,.
我太太昨天做的麵包,.
加了一種新的材料,.
如果我如常慣性地吃麵包,.
吃了也不知道吃了甚麼,.
太太一口心機,你也看不見,.
我們太多時候如常工作,.
完全沒有看到內宗的變化,.
一些恩情,.
我問一個埔口神學院的姐妹,.
信主這麼多年,.
你有甚麼特別的屬靈經歷?.
那姐妹想了一會兒,繼續說,.
她說她跟丈夫以前的關係很惡劣,.
丈夫有婚外情,.
他們很多次漏離婚,.
自己也在一個嚴重的情緒困擾中,.
信主之後,夫妻兩個人關係大幅改善,.
家庭和睦,.
她自己的精神狀態也恢復正常,.
她說,這個對她來說是上帝奇妙的作為,.

$^{321}$很真實,也是切中她的需要,.
那姐妹說到一個最少她可以有渣拿,.
實現上帝的恩典,上帝的作為,.
你呢?你會說甚麼?.
耶穌基督問安德烈兩個人,.
你們要甚麼?.
這是他們跟隨耶穌面對的第一個問題,.
You have to articulate what you really want in life,.
你想要甚麼?.
不過,安德烈這個時候也不知道自己想要甚麼,.
或許他沒有任何urgent need,.
他也不覺得有甚麼迫切的需要,.
又沒有病,又沒有缺錢用,.
又沒有甚麼特別問題要等待處理,.
但也可能是他們也不認識耶穌,.
他們也不知道耶穌有甚麼實力,.
耶穌能夠幫助他們甚麼?.
例如你逛街,看到一間裝修得很特別,.
招牌很特別的店舖,.
你便好奇,便進去看看,.
店員問你想買甚麼,.
我立即問你,你有甚麼賣?.
我也不知道你有甚麼賣,.
我怎樣說我想買甚麼?.
我只是好奇,看看你的店舖是甚麼,.
他們也不知道耶穌有甚麼供應,.
我怎樣說他們有甚麼需要?.
所以安德烈傻乎乎地反問耶穌,.
老師,你住在哪裡?38節,.
這個回應似乎是答非所問,.
不過意思仍然很清楚,.
他們是計劃跟隨耶穌,.
那耶穌住哪裡,他們便跟隨到哪裡,.
他們是很想知道耶穌住哪裡,.
其實是很想知道他們將來住哪裡,.
很想知道他們將來會怎樣,.
你們住在哪裡,.
這個動詞,manual這個字可以譯成.
你通常在哪裡,你通常做甚麼?.
安德烈他們知道他們要離開.

$^{361}$是世若寒,離開生活多時的約旦河邊.
那個修道的群體,.
但他們不知道跟隨耶穌要去哪裡,.
跟隨耶穌要進入一個全新的世界,.
我們的生活會徹底改變,.
我們的價值觀,我們的世界觀會改變,.
我們的想法,我們的做法會改變,.
至於我們要變成怎樣,.
這個確實是我們在跟隨耶穌之前,.
信主的時候,我們是無法預知的情況,.
我們不一定知道自己的真實需要,.
也不知道耶穌能夠怎樣幫助我,.
無知,這個不是一些很可憐的情況,.
是我們信仰實踐一個上上一個重要的標記,.
信仰很多時候是包括未知的部分,.
所以我常說信仰是一個adventure,是一個歷險,.
好像亞伯拉罕離開吾爾,.
不知道將來要去的哪裡,叫做應許之地,.
跟隨耶穌學習信服上帝,過一個順命的生活,.
放棄虛假的自信和自由,.
放棄自編自導自演的生活,.
因為無知,因為我不知道,.
所以才需要信仰,信仰就是信任,.
信任是一種依賴的關係,.
我信對方,用善意期待對方,.
並且甘願承受失去自我保護,.
乃至傷害的風險,.
信任是放棄對現在對未來的充分掌握,.
將未來交付給對方,.
相信對方有意圖,也有能力給自己最好的將來,.
好像古代的社會,婦女出嫁,.
嫁到夫家,跟夫家同住,.
可以是另一條村,可以是另一個省,.
甚至是過剖新娘去另一個國家,.
妻子的命運從此跟丈夫綁在一起,.
同樣信耶穌就是跟隨耶穌,.
耶穌在哪裡,我們就去哪裡,.
耶穌所在的地方就是我們最終的歸宿,.
是我們最好可能的將來,.
耶穌沒有直接回答安德烈的提問,.

$^{401}$他只是說你們來旱,.
然後他們就去到耶穌所住下的地方,.
就是去住下來,.
跟隨耶穌就成為了他們整個信仰的開始,.
聖經提到這個時候是新正,.
就是下午四點之後,.
為什麼要特別提一個時間呢?.
我們不知道,.
有些識經者企圖去找一些特別的象徵含義去解釋,.
不過解釋來解去都是沒有什麼意思的,.
總之四點,就四點吧,.
我們來看,我們看到耶穌,.
我們看到耶穌的生活和事奉,.
我們在耶穌身上看到自己的需要,.
看到自己要尋找什麼,.
我們看到自己的將來,.
各位如果還沒信耶穌的朋友,.
你先看看你身邊的那些跟隨耶穌的人,.
看看耶穌在他們身上成就了什麼奇妙的恩典,.
我相信只要你們肯信耶穌,.
你們的生命都會同樣發生這樣生命改變的故事,.
各位弟兄姊妹,.
我們大部分人是信了耶穌的,.
我們都要親自去看看,.
不僅僅看看跟隨耶穌之後,.
他會在我們身上成就什麼作為,.
更加要看耶穌要藉著我們,.
去成就他託付給我們一個更偉大的作為,.
正如安德烈在耶穌身上看到的塞亞,.
這個超出了他個人的追求和期望,.
他只是想找一個師父,.
但耶穌卻是比他所期望的師父的角色大,.
是不是?.
跟隨耶穌會讓我們擺脫我們既有,.
包括我們自己所期望的生活框框,.
耶穌滿足我們的需要,.
耶穌改變我們的生命,.
耶穌承傳我們的現在,.
耶穌要帶引我們走向未來..
在早期,.

$^{441}$耶穌還沒有建立一個緊密的信仰群體,.
還沒有正式招收門徒,.
沒有拜師的儀式,.
所以跟隨耶穌的人,.
雖然有很多時間跟他住在一起,.
但他們偶爾還可以回家,.
他們還可以承擔家裡的責任,.
甚至有些人還在繼續他們原來的工作,.
原來的職業..
安德烈沒多久就回家,.
跟他的兄長彼得見面,.
他馬上跟彼得說,.
「我們遇見彌賽亞了」,.
才是一節..
上一次其實我已經講過,.
安德烈發現耶穌是彌賽亞,.
這是有些古怪的..
首先,耶穌跟安德烈,.
在約翰簡短的交談中,.
根本沒有提到自己的身份和使命,.
更加沒有提到彌賽亞這個稱呼..
所以安德烈是不可能這麼早,.
還沒見到耶穌做了什麼,.
就發現耶穌的身份的,.
是不是?.
按照馬太和其他婦女夫音的技術,.
彼得是第一個門徒,.
確立耶穌是彌賽亞的身份,.
而他要在.
在該撒尼亞肥納比,.
然後他才能宣認耶穌是基督,.
是不是?.
這是馬太和其他婦人十六章,.
十三到二十節的技術,.
這是很重要的技術,.
這個叫做Petrine text,.
這個叫做彼得經文,.
天主教是用這段經文去確立,.
彼得是教會的元首的一段經文,.
要去到馬太十六章才能確定,.

$^{481}$為什麼無端端約翰將安德烈.
變成第一個說耶穌是彌賽亞的人呢?.
有很多解經家認為,.
約翰夫音是錯配了時序,.
所以這是一個非歷史性的技術,.
約翰在最後寫的時候,.
是將時序插來插去,.
插錯了,.
這是Gospel的非歷史性的技術,.
我的解說就是,.
約翰不是亂錯配時間,.
安德烈自己沒有發現耶穌的身份,.
我相信是他的師父,.
師弟約翰告訴他的,.
雖然他是第一個說耶穌是彌賽亞的人,.
但不是他自己發現的,.
我們甚至相信,.
由於師弟約翰一早通水給兩個門徒,.
這兩個門徒也成為了耶穌的門徒,.
所以在耶穌的門徒中間,.
其實一直都有咬耳朵流傳著,.
耶穌是彌賽亞的說法,.
只不過因為這個身份,.
太過非同小可,.
耶穌日後的行事表現,.
又跟猶太人一直理解的彌賽亞形象,.
不是完全一樣,.
所以門徒在跟隨耶穌之後,.
一直都是半信半疑,.
直到彼得在馬第十六,.
做一個清楚的宣認,.
然後才確定耶穌是彌賽亞,.
聖經沒有提到安德烈跟彼得說,.
我們預見彌賽亞這句話的反應,.
沒有提到他是不是相信他弟弟的話,.
只不過提到彼得跟隨安德烈去見耶穌,.
沒有人因為別人的口才,.
別人的智慧去相信耶穌,.
我們永遠是將人帶到耶穌面前,.
耶穌親自去顯現,感化他,.

$^{521}$你開放你的心靈,.
直面耶穌,.
這個才是最重要的,.
當彼得兄弟來到耶穌面前,.
聖經也沒有詳細講到,.
耶穌跟彼得說了什麼道理,.
感召他成為門徒,.
不過就記述了耶穌為彼得做的第一件事,.
就是替他改一個名字,.
你是約翰的兒子西門,.
你要稱為基發,.
約翰這個字跟約拿是同一個字,.
這個是亞蘭文的short form,.
所以約翰的兒子,.
聖經有說是約拿的兒子,.
是一樣的字,.
基發這個字就是石頭盤石的意思,.
盤石這個字是用亞蘭文寫出來的,.
如果將它transliterate變成希臘文,.
就是彼得了,.
所以耶穌替彼得改一個名字,.
顯示他預見,.
要刻意安排彼得在他的救贖計劃裡有特殊的位置,.
這是為日後馬第法律16章鋪路的,.
是這樣的,.
耶穌對彼得不僅發出成為門徒的呼召,.
更加給了他一個很特別的期許,.
一個還沒說明白,.
但已經說了的使命,.
約翰福音沒有記載馬第法律16章這個事件,.
不過他在整本書最後一章,.
我們叫做附錄的一章,.
即是約翰福音21章,.
記述了耶穌對彼得的召命,.
這個不是給所有使徒一般性的召命,.
是特別對彼得說的,.
所以就算我們不是天主教徒,.
我們也要承認彼得在眾使徒中間有特殊的地位,.
摩洛因學,.
耶穌為彼得改名的事件,.

$^{561}$說明耶穌其實是很清楚認識他所揀選的呼召的門徒,.
他不僅看到他們的外貌,.
更加看到他們的內心,.
不僅看到他們的現在,.
更加看到他們的將來,.
前面說第一批門徒來找耶穌,.
不算是正式回應耶穌的呼召,.
耶穌沒有呼召,.
他們又不是回應呼召,.
純粹是司祭約翰的吩咐,.
在中間你會看到有很多偶然性的成分,.
大家都夢察察,.
不過在耶穌呼召彼得的事件裡面,.
你會看到其實在偶然性裡面,.
卻是有必然性在其中,.
耶穌其實早就知道每一個門徒,.
知道他們的過去,.
知道他們的未來,.
在這段聖經最後一部分,.
有說到再過一天,.
耶穌出發前往加利利,.
就遇到菲力,.
然後向他發出呼召,.
來跟從我,.
這是約翰福音裡面,.
耶穌第一次對人作出的呼召,.
看完這段聖經,.
好像很瑣碎,.
在這段聖經裡面,.
我們看到跟隨耶穌的行動有兩個特徵,.
第一,跟隨的人要知道自己想怎樣,.
你們要些什麼,.
第二,跟隨耶穌的初期,.
其實我們是不知道要去哪裡,.
你住在哪裡,.
兩個對白,.
你們要些什麼,.
你們住在哪裡,.
我要說,這兩個都是重要的特徵,.
日後你會看到耶穌一再重新跟隨耶穌這兩個特徵,.

$^{601}$第一,我們要些什麼,.
我們要知道我們要些什麼,.
我們期望從耶穌得著些什麼,.
跟隨耶穌不可以是一個貿然的決定,.
不可以是一個傻乎乎的決定,.
最少我們要知道我們想怎樣,.
但是另一方面,.
我們知道我們想怎樣是一件事,.
我們同樣要慢慢知道我們的主想怎樣,.
他要帶引我們去哪裡,.
他在哪裡,.
他要我們去哪裡,.
這個不是我們今天就我們自己的想望就能夠知道,.
我在這裡思考,.
我們的信仰很有趣,.
一方面,.
我們相信上帝的計劃,.
相信上帝的傳之,.
相信上帝的預定,.
我們在生活裡面確認冥冥之中有必然性,.
上帝預知一切,.
上帝確定我一生,.
但在每一個當下,.
我們又充滿著很多很多的未知之數,.
我們要做一個又一個的決定,.
我媽媽走之前差不多一個多兩個月,.
總共報了不知七次還是八次的案,.
每一次送入急症室,.
醫院打來,.
你做什麼無論是白天還是晚上,.
你決定立刻去,.
不要想,立刻去,.
去到之後沒事,.
然後又回家,.
然後又再打電話,又再去,.
我猜你們也試過,.
很多人也試過這些故事,.
一次又一次,.
總之是的,.
很多時候你覺得很勞累,.

$^{641}$很多時候你會問會不會又是假案,.
不過總之就是你不要想那麼多,.
你快點去,.
你不去將來可能會悔恨一生,.
去了之後,.
到時候你是一個交代,.
就是這樣,.
我們照顧年老的傷親,.
我們照顧家裡體弱多病的人,.
每一次你都覺得很勞累,.
重複又重複,.
沒完沒了的小故事,.
不過你不知道的是,.
當你過了一個段落之後,.
你回頭看,.
原來就是這麼多這些瑣碎的小決定,.
最少成就了一生的母子關係,.
成就了一生的某一個姻緣,.
某一段的感情,.
基督徒我們有永恆,.
有預定的觀念,.
這個為我們提供一個安穩,.
踏實的感覺,.
我生命不是由一連串的意外組成,.
我知道我安插在我的主,.
在他的計劃裡面,.
但是生活仍然充滿了片段的偶然性,.
偶然性和我做的每一個大大小小的個人決定,.
卻為我在生活裡面,.
特別是那些不斷重複沉悶的日子裡面,.
我去迫使自己做一個負責任的人,.
你一生的愛情故事,.
你一生的信仰故事,.
都是由很多細碎的決定組成,.
你知道上帝知道你的結局,.
所以我們求求就是每一個刻,.
今天我定義做一個怎樣的人,.
我定義怎樣善待我身邊的人,.
我定義怎樣用什麼態度去承擔我做的細碎決定,.
包括一會兒和爸爸媽媽去吃午飯喝茶的時候,.

$^{681}$定義過一個快樂的時光,.
定義讓他覺得開心,.
信仰是由很多細碎的決定所組成,.
我們知道上帝為我們鋪設了一生,.
為我預定了一切,.
不過我們不知道,.
我們日常生活常常說,.
贏在起跑線,.
在人間讀書做事,.
起點確實是很重要,.
不過信仰剛剛好就是相反的,.
起點不重要,.
信仰是贏在終點線,.
終點遠遠重要過起點,.
在前的可以變成在後,.
在後的會變成在前,.
聖經不停強調,.
我們的信仰之路一定要一走到底,.
堅持到底,.
然後才可以以勝利者的姿態進入永生,.
忍耐到底的才能得救,.
唯有致死忠心的才能得到生命的冠冕,.
啟示錄十四章十三節說,.
在主裡面而死的人有福了,.
你死的時候是基督徒,.
比你再生的時候是基督徒重要,.
比你一出生的時候重要,.
有些教會又幫你做嬰孩洗禮,.
死是怎樣,.
比你生的時候任何一點是怎樣重要,.
所以弟兄姊妹,.
我們已經有了不是很壞的起步,.
我們已經是一個基督徒,.
我們有人信主已經有十幾年,.
有幾十年了,.
求主幫助我們,.
堅持下去,.
謹守我們信仰的身份,.
並且在我們生活面前的,.
不要說過去,.

$^{721}$每一刻確立我們的身份,.
我們應有的大大小小的責任,.
直到我們見主的面,.
然後就完成我們一生的得救故事,.
信仰故事..
\newpage



\section{撒母耳記上 17:26-40}
\label{sec:PSHqqITZkIo}
\textbf{2022年10月2日崇拜直播｜陳明泉牧師｜以石為記:擊打巨人的五塊小石｜撒母耳記上17\_26-40}
\newline
\newline
連結: \href{https://youtube.com/watch?v=PSHqqITZkIo}{\texttt{https://youtube.com/watch?v=PSHqqITZkIo}} ~~~~ 語音日期: 2022-10-02
\newline
\newline
\hyperref[sec:uD1kVchCRPU]{\small{< < < PREV SERMON < < <}}
~
\hyperref[sec:index_chronic]{\small{[返順時目]}}
~
\hyperref[sec:x9IT18Te0VA]{\small{> > > NEXT SERMON > > >}}
\newline
\newline
撒母耳記上 17:26-40
\newline
\begin{longtable}{cl}
\hline
\hline
章節 & 經文 (和合本修訂版)\\
\hline
17:26 & \begin{tabularx}{0.7\textwidth}{X} 大衛對站在旁邊的人說：「若有人殺這非利士人，除掉以色列人的羞辱，他會怎樣呢？這未受割禮的非利士人是誰，竟敢向永生神的軍隊罵陣！」 \end{tabularx} \\ \\ \relax
17:27 & \begin{tabularx}{0.7\textwidth}{X} 百姓照同樣的話對他說：「若有人殺了那人，必這樣待他。」 \end{tabularx} \\ \\ \relax
17:28 & \begin{tabularx}{0.7\textwidth}{X} 大衛的長兄以利押聽見大衛與他們所說的話，就向他發怒，說：「你下來做甚麼呢？在曠野的那幾隻羊，你交託誰了呢？我知道你的驕傲和你心裡的惡意，你下來只是為了看戰爭！」 \end{tabularx} \\ \\ \relax
17:29 & \begin{tabularx}{0.7\textwidth}{X} 大衛說：「我現在做了甚麼呢？只是問一句話也不可以嗎？」 \end{tabularx} \\ \\ \relax
17:30 & \begin{tabularx}{0.7\textwidth}{X} 大衛離開他轉向別人，問了同樣的事，百姓也照先前的話回答他。 \end{tabularx} \\ \\ \relax
17:31 & \begin{tabularx}{0.7\textwidth}{X} 有人聽見大衛所說的話，就在掃羅面前報告；掃羅就派人叫他來。 \end{tabularx} \\ \\ \relax
17:32 & \begin{tabularx}{0.7\textwidth}{X} 大衛對掃羅說：「人不必因那非利士人灰心。你的僕人要去與他決鬥。」 \end{tabularx} \\ \\ \relax
17:33 & \begin{tabularx}{0.7\textwidth}{X} 掃羅對大衛說：「你不能去與那非利士人決鬥，因為你年紀太輕，他從小就是戰士。」 \end{tabularx} \\ \\ \relax
17:34 & \begin{tabularx}{0.7\textwidth}{X} 大衛對掃羅說：「你僕人為父親放羊，有時獅子來了，有時熊來了，從群中抓走一隻羔羊。 \end{tabularx} \\ \\ \relax
17:35 & \begin{tabularx}{0.7\textwidth}{X} 我就追趕牠，擊打牠，把羔羊從牠口中救出來。牠起來攻擊我，我就揪牠的鬍子，打死牠。 \end{tabularx} \\ \\ \relax
17:36 & \begin{tabularx}{0.7\textwidth}{X} 你僕人曾打死獅子和熊，這未受割禮的非利士人必像獅子和熊一樣，因為他向永生神的軍隊罵陣。」 \end{tabularx} \\ \\ \relax
17:37 & \begin{tabularx}{0.7\textwidth}{X} 大衛又說：「耶和華救我脫離獅子和熊的爪，他必救我脫離這非利士人的手。」掃羅對大衛說：「你去吧！耶和華必與你同在。」 \end{tabularx} \\ \\ \relax
17:38 & \begin{tabularx}{0.7\textwidth}{X} 掃羅把自己的戰衣給大衛穿上，將銅盔戴在他頭上，又給他穿上鎧甲。 \end{tabularx} \\ \\ \relax
17:39 & \begin{tabularx}{0.7\textwidth}{X} 大衛佩刀在戰衣上，試著走走看。因大衛沒有試過，就對掃羅說：「我穿戴這些不能走路，因為我沒有試過。」於是他脫下身上的這些軍裝。 \end{tabularx} \\ \\ \relax
17:40 & \begin{tabularx}{0.7\textwidth}{X} 他手中拿杖，又在溪中挑選了五塊光滑的石子，放在袋裡，就是牧人帶的囊裡，手裡拿著甩石的機弦，迎向那非利士人。 \end{tabularx} \\ \\
[1ex]
\hline
\hline
\end{longtable}
$^{1}$《十七章二十六至四十節》 - 全體委員會審議階段 - 第359至360頁.
「大衛問站在旁邊的人說:.
『有人殺者非利是人,.
除掉以色列人的恥辱,.
怎樣待他呢?』.
『這未受國禮的非利是人是誰呢?.
竟敢向永生神的軍隊罵陣嗎?』.
百姓照先前的話回答他說:.
『有人能殺者非利是人,.
必如此如此待他.』.
大衛的長兄以利亞.
聽見大衛與他們所說的話,.
就向他發怒說:.
『你下來作甚麼呢?.
在曠野的那幾隻羊,.
你交託了誰呢?.
我知道你的驕傲和你心裡的惡意,.
你下來特為要我,要看真戰.』.
大衛說:.
『我作了甚麼呢?.
我來豈沒有緣故嗎?』.
大衛就離開他,轉向別人,.
照先前的話而問:.
『百姓仍照先前的話回答他.』.
有人聽見大衛所說的話,.
就告訴了蘇羅,.
蘇羅便帶法人叫他來..
大衛對蘇羅說:.
『人都因那非利是人膽怯,.
你的僕人要去那非利是人戰鬥.』.
蘇羅對大衛說:.
『你不能去與那非利是人戰鬥,.
因為你年紀太輕,.
他自幼就作戰事.』.
大衛對蘇羅說:.
『你僕人為父親放羊,.
有時來了獅子,有時來了熊,.
從群中行一隻羊羔去,.
我就追趕他,擊打他,.
將羊羔從他口中救出來..

$^{41}$他起來要害我,.
我就揪著他的鬍子,.
將他打死.』.
你僕人曾打死獅子和熊,.
這未受國禮的非利是人.
向永生神的軍隊罵陣,.
也必將獅子和熊一半..
大衛又說:.
『耶和華救我脫離獅子和熊的爪,.
也必救我脫離這非利是人的手.』.
蘇羅對大衛說:.
『你可以去吧,.
耶和華必與你同在.』.
蘇羅就把自己的戰衣給大衛穿上,.
將銅盔給他戴上,.
又給他穿上海甲..
大衛把刀誇在箭之外,.
試試能走不能走,.
因為素來沒有穿慣,.
就對蘇羅說:.
『我穿大者,些不能走,.
因為素來沒有穿慣.』.
於是就脫了..
他手中拿著,.
又在溪中挑選了五塊光滑石子,.
放在袋裡,.
就是僕人帶來的囊裡,.
手中拿著脫石的基燃,.
迎拿非利是人聖經獨筆..
弟兄姊妹平安!.
平安!.
今天我們一起來看.
繼續以石為記的講道系列..
上星期,.
波木跟我們說,.
一個圖像,.
或者一個石柱,.
就成為了整個以色列經歷上帝帶領他們離開圍爐之地,.
進入應許之地的一個很重要的見證..
這個見證除了幫他們成為家庭裡面的教導之外,.

$^{81}$也都是提醒他們這個民族,.
整個的國家,.
上帝是看顧他們,保守他們,.
他們要帶著信心走面前的路..
今天我們就將這個鏡頭,.
我們就退後一段日子,.
我們就看一下整個以色列民族,.
或者我們這次將焦點放在他的君王身上,.
一顆石子可以成為見證,.
也都是成為了一個君王,.
一個新世代君王,.
站上他們歷史巔峰裡面的一個很重要的一個名證..
一顆石子竟然成就了一件很大的事,.
也都是成就了一個權力的更替..
今天我們看的《撒母耳記》上第十七章,.
就是在講大衛用這塊石頭來鋪排他,.
上帝的恩典淋到他身上,.
讓他怎樣成為一個新一代的領袖,.
成為新一代的君王,.
就由這一顆石子開始..
希望今天我們看的時候,.
我們找回我們生命裡面也有這一顆石子,.
究竟是什麼呢?.
我們今天一邊聽講道的時候,.
可能會幫我們整合得到..
我們一起求上帝幫助我們,.
我們一起禱告,同心領受上帝的話語..
阿爸天父願你在我們當中帶領我們,.
讓我們能夠在你的話語當中,.
找到很多生命的亮光,.
賜福給我們眾人,.
讓我們聆聽你話語的時候,.
讓我們的心被激勵,被振奮,.
又讓我們帶著謙卑受教的心在你面前,.
願你親自對我們講說話,.
幫助聆聽的每位弟兄姊妹,.
我們真的收到上帝對我們的照明,.
也幫助宣講的講得清楚,.
將你話語中心來到向弟兄姊妹展明,.
我們祈求你與我們同在,.

$^{121}$獻上禱告奉靠基督耶穌得勝的聖名祈求..
Amen.
今天我們看的《三寶衣機》上第十七章,.
是一個很重要的技術重點,.
不過我想大家了解一下,.
上文下你講的是什麼,.
在《三寶衣機》上第十六章,.
那裡是講上帝厭棄掃羅作王,.
在第十六章裡面特別講到,.
耶和華上帝直著先知,.
三寶衣,.
向掃羅來說話,.
第十六章第一節裡面講到,.
我厭棄掃羅作以色列的王,.
你為他悲傷到幾時呢?.
你將高遊城門遼閣,.
我差遣你,.
往伯利行的耶西嘉里,.
因為我在眾子之內,.
預備一個作王的..
在《三寶衣機》上第十六章,.
他首先在講一件事,.
當代在《三寶衣機》的日子,.
有掃羅作王,.
但是掃羅這個王似乎不合格,.
在他生命當中有很多地方,.
甚至是不聽上帝的吩咐,.
來管理以色列民,.
慢慢這個本來高人一籌的君王,.
就走火入魔,.
所以上帝就厭棄他,.
要燉他冬菇,.
讓他不會成為以色列的君王,.
所以《三寶衣機》上其實是在講,.
厭棄這個君王,.
然後改朝換代,.
改換一個新的君王,.
就在耶西,.
即是大衛的父親家裡,.
預備了一個新生的王,.

$^{161}$來接替他的王位..
當我們了解了這樣的情況之後,.
我們知道大衛他接下來的行事,.
其實已經是在講,.
藉著上帝透過三寶衣,.
其實他做了一些動作,.
在第十六章第六節開始,.
就講到三寶衣去找耶西,.
逐個耶西的兒子見,.
頭幾個又高大,.
那些大哥,.
三寶衣說不是那些,.
反而那個小孩,.
面色光紅,.
雙目清瘦,.
容貌俊美,.
三寶衣說我就要告他了,.
這是十六章第十二節,.
耶和華說,.
這就是他,.
你起來告他,.
三寶衣就用國女的高柔,.
在他朱輕中告了他,.
從這日開始,.
耶和華的靈就大大感動大衛,.
三寶衣起身回拉瑪去了..
在這一段的記載裡面,.
很簡單的,.
好像一個儀式那樣,.
但是在大衛都未出道,.
未成名,.
甚至是一個牧羊仔,.
但是三寶衣說,.
上帝要揀選這個人,.
將來成為新一代,.
取代蘇羅的君王,.
經文就是這樣記載,.
接下來的發展,.
第一件事就是第十七章裡面,.
就是講到有一些非利士人來罵陣,.

$^{201}$在經文裡面,.
第十七章,.
因為比較長,.
我們跳來讀,.
第十七章第四節,.
到第七節,.
就講到在非利士的境內裡面,.
從非利士營中出來一個土戰的人,.
這個人的名字叫做哥利亞,.
是加特人,.
身高六爪,.
零一虎口,.
頭戴灰盔,.
身穿海甲,.
甲重五千尺赫勒,.
腿裡面有銅護膝,.
兩肩之中背負銅戟,.
槍桿粗如織布的機軸,.
槍頭重六百尺赫勒,.
有一個拿盾牌的人在攤前面走,.
這個是撒布爾基上十七章四至七節裡面記載著,.
當以色列,.
其實他們雖然有自己的地方,.
但是外族或者迦南地裡面有很多不同的民族,.
經常會犯警的,.
其中一個宿敵就是非利士人,.
這裡記載的非利士人有一個人叫做哥利亞,.
他的高度有多少呢?.
六爪,零一虎口,.
兌換今天的道量行,.
這個人大概是,我查過,兩米八九,.
九尺,九點四八尺高,.
九尺多,我們這裡的樓底都不夠他進來,.
即是教會的樓底都沒有,.
這個加上假天花,進不了的狀態,.
即是他搭電梯都搭不到的,.
去到今天一個九尺高的一個巨人,.
在經文裡面,.
他特別的很細緻描繪著這個巨人,.
這個向以色列軍隊罵陣的某程度上的一個代表人物,.

$^{241}$他講到很細緻地描述他一些特徵,.
除了他高大之外,.
他身體有大量的裝備,.
戴銅盔,盔甲,新村海甲,.
盔甲的重量有五千射赫勒,.
五千射赫勒又是用今天的道量行,.
就是講百幾磅,就是那件盔甲,那件衣服,.
另外還有護膝,又有銅的裝備等等,.
你看到整個身,一個高大的巨人,.
身上還有穿上整套的戰備,.
到今天就是鋼鐵人,.
你看到一個高大的人穿上鋼鐵人,.
那個勢已經很強了,.
再加上他經文裡面講到鐵槍頭重六百射赫勒,.
那個重量都是值得留意的,.
不過他最想告訴你甚麼呢?.
是那個槍頭是用鐵製造的,.
牧師鐵製造很普遍,.
在大衛當代就不是了,.
有很多考古的發現,.
其實非利士人都頗兇的,在當代裡面,.
為甚麼這麼兇呢?.
因為他們能夠拿到一個技術,就是煉鐵的一個技術,.
所以他們能夠將一些很多的礦物提煉出來,.
成為了當時的鐵器,.
當以色列民族去打仗的時候,.
用石頭,用火把等等那些的時候,.
非利士人他們已經有精良的金屬,.
就是那個槍頭,.
他製造一個很硬的金屬,.
來刺穿敵人的敵方,.
所以在經文裡面,.
他想強調的就是,.
不單止他的身形巨大,.
而且更加講到他有精良的武器,.
非利士人有鐵的技術,.
基本上在牌面裡面,.
面對著一個這麼強的敵人,.
是不能打的,.
在以色列民族裡面,.

$^{281}$面對著一個這麼強的一個人出來討戰,.
一個這樣的罵陣的一個人,.
在經文裡面,.
即使這個這麼強悍的巨人,.
整天都在講,.
你們出來擺列隊伍做甚麼?.
我不是非利士人嗎?.
這是第八節講的,.
你不是數落的僕人嗎?.
可以在你們中間,.
選一個人出來,.
使他來到我當中,.
他如果贏了我,將我殺死的,.
我們非利士人做你的僕人,.
如果我贏了他,將他殺死,.
你們就要做我的僕人了,.
經文裡面就是這樣講到這個巨人哥利亞,.
帶著精良的裝備,.
來到去罵陣,.
當他做這些罵陣的時候,.
前線的軍隊,.
包括大衛的那幾個哥哥,.
因為跟了數落來打仗,.
他們有甚麼反應呢?.
經文裡面就特別講到,.
十一節,.
數落和以色列眾人聽見非利士人的這些話,.
就驚慌,.
極其害怕,.
害怕,.
輸定了,.
但又要擺一個隊形出來,.
你可以見到那種張力,.
其實大家都想著,.
當這個狂人發狂的時候要攻打,.
大家可能心裡想著要怎樣走,.
但樣子還要扮得很英明神武,.
因為是軍隊,.
但面對一個這麼強悍的敵人,.
有這麼精良的裝備,.

$^{321}$大家就很害怕,.
於是就有一個暗暗的指令就來了,.
即是說,.
數落說如果誰夠膽出去迎戰,.
就可以得到做駙馬爺的身份,.
但經文第一件事,.
未必講到那些群眾立即反應,.
作為描述的主調,.
在第十二節裡面,.
當這個強人這麼強的時候,.
鏡頭一轉就講到他,.
之前高立了的牧羊仔大衛,.
大衛聽他爸爸吩咐,.
要將那些食物給他,.
上去前線裡面有三個哥哥,.
即是以利亞,.
阿比,.
拿達和沙馬,.
大衛是拉仔,.
於是聽爸爸吩咐,.
將食物送上前線,.
希望可以供應他的哥哥所需,.
一路走的時候,.
大衛就聽到一個這樣的說法,.
他說如果哪一個人可以上前去,.
將這個非利士人擊殺,.
數落這個王,.
就會將自己的女兒嫁給他,.
成為妻子,.
而且在以色列當中免去父家納娘當妻,.
簡單來說,.
打贏了這個巨人,.
有這樣的指令,.
就可以做駙馬,.
還要不用納娘之外,.
還會賜你黃金萬兩,.
這些是坊間裡面,.
大衛一路送食物去比個哥的時候,.
沿路聽到百姓所關心的事情,.
你可以想像的,.

$^{361}$你投入到整個情景,.
三毛二寫的東西是很戲劇化,.
那些人在說,.
那個巨人多厲害,.
王說如果誰打贏,.
打死他,.
做駙馬之外,.
這次還有黃金萬兩,.
這次發達,.
那些人就說,.
那你去吧,.
那個人說,.
有錢沒命享,.
而且這麼大個人,.
一爭一的,.
我就好像炸猛死了,.
所以那些人只是說,.
沒人敢去上他錢,.
不過大衛聽到這些人的說法,.
在第十七章,.
二十六至二十七節,.
就記載著大衛和這些市民百姓,.
那個關注點不同,.
大衛問旁邊的人,.
有誰殺這個非利士人,.
除掉以色列人的恥辱,.
怎樣對他,.
怎樣待他,.
這些未受國禮的非利士人是誰,.
竟然敢向永生神的軍隊罵,.
百姓就將之前的說話,.
做什麼,.
殺誰,.
就可以黃金萬兩,.
做駙馬等等這些,.
大衛的關注點,.
和這些百姓的關注點,.
是完全不同的,.
百姓的關注點,.
老百姓想不用納糧,.

$^{401}$黃金萬兩,.
飛黃騰達等等,.
想這些,.
但大衛他的關注點,.
在經文裡多次又一次提及,.
以色列的恥辱,.
永生神的軍隊,.
他怎樣對這班未受國禮的非利士人,.
即是說未受國禮,.
即是和我們不同類,.
這班低下的,.
和我們這些神國子民,.
是不能揮的,.
不過在眼前的形勢,.
實際上,.
完全實力懸殊,.
但是在大衛的心裡,.
他在想的,.
他看到的,.
這班人雖然好像很懼怕,.
很害怕,.
但是他稱呼這一班,.
面對著,.
面臨大敵的,.
這班人是永生神的軍隊,.
和以色列的恥辱,.
當我們仔細去看的時候,.
你會發覺,.
大衛的關注點,.
從來不是在說金銀財寶,.
或者那些百姓心裡覺得,.
打贏了之後就得到的結果,.
反而他心裡面,.
他有一份對於自己的民族,.
對於上帝的那份認同,.
當一個外敵,.
在面前罵陣,.
來抵毀也好,.
或者去笑,.
癡笑的那些軍隊的時候,.

$^{441}$他不覺得是笑那些軍隊,.
那些人都廢廢地,.
不同大家有時看球,.
球員看得不好,.
然後說,.
這麼差勁,.
你覺得有距離,.
那個現在不是在笑他和他有距離,.
而是我們,.
他在笑我們,.
他覺得我們不濟,.
他心裡面有一份憤怒,.
而且他覺得自己有責任,.
來回應這個敵人,.
即使很多百姓在說很多的好處,.
但是在大衛心裡面,.
他看見的就是一班永生神的軍隊,.
面對著這個,.
人家罵陣裡面帶來的恥辱,.
在這裡面我們稍微了解,.
或者我們再進深來看,.
大衛的思考,.
大衛對時局的看法,.
對於我們今天很多弟兄姊妹基督徒,.
都是一個很重要的提醒,.
第一方面,.
當我們身邊有些基督徒,.
或者我們有一些弟兄姊妹,.
面對生命困難的時候,.
我們會不會和他劃清一條很遠的界線,.
覺得他這麼不濟,.
如果比起我,.
我聰明很多,.
自己當八卦的說,.
另一方面,.
覺得那件事不關自己的事,.
不過在大衛的關懷裡面,.
當軍隊,.
當他眼前的這個族,.
被人家罵陣的時候,.

$^{481}$他不覺得整件事和他是抽離的,.
他反而為了這件事,.
他自己覺得心裡面有一份憤憤不平,.
有一份覺得切膚之痛,.
甚至有一份義憤在心裡面,.
他希望自己可以回應,.
在一個這樣的時代,.
或者一個這樣的處境當中,.
他們會不會有一些力量,.
去回應一個這麼大的敵人,.
在經文裡面更加提到,.
當大家恐懼,.
沒有士氣,.
當整個軍隊表面是很有隊形,.
心裡面是恐懼,.
整個以色列民族,.
面對著一個這樣的大敵,.
當大家很意志消沉的時候,.
大衛和這班人很不同,.
看見的不是利益,.
或者是眼前有多困難,.
他看見的是上帝在他當中,.
他看見上帝沒有放棄他這個民族,.
他看見上帝要藉著一些人,.
面對著一個這麼大的巨人,.
他們都有力量去增盛,.
有機會和一些同道傾談,.
我不知道大家可能在生活裡面,.
都會遇到一些基督徒也好,.
或者不同的人士,.
不論留在香港或者移民,.
講到有一個很大的心理的想法,.
覺得香港人我們自己不行,.
我不知道你們有沒有這樣想過,.
覺得我們怎樣行,.
又被新加坡搶了,.
經常都說這些,.
或者我們又沒有這些東西,.
我們不行了,.
什麼都不行了,.

$^{521}$當然客觀環境我們真的面對著很多困難,.
但心底裡面,.
我們面對著一個這麼大的難題,.
我們是否純粹採取一種捱打的狀態,.
在我們面對著一個這麼大,.
這麼困難的處境裡面,.
是否我們連一點點回應的能力都沒有,.
我們除了害怕,.
除了沮喪之外,.
我們還有什麼可以做,.
有一句名言是這樣說的,.
可能我們直播去到外國,.
我也要說,.
有些移民到外國的弟兄姊妹,.
人家第一件事問她們,.
settle down 了沒有,.
那些移民了的弟兄姊妹說,.
我們已經settle了,.
settle了之後,.
我們就很down了,.
明白我的意思吧,.
真的他們想著移民了,.
或者處理了之後,.
就會解決,.
不過原來當他們去到轉換環境,.
生活裡面都有很多的低沉,.
未必一定移民的,.
香港裡面都是,.
我們都有一種很低下,.
很屈不得志的那種情感,.
不過求主幫助我們,.
在這裡我們看到,.
原來這個時代裡面,.
上帝預備了大衛這個這麼有趣的人,.
當大家覺得不行的時候,.
誰向永生神的軍隊來罵陣,.
誰來為我們這個民族帶來羞辱,.
我們有力量可以面對的,.
所以當大衛一路一路說的時候,.
在第17章33節記載著,.

$^{561}$大衛就要向自薦,.
來向蘇羅說,.
我要和這些非利士人爭戰,.
我漏了一句,.
我想讀給大家聽,.
大衛的豪情壯語,.
人都不必因非利士人膽怯,.
你的僕人要去與那非利士人戰鬥,.
這是32節裡面說的,.
在大衛和蘇羅當大家很害怕的時候,.
大衛說,.
個個都不用怕那些非利士人,.
現在你的僕人是我,.
他很謙稱自己是僕人,.
我現在就要和這些非利士人,.
和這個非利士人來戰鬥,.
當他說這件事的時候,.
33節,蘇羅對大衛說,.
你不能夠和這個非利士人決鬥,.
你還這麼年輕,.
他們從小就是戰士,.
某程度上是年輕人,.
人家小時候已經是打仗的,.
你小時候去牧羊,.
你年紀太輕了,.
大衛對蘇羅說,.
你僕人為父親放羊,.
有獅子來到的時候,.
有時候熊來到的時候,.
在群羊當中捉了一隻羊羔,.
我就追趕他,.
擊打他,.
在羊羔從他的口裡面救出來,.
他起來攻擊我的時候,.
我就抽著他的鬍鬚,.
打死他,.
你僕人曾經打死獅子和熊,.
那些未受國禮的非利士人,.
好像獅子和熊一樣,.
因為他向永生神的軍隊罵陣,.

$^{601}$大衛又說,.
又說救我脫離獅子和熊的爪,.
他必救我脫離這非利士人的手,.
嘩,豪情壯語,.
一個牧羊仔,一個小孩,.
原本是送飯的,.
當他想送飯的時候,.
看到這些狀況,.
那些哥哥說,.
你來幹什麼?.
你來這裡八卦,.
看燈,看打架?.
回去吧,.
你那幾隻羊看好了沒有?.
之前你哥哥是這樣跟他說的,.
但是大衛不理他們,.
直接跟素羅對話,.
獅子和熊夠兇猛,.
牧羊的時間,.
上帝已經訓練他,.
上帝昔日怎樣去救他,.
如今他都要知道,.
上帝救他脫離這個非利士人的手,.
在大衛的心目中,.
他不否定自己的能力,.
他不否定自己的經驗,.
可以戰勝獅子和熊,.
不過他最心底裡面,.
有一份這麼強烈的自信,.
因為他知道,.
耶和華救他脫離這些兇惡,.
耶和華上帝與他同在,.
給他有力量去打贏,.
昔日的獅子和老虎,.
接下來的非利士人,.
上帝的靈都是這樣,.
給他有這份的確信,.
去戰勝他們,.
當然聽到這些豪情壯語,.
有時候我們會很害怕,.

$^{641}$年輕人真的有一份激情,.
覺得這個世界完全可以改變,.
我們有時候會淋一淋年輕人的冷水,.
不過我想講一個事實給大家聽,.
原來當我們覺得一些很穩陣的事,.
在年輕人身上講的時候,.
原來時代真的會被這些說話改變,.
早幾輪我有機會跟一些,.
將要入神學院蒙照的,.
有些年輕人聊天,.
他跟我說,.
我問他為什麼你打算讀神學呢?.
他說我都不知道,.
我不知道可以做什麼,.
然後我說你不知道做什麼,.
打算讀神學?.
這裡不是職業先修,.
不過他說他面對很多艱難,.
情緒很苦,.
很悲哀,.
他想想可以靈修幫他,.
他更加說一些說話,.
說到學會處理一些自己的不濟,.
學Richard Nabor的禱告,.
這樣就很豐厚了,.
我聽見那是一個二十多歲的年輕人,.
我聽完他說,.
他說你說的話很有深度,.
不過不是你的年紀說的,.
在你的年紀裡,.
我好像你這個年紀,.
我是不識天高地厚,.
我覺得這個世界我可以有力量改變,.
但是你現在,.
你二十多歲說一個六十多歲的說話,.
你是一個老人精,.
老人精帶來的結果是什麼呢?.
你都沒有青春過,.
你已經蒼老了,.
今天這個世代,.

$^{681}$很多人都沒有經歷青春期,.
已經一條條的成為一個蒼老,.
沒有動力,.
失去了盼望的一個人,.
今天我們香港,.
我們一班弟兄姊妹,.
會不會我們都在經歷一個這樣的狀況,.
未經歷青春,.
沒有經歷過年少的鬥志,.
很快已經消磨了自己,.
坊間裡面叫他想清楚,.
但到最後其實是在說叫他放棄,.
做一個沒有改變現實的那種鬥志和期盼,.
留下大衛這個人給我們看到,.
這個年輕人,.
他沒有妄自菲薄,.
他知道上帝與他同在,.
他知道在他的能力範圍裡面,.
他可以改變,.
即使如臨大敵,.
他都不會純粹用懼怕來回應,.
眼前的那種令人很恐怖,.
很驚惶的處境,.
反而他自告奮勇,.
他自己向蘇羅來說,.
他要與這個非利士人爭戰,.
當蘇羅聽見大衛說完這番說話之後,.
蘇羅雖然說是被取代的一個君王,.
但他似乎都是眷顧這個年輕人的,.
三十八節,蘇羅就把自己的戰衣給大衛穿上,.
將自己的銅頭盔給他戴上,.
又給他穿上海甲,.
大衛把刀誇在戰衣外,.
是是能不能走,.
因為蘇羅沒有穿冠就對蘇羅說,.
我穿戴這些不能走,.
因為蘇羅沒有穿冠,.
於是擲脫了,.
很簡單的記載,.
但很有意思,.

$^{721}$蘇羅看到這個年輕人竟然這麼有自信,.
又有些check record,.
獅子老虎,紅癮都打死了,.
於是讓他試,.
不過怕他不夠裝備,.
蘇羅將自己最漂亮的,.
做君王的裝備一定是最漂亮,.
將他的頭盔,海甲和刀,.
給大衛戴在身上,.
但你記得勝經說,.
蘇羅是高人家一個頭,.
在他當代,一個高大的,.
雖然高利亞九呎高,.
蘇羅都不矮,都有七呎左右,.
將一個七呎的盔甲給大衛五呎左右,.
擺在身上,.
整件事大衛試穿上去,.
那把劍很漂亮,.
不過拖著他們,.
如何行,如何打仗,.
於是幾個字就說到大衛的做法,.
摘脫了,.
全部脫下,.
將所有精良的裝備,.
一一放下,.
大衛為何這樣做呢?.
他直接說,.
因為素來沒有穿慣,.
穿不慣,.
他打獅子老虎不是穿這些裝備,.
所以全部放下,.
不容讓這些東西成為自己的阻攔,.
這些東西是最好的東西,.
蘇羅最厲害的軍王盔甲,.
還有最漂亮的配件,.
不習慣,不穿了,.
不合適,所以不用,.
在這裡,.
經文裡面似乎給我們一些比較,.
菲律賓人穿著鋼鐵人,.

$^{761}$九尺多高,.
穿著整套盔甲,.
蘇羅很害怕,.
但是他也穿著整套盔甲,.
整裝待發,.
但是不敢打仗,.
相反來說,這個牧羊仔,.
大衛什麼都沒有,.
到最後有些什麼呢?.
在經文裡面四十節,.
手拿著杖,.
在河裡挑選五粒光滑的石子,.
放在袋裡,.
就是牧人帶的囊裡,.
手中拿著掉石的機緣,.
就去迎那菲律賓人,.
整件事是一個很大的諷刺,.
那些整裝待發的在那裡罵陣,.
全新裝備,最漂亮的了,.
想著要打勝仗的了,.
想著要贏,.
又或者蘇羅要用最強的裝備,.
來迎戰這個巨人,.
不過大委身上用最簡單,.
他最用慣的東西,.
來面對一個最強悍的敵人,.
在經文裡面他特別說到,.
大衛這個小子去打仗,.
他不是用人家覺得好的裝備,.
用人家覺得好的東西,.
用人家覺得最強的東西,.
放在自己身上,.
就令自己變得強勁,.
而是用他素來用慣的東西,.
他一直怎樣去對抗獅子,.
紅人的那些方法,.
如今他就在這一刻裡面,.
就用這些東西,.
來擊退這個巨人,.
在經文裡面,.

$^{801}$如果將這一段經文,.
拍為一套戲,.
我想說是不好看的,.
為什麼呢?.
因為開頭說的那些裝備,.
那些對峙,.
場面很大,.
但到打的時候,.
兩句就搞定了,.
《三武二記》是真的,.
令讀者真的很好笑,.
幾句話已經搞定了這場仗,.
經文怎樣說呢?.
48至51節,.
4節經文,.
我讀給大家聽,.
「那非利士人起來迎戰大衛,.
走近前來,.
大衛急忙往戰場,.
迎向非利士人跑去,.
大衛伸手入囊中,.
從裡面逃出一塊石子來,.
用勒石機鉛勒出去,.
擊中非利士人的前額,.
石子進入額內,.
他就撲倒,.
面伏於地,.
這樣大衛用機鉛和石子勝了那非利士人,.
擊中他,把他殺死,.
大衛手中沒有刀,.
大衛跑去站在那非利士人身旁,.
把他的刀從肖中拔出來殺死他,.
用刀割下他的頭,.
非利士眾人看見他們的勇士死了,.
就都逃跑.」.
幾節經文,.
我以為大家在私塞,.
大家在拉鬍子,.
有些埋身肉搏,.
好像成龍片一樣,.

$^{841}$你以為沒有,.
大衛拿了五顆石,.
記不記得我剛才說過,.
他的裝備,一支棍,.
五塊石頭,.
那個非利士人就說,.
「你打狗呀?拿支杖來跟我打仗?」.
在取笑他,.
大衛不理那麼多,.
打到他前線要衝出去的時候,.
在他那五塊石頭裡拿一顆石,.
就用「甩石機元」就發出去了..
在這裡再值得我們去想,.
大衛那個信心大到,.
在拿那幾顆石,.
他在河裡挑選五塊光滑的石子,.
如果我作為打仗之前,.
我一定想那些是子彈,.
那些「甩石機元」的石仔,.
如果我真的想面對一個那麼強悍的敵人,.
如果我子彈多一點,.
我就穩陣一點..
但是他只拿了五顆石,.
五顆他覺得已經夠了..
然後到最後,.
真的打仗的時候,.
他只用了五分一,.
只用了第一顆石頭,.
用一個「甩石機元」,.
就一下擊斃了這個非利士人..
當這個非利士人撲倒在地上,.
死了,.
大衛手上,.
剛才說的「素羅比」那把劍他沒有帶,.
沒有劍,.
斬他的頭下來,.
他要走到那個非利士巨人那裡,.
拔掉他的配劍,.
一下子就將這個非利士人的頭斬了出來,.
然後就示眾了..

$^{881}$整件事,.
如果你說真的打仗,.
你真的打鬥場面,.
欠奉,.
因為薩姆爾記他要告訴你的,.
不是說那個廝殺有多血腥,.
或者要做些什麼,.
他想說整件事是什麼呢?.
耶和華上帝使人憎勝,.
耶和華上帝令到這個小子,.
好像沒有什麼裝備,.
沒有什麼精良的武器配搭之下,.
他都可以戰勝一個.
令到人聞風喪膽的巨人..
一個很簡單,.
直接了當的一擊,.
就已經成為了這個巨人最後的一擊,.
已經成為了他葬身在沙場裡面的一粒機緣的石子..
這塊石子代表著上帝與大衛同在,.
藉著薩姆爾的高立,.
這個年輕人知道他將會成為一個新一代的領袖,.
但心裡面不只是純粹沾沾自喜,.
而是他為著整個民族,.
為著上帝的同在,.
他有一份莫名其妙,.
一份與世代不同的信心和鬥志,.
藉著這塊石子將這個巨人擊退..
一粒石子表面上好像一粒子彈,.
一個武器,.
但背後裡面卻是隱藏著上帝的計劃,.
這塊石子背後也代表著.
這個牧羊人其實上帝早就預備他,.
讓他成為明日之君..
我們去到第17章第45至47節,.
某程度上是我們今天結束的一個結語..
大衛對非利士人說:.
「你來攻擊我,是靠著刀槍和矛,.
但我來攻擊你,是靠著萬軍之耶和華的名,.
就是你所辱罵帶領以色列軍隊的神.」.
47節:.

$^{921}$「有使這類的傳媒眾知道,.
耶和華使人得勝,.
不是用刀用槍,.
因為戰爭存在乎耶和華,.
他必將你們交在我們手裡.」.
在大衛的心目中裡面,.
真正的致勝被決,.
不是說外在的裝備或敵人有多強悍,.
在大衛心目中,.
真正得勝的是來自上帝他自己..
在《撒慕爾記》裡面,.
他有很多篇幅,.
他講到這個巨人有多恐怖,.
講到眼前以色列民族裡面面對著這個宿敵,.
有多精良的武器,.
有多巨大,.
多令人生畏,.
其實這是經文裡刻意要告訴讀者,.
真是很恐怖的,.
很令人害怕的..
這種呈現,.
同樣在我們讀聖經的時候,.
我們都可以在腦裡想像,.
每個人生命裡面都有一個巨人,.
有一個這樣的敵人,.
不過今日我們的巨人未必是一個九尺高的男人,.
來到在我生命裡面罵陣,.
但是我人生裡面可能有些問題,.
就好像這個巨人一樣,.
在我們面前威嚇我們,.
令到我們驚惶失措,.
令到我們不知進退,.
甚至令到我們自己看不起自己,.
我們覺得自己已經是一敗塗地,.
他都未用打,.
你已經輸了在你的心靈裡面,.
是一種這樣的形態出現,.
不過上帝為我們預備了這些石頭,.
上帝這些預備是甚麼呢?.
這些石頭代表著甚麼呢?.

$^{961}$可能在大衛自己自幼牧羊的時候,.
他去和獅子,紅癮搏鬥,.
他都沒有想過自己這些能力.
可以將來成為君王的一個技能,.
他學會了,.
他面對過這些處境,.
就在當下裡面,.
上帝用他這些訓練,.
鋪排他未來所走的每一步,.
這些不起眼,.
或者生活日常這些能力,.
其實可能久而久之,.
就成為了我們戰勝巨人,.
一些很重要的力量,.
亦都這些能力訓練到我們,.
我們其實不需要面對生命裡面很多問題,.
不需要怕,.
一個比較輕和的例子,.
這樣說好像很厲害,.
輕和一點,生活一點,.
早一陣子我女兒在班裡面開始要演講,.
然後就很害怕,.
出手汗,又腳震,又睡不著,.
她說:Daddy你經常說話,.
你怎樣克服這些害怕?.
我就說:我的東西怎樣適合你?.
你怎樣克服我的恐懼?.
就好像素羅那套規格,.
我不可以套給你身上,.
不過那一段我告訴她,.
其實你都有足夠能力可以講的,.
你想想你在皇寵裡面,.
你都曾經對著百多二百人,.
至少一個祖堂,五堂,.
對著百多二百人唱詩歌,彈琴,.
最醜怪的事,.
最需要膽量的事你都做了,.
對著同學幾十個,.
大家跟你這麼高,.
你講幾句話你就害怕,.

$^{1001}$其實上帝早已預備了你,.
在教會裡面訓練了你,.
你已經有這樣的能力,.
來應對一些大一點的挑戰,.
只不過場景不同了,.
那些人不同,.
但是用的方法是一樣的,.
我女兒笑著笑著,.
她說:我都害怕,.
不過你講完這句話,.
好像心靜一點,Daddy,.
然後我另外一個例子,.
有機會我和一個商界精英,.
來跟他說他夫妻之間的關係,.
因為這個商界精英,.
他做生意很厲害,.
做工作,職場很厲害,.
我說:搞不定,.
不知道怎樣跟老婆聊天,.
然後我就說:.
在你公司裡面.
有沒有對過一些最難對的客人?.
然後他說:當然有,.
他們經常都是這樣對,.
你猜老婆難受一點,.
還是客人難受一點?.
我說:最難搞的客人你都搞得定,.
沒理由搞不定你老婆,.
然後那個商界精英說:.
對,我從來沒有想過,.
原來對老婆就好像待客一樣,.
我說:對,你搞得定她,.
你又用你對老婆的方法來對客,.
你又可以,.
兩樣東西是可以的,.
不過不是要哄她,.
而是你要告訴她,.
你要對老婆講真誠,.
你想有誠意,.
希望待她好,.

$^{1041}$這樣來溝通,.
頂尖我要講什麼呢?.
在我們人生裡面,.
其實上帝已經為我們預備了很多,.
好像大衛一樣,.
擊退獅子,老虎的那些小石,.
那粒凹石機緣,.
日復如日,.
我們不要忘記上帝已經為我們.
預備了這些能力,.
面對著生命裡面有很多困難,.
我們第一件事不是害怕,.
而是想一想,.
在我這個牧人的袋裡面,.
我有多少彈藥,.
是上帝為我預備,.
昔日了這些能力,.
當在這些困難的日子裡面,.
我們就將這些能力改換,.
我們用同樣的思維,.
來克服眼前這些困難,.
今日這一段經文,.
希望大家繼續回去,.
再一路回答,.
《三寶二記》上第十七章,.
我相信當我們去細味.
這一段經文的時候,.
特別是這個好像沒有什麼盼望,.
沒有什麼力量,.
在大家很絕望的世代裡面,.
找回耶和華上帝在我們生命裡面,.
為我們所預備的恩典,.
也找回我們生命裡面,.
那顆光滑的石頭,.
迎戰這個我們覺得有大有難的巨人,.
求主繼續幫助我們..
法定人數不足.
\newpage



\section{以賽亞書 28:14-22}
\label{sec:x9IT18Te0VA}
\textbf{2022年10月9日崇拜直播｜陳冠成牧師｜以石為記:試驗過的房角石｜以賽亞書28\_14-22}
\newline
\newline
連結: \href{https://youtube.com/watch?v=x9IT18Te0VA}{\texttt{https://youtube.com/watch?v=x9IT18Te0VA}} ~~~~ 語音日期: 2022-10-10
\newline
\newline
\hyperref[sec:PSHqqITZkIo]{\small{< < < PREV SERMON < < <}}
~
\hyperref[sec:index_chronic]{\small{[返順時目]}}
~
\hyperref[sec:862eUMKnVnc]{\small{> > > NEXT SERMON > > >}}
\newline
\newline
以賽亞書 28:14-22
\newline
\begin{longtable}{cl}
\hline
\hline
章節 & 經文 (和合本修訂版)\\
\hline
28:14 & \begin{tabularx}{0.7\textwidth}{X} 因此，你們這些傲慢的人，就是管轄住耶路撒冷這百姓的，要聽耶和華的話。 \end{tabularx} \\ \\ \relax
28:15 & \begin{tabularx}{0.7\textwidth}{X} 你們曾說：「我們已與死亡立約，與陰間結盟，不可擋的鞭子揮過時，必不臨到我們；因我們以謊言為避難所，靠虛假來藏身」； \end{tabularx} \\ \\ \relax
28:16 & \begin{tabularx}{0.7\textwidth}{X} 所以，主耶和華如此說：「看哪，我在錫安放一塊石頭作為根基，是衡量的石頭，是寶貴的房角石，穩固的根基；信靠他的人必不致驚恐。 \end{tabularx} \\ \\ \relax
28:17 & \begin{tabularx}{0.7\textwidth}{X} 我以公平為準繩，以公義為鉛垂線；冰雹必沖去謊言的避難所，大水必漫過藏身之處。 \end{tabularx} \\ \\ \relax
28:18 & \begin{tabularx}{0.7\textwidth}{X} 你們與死亡所立的約必廢除，與陰間所結的盟不得堅立；不可擋的鞭子揮過時，你們必被踐踏。 \end{tabularx} \\ \\ \relax
28:19 & \begin{tabularx}{0.7\textwidth}{X} 每逢它揮來，必將你們擄去；每早晨它必揮過，白晝黑夜都是如此。明白這信息的都必驚恐。」 \end{tabularx} \\ \\ \relax
28:20 & \begin{tabularx}{0.7\textwidth}{X} 床榻短，人不能伸展；被子窄，人無從裹身。 \end{tabularx} \\ \\ \relax
28:21 & \begin{tabularx}{0.7\textwidth}{X} 耶和華必興起，像在毗拉心山，他必發怒，如在基遍谷；為要做成他的工，就是非常的工，成就他的事，就是奇異的事。 \end{tabularx} \\ \\ \relax
28:22 & \begin{tabularx}{0.7\textwidth}{X} 現在你們不可傲慢，免得捆綁你們的繩索更結實，因為我從萬軍之主耶和華那裡聽見，在全地施行滅絕的事已定。 \end{tabularx} \\ \\
[1ex]
\hline
\hline
\end{longtable}
$^{1}$今日的經訓是《以賽亞書》第十四章.
不是 第二十八章十四節至二十二節.
在唐用的聖經九百四十四頁.
是舊約的《以賽亞書》.
請聽我讀出.
所以你們這些洩漫的人.
就是核管著在耶路撒冷這百姓的.
要聽耶和華的話.
你們曾說.
我們與死亡立約.
與陰間結盟.
敵軍如水漲漫經過的時候.
必不臨到我們.
因我們以謊言為被鎖.
在虛假以下藏身.
所以主耶和華如此說.
看啊!我在石安放一塊石頭作為根基.
是試驗過的石頭.
是穩固根基.
寶貴的防角石.
信靠的人必不著急.
我必以公平為準繩.
以公義為善途.
兵駁必衝去謊言的被鎖.
大水必漫過藏身之處.
你們與死亡所納的約必然廢掉.
與陰間所結的盟.
必納不住.
敵軍如水漲漫經過的時候.
你們必被他踐踏.
每逢經過必將你們擄去.
因為每早晨他必經過.
白晝黑夜都必如此.
明白傳言的必受驚恐.
原來藏塔短使人不能書身.
鐵鍋窄使人不能遮體.
耶和華必興起.
葬在比拉森山他必發怒.
葬在雞片谷.
好作成他的功.

$^{41}$就是非常的功.
成就他的事.
就是奇異的事.
現在你們不可洩漫.
恐怕捆你們的綁索更結實了.
因為我從主萬軍之耶和華那裡聽見.
已經決定在全地上施行滅絕的事.
聖經讀不願神賜福常讀句話與並且遵行的人.
請姐妹平安.
相信剛才我們所聆聽的經文.
特別提到「防國石」這三個字.
我想我們對這三個字不陌生.
特別留意新教聖經常常出現.
與「防國石」相關的教導.
我們明白在建築層面來說.
「防國石」就是整棟建築物裡.
最大塊最堅固的石頭.
並且它是第一塊.
安放在建築物角落的基石.
用來做什麼呢?.
除了根基外.
當我放了這塊防國石作為基石的時候.
是用來提醒和指導建築工人.
之後每一塊石頭都要靠著這塊防國石堆疊.
要與這塊防國石對準和拍齊.
所以其實「防國石」不單只是建築物的根基.
其實它同時制定了整棟建築物的發展方向.
向著哪裡來建.
它的定位到底是怎樣.
這是我們對「防國石」很熟悉的地方.
不過同樣我們剛才聽這段經文.
其實也有一些我們陌生的部分.
我們不是太熟悉當時這段經文的背景.
所以需要花點時間來做一些了解.
其實講的是公元前八世紀的後期.
北角以色列已經被亞述帝國消滅了.
列王記夏解釋.
北角以色列之所以敗亡的真正原因.
是所有的君王都行耶和華眼中看為惡的事.
包括什麼罪呢?.

$^{81}$包括在宗教上他們得罪了上帝.
又敬拜上帝.
但同時又拜偶像.
我們說是心懷異議.
這是第一方面.
此外也有道德上他們得罪上帝.
在先知書裡很多時責備他們屈往正直的人.
以及沒有看顧孤兒寡婦.
甚至欺壓弱勢的社群.
第三方面他們在政治上也得罪上帝.
就是當面對外敵入侵時.
他們應該第一時間尋求上帝的指引.
不過很可惜.
他們首先考慮的是如何用外交和結盟方法解決問題.
這是他們三方面得罪上帝的地方.
其實上帝已經多次多方.
差派不同的先知來警告他們.
提醒他們要回頭.
不過正如先知書裡說.
高傲自誇的以色列.
總是勸不聽.
甚至以為自己仍然是形勢大好.
於是繼續用人的方法.
用人的想法來主導整個國家的發展.
當上帝的子民不再按照上帝的心意.
來作為建立國家的防國石時.
上帝會怎樣?.
正如你聽過的新聞.
當你發現那棟樓的樁腳有問題.
那你會繼續住下去嗎?.
繼續住下去嗎?.
不行.
在上帝的角度.
既然有問題的話.
就拆掉它.
重新建造.
總比住了之後倒塌的傷亡更嚴重.
所以當北角以色列已經偏離了.
上帝的心意並且屢勸不聽的時候.
上帝就透過他怒氣的象.

$^{121}$亞述帝國出手.
認罰管教他的百姓.
正如剛才我們看28章第14至21節.
28章一開始是以色列對北角的宣判.
第一至第二節.
在程序表上也印了第一段經文.
「禍災以法連的酒徒.
住在肥美谷的山上.
他們心裡高傲.
以所誇的為冠冕.
猶如樟蟬之花.
看!主有一大能大力者.
將一陣冰雹.
將毀滅的風暴.
將將日的大水.
他必用手將冠冕摘落於地」.
這裡是上帝透過亞述管教以色列.
將他們滅亡.
眼見自己的兄弟國.
以色列已經出事了.
身為猶大的兄弟國應該怎樣做?.
到底是引以為鑒.
避免重蹈覆轍得罪上帝.
又或者耶路撒冷的領袖.
見著北角已經敗亡.
他們應該選擇.
用什麼來作為建立國家的防國石?.
到底是繼續用人的方法.
還是順應上帝的指引.
帶領百姓面對亞述將來的入侵?.
這正是我們今天所看的經文.
裡面提到南國猶大王希西加.
他要面對的處境.
事實上我們對希西加王也有印象.
他25歲就登基作猶大王.
在耶路撒冷作王29年.
他是一個很好的君王.
因為王紀形容他.
行耶和華眼中看為正的事.
效法他的先祖大位一切所行的.

$^{161}$其實不簡單.
為什麼呢?.
在他的上一代.
即是在他父親亞克斯王作王期間.
其實引入了很多外邦偶像敬拜的風俗.
甚至私下更改了祭壇祭禮的規矩.
改變了上帝所訂立的法則.
還要將聖殿的門封鎖.
即是說他的上一代.
他的父親.
令整個猶大國陷入屬靈的最低谷.
所以你看到希西加為什麼被.
廖王紀平為行上帝眼中看為正的事.
是因為他一上場之後.
他決心擺脫上一代所帶來的厄運.
他推動了一連串的改革.
包括在宗教方面.
他徹底剷除當時所有的休壇.
廢掉了那些木偶,柱像.
並且打碎了摩西當時所建造的銅蛇.
是在文素所立的銅蛇.
因為到了希西加年代.
以色列人已經把銅蛇當作偶像來拜.
失去了原本的意思.
所以你看到希西加做了.
在宗教方面很大的改革.
將所有的偶像崇拜徹底粉碎.
另一方面他重開聖殿.
結正聖殿.
重新按照祭司利未人的班次.
來安排他們事奉.
恢復聖殿敬拜的秩序和功能.
而另一方面隨著北角以色列的殞落.
出現了很多難民.
希西加通過全國.
包括由南至北,全國以色列.
來呼籲散落的同胞.
你們可以一起回到耶路撒冷.
向耶和華守逾越節.
你看到他很想.

$^{201}$將分裂的南北兩國重新統一起來.
聯合十二個支派.
亦讓亡國的同胞可以回到上帝的身邊.
得到上帝的醫治和盼望.
而另一方面.
聶王頡亦記載在軍事方面.
希西加亦很蒙上帝的祝福.
收服了很多他父親.
阿哈斯王在位的時候.
被菲律斯人奪去的土地.
令猶大的西南部局勢得到安穩.
至於在政治層面方面.
阿哈斯王在位的時候.
為了對抗北面的亞蘭國.
和以色列的聯合攻擊.
於是當時的猶大.
與亞述帝國結盟.
亦甘願成為亞述的附庸國.
需要連連向亞述王進攻.
但希西加上場後.
他們都決心擺脫亞述的控制.
當時的亞述王一死之後.
希西加就趁機背叛亞述.
拒絕向新任亞述王.
西拿基納進攻.
於是當然.
這舉動惹來當時新任亞述王的不滿.
他就開始起兵攻打猶大.
聖經記載他亦都攻取了一些猶大的城邑.
在那個時候.
當希西加面對敵人的強勢.
於是他就唯有向亞述王認罪.
並且賠償巨款.
為了籌這筆經費.
聖經記載他甚至.
要將聖殿門上的金.
和包在柱子上的金紙.
刮下來.
刮下來送刀給亞述賠款.
不過.

$^{241}$希西加不單止這樣.
一方面他.
緩和亞述王的強勢.
但同時他又和埃及帝國.
當時埃及是另一邊很強大的國家.
他和埃及結盟.
試圖拉攏埃及對抗亞述王.
這是當時經文的歷史背景.
你會見到希西加是一個很好的君王.
他成功擺脫了前朝所帶來的枷鎖.
行上帝眼中看為正的事.
不過.
當逆境臨到的時候.
他又不知不覺重蹈了前朝的覆轍.
昔日他的父親阿克斯王.
為了對抗亞蘭和北角的入侵.
不惜不聽先知爾賽亞的警告.
堅持要和亞述來到結盟.
結果得罪上帝.
同樣希西加當他面對急難的時候.
他又效法了他的父親.
用外交的方式來處理亞述的問題.
他似乎相信比起在上帝裡面尋求出路.
埃及的軍事幫助可能更加有效.
結果當然就像我們今天所看到的爾賽亞書.
他惹來上帝的責備.
上帝如何責備他呢?.
我想大家一起讀一段經文.
印在程序表裡.
爾賽亞書第30章.
一至三節.
第二段經文.
預備起.
.
謝謝.
上帝很清楚他斥責猶大是甚麼.
是活逆的兒女.
為何說他們背叛了上帝.
是因為猶大唯一的盟友只有上帝一個.
他們只可以和上帝結盟.

$^{281}$換句話來說.
猶大現在選擇和埃及結盟.
不再能和上帝保持專一和聖結的關係.
其實就等同於背叛了.
他們昔日和上帝所納的約.
更嚴重的是.
其實這段經文說得很清楚.
他們繞過了上帝.
沒有事先求問上帝的心意.
就私下達成這單交易.
來和埃及同謀結盟.
其實他們明知上帝很憎恨這件事.
為甚麼.
他們明知上帝很憎恨和埃及結盟.
因為在這個民族歷史.
上帝昔日將他們從埃及遺牢之家拯救出來.
拯救了之後就和他們立約.
所以現在上帝怎會喜歡他們.
回到舊路.
和埃及這個遺牢之地結盟.
他們其實應該很清楚.
上帝是憎惡.
所以上帝在這裡表明.
猶大的計謀是不會得逞.
埃及不會成為他們的力量.
只會成為他們的羞辱和慚愧.
在這裡讓我們看到.
即使希西加是一個很好的君王.
幸而我眼中看為正的事.
都有他閃失的時候.
正如人為了盡快解決出現在眼前的危機.
有時都會怎樣.
明知不應該做的事都會去做.
我曾經聽過一位中學校長分享.
他說在疫情期間要考試.
考試期間曾經接到一位家長打電話來學校.
說他的兒子身體不適.
他有點作嘔但沒有發燒.
所以需要遲一點才能回到學校.
學校明白這個情況.

$^{321}$但因為疫情緣故都緊張.
所以馬上做應變安排.
特意找一個獨立的房間.
讓這位遲一點回來身體不適的同學考試.
這是一個很好的安排.
結果同學就遲了回來.
當老師問他哪裡不適時.
他就說「沒事,我只是睡懶了」.
於是乎.
事後校方聯絡家長了解清楚情況.
那位家長說他只是擔心.
因為怕兒子遲到會影響考試分數.
所以情急之下出此下策.
其實人有時就是這樣.
明知不應該這樣.
但情急之下就犯了像希西加的錯誤.
用了一些不應該用的方法.
或者說錯誤的方法.
甚至上帝不喜歡的方法.
來解決眼前的危機.
有時甚至用惡的手段來制裁惡的事情.
就好像我最近.
吃晚飯時在家.
打開電視.
我想大家有時會用電視撈飯.
看到一套劇集.
當時的主角的家人.
因為在網絡上玩網絡遊戲.
結果被騙了錢,課金.
或者說騙了感情.
於是主角就想替家人出頭.
要替他討回公道.
他就設局.
騙騙子.
結果當然成功.
成功騙騙子.
令他有財物損失.
可以替家人討回公道.
家人亦大團圓結局.
雖然是一個搞笑的劇集.

$^{361}$但不知不覺原來.
把不好的訊息.
植入了家庭內.
某程度上是說.
原來很多事情.
我們可以以惡制裁惡.
這樣是合理的.
這樣是有好的結果.
最慘是那時是.
合家歡事段.
所以我要一邊看.
一邊替兒子做家長指引.
因為很容易入口.
笑著笑著就一集.
但原來把很多.
上帝不起悅的做事方式.
不知不覺植入了我們腦袋.
同樣回到這段經文.
現實生活中.
我們會否有時都像猶大一樣.
用了一些錯誤.
上帝不起悅的方法.
來應付一些生活危機.
或者說以上帝憎惡的手段.
作為一個避難所.
問題可能真的解決到.
甚至很快解決到.
但卻是衍生了.
另一個更嚴重的問題.
是什麼?.
你得罪了上帝.
那是更嚴重的問題.
第十四節.
你會看到.
先知對耶路撒冷領袖.
所發出的警告.
其實你看到.
先知兩次來諷刺.
猶大和埃及結盟.
來對抗亞述.

$^{401}$就形同什麼?.
和死亡立約.
與陰間結盟.
用廣東話說是什麼?.
死路一條.
實死無三.
猶大以為埃及很穩陣.
靠得住.
可以成為避難所.
抵禦亞述的攻擊.
但上帝卻拆穿了.
他們的謊言和虛假.
上帝表明他們所依靠的.
其實是豆腐渣工程.
會被亞述一沖就散了.
他們以為埃及.
可以讓他們安枕無憂.
但先知這裡諷刺他們.
第十節說什麼?.
原來床塔短.
使人不能書身.
被窩窄.
使人不能遮體.
好像有幫助.
但卻是怎樣?.
令他們睡得不好.
睡得.
一來可能會著涼.
二來可能書枕不著.
沒辦法安睡.
相反.
上帝在第十六節.
說他會安放一塊石頭.
上帝說這塊石頭是試驗過的.
經得起考驗.
是寶貴堅固的防角石.
它會成為建立.
猶大國的準則和垂線.
任何不符合上帝規格的建造工程.
都會怎樣?.

$^{441}$基於上帝這塊防角石的準則.
一方面被揭露出來.
另一方面會怎樣?.
拆走它.
所以.
先知這裡提醒猶大的領袖.
到底今後猶大的發展.
是堅固穩妥.
還是一衝就散?.
很在乎於.
人怎樣去看.
上帝安放這塊防角石.
如果人懂得靠這塊防角石.
他就會懂得.
好像經文裡所說.
他不是先急於向外.
亂來幫手.
來幫自己扭轉敗局.
而是怎樣?.
他先做好他份內的事.
他應該對內.
持守著上帝公平公正的吩咐.
先知說這才是活路.
很多時和人的想法相反.
人遇到危難時.
就急忙向外找方法.
來解決自己的問題.
很想盡快扭轉敗局.
上帝說不是.
你先看一看它.
你先靠近它.
先對準它.
你先做好.
上帝吩咐你的事.
其他的事交回給上帝.
所以弟兄姊妹.
值得我們思考的是.
特別在某程度上.
我們身處在一個動盪的環境.
很多不安感的環境.

$^{481}$特別可能.
有時你會覺得.
辛苦經營下來的事.
很可能會怎樣?.
慢慢不見感.
好像被人搶走一樣.
你會怎樣做?.
你會靠甚麼作為避難所和防擴石?.
你所依靠的.
到底是令你越來越輕散,平安.
還是越來越疲倦.
甚至是挫敗?.
又或者換另一個方法來說.
當我們面對困局.
真的要做決定的時候.
你的決定是令你更加靠近上帝.
還是遠離它?.
正正是這段經文.
給我們思考的空間.
特別可能.
你會收到每個月的強積金.
好像一直怎樣?.
正在下跌.
你辛苦過往儲蓄十多年的那些怎樣?.
真的好像如浮雲一樣.
轉眼就飄走了.
你捉都捉不緊.
又或者你的樓價會怎樣?.
一路跌,一路跌,一路跌,一路跌.
你的身家好像越來越貶值.
你又搞不定.
有人會選擇盡快賣掉手上那層樓.
轉移你的資產.
或者甚至移民去其他地方.
來繼續保住你所擁有的.
又或者有人覺得.
先看清楚.
不是那麼壞的.
環境.
按兵不動繼續觀望.

$^{521}$再想.
其實真的很難說.
什麼才是最明智的決定.
畢竟我們手上沒有那個.
可以預知未來的水晶球.
可以為自己籌算到一份百分之一百的保障.
其實作為上帝的管家.
我們一方面真的要.
為未來做好籌劃和計算.
如何善用上帝給我們的資源.
對不對?.
但是有一樣東西.
我們更加要明白的就是.
哪一個才是我們真正的保障.
在籌算的過程.
在計劃的過程.
這段經文就提醒我們.
什麼才是我們真正的保障.
世上到底有沒有其他東西.
會比上帝更加可靠.
更加穩陣呢?.
沒錯,上帝沒有允許每一個.
投靠,信靠祂的人.
一定會時時順境.
凡事成功.
上帝亦沒有允許你投靠祂.
祂會不斷增加.
保存著你手上的財產.
但是上帝允許.
祂一定夠用.
祂一定有足夠的平安保護.
來承載著我們.
所以等於姐妹.
當上帝祂表示.
祂就是那一塊.
經得起考驗的防角石.
祂就是那一塊.
其實珍貴穩固的防角石.
並且呼籲我們.
去投靠祂的時候.

$^{561}$我們懂不懂得選擇呢?.
我們到底懂不懂得.
真的去投靠祂.
靠著祂.
以祂作為人生的基準呢?.
還是我們總是很想.
在上帝之外.
找一些額外的保障.
總是覺得上帝不夠穩陣.
所以多買幾份保險.
經文提醒我們.
無論你面對什麼人生考驗.
昔日的希西加也好.
今日的我們也好.
首先保住的.
不是地上所擁有的.
經文提醒無論什麼時候.
我們首先要保住的.
是你和我的屬靈財產.
就是我們和上帝的關係.
永遠都是在考驗那一件事.
當危難臨到.
你會在乎誰?.
什麼是你的保障?.
什麼才是你真正的需要?.
通通都是這段經文.
給我們的提醒.
到底我們願不願意.
去靠著上帝這個防國石.
順著祂隨時的帶領和幫助.
來生活呢?.
上帝在問我們這件事.
我們又會怎樣去回應呢?.
和大家一起看最後那段經文.
我們一起讀.
最底那段經文.
第23至29節.
今天我們沒有讀到.
現在可以讀一讀.
好.

$^{601}$預備起.
你們當質疑聽我的聲音.
留心聽我的言語.
那為殺種而耕地的.
會不停地耕地.
送土爬地嗎?.
他剷平了地面.
豈不就種小茴香.
播種大茴香.
按行列作小麥.
在定處種大麥.
在田邊種粗麥嗎?.
他的神指教他.
指導他合宜的方法.
原來打小茴香.
不用尖利的器具.
紮大茴香.
不是用車輪.
卻要用杖打小茴香.
用棍打大茴香.
谷要打.
但不能持續地倒.
用車輪和馬紮.
卻不紮碎它.
這也是出於因之耶和華.
他的謀略奇妙.
他的智慧廣大.
很難聽吧?.
我們很少讀這段經文.
但是他的道理不難明白.
正如先知呼籲當時猶大的百姓.
面對上帝的管教.
首先要學習農夫.
向他們學習生活的智慧.
那些當時被譽為.
沒有高深學問.
沒有上過農業學的農夫.
他們為何會懂得這麼多.
為何會懂得耕田.
為何會懂得用不同工具.

$^{641}$處理不同的穀物收成.
先知說得很清楚.
答案是上帝教他的.
上帝指導他.
用合宜的方法做合宜的事.
這樣就會有收成.
是因為上帝將他自己的謀略智慧.
指示給農夫知道.
其實不難理解.
即使我們今天.
農民我們說這個行業是.
靠天吃飯.
看天做人一個的行業.
他既需要努力工作.
但更需要.
順應上帝所命令的法則.
一些農業的法則.
一些節期.
要在適當的時候.
按照上帝的安排.
做適當的事情.
農夫就是這樣.
每天依靠上帝.
日復日累積到.
在他心目中的人生智慧到底是怎樣.
他明白無論我多麼努力.
多麼精明也好.
上帝才是我真正的策士.
因為上帝的策略指引.
簡單易明最有效.
順著上帝的帶領.
生活工作自然能夠更加的輕散穩妥.
所以.
先知很想藉著.
農夫這一份的智慧.
來一方面凸顯.
其實尤大領袖當時.
那種愚昧無知.
也勸勉他們.
與其在混亂的情況下亂來.

$^{681}$亂找辦法解決.
倒不如.
單純地仰望上帝.
給你當下的指引.
順服他的管教.
這樣才是上策.
正如我以前有一位神學院老師.
他事實上這樣教導他的學生.
他說我們每一天不多不少.
出門前都要看天色做人.
正如今天你不會戴傘.
今天你看到應該不會下雨.
所以我看大家應該沒有傘在身.
不單看天色做人.
還要看什麼呢?.
有沒有看上帝的面色做人?.
凡事如果學習先仰望上帝的面色.
其實有時不再需要看.
那個界懷的局勢到底變化是怎樣.
不再需要計較別人的面色是怎樣.
確實是一個很有智慧的道理.
每天我們都先看看上帝的面色做人.
凡事做決定的時候.
先看看上帝是否對我們和顏悅色.
還是開始憤怒我們.
其實真的.
我們都明白有時太過.
看別人面色做人其實舒不舒服.
很辛苦.
其實某程度上是被人的需要和想法牽著被指走.
更加會讓我們忽略了.
上帝在我們生命裡的帶領和旨意.
所以在《鼎姐妹經文》裡面.
最後勉勵我們能否保持著一顆.
單純依靠上帝生活的心.
事實上我們喜不喜歡單純.
在這個社會生活.
有沒有賣點呢?單純.
單純不是在說我們空無大志.
而是在說我們怎樣去專注.

$^{721}$執著我們人生首要的目標.
怎樣去榮耀上帝討他的喜悅.
這叫單純.
單純亦不是在說我們遠離世俗.
而是在說我們面對惡人惡事的時候.
我們仍然可以在上帝面前.
保持一個分別為勝的心.
單純亦不是在說我們對人完全不設防.
而是在說我們懂得自保之餘.
我們不和人爭敬.
我們總是持守良善的心.
設法與人和睦.
單純更加不是在說.
我們一點也不為未來去籌算.
而是我們確信上帝的意念安排.
總是怎樣?.
最穩妥.
比人的想法超越.
所以我們願意隨時配合.
上帝的帶領.
願意被祂糾正.
接受祂的管教.
來調教我們做人的方向和目標.
確實按照社會的遊戲規則.
做人單純正直.
可能真的會很吃虧.
特別身處一個動盪的環境.
當人人都在說要怎樣去自保的時候.
做人單純正直.
好像真的沒有甚麼保障.
所以你看到.
這個世界有不少比我們聰明.
有學識有地位的人.
他們為了預約他們的成功人生.
就會怎樣?.
他們選擇的是處心積慮.
他們選擇的是用一些手段來上位.
增加他們的財富.
或者保住他們所擁有.
他們相信這才是生存的法則.

$^{761}$這才是自性之道.
但今天這段經文.
它說的是相反.
原來在上帝的生命之道裡.
持守單純正直.
一點都不吃虧.
因為當你在上帝面前.
持守純全正直.
其實本身.
這就是最大的保障.
持守純全正直.
本身就有它的保護功能.
因為誰保護你?.
上帝應許了.
他會看顧那些在他面前.
單純投靠他的人.
上帝應承了.
所以.
很盼望今天我們藉著上帝的話語.
幫我們去重新審視調教.
到底我們現在在走一個什麼方向?.
我們是否真的對準上帝安放在我們生命裡面.
那一塊防角石?.
以致我們真的可以穩行在上帝的心意裡面?.
讓我們同心祈禱.
有如希西加般的軟弱.
特別當急難困境來到.
我們很想盡快解決的時候.
我們就會依靠其他勢力.
來成為人生的防角石.
上帝的意志.
我們真的輕慢了你.
甚至乎得罪了你.
但上帝我們知道.
你有恩典,有憐憫.
每當我們真的走迷的時候.
你總是會用不同的方式.
用我們明白的方式.
來挽回,提醒我們.
以致我們慢慢學習.

$^{801}$你才是我們堅固的保障.
在任何時候.
即或在困境臨到的時候.
我們首先要保住的.
仍然是和你的關係.
在困難的時候.
我們首先要審視的.
我們依靠的是甚麼.
到底是不是上帝你.
以致即或我們.
知道世界上有很多事情.
不是我們能夠掌管.
但我們知道上帝.
你的恩典總是夠用.
當我們願意靠近你.
單純地信靠你的時候.
上帝你就會帶領我們.
走過一個一個.
又一個的難關.
讓我們和你的關係.
更加堅固,更加穩妥.
上帝為此幫助我們.
讓我們能夠藉著你的話語.
藉著聖靈的提醒.
重新調教我們往後的生活方向.
以致無論我們或信或抑.
我們知道我們都是上帝你的兒女.
你是我們的主,我們的神.
我們生命的唯一目標.
是榮耀你,討你的喜悅.
進行你的旨意.
為此我求你幫助我們.
誰聽我們在你面前的禱告.
奉基督耶穌你得勝的命旨.
來祈求.
阿們.
\newpage



\section{馬太福音 4:1-11}
\label{sec:862eUMKnVnc}
\textbf{2022年10月16日崇拜直播｜梁煒傳道｜以石為記:試探之石｜馬太福音4\_1-11}
\newline
\newline
連結: \href{https://youtube.com/watch?v=862eUMKnVnc}{\texttt{https://youtube.com/watch?v=862eUMKnVnc}} ~~~~ 語音日期: 2022-10-18
\newline
\newline
\hyperref[sec:x9IT18Te0VA]{\small{< < < PREV SERMON < < <}}
~
\hyperref[sec:index_chronic]{\small{[返順時目]}}
~
\hyperref[sec:YIG0i1d0dV8]{\small{> > > NEXT SERMON > > >}}
\newline
\newline
馬太福音 4:1-11
\newline
\begin{longtable}{cl}
\hline
\hline
章節 & 經文 (和合本修訂版)\\
\hline
4:1 & \begin{tabularx}{0.7\textwidth}{X} 當時，耶穌被聖靈引到曠野，受魔鬼的試探。 \end{tabularx} \\ \\ \relax
4:2 & \begin{tabularx}{0.7\textwidth}{X} 他禁食四十晝夜，後來就餓了。 \end{tabularx} \\ \\ \relax
4:3 & \begin{tabularx}{0.7\textwidth}{X} 那試探者進前來對他說：「你若是神的兒子，叫這些石頭變成食物吧。」 \end{tabularx} \\ \\ \relax
4:4 & \begin{tabularx}{0.7\textwidth}{X} 耶穌卻回答說：「經上記著：『人活著，不是單靠食物，乃是靠神口裡所出的一切話。』」 \end{tabularx} \\ \\ \relax
4:5 & \begin{tabularx}{0.7\textwidth}{X} 魔鬼就帶他進了聖城，叫他站在聖殿頂上， \end{tabularx} \\ \\ \relax
4:6 & \begin{tabularx}{0.7\textwidth}{X} 對他說：「你若是神的兒子，就跳下去！因為經上記著：『主要為你命令他的使者，用手托住你，免得你的腳碰在石頭上。』」 \end{tabularx} \\ \\ \relax
4:7 & \begin{tabularx}{0.7\textwidth}{X} 耶穌對他說：「經上又記著：『不可試探主—你的神。』」 \end{tabularx} \\ \\ \relax
4:8 & \begin{tabularx}{0.7\textwidth}{X} 魔鬼又帶他上了一座很高的山，將世上的萬國和萬國的榮華都指給他看， \end{tabularx} \\ \\ \relax
4:9 & \begin{tabularx}{0.7\textwidth}{X} 對他說：「你若俯伏拜我，我就把這一切賜給你。」 \end{tabularx} \\ \\ \relax
4:10 & \begin{tabularx}{0.7\textwidth}{X} 耶穌說：「撒但，退去！因為經上記著：『要拜主—你的神，惟獨事奉他。』」 \end{tabularx} \\ \\ \relax
4:11 & \begin{tabularx}{0.7\textwidth}{X} 於是，魔鬼離開了耶穌，立刻有天使來伺候他。 \end{tabularx} \\ \\
[1ex]
\hline
\hline
\end{longtable}
$^{1}$今日主教崇拜要讀的經文是記載在馬太福音第四章第一至十一節.
我們可以打開我們手上的聖經.
新約聖經的馬太福音第四頁.
馬太福音第四章第一節.
當時耶穌被聖靈引渡曠野受魔鬼的試探.
他禁食四十晝夜後來就餓了.
那試探人的進前來對他說.
你若是神的兒子可以吩咐這些石頭變成食物.
耶穌卻回答說.
經上記載說.
人活著不是單靠食物乃是靠神口裡所出的一切.
魔鬼就帶他進了城城.
叫他站在墊頂上.
對他說.
你若是神的兒子可以跳下去.
因為經上記載說.
主要為你吩咐他的使者用手托著你.
免得你的腳痛在石頭上.
耶穌對他說.
經上又記載說.
不可試探主你的神.
魔鬼又帶他上了一座最高的山.
將世上的萬國與萬國的榮華都致給他.
對他說.
你若苦復敗我.
我就把這一切都賜給你.
耶穌說.
撒旦.
退去吧.
因為經上記載說.
當拜主你的神.
丹要事奉他.
於是魔鬼來了耶穌.
由天使來侍候他.
靈姐妹平安.
多謝靈姐妹經常去問候.
關心我們家庭.
以及問候小女的健康情況.
現在小女讀M班一間特殊學校.
上學後一個多月.

$^{41}$感恩這間學校其實很有愛心.
家長都可以陪她上課.
小女在學校裡的適應都很好.
上課都很專心.
她在同學裡有六個.
她是懂得坐直的那個.
其他人都需要其他輔助.
因為特殊學校有很多困難的小朋友.
這幾個月裡都看到她.
有很多的成長和進步.
接下來小女都會開始一些新的治療.
無論是言語治療.
或者一些其他方法.
推拿或者感同的刺激.
其實很想幫助她.
能夠更加努力地成長.
因為小女這個病.
我們家裡一家人裡.
我們的起居飲食都要特別小心.
特別現在就是新冠肺炎的一個時期.
這個病毒其實對長期病患的人.
都是有影響的.
所以從年頭開始的時候.
我們一家人已經是不出外.
不在餐廳裡面吃飯.
減少感染的機會.
不過雖然不出去和人吃飯.
或者是不脫下口罩和人說話.
不過都沒有難了我去服侍.
我仍然都是戴著口罩和年輕人入營.
三十多天或者一夜都好.
晚上我都是不脫下口罩.
戴著口罩睡覺都可以的.
和其他人吃飯之前.
我就預先消失十分鐘.
簡單地吃了它.
然後繼續和人吃飯裡面的交談.
記得有一次和一班年輕人吃火鍋.
哇.
桌子上擺滿著肥牛.

$^{81}$肉丸.
海鮮.
烏冬.
這些東西都很吸引.
還要火鍋的味道都很強.
我已經和大英姐妹說.
其實我都不會和大家吃的.
你們OK的你們自己放心吃.
大英姐妹都很好嚇她.
她再問.
喂你真的不用都很好吃的.
又過了一陣子.
她說真的不吃你OK的.
你們吃吧我夾給你們.
服侍你們.
都很好.
幸好我都忍得過.
上個星期的時候.
我剛剛浸大完了一些after course.
和一些大學生去樂富附近.
有一個比較有名的大排檔一起吃飯.
那天因為忙了一點.
沒機會吃飯之前吃東西.
所以肚子都很餓.
大學生都很會叫東西.
叫了一堆蠔餅.
知不夠吐魚.
知足火腩.
嘩.
那時候我都很餓.
竟然不聞到這些.
我都很久沒吃過了.
我都真的很想吃.
不過一想起.
這個食物的誘惑都很吸引.
但我想起小女的健康.
我就忍住.
我堅持下去.
不可以脫口罩.
或者最多早點走.

$^{121}$回家吃個飯.
在生活裡面.
不知道你會不會遇到.
這些類似的掙扎.
困難.
試探.
誘惑在我們裡面.
在我們生活裡面的確是充滿著.
各式各樣的試探和誘惑.
我們沒有辦法.
沒有一個人能夠避免.
沒有一個人能夠免疫.
今天很想透過.
一個role model.
我們的耶穌基督.
一起去看看.
看看如何面對著.
殺蛋的試探.
成為我們的幫助.
我們開始.
在聽完之前我們一起祈禱.
今天生活裡面.
有各式各樣的試探.
金錢的誘惑.
一些犯罪的機會.
一些關係裡面的試探.
最後我們都陷入在.
這些眼目裡面的情慾.
甘心裡面的驕傲裡面.
你留下了很寶貴的聖經.
你都將最寶貴的聖靈賜給我們.
神我們求你.
開我們的眼.
開我們的心.
開我們的耳朵.
更加幫助我們明白你的話語.
讓這些話語都成為我們有力.
繼續面對著試探.
我們都能夠學習耶穌基督.
一起同樣去勝過.

$^{161}$聽我們祈禱.
奉耶穌得勝之名而求.
阿們.
對魔鬼殺蛋來說.
他知道耶穌來到這個世界的目的.
耶穌來到這個世界的目的.
就是要拯救人.
從罪惡裡面能夠拯救出來.
然後擊敗魔鬼.
殺蛋罪惡裡面的權勢.
如果魔鬼在這個階段.
耶穌還沒有出動傳道的時候.
能夠成功誘惑耶穌.
令到他犯罪.
令到他不遵行上帝的話語.
這樣耶穌上帝.
他拯救的工作.
就要提早收工.
就要面臨失敗.
所以耶穌基督.
面對著這三個試探.
這三個試探是魔鬼殺蛋.
特別精心策劃.
在耶穌最軟弱的時候.
在他最困難無助的裡面.
在他沒有任何的幫助裡面.
他很想設計三個試探.
讓耶穌跌倒.
耶穌受試探.
在三權福音書裡面.
其實都有記載.
不過馬可只是用了兩字經文裡面記載.
沒有提及三個試探裡面的內容.
而路加福音.
都有記載到三個不同試探裡面的內容.
只不過次序是不一樣.
三個的試探.
好像很合理.
一個人肚餓.
找食物.

$^{201}$有東西吃都很正常.
遇到危險.
求神的保護都好.
在我們人生裡面.
我們為什麼努力奮鬥呢.
都是要追求可能成就或者榮譽.
但是為什麼這三個的場景.
耶穌或者魔鬼.
就說這件事.
聖經的作者.
形容這三件事是一個試探.
第一個試探.
其實是要耶穌利用他自己的權能.
去滿足他的需要.
在經文的上下文裡面.
第三章.
第三章裡面.
他記載了耶穌受施浸約翰的真理.
天開了.
聖靈好像甲子一樣降臨在耶穌的身上.
有聲音跟耶穌和施浸約翰說.
這是我的愛子.
我所喜悅的.
下文耶穌經歷了這個試探之後.
他開始去傳講天國裡面的福音.
叫人去悔改.
去認識天國.
因此.
在耶穌準備出發傳道之前.
在得到上帝的認可.
神認愛子身分之後.
要面臨這個試探.
更加確定他出發的時候.
他的心態有沒有純正.
他的方向是否正確.
在《馬太福音》四章一至到第四節.
我們邀請弟兄姊妹一起讀一至四節的經文.
.
在第一個試探裡面.
試探的特點就是.

$^{241}$這個試探發生在曠野.
在猶太人的心目中裡.
曠野其實是一個魔鬼撒旦活動的地方.
世界上的這些以色列人.
就在曠野裡漂流.
經歷各樣的試探和困難.
在曠野裡面.
他們面對著孤單.
在曠野裡什麼都沒有.
面對著內心裡真實的需要.
這個曠野是一個人的時候.
這個曠野是沒有人的地方.
可能是一個失縮的角落.
可能是自己一個人在房間的時候.
可能是自己一個人在用電話的時候.
這些試探來得特別真實.
而這個試探發生在耶穌禁食了四十週夜.
猶太人的習慣裡面.
當他們要去領受上帝的誡命的時候.
他們會先去禁食.
去潔淨自己幫助自己的心清意潔.
領受從上帝而來的誡命.
就像摩西領受上帝的十誡的時候.
他先前都會去禁食.
禁食四十週夜其實都挺辛苦的.
禁食過後其實肚子都應該是很餓了.
其實人裡面有很大的慾望.
特別是吃肉.
有很大的很想吃東西.
而這個魔鬼撒旦的試探.
經文裡面就說到.
如果你是神的兒子.
這個試探不是說魔鬼撒旦不知道.
耶穌是神的兒子.
如果你是.
其他日本說既然你是神的兒子.
其實他是知道耶穌的身份是神的兒子.
不過他考驗的是.
究竟耶穌會不會利用他超自然的能力.
他的權柄.

$^{281}$用在自己身上.
滿足自己.
將這些石頭變成食物.
這個才是真正的試探.
耶穌絕對有能力將這些石頭變成食物.
因為他能夠將水變成酒.
他能夠將五個餅兩條魚餵飽超過五千人.
不過生存的確是人最基本的關注.
食物,工作,賺錢,休息.
或者是各樣人最基本的需要.
將石頭變成食物.
絕對有需要.
因為耶穌來到這個世界上有更加重要的使命.
如果他餓死了怎麼辦?.
他怎麼能夠完成上帝的拯救工作?.
不過耶穌面對魔鬼撒旦這個試探.
他引用了《申命記》第八章第一節.
《申命記》的經文.
我們一起去看看《申命記》.
其實這節經文的上下文是說什麼?.
《申命記》第八章第一節.
「我今日所吩咐的一切誡命,你們要謹守遵行,好叫你們傳閱,人數增多..
且進去得耶和華向你們列祖所建的那地..
你們也要紀念耶和華你們神在曠野引導你這四十年,.
是要苦練你,試驗你,要知道你內心如何,.
狠守他的誡命不肯..
他苦練你,任你飢餓,.
將你和你列祖所不認識的猛辣似給你吃,.
使你知道人活著不是單靠食物,.
乃是靠耶和華口裡所出的一切話.」.
當我們將《馬太福音》第四章的場景.
和《申命記》第八章的場景.
拼合在一起.
我們看到很多相似的地方.
曠野,四十,苦練,試驗.
人活著不是單靠食物.
乃是靠耶和華口裡所出的一切話..
耶穌基督就成為了新的一班的以色列人.
面對著這個試探..
其實曠野的環境其實都很艱鉅.

$^{321}$在白天的時候太陽很猛烈.
在晚上的時候天氣都很冷.
而且曠野有很多野獸,猛獸夜間裡面的行動.
曠野裡面都沒有食物.
也沒有水源.
是一個了無人煙的地方.
曠野也沒有任何的保障.
如果忽然間有人晚上打你的時候.
你無法防備.
但是《申命記》第29章第5節.
「我領悟你們在曠野四十年.
你們身上的衣服並沒有穿破.
腳上的鞋也沒有穿壞」.
你一雙鞋子應該不會穿四十年吧.
但是神很奇妙.
他在說的是上帝的保守.
上帝帶領以色列人離開了埃及.
在曠野裡面漂流了四十年.
上帝保守他們.
供應他們的食物,供應他們的水源.
讓雲柱,火柱成為他們生命裡面的導引.
能夠供應他們的所需.
上帝是允許這群以色列人進入應許之地.
不過在進入這個應許之地,加林地之前.
上帝藉著《申命記》這段經文.
提醒這群以色列人.
這四十年裡面.
其實他們在學什麼功課.
這些功課就是當他們願意遵行上帝的說話.
上帝會供應他們的需要.
祝福他們的人數.
讓他們帶領他們的生命.
能夠進入這個應許之地.
上帝要這群以色列人明白.
遵行上帝的說話.
能夠行使上帝的使命.
比食物,比其他人一切的需要更加看重.
當生命的需要和上帝的使命.
相違背,相衝突的時候.
這個其實是一個試探.

$^{361}$要人去學習.
如何去信服,如何去寫記.
耶穌在《約翰福音》四章三十四節.
說到「我的食物就是遵行差我來者的旨意.
造成他的功」.
耶穌清楚知道祂來到這個世界的目的.
就是要拯救人從罪惡中得釋放.
為了這個目的.
祂願意放下肉體裡面的需要.
在耶穌傳道的日子.
祂由早上天還沒亮的時候.
已經開始出來服侍.
到晚上.
合乘的人都去找耶穌.
耶穌在當中去接觸人.
教導人,醫治人.
祂放下了休息的需要.
在《路加福音》九章五十八節.
祂說「狐狸有洞,天空的飛鳥有窩.
只是人子沒有枕頭的地方」.
耶穌放下了安窩的需要.
耶穌沒有恃著自己是上帝.
而自己的身份有特權.
祂沒有為自己謀求福利.
當祂知道生存的需要和祂的使命的時候.
祂選擇寫記.
選擇願意放下祂的需要.
不過在聖經裡面.
耶穌因為堅餓.
不惜用長子的名份和雅各換取了紅豆湯.
在《士師記》裡面.
參孫因為貪愛.
即妓女.
違背了那世二人的立約和使命.
最後付上生命的代價.
在新約裡面.
加略人猶大.
他因為三十萬銀錢的緣故出賣了耶穌.
放棄了成為門徒的身份和使命.
當我們打開新聞.

$^{401}$當我們看到教圈子裡面.
其實都有很多.
牧者,傳代人,教徒.
都會因為各種名利,各種利益的原因.
他們違背了在自己的崗位裡面.
為了自己的利益去行使了不應該行使的權利.
的確.
食物.
一些的情慾.
金錢,權力.
其實都是很重要的.
都是很多人很努力追求.
不過當這些東西和我們的身份.
有違背,有衝突的時候.
其實我們要學習如何去取捨.
耶穌拒絕摩訶撒旦的試探.
將石頭變成食物.
因為祂知道祂因此不是滿足祂自己.
祂知道只要當祂願意去信服上帝的說話和使命的時候.
做上帝應要祂做的事情.
上帝必定會供應祂的需要.
在第四章第十一節.
第十一節經文裡面.
他說到.
「於是魔鬼離開耶穌,有天使侍候他」.
當魔鬼走了之後.
天使侍候.
侍候的意思是有照顧.
而且這個字更加和食物有關係.
當耶穌做了一些對的事情.
對的決定的時候.
上帝自然會為他安排.
供應他肉體裡面的需要和滿足.
因為上帝是耶華爾勒.
他是供應的上帝.
他是漢見的上帝.
他也是必有預備的神.
他不會眼巴巴看到我們的需要.
他不理我們.
《馬太福音》六章三十三節裡面說到.

$^{441}$「你們要先求他的國和他的義」.
「這些東西都要加給你們了」.
當我們願意去遵行上帝的說話和使命的時候.
上帝會供應我們一切所需要的一切.
我分享一下.
在八年前我夢見要去讀神學.
當時讀神學的過程其實也挺奇妙的.
那時候父親和家人都還沒有相信.
但是他們反而主動地開綠燈.
說既然你將來工作有需要.
那你不如去讀讀神學進修一下.
我覺得已經很奇妙.
當我想讀神學這個想法.
和當時的女朋友.
現在的太太分享的時候.
女朋友也說.
也好不如一起去讀.
而且還是一起讀四年制的課程.
當時我又會在教會裡面工作.
和教會的弟兄姊妹分享.
他們都表示很認同和支持.
和會幫忙做推薦.
於是乎我和舞會的牧者.
去分享上帝裡面的帶領和感動.
告訴他們家人已經開了綠燈.
伴侶也好像開了綠燈.
其他人也有同感一零.
去驗證了上帝.
其實也很想我去讀神學.
再去下一階段的服事.
牧師聽完之後也覺得很受感動.
因為一個年輕人願意去服事神.
牧師跟我說.
你不用擔心.
教會一定會全力支持你.
舞會過往的做法.
如果有人讀神學.
也會有三分之一的支持.
牧師很好.
他也說你這麼有心.

$^{481}$而且你年輕也未必有很多金錢的積蓄.
他說我試試幫你爭取多一點.
希望有一半以上.
牧師也很好.
很愛成年輕人.
但過了一段時間.
牧師打電話給我.
阿偉呀.
關於你報讀神學.
你要記得.
我會支持你.
教會全部人都支持你.
不過因為你報讀的神學院.
和我們本身的宗派不一樣.
所以教會是不會用現金支持你讀神學的.
不過他說相信上帝會預備.
你放心繼續去.
走在上帝的道路裡.
那時候讀建道.
我想去報讀的學校.
建道是要學費.
要住在裡面.
住在裡面的三餐.
我昨天也特意上網找找.
現在的費用是一年的學費,宿費和膳食費.
大約是八萬元.
要全時間在裡面讀.
你沒有收入.
如果你生活費節省一點的話.
一年大約是兩萬元.
總的來說.
一年可能要學費,膳食費,宿費和生活費加起來.
一年要十萬.
如果我讀的課程是四年的東西.
那條數也挺大的.
八年前如果有幾十萬在手的時候.
我覺得我可以做其他東西.
但是那時候我因為.
剛剛大學畢業.
在基督教機構裡面工作.

$^{521}$在教會裡面做科研幹事.
其實薪金也比較低.
那時候也有很大的掙扎.
這麼大的數目.
我沒有錢.
我怎樣去讀呢?.
我祈禱過後也相信這是上帝的帶領.
上帝感動我的.
上帝會為我去預備.
於是我就帶著戰驚的心.
很擔心那時候.
擔心不夠錢.
擔心讀著讀著書沒法繼續讀.
但是上帝很奇妙.
上帝差派了很多天使.
上帝是一個很信實的神.
在我的神學幾年的過程裡面.
他用了不同的方法.
讓我繼續能夠去完成神學的課程.
而且也帶領我來皇浸裡面去服侍.
上帝的帶領很奇妙.
今天我們在四樓這個地方聚會.
是因為六樓要裝修.
六樓裝修也要錢.
我們不單止要為了裝修.
我們更加是為了上帝的國.
為了一些侍宮的服侍.
我們希望能夠有三百萬裡面的籌款.
這個是上帝的工作.
上帝的心意.
上帝必定會在我們心裡面做感動的工作.
我們只要按著上帝的感動.
他就會供應我們一切所需.
經文的第二個部分.
第五到第七節.
不如邀請大家一起讀吧.
第五到第七節.
一二三.
(念經).
在抗疫的地方不行了.

$^{561}$他帶耶穌去到另一個場景.
就是去到耶路撒冷的聖殿頂.
最高的位置裡面.
他讓耶穌置身在危險的當中.
在危險的時候.
人都會很慌張很慌亂.
覺得危險迫在眉睫.
覺得危險有性命裡面的威脅.
是一個很緊急的困難.
令到人在困難裡面.
懷疑上帝的說話.
很想迫使上帝去信服人裡面的需要.
用神蹟的旗幟.
去顯示他的慈愛和信實.
而魔鬼撒旦試探的地方.
他都很厲害.
他應該比我們每一個更加熟悉聖經.
因為他看了聖經.
上帝的說話很多年.
他引用聖經.
他引用聖經裡面的說話.
不過他是斷章取義.
錯誤地去形容.
魔鬼撒旦將耶穌帶到耶路撒冷的聖殿.
他引用了詩篇91篇第12節的經文.
叫耶穌跳下去.
不用擔心.
因為聖經應許過.
有天使會保護你.
根據聖經裡面的考古.
當時耶穌的年代.
在大希律所興建的聖殿.
大約是樓高170呎.
即是52米.
即是當十幾樓高.
其實十幾樓高都很高.
跳下去一定會有生命危險.
詩篇91篇裡面的記載.
耶穌說是我的保護.
是我的避難所.

$^{601}$他能夠叫一切災害都遠離.
屬於他的人.
的確上帝會保護人.
但是詩篇91篇第14節.
裡面也會說一些條件.
類似是.
神說因為他專心愛我.
我要答救他.
他知道我的命.
我要把他安置在高處.
是因為這班人愛他.
遵守上帝的誡命.
遵守上帝的說話.
是一班信服的人.
所以上帝願意去保護他.
上帝願意保護一班愛他的人.
遵從他說話的人.
當他們遭受危險的時候.
上帝必定會拯救.
這個是上帝的應許.
不過魔鬼撒旦厲害的地方.
他將上帝的話語.
說一點不說一點.
上帝的話語不是說人.
能夠在任何的試探.
困難甚至危險.
那些危險可能是自己罪的緣故.
然後引自己的危險.
在困難裡面.
叫人去懷疑上帝.
上帝也不是任由一個人.
濫用上帝的恩典和慈愛.
不斷犯罪.
期望上帝去做拯救和原諒的工作.
上帝是會去拯救.
不過是拯救愛他的人.
所以耶穌回答說.
在他引用的申命記第六章十六節.
他說:你們不可試探耶華你們的神.
像你們在瑪撒那樣試探他.

$^{641}$瑪撒是一個什麼地方呢?.
發生什麼事情呢?.
瑪撒其實是記載在出埃及記的第十七章.
出埃及記的第十七章裡面第七節.
他去那地方名叫瑪撒.
又叫米利巴.
因為以色列人爭鬧.
又因他們試探耶華說.
耶華是在我們中間不是?.
在出埃及記的頭十五章裡面.
上帝用了很多神蹟騎士.
用了十災.
帶領以色人過紅海.
帶領他們出埃及逃離法老的手.
而在第十六章在曠野裡面的時候.
以色人因為沒有食物.
上帝賜下了瑪撒.
賜下了鵪鶉.
上帝一直都在.
上帝一直保守住這班以色列人.
沒有上帝的時候.
以色列人仍然在埃及裡面過著圍爐的日子.
沒有辦法脫離在埃及裡面的追兵.
是上帝一直的同在.
他用神蹟騎士幫助這班以色列人.
但是當這班以色列人去到瑪撒的時候.
因為沒有水.
他們和摩西去爭辯.
叫摩西給他們水喝.
他們埋怨摩西為甚麼要帶他們出來.
帶他們出埃及.
他們埋怨摩西.
是否要把他們嚇死.
令他們的女兒子女和親戚都被嚇死.
他們幾乎要用石頭打死摩西.
以色列人忘記了上帝的恩典和神蹟.
在過往的日子裡面怎樣伴隨著他們.
以色列人因為面對著今天所面對的困難.
他們因為沒有水.
他們很想試驗究竟耶華是否真實的同在.

$^{681}$其實這是一個試探神的想法.
是一個不相信上帝.
懷疑上帝.
質疑上帝的一個行動.
耶穌不願意這樣做.
並不是因為他不知道天父有沒有能力.
而是因為他知道上帝不允許.
上帝禁止人去試探他.
對以色列人來說.
其實要求神蹟的保護並不是一個大的問題.
問題在於他們想試探神.
當耶穌被猶太人出賣.
在赫西瑪利亞人祈禱.
祭司和一班長老帶了很多人來捉拿耶穌.
耶穌說他可以差派十二型的天使來拯救他.
不過耶穌沒有這樣做.
因為他知道他要去信服上帝的拯救計劃.
他也沒有質疑過上帝.
因為他知道上帝和他同在.
不需要靠神蹟騎士救他脫離的危險.
而上帝的信神一直都在.
的確是.
人面對著苦難,困難.
或者我們一些不愉快的事情.
我們很想很快地在這些困難中走出來.
所以很想要求上帝.
神啊!你即時即刻用這個方法解決我的問題.
我希望上帝按照人的方式去解決問題的方法.
因為這個問題一天沒有在的時候.
可能好像這班以色列人一樣.
會去埋怨.
會懷疑上帝.
將這個問題歸咎在上帝身上.
甚至乎向上帝發怒.
好像整個問題裡面.
上帝都要去負責任.
特別在過去幾年的時候.
香港都面對著各種挑戰困難.
或者社會上.
或者疫情裡面.

$^{721}$其實有很多人都在懷疑上帝.
懷疑上帝的公義.
懷疑上帝的信實.
上帝為什麼不出手去拯救.
上帝明明是公義的神.
這班人.
做這些事情你為什麼不出手.
上帝你是一個憐憫的神.
這班人患病.
你為什麼沒有去做拯救的工作.
阻止悲劇發生.
所以有一些人對上帝失望.
因為發現上帝的應許和現實之間.
有很大的落差.
最後他們選擇離開上帝.
拒絕上帝.
人在困難的時候.
真的很容易懷疑神.
我們家庭都試過懷疑神.
因為小女病患的過程.
我們發現這個病已經一年多了.
小女患有一個罕有的病.
這個病傳到香港.
只有幾十個人.
我們去做一些基因排序.
將這些基因排序傳到外國.
發現基因庫裡面.
全球沒有相同的基因出現.
就是說沒有一個醫治的方法.
我們每一天都去為小女的健康祈禱.
無論是我們還是大女.
每晚都祈禱.
祈禱的內容都是求主醫治妹妹.
給妹妹健康.
給妹妹懂得走路.
懂得走路.
陪她玩.
我們每一天祈禱的時候.
第二天早上起床.
我們都在祈望.

$^{761}$會不會有神蹟出現呢?.
會不會忽然間她懂得走路.
懂得跳 懂得走.
唱歌去讚美神呢?.
如果是的話.
真的很震撼 很厲害.
不過上帝沒有容許這些事情發生.
而且當我們看醫生的時候.
醫生每一次都會告訴我們.
這個基因突變會有新的情況出現.
我們每一次見醫生都會覺得很無力.
不知道怎麼走下去.
都會對上帝有些懷疑在當中.
不過上帝卻是信實.
在很多的位置裡面.
祂都很奇妙地安排.
祂為我們安排了一個超級好的醫生.
在香港研究這個病的裡面.
其實也是一個很出色的人.
在妹妹的醫治方面.
又給了我們一個很好的物理治療師.
這個物理治療師也是一個基督徒.
她丈夫也是一個牧師.
每次收費都便宜了給我們.
很重要.
她在小女兒訓練的過程裡面.
由她不懂得爬.
開始教她四點跪 爬.
慢慢地學習站立 學習平衡.
然後在這個階段裡面去學走路.
每個星期.
這個治療師見我見小女兒的時候.
他都在稱讚小女兒.
他說她進步得很快.
每個星期都看到她一點一滴的進步.
相信她來走的日子不遠.
上帝也為我們預備了一個很好的幼兒導師.
這個幼兒導師就是上門去幫小女兒做訓練.
記得年頭的時候.
其實疫情都很嚴重.

$^{801}$經常都過萬宗.
其實都很害怕.
這個幼兒導師住在深井.
我們家住在馬鞍山.
其實都很遠.
深井到馬鞍山.
她其實不是接很多宗.
她都是選擇來做.
因為她也是半退休人士.
她分享說她每次來馬鞍山.
她都是叫Uber來.
她盡量避免搭很多交通工具.
如果有一些病毒感染了就不好.
她很細心.
她很有耐性.
她教小女兒.
將那些幣錢.
進入到小豬錢罐裡.
每一次小女兒都做不到.
我們見到的時候都很灰心.
但是她說又再近了一點.
找些地方讚她.
每一次都很雀躍.
很有耐性地去.
很欣賞小女兒的表現.
原本我們都很洩氣.
很灰心.
不過她說她的出現.
成為了我們很大的力量.
她每次臨走之前.
都會很客氣地問我們.
不知道你介不介意.
將你小女兒資訊的情況.
告訴她教會裡面的小組人聽.
她很想為她祈禱.
她說原來每一次治療裡面.
不單單是用她的技能和知識.
她更加相信上帝的奇妙工作.
這個職業也是很奇妙.
我都很感動.

$^{841}$每一次她上來.
我在家裡服事治療之後.
我覺得她很有力量.
不單單是因為她的治療幫助了小女兒.
更加鼓勵了我們.
上帝說她不會即時.
用神蹟幫人脫離困難危險.
和災難.
不過上帝一直都在.
他一直都會有奇妙的安排.
出人意外的方式.
安排很多天使成為人的幫助.
上帝沒有應許天色常藍.
不過他卻是應許.
無人遇到任何困境.
他都會和我們一起面對.
第三個試探.
也是最大的試探.
第八到第十節.
邀請大家一起讀.
預備一二三.
.
魔鬼撒旦最後一個試探的場景.
將耶穌帶到世界裡.
最高的山的裡面.
大地在我腳下.
他利用人的虛榮心.
利用人的慾望.
將萬國展現在耶穌的面前.
他說只要耶穌做一個很簡單的動作就行了.
這一切都是屬於耶穌.
撒旦在這個世界.
以弗所書裡面說.
撒旦是這個世界裡的王.
當願意去拜祂.
去俯伏在祂面前的時候.
能夠享受榮華富貴.
他也能夠有超自然的能力.
只要人願意俯伏拜祂的時候.
他有預知的能力.

$^{881}$他有交軌的能力.
去讓人有其他的方法.
很直接地得到權利.
金錢和一些慾望裡面的滿足.
特別是在香港一個物質社會的環境裡面.
這些物質的誘惑顯得特別吸引.
香港寸金尺土.
如果忽然有一天有人帶你去太平山.
然後叫你看著整個維港.
全部的大廈.
他告訴你.
只要你和我一起工作.
我將整個香港送給你.
整個香港.
一個香港的單位也要幾百萬.
是不是?.
整個香港不知道有多少萬億的價值.
更何況摩根士丹利說到.
是一張萬國.
全世界很重要的.
一些很出名的城市.
無論是倫敦,紐約,悉尼,北京,東京.
都說只要你拜一拜.
你就能夠擁有這一切.
真的很吸引.
早前在他的書卷裡面.
他說到面對著肉體的情慾.
啞木的情慾.
並今生的驕傲.
其實這些東西真的沒有人能夠抵受得住.
可能我們今天能夠抵受得住.
可能因為那些吸引.
不夠吸引.
使我們沒有去陷入那些試探.
但是如果那些試探再吸引一點.
利潤再豐厚一點.
再吸引.
再主動靠近你.
隔離的時候.
其實真的很難抵受得住.

$^{921}$不過耶穌的回答.
他是在引用《申命記》.
《申命記》第六章第十三節他引用.
我讀上下文給大家聽.
《申命記》第六章第十節.
(1)你的神令你進他向你列祖亞伯拉罕以撒瓦戈.
起誓應許你的地.
那裡有城邑又大又美.
非你所建造的.
有房屋裝滿各樣的美物.
非你所裝滿的.
有作成的水井.
非你所作成的.
還有葡萄園,橄欖園.
非你所栽種的.
你吃了並且飽足.
那時你要謹慎.
免得你忘記將你從埃及地.
遺留之家領出來的耶和華.
你要敬畏耶和華你的神.
事奉他.
指著他的名起誓.
不可隨從必神.
就是你們四位國民的神.
因為在你們中間.
耶和華.
你神是既邪的神.
唯恐耶和華你神的怒氣向你發作.
就把你從地上除滅.
昔日這群以色列人.
臨入迦南地之前.
摩西提醒這群以色列人.
迦南地上帝已經為你預備了.
很豐富的筵席.
有很豐富的資源.
有美麗的房屋.
或者很豐富的土產.
有很多你一進去.
能夠享受的一些資源.
不過你享受這些.

$^{961}$豐富的資源的時候.
記住不要忘記.
你擁有的一切.
是由神而來的.
不要隨從外邦的人的習慣.
習俗去敬拜那些偶像.
單單要事奉耶和華.
不過當這群以色列人.
進入迦南地.
我們看到士師記的時候.
他們跟隨了其他外族的人.
敬拜偶像.
隨從當地人的習俗.
恨一些邪人.
離棄了上帝.
恨耶和華眼中看為惡的事情.
當敬拜的對象錯誤的時候.
他們的生命.
他們的行為.
越來越走偏.
離開上帝的心意.
對耶穌來說.
當拜主來的神.
單要事奉他.
耶穌拒絕魔鬼撒旦的計劃.
定要跟從上帝的心意.
上帝的計劃是要耶穌.
來到這個世界為人的罪.
走在十字架的苦路裡面.
過程裡面要經歷人的背叛.
鞭打羞辱.
最後要死在十字架.
被埋葬.
三天之後復活升天.
然後等待第二次來到這個世界裡面.
坐在榮耀的寶座上面.
萬國萬民都要去敬拜他.
上帝的計劃.
是要讓耶穌基督先受苦.
然後能夠讓人在罪惡裡面拯救出來.

$^{1001}$最後萬國萬民一起去敬拜他.
不過魔鬼撒旦提出的試探.
就是叫耶穌直接拜他.
萬國萬民.
所有榮華都是屬於他.
這是一個捷徑.
叫耶穌不用上十字架.
就能夠直接享受了.
彌撰應有的權柄.
他想耶穌放下上帝的心意.
不要受苦.
這麼辛苦做什麼呢?.
十字架在其他人眼中.
捍衛一個如注的拯救.
為什麼耶穌一定要受苦呢?.
為什麼不直接赦免別人的罪?.
其實死有很多方法.
為什麼要選擇十字架.
這麼羞辱的痛苦的方法呢?.
或者選擇門徒也有很多種.
你不選擇尖子.
為什麼要選擇這些漁夫.
或者這些在社會邊緣的人呢?.
當耶穌去拜魔鬼的時候.
表達他認同魔鬼的方法.
成為魔鬼計劃的一部分.
神安排耶穌來到這個世界.
是要拯救人的罪.
這個方法在其他人眼中.
好像要用很長的時間.
要經過很長的掙扎.
而這些效果好像很不明顯.
當耶穌升天的時候.
這群門徒還要四散.
沒有一個信任可靠的一個人.
好像要用時間一點一滴.
那一刻又很慢.
魔鬼撒旦的計劃就不一樣.
果效很大.
能夠很簡單很快直接享受成果.

$^{1041}$更重要的是不用經歷痛苦.
不用經歷掙扎.
你用他的方法.
你很容易上位.
你很快地直接去享受到各式各樣的榮華富貴.
都可以給你.
如果你是耶穌你會怎麼選擇呢?.
耶穌在傳道之前受魔鬼的試探.
其實這個試探.
第一節裡面他講到受聖靈的引導.
是聖靈容許.
然後在天使裡面又侍候他.
這個試探不單只是魔鬼的試探.
更加是上帝的考驗.
他重新問耶穌幾個很重要的問題.
耶穌說你沒有能力.
沒有權並沒有恩賜.
但是你的恩賜能力權並.
是用在服侍其他人身上.
還是滿足自己的需要?.
當耶穌走在十字架傳道的路上.
面對各式各樣的挑戰困難.
有人會挑機.
有人會不滿.
有人會將他陷入苦難.
更加重要的是走在十字架路上.
在這些苦難中.
耶穌你會不會懷疑上帝的同在?.
耶穌有能力.
耶穌有聰明有智慧.
他能夠用一些不同的方式.
特別是很吸引人的目光的方法.
去用一些很戲劇性的方法.
去讓人有改變.
究竟耶穌在用的是.
魔鬼撒旦快捷方面的方法.
還是在用上帝的方法?.
其他人好像很蠢.
一步一步踏踏實實.
好像中中夕夕一樣.

$^{1081}$走在十字架的路上.
三個試探.
不單只是試探耶穌.
其實也在問今天的我們.
神給恩賜我們.
是要做什麼的?.
在困難的時候.
我們會不會懷疑神?.
在服侍上帝的裡面.
我們會用自己的方法.
魔鬼方法.
還是上帝的方法?.
其實三個試探.
最終可能都是在問一個試探.
一個問題.
我們生命最重要的主權.
在誰手上?.
誰是我們生命的主事人?.
是我們自己?.
是魔鬼?.
還是上帝?.
我們願不願意將我們手中.
所有的計劃.
所有的想法.
傳言.
去奉獻.
交在上帝的手中.
告訴上帝.
神啊,我也是一個無用的僕人.
求你用我.
用在你的手中.
成為其他人的祝福.
最後的階段.
我想分享一個宣教士的見證.
這個宣教士叫做斯可福.
他的譯名是.
他在18,19世紀.
跟隨著大質生中國內地的警察.
來到中國服侍.
當時斯可福.

$^{1121}$他很聰明.
他小時候.
他15歲已經入大學.
入大學的過程中.
他讀一些最優秀.
其他人眼中的項目.
最羨慕的一個學科.
就是醫科.
他讀完醫科.
成為了醫生.
在不同的地方.
他有些研究.
有些成果.
和獎學金.
當他在牛津最高的學府.
畢業的時候.
他同時畢業的時候.
都有七個獎牌.
在他的手中.
他在醫學中.
都有很大的成就.
有些教授告訴他.
不如你.
會不會留下來.
在大學中去.
去做一些研究.
能夠帶來更多的祝福.
但是在1880年的時候.
當時的斯可福.
29歲.
他聽到中國內地的需要.
他放下了.
很優厚的條件.
他去到中國內地.
去成為宣教士.
因為他的聰明.
他很快學會了當地的語言.
我不知道是廣東話還是普通話.
但是很快地學會了語言.
他在太原的第一年裡面.

$^{1161}$他醫治了超過千個病人.
在門診裡面.
而且當中.
有些人是施行了麻醉手術.
在當時醫療環境很缺乏的情況下.
他施行這些手術.
而且能夠成功.
在第二年的服侍.
就是醫病的過程裡面.
他醫治了超過六千多個病人.
而當中有很多人都是被狼咬傷.
他進行了過百個手術.
而且這些手術.
有超過四十個.
在麻醉的情況下進行.
斯可福在太原的服侍期間.
他很能夠使用上帝的恩慈.
他用他的才幹能力.
他也很尊重當時.
中國老百姓裡面的文化習俗和需要.
他看到當時中國裡面最大的困難.
就是受著鴉片毒癮裡面的煎熬.
他成為了幫助人.
戒毒的一個先驅.
他設立了一些戒毒所.
很希望用醫療的方式.
去幫人戒毒.
然後在當時的山西去服侍的地方.
鴉片的毒害都是很嚴重.
當時斯可福覺得.
不是辦法.
他覺得這個需要很大.
這個禍場很大.
單靠他一隻手.
做不完所有的事情.
所以他去告訴戴德新.
戴德新有沒有機會呢.
你去呼籲叫多一個七十個人.
七十個人來到當中.
能夠幫助這幫人離開水深火熱的環境裡面.

$^{1201}$在1883年的一個春天的晚上.
斯可福去祈禱.
他跪下.
他去求神.
求神去差派更多的高材生.
占子.
來到中國裡面去服侍中國人.
他跪在床邊裡面.
求神去差派更多的工人來到當中.
幾個月過去.
在一次醫病的過程裡面.
他不幸感染了白喉病.
然後很快.
三天之後就死了.
他當時死的年紀是32歲.
當他死的時候.
很多人都會覺得很可惜.
都會問很多問題.
32歲.
他原本可能可以在牛津享受最優厚的環境.
或者是條件.
或者生活裡面的需要都是不缺乏的.
為什麼要來到中國.
上帝用人.
為什麼不用盡他.
為什麼只在中國服侍了三年多.
這麼快就殉道呢.
為什麼他的服侍是這麼短呢.
可能都會問很多問題.
上帝在苦難裡面為什麼不救他呢.
不過上帝的計劃.
卻是很奇妙.
上帝垂聽了斯可福的祈禱.
因為斯可福是牛津大學畢業.
他帶動了大學裡面的基督徒運動.
在英國.
在美國裡面.
有過千人願意參與在中國服侍.
而且當中有劍橋七傑.
都是因為斯可福的見證原因.

$^{1241}$他們願意投身在中國服侍中國人.
今天我們生命裡面可能會遇到很多挑戰困難試探.
可能在問我們也是一個問題.
我們願不願意將我們生命的主權.
交在上帝的手裡.
給上帝用呢.
因為這個世代裡面真的有太多需要.
當我們打開新聞.
看到一個國家打仗.
看到某一個國家貧富懸殊.
或者不法的事很多.
我們的心又是怎麼想呢.
神有沒有感動我們去回應這個世代裡面的需要呢.
求神幫助我們.
讓我們靠著上帝的恩典勝過這些試探.
更加在生命裡面成為.
被上帝在黑暗世代裡面成為一個明亮的燈台.
我們一起祈禱吧.
親愛的主我們的神.
你是一個聖潔的神.
你是一個榮耀的神.
你是一個充滿著慈愛的神.
當我們去到你面前的時候.
主要我們看到我們自己的軟的.
看到我們的不堪.
看到我們每一次面對著試探.
我們都很容易跌倒.
神我們祈求你.
這些試探在我們生命裡面很真實.
這些真實的試探可能是我們肉體裡面的需要.
可能是我們信心裡面的需要.
甚至乎可能是我們願不願意將我們所擁有的降服在裡面.
主在我們面前我們承認.
我們都是一群沒辦法自救的人.
離開了你我們什麼都不能做.
主我們祈求你.
讓我們用力地抓緊你.
讓我們用力地在每一次面對生活裡面.
在各種試探裡面的時候.
我們想起你.

$^{1281}$以致我們能夠很體寵我們.
基督徒很尊貴的上帝兒女的身份.
我們不會做得罪你的事情.
我們更加不單止不做得罪你的事情.
我們更加要去做討你氣的事情.
今天在這個五花.
就是花花世界裡面.
很多事情都很吸引人的眼目.
有很多人都在問.
就是傳福音要慢慢帶人順處.
要關心人一步一步地去教他們看聖經.
我說這些方法很慢.
還有沒有效的.
基督上帝是不是還在這個世界.
主阿大我們承認你是這個世界裡面的主.
你仍然是這個世界裡面的王.
只有你才能夠拯救人的生命.
上帝你仍然會在這個世代裡面.
用神蹟祈禱.
大神我們仍然願意配合你的工作.
用你的方法.
可能過程裡面會慢一些.
過程裡面會經歷掙扎.
會經歷痛苦.
大神我們願意降服在你裡面.
求你親自去幫助我們.
給我們力量回應這個世界的需要.
求你親自在當中幫助我們.
在我們祈禱.
奉耶穌得勝之名求.
\newpage



\section{約翰福音 8:1-11}
\label{sec:YIG0i1d0dV8}
\textbf{2022年10月30日崇拜直播｜吳真真牧師｜以石為記:誰敢拿起審判之石｜約翰福音8\_1-11}
\newline
\newline
連結: \href{https://youtube.com/watch?v=YIG0i1d0dV8}{\texttt{https://youtube.com/watch?v=YIG0i1d0dV8}} ~~~~ 語音日期: 2022-10-31
\newline
\newline
\hyperref[sec:862eUMKnVnc]{\small{< < < PREV SERMON < < <}}
~
\hyperref[sec:index_chronic]{\small{[返順時目]}}
~
\hyperref[sec:ucA_hp7yY_0]{\small{> > > NEXT SERMON > > >}}
\newline
\newline
約翰福音 8:1-11
\newline
\begin{longtable}{cl}
\hline
\hline
章節 & 經文 (和合本修訂版)\\
\hline
8:1 & \begin{tabularx}{0.7\textwidth}{X} 耶穌卻到橄欖山去。 \end{tabularx} \\ \\ \relax
8:2 & \begin{tabularx}{0.7\textwidth}{X} 清早，他又回到聖殿裡。眾百姓都到他那裡去，他就坐下，教導他們。 \end{tabularx} \\ \\ \relax
8:3 & \begin{tabularx}{0.7\textwidth}{X} 文士和法利賽人帶著一個犯姦淫時被捉的女人來，叫她站在當中， \end{tabularx} \\ \\ \relax
8:4 & \begin{tabularx}{0.7\textwidth}{X} 然後對耶穌說：「老師，這女人是正在犯姦淫的時候被捉到的。 \end{tabularx} \\ \\ \relax
8:5 & \begin{tabularx}{0.7\textwidth}{X} 摩西在律法書上命令我們把這樣的女人用石頭打死。那麼，你怎麼說呢？」 \end{tabularx} \\ \\ \relax
8:6 & \begin{tabularx}{0.7\textwidth}{X} 他們說這話是要試探耶穌，要抓到控告他的把柄。耶穌卻彎下腰，用指頭在地上寫字。 \end{tabularx} \\ \\ \relax
8:7 & \begin{tabularx}{0.7\textwidth}{X} 他們還是不住地問他，耶穌就直起腰來，對他們說：「你們中間誰沒有罪，誰就先拿石頭打她！」 \end{tabularx} \\ \\ \relax
8:8 & \begin{tabularx}{0.7\textwidth}{X} 於是他又彎著腰，用指頭在地上寫字。 \end{tabularx} \\ \\ \relax
8:9 & \begin{tabularx}{0.7\textwidth}{X} 他們聽見這話，從老的開始，一個一個都走開了，只剩下耶穌一人和那仍然站在中間的女人。 \end{tabularx} \\ \\ \relax
8:10 & \begin{tabularx}{0.7\textwidth}{X} 耶穌就直起腰來，對她說：「婦人，那些人在哪裡呢？沒有任何人定你的罪嗎？」 \end{tabularx} \\ \\ \relax
8:11 & \begin{tabularx}{0.7\textwidth}{X} 她說：「主啊，沒有。」耶穌說：「我也不定你的罪。去吧！從今以後不要再犯罪了。」〕 \end{tabularx} \\ \\
[1ex]
\hline
\hline
\end{longtable}
$^{1}$以下我們以聆聽聖道繼續去敬拜我們的主.
首先由陳王惠敏姐妹為我們讀出《讓福音》第八章一到十一節.
然後由吳真真牧師為我們宣講.
題目是「以石為記,誰敢拿起審判之石」.
《讓福音》第八章一到十一節.
「於是各人都回家去了,耶穌卻往橄欖山去,清早又回到殿裡,眾百姓都到他那裡去,他就坐下教訓他們,問是何法理賽人帶出一個行淫事被拿的婦人來,叫她站在當中,就對耶穌說:.
夫子,這婦人是正行淫事被拿的,若在法律上吩咐我們,把這向的婦人用石頭打死,你說該把她怎麼樣呢?.
他們說者說,乃他試探耶穌,要得著告他的把柄.耶穌卻彎著腰,用指頭在地上畫字.他們還是不住地問他,耶穌就直起腰來對他們說:.
你們中間誰是沒有罪的,誰就可以先拿石頭打他?.
於是又彎著腰,用指頭在地上畫字.他們聽見者說,就從老到少,一個一個的都出去了,只剩下耶穌一人,還有那婦人仍然站在當中..
耶穌就直起腰來對她說:婦人,這些人在哪裡呢?沒有人定你的罪嗎?.
她說:主啊,沒有.耶穌說:我也不定你的罪.去吧,從此不要再犯罪了..
請獨筆..
各位大英姐妹平安..
《以石為記》系列由九月開始,到今天是第七講,我們總共有十講..
前面的五講是《九曰》,上次是兩位傳道和我們講耶穌受試探..
今天來到第七講,《約翰福音》..
剛才讀經為我們讀出之後,我們發覺這段經文大家很熟悉..
這次我們講的這塊石頭,其實是關於舊約時代摩西律法審判下處死犯罪的人的一個刑具..
這塊石頭可以代表審判,可以代表公義..
今時今日,在英文世界裡,我不擅長英文,不過我查到這句諺語,叫做"Cast the first stone",摘出第一塊石頭..
原來這句諺語,就像我們的成語,會說典故在哪裡..
這句典故是來自聖經第八章第七節,耶穌問的一條問題:.
你們中間誰沒有罪的,誰就可以拿石頭打他?.
誰就可以先拿石頭打他?.
所以這句諺語往往在現今世界用這個字,代表了第一個作出批評,責怪的批判的一個字..
在今時今日,其實你有沒有發覺,我們人人都準備向人摘石頭..
我們可能每一天都對不同的人,頻道品足,加以評論,甚至論斷他們..
我們有時聽到有些人說有些毒舌,嘴巴一張出來都是責怪人,罵人,責備人..
有時不是當面說的,而是在心裡面,對人做批評,在心裡面有很多挑剔,甚至定罪..
所以有些時候我們都做了摘石頭的人,不過更有可能的就是我們都會被人摘石頭..
今天這段經文對我們來說到底在這個時代裡面有甚麼提醒呢?.
不如讓我們一起祈禱吧!.
求你預備我們的心,光照我們,虛心聆聽,讓你的光照亮我們內心的黑暗,.
讓我們以真誠去回應你今天對我們的提醒,.
又讓你的恩典和愛浮避著我們,讓我們真的在崇拜的裡面,能夠去領受主的更多恩惠..
禱告教徒,奉主名求,阿們..
我們來看看第一點,究竟這些摘石頭的人的內心是怎樣想的呢?.
我們先看第一至六節上,第一,第二節就形容了這件事的背景是怎樣的呢?.
原來這件事發生在主白銜節期間,或者再推後一點,就是耶穌受難之前,.

$^{41}$祂每一天早上都會去到聖殿裡面去教訓人..
所以當時有很多群眾在聽耶穌講解聖經的真理..
不過在這個時候,有一些文士和法理才人走來找耶穌,.
但是他們來不是要聽耶穌的教導,而是他們帶著一個被人捉姦在床的婦人來到祂面前,.
去到一大群人的中間,向耶穌發出一條問題,按照摩西的律法,書的命令,.
說要將這個犯姦淫的婦人用石頭打死,.
你說我們應該怎樣做?.
其實他們好像表面上是按照神的公義去執法,.
但是你總是覺得有點奇怪,疑點重重..
首先他們問的問題,他們是文士,文士是解經的學者,.
他們是法理才人,他們是教人怎樣行律法的,.
他們沒有理由不知道聖經要怎樣處置的..
在《新命紀》22章23至24節說,若果有處女已經許配丈夫,.
有人在城裡遇見他,遇他行淫,.
你們就要將這兩個人帶到城門,用石頭打死..
他們一定會回答,不用問耶穌了..
第二,當時這種犯姦淫要去定罪,往往會拉到猶太的公會去審判他們..
何時輪到這個在他們眼中,完全沒有受過正統訓練的耶穌去做判決?.
這是第二個疑點..
第三個疑點,你試想一下,行淫的時候會不會只有一個人?.
不會的,那為何只帶了一個女的來?.
同時行淫的男人去了哪裡?.
第六節很快就告訴我們,其實因為這班人來者不善,.
他們來者要試探耶穌,要找他的把柄去指控他..
你可能會問,為何問這條問題,就可以找到把柄?.
他問這條問題是,按照摩西律法書說,行淫的人要用石頭打死,.
耶穌,你說我們應該怎樣做?.
就是這條問題..
問題在於,耶穌回答,好啊,用石頭打死,.
或不要用石頭打死,.
回答是又死,不是又死..
假如耶穌說是,.
閩氏法利賽人就可以在羅馬官府面前去控告耶穌,.
因為猶太人是被羅馬帝國去管治的,.
猶太人是沒有執行死刑的權利的..
你想想,約翰福音十八章說到耶穌被審,.
被人拉到羅馬巡撫彼拉多面前,.
彼拉多對猶太人說,你們就按照猶太人的律法去審判他..
但是猶太人怎樣回答?.
他們回答,我們沒有殺人的權柄..

$^{81}$所以,如果耶穌按照摩西的律法,.
用石頭打死這個婦人,.
他們就可以拿到把柄,.
可以指控耶穌,說耶穌執行私刑,.
甚至再進一步指控耶穌,.
是有心推翻羅馬王國,要自立為王..
答是就死了,.
答不是可以嗎?.
假如耶穌答不是,.
就不要用石頭打死這個婦人..
那麼,文士和法理賽人就可以趁機煽動群眾,.
指控耶穌公開教導人不守摩西的律法..
你想想,猶太人亡國以來,.
摩西的律法不單單是宗教規條,.
對於他們來說,.
簡直是他們這些猶太人的身份代表,.
是這些屬於猶太人,.
猶太人能夠得以維繫的一個中心..
而耶穌當時就在這些群眾裡面,.
去講解摩西的律法,.
他是一個可能是新冒起,.
很有影響力的律法教師,.
他竟然叫人不守摩西的律法?.
這樣他們就會令到耶穌在群眾裡面失去公信力,.
甚至成為群眾的敵人..
所以,原來問這條問題,.
真是有玄機的..
法理賽人和文士,.
原來真是在這裡想陷害耶穌,.
拿他的把柄..
其實這群法理賽人和文士,.
就是一群經常預備隨時砸石頭的人,.
他們好判斷,好審判,.
隨時向人攻擊行刑..
我們看看日常生活裡面,.
他們會有什麼表現呢?.
第一個表現就是,.
會經常抓人的錯處,.
不知道當中有沒有老師呢?.
我掃一下,沒有人回應我..

$^{121}$老師,你有沒有試過有些小朋友,.
尤其是小學生,.
很多時候很流行,.
老師,老師,陳小明今天沒帶書..
老師,張小玲你怎麼還在這裡?.
張小玲今天拿支鉛筆撩我..
是不是這樣?.
你記不記得你小學時有這樣的同學?.
沒有人回應我,.
今天你們怎麼沒有反應?.
記得,記得,多謝鄭鵬叔..
其實他們是隨時找人的錯處,.
預備去告他,預備去提案..
這群法理賽人就是這樣,.
聖經裡面經常提到,.
要找耶穌把柄..
經文剛才去到第七章,.
我們沒有看,.
但第七章尾其實是描述到,.
那群法理賽人在質疑耶穌,.
去說摩西律法的權柄,.
質疑他在當中影響了人,.
於是他們就找機會,.
找把柄,.
嘗試從他說的話裡面,.
看看有沒有抽籤的地方,.
然後就找警察去捉拿他..
但是,第七章尾說,.
警察去捉拿他的時候,.
覺得這個人說的道很特別,.
很有力,.
於是他們都找不到把柄去捉拿他..
你想像一下,.
接著怎麼樣?.
八章就來了,.
他們的心水不寧,.
晚上睡不著,.
於是就去捉拿耶穌的把柄,.
竟然被他們找到一個行淫的婦人,.
捉她去到耶穌面前,.

$^{161}$找耶穌的把柄..
但是我覺得很奇妙,.
你再看聖經,.
你看第八章第一節,.
耶穌在講完道,.
耶穌在他向群眾當中教導完之後,.
他去橄欖山,.
他面對四面的攻擊,.
面對人往往想抽清他,.
耶穌是去橄欖山,.
他安靜下來,.
他去禱告,.
他去親近神,.
他去享受這個休息..
然後第二天一早,.
他又精神抖擻,.
去到群眾當中去做教導..
你想不想到,.
他們很忙,.
那群法利塞人,.
和耶穌都是很忙,.
不過他們做的事情,.
他們的心靈成為一個很強烈的對比..
第二,我們發覺,.
這些喜歡砸石頭,.
指控人的法利塞人,.
他們背後的動機是什麼?.
原來他們的動機,.
是為了要武裝好自己,.
去保護自己..
為什麼我這樣說呢?.
我們看到,.
其實這群法利塞人,.
在《約翰福音》第七章說到,.
他們去質疑耶穌,.
你是否一個真真正正從神而來的教師?.
耶穌怎樣說他們?.
耶穌在七章第十九節說,.
「摩西豈不是傳律法給你們嗎?.
你們卻沒有一個人守律法,.

$^{201}$你們為何要殺我?」.
這群宗教領袖,.
被耶穌看穿了,.
而且耶穌毫不留情地指責他們..
自然而然,.
他們為了捍衛自己的地位,.
他們就產生了一些防衛的機制,.
去保障自己..
他們知道耶穌會遙控他們的地位,.
和身份,.
所以他們力證自己是對的,.
而耶穌是錯的..
耶穌錯,他們對,.
他們才有位置..
結果他們就找了這個行刑的婦人,.
來公審,.
問耶穌是否守摩西的律法,.
就是要證明他們是真正的基督徒,.
而不是守摩西的律法,.
就是要證明,.
「你說我不守,.
那我又看看你怎樣守.」.
心裡就越走越歪,.
甚至到了一個冷酷無情的地步,.
去當眾收入一個.
手無寸鐵的婦女,.
而背後的動機,.
並不是為了伸張正義,.
並不是為了痛恨罪惡,.
而是為了自己的利益..
他們這麼冷酷無情,.
其實都真的很可恥..
而事實上,.
當他們帶著這個行刑的婦人,.
來到會眾面前,.
表面上他們要向婦人定罪,.
向婦人砸石頭,.
但是真正要砸石頭的對象,.
他們是想砸向耶穌..
弟兄姊妹,.

$^{241}$我們都要留意,.
我們有些時候都要小心,.
我們會不會都會.
好像這班法利賽人一樣,.
向人砸石頭呢?.
我們要留意,.
其實有時我們會去.
將人去指責,.
去批評,.
甚至將人去dehumanize,.
去人性化..
那個目的是什麼?.
看不到他的需要,.
他成為你的工具..
那個目的就是試圖.
將自己變得更加優越,.
變得我企高人一格..
我們有時真的會不經意地流露,.
其實我比別人有價值..
然後我們就會有很多理由.
去證明這一點,.
我們沒有犯罪,.
我們不是通姦,.
我又不是吸毒,.
我又沒有立心不良,.
我們還要是教會的成員,.
我很積極行十一奉獻,.
我很積極去事奉,.
參與教會活動,.
我比其他人更有價值去活..
我記得基督教今日報.
曾經跟我們提過,.
法理賽人有五個特點,.
不如讓我們去聽一聽,.
審視一下,.
自己會不會在我們信仰的歷程裡,.
不經意地流露了這些特質..
第一,就是盲目地高舉傳統,.
例如他們堅守禁食,.
十一,各種宗教儀式,.

$^{281}$不要緊,.
但卻忘卻了信仰的真正本質..
第二,很多時候用人的行為去評價他的信仰,.
例如他是否積極地每星期來送拜,.
或者參與教會活動,.
來評價這個人的信仰,.
而不是真真正正了解他的性格,.
他的生命..
第三,如果法理賽人有些特質,.
就是經常盲目地看別人的信仰,.
而不是盲目地看自己的信仰,.
他們會批評約翰,.
說他不吃不喝,是否被鬼附,.
他們批評耶穌,.
說他又大飲大食,.
是否很愛吃喝,.
他們會用自己的尺度去量度別人..
第四,就是會變得很僵化,.
沒有甚麼彈性,.
在一個狹隘的宗教世界裡生活,.
去排除一些新的思維,.
而且我們說是很律法主義..
第五,就是非我族類的狹隘思想,.
他很看重組織的認同,.
懷著一個叫「一旦你不是我的族類,.
我就要排斥你的」的濃厚觀念..
這也是值得我們深思的..
第六,就是我們知道.
拋砸石頭的人的思維後,.
我們看看耶穌如何回應他們..
我們發覺在第六節下至第九節上,.
我們看到耶穌問了一條問題,.
於是就沒有人敢拿起這個石頭..
雖然法利賽人和文士將一個行淫的婦人的死活.
作為這個場景裡的焦點,.
他將她放在聖殿的中間,.
但是耶穌卻很巧妙地將這個焦點,.
移到群眾當中..
我們看看耶穌如何回應那條.
「答是又死,答不是又死」的問題..

$^{321}$首先,耶穌彎腰,.
用指頭在地上畫字,.
即是說耶穌不是用說話回答他們的問題..
有些解經學者告訴我們,.
不要以為耶穌是在「買時間」,.
什麼是「買時間」?.
他想不到如何回答,於是就在拖延時間,.
然後就自己在想如何回答..
不是的,耶穌已經啟動了回答的模式,.
因為他的回答帶來了一個很好的回應..
耶穌彎腰,人們沒有放過他,.
繼續不斷地問,不斷地死纏砸打,.
耶穌於是回答他們,.
說了一句話:.
「你們中間誰沒有罪,.
就可以先拿石頭打他.」.
說完這句話,又彎腰,又繼續畫字..
有很多人問,到底要畫什麼字?.
我看到釋經書有不同的演繹,.
有寫經文,有罪字,.
不過我不探究了,.
原因是我覺得作者不寫,.
就不重要了,.
沒有影響我們去理解這段經文,.
所以我們不去問..
不過重要的意識是什麼呢?.
重點是耶穌蹲下寫字,.
其實是向群眾表達了什麼呢?.
耶穌的舉動表明了他支持摩西的律法,.
不過同時他也避免了行私刑,.
即是他完全沒有被法利賽人和民事陷阱打中他,.
害到他..
耶穌彎腰寫字,.
其實目的就是將人群情洶湧的情緒,.
慢慢舒緩下來..
當一些人群情洶湧,.
但被他質詢,被他圍攻的對象,.
沒有動氣,.
而且還不拙不拙地在地上畫字,.
耶穌在等候一個適當的時機,.

$^{361}$等候群眾的心理狀況預備好,.
向他們心靈的最深處,.
問一條最針對他們的問題,.
你們中間誰沒有罪,.
就可以拿石頭打他..
在原文中有些希臘文的譯本,.
第九節有這樣說,.
「聽見者說」,.
本來就是從老到少,一個一個走,.
有些譯本中間加了一句,.
「並且被良心責備而後」,.
由老到少,一個一個走..
有時我們在輔導室裡,.
大家想我們做輔導就不斷說話,.
室內充斥著很多說話,.
但在輔導中有一種很厲害的技巧,.
就是沉默的力量,.
真的不可以忽略,.
因為沉默可以使當事人有空間,.
不受干擾地經驗和處理他們的情緒和想法,.
可以向他們心靈的最深處進發,.
面對真正的自己..
所以耶穌問的這條問題,.
你們中間誰沒有罪,.
就可以拿石頭打他,.
是一個非常非常有智慧的問題..
成功地將這班文士和法理賽人,.
由怨告變成被告,.
由平時高人一格,.
拉下來,.
來到上帝的審判寶座面前,.
他們基本上和犯罪的奸淫婦人是同等的,.
是沒有分別的..
在《申命記》第17章第6至7節說,.
「要憑兩三個人的口作見證,.
將那當死的人致死,.
不可憑一個人的口作見證將他致死,.
見證人要先下手,.
然後眾民也下手將他致死,.
這樣就把那惡從你們中間除掉.」.

$^{401}$這裡提到原來舊約的律法,.
誰是行刑者?.
是見證人..
是第一個,有兩個人,一個都不行,.
要兩個見證人,.
同時見證那個人是犯罪,.
然後到真的要執行刑罰的時候,.
就由第一個和第二個見證人先用石頭打他,.
然後其他人才可以用石頭..
這一班的領袖,.
沒有辦法去做第一個和第二個.
拿起石頭去打這個淫婦的見證人,.
因為當耶穌在這麼寧靜,.
追問他們內心深處的問題下,.
他們不能逃避一個真實的自己..
當他們想到他們這麼不堪的動機,.
他們想到他們用一個侮辱婦女的方法,.
作為他們陷害耶穌的工具,.
他們逃不過良心的責備,.
於是他們要離開,.
他們不敢走出來認自己沒有罪,.
唯有毫無顏面,一個一個靜悄悄地流走,.
離開這個被他們變成公審法庭的聖殿,.
這個敬拜神的地方..
你會問,為何會從老到少呢?.
我覺得很簡單,.
年紀小的當然是看老人家的頭,.
那些領袖,那些文士,法理人,是長輩,.
我們經常聽他們德高望重地教導我們聖經真理,.
我當然是看他們的頭,.
他們離開,我就跟著離開..
但同樣有一樣我覺得很真實的,.
就是年紀大的人會較容易意識到自己的罪,.
年輕人沒有那麼強烈意識,.
為何這樣說呢?.
我用自己做例子,.
我記得我初中的時候信了主,.
導師教我每晚要認罪祈禱,.
你們有沒有人教你們?.
有,是吧?.

$^{441}$Maggie,你點頭,其他人沒有嗎?.
我很記得,我小時候認罪祈禱,.
有一段時間我覺得很困擾,.
因為我低頭祈禱,.
然後就想自己有甚麼罪,.
我想了很久都想不到,.
然後我就跟神說,.
「神啊,我又想不到今天有甚麼罪,.
求你寬恕我隱而未然的罪.」.
你們試過嗎?.
試過,這班人回應得很好..
上帝是真的,我們年輕的時候很純真的,.
很單純的,.
我們會覺得自己沒有甚麼罪,.
不過很坦白說,.
到我們年紀大了,.
在信仰裡浸久了,.
聖的真理明多了,.
我就會發覺,.
其實我認真探究自己的心思意念的時候,.
我都有很多得罪神,得罪人的地方,.
是吧?.
各位靈姐妹,.
其實你不要以為我今天跟你分享這段經文,.
是因為我已經得罪了,.
我不是法理賽人,.
所以我教你們千萬不要做法理賽人,.
不是的,不是的..
我反而記得我小時候,.
剛剛信主沒多久,.
我家裡有人做了一件很差的事,.
我真的很生氣,.
我要跟他割席,.
我不想再見到他,.
我沒辦法寬恕他,.
但是我記得有一個姐妹,.
她很好,.
她問了我一句,.
她說你要跟他割席,.
你有沒有留意到,.

$^{481}$其實你也是一個罪人,.
耶穌有沒有跟你割席?.
耶穌怎樣對你?.
你願不願意用耶穌對你的愛,.
去愛這個親人?.
當我聽到這個姐妹這樣說的時候,.
我真的立即悔改,.
我真的很乖,.
我默默守候這個家人,.
直到他信主為止,.
但是,.
今天祂讓我預備這段經文,.
向大家宣講,.
祂好像蹲在地上,.
問了我一條很法人心聖的問題,.
植於我內心的深處,.
我不禁去問自己,.
我作為一個回應上帝呼召,.
成為主的僕人,.
我是一個傳道人,.
我是一個牧師,.
但是苦心自問,.
我有時也是一個砸石頭的人,.
我會為在人際關係裡的不舒服,.
去對人批評,.
作出控訴,.
有時也會為這些情緒困擾而傷害別人,.
我有時也會站高一格,.
用手指著別人,.
說別人不對,.
我經常教夫婦要清垃圾,.
不是家裡的垃圾桶,.
是兩個人之間吵架的垃圾,.
大家的不滿,.
大家的不開心,.
我經常勸夫婦清垃圾,.
但是我發覺,.
其實我們也要向神清自己的垃圾,.
我們應該學校,.
這些人,.

$^{521}$耶穌製造了這樣的情景,.
讓我們安靜,.
我們這些忙碌的人,.
我們在工作裡,.
我們在事奉裡,.
這麼密切地處理事奉的人,.
我們有沒有真的撥出一些空間時間,.
認真的同神親近,.
在主的裡面,.
讓祂幫我們清垃圾?.
第三,.
我們發覺沒有人敢砸出這個石頭,.
最後,.
連有資格砸出這個石頭的耶穌,.
祂也沒有,.
第九節下至第十一節,.
最後所有人離開了,.
只有耶穌和被捉拿的行淫婦人,.
只剩下他們,.
耶穌問他們,.
「那些人在哪裡?.
他們沒有定你的罪嗎?」.
婦人說,.
「主啊,沒有.」.
耶穌說,.
「我也不定你的罪,.
去吧,從此不要再犯罪了.」.
對於他們的對話,.
我想有幾個重點,.
我很想提出,.
我們看到其他人,.
因為看到自己的罪就離開,.
但是誰沒有離開?.
耶穌,.
因為耶穌沒有犯罪,.
祂不需要離開,.
第二,.
耶穌有權柄去審判,.
約翰福音第五章第二節這樣說,.
「父不審判任何人,.

$^{561}$即是天父不審判任何人,.
而將審判的事交予給子.」.
說明審判的權柄在耶穌那裡,.
祂有權柄去審判,.
祂也有權柄去赦免,.
馬可福音二章說到,.
有四個朋友帶著一個毯子,.
在屋頂上將毯子脫下,.
耶穌說了一句,.
「小子,你的罪赦了.」.
然後祂對法利賽人說,.
「我要讓你們知道,.
人子在地上有赦罪的權柄.」.
耶穌有權去審判,.
祂也有權赦罪,.
祂這次很清楚對面前的婦女說,.
祂不定祂的罪,.
祂不會用石頭打死她,.
原因是什麼?.
原因是約翰福音三章十七節,.
因為神差祂的兒子幹世,.
不是要定世人的罪,.
而是要世人因祂得救,.
耶穌可以審判世人,.
不過祂最終的目的是要呼召罪人,.
請他們認罪悔改,.
因而得到新的生命..
約翰福音一章十四節說,.
「道成了肉身.」.
耶穌成為肉身來到這個世界,.
住在我們中間,.
充充滿滿地有恩典,有真理..
因為全是恩典,.
因為上帝的愛,.
祂包容這個人,.
你要記住,.
耶穌在祂的對答中告訴我們,.
祂包容這個人,.
但是祂沒有包容罪,.
所以祂給予他悔改的機會,.

$^{601}$給予他一個不用再被罪捆綁的新生命..
耶穌在這裡說,.
叫他不再犯罪,.
不是只是說不再犯奸淫,.
而是希望他能夠請他離開罪惡..
弟兄姊妹,.
耶穌基督對於這個婦人的罪的處理,.
都是很值得我們認真的思考的..
我們在做,.
當我們要向人作出判斷,.
作出批評的時候,.
我們不代表我們要當她沒有罪,.
不代表我們要隻眼開隻眼閉,.
反而很多時候我們要提醒的是,.
其實我們都是蒙恩的罪人,.
我哪有資格去向她拋石頭呢?.
我哪有資格去定人的罪呢?.
但是我們不隻眼開隻眼閉,.
不是說我們不去對人作判斷,.
坦白說,.
我們做教睦,.
做領袖,.
我們不可以看著弟兄姊妹錯,.
然後當沒事,.
然後不作勸戒,.
如果是這樣,.
我們是沒有盡我們的本份,.
但是我們要留意.
《日漢福音》一章十七節,.
「律法本是藉著摩西傳的,.
恩典和真理都是由耶穌基督來的.」.
所以,領子妹,.
我們在真理的面前,.
上帝會用真理的標準來判別我們,.
但是背後,.
最終,.
祂是用恩典去承載..
我聽過一個這樣的故事,.
有一個美國的牧者,.
他星期日晚的晚堂崇拜,.

$^{641}$就問,.
「在我們當中,.
有沒有人想分享?」.
這個在外國很流行的,.
崇拜完之後,.
就請人分享..
突然有一個中年人站起來,.
他是一個很莊重,.
好像給人一個很得體的中年人,.
他就舉手,.
「我這次有些事情想分享.」.
這個中年人就開口說,.
「我三十五年以來,.
第一次再回到這個教堂裡面,.
我小時候,.
我上小學的時候,.
我是這裡的主要學生,.
今晚我來這裡,.
是想見我的主要學老師,.
Mr. Orton.」.
然後他就向著後面,.
轉過去,.
然後就說,.
「Mr. Orton,你好啊,.
你還記得我嗎?.
我是Hurley.」.
Mr. Orton最初就皺眉頭,.
但是看著看著,.
兩分鐘後,.
他就開始認得這個Hurley,.
然後他的臉上就展現了一個很燦爛的笑容..
他轉身面對Hurley,.
就繼續說,.
「Mr. Orton,.
我今天來這裡是要告訴你一件讓你很開心的事情,.
我是特地從很遠的地方開車來這裡,.
參加這裡的崇拜,.
因為我很想求你的原諒,.
你還記得嗎?.
當年我上主要學的時候,.

$^{681}$我們一群人,.
從頭到尾,.
卑鄙到極,.
做盡很多搗蛋的事情,.
搞盡很多惡作劇,.
而我就是那個帶頭的首領,.
我是真正的罪魁禍首,.
你還記得嗎?.
在我剛搬來這裡,.
去其他社區之前的幾個星期,.
我是特別的頑皮,.
特別的討人厭,.
你還記得嗎?.
你知不知道,.
最後回來的那一次,.
我下課了,.
我已經走了,.
但是我醒來,.
咦,我有些東西還沒有拿走,.
於是我回頭,.
因為我不再回來了,.
我回來的時候,.
我發現到,.
你坐在課室的面前,.
最前面,.
你在哭,.
我聽到你在那裡祈禱,.
我聽到你為我祈禱,.
我聽到你向上帝求幫助,.
你知不知道,.
那天其實我站在那裡很久,.
不過我沒有勇氣進來,.
於是我走了,.
接著那五年,.
每當我想起你那個禱告的樣子,.
我的心就覺得很窄心,.
有一天,.
突然間,.
我覺得我站在神的面前,.
神讓我去看到我的劣行,.

$^{721}$我在班裡面多麼的搗蛋,.
多麼的傷害人,.
這件事令我非常痛苦,.
於是我向神祈求,.
將我的生命獻上給他,.
今天,.
我已經是社區裡面的一間教會的長老,.
並且一直為我祈禱,.
所以我今天特地來,.
請你原諒,.
也告訴你,.
我在神裡面很生性地長大,.
弟兄姊妹,.
我們都是蒙恩的罪人,.
其實我們很多時候,.
都需要別人對我們寬容接納,.
是嗎?.
同樣,.
我們需要以耶穌基督的愛和包容,.
有真理,.
有恩典地去寬容和接納我們身邊的人,.
今天我們說到納石,.
以石為記,.
我們拿著這塊石頭,.
這塊我們說是審判公義的石頭,.
我們要像法利賽人一樣,.
用這塊石頭坐在高高在上的審判座上,.
去帶來的,.
是一些腐敗的信仰生活,.
還是一些對人的絕望,.
還是我們要像耶穌一樣,.
耶穌對我們的審判之石,.
是要叫我們因為律法而得生,.
以一個沒有秩序的混沌生命,.
轉化成為一個豐盛,.
有目標,.
被耶穌掌管的生命,.
我們或許,.
今天聽這篇道,.
我們都擇心,.

$^{761}$我們會不會有時都會站在這個擇石頭的位置呢?.
同樣,.
耶穌告訴我們,.
祂都會用永遠的恩慈,.
永遠的愛去接納我們,.
我們不用用我們自己的能力,.
在神的裡面得到肯定,.
我們其實需要的是上帝的恩典,.
去包容,.
接納我們,.
讓我們一起祈禱,.
讚美我們天上的父,.
我們多謝你,.
你的愛,.
你的心恩,.
在我們每一個人的生命裡面,.
就是這樣發生,.
上帝啊,.
我們要向你求寬恕,.
有些時候,.
我們都會站在擇石頭的位置,.
主啊,如果是這樣,.
我們求你寬恕我們,.
想起你對我們憐憫包容的心,.
好使我們能夠避免走在這個生命裡面的困局,.
我們很願意,.
很樂意,.
你給人帶來生命的希望,.
我們很願意幫助人,.
在生命裡面,.
能夠在罪裡面得到釋放,.
就求你先以恩慈待我們,.
好讓我們能夠更加懂得,.
怎樣去愛我們身邊的人,.
對人更多的包容,.
多謝你,.
聽我們的祈禱,.
禱告交托,.
是奉主耶穌基督得勝的明智祈求,.
阿們.

$^{801}$謝謝大家收看 再見.
\newpage



\section{羅馬書 9:30-33}
\label{sec:ucA_hp7yY_0}
\textbf{2022年11月6日崇拜直播｜陳明泉牧師｜以石為記:跌人的磐石｜羅馬書9\_30-33}
\newline
\newline
連結: \href{https://youtube.com/watch?v=ucA_hp7yY_0}{\texttt{https://youtube.com/watch?v=ucA\_hp7yY\_0}} ~~~~ 語音日期: 2022-11-08
\newline
\newline
\hyperref[sec:YIG0i1d0dV8]{\small{< < < PREV SERMON < < <}}
~
\hyperref[sec:index_chronic]{\small{[返順時目]}}
~
\hyperref[sec:y7hPaMojRiM]{\small{> > > NEXT SERMON > > >}}
\newline
\newline
羅馬書 9:30-33
\newline
\begin{longtable}{cl}
\hline
\hline
章節 & 經文 (和合本修訂版)\\
\hline
9:30 & \begin{tabularx}{0.7\textwidth}{X} 這樣，我們要怎麼說呢？那不追求義的外邦人卻獲得了義，就是因信而獲得的義。 \end{tabularx} \\ \\ \relax
9:31 & \begin{tabularx}{0.7\textwidth}{X} 但以色列人追求律法的義，反而達不到律法的義。 \end{tabularx} \\ \\ \relax
9:32 & \begin{tabularx}{0.7\textwidth}{X} 這是甚麼緣故呢？是因為他們不憑著信心，而是憑著行為，他們正跌在那絆腳石上。 \end{tabularx} \\ \\ \relax
9:33 & \begin{tabularx}{0.7\textwidth}{X} 就如經上所記：「我在錫安放一塊絆腳的石頭，使人跌倒的磐石；信靠他的人必不蒙羞。」 \end{tabularx} \\ \\
[1ex]
\hline
\hline
\end{longtable}
$^{1}$讓我們以聆聽聖道敬拜.
今日的經訓是記載在.
羅馬書9章30至33節.
在大家聖經的新約第216頁.
這樣我們可說什麼呢.
那本來不追求義的外邦人.
反得了義.
就是因信而得的義.
但以色列人追求立法的義.
反得不著立法的義.
這是什麼緣故呢.
是因為他們不憑信心求.
只憑著行為求.
他們正跌在那半腳石上.
就如經上所記.
我在石安放一塊半腳的石頭.
跌人的盆石.
信靠他的人必不至於羞愧.
願神的話語今日光照我們的內心.
(聖詩).
弟兄姊妹平安.
平安.
我們用石頭在聖經裡面的意象.
來將上帝的計劃.
將上帝在我們的心意.
我們一齊來領受.
過往大家去聽關於石頭的意象.
在聖經大部份的技術.
都是一些比較正面的技術.
在若但河裡面立石為記.
這條石柱代表著上帝世世代代.
看顧以色列人.
又或者用石頭來比喻上帝的堅固.
上帝是我們的岩石.
是我們的盆石.
是我們的拯救.
我們用不同的詩歌來講.
石頭就好像正面的說.
上帝在一群被選召的子民身上.
不過今日我們所說的石頭這個意象.

$^{41}$就和之前的說法有些不同.
今日某程度上比較負面.
那個負面就是.
剛才你一邊聽聖經的時候.
你會知道.
這塊石不單止是說.
很堅固.
很歌頌上帝計劃.
同時是一塊.
棘以色列人的一塊絆腳石.
一路走的時候你會棘跌.
會跌倒的.
甚至這塊石頭令到.
那個人不經意地.
進入了一個跌倒的狀態.
究竟今日這段經文.
或者這個絆腳石頭.
說的是一個什麼訊息呢.
希望今日我們一起去聽的時候.
我們一起領受.
特別是說我們這一群.
已經信了一段日子的基督徒.
不論你初信或者信了一段日子.
今日這段經文我們要拿來用.
要成為我們生命的一部分.
我們一起祈禱.
求聖靈幫助我們.
讓我們明白話語.
同心祈禱.
親愛主我們將我們每一個.
心思意念交代給你.
讓我們能夠明白.
你對我們所說的每一句話.
求你在我們當中.
你的聖靈與我們眾頂.
展回同在.
讓我們每一個.
恭謹在你面前聆聽的時候.
我們能夠明白.
你在這個世代的心意.

$^{81}$又幫助宣講的講得清楚.
忠心的將你的話語.
來到向頂展回解明.
讓我們不單止.
腦袋裡面仍知道.
更加在我們的內心裡面.
帶著你的話語.
奔走我們人生的道路.
恭敬將來的時間交託.
願你賜福.
我們獻上禱告.
奉靠基督耶穌得勝的聖名.
Amen.
剛才剛才說過了.
聖經裡面在這段經文.
我們用一個叫做「半腳的石頭」.
來講述上帝的工作.
其實這個「半腳的石頭」的出處.
本身是什麼?.
它在講述一個什麼樣的處境呢?.
在羅馬書第九章裡面.
其實保羅在回答一些很重要的問題.
回答以色列民.
一群被上帝選召的.
一群神國的子民.
這一群人是一群很謹守律法的一群人.
這一群我們看為一群道德的分子.
一群很decent.
很謹守的一群人.
竟然在上帝的永恆救恩當中.
竟然失落.
相反來說.
那些外邦人.
那些未認識主.
甚至乎之前是洩漫上帝的這一群人.
竟然得到上帝永恆的救恩.
在一個這樣的撥論當中.
保羅就要藉著上帝的話語.
來去指明究竟是一回什麼事.
保羅用什麼來應答這些人對保羅的質問呢?.

$^{121}$保羅用舊約裡面的以賽亞書.
來回答他們的提問.
在以賽亞書裡面的八章十四至十五節.
裡面就說到上帝的救贖計劃.
他必作為聖所.
卻向以色列兩家作半腳的石頭.
跌人的盤石.
向耶路撒冷的居民作為圈套和網羅.
許多人必在其上.
半腳跌倒而且跌碎.
並陷入網羅.
被纏住.
在這裡面.
以賽亞書八章十四至十五節.
你會看到大量負面的方式.
說的是上帝對以色列.
說的是兩個家.
即是以色列和猶大.
即是整個民族.
對他們整個神國.
被上帝選取的這些子民.
上帝要踢一下他們.
接下來的意象更加特別.
他說到不單止踢一下他.
踢他腳.
而且要整跌他.
要成為他們的圈套.
成為他們的網羅.
許多人因為這樣跌倒.
跌碎.
跌碎不是一跌就這麼簡單.
跌碎.
陷入網羅.
被纏住.
整件事說得很負面.
這塊石頭竟然帶來的.
整個以色列民族.
都陷入一種很負面的經歷當中.
而且被扯住被纏住.
在經文裡面說到.

$^{161}$上帝的計劃.
他要和以色列民.
他們將要進入一種這樣的狀態.
當然後面是說負面的.
前面他很輕輕說過一個字.
就是他必作為聖所.
同樣一塊石頭.
對於某些人來說是一個聖所.
但對於另外一些人來說.
是一個踢腳石.
是會跌碎陷入網羅.
究竟這塊石頭是什麼原因呢.
保羅在羅馬書裡面.
他就要用這句經文來去展述.
那個福音的工作.
所以用這樣的這段經文.
保羅就繼續來解釋.
剛才說過九章三十節到三十一節.
這樣我們可說什麼呢.
原來那本來不追求義的外邦人.
反成了義.
就是因信而得的義.
但以色列人追求律法的義.
反得不著律法的義.
這是什麼緣故呢.
是因為他們不憑信心求.
只憑行為求.
他們正跌在那盤石上.
在第三十節到三十二節.
其實保羅很直接了當地說.
那個踢腳石.
或者令到那些以色列人.
為什麼會被踢倒呢.
原因就是因為他們以為藉著律法.
藉著遵守上帝的教訓.
他們就可以乾淨.
他們就可以在上帝面前被稱為義者.
用一個完全的身份.
用一個強人的姿態.
來進入上帝所命定.

$^{201}$或者上帝所預備的.
應許的永生當中.
值得留意的是.
在這裡他在說.
本來不義的人.
都不知道做了什麼.
就被稱為義.
本來一直一心追求義的那些以色列人.
反而在這裡被追極都追不到.
他們的問題在於什麼呢.
三十二節回答了給我們知道.
原來這一班以色列民.
他們追求的.
或者他們所做的.
不是憑著信心來行義.
他們只是一心覺得.
跟隨來做.
跟著這條規則.
跟著律法的軌跡.
一路做.
跟著走就不會死.
所以他們就在這樣的軌跡當中.
他們就一路一路.
得到永生的道路.
不過在這裡.
保羅特別說到.
他們這一班人.
是得不到律法的義的.
畢生追求律法.
但是卻是.
終生都沒辦法得到.
律法上的正義.
就是在永生的道路裡面.
沒有他們份.
究竟是什麼問題呢.
其實不是因為上帝.
跟這一班以色列人有仇.
不是會特意做一些事情.
來作弄他們.
上帝的心意是要讓.

$^{241}$普世的子民.
上帝所選取的這一班人.
都得著永生.
不過在聖經裡面說到.
這班以色列人.
他們信的信什麼呢.
他們不是信上帝.
他們只是信有這個規條.
他們跟著來做.
他們就可以得到永恆的生命.
律法本來是要導人.
走在上帝的心意當中.
不過律法也有它終極的限制.
律法沒辦法.
建立一個義的人.
在加拉太書裡面.
保羅再繼續補充.
他這樣說.
這律法是我們糞蒙的師父.
引我們去基督那裡.
使我們恩順清義.
律法的作用.
或者上帝教導的規則.
這些人神之間.
人如何去回應上帝.
這些作法,這些行為.
其實是要訓練我們.
但不是一個永生的道路.
到最後是要引領我們去基督那裡.
不是直接做這些事就得到永生.
這樣說好像很抽象.
不過用一個生活例子.
可能大家會明白一點.
早前有機會去跟老師聊天.
我是借街醫的.
有時候吹吹水問問老師.
我說現在的世代.
最差的是什麼學生最難教的呢.
輕描淡寫想著說一個這樣的問題.
但老師很認真地回答.

$^{281}$他說最差的學生.
通常都不是那些常常破壞規矩的學生.
最差的反而是上課.
最循規蹈矩.
最懂得看老師眉頭眼額的學生.
在你面前最恭敬.
但背地裡他所做的事通常都很嚴重.
通常最容易闖禍的就是這些學生.
這個老師說的不是一個負面標籤.
說這些學生.
不過他講到一些心態.
在學校的規矩裡面.
可以掩飾一個人真正的心理.
在人的焦點之下.
或者在一些規則底下.
人可以扮得很端莊.
但心底裡或做的事.
可以是很差的一個人.
我想起我小時候.
這個老師一講的時候.
我想起我小時候.
我小時候高高大大.
又懂得看眉頭眼額.
見到老師通常最恭敬的是我.
第一個鞠躬的是我.
對於校規.
都不知道為什麼這麼熟.
我看過一些字.
佛抄過一次我記得.
所以現在你看到我寫講章.
我不看的.
我抄過一次我記得.
天賦給我的能力.
以前可能佛抄過一次.
我就記得全部校規.
又這麼熟校規.
又高高大大.
又懂得看眉頭眼額.
這些通常在學校做什麼呢.
做風紀.

$^{321}$做風紀隊長.
那時候拿著風紀隊長帶.
作威作福.
從小到大.
特別是高年級的時候.
通常都是做這些事情.
有一次拿著濫用職權.
喜歡作威作福.
經常作弄女同學.
從小到大我都是這麼衰格.
我都知道.
見到女生經常作弄.
有一次見到一個女生.
有一點點行為上.
打了小色鐘之後就吃東西.
然後就想抓她.
摘她名字.
然後我就想作弄她.
然後就好像追逐一樣.
追逐她.
追她來摘名字.
就好像兵捉賊一樣.
有意無意之間.
那女生在跑.
然後因為追逐的時候.
她就一頭撞到學校那些.
百頁窗的玻璃那裡.
就撞穿了頭.
起了一個大瘤.
我是風紀來的.
但是抓一個女同學吃東西.
弄到頭部被撞穿了.
老師就大興問罪之師.
他就問陳明全.
為什麼會發生這樣的事情.
那個女同學無端端.
為什麼會這樣走.
撞穿了頭.
你見到這件事的發生.
究竟是什麼.

$^{361}$然後我就在掩飾.
就說因為校規不准吃東西.
她又頑皮.
我想摘她名字.
她又拒捕.
那時候我在說.
然後她又追逐之間.
她自己不小心有什麼.
就撞到頭撞穿了.
然後老師說.
一般女同學是不會這樣追逐的.
就是你一定有些事情要做.
然後還要掩飾說.
沒有啊.
人就是這樣的.
又達到校規.
到最後老師就一聲.
你會不會濫用職權.
你是在說用你這段時間.
去糟蹋那些同學.
然後臉都青了.
說這些.
然後否認.
不過到最後是那個女同學.
只好直接說.
是我在欺凌她.
所以她要走.
整件事真相都知道了.
真相知道了之後.
老師就馬上撕掉.
撕掉了馬上沒得做.
這件事小時候.
我還覺得自己很正義.
她只是拒捕而已.
我擇名不聽.
覺得什麼.
不過當我再聽老師說這件事的時候.
我想起.
其實我心裡基本上不是想忠實.
去執行我的工作.

$^{401}$我只不過是運用我的權力.
來欺壓.
作弄一些人.
表面上可以是很正氣凜然的一種表達.
但生命裡面可以藏於很多污穢.
所以用這個角度來看.
在耶穌基督.
在福音書裡面.
祂對於那些法利賽人.
耶穌基本上.
很多時候所做的言行.
好像衝淨他們一樣.
耶穌特別稱呼這些人叫做假冒為善.
這群人其實在眾人眼中.
不是假冒為善的人.
是一群很高尚.
很道德高地的一群人.
但是卻是在真理的亮光之下.
其實祂是用這些律法.
來遮蓋自己.
包裝自己.
讓自己成為一個高人一等的人.
祂所做的好事.
目的是什麼呢.
不是因為祂想接近上帝.
而是祂很想顯出別人的不義.
祂所做的好事.
為什麼是很純粹地.
因為覺得這件事是好事呢.
不是.
祂只是覺得.
跟著做.
是律法.
很功利地跟隨這條毒路.
不過在基督耶穌.
去輾明的救恩.
祂在說這些法利賽人也好.
這些以色列人也好.
祂走在這條.
用律法來掩飾自己生命的道路.

$^{441}$到最後不能夠引領他.
走向一個對的人生道路當中.
反而他會面對著.
他人生走向一個很大的困境.
在經文裡面.
他在說一些不追求義.
或者外邦人.
因為信.
單單因為信而得到的義.
因為基督耶穌.
祂說到在律法的限制下.
人沒辦法做每一樣事情.
滿足了律法的要求.
以致在人力當中.
可以得到這個永生的道路.
人的盼望不是在於.
人有力量實踐.
律法的義.
人的盼望是在於.
因為有一個上帝可以讓我們信靠.
我們單單信靠祂.
以致這些本來不濟的外邦人.
本來這些生命裡面很糜爛的.
或者以前昔日褻瀆上帝的外邦人.
因為相信.
以致祂走另外一條路.
一條確保他可以得到永生的路.
這就是信靠上帝的意思.
九章三十三節.
「就如經上所記,我在石安.
安放一塊半腳的石頭,跌人的盤石.
信靠他的人必不至於羞愧」.
在羅馬書九章三十三節裡面.
其實這一節的經文.
其實是引用兩段以賽亞書的.
頭半部.
即是半腳的石頭.
跌人的盤石.
剛才說過了.
就是以賽亞書第八章.

$^{481}$第十四至第十六節.
特別是第十六節.
就是說他們是倚靠上帝.
另一個部分就是在以賽亞書第二十八章十六節.
經文是這麼寫的.
「所以主耶和華如此說.
看,我在石安.
安放一塊石頭作為根基.
是試驗過的石頭.
是穩固的根基.
寶貴的防割石.
信靠的人必不著急」.
以賽亞書這裡說的是什麼呢?.
理解了這裡再說羅馬書.
在以賽亞書第二十八章.
它說的是一個很重要的時刻.
就是上帝的審判.
上帝的審判就像洪水一樣.
氾濫著整個大地.
人人沒有辦法自救.
面對洪水的威猛衝擊之下.
每一個人都是危險的.
不過在一個洪水氾濫的日子.
在以賽亞書裡面說的審判之日.
就是上帝的子民.
他還有唯一的救法.
就是抓住一些在洪水裡面衝不掉的石頭.
經文是說一個這樣的意象.
在洪水.
在好像審判的日子裡面.
人人都抓住一些東西.
抓住那些東西被衝擊的.
那些東西沒有用的.
是要抓住一塊穩固的石頭.
是試驗過根基穩固.
被洪水不會衝掉的一塊石頭.
如果當我們抓住這一塊石頭.
在以賽亞書裡面說的.
上帝耶穌基督就是那個防角石.
抓住這一塊防角石.

$^{521}$信靠的人就不用著急.
就不用怕了.
所以經文裡面就說到.
在一個末日或者審判的日子裡面.
每一個人都要找一些東西抓住.
最重要抓住什麼呢?.
不是抓住身邊的.
你覺得重要的金銀財寶.
而是抓住一個經得起考驗.
而且是上帝保證的那塊防角石.
那個就是耶穌基督.
所以在經文裡面.
保羅引用以賽亞書.
說的是既是跌人.
同時又是穩固的防角石.
他將兩樣東西放在一起的時候.
整句的解釋就是說.
對於那些以色列人來說.
他們以為抓住一塊好的石頭.
抓住做律法.
他就得救了.
那個不是的.
那個是跌死你的.
那塊是跌人的盤石.
一定要抓住什麼呢?.
抓住基督耶穌.
單單順靠他.
不要太過相信自己的行為.
不要相信人力可以按著律法.
按部就班.
做好人你就得救.
在聖經裡面說是不行的.
律法已經將我們放在一個失敗的位置.
我們不要再跟著這件事.
我們反而要將我們的信心.
放在基督耶穌的身上.
順靠他.
以致我們得著這份永恆的救恩.
所以在經文裡面說.
我們不需要怕.

$^{561}$不需要著急.
他引用這段經文.
不過在三十三節裡面.
他再用一點字眼.
有一點改變的.
因為保羅身處的狀況.
不是即時的審判.
他將那個不至於.
必不著急.
改為必不至於羞愧.
不會覺得蒙羞.
你讀上去還是覺得奇怪.
有些的聖經學家就在說.
不至於失望.
你圍著上帝.
你人生是不會失望的.
在人生無論什麼處境當中.
上帝的應許是永不落空的.
上帝答應過救你.
他一定會幫助你.
所以在整句九章三十三節裡面.
他強調的就是.
那些人走錯的路.
走錯的路.
那塊石頭就卡腳了.
走對你的道路.
那塊就是你生命救不你的防角石.
你圍著它的時候.
你的生命就走在上帝命定的.
永生的道路當中.
我想再看多一點.
就是那個信靠的字.
信靠的人必不著急.
信靠是什麼意思呢?.
我們今天很多丙姐妹以為信靠.
就是純粹在小學階段.
老師呼召舉手.
那個就叫做信靠的生活.
我們純粹將信靠.
以為是.

$^{601}$信以為真.
又或者是很簡單的.
知道有上帝存在.
我們叫這個叫做信.
不過在聖經裡面所說的信靠.
從來不是說一個腦袋裡面的認知.
那個信靠是說.
將你整個人的生命.
完全依靠它.
信任它.
讓它帶領你走人生的道路.
說一個比較容易理解的例子.
你今天坐巴士.
你坐車來到教會.
你坐的車為什麼會上這輛車呢?.
可能你正在上某個號碼的巴士.
它說它的目的地在這裡.
你知道它會將你.
由你身處的位置帶你去到目的地.
你相信它.
所以你就上這班車.
你上去的時候.
你投放你自己.
你給錢.
但是如果你坐人生的巴士.
你投放生命裡面的力量.
然後安然地等待這班車.
送到目的地當中.
信靠的意思就類似這樣.
上帝告訴你.
你人生的目的地在這裡.
你上這班車.
你讓它帶領你去到那裡.
在你去的過程裡面.
你不要騷擾司機.
你不會說你走錯路.
司機決定了你.
這樣去到目的地.
所以你就安然地.
繼續接受.

$^{641}$等待司機載你去到目的地.
同樣在我們人生道路裡面.
耶穌基督會帶領我們去走.
我們人生的道路.
我們將我們行走的主權交給祂.
安心地信任祂.
而且有百分百的信心.
我一定去到那裡.
是快與慢的不是我們控制的.
但是我們知道.
我們一定會去到那裡.
信靠就是一個這樣的意思.
所以,弟兄姊妹.
不要純粹停留在小學階段.
中學階段我舉手缺個字.
又或者我知道上帝是真的.
純粹知道祂是真的不足夠.
我們要讓我們的生命.
被上帝帶領我們走人生的道路.
按著祂給我們的方向.
我們行事為人.
由祂來做帶領.
這個就是信靠的意思.
當我們去看這段經文.
其實對於以信主的弟兄姊妹.
我相信理解不是很複雜.
我們大概都知道福音的意思.
不過很多時候弟兄姊妹讀這段經文.
就進入一個誤區.
就是我們覺得.
我們把自己視為外邦人.
所以我們就恩信稱義.
相反來說.
我們就批判猶太人或者以色列人.
他們不信.
所以他們就不得救.
我們就站在我們的位置.
我們就沾沾自喜.
我覺得一定會得到上帝的救恩.
不過如果我們純粹是在說.

$^{681}$讀這段經文.
我們覺得自己沾沾自喜.
大概我們失去了這段經文.
對我們的營養價值.
真相是.
我們很多時候.
我們很多個都好像那些以色列人一樣.
在我們的信仰裡面.
很多時候很發力賽.
我們當我們沾沾自喜.
以為自己一定得救的時候.
可能我們有意無意之間.
其實我們是在走那些被擊跌的那些以色列人.
或者那些發力賽人的道路.
我們用聖經包裝自己成為一個好人.
我們用上帝的話語.
裝扮自己成為一個上帝喜悅的人.
我們有意無意之間.
我們都是在走舊約的那些律法主義.
以致我們以為這樣就可以得救.
近來有機會看一本書.
《我為什麼要信》.
一個作者叫Tim Keller.
他的書很好看.
他說裡面說一個很重要的.
一個今天世代裡面面對的處境.
他特別對基督徒說.
他說今天罪惡有兩種呈現的方式.
第一種呈現方式就是很壞的.
破壞一切的規則.
糜爛的生活.
很差的人生.
罪惡不難理解.
他說罪惡很多時候用這種方式來表達.
我們認同.
不過他說第二種罪惡呈現的方式.
就是讓你覺得自然.
覺得自己很理想.
覺得自己是一個上帝很喜悅的一個人.
我們用聖經,用律法來包裝自己.

$^{721}$讓上帝的話語成為了.
遮蓋自己生命真相的那種生命表現.
Tim Keller繼續說.
如果是魔鬼.
他容易用的方法.
他會用那些自然的人的方法.
因為那些人不由自主地.
其實是呈現罪惡生命的那一部份.
他們外表可以很好.
但心底裡可能離上帝很遠.
他再繼續說.
他說他引用一個叫做Franny.
就是歐康娜的一個作家.
他說避開上帝最好的方法是什麼呢?.
就是避開罪.
如果你正在避開罪.
過著一個道德的生活.
以為這樣上帝就會祝福你.
我們基督教信仰只是這樣的時候.
可能你看耶穌只是一個道德的老師.
是一個幫助者.
是一個智者等等.
不過你沒有打算讓祂成為你的救主.
聖經或者恩順清二的道理是在說什麼呢?.
是在說我們不能夠透過人世間的努力.
讓我們成為異議的一部份.
零啊!沒有辦法做.
律法目的是要告訴我們上帝的心意.
但按著人的力量.
你打算走這條路.
來被上帝稱讚成為一個潔淨的人.
去到上帝的底層.
No way!不是的.
這條不是路.
律法是要把我們全人類都關在罪當中.
我們只是體會到自己生命的本相.
因為我們不能夠用人力得救.
所以我們單單信靠上帝.
恩順清二提醒我們.
我們不是想像中那麼容易.

$^{761}$我們只不過是一個罪人.
我直白一點去說.
有時候我們做一些好的事.
可能我們想為自己做一些包裝.
讓自己可以在Facebook.
在IG上有些相片登下.
讓別人讚許一下自己.
讓自己感覺超越了別人.
我比別人更加優勝.
又或者我所做的好行為.
我所做的義.
我所做的每個人都覺得好的事.
其實我是想反襯那些人.
那些人多麼不義.
那些人多麼差.
我和他們不同.
我站高一個梯級.
他們就成為這樣的人.
求主幫助我們.
如果我們的生命裡.
我們的義.
我們去實踐聖經的教導.
如果我們帶著這種心腸.
我們其實和昔日的法理賽人.
沒有什麼分別.
恩順清二卻是叫我們打回原形.
回想一下自己.
我們生命本相是怎樣.
我們本來是一個什麼人呢?.
我們的生命按著我們自己.
根本上不是想像中那麼理想.
不是那麼令人羨慕.
甚至很多生命的污點.
早前有機會.
我教給同工知道.
看著整件事發生.
話說三十年前.
有一個暑期的課程.
出去上課.
認識了一些學校的同學.

$^{801}$其他同學.
那時候我中六.
上一個課程是補英文的.
因為跨學校上課.
認識了很多其他學校.
都是那樣東西.
都是微信主.
很多的職集.
那些雙雙格格.
從小到大都是這樣.
找一些女孩子來作弄.
經常都是這樣.
為什麼要說這件事呢?.
早兩三個星期前.
有一個女同工.
一個機構的女同工.
敲我房門.
就和我談一談.
說未來有一個聚會.
陳牧師想和你談一談.
那天會講福音訊息.
我不知道她會進來.
但我認到她的樣子.
這個樣子就是我三十年前.
上過英文課程.
經常遭殃的女孩子.
我認到她.
因為都幾容易認.
但她認不到我.
在那裡說的時候.
我當然英明神武.
要這樣做.
整件事就是這樣.
說完之後.
我忍不住.
我說姐妹.
我很蒙昧.
我想問你記不記得.
大概多少年的時候.
在你上夏季課程.

$^{841}$上過一個英文課程.
你記不記得有這樣的事.
姐妹說是.
都記得.
為什麼你會問起呢?.
我說你記不記得.
有些男同學.
特別有一個.
經常喜歡糟蹋女孩子.
她說.
男孩子個個都是這樣.
我說特別有一個.
很口臭的那個.
她的眼睛好像叮一叮.
我說很口臭.
她說是呀.
真的有一個人很口臭.
我說那個很口臭的那個.
就坐在你面前.
那個是我.
她就大聲在房間大笑.
笑到整個辦公室都聽到.
在我房門口.
她說這個世界真是有神蹟.
一個口臭男孩.
竟然信主.
我還說是呀.
她還說你不單做牧者.
還要做主任牧師.
你不要害會友.
我就在我房門口.
我說畢竟給點面子.
剛才你也覺得我英明神武.
但你一刻一說我真.
舊時的樣子.
你就這樣糟蹋我.
反過來.
打回原形.
在她面前.
真是舊時做的事.

$^{881}$中什麼收什麼的.
我知道.
當那天她離開了.
我坐車回家的時候.
一直坐在巴士上經過旺角.
我回想這幾十年發生什麼事.
我只是在旺角裡.
一個死小子.
一個口臭男孩.
幾十年後.
竟然在一個旺角的教會裡.
成為一個主任牧師.
我有多厲害呢?.
我自己知道.
我有多勤奮呢?.
我更加自己知道.
一邊數.
我跟上帝說.
其實今天我能夠.
繼續幫助弟兄姊妹.
幫助很多人成長.
幫助人們讀聖經.
其實也不是我有什麼才智.
或者我有什麼優越的背景.
我某程度上也是爛命一條.
是因為上帝的恩典.
將我從這個口臭男孩.
成為今天用話.
用上帝的話語.
牧羊的一個牧者.
我多謝上帝.
我們生命需要打回這些原形.
找回自己.
當我明白了這個道理的時候.
我心裡好像忽然間釋懷.
我對於有時候有些人與事.
我覺得耿耿於懷.
我懂得放開.
每一個人都有自己的過去.
每一個人都有自己的個性.

$^{921}$有些人都是很不理想.
上帝沒有嫌棄我.
同樣上帝也不是想我.
變得更加敏感地對身邊的人.
反而是更加開放地對待身邊.
有些不是很理想的那些人.
上帝的恩典.
都淋在那些人身上.
恩順清義幫助我們打回原形.
看回我們自己.
根本就沒有自義的本錢.
我們只是用罪人的心態.
求上帝的恩典在我們生命裡訴訟我們.
第二個最重要的重點.
恩順清義提醒我們.
世界上有很多東西.
鼓吹著你去相信.
不過當你做人越來越久的時候.
你發覺這個世界有很多東西.
是你不值得相信的.
今天我們批判很多人是懷疑主義.
不過事實上是真的.
你有什麼可以完全相信呢?.
恩順清義這個很重要的教義.
提醒我們.
在它改教時間裡.
它告訴當時的人.
你不要相信建築物的教會.
你不要相信教廷的制度.
你的信心要放在哪裡?.
唯獨聖經.
唯獨基督.
唯獨恩典.
唯獨信心.
你的信心是放在那些值得你信靠的身上.
保羅在說.
這麼多東西相信.
但是你要扶持著的是.
試驗過的防角石.
那個可以讓你抵禦洪流的那塊石.

$^{961}$基督耶穌.
扶持著基督耶穌.
由這個作為你生命的起點.
也會成為你生命的終點.
今天很多東西鼓吹你去相信.
叫你去相信.
你的WhatsApp.
你的媒體等等.
你想避都不能避的.
不過.
那件事不是聖經所說的.
那是引你去到聖經.
還是拉你離開聖經.
離開基督的那些信息呢?.
求主幫助我們.
今天我們需要.
回到最基本的信心.
找回我們信靠的對象.
由這個成為我們生命.
實踐信仰的起始點.
當我們在恩順清義的光照下.
看到自己的本相.
也看到自己生命的起點的時候.
我們更加知道.
我們追求一生的對象.
就是那個基督耶穌.
大概我們人生.
我們都不會怎樣行差踏錯.
求主繼續幫助我們.
今天我們需要.
這一塊防割石.
拯救我們.
在很多自然的人當中.
耶穌基督的生命是伴腳石.
但在我們心目中.
祂是我們生命的支柱.
依靠讓我們不會走失.
走錯路的.
不會被洪水沖去的防割石.
今天怎樣去選擇.

$^{1001}$希望今天這篇道幫助我們.
讓我們懂得怎樣行人生的路.
謝謝大家收看 再見.
\newpage



\section{馬太福音 16:13-19}
\label{sec:y7hPaMojRiM}
\textbf{2022年11月13日崇拜直播｜陳冠成牧師｜以石為記:建造教會的磐石｜馬太福音16\_13-19}
\newline
\newline
連結: \href{https://youtube.com/watch?v=y7hPaMojRiM}{\texttt{https://youtube.com/watch?v=y7hPaMojRiM}} ~~~~ 語音日期: 2022-11-15
\newline
\newline
\hyperref[sec:ucA_hp7yY_0]{\small{< < < PREV SERMON < < <}}
~
\hyperref[sec:index_chronic]{\small{[返順時目]}}
~
\hyperref[sec:esPM0k_dZBY]{\small{> > > NEXT SERMON > > >}}
\newline
\newline
馬太福音 16:13-19
\newline
\begin{longtable}{cl}
\hline
\hline
章節 & 經文 (和合本修訂版)\\
\hline
16:13 & \begin{tabularx}{0.7\textwidth}{X} 耶穌到了凱撒利亞．腓立比的境內，就問門徒：「人們說人子是誰？」 \end{tabularx} \\ \\ \relax
16:14 & \begin{tabularx}{0.7\textwidth}{X} 他們說：「有人說是施洗的約翰；有人說是以利亞；又有人說是耶利米或是先知中的一位。」 \end{tabularx} \\ \\ \relax
16:15 & \begin{tabularx}{0.7\textwidth}{X} 耶穌問他們：「你們說我是誰？」 \end{tabularx} \\ \\ \relax
16:16 & \begin{tabularx}{0.7\textwidth}{X} 西門‧彼得回答說：「你是基督，是永生神的兒子。」 \end{tabularx} \\ \\ \relax
16:17 & \begin{tabularx}{0.7\textwidth}{X} 耶穌回答他說：「約拿的兒子西門，你是有福的！因為這不是屬血肉的啟示你的，而是我在天上的父啟示的。 \end{tabularx} \\ \\ \relax
16:18 & \begin{tabularx}{0.7\textwidth}{X} 我還告訴你，你是彼得，我要把我的教會建造在這磐石上，陰間的權柄不能勝過它。 \end{tabularx} \\ \\ \relax
16:19 & \begin{tabularx}{0.7\textwidth}{X} 我要把天國的鑰匙給你，凡你在地上所捆綁的，在天上也要捆綁；凡你在地上所釋放的，在天上也要釋放。」 \end{tabularx} \\ \\
[1ex]
\hline
\hline
\end{longtable}
$^{1}$經法會議事項.
《馬太福音》第十六章十三至十九節.
《新約聖經》第二十四頁.
《馬太福音》第十六章十三節.
耶穌到了埃塞尼亞肥納比的境內.
就問門徒說:人說我人子是誰?.
他們說:有人說是施浸的約翰.
有人說是以利亞.
又有人說是耶利米或是先知裡的一位.
耶穌說:你們說我是誰?.
西門彼得回答說:你是基督.
是永生神的兒子.
耶穌對他說:西門巴約拿,你是有福的.
因為這不是屬血肉的子是你的.
乃是我在天上的父子是的.
我還告訴你,你是彼得.
我要把我的教會建造在這盤石上.
陰間的權柄不能勝過他.
我要把天國的約事給你.
凡你在地上所捆綁的.
在天上也要捆綁.
凡你在地上所釋放的.
在天上也要釋放.
弟妹平安.
相信我們都熟悉剛才這段經文.
特別我們記得彼得對耶穌的認訊.
認耶穌是基督,是永生神的兒子.
這件事件.
事實上這件記錄.
可以說是整卷馬太福音的一個分水嶺.
如果你有留意.
在這個宣告出現之前.
其實耶穌從未透露過.
祂要上耶路撒冷來受苦犧牲.
並且耶穌從未告訴門徒.
祂要建立教會.
現在一群人來到.
凱撒利亞菲勒比這個地方.
耶穌知道自己離世的日子也將近.
所以祂也準備把福音的棒.

$^{41}$交給門徒.
託付他們來建立教會.
於是你見到耶穌就開始出題目.
考問門徒對耶穌的認識到底是怎樣.
為何會選凱撒利亞菲勒比這個地方呢?.
首先我們知道.
這是一個外邦人居住的地方.
位於當時加利利海北邊的四十公里.
是當時的昏風王希律肥利.
為了討好羅馬帝王而興建的一個城市.
我們不妨看看那張圖片.
看看它的外貌是怎樣.
我們說其實這個城市.
除了敬拜當時的羅馬帝王凱撒.
有他的廟宇之外.
其實還有各式各樣的神鵪.
就像一個鵪鶉一樣.
供奉著不同當時的人膜拜的神明.
當地的人主要去膜拜一個叫做Pang.
我們說上半身是人.
下半身是馬的一個希臘神奇.
傳說Pang有一個很獨特的情況.
就是他很喜歡嚇人.
突然間躲在陰暗的地方.
跳出來就會發出一些很恐怖的叫聲.
令人毛骨悚然.
所以英文有個字叫Panic.
驚恐.
其實這個字P-A-N-I-C.
也是來自他們對Pang的一種恐懼.
至於在這些神鳥之上.
大家看到可能有一大片的懸崖.
當地人叫它做神奇之石.
在裡面有一個巨大的山洞.
在圖中的凱撒奧古士督神殿上.
有一個很大的黑影.
其實那是一個山洞.
下一張就會很清楚看到.
這個山洞.
人們傳說是通往陰間的入口.

$^{81}$所以給了它一個名字.
叫做陰間之門.
你想像一下整幅圖畫是怎樣的.
耶穌和祂的門徒.
身處在這片陡峭的石塊之下.
面對著這麼多宏偉的建築物.
以及他們背後所代表的不同勢力.
耶穌和門徒是否顯得格外渺小.
其實很渺小的.
耶穌和祂的門徒.
在這麼多宏大的事件裡.
不過你看到耶穌是專門選擇.
在凱撒里亞肥勒比這個地方.
我們說這個充滿了.
異教廟宇林立的地方.
這個遠離了耶和華信仰的氛圍的地方.
並且在這一切.
代表著君王,賈臣.
甚至陰間的勢力面前.
耶穌要說一個宣告.
就是說.
祂即將要透過被人看為很微小的教會群體.
來為當時的世界.
帶來翻天覆地的改變.
並且要將那些在不同權勢管轄下的人.
釋放出來.
這可以說是整段經文.
為什麼耶穌要選擇在凱撒里亞肥勒比.
宣告一個原因.
你看到耶穌首先問第一個問題.
祂問祂的門徒.
那些人怎麼看我.
他們說我是誰.
門徒於是回答.
有人說.
耶穌你不是被希律王斬首.
不過已經復活了的施浸約翰.
詳情可以翻前兩章.
《馬太福音》十四章就有交代.
為什麼當時的人會覺得.

$^{121}$耶穌是復活了的施浸約翰.
也有人說.
耶穌是以利亞.
就是舊約最後一卷.
《馬拉基書》裡.
預告上帝將要差派來.
復興以色列的一個先鋒.
當然也有人說.
耶穌你像耶利米.
大概是因為.
兩個人的遭遇其實都很相似.
他們的信息都是受到.
當時一些宗教領袖來去嫌棄.
並且他們經常也都被人針對.
甚至乎攻擊.
但是無論如何.
你會看到.
耶穌的宣講.
和祂所行的各樣神蹟.
其實也都令大部分的群眾.
認定祂是一位先知.
是一位上帝差派來的代言人.
就連當時猶太的宗教領袖.
其實他們都不敢排除這種可能.
事實上耶穌.
祂又怎會不知道.
身邊的人如何去議論祂呢?.
所以其實祂問頭一條問題.
可以說是一碟前菜.
實際的作用是用來測試.
到底那些跟隨祂的門徒.
會不會都是人魂逆魂.
純粹當祂是先知呢?.
還是他們已經緊密地.
近身地跟隨了耶穌一段日子.
他們對耶穌其實.
會不會有另一種的體會呢?.
這是耶穌更想他們去回答.
所以你看到.
他單獨直入.

$^{161}$他問第二個問題.
那麼你們說.
我是哪一個?.
每一個門徒.
他們都要親自回答.
其實在過去不同的年代.
不同地方的人.
不同文化的人.
你有沒有發覺.
其實都在問這條問題.
就是耶穌到底是哪一個?.
其實今天還在問.
如果你有機會.
上網搜尋一下.
你會發現.
其實每隔幾年.
就有些民意調查出來.
問一下耶穌到底是哪一個.
收集一下他們對耶穌的看法.
有人覺得耶穌是一個.
很王愛可親的人.
有些人覺得耶穌是一個.
道德很完美的人.
又或者有些人覺得.
耶穌好像一個社工.
社會工作者.
專門去幫助貧苦大眾.
有些人說他超級巨星.
去到哪裡.
總會有些粉絲跟隨他.
當然也有人對耶穌.
有另一種的看法.
覺得他說的話都不落地.
是一個理想主義者.
又或者覺得他是一個.
革命的先驅.
總是打破常規.
甚至有人覺得.
他是一個悲劇的英雄.
雖然開頭很厲害.

$^{201}$但最後失敗收場.
無論如何.
對於基督徒來說.
耶穌是誰這個問題.
其實我們每一個都需要去回答.
因為他一方面.
很實際地考驗.
你和我對耶穌基督的身份.
有沒有正確的理解.
同時也在測試我們.
假如有人真的去問你.
這條問題的時候.
你是否已經預備好.
如何去回答呢?.
當有人問我們.
耶穌是誰的時候.
我們能否預備好.
將耶穌介紹給對方知道.
頂子妹這條問題.
其實沒有人可以代我們自己回答.
你不可以找.
回得久一點的頂子妹代你回答.
你不可以找教會的牧者代你回答.
你必須要.
好像這裡的門徒.
要親自回答.
事實上耶穌期望我們.
當問到耶穌是誰的時候.
我們都需要很誠實地去檢視回應.
其實我們和耶穌的關係是怎樣.
他在我們心目中.
有多大的影響力.
這才是這條問題背後.
一個重要的意義.
第十六節到此.
彼得的回應.
耶穌你就是基督.
你就是永生神的兒子.
我們對這兩個的清涵其實不陌生.
基督我們知道是來自希臘文.

$^{241}$等於希伯來文我們常說的彌賽亞.
也就是受高者的意思.
特別在舊約裡.
我們知道有幾個身份是需要被高納的.
需要受高的儀式.
例如祭司,先知或當時的君王.
他們都需要通過高納儀式來承接職份.
來顯示他們的職份是受託於上帝.
並且背後真正服侍的是上帝他自己.
當以色列人經歷了亡國被擄之後.
其實在舊約和新約這四百多年裡.
你知道他們不停被戰火摧殘.
被不同的外邦政權統治.
所以令他們對基督.
對彌賽亞的觀念慢慢演變.
他們很渴望特別在耶穌的時代.
他們很希望上帝興起一位拯救者.
來幫助他們脫離羅馬政府的手.
真的幫他們復興以色列.
所以彼得在這裡認耶穌是基督.
是永生神的兒子.
其實是在說他承認耶穌的職份來自上帝.
是上帝揀選了耶穌.
來將拯救,復興以色列的工作.
托付給耶穌.
至於永生神兒子這個稱呼.
其實按照當時猶太人的傳統.
他本身不一定是一個神聖的身份.
事實上當時以色列人或大衛的子孫.
其實通通都可以泛稱為上帝的兒子.
不過基於耶穌和門徒.
那種朝夕相對的緊密相處.
他們見識過耶穌所行的不同神蹟奇事.
並且你有留意耶穌經常在門徒面前.
透露他和天父的關係.
那種父子緊密的關係.
所以雖然未必完全理解.
但彼得已經意識到.
眼前這一位的耶穌.
他不單單只是上帝差派來的救主.

$^{281}$他更加擁有和上帝一樣的神聖身份.
所以第十七節.
耶穌就回應他的宣認.
耶穌對彼得說.
西門巴若拿.
其實是什麼意思呢?.
「巴」就是兒子的意思.
「巴若拿」就是若拿的兒子.
簡單來說西門巴若拿.
就是若拿的兒子西門.
你留意其實科林書.
只有這裡很鄭重地說彼得的全名.
很明顯你見到耶穌.
即將要對西門彼得說一些很重要的事情.
所以才會說全名.
意思是什麼呢?.
譬如我平時見到.
量瑞倫弟兄.
我會叫他「細倫」.
一個稱呼.
我不會特別叫他「量瑞倫弟兄」.
如果我這樣叫他.
代表什麼呢?.
可能他欠我一頓飯還沒還.
又或者是一些很重要的場合.
要說一些很重要的事情.
所以才會說全名.
是一些很重要的宣佈.
耶穌首先稱讚西門彼得.
你是有福的.
其實等於承認了.
彼得他剛才對耶穌的宣認是怎樣呢?.
正確無誤.
不過耶穌說.
你之所以能夠宣認我的真正身份.
不是靠你自己本事.
不是靠人的本領學識.
來正確推論到耶穌到底是誰.
而是天父來賜福指引給他.
透過什麼場合呢?.

$^{321}$很大機會就是彼得一直跟隨著耶穌.
在裡面的互動場景.
譬如說當日耶穌怎樣呼召西門彼得.
而彼得又真的肯回應.
放下所有跟隨耶穌.
又或者說他親自見到自己的岳母.
怎樣給耶穌治好他的熱病.
又或者最記得有一幕很經典的.
就是耶穌怎樣邀請彼得在水面上行走.
這是他唯獨一手的親身經歷.
並且當彼得掉下水的時候.
耶穌怎樣即時把他拔起.
從水上面.
就是透過這種種的互動經歷.
其實上帝就將耶穌基督的真正身份.
一路一路啟示揭示出來.
讓彼得他對耶穌是基督是救主.
是永生神兒子.
這個信念就變得越來越肯定和清晰.
由最初可能他最初跟隨耶穌的時候.
他期望想像這位應該是救主.
應該是永生神的兒子.
所以我跟隨他.
到他累積了大大小小的信心.
累積了大大小小的經歷的時候.
他不只是期望他是.
他是認定眼前這位的耶穌.
就是救主.
就是永生神的兒子.
所以決心跟隨他.
換句話來說.
靈姐妹.
彼得他之所以獲得耶穌的稱讚.
不是因為他答對了耶穌的身份.
不是因為他答對了耶穌是誰.
而是因為他跟隨耶穌.
在耶穌訓練他的過程裡.
他願意不斷去回應上帝給他的啟示.
他願意不斷去回應.
耶穌給他的一些挑戰和新嘗試.

$^{361}$來累積對耶穌的信心.
所以來到此時此刻.
當耶穌來問.
你覺得我是誰的時候.
他就能夠夠膽宣認.
你就是永生神的兒子.
你就是我的救主.
所以弟兄姊妹.
很值得我們停下來去想一想.
其實我們知道.
要答對耶穌是誰難不難?.
難不難?.
如果純粹只是答對答中耶穌的身份.
其實不是很難.
一個不相信的人都可以答對.
只要他通過一些看書的研究.
其實他可以答出耶穌的真正身份.
是基督是永生神的兒子.
甚至比你和我.
可以答得更加仔細詳盡深入.
不過答對耶穌是誰.
和認信耶穌是誰.
其實我們知道是兩回事.
後者才能夠真實反映.
我們和耶穌的關係.
和他的距離到底是遠還是近.
耶穌對我們的生命.
到底有多大的影響力.
你看到耶穌真正喜悅稱讚的.
其實不是答對他是誰的人.
不是那些純粹頭腦上.
知道他是誰的人.
他期望我們就像彼得那樣.
成為一個認信他的人.
所以今天這段經文.
其實都在考問我們.
當我們宣認耶穌是基督.
耶穌是永生神的兒子的時候.
我們是否真的認信他?.
是否真是一個課真價實的門徒?.

$^{401}$我記得曾經有位弟兄.
他打算全職事奉布督神學.
但他又不太確定.
到底上帝是否想我這樣做呢?.
於是他開始有很多困惑.
他想找牧師問清楚.
於是牧師就問.
「你想全時間服侍上帝,是嗎?」.
「你不是在做壞事吧?」.
「既然是一個好的意念」.
「為何有這麼多擔憂,憂慮?」.
「你一心想事奉,又不是搞事」.
「為何要為首為尾?」.
這位弟兄就說.
「我害怕,我怕上帝」.
「如果不是想我走這條路,怎麼辦呢?」.
「我們都經常有這些疑惑」.
「我這樣選擇,為上帝擺上」.
「上帝是否喜歡呢?」.
「我想我們多少都有個疑惑」.
於是牧師就問這位弟兄.
「耶穌到底是否你的救主呢?」.
「當然是,耶穌當然是我們的救主」.
「既然如此」.
「就算你所說的可能行錯,選錯」.
「難道在你整個申請布督的過程中」.
「耶穌沒有能力拯救你?」.
「沒有能力將選錯,行錯的你」.
「帶回一條正路嗎?」.
「你剛才不是認信祂是你的救主嗎?」.
弟兄姊妹.
我們確實認耶穌是基督.
是永生神的兒子.
但在現實生活中.
我們會否總覺得.
有些事情.
連耶穌這位救主.
都不能擺平.
都不能拯救.
以致我們都覺得.

$^{441}$自己出手解決比較好.
又或者我們.
口稱耶穌是永生神的兒子.
但在現實生活中我們對焦的.
是否只是眼前的成敗得失?.
我們有沒有願意被耶穌塑造.
改變我們的身份.
來配合上帝永恆的計劃?.
這是我們宣認耶穌是基督.
耶穌是永生神的兒子.
背後真正一個實際操作的意義.
事實上當我們宣認的時候.
並不是表示純粹在得救的事情上.
在天堂的入場券裡.
我們需要依靠耶穌.
當我們宣認耶穌是永生神的兒子.
是基督是救主的時候.
其實背後生活中大大小小的事情.
我們沒有辦法靠自己做得滴水不漏.
人生總會有所謂計錯數.
總會有選錯.
甚至犯錯的時候.
否則我們就不需要一位救主.
祂真是願意成為你和我生命的主.
並且祂說了會拯救我們到底.
所以其實你可以放膽為耶穌冒險.
基督徒有時可能很擔心.
做很多特別是事奉上的決定.
會畏首畏尾.
覺得這樣可能是上帝的心意.
好像我又沒有什麼能力.
又沒有什麼經驗.
但上帝之所以透過耶穌拯救我們.
其實就是想我們放膽為祂冒險.
因為祂會拯救帶領我們.
所以弟兄姊妹其實.
你可以放膽依靠祂.
面對生活不同的挑戰.
放膽為耶穌冒險.
突破一些可能在你成長中.

$^{481}$必然會遇到的瓶頸位.
並且因為這位永生上帝的兒子的緣故.
你願意開始改變一些固有的價值觀念.
既然他是一位永生上帝的兒子.
你願意跟隨他.
你的目光需要重新調教.
你生活的優先次序也需要改變.
學習將你的人生.
放在具有永恆價值的投資上.
這是我們對他的宣認.
一種實際的運作.
當然耶穌沒有允許過.
當我們這樣做的時候.
過程一定是一帆風順.
甚至乎真的會跌跌碰碰.
會損手爛腳.
但耶穌說了什麼?.
祂會包底.
祂會拯救我們到底.
既然我們宣認祂是救主.
是永生神的兒子的時候.
即使在現階段我們會跌跌碰碰.
但祂會確保我們.
可以安穩在上帝永恆的裡面.
所以無論在這一刻你遇到什麼情況.
有什麼事情.
假如我們真的認定耶穌是基督.
是救主的時候.
你真的可以放心將這一切交託給祂.
祂真的願意拯救我們.
幫助我們.
並且好像聖經所說.
將永生的道路指示給我們.
單單這樣.
上帝救我們其實是要用我們.
祂要使用我們成為合用的器皿.
就好像祂即將要使用彼得.
來任命祂建立教會一樣.
所以我們繼續看.
耶穌是怎樣回應的?.

$^{521}$耶穌說你是彼得.
其實要留意.
在福音書的一開始.
當耶穌呼召西門的時候.
他已經幫他起了一個名字.
阿南文叫基法.
也等於彼得的意思.
所以一開始已經這樣叫他.
為什麼這裡又要再說你是彼得呢?.
其實不是再一次和他重新學改名.
不是和彼得重新改名.
而是要解釋為什麼要叫他做彼得.
希臘文彼得這個字是盤石石頭的意思.
耶穌稱他為盤石.
不是純粹因為彼得夠硬朗.
性格剛烈,能夠站得住.
稍後你知道其實他有很多軟弱.
甚至乎他不認耶穌.
叫他做彼得.
叫他做盤石.
是表明他稍後在上帝的角度.
到底角色,到底功用是什麼.
耶穌說你是彼得.
其實有沒有留意跟彼得宣認.
你是基督.
其實是互相呼應的.
你既然認信我做你的救主.
是永生神的兒子.
你不單單是有福.
我也要任命你成為建造教會.
一塊很重要的材料.
互相呼應.
當然我們知道.
特別在加拉太書或者以弗所書告訴我們.
教會的儲石和根基.
不限於彼得一人.
也包括其他使徒.
例如約翰或雅各.
跟彼得一樣.
他們也要成為教會重要的一塊石材.

$^{561}$靠著耶穌這一塊最重要的防角石.
來一起參與教會的建立.
當然最終.
我們說是耶穌基督親自建立教會.
但是彼得在初期教會.
扮演著一個很獨特的角色和功用.
這確實是毋庸置疑.
當日在凱撒利亞菲勒比.
他宣認耶穌是基督.
是永生神的兒子.
日後他要繼續發揮盤石的功用.
在《使徒行傳》裡面.
你繼續看到.
他願意在群眾面前公開宣揚.
那一位已經從死裡復活的主.
他是基督.
是人類的救主.
是永生神的兒子.
你看到上帝大大使用彼得.
藉著他的宣講.
帶領了數千人信主.
也讓教會群體.
從那一刻開始被建立起來.
第十九節.
耶穌繼續對彼得說.
「我要把天國的約時給你.
凡你在地上所捆綁的.
在天上也要捆綁.
凡你在地上所釋放的.
在天上也要釋放」.
「捆綁」和「釋放」這兩個詞.
是當時猶太的拉比.
很常用的術語.
「捆綁」其實是對某人加以制裁.
至於「釋放」是對某人宣告的寬容.
耶穌不單止將「捆綁」和「釋放」的權柄.
交給彼得.
稍後你會看到.
結後兩章第十八章.
同一番說話也出現.

$^{601}$不過這次是說.
耶穌要將「捆綁」和「釋放」的權柄.
賜給所有的門徒.
授權他們來建立教會.
包括擬定教會的準則.
又或者執行教會的紀律.
最明顯的例子.
你會在《荃唐傳》看到.
彼得如何實行「捆綁」的權柄.
反映在阿拿尼亞撒菲拉事件.
他們欺控聖靈的時候.
彼得宣告他們的罪狀.
於是聖靈就馬上擊殺了二人.
又或者彼得如何在《荃唐傳》.
行使他「釋放」的權柄.
在他帶領外邦人.
哥尼留一家信主的事情上.
聖靈感動他.
告訴他福音不單止給自己人.
不單止給猶太人.
在上帝的心意.
是要擴展到給猶太人.
認為不潔的外邦人身上.
就是這樣.
彼得帶領上帝耶穌.
給他「釋放」的權柄.
打開了外邦人信主的福音的門.
而耶穌最後也說.
他說「陰間的權柄不能勝過教會」.
權柄這裡.
如果你看看畫本.
特別是他有一個括弧.
印了權柄的原文是「門」.
還記得我們剛才一開始提到.
凱撒利亞的菲勒比有個「陰間之門」嗎?.
有印象嗎?.
耶穌其實這裡是說.
陰間的門不能勝過教會.
那個大山洞.
那個陰間的門.

$^{641}$當時的人覺得.
這裡就是直通陰間.
耶穌為什麼要這樣說.
陰間的門不能勝過教會.
是什麼意思呢?.
你知道門是否一個用來攻擊的武器.
不要把它聯想成摺椅.
你不會用門去攻擊別人.
門是用來做什麼的?.
是防禦的.
我關了門.
外面的人又如何?.
進不來拿東西.
門是防禦的.
但耶穌說這裡陰間的門.
不能勝過教會.
意思是什麼呢?.
其實在說教會.
不僅是被動防守.
敵人或死亡的攻擊.
反而教會可以主動出擊.
你去踢爛攻破陰間的門.
進去拿東西.
進去奪回那些.
被死亡權勢所限制.
但原本屬於耶穌基督的迷洋.
你衝進去拿回.
那扇門難阻不了.
耶穌基督的應許.
所以這段經文其實很好提醒我們.
特別我想在這幾年.
其實大家看到香港的教會.
都面對著成員流失的情況.
對不對?.
我想今天這段經文很好提醒我們.
很多時候最好的防守是什麼呢?.
進攻.
進攻往往有時是最好的防守.
一方面面對著今天教會的剛勁.
我們當然要守住我們已有的陣地.

$^{681}$守住一些每間教會.
它很寶貴的屬靈傳統.
守住弟兄姊妹的相交關係.
以至不會內部出現混亂分裂.
穩住了教會內部的陣腳.
但正所謂「九守」有時會怎樣呢?.
「九守」有時真的會必失.
舉個例子.
譬如我們唱詩.
我們有時唱到每一句尾句都要保持長拍.
可能要唱四拍,六拍甚至再唱.
要怎樣將長音唱得穩定.
原來在意識上.
你不可以只是保持住它.
不要讓它跌.
當你純粹說不要讓它跌.
其實就會慢慢跌.
你要保持長音唱得穩定.
其實你要在每一拍.
譬如四拍.
每一拍都要用適當的力度.
托起它.
聽起來整個音就會很穩定地延續.
相反如果只是保持住它.
其實會慢慢跌.
然後就會有走音的情況.
同樣教會今天的成長.
不可以單靠內部維繫.
這固然重要.
但耶穌藉著今天的經文提醒我們.
要主動出擊.
我們需要走入人群裡.
開發不同的戰線或陣地.
更加拉闊我們服事或福音的工作.
要在這個時間讓更多人接觸到教會.
所以真正的弟兄姊妹.
歷史告訴我們.
教會是否穩定.
能夠健康發展.
和外面的環境是否艱難.

$^{721}$其實沒有必然關係.
教會能夠維持健康的成長.
與甚麼有關?.
與教會的群體有沒有承傳.
承接福音的使命有關?.
事實上客觀來看.
假如我們是初期教會的信徒.
回到耶穌的年代.
你會否打算在耶路撒冷這個地方.
建立初期教會?.
你會否這樣計劃?.
在耶路撒冷這個地方.
在第一世紀耶穌的時代.
在那裡建立教會?.
你覺得這是否一個明智之舉?.
其實不是一個明智之舉.
為甚麼?.
耶路撒冷是哪個宗教的根據地?.
是猶太教的根據地.
門徒在耶路撒冷.
是否被人尊重?.
不被尊重.
被視為鄉下人.
沒有見過世面的.
沒有學識的小民.
那麼耶穌呢?.
耶穌在耶路撒冷受歡迎嗎?.
也受歡迎的.
不過最終他的下場是怎樣?.
是被人嫌棄.
甚至被殺害的地方.
所以如果純粹用現代的角度.
用一個商業管理的角度.
你們的教主剛剛被殺害.
你們竟然敢在這裡插旗開教會?.
明不明智?.
其實沒有人會這樣做.
因為在一個當時充滿.
對基督信仰充滿不利條件的環境.
來建立教會.

$^{761}$其實沒有人覺得會成功.
甚至乎輸定.
但是初期教會的發展.
他又讓我們看到.
他的成長真是見證著.
耶穌所說的才是真理.
就是陰間死亡的權柄.
哪怕是迫不及待苦難.
都不能夠阻止耶穌的教會.
來建立和擴展.
所以有人說得很好.
其實我們的信仰離不開兩個字.
知不知道是哪兩個字?.
離不開受苦兩個字.
基督信仰和受苦是分不開的.
我們的最高負責人耶穌基督.
他本身是一位怎樣的救主?.
是一位受苦的救主.
作為教會的防國石.
耶穌都是通過苦難.
成就了上帝的拯救計劃.
承傳了教會的開展.
所以作為教會群體一份子的我們.
自必然也會走這條路.
來建立教會.
在我們的信仰歷程裡.
在實踐愛.
實踐福音使命的過程裡.
我們一定會像耶穌一樣.
經歷不同程度的苦.
不過這些苦.
耶穌說是有益的.
可以與他連結.
所以我想起.
以前在小組有位弟兄.
他有個分享.
他說他有一段時間.
在信仰上沒有突破.
他覺得很閹悶.
於是他很想衝過這個位置.

$^{801}$於是他很想向上帝祈求.
你猜他求什麼?.
他求上帝讓他經歷一些苦難.
奇怪嗎?.
你和我都不會這樣祈.
哪有人這樣祈苦難.
沒有人這樣做.
所以我們一群人都在笑.
怎麼搞的?.
求苦難臨到你身上?.
事後回想這段話.
有沒有智慧?.
其實很有智慧.
我們都希望與耶穌同行.
我們都希望跟隨他.
分享到耶穌的生命.
自然就會怎樣?.
分享到耶穌受苦的生命.
分享到耶穌的苦難.
分享到耶穌的艱難.
背後就是提醒我們.
我們要隨時有為主.
吃苦受苦的心智.
當我們表示認耶穌是基督.
認祂是永生神的兒子的時候.
但好消息.
我們不只是因為受苦捱打.
因為耶穌應許我們.
雖然世上會有苦難.
但祂也保證我們會勝過世界.
因為祂已經勝過死亡.
勝過苦難.
成為了我們今天的安慰力量和盼望.
所以弟兄姊妹很盼望.
藉著今天的經文.
我們互相勉勵.
一方面我們樂意承擔.
耶穌基督託付給我們的科興使命.
我們願意面向服侍我們的社群.
讓教會成為一個有活力.

$^{841}$有凝聚力和有歸屬感的群體.
很真實的我們如何建立歸屬感.
就是透過你的參與.
透過你去實踐耶穌基督的吩咐.
我們一起對準耶穌的時候.
歸屬感,凝聚力,活力都會產生.
盼望藉著這段經文.
我們能夠幫助更多的人認識耶穌.
就好像我們這樣.
幫助他們認耶穌是基督.
是永生神的兒子.
我們同心禱告.
你的心恩厚愛拯救我們.
讓我們真的可以脫離死亡陰間的核制.
在你裡面過著不一樣的生命.
上帝幫助我們.
藉著今天的經文提醒我們.
我們不單止知道你是基督.
是永生神的兒子.
我們更加宣認你就是基督.
就是我們的救主.
就是永生神的兒子.
以致我們的行事為人.
我們不再純粹著眼於今世有限的事情上.
我們更加願意投放在上帝的喜悅.
有永恆價值的事情上.
上帝在過程中我們會有膽怯.
我們會有很多懷疑.
但耶穌求你的救恩幫助我們.
既然你願意拯救我們到底的時候.
就讓我們放膽.
放手地跟隨你.
哪怕為耶穌你的冒險.
我們確信耶穌你的保護.
你的供應.
你的看顧是足夠的.
以致我們在人生裡.
無論或順或逆.
當我們在捉緊你的時候.
我們就有力量去面對.

$^{881}$上帝你給我們的不同挑戰.
和每一個成長的機遇.
我們知道最終.
上帝你一定會拯救我們.
進到你永恆的裡面.
為此幫助我們.
幫助教會.
在艱難的日子.
我們更加堅守你託付給我們的使命.
以致我們真的能夠幫助更多的人.
來認識你.
將更多失散的迷樣.
從陰間的門裡面去拯救釋放出來.
以致我們真的能夠課真價實地.
承認耶穌.
你就是我們的救主.
是永生神的兒子.
求你帶領教會.
求你繼續去使用我們每一位.
願你聽我們在你面前的禱告.
奉基督耶穌你得勝的名字.
來祈求.
阿們.
其他事項.
\newpage



\section{詩篇 71:5-18}
\label{sec:esPM0k_dZBY}
\textbf{2022年11月20日長者主日崇拜直播｜羅志雄牧師｜恩典的年歲｜詩篇71\_5-18}
\newline
\newline
連結: \href{https://youtube.com/watch?v=esPM0k-dZBY}{\texttt{https://youtube.com/watch?v=esPM0k-dZBY}} ~~~~ 語音日期: 2022-11-22
\newline
\newline
\hyperref[sec:y7hPaMojRiM]{\small{< < < PREV SERMON < < <}}
~
\hyperref[sec:index_chronic]{\small{[返順時目]}}
~
\hyperref[sec:4hDfY5nlUxk]{\small{> > > NEXT SERMON > > >}}
\newline
\newline
詩篇 71:5-18
\newline
\begin{longtable}{cl}
\hline
\hline
章節 & 經文 (和合本修訂版)\\
\hline
71:5 & \begin{tabularx}{0.7\textwidth}{X} 主耶和華啊，你是我所盼望的；自我年幼，你是我所倚靠的。 \end{tabularx} \\ \\ \relax
71:6 & \begin{tabularx}{0.7\textwidth}{X} 我自出母胎被你扶持，使我出母腹的是你。我要常常讚美你！ \end{tabularx} \\ \\ \relax
71:7 & \begin{tabularx}{0.7\textwidth}{X} 許多人看我為異類，但你是我堅固的避難所。 \end{tabularx} \\ \\ \relax
71:8 & \begin{tabularx}{0.7\textwidth}{X} 我要滿口述說讚美你的話終日榮耀你。 \end{tabularx} \\ \\ \relax
71:9 & \begin{tabularx}{0.7\textwidth}{X} 我年老的時候，求你不要丟棄我！我體力衰弱時，求你不要離棄我！ \end{tabularx} \\ \\ \relax
71:10 & \begin{tabularx}{0.7\textwidth}{X} 我的仇敵議論我，那些窺探要害我命的一同商議， \end{tabularx} \\ \\ \relax
71:11 & \begin{tabularx}{0.7\textwidth}{X} 說：「神已經離棄他；你們去追趕他，捉拿他吧！因為沒有人搭救。」 \end{tabularx} \\ \\ \relax
71:12 & \begin{tabularx}{0.7\textwidth}{X} 神啊，求你不要遠離我！我的神啊，求你速速幫助我！ \end{tabularx} \\ \\ \relax
71:13 & \begin{tabularx}{0.7\textwidth}{X} 願那與我為敵的，羞愧滅亡；願那謀害我的，受辱蒙羞。 \end{tabularx} \\ \\ \relax
71:14 & \begin{tabularx}{0.7\textwidth}{X} 我卻要常常仰望，並要越發讚美你。 \end{tabularx} \\ \\ \relax
71:15 & \begin{tabularx}{0.7\textwidth}{X} 我的口要終日述說你的公義和你的救恩，因我無從計算其數。 \end{tabularx} \\ \\ \relax
71:16 & \begin{tabularx}{0.7\textwidth}{X} 我要述說主耶和華的大能，我單要提說你的公義。 \end{tabularx} \\ \\ \relax
71:17 & \begin{tabularx}{0.7\textwidth}{X} 神啊，自我年幼，你就教導我；直到如今，我傳揚你奇妙的作為。 \end{tabularx} \\ \\ \relax
71:18 & \begin{tabularx}{0.7\textwidth}{X} 神啊，我年老髮白的時候，求你不要離棄我！等我宣揚你的能力給下一代，宣揚你的大能給後世的人。 \end{tabularx} \\ \\
[1ex]
\hline
\hline
\end{longtable}
$^{1}$《聖經》第七十一篇 五至十八節.
主耶和華,你是我所盼望的,.
從我年幼,你是我所倚靠的..
我從出母胎被你扶持,使我出母胎的是你,.
我必常常讚美你.許多人以我為怪,.
但你是我堅固的避難所.你的讚美,你的榮耀,.
終日必滿了我的口.我年老的時候,求你不要丟棄我;.
我力氣衰弱的時候,求你不要離棄我..
我的仇敵議論我,那些窺探要害我命的彼此相議,說:.
神已經離棄他,我們要追趕他,捉拿他吧,因為沒有人搭救..
神啊,求你不要遠離我;我的神啊,求你速速幫助我;.
願你與我性命為敵的羞愧避滅;願你謀害我的受辱蒙羞..
我卻要常常盼望,並要越發讚美你..
我的口終日要述說你的公義和你的救恩,因我不計其數..
我要來說主耶和華大能的事,我單要提說你的公義..
神啊,自我年幼時,你就教訓我,直到如今,我傳揚你奇妙的作為..
神啊,我到年老發白的時候,求你不要離棄我,.
讓我將你的能力指示下代,將你的大能指示後世的人..
聖經獨筆..
等姐妹平安..
今日是長者日,祝願可愛的長者平安喜樂..
我也是在這個行列,包括我自己也要祝願..
今日在坊間都有很多優惠,不知道大家知不知道..
這些是對長者的優惠和回報..
大臣給年長者厚厚的福氣..
剛才聽到師班出來獻頌,那裡有十多位年長者,.
你知不知道那裡加起來有多少歲呢?.
以我粗略平均計算,其實已經過了千歲了..
所以剛才的獻頌可以說是千歲獻頌,.
讓他們的歌聲頌讚歸給上帝,.
因為可以帶領他們經歷人生的歲月..
我們一起來禱告吧..
生命的旅程是在上帝你自己的掌管裡面,.
無論順逆,起跌高低,.
主我們深信你託戴我們一生..
主願上帝,你的恩典厚厚的在我們身上,.
我們不怕失迷,不怕在艱難,風雨之中,.
因為你常常的與我們同在..
主祈求你的話語成為我們的幫助,.
讓我們剛才所讀的聖經都能夠有從你而來的亮光,.

$^{41}$在我們的人生路裡面,每一天都能經歷到上帝你自己的同行同在..
願你臨到我們,聖靈來啟迪我們,開我們的心竅,.
讓我們能夠裝載你自己的聖言,.
你引導我們行走前面的道路..
阿們..
不知道大家有沒有看過這齣電影,.
這齣是甚麼電影呢?.
沒有人讀的,早課有人反應的..
《爸爸可否不要老》,英文是The Father,.
其實這套電影取材自舞台劇,.
當中的主角是八十多歲的,.
他飾演一個患上智力障礙,或者叫老退化的一位年長者..
如果是這樣的話,整套電影有甚麼好看呢?.
今天我也不會劇透這套電影,.
但是這齣電影,.
他是以一個反常聚事的進路,.
帶出年長者在他罹患老退化的真實故事..
在當中,他記錄了混亂,無記性,.
情緒化,甚至是疑患似真的很多生活場景..
這種電影手法,不知道大家看了之後有甚麼想法呢?.
本會帶給大家很多剝洋蔥,流眼淚,.
但我相信導演是想讓觀眾感同身受,.
帶出一個很重要的訊息,.
就是要關心家中或身邊的年老長者,.
怎樣關心呢?.
是否只給提供物質,.
這些需要呢?.
還是更加關心他們的情感,.
怎樣理解他們,怎樣認同他們,.
能夠更加了解他們生理的變化呢?.
的的確確,在這些真實的場景中,.
其實我們未必每一個家庭或每一個人都能夠遇到,.
但我們的人生,.
由出生到年幼成長,成年,壯年,.
甚至到年老發白的階段,.
當然生活上會經歷很多不同的際遇,.
我們會面對很多風風雨雨,.
但我們的生活都有上帝的保守,.
我們今天可以來到這裡一起敬拜,.
一起來讚美,.

$^{81}$一起來領受神的話,.
相信全都是上帝的恩典,.
剛才讀的聖經71篇,.
整篇的詩篇,.
如果有留意經文的內容,.
都知道是一位年老發白的人,.
一首祈禱的詩篇,.
在第8節和18節已經清楚說到,.
作者就是年老發白的人,.
但作者是誰沒有說到,.
有古抄本將詩篇70篇和71篇合併成一篇,.
以70篇為71篇的序言,.
70篇如果了解的話,.
標題裡面也有寫了大衛的紀念詩,.
整篇71篇可以分為三個段落,.
1至3節就是.
這位作者向神呼求,求神幫助,.
4至13節,.
14至24節,.
如果大家有留意整篇詩篇,.
當中出現了中日,.
這裡有三段,.
所以有悉經學家也會將它分段,.
1至8節,9至15節,.
16至24節,.
剛才永祥弟兄和我們讀的5至18節,.
介乎兩者之間,.
我發現到詩人就是.
以整個人一生的階段,.
由母覆到年幼,.
直到年老發白的時候,.
所說的人生信仰的一個歷程,.
我們來看詩篇71篇1至5節,.
好,想邀請弟兄姊妹來讀1至5節,.
預備起.
(念書).
好,謝謝.
詩人祈禱的時候,.
一開始就以神為耶和華為主,.
講明了神在詩人生命裡有著很重要的位置,.

$^{121}$所以詩人呼喚耶和華的名字,.
這就是她唯一專心一意信靠的對象,.
詩人一生的際遇並不是事事順利,生活安定,.
相反的,在經文裡已經透露,.
她應該經歷了生活處境的艱難,.
受到不同程度的逼迫,.
甚至在她生活裡都是步步為營,.
期間甚至遇上風高浪急的危機,.
甚至有很多惡人不斷地攻擊她,破壞她,.
甚至趁詩人軟弱的時候,.
更加以兇殘的方式來攻擊,.
令她不單單是身體,.
甚至精神上更加惶恐,因為擔憂生命的安危,.
不過,詩人在禱告裡,.
她仍然很肯定一件事,.
就是神一定會出手幫助她,.
其實已經發生了,.
所以詩人不斷地認定神已經來解救她,.
詩人更加表達,.
這位她所謂一頭靠的神,.
是答救她,救拔她,拯救她,.
不單單是這樣,更加給她保護,.
這個保護方式是以盤石,岩石,山寨,.
是當時候一個最高級別的保護,.
可以來為她阻擋一切的危險,.
現在這位詩人進入去年老另一個人生階段,.
在她的禱告裡,她回憶了剛才所描述的一些惡劣的遭遇,.
雖然是這樣,但她仍然再次向神去呼求,去禱告,.
因為她確實領受到神的福氣,.
這是上帝給她的恩典,.
當她呼求的時候,.
她呼求神,要求神投靠她,.
求神叫她永不羞愧,.
不羞愧在呂鎮忠的譯本中,.
就譯為不失望,.
不失愧,不失望,.
是表示詩人在當下仍然心存盼望,.
所以在第五節的時候,.
主耶和華說:你是我所盼望的,.
她雖然是一把年紀,身體越來越多毛病,軟弱,.

$^{161}$甚至可能出現了無法避免的身體衰殘,甚至死亡,.
她仍然有盼望,仍然堅信,.
因為幫助是從耶和華而來,.
信就是對所盼望的事的一種把握,.
因此信心加上盼望,.
兩者連結的時候,.
對神的那種倚靠,.
詩人是沒有動搖的,.
她不會因為周遭的環境,.
而對神產生任何的懷疑,.
使徒彼得在耶穌的門徒裡面,.
他是一個領袖,.
他仍然要學習信心的功課,.
其中有一幕聖經裡面福音書記載,.
彼得在風浪裡面,.
他仍然求主准許他在海面上走到耶穌那裡,.
雖然在過程裡面,.
他面對一些軟弱,.
但耶穌來拯救他,.
到他年老的時候,.
他繼續來信靠這位主,.
相信這個經歷,.
在他年輕的時候,.
他領受到耶穌給他的這一份確據,.
所以他能夠勉勵信徒學習信心,.
信靠的功課,.
我服侍長者的群體,.
在皇浸裡面有十年的時間,.
每逢去探訪一班長者的時候,.
他們都很分享他們,.
他們如何能夠去到這個年歲,.
甚至會講他們過往人生的經歷,.
當然在那個年代,.
他們實在生活很艱苦,.
又會面對戰爭的危險,.
甚至他說不單止這樣,.
在窮困的階段裡面,.
他們要打幾份工,.
因為他們所養育的,.
不只是一兩個兒女,.

$^{201}$而可能是很多的兒女,.
他們捱得很辛苦,.
甚至有時候在工作期間,.
都面對很多的危險,.
甚至去到生死的邊緣,.
當他們這樣分享的時候,.
當時雖然他們未信主,.
但他們經歷了上帝,.
在現在能夠認識耶穌的時候,.
他還分享,.
如果不是當時候有上帝的保守,.
他們不能夠活到今天,.
白髮兩鬢,.
可以回來教會,.
可以跟弟兄姊妹一起分享,.
一起讚美,.
他說這些都是傳示上帝的恩典,.
私人的禱告,.
從盼望,倚靠,.
轉去另一個方向,.
就是讚美的主題,.
經文在第六節是這樣說,.
「我從母胎被你扶持,.
使我出母腹的是你,.
我必常常讚美你,.
許多人以我為怪,.
但你是我所堅固的避難所,.
你的讚美,你的榮耀,.
終日滿了我的口.」.
私人來到這個年老的光景,.
他自問一生倚靠神,.
經歷過任何信逆際遇,.
在荒漠裡,.
神都拯救他,幫助他,.
所以他必然發出讚美,.
因為他知道生命的源頭是從神而來,.
生命是非常美妙,.
未出母胎,.
上帝已經知道他,.
也知道他未來的人生路,.

$^{241}$詩篇139篇裡.
都很詳細地描述整個.
由母腹到成長的階段,.
實在看到人生裡的生命氣息.
全是上帝所賜予,.
試問一下,弟兄姊妹,.
如果沒有生命的時候,.
如何能夠走到世上.
現在的生命歷程呢?.
又或者大家可以安坐這裡,.
去經歷上帝呢?.
甚至,.
會不會還有機會來到.
年老發白的日子呢?.
所以詩人在這裡.
他必然要來讚美,.
有人會說,.
人自出生,年幼,成長,成年,.
壯年,到年老發白的階段,.
就叫做一條線,.
線性的發展進程,.
於是有很多的生理學家,.
心理學家,.
不斷地研究人生命奧妙的成長,.
他們這樣研究,.
這樣去做很多的.
一些資料的搜集的時候,.
其實他們想做些什麼呢?.
他們企圖在生理醫學,.
心理學,.
對人做很多的預測,.
嘗試去了解,.
在這麼複雜的一種階段裡面,.
他們人生的互動是怎麼樣,.
甚至,.
如何有這樣的能量,.
來走他們的人生路,.
又或者去了解一下,.
他們用什麼的態度,.
來做生命的反省,.

$^{281}$很可惜,.
坊間裡面沒有任何精闢,.
講到關於對生命的重責,.
不知道大家有沒有想過,.
如何去回應這個提問呢?.
唯有私人,.
就是以信仰的角度,.
來講出神對生命的體重,.
對人施行的慈愛,.
一次又一次地去拯救人,.
所以向他發出頌讚,.
因為他的生命,.
私人的生命,.
完完全全地靠賴,.
信靠神的信實,.
神的公義,.
神的榮耀,.
所以私人要終日去頌讚他,.
在幾天之前,.
我在快餐店下午茶,.
突然間出現一位舊同事,.
她是一位姐妹陪她看醫生覆診,.
於是就在快餐店,.
我就覺得奇怪,.
她告訴我,.
她剛剛在QE覆診完畢,.
為什麼會無端來到旺角快餐店呢?.
她說有些事要做,.
但我通常都會遲到快餐店,.
但當時我還在想去哪裡,.
不過終於去了那裡,.
她就在那裡說,.
「我在二十年前在病房認識你,.
現在在這個地方見到你,.
我相信不是巧合,.
是上帝的恩典.」.
這位舊同事在念書上,.
其實也有很多訊息,.
她退休幾年,.
她患了癌症,.

$^{321}$已經有一段日子,.
她還在接受治療,.
她在念書上也知道我做了傳道牧者,.
她說亦感恩,.
但重點是,.
我在她的Facebook上看到,.
在她患病的過程中,.
她仍然讚美和感恩,.
因為她說她有生命氣息,.
是上帝讓她在這段日子,.
可以陪伴她的家人,.
特別是她可愛的小孫女,.
相信現在我看到的時候,.
她應該有兩歲多,.
弟兄姊妹,.
一個未信主的人,.
在退休之後,.
應該是很享受生活,.
但突然患病,.
可能打亂了她整個退休計劃,.
怎可以讚美和感恩呢?.
可能會擔心自己的病情,.
不知道如何治療,.
甚至治療的時候會否很辛苦,.
考慮這些因素,.
確實是一件不幸的事,.
但當信主的人,.
他們真的聽到要讚美和感恩,.
是一件很難想像的事,.
甚至是奇怪的事,.
就像第七節那樣說,.
許多人以我為怪,.
但你是我堅固的避難所,.
以我為怪,.
「為怪」這個字可以解釋為.
歧視神跡,.
又甚或可以指.
神的能力是一種特殊的顯現,.
所以如果原文成句直譯,.
對許多人而言,.

$^{361}$我是奇怪的,.
但你是我堅固的避難所,.
如果這篇詩是大衛所作的時候,.
人人都會逃避大衛,.
甚至會嘲笑他,.
為何你淪落到這個光景?.
但大衛這位詩人說,.
神仍然是他的避難所,.
「為怪」這個字.
第一次出現在《出埃及記》,.
那裡翻譯為神跡歧視,.
第一是指摩西的象變蛇,.
第二是降十災,.
也在《申命記》中,.
摩西大理以色列人回顧.
曠野這段歷史的時候,.
每次提到的記號.
其實就是「為怪」這個字,.
是指舊因歷史中的記號,.
在曠野中,.
他們沒有東西吃,.
有嗎哪,.
沒有水喝,.
盆石可以出水,.
日間有雲柱,.
夜間有火柱,.
這些都是記號,.
所以這一節的經文,.
詩人提及「為怪」,.
我相信是一個笑頭或記號,.
意思是要向人宣揚神的大能,.
神的美善,.
神的慈愛和拯救,.
我想起我自己.
2006年在夏令營.
蒙召讀神學的事,.
其實是預備2007年.
會報讀神學院,.
但是在2007年初的時候,.
太太確診了胃癌,.

$^{401}$其實發現的時候已經很後期,.
之後安排做手術,.
第二天我去探望她,.
太太,我都問她,.
你做完手術痛不痛?.
情況如何?.
她都不是這樣說,.
她問我,.
你報讀神學的事,.
正在做嗎?.
或者開始了嗎?.
這也令我….
為何太太會這樣說呢?.
反而她不是擔心自己的病情,.
而是神的照明,.
如此之後,.
太太大概在10月份離世,.
之後我很多時候都去禱告,.
尋求神的心意,.
最終在2009年的時候.
就入讀神學院,.
但在期間仍然在病房工作,.
當時上帝的恩典,.
讓我能夠升職,.
當我說辭職的時候,.
同事都覺得很奇怪,.
你剛剛升職,.
怎會這樣離任?.
沒錯,有時我們跟隨耶穌的人,.
跟隨神的人,.
我們所做的,.
可能真的別人不了解,.
會怪怪的,.
丁姐妹,今天有沒有一些.
你覺得在你的生命上,.
別人覺得以為怪的,.
是一些記號,.
是一些申命記號,.
所以來到講述.
神在你身上的生命故事,.

$^{441}$可以成為讚美神的一些記號嗎?.
來到經文第9至18節,.
很想邀請丁姐妹一起來中讀,.
第14至18節,.
預備起,.
.
好,謝謝,.
詩人在這裡承認自己是一位.
白髮蒼蒼,.
年紀老邁的一個人,.
他又發現自己力氣衰弱,.
他是確確實實在神面前,.
去坦誠,.
去表白自己這一種階段,.
詩人他沒有自棄,.
甚至躺平下來,.
什麼都不做,.
相反,.
他是繼續迫切來禱告,.
向神尋求,.
並且他不斷地向神呼求,.
不要離棄他,.
不要遠離他,.
詩人這一種做法,.
不是我們今天所說的那種隱蔽長者,.
反而他是更加積極地生活,.
究竟隱蔽長者的心態是怎樣的?.
不知道大家有沒有見過,.
很多的老年學家,.
老年精神健康專家,.
又或者社工,.
他們研究,.
做了很多這樣的調查,.
發現原來隱蔽長者的心態,.
是有以下幾種情況,.
可能會受到生理,心理,.
甚至現實環境所影響,.
他們覺得老了,.
體弱,多病,.
想法通常都是消極的,.

$^{481}$孤僻,寂寞,.
甚至他們覺得沒辦法學一些新的東西,.
因為學新東西感覺會有很大的壓力,.
因為會沒有記性,.
甚至也覺得自己抵抗疾病的能力,.
都是大不如前,.
就算有的,.
他們都覺得發現了,.
還是治不了的,.
退休之後收入減少,.
積蓄遲早都會用完,.
他們這樣的想法,.
就是吃老本肯定會坐食山崩,.
怎樣可以安老呢?.
實在很艱難,.
他們想起就有心無力,.
於是會自覺是廢老,.
在這裡等死,.
於是在坊間裡面,.
很多學者探討,.
有什麼理論或者學說,.
可以讓一班長者,.
能夠過渡老化進程,.
於是又有很多工作坊,.
或者研討會,.
去了解安度晚年應該怎樣做,.
他們通常都是用他們的學術,.
理論基礎,.
來想一些能夠安度老化的進程,.
不知大家聽完這些老化,.
什麼老化的時候,.
你會覺得他們用什麼進路,.
用什麼進程,.
讓年老的長者,.
可以安度晚年呢?.
他們研究出有什麼呢?.
正向老化,.
健康老化,.
有生產力老化,.
活躍老化,.

$^{521}$不知大家覺得是怎樣呢?.
當我們翻開整本聖經,.
又或者我們看到詩篇的時候,.
其實沒有任何指引,.
教我們信徒弟兄姊妹,.
怎樣度過老化的過程,.
不過聖經裡面有很多例子,.
我們熟悉的人物,.
例如摩西,亞伯拉罕,迦勒,.
甚至在新約的保羅,.
使徒彼得,約翰,.
他們在晚年的時候,.
仍然面對很多困難,.
但他們沒有自棄,.
沒有躺平,.
堅持對神的信靠,.
他們憑著信心,.
帶著盼望,.
繼續走他們餘下的人生,.
所以一個不自棄的人,.
對神絕對忠心到底,.
向神不斷懇求,.
懇切的呼求,.
神怎會離棄他們呢?.
有一位牧者,.
神學家,靈修學家,.
他的名字叫侯士廷,.
James Houston,.
他對於牧養關懷長者,.
有很豐富的經驗,.
他確立了很多研究,.
他說長者在晚年生活中,.
如果是參加教會活動,.
得出來的,.
他們有較長的壽命,.
如果剛才聽過長者的獻詩,.
你就會知道,.
當中有90歲的長者獻祭,.
另外他們亦較少患疾病,.
或殘障的情況,.

$^{561}$反而他們的認知免疫功能,.
還會有向好的改變,.
更能對他們的生命,.
起了正面積極的意義,.
所以在這裡,.
我也想推介旺角浸信會,.
長者部的一些活動,.
除了平常的聚會,.
如何認識聖經,.
如何學習禱告,.
我們還有很多其他活動,.
大家知不知道,.
我們七月份開始有五十家,.
一提到五十家,.
有位弟兄跟我說,.
五十家,五十都算長者嗎?.
不知道大家有何看法,.
有沒有聽過,.
「輕老,中老,老老」,.
「Young old, old old old」,.
有沒有聽過?.
你會知道,.
現在人在很多研究中,.
去到一定的年齡,.
其實仍然在年齡階段中,.
會這樣被界定,.
在平常的活動,.
我們一至五都有不同的活動,.
剛剛我收到一個很新的消息,.
就是昨天行山,.
昨天是一個進森的行山活動,.
所以進森只有二十多人去,.
但當中我收到的消息,.
也不是少得多,.
原來有兩個高齡的參加者,.
一位八十歲,.
一位八十九歲,.
他們是行三點五星難度,.
他們行不全程,.
當時之前我收到消息,.

$^{601}$有人去行山,.
在同工會分享,.
牧師說你真的勸勸他不要去,.
很危險的,.
我自己也想想,.
我們當中有很多行山經驗的人,.
有熊爸爸在,.
有田先生在,.
今早完了早課,.
這位先生也有回來,.
我也跟他打招呼鼓勵他,.
很想他繼續回來教會聚會,.
這個教會的活動,.
能夠幫助一班年紀長的人,.
他們參與活動的時候,.
增加人與人之間的接觸,.
有很多的互動,.
可以在分享的裡面,.
在笑聲的裡面,.
大家可以開懷,.
心靈裡面的空間,.
對他們來說,.
一定是擴闊了很多,.
所以當眼界開闊,.
我們就能夠以新的目光,.
看我們的人生,.
詩人更加肯定,.
我卻要常常盼望,.
並要越發讚美,.
因為他一直持守對神的盼望,.
這一股動力驅使他去讚美神,.
因為神對他的眷顧,.
和施行慈愛,.
他就更加禱告,.
更加祈求,.
更加渴望終日的去讚美神,.
基於盼望是源自神,.
並且神的拯救,.
所以詩人能夠以樂觀,.
正向的人生觀,.

$^{641}$來過他每一天的生活,.
他這個生活,.
雖然我們看到祈禱裡面的詩篇,.
好像很艱難,.
我相信,.
他會以一個開開心心的情懷,.
過他每一天,.
因為他知道他的拯救,.
是從哪裡而來,.
也知道他有什麼地方,.
是他生命裡面的保障,.
所以在長者團契裡面,.
其實每一個長者回來聚會,.
他們不是星期日回來,.
是星期三或星期六,.
他們帶著一個很開懷,.
很舒暢的心,.
出席聚會,.
就算他們有一點點病痛,.
他們都會出席聚會,.
如果真的病得不能走路,.
他們就會通知,.
我今天不能回來,.
在他們聚會裡面,.
他們常常禱告,.
我聽到他們的禱告,.
最多的禱詞是什麼呢?.
他們說,.
求主讓我們每天都開開心心過生活,.
有一對夫婦,.
他們真的很開心,.
平時回來真的很開心,.
他們有生命氣息,.
我去探訪他們的時候,.
我就問他們,.
為什麼他們經常都這麼開心?.
可能多點時間跟他們聊天,.
他們說,.
我一起來就已經先背今句,.
什麼都不做了,.

$^{681}$先在這裡讀聖經,.
於是他們在房間裡拿出一些很長的,.
寫得很漂亮的字出來讀給我聽,.
原來他們的秘訣就是要常常感恩,.
常常讚美,.
常常讀神的聖言,.
這對夫婦其實不是住在我們旺角附近,.
是住在天水圍,.
所以他們如果要有聚會,.
那星期就要來教會三次,.
這就是上帝給他們的恩典,.
這種積極活在當下的人生,.
弟兄姊妹,.
可能大家還沒有去到這個階段,.
你會不會也分享你生命裡的一些故事,.
傳承下去呢?.
詩人就是這樣禱告,.
神啊,自我年幼時,.
你就教訓我,.
直到如今,.
我傳揚你奇妙的作為,.
神啊,我到年老發白的時候,.
求你不要離棄我,.
讓我將你的能力指示下代,.
將你的大能指示後世的人..
詩人所表達的就是.
將他的生命歷練,.
包括神的拯救,.
神的公義,.
神的公義在整篇詩篇71篇出現了很多次,.
神的公義帶有審判,.
而在審判裡面,.
上帝也拯救詩人,.
我相信如果是大衛所寫的71篇詩篇,.
神對大衛的公義,.
大家一定會記得很清楚..
這些通通都成為活生生的教材,.
是他的樣板,.
要指示神的能力,.
指出神的作為,.

$^{721}$在他的身上怎樣彰顯一個奇妙,.
一個人們覺得為怪的,.
他要宣揚出去..
經文指出,.
我們要將我們自己.
與神互動的生命,.
分享傳承,.
能夠將基督徒生活的準則分享,.
相信是一份恩典,.
也是值得感恩和讚美..
在《長者團契》裡面,.
一位身為奶奶的姐妹,.
她不時也會分享,.
她有兩個孫女,.
小時候就說帶她認識耶穌,.
她在家裡會以一些信仰,.
讀聖經,講故事,祈禱,.
來分途他們年紀小小的生命,.
甚至講關於神的公義,.
來處事為人..
她現在的孫女已經在讀大學,.
但又繼續聽她分享,.
在大學裡面除了要學有關的學科之外,.
都會有很多不同的價值,.
人生觀來衝擊她,.
究竟她和她的信仰,.
會是怎樣,.
可不可以銜接到呢?.
這位奶奶繼續鼓勵她,.
你不如回來教會的年青團契,.
可以認識多一點的信仰,.
令信仰的根基更加紮實..
可見長者在信仰上的指引,.
分途是非常重要,.
提摩太也是一樣,.
有她的外祖母,.
也有她的媽媽,.
在提摩太小的時候,.
從小就教她認識聖經,.
以致當保羅揀選找到她做拍檔的時候,.

$^{761}$她在一個中心的宣教裡面有份,.
甚至保羅也猜派她來牧羊教會,.
教養孩童使她走當行的道,.
就是道路也不偏離..
另外,我們過去長者的團契,.
我們也很想讓年青人在當中參與,.
所謂cross over,.
用不同的世代去幫助,.
可以說是照顧一下年長者,.
但原來也不只是這樣,.
可以讓年長者信仰那種生命,.
信仰那種依然有的活力,.
和那種信心,.
可以承傳給年青的一代,.
讓他們在年青的信仰裡面,.
面對現在很複雜,彎曲薄謀的世代裡面,.
讓他們更加知道,.
如何在信仰上有更好的紮根..
詩篇71篇是一篇哀討詩,.
內容非常生活化,.
詩人的禱告其實提醒我們,.
無論是什麼環境,.
我們的生活都總是要過,.
不過我們過生活的時候,.
是要定精在神上面,.
在神的身上我們就能夠更有信心,.
更有盼望去投靠祂,.
和要常常的去讚美祂,.
因為祂在我們身上做了很多奇妙的工作,.
所以我們的人生,.
我們要帶著神的榮耀作為,.
去見證祂,成為別人的榜樣,.
教養下一代..
所以年輕人,.
如果你覺得你是年輕的時候,.
你就要在你壯年的日子,.
要紀念做你的主,.
我們年歲的日子,.
要在神的恩典裡面,.
我們要開開心心,.

$^{801}$要以一個喜樂的心,.
每天來活出我們的信仰,.
這個就是恩典的記號..
有一節的經文,.
我相信大家都很熟悉,.
在白髮人的面前,.
你要站起來,.
也要尊敬老人,.
又要敬畏你的神,.
我是耶和華..
可見尊敬老人家,.
其實就等於去敬畏上帝,.
去敬畏神..
最後,很想以一段禱告來結束,.
這一段禱告,.
就是一位傳道人,.
他去另一間教會,.
探望另一間教會的傳道人,.
他在禮堂的公告欄,.
發現了一篇祈禱文,.
這篇祈禱文,.
是年長者寫的祈禱文,.
他是這樣說的,.
「主啊,你知道我年歲漸長,.
讓我不要無視這個事實,.
讓我不要變得令人討厭,.
也不要自命不凡,.
在任何場合裡,.
要發表意見令人討厭,.
同時,也讓我不要.
只是鉅細無遺地.
說自己過去的經驗,.
封起我的嘴巴,.
不要哭這裡痛,那裡酸,.
我不求更好的記憶力,.
但求少管別人的閒事,.
求,我就是要承認有時候被誤解,.
求你使我隨著年歲增長,.
變得較為甜美成熟,.
讓我永遠不變老,而是成長,.

$^{841}$但願聖經的說話提醒我,.
我們外體雖然毀壞,.
內心卻一天生持一天,.
求讓我的內心外表都保持甜蜜,.
求讓我不要為生命的盡情煩躁心焦,.
求讓我盡情享受那永生的歡樂,.
求讓我在變老時也能心宜歡笑,.
求神使年歲漸長的生命更加豐盛..
\newpage



\section{提摩太後書 1:1-7}
\label{sec:4hDfY5nlUxk}
\textbf{2022年12月4日主餐暨堂慶主日崇拜直播｜梁家麟牧師｜遺產.使命.恩典｜提摩太後書1\_1-7}
\newline
\newline
連結: \href{https://youtube.com/watch?v=4hDfY5nlUxk}{\texttt{https://youtube.com/watch?v=4hDfY5nlUxk}} ~~~~ 語音日期: 2022-12-08
\newline
\newline
\hyperref[sec:esPM0k_dZBY]{\small{< < < PREV SERMON < < <}}
~
\hyperref[sec:index_chronic]{\small{[返順時目]}}
~
\hyperref[sec:6BkHUdPWvFs]{\small{> > > NEXT SERMON > > >}}
\newline
\newline
提摩太後書 1:1-7
\newline
\begin{longtable}{cl}
\hline
\hline
章節 & 經文 (和合本修訂版)\\
\hline
1:1 & \begin{tabularx}{0.7\textwidth}{X} 奉神旨意，按照基督耶穌裡所應許的生命，作基督耶穌使徒的保羅， \end{tabularx} \\ \\ \relax
1:2 & \begin{tabularx}{0.7\textwidth}{X} 寫信給我親愛的兒子提摩太。願恩惠、憐憫、平安從父神和我們的主基督耶穌歸給你！ \end{tabularx} \\ \\ \relax
1:3 & \begin{tabularx}{0.7\textwidth}{X} 我感謝神，就是我接續祖先用純潔的良心所事奉的神，在祈禱中晝夜不停地想念你。 \end{tabularx} \\ \\ \relax
1:4 & \begin{tabularx}{0.7\textwidth}{X} 我一想起你的眼淚，就急切想見你，好讓我滿心快樂。 \end{tabularx} \\ \\ \relax
1:5 & \begin{tabularx}{0.7\textwidth}{X} 我記得你無偽的信心，這信心先存在你外祖母羅以和你母親友妮基的心裡，我深信也存在你的心裡。 \end{tabularx} \\ \\ \relax
1:6 & \begin{tabularx}{0.7\textwidth}{X} 為這緣故，我提醒你要把神藉著我按手所給你的恩賜再如火挑旺起來。 \end{tabularx} \\ \\ \relax
1:7 & \begin{tabularx}{0.7\textwidth}{X} 因為神賜給我們的不是膽怯的心，而是剛強、仁愛、自制的心。 \end{tabularx} \\ \\
[1ex]
\hline
\hline
\end{longtable}
$^{1}$今天讀經是在提摩太后書一章一至七節.
是在新約聖經二百九十九頁.
提摩太后書一章一至七節.
新約聖經二百九十九頁.
請聽我讀速.
「奉神旨意照著在基督耶穌裡生命的應許.
作基督耶穌仕途的保羅.
寫信給我親愛的兒子提摩太.
願因慰憐憫平安.
從父神和我們主基督耶穌歸於你.
就是我接續祖先用清潔的良心所事奉的神.
祈禱的時候不住的想念你.
紀念你的眼淚.
晝夜徹徹的想要見你.
好叫我滿心快樂.
想到你心裡無謂之信.
這信是先在你外祖母羅爾和你母親有離基心裡的.
我深信也在你的心裡.
為此我提醒你.
使你將神藉我按手所給你的恩賜.
再如火挑旺起來.
因為神賜給我們.
不是膽怯的心.
乃是剛強仁愛謹守的心.
你不要以給我們的主作見證為此.
也不要以我這位主為酬的為此.
總要按神的能力與我為福音同壽.
我們讀到「謹守的心」.
好的.
今天是黃國俊信會六十週年堂慶.
我們懷著感恩的心.
迎接我們所有人都有份的大壽.
生日快樂.
德國十八世紀大哲學家康德.
以一生智力問一個問題.
就是「人是甚麼」.
人是甚麼來的呢?.
他將這個大哉問.
一個如此重大的問題.
分為三個子問題.

$^{41}$三個世的問題.
我能夠知道甚麼.
我要做甚麼.
我可以盼望甚麼.
這個分別就是.
知識,倫理和宗教三個領域的問題.
香港二十一世紀的一個小傳道梁家倫曾經說過.
基督信仰關乎上帝與人的關係.
我們主要問四個問題.
上帝是誰?.
上帝為我們做了甚麼?.
我們是誰?.
我們要為上帝做甚麼?.
上帝是誰?.
和祂為我們做了甚麼?.
這是關乎科興信息.
而我們是誰?.
和我們要做甚麼?.
這關乎我們的身份和使命.
福音為我們塑造新的身份.
賦予新的使命.
先確定福音.
然後確定我們.
我們是誰?.
我們要做甚麼?.
身份和使命是我們生命息息相關的兩大元素.
今天我們借用這段短短的聖經.
來做一些思考.
我們所選讀的經文.
記載保羅和他年輕的同工提摩泰的眷戀.
如所周知.
提摩泰是保羅在晚年一個傳道的助手.
保羅猜排他到伊弗所.
處理當地教會嚴重的人事問題.
這間教會問題非常扭曲.
提摩泰似乎不一定很年輕.
三十多歲.
不過總之不是很有能力,魄力,權威去處理.
所以處理得很淪盡.
所以教會問題一直沒有解決的跡象.

$^{81}$保羅寫信給他.
一方面給他一些具體的指導.
另一方面也為他打氣.
希望鞏固他的信心.
在崗位上堅持下去.
保羅沒有給提摩泰一些空洞的安慰.
明知他做不來.
只說「我對你有信心」.
「你可以的,你行的」.
他不是這樣說.
保羅是就提摩泰已經有的東西說起.
已經確定的說起.
那些「Given」說起.
強調提摩泰是上帝所揀選的工人.
「我是誰?」.
「我要為上帝做些什麼?」.
身份和使命息息相關.
保羅幫助提摩泰確認他的身份.
然後才引伸鞏固他的使命.
正如前面所說.
教徒的身份是由福音所塑造的.
上帝是誰?.
上帝為我們做些什麼?.
這兩點是影響深遠的.
保羅要提摩泰確認一件事.
他是上帝所體重的人.
並且上帝不是隨便去花墟選一束花.
「這束吧」.
不是的.
上帝早在提摩泰生命裡有長久的作為.
是的.
今天我們說要為上帝承擔這個工作.
但是我們要確認.
上帝在我們生命裡已經長久地有工作.
他的工作早於我們的工作.
他的工作也遠遠大於我們的工作.
嚴格地說.
我們說我們為上帝工作.
不過是在某方面配合他在我們身上的工作意義.
我們的工作也是上帝在這裡工作.

$^{121}$可能你覺得我太極端了.
舉個例子.
我在見道神學院事務院工作了三十七年.
從某個角度來說.
我是在見道神學院工作的.
但是骨子裡.
是上帝借用見道這個平台.
繼續完成他在我生命裡的工作.
上帝幫我完成我整個救贖的生命.
是不是?.
就好像你們在旺鎮服侍.
你們同時在這裡經歷他的恩典.
同時在這裡成長.
同時在這裡見證上帝奇妙的作為.
是不是?.
所以你們在這裡事奉.
但你們也在這裡經歷上帝的工作.
上帝在你生命裡的工作.
他要繼續模造你.
雕琢你成為他眼中的美好.
弟兄姊妹.
我們都是上帝的創造.
但上帝不是一下子就完成我們這件產品.
我們都是發展中國家.
都是在繼續發展中.
Developing.
But not developed.
不是已經發展了.
是發展中.
也就是說.
是繼續讓上帝的創造和救贖計劃.
在我們生命裡完成.
我們今天的生活和事奉.
都是逐步展開.
將上帝的旨意實踐出來.
彰顯出來.
工作尚未完成.
使命尚未完成.
我們的生命也尚未完成.
上帝在我們生命裡有很多工作.

$^{161}$上帝如何在我們生命裡開始工作呢?.
他有很多不同的方法.
藉著不同的際遇.
不同的人和事的發生.
在我們生命裡各種的安排.
其中最不需要多說的.
是我們的家庭和父母.
我今天的題目是.
遺產,使命和恩典.
我們先說遺產這部分.
家庭對我們的影響巨大.
不僅僅是我後天出生成長的環境.
我是在橫頭磡出生.
我在慈雲山紅番區出生.
那個環境對我當然有影響.
更加重要的是.
我們從父母和家族.
領受各種遺傳基因.
那些都是遺產.
家庭對我們的影響.
很多研究都告訴我們.
遺傳對人的影響很大.
我們在生理,心理各方面.
都受遺傳基因影響.
我比較胖.
極有可能是遺傳道致的.
「你擦汽器吧,你餵食而已」.
餵食也是一個遺傳基因.
所以我是不能不餵食的.
基因決定了我餵食.
我們的智力,我們的性格.
其實都很受遺傳影響.
父母很清楚.
同一對父母生兩個子女.
性格不一定一樣.
有些人不是後天的.
先天決定了.
所以小朋友的成長.
很大程度上真的受先天遺傳所影響.
後天教養能夠影響的只是短時間.

$^{201}$為什麼這麼多輔導員.
這麼喜歡將我們的問題歸咎.
我們的原生家庭.
我們FOO.
Family of Origin.
說原生家庭.
通常給人一個性賴的印象.
通常為我們一些不好的事情.
做一個歷史的解釋.
指出我們的問題都是父母的錯.
祖宗的錯.
藉此去推卸我們的責任.
我卑鄙,因為我爸爸卑鄙.
我軟弱,因為我媽媽軟弱.
諸如此類.
今天我們說正面的話.
我們所有的優點.
所有的長處.
如果有的話.
豈不是都是從我們的父母.
我們的家庭遺傳得來的嗎?.
弟兄姊妹.
我很相信你們在香港.
都算是Above Average.
比較好的一班人.
基督徒的男人多數都是好男人.
最少是負責任顧家.
照顧妻兒.
如果我們有任何的好處.
我相信都和你的父母.
在你過去成長的階段.
對你悉心的教導,提攜.
有很重要的關係.
我媽媽幾個月前逝世.
如果我有任何的好處.
我都相信是和我媽媽有關.
是她幫我成為今天這樣的人.
給我們後天的教養.
所預先做的工作.
在剛才我們所讀的經文裡面.

$^{241}$保羅提醒提摩泰.
他今天所持定的基督信仰.
不是由他個人開始的.
而是由他在某處不小心撿回來的.
我們很多人是第一代的基督徒.
我們確實是不小心撿回這個信仰.
我在街上撿這個信仰回來的.
我去過報道會聽了耶穌.
不過提摩泰不是.
他的信仰是出自他的母親和外祖母.
第五節也說到.
「想到你心裡無偽的信」.
「才是在你外祖母羅爾」.
「和你母親有離基心裡的」.
「我深信也在你心裡」.
這是一個延續了三代的信仰.
雖然也有可能.
他母親和他外祖母是同時間信主的.
不過如果說世代來說.
已經是三代了.
三代的信仰.
我不知道外在當中有多少人.
是出生在基督徒家庭.
有多少是第一代基督徒.
出生在基督徒家庭的.
你們屬於第幾代呢?.
我們中間有沒有超過三代的?.
即是數上兩代.
有沒有?.
即是你爺爺,阿麻.
或者阿公,阿婆信主.
有多少個?.
一個?兩個?.
也不是吧?也有五,六個.
出生在一個基督徒的家庭.
這意味著信仰的獲得.
是不需要經過一個尋索和轉向的歷程.
生命不需要出現一個斷層.
當然,仍然不是說.
我們信仰可以純粹由父母繼承過來.

$^{281}$不用個人的抉擇.
沒有這樣的事.
嚴格地說.
每個人都要自己獨自面對上帝.
都要做跟隨基督的個人決定.
個人的決志.
沒有自動轉向這件事.
但是,在基督徒家庭成長.
從小回教會.
接觸信仰的機會就是百分百.
不用在街頭等人向你派單張.
才能接觸到信仰.
在家裡,從小張開眼睛.
或是嬰兒時就已經手抱回教會.
你們的決志跟隨主.
要付的社會代價.
那個社會代價也少很多很多.
保羅提到提摩太有這樣的家庭優勢.
這正好跟他自己出生在一個嚴謹的猶太教.
法利賽派的家庭構成極其強烈的對比.
保羅由迫害基督徒.
到成為基督徒.
期間真的出現一個嚴重的斷層.
他不僅要跟原來他的家庭.
跟社區斷絕關係.
負上沉重的社會代價.
如果他是一個法利賽派的人.
他不可能過了二十歲還未結婚.
但你在經文十三封書信.
有沒有提過保羅有老婆?.
有沒有提過他老婆?.
完全沒有.
他還酸溜溜地說.
給你就好了.
可以帶老婆到處去.
我就自己一個孤家假人.
可能他老婆早死.
不過更可能的是.
他因為信耶穌的緣故.
被驅趕出整個法利賽派的家族.

$^{321}$所以他老婆也跟他分開了.
這就是為什麼他在教導的時候說.
你做了基督徒.
如果你的配偶未信.
他不想跟你分開.
就不要跟他分開.
但如果他堅持要跟你分開.
你也要由他走.
信仰先行.
我相信保羅說這些的教導.
是跟他個人處境或密切的關係.
所以保羅極有可能是因為.
信了耶穌的緣故.
跟他整個家族.
整個成長的社群徹底的割裂.
並且不止這樣.
他棄暗投明.
棄暗當然辛苦.
投明也不自在.
他最初要加入基督徒群體.
也被新的群體成員排斥他.
他們懷疑保羅是否真誠信主.
是否奸的.
這是否無奸道.
滲入教會顛覆我們.
如果不是因為有巴拿巴的信任和幫助.
保羅要取信於教會.
特別是以耶路撒冷為主的教會.
是極其不容易.
不過保羅始終在他一生.
都無法融入耶路撒冷的信徒群體.
他只能夠另起爐灶.
落直安提阿.
並且專注向外邦人傳福音.
你做這個演奏一看.
提摩太真是自在.
在基督教家庭長大.
信仰沒有斷層.
並且很早就被接納.
保羅是碰到人們懷疑他不接納他.

$^{361}$提摩太不碰都被人硬捧上去.
差多遠.
真是好著數.
所以保羅做這個對照.
告訴提摩太.
你真是很幸福.
我們知道教會最基本的單位是家庭.
不是個人.
家庭是傳福音的第一戰線.
一個人信了耶穌.
就會將家人和朋友一個個帶回基督.
除了傳福音之外.
家庭也是栽培和訓練的場所.
我以前也說過.
宗教教育主要在家做.
不是在教會做.
不是教會主要學不重要.
兒童時空不重要.
但我希望各位家長.
各位父母一定要認定.
宗教教育是在家做的.
這是很符合浸信會的主張.
浸信會是強調家庭教育.
非常強調家庭教育.
我的子女的信仰.
是由我和我太太建立的.
我不會將責任推給教會的主教老師.
傳福音栽培在家做.
家庭也是實踐信仰的主要場所.
保祿在《伊弗所書》四到五章.
提到的三種主要關係.
夫妻,父母,子女和主僕.
都是家庭的關係.
我們要在家努力成為基督徒.
實踐基督徒的信仰.
兌現基督徒對我們要求的責任.
從傳福音到栽培到信仰實踐.
子女如果在健康的基督教家庭長大.
就能夠明白和經驗信仰.
奠定一個鞏固的根基.

$^{401}$父母對子女有重要的宗教教育責任.
弟兄姊妹又說了一些題外話.
我們為子女其實已經付出一切.
轉校區搬家,搞移民,傾家蕩產.
讓子女去外國讀書.
如果不是非常富有的家庭.
供一個子女去英國,美國讀書不便宜.
現在就不可以了.
以前隨時供她讀完大學學位就有一層樓.
現在就半層了.
很貴很貴.
不過我們要認定.
努力為子女建立一個穩固的信仰基礎.
讓他們從小學習聖經.
投入教會生活.
並且學習用信仰的角度去理解.
去評估這個世界.
不管子女學問有多高.
成就有多大.
財富有多豐厚.
如果他們放棄了信仰.
我們都有很大的遺憾.
無論如何我們為子女創造財富.
給他們最多的遺產.
但我們也要將信仰傳遞給他們.
使他們不僅成為富二代,富三代.
更要成為信二代,信三代.
家庭要建立一個所謂的遺產.
正面說,不要說負面的事.
我們中間有很多信二代,信三代.
如果你們有一個比較鞏固的信仰.
我們常常說負面的事.
但今天很多信二代,信三代離開教會.
這是令人難過的現象.
但我們同時要看到.
有不少信二代,信三代.
仍然是積極參與教會的服事.
父母對他們的影響是真真實實.
我記得有一年.
我去一間很古老的教會.

$^{441}$名叫禮賢會香港堂.
其中一堂的主席是姓黃的長老.
他告訴我.
他祖父是這間教會的牧師.
當他說到他祖父的時候.
這位長老已經七十多歲.
他的表情既莊重又崇敬.
家庭的屬靈遺產不是一個抽象的觀念.
絕對不是抽象的觀念.
如果教會是我們屬靈的家庭.
教會也有屬靈的遺產.
今天我們來參加望朕聚會.
我們知道這間教會.
是由很多前輩努力奮鬥打拼出來的.
甚至是直接由他們傳福音和栽培而成長的.
教會的遺產除了我們的建築物.
於是我們有這幾層的房子.
更重要的遺產是無形的.
屬靈的傳統.
我們說望朕的精神.
望朕的DNA.
是什麼構成我們和其他教會不同的地方.
今天六十週年.
我們一定要問這個問題.
望朕的DNA是什麼.
什麼是我們.
構成我們非是這樣不可.
不是這樣就不是望朕的東西.
有什麼是我們引以為傲的東西.
從遺產說到使命.
家庭是教會的基本單位.
家庭的傳承影響教會的傳承.
信仰必須一代一代的傳遞下去.
教會的使命也得由一代再傳遞去.
另一代代代相傳.
一間擁有六十年歷史的教會.
總是不停經歷世代交替的故事.
我們的服侍是接觸前人的服侍.
不是自己獨力開創的.
保羅以他的經驗在第三節說.

$^{481}$就算我接觸祖先.
用清潔的良心所事奉的上帝.
你要注意.
這個說的祖先一定不是他自己的血緣祖先.
他第一代基督徒有什麼祖先.
是不是.
不關他的事.
不是他的祖先的事.
他還跟家族割裂了.
有什麼祖先.
這裡說的祖先不是他血緣家庭的祖先.
是說他的肅靈家庭的祖先.
那些先賢先成.
就好像對望針來說.
我們的祖先是誰.
是子強.
是馬美.
是陳太.
還有很多我不認識的前輩.
他們是我們肅靈的祖先.
我經常說.
我們在座的很多都是我很尊重.
我相信將來我們會被傳頌.
不過對仍然在生的人我們不會說話.
只是說已經被主接去.
我們認定他們是我們的祖先.
他們是我們的祖先.
這個繼承的治國.
延續到保羅的下一代.
他在第二節稱提摩泰為他的兒子.
說明這是一個肅靈的血緣關係.
世代相傳.
提摩泰要接續保羅的服侍.
薪火相傳.
我知道古代的社會資源不足夠.
社會流動也不高.
子女通常是從父母繼承職業.
創業的空間很狹窄.
父母耕田 兒子耕田.
父母打魚 兒子打魚.

$^{521}$父母是木匠.
做兒子的通常都是木匠.
到今天子承父業.
這不再是香港的必然.
但你說完全沒有也不是.
特別是做生意的.
父母很想子女接他的生意.
將他發揚光大.
還有很多專業也有類似情況.
例如做父母是醫生.
很多子女會做醫生或醫護界.
會計師的子女做會計師.
我兒子是做律師.
他告訴我 當然他不是出身律師世家.
他告訴我 讀法律的很多是法律世家.
特別你不是法律世家.
你想讀法律的很難.
基本上很難.
你哪有生意 怎樣開始 怎樣起步.
很多都是法律世家.
西叔伯照著 很多很多.
香港這些老字號近年紛紛結業.
做舖租貴 成本高.
這是其中一個因素.
更重要的是年輕一代出路比較多.
不喜歡從事飲食業.
需要長時間工作的行業.
所以當店主老去 店舖後繼無人.
台灣和日本的情況很不同.
他們的社會晉升沒有我們那麼多.
沒有那麼流動.
當然他們沒有我們那種萬般佳下品.
唯有讀書高的觀念.
所以很多古老店舖甚至街頭小攤販.
都是幾代經營.
並且他們強調 他們遵從古法製造.
台灣叫古早味.
總之不變的 幾十年不變 幾代不變.
初期家會主要在家裡聚集.
家會因此像一個家族企業.

$^{561}$由父傳子一代一代繼承下去.
幾個核心家庭拱衛了教會的發展血脈.
今日香港很少這樣的情況.
除了一些方言教會 方言地域教會.
例如汶南教會 潮州教會.
這些比較強調方言和地域.
家族的影響比較大.
其他教會這方面不多.
不過教會主要仍然由家庭組成.
我記得有一個友會跟我說.
他們有二三千人.
不過三個家族都可以串到50\%的人.
兒女交叉結婚 很多的.
交叉幾代就很大群的.
家族是有很大影響力的.
信仰由家庭一代一代傳遞下去.
猶太人比我們更加注重這個.
為何他們亡國千多年.
民族和信仰仍然得以鞏固.
基督教會一定要重視宗教教育.
亦要注重將我們的意象和使命傳遞給下一代.
並且盡量維持組織的開放態度.
讓下一代願意和有能力去接行.
一間教會如果做來做去都過幾十年.
幾十年不變 掌權的 做掌櫃的不肯放權.
其實不是很健康的情況.
說一些不應該說的例子.
昨天看報紙 公民黨因為沒有人肯選任黨主席.
宣佈進入公司結束程序 要拉閘.
這是今日香港特殊的政治生態環境使然.
不理我們的政治立場是甚麼.
聽到這些消息我都有些黯然的感覺.
今日教會老化 缺乏年輕人.
這是一個嚴重的問題.
如何栽培接班人更是教會最嚴峻的挑戰.
香港一千四百間教會缺主任牧師的.
我估計超過三百間.
很多教會掌執都有空缺.
原因不用多說 不分析這些.
我們只能夠咬緊牙關.

$^{601}$做栽培 輔助年輕人上位的工作.
不要太多說 今日的年輕人不夠理想.
提摩太太的才華遠不及保羅.
但保羅仍然硬食.
甚麼都自己做 你做了幾年.
無論如何都要扶提摩太太上位.
將責任托付給他們 努力幫他們完成任務.
千萬不要我們這一代做得好.
但我們下台後 後繼無人.
教會亦要解散 這是很糟糕的事.
我們不想人亡政息.
希望教會能夠繼續未竟的使命.
一直屹立發展到主在內的日子.
沒有年輕人 沒有接班人.
教會就沒有明天.
保羅自覺從祖先領受義將和使命.
亦銳意訓練像提摩太太般的接班人.
這是成仙啟後的精神.
當然成仙啟後不等於我們要求下一代重複我們所做的所有事.
上帝不需要我們做這樣的重複.
上帝是獨行其事的上帝.
我多提幾句.
希伯來書的作者.
我覺得很有趣.
他一方面確認歷代的前輩.
在信仰和使命上樹立了光輝的榜樣.
希伯來書十一章二節說.
古人在這信上得到美好的證據.
但他同時要求鼓勵信徒不要往後看.
而是往前看.
不是搬前人的做法.
而是仰望在我們前頭帶領我們的耶穌基督.
我們既有許多見證人如同雲彩圍著我們.
就當放下各樣的重擔.
脫去容易纏累我們的罪.
全心忍耐.
不那擺在我們前面的路程.
仰望為我們信心創始成終的基督.
基督才是旺朕的上訴.
阿叔建立的傳統我們要尊重.

$^{641}$但阿叔不是教會的權威.
我們的權威只有一個.
剛才師班大師跟我們說了.
耶穌基督才是永遠的創始成終的那一位.
對父母來說.
子女有些地方跟自己相似.
這是很美妙的事.
不過我們每個人都想子女比我們好.
比我們厲害.
成就比我們大.
跟我們一樣有什麼好.
要厲害一點.
從遺產到使命將我們說到恩典.
我們的過去是由上帝的恩典.
和身邊的人的恩情構成.
在第六節.
保羅圈面提摩太將恩賜如仿挑旺起來.
這裡說恩賜不是提摩太個人的某些特殊的才能.
我們常常說恩賜是說這件事.
你有什麼恩賜.
我口才好.
我的社會關係好.
我的什麼什麼好.
我們說恩賜是說我們個人的能力.
我會彈琴.
我八級樂理.
這是我的恩賜.
保羅說恩賜不是這樣說的.
恩賜是說你有的東西.
但你有的東西不是純粹由你的才能決定.
保羅說的是什麼.
保羅告訴提摩太.
他的恩賜是指著他所有的經歷和關係.
例如他有一個虔誠的外祖母.
有一個虔誠的媽媽.
他是在一個信仰群體成長.
他被認識了保羅.
被保羅差派.
藉著一個按手禮成為傳道人.
所有這些經歷和關係.

$^{681}$都不是他賺回來的.
是上帝無形的手冥冥中安排給他的.
所以是注定的.
不是他賺回來的.
就是恩典.
就是恩賜.
是上帝給他的禮物.
想想我們的過去.
有很多很特別的經歷.
有很多很特別的遭遇.
你發覺有一個無形的手一直拖帶著我們.
那隻無形的手就是我們生命的大導演.
就是我們生命的編劇.
那些構成我今天最重要的元素.
這是恩賜.
在保羅眼中.
提摩太所有先天和後天的條件.
他的出身.
他的家庭背景.
通通都是恩賜.
今天他領受傳道的工作.
都是一個恩賜.
這個恩賜是上帝藉著保羅所做的按手禮而給他的.
保羅勸提摩太將恩賜如火挑旺起來.
意思就是你要維持奮鬥心.
保持積極.
保持戰意.
將奮鬥心如火挑旺.
心中要有火.
你才能夠承擔工作.
提摩太不是一個勇敢的人.
上上缺乏勇氣和毅力.
保羅沒有教導他建立一個正確的自我形象.
教導他減壓.
教導他如何處理當前的困難.
而是讓他看到一個比他眼前的困難更重要的事實.
你不是你的困難.
你不是你的問題.
你是上帝花幾十年時間鋪造出來的一位.
當你看見上帝的作為.

$^{721}$看見恩典大於你眼前的問題的時候.
你就可以醒定做人.
生命是一個恩典.
生命裡面有很多的恩賜.
恩典能夠使人活下去.
活得有意義.
活得積極.
唯有在生命裡面發現恩賜.
我們才可以得到生之勇氣.
The courage to be.
然後才能夠建立剛強.
仁愛和謹守這些美德.
我常常說.
生命最重要的東西.
都是不會失去的.
因為都不是我們賺回來的.
上帝的恩典.
身邊人的恩情.
這個構成了我最重要的本相.
我依次著這個.
我不需要不停為自己經營這樣那樣.
彷彿一旦我鬆懈了.
我就一文不值.
一無所有.
我輸不完.
永遠重要的東西都輸不到.
這是我們最大的福氣.
弟兄姊妹.
我們為我們自己的過去感恩.
為我們的教會感恩.
為我們的家庭感恩.
或者過去我們曾經有一些不開心的經歷.
不過將那些東西忘記.
認定能夠在這樣的環境成長.
這是我們的福氣.
你越為你的過去感恩.
你就越能夠看到今天你的位置.
也對前景有信心.
家庭和父母給我們很多信仰的遺產.
家庭和父母也都傳遞給我們信仰的使命.

$^{761}$家庭和父母也是我們經歷上帝恩典的渠道.
我們的父母.
我曾經不止一次對神學生和傳道人.
我和他們說話的時候說過.
我們在講道的時候.
說過很多不同人的故事.
但我們留在信徒心中的.
永遠是我們自己的故事.
不是我在網上剪輯下來的故事.
這番說話不僅適用於傳道人.
適用於所有弟兄姊妹.
今天我們會提到叔叔的故事.
他日下一代會提到我們的故事.
我們有各自的故事.
也有共同的故事.
看看我們能夠留下什麼故事.
告訴我們的下一代.
確認我們所領受的遺產.
確認我們繼承的使命.
發現生命裡上帝的恩典和身邊的人的恩情.
讓我們清楚明白我們的身份和使命.
於是我們就可以有.
保羅對提摩太的總結要求.
建立一個剛強,仁愛,謹守的心.
迎接未來的日子.
\newpage



\section{馬太福音 2:1-12}
\label{sec:6BkHUdPWvFs}
\textbf{2022年12月11日崇拜直播｜陳明泉牧師｜看見星星的你｜馬太福音2\_1-12}
\newline
\newline
連結: \href{https://youtube.com/watch?v=6BkHUdPWvFs}{\texttt{https://youtube.com/watch?v=6BkHUdPWvFs}} ~~~~ 語音日期: 2022-12-13
\newline
\newline
\hyperref[sec:4hDfY5nlUxk]{\small{< < < PREV SERMON < < <}}
~
\hyperref[sec:index_chronic]{\small{[返順時目]}}
~
\hyperref[sec:1SOWC0lwIU0]{\small{> > > NEXT SERMON > > >}}
\newline
\newline
馬太福音 2:1-12
\newline
\begin{longtable}{cl}
\hline
\hline
章節 & 經文 (和合本修訂版)\\
\hline
2:1 & \begin{tabularx}{0.7\textwidth}{X} 在希律作王的時候，耶穌生在猶太的伯利恆。有幾個博學之士從東方來到耶路撒冷，說： \end{tabularx} \\ \\ \relax
2:2 & \begin{tabularx}{0.7\textwidth}{X} 「那生下來作猶太人之王的在哪裡？我們在東方看見他的星，特來拜他。」 \end{tabularx} \\ \\ \relax
2:3 & \begin{tabularx}{0.7\textwidth}{X} 希律王聽見了，就心裡不安；耶路撒冷全城的人也都不安。 \end{tabularx} \\ \\ \relax
2:4 & \begin{tabularx}{0.7\textwidth}{X} 他就召集了祭司長和民間的文士，問他們：「基督該生在哪裡？」 \end{tabularx} \\ \\ \relax
2:5 & \begin{tabularx}{0.7\textwidth}{X} 他們說：「在猶太的伯利恆。因為有先知記著： \end{tabularx} \\ \\ \relax
2:6 & \begin{tabularx}{0.7\textwidth}{X} 『猶大地的伯利恆啊，你在猶大諸城中並不是最小的；因為將來有一位統治者要從你那裡出來，牧養我以色列民。』」 \end{tabularx} \\ \\ \relax
2:7 & \begin{tabularx}{0.7\textwidth}{X} 於是，希律暗地裡召了博學之士來，查問那星是甚麼時候出現的， \end{tabularx} \\ \\ \relax
2:8 & \begin{tabularx}{0.7\textwidth}{X} 就派他們往伯利恆去，說：「你們去仔細尋訪那小孩子，找到了就來報信，我也好去拜他。」 \end{tabularx} \\ \\ \relax
2:9 & \begin{tabularx}{0.7\textwidth}{X} 他們聽了王的話就去了。忽然，在東方所看到的那顆星在前面引領他們，一直行到小孩子所在地方的上方就停住了。 \end{tabularx} \\ \\ \relax
2:10 & \begin{tabularx}{0.7\textwidth}{X} 他們看見那星，就非常歡喜； \end{tabularx} \\ \\ \relax
2:11 & \begin{tabularx}{0.7\textwidth}{X} 進了房子，看見小孩子和他母親馬利亞，就俯伏拜那小孩子，揭開寶盒，拿出黃金、乳香、沒藥，作為禮物獻給他。 \end{tabularx} \\ \\ \relax
2:12 & \begin{tabularx}{0.7\textwidth}{X} 因為在夢中得到主的指示，不要回去見希律，他們就從別的路回自己的家鄉去了。 \end{tabularx} \\ \\
[1ex]
\hline
\hline
\end{longtable}
$^{1}$今日的經訓是馬太福音第二章一至十二節.
新約聖經第二頁.
請聽聖經的記載.
當希律王的時候.
耶穌生在猶太的伯尼行.
有幾個博士從東方來到耶路撒冷說.
那生下來作猶太人之王的在哪裡?.
我們在東方看見他的星,特來拜他.
希律王聽見了就心裡不安.
耶路撒冷合乘的人也都不安.
他就召齊了祭司長和民間的文士.
問他們說:基督當生在何處?.
他們回答說:在猶太的伯尼行.
因為有先知記載.
說:猶太地的伯尼行啊.
你在猶大諸城中並不是最少的.
因為將來有一位君王要從你那裡出來.
牧養我以色列民.
當下希律暗暗的召了博士來.
細問那星是什麼時候出現的?.
就猜他們往伯尼行去.
說:你們去仔細尋訪那小孩子.
尋到了就來報信.
我也好去拜他.
他們聽見王的話就去了.
在東方所看見的那星.
忽然在他們前頭行.
直行到小孩子的地方.
就在上頭停住了.
他們看見那星就大大的歡喜.
進了房子.
看見小孩子和他母親瑪利亞.
就俯伏拜那小孩子.
揭開寶盒.
拿黃金,乳香,抹藥為禮物.
獻給他.
博士因為在夢中被主指示.
不要回去見希律.
就從別的路回本地去了.
各位頂展會平安.

$^{41}$突然間螢幕看到自己這麼大.
真人也沒有這麼大.
嚇到.
我們來到新裝修的堂子.
我們由今個星期開始.
逐個星期來思考聖經.
好像倒數一樣.
來紀念基督耶穌降生.
今天我們一起來看《瑪達福音》.
這一段經文.
剛才梓梅為我們讀出的這段經文.
對我們來說也是耳熟能詳.
大家在每年的聖誕節.
可能讀過幾十年.
我們都覺得是一段很熟悉的經文.
我們唱詩歌都會唱東方三博士.
但是當你想深一層.
其實為什麼在聖經裡面.
會記載著一段這樣的經歷呢.
究竟東方這幾位博士.
我們今天也做.
其實他們記載在聖經裡面.
對整個基督耶穌降生的經歷.
或者對於我們今天來說.
有什麼重要的信息是我們知道的呢.
聖經裡面記載的說話.
我們不會純粹很負面地去做結論.
很多頂展會讀這段經文.
或者有些同道去處理這段經文.
很容易就會下了一個結論.
就是希律很壞.
是一個壞皇帝.
殘害當代的伯利恆的孩子.
我們很快就會下了這樣的結論.
如果純粹是為了一個結論.
去處理這段經文.
就好像錯失了這段經文裡面.
還有更加豐富的意涵.
希望今天我們從這段經文裡面.
我們嘗試去深入了解裡面的內容.

$^{81}$究竟來自東方的這幾位仁兄.
究竟對我們整個理解聖經裡面.
有什麼重點.
更加重要的是.
我們如何拿著這些重點.
來成為我們屬靈的養分.
求主幫助我們.
我們一起祈禱.
求聖靈開啟我們的心竅和心祈禱.
阿爸天父我們求你.
在我們當中繼續來到瑞嶺中頂展會.
讓我們去讀你話語的時候.
讓我們能夠認清你的心意.
開啟我們的心思念念.
讓我們能夠.
當我們讀到明白.
讓我們知道我們如何行事為人.
我們恭敬將今天早上.
我們一起領受這段話語.
我們全然交託給你.
將我們的心交託給你.
願你賜福.
獻上我們的禱告.
奉靠基督耶穌得勝的聖名祈求.
Amen.
我們一起看.
在經文裡面特別提到.
在希律王的時候.
耶穌生在猶太的伯尼行.
有幾個博士.
從東方來到耶路撒冷.
特別留意一個字眼.
就是在說幾個博士.
博士這個字就是和合本的翻譯.
在我們看回原文聖經裡面.
其實不是博士這個意思.
不是Doctor.
他用一個字眼.
英文聖經也會用這個字.
叫做Megui.

$^{121}$Megui這個字其實是來自當時的波斯話.
即是說我們要丟下一些概念.
不是那些人讀了一個博士學位.
或者認識了很多學問的人.
可能他們也有些學問.
不過在經文裡面重點不是說這些知識分子.
或者像我們今天所說的在大學裡面做教授的博士.
原文的Megui這個字是在說占星師.
如果你這樣看你會知道.
以前的占星.
類似我們今天看的紫微斗數.
朱國亮也做過這些東西.
我們中國人.
但是他們是研究天文.
藉著天文的變化.
來看到當時的政治,天氣.
或者他們國家的處境.
而他們這一班人.
他們藉著夜觀星象.
或者看天文現象.
來洞悉他們國家的命運.
如果用中國人的說法.
那班可能高級一點的叫國師.
低級一點的叫占卜家.
經文裡面特別只是說他們是一班Megui.
但是他們在他們自己身處的位置.
是站在一個什麼角色呢.
經文沒有特別去說.
所以我們不需要猜測.
不過在經文裡面再繼續說.
這幾個人是來自什麼地方呢.
來自東方.
我們看看地圖就知道了.
耶路撒冷的東邊是什麼地方呢.
舊時入侵的國家.
最近那個就是巴比倫.
再遠一點就是波斯.
其實巴比倫已經滅了.
很多的色經學家都相信.
這幾個來自東方的.

$^{161}$其實就是來自波斯的占星師.
你明白了吧.
你見到原來在經文裡面的字眼這麼微小.
但是他想告訴你.
這班人不是一班知識分子.
或者讀很多書.
和合本怕你因為有些偏差.
怕大家覺得占星師在聖經裡面這樣出.
好像不太屬靈.
所以和合本為了要在建立大家屬靈上的正面.
就用博士這個字.
但其實意思本來就是來自波斯.
來自巴比倫.
巴比倫最厲害占星術.
由這些人看到夜觀星象.
由東方的波斯來到耶路撒冷.
要尋找生下來的猶太王.
他們跟著這顆天上的星.
有些說法這顆叫做亞國之星或者大衛之星.
看到這顆星很奇怪地突然明亮了.
跟著他們跟著這個星光.
來到猶太的地方.
他們一找.
他們一想到君王.
第一樣要找的是什麼呢.
人口最沉默.
政治中心.
所以他們第一個到達.
或者他們主要的目的地去到哪裡呢.
幾個博士從東方來到耶路撒冷.
當時猶太人集中的地方.
聖殿在那裡敬拜.
在那個地方裡面.
最主要就是在耶路撒冷裡面.
他們覺得人口最沉默.
君王應該在這裡出生.
不過當他們去到這個地方的時候.
他們問了一條問題.
這條問題其實可以招惹殺身之禍.
無有識死.

$^{201}$他們問什麼問題呢.
那生下來作猶太人之王在哪裡.
我們在東方看見他的星.
特來拜他.
我們小朋友讀的時候.
我們覺得這句話很浪漫.
我們來找那個王的.
不過你留意一下.
當時的耶路撒冷.
其實猶太人已經沒有了自己的國家.
他們只是一個自治的地方.
他們的政權其實就是在羅馬帝國的管治之下.
因為福源很大.
所以讓他們不同的民族有自己的管治的權力.
當這班來自東方的占星師來到他們當中的時候.
他問一條生下來的猶太人的王在哪裡.
其實這個問題也是在猶太人心裡面不停想.
因為他們引頸以待.
等待一個政治的領袖.
從他們那裡出生.
希望有一招可以推翻他們不滿意.
或者一個好像由外邦管理他們自己國家的君王一樣.
他們很想自己有自己血脈的人.
來跟正苗紅地帶領他們自己的國家.
就好像之前的馬卡比時代.
他們想造反.
自己治理自己的地方.
所以這班人問一條很政治敏感的問題.
這班人其實這幾個占星師來到.
沒有有色死.
去到別人的地方問一條政治最敏感的題目.
接著他們說他們因為看到星星.
來到要拜他們.
值得留意的是.
剛才說的.
他找的是跟正苗紅的猶太人的王.
他來到耶路撒冷.
不過為什麼他們要問這個問題呢.
星光其實去到他們這裡的地方.
因為去錯地方.

$^{241}$所以星光就好像找來找去都找不到.
做不到他們心目中他們所做的目的.
當他們一直問的時候.
這些呼聲就去到當時管理著猶太人的王.
叫做希律王.
聽到這些人的問題.
他們的心就非常不安.
這個希律王說多一點.
待會我們會再仔細去看.
不過希律王跟當時生下來的猶太王有點不同.
希律這個王.
其實他本身的血統不是猶太人.
他只不過是他爸爸歸附了進去猶太人的血統.
稱自己為猶太人.
其實本身他是以東人.
不是那個鄉下.
但是說我現在認了.
認了他.
所以就成為了一個入籍的.
歸化了的一個猶太人.
而在他一直成長的過程裡.
因為對羅馬政權.
當時的羅馬帝國.
他救了他一些皇室成員.
所以後來那個羅馬政府就給他.
給這個地方成為他的食泣.
類似是這樣理解.
在一個這樣的狀態之下.
問一個根正苗紅的猶太人.
問一個非猶太裔.
不是根正苗紅的希律聽到這些呼聲.
他心裡一定有很多不安感.
他就開始做一些工作.
希望可以找到真正的猶太人的王.
來去鏟草除根.
我們稍微在這裡.
我們理解了這一部分之後.
我想有些地方值得我們思考.
剛才我所說的這幾個.
所謂博士其實是占星師.

$^{281}$幾個而已,沒有說三個.
無三不成幾.
所以我們就用三個.
不過千里迢迢.
就由波斯來到耶路撒冷.
或者來到猶太人聚居的地方.
聖經為什麼要記載一個這樣的故事呢.
如果你再認真一點.
你再理解一下.
其實在聖經裡面對於占卜聖上.
都採取一個很負面的態度.
聖經裡面從來都沒有叫人夜觀星象.
來尋求上帝的旨意.
聖經舊約的教導.
甚至在《但爾里書》裡說.
但爾里他們從小學到大.
但是他對於上帝的心意.
從來都不是在說.
在巴比倫裡面學的術數.
來幫助當時的王.
或者幫助他們自己解決眼前的困境.
他們只是尋求天上的符.
在聖經裡面即使.
但爾里這些領袖人才熟讀了這件事.
但是他們是不用的.
甚至在聖經裡面說到.
有一個很負面的對於占星術.
這種說法,這些是外邦人才用的.
為什麼會記載在耶穌出生.
還要引導外邦人來敬拜神呢.
近代的神學家幫我們解決了這個疑難.
在聖經裡面.
雖然反對占星術.
但是對於一些外邦人.
他們沒有聖經來看.
對於一些不是猶太人.
或者不是以色列人.
他們這班人不是上帝放棄的對象.
這班人他們怎樣去洞悉到.
知道這個世界有一個真正的神呢.

$^{321}$上帝往往透過一些天文現象.
來幫助那些外邦人.
來明白真理.
在神學裡面我們叫這個普遍啟示.
普遍啟示的意思是什麼呢.
就是說全世界所有人都可以知道的.
不過普遍啟示有一個最困難的地方就是.
普遍啟示只是告訴你這個世界有一個上帝.
但是上帝的工作是怎樣.
上帝具體是誰.
或者耶穌基督是怎樣的一個救贖工作呢.
普遍啟示是不知道的.
所以在聖經裡面.
今天的神學家都是在說.
上帝藉著天文現象.
將一個外邦人.
外邦的占星師.
透過耶觀星像.
似乎這件事都好像.
襯托著.
引證著基督耶穌的降生.
但是這班人因為沒有聖經.
沒有一個點注在那裡.
所以他們都是盲從中.
看他們最熟練的東西.
以致他們可以尋求到真正的耶穌基督.
如果不說你不知道.
在今天這個世界裡面.
今天都有很多人都是面對一個這樣的狀況.
特別是在一些傳福音很困難的國家.
一些很保守的地方.
或者一些很反對基督教.
其他宗教的地域.
很多人會自動無端端摸進去尋找信仰.
為什麼呢.
因為他們做夢.
他們做夢見到耶穌.
很多都是歸順主.
他們因為做夢見到耶穌.
所以聽到耶穌這個名字.

$^{361}$他們就會去尋訪.
而在他們的宗教信仰裡面.
他們知道很嚴重的.
他們這樣做.
但是他們因為很真實的經驗.
很多歸順者都是在夢中尋求.
接著再接受真正的信仰.
找教會來幫他們認識基督耶穌.
到了我們今天這個世界.
還是有這樣的事情發生.
不過在聖經裡面很強調的.
那個也是起始點.
真正認識耶穌.
認識信仰那個.
始終都是要靠聖經的點注.
接著經文裡面就繼續說這個主題.
在第三節到第六節.
經文就這樣說.
剛才說到希律聽見了就心裡不安.
耶路撒冷合成的人都不安.
他就招齊了祭司和民間的民事.
問他們說基督當生在何處.
他們回答說在猶太的伯利恆.
因為有先知記著說.
猶大地的伯利恆.
你在猶大諸城中並不是最少的.
因為將來有一位君王要從你那裡出來.
牧養我以色列民.
希律聽到外面的波斯.
或者東方人來尋找生來的猶太王.
希律因為他心裡不是根正苗紅的猶太人.
他爸爸是歸附猶太.
以至他坐順風車.
這個王其實在他的證跡裡.
在猶大的約瑟夫歷史學者裡記載著.
原來他是一個極度疑心重的君王.
他曾經疑心大到什麼地步呢.
他見到他的部下.
他有一個很漂亮的老婆.
他有些屬下見到他老婆.

$^{401}$見到她眼睛超纖.
心花怒放.
見到美女.
然後見到他的部下.
一次過殺了八個人.
就說看我老婆.
看了就殺了他們.
再加上他自己.
因為他是一個不是根正苗紅的猶大人.
所以當有任何的風吹草動.
猶太人想作亂或者有些反抗.
事實上猶太人自己都不想他做君王.
聽到有任何風聲的時候.
他就採取一種比較極端的手段.
來殺害拉比和先知.
雖然這一方面是殺害.
另一方面他就很豪華.
做很多善事.
當時猶太人的聖殿都是由他來建的.
所以耶穌進去的聖殿.
那個金碧輝煌的猶太人聖殿.
其實就是希律裡面的德政.
對於猶太人來說.
建得很漂亮.
讓他們有聚焦敬拜的地方.
不過希律聽見耶穌基督的出生.
心裡面有一種被取代的不安感.
當他不安的時候.
身邊的人.
整個耶路撒冷的居民.
死了.
希律這麼不安.
身邊的人一定會充滿不安.
所以整個城市.
整個耶路撒冷城.
都會感覺一種不安的感.
來籠罩在他們當中.
這種不安感.
這種感受或者這種情緒.
其實真的會傳染.

$^{441}$在你家裡有時候都會這樣.
你父母或者你工作裡面的老闆.
有一些不安感.
那種不安的情緒.
就好像傳染病一樣.
會直捲你身邊的人.
和你身邊的人一起生活的人.
會經歷到那種類似氣場.
那種不安的氣場.
你不安的時候.
你做出的決定是不安的決定.
你不安的時候.
你會帶著很多的情緒.
憤怒,窒道.
來對待你身邊的人.
在耶路撒冷裡面.
你會看到希律的不安.
大概就會將這種不安.
因為他是當代.
當時的王.
他就會將這種糾纏著權力.
和情緒的管治.
就好像擴散了一樣.
整個城市都會瀰漫著一種不安的情感.
當他不安的時候.
第一件事就是找資料.
問一些什麼呢?.
耶穌.
因為生來基督的王.
他就問那些熟讀聖經的人.
祭司長和文士.
問他們.
基督生在哪裡?.
記不記得剛才說的.
那些東方的占星師.
他們見到天文現象.
但是他們不知道耶穌是誰.
也不知道耶穌應該是怎樣.
就用一個大圖畫來問希律.
大家對於耶穌在哪裡出生.

$^{481}$所有東西都好像一個謎一樣.
這個謎.
或者到具體指清楚在哪裡呢?.
其實就來自聖經.
具體的來指出.
那些文士和祭司長就引用尼加書.
他們所說的那句話是尼加書五章二節.
他說到猶大地的伯利恆.
你在諸城中並不是最少的.
因為將來有一個君王要從你那裡出.
牧養我以色列民.
回到剛才那裡.
這是一個引用舊約聖經.
用聖經的聚焦.
讓人可以知道神是誰.
在哪裡出生.
藉著聖經就把人拉進來.
回到那個焦點.
不過希律很聰明.
他問了這條問題之後.
他得到這個答案.
某程度上他想藉助那幾個占星師.
來找回耶穌出來.
不過在經文裡真的很巧妙.
在這個狀態下.
某程度上也幫了那些.
蒙查查那些看星象的占星師.
他可以能夠.
不是在耶路撒冷找耶穌.
而是回到伯利恆.
距離耶路撒冷很近的一個小城市.
伯利恆裡面尋找生來的猶太王.
在這個經文裡.
他就要告訴我們今天的讀者知道.
耶穌的降生.
面對著一個好像不安的希律.
他好像在做一些破壞的工作.
但是又藉著他這種好像很壞的動機.
做了一些這麼壞的事.
某程度上也幫了一個外邦人.

$^{521}$能夠找到真正生來的猶太王.
這個是留下來的一些伏線我們再去看.
不過我想在這裡再看兩個字.
首先就是希律在聽見.
這群人就是很不安的那種管治.
然後到經文裡說到生來的猶太王.
將來有一個君王從尼那裡出.
穆養我以色列民.
在經文裡在舊約聖經.
他用穆養這個字來說出那個君王.
那個猶太人的王那種的.
那種的心腸和那種的角色.
或者那個作法.
在希律的手段裡面剛才說過了.
疑心很重用強權.
還有又威逼又利誘.
令到那些百姓就被迫要信佛.
不過在經文裡面真正的王.
猶太生出來的君王.
不是用這些方法.
他是用穆養這個.
好像穆養人的這種心態.
來帶領他整個民族的.
值得注意的是.
同樣在舊約裡面.
用同樣穆養這個字根的那個字.
在詩篇第二章第九節.
都是用這個字.
那個穆養的字.
其實和合制的字是同一個字.
字根是一樣的在舊約.
穆養和合制.
負面的就是合制.
正面的就是穆養.
穆養是說到上帝好像一個穆羊人.
一直從舊約.
你讀詩篇你知道.
上帝對子民的心態.
走了穆羊人.
帶領羊仔這樣.

$^{561}$很浪漫的.
有時候會惡.
但是目的不是要傷害羊的.
他要帶領他去清草地.
在水旁能夠領受生命的恩典.
這個是穆養的角色.
不過去到如果是強權管治之下.
那個就是合制.
就是用威逼.
用利誘.
又或者利用當時希律所擁有的權力.
來去合制住當時猶太的百姓.
這個聖經裡面他很想說.
某程度希律不是真正上帝.
真是在尼加書裡面所啟示出來.
或者是說將來要應許出現的君王.
他的不安感是很正常的.
因為他很擔心將來有一個王要取代他.
以致他在一種擔心的狀態之下.
他還更加多一重的去負面.
去做一些工作在那些百姓身上.
不過去到今天我們去拿這段經文去理解的時候.
合制本身的意思就是管理.
去帶領.
本身不是負面的.
不過在聖經裡面他藉著耶穌基督的降生.
他講到我們今天的身份.
那個君王是在牧養我們.
同樣在我們生命裡面.
他學習基督耶穌的樣式.
或者我們的弟兄姊妹.
在家庭當中的那種關係.
都是用一種牧養的那種關係.
彼得前書他都是用這個字.
用合制和牧養.
彼得在講作長老的教會長老.
不是合制著你的羊群.
而是成為一個牧養人.
因為我們有一個大牧人來牧養我們.
所以我們作領袖也好.

$^{601}$作長輩也好.
我們就要牧養身邊的人.
在彼得前書裡面.
五章二至三節再繼續的去講.
具體的牧養是甚麼意思呢.
凡事作榜樣.
牧養就是凡事作榜樣.
就是叫那些東西去做.
自己去做先.
你去做先.
新鮮事出的時候.
其他被你帶著的後輩也好.
家庭成員也好.
其他在教會裡面服侍的也好.
他們就會跟著來做.
如果你叫他做.
自己退後.
那不是牧養.
曾經有一個這樣的比喻.
我聽過一些老牧者說.
牧羊的人通常.
你猜牧羊群的前面還是後面呢.
趕羊的那些懶的牧童.
就是找一隻牧羊犬走前去.
然後在後面掃一掃.
這樣就有一天了.
那些叫做趕羊人.
在後面而已.
真正的牧羊人.
是拿著那隻牧羊的象.
在那些象群面前.
一直帶領著羊仔.
牧羊的象上面有個勾.
知道有什麼用嗎.
那些羊其實很蠢的.
走著走著會自己偏行幾路.
衝下山崖也不知道.
一滾一隻.
你不要想著一隻跌死.
羊群心態.

$^{641}$一起的就滾下山了.
那牧羊人的象.
那隻象的勾是做什麼呢.
勾著羊的角.
當牠在那裡蠢蠢蠢蠢衝下去的時候.
他就找一隻牧羊象.
那隻勾勾著牠的角.
拉牠回來.
那隻羊一定會說一句什麼呢.
「咩呀」.
走回來.
這樣就走回正路.
不要滾下山崖.
牧羊人的角色就是.
當那些人「咩呀咩呀咩呀」的時候.
牧羊人的象大力一搖.
搖回來走回正路.
今天我們也是一樣的.
我們的長輩也好.
我們家裡也有牧羊的角色.
偏行幾路.
我們首先第一件事.
不是在後面看牠怎麼坐.
走在牠的前面.
然後用.
令牠有時候不舒服.
但是要令牠幫助牠走回正路的方法.
用一隻象勾回去.
即使牠說什麼也好.
你也要拉牠回來.
讓牠不要衝下山崖.
不要走向死路.
牧羊的意思就是走在羊群的前頭.
凡事作榜樣.
今天對我們來說.
這也是一個很重要的訊息.
以色列的君王.
是牧羊整個以色列文.
這個君王會身先至尊.
耶穌基督.

$^{681}$在希伯來書裡就說到.
凡事就作了我們的元首.
他打開的那條信仰的路.
是他自己帶領我們這些屬神的子民.
來出死入生.
我們只是跟隨基督.
這個君王的腳蹤.
同樣我們都是領受.
這個呼召.
我們每一個.
我們生命裡總會有一些影響力.
我們身邊的家人.
我們的後輩.
我們在教會裡的羊仔.
我們都是走著基督的路.
我們也是引領著我們後面的.
那些羊群.
那些人.
走在基督的路上.
凡事作榜樣.
就是這個新生的王的榜樣.
也是對我們今天.
需要理解.
相反來說.
那個希律.
他用強權.
用其他的方法.
來去合制住那些人.
到最後那條不是真正要走的路.
希律到最後.
他自己都要走.
到最後他自己都要承受自己的惡果.
我們再繼續看第七節到第十二節.
剛才說到希律.
知道了.
原來這個生下來的基督.
就在伯利恆城.
接著第七節.
當下希律暗暗地召了博士來.
叫那些占星師.

$^{721}$細問那星是什麼時候出現.
就猜他們往伯利恆去說.
你們去仔細尋訪那孩子.
尋到了就來報信.
我好去拜他.
他們聽見黃的說就去了.
在東方所看見的星.
忽然在他們前頭行.
直行到小孩子的地方.
就在上頭停住了.
他們看見那星.
就大大的歡喜.
進了房子.
看見小孩子和他母親瑪利亞.
就俯伏在那孩子揭開寶盒.
拿黃金與香活躍.
為禮物獻給他.
博士因為在夢中被主指示.
不要回去見希律.
就從別的路回本地去了.
在這裡你會看見.
這是做一個另外的比較.
剛才說到希律.
心裡有很多不安.
合成的人都很不安.
那種情緒大家都爆錶了.
但另一方面.
這個希律.
是雙面人.
不是壞這麼簡單.
雙面人.
不是壞這麼簡單.
他不是一個.
除了疑心重他也有些手段.
他知道那顆星是指向伯利恆.
於是就叫那些.
占星師.
問候何時.
然後就猜險他們.
命令他們去到伯利恆.

$^{761}$然後就跟他說.
你找回那個小孩子出來.
你找到的時候.
就回來告訴我.
然後他就說.
你也去拜祭他.
他的意思是這樣的.
其實如果他是一個.
很有威勢的王.
他要怎樣做呢.
如果沒有那麼.
工於心計.
或者不是那麼.
他想低調.
不要令猶太人生亂.
所以他第一件事做.
就是避地裡來做.
第二件事就是當他找到這個孩子的時候.
我們知道.
這個孩子的目的不是希律.
想去拜祭他.
是想殺了他.
後來經文也知道.
找不到就將整個伯利恆城的孩子.
兩歲以下全部殺了.
目的就是要殺人滅口.
剷除那些.
威脅他自己的.
未來的君王的可能.
不過在這裡.
他猜險.
這些.
那些人來到去.
某程度上就像一個負面的幫助.
剛才說過.
幫助了這些.
尋找這顆星.
尋找耶穌基督的這班人.
去到伯利恆找到耶穌.
所以這些人.

$^{801}$其實是糊塗的.
猜他去就去.
他也不知道.
其實是非常危險的一件事.
不過他們聽見他們的話.
就去了.
被這顆星引導.
在位置上.
收去光芒.
接著.
這幾個人.
和當時希律.
和耶路撒冷的城.
裡面的人的情緒反應.
是完全兩回事.
當那些人嚴重不安.
很害怕的時候.
這幾個人糊塗.
但見到這個孩子.
聖經裡面用.
用了很多字眼.
大大歡喜.
面對著很不安的世界.
所有人都.
不知道發生了什麼事.
但這幾個人.
只是很專注找耶穌.
一個專注找耶穌的人.
外面的事他也不太理會.
反而他們的目的.
就是要尋訪.
這個生下來的.
耶穌基督.
所以當他們找到的時候.
他們就快樂到不得了.
大大的歡喜.
他們將他們整個行程裡面.
最重要的.
最高峰的.
就是要向這個.

$^{841}$生下來的猶太君王.
來下拜.
大人大者.
從東方千里迢迢.
其實某程度.
不太關他們事.
他們.
在波斯.
有他們自己國家的生活.
不過.
這顆星他們就是尋找這顆星.
他們知道猶太人的王.
反而他將.
這顆星成為他生命的.
目標的一部分.
即使外在的環境.
多麼恐怖.
希律裡面的手段有多麼凌厲.
但是呢.
他都是按著這顆星的指示.
找到他們心目中.
所要找到的這個.
小童的耶穌.
向他獻上最.
最名貴的禮物.
黃金乳香活藥.
當然今天有些聖經.
或者有些童宮會用.
黃金乳香活藥.
為什麼送這些呢?可能用靈異化的解釋.
來處理.
不過順便問問你.
沒有特別再說黃金乳香活藥.
即使耶穌釘十字架的時候.
用的活藥.
都不會是這幾個博士.
送的活藥.
其實是沒有.
直接關係的.
唯一的解釋就是.

$^{881}$這些是最名貴.
在東方裡面最好的東西.
那時候沒有保應丹送.
沒有珍珠末那些東西送.
送一些藥.
一些值錢的東西.
由東方送到來.
更加重要的就是.
向著這個小孩耶穌.
來下拜.
值得重用一個更加有趣的地方.
是.
剛才這幾個.
來自東方的人和希律.
都有些對話.
經文裡面看到這個.
在做管治的王.
聖經沒有特別記載.
沒有向這個王下拜.
沒有特別去做一些.
俯首稱臣的表現.
甚至乎他做的事.
都好像某程度是.
阿王啊,請你幫幫我.
我找一個人來取代你的.
我會和他下拜.
其實說這件事很冒犯.
不過那幾個人好像.
懵聖城,什麼都不知道.
他們就是單一.
要找這個小孩的耶穌.
要向他下拜.
要向他獻上禮物.
在聖經裡面就說到這幾個人.
就是萬萬壯壯.
來自東方.
但是將最好的禮物.
獻給基督耶穌.
經文裡面.
我們大致看到.

$^{921}$一個這樣的脈絡.
那個故事或細節我們知道.
在這裡我們有兩個地方值得.
用這段經文來幫助我們.
當基督耶穌出生的時候.
在聖經裡面記載.
一個小孩子他都未能夠有.
行神蹟的能力.
他都未做任何事.
但是已經扭轉了當代.
希律.
或者在猶太人社會裡面.
已經製造了一些.
一些的動盪.
耶穌的出生.
改變了.
在當代裡面.
有一些人的.
心目中的一些想法.
可能他的出現.
會挑戰著.
或者會將心裡面.
潛藏起來的那些罪.
那些.
不是合乎上帝心意的東西.
會揭露出來.
當這些東西被揭露的時候.
那種不安感.
是會呈現出來.
不過.
耶穌基督的降生.
就是要將.
這個混亂的心態.
未必要說.
我不會無限上綱.
講任何政治.
講我們的生命.
當基督耶穌在我們生命裡面出現的時候.
我們有很多的不安感.
有一些剛剛認識主的.

$^{961}$一信耶穌.
第一件事要看到的是什麼呢?.
除了看到耶穌的愛之外.
也看到自己的罪.
這些東西一揭的時候.
其實是.
不容易的,挺痛苦的.
但是正正是這樣的時候.
當這些東西.
被處理的時候.
就更加有一個很重要的.
那個價值.
真的要面對.
說回這個堂子.
我們裝修完之後.
大家不知道,有些地方你不察覺的.
油漆很漂亮,坐在角落那裡.
其實那次我看著.
那些天花頂拆了.
大家知不知道?沒有天花頂了.
以前有天花頂遮住的.
那些冷氣.
又有天花頂包圍住.
所有東西.
很漂亮的,是吧?.
不過那次一揭開.
嘩!.
那些生鏽水.
你們有沒有想過.
那些水不是清的.
大家不要覺得噁心.
你們感覺一下.
想像一下,一揭開那些水.
有些鐵鏽色.
黏的,好像一些痰.
然後冷氣因為.
太久,或者什麼的.
不通的,然後那些師傅說.
下一點化痰藥就行了.
化痰藥?冷氣機.

$^{1001}$那些水太久了.
黏了,所以痰那樣.
他們說下一些化痰藥.
但是遮住了,看不到的.
但一揭開的時候.
所有的東西.
呈現了出來.
我們信仰都是一樣的.
當基督耶穌出現的時候.
我們生命裡邊.
那些污穢.
那些黏黏的.
除了看到上帝的恩典之外.
我們也會看到自己的庇護.
那些東西.
看見我們會不安.
我們以前是避開的,找東西遮住.
但是正正是這些.
藏污納垢.
卻是糾纏在我們的生命當中.
以致我們經常都會帶著.
不安感.
如果我們還要做一個好樣子.
或者找些東西遮住它的時候.
我們這種不安.
其實是散播給身邊的人.
其實你掩飾著你生命的罪.
掩飾著你生命的問題.
身邊的人是知道的.
身邊的人其實都是在吃.
你那些不安.
那些污穢.
將這件事繼續來發大.
當我們今天.
我們一起預備.
慶祝基督耶穌降生的時候.
我們也容讓基督耶穌的生命.
光照我們.
我們也容讓.
基督耶穌的生命光照我們.

$^{1041}$看看我們的心靈.
心處裡面會不會有這些.
很多的不安.
很多東西我們遮掩著.
又或者有很多的方式.
來令自己好樣一些.
令自己的污穢不要被別人看到.
但是聖經.
它是在挑戰我們.
當我們讀的時候我們要問一問.
我們的心靈心處會不會有這些不安.
呈現在這些.
在我們的生活當中.
是我們第一個思想.
第二個.
記得剛才那幾個來自東方.
波斯的占星師.
我們叫博士.
他們帶著.
最純正的動機.
他們.
夢察察.
去說一些.
當時最尖銳.
最大加最卑微的.
猶太人的皇生在哪裡.
其實這個.
這動物本身是一個.
今天如果你稍微明智一點.
你明白了那個處境.
你不會這樣做的.
不過他們就是.
什麼都不知道.
他只是帶著最單純的動機.
尋找耶穌.
正正就是因為這種.
單純的動機.
所以身邊裡面所有的不安感.
都不能夠侵擾到他.
正正就是這一份.

$^{1081}$好像.
去到一個.
傻勁的字來形容他們的時候.
這一種.
單純的心.
令到這些人可以排除萬難.
令到這幾個人.
可以真正去到小孩.
耶穌面前獻上最好的禮物.
今天對我們來說.
我們今天生命太多掛慮.
我們今天很多東西.
很多東西都變得成黃成紅.
我們怕這些.
怕那些.
又或者看環境.
做人看得太多.
以致我們失去了焦點.
我們所有東西都覺得來勢洶洶.
所有東西都很不友善.
不過今天.
當我們看回這幾個.
來自東方的博士.
他們單純的心.
都提醒我們.
我們生命裡面有很多東西其實不需要考慮太多.
你考慮太多反而你會怕多一點.
反而你專注在基督耶穌.
你一生.
你去尋求祂的時候.
你的生命就可以舒坦.
上帝會保守你.
這幾個人懵聖城.
到他們拜完耶穌之後.
上帝都報夢給他們.
不要回去.
有殺身之禍.
然後就按著那條路回到自己的本地.
即使外面的環境.
這麼兇猛.

$^{1121}$尋求上帝的人.
上帝會保守他.
今天我們都需要這一份這樣的所敬.
需要這一份這樣的單純.
單單仰望那位基督耶穌.
我們的生命就可以舒坦.
即使外在環境有多不安都好.
我們都依然可以很專注的.
來找到這個.
為我們生命帶來改變的基督.
這個生下來的猶太君王.
他的影響力已經不再是在猶太民族.
或者這個城市裡面.
他的平安已經覆蓋著.
由這幾個東方的波斯人開始.
這個福音.
亦都延展到我們今天每一個弟兄姊妹當中.
求主繼續幫助我們.
讓我們這一段經文繼續拿回家.
踏真了它.
經歷耶穌生命裡面帶來的那種豐富和預備.
願主繼續來幫助我們.
\newpage



\section{路加福音 2:1-7}
\label{sec:1SOWC0lwIU0}
\textbf{2022年12月25日崇拜直播｜陳冠成牧師｜展現脆弱的一面｜路加福音2\_1-7}
\newline
\newline
連結: \href{https://youtube.com/watch?v=1SOWC0lwIU0}{\texttt{https://youtube.com/watch?v=1SOWC0lwIU0}} ~~~~ 語音日期: 2022-12-29
\newline
\newline
\hyperref[sec:6BkHUdPWvFs]{\small{< < < PREV SERMON < < <}}
~
\hyperref[sec:index_chronic]{\small{[返順時目]}}
~
\hyperref[sec:HUos2vHhz9s]{\small{> > > NEXT SERMON > > >}}
\newline
\newline
路加福音 2:1-7
\newline
\begin{longtable}{cl}
\hline
\hline
章節 & 經文 (和合本修訂版)\\
\hline
2:1 & \begin{tabularx}{0.7\textwidth}{X} 在那些日子，凱撒奧古斯都降旨，叫全國人民都登記戶籍。 \end{tabularx} \\ \\ \relax
2:2 & \begin{tabularx}{0.7\textwidth}{X} 這第一次登記戶籍是在居里扭作敘利亞總督的時候行的。 \end{tabularx} \\ \\ \relax
2:3 & \begin{tabularx}{0.7\textwidth}{X} 眾人各歸各城，辦理登記。 \end{tabularx} \\ \\ \relax
2:4 & \begin{tabularx}{0.7\textwidth}{X} 約瑟也從加利利的拿撒勒城上猶太去，到了大衛的城名叫伯利恆，因為他是大衛家族的人， \end{tabularx} \\ \\ \relax
2:5 & \begin{tabularx}{0.7\textwidth}{X} 要和他所聘之妻馬利亞一同登記戶籍。那時馬利亞已經懷孕。 \end{tabularx} \\ \\ \relax
2:6 & \begin{tabularx}{0.7\textwidth}{X} 他們在那裡的時候，馬利亞的產期到了， \end{tabularx} \\ \\ \relax
2:7 & \begin{tabularx}{0.7\textwidth}{X} 就生了頭胎的兒子，用布包起來，放在馬槽裡，因為客店裡沒有地方。 \end{tabularx} \\ \\
[1ex]
\hline
\hline
\end{longtable}
$^{1}$主席.
問問你自己.
你又願不願意經常在別人面前展示自己的脆弱?.
你又願不願意夠膽向人真情流露?.
你又舉手給我看.
少了很多.
弟兄更少了很多.
這是一個特別的警況.
也可以說是一種矛盾.
正如美國休斯頓大學有一位很出名的社工系教授.
布朗教授Bernie Brown.
他說得很好.
我們很渴望看到別人真實和真誠.
但是我們很害怕被人看到我們自己的真情流露.
你很想看到對方是真誠待你的.
看到他真實的一面.
哪怕是脆弱的一面.
但自己卻會收得很近.
不想被人看到.
或者很害怕被人看到自己脆弱的一面.
原因是什麼呢?.
他說因為你向別人赤露闖開的時候.
當你向別人展示自己的脆弱的時候.
其實是很冒險的.
其實是一個很不安全的行為.
因為你不知道對方有什麼反應.
怎樣去看你.
怎樣去回應你.
不過如果我們真的要和別人產生一個真誠的連結的話.
我們真的很想拉近和對方之間的距離的話.
其實我們確實又要學會擁抱這種展現脆弱的力量.
就正如剛剛世界盃上星期打完了.
哪隊贏了?.
阿根廷了.
無論你捧法國也好.
捧阿根廷也好.
有沒有人喜歡日本隊的?.
有了.
看吧.
日本隊.

$^{41}$沒有了?這麼少人?.
或者你聽過了.
其實很多人喜歡日本隊的.
如果你有留意新聞報導.
特別今屆世界盃.
其中有一位球員.
他叫做三尖勳.
有沒有聽過?.
25歲的.
很年輕的.
他是其中一位日本國家隊的代表球員.
他有份對克羅地亞那場賽事.
最後階段主射12碼.
雖然很可惜日本隊未能晉身八強.
但三尖勳這位球員其實在很多人心目中留下了一個很深刻的印象.
雖然他在最後階段互射12碼的時候.
他射失了.
被對方克羅地亞的球員撲出了.
令到最後日本隊在互射12碼的結果.
以1比3出局.
亦令到這位首次在世界盃出局的三尖勳.
哭到你真的想都想不到.
他由球場射失了開始哭.
哭到完場都繼續哭.
哭到接受訪問都繼續哭.
哭成淚人.
他表示自己其實已經在之前一天做好準備.
練習了很多.
他覺得自己是有能力承擔那一個時刻.
他已經做好準備.
所以他自動請英舉手.
向教練說我去主射其中一球12碼.
他以為自己可以完成任務.
但結果又怎樣?.
留下了人生的遺憾.
他很真誠地對傳媒說.
我今次覺得自己是闖了大禍.
而被問到為何哭成這樣.
他就說為了那些比我有更強感受的人來難過.
你看到很特別.

$^{81}$三尖勳這個球員在接受訪問時.
他沒有所謂像我們今天這樣.
做一個公關show.
做一場公關talk.
寫好稿子在傳媒面前很大方地說.
我很對不起大家.
我辜負了所有人的期望.
我之後會繼續努力成為一位出色的球員.
我會為下一屆做好準備.
然後起身鞠躬離開.
很冠冕堂皇.
很見得人.
他沒有這樣做.
反而是他很勇敢地在大眾面前流露他的真性情.
他淚流滿面承認自己的失敗.
結果很奇妙.
大家沒有人記得他射失了十二碼.
反而很欣賞有一個活生生的人.
在他們面前展露自己的脆弱.
確實是的.
不知你是否也有這樣的想法.
比起聽很多成功例子的見證.
其實有時我們更加渴望聽到.
好像三尖分這些真實.
有血有肉的分享.
因為它代表著我們.
因為我們都很明白.
其實沒有一個人是完美的.
我們都會面對挫敗.
都有難過掙扎的時候.
不過我們未必選擇和別人分享釋放出來.
但我們確實需要像三尖分一樣.
適當的時候將你壓抑的情緒.
用一個正確的方法去呈現出來.
所以,亭姐妹.
人總是有堅強和脆弱的一面.
其實兩樣東西沒有分對錯.
不過我們會否只是習慣展示好一面.
堅強的一面.
硬朗的一面.

$^{121}$但對於脆弱,挫敗,掙扎的一面.
我們就藏起來不讓別人知道.
這正好是我們今天很值得反省.
我曾經聽過有人這樣說.
「為什麼你那麼怕向人展露你的脆弱?」.
對方就回應.
「如果要我赤裸裸地讓人看到我的脆弱」.
「我寧願扮作堅強」.
「怎樣也不讓人知道我這些掙扎」.
因為我很害怕.
我害怕別人小看我.
我害怕別人討厭我.
又或者我很害怕.
當我分享這些軟弱掙扎的時候.
其實我會否增添別人的麻煩?.
我會否反而令到我身邊的人為我擔憂.
但又幫不到我?.
我們會否也有這些想法?.
不過其實很多研究告訴我們.
原來能夠合適地展現自己的脆弱.
其實是一件很迷人的事.
是一件很輕鬆的事.
我們的脆弱其實比我們的想像吸引人很多.
所以你看回耶穌.
耶穌的出生不是以一個強大的領袖.
不是以一個硬漢的姿態登場.
而是透過一個最脆弱.
脆弱到根本沒能力照顧自己.
很無助的.
耶穌作為嬰孩.
祂一定是需要別人照顧.
唯有這樣祂才可以存活.
耶穌是用一個這樣的嬰孩形象.
來吸引人靠近祂.
嬰孩耶穌不單單向我們表示.
祂是一位可親的上帝.
祂是一位願意與人連結交往的上帝.
其實祂同時教曉了我們一件事.
就是原來我們讓別人看到自己的脆弱.
甚至無能無助.

$^{161}$其實不是一件很羞愧的事.
耶穌祂自己也不覺得羞愧.
不是一件那麼可怕的事情.
相反.
耶穌告訴我們.
是一個勇敢,愛的表達.
意思是什麼呢?.
在座有結婚的人士.
有沒有人在拍拖?.
應該有吧?.
不敢說話吧?.
有沒有試過示愛?.
沒有嗎?.
那怎樣拍拖呢?.
抽籤嗎?.
還是導師幫你匹配呢?.
應該不是吧?.
應該有示愛的.
不過不方便說.
無論男女示愛或相反.
有婚姻,有拍拖就有示愛.
有向人表達過愛意的經驗.
其實需要勇氣嗎?.
需要的.
為什麼要很大的勇氣?.
因為你不知道對方有什麼反應.
對方高興接受當然好.
但對方強烈拒絕怎麼辦?.
最慘是兩者都不是.
不表態.
吊著你忍.
那一刻你還不心腸.
所以跟對方表達.
無論是你心儀的對象.
是你的伴侶或家人,朋友.
其實等於你毫無保留.
赤裸裸揭開自己給別人看.
並且預計你的情感.
有機會受到傷害.
我還記得當年追求太太的時候.

$^{201}$我準備向她表白的時候.
其實心情很忐忑.
即使總覺得自己有機會成功.
條件也不差.
但同時又會很害怕.
萬一她拒絕怎麼辦?.
其實很尷尬.
她還要是我的同事.
很尷尬.
我會否嚇走她辭職不做?.
以後就沒得見面?.
那種心情其實不難理解.
不知道大家有沒有看到圖畫.
當你很想跟對方談戀愛的時候.
你會撕花瓣.
她喜歡我,不喜歡我.
成功嗎?.
原來就是這樣.
很真實.
結果我去表白.
是否一次就成功?.
當然不是.
我第一次是失敗了.
硬生生地拒絕.
因為當時對方是基督徒.
他送了一段經文給我.
雖然我不懂解釋.
但我知道我沒有行.
因為他說「信與不信不能同乎你呃」.
成功擊退了我.
很傷心的.
失戀的感覺.
好像被針刺,刀插.
很辛苦.
但我沒有放棄.
後來再鼓起勇氣.
再一次冒著被受傷的風險.
再去表達愛意.
那時後來已經信了主.
過了很久.

$^{241}$最終打動了對方.
今年十一月剛好.
慶祝二十週年結婚紀念.
真正的弟兄姊妹.
其實你想想.
每一段情感的建立和發展.
其實是有風險的.
無論是友情,親情和愛情.
都會存在多多少少風險.
當然我們不去表達愛意.
不肯去表達心事.
你當然不會受傷.
因為對方都不知道.
但很大機會我們永遠體驗不到.
什麼是真實的愛.
所以有人說得很好.
有愛就一定有風險.
因為你永遠不知道.
對方會是什麼狀況下一刻.
又或者對方有什麼反應.
有可能會不辭而別.
有可能不再愛你.
甚至背叛你.
令你傷心欲絕.
你看看耶穌也是這樣.
耶穌是因為愛的緣故.
所以他選擇來到這個世界.
還要選擇以一個英雄的方式.
來到這個世界.
甚至乎之後為世人.
來到犧牲受死.
風險大不大?.
極之大.
因為耶穌其實自己都知道.
原來總會有人不領情.
總會有人拒絕他的恩典.
他的愛.
他的拯救.
你知道在科幻書看到很多.
他出世沒有領情又如何?.

$^{281}$有人想要他的命.
又或者有人對他的救恩嗤之以鼻.
也有些人.
起初是跟隨耶穌的.
但後來發現耶穌好像幫不了忙.
不是我想的.
就慢慢疏遠他.
也有些人曾經說至死不渝.
要愛耶穌跟隨耶穌.
不過到了危急關頭.
禍患臨到的時候.
他們又拋棄了耶穌.
但即使如此.
你看到耶穌仍然願意愛.
願意冒險.
他選擇赤裸裸來到這個世界.
其實就是毫無保留告訴我們.
他有多麼愛我們.
愛到一個地步.
願意為我們冒風險.
並且不單止停留在這裡.
耶穌其實也很願意幫我們.
打開我們的心扉.
意思是從前我們可能.
跟別人交往的時候.
跟別人建立關係的時候.
我們曾經可能在情感裡受到傷害.
於是我們就將自己的內心收起來.
不想再分享自己的內心世界.
讓你身邊的人知道.
可能是你的伴侶.
可能是你的父母.
可能是你的兄弟姐妹.
或者你的弟兄姊妹.
你將自己的心收起來.
因為你很怕再受傷害.
我再分享自己內心的事.
會不會被人取笑?.
會不會他看不起我?.
會不會他到處說?.

$^{321}$你曾經有過這樣的傷痕.
所以你下次再不想分享.
但是耶穌的愛.
要再一次釋放我們.
要再一次融化解開了.
我們被冰封的心.
祂很想幫我們.
你嘗試再次放開懷抱.
相信祂.
相信你身邊的人.
你再一次鼓起勇氣.
邀請人進入你的內心世界.
讓他們有機會.
分享,守望你生命中的脆弱位置.
也給機會你身邊的人.
關心,幫助你.
所以不妨停下來想一想.
會不會在過往的日子.
我們曾經受到傷害.
所以我們很怕再流露自己的軟弱呢?.
而久而久之.
也窒礙了我們和別人之間的友誼.
我們和別人之間的關係.
甚至是夫婦關係.
唯有重新開放我們自己.
慢慢一點一滴去學習.
分享一些感受.
一些內心的需要和想法.
其實你才可以重新去真正投入那段關係.
無論和別人的關係是怎樣.
其實和上帝的關係也是這樣.
一旦我們意識到.
我們本身是脆弱的.
我們本來就像脆弱的耶穌.
像嬰孩般渺小的時候.
其實你就自然願意靠近上帝.
你就自然願意和身邊的人重新連結.
開始發展正常的相交.
所以弟兄姊妹.
今天這段經文雖然很短.

$^{361}$但是耶穌以嬰孩的樣式來到世上.
其實祂在告訴我們一個很重要的訊息.
你發覺耶穌來到世上.
祂第一時間不是開口教訓人.
是用祂自己一個很純粹的生命.
一個赤裸,嬰孩,脆弱的生命.
告訴我們幾件很重要的事.
第一方面就是祂和我們原來一樣.
祂經歷過人世間那種無助無能.
祂也需要別人的照顧.
需要別人的保護和同行.
祂經過這些時刻.
所以祂很明白我們.
其實不單止這樣.
如果你有留意.
耶穌將離世之前.
祂在喀西瑪利園.
同樣展示了祂脆弱的一面.
祂面對上帝給祂的任務.
祂的軟弱.
祂多麼想逃避.
但是藉著祂和上帝的分享,禱告.
祂慢慢去跨性.
祂也有能力幫我們去跨性.
我們現在這一刻所面對的掙扎和軟弱.
事實上這位向我們赤裸闖開的主.
祂也想我們以這樣的方式回應祂.
祂很想我們同樣對著祂.
來打開我們的心扉.
你可以隨時找祂傾心吐意.
哪怕是你生命中覺得最醜怪.
最不見得人最幽暗的事情.
其實耶穌願意聽.
耶穌願意接納.
並且祂能夠有安慰和盼望.
並且耶穌祂很想我們.
不單止向祂傾心吐意.
耶穌想我們倚靠祂.
你也可以效法祂.
來接納身邊同樣有軟弱的人.

$^{401}$所以,弟子們為此.
求耶穌幫我們.
我們不單止開放自己的情感.
同時我們也學習.
從主裡面支取更多的愛心和耐性.
你去承載別人的脆弱.
你去和那些掙扎有挫敗的人同行.
我還記得幾年前.
一群弟兄在分享婚姻經歷.
一些掙扎.
其中一位弟兄分享到.
他和太太已經冷戰了幾年.
其實不知道再如何溝通下去.
於是其他弟兄出於愛心.
提供很多建議給他.
很希望幫他解決即時的困難.
你嘗試寫信給他.
買一份禮物給他.
或者約他吃燭光晚餐.
甚至給他WhatsApp.
每天慢慢感動他.
給很多建議.
希望幫他解決這件事.
不過這位弟兄很沮喪地說.
「你以為我沒有試過這些嗎?」.
對方一點反應都沒有.
「我試過了,還可以做什麼?」.
這位弟兄很灰心.
其他人似乎有些計窮.
不知道如何回應.
「丁姐妹,如果你是你」.
「你怎麼辦?」.
「繼續給答案?」.
「繼續提供意見?」.
「還是很簡單聽他說」.
「認同他這一刻的情況」.
「我體會到你真的很艱難」.
「用盡所有方法」.
「你很愛對方」.
「但對方此時此刻好像不領情」.

$^{441}$「你也不知道該怎麼做」.
「你告訴對方」.
「你是認同他」.
「你是接納他」.
「甚至一個擁抱」.
「一個禱告」.
「很多時勝過千言萬語」.
確實丁姐妹.
假如真的有人在你面前.
願意分享他的掙扎.
軟弱的時候.
甚至哭泣的時候.
不要錯過上帝給你服侍的機會.
他信得過你才會這樣做.
他當你是朋友.
當你是知己.
他才願意讓你看到他最脆弱的一面.
所以在這個時候不要害怕.
多數我們都會害怕.
我們一害怕就會.
很想快些寫好答案.
於是給他很多答案.
「你這樣吧」.
「希望可以幫到對方」.
「然後對方就不再出聲」.
「那你自己就平安了」.
但其實反過來.
你不要害怕.
你也不要介意他哭泣.
反而你心裡祈禱.
讓你遇見這個時分.
你可能會很擔心.
但你求上帝幫助你.
不要走開就行了.
讓你和他和上帝.
一起三個的互動.
你珍惜這個時刻.
因為是上帝給你的.
你也不要急於提供答案.
重點是你留下.

$^{481}$靜靜地聽他講.
假如對方哭泣的時候.
千萬不要阻止.
你要相信.
眼淚其實是情緒的淚水.
他能夠哭泣是有益.
因為多多少少.
他能夠抒發到他的情懷.
假如你覺得適合的.
你可以拍拍對方的肩膀.
如果是同性的.
甚至乎給一個很真誠的擁抱.
向對方表示.
你知道你接納他脆弱.
甚至乎你不懂得回應的時候.
你真的可以很真誠地說.
聽到你這樣說.
我也很難過.
其實我也不知道.
怎樣幫你,怎樣回應.
但我很願意留下.
陪伴你.
你想說你可以繼續說.
你很真誠地說出自己的狀況.
也很真誠地表示.
你會陪伴他.
讓對方知道.
即使在他人生最糟糕的時候.
原來他不孤單.
原來仍然有人.
關注他的生命.
願意接納他,關心他.
即使問題解決不了.
但情感就連結了.
所以剛才提到的Bernie Brown.
她說得很好.
她說很多時候.
我們的事情之所以變好.
不是因為問題解決了.
是因為你願意和對方的情感連結.

$^{521}$你願意讓對方知道.
原來他是不孤單.
這樣事情就會變好.
所以弟兄姊妹.
其實今天在教會群體裡面.
我們很多時候強調.
我們要有策動事工的能力.
解決難題的能力.
我們要有很多做事的能力.
這一切都是好的.
但我們不能缺少的就是.
我們願意盛載別人情感的能力.
我們願意聆聽別人的恩賜.
我們願意接納對方軟弱的情懷.
事實上.
教會不只是分享開心的地方.
教會也可以分享掙扎,挫敗,軟弱的地方.
因為無論是信或逆.
上帝還在嗎?.
上帝還在.
即是上帝的恩典仍然在.
在上帝的手中.
這一切都可以變成祝福.
變成每一個人成長的途徑.
所以弟兄姊妹千萬不要錯過.
別人向你真誠流露的時刻.
可能我們真的會不習慣去面對.
但求主幫助我們.
一方面我們願意開放自己的心靈.
和對方分享.
亦都願意盛載別人的情緒.
我相信這樣的時候.
人與人之間的牆壁隔膜.
就會慢慢鬆開.
亦都令很多曾經受傷的心靈.
可以經歷到醫治和安慰.
又或者破損的關係.
可以漸漸得到修補.
很盼望今天這段經文.
今天的分享.

$^{561}$能夠造就我們整個群體.
讓我們不只是用我們的理性.
用我們的知識來經歷信仰.
我們亦都同時用我們的情感.
用我們的心靈來經歷信仰.
好嗎?.
讓我們同心一起祈禱吧.
求上帝幫助我們.
因為你親自降世為人.
以一個脆弱英骸的樣式來到世上.
你帶頭向我們所有人.
展示你最脆弱的一面.
讓我們知道其實.
你願意和我們認同.
你願意接納我們作為人.
在生活裡面面對種種軟弱和掙扎.
並且我們可以在你的面前.
坦誠地分享.
甚至誇耀我們的脆弱和挫敗.
因為我們深信.
你的能力在我們的軟弱上.
能夠完全彰顯.
聖靈啊為此.
我求你幫助王浸每一位.
幫助我們每一個.
讓我們真的可以靠著主你.
慢慢放開我們的心懷.
我們能夠和身邊的人彼此坦誠.
我們可以互相扶持.
互相守望.
也讓我們像耶穌基督一樣.
有一個柔和謙卑的心.
有更多的愛心,耐性.
去守望身邊的人.
去承載別人的軟弱.
以致我們真的能夠.
和別人建立真誠的關係.
一起在主的裡面.
在靈的裡面.
經歷成長,經歷更生.

$^{601}$為此我們同心獻上祈禱.
奉基督耶穌你得勝的名字.
來祈求.
阿們.
謝謝大家今天來看我們節目 拜拜.
\newpage



\section{約翰福音 1:43-51}
\label{sec:HUos2vHhz9s}
\textbf{2023年1月1日崇拜直播｜梁家麟牧師｜掌握召命｜約翰福音1\_43-51}
\newline
\newline
連結: \href{https://youtube.com/watch?v=HUos2vHhz9s}{\texttt{https://youtube.com/watch?v=HUos2vHhz9s}} ~~~~ 語音日期: 2023-01-05
\newline
\newline
\hyperref[sec:1SOWC0lwIU0]{\small{< < < PREV SERMON < < <}}
~
\hyperref[sec:index_chronic]{\small{[返順時目]}}
~
\hyperref[sec:WHVDHKq51Y4]{\small{> > > NEXT SERMON > > >}}
\newline
\newline
約翰福音 1:43-51
\newline
\begin{longtable}{cl}
\hline
\hline
章節 & 經文 (和合本修訂版)\\
\hline
1:43 & \begin{tabularx}{0.7\textwidth}{X} 又過了一天，耶穌想要往加利利去。他找到腓力，就對他說：「來跟從我！」 \end{tabularx} \\ \\ \relax
1:44 & \begin{tabularx}{0.7\textwidth}{X} 這腓力是伯賽大人，是安得烈和彼得的同鄉。 \end{tabularx} \\ \\ \relax
1:45 & \begin{tabularx}{0.7\textwidth}{X} 腓力找到拿但業，對他說：「摩西在律法書上所寫的，和眾先知所記的那一位，我們遇見了，就是約瑟的兒子拿撒勒人耶穌。」 \end{tabularx} \\ \\ \relax
1:46 & \begin{tabularx}{0.7\textwidth}{X} 拿但業對他說：「拿撒勒還能出甚麼好的嗎？」腓力說：「你來看。」 \end{tabularx} \\ \\ \relax
1:47 & \begin{tabularx}{0.7\textwidth}{X} 耶穌看見拿但業向他走來，就論到他說：「看哪，這真是個以色列人！他心裡是沒有詭詐的。」 \end{tabularx} \\ \\ \relax
1:48 & \begin{tabularx}{0.7\textwidth}{X} 拿但業對耶穌說：「你從哪裡認識我的？」耶穌回答他說：「腓力還沒有呼喚你，你在無花果樹底下，我就看見你了。」 \end{tabularx} \\ \\ \relax
1:49 & \begin{tabularx}{0.7\textwidth}{X} 拿但業回答他：「拉比！你是神的兒子，你是以色列的王。」 \end{tabularx} \\ \\ \relax
1:50 & \begin{tabularx}{0.7\textwidth}{X} 耶穌回答他說：「因為我說在無花果樹底下看見你，你就信嗎？你將看見比這些更大的事呢！」 \end{tabularx} \\ \\ \relax
1:51 & \begin{tabularx}{0.7\textwidth}{X} 他又說：「我實實在在地告訴你們，你們將要看見天開了，神的使者在人子身上，上去下來。」 \end{tabularx} \\ \\
[1ex]
\hline
\hline
\end{longtable}
$^{1}$(教育)今天要讀的經文是記載在.
約翰福音一章四十三至五十一節.
我們可以翻開程序表的內頁.
又或者是教會的聖經新約一百二十六至一百二十七頁.
約翰福音一章四十三至五十一節.
「有次日,耶穌想要往加里尼去,.
於見腓力就對他說:來,跟從我吧!.
這腓力是伯塞大人,和安德烈彼得同成..
腓力找著拿達葉,對他說:.
莫西在律法上所寫的,和眾先知所記的那一位,.
我們遇見了,就是約瑟的兒子拿瑟勒人耶穌..
但拿達葉對他說:拿瑟勒還能出什麼好的魔?.
腓力說:你來看!.
耶穌看見拿達葉來,就指著他說:.
看啊!這是個真以色列人,他心裡是沒有鬼的..
拿達葉對耶穌說:你從哪裡知道我?.
耶穌回答說:腓力還沒有招呼你,.
你在無花果樹底下,我就看見你了..
拿達葉說:拉比,你是神的兒子,你是以色列的王..
耶穌對他說:因為我說,我在無花果樹底下看見你,你就信嗎?.
你將要看見比這更大的事..
又說:我實實在在地告訴你們,你們將要看見天開了,.
神的使者上去,下來,在人子身上..
聖經讀筆,我們繼續以聆聽聖道敬拜神..
(主持人:各位姐妹,新年快樂!).
(各位姐妹:新年快樂!).
我們過了一個比較艱難的2022年,.
我們過了一個更艱難的2021年,.
我們現在迎向2023年,.
我們立志要在2023年,在香港這個地方,.
我們一起去見到上帝奇妙的作為,.
我們會在這裡經歷祂,服侍祂..
我去年都出了好幾次場門,每次去不同的地方,.
有人聽到我說,我是香港來的,都用一個比較憐憫的眼神看著我,.
問怎樣怎樣,我要很努力地用微笑去回應他們,.
告訴他們,我們不是被遺棄在這個地方,.
我們真的很愛很愛這個地方,.
香港從來不是一個完美,不是一個好地方..
我記得有一個台灣的領隊,曾經嘲笑香港人的居住環境的惡劣,.
他說他每一次帶團來香港,都會將團友,.

$^{41}$專登帶他們去走黃大仙樂府一帶,.
然後叫他們去猜有多少個窗是屬於一戶人,.
當然答案一定是嚇他們一跳,.
因為每一個窗戶就等於一戶人..
香港人居住環境的狹窄是舉世知名的,.
我們上一輩生活在一個很艱苦的環境裡面,.
我不知道我們中間有多少人在屋村長大,.
我不是在屋村長大,是因為我無幸在屋村長大,.
我由我出生到十八歲,我都是住在一間五百呎的唐樓,.
建了六個單位,板間房,.
我住在一家七口住在一個幾十呎的洗澡房改建的板間房,.
我是到十八歲才有幸升級上到公共屋村,.
所以我們都是艱難的日子,.
不過我們回望我們的過去,.
我們看到高高低低的日子裡面,.
上帝一直恩眷我們,.
我從來沒有像現在那樣為香港人自豪,.
為我們的激情,為我們的勇敢,.
為我們的沮喪,為我們的膽怯,.
為我們的挫折,為我們的奮鬥,.
我引以為榮..
2023年開始,我們先講照命,.
不過你知道今天本來不是我講的,我是吃雞的,.
今天本來應該是陳明全牧師,主任牧師,.
和大家講教會新一年的方向,.
不過因為他病了,我臨時替上,.
所以我講照命,我就不會講群體性的,.
我會講個人性的,整個教會的大方向就留給我們的總舵主,.
等他在下個星期,或者是農曆新年第一個星期,.
就是年初一,等他講吧,.
我們繼續看《約翰福音》,.
我們看耶穌呼召門徒的第二部分,.
剛才讀的聖經記載了耶穌呼召那丹葉的故事,.
我們在這個故事裡面講照命,.
我們講照命不是講偉大的東西,.
不是講摩西在荷列山上見到燃燒的荊棘,.
不是講以賽亞在祈禱的時候見到坐在寶座上的上帝,.
這些照命太崇高,太偉大,.
最少我很抱歉地講,我講不到的,因為我沒有經歷過,.
我順處視同幾十年,我未見過燃燒的荊棘這麼厲害的東西,.

$^{81}$如果你看過的話,應該出來講,.
我見到的都是微思細眼的東西,.
不過那丹葉的故事就是一個微思細眼的故事,.
耶穌呼召那丹葉的故事有些婉言曲折,.
聖經首先記載耶穌呼召肥力,.
肥力和彼得一樣都是白菜大人,.
耶穌在傳道初期,由於他家鄉拿撒勒不歡迎他,.
所以他很多時間都停留在白菜大,.
在加里尼湖邊的一個小村落,一個小鎮,.
這個時候耶穌還沒有到處周遊報道,.
不少他所呼召的第一批門徒仍然去Day Camp,.
不是去Full Camp,不是Boarding,.
每天回家,不是跟著耶穌到處去,.
是白天跟著他,晚上回家,.
甚至白天都不是時時跟著耶穌,.
是Part Time(兼職)而不是Full Time(全職).
他們還會繼續做他們自己原來的工作,.
繼續營生,繼續耕田,繼續打魚,.
這是早期的情況,.
在跟隨耶穌之後,.
馬上就去找他的一個閨中密友,一個發小,.
就是他的好朋友拿丹葉,.
告訴拿丹葉這個消息,.
他說我們蘇常盼望的,.
也是舊約的眾先知所預言的那個彌賽亞,.
我們今天終於見到他了,.
他就是拿撒勒人耶穌,.
拿丹葉聽到這番說話,反應很冷淡,.
隨便回應一句,.
拿撒勒這個窮鄉僻壤,.
怎會產生這些好東西來呢?.
山東街,山東街有什麼好東西?.
山東街只見到人們扔油,.
見到人們賣粉,.
山東街有好東西嗎?.
不會的,他們見不到我們有旺角浸信會,.
所以見不到好東西,.
見到都是,哇,這樣的地方,.
怎可能出到好東西呢?對不對?.
這個耶穌如果是拿撒勒人,.

$^{121}$大概他不可能是彌賽亞,.
肥力是一個很會傳福音的人,.
他沒有和拿丹葉做口舌之爭,.
糾纏於拿撒勒能不能夠出產彌賽亞這個神學問題,.
他只是簡單說一句,.
你自己來看,你自己來看,.
真的,沒有人聽到別人說有關耶穌的故事,.
就貿然願意跟隨主的,.
這個故事太奇怪,.
耶穌出生太貧寒,.
道成肉身的故事太放誕,.
就這樣聽,你不會接受的,.
唯有是他親自去到耶穌的面前,.
見過祂,聽過祂,摸著祂,.
亦都應該說被耶穌動過祂,.
他才會向耶穌折服,.
承認耶穌果然是上帝的意志,.
所以傳福音的人最大的使命不是在理論層面去搏倒對方,.
而是將人帶去耶穌的面前,.
讓耶穌親自動他,親自改變他,.
我不是想用我的言語去說服你,.
我是想你見到耶穌,.
你見到耶穌就什麼都可以搞定,.
我不用說這麼多,.
費力帶著拿丹葉去找耶穌,.
因為北塞大是很小的地方,.
耶穌離他們不遠,.
很快就找到耶穌,.
耶穌遙遠看到他們,.
就大聲指著拿丹葉說,.
看,這是一個頂天立地的以色列人,.
他的心地純樸,.
沒有鬼詐,.
沒有什麼陰暗面,.
我們可以相信拿丹葉真是一個心地善良,.
人格正直的人,.
耶穌沒理由看錯人,.
耶穌不會擄作的,.
耶穌說他是好人,.
一定是好人,.

$^{161}$我們也不需要浪漫化地說,.
拿丹葉是一個性格完全毫無瑕疵的人,.
如果是一個完美的人就不用相信耶穌,.
他已經上天堂了,.
沒有一個這樣的人,.
拿丹葉應該是有好有壞的,.
不過可能性格的優點多於缺點,.
或者六四,.
或者七三,.
或者八二,.
不知道,.
不無論何,.
最少你見到的是,.
耶穌刻意去點出拿丹葉的優點,.
而不是缺點,.
是指出他性格上有可誇之處,.
公開去讚美他,.
耶穌沒有捉住拿丹葉的缺點來對他說,.
沒有指責你是罪人,.
你必須盡快悔改,.
否則你會落地獄,.
不是,.
耶穌只是說他的優點,.
我兒子,.
我記得他在中學階段尾,.
中五,中六,中七的時候,.
他曾經想過想去讀拜柯,.
為什麼想讀拜柯呢?.
他後來沒有,.
因為我也知道拜柯根本不適合他,.
完全不是他的性格,.
只不過因為拜柯阿Sir很疼他,.
他很喜歡那個拜柯阿Sir,.
很多人是這樣的,.
我喜歡那個先生,.
或者那個先生喜歡我,.
我就喜歡他那科,.
我就自然想跟他那科讀,.
是這樣的,.
你看得起我,.

$^{201}$你喜歡我,.
你讚賞我,.
我就會為你拼命,.
我就會用心讀好這一科,.
通常我們人就是這樣,.
肯定人,.
找出人生命裡面的優點,.
然後到位的去讚賞他,.
他會知道他是誰,.
並且他會繼續奮鬥,.
繼續努力,.
弟兄姊妹,.
不要說我們去讚賞別人,.
我們就算去看我們自己也是這樣,.
多一點從自己有的,.
而不是沒有的,.
從未有的角度,.
去看上帝我們個人生命裡面的心意,.
很多時候我們想我們的未來,.
想我們要做什麼,.
總是想一些我們未做到,.
甚至可能做不到的事,.
見到別人上太空,.
我很希望上太空,.
見到李麗珊,.
我很想游得快過她,.
我羨慕別人有,.
而我沒有的事,.
我期望做到她,.
翻天覆地,.
改變自己,.
做到她,.
但是,.
想一想,.
有什麼是我有的,.
有什麼是上帝給了我的,.
上帝在我們生命的劇本已經開展了,.
那問,.
如果今天耶穌要來讚我,.
祂會讚我什麼呢?.

$^{241}$你不要以為這樣很自誇,.
不是的,.
保羅常常說祂只是誇口祂的軟弱,.
但是,.
他在提摩太前書1章12節是這麼說的,.
基督耶穌,.
因祂以我有忠心,.
派我事奉祂,.
保羅是不是口面屁,.
不是,.
他是真誠的說,.
我,.
至少我是很忠心我的主,.
所以我的主因著這樣的緣故事奉我,.
上帝,.
我常常說,.
不會為一件垃圾而死的,.
祂為我死,.
證明我是有價值的人,.
我一定是有價值的人,.
我有什麼價值呢?.
我不常常看見,.
我求主開我的眼叫我見到,.
那丹葉對耶穌讚美的說話非常震動,.
本來他對耶穌還有一些狂傲,.
有一些對立的態度,.
想挑戰耶穌,.
否定耶穌,.
那丹不會出好東西,.
所以耶穌不是好東西,.
但一下子,.
你見到他,.
馬上謙卑下來,.
在一個認識你底蘊的人面前,.
你最好謙卑一點,.
對方認識你,.
你不認識他,.
你知不知道你輸了,.
你很害怕,.
他認識你,.

$^{281}$你不認識他,.
他就問耶穌,.
你怎麼認識我?.
耶穌再大他,.
耶穌回答,.
肥力還沒招呼你,.
你在無花果樹底下,.
我就看見你了,.
耶穌不是在肥力帶那丹葉來找他的時候,.
才看到那丹葉,.
他卻在肥力還沒找那丹葉,.
沒有邀請他見耶穌之前,.
耶穌已經看到他,.
耶穌特別提到那丹葉,.
你那時候是否站在無花果樹底下呢?.
證明我真的看到你,.
早在那丹葉從肥力口中認識有耶穌這個人之前,.
耶穌已經認識那丹葉,.
實際上我們知道,.
耶穌不是在肥力找那丹葉之前,.
那一剎那才認識那丹葉,.
他是早在那丹葉出生之前就已經認識他,.
耶穌認識那丹葉,.
也認識我們每一個,.
早在我們認識他之前,.
耶穌已經認識我們,.
早在我們看見耶穌之前,.
耶穌已經看見我們,.
我們一直活在他的目光之下,.
從來沒有離開過,.
我以前寫過一本小書叫《信定一生》,.
今天講的不是新的東西,.
我在十八歲那一年成為基督徒,.
在那一年我承認上帝是上帝,.
但上帝何時是我的上帝,.
何時認識我,.
他從來都是我的上帝,.
從來沒有一刻不是,.
他一直認識我,.
我的基督徒生涯在十八歲那一年開始,.

$^{321}$但我的救贖故事,.
早在我的骨頭還未形成之前就已經開始,.
我在未信主前的十八年,.
我在人間飆盪的十八年,.
我經歷各種挫折,.
困乏的時候,.
他已經是我的上帝,.
他一直是我的上帝,.
那丹爾被耶穌基督的認識徹底震撼,.
他立即在耶穌面前伏手稱臣,.
拉比,你是上帝的兒子,.
你是以色列的王,.
他本來是要來看看耶穌是誰,.
認識耶穌,.
發現耶穌,.
判定耶穌是不是肥力所說的救約所允許的彌賽亞,.
但如今,.
他第一個發現是甚麼,.
他發現耶穌,.
發現他,.
他認識甚麼,.
他認識耶穌,.
他認識他,.
耶穌認識他,.
比他認識耶穌更重要,.
這是因為被發現而帶來的震撼,.
信仰除了幫助我們認識上帝,.
知道他為我們做了甚麼之外,.
也更重要的是幫我們認識我們自己,.
我們在上帝眼裡,.
發現真正的自己,.
就是被上帝創造,.
被他救贖,.
被他作為一席無數的自己,.
也是要開展的救贖故事裡的第二主角,.
我的救贖故事的第一主角當然是耶穌,.
是耶穌和我,.
以馬來尼,.
上帝和我的故事,.
很多時候,.

$^{361}$我們對自己有很多假象,.
特別是我們對眼前的我們自己,.
很多不滿意的時候,.
我們總是閉上眼睛,.
幻想,.
預測我是怎樣怎樣的人,.
我經常和人格價,.
和人比較,.
希望我可以做到別人的樣子,.
希望可以實現某一個境界,.
然後我們東奔西跑,.
就是蠻頭蒼蠅到處去嘗試,.
直到我們碰了很多釘,.
筋疲力盡之後,.
我們才乖乖安頓下來去問上帝,.
到底在你眼中,.
我是一個怎樣的現實,.
我有一個怎樣的將來,.
我常常想成為somebody,.
常常想do something,.
但是上帝告訴我,.
我要成為他創造計劃裡面的我,.
Be yourself,.
當然不是如今爛爛的self,.
而是未來最大可能的self,.
師父需要的第一照明,.
永遠不是宏大的外在需要,.
而是上帝在我們生命裡面的計劃,.
常常問客觀的需要其實沒有意思,.
這個世界當然有需要,.
好像很厲害,.
我可以做些什麼,.
我可以幫這個世界什麼,.
不用問,.
每樣都需要,.
二百個國家,.
每個國家都需要,.
柬埔寨有沒有需要,.
菲律賓有沒有需要,.
泰國有沒有需要,.

$^{401}$緬甸有沒有需要,.
不用問,.
每個國家都有需要,.
沒有一個世界的地方是以德之地,.
全部都不是,.
歐美都不是以德之地,.
很多人不相信耶穌,.
很多人放棄信仰,.
對不對?.
要問旺朕有什麼需要,.
兒童侍工有沒有需要,.
青少年侍工有沒有需要,.
成人侍工有沒有需要,.
省一點,.
不用問,.
每樣都有需要,.
只是問需要,.
然後我可以怎麼做,.
我可以怎麼做,.
一方面叫你自己有個假象,.
以為你很偉大,.
另一方面,.
實際上你做的事情都很少,.
問一問,.
上帝預備一個怎樣的你,.
你是怎樣的人,.
正如我以前所說,.
不用問誰是我的鄰舍,.
所有有需要的人,.
都是我潛在的鄰舍,.
我要問的是,.
我是誰的鄰舍,.
先講感情上,.
什麼是我情之所終,.
心之所繫,.
我見到我感動,.
我見到我很想抱他的那個,.
那個就是我會成為他的鄰舍,.
但同時,.
更重要的是,.

$^{441}$我要問我自己,.
我可以成為一個怎樣的事奉者,.
正如我從來沒想過,.
我會做婦女工作,.
限電,.
頻道都完全不對,.
根本上沒得談,.
都不知說什麼好,.
另外,.
我們都已經放棄做精緻人工很多年,.
那個generation gap(世代差距)這麼大,.
不用說了,.
不會做,.
做不到,.
我比較有能力做的,.
是做一些中坑老坑的工作,.
那個跟我對到嘴型,.
我們大家很容易,.
不用扮,.
就已經可以合得來,.
我有什麼,.
我是誰,.
我是一個怎樣的事奉者,.
我就做什麼事奉,.
我在《信徒一生》的書中,.
也用過,.
以法所書四章七節這樣講,.
我們各人蒙恩,.
都是照基督量級各人的恩賜,.
你要注意這個字,.
這個恩賜這個字,.
指的不是那個gift(禮物),.
不是我們常用的那個gift,.
是在講Christ apportioned(基督賜予的),.
基督派定的,.
基督給的,.
那些就叫做恩賜,.
恩賜不是指我們某些知識,.
某些技能,.
我會彈琴,.

$^{481}$我會唱歌,.
我有某些技巧,.
我們沒有人純粹用這些知識技能技巧去事奉,.
他攤派他整定的東西,.
我成長的背景,.
我過去的經歷,.
快樂的不快樂的,.
成功的失敗的,.
受傷的醫治的,.
被建立的被剝奪的,.
我所有的人際關係,.
各種悲歡離合,.
如果我用信心的眼睛,.
看見這些通通是上帝給我的恩典,.
既然是恩典,.
就可以轉化成為我事奉的資本,.
我這樣跌倒過,.
我這樣病過,.
我好知道跌倒的人有什麼感受,.
我好知道那些有病的人有什麼需要,.
我用我的成功去事奉,.
我用我的失敗去事奉,.
我用我的榮耀去事奉,.
我用我的傷痕去事奉,.
你回頭以後,.
要兼顧你的弟兄,.
這是耶穌對彼得的吩咐,.
在你跌倒之後,.
你用你跌倒的經歷去兼顧你身邊的人,.
最重要的是,.
我們用信仰的眼睛,.
去看我這一生的救贖歷史,.
確認上帝在我生命裡已經有的作為,.
不是個別事件,.
我求這件事,.
靈驗,.
得到,.
求不到,.
就動它,.
不妥它,.

$^{521}$不是那些事,.
是面現的故事,.
我們都是上帝的工作,.
在基督裡面完成,.
我將我過去幾年在香港的經歷,.
加在我的香港人故事裡,.
加在我的救贖故事裡,.
我立志我在這裡,.
我會見到我的主,.
我會服侍我的主,.
恩典是藉著回顧而確定,.
弟兄姊妹,.
你認不認識你自己?.
你認不認識上帝眼裡的你自己?.
你看不看到上帝在你生命裡的心意?.
找到你現在的位置,.
先確定你現在要做的事,.
這是最重要的,.
我經常說這些很老套的話,.
重要的事要說三萬次,.
做個好樣的老公,.
做個好樣的老婆,.
做個好樣的父母,.
做個好樣的子女,.
先在我那些不能推,.
上帝派給我的關係位置裡,.
我先做好它,.
努力做好它,.
我先成為我身邊的人的和平之子,.
我給它四萬個樣子,.
讓它開心,.
讓它覺得被光拜,.
先做好我眼前上帝派給我的東西,.
然後你會見到上帝有更大的東西,.
更多的可能在你生命裡出現,.
重要的不是在於我見到我的可能性,.
而是我在我的生命裡見到耶穌,.
意象從來不是某些外在的意象,.
而是上帝在我生命裡的意象,.
我很喜歡那首古老的歌,.

$^{561}$"Be Thou My Vision",.
讓你成為我的意象,.
我見到我的主,.
我就見到我將來的可能性,.
我的主和我在一起,.
我很清楚我可以做些什麼,.
當我一心想著我很愛我的耶穌,.
我很想擁抱祂,.
我很想榮耀祂,.
很想哄祂開心,.
我從這個角度去想,.
我就很清楚我面前可以做些什麼,.
當我認定我要榮耀我的主,.
我很愛很愛很愛我的主,.
很想很想哄祂開心的時候,.
我就看到我有很多很多不同的可能性,.
那丹妮還沒有認識耶穌,.
因為耶穌根本沒有告訴她其他東西,.
沒有說過什麼基督教道理,.
她只知道耶穌充分認識她,.
她也沒有問耶穌是誰,.
但知道耶穌認識她之後,.
她就確定耶穌是彌賽亞,.
拉比,你是上帝的兒子,.
你是以色列的王,.
耶穌怎麼回答,.
耶穌再給她一個挑戰,.
因為我說在無花果樹下見到你,.
你就相信了,.
這麼容易,.
不是這麼簡單的,.
你將要看見比這個更大的事,.
她接著說,.
我實實在在的告訴你們,.
你們要看見天開了,.
上帝的使者上去下來在人子身上,.
哇,.
那丹妮看到耶穌看見她就很震撼,.
她因此就相信了耶穌,.
但耶穌告訴她,.

$^{601}$就這樣相信我並不重要,.
不是最重要,.
你會看見更大的事,.
什麼事?.
不是一般自然的事,.
是有關超自然的事,.
天開了,.
自然和超自然的界線吻隨了,.
好像有一道天梯,.
踏落人間,.
踏在耶穌身上,.
好像亞國在伯特利所見到的異象,.
上帝的使者上去下來,.
你記不記得,.
異象是什麼?.
上帝的旨意會被頒佈,.
上帝的旨意會被實現,.
上帝的工作會成就,.
天使是派落人間去成就上帝的工作,.
耶穌預告,.
你會看見一條天梯,.
從天上跌落到耶穌身上,.
上帝的旨意,.
上帝的作為,.
會一直一直成就在耶穌基督身上,.
而那個耶穌是在拿丹葉眼前的耶穌,.
在我們每一個人生命裡的耶穌,.
兄弟姐妹,.
這是我們大家的期望,.
這是我們大家的應許,.
我們個人的生命裡,.
如果耶穌在我們生命中,.
你會看見上帝更多奇妙的作為,.
會在你生命裡的耶穌實現,.
所以你永遠不是看見你自己的作為,.
好像很厲害,.
我在我生命裡,.
我不斷經歷耶穌基督奇妙的作為,.
祂藉著我去彰顯祂的恩賜,.
祂的慈愛,.

$^{641}$不是我,不是我,.
同樣在我們皇寺也一樣,.
我們會看見耶穌基督在我們中間,.
天開了,.
上帝的使者會上去下來,.
在我們皇寺之中,.
耶穌基督會彰顯祂奇妙的作為,.
耶穌基督會彰顯祂更大的榮耀,.
弟兄姊妹,.
求主奪回我們的心意,.
使我們可以徹底的信服基督,.
這是保羅在哥林多後書十章五節所講的,.
今天我們要講的其實很簡單,.
弟兄姊妹,.
我們要掌握上帝給我們的召命,.
那個召命在我們專注耶穌,.
專注現在開始,.
我的祈禱是,.
好像紹約翰在啟示錄對伊弗所教會的責備,.
這間教會樣樣都good perfect,.
既認定真理,.
踐行真理,.
抵抗異端,.
又忠心又勞苦,.
努力事奉,.
不過他們只需一樣東西,.
就是他們丟棄了起初的愛心,.
結果耶穌藉著約翰口講一個怎樣的診斷,.
你如果不馬上改變,.
我是會撤去你的金燈台,.
即是說這是致命的缺失,.
就這件事可以使他由教會變成不是教會,.
我希望我們弟兄姊妹,.
大家在面前一年,.
我們立志更加信基督教,.
更加堅持我們的信仰,.
更加努力去實踐我們的信仰,.
當然這些都重要,.
最重要的是我們立志更加愛耶穌基督,.
愛祂,.

$^{681}$切實地愛祂,.
在我們的生命裡發現祂,.
經歷祂,.
只要我們執著我們對耶穌的愛,.
耶穌對我們的愛,.
你會看見在面前一年,.
我們會經歷更多上帝奇妙的作為,.
上帝的召命會在我們生命裡得到充分的兌現,.
上帝恩典我們每一個,新年快樂!.
\newpage



\section{馬太福音 4:18-25}
\label{sec:WHVDHKq51Y4}
\textbf{2023年1月8日崇拜直播｜陳明泉牧師｜給當代門徒的15個忠告:真正的追隨｜馬太福音4\_18-25}
\newline
\newline
連結: \href{https://youtube.com/watch?v=WHVDHKq51Y4}{\texttt{https://youtube.com/watch?v=WHVDHKq51Y4}} ~~~~ 語音日期: 2023-01-09
\newline
\newline
\hyperref[sec:HUos2vHhz9s]{\small{< < < PREV SERMON < < <}}
~
\hyperref[sec:index_chronic]{\small{[返順時目]}}
~
\hyperref[sec:Z3AMKkvW_mc]{\small{> > > NEXT SERMON > > >}}
\newline
\newline
馬太福音 4:18-25
\newline
\begin{longtable}{cl}
\hline
\hline
章節 & 經文 (和合本修訂版)\\
\hline
4:18 & \begin{tabularx}{0.7\textwidth}{X} 耶穌沿著加利利海邊行走，看見兩兄弟，就是那叫彼得的西門和他弟弟安得烈，正往海裡撒網；他們本是打魚的。 \end{tabularx} \\ \\ \relax
4:19 & \begin{tabularx}{0.7\textwidth}{X} 耶穌對他們說：「來跟從我，我要叫你們得人如得魚一樣。」 \end{tabularx} \\ \\ \relax
4:20 & \begin{tabularx}{0.7\textwidth}{X} 他們立刻捨了網，跟從他。 \end{tabularx} \\ \\ \relax
4:21 & \begin{tabularx}{0.7\textwidth}{X} 耶穌從那裡往前走，看見另外兩兄弟，就是西庇太的兒子雅各和他弟弟約翰，同他們的父親西庇太在船上補網，耶穌就呼召他們。 \end{tabularx} \\ \\ \relax
4:22 & \begin{tabularx}{0.7\textwidth}{X} 他們立刻捨了船，辭別父親，跟從了耶穌。 \end{tabularx} \\ \\ \relax
4:23 & \begin{tabularx}{0.7\textwidth}{X} 耶穌走遍加利利，在各會堂裡教導人，宣講天國的福音，醫治百姓各樣的疾病。 \end{tabularx} \\ \\ \relax
4:24 & \begin{tabularx}{0.7\textwidth}{X} 他的名聲傳遍了敘利亞。那裡的人把一切病人，就是有各樣疾病和疼痛的、被鬼附的、癲癇的、癱瘓的，都帶了來，耶穌就治好了他們。 \end{tabularx} \\ \\ \relax
4:25 & \begin{tabularx}{0.7\textwidth}{X} 當時，有一大群人從加利利、低加坡里、耶路撒冷、猶太、約旦河的東邊，來跟從他。 \end{tabularx} \\ \\
[1ex]
\hline
\hline
\end{longtable}
$^{1}$今日我們所選的經訓.
在新約的《馬太福音》.
四章十八節至二十五節.
是教會的聖經第四至五頁.
聖經是這樣說的.
耶穌在加利利海邊行走.
看見弟兄二人.
就是那稱呼彼得的西門和他兄弟安德烈.
在海裡撒網.
他們本是打魚的.
耶穌對他們說.
來 跟從我.
我要叫你們得人如得魚一樣.
他們就立刻捨了網.
跟從了他.
從那裡往前走.
又看見弟兄二人.
就是西比太的兒子雅各和他兄弟約翰.
和他們的父親西比太在船上保網.
耶穌就招呼他們.
他們立刻捨了船.
別了父親.
跟從了耶穌.
耶穌走遍加利利.
在各會堂裡教訓人.
傳天國的福音.
醫治百姓各樣的病症.
他的名聲就傳遍了敘利亞.
那裡的人把一切害病的.
就是害各樣疾病.
各樣疼痛的.
和被鬼附的.
癲癇的.
癱瘓的.
都帶了來.
耶穌就治好了他們.
當下有許多人.
從加利利.
底加波利.
耶路撒冷.

$^{41}$猶太.
約旦河外來跟著他.
弟兄姊妹平安.
平安.
之前中了COVID.
還有一點鼻子有些嗡聲.
有些爛聲.
大家多多忍耐.
希望今天的聲音可以令我不會咳嗽.
可以完成今天的宣講.
在接下來我們踏入新的一年.
2023年.
我們教會的主題就叫做門徒不斷.
或者門訓不斷.
你訓練就有門徒了.
你不訓練就不會有的.
所以.
門訓也好.
或者門徒也好.
這個生生不息的門徒生活.
是經歷過教會裡面.
刻意訓練.
以致讓弟兄姊妹可以起來.
成為上帝所使用的器皿.
不過說起門徒訓練或者門訓.
其實坊間有很多書.
不同的書.
或者不同教會在做的門徒訓練.
大家都是同樣的字眼.
但是內涵或者大家所看的東西可以很不同.
有些教會藉著門徒訓練.
其實就是揀卒.
希望可以將新一代的領袖.
藉著一些訓練.
一些層次比較高一點的.
或者不同的密集式的訓練.
希望可以讓一些人可以承接教會的棒.
這是一個將教會的領袖權交給一些門徒.
所以這種訓練某程度上比較隱蔽.
因為看到這些弟兄姊妹.

$^{81}$但是否真的可以將來成為教會的領袖.
都要觀察一段時間.
如果一下子公開了.
如果她還未去到那個位置就很尷尬了.
所以就有一些很隱蔽,神神秘秘的.
來做一些秘密訓練.
希望揀卒來建立一些新一代的門徒.
這是一些想法.
第二就是一般教會聚會的內容.
已經有一定的熟練.
有些弟兄姊妹追求的就是更加高層次.
或者更加有深度的分享.
所以覺得門徒訓練是某程度上.
是一個很追求性的訓練.
希望可以將更加高階,更加深奧.
或者更加看得遠一點的內容.
透過這些訓練來給予一些弟兄姊妹.
但是未必一定和你的崗位.
或者未來教會的目的有些掛勾.
只不過是令弟兄姊妹的滿足感更加大.
這是第二種說法.
第三種就是有一些使徒濟做多一些的.
教會可能有些事工.
如果沒有人去接棒.
可能短時間就會沒有了.
所以就說怎樣去做一件事.
怎樣去實踐事工的內容.
透過使徒濟.
一隻撚隻帶一個.
來將這個事工實踐出來.
而且還可以確保這個事工可以.
源源不盡地有人來做.
坊間說的青少年事工.
或者全福音,三福等等這一類的事工.
就是用這個角度.
就是落手落腳.
怎樣去做都是有章法的系統來做.
五花八門大家聽下去都知道.
原來對於門徒的看法.
有些可能門徒等於領袖.

$^{121}$有些可能門徒等於前線的事奉者.
又或者有一些門徒代表著.
一些對聖經或者對神學.
有深入理解的一些弟兄姊妹.
在這種想法裡面.
有意無意之間都塑造了.
弟兄姊妹原來有一些比較高階一點.
有一些就比較顯淺一點.
或者有些叫做「平信徒」.
即是在說不同的層次.
不過如果在《旺朕》裡面我們看.
我們不相信.
聖經裡面沒有特別講到.
將信徒劃分了不同的層次.
當我們真心信主的時候.
其實在聖經裡面講到真心相信主的人.
已經是進入門徒訓練的門檻.
原則上在《旺朕》.
或者在我們自己的理解裡面.
其實我們的牧羊和門徒訓練是掛勾的.
我們的牧羊其實都是門徒訓練的一部分.
而且是一個很重要的那一環.
我們也沒有刻意要將一些信徒.
變成一個高階信徒.
我們有機會可能會培訓多一些的領袖.
但這些領袖是僕人式的領袖.
是承擔更加大.
而不是掌握一些更加深奧的道理.
或者掌握一個其他人不知的信息.
來讓他成為一個秘密的領袖.
反而在教會裡面我們希望.
每一個都成為門徒.
領袖可能是一小撮人.
但門徒就是所有弟兄姊妹.
當他融入了在教會裡面的牧羊.
他就成為門徒了.
我們帶著這個角度.
在接下來這十多次的講論裡面.
我們每一個都是被塑造者.
被上帝的話語.

$^{161}$被上帝的真理來塑造.
我們每一個都是在走門徒的路.
接下來的講論篇幅.
我們集中會用《馬太福音》裡面的登山寶訓.
以他們的內容.
作為一個門徒可以的生命怎樣呈現出來呢?.
每一次都有不同的面向.
今日我所講的這段經文.
其實就不是登山寶訓的內容.
而是登山寶訓之前.
耶穌對著門徒和群眾.
耶穌講的對象是甚麼人呢?.
今日我們會處理這段經文.
不過我希望我們今日.
當我們去理解門徒訓練.
或者我們每一個被塑造的時候.
不要覺得是隔壁那個人.
或者那個熱心的弟兄姊妹.
這個訊息或者未來這半年我們所看的經文.
其實是關乎我們自己的屬靈生命.
是關乎我們怎樣被上帝.
來打造這個世代上帝所使用的一班人.
求主繼續來幫助我們.
在我進入經文之前.
我想講一個不知道大家有沒有.
一些這樣的經歷.
一些我自己家裡的香港情懷.
話說我家裡我爸爸是做製衣的.
我爸爸做製衣的時候.
他偷渡下來香港.
所以穿著一條游水褲來到香港.
謀生的時候要找一門手藝.
於是他就進入製衣行.
走去那些廠裡面做學師.
由最簡單的裁剪,摺衫,剪線頭.
一直學到他做到一件衣服.
然後看到別人怎樣做就接生意.
然後大概六七十年代.
他就自己開了一間山寨廠.
不知道大家有沒有聽過山寨廠.

$^{201}$就是這樣做的了.
就是有一個山寨廠.
我就見證不到我爸爸開山寨廠的創業過程.
不過我見到我爸爸收徒弟的過程.
話說有一些家裡比較遠房的親戚.
有一個住在香港仔.
有一個住在長洲那裡.
他們那時候那些地方比較窮鄉僻壤.
或者出去城市都不知道找些什麼一技之長.
十多歲的年輕人.
於是乎就透過親戚所託.
就去到我家找我爸爸來拜師學藝,學製衣.
我見到那個情景很有趣的.
就好像看電影一樣.
見到他來拜師的時候要給我爸爸倒茶.
接過這個茶之後我爸爸給他封利是.
就成為了我爸爸的徒弟了.
他叫我爸爸陳生.
不會叫師父這麼難看的.
他們都是叫陳生.
不過我見到的那個狀況.
去到現在都很難忘.
那個師父就是我爸爸.
出去接貨也好.
或者做工趕貨也好.
這些徒弟仔都一定跟著他的.
十多歲的年輕人.
其實有時候不識天高地厚.
有時候會撞板的做事.
做著做著做錯了就被我爸爸罵.
但是有時候做了一些事很幫得上忙.
我爸爸就會搭著肩膀.
稱讚他做得好.
我很記得有一次.
有一個年輕人.
受不了被我爸爸罵.
發脾氣.
就扔掉電話.
那時候用那些牛龜電話.
黑色的.

$^{241}$發脾氣.
吸線打電話.
我爸爸就氣得飛起來.
然後走去錄影廠找他.
然後害怕.
躲在衣服那裡.
害怕被師父罵.
被師父罰.
他們十多歲的年輕人.
這樣來學師.
由最基本一步一步.
由一個窮鄉長洲.
或者在香港仔的漁上.
水上人.
來轉換一條跑道學製衣.
我還記得他們經歷過十幾年之後.
他大我十年八載.
然後我記得有一年的新年.
類似現在這些時間.
農曆新年之前.
我見到有一個爸爸的徒弟.
就跟我爸爸說.
陳先生.
我今年就不做了.
一直當我自己.
好像弟弟一樣.
來看大我.
教我很多東西.
接下來我都要離開.
我希望自己開新的廠.
我記得那個情景.
這個學師的.
我爸爸的徒弟.
一路說的時候.
自己的眼淚都流出來了.
很不捨得.
但是他不想繼續屈在這裡.
他想自己可以獨當一面.
於是乎就說他自己會開廠.
接著他會在屯門那裡開.

$^{281}$不會和我們爭生意.
一路一路這樣.
來扶持他們.
我爸爸聽到這個消息之後.
就給了他一封利是.
接著這個學師就給了他一杯茶.
滿師了,就離開了.
我見到我爸爸眼濕濕.
我爸爸平時很少會這樣.
我見到我爸爸眼濕濕.
我就不惜天高地下.
「老爸你不捨得人家嗎?」.
這樣問他.
我爸爸就真的說.
「真的有些不捨得的」.
「這個徒弟」.
「他是我在香港來到」.
「我隻身穿條牛水褲來」.
「這些是我第一批所收的徒弟」.
「不捨得你不留下他?」.
這樣問他.
我爸爸說.
「學師為了什麼呢?」.
「學師」.
「師父不是要」.
「讓徒弟永遠在自己身邊的」.
「如果真正成功的師父」.
「是有令到自己的徒弟可以獨當一面」.
「雖然我很不捨得」.
「但是我覺得我很滿足」.
「因為我可以令到一個人」.
「可以自己滿師」.
「來獨當一面開創自己的事業」.
我爸爸又借了一句.
「你將來都要是這樣」.
我不知道他說什麼.
不過在我預備這篇講章的時候.
我心裡想.
是的.
今天教會的門徒訓練.

$^{321}$不是要令到那些頂子妹.
逐擁那個師父.
逐擁在一個人身上.
應該要令到每一個人.
在真理的栽培之下.
每一個人都能夠獨當一面.
今天我們看的這段經文.
耶穌都是帶著這一份心腸.
來呼召一些門徒.
希望我們從這段經文中.
找到一些養分.
我們先有一個祈禱.
求主幫助我們.
耶穌求你在我們當中對我們說話.
藉著你怎樣去呼召門徒.
這個訊息同樣在我們心靈裡.
繼續來動功.
遙控我們對你的認知.
改變我們.
可能我們過往有很多對你錯誤的理解.
求你在我們當中繼續對我們說話.
恭敬將來領受你話語的時間交托.
幫助宣講的說得清楚.
幫助每一位聆聽的弟兄姊妹聽得明白.
奉靠基督耶穌得勝的聖名.
Amen.
我們一起看.
《馬太福音》第四章十八節到二十五節.
這個段落是在登山補糞之前的一個技術.
在之前的技術講到耶穌受過試探.
接著耶穌一路一路開始.
受完試探之後就開始.
走遍不同的加利利,拿撒勒的地方.
開始講天國的道.
在第十七節講到.
耶穌就傳起道.
天國近了,你們應當悔改.
耶穌很簡單.
大概和施贈約翰所講的那個訊息是一致的.
不過他講完這一篇道.

$^{361}$講完一個很簡單.
天國近了,你們應當悔改的訊息之後.
耶穌不是選擇以獨行俠的方式.
來完成這個天國召命.
接著在第十八節開始.
就講到耶穌一路走遍加利利.
見到有一些捕魚的兩兄弟.
彼得和安德烈.
耶穌就直接呼召他們.
當他們在船上在游泳.
在休息的時候.
耶穌見到他們兩兄弟.
就叫他們跟從他.
耶穌應許叫他們得人如得魚一樣.
在二十節他們就立刻寫了網.
跟從了他.
十八至二十節其實是一個很靈快的記載.
裡面沒有描述彼得和安德烈.
寫下網的內心掙扎.
在這裡其實很明快.
耶穌很清晰地說.
叫這兩兄弟跟從他.
就好像昔日打魚一樣.
耶穌用一些他所熟悉的.
這兩兄弟所熟悉的語言.
就是打魚.
這個概念套用於在未來的日子.
這兩兄弟將來就好像漁夫一樣.
不過得的不是漁獲.
而是得到人的生命.
當然耶穌說這番話.
是在一個天國福音的狀況下.
得魚的意思不是說這兩個漁夫.
得人就好像做了教主.
然後人們就用心擁戴他.
大概是在一個天國的福音下.
耶穌說這兩個人是被使用的.
他得到的人就好像昔日捕魚得到的漁獲一樣.
有滿足感.
而且在過程裡面.

$^{401}$他充滿了那種踏實.
在經文裡面.
一個很簡單的呼召.
這兩個人立刻寫了網.
跟從他.
經文裡面說到彼得和安德烈.
對於耶穌的呼召其實是很爽快的.
經文怕我們有錯誤的詮釋.
所以在這裡要特別寫到.
是立刻寫了網.
聽完之後立刻做的.
不是嘔很久.
又要祈禱一下.
沒有做這件事.
直接了當.
耶穌呼召.
立刻就寫下了網.
跟從他.
同樣的公式在二十一到第二十二節又再重複.
耶穌往那裡往前走看見另外兩兄弟.
就是西比泰的兒子雅各和弟弟約翰.
這次他還多一個人物.
就是他們的爸爸西比泰.
在船上來保望.
耶穌就呼召他們.
他們就立刻寫了船.
辭別了父親.
跟從了耶穌.
到了二十二節的公式是一樣的.
不過人物關係就複雜了一點.
首先就說到這兩兄弟除了他自己之外.
還有第三者就是西比泰這個爸爸.
這個爸爸都是.
如果你明白當代社會的社會流動性很低.
所以爸爸做什麼兒子就做什麼.
而這次是爸爸都聽到.
或者親歷耶穌對這兩兄弟的呼召.
就是要放棄這個世世代代相傳的行業.
但是放下了之後.
他們是跟從直接呼召他的那位耶穌.

$^{441}$在經文裡面說到他們放下了什麼呢?.
捨棄了以前吃日餬的那艘船.
那隻漁船.
辭別了父親.
跟爸爸說要跟從一個師父.
辭別了爸爸既是代表著.
可能要離開家庭.
同時間也是離開爸爸的祖業.
離開爸爸一直努力的.
教導的從小到大的那個家業.
由這個家業學一門新的做法.
跟隨耶穌傳天國的道理.
當然有很多《釋經》學家就說.
為什麼他們這麼立刻.
來回應耶穌的呼召呢?.
大概可能是猶太人.
心裡面一直有一個.
引頸以待的一個根深蒂固的概念.
他們期待一個新的領袖.
一個彌賽亞.
會繼續帶領他整個民族.
所以他們一直在等.
當他們見到耶穌的時候.
聽到耶穌傳講天國的福音.
引用先知的道理.
於是他們就覺得是他了.
所以就很爽快地跟從.
不過同樣道理.
其實當時有很多先知.
有人跟隨施浸若寒.
又或者有不同的人.
在他們的思路裡面.
在《聖經》裡面其實就保持緘默.
只不過可以肯定的就是.
他回應耶穌的呼召.
這兩對兄弟.
他們都是立刻捨棄眼前的需要.
捨棄眼前謀生的工具.
轉換跑道.
跟隨耶穌.

$^{481}$成為一個得人的漁夫.
打人漁夫也很恐怖.
不要記錯.
通常錯的事情大家一定記得很清楚.
好,接著第23至25節.
剛才我們看了那個門徒.
他們被稱為門徒去跟從.
跟從這個字.
其實在第23至25節都是重複的.
那些群眾都是跟從耶穌的.
首先我講講跟從這個字很有意思.
跟從這個字.
其實是來自一個複合的字眼.
Aquiluteo.
這個字是怎麼來的呢?.
那個跟從.
跟從其實前頭那個字.
Aquio.
就是聽到的意思.
我聽.
聆聽.
Aquiluteo.
後面那個字就是同行.
如果你說跟從這個字.
將兩個字放在一起.
意思就是說.
聆聽得到.
接著就同行了.
所以那個跟從的意思.
其實不是很複雜.
聽到耶穌的訊息.
聽到耶穌的呼籲.
接著就與耶穌同行.
這個就是跟從.
不過除了這些門徒.
很快地跟從耶穌之外.
23節到第25節.
他們也有另外一群人.
來跟從耶穌.
23節.

$^{521}$耶穌走遍加里尼在各會堂裡教導人.
宣講天國的福音.
醫治各樣百姓的疾病.
他的名聲傳遍了敘利亞.
那裡有些人把病人.
各樣疾病,疼痛,被鬼附,癲癇,癱瘓.
都帶到耶穌面前.
耶穌就治好他們.
當時有一大群人.
從加里尼,底加坡里,耶路撒冷.
猶太,約旦河的東邊.
都來跟從他.
在23節到第25節裡面.
你看到.
其實這一群人都是來跟從耶穌的.
到了23節到25節.
耶穌對這群人做的事情就多了.
開始的時候.
耶穌呼召門徒.
耶穌說了一個單一的訊息.
「你來跟從我」.
「我要叫你得人如得魚」.
很簡單而已.
但是到了23節到25節.
耶穌走遍不同的地方.
在會堂裡教訓人.
傳天國的福音.
除了講之外.
再加多一重.
就是醫治各樣的病症.
有些奇難雜症.
可能患了重病.
各樣的疼痛.
又或者當時人力無法做的被鬼附.
癲癇,癱瘓.
注定了沒得走的了.
所有的病人,病患者.
人力無法解決的問題.
所有這些人就將自己.
也將這些問題帶到在耶穌面前.

$^{561}$耶穌就一一來醫好他們.
很正常的.
以前看病.
你不要想著今天.
找到好的醫生是很難的.
找到好的藥更加難.
在當代裡.
能夠得到醫治是很奢侈的.
一種很大的神蹟.
所以這群人真的慕名而至.
遍及整個敘利亞.
於是所有的人.
聽到耶穌的名聲.
他們就跟從耶穌.
在這群人裡.
他們不單想聽天國的福音.
我想這裡有一個很重要的百分比.
不過這群人除了想聽天國福音之外.
更加多的是想得到上帝的醫治.
想得到從耶穌而來的好處.
所以這群人的跟從.
不是那麼單純的.
他們想拿些東西.
在耶穌的醫治.
或者耶穌的講論當中.
他們想得到一些.
他們過去得不到的好處.
所以這群人就跟從耶穌.
值得留意的是.
在十八節到二十二節.
剛才耶穌呼召的兩對兄弟.
和這些群眾.
耶穌做的事是很近似的.
都是宣講.
後面加了一些醫治.
然後人去跟隨.
不過界定了第一班的門徒.
和這些群眾.
最不同的地方就是.
這些群眾沒有為了聽到.

$^{601}$耶穌所講的內容.
而捨下任何.
或者付出任何代價.
或者放棄眼前的東西.
但是在耶穌所呼召的兩對兄弟.
當耶穌直接了當.
講了一個訊息的時候.
這兩對兄弟.
就將自己眼前的生活.
將漁網或者自己的船.
放下,捨下.
然後就去跟從主.
在《馬太福音》裡面很低調.
但是其實在這裡.
已經隱含著.
原來這個世界裡有些人很喜歡耶穌.
很喜歡認定耶穌能夠為他解決問題.
不過門徒和群眾之間的分別在於什麼呢?.
捨下.
放下你眼前覺得重要的東西.
你說牧師就是這麼簡單.
捨下這個字.
這個釋經很簡單.
但是人生要做這件事很難.
其實在後來的篇幅.
如果大家記得.
有一個青少年,一個少年的官.
一個如日方中.
一個很有前景的年輕人.
走來問耶穌.
良善的夫子.
我應該作什麼可以承受永生?.
耶穌見到他就愛上他.
就覺得他有為青年.
不過耶穌就給他一個挑戰.
他說如果你想承受永生.
你就變賣你一切.
來變周濟窮人.
而且還要跟從.
這個年輕的官.

$^{641}$聽到耶穌這個挑戰.
叫他放下他手頭上擁有的東西.
周濟窮人.
周濟窮人是很樂意的.
不過把他的家產都賣了.
他就有點不願意.
聖經裡就說到.
這個年輕的青少年才信.
聽到耶穌這個挑戰.
就悠悠愁愁離開.
因為他擁有的東西是多的.
捨下,放下.
《釋經》裡是很容易的.
但是原來在我們的生活.
在我們的人生裡.
我們要能夠成為真正門徒.
原來這個功課我們必須要學.
如果我們不能夠在真理的亮光下學習.
捨下,放下.
大概我們都沒有辦法經歷到門徒的生活.
我們都沒有辦法經歷到.
從神而來那種豐盛的恩惠.
其實今天那段經文.
其實很簡單.
其實就是把門徒和這群群眾.
做一個很簡單的對比.
接下來所說的天國的信息.
其實都是用一個捨下.
作為一個很重要的主軸.
你要經歷豐盛.
如果你不放下眼前有一些你很執著的東西.
你經歷不了.
所以未來接下來在登山部分裡.
所說的很多內容.
其實都是圍繞著這樣的主題.
不過我很想我們今天再去踏真一點.
這個成為門徒.
在我們生命裡有三個很重要的.
我們需要今天留意的三部曲.
我們希望我們在這個世代裡.

$^{681}$不是純粹做一個聽道.
或者很喜歡耶穌的一個粉絲.
而是成為一個真正跟從主的門徒.
這三部你記住.
第一部分是「傾聽」.
那個跟從的字記不記得是「Aquo」.
「聆聽」.
「傾聽」的意思.
是代表著我們真的聽得到耶穌對我們所發出的呼聲.
我們今天很容易水過鴨背.
或者我們今天很容易覺得有些事情不關自己事.
選擇性的聆聽.
選擇性接收一些我們喜歡好聽的東西.
不過我們如果真正認真聆聽上帝對我們所說的話.
其實每一句對我們都是鏗鏘有力.
每一句都要直撼我們內心深處.
怎樣才算傾聽?.
怎樣才算有效的聆聽?.
每一次你聽.
每一次你去讀聖經的時候.
你想起每一句說話都對你有關.
曾經有些弟妹說.
牧師聽上帝的話.
聽到上帝說話好像很神秘.
我又不懂得那些超自然的祈禱.
有些人祈禱很厲害.
他說我不是的,我很拙劣.
或者我沒試過在天上突然靈修的時候.
寫了幾個字.
然後有個雲就報出一些字.
我都沒試過.
那我怎樣聽上帝所說的話?.
我說剛才說的那些很神秘.
那些可能會有的.
但出現一些字的那些.
我基本上在荷里活片才見過的橋段.
原則上今天在我們的信仰實踐是沒有的.
不過其實上帝已經對我們說了話.
上帝所說的話其實在整本聖經裡.
有一個猶太拉比說得很好.

$^{721}$我們今天很著重祈禱來聽上帝的聲音.
那是比較拙劣的做法.
那是很片面的聽那個聲音.
真正要聽上帝的聲音.
其實不是純粹我們的祈禱.
因為祈禱大部分都是你多說過上帝說.
真正的聆聽是在於你查經.
上帝要說的話在六十六卷聖經裡面.
當我們在上帝面前.
來查考祂的話語.
其實上帝已經對我們說了話.
怎樣才算是對我們說話呢?.
剛才說到.
怎樣去將上帝的話語和我們產生關係呢?.
你每一次查經.
你看聖經的時候.
你都要問上帝一個問題.
你今天藉著聖經對我說些甚麼呢?.
你都要聚焦在我自己的身上.
不要說今天對我老公說些甚麼.
或者對個人說些甚麼.
如果你這樣想.
就會將聖經成為武器來罵人.
但你用聖經來聚焦在自己上帝.
今天對我說些甚麼呢?.
你會自己個人拿到聖經的內容.
成為你生命的內容.
再舉多一個例子你就明白了.
如果你有機會去到機場.
坐飛機的時候.
那些廣播不是經常都會說的嗎?.
你聽完之後當耳邊風的.
你不會記得的.
你不會覺得那件事跟你有關.
但當你坐在機場很久.
在滯後.
你坐那班機.
何時要開出呢?.
你一直在等待.
豎起耳朵來聽的時候.

$^{761}$你那個就是積極的聆聽.
當你聽到那班機的號碼.
你一想起CX是甚麼號碼.
你一聽到這個.
你馬上就為之一震.
接著你聽著那個廣播員說的內容.
接著就來做.
那個就是主動的側耳而聽.
同樣道理.
在我們今天看聖經都是一樣.
我們每一樣查考的內容的時候.
找一找我們生命.
我們需要今天說的東西.
對我們個人有甚麼幫助.
由個人和上帝的話語重新建立.
如果我們聽不到上帝的聲音.
我們不願意尋求上帝的話語.
我們做不到門徒.
唯有傾聽.
仔細主動的聆聽的時候.
我們才能夠進入到被訴作門徒的身份當中.
這是第一個.
第二個就是同道.
我們信仰的路不是孤身作戰.
我們不是孤島式經歷我們的信仰生活.
我們需要同行者.
我們需要有人和我們一齊來行.
同行者不是純粹崇拜裡面.
歡迎民安握手.
旁邊那個見面熟口熟面.
而是知道你願意打開心扉.
和他一起行.
他又和你一起知道你生命的經歷.
為你互相禱告的那個同行者.
在教會裡面那個氛圍就是小組.
很多弟兄姊妹興翰小組的模樣.
有些覺得聽到已經很上乘.
很好了.
很足夠了.
不過我想說.

$^{801}$如果我們只是聽到.
我們的信仰缺乏了最精彩的部分.
有一個很不光彩的過去例子.
讓大家說一下.
最初回教會的大學二年級信主.
大學生自命不凡.
覺得自己掌握了整個基督教信仰.
才是決志信主.
那時候我是這樣想的.
回到教會的時候.
有些弟兄姊妹.
有些年輕的邀請我一起.
來參與他們的團契小組.
叫了幾次都不肯.
都拒絕.
藉口說自己沒空.
心底裡是怎樣想的呢?.
我看到邀請我那個.
跟我同年紀的.
大學二年級.
你大概想想大學二年級是怎樣的.
有一個我覺得是導師.
他邀請我去參與團契.
我以為他是導師.
原來他跟我同年.
他的樣子比我大二十年.
穿著Polo Shirt.
扣上兩粒頸後鈕扣.
穿著西褲.
穿著涼鞋.
穿襪子那種.
整件事很帥吧?.
我心想如果我跟他同一個團契.
我真的會死的.
我會不會變成他這樣.
我心裡覺得他老套.
所以我看不起他.
然後我就不想.
還有另外一個.
肥肥等等.

$^{841}$我看到他這麼不修邊幅.
如果我每個星期都看著他.
我真的不知道怎樣.
我不知道回到教會一段日子後.
我真的變成另一個大叔.
心裡看不起他們.
覺得他們老套.
覺得他們不配合我.
我覺得我自己跟他們一起.
我覺得我自己會拖垮自己.
所以我說.
很飄逸地回到教會.
孤身上路.
挺好的.
聽完道之後.
就可以消消灑灑地離開.
不過後來有一個姐妹再跟我說.
姐妹說得多一點心底話.
經過包裝.
將剛才的想法告訴她.
或者說我怕我將來.
在教會長久了就像他們一樣.
她說我怕我老套.
那個姐妹罵我.
她說.
你不想跟別人在一起.
不是因為別人老套.
而是因為你驕傲.
你覺得自己比別人更高一籌.
你覺得自己比別人聰明.
你覺得自己不需要別人幫你.
所以你不願意入組.
如果你說這些手法.
叫別人入組.
其實是很失敗的.
但這個姐妹說的.
正是我內心最污穢的一面.
這句說話令我改變.
我要回去.
我要找一個團契.

$^{881}$去到今天.
我回到教會二十多三十年.
二十多年.
剛才那些老套.
有一個胖胖的那個.
去到今天也是我生命中.
很重要的好友.
我生命中最不起眼.
最污穢的.
她一定知道.
同樣她生命的問題.
我也一定知道.
這些這麼好的同行者.
所以弟兄姊妹.
如果你來到教會.
你打算聽偏道.
就可以成為門徒.
我想說.
你錯失了最重要的那部分.
同行是最重要的.
教會席著一群人.
一群未必是最英俊.
最美麗.
最有體面的人.
甚至可能大家都充滿了不同的問題.
但這一群人.
上帝卻是用這群人.
來塑造我們自己的屬靈生命.
甚至乎坊間經常聽到.
爭論教會的不事.
說教會都是一群很不完美的人.
很差的人.
說很多這些論調.
你上網可以聽到.
不過你可以得意一點來說.
慘得上帝用這群人.
來塑造我們.
這群不起眼,不完美的人.
卻成為上帝家裡的一份子.
我們很厲害嗎?.

$^{921}$我們只不過是蒙恩的罪人.
我們在一個不完美的群體裡.
被塑造,被模樣,被建立.
我們認定我們是不完美.
我們不需要找一個完美的群體來參與.
我們只需要群體.
來成為我們的同行者.
帶領我們.
所以弟兄姊妹.
我鼓勵大家.
如果在新的一年裡.
或者你回到旺鎮一段日子.
你未有自己的群體.
你要找一個屬於你的小組.
倘若你在小組裡.
你生活得一般般.
你弱積弱離.
我鼓勵你.
全身投入在這個小組裡.
這是上帝用來塑造你生命的一個群體.
不要弱積弱離.
不要在這個群體裡.
不完美的群體裡再挑骨頭.
反而是珍惜這個群體.
上帝藉著這個群體.
塑造你和我的生命.
第三方面.
舍下.
剛才我已經說過了.
西比太的兒子也好.
或者西門彼得.
或者安德烈.
為什麼他們願意舍下呢?.
我們以為他們很犧牲.
甚至可能西比太也覺得.
你把這個祖業忘卻了.
不過.
如果我們經常有種心態覺得.
我為上帝犧牲很多東西.
我們不用走的.

$^{961}$其實舍下.
放下.
一個很重要的意涵是.
我找到一個更加有價值的東西.
以致我把眼前.
我手執著的次要價值放下.
其實是懂得選擇.
舍下的意思是懂得選擇.
選擇一些重要的東西.
把次要的東西放下.
今天我們生命裡.
充斥著大量次好的東西.
在我們的內心.
以致我們沒有辦法.
選擇最好,上好的東西.
不是純粹講錢.
我只是講時間.
我們今天的時間.
被大量次好的東西佔據.
我不是否定那些價值.
不過我想講是次好.
即是次好的東西.
佔據了的時候.
你只有24小時.
你沒有辦法選擇上最好的.
容許我在這裡講.
我不是自誇.
我只是覺得我是教子女如何選擇.
我今天有些子女.
他們都很喜歡回到皇寺.
為何會這樣做呢?.
一方面我太太很努力.
要建立他們的紀律.
不過如果你強迫他們紀律他們.
其實他們不會喜歡回到教會.
我們在家裡經常強調.
講教會的過癮事來聽.
快樂事來聽.
每一個星期回到教會.
講一些好笑的事.

$^{1001}$當一個群體可愛.
在這裡有些傻乎乎的事.
但那是我們生命的一部份.
子女就是這樣數著數著.
他們就愛上了教會.
另一方面就是.
我們不是放高調.
我也給我的子女學東西.
學音樂,學畫畫.
我都會給的.
不過我們的金科玉律是.
安排這些放外活動.
不可以和自己的教會聚會撞時間.
我不會刻意.
甚至我不會容許.
我的子女的放外活動.
安排在星期六下午團契.
或者星期日.
回崇拜事奉的時間.
來參與他們的興趣班.
他們喜歡的事.
我都是講講.
那些是次好的事.
不是差的事.
不過如果你被這些佔據了.
將來他們就沒有辦法選上最好的.
我要告訴他們.
我們的生命,我們的時間.
是用最好的時間.
來選上最好的方式.
最好的方式.
未必一定是那次的茶經蕩氣回腸.
單單回來.
單單一齊.
單單敬拜神.
都已經足夠我們稱為最好的.
今日的兄弟姐妹.
你怎樣去選擇呢?.
願不願意將你自己覺得.
多好,多重要的事.

$^{1041}$學會卸下.
選擇上好的福分.
以上帝成為你生命裡.
選上的時間表.
重新規劃下你自己的生命.
當你明白了時間怎樣運作的時候.
你的金錢.
你的人生擺放.
你的學習.
都會有很不同.
求主幫助我們.
昔日耶穌很簡單地來到呼召.
這班人看見基督耶穌的份量和價值.
他們就安然快樂地.
沒有掙扎地將眼前的願望卸下.
同樣我們都希望.
當今日耶穌的話語來到我們當中.
我們別無異心.
我們輕省地選擇耶穌.
門徒訓練是甚麼呢?.
是傾聽.
是有人同道.
和學習卸下.
盼望在這個世代裡.
上帝打造我們.
成為祂所使用的門徒.
我希望我們往心.
我們這個群體.
我們繼續門徒不斷.
\newpage



\section{馬太福音 5:13-16}
\label{sec:Z3AMKkvW_mc}
\textbf{2023年1月15日崇拜直播｜陳冠成牧師｜給當代門徒的15個忠告:成為世上的鹽和光｜馬太福音5\_13-16}
\newline
\newline
連結: \href{https://youtube.com/watch?v=Z3AMKkvW_mc}{\texttt{https://youtube.com/watch?v=Z3AMKkvW\_mc}} ~~~~ 語音日期: 2023-01-17
\newline
\newline
\hyperref[sec:WHVDHKq51Y4]{\small{< < < PREV SERMON < < <}}
~
\hyperref[sec:index_chronic]{\small{[返順時目]}}
~
\hyperref[sec:AMjA5Kj3ujM]{\small{> > > NEXT SERMON > > >}}
\newline
\newline
馬太福音 5:13-16
\newline
\begin{longtable}{cl}
\hline
\hline
章節 & 經文 (和合本修訂版)\\
\hline
5:13 & \begin{tabularx}{0.7\textwidth}{X} 「你們是地上的鹽。鹽若失了味，怎能叫它再鹹呢？它不再有用，只好被丟在外面，任人踐踏。 \end{tabularx} \\ \\ \relax
5:14 & \begin{tabularx}{0.7\textwidth}{X} 你們是世上的光。城造在山上是不能隱藏的。 \end{tabularx} \\ \\ \relax
5:15 & \begin{tabularx}{0.7\textwidth}{X} 人點燈，不放在斗底下，而是放在燈臺上，就照亮一家的人。 \end{tabularx} \\ \\ \relax
5:16 & \begin{tabularx}{0.7\textwidth}{X} 你們的光也要這樣照在人前，叫他們看見你們的好行為，把榮耀歸給你們在天上的父。」 \end{tabularx} \\ \\
[1ex]
\hline
\hline
\end{longtable}
$^{1}$經內容記載在馬太福音第五章第十三至第十六節.
亦是大家熟悉的章節.
馬太福音第五章第十三節.
他是這樣說的.
你們是世上的鹽 鹽若失了味 怎能叫它再鹹呢.
以後無用 不過凋在外面 被人踐踏了.
你們是世上的光 成灶在山上 是不能隱藏的.
人點燈 不放在斗底下 是放在燈台上.
就照亮一家的人 你們的光也當這樣照在人前.
叫他們看見你們的好行為 便將榮耀歸給你們在天上的父.
弟兄姊妹平安.
今天我們會看耶穌基督在登山寶訓上的教導.
其實前兩年 二零二一年.
大概在這個時候 第一季.
我們也分享過八福這個系列.
大家還有印象吧.
或者如果有需要可以在網上重溫.
其實八福是耶穌要求門徒應有的品格.
而緊接下來我們剛才所聽的光和鹽的教導.
其實耶穌進一步說明門徒和這個世界的關係.
對這個世界到底會發揮什麼影響力.
如果我們有留意剛才所讀的經文.
你會發現其實耶穌宣講的對象是哪些人?.
是門徒 或者一些願意跟隨耶穌.
喜歡聽他的教訓的人.
絕大部分都是學歷或社會地位不高的人.
門徒甚至乎在當時的社會是少數和異類.
其實這麼微小,這麼薄弱的力量.
真的可以改變這個世界?.
其實聽起來有點令人難以置信.
於是你看到耶穌就舉了兩個跟日常生活息息相關的比喻.
來肯定門徒確實對世界是有影響力的.
祂表示門徒是世上的鹽,世上的光.
鹽和燈光可以說是當時每家每戶必備的.
其實是很普通,很平凡的事情.
不過卻是很有價值.
可以為當時的人帶來許多益處.
特別在古羅馬的時代.
鹽可以說是當時羅馬軍隊的糧餉.
就好像今天我們說的英文的薪金.

$^{41}$其實是甚麼英文字?.
Salary,是嗎?.
其實是來自古羅馬的拉丁文Salarium.
意思其實是Short Money.
即是鹽的貨幣.
反映當時鹽在當時的社會是相當有價值的.
當然今天我們出糧.
沒有人會給你一包鹽當作是出糧.
不過鹽同樣對我們有許多的用處.
首先我們會明白鹽裡面有鈉的成份.
Sodium,是我們人體必需的營養素.
為我們提供變解質.
鹽亦可以作為一個防腐劑.
特別在以前未有雪櫃,冰箱發明之前.
我們會用鹽來做甚麼?.
來醃製食物.
令它不容易變爛.
又或者可以增加保存度.
可以放得久一點.
其實今天我們都沿用這個方法.
你在生活上都接觸到很多用鹽醃製的食物.
可以放得久一點,存夠一點.
又或者喜歡煮食的弟兄姊妹.
在你的廚房一定會有甚麼?.
一定有鹽,用來做甚麼?.
用鹽煮食物,醃肉,調味.
還有如果我們有急凍雪藏的肉類.
拿出來有雪櫃的氣味.
如何去除它?.
其實可以用鹽水浸泡.
去除雪味.
令食物更加滋味和可口.
當然你會看到耶穌表示.
鹽要發揮功效.
其實是有一個條件.
就是必須保持它應有的鹹味.
耶穌這裡說.
鹽若失去鹹味.
就會變成沒用.
亦會被人調戲.

$^{81}$其實祂在說甚麼意思呢?.
原來在古代巴勒斯坦.
他們吃到的鹽.
和我們今天在超市,雜貨店.
買到的食鹽.
其實不是那些.
絕大部分都是從礦場.
或是岩洞.
挖出來的岩鹽或礦鹽.
挖出來的時候.
像一塊石頭般大.
而鹽份就包含在岩石的表面.
換句話來說.
當猶太人吃飯的時候.
他們可能會用右手拿著食物.
左手拿著甚麼?.
就是拿著那塊鹽的石塊.
如果想食物更加滋味,更加鹹.
可以怎樣做呢?.
吃一口食物.
然後舔一舔它.
舔一舔它.
很不衛生吧?.
不過你放心.
不會大家共用一塊的.
理論上可能家裡每個成員.
都有自己那塊鹽塊.
每次吃飯出來.
就拿出來舔一舔.
但你想想.
石頭上的鹹味會怎樣?.
越來越少了.
石頭就會越來越淡.
甚至失去鹹味.
表面看來它還是原先那塊.
包著鹽份的石頭.
但實際已經是一塊普通石頭.
已經失去鹹味.
結果就像耶穌所說.
沒用的石頭.

$^{121}$普通的石頭就怎樣?.
扔掉它.
扔在外面.
被人踐踏.
這就是耶穌所說的比喻典故.
提醒門徒不可以成為.
一塊失去味道的鹽.
到底怎樣才能保持.
作為鹽的效果呢?.
嚴格來說.
鹽本身的化學成分是什麼?.
還記得你中學上過的化學課嗎?.
NaCl.
嚴格來說.
經過提煉純度很高的鹽.
其實是頗穩定的.
不容易失去味道.
除非你將它和其他雜質混合.
就可能讓它受到污染.
甚至失去效用.
同樣的.
作為基督徒我們也有這個情況.
耶穌猜我們去改變這個世界.
影響這個世界.
但與此同時.
這個世界又想怎樣?.
又想混亂你.
又想影響你.
所以我們有時會很容易.
落入一種試探.
就是不自覺被社會的價值觀.
來參雜.
混亂了我們的信仰.
特別是教會的歷史.
悠長的歷史裡.
總離不開一種情況.
就是.
作為基督徒有時.
很容易和世俗的文化妥協.
有時會靠過去.

$^{161}$不知不覺.
相反有時我們也會被他們溝淡了.
沒有自力去傳遞基督信仰的美善.
久而久之.
我們有時會變得.
和未認識耶穌的人.
看上去沒什麼分別.
就好像那塊失去了鹹味的石頭.
弟兄姊妹.
耶穌這裡很清楚說.
唯有我們將.
祂所吩咐門徒的要求.
或者說八福裡的品行.
實踐出來.
唯有我們持續為耶穌基督作見證.
才能夠成為世上的鹽.
能夠確保發揮到我們的鹹味.
相反.
如果我們不斷被世界的價值觀同化.
不斷妥協的時候.
被世俗的風氣污染我們的時候.
其實我們的鹹味.
就會漸漸淡化,溝淡.
甚至乎.
好像耶穌這裡說.
沒有了.
失去了.
最後也就像耶穌所說.
以後無用丟在外面.
被人踐踏.
沒什麼人會欣賞我們的信仰.
甚至看不起我們.
不單止失去了對世界的影響力.
甚至反過來被世界掌控我們.
所以耶穌這裡說.
我們作為世上的鹽.
其實.
有一件事很清晰的就是.
鹽怎樣才能發揮到它的價值?.
鹽放在這裡會不會發揮到價值?.

$^{201}$如果鹽就這樣放在這裡.
有沒有收過一些禮物.
拿一個很漂亮的瓶子.
放著一瓶鹽給你?.
有沒有收過這些?.
捨不捨得用?.
通常你不會用的.
那瓶鹽還是不是鹽?.
看上去好像是.
但其實.
它只是一個擺設.
發揮不了應有的功效.
所以作為世上的鹽.
那瓶鹽是不能拿來放的.
一定要怎樣?.
是要拿來用的.
一定要拿來接觸身邊的人.
造就到身邊的人.
那你才有所謂鹽的價值.
正如.
歌羅西書四章六節這裡說.
我們的言語要常常帶著和氣.
好像鹽調和.
我們有沒有世上.
帶著耶穌基督給我們的福氣.
帶著他的話語.
在我們不同的人際關係.
接觸人的時候.
成為一個調和的作用呢?.
我還記得.
剛剛去到平安夜.
教會的工作都完成了.
於是就回家.
做到很累.
就想著買外賣.
結果去了一間店舖.
他賣炸雞和海南雞飯.
我點了海南雞飯.
你知道海南雞飯.
通常配一個油飯.

$^{241}$比較好吃.
但那天想吃得清一點.
要白飯.
結果最後.
他又舀了油飯給我.
我看到那兩個店員很生氣.
弄錯了東西.
因為那天很多訂單.
你知道平安夜很多人叫外賣.
炸雞,烤雞.
我看到他們這麼嘈雜.
我就覺得算了.
不要緊.
胖一點也沒有所謂.
不要緊.
照樣油飯.
結果他們就沒事了.
很開心.
臨走前我說了一句.
聖誕快樂.
其實弟兄姊妹.
有時我們很簡單一句.
俠兒的說話.
你會否發覺.
其實可以改變人與人之間.
一些相處的氣氛.
本來可能是嘈吵.
很露露的.
但你送上一段俠兒的話.
有時可能能夠替人造就和睦.
讓我們看到其實.
鹽有時發揮的功效.
其實相對低調.
你不需要敲鑼敲鼓.
你想想.
你將鹽灑在食物上的時候.
鹽又會怎樣?.
靜悄悄地自己融化.
自己滲入肉裡.
不知不覺就改變了食材的質感和味道.

$^{281}$同樣我們有時為耶穌作見證.
其實亦可以不經意自然流露.
我們代表耶穌向你身邊的人.
譬如你出去食飯的時候.
送上一句關心鼓勵的話.
或者即使對方有做錯的時候.
你送上一句體諒的話.
其實就好像不經意在那些場合.
在人身上灑少少鹽.
將那些未善的事物.
可以保存得久一點.
同時亦將那些可能不太理想.
充滿嘈吵或者充滿怨氣的局面.
你不要讓它再難下去.
不要讓它再腐化下去.
在形形異異的奴隸生活裡.
你為你身邊的人灑少少鹽.
加添生活的滋味和樂趣.
這是耶穌給我們的提醒.
事實上祂將你和我放在世上.
正正是要阻止社會這個世界.
不要再難下去.
並且祂很想我們慢慢滲入.
社會不同的階層角落.
你去改變裡面的氣氛.
你去祝福多些那裡的人.
所以有一個問題很值得我們去想.
當你在身處的環境.
無論是工作的地方.
或者是你的家.
或者是人際關係裡面.
甚至是教會裡面.
當你看到一些東西不太好.
差一點點.
甚至開始爛的時候.
我們除了表達不滿,投訴.
或者埋怨.
其實這段經文很值得我們去想.
究竟我們的尊嚴去了哪裡?.
我們的尊嚴去了哪裡?.

$^{321}$我們有沒有及時在適當時候.
灑上適當的鹽.
來發揮防止腐爛.
發揮潔淨或者保鮮的功效.
還是我們尊嚴還藏起來沒有用過呢?.
很值得我們去想.
求主亦給予我們勇氣.
我們在致力守護那些美善事物的同時.
我們亦敢於正視那些罪惡不好的局面.
特別是當我們為主作見證.
或者你去分享福音的時候.
其實有時在所難免.
我們都要幫我們的對象.
我們身為罪人本質的對象.
要正視他的本相.
有時難免好像是將少許鹽灑在傷口裡.
有沒有試過用鹽去灑在傷口裡.
如果你的傷口剛好有傷口.
或者有粒飛脂的時候.
會怎樣呢?.
灑一灑那一下.
很不舒服的.
因為有刺痛的感覺.
當我們有時要指出一些盲點的時候.
就好像灑少許鹽在傷口裡.
會令對方叫痛.
但你知道這個過程是必須的.
亦都給予他好處.
所以求主給我們真是有愛心.
智慧和耐性.
無論我們是做一些人與人之間調和的工作.
又或者做一些撥亂反正的事情.
幫我們在適當的時機.
將適當份量的鹽灑在對方身上.
讓他真是得到福音的好處.
得到聖經真理的光照.
繼續看下去.
十四至十六節.
耶穌說門徒是世上的光.
要被人看見.

$^{361}$不能夠隱藏.
耶穌舉了一個比喻.
祂表示人在夜晚.
或者在很黑的地方.
很自然就會點燈.
點了的燈一定會放在燈台上.
使全家人可以得到光亮.
提供照明.
所以沒有人會將點了的燈放在斗底下.
斗底即是當時的一些糧器.
用來糧食和裝飾.
沒有人會這麼奇怪.
因為點了的燈等於沒有點.
因為沒有人看到.
亦發揮不了原有的照明功效.
所以耶穌這個比喻已經假定了答案.
就是沒有人會這樣做.
即是基督徒的行為是不會不被人看見.
不合理的.
所以比起世上的炎.
我們剛才說比較低調.
我們靜悄悄地將自己融化.
將自己寫寫.
滲入別人的生命裡.
讓別人得到受惠.
這種相對比較低調的表現.
世上的光就顯然高調很多.
特別耶穌這裡很清楚說.
什麼是我們的光.
我們的好行為.
我們的好行為就是.
我們的光要具體呈現出來.
要被人看得見.
注意到.
換句話來說.
原來我們不可以只讓人知道我們是基督徒.
我們更加要讓人看到.
我們的好品行.
我們的好行為.
來隨時為上帝作見證.

$^{401}$我還記得很久以前.
坐電梯在商場.
進了電梯後.
有一對母子進來.
媽媽很生氣.
未進電梯已經聽到她在罵兒子.
一直罵.
內容我不知道是什麼.
只聽到她不停地罵.
可能是扭賣玩樂.
或者其他東西.
電梯很嘈吵.
又有其他人在.
我又想有什麼可以做呢.
你會怎樣做?.
可以怎樣做?.
給我一些辦法.
叫媽媽不要罵.
不要罵了,太太.
我沒有這樣做.
因為我怕對方不知道會有什麼反應.
我當時只說了一句話.
太太,你去哪一層?.
很奇妙的一句話.
原本罵人的太太.
立即靜了.
我想大概是因為這句很簡單的說話.
她發覺原來.
有人注意到她.
有人關心她的需要.
那一刻很微妙.
她立即由憤怒的情緒.
出來.
靜了.
直至到她的樓層.
很奇妙,兒子突然說.
媽媽,你剛才是否因為這樣.
所以罵我?.
媽媽說,是的,就是這樣.
開始變得很溫柔.

$^{441}$開始有正常的對話.
很奇妙.
正如上帝在詩篇裡所說.
祂的話語.
就是人腳前的燈.
路上的光.
是否只有基督徒適用?.
不是的.
上帝的話語.
適用於所有人.
對所有人都能夠.
提供幫助,都是有益的.
所以,弟兄姊妹.
當我們願意.
讓上帝的話語.
來掌管我們言行的時候.
你要確信.
上帝是會使用的.
是會幫助身邊的人.
是會帶來一些正面的影響.
其實是真的.
耶穌說我們是炎,是光.
其實背後意味著.
祂是很願意使用我們.
一方面,我們可以光照.
引導別人來走進.
上帝喜悅的道路裡.
同時,其實.
上帝的話語也作為.
一個光.
作為一個警惕的訊號.
就好像.
我們開車或逛街.
會看到有維修道路.
通常那些維修的地方.
會怎樣?.
凍了一些欄杆或水馬.
還有加一些燈.
那些黃黃的圓形的一盞.
晚上就會斬下斬下.

$^{481}$有什麼作用?.
叫你不要走過去.
前面有個洞.
不要走,危險.
再走過去就會受傷.
出事.
同樣,我們也可以運用.
上帝的話語.
勇敢地指出一些.
不恰當的事情.
或是一些上帝不起悅的事情.
來勸對方.
不要繼續.
以免.
他得罪上帝.
最後回不了頭.
所以.
耶穌也為這個緣故.
祂吩咐我們要成為.
世界的光.
透過我們的行為幫助別人.
可以體會到聖經真理的光照.
我們也是因為這個緣故.
被夢照作為門徒.
致力來去活出一個.
我們的行為.
和福音相稱的生活.
所以,弟兄姊妹.
我們需要操練.
我們需要很努力地.
用行動.
來表達我們的信仰.
好像.
我們現在很流行說一句話.
「唱好,說好香港故事」.
我們也要用我們的行動.
用我們的品格.
說好我們的信仰故事.
並且我們也要刻意地.
提醒自己.

$^{521}$謹慎自己的言行.
小心不知不覺.
將我們的信仰唱衰了.
不知道大家.
最近有沒有看一套電影.
《正義迴廊》.
可能很多人都看了.
裡面有一個角色.
扮演一個虔誠的基督徒.
而飾演這個.
虔誠基督徒的演員.
早前曾經接受一個訪問.
他表示自己.
他真的相信.
基督教的神是存在的.
他也很認同.
基督教所說的道理.
只不過.
他不想做一個基督徒.
他透露.
在現實生活裡.
他自己也曾經舉手決志.
也有參加過教會的查經班.
不過那位.
帶他回教會的朋友.
之後他再也找不到他.
感覺就像被人帶回教會後.
就丟下了.
當他面對很多.
信仰的疑惑.
找不到答案的時候.
他發覺也沒有人可以幫助他.
陪他一起成長.
所以最後他選擇.
離開了教會這個信仰群體.
最初他也有.
在網上自修.
自己找一些查經的資料.
但後來他終於決定.
將這個信仰放下.

$^{561}$原因也和.
基督徒的行為有關.
他表示自己.
曾經和一位對外.
說自己是一個虔誠基督徒的人.
但對方的所作所為.
卻令他大跌眼鏡.
弟兄姊妹.
我想這個例子多多少少反映.
世人是怎樣看我們.
怎樣看這個信仰.
當他們看到.
基督徒的行為不合乎體統的時候.
就會對我們的信仰.
產生疑惑.
甚至反感.
雖然我們.
可以很努力地解釋.
辯護.
基督徒也是罪人.
也是軟弱的.
應該將信心.
放在上帝身上.
不應該放在人的行為上.
我們也明白這個道理.
不過現實裡.
很多時候.
影響別人不信耶穌的.
確實是.
我們不好的見證.
不好的行為.
曾經聽過有人這樣分享.
他對基督教信仰的體會.
他說.
信耶穌是很好的.
不過和信耶穌的人.
在一起就不是很好.
弟兄姊妹聽完之後.
我們會不會也有種.
心鬱無奈的感覺.

$^{601}$但我們甘不甘心.
我們甘不甘心.
讓耶穌基督的命.
受虧損.
既然耶穌說.
我們是世上的光.
其實就表明.
我們一定會接觸.
面向這個世界.
所以我們的言行.
不只是為了回應.
耶穌的吩咐.
不只是回應耶穌的愛和期望.
某程度.
耶穌也要我們回應.
世人的期望.
正如耶穌說.
叫人看見我們的好行為.
便將榮耀.
歸比在天上的附近.
我們要面向.
接觸這個世界.
這是我們無可推諉的責任.
除非我們不信耶穌.
並且你留意到.
耶穌說.
你們是世上的光.
世上的燕.
耶穌沒有說一個單數只有你.
只有隔離那個.
意味著.
這是所有人的責任.
只靠個別人的努力.
其實不足夠.
畢竟當我們為主作見證的時候.
其實總會遇到挑戰.
遇到試探和挫敗.
單打獨鬥.
有時會很氣餒.
走不遠的我們這條信仰路.

$^{641}$如果只有我們孤單一個人努力.
我們很需要有夥伴.
需要人聽道行道.
需要動員教會.
一起來成為耶穌所說的.
世上的燕,世上的光.
持續不斷的.
為世界帶來正面的改變和祝福.
換句話來說.
其實耶穌提醒我們.
沒有一個人可以獨善其身.
沒有一個人可以自享其福.
每個人.
上帝說他都有能力成為光和燕.
只欠我們肯不肯去.
並且.
我們是要去面對.
一個罪惡的世界.
可能會傷害你.
可能會得罪你.
令到你不舒服.
但是.
經文提醒我們.
仍然要對未信的人.
甚至抗拒的朋友.
你要心存歎望.
耶穌沒有放棄過他們.
所以我們也不能放棄.
不接觸他們.
因為我們被召.
夢召了門徒.
就是要服侍那些不信的人.
事實上.
在耶穌時代.
當時死海一帶.
住了一個社群.
叫做愛悉尼派.
他們有一個生活取向.
就是覺得當時的世界.
凹凹亂亂烏煙瘴氣.

$^{681}$於是怎樣?.
自己組成一個群體.
我們去到死海附近的曠野.
在山洞裡面生活.
我們選擇禁慾主義的方式.
來過每一天.
確實是一塵不染.
確實很乾淨.
完全抽離這個世界.
他們並稱自己為真光之子.
但是很抱歉.
如果按照耶穌的教導.
他們並不是世上的鹽.
並不是世上的光.
他們的乾淨.
他們的聖潔.
只是自享其成.
自得其樂.
所以弟兄姊妹.
今天上帝可能將你和我.
放在某個場景裡面.
未必是一帆風順.
可能都是烏煙瘴氣.
有很多事我們都擺不平.
好像無能為力.
但是有沒有想過.
其實耶穌說.
對的.
那個地方就是將你成為鹽.
成為光.
一個很操練實踐的場景.
耶穌說可以.
你回應祂的呼召.
承擔祂給你的責任.
來面向問題多多的世界.
可能有時我們會.
鬼鬼地.
怪怪地.
覺得自己很渺小.
很無助.

$^{721}$很無力.
是否真的可以.
祝福到身邊的人呢?.
弟兄姊妹.
你相不相信耶穌基督的應許?.
正如馬丁路德.
他曾經在宗教改革的時期.
他都過著一條艱辛的路.
他面對很多的敵黨挑戰艱難.
他曾經說過一句話.
就是憑著耶穌基督一句話.
我就能夠和那些掌握大權的人.
手握刀劍和槍械的人.
我可以比他們更加無畏懼.
比他們更加有ye 可誇.
真的現實在裡面.
我們有時都無能為力的感覺.
但不要忘記.
我們有耶穌基督.
有祂所賜的福音使命和應許.
所以.
即是說可能有人會拒絕你.
甚至反對你所分享的信仰.
甚至和你唱反調.
但不要因為這樣而意志消沉.
更加不要因為這個世界.
有很多惡人,衰人.
所以我們都擺爛自己.
來將自己變差.
漸漸失去了光和炎的功效.
相反.
耶穌說你可以學習感恩.
可以學習歡喜快樂.
因為當你遇到這些挑戰逼迫的時候.
你就體驗到.
祂所說的話是真的.
是很可靠的.
你就體驗到.
甚麼是真實和耶穌基督同行.
就正如祂在八福裡面.

$^{761}$一早已經說明了.
是會這樣.
所以我很想邀請大家讀最後一段經文.
在馬太福音第五章十至十二節.
預備起.
當我們遇到困難的時候.
這段經文盼望成為我們的幫助和激勵.
耶穌自己也是這樣.
總會遇到挑戰和敵黨.
當我們為主作見證的時候.
不過耶穌同時說.
這是蒙福土上帝喜悅的方法之一.
事實上.
耶穌自己也帶頭做了榜樣.
祂第一個走出來.
服侍那些討厭祂.
反對祂.
不信祂.
迫害祂的人.
祂也吩咐我們.
我們也一起這樣吧.
就是這樣.
\newpage



\section{馬太福音 5:17-20}
\label{sec:AMjA5Kj3ujM}
\textbf{2023年1月29日崇拜直播｜吳真真牧師｜給當代門徒的15個忠告:遵行耶穌所成全的律法｜馬太福音5\_17-20}
\newline
\newline
連結: \href{https://youtube.com/watch?v=AMjA5Kj3ujM}{\texttt{https://youtube.com/watch?v=AMjA5Kj3ujM}} ~~~~ 語音日期: 2023-02-04
\newline
\newline
\hyperref[sec:Z3AMKkvW_mc]{\small{< < < PREV SERMON < < <}}
~
\hyperref[sec:index_chronic]{\small{[返順時目]}}
~
\hyperref[sec:iHCyyrWPaj4]{\small{> > > NEXT SERMON > > >}}
\newline
\newline
馬太福音 5:17-20
\newline
\begin{longtable}{cl}
\hline
\hline
章節 & 經文 (和合本修訂版)\\
\hline
5:17 & \begin{tabularx}{0.7\textwidth}{X} 「不要以為我來是要廢掉律法和先知。我來不是要廢掉，而是要成全。 \end{tabularx} \\ \\ \relax
5:18 & \begin{tabularx}{0.7\textwidth}{X} 我實在告訴你們，就是到天地都廢去，律法的一點一畫也不能廢去，直到一切都實現。 \end{tabularx} \\ \\ \relax
5:19 & \begin{tabularx}{0.7\textwidth}{X} 所以，無論誰廢掉這誡命中最小的一條，又教導人也這樣做，他在天國裡要稱為最小的。但無論誰遵行並如此教導人的，他在天國裡要稱為大。 \end{tabularx} \\ \\ \relax
5:20 & \begin{tabularx}{0.7\textwidth}{X} 我告訴你們，你們的義若不勝過文士和法利賽人的義，絕不能進天國。」 \end{tabularx} \\ \\
[1ex]
\hline
\hline
\end{longtable}
$^{1}$經文記載在馬太福音第五章十七至二十字.
莫想我來要廢掉律法, 和先知我來不是要廢掉, 乃是要成全..
我實在告訴你們,就是到天地都廢去了, 律法的一點一劃也不能廢去,都要成全..
所以無論何人廢掉這誡命中最少的一條, 又教訓人這樣作..
他在天國要稱為最少的, 但無論何人遵行這誡命, 又教訓人遵行, 他在天國要稱為大的..
我告訴你們,你們的義若不勝於民事和法理裁人的義, 斷不能進天國..
承認獨白..
各位弟兄姊妹,今天是農曆年初八, 所以我都在說一些恭賀的說話,.
恭祝大家身心靈健壯,靈命步步升!.
我們除了套年開始, 現在是一月,大家記不記得今年教會的主題是什麼?.
門徒不斷!.
在一月開始,我們有一個主題的講道系列, 叫做「給當代門徒十五個忠告」..
今天我們來到第三講, 講登山寶訓的教導..
第一個教導是陳明荃牧師提醒我們, 要求我們要作一個真正追隨基督的門徒..
然後上一講是Ray牧師跟我們分享, 要我們去作為世上的炎和光..
今天我們會講到《馬太福音》五章十七至二十節..
上次我們講到五章十六節的時候, 耶穌吩咐我們,.
「你們的光也當這樣照在人前, 叫他們看見你們的好行為,.
便將榮耀歸給你們在天上的父. 」.
門徒的好行為可以照耀這個世界, 好像世上的明燈..
使人將榮耀歸給天父. 但到底什麼是好行為呢?.
對於當代的猶太人來說, 耶穌必須宣講一個很重要的基礎,.
讓那些天國的子民能夠理解到 律法命令的重要本質,.
以及用什麼樣的態度去回應. 所以我們今天會去聽或學習第三個忠告,.
就是去遵行耶穌所承傳的律法, 讓我們同心祈禱..
「親愛的恩主,求你幫助我們, 讓我們謙卑自己的心,.
倒空自己,讓我們聆聽你的道, 一會離開禮堂之後,.
我們的心裡裝滿滿, 因為你自己使我們滿足,.
因為你的道叫我們很紮實地活出來, 求你幫助我們,.
聽我們的祈禱,禱告交托,奉主明求, 阿們. 」.
對於我們這一班基督徒來說, 我們很清楚,.
登山寶訓就是我們日常生活裡面 耶穌基督對我們的命令,.
但是對於當代的猶太人來說, 他們就很不同了..
讀一讀17至18節, 耶穌說:.
「莫想我來是要廢掉律法和先知, 我來不是要廢掉乃是要成全..
我實在告訴你們,就算天地都廢去了, 律法的一點一劃也不能廢去,都要成全.」.
這裡說到耶穌是要來成全律法和先知,.
對於當時的猶太人來說, 他們去回應,去思考,.
究竟他們說的好行為是什麼呢?.
原來就是去遵行律法和先知的命令,.

$^{41}$所以耶穌在這裡開始就要告訴我們, 究竟什麼是律法和先知..
當時的猶太人,他們以摩西五經, 即是我們說的「妥拉」,.
先知書和聖卷,作為他們舊約的希伯來文的聖經教導,.
就是他們的行事為人的標準..
在「妥拉」裡面,他們列了613條的誡命,.
成為他們在人生裡的道德標準,.
不過法利賽人在這些誡命上, 再加一些解釋和補充,.
他們希望能夠達到所有的民眾,.
都不會因為無知或意外而干犯主的律法..
這613條的律法是什麼呢?.
舉個例子,就好像他們其中一條, 就是「在安息日,民遂安息」,.
來自《初一及記》23章12節,.
而法利賽人在上面再加一些注釋,.
就是清楚地列明, 怎樣才算是安息呢?.
他們列舉了39條主要的動作, 是不可以做的,.
包括什麼叫作「作供」呢?.
就是收割,打谷,準備飯食,.
這些都是在安息日禁止的..
而且他們也清楚地列明,.
安息日最多可以走960米的路,.
這就是他們定的規矩..
當時的猶太人認為, 只要實行這些規矩,.
他們就能守法律及先知的律例,.
他們就能做好行為..
耶穌一開始就說,.
「我來不是要廢掉律法.」.
為什麼耶穌要這樣說呢?.
為什麼無緣無故一開始就說自己不是廢掉律法呢?.
因為在耶穌的傳道生涯一開始,.
祂的言行都是很破格的..
例如, 祂的言行是怎樣呢?.
馬可福音一章二十二節, 一開始就說,.
「耶穌走遍各城各鄉, 去會堂裡教訓人,.
眾人稀奇祂的教訓,.
因為祂教訓他們的, 正像有權柄的人,.
不像民事.」.
耶穌就算是說的道理,.
都是和民事, 他們的律法, 宗教的領袖,.
是很不同的..
另外, 耶穌的行事,.

$^{81}$對於當時的法利賽人來說,.
也是很犯禁的..
大家應該還記得,.
耶穌在安息日治病,.
祂的門徒在田間裡摘蜜水,.
又試過在吃飯的時候不洗手,.
結果引來很多法利賽人對祂的批評,.
但耶穌每一次都用真理去駁斥他們的指控,.
所以耶穌的言行都會令人有個印象,.
「你來是否要廢掉律法?.
你來是否來踢館?」.
其實廢掉這個字,.
在第17節的廢掉這個字,.
是指在建築物裡拆東西下來,.
去毀壞它, 或者去廢除它, 終止它..
耶穌在這裡,.
祂要除去當時的人對祂言行的質疑,.
認真地宣告,.
「我來不是要將律法剷除, 另起爐灶.」.
反而在第18節,.
祂說「我實在地告訴你們,.
實在的這個字是告訴我們,.
祂是肯定的,千真萬確的,.
甚至在約翰福音一說到「實實在在」這個字,.
就是說話的那個人的地位,.
是比一般人超越的..
祂要很清晰,.
無可置疑地宣告,.
兩個很重要關於律法的本質..
第一, 祂怎麼說呢?.
就是「天地都廢去了,.
律法的一點一劃都不廢去.」.
其實耶穌用了兩個誇張的修辭手法.
去表達律法的本質..
第一, 祂將律法的存在.
和天地的存在做了一個對比,.
代表了律法的永恆性,.
律法的本質是永恆的,.
律法是從創天造地的神所設立的,.
祂命立就立..

$^{121}$然後祂說「一點一劃都不廢去」,.
一點一劃就是聖經的文字.
那些最微細的標記,標點符號,.
這麼微小的標點符號,.
都不能夠廢去,.
都不能夠終止,.
都不能夠剷除,.
這樣就代表了律法的整全性..
耶穌不是要駁斥律法,.
不是要廢止律法,.
反而是要表明律法的永恆性和整全性..
這位獨一無二,.
道成肉身,來到世界的基督,.
才是那位承傳永恆而整全律法的那位..
他要為了律法,.
去作出重新的演繹,.
以及要帶出律法的真正意思..
「承傳」這個字,.
在馬太福音是非常重要的用詞,.
這個字在新約出現了86次,.
在馬太福音就已經出現了16次,.
而有13次是說到.
「承傳」這個字是說到應驗舊約先知的預言..
他告訴我們,.
這個舊約的律法和先知的預言,.
根本就是指向了耶穌基督,.
唯有耶穌基督才可以親自來到這個世界,.
去承傳這個律法..
他要帶出律法的真正目的和定義,.
他是創造律法的人,.
創造律法的那個,.
他所說的是帶著權柄,.
他所反映的是比摩西的律法精神更高的要求,.
他的行事為人本身就可以滿足律法的要求,.
甚至可以提升到另一個層面..
他如何重新演繹律法呢?.
耶穌超越了當時宗教領袖對於律法和先知命令的解釋,.
也糾正了他們傳遞的一些律法傳統的見解,.
耶穌是唯有他才能真正教導天國的門徒..
接下來我們會有六個宗谷,.

$^{161}$是21至48節,.
講到六個律法或誡命的教導,.
他指出了舊約裡面,.
講到不可以殺人,不可以姦淫,.
講到丘妻不可以背誓,.
講到以眼還眼,愛鄰舍,恨仇敵,.
這些舊約所指明的律法誡命,.
他將新和舊做出全息對比,.
耶穌鏗鏘有力地鋪陳摩西律法的完整含義,.
對於當時來說是一個好革命的表達,.
不單單延續了律法的有效性,.
而且比起文士和法理賽人的要求更加嚴格,.
告訴我們如何將這些律法落實在生活裡面,.
才是真真正正符合上帝對義的要求..
而且在耶穌的傳道生涯裡面,.
他不單單有這樣的教導,.
他更加指明了律法的真正意義和本質,.
其實就是我們和立約的上帝的關係..
《約翰福音》14章21節說:.
「有了我的命令就遵守的,.
這人就是愛我的,.
愛我的,便蒙我父愛他,.
我也要愛他,並且要向他顯現.」.
我們遵循律法,反映的是我們對上帝的愛,.
所以耶穌更加清楚地說明了律法和教命的總綱是什麼..
《馬太福音》22章37-40節說:.
「你要盡心盡性盡義愛主你的神,.
這是最大的,且是第一個誡命,.
第二條也是如此,.
就是愛,倫,捨,如同自己..
這兩條誡命是一切律法和先知書的總綱.」.
耶穌說明了,.
祂不是要來另起爐灶,.
祂只是很清楚地告訴我們,.
天國子民,你要知道律法的永恆性和整全性,.
律法的真正意義和它的總綱,.
然後這個律法才可以繼續延續下去..
究竟天國的子民應該怎樣回應呢?.
就要看第19-20節,.
其實耶穌就是叫我們進行律法..

$^{201}$耶穌在這裡告訴我們,.
祂怎樣去承傳律法,.
怎樣去延續律法,.
然後祂就說明了,.
作為天國的子民,.
對於律法的態度,.
原來就會影響了我們在天國的地位..
第19節,.
「無論何人廢掉這誡命中最少的一條,.
又教訓人這樣做,.
他在天國要稱為最少的..
但無論何人遵行這誡命,.
又教訓人遵行,.
他在天國要稱為大的.」.
這裡對我們有個很好的提醒,.
我們要小心,.
我們不可以廢掉律法..
第19節其實是衝著當時的.
法理賽人和文士而來的,.
因為主耶穌認識他們,.
因為他們的教導.
是將律法分等級的,.
他們將罪分類,.
有些是嚴重到必死,.
有些是很輕微的,.
他們會教導在律法裡,.
有些是比其他重要,.
違反最少的命令,.
是小事而已..
但是耶穌在這裡告訴我們,.
只要你廢掉戒命中最少的一條,.
還要教人這樣行,.
你原來在天國裡是小的..
第19節的廢掉是一個很特別的字,.
這個字用得很有意思,.
它和第17節的廢掉是不同的,.
這個廢掉是何時會用呢?.
這個廢掉是當死人的身體分解的時候,.
分解的字就是廢掉的意思,.
解容解,分解,.

$^{241}$又或者你起了一個帳幕,.
你的帳幕拆下來,.
就好像放鬆它,.
放鬆它的意思,.
在這裡對於門徒來說,.
是一個很好的提醒,.
我們千萬千萬不要放下律法,.
不要鬆開我們的律法,.
其實每一個時代的基督徒,.
都會受到當代文化的洪流所影響,.
時代的變遷都會衝擊著我們對真理的持守態度,.
有時會令我們不是用真理去看世界,.
反而是用世界去幫我們運用真理,.
令我們鬆開律法,.
放下律法,視律法為綑綁,.
所以我們會有個想法要鬆開它,.
特別我們這班蒙恩的罪人,.
我們都是身處在恩典的時代,.
羅馬書六章第十四節,.
保羅這樣說,.
「罪必不能作你們的主,.
因你們不在律法之下,.
而是在恩典之下.」.
如果單看這一句,.
很多人就會說,.
其實我們都是在恩典之下,.
在這個後現代的時代,.
鼓吹的是沒有絕對的真理,.
如果你很嚴謹地遵守真理,.
可能你會被人笑,.
「你怎麼這麼律法主義?」.
很多時候,.
這個時代都會用恩典蓋過其他東西,.
我相信耶穌,.
我已經有救恩了,.
耶穌體恤我們的軟弱,.
不要緊的,.
只要認罪,.
神就會寬恕,.
於是,.

$^{281}$當我們看登山寶訓,.
當我們看耶穌所吩咐我們,.
在天上,.
在地上行事為人的標準的時候,.
我們就會說,.
「哎,聖人都做不到,.
只是說而已,.
難道你被人打左邊臉,.
你真的被人打右邊臉嗎?」.
我們真的要提醒自己,.
我們會否鬆綁了律法呢?.
我們會否放下了律法呢?.
我們要認真去想,.
保羅在他的《羅馬書》裡面,.
他繼續說,.
他不是反對律法,.
他反對猶太人的律法教師,.
錯誤地將這班外邦的歸信者,.
帶到他們舊約的律法和傳統之下,.
反而他在《羅馬書》八章一至四節,.
很清楚地說,.
「如今那些在基督耶穌裡的人,.
就不被定罪了,.
因為似生命的聖靈的律,.
在基督耶穌裡從罪和死的律中,.
把你釋放出來,.
律法的因肉體軟弱,.
而無能為力,.
神就差遣自己的兒子成為罪身的樣子,.
為了對付罪,.
在肉體中定了罪,.
為要使律法要求的義,.
實現在我們這不隨從肉體,.
只隨從聖靈去行的人身上,.
律法透過耶穌去承傳,.
我們不再被定罪,.
但保羅同時也告訴我們一個更重要的道理,.
就是我們這群信徒,.
要信著聖靈,.
信服耶穌,.

$^{321}$這就是上帝對於我們承傳律法的義的要求,.
而且在第二十節這樣說,.
我告訴你們,.
你們的義若不勝於民事和法理賽人的義,.
斷不能進天國..
馬太福音經常說的義,.
其實是對神的話語和神的旨意遵守,.
這樣的意思,.
但是法理賽人的義是怎樣的義呢?.
馬太福音六章第一節說,.
你們要小心,不可將善事行在人的面前,.
故意叫他們看見..
善事這個字在希臘文就是義,.
法理賽人和民事的義,.
只是講求他們外表的敬虔,.
卻沒有內心敬虔的真實意思..
你們能否勝過法理賽人和民事的義,.
是在乎一個人的內心..
耶穌正正就指出了法理賽人最大的問題,.
他們驅離於律法的傳統,.
以進行自己,進行這些的誡命來自誇,.
叫人看見的是他們自己,.
他們並沒有從他們的心裡,.
驅動到以愛神,愛人的大誡命,.
所以我們必須要勝過民事,.
他們只求有敬虔的外表,.
將好行為,.
是帶人去看見他們的好行為,.
以進行誡命有多少來自誇,.
我們要勝過這個義..
但我們要強調的不是好行為才能進入天國,.
而是一個真正的天國子民,.
他的身份,他的生命的展現,.
就會展現他們如何進行上帝的教導,.
其實就是他內心信服上帝的主權,.
並且一心一意追求一個公義的生活,.
以至剛才五章第十六節所說,.
彰顯天父的榮耀..
所以弟兄姊妹,.
在這個時代,.

$^{361}$我們需要信服耶穌基督的主權,.
基督的真理是活的,.
不是死胡胡的,.
不是不能動搖的傳統,.
我們必須要每一天好好地.
去將自己立定心智,.
信服主..
有一位事奉主超過五十年的一位牧師,.
他叫做雷斯牧師,.
他在道行會門訓的訓練下成長,.
他曾經是葛培理報道團裡.
專負責栽培和牧養的一位牧者,.
他每一天開始會做兩件事,.
第一,.
他說,.
「我每一天會將身體獻上當作活祭,.
在神面前降服,.
將自己的主權交給祂,.
讓祂做主.」.
第二,.
他求聖靈管理他那一天的生活,.
他願意信服聖靈行事,.
不是滿足自己肉體的私慾..
每一天祈禱之後,.
他就讀經,.
然後散步,.
一邊散步,一邊祈禱,.
他視他每一天早上起床親近神,.
就代表了每一天與神同行的起步,.
所以他每一天都依靠神,.
求神引導和賜下智慧能力,.
去應對每一天的生活..
弟兄姊妹,.
其實我們知道這些訊息,.
我們知道律法的永恆和整全性,.
我們知道耶穌對我們的要求,.
但我很想在今天對大家提出三個建議,.
是在我們日常生活中可以實踐出來的..
第一,.
認真去對待神的話,.

$^{401}$認真去進行祂的教導,.
當你每一次去聽道,.
每一次去查經,.
去認識主的真理的時候,.
你去想想如何確切執行..
我記得在我成長中,.
我參加過道行會的門訓團契,.
我們的團契很簡單,.
每個月都是這樣,.
三個星期查經,.
一個星期分享和祈禱,.
如果有第五個星期,.
我們就會大早活動,.
每個月有一個星期六下午出去傳福音,.
最初我們的團契只有十多人,.
但四年後,.
我們的團契有八多人..
我想說,.
這個團契的增長,.
並不是因為我們有這樣的制度,.
而是我很深的認同一件事,.
就是大家對查考聖經的認真,.
以及對進行聖經的真理質素..
我最初很深刻的印象,.
就是每次和團契的領袖分享的時候,.
經常聽到他們掙扎和分享,.
他們如何應用聖經在他們日常生活的各個層面,.
包括他們的工作,.
包括在他們的家庭,.
包括在他們的事奉,.
令我們很深刻..
然後我們每一次查經,.
都會刻意預留一段比較長的時間,.
去討論我們如何應用經文,.
我們如何實踐..
他們的要求是我們要寫得仔細,.
要有計劃,.
不可以敷衍,.
查完「愛受敵」,.
你的應用不是「我愛受敵」,.

$^{441}$他們會問你,.
「說到愛受敵,你想起誰?」.
然後再問,.
「好吧,那你如何實踐?」.
「我為他祈禱.」.
「那你如何祈禱?」.
「一個星期祈一次?」.
「天天祈?」.
「我去關心他.」.
「如何關心?寫信,還是買東西給他吃?」.
每一次我們的應用,.
我們的組長都記得很清楚,.
他寫下,他抄下,.
然後下星期,.
就像追問劇情一樣,.
「怎樣?你如何實踐上星期的教導?」.
我記得當中,.
我們一起流過很多眼淚,.
也一起很歡喜,.
回來分享原來實踐主義教導的時候,.
那種和耶穌的同行和蒙福..
我們當時經常問,.
「上帝,你想我怎樣?」.
所以實踐的行動,.
第一,我呼籲大家,.
想想,.
你讀聖經的時候,.
想想如何應用..
第二,.
每一天晚上的時候,.
你去審查自己..
鄺炳釗博士在《如何經歷神童在》那本書裡,.
他有一段文字,.
很鼓勵我..
他說每一天結束的時候,.
我們就去審查自己當日有沒有信服神..
第一,.
聖經有清楚指示我的,.
我有沒有做..
第二,.

$^{481}$我有沒有做了一些聖經叫我不要做的事情..
第三,.
聖經沒有特別提及,.
我有沒有藉著禱告,.
用四個問題去確定,.
我應不應該做..
那四個問題是,.
這件事對我身心靈有沒有益?.
這件事會不會克制我?.
這件事對其他人有沒有害處?.
這件事是不是榮耀神?.
鄺炳釗博士說,.
願意遵行神命令的人,.
神就與他同在..
聖靈保衛師就跟你永遠在一起,.
他會指教你,.
教導你如何去行上帝的吩咐..
第三,.
在我們的生活裡,.
聖靈保衛師會提醒你..
當他提醒你的時候,.
我勸大家快快降伏,.
不要讓自己的意思,.
心思意念帶動你去回應,.
求主加力,.
給你勇氣,.
走出祂指引的第一步..
我最近有一個學習,.
家人吃團年飯,.
我口快快地說了一句.
對家人批評的話,.
結果引起了家人的爭執..
我當時馬上辯駁,.
覺得自己說的話是真的,.
很有道理..
但是當晚上你安靜下來的時候,.
聖靈就提醒我,.
剛剛那星期,.
Raymond不是教我們要作炎作光嗎?.
不是要調和嗎?.

$^{521}$那一刻,.
我立刻降伏,.
向主認錯,.
並且求主賜勇氣,.
讓我在我的家人當中,.
去說一些和解的話,.
去關心他們不同人的感受..
弟兄姊妹,.
我很希望大家能夠確切執行.
剛才那三個建議,.
認真地去思考,.
我們如何去行上帝的道,.
寫下你的應用,.
確切地去行..
每一天,.
要去審查自己,.
有沒有行上帝的道,.
有沒有得罪神,.
得罪人的地方..
當聖靈感動我們,.
提醒我們的時候,.
快快降伏,.
讓我們能夠.
懷著這個心,.
繼續去行主的吩咐,.
目的不是讓人見到我的好行為,.
目的是叫人將榮耀歸給天父,.
讓我們一起祈禱..
主,我們多謝你,.
我們實在有很多不明白,.
若果我們憑著自己的努力行為,.
根本我們沒有辦法.
符合上帝對我們這群天國子民.
行義的要求,.
因為耶穌基督擔當了我們的罪,.
祂讓我們活在恩典之中,.
但祂更願意讓我們過豐盛的生命,.
並且在天國得到賞賜,.
求你幫助我們有一個願意成長,.
得賞賜,.

$^{561}$夢想地閱立的生命,.
讓我們能夠認真去明白上帝的道,.
並且確切去實踐,.
求你幫助我們,.
我們是何等軟弱,.
我們需要上帝親自的加力,.
我們更需要弟兄姊妹彼此提醒,.
彼此激勵,.
好讓我們能夠一同在天國的跑道裡努力,.
更願意叫人因為看見我們的好運,.
好行為,.
去將榮耀歸給我們的天父,.
我們禱告,交託,奉主耶穌基督得勝的明智祈求,.
阿們..
\newpage



\section{馬太福音 5:21-26}
\label{sec:iHCyyrWPaj4}
\textbf{2023年2月5日崇拜直播｜陳啟興牧師｜給當代門徒的15個忠告:處理怒氣｜馬太福音5\_21-26}
\newline
\newline
連結: \href{https://youtube.com/watch?v=iHCyyrWPaj4}{\texttt{https://youtube.com/watch?v=iHCyyrWPaj4}} ~~~~ 語音日期: 2023-02-07
\newline
\newline
\hyperref[sec:AMjA5Kj3ujM]{\small{< < < PREV SERMON < < <}}
~
\hyperref[sec:index_chronic]{\small{[返順時目]}}
~
\hyperref[sec:N8TX9PAocpk]{\small{> > > NEXT SERMON > > >}}
\newline
\newline
馬太福音 5:21-26
\newline
\begin{longtable}{cl}
\hline
\hline
章節 & 經文 (和合本修訂版)\\
\hline
5:21 & \begin{tabularx}{0.7\textwidth}{X} 「你們聽過有對古人說：『不可殺人』；『凡殺人的，必須受審判。』 \end{tabularx} \\ \\ \relax
5:22 & \begin{tabularx}{0.7\textwidth}{X} 但是我告訴你們：凡向弟兄動怒的，必須受審判；凡罵弟兄是廢物的，必須受議會的審判；凡罵弟兄是白痴的，必須遭受地獄的火。 \end{tabularx} \\ \\ \relax
5:23 & \begin{tabularx}{0.7\textwidth}{X} 所以，你在祭壇上獻祭物的時候，若想起有弟兄對你懷恨， \end{tabularx} \\ \\ \relax
5:24 & \begin{tabularx}{0.7\textwidth}{X} 就要把祭物留在壇前，先去跟弟兄和好，然後來獻祭物。 \end{tabularx} \\ \\ \relax
5:25 & \begin{tabularx}{0.7\textwidth}{X} 你同告你的冤家還在路上，就要趕快與他講和，免得他把你送交給法官，法官交給警衛，你就下在監裡了。 \end{tabularx} \\ \\ \relax
5:26 & \begin{tabularx}{0.7\textwidth}{X} 我實在告訴你，就是有一個大文錢還沒有還清，你也絕不能從那裡出來。」 \end{tabularx} \\ \\
[1ex]
\hline
\hline
\end{longtable}
$^{1}$今日的經文是在新約聖經.
馬太福音第五章二十一到二十六節.
新約聖經第六頁.
馬太福音五章二十一到二十六節.
二十一節.
你們聽見有吩咐古人的話.
說不可殺人.
又說凡殺人的.
難免受審判.
只是我告訴你們.
凡向弟兄動怒的.
難免受審判.
凡罵弟兄是拉加的.
難免公會的審判.
凡罵弟兄是摩利的.
難免地獄的火.
所以你在祭壇上獻禮的時候.
若想起弟兄向你懷願.
就把禮物留在壇前.
先去和弟兄和好.
然後來獻禮物.
你同告訴你的對頭還在路上.
就說著與他和諧.
恐怕他把你送給審判官.
審判官交付遺役.
你就下在監裡.
我實在告訴你.
若有一文錢沒有還清.
你斷不能從那裡出來.
今天我們很高興邀請到.
金融業福音團契總幹事陳啟興牧師.
為我們正道講題是.
「給當代門徒的十五個忠告」.
處理怒氣.
請陳牧師.
好,弟兄姊妹午安.
我們先同心禱告.
親愛的主我們.
懷著很敬畏.
很謙卑的心.

$^{41}$來到你的面前.
我們懇求你的靈去引導我們.
進入真理.
我們深深知道.
如果沒有你的聖靈去引導.
我們沒有辦法去明白.
亦都活出真理.
求你幫助我們.
我們是愚鈍的.
我們需要你.
求你開我們的耳朵.
聽得明你的真理.
求你讓我們的手.
我們的腳.
能夠行出你的道.
改變我們.
讓我們在地上的生命更加土地起源.
同心禱告仰望.
奉耶穌基督的聖名.
阿們.
在幾年前.
我有一次去赤柱監獄.
為幾個信主的所願洗禮.
洗完禮之後.
我跟身旁一些其他的所願去聊天.
我也知道這些是甲級犯.
多數都是終身監禁.
我跟其中一個所願去聊天.
我覺得他怪怪的.
他的言談古古怪怪.
我也覺得可能有些問題.
我再看一看他衣服上.
有一個名牌.
有一個名字.
他姓周.
這個名字很熟悉.
很有印象.
但是我又想不起來.
在哪裡見過這個名字呢?.
回家上網查.

$^{81}$原來這位周先生.
就是剛剛十年前.
2013年.
有一宗很轟動的案件.
就是大角咀的碎屍案.
這位周先生.
他很生父母的氣.
謀殺了父母.
將他的碎屍.
將父母的人頭.
切下來.
放在雪櫃裡.
當時很轟動的一宗案件.
如果大家最近有看一部電影.
叫做《正義迴廊》.
就是以這宗案件作為藍本.
我再看一看這位周先生.
為什麼會做出這樣的慘案呢?.
原來他.
至少被父母迫他學琴.
學琴又不能做運動.
他很生父母的氣.
迫他學琴.
做不到運動.
人就比較受虐.
可能追不到女孩子.
後來父母送他去澳洲讀書.
去到澳洲被同學欺凌.
他很生父母的氣.
懷恨在心.
經過測試.
他的智商是很高的.
120多.
即是高於常人.
所以他整宗案件.
他不是一時衝動去殺父母.
他是處心積慮去謀殺父母.
當然根據一些專家分析.
他的性格障礙.
他的精神問題.

$^{121}$都是這宗慘案的成因之一.
不過令我很多感觸的就是.
原來憤怒會引致這麼嚴重的後果.
你想想.
學琴而已.
送他去外國讀書.
在座有多少父母這樣做?.
你有否給你的子女學琴?.
有沒有送到外國讀書?.
都有啊.
很多都有.
用不著憤怒成這樣呢?.
當然在座也有小朋友在.
千萬不要這樣.
不過原來憤怒會演變成仇恨.
仇恨會掩蓋你的人性.
有時候會造成一些無可挽回的大錯.
所以憤怒是要處理的.
其實聖經早就提醒了我們.
如果你熟悉聖經.
在人類第一宗的兇殺案.
就是該隱殺了阿鵬.
是兩兄弟.
人類第一宗兄弟.
第一對兄弟.
為什麼要殺了弟弟呢?.
因為該隱獻祭.
耶和華看不中.
不閱立他和他的共物.
但是就閱立他弟弟的共物.
他就很生氣.
在聖經創世記四章五節.
那裡說他就是該隱.
非常生氣.
他很生氣.
生氣到一個地步.
要殺了弟弟.
所以原來聖經早就告訴我們.
憤怒的後果很嚴重.
可以無可挽回的大錯.

$^{161}$憤怒仇恨的確會掩蓋你正常的人性.
引致殺人.
親人也無情講.
防範於未然.
我們憤怒應該要及早處理.
所以耶穌在登山寶訓的這一段.
就提醒我們.
憤怒是要及早處理.
這一段講一點背景.
因為我知道王震是有一個系列做門徒.
給當代門徒的十五個忠告.
在一月中.
上個月中.
Ray 牧就講了第一個忠告.
用了五章的十三到十六節.
提醒我們做門徒的第一個忠告.
有沒有人記得是什麼?.
沒有什麼印象.
是叫我們做世上的炎和光.
我們的信仰要行出來.
我們要照耀於人前.
上個星期應該記得.
上個星期真真牧師提醒我們.
他用了上一段.
五章十七到二十節.
給我們第二個忠告是什麼?.
不要緊.
大家回去可以重溫一下.
聽到我們要記得才重要.
還有記得才能夠行道.
上個星期真真牧師講什麼呢?.
我們是要行出來.
不過不是行舊約的那些律法規條.
那些不是錯的.
不過要行耶穌承傳了的律例.
耶穌提出更高的義.
就是我們新約的教導.
所以我們要行.
耶穌所提出重新演繹的更高的義.
所以耶穌接著在五章的二十一到四十八.

$^{201}$講了六個對比.
舊約的律例就是這麼說.
古人就這麼說.
有話說了.
但是耶穌就這麼說了.
有另外一些更高的教導.
祂要求我們要達至更高的義.
目的就是在這段的最後一節.
就是五章的最後一節.
第四十八節.
就是讓我們完全.
好像我們天父完全一樣.
這就是大段落的一個概覽.
今天我只會分享五章二十一到二十六.
就是耶穌所說的第一個對比.
律法叫我們不要殺人.
不過耶穌更加叫我們不要恨一個人.
我們要處理心裡面的憤怒.
進入經文在第五章二十一節.
「你們聽見有吩咐古人的話說不可殺人」.
在你的程序上也印了.
出自《初一及記》二十章十三節.
或者在《申命記》五章十七節.
也重申了一次.
不可殺人.
這是十誡中的第六誡.
很清楚.
這是在舊約的很重要的律例.
然後又說了.
「凡殺人的難免受審判」.
這可能是出自《初一及記》二十一章十二節.
或者《利未記》二十四章十七節.
提到如果你打死人.
就要被處死.
就是你殺人是有審判有後果的.
不可殺人是律法的義.
是對的.
耶穌完全沒有否定.
也沒有推翻祂.
耶穌沒有說容許我們殺人.

$^{241}$但是耶穌要求我們的義.
要怎樣呢?.
要勝過民事和法理裁人的義.
要勝過舊約律例的義.
所以耶穌接著提出一個更高的義.
這就是今天耶穌希望.
對於每一個跟從耶穌的人的要求.
在第二十二節.
祂說「只是我告訴你們」.
你經常看到這個模式.
舊約有古人是這樣說.
「只是我告訴你們」.
就是耶穌要求更高的義.
「凡向弟兄動怒的難免受審判」.
耶穌用了弟兄.
整個段落都是用弟兄.
你不要罵弟兄是拉加.
不要罵弟兄是摩利.
耶穌用弟兄是指屬於同一個信仰群體.
在這個群體裡面.
當然這個信仰群體可以是教會.
可以是你的家庭.
有同一個類別.
同一個群體.
我們不要這樣彼此對待.
我們不要生一個人的氣.
你生他氣就難免受審判.
上一節說到你殺人就難免受審判.
這一節說到你生人就難免受審判.
耶穌很明顯將兩件事放在一起.
就是說你生人就好像殺人一樣.
弟兄姊妹,你有沒有生過人?.
有小朋友轉過頭來.
你從來沒有生過人的嗎?.
你說不是.
OK.
如果不是的話,我真的會….
我沒有轉過頭來.
OK.
如果不是的話,完了崇拜要請教你.

$^{281}$能夠不生人.
每個人都生過人.
但在在應該沒有人殺過人.
我相信.
你有沒有這麼嚴重?.
生人就好像殺人一樣.
不過剛才我已經說了.
該人殺那百就是這樣.
他只是生弟弟的.
那位周先生殺父母一樣.
他是生父母的.
耶穌一針見血.
祂提出更高的義.
不要生人,不要恨人.
很重要.
因為如果我們不去處理裡面的怒氣.
裡面的憤怒.
祂一直累積發酵.
可能引致.
我相信是一個非常非常不理想的後果.
甚至是無可挽回的後果.
耶穌說.
你凡罵弟兄是拉加.
難免公會的審判.
拉加是原文的一個音譯.
拉加這個字就是廢物.
廢柴.
飯桶的意思.
然後.
他說如果你罵弟兄是摩利.
摩利也是原文的音譯.
摩利這個字就是白癡.
蠢材的意思.
如果你罵弟兄是摩利.
就難免地獄的火.
我想問問在座.
你有沒有罵過別人是不飯桶的?.
你有沒有罵過別人是廢柴的?.
有人說髒話了.
你有沒有罵過別人是白癡的?.

$^{321}$我經常被我兒子罵.
問題是要受公會的審判.
以及受地獄的火.
耶穌有沒有誇張一點?.
你罵別人是飯桶.
你就要去地獄.
當然在這裡不一定要照字面的解釋.
這裡的意思其實是.
後果是很嚴重的.
憤怒,仇恨是可以引致殺人的.
後果很嚴重.
我希望大家有耶穌的眼光去看待這件事.
今天的重點.
我不說這段篇幅.
你也明白.
你看你也明白.
不過希望大家很認真去看待.
耶穌很重視的事情.
這也是我們人生的態度.
耶穌重視的事情.
我們應該重視.
耶穌不重視的事情.
我們也同樣應該不要這麼重視.
耶穌有一些很外顯的罪行.
外顯的行為.
就是殺人.
現在教導我們轉移.
看我們內心的心態.
內在的心態.
就是生氣.
這就是更高的義.
一般的大方向也是這樣.
是處理內在的心態.
我們不要單單處理外在的行為.
一些神不喜歡的事情.
其實聖經也有說.
詩篇九十篇.
有罪孽.
但是也有隱惡.
可能你能夠看見罪孽.

$^{361}$你打人.
殺人.
隱惡在你內裡.
你會看不到一種仇恨,嫉妒.
可能在教會.
你坐在身邊的人.
可能內心看不到.
我們對神不喜歡的事情.
我永遠都有這個心態.
要消滅在他萌芽的階段.
不要讓他在內裡.
可以長大.
可以發芽.
要一早處理.
一早消滅.
那如何處理呢?.
耶穌在第23-26節.
用了兩個例子.
給我們一些指引.
第23節.
他說一開始用「所以」.
就是表示23-26節的段落.
這兩個例子的教導.
就是22節的教導.
一個很自然的結論.
我們應該要這麼做.
我們先看23-24.
第一個例子.
「你們在祭壇上獻禮物的時候.
如果想起弟兄向你懷怨.
就把禮物留在壇前.
先去和弟兄和好.
然後來獻禮物」.
雖然整段都是弟兄.
從頭到尾都是弟兄.
但當然一定包括姐妹.
你不要以為姐妹可以脫身.
這是一個涵蓋性的名詞.
弟兄姊妹.
你留意細心看.

$^{401}$這裡不是說你去惹人生氣.
這裡是說.
你知道別人生氣.
你知道有人向你懷怨.
或者和修羅版用「懷恨」.
有人恨你,生氣你的時候.
你要先去尋求和好.
然後才去獻祭.
意思就是.
如果你不先和人和好.
你獻祭.
神也未必悅立.
神也未必喜歡.
所以你看看神重視的是什麼?.
是我們內在.
弟兄姊妹之間一個和好和諧的關係.
是多於外在的一些表面上的禮儀.
不是說那些不重要.
但神更加重視什麼呢?.
如果你想起這一陣子.
在旺鎮裡面.
有弟兄,有姐妹在生你的氣.
你會不會想想.
先和他們和好.
再回到旺鎮崇拜呢?.
我們可能有很多顧慮.
你問我,我也有很多顧慮.
如果要和人和好.
然後才去敬拜,獻祭.
我會想.
我只是想和好而已.
他會不會原諒我?.
我又會想.
我費盡唇舌.
行不行呢?.
是不是一定達至和解呢?.
真的不知道.
我們也寫不到保單.
我也試過想和人和解.
對方是不願意和解.

$^{441}$不願意和你和好.
他要是生你的氣就生你的氣.
或者生你的氣一段時間.
不過我的態度就是.
我們不做.
你一定不知道行不行.
但你做了嘗試和解.
你就知道行不行.
我相信如果神喜悅的東西.
神的恩典一定夠用.
一定能夠幫助你.
可能不是即時和解.
可能怨恨積累了很久.
你也要花相當的努力.
相當的日子.
可能花幾年.
甚至十幾年才做到.
不過.
這是上帝的心意.
這也是王浸很想給各位門徒的一個忠告.
我們要及早處理.
我覺得和解的機會是大一點的.
成功的機會是大一點的.
你越遲和解.
怨恨積得越深.
大家都明白這個道理.
不過最重要的不是成數大與小.
最重要的是耶穌基督喜悅我們這樣去做.
神是不喜悅我們不和好.
神是很重視和好.
所以為什麼耶穌要來到這個世界呢?.
祂是要令人和神和好.
弟兄姊妹.
耶穌舉了另外一個例子.
在第25和26節.
第25節.
你和告你的對頭還在路上.
告你的對頭就是你的控訴者.
他在控告你.
就在說與他和息.

$^{481}$如果你看原文.
和息這個字是放在這一句.
第25節的最開頭.
是一個強調的作用.
原來和好是耶穌在這裡想強調.
你要和好.
跟你告你的那個人.
如果不和好.
有什麼後果呢?.
恐怕他把你送給審判官.
審判官交付牙役.
你就下在監獄.
我實在告訴你.
若有一文錢沒有還清.
你斷不能從那裡出來.
如果我們不和好.
你要承擔你虧欠人.
又不去和好的後果.
下在監獄還清才出來.
一文錢在原文的音譯就是科特拉.
科特拉是一個羅馬的閉籍.
科特拉有多少呢?.
大家可能比較多聽就是德拿利.
在聖經.
一個德拿利就是一個普通工人一天的工資.
如果今天可能就是幾百塊.
科特拉就是六十四分之一的德拿利.
如果今天可能就是幾塊.
你幾塊都不還清.
沒有寫出來.
這裡的意思是什麼呢?.
就是我們對人不要有絲毫的虧欠.
要盡量還清.
消除怨恨.
及早和解.
有一個故事.
我印象也深刻.
我也是聽回來的.
未必真的有這個故事.
話說有一個小學生.

$^{521}$他叫小強.
我給他起名字.
他很瘦很瘦.
在學校經常被大隻的同學欺凌他.
特別是有一個大文.
大大隻隻經常欺負他.
有一天小強放學回來就哭了.
跟爸爸說.
我又被同學欺負.
這個大文經常欺負我.
他爸爸當然很心痛.
我跟他想辦法.
如何處理這個情況呢?.
現在雙方都很敵對.
爸爸就建議.
不如這樣.
我們就請這幾個經常欺負你的小朋友.
來我家開派對.
大家開派對.
好像和解一樣.
小強就說.
不可能.
這些根本不是朋友.
他不會來我家.
不會應邀參加這個派對.
這是仇人.
特別是這個大文.
這是世仇.
爸爸就笑笑口說.
不如我們開一個仇家派對.
既然是仇人.
爸爸不僅說他真的做.
他努力安排.
努力跟他聯絡.
邀請了這幾個小朋友.
來到小強家裡玩.
結果很出人意外.
他們玩得很開心.
小強拿了玩具跟大文分享.
跟幾個以前的仇家分享.

$^{561}$平時大家的舌劍唇槍.
拳來腳往.
都沒有了.
沒有了這些場面.
大家玩得很開心.
真的打成一片.
從此以後小強就沒有再哭著回家.
他最大的仇敵.
大文.
就變成最忠實的夥伴.
沒有人再敢欺負小強.
因為有一個高高大大的大文.
現在照著他.
沒有人敢欺負他.
接著呢.
有趣的事情就是.
大文的爸爸.
事後打電話給小強的爸爸.
我告訴你一個秘密.
其實我兒子大文.
因為他太討人厭了.
從來沒有人請他去派對.
是因為他是仇敵而不被邀請.
還是因為他不被邀請.
而變成仇敵呢?.
我們不知道.
但是這個故事教導我們.
原來怒氣.
大家敵對,憤怒,仇恨.
是可以及早處理的.
也有方法處理.
我們要用神給我們的智慧去處理.
不要當這些是沒有的.
生氣就生氣.
還擊.
報復.
這是我們最常用的方法.
但是弟兄姊妹.
在這一段.
耶穌沒有這樣教導.

$^{601}$在其他.
你看看登山寶訓.
耶穌更加教導我們怎樣呢?.
要愛我們的仇敵.
弟兄姊妹.
這是很顛覆性的教導.
在我們今天.
香港社會.
一個資本主義很競爭的社會.
一點都不像我們的文化.
我們要報仇.
我們要以牙還牙.
以眼還眼.
但是耶穌不是這樣教導我們.
起碼一開始在這六個對比.
第一個已經教導我們.
要處理怒氣.
消除仇恨.
到了後面更加很白.
在其中的對比.
教導我們怎樣對仇敵.
要愛你的仇敵.
要你走一里路.
走兩里路給他.
拿你的外衣.
還給他你的內衣.
弟兄姊妹.
這個在下幾次.
應該有牧者會跟你分享.
很簡單的一段經文.
我再重申.
不說你也明白.
這段經文.
這段經文是難做的.
當我預備這段經文.
我用了最多時間想.
我怎樣做呢?.
怎樣做得到呢?.
我就用以下三個學習.
去總結這段經文的教導.

$^{641}$我也寫在結語的部分.
你參考程序.
第一個學習.
在這段經文給我自己.
就是我要追求更高的義.
我們的信仰.
我們信基督教的信仰.
肯定不是律法主義.
我們不是要來守一條條的規條.
信耶穌不要打麻將.
不要賭馬.
不要說髒話.
這些全部都是對的.
我覺得.
我們不是守這些規條的信仰.
我們是守這些信仰的精義是什麼呢?.
所以耶穌為什麼引用六條的規條.
說出更高的義.
我們要守的是更高的義.
耶穌重新演繹的.
更高的義.
我舉一個例子.
很多弟兄姊妹真的這麼想.
奉獻要給多少呢?.
十分之一.
他們很嚴守這個十分之一.
少一個「崩」不行.
不過多一個「崩」也不行.
百分之十一不對.
只要十分之一.
沒有多沒有少.
弟兄姊妹是不是真的?.
是不是這樣教的?.
新約聖經有沒有教十一奉獻的?.
沒有.
保羅不是這樣教的.
「捐得甘心樂義,是神所喜悅」.
不要作難.
不要勉強.
弟兄姊妹.

$^{681}$我們不要嚴守規條.
要捐十分之一奉獻給皇寢.
那麼十分之九是誰的呢?.
那個十分之九應該怎麼用呢?.
那個十分之九也是上帝的.
所以.
其實我自己很少.
我信了主.
很快就沒有跟十分之一的奉獻.
弟兄姊妹.
我想說的是.
我們不要用律法主義.
去過我們的信仰人生.
我們永遠都是在想.
十一奉獻的精神在哪裡.
我們要遵守.
要去走更高的義.
第二個學習.
我們要注重內在的生命.
外顯的行為是重要的.
我沒有否定.
我們外顯的行為是重要的.
你不要說我們只是講精神而已.
想我回來用粗口問候.
沒有問題的.
我很愛你.
不過我習慣講粗口.
這個我不是這樣去演繹.
不過我們更加要由外而內.
去留心我們內在的動機和心態.
我不知道你有沒有這些經歷.
我在教會一段時間.
我也有參加過會議.
我真的看到不少這些情況.
有時候弟兄姊妹之間.
表面是很和氣.
很客氣的.
甚至對你笑容滿面.
但是內心可能是在生你的氣.
可能是在咒罵你的.

$^{721}$你怎麼知道呢?.
最後你從第三者知道.
他跟第三者說.
他是在生陳牧師的氣.
我有這些經歷.
弟兄姊妹.
但是他跟你面對面的時候.
他仿佛很尊敬你.
對你很和善.
很客氣的.
我們內心的污穢要處理.
由外而內.
最重要處理的是裡面.
你處理到裡面.
你就處理到外面.
在我的媒體我看到更多的就是.
很多事奉的人.
裡面的動機是不是純正的呢?.
或者我自己也出過問題.
有時候我們的事奉.
是想別人欣賞你.
想別人肯定你.
還是藉著事奉你肯定自己.
在這個群體的價值呢?.
釐清一下.
皇帝這些全部.
上帝不喜悅.
是沒有用的.
我的意思不是我們事奉.
不理別人的看法.
不是這樣的意思.
我們謙卑聽別人對我們的意見.
但我們不是因為別人的看法.
拖著我們的事奉.
永遠不要被人拖著你走.
我們重視上帝怎麼看我的事奉.
這就是你要檢視的動機.
有時候事奉很厲害的.
上來講大家很多人信主.
說你的聚會很多人參加.

$^{761}$很成功.
但檢視是一個動機.
可能很污穢.
所以我們的信仰生命.
是注重內在.
給多點時間去檢視自己.
我的動機,我的態度,我的心態.
是怎麼樣的.
第三個學習.
就是處理怒氣要及早處理.
「早」比「遲」好.
我剛才說.
殺虛於萌芽的階段.
如何及早處理呢?.
通常處理怒氣.
我不是心理學家.
但我聽心理學家講.
有兩個大的方向.
很籠統地說.
一個方向就是壓抑.
壓抑怒氣.
適度的壓抑是對的.
我舉一個例子.
如果我回到家.
我兒子罵我一句.
接著我打他一頓.
我火都來了.
用刀砍他.
這樣不是太好.
壓抑會不會他也有難處呢?.
會不會他也有被人可能.
令他很生氣.
他發洩在我身上呢?.
如果做打工的應該知道.
職場的怒氣.
通常是轉嫁給誰呢?.
家人.
在公司上司的戲.
很多時候回到家.
跟家人鬥一輪.

$^{801}$很多時候都是這樣.
所以適度的壓抑不是錯的.
第二個大方向.
不是壓抑就是發洩.
發洩有時候.
我們找人談判.
發牢騷.
鬥一輪就舒服了.
約一個好朋友.
好兄弟姐妹.
摸摸酒杯底.
談一輪就舒暢了.
問題沒有處理.
所以適度的發洩.
適度的壓抑是好的.
我們千萬不要走極端.
過度的壓抑.
或者過度的發洩.
這就不是太好.
當然我們不走極端.
在中間落寞的時候.
我們會做什麼呢?.
我就很鼓勵.
所謂及早處理.
怎麼做呢?.
耶穌沒有教的.
這些全部都是我一直在預備.
我自己在想.
或者我正在實踐.
我常常都這麼想.
我有一個人得罪我.
我很生氣.
我經常想.
我有沒有權.
按照真理的標準.
我認識的聖經.
我有沒有權去生他的氣呢?.
我應不應該生他的氣呢?.
如果應該的話.
我就很心安理得.

$^{841}$如果不應該的話我就要處理.
有些生氣也是應該的.
我經常都用這個例子.
如果我包二奶.
我老婆生我的氣.
應不應該呢?.
她生我的氣應該.
她不生我的氣就奇怪了.
是不是很正常?.
我們人應不應該讓神生氣呢?.
應該.
我們太過活躍.
我們太過不濟.
上帝沒有小器.
是我們該死的.
如果你看完整本聖經.
你應該很明白.
上帝已經寬容到不得了.
弟兄姊妹.
我們人與人之間.
我會鼓勵如何及早處理怒氣呢?.
在過分壓抑與不過分發洩的中間.
我們要想想.
按真理我有沒有權去生他的氣呢?.
應不應該生他的氣呢?.
如果真的不應該的話.
我真的要做事了.
這是我要開始行動的起點.
我怎麼行動呢?.
我第二個問題.
我通常問自己.
我能不能夠寬恕他.
如果我回答我不能夠寬恕他.
我多數都沒有事情做.
因為我知道我再找他談.
可能火上加油.
當我回答自己.
我也應該可以寬恕他了.
那麼我就行動了.
我可能約他需要見面.

$^{881}$坐下來喝杯茶喝杯東西.
大家說說前因後果.
能夠互相原諒.
很多時候都是互相原諒的.
不是單方錯的.
這個世界很少一方錯的.
弟兄姊妹.
及早處理有很多方法的.
這是其中我正在用的方法.
我最近有一個經歷.
這個經歷也不是第一次.
我之前在疫情很嚴重的時候.
我有一次去一家快餐店吃東西.
有一個人咳嗽.
咳嗽是很正常的.
不過特別之處是除了口罩才咳嗽.
他咳嗽完就戴上口罩.
我心想很奇怪.
坐在我旁邊除了口罩咳嗽.
這是最近發生的.
我最近坐巴士.
同樣的有一個人也是咳嗽.
他除了口罩咳嗽.
咳嗽完之後就戴上口罩.
我也是在想.
我正在預備這段經文.
究竟是壓抑還是發洩呢?.
千萬不要過度.
如果過度壓抑.
你會很緊張.
會不會我中招.
然後可能心跳加速.
然後頭暈.
這是過度壓抑.
過度發洩就是罵他.
真的有一個乘客真的罵他.
幸好他罵得不是很激烈.
我也感恩.
我又不是站好.
我會怎麼做呢?.

$^{921}$一就是你不行就調位.
我覺得.
你可以去.
你坐樓上還是坐樓下.
你坐樓下就是搬樓上.
一就是你可以勸他.
如果你覺得可以勸喻.
勸他先生.
或者如果你真的咳嗽.
下次就戴上口罩.
有時候不用調位.
坐在旁邊也可以的.
當然勸人戴口罩.
我幾乎是屢試失敗的.
不是這一年.
在去年疫情很高峰的時候.
連做運動也要戴口罩的時候.
我試過勸一個人戴口罩.
他不戴口罩.
沒有人聽.
一個也沒有.
不是有一個人聽.
有一個老伯肯聽.
不過我猜他看我這麼高大.
他是怕我.
為什麼呢?.
因為他戴了一會就馬上脫下來.
你又不是警察.
我戴不戴關你什麼事呢?.
所以我也反省.
罪就是這樣.
勸也未必聽.
不過我們不要靠得太遠.
在這個例子告訴我們.
我們可以處理.
如果你.
因為這樣的場合.
你跟陌生的一位乘客.
暫時有一個敵對的心態.
也要處理的.

$^{961}$是不是?.
我就建議.
一邊搬位一邊勸他.
一邊不出聲.
不要過激.
這些很小事處理不了.
因為兩個站之後.
大家都下車.
各散東西.
你又不知道他什麼名字.
也不會認得他的樣子.
他又不會認識我.
大家沒事的.
但如果是你的親人.
是旺鎮的弟兄姊妹.
是你的同事.
朝見口晚見面的.
我們為什麼不這麼做呢?.
最傷氣的不是被氣的.
是氣人的.
他也不知道你氣了他十年.
你自己又無緣無故.
辛苦了十年.
更加重要的原因就是.
耶穌不起怨.
耶穌不起怨我們這樣的關係.
所以耶穌今天這段經文很實際.
也很簡單提醒我們.
祂希望大家記得第三個忠告.
及早處理你的怒氣.
你的一生也會平安喜樂很多.
我們重新禱告.
你的智慧不是我們能夠想像得到.
你每一個的教導.
對我們都是畢生的受用.
求主幫助我們.
我們願意去處理我們的內心.
我們願意去檢視.
連我們最陰暗的地方.
都願意被你光照.

$^{1001}$我們願意你幫助我們.
去直接正面地面對我們.
對人的嬲怒.
或者處理別人對我們的恨意.
讓我們和好.
好像與你和好一樣.
當我們彼此和好.
我們整個群體就能夠照耀於人前.
成為一個美好的見證.
讓更多人想認識你.
追求你.
渴慕你.
同心禱告奉耶穌基督的聖名.
\newpage



\section{馬太福音 5:27-32}
\label{sec:N8TX9PAocpk}
\textbf{2023年2月12日崇拜直播｜陳冠成牧師｜給當代門徒的15個忠告:持守聖潔｜馬太福音5\_27-32}
\newline
\newline
連結: \href{https://youtube.com/watch?v=N8TX9PAocpk}{\texttt{https://youtube.com/watch?v=N8TX9PAocpk}} ~~~~ 語音日期: 2023-02-13
\newline
\newline
\hyperref[sec:iHCyyrWPaj4]{\small{< < < PREV SERMON < < <}}
~
\hyperref[sec:index_chronic]{\small{[返順時目]}}
~
\hyperref[sec:VsBsdArDKAI]{\small{> > > NEXT SERMON > > >}}
\newline
\newline
馬太福音 5:27-32
\newline
\begin{longtable}{cl}
\hline
\hline
章節 & 經文 (和合本修訂版)\\
\hline
5:27 & \begin{tabularx}{0.7\textwidth}{X} 「你們聽過有話說：『不可姦淫。』 \end{tabularx} \\ \\ \relax
5:28 & \begin{tabularx}{0.7\textwidth}{X} 但是我告訴你們：凡看見婦女就動淫念的，這人心裡已經與她犯姦淫了。 \end{tabularx} \\ \\ \relax
5:29 & \begin{tabularx}{0.7\textwidth}{X} 若是你的右眼使你跌倒，就把它挖出來，丟掉。寧可失去身體中的一部分，也不讓整個身體被扔進地獄。 \end{tabularx} \\ \\ \relax
5:30 & \begin{tabularx}{0.7\textwidth}{X} 若是你的右手使你跌倒，就把它砍下來，丟掉。寧可失去身體中的一部分，也不讓整個身體下地獄。」 \end{tabularx} \\ \\ \relax
5:31 & \begin{tabularx}{0.7\textwidth}{X} 「又有話說：『無論誰休妻，都要給她休書。』 \end{tabularx} \\ \\ \relax
5:32 & \begin{tabularx}{0.7\textwidth}{X} 但是我告訴你們：凡休妻的，除非是因不貞的緣故，否則就是使她犯姦淫了；人若娶被休的婦人，也是犯姦淫了。」 \end{tabularx} \\ \\
[1ex]
\hline
\hline
\end{longtable}
$^{1}$《新約聖經》第6頁.
第27節.
「你們聽見有話說,不可姦淫,只是我告訴你們,凡看見婦女就動淫念的,這人心裡已經與她犯姦淫了.若是你的右眼叫你跌倒,就腕出來丟掉,令可失去八體中的一體,不叫全身丟在地獄裡.若是右手叫你跌倒,就撼下來丟掉,令可失去八體中的一體.」.
「若是右手叫你跌倒,就撼下來丟掉,令可失去八體中的一體,不叫全身下入地獄.又有話說,人若休妻,就當給他休戍,只是我告訴你們,凡休妻的,若不是為淫亂的緣故,就是叫他作淫婦了.人若娶這被休的婦人,也是犯姦淫了.」.
聖經獨筆.
弟兄姊妹平安,上次Paul木託問大家,大家已經有所準備了,他上次透過耶穌基督的教導提醒我們什麼?處理怒氣,還要是及早處理怒氣,耶穌表示上帝不可殺人的吩咐,其實不單止針對人外在的行為,.
其實也要人正視他內心,因為殺人往往由於人沒有好好處理自己的怒氣,任由他積存甚至爆發,最終明明只是憎恨對方,最後就變成殺害對方..
這就是耶穌在登山寶訓裡面,關於承傳律法的首個實際應用,今天我們會繼續看之後的兩個和耶穌承傳律法有關的吩咐..
剛才所讀的經文,一開始第27節說,不可姦淫,其實是來自舊約十誡第七條的誡命,猶太人很看重這個吩咐,認為犯姦淫是一個很嚴重的罪行,被捉拿的人會被拉去致死,不過他們往往著重的是這條條文裡面的字面意思,.
換句話來說,他們覺得只要沒有姦淫的行為,就不會抵觸不可姦淫這個上帝的吩咐,不過耶穌表示上帝要求的是一個更高的義,耶穌說,凡看見婦女動淫念的,這人心裡已經與他犯姦淫了..
當然,耶穌並不是針對婚姻中正常的性關係,也不是要禁止所有人以後不可以看異性,看異性,耶穌說,看見就動淫念,其實原本有兩個意思,.
第一個就是當看見異性的時候,故意讓自己幻想出色情的畫面,這是第一個,而另一方面就是故意向對方拋眉弄眼,故意用你的眼神挑起對方的情慾,引誘他犯罪,.
這樣做的話,即使沒有發生任何肉體上的姦淫,但是按照耶穌的標準,那人心裡已經犯姦淫了..
我們看到耶穌在這裡再一次表明,上帝律法的原意不只是針對外在的行為,而是要人進一步檢視他的內心想法和動機,既然不可殺人是提醒人要及早處理對人的憎恨和憤怒,.
換句話來說,不可姦淫也是提醒人要及早正視對人的非憤之想,耶穌在這裡說,若是你的右眼叫你跌倒,就怎樣?.
「寧可失去白體中的一體,不叫全身丟在地獄裡.」又或者說,若是右手叫你跌倒,就怎樣?.
「寧可失去白體中的一體,不叫全身下入地獄.」.
過往確實在教會歷史裡曾經有基督徒真的照著這幾節經文的字面意義去做,.
為了禁止自己動淫,唸甚至犯姦淫,將自己弄至傷殘..
最極端的例子,可以說是公元第三世紀,可能大家也聽過他的名字,教父何利根,有沒有聽過他的名字?.
他是一個很主張極端禁慾主義的信徒,他不單止放棄自己的財產,過著禁食,不讓自己睡覺,不眠的苦行生活,.
甚至他會過度用字面意義來解釋剛才我們所讀的那段話題的復音,.
他為了不讓自己犯姦淫甚至動淫,唸,你知道他怎樣做嗎?.
他將自己變成一個閹人,很可怕吧?.
所以在以後的尼西亞會議就禁止這種極端的做法..
當然,頂子妹耶穌也不是叫我們按照字面意義來理解,.
將那些可能讓你陷入情慾的,試探的眼,手,甚至腳,通通從你身上撤除,.
如果這樣會怎樣?很大件事,可能你身邊看到的不單是男士或女士都會怎樣?.
都會有殘障,可能沒有眼,手,甚至腳..
再說,就算真的將眼,手,腳劃走,.
其實都阻止不了人用腦袋去犯姦淫,對不對?.
所以耶穌這裡說「劃眼劈手」,其實是一種誇張的表達手法,.
是很想帶出那種迫切性和嚴重後果,.
提醒門徒要及早將淫念消滅於萌芽,盡你一切努力,.
想盡辦法,不要讓它有機會在你身上出現,.
不要讓它留下,不要讓它拖累你,這是耶穌的意思..
套用在我們今天的生活場景,耶穌可能會這樣吩咐我們,.
假如你每次路過旺角,油麻地,你見到那個招牌就想起壞事,那怎麼辦?.
下次提醒自己不要抬頭望,情況就好像你畫了隻眼出來,.
又或者假如那個網站會令你聯想起很多色情畫面,.

$^{41}$那怎麼辦?下定決心控制你的手不要揮進去,情況就像你劈了手,.
又或者你每次路過某個報紙檔擺放的雜誌,.
有些會令你挑起一些對性的幻想,幫自己下次繞路走,不要經過,.
就好像你連腳也斬了一樣,耶穌是這個意思,.
簡單來說就是祂要求我們早一步拒絕那些構成情慾試探的事物,.
不要讓它靠近,這樣做很可能這個世界會把我們當成傷殘人士,.
意思是今天很多潮流文化裡覺得沒問題的事,可以放膽去做的事,.
例如無論你想不想,你都會發現今天很多的電視節目或者網上資訊都是用什麼來做賣點,.
吸引你去看,可能用小鮮肉或者性感來吸引你去看,.
展示它的身材,讓你去瀏覽,基督徒卻要在這種文化處境裡,.
你要持守聖潔,你要接受很多的限制,選擇不去做,不去看,.
因為耶穌提醒我們為了得到真正的生命,得到上帝的喜悅,.
寧可摒棄在今生的罪終之樂,總好過將來要承擔被毀滅的後果,.
耶穌說得很清楚,所以為此緣固定,這段經文提醒我們,.
其實是需要在生活裡設立一條界線,留意,如果我們細看這段經文,.
耶穌沒有說所有人都要將他的右眼右手,那些很可能令你犯罪跌倒的事情,.
全部一刀切,意思是,譬如對A來說,他可能見到對方穿背心,穿緊身褲,.
見到這樣的異性的時候,其實他一點反應都沒有,但是對B來說,.
同一件事情可能會令他很快上入非非,又或者對有些人來說,.
他一個人在房間上網,不會出事,但是對另一些人來說,只要一上網,就立刻墮網,不能夠止拔,.
耶穌的意思是,原來不同的信徒都會面對不同的情慾試探的情況,.
所以做法不是一刀切,禁止所有的基督徒不准你上網,不准你看電視,不准你看電影,.
教會從此禁止我們接觸任何的社會文化,耶穌不是這個意思,.
耶穌是很想我們依靠祂,來到適當地為自己立一條安全線,確保我們不會出了界,越軌,得罪了上帝,.
所以說到如何去立界線,其實這個經歷不難理解,我曾經分享過一個例子,.
大家有沒有走過行人天橋,行人天橋兩邊一定有圍欄,有牆擋住,是不是要管轄你?.
那兩面牆是想保護你,你試想想,以前有沒有走過沒有頂的行人天橋,兩邊都扶手,.
你想想如果那條天橋沒有兩邊的圍欄,你夠不夠膽走到哪裡?還夠不夠膽?不夠膽,.
你多大都保持起碼一米距離,你怕跌下去,但是如果有了兩邊的圍欄,你靠著它,扶著它都夠膽走,.
你走闊了,因為你知道那條圍欄是保護你,不是攔阻你,你走近了都沒事,.
同樣為我們設立一條界線其實是一個保護,確保你可以在界線裡面很安全很自由地行走,.
不會跌死,不會失足,所以這段經文提醒我們,你真的要為你,可能你在某方面有情慾試探,.
你真的要設立一條界線,我記得大概十多年前,那時候我住在荃灣,你知道荃灣有些紅Van直接來旺角,.
每次可能七點多就上小Van,來到旺角其實都是早上八點多,很早,.
那輛車通常經上海街,山東街或者豉油街交界,我每次叫下車都會在上海街,山東街和豉油街交界,.
我一下車那一次,你知道那裡有很多唐樓,我走回教會,突然間在唐樓裡面有一男一女走下來,.
那男的問「先生,需要服務嗎?」,原來這麼早都有這些情況,我當然不回應,馬上急急走回教會,.
那以後怎樣?是不是不坐小Van,以後坐地鐵,又或者你老早,亞皆老街都未到就叫下車,.
你確保你不會再繞到那裡又遇到這些試探,所以同樣的,姐妹,譬如又說,.
可能一會兒你無論男女都好,崇拜完你下街,一落街就見到一位很性感打扮的異性,.

$^{81}$向你迎面而來,你不看都看了,那怎樣設立界線?怎樣算好?.
很值得去想的,因為你避不了的,第一件事你不要再繼續看,.
還有你跟耶穌說什麼?耶穌眼前所見的真的很吸引,不過求你幫助我,停了,夠了,不要再看下去,不要再想多了,.
事實上原來你帶著耶穌去是差很遠的,很多事情,有時我們不單單只需要設立界線,.
其實單靠人的意志你很難去守住,特別我們曾經有些經驗,當你提醒自己不要做這件事,不要做那件事,不要繼續的時候,.
你越是說不要做,你的腦袋就會怎樣?你越是想起那件事,所以你帶著耶穌去才是最安全最穩陣,.
當你遇到試探的時候,你馬上跟耶穌說救我,讓祂監察你,讓祂掌管你的思想,你的眼,你的手,你的腳,讓我們不要再亂來,.
當然你更加可以邀請你身邊的人幫你守住你的界線,譬如說家裡有子女,盡量鼓勵他們上網的時候,在哪裡上?.
在客廳上,你電腦設在客廳上,每次上網都要面壁看著牆,那部電腦是向著所有人的,那他就不會亂來,他就知道有人看著他,.
自然會謹慎很多,減低他墮網的風險,又或者你可以找一些你信得過的人,找牧者,或者找你信得過的弟兄姊妹,你跟他們分享你在面對這些試探的時候,.
你的掙扎和難處,讓他們可以定期跟你傾談,為你代禱,守望你,真的有效的,因為你會發覺原來有人願意撐你,跟你同行的時候,其實你會有很多力量來抵抗情慾的試探,.
這是耶穌關於不可姦淫所引伸的正視內心的教導,我們繼續看下去,耶穌從不可姦淫引伸到有關婚姻甚至離婚的教導,.
第31節說人若休妻就當給他休書,其實同樣是來自舊約的申命記,第24章1節,程序表裡面也印了這段經文給大家,.
這裡說人若娶妻以後,見他有什麼不合理的事,不喜悅他,就可以寫休書交在他手中,打發他離開夫家,.
特別是在男性主導的社會,當丈夫認為妻子有不合理的事情,不喜歡他,原來他有權讓他簡單地寫幾頁字立一封休書,然後在兩個見證人面前將休書遞給妻子,.
這樣就成為了休妻,過程是不是超級簡單?真的很簡單,但是問題來了,什麼是這段經文所說的不合理的事?.
當時猶太人有兩派的解說,屬於嚴謹的薩米拉比學派,他們認為不合理的事第一件事就是不忠,換句話說,他們主張除非妻子對丈夫不忠,否則不可以休妻,.
至於另一派,屬於寬鬆的希烈拉比學派,他們認為不合理的事包括不足掛齒的小事過失,例如一餐太太煮飯加了鹽,不好吃,糟蹋了食物,他就可以休妻,.
老婆說話大聲一點,又可以休妻,又或者覺得老婆不夠吸引,其實是他認識了一個更加吸引的,他又可以提出休妻,換句話說,這麼寬鬆的學派主張,幾乎等於丈夫可以用任何理由撇開太太,.
所以,問問大家,你猜耶穌時代哪一個學派對休妻的說法比較有影響力?哪一派?寬鬆那一派還是嚴謹那一派?.
寬鬆那一派比較受歡迎,所以大家讀一段經文,印給大家,中間那段馬太福音的十九章三至九節,預備起.
(教學).
你看到法利賽人專注在離婚的理由,何謂合法離婚?但耶穌是專注在上帝設立婚姻的原意是怎樣?.
所以耶穌反問法利賽人有沒有看過《創世記》,上帝起初做男做女,並且說人要離開父母,與妻子聯合,二人成為一體..
當人企圖在律法中捐瓶捐瓶,找漏洞轉空子,希望將休妻合理化時,耶穌卻重新為婚姻關係定調..
耶穌說得很清楚,夫婦不再是兩個人,乃是一體了.婚姻的目標不是分開,而是合一..
第二方面,耶穌說上帝所配合的人不可以分開,婚姻是神撮合的,人不能擅自用任何手段解除這個關係..
耶穌斬聽鐵鐵說了這兩點,但你看到法利賽人不同意,他們反駁,說不是,休妻是合法沒問題的,因為摩西吩咐我們這樣做,.
我們休妻是遵守律法的.耶穌指正他們說,其實這條規條不是鼓勵人休妻,因為上帝根本不期望人這樣做,.
只不過是人的心硬,摩西很無奈才作出讓步.所以你看到耶穌在馬太福音5章的第32節和第19章的9節,總結了他對休妻的看法到底是怎樣..
第5章32節說:「只是我告訴你們,凡休妻的若不是為淫亂的緣故,就是叫他作淫婦了,人若娶這被休的婦人,也是犯姦淫了.」.
第19章9節說:「凡休妻另娶的,若不是為淫亂的緣故,就是犯姦淫了,有人娶那被休的婦人,也是犯姦淫了.」.
首先我們要處理的問題是,為何耶穌說,無論是休妻另娶的,又或者娶了被休的婦人,甚至是無辜被休的婦人,三方面都是犯姦淫呢?很奇怪..
首先我們要明白當時的背景,當時有沒有贍養費?沒有.當時對被休的婦人是沒有社會保障,她不會得到贍養費,她甚至不能回娘家..
被休的婦人往往需要通過一個方法來繼續獲取社會的保障,就是再婚.她唯有依靠另一個男士成為丈夫,她才可以在社會生活..
所以耶穌的意思是,休妻其實是變相逼迫被休婦人再婚,她無法選擇.另一方面,剛才我們已經說了,夫婦的關係是上帝所配合的,人不能擅自用方法來解除..
即使律法容許人離婚,好像讓某一對夫妻可以結束他們的婚姻關係,但是在上帝的眼中,已經完結了嗎?即使律法容許他們這樣做,但是在上帝眼中,那對夫婦的婚姻關係從來沒有終止過..
他們仍然是對方的配偶,就算寫了休書.換句話來說,耶穌的意思是,假如跟配偶離婚之後再婚,又或者跟離婚的人結婚的話,在上帝眼中其實是什麼問題?是重婚.所以耶穌說是犯姦淫..
而唯一的例外,就是當二人成為一體的關係,他遭到破壞,譬如婚外情不忠,甚至我們今天所說的虐待,虐打,這個二人成為一體的盟約遭到破壞受損..
在這個情況下,耶穌容許,但也未必是鼓勵他們,他容許他們選擇離開,甚至再婚,而不構成舊約犯姦淫的問題..

$^{121}$所以你看到,弟兄姊妹,上帝設立的婚姻制度,在古今中外的社會,其實都面對不同情況的挑戰,衝擊,甚至是瓦解..
特別你看看最近的調查,在疫情之前,2019年的結婚數字,有大約4萬4千宗結婚的個案,離婚的有多少?你猜猜,樹根可能清楚一點做統計,有2萬1千多宗..
換句話說,離婚率有48\%,即是接近每兩對結婚的,就有一對離婚,當中亦有不少基督徒夫婦..
當夫妻在婚姻關係上覺得不滿足,對對方有很多不滿,開始嫌棄對方,不喜歡對方,甚至想離開對方的時候,.
特別是剛好在對方身邊出現了一個所謂的真命天子,對他更加關心,更加懂得欣賞他,對他更好..
在這個時候,世界往往會教導甚至鼓勵你怎樣做?跟他分開,離婚,為何要讓他拖累你的幸福?.
正當很多人選擇用律法來解決婚姻關係的時候,耶穌卻是重申,其實婚姻沒有什麼退路,但可以靠上帝有很多出路和盼望..
耶穌其實背後是這個意思,耶穌說「二人成為一體」,其實就是提醒夫婦雖然結婚之後會住在一起,但住在一起不會自動成為一體..
始終兩塊不同質地的物料要把它們混合成一塊,裡面是一個過程,充滿了藝術和挑戰,一定會有不順暢,有掙扎的時候,.
因為大家在性格上,習慣上,甚至很多價值觀上都有不同的看法,需要很多時間來適應和對方磨合..
舉個例子,大家有沒有試過打磨?男士可能多一些,去磨一件物件,磨好它..
你需要很多技巧,很多耐性,要一直轉沙紙,由粗轉到幼,再去拋光,整個過程還要不斷加水和潤滑劑,不少得這兩樣東西..
因為要確保整個打磨磨合的過程要暢順,不會磨爛那件物件,也不會發過熱..
同樣,夫婦關係要磨得亮麗,其實很需要上帝的幫助,需要常常透過上帝的愛來滋潤這段關係,.
讓上帝的話語來調教改變我們的言行,讓上帝的恩慈可以包住我們,以致我們可以用恩慈彼此相待..
這樣婚姻的磨合才會暢順,不會經常卡住或是絕住了..
耶穌這裡說,神所配合的不可以分開,祂不是要消極地讓每一對基督徒夫婦,我的底線就是不和他們離婚就算了,因為上帝說不可以分開,上帝的意思不是這樣..
上帝的意思是很正面的,祂基本上替每一對在祂面前立了婚盟的夫婦,祂幫我們封了分開這條路,此路不通,但是有另一條路,祂開了,祂會幫我們成為一體..
所以,祂強調的是不要純粹不離婚就成為我們的婚姻觀,祂強調的是很正面地去竭力地黏在一起,祂會幫我們..
當我們遇到問題的時候,祂會幫我們走在一起,所以當婚姻遇到問題的時候,當彼此都不想靠近對方的時候,我們首先要怎樣做?.
我們先靠近上帝,靠近那位設立婚姻的主,靠近那位撮合你們成為夫婦的主,這個才是重點..
當我們這樣做的時候,上帝的愛就會充滿我們,上帝的話語就會掌管保守我們,首先不會再做那些魯莽,衝動,羞愧的事..
並且也開始提醒我們,耶穌不是在登山寶訓裡面說過八福,我們要使人和睦,我們要溫柔,我們要謙卑,我們要實踐彼此相愛,彼此饒恕..
在那些爭吵不和的處境裡面,靠上帝成為和平之子..
事實上,當我們願意遵守耶穌教導,願意謙卑柔和,接納別人的脆弱的時候,當我們成為一個這樣的基督徒的時候,其實也會自然成為一個成熟稱職的配偶..
當我們先靠近上帝,上帝也會幫助我們和配偶彼此靠近..
記得那個三角形,上帝,夫婦,當夫婦二人想走近的時候,其實是怎樣呢?.
是先走近上帝,夫婦就會自然走近..
所以,頂尖會盼望我們及早讓上帝,或者讓牧者,又或者是專業的輔導,介入我們的婚姻,幫我們處理困難..
不要讓問題積累得太多,太過久..
與此同時,其實我們要想辦法為我們的婚姻來儲蓄增值..
有一句名言是這樣說的,「結婚之前,努力選擇你所愛的.」.
下一句是什麼呢?「結婚之後,努力愛你所選擇的.」.
其實與上帝是一樣的,與上帝的關係,你選擇上帝,就努力去愛他..
你選擇了對方做另一半,就努力去愛他..
上帝是這樣鼓勵,勉勵提醒我們..
雖然好像沒有什麼退路,但是上帝有很多恩典,很多出路給我們..
順帶一提,過兩天是什麼日子呢?.
大家都記得這一天,但不只是記得這一天..
預備好了沒有?有人轉頭了..

$^{161}$預備好向你的愛侶,情人表達愛意了沒有?.
可能有些頂尖的說,老夫老妻,不用了,不要浪費錢,不要被人拋棄,出去,是嗎?.
但不是的,你想想,如果你真是精心為對方去揀選一份不一定要很貴的禮物,.
你為對方一些需要,你平時有留意,就趁這一刻,你為他揀選,送給他..
如果這樣能夠為婚姻增值,其實為什麼不做呢?.
又或者說,出去吃飯很貴的,月十四,沒錯,真的很貴,那就自己做吧..
如果兩夫妻可以親手製作一頓的燭光晚餐,這樣可以增潤彼此的關係,為什麼不做呢?.
還有多少天?還有兩天,很多時間補飛,快點去吧..
盼望我們都能夠靠著上帝的愛和祂的話語,來增潤我們和上帝的關係,增潤我們和對方的關係,.
多花一些心思來為這段關係投資和增值,值得我們這樣去做,我們同心同意..
挑起上帝,我們每一位仰望,示納在你面前,願上帝你都去幫助我們,.
讓我們今天這段經文成為我們生命的提醒和幫助,讓我們首先和你結連,.
以致我們真的能夠有更多的愛,更多的恩典來和我們身邊的人,特別是我們的配偶來相處..
讓我們都有力量來靠著你,好好為我們的婚姻增值,保值,.
因為這是你對我們的吩咐和期望,即使有難處,但我們深信上帝你的恩典一定夠用..
我又求主保守我們,讓我們真的面對生活的試探,特別是情慾的試探的時候,.
讓我們真的有力量去勝過,讓我們設立合宜的界線,讓我們真的常常持守聖潔,.
保守一個清潔的心,成為合用的器皿,讓耶穌你去使用..
上帝我們恭敬,將每一位弟兄姊妹,我們仰望在你的恩手,.
我們在生命裡面都會面對不同的掙扎,不同的難處,我們都會有很多的軟弱,.
但我們知道上帝你已經接納我們,你已經寬恕我們,你也會幫助我們走近你,.
也會幫助我們和身邊的人走近,維持我們將每一個不同的境況都交在上帝你的恩手,.
願上帝你拖大我們,幫助我們能夠過一個土地喜悅,榮神益人的生命..
我們這樣的禱告交託,奉基督耶穌你得勝的名字,你也可以祈求..
Amen.
阿們.
(影片完結).
\newpage



\section{路加福音 10:25-37}
\label{sec:VsBsdArDKAI}
\textbf{2023年2月26日崇拜直播｜李盛林牧師｜誰是你的同行者｜路加福音10\_25-37}
\newline
\newline
連結: \href{https://youtube.com/watch?v=VsBsdArDKAI}{\texttt{https://youtube.com/watch?v=VsBsdArDKAI}} ~~~~ 語音日期: 2023-03-02
\newline
\newline
\hyperref[sec:N8TX9PAocpk]{\small{< < < PREV SERMON < < <}}
~
\hyperref[sec:index_chronic]{\small{[返順時目]}}
~
\hyperref[sec:O_RnYw4oPiY]{\small{> > > NEXT SERMON > > >}}
\newline
\newline
路加福音 10:25-37
\newline
\begin{longtable}{cl}
\hline
\hline
章節 & 經文 (和合本修訂版)\\
\hline
10:25 & \begin{tabularx}{0.7\textwidth}{X} 有一個律法師起來試探耶穌，說：「老師！我該做甚麼才可以承受永生？」 \end{tabularx} \\ \\ \relax
10:26 & \begin{tabularx}{0.7\textwidth}{X} 耶穌對他說：「律法上寫的是甚麼？你是怎樣念的呢？」 \end{tabularx} \\ \\ \relax
10:27 & \begin{tabularx}{0.7\textwidth}{X} 他回答說：「你要盡心、盡性、盡力、盡意愛主—你的神，又要愛鄰如己。」 \end{tabularx} \\ \\ \relax
10:28 & \begin{tabularx}{0.7\textwidth}{X} 耶穌對他說：「你回答得正確，你這樣做就會得永生。」 \end{tabularx} \\ \\ \relax
10:29 & \begin{tabularx}{0.7\textwidth}{X} 那人要證明自己有理，就對耶穌說：「誰是我的鄰舍呢？」 \end{tabularx} \\ \\ \relax
10:30 & \begin{tabularx}{0.7\textwidth}{X} 耶穌回答：「有一個人從耶路撒冷下耶利哥去，落在強盜手中。他們剝去他的衣裳，把他打個半死，丟下他走了。 \end{tabularx} \\ \\ \relax
10:31 & \begin{tabularx}{0.7\textwidth}{X} 偶然有一個祭司從那條路下來，看見他就從另一邊過去了。 \end{tabularx} \\ \\ \relax
10:32 & \begin{tabularx}{0.7\textwidth}{X} 又有一個利未人來到那裡，看見他，也照樣從另一邊過去了。 \end{tabularx} \\ \\ \relax
10:33 & \begin{tabularx}{0.7\textwidth}{X} 可是，有一個撒瑪利亞人路過那裡，看見他就動了慈心， \end{tabularx} \\ \\ \relax
10:34 & \begin{tabularx}{0.7\textwidth}{X} 上前用油和酒倒在他的傷處，包裹好了，扶他騎上自己的牲口，帶他到旅店裡去，照應他。 \end{tabularx} \\ \\ \relax
10:35 & \begin{tabularx}{0.7\textwidth}{X} 第二天，他拿出兩個銀幣來，交給店主，說：『請你照應他，額外的費用，我回來時會還你。』 \end{tabularx} \\ \\ \relax
10:36 & \begin{tabularx}{0.7\textwidth}{X} 你想，這三個人哪一個是落在強盜手中那人的鄰舍呢？」 \end{tabularx} \\ \\ \relax
10:37 & \begin{tabularx}{0.7\textwidth}{X} 他說：「是憐憫他的。」耶穌對他說：「你去，照樣做吧！」 \end{tabularx} \\ \\
[1ex]
\hline
\hline
\end{longtable}
$^{1}$接著是一段經訓時間.
經文是《路加福音》第十章25.37.
我們請到林永祥弟兄.
接著今天的講道.
是由李成林牧師考慮中間.
大家可能對他都有多少認識.
他除了是慕織臣學院的前院長之外.
大家如果有留意我們的程序.
他不停為了他的兒子Moe的身體祈禱.
我會將時間和大家分享訊息.
主題是「誰是你的同行者」.
我們先把時間交給林永祥弟兄.
.
今天的經訓是《路加福音》第十章25-37.
請聽我讀出.
《路加福音》第十章25節.
.
耶穌說:你回答的事.
你這樣行,就必得永生.
那人要顯明自己有理.
就對耶穌說:誰是我的鄰舍呢?.
耶穌回答說:有一個人.
從耶路撒冷夏耶利哥去.
落在強盜手中.
他們剝去他的衣裳.
把他打個半死.
就丟下他走了.
偶然有一個祭司從這條路下來.
看見他就從那邊過去了.
又有一個利未人來到這地.
看見他也照樣從那邊過去了.
唯有一個撒瑪利亞人.
行路來到那裡.
看見他就動了持心.
上前用油和酒倒在他的傷處.
包裹好了.
扶他騎上自己的身後.
帶到店裡去照應他.
第二天拿出二錢銀子來.
交給店主說:你且照應他.

$^{41}$此外所費用的.
我會來必還你.
你想.
這三個人.
哪一個是落在強盜手中的鄰舍呢?.
他說:是憐憫他的.
耶穌說:你去照樣行吧.
聖經獨筆.
頂主每早晨.
和大家在今天主日的分享.
我也是在過去在浸信會事奉.
不過在加拿大多倫多那邊.
今天和大家分享這一段經文.
耶穌說的比喻就是這樣.
祂面對民事發理賽人.
他們並不是同行者.
他們只不過是指手畫腳的人.
沒錯,大家都是信仰.
所以耶穌說:他們是將難擔的擔子.
放在人的身上.
但是他們是一個指頭都不動.
弟兄姊妹.
同行者是需要下水的.
在信仰上,今天我常常說.
基督徒在這個世代很少有屬靈的經驗.
很少有屬靈在每一天裡面的觸覺.
是因為我們在信仰上很少與神同行.
在信仰上更加少與人同行.
我們活在自己的圈子,活在自己的世界裡面.
和大家看看一段聖經.
耶穌所說的那個同行者.
他有兩方面.
同行者,我們最懷疑的.
就是用了一個很僵化的屬靈信仰.
以為很屬靈,但是僵化的.
沒有那種感受,沒有那種讓人看見.
當你與人接觸的時候,人們體會到那種真誠.
一種僵化的屬靈,僵化的信仰.
你接觸任何人,你都不能夠.
或者不會讓人看到你那種真誠.

$^{81}$人很容易的,當你與人接觸的時候.
你深感受到.
民事法理賽人,這些律法師是沒有的.
他們只不過是按照傳統,按照本質,按照責任去辦事.
但是耶穌說的這個比喻.
他就是要說出來一個真正的同行者.
是有著一個活潑屬靈的信仰.
你與人接觸的時候,你讓人感覺到那種生命的真誠.
你讓人感覺到,體會到,經歷到.
那種信仰的能力在裡面.
所以在這段經文.
我們可以看見的.
很多時候,教會裡面,這個世界上.
都是羊群的心態.
人怎樣做,我就怎樣做.
人這樣走,我就照大衛這樣走.
所以耶穌曾經在福音書裡面說過.
他說闊的路很多人走.
窄的路很少.
因為當你走窄的路的時候.
你要去面對著群眾.
又或者面對著絕大多數的人這樣去做.
為何你不這樣做呢.
你要有理由,理據和理能去支持.
你要走窄路的心態和意志.
在這個時代很少了.
最舒服的是大家一起走.
大家一起走,不需要任何的解釋.
因為每個人都是這樣.
當時在這段經文裡面.
我們可以看見.
比喻的基礎出現.
是因為有一個律法師起來要試探耶穌.
可能那個版本不同了,字就不齊了.
還有,在比喻裡面有一個主角.
有一個人.
耶穌著意沒有說任何人.
有一個人.
但是大家留意.
在聖經裡面所說的比喻.

$^{121}$不是一種抽象虛構的故事.
又或者是無端端想出來的.
不是.
聖經裡面所說的比喻.
是一個在現實的環境.
當時的人.
可以在實際生活裡面體驗到的.
所以在比喻裡面.
耶穌是將當時的一個景象抽出來.
然後就和這個律法師講.
除了這個人之外.
還有強盜.
真的有強盜.
為什麼有這麼多強盜在這裡呢?.
一會兒我會和大家看這段經文再解釋.
有祭司有利未人.
還有撒瑪尼亞人.
大家看得辛苦一點.
但是不要緊.
引發這個比喻的.
就是最後這個律法師問耶穌一個問題.
誰是我的鄰舍?.
深入一點和大家看看.
這個同行者.
耶穌在說什麼?.
一個同行者如果他單是有信仰.
但是這個信仰是沒有經過.
真的投入那種熱誠.
深思為別人著想的.
一個這樣的信仰.
你並不是同行者.
你只不過好像文氏,法利,賽人那樣.
一套的東西加進去.
難怪有很多人不信耶穌.
就是因為我信了耶穌有很多規矩.
我信了耶穌有很多東西是不行的.
我們很容易被人看見信耶穌就是這樣.
信耶穌.
如果我們能夠被人.
親身心底裡體會,經歷到.

$^{161}$那種人性.
和那種在耶穌基督裡願意的付出.
他不會感受到這是僵化.
不會感受到這是一套規矩.
有一個律法師.
當時在耶路撒冷裡.
有六千個律法師,六千個拉比.
這些律法師,拉比.
我們可以了解到.
在以色列人被擄到巴比倫時.
就開始有了.
這些律法師.
原本在以色列裡是沒有的.
但是被擄到外邦.
特別是巴比倫時.
沒有聖殿,沒有敬拜,沒有節日.
沒有一個好像在耶路撒冷.
有聖殿那種屬靈的環境.
所以被擄到外邦的祭司.
他將那些律法.
濃縮成為613條.
容易記.
然後讓被擄的以色列人知道.
其實我們在外邦的地方.
不再像以前一樣.
在聖殿裡面那種屬靈的環境.
周圍的人有很多不信耶穌的.
或者不信神的.
我們可以怎樣走出來呢?.
當時還好的.
當時很純正.
就是這些祭司講解律法.
將它顯化成為一些可以實踐的信仰.
這些律法師,文士,法利,賽人.
就在那段時間出現.
一直來到耶穌的年代.
有三,四百年的日子.
他們記錄了很大一套的.
律法的文稿是怎樣實行出來的.
然後他就跟耶穌說.

$^{201}$他要試探耶穌.
「夫子,我該做甚麼才可以承受永生?」.
大家留意這個比喻.
在記載下不是一種很立體.
也不是一種很現場讓我們看到.
周圍有些甚麼.
其實如果我們能夠明白當時的處境.
就會更加立體一點.
為何律法師跟耶穌有這樣的對話.
其實以色列人.
他們十二歲過了一個成人禮.
他們要在信仰,宗教上高.
履行一個責任.
今天也有.
如果你有機會去到以色列.
耶路撒冷的哭場.
剛好遇到他們有成人禮的時候.
你先不要走.
你跟領隊說.
我可不可以自己回酒店.
你站在那裡看一看.
很隆重的.
全個家族出現.
十二歲.
由以色列救約的時候開始.
他們在成人禮之後.
就要找一個生命的師傅.
那些生命的師傅.
往往就是拉比.
律法師.
所以.
律法師看到耶穌問這個問題.
夫子我該穿甚麼才可以承受永生.
大家要留意.
這是拉比與拉比之間.
大家切磋的問題.
不是他突然間想到.
耶穌我要問你一個問題.
他知道答案.
他有一個很圓滿的答案.

$^{241}$不過在律法師與律法師.
大家見面的時候.
通常是這樣切磋.
因為不單止是.
律法師與耶穌在現場.
其實律法師的背後.
與耶穌的背後.
是有他們跟隨的門徒在.
這些門徒的再外圍.
就有一些圍觀的猶太人.
當中.
他們是有.
已經跟隨一些生命的師傅.
當中是沒有的.
所以通常他們見到.
律法師之間是互相切磋的時候.
他們一定是圍在一起的.
因為他們要看哪個律法師是最厲害的.
他們可能改變.
去跟隨那個師傅.
又或者那些沒有生命師傅的.
他們在聽這兩個對話切磋的時候.
他們自己選擇.
完了之後.
我可不可以跟你.
你成為我生命的師傅.
當時耶路撒冷有六千個這樣的生命師傅.
所以這個律法師.
他特意來面對耶穌的時候.
其實耶穌當時已經在地上高傳道.
傳福音.
傳天國有三年的時候.
所以不少人跟隨了耶穌.
他想著如果是撞倒耶穌的話.
真是不得了.
耶穌的門徒全部過當在他那裡.
所以他說.
夫子我該著什麼才可以承受永生呢?.
耶穌回答他二十七節.
你要盡心盡性盡力盡意愛主你的神.

$^{281}$又要愛倫舍如同自己.
這兩段經文.
前段是生命期.
愛倫舍如同自己是在利微期出現.
耶穌是將那六百十三條的律法更加濃縮.
成為了一條.
就是跟神有關係.
跟人有關係.
從你那裡反映出來.
那裡很厲害.
一次一個答案.
仍能夠壓制著這個律法師.
所以這個律法師.
他就要顯明自己有理.
如果被耶穌打壓了.
全班門徒都在他後面.
周圍的人還在看著.
如果我不能夠再回應或還擊的時候.
周圍的人怎麼看我?.
我的門徒會不會走呢?.
所以他問耶穌一個問題.
就是誰是我的倫舍?.
大家留意.
這裡所說的倫舍.
並不是住在你隔壁屋的人.
這裡所說的倫舍.
凡是和我有關係的人.
就是說.
耶穌.
誰是和我有關係的?.
他問這個問題.
是要對應耶穌剛剛先前在第九章做的事.
對不起,我喝了一杯水.
他要對應耶穌剛剛所發生的事情.
因為這個是能夠對耶穌最致命的打擊.
剛才發生了什麼事?.
在第九章.
因為路加福音第九章.
一直到十九章,十一章的聖經.
是記載耶穌最後幾個月.

$^{321}$由北面的加利利.
到南面的耶路撒冷.
因為離開的時候.
耶穌對那班門徒說.
我要上耶路撒冷去.
受長老,祭司,長民氏許多的苦.
並且避殺第三天復活.
這班門徒.
他們接收了這個訊息.
但是焦點並不是耶穌.
他要受死和復活.
他們所接收的訊息.
想到第二邊就是.
誰去承繼大亞哥的位置.
所以從第九章開始.
一直到十九章.
其實在這個過程裡.
他們一直在討論誰是大佬.
聖經說:爭論誰偉大.
所以一到耶路撒冷.
那個月節的時候.
進屋之前沒有人肯洗腳.
沒有人肯為其他人洗腳.
過往很多時候是可以給的.
但是可以給完全不動.
大家當沒事發生.
因為誰洗腳就誰認輸.
因為這幾個月大家都在討論.
沒有答案誰是大佬.
所以耶穌先正.
拿起一盆水,一條巾.
第一個走到彼得面前.
就是這個原因.
誰是我的淪舍.
因為在第九章.
當他們一開始下去耶路撒冷.
就經過撒瑪利亞.
在那裡.
不像三年前.
邀請耶穌進入撒瑪利亞.

$^{361}$去港島.
然後很多人信著.
三年面目全非.
當耶穌帶著那些門徒.
進入撒瑪利亞的時候.
沒有人肯接待他們.
那些門徒就很生氣.
三年前我們知道.
你們有這樣的反應.
但是今天你們竟然一反常態.
於是他們就跟耶穌說.
我們祈求天上的火.
將他們燒死.
將整個城都燒死.
當然耶穌就禁止他們.
耶穌這個表現.
是要表達出.
什麼叫做真的淪舍.
而不是帶著那種仇恨.
所以這個律法師.
他就是要捉住耶穌這個焦點.
誰是我的淪舍.
因為所有的猶太人.
他們對撒瑪利亞人那種憎恨.
已經到了極點.
因為他們在被擄之後.
是跟外邦人通婚.
無論他們主動還是被動.
因為亞述的政策.
是要將所控制的民族.
將他們通婚.
擾亂他們的獨有民族性.
他們不可以再造反叛亂.
所以留下來的猶太支派.
辯亞敏支派.
利美支派.
就對這些其他支派.
他們變成了混雜了的猶太民族.
他們很憎恨.
他們不持守這個信仰.

$^{401}$他們不能夠堅守著亞伯拉罕.
傳下來的神所應許的子民的血統.
所以他們很憎恨.
但是今天.
律法師知道耶穌這樣的反應.
他就盼望藉著這個問題.
耶穌再回答的時候.
就喚起當時那一班猶太人.
對耶穌的拒絕.
你為什麼要維護那班撒瑪利亞人.
燒死他們就燒死他們.
但是你阻止.
他們很希望能夠產生這樣的結果.
因為他們第一個問題已經輸給了耶穌.
所以聖經說.
他們要顯明自己有力.
當時我們就了解到.
耶穌就立刻說一個.
在現實上他們常常看到的現象出現.
耶穌的回答就是.
針對他們只有屬靈的知識和憎恨.
針對律法師的那種心態.
針對周圍的那一班圍觀的猶太人.
他們那種心態.
誰是我的靈性.
連那一班跟隨耶穌的門徒.
跟了三年都是動不動就燒死他們.
動不動就不理會他們.
跟他們沒有關係.
頂著門.
這個淪陷.
如果正確地表達.
是一個跟我有關係的人.
今天什麼人跟你有關係.
除了你的家人,鄰舍,同學,同事.
還有跟你接觸的人.
很多時候人與人之間的關係很微妙.
大家的性格不同.
大家有時在說話上.
是有意無意地得罪了.

$^{441}$言者無心聽者猶疑.
你又聽了.
頂著門.
淪陷與你有關係.
是不是像那一班門徒一樣.
燒死他們就行了.
免得看到他們.
耶穌並不是這樣.
所以耶穌說.
有一個人從耶路撒冷到耶利哥.
耶路撒冷在聖安山上高.
大概740米.
耶利哥在水平線以下.
在地理環境上高.
向下350多米.
從耶利哥到聖安山的聖殿.
距離有千多米.
一千多米有多高.
我嘗試給大家一個印象.
廣州很多時候香港人都去.
三百多米.
但是香港的太平山五百多米.
即是大概兩個太平山的高度.
由耶利哥一直上去.
然後距離有多遠.
大概22公里.
22公里是怎樣.
大概馬鞍山一呎.
但不是坐車.
是走路.
經過海底隧道走過去.
22公里.
平均一個人走5至6公里.
是一個小時.
所以大概要走5小時.
而且不是平路.
平路每小時是5至6公里.
所以不止一個路程.
讓大家立體看看.
當時耶穌所說的比喻.

$^{481}$這裡是耶路撒冷.
這裡是耶利哥.
不是平的.
是有距離的.
這件事發生在.
從耶路撒冷到耶利哥的一個場景.
當我們再看的時候.
我們就了解到.
當中有強度的出現.
剛才我說為何有強度.
是否作外?不是.
的確在耶穌的時代.
在這條路上有很多強度.
所以他們過節的時間.
不會是一個人.
整個家族一起走上去.
因為過節,守節.
到聖殿要帶牛,羊,甲子,班溝.
還要帶錢.
因為要到聖殿納電稅和汀稅.
強度是怎樣的呢?.
是因為希律王當時在耶利哥.
耶利哥整個範圍被譽為沙漠的綠洲.
因為周圍都是沙漠.
唯有那裡有水源.
有很清綠的植物.
就算你今天去耶利哥.
或去聖地.
都會帶你去那裡看看.
周圍沙漠.
在那裡.
所以希律王在那裡建造一個行宮的時候.
就需要很多人.
完成了.
他們沒有遣散費,沒有公積金,沒有強積金.
沒有任何做完了的花紅.
有些人當然回到自己的老家.
但是有一部分就流落在這條路上.
為什麼流落在這條路上?.
流落在這條山路上?.

$^{521}$因為他們知道.
猶太人起碼一年三次.
上去聖殿過節的時候.
一定帶著很多物資.
他們就去搶.
看準機會就去搶.
聖經形容到有一個人.
這個人要過這條路.
是有他的目的在的.
一會兒大家就知道.
聖經說.
這班強盜把他的衣服剝去.
大家留意.
當時這些衣服.
就是代表著一個民族的特徵.
就好像我們中國有五十六個少數民族.
每一個少數民族都有特別象徵著.
他們民族的衣服.
一看就知道是屬於什麼族人.
耶穌特別說出.
「誰是我的鄰舍?」.
是要表明這個人.
你不知道他是什麼背景.
你不知道他是什麼民族.
這才顯示出.
你會不會去做.
你不知道是什麼背景.
你會不會去看他的情況.
聖經說他打得半死.
打得半死是鋪排.
跟著說的祭司和利未人有關.
打得半死是昏迷在那裡.
這跟潔淨條例有關.
做祭司.
做利未人.
特別是你事奉.
在潔淨條例底下.
不是摸那個不潔淨.
或摸那個死屍.
才叫不潔淨.

$^{561}$在他們的條例底下.
五尺就是在我這個範圍.
去到陳牧師那裡.
我不能夠經過這些.
不能再埋的.
我只能在這個五尺範圍底下.
去看那個被打得半死.
昏迷的人是怎樣的.
所以他不知道這是什麼人.
他不知道他是生還是死.
因為如果他能夠走過去.
可以看到摸到他有沒有呼吸.
他的情況是怎樣.
他不能夠走過去.
因為再走過去.
半寸他都不潔淨.
所以因為這樣的限制底下.
聖經形容到這個祭司.
他從這條路下去.
他剛剛完成了聖殿的一個值班.
他是潔淨之軀.
他回到耶利哥祭司城.
可能他的家就在那裡.
你潔淨的.
你回到家裡不潔淨.
很多問題出現.
你是回到聖殿再走上去.
你再潔淨自己.
但問題是你先要處理這個人是怎樣.
所以很多這樣的原因.
他沒有做,他走了.
聖經形容到.
他從那裡走過去.
然後另外一個人出現.
就是利未人.
他都同樣好像一個祭司一樣.
背著一個歷史信仰的包袱.
他很掙扎.
我去幫這個人,去看這個人.
不要說幫不幫.

$^{601}$我去看一看他.
我就不潔淨了.
如果我看到這個人有需要.
我又不知道他是什麼背景.
如果他是薩馬利亞人,那怎麼辦.
我又要扔下他.
就算他有生命的危險.
我都不能夠理他的.
有很多這樣的限制.
但是信仰是怎樣的信仰.
特別我們基督徒.
我們要與人同行.
我們是怎樣.
耶穌要說出.
沒錯,我們有很多歷史的包袱.
甚至我們有很多信仰的傳統.
我在浸信會剛剛的事奉.
我教會有十幾二十個牧師.
有七十幾個執事.
我常常跟他們說.
我們浸信會,沒錯.
有很多傳統.
傳統裡面有很多很好的.
但是有一些我們要.
真的常常去反省,反思.
究竟這些條例.
是怎樣影響我們事奉的.
是怎樣影響我們對弟兄姊妹的.
究竟這些規矩.
我們能不能夠去了解到.
去建立到人的生命.
還是只不過給人看見.
這些規矩教會是這樣定的.
是這樣做的.
弟兄姊妹.
很多時候我們就有這些信仰的包袱.
傳統的包袱.
歷史的包袱.
這個利未人一樣從那裡下來.
他也同樣是事奉完了.

$^{641}$然後回家.
聖經說他也是這樣.
從那條路經過.
活培的信仰不在乎.
你是甚麼角色.
而是處理那個歷史包袱.
是怎樣處理的.
也不是單單在人面前所做.
雖然利未人在聖殿裡的角色.
不像祭司那麼明顯.
他只不過是一個輔助的角色.
但不因為他這個輔助的角色.
他就覺得我只是小人物.
小人物.
不需要做了.
算了.
耶穌著意用利未人在比喻上.
他就要點明無論你那個角色是怎樣.
這個不影響你與別人同行.
影響你與別人同行的.
就在於你那個生命的流露和表達.
你那個真誠的信仰是怎樣.
第三.
一個角色的出現.
撒瑪利亞人.
當耶穌說撒瑪利亞人的時候.
其實周圍圍觀的那班已經打了個特.
我們都不用這幾個字眼.
說出來的.
耶穌你夠膽這樣去說.
所以圍觀的人.
他們內心都已經對耶穌打了個問號.
當然.
對耶穌在剛才第九章帶著那班門徒.
特意走經過撒瑪利亞城市.
他們已經感覺到有少許的遺憾.
但耶穌不怕說.
是因為這個撒瑪利亞人.
他所做的事.
撒瑪利亞我們可以了解到.

$^{681}$他們都有五經.
他們都有創出利民生.
當然他們的五經叫撒瑪利亞五經.
是和猶太人.
即是和我們現在聖經.
創出利民生是一模一樣的.
只不過.
他們的節日是比猶太人晚一個月.
舉例說,逾越節.
在猶太的曆法上.
猶太人是正月.
撒瑪利亞人是二月.
他們照樣有這樣做.
所以這班撒瑪利亞人.
他們都是一直在走律法.
他們當然.
因為.
以色列人的緣故.
對他們的態度.
他們的反射反擊.
同樣是對猶太人沒有好感.
即是兩個互相對立.
但這兩個互相對立.
很明顯在比喻上表達出來.
你看,祭司和利尾人.
原本他們是在整個信仰上做帶領.
誰知他們是這樣的帶領.
撒瑪利亞人和猶太人是互相對立的世仇.
但他們面對一個.
他都不知道是什麼背景.
被打得半死的人.
在路邊,他沒有計較這些.
生命最重要.
他見到那個人的需要是最重要.
他就先走了.
所以他走過去.
我特別提到.
他們有行律法.
他們自己本身都有潔淨的條例.
他們自己本身都不能夠摸死屍.

$^{721}$這個撒瑪利亞人.
他動了慈心.
這個動了慈心,我們可以很清楚了解.
他不只是一個信念.
信仰.
特別在近幾個月.
我常常分享,鼓勵.
弟兄姊妹和我家人.
信仰不是一個.
我們的信仰不是一個道德的信仰.
做好人多好.
所以當你傳福音的時候.
很多人回應.
哪個教一樣.
都是導人向善.
是道德的.
我們的信仰並不是道德的信仰.
是一個關係的信仰.
你看每一個信仰.
在我們今天所了解的.
不會和你講神和人之間的關係.
有,只是一個利害的關係.
你求我嗎?.
你對我有所求嗎?.
你燒一隻豬給我.
或者在人的角度.
我燒一隻豬給你.
盼望你有相應的回應給我.
但是我們的信仰.
是講關係的信仰.
是耶穌基督.
他甘心主動.
道成肉身在地上高.
去體驗我們那個需要.
為我們付出.
建立和我們的關係.
道成肉身很重要.
在照顧阿Mo的時候.
我更加體驗到.
很多時候我們是從自己的角度去.

$^{761}$為什麼呢?.
因為我們沒有站在對方的角度.
譬如他一個例子.
他說:你可不可以把我移到左邊一點?.
我們真的移到左邊.
但是對他來說是右邊.
你這樣對著他.
移到左邊.
我們很自然就是.
我們的左邊就是這裡.
其實他的左邊就是那邊.
很多時候我們形態的觀念就是這樣.
就是站在我們自己出發.
所以你照顧開病人.
或者有需要的人.
如果你站在他的角度去看.
你就會懂得好好照顧他.
耶穌基督道成肉身就是這樣.
他不是高在上.
他來到地上站在我們的角度.
來看我們的需要.
頂曉.
這個撒米拉人.
他動了慈心.
他站在被打得半死的人的角度.
他的信仰.
化成了對對方的認同,體會.
所以他拿出油和酒.
倒在他的傷處.
油和酒.
你不要想到今天.
你廚房用的花生油,橄欖油.
你倒這樣的油在我身上.
很絲粒的.
不是.
這些油是橄欖油.
不是我們在超級市場.
只有兩種橄欖油.
一種初榨,一種第二浸.
他們有幾種不同的製法.

$^{801}$有不同程度的橄欖油.
這些通常在遠行的時候.
用比較稀釋的油.
用來清洗.
因為當時沒有水源.
又不像你那麼方便.
你去買一瓶水.
不是.
他用來潔淨,清潔.
甚至是用一點點來頂住.
酒是用來消毒.
所以薩馬利亞人把他的傷處洗淨.
把酒消毒.
然後包裹.
這個包裹讓我們更加了解.
他不是在附近住的.
因為當遠行的時候.
通常他們帶油和酒.
還有一條長的帶.
是用來救命的.
救命是甚麼呢.
就是他們遇見一些強盜或跌傷.
馬上用來包紮.
是遠行的人才會有.
所以從這些裝備.
我們可以了解到.
因為耶穌是用一個很現實的處境.
讓那些圍觀的人知道.
他不是住得很近的.
他是一個遠行的人.
而且他是有目的地.
有目的地在這條路出現.
特別是把傷者扶上自己的身後.
這個身後是一個雷仔.
在當時主要是運送貨物.
不是載人.
即是他只是做生意載貨.
但現在他暫且放下自己的生意.
因為做生意時機很重要.
不過他放下.

$^{841}$他見到這個人的需要.
他先處理幫這個人.
然後把他放在雷仔身上.
然後把他帶到客店.
放下二錢銀紙.
如果你是店主.
你放下錢就可以.
你扔下傷者.
你何時回來.
其實在當時我們了解到.
一錢的銀紙.
等於一個星期的住宿.
即是我扔下兩隻銀紙.
已經是隱喻到.
我兩個星期後會回來.
其他在這兩個星期內.
如果你要再照顧而有的花費.
我回來再跟你算帳.
還給你.
然後耶穌問那一個律法師.
這三個人.
哪一個落在強盜手裡的淪陷.
那一個律法師.
他回答得很有技巧.
他只說是憐憫他的人.
他不用撒瑪利亞人.
耶穌說的那三個比喻的主角.
是祭司利未人和撒瑪利亞人.
不是祭司利未人和憐憫人.
為什麼?因為他知道.
一說撒瑪利亞的時候.
周圍又會起哄.
那些門徒可能會立即飛起.
所以他回應耶穌就是.
是憐憫他的.
對他來說是一個突破.
因為他是接受撒瑪利亞人所做的事.
而不是否定他所做的事.
但是這只是開始.
不完全.

$^{881}$但是耶穌說你照樣行拜.
雖然不完全.
但是耶穌接納他,鼓勵他.
你繼續去做.
你繼續這樣去想.
你繼續這樣去教導就可以了.
頂止不去.
同行者是需要突破很多人際的隔膜.
是需要付出的.
當我們看回耶穌說這個比喻的時候.
第一方面我們就要去思想反思.
究竟今天我們說信耶穌.
這是一個什麼信仰?.
不要想著是一個傳統的信仰.
不要想著是一個傳遞經書的信仰.
不要想著是一個道德的信仰.
不要想著只是我有難的時候.
可以捉住他.
林吉抱佛腳的信仰不是.
你為什麼信耶穌?.
在這些代祖信一出了兩三個星期.
我收到不少回應.
牧師.
我不是基督徒.
但我一收到你的代祖信.
我就為你們祈禱.
為你們祈禱.
我從來沒聽過.
因為試用了四十年.
信了耶穌超過半個世紀.
我從來沒聽過一個非基督徒.
為一個基督徒來祈禱.
頂止不去.
有一些.
牧師.
我因為教會.
給我這個信仰上.
有一些陰影.
我已經十幾年沒有回教會.
但因為你們的經歷.

$^{921}$我回到教會.
我慢慢開始有時服.
我慢慢開始投入一個.
什麼叫信仰.
頂止不去.
你信耶穌為什麼?.
沒錯.
有很多歷史的包袱.
但是.
我怎樣去看呢?.
我盼望.
我們不是像這個律法師.
有這麼大的歷史包袱.
他看不見那種生命的流露.
一個活潑的信仰同行者.
是需要付出.
行動.
甘心樂意.
不是看你是哪個人.
不是看你是否撒瑪利亞人.
不是看你是不是我所憎惡的人.
耶穌著意.
去用這個比喻.
盼望他能夠解開.
律法師他的心結.
又盼望解開周圍圍觀的人.
他那種思想.
我盼望.
今天我們都在神的面前.
藉著聖靈在我們心裡面的工作.
能夠活在我們現實的一個環境.
誰是我的鄰舍.
我的同行者是怎樣.
請我們一起祈禱.
因為如果沒有聖經.
我們只不過仍然活在自己.
靠著我們的經驗.
靠著我們所見到的環境.
來判斷我們的生活模式是怎樣.
因為耶穌你道成肉身.

$^{961}$站在我們的角度.
並且有神你自己的話語.
有真理的聖靈提醒我們.
與神你繼續的工作.
我們的禱告.
是奉基督耶穌名字.
阿們.
\newpage



\section{馬太福音 5:33-37}
\label{sec:O_RnYw4oPiY}
\textbf{2023年3月5日崇拜直播｜陳明泉牧師｜給當代門徒的15個忠告:成為誠信可靠的人｜馬太福音5\_33-37}
\newline
\newline
連結: \href{https://youtube.com/watch?v=O-RnYw4oPiY}{\texttt{https://youtube.com/watch?v=O-RnYw4oPiY}} ~~~~ 語音日期: 2023-03-12
\newline
\newline
\hyperref[sec:VsBsdArDKAI]{\small{< < < PREV SERMON < < <}}
~
\hyperref[sec:index_chronic]{\small{[返順時目]}}
~
\hyperref[sec:DeJAKAGjSug]{\small{> > > NEXT SERMON > > >}}
\newline
\newline
馬太福音 5:33-37
\newline
\begin{longtable}{cl}
\hline
\hline
章節 & 經文 (和合本修訂版)\\
\hline
5:33 & \begin{tabularx}{0.7\textwidth}{X} 「你們又聽過有對古人說：『不可背誓，所起的誓總要向主謹守。』 \end{tabularx} \\ \\ \relax
5:34 & \begin{tabularx}{0.7\textwidth}{X} 但是我告訴你們：甚麼誓都不可起。不可指著天起誓，因為天是神的寶座。 \end{tabularx} \\ \\ \relax
5:35 & \begin{tabularx}{0.7\textwidth}{X} 不可指著地起誓，因為地是他的腳凳；也不可指著耶路撒冷起誓，因為耶路撒冷是大君王的京城。 \end{tabularx} \\ \\ \relax
5:36 & \begin{tabularx}{0.7\textwidth}{X} 又不可指著你的頭起誓，因為你不能使一根頭髮變黑變白。 \end{tabularx} \\ \\ \relax
5:37 & \begin{tabularx}{0.7\textwidth}{X} 你們的話，是，就說是；不是，就說不是。若再多說，就是出於那惡者。」 \end{tabularx} \\ \\
[1ex]
\hline
\hline
\end{longtable}
$^{1}$接下來我們繼續聆聽聖道.
我們繼續我們的敬拜.
今天我們選讀的經文是新約聖經《瑪提福音》.
瑪提福音第五章三十三到三十七節.
請聽我讀出.
你們又聽見有吩咐古人的話說.
不可背誓.
你起的誓總要向主謹守.
只是我告訴你們.
什麼誓都不可起.
不可指著天起誓.
因為天是神的座位.
不可指著地起誓.
因為地是他的腳墩.
也不可指著耶路撒冷起誓.
因為耶路撒冷是大君的京城.
也不可指著你的頭起誓.
因為你不能使一根頭髮變黑變白了.
你們的話.
是就說是.
不是就說不是.
若再多說就是出於那惡者.
弟兄姊妹平安.
戴口罩好看一點.
但是能夠可以.
好像不用戴著面罩.
可以比較見到整全人的面貌.
其實對我們來說.
好像這幾年經歷過疫情之後.
現在好像變得很奢侈.
但也很珍惜這些時間.
可以面部坦蕩蕩去做人.
不過人生有時候未必可以完全坦蕩蕩.
我記得我小時候讀書不太好.
有時候默書通常都不高分.
但很容易每次都過關.
為什麼呢.
因為第一次默書一定不錯.
媽媽就會簽名.
我以前用印水紙印了媽媽的簽名.

$^{41}$之後不合格的.
我趁媽媽不留意.
用紙套著來冒簽媽媽的名字.
不合格的就我冒簽.
合格的拿100分的就給媽媽簽.
有一次是這樣的.
那次有100分就給媽媽簽.
我用萬字夾給媽媽簽.
媽媽簽完就拿萬字夾去看前面.
我簽過這次嗎.
那次零分為什麼我簽了.
我說媽媽你簽了你不記得.
不過媽媽一直揭的時候.
每次這麼多不合格.
零分的一定沒有簽過.
然後就氣我一個星期.
媽媽氣我以為我做什麼呢.
我以為媽媽氣我不合格.
不勤奮去做.
氣了一個星期都不知道發生什麼事.
板著臉直至到一個星期後.
媽媽就跟我說.
誰知道我氣你什麼.
不是氣你拿零分.
是氣你不坦白.
氣你面對自己的錯誤.
不肯真誠承認.
還要沒有簽我的簽名.
簽得這麼噁心.
類似這樣.
你過到別人的關.
但過不到自己的關.
這件事小朋友可能很多在座的.
有些弟兄姊妹都會經歷過.
不過原來人生裡面.
當我們靜下來去想.
原來在我們人生裡面.
很多時候我們都經歷過一些不坦白.
又或者面對人生命黑暗的地方的時候.
我們會用一些東西來遮掩自己.

$^{81}$今天我們一起讀的這段經文.
其實不是一個很複雜的經文.
但希望我們從這段經文.
從耶穌基督對一班群眾.
一班弟兄姊妹對他的門徒所說的這番說話.
成為我們警惕.
亦都希望比我們今天.
我們更加著意自己生命裡面.
是否成為一個坦誠有誠信的一個人.
我們求主幫助我們.
我們先祈禱.
求聖靈幫助我們.
我們同心祈禱.
讓我們今天一起來聆聽上帝的話語.
願你的靈願我們每一個弟兄姊妹同在.
讓我們能夠聽得明白.
能夠將你的話語藏在心.
在生活當中我們應用實踐出來.
驅除我們生命的黑暗.
你的話語來到的時候.
讓我們生命的黑暗被驅走.
幫助宣講的講得清楚.
讓你的話語忠心地被講解出來.
讓我們繼續按著你的話語.
繼續行事為人.
恭敬將來的時間交託.
願你賜福.
天上禱告奉靠基督耶穌得勝的聖名.
Amen.
當我們去看這個登山寶訓.
其實都是講連續這麼多次.
十五次就在講耶穌對門徒的忠告.
在這一堆忠告裡面.
有一個很重要的價值觀念.
或者在耶穌對一群群眾講論的時候.
有一個很重要的核心理念.
在五章二十節.
耶穌說:我告訴你們.
你們的義若不勝於民事和法理賽人的義.
斷不能進入天國.

$^{121}$在耶穌講論這一番登山寶訓的時候.
某程度上是在講天國的倫理觀念.
天國的子民的生命的呈現.
不過耶穌特別的告訴那些聽的人.
如果你們的義不能勝過民事和法理賽人的義.
你沒有辦法進入天國.
某程度上來說.
這也是一個很高的要求.
你再細心想想.
為什麼呢?.
因為民事和法理賽人當時來說是道德的專家.
他們解釋律法.
即使猶太人失去了自己國家的主權.
在一個羅馬帝國之下.
但是他們維繫著整個文化.
靠著律法.
透過民事和法理賽人來做教導.
他們是道德的典範.
耶穌說你們的義要勝過他們.
要比一群專家厲害.
是不是真的很困難呢?.
所以這群民事和法理賽人.
對於當代的人來說.
根本上某程度上已經是代言人.
KOL.
耶穌說叫他們要勝過他們.
某程度上這是一個更加難擔的擔子.
另一句話來說.
耶穌又說過那些民事和法理賽人.
將一些難擔的擔子放在群眾身上.
令到那些人就像一個鴉一樣.
扼殺他們.
他們的生命是不自由的.
他們就像死守規條一樣.
來活出信仰的真諦.
但是是不是真諦呢?.
真是不知道的.
耶穌說這些是一些鴉.
令到他們的生命不自由.
耶穌說這番話好像有一點前後矛盾.

$^{161}$不過那個勝過的意思.
不是耶穌要將他所說的講論.
和這些法理賽人和民事做一些較勁.
大家比對一下誰比較勁.
他要將真正的生命之道.
將真正的天國倫理.
告訴那些聽的子民身上.
耶穌基督他希望這班人聽的時候.
他們不只是像以前死守律法.
按著字面.
或者不同的規條.
來守著得到永生.
或者得到天國的盼望.
而是真正得到上帝的真理.
所以在這裡.
五章二十節.
耶穌說要勝過.
那個勝過除了說要超越他之外.
更加不要用他們的思維.
來經歷信仰.
用這個角度來看五章三十三至三十四節.
你們又聽見有吩咐古人的話說.
不可背誓.
所起的誓總要向主謹守.
只是我告訴你們.
什麼誓都不可起.
這個是五章三十三至三十四上半節.
當我們讀到這裡的時候.
耶穌好像加多一條新的規條.
對於當代的子民百姓一樣.
不可什麼好像十戒的語法.
去到這裡.
在他們讀的律法.
或者古人的教導是怎樣教呢.
不要背誓.
違背誓言.
所起的誓要向主謹守.
其實是一個很正路的教導.
不要背誓.
言而有信.

$^{201}$你向主所答應的事就要去謹守.
不過在這裡一講不可起誓.
當耶穌說要勝過法理錯人的義的時候.
我們讀上去就會覺得有少少納悶.
因為在我們生活當中.
其實我們不多不少都會做起誓這個動作.
或者做這個表達出來的宣誓.
最簡單的結婚.
是誓言.
還有就是如果你萬萬物業.
要轉名.
甚至有些產業的轉換.
有時候都要經歷過這個在政府裡面做宣誓.
甚至你要面對法庭的訴訟.
按著聖經.
這個起誓.
其實在我們生命裡面是成為我們的一部分.
再看聖經.
其他的篇幅.
上帝自己都起誓.
耶和華起誓.
他直接來講.
他指著自己的起誓來將福氣賜給亞伯拉罕.
在聖經裡面.
上帝通常都是指著自己的起誓.
來向那些以色列人講.
上帝是守約斯前外的神.
上帝自己都起誓.
甚至耶穌基督的設立成為大祭司.
在希伯來書裡面.
特別講到.
他是上帝命定的大祭司.
耶和華上帝起誓而立的.
相對於當代的祭司按照制度制籤出來的.
耶穌的身份是超然的.
所以上帝自己起誓.
耶穌也是起誓而立.
你說上帝可以做.
我們人不能做.
但是到了新約再後一點.

$^{241}$保羅自己.
保羅自己講道.
或者對人講論.
他都起誓.
羅馬書一章九節那裡.
他說我在他兒子的福音上.
用心靈所事奉的神可以見證.
我怎樣不住地提到你們.
那裡向著神.
用心靈所事奉的神可以見證.
這個字不只是純粹見證這麼簡單.
而是起誓.
那個字眼就是這個意思.
掃尊於上帝所講的東西是真的.
另外一個就是九章一節羅馬書.
很容易記.
一章九和九章一.
都是保羅起誓的.
我在基督耶穌裡說真話並不謊言.
有我良心被聖靈感動.
給我作見證.
給我作見證這個字就是起誓.
你看到其實在聖經裡頭一尾都會有起誓.
我們生活裡也有起誓.
如果我們面對在耶穌的講論當中.
我們就全部違背了聖經.
我們是不是做錯了事呢.
有一些很極端的說法是真的.
甚至乎買賣也不應該這樣做.
婚禮也不應該起誓.
有一些很極端的教派真的這樣說.
不過我們看回聖經的述文.
那個不可起誓.
其實不是在說字面裡.
不准做這個禁戒.
或者做一條新的誡命.
耶穌要頒佈給子民當中.
不可起誓其實是指向一個.
更加內裡的人與人之間的關係.
背後裡面不要動諜地用誓言.

$^{281}$來證明你所說的話是真的.
那個誓言背後其實訴諸一個.
最重要你生命的特質.
你是不是成為一個誠信和可信任的一個人.
你成為一個可信的人.
你不用做每一樣東西.
或者做任何的表達.
都要用誓言來增加你自己的可信性.
反而是你的言行舉止一落千金.
你的生命成為別人裡面很信任的一個人.
相反來說如果你是一個沒有信用的人.
你起了多少誓也好.
你都沒有辦法來增加你自己的信用.
有一本書我很喜歡看的.
不是一些好心愛的.
是一些管理書.
作者叫做Steven Covey.
可能有些頂展會讀過.
有一本書叫做《高校信任力》.
這本書說什麼呢?.
這本書不是說教會的.
不過是說在職場裡面.
如果在職場人與人之間裡面.
建立了一個真正的信任.
建立了一個人與人信任的制度.
這個信任會成為一種力量.
可以令到在運作裡面.
減低了很多那些孜孜孜孜.
或者不信任帶來的額外的成本.
而且當你建立了信任之後.
你的工作關係.
你的團隊成為一個快樂的團隊.
相反來說如果你沒有了這種信任.
人與人之間就有很多猜忌.
你做事不安心.
經常都不知道我做這些事上司喜不喜歡.
我的同事會不會這樣看我.
或者我跟外面的人合作做生意的時候.
那個供應商信不信得過.
他會不會騙我等等.

$^{321}$這本書其實就是說.
真正的信任是帶來更加好的果後.
而且會帶來更加快樂的一個環境.
不過這本書不是停在這裡.
這本書最重要的不是說信任的好處.
而是說真正能夠建立你的信任.
你真的做些什麼呢.
它說幾樣東西.
除了你說的東西能夠兌現之外.
更加重要的是你平時做事的實力.
你做事是怎樣做的.
當有些處境有任務交給你的時候.
你是不是真的完完全全.
百分之一百地做好這件事.
這就成為了一個很簡單的倫理.
不過卻是今天竟然商業社會裡面最缺乏的一件事.
斯蒂芬·科斐所說的.
其實跟耶穌基督所說的是一脈相承的.
甚至乎其實他是一個基督徒的管理學家.
他用基督教的價值觀念來放在職場的運作當中.
他特別跟耶穌基督所說的互相呼應的就是.
今天我們要建立一個彼此信任的一種關係.
而這種信任不是在乎你說的話說得多動聽.
又或者你用一些誓言來支持你所說的話.
反而是用你的實力.
用你眼前完成任務的能力.
來證明你是一個可信的人.
回到經文那裡.
耶穌基督特別說到不可起誓.
或者那些古人吩咐人們不要違背誓言.
那個思路本身不是錯的.
不過如果人與人之間的關係只是建基了.
純粹是在你起誓的那些說話當中.
誓言以外的那些話語.
不要違背上帝.
但是有一些情況你未必所有東西都會訴諸於上帝的.
你平時人與人之間的交往.
如果你沒有一個誓言來做一個支持的時候.
人與人之間會不會成為一個不可信.
或者大家都沒有辦法有誠信建立的一個世代.

$^{361}$所以耶穌基督特別指出.
起誓根本就沒有辦法成為一個人建立誠信的一個手段.
反而是動諜地起誓.
會製造很多人與人之間的不信任.
好幾年前我聽過一個戰亂的國家.
因為那裡長期戰亂.
所以宣教士說那些人不會說真話.
那些人在戰亂國家生活的那些人.
第一眼見到你很友善.
他永遠都是笑笑口請你喝茶.
第一句你問他今天吃飯了沒有.
他永遠都是答吃了.
其實不是真正的答案.
他說這個民族因為經歷長期的戰亂.
你要問他五次同樣的問題.
第五次才說真話給你聽.
所以他說他要傳福音.
你問他信不信耶穌呢.
他會信的.
但根本心不是信的.
又或者你對信仰對基督徒有不友善.
他第一句到第四句都說友善.
到第五句才說真話.
不過你可不可以熬到第五句.
或者第一至第四句已經信了他.
你就很多時候會面對生活的遭殃.
他說他曾經分享過.
有一些被出賣的經歷.
都是輕易地相信頭那四句.
所以到最後他是被人騙了.
不過那個簡單沒有信任.
令到人與人之間成為互相猜忌.
這個國家可能是戰亂.
令到人變得不真實不誠實.
但亦可能因為這些人的不誠實.
加劇了這個國家彼此的誤會.
有很多的戰爭在當中.
好像惡性循環一樣.
耶穌基督他要針對的就是.
我們今天是否成為一個誠信的人.

$^{401}$起誓沒有辦法幫助我們成為一個可信的人.
究竟我們今天的生命.
頂子們你覺得你的生命可不可信呢.
你覺得你自己是否在別人的眼中.
別人會信任你呢.
在我預備這段經文的時候.
我自己也在問自己這個問題.
雖然我是教會的主任牧師.
我經常都問自己.
究竟我是否一個可信的人呢.
我講的道大家又信多少呢.
又或者我平時從別人呈現出來的樣貌.
我是一個怎樣的人呢.
我身邊的人信不信任我呢.
我回答自己.
我有一段時間有過這樣的經歷.
話說在我事奉早期的時候.
我不是特意去說謊的.
不過我發覺很多時候我答應了家人.
我會出席家庭的家聚.
又或者答應了家人會做某些事.
答應了老婆做這些事.
到最後不是忘記了.
而是說忘記了.
到最後是不了了之的.
我記得有一次答應了子女.
那晚一定要陪他們.
跟他們慶祝有一個特別聚會.
我忘記了是什麼聚會.
不過我答應了他們.
我打算回去.
但當我答應了之後.
我發現原來我答應了別處.
我要帶他們去教會的一個聚會.
到最後我打算很快搞定就馬上回家.
但最後是搞到晚上十一點的十二點.
子女基本上都睡了.
他們其實預備了所有的食物.
打算我回去.
我太太也睡了.

$^{441}$我自己就在桌面上吃蛋糕.
吃他們預備的東西.
我自己在問自己在做什麼.
我為什麼會成為這樣的人.
我為什麼要令他們失望.
我明明早就答應了他們會這樣.
但卻突然之間有些事情出現的時候.
我就將他們的重要性放到最低.
我打算我自己的能力可以做到.
但到最後我失去了跟他們一起約會的時間.
我心裡很不安樂.
隔了幾天之後.
我說我補償一下.
我做回這些事情.
跟子女和老婆說.
但他們說先看看怎樣.
爸爸你很忙的.
這句話大家明白吧.
他們不是潤我的.
不過失去了信任.
我未必是存心要騙他們.
不過在我心裡我知道.
我沒辦法兌現我口裡所說的話.
即使我再繼續說多少句也好.
我都要再用時間去重建這種信任.
在我們人生裡.
你是否一個可信的人呢.
不是純粹問陌生人可不可信.
問問你身邊和你共處的人.
你的家人最清楚你是否一個可信任的人.
從你的生活裡呈現你是否一個可信的人.
我相信這已經成為我們生命裡.
進入耶穌這番說話裡.
一個很重要的特徵.
我們繼續看經文.
34-36節.
不可指著天起誓.
因為天是神的座位.
不可指著地起誓.
因為地是他的腳凳.

$^{481}$也不可指著耶路撒冷起誓.
因為耶路撒冷是大君的京城.
也不可指著你的頭起誓.
因為你不能使一根頭髮變黑變白了.
耶穌基督在這裡再繼續說.
剛才說的不可起誓.
訴諸於人要建立誠信.
不過當時的猶太人或怎樣教他們.
他們指著一些比自己更大的來起誓.
所以耶穌要顛覆他們.
你不要指著這些.
以為有人來做你的抵押.
或有一些物件成為你的後台.
以至你的說話更加可信.
耶穌在這裡完全用否定的方式.
不要指著那些.
他告訴你指著那些是沒有用的.
也沒有必要用這些方法增加自己說話的實力.
更加重要是因為當時的古人.
是這樣教不可背誓指著神的就要對言.
有些人為了增加自己可信性.
但又不想指著神.
既有忌諱.
但同時又怕這種忌諱.
萬一對言不成諾帶來一些後果.
所以就有一種技術性的起誓.
這個聽起來很荒謬.
有些色經學家看猶太早期文獻就這樣教.
技術性起誓就是指著耶路撒冷起誓.
萬一承諾不對言就不需要擴散後果.
但如果面向耶路撒冷起誓.
就要承擔後果.
是不是很荒謬.
同一句說話.
指著他就不需要承擔後果.
面向耶路撒冷就要承擔後果.
所以當時每一個人做這些起誓.
既然是這樣就指著耶路撒冷.
表面上牌面上他很可信.
但實際上萬一他爽約.

$^{521}$或不能對言承諾的時候.
他就不需要承擔後果.
借助一個比他更大更有權威的東西來幫助他.
但實際上失敗的時候他就逃之夭夭.
很多時候當時的人就用這種技術性起誓.
來幫助自己令自己可信.
耶穌基督在馬太福音第23章裡面.
他再繼續說.
他說有些人這樣教起誓.
建立自己的誠信是有禍的.
馬太福音23章16-22節.
那段經文是很長的.
是在說法利賽人的禍患.
他說法利賽人.
你們那些核眼領路的有禍了.
你們說凡指著電起誓的.
這算不得什麼.
只是凡指著電中金紙起誓的.
他就該謹守.
後面說的是輕與重.
指著大的那些因為很空泛的.
所以就算不得什麼.
沒什麼的.
只要那些金紙.
底壓的那種狀態之下.
你就要謹守了.
入獄的.
有結果的.
耶穌說這一班人是一派胡言.
用這些方法.
用這些的起誓.
或者哪個大一點哪個小一點.
哪個重要哪個不重要.
既是顛倒了重要和重輕之分.
基本上顛倒了.
更加不理想的地方是.
以為用起誓就可以建立了.
人與人之間的信任.
耶穌說這一班人是一派胡言.
製造更加多的荒謬.

$^{561}$在這個世界裡面.
不知道你相不相信.
其實我們某程度都是活在.
這些虛荒當中.
有時候為了證明自己有道理.
或者我們證明自己所說的話.
我們有時候未必是起誓的.
不過我們會訴諸於某一些的說法.
片面的去拿某一些.
我們覺得適合我們用的一些權威.
然後套入放入在我們說話當中.
或者拿人家小小的地方.
截出來.
然後不是看上文下理不是整個意思.
拿那個部份.
將部份的真理引用.
然後就成為了我自己說話的內容.
這些是我們生命裡面很常見的.
雖然我們未至於到好像起誓那種方式.
不過我們的心態大致是要證明自己有道理.
我們找一些更大的來證明自己.
或者用人家片面的方法.
來幫自己令自己可以可信.
最常見的一些例子.
你當笑話聽.
頂指會有些茶經.
我不是介意引用我說的話.
不過有時候明明我不是說那些.
他就會說牧師說的.
陳牧師說的.
還要錄下音的那句.
有時候就說下.
但我的意思也不是這樣的.
有時候會這樣.
什麼人都會用的.
有時候會說珍珍牧師這樣說.
Raymond這樣說.
梁院長這樣說.
找個院長再大一點.
令到自己充滿了有backup.

$^{601}$有時候除了人之外.
我們甚至會說聖經這樣說.
我們片面地理解聖經.
不是整全地理解上帝的話語.
我們就抽取幾個字.
合用的.
我們就拿來放在我們要說的話當中.
這樣做是否增加了我的說服力呢.
對於一些不心夠的人可能是.
不過久而久之.
當我們成為了一個.
只是以別人或比自己更大.
來做一些後盾.
但沒有真正自己去查究一個真相的時候.
去說的這些說話.
其實到最後.
我們都是成為一個面目模糊的人.
我們都沒有辦法有自己的話.
我們也不是在說我們自己的心底話.
求主幫助我們.
在這裡耶穌基督.
他要拒絕法利賽人帶來的荒謬.
你不可以指著天起誓.
天是上帝的座位.
你不可以指著地起誓.
地是上帝的腳墩.
你不要指著耶路撒冷起誓.
耶路撒冷是上帝作為大君王的經承.
不可以指著你的頭起誓.
因為你不能使一斤的頭髮變黑變白了.
全部耶穌所引用的所說的那個不可沒有.
因為他的背後就是在說你所說的每句說話.
你不要想著指著那裡就不用負責任.
到最後指著什麼的對象.
訴諸於都是上帝你口裡所說的話.
即使你指著自己的頭.
指著自己的頭就行了.
你的頭髮也是上帝數算過的.
無論你禿頭也好.
或者你的頭髮很濃密也好.

$^{641}$上帝有最終的訴諸權.
所以我們所說的每句說話.
指著的每一個對象.
用什麼來背負自己.
我們都離不開上帝的主權底下.
既然是這樣.
耶穌基督其實推到極致.
你可以用這個角度來理解.
你當是你口裡所說的每句說話.
都好像誓言一樣.
每句說話你都有結果的.
每句說話都帶著最終上帝.
將來用這句說話來看你成為一個怎樣的人.
在你人際關係裡面.
你的人生命的誠信.
都是憑著你生活的每一句說話.
來表達出來.
不論你指著什麼對象.
全部都和上帝有關.
所以耶穌基督在這裡.
要特別將當時的古人.
花里草人所說的話顛覆了.
耶穌的應用是什麼呢.
很簡單的.
第37節.
你們的話.
是就說是.
不是就說不是.
再若多說.
就是出於那惡者.
其實理解不複雜的.
是就說是.
不是就說不是.
你按著你所知道的.
或者你所經歷的實況.
是就這樣說的.
不是就不是.
簡單來說就是.
按著你自己的經歷.
按著你眼前的.

$^{681}$來將真相表達出來.
你的溝通就是這樣表達出來.
不過後面再加多一句.
若再多說就出於那惡者.
這裡就很奇怪.
為什麼會再多說就出於惡者呢.
其實你用一句語意分析就不難了.
簡單來說.
一句說話.
開頭的時候就會說了你自己的立場.
你平時溝通也是.
譬如我沒有偷東西.
我沒有偷東西.
第二次再說.
我沒有偷東西.
人家好像有點疑惑.
第三句就說.
如果你不相信.
我就指著什麼是起誓.
如果覺得人家好像有些質疑你.
有些懷疑你的時候.
再多說.
那你為了要證明自己.
是真的.
所以你就會用更加多的力量.
用其他起誓的方式.
來證明自己有道理.
將這句說話演繹的.
最清楚的是什麼呢.
彼得不認耶穌的段落.
當耶穌面對審判受難的時候.
門徒四散.
彼得就離開了.
有些人看到彼得在這裡.
在這裡團結一些火.
在這裡取暖.
有些人認出.
你不是跟隨耶穌的門徒嗎.
彼得第一次說是什麼.
我不是.

$^{721}$第二次就說.
我不認識那個人.
第三次就說.
經文裡面說.
彼得就發奏起誓說.
我不認識這個人.
看到了嗎.
其實整個語意分析.
在彼得裡面就表達出.
第三句說話.
說的錯誤的信息.
為了要證明自己有道理.
就會引入去.
發奏起誓.
來告訴別人.
證明他自己是對的.
或者證明自己是無辜的.
不過我們讀聖經.
我們知道.
其實他是耶穌基督的門徒.
他是彼得.
他是說謊.
他是不認主.
在我們人生裡面.
有時候可能.
我們為了證明自己有理.
或者我們想告訴別人.
我是無辜.
我們會訴諸於.
或者我們會再繼續多說.
用更加多的話.
來背後自己.
以致我成為一個可信.
或者幫自己脫身.
在這些困難當中.
耶穌基督說.
再多說就出於那惡者.
雖然你奉神的名來起誓.
但那句話帶著的就是一個.
狡猾的動機.

$^{761}$惡者就是撒旦.
就是一個欺控上帝.
萬惡之者.
就是將所有的罪惡.
藏在他身上.
當耶穌基督教導門徒.
就是要脫離世上的惡者的時候.
我們由說真話.
來成為我們的起點.
我們生命裡面能夠.
說是就說是.
不是就說不是.
別人不相信你.
你不要再強行.
用不同的方法來逼別人相信.
用你的行為告訴你.
你的生命是可信的.
這是耶穌基督對門徒.
一些很重要的免禮.
經文就大致解釋在這裡.
不是很複雜.
有一點典故.
我們解開了.
我們大概明白.
當時的人為了證明自己可信.
做了一些很荒謬的事情.
不過這裡再深入一點去看.
我們的生命.
我們是不是真的說得出光明磊落.
一句大話都沒有說過.
又或者我們自己的生命.
是一個完全無懈可擊的義人呢?.
負心自問我們絕對不是.
我作為牧者.
我自己都有很多悲污.
我很坦誠地跟大家說.
我自己都有很多不濟的地方.
在我們生命覺得最不光彩的那一面.
上帝依然施行拯救.
對於中國人這個文化來說.

$^{801}$一個人失去了信用.
就等於失去了生命所有.
有一句說話.
一字不忠百字不容.
就是你失敗一次.
你就永遠沒有翻身之地.
不過不是代表中國人很誠實.
中國人很強調誠實這個價值.
但是中國人其實都不是很誠實.
那個一字不忠百字不容.
關鍵在於什麼呢?.
你失敗了就沒有人看到你不誠實的地方.
所以騙人騙到底.
中國人有時文化有這種特性.
不過聖經裡面.
如果我們成為基督徒.
我們去讀聖經或讀自己的人生.
或自己的文化層面.
我們如果有不理想的地方.
上帝是原諒的.
聖經文在約翰一書一章九節.
我們若認自己的罪.
神是信實的是公義的.
必要赦免我們的罪.
洗淨我們一切的不義.
我們若說自己沒有犯過罪.
便是以神為說方的.
他的道也不在我們心裡了.
約翰一書在這裡說什麼呢?.
他同樣用誠實這個字或信實這個字.
他如果放在上帝那裡.
便成為了我們基督徒的出路.
因為上帝答應了人的事.
上帝答應我們一定會赦免的心意.
他不會改變的.
人沒有一樣的罪.
不可以被上帝赦免的.
我們反而苦心自問.
我們是否成為一個坦誠坦蕩蕩的人.
如果我們不是的時候.

$^{841}$面向上帝.
坦誠面向上帝.
他必要赦免我們的罪.
洗淨我們一切的不義.
如果我們反而說自己沒有罪.
上帝便說謊了.
究竟是你說謊還是上帝說謊呢?.
約翰一書說.
乖乖的回到上帝那裡.
上帝是一個信實的主.
上帝願意赦免我們.
幫助我們.
以致我們可以生命重新建立起來.
信任在上帝的恩典當中.
中國人覺得沒有信.
基本上便沒有價值.
但在基督教裡.
一個人曾經失足.
大至好像耶穌基督的門徒彼得一樣.
失信於人,失信於基督.
但在基督的恩典裡.
他都可以重新做人.
信任是可以重新建立的.
當我們過去有些不理想的地方.
我們祈求上帝幫助我們.
讓我們重新走回正路.
有些研究這樣說.
如果做了一些似乎那件失信的事有多大呢?.
但重新建立信任.
不是一件很難的事.
如果你在小事裡失信.
你用十倍的時間.
你便可以重新建立.
不過有些大的.
可能在你的家裡曾經有失信的行為.
一般來說.
關係的破裂.
如果你要重新建立.
兩至五年.
你真的很坦誠地努力.

$^{881}$重新建立的時候.
信任是可以建立的.
用一段時間.
兩年至五年.
你真的好好地去活現.
事就說事.
不是就說不是.
你的生命最終會被人接納.
第二個方面.
我想作為一個應用.
在信任這個題目裡.
我想不單止是說.
曾經失信於人的人要真心改過.
我們也要建立一個可以可信任的環境.
或者我們也要致力去相信身邊的人.
有時候可能我們會因為一些小成見.
或者那個人有些不起眼.
但是你看下去.
刺一刺你的一些.
你覺得失信的行為.
但我們都勉勵自己.
上帝沒有放棄他的時候.
我們也不要放棄他.
我們也給身邊的人.
有足夠的空間來改過自身.
讓我們這個群體.
在我們的家裡裡.
我們也製造一些空間.
來幫助曾經可能令你失望,失信.
或者令你陷入情感裡面很困難的人.
這個過程是很不容易的.
我也明白.
有時候我們會生氣的.
不過求主的恩典.
臨到我們生命當中.
幫助我們.
給我們有足夠的空間.
讓人生命可以真正改過.
你可以不用這麼快完全接納他.
但你不要第一時間拒絕了他.

$^{921}$我相信這個就是今天.
當我們重新建立.
誠信,信任的一個環境裡面.
大家雙方面一起互相努力.
求主繼續幫助我們.
今天這段經文不是很複雜.
但在實踐裡面.
確是要大家每一句說話.
在我們生命每一個小節當中.
加緊的來留意.
就讓我們成為一個有誠信的人.
在這個世代裡面.
不需要再細神疲軟.
來證明自己是對的.
而是我們活出一個對的生命.
求主繼續帶領我們眾人.
謝謝.
\newpage



\section{馬太福音 5:38-48}
\label{sec:DeJAKAGjSug}
\textbf{2023年3月19日崇拜直播｜陳冠成牧師｜給當代門徒的15個忠告:像天父完全一樣｜馬太福音5\_38-48}
\newline
\newline
連結: \href{https://youtube.com/watch?v=DeJAKAGjSug}{\texttt{https://youtube.com/watch?v=DeJAKAGjSug}} ~~~~ 語音日期: 2023-03-19
\newline
\newline
\hyperref[sec:O_RnYw4oPiY]{\small{< < < PREV SERMON < < <}}
~
\hyperref[sec:index_chronic]{\small{[返順時目]}}
~
\hyperref[sec:XpqKURe4134]{\small{> > > NEXT SERMON > > >}}
\newline
\newline
馬太福音 5:38-48
\newline
\begin{longtable}{cl}
\hline
\hline
章節 & 經文 (和合本修訂版)\\
\hline
5:38 & \begin{tabularx}{0.7\textwidth}{X} 「你們聽過有話說：『以眼還眼，以牙還牙。』 \end{tabularx} \\ \\ \relax
5:39 & \begin{tabularx}{0.7\textwidth}{X} 但是我告訴你們：不要與惡人作對。有人打你的右臉，連另一邊也轉過去由他打。 \end{tabularx} \\ \\ \relax
5:40 & \begin{tabularx}{0.7\textwidth}{X} 有人想要告你，要拿你的裡衣，連外衣也由他拿去。 \end{tabularx} \\ \\ \relax
5:41 & \begin{tabularx}{0.7\textwidth}{X} 有人強迫你走一里路，你就跟他走二里。 \end{tabularx} \\ \\ \relax
5:42 & \begin{tabularx}{0.7\textwidth}{X} 有求你的，就給他；有向你借貸的，不可推辭。」 \end{tabularx} \\ \\ \relax
5:43 & \begin{tabularx}{0.7\textwidth}{X} 「你們聽過有話說：『要愛你的鄰舍，恨你的仇敵。』 \end{tabularx} \\ \\ \relax
5:44 & \begin{tabularx}{0.7\textwidth}{X} 但是我告訴你們：要愛你們的仇敵，為那迫害你們的禱告。 \end{tabularx} \\ \\ \relax
5:45 & \begin{tabularx}{0.7\textwidth}{X} 這樣，你們就可以作天父的兒女了。因為他叫太陽照好人，也照壞人；降雨給義人，也給不義的人。 \end{tabularx} \\ \\ \relax
5:46 & \begin{tabularx}{0.7\textwidth}{X} 你們若只愛那愛你們的人，有甚麼賞賜呢？就是稅吏不也是這樣做嗎？ \end{tabularx} \\ \\ \relax
5:47 & \begin{tabularx}{0.7\textwidth}{X} 你們若只請你弟兄的安，有甚麼比別人強呢？就是外邦人不也是這樣做嗎？ \end{tabularx} \\ \\ \relax
5:48 & \begin{tabularx}{0.7\textwidth}{X} 所以，你們要完全，如同你們的天父是完全的。」 \end{tabularx} \\ \\
[1ex]
\hline
\hline
\end{longtable}
$^{1}$今段經文是在《馬太福音》第五章三十八到四十八.
是在《聖經》新約第七頁,請聽我讀出.
你們聽見有話說,以眼還眼,以牙還牙,.
只是我告訴你們,不要與惡人作對,.
有人打你的右臉,連左臉也要轉過來由他打..
有人想要告你,要拿你的女衣,連外衣也由他拿去..
有人強迫你走一里路,你就和他走二里..
有求你的,就給他;有向你借太的,不可推辭..
你們聽見有話說,當愛你的淪陷,恨你的仇敵,.
只是我告訴你們,要愛你的仇敵,也為那逼迫你們的禱告..
這樣就可以作你們天父的兒子,.
因為他叫日頭照好人,也照大人,.
降雨給異人,也給不異的人..
你們若單愛那愛你們的人,有什麼賞賜呢?.
就是遂理,不也是這樣行嗎?.
你們若單請你們弟兄的安,給人有什麼賞賜呢?.
就是愛邦人,不也是這樣行嗎?.
所以你們要完全將你們的天父完全一樣..
請姐妹平安,我們繼續登山寶訓的宣講..
剛才所誦讀的經文,對我們不陌生吧?.
我們可能聽過很多次,.
只不過很特別的是,我們對耶穌的某幾個吩咐,.
其實會有保留,.
覺得怎麼可能要我這樣做呢?.
有人打我的右臉,難道我連左臉也要主動轉給他打嗎?.
又或者有人要告我,要連我的內衣也要拿走,.
我反而還要給他外衣?.
耶穌的要求似乎和普世所倡議的公平公道,.
要保護自己或別人應有的權益,.
這種社會價值觀念好像相反,.
怎麼可能要我們完全按照字面的意思來遵守呢?.
特別是當自己或家人真的受別人欺負,侮辱,被人傷害的時候,.
難道還要我完全不設防,不保護自己?.
盲目任由對方在自己或家人身上佔便宜,作惡?.
怎麼說也好,我們心裡總覺得很荒謬,很離地,.
於是我們很自然地會靜悄悄地將這幾節經文放在一旁,.
還是不要處理它,可能這樣會好一點,.
甚至我們骨子裡也渴望,耶穌其實不是真的叫我這樣做的,.
弟兄姊妹,其實耶穌想跟我們說什麼呢?.
首先我們看回上文,跟大家重溫一點,.

$^{41}$我們之前也看過第五章第十七至二十節,.
耶穌說了兩個很重要的原則,.
就是說他來並不是要廢掉律法和先知的教導,.
而是要成全對,.
他也說過,門徒的義要勝過文士和化理卓人的義,.
才可以進天國,我們必須謹記這兩個前提,.
因為耶穌是根據這兩個大前提,.
來提出六個的教導,.
每一個都是用「你們聽見有話說」,.
又或者說「古人說」這一類的字眼來作為開始,.
藉以糾正當時的社群對舊約律法的曲解,.
從而說明了律法原本的意思到底是怎樣,.
又或者說什麼才是符合上帝旨意的義,.
我們必須謹記這個大前提,.
也因著這個大前提,我們之前已經說了四個教導,.
還記得嗎?.
不可殺人,不可姦人,不可休妻和不可起誓,.
現在我們就分享第五和第六個教導,.
三十八節,剛才也讀了,.
「你們聽見有話說,以眼還眼,以牙還牙」,.
如果特別看《和學本》,其實有個括號,.
它也說明這八個字是引述自《昭及記》,.
《利未記》,《新命記》都有的律法教導,.
驟眼看來,似乎真的容許人有權去報復,.
你來到去以暴易暴,.
事實上這個教導,你想想,.
在當時,其實也挺迎合社會貧苦大眾的口味,.
畢竟你想想,特別是猶太人過往幾百年,.
他們怎樣?.
經歷了國破家亡,.
又成為了俘虜,被擄到外地,成為二等公民,.
終於可以回歸耶路撒冷,有自己的地方,有聖殿,.
但其實還是活在不同的外邦政權的欺壓下,.
所以,普遍百姓心目中,自然積聚了不少怨氣,.
不少的苦毒,.
所以不難想像,其實有關報仇的教導,.
其實很有市場,.
透過報仇,人往往可以怎樣?.
將你內裡的苦毒,那些憤怒,多多少少發了出來,.
又或者透過報復,.

$^{81}$也讓欺負你的人知道怎樣,.
你不要再來,我不是很欺負的,.
我會還手的,.
所以,我們不難想像,當時是一個這樣的情況,.
事實上你會看到,即使今天報復的心理,.
都活躍在人類的生活裡面,.
有沒有留意?.
如果你留意,特別可能外國很多的調查都會研究,.
覺得有兩成以上的殺人案,.
或者六成的校園槍擊案,.
和什麼有關?.
和報仇有關,.
甚至這種報復心理,.
你會發覺,其實在政治層面都會出現,.
曾經看過一篇報道,.
BBC,英國廣播公司,.
有一篇文章分析,.
當年特朗普,Donald Trump,.
他為什麼會贏得一個總統的選舉呢?.
其中一個原因就是,.
由於當時美國的藍領白人,.
特別是住在鄉村的選民,.
他們長期覺得自己被主流社會遺棄,.
他們的訴求,他們的聲音得不到聆聽,.
長期被打壓,.
於是就一面倒轉去支持,.
當時被主流社會所不齒的特朗普,.
其實就是報復,.
用選票來報復,.
發洩他們心裡的怨憤,.
事實上,說得更近一點,.
我們都看荷里活電影,.
你有沒有發覺不少荷里活票房大賣的影片,.
和什麼題材有關?.
復仇,.
特別幾年前,.
在座很多人都看過,.
聽過復仇者聯盟系列的影片,.
在全球大賣,.
你想想,.

$^{121}$如果用原班的演員,.
原班的製作單位,.
同樣的投資,.
但是改了個名字,.
不叫復仇者聯盟,.
叫復和者聯盟,.
他們穿上裝備,.
但不是和敵人拼命,.
而是和他們和解,.
大家抱著頭抱著頸吃頓飯,.
握手做朋友,.
你會不會想看?.
看一集都夠了,.
還要看整個系列,.
不會大受歡迎,.
事實上,.
電影姐妹回到我們自己,.
苦心自問,.
有沒有試過被欺負?.
不說就當有,.
我們多少試過被欺負,.
被人傷害,.
或者遭遇不公平的對待,.
可能你被人無理解僱,.
或者同事陷你一鑊,.
偷了你的主意,.
又或者你很熟的朋友,.
借了你錢不還,.
你的家人很疼你,.
但有一天因為這件事,.
互相指責,.
又或者你的愛人,.
移情別戀,.
令我們感到委屈,.
很受害,.
很憤怒,很憎恨他的時候,.
有沒有曾經閃過一個念頭,.
要以牙還牙,.
以眼還眼?.
來到反擊地方,.

$^{161}$怎樣都要令對方感受到,.
一樣的痛,.
覺得這樣才是公道,.
覺得這樣才能下那口氣,.
就正如我們經常說,.
俗語有云,.
善有善報,.
惡有惡報,.
若然未報下一句是什麼?.
時辰未到,我們搞不定他,.
希望上天幫我們收他,.
但現實不是這樣的,.
現實很多人奉行的,.
若然未報,.
由我來報,.
我出手看天行道,.
就像耶穌時代,.
普遍人都贊同.
以牙還牙,.
以眼還眼的做法,.
不過我們首先要,.
如耶穌所說,.
回到這個律法原意,.
其實這條律法原意,.
不是教人報仇,.
相反,.
是教導人不要報仇,.
為什麼這樣說呢?.
首先我們要明白,.
在古代社會,.
在舊約的世界裡,.
很多時候有所謂的集體報仇的觀念,.
例如A部落,.
有人被B部落某人傷害時,.
A部落自然會向B部落所有人報復,.
見一個就打一個,.
甚至要置對方於死地,.
你會看到,.
如果這樣不斷循環的時候,.
報仇只會最終造成更多更大的仇恨,.

$^{201}$所以律法這裡說,.
以牙還牙,以眼還眼,.
其實目的是要減少再多的傷害出現,.
加害者當然有責任向受害者提供相應的補償,.
但受害者不能夠過分地索償,.
對方如果真的令你損失了一隻眼,.
你最多只可以要求他賠一隻眼給你,.
賠兩隻都不行,賠一隻,.
亦不能要求對方因為你失去眼的緣故,.
給你一條命,.
更加不可以傷害你家人來發洩,.
否則,我們常說冤冤相報何時了呢?.
而另一方面我們必須留意的是,.
其實以眼還眼的律法,.
其實只適用於公共審判的場景,.
你當今天法庭那樣,.
其實是用在那裡的,.
不是用在私人糾紛解決私人糾紛的場景裡面,.
審判官他依據著以牙還牙的原則,.
來量刑指引,.
讓加害的和受害的雙方,.
他們都得到合理公道的裁決,.
但是到了新約時代,.
我們不難想像在古人權威教導這種加持的底下,.
其實以牙還牙,以眼還眼,.
這個佛法的原意已經被曲解為,.
容許人你有權去報復,.
你有權自己按著這八個字,.
自己執行私刑,.
變成了一個這樣的藉口和歪理,.
於是耶穌必須出手撥亂反正,.
他在第39至42節怎麼說,.
他列舉了四個標題,.
這四個例子有一個標題就是,.
不要和惡人作對,.
根據這個標題,他講了四個例子,.
你留意,其實和之前我們已經說過,.
耶穌也說,那些如果讓你陷入情慾方面的罪行,.
那些眼和手要怎樣,.
通通挖出來斬了它,.

$^{241}$但其實我們上次說過,.
不是叫你用字面意義去理解,.
這個只是一個誇張的表達手法,.
目的是要人怎樣,.
不是一文也不要想,.
就照著100\%的字面意義去做,.
而是藉著這個看上去,聽上去好像很荒謬,.
不可行的吩咐,.
來引起所有聽到的人的注意,.
你聽到的時候,你馬上說不是吧,.
讓你思考一下,追問耶穌,.
其實耶穌究竟想怎樣,.
其實他是最大發揮這個作用,.
所以其實,.
耶穌不是要我們畢加斯索死守著經文字面的意義,.
假如我們不幸被人欺負,.
打了,刮了你的右臉,.
千萬不要把你的左臉也送給別人刮,.
因為連耶穌自己,.
當他面對無理的審判,.
他被人打臉的時候,他有沒有這樣做?.
他沒有,.
他沒有說你打我這邊,任你打,.
耶穌沒有這樣做,.
再者你會看到,.
耶穌這裡說,.
不與惡人作對,並非要求我們縱容別人作惡,.
或者盲目地息事寧人,.
你留意耶穌經常公開地指斥那些違背上帝心意,.
作惡的人,.
包括宗教領袖,.
亦都事實上對他們怎樣?.
對著幹,.
不會縱容,姑息他們,.
只不過,.
耶穌絕對不會採用以牙還牙,.
以眼還眼,.
這種報復私了的方式,.
來回應對方的惡行,.
即使有人攻擊他,侮辱他,.

$^{281}$或者試探他,.
所以我們基於這個原則,.
試試看那四個例子,.
首先第一個,按照當時的文化,.
我們說是右手主導的文化,.
右手世界的文化,.
左手只是處理那些不潔淨的東西,.
右手代表力量,權威,.
用有人打你的右臉,.
其實意思就是說,.
他不是這樣打,.
右手這樣打不到對方的右臉,.
唯有怎樣?.
他唯有反手用手背,.
這樣才打到對方的右臉,.
目的不是要傷害對方的身體,.
而是要怎樣?.
羞辱他,侮辱他,.
如果按照以眼還眼的方式,.
你覺得應該怎樣做?.
如果准你以眼還眼,.
合法的,.
你會怎樣?.
不敢說是吧?.
你可能真的會把他升級,.
是吧?.
但是,耶穌這裡說什麼?.
你不可以他做初一,.
你就做回十五,.
你不可以因為他羞辱你,.
所以你也不跟他客氣,.
羞辱他,.
耶穌這裡說,.
有人打你的右臉,.
你連左臉也拿給他打?.
其實意思是,.
你想像一下,.
當我被人用右手,.
刮了我的右臉一下,.
我的頭是怎樣?.

$^{321}$是這樣,.
我刻意把左臉轉過來,.
其實是什麼意思?.
他還可以這麼容易打我的右臉嗎?.
不容易了,.
其實把右臉轉過來,.
對方就不再用他的右手,.
繼續羞辱你的右臉,.
其實也是間接用一個行動,.
去保護上帝給我們的尊嚴,.
向對方表明,.
你沒有權再羞辱我,.
我也不會再讓你這樣做,.
因為我也有上帝的形象,.
第二個例子,.
有人打官司告你,.
要求你賠償,.
本來很合理的,.
如果你真的欠人錢,.
或者令人有損失,.
但是你已經一貧如洗了,.
身上只剩下兩樣東西,.
就是你的外衣,.
和裡面的女衣,.
按照當時的律法規定,.
對方是不可以拿你這件價值錢的外衣,.
來做抵押,.
或者說抵償,.
因為他要晚上給你,.
來當被擒,.
你知道巴勒斯坦早一夜的溫差很大,.
他不可以拿走你的外衣,.
所以其實,.
裡面那件不值錢,.
不能被擒,.
他拿來也沒用,.
他要怎樣?.
你真的不能賠償,.
沒有錢身,.
他要先放過你,.

$^{361}$再想想有什麼辦法,.
但是假如他,.
真的把你抓乾,.
他有權向你索償,.
但是連你那件很不值錢的內衣,.
也不放過,.
要拿走的時候,.
其實對他一點好處也沒有,.
不過他可以,.
可能發洩他的不滿,.
或者拿回財,.
在這個時候,.
你如果以牙還牙,.
你會怎樣?.
你可能會跟他拼命,.
不是吧?.
你這樣也做得出來?.
我什麼都沒有了,.
還要我把那件給你,.
不值錢,.
我給你也沒用,.
你死都不給他,.
但這個時候很特別,.
耶穌說你主動把外衣脫掉,.
其實想像一下,.
在一封什麼圖,.
他要你裡面那件,.
你把外面那件也給他,.
你主動讓自己在眾目睽睽下,.
變成赤裸上身離開,.
其實也是間接令對方感到羞辱羞愧的行動,.
因為在場不止你和他,.
有審判官,.
甚至有其他百姓,.
當所有人看到你搞成這樣的時候,.
其實責怪誰?.
其實大家就會知道,.
你不應該把他趕盡殺絕,.
怎樣也好,.
他欠你錢也好,.

$^{401}$什麼也好,.
他什麼也沒有,.
你還要趕盡殺絕?.
從而希望對方可以回心轉意,.
不要逼得太盡,.
可以心平氣和,.
大家坐下,.
有什麼解決方法?.
第三個例子,.
按照羅馬人的律法,.
他們是准許士兵有權強迫百姓,.
替他們背行李,.
行一里路,.
同樣可能貧苦百姓聽到這條規例,.
特別當有羅馬士兵強迫他們,.
勞役他們的時候,.
會不會生氣?.
也會的,.
特別你想像一下,.
極端一點,.
當普通百姓要向北走的時候,.
羅馬士兵說,.
不行,你要先向南走,.
走完再做自己的事,.
會不會生氣?.
也會的,.
如果以牙還牙的時候,.
你被他勞役完,.
你還會不會跟他這麼客氣?.
不會了,.
你一定會罵他,.
或者就坐他,.
甚至會怕他再得寸進尺,.
冤枉你,.
貪得無厭,.
行多一里路,.
這個時候耶穌說,.
你主動告訴對方要陪他,.
行多一里路,.
其實也是間接提醒對方,.

$^{441}$不可以再勞役你,.
除非你自願幫他,.
陪他,.
當對方聽到這樣說的時候,.
其實他也會知難而退,.
你間接提醒他,.
法例規定行一里路,.
下一里路就是你犯法,.
就算我自願也好,.
所有人看來,.
其實都是你犯了法,.
到最後一個例子,.
我們不要忘記,.
耶穌當時的聽眾,.
絕大部分都是貧苦大眾,.
找招不得萬,.
你已經夠窮夠慘,.
突然有一天,.
還有人求你,.
挖苦你,.
要求你求助,.
問你借錢,.
你會怎樣做?.
你會不會以眼還眼?.
喂,.
我比你還慘,.
你說你慘,.
我昨天飯也沒有吃,.
我也想找人救濟,.
好像你比我還好,.
不如你救我,.
以牙還牙,.
然後打發對方走,.
耶穌說,.
不要立刻那麼冷漠,.
想想有沒有辦法,.
幫得多少,.
得多少,.
幫得一里路,.
一里路,.

$^{481}$耶穌沒有叫你照顧他一輩子,.
耶穌也沒有叫他,.
不用自己想辦法,.
你一個餅,.
撕一塊給他,.
比他悠悠愁愁,.
孤立無援地走好,.
弟兄姊妹,.
耶穌舉這四個例子,.
就是要跟隨他的人,.
活出一種不一樣的做人態度,.
就是見他用一種以眼還眼的方式,.
來維護自己的權益,.
死都不放,.
又或者用一種以牙還牙的心態,.
死都要為自己討回公道,.
你留意,.
耶穌舉這四個例子,.
其實有一個共通點,.
他都是要求當事人去付出,.
有沒有留意,.
特別頭三件事,.
被人打臉,.
被人拿了女衣外衣,.
被逼走一里路,.
其實付出的,.
往往是比對方要求的更多,.
有沒有看到,.
目的其實是突顯一種,.
默默地一次又一次,.
無條件地付出,.
甚至超乎對方的想像,.
並且沒有要求對方,.
因為自己受虧損,.
而要對方賠償這個形象,.
弟兄姊妹,.
誰有這種形象?.
誰像這種形象?.
如果我們一開始說,.
原來耶穌的身份使命,.

$^{521}$就是承傳律法,.
他對舊約律法,.
所有的傳息和教導,.
其實就是在說耶穌自己,.
其實這四個例子,.
最終是指向耶穌基督本人,.
這四個例子的焦點,.
其實是在描寫耶穌的生命,.
他完全在演繹上帝律法的精髓,.
在人的面前彰顯上帝愛人的心腸,.
到底是怎樣?.
而耶穌也期望我們透過這四個例子,.
可以看到,.
其實他就是這四個例子的主角,.
一個這樣被人冒犯欺負的人,.
最終竟然會成為眾人的給予者,.
一個這樣被人傷害,.
他絕對有權為自己討回公道,.
向所有人追討賠償損失的人,.
他最終反而成為眾人的施恩者,.
聽起來是不是很荒謬?.
是不是很不可能?.
但又真的出現在耶穌的生命裡面,.
換句話來說,.
耶穌其實在這裡提醒我們,.
不要永遠以一種好像債主思維的心態,.
來對待你身邊的人,.
當別人對不起你,.
虧欠你,欺負你的時候,.
你就自動想都不用想,.
債主模式上身,.
向對方追究到底,.
用盡你一切的手段,.
甚至不惜報復,.
攬炒玉石俱焚,.
你敢害我,.
我都不讓你好過,.
勢要為自己取回公道,.
怎樣都要取回自己應得,.
弟兄姊妹,.

$^{561}$一旦我們真的有這些念頭閃過,.
這段經文正好提醒我們,.
你看一看那位為你釘十架的主,.
他會不會這樣對我們?.
停一停想一想,.
這位作為全人類最大的債主,.
他有沒有這樣惡待他的債主?.
就是那些經常得罪虧欠他的人,.
包括我們,.
包括我們的受敵,.
如果當我們這樣想的時候,.
我們就開始明白,.
我們不能只祈求上帝放過我,.
饒恕我的過犯,.
但自己又不肯放過那些得罪虧欠我們的受敵,.
所以很想請大家,.
一起看最後一段經文,.
43至48節,.
請大家一起讀出來,.
預備起!.
.
如果說我們一開始看,.
以牙還牙,.
以眼還眼,.
是一種斷章取義,.
引用律法的做法,.
現在在這裡,.
當愛你的鄰舍,.
恨你的受敵,.
就是公然篡改了律法,.
故意將恨受敵,.
這個上帝沒有的吩咐,.
添加了上去,.
來混淆當事人的視聽,.
將上帝吩咐愛人如己,.
那個對象範圍,.
大大收窄了,.
愛人如己本來沒有限,.
是一些什麼人,.
包括我們的受敵,.

$^{601}$但現在卻收窄為,.
你只需要愛你的自己友就行了,.
愛那些對你好的就行了,.
但耶穌說,.
上帝的意思不是這樣,.
因為在上帝眼中,.
我們和受敵,.
其實都是上帝私恩的對象,.
經文這裡說得很清楚,.
無論是好人壞人,.
異人不異的人,.
天父同樣有一些普遍的恩典,.
譬如像陽光與水,.
給過他們,.
使他們可以存活,.
既然如此,.
上帝的兒女就要想一個問題,.
我是否仍然要堅持去憎恨那壞人,.
那一個欺負我的人,.
傷害我的人,.
天父說都是他所愛,.
是不是很矛盾?.
是不是很不公平?.
是,.
我們對上帝和耶穌都是很不公平,.
都是很矛盾,.
上帝愛我們,.
但我們又經常惹祂生氣,.
耶穌當然知道我們的難處,.
但他同時也不想我們遠離了上帝,.
和上帝的旨意背道而馳,.
於是他提供了一條出路給我們,.
去愛你的受敵,.
愛你的對頭人,.
為那些傷害過你的人,.
試試為他們禱告,.
不知道當我們聽到耶穌的吩咐時,.
會有什麼反應,.
可能會有人覺得,.
不是吧,.

$^{641}$哪有可能,.
我死都不願意,.
還是有人試過,.
我試過,.
其實是難,.
但是可以的,.
是有效的,.
有用的,.
弟兄姊妹,.
當我們一想到,.
叫我們愛受敵,.
之所以感到反感,.
很不想做的時候,.
是因為我們有個前設,.
我們覺得愛受敵,.
就等於便宜了他,.
這麼便宜他?.
明明我已經很傷,.
被他佔便宜,.
還要我便宜他,.
還要我愛他,.
為他禱告?.
耶穌為什麼要這麼為難我?.
特別在我們很生氣的時候,.
我們真的很難不這樣想,.
但當你冷靜下來的時候,.
你試想一想,.
耶穌吩咐我們愛受敵,.
為迫迫我們的禱告,.
其實首先便宜了誰?.
首先便宜了誰?.
是我們自己,.
因為愛受敵,.
為迫迫我們的禱告,.
首先便幫我們回到上帝身邊,.
你進行上帝的旨意,.
走在上帝的路,.
你自然便住在祂的裡面,.
你便可以體驗到,.
上帝更多對你的愛,.

$^{681}$對你的恩慈,.
包容和接納,.
這樣你內心的苦澀,.
你的毒,.
才可以慢慢因為上帝的愛的緣故,.
慢慢排走,.
相反,.
我們選擇堅持憎恨對方,.
不放過他,.
你猜最後毒發身亡的是誰?.
是對方還是自己?.
所以有人說得很好,.
如果你嘗試透過不斷憎恨對方,.
以為對方會知道,會內疚,.
會難受,.
其實就等於你一直在喝毒藥,.
但你以為死的是對方,.
不會的,.
死的是你,.
沒錯,.
愛受敵,.
為迫迫你的禱告,.
確實不一定會令那個壞人變好,.
不一定,.
當我們釋出善意的時候,.
對方也不一定會大受感動,.
願意和我們化敵為友,.
不過,.
老土一點說一句,.
你不試過又怎知道?.
你未試過,又怎知道不行?.
正如耶穌吩咐我們這樣做,.
你會否覺得他是在點黑路我們走?.
你會否覺得耶穌叫我們愛受敵,.
為受敵禱告,.
是在點黑路我們走?.
你有沒有看十字架上,.
他就帶頭示範,.
怎樣愛受敵,為受敵禱告?.
你說耶穌這樣做有沒有用?.

$^{721}$你不會說沒有用,一定有用,.
如果不是,.
我們今天也不能和上帝化敵為友,.
可以一起在這裡敬拜,服侍祂,.
話說,有一位猶太拉比,.
有一天,很悠閒地,.
騎著自己的蘿花騎,.
他臉上流露著很幸福的笑容,.
因為他剛剛和其他拉比一起鑽研摩西五經,.
所以有很多得著,.
令他有很多飽足,滿心歡喜,.
在這個時候,迎面出現一個很醜陋的男人,.
這個男人很有禮貌地和拉比打招呼,.
拉比你好,願你平安,願上帝賜福你,.
但拉比看到這個人實在太醜陋了,.
他美好的心情馬上消失了,.
於是他就口沫遮攔地嘲諷對方,.
你真的很醜,.
你們城裡的人是不是也像你一樣醜?.
受到這樣的侮辱,.
這個面貌醜陋的人他恍了一恍,.
於是他就回答,.
我也不知道,.
不如麻煩你去問一問那位創造我的人,.
為什麼把我做得這麼醜呢?.
這個拉比一聽到就知道自己失敗了,.
得罪了人,.
兼且得罪了上帝,.
於是他就隨即從牢身上下來,.
很謙恭地年年向這個人道歉,.
對不起,.
我得罪了你,侮辱了你,.
請你原諒我..
弟兄姊妹,.
這個故事裡那位容貌醜陋的人,.
其實他真的可以以惡還惡,.
他可以惡言相向,.
你自己又不照鏡子,.
你不是更醜?.
像不像我們小學生吵架?.

$^{761}$當你被人說是豬的時候,.
你畢金就說,.
你才是豬,.
對方就說你蠢豬,.
你臭豬,.
你死豬,.
豬來豬去,.
最後小學生吵架,.
但是這個容貌醜陋的人,.
他是選擇為上帝多走一步,.
用恩慈,.
很有智慧地將他和對方帶到創造他們的上帝面前,.
因為他很明白,.
無論是美還是醜,.
好人或歹人,.
其實都是上帝所愛,.
裡面都有天父的形象,.
同樣耶穌說我們是天父的兒女,.
擁有天父的形象,.
其實是期望我們遇到惡人惡事的時候,.
可以為他多走一步,.
向罪人,.
向那些得罪我們的人,.
彰顯上帝的美善,.
彰顯上帝的心意,.
所以耶穌也是因為這個緣故,.
他舉了一個例子,.
就是用當時的聽眾猶太人,.
他們最憎恨覺得差人一等的瑞利和外邦人來做比較,.
耶穌說如果只是愛那些對我們好的人,.
只是愛自己的,.
向他們說早晨,.
說Hello,.
這樣其實我們和那些我們覺得很壞,.
覺得比我們差一截的人,.
其實真的沒有什麼大分別,.
耶穌自己這樣說,.
在他眼中沒有什麼優點,.
沒有什麼長處,.
也沒有上次,.

$^{801}$所以耶穌他吩咐我們要愛仇敵,.
為迫伯尼的禱告,.
是為了我們的益處,.
他很想我們完全像我們的天父一樣,.
他這個說話就好像一面用來影照我們和對方的鏡子,.
即使在我們的生命裡面,.
好像這個故事的容貌醜陋的人,.
在你成長過程中總會遇到,.
突如其來的羞辱,.
侮辱,.
甚至欺負,.
你不想的,.
但總會遇到,.
但正是因為你深知自己擁有上帝天父尊貴的形象,.
上帝看你為寶貴,.
上帝很疼你,.
所以你可以不悲不亢地面對,.
你可以依靠耶穌來努力試下以善成惡,.
完全將八福的吩咐,.
將耶穌這個使命實現出來,.
這樣耶穌說,.
你的生命就能夠越來越反映到天父美善完全的形象,.
而另一方面,.
我們要明白每一個人,.
包括我們心目中的那個壞人,.
壞人,.
都是天父手所做,.
和我們一樣,.
裡面都有上帝美善的形象,.
不過被他的罪惡,.
被他的罪行包住了,.
所以我們都需要學習在你的生命裡面,.
你嘗試為某個壞人禱告,.
既然耶穌沒有將他置諸不理,.
救恩同樣為他預備,.
耶穌要拯救他,為他犧牲,.
這樣我們也不能太過心硬,.
永遠採取視而不見,.
或者冷眼旁觀的態度,.
但確實為壞人禱告,.

$^{841}$最初是很難的,.
因為你心裡面總是有另一個律出現,.
我們很想看到惡人衰,.
很想看到他跌倒,.
甚至想他不能翻身,.
但弟兄姊妹,.
假如我們覺得很困難,.
為你生命中某個壞人禱告的時候,.
我鼓勵你,.
你為你自己禱告,.
求耶穌幫你,.
將你的情緒,.
你的想法,.
哪怕是一些很惡的想法,.
你的掙扎通通告訴他,.
求他梳理你的思路,.
安撫你的情緒,.
幫你放下以眼還眼報復的心態,.
另一方面,.
也求上帝幫你好好管理你的情緒,.
將你未處理到的怨恨,.
放在一個適當安全的位置,.
不要任由它爆發,.
甚至轉嫁在別人,.
特別是無辜的人身上,.
並且我們也求神慢慢幫我們多走一步,.
試試將你和你的對頭,.
一起帶到耶穌面前,.
一方面求上帝阻止他的惡行,.
不要讓他再傷害你,傷害別人,.
同時也引導他,.
求聖靈感動他可以悔改,.
甚至像你一樣認識耶穌基督,.
假如再碰面的話,.
求上帝也給你智慧,.
給你勇氣和忍耐,.
你懂得如何面對他..
弟兄姊妹,.
不知道在這一刻有沒有一些面孔,.
有沒有一些以往不快的事情出現在你腦海?.

$^{881}$我很想你放手,.
你試試交給上帝,.
交給耶穌幫你走出困局..
很多時候真的有利說不清,.
我們也很難靠別人的EQ來擺平,.
但是耶穌說祂可以幫你,.
我們也很願意為你禱告,.
因為我們也相信.
耶穌一定可以幫我們走近天府,.
讓我們變得像天府完全一樣,.
讓我們同心祈禱..
因為你的訊息真的將人的心浮開,.
而我們內心深處,.
我們很不想說出口,.
可能和對頭人的關係,.
一些的苦毒,一些的傷害,.
上帝其實你知道,.
上帝也知道我們沒有能力靠自己去處理,.
唯有靠你,.
靠那位十字架上的恩主,.
當我們仰望你的時候,.
當我們回到天父懷你的時候,.
我們才知道怎樣自處,.
我們才知道怎樣去面對那些傷害過我們的人,.
上帝我們求你去幫助每一位,.
當我們真的想起那位曾經傷害過我們,.
當我們曾經想起怎樣被人加害,.
又或者我們曾經傷害過別人,.
求主你都在我們中間,.
使我們可以歸向你,.
而自我們的心思意念,.
我們的情緒,我們的行動,.
我們首先交給你管理,.
讓我們知道上帝你在這裡,.
你是公義的上帝,.
你同樣是慈愛的上帝,.
你不會姑息罪惡,.
你也正視罪惡,.
但你同時對人有恩典,有憐憫,.
你有等待,有恩慈,.

$^{921}$希望人在當中迴轉,.
無論是我們,無論是我們的受敵,.
所以上帝求你藉著今天的經文,.
幫助我們放下以牙還牙,.
以眼還眼的觀念,.
讓我們學習恩慈,.
學習去愛上帝你,.
以致我們有力去愛我們身邊的人,.
愛自己,並且愛我們的受敵,.
讓我們作為天父的兒女,.
我們同樣將天父的樣式,.
你的美善向人前映照出來,.
主啊,我們願意這樣做,.
我們其實心底裡也很想這樣做,.
我們也知道唯有靠你,.
我們才能夠這樣做,.
所以求上帝你閱立我們這份心願,.
幫助我們,.
我們這樣禱告,奉耶穌基督,.
你得勝的名字,來祈求,.
阿們..
請郭城帶領我們唱回應詩..
主啊,我願意..
\newpage



\section{馬太福音 6:1-4}
\label{sec:XpqKURe4134}
\textbf{2023年3月26日崇拜直播｜余嘉敏傳道｜給當代門徒的15個忠告:擁有一顆單純服侍主的心｜馬太福音6\_1-4}
\newline
\newline
連結: \href{https://youtube.com/watch?v=XpqKURe4134}{\texttt{https://youtube.com/watch?v=XpqKURe4134}} ~~~~ 語音日期: 2023-03-27
\newline
\newline
\hyperref[sec:DeJAKAGjSug]{\small{< < < PREV SERMON < < <}}
~
\hyperref[sec:index_chronic]{\small{[返順時目]}}
~
\hyperref[sec:F3PyAVrDKHQ]{\small{> > > NEXT SERMON > > >}}
\newline
\newline
馬太福音 6:1-4
\newline
\begin{longtable}{cl}
\hline
\hline
章節 & 經文 (和合本修訂版)\\
\hline
6:1 & \begin{tabularx}{0.7\textwidth}{X} 「你們要謹慎，不可故意在人面前表現虔誠，叫他們看見，若是這樣，就不能得你們天父的賞賜了。 \end{tabularx} \\ \\ \relax
6:2 & \begin{tabularx}{0.7\textwidth}{X} 「所以，你施捨的時候，不可叫人在你前面吹號，像那假冒為善的人在會堂裡和街道上所做的，故意要得人的稱讚。我實在告訴你們，他們已經得了他們的賞賜。 \end{tabularx} \\ \\ \relax
6:3 & \begin{tabularx}{0.7\textwidth}{X} 你施捨的時候，不要讓左手知道右手所做的， \end{tabularx} \\ \\ \relax
6:4 & \begin{tabularx}{0.7\textwidth}{X} 好使你隱祕地施捨；你父在隱祕中察看，必然賞賜你。」 \end{tabularx} \\ \\
[1ex]
\hline
\hline
\end{longtable}
$^{1}$今次經訓時間.
今日讀經的.
是記載在馬太福音第六章一至四節.
新約的第七頁.
唐用聖經新約的第七頁.
請丙姐妹翻開聖經聽我讀出.
你們要小心.
不可將善事行在人的面前.
故意叫他們看見.
若是這樣就不能得你們天賦的賞賜了.
所以你施捨的時候.
不可在你面前.
你前面吹號.
將那假冒為善的人.
在會堂裡和街道上所行的.
故意要得人的榮耀.
我實在告訴你們.
他們已經得了他們的賞賜.
你施捨的時候.
不要叫左手知道右手所作的.
要叫你施捨的事.
行在暗中.
你負在暗中察看.
必然報答你.
聖經讀不.
元神賜福常讀他話語並且進行的人.
讓我按一下.
沒有了第一頁.
不要緊.
我問問你.
現在說的當代門徒是第十五宗誥.
先考考你們.
你知不知道今天是第幾講?.
(七).
七?.
很接近.
原來已經過了一半.
時間過得很快.
已經到第八講了.
不如我們快速重溫.

$^{41}$我們之前學過什麼.
第一個是教我們如何做一個真正追隨基督的門徒.
然後叫我們要作細信的光和炎.
目的是什麼?.
叫我們要讓別人看到我們的好行為.
而且要將榮耀歸給我們的天父.
第三講是.
耶穌說不要廢掉律法.
而是要去承傳律法.
叫我們要遵行祂的義.
第四講是.
教我們不要發怒.
要先和弟兄和好.
才可以獻上祭物.
第五講.
提醒我們要保持聖潔.
不要奸人.
不要丟妻.
第六講是.
叫我們不要起誓.
要做一個真誠誠實的人.
第七講.
就是要愛我們的仇敵.
要我們考法我們的天父.
要將天父完全一樣.
過去的七個忠告.
覺得容易做嗎?.
不容易?.
很難吧?.
我想問我是否要拆聲?.
是嗎?.
我用這個會好一點嗎?.
好.
過去的七個忠告.
其實在第五章裡.
在第五章裡.
應該是說耶穌在第五章裡的講述.
在第五章第一節是說登山寶訓.
直至第五章尾的時候.
就開始說仇敵.

$^{81}$中間的位置帶出了一個非常重要的主題.
就是第五章二十節說.
我告訴你們.
你們的義若不過聖幹民事.
或法利賽人的義.
絕不能進入天國.
如果我們是當時耶穌時代的人.
我們會覺得這個教導是非常震撼的教導.
因為當時的人在他們眼中.
法利賽人是誓死去遵循律法的人.
我們之前Anders陳牧師也提到.
民事和法利賽人.
他們很喜歡嚴守一些律法.
一些外在的規條.
以致達至上帝的義.
不單止他們做得很足.
而且他們更加做多.
他們將律法分成很細條.
千千萬萬的規條.
例如律法上說.
不可在安息日不可以工作.
於是民事就辛辛苦苦將工作下定義.
安息日不可以做什麼.
甚至走路也要規定.
不可以超過一九百六十米的距離.
他們跟得那麼足.
要我們去勝過他們的義.
很難的.
究竟耶穌所說的義是什麼意思呢?.
民事和法利賽人所說的義.
究竟出了什麼問題呢?.
我們作為信徒.
又應該怎樣擁有一些什麼特質.
才可以勝過他們的義呢?.
首先我們一起有個禱告.
好嗎?.
親愛的三一佛上主.
求聖靈賜給我們一份柔軟的心.
去聆聽.
並且去實踐你的教導.

$^{121}$叫到我們今天踏出禮堂門口的時候.
我們的生命不再一樣.
多謝你.
祈禱乃奉我主耶穌基督的聖名祈求.
阿們.
近期熱話.
Netflix有個紀錄片.
《以父之名》.
In the name of God.
揭發了韓國有四大異端.
我今天不是說最有名的那個攝理教.
我是想說第二個.
名為五大羊的異端.
當中有一個圖中的教主.
叫做朴信子.
他在1974年.
他有一場大病.
到處求醫.
但沒有辦法得到醫治.
但有一天他竟然無緣無故病好了.
他認為神奇.
他覺得得到上天的救贖.
於是他開始接觸基督教.
在1986年.
他開始經營了一間公司.
叫做五大羊股份公司.
專做一些工藝品.
他很重視員工的福利.
他們好像一家人初期教會.
一起居住.
公司內又會經營一間幼兒院.
幫助很多孤兒.
所以很多貧苦的人.
都非常感激朴信子.
亦很信任他.
後來五大羊的員工.
入了教的信徒.
就認為追隨聖靈的路.
就是要借錢給五大羊公司.
讓他們可以幫助更多貧苦的百姓.

$^{161}$所以信徒因為朴信子的義行.
就把自己甚至家人的家產.
都巨額借出來.
你猜他累積了多少錢?.
當時累積了一百億韓幣.
如果用現在的市價.
應該是幾千億韓幣.
是否很大? 即是很多錢.
最終被揭發.
五大羊的信徒和教主一起逃走.
最終他們被發現.
三十二人集體自殺或被殺.
其後更揭發這間公司.
原來是虛有其表.
沒有實質的生意.
甚至乎孤兒院的孤兒.
其實都是假裝出來的.
所有人都蒙在鼓裡.
全都掉眼鏡.
其實憑一個人的外表.
很難看到他的內心究竟發生了甚麼事.
教主的義行.
所謂的善事.
用我們今天的基督徒來說.
就叫做好行為.
在當時的人眼中.
他是一個非常有愛心.
非常好.
亦想不到他內裡.
原來全部都是虛偽的東西.
當時耶穌的時代.
法利賽人其實都行了很多.
我們所謂的義事,好行為.
但耶穌卻斥責他們.
說他們表裡不一.
耶穌更加警告我們.
叫我們不要效法民事和法利賽人.
在今天的《絕經文》裡也有提到.
就是要指出他們這些錯誤的行善心態.
我們看看第一節.

$^{201}$這裡說了一句「小心」.
其實第一節是一個總綱.
是一個principle.
告訴我們一件很重要的事.
是甚麼呢?.
「小心」這個字.
在中文標準的譯本裡.
解作「謹慎」的意思.
耶穌很想吸引我們的注意.
留意甚麼呢?.
就是要告訴我們.
下面的吩咐很重要.
一定要來聽.
要我們留心.
他告訴我們甚麼呢?.
就是叫我們不要將善事行在人的面前.
故意叫他看見.
究竟甚麼是善事呢?.
原文是「你的義」.
英文是「righteousness」.
這個字其實跟剛才我們說的五章二十節.
要我們勝過民事和法利賽人的義.
那個「義」字其實是同一個字.
意思是甚麼呢?.
就是要告訴我們要做正確的事.
不單是在一般人眼中.
覺得是合乎道德.
正確的事之外.
而且特別是指神面前合義的行為.
即是要合神心意.
神喜悅的才叫做義事.
好像我們第一節說的「義」.
就是善事這個字.
所謂要做一些合乎神心意的事.
正呀!是好事.
當然是對的.
問題在哪裡呢?.
問題在於背後的動機.
剛才我們說過.
民事和法利賽人會嚴守樽格的外在規條.

$^{241}$認為做到足.
這些在律法上所謂正確的事.
就會合乎神心意.
神所喜悅.
但很可惜.
他們只是越著重外在的時候.
他們就越忽略了自己的內心.
以為自己做得越多就越虔誠.
甚至用自己所謂的行為.
所謂的善事去衡量.
究竟是否配得上帝的恩典和稱讚.
自以為義,虛偽.
做給人看.
才是錯誤的動機.
耶穌在第一節當頭棒喝地警告.
如果你們故意讓其他人看見的話.
你們就不能得到天賦賞賜.
耶穌在第二節以下.
用了三個例子.
是猶太人著重表現自己虔誠的善事例子.
我們今天會說的是施禦.
即是給予的例子.
另外兩種在下兩節會說到的.
就是禱告和禁食.
我們一起看看第二節.
這裡說到.
所以你施捨的時候.
不可在你面前吹號.
盡那假冒為善的人.
在會堂和街道上所行.
是要故意讓人看見.
得到榮耀.
施捨其實是舊約裡面很重要的教導.
其實是神吩咐我們.
要周祭身邊的貧窮人.
在律法裡面.
如果細心看.
其實有很多是說到.
要照顧到貧窮的條例.
例如安息日.

$^{281}$安息日是想叫有地耕種的人.
留下出產的產物.
給有需要的人吃.
而在利未記我們看到.
第25節.
我們也有說到.
如果你看見弟兄越來越貧窮的時候.
那你要怎樣?.
要幫他.
還要和他同住.
甚至乎可以借錢給他.
或者借糧給他.
但是可不可以收得多?.
不可以收得多.
神很眷顧貧窮人.
而在詩篇41.
這一節也有說到.
原來眷顧貧窮人是有福的.
原來箴言更加說到.
憐憫貧窮人就是借給他.
這件事很正.
你施與給窮人的時候.
等同於借給耶和華.
好像代神施與給窮人.
施與給窮人.
有需要的人做善事是必須的.
我們做門徒.
神是很喜悅的事.
不過耶穌要針對的.
就是施捨背後的動機.
第二節說到.
施捨的時候.
不可在你面前吹號.
如果我們嘗試用字面去解釋.
其實也有可能.
當時有緊急需要奉獻的時候.
原來當時的聖殿會吹起號角.
你有沒有見過那個.
喇叭形的號角呢?.
召集一些當時的市民.

$^{321}$到聖殿的時候去捐獻.
我們又試試代入那個情景.
當號角聲吹起的時候.
我們每個人都聽到.
放下手上的工作.
然後跑去聖殿做善事.
是不是很震撼?.
每個人都看到我們做善事.
原來除了用字面去解釋這段經文.
也有可能是一個隱喻.
看回當時的歷史.
放在聖殿的奉獻箱.
形狀設計就像一個號角.
有研究說.
可能有兩種.
好像這個.
有兩種設計.
他們的口都是用青銅的金屬去製造.
下面的口可能是向上很大.
但會收窄.
或者是嘴巴很小.
然後慢慢引導一條管路.
把錢放下去.
目的是為了什麼呢?.
目的是防止偷竊.
我以前服侍的教會.
奉獻箱經常被偷錢.
所以我們就要…….
不好意思.
所以我們就要在奉獻箱上.
加多了很多鐵板.
才能最終杜絕偷竊事件.
所以在耶穌時代.
奉獻箱這樣用金屬設計.
其實是很合理的.
我們有試想像.
金屬是金屬.
放在奉獻箱上.
聲音會怎樣呢?.
「叮叮噹噹」.

$^{361}$如果刻意大力一點.
或者放很多.
會怎樣呢?.
更大聲.
好像在這個堂.
突然大聲一點.
電話聲也好.
我們會怎樣呢?.
立即看過去.
你覺得它的動機是什麼呢?.
想人們看到.
我們看到原來放了很多錢.
你看我多好.
值得去榮耀自己.
炫耀自己是一個好機會.
這個情境讓我想起.
小時候我看歡樂滿東華.
大家有沒有看過.
「我可以去認捐嗎?」.
打電話去.
鏡頭前拍到.
電話聲「叮叮噹噹」.
很大聲.
然後螢幕上會出現捐贈人.
和多少錢的金額.
例如小朋友捐一百.
吳珍珍捐五百.
小時候我坐在電視上.
看著字幕.
想看誰捐得最多.
如果看到自己的名字.
心裡會飄飄然.
當我分享這個例子.
給我們的堂主任聽.
他也說「我也有」.
「我也有打過去認捐」.
他是不是小時候也很有愛心呢?.
當時他告訴我.
「我還留下了一個名字」.
「你們幫我數數有多少字」.

$^{401}$他說「利傳」.
「製衣廠陳明全」.
「妹妹一百」.
「你數了多少字?」.
「不算數字的話」.
「十個?」.
「差不多了」.
「是十一個字」.
「水蛇春」.
堂主任說.
「這樣才夠喜眼」.
「一百元就可以出自己的名字」.
「而且又可以幫公司賣廣告」.
如果他不說.
我也不知道他有這個動機.
我又想「這麼有趣」.
「為什麼只出自己全名」.
「為什麼不出妹妹全名呢?」.
做善事固然是一件好事.
但在第二節.
他在經文中警告信徒.
「不要像這些假冒為善的人」.
「在會堂或街道上」.
「故意得到別人的榮耀」.
「假冒為善」這個字.
原文是表演者.
即是演員的意思.
即是本身不是這樣的.
只是演繹一個角色.
而假冒為善的人.
其實是在表演虔誠.
不是真實的.
是故意在你面前表演.
裝模作樣.
目的是要得到你們的榮耀.
榮耀即是甚麼?.
有人稱讚他.
頭上有光環那種.
我們嘗試說說有三種假冒為善.
讓大家可以反思一下自己.

$^{441}$是否有時也落入了這些圈套.
第一種是本身是邪惡的.
但他扮善良.
就像我們剛才說的教主一樣.
他是邪惡的.
魔鬼也是很厲害的.
扮演一切.
但讓人覺得他欺騙其他人.
第二種開始比較細膩.
就像我們《馬太福音》第七章.
三至五節說到.
這種假冒為善.
明明自己眼中有良目.
但要批評別人眼中有痴.
因為他們看不見自己的問題.
自以為是.
於是他陶醉在自己的演技當中.
外面的人.
你猜他們看不看到?.
也有可能看到.
這麼假.
有時我們也會這樣.
很想跟別人比較.
想告訴別人我多厲害.
我做得到.
但你做得不及我好.
所以有時也會有這些動機.
你們猜這些動機是想做甚麼?.
要抬高自己.
就是想被人稱讚.
第三種.
我們形容為最不顯眼.
最細膩.
但很值得我們留意的.
以為自己所做的事.
是為神為人.
我拼命虔誠地做好事.
去服侍.
獻上我的金錢.
但當我們細心測驗內心時.

$^{481}$原來我們是為了自己.
滿足自己的內在.
原來行善的動力.
是來自別人的讚賞.
誰是一個虔誠的人.
他一定是一個很愛耶穌的人.
於是引發我們.
越做越多.
很想追求頭上的光環.
我們要很小心.
因為我們很容易.
代入這種假冒為善的圈套.
反思我們教會的生活.
有很多機會.
讓我們被人看到.
我們好的行為.
例如我們信徒.
很多建堂,礦堂,裝修等等.
獻上的金錢.
還有我們每個月的十一奉獻.
金額的大小.
其實很容易被人知道.
甚至我們有些事奉崗位.
很容易被我們.
帶上了一些光環.
跌入了這些陷阱當中.
我們近年有很多社區服務.
例如我們去探望無家者.
基層的家庭.
派一些物資給他們.
反思我們內在的心態.
究竟我們想拍張照.
發上社交平台.
希望別人讚好我的動作.
還是我們出於感動.
很真誠地支援他們.
其實我很介意別人怎樣看我.
即使我有單純的動機.
但我也很想別人稱讚我.
肯定我.

$^{521}$有一個姐妹邀請我到一間礦堂的教會.
一間學校教會.
她一邊帶我看礦堂的設施.
一邊講解給我聽.
將來你捐獻了這麼多.
這裡就會有這些設施.
又講了一個很有趣的事.
如果你捐獻到某個金額.
課室裡面就會有一個金牌.
寫上你的名字.
還會邀請你去開幕典禮.
我一邊聽的時候.
一邊計算金額.
金額也不算很多.
捐多一點.
又可以為神.
為學生的福音工作.
獻上我的金錢.
也算是一件好事.
但又想一想.
如果我捐得太少.
不知道這個姐妹會怎樣看我.
於是就覺得不太好意思.
就在糾結.
但當我坐下來想一想.
其實捐錢就是捐錢.
為什麼那麼介意姐妹怎樣看我.
這些光環我要來做什麼.
說到底.
我內心的動機就是想被人看見.
得到這位姐妹的稱讚.
這種施予的善行其實是扭曲了.
因為我的焦點.
是放錯了在別人稱讚我.
或者是我自己得著的榮耀上.
你會不會跟我一樣.
也會有時擁有這份錯了的動機.
同樣這個民事和法理錯人的問題.
都是焦點放錯了.
放錯了喜歡愛人稱讚.

$^{561}$多於愛神稱讚.
所以這裡告訴我們.
我們要很小心.
看回這個經文.
這裡說到.
我實在告訴你們.
這些追求光環的人.
已經得到其他人的賞賜.
所以你不會再得到天賦的賞賜.
試想想人的賞賜是怎樣的.
短暫的.
唯有天賦的賞賜才是永恆.
所以我們要很小心.
檢視我們有沒有代入假冒為善的圈套.
但其實我們很多時候行善的動機.
本身是可以純正的.
就好像法理錯人.
其實他之所以在律法上設一些限制.
這些界限.
其實他們就是不想越過上帝的界線.
他們不想得罪神.
所以才定這些律法.
這些細條出來.
我們嘗試回顧一下昔日猶太人的歷史.
他們是怎樣的?.
利器上帝.
又敬拜偶像.
結果他們經歷亡國.
被擄到巴比倫.
甚至到後面.
就說到敘利亞的暴君.
安提阿四世.
他更加毀壞耶路撒冷.
然後褻瀆聖殿.
等等這麼多的慘痛歷史.
所以耶穌年代的猶太人.
其實他們是痛定思過.
他們堅決守住律法.
和堅持維持猶太教的傳統.
所以他們才會加這麼多條例.

$^{601}$目的是想幫自己.
也幫人民去遵行義.
可以做出來.
但很可惜他們越做得多.
就越故意想別人看見.
很想為自己頭上多了一些光環.
更慘的是.
他們把這些條例.
加了這麼多瑣碎的規則.
令到自己和其他人.
都加了很多重擔重壓.
這只不過是自以為宜.
並不是耶穌所說的義行.
因此我們真正要遵行神的義.
我們怎樣去勝過.
民事和法理賽人的義.
我們必須要知道.
究竟我們什麼才是.
神眼中正確的心態.
我們看最後一段.
第三節說到.
耶穌吩咐我們施捨的時候.
不要讓左手知道.
我們右手做了什麼.
什麼意思呢?.
意思是要我們暗中施捨.
隱密到幾乎.
我們右手做了的事.
左手也不知道.
這麼低調.
雖然沒有人看見.
但耶穌很肯定地告訴我們.
我們的天父爸爸.
我們一定會看見.
而且經文說.
祂會在暗中察看我們.
必定報答我們.
這裡又引申一個問題.
之前我們記得Raymond說過.
不是Raymond說的.

$^{641}$是耶穌說的.
就是要我們做世上的炎和光.
要人看見我們好行為.
現在耶穌叫我們暗中做.
人們怎會看見我們好行為.
讓我們的生命吸引到認識主呢?.
我們又看回之前的經文.
五章十六節說.
耶穌吩咐門徒.
要我們做出什麼好行為.
動機是要讓人稱讚我們天上的父.
將榮耀歸給他.
而不是歸給我們自己.
所以如果你想炎耀自己.
人們只會將光環放在頭上.
而頭上所發出的光.
根本不能反照出基督的香氣.
令人去榮耀主.
試試反思一下.
作為施與者.
若我們用一種高高在上的心態.
去展示出來.
告訴別人.
我有能力來幫你.
這只不過是一個同情.
不是同情.
同情只是說.
我們從高處看下來.
你很可憐.
我來幫你.
而不是同情.
同情是從高處走下來.
到人群跟他說.
我們打從心底.
同理對方的感受和需要.
而真心出於愛人如己的心腸.
這個心腸是符合神心意的憐憫.
唯有真誠出於愛人如己的心腸.
才能吸引人認識我們的主.
我們的父神.

$^{681}$我曾經去過一個長者的喪禮.
多年來.
我見這位長者.
每星期回來都會幫忙接情序表.
不算是喜愛.
但又很可愛的長者.
喪禮裡面.
他的兒子分享.
當他收拾爸爸的遺物時.
他發現原來爸爸.
在養了幾個山區小朋友.
而且他很多年來.
一直捐獻宣教機構.
發展內地的福音工作.
十年如一日.
從來沒有間斷過.
我聽到這個分享.
對我來說.
其實我內心很震撼.
我很被他那份謙卑的生命所為感動.
他的動機.
很肯定不是想別人讚賞.
否則連兒子也不知道.
雖然他在世時沒有人知道.
沒有人知道他所做的善事.
但天父是知道的.
因為從今天的經文.
耶穌是應許告訴我們.
天父爸爸是知道的.
並且他一定會欣賞和讚賞他.
這份單純服侍主的心.
你想不想擁有這樣的信仰生命?.
我有兩個應用.
想帶回家一起實踐.
第一個是求聖靈幫助我們.
幫助我們懂得自省.
時常檢視我們內心的內在動機.
究竟我想取悅人.
還是我真的想取悅上帝?.
有沒有一些事情我很介懷的?.

$^{721}$我經常介懷我事奉的得失成敗.
我看重回報.
究竟我事奉的奉獻.
有多大有多少.
或是我事奉的崗位.
角色有多高有多低.
有沒有印象別人讚賞我.
當別人讚賞我時.
我就感覺到飄飄然.
或是當沒有人讚賞你時.
你會否有心有不甘的感覺?.
記住耶穌叫我們.
要勝過民事和法利賽人的義.
我們要時刻警醒.
不要代入假冒為善的圈套.
第二就是要暗中行善.
全面地服侍主.
其實施予幫助給有需要的人.
其實是我們信仰.
做信徒的時候必須要做的.
耶穌吩咐我們.
不是不做善事.
我們一定要做善事.
但祂叫我們要暗中行善.
免得我們自誇.
其實行善的主要動機.
是我們回應上帝的愛.
我們所擁有的一切.
包括金錢,能力,恩賜.
各樣都是神賜予我們.
有什麼好誇后的.
如果我們真心愛天父爸爸.
自然會怎樣?.
會疼愛他,尊敬他.
自然地愛我這個家庭.
自然地不計較.
有什麼需要我們便幫忙.
我們會甘心樂意獻上.
作為父母的我們.
我們會很明白.

$^{761}$其實爸爸不是很著重我們做多少.
我們做什麼.
除非做了很壞的事.
而是他很著重看我們的心.
天父爸爸其實很喜悅我們.
擁有一份單純事奉他的心.
而這份單純土神喜悅的行義心態.
正正就是聖歌,文士和法利賽人的自以為是.
求主幫助我們.
我們一起有個禱告.
厚似白麥的主.
我們要謙卑彌度你的金錢.
向你躺開心扉.
去認罪悔改.
求你用你的真光去光照我們內心.
最深最深的深處.
願意讓我們被你的生命.
被你的教導去更生改變.
其實我們不得不承認.
我們心裡有很多不純正的雜質.
或多或少.
好像文士和法利賽人.
我們都想要威風.
想贏得別人的稱讚.
我求你幫助我們.
聖靈求你幫助我們.
時刻幫助我們去檢視我們的內心.
去指出我們這些假冒為善的動機.
好讓我們能夠單單地去逃離喜悅.
清心地單純地去事奉你.
去回應你的愛.
又但願望朕的弟兄姊妹.
能夠一起去實踐愛人如己.
能夠清心地為我們這個屬靈的家.
和我們身邊的鄰舍.
甘心樂意,全然地去獻上.
以上不完全的禱告.
乃奉我主耶穌基督的聖名祈求.
阿們.
\newpage



\section{馬太福音 6:5-15}
\label{sec:F3PyAVrDKHQ}
\textbf{2023年4月2日崇拜直播｜陳明泉牧師｜給當代門徒的15個忠告:禱告無添加｜馬太福音6\_5-15}
\newline
\newline
連結: \href{https://youtube.com/watch?v=F3PyAVrDKHQ}{\texttt{https://youtube.com/watch?v=F3PyAVrDKHQ}} ~~~~ 語音日期: 2023-04-02
\newline
\newline
\hyperref[sec:XpqKURe4134]{\small{< < < PREV SERMON < < <}}
~
\hyperref[sec:index_chronic]{\small{[返順時目]}}
~
\hyperref[sec:HNqZXsZMfMU]{\small{> > > NEXT SERMON > > >}}
\newline
\newline
馬太福音 6:5-15
\newline
\begin{longtable}{cl}
\hline
\hline
章節 & 經文 (和合本修訂版)\\
\hline
6:5 & \begin{tabularx}{0.7\textwidth}{X} 「你們禱告的時候，不可像那假冒為善的人，愛站在會堂裡和十字路口禱告，故意讓人看見。我實在告訴你們，他們已經得了他們的賞賜。 \end{tabularx} \\ \\ \relax
6:6 & \begin{tabularx}{0.7\textwidth}{X} 你禱告的時候，要進入內室，關上門，向那在隱祕中的父禱告；你父在隱祕中察看，必將賞賜你。 \end{tabularx} \\ \\ \relax
6:7 & \begin{tabularx}{0.7\textwidth}{X} 你們禱告，不可像外邦人那樣重複一些空話，他們以為話多了必蒙垂聽。 \end{tabularx} \\ \\ \relax
6:8 & \begin{tabularx}{0.7\textwidth}{X} 你們不可效法他們。因為在你們祈求以前，你們所需要的，你們的父早已知道了。」 \end{tabularx} \\ \\ \relax
6:9 & \begin{tabularx}{0.7\textwidth}{X} 「所以，你們要這樣禱告：『我們在天上的父：願人都尊你的名為聖。 \end{tabularx} \\ \\ \relax
6:10 & \begin{tabularx}{0.7\textwidth}{X} 願你的國降臨；願你的旨意行在地上，如同行在天上。 \end{tabularx} \\ \\ \relax
6:11 & \begin{tabularx}{0.7\textwidth}{X} 我們日用的飲食，今日賜給我們。 \end{tabularx} \\ \\ \relax
6:12 & \begin{tabularx}{0.7\textwidth}{X} 免我們的債，如同我們免了人的債。 \end{tabularx} \\ \\ \relax
6:13 & \begin{tabularx}{0.7\textwidth}{X} 不叫我們陷入試探；救我們脫離那惡者。因為國度、權柄、榮耀，全是你的，直到永遠。阿們！』 \end{tabularx} \\ \\ \relax
6:14 & \begin{tabularx}{0.7\textwidth}{X} 你們若饒恕人的過犯，你們的天父也必饒恕你們； \end{tabularx} \\ \\ \relax
6:15 & \begin{tabularx}{0.7\textwidth}{X} 你們若不饒恕人，你們的天父也必不饒恕你們的過犯。」 \end{tabularx} \\ \\
[1ex]
\hline
\hline
\end{longtable}
$^{1}$今天的經份是記載在.
《馬太福音》六章五至十五節.
是你們手上聖經新約部分七至八頁.
聖經這樣記載.
「你們禱告的時候.
不可像那假冒為善的人.
愛站在會堂裡和十字路口上禱告.
故意叫人看見.
我實在告訴你們.
他們已經得了他們的獎賞.
你禱告的時候.
要進你的內屋關上門.
禱告你在暗中的父.
你父在暗中察看必然報答你.
你們禱告.
不可像外邦人.
用許多重覆話.
他們以為話多了必蒙誰聽.
你們不可效法他們.
因為你們沒有祈求以先.
你們所需用的.
你們的父早已知道了.
所以你們禱告要這樣說.
我們在天上的父.
願人都尊你的名為聖.
願你的國降臨.
願你的旨意行在地上.
如同行在天上.
我們日用的飲食.
今日賜給我們.
免我們的債.
如同我們免了人的債.
不叫我們遇見試探.
救我們脫離凶惡.
因為國道,權柄,榮耀.
全是你的.
直到永遠.
阿們」.
你們要恕人的過犯.
你們的天父也必要恕你們的過犯.

$^{41}$你們不要恕人的過犯.
你們的天父也必不要恕你們的過犯.
聖經讀得.
回庭安.
答進這個禮拜是我們基督教傳統裡面.
叫聖咒的一段時間.
很多教會的禮儀裡面.
今天是忠子主日.
很多教會傳統裡面.
弟兄姊妹會拿著忠說之.
迎接耶穌和撒勒.
應當稱頌的.
來歌頌耶穌.
不過在這個禮拜裡面.
亦都說到耶穌由最受民眾的歡迎.
去到受苦節當中.
被百姓唾棄.
藉著一個最謙卑.
最一種最痛苦的樣式.
來去成就這個救恩.
在這個禮拜裡面.
其實我們有很多的祈禱.
我們藉著思念基督耶穌的生命.
藉著思念基督耶穌為我們所做的一切工作.
我們很深刻地.
重新將自己整理.
思想基督耶穌在我們生命裡面的工作.
但當我們一起去禱告的時候.
有些弟兄姊妹會問.
牧師,其實我怎樣的禱告才是合宜的呢?.
我要怎樣的方式.
我用甚麼字眼.
可以成為我更加深刻的祈禱.
於是就會問.
你怎樣令我禱告更加盡心.
不過今天這段經文.
雖然是耶穌基督教我們怎樣去祈禱.
但是祂教的內容.
不是教我們怎樣用字變得更加優雅.
或者讓我們在禱告裡面.

$^{81}$帶著更加讓人覺得.
我的態度更加誠懇和謙卑.
以致我的祈禱得到更加好.
上帝的垂聽.
今天這段經文其實耶穌都頗顛覆.
當時人們對禱告的一些觀念.
希望我們今天帶著這段經文.
特別是在耶穌裡面用幾場篇幅.
來教我們怎樣去禱告.
讓我們可以進入一個更加深的一種.
即是和祂之間的關係.
其實老實說,怎樣用字眼呢?.
其實今天不用問我.
問電腦已經可以了.
我早幾天用上網看過一個人工智能.
checkGPT.
幫我寫一個圖文.
他的用字比我好很多.
我想我們今天不是學甚麼用字.
學學我們怎樣去建立一個正確的心態.
向神祈禱.
我們先交托,求神幫助我們和心祈禱.
讓我們今天一起學習你的話語.
幫助我們,讓我們能夠在你的教導當中.
我們更加明白禱告的心靈.
撇除我們過往很多時候.
我們想錯了或誤解了你的地方.
求你在今天早上對我們每一個人說話.
願你的話語光照我們.
恭敬將來的時間交給你.
幫助每位弟妹聽得明白.
宣講的講得清楚.
獻上我們禱告.
奉靠基督耶穌得勝的聖名祈求.
Amen.
我們要知道.
我們一直看整個講道的系列.
這十五講.
講的是耶穌登山寶訓的內容.
雖然我們分了不同的段落.

$^{121}$分了不同的.
每一次有不同的主題.
但是這是講耶穌整體的講論.
所以在我講的內容裡面.
耶穌因為一氣呵成地講.
他要講的內容.
所以觀點是一脈相承的.
不會跟前後有矛盾.
甚至講不同的內容.
會回應之前所講的內容.
所以大家不要覺得憎氣.
反而我們一直按照耶穌的思路.
去想我們今天怎樣去作門徒.
今天所講的祈禱.
頭半部分就是按照上星期.
余姑娘所講的內容.
怎樣去作善事.
善事的心態.
來應用在禱告當中.
在六章第五至第六節裡面這樣講.
你們禱告的時候.
不可以像那些假冒為善的人.
愛站在會堂裡和十字路口上禱告.
故意叫人看見.
我實在告訴你們.
他們已經得了他們的賞賜.
你禱告的時候.
要盡你的內屋關上門.
禱告你在暗中的乎.
你負在暗中察看.
必然報答你.
你留意到在這裡所寫的內容.
如果你看回聖經就知道.
或者你記得上星期講.
其實很多字眼的句式.
和上星期講的很近似.
特別講到有些字眼是重覆又重覆的.
故意叫人看見.
然後再講.
他已經得到他們的賞賜.

$^{161}$然後再講禱告暗中的乎.
暗中察看就必然報答你.
故意,賞賜,暗中和報答.
這個句式其實在上星期出現過.
簡單來說.
耶穌說你禱告的時候.
你不要特意要讓人看見你怎樣祈禱.
不要告訴人你是用什麼方式.
表達你的敬虔.
我們也有公禱的時間.
不過重點在於.
禱告的表達.
你和神的聊天.
不是一種表演.
值得注意的是.
耶穌基督講這番說話.
其實是回應當時的法利賽人和民事怎樣教人祈禱.
或者他們禱告的形態.
他們喜歡站在十字路口.
最多人走過的路口.
或者在會堂裡.
人才濟濟坐在這裡.
最多人聆聽.
在最多人面前表演他們禱告的敬虔.
用語.
那種對神忠誠的表達.
不過耶穌說你不要學他們.
首先用法利賽人和民事成為了他要顛覆的對象.
當時很多人說.
祈禱就是一個宗教的禮儀.
是一個表達我對神之間的關係的平台.
不過耶穌說.
祈禱不是這樣的.
祈禱不是一個表演你自己有多屬靈.
或者你自己的靈性有多高的一個指標.
反而祈禱是一個很秘地裡和上帝之間的關係.
值得留意另外一個字眼.
我們很多時候很快就跳進去.
就是假冒為善這個字.
假冒為善即刻說就是裝模作樣.

$^{201}$裝好.
當然字意是對的.
不過原文裡面用這個字.
hypocrite.
是用什麼意思呢.
原來這個字本身是說戴面具的意思.
什麼人會戴面具呢.
當時希臘裡面做的舞台劇.
那些人不是用你的面形來演戲的.
是會戴著不同的面具.
藉著戴著這些面具就扮演不同的角色.
因著劇情的需要.
可能戴個快樂的面具.
或者戴個憂愁流眼淚的面具.
用這個面具來呈現他要演繹的角色的內容.
耶穌用這個字來形容.
法利賽人的祈禱.
本身戴面具在當時社會裡面.
不是一個特別貶義的字眼.
但如果用在祈禱當中的時候.
這個禱告就好像成為了一個表演的平台.
那些法利賽人和民事.
藉著這個跟神本身溝通的機會.
就成為了表達自己表演演戲的一個平台.
所以那個假冒為善.
就是那些人的動機錯誤了.
用了這個空間.
用了這個時間.
來呈現他自己的敬虔.
讓這件事成為自己展示肌肉的一種方式.
當我們明白這個道理.
其實裡面所說的東西是不深奧的.
如果用上星期的內容來說.
耶穌呼籲我們.
其實真正的禱告.
不是一些所謂表演式的.
或者公開式的.
來將一些很漂亮的字眼堆砌出來教你.
或者引導你去祈禱.
相反來說.

$^{241}$耶穌基督所強調的那個是一個很個人化.
某程度是一個你與上帝之間很親密的關係.
在經文裡面特別說到.
禱告你暗中的負.
你負在暗中察看.
必然報答你.
暗中的意思不等於一定要有一個黑房.
你不用在家裡最暗黑的角落.
你在那裡祈禱才算最敬虔.
不是這個意思.
那個暗中的意思就是說.
你不要將禱告成為一個平台.
反而是成為你生活的一部份.
很低調的,沒有人知道的.
你祈禱也沒有人知道.
你是為了很多事情來守望,在祈禱.
那個暗中的意思.
其實是很低調的意思.
想一想,弟兄姊妹.
有時候我們的祈禱.
我們今天反而相反.
我們不太想祈禱.
為什麼呢?.
因為我們有一種心態.
在小組裡面,或者弟兄姊妹當中.
或者在家裡面,借飯祈禱.
你會覺得很尷尬的.
你不知道怎麼做的.
最尷尬的地方就是你怕自己.
祈禱的時候那些字眼很刺眼.
用詞不夠好.
所以就很擔心.
不知道你帶領的時候會不會妨礙了人家.
禱告的那個心.
不過在這裡,經文裡面提醒我們.
其實那個禱告本身是一個很自然不過的.
一種的表達.
不需要太過介懷別人怎麼看你的祈禱.
當我們很介懷的時候.
我們有意無意之間就進入了法理賽人的思路.

$^{281}$就是很想讓人知道我是一個怎樣禱告的人.
但是經文正正就是要針對.
其實我們不需要太過介意.
我們的用詞顯如.
我們是不是說得優雅呢?.
我們的內容是不是合邏輯呢?.
是不是令到別人很擾樑三日的去祈禱呢?.
別人會不會記住我的祈禱呢?.
這些通通都不是我們需要考慮的地方.
我們只是帶著單純的心.
去到上帝面前來祈禱.
有弟兄姊妹曾經跟我說過.
牧師,其實平時跟同事吃飯.
那個借飯祈禱是一個見證.
這麼有趣,為什麼這麼說呢?.
他說當大家把碟子放在那裡的時候.
菜都來了,第一件事不是拍照.
而是在別人面前這樣祈禱.
還要很敬虔的.
帶領同事們祈禱.
或者自己自浸地祈禱.
告訴別人我很尊主偉大的.
那個祈禱就很棒了,有一個見證.
我說你這麼說.
表面上我知道你是敬虔的.
但實際上,到最後你將禱告成為一個表演自己.
即使你的動機是希望別人知道你是基督徒.
或者上帝在你的生命裡面作主.
但如果你做一些刻意的事情.
反而你還弄巧反拙.
你還達不到祈禱的真正目的.
所以當你有食物來的時候.
你帶著這樣的心領受.
你看到食物很好吃.
你多謝上帝.
不用將食物變得很儀式化.
其實那樣東西進入生活當中.
那個才是真正禱告的那種經歷.
第二件事,很多時候有些弟兄姊妹.
將禱告成為了.

$^{321}$今天有一種感覺叫做儀式感.
我不知道大家有沒有什麼.
即是做每一樣東西.
將生活所有的細節.
變成一個很有儀式的表達.
其實本身是很高尚的處理.
不過過份的儀式感.
其實你被人感覺到你是繁文績節到不得了.
即是你變成了一個沒有儀式.
或者徒有儀式.
但沒有實在的一個.
沒有什麼踏實感的一個人.
在禱告的教導裡面.
耶穌提醒我們.
將禱告成為你生活的一部份.
低調,即使用字很拙劣也好.
你用這樣的形式來表達.
你的天父會報答你的.
如果你在別人面前表達.
即是呈現你很敬虔的樣貌.
你已經拿到你自己地上的賞賜.
你要得到的存在感已經拿到了.
反而天上或上帝親自報答你的東西.
已經被地上的東西取代了.
求主幫助我們.
所以在第一件事我們要知道.
禱告的用詞顯然不是最重要的.
反而在禱告的過程裡面.
我是否真誠和低調地向上帝表達.
就著這一點我想再引用一個神學家所說的內容.
潘霍華,大家看過了.
可能有些頂尖姐妹經常聽這個名字.
他是德國人,不是中國人.
他姓潘,只是一個譯名.
他在一本書叫作《作門徒的代價》裡說到.
作門徒一個很重要的特徵.
特別是在禱告的內涵裡.
就是要變得forgetfulness.
forgetfulness是一個比較貶義的字眼.
是說有很多東西忘記了,善忘.

$^{361}$這個字是一個貶義.
潘霍華說在禱告或實踐信仰裡作為門徒.
我們要成為一個forgetfulness的一個人.
他繼續說的意思是.
在你的生活裡,如果真的要操練禱告的時候.
你要盡量讓自己變得隱藏.
讓你的愛心也好,讓你為別人的代禱也好.
將這一切的工作低調地進行.
他甚至說到極致.
真正的愛心是低調的,是隱藏的.
顯露出來的愛心,顯露出來的祈禱.
其實已經拿了地上的賞賜了.
但真正為別人帶來祝福的.
永遠做了也沒有人知道的.
這種禱告的生活.
其實就是帶來一種真正的操練.
再退一步來說.
其實禱告,告訴別人我為了別人.
或者我做了一些事情來禱告.
其實你越低調,才更加適合那個原理.
因為其實你禱告的時候,你做不到的事情.
你就求神來做.
簡單來說,那些事情不是你做的.
你就很大聲告訴別人,我叫神幫你做事.
其實在概念上,沒有什麼值得誇口的.
沒有什麼值得表演的.
因為做公的那個是上帝.
真正要改變,更新的那個是上帝.
人只不過能夠有份參與在當中.
你愛這個弟兄,你愛這個姐妹,你為他守望.
將他放在你的禱告裡面.
求上帝有恩典臨到在他身上.
其實做這些事情,沒有值得可以誇口的地方.
反而越隱藏,將你自己變得默默的成為別人的祝福.
那種守望,那種對自己的存在感不再那麼著緊的時候.
那種祝福,那種心.
反而是成為一個真正最謙卑與人同行的.
一個最重要的那種態度.
所以你記住了,在禱告裡面.
第一個,耶穌基督教我們的功課.

$^{401}$禱告不是用來測量我們的存在感.
反而是透過很低調默默的守望.
讓我們不是成為一個要將焦點,將聚光燈放在我們身上的人.
反而我們不斷地學會.
Forgiveness,將自己放在離開中心的位置.
將主放在中央就可以了.
這個是第一點.
第二點,第六章七至第八節.
你們禱告不可像外邦人,用許多重覆的話.
他們以為話多了就必蒙誰聽.
你們不可效法他們.
因為他們沒有祈求以先.
你們所需用的,你的父已經知道了.
剛才說的是法利賽人成為顛覆的對象.
接著去到第二個部分.
耶穌以外邦人的祈禱.
其實外邦人也不是向上帝祈禱.
外邦人向他自己的神明祈禱.
用外邦人成為顛覆或某程度要解構的對象.
叫那些聽著登山寶訓的百姓.
不要學那些外邦人方式的禱告.
讓這些東西成為你的祈禱.
所以在這裡首先他講的是外邦人.
外邦人怎樣祈禱呢?.
經文裡面說重覆的話.
我們不太明白什麼叫重覆的話.
原來重覆的話本身的意思.
就是喃喃自語.
沒有意思地,不如說那些叫做Magic Word.
一些咒語.
說某一些字眼.
唸來唸去,不停地唸的時候.
就可以操控或祈求神明不要忘記自己.
以致你的祈願可以達成.
聖經原文很有趣.
他翻譯這個字類似blah blah blah.
重覆話就是這個意思.
你在上帝面前不要說blah blah blah.
不停地說沒有意思的重覆字眼.
你不要想今天怎會這麼傻.

$^{441}$如果我們留意一些其他的宗教.
坊間裡面有些所謂的操練方式.
其實是同一句話不停地唸.
不停地唸的時候.
就好像進入了一種境界一樣.
甚至今天有一些靈修.
基督教的靈修.
可能捉住聖經某一句.
不停地唸.
唸一千次一萬次.
唸了之後你就頓辣明白那個句子的意思.
聖經裡面不是這樣教我們祈禱的.
聖經不是叫我們要元神出竅.
唸同樣的字眼.
沒有意義的.
即使你覺得後來你賦予他意義的.
重覆又重覆的形式.
那種禱告.
上帝必蒙誰聽.
在這裡耶穌用了幾串的說話來說.
他們以為說多了就必蒙誰聽.
重覆的字眼.
重覆的話不停地說.
說幾十萬次.
其實功效只是一次.
一樣的.
跟你說一次是沒有分別的.
不是說得多.
就必蒙誰聽.
耶穌說你不可效法他們.
因為你們沒有祈求以先.
你們所需用的你們的父已經早已知道了.
耶穌後面再引述.
我們基督徒或者我們作為門徒的人.
和當代裡面那些外邦人所信奉的神明是不同的.
那些外邦人信奉的神明.
因為他們某一種安全感.
他們擔心他們的神會忘記他們.
他們覺得要喃喃自語.
不停地上口一些咒語.

$^{481}$一些說話.
以致上帝所信奉的神明會記得他們.
知道他們.
當他們很忙碌的時候.
聽到他們的聲音.
他們就會按著他們的恩典.
來報答這些外邦人.
外邦人是這樣想的.
所以他們看為何得不到祝福.
因為天上的神明聽不到他們的祈禱.
或者忘記了他們.
不過耶穌基督說.
我們所信奉的上帝不是這樣的神.
我們的神不是一個.
你不提他他不知道的上帝.
反而是在我們未祈求之先.
我們都不知道自己祈求什麼.
上帝已經知道我們的需要.
所以在經文裡.
他不單是顛覆祈禱.
不要用什麼字眼.
不是用方法的錯誤.
來解拆這些想法.
他更加要糾正那些人對神的觀念.
那個神觀.
我們的神是一個知道人的需要的上帝.
人只需要帶著信心.
在上帝面前祈求.
我們每一次去祈求的時候.
在我們覺得自己需要的東西當中.
可能我們眼前覺得這些東西很重要.
但上帝總會將最好的賜給我們.
有時候上帝不按我們的心意.
給我們一些東西.
不是上帝要作弄我們.
或者上帝喜歡糟蹋地上的人.
反而是透過我們在過程裡表達.
調教自己的心意.
更加知道什麼對我們生命是好的.
其實同樣的主題.

$^{521}$或者繼續上帝將好的東西賜給我們.
其實在六章第25至34節會繼續說.
我在這裡不說了.
就跟公務員說了.
不過我想引用一句.
在經文裡面說到.
我們在生命裡面.
上帝的旨意.
是按著他的信實,慈愛來賜給我們.
即使那些小信的人.
小信的人在野地裡面的草.
今天還在.
明天就丟進爐裡燒了.
神還給他們這樣的裝飾.
何況你們呢?.
你們要憂慮吃什麼,喝什麼,穿什麼.
這都是外邦人所求的.
你們需要的這一切東西.
你們的天父是知道的.
外邦人的祈求.
嘴裡面淋淋自雨.
希望拿到他們自己想要的東西.
不過在這裡面耶穌說.
你知道的東西.
你想要的東西.
其實上帝都知道.
在你未祈求以前.
上帝已經賜給你.
有祂的美善旨意.
我們只是透過禱告.
慢慢更新一下.
有時候上帝不給你這些東西.
可能有祂美好的心意.
細緻到你追求的女孩.
那個女孩你追不到.
年輕的弟兄.
你追不到那個女孩.
焉知非福.
十多二十年後.
你已經看到真章.

$^{561}$又或者你眼前想得到.
某一些物質.
或者你要追求的某一些生活形態.
可能你覺得這一刻上帝不給你.
你忿忿不平.
不過日子有功.
你就會看到上帝的心意是好的.
上帝會將最好的東西.
放在你的身上.
所以在禱告裡面.
我們要學習第二要功課.
我們要有信心.
上帝在我們未祈求之前.
祂已經知道我們的需要.
上帝會將好的東西賜給我們.
上帝是一個美善的上帝.
祂會帶領地上跟隨祂的人.
將最好的東西放在祂的身上.
眼前你未必明白.
不過透過禱告.
一路察驗.
一路洞察上帝的心意.
以致到你一路走人生的路的時候.
你就發現上帝美好的帶領.
這個是第二點.
第二個禱告裡面我們要學的.
就是信心.
第三個.
六章第九至十一節.
所以你們禱告要這樣說.
我們在天上的父.
願人都尊你的名為聖.
願你的國降臨.
願你的旨意行在地上.
如同行在天上.
我們日用的飲食.
今日賜給我們.
這個是主禱文其中一部分.
牧師你為什麼停在這裡呢.
當然為了有效去表達.

$^{601}$我停在中間這個位置.
不過我想大家留意一下.
當耶穌在說那些發理賽人.
或者外邦人的禱告.
你不要學的時候.
你們應該要怎樣祈禱呢.
所以你們要這樣說.
我們以為這樣說就要跟足那些字眼.
其實不是那個字眼要跟足.
而是跟耶穌基督所表達裡面.
你關懷的內容.
關懷的內容一開始禱告的時候.
說什麼呢.
是說上帝的名.
願祂的名為聖.
願祂的國降臨.
願祂的旨意行在地上.
如同行在天上.
先求上帝的國.
求上帝的榮耀.
和上帝的名被彰顯.
然後再求自己日用的飲食.
你看到嗎.
在這個祈禱裡面.
那個次序也好.
或者那個關懷內容.
是由天上開始到地上.
我們今天祈禱的是什麼呢.
第一件事.
雖然我們以前在初信班教會有不同的步驟.
不過我們最終是想什麼呢.
我們最終是想說上帝你給我.
我想要的東西.
我們祈求的東西.
經常都求上帝給我們的.
所以我們就有意無意之間.
將重點側重了.
在自己想要的東西當中.
在這裡面耶穌基督.
他告訴你那個foremost.

$^{641}$那個重點.
你祈求你禱告的內容.
你嘗試一下.
不要以你自己的關懷.
作為你禱告主要的內容.
以你的父的事為念.
所以耶穌示範的禱告.
是首先告訴我們.
我們先祈求神的國降臨.
願上帝的名被尊為聖.
願祂的旨意行在地上.
如同行在天上.
這句話不是禱告純粹是技術性的開手.
而是心態上的調教.
相反來說.
如果我們祈求一句.
我們每天五天都吃飯.
今天才給我們.
那就是自我為中心.
內在求自己福氣的祈禱.
這不是錯的.
不過我們如果永遠都帶著這樣的心態.
只是以自己的需要.
來成為我們禱告主要的內容.
有時候我們越祈禱的時候.
我們的心胸會越來越變得狹窄.
我不知道你們有沒有經過.
有沒有見過.
有些弟兄姊妹.
經歷了幾十年的信仰.
他們都是如常不住地來祈禱.
不過很多時候.
越祈禱越苦毒.
越自我為中心.
為什麼呢?.
因為他們永遠都是以自己眼前的事.
眼前自己的需要.
眼前別人得罪他們的情緒.
來作為祈禱的主要中心點.
當我們的事只是放在我們自己.

$^{681}$整個宇宙的中心裡的時候.
我們會發覺這個世界.
其實是一個很殘酷的世界.
很痛苦的世界.
很沒有人情味的世界.
因為所有事情都不是如我心願.
這個世界不是按著我心裡的軌跡來進行.
但當我們拉闊一點.
將我們自己的眼光拉闊了的時候.
我們的禱告.
或者令我們對生命的操練.
可以有一個截然不同的境地.
我認識一位姐妹.
她和她先生結婚幾十年.
她和我很坦白地說.
牧師我祈禱的時候.
你不知道我祈什麼.
因為我和我老公經常吵架.
幾十年.
每次祈禱我都在上帝面前罵我老公.
她覺得上帝幫她出一口氣.
我知道她可能正在面對一些痛苦的狀況.
我都明白的.
不過我鼓勵她.
如果永遠都將你的思維放在這裡.
你的老公依然是這麼刺眼的.
你的生活依然是這麼痛苦的.
嘗試一下改變一下你的角度.
為著這個世界有不同的需要.
來擺放代禱的功夫.
你將你的眼光拉闊一點.
我相信在這個禱告操練你的時候.
你的老公的樣子都會好一點.
你和生命裡面很多糾纏的關係.
慢慢都會言喻意解.
求主幫助我們.
耶穌在這裡提醒我們.
禱告的人慢慢都要將自己的關懷變得廣闊.
以上帝的事為念.
願上帝的名為聖.

$^{721}$願上帝的國降臨.
願祂的旨意行在地上.
如同行在天上.
上帝的說話在這裡也不是叫我們技術性地開頭.
每一次祈禱你要求自己的時候.
就說完這一句.
然後就說自己所需要的東西.
先和後不是指次序的先後.
而是基礎上.
你的眼界,你的心態拉闊了.
以至這個禱告成為了你生命的土壤.
改變你更生.
令你自己不再陷在.
你經常覺得自己以自我為中心的信仰觀念.
我們再看,六章十二至十五節.
這是最後一個段落.
免我們的債,如同我們免了人的債.
不叫我們遇見試探.
救我們脫離凶惡.
因為國度權柄榮耀全善利的.
直到永遠,阿們.
你們饒恕人的過犯.
你們的天父也必饒恕你們的過犯.
你們不饒恕人的過犯.
你們的天父也必不饒恕你們的過犯.
當你讀這一句的時候.
通常我們教禱告.
或者外面有一些的分段.
它的停止到哪裡呢?.
停到因為國度權柄.
全善利的直到永遠,阿們.
我們讀經就讀到這裡停止了.
然後就開新段落.
不過你要看到.
接下來的第十四第十五節.
如果你跳出來之後.
第十六節你們禁食的時候.
好像上下文不是很關係.
無端端說饒恕,然後說禁食.
是不是這兩個是孤句呢?.

$^{761}$原來最新有一些.
很多的《釋經學家》或者你的《聖經》裡面都會看到.
其實說這一句話.
特別是說那些古抄本.
因為國度權柄榮耀全善利的直到永遠,阿們.
這一句很多的《釋經學家》都同意.
這一句是後加下去.
編輯成為一個主土文的內容.
其實因為國度榮耀權柄全善利的直到永遠.
是主要第二世紀才流行說的一句用語.
如果用回古抄本.
如果用回耶穌講論或者登山寶訓裡面所說的.
就是免我們的債如同我們免了人的債.
不叫我們遇見試探.
救我們脫離凶惡.
然後就直到第十四節.
你們饒恕人的過犯.
你們的天父也必饒恕你們的過犯.
你們不饒恕人的過犯.
你們的天父也必不饒恕你們的過犯.
你會看到如果沒有這一句的時候.
你就會解得通了.
一頭一尾其實是在說一個免債和饒恕的主題.
中間夾住的字眼就是.
不叫我們遇見試探.
救我們脫離凶惡.
那個試探和凶惡代表著什麼呢.
原來是回應前面免人的債.
如同我們免人的債.
和後面對人的饒恕.
如果整個句式其實是在說.
耶穌基督教我們禱告的時候.
可能當我們閉上眼睛.
我們想起有些人我們得罪了你.
虧欠你.
令到你很生氣.
禱告的人慢慢要學習怎樣去饒恕.
不容易的功課.
都很違逆人的本性.
我們的本性很想舊時說的施行公義.

$^{801}$正如之前Raymond所說.
我們舊約裡面以眼還眼.
以惡還惡.
是公道的.
不過禱告的人最終要學習的是.
要學習免債.
免人的債.
饒恕和免債是一個對等.
是一個令我們不進入惡者.
或者令自己進入罪惡當中的.
一個很重要的自我覺醒.
當我覺得我們自己不願意饒恕.
我們心裡面帶著憤恨.
來求上帝報應曾經得罪我們的人.
我們很容易走火入魔.
即使我禱告也好.
我們會進入心靈裡面填滿了苦毒的人.
當我們心靈充滿了苦毒的時候.
我們的祈禱大概都是將一些沒有意義的憤恨.
重覆又重覆.
日復如日 年復如年.
即使我不斷禱告.
那種痛苦.
那種仇恨在內心裡面會加增.
耶穌基督提醒我們.
當我們面對著這些人得罪了你的時候.
我們要自我覺醒.
每次祈禱的時候.
想起求天父免我們的債.
那種免債好像一個條件.
上帝你看看我怎樣免別人的債.
按著我所做的來免我們的債.
當然上帝的赦免當然比我們多很多.
不過當我們用這個方式來覺醒來祈禱的時候.
我們避免自己落入了一種苦毒的狀態.
我有過這樣的經歷.
這是我們整個家庭裡面.
曾經經歷過一個很痛苦的經歷.
話說我16歲我爸爸去世.
家裡是做生意的.

$^{841}$在我家裡做生意裡面.
有一些合作的拍檔孤兒寡婦.
但是他們又很不濟.
沒有想過我們孤兒寡婦.
然後騙了我們很多錢.
我媽媽曾經在生前跟我說過.
有一段時間很艱辛.
打算結帳的那些不是真正要結帳的帳.
是被騙的一筆錢.
而這些錢是給我們讀書的.
給我們三兄妹.
很生氣.
找不到這個人.
這個人失蹤了.
到後來我們要賣樓.
或者要做一些事情.
可以周轉到給我們供書教學.
繼續學業.
我媽媽其實一直都在說這件事.
直到我媽媽去世之前.
我曾經問過她.
這件事你怎樣看呢?.
我媽媽相信了之後.
她跟我說.
她說很簡單.
算了.
她跟我說就算了.
我說算的意思是不是一種無奈?.
她說不是.
信主就要教我們去算.
不要再.
她騙了的.
有一天上帝會有公道.
今天你們可以很生性地長大成人.
他們騙了你.
但騙不走恩典的.
上帝看到他們所做的.
也上帝看到我們這個家是怎樣熬下去的.
我想我們學祈禱學遙書.
有時候我們都要學習.

$^{881}$算了.
生命裡面有些事情是痛苦的.
是不容易的.
當我們祈禱的時候.
我們可以選擇繼續就在上帝面前.
就在你所憎恨的人.
不過當我們知道不停地這樣做.
我們很容易墮入那種試探.
令我們成為那惡者不停控訴的對象.
我會令自己停止.
在眼前有很多恩恩怨怨.
我們可以祈求上帝幫我們.
說一句算了.
慢慢由一個意念開始.
逐步逐步將這些重擔一路一路放走.
讓我們真正可以饒恕.
這些曾經得罪我們的人.
求主幫助我們.
如果我們真的要學禱告.
最後要學的就是「原諒」.
耶穌在十字架上的祈禱也是這樣.
「父啊,赦免他們.
因為他們所作的,他們不知道」.
同樣在我們生命裡.
我們也學習了基督耶穌的祈禱.
生命裡充滿了張力.
但當我們成為一個禱告的人的時候.
我們要學的就是學懂.
饒恕身邊曾經得罪你的人.
今天教禱告的內容不是很複雜.
記住這四個F.
「忘記」.
不要用禱告來測試你自己的存在感.
禱告不是你的平台.
表演你屬靈的肌肉.
反而是一個很純粹向上帝表達.
真心吐涼的一個機會.
越將自己低調隱藏就越理想.
讓自己成為一個默默祝福別人的人.
成為一個默默的守望者.

$^{921}$第二個F.
「信任」.
我們每一個禱告的人都知道.
上帝已經將最好的給我們.
眼前有些事情我們所求.
可能上帝未即時給我們.
但上帝不是作弄我們.
上帝總會將最好的給他的兒女.
我們就帶著這份信心去到上帝面前.
甚至有時候我們未祈求以先.
他已經知道我們所需.
我們信得過上帝.
相信上帝會為我們預備最好的.
第三.
我們祈求不要只是純粹以自己的念.
成為我們禱告的主要內容.
拉闊你的心胸.
不要將自己成為宇宙的焦點.
以上帝的旨意.
上帝的名被高舉.
上帝的心意如何實施在我們生命當中.
拉闊你.
以上帝的事為念.
最重要.
以上帝的事.
作為我們優先和基礎來禱告.
拉闊我們.
最後.
我們生命中總會有人得罪你.
又或者你得罪人.
學習.
饒恕.
越禱告.
越能夠學習饒恕.
那就是這個世代裡面.
禱告的人所需要學習的功課.
求主繼續幫助我們.
Check GPT.
看不到我們這些作門徒禱告的心法.
但我們今天跟著這四個F.

$^{961}$跟著耶穌基督的教導.
我們要和上帝的距離會越來越近.
上帝會聆聽你的祈禱.
上帝會越立聽我們的心聲.
求主繼續幫助我們.
\newpage



\section{約翰福音 21:1-17}
\label{sec:HNqZXsZMfMU}
\textbf{2023年4月9日復活節崇拜直播｜陳冠成牧師｜復活主的顯現｜約翰福音21\_1-17}
\newline
\newline
連結: \href{https://youtube.com/watch?v=HNqZXsZMfMU}{\texttt{https://youtube.com/watch?v=HNqZXsZMfMU}} ~~~~ 語音日期: 2023-04-09
\newline
\newline
\hyperref[sec:F3PyAVrDKHQ]{\small{< < < PREV SERMON < < <}}
~
\hyperref[sec:index_chronic]{\small{[返順時目]}}
~
\hyperref[sec:Dcz9krAW4uI]{\small{> > > NEXT SERMON > > >}}
\newline
\newline
約翰福音 21:1-17
\newline
\begin{longtable}{cl}
\hline
\hline
章節 & 經文 (和合本修訂版)\\
\hline
21:1 & \begin{tabularx}{0.7\textwidth}{X} 這些事以後，耶穌在提比哩亞海邊又向門徒顯現。他怎樣顯現記在下面。 \end{tabularx} \\ \\ \relax
21:2 & \begin{tabularx}{0.7\textwidth}{X} 西門‧彼得、叫低土馬的多馬、加利利的迦拿人拿但業、西庇太的兩個兒子，和另外兩個門徒，都在一起。 \end{tabularx} \\ \\ \relax
21:3 & \begin{tabularx}{0.7\textwidth}{X} 西門‧彼得對他們說：「我打魚去。」他們對他說：「我們也和你一起去。」他們就出去，上了船；那一夜並沒有打著甚麼。 \end{tabularx} \\ \\ \relax
21:4 & \begin{tabularx}{0.7\textwidth}{X} 天剛亮的時候，耶穌站在岸上，門徒卻不知道他是耶穌。 \end{tabularx} \\ \\ \relax
21:5 & \begin{tabularx}{0.7\textwidth}{X} 耶穌就對他們說：「孩子們！你們有吃的沒有？」他們回答他：「沒有。」 \end{tabularx} \\ \\ \relax
21:6 & \begin{tabularx}{0.7\textwidth}{X} 耶穌對他們說：「你們把網撒在船的右邊，就會得到。」於是他們撒下網去，竟拉不上來了，因為魚很多。 \end{tabularx} \\ \\ \relax
21:7 & \begin{tabularx}{0.7\textwidth}{X} 耶穌所愛的那門徒對彼得說：「是主！」那時西門‧彼得赤著身子，一聽見是主，就束上外衣，跳進海裡。 \end{tabularx} \\ \\ \relax
21:8 & \begin{tabularx}{0.7\textwidth}{X} 其餘的門徒因離岸不遠，約有二百肘，就坐著小船把那網魚拉過來。 \end{tabularx} \\ \\ \relax
21:9 & \begin{tabularx}{0.7\textwidth}{X} 他們上了岸，看見那裡有炭火，上面有魚和餅。 \end{tabularx} \\ \\ \relax
21:10 & \begin{tabularx}{0.7\textwidth}{X} 耶穌對他們說：「把剛才打的魚拿幾條來。」 \end{tabularx} \\ \\ \relax
21:11 & \begin{tabularx}{0.7\textwidth}{X} 西門‧彼得就上船，把網拉到岸上，網裡滿了大魚，共一百五十三條；雖然魚這樣多，網卻沒有破。 \end{tabularx} \\ \\ \relax
21:12 & \begin{tabularx}{0.7\textwidth}{X} 耶穌對他們說：「你們來吃早飯。」門徒中沒有一個敢問他：「你是誰？」因為他們知道他是主。 \end{tabularx} \\ \\ \relax
21:13 & \begin{tabularx}{0.7\textwidth}{X} 耶穌走過來，拿餅給他們，也照樣拿魚給他們。 \end{tabularx} \\ \\ \relax
21:14 & \begin{tabularx}{0.7\textwidth}{X} 耶穌從死人中復活後向門徒顯現，這是第三次。 \end{tabularx} \\ \\ \relax
21:15 & \begin{tabularx}{0.7\textwidth}{X} 他們吃完了早飯，耶穌對西門‧彼得說：「約翰的兒子西門，你愛我比這些更深嗎？」彼得對他說：「主啊，是的，你知道我愛你。」耶穌說：「你餵養我的小羊。」 \end{tabularx} \\ \\ \relax
21:16 & \begin{tabularx}{0.7\textwidth}{X} 耶穌第二次又對他說：「約翰的兒子西門，你愛我嗎？」彼得對他說：「主啊，是的，你知道我愛你。」耶穌說：「你牧養我的羊。」 \end{tabularx} \\ \\ \relax
21:17 & \begin{tabularx}{0.7\textwidth}{X} 耶穌第三次對他說：「約翰的兒子西門，你愛我嗎？」彼得因為耶穌第三次對他說「你愛我嗎」，就憂愁，對耶穌說：「主啊，你無所不知，你知道我愛你。」耶穌說：「你餵養我的羊。 \end{tabularx} \\ \\
[1ex]
\hline
\hline
\end{longtable}
$^{1}$多謝師班的衷心獻情.
我們繼續以聆聽聖道去敬拜我們的主.
首先由關王桂民姐妹為我們讀出約翰福音第21章1-17節.
然後陳冠誠牧師會告訴我們復活主的顯現.
請桂民.
《約翰福音》第21章1-17節.
「這事以後,耶穌在提比利亞海邊又向門徒顯現,他怎樣顯現,記在下面..
有西門彼得(和稱為底土馬的多馬)並加利利的加拿人(拿旦葉).
還有西比太的兩個兒子,又有兩個門徒都在一處..
西門彼得對他們說:「我打魚去.」他們說:「我有和你同去.」.
他們就出去上了船,那一夜並沒有打著什麼..
天亮的時候,耶穌站在岸上,門徒卻不知道這是耶穌..
耶穌就對他們說:「小子,你們有吃的沒有?」他們回答說:「沒有.」.
耶穌說:「你們把網紮在船的右邊就必得著.」.
他們便紮下網去,竟拉不上來了,因為魚甚多..
耶穌所愛的那門徒對彼得說:「是主.」.
那時西門彼得赤著身子,一聽見是主,就束上一件外衣,跳在海裡..
其餘的門徒離岸不遠約有二百周,就在小船上把那網魚拉過來..
他們上了岸,就看見那裡有炭火,上面有魚又有餅..
耶穌對他們說:「把剛才打的魚拿幾條來.」.
西門彼得就去,把網拉到岸上,那網滿了大魚,共一百五十三條..
魚雖這樣多,網卻沒有破..
耶穌說:「你們來吃早飯.」.
門徒中沒有一個敢問祂:「你是誰?」.
因為知道是主,耶穌就來拿餅和魚給他們..
耶穌從死裡復活以後,向門徒顯現這是第三次..
他們吃完了早飯,耶穌對西門彼得說:.
「約翰的兒子西門,你愛我比這些更深嗎?」彼得說:「主啊,是的,你知道我愛你.」.
耶穌對他說:「你餵養我的小羊.」.
耶穌第二次又對他說:「約翰的兒子西門,你愛我嗎?」彼得說:「主啊,是的,你知道我愛你.」.
耶穌說:「你牧養我的羊.」.
第三次對他說:「約翰的兒子西門,你愛我嗎?」彼得因為耶穌第三次對他說「你愛我嗎?」就憂愁..
對耶穌說:「主啊,你是無所不知,你知道我愛你.」.
耶穌說:「你餵養我的羊.」.
請姐妹平安.
剛才所讀的《讓福音》第21章.
記載了耶穌從死裡復活後第三次向門徒顯現.
這個段落由兩個大部分組成.
首先是耶穌帶領他的門徒如何獲取豐厚的餘禍.
然後記錄了耶穌基督和彼得的單獨對話.

$^{41}$這兩個部分其實離開世體有一個相同的主題.
就是在說復活的耶穌如何給予那些經歷挫敗的門徒再來一次的機會.
我們首先看一至十四節.
最小的場景是提比利亞海,也就是我們熟悉的加利利海.
耶穌基督曾經在他復活後.
吩咐那些見證他從死裡復活的婦女.
要回去告訴門徒.
叫他們回到加利利等候耶穌.
不過耶穌沒有說清楚到底要等多久.
所以在這個情況下.
當門徒一日復一日等候.
他們總要解決一個生活中最基本的問題.
就是開飯.
總要去解決肚腹的需要.
所以漁夫出身的彼得自然會去打魚.
這是他一向的專長和技能.
可以用來解決當下食物的需要.
於是他就說我要去打魚.
而其餘的六個門徒也一起去.
經文記載他們打了一整晚魚.
你不難想像那一晚他們一次又一次地將網殺在海上.
不過經文說他們直到天快亮的時候.
仍然一無所獲.
而在這個時候鏡頭一轉.
耶穌就在岸邊顯現.
他對門徒說將網殺在船的右邊.
就會有所得著.
耶穌的意思是相信我.
聽我的吩咐再試一次.
於是門徒就聽從耶穌的吩咐.
結果就像經文所說.
他們大收漁獲.
熟悉福音書的讀者會在此發現此情此景.
和路加福音第五章.
彼得夢召的經歷很相似.
不過你看到約翰福音的作者.
在此沒有意義再說彼得如何夢召.
因為在這個段落.
他自己是一頭一尾.
第一節和十四節.

$^{81}$他很講清楚目的是什麼.
他要描述記載耶穌基督復活後的顯現.
目的是要我們知道.
到底這位復活的主.
他今次對門徒做了什麼.
從而顯明他又是一位怎樣的主.
事實上聖經學者留意到.
一至十四節的內容.
其實反而和另一段經文有很密切的關係.
就是同樣在約翰福音第六章.
五餅二魚的神蹟.
因為留意到.
二十一章所出現的.
其實有很多和食物和飲食吃相關的字眼.
在程序表講到大都印了給大家.
包括耶穌問門徒有沒有打到魚.
又或者他專程在岸上.
為門徒預備了魚和餅.
然後他又吩咐門徒.
把打到的魚拿幾條當早飯.
之後他又邀請門徒一齊吃早飯.
最特別的是第十三節.
耶穌親自把魚和餅遞給門徒吃.
此情此景很明顯是想表達.
耶穌基督在第三次復活顯現的時候.
他首先要解決處理的.
是門徒一個生活裡最基本的需要.
從而顯明他就是生命的糧.
所以你看到.
即使西門彼得和西比太的兩個兒子.
其實他們都是經驗老道的漁夫.
不過他們努力打魚了一整晚.
仍然是毫無收穫.
他們在這個情況唯一能夠做的是什麼?.
就是聽耶穌的吩咐.
配合耶穌為他們預備的時機.
同樣讓我們看到.
人有時會以為我們只要從事自己熟練的事情.
再加上我們的努力.
就一定會有果效.

$^{121}$多多少少都會解決眼前的危機困難.
不過這段經文剛好告訴我們.
不一定是這樣.
譬如可能有人擁有不錯的學歷.
他做事也很能幹,勤快.
但可能因為突如其來的世事變遷.
譬如可能在早幾年疫情下.
公司沒有生意,結果怎樣?.
裁員甚至倒閉.
又或者可能突如其來.
那個很愛惜他,很愛他的上司離開了.
換了一個新的老闆.
我們都會理解一朝天子是怎樣?.
一條神.
不一定是你不能幹.
不一定是你辦事能力不好.
不過新人士就有新作風.
結果他不再用你.
甚至將你辭退.
你突然間失去了戶口甚至養家的機會.
多努力都沒有用.
又或者我們都試過.
或者你身邊的人都出現過.
你的工作表現非常好.
於是得到公司的賞識.
升你職,加你人工.
當然與此同時就會加你甚麼?.
加添你的工作量.
你開頭覺得雖然有少少吃力.
不過你跟自己說,沒問題的.
能者就多勞.
我撐得住.
只要我加班就可以解決了.
結果通常是怎樣?.
越加班就會越多工作.
不知為何,很特別.
越加班反而你解決不了手頭上的工作量.
反而不知為何會更多工作量.
於是公司見到你應付得到.
就怎樣?.

$^{161}$再加多些工作給你.
結果從此你每晚可能沒有晚上十點.
你都走不了.
你都回不了家.
原本以為能者多勞.
靠自己的努力,勤奮,經驗,才幹.
就可以解決.
結果卻是織奴成質.
弄壞了自己的身體.
沒錯,其實上帝賜予我們的才幹.
給我們的經驗,甚至是努力.
是有它的價值的.
是有益的.
但聖經亦告訴我們.
原來生命裡有很多事情.
其實我們沒辦法靠自己的努力去解決.
就連生命裡最基本的事情.
食物,肚腹的需要.
原來不在你和我的掌握裡.
更何況生命裡更重要的.
與永生有關,與事奉上帝的事情.
經文很想我們這樣去想一想.
所以弟兄姊妹.
其實耶穌向門徒的提問.
同樣向我們提問.
很值得我們去深思.
當耶穌對著這班打魚的門徒說.
小子們,你們有沒有吃的沒有?.
其實耶穌早已知道門徒的答案是甚麼.
是沒有.
換句話來說,其實耶穌有帶諷刺的口吻.
來說,小子們.
是否忙了一整晚,都只有空腹呢?.
事實上,我們在事奉上帝的路途上.
多多少少都遇過類似的困局.
可能你花了很多時間.
用盡你的全力.
熬了很多晚夜.
結果你發覺事情都沒有進展.
沒有果效.

$^{201}$甚至你想服侍動員的對象.
好像對你所付出的沒有反應.
你會發覺,單靠自己的才幹,經驗和勤奮.
我們遲早會面對門徒所面對的困惑.
我們同樣會被耶穌提問.
小子們,是否覺得只有個做字呢?.
是否覺得靠自己的努力解決困難.
結果發覺不是走得通呢?.
耶穌不是要小看我們.
也不是責怪我們.
祂是想我們回想一個很基本的事情.
到底我們是依靠誰的供應和帶領.
是耶穌基督本人.
還是我們骨子裡依靠自己的才華.
靠自己的努力.
靠自己的經驗.
當門徒上岸時.
他們發現耶穌已經為他們準備了早飯.
烤好魚和餅.
這些食物不是來自門徒的餘渴.
是耶穌親手包辦的供應.
相比起豐富的餘渴.
其實門徒獲得一個更寶貴的經驗.
就是耶穌會在適當時候親自餵養和服侍他們.
不需要他們開口.
不需要他們祈求.
這段經文告訴我們一個原則.
當上帝把使命託付給每一個門徒.
包括我們之前.
其實耶穌已經預先為我們預備.
供應了一切.
所以只要我們每一個人.
都願意時常來到耶穌面前.
來跟祂結連.
向祂支取力量.
其實這位大牧人.
一定會在適當的時候.
提供所需給我們.
祂會餵養我們.
祂會幫助我們.

$^{241}$讓我們首先經歷到從神而來的滿足.
然後我們也有力量.
有足夠的供應.
來餵養服侍我們身邊的人.
就是上帝託付給我們照顧的人.
所以這段經文首先提醒我們.
都可以像那些打魚的門徒.
你來到復活主的面前.
你坦然承認.
原來生命不是由我們掌握.
即使最小的事情.
當我們以為可以靠經驗.
靠努力搞定的時候.
其實不是.
是靠主.
甚或你覺得.
主啊,其實我在工作裡.
在事奉裡,或者同人的關係裡.
其實我都一般般.
我很努力,我嘗試去搞好它.
但我發現我靠自己搞不定.
耶穌,你可不可以幫我?.
你可不可以再給我一次機會?.
你讓我可以修正改變.
耶穌,我知道你的時機.
你的預備才是最好的.
我願意去聽從你的吩咐.
順應你的大靈.
投主你給我再活一次.
並且這一次.
我很想為耶穌你而活.
你嘗試向耶穌祈求.
當然謹記的是.
這段經文沒有叫我們.
可以放棄你所有的才華.
你所有的經驗,技能.
甚至乎努力.
什麼都不需要做.
等耶穌搞定,等祂出手.
並不是這個意思.

$^{281}$其實你看到.
沒錯,耶穌是真的有能力.
祂包辦所有的東西.
真的不需要我們出手.
祂可以一手供應,牧養所有的人.
但是,耶穌卻很樂意.
邀請你和我參與當中.
你留意耶穌叫門徒.
將剛才打到的那幾條魚.
拿幾條來.
並不是耶穌預備得不足夠.
所以要靠門徒協助.
拿幾條魚來加sung .
不是這個意思.
而是要藉著我們微小的參與和擺上.
耶穌要建立我們的生命.
耶穌祂很願意和我們同工.
祂很期望我們和祂的使命.
和祂的工作.
和祂的生命有份.
這是經文首先提醒我們.
而繼續看下去.
你會發現.
當處理了門徒討福的需要.
接著就要處理另一件事.
耶穌開始處理.
祂和彼得之間的關係.
如果細心看,你會發現.
自從彼得三次不認主之後.
這次是耶穌首次和他進行的單獨對話.
耶穌問彼得:你愛我嗎?.
對彼得來說.
一定是百般滋味在心頭.
因為你留意經文刻意提到.
岸上有炭火.
你想像一下.
當彼得看到這些炭火和耶穌.
他想起什麼?.
想起什麼?.
對,想起大祭司的原子.

$^{321}$當耶穌受難受審的時候.
彼得看著這些炭火.
三次不認主.
所以你不難想像.
這一餐的早飯.
對彼得來說一點都不自在.
吃得很辛苦.
一邊吃一邊想.
糟了,耶穌會不會和我算帳?.
會不會質疑我.
為什麼當日要這樣對他?.
我當日不是曾經在耶穌.
在眾人面前高言大志.
說要愛他,跟隨他到底嗎?.
為什麼當禍患臨到的時候.
我又退縮.
我又不認他,不愛他.
甚至遠離他.
耶穌會怎麼看我?.
正所謂一次不忠,百次不容.
耶穌還會不會寬恕我?.
他還會不會相信我?.
弟兄姊妹,其實我們不難代入.
西門彼得的心情.
假如我們每一個人.
都有一個履歷表.
上面寫著你愛不愛主.
我們每個人都有一個履歷表.
記錄著你愛主或不愛主的事件.
當拿著這份履歷表.
在耶穌面前,耶穌問.
你愛不愛他的時候.
其實我們可能會像彼得那樣.
在耶穌面前詞窮.
不知道該怎麼說.
甚至像彼得那樣.
被耶穌問得憂愁起來.
沒錯,我們確實有不少.
愛主的證據.
只不過不愛主的紀錄.

$^{361}$也不少.
甚至更加多.
所以我們也會像西門彼得那樣.
主啊,經歷了這麼多風風雨雨.
我不再敢在你面前誇口.
因為我真的拿不到什麼實際的憑據.
來證明我還是愛你.
所以,耶穌你是無所不知.
你知道我的心現在是怎樣.
都是由你來評價.
我是不是愛你.
弟兄姊妹.
耶穌當然知道彼得的心是怎樣想.
其實他也沒有懷疑過彼得對他的愛.
否則你不難想像的就是.
無論彼得回答多少次.
有沒有用.
當耶穌還在懷疑彼得的時候.
無論彼得怎樣回答.
回答多少次.
其實耶穌也不會收貨.
正如當我們懷疑不信一個人的時候.
無論對方指天篤地發誓.
說要在你面前自盡也好.
你也不會收貨.
所以你會看到其實.
耶穌之所以三次問彼得你愛不愛我.
不是為了滿足耶穌的需要.
更不是因為耶穌.
如坊間的人說有本小氣堡.
把對方對不起的紀錄都記下.
不是因為耶穌小氣.
所以故意要為難彼得.
當日三次不認我.
所以我要為難你三次.
問你三次.
問到你口啞為止.
不是這回事.
問三次完全是為了彼得的益處.
耶穌要徹底修復他和彼得的關係.

$^{401}$弟兄姊妹你試想像.
面對一個曾經傷害過你.
離棄過你的人.
如果你仍然願意來到他面前.
主動跟他說.
你還愛不愛我.
你還想不想在一起的時候.
其實背後代表什麼.
你想一想.
假如你仍然很憤怒他.
很憎恨他.
很懷疑他.
甚至很怕他再傷害你的時候.
大概你不會想跟他有什麼瓜葛.
更加不會主動去開口.
問對方還愛不愛你.
所以你看到耶穌之所以問彼得.
你愛不愛我.
其實是因為耶穌對彼得的愛.
從來沒有改變.
他仍然很重視彼此之間的關係.
所以耶穌其實是主動.
首先向彼得示愛.
意思就是彼得.
沒錯,在那三次你雖然沒有回應我.
你雖然不愛我.
但是我從來沒有不回應你.
我從來沒有不愛你,彼得.
到現在都是一樣.
我不單單只要恕你.
接納你一次,兩次,三次.
即或在將來事奉的日子.
假如你仍然有軟弱跌倒的可能.
甚至好像那一晚不回應我的話.
你放心,我不會捨棄你.
任何時候你經歷挫敗,氣餒.
我都會主動愛你,挽回你.
我仍然要對你說.
彼得,我愛你.
我願意給你機會愛我.

$^{441}$將天國的使命繼續託付給你.
所以,西門彼得.
你願意愛我嗎?.
你願意和我同行.
牧養我交託給你的羊嗎?.
弟兄姊妹.
這才是耶穌愛彼得的心腸.
話說.
IBM,大家都認識我這間公司.
大部分電腦都是用IBM的電腦.
當時有一位年青有潛質的主管.
因為參與一項很高風險的投資項目.
而遭遇公司損失一千萬美金.
當時的創辦人.
領導了IBM長達四十年的商界傳奇.
叫做湯瑪士華生.
他要召見這個人.
如果你是老闆,你會怎樣處理這件事?.
他不是幫你賺了一千萬.
他幫你虧了一千萬.
你會怎樣處理?.
大家沒有反應.
可能會辭職.
總要找人來祭旗.
所以這位年青人.
當時百感交集.
他主動和湯瑪士華生說.
我想你叫我來.
都是要散我一餐.
都是要我辭職.
所以我決定自動請辭.
這封是我的辭職信.
湯瑪士華生回答.
開玩笑嗎?.
我剛剛在你身上投資了一千萬美金的學費.
怎會這麼容易讓你走?.
弟兄姊妹,你可以想像.
類似的對話.
很可能都出現在耶穌和西門比德的過程中.
他們的互動中.

$^{481}$譬如當西門比德說.
主啊,我曾經在風浪裡踏出船身.
在水面上行走.
我很想走到你的身邊.
但是因為我小信.
因為我疑惑驚慌.
當我沒有定睛仰望你的時候.
我就下沉到水中.
所以你責怪我.
說我小信.
主啊,我想我還是做不成你的門徒.
我還是辭職吧.
耶穌可能會說.
你開玩笑嗎?.
我剛剛在水上救你回來.
我這麼愛你.
我怎會讓你辭職.
你繼續跟從我.
又或者.
主啊,在你最需要有人禱告警醒守望你的時候.
我睡著了.
在你最需要有人分擔你的憂患的時候.
我不認你.
還不認你三次.
你一定是對我的表現極度失望.
我還是辭職吧.
耶穌於是乎會說.
你開玩笑嗎?.
我剛剛才為你受死.
為你復活.
我用重價把你從罪惡裡拯救回來.
我這麼愛你.
怎捨得讓你離開.
你繼續跟從我.
做我的門徒吧.
弟兄姊妹,其實由我們認順耶穌為主那一刻.
我們就是他的門徒.
有沒有辭職?.
沒有.
可能你可以真的有機會退休.

$^{521}$休息一下.
暫時退下來.
你某個事奉崗位.
但是門徒這個身份是終身制的.
不能請辭.
在訓練的過程裡.
耶穌總是願意寬恕我們的軟弱過範.
總是給我們機會一次又一次.
繼續這樣愛他.
愛他所愛的人.
跟西門比特一樣.
唯有我們先被耶穌基督的愛充滿.
我們才有力量.
來愛神.
愛人.
去展現愛心.
所以.
耶穌在這裡不單要求西門比特.
也要求我們在承接天國使命的從事.
千萬不能忘記失去了那份愛主的初心.
因為.
當我們要承擔更多更大的使命和毀身的時候.
我們就需要更多從上帝而來的愛.
這樣才可以推動我們.
繼續去堅持跟隨耶穌.
確實唯有主的愛.
才可以幫你和我的事奉.
的深度和溫度有所提升.
單靠人的努力.
單靠人的血氣.
我們的事奉慢慢.
會變成例行公事.
甚至連AI都不如.
唯有我們持續去愛耶穌.
持續經歷耶穌的愛.
我們才有力量繼續去愛祂.
愛身邊的人.
所以我想大家繼續讀.
我們有幾節經文未讀.
第十八至二十二節.

$^{561}$請大家一起讀.
預備起.
(唸經).
好,謝謝.
耶穌祂表示.
西門彼得將來會透過受苦.
甚至是殉道的方式.
來榮耀神.
顯明他對耶穌基督的愛.
正如英文.
Passion這個字.
PASSION.
意思是解作.
如果是世界的時候.
世界P的時候.
熱情或者熱愛.
但當這個字變成.
大街passion的時候.
depression的時候.
就是指什麼?.
耶穌基督的受苦.
我們看到.
其實耶穌基督對父神的愛和熱情.
原來和他的受苦.
被綑綁在一起.
分不開.
所以耶穌不是要嚇怕彼得.
而是要每一個愛主的人.
包括我們.
要預先有一個覺悟.
因為在實踐.
愛神,愛人的吩咐上.
其實有可能意味著.
我們會失去了自由.
好像彼得這樣.
被帶到一些.
自己不願意去事奉的場景.
甚至有可能面對那些.
不可愛,反對你.
甚至攻擊你的人.

$^{601}$會令到你沒辦法.
再隨心所欲來行事.
但不是因為我們做錯了什麼.
所以遭受了苦難.
相反.
正因為我們好像彼得這樣.
你選對了.
你做對了.
你願意跟隨耶穌.
好像彼得這樣.
決心愛耶穌的時候.
你就會自然和耶穌基督的.
捨己,受苦.
聯合在一起.
這是必然的事.
所以耶穌基督和西門彼得說.
免禮他跟從我吧.
因為這條不是悲慘的路.
正如耶穌所說.
是榮耀上帝的道路.
耶穌基督要藉著受苦犧牲.
來榮耀神.
同樣他的門徒也要通過.
這個方式來榮耀神.
提醒我們每一個.
不要太在意個人榮辱.
和眼前的得失成敗上.
事奉的路上.
如果有太多比較和計較的話.
其實一定會影響或溝淡了.
我們跟隨主的信心和決心.
正如彼得.
你會看到.
當我們聽到耶穌.
如何預告自己的結局的時候.
並呼召他跟隨主的時候.
你看到彼得的反應是什麼嗎?.
他並不是對耶穌說.
主啊,好,我願意.
你差遣我.

$^{641}$他很人性地.
將他的目光.
放在主所愛的門徒身上.
耶穌問彼得.
主啊,這個人將來會怎樣?.
很明顯他並不是純粹.
出於八卦和好奇.
想知道對方的結局.
而是他呈現了.
一種很想計較和比較的想法.
主啊,你要我為你犧牲捨己.
那他呢?.
你所愛的門徒.
他又需不需要?.
主啊,你要我為你受苦.
沒有問題,我願意.
但是他也應該要這樣.
這樣才公道.
否則,耶穌似乎偏心.
你捨他多於捨我.
又或者,主啊.
如果他不需要受苦.
那我可否用一個比較舒服的方法.
來愛你榮耀神?.
我也想像他這樣.
弟兄姊妹.
不知道我們有沒有閃過這些念頭.
如果有的話.
盼望耶穌基督對彼得的回答.
可以幫到你.
耶穌怎樣回答彼得?.
四個字,是什麼?.
與你何干?.
用廣東話說一次.
啊.
關你什麼事?.
其實是很好的提醒.
真的關你什麼事?.
耶穌其實表示.
每一個門徒榮耀神的方式是獨特的.

$^{681}$是不相同的.
因為是出於上帝的主權和安排.
所以耶穌對彼得說.
這點關你什麼事?.
你又是榮耀上帝,他又是榮耀上帝.
只不過方式不同.
沒有對錯,沒有高低.
我愛你和愛他都是一樣的.
所以真的,弟兄姊妹.
沒有比較就沒有什麼?.
就沒有煩惱.
做人,做基督徒,為什麼經常不開心?.
你想想,其實多數與計較和比較有關係.
耶穌正正提醒我們.
不要拿別人的際遇來做比較.
不要總覺得別人給自己幸運和幸福.
覺得上帝總是欠了我一些東西.
對我不是很好.
又或者懷疑.
耶穌你是不是真的愛我,為我預備最好的?.
結果又如何?.
結果你越事奉,就會越多埋怨.
你身邊的人就會漸漸覺得.
這樣不妥,那樣不順眼.
總是找到一些東西來挑剔.
漸漸你開始抽離這個群體.
甚至與上帝的關係疏遠.
沒錯,雖然仍然是崇拜和事奉.
但沒有了以前那份與上帝親密的關係.
沒有了以前那種喜樂滿足.
為什麼?.
是因為比較和計較.
所以耶穌對西門比德.
亦對我們說.
不要懷疑上帝.
不要比較和計較.
專心只管跟隨耶穌.
祂不會騷擾你.
唯有我們定睛仰望那位死而復活的主.
回想起耶穌,祂也沒有跟我們計較.

$^{721}$祂也沒有向天父投訴.
有沒有搞錯?這麼不公平.
明明我聖潔的.
要我這麼吃虧.
為了一群罪人犧牲釘十架.
耶穌沒有這樣計較.
耶穌反而處處為上帝的益處.
上帝的榮耀.
作為優先的考慮.
所以當我們仰望這位復活的主的時候.
你就不容易迷失.
因為有耶穌作為你的榜樣.
並且你知道祂一定會扶正你.
讓你穩行在一條榮耀上帝.
跟隨上帝的路上.
所以弟兄姊妹.
來到這裡我們一起祈禱.
無論我們面對的光景是什麼.
特別是剛才經文對你有提醒.
可能你就像那些正在打魚的門徒.
有很多事情你很想為上帝做好.
但總是努力也達不到效果.
你求上帝幫你.
又或者你曾經立志.
好像彼得那樣要愛主愛人.
但卻做了一些上帝虧欠了人的事情.
其實你很想修補很想改正.
但不知如何入手的時候.
求主幫你.
又或者你生命裡有很多計較比較的念頭.
影響了你愛主愛人的熱心.
求主幫你.
無論如何.
耶穌他認識我們.
無論如何他知道我們的軟弱.
並且他很願意愛我們.
挽回我們.
他很樂意幫我們重新跟隨他.
所以弟兄姊妹.
你愛這位為你死為你復活的主媽.

$^{761}$我們同心祈禱.
因為你為了我們的罪惡過犯緣故.
你犧牲在十字架上.
你為了我們永生盼望的緣故.
你復活過來.
今天你已經坐在上帝寶座的右邊.
為我們代求.
上帝幫助我們.
正如經文裡所說.
門徒的軟弱都是我們的軟弱.
他們的盲點都是我們的盲點.
我們很多時候.
想為上帝你成就美事.
卻是有心無力.
又或者我們盡力了.
好像徒勞無功.
耶穌讓我們重新歸向你.
讓你再給我們機會.
讓我們敏銳你的時機.
順服你的大靈.
以致我們靠著主你.
可以結到果子.
經歷到服侍你的美好和甘甜.
主也或者我們過往曾經很虧欠你.
我們雖然口裡愛你.
但我們的心卻遠離你.
甚至我們的行事為人相反.
得罪你.
主也求你愛我們.
我們知道你願意幫助我們.
因為你已經用重價將我們贖回.
你要繼續使用我們.
讓我們一生跟隨你.
求主你再給我們機會.
讓我們可以修正改過.
讓我們可以順服你.
跟隨你.
過一個榮神益人的生活.
主也給我們充滿你的愛.
願你以你的愛懷抱我們.

$^{801}$當我們看見主對我們的愛.
是何等深何等大的時候.
我們便已經滿足.
我們便不再計較.
不再和別人比較.
我們的際遇.
我們的眼前得失.
無論如何.
我們知道上帝很愛我們每一個.
你對我們一定有美善的安排.
主呀.
讓我們每一個重拾對你的熱情.
讓我們每一個回到跟隨你的初心.
讓我們就好像你所說.
只管跟隨你.
榮耀你.
求上帝你閱立我們在你面前的禱告.
奉復活的耶穌基督.
得勝的名字來祈求.
Amen.
\newpage



\section{馬太福音 6:16-23}
\label{sec:Dcz9krAW4uI}
\textbf{2023年4月16日崇拜直播｜羅志雄牧師｜給當代門徒的15個忠告:禁食為何｜馬太福音6\_16-23}
\newline
\newline
連結: \href{https://youtube.com/watch?v=Dcz9krAW4uI}{\texttt{https://youtube.com/watch?v=Dcz9krAW4uI}} ~~~~ 語音日期: 2023-04-17
\newline
\newline
\hyperref[sec:HNqZXsZMfMU]{\small{< < < PREV SERMON < < <}}
~
\hyperref[sec:index_chronic]{\small{[返順時目]}}
~
\hyperref[sec:zKeckdMvwPY]{\small{> > > NEXT SERMON > > >}}
\newline
\newline
馬太福音 6:16-23
\newline
\begin{longtable}{cl}
\hline
\hline
章節 & 經文 (和合本修訂版)\\
\hline
6:16 & \begin{tabularx}{0.7\textwidth}{X} 「你們禁食的時候，不可像那假冒為善的人，臉上帶著愁容；因為他們蓬頭垢面，故意讓人看出他們在禁食。我實在告訴你們，他們已經得了他們的賞賜。 \end{tabularx} \\ \\ \relax
6:17 & \begin{tabularx}{0.7\textwidth}{X} 你禁食的時候，要梳頭洗臉， \end{tabularx} \\ \\ \relax
6:18 & \begin{tabularx}{0.7\textwidth}{X} 不要讓人看出你在禁食，只讓你隱祕中的父看見；你父在隱祕中察看，必然賞賜你。」 \end{tabularx} \\ \\ \relax
6:19 & \begin{tabularx}{0.7\textwidth}{X} 「不要為自己在地上積蓄財寶；地上有蟲子咬，能鏽壞，也有賊挖洞來偷。 \end{tabularx} \\ \\ \relax
6:20 & \begin{tabularx}{0.7\textwidth}{X} 要在天上積蓄財寶；天上沒有蟲子咬，不會鏽壞，也沒有賊挖洞來偷。 \end{tabularx} \\ \\ \relax
6:21 & \begin{tabularx}{0.7\textwidth}{X} 因為你的財寶在哪裡，你的心也在哪裡。」 \end{tabularx} \\ \\ \relax
6:22 & \begin{tabularx}{0.7\textwidth}{X} 「眼睛是身體的燈。你的眼睛若明亮，全身就光明； \end{tabularx} \\ \\ \relax
6:23 & \begin{tabularx}{0.7\textwidth}{X} 你的眼睛若昏花，全身就黑暗。你裡面的光若黑暗了，那黑暗是何等大呢！」 \end{tabularx} \\ \\
[1ex]
\hline
\hline
\end{longtable}
$^{1}$以下我們要聆聽神給我們的說話.
我們要聽神的說話.
記載在馬太福音第六章十六至二十三節.
新約聖經的八頁.
馬太福音第六章十六節.
你們禁食的時候.
不可仗那假冒為善的人.
點上呆雀壽容.
因為他們把臉弄得難看.
故意叫人看出他們是禁食.
我實在高騷你們.
他們已經得了他們的賞賜.
你禁食的時候.
要梳頭洗臉.
不叫人看出你禁食來.
只叫暗中的乎看見.
你乎在暗中察看.
必然報答你.
第十九節.
不要為自己積贊財寶在地上.
地上有蟲子咬.
能受壞也有賊挖洞來偷.
只要積贊財寶在天上.
天上沒有蟲子咬.
不能受壞也沒有賊挖洞來偷.
因為你的財寶在那裡.
你的心也在那裡.
聖經獨筆.
弟兄姊妹主內平安.
有請梁美芬來做禱告.
天上的包袱.
我們獻上感恩.
獻上讚美.
是上帝你自己的恩典.
招聚眾弟兄姊妹一齊去敬拜.
去聆聽上帝你自己的教導.
主幫助我們今天.
能夠專心聽你自己的教導.
無論講的聽的.
都同蒙上帝你自己的賜福.

$^{41}$幫助我們在生活當中.
去見證你.
去榮耀你的名.
我們來公敬的.
向你獻上禱告.
奉主耶穌基督的名.
阿們.
今日的題目是.
給當代門徒的十五個忠告.
禁食為何.
一開始很想說說.
疫情過了.
很多人會去旅行.
我相信都會成為大家吃喝的日子.
時刻又到.
品嘗不同地方的美食.
去很多不同的地方.
我相信已經是非常普遍.
亦都是一件令我們很滿足的事情.
我不知道當你有飲食的時候.
你有什麼心情呢.
會不會第一時間都不是自己先吃.
你會怎樣做.
手機先吃.
我都知道.
手機先吃之後.
還會吸不及待.
第一時間給一些親朋好友.
讓他們在上面看看.
試試一下.
這是無可厚非.
人人都喜歡這樣的美食.
聖經裡面亦都沒有禁止.
或責備我們飲食.
辛苦賺來自在的吃.
我相信大家都很熟悉這句話.
但聖經裡面傳道書提醒我們.
人最好吃喝.
在自己的奴隸裡自得其樂.
我看這也是出於神的手.

$^{81}$的確是.
我們知道.
如果我們更加清楚明白.
吃喝能夠帶給我們無比的歡愉.
同樣是神厚賜.
亦都是恩典的時候.
我們更加應該心存感恩.
不過有很多不同的情況.
或其他原因.
可能不會去飲食.
甚至要停一段時間.
或身體疾病的緣故.
又或近來好成風氣.
剛才在等候的時候.
都聽到說.
十六八節食.
減糖.
還有生酮.
不知道大家有沒有試過.
我第一次聽.
生酮.
他們去時上稱之為.
這些身體的緣故.
都可能叫節食.
或禁食.
但在我們信仰上.
我們都有這種禁食.
這些就是宗教生活上的禁食.
好像是伊斯蘭教.
天主教.
猶太教.
東正教.
甚至我們跟隨耶穌基督.
基督新教.
亦都是在指定的日期裡面.
舉行禁食.
不過弟兄姊妹.
我們一班信徒群體.
特別在這個時候.
我們都很少提及禁食.

$^{121}$反而其他靈性的操練.
我們都會常常去做.
好像是讀經.
祈禱.
又或者一些幫助我們.
可以操練我們屬靈生命的.
我們都會去做.
但偏偏在禁食這個課題上.
我們比較少提到.
很希望藉著今天這段經文.
來幫助我們.
讓我們深入背後的意義.
讓我們能夠增加.
我們在信仰生活裡的活力.
弟兄姊妹.
究竟禁食是甚麼呢?.
是否只是為了身體的緣故呢?.
我們知道禁食是完全不吃東西的.
甚至水都不喝.
甚至可能你會用一段很短的時間.
短的時間可能是以日計算.
甚至是以日計算.
大家又知不知道.
我們經常都有禁食.
禁食英文叫fast.
但我們早餐叫甚麼呢?.
英文叫breakfast.
即是告訴我們.
其實我們每一天早餐.
一開始早餐的時候.
其實我們就解除了禁食.
在舊約聖經裡面.
關於禁食的一些條例.
其實是出現在利未記裡面的.
利未記裡面有這樣的記載.
每逢七月的初十日.
你們要刻苦己心.
無論本地人是寄居的.
在你們中間的外人.
甚麼功都不可作.

$^{161}$只要作為你們永遠的定例.
在這天要為你們贖罪潔淨.
你們要在耶和華面前得以潔淨.
脫盡一切的罪孽.
這裡提到的是刻苦己心.
可能沒有提到禁食.
但原來這個字.
字面是刻苦己心.
但其實是有一個禁食的意思.
禁食是一種受苦.
面對罪的一種受苦.
從而我們能夠謙卑在神的面前.
有一個態度.
當然刻苦己心.
其實是告訴我們.
除了不吃之外.
其實我們還有一些可以禁戒.
這就是利未記裡提醒我們.
另外一段經文亦提到.
指明要在贖罪日裡.
守禁食.
我們都知道.
猶太人面對很多種種的艱難困苦.
面對很多的狀況.
受了很多民族的苦況.
他們更加會犯罪.
得罪神.
常常都會受到神的懲罰.
最後他們都成為國家滅亡.
甚至整個民族都被迫移民到被擄.
在這種情況下.
他們得罪神.
他們又悔改.
在悔改的時候.
其中他們要做的是什麼呢?.
就是禁食.
以色列人在被擄的時期.
他們也制定了很多關於禁食的情況.
他們的禁食原意是為了悔罪.
甚至希望有一天.

$^{201}$能夠在祈禱和禁食的情況下.
上帝憐憫他們.
他們能夠回到自己的故土耶路撒冷.
在被擄回歸的那段時間.
我相信大家都聽過.
從舊約的聖經裡面.
都知道有很多的情況.
猶太人整個民族.
雖然在異地生活.
他們都會禁食.
最明顯的是.
以斯帖皇后.
當時有一個詭計.
要將整個以色列民族.
散居在巴比倫這些國土的時候.
所有的以色列人都要滅亡.
這位皇后知道的時候.
她也要呼籲百姓.
來禁食禱告.
以致她能夠從上帝而來.
進去見皇帝.
在這個過程裡.
最終能夠解決.
滅絕民族的情況.
又或者在尼希米的時候.
當他回到耶路撒冷.
重建城牆的時候.
他也發起了禁食.
今天我們看這一段經文.
六至十八節.
第一段經文的時候.
整段經文.
如果我們有上文下理看的時候.
就是由第一節到第十八節.
在第六章裡面.
當中耶穌有提及三個宗教敬拜生活.
有施捨,禱告.
這兩個分別已經.
余姑娘和陳牧師.
在崇拜的時候.

$^{241}$在他們的信息裡面.
都有很多的教導和提醒.
今天我們嘗試看第三個.
有關猶太人在宗教生活裡面.
的禁食.
大家有留意.
這段經文.
其實與施捨和禱告的句式.
是一模一樣.
我相信耶穌在這裡.
刻意地將它這樣教導的時候.
一定有它的用意.
我們先來看看.
關於耶穌禁食的教導.
在新約裡面.
這是唯一的記載.
在以後就算在保羅的書信.
同樣都有提及禁食.
但沒有很詳細的教導.
雖然如此.
但我們在禁食的傳統裡面.
究竟對我們今天.
每一個跟隨耶穌的人.
在我們的信仰生活裡面.
又或者在我們屬靈生命裡面.
有什麼可以反思的地方呢?.
第十六節一直至到.
經文裡面.
第一節十六節是這樣說的.
你們禁食的時候.
不可像假冒為善的人.
念上戴著愁容.
因為他們把念弄得難看.
故意叫人看出他們是禁食.
在耶穌的時代.
禁食是非常盛行的.
除了是律法所指定要求的.
在贖罪日必須禁食之外.
還有很多其他種種的禁食.
甚至有些是什麼呢?.

$^{281}$是自願選擇禁食.
所以耶穌在這裡看到這樣的現象.
在他的登山寶訓.
來到第六章.
來到這裡的時候.
講了施捨,祈禱之後.
繼續的講出關於禁食的教導.
你們禁食的時候.
不要像那些假冒為善的人.
什麼是假冒為善呢?.
相信聽了施捨和祈禱的時候.
都知道禁食和假冒為善.
其實是相關語言.
當然一定是指那些.
法利賽人,文士這些宗教領袖.
耶穌沒有反對禁食.
禁食這個行為是好的.
但是為什麼耶穌要在這裡講.
不要像那些法利賽人,文士.
宗教領袖這樣的禁食呢?.
耶穌在這裡觀察到.
這個現象是怎樣的.
他們有什麼特色呢?.
他們原來進行禁食的時候.
普遍是怎樣呢?.
是要故意叫人看到.
還要在一些最多人走的地方.
讓人看到他們在這樣的狀態.
在這樣的情況下.
他們是禁食的.
不單止是這樣.
他們會怎樣呢?.
經文告訴我們.
耶穌說他們故意令自己.
弄得很難看.
難看到一個地步.
甚至要將一些灰泥.
抹在臉上.
這個有什麼作用呢?.
面帶受用.

$^{321}$這個說法是形容一些親人死了.
他很悲哀地表達.
而弄得難看的時候.
他抹了很多那些灰泥.
其實就是讓人看不出.
究竟這個是誰來的.
但是在這樣的情況下.
更加顯出他是一個特別聖潔.
一個這樣的情況.
耶穌看到當時禁食.
人人都這樣禁食.
當時的現象就是這樣.
究竟耶穌有什麼評價呢?.
就說他們是假冒為善.
原來他們進行禁食.
就算不是法利賽人.
不是一些宗教領袖也好.
可能他們都會去禁食.
禁食是怎樣呢?.
就成為假冒為善.
耶穌說他們是在演戲.
假冒為善這個字其實是演員.
另一個意思是要做演員.
那他們怎樣做呢?.
就戴面具.
戴什麼面具呢?.
就是弄得自己很難看.
甚至在臉上塗了一些灰泥.
弟兄姊妹.
我想在這裡停一停.
想想究竟我們的信仰生活.
是否為了.
好像假冒為善的人.
只是為了回來教會聚會.
崇拜參加小組.
讓你認識的家人朋友覺得.
你回來教會很忙碌.
你真是虔誠了.
又或者會不會覺得.
你做這些活動.

$^{361}$例如禁食.
做這些動作成為你的主菜呢?.
回來教會.
其實最終想做些什麼呢?.
我相信弟兄姊妹.
藉著耶穌再三提醒我們.
假冒為善.
其實耶穌來到這裡.
我相信看到人真正表演的面目的時候.
耶穌在這裡淋漓盡致.
用這個字.
因為實在有行為可以做出來.
禱告最多是大聲喊.
站在街角.
但原來你化妝的時候.
讓人更加體會到.
所以耶穌就在這裡.
說他們這種做法.
其實就是在演戲.
在演戲地做.
他們有什麼賞賜呢?.
耶穌就說.
耶穌對他們的評論.
就是我實在告訴你們.
他們已經得到他們的賞賜.
這個賞賜.
也是一個商業的用語.
如果大家下了單.
交收了一個物品.
付了錢.
交了貨.
收了貨.
如果大家知道.
以前有些單據.
會蓋上一個印.
叫做課銀兩銀.
就是說沒有拖沒有欠.
從此就不會再有麻煩出現.
耶穌在這裡這樣說.
他們得到賞賜.

$^{401}$就好像做交易一樣.
課銀兩銀.
就是說如果你的禁食.
成為一個宗教行為的時候.
其實你已經收了你應該有的賞賜.
不過不是從上帝而來的.
禁食原來的意義.
告訴我們是謙卑.
悔改這種心智.
但是法利賽人.
他就成為一個表演.
沽名釣譽.
去得人稱讚的途徑.
就算在宗教的節期.
特別是禁食日.
就是贖罪日裡面.
他們要舉行禁食.
可能這些法利賽人.
教會領袖.
甚至是一班信徒.
都跟風的時候.
可能只是一個守禮如儀.
甚至是履行宗教責任.
弟兄姊妹說到宗教責任的時候.
不知道大家如何去理解.
耶穌在這裡提醒我們.
我們究竟跟隨耶穌.
還是跟隨我們所說的耶穌教呢.
昔日如果大家有聽過.
不是叫基督教.
叫耶穌教.
甚至我們去想一想.
究竟我們是跟隨一個宗教.
還是跟隨耶穌.
是耶穌真正的門徒呢.
我們會不會是戴上了一些面具.
回來教會.
出席各種的活動.
甚至有一些宣教.
又或者一些服事的活動.

$^{441}$我們會去探訪.
探訪一些有需要的人.
我們會去探訪一些.
真是他們有需要幫助的.
好像過去的日子.
我們有步兵.
組織了步兵.
來探訪一些.
他們感染了新冠病毒的.
在那一刻有沒有想過.
我只是履行一種責任.
聖經是這樣教.
耶穌是這樣教.
那我們就這樣去做.
還是說.
我們內裡如一.
又或者連我們自己都不知道.
原來我們真正回教會.
那個敬拜.
變成了一個例行公事.
自己騙自己.
甚至連上帝都欺騙了.
我相信這個.
更加讓我去想清楚.
不是的.
我回來教會.
真是實際做了很多事.
最重要的是什麼.
十一奉獻.
最近教會又呼籲.
我們有裝修.
令到禮堂.
甚至我們在五樓.
也有裝修.
是不是就是這樣.
我做足了這些要求.
或者是吩咐.
在這裡讓我們去思想一下.
求主幫助我們.
讓我們在信仰的生活裡.

$^{481}$不是演戲地去做.
究竟什麼是真正的禁食呢.
怎樣做才是禁食呢.
接著耶穌教導門徒.
一種另類的禁食.
不知道大家一聽到.
另類的禁食是什麼呢.
很想知道.
耶穌沒有其他的方法.
其實都是禁食.
但是那個方法很不同.
耶穌就告訴我們.
他說不要叫人見到.
要暗中.
要靜悄悄.
不要被人發現.
就算是日常生活也好.
逛街也好.
在大街大巷都沒有人知道你禁食.
但是只有誰會知道呢.
就是暗中的符.
就是你看不到的那位符上帝.
成為你可以看到你禁食的狀況.
耶穌在這裡教導的門徒.
這一種禁食.
其實是隱含著一個情況.
就是可以體會到.
那些禁食的人.
他們去禁食是怎樣呢.
假冒為善.
假冒為善其實背後.
隱含著一連串的想法.
其中就是他的動機.
他的動機是怎樣呢.
這個耶穌已經講得最清楚.
但是最清楚我們的動機是誰呢.
其實天父是會看到我們人內心的深處.
所以耶穌的教導.
我們的禁食.
只是給我們的符上帝知道.

$^{521}$我們不需要披麻蒙灰.
樣子做得很難看.
但是反而要怎樣呢.
耶穌在這裡教導我們的時候.
就是要梳頭洗臉.
梳頭洗臉是怎樣的呢.
我們每天都做的.
你會不會睡醒就這樣出門呢.
就算很忙碌.
過了一小時睡醒.
原來過了一小時.
那你會不會過了一小時之後.
就這樣穿上衣服.
隨便就出門呢.
也不會的.
就算過了一小時.
多趕時間也好.
你都會怎樣呢.
都是用十分鐘八分鐘.
以前可能會用三十分鐘來做的.
但你都會做一個這樣的梳頭洗臉.
耶穌在這裡提醒我們.
就是說我們一般的梳洗.
好像我們平時這樣的裝束.
這樣的裝扮.
當然一定是整齊清潔.
是我們如常生活這樣.
來做這個禁食.
弟兄姊妹.
耶穌這樣教導我們的時候.
耶穌本身已經身體力行.
這個禁食.
這個另類的禁食.
耶穌在傳道之前.
祂就獨自去到曠野.
去到一些杳無人煙的地方來禁食.
而且在曠野的地方.
祂禁食了四十晝夜.
我相信聖經的記載.
這裡是耶穌最長的禁食時間.

$^{561}$在祂這個禁食的裡面.
祂謙卑自己.
預備祂的心智.
作傳道.
將天國的信息.
宣揚一個的準備.
當祂禁食完畢.
聖經裡面也記載著.
祂還要受靈性和肉體的試探.
很明顯.
耶穌不是一下子進入四十天的禁食.
而是在祂平常的時候.
祂一直都有禁食的操練.
只不過祂不是在人群那裡.
祂往往禁食和禱告.
是同時進行.
遠離人煙.
去到曠野或者山上.
比耶穌更早之前.
進行禁食的.
聖經裡面有一個很特別的記載.
也是在耶穌的時期.
有一位年老的女先知.
叫做阿娜.
她長期禁食.
祈禱和祈求.
晝夜事奉神.
在她年老的日子.
上帝給她有回報.
她得到能夠在聖殿裡面.
看到嬰孩的耶穌.
這就是她長期盼望.
以色列的聖者.
所以這位年老的女先知阿娜.
在她人生最後的階段.
或者在她人生渴求的裡面.
上帝給她回報.
又或者在初期教會.
耶穌的門徒彼得約翰他們.
當他們經歷了聖靈降臨的時候.

$^{601}$他們就大有膽量.
他們去傳道.
當面對種種困難的時候.
他們依然守住.
耶穌給他們的教導.
就是禁食.
也有保羅的禁食.
又或者在早期的教會.
一些清教徒.
他們常常都刻苦己生.
可能在初期的教會.
又或者初代的教會.
通常都很著重這種禁食.
可能一路演變的時候.
一路去到某一個狀況之下.
今天都比較少做這些操練.
我自己也有一個禁食的經歷.
一個學習.
但可能說不上操練.
我想起預備的時候.
我曾經向上帝禱告.
聖經裡不時都說禁食.
究竟是怎樣的呢?.
我都去祈禱.
可不可以讓我學習一下.
究竟禁食是怎樣的呢?.
禁食整個過程是怎樣的呢?.
我就去祈禱.
想想是何時開始.
禁食的方法是怎樣的呢?.
都跟一些導師傾談過.
但當時我只想起禱告.
好像沒有實際行動.
但後來就真的很機緣巧合.
我想都是上帝了解了我的需要.
了解了我這種渴求.
於是就參加了浸信會聯會.
即是我們自己的宗派.
舉辦的禁食觀摩團.
我在那裡祈禱.

$^{641}$想想在什麼地方都在自己的.
會不會在家裡.
或是在什麼地方禁食.
原來這個禁食觀摩團去哪裡呢?.
去韓國.
真的很特別.
當時大家都知道.
為什麼很多都會去韓國呢?.
原來很多不同宗派的弟兄姊妹.
傳道人都會帶他們去韓國的地方禁食.
也可以跟韓國的弟兄姊妹.
有一個信仰的交流.
原來真的有後著.
之後家人患病.
我除了要照顧他們.
還有教會的服侍之外.
我仍然如常在病房工作.
當我在想.
我之前好像去過一個禁食交流團.
我會不會學習一下禁食呢?.
於是我在病房的時間.
放飯的時候.
我就學習禁食.
如何禁食呢?.
走進病房空置出來的房間.
我就不是去放飯.
學習一下不吃這一餐.
是做什麼呢?.
原來可以學習安靜禱告.
甚至可以嘗試.
雖然這一段禁食的經歷不是很長.
但原來讓我可以將病患的家人交託給主.
甚至在這段時間.
都是操練我.
究竟我自己的信仰.
在這個艱難的時刻.
是否只要家人可以康復.
還是怎麼樣呢?.
當我回想有這個經歷的時候.
讓我體會到.

$^{681}$在艱難的時刻.
如何信靠主.
更讓我體會到.
究竟我的心智.
一直跟隨耶穌有沒有改變呢?.
的確是的.
耶穌教導門徒的禁食操練.
是叫我們不要有得望回報的賞詞.
就是別人的讚賞.
你真的虔誠了.
其實耶穌在這裡教導我們.
在沒有人見到.
不顯眼的時刻.
我們願不願意去選擇.
做這樣的操練呢?.
就好像我們日常生活一樣.
淺行出來.
我相信上帝會看到.
如果弟兄姊妹願意做禁食.
其實上帝會察看.
因為上帝會看到你內心的動機.
你是否內外如一.
上帝都會知道.
當我去預備的時候.
我看到禁食不只是一種苦修.
認罪悔改的一個場景.
在撒迦利亞書裡也有提到.
原來在禁食的時候.
可以將它轉變過來.
那裡怎麼提呢?.
耶和華如此說.
四月五月禁食的日子.
七月十月禁食的日子.
必變為由大家歡喜快樂的日子.
和歡樂的節日.
所以你們要喜樂.
喜愛誠實與和平.
弟兄姊妹.
禁食操練.
很多時候都會被認為.

$^{721}$甚至我們今時今日.
都會覺得是一個虔誠.
不單只是我們信仰對上帝的虔誠.
甚至會被人說是宗教虔誠.
你有好明星.
必定你事奉會有好的效果.
甚至可能在一些群體裡.
說我是教會裡的一個領袖.
我都要禁食.
以致到他可能只是為了事奉.
為了自己是一個領袖.
而禁食.
其實這樣的想法.
這樣的想法.
其實在我們的信仰生活裡.
可能把焦點放錯了.
令到我們的動機.
都出現了不純正的地方.
求主在這裡監察我們.
禁食就好像其他一些.
屬靈操練一樣.
其實我們要專注在上帝那裡.
在上帝的身上仰望耶穌.
去尋求祂.
當然禁食不單只是.
想著吃或不吃.
其實有很多方式.
來表示我們.
是否以一種禁食的態度.
來面對.
特別是一些我們最喜愛.
一些要做的事情.
又或者一些我們最喜歡吃的東西.
當你願意放下.
這些你最喜愛.
最愛不釋手的.
然後你集中朝向上帝那裡.
在神那裡.
特別是現在最流行的是什麼.
拿著手機就在那裡玩遊戲機.

$^{761}$甚至是看劇集.
有很多的方式.
來看一些劇集.
不一定是看電視.
明珠台,翡翠台.
大家可能有很多的方法.
會不會透過這些禁戒.
戒除一些自己很喜愛的事情.
來減低對這些物質的誘惑.
對你的誘惑.
從而將你的聚焦點.
回到耶穌的身上.
更加敏銳,聖靈.
帶領你原本不想做.
不願意做的事情.
因為你有禁食的心智.
而可以來做一些.
原本上帝令你有益的事情.
或者討他喜悅的事情.
甚至禁戒一些.
成為你自己沉溺的習慣.
或者成癮的事情.
我不知道這些是什麼.
但當你本著禁食.
這種回歸上帝那裡的時候.
這樣可以改善你的工作.
你的學習.
甚至你的人際關係.
我相信這個.
可以讓你有一個嶄新的.
是一個沒有外表敬虔.
但有內而外謙卑的信仰生活.
相信這個就是耶穌很渴望.
來提醒我們.
這段經文.
好了,耶穌說完.
教導了門徒.
關於禁食這個.
屬靈生活美好的敬虔.
我們不要來把它做壞了.

$^{801}$隨之而來.
耶穌教導我們.
信仰生活可以結合在.
我們真實的場景裡.
這段經文我相信.
大家都很熟悉.
我很想請兄弟姐妹.
一起來讀這段經文.
19至20節.
預備起.
(念經).
好,謝謝.
你的心也在哪裡.
謝謝弟兄姊妹.
我們這種信仰的生活裡.
當然包括了.
施捨,禱告和禁食.
但耶穌所關心的.
不單單是這些.
更加關心我們每一個跟隨耶穌的人.
究竟我們在世上.
在這個社會裡.
怎樣真實地生活.
當然包括了.
我們對金錢,財富,飲食.
和衣著的範疇.
這一部分.
怎樣可以體現到呢.
在經文裡.
登山寶訓第六章.
第19至34節.
反映出來.
耶穌在這裡繼續.
教導和提醒.
跟隨耶穌的人.
但很可惜的是.
很多時候我們就將了.
耶穌教導的時候.
提及地上的生活狀況.
可能有意或無意.

$^{841}$將耶穌的教導.
劃分為屬靈或屬世.
兩個面向.
一位牧師.
史托德牧師提醒我們.
原來我們這樣劃分.
很多時候.
都是誤導了我們.
信耶穌的人.
特別是我們一群.
信耶穌的.
耶穌基督的門徒.
因為史托德牧師說.
這樣的劃分.
是會將聖俗.
這種分割.
帶來很多.
我們難以想像的後果.
當我們今天.
學習經歷的時候.
有時候我們都很模糊.
甚至會將一些情況.
在我們生活當中.
分為屬靈或屬世.
基督徒在做每一件事上.
我們大家都知道.
我們會開車.
我們會煮飯.
在辦公室上班.
又或者一些娛樂.
甚至我們會.
做很多其他事情.
但在這裡.
按照史托德牧師的提醒.
其實這些都是屬靈的.
為什麼呢.
因為我們這些事情.
不單是做在人的面前.
或者做在自己身上.
在這裡告訴我們.

$^{881}$我們其實都是做在神的面前.
我們也是遵照著神的旨意.
教導我們怎樣在地上.
實現實際的生活.
因此在登山寶訓第六章裡.
來到這裡的時候.
其實是概括提醒我們.
我們要注意兩種人生觀.
就是世人的價值觀.
和基督裡面的價值觀.
耶穌正正針對著這一方面.
提醒我們.
讓我們知道.
同樣耶穌是關心我們.
我們怎樣面對地上的種種誘惑.
正如耶穌在傳道之前.
他去了曠野.
四十晝夜.
禁食之後.
所帶來一連串.
接受魔鬼的挑戰.
耶穌就在這裡提醒我們.
我們在地上生活.
會面對很多不同的誘惑.
挑戰.
究竟在這個情況之下.
我們和基督的價值觀.
是否一致呢?.
還是出現了很多張力的時候.
我們怎樣生活呢?.
在這裡耶穌就是這樣提醒我們.
我們看到十九至二十三節裡.
這一段落.
如果我們再看回頭三節的時候.
其實耶穌就以一個智慧格言的方式.
來對門徒提醒.
究竟我們在生活裡.
我們怎樣來分.
地上的財寶和天上的財寶呢?.
來對應在上文所說的.

$^{921}$地上的報答和天上的賞賜.
在這三節的描述裡.
就是告訴我們這兩個的價值觀.
解經家也繼續來看.
他說這三節裡.
是重複了.
重複了什麼呢?.
重複了兩種的錢財.
兩種的眼光.
兩種的主人.
如果我們再看第二十四節.
就說一個人不能夠事奉兩個主.
但是我們先看錢財和目光.
這就成為解經家,釋經家.
告訴我們一個很重要的主題.
就是我們不要追求物質.
而是要專一地追求天父的旨意.
耶穌要求他的門徒.
不要積賺財寶在地上.
然後就說要積賺財寶在天上.
耶穌在這裡說的.
其實是一個命令.
大家都知道.
積賺財寶.
是要找一個地方.
那個地方放在哪裡呢?.
無論那個地方放在哪裡也好.
耶穌說也會有蟲來咬.
也會受害.
也會有賊在房屋裡.
在牆上鑿洞來偷走.
耶穌這樣的描述.
當時的財產是用什麼呢?.
用穀物,很名貴的衣物.
甚至是如果有加速的時候.
就用這些來做他的財產.
蟲子是咬什麼呢?.
咬這些穀物,農作物.
如果你積得太多.
你又用不完.

$^{961}$你又沒有施捨.
你會怎樣呢?.
蟲子會吐出來咬掉.
受害就是很寶貴.
很值錢的衣服,衣物.
當你不穿著.
一來會褪色.
二來也會受害.
當時猶太人所住的房屋.
是用泥牆來做的牆.
所以很容易被賊挖洞.
把你的財物偷走.
耶穌就是用這種.
人能夠明白的.
一個智慧,真言的方式.
來教導.
猶太人一定知道.
他的特定意思.
更加猶太人.
常常要追求天上的事.
什麼是天上的事呢?.
猶太人會很清楚.
天上的事是怎樣的呢?.
天上的財寶.
是不會變質.
是不會貶值.
也是永恆.
是最安全的.
按照耶穌的表達.
接著在22至23節.
耶穌再以一個比喻.
來提醒他的門徒.
要作出明智的選擇.
究竟你想要地上的財寶.
還是天上的財寶呢?.
在這裡耶穌提醒我們.
我們的選擇取決於.
我們身上一個.
很微小但很重要的器官.
就是我們的眼睛.

$^{1001}$耶穌在這裡說.
究竟你有瞭亮的眼睛.
還是有昏暗的眼睛.
因為眼睛就是身上的燈.
意思就是告訴我們.
眼睛成為我們的窗戶.
讓光線進入我們的人體.
耶穌用這個比喻是很生動.
甚至常常會將眼和心.
互為通用.
所以當我們說眼和心的時候.
即是表示我們要專心.
要集中.
要專注去做某一件事.
耶穌在這裡登山寶訓.
來提醒我們.
我們的心要放在正確的位置上.
我們要怎樣去選擇呢?.
所以耶穌有一個提醒.
祂說你的財寶在哪裡?.
你的心就在哪裡?.
耶穌曾經在祂的.
講道關於錢財財富的時候.
有很多的教導.
我們在福音書裡都看到.
耶穌有一次.
來到全道做教導的時候.
其中有一個少年人.
有一個富有的人.
要求耶穌為他和他的兄弟分家產.
耶穌沒有答應.
祂只是說了一個比喻.
祂說有一個很有錢的農夫.
他有很多豐富的農作物收穫.
當他為了農作物怎樣收藏.
怎樣擺放得好也好.
他不斷拆舊的倉房.
而是搭建更大的倉房.
來收藏他的財產.
是為了什麼呢?.

$^{1041}$為了吃喝快樂.
只是為了自己.
為了滿足自己的需要.
為了滿足自己的慾念.
但耶穌在這個比喻中說.
神今夜就要你的靈魂.
你所積存的去了哪裡呢?.
能否帶走祂呢?.
這就是自我自私.
他的視野只能看到自己.
就像耶穌所說的.
他的眼睛是昏暗的.
另外有一次.
有人來問耶穌.
他說:耶穌.
我做什麼善事.
才能得到永生呢?.
這個問耶穌的少年人.
他告訴耶穌.
他說:所有的敬虔的事.
律法的條例.
我已經守住了.
連禁食我都做了.
不過耶穌告訴他.
他說:你還欠缺一件事.
就是要將你所有的變賣.
消滯窮人.
然後來跟隨我.
原來這位少年人.
聽到耶穌所說的之後.
他就面帶愁容.
悠悠愁愁地走了.
這豈不是耶穌所說的.
他的眼睛也是昏暗的嗎?.
他只能看到短暫的利益.
卻失去了長遠看的視野.
還有一位弟兄.
他是一位成功的商人.
在他地上真的積賺了很多財寶.
很多物業.

$^{1081}$但是這位成功商人.
這位弟兄.
他做了很多的捐助.
很多熱心公益的事.
甚至變賣他的物業.
來做一些公益.
一些慈善事業.
這位是恩實商人.
這位弟兄.
他不求任何的名利.
他只是默默耕耘他的善事.
成為他的志業.
但是他的家人.
很想他來信耶穌.
但是他聽了很多年的福音.
他總覺得.
他說我行善已經積德.
中國人是這樣說的.
行善積德是好事.
我已經做了很多善事.
我又沒有積賺所有的財寶.
在地上留給自己.
所以他的價值觀是很清晰的.
他的眼光也很好.
很嘹亮.
但是.
都是.
牧師不斷跟他說.
傳福音.
都是說他欠缺了一件事.
但是到最後.
這位恩實成功的商人.
他真的領受到耶穌的福音.
給他的呼喚.
他就把他的視線.
把他的目光教正.
結果他的眼睛.
由朦朦朧朧.
或者叫做明亮也好.
他真的變為完全的明亮.

$^{1121}$來跟隨耶穌基督的價值觀.
到最後他真的.
決志來相信耶穌.
得到天上的賞賜.
讓姐妹.
明亮的眼睛.
會帶來我們身體的光明.
為我們的人生.
為我們的信仰生活.
更加可以有更好辨別的方向.
以致我們能夠在服侍神.
親近人.
我們能夠做得.
在上帝.
在耶穌基督裡面所說的.
我們施捨,禱告,禁食.
會帶來生活上.
更新的意義.
以致朝向積聚財寶在天上.
我深願弟兄姊妹.
能夠經歷.
把你的目光.
注視在積賺財寶在天上.
聖經裡面.
沒有否定或禁止我們擁有財物.
也沒有教導我們.
輕看我們所擁有的財物.
我們都會知道是上帝所賜的.
但更重要的是.
我們如何能夠.
享受上帝創造.
這位創造主厚賜給我們.
各樣豐富的時候.
我們不僅僅顧自己.
看自己.
如何好好地生活.
享受生活.
又或者任意揮霍.
又或者積存很多.
成為守財奴.

$^{1161}$而不願意施捨.
不關心別人的需要.
又或者更甚推演到.
貪得無厭.
以致走出一些.
不合乎耶穌教導的原則.
去尋找一些錢財.
這段經文提醒我們.
不要這樣做.
弟兄姊妹.
讓我們好好地反思.
去審查.
我們用什麼目光.
來看待地上的事情.
以致我們能夠.
作出正確的人生價值觀的取向.
為愛你而受苦.
受死復活的耶穌.
作出回應.
將你的生命.
將你的人生.
投資在永恆.
屬天的事情上.
我盼望你們願意.
來尋求神.
祈求祂賜下明亮的眼睛.
看到耶穌所賜.
屬天的賞賜.
我們一起來禱告.
讓你的教導成為我們的幫助.
好叫我們在今生.
在地上生活的時候.
能夠辨別方向.
更加看清楚.
我們所渴求的.
就是積贊財寶在天上.
願上帝你幫助我們.
我們會面對很多.
不同的挑戰和張力.
但願上帝你自己的教導.

$^{1201}$華語藏在我們心裡.
我們時刻靠著你.
在我們的生活裡見證出來.
誰聽我們在你面前的禱告.
奉主耶穌基督的名.
阿們.
\newpage



\section{馬太福音 7:1-6}
\label{sec:zKeckdMvwPY}
\textbf{2023年4月23日崇拜直播｜陳明泉牧師｜給當代門徒的15個忠告:別錯放辨別力｜馬太福音7\_1-6}
\newline
\newline
連結: \href{https://youtube.com/watch?v=zKeckdMvwPY}{\texttt{https://youtube.com/watch?v=zKeckdMvwPY}} ~~~~ 語音日期: 2023-04-24
\newline
\newline
\hyperref[sec:Dcz9krAW4uI]{\small{< < < PREV SERMON < < <}}
~
\hyperref[sec:index_chronic]{\small{[返順時目]}}
~
\hyperref[sec:PtF_7l7B_qs]{\small{> > > NEXT SERMON > > >}}
\newline
\newline
馬太福音 7:1-6
\newline
\begin{longtable}{cl}
\hline
\hline
章節 & 經文 (和合本修訂版)\\
\hline
7:1 & \begin{tabularx}{0.7\textwidth}{X} 「你們不要評斷別人，免得你們被審判。 \end{tabularx} \\ \\ \relax
7:2 & \begin{tabularx}{0.7\textwidth}{X} 因為你們怎樣評斷別人，也必怎樣被審判；你們用甚麼量器量給人，也必用甚麼量器量給你們。 \end{tabularx} \\ \\ \relax
7:3 & \begin{tabularx}{0.7\textwidth}{X} 為甚麼看見你弟兄眼中有刺，卻不想自己眼中有梁木呢？ \end{tabularx} \\ \\ \relax
7:4 & \begin{tabularx}{0.7\textwidth}{X} 你自己眼中有梁木，怎能對你弟兄說『讓我去掉你眼中的刺』呢？ \end{tabularx} \\ \\ \relax
7:5 & \begin{tabularx}{0.7\textwidth}{X} 你這假冒為善的人！先去掉自己眼中的梁木，然後才能看得清楚，好去掉你弟兄眼中的刺。 \end{tabularx} \\ \\ \relax
7:6 & \begin{tabularx}{0.7\textwidth}{X} 不要把聖物給狗，也不要把你們的珍珠丟在豬面前，恐怕牠們踐踏了珍珠，轉過來咬你們。」 \end{tabularx} \\ \\
[1ex]
\hline
\hline
\end{longtable}
$^{1}$今日的經訓是記載在.
《瑪太福音》七章一至六節.
是在你們手上的聖經新約第九頁.
上帝的話語這樣記載.
「你們不要論斷人,免得你們被論斷..
因為你們怎樣論斷人,也必怎樣被論斷..
你們用甚麼良器量給人,也必用甚麼良器量給你們..
為甚麼看見你弟兄眼中有刺,卻不想自己眼中有良目呢?.
你自己眼中有良目,怎能對你弟兄說:.
「容我去掉你眼中的刺呢?.
你這假冒為善的人,先去掉自己眼中的良目,.
然後才能看得清楚去掉你弟兄眼中的刺..
不要把聖物給吸走,也不要把你的珍珠掉在珠前,.
恐怕他踐踏了珍珠,轉過來咬你們.」.
.
今天原本不是我講的,但是梁偉中了COVID,.
原本今天是青少主日,年輕人繼續領詩,.
但是我們的青少主日會調到五月第一個星期,.
我臨時拉夫吃了這隻雞,去預備這篇的信息,.
但其實信息很早之前我們跟公一起談就預備了,.
大概有個框架,希望今天我們調一調次序,.
先講第七章,下一次梁偉再講的時候,.
講第六章最後的部分..
當我們去看整個系列,.
當代門徒十五個忠告,.
去到第七章開始的時候,.
某程度上就是耶穌基督講登山部分的結論部分,.
這個部分特別有一個很重要的思想,.
或者一個很重要的主調就是,.
你怎樣對人,人家將來會怎樣對你,.
更加重要的是天父將來會按著你怎樣對人的態度,.
來對待你,在你的身上,.
所以有一個類似因果循環的思路在這裡,.
而今天所講的論斷,.
或者關於怎樣去做判斷力,分辨能力,.
這種態度也成為了將來天父怎樣對待我們,.
或者我們身邊的人怎樣對待我們的一種態度,.
希望今天我們去領受這段經文的時候,.
我們都能夠為我們自己帶來一些生命裡的養份和反省,.
事實上這段經文很多的童工都講過,.

$^{41}$但實際上當我仔細看,不容易處理的,.
希望今天大家一起,我們專心的一起來看看上帝的話語,.
對我們今天有什麼訊息,.
我們先祈禱,求上帝幫助我們,同心祈禱,.
親愛主我求你在我們當中帶領我們,.
讓我們能夠在你的經文當中,.
讀到真正對我們門徒生命有意義的一些訊息,.
幫助我們,讓你的話語銘記在我們的內心當中,.
幫助此時此刻,中弟兄姊妹能夠清晰的領受你的教導,.
又讓孫講的講得清楚,.
讓我們銘記你的話語,讓我們行事為人,.
被你所稱他,也成為這個世界的光和炎,.
求你在我們當中繼續率領我們下來領受你話語的時間,.
獻上禱告奉靠基督耶穌得勝的聖名,Amen.
我們看第七章第一至到第六節,.
你們不要論斷人,免得你們被論斷,.
因為你們怎樣論斷人,也必怎樣被論斷,.
你們用什麼良器量給人,也必用什麼良器量給你們,.
這是第七章第一至到第二節,.
這是第一個小段落,.
在這個第一個小段落裡,.
我們一定要把焦點放在論斷這個字裡面,.
論斷這個字,原文的字眼用的就是quino,.
這個字其實有很多意思,.
這個字在不同的上下文裡面有不同的翻譯,.
在這裡面的論斷,是在講一種批評,挑剔,.
一個比較負面的意思,.
所以論斷在這個上文下理裡面,.
有一個解法就是一種批評,.
有一些更加極端的翻譯,.
那個意思就定為定罪,condemn,.
論斷這個字如果用在上帝對人身上,.
或者用在一個法官身上,.
這個字就代表定罪的意思,.
當然有時候這個字未必去到最極致,.
去到審判或者定罪,.
在生活層面裡面都用這個字,.
分辨,明辨是非,判斷能力,.
都是用這個字的,.
所以論斷這個字的意思或者內容就要看上文下理,.

$^{81}$更加重要有時候在聖經裡面都會玩work play,.
就是在講一字多義,.
一個字同樣的意思,.
但上面是這個意思,.
後面有另外一種意思,.
所以我們去讀這段經文的時候,.
首先要留意一下耶穌基督用這個字眼,.
究竟要講什麼內容,.
我們掌握到之後,.
我們就會容易理解經文當中所講的信息,.
在這裡其實很清楚,.
大概在這裡耶穌基督呼籲那些百姓或者門徒,.
你們不要,如果用論斷覺得不清楚,.
不要肆意批評,不要挑剔人家,.
免得你們被論斷,.
論斷就不是被人批評,.
免得你們被審判,.
那是講上帝對人的內容,.
因為你們怎樣挑剔人,.
也必怎樣被審判或者被定罪,.
你們用什麼良器良給人,.
也必用什麼良器良給你們,.
為什麼會用審判這個字,.
或者用挑剔這個字,.
其實記不記得剛才所講,.
在第七章開始,.
耶穌基督要講登山寶訓,.
他有一個神學很重要的主調,.
七章十二節,.
在七章十二節裡面有一個很重要的主調,.
告訴我們耶穌基督的思路,.
所以無論何時,你們願意人怎樣待你們,.
你們也要怎樣待人,.
因為這就是律法和先知的道理,.
你想怎樣別人對你,.
你用什麼方法對人,.
根本上後面所講的,.
別人怎樣待你,.
也講到上帝怎樣待你,.
所以你的人倫關係和天父將來和你之間的關係,.

$^{121}$兩者是有關連的,.
那不是純粹講人際關係這麼簡單,.
這是講到將來天父會按著我們在地上的對人方式,.
來對待我們,.
當然天父是很慈愛的,.
但不要覺得天父是很慈愛,.
我們做的所有錯事,.
我們就放肆或者任由我們自己,.
耶穌基督在這裡面,.
他警醒我們,.
他叫我們今天怎樣待人接物,.
和將來上帝怎樣對我們,.
兩者是有關係的,.
所以我們生活裡面很多事情是不能輕忽的,.
不能肆意地覺得自己有恩典,.
我們喜歡怎樣就怎樣,.
舉個例子,.
不知大家在公司也好,.
在家裡面有沒有見過,.
接觸過一些老佛爺,.
知道甚麼是老佛爺嗎?.
做甚麼他都不對,.
然後他就會挑剔,.
杯水這麼熱,.
倒另一杯,.
冷一點,.
杯水這麼冷,.
都不知你覺得多少度為之對,.
或者做每一樣事,.
他都會在你身上找到一些不理想的地方,.
你覺得如果和這些老佛爺,.
未必一定是女士,.
男士都會有老佛爺感覺,.
面對這一類的弟兄姊妹,.
你辛苦嗎?.
很辛苦的,.
在你的人際關係裡面,.
如果喜歡挑剔,.
或者肆意批評,.
對我們的人際關係有很大的困難,.

$^{161}$換個角度來說,.
其實很喜歡論斷,.
或者批評人的人都很慘,.
為甚麼呢?.
因為喜歡批評的人,.
大概內心不是一個快樂的人,.
他很多時看見身邊所有事都不合心意,.
所有人與事都沖著他自己,.
他看每一樣事都不順眼,.
他是一個不快樂的人,.
他的人際關係大概都不是很理想,.
因為沒有甚麼人可以承受一個經常論斷的關係,.
偶爾為之可能有些近似,.
都是很喜歡論斷的人走在一起成為一個團伙,.
最初可能是論斷外面的人,.
但到沒有人在旁邊的時候,.
論斷自己的朋友,.
所以在這個過程裡面,.
他的人際關係,.
一種論斷的關係,.
是沒有辦法為這個人帶來真正親密的關係,.
一個論斷慣了的人,.
他是一個孤單的人,.
論斷的人可能都是一個完美主義者,.
他看身邊所有事,.
總是有一種不完美的感覺,.
而這種不完美的感覺,.
就會令到自己的心裡面有一種,.
一團亂,刺著,.
不過這個完美主義者不單止是對事物,.
他的神的觀念,他的神觀,.
他都會覺得天父是一個很嚴苛的上帝,.
對於世上所有瑕疵,.
對於人生命裡面有很多不理想的地方,.
可能他覺得天父是不會接納的,.
他覺得天父容不下,.
生命裡面有一點點瑕疵,.
或者不完美的地方,.
他極力成為一個完美主義者,.
因為他盼望他這樣才可以得到天父的接納,.

$^{201}$或者人對他的接納,.
論斷的人也是一個沒有安全感的人,.
他生命裡面沒有一件事是可以靠賴的,.
他藉著論斷,.
希望可以透過這些批評的話語,.
找到自己的價值,.
不過當他身邊沒有人沒有物,.
或者在論斷的過程裡面,.
他找不到真正的對象的時候,.
他就好像失去了身邊所有.
令他站得穩陣腳的一些東西,.
以致他一個很痛苦的經歷,.
有一本書叫做《建立高EQ的靈性》,.
Peter Cazoro那本書,.
這本書是曾真牧師介紹給我們看的,.
一本很好看的書,.
他其中講到一個很重要的一點,.
他說屬靈裡面,.
在靈性裡面有一個很重要的困難,.
就是慣性去做一些論斷和批評,.
當一個人慣性進入了這種論斷批評,.
或者群體之間習慣了一種這樣的關係的時候,.
群體沒有辦法經歷那種親密感,.
更加重要的是在這種批評,.
或者挑剔的視野裡面,.
其實去看自己比其他人更高一線,.
他覺得自己的標準,.
是成為世界上絕對的標準,.
他是一個驕傲的人,.
是一個沒有辦法過著融洽相處的生命,.
這本書的作者講到,.
其實論斷講到尾是什麼呢?.
其實是將世界的不滿投射在其他人的身上,.
將一種不滿感,.
用一種自己不滿的憤怒,.
指向身邊某一個人,.
可能對世界不滿,.
又或者對他自己曾經的經歷不滿,.
而這種不滿感,.
他要找一些對象來發洩,.

$^{241}$論斷就成為了一種這樣的狀態,.
令到人破壞了這種關係,.
這樣講我們知道,.
耶穌基督叫門徒不要進入一種論斷的關係,.
這是一個很顯明意見的道理,.
論斷不是一種人與人之間理想的關係,.
即使我們很想幫助別人,.
鞭策別人去做得更加好,.
但如果用著論斷的態度,.
我們沒有辦法幫助別人成長,.
耶穌基督叫我們要棄絕這種態度,.
在羅馬書裡面,.
保羅也是承接著耶穌基督的教導,.
所以我們不可再彼此論斷,.
寧可定以誰也不給弟兄放下半腳跌人之物,.
簡單來說,.
羅馬書裡面說這個人吃這些藥,.
或禁戒自己做禁食,.
人們就用別人的行為來互相指控姐妹之間的關係,.
這種論斷關係,.
在教會裡面,.
在我們的人生裡面,.
我們要禁戒的,.
要令到自己不要進入這種狀況裡面,.
其實經文講完這裡,.
可以講完的了,.
今天可以完結了,.
可以收工了,.
不過耶穌還未收工,.
講完這裡之後,.
接下來更多的篇幅,.
告訴我們,.
既然論斷是這麼差,.
我們應該怎樣去做改變呢?.
答案就是,.
論斷本身這種能力,.
記不記得剛才說的這個字,.
是一字多二的,.
這個判斷能力,.
這種某程度很細緻地去看,.

$^{281}$人們生命的不是,.
這種能力,.
其實如果放在對的位置,.
或者放在對的位置,.
是成為我們的出路,.
放錯位置就是罪,.
放在對的位置,.
那件事是成為生命一個很重要的幫助和養分,.
不過耶穌基督就用一個比較負面的角度來講,.
我們繼續看,.
七章第三節至第五節,.
為什麼看見你弟兄眼中有刺,.
卻不想自己眼中有良目呢?.
你自己眼中有良目,.
怎能對你的弟兄說,.
用我去掉你眼中的刺呢?.
你這假冒為善的人,.
先去掉自己眼中的良目,.
然後才能看得清楚,.
去掉你弟兄眼中的刺..
在這裡,.
耶穌繼續講這個講論,.
耶穌繼續用一些很絕核的方式來講一個活生生的比喻,.
他講到對著那些經常論斷的人,.
又或者要警惕自己不要論斷的人,.
他就用一個活靈活現的比喻,.
為什麼你經常看見弟兄姊妹眼中,.
那裡有條刺,.
有條幼的木舌,.
眼睛都很小,.
但在論斷的人眼中裡面,.
可以細緻到別人眼裡面有條刺,.
他都看得到的,.
是不是很眼利?.
這個很眼利的人,.
耶穌說,.
但你看不到有條大的良目在你的眼前,.
你視力那麼好,.
但看不到你自己原來眼前有一條大的良目,.
阻礙了自己來看人,.

$^{321}$看事物..
耶穌用這個那麼有趣的,.
或者一個很絕核的比喻,.
就誇張了,.
原來很多時候我們喜歡批評的人,.
我們就在別人身上找一些小錯處,.
但我們很少用這種能力來審視自己生命裡面的大錯處,.
所以耶穌基督說,.
你這假冒為善的人,.
先去掉自己眼中的良目,.
然後才能看得清楚,.
去掉你弟兄眼中的刺,.
你留意一下,.
講這個論斷的時候,.
他最後某程度上有一個呼籲,.
他稱呼那些人叫做假冒為善,.
之前說過,.
假冒為善的人,.
其實在登山部分,.
順民夏利都在用的,.
在祈禱裡面戴著面具,.
或者做善事的時候,.
扮演有一個好的角色,.
但論斷的人都是假冒為善,.
為什麼呢?.
假冒為善除了是找一些東西來表演自己的敬虔之外,.
原來假冒為善也有一種心態,.
就是進攻就是最佳的防守,.
什麼意思呢?.
就是說去挑剔別人生命的錯處,.
就是為了保障自己生命不需要揭露的一個手法,.
用挑剔別人來遮掩自己生命的錯處,.
所以耶穌基督在這裡,.
他很尖銳地將論斷人的生命呈現出來,.
論斷到最後,.
其實是一種假冒為善的表現,.
我們不停地指出別人生命的不理想,.
或者挑剔身邊的人,.
其實這種舉動是在掩飾自己生命有很多不理想的地方,.
所以耶穌基督在這裡,.

$^{361}$你有這種細緻看到別人錯處的能力,.
這種能力不要再放在別人身上,.
來查找任何一些不理想的地方,.
反而你要用這種能力來看回自己生命裡面,.
找回自己生命裡面那條大的良木,.
別人有錯不等於自己對的,.
別人的不是不等於自己一定是站得對的,.
整個過程是可以別人錯,自己也會錯的,.
耶穌基督盼望他的門徒,.
用這種審視別人的能力,.
同樣應用在自己的身上,.
所以在這裡,.
其實也是在回應第一節和第二節,.
你要用相同的良器來量度別人,.
如果你看到別人生命裡面的錯處,.
在第二節裡面,你用什麼良器量級別人,.
也必用什麼良器量級你們,.
這個既是說那個結果,.
同時他也是在呼籲這麼喜歡論斷批評的人,.
請你看看,用你同樣的標準來應用在自己身上,.
你不可以說你自己對自己很寬鬆的,.
看不到自己的東西,.
完全超然在那個狀態,.
純粹用一個最mean的標準來量度身邊的人,.
既然那個標準是可以量給別人,.
同樣也要量給自己,.
用同樣的標準來看,.
最經典的例子,.
看球賽,.
如果男士喜歡看球賽,.
通常會說那是球史,.
罵得最厲害,.
有時候不一定要看英超,.
在麥花臣,我那時候有空去看人生百態,.
有些大叔坐在樓梯級看球賽,.
下面的人都是像我們這些傻子一樣踢球,.
踢球不是傻子,我才是傻子,.
踢球踢得不太好,.
有個人在樓梯級說那是球史,.
罵人家,.

$^{401}$但當球在他身上踢的時候,.
他也是差不多一樣,.
但很多時候罵也罵得很凌厲,.
但如果同樣的東西應用在自己身上的時候,.
自己就知道其實不是想像中那麼簡單,.
還有各位大英姐妹,.
有時候當我們去看別人不理想的地方,.
我們也用同樣的標準來看看我們自己,.
如果我們設身處地,.
我們做同樣的事,.
我們是否同樣也能有這樣的表現呢?.
唯有我們真的比較客觀一點,.
用那麼嚴厲的標準套用在自己身上,.
我們看人,我們對人的判斷才比較公道,.
聖經裡面說的分辨就是這個意思,.
分辨的意思同樣是用quino這個字,.
不過這個字不是純粹按自己的喜好,.
按自己的標準來看身邊的人,.
而是按一個客觀和上帝啟示出來的標準,.
來應用在別人,也應用在自己身上,.
當我們能夠客觀地去看的時候,.
我們看事物,我們對自己,對別人,.
我們就會變得公道一些,.
而不是純粹站在我們的位置,.
就指點別人的江山,.
經文裡面再繼續說,.
這種如果能夠挑剔別人的態度,.
來檢視自己的生命,.
我相信那是為我們自己的生命帶來革命性的改變,.
有一個沙漠教父,.
他在說一個論斷的主題,.
他說如果當人用大量的時間來看自己,.
用論斷的能力來看自己的生命,.
根本就沒有時間去看其他人生命的不是,.
因為我們生命千瘡百孔,.
我們都沒有辦法去兼顧別人的生命,.
有什麼不順眼或不順眼的地方,.
用同樣的東西來看自己的生命狀態,.
如果我們不看自己的生命,.
而是盲目地去還給別人的時候,.

$^{441}$耶穌基督在這裡說,.
你就沒有辦法看得清楚,.
我們眼睛如果有一條大的良目,.
我們說多細緻地去看別人也好,.
我們看的東西都是帶著偏頗,.
都是帶著自己的成見來看身邊的人,.
我們就沒有辦法好好地去活出一個好好的生命,.
在這裡我有一個這樣的例子,.
這是上帝教導我一個很重要的例子,.
話說我認識一個年輕人,.
我自己也很年輕的時候,.
覺得這個人有點成見地對待他,.
小時候看著他做事不是很規矩,.
又很自我中心,.
覺得這個人難成大器,.
所以我那時候還沒有相信他,.
覺得這個人這樣,.
我小時候很壞,.
喜歡欺凌比我小的人,.
經常糟蹋這些比我小的後輩,.
覺得他沒有用,.
然後就打他,擠他,.
未至於很重的欺凌,.
不過就,.
總之這些壞東西都出來了,.
大家明白了,.
不用我去畫公仔畫出腸,.
後來自己一路信了主,.
自己一路改變,.
覺得原來自己以前那麼壞,.
這樣欺負自己的後輩,.
糟蹋別人,.
但是有一種成見,.
心裡面也在,.
覺得那些人都是難成大器,.
不過後來這個後輩,.
一路長大的時候,.
這個人也成長了,.
跟我小時候看他的樣貌,.
或者做人處事的方式,.

$^{481}$完全脫胎換骨,.
然後看到他一路成長長大的時候,.
然後我就不好意思說,.
不過我心裡面有一種很大的悔疚,.
我跟上帝認罪,.
我問上帝為什麼我以前,.
我看人,.
我那麼有成見,.
為什麼我判定了這個人是差的,.
我就覺得他永遠是差,.
為什麼我不覺得這個人可以,.
在你的恩典當中會有成長呢?.
但是上帝改變了我,.
上帝也改變了這個人,.
以致我發覺生命是可以改變的,.
更加重要的就是,.
我們作為長輩也好,.
或者我們成為一個基督徒,.
不要再用成見,.
不要再用我們自己覺得對的標準,.
套在別人身上,.
肆意批評身邊的人,.
這個人後來成為了我一個很好的朋友,.
後輩也好,.
朋友也好,.
我有心事都會跟他談,.
我覺得他生命是成熟,.
可以承托著人生很多的故事,.
弟兄姊妹,我不知道你會不會是這樣呢?.
在教會裡面,.
人與人之間的相處,.
這個人這樣就沒得救了,.
又或者覺得某些人,.
你已經覺得他不行,.
你就用成見的態度,.
跟他有一種距離感,.
隔絕自己,.
覺得自己不會跟這些人在一起,.
我們喜歡跟自己談得來的人走在一起,.
而跟我們不同的,.

$^{521}$或者我們看不上眼的人,.
我們就將他分門別類,.
敬而遠之,.
甚至乎不敬而遠之,.
這種態度,.
耶穌基督是很嚴厲地批評,.
耶穌基督甚至乎說,.
當代的法利賽人,.
對於那些稅吏,.
對於那些妓女,.
低下階層,.
他們就是這樣,.
作為一個門徒,.
如果我們都好像法利賽人,.
我們勝不過他們的義,.
我們怎樣進入天國呢?.
所以經文裡面,.
耶穌基督就叫我們,.
你要用這種判斷的能力,.
用這種看別人的挑剔的能力,.
來挑剔自己,.
將你自己眼中的良目整走,.
以致你對人,.
對事,.
你才能夠看得清楚,.
所以在這裡所說的,.
一字多二,.
就在這裡應用,.
你要分辨,.
分辨生命,.
你要運用這種能力,.
而這種能力就是幫你自己,.
看你自己的罪,.
能夠好好地去看,.
經文最後一個小節,.
七章第六節,.
這個也是一個很難解釋的經文,.
先聽我讀,.
不要把聖物給狗,.
也不要把你們的珍珠丟在豬前,.

$^{561}$恐怕牠踐踏了珍珠,.
轉過來咬你們,.
當你讀這裡的時候,.
很多色經學家,.
或者很多朋友,.
停在七章第五節裡,.
就說完了,.
不過耶穌說的語氣還沒完,.
無端端說不要把聖物給狗,.
不要把珍珠丟在豬面前,.
那隻狗和豬,.
那隻狗不是你家養的那隻,.
那隻是野狗,.
那隻豬不是你家吃的那隻,.
或者那些貴價黑毛豬,.
如果你知道當代猶太人養豬,.
和放羊一樣是放牧的,.
某程度上是野性的豬,.
你見到下面有些牙那些,.
比較有力量的豬,.
會有攻擊性的,.
在這裡,.
耶穌基督用一個很奇怪的比喻,.
有些悉經學家說,.
不關事的,可能關下面的事,.
但下面更加不關事,.
你們祈求的,就給你們,.
尋找就尋見,.
上文下理好像都不通,.
有些索性不解,.
姑懼,.
不處理這段經文,.
用其他東西來說,.
不過其實在經文裡,.
似乎耶穌基督用這個比喻,.
我要再引申判斷力的問題,.
記不記得剛才所說,.
那個論斷的字,.
是有很多的意思,.
而下面這個,.

$^{601}$將聖物給狗,.
或者將珍珠掉在豬的面前,.
這個是一個不加思索,.
不判斷,.
將好的東西,.
就掉,.
隨意給那些野生的動物,.
這個過程是一個不用腦的表現,.
一種不判斷,.
不想那樣東西對不對,.
那些是好東西,.
我想說,.
將珍珠,.
你不要想著是一顆碎石,.
是一顆珍珠,.
是有價值,.
是好東西,.
不加思索地將好東西,.
給那些不識貨的,.
或者他都不明白,.
什麼叫珍珠的動物身上的時候,.
反過來那隻豬,.
不會覺得你善待牠,.
牠覺得你只不過是玩牠,.
甚至乎牠覺得牠有攻擊,.
所以牠就即使已經咬一顆珍珠,.
硬的,.
然後牠就走過來,.
就咬你回頭,.
覺得你是對牠不懷好意的,.
耶穌基督說這個這樣的比喻,.
如果用剛才說的判斷,.
你就會解得通了,.
那個判斷,.
或者那個分辨的智慧,.
不是去看別人的挑剔,.
第一步是要看自己,.
檢視自己,.
生命裡面眼中有沒有良目,.
第二件事就是要檢視自己所做的行為,.

$^{641}$是不是合宜的,.
中國人有句話說道,.
「宜者宜也」,.
什麼叫宜呢?.
中國人裡面說,.
除了說公義之外,.
更加是說在對的場合做對的事,.
那個就為之宜了,.
同樣道理在這裡裡面說,.
什麼叫判斷呢?.
除了說那種能力來檢視自己生命的靈性之外,.
更加要看自己所做的,.
基於一個好的心腸,.
但是做的事是不是帶來好的效果,.
你身邊的人是不是接納你所說的內容,.
或者你所做的那件事,.
如果你不能夠做到這件事,.
失去了這種辨別能力,.
你就會糟蹋了一些好的東西,.
那些聖物,.
比如說狗,.
把神聖的東西給了一隻不神聖的狗來吃,.
一顆好的珍珠就給了一隻野豬來踐踏,.
再加上那隻豬不會覺得你的好意,.
反過來就咬你,.
整個表達的判斷能力,.
其實基督徒是要運用的,.
只不過是放在一個位置上,.
不要覺得論斷,.
或者判斷是與非,.
這種能力,.
做任何事,.
然後就覺得不要論斷,.
所以就把這種能力扔掉,.
反而我們要放在一個位置上,.
這個世界依然要分辨的,.
我們讀聖經,.
其實就要有分辨的能力,.
所以在哥林多前書裡面,.
保羅同樣用這個字,.

$^{681}$用一個更加強力的字,.
他就說每一次領主餐的時候,.
如果你不是分辨主的身體,.
你就吃喝自己的罪,.
那個分辨就是用quino這個字,.
我們人生裡面,.
還有你聽道的時候,.
三兩個人都要分辨講的內容,.
是不是合乎聖經,.
這個是在哥林多前書第十四章裡面,.
保羅是這樣教一個群體,.
分辨能力,.
或者論斷的能力,.
放對了位置,.
這個是我們生命的寶藏,.
如果你不去不加思索,.
不用,.
你就會帶來不好的效果,.
如果這個裡面,.
用回這種應用,.
我們今天覺得好的福音,.
我們宣講真理,.
我們是不是用合宜的態度來表達,.
有時候我們會用令到人反感的方式,.
來將這個真理傳開他,.
在我還沒有講真正的福音之前,.
我已經令到身邊的人害怕了,.
有一個這樣的負面例子,.
這個是我們教會,.
不單止我們教會,.
是整個旺角的教會做的一個例子,.
我這樣公開批評好像不太禮貌,.
不過這個是讓我自己學一個功課,.
我不知道大家記不記得,.
在大概沙士2003年的時候,.
有些老黃浸就知道了,.
我們為了要為沙士的弟姐妹打氣,.
就做了一個叫做生命加油站,.
有些印象了,.
我們那個做法,.

$^{721}$因為那時候我們下面那條街,.
行人專用區,.
整條路是禮拜日可以封了街,.
沒有車可以走的,.
就可以讓人做一個專用區,.
讓我們走得鬆動一點,.
我們一班有時候有些旺角區的,.
原本打算找一個角落,.
找一個角落傳好訊息,.
傳福音,.
有一些好訊息告訴大家,.
這個動機很好的,.
一直做的時候,.
到我們實習的時候,.
2005年做了兩年多,.
兩年多三年了,.
都還繼續做,.
舊時是說一個街裡面一個角落來做的,.
越做越大型,.
在我來實習的時候,.
我看到半條行人專用區,.
擺了那些椅子,.
又用發電機,.
又搭過台在那裡,.
大家有沒有印象,.
很大場景,.
整個旺角行人專用區,.
就是一個舞台,.
一個講福音,.
一個唱詩歌,.
有人坐在那裡的一個很大型的街頭製作,.
我們覺得很棒,.
後來那次我有機會下去傳福音,.
打算做完這個大製作之後,.
跟著下去跟商戶傳福音,.
那時候第一次,.
我走到商舖那裡,.
給他一張單張,.
商舖裡面那個,.
進去某間電器舖,.

$^{761}$他就說,.
走啦你們,.
很煩呀你們,.
很不禮貌的,.
來對待,.
我就說,.
我什麼都沒做,.
我只是給你一張單張,.
我想問一下,.
為什麼這麼反感,.
然後那個店員就很清晰說,.
以前你們只是一個角落,.
小聲一點,.
但是現在,.
你每個星期日都吵得很厲害,.
這段時間我們是做生意的,.
我們有時候要講解,.
我的產品是怎麼用,.
你們的音樂比我們還大聲,.
所以我要跟你們的音樂,.
跟你們的詩歌鬥大聲,.
然後,.
整條路被你們霸佔了,.
我們真的被行人走的,.
都沒地方走,.
全部就是你們教會的人,.
就霸佔了整條路,.
我見到你們就覺得你們很討人厭,.
我第一次聽真話,.
我們第一次在這裡,.
我們覺得我們做的事情,.
是祝福街坊的,.
是好事來的,.
但是我發覺,.
不合宜的,.
在這件事裡面,.
我放在心裡幾年了,.
十幾年才跟大家說,.
做了主任牧師,.
這件事沒了,.

$^{801}$才總結這個經驗,.
我們很多時候說,.
這個世界敵黨福音,.
我們覺得這個世界是抗拒基督耶穌,.
不過,.
可能有時候那種敵黨,.
不是因為敵黨這個信仰,.
可能是敵黨一些,.
其實我們覺得好,.
但是將聖物給狗那種,.
我們不加思索,.
很懶,.
人家做開覺得第一次好,.
然後就不停地重複又重複,.
那件事,.
我們成為了福音的難阻,.
我們帶著善良的動機,.
來做這件事,.
但是當我們不去再用分辨的能力,.
不再去看那件事合不合宜的時候,.
可能我們所做的一切努力,.
反而成為福音的難阻,.
在人家沒接受,.
沒聽福音之前,.
他已經拒絕了我們這班人,.
拒絕了我們自己,.
求主幫助我們,.
如果我們這個世代裡面,.
你說論斷,.
是不是說這件事就是論斷呢?.
如果人家覺得是的,.
我覺得我們都要接受,.
不過我們更加要負心自問,.
今天在我們的場景,.
我們有沒有用這種,.
說人家,.
批評人家的能力,.
都應用在自己的檢視身上,.
檢視自己,.
我們是不是用同樣的標準,.

$^{841}$來看我們自己的生命,.
檢視自己,.
我們說人家的時候,.
我們會不會都看看自己的生命,.
是不是想像中那麼理想,.
我們更加要檢視,.
有一些我們的做法,.
我們的行為,.
我們會不會帶著善良的動機,.
但是做了一些不合宜的工作,.
求主幫助我們,.
當我們如果能夠善用這種判斷能力,.
我們一銀兩面的,.
我們會成為別人的祝福,.
我們會成為我們自己靈性的提升,.
我們更加看到自己生命的真相,.
我們更加踏實地,.
來經歷門徒的生活,.
求主繼續幫助我們,.
今天基督耶穌叫我們不要論斷,.
但是不要失去了判斷的能力,.
這個就是今天我們要領受的訊息,.
求主繼續對我們每一個人說話..
拜拜.
\newpage



\section{詩篇 73:1-28}
\label{sec:PtF_7l7B_qs}
\textbf{2023年4月30日崇拜直播｜區祥江博士｜如何面對嫉妒的困擾｜詩篇73\_1-28}
\newline
\newline
連結: \href{https://youtube.com/watch?v=PtF_7l7B_qs}{\texttt{https://youtube.com/watch?v=PtF\_7l7B\_qs}} ~~~~ 語音日期: 2023-05-01
\newline
\newline
\hyperref[sec:zKeckdMvwPY]{\small{< < < PREV SERMON < < <}}
~
\hyperref[sec:index_chronic]{\small{[返順時目]}}
~
\hyperref[sec:38XteBMXhDk]{\small{> > > NEXT SERMON > > >}}
\newline
\newline
詩篇 73:1-28
\newline
\begin{longtable}{cl}
\hline
\hline
章節 & 經文 (和合本修訂版)\\
\hline
73:1 & \begin{tabularx}{0.7\textwidth}{X} 神實在恩待以色列那些清心的人！ \end{tabularx} \\ \\ \relax
73:2 & \begin{tabularx}{0.7\textwidth}{X} 至於我，我的腳幾乎失閃，我的步伐險些走偏； \end{tabularx} \\ \\ \relax
73:3 & \begin{tabularx}{0.7\textwidth}{X} 因為我嫉妒狂傲的人，我看見惡人享平安。 \end{tabularx} \\ \\ \relax
73:4 & \begin{tabularx}{0.7\textwidth}{X} 他們的力氣強壯，他們死的時候也沒有疼痛。 \end{tabularx} \\ \\ \relax
73:5 & \begin{tabularx}{0.7\textwidth}{X} 他們不像別人受苦，也不像別人遭災。 \end{tabularx} \\ \\ \relax
73:6 & \begin{tabularx}{0.7\textwidth}{X} 所以，驕傲如鏈子戴在他們項上，殘暴像衣裳覆蓋在他們身上。 \end{tabularx} \\ \\ \relax
73:7 & \begin{tabularx}{0.7\textwidth}{X} 他們的眼睛因體胖而凸出，他們的內心放任不羈。 \end{tabularx} \\ \\ \relax
73:8 & \begin{tabularx}{0.7\textwidth}{X} 他們譏笑人，憑惡意說欺壓人的話。他們說話自高； \end{tabularx} \\ \\ \relax
73:9 & \begin{tabularx}{0.7\textwidth}{X} 他們的口褻瀆上天，他們的舌毀謗全地。 \end{tabularx} \\ \\ \relax
73:10 & \begin{tabularx}{0.7\textwidth}{X} 所以他的百姓歸到這裡，享受滿杯的水。 \end{tabularx} \\ \\ \relax
73:11 & \begin{tabularx}{0.7\textwidth}{X} 他們說：「神怎能曉得？至高者哪會知道呢？」 \end{tabularx} \\ \\ \relax
73:12 & \begin{tabularx}{0.7\textwidth}{X} 看哪，這就是惡人，他們常享安逸，財寶增多。 \end{tabularx} \\ \\ \relax
73:13 & \begin{tabularx}{0.7\textwidth}{X} 我實在徒然潔淨了我的心，徒然洗手表明我的無辜， \end{tabularx} \\ \\ \relax
73:14 & \begin{tabularx}{0.7\textwidth}{X} 因為我終日遭災難，每日早晨受懲治。 \end{tabularx} \\ \\ \relax
73:15 & \begin{tabularx}{0.7\textwidth}{X} 我若說「我要這樣講」，就是愧對這世代的眾兒女了。 \end{tabularx} \\ \\ \relax
73:16 & \begin{tabularx}{0.7\textwidth}{X} 我思索要明白這事，眼看實係為難， \end{tabularx} \\ \\ \relax
73:17 & \begin{tabularx}{0.7\textwidth}{X} 直到我進了神的聖所，思想他們的結局。 \end{tabularx} \\ \\ \relax
73:18 & \begin{tabularx}{0.7\textwidth}{X} 你實在把他們安放在滑地，使他們跌倒滅亡； \end{tabularx} \\ \\ \relax
73:19 & \begin{tabularx}{0.7\textwidth}{X} 他們轉眼之間成了何等荒涼！他們被驚恐滅盡了。 \end{tabularx} \\ \\ \relax
73:20 & \begin{tabularx}{0.7\textwidth}{X} 人睡醒了，怎樣看夢，主啊，你醒了也必照樣輕看他們的影像。 \end{tabularx} \\ \\ \relax
73:21 & \begin{tabularx}{0.7\textwidth}{X} 因此，我心裡苦惱，肺腑被刺。 \end{tabularx} \\ \\ \relax
73:22 & \begin{tabularx}{0.7\textwidth}{X} 我這樣愚昧無知，在你面前如同畜牲。 \end{tabularx} \\ \\ \relax
73:23 & \begin{tabularx}{0.7\textwidth}{X} 然而，我常與你同在；你攙扶我的右手。 \end{tabularx} \\ \\ \relax
73:24 & \begin{tabularx}{0.7\textwidth}{X} 你要以你的訓言引導我，以後你必接我到榮耀裡。 \end{tabularx} \\ \\ \relax
73:25 & \begin{tabularx}{0.7\textwidth}{X} 除你以外，在天上我有誰呢？除你以外，在地上我也沒有所愛慕的。 \end{tabularx} \\ \\ \relax
73:26 & \begin{tabularx}{0.7\textwidth}{X} 我的肉體和我的心腸衰殘；但神是我心裡的力量，又是我的福分，直到永遠。 \end{tabularx} \\ \\ \relax
73:27 & \begin{tabularx}{0.7\textwidth}{X} 看哪，遠離你的，必要死亡；凡離棄你行淫的，你都滅絕了。 \end{tabularx} \\ \\ \relax
73:28 & \begin{tabularx}{0.7\textwidth}{X} 但我親近神是於我有益；我以主耶和華為我的避難所，好叫我述說你一切的作為。 \end{tabularx} \\ \\
[1ex]
\hline
\hline
\end{longtable}
$^{1}$(教育)接著是讀經的時間,請翻開舊約聖經748頁,經文是詩篇73篇,翻開之後請聽我讀..
(教育).
(陳健波) 大家早安,剛才打出來嚇了我一跳,這個鍋很美麗,值得人羨慕..
(陳健波) 今天跟大家分享一段經文,詩篇73篇是很多人都熟悉的詩篇,一般是處理義人與惡人之間,見到義人風生水起,.
(陳健波) 反過來說,見到惡人風生水起,義人心中如何疏解自己的內心呢?.
(陳健波) 今天我選擇了另一個角度跟大家分享,就是「疾度」這個主題..
(陳健波) 不知道大家在生活中有沒有羨慕或妒忌一些人呢?.
(陳健波) 生了一個兒子,那麼乖,那麼聰明,去他家探訪,看到他家那麼美麗,我就住在這些地方..
(陳健波) 原來妒忌這種心情,在日常生活中都經常出現..
(陳健波) 不過如果輕度一點,我們會說你羨慕還是妒忌?.
(陳健波) 羨慕就OK,妒忌其實內心會有種不太舒服的感覺..
(陳健波) 通常妒忌這種情緒都會影響我們,會令我們有些自卑..
(陳健波) 你有我沒有,你有我沒有..
(陳健波) 不過在我們中文翻譯,妒忌可以有兩個英文字..
(陳健波) 不知道大家有沒有一個叫jealous,一個叫envy..
(陳健波) 其實這兩個字是有少少分別的,不知道大家懂不懂得分辨..
(陳健波) 今天是講envy,不是講jealous..
(陳健波) 很多年前看過一套電影,如果你看過你應該都年紀不輕了..
(陳健波) 叫莫扎特之死,有沒有人看過這套電影?.
(陳健波) 看過的都是跟我差不多年紀的..
(陳健波) 這套電影簡單來說是講一位音樂家,他同期莫扎特在,很年輕,三十多歲..
(陳健波) 天才橫刃,作曲很美麗..
(陳健波) 連這個音樂家也是作曲的人,都會被莫扎特的作品感動..
(陳健波) 不單止感動,他還出了一種妒忌..
(陳健波) 為什麼他這麼厲害,這麼有才華,而我沒有呢?.
(陳健波) 妒忌嚴重起來是很厲害的,因為這套莫扎特之死是一個野史..
(陳健波) 莫扎特為什麼會死呢?會不會是因為妒忌而殺死了這個音樂家呢?.
(陳健波) 妒忌是兩個人可以發生的,你有我沒有,所以就是妒忌..
(陳健波) 但妒忌就有一點點醋味,跟客醋有一點點關係..
(陳健波) 那個音樂家其實是信神的,他的情緒不單止是憤怒,他也有妒忌..
(陳健波) 妒忌就是說,為什麼上帝不給我這些才幹,你把這些才幹給了那個人,給了莫扎特..
(陳健波) 這樣他就有一種客醋,所以妒忌是三個人才會出現的..
(陳健波) 因為有一個人,你很想得到他的愛,得到他的恩寵,但他偏偏沒有給你,就給了別人..
(陳健波) 所以就令你有種妒忌的感覺..
(陳健波) 今天跟大家分享的主題是Anthe..
(陳健波) 我選擇了和修版跟大家分享..
(陳健波) 大家讀經就用和合本,我跟大家分享的時候就用和修版..
(陳健波) 我們一起看看,這裡說:神實在,恩代以色列那些清心的人..
(陳健波) 他這樣說,其實就是說,他曾經自己不夠清心..
(陳健波) 可不可以這樣說?一首詩篇的起頭,其實可以說是一個總結,或者一個回顧..

$^{41}$(陳健波) 他覺得自己曾經不夠清心,所以就掉進了一種情緒的困擾裡面..
(陳健波) 今天我會帶大家透過詩人走一圈,就是說一個妒忌的人的內心是怎樣..
(陳健波) 他怎樣才能夠疏解自己的妒忌..
(陳健波) 大家一邊看著聖經,其實這篇詩篇可以說是他自己的一個醒覺,一個心路歷程..
(陳健波) 他起初就覺得自己不夠清心,差點滑跌,差點跌倒..
(陳健波) 特別想跟大家看第三節,因為我嫉妒狂傲的人,我看見惡人享平安..
(陳健波) 如果大家看和修本,他翻譯成心懷不平,如果大家看回秩序表讀經的時候,心懷不平..
(陳健波) 比較準確的就是妒忌,所以和修本是修正了一些字眼上提到妒忌..
(陳健波) 妒忌這個主題我們可以從幾方面去看,第一就是妒忌什麼..
(陳健波) 你對妒忌的對象有什麼反應?我們先看看他妒忌什麼..
(陳健波) 第三節說我妒忌狂傲的人,我看見他們享平安..
(陳健波) 第四節說他們力氣強壯,死的時候沒有疼痛..
(陳健波) 第五節說他們不像一些人會受苦,他們沒有遭災害..
(陳健波) 這是妒忌的內容,就是說這些人生活很平安,個人很強壯,死的時候也不痛..
(陳健波) 還有就是沒有災難淋到他們..
(陳健波) 還有一段是說妒忌的內容,很快跳到第十二節說,看下這就是惡人他們常享安逸..
(陳健波) 常享安逸,就是生活很舒服..
(陳健波) 然後財寶增多,又有錢,生活無憂,又沒有災害,又沒有病痛..
(陳健波) 是否令人羨慕呢?.
(陳健波) 其實在人生裡,我們妒忌的東西也很多,不知道大家有沒有浮出一些妒忌的心情..
(陳健波) 有沒有妒忌澳門人拿的錢比我們多?我們只有五千元,好像少了一點..
(陳健波) 多一點就好了..
(陳健波) 有沒有妒忌有些人神給他很多恩賜,我唱歌就走音,其實這些心情很容易出現..
(陳健波) 看下你妒忌什麼內容..
(陳健波) 好了,我想和大家分享,當我們真的妒忌一個人的時候,其實我們的內心會和那個人有一種距離..
(陳健波) 這裡明明是惡人,沒什麼好說的..
(陳健波) 但如果放在我們日常生活,其實妒忌這種情緒有時候也會破壞我們人與人之間的友情..
(陳健波) 你本來和那個朋友很好,但你見到他風生水起,如此都犯,因為他有你沒有,就酸溜溜..
(陳健波) 如果那些有財富又順風順水的人,如果他們招職的話,我們就更加難堪..
(陳健波) 曬爛,說自己多威風..
(陳健波) 其實一個妒忌的人,我們會對這個妒忌的對象起了一些內心的變化..
(陳健波) 我們有時候會想,他很聰明,恃才在表現高級..
(陳健波) 我不知道你有沒有試過對一些人生氣,可能平常人看他沒什麼,但因為你妒忌他,你會說他不好..
(陳健波) 你以為你很聰明嗎?恃著這些東西嗎?.
(陳健波) 當然我們看對象是惡人,但他的內心也很被惡人糾結..
(陳健波) 我們一起看回第六節..
(陳健波) 你會發覺他描述惡人的情況也用了很多筆墨..
(陳健波) 他用了很多筆墨,表示他內心經常前思後想著這些事情,我們說preoccupy..
(陳健波) 他說驕傲好像鏈子帶在他的頸上,很囂張.殘暴好像衣裳蓋在他身上,很暴躁,對人很差..
(陳健波) 第七節更離譜,他說他的眼睛因肥胖而凸出,這不是妒忌的對象..

$^{81}$(陳健波) 只不過是這個人很壞,不斷吃,吃大了眼睛也凸出,看他不順眼..
(陳健波) 因為妒忌對象的外貌和情況,會令你抗拒或講壞他,用來平復自卑..
(陳健波) 為何他有我沒有呢?.
(陳健波) 內心放任不羈,說他的人際關係,他經常笑人,第八節欺壓人,又自高..
(陳健波) 不單止,他還褻瀆上天,毀謗全地..
(陳健波) 可以是一個惡人的表現,但我覺得也可以是一個我們妒忌對象的緣故..
(陳健波) 我們用了很多筆墨,用了很多心機去想他有多壞,這影響了我們的關係..
(陳健波) 第十節想和大家釋經上的問題..
(陳健波) 如果我們看和合本,第十節說:.
(陳健波) 「所以神的民歸到這裡,喝盡了滿杯的苦水.」.
(陳健波) 說到義人來到這個位置,見到惡人風生水起,義人其實很痛苦,在喝苦水..
(陳健波) 但和修本有一個中性一點的翻譯..
(陳健波) 和修本說:「所以他的百姓歸到這裡,享受滿杯的水.」.
(陳健波) 有一點不同,和合本說苦水,苦水是很難啃的,但這裡用享受的字..
(陳健波) 所以新一點的釋經,釋經採取了這班百姓不是講神的子民,中性一點..
(陳健波) 而是說這些惡人其實有很多人擁戴,擁戴這些惡人的生活方式,令人羨慕,很多人崇拜他..
(陳健波) 都會喜歡他那套,很想喝到他的水,重點在這裡..
(陳健波) 這是第十節的解釋..
(陳健波) 用這個解釋是比較順理成章,因為第十一節開始,他們說,如果是指那批擁戴惡人的人,那些百姓,他們怎麼說呢?.
(陳健波) 怎能曉得,至高者怎會知道呢?看那,這就是惡人,他們常享安逸,財寶增多..
(陳健波) 有一個總結,惡人就是這樣..
(陳健波) 私人用了第十二節到第三節的篇幅,去看他疾度的對象,他的描述,反應,你說他痛不痛苦呢?.
(陳健波) 其實都很痛苦.所以,如果你有一些疾度的對象,其實你都可以想想,你有沒有被這件事影響了你的情緒,困擾,覺得自己很慘,神為什麼對這個人這麼好,對我這麼差,有這樣的感受都未必..
(陳健波) 接著,私人因為這種疾度的緣故,落了一種自憐的狀態,我們一起看第十三節..
(陳健波) 十三節,他說,我實在徒然潔淨了我的心,徒然洗手表明無辜..
(陳健波) 就是說,我做好人幹嘛呢?.
(陳健波) 我落在這樣的處境..
(陳健波) 比如說,我們不信耶穌的人,順風順水,我信耶穌幹嘛呢?守好規條都這麼慘..
(陳健波) 他有一種自憐,覺得自己這麼慘,值不值得呢?.
(陳健波) 為什麼呢?他的自憐是來自,這裡第十四節可以說是有一個對比,因為我終日遭災難,相對那些惡人是無災難的..
(陳健波) 這是一個對比.然後他說,每天早上受情治,其實經常都受到情緒困擾..
(陳健波) 而那些惡人是怎樣呢?常享安日,生活很容易,成為一種很強的對比..
(陳健波) 第十五節,十五節是一個轉機,這一點我覺得值得跟大家分享..
(陳健波) 這裡說,我要這樣說,就是愧對這世代的眾兒女調..
(陳健波) 就是說,如果我這樣說,怎樣說呢?就是說第十三節那番話,我實在徒然..
(陳健波) 就是說,我們基督徒有時候都會這樣,當我們遇到艱難的時候,我們都會慎,是吧?.
(陳健波) 唉,我真是苦了,信來做什麼呢?搞到今天這樣..
(陳健波) 如果你對著一群弟兄姊妹這樣說這些負氣話,其實是會叛倒人的,會影響人的..
(陳健波) 私人有個好處,他突然想起,如果我這樣亂說話,不斷發怨言,自憐,其實是會影響其他眾子女..
(陳健波) 就是說,教會的群體,我講這些負面說話,其實是會影響初信的,或者是人的..

$^{121}$(陳健波) 這種現象,我覺得在神的家裡面,屬神的群體,應該有這種彼此看守的力量..
(陳健波) 但是很可惜,我發覺現在我們的信仰都是來得比較個人主義一點..
(陳健波) 我喜歡這樣就這樣,我想說什麼就說什麼,我們有時候都會發一些負氣的話,說就說..
(陳健波) 沒有想過,我說了之後,原來可能會影響我身邊的弟兄姊妹..
(陳健波) 說就好像理所當然地表達了出來,但是原來其他人怎樣呢?群體怎樣看呢?.
(陳健波) 私人有了一個醒覺之後,就有一個轉化了..
(陳健波) 不過內心還是覺得想不通的,自憐,不想影響弟兄姊妹,但是裡面仍然有一個困擾..
(陳健波) 困擾是什麼呢?第十六節,「我思索要明白者事,眼看實是為難.」.
(陳健波) 想來想去都想不通,即是說現實看到的景象,惡人真的不明白..
(陳健波) 怎樣才能明白呢?第十七節是轉捩點,我們一起讀第十七節吧..
(陳健波) 「直到我進了神的聖所,思想他們的結局.」.
(陳健波) 這個是私人的轉捩點..
(陳健波) 我想大家想想,進入神的聖所,是一個去敬拜神,對準神的時候..
(陳健波) 進入神的聖所其實是很重要的,因為疫情的緣故,有幾年我們都在用Zoom..
(陳健波) 但我自己覺得,若是能夠回教會,會比用Zoom好..
(陳健波) 來到神的聖所,有沒有一種覺得神是聖潔,神聖的情感呢?.
(陳健波) 屬靈上也有種情感,跟你在家坐在沙發上,穿著睡褲,還未換衣服,跟你今天這樣整理自己來朝見神,是不同的..
(陳健波) 所以很鼓勵,正在上網的,若是可以,你也要回來神的聖所..
(陳健波) 因為當你對準神的時候,我們的內心是會有變化的..
(陳健波) 就好像禱告,禱告之前,內心有很多憂慮和掙扎..
(陳健波) 但你進入禱告之後,你會慢慢用神的眼光來看待事情.禱告發揮的作用就是這樣..
(陳健波) 同樣地,我們來到神的聖所,是開始用神的眼光來看待事實的真相..
(陳健波) 事實的真相,第17節最重要的是,不要看一個人,而是看他中間發生了什麼..
(陳健波) 最重要的是看什麼?看結局..
(陳健波) 我們中國人有一句話叫做「蓋棺定論」,就是蓋上棺材就可以做出定論..
(陳健波) 其實我們看到惡人風生水起,只不過我們還不夠長命,看他如何衰收尾..
(陳健波) 一個人的結局和收尾是最重要的..
(陳健波) 例如我們內地也抓了很多貪官,《內蒙內新聞》也有報告..
(陳健波) 其實那些貪官在出事之前是很輝煌的,但被抓之後就像落水狗一樣..
(陳健波) 其實很在乎你看的時候是看哪一段,他人生裡面哪一段..
(陳健波) 如果你看完結局,其實你就不會這麼不服氣..
(陳健波) 結局是怎樣呢?.
(陳健波) 詩人第十八節:.
(陳健波) 「你實在把他們安放在滑地,使他們跌倒滅亡,他們轉眼間成為何等的荒涼,他們被驚恐滅盡.」.
(陳健波) 不知道大家看不看得到那幅圖畫呢?.
(陳健波) 就是說,我們以為那些惡人很順風順水,其實神是把他們放在滑地,在滑他們,其實是香蕉皮..
(陳健波) 這裡有一點吊詭的地方,因為詩篇的起頭,詩人說「我差點失禪,我也差點滑倒」..
(陳健波) 其實有兩個「滑倒」,一是詩人自己差點滑倒,然後這裡說惡人「滑倒」..
(陳健波) 其實惡人「滑倒」其中一個原因,我們很多時候都會有這樣的觀察,.
(陳健波) 惡人為何不悔改呢?是神任憑他們,任憑他們是什麼呢?.

$^{161}$(陳健波) 他們有東西依靠著,他們的生活很平順,很多錢,很順利..
(陳健波) 其實什麼人才會尋求神呢?就是覺得自己不行的時候才會尋求神,是不是?.
(陳健波) 否則,有什麼需要神的救恩呢?.
(陳健波) 其實惡人煽動他們的地方就是他們太順暢,他們的順暢是一個滑地,明白我的意思嗎?.
(陳健波) 但問題在哪裡呢?其實現在詩人經常去羨慕惡人的順暢,明白嗎?.
(陳健波) 其實他們在羨慕,妒忌的東西是香蕉皮,明白我的意思嗎?.
(陳健波) 其實詩人透過來到神的聖所,看到那些鐘鼓,就醒來了..
(陳健波) 其實那些惡人就像神化一場夢,夢醒之後什麼都沒有,那些人什麼都沒有..
(陳健波) They are nothing,沒有的..
(陳健波) 我們為了這些夢幻的東西,搞到自己妒忌,內心有很多掙扎,困苦,值不值得呢?.
(陳健波) 不值得了,所以有了這個醒覺之後,詩人就很痛苦了..
(陳健波) 質度的疏解,心路歷程是很厲害的,起初很羨慕,抹黑那些質度的人,覺得自己很自憐,然後去到神的聖所,有了轉變..
(陳健波) 看到他們的結局,然後就看到自己有多蠢,我真是蠢,我妒忌這些東西,做什麼呢?.
(陳健波) 所以我們一起看第21節,因此我心裡苦惱,肺苦鼻刺,我這樣愚昧無知,覺得自己很蠢,很無知,羨慕香蕉皮..
(陳健波) 善妒惡人的東西,我羨慕,是不是蠢呢?.
(陳健波) 所以他就給自己一種羞愧,很羨慕,覺得自己很羞愧,很痛苦..
(陳健波) 不過,第23節他就抓緊神,然而,我常與你同在,你參乎我的右手,你要以你的訓言引導我,以後你必接我到榮耀裡..
(陳健波) 來到神的聖所,雖然有懊悔,有羞愧,但是對著神,你還可以抓住神,不怕的..
(陳健波) 第25節,待會回應詩歌,我會唱這節經文,這節經文可以說是解決我們妒忌的一個最重要的入手,是什麼呢?.
(陳健波) 我們一起讀25,26兩節,好嗎?開始!.
(陳健波) 除你以外,在天上我有誰呢?除你以外,在地上我也沒有所愛無的.我的肉體和我的心常需,但上帝是你的力量,也是我的福分,直到永遠..
(陳健波) 妒忌是人有我無,有我無,到最後,私人得到的秘訣是,我有了神,其實就無憂了..
(陳健波) 最重要是有了神,耶和華是我的牧者,我必不自缺乏.你有了耶和華是你的牧者,你還擔心什麼呢?.
(陳健波) 我們有沒有以擁有神,有了神,成為我們心裡最大的滿足呢?.
(陳健波) 不需要到處跟別人比較,不容易的,要在屬靈上實踐才行..
(陳健波) 你內心是否真的相信,有了神,萬事都無憂..
(陳健波) 有了神,不等於你跟著順風順水,因為第26節,我覺得精彩的地方在第26節,我的肉體和我的心常需殘..
(陳健波) 即是說,他的肉體和心常的狀態是沒有變的,他仍然是需殘的,但是他的內心已經更新了..
(陳健波) 他的內心的更新,令他有這樣的宣稱,但是神是我心裡的力量,是我的福分..
(陳健波) 每個人都有一份,傳道書教我們,我們滿足於神給我們的福分,滿足於我們擁有的東西,我們就會開心快樂,直到永遠..
(陳健波) 通常在結尾,這首詩篇是群體的敬拜,也是一種宣稱,疏解了自己的內心..
(陳健波) 很想一齊敬拜的人,同樣有這樣的體會,一個什麼體會呢?.
(陳健波) 看那遠離你的必要死亡,犯離棄你的行淫的,你都滅絕.但親近上帝是於我有益,我以主耶和華為我的避難所,好叫我述說你一切的作為..
(陳健波) 親近上帝是於我有益..
(陳健波) 今天我們有沒有羨慕或妒忌身邊其他人呢?如果那種妒忌強到令你憎恨那個人,或者內心有很多糾結,不得不回到神的聖傻,看神是怎樣看待這些人..
(陳健波) 當我們有神的眼光去面對這些情況後,我們會體會到,原來擁有神是最寶貴,有神勝過擁有其他的一切..
(陳健波) 不等於我明天突然間飛黃騰達,順風順水,最大的滿足是來自內心,不要再有那麼多糾結,信神會將最好的東西給我們,我們一起低頭祈禱..
(陳健波) 今天謝謝你讓我們看到,詩人經歷了一個窒道的過山車,透過詩人很細緻地描述她的內心掙扎,也讓我們體會到窒道對我們沒有益處..
(陳健波) 我們很容易會掉進一種自憐和困擾之中,直至我們來親近你,從你的眼光去看事情的時候,我們才豁然開朗,知道你為我們所預備的,是最好的..
(陳健波) 我們知足常樂,謝謝你,我們這樣禱告,鳳蘭你名及..

\newpage



\section{馬太福音 6:25-34}
\label{sec:38XteBMXhDk}
\textbf{2023年5月7日青少主日崇拜直播｜梁煒傳道｜給當代門徒的15個忠告:不要憂慮｜馬太福音6\_25-34}
\newline
\newline
連結: \href{https://youtube.com/watch?v=38XteBMXhDk}{\texttt{https://youtube.com/watch?v=38XteBMXhDk}} ~~~~ 語音日期: 2023-05-08
\newline
\newline
\hyperref[sec:PtF_7l7B_qs]{\small{< < < PREV SERMON < < <}}
~
\hyperref[sec:index_chronic]{\small{[返順時目]}}
~
\hyperref[sec:iCj_rgLUlC4]{\small{> > > NEXT SERMON > > >}}
\newline
\newline
馬太福音 6:25-34
\newline
\begin{longtable}{cl}
\hline
\hline
章節 & 經文 (和合本修訂版)\\
\hline
6:25 & \begin{tabularx}{0.7\textwidth}{X} 「所以，我告訴你們，不要為你們的生命憂慮吃甚麼喝甚麼，或為你們的身體憂慮穿甚麼。生命不勝於飲食嗎？身體不勝於衣裳嗎？ \end{tabularx} \\ \\ \relax
6:26 & \begin{tabularx}{0.7\textwidth}{X} 你們看一看那天上的飛鳥，也不種也不收，也不在倉裡存糧，你們的天父尚且養活牠們。你們不比飛鳥貴重得多嗎？ \end{tabularx} \\ \\ \relax
6:27 & \begin{tabularx}{0.7\textwidth}{X} 你們哪一個能藉著憂慮使壽數多加一刻呢？ \end{tabularx} \\ \\ \relax
6:28 & \begin{tabularx}{0.7\textwidth}{X} 何必為衣裳憂慮呢？你們想一想野地裡的百合花是怎麼長起來的：它也不勞動也不紡線。 \end{tabularx} \\ \\ \relax
6:29 & \begin{tabularx}{0.7\textwidth}{X} 然而我告訴你們，就是所羅門極榮華的時候，他所穿戴的還不如這些花的一朵呢！ \end{tabularx} \\ \\ \relax
6:30 & \begin{tabularx}{0.7\textwidth}{X} 你們這小信的人哪！野地裡的草今天還在，明天就丟在爐裡，神還給它這樣的妝飾，何況你們呢？ \end{tabularx} \\ \\ \relax
6:31 & \begin{tabularx}{0.7\textwidth}{X} 所以，不要憂慮，說：『我們吃甚麼？喝甚麼？穿甚麼？』 \end{tabularx} \\ \\ \relax
6:32 & \begin{tabularx}{0.7\textwidth}{X} 這都是外邦人所求的。你們需要這一切東西，你們的天父都知道。 \end{tabularx} \\ \\ \relax
6:33 & \begin{tabularx}{0.7\textwidth}{X} 你們要先求神的國和他的義，這些東西都要加給你們了。 \end{tabularx} \\ \\ \relax
6:34 & \begin{tabularx}{0.7\textwidth}{X} 所以，不要為明天憂慮，因為明天自有明天的憂慮；一天的難處一天當就夠了。」 \end{tabularx} \\ \\
[1ex]
\hline
\hline
\end{longtable}
$^{1}$《新約聖經》第8頁.
馬大福音第6章25至34節.
請聽我讀出.
所以我告訴你們.
不要為生命憂慮吃什麼 喝什麼.
為身體憂慮穿什麼.
生命不勝於飲食什麼.
身體不勝於衣裳什麼.
你們看那天上的飛鳥.
也不種也不收.
也不積蓄在倉裡.
你們的天父尚且揚活他.
你們不比飛鳥貴重得多麼.
你們哪一個能用思慮使壽數多加一克呢.
何必為衣裳憂慮呢.
你想野地的百合花怎樣長起來.
它也不勞苦也不枉善.
然而我告訴你們.
就是所羅門極榮華的時候.
它所穿戴的還不如這花一朵呢.
你們這小信的人啊.
野地裡的草今天還在.
明天就凋得在爐裡.
神還給它這樣的裝飾.
何況你們呢.
所以不要憂慮說吃什麼.
喝什麼穿什麼.
這都是愛邦人所求的.
你們需要的這一切東西.
你們的天父是知道的.
你們要先求他的國和他的義.
這些東西都要加給你們了.
所以不要為明天憂慮.
因為明天自有明天的憂慮.
一天的難處一天當就夠了.
原本我的港獨是四月尾的那段時間.
不過兩個星期前.
我們家裡一家人就感染了COVID-19.
感恩大人不是很嚴重.
不過家裡的小朋友.

$^{41}$特別是大女兒.
其實連續發了三日瘋.
我猜大家都關心我小女兒.
因為小女兒長期病患.
感恩她在病患過程中.
除了咳嗽之外.
其他都不是很嚴重.
能夠正常地生活.
多謝你們的代禱和問候.
今天是清少主日.
由我在教會服侍了五年的青少年的傳道.
跟大家分享.
今天是我最後一次.
能夠在這裡和大家分享上帝的話語.
六月份的時候.
我們一家人會到英國牛津.
服侍當地的橫教會.
主要是去做一些大學生的科研工作.
其實充滿了很多不同的挑戰和困難.
但當我們回望的時候.
其實發現上帝的恩典.
上帝一直都在這裡.
其實我自己都是一個比較容易焦慮的人.
我記得2018年4月.
我最初來到旺鎮的團隊裡.
我都帶著不同的擔心和憂慮.
我覺得自己對自己沒有什麼大的信心.
雖然是神學畢業.
雖然有一些事奉的經驗.
但都會擔心自己不夠能力.
都會擔心可能時代變遷.
擔心其他人的眼光.
擔心這個團隊能不能夠合作.
除了擔心服侍之外.
因為家裡其實都有小朋友.
都會擔心他們的成長.
會不會因為想著服侍而忽略了他們呢.
中國人有句話叫養兒一百歲長憂九十九.
其實做家長都會擔心小朋友的成長.
人越大要去兼顧的事情越多.

$^{81}$發現很多事情都會很擔心.
特別是面對著未知的將來.
其實心裡有一份不安全感.
這份不安全感令我們很想捉住更多東西.
捉住更多資源.
捉住更多的人.
但當我們越捉得多的時候.
原來我們發現心裡那份憂慮好像沒有消失.
好像心裡每個人的慾望無窮無盡.
當我們越看得多的時候.
反而我們的憂慮都越多.
青少主義特別對於一班青年人來說.
其實成長裡面都充滿著各種迷思和憂慮.
升學就業人際關係.
家人的關係.
拍拖戀愛.
倫理道德.
人生的價值觀.
青春期裡面的新心靈的變化.
其實很多東西都很值得產生憂慮.
各方這幾年的社會裡面.
發生了很多不同的情況.
無論是社會事件或者疫情.
都對年青人.
其實有很多新聞都說.
年青人對將來都充滿著很多憂慮和恐懼.
其實他們都很需要人去跟他們互動.
聆聽他們的需要.
不過因為疫情的緣故.
阻隔了人與人之間的互動.
青年都少了一些學習群體的生活.
有問題有心事.
可能只是在網上.
不論是Google或者IG.
甚至是TikTok.
裡面去找一些答案.
其實年青人很需要.
花大量的時間.
去陪伴他們.
去跟他們一起去分享.

$^{121}$去減輕他們的憂慮.
我相信不單止年青人憂慮這件事.
其實成年人.
甚至是老人家.
我們整個社會.
因為憂慮的原因.
產生了不同的問題.
面對著憂慮這麼大的困難.
我們有時候面對著憂慮.
都會花更多的時間.
會不會想多一點.
會不會周詳多一點.
但明明很簡單的事情.
想得太多的時候.
就會複雜化.
將原本很簡單的事情.
我們將它變得越來越困難.
也會將很多的精神.
放在胡思亂想的時候.
令自己更加害怕.
我分享前兩個月.
我自己考車.
考車的過程中.
前一晚都很緊張.
緊張到睡不著.
所以第二天早上.
一泊車的時候.
就砰.
沒有撞到.
不過從來都沒有試過.
泊車撞到我的肩膊.
不過因為真的太過緊張.
和焦慮.
因為擔心很多事情.
影響查犯不思.
影響情緒.
影響我們的生活.
有很多事情牽掛著.
令我們產生疑惑.
甚至嚴重的.

$^{161}$影響我們的生活.
工作.
我們的學業.
影響我們的身體.
社交的活動.
令我們有壓力.
其實很多病.
都和壓力有關係.
我們的胃.
我們的頭.
都反映著一些狀況.
更何況現代人.
其實睡眠的質素.
都是一般般.
有很多人有失眠的問題.
這些都和憂慮有關係.
作為處境門徒.
究竟我們面對著這些憂慮.
我們和未信者有什麼分別呢?.
我想這段經文.
其實大家應該是聽過的.
不過我想實踐起來的時候.
真的有些挑戰.
但很想繼續透過這段經文.
大家互相去分享.
我們繼續給當代門徒.
15個忠告.
其實今天已經是第12個了.
經文是記載在.
《馬太福音》六章25到34節.
主題就是不要憂慮.
不如我再將這段經文.
分開了兩大部分.
第一個部分就是25到29節.
不如再邀請大英姐妹.
一起讀25到29節的經文.
25.
預備.
1,2,3.
(經文).

$^{201}$(大英).
「我臣望為人體有慮全尋望」.
「生命不勝於人成望」.
「生命不勝於人之常望」.
「你會看眼天上的奇妙」.
「假不中 假不收 假不正宗 作錯了」.
「你會看天空上的妖奸」.
「你會不記奇妙的一種的過往」.
「你會看我能用此類」.
「是受數國家的合力」.
「我願為你以上有慮」.
「你將惹記於你的媽媽」.
「甚至將其女兒」.
「她有多辜負 也不妨事」.
「顯然我高掃門」.
「就使所門即會垮的時刻」.
「她所穿的衣 反而越是她的慮」.
耶穌在登山的裡面.
教導了這班門徒和會眾.
關於天國的價值觀.
雖然他們和我們一樣都是人.
但當我們用不同的眼光時.
同樣面對同一件事.
都有不一樣的心態.
經文中有幾個字眼不斷重複.
當然最重要的字眼就是「憂慮」.
因為這字的原文是說焦慮.
一些無謂的擔心.
特別是其他經文中可能是說思慮煩擾.
特別是路加福音.
思慮煩擾.
特別是路加福音的十章四十一節.
路加福音記載了.
馬太和瑪利亞兩姐妹.
這兩姐妹接待著耶穌.
瑪利亞在耶穌面前專心聽耶穌講道.
馬太事後很多事.
心裡忙亂.
到了耶穌的腳前跟他說.
主啊!我的妹子留下我一個侍候.

$^{241}$你不在意嗎?.
請吩咐她來幫我.
耶穌回答說.
馬太!馬太!.
你為許多事情思慮煩擾.
不是.
我不是故意讀錯.
不過我真的有點困難.
但只不過缺少的只有一件.
瑪利亞已經選擇.
那上好的福分是不能奪去的.
是的!為了很多事情去忙碌.
而忘記了最重要的事情.
外國有位作家形容.
憂慮是什麼呢?.
他說憂慮就像一張搖椅.
它能給你做些什麼.
但它不會給你任何的珍貴.
他說憂慮就像一張搖椅.
當你搖來搖去的時候.
好像很忙.
搖來搖去做了很多事情.
但原來只是原地踏步.
這經文講到三個憂慮.
憂慮吃什麼.
喝什麼.
和囂張什麼.
吃喝.
我們將它歸類成同一類.
而囂張什麼.
而歸納成另一類.
吃喝其實是人生中的生存問題.
特別對於耶穌的年代.
食物和水源.
其實都不是必然的事情.
平常的百姓.
其實他們都餵著兩餐.
他們很努力.
就算他們很努力地耕田.
但都會受天氣的影響.

$^{281}$經常會擔心農作物失收.
經常會擔心.
因為天氣的緣故.
因為有其他動物的緣故.
他們的生活沒有保障.
所以食物對他們來說.
真的很重要.
不過我想對於現代人來說.
吃什麼穿什麼.
未必是我們很擔心的事情.
但可能我們會擔心的.
吃什麼種類的食物.
甚至那些食物.
吃得好不好吃.
對於現代人來說.
吃和喝.
這個其實是生存裡面最重要的.
(咳嗽).
.
我們心靈裡面應該要追求什麼?.
信心的問題.
其實是一個屬靈裡面的需要.
其實這段經文25節之前.
應該是24節.
應該是羅木斯上次的講道.
經文裡面有講到.
"一個人不能夠事奉兩個主".
"不事護者過,愛那個,就是中者過,興那個".
"你們不能又事奉神,又事奉馬門".
其實在說這一群的人.
就是心靈裡面不能夠有兩個主.
當一群外邦人去追求馬門的時候.
他們心靈裡面沒有這個天賦的看顧.
所以他們求的就是他們的憂慮.
吃的,住的,穿的.
但是當我們能夠這一群每一個門徒認識他的人.
我們每一個去祈求上帝的時候.
去追求上帝.
上帝能夠供應我們一切.
叫我們不需要去憂慮.

$^{321}$其實要成為上帝的門徒.
是要對上帝有信心.
相信上帝會供應我們.
特別是當上帝給了我們使命的時候.
無論遇到任何困難和挑戰.
上帝都能夠給我們足夠的恩典.
憂慮難阻我們經歷神德.
憂慮也讓我們自己籌算自己的計劃.
最後忘記了上帝的恩典.
聖經教導我們.
「焉綱一無掛慮,只要凡事藉著禱告和祈求」.
將我們所有的交託給上帝.
上帝所賜的是出人意外的平安.
迫在基督耶穌裡保守你們的心懷意念.
我們每一個服侍青年人的影片.
我們會提供一些場景.
幫助年輕人操練什麼叫信心.
每個星期年輕人都會去負責不同的服務.
敬拜隊,音響.
每個月的時候我們都會走出街頭.
去探訪一些有需要的人.
無論是露宿者,清潔工.
突破信心的界線.
我記得有一次.
我們被一間中學邀請一起去舉辦報道會.
一些大專生負責詩歌敬拜.
分享見證.
我負責去港島分享呼嘲.
年輕人去分享.
其實面對八百多個中學生.
心裡面都有點緊張和怯.
他們都有很多擔心.
但當我們一起去祈禱.
交託給上帝的時候.
上帝給他們勇氣.
見到這班中學生在敬拜過程.
在報道會裡面都很投入.
這班大專生他們說.
他們很正在經歷上帝的同在.
正在經歷上帝給他們力量.

$^{361}$幾年前.
教會都有機會和一班大專生.
和一些不同的弟兄姊妹.
一起有機會去越南的短孫.
其中兩天的時候.
我們有機會去到山區裡面.
去探訪一些少數民族.
當時的天氣都下著一些大雨.
當地的一些當地人.
他們用電單車載我們上山.
你可以想像.
我們有十幾人.
他們用十幾輛電單車載我們.
車程大約是半個小時.
他們開著電單車.
載著我們穿梭在大小的山路裡面.
而山路因為下雨.
石頭崎嶇凹凸.
兼且有些泥路在當中.
其實都挺危險的.
很濕滑.
增加了意外的難度在裡面.
也有些短孫隊員.
不小心滑倒.
全身都是泥.
加上上山的路裡面.
有很多很滑的.
很斜的一個環境在裡面.
所以當中要很大的信心去突破.
雖然我們不能每一個人都去到越南.
不過我相信你們的參與和支持是很重要的.
不知道你們有沒有留意.
外面洗手間旁邊其實有一塊碧波板.
碧波板裡面鼓勵大家一起去參與和支持.
用金錢.
用祈禱.
現在暫時支持神學生的財產.
大約兩萬元.
我們都很希望目標是七萬五千元.
鼓勵大家一起去經歷上帝.

$^{401}$經歷信心的功課.
上帝去保守我們.
上帝去看顧我們.
在不同的場景裡面.
都挑戰我們信心裡面的突破.
經文裡面第三十三節.
你們要先求他的國和他的義.
這些東西都要去加給你們.
不要憂慮吃和喝和穿的問題.
不是說我們不需要為了生命的所有事情.
我們拋諸腦後.
不需要對生命去負責.
反而耶穌教導我們.
我們生命的優先是要為了什麼的籌算呢?.
他說他教導我們要為了神的國和神的義.
去求上帝的國.
去求上帝的公義臨到我們當中.
求一些屬靈裡面的事情.
特別是在馬太福音第六章裡面.
耶穌教導我們的主禱文.
他說.
你們祈禱的時候要這樣說.
我們在天上的父.
願人都尊你的名為聖.
願你的國降臨.
願你的子兒行在地上.
如同行在天上.
為人用的飲食.
今日賜給我們.
為上帝的國去祈求.
為上帝的子去祈求.
耶穌在地上的服侍.
他也會擔心一群失喪的靈魂.
他常常為著這群門徒去憂心.
他每次服侍完之後和服侍之前.
都會去到曠野裡面去祈禱.
將難處和憂慮交給天父.
保羅更加說.
他服侍的時候好像欠了福音裡面的債.
他為了愛國人信主.

$^{441}$他常常去憂心.
他擔心上帝的國.
擔心人得救的事情.
你們要先求上帝的國和上帝的兒.
這些東西都要賜給你們.
這個更加是上帝的一個應許.
當我們這樣去祈求的時候.
這樣去做的時候.
上帝說他自然會看顧我們.
既然我們都有憂慮.
解決憂慮最主要的方法.
最有效的.
當然是尋求上帝的幫助.
上帝告訴我們.
當我們願意尋求上帝的國和他的兒的時候.
他會供應我們的生命.
最後用自己最近的經歷.
和大家分享.
剛才說到這個主日.
其實是我們家庭最後一個的主日.
之後我們就會.
下個六月份就會去到英國牛津.
叫做移民.
移民這個想法其實是什麼時候開始的.
我們家庭裡面.
其實是從小女兒患病開始.
我們開始去為小女兒的事情.
去在不同的地方找很多資料.
特別當時.
其實去年的時候.
我們去申請香港一些特殊學校的環境.
我們發現其實香港要申請特殊學校的資源.
都是挺有限的.
排隊其實最少也要排兩三年.
不過後來很感恩.
因為有一間特殊學校是自行修身.
所以小女兒能夠順利地入讀這間學校.
現在讀得挺開心.
不過我們也會想.
讀完這間特殊學校.

$^{481}$假如一年兩年之後.
下一步會是怎麼樣呢.
我們也會找相關的資料.
我們也發現.
在香港特殊學校裡面的支援.
或者人對特殊的小朋友的看待.
好像也會有很多大家要去適應的地方.
於是乎我們也去想.
會不會外國的特殊教育會比較完善呢.
由去年年中的時候.
我們家庭開始有這個想法.
我們將這個想法告訴同工隊.
告訴執事他們.
除了叫他們祈禱之外.
更加告訴他們有沒有人物網絡.
介紹一下.
搜查一下我們家庭.
我們發現原來移民其實也挺大困難的.
第一個最大的困難.
原來就是經濟的問題.
因為我想太太.
我們兩夫妻.
太太最主要其實是要照顧小女兒的成長.
所以能夠上班的其實只有我一個人.
如果我出去工作的時候.
我去外國的橫教會.
其實相對的待遇.
其實未必有香港那麼理想.
當中我也在想一個問題.
如果我真的要為了小女兒的緣故.
要去移民的時候.
我究竟要做什麼工作呢.
做什麼工作可以維持我們的生活呢.
我要不要轉行呢.
我在大學的時候是讀computer science的.
我在想我要不要轉做IT的工作呢.
或者要不要用我傳福音的口才.
去告訴其他人.
sell(出售)一件產品.
就是要成為一個sales(售貨員).

$^{521}$經過我們夫婦的討論之後.
我覺得上帝的呼召其實也是很清晰的.
我們夫婦仍然是很想很想去服侍上帝的.
之後我們開始去探索海外的工場.
我們去探索加拿大.
探索美國.
探索澳洲.
英國裡面的不同的工場.
我們自己上網去發求職信.
然後開始面試.
面試的過程中.
其實也有些教會也覺得挺友好.
不過透過一些人物網絡.
我們打探到一些薪酬裡面的待遇.
計一計那筆錢.
再計一計當地的人.
就是租金.
然後稅收.
兩個小朋友的教育.
醫療等等.
原來發現都生存不了.
真的只剩下.
扣完之後.
我想我們只剩下兩三千塊.
給我們家裡用.
一個月根本上.
原來如果維持一全人的身份去移民的話.
其實也挺困難的.
不過我們仍然也很想繼續去事奉上帝.
之後下來繼續剩下一間教會.
就是英國牛津這間教會.
他繼續用Zoom跟我們面試.
面試了幾個月.
三四個月.
在當中繼續互相祈禱.
大家都覺得很好.
上帝也是這樣去長大.
用一兩個月的時間大家一起祈禱.
在祈禱的裡面.
神也讓我看到.

$^{561}$原來英國牛津這個地方.
其實也是一個很淑齡豐富的地方.
上帝在當地裡面興起了.
特別大學生興起了很多宣教士.
這些宣教士甚至乎來到中國裡面去傳福音.
我們夫婦也很想能夠有機會服侍這些大學生.
點燃年輕人服侍上帝的心.
之後我們跟教會分享.
我們也很正面.
你知道傳達人的工作通常已經很久了.
然後到了最後才問薪金.
跟外面的工作不同.
外面的工作一開始就問薪金.
但是傳達人的環境就是這樣.
大家的文化就是這樣.
到了最後才會說.
其實多少錢呢?.
我也不是很重視錢的.
不過我也要算一下數的.
是不是?.
就是要算一下.
原來薪金真的不是很高.
不過很奇妙的就是.
原來教會裡面有宿舍.
當有宿舍的時候.
真的解決了我們家庭裡面的問題.
教會的弟兄姊妹也知道我們小女兒的情況.
心裡面也覺得很難過.
很想去幫助我們.
經過一起祈禱.
一起去分享.
他們也願意等候我們大半年.
接下來下個月.
六月的時候.
我們一起去英國服侍.
我們從來沒有想過.
我們家庭會去外地的原因.
我們也沒有想過用Zoom interview.
其實也有點像.
那些人口販賣的KK樂園.

$^{601}$就是不知道去到哪裡.
那個人是什麼.
因為我們家庭連英國也沒去過.
但是我們學習先求上帝的國.
還有上帝的義.
我們相信上帝會公認我們.
我們相信上帝會保守我們.
當我們接受這個offer的時候.
我們發現原來這間宿舍附近.
原來走二十分鐘.
還有一間很好的兒童醫院.
家裡附近也有一間很好的幼稚園和小學.
上帝的公認超出我們所想所求.
當我們願意去尋求祂的時候.
上帝會公認我們.
第34節最後一節就是.
所以不要為明天憂慮.
因為明天自有明天的憂慮.
一天的難處.
一天擔當就夠了.
憂慮的產生.
可能是因為我們對於未知的將來.
經常有很多的估計.
有很多的幻想.
我們花很多的時間去想.
這個事情會不會有很惡劣的情況出現呢.
於是我們害怕.
面對著不肯定的未來.
我們生活在懼怕的裡面.
經文裡面他講到明天自有明天的憂慮.
直譯就是說.
明天會為到明天自己去憂慮.
就是明天的憂慮明天才去處理.
經文更加告訴我們.
一天的難處.
一天擔當就夠了.
他沒有告訴我們.
不用擔心無憂無慮.
而是憂慮必定會有.
每一天都有.

$^{641}$當我們生活的時候一定都會有.
是一定必然的.
不過一天去處理一天的憂慮已經足夠了.
這個憂慮提醒我們.
不單止不要為明天明天去憂慮.
更加告訴我們.
我們不要為過去昨天去憂慮.
可能有時候我們都會很後悔.
活在自責的裡面.
如果過去我有這樣做的時候.
那就好了.
我們會有很多擔心憂慮自責在裡面.
英文就是今天就是Present.
Present這個字其實也可以解作禮物.
神給我們今天.
就好像給我們一份禮物.
叫我們好好活在今天.
當如果我們今天活在這裡的時候.
但我們的心竟然在將來.
竟然活在過去.
其實我們沒有辦法享受今天這一刻.
今天這段經文其實對.
我不知道對你來說大不大提醒.
但是對於我自己.
對於我們家庭來說其實很寶貴.
接下來會有前面有很多的挑戰和困難.
我相信都是在迎接我們.
有很多不同的適應.
加上我們有病患的小朋友.
小朋友未必會這麼快有書讀.
我們有很多的調節在當中.
我們都會有很多的計算.
很多的很怕做錯決定.
影響小朋友的發展和成長.
但是當我們有一天的時候.
我們看著我們的小女孩.
其實平時我們都會很憂心.
哎呀為什麼她說話還是這麼落後.
為什麼她走路還是拐彎抹角.
為什麼她一些技能都不能夠做到.

$^{681}$我們都會擔心其他人對她的眼光和想法.
不過這個小女孩都很可愛.
上帝做她都是很奇妙.
她是一個很喜歡笑的小朋友.
她好像不用擔心的.
她就活在現在.
知道我們和她在一起.
知道我們會陪伴著她.
她心裡面有一份的滿足.
她經常都會笑著告訴我們.
神都很想藉著她去提醒我們.
活好今天.
不要活在憂患的裡面.
如果不是的話.
我們的生命就會浪費了.
和妹妹相處和小女孩相處.
我們覺得不是必然.
其實每一天都是賺回來的.
上帝隨時可以拿走她的生命.
你會不會這樣去珍惜上帝給你的一切.
珍惜上帝給我們的家人.
珍惜上帝給我們所擁有的資源.
所擁有的一切.
我們懷著感恩的心.
細味上帝.
更加相信上帝會去保守.
會去拯救我們.
好好活在今天.
不要被憂慮.
偷走我們和上帝一起經歷的美好時光.
讓我們一起祈禱.
我們承認我們都是很多計算.
很多小信的人.
我們活在自己的憂患的裡面.
我們活在自己的恐懼裡面.
可能在聖經裡面說到.
其實我們和未信者外的人有什麼分別呢.
最後我們去祈求你.
這段經文聽過很多次.
但是有沒有改變我們的心.

$^{721}$最後我們求你.
你的話語是大有能力的.
你的應許是真實的.
最後我們祈求你.
讓我們有一個熟天的眼光.
我們相信你就是這個世界創造偉大的上帝.
你會供應我們生命裡面一切所需要的.
我們想到我們沒有想過的.
你都會去照顧我們.
特別當我們願意去尋求你.
按著你的話語去走的時候.
上帝你會保守我們.
求你幫助我們.
生活已經有很多挑戰,很多困難.
叫我們不要再去加添很多的.
沒有必要的憂慮在裡面.
叫我們每一天活好現在.
每一天經歷上帝你和我們同在.
世美我們的人生.
世美你裡面的甜酸苦辣.
看在我們面前的.
逢耶穌得勝之面.
Amen.
\newpage



\section{撒母耳記上 1:4-28}
\label{sec:iCj_rgLUlC4}
\textbf{2023年5月14日母親節主日崇拜直播｜余嘉敏傳道｜苦媽的最佳選擇｜撒母耳記上1\_4-28}
\newline
\newline
連結: \href{https://youtube.com/watch?v=iCj-rgLUlC4}{\texttt{https://youtube.com/watch?v=iCj-rgLUlC4}} ~~~~ 語音日期: 2023-05-15
\newline
\newline
\hyperref[sec:38XteBMXhDk]{\small{< < < PREV SERMON < < <}}
~
\hyperref[sec:index_chronic]{\small{[返順時目]}}
~
\hyperref[sec:qRN_c71zY24]{\small{> > > NEXT SERMON > > >}}
\newline
\newline
撒母耳記上 1:4-28
\newline
\begin{longtable}{cl}
\hline
\hline
章節 & 經文 (和合本修訂版)\\
\hline
1:4 & \begin{tabularx}{0.7\textwidth}{X} 每逢獻祭的日子，以利加拿把祭肉分給他的妻子毗尼拿和毗尼拿所生的兒女。 \end{tabularx} \\ \\ \relax
1:5 & \begin{tabularx}{0.7\textwidth}{X} 他給哈拿的卻是雙分，因為他愛哈拿。耶和華卻不使哈拿生育。 \end{tabularx} \\ \\ \relax
1:6 & \begin{tabularx}{0.7\textwidth}{X} 她的對頭毗尼拿因耶和華不使哈拿生育，就常常惹她發怒，要使她生氣。 \end{tabularx} \\ \\ \relax
1:7 & \begin{tabularx}{0.7\textwidth}{X} 年年都是如此。每當她上到耶和華殿的時候，毗尼拿就這樣惹她發怒，以致她哭泣不吃飯。 \end{tabularx} \\ \\ \relax
1:8 & \begin{tabularx}{0.7\textwidth}{X} 她丈夫以利加拿對她說：「哈拿，你為何哭泣？為何不吃飯？為何傷心難過呢？有我不比有十個兒子更好嗎？」 \end{tabularx} \\ \\ \relax
1:9 & \begin{tabularx}{0.7\textwidth}{X} 他們在示羅吃喝完了，哈拿就站起來。祭司以利坐在耶和華殿門框旁邊的位子上。 \end{tabularx} \\ \\ \relax
1:10 & \begin{tabularx}{0.7\textwidth}{X} 哈拿心裡愁苦，就痛痛哭泣，向耶和華祈禱。 \end{tabularx} \\ \\ \relax
1:11 & \begin{tabularx}{0.7\textwidth}{X} 她許願說：「萬軍之耶和華啊，你若垂顧你使女的苦情，眷念不忘你的使女，賜你的使女一個子嗣，我必使他終生歸給耶和華，不用剃刀剃他的頭。」 \end{tabularx} \\ \\ \relax
1:12 & \begin{tabularx}{0.7\textwidth}{X} 哈拿在耶和華面前不住地祈禱，以利注意她的嘴。 \end{tabularx} \\ \\ \relax
1:13 & \begin{tabularx}{0.7\textwidth}{X} 哈拿心中默禱，只動嘴唇，聽不到她的聲音，因此以利以為她喝醉了。 \end{tabularx} \\ \\ \relax
1:14 & \begin{tabularx}{0.7\textwidth}{X} 以利對她說：「你要醉到幾時呢？不要再喝酒了！」 \end{tabularx} \\ \\ \relax
1:15 & \begin{tabularx}{0.7\textwidth}{X} 哈拿回答說：「我主啊，不是這樣。我是心裡愁苦的婦人，清酒烈酒都沒有喝，只在耶和華面前傾心吐意。 \end{tabularx} \\ \\ \relax
1:16 & \begin{tabularx}{0.7\textwidth}{X} 不要將你的使女看作不正經的女子。我因極其難過和生氣，所以一直禱告到如今。」 \end{tabularx} \\ \\ \relax
1:17 & \begin{tabularx}{0.7\textwidth}{X} 以利回答說：「平平安安地回去吧。願以色列的神允准你向他所求的！」 \end{tabularx} \\ \\ \relax
1:18 & \begin{tabularx}{0.7\textwidth}{X} 哈拿說：「願你的婢女在你眼前蒙恩。」於是婦人上路，去吃飯，臉上不再帶愁容了。 \end{tabularx} \\ \\ \relax
1:19 & \begin{tabularx}{0.7\textwidth}{X} 他們清早起來，在耶和華面前敬拜，就回去，往拉瑪自己的家裡。以利加拿和妻子哈拿同房，耶和華顧念哈拿。 \end{tabularx} \\ \\ \relax
1:20 & \begin{tabularx}{0.7\textwidth}{X} 時候到了，哈拿懷孕生了一個兒子，給他起名叫撒母耳，說：「這是我從耶和華那裡求來的。」 \end{tabularx} \\ \\ \relax
1:21 & \begin{tabularx}{0.7\textwidth}{X} 以利加拿和他全家都上去，要向耶和華獻年祭和還願祭。 \end{tabularx} \\ \\ \relax
1:22 & \begin{tabularx}{0.7\textwidth}{X} 哈拿卻沒有上去，因為她對丈夫說：「等孩子斷了奶，我就帶他上去朝見耶和華，讓他永遠住在那裡。」 \end{tabularx} \\ \\ \relax
1:23 & \begin{tabularx}{0.7\textwidth}{X} 她丈夫以利加拿對她說：「就照你看為好的去做吧！可以留到兒子斷了奶，願耶和華應驗他的話。」於是婦人留在家裡乳養兒子，直到他斷了奶。 \end{tabularx} \\ \\ \relax
1:24 & \begin{tabularx}{0.7\textwidth}{X} 斷奶之後，她就帶著孩子，連同一頭三歲的公牛，一伊法細麵，一皮袋酒，上示羅耶和華的殿去。那時，孩子還小。 \end{tabularx} \\ \\ \relax
1:25 & \begin{tabularx}{0.7\textwidth}{X} 他們宰了公牛，就領孩子到以利面前。 \end{tabularx} \\ \\ \relax
1:26 & \begin{tabularx}{0.7\textwidth}{X} 婦人說：「我主啊，請容許我說，我向你，我的主起誓，從前在你這裡站著祈求耶和華的那婦人就是我。 \end{tabularx} \\ \\ \relax
1:27 & \begin{tabularx}{0.7\textwidth}{X} 我祈求為要得這孩子，耶和華已將我向他所求的賜給我了。 \end{tabularx} \\ \\ \relax
1:28 & \begin{tabularx}{0.7\textwidth}{X} 所以，我將這孩子獻給耶和華，使他終生歸給耶和華。」他就在那裡敬拜耶和華。 \end{tabularx} \\ \\
[1ex]
\hline
\hline
\end{longtable}
$^{1}$今日選讀的經文是記在《撒母耳記》上一章4至28節.
在《舊約聖經》337頁.
經文是記載一個婦人向神求得兒子的故事.
「耶利加拿每逢獻祭的日子.
將祭肉分給他的妻比利娜和比利娜所生的兒女.
給哈娜的卻是傷份.
因為她愛哈娜.
無奈耶和華不使哈娜生育.
比利娜見耶和華不使哈娜生育.
就作她的對頭.
大大激動她要使她生氣.
每年上到耶和華殿的時候.
耶利加拿都以傷份分給哈娜.
比利娜仍是激動她.
以致她哭泣不吃飯.
她丈夫耶利加拿對她說.
「哈娜,你為何哭泣不吃飯.
心裡受悶呢?.
有我不給十個兒子還好嗎?」.
他們在室內吃喝完了.
哈娜就站起來.
細思以理在耶和華殿的門框旁邊.
坐在自己的位上.
哈娜心裡受苦.
就痛痛哭泣.
祈禱耶和華許願說.
「萬君知耶和華啊.
你若誰顧婢女的苦情.
眷念不忘婢女.
賜我一個兒子.
我必使他終身歸於耶和華.
不用剃頭到剃他的頭」.
哈娜在耶和華面前不住的祈禱.
以理定睛看她的嘴.
原來哈娜心中脈搏.
只動嘴唇.
不出聲音.
因此以理以為她喝醉了.
以理對她說.
「你要醉到幾時呢?.

$^{41}$你不應該喝酒」.
哈娜回答說.
「主啊!不是這樣.
我是心裡受苦的婦人.
清酒濃酒都沒有喝.
但在耶和華面前傾心圖義.
不要將婢女看作不正經的女子.
我因被人激動.
受苦太多.
所以祈求到如今」.
以理說.
「你可以平平安安的回去.
願以色列的神.
允許你向他所求的」.
哈娜說.
「願婢女在你眼前蒙恩」.
於是婦人走去吃飯.
面上再不大受用了.
次日清早.
她們起來.
在耶和華面前敬拜.
就回拉瑪.
到了加里.
以利加拿和妻哈娜同房.
耶和華眷念哈娜.
哈娜就懷孕.
日子滿足.
生了一個兒子.
給他起名叫薩姆爾.
說.
「這是我從耶和華那裡求來的」.
以利加拿和他全家都想試下去.
要向耶和華獻蓮濟.
並還所許的願.
哈娜卻沒有上去.
對丈夫說.
「等孩子斷了奶.
我便帶他上去朝見耶和華.
使他永遠住在那裡」.
她丈夫以利加拿說.

$^{81}$「就隨你的意行吧.
可以等兒子斷了奶.
但願耶和華應驗他的話」.
於是婦人在家裡預養兒子.
直到斷了奶.
既斷了奶.
就把孩子帶上時來.
到了耶和華的殿.
又帶了三隻公牛.
一衣髮細棉.
一皮袋酒.
那時孩子還小.
載了一隻公牛.
就領孩子到以利面前.
婦人說.
「主啊.
我敢在你面前起誓.
從前在你這裡.
站著祈求耶和華的那婦人.
就是我.
我祈求為了得這孩子.
耶和華已將我所求的賜給我了.
所以.
我將這孩子歸於耶和華.
使他終生歸於耶和華」.
接下來有傳道人.
余嘉敏姑娘為我們宣講神的訊息.
「虎媽」的最佳選擇.
交給余姑娘.
余:各位早晨.
我相信在座應該有很多媽媽.
不如舉手讓我們認一下你們.
歡迎.
我們也跟他們說「母親節快樂」.
(笑).
其實我也是一個媽媽.
我有兩個兒子.
除了曾睦.
原來這教會裡.
我就是第二個很珍貴的教睦.

$^{121}$女童工.
是媽媽.
所以今天很幸運.
可以在此分享關於媽媽的心聲.
不如我們一起開始禱告.
親愛的恩主.
你留下很寶貴的話語給我們.
教導我們如何作為你的門徒.
求聖靈你親自賜我們受教的心.
願意讓我們聽完道之後.
我們會行道.
藉著今天的訊息.
更新我們的生命.
求主你幫助我們.
祈禱交託來奉主耶穌基督成名祈求.
阿們.
好,現在呢.
你們知道嗎?.
我的媽媽應該清楚一點.
其實在三月開始.
是爸爸媽媽都挺辛苦的日子.
知道為什麼嗎?.
因為三月到七月.
在各個幼稚園到中學.
他們會舉辦很多簡介會.
然後你們要決定報名的時間.
六月到八月.
更加是小學和中學的面試.
面試的高峰期.
所以回想起那段時間.
都挺辛苦的.
因為壓力很大.
要選擇自己的前路已經不是小事.
何況我們要幫別人去選擇.
加上今天的媽媽.
不單是全職媽媽那麼簡單.
我們是雙職媽媽.
要上班,但又要照顧家庭.
所以我們的身心靈.
其實是年終無憂的.

$^{161}$所以我們都要稱讚.
或慰勞一下身邊的媽媽.
其實爸爸也很勞苦功高.
也很辛苦.
所以我們不知道.
我只是想像而已.
不知道是否母親節放在五月.
然後父親節放在六月.
就是要慰勞一下.
我們爸爸媽媽的辛苦.
對於很多媽媽來說.
我這一隻.
你知道她叫甚麼媽媽嗎?.
就是「搬上岸」的媽媽.
為甚麼?.
因為我大兒子已經上了中學.
去中三.
小兒子已經考進了一間小學.
而且他和我哥哥的學校有關係.
所以我其實不太擔心他上中學.
那些甚麼叫「上岸」的媽媽呢?.
當然就是送子女去讀大學.
出來社會工作.
甚至最好找老婆.
生兒子.
這樣就直接上了內陸.
在內陸確保自己不會水浸眼眉.
做媽媽其實都挺辛苦的.
很多事情要擔心.
以前在沒有那麼多資源的時代.
我們不只是擔心他們的衣食住行.
他們成家立室.
我們就覺得很放鬆.
但今天的媽媽.
除了要擔心這些之外.
她還要擔心子女讀書.
如果能分辨到.
我先考考你們.
分辨到甚麼是官校,進校,直資.
私校,國際學校.

$^{201}$哪一間你能處理好.
其實你已經很厲害了.
我還要看這堆學校裡.
你還要分辨.
甚麼是DSE,IB,GCSE,IGCSE和IAL.
你還要分辨到.
還要知道哪一間中學的公開考試制度.
適合我兒子去讀的話.
你已經是一個超級厲害的媽媽.
其實我很少認識的爸爸會為這些事情很煩.
他經常跟我們說.
不用擔心,可以的.
但我真的很想說.
其實我們之所以擔心.
是因為我們看到到處都是大學生.
我們很擔心他沒有了那張沙紙之後.
他就會失去那種競爭力.
再加上我們沒有辦法預知.
我們的子女是甚麼材料,甚麼質地.
甚麼選擇是最好的.
所以我們漸漸地想為他做更多.
為他預備更好.
其實這份焦慮也令我們很苦惱.
而往往這份焦慮.
也令我們越來越想控制.
越來越做越多.
很想為我們的兒子預備最好.
漸漸地我們就變成了這個.
虎媽.
有時我們不想做那些虎媽.
好像很有目標.
然後在追逐社會上的東西.
我有個經歷就是.
我昔日幫我兒子大兒子時考小學.
我覺得他不適合讀傳統學校.
所以我就幫他報了八間小學.
我覺得八間要面試八次.
其實這件事已經很恐怖了.
但當我跟其他家長說的時候.
我是被家長罵的.

$^{241}$他們跟我說:你報得這麼少.
我們每個至少都13間.
我心裡立即想一想.
是不是我不夠進取呢?.
我的兒子會不會落後於人呢?.
於是我不知不覺.
我從一個很苦的媽媽.
變成一個很「嘩」的虎媽.
其實幫兒子選擇得多.
又會被人說是虎媽.
但如果不計劃的話.
他又會擔心我們寫詞沒有多敗意.
究竟做虎媽.
這隻虎媽「嘩」的虎媽.
是否一定有問題呢?.
作為當代的媽媽.
特別是我們作為基督徒的媽媽.
究竟我們應該怎樣幫下一代.
做最好的選擇呢?.
我想起聖經裡有一個.
苦苦的媽媽.
可以為我們成為借鏡.
就是我們今天所說的.
撒母耳的媽媽哈娜.
她出生於一位.
一千多年前的舊約裡.
生不到孩子的媽媽.
通常大家提起她.
就會想起她很苦.
被人欺負.
又會提起她的信心.
但當我看這段經文的時候.
我不知道為什麼.
我自己說的.
我真的覺得.
在我的眼中.
她好像港媽一樣.
她很有目標.
很有想法.
甚至她身處的處境.

$^{281}$背景和性情.
其實很多地方.
跟我們香港現今的媽媽.
都很有共鳴.
所以不如我們一起.
藉著看哈娜的背景.
她經歷的事.
讓我們成為借鏡.
我們一起看下一段.
第一段.
第四至第八節.
我們看回今天的經文.
如果我們看排序.
其實我們是根據.
七十事譯本去排序的.
就是撒母耳的上下.
其實在路德.
他上一本書是甚麼?.
路得記的,是嗎?.
但原來在希伯來的.
《舊約聖經》裡面.
原來撒母耳的上下.
是列為先前的先知書.
所以他是緊接著士師記.
哈娜的生活.
就是生於士師記的後.
比較末期的時間.
而她所生的孩子.
撒母耳.
其實都是一位士師.
我們今年一起查士師記.
你們是否記得.
士師年代是一個怎樣的年代?.
非常敗壞.
道德,是嗎?.
社會民心.
是完全混沌不安的一個時代.
大家記得.
以色列民在當時的背景.
就是約書亞帶領他們.

$^{321}$去到.
終於定居在迦南地.
結束了很多年.
他們在曠野漂泊的日子.
他們分了家產.
各人.
於是住在迦南的外邦之地.
其實生活比起以前.
是富足了很多.
很安穩的.
但是有一句話.
不是經常說的.
人一安定又怎樣?.
我改了它叫做「一心就起」.
他們很快跟隨著外邦人的生活.
漸漸就落入了.
他們去敬拜偶像的圈套.
他們慢慢遠離神.
過著一些很淫邪的生活.
以致他們.
結果被落入外邦人的手中.
上帝很憐憫.
馬上興起一些士師去拯救他們.
但過了不久.
他們又再犯罪.
然後又落入外邦人的手中.
於是上帝又興起士師.
所以你會看到整卷的士師記.
在無限輪迴地不斷地發生.
在士師記裡面最後那一句是這樣寫的.
「那時以色列中沒有王.
任人……國人任意而行」.
當人們沒有上帝在心裡的時候.
他們甚麼敗壞的事都做得出來.
因此上帝要拯救以色列民.
他興起薩姆爾.
幫他選立君王.
取消士師的制度.
令到這些以色列人進入一個新的時代.
從士師的年代過渡到君王的時代.

$^{361}$薩姆爾的母親哈拿出身於這樣的背景.
一個不太好的年代.
不單這樣.
她自身的家庭也困難重重.
我們一起看看第四至第五節.
她提到哈拿的丈夫伊利加拿有兩個妻子.
每逢憲制的時候.
他會把妻子比利娜和兒女分給他.
但哈拿會分相分.
因為他愛哈拿.
我們會問.
伊利加拿是甚麼人?.
經文說他那麼愛他的妻子.
為何要有兩個妻子?.
我們一起看回前一節.
他說伊利加拿是一個以法連人.
住在以法連的山地裡.
我們看回經文.
其實伊利加拿不是以法連人.
按著歷代字第六章的續譜.
我們發現伊利加拿是利未人.
他還是利美的哥核後代.
他只是一個住在以法連地的利未人.
既然他是利未人.
即是以色列人.
為何要娶兩個妻子?.
在舊約的年代.
的確猶太人和中國人很相似.
傳宗接代對他們來說非常重要.
雖然舊約在一般情況下.
是實行一夫一妻制.
但當妻子不能生育為傳宗接代時.
才可以為自己的婢女.
交給丈夫成為妾侍.
幫你傳宗接代.
就像昔日的撒拉.
她不能生育.
於是亞伯拉罕就娶了夏婚為妾侍.
生了二十萬個女兒.
而且我們看回第二節.

$^{401}$第二節說到.
這裡這樣翻譯.
就說他有兩個妻子.
一名叫哈娜.
一名叫比利娜.
但原文直譯的意思是.
他有兩個妻子.
第一個是名字叫哈娜.
第二個才是比利娜.
所以我們看到.
以利加拿的第一任妻子是誰?.
是哈娜.
所以很有可能.
因為她不能生育.
丈夫就娶了一個伊拉.
為她傳宗接代.
原來現代如果不能生育.
都會不開心.
可能會被人問你為什麼.
但原來當時的婦女不能生育.
對於以色列人來說.
是一個極大的羞辱.
其實已經被人白眼.
背著羞辱的名字已經夠慘.
雖然丈夫每逢特別愛惜她.
給她雙份肉.
表示她對她的尊重和愛意.
但這個哈娜似乎未甘心.
她仍然很仇苦.
我們看第六節提到.
比利娜更加因為哈娜不能生育.
大大激動她,使她生氣.
大大激動是什麼意思呢?.
原來她用不當的字,言語和行為.
故意羞辱哈娜.
令她每年到耶和華的殿時.
哈娜都哭著吃不下飯.
而生氣並非像生氣或不開心般簡單.
而是有不斷發怨言的意思.
所以哈娜生氣得像打雷一樣.

$^{441}$丈夫很愛她.
但丈夫似乎不太明白.
第八節提到.
為何哭泣,不吃飯,心裡仇苦.
其實和合本的譯文已經很仁慈.
把為何縮短成一句.
但如果我們看原文或英文譯本.
會看到三個為何.
其實是三次的反問句去問老婆.
男人們,男人們,我問問你們.
當老婆哭泣,吃不下飯時.
你還用三個反問句去問她.
你們是否想跪玻璃呢?.
明顯地,丈夫不太明白她.
你試想像吧,男人.
如果你發現你最深愛的絕版高達.
突然間破爛了.
老婆便跟你說.
不用擔心,這麼不開心吧.
有我,不是比擁有十個高達好嗎?.
你那一刻想怎樣?.
你可能想掐死她吧.
你覺得自己不明白嗎?.
但仔細想想.
Hannah的丈夫給她雙份肉有甚麼用呢?.
一個女人能吃兩份肉嗎?.
反而更加能幫忙.
令Hannah和她的二奶關係如何?.
會越來越多嫉妒.
而且每年都要凸顯.
「我生不了孩子」.
「所以你要專程彌補我」.
面對這些苦情.
其實她的丈夫.
似乎在意大利做甚麼都會死.
男人應該很明白.
人做甚麼都很難彌補.
我們心中的苦況.
這位虎媽很有智慧.
她知道掌權的神是可以的.

$^{481}$她主動尋求掌權的神.
我們一起看下一段.
面對受苦.
Hannah做了一個非常智慧的選擇.
在第九節的時候.
她說見到Hannah站起來.
其實這是一個非常主動的動詞.
她不是坐著甚麼都不做.
她這麼主動做甚麼呢?.
原來她站起來祈禱.
她去找神.
第十節說Hannah心裡受苦.
痛痛地哭泣.
去禱告耶和華.
她向神許願.
「萬君子耶和華.
你若垂憐誰姑婢女的苦.
眷念她的婢女.
賜給我一個兒子的話.
我就怎樣怎樣」.
現代如果不辱.
我們會怎樣?.
我們第一時間會看中醫.
我們會看醫生有沒有方法.
如何用科學方法解決這個問題.
在Hannah的時代.
雖然沒有這樣的方法.
但仍然有很多方法.
可以幫她解決生育問題.
例如我想到.
如果她找另一個二奶.
是自己的心腹.
之後再幫她生一個.
二奶的地位削弱了.
不保了.
心也涼了.
又或者可以恃著丈夫對她的愛.
所以耍手段.
佔據她的心.
令她用這些方式報復.

$^{521}$還有.
第五節說.
其實耶和華使Hannah不生育.
是易作神關閉了她的子宮.
試想像一個神關閉了子宮.
雖然我們看的時候覺得.
神有神的心意.
但如果你是Hannah的話.
你會怎樣?.
可能最直接的方法就是.
我離開上帝.
去拜其他神.
那時很流行.
拜巴力.
可以幫我求兒.
為何Hannah不做這件事?.
Hannah的確沒有選擇.
其他方法.
即使她知道不讓她生育.
是神當刻的安排.
但她仍然選擇.
主動尋求神的幫助.
因為她深信.
掌管生命的神.
會賜予她生命.
即是孩子.
給她的福祿.
只有她一個可以彌補.
她心靈上的苦.
各位弟兄姊妹.
當我們覺得很無助.
沒有人明白你苦無出路的時候.
我們不一定要用很多人的方法去挽救.
Hannah教會了我們一件很重要的事.
就是她會主動尋求.
讓上帝介入我們的困境之中.
你願不願意相信這個神.
仍然掌權.
祂是一個向你施慈愛的主.
在第十節和第十一節.

$^{561}$我們看到Hannah的禱告.
一點也不簡單.
除了求子之外.
她更加對神說.
如果她真的能生育.
就使她終身歸於耶和華.
不用剃頭都剃他的頭.
這是跟神立約的許願.
根據《民數記》第六章.
如果不是出於祭司家族的猶太人.
想自己的兒子服侍耶和華.
就要在那一段特定的時間許願.
去做拿世爾人.
拿世爾人是甚麼意思呢?.
就是獻身歸主的人.
在那段特定的時間.
他有些事情是不能做的.
不可以喝酒,剃頭髮.
不可以吃鮮葡萄甚至乾葡萄.
更加不能靠近店內的死屍.
Hannah按照舊約的條例.
為兒子許願成為拿世爾人.
但她的許願不是在特定的時間.
在經文中看到是一輩子.
我們試想想.
如果幫助兒女找學校.
我們已經夠頭痛.
何況要幫助他們找.
就是找一份終身的職業.
走那條路更加難.
Hannah的許願就像我們今天.
在你們的兒女還未出世之前.
已經幫助他們決定了.
將來要全職事奉.
然後做神父.
甚至像我們堂主任陳明全牧師.
終身服侍主.
但她並非幫助兒女鋪排一條.
很順利在社會上名成利就的路.
為何Hannah如此肯定.

$^{601}$去做這些抉擇呢?.
我們一起看看下一段.
看第十二節.
第十三節Hannah心中不住的禱告.
其實她旁邊的屬靈領袖.
二利誤解了她喝醉酒.
但我看到Hannah的性情是怎樣.
她不像我們傳統的女人.
沒有地位,不太會說話.
反而她像我們今天的女性.
她站起來為自己發聲.
她對二利說.
你不要當我是不正經的女人.
我只是受苦太多.
所以才祈禱,直到現在.
Hannah勇敢為自己發聲.
令二利明白.
其實她誤解了Hannah.
第十七節二利說.
祝你平平安安.
願上帝聽你的禱告.
英文禱告.
其實祈禱後Hannah的反應是怎樣.
很有趣.
她的臉上再沒有帶著仇容.
還回家吃飯.
其實二利的回應沒有神奇功能.
不要誤解了.
只不過是因為Hannah的心態.
直著禱告後被神改變.
從小到大我們看到Hannah很深信.
上帝是她唯一的倚靠.
她倚靠上帝是最好的福氣.
因此她選擇將兒子獻給神.
你有沒有這份像Hannah一樣對上帝的信心?.
她很相信.
當人虎無出路的時候.
掌權的神仍然可以依靠.
可以信靠.
她也祈求神能夠去眷念她.

$^{641}$眷念的意思是記得她.
希望神記得Hannah所祈求的事.
為Hannah做一些行動.
神真是一個信實的神.
所以她沒有忘記Hannah.
在最後一段.
我們會看到神的行動非常之快.
第十九節.
耶和華很顧念Hannah.
第二天清早回家的時候.
兩夫妻行房.
Hannah就懷孕了.
足月之後她就生了撒母耳.
如果我們看經文.
第21至28.
我們會看到Hannah.
就像香港的媽媽.
很有主見.
也很堅定地做了一些決定.
首先在第20節.
她選擇了兒子的名字.
叫做撒母耳.
意思是.
這是我從耶和華所求回來的.
是神恩典的次語.
她很清楚.
這個兒子是從耶和華而來的.
這個兒子是屬於她的.
因此她很清楚地做了第二個選擇.
就是當她的丈夫想回去獻祭還願的時候.
Hannah就選擇不跟大隊.
她對她的丈夫說.
等到兒子斷奶的時候.
我就會帶他去朝見耶和華.
使他永遠住在那裡.
你試試代入Hannah的世界去想.
這麼辛苦才得到一個兒子.
我不知道我之後有沒有機會再生.
這麼辛苦才除去恥辱.
現在竟然要獻給上帝.

$^{681}$一年只見一次或兩次.
有兒子等於沒有兒子.
一定會很掛念他.
但這個婦孖Hannah不是想擁有這個兒子.
她不是為了留這個兒子在身邊.
去榮耀自己.
或者對付那個姨母.
而是她很想這個兒子永遠住在神的聖殿裡.
專心地服侍上帝.
她很決心.
決定自己用幾年的時間.
如果用當時的文化來說.
斷奶大概是三至六歲左右.
所以到那個時候.
Hannah就選擇繼續堅持她對上帝的承諾.
在第24至25節說到.
她去到侍羅做了一些奉獻.
她說她帶了三隻公牛.
又有一衣的細化面型和一皮的袋酒.
又帶了一隻公牛.
但我們回顧舊約的要求.
摩西哥律法只說要獻一隻公牛就足夠了.
所以我們可以看到.
Hannah獻給神的獻祭非常貴重.
而且是豐富.
可以說Hannah獻上最好的給我們的主.
其實這個婦孖是很值得我們學習.
她的信心是很好的.
更重要的是她為了兒子獻上了最好的選擇.
其實在聖經裡記載了七個不辱的媽媽的例子.
通常都是說關於舊屬的時代的開展.
例如我們會在舊約看到亞伯拉罕的太太薩拉生了二兒子.
新約的時候我們會看到撒迦利亞的太太以利沙白生了斯洗約翰.
都是非常重要的人物.
其實在他們不辱的背後.
神有一個更大的舊屬計劃在當中.
我相信當Hannah為孩子做這個選擇的時候.
她一定沒有想過薩姆爾將會成為以色列歷史上第一個先知.
或最後一個的事師.
而且她將會為神高納兩個很重要的王.

$^{721}$就是蘇羅和大衛.
薩姆爾的一生我們知道他是很倚靠神.
他得到最好的福氣.
而且他的名字一直在聖經裡流傳和記錄下來.
其實有什麼彼得到子女的救恩更好?.
得到最好的福氣?.
我一開始就說過.
我替我的兒子報了八間學校都被人罵.
很可惜我報了八間.
但我兒子一直都沒有學校收.
甚至我們去到派位的時候都失敗了.
所以要求他去購門.
當時我只有一間我心目中覺得好的學校去購門.
其實在裡面都很絕.
在面試前一天.
我兒子當時是讀BSF的茶經班的老師.
在一個機緣巧合的情況下.
他知道我們的情況.
於是他膽小膽大地去向他的老師.
或是向校長面前推薦.
安排我們在第二天面試.
希望我們可以多一間學校選擇.
但很吊詭的.
當我滿心歡喜的時候.
發現這兩間學校都在同一個早上面試.
我心裡想.
我這個香港虎媽會做什麼?.
當然做了很多功夫.
詳細我不說了.
總之做了很多功夫令我能夠一天走兩場.
幫我兒子面試.
以為我安排得非常好.
於是我們先去一間被插入隊的學校.
是老師介紹的那間學校.
進去面試.
怎料我坐在那裡.
等到沒有人喊咪叫我上去.
我越等越生氣.
時間夠了.
我們想趕快去另一間學校面試前.

$^{761}$去行政問發生什麼事.
但行政說.
「不好意思,有些事情弄錯了」.
「你是排在最前的」.
「不好意思」.
立即抓住他們.
帶他們上房間.
我心想還趕得及.
怎料他們進入房間後.
就像現在進入了這麼大的房間裡.
坐著很多個家長.
我們就困在裡面.
很無奈地坐著.
看著電視播放詩歌.
正正因為我坐著聽著詩歌的時候.
我跟上帝有些對話.
上帝讓我反思.
究竟我抓得這麼緊.
發生了什麼事呢?.
我一邊祈禱一邊聽著詩歌.
上帝將我的心改變下來.
我一邊想.
既然我已經放在這裡.
這個機會是上帝賜予我的.
我為什麼要追求這麼多.
要抓這麼多呢?.
於是我就決定留在這間.
在我心目中這麼好的.
比較二線的學校裡面去面試.
過了兩天.
校長就親自說要見我和兒子.
那次是我人生中第一次.
有機會陪兒子一起面試.
我終於體會到發生了什麼事.
原來一個六歲的兒子.
你試想想六歲的兒子.
他面對的壓力非常大.
所以他不敢跟別人有眼神的接觸.
甚至他不敢回答他所認識的東西.
所以他就像傻乎乎的坐在那裡.

$^{801}$面試之後.
我感到很感動.
我覺得勝靈感動我.
於是我跟兒子說.
媽媽終於明白了.
原來過往一年.
你承受著很大的壓力.
無論你這間學校是否收你.
或者你去哪間學校.
在媽媽心目中.
你都是一個最厲害,最勇敢.
做得最好的孩子.
我們在街上就像兩個傻子.
抱在一起哭.
每次說起這件事.
我都很感動.
其實神就像是專門安排了這個面試給我.
讓我親自去看.
甚至是安慰我.
其實當時那一刻.
我一切的勞苦,內疚.
甚至我對上帝的憤怒.
都消失了.
最意外的是這間學校收了我兒子.
後來我聽查經班的老師說.
原來中間有很多小插曲.
我以為我安排了很多.
但原來上帝在後面.
才是精心策劃整件事的人.
完全不是我能力所能做到.
當我回顧過往.
我兒子畢業後的六年.
我發現這間學校是非常重視.
學生福音教育的一間學校.
我兒子更在裡面認識了兩個.
他見過一生很多的基督徒老師.
幫助他在信仰裡扎根成長.
甚至其中一個老師.
在兒子畢業後的今天.
還發訊息祝我母親節快樂.

$^{841}$有什麼彼得上.
兒子得到上天賜給他的最好的福氣好呢?.
所以我們今天在經文裡.
有兩個給當代基督徒媽媽的忠告.
第一個是處身於亂世困局的時候.
我們仍然堅心地尋求神.
其實比對於希臘的年代.
跟我們的年代也有點相似.
我們也有很多內戰.
國與國的戰爭.
我們人心很腐敗.
我們的家庭倫理很扭曲.
甚至我們見到很多異端邪說.
嚴重地挑戰我們和我們下一代的信仰.
作為基督徒的媽媽.
我們會怎樣選擇呢?.
我們會否仍然堅守.
持守真理.
繼續帶我們的子女回到教會當中.
我們的教養會採取積極的態度.
還是被動消極呢?.
我們會否像哈拿一樣.
仍然抓緊上帝.
將我們的困境.
四周充滿變數的環境.
甚至為我們子女計劃的將來.
一一用禱告交給天父.
我們要依靠神.
而不是依靠自己的聰明.
第二個大家要學習的.
就是要放手.
為我們的子女選擇最好的福氣.
其實哈拿教會了我一件很重要的事.
原來一個媽媽的決定.
是影響子女一生的.
薩姆爾一生那麼忠心地事奉主.
其實是由她的媽媽.
幫她做這個決擇的開始.
其實哈拿除了很苦惱苦之外.
其實也挺苦的.

$^{881}$其實她有目標地為薩姆爾.
做了很多很多主導性的決擇.
做苦媽其實也不是一定錯的.
最重要的是我們的目標放在哪裡.
我們是放在神的道路上.
還是其他地方.
我們有沒有好好引導子女的信仰.
向著正確的方向.
飆江直跑.
疫情過去.
我知道很多媽媽.
其實都想幫子女追趕學習進度.
有時甚至放棄了主日學的時間.
幫他們補習.
甚至參加比賽.
其實這種焦慮.
只會慢慢地將我們的子女.
帶得越來越遠離神.
在《箴言》說過.
我們要好好地管教我們的孩子.
使他走當行的道路.
讓他一生都不偏離正道.
所以放手不是什麼都不做.
讓子女自由發揮.
不是.
而是像哈拿一樣.
她經常提醒自己.
這個兒子.
我們的子女.
是神所賜的.
是屬於神的.
所以我們可以放心交託給上主.
我們會怎樣選擇呢?.
我們一起同心圖告.
親愛的主.
我們要高舉讚美你.
因為你賜給我們生兒育女的福分.
又賜給我們有家庭.
讓我們經歷那份溫暖和保護.
願你在今天母親節的日子.

$^{921}$我特別求你.
去紀念我們在在的每一個母親.
紀念每一個母親的勞苦.
特別在現今的世代裡.
去教養孩子一點都不容易.
很容易被世俗的歪理.
很多歪曲了的價值觀所影響.
你藉著哈拿去教曉我們.
我們要為我們的孩子選擇最好的.
是最好的上帝賜予我們的福氣.
我求你幫助我們.
去學習去認定你是我們生命的救主.
我們可以放心將我們的子女交託給你.
使他們走當行.
在上帝你的道路裡.
一生都不偏離主你的道.
求主你幫助我們.
祈禱交託乃奉主耶穌基督的聖名祈求.
阿們.
\newpage



\section{馬太福音 13:1-9}
\label{sec:qRN_c71zY24}
\textbf{2023年5月28日崇拜直播｜張培良傳道｜成為不合理的撒種人｜馬太福音13\_1-9}
\newline
\newline
連結: \href{https://youtube.com/watch?v=qRN-c71zY24}{\texttt{https://youtube.com/watch?v=qRN-c71zY24}} ~~~~ 語音日期: 2023-05-29
\newline
\newline
\hyperref[sec:iCj_rgLUlC4]{\small{< < < PREV SERMON < < <}}
~
\hyperref[sec:index_chronic]{\small{[返順時目]}}
~
\hyperref[sec:bwknUBKFdWs]{\small{> > > NEXT SERMON > > >}}
\newline
\newline
馬太福音 13:1-9
\newline
\begin{longtable}{cl}
\hline
\hline
章節 & 經文 (和合本修訂版)\\
\hline
13:1 & \begin{tabularx}{0.7\textwidth}{X} 就在那天，耶穌從房子裡出來，坐在海邊。 \end{tabularx} \\ \\ \relax
13:2 & \begin{tabularx}{0.7\textwidth}{X} 有一大群人到他那裡聚集，他只好上船坐下，眾人都站在岸上。 \end{tabularx} \\ \\ \relax
13:3 & \begin{tabularx}{0.7\textwidth}{X} 他用比喻對他們講了許多話。他說：「有一個撒種的出去撒種。 \end{tabularx} \\ \\ \relax
13:4 & \begin{tabularx}{0.7\textwidth}{X} 他撒的時候，有的落在路旁，飛鳥來把它們吃掉了。 \end{tabularx} \\ \\ \relax
13:5 & \begin{tabularx}{0.7\textwidth}{X} 有的落在土淺的石頭地上，因為土不深，很快就長出苗來， \end{tabularx} \\ \\ \relax
13:6 & \begin{tabularx}{0.7\textwidth}{X} 太陽出來一曬，因為沒有根就枯乾了。 \end{tabularx} \\ \\ \relax
13:7 & \begin{tabularx}{0.7\textwidth}{X} 有的落在荊棘裡，荊棘長起來，把它擠住了。 \end{tabularx} \\ \\ \relax
13:8 & \begin{tabularx}{0.7\textwidth}{X} 又有的落在好土裡，就結出果實，有一百倍的，有六十倍的，有三十倍的。 \end{tabularx} \\ \\ \relax
13:9 & \begin{tabularx}{0.7\textwidth}{X} 有耳的，就應當聽！」 \end{tabularx} \\ \\
[1ex]
\hline
\hline
\end{longtable}
$^{1}$讓我們繼續以聆聽聖道的敬拜.
接下來是經訓時間.
今天選讀的經文是在新約聖經《馬太福音》.
新約聖經《馬太福音》第十三章一至九節.
請聽我讀出.
當那一天耶穌從房子裡出來.
坐在海邊.
有許多人到他那裡聚集.
他只得上船坐下.
眾人都站到岸上.
他用比喻對他們講許多道理.
說:.
有一個殺種的出去殺種.
殺的時候有落在路旁的.
飛了來吃盡了.
有落在土淺石頭地上的.
土既不深發苗最快.
日頭出來一篩.
因為沒有根就枯乾了.
有落在荊棘裡的.
荊棘長起來把它擠住了.
又有落在土土裡的.
就結實.
有一百倍的.
有六十倍的.
有三十倍的.
有耳可聽的.
就應當聽.
各位我是張培良傳道.
來自澳門百甲巢浸信會.
今天我是第一次在舞會上講道.
所以對我來說是有點緊張.
但又有點興奮.
看到大家有一些熟悉的面孔.
可能有些可能令我看起來不太清楚.
但不要緊.
因為我們今天來到.
我想跟大家分享神的恩典.
我記得我自己.
算上我讀神學到回教會服侍.

$^{41}$前後大概是踏入第九年.
差不多十年的時間.
十年後我回到這裡.
我覺得很感恩也很感動.
也看到舞會不斷發展.
不斷鼓勵我們.
所以今天來到這裡.
我也希望代表教會.
再次祝福弟兄姊妹生命健壯.
繼續在上帝的過路中一同努力.
今天我們分堂在澳門.
希望藉著年青人.
剛剛大家聽到他們大敬拜.
我們也不斷學習成長當中.
其實是給上帝的.
如果不是上帝.
他們是沒有今天的.
你看到帶領敬拜的弟兄.
比我更加像傳道人.
看聖經好像比我更加像傳道人.
教會是這樣的.
是很有盼望的.
我們今天來到當中.
雖然我知道今天有很多傳道同工.
今天不在當中.
有很多不同的事奉.
很多不同的事情.
在這麼多年來.
陪著百甲巢浸信會的成長不離不棄.
特別當中有幾位的執事.
牧者傳道.
今天在莊中也有.
其實真的很辛苦.
不斷港澳兩邊走.
陪著我們長大.
我自己經常覺得.
當時也不是很意會到那種艱辛是多少.
當我自己在服侍的當中.
有機會出外.
或者有機會在不同地方服侍的時候.

$^{81}$我深深體會到他們的辛勞是多麼的大.
那一程傳道也很要命的.
你知道交通上也是很浪費精神的.
但是他們每一次來到.
都帶來很多的教導.
鼓勵,供應.
帶來很多母會弟兄姊妹的祝福.
就是在座的我們認識的.
包括羅牧師.
有些執事我們以前一起同工會開過會.
包括甚至在音響房裡的一些.
我認識的弟兄.
我們以前曾經在一些宣教的事工上.
有些弟兄姊妹我們一起.
在當中參與過.
這些弟兄姊妹讓我們的信仰更加豐富.
讓我們更加.
或者包括我自己.
更加明白應該怎樣來服侍上帝.
就是在座的弟兄姊妹你們每一位.
因為你們很厲害.
你們甚至連.
可能有些當中我不知道.
有些可能連百合巢浸信會都沒去過.
連那個百合巢浸信會的張培良傳道的樣子.
是什麼都沒見過.
連百合巢弟兄姊妹的樣子你都不認識.
但你竟然憑著上帝的信心.
你們繼續的奉獻.
繼續的為著教會來祈禱.
如果不是大家這樣來用信心擺上.
你們就不會見到我們今天在當中.
教會的成長.
和不斷的來見證上帝.
我們中間雖然不認識.
但我想我們有一件事是有關連的.
就是什麼呢.
就是我們都是因為上帝的國度.
因為天國的福音.
所以我們今天坐在這裡.

$^{121}$所以我剛剛也跟敬拜隊說.
我們帶敬拜的時候.
可能面對的弟兄姊妹不太熟悉.
不要緊的.
看上去有些生寶也不要緊.
你記得.
那個就是主裡面的弟兄姊妹就行了.
當我們一敬拜.
身份背景的東西就會飛走了.
我們就是主裡面的弟兄姊妹.
所以不需要緊張.
我們就一起去敬拜.
今天我們當說到.
如果我們中間的連繫.
就是天國的福音.
就是耶穌基督的福音的話.
所以我今天正正要來分享的就是.
耶穌基督的福音.
天國的比喻.
很想跟大家分享這段經文.
讓我們一起在主裡面.
再一次明白這個我們很熟悉.
可能已經聽了很多次的比喻.
到底今天還可以給什麼提醒和鼓勵我們呢.
在開始之前我想帶大家做一個簡單的祈禱.
我們祈禱.
主你幫助我們去聽到你的聲音.
是因為你的福音 是因為你的救贖.
今天我們一同可以這樣敬拜.
我們讚美你.
願你讓我們繼續聆聽你的聖言.
也讓我們今天.
可以再一次在你的聖言當中去回應你.
帶領以下的時間打開我們的心.
我們同心禱告.
逢禮基督 得勝名至極 阿們.
今天的主題叫做.
成為一個不合理的殺種的人.
我一會兒會說到底有多不合理.
首先我想問大家一個問題.

$^{161}$進入經文之前.
我不知道在座的弟兄姊妹.
包括白鴿巢.
有沒有想過.
你有沒有期待過.
教會要成為一個怎樣的教會.
有沒有想過的.
當我在教會聚會的時候.
有沒有想過這間教會要變成怎樣的呢.
有些人可能說我要變成一間千人的教會.
我要變成一間很大的教會.
然後熱切傳福音.
非常好.
有些人說不是.
我們不一定很大.
但是我們要成為一間很落地.
去做到一些服務社會.
服務基層的一間教會.
有些教會不斷在做這些工作.
我知道母會也做了很多這些.
有些人可能覺得我們教會不理大小.
總之我們要成為一間.
時常都可以讚美到神的一間教會.
每一間教會都不同.
我不知道旺角浸宿的弟兄姊妹.
你們是怎樣期待.
但不要只是期待.
如果你要你的教會變成怎樣呢.
你怎樣才可以變成.
可以看到那間教會出現.
就是由你開始做的那件事.
如果你想教會成為一間傳福音教會.
怎樣呢.
你成為一個傳福音的人.
那你教會就會慢慢成為一間傳福音的教會.
如果你想你的教會可以有很多的服務基層.
可以服務社會.
那怎樣呢.
由你開始去服務社會.
由你開始起步.

$^{201}$那教會就能夠去成就.
那我自己說一下我們教會.
我都可能.
他們都聽一下.
我很希望百甲巢浸宿會成為一間.
不是說大小.
而是說無論這個世界這個環境.
怎樣變化都好.
在不同的環境下都可以成為一間.
不斷去傳揚福音的教會.
為何我會這樣說呢.
因為百甲巢就是一間這樣來成立的教會.
是因為有很多人不計得失.
不計後果.
甚至不計付出.
而來去供應.
去當中服侍.
才有今天百甲巢浸宿會.
來到今時今日這個日子.
所以我希望這個種子.
能夠繼續在教會當中生根發芽.
無論將來的社會變化怎樣.
我們不知道.
但是無論怎樣.
得時不得時.
總要傳主的福音.
這是我自己的期望.
當然今天有很多人.
對於教會有不同的期望.
不過今天在這段福音裡面.
我們看到耶穌基督正正就是來跟我們說.
到底天國是怎樣來開展的呢.
到底天國是怎樣來發大的呢.
或者這麼說.
在經文之前我們看到.
這一段經文我們很熟悉的.
我相信在座很多人都聽過.
第一次讀.
我們通常注意的位置在哪裡.
通常聽到100倍.

$^{241}$60倍.
30倍.
當然是最正的這裡.
我希望自己教會100倍.
希望教會60倍.
或最少也30倍.
或者有些人的注意力會放在.
不同的土壤那裡.
然後問一下我自己.
我是不是很肚餓呢.
有沒有讓上帝的種子.
在我心裡面生根發芽呢.
也有時候我們會這樣問的.
如果覺得自己現在是一些荊棘.
拜託幫我除去那些荊棘.
讓我們生命成長得再好一點.
但今天我想由一個第三個的角度.
來看這段經文.
這段經文裡面有一個人.
一開始說到有一個殺種的人.
這個殺種的人是怎樣的呢.
他有一些很不合理的一種殺種的方式.
首先他說這個人出去殺種.
他殺的時候.
有一些就殺在路旁.
有一些就殺在土踐石頭地.
有一些就殺在荊棘裡面.
你這樣聽就覺得.
那就殺在那裡吧.
有什麼特別呢.
但你想一下.
如果一個正常的殺種的人.
他希望有什麼呢.
他希望有收成.
他希望種子是發芽生根的.
如果一個正常的人.
我不知道在座當中有沒有家裡有養殖物.
甚至以前有耕種過.
我不懂得耕種的.
但我聽那些人說.

$^{281}$如果你想種子發芽發得好.
首先你要找一些好的土壤.
你首先要翻土.
然後你又要除去泥土的雜質.
還有現在有些很厲害的肥料.
可以放在泥土上.
令到那些泥土.
不知道那些什麼酸鹼值.
是很適合栽種的.
我不認識.
我是聽有一個姐妹跟我分享.
哇,這麼複雜.
種東西而已.
她說是的.
如果你想那些植物發芽發得好.
是有很多工序要做的.
不是這麼簡單的.
我去想一想.
如果一個正常的農業社會.
很清楚耕種的需要.
他們這麼多工序.
那麼這個殺種的人.
為什麼會這麼隨意地殺呢.
殺路旁,荊棘旁,石頭旁.
難道他不著重收成嗎.
不可能的.
殺種的人一定著重收成.
但今天我們看到馬太.
當他描述.
我們一路看聖經.
整個十三章有很多的比喻.
去描述天國的發展.
你發覺.
他告訴我們.
上帝好像用了一個.
跟我們所理解的不太一樣的思維.
去講這個故事.
這個也是我們每一個人.
我們每一次來看耶穌基督的比喻.
我們一定要常常戴著這副眼鏡.

$^{321}$耶穌基督是很厲害的.
用一些大人小孩都聽得懂的比喻.
但是他在中間.
要反轉你平常的思維.
在這裡我們首先看到什麼呢.
我今天講的第一點就是.
這個殺種的人.
你發覺他不是不重視收成.
因為沒有可能的.
沒有一個殺種的人不重視收成.
但為什麼他會這樣殺種呢.
是因為這個殺種的人.
很明顯他對於他的種子非常有信心.
他相信這些種子.
只要落到一個很土的地方.
就會發生什麼事呢.
就會有一個爆炸性的成長.
一百倍,六十倍,三十倍.
他對他的種子非常有信心.
在這裡他告訴我們.
那顆種子代表著什麼呢.
天國的福音.
這個殺種的人代表著天國的福音.
他殺去種.
但這個殺種的人是誰呢.
很多人呢.
剛剛我們看了比喻經文講.
有很多人圍著耶穌.
但是如果你看經文的第十節.
或者你程序裡面應該都有.
你會發覺原來這一段經文.
不是全部講給所有人聽的.
是講給誰聽的呢.
是講給門徒聽的.
因為那些門徒當耶穌講完比喻之後.
就來抓耶穌一邊.
然後問耶穌.
耶穌啊,為什麼講這麼難明白的比喻呢.
那耶穌說什麼呢.
這些比喻這麼難明白.

$^{361}$就是為了讓他們聽不明白.
讓你們聽得明白.
然後耶穌就跟他們解釋.
整個比喻是什麼意思.
耶穌沒有向門徒隱瞞.
他要講的東西.
所以今天這個比喻.
這個關於天國的開展.
是講給那些門徒聽的.
我知道今年.
王浸的一個主題叫做門徒不斷.
按我理解.
在座的每一位都是門徒.
我們要不斷建造門徒.
門徒不是只是說.
我相信耶穌的問答.
而是說一個真真實實用生命去跟隨耶穌的人.
今天這一段經文.
就是講給這些門徒聽的.
這些門徒帶著天國的福音.
不斷的來傳播.
而天國的福音.
卻是有著爆炸性.
有著那種你想像不到的能力.
可以展現出來.
所以這個殺種的人.
他去到哪裡都殺.
不是因為他不理那個修成.
而是因為他相信福音的種子帶著能力.
只要你願意為主去傳.
總有一天.
你不知道什麼時候.
那顆福音的種子.
會有一天下到好土.
然後那顆好土.
就會突然間一百倍.
六十倍.
三十倍的成長.
然後當你看到那些門徒.
可能甚至你沒有想過.

$^{401}$你沒有想過原來在那裡.
上帝這樣來工作.
你沒有想過.
在這裡我想起一個故事.
跟我自己的經歷有關.
我記得當時我在大學團契.
我當時有上大學團契.
其實我那時候剛剛信主沒多久.
剛剛上了教會什麼都不懂.
但是因為當時我換屆.
所以進去之後我做了職員.
我記得當時剛剛換屆.
所以人數不是很多.
我記得當時有一個弟兄.
那個弟兄是一個運動健將.
他真的經常跑步.
你知道我不知道是不是運動健將的人.
性格都比較直男一點.
那弟兄進來團契.
第一天就說什麼.
我想認識女孩子.
這裡有沒有女孩子.
哇 這麼直接的弟兄.
他當時還沒有信主.
他就說聽說教女孩子挺好的.
想進來團契找女孩子.
弟兄你是不是找錯門口了.
你去旁邊那裡找吧.
你去聯誼的地方吧.
這個弟兄很有趣 很搞笑.
我經常覺得他很直.
做運動的直男.
我跟他聊天的過程當中.
有一次有機會.
當時的團長邀請我.
就說這個弟兄剛剛來的.
你不如就跟他做一些.
當時我來沒多久.
你帶他去信主吧.
或者跟他分享福音吧.

$^{441}$沒所謂吧.
團長叫了我就帶他去.
我說弟兄除了女朋友.
我們還有耶穌.
不如我就介紹耶穌給你認識吧.
弟兄就說好啊.
他覺得好像一些入會儀式.
聽我說完一輪初信的介紹之後.
我說弟兄你願不願意相信耶穌.
弟兄就說.
相信耶穌之後.
是不是可以認識女孩子.
我心裡面一沉再躁.
弟兄就說你真的搞不定.
上帝幫到你.
就在這個時候.
他願意相信.
那一刻我帶著他去做了這個祈禱.
這個弟兄就一直回到團契.
但是很遺憾地.
當時我們因為換屆.
當時我們團契是沒有女孩子的.
全部回來都是男孩子.
這個弟兄當然很沮喪.
他回來之後.
他覺得這裡可能是一個騙局.
明明這裡有很多女孩子.
又說好女孩子.
誰知道全部都是男孩子.
搞什麼呢.
我記得有一次.
那一次團契.
不知為什麼.
只有我們兩個回來.
你眼看我眼.
你說我們真的很想多些女孩子.
當時我還沒有拍拖.
我都說那真的需要女孩子.
團契不是為了追女孩子.
但都需要姐妹的.

$^{481}$我說不如弟兄這樣.
我們一起祈禱.
這個弟兄剛剛順著母奶.
你記得吧.
然後他竟然和我一起.
就在團契室裡.
慰勞團契接下來多些姐妹.
他很懇切地祈禱.
他真的很懇切地祈禱.
他不是很會祈禱.
煮多些女孩子.
阿們.
當然我再補充.
我再豐富一點.
但奇妙的事情就在這裡發生.
在那一年年尾結束的時候.
我們團契裡多了三倍的女孩子.
最後只剩下我們兩個男孩子.
這個弟兄到今天還在說這件事.
他說.
他和我說什麼.
他說我以後不敢亂祈禱.
他說以前想多些女孩子.
現在多到不知道怎麼選擇.
他說上帝真的很信實.
但神奇的事情不止這件事.
神奇的事情是這位弟兄.
這位來追女孩子的運動直男.
最後被上帝呼召了.
去了新加坡讀神學.
我早兩年的時候.
我很感恩.
因為我收到他寄給我的一封信.
不是,是網上Facebook寄過來的.
那封信是什麼呢.
不是問候我.
也不是說一些寒暄的話.
那封信是他自己在大學寫的一篇論文.
論文的題目是關於.
希伯來文其中一隻字的分析.

$^{521}$可能現在這樣說.
你不知道我當時多麼震撼.
我沒有想過這個弟兄.
這個為了追女孩子而來到這裡.
懇切祈禱的弟兄.
他今天竟然.
他開始學習聖經原文.
他去到新加坡準備讀公道博士.
他最後成為了一個服侍上帝的人.
在舊約裡面不斷在研讀.
我說上帝這個不是你的工作.
有什麼可能成就.
當上帝的種子落到對的土壤裡面的時候.
那種爆炸力是你想不到的.
今天我跟弟兄說.
我說弟兄.
你快點有機會多教我的原文.
我現在搞不定原文.
你多點跟我分享.
上帝在你生命當中的工作.
他現在很開心.
在新加坡找了老婆.
收了小朋友.
在新加坡定居.
繼續在當中服侍.
每一次我上帝兄.
我就想起這篇經文裡面.
那個殺種人.
為什麼能夠對方的種子這麼有信心.
是因為上帝真的讓我們看到.
原來上帝的道.
是真的有能力的.
我相信今天在旺朕裡面.
你們也經歷過這樣的經歷.
可能當初還沒見到.
但是有一天你總會見到.
但是在美好的畫面裡面.
先不要太快開心.
因為這段經文.
不只是100倍60倍30倍.

$^{561}$你記得那三種土壤.
還有三種卡住的土壤.
當經文來講這三種土壤的時候.
你發覺在經文的.
如果你有聖經.
你看看第11節到第17節裡面.
他引用了一段以賽亞書的經文.
這段以賽亞書的經文.
講的背景是什麼呢.
就是以賽亞書六章的故事.
以賽亞書六章的故事.
是上帝呼召先知以賽亞的故事.
如果你還記得上帝呼召以賽亞是什麼故事.
是很有趣的.
當上帝說我在這裡.
我可以妻獻誰呢.
以賽亞很厲害.
以下立刻衝出來.
妻獻我,我才是姐女.
請妻獻我.
如果今天我們教會.
陳牧師出來.
我們可以妻獻誰呢.
第一節節目一起舉手.
請妻獻我.
陳牧師應該很開心.
耶穌也很開心.
但接著說什麼呢.
陳牧師或者那段六章說什麼呢.
上帝說什麼呢.
好啊,你參與服事,沒問題.
我告訴你.
你接著服事的那一班是什麼人呢.
就是那些聽完又不會做.
做都不會做得好.
還會埋怨你的人.
哦,是嗎.
原來是這樣的.
原來服事是這樣的.
那要到什麼時候呢.

$^{601}$不是很久.
做到香港淪陷為止.
六章差不多是這個意思.
做到房屋荒蕪.
是嗎.
哇,原來.
當以賽亞這麼興奮.
來到說主啊,我才是姐女.
請妻獻我的時候.
上帝這樣一盤冷水潑過來.
不是應該說很憤慨嗎.
應該說,行了,你做就去吧.
不是啊,原來這裡.
經文作者或者聖經都很真實.
告訴你原來在服事裡面.
為什麼要將100倍60倍30倍放在最後.
因為你遇到的通常都是那些石頭啊.
淺肚啊.
荊棘啊.
說完又不聽.
聽完又不做.
做完又不認.
認完又不理.
理完下一次又不會再來.
原來不斷經歷著這些狀況.
所以當經文去引用這些土壤的時候.
你要看到100倍60倍30倍.
但同樣也會看到.
原來你現在會面對很多困難.
面對很多荊棘.
我相信在.
黃浸的裡面.
當.
黃浸在當初.
十幾年前.
來到澳門,你們帶著宣教的心智.
來到澳門這個地方做宣教.
經過輾轉.
接納了.
接受了澳門白頭陣信會的服侍.

$^{641}$我相信歷史在有當中.
很多定律會比我更加熟悉.
但我在想.
在當初十幾年前.
當舞會的牧者,當舞會的執事.
看著我們十幾年前.
還是失鬥窿.
還在流鼻涕的學生.
什麼都不懂的時候.
他是不是已經看到我們今天會.
站在這裡跟大家分享聖言.
是不是看到我們今天會跟大家大領敬拜.
是看到我們開始像基督徒.
一點點的樣子.
是不是已經看到這個畫面.
所以他們做這個工作呢.
我相信不是.
其實如果用人的角度.
我想說.
我們當時這班失鬥窿.
即是小孩子.
如果再壞一點.
早幾年可能要交死者在前面.
這班小孩子在面前.
其實說真的我們又沒有奉獻.
我們也不是很擅長服侍.
在教會就是有破壞.
沒有建設.
在服侍的時候面對這班年青人.
當然我們可以很理想地告訴他們.
不是的,年青人將來.
很有盼望的.
我們也會這樣說的.
但你現在眼前看著我們,那個樣子不是那個樣子.
你看著我們說話.
還在質疑你,是嗎?.
是這樣的嗎?信得是這樣的嗎?.
即是.
你明白嗎?我體會到.
當初舞會的牧者.

$^{681}$他們絕對不是因為.
已經看到100倍,60倍,30倍.
所以他們繼續這樣來服侍.
而是因為他們相信上帝.
福音種子的能力可以改變人.
同樣的,今天.
我想說,我自己.
有機會回教會服侍.
我服侍一班,進到學校服侍一些中學生.
那些中一中二的那些.
我想說,我真的嘔血的.
因為我現在進到學校.
我覺得我自己好像跟那些話題.
已經完全連結不到.
我不知道你們對著年青人會不會這樣.
我不知道,那些潮語我不明白.
好像暗傳的,我又不知道.
是嗎?是嗎?.
笑著說,他們覺得很好笑.
不知道說什麼,一直都有笑的.
不知道,這麼多笑點,但我不知道他們笑什麼.
只能說,是啊,耶穌都很愛你.
或者有時候.
這次出現,很開心,很乖.
下一次,突然間,一班人不見了.
食物就預備在這裡.
然後那一班人不見了,那怎麼辦?我吃光他們.
那你開始明白為什麼傳道人的身形.
開始變化了,是嗎?.
就是這些服侍裡面的經歷.
今天我看著這班.
中一中二的學生.
我真的可能能夠體會到.
當初,舞會執事.
或者他們看著我們那種感受.
哇,真的,如果不是相信.
上帝那種福音的能力.
的改變,我真的.
很想棄你們於不顧.
我不如找一些好像更加.

$^{721}$有能力,或者好像.
更加做到我想做的.
那些人.
但上帝在這裡,用這一段經文.
這些不同的土壤,告訴那班門徒.
是,你會遇到這些情況.
你會遇到這些荊棘.
你會遇到一些令你沮喪的時間.
但請你不要忘記.
上帝福音的種子是有能力的.
當你今天.
你願意繼續為上帝.
來奔走,去擺上.
的時候,你會知道.
有一天,你會.
看到那個修成.
我也是這樣相信,來在繼續當中.
的服侍,所以.
以賽亞書裡面這一段.
的經文,請你不要忘記.
我剛剛說的故事.
六章,聽起來好像有點悲情.
但以賽亞書不是只有六章.
以賽亞書有六十六章.
去到六十六章,告訴我們.
什麼?原來去到最後,以賽亞.
不是失敗,以賽亞.
是來成就了.
上帝的心意,或者.
整本以賽亞書的主角.
根本不是以賽亞.
根本是誰?是上帝.
告訴我們,原來.
人的歷史,人的故事.
會起起跌跌,會有成功.
有失敗,是我們.
掌管不到的,即是喜福.
但最終,上帝.
在這裡,上帝掌權.
弟子們,今天,我相信.

$^{761}$我想在香港當中,我們都面對很多.
高高低低的喜福.
特別經過疫情.
經過很多.
香港這幾年發生的事.
我相信香港都有很多要面對的艱辛.
教會都有很多的挑戰.
這些我們控制不到.
但我們可以.
控制的是什麼?繼續對上帝有信心.
繼續的來到堅定在主的裡當中.
繼續的去做我們應該要做的工作.
繼續的去殺那些.
上帝天國福音的種子.
繼續的去殺那些上帝天國福音的種子.
在那些你殺到的地方當中.
第二樣.
我想跟大家分享的.
殺種的人,當他殺完種之後.
他不是殺完就算的.
他要等,你種的東西是知道的,要等的.
等的過程是不知道要等多久.
當然有些人可能你對那個植物.
越來越認識,可能你會掌握到.
大概多久,但通常你都掌握不到.
可能因為天氣.
可能因為你給的養分等等.
你不知道什麼時候才會看到.
發芽成長,在這個過程當中.
作為殺種的門徒.
我們需要來到堅持.
去忍耐,以至有一天.
我們才能夠看到.
那個一百倍,六十倍.
或者三十倍的收成.
堅持忍耐,我相信.
對在座可能.
是比較容易一點,對於新一代.
我特別在說我現在面對的.
Form1,Form2的學生.

$^{801}$堅持忍耐,多忍兩秒.
對他們來說.
是一個災難.
這個功課是很難的.
當然我自己.
年輕的時候,我都覺得忍耐是一個難的功課.
要我坐在這裡.
聽道聽三四十分鐘.
其實是一個很大的操練.
所以我很明白.
當中有時候我進學校.
看到學生睡覺,我很諒解.
我說主啊,你在夢中對他說話.
你有你的方法.
我很明白,但是.
上帝在教會中不斷操練的時候.
你多了教會,你會發覺.
你的耐性多了,因為聽久了.
你會發覺三十分鐘,四十分鐘.
OK,還未睡著.
還可以忍耐的.
開始我們的耐性,聽上帝的耐力是不是多了.
所以今天在我們生命裡.
做上帝的工作.
我們不可以想一下就立刻到位.
我們都有那個堅忍.
有那個操練,有那個成長的功課.
你和我今天都是.
無論什麼年紀,都是.
我們都要繼續的持守堅忍.
第三點.
我們分享的就是.
經文有一個很重要的畫面.
就是由小到大.
你會發覺不只是說種子.
裡面有說麵糰.
麵粉,或者一些.
很小的東西慢慢變成大.
這個其實只是說.
上帝的道.

$^{841}$聽起來好像很小的東西.
好像很簡單.
只是分享耶穌基督的十字架就行了.
原來正正就是.
透過很小的事.
分享福音的種子.
它是要帶來深遠的一些影響.
上帝的道.
今天很多人都覺得很微小.
你還說耶穌.
你還說耶穌基督的十字架.
今天都說元宇宙.
還說這些.
但今天在經文裡.
我們看到原來上帝的道.
就好像一顆種子.
埋在土裡的時候.
它所成長出來的.
是你無法估計的.
而且這裡說的大和小.
不只是說數字上的大和小.
經文在31至32節.
就說那些種子.
長成一棵大樹.
那些大樹可以讓飛鳥.
在當中安居.
提供了一個什麼呢?.
幫助和安慰.
原來今天教會我們看到的那種變大.
不是說白鴿巢.
或者黃浸變成更大的教會.
不是人數變大了.
是說這間教會.
成為更加有影響力.
可以帶來安慰.
帶來社區改變的一間教會.
這才是最重要.
那種福音的改變.
所以今天.
弟兄姊妹我們要讓教會.

$^{881}$成為一間有影響力的教會.
不單止數字要有.
不斷的變化.
而是我們的生命.
要讓人看到.
那種改變和影響力.
當我們這樣做的時候.
雖然會遇到經濟.
會遇到沙石.
會遇到那些可能不太生長的地方.
甚至有時候我們會.
沮喪可能會卡住.
但是你相信有一天.
當你不放棄的時候.
你會看到100倍.
60倍 30倍.
我整理一下.
今天我們所聽的內容.
這個殺種的人.
就是那個門徒.
他用一種很不合理的方式.
或者對於當時來說很不合理的方式.
為什麼會這樣做.
因為他相信那顆種子本身的潛力.
相信福音本身帶著能力.
這個殺種的人.
甚至沒有除去所有的雜質.
沒有將種子放在.
所謂最好的環境當中.
他只是無論去到哪裡.
他就殺到哪裡.
無論他去到哪裡.
他就不斷的繼續殺種.
因為他相信有一天.
這顆種子總會去到.
好土的當中.
然後殺種的人可以經歷.
100倍 60倍 30倍的收成.
所以整個殺種的過程.
不是說我們人到底有多精心的策劃.

$^{921}$不是說我們有多世界最好的計劃.
或者當然我們在教會裡面.
我們都會說計劃.
我們都會說一些籌備.
但不只是這些策劃籌備.
令到那件事成就.
是什麼令到這件事成就呢.
是福音本身的能力令到這件事成就.
但是當然.
如果你要讓這顆種子能夠成長得好.
那個農夫.
那個殺種的人用心的栽種需不需要.
當然需要.
那個農夫每一天來去擦旱.
努力擺上.
去保護那顆種子的行動需不需要.
今天同樣在教會.
上帝的福音帶著能力.
但是都需要門徒.
悉心的擺上.
努力的擺上.
讓福音的種子.
可以能夠不斷的擴大.
所以我會說.
我們苦幹不一定成功.
但是成功我們必須要苦幹.
要擺上.
你不能夠想著坐在這裡.
就等上帝的事成就.
去到最後結尾.
我想分享一段關於Facebook的故事.
或者我們請勁拜隊可以有些預備.
Facebook我想很多人都有.
或者我不知道在座有多少人有Facebook.
Facebook如果你知道故事.
是有一個人.
他叫做Sukerberg.
這個人是他的總裁.
為什麼現在會這麼流行.
甚至成為一個很出名的社交媒體.

$^{961}$其實當初.
他是一個和同學之間開玩笑.
平時你們讀書中學有沒有寫紀念冊.
大家臨畢業.
幫我寫個名字.
貼張相.
其實當時他們Facebook開始就是這樣.
在學校有些所謂紀念冊的東西.
經常說不如將這些放上網.
看上去好像比較好.
而且方便看.
然後同學說試試.
自己圍內幾個人試試.
試多幾個.
朋友說幾個好像浪費了.
不如連同隔壁班一起玩.
隔壁班好像挺好玩.
多了人參加.
連同隔壁學校一起玩.
幾間美國大學一起之後.
慢慢慢慢他就開始退學.
然後就找人集資.
然後就成就了整個Facebook的發展.
去到今時今日.
發展到今天已經是.
不知道多少億了.
我想用這個故事告訴大家.
其實Facebook是一個很簡單的概念.
開始就成為了今天這麼大的影響力.
等於我們福音難度不比Facebook更加強大.
可能只是我們今天.
來相信福音的能力.
可能是一個很簡單的意念.
可能是一個很簡單的嘗試.
你不會知道他朝.
在黃浸的工作.
或者在你的生命當中.
上帝可以怎樣來使用.
所以我相信我們福音的能力比Facebook更強.
所以Facebook要成就一個全球連繫的網絡.

$^{1001}$今日我們的福音.
都要成就一個全球連繫的網絡.
讓每個人都能夠聽聞到福音.
最後我想分享教會.
都要講一講我們教會.
最近拍了一張照片.
我們教會有很多年青人.
當初我們教會.
當舞會的牧者執事來到的時候.
我們教會當時只有十多人.
但今日教會.
我們還是很虧欠.
因為我們今天可能是十倍左右.
十倍多的增長.
我們教會現在有百多人.
但我覺得不要看輕我們只是百多人.
因為如果你有機會去到澳門.
你問一問一些基督徒.
在澳門裡面.
做年青人教會的哪一間教會.
我相信基本上有九成的教會都能夠說出.
就是百甲巢浸信會.
很多年青人一來到澳門.
他說有些年青人介紹來澳門.
介紹去百甲巢.
我們不知道的.
其實我們不知道這件事.
後來外面教會的人告訴我們.
但我講這件事不是想講.
有多驕傲我們做得多成功.
不是,我們百多人.
有什麼值得驕傲.
離三十倍六十倍一百倍還差很遠.
但我想講這件事.
是上帝的能力.
彰顯在教會弟兄姊妹裡面.
讓百甲巢有一間.
基本上準備倒閉式微的教會.
來到今天成為一間有活力.
可以有上帝的道.

$^{1041}$有真理.
甚至繼續祝福這一代年青人的教會.
弟兄姊妹這件事.
全部都是上帝福音的能力.
所以我很想藉著今天.
分享我們教會.
亦都分享這段經文.
鼓勵大家.
未來的日子.
還有很多荊棘.
沙土.
很多這些挑戰.
但不要只看衝擊.
那些難阻.
你先看看我們.
看看我們.
你們的付出.
是有收成的.
然後你再看自己的教會.
現在做的事.
可能是難阻.
可能有挑戰.
亦都有難關.
但當你用上帝的愛.
當你用上帝的福音繼續去做的時候.
你會有一天見到一百倍.
六十倍.
三十倍的收成.
讓我們一起祈禱.
欣主多謝你.
賜下了美好的福音給我們.
讓我們知道如何來跟從你.
今天是上帝你自己的愛的揀選.
讓我們一班不配的人.
能夠成為你的門徒.
甚至我們可以成為一個傳揚福音的人.
大家知道傳揚福音裡面.
有很多的挑戰.
有很多的艱難.
有很多我們自己的困難.

$^{1081}$外在環境的困難.
但求你幫助我們.
成為一個得時不得時都傳揚福音的人.
成為一個面對經濟.
面對難阻.
都繼續傳揚福音的人.
因為我們對福音要有信心.
不是我們自己的能力有多大.
而是福音本身就帶著能力.
我們相信當我們這樣來到.
無論環境如何變化.
我們都繼續傳揚的時候.
有一天這些種子會去到好土的裡面.
他們會結成三十倍.
六十倍.
一百倍.
然後當我們回望的時候.
我們心裡面就有大大的感恩和讚美.
是上帝你的愛.
讓我們可以經歷這麼奇妙.
這麼祝福的畫面.
願主你祝福我們舞會.
旺角浸訴會的眾弟兄姊妹.
讓我們繼續在我們門徒的生命當中.
為你擺上.
為你踏出那一步.
為你繼續的堅忍.
為你繼續的有盼望.
願主你使我們眾人.
將每個人都擺上.
逢難而獲得的明智祈求.
Amen.
讓我們一起用詩歌.
來回應我們的主.
讓上帝的愛.
來到我們的生命裡面.
也不單止是在我們生命裡面.
而是感染我們身邊的人.
來活出上帝的愛.
\newpage



\section{馬太福音 7:13-23}
\label{sec:bwknUBKFdWs}
\textbf{2023年6月4日主餐崇拜直播｜吳真真牧師｜給當代門徒的15個忠告:走窄路｜馬太福音7\_13-23}
\newline
\newline
連結: \href{https://youtube.com/watch?v=bwknUBKFdWs}{\texttt{https://youtube.com/watch?v=bwknUBKFdWs}} ~~~~ 語音日期: 2023-06-05
\newline
\newline
\hyperref[sec:qRN_c71zY24]{\small{< < < PREV SERMON < < <}}
~
\hyperref[sec:index_chronic]{\small{[返順時目]}}
~
\hyperref[sec:UnHN8_G394I]{\small{> > > NEXT SERMON > > >}}
\newline
\newline
馬太福音 7:13-23
\newline
\begin{longtable}{cl}
\hline
\hline
章節 & 經文 (和合本修訂版)\\
\hline
7:13 & \begin{tabularx}{0.7\textwidth}{X} 「你們要進窄門。因為通往滅亡的門是寬的，路是大的，進去的人也多； \end{tabularx} \\ \\ \relax
7:14 & \begin{tabularx}{0.7\textwidth}{X} 通往生命的門是窄的，路是小的，找到的人也少。」 \end{tabularx} \\ \\ \relax
7:15 & \begin{tabularx}{0.7\textwidth}{X} 「你們要防備假先知。他們到你們這裡來，外面披著羊皮，裡面卻是殘暴的狼。 \end{tabularx} \\ \\ \relax
7:16 & \begin{tabularx}{0.7\textwidth}{X} 豈能在荊棘上摘葡萄呢？豈能在蒺藜裡摘無花果呢？憑著他們的果子，就可以認出他們來。 \end{tabularx} \\ \\ \relax
7:17 & \begin{tabularx}{0.7\textwidth}{X} 這樣，凡好樹都結好果子，而壞樹結壞果子。 \end{tabularx} \\ \\ \relax
7:18 & \begin{tabularx}{0.7\textwidth}{X} 好樹不能結壞果子，壞樹也不能結好果子。 \end{tabularx} \\ \\ \relax
7:19 & \begin{tabularx}{0.7\textwidth}{X} 凡不結好果子的樹就砍下來，丟在火裡。 \end{tabularx} \\ \\ \relax
7:20 & \begin{tabularx}{0.7\textwidth}{X} 所以，憑著他們的果子就可以認出他們來。」 \end{tabularx} \\ \\ \relax
7:21 & \begin{tabularx}{0.7\textwidth}{X} 「不是每一個稱呼我『主啊，主啊』的人都能進天國；惟有遵行我天父旨意的人才能進去。 \end{tabularx} \\ \\ \relax
7:22 & \begin{tabularx}{0.7\textwidth}{X} 在那日必有許多人對我說：『主啊，主啊，我們不是奉你的名傳道，奉你的名趕鬼，奉你的名行許多異能嗎？』 \end{tabularx} \\ \\ \relax
7:23 & \begin{tabularx}{0.7\textwidth}{X} 我要向他們宣告：『我從來不認識你們，你們這些作惡的人，給我走開！』」 \end{tabularx} \\ \\
[1ex]
\hline
\hline
\end{longtable}
$^{1}$也是經訓的時間.
今天所讀的經文是記載在.
新約《馬太福音》第七章十三至二十三節.
是在本堂的聖經.
唐庸聖經是在新約的九頁和十頁.
請聽我讀出.
「你們要進閘門,因為引渡滅亡,那門是寬的,路是大的,.
進去的人也多,引渡永生那門是窄的,路是少的,.
找著的人也少.你們要防備,假先知,.
他們到你們這裡來,外面披著羊皮,裡面卻是殘暴的狼..
憑著他們的果子,就可以認出他們來..
荊棘上豈能摘葡萄?.
窒泥裡豈能摘無花果?.
這樣,凡好樹都結好果子,.
唯獨壞樹結壞果子..
好樹不能結壞果子,壞樹不能結好果子,.
凡不結好果子的樹,就撼下來,凋在火裡..
所以,憑著他們的果子,就可以認出他們來..
凡稱呼我主呀,主呀的人,不能都進天國,.
唯獨遵行我天父旨意的人,才能進去..
當那日必有許多人對我說,主呀,我們不是奉你的名傳道,.
奉你的名趕鬼,奉你的名行許多義能,.
我就明明白白地告訴他們說,.
我從來不認識你們,你們這些作惡的人,離開我去吧..
經文獨筆,願神賜福上讀他話語並且進行的人..
各位弟兄姊妹,今天我們要講的經文,.
記載在《馬太福音》七章十三至二十三節..
今天已經來到我們由一月開始的講道主題,.
耶穌對門徒的十五個忠告的第幾講?.
十三還是十四?十四,很清楚,是第十四講..
今天我們的主題是「走澤路」..
前幾天我看報紙,看到有一個名人,.
他正在準備出一本書,.
是講關於對他五十句對他影響心軟的話..
其中一句是這樣說的,.
他說他的爸爸教他,.
如果你迷茫的時候,就要選擇一條最難的路,.
因為就算不成功,你從中也會學到很多,.
以及可以成長..
這一句話對於那位名人來說,.

$^{41}$心裡產生了極大的迴響,.
他自己不單單常常思考,.
在他的生活中他也常常應用,.
甚至他的兒子也是這樣教導..
這句話是不是很有道理?.
這句話是不是很有智慧?.
但是我今天要和你們講的訊息,.
耶穌告訴我們要走澤路,.
這一句話同樣會幫助我們成長,.
也會幫助我們有很多學習,.
但是不同的是,.
選擇難走的路不是在迷茫的時候才選擇,.
而且這條路是一定成功的..
我們一起來看看,.
究竟耶穌所說的選擇難走的路,.
走澤路,走小路,究竟是怎樣?.
讓我們一起祈禱..
親愛的主,幫助我們,.
在人生的路上我們認定你,.
我們選擇一條走向永生的路,.
求你幫助我們,.
打開我們的心去認定你的話,.
以致我們人生的路上沒有選錯,沒有走錯,.
我們這樣禱告交托,奉主明求,阿們..
上次第十三個忠告,.
是Ray牧師和我們分享,.
七章第十二節說,.
所以無論何時,你們想要人怎樣待你們,.
你們都要怎樣待人,.
因為這就是律法和先知的道理,.
這可以說是登山寶訓主體的一個總結,.
耶穌一開始告訴我們,.
祂不是要廢除律法,.
祂是要承傳律法和先知,.
而待人如己,.
以恩慈去待人,.
就是這個律法的總結,.
是我們踐行律法的目標,.
耶穌希望祂的門徒能夠竭力地去執行,.
到了第十三至二十七節,.

$^{81}$祂會鄭重地向他們警告,.
我用警告這個字,不是用忠告,.
是因為在這十三至二十七節的教導,.
帶來的後果是嚴重的,.
所以其實這個段落就是告訴我們,.
我們聽了這麼多登山寶訓裡,.
這麼寶貴的教訓,.
律法和先知的教訓之後,.
我們應該如何去回應?.
首先我們要看看,.
其實這裡是呼籲我們,.
呼籲我們怎樣可以走進天國?.
在十三至二十七節這一連串的教導裡,.
總共有五個圖像,或者我們稱為五個比喻,.
裡面都有一個對比,.
都會讓群眾看到,其實是有兩個選擇,.
你會看到這些選擇裡面是很分明,.
很對立,.
而且是彰顯了究竟一個門徒和一個非門徒的特徵,.
你會看到有哪幾個比喻?.
有兩道門,.
門徒的特徵就是,.
他們走的會是進閘門,走小路,.
非門徒的特徵就是,.
他們會走寬門,走大道,.
然後他又有兩種動物的比喻,.
雖然這裡沒有說到門徒是羊,.
不過他暗示了其實門徒是羊,.
然後那群非門徒是狼,.
而且他們披著羊皮,.
然後有兩種果樹,.
在第七章十六至二十節,.
他提到兩種果樹,門徒的特徵就是一棵好的樹,.
而且他們會結好果子,.
但是非門徒的特徵就是,.
他們是一棵壞的樹,會結壞果子,.
另外我們會看到一個對比,.
就是門徒會進行父的旨意,.
但是非門徒只是在口裡稱呼耶穌為主,.
而沒有進行到天父的旨意,.

$^{121}$到二十四至二十七節,.
就是用兩個蓋房者的比喻,.
究竟那個房子是建造在盤石還是沙土上?.
其實這種對比對猶太人來說一點也不陌生,.
你馬上會想起詩篇第一篇,.
「不寵惡人的計謀,.
不站惡人的道路,.
不坐舌萬人的座位.」.
然後「唯有喜愛耶和華的律法,.
便為有福.」.
是一個對比,.
然後詩篇第一章第六節就說,.
「因為耶和華知道義人的道路,.
而惡人的道路必自滅亡.」.
是一個這樣的對比,.
告訴你有兩個選擇,.
而這兩個選擇的結局是甚麼?.
我們看看這個圖表,.
這個圖表就是將這兩個比喻,.
在我們一會要說的段落裡,.
不同的比喻,不同的圖像,.
兩種的人,門徒和非門徒,.
他們最終走到的結局是甚麼?.
兩道門的結局是,.
門徒的結局是去永生,.
兩棵樹門徒的結局是去永生,.
兩個稱呼門徒的結局是可以進天國,.
蓋房的比喻是,.
門徒是經得住大水的衝擊,.
如果不是真正的門徒,.
他們的結局是怎樣?.
兩道門的比喻就是,.
他們會去滅亡,.
壞樹的比喻,.
那群人最終會去審判,.
單單稱呼耶穌主,.
而不遵行祂的旨意的人,.
主耶穌說:我根本不認識你,.
請你離開,.
見在沙土上大水會直接衝走你,.

$^{161}$這些對比在在表達一個結果,.
結果是關乎生和死,.
它是挑戰我們,.
到底你在天國是有份,.
還是走向滅亡?.
它告訴我們,.
如果我們是一個門徒,.
我們的生命是走向永生的道路上,.
有一天你和我都要面對死亡,.
我們的生命最終究竟可以上天堂,.
可以和上帝永遠在一起,.
在一個沒有疾病,沒有痛苦,.
沒有死亡,沒有眼淚的地方,.
還是你要選擇永遠和主隔絕,.
朝至滅亡的地方?.
通向永生的道路,.
不僅是死後才開始,.
而是一開始我就選擇了跟從耶穌,.
就是去進閘門,.
就是去走閘路,.
這條路上,.
你會經歷很多恩典和豐盛,.
因為耶穌說:.
「我來是要叫人得生命,.
並且得的更豐盛.」.
耶穌在這裡清清楚楚地展現了,.
層層地推進了人的選擇,.
會帶來什麼樣的結局?.
弟兄姊妹,.
難道你想在審判的日子裡,.
不讓基督承認,.
讓祂離開我們,.
要將我們像枝枝般撼下來,.
永遠和祂分開,.
盡到滅亡嗎?.
所以我會說,.
第十三節開始,.
到二十七節,.
是一個很嚴厲的警告,.
祂告訴你,.

$^{201}$有兩條路,.
但事實上,.
你應該是沒有選擇的,.
我們門徒要選擇的,.
很清楚,.
就是走這條閘路,.
進閘門,.
走閘路,.
但我們也很清楚知道,.
第十三至第十四節,.
祂告訴我們,.
走這條閘路,.
是很艱苦的,.
讓我們一起讀出,.
第十三至第十四節,.
你們要進閘門,.
因為引致滅亡,.
那門是寬的,.
路是大的,.
進去的人也很多,.
引到永生,.
那門是窄的,.
路是小的,.
找著的人也少,.
雖然耶穌的警告,.
是這麼清楚,.
我們應該會選擇,.
去永生那條路,.
但是弟兄姊妹,.
我不知道當你聽完,.
之前的十三個忠告,.
你覺得怎樣?.
馬太很清楚,.
這樣指明登山寶訓裡,.
耶穌說,.
我們要作門徒的時候,.
我們要走真理和教訓,.
我覺得一點也不容易,.
講了這麼多堂,.
有多少的吩咐,.

$^{241}$有多少的教導,.
我們回去認真去想,.
怎樣信服耶穌,.
跟從祂實踐,.
第十三節,.
一開始,.
他講到進閘門,.
進閘門,.
他說要進閘門,.
這時是一個命令的動詞,.
他講得很清楚,.
我們選擇的路,.
一定是窄的那條,.
一定是小的路,.
其實相對地來說,.
寬闊的路,.
是容易走很多的,.
寬闊的路,.
是可以讓你不用怎樣看到界線,.
可以讓你找到很多走棧,.
就是跟從這個世界的價值去走,.
所以是容易很多的,.
舉個例子,.
如果你是一個售貨員,.
跟著這個世界的價值去走,.
你的目標就是賺錢,.
你會按著世界的處事方式去行,.
在工作上,.
你可能有時要說謊,.
踩著別人上位,.
然後你就得到你所得,.
而不需要跟從耶穌的真理,.
很容易你就會達到目的,.
又或者如果你是老闆,.
你為了希望你的下屬可以達到目標,.
你可以用這個世界的方法,.
用威權去鎮壓他,.
你不跟上帝的吩咐,.
不跟耶穌的吩咐,.
去告訴我們,.

$^{281}$我們為首的,.
做領袖的,.
我們就是做僕人,.
的確,.
有些人說,.
路是小的,.
其實小的意思是什麼?.
小的意思,.
原文是被逼迫,.
被擠壓,.
它是由一個動詞「壓」,.
來衍生出來的,.
它的同源詞就是「逼迫」和「苦難」,.
所以其實現在說的路是小的,.
就是說我們會因為這些困難,.
而帶來我們很多的苦惱,.
新漢語譯本譯做路是小的,.
它譯做路是多麼崎嶇,.
走窄路是很崎嶇,.
是很艱難,.
其實窄路是什麼?.
窄路是一條順服上帝的道路,.
但是提摩太后書三章十六節說得很清楚,.
如果你立志在耶穌基督裡面敬田度日的話,.
你就會受到逼迫,.
很多時候門徒要展現愛心,.
很容易,.
也很受歡迎,.
但是如果門徒要指出世界的不是,.
活出和世界不同的價值觀的時候,.
很可能是不受歡迎,.
而且甚至要受苦,.
這個時代,.
我們面對世界價值觀的挑戰和衝擊,.
去做合神心意的事,.
有時候反而會被譏諷被欺凌,.
我認識一個小六的學生,.
一個女孩子,.
她是一個很有愛心,.
很樂於幫助人的小基督徒,.

$^{321}$她因為見到班裡面有一個同學,.
平時的校服比較舊,.
不整齊,.
沒什麼人和她玩,.
於是她主動和他玩,.
和他聊天,.
希望能夠成為他的朋友,.
但後來的結果是什麼呢?.
後來的結果就是她被班裡的其他同學排擠和杯葛,.
是的,我們走這條路,.
我們要認清一件事,.
你不要以為走天國的路是簡單而易,.
走天國的路並不容易,.
你必須要有受苦的心智,.
我還很記得,.
17年前,.
我先生本來在港交所做了接近17年的工作,.
一直以來生活無憂,.
年年晉升,.
亦享有高薪厚職,.
很多時候他都會出差,.
有很多免費旅遊的機會,.
去到不同的地方,.
就受到高級的禮遇,.
其實真的很舒服,.
走的路真的容易,.
但是當神在17年前去呼召他,.
獻身全時間服侍的時候,.
他就要放棄這個舒服容易走,穩定的生活,.
我還記得當他在讀神學的時候,.
有很多的挑戰,.
腦袋已經塞不進希臘文,.
希伯來文,.
還要這麼多的背誦,.
家裡一個個的難關要去過,.
沒有收入,.
而兩個兒子又出世,.
我們的確是經歷了一個.
遇到受苦艱難的路,.
兒子今年17歲了,.

$^{361}$回想起來,.
原來真的選擇放棄走寬的路,.
走窄路,.
是預料到要艱難,.
有時還很孤單,.
我們真的要問,.
為什麼走窄路會這麼難呢?.
因為有些時候,.
我們走窄路的時候,.
是會有假先知去搞擾我們,.
我們一起看第15至20節,.
讓我們一起讀出,.
你們要防備假先知,.
他們到你們這裡來,.
外面披著羊皮,.
裡面卻是殘暴的狼,.
憑著他們的果子,.
就可以認出他們來,.
荊棘上豈能摘葡萄呢?.
窒泥裡豈能摘無花果呢?.
這樣凡好樹都結好果子,.
唯獨壞樹結壞果子,.
好樹不能結壞果子,.
壞樹不能結好果子,.
凡不結好果子的樹,.
就撼下來,.
凋在火裡,.
所以,.
憑著他們的果子,.
就可以認出他們來..
同樣來到第15節,.
防備這個動詞,.
也是一個命令的動詞,.
耶穌告訴我們,.
我們必定必定要去防備假先知,.
然後,經文教導我們如何防偽,.
如何防偽,.
他用了動物和果樹來做比喻,.
他告訴我們,.
壞的樹,.

$^{401}$例如荊棘和窒泥,.
是結不到好果子,.
結不到像葡萄或無花果的好果子,.
其實,究竟當時的假先知是指哪些人呢?.
那些假先知,.
往往是一些宗教領袖,.
他們可能是法利賽人,.
可能是文士,.
甚至在當時的《十二使徒遺訓》裡,.
提到一些基督徒的先知,.
會巡迴各地,.
而這些假先知,.
是披著羊皮的狼,.
披著羊皮的狼是怎樣的呢?.
他們表面純良,.
但內心是殘暴的,.
殘暴這個字有兩個解釋,.
一個解釋是說,.
他們是殘暴的意思,.
為甚麼呢?.
因為他們引人走進歡歡,.
走歡路,.
引導人走向滅亡,.
馬太福音24章11至12節這樣說,.
「且有好些假先知起來,.
迷惑多人,.
只因不法的事增多,.
許多人的愛心才漸漸冷淡了.」.
他說假先知會迷惑人,.
假先知會令人離開基督,.
令這個世界的罪惡越來越多,.
就算信徒有著愛心,.
他們的愛心的火也會漸漸熄滅,.
他們會灰心,.
他們會越來越武裝自己,.
保護自己,.
以至沒有實踐主的愛,.
這就是披著羊皮的狼,.
假先知的殘暴,.
帶人去死,.

$^{441}$但另一個解釋,.
他們貪婪,.
他們可以藉著基督的名義,.
去謀取自己的利益,.
他們可以令自己很有魅力,.
令很多人跟隨他,.
以至他謀取自己的利益,.
這裡說我們要防備假先知,.
我們如何防備?.
是因為我們可以從他的果子認出他,.
壞樹會結出壞果子,.
他們是貪婪自私的,.
我們會看得到,.
好牧羊,.
好牧人是會為羊捨命的,.
一個壞的領袖,.
他注重的是自己的榮耀,.
自己的好處,.
令自己得益,.
但他們的道德觀,.
他們的價值觀,.
他們的人生,.
跟耶穌所提的並不一致,.
他們的生命並不能活出與悔改的.
心相稱的行為,.
美國有一個電視報道家,.
他叫Jim Barker,.
他在七八十年代,.
上世紀七八十年代,.
他是一個有名的電視報道家,.
但在一九八八年,.
他被控告一些關於詐騙的罪名,.
而在過程中揭發了他有很多性醜聞,.
最終他囚禁了六年,.
他承認他所作的一切,.
他口說耶穌的道,.
宣揚耶穌的福音,.
但他的生命完全是一個.
與耶穌的價值背道而馳的人,.
他在他的自傳中,.

$^{481}$I was wrong,.
承認了自己的錯,.
從他的果子,.
我們知道他是一棵壞樹,.
因為他走出來的,.
絕對是與得救的恩相違背的品格,.
弟兄姊妹,除了我們要去防備先知的搞擾之外,.
有一樣東西是極之極之重要的,.
就是我們要防備自己成為假先知,.
我們的生命有沒有活出好的果子,.
我們的生命會不會活出基督的樣式,.
我們要小心自己都會成為一隻披著羊皮的狼,.
在迷惑人離開基督,.
弟兄姊妹,.
門徒是要走窄路,.
同樣我們都要結出好果子,.
我們的生命有沒有聖靈所結的果子,.
接下來的經文就告訴我們,.
其實我們要結果子,.
一定要留意一樣很重要的東西,.
就是我們的生命有沒有遵行神的旨意,.
讓我們一起讀21至23節,.
「凡稱呼我主呀主呀的人,.
不能到盡天國,.
唯獨遵行我天父旨意的人,.
才能進去..
當那日必有許多人對我說,.
主呀主呀,.
我們不是奉你的名傳道,.
奉你的名講鬼,.
奉你的名行許多義能嗎?.
我就明明地告訴他們說,.
我從來不認識你們,.
你們這些作惡的人,.
離開我去吧.」.
耶穌說,.
有些人的口會叫他主呀主呀,.
還會奉他的名講鬼,.
傳道行神跡,.
但是耶穌怎麼稱呼他們?.

$^{521}$耶穌稱呼他們做惡人,.
但最重要的,.
不單稱呼他們做惡人,.
耶穌說,.
「我不認識他們,.
我不承認你們.」.
意思就是說,.
我根本和你們一點關係都沒有,.
因為耶穌關注的是,.
究竟他們是否正在遵行天父的旨意..
耶穌在此指出一個實況,.
即使有些人信了耶穌基督很久,.
甚至是一些很有事奉機會的人,.
或者領袖,.
但在他們的生活和行為上,.
是沒有實行天父的旨意..
很多人會說,.
保羅在羅馬書說,.
「因信稱義.」.
只需要信耶穌,.
就可以上天堂,.
於是就可以翹起雙手,.
只要受浸主日的回教會,.
每個月十一奉獻,.
你就可以做你喜歡的事情..
有些信徒只是求耶穌的愛和恩典,.
卻拒絕耶穌的責備和管教..
這就像現在的當代,.
有些人信了主二十年,.
三十年,.
四十年,.
我們都會稱呼他為「知心的初信者」,.
因為他的生命,.
他的個性,.
他的行為,.
他的態度,.
一點都沒有改變..
但雅各書很清楚地告訴我們,.
信心沒有行為是死的,.
沒有帶來行為改變的信心,.

$^{561}$是假的信心..
所以你說你恩信清義,.
我會質疑你,.
你的信到底是什麼信?.
這個信是不是真真正正的信?.
所以,.
弟兄姊妹,.
要進天國,.
要走擇路,.
我們必須要遵行神的旨意,.
才能結出好的果子..
其實到底我們遵行神的心意,.
那個原則和心態應該是怎樣呢?.
我們說我們要遵行天父的心意,.
我們究竟知不知道天父的心意是怎樣?.
其實,.
星期三那一晚,.
誠心所願,.
楊叔教授跟我們分享了一樣很重要的事,.
在這個世代裡,.
小心讀科技,.
我們隨手接觸手提電話,.
或者一些科技的產物,.
很容易,.
於是我們提到上網,.
上YouTube,.
玩遊戲,.
看IG,等等..
會用了我們日常生活很多很多的時間,.
我們接觸這些時間非常多,.
但我很想問大家,.
我們去明白主日道,.
和神親近的時間有多少?.
大家現在的靈修生活是怎樣?.
我們有沒有想過,.
不只是看完一次靈修APP就叫做靈修,.
我們有沒有試過放下,.
我們常常用電話來做娛樂,.
消磨時間,.
我們刻意節制,.

$^{601}$減少這些時間,.
去讀聖經親近神,.
去問上帝的心意,.
然後當我們讀到聖經的說話的時候,.
我們切實去思考,.
究竟耶穌今天跟我說的話,.
對我有什麼提醒?.
我們需要去背今句,.
真門徒的生命裡面,.
是有聖靈的提醒,.
當聖靈向你啟示他的真理,.
你的心有擇心,.
你不要白白流失了這個機會,.
因為這個就是主耶穌要去教導你,.
要讓你在生活裡面實踐他的心意的時候..
另外,我很想告訴大家,.
我們要在這個世代走擇路,.
順服主心意的時候,.
有時候真的很艱難,.
但是,這段經文告訴你,.
耶穌會認我們的,.
耶穌會認我們,.
他會跟我們同行,.
他跟我們有親密的關係,.
我們深信最寶貴的是,.
我們是不完美的,.
但是,當我們願意去明白他的心意的時候,.
他一定會跟我們同行,.
他一定會給我們很多恩典,.
給我們很多力量去面對,.
另外,我想說,.
進行神的旨意的時候,.
我們會很孤單,.
我鼓勵大家,.
不要只是自己跟神說,.
去跟你身邊屬靈的同行者分享,.
你有沒有土伴,.
你有沒有屬靈的密友,.
你的伴侶呢?.
你可不可以跟他分享你的掙扎?.

$^{641}$千萬不要單打獨鬥,.
唯有你經歷一個,.
不是靠自己隨己心意的生命,.
不用自己的方法去走人生路的時候,.
你在地上,.
你已經經歷與神同行,.
在耶穌裡面的豐盛,.
所以,弟兄姊妹,.
如果我們說,我們要得生命,.
一個永遠的生命,.
不要滅亡的話,.
我們一定要走窄路,.
是,我們要有一個提醒自己的心,.
走窄路是艱苦的,.
走窄路是孤單的,.
但是在這條走窄路的上面,.
我們要提防才對我們有所搞擾,.
不要被他們影響,.
引我們離開了基督,.
而且自己的生命,.
也要結出生命的好果子,.
而這個結出生命的好果子,.
最重要的是,.
我們必須去遵行天父的旨意,.
其實當我去想,.
我們走這個生命的窄路的時候,.
我想到一個完美的榜樣,.
就是耶穌基督,.
耶穌基督,他放棄了天上的尊榮,.
來到這個世上,.
經歷了很多的艱苦,.
他被逼迫,.
他遇上法利賽人的宗教領袖,.
他被迫害,.
但是,在他的生命路上,.
他是滿結果子,.
在這段窄路上,.
他很孤單,.
沒有人明白,.
只有天父知道,.

$^{681}$在帕西瑪利園,.
門徒都睡著了,.
彼得也不認他,.
但是他選擇遵行天父的旨意,.
釘在十字架上,.
成就救恩..
弟兄姊妹,.
讓我們一起,.
和我們這位釘身十架的基督,.
去走這條通往永生的窄路,.
我們一起祈禱..
主耶穌,我們多謝你,.
多謝你,.
這麼清楚地向我們作出警告,.
告訴我們,.
我們的人生,.
要怎樣才可以通向永生的道路,.
求你幫助我們,.
立志要在這條艱苦和孤單的窄路上,.
行出你的心意,.
結出好的果子,.
進行你的使命,.
讓我們去經歷你豐盛的同行,.
主耶穌,我們會軟弱,.
我們會有時跌倒,.
但是你答穩我們,.
我們去走這條路的時候,.
恩典不會少,.
我們這樣禱告教堂,.
奉主耶穌基督得勝的明智祈求,.
阿們..
\newpage



\section{哥林多前書彼得前書 11:3-12}
\label{sec:UnHN8_G394I}
\textbf{2023年6月18日父親節崇拜直播｜梁家麟牧師｜作頭的男人｜哥林多前書11\_3-12;彼得前書3\_7}
\newline
\newline
連結: \href{https://youtube.com/watch?v=UnHN8_G394I}{\texttt{https://youtube.com/watch?v=UnHN8\_G394I}} ~~~~ 語音日期: 2023-06-20
\newline
\newline
\hyperref[sec:bwknUBKFdWs]{\small{< < < PREV SERMON < < <}}
~
\hyperref[sec:index_chronic]{\small{[返順時目]}}
~
\hyperref[sec:rlyrCSM_7W8]{\small{> > > NEXT SERMON > > >}}
\newline
\newline
$^{1}$今次經訓有三段.
哥林多前書十一章三節.
八到十二節.
還有彼得前書三章七節.
請聽我讀出.
我願意你們知道.
基督是國人的頭.
男人是女人的頭.
神是基督的頭.
起初男人不是由女人而出.
女人乃是由男人而出.
並且男人不是為女人做的.
女人乃是為男人做的.
因此女人為天使的緣故.
應當在頭上有伏權病的記號.
然而照主的安排.
女也不是無男.
男也不是無女.
因為女人原是由男人而出.
男人也是由女人而出.
但萬有都是出乎神.
你們作丈夫的.
也要按情理和妻子同住.
因她比你軟弱.
但你們也要滿一桶承受生命的恩德.
所以要敬重她.
這樣便叫你們的禱告沒有阻礙.
.
各位節目好,早安.
聽到嗎?.
很久沒見大家了.
很不好意思.
再次見大家的場合是父親節.
我連續三年和王浸說父親節主日.
我自己也有想死的感覺.
我可以想像你們作為聽眾的困苦.
為何年年都擺脫不了這個老頭子.
為了避免悶死你們.
今天我稍稍轉一下牛角尖.
找一段非常不好處理的經文來講論.

$^{41}$這是有關男人作為女人的頭的經文.
這些經文通常試探討女性在教會中.
是否能夠事奉角色才被引述.
如果男人是女人的頭.
女人能否在教會中扮演領導的角色呢?.
女人能否講道,能否按木,能否做堂主任呢?.
這些問題在香港的爭論性已經不太強.
不過在北美仍然是一個繳錢不休的課題.
你不要以為北美教會比香港前衛.
不,通常是更保守,非常保守.
中國教會早在1930年代.
這是我論證出來的.
已經按納了女性做牧師.
不過時至今日.
美國仍然是一個極富爭論性的課題.
不久之前,美國以前最大的浸信會.
擁有5萬7千個信徒的馬鞍峰教會.
Settle Back Church.
我們熟悉的Vic Warren的主任牧師.
寫了《目標干人生》的牧者.
由於他教會按納了幾個女性做牧師.
結果美男浸信會將這教會趕出會.
Vic Warren在他的網頁指出.
如果我們讓一半的會友坐著不做事.
我們永遠無法完成大使命.
如果一半的會友坐著.
然後又說.
我相信美男浸信會有數以百萬計的姊妹.
她們的恩賜和才幹是被浪費掉的.
我相信數以百萬計的美男浸信會的恩賜和才幹是被浪費掉的.
今天我不打算在這裡和大家說.
婦女在教會裡事奉的角色的問題.
我也不會和大家討論.
婦女在聚會裡是否需要蒙頭的問題.
不是我想迴避爭論.
而是香港的浸信會在婦女按納的問題上.
很早就已經止息了爭論.
理論上是沒有爭論的.
旺角浸信會早在2007年.
我們教會慶祝45週年的時候.

$^{81}$我們按納了林艷芳和陳信錦.
作為牧師.
今天我們也有吳珍珍牧師.
除非你們想翻案.
否則我們已經解決了這個問題.
旺浸已經沒有這個問題存在.
我們還有Viola作為我們的執事會主席.
所以我說按納女性.
女性做領袖不再是旺浸的問題.
你們應該也同意的.
今天我們也接受女性講道.
女性按木出任堂主任.
不過保羅仍然說男人是女人的頭.
這句話是不是已經過時失效呢?.
如果沒有過時,聖經仍然是聖經.
我們要問這個教訓.
今天可以怎樣理解和應用呢?.
很明顯再保守的基督徒.
也不可能不會明目張膽的.
反對男女平等這個現代社會的.
差不多是最高原則.
平等,種族平等,性別平等.
我猜沒有人敢挑戰這些.
任何在公共領域.
禁止婦女讀書,工作,晉升,做領導的主張.
都是作死的,你同意嗎?.
死定的.
不過唯有在教會裡.
我們仍然有人提出很多古怪的理論.
見到一個很普遍提出的理論.
就是所謂的互補主義.
Complementarianism.
用它來作為對平等主義.
和Nationalism相對的主張.
什麼叫做Complementarianism呢?.
互補主義呢?.
就是宣稱男女的地位其實是一樣的.
不過角色不需要相同.
兩性不同的角色.
正好是互相補充對方的不足.

$^{121}$所以不准婦女講道.
不准婦女安慰.
不准婦女做領袖.
那是無損男女平等的精神.
有沒有聽過這個理論?.
必須很不客氣地說.
這麼荒唐,這麼拙劣的狡辯.
真的只能在教會說出來.
在外面世界說這樣的理論已經死了很久.
我不知道為什麼有時候.
做了教徒就變笨了.
這麼拙劣的理論竟然在教會可以站得住腳.
無論在政治,社會,經濟,文化任何領域.
我們都沒有人可以引用什麼互補主義的理論.
禁止婦女做行政主管.
做教導者,做總統.
或者在我們要作聲一個員工的時候.
考慮的不是才幹,不是表現.
而是性別.
我們肯定婦女平等的地位.
不過出於互補的原則.
做總統這些粗重的事情.
就讓男人繼續做吧.
你們留在茶水間做大姐吧.
你嘗試說出這樣的說法.
今天出於互補的原則.
女性還是不要做醫生好了.
繼續做護士吧.
做護士比較好.
這個叫互補.
你嘗試在醫管局說出這樣的說法.
我保證你死了很久.
今天在社會上還有誰敢設禁區.
限制婦女擔當某些角色.
禁止她們承擔領導地位.
要求她們永遠注定做從屬的角色.
然後宣稱這是為了良性互補.
平等的互補.
這些說法只能套用在地位完全相等的崗位上.
你搓麵,我包餡.

$^{161}$這就叫互補.
卻不能用在權力和權威完全不同的崗位上.
我永遠坐高位,你永遠坐低位.
這就叫平等中的互補.
你還說什麼?.
怎麼可能?.
原則上禁止婦女涉足某些崗位和職級.
這百分百是歧視.
不要跟我胡扯什麼互補主義.
男人賺十萬,女人賺一萬.
不過兩性仍然是平等的.
當我們爭取兩性平等的時候.
我們著眼的是機會.
是待遇這些實際的東西.
同工同酬.
而不是精神上是平等的.
這些空洞無物的口號.
所以男人是女人的頭.
這個教導如果仍然有效.
這個有效範圍在公眾生活上.
基本上沒有人能夠說得出.
只能局限在教會或家庭之內.
如果如我剛才所說.
我不同意在教會裡做禁止婦女出頭的原則性設限.
其實現在只剩下.
唯一可以應用的範圍就是在家裡.
在我們的私人世界裡.
不涉及公眾生活.
在家庭裡.
男人就是丈夫.
女人就是妻子.
所以男人是女人的頭的意思.
就是丈夫是妻子的頭.
這是我對保羅這個教導的認訊.
他是圈面性的.
做正面的鼓勵的.
但不是原則上的規限.
他基本上只適用於家庭私人領域裡的夫妻關係裡.
我們看看在家庭裡.
可以怎樣兌現這個聖經教導.

$^{201}$我們注意到在哥林多前書十一章.
在伊弗所書.
保羅強調男人是女人的頭.
他有找到一些舊約論據.
這句理由是創造了秩序.
就是上帝先做男後做女.
另一個理由就是.
做女的原因是為男人找一個伴侶和幫助.
第一是先後.
第二就是幫助.
如果我們相信創世記這個技術的事實.
這樣創造秩序上.
男先女後是毋須爭論的.
確實先有亞當後有夏娃.
不過如果我們說亞當是先.
所以亞當就是夏娃的頭.
夏娃就要順服亞當的權柄.
這樣的引伸討論是非常勉強的.
你知不知道上帝先做豬後做人的?.
上帝先做禽獸後才做人.
人是否須順服豬的權柄呢?.
創造先後就變成權威的高低.
我不知道這個論據是什麼意思.
至於夏娃被創造的目的.
是要成為亞當的幫助.
這也是歷史的事實.
我們同樣不能簡單地將幫助者視為次等的.
從屬的角色.
因為聖經上上強調.
上帝是人的幫助.
而又說是我的幫助.
是同一個字下.
不是兩個字.
所以我們不能推論說.
上帝是人的幫助.
所以上帝從屬於人.
老師是學生的幫助.
所以老師的地位低於學生.
不是這樣吧.
幫助者是低位嗎?.

$^{241}$所以不要期望單從創世記的技術.
就可以將羅出一個男尊女卑的充分理據.
我們看保羅的新論.
與其說他是直查考創世記.
來確定男尊女卑的原則.
不如說.
他是就當時社會上一個普遍存在.
普遍認受男尊女卑的事實.
去將羅一個屬靈的理由.
保羅不是宣揚.
保羅只是確認當時的事實.
然後找一個聖經的根據.
去合理化這個存在很久的事實.
保羅不是從無到有.
去確立男尊女卑的原則.
只是為這個普遍存在的事實.
去找一個屬靈的理由.
就好像.
你可能不知道.
一直到最少.
有文獻記載.
直到十四世紀,十五世紀為止.
男人都比女人長壽.
你不知道這件事吧.
我們一直以為是女人長壽過男人.
直到十四世紀前.
不是的.
男人普遍比女人長壽.
這個主要原因是因為古代社會.
食物供應是很不足夠的.
營養也不足.
所以.
女性因為生育.
因為月經等等的原因.
往往是最受營養不良.
醫療條件不好的影響而早逝.
你看中國的古代.
常常你見到那些.
這個名人那個名人又續延又續延.
就好像他們老婆死得很快.

$^{281}$死得很容易.
現在就少一點.
現在通常都是男人先死.
不過以前真的是.
女人先死.
這是一個普遍的事實.
不是個別男人長命的女人.
而是極普遍.
普遍到不用討論的現象.
古代的人不明白為什麼.
但要解釋為什麼男人更長壽的現象.
所以就有很多理論說出來了.
包括亞里斯多德.
為什麼男人比女人長壽呢.
是因為男人比女人暖.
Man is a warmer animal.
男人是暖.
你牽牽老婆的手就知道.
她通常都手冷腳冷.
不知道為什麼女人的手會冷一點.
因為男人暖一點.
暖一點就耐一點.
暖一點就長壽一點.
就是這樣.
今天我們不用討論亞里斯多德.
這個解釋是否合理.
因為男人更長命這個現實都不存在了.
有沒有討論呢.
討論是否合理都不用討論.
因為現在女人比男人長命.
這是普遍的現實.
我們今天引用的另一處經文.
彼得前書三章七節這樣說.
「你們作丈夫的.
也要按情理和妻子同住.
因她比你軟弱.
與你一同承受生命之恩的.
所以要敬重她.
將便叫你們的禱告沒有阻礙.
妻子比丈夫軟弱」.

$^{321}$這個不是彼得提出的一個理論.
而是當時一個普遍的現實.
當時多數女性都沒有受過教育的機會.
社會資源有限.
各方面的能力都比較低.
通常她們比較感性.
認知分析能力不高.
所以常常給人一個軟弱的感覺.
動不動就哭.
動不動就不知如何是好.
手足無措.
不過今天我們能否仍然說.
女性必然比男性軟弱呢?.
女性多數都比男性軟弱呢?.
當然仍然有軟弱的女性.
我見過不少.
不過軟弱的男性我都見不少.
真可憐.
只有女性才軟弱嗎?.
不覺得.
所以我們不要將聖經的實言技術.
將它變成一個應言的論斷.
聖經沒有定規男尊女卑.
也沒有定規女性要比男性軟弱.
所以不軟弱的女性.
請你自己扮軟弱.
否則我都要弄到你變得軟弱為止.
不是那麼變態吧?.
聖經沒有將它變成一個必須的原則.
聖經只是道出當時的社會現實而已.
今天的世界變了.
再沒有這樣的現實.
變了就是變了.
保羅在為男尊女卑的現實.
尋求屬靈的解說之餘.
仍然不忘記指出.
聖經最根本的平等原則.
臨前十一章十一到十二節.
然而照主的安排.
女也不是無男.

$^{361}$男也不是無女.
因為女人原是由男人而出.
男人也是由女人而出.
但萬有都是出乎上帝.
男與女是互相依存的.
男人不可以離開女人.
女人不可以離開男人.
這就像中國傳統說.
孤陰不生獨陽不長.
陰陽協調才可以化玉萬物.
另外保羅從創造故事尋找靈感再指出.
除了夏娃的被造是由男人而出.
因為男人把肋骨拔出來.
造了女夏娃.
日後女人懷孕生子.
也說明男人也是由女人而出.
女由男而出.
男也由女而出.
兩性因此是平等的.
所以保羅最後說.
萬有都是出乎上帝.
這是男女平等的最終依據.
上帝是人類平等的依據.
正如在這份書中.
保羅說無論是主人或僕人.
他們都同有一個在天上的主人.
所以主人和僕人是沒有本質上的分別.
上帝是男是女.
男和女沒有本質上的分別.
在上帝面前眾生平等.
我是我兒子的父親.
我是學生的老師.
不過在天父上帝面前.
我的兒子和我的學生.
同樣都是我的弟兄.
上帝的絕對性.
將人間一切事物都相對化.
人間不存在必然的尊與卑.
不過我也要補充說.
聖經雖然沒有男尊女卑的要求.

$^{401}$不過聖經也沒有直接鼓勵我們.
在人間爭取男女平等.
在上帝信仰的前提下.
我們確認平等的原則.
男女平等,性別平等的原則.
不過聖經沒有指引我們在社會上.
去推動平權運動.
聖經作者對現實的關懷和追求.
其實不是很多.
正如聖經要求主人要善待僕人.
但沒有要求基督徒起來.
去推翻奴隸制度.
聖經沒有直接這樣呼籲.
我們拒絕用一些經文.
去合理化奴隸制度.
或者合理化任何歧視制度.
這些荒謬的神學理論.
但也不要將耶穌或保羅.
裝扮成為解放奴隸的社會革命家.
或者是婦解分子.
保羅不是,耶穌也不是.
今天我們如何去理解.
男人是女人的頭.
這個聖經的說法呢?.
我有兩個解說.
第一,我會將它看為一個實現的描述.
正如保羅所說的.
男人是女人的頭.
是對當時的社會一個現實的描述.
今天雖然男女平等已經成為共識.
但在多數的社會.
男女幾乎仍然不是一個完全平等的位置.
不要說在伊斯蘭世界.
對婦女在學習,就業,結婚有很多限制.
就算在香港.
香港其實已經是一個相當男女平等的社會.
特別現在.
在20到29歲的年齡層.
擁有大學學位的女性比男性多.
不過因為有些歷史的累積.

$^{441}$上一輩多數是男人讀書多一點.
女人少一點.
所以整體而言還是男人讀大學多於女性.
但現在這一代是女人多於男人.
另外,根據政府統計處的數字.
全港目前有超過60萬全職的家庭主婦.
所以說男主外,女主內的原則已經不再適用.
其實也不是.
60萬也不少.
全職做家務的男人.
不是沒有,有38000多人.
所以有些傳統的觀念做法不是完全失效.
我相信基督徒的價值觀.
通常我們都是比較保守的.
也更加注重家庭,更加注重子女教育.
所以我相信,我沒有數字.
基督徒中間全職的家庭主婦的比例.
是會高於外面的世界.
你同不同意嗎?.
我相信價位要高一點.
我認識不少傳道人.
不少弟兄姊妹都是太太不做事的.
帶小孩的,專心.
所以,這是一個事實的描述.
另外,我們要留意的是.
聖經教導我們做男性的.
要做女人的頭.
聖經不是要求所有男人一定要做頭.
是教導如果我們要做頭的話.
就要這樣這樣做頭.
保羅是跟我們說明.
做頭是什麼一回事.
我們知道,無論是耶穌還是保羅.
都有非常激進的社會教導.
幾乎是全面顛覆當時的想法和做法.
例如,我們的主.
那個僕人領袖的教導.
我非常喜歡這段經文.
《馬福音》十章42-45節.
你們知道外邦人有尊為君王的治理他們.

$^{481}$有大臣操權管束他們.
只是在你們中間不是這樣.
你們中間誰願為大.
就必作你們的僕人.
在你們中間誰願為首.
就必作眾人的僕人.
因為人子來並不是要受人的服侍.
而是要服侍人並且要寫命作多人的屬價.
做大的要服侍做小的.
耶穌不是空口說一個原則.
耶穌是當時期先做一個洗腳的示範.
但他做一個最終極的示範是什麼.
寫命作多人的屬價.
耶穌基督在十字架上.
將他說的這個僕人領袖的原則完全活現出來.
說到男人怎樣做女人的頭.
我們知道保羅在哥林多前書和這份所書.
都有強調男人是女人的頭的事實.
而在他談論做頭的道理的時候.
他都將丈夫和妻子的關係.
比附為基督與教會的關係.
丈夫是妻子的頭.
正如基督是教會的頭.
這是基督論.
是一個救贖論的比輔.
比他之前說創造故事的解說更加重要.
是的.
基督與教會從來都不是一個平等的關係.
是一個從屬的關係.
所以教會要信服基督.
我們卻不會說基督會信服教會.
妻子如果要像教會信服基督那樣去信服丈夫.
就不要期望丈夫反過來都信服妻子.
這個信服是單向性的.
我經常說解經的時候.
不要將二弗所書五章二十一節.
彼此信服的命令夾纏在.
跟著他討論的夫妻關係那裡.
因為教會與基督不存在彼此信服的關係.
那怎麼死呢?.

$^{521}$基督信服教會.
保羅的說法.
二弗所書五章二十四節是很硬朗的.
教會怎樣信服基督.
妻子又要怎樣凡事信服丈夫.
是單向的.
不過你們一定聽過我不止說過一次.
聽到有點悶.
這個老頭子有些事情是說來說去.
沒有什麼新的東西.
不勝經的道理沒有什麼新的東西.
保羅這個夫妻關係的教導.
與其說是對女性不公平.
不如說是對男性非常不公平.
丈夫與妻子的關係.
如果等同基督與教會的關係.
那你告訴我.
基督付出得多還是教會付出得多.
不用說了.
基督付出的是壓倒性.
基督是為教會寫明.
所以如果按照基督與教會這個關係的比喻.
保羅是要求做丈夫的.
為妻子全然擺上犧牲自己.
不計代價成全太太.
二弗所書五章二十八二十九節說.
丈夫也當照樣愛妻子.
如同愛自己的身子.
愛妻子就是愛自己了.
從來沒有人恨惡自己的身子.
總是保養固息.
正如基督待教會一樣.
保羅要求妻子信服丈夫.
同時要求丈夫為妻子犧牲一切.
你告訴我這個教導對男人公平還是對女人公平?.
弟兄姊妹.
這是保羅眼中男人是女人頭的道理的具體實踐.
要做頭就去死吧.
要做頭就準備犧牲一切吧.
誰說因為男人是女人的頭.

$^{561}$所以丈夫可以打老婆?.
發神經.
誰說因為男人是女人的頭.
所以丈夫在家裡可以攪拌.
什麼家務都不做?.
誰說因為男人是女人的頭.
所以做男人做丈夫的就可以做皇帝.
做惡霸在家裡哈哈爸爸?.
沒錯,男人是女人的頭.
誰敢將人間尊卑高低的標準和做法.
偷混到聖經裡.
耶穌清楚將聖經的領袖官告訴我們.
你們中間卻不是這樣.
你們中間不是這樣.
不是這樣.
你們中間誰願為大就必作你們的僕人.
男人是女人的頭.
所以男人在家裡要做大夫.
回到這篇講道的真心信息.
我們所處的社會.
如果仍然是男人為主導.
如果男人仍然是頭.
保羅就教導我們.
做好作為頭的角色.
保羅不是定規男性一定要做頭.
他教導如果我們做頭.
就要做出聖經的做頭的.
那個僕人領袖的樣子.
我再說,在今天的社會.
男人不一定再扮演頭的角色.
正如剛才所引的數字.
全職主婦有六十萬.
是很多的,六十多萬.
但是全職的主公.
做老公在家裡做家務的.
都有三萬多.
在我認識的人裡面.
我能夠馬上說出的.
有五個,我數到五個家庭.
很多年前已經是.

$^{601}$丈夫在家裡做家務.
照顧小孩.
太太在外面工作.
其中一個是我的家人.
做老公的.
在大學的時候是讀EE.
不過在八十年代已經.
公眾移去大陸.
他們決定不想老公和老婆分開.
在大陸工作.
所以決定老公在家裡做家務.
老婆一個人出去賺錢.
另外,我也認識一對.
我以前教會的.
太太是一個accountant.
不過做了一個英國的跨國公司的.
高級外派員.
派去大陸很多不同的城市.
做老公的只能跟著老婆.
去不同的地方工作.
如果太太是家裡收入的來源.
如果在現實上太太做了頭.
我們強行要求.
老公仍然是頭.
怎麼做?.
你告訴我怎麼做?.
在我看來.
出於聖經教導常有性別互換的原則.
如果太太是家裡的頭.
保羅對丈夫做頭的要求.
對太太有效.
誰做頭的.
都要做耶穌所吩咐的個人領袖.
不可以對弱勢的一方指手畫腳.
哈哈爸爸.
正如彼得的教導.
你們作丈夫的.
要按情理和妻子同住.
因他比你軟弱.
與你一同承受生命之恩.

$^{641}$所以要敬重他.
叫你們的禱告沒有阻礙.
如果太太比丈夫軟弱.
做丈夫的都不可以輕看他.
反倒要敬重他.
認定你不可以撇下他.
擺脫對方自行去做信仰的追求.
你要與對方攜手.
一同承受生命之恩.
這是保羅對在強一方的男人的吩咐.
我們讀聖經時常常選擇性地讀.
所以當讀到臨前十一章.
有關男人是女人的頭的時候.
我們總是著眼於女性要服權病.
在聚會的時候要蒙頭這些細節.
我們忽略到聖經.
如何要求男人做頭.
做頭是什麼意思.
要怎樣做頭的表現.
男人如果沒有做頭的樣子.
卻勉強女人服從他們.
這是超級混帳的事.
耶穌基督首先為教會寫記.
然後才要求教會信服.
你要人信服.
你值得對方信服嗎.
你要人當你是頭.
你做到頭的表現嗎.
親愛的弟兄姊妹.
你認識我這麼久.
雖然我主張男女平等.
不過我心底裡.
我是百分之二百的大男人主義者.
我是一個大男人.
非常大男人.
我對男性.
對弟兄.
對丈夫.
總是體高一線.
有格外的要求.

$^{681}$做男人要有男人的樣子.
做丈夫做父親.
要有俠如其分的表現.
我常常說大男人不可怕.
小男人才可怕.
一對夫妻鬧離婚.
當然雙方都有責任.
不過我過去的經驗.
所形成的一個偏見是.
丈夫通常要負較多的責任.
大概是六四或七三的比.
丈夫好.
家庭通常都容易變好.
我是家庭的頭.
我立志要做我兒子的榜樣.
我立志要做我學生的榜樣.
以下這段說話只是跟弟兄說.
父親節也不跟弟兄說話.
我們都被邊緣化很久了.
那位弟兄.
我們被命定去做男人.
我們要承擔做男人的命運.
責任 責任 責任 責任.
付出不計較.
肯吃虧 成全對方.
不抱怨 不畏過於人.
不退縮 不逃避.
有錯就認.
正視現實 迎向挑戰.
不要將那些女性化的輔導理論.
污染我們的男性化的理論.
污染我們的男性化.
我常說聖經要求我們做大丈夫.
你數一下有多少個聖經對大丈夫的說法.
《三會記》上四章九節.
《三會記》下三章三十八節.
十章十二節.
《聶王記》上二章二節.
《歷代志》上十九章十三節.
《以賽亞書》四十六章八節.

$^{721}$《哥林多前書》十六章十三節.
Be Man.
Be a Great Man.
翻譯成NIV就是這麼翻譯的.
我從來不在家發洩我心中的壓力和情緒.
我在家不會發瘋.
不過我在教會 在神學院我幾乎都不發瘋.
我不覺得.
因此我心靈就變得很不健康.
老男人 臭男人.
我唱的歌仍然是.
男兒一生要經過.
世上磨練共多少.
男兒一生要幾次做到失望與心照.
我有無邊毅力.
捱盡困難考驗.
立誓披荊斬棘.
心裡更怨 永不折腰.
你不要說這些.
忍耐不作聲的態度.
是落伍過時的臭男人的做法.
我告訴你.
這是耶穌基督為我們樹立的榜樣.
你自己看二頌書53章7節.
就像很多年前我在上帝面前.
所立的一個小男人的字.
我希望我太太覺得幸福.
我兒女覺得幸福.
我身邊的人因著我覺得幸福.
各位弟兄.
我們做不做得到.
立誓讓你身邊的人覺得幸福.
我們就此致這樣做.
如果我們真的做頭做到這樣的樣子.
你想想我們的太太或其他異性.
會不會很抗拒我們.
會不會大罵你這個死人頭.
上帝恩典我們.
父親節快樂.
\newpage



\section{約翰福音 2:13-25}
\label{sec:rlyrCSM_7W8}
\textbf{2023年7月9日崇拜直播｜梁家麟牧師｜潔淨聖殿｜約翰福音2\_13-25}
\newline
\newline
連結: \href{https://youtube.com/watch?v=rlyrCSM-7W8}{\texttt{https://youtube.com/watch?v=rlyrCSM-7W8}} ~~~~ 語音日期: 2023-07-12
\newline
\newline
\hyperref[sec:UnHN8_G394I]{\small{< < < PREV SERMON < < <}}
~
\hyperref[sec:index_chronic]{\small{[返順時目]}}
~
\hyperref[sec:CAkR8_FGcmQ]{\small{> > > NEXT SERMON > > >}}
\newline
\newline
約翰福音 2:13-25
\newline
\begin{longtable}{cl}
\hline
\hline
章節 & 經文 (和合本修訂版)\\
\hline
2:13 & \begin{tabularx}{0.7\textwidth}{X} 猶太人的逾越節近了，耶穌上耶路撒冷去。 \end{tabularx} \\ \\ \relax
2:14 & \begin{tabularx}{0.7\textwidth}{X} 他看見聖殿裡有賣牛羊和鴿子的，還有兌換銀錢的人坐著， \end{tabularx} \\ \\ \relax
2:15 & \begin{tabularx}{0.7\textwidth}{X} 耶穌就拿繩子做成鞭子，把所有的，包括牛羊都趕出聖殿，倒出兌換銀錢之人的銀錢，推翻他們的桌子， \end{tabularx} \\ \\ \relax
2:16 & \begin{tabularx}{0.7\textwidth}{X} 又對賣鴿子的說：「把這些東西拿走！不要把我父的殿當作買賣的地方。」 \end{tabularx} \\ \\ \relax
2:17 & \begin{tabularx}{0.7\textwidth}{X} 他的門徒就想起經上記著：「我為你的殿心裡焦急，如同火燒。」 \end{tabularx} \\ \\ \relax
2:18 & \begin{tabularx}{0.7\textwidth}{X} 因此猶太人問他：「你能顯甚麼神蹟給我們看，表明你可以做這些事呢？」 \end{tabularx} \\ \\ \relax
2:19 & \begin{tabularx}{0.7\textwidth}{X} 耶穌回答他們說：「你們拆毀這殿，我三日內要把它重建。」 \end{tabularx} \\ \\ \relax
2:20 & \begin{tabularx}{0.7\textwidth}{X} 猶太人說：「這殿造了四十六年，你三日內就能重建嗎？」 \end{tabularx} \\ \\ \relax
2:21 & \begin{tabularx}{0.7\textwidth}{X} 但耶穌所說的殿是指他的身體。 \end{tabularx} \\ \\ \relax
2:22 & \begin{tabularx}{0.7\textwidth}{X} 所以他從死人中復活以後，門徒想起他曾說過這事，就信了聖經和耶穌所說的話。 \end{tabularx} \\ \\ \relax
2:23 & \begin{tabularx}{0.7\textwidth}{X} 耶穌在耶路撒冷過逾越節的時候，有許多人看見他所行的神蹟，就信了他的名。 \end{tabularx} \\ \\ \relax
2:24 & \begin{tabularx}{0.7\textwidth}{X} 耶穌自己卻不信任他們，因為他認識所有的人， \end{tabularx} \\ \\ \relax
2:25 & \begin{tabularx}{0.7\textwidth}{X} 也用不著誰來證明人是怎樣的，因為他自己認識人的內心。 \end{tabularx} \\ \\
[1ex]
\hline
\hline
\end{longtable}
$^{1}$接下來我們繼續用聖道去敬拜.
今天我們選讀的經文是在新約聖經約翰福音.
約翰福音第二章十三到二十五節.
請聽我讀出.
猶太人的愉悅節近了.
耶穌就上耶路撒冷去.
看見店裡有賣牛羊甲子的.
並有兌換銀錢的人坐在那裡.
耶穌就拿繩子作成鞭子.
把牛羊都趕出店去.
倒出兌換銀錢之人的銀錢.
推翻他們的桌子.
有對賣甲子的說.
把這些東西拿去.
不要將我付的電當作買賣的地方.
他的門徒就想起經上記著說.
我為你的電心裡焦急如同火燒.
因此猶太人問他說.
你記著這些事.
還顯什麼神蹟給我們看呢.
耶穌回答說.
你們拆毀這電.
我三日內要再建立起來.
猶太人便說.
這電是四十六年才造成的.
你三日內就再建立起來嗎.
但耶穌這話是以他的身體為電.
所以到他從死裡復活以後.
門徒就想起他說過的話.
便信了聖經和耶穌所說的.
當耶穌在耶路撒冷過雨節的時候.
有許多人看見他所行的神蹟.
就信了他的名.
耶穌卻不將自己交託他們.
因為他知道萬人.
也用不著誰見證人怎樣.
因他知道人心裡所存的.
多謝各位早安.
我們看一個很有趣的故事.
耶穌結證聖殿.

$^{41}$猶太人有三個主要的節期.
駐朋節,雨節和七七節.
七七節又叫五旬節.
如果是可以的話.
他們會從住處徒步到耶路撒冷朝聖.
並且將祭物帶到聖殿獻上.
在約翰福音裡面.
只是約翰福音說.
記載了最少三個雨節.
分別是二章十三節,六章四節和十一章五十五節.
或者還有第四個.
就是五章一節.
因為那裡沒有提到什麼特別的節期.
不過推測也可能是雨節.
這就吻合了耶穌大概三年半公開傳道生涯.
三年半就有三個或四個雨節.
所以福音書主要記載這段時間.
大概就是這樣.
耶穌是一個遵守猶太律法的人.
每年都不止一次上耶路撒冷朝聖.
雨節是他幾乎一定去耶路撒冷的時間.
雨節是猶太年曆裡面的尼撒月十四日.
就是我們陽曆大概三月到四月之間.
我們知道復活節總是在雨節後面.
所以我們大概知道復活節的時間.
就猜到雨節是什麼時間.
聖經一開始就說.
猶太人的雨節近了.
耶穌就上耶路撒冷去十三節.
為什麼要強調這是猶太人的雨節呢?.
我們剛剛過了端午節.
我們不會說中國人的端午節過得這麼噁心.
我們都是中國人.
為什麼要聲高聲小說中國人呢?.
表明猶太人的雨節.
證明這段說話是針對外邦人說的.
約翰福音的作者面對以外邦人為主的讀者.
所以特別著名.
這是猶太人的雨節.
耶穌不是一個人上耶路撒冷.

$^{81}$這裡說他是跟門徒一起去的.
到了聖殿.
聖經作者沒有記述他們如何獻祭.
但記載了耶穌進行了一次潔淨聖殿的行動.
十四到十六節.
看見殿內有賣牛羊甲子的.
並有兌換銀錢的人.
坐在那裡.
耶穌就拿繩子作成鞭子.
把牛羊都趕出去.
倒出兌換銀錢的人的銀錢.
推翻他們的桌子.
從兌換甲子的雪.
把東西拿去.
讓他們把付的電當作買賣的地方.
這裡首先要討論的是.
其他三卷婦女福音.
即馬可路加.
都將潔淨聖殿的事件.
放在耶穌事奉生涯的最末段.
即是在上十字架前不久.
潔淨聖殿的行動意味著.
他跟猶太教的領袖公開對決.
彼此勢不兩立.
再無轉圜餘地.
但是約翰將潔淨聖殿的事件.
放在耶穌事奉生涯的最早期.
第二章.
放在最早期.
技術上的差異其實不容易把他和諧化.
有人甚至建議.
不如估計耶穌有兩次潔淨聖殿的行動.
一次在前一次在後.
不過我覺得這樣的說法有點奇怪.
這樣的行動可以搞兩次.
我們寧願相信聖經作者的編排.
有特別的深意在其中.
因為聖經有說到兩個可能的時間.
我們不能百分百確定.
潔淨聖殿的事件實際發生在什麼時候.

$^{121}$不過我估計.
發生在耶穌上十字架之前.
即是說在他事奉生涯的末段.
應該是比較合理的.
那為什麼約翰.
寧寧懾懾要將這個事件.
搬到耶穌事奉生涯的前期呢?.
原因很可能是他認定.
耶穌整個的救贖使命.
目的就是要將猶太人的聖殿徹底潔淨.
甚至是將它拆毀重建.
潔淨聖殿的事件.
是標誌著耶穌的救贖使命的目標.
耶穌藉著這次行動.
公開宣示.
我來便是要潔淨聖殿.
將之還原為真正敬拜上帝的地方.
我要恢復上帝指引.
一個純潔的敬拜.
我們接著看第四章.
後面四章.
耶穌與聚家城的婦人談道的時候.
就強調什麼叫做真正的敬拜.
上帝期望我們如何去敬拜祂.
我們後面會說的.
我們看看潔淨聖殿行動的細節.
從下面吸輪谷.
一路上橄欖山.
沿路都有很多攤販.
販賣憲制用的牛羊甲子.
這是為了方便.
那些從遠方來憲制的人.
而做的設施.
你知道人很難拉著牛羊.
走幾天的路.
牛羊走很慢.
拉著牠走.
最少多走一倍時間.
兩倍時間.
不止,對不對.

$^{161}$與其在家裡帶一隻牛來憲制.
不如來到這裡買一隻牛來憲制.
與其帶一隻羊來憲制.
走三天路.
由加利利走下來.
不如在這裡買一隻.
是不是簡單很多.
所以這些設施不是壞的設施.
是方便的設施.
這些攤販.
通常都是設在聖殿的外面.
好像今天不同的宗教寺廟.
外圍沿路都有很多.
販賣那些制制用品.
或者宗教紀念品的店舖.
黃大仙祠.
或者車公廟.
門口都有很多這些攤販.
如果你去日本或泰國的廟宇.
都會看到這樣的情況.
聖經這裡說.
有人要佔地來做生意.
竟然將檔口開到聖殿裡面.
大概應該不是在聚裡面.
應該是在外邦人院.
沒有那麼誇張.
不過開在外邦人院.
都會有些混帳.
為什麼要開到圍牆裡面.
管聖殿的人為什麼不干涉這件事.
幾個可能.
或者他們管不住.
武力控制一些惡勢力.
或者中間可能有人收受了利益.
所以睜開眼閉.
販賣生畜的攤檔開在聖殿裡面.
你知道通常都要討價還價.
又要選生畜.
你們去街市都要選來選去.
你會發覺.

$^{201}$整個聖殿本來都已經夠擠迫.
第二聖殿不是很寬大.
如今可以看到更擠迫.
並且有烏煙瘴氣的感覺.
除了販賣祭祀的攤檔.
還有兌換錢的攤檔.
律法規定.
這是在《出乃及記》的30章11-16節說的.
每個猶太成年男人都要獻上半赦克勒的銀紙.
聖殿稅.
這是聖殿稅.
從各地來的人.
他們來自不同的國家.
所以他們其實帶著不同的貨幣.
貨幣要兌換.
要等價.
半赦克勒.
還有不同的貨幣.
當時期不像今天.
銀的含銀量不一定一樣.
純度不一樣.
所以當時規定.
要用推羅的半赦克勒的銀幣.
希臘的銀幣也可以用.
其他就不可以用.
所以如果剛好帶著港幣去.
就要兌換推羅的貨幣.
才可以知道等價是多少.
才去獻祭.
通常這些攤檔.
不是任何時候都在開在聖殿裡面.
據說在月前十天.
兌換銀錢的人.
可以將攤檔搬到聖殿外邦人院.
方便去敬拜那些猶太人.
我們知道聖經裡面說.
耶穌沒有原則上反對繳付聖殿稅.
你記得他曾經指示彼得.
用魚擔著的銀幣去交稅.
所以耶穌不會公然反對.

$^{241}$是律法規定要交聖殿稅.
耶穌沒理由反對.
有人估計.
耶穌之所以對那些攤販有激烈反應.
是因為他們做生意不老實.
收過高的手續費或利潤.
或與管聖殿的人勾結貪污.
我們不能完全排除這個可能性.
不過聖經沒有這樣說.
從耶穌的口中.
我們看不到他對這方面的暴利或貪污有任何不滿.
耶穌只說.
把這些東西拿去.
不要將我付的電當作買賣的地方.
所以耶穌只不喜歡在聖殿裡經營買賣做生意.
二章十五節記述.
耶穌說.
繩子作繩鞭子把牛羊都趕出殿去.
倒出兌換銀錢.
指人的銀錢推翻他們的桌子.
在聖經裡.
耶穌用暴力行為.
即動粗.
在福音書的記述裡.
大概就是這一次.
只一次.
就算被抓.
也只能被納加山動粗.
耶穌也沒有動粗.
耶穌動粗只有一次.
當然.
這個看來是暴力的行動.
其實也不是很暴力.
不是很血腥.
因為耶穌只是用一根繩子作為鞭子.
這不是真正的武器.
不具殺傷力.
耶穌沒有打人.
只是推翻他們的桌子.
並且.

$^{281}$雖然聖經沒有詳細說明.
但看來.
耶穌的驅趕行動.
並沒有招來激烈的反抗.
那些貪蕩的主人.
沒有和耶穌繳錢不休.
堅決拒絕離場.
更加奇怪的是.
聖殿的管事和即日的祭司.
也沒有出面攔阻.
沒有制止耶穌的清場行動.
保護貪犯的利益.
他們似乎也覺得.
在聖殿擺賣是有違祖制.
在情在理也說不過去.
所以不敢公然支持.
那些貪犯的商業行動.
只好任由耶穌在聖殿裡放肆.
在聖經的技術裡.
我們看到耶穌的潔淨行動.
沒有看到貪犯和祭司的反應.
不好意思.
也可以這麼說.
希律王在聖殿山.
建造了安東尼堡.
派駐了羅馬軍隊駐守.
用那裡監察聖殿的行動.
那時候是羅馬統治的時間.
所以如果耶穌的潔淨聖殿.
招來暴力反抗.
甚至釀成暴動,流血事件.
後果是很嚴重的.
如果這次羅馬軍隊沒有出動.
證明這次聖殿的行動.
是耶穌單方面的行動.
並沒有造成肢體衝突.
並沒有造成大事件.
我們想一想.
耶穌潔淨聖殿行動的意義.
在聖殿裡有貪犯,賣祭司,兌換銀錢.

$^{321}$就算不是由來已久.
我也相信沒有理由是今次才開始的.
耶穌每年都來聖殿最少一次.
祂不應該是第一次看到這樣的現象.
為什麼今次祂有這麼激烈的反應呢?.
有兩個可能的解釋.
第一,祂純粹心血來潮.
這個行動是即興的.
耶穌其實之前沒有什麼想法.
祂只是看到聖殿裡的嘈亂巴閉.
一時火起,就作出暴力行為.
第二,事先部署.
耶穌是有計劃做這次行動的.
祂要藉著這次行動.
宣示一個關乎祂使命的訊息.
我不能排除第一個可能性.
耶穌是即興的,突然間火起.
然後就作出一個火爆的行動.
不過我傾向相信第二個解釋.
耶穌是有計劃作出一個結聖聖殿的行動.
我們知道結聖聖殿.
是舊約聖經的一個應許.
撒迦利亞書最後一節的聖經說.
當那日在萬君子耶華的殿中.
必不再有迦南人.
這是撒迦利亞書14章21節.
迦南人這個字原文也可以解作商人.
就是說上帝在所預定的那一天.
要剪除所有在聖殿裡營商謀利的人.
另外,《馬拉基書》也提到.
上帝差遣使者結聖聖殿.
以及在聖殿裡服侍的人.
作出了一個這樣的應許.
《馬拉基書》三章一節和三節說.
萬君子耶華說.
我要差遣我的使者在我面前預備道路.
你們所尋求的主必忽然進入他的殿.
立約的使者就是你們所仰慕的快將快要來到.
去到第三節.
他必坐下如煉證銀紙的.

$^{361}$必潔淨利未人.
熬煉他們像金銀一樣.
他們就憑公義獻供物給耶華說.
聖殿敬拜是救藥以色列民的宗教生活的中心.
耶穌基督潔淨聖殿的行動.
於是祂要更新以色列民的宗教生活.
藉此修補他們和上帝破損了的關係.
我們在耶穌接著和猶太人的對話裡面.
看到這一方面的端倪.
先是第十七節.
使徒約翰記述這個行動之後.
門徒的反應.
他的門徒就想起經上記著說.
我為你的殿心裡焦急如同火燒.
這個引經文是詩篇69篇第九節.
不過如果你看詩篇69篇的話.
你會知道其實是有些硬.
斷章取義地去拔這節經文出來的.
因為跟69篇所說的東西沒有什麼關係.
而約翰斷章取義拔這節經文.
目的是要說明耶穌對聖殿是很著緊的.
因為祂著緊.
祂期待聖殿是乾乾淨淨的.
期待聖殿有合上帝心意的敬拜.
是這樣祂才做出這麼激烈的行動.
因為祂在意.
所以祂作出這樣的行動.
值得注意的是.
門徒不是在當下現場就有這方面的聯想.
他們是事後回想.
才明白耶穌這個行為背後的心意.
22節說.
所以到他從死裡復活以後.
門徒就想起他說過的話.
便信了聖經和耶穌所說的.
門徒是要到耶穌死而復活之後.
才逐漸明白耶穌一言一行的心意.
他們真是後知後覺.
我們可以推想.
彼得,雅各,約翰等門徒.

$^{401}$在當場他們其實不明白耶穌在做什麼.
直到約翰寫這本書的時候.
他努力要去整理耶穌的生平和言行.
然後才發覺他的夫子做每一件事都有心意.
所以約翰將耶穌結證聖殿的故事.
放在耶穌事奉生涯的最早段.
目的也是要將這件事件作為一個標誌.
反映耶穌救贖使命的本質.
耶穌結證聖殿.
耶穌重建聖殿.
這是耶穌整個事奉最關鍵的含義.
好了,我們來看18節.
18節開始有猶太人出場.
他們或是管聖殿的人.
或是工會成員.
不要說他們是祭司,不過也可能是祭司.
他們質問耶穌憑什麼權柄來做結證聖殿的行動.
干涉聖殿的事務.
「你既作這些事,還顯什麼神跡給我們看?」.
你憑什麼做這些事?.
你是誰?.
你如何證明你有這個權柄?.
管聖殿的人質問這樣的問題其實有點奇怪.
不過我將它放在現代處境.
猜一猜他們在做什麼.
就好像有兩個金毛仔在後面撥欄街,酒吧醉酒生事.
看搶的社團分子就打殘他們.
根本不想討論,打打打打,然後趕走他們.
但假如有一個人像我這樣.
雖然老老的,不過也黑黑實實的.
跑去澳門葡京酒店.
就這樣打爛他一件雕塑.
那些管事的人一定不敢掉以輕心.
隨便派幾個馬仔就幹掉我.
一定會先問我兄弟,什麼來的?.
你做這樣的行動,你當然可能有背景.
不是猛龍不過江,你竟然跑去聖殿搞事.
所以那些在聖殿管事的人不敢立即對耶穌採取行動.
就去摸摸他的底,探探他的路.
Sodicore顯個神跡來看.

$^{441}$猶太人問這樣的問題.
顯示他們不是對舊約的應許一無所知.
如果眼前這個人潔淨聖殿.
他是不是舊約聖經所應許要來的上帝使者呢?.
你是不是那個上帝使者?.
你是不是真的要來潔淨聖殿呢?.
耶穌毫不猶豫就說了.
十九節,你們拆毀這殿.
我三天內要再建立起來.
由於猶太人問的是耶穌對聖殿有什麼權柄.
耶穌就強調他有建造聖殿的能力.
並且是超強能力.
你們試試整垮這個聖殿.
我三天內就可以將它重建起來.
這番說話是不是耶穌的話呢?.
我們可以再想想.
馬可福音的技術是.
當耶穌接受大祭司的審判.
就是要查證他犯過什麼罪的時候.
有人就這樣指證.
馬可福音14章58節.
他們指證耶穌曾經公開宣告.
我要拆毀這人手所造的殿.
三天內另造一座不是人手所造的.
如果那些人轉述耶穌的話是對的.
那麼耶穌在聖殿說的不是.
你們拆毀聖殿,我將它重建.
而是說我會拆毀聖殿,然後我將它重建.
到底耶穌說的是你拆毀還是我拆毀呢?.
我不知道.
總之我將聖經的資料放在一起給大家看看.
可能耶穌說的是.
我可以拆毀聖殿,然後將它重建起來.
這更顯出耶穌的權柄.
耶穌的狂言當然不為猶太人所相信.
20節他們質疑說.
聖殿是46年才造成的.
三天內就能夠重建起來嗎?.
他們當然聽不明白耶穌所說的話.
耶穌說祂三天重建聖殿.

$^{481}$指的是祂自己的身體.
耶穌的身體是聖殿.
這不是人手所造的,是屬天的.
聖殿的含義是上帝與人相遇.
上帝在聖殿裡與人相遇.
如果耶穌的名字叫做「以馬來利」.
我們說過了.
「以馬來利」的意思就是上帝和我們.
同在是加上去的.
上帝和我們.
這是上帝與人有真正聯繫的地方.
所以耶穌就是聖殿.
就是上帝與人聯繫.
如果聖殿是人向上帝獻祭的地方.
我們也知道.
耶穌一次過為人獻上了挽回祭.
並且是最美的獻祭.
所以耶穌在三天內重建聖殿.
指的是祂在三天內復活.
耶穌不僅僅潔淨了聖殿.
更是取代了聖殿.
並且將聖殿的功能變得完全.
耶穌的說話在那時候.
不僅僅那些猶太人聽不懂.
門徒也聽不懂.
他們要到耶穌復活之後.
將祂的說話與舊約的寓言連結起來.
才逐漸明白耶穌的意思.
22節說.
所以到祂從死裡復活以後.
門徒就想起祂說過的話.
便信了聖經和耶穌所說的.
好,耶穌潔淨聖殿的事件.
我們就講到這裡.
有多少聖經我們就講多少.
這段聖經沒有提到猶太人對耶穌採取進一步的行動.
雖然他們聽不懂耶穌的說話.
但似乎也沒有打算追究到底.
就讓耶穌和門徒離去.
整件事件就這樣不了了之.

$^{521}$23-25節.
約翰反而說了以下的說話作為結束.
當耶穌在耶路撒冷過御節的時候.
有許多人看見祂所行的神蹟就信了祂的命.
耶穌絕不將自己交託他們.
因為祂知道萬人也用不著誰見證人怎樣.
因為祂知道人心裡所存的.
耶穌一直在眾人的誤會中間服侍.
沒有人完全明白.
真正明白耶穌說什麼做什麼.
不僅是他的敵人不明白.
連崇拜他的人其實也不明白.
所以耶穌甘於寂寞.
孤獨前行是我們知道的.
祂不被崇拜祂的人覺得很開心.
也不被攻擊祂的人激怒.
不管是崇拜祂還是攻擊祂.
對耶穌來說.
在耶穌看來都是出於誤會.
如果當時有新聞記者的話.
對潔淨聖殿的行動或者會有這樣的技術.
昨天有一個人職神失常.
在聖殿鬧事.
推倒了很多攤檔.
大概就是潔淨聖殿的行動.
就是這樣的報導.
我們來到應用部分.
耶穌潔淨聖殿的行動對我們有什麼啟示呢?.
很多人會馬上聯想.
潔淨聖殿.
聖殿就是我們新約的教會.
即是我們教會也需要被潔淨.
我不知道大家對你所參與.
所服侍的教會有什麼看法.
我猜在我們所在的處境世界.
有不少人特別是年輕一代.
對教會有很多很多的看法.
很多負面的評價.
這是真的.
我聽過不少年輕人跟我說.

$^{561}$他期待耶穌再來潔淨教會.
推倒我們的祭壇.
推倒我們的講壇.
推倒桌椅.
我不知道是不是這樣.
希望在教會來一次大掃除.
大潔淨.
沒有人在原則上.
反對教會需要被潔淨和更新.
坦白說.
我們也不會說我們的教會是完美的.
就算現在我們不覺得有問題.
是我們看不到的問題.
教會任何時候都需要.
不停被上帝更新和潔淨.
這是永遠不會錯的.
不過關鍵問題是.
要潔淨什麼呢?.
我們教會如今藏污納垢.
包庇罪惡.
我們教會裡面有人男盜女娼.
滿口仁義道德.
不過比不信的人更腐敗.
教會裡面有人競逐名利權勢.
沒有真實的信仰.
沒有真實的道德.
我們沒有人可以否認.
我講廣義的教會.
我不是在說我們王浸.
有失信,失德的事件發生.
基督徒犯罪是不可避免的.
但是.
剛才所說的那些.
什麼男盜女娼,包庇罪惡.
是不是真是我們教會.
普遍而真實的寫照呢?.
這是我覺得要問一問的.
我看到的那些對教會的批評.
通常不是說教會有大量.
藏污納垢的證據.

$^{601}$事實上我甚至會說.
年輕一代對道德操守.
對教義的執著其實不是很高.
對權威,對規條通常持一個嘲諷的態度.
對人的軟弱.
包括對自己的軟弱通常是很包容的.
所以我聽到幾個年輕人.
他們批評教會多數是說教會.
太dogmatic(妒忌).
太排他.
特別是對同性戀.
甚至吸毒的態度.
說我們太惡邪.
他不喜歡教會過分高舉道德.
特別是堅持某一套.
牢固的道德表達方法.
他們覺得這是一個矯情,虛偽,律法主義.
妨礙了人的真性情的流露.
他們最常對教會的批評.
不是教會道德,律法.
是說教會的生活則外過盛.
和生活脫節.
讓人沒辦法活出真我.
我們必須要承認.
我們所處的時代急劇轉變.
兩代之間的文化差距不斷擴大.
年輕一代沒辦法.
適應上一代所形塑的教會文化.
包括敬拜模式,真理和道德的教導.
行政體制,人際關係.
以及教會與社會的關係等等.
這是很真實的事.
跟大家說一些無聊事.
有人就….
我不用Facebook的.
不過有人轉了一段.
我兒子在Facebook上所貼的文字.
我就讀這段文字.
他說早月做了一個關於PK.
PK就是Pastor Kids的分享.

$^{641}$即是我兒子被邀請去一個聚會做分享.
我不知道是什麼聚會.
說一下PK的牧師仔的情況.
分享結束後提問環節.
有人問我.
會怎樣叫已經離開教會.
(堂會)的PK回來.
我沒有想過這個問題.
打了個突,唯有誠實回答.
我兒子在那個聚會的答案.
我不會叫他們回來.
因為我得承認PK作為最了解教會.
運作和黑暗面的人.
離開總有合理的原因.
對不起,以下說的話.
他寫的,很噁心.
他說每家堂會都好像一坨屎.
有咖喱味,有水果味.
只有不同的味道.
而沒有不是屎的教會.
在這前提之下.
選擇不再吃這坨屎是再合理不過的決定.
既然吃屎是不合理的.
我就沒有道理可以說服他們回來.
要他們回來吃屎.
就需要他們先自己找到足夠的理由.
然後甘心情願吃這坨屎.
討論的重點從來都不是.
這是不是屎.
而是有沒有甘心吃屎的理由.
有就留,沒有就走.
很明顯,你看到我兒子比我偏激很多.
這不用說.
不過幸運的是他找到繼續吃屎的理由.
所以他仍然在教會服侍.
他還在讀神學.
他在這篇文章.
其實我不知道頭不知尾.
不知道還有什麼.
因為我沒有看原文.

$^{681}$說到PK離教的問題.
這個比較複雜.
我不在這裡和大家處理.
牧師的二代.
牧二代離教的問題.
這個不是我們今天要說的東西.
我要說的到頭來.
是真的.
信二代離教的情況其實是相當普遍.
如果用北美的數字的話.
通常信徒的第二代不再回教會.
是超過60\%.
華人可能沒有這麼多.
不過也不見得很少.
我其實今早才從加拿大回來.
幫一個教會講夏季營.
在最後一個下午.
也有一對夫婦.
很愛主的弟兄姊妹.
來和我說他們一對子女.
十幾歲之後不再回教會.
說他們的困擾.
我要說這些情況是很普遍的.
如果我們中間沒有這個情況.
大家不知道我在說什麼.
你不需要明白.
沒有就簡單了.
不用討論.
我要說很多人宣稱他們討厭教會.
宣稱他們在教會裡找不到耶穌.
意思就是找不到他們心目中想要的耶穌.
那個可以凡事包容接納.
絕對不批評.
不干涉他們的想法和做法.
可以讓他們做什麼都可以.
或者不做什麼都可以.
這樣的耶穌.
我要說我們大家對教會有不同的期望.
不同的年代對教會有不同的期望.
這是我們必須要正視,承認的事情.

$^{721}$也沒有什麼好辯論.
不過如果將我們對教會的某些挫折或失望.
連接到耶穌基督.
說潔淨聖殿的事件.
我就說我們需要進一步多想一點.
我們對教會到底有什麼不滿呢?.
我們要推倒什麼攤販?.
要趕走什麼買賣行動呢?.
我們可以說我們不喜歡教會的制度.
我們可以說我們不喜歡教會的傳統.
甚至我們不喜歡教會的教會政治.
什麼叫教會政治呢?.
這是很複雜的問題.
可能個別有教會有終身制.
他們覺得有人把持教會的局面.
不過如果多數教會其實都沒有終身制.
那個是一個民主選舉的制度的話.
你就要元祖復蘇了.
你可以說我不喜歡現在的領袖.
不過問題是那些領袖是選出來的.
所以其實關鍵的問題是.
我們還不是要換走什麼.
我們要關心的問題是.
換上什麼.
所有制度都和人事有密切的關係.
制度是次要,人是主要的.
我們常說上帝的方法是人.
我們可以請上一輩的人離開領導層.
不過我要問.
什麼人會接上去.
繼續承擔教會的呼應使命和牧養使命呢?.
如果我們只是批評.
我們不願意毀身.
我們不願意承擔責任.
我們不願意競選執事.
我們將現在的執事拉下馬之後.
教會不如就拉閘了.
我們也說過我們喜歡耶穌的反台行動.
刺激和痛快.
但我們分享了耶穌對聖殿的關切美.

$^{761}$他的門徒就想起經上記著說.
「我為你的殿心裡焦急如同火燒」.
弟兄姊妹.
我想問.
包括我們上一代.
中間的一代.
或者是年輕的一代.
我們對教會有多少關懷.
多少迫切.
多少承擔.
我們如果看到教會有不理想的地方.
或者教會在今天實踐福音使命和牧養使命.
有不到位的地方.
我們的福音傳不出去.
我想問.
我們有多少關切.
我們有多願意真的毀身下去.
我們來.
我們來參與.
上帝幫助我們.
我們可以為教會鋪開一個新的局面.
我們都關心教會的前景.
我們都確認.
今天的情況不算是很理想的情況.
我們都需要大幅度去更新和復興.
不破不立.
改革目前的一些做法是必須的.
我的祈禱是.
我們有更多為教會的前景.
心照如焚的人.
有更多願意參與.
毀身.
承擔.
開放自己.
成為上帝在這個世代.
施展作為的平台的人.
拆毀不是神蹟.
將拆毀了的聖殿重建.
這才是真真正正的神蹟.
我們祈求.

$^{801}$看到這樣的神蹟在我們的世代發生.
我們一起為這件事祈禱.
\newpage



\section{士師記 2:6-3:6}
\label{sec:CAkR8_FGcmQ}
\textbf{2023年7月16日崇拜直播｜吳真真牧師｜「家醜說出口」系列:信仰傳承失效｜士師記2\_6-3\_6}
\newline
\newline
連結: \href{https://youtube.com/watch?v=CAkR8-FGcmQ}{\texttt{https://youtube.com/watch?v=CAkR8-FGcmQ}} ~~~~ 語音日期: 2023-07-17
\newline
\newline
\hyperref[sec:rlyrCSM_7W8]{\small{< < < PREV SERMON < < <}}
~
\hyperref[sec:index_chronic]{\small{[返順時目]}}
~
\hyperref[sec:iklPHMTALSA]{\small{> > > NEXT SERMON > > >}}
\newline
\newline
士師記 2:6-3:6
\newline
\begin{longtable}{cl}
\hline
\hline
章節 & 經文 (和合本修訂版)\\
\hline
2:6 & \begin{tabularx}{0.7\textwidth}{X} 約書亞解散百姓，以色列人回到自己的地業，佔各自的地。 \end{tabularx} \\ \\ \relax
2:7 & \begin{tabularx}{0.7\textwidth}{X} 約書亞在世的日子和他死了以後，那些見過耶和華為以色列所做一切大事的長老還在世的時候，百姓都事奉耶和華。 \end{tabularx} \\ \\ \relax
2:8 & \begin{tabularx}{0.7\textwidth}{X} 耶和華的僕人，嫩的兒子約書亞死了，那時他一百一十歲。 \end{tabularx} \\ \\ \relax
2:9 & \begin{tabularx}{0.7\textwidth}{X} 以色列人把他葬在他自己地業的境內，以法蓮山區的亭拿‧希烈，在迦實山的北邊。 \end{tabularx} \\ \\ \relax
2:10 & \begin{tabularx}{0.7\textwidth}{X} 那世代的人也都歸到自己的列祖。後來興起的另一世代不認識耶和華，也不知道他為以色列所做的事。 \end{tabularx} \\ \\ \relax
2:11 & \begin{tabularx}{0.7\textwidth}{X} 以色列人行耶和華眼中看為惡的事，去事奉諸巴力。 \end{tabularx} \\ \\ \relax
2:12 & \begin{tabularx}{0.7\textwidth}{X} 他們離棄領他們出埃及地的耶和華－他們列祖的神，去隨從別神，就是四圍列國的神明，向它們叩拜，惹耶和華發怒。 \end{tabularx} \\ \\ \relax
2:13 & \begin{tabularx}{0.7\textwidth}{X} 他們離棄了耶和華，去事奉巴力和亞斯她錄。 \end{tabularx} \\ \\ \relax
2:14 & \begin{tabularx}{0.7\textwidth}{X} 耶和華的怒氣向以色列發作，把他們交在搶奪他們的人手中，又把他們交給四圍仇敵的手中，以致他們在仇敵面前再也不能站立得住。 \end{tabularx} \\ \\ \relax
2:15 & \begin{tabularx}{0.7\textwidth}{X} 他們無論往何處去，耶和華的手都以災禍攻擊他們，正如耶和華所說的，又如耶和華向他們所起的誓；他們就極其困苦。 \end{tabularx} \\ \\ \relax
2:16 & \begin{tabularx}{0.7\textwidth}{X} 耶和華興起士師，士師就拯救他們脫離搶奪他們之人的手。 \end{tabularx} \\ \\ \relax
2:17 & \begin{tabularx}{0.7\textwidth}{X} 然而，他們卻不聽從士師，竟隨從別神而行淫，向它們叩拜。他們列祖所行的道，所聽從耶和華的命令，他們都速速偏離了，並不照樣遵行。 \end{tabularx} \\ \\ \relax
2:18 & \begin{tabularx}{0.7\textwidth}{X} 耶和華為他們興起士師，耶和華與士師同在。士師在世的一切日子，耶和華拯救他們脫離仇敵的手。耶和華因他們受欺壓迫害所發出的哀聲，就憐憫他們。 \end{tabularx} \\ \\ \relax
2:19 & \begin{tabularx}{0.7\textwidth}{X} 但士師一死，他們又轉去行惡，比他們祖宗更壞，去隨從別神，事奉叩拜它們，總不放棄他們的惡習和頑梗的行為。 \end{tabularx} \\ \\ \relax
2:20 & \begin{tabularx}{0.7\textwidth}{X} 於是耶和華的怒氣向以色列發作，說：「因為這國違背我吩咐他們列祖當守的約，不聽從我的話， \end{tabularx} \\ \\ \relax
2:21 & \begin{tabularx}{0.7\textwidth}{X} 約書亞死的時候所剩下的各國，我必不再從他們面前趕出任何一個， \end{tabularx} \\ \\ \relax
2:22 & \begin{tabularx}{0.7\textwidth}{X} 為要藉此考驗以色列是否肯謹守遵行耶和華的道，像他們列祖一樣地謹守。」 \end{tabularx} \\ \\ \relax
2:23 & \begin{tabularx}{0.7\textwidth}{X} 耶和華留下那些國家，不將他們速速趕出，也不把他們交在約書亞的手中。 \end{tabularx} \\ \\ \relax
3:1 & \begin{tabularx}{0.7\textwidth}{X} 耶和華留下這些國家，為要考驗所有未曾經歷迦南任何戰役的以色列人， \end{tabularx} \\ \\ \relax
3:2 & \begin{tabularx}{0.7\textwidth}{X} 只是為了要以色列人的後代認識戰爭，教導他們，尤其那些未曾認識這些事的人。 \end{tabularx} \\ \\ \relax
3:3 & \begin{tabularx}{0.7\textwidth}{X} 留下的有非利士的五個領袖，所有的迦南人，西頓人，以及從巴力‧黑門山到哈馬口，住黎巴嫩山的希未人。 \end{tabularx} \\ \\ \relax
3:4 & \begin{tabularx}{0.7\textwidth}{X} 他們是為了要考驗以色列，好知道他們是否肯聽從耶和華藉摩西吩咐他們列祖的命令。 \end{tabularx} \\ \\ \relax
3:5 & \begin{tabularx}{0.7\textwidth}{X} 以色列人住在迦南人、赫人、亞摩利人、比利洗人、希未人、耶布斯人中間， \end{tabularx} \\ \\ \relax
3:6 & \begin{tabularx}{0.7\textwidth}{X} 娶他們的女兒，將自己的女兒嫁給他們的兒子，並事奉他們的神明。 \end{tabularx} \\ \\
[1ex]
\hline
\hline
\end{longtable}
$^{1}$因為「著崇拜」要讀的經訓.
是記載在我們舊約聖經裡的.
《士師記》第二章第六節至三章六節.
我們可以打向我們的聖經.
舊約聖經第299頁至300頁.
《士師記》第二章第六節內容是這樣說的.
「從前約書亞打發以色列人百姓去的時候.
他們各歸自己的地業.
佔據地土.
約書亞在世和約書亞死後.
那些見耶和華為以色列人所行大事的長老.
還在的時候.
百姓都事奉耶和華.
耶和華的僕人.
亂的兒子約書亞.
正一百一十歲就死了.
以色列人將他葬在他地業的境內.
就是在以法連山地的亭南基列.
在加什山的北邊.
那世代的人也都歸了自己的列祖.
後來有別的世代興起.
也不知道耶和華.
為以色列人所行的事.
以色列人行耶和華眼中看來惡的事.
去事奉朱巴力.
彌起了領他們出埃及地的耶和華.
他們列祖的神去後拜別神.
就是四位列國的神.
惹耶和華發怒.
並彌起耶和華去事奉巴力和阿斯他祿.
耶和華的怒氣向以色列人發作.
就把他們交在搶奪他們的人手中.
又將他們賦予四位仇敵的手中.
甚至他們在仇敵面前再不能站立得住.
他們無論往何處去.
耶和華都以災禍攻擊他們.
正如耶和華所說的話.
「由於耶和華向他們所起的誓.
他們必極其困苦」.
耶和華興起示師.

$^{41}$示師就拯救他們脫離.
搶奪他們人的手.
他們卻不聽從示師.
竟隨從叩拜別神.
行了邪人.
速速的偏離他們列祖所行的道.
不如他們列祖順從耶和華的命令.
耶和華為他們興起示師.
就與那示師同在.
示師在世的一切日子.
耶和華拯救他們脫離仇敵的手.
他們因受欺壓.
繞開就哀聲嘆氣.
所以耶和華後悔了.
及至示師死後.
他們就轉去行惡.
被他們列祖更神.
去事奉叩拜別神.
總不斷絕還梗的惡行.
於是耶和華的怒氣向以色列人發作.
他說:因者民違背我吩咐他們列祖所守的約.
不聽從我的話.
所以約書亞死的時候.
所剩下的各族.
我必不再從他們面前趕出.
惟要藉此試驗以色列人.
看他們肯照他們列祖謹守遵行我的道不幸.
這樣,耶和華留下各族.
不將他們速速趕出.
也沒有交付約書亞的手.
耶和華留下這幾族.
惟要試驗那不曾知道與迦南征戰之時的以色列人.
好叫以色列的後代.
又知道又學習.
未曾曉得的戰事.
所留下的就是.
非理士的五個首領和一切迦南人,西頓人.
並住黎巴嫩山的希美人.
從巴力克門山直到哈馬口.
留下這幾族惟要試驗以色列人.

$^{81}$知道他們肯聽從耶和華.
藉摩西吩咐他們列祖的誡命不幸.
以色列人竟住在迦南人,黑人,阿摩尼人,比利西人,希美人,耶布斯人中間.
取他們的女兒為妻.
將自己的女兒嫁給他們的兒子.
並事奉他們的神.
(字幕).
各位弟兄姊妹平安.
早前我看了一篇由崇基神學院副教授.
伯德.培木斯所寫的文章.
那篇文章的題目是「基督教在歐洲的衰落」.
當中提到他帶一些學生到歐洲遊學時.
看到教堂已經轉型成為其他機構和服務場所.
他們都感到很可惜.
他提到一樣不爭的事實.
他說淺行信仰的基督徒人數正在減少.
「位於鄉村或城市社區中心的大教堂內.
參加主日崇拜的會眾中.
都是由白髮滄滄的老人家組成.
他們的虛弱聲音被淹沒在管風琴的轟鳴之中」.
在世界基督教資料庫的統計數字.
可以看到基督教的急劇衰落.
在德國和英國都很明顯.
在1900年,德國人口中有98.56\%信奉基督教.
但到了2015年,只有67.29\%.
英國人口也由1900年的97.44\%降至69.4\%.
甚至統計學專家預測.
到2050年,這兩個國家的基督徒人數會跌至50\%.
一個是16世紀宗教改革的發源地.
一個是19世紀海外宣教運動的基地.
當時的信徒沒想過在21世紀的今天.
他們的信徒人數會大跌至這個地步.
似乎在他們當中,信仰的傳承已經失效.
沒有了一個有效的延續.
當我們身處現在的香港.
會否想到幾十年後,教會是衰落還是復興呢?.
經過幾年的疫情,社會運動.
移民潮對香港,特別是香港的教會都造成很大的衝擊.
我們都在想,領袖離開,信徒稱王不接.
我們有沒有認真想過我們的承傳是怎樣呢?.

$^{121}$今天這段經文正正就是從一個時代走進另一個時代.
當中我們看到以色列人失敗的經歷.
讓我們在他們的失敗當中去反省和學習.
讓我們一起祈禱.
因為我們蒙你的恩典,成為你的子民.
但主啊,我們知道面對世代的確有很多挑戰和衝擊.
到底我們可以怎樣將你在我們身上的這份恩典.
我們對你的確信能夠承傳下去.
我們實在需要你督澤我們,提醒我們.
願意我們預備好我們的心去聆聽你的教導.
願意我們生命有轉化去行出你的心意.
我們這樣禱告交託,奉主耶穌基督的聖名自祈求.
阿們.
我們繼續加醜要說出來的第二講.
我們說到信仰承傳的失效.
事司紀有兩個序言.
上星期Ray牧師已經和我們分享到第一章的第一個序言.
第一個序言是用一個歷史性的基礎告訴我們.
以色列人在約書亞這位領袖死後.
在戰事上,他們怎樣和迦南人爭戰.
攻取地為業.
並且他們容許了迦南人在他們當中.
而今天分享的這段經文.
二章六至三章六節一段很長的經文.
是從一個屬靈的角度.
去說出以色列人怎樣從一個黃金的年代.
走進一個衰敗墮落的惡性循環.
為了接下來的三章七節至十六章三十一節.
做一個蓋覽,做一個神學的註腳.
讓我們一起看看這兩代是怎樣.
大家一起看看第六至十節.
很特別的,這段經文的頭四節.
和約書亞記二十四章二十八至三十一節幾乎一模一樣.
有很多人說士斯記的序等於約書亞記的拔.
是一個共同的連接.
將它特別提到約書亞的時代的延續.
約書亞的時代可以說是一個黃金的時代.
他們能夠走進迦南地,得地為業.
而約書亞也被稱為神的僕人.
下一章的頭影片可以看到.

$^{161}$這個僕人到最終是葬在自己的地獄裡.
連偉大的摩西領袖都做不到的事.
約書亞他做到了.
他打過美好的仗.
他們的這個年代是一個黃金的年代.
是功德圓滿的年代.
不過,我們看到第七節和第十節.
我們看到約書亞死了之後.
仍然有一群見證過耶和華為他們攻取迦南地.
行其事神職的長老仍然在他們當中.
這群以色列民仍然是事奉耶和華的.
不過,當這群人都死了.
這個時代就終結了.
聖經形容有別的世代興起.
是別的世代興起.
這是事司的時代.
就是我們一直以來的事司系列會說下去.
這個時代有一個特點.
去強調這個時代是不知道耶和華.
也不知道耶和華為以色列人所行的事.
以色列人要知道神耶和華的作為.
要知道耶和華是誰.
一定要靠上一代的教導.
在《申命記》六章一至九節.
那裡很清楚地記載了摩西在還沒進入迦南地的時候.
已經要祝福那群以色列人進入迦南地之後.
要千萬千萬地教導他們的後裔.
經文裡說.
「這是耶和華你們的神所吩咐.
要教導你們的誡命」.
《律例典章》.
叫你們在所有過去得遺孽的地上遵行.
要謹守遵行敬畏耶和華.
謹守一切的律例典章.
就是我吩咐你們的.
使你們的日子得以長久.
還教他們第六節.
我吩咐你們將這些話要記在心上.
要殷勤教導你們的兒女.
無論坐在家.

$^{201}$走在路上.
躺在地上.
起來都要淫訟.
要繫在手上作記號.
戴在額上作經鑑.
又要寫在你房屋的門框上.
和你的城門上.
這是摩西在未入迦南地之前.
已經吩咐他的子民.
要一代一代這樣去述說.
耶和華的作為.
要教導他們謹守遵行耶和華的吩咐.
但是顯然這個別的世代.
這個新的世代的以色列民.
沒有按這個吩咐去做.
我想原因有兩個.
第一就是.
當這個別的世代的人.
不需要再爭戰.
再沒有目標.
進入了一個安逸的時代.
又有其他的神明在他們其中的時候.
他們漸漸被同化.
他們忘記了上帝.
這個在《史詩·第三章》七節.
很清楚地說.
他們是忘記了上帝.
第二.
經文亦都強調.
那些見證耶和華恆大事的長者.
長老仍然在的時候.
他們會事奉上帝.
這些人.
都是曾經大大經歷過耶和華.
信仰對於他們來說.
是第一身.
是第一手.
他們確實知道神的信實.
所以他們仍然會鼓勵他們的家.
他們的族.

$^{241}$要這樣事奉神.
但是接下來的第二代.
他們沒有這些親身的經歷.
沒有這些第一手的信仰.
於是就沒有再去教導.
其實這是一個很好的警惕.
值得我們引以為鑑.
有時我們都發覺.
我們的信義代.
所信的是二手的信仰.
他們並不是親身去經歷神.
或者是因為他們的生命.
他們的生活很順利.
資源很豐富.
沒有什麼困難危機.
或者他們有困難危機.
父母.
我們上一代有足夠的能力和資源.
幫他們排難解分.
擋開一切困難.
變成信義代.
我們信耶穌的.
信上帝的父母.
是怎樣供應和愛我們.
但是我們有沒有帶領這班孩子.
去看到其實是上帝的恩手.
我們有沒有幫這班孩子.
和我們的上帝去結連呢?.
剛剛的星期一早上.
也是加拿大時間的晚上.
是寧靈學院院長.
Tim Sir的追思禮.
大家都記得他吧.
他來到我們當中分享.
怎樣去和神親近.
怎樣在靈修中得力.
當日他的太太.
也是一個傳道人.
Pastor Jessica.
她分享到.

$^{281}$原來Tim Sir的家.
每一天都會在出門前.
一個小時.
一家人一起吃早餐.
一起分享.
一起祈禱.
Tim Sir是最早起床的人.
他五時多已經起床.
因為他自己.
首先享受和上帝.
獨處安靜的時間.
然後為家人預備早餐.
然後在早餐的過程中.
分享上帝的話.
他在當中的領受和得著.
帶動整個家庭的屬靈氣氛.
讓他兩個女兒.
讓他的太太.
一起去認識上帝.
Pastor Jessica提到.
很多時候兩個女兒會問.
「爸爸,我可以怎樣做?」.
「我應該怎樣決定?」.
通常Tim Sir會跟她說.
「不是我要告訴你怎樣決定」.
「最重要的是上帝想你怎樣」.
「你問上帝吧」.
Tim Sir的榜樣.
的確是一個很好的提醒.
他是一個示範.
告訴我們.
我們要怎樣.
將我們所信靠的上帝.
去教導他們.
讓我們的下一代認識他.
知道信實的上帝是怎樣.
讓他們親自去經歷上帝.
不過很可惜.
回看聖經.
我們看到的是以色列人.

$^{321}$當他們不知道上帝.
不知道耶和華的時候.
帶來了一個很嚴重的後果.
我們看到第二章的十一至十九節.
我們說的是這個年代.
是一個degeneration的年代.
是一個熟齡衰退的年代.
是一個向著沉淪的漩渦.
一直下陷的年代.
主要的原因是一個.
我們回看第十一至十三節.
這裡說.
以色列人行耶和華.
眼中看為惡的事.
這個說話很關鍵.
因為在接下來的士司紀裡.
每當一個士司出現的時候.
你就會看到這一句說話.
甚麼是行耶和華眼中看為惡的事呢?.
這裡很清楚地說.
就是他們離棄上帝.
叩拜偶像.
你會看到和修版有一個很特別的.
就是他逆出一個字.
那個字是王合本沒有寫的.
就是他們去事奉其他的偶像.
朱巴力,阿斯他祿和巴力.
他們隨從.
跟著這些偶像.
事奉這個字很重要.
事奉這個字是這兩個世代.
以色列人興衰成敗的關鍵字.
在二章六字到三章六字.
我們今天讀的這段經文.
事奉這個字出現了五次.
這五次裡只有我們剛才讀的第七節.
事奉的對象是耶和華.
是約書亞那一代.
見證著耶和華為以色列人行大事的領袖.
這些人事奉耶和華.

$^{361}$但接下來的篇幅有四次提到事奉.
事奉的對象都是外邦的神.
是別神.
這就是關鍵.
什麼是事奉呢?.
事奉的原文.
我們在舊約裡的翻譯.
我們會看到.
特別在約書亞記和士師記.
會翻譯成「瀑人」.
「獻」字是「獻」字.
「服侍」.
事奉是最多的.
這裡告訴我們.
當我們事奉一件事.
事奉一個神明.
就代表了對他的效忠.
你是敬拜他.
你當他是你的主.
你就成為他的瀑人.
你就為了哄他開心.
去服侍他.
其實約書亞記23至24章.
是約書亞和以色列人.
最後臨終的遺命.
他在其中14至24章.
第24章14至24章.
那裡短短的11節經文.
就出現了14個事奉.
在這短短的經文裡.
事奉就是他們與神的立約.
當我們提到這群以色列人.
離棄上帝.
不再事奉他.
其實就是與神之間的約.
他們主動毀這個約.
離棄曾經為他們行騎士神蹟的耶和華.
那個在埃及裡.
救他們從奴役的生涯.
而得到釋放的耶和華.

$^{401}$這一代的以色列人.
他們選擇誤再事奉耶和華.
他們事奉那些假神.
做假神的僕人.
去哄這群神明開心.
其實當他們進入迦南地.
他們沒有趕走迦南人.
由得他們在各個支派裡生活.
當他們由曠野約旦河東.
要進入這個流乃與物之地之前.
他們過著的是束木的生活.
但當他們進入這個新的地方時.
他們就變成一個霧濃的生活.
他們經歷霧濃的生活.
就受到與他們同住在一起的.
迦南人的文化所影響.
他們漸漸被同化.
這群迦南人是霧濃的.
他們會服侍,敬拜一些神明和偶像.
而這些偶像其實是因應他們實際的環境需要.
而想像出來的神明.
巴力.
本來的意思就是主人.
巴力本來就是主人.
而他就是眾多的迦南人敬拜的神.
你們看到經文裡說的朱巴力.
其實他們是拜很多不同的神明.
按照他們的需要拜不同的神明.
風暴的神,新玉的神,賜豐收的神.
因為他們是霧濃的.
所以他們奉巴力作為農作物豐收的施與者.
所以以色列人當他們走入一個霧濃的世界的時候.
他們就像他們一樣去拜這些偶像.
確保他們能夠在農作物裡有豐收.
再加上他們進去之後.
可能他們在《士師記》之後也告訴我們.
其實他們完全不知道耶和華.
不過《士師記》裡告訴我們.
他們所認識的耶和華對於他們來說是知不知道的.
他們只當耶和華是諸神裡其中的一個.

$^{441}$最重要的目的就是可以祝福到自己.
所以以色列人在加拿大眾多偶像當中.
他們去朝拜諸巴力不同的神明的時候.
他們其實以他們為主人.
不過有些解經學者告訴我們.
當他們不斷去拜其他不同的神明的時候.
其實真正的主人是自己.
因為是為了滿足自己的需要.
為了滿足自己的慾望.
其實我們反省到基督徒未必口中會說他不信上帝.
我們不會說我們不信上帝.
我們不會說我們拜其他東西.
不過有些時候我們要小心.
我們雖然口中說我們是敬拜神.
但心裡其實是敬拜自己.
我們活在一個這麼多元.
這麼吸引,五光十色的社會裡.
有很多東西吸引我們.
Richard Forster說得很對.
他說基督徒有三大引誘.
Money, Sex, and Power.
金錢,權力和性.
這些引誘在我們當中.
久而久之我們就會以這些東西成為主人.
去事奉這些東西.
我們會不知不覺成為物質的奴隸.
享樂的奴隸,金錢的奴隸,榮譽的奴隸,性的奴隸.
人生不知不覺地以這些為目標.
當我們人生的優先次序選擇的時候.
我們就會以這些東西為先.
反而將我們的上帝放在一旁.
這真是值得我們好好地去警惕.
然後我們看到下一段經文.
我們看到第十四節和二十節.
經文告訴我們.
因為以色列人的背弱引起了神的怒氣.
以至耶和華跟以前不同了.
他以前會為以色列人爭戰.
為以色列人剷除他的敵人.
現在他不會了.

$^{481}$他還要將以色列人交在敵人的手中.
任由他們搶奪,攻擊.
甚至他自己也降災禍給這些以色列人.
有很多人接受不到.
這位充滿慈愛的神會有憤怒.
不過我覺得我們要留意一件事.
這個憤怒是基於他愛這群子民.
跟他立的約.
我們搞錯了的原因.
不讓神憤怒的原因.
是因為我們往往會將憤怒和破壞力.
連結在一起.
但憤怒不一定是破壞.
不一定是無理的發洩.
當面對跟神立約的子民.
作出不忠的背叛.
憤怒絕對可以理解.
但神的憤怒所作出的舉動.
他要放敵人在他們當中.
對他們做搶奪.
有災禍淋到他們.
令他們站不穩.
這是以色列人的選擇和結果.
就像他們15節所說.
「就如耶和華向他們所起的誓」.
耶和華在《申命記》27章15至26節.
很清楚地說.
背若的要懲罰他們.
這就是根據.
所以耶和華的憤怒的發作.
他寫的「發作」.
其實是以色列人背若的結果.
就像一個學生.
說明了在學校的校規.
他不守這個校規.
最終他就要承擔這個校規.
清清楚楚列明的處罰.
並不是說.
「哦,訓導主任發脾氣,心情不好」.
然後就亂罰你一樣.

$^{521}$當耶和華怒氣發作.
以色列人自己走進了一個.
萬劫不復的循環.
就像在第16至19節所說.
他們離棄上帝,師封必慎.
導致他們走進惡性的循環.
第16至18節.
形容他們.
因為神將以色列人交在敵人的手中.
他們就很困苦.
耶和華興起事司救他們.
但以色列人沒有珍惜.
很快就叛逆.
後來又受迫.
他們哀聲嘆氣.
結果到最後.
神還是拯救他們.
不過聖經裡告訴我們.
第19節.
他們並沒有因為神不斷的拯救而回轉.
最終形容他們.
比他們列祖更甚.
去事奉叩拜別神.
總不斷絕還緊的惡行.
其實這裡說的列祖.
並不是約書亞時代的列祖.
而是事司時代裡.
一代又一代的先祖.
他們越來越差.
越來越沉淪.
這短短幾節的經文.
就是事司紀三章七至十六章三十一節的寫照.
所以他們步入一個背若.
受欺壓,悔改,拯救.
然後又再背若.
受欺壓,悔改,拯救的循環裡.
弟兄姊妹.
這一群不知道神的以色列人.
越來越沉淪的警方.
有一個美國的牧者叫Jerry Vines.

$^{561}$他說到這些就是以色列人進入.
degeneration.
就是一個衰退的情況的局面.
進入一個靈性越來越衰退的情況.
但是他提到.
其實我們可以從事司紀這一段經文裡學習.
以至我們不用再degenerate.
我們不用衰退.
而是我們可以regenerate.
我們可以有一個屬靈的恢復.
我們有一個屬靈的再生.
我們可以不用被這個世代同化.
而是我們反而可以影響這個世界.
我們再看一個很關鍵的章節.
就是在第十八節.
他們提到一個字.
就是當我們看到.
為什麼我們可以有regenerate的機會.
其實是因為我們有一位恩慈的上帝.
如果不是這個恩慈的上帝.
我們只會仍然看到.
我們看到這個恩慈的上帝.
我們仍然看到他可以給世代的意識連人機會.
如果不是我們剛才說的惡性循環.
根本不會發生.
早就已經斷了纜.
關鍵就在於第十八節這個字.
後悔.
因為耶和華看到他們極其痛苦.
他就後悔.
後悔這個字.
原文是breathing deeply.
深深地吸一口氣.
這表達什麼?.
這表達神極其悲傷.
因為意識連人的不忠.
所以上帝很悲傷.
神很愛他們.
雖然知道他們罪有應得.
不過他仍然改變心意.

$^{601}$去派事司去拯救他們.
其實有時真的很難理解.
神對他們的愛.
因為一群如此頑劣的人.
他仍然盡方法去挽回他們.
我們看接下來的二章二十一節.
到三章四節.
這段經文裡.
二章二十一節.
提到「試驗」.
其實這裡作出了一個解釋.
就是為什麼神不親自趕走.
以色列的仇敵.
迦南的仇敵.
這裡很清楚地說.
其實是為了一件事.
就是「試驗」.
就是為了「試驗」.
在二章和三章.
總共出現了三個「試驗」.
這三個「試驗」的目的是什麼呢?.
我們會看到.
第一個「試驗」所說的.
是二章二十二節.
和三章四節裡說的「試驗」.
「試驗」就是他們.
是否願意像列祖一樣.
去謹守,遵行祂的道.
意思是什麼呢?.
意思就是.
其實對於這個世代的人的屬靈質素.
是需要讓他們遇到一些困難和逆境.
才是真正可以「試驗」出來的.
所謂「真金不怕紅爐火」.
純正的金是不怕火煉.
反而是越煉越正.
在這個情況下.
如果所有的迦南人都被趕出.
根本以色列人不用選擇.
只好跟從耶和華.

$^{641}$但是迦南人的存在.
就讓以色列人去作一個選擇.
究竟去跟還是不跟上帝.
而跟不跟上帝.
最重要的是.
他們是否願意聽神的話.
他們究竟是否願意在任何情況下.
在順境,逆境,艱難的情況下.
都願意信服神的帶領和心意呢?.
很多時候人會在他們遇到困難的時候.
去找神明.
在我們今天的社會裡.
當人失去的時候.
失去健康,失去親人,失去工作.
失去金錢,失去尊嚴.
很多時候就會找東西來填補.
但是一個真正跟從神的人.
他看萬事為糞土.
他以基督耶穌為至寶.
其實當我們.
這位又愛我們.
又滿有大能,又滿有智慧的上帝.
成為我們生命中的中心的時候.
祂的同在,祂的引導.
就已經成為我們最大最大的滿足.
成為最大最大的安全感.
其實這就是考驗我們.
是否緊緊跟隨神的時候.
然後第二個試驗.
第二個試驗就是在三章一至二節.
因為他們不知道當時耶和華.
是怎樣帶以色列人和迦南人征戰.
所以上帝就用了這個考驗.
去幫他們知道和學怎樣去征戰.
這次這個試驗.
其實是不同於剛才.
剛才的試驗是試他們.
究竟是否真的跟隨神.
這個試驗有第二個目的.
這個試驗的目的.

$^{681}$就是要教導他們一些真理.
讓他們除了在頭腦上知道.
變成真知道.
就好像臨到西乃山.
百姓見到雷轟,閃電,覺醒.
山上無煙.
他們就發戰.
但摩西跟他們說.
不用懼怕.
因為神降臨是要試驗你們.
叫你們時常敬畏祂.
不致犯罪.
這個試驗是耶和華作出教導.
教導以色列民要敬畏祂.
不要犯罪.
同樣在這個試驗裡面.
有一個目的.
就是這班以色列民.
在安迪的以色列民.
從來沒有實戰經驗的以色列民.
從以前不懂.
沒有人教.
沒有實踐過.
可以學習怎樣在戰場裡面打仗.
真真正正讓他們從以前的不知道.
變成真知道.
不過很可惜.
三章五至六節告訴我們.
神那麼愛他們.
一次又一次給他們機會.
到最終他們仍然選擇偶像.
跟他們迦南人通婚.
最重要的是他們通婚後.
跟隨他們去拜他們的偶像.
這個就是他們的結局.
但是今世的基督徒.
我們千萬不要學他們.
我們應該怎樣去回應.
神這份不離不棄的愛呢?.
剛才已經說過.

$^{721}$兩個試驗的目的.
其中一個基督徒跟從神.
事奉神的人的標記就是.
聽命和信服.
聽命和信服很難學.
在這麼多東西要選擇的世代.
你不去選擇世界.
而是去選擇神.
真的不是一件容易的事.
但我認為福塞特.
去呼應馬丁路德的一個觀點.
很有意思.
他這樣說.
上帝兒女最高的呼召.
是在這個世界.
但我們不屬於這個世界.
正如船在海上是安全的.
只要海水不在船上.
上帝的兒女在這個世界是安全的.
只要世界不在上帝的兒女的心裡.
弟兄姊妹.
我們要的不是讓這個世界影響我們.
這個世界是我們事奉的工場.
當我們心裡有上帝.
走他的道.
其他未認識神的人.
反而會藉著我們去認識這位上帝.
他會給我們影響.
他會想認識我們.
認識我們所信靠的上帝.
然後學我們一起去事奉他.
敬拜他.
最近認識一個弟兄.
他本來是一個很有資歷的專業人士.
他可以賺很多錢.
也成為他業界舉足輕重的人物.
但當他有一個重新的機會.
他去思想道.
被神提醒了他.
去思想道的就是.

$^{761}$其實更重要的是他要事奉神.
要將人靈魂的需要.
放在他生命裡最重要的優先次序.
於是他就用他的專業知識去服侍人.
不是為了賺錢.
不是為了名譽.
反而是用一個極商宜的價錢.
去服侍一些有需要的人.
因為他心裡最重要的優先次序有兩個字.
就是福音.
其實我覺得真正認識神.
我們一定要捲起衣袖.
絕對不可以坐著坐享其成.
如果我們單單去崇拜.
單單聽到.
聽別人說道.
我們很容易忘記.
我擔保下星期你已經忘記我今天說過什麼.
但是如果我們不是單單坐著.
我們是捲起衣袖.
我們去毀身.
我們去投身.
我們去訓身.
我們才有機會真知道神.
我們才會不忘記祂.
所以重生.
不會我們坐著就不用捱打.
我們需要站起來去回應這個世界.
我們要在當中去事奉神.
經歷神.
去影響這個根本站不住腳的世界.
剛剛那個星期五.
我有機會和桂慧即時聊天.
她說到很多年前.
我們當中有一群肢體.
去到一個叫雲嶺的地方做短孫.
她說那個短孫的環境極之艱難.
很刻苦.
去廁所和洗澡的那一格.
其實是那麼的丁方.

$^{801}$但是他們一起.
有同一的心智.
要去宣揚主的愛.
經歷上帝的信實和恩慈.
和祂一起打仗.
一起去經歷.
在整個歷程裡.
她說了很多次很多次那個字.
她說不會忘記.
不會忘記.
她說了很多次.
我相信不單單是桂慧姐妹.
不單單是桂慧執事.
我相信當年和她一起去短孫的弟兄姊妹.
都不會忘記.
因為我現在還看到他們每一個.
很忠心地在我們當中.
去服侍神,見證神.
還是很積極的為神去事奉神.
所以弟兄姊妹.
我自己在這裡呼籲大家.
如果我們不想.
我們屬靈的世代一直衰落.
反而我們是復興重生.
關鍵就是我們有沒有將.
上帝的信實告訴別人.
我們是否事奉我們的神.
我們是否以祂為中心去服侍祂.
在這麼多姿多彩的世界裡.
我們選擇神.
去行祂的信仰.
宣揚祂.
去敬拜祂.
去回應這個極度自我中心的世代.
弟兄姊妹.
我呼籲黃浸的弟兄姊妹.
一起捲起衣袖.
全民皆兵去服侍我們的上帝.
我們一起祈禱.
天父,我們要在你面前立志.

$^{841}$我們不要單單坐在這裡.
過安逸的生活.
我們要的是.
我們全心全意奮身去事奉敬拜你.
否則我們就會忘記你.
忘記你的慈愛.
忘記你的信實.
我們就會跟隨這個世代的影響去走.
天父,我求你幫我們.
我們希望我們的信仰能夠得以承傳.
我們希望我們能夠做一個.
奮勇去事奉你的人.
求你親自去引領.
求你興起我們當中每一個.
我們這樣禱告,交託.
奉主耶穌基督的聖名字祈求.
\newpage



\section{士師記 8:22-35}
\label{sec:iklPHMTALSA}
\textbf{2023年8月6日崇拜直播｜陳明泉牧師｜「家醜說出口」系列:竟為下代設下網羅｜士師記8\_22-35}
\newline
\newline
連結: \href{https://youtube.com/watch?v=iklPHMTALSA}{\texttt{https://youtube.com/watch?v=iklPHMTALSA}} ~~~~ 語音日期: 2023-08-07
\newline
\newline
\hyperref[sec:CAkR8_FGcmQ]{\small{< < < PREV SERMON < < <}}
~
\hyperref[sec:index_chronic]{\small{[返順時目]}}
~
\hyperref[sec:Fr77s7J2G5I]{\small{> > > NEXT SERMON > > >}}
\newline
\newline
士師記 8:22-35
\newline
\begin{longtable}{cl}
\hline
\hline
章節 & 經文 (和合本修訂版)\\
\hline
8:22 & \begin{tabularx}{0.7\textwidth}{X} 以色列人對基甸說：「你既然救我們脫離米甸的手，願你治理我們，你的兒子孫子也治理我們。」 \end{tabularx} \\ \\ \relax
8:23 & \begin{tabularx}{0.7\textwidth}{X} 基甸對他們說：「我不治理你們，我的兒子也不治理你們，耶和華會治理你們。」 \end{tabularx} \\ \\ \relax
8:24 & \begin{tabularx}{0.7\textwidth}{X} 基甸又對他們說：「我有一件事求你們，請你們各人把所奪的耳環給我。」因敵人都戴金耳環，他們是以實瑪利人。 \end{tabularx} \\ \\ \relax
8:25 & \begin{tabularx}{0.7\textwidth}{X} 以色列人說：「我們情願送給你！」他們就鋪開一件外衣，各人將所奪的耳環丟在上面。 \end{tabularx} \\ \\ \relax
8:26 & \begin{tabularx}{0.7\textwidth}{X} 基甸所要求的金耳環，重一千七百舍客勒金子。此外還有米甸王所戴的月牙圈、耳環，和所穿的紫色衣服，以及駱駝頸項上的鏈子。 \end{tabularx} \\ \\ \relax
8:27 & \begin{tabularx}{0.7\textwidth}{X} 基甸以此造了一個以弗得，設立在他的本城俄弗拉。全以色列就在那裡拜這以弗得行淫，這就成了基甸和他全家的圈套。 \end{tabularx} \\ \\ \relax
8:28 & \begin{tabularx}{0.7\textwidth}{X} 這樣，米甸就被以色列人制伏了，再也不能抬頭。基甸還在的日子，這地太平四十年。 \end{tabularx} \\ \\ \relax
8:29 & \begin{tabularx}{0.7\textwidth}{X} 約阿施的兒子耶路巴力回去，住在自己家裡。 \end{tabularx} \\ \\ \relax
8:30 & \begin{tabularx}{0.7\textwidth}{X} 基甸有七十個親生的兒子，因為他有許多妻子。 \end{tabularx} \\ \\ \relax
8:31 & \begin{tabularx}{0.7\textwidth}{X} 他在示劍的妾也為他生了一個兒子，基甸給他起名叫亞比米勒。 \end{tabularx} \\ \\ \relax
8:32 & \begin{tabularx}{0.7\textwidth}{X} 約阿施的兒子基甸年紀老邁而死，葬在亞比以謝族的俄弗拉，他父親約阿施的墳墓裡。 \end{tabularx} \\ \\ \relax
8:33 & \begin{tabularx}{0.7\textwidth}{X} 基甸死後，以色列人又去隨從諸巴力而行淫，以巴力‧比利土為他們的神明。 \end{tabularx} \\ \\ \relax
8:34 & \begin{tabularx}{0.7\textwidth}{X} 以色列人不記得耶和華－他們的神，就是那位拯救他們脫離四圍仇敵之手的， \end{tabularx} \\ \\ \relax
8:35 & \begin{tabularx}{0.7\textwidth}{X} 也不照著耶路巴力，就是基甸向以色列所施的恩惠善待他的家。 \end{tabularx} \\ \\
[1ex]
\hline
\hline
\end{longtable}
$^{1}$我們以下要聽神給我們的說話.
今天我們要聽神的說話.
是記載在《釋詩記》第八章22至35節.
神是這樣說.
以色列人對基澱說.
你救我們脫離美顛人的手.
願你和你的兒孫管理我們.
基澱說:我不管理你們.
我的兒子也不管理你們.
唯有耶和華管理你們.
基澱又對他們說.
我有一件事求你們.
請你們各人將所奪的耳環給我.
原來受的是以實瑪利人.
都是戴金耳環的.
他們說:我們情願給你.
就鋪開一件外衣.
各人將所奪的耳環雕在其上.
基澱所以出來的金耳環.
重一千七百四克.
戴金紙.
此外還有美顛王所戴的月環.
耳墜和所穿的紫色衣服.
並落駝項上的金鏈紙.
基澱以此製造了一個伊弗德.
設立在本城和弗拉.
後來以色列人拜那伊弗德.
幸了邪人.
就作了基澱和他全家的網羅.
這樣米顛人被以色列人制服了.
不再抬頭.
基澱還在的日子.
國中太平四十年.
二十九節.
約阿斯的兒子耶路巴力回去.
住在自己家裡.
基澱有七十個親生的兒子.
因為他有許多的妻.
他的妾住在士劍.
雅吉他生了一個兒子.

$^{41}$基澱與他起名叫阿比米勒.
約阿斯的兒子基澱.
年紀老邁而死.
葬在阿比爾籍的俄弗拉.
在他父親約阿斯的墳墓裡.
基澱死後.
以色列人又去跟隨朱巴力行邪淫.
以巴力比利圖為他們的神.
以色列人不記念耶和華他們的神.
就是拯救他們脫離四位仇敵之手的.
也不照著耶路巴力.
就是基澱向他們所施的恩惠.
厚待他的家.
聖經獨筆.
弟兄姊妹平安.
平安.
我們一齊來到再看時事記裡面的記載.
剛才一路讀的時候.
你會看到基澱是一代的領袖.
但是在他的結束.
或者在時事記的作者對他一生的評論.
你會看到是一個比較負面的評論.
一代的偉人.
一個戰績彪炳的時事.
但是在歷史公允的審視之下.
他留下一個不光彩的生命故事.
在時事記的記載裡.
我們很多時候雖然會集中在時事如何打仗.
如何擊敗敵人.
這是其中一個很重要的.
整個以色列家發展的命脈.
不過我們也不能忽視的一點.
在時事記的記載裡.
他不單止在說以色列國如何去爭戰.
不論是打贏還是打輸.
他更加說到時事會進入他的心靈空間.
我們之前讀過耶路巴力的講章.
我們看到基澱的生命的掙扎.
那種懼怕和憤怒.
他如何被疏解.

$^{81}$如何踏前一步.
但這些懼怕其實還是伴隨在不同的戰役當中.
當他打勝仗之後.
他還有很多問題在他的家庭裡都呈現了出來.
希望我們今天較長的篇幅.
我們由第七章開始看.
看到第八章尾.
我們從一個比較整全的經歷來看.
這麼長的經文很難完全仔細去解.
不過我們就集中看看他的家庭故事.
看看他的掙扎.
甚至乎看看他講什麼話.
到他能否兌現.
我們就看看基澱的故事.
我們一起禱告.
求聖靈幫助我們.
讓我們明白話語和心祈禱.
親愛的天父,我們願你在我們當中.
求你指教我們.
藉著舌絲記.
好像一面鏡子一樣.
照出我們生命的真相.
幫助我們.
不論正面或負面的記載.
都能夠為我們生命帶來正面的造就.
我恭敬將我們領受你話語的時間交託給你.
幫助每位弟兄姊妹.
我們一起聆聽的時候.
我們心存敬畏的心.
來看待你的話語.
有樣宣講的.
我們謹慎.
我們按著你的真意.
將你的話語辨明.
求你指教我們.
獻上禱告奉求基督耶穌得勝的聖名.
Amen.
我們繼續看.
我想大家幫我.
如果你有聖經或者方便的話.

$^{121}$如果你記得之前基甸的別名.
一個綽號叫耶路巴力.
記得了嗎?.
上一次第六章裡面說到耶路巴力的意思.
就是說讓巴力與他爭辯.
後來基甸用這個名貫穿了.
整個舌絲記裡面對他的記載.
因為之前他是沒有膽量又憤怒的人.
上帝給他一個空間.
藉著憲制來驗證.
跟他說話的那個就是上帝.
然後基甸就有多一點的勇氣.
一步又一步.
上帝甚至要拆毀了他父家的偶像.
之後就拆了巴力的偶像.
將偶像當柴燒.
然後獻祭給耶和華上帝.
所以他有這個綽號.
你記得這個名字了.
幫大家有一點記憶.
耶路巴力後來就成為了基甸.
不單止是他在民族當中的綽號.
也成為了在經文裡面多次提及.
一講基甸的時候就講耶路巴力這件事.
來提醒我們每一個讀者.
他有這個身份的.
他有這個綽號的.
他怎樣演繹他生命的這個耶路巴力的故事.
如果你見到第七章第一節.
在這裡裡面他第一件事就說.
耶路巴力就是基甸.
他用這個綽號.
用這個身份來講他自己的故事.
在耶路巴力這個故事裡面.
我想大家再跳前一點.
在《士師記》裡面.
在六章第23至24節.
寧光幕你會看到.
有一個經文的六章23至24節.
在經文裡面.

$^{161}$基甸因為要驗證耶和華.
對他說的是不是耶和華上帝.
所以耶和華對他說.
你放心不要懼怕.
你必不致死.
於是基甸在那裡為耶和華築了一座壇.
起名叫耶和華沙龍.
然後括號這個字是很值得大家留意的.
這壇在阿比爾籍的俄弗拉直到如今.
為什麼會提及一個這樣的祭壇.
叫耶和華沙龍呢.
很正面的耶和華沙龍.
上主賜下平安.
Shalom.
很正面的.
他特別告訴後代.
不過這個祭壇不是在敬拜的中心.
而是在俄弗拉這個地方.
大概就是基甸的根據地.
在真正的敬拜中心裡.
還有另設一個敬拜中心.
而這個敬拜中心.
是耶路巴黎或基甸做的事來驗證上帝.
但這個壇某程度上也反映了基甸.
他不單止驗證上帝.
他也有意將敬拜中心.
由原本的敬拜中心或祭司的制度轉移.
去到他自己方便的地方.
待會你再讀就會明白.
為什麼有這樣一句說話.
不過你記住.
這個壇是基甸所建的.
由阿比爾族的俄弗拉直到如今.
讀者那時候知道的.
還有這樣的一個祭壇.
接著我們跳回耶路巴黎.
耶路巴黎就是基甸七章一至三節.
他和一切跟隨他的人早晨起來.
在哈律全旁的安營.
米顛營在他們北邊的平原靠近摩利江.

$^{201}$耶路巴對基甸說.
跟隨你的人過多.
我不能將米顛人交在他們手中.
免得以色列人向我誇大說.
是我們自己的手救了我們.
現在你要向這些人宣告說.
凡懼怕膽怯的可以離開基列山回去.
於是有二萬二千人回去.
只剩下一萬.
在第七章第一至第三節就交代了.
試驗了斬了巴力將之後.
上帝不是純粹叫他拆了偶像.
就完成了任務.
作為一個事司.
他就要為當時壓迫他的米殿人開始要征戰.
在經文裡面一路交代.
基甸就一路招兵買馬.
或者一路招聚不同的宗族.
來組成一隊軍隊.
去到不同的地方.
在第六章後已經說到不同的地方.
招聚軍隊.
但去到第七章某程度.
軍隊已經齊營.
如果你計算一條數.
二萬二千人回去再加一萬.
有三萬多人一次過去催雞.
有三萬多人的軍隊.
可以組成來對抗米顛的軍隊.
不過當他齊了腳.
大概他招聚了三萬多人的軍隊之後.
上帝對基甸說.
跟隨你的人實在太多了.
如果你在座裡面打仗的人也好.
事奉的人也好.
人多就越多越好.
做任何事.
由我們第一天開始.
我們的直覺.
我們知道人多就好做作.

$^{241}$所以人多好辦事.
我們經常都希望集結更多的人.
來做更大的事.
但在上帝的角度裡面說.
不是.
面對米殿的外族.
你人太多了.
所以上帝就說叫基甸勸退有一部份的人.
不過上帝的意思是.
為什麼要叫他勸退那些願意從軍的人呢.
因為耶和華對基甸說.
跟隨你的人太多.
我不能將他們交在你手裡.
免得以色列人向我誇大.
上帝看重的不是兵力的懸殊問題.
而是說.
如果你叫那麼多人去打聖仗.
你就當作是自己的能力.
當作是人自己覺得.
人可以有辦法去打平外族.
所以上帝就叫他們勸退.
所以基甸很聽話.
當基甸聽到上帝說這番話之後.
然後就向那些人說.
怕的就回家.
原來三萬多人裡面有二萬多人是怕的.
三分二的人是怕的.
所以有二萬二千人就回家.
剩下一萬一萬的軍隊.
一萬軍隊面對著米顛人的軍隊.
大概有多少人呢.
在第八章就說到.
其實米顛的軍隊好像蝗蟲一樣.
遍滿全地.
如果按數字在第八章有記載.
大概十三萬人.
一萬人對十三萬人多不多.
一萬人對十三萬人都很實力懸殊.
都是很不足夠的.
原本三萬多對十三萬.

$^{281}$都好像可以打的.
有些戰術.
然後上帝說太多了.
剩下一萬人.
剩下一萬人之後.
面對十三萬大軍.
上帝都說都是太多了.
七章第七節至第八節.
耶和華對基甸說.
我要用這點水的三百人拯救你們.
將米顛人交在你手.
其餘的人可以各歸各處去.
這三百人就帶著食物和割.
其餘的以色列人基丁都打發他們各歸各的帳篷.
只留下這三百人.
米顛營在他們下面的平原裡去.
上帝說上帝的減法.
上帝的打仗方式.
他要告訴以色列和以色列的全軍知道.
不是靠人多.
不是靠人體的實力.
一萬大軍他都覺得嚴多.
於是他就從喝水的動作.
有人用手捧水的有三百人.
就用這三百人來.
有些人喝水就按地.
舌頭舔舌頭像小狗一樣.
他就這樣分配.
當然有些解經學家解多了.
舔舌頭像小狗一樣喝水.
是這樣形容的.
那些人打仗不厲害.
所以上帝就撥走他們.
只剩下三百最精良的軍隊.
這些是解多了的.
因為上帝從頭到尾.
他不是說要人的實力的問題.
那三百人為什麼會被揀選出來呢.
單單是因為喝水的動作.
不是因為他有什麼戰爭的實力.

$^{321}$我又不知道這樣拿著喝水的人的戰鬥力.
會強過這樣喝水.
你在軍隊在軍校裡面.
你不會見到有這樣的現象.
不是這樣來判斷一個人.
大概只是上帝說.
就好像包踏小卒一樣.
誰少就揀選誰.
就揀選這三百人來打仗.
在經文裡面.
其實第一件事.
上帝要以色列人明白的一個功課.
或者幫基奠要明白的一個功課就是.
根本打勝仗的關鍵不是在於自己自身的實力.
關鍵在於上帝是否與這班人同在.
即使實力超級懸殊.
面對十三萬大軍.
三百人如果有上帝同在的時候.
三百人都可以打一個徹底令敵軍潰敗的勝利之戰.
不過今天我們看事情不是這樣看的.
我們今天看事情.
我們習慣用自家的實力去看眼前的狀況.
自家實力並不代表是錯的.
我們也不用刻意令我們捉襟見肘.
什麼都做不到.
讓上帝帶領我們.
這些就是反過來的經文.
經文的重點就是.
我們無論作事行事為人.
或者我們要去擴展.
眼前可能有些條件不是太有利於我們.
但當我們知道是上帝叫我們去做的.
上帝必然會與我們同在.
今天我們很多時候.
我們姐妹有一種心態是我們沒有實力.
講得直白一點.
例如開荒工作.
傳福音去到不同的地方.
或者我們要做一些很困難的使命.
其實每一個使命都是困難的.

$^{361}$但我們不敢去的原因是.
我們實力不夠.
我們不夠錢.
奉獻不夠.
我們老人家多一點.
或者年輕人年紀小的人很多.
等等有很多人世間我們所想的問題.
不過在上帝的眼中.
這些通通都不是問題.
問題是上帝有沒有與我們同在.
我們信不信得過上帝.
在經文第七章裡.
基顛就經歷了這個信心的功課.
其實基顛去到要打這十三萬大軍.
經文第七章第四節.
對基顛說.
人還是過多.
你要帶他們下到水旁.
我在那裡為你們試試他們.
我指點誰說.
這人可以跟你同去.
他就可以跟你同去.
我指點誰說.
這人不可跟你同去.
你就不可跟你同去.
接著第七節.
耶和華對基顛說.
我要用這添水的三百人.
拯救你們.
將米顛人交在你的手中.
其餘的人可以各歸各處去.
這三百人就帶著食物和各.
其餘的以色列人.
基顛都打發他們各歸各的帳篷.
只留下這三百人.
米顛仍在他們下邊的平原裡.
接著第九節.
當那耶和華吩咐基顛說.
起來.
下到米顛營去.

$^{401}$因為我已經將他們交在你手裡.
倘若你怕下去.
就帶著你的僕人普拉.
下到那營去.
你必聽見他們所說的.
然後你就有膽量下去攻營.
在第九節裡面首先說.
三百人面對十幾萬.
好像蝗蟲那麼大的軍隊.
順帝知道基顛是害怕的.
人之常情.
你在山那裡偷偷地看下去.
下面軍隊是片滿了.
而你看看後面.
三百人在那裡喝水.
整件事是實力懸殊.
是害怕的.
上帝也知道基顛是害怕的.
當他害怕的時候.
上帝說你帶著你的僕人去.
當然不單止是普拉.
普拉為首.
好幾個僕人.
但多極都有限度.
你必聽見他們所說.
然後你就有膽量下去攻營.
聽見他們所說的.
你就有膽量去攻營.
其實說什麼呢.
於是基顛帶著他的僕人普拉下到營旁.
米顛人,阿瑪力人和一切東方的.
都佈在平原上.
如同蝗蟲那樣多.
他們的駱駝無數.
多如海邊的沙.
基顛到了.
聽見一人將夢告訴同門說.
我作了一個夢.
夢見一個大麥餅滾入米顛營中.
到了帳篷.

$^{441}$將帳篷撞倒.
帳篷就翻轉傾覆了.
那同伴說.
這不是別的.
乃是以色列人約阿斯的兒子基顛的刀.
臣已將米顛和全軍都交在他的手中.
在這裡.
基顛帶幾個同伴.
窺探敵方的軍情.
當他知道敵方有很大量.
其中一個同伴就會說.
我作了一個夢.
我見到一個大麥餅衝入米顛營中.
到帳篷.
帳篷就全推倒.
然後那隊軍隊就傾覆了.
我作了一個這樣的夢.
我們經常都做夢.
然後解夢的人就說.
這不是別的.
那不是大麥餅.
那是以色列人約阿斯的兒子基顛的刀.
臣已將米顛和全軍都交在他的手中.
當你聽到這個時候.
如果你真的打仗.
為何打?為何進攻?.
因為其中一個探聽的人作了一個夢.
做了一個夢之後解夢.
就知道這隊軍隊就交付在基顛的手中.
如果按照人理上程.
你會覺得很奇怪.
而且是一個很大的舉動.
是一個性命是否可以保存的進攻行為.
但是基顛聽到這個夢和解夢之後.
基顛有什麼反應呢?.
基顛聽到這個夢和夢的講解.
就敬拜神.
回到以色列營中說.
起來吧.
也就是說已將米顛的軍隊交在你們的手中了.

$^{481}$在這裡.
如果你按照那個夢.
或者按照那個人解夢.
其實那件事是不足夠有理據的.
整件事是覺得好像很任意.
不過整個舉動是誰叫他這樣做的呢?.
基顛很清楚.
是上帝叫他帶這些同伴過去.
上帝藉著這些同伴.
來將這個壯膽的訊息告訴基顛知道.
其實要打這段戰役.
基顛已經知道了.
甚至乎基顛也知道.
上帝是很實在地與他同在.
不過面對著將要進攻的時候.
實在有很多很擔心的.
有很多懼怕的地方.
而這兩三個同伴就成為了.
一個好像在過程中既了解敵人.
同時間也再一次肯定基顛.
他可以上前去去做的這些同伴.
所以基顛就下了決心.
就要相信耶和華已經將米顛的軍隊.
交付在他們的手中.
仗都是這樣打的.
不過最重要的是在這一刻.
上帝藉著這些同伴幫助他.
讓他有更加大的決心和信心.
來打這場的仗.
值得注意的是.
基顛聽到這個夢和講解的時候.
他就在那裡去敬拜神.
基顛的敬拜.
與之前他建立了一個.
耶和華沙龍的祭壇有一點不同.
上一次的耶和華沙龍耶路巴力.
他想知道那個是不是上帝對他說話.
然後他就拿一些禮物來獻祭的那個祭壇.
他某程度上試探聽.
究竟跟他說的是不是上帝.

$^{521}$不過這次的敬拜.
就跟上一次的祭壇有一點不同.
因為這一次的敬拜.
他知道是上帝藉著這些人.
肯定他要攻打這些米顛人.
這個敬拜也代表著.
他對上帝一份的決心.
他去敬拜上帝.
也將他這一次的戰役.
整個打仗他就要交付給上帝.
求上帝幫助他.
在經文裡面你值得留意.
如果自發地去做敬拜.
在《釋師記》裡面.
其實是很少的一些記載.
不過可以肯定的一點就是.
基澱他很邀請上帝與他同行.
即是說當時已經建立了一個敬拜的制度.
或者獻祭都是由利未人.
或者由祭司來做.
他本身就不是那個支派.
馬來西亞的支派.
他不是那個屬於獻祭的支派.
但是在打仗行軍之前.
他只是邀請上帝同行.
上帝某程度默許他做這件事.
他沒有說他做錯了一件事.
所以在這裡面.
他的敬拜成為了機電壯膽的一個歷程.
接下來其實經文很快.
我們要看了.
因為沒有時間去看.
經文接著就說.
這三百人進軍去打這些米顛人.
打的方式是怎樣呢?.
其實很原始的.
他們用一些瓶子.
然後裡面有一些火把.
類似今天所說的那些炸彈,汽油彈等等.
他們放了一些火進去.

$^{561}$然後就扔進去敵方的軍營裡面.
然後那些人一起火的時候就內訌.
經文裡面說到耶和華派那些敵軍自己互相廝殺.
整件事不是那三百大軍衝進去.
接著就進去十幾萬大軍裡面進行殺戮.
不是這樣.
而是用一個很原始.
某程度是以小勝大的方式.
就是用一些瓶子帶著一些火把.
接著就令到那裡失火.
令到敵軍很慌忙.
接著就說到這十幾萬大軍.
一直在敗陣.
直到那兩個領袖被捉拿.
接著就將他們殺了.
而在打仗的時候.
跟隨他們的人.
見到時機到了.
在七章十八節.
我和一切跟隨我的人摧國的時候.
你們要在營的四處摧國.
喊叫說.
耶和華和機電的刀.
機電和跟隨他的一百人在三更之初.
才換更的時候來到營國就摧國.
打破手中的瓶.
三對人就摧國.
打破瓶子.
左手拿著火把.
右手拿著割.
就喊叫說.
耶和華機電的刀.
他們在營的四處各站各地方.
全營都亂鼠.
三百人立陷.
使他們逃跑.
三百人就摧國.
耶和華使全軍的人用刀互相擊殺.
逃到西尼拉的白哈士塔.
直到靠近塔伯亞伯米荷拉.

$^{601}$經文很簡單很明快.
某程度上是說.
用一個以小勝大的方法.
每人帶著一百人.
就打敗了圍剿著這十幾萬.
軍心賄亂的米電聯軍.
還有其他外邦.
這十幾萬人就被擊殺了.
擊殺之後.
打勝仗了.
我們那裡跳了.
因為沒辦法仔細讀戰爭是怎樣.
不過我想我們再看多一點.
就是打完仗之後.
八章二十二節到二十七節.
當打勝仗了.
敵方米電的首領已經被處決了.
這個外族某程度的威脅.
昔日欺壓的事已經改變了.
整個局面已經平定了的時候.
八章二十二節.
以色列人對著基甸說.
你既救我們脫離米電人的手.
願你和你的兒孫管理我們.
基甸說.
我不管理你們.
我的兒子也不管理你們.
唯有耶和華管理你們.
基甸又對他們說.
我有一件事求你們.
請你們各人將所奪的耳環給我.
原來仇敵是以實瑪利人.
都是戴金耳環的.
他們說.
我們情願給你.
就鋪開一件外衣.
將各人所奪的耳環雕在其上.
基甸所要出來的金耳環.
重一千七百射克勒金紙.
此外還有米電王所戴的月環.

$^{641}$耳兌和所穿的紫色衣服.
並落在頂上的金鏈紙.
基甸以此製造了一個爾弗德.
設在本城俄弗拉.
後來以色列人拜那爾弗德.
行了邪人.
這就作了基甸和他全家的網絡.
在經文裡面.
說到基甸打勝了仗.
三百人軍隊打贏了這場仗之後.
本來我們可以很容易就有一個結論.
上帝說免得你們自誇.
免得那些以色列人自誇.
靠自己的實力.
來打贏了米電.
一直欺壓他們的外邦.
不過一銀有兩面.
正正因為人少.
基甸帶著這三百人.
打贏十幾萬大軍.
結果是基甸很厲害.
明不明白我的意思.
原本是實力懸殊之戰.
人力無法打贏.
但結果是覺得.
這班以色列人覺得.
原來基甸這麼厲害.
三百人都打勝了.
贏了十幾萬大軍.
這簡直是驚世人才.
他們想的.
所以他們覺得.
應該由這麼強勁的軍事領袖.
來管理他們整個國家民族.
以後他們就有好日子過了.
他們看基甸成為了他們一個很重要的領袖.
經文裡面說.
叫他和他們的兒孫來管理他們.
在這裡面很隱含的一句說話.
事先不是代傳的.

$^{681}$如果你看事先的上下文.
每一個事先是獨立的.
是上帝的挑選.
上帝的揀選.
令這個人可以成為整個以色列民族的領袖.
但去到基甸的時候.
那些以色列民族說.
由你和你的子孫來管理.
隱含了就是一個帝制.
他要將基甸封為.
他要請基甸成為他們的王.
用世襲的方式.
來管理他們這一班以色列整個民族.
尤其是這個人有驚人的實力.
三百人就可以打贏十幾萬大軍.
所以他懇求基甸和他的子孫來管理他.
不過基甸在這裡面很清晰.
在這一刻裡面的答案.
這一刻很勤調.
我不管你們.
我的兒子也不管你們.
因為管理他們的只是耶和華上帝.
唯有耶和華管理你們.
其實這是一個很正確的答案.
很正確的.
基甸的心也不在於要做一個王.
不過基甸的心想做什麼呢.
他不是想做王.
他叫其他人將他的耳環.
金耳環,金器.
找件衣服.
那些人就自發鋪在地上.
鋪在地上.
將他的金耳環都脫下來.
基甸就說用了這些金.
來做一個爾忽德.
基甸的心不是想做王.
不過做爾忽德他想做什麼.
大祭司.
他想做傳道人.

$^{721}$簡單來說.
他想做主任牧師.
你這樣理解吧.
他不是想做執事會主席.
他是想做主任牧師.
用現代來理解.
於是當時的人就將所有的金器.
貴重的金屬都脫下來.
鑄造了一個爾忽德.
爾忽德是什麼呢.
其實是當時大祭司裡面的大祭司的袍.
大祭司的袍就是用細麻布來製造的.
類似一個圍腰的長條狀.
代表著一種尊貴的祭祀身份.
基甸其實由最初.
他根本就不是喜歡打仗的人.
他是想做一個祭司.
他在這裡打勝仗.
外族就停定了.
他就想去想這件事.
不過他想做大祭司.
如果按照他自己的宗族是沒有辦法做的.
如果按照舊時的律法規定.
是利未人才可以做的.
而他的支派是馬來西亞支派.
原則上這件事是距離他很遠的.
不過就是因為他打勝了仗.
他得到這樣的榮譽.
其他人都很樂意將這些金器脫下來.
幫他們完成了基甸的想法.
基甸就用錢來做了一個以弗德.
將這個以弗德設立在本城的俄弗拉.
後來以色列人拜那以弗德行了邪人.
這就作了基甸和他全家的網絡.
經文裡面對於基甸的想法.
帶著一個比較強烈的批判.
基甸上帝呼召他去打仗.
呼召他成為一個行軍遣將.
整個民族的一個領袖.
他沒有膽子,上帝給他膽子.

$^{761}$他的敬拜是幫他壯膽.
第一次的敬拜是驗證上帝.
第二次的敬拜是幫助他.
設立為他有信心去打仗.
但這個敬拜不是要肯定他成為一個祭司.
但基甸打勝仗後.
他將他的焦點轉移.
他想打造一個以弗德.
在他自己的根據地俄弗拉這個地方.
建立一個以弗德.
藉著這個以弗德.
希望得到一種祭司的尊榮.
不過這件事.
他這樣帶來的後果是什麼呢?.
這個後來用金打造的以弗德.
就成為了一件聖物.
一個神聖之物.
這件神聖之物背後有故事的.
就是上帝怎樣打勝仗.
然後打贏了米顛人.
然後人們願意奉獻成為了一個這樣的聖物.
但這件事就成為了一個新的偶像.
昔日基甸有個別名.
你記不記得是什麼?.
耶路巴力.
他打爛了愛邦人巴力的偶像.
不過因為這件事.
他鑄造了一個新的偶像.
一個金造的以弗德.
也成為了他和以色列人後來的網絡.
簡單來說.
他因著他自己的想法.
打造了一個新一代的偶像.
昔日他以打破巴力這個偶像.
作為他最大的標榜的一個業績.
人生最光彩的一頁.
但這個名字打勝之後.
他晚年就成為了他生命裡面一個很重要的網絡.
也成為了他第一個的污點.
就是這個金造的以弗德.

$^{801}$在這裡我們也有去想.
在我們的人生裡面.
我們有一些很光彩,很美麗的經歷.
我們有一些熟悉某些經驗.
可能這些經驗令我們每個人拍爛手掌.
但當我們人生裡面有些焦點迷失了的時候.
可能我們會將這些經驗.
構定成為教會裡面.
另一個必須要有的傳奇.
有一些神聖不可逾越的規範.
不過這可能成為了我們敬拜另一個偶像的開始.
頂子們,我們今天信仰裡面有沒有這些雜質呢?.
有些東西,我們過去的經驗是很好的.
上帝某程度給予我們.
直到極致,我將這東西捧到最高.
當作不可失去的一件事.
成為了最高價值的傳統.
可能我們有機會步向基甸所走的失敗的那條路.
教會傳統裡面有大量這些事.
之前我們一班頂子妹去過羅馬.
去到天主教的禮拜堂.
我們看到那些地方很美麗,很滑.
不過有些東西,我們作為基督徒都很冷淡.
我們都覺得不太認同.
譬如有些人說那口釘是比彼得釘的.
後來成為了一個聖物.
找了一個很美麗的金壇放在那裡.
然後看著那口用木做的釘.
那口釘是提醒你為義而受逼迫.
有些人將它成為了神聖的源頭來祝福你.
當時很多人覺得這些可以治病的聖物.
就去摸摸,或者去做不同的敬拜.
可能將我們的焦點擺錯了.
我們今天基督教沒有這些東西.
不過可能有些無形的偶像潛藏在我們的內心當中.
基甸的經歷提醒我們.
我們要撇除,不要讓新的伊弗特成為我們生命裡面新一代的網絡.
基甸不單止做了一個金伊弗特.
它還有第二個很不光彩的地方.
29至31節.

$^{841}$約阿斯的兒子耶路巴力回去.
再提一次,耶路巴力.
住在自己的家裡.
基甸有70個親生的兒子.
因為他有許多的妻.
他的妾住在事件.
也給他生了一個兒子.
基甸與他起名叫阿比米勒.
這裡只是輕輕的一句說話.
不過基甸一個這麼成功,這麼彪炳的人.
他製造了第二個網絡為他們自己的民族.
首先,經文裡面有一點諷刺.
這個叫耶路巴力的人.
他有多少個兒子?.
70個兒子.
也不知道怎樣生的.
經文裡面補充給你聽.
一個妻子生不了70個.
豬也生不了這麼多.
因為他有很多的妻.
大量的妻.
你見到和昔日基督教和猶太人所屏行的傳統已經有些偏離.
然後生了大量的兒子之後.
他還有一個妾.
這個就是可圈可點.
正室也夠了.
他還有一個妾.
就在事件這個地方.
而這個妾生了一個兒子.
名字叫阿比米勒.
阿比米勒是什麼意思?.
阿比米勒要回頭.
剛才那些以色列人求基甸.
你既然救我們脫離米殿人的手.
求你和你的兒孫管理我們.
做王嘛.
基甸說我不會管理你.
唯有耶和華管理你.
Modal answer.
不過到妾所生的兒子.

$^{881}$基甸給他的名字叫阿比米勒的時候.
原來阿比米勒的意思.
阿比就是說爸爸.
米勒就是君王.
就是說這個兒子每一次叫阿比米勒.
就是叫什麼?.
叫父王.
基甸是不是真的.
心裡面純粹想著.
由耶和華管理他.
在那一刻他說了一個modal answer.
但是到了這一刻.
似乎他也想.
他成為一個新一代的王.
又或者讓他的兒子.
繼續接在管理的位置.
繼續管理這群以色列人.
不過在經文裡面.
第八章的結尾.
他只是用阿比米勒.
成為了他的結束.
不過阿比米勒這個人.
和耶路巴力這個綽號.
在接下來的第九章會再出現.
新一個新的網絡.
又再在以色列人當中.
究竟這個網絡是什麼呢?.
下一次我們再說.
不過今天這段經文提醒我們.
很多人說失敗乃成功之母.
我們覺得成功是我們人生的終點.
不過這段經文告訴我們.
可能成功是失敗之母.
成功不是我們生命的終點.
可能我們今天的成功.
隱藏了很多未來的失敗.
又或者未來我們後一代的失敗.
今日弟兄姊妹我們要好好思想.
從機電的身上.
從這個很彪炳.

$^{921}$但充滿污點的家庭裡面.
提醒我們怎樣行事為人.
求主繼續帶領我們.
對我們說說話.
\newpage



\section{士師記 11:1-40}
\label{sec:Fr77s7J2G5I}
\textbf{2023年8月13日崇拜直播｜陳冠成牧師｜「家醜說出口」系列:傷痛不代傳｜士師記11\_1-40}
\newline
\newline
連結: \href{https://youtube.com/watch?v=Fr77s7J2G5I}{\texttt{https://youtube.com/watch?v=Fr77s7J2G5I}} ~~~~ 語音日期: 2023-08-14
\newline
\newline
\hyperref[sec:iklPHMTALSA]{\small{< < < PREV SERMON < < <}}
~
\hyperref[sec:index_chronic]{\small{[返順時目]}}
~
\hyperref[sec:bCwRlVK9wDE]{\small{> > > NEXT SERMON > > >}}
\newline
\newline
士師記 11:1-40
\newline
\begin{longtable}{cl}
\hline
\hline
章節 & 經文 (和合本修訂版)\\
\hline
11:1 & \begin{tabularx}{0.7\textwidth}{X} 基列人耶弗他是個大能的勇士，是妓女的兒子。基列生了耶弗他。 \end{tabularx} \\ \\ \relax
11:2 & \begin{tabularx}{0.7\textwidth}{X} 基列的妻子也給他生了幾個兒子。他妻子生的兒子長大後，就把耶弗他趕出去，說：「你不可在我們父家繼承產業，因為你是別的女人生的兒子。」 \end{tabularx} \\ \\ \relax
11:3 & \begin{tabularx}{0.7\textwidth}{X} 耶弗他就逃離他的兄弟，住在陀伯地。有些無賴的人聚集在他那裡，與他一同出入。 \end{tabularx} \\ \\ \relax
11:4 & \begin{tabularx}{0.7\textwidth}{X} 過了些日子，亞捫人攻打以色列。 \end{tabularx} \\ \\ \relax
11:5 & \begin{tabularx}{0.7\textwidth}{X} 亞捫人攻打以色列的時候，基列的長老去請耶弗他從陀伯地回來。 \end{tabularx} \\ \\ \relax
11:6 & \begin{tabularx}{0.7\textwidth}{X} 他們對耶弗他說：「請你來作我們的指揮官，好讓我們跟亞捫人打仗。」 \end{tabularx} \\ \\ \relax
11:7 & \begin{tabularx}{0.7\textwidth}{X} 耶弗他對基列的長老說：「你們不是恨我，把我趕出父家嗎？現在你們遭遇急難，為何到我這裡來呢？」 \end{tabularx} \\ \\ \relax
11:8 & \begin{tabularx}{0.7\textwidth}{X} 基列的長老對耶弗他說：「現在我們回到你這裡，是要請你同我們去跟亞捫人打仗，作基列所有居民的領袖。」 \end{tabularx} \\ \\ \relax
11:9 & \begin{tabularx}{0.7\textwidth}{X} 耶弗他對基列的長老說：「若你們請我回去跟亞捫人打仗，耶和華把他們交給我，我就作你們的領袖。」 \end{tabularx} \\ \\ \relax
11:10 & \begin{tabularx}{0.7\textwidth}{X} 基列的長老對耶弗他說：「有耶和華在你我之間作證，我們必定照你的話做。」 \end{tabularx} \\ \\ \relax
11:11 & \begin{tabularx}{0.7\textwidth}{X} 於是耶弗他與基列的長老同去，百姓就立耶弗他作他們的領袖和指揮官。耶弗他在米斯巴將他一切的事陳述在耶和華面前。 \end{tabularx} \\ \\ \relax
11:12 & \begin{tabularx}{0.7\textwidth}{X} 耶弗他派使者到亞捫人的王那裡，說：「你與我有甚麼相干，竟來到我這裡攻打我的地呢？」 \end{tabularx} \\ \\ \relax
11:13 & \begin{tabularx}{0.7\textwidth}{X} 亞捫人的王對耶弗他的使者說：「因為以色列從埃及上來的時候佔據我的地，從亞嫩河到雅博河，直到約旦河。現在你要和平歸還這些地方！」 \end{tabularx} \\ \\ \relax
11:14 & \begin{tabularx}{0.7\textwidth}{X} 耶弗他又派使者到亞捫人的王那裡， \end{tabularx} \\ \\ \relax
11:15 & \begin{tabularx}{0.7\textwidth}{X} 對他說：「耶弗他如此說：以色列並沒有佔據摩押地和亞捫人的地。 \end{tabularx} \\ \\ \relax
11:16 & \begin{tabularx}{0.7\textwidth}{X} 以色列人從埃及上來，是經過曠野到紅海，來到加低斯。 \end{tabularx} \\ \\ \relax
11:17 & \begin{tabularx}{0.7\textwidth}{X} 那時，以色列派使者去以東王那裡，說：『求你讓我穿越你的地。』以東王卻不聽。以色列又照樣派使者去摩押王那裡，他也不肯。於是以色列人就住在加低斯。 \end{tabularx} \\ \\ \relax
11:18 & \begin{tabularx}{0.7\textwidth}{X} 他們又經過曠野，繞過以東地和摩押地，到摩押地的東邊，在亞嫩河邊安營，並沒有進入摩押的境內，因為亞嫩河是摩押的邊界。 \end{tabularx} \\ \\ \relax
11:19 & \begin{tabularx}{0.7\textwidth}{X} 以色列派使者去亞摩利王，就是希實本王西宏那裡；以色列對他說：『求你讓我們穿越你的地，到我自己的地方去。』 \end{tabularx} \\ \\ \relax
11:20 & \begin{tabularx}{0.7\textwidth}{X} 西宏卻不信任以色列，不讓他們穿越他的疆界。他召集了他的眾百姓在雅雜安營，與以色列爭戰。 \end{tabularx} \\ \\ \relax
11:21 & \begin{tabularx}{0.7\textwidth}{X} 耶和華－以色列的神將西宏和他的眾百姓都交在以色列手中，以色列人就擊殺他們，佔領了那地居民亞摩利人的全地。 \end{tabularx} \\ \\ \relax
11:22 & \begin{tabularx}{0.7\textwidth}{X} 他們佔領了亞摩利人所有的疆土，從亞嫩河到雅博河，從曠野直到約旦河。 \end{tabularx} \\ \\ \relax
11:23 & \begin{tabularx}{0.7\textwidth}{X} 耶和華－以色列的神如今從他百姓以色列面前趕出亞摩利人，你竟要佔領它嗎？ \end{tabularx} \\ \\ \relax
11:24 & \begin{tabularx}{0.7\textwidth}{X} 你不是已經得了你的神明基抹賜給你的地為業嗎？耶和華－我們的神在我們面前所趕出的，我們也要得它為業。 \end{tabularx} \\ \\ \relax
11:25 & \begin{tabularx}{0.7\textwidth}{X} 現在你比西撥的兒子摩押王巴勒還強嗎？他真的曾與以色列爭執，或是真的與他們爭戰了嗎？ \end{tabularx} \\ \\ \relax
11:26 & \begin{tabularx}{0.7\textwidth}{X} 以色列人住希實本和所屬的鄉鎮，亞羅珥和所屬的鄉鎮，以及沿著亞嫩河的一切城鎮，已經有三百年了。在這期間，你們為甚麼不取回呢？ \end{tabularx} \\ \\ \relax
11:27 & \begin{tabularx}{0.7\textwidth}{X} 我並沒有得罪你，你卻要攻打我，加害於我。願審判人的耶和華今日在以色列人和亞捫人之間判斷是非。」 \end{tabularx} \\ \\ \relax
11:28 & \begin{tabularx}{0.7\textwidth}{X} 但亞捫人的王不聽耶弗他傳達給他的話。 \end{tabularx} \\ \\ \relax
11:29 & \begin{tabularx}{0.7\textwidth}{X} 耶和華的靈降在耶弗他身上，他就經過基列和瑪拿西，經過基列的米斯巴，又從基列的米斯巴過到亞捫人那裡。 \end{tabularx} \\ \\ \relax
11:30 & \begin{tabularx}{0.7\textwidth}{X} 耶弗他向耶和華許願，說：「你若真的將亞捫人交在我手中， \end{tabularx} \\ \\ \relax
11:31 & \begin{tabularx}{0.7\textwidth}{X} 我從亞捫人那裡平平安安回來的時候，無論誰先從我家門出來迎接我，就要歸給耶和華，我必將他獻上作為燔祭。」 \end{tabularx} \\ \\ \relax
11:32 & \begin{tabularx}{0.7\textwidth}{X} 於是耶弗他往亞捫人那裡去，與他們爭戰。耶和華將他們交在他手中， \end{tabularx} \\ \\ \relax
11:33 & \begin{tabularx}{0.7\textwidth}{X} 他就徹底擊敗他們，從亞羅珥到米匿，直到亞備勒‧基拉明，攻取了二十座城。這樣，亞捫人就在以色列人面前被制伏了。 \end{tabularx} \\ \\ \relax
11:34 & \begin{tabularx}{0.7\textwidth}{X} 耶弗他回米斯巴去，到了自己的家，看哪，他女兒拿著手鼓跳舞出來迎接他。她是耶弗他的獨生女，除她以外，沒有別的兒女。 \end{tabularx} \\ \\ \relax
11:35 & \begin{tabularx}{0.7\textwidth}{X} 耶弗他一看見她，就撕裂衣服，說：「哀哉！我的女兒啊，你使我非常悲痛，叫我十分為難了。因為我已經向耶和華開了口，不能收回。」 \end{tabularx} \\ \\ \relax
11:36 & \begin{tabularx}{0.7\textwidth}{X} 他女兒對他說：「我的父親啊，你既向耶和華開了口，就當照你口中所說的向我行，因為耶和華已經在你的仇敵亞捫人身上為你報了仇。」 \end{tabularx} \\ \\ \relax
11:37 & \begin{tabularx}{0.7\textwidth}{X} 她又對父親說：「我只求你這一件事，給我兩個月，讓我和同伴下到山裡，好為我的童貞哀哭。」 \end{tabularx} \\ \\ \relax
11:38 & \begin{tabularx}{0.7\textwidth}{X} 耶弗他說：「你去吧！」他就讓她離開兩個月。她和同伴去了，在山裡為她的童貞哀哭。 \end{tabularx} \\ \\ \relax
11:39 & \begin{tabularx}{0.7\textwidth}{X} 過了兩個月，她回到父親那裡，父親就照所許的願向她行了。她從來沒有親近男人。於是以色列中有個風俗， \end{tabularx} \\ \\ \relax
11:40 & \begin{tabularx}{0.7\textwidth}{X} 每年按著日期以色列的女子要去為基列人耶弗他的女兒哀哭四天。 \end{tabularx} \\ \\
[1ex]
\hline
\hline
\end{longtable}
$^{1}$今日選讀的經文是.
《士師記》第十一章一至十一節.
聖經舊約第三八一十四頁.
這裡記載了耶乏他士師的一些事蹟.
請聽聖經的記載.
「激烈人耶乏他是個大能的勇士.
是妓女的兒子.
耶乏他是激烈所生的.
激烈的妻也生了幾個兒子.
他妻所生的兒子長大了.
就趕逐耶乏他.
說:你不可在我們父家承受產業.
因為你是妓女的兒子.
耶乏他就逃避他的弟兄.
去住在陀伯地.
有些匪徒到他那裡聚集.
與他一同出入.
過了些日子.
阿們人攻打以色列.
阿們人攻打以色列的時候.
激烈的長老到陀伯地去.
要叫耶乏他回來.
對耶乏他說:.
請你來作我們的元帥.
我們好與阿們人爭戰.
耶乏他回答激烈的長老說:.
從前你們不是恨我.
趕逐我出來父家嗎?.
現在你們遭遇急難.
為何到我這裡來呢?.
激烈的長老回答耶乏他說:.
現在我們到你這裡來.
是要你和我們去與阿們人爭戰.
你可以作激烈一切居民的領袖.
耶乏他對激烈的長老說:.
你們叫我回去與阿們人爭戰.
耶和華把他交給我.
我可以作你們的領袖嗎?.
激烈的長老回答耶乏他說:.
有耶和華在你我中間作見證.

$^{41}$我們必定照你的話行.
於是耶乏他和激烈的長老回去.
百姓就納耶乏他作領袖.
作元帥.
耶乏他在米斯巴將自己的一切話.
陳明在耶和華面前.
各位姐妹平安.
今天這個講題其實來自一本書.
講關於家庭治療.
《傷痛不代存》.
已經出版了很多年.
這本書頗為經典.
裡面提及到.
我想每個人都認同.
家庭本來是一個很溫暖.
充滿愛.
來培育下一代孩子.
健康成長的地方.
我們都很認同這一點.
不過對很多人來說.
在他小時候.
或者在他青少年成長階段的時候.
裡面很多負面的家庭經驗.
其實都令他們飽受傷痛.
我們都聽聞過.
或者甚至經歷過.
不能甩掉.
亦都忘記不了.
就算他日後長大成人.
甚至結婚生子.
有自己的家庭.
這些上一代遺留下來的傷害.
很奇怪都會繼續用不同的方式.
來延伸到他們下一代的身上.
就好像我們今天看到.
耶佛他的故事一樣.
其實和其他的事司一樣.
耶佛他的興起.
始於以色列人.
又行上帝眼中看為惡的事.

$^{81}$他們敬拜事奉列邦的神.
於是自必然惹來上帝的法路.
將他們交在.
菲律斯人和阿滿人的手上.
當他們被仇敵來到碧巴.
感到極其困乏的時候.
自然來到上帝面前哀求他.
承認自己的過犯.
不過很奇怪你會看到.
上帝不再像我們之前看到的.
事司的故事一樣.
主動為百姓興起一位拯救者.
上帝這次沒有這樣做.
上帝沒有幫他們.
揀選一位領袖.
來帶領他們脫離敵人的核制.
上帝說.
印在這裡的經文.
在你的秩序表講到大綱.
我們今天還未看的.
就是第十章前一章.
上帝在那裡表示.
其實我已經拯救了你很多次.
從不同的外邦人手上幫過你們.
不過每一次試過警籤之後.
你們又固態服民.
又再一次捨棄我.
去敬拜外邦的偶像.
上帝看到這班百姓.
只是想上帝提供服務.
不想他們專心跟隨上帝.
於是上帝這次把心一橫.
狠心一點.
他決定不出手.
任憑他們承擔罪所帶來的惡果.
並且上帝在這裡諷刺他們.
你們去拜別神.
好呀.
這次你們就試試.
你們所拜的別神.

$^{121}$求祂幫你們.
這是耶律發興起的背景.
而在這個時候.
阿摩人的軍隊已經結集在.
基列這裡準備攻打以色列人.
在沒有上帝為百姓興起事先的情況下.
百姓只有自己選擇.
自己用自己的方法.
選一位軍事領袖來對抗阿摩人.
於是他們派人到陀伯地.
招募一位既懂得打仗.
又懂得帶兵的大能勇士.
希望他出手解決燃眉之急.
這就是耶律發興起.
在十一章一開始.
他會先交代耶律發興起的身世.
經文這裡說.
剛才我們也一開始讀了.
他是父親和一位妓女所生的兒子.
他的父親也跟他的妻子.
即是正室.
生了幾個兒子.
當這班同父異母的兄弟長大後.
很大機會他們的父親已經離開了人世.
聯手趕逐耶佛他.
不准他繼續留在父家承受產業.
理由是他出身不好.
他是妓女的兒子.
不過按照摩西的律法.
耶佛他其實是可以承繼到.
他父家的產業.
因為承繼權只會受父親的影響.
跟誰是他的母親沒有關係.
所以就算他是妓女所生的.
耶佛他仍然符合承繼家業的條件.
所以很大機會當時他和父異母的兄弟.
需要打官司.
透過訴訟才能夠.
強行剝奪耶佛他的承繼權.
並且將他趕走.

$^{161}$至於當時負責裁決的.
很大機會是現在.
找他出手.
找他打仗的長老.
至於耶佛他被家人遺棄.
強行將他應有的權益拿走.
沒辦法.
自然要逃跑.
逃到塘北地這個地方.
聽聞說他和一些匪徒.
我們今天好像地痞流氓.
或者說社會邊緣人.
經常和他們出入為伍.
說得俗一點.
就是出來逛.
亦因為這樣.
你會看到其實.
耶佛他他整個人生.
他有些東西很厲害.
他很懂得和人討價還價.
他很懂得為自己爭取應有的權益.
所以你留意.
剛才我們所聽到的那一段.
在十一章.
他和長老的對話.
像不像談判和談判?.
其實都很像.
我們再看一遍.
你留意激烈的長老首先出價.
他們開出的條件是什麼?.
耶佛他他來做我們的元帥.
幫我們和亞滿人爭戰.
元帥的職位吸不吸引?.
不錯啊.
聽下去.
挺好.
不過.
耶佛他他很聰明.
他不是省油的燈.
他知道元帥的權力.

$^{201}$只限於什麼場面?.
打仗而已.
打完仗怎麼辦?.
打完仗.
其實作為一個元帥.
他其實沒有什麼實權.
沒有什麼地位.
說得難聽一點.
其實.
就是一個可以用完即棄的爛頭卒.
找你去打仗.
給你權力.
但打完仗之後.
沒有了.
完了.
所以你會很明白.
耶佛他他為什麼不會接受.
試問一個曾經被人遺棄的人.
他想不想.
這個被遺棄的經歷.
再次發生在自己身上?.
沒有人想這樣.
沒有人喜歡這樣.
再被人遺棄.
他也知道.
在這一刻.
他特別在長老的面前.
他是最有市場價值的.
他是最有爭議力.
所以他首先回應的是什麼?.
和這班長老翻舊帳.
他說.
你們以前不是很憎恨我嗎?.
不是憎恨到要將我趕出父家嗎?.
現在你們有難.
為什麼來找我?.
言下之意是什麼意思?.
想找我幫忙.
幫你們打仗?.
可以.

$^{241}$先搞定我們之間的舊帳.
所以你看到.
雙方都有各取需要.
耶佛他他很知道.
他可以把握這個機會.
讓有能力說話的長老.
幫他平反洗底.
除去昔日被人趕走.
被人隔斜.
褫奪了承繼權的污名.
而這班長老也看中他什麼?.
看中耶佛他有很強的打仗作戰能力.
可以幫他們趕走阿滿人.
如果談得攏就怎樣?.
各取所需.
皆大歡喜.
於是乎你看到.
這班長老也很聰明.
馬上加碼.
開一個更加優效的條件.
耶佛他只要你肯回來.
跟我們阿滿人爭戰.
你可以做激烈一切居民的首領.
不同了.
做激烈一切居民的首領.
代表什麼?.
領袖.
打完仗他還是領袖.
還可以維持到.
他在眾人面前的權力地位和利益.
對耶佛他來說.
這一定是一個可遇不可求的機會.
可以說是翻身了.
由一個不務正業的流氓搖身一變.
變成什麼?.
變聖民族的英雄和領袖.
其實怎會不心動?.
不過你會看到.
耶佛他有沒有馬上答應?.
他沒有.

$^{281}$其實你看到他信不過這班.
曾經放棄了他.
踢走他.
傷害過他的人.
他很怕他們.
怎知道你之後會不會反口.
怎知道到時打完仗.
你會不會不認帳.
又踢我走.
於是乎要求他們.
你再三保證.
我想請大家一起讀.
第九至十一節.
他怎樣要求這班長老再三保證.
請大家讀.
九至十一節.
預備起.
.
二至四.
.
記念地長老回來.
八星就立耶佛他作領袖.
作元帥.
耶佛他在米斯巴將自己一切畫.
曾命在耶和華面前.
出現了多少次.
上帝的名字「耶和華」?.
多少次?.
三次.
這三節的經文出現了三次.
不過你留意.
上帝是一直緘默的.
上帝說明了不幫忙.
上帝沒有參與.
不像以前那樣.
上帝很主動地.
向眾人顯明他的旨意.
甚至上帝親自選擇一位士師.
這次完全是人為操作.
所以你看到.

$^{321}$表面上大家好像很尊重上帝.
向上帝稟報一切.
並且請他在眾人面前做見證.
但你覺不覺得整件事.
與拜關公沒有分別?.
你覺不覺得?.
你試試把關義哥的名字.
換成「耶和華」.
再讀一次第九至十字.
例如.
你們叫我回去與阿滿人爭戰.
關義哥如果將他們交在我手上.
我是否可以做你們的領袖?.
激烈的長老回答.
有關義哥在我們中間.
做見證.
我們一定照你們所說的行.
是否像電視劇集?.
其實都很像.
某程度上你看到.
在耶和華和長老的心目中.
上帝只不過是陶章.
橡皮陶章.
大家談好雙方利益.
那又如何?.
敦請上帝出來.
麻煩你蓋個印.
麻煩你幫我們背書.
防止雙方日後反口不認數.
所以你有沒有發覺.
整件事讓我們看到.
耶和華在人生中.
特別當他遇到一些.
他覺得重要的時刻時.
他就會很厲害.
做一件事.
搬上帝的名字出來.
好像我們現代人.
有時都會這樣.
我們喜歡找一些有知名度的人.

$^{361}$有權威的人士.
來加持自己的想法或需要.
增加認受性和可信性.
耶和華和長老的做法.
其實是一個警惕.
我們和上帝的關係.
會不會只是口舌上的尊敬?.
回想我們有不少的人生經歷.
我們都奉主的名來進行.
我們都以上帝的名來做事.
甚至是事奉.
我們甚至求他的幫助.
不過這段經文.
值得我們提醒.
他們都是這樣.
但是我們有沒有真正.
在裡面尊上帝為大?.
有沒有真正一心.
尋求信服他的帶領?.
還是只是叫上帝出來.
承傳自己的想法?.
上帝會不會不知不覺.
在我們的橡皮圖章.
有需要緊急的時候.
上帝來幫幫我.
上帝求你承傳我.
求上帝扔個印.
希望他幫手保佑我們.
為我們提供服務.
這是我們有時會落入的試探.
你看到耶復他他很厲害.
他以上帝的名來滿足自己的需要.
無論他和這班長老講述.
稍後他和阿滿王講述.
其實都是這樣.
當然我們今天沒有機會去看.
之後第十二至二十八節的經文.
甚至乎我們會看的另一段經文.
提及到他準備打仗.
在他人生生死關頭的時刻.

$^{401}$他也會很習慣.
把上帝的名字搬出來.
來幫助自己.
不過雖然如此.
你看到上帝很奇妙.
他默默地仍然掌管一切.
他控制著大局.
並且他會使用.
耶復他這些軟弱.
這些瑕疵.
來完成上帝自己的旨意.
來拯救百姓.
脫離阿滿人的管轄.
當然.
當我們看到最後的時候.
你會發覺.
耶復他仍然要為他自己的行為.
為他自己所講過的話.
負上責任.
正所謂.
人種的是什麼.
最終收到的也是什麼.
我們繼續看下去.
跳到第二十九節.
二十九節三十三節.
這裡讓我們看到.
一直保持緘默的上帝.
現在主動採取行動.
將他的靈降臨在耶復他的身上.
雖然這位是人所設立的領袖.
但是上帝對他的子民仍然有恩典.
有憐憫.
沒有完全放棄他們.
上帝仍然願意.
與這班子民同在.
暗暗地幫他們脫離險境.
你會看到.
因為上帝的靈同在.
耶復他就像之前的一些事司.
例如基甸.

$^{441}$他很順利號召到.
無論是基列,馬來西亞.
或是米斯巴的地方的以色列人.
來到這裡起來.
準備與阿們人爭戰.
三十二至三十三節.
這裡很清楚強調.
勝負的關鍵是什麼.
是出於上帝的工作.
是上帝將阿們人交給耶復他手中.
而不是耶復他起誓.
所以你會看到.
就算抽走了耶復他.
在三十至三十一節.
特別深色了那段經文.
其實抽走了這段誓詞.
你會發覺上帝的旨意依然會成就.
依然會讓耶復他來戰勝敵人.
所以耶復他這次的起誓.
他表示如果上帝你幫我打勝仗.
平平安安回來的話.
誰第一個出來門口迎接我.
我就將他獻為梵祭.
獲給上帝你.
耶復他這個起誓其實是多餘的.
說多了說大了.
如果他沒有這樣做的話.
很大機會他的故事.
其實可以喜劇收場.
不過我們要集中去看.
為什麼他很著緊.
他要向上帝發這個誓.
為什麼他要很強調.
上帝如果我平平安安回來.
怎樣也好誰一個出來迎接我的.
我就將他獻為梵祭給你.
很明顯反映他著眼的是什麼.
他最緊張的是什麼.
不單單是要打贏阿們人.
更加要怎樣.

$^{481}$你試想像如果他只是打贏敵人.
但弄到自己重傷殘廢.
甚至戰死沙場.
那又怎樣.
那就是白做.
他不能享受之後所期望的終極成果.
耶復他需要的不單單是打贏巡將.
他要的是平平安安回來.
他才能成為所有百姓的領袖.
他才能站在昔日踩低他的人.
他才能消除他人生最大的污點和傷痛.
補償自己所失去的.
所以他要平安回來.
不單單是打贏巡將.
耶復他做這麼多事.
其實都是要達成他的心願.
而為保萬無一失.
他需要怎樣做.
他發覺獻生束不夠.
不夠誠意.
不足以打動上帝.
所以怎樣.
就像講數一樣.
加大注碼.
他仿效愛邦人.
獻人為繁祭的方式.
這樣就行了吧?上帝.
來到很希望來換取上帝的幫助.
所以等於子沫從耶復他這一個誓詞.
充分反映到上帝在他心目中是怎樣.
某程度上和其他的偶像假神.
其實真的分別不大.
都是可以收買到的.
都是可以和他做交易.
可以討價還價的對象.
頂智媒.
這裡給我們看到.
另一個提醒.
原來就算有上帝的靈同在.
也不確百分百保證.

$^{521}$那個人的心思意念.
必定對準上帝.
必定做出一些上帝喜悅的事情.
你會看到.
基典他晚年就是這樣.
耶復他現在又是這樣.
到最後.
遲些我們看到參宣這個事司.
甚至乎在撒母耳基的《素羅王》.
通通都有上帝的靈同在.
甚至乎上帝會使用他.
成就上帝的旨意.
不過不代表他們的生命.
一定討上帝喜悅.
耶復他.
他以人獻為凡祭.
他試圖操弄收買上帝.
其實很明顯.
我們都知道.
違反了摩西律法的吩咐.
耶復他.
他知道有上帝在這個世界上.
但很可惜的是.
他不認識上帝的話語.
他的情況也提醒我們.
其實你發不發覺.
人生很多時候遇到問題.
都是因為什麼?.
我們不認識.
沒有捉緊上帝的話語.
大家有坐巴士,地鐵的經驗.
習慣捉扶手嗎?.
有沒有發覺.
有些人是不會捉扶手的.
有很多原因.
可能他有手拿著東西.
或者他沒有東西拿.
但他也不捉.
很特別.
可能他覺得沒事的.

$^{561}$靠自己雙腳.
靠自己的平衡力.
都可以站穩.
其實很多時候都可以.
分別在哪裡?.
坐車的時候就看到.
當你遇到人生考驗.
有捉緊扶手和沒有捉緊扶手.
是完全兩回事.
沒有捉緊扶手會怎樣?.
不單自己跌倒或滑倒.
還會怎樣?.
最慘是你連累到身邊那個.
無辜被你撞一撞.
就不知如何是好.
某程度上耶乏他的故事.
都是這樣.
沒有捉緊.
沒有捉實上帝的話語.
覺得靠自己可以處理.
結果不單止連累自己.
連累他家人.
上帝允許我們有豐盛的生命.
但為什麼人生總是問題多多.
其實很值得我們透過這段經文.
來深思.
當然一定會找到很多原因理由.
來解釋.
甚至乎都可以合理.
不過歸根究底.
往往是我們怎樣看待上帝的話語.
我們有沒有認真去讀上帝的話語.
有沒有將上帝的吩咐.
真是實踐出來呢?.
如果我們真的不太清楚.
到底上帝對我們的要求是什麼.
到底上帝的吩咐是什麼.
只是有很多宗教的活動.
或者儀式進行.
其實我們很容易變成了耶乏他.

$^{601}$我們會將很多世俗的想法手段.
加進你的信仰裡.
好像還在信上帝.
還在參與這個信仰.
但實情.
我們可能已經離上帝很遠.
所以你細心看到.
上帝有沒有要求耶乏他起誓.
從來都沒有.
所以上帝也沒有回應他的誓詞.
這只是他一廂情願的想法.
所以上帝也怎樣做?.
就由你吧.
由你按司意來行事.
不過你要承擔後果.
所以.
沒錯.
耶乏他最後打贏了長征.
他不單打贏了阿們人.
還打贏了誰?.
打敗了自己.
打散了自己的頭家.
我們再看最後一段經文.
我想請大家一起讀.
34至39節.
預備起.
(念經).
謝謝.
耶乏他大概預計.
出來迎接他海船回歸.
很大機會.
應該是他的手下或是僕人.
所以他才會起這個誓.
但萬料不到.
第一個出來的是自己的獨生女.
其實挺慘的.
耶乏他一心想.
自己這次終於可以擺脫厄運.
殊不知厄運是怎樣?.
轉移到自己的女兒身上.

$^{641}$沒有擺脫.
他可以拯救百姓脫離阿們人的手.
但卻拯救不到女兒脫離自己所發誓所種下的網羅.
女兒現在成為了.
滿足自己心願下的犧牲品.
我們會覺得耶乏他很討人厭.
他還要怪責女兒.
為何還要出來?.
不出來就沒事.
但從另一個角度去看這悲劇.
就不會難理解耶乏他.
由天堂掉落地獄.
多麼難受痛苦的心情.
其實根本不懂得如何面對自己.
如何面對女兒.
很憎恨自己為何會令女兒這樣.
另外自己又是大男人.
身為父親.
面子尊嚴的問題.
根本沒有勇氣.
不敢在女兒面前認錯.
很痛.
他知道自己才是罪魁禍首.
但不可以承認.
承認了更痛.
唯有怪責女兒.
女兒為何第一個衝出來?.
不出來就沒事.
其實實際是怪誰?.
怪自己.
為何自己這麼壞?.
不發這個誓就甚麼事都沒有.
你害死自己的女兒.
你如何處理?.
沒辦法處理.
正在耶律發他覺得很為難.
很糾結的時候.
他自己也沒有想過.
是誰為他解圍.
是他自己的女兒.

$^{681}$他說:爸爸.
你既然向上帝開口.
就當照你口裡所說的去行.
在我的身上.
因為耶律說已經在仇敵阿們人身上.
為你報仇.
報的這個仇.
當然不是單單只.
阿們人管轄以色列人的仇.
報的仇更加是.
他的父家帶給他的傷痛.
帶給他的仇.
有沒有留意.
耶律發他兩父女很特別.
他們有一個相同的核心價值.
他們很重視承諾.
夠不夠?.
大概和耶律發他的成長背景.
都有些關係.
試想像一個曾經被父家遺棄.
傷害過的人.
自然對人沒有安全感.
沒有信任.
很害怕怎樣?.
再被人遺棄.
再被人傷害一次.
因為這樣.
他很重視承諾這兩個字.
無論他和長老的講述.
甚至和上帝的喜事.
其實都是強調.
承諾一定要做.
承諾是很重要.
對人是這樣.
對上帝更加應該是這樣.
相信耶律發他.
也是一直教導女兒這樣做.
確實遵守承諾.
本身是一個很好的做人態度.
但遵守一個錯誤.

$^{721}$上帝不起悅的承諾.
很可能會造成挽救不到的後果.
事實上耶律發他這個路不可破的承諾.
並不是完全沒有彎轉.
關鍵是要有人.
一個很熟悉上帝話語心意的人.
及時提醒他.
喂!摩西律法已經講明了.
不可以獻人為祭.
你這樣做是得罪上帝.
上帝最憎恨這樣.
所以耶律發他要獻的是贖罪祭.
為了你自己誣蔑,亂說話,得罪神.
害死你的女兒.
來徹底向神認罪悔改.
停止你這樣做.
如果當時.
有一個熟悉,重視上帝話語的人出現.
肯站出來立即煞停耶律發他.
相信就可以阻止悲劇發生.
我還記得認識一位弟兄姊妹.
很多年前.
她是一位很虔誠,很愛主的基督徒.
不過她的媽媽未信主.
她自然很希望媽媽決自信耶穌.
所以有機會她就將耶穌介紹給她認識.
甚至邀請她回教會.
但她的媽媽死都不肯.
堅持不接受.
自然會問她原因.
她說沒有原因.
總之不去.
到後期當她的媽媽患了重病.
即將離世之前.
這位弟兄姊妹就穩牧者.
去探望她的媽媽.
很希望盡最後的努力.
解開她一些我們都不知道的心結.
希望她可以信耶穌.
結果在聖靈的引導下.

$^{761}$這位弟兄姊妹的媽媽終於說了一件事.
她說她曾經信過一個信仰.
當然不是基督教信仰.
這個信仰有一個特點.
就是你不可以告訴別人你信了這個教.
你要發誓.
否則,如果你洩露了.
你和你的家人都會有危險.
毒不毒?.
很毒.
你可以想像.
這位弟兄姊妹的媽媽為了保護自己.
保護家人.
死都不說.
死都不信耶穌.
但最終家人的愛打動.
上帝的恩典憐憫.
說出這個不能說出的秘密.
結果在牧者的帶領下引導她.
真正的上帝不會這樣害人.
不會這樣要人發這樣的毒誓.
慢慢地鬆開了她的心.
最終她願意決志歸信耶穌.
最後回到天家的裡面.
在那個時候幸好有人不斷介入.
有牧者不斷提醒.
她願意接受聖靈的引導.
不過在耶佛他的故事裡.
我們看不到這一幕.
我們看不到祭司.
看不到長老.
這些應該是重視上帝的話語.
為上帝把關的人出來及時阻止.
耶佛他你不要再錯下去了.
某程度上也是司祭向我們呈現出來的.
那種屬靈敗壞光景.
沒有人願意按上帝的吩咐而行.
但同時也提醒今天的我們.
不要這樣.
不可以這樣.

$^{801}$教會群體你和我都有責任.
來阻止減少這些家庭悲劇再發生.
我們要切實在裡面將那些傷痛切斷.
不要讓它傳播下去.
避免上一代的問題.
繼續延禍到下一代的身上.
所以弟兄姊妹今天這段經文.
給了我們幾個回應方向.
第一個就是我們要謹慎自己的言語.
聖經也有說舌頭是最難搞.
有時一句冒失的話.
特別父母對子女一句無心之失的話.
可以很大傷害.
你和我都會明白.
就像耶乏他那樣.
有機會影響到子女的幸福.
事實上我們也有這些經歷.
可能你曾經在上一代聽過這樣的說話.
或者你現在也不知不覺.
告訴你的下一代.
譬如說.
你這麼頑皮.
我就不要你了.
你不乖.
我就不疼你了.
你經常做錯.
你長大了一定會被人罵.
有沒有聽過?.
有沒有說過?.
也有求主赦免.
其實很奇怪.
特別對一個信耶穌的人來說.
他這三句說話.
其實完全是違反聖經教導的.
你覺不覺得?.
我們雖然不乖.
但上帝有沒有不要我們?.
沒有.
繼續抱著你和我.
我們雖然頑皮.

$^{841}$但上帝有沒有不疼我們?.
上帝會管教.
但也會疼我們.
他的管教就是疼我們.
我們現在經常做錯事.
上帝會不會看死我們.
說沒有鬼用?.
確實鬼是不會用我們.
但上帝看我們為寶貴.
他會使用我們.
所以突然間.
假如你曾經被人說過.
或者你現在不知不覺.
在你憤怒的時候.
或者你無室的時候.
這樣對你的下一代說.
請你求上帝幫你.
下次不要了.
上帝既然不是這樣對待我們.
為何我們會這樣對待別人.
甚至乎我們的下一代呢?.
而另一方面.
耶伏他他呈現的某程度上.
是將自己未了的心願.
透過下一代來實現.
而犧牲了下一代的幸福.
這個例子我想不需要多說.
大家都會經歷過.
譬如上一代很想子女怎樣?.
你學琴.
學樂器.
你要做醫生.
做律師.
甚至乎要承繼父業.
爸爸做甚麼生意你都要做.
到底是想他好.
還是只是彌補自己一些人生遺憾.
很值得我們去想.
就好像耶伏他他這樣.
我曾經認識一位弟兄.

$^{881}$他的家人是做生意的.
他的父親很想他承繼這生意.
但他知道自己根本不喜歡做生意.
他喜歡與人聊天,關顧.
你要他每天對著數字.
他就要去讀帳戶,讀法律.
結果令他痛苦不堪.
到最後他還是要承擔父親的生意.
開心與否?其實不開心.
但沒有辦法.
弟兄姊妹有時我們會否不知不覺.
因為我們自己一些未了的心願.
而影響下一代.
一些成長發展的機會.
影響了他們的幸福.
甚至影響了他們遠離上帝.
最後.
我們很需要有人陪伴我們成長.
我們也很需要守望別人成長.
這兩樣都是耶伏他他缺乏的.
沒錯,原生家庭對我們很重要.
在我們的成長中扮演一個不可或缺的角色.
但這不是唯一的因素.
有些人可能沒有一個完整的家庭.
不過他有教會.
他有目姐,有弟兄姊妹.
來守望同行的時候.
雖然他的家帶給他很多成長壓力和傷痛.
但他知道天父很愛他.
他知道有一班弟兄姊妹.
願意關心他,對他不離不棄,撐他.
看好他的時候.
他的生命仍然可以很亮麗.
他的見證亦都可以感動.
幫助到很多人來改變.
所以弟兄姊妹.
求上帝幫我們.
藉著他的經文,話語,恩典.
引領我們成為別人的祝福.
擺脫一些來自上一代,來自家庭的咒咒和傷痛.

$^{921}$我們一起禱告.
主人我們願意將我們每一位交託給你.
主人我們每一個都有上一代.
亦都有下一代.
上帝這一切本來是你命定.
一個傳播祝福.
承接你自己福氣的途徑.
但因著人的軟弱,人的罪.
因著我們不聽你講,不重視你的話語.
我們所傳遞,蔓延開去的.
往往是人與人之間的傷痛和苦痛.
就好像我們今天看到的經文.
但上帝我們知道.
這不是我們的結局.
這亦不是你期望發生在我們身上.
我們知道到今天.
你仍然有恩典,有憐憫.
幫我們走出這些困局.
幫我們來幫助別人擺脫這些傷痛.
主人藉著今天的經文.
我首先求你幫我們.
倘若我們都像耶佛他一樣.
承受上一代傷痛的時候.
願你幫我們擺脫它.
我們不一定能夠忘記它.
但求你幫我們有智慧,有力量.
擺在一個適當的位置.
不要影響我們與你的關係.
不要影響我們與其他人的關係.
讓我們仍然可以靠上帝的恩典.
靠弟兄姊妹的守望.
慢慢地將它化解.
讓它不再在我們生命裡生根建造.
影響我們靠近你.
祝福更多的人.
我又願你幫助我們.
能夠成為別人的守望者.
守望我們身邊的人.
守望教會.
守望我們的下一代.

$^{961}$讓我們真的可以切斷這些傷痛的傳播鏈.
讓我們因著上帝你的緣故.
我們的生命可以出死入生.
成為更多人的祝福和幫助.
為此我恭敬將我們每一位仰望在你的恩手.
願你藉著你的話語.
藉著你的帶領.
幫我們成為美好,輕香的祭.
成為你合用的器皿.
來榮神益人.
願你聽我們在你面前的禱告.
奉耶穌基督你的名來祈求.
阿們.
\newpage



\section{士師記 9:1-6}
\label{sec:bCwRlVK9wDE}
\textbf{2023年8月20日崇拜直播｜陳明泉牧師｜「家醜說出口」系列:寵出個皇兒｜士師記9\_1-6}
\newline
\newline
連結: \href{https://youtube.com/watch?v=bCwRlVK9wDE}{\texttt{https://youtube.com/watch?v=bCwRlVK9wDE}} ~~~~ 語音日期: 2023-08-22
\newline
\newline
\hyperref[sec:Fr77s7J2G5I]{\small{< < < PREV SERMON < < <}}
~
\hyperref[sec:index_chronic]{\small{[返順時目]}}
~
\hyperref[sec:sDc2iE_sQ_g]{\small{> > > NEXT SERMON > > >}}
\newline
\newline
士師記 9:1-6
\newline
\begin{longtable}{cl}
\hline
\hline
章節 & 經文 (和合本修訂版)\\
\hline
9:1 & \begin{tabularx}{0.7\textwidth}{X} 耶路巴力的兒子亞比米勒到示劍他的母舅那裡，對他們和他外祖父全家的人說： \end{tabularx} \\ \\ \relax
9:2 & \begin{tabularx}{0.7\textwidth}{X} 「請你們問示劍所有的居民：『是耶路巴力的眾兒子七十人都治理你們好，還是一人治理你們好呢？』你們要記得，我是你們的骨肉。」 \end{tabularx} \\ \\ \relax
9:3 & \begin{tabularx}{0.7\textwidth}{X} 他的母舅們為他把這一切話說給示劍所有的居民聽，他們的心就傾向亞比米勒，因為他們說：「他是我們的弟兄。」 \end{tabularx} \\ \\ \relax
9:4 & \begin{tabularx}{0.7\textwidth}{X} 他們從巴力‧比利土的廟中取了七十銀子給亞比米勒，亞比米勒用這些錢雇了一些無賴匪徒跟隨他。 \end{tabularx} \\ \\ \relax
9:5 & \begin{tabularx}{0.7\textwidth}{X} 他來到俄弗拉他父親的家，在一塊磐石上把他的兄弟，就是耶路巴力的七十個兒子都殺了，只剩下耶路巴力的小兒子約坦，因為他躲了起來。 \end{tabularx} \\ \\ \relax
9:6 & \begin{tabularx}{0.7\textwidth}{X} 示劍所有的居民和全伯‧米羅都聚集在一起，到示劍橡樹旁的柱子那裡，立亞比米勒為王。 \end{tabularx} \\ \\
[1ex]
\hline
\hline
\end{longtable}
$^{1}$接下來是經文的時候.
今天的經文記載在.
《士司紀》九章一到六節.
《舊約聖經》三百一十二.
請聽我讀.
耶路巴黎的兒子阿彼米勒.
到了事件見他的眾母寇.
對他們和他外祖全家的人說.
請你們問事件的眾人說.
是耶路巴黎的眾子七十人.
都管理你們好呢.
還是一人管理你們好呢?.
你們又要紀念我是你們的骨肉.
他的眾母寇便將這一切話.
為他說話說給事件人聽.
事件人的心就歸向阿彼米勒.
他們說:他原是我們的兄弟兄.
就從巴黎比利土的廟中取了七十射克勒銀子給阿彼米勒.
阿彼米勒用耳顧了西夫土費跟隨他.
他往俄佛拉到他父親的家.
將他弟兄耶路巴黎的眾子七十人.
都殺在一塊盆石上.
只剩下耶路巴黎的小兒子約旦.
因他躲藏了.
事件人和米羅人都一同聚集.
往事件像樹旁的處子那裡.
納阿彼米勒為王.
聖經獨筆.
頂子會平安.
我們今天繼續用舌詩記來思考我們生命真實的狀況.
如果你記得一個星期前.
我們繼續看機電的生命故事.
在機電的一個很成功的生命當中.
我們看到其實他生命裡面有些失敗因子.
在他人生裡面戰績很彪炳.
他帶著三百多人戰勝了米電人十幾萬的大軍.
當有一個這麼彪炳的人生.
但是在他的生命裡面留下一些污點.
其中就是他用以弗德.
就是用以色列人獻上的金耳環打造了以弗德.

$^{41}$昔日他的名字 綽號叫做耶路巴黎.
就是說用巴黎和他來爭戰.
他的生命就是和那些外邦的邪靈邪神來做一些對抗.
但是到了他戰勝之後.
他打造了一個新的以弗德.
成為了好像新一個的偶像一樣.
這個成為了他生命裡面一個很重要的污點.
另外第二個污點就是他有七十個兒子.
因為有很多的老婆.
還有一個妾侍在事件那裡.
這個妾侍生了一個兒子叫做亞比米勒.
亞比米勒的名字的意思就是.
我的爸爸 我的父 是王者.
在基典的生命當中.
雖然以色列文稱他.
或者甚至想他成為一個國家的管理者.
成為王.
不過他在當時來說是拒絕.
但是當生了亞比米勒這個兒子的時候.
其實他心裡面就有一些隱藏著一種沾沾自喜的因子.
就是這一種作王的心態.
不單止在他生命裡面留下了這個落印.
而在他下一代裡面.
今天我們所看的亞比米勒.
基典其中一個兒子.
就將他的生命一些污點呈現了出來.
今天我們去看這段經文.
在時詩記裡面有時候很難看到一些很正面.
很彪炳的例子.
有時候我們看的時候會覺得很納悶.
甚至乎覺得好像很負面.
不過任何負面的經歷.
在聖經裡面啟示一部分我們都要注意的.
我們不是純粹看正面或者看得很亮麗的生命故事.
有時候我們都要進入黑暗.
或者映射我們生命裡面黑暗的地方.
今天這一段經文都照到我們生命裡面.
一些不是很理想的那一面.
希望我們虛心地來聆聽.
在我們聆聽之前.

$^{81}$我想叫大家注意一個字眼.
就是管理這個字.
管理這個字其實在和合本裡面翻譯.
《時詩記》八章二十三節.
其實基典打勝仗的時候.
那些以色列民就說.
那些以色列民叫基典管理他.
不過基典的答話就是.
我不管理你,我的兒子也不管理你們.
唯有耶和華管理你們.
接著到第九章.
管理這個字再出現.
不過那個層次或者用字的深入度提升了.
待會我們用管理.
這個字成為了解釋亞比米勒.
他生平故事的一個很重要的字眼.
究竟亞比米勒怎樣演繹他父親說.
不管理你這件事.
還是他偏行幾路呢?.
我們待會讀的時候繼續去理解.
我們先祈禱求聖靈幫助我們.
讓我們明白他的話語.
同心祈禱.
阿爸天父願你在我們當中帶領我們.
幫助我們眾弟兄姊妹.
讓我們能夠從《時詩記》當中.
看到我們自己生命需要正視的地方.
幫助每一位弟兄姊妹.
在我們生命的時候.
我們打開我們的心扉.
虛心地求你指教我們.
讓孫講的講得清楚.
將你的話語中心告訴弟兄姊妹知道.
讓我們聽見你的訊息.
我們立志要更新改變.
復興我們.
或者照亮那些黑暗的角落.
在我們心靈當中.
被你的話語光照.
恭敬將來領受你話語的時間交托.

$^{121}$願你賜福.
獻上禱告奉求基督耶穌得勝的聖名.
Amen.
剛才我說過.
在聖經裡面.
第八章和第九章.
其實是一個整體.
在《時詩記》裡面.
亞比米勒基澱這個兒子.
聖經的作者不會說一代歸一代.
他用一個字眼去貫穿.
就是用耶路巴力的兒子亞比米勒.
在第九章第一節裡.
聖經作者特別刻意地.
要將基澱的下一代.
連接在基澱去講.
他要我們特別注意的就是.
上一代的故事和下一代是一個延續.
在《時詩記》裡面.
雖然那個時詩是沒有世襲制的.
每一個時詩都是獨立地被上帝的靈興起.
但唯獨是亞比米勒開始有一點變奏.
在經文裡面.
在第九章第一節裡面這樣說.
耶路巴力的兒子亞比米勒到了事件.
見他的眾母舅.
對他們和外祖全家的人說.
請你們問事件眾人說.
是耶路巴力的眾子七十人都管理你們好.
還是一人管理你們好呢?.
你們又要紀念我是你們的骨肉.
為什麼亞比米勒無端端說了這句話呢?.
其實因為前面的經文已經交代了.
因為基澱年紀漸長了.
老邁了.
而且他又去世了.
他死後那些以色列百姓.
又再走回舊的循環裡面.
在八章第三十三節裡面說到.
基澱死後以色列又回去隨從.

$^{161}$諸巴力行邪淫.
以巴力比利圖為他們的神.
以色列不紀念耶和華他們的神.
就是拯救他們脫離四位仇敵之手的.
也不照著耶路巴力.
就是基澱向他們所施的恩惠.
後代他的家.
在經文裡面說到.
整個時代當基澱死後的時候.
那些百姓本來是經巴耶和華上帝的.
基澱一死了之後.
那些人又走回舊路.
拜另外一個新的神叫做巴力比利圖.
巴力比利圖這個字眼.
或者這個外邦神明.
稍後會再出現的.
這樣很絕鬱的.
不過在這裡面說到.
當他死後的時候.
就開始產生權力鬥爭.
你記得了.
剛才說到基澱有多少個兒子.
七十幾個兒子.
還有一個妾侍在士劍這個地方.
生出來的.
亞比米勒某程度上是妾侍所生的私生子.
原則上如果按照世襲的管理位置.
怎樣排隊都排不到他.
因為他是妾侍所生的.
不是自己父家裡面所生的.
名門正取的一個兒子.
所以亞比米勒.
他父親給他這麼霸氣的名字.
我父親是王者.
於是他採取了一種王者的姿態.
他就問他家人.
那些舅父.
他媽媽.
外祖父家等等.
外家.

$^{201}$他就在士劍的地方提出一個問題.
問他那些人.
基澱死後之後.
你想耶路巴力那七十個兒子管理你們好.
還是一個人管理你們好呢.
言下之意.
他不是在說一個很客觀的分析.
是在說數字上的問題.
因為後面他再加多一句.
你們又要紀念.
我是你們的骨肉.
在經文裡面其實.
講到亞比米勒.
他很想將一個管理權.
記不記得管理這個字.
將管理權.
希望他的外家.
他的根據地士劍.
那些的鄉親父老都能夠支持他.
讓他成為了新一代的領袖.
他用的方法.
首先某程度上是煽動一下.
或者挑動一下大家的神經.
令到大家覺得不安.
將來我們由七十個人來管理.
都很嚴重.
都很混亂.
所以就由一個人管理好一點.
但是那個人是誰呢.
亞比米勒說.
你要記住.
我和你是同鄉.
我是一件人.
所以這一班的鄉親父老.
眾母舅.
即是那些舅父.
他媽媽.
等等的外家那些人.
就將這些話.
就告訴他的同鄉士劍人聽到.

$^{241}$士劍人聽到這一番話.
他們的心就歸向亞比米勒.
所以你看到.
在講論裡面.
亞比米勒都很心思熟慮.
也都攻於心計.
首先挑動了那些人.
對於管理的那種恐慌.
第二件事就是.
講到自己是一個合適的人選.
那個合適就是因為在於.
他是自己人.
自己人.
是士劍人.
由他來管理.
就最正常不過了.
也是最好的選擇.
他的母舅聽到這番話.
除了傳向之外.
他們做了一個舉動.
在九章第四節.
記不記得剛才所說的.
那些以色列人.
隨從巴力.
他們就在巴力廟.
第四節.
就從巴力比利土的廟中.
取了七十赤赫勒銀子給亞比米勒.
拿了一些香油錢.
拿了香油錢.
就拿了七十赤赫勒.
就給了亞比米勒.
亞比米勒拿這些錢來做什麼呢.
經文再繼續說.
亞比米勒就用以.
僱了一些匪徒來跟隨他.
請了一班僱傭兵.
這些僱傭兵是什麼人呢.
不是一些志願軍.
這些是為了錢而打仗的.

$^{281}$本身是匪徒.
貪錢的一些人.
用一些落草為寇的.
不務正業的這些人.
成為了他的士兵.
用了七十赤赫勒就請了這班人.
請了這班人就壯大了自己的軍事力量.
當他組成了這個軍事力量之後.
他們做了什麼呢.
在第五節.
然後就浩浩蕩蕩.
整隊人馬.
他往俄弗拉.
俄弗拉就是基定自己的鄉下.
到他父親的家.
將他弟兄耶路巴力的眾子.
七十人都殺在一塊盆石上.
只剩下耶路巴力的小兒子約坦.
因為他躲藏了.
士劍人和米萊人都一同的聚集.
往士劍匠樹旁的處子那裡.
立亞比米勒為王.
整件事你會看到.
很明快地簡單記載.
亞比米勒說.
因為他是自己人幫他管理.
所以整個鄉親父老就給錢他.
供應他一些養分.
拿了這些錢之後.
他就僱傭了這些匪徒成為他的士兵.
這些士兵是做什麼呢.
去打仗.
去打什麼人呢.
其實是打反以色列自己的人.
是他父家的兄弟.
他父親所生的七十個.
有機會問鼎世襲領袖位置的一班人.
在經文裡面特別強調.
將這七十人殺在盆石上.
如果你記得.

$^{321}$其實聖經的作者有意要你回憶.
如果你說盆石.
其實在第七章.
士司紀第七章.
基澱就在盆石那裡獻祭.
要驗證那個是耶和華上帝的使者.
那塊盆石昔日是他父親用來驗證上帝的一塊盆石.
他未必是同一塊石.
不過同樣的盆石出現的時候.
到他的兒子的時候會做什麼呢.
就是兄弟互相殘殺.
或者這麼說.
一個殘殺七十個.
全部死在盆石上.
整件事是很血腥的.
令你一邊想的時候.
一代就可以轉換成一個這麼恐怖的經歷.
經文裡面繼續說.
為什麼亞比米勒會有這樣的能力呢.
亞比米勒有這個能力.
原來其實是來自他的家族.
是來自士兼人.
是他母親的家.
來製造了一些養分.
來煉成了一個這樣的霸王出來的.
在經文裡面說到.
這一班的鄉親父老.
在巴黎比利土的廟裡面拿到一些香油錢.
來製造一個機會給他.
讓他可以組織成為一個軍事力量.
讓他成為新一代的霸主.
在士兼人來說.
他覺得可能是為了他們自己的家鄉.
成為了一個比其他鄉族高一等的地方.
不過在聖經的作者.
在這裡看.
這件事一點都不光彩.
他們所用的錢是巴黎廟裡面的香油錢.
不是來自上帝的靈充滿.
這個領袖不是由上帝的靈充滿.

$^{361}$在一個人身上.
上帝揀選的一個人成為他們的領袖.
而是透過血腥.
用暴力的方法.
用殺戮的方法.
來將一個位置加諸在亞比米勒的身上.
經文另外刻意要我們注意.
整個顧名思義.
整個時詩記是在說時詩.
所以領袖是時詩.
時詩這個身份.
是上帝給的一個判斷.
或者是一個軍事領袖.
是打外邦人的.
不過去到這裡的時候.
亞比米勒就不是做時詩了.
去到橡樹旁.
是那些人立亞比米勒為王.
這個是時詩記裡面一個很重要的變奏.
很多人想找自己心目中的領袖.
很想一個人有效的來管理他.
不過在這裡他們用了人的方法.
用了人的暴力.
用了血腥的方法.
來建立一個領袖.
所以這個領袖就成為了.
他們之後的整個事件.
或者整個以色列的網羅.
在這裡我們稍微停一停.
想一想.
小小題外話.
不過都值得我們去注意.
亞比米勒這個人怎樣煉成出來的呢?.
在經文裡面.
首先他講到機電是有責任的.
所以在經文上上下下.
他不斷重複耶路巴力這個字.
是在講耶路巴力.
他生命裡面的污點.
這種70個老婆.

$^{401}$兒子之間有很多的張力.
這個家庭變得很複雜.
耶路巴力他在事件地生了這個兒子.
為他的家族產生了一些不可或缺的網羅.
你又會問為什麼無端端會講事件和俄弗拉.
這兩個地方.
俄弗拉就是機電的自己地方.
事件為什麼會在這裡生了一個兒子呢?.
如果你看回聖經.
或者你理解米電人之前打的.
其實之前打米電人.
集結就是在事件這個地方.
打完仗之後.
那個將領很浪漫的.
然後結識了一個女子.
有些的聖經者說.
那個妾侍其實是一個妓女.
在那裡打完仗.
就和一個妓女發生關係.
她又因為是妓女.
所以不可以明媛正娶.
所以就成為了她另外一頭家.
在距離她自己俄弗拉.
大概三四十公里的一個地方.
裡面建立了另外一頭家.
耶路巴黎在這裡是有責任的.
在他死後這些問題浮現.
而這個問題繼續惡化的就是.
他身邊的人.
對他的那種.
好像期望也好.
又或者被他煽動也好.
他這個很想作王.
很想作管理的這個人.
帶來了一些很多的養分.
今天我們留意一下.
我們今天當然家裡不會有亞比米勒.
不過在我們家裡.
或者在很多的弟妹的家庭裡.
可能一兩個小霸王都會出現.

$^{441}$這個小霸王我們說那個小孩很頑皮.
我們會將矛頭指向.
屬於那個小孩身上.
不過更加多的其實未必是那個小孩.
天生出來是有問題.
更加多的問題是在於可能.
在他的家庭裡.
在他的外圍的外祖家.
父家等等那裡.
煉成了一個小霸王出來.
這個霸王不是純粹.
他的性格使然會這樣.
可能這裡也有一小部分.
但更加多的部分可能是在家裡的教養.
身邊的人的驕寵.
甚至製造很多的養分.
給很多的資源.
來哄他開心.
或者出於自己的擔心.
所以什麼都要順著霸王的心裡面.
他想要什麼就給什麼.
所以就製造了一個這樣的小霸王.
今天我也要很坦誠地說.
今天很多的坊間裡.
有一些教導孩子的方法.
有時候不是很適合聖經.
甚至有時候坊間裡有些教導方法.
可能是著重孩子怎樣發展.
但未必是將真理.
將對的東西教導給他.
我們今天很著重的就是.
用什麼方法令小孩子開心.
但很少去想.
我怎樣令一個孩子正直.
或者是一個對的環境裡長大.
有時候上一代或者父母.
想用哄他快樂的方法.
來讓兒子快樂.
但甚少去明白.
或者去貫徹的.

$^{481}$讓孩子在對的情境裡成長.
我自己做父母也知道.
有時候兒女有意無意之間.
會將自己生命的缺點放大.
先說一些秘密.
前幾天我和女兒聊天.
她跟我說.
爸爸其實我也有幾分像你.
她說放心不是樣子像你.
是很嚴重的.
像我樣子很嚴重我明白的.
她說性格有點像我.
我說有什麼像我呢.
她說有一些擺款.
她跟我說.
她說有些擺款喜歡指點江山.
她覺得自己有一點領導的能力.
我說不是領導能力.
是指點人做事的能力.
她說是呀,我有些很像你的這方面.
然後一邊說的時候.
開始的時候我也沾沾自喜.
自己的子女像自己多開心.
但越說越時候我就越害怕.
我想如果我女兒將來朝著.
她覺得像我.
將我某一面她覺得很棒的東西.
但原來可能不是很理想的.
性格缺陷放大了.
用她人生幾十年來放大.
其實可能是一個很大的問題.
然後她說完之後我說.
女兒有些東西你應該要像爸爸.
但有更多的東西你要像耶穌.
你不要像我.
像爸爸有時候很嚴重.
爸爸也不是一個完美的人.
弟兄姊妹我們有時候會不會這樣呢.
我們生命裡面可能整體是好的.
不過我們生命裡面有一些失敗的因子.

$^{521}$可能你的兒女看了不順眼.
如果你繼續覺得引以為傲.
而且給她養分.
可能你會久而久之製造一個新的阿比米勒.
機電的生命本身.
不是一個敬虔的人去驗證上帝的人.
即是說他有這些污點.
但通盤來說他也是上帝重用的僕人.
但這些污點如果沒有處理的時候.
就遺棄到下一代的時候.
這些污點就會放大.
成為了其他人的網絡.
在這裡我們看到阿比米勒的生命.
他的出現就是因為他的父家.
他的母系那邊來為他製造養分.
以致他製造了一個這麼大的悲劇.
將他其他的七十個兄弟殺在盤石上.
經文再繼續說.
到了第十六節至二十節.
經文裡面就說到.
七十個兒子死了.
唯有一個逃脫了.
那個叫約坦.
在耶路巴黎其中一個兒子.
小兒子.
大概他差不多是最小的.
他就逃脫了.
逃脫的時候.
後來有人就告訴約坦.
比米勒已經成為了王.
已經被納為王.
這個小兒子他很勇敢.
在經文裡面就交代了他.
在第七節就說到.
他向所有的以色列百姓和使檢人.
特別是使檢人說了這一番話.
第七節到十五節我讀給大家聽.
使檢人啊!你們要聽我的話.
神也就聽你們的話.
有的時候樹木要高一樹為王.

$^{561}$管理他們就去對橄欖樹說.
請你作我們的王.
橄欖樹說.
回答說.
我豈能指著供奉神和尊重人的柔.
飄搖在種樹之上呢?.
樹木對無花果樹說.
請你來作我們的王.
無花果樹.
回答說.
我豈肯指著所結甜味的果子.
飄搖在種樹之上呢?.
十二節樹木對葡萄樹說.
請你來作我們的王.
葡萄樹.
回答說.
我豈肯指著使人和人喜樂的新酒.
飄搖在種樹之上呢?.
種樹對荊棘說.
請你來作我們的王.
荊棘.
回答說.
你們若誠誠實實地高我為王.
就要投在我的任下.
不然願火從荊棘裡出來.
消滅黎巴嫩的香柏樹.
七至十五節就是走掉的藥毯.
用一個預言的方式.
來交代亞比米勒作王之後的事.
這件事用這個比喻說什麼呢?.
很隱晦的.
這個比喻說什麼呢?.
他用樹木來做一個比喻.
他說有幾種樹.
有橄欖樹,無花果樹,葡萄樹,荊棘和香柏樹.
這幾種樹.
他就說其中有一段時間.
樹木說我要高一棵樹作為樹中之王.
對橄欖樹說.
橄欖樹,現在我請你成為我們樹中之王.

$^{601}$橄欖樹接不接受這個邀請呢?.
橄欖樹拒絕.
他說我怎麼可以停止了供奉神和尊重人.
來製造油出來給別人.
我的本份是製造橄欖油.
我怎麼可以停止了這個本份.
然後就在種樹之上來飄搖.
指點別人降山呢?.
所以這個答法就拒絕了.
橄欖樹不肯了.
於是就問無花果樹.
無花果樹就說你做我們的王吧.
無花果樹說我不會.
我怎麼可以停止結出甜美的果子.
吃無花果.
我怎麼不結無花果來做王呢?.
所以我就不做了.
接著十二節.
其他樹見他拒絕就去找葡萄樹.
葡萄樹說你做我們的王吧.
葡萄樹又拒絕.
他說怎麼可以停止了產生神和人喜樂的新酒呢?.
然後就飄搖在種樹之上.
我的職責是結出葡萄子.
所以我是做新酒的.
我不會做王的.
之前的幾種樹全部都有出產的.
橄欖油,無花果和葡萄樹結出果子.
全部都有收成結果的.
於是就找個比較空閒的.
就找荊棘.
大家知道荊棘是什麼植物吧.
很多刺的,爛藤的那些.
荊棘就沒有果子的.
荊棘就請他來做王.
荊棘的回答就跟其他樹很不同了.
荊棘就說.
如果你是誠誠實實的高我為王.
就要投在我的蔭下.
不然願火從荊棘裡出來.

$^{641}$消滅黎巴嫩的香柏樹.
在這個說法裡面.
你會覺得很荒謬的.
荊棘其實是爛藤和一些尖銳的刺在那裡.
他的蔭下是不是可以保障人的呢?.
他的蔭下是沒有的.
只是用荊棘的刺刺人而已.
荊棘本身是不能夠製造出產的.
沒有成果的,沒有果子的樹.
他不用去結果.
不過找到的時候.
他就成為了一個刺痛別人的植物.
再加上荊棘是很易燃的.
很容易燒著的.
所以荊棘就用威嚇的話說.
如果你不誠實地尊敬我,高我.
投在我的蔭下.
就會有火在我裡面燒著你們.
消滅黎巴嫩的香柏樹.
為什麼會說黎巴嫩的香柏樹呢?.
當時真是樹中的王.
就是香柏樹.
最好的木材就是香柏樹了.
因為香柏樹就是做宮殿,聖殿.
甚至未來宮殿裡面樹中之王就是香柏樹.
整件事在這個比喻裡面.
約坦就用一個很荒謬的說法.
就是暗指亞比米勒.
就好像荊棘一樣要作王.
他既不是帶來真正為人們的管理.
帶來一個真正的福祉.
反而是帶來大量的禍患.
他的火成為了眾人的網絡.
後來約坦就在這裡.
其實他再解釋19節.
你們如今若按誠實,正直地.
待耶路巴黎和他的家.
就可以因為亞比米勒而得歡樂.
他也可因你們得歡樂.
不然,願火從亞比米勒發出.

$^{681}$消滅事件人和米羅眾人.
又願火從事事件人和米羅人出來.
消滅亞比米勒.
約坦在說什麼比喻呢?.
是在說一個預言.
一個很荒謬的人.
透過軍事力量成為王.
一個很荒謬的人.
他想著擁有權力.
就能夠管理身邊的人.
但這是一個荊棘作王的比喻.
他表面上好像說.
事件人如果正直誠實地對待他.
大家就好了.
不過在約坦的口中.
不是這樣的.
基本上這班人.
是沒有辦法和他和諧共處.
到最後.
亞比米勒.
會有火從亞比米勒發出.
消滅事件人和米羅眾人.
後來亞比米勒也會被消滅.
這是一個預言.
是一個和幻的預言.
也是說昔日這些事件人.
為什麼要擁納他為王.
約坦也特別針對.
這些事件人.
昔日給錢亞比米勒.
來殺害其他兄弟.
將來他們也會面對後患無窮.
說完這句話之後.
約坦就消失了.
在聖經裡再沒有記載他的出現.
或者說這個人好像不在當中.
但這個就好像一個緊箍咒.
一個咒詛在這班人身上.
然後聖經就繼續記載整件事的發展.
九章二十二節至第二十九節.

$^{721}$就記載著這件事件是怎麼發生的.
亞比米勒管理以色列人三年.
神使厄摩降在亞比米勒和事件人中間.
事件人就以鬼咒代亞比米勒.
這是要叫耶路巴力七十個兒子.
所受的殘害歸於他們的哥哥亞比米勒.
又叫那流他們血的罪歸於幫助他殺弟兄的事件人.
事件人在山頂上設埋伏.
等候亞比米勒.
凡從他們那裡經過的人.
他們就搶奪.
有人將這事告訴亞比米勒.
以別的兒子迦勒和他的弟兄來到事件.
事件人都信求他.
事件人出城到城填間去摘下葡萄.
穿走設擺筵席.
進他們神的殿中吃寒.
咒詛亞比米勒.
二十九節.
惟願者民歸我的手下.
我就除掉亞比米勒.
迦勒又對亞比米勒說.
增添你的軍兵出來吧.
這件事交代什麼呢.
整件事很人為操縱的.
權力鬥爭.
荒謬絕倫.
上帝在哪裡呢.
在《士師記》裡面.
是不是上帝任由這些亂局發生呢.
在作者裡面.
聖經的作者不會.
上帝是掌管整個的局面的.
不過在經文裡面.
他叫我們特別要注意的就是.
亞比米勒真的管理.
以色列人只有三年的時間.
一個短的時間.
對於受苦的人來說很長的.
但其實在《史詩》的時間裡面.

$^{761}$不是很長的.
三年.
第二個字需要特別注意的.
就是管理這個字.
待會我會再說的.
不過管理這個字就升級了.
他用統治這個字.
聖經裡面只有荷西亞書.
在說統治天使.
用這個字出現兩次.
由本身的管理變成統治.
待會我會再解釋這個字眼.
不過在他統治了幾年間.
上帝做了什麼工作呢.
上帝是一個.
統盤處理的上帝.
惡事都成為上帝的服役.
上帝將惡魔降在.
亞比米勒和事件人中間.
這是一個很形象化的說法.
昔日的事件人給錢.
來遵座亞比米勒.
做王者.
就是事件人.
但當他做了王之後.
上帝就不祝福整件事.
上帝甚至將惡魔.
將和環降到亞比米勒.
和事件人中間.
在他們中間做了什麼呢.
上帝的工作令到.
任憑他們.
他們自己打生打死.
互相來到差界.
那些事件人不相信亞比米勒.
同時亞比米勒也更不相信事件人.
所以兩者就製造了一個.
互相鬥爭的局面.
在二十四節說到.
這是要應驗.

$^{801}$要報昔日的血債血仇.
就是七十個耶路巴黎的兒子.
被殘害的這件事.
二十五節就說到.
事件人做了什麼來頂撞亞比米勒.
亞比米勒有任何的人經過.
有什麼軍隊也好.
或者有人員走過某些地方.
其他的事件人就切埋伏.
要把那些人搶奪.
搶走他們的東西.
甚至殺了他們的人.
矛頭是直指亞比米勒.
另一方面那些事件人.
也希望擁立一個新的領袖.
來取代亞比米勒.
在經文裡就交代了一個人叫加勒.
其實某程度上在經文裡是Cathefer.
沒什麼特別去說.
不過他想製造一個.
讓你知道的圖畫.
亞比米勒用強權成為了一個管理者.
但是否真的得到那些事件人的馴服呢?.
是否得到那些所有人的愛戴呢?.
不是.
而是製造了一個互相紛爭.
鬥爭的局面.
當這種鬥爭出現的時候.
亞比米勒也不是省油的燈.
他也會報應這些人所作的行為.
其中剛才所說的一個加勒這個Cathefer.
他就被那些事件人又再擁立了.
然後就說他心紅的第二度.
來到這個事件的地方.
於是就希望挑動他.
發起一個軍事力量.
讓加勒成為新一代的霸主.
所以在二十九節裡面說.
為願者們歸到我的手下.
我就除掉亞比米勒.

$^{841}$加勒又對亞比米勒說.
增添你的軍兵出來吧.
挺豪氣的.
一個好像很有霸氣的梟雄.
跟他決一死戰.
在這裡我們看到.
上帝看到一個這麼混亂的局面.
上帝不是完全沉默的.
甚至在神學裡面的角度來說.
那種最混亂的狀態.
是上帝促成的.
那種混亂的本身.
不是上帝的心意.
而是這群人在荒謬的處境當中.
擁立亞比米勒為王.
事件人都有角色的.
所以上帝就將惡魔降在他們中間.
令到他們互不信任.
製造了一個極之鬼打鬼的局面.
我們看到這些工作.
我們覺得很納悶的.
我們看上去.
上帝你容許這些亂事發生.
我們最期望的就是上帝.
從高而下.
然後有一個強人的領袖.
來平息所有的混亂.
不過在聖經的作者裡面.
講到事思的時代.
各人任意而行.
上帝有時候任憑他們自己自食其果.
他不想過一堂.
知道有什麼混亂帶來的功課.
他們是不會學得精.
有時候我們人生也是.
我們偏行幾路.
我們進入混亂的狀態.
有時候上帝可能會透過一些.
不是太好的東西.
來讓我們自己吃回自己自食其果.

$^{881}$以致我們掌握到一些功課.
如果我們不通過.
我們是不會懂得去尊崇上帝.
或者會懂得去迴避.
在聖經裡面也講到.
這群人繼續執意而行.
製造了一個新的混亂的局面.
另外一個字眼.
剛才說的管理這個字.
管理這個字.
跟剛才說的八章二十二二十三節.
說的以色列人擁戴機電.
管理他那個字是不同的字.
在這裡九章二十二節出現的管理.
這個是統治這個意思.
是一個超然極高的權兵.
用權力來管治.
統治一個很極端的字眼.
聖經作者用這個字.
不是尊崇阿彼米勒.
他真的運籌為握.
掌握權力.
而是在諷刺他.
阿彼米勒管你.
這一群人的三年.
上帝就將惡魔降在他們中間.
當他覺得擁有了權力.
等於能夠管理的時候.
上帝就要打破這些想法.
真正的管理從來不是製造了更多的混亂.
真正的管理也不是製造了人互不信任.
如果我們覺得這種統治是來自上帝.
大概都是一廂情願的想法.
經文的作者就用這個字來諷刺阿彼米勒.
那三年是混亂的.
管理其實是上帝給人的創世裡面.
一個很重要的責任.
你和我每個人都有管理.
不過管理和統治是兩碼子的事.
我們要時間管理.

$^{921}$我們要關係的管理.
我們要管理我們眼前做的不同的任務.
調配不同的資源.
這是上帝給我們的智慧和職份.
來做好眼前的工作.
這就是真正的管理.
在薩斯的時代.
上帝委託某一些領袖攻打其他外邦.
不容讓外邦的文化侵入以色列文化當中.
能夠按照正義來做判斷.
這就是薩斯的管理責任.
雖然用管理這個字眼來看.
不過在薩斯記裡面從來沒有用管理這個字.
套用在時事身上.
唯獨是阿彼米勒用管理這個字來說他.
其實是某程度上潤他.
他所說的管理其實一點都不是管理.
只是帶來混亂.
帶來更加多的悲瘡.
弟兄姊妹今天我們要負心自問.
可能在職場裡面有些人說.
我們管理階層.
我們是為人帶來混亂.
還是實踐上帝給我們生命的召命.
很多人今天在職場也好.
在生活裡面也好.
他覺得以為掌握一些權力.
或者製造一些操控的能力.
就能夠很好的來管理.
不過在聖經的作者裡面不是這樣的.
真正的管理從來都不是要成為一個統治者.
一個操控別人的人.
真正的管理反而是管理我們自己的生命.
在新約的作者裡面說.
我們真正的管理是要苦乞己生.
叫身服我.
制服你生命裡面的那隻怪獸.
面對你自己生命的黑暗.
不要容讓你的肉體成為你生命的王.
那個才是真正的管理.

$^{961}$真正的管理不是要嚴苛地對待身邊的人.
真正的管理是能夠發揮上帝在你生命裡面所擁有的本質.
好好地實踐上帝的照明.
今天我們的管理帶來什麼結果呢?.
在亞比米勒裡面.
他的管理帶來了操控.
他的操控帶來了猜忌.
帶來一個像剛才約坦所說的語言那種苦果.
經文裡面再繼續說.
亞比米勒聽到或知道那些村民.
那些事件人對他的那種反抗.
在第36節.
接下來就記載著.
一路整個段落就說到亞比米勒如何去回敬那些對待.
反對他的人.
亞比米勒記載著.
知道有人這樣對待他.
他先收服對待興起上來很囂張不護的加勒.
在36節裡就說到他運用自己的軍事力量埋伏.
就要驅散加勒.
凡是見到加勒的那些人.
就要埋伏地打仗.
去針對他 驅散他.
經文裡面很仔細.
我們當然沒有時間仔細地說.
我跳過了.
大概你知道那個結果就是.
加勒剛才這個叫聲就沒有了.
不過我想多說一點.
在後一點的.
在第46節開始.
就說到那些事件人的結局.
46節事件人.
事件人聽見了.
那些人知道亞比米勒的大軍.
那些軍事力量殺了加勒之後.
他就乘勝追擊追趕其他的事件人.
那些事件人聽見了.
就躲進巴黎比利土的廟裡的圍鎖.
即是廟旁邊住的那些圍鎖.

$^{1001}$有人告訴亞比米勒說.
事件人都聚集在一處.
亞比米勒和跟隨他的人就上到薩門山.
亞比米勒手拿斧子砍下一根樹枝.
供在肩上.
對跟隨他的人說.
你們看我所行的.
當趕緊照樣行.
各人就各砍一枝.
跟隨亞比米勒把樹枝堆在圍鎖的四周.
放火燒了圍鎖.
以致事件留的人都死了.
男女約有一千.
在46節交代什麼呢.
應驗約旦所說的.
荊棘是會著火的.
亞比米勒要追殺那些反對他的事件人.
知道那些人全部躲在巴黎比利土.
昔日拿香油錢的那間廟.
那些人聚集在那裡.
於是亞比米勒就砍了一棵樹.
剁了一根樹枝.
然後叫全部人跟隨我這樣做.
全部人拿了樹枝之後就圍著圍鎖.
用火燒了.
當時燒死了一千個男女.
很殘酷吧.
然後見到這件事之後.
死了一千人.
還有一些人還沒死.
亞比米勒再去提比斯.
打算在那裡安營.
再繼續攻取另一個城.
如果誰不聽他的話.
他繼續殺戮.
城裡面有一個很堅固的城樓.
亞比米勒見到這個堅固的城樓.
都很難攻的.
之前用火燒巴黎比利土的圍鎖.
如今他又再用照板煮碗.

$^{1041}$走到城樓那裡.
拿那些樹枝圍著柴燒城樓.
因為城樓很高.
遠火燒不著.
所以要走近一點.
走到城樓的牆上.
堆那些樹枝.
打算堆高木柴來攻城.
用火燒城.
在他走近城樓的時候.
忽然間在城樓裡.
萬不經意.
或者那些人都是奮力抵抗.
其中一個女人.
有個婦女.
拿著一個上磨石.
就在城樓那裡扔下去.
因為城樓很高.
有險可守的.
見到那些人用火來圍攻的時候.
那個女人.
所有人就找東西扔下去.
而這塊石就是.
剛剛那個女人拿上磨石.
磨的那些石.
很重的.
用盡九牛二虎之力.
一扔.
一扔就扔中了.
亞比米勒的頭.
扔穿了他的頭.
頭破血流.
其實都是一個重傷.
當他的頭.
亞比米勒的頭受傷的時候.
他就說.
叫身邊帶著兵器的年輕人.
拿把刀出來.
用刀刺繩.
為什麼他這樣說呢?.

$^{1081}$免得我被人家後世來取笑.
我是死在一個女人的手下.
他自己知道那塊石.
是一個女人扔下來的.
他就要叫士兵插死她.
英勇而死.
至少戰死.
也有一些傳統就是說.
被女人殺死是永不超生的.
在中東有一些這樣的傳統.
無論如何.
亞比米勒就被這個少年人.
一刀刺下去就死了.
整件事.
你會看到.
是一個悲劇收場.
在聖經裡面.
56-57節9章.
這樣.
神報應亞比米勒.
向他父親所行的惡.
就是殺了弟兄七十個人的惡.
是劍人一切的惡.
神也都報應在他們頭上.
耶路巴力的兒子約旦的就坐.
歸到他們的身上.
整件事是一個就坐的結束.
耶路巴力這個名字.
本來是代表著.
驗證上帝的同在.
和巴力來到去鬥過.
但是隨著亞比米勒.
隨著士劍人.
全部的死亡.
這個就好像一個悲劇.
為機電的事跡作一個落幕.
當我們讀的時候.
我們心裡面會覺得不舒服.
為什麼好好的一個時代.
會搞成這樣呢?.

$^{1121}$在聖經裡面.
今天他告訴我們.
很多時候一個時代也好.
或者一個悲劇的產生.
很多時候是源自一個人生.
生命裡面一個罪惡因子.
作為父母上一代.
我們要留意我們的生命.
有沒有這樣的罪惡因子.
如果你不好好對付它.
這些會在你的下一代裡面放大.
它會用一生來將你的因子.
將你的問題放大.
當你繼續身邊的人給養分給它的時候.
這些罪惡會越來越糾纏.
製造的悲瘡會越來越大.
在這裡提醒我們作父母.
作上一代.
今天除了要學親子要怎樣溝通.
講一些技巧上的東西.
聖經在說.
真正需要注意的就是我們的生命.
我們是一個怎樣的上一代.
我們是一個怎樣的人.
在家裡我是一個怎樣的父母.
今天如果你純粹將你的注意力放在技巧上.
你會有什麼影響呢?.
如果你純粹將你的注意力放在技巧上.
你到最後還是不能幫助下一代正直不阿地去成長.
在基澱的事蹟裡面.
很清楚地告訴我們.
這些罪惡因子.
去到亞比米勒無限地放大.
第二件事是需要我們注意的.
什麼叫做管理呢?.
管理從來不是說擁有權力就等於管理.
在你的生命裡面.
你的地界,你的任務.
你要做的事你有沒有好好做好.
今日坊間裡面很多人以為擁有權力.

$^{1161}$掌握一些實力.
就能夠統治身邊的人.
在聖經裡面真正統治.
真正管理的不是我們自己.
是只有耶和華上帝.
唯有基督耶穌成為我們的主.
主宰我們的生命.
我們不要固執我們自己擁有的權力.
我們覺得重要的東西.
我們真正最重要的不是這些外在的東西.
真正最重要的是基督耶穌.
我們是不是按照真理來管治我們自己的生命.
我們要攻克己身.
叫身服我.
不要我們將福音傳給身邊的人.
自己反落在黑暗裡面.
這是經文裡面第二個提醒我們.
第三件事.
我們不要輕忽上帝的管教.
上帝有時候製造一些混亂的局面.
其實告訴我們我們種什麼收什麼.
有時候上帝任憑一些東西出現在我們生命當中.
上帝的任憑其實是一個很大的刑罰.
我們不要想著我們做某一些惡事.
沒有人知道.
順風順水.
我們就瞞天過海.
我們時刻都要求上帝監察我們的內心.
讓我們不要執意而行.
不要再偏行己路.
今日亞比米勒的故事.
耶路巴黎兒子的故事.
成為了我們生命一個很重要的警惕.
求主繼續來對我們每一個人講說話.
拆亮我們監察我們的內心.
這是我對他的內心.
\newpage



\section{士師記 13:8-23}
\label{sec:sDc2iE_sQ_g}
\textbf{2023年9月3日崇拜直播｜陳明泉牧師｜「家醜說出口」系列:錯焦的敬虔｜士師記13\_8-23}
\newline
\newline
連結: \href{https://youtube.com/watch?v=sDc2iE_sQ-g}{\texttt{https://youtube.com/watch?v=sDc2iE\_sQ-g}} ~~~~ 語音日期: 2023-09-04
\newline
\newline
\hyperref[sec:bCwRlVK9wDE]{\small{< < < PREV SERMON < < <}}
~
\hyperref[sec:index_chronic]{\small{[返順時目]}}
~
\hyperref[sec:Up7U5HlOZfo]{\small{> > > NEXT SERMON > > >}}
\newline
\newline
士師記 13:8-23
\newline
\begin{longtable}{cl}
\hline
\hline
章節 & 經文 (和合本修訂版)\\
\hline
13:8 & \begin{tabularx}{0.7\textwidth}{X} 瑪挪亞祈求耶和華說：「主啊，求你再差遣那神人到我們這裡來，指示我們對這將要生的孩子該怎樣作。」 \end{tabularx} \\ \\ \relax
13:9 & \begin{tabularx}{0.7\textwidth}{X} 神垂聽了瑪挪亞的聲音。那婦人坐在田間的時候，神的使者又到她那裡，但是她的丈夫瑪挪亞沒有同她在一起。 \end{tabularx} \\ \\ \relax
13:10 & \begin{tabularx}{0.7\textwidth}{X} 婦人急忙跑去告訴丈夫，對他說：「看哪，那日到我這裡來的人又向我顯現了。」 \end{tabularx} \\ \\ \relax
13:11 & \begin{tabularx}{0.7\textwidth}{X} 瑪挪亞起來，跟隨他的妻子來到那人那裡，對他說：「你就是跟這婦人說話的那個人嗎？」他說：「是我。」 \end{tabularx} \\ \\ \relax
13:12 & \begin{tabularx}{0.7\textwidth}{X} 瑪挪亞說：「現在，願你的話應驗！這孩子該如何管教呢？他當做甚麼呢？」 \end{tabularx} \\ \\ \relax
13:13 & \begin{tabularx}{0.7\textwidth}{X} 耶和華的使者對瑪挪亞說：「我告訴這婦人的一切事，她都要遵守。 \end{tabularx} \\ \\ \relax
13:14 & \begin{tabularx}{0.7\textwidth}{X} 葡萄樹所結的不可吃，清酒烈酒都不可喝，任何不潔之物也不可吃。凡我所吩咐的，她都當遵守。」 \end{tabularx} \\ \\ \relax
13:15 & \begin{tabularx}{0.7\textwidth}{X} 瑪挪亞對耶和華的使者說：「請容許我們留你下來，好為你預備一隻小山羊。」 \end{tabularx} \\ \\ \relax
13:16 & \begin{tabularx}{0.7\textwidth}{X} 耶和華的使者對瑪挪亞說：「你雖然留我，我卻不吃你的食物。你若預備燔祭，就當獻給耶和華。」因瑪挪亞不知道他是耶和華的使者。 \end{tabularx} \\ \\ \relax
13:17 & \begin{tabularx}{0.7\textwidth}{X} 瑪挪亞對耶和華的使者說：「請問大名？好讓我們在你的話應驗的時候尊敬你。」 \end{tabularx} \\ \\ \relax
13:18 & \begin{tabularx}{0.7\textwidth}{X} 耶和華的使者對他說：「你何必問我的名字呢？我的名字是奇妙的。」 \end{tabularx} \\ \\ \relax
13:19 & \begin{tabularx}{0.7\textwidth}{X} 瑪挪亞取一隻小山羊和素祭，在磐石上獻給耶和華。他行奇妙的事，瑪挪亞和他的妻子觀看， \end{tabularx} \\ \\ \relax
13:20 & \begin{tabularx}{0.7\textwidth}{X} 火焰從壇上往上升，耶和華的使者也在壇上的火焰中升上去了。瑪挪亞和他的妻子看見，就臉伏於地。 \end{tabularx} \\ \\ \relax
13:21 & \begin{tabularx}{0.7\textwidth}{X} 耶和華的使者不再向瑪挪亞和他的妻子顯現了。那時，瑪挪亞才知道他是耶和華的使者。 \end{tabularx} \\ \\ \relax
13:22 & \begin{tabularx}{0.7\textwidth}{X} 瑪挪亞對他的妻子說：「我們一定會死，因為我們看見了神。」 \end{tabularx} \\ \\ \relax
13:23 & \begin{tabularx}{0.7\textwidth}{X} 他的妻子卻對他說：「耶和華若有意要我們死，就不會從我們手中接受燔祭和素祭，不會將這一切事指示我們，這時也不會讓我們聽到這話。」 \end{tabularx} \\ \\
[1ex]
\hline
\hline
\end{longtable}
$^{1}$今日的經份是我們會一起看舊約聖經.
《釋詩記》十三章八到二十三節.
大家可以打開舊約聖經三百一十七頁.
請聽我讀出.
嗎哪亞就祈求耶和華說:主啊,求你再差遣那神人到我們這裡來,好指教我們怎樣待這將要生的孩子..
神應許應允嗎哪亞的話:婦人正坐在田間的時候,神的使者又到她那裡,她丈夫嗎哪亞卻沒有和她一在一處..
婦人急忙跑去告訴丈夫說:那日到我面前來的人,又向我顯現..
嗎哪亞起來跟隨他的妻來到那人面前對他說:與這婦人說話的就是你嗎?.
她說:是我.嗎哪亞說:願你的話應驗,我們當怎樣待這孩子,他後來當怎樣呢?.
耶和華的使者對嗎哪亞說:我告訴婦人的一切事,她都當謹慎,葡萄樹所結的都不可吃,青酒,濃酒都不可喝,一切不潔之物也不可吃,凡我所吩咐的她都當盡來盡守..
嗎哪亞對耶和華的使者說:求你用我們款留你,好為你預備一隻山羊羔..
耶和華的使者對嗎哪亞說:你雖然款留我,我卻不吃你的食物,你若預備繁祭,就當獻與耶和華..
原來嗎哪亞不知道她是耶和華的使者..
嗎哪亞對耶和華的使者說:請將你的名字告訴我,到你說應驗的時候,我們很尊敬你..
耶和華的使者對她說:你何必問我的名,我名是奇妙的..
嗎哪亞將一隻山羊羔和素製到盆石上獻與耶和華..
耶和華的使者在壇上的火焰中也升上去了..
嗎哪亞和他的妻子看見,就附伏在地..
耶和華的使者不再向嗎哪亞和他的妻顯現,嗎哪亞才知道她是耶和華的使者..
嗎哪亞對他的妻子說:我們必要死,因為看見了神..
他的妻卻對他說:耶和華卻若要殺我們,必不從我們手裡收納凡祭和素祭,並不將這一切事指示我們,.
今天也不將這些話告訴我們..
頂子回庭安!.
我們今天繼續以舌詩記重新思考一下,我們今天信仰生命是一個怎樣的狀況..
有一個假設性的處境,如果有一個朋友叫你幫他家拆了一個偶像,.
很快很成功地將那個偶像拆除了,這個朋友可能是一個比較年紀大的,.
他覺得舊時擺的偶像空空如也,感覺上很不安樂,.
於是你送了一個基督是我的主的牌給他,放在舊時擺偶像的位置,.
我們就想著一切就完成了整個任務,.
不過第二次之後你再去探望這個朋友,.
你會看到這個比較上了年紀的長者,對著這個牌,.
用舊時的裝香,或者放三個杯子在這裡,.
他用的牌是基督是我家的主,.
他敬拜的,他說我是敬拜耶穌基督,.
不過用的方法是舊時拜偶像的方法,你覺得有沒有問題?.
有的,不過有時候我們會容忍的,.
他也是拜神,拜回正確的神,.
不過我們知道一直聽下去,我們知道有些不妥的,.
今天我們看的事事記某程度都是一個類似這樣的處境,.
我們今天看的這個事事就是三宣,.

$^{41}$其實在事事記裡面,技術的事事,最後一個就是參宣了,.
在參宣裡面相對其他的事事來說,他的能力是最高的,.
甚至他被揀選的,只是推到最盡的,在武福裡面已經被揀選了,.
一個沒有能力,上帝的靈充滿的一個年輕人,.
一個沒有智慧能力的一個事事,他的家庭狀況是怎樣呢?.
如果我們大概明白聖經就知道事事不是我們一個要跟從的榜樣,.
不過在他生命裡面又有某程度的從神而來的能力,.
但又不是我們跟隨的榜樣,甚至他有些生命的質素是我們批判的,.
究竟我們今天怎樣去看這個事事呢?.
在聖經裡面他所說的這個事事就不是由他,.
耶和華的靈充滿他開始說起,是由他的家庭,他的父母,.
怎樣去領受上帝對他說的話開始,.
接下來這三個星期我們會分不同的階段來讀參宣的故事,.
今天我們先看看他的父母怎樣領受上帝的吩咐,.
在參宣出世之前發生了什麼事,.
我們今天一起去讀的時候看看我們會不會明白上帝.
或者在聖經裡面告訴我們一個什麼樣的訊息,.
我們先有個祈禱,求上帝幫助我們,同心祈禱,.
親愛的伯父,求你在我們當中,願你指教我們,.
讓我們能夠明白你的心意,讓我們在我們的人生當中,.
我們不是執意偏行幾路,我們希望你調教我們,.
督促我們,讓我們行在正路當中,.
恭敬將來領受你話語的時間交託,幫助每位的弟姐妹聽得明白,.
讓宣講的講得清楚,恭敬將來的時間交託給你,.
願你賜福,獻上我們的禱告,奉基督耶穌得勝的聖名,.
我們一起看《士師記》第十三章,.
今天看的經文會比剛才我們一起聆聽的段落寬闊,.
因為篇幅也比較長,.
在第一節到第七節,其實就講到耶和華的使者,.
向著一個不育的婦人,就是參孫的媽媽,來到她顯現,.
值得留意的是在第十三章第五節,.
這個使者就特別講到這個孩子一出胎,就要歸神作拿世爾人,.
拿世爾人先解釋一下什麼意思,然後我們再下那個段落,.
拿世爾人其實在聖經裡面,是講一幫很特別的,.
一幫歸耶和華為聖的,一幫願意操練自己的一幫人,.
作為拿世爾人有兩種方式,.
一種就是按著那個成年人,或者那個人自己的意願,.
他願意守上帝一個很嚴謹的誡命,.
什麼誡命呢?首先就是清酒濃酒都不可嘗,.
或者在民數記的記載裡面,特別講到連那些葡萄樹的果子,.

$^{81}$葡萄的果子,他是完全碰都不碰的,.
不單止不喝酒,而是釀酒的葡萄都不會碰的,.
第二樣東西就是他不會剪頭髮的,包括那些鬍子他都不會剃的,.
身上的毛他都不會剃的,相對於當時的外邦剃得乾乾淨淨,.
埃及的王剃得頭很乾淨,在這裡講到身上的毛髮他是不會剃走的,.
第三樣東西就是他不會做一些不潔的行為,特別是在講埃及的屍體,.
在民數記裡面講到極端,甚至乎家人有死人,.
如果他是成為一個拿世爾人,他都成為了一個不潔,.
他都沒有辦法完成拿世爾人這樣的心願,.
做這件事全部都在聖經裡面,特別是在民數記裡面講,.
是那些百姓自願歸耶和華為聖,所以是他們個人的意願來守一個這樣的很嚴謹的誡命,.
而這個誡命大概和大祭司的規格是很近似的,.
他不是大祭司,不過是以大祭司的規矩來過他自己的生活,.
不過這種自願性質的歸耶和華為聖成為拿世爾人是有一個結束的日期的,.
他覺得日期滿了,他就可以還俗去回他自己普通人的生活當中,.
這個是第一種,第二種就是一些在母胎裡面已經被上帝選照成為拿世爾人,.
這些人其實是很少的,在聖經裡面真的記載著一出生就已經成為拿世爾人,.
參孫就是其中一個,第二個就是薩姆爾,第三個如果再嚴格來說他的新約就是斯朕約翰,.
在胎裡面媽媽懷著胎,甚至參孫現在連胎都還沒有懷,.
都已經是在講上帝有一個任命給一個將要出生的孩子,.
他們就守這些規矩,在肚子裡面的孩子就不是自己做決定,.
而是上帝命定他,上帝揀選他,讓他這一生都要遵守這樣的一種守誡命的身份,.
不過如果上帝選照的話,這個就沒有一個限期來結束,.
某程度來說是一生他都是被上帝揀選成為拿世爾人,成為了一個與世俗不同的人,.
在《創世記》第49章26節其實最早已經記載著拿世爾人這個名字,.
不過在《創世記》裡面講的拿世爾人他就不是這樣譯的,.
那句經文就說到最後雅各給他十二支派,他十二個兒子,.
包括約瑟的祝福的時候,他最後祝福約瑟就說一句這樣的說話,.
臨到那與弟兄詭別之人的頂上,與弟兄詭別之人的字眼就是拿世爾人的原文,.
在原文就用這個字,不過在《和合本》的《創世記》裡面就譯與弟兄詭別,.
就是與弟兄有別於一般的弟兄,他是分別了出來的,.
所以你明白了拿世爾是一個分別為性的一種身份,.
由《創世記》開始已經有這種記載了,去到參孫出世之前,.
上帝就透過他的使者向他的媽媽顯現告訴他,他的兒子就是一個這樣的狀況,.
在聖經裡面第十三章第一節裡面,他的起點都是用一個《士師記》裡面以色列人行惡的起點,.
十三章第一節裡面說到以色列人又行耶和華眼中看為惡的事,.
耶和華將他們交在非利士人手中四十年,.
這個是一個很悲壯的一段歷史,因為被外族欺侮之前說的是三兩年,十年,.
現在去到這裡就說的是四十年,被非利士人欺壓,然後上帝就興起《史詩》,.
在興起的時候,在第二節裡面就說,那時候有一個所拉人是屬彈族的,名叫馬羅亞,.

$^{121}$他的妻不懷孕不生育,耶和華的使者向那婦人顯現對他說,.
向來你不懷孕不生育,如今你必懷孕生一個兒子,.
所以你當謹慎,清酒濃酒也不可喝,一切不潔之物也不可吃,.
你必懷孕生一個兒子,不可用刀剃他的頭,因為這孩子一出胎就歸神作那世人,.
他必起手拯救以色列人脫離非利士人的手..
在這句話是耶和華的使者對那個婦人說的,.
你值得留意整句話首先是說他要成為那世人,.
然後有一個目標是做什麼?.
就是他要起手拯救以色列人脫離非利士人,.
脫離非利士人的手,他的起點是說因為非利士人欺壓四十年,.
所以這個孩子出世被分別為聖,然後他的目的就是要脫離非利士人的手,.
整個的意思是這麼完整的,.
當這個婦人聽到這個使者說完這番話之後,.
這個婦人有什麼反應呢?.
接著第六節就立刻說,這個婦人立刻回去找老公,.
婦人就回去對丈夫說,有一個神人到我面前來,.
他的相貌如同神使者的相貌,甚是可畏,.
我沒有問他從哪裡來,他也沒有將他的名告訴我,.
卻對我說,你要懷孕生一個兒子,.
所以清酒濃酒都不可喝,一切不潔之物也不可吃,.
因為這孩子從出胎一直到死,必歸神作拿世以人..
在這裡聽下去,驟眼聽下去都很興奮,.
不孕的婦人,有個神人應許她,有個嬰兒將要出世,.
這個婦人她的著重點,她在講福術這件事,她有兩個重點的,.
第一,她講的重點是甚麼呢?她說跟她講話的神人是怎樣的?.
從他的樣子開始講起,他的相貌,有神使者的相貌,.
樣子都幾像樣,幾好樣,甚是可畏,.
即是從外貌來到品評跟她講話的神使者,耶和華的使者,.
第二部分就是講到這個孩子將要出世了,.
基本上有八成都講了,要歸作拿世以人,.
清酒濃酒不可嘗,不潔之物也不可吃,.
這個孩子一出胎直至死都歸作拿世以人,OK的,齊了的,.
不過最後一句沒有講,沒有講甚麼呢?.
沒有講這個孩子起手將要作以色列人脫離非利人手的拯救者,.
她沒有講這句話,所以這個婦人心裡面,.
第一個如果你代入進去,她在想的就是,.
我將要不孕的將要懷孕了,我將要有一個兒子來承繼家業,.
我當代的婦人看不孕是一個羞辱,從此她的羞辱就沒有了,.
而且這個神人講話是可信的,因為她的樣貌很可敬,很可畏,.
從她的外貌特徵來判斷這句話是不是真的,.

$^{161}$第二個重點就是,不過這個兒子將來是要歸耶和華為姓的,.
所以在一個很嚴格的規條之下,她一生都要這樣守的,.
不單止這個兒子開始守,這個婦人其實這個使者都吩咐她,.
她自己都要守的,媽媽為了生孩子,戒口在中國人的文化裡面是很平常的,.
大概這個婦人為了要令這個兒子順利誕生,.
她要戒口都在所不辭的,她的著重點就是,.
接下來這個兒子將要出生了,很興奮地跟丈夫說,.
而且這個是一生要守戒命的一個兒子,.
不過沒有想過她也沒有將重點放在成為拯救以色列人脫離非利人的守的拯救者,.
當這個婦人對著她丈夫馬羅亞說了這番話之後,.
在第八節,馬羅亞聽到他的太太這樣說,.
第八節馬羅亞就祈求耶和華說,主啊求你再差遣那個神人在我們當中,.
這類來好指教我們怎樣待這將要生的孩子..
神應問馬羅亞的話,婦人正坐在田間的時候,.
神的使者又坐在她那裡,她丈夫馬羅亞卻沒有與她同在一處..
婦人急忙跑去告訴她丈夫說,.
那一天到我面前來的人又向我顯現,.
馬羅亞起來跟隨他的妻子來到那一天面前對她說,.
與這婦人說話的就是你媽,她說是我..
我們由八節到第十一節看到,丈夫其實他要查證,.
在他老婆說完這番話之後,.
馬羅亞某程度上帶著一個很勁竭的態度來祈求耶和華上帝,.
祈求上帝再一次將這個神人顯現給他知道,.
因為他老婆說一點不說一點,.
作為丈夫有時候聽到老婆說話都是說一點不說一點,.
我很怕我得罪我老婆,所以我不敢說得太多,.
不過很多時候女士是情感主導的,將最興奮的先說了,.
其他的細節他不太理會,.
丈夫聽到這番話之後,於是他求上帝再一次讓這個神人顯現在他面前,.
好查證究竟這個講的是一個怎樣的訊息,.
他特別注意的就是他要很想神人指教他,.
告訴他怎樣去對待這個將要生出來的孩子,.
不過如果你再去回想一下,其實這個神人之前對著婦人已經說了,.
要將他歸耶和華為聖,.
所以要讓他除了戒命之外,.
如果真是嚴格的那西爾人,如果你比對薩姆爾,.
他要放在殿裡面來學習事奉耶和華上帝,.
不過在他的眼中裡面,或者在他婦人眼中裡面,.
他就覺得要禁戒這些條例,.
但是怎樣敬拜上帝,怎樣去實踐,.

$^{201}$其實某程度上那個使者已經跟這個婦人說了,.
不過馬羅亞似乎不是很收到這個訊息,.
所以他祈求上帝再一次差派這個使者去當中,.
於是上帝就應運馬羅亞的說,.
當這個婦人坐在這裡的時候,她老公不在她身邊,.
顯現的時候,老婆就馬上找她老公來,.
就去到這個使者面前,就作出這種說話,.
不過這個老公第一樣東西見到這個使者,.
他第一句問她的是什麼呢?.
他說那天在我老婆或者在這個婦人面前顯現的那個是不是你呀?.
這句說話如果你留意,其實馬羅亞那種心態,.
就好像面對著今天我們那些,面對一些驚詐騙電話的那種感覺,.
是不是真的?是不是你跟我老婆說這番話?.
是不是你打電話給我的?.
要某程度上都帶著謹慎,小心翼翼的態度,.
來去問那個當時的人的,這個使者很清晰,.
不過和合本就說是示我,感覺上就讀出來就好像很被動,.
原文他做的那句話就是我示,我就是,.
那個我強調我多一點,有些霸氣說多一點這句話,.
而不是說這個是我,不是這樣的,.
是一個很肯定,告訴他他是和他太太見面的使者,.
然後經文裡面第十二節就說了,.
願你的話應驗,我們當怎樣待這個孩子,.
他後來會怎樣呢?老公問的問題反而對板一點,.
因為那個神人或者那個使者之前說,.
他要歸耶和華為聖,然後成為以色列的拯救者,.
脫離非利士人的手,整個句式是這麼整齊的,.
不過當馬萊亞這樣問那個使者的時候,.
那個使者再重複的就是說,我告訴婦人的一切時,.
她都當謹慎,葡萄樹所結的都不可吃,.
清酒濃酒都不可喝,一切不潔之物也不可吃,.
凡我所吩咐的她都要遵守,某程度其實他就強調了,.
拿西爾人前面的守戒,然後後面所說的這個婦人,.
吩咐她所有事都要遵守,這樣就完了,就全部回答了,.
馬萊亞好像收到,又好像收不到,.
他覺得他老婆說的話就等於事實的全部,.
因為守戒名,也沒有想過自己的兒子將要成為一個士師,.
被分別出來拿西爾人的士師來拯救以色列人,.
然後經文馬萊亞就說了這麼偉大的說話,.
這麼好的一個消息,於是馬萊亞就對那個使者說,.

$^{241}$不如這樣吧,讓我款留你,款留你這個字稍後再解釋一下什麼意思,.
好為你預備一隻山羊羔,他知道神人給了他一句這樣的說話之後,.
他就要款待他,款待他用什麼來款待他呢?.
就是用一隻山羊羔,如果記得剛才說的這個婦人說有一個神人,.
那個神人也是帶著憑外貌來認人的那種感覺,.
然後到了這裡就說,老公聽到這樣的消息,.
我現在tong 一隻雞來款待你,類似中國人這樣,.
一般來說你對朋友,有個朋友做了一些恩惠給你,.
或者不是tong 雞,你今天我就tong 一隻燒豬來款留你,.
其實在當代裡面用山羊羔來款留,.
就是對於當時的神或者當時的偶像敬拜,.
對於那些半人半神的供奉就用這麼高規格的方式,.
來接待這個超越了他想像的賓客,.
某程度馬羅亞當這個人是神,神人之間,.
但是用最高規格,他的態度是怎樣的?.
是敬虔的,我依然在說,他是一個敬虔的人,.
他是一個很小心翼翼的人,.
不過他用的方法是用當時民間信仰的方式,.
來接待耶和華的使者,.
當耶和華的使者知道他這樣做的時候,第十六節,.
耶和華的使者對馬羅亞說,.
你雖然款留我,我卻不吃你的食物,.
你若遇被凡濟就當獻與耶和華,.
經文再繼續加一個註腳,原來馬羅亞不知道他是耶和華的使者,.
整句經文其實已經告訴你了,.
使者說我不會吃你的食物,.
你用這個方式來款留我,我是不會接納的,.
應該用凡濟的方式來獻給耶和華上帝,.
這個使者其實某程度是調教馬羅亞,.
他對於擺錯了敬虔的心態,.
一個敬虔的人在使者引導之下,.
他就做適當的行為,.
經文裡面特別說到馬羅亞不知道他是耶和華的使者,.
這句話很多時候有很多不同的解法,.
究竟他是基於敬拜神,一個有神的概念之下,.
不知道上帝差遣一個使者在他當中,.
還是他不知道耶和華的使者擺錯位,.
他用著二邦的方法來款留耶和華的使者,.
他不知道他是對錯位的,.
我覺得偏向後者會比較適合上文下理,.

$^{281}$馬羅亞這一刻他只知道很敬虔,.
有人應許他要生一個兒子,.
清酒濃酒不可嘗,.
即使他要成為那西爾人在耶和華上帝那條線,.
他還是用了二邦的敬拜方式來對待他,.
所以這個使者拒絕馬羅亞的款留,.
馬羅亞繼續問,.
你把名字告訴我,.
到這句話應驗的時候,.
我很尊敬你,.
很尊敬你又是還神,.
做回那件事,.
你告訴我,神也有個名字,.
或者你有個名字告訴我,.
雖然好像調教了他,.
但他好像也不是很明白,.
於是這個神人對於馬羅亞的問題,.
他怎樣回答呢?.
第十八節,耶和華的使者對他說,.
你何必問我的名字呢?.
我的名字是奇妙的,.
我們一解經以為使者的名字叫奇妙,.
其實不是這個意思,.
整句話就是說,.
你不用問我什麼名字,.
說你也不明白這個意思,.
就是我的名字叫奇妙,.
其實是超越了你的理解能力,.
所以你不用問我,.
說給你聽也是白費勁,.
你不明白的,.
所以你不要問,.
在不問和這些規限當中,.
馬羅亞好像有點懵,.
聽著他老婆說一知半解的話,.
他就好像全盤接收,.
然後這個使者因為全部告訴他老婆,.
所以他就好像懵懵懂懂,.
蒙在鼓裡用了二幫的方法,.
來去款留這個神人,.

$^{321}$後來所有事情都調教對了,.
十九節,馬羅亞將一隻山羊羔和素祭放在盤石上獻於耶和華,.
使者行奇妙的事,.
馬羅亞和他的妻觀寒見火焰從壇上往上升,.
耶和華的使者在壇上的火焰中也升上去,.
馬羅亞和他的妻漢見,.
就俯伏於地,.
在這裡真的做對了事情,.
做對了一個凡祭,.
見到這個凡祭和素祭在盤石上獻於耶和華上帝,.
這個使者就做了一些很奇妙的事,.
用了奇妙的名字,.
說你都不明白的,.
同樣做了一些說你都不明白的事,.
就是當獻祭的時候,.
那個祭身被火燒上去,.
而這個使者也隨著這些火焰往上升到天上,.
然後就消失了,.
去到這樣的狀況之下,.
馬羅亞見到一個這樣的處境,.
他們就俯伏在地,.
然後在第21節,.
耶和華的使者不再向馬羅亞和他的妻顯言,.
馬羅亞才知道她是耶和華的使者,.
那時候才知道她是使者,.
馬羅亞對他的妻說,.
我們必要死,.
因為我們見到神了,.
他的妻卻對他說,.
耶和華如果要殺我們,.
必不從我們的手裡面收納你的凡祭和素祭,.
並不將這一切的事指示我們,.
今天也不將這些話告訴我們,.
他老婆很淡定,.
但馬羅亞就很害怕,.
在那一刻才害怕,.
在那一刻相會,.
遇見才如夢初醒,.
他當時向耶和華祈求的說話已經應驗,.
而眼前這個人,.

$^{361}$也是眼前消失的使者,.
是來自耶和華上帝,.
之前他是敬畏上帝的,.
很有態度,.
很誠懇的,.
不過是用義邦的方式,.
來敬拜上帝,.
不過到了後面,.
當看到真的奇妙超乎他所想像的事的時候,.
兩夫妻才俯伏在地,.
才將最大的恐懼,.
主權,.
伏在上帝的面前,.
來敬拜上帝,.
這個段落其實很值得我們去深思,.
今天我們的信仰觀念,.
其實可能很多弟兄姊妹也會有,.
我們信了主,.
我們成為第一代的基督徒,.
但是在我們的信仰當中,.
有時候我們也會將一些,.
昔日我們拜偶像的民間習俗,.
參雜了進去,.
放進去我們自己的信仰觀念裡面,.
中國人有五花八門的,.
有不同的方式,.
來敬拜上帝,.
敬拜當時民間覺得是神的神明,.
不過這些心態,.
當他進入基督,.
或者成為一個基督徒的時候,.
他也會有意無意之間,.
將這些東西呈現出來,.
我記得不是嘲笑,.
而是我覺得有時候對這些想法很納悶,.
有些弟兄姊妹家裡新屋入伙,.
然後她就說牧師,.
你來為我新屋入伙,.
通常就說,.
為什麼呢?.

$^{401}$其實你們自己祈禱也可以,.
我初時請我去吃一頓飯,.
她說不是,.
你來祈個靚肚,.
我們才入伙,.
這樣會好一點,.
我說其實你們自己祈禱也可以,.
我其實不是不樂意去做這件事,.
很快樂的,.
新屋入伙,.
為弟兄姊妹來祈禱,.
新屋入伙,.
不過如果我心裡面知道,.
她是無法脫離舊時的民間信仰,.
我就會覺得很納悶,.
究竟我幫不幫她好呢?.
如果我幫她,.
我就好像助長了她繼續帶著這種觀念,.
來敬拜上帝,.
但是不幫她,.
我就覺得她沒有愛心,.
為什麼我會這樣做呢?.
令到她的心很單純,.
基於敬虔的動機,.
她希望這間屋歸炎和華為勝,.
在這些位置裡,.
有時候真的很難去拿捏的,.
不過有時候我們都要問自己,.
我們在這些位置裡,.
我們會不會都沾染了一些民間習俗,.
再講一個例子,.
這是一個老基督徒告訴我的,.
他說,.
因為他以前經歷過在他自己的鄉下裡,.
有一些習俗,.
或者有一些曾經受害過,.
話說他們鄉下裡曾經有些人,.
拿了家裡人的頭髮,.
枕頭總會有幾條的,.
拿了去,.

$^{441}$然後就下降頭,.
在那些南洋的位置,.
然後家裡人就受害了,.
有病了,.
然後就有一些印傭來到,.
離開的時候,.
那個老基督徒就說,.
千萬不要給印傭繞你的枕頭,.
他拿了你幾條頭髮,.
拿回印來下降頭你就知道了,.
通常不是happy ending,.
炒印傭那些,.
他就說不要給他碰枕頭的東西,.
所有東西,頭髮什麼都不准給他碰,.
小朋友枕頭底下的都不要給他抹,.
就這樣給他走就行了,.
我說,你基督徒有時候不用這麼害怕,.
他說,信我,.
我以前經歷過這些事情,.
有時候會有一些很沒信心的行為,.
去到這些位置我就會說,.
你放心,.
我們被上帝保護的人,.
我們心裡面有聖靈保護,.
我們是不需要害怕的,.
反而我們要更加有信心走這條信仰的路,.
不要再將舊時的信仰觀念放進去,.
這種敬拜,.
有另外一個基督徒就說,.
七月夠漂亮,.
適合他那幾天,.
走路不要靠近牆,.
貼著牆走,.
我說,為什麼?.
你不要說,小朋友不要問這麼多,.
這樣說,.
我知道他問什麼,說什麼,.
不過有時候都會帶著這樣的觀念,.
我們在經歷我們自己的信仰生活,.
我們表面是一個基督徒,.

$^{481}$但是暗處裡面其實我們都帶著很多的恐懼,.
迷信,.
不知道怎麼做,.
放進去我們的基督教信仰裡面,.
我們依然是一個敬虔的人,.
我們都是以上帝為主,.
我們很尊重上帝的人,.
不過我們沒有辦法將這種生命徹底改變的時候,.
我們的信仰生命都是成王成孔,.
我們和拜偶像其實只不過是那個標誌改變了,.
但是我們的心態是一樣的,.
求主幫助我們,.
在這裡馬羅亞就成為了一個,.
或者她的妻子,.
是敬虔的,.
是守戒命的一個人,.
但是心底裡面他們經歷不到,.
從神而來,.
敬拜耶和華真理的釋放,.
這種心態你不要小看,.
你不要以為都是敬虔就可以,.
或者我們不去做一些調教,.
如果這個心態不去處理的時候,.
很多時候這些信仰就會有後果來承接出現,.
在經文裡面,她的兒子參孫,.
就將這種信仰觀念,.
在這樣的家庭狀況下表露無遺,.
我們再繼續看,.
我們看在經文裡面第十四章,.
或者我們再看多一點,.
我們看回第二十四節那裡,.
十三章二十四節,.
婦人見到這件事,.
或者見到整個狀況,.
這個兒子出生了,.
後來婦人生了一個兒子,.
給他起名叫參孫,.
孩子長大,.
耶和華賜福於他,.
在瑪哈尼旦,.

$^{521}$就是所拉和以實圖中間,.
耶和華的靈才感動他,.
很簡單的一句話,.
說的那個兒子叫什麼名字,.
這裡裡面特別有些記載的,.
這個兒子的名字叫參孫,.
參孫這個名字的意思是什麼意思呢,.
在原文裡面,.
Samson這個字其實是太陽之子的意思,.
在聖經裡面,.
基本上,.
這些如果上帝很清晰,.
感召那個父母,.
不育的婦人生兒子,.
特別是比對,.
文本互涉,.
我們叫比對其他經文,.
用其他經文來講故事,.
或者有很多並排的那些說法,.
撒母耳由一個不孕的哈拿婦人,.
來生兒子,.
他生的兒子,.
他會按著他這些經歷,.
對神的敬拜,.
他會改兒子是和上帝有關的,.
所以撒母耳的名稱,.
就叫做上帝聆聽了,.
上帝垂聽了,.
Samuel,.
就是在說上帝垂聽了哈拿婦人的祈禱,.
所以生了這個兒子,.
而太陽之子,.
其實不是按著他的經歷,.
而命名他的兒子,.
為什麼要用太陽之子呢?.
當時的義邦是一個農業社會,.
敬拜太陽就是他們最重要敬拜的神,.
敬拜太陽令他們風調雨順,.
所以他們會視太陽成為他們的神,.
太陽神有時候看考古都知道,.

$^{561}$在這裡他將他的兒子視為太陽之子,.
不是視為耶和華之子,.
所以整個過程裡面,.
你會看到他們父母的心是帶著敬虔,.
但骨子裡他也是帶著義邦的敬拜,.
來養育自己的孩子,.
很諷刺的地方是馬來西亞問神人,.
這個兒子要怎樣幫他成長,.
他們做了一些表面,.
就是要幫他清酒濃酒不可嘗,.
不可靠近死屍,.
不用剃刀剃身體上的毛,.
這些外在的事情都做了,.
但骨子裡他是一個愛邦敬拜神的方式,.
來教育他的兒子,.
在這樣的家庭教育下,.
出現了一個什麼樣的兒子呢?.
《蝕師記》第十四章一至三節,.
參孫下到亭那,.
在那裡看見一個女子,.
是非利士人的女子,.
參孫上來稟告他的父母說,.
我在亭那看見一個女子,.
是非利士人的女子,.
願你們給我娶來為妻..
他父母說,.
在你弟兄的女兒中,.
或在本國的民中,.
豈沒有一個女子,.
何知你未受國禮的非利士人中娶妻呢?.
參孫對他父親說,.
願你給我娶那女子,.
因我喜悅她..
在一個很小的交代,.
說的是參孫從小到大,.
很嚴格的聖殿規例,.
拿世以人的規條長大,.
到大了,.
娶妻的年紀了,.
看見外邦的女子,.

$^{601}$就喜歡,.
他父親對他父親說,.
為何要找一些未受國禮,.
非利士人外邦女子為妻,.
找自己本族人,.
參孫的答案是甚麼呢?.
我喜歡..
今天,我不知道你會否覺得這句說話很有趣,.
今天這一代經常聽,.
我喜歡,.
我們可能在基督教家庭長大的孩子,.
但有時偏行幾路,.
到某一刻,父母知道不對辦,.
但在一個只有規範,.
但沒有內涵的信仰生命裡,.
我們會孕育出新一代,.
我行我素的兒女,.
繼續他們自己的信仰觀念,.
我喜歡,.
成為了當代人的寫照,.
或成為了我們這一代人的寫照,.
在《士司記》裡,.
他不是純粹描述一個人,.
有多不理想,.
在《士司記》裡,.
他用一把聖靈的寶劍否開了,.
這個人的出世,.
他家裡某程度裡裡都製造了一些因子在我們當中,.
有時我們讀書時,.
我們自己都會痛心,.
或我們自己都謹慎,.
我們作為一個基督徒父母,.
自問自己是一個敬虔的人,.
我們對上帝都有很多尊敬,.
放到很高的位置,.
有很多規條禁止我們的兒女來做,.
不過當他們長大的時候,.
他們經常偏行幾路,.
問題在於,.
可能不是純粹下一代的問題,.

$^{641}$而可能在我們的信仰傳承裡,.
我們出了事,.
我們只是將表面的事告訴我們的孩子,.
在我們的孩子面前,.
我們可能表面上成為一個敬虔的人,.
但骨子裡我們依然帶著很多異教,.
又或者我們心底裡都偏離了上帝,.
我們的孩子可能看在眼中,.
以致他們成為了反叛的一代,.
成為了一個我行我素,.
目空上帝的一代..
在聖經裡面,.
舍斯記最後的一句說話,.
第25章,.
他經文裡是這樣說的,.
我讀給大家聽,.
在21章第25節,.
舍斯記第21章25節,.
那時以色列人中沒有王,.
各人任意而行..
這句說話說的是什麼呢?.
任意而行的意思不是說喜歡做什麼就做什麼,.
以色列沒有王,.
各人任意而行,.
是各人行自己覺得認為對的事,.
自己覺得好的,.
自己覺得對的就去做,.
各人有各人不同版本的敬虔,.
在他當代的世代來活現,.
到最後就製造了整個時代,.
一路越走下坡的一種狀況..
弟子們,.
今天我們敬拜上帝,.
我們都要闊心自問,.
究竟我們會不會只有外在的敬虔,.
我們心底裡是一個怎樣的人呢?.
未必一定是一個父母,.
但在我們的屬靈生命裡,.
究竟我們會不會只有一個框框,.
有大量的規條不可做,.

$^{681}$不可以做這些事,.
是出於敬虔的動機,.
但我們心裡不會細想,.
我們不會真的進入真理當中,.
以致我們只是成為了一個敬拜神的異邦基督徒,.
我們只是敬拜神,.
但我們心裡都是敬拜的偶像,.
事實既提醒我們,.
我們這個世代不要再這樣繼續走下去,.
我們生命的敬虔,.
我們真的要從心底裡明白上帝的話語,.
知道上帝的心意,.
用上帝的心意,.
用上帝的曉喻方式來敬拜神,.
那個敬虔才是徹頭徹尾的敬虔,.
唯有這樣的敬虔才可以帶來真正有效的傳承,.
求主繼續幫助我們,.
繼續對我們說話..
\newpage



\section{士師記 14:1-9}
\label{sec:Up7U5HlOZfo}
\textbf{2023年9月10日崇拜直播｜陳冠成牧師｜「家醜說出口」系列:欲望阻不了｜士師記14\_1-9}
\newline
\newline
連結: \href{https://youtube.com/watch?v=Up7U5HlOZfo}{\texttt{https://youtube.com/watch?v=Up7U5HlOZfo}} ~~~~ 語音日期: 2023-09-10
\newline
\newline
\hyperref[sec:sDc2iE_sQ_g]{\small{< < < PREV SERMON < < <}}
~
\hyperref[sec:index_chronic]{\small{[返順時目]}}
~
\hyperref[sec:rllwPQsYaK0]{\small{> > > NEXT SERMON > > >}}
\newline
\newline
士師記 14:1-9
\newline
\begin{longtable}{cl}
\hline
\hline
章節 & 經文 (和合本修訂版)\\
\hline
14:1 & \begin{tabularx}{0.7\textwidth}{X} 參孫下到亭拿，在亭拿看見一個女子，是非利士人的女兒。 \end{tabularx} \\ \\ \relax
14:2 & \begin{tabularx}{0.7\textwidth}{X} 他上來告訴他父母說：「我在亭拿看見一個女子，是非利士人的女兒，現在請你們給我娶她為妻。」 \end{tabularx} \\ \\ \relax
14:3 & \begin{tabularx}{0.7\textwidth}{X} 他父母對他說：「在你弟兄的女兒中，或在本族所有的人中，難道沒有女子嗎？你何必在未受割禮的非利士人中去娶妻呢？」參孫對他父親說：「請你給我娶那女子，因為我喜歡她。」 \end{tabularx} \\ \\ \relax
14:4 & \begin{tabularx}{0.7\textwidth}{X} 他的父母並不知道這事是出於耶和華，因為他在找機會攻擊非利士人。那時，非利士人轄制以色列人。 \end{tabularx} \\ \\ \relax
14:5 & \begin{tabularx}{0.7\textwidth}{X} 參孫跟他父母下亭拿去，他們到了亭拿的葡萄園。看哪，有一隻少壯獅子對著他吼叫。 \end{tabularx} \\ \\ \relax
14:6 & \begin{tabularx}{0.7\textwidth}{X} 耶和華的靈大大感動參孫，他就手無寸鐵撕裂獅子，如撕裂小山羊一樣。他做這事，並沒有告訴他的父母親。 \end{tabularx} \\ \\ \relax
14:7 & \begin{tabularx}{0.7\textwidth}{X} 參孫下去跟那女子說話，看著就喜歡她。 \end{tabularx} \\ \\ \relax
14:8 & \begin{tabularx}{0.7\textwidth}{X} 過了些日子，他回來要娶那女子，繞道去看獅子的殘骸，看哪，有一群蜜蜂在獅子的屍體內，也有蜜在裡面。 \end{tabularx} \\ \\ \relax
14:9 & \begin{tabularx}{0.7\textwidth}{X} 他就取了蜜，放在手掌上，邊走邊吃。他到了父母那裡，給他們蜜，他們也吃了。但他沒有告訴他們，這蜜是從獅子的屍體內取來的。 \end{tabularx} \\ \\
[1ex]
\hline
\hline
\end{longtable}
$^{1}$十四章一至九節.
舊約時詩記第十四章一至九節.
在舊約聖經三百一十八頁.
舊約聖經三百一十八頁.
請聽我讀出.
第十四章一節.
參孫下到亭拿.
在那裡看見一個女子.
是非利士人的女兒.
參孫上來稟告他父母說.
我在亭拿看見一個女子.
是非利士人的女兒.
願你們給我娶來為妻.
他父母說.
在你弟兄的女兒中.
或在本國的民中.
豈沒有一個女子.
何至你去.
在未受國禮的非利士人中娶妻呢?.
參孫對他父親說.
願你給我娶那女子.
因我喜悅她.
他的父母卻不知道.
這是出於耶和華.
因為他找機會攻擊非利士人.
那時非利士人.
核製以色列人.
參孫跟他父母下亭拿去.
到了亭拿的葡萄園.
見有一隻笑臟獅子.
向他嘲笑.
耶和華的靈大大感動參孫.
他雖然身無氣械.
卻將獅子撕裂.
如同撕裂山羊羔一樣.
他行者是並沒有告訴父母.
參孫下去與女子說話.
就喜悅她.
過了些日子.
再下去要娶那女子.

$^{41}$轉向道旁要看死屍.
見有一群瘋子.
黃脈在死屍之內.
就用手取脈.
扯起扯走.
到了父母那裡.
給他父母.
他們也吃了.
只善沒有告訴者物.
是從死屍之內取來的.
弟兄姊妹平安.
上星期陳明全牧師已經看了.
參孫的夢照和出生.
今天我們繼續看下去.
話說回來.
假如我們對參孫這個人物.
不太認識的話.
假如我們沒有看到今天的經文.
不知道參孫長大後.
到底會發生什麼事.
其實我們相信他和.
薩姆爾或斯贈約翰.
都會成為一番作為的祭司.
因為他確實有潛質和條件.
畢竟他在母胎成形之前.
上帝已經揀選了他.
成為一個拿洗人終身的.
來給上帝使用.
所以在我們的期望裡.
參孫將會在上帝的計劃裡.
大有一番作為.
不過很可惜.
今天我們看到的經文.
你發覺參孫有一個致命的缺點.
是什麼?.
他很任我行.
很惡行惡掃.
你看到他的行為模式.
基本上是受到自己的欲念支配.
他看到想要的東西.

$^{81}$就會伸手去拿.
想到想做的事情.
就會馬上去做.
你看到他為了滿足得到的事物.
他可以超越很多的規範.
甚至可以無視上帝的吩咐.
當然也不會理會身邊的人怎麼看.
或者會否介懷自己的行為.
會連累到其他人.
就好像剛剛樹錦執事.
和我們所讀的經文.
你看到這裡.
經文記載參孫他下到庭拿的地方.
經文說他看見.
或者說看中了.
一個菲律賓人的女兒.
於是就回去稟報他的父母.
跟他們說.
「爸媽,我在庭拿見到一個女子」.
「是菲律賓人的女子」.
「願你們幫我娶她為妻」.
這句說話很簡單.
不過卻暴露了.
參孫也挺受他眼目情慾的支配.
只是看見而已.
作者刻意說只是看見而已.
其實還沒聊天.
讀到第七節才開始聊天.
只是看見那個女子.
連聊天也沒聊.
就已經.
很大機會被他的外表.
或者說看太吸引.
於是就立即有行動.
不需要通過我們今天說的.
拍拖,交往.
不需要認識對方的底蘊背景.
由我很喜歡這個女子.
直接跳到怎樣.
「我要跟她結婚」.

$^{121}$你做爸爸媽媽怕不怕?.
如果你的子女有一天跟你說.
「我今天在街上看見一個女子」.
「我覺得她很漂亮」.
「我要跟她結婚」.
說真的.
不是拍拖.
不是追求她.
是要跟她結婚.
你發覺參孫比我們今天說的閃婚更厲害.
跳了很多步驟.
一跳就跳到差不多終極.
所以我們與其說參孫.
她在談戀愛.
不如說她更像甚麼.
像不像動物的狩獵.
Hunting.
看見獵物.
認定是牠.
用最快的速度全速撲過去.
基本上就像Hunting.
純粹是隨著自己的慾望.
做一個即時的反應.
不會很理性地事先計劃.
應該怎樣做.
這樣做好不好.
或者我這樣做的時候.
會否有甚麼後果要承擔.
總之.
我喜歡就是最大.
我想要的.
就一定要得到手.
你看到其實.
她有一個甚麼身份.
拿世人.
這個與拿世人應該擁有一種.
小心謹慎.
自我約束的.
一種行為.
一種態度.

$^{161}$其實基本上是反轉的.
不過話說回頭.
到底我們說參孫這種.
我們說如此原始.
如此衝動的個性.
我一看中就馬上要得到.
這種衝動.
其實是怎樣形成的呢?.
怎樣形成的呢?.
值得我們想一想.
是否天生的呢?.
當然經文她沒有直接交代.
不過我們不難看出.
她父母有沒有幫忙.
她父母有幫忙.
練成參孫這樣的.
參孫這種任意行為的表現.
她父母其實責無旁貸.
為何這樣說呢?.
你留意她父母.
對參孫這個請求的回應是甚麼呢?.
在你的弟兄的女兒中.
或者在本國的民中.
豈沒有一個女子.
為何你去在未受國禮的.
非理士人中娶妻呢?.
有沒有見到耶和華三個字?.
有沒有見到上帝的名字出現?.
沒有.
參孫的父母只是基於當時.
一些文化,一些習俗.
或者說種族的觀念.
不應該和外族通婚.
應該和自己人結親.
那才怎樣呢?.
門當戶對.
他們嘗試想用這個理由.
來反對參孫的想法.
但卻不訴諸於一個最高的權威.
應該是甚麼呢?.

$^{201}$應該是用上帝的旨意.
嚴肅制止自己的兒子的行為.
他們應該和參孫說.
兒子.
上帝吩咐過我們.
不可以和家男人通婚.
否則他們必定會引誘我們.
離棄上帝事奉別臣.
《五經》裡面說了.
又或者更重要的是.
其實上帝.
第一,你這樣做一定會得罪上帝.
你還記得.
在你未出生之前.
上帝已經怎樣?.
上帝已經說明了.
你要分別為聖.
是要用你來拯救以色列人.
脫離非利士人的手.
換句話來說.
其實我們一定不可以.
和非利士人走得那麼近.
現在你還說和他交往,結婚.
這件事不可能.
正路應該這樣勸參孫.
不過他們只是想單憑一個文化.
或者種族觀念這麼薄弱的理據.
來勸退那個血氣方剛.
對非利士女子充滿慾念的參孫.
其實當然一定不奏效.
所以你見到參孫只是很簡單一句.
你幫我娶她為妻吧.
因為我喜歡.
因為我喜悅.
直譯這句說話就是.
因為她在我眼中是看為正的.
其實參孫就是代表.
當時以色列人對上帝的態度.
覺得正的事.
不理三七二十一.

$^{241}$不理上帝.
先做了.
不需要理後果.
你看到.
參孫就是這樣.
總之我喜歡.
我就要一句簡單的說話.
就終結了他和父母這場對話.
你看到他的父母.
也沒有繼續去堅持.
反而是怎樣?.
你看到之後的經文.
他們是默許.
甚至配合這頭的親事.
當然.
我個人很相信.
他的父母應該是愛惜參孫.
因為照道理.
一個難得.
本來不育的家庭.
現在竟然沒有上帝眷顧.
賜給他們一顆兒子.
沒理由不愛惜.
但很可惜.
他們大概沒有將上帝.
賜予給參孫.
這麼尊貴的身份和使命.
從小就灌注在參孫身上.
讓他明白到.
其實忠義並不是最大.
人是應該接受上帝的監管.
如果參孫的父母.
有將上帝的吩咐.
落印在參孫的生命裡.
我相信參孫應該不會停留.
甚至見到菲律賓人.
他也會掉頭走.
不要說和他們交往.
甚至結婚.
同樣地,我們也愛惜身邊的人.

$^{281}$你的親人,子女,配偶.
我們也想將最好的給他們.
來展示我們對他們的愛.
我們有沒有將上帝的話語.
從聖經以來的價值觀念.
做人態度.
再重在他們身上呢?.
參孫的家庭其實.
很值得我們去思考這方面的事情.
讓我們看到.
我們有沒有用上帝的話語.
引導我們所愛的人.
行在上帝的旨意當中.
又或者假如他們真的.
走錯路走迷了.
遠離了上帝的話.
其實我們有沒有勇氣站出來?.
夠不夠膽在他們面前.
冒著可能破壞關係的風險.
仍然要按上帝的吩咐.
按聖經的教導.
指出對方一些生命上的盲點.
設法去勸他們回頭呢?.
很值得我們去思考.
特別作為父母.
可能在在都有些父母.
可能當你的小朋友.
開始升中的時候.
你都要為他預備一部智能手機.
我兒子是這個情況.
我兒子升中的時候.
要用智能手機.
以前小學還是用2G.
上不到網.
基本上只可以打電話.
但是上到中學沒辦法.
他要開WhatsApp群組.
來聯絡做項目.
你知道一有智能手機就會怎樣?.
就會進入網絡世界.

$^{321}$那裡進去之後.
回不出來.
每天做完功課就會怎樣?.
就會拿手機日看.
你做父母自然怎樣?.
你自然想制止他.
停止他.
在那些時候.
當他覺得我的自由被剝奪.
我的享樂被剝奪的時候.
受到規範的時候自然會反彈.
那怎樣?.
要不要站硬?.
當然要站硬.
但站硬又不可以破壞關係.
那怎樣做?.
唯有溫柔而堅定.
每一次都要停.
要停.
不過不是純粹為了停而停.
不是純粹因為你玩了很久.
所以要停.
而是因為要保護你的眼睛.
不想像爸爸一樣.
戴眼鏡近視的話.
就少一點.
看屏幕.
爸爸以前沒人看.
不過你現在有.
你要像媽媽一樣.
保護眼睛.
那就不用戴眼鏡.
他開始明白多一點.
不過依然很不滿.
那怎樣?.
唯有再說一件事.
其實你覺得.
爸爸媽媽很喜歡專程跟你對著幹嗎?.
不會的.
我們那麼疼你.

$^{361}$我們都想省一口氣.
但為何仍然要冒著破壞風險的關係.
來制止你呢?.
是因為愛你.
因為想你好.
想保護你.
不是要攔阻你.
他開始明白多一點.
然後我們再說.
更重要的是.
你的身體是上帝的.
你要保護你的身體.
來事奉上帝.
他似明非明.
但他知道.
是出於愛的緣故.
他願意選擇放下.
當然這個過程不會一下就搞定.
這個角力是經常出現的.
很考驗父母的耐性.
但你不做的話.
錯過了那個時機又怎樣?.
他將來就回不了頭.
所以弟兄姊妹.
參孫的父母.
雖然有一個很好的背景.
有上帝給他們一個獨生兒子.
並且上帝呼召他們成為那四人.
但他們的父母卻沒有運用.
屬靈的權柄教導參孫.
來堅持拒絕.
他們的兒子一些錯誤的要求.
其實很值得我們去想.
到底是愛他.
還是害死了他呢?.
不過.
回過頭來.
其實.
上帝既然呼召參孫終身成為那四人.
賦予他一個這麼聖潔尊貴的身份和使命.

$^{401}$很應該怎樣?.
亦為他預備一個.
在信仰上,屬靈上.
匹配的家庭.
應該是這樣.
好像薩姆爾,施浸約翰.
有那四人.
但亦有經乾的父母.
但很奇怪.
你會看到.
上帝是否錯配呢?.
竟然.
讓參孫在一個這樣的家庭.
似乎不太屬靈.
從小沒有教導兒子.
遵守上帝的吩咐.
到底發生了甚麼事?.
又或者在這件事上.
就算參孫的父母沒有能力阻止.
參孫這種任意妄為.
其實上帝可以直接干預.
上帝沒有理由容許一個那四人亂來.
為甚麼上帝不出手.
要容許參孫娶非利士人的女子?.
作者在第四節就開始告訴我們原因.
他在這裡說.
「他父母卻不知道這是出於耶和華」.
作者說.
「因為他找機會攻擊非利士人」.
「那時非利士人核制以色列人」.
我最初讀這段經文覺得很有趣.
我有個疑問.
上帝其實是全能的.
不是隨時隨地.
任何時候都可以攻擊.
非利士人的嗎?.
為甚麼他要找機會?.
連上帝都要找機會才能攻擊非利士人?.
難度在哪裡?.
這件事很困難嗎?.

$^{441}$為甚麼上帝要找機會?.
他不能像以前那樣.
興起一位士師.
帶領百姓向非利士人爭戰.
很簡單的.
一向都是這樣做.
成功的.
為甚麼這次.
好像困難了.
要找機會才能攻擊非利士人?.
到底發生了甚麼事?.
首先我們必須回到參孫故事的開頭.
十三章一開頭.
以色列人遊行.
耶和華眼中看為惡的事.
耶和華將他們交在.
非利士人的手中40年.
如果我們熟悉士師記.
我們說了這麼多.
按照慣常在士師記出現的循環.
下一步會是怎樣?.
百姓因為受到外敵的欺壓.
然後就.
喊苦喊罰.
叫苦連天.
然後就呼求上帝.
你出手拯救我.
最後上帝憐憫.
興起一位士師.
帶領百姓與外敵爭戰.
脫離受敵的日制.
這是我們說的基本的士師記循環.
但有沒有發覺參孫故事中.
這個循環已經改變了?.
我們再看不到百姓因為受敵人核制.
而叫苦連天的描寫.
取而代之是甚麼?.
是百姓已經甘心接受外敵的管轄.
讀一段經文我們就會明白.
在秩序表也印了那段經文.

$^{481}$跳到下一章.
十五章的九至十二節.
這裡記載了之後.
參孫與以色列人.
與非利人三方面一起所起的衝突.
經文說甚麼呢?.
它說.
非利人想安營在猶大.
布散在尼西.
猶大人說.
你們為何上來攻擊我們?.
他們說.
我們上來是要捆綁參孫.
他向我們怎樣行.
我們也要向他怎樣行.
於是有三千猶大人.
下到以坦盤的院內.
對參孫說.
非利人管轄我們.
你不知道嗎?.
你向我們行的是甚麼事呢?.
他回答說.
他們向我怎樣行.
我也要向他們怎樣行.
猶大人對他們說.
我們下來是要捆綁你.
將你交在非利人手中.
弟兄姊妹我們開始會明白.
為何.
非利人可以管轄.
以色列人長達多少年?.
四十年.
是否因為敵人太過勢力?.
不是.
而是因為.
神的子民已經.
甘心接受敵人的管治.
習慣了其實不想改.
其實覺得沒甚麼問題.
不再需要有意欲.

$^{521}$來呼求上帝幫忙.
所以你看到他們有多少人?.
有三千人.
三千個猶大人.
再加上很能打的參孫.
他們也不打算征戰.
他們只想.
息事寧人.
交參孫這個.
野士的人出來.
就這樣算了.
他們選擇交了參孫出來.
繼續在.
非利人的管轄下生活.
沒有想過要征戰.
換句話說.
其實你不難想像.
就算上帝.
興起一位比參孫更加純品.
更加正直的人來做事司.
很希望像從前一樣.
透過一位事司.
帶領現在的百姓.
來與非利人征戰.
百姓也不會再有意欲配合.
習慣了在罪的裡面.
上帝怎樣說也沒有反應.
上帝怎樣做.
他們也覺得算了吧.
他們寧願做非利人的順民.
也不會選擇順從上帝的帶領.
所以問題.
不是在上帝那裡.
是百姓已經怎樣了.
廣東話.
爛泥又怎樣了.
扶不上壁.
不過.
我們的信仰不是那麼差.
因為上帝不會這樣放棄.

$^{561}$上帝沒有放棄他的計劃.
所以你看到.
作者這裡說.
上帝要找機會.
來攻擊以色列人.
來拯救他的百姓.
他選擇在這一刻.
他不制止參宣與非利人的結合.
反過來使用參宣.
這種任性衝動的舉動.
來亂衝亂撞.
我們說在非利人與以色列人.
這個主僕關係裡面.
橫衝直撞.
其實你想像參宣.
就好像上帝手上一支鐵筆.
把兩塊黏得很緊的木板.
一塊是非利人.
一塊是以色列人.
黏得很緊.
用釘釘緊的.
一般方法是不行的.
唯有怎樣呢?.
上帝拿著參宣這支鐵筆.
透過他的橫衝直撞.
又怎樣呢?.
使勁攔勁.
撬鬆那段不應該有的關係.
拆散那段不應該有的關係.
這一段不合上帝心意的關係.
上帝說不可以讓它繼續發展下去.
所以等於姐妹.
當我們讀《士師記》的時候.
我們發覺好像越看越灰.
越看就越多人性黑暗.
污穢的面出來了.
我們亦會慨嘆為何.
身為上帝的子民.
他們的品格.
好像連一個外邦人.

$^{601}$一個連未信上帝的人都不如.
確實這是作者要傳遞其中一個信息.
來警惕我們每一位.
或每一個世代的基督徒.
不要重蹈覆轍.
有前車可鑑.
不過這不是《士師記》唯一的信息.
作者更加想我們明白.
在一片敗壞黑暗.
沒有盼望的光景裡面.
上帝從來沒有打算放棄.
上帝從來沒有打算.
因為人的軟弱.
推倒自己的旨意和計劃.
所以如果從「趙聖」上看.
以色列人至今墮落到這個地步.
其實應該怎樣?.
你還救不救他們?.
你是人.
不救了.
任由他們.
玩完後確認了.
沒救了.
但上帝說不是.
是有救了.
是有盼望.
因為他們不放棄.
他們不放棄.
所以以色列人仍然有生機.
因為上帝堅持對他們守約思慈愛.
即使表面上看來.
已經沒有人願意起來.
配合上帝的旨意.
就連事事的人格都大有問題.
上帝仍然默默地獨行其事.
他不需要人知道.
他亦有辦法使用人的軟弱和黑暗.
來配合他的工作.
來成就他的旨意.
無論參宣有多任意妄為.

$^{641}$一切仍然在上帝的掌握裡.
這是我們看事事記時.
何時都要捉緊的盼望.
繼續看下去.
我想請大家讀讀.
第三段經文.
記載在.
禪罪表裡第十四章五至九節.
讀多一次這段經文.
麻煩大家預備起.
參選特安府部下屬的立法會.
都邀請立法會普通人.
現有一隻小黃之子.
他怕黑暗.
要話令大大分工參選.
下屬仍受無欺害.
得獎時之私利.
亦忘之三日無一日.
下屬之子並沒有報訴父母.
參選立法會議員之說話.
就應業客.
過了這日子.
在下屬要取回原詞.
轉向道學 安可四次.
現有一筆公知和物.
在四次之後.
就用手取物.
且行且走.
保了父母家屬.
及家父母.
阿們也一票.
以示沒有報訴者.
是從四次之後取來的.
謝謝.
經文很清楚記載.
參選主動走近葡萄園.
獅子的屍體.
亦伸手取獅子屍體內的蜂蜜.
三個舉動.
請問這些舉動有甚麼問題.

$^{681}$參選可否走近葡萄園.
其實不可以.
一般以色列人可以.
但他不可以.
因為他是拿西人.
如果我們還記得.
拿西人其中一個重要要求.
不單止清酒濃酒不可喝.
酒是用葡萄製的.
就算葡萄乾.
一切與葡萄有關的製品.
他碰也不可以碰.
如果你是拿西人.
你明知這樣的時候.
其實你應該把防線設得更遠.
你不會主動走近葡萄園.
因為你知道很容易中招.
特別以參選這種性格.
真的隨時會出事.
第二方面.
他不可以走近死屍.
動物的死屍都不可以.
因為會令自己構成不潔.
他亦不可以拿死屍內的蜂蜜.
因為同樣會令自己不潔.
當然本質上他不應該接觸.
非理士人的女子和他結婚.
因為從相應的角度來說.
外邦人亦被視為不潔.
所以參選本來是要設法.
避開這一切的試探.
但現在他主動走近.
他主動接觸.
他可以任由自己.
落入一個不潔的狀況.
為了滿足自己的慾望.
他可以無視一切.
拿西人應該守的規範.
一切的界線都可以越過.
事實上你會發覺.

$^{721}$你覺得參選知不知道.
是不應該這樣做的嗎.
其實他知道的.
你覺不覺得.
他是故意知道.
如果這樣做的時候會有收尾.
因為如果一個拿西人.
他不潔的時候.
按照律法規定.
他要做很多禮儀獻祭.
使自己回復聖潔的狀態.
但他故意選擇不告訴別人.
不告訴他父母.
就是不想有收尾.
不想被人煩惱.
不想再接受規範.
又要搞一大輪.
才能讓自己回復聖潔.
所以你看到.
參選就是這樣的人.
為了滿足自己的需要.
他不介意身邊的人受牽連.
告訴你 無所謂.
結果他的父母被蒙在鼓裡.
被牽連成為不潔的人.
等於我們看到.
不知你覺不覺得參選.
好像一個一直在走的炸彈.
你不知道它何時爆.
它可能突然走到你身邊.
不單是炸到你.
最慘是炸到你.
參選就是這樣.
突然引爆.
不單是隨時禍及自己.
亦連累身邊的人.
但最慘的是.
他不覺得是甚麼問題.
無所謂 隨便.
總之我需要的東西先滿足了.

$^{761}$為了得到自己想要的東西.
他一點也不介意.
打破上帝給他的規範.
所以其實參選這種情況.
當然我們未必像他那樣.
去到極端.
但某程度上我們同樣.
上帝給了我們聖潔的身份.
這個聖潔身份隨之而來.
在生活裡我們需要.
接受上帝話語的規範.
我們的品格.
行事為人.
和我們的身份是否相稱.
亦很值得我們去想.
我們是否持守著.
基督徒的身份.
為了這個緣故.
我們願意接受上帝的管教呢.
當然我們會明白.
很多時候人不太喜歡.
管教規範這兩個字.
因為一聽到就覺得.
我的自由被限制.
我無法作主.
我無法隨心所欲.
這樣會很不爽.
很不暢快做人.
不過我們試試反過來.
規範真正的用途是甚麼.
是保護.
是保護你.
不要出事.
不要危險.
我很喜歡用這個比喻.
之前也和弟兄姊妹說過.
行人天橋有沒有走過.
有走過.
但不是行人天橋.
現在那些好些.

$^{801}$現在那些會包頂.
但舊式的行人天橋.
沒有頂.
只有兩邊的圍欄.
你覺得那是限制你.
還是保護你.
兩邊的圍欄.
好像很難理解.
如果拆走兩邊的圍欄.
你覺得那條行人天橋.
有沒有那麼安全.
你會否夠膽走到問哪些.
不會吧.
如果沒有兩邊的圍欄.
橋這麼闊的話.
這裡就只有這麼多.
你不會夠膽走過.
因為你怕掉下去.
但反過來有了兩邊的圍欄.
其實你走得闊了.
因為你知道.
無論問哪邊.
兩邊都有保護.
不會死不會掉.
所以其實圍欄.
到底是限制了我們的自由.
還是確保我們.
安全自由.
其實看你怎樣看.
看你怎樣看.
事實上我們知道.
上帝的規範.
是重要的.
上帝給我們的管教.
是重要的.
是為了保障我們的生命.
不會越過了.
上帝給我們的界線.
以致我們得罪了上帝.
破壞和上帝的關係.

$^{841}$不過你看到的是.
參孫一再無事.
上帝給他的規範.
任由自己的私慾.
來到他無限放大.
無限滿足的時候.
就快要跌了.
下一次你會看到他跌.
如姑娘所說.
其實你看到.
參孫現在好像沒事.
但其實距離崩壞的結局.
其實不遠了.
他對我們的生命.
其實是一種警惕.
其實你看到.
頗諷刺的.
剛才我們讀的經文.
他可以靠上帝的幫助.
徒手制服強大的獅子.
但他卻無法勝過獅子.
獅體裡面那些.
微小的蜂蜜誘惑.
很奇怪.
相當的諷刺.
正如他之後.
他可以靠上帝的幫助.
以一人之力.
來制服強大的非利士人.
但他同時又無法勝過.
一個弱小的非利士女子的引誘.
弟兄姊妹其實.
我們的情況會否和他相似.
在某些範疇裡面.
我們很有能力.
我們揮灑自如.
可以好好處理面對.
但在一些很細小的引誘裡面.
我們卻是軟弱跌倒.
很容易被他搞定.

$^{881}$某程度上其實.
參孫有慾望的試探.
其實我們也有.
我們可能有購物慾.
佔有的慾.
權力慾.
或是現在上網的慾.
又或是像參孫一樣.
有情慾的試探.
其實這些試探.
都會好像在參孫身上.
其實成為了你或我的死穴.
一戳下去就中.
然後你就好像無力反抗.
被牠牽著走.
到底是及早被上帝介入.
幫我們好好去管理.
這些慾望試探.
還是好像參孫一樣.
任由自己的慾念無限放大.
我還記得最初出來工作.
那時是一九九幾年.
熟悉我的弟兄姊妹.
知道我很喜歡玩高達.
玩玩具.
我記得最初工作月入萬多元.
每個月花兩千至三千元.
即是多少百分比.
三十吧.
大概.
買我的心頭好動漫的產品.
瘋狂是甚麼呢.
舉例.
我不喜歡Hello Kitty.
不過我聽過一些弟兄姊妹.
她喜歡Hello Kitty.
總之袋蓋了Hello Kitty.
總之袋是粉紅色她就買.
去到這個地步.
不過我不是喜歡Hello Kitty.

$^{921}$但我喜歡一些動漫機械人產品.
總之差不多.
蓋了她的頭.
我可能差點就買了.
只要她漂亮我就會買.
買了買了.
我發覺這樣不知不覺就買了三年.
你計算一下.
花了多少錢.
可能花了十幾二十萬.
是說多少年前.
三十年前.
都不少錢.
其實買一買你會覺得很無謂.
因為一開始買回來會拆出來看.
後來買了買了都放在一邊.
然後堆了一堆.
你都不再記得買了甚麼.
但就好像一個印.
撕不掉的.
很特別.
它又不是殺人放火.
又不是傷天害理.
又不是不良嗜好.
更加不會影響到其他人.
我又負擔到.
有甚麼問題.
確實好像真的沒有問題.
何時開始改變呢.
信了主就改變了.
特別是信了主之後.
其實現在回頭想.
當你原來發覺.
你怎樣去撕掉這些人.
你要找到一個更大.
更值得追求的目標和對象.
其他的就變得次要.
我先理生.
我現在不是不買玩具.
我現在不是無欲無求.

$^{961}$我都會買.
但節制了很多.
那件事不是很重要.
更加重要的是.
你將焦點放在上帝身上.
很多慾念慾望.
你就會放得低.
真的,弟兄姊妹.
我想我們每一個人.
都多多少少有所謂慾望慾念.
有些甚至乎不是.
十惡不捨的事.
聖經亦沒有主張.
我們要行禁慾主義.
無欲無求,甚麼都不要想.
聖經教導我們的是.
功克己身.
好好管理我們的慾望.
正如史祿寶樂所說.
我們都熟悉這段經文.
哥林多前書九章二十七節.
我是功克己身.
叫身服我.
恐怕我全福音給別人.
自己反被棄絕了.
其實幾次參宣的情況.
他被上帝使用.
亦都有一些果效.
但最後他和上帝的關係.
是離寒離辣的.
很淒涼的,很慘的.
來提醒我們.
我們都需要靠主.
幫我們努力.
在生活上大大小小的.
慾望試探.
否則我們真的有可能.
變成另一個參宣.
雖然上帝有使用他.
甚至他的事奉.

$^{1001}$亦都有些果效.
但他自己都不知道.
原來已經偏離上帝.
還要越行越遠.
弟兄姊妹.
參宣這個故事.
其實提醒我們.
一位理想的上帝僕人.
不只是有能力就行.
上帝更加體重的是.
他僕人的品格.
品格是怎樣.
才是更加重要.
盼望今天我們藉著.
上帝這段話語.
來光照我們的生命.
來潔淨我們的一些心思意念.
讓我們可以成為.
聖潔合用的器皿.
獻給上帝使用.
我們一起禱告吧.
天主上是願你的話語光照我們.
成為我們腳前的燈.
路上的光.
確實我們未必像參宣一樣.
去到極端的衝動為所欲為.
但多多少少我們都有他的縮影.
我們都會有時被.
我們的慾望慾念所控制.
而將上帝和你的話語.
你的吩咐放在一邊.
甚至我們可能有教會生活.
有事奉.
但今日的經文.
很值得我們重新檢視.
我們的品格.
我們和上帝的距離.
我們是否真的願意.
以上帝為居首.
將其他事情捍衛次要.

$^{1041}$盼望上帝你藉著這段話語.
幫助我們.
讓我們重新回到你的裡面.
讓我們真的靠上帝你.
來好好管理我們的人生.
或他有一些不合上帝你心意的.
一些想法,一些行為.
求上帝幫我們去撕掉它.
求上帝幫我們好好管理.
我們的心思意念.
讓我們可以以你為首.
讓我們全人都可以被你潔淨.
成為合用的器皿來榮耀你.
為此我們同心獻上祈禱.
奉耶穌基督你得勝的名字來祈求.
阿們.
\newpage



\section{士師記 16:1-31}
\label{sec:rllwPQsYaK0}
\textbf{2023年9月17日崇拜直播｜余嘉敏姑娘｜「家醜說出口」系列:過度自信者的求告｜士師記16\_1-31}
\newline
\newline
連結: \href{https://youtube.com/watch?v=rllwPQsYaK0}{\texttt{https://youtube.com/watch?v=rllwPQsYaK0}} ~~~~ 語音日期: 2023-09-18
\newline
\newline
\hyperref[sec:Up7U5HlOZfo]{\small{< < < PREV SERMON < < <}}
~
\hyperref[sec:index_chronic]{\small{[返順時目]}}
~
\hyperref[sec:jXXnWoB7EBk]{\small{> > > NEXT SERMON > > >}}
\newline
\newline
士師記 16:1-31
\newline
\begin{longtable}{cl}
\hline
\hline
章節 & 經文 (和合本修訂版)\\
\hline
16:1 & \begin{tabularx}{0.7\textwidth}{X} 參孫到了迦薩，在那裡看見一個妓女，就與她親近。 \end{tabularx} \\ \\ \relax
16:2 & \begin{tabularx}{0.7\textwidth}{X} 有人告訴迦薩人說：「參孫到這裡來了！」他們就包圍起來，整夜在城門埋伏等著他。他們整夜靜悄悄地，說：「等到天一亮我們就殺他。」 \end{tabularx} \\ \\ \relax
16:3 & \begin{tabularx}{0.7\textwidth}{X} 參孫睡到半夜，在半夜起來，抓住城門的門扇和兩個門框，把它們和門閂一起拆下來，扛在肩上，抬到希伯崙前面的山頂上。 \end{tabularx} \\ \\ \relax
16:4 & \begin{tabularx}{0.7\textwidth}{X} 這事以後，參孫在梭烈谷愛上了一個女子，名叫大利拉。 \end{tabularx} \\ \\ \relax
16:5 & \begin{tabularx}{0.7\textwidth}{X} 非利士人的領袖上去，到那女子那裡，對她說：「請你哄騙參孫，探出他為何有這麼大的力氣，以及我們要用甚麼方法才能勝他，將他捆綁制伏。我們就每人給你一千一百塊銀子。」 \end{tabularx} \\ \\ \relax
16:6 & \begin{tabularx}{0.7\textwidth}{X} 大利拉對參孫說：「請你告訴我，你為何有這麼大的力氣，要用甚麼方法才能捆綁制伏你。」 \end{tabularx} \\ \\ \relax
16:7 & \begin{tabularx}{0.7\textwidth}{X} 參孫對她說：「若用七條未乾的新繩子捆綁我，我就像平常人一樣軟弱。」 \end{tabularx} \\ \\ \relax
16:8 & \begin{tabularx}{0.7\textwidth}{X} 於是非利士人的領袖拿了七條未乾的新繩子來，交給她，她就用繩子捆綁參孫。 \end{tabularx} \\ \\ \relax
16:9 & \begin{tabularx}{0.7\textwidth}{X} 當時，埋伏的人正在她的內室等著。她對參孫說：「參孫，非利士人來捉你了！」參孫就掙斷繩子，繩子如遇到火的麻線斷裂一樣。這樣，人還是不知道他的力量從哪裡來。 \end{tabularx} \\ \\ \relax
16:10 & \begin{tabularx}{0.7\textwidth}{X} 大利拉對參孫說：「看哪，你欺騙我，對我說謊。現在請你告訴我，要用甚麼方法才能捆綁你。」 \end{tabularx} \\ \\ \relax
16:11 & \begin{tabularx}{0.7\textwidth}{X} 參孫對她說：「若用未曾用過的新繩子捆綁我，我就像平常人一樣軟弱。」 \end{tabularx} \\ \\ \relax
16:12 & \begin{tabularx}{0.7\textwidth}{X} 大利拉就用新繩子捆綁他，對他說：「參孫，非利士人來捉你了！」當時，埋伏的人在內室等著。參孫掙斷手臂上的繩子，如掙斷一條線一樣。 \end{tabularx} \\ \\ \relax
16:13 & \begin{tabularx}{0.7\textwidth}{X} 大利拉對參孫說：「你到現在還是欺騙我，對我說謊。請你告訴我，要用甚麼方法才能捆綁你。」參孫對她說：「只要用織布的線將我頭上的七條髮綹編織起來就可以了」。 \end{tabularx} \\ \\ \relax
16:14 & \begin{tabularx}{0.7\textwidth}{X} 於是大利拉用梭子將他的髮綹釘住，對他說：「參孫，非利士人來捉你了！」參孫從睡中醒來，將織布機上的梭子和織布的線一齊都拔出來了。 \end{tabularx} \\ \\ \relax
16:15 & \begin{tabularx}{0.7\textwidth}{X} 大利拉對參孫說：「你既不與我同心，怎麼能說『我愛你』呢？你這三次欺騙我，不告訴我，你為甚麼有這麼大的力氣。」 \end{tabularx} \\ \\ \relax
16:16 & \begin{tabularx}{0.7\textwidth}{X} 大利拉天天用話催逼他，糾纏他，他就心裡煩得要死， \end{tabularx} \\ \\ \relax
16:17 & \begin{tabularx}{0.7\textwidth}{X} 終於把心中的一切都告訴她。參孫對她說：「從來沒有人用剃刀剃我的頭，因為我一出母胎就歸給神作拿細耳人。若有人剃了我的頭髮，我的力氣就會離開我，我就像平常人一樣軟弱。」 \end{tabularx} \\ \\ \relax
16:18 & \begin{tabularx}{0.7\textwidth}{X} 大利拉見他說出了心中的一切，就派人去召非利士人的領袖，說：「請再上來一次，因為他已經說出了心中的一切。」於是非利士人的領袖手裡拿著銀子，上到她那裡。 \end{tabularx} \\ \\ \relax
16:19 & \begin{tabularx}{0.7\textwidth}{X} 大利拉哄參孫睡在她的膝上，叫一個人來剃掉參孫頭上的七條髮綹。於是大利拉開始制伏參孫，他的力氣就離開他了。 \end{tabularx} \\ \\ \relax
16:20 & \begin{tabularx}{0.7\textwidth}{X} 大利拉說：「參孫，非利士人來捉你了！」參孫從睡中醒來，說：「我要像前幾次一樣脫身而去。」他卻不知道耶和華已經離開他了。 \end{tabularx} \\ \\ \relax
16:21 & \begin{tabularx}{0.7\textwidth}{X} 非利士人逮住他，挖了他的眼睛，帶他下到迦薩，用銅鏈鎖住他，叫他在監獄裡推磨。 \end{tabularx} \\ \\ \relax
16:22 & \begin{tabularx}{0.7\textwidth}{X} 然而他的頭髮被剃以後，又開始長起來了。 \end{tabularx} \\ \\ \relax
16:23 & \begin{tabularx}{0.7\textwidth}{X} 非利士人的領袖聚集，要向他們的神明大袞獻大祭，並且慶祝，說：「我們的神明把我們的仇敵參孫交在我們手中了。」 \end{tabularx} \\ \\ \relax
16:24 & \begin{tabularx}{0.7\textwidth}{X} 眾人看見參孫，就讚美他們的神明說：「我們的神明把那毀壞我們的地、殺害我們許多人的仇敵交在我們手中了。」 \end{tabularx} \\ \\ \relax
16:25 & \begin{tabularx}{0.7\textwidth}{X} 他們心裡高興的時候，就說：「叫參孫來，逗我們歡樂。」於是他們把參孫從監獄裡提出來，在他們面前戲耍。他們叫他站在兩根柱子中間。 \end{tabularx} \\ \\ \relax
16:26 & \begin{tabularx}{0.7\textwidth}{X} 參孫對牽他手的僮僕說：「讓我摸摸支撐這廟宇的柱子，我要靠一靠。」 \end{tabularx} \\ \\ \relax
16:27 & \begin{tabularx}{0.7\textwidth}{X} 那時廟宇內充滿男女，非利士人的眾領袖也都在那裡，屋頂上約有三千男女觀看參孫逗他們歡樂。 \end{tabularx} \\ \\ \relax
16:28 & \begin{tabularx}{0.7\textwidth}{X} 參孫求告耶和華說：「主耶和華啊，求你眷念我。神啊，就這一次，求你賜給我力量，使我向非利士人報那挖我雙眼的仇。」 \end{tabularx} \\ \\ \relax
16:29 & \begin{tabularx}{0.7\textwidth}{X} 參孫抱住中間支撐廟宇的兩根柱子，左手抱一根，右手抱一根。 \end{tabularx} \\ \\ \relax
16:30 & \begin{tabularx}{0.7\textwidth}{X} 然後他說：「讓我與非利士人一起死吧！」他盡力彎腰，廟宇就倒塌了，壓住領袖和廟宇內的眾人。這樣，參孫死的時候所殺的人比活著所殺的還多。 \end{tabularx} \\ \\ \relax
16:31 & \begin{tabularx}{0.7\textwidth}{X} 他的兄弟和他父親的全家都下去收他的屍首，抬上去，葬在瑣拉和以實陶中間、他父親瑪挪亞的墳墓裡。參孫作以色列的士師二十年。 \end{tabularx} \\ \\
[1ex]
\hline
\hline
\end{longtable}
$^{1}$今天選讀的經訓是在《士師記》十六章第一至四節.
以及二十五至三十一節.
在舊約聖經的三百二十二開始.
請聽我讀出.
參孫到了加薩.
在那裡看見一個妓女就與她親近.
有人告訴加薩人說.
參孫到這裡來了.
他們就把她團團圍住.
中夜在城門悄悄埋伏.
說等到天亮我們便殺她.
參孫睡到半夜起來.
將城門的門扇,門窗,門衫.
一起拆下來.
剛在肩上.
剛到希伯倫前的山頂上.
後來參孫後來後來後來後來….
(參孫的名字).
(參孫的童子說).
參孫求告耶和華說.
主耶和華求你眷念我.
神啊,求你賜我這一次的力量.
使我在非理士人身上.
布那碗我雙眼的仇.
參孫就抱著托房的那兩根柱子.
左手抱一根.
右手抱一根.
說我情願與非理士人同死.
就盡力屈身.
防止倒塌.
壓在首領和房內的眾人.
這樣參孫死時所殺的人.
比活著所殺的還多.
參孫的弟兄和他父的全家.
都下去取他的屍首.
抬上來.
葬在所拉和耳實圖中間.
在他父嗎哪亞的墳墓裡.
參孫作以色列的事司二十年.
各位弟兄姊妹,平安.

$^{41}$上星期Ray Mook帶我們分享了.
參孫未出世已被神揀選.
終身跟從他成為拿世二人的事司.
被委以重任.
被神揀選.
對很多人來說.
應該是一個很大的祝福.
應該按照上帝的旨意.
和分別為聖.
為上帝爭戰.
但相反.
參孫的行為非常隨心所欲.
他按自己的慾望行事.
多於按照上帝的旨意.
最納悶的就是.
上帝仍然繼續使用他.
順應參孫的隨心所欲行為.
去成就拯救以色列人的旨意.
不過他這種隨心所欲的行為.
並非沒有後果.
今天我們一起看看.
參孫會有什麼終局.
我們一起禱告.
親愛的三一上帝.
去教導我們.
求聖靈你光照我們的內心.
幫助我們有一份柔軟的心去聆聽.
也幫助我們去實踐你的真理.
叫我們今天聽完道之後.
我們走出禮堂門口.
我們的生命更加貼近你.
多謝你聽我們的禱告.
奉主耶穌基督的聖名祈求.
阿們.
認識我的人應該都知道.
我很喜歡一件事.
就是我很喜歡開車.
特別是我喜歡開車到處去.
去新路 探新的地方.
我覺得很爽的感覺.

$^{81}$很無拘無束.
享受那份開車的自由.
我覺得自己的體位應該都不錯.
為什麼呢?.
因為我通常開車.
都是一至兩手tie .
我就已經停好車了.
但當回想我剛考到牌照的時候.
我絕對不是像今天這樣.
那麼有自信.
當時我開車.
我問我老公就知道.
我是震驚的.
我最怕的就是泊位.
每一次我開車之前.
你猜我做什麼?.
不是TASI.
我是會先祈禱.
我祈求懇切地求神.
你給我智慧.
你給我看位準確.
給我保存我和途人的生命.
不過日子久了.
我的技術就成熟了很多.
連油麻地 佐敦這些.
那麼窄的停車場我都泊到.
我試過自駕遊上雪山.
也去過台灣出名比較危險的太老角.
一個女孩帶著兩個朋友.
上到太老角玩.
自此之後.
我就覺得自己駕馭得到.
我行的.
特別是之後我就會更加有信心.
慢慢就喜歡開快一點.
特別是當我兩個兒子.
跟我說 讚我.
媽媽明明我們同一個時間出門上學.
但你送我上學的時候.
我就不會遲到.

$^{121}$但爸爸送我上學就會遲到.
我聽完之後.
我心裡就沾沾自喜.
很有自信.
覺得自己行的.
我都很有效率.
超小小速我又能控制.
駕得又穩又穩.
有什麼問題.
我完全忘記了.
當初神賜給我的智慧去駕車.
也忘記了我每一次平平安安去到目的地.
其實都是神的恩典.
久而久之.
我現在都沒有了.
沒有了上車祈禱的習慣.
我再沒有依靠神駕駛.
我覺得自己有信心的地方.
其實做人有信心有自信.
其實是很重要的.
因為自信可以幫助我們.
面對挑戰的困難.
如果沒有自信的話.
我駕車連小路都不敢出大路.
不過當我們過度自信的時候.
往往都很容易失去理智.
以為自己行的.
很可能我們會闖大禍.
甚至乎後果可以不堪設想.
其實什麼叫做過度自信.
Overconfidence在心理學裡有個定義.
就是受到信念,情緒,偏見,感覺等.
這些主觀的心理因素所影響.
人往往都會過度自信.
高估自己的判斷能力.
亦都高估成功率.
就是百分比.
以為自己一定行.
簡單來說.
過度自信會使我們失去對現實的清晰認知.

$^{161}$高估自己的能力.
低估外面的形勢.
甚至我們很容易誤踏輕敵的陷阱.
尤其是在界線方面.
容易高估自己能力.
界線能力可以應付.
所以我們很容易踩界.
以為自己有能力招架一切.
參宣就是這樣的人.
他任意妄為,隨心所欲.
源於他過度自信.
他輕視神給他的界線.
甚至他喜歡專門踩界.
他走在踩界的位置挑釁敵人.
目的是表現自己.
而且他與其他士師不同.
他背後沒有軍隊.
亦不需要有人支援他.
他只用自己的方法.
向非理士人報仇雪恨.
這種性格在今日的經文.
表露得淋漓盡致.
而他這份過度自信的生活態度.
將會帶來的悲劇收場.
對於我們今日的我們來說.
究竟參宣的終局.
又對我們有甚麼借鏡呢?.
我們一起看看第一至第三節.
這段經文很有趣.
這三節看來可有可無.
我們看上文第十五章.
和下理第四節打後.
看來好像沒有直接關係.
上一段說到.
參宣單人匹馬橫擊非理士人.
結果擊殺了一千個.
打算與他們打仗的非理士人.
做完這個壯舉後.
他覺得很口渴.
很想喝水.

$^{201}$他的反應很有趣.
他走去像一個未懷孕的嬰兒.
去投訴父親.
為何不給他喝水.
但很吊詭的地方是.
神又聽他的.
神給他行了一個神跡.
叫利希的蛙柱湧出水來.
參宣就有水喝了.
而在第二十節.
這裡有提及.
二十節說.
參宣因為這個神跡.
成為了以色列人認可.
作他們事先二十年.
我們又看看夏理.
第四節的打後.
記載非理士人.
一直覺得大巨人.
想整治他.
於是他找了一個婦人.
戴莉拉引誘他.
這件事.
我們看回第一至三節.
這段經文時.
我們發現抽起這段.
其實無損我們理解.
故事的發展.
不過這段經文非常重要.
記載了參宣去加薩的事跡.
這段事跡讓我們看清楚.
參宣是一個怎樣的人.
原來是充滿過度自信的表現.
我們一起讀第一節.
參宣到了加薩.
眼裡看見一個女兒.
與他親近.
加薩是一個怎樣的地方呢?.
原來加薩這個名字.
代表著強壯.

$^{241}$現在叫做易造加薩.
加薩是當時.
非理士人的首都.
是屬於他們五大.
最大的城市之一.
如果看地圖.
你會發現.
它是在地中海的沿岸.
你有沒有想起地中海旅行?.
陽光普照.
很美的.
我相信那裡是一個.
很受歡迎的度假勝地.
參宣是否去那個度假呢?.
還是他有什麼.
要去非理士人的地方呢?.
聖經沒有交代.
但作為神所差派.
去消滅非理士人的事司.
他竟然大無私讓去那裡.
竟然要和一個妓女親近.
其實不用多解釋.
我們都知道.
其實都不是很恰當.
他沒有遵守到.
他作為拿世蟻人.
要該守的聖潔.
不單止這樣.
他這個行為.
還極富挑戰性.
在上兩章才剛剛說到.
他說他要想娶一個.
庭娜女子.
即是非理士人的女子.
他就殺了三十個非理士人.
奪取了他們的衣服.
然後給了一些猜出謎語的人.
之後他的岳父.
就把他剛剛認取的婦人.
給了作為他的朋友的男人.

$^{281}$於是參宣大發雷霆.
報復式地將非理士人的和將.
和這個橄欖園都燒了.
不過他都真的頗大力.
如果我們看第十五章的四至五節.
他用的方法都很神奇.
他捉了三百隻狐狸.
將狐狸尾巴一對一對綁住.
中間放一個火把.
然後夾住了.
捉了三百隻.
你叫我捉一隻.
我都死了.
不單止這樣.
他還點著火.
將狐狸放進田裡.
做什麼?.
農田來的.
滅了人家的生計.
這種報復式的行為.
演變成為以色列的危機.
整個國家的危機.
第十五章第九節開始說.
這班非理士人要上安寧在猶大.
散佈在尼西.
原來他們想準備攻打以色列.
以色列知道這件事之後.
他們就自保.
於是他們就去找參宣.
結果就將參宣.
用兩條繩子綁住了.
參宣就交給非理士人.
不過參宣真的很厲害.
他能夠殺出蟲蟻.
第十五節說.
他用了一個未乾的蘆筛骨.
擊殺了一千個非理士人.
一個這樣和非理士人為敵的士師.
在我們今天的經文.
他竟然單人匹馬.

$^{321}$去了加薩這個地方.
簡直送羊入虎口.
難道只是因為女色.
傻乎乎地去了加薩嫖妓?.
我相信應該不會這麼簡單.
以參宣的性格.
他一定不會覺得自己是羊仔.
我們一起看看今天的經文.
第二節說.
果然不出所料.
加薩人知道他來了.
他們就將他團團圍住.
中夜在城門悄悄埋伏.
想做什麼?.
等到天亮的時候就殺了他.
明顯非理士人過了二十年.
心頭大恨都未能放下.
他們無時無刻都想消滅.
報這個仇.
但第三節就說.
看到參宣睡到半夜.
他就起床.
將城門的門扇,門窗,門衫.
一齊都拆下來.
剛在他的肩上.
不知道背到哪裡.
就是去到希伯倫山.
回到自己的地頭的山頂上.
其實參宣這個舉動.
都很難明白.
我看了幾次這段經文.
離開城門就離開城門.
不用拆了城門.
然後背到對面山這麼誇張.
很可能他知道.
非理士人的計謀.
所以他沒有睡到天亮才醒.
還將城門拆下來.
背在肩上.
我們先說說這個城門有什麼特別.

$^{361}$原來有些專家說.
古代的城門有多重?.
我都驚訝了.
原來有五千至一萬磅.
這麼重的重量.
如果以壯男.
即是那些健碩的男人.
以二百磅計算的話.
即是等於多少個男人?.
計算很厲害的那些.
即是大概二十五至五十個.
這麼大的壯男.
將他搬到對面的山.
而希伯倫距離加薩.
有四十至五十公里路程.
所以路程一點都不近.
如果單靠走路.
其實至少要走十八個小時.
所以如果我們這樣去理解.
這段經文的時候.
有可能我們想像到.
參孫就像李小龍.
在精武門裡單人匹馬.
走去踢館.
然後將人家武館的牌匾拆下來.
參孫就像要向非利人士宣示.
我多厲害.
你別指望走近我身.
而加薩.
剛才我們說到.
是代表著強壯的意思.
參孫的舉動就像.
告訴非利人士.
我在你覺得最強壯的城市裡.
連防禦門都拆下來.
想告訴他們.
我除了展示自己的力量之外.
其實是要向敵人宣戰.
對於非利人士來說.
這是一個極大的羞辱.

$^{401}$參孫的確力大無窮.
正如我剛才所說.
他不像其他士師.
有一隊軍隊和他一起迎戰.
他總是獨來獨往.
一個人去對戰非利人士.
就好像一部電影《殺神》.
John Wick 裡面.
每一個人在打以一敵百.
好像沒有人可以抵擋他.
阻擋他向前報仇.
其實他的力量是非常沉上.
但是他的力量從哪裡來呢?.
我們看回第十四章十九節.
當他擊殺三十個非利人士時.
經文說是耶和華的靈大大感動參孫.
他才可以殺死三十個人.
第十五章第十四至第十五節.
當非利人士和他開戰時.
是神的靈大大感動他.
他才可以用盧塞骨.
擊殺一千個非利人士.
由此可見.
因著神的靈大大感動他.
他才有超乎的力量擊敗非利人士.
你猜參孫知不知道?.
我覺得他是知道的.
但是你看看他擊打一千人之後.
他說的一句話.
他說我用盧塞骨殺人成堆.
用盧塞骨殺了一千人.
他的主體是誰?.
是我.
我們可以看到.
他沒有把神的感動.
和神在他身上的作為放在心裡.
他反而過度自信地.
引伸了一些驕傲出來.
他把神所成就的榮耀.
歸在自己的身上.

$^{441}$參孫完全忘記了.
這完全是神的恩典.
他也沒有把榮耀歸給.
次能力給他的上帝.
他認為我擁有這些能力.
理所當然.
除了沒有感恩之心.
彷彿他覺得自己做什麼都好.
我這些力大無窮的恩賜能力.
是不會被人拿走的.
再者我們看看.
當我們過度高舉自己的時候.
我們不要妄想.
其他人會從我們的生命裡.
反照到神的榮耀.
非利士人一直看著參孫.
他只看到參孫的力大無窮.
他從來沒有看到.
他在他的身上看到神的榮耀.
我們看看第四節和第五節.
非利士人叫了一個婦人.
達莉拉去看守參孫.
她要設法找出.
參孫背後的神奇魔力.
很可能她認為.
參孫的超自然能力.
就像我們現在戴著的護身符.
或能量手鏈.
只要拿走那件物件.
我們就可以打敗參孫.
看回我們的生命.
我們應該沒有在座的人.
像參孫一樣.
有這個力大無窮的恩賜.
也沒有被上帝呼召我們.
去拯救萬民.
為他爭戰這麼大的恩賜.
但毋庸置疑.
我們每一個人.
都是上帝賜我們每一個人.

$^{481}$有些很獨特的恩賜.
有些很獨特的使命.
讓我們在我們的家庭裡.
在我們的教會當中.
我們的工作場所.
甚至我們的社區裡.
去做一些神想我們做的工作.
但我們苦心自問.
我們有沒有認定.
這一切是神賜予給我們的呢?.
是他授權給我們.
命定我們去完成的呢?.
還是你將這些恩賜才幹.
智慧.
甚至成功的一些個案例子.
那些恩惠據為己有.
我們有沒有覺得.
這個範疇是我熟悉的.
我也是專家.
我是專家.
所以這一切的東西.
這個領域,這個範疇.
上帝你不要搞我.
你不要插手.
我能搞定的.
覺得這些功德,成功.
全部都是出於自己.
我們有沒有用上帝給我們的恩賜.
去榮耀我們上帝呢?.
還是我們會用錯了我們的恩賜.
去榮耀我們自己.
甚至是將上帝的名蒙羞.
當過度自信所引出來的驕傲.
其實會將我們.
將神所做的工作.
歸落我們自己那裡.
所以弟兄姊妹.
不要小看這個過度自信.
它不單止讓我們.
看不見上帝的榮耀.

$^{521}$它也會削弱我們對罪的敏感度.
令我們越來越離開上帝的旨意.
我們一起看看下一段.
這裡說到我們會見到.
參孫喜歡上一個女人.
剛才說的那個大利拉.
又是女人.
這個大利拉很努力.
起碼經文裡至少三次以上.
她出一些「lum 宮」去批評參孫.
要她說出「你這麼有能力」.
「還有什麼才可以制衡你」.
「捆綁你」.
當我們細心看大利拉的說話時.
作為讀者的我們.
其實也不難發現.
這個女人有一點「境廣」.
但參孫好像不太有「境覺性」.
我們看看第六節這裡說.
大利拉第一步直接問參孫.
「你這麼有能力」.
「還有什麼才可以捆綁你」.
「克制你」.
第二次在第十一節裡說.
第二次和第十三節.
開始用了另一招.
我想在座每一個女人都認識.
就是撒嬌.
發脾氣.
然後撒嬌參孫.
沒有什麼特別伎倆.
應該每一個人都認識.
所以如果要撒嬌.
或者要拆穿她.
其實也不難.
但我們看參孫的回應.
她沒有什麼防備心.
而且她一次次.
比一次次的答案更有創意.
她好像在玩遊戲.

$^{561}$我們一起看看第六至第七節.
她說如果用七條未乾的清繩子.
捆綁我.
我就會軟弱.
好像普通人一樣.
七條未乾的清繩子是什麼呢?.
其實是七條未乾的繩子.
清繩子是當時普通人家都會用的繩子.
是遊牧民族很喜歡用的繩子.
所以其實沒什麼特別.
到第二次她再回應.
她回應說用過沒用過的新繩子.
不就是乾的.
第三次她更有創意.
她說把頭髮上的髮柳.
七條髮柳和尾線織在一起.
意思是把七條的辮子.
和織布機橫中橫直.
不懂得怎麼說.
總之就是織在織布機上.
然後釘死它.
之後就可以了.
聽起來答案都越來越離譜.
都不太合理.
我覺得好像挺好玩的.
還沒察覺到自己已經在危險當中.
但是大利拉其實一次又一次.
就像第八章和第九節那樣說.
她照樣捉住了.
用方法用繩子捆綁參孫.
而且每一次都會預備一些人.
這裡的人是眾數.
所以是多過一個人.
去埋伏大利拉的內室.
準備捉住她.
就橫擊了.
雖然我明白.
但想也想不通.
其實那個地方是不是很大呢.
為什麼多過一個人埋伏在內室.

$^{601}$參孫都發現不到.
就當她真的發現不到.
當枕邊人不斷問你什麼弱點.
之後又會照樣捉住你.
好像想弄你一樣的時候.
其實足以令我們意識到.
這個女人可能有不懷好意.
作為讀者我一邊看的時候.
心裡會有一句.
你為什麼這樣都看不到.
有可能參孫其實早就知道.
但她很享受這種刺激感.
自信地覺得即使有危險.
那又怎樣呢.
我覺得自己可以駕馭這些危險的場面.
踩下界 玩下又怎樣呢.
只要我凌駕得到 招架得到.
以我的聰明才智 以我的體力.
即使有人來捉我.
我都絲毫無損.
而這種心態其實不是第一次.
出現在參孫的心裡.
其實一點都不奇怪.
她作為非利士人的頭號敵人.
她竟然多次自由出入非利士人的重地.
毫無戒心.
明明知道自己親近女色是她的超級弱項.
但她每一次的危機都發生在女人身上.
但她毫無忌諱.
沒有學習過她的教訓.
她總是親近女人.
甚至走近妓女身邊身犯險境.
很自信地覺得我有能力招架所有危險危機.
她覺得自己能夠抵禦罪的能力很高.
對於開車很有自信的我.
剛才不是說我有自信的我.
其實很明白參孫的心態.
因為我覺得自己招架得到.
我駕馭得到.
不會出意外的.

$^{641}$開快一點沒問題.
只要我在拍快照的位置.
收油就行了.
亦因為這份過度的自信.
我明明看到紅燈快要轉紅燈.
我仍然會踩油衝過去.
我不會那麼差的.
我不會衝紅燈的.
但大家都知道.
我應該高估了自己.
結果某年某月某日.
我累積了一個為例分數.
超過了十分.
所以要參加一個強制駕駛改進課程.
這個經驗是我人生中.
很尷尬的經驗.
坐在我身邊的.
你猜他們是甚麼人?.
全部都是專業的貨車司機,的士司機.
我們十多人.
被困在一間類似G房的房間.
我們要在那裡坐八個小時.
然後那個司機.
好像我們應該坐牢.
像懲教署的職員.
不斷在那裡囚禁我們.
又在那裡.
不知怎樣說.
遊蕩了我們一段時間.
八個小時.
被囚禁不止.
我們聽完這些後.
還要考一個試.
即場考.
你合格才可以離開.
自此之後.
這個慘痛的經驗.
我就不敢高估自己的能力.
去為例駕駛.
其實過度自信.

$^{681}$弟兄姊妹.
不是沒有後果的.
原來會削弱我們對罪的敏感度.
參選是一個很好的例子.
我們反思我們的生命.
其實我們生活裡.
每每都犯同樣的錯.
總覺得我們自己可以駕馭這些試探.
例如你可能會想.
我對我的婚姻很有信心.
我和同事之間去應酬.
我去上面玩玩.
OK的.
工作上.
我對條例很熟悉.
偶爾踩踏灰色位.
不代入法網.
有什麼問題?.
OK的.
我對我的信仰很有信心.
我老闆.
拜滿天神佛.
帶我去看風水.
又要我設置.
辦公室的位置.
要我設置風水陣.
OK的 影響不到我.
我和我男朋友.
都是很好的基督徒.
我們兩個一起去旅行.
同一間房.
OK的.
指圖保羅曾經在哥林多前書五章六節.
告誡哥林多的弟兄姊妹.
他說:你們這自誇是不好的.
豈不知一點面孝.
能使全團發起來嗎?.
其實罪就像面孝.
那些孝母一樣.
只要放一點點進去.

$^{721}$一個小小的麵團.
就可以發到很大.
因此我們不要過度自信.
覺得我們自己可以凌駕一點點的罪.
的人有.
過度自信的參算.
就是容許自己的生命.
這樣發孝.
那些罪在裡面發孝.
最終他抵不過.
亞特里拉天天這樣去扯他.
在第十七節說.
他真的將上帝的交付.
給他的事實.
全盤托出.
他說:若果你剃我的頭髮.
那些力氣就會離開我.
我會變成和普通人一樣.
參宣闖下了一個難以逆轉的慘痛經歷.
我們一起看看.
第20至21節說.
這裡說到亞特里拉.
真的這麼做.
過度自信的參宣.
以為即使他頭髮被剪了.
他仍然能力大無窮.
所以當他醒了的時候.
他在第20節說.
我心想.
我要像之前那幾次.
去活動一下身體.
他不知道原來耶和華已經離開了他.
他的能力其實與頭髮無關.
而是神賜給他的.
所以當神的靈離開他的時候.
他這些特別能力就沒有了.
在第21節.
菲律賓人就抓住他.
然後刮了他的雙眼.
帶他去加薩.

$^{761}$用一個銅鏈綁住他.
然後就把他放在監獄裡推磨.
原本很有自信.
力大無窮的參宣.
現在時移世易.
形勢大逆轉.
推磨是一個什麼工作?.
其實推磨是一個很辛苦.
而且是當時最低等的奴役工作.
一個大男人.
做一些女人平時做的工作.
而參宣也被帶到加薩.
被稱為強壯的地方.
被人任人魚肉.
戲弄.
變得軟弱無能.
其實對參宣來說是一個極大的諷刺.
第25節就說菲律賓人.
在那裡宴樂的時候.
就叫參宣出來.
就說要把參宣帶出來.
好像是娛樂他們.
戲弄他們.
給他們一點開心.
當時這班菲律賓人的領袖.
其實包括首領.
聚集在一個.
當時的神明叫大滾.
Dagon一個寺廟裡.
他們很開心.
因為捉到一個大敵人.
於是他們就好像在開派對.
屋頂大概.
不止這些.
下面有人.
第28節就記載著.
參宣在聖經裡的第二次祈禱.
他這次很特別.
他稱呼神為主耶和華.
以往他從來都沒有這樣去稱呼.

$^{801}$用這麼敬畏的稱呼去敬畏耶和華.
去呼求神.
他就在說.
求你眷顧我.
求你賜我這一次的能力.
其實參宣去到窮途末路.
他終於都清楚知道.
他的能力完全是出於神.
他不再自以為是.
他終於都學習去謙卑.
倚靠耶和華.
祈求神再一次給他能力.
不過我們要細心留意一下.
其實他的動機是什麼?.
其實他是想求神.
使他在非利人身上報仇.
報他被人挖了雙眼的仇.
最後29至30節說.
參宣左右手.
各抱著一支柱.
托著房的柱子.
他屈身將房子倒塌下來.
使他與非利人同歸於盡.
這就是參宣的一生.
參宣最後所關心的.
仍然都是為了報仇.
他從來都沒有真正.
履行過神所託付給他的使命.
雖然上帝是無所不能的.
神的工作最後都能夠達成.
不過參宣的終局.
不禁令人感到很慨嘆.
他原本是神所設立.
選定他成為終生的拿撒爾人.
背負著以色列人脫離非利人.
的一個重大使命.
本應得到上帝的祝福.
可惜他沒有珍惜這份恩典.
他反而恃著自己力大無窮.
過度自信.

$^{841}$最終悲劇收場.
有時我會想.
如果參宣能夠早一點回轉的話.
他的生命,人生的終局.
會不會這麼悲慘呢?.
參宣的人生我們無法掌握.
我們控制不了.
但我們自己的人生.
我們的選取.
是在我們手中.
神賜給我們各人各種恩賜.
有各種能力.
讓我們活出榮耀上帝的生命.
今天神給了你們什麼恩賜呢?.
他又託付給你們什麼使命呢?.
你願不願意謙卑順服上帝的心意.
看自己合乎忠道.
在我們各人的事奉崗位上.
敬虔道逸.
去全著敬畏的心去榮耀.
我們榮耀上帝的生活.
我們一起同心禱告.
全行的恩主.
當我們還作罪人的時候.
你就已經命定要拯救我們.
你洗淨我們的污穢.
引導我們走正義的路.
聖經告訴我們.
人若離開了你之後.
其實我們不能作什麼.
我們不得不承認.
我們只是一群過度自信的罪人.
經常忘記了.
上帝你才是賜予恩典的神.
不是我們有多厲害.
不是我們多有才幹.
而是你白白的恩典.
賜予我們能力可以和你同功.
求你賜我們一份謙卑受教的心.
幫助我們去珍惜.

$^{881}$上帝賜給我們的每一份恩典.
不要當作理所當然.
好叫我們的生命成為你美好的活祭.
見證你 榮耀你.
以上坦誠真誠的禱告.
乃奉我恩主.
主耶穌基督得勝的名字.
其他事項.
\newpage



\section{士師記 17:1-13}
\label{sec:jXXnWoB7EBk}
\textbf{2023年9月24日崇拜直播｜陳明泉牧師｜「家醜說出口」系列:私有化的祭司｜士師記17\_1-13}
\newline
\newline
連結: \href{https://youtube.com/watch?v=jXXnWoB7EBk}{\texttt{https://youtube.com/watch?v=jXXnWoB7EBk}} ~~~~ 語音日期: 2023-09-26
\newline
\newline
\hyperref[sec:rllwPQsYaK0]{\small{< < < PREV SERMON < < <}}
~
\hyperref[sec:index_chronic]{\small{[返順時目]}}
~
\hyperref[sec:_4zbsWU6e2A]{\small{> > > NEXT SERMON > > >}}
\newline
\newline
士師記 17:1-13
\newline
\begin{longtable}{cl}
\hline
\hline
章節 & 經文 (和合本修訂版)\\
\hline
17:1 & \begin{tabularx}{0.7\textwidth}{X} 以法蓮山區有一個人，名叫米迦。 \end{tabularx} \\ \\ \relax
17:2 & \begin{tabularx}{0.7\textwidth}{X} 他對母親說：「你的一千一百塊銀子被人拿走了，為此你發咒起誓，也說給我聽。看哪，銀子在我這裡，是我拿的。」他母親說：「願我兒蒙耶和華賜福！」 \end{tabularx} \\ \\ \relax
17:3 & \begin{tabularx}{0.7\textwidth}{X} 米迦把這一千一百塊銀子還他母親。他母親說：「我把這銀子分別為聖，親手獻給耶和華，為我兒子造一尊雕刻的像，以及一尊鑄成的像。現在我把銀子交給你。」 \end{tabularx} \\ \\ \relax
17:4 & \begin{tabularx}{0.7\textwidth}{X} 米迦把銀子還他母親，他母親把二百塊銀子交給銀匠，去造一尊雕刻的像，以及一尊鑄成的像，安置在米迦的房子裡。 \end{tabularx} \\ \\ \relax
17:5 & \begin{tabularx}{0.7\textwidth}{X} 米迦這個人有了神堂，又造了以弗得和家中的神像，派他的一個兒子作祭司。 \end{tabularx} \\ \\ \relax
17:6 & \begin{tabularx}{0.7\textwidth}{X} 那時，以色列中沒有王，各人照自己眼中看為對的去做。 \end{tabularx} \\ \\ \relax
17:7 & \begin{tabularx}{0.7\textwidth}{X} 猶大的伯利恆有一個年輕人，是猶大族的人。他是利未人，寄居在那裡。 \end{tabularx} \\ \\ \relax
17:8 & \begin{tabularx}{0.7\textwidth}{X} 這人離開猶大的伯利恆城，要找一個可住的地方。他來到以法蓮山區米迦的家，還要往前走。 \end{tabularx} \\ \\ \relax
17:9 & \begin{tabularx}{0.7\textwidth}{X} 米迦對他說：「你從哪裡來？」他說：「我從猶大的伯利恆來。我是利未人，要找一個可住的地方。」 \end{tabularx} \\ \\ \relax
17:10 & \begin{tabularx}{0.7\textwidth}{X} 米迦說：「你就住在我這裡吧！我以你為父為祭司，每年給你十塊銀子和一套衣服，以及生活所需的食物。」利未人就來了。 \end{tabularx} \\ \\ \relax
17:11 & \begin{tabularx}{0.7\textwidth}{X} 利未人願意和這人同住；他待這年輕人如自己的兒子一樣。 \end{tabularx} \\ \\ \relax
17:12 & \begin{tabularx}{0.7\textwidth}{X} 米迦授這年輕的利未人祭司的職任，他就住在米迦的家裡。 \end{tabularx} \\ \\ \relax
17:13 & \begin{tabularx}{0.7\textwidth}{X} 米迦說：「現在我知道耶和華必恩待我，因為我有利未人作我的祭司。」 \end{tabularx} \\ \\
[1ex]
\hline
\hline
\end{longtable}
$^{1}$經訓時間.
陳偉業議員.
第一節.
二百年山地有一個人名叫米家.
他對母親說.
你那一千一百赦赫勒銀子被人拿去.
你因此就坐.
並且告訴了我.
看啊!這銀子在我這裡.
是我拿去了.
他母親說.
我兒啊!願耶和華賜福與你.
米家就把這一千一百赦勒銀子還他母親.
他母親說.
我分出這銀子來為你獻給耶和華.
好雕刻一個像.
鑄成一個像.
現在我還是交給你.
米家將銀子還他母親.
他母親將二百赦勒銀子交給銀像.
雕刻一個像.
鑄成一個像.
安置在米家的屋內.
這米家有了神堂.
又製造以弗得和家中的神像.
分派他一個兒子作祭司.
那時以色列中沒有主.
各人任意而行.
猶大的伯利恆有一個少年人.
是猶大族的利未人.
他在那裡寄宿.
這人離開猶大伯利恆城.
要找一個可住的地方.
行路的時候.
到了以法連山地.
走到米家的家.
米家問他說.
你從那裡來?.
他回答說.
從猶大伯利恆來.

$^{41}$我是利未人.
要找一個可住的地方.
米家說.
你可以住我這裡.
我以你為父為祭司.
我每年給你十赦赦勒銀子.
一套衣服和度日的食物.
利未人就進了他的家.
利未人情願與那人同住.
那人看這少年人.
如自己的兒子一樣.
米家分派這少年人的.
利未人作祭司.
他就住在米家的家裡.
米家說.
現在我知道耶和華.
必賜福於我.
因我有一個利未人作祭司.
願神賜福常讀他話語.
並且遵行的人.
丁姐妹平安.
平安.
我們今天繼續來看時事記.
之前就說過.
在時事記裡面記載著.
有不同的時事.
不同的領袖興起.
到參宣就是時事記裡面.
記載的最後一位時事.
但時事記並不是停止了.
在參宣之後就停止了.
沒有時事興起.
就等於整個書卷的結束.
到第十七章.
就開始進入一個新的時代.
就是後時事時代.
後時事時代是甚麼意思呢?.
就是沒有時事興起.
如果你明白以色列.
或舊約裡面.

$^{81}$有三個角色身份是很重要的.
首先第一個.
利用時司或君王.
即是成為統治帶領.
或打仗的領袖.
先知就是上帝的代言人.
上帝直著他.
將上帝的話語帶到整個群體當中.
祭司或大祭司.
就是帶領以色列民.
一來教會他們如何做教導.
他們如何敬拜上帝.
另一方面就是主理整個聖殿.
在會幕裡面的敬拜.
上帝選召利未人成為了這個角色.
在這個支派裡面.
選召一些人來參與.
聖殿或祭祀的工作.
除了這個祭祀的宗教儀式之外.
利未人亦要負責平時律法的教導.
都要幫助以色列整個群體.
都要明白上帝的心意.
這三個角色是舉足輕重的.
不過當你一邊看的時候.
事司已經是最後一個了.
當時以色列又末又王.
所以就沒有了這樣的領袖.
當以色列人咬牙切齒.
或者我們看上去.
搞錯呀,這樣的事司都可以的時候.
接著就連事司都沒有了.
上帝的說話已經越來越稀少.
所以不是那麼多先知來代表上帝.
或者從上帝直掌一個人來說話給整個群眾.
整個以色列民知道.
最後剩下的就是祭祀.
而在接下來的篇幅就記載著兩個利未人的故事.
在後世詩時代裡面.
最後把關了.
在以色列民族裡面的這些祭祀或者利未人.

$^{121}$他們在他們的群體當中.
究竟能否發揮作用呢?.
整個以色列民族走向一個什麼樣的地步呢?.
到了第十七章我們開始.
接下來這幾次的講道.
我們都是用這幾個.
用兩個祭司的故事和最後的結語.
來講究竟整個以色列走向一個什麼樣的地步.
讀上去都是有一點負面的.
有一點暗黑的.
不過負面的經歷也都希望為我們帶來正面的提醒.
我們先有個禱告.
求上帝幫助我們開啟我們的心靈.
同心祈禱.
天父上帝願你在我們當中帶領我們眾人.
讓我們能夠明白你的話語.
願你此時此刻.
藉著你的話語.
將你的心意對我們來說話.
不單止讓我們投入認識.
更加對著我們的心靈說話.
調教我們.
更新我們.
讓我們所作的被你閱立.
我們祈求你在我們當中繼續來到瑞領.
獻上我們的禱告.
奉靠基督耶穌得勝的聖名祈求.
Amen.
我們看第十七章.
剛才玉峰弟已經幫我讀了.
其實第十七章第一節到第二節.
你讀上去.
如果就這樣平鋪直敘讀.
你讀不到那個味道.
你想像一下周星馳的畫面.
你就會讀得這麼奇怪.
這段經文開頭講什麼呢.
以化連山地裡面有個人叫米加.
有一天米加跟他媽媽說.
媽媽那一千一百射克勒銀子被人偷了.

$^{161}$你不是要詛咒那個人嗎.
我現在告訴你.
你詛咒的人就在你眼前.
那個就是我.
那一千一百射克勒就在我那裡.
整件事你覺得是不是.
爛爛的.
就像周星馳一樣.
你詛咒那個人其實是你的兒子.
所以現在錢就在我那裡.
他媽媽看到這個情景.
更加周星馳.
媽媽說了什麼話呢.
兒子啊,願耶和華賜福給你.
開頭是說之前就坐的兒子.
現在就說願耶和華賜福給你.
他想著可以傷底.
可以level(等級)了.
將祝福可以平減了那個就坐.
接著最荒唐的地方就是.
米加就將這一千一百射克勒還給他媽媽.
他媽媽就說了.
我分出這些銀子幫你獻祭.
獻給耶和華.
那你怎麼做呢.
我做一個像給耶和華.
雕刻一個像,鑄成一個像.
那我就還給你了.
那一千一百射克勒已出之物.
不見了,就走了就算了.
現在我就用這些錢來做一個像給你.
米加就不要.
於是他媽媽就強行這樣做.
於是就將二百射克勒交給一個銀像.
用銀或者貴金屬來鑄造一個耶和華的像.
經文沒有特別說為耶和華鑄像.
不過你上文下理就知道了.
為耶和華鑄成一個神像.
安置在米加的家裡.
所以米加家裡就有個神堂.

$^{201}$米加要擴大這個神堂.
在存放這些神像.
就會製造一些以弗德.
就是祭司的袍.
還有在家裡放其他的神像.
另外在他家裡分派一個兒子.
作為打理這個神堂裡的祭司.
這個就很周星馳.
一個你覺得讀上去很奇怪的.
為什麼以色列出現這些東西呢.
首先第一件事就是.
耶和華不可以做像.
敬拜上帝.
如果你真的做錯了事.
不是你做一個像就可以平減了.
或者抵消了之前做錯事的方式.
但是這個媽媽就說做一個像.
就好像做一些功德.
給米加有一個像來拜拜.
就令到他可以平息了上帝的怒氣.
律法上說到不可為耶和華做像.
第二件事.
就坐和祝福.
很輕易的.
隨口說的.
在他媽媽裡面.
就坐罵到死的時候.
罵正就坐自己的兒子.
然後說一句願耶和華賜福給你.
就好像搞定了整件事.
第三件事.
有些很低調的你不知道.
我們以為後契斯時代沒有一個領袖.
整個以色列文就連耶和華上帝都忘記了.
不過你看上文下理.
從第17章開始.
不停提.
提上帝的耶和華的名稱.
甚至乎剛才米加的媽媽改名.
都很神聖的.

$^{241}$米加的名字你知不知道什麼意思.
瑞章耶和華.
這個是一個很謙卑.
一個很神聖的名字.
在聖經裡面有第二個人.
都是用這個名字.
是先知來的.
叫做梅加書裡面的梅加先知.
同樣一個字.
你想一下米加.
他的名字是他媽媽覺得.
誰像耶和華上帝一樣.
是讚嘆上帝.
謙卑苦伏的一種.
稱頌上帝的表現.
在他的生活當中.
不斷的提及耶和華的名字.
所以即使兒子做錯事.
他的心是宗教性的.
是想做一些事情為上帝而做的.
不過他沒有想過.
他做的那些事情完全違反律法.
在第六節.
聖經的作者都怕我們有誤會.
他告訴你.
那時以色列沒有王.
各人任意而行.
之前說過了.
任意而行是什麼意思呢.
那個任意的意思不是說.
我感覺到了.
好像藝術家.
好像教師一樣.
做一些行為藝術的事情.
任意的意思是.
那些人做他們眼中看為正的事.
那個正或者邪.
善與惡.
是從做的人的角度去看的.
不是從律法角度去看.

$^{281}$所以在以色列當中.
後事詩時代.
所有人都提及上帝的名字.
提及耶和華的名字.
按著自己的心意去做敬拜.
甚至他們建立自己的神堂.
在自己的家中擺一個偶像.
多方便.
隨時隨地.
傳天后敬拜上帝.
一想起就不用去示羅.
去回望當中敬拜上帝.
在家中首到拿來.
為了更方便他們整個祭祀.
或者整個祭祀儀式.
更加家庭化.
米加就分派他自己的兒子.
祭祀本來是上帝選召的一班子民.
指定要利未人來做的.
如果說米加大概是二法蓮支派的一個人.
住在二法蓮山地那裡.
他們的宗派不是擔演這個角色的.
不過他們為了.
他們口中常常敬拜上帝.
他們不跟律法去做.
就按自己的心意.
按自己覺得敬虔的方式.
去做敬拜上帝的一些事宜.
在這裡就有一個很荒謬的開始.
接著去到第七節.
再繼續說.
猶太裡面有一個伯利恆的少年人.
是猶太族的利未人.
他在那裡寄居.
這個人離開了猶太的伯利恆城.
要找一個可住的地方.
行路的時候去到二法蓮山地.
走到米加的家.
我們看到第八節.
第七節第八節.

$^{321}$在這裡你值得留意的地方.
他講到一個真是負責上帝選照的利未人.
祭司的這個支派.
他們就四散.
住在不同的地方.
當然在這裡所講的四散.
或者他們寄居在不同的地方.
為什麼會這樣做呢.
失業.
當時的人因為按著自己的心意.
任意妄為.
所以基本上回望的敬拜是沒有制度的.
而舊時專職在回望裡面敬拜的那些利未人.
沒事做.
沒人願意奉獻.
沒人願意承擔整個敬拜中心的方法.
那些利未人舊時都是靠著那些祭肉.
或者那些祭身.
或者那些餅.
來支持他的生計.
當沒人去這樣做.
每個人自己在家裡建立自己的神堂的時候.
利未人就沒辦法來生存.
於是他們就四處漂泊.
找一個可住的地方來寄居.
這個是利未人在當時已經失業.
更加重要的是利未人跟其他支派不同.
在《申命記》第十八章說到.
利未人跟其他支派有些不同的地方就是.
祭司利未人和利美全支派.
必在以色列中.
無份無業.
這班人是無法分地的.
他們所吃用的就是獻給耶和華的火祭和一切所捐的.
他們在弟兄中活有產業.
耶和華就是他們的產業.
正如耶和華所應許他們的.
利未人進入迦南地.
《申命記》裡面摩西很清晰說到.
這個支派的人.

$^{361}$他們的生存就是靠賴整個以色列民.
大家真的願意來獻祭.
來供養這班人.
當這班人各自已經.
雖然自己很強調自己的信仰.
但是他們對敬拜.
對於律法的要求.
他們完全既一無所知.
也沒有想過要實踐.
所以一班本身沒有產業的利未人.
他們就要四處漂泊.
經文裡面再繼續說.
這個年輕人.
有些色經學家說25至30歲.
年輕人來的.
一個少年人來的.
一直走的時候.
去到第九節.
米加就問他.
你在哪裡來的.
他回答說.
從猶大的伯利行來.
利未人要找一個可住的地方.
即是寄居.
行乞的方式.
來過他的生活.
米加說你可以住在我這裡.
我以你為我的父.
為祭司.
每年給你十射克勒的銀紙.
一套衣服.
和到日的食物.
利未人就盡了他的家.
利未人情願與乃人同住.
乃人看這少年人.
如同自己的兒子一樣.
米加分派這少年的利未人作祭司.
他就在米加的家裡.
米加說現在我知道.
耶和華必賜福於我.

$^{401}$因我有一個利未人作祭司.
在米加的眼中.
這個年輕利未人漂泊流離.
沒有地方住.
收留了他.
不過當他講到他是一個利未人.
聖經裡面說的是專職負責祭祀的角色身份.
於是就提出一些條件.
你可以住在我們這裡.
還要說得很重金禮聘.
你可以作我的父.
為祭司.
我每年給你十射克勒銀紙.
和一套衣服.
和到日的食物.
在這裡.
你是不是最優厚呢?我不知道.
不過總比四處遊蕩好.
對於這個流離失所的少年利未人來說.
能夠有一戶收留他.
又有人工鬥.
又有衣服.
又有食物.
總是一定比四處流浪好.
他沒有想過.
他的身份.
他的支派.
是供奉上帝聖所的支派.
他只想眼前的生活.
於是這個利未人就很甘心情願.
和這個人住了.
米加真的對少年的利未人很好.
在經文裡面第十二節.
用同樣的字眼.
米加分派這個少年人.
分派這個利未人作祭司.
住在米加的家裡.
同樣的字眼.
剛才說到.
米加分配他其中一個兒子做祭司.

$^{441}$現在就是分派一個利未人.
作為他家神堂的祭司.
然後米加有一個結論.
為了整個安排.
我現在才知道.
媽媽說的那個祝福是什麼.
知道耶和華必賜福給我.
我知道什麼原因.
因為有一個利未人.
在這裡做祭司.
多幸福.
最方便的位置.
又有他自己最合乎心意的敬拜方式.
所有東西都是朝著他的心意.
所有東西都是按著他的心意.
整個敬拜.
按著他的心意來建立起來.
好像打麻將一樣.
儲牌一樣.
儲著儲著.
快要吃糊了.
吃了糊的時候.
哇 我知道.
上帝真的賜福給我了.
我有一個利未人.
現在如虎添翼.
所有東西都齊了.
很開心.
不過在這裡.
我們值得留意.
利未人是不是成為一個家裡面的祭司呢?.
利未人是成為祭司.
不過這個祭司不是屬於米加的家裡.
利未人成為的是神的家裡的祭司.
他的生命.
他的供應.
誰是他出糧給他的呢?.
在申命記中說的是耶和華出糧給他.
耶和華是利未人的產業.
上帝成為了他的老闆.

$^{481}$不過在這裡.
一個按著自己的心意來建立的神堂.
在敬拜上已經是錯誤了.
在一個這樣的錯誤基礎裡面.
再接受了這個家庭裡面的供養.
成為這個家庭裡面的私有化了的祭司.
其實在本質上他已經將自己的即時扭曲了.
經文裡面不是停在這裡.
我們還要看的就是第十八章.
那個才是今天裡面我們最需要看的重點.
在第十八章一至六節.
經文再一次提醒我們.
那時以色列沒有王.
但之派的人仍是尋地居住.
因為到那日子.
他們還沒有在以色列之派中得地為業.
但人在所拉和以實圖打發本族的五個勇士.
去仔細窺探那地.
吩咐他們說.
你們去窺探那地.
他們來到以法連山地.
進了米家的住宅.
就在那裡住宿.
他們臨近米家的住宅.
聽出那少年你美人的口音來.
就進去問他說.
誰令你這裡來?.
你在這裡作什麼?.
你在這裡是什麼?.
他回答說.
米家待我如此如此.
讓我作祭司.
他們對他說.
請你求問神.
使我們知道所行的道路通達不通達.
祭司對他們說.
你們可以平平安安地去.
你們所行的路是在耶和華面前的.
又再出現耶和華的字眼.
尋問耶和華.

$^{521}$不過這次的尋問不是米家來尋問.
當時十二支派中其中一個支派叫做旦.
如果你記得.
我們很早的時候說過.
第一章第34節.
說不同支派進入郊民地打仗的時候.
各個支派有些打贏仗.
有些打贏但不趕出那些外邦人.
而成績最差的支派就是旦支派.
他怎樣差呢?.
他想趕出阿摩利人.
但阿摩利人反而趕不走.
還要趕他們下山腳.
趕他們下山.
總之趕盡殺絕的就是這班旦的人.
去到十二支派某程度.
去到舌斯記這個時間.
大家都已經落地生根的時候.
經文說到旦這個支派.
某程度還是覓地居住.
尋找一個可住宿的.
一個整個支派可以安身立命的地方.
於是他們就四出來窺探.
看看有什麼地方可以做.
昔日的窺探是說.
以色列人的窺探外邦地.
現在他們看什麼呢?.
看自己有.
看其他支派的地方.
想想有沒有地方可以讓他們住.
這五個勇士就去到米加的家.
聽到這個少年利未人的口音.
覺得他不是本地人就問他.
你來這裡做什麼?.
你是誰?.
你在這裡做什麼?.
然後很簡單的一句回答.
米加待我如此如此.
讓我在這裡作祭司.
他成為這個家庭裡面的祭司.

$^{561}$以色列人不是不知道.
有時來有一個回眸.
在米加和時來其實兩者相距的步程.
不是太遠.
不過就在一個這樣的狀況之下.
這五個窺探這個地方的勇士就問這個祭司.
問這個少年人.
喂,幫我問一下耶和華上帝.
求問神.
讓我們知道.
走我們的道路通達不通達.
其實簡單來說就是說.
他沒有告訴我們我們要做什麼.
他只是想說他做的事情是否順利.
於是這個祭司不問情由.
他其實說得比較俗一點.
找飯吃而已.
沒有必要結怨.
所以有人來問他.
人家問他信不信你.
諸事大吉.
平平安安回去.
耶和華的道路就在你面前.
你所行的道路就在耶和華面前.
但是他沒有想過.
其實這五個人就是想吞併米加和他所住的這個二法蓮山地的地方.
糊里糊塗.
盲目的說平安的話.
就希望打發了這些人離開.
也不用為自己製造一些麻煩.
而這個祭司也不是存心來尋問上帝.
上帝在他的口中裡.
只不過是一個口頭禪.
他不停提及上帝.
但是他的生命內容.
他的信仰內容是沒有上帝的.
在這裡我們稍微思想一下.
整個去到這一部分的時候.
我們想我們今天的信仰觀念.
或者我們今天實踐的信仰.

$^{601}$有時候會不會也類似在後世詩時代的一群以色列人.
甚至是祭司呢?.
我們口中可能說很多上帝.
提及很多上帝.
但是我們心裡只是希望上帝按照我們自己的議程.
按照我們自己的心意.
保佑我.
賜福給我.
將我所要的東西加給我.
我覺得這個就是上帝的祝福.
曾經有一個我聽過一個最荒誕的.
不是旺朕的,大家放心.
一個很荒誕的說法是什麼呢?.
有個弟兄姊妹買六合彩.
祂祈求上帝.
上帝啊,你給我六個號碼.
六個就夠了.
就是給我能夠拿到這六個號碼.
最多我十分之三來奉獻給神的家.
在神的家風足有餘.
養活陳明泉牧師一家.
教會現在裝修沒錢.
我求六個號碼吧.
當然這個你覺得很荒謬.
不過有時候我們的信仰運作是這樣的.
調換了不是說六合彩.
可能是說其他你生命裡很想追求的東西.
你希望上帝幫助你.
你希望上帝滿足你.
你不會問那件事.
最終你所追求的東西是否合乎聖經.
你只是很想得到.
上帝就成為了一個你想求祂.
威逼你有賄賂的方式.
來滿足自己的心願.
其實這種信仰心態.
和坊間民間信仰相差無幾.
只不過是將標籤轉換了.
由昔日的某個神明撕掉.
現在就貼了另一個.

$^{641}$就是耶和華上帝.
但是敬拜的態度.
或者那種信仰的生命.
內涵是一樣的.
不過在聖經裡面提醒我們.
我們敬拜上帝從來不是說.
要上帝滿足不是天從人願.
人心裡面所想到的祝福.
上帝要滿足就代表著一個敬拜的生活.
相反來說.
我們以上帝為主.
上帝要我們所作的.
我們按著上帝心意而行.
我們如何去明白上帝呢?.
既需要有人來解釋.
或者將上帝的話語傳遞給我們.
但同時間我們也要看.
成為點注的聖經.
聖經很清楚將上帝的心意.
來表明的時候.
我們按著這個心意.
來實踐我們的信仰.
而不是在說我們.
貿貿然我們自己喜歡怎樣就怎樣.
來過我們信仰的生活.
我們繼續看經文.
在經文裡面.
第十八章第七節之後.
剛才說的五個勇士離開了.
去到拉翼.
見那裡的人安居無慮.
如同西頓的人安居一樣.
那裡是很有錢.
安居樂業的地方.
那裡也沒有掌權擾亂他們.
他們離西頓的人也遠.
與別人沒有來往.
於是這五個人就回到.
所拉和以實圖.
見他們的弟兄.

$^{681}$簡單來說就是.
回到自己的根據地.
與他們自己的村民.
以實圖和所拉.
他們自己的根據地.
來說現在要攻打這個地方.
我們窺探過那個地方.
一切都很好.
大家不要靜態地坐著不動.
現在就速速前往.
於是他們就招聚了多少人呢.
十一節裡面.
他們就招聚了民族.
氮族裡面的六百人.
各帶兵器.
從所拉到以實圖.
來到前往.
去到猶大的基烈耶林.
在基烈耶林那邊安營.
然後就一路一路.
逆步逆趨.
去到以法連山地.
去到米家的住宅.
然後經民的發展.
其實是荒謬中的荒謬.
記載著一個更加荒誕的故事.
當現在不是五個勇士.
而是六百人.
去到米家的家裡.
之前的那五個勇士.
進到米家的家裡.
他第一件事要做什麼呢.
第十六節.
六百人各帶兵器.
站在門口窺探地的五個人走進去.
將雕刻的像.
以佛德家中的神像.
並著成的像都拿去.
祭司和帶兵器的六百人.
一同站在門口.

$^{721}$你看到這六百人第一件事拿什麼呢.
他拿米家自己.
周星馳式的那幾個神像.
當然當時的神像是值錢的.
是貴的東西.
不過他第一件事要拿的就是.
拿這幾個神像.
然後看到少年人的祭司.
這個祭司正在吃米家的糧.
看到他拿了那些像.
將他的生財工具都拿走了.
然後這個利未人就問.
你們在做什麼呢.
那五個勇士就回答他.
不要出聲.
掩著嘴巴.
跟我們一起走.
第十九節那裡.
我們以你為父為祭司.
你作一家的祭司好嗎.
還是作以色列一族一支派的祭司好呢.
一個極之荒誕的提請.
以色列人偷神像.
已經是一個很荒誕的舉動.
然後還要高薪劃國.
那個祭司.
他還要給一個條件.
在這裡沒有前途.
作一家的祭司.
現在是說一族一支派的祭司.
整個支派的祭司更加有榮耀.
這個祭司聽到之後有什麼反應呢.
二十節.
祭司心裡喜悅.
便拿著以弗德和家中的神像.
並雕刻的像進到他們中間.
一個極大的諷刺.
讀上去會覺得更加咬牙切齒.
簡單來說一個沒有腰骨的祭司.
哪裡有飯吃就是哪裡.

$^{761}$他就去哪裡了.
當有人來規探.
接著要佔領這個地方.
接著有好路數.
於是乎他就情願跟著這班人.
即是說他偷了米家之前老闆的神像.
又如何呢.
現在有一個更加好路數.
他幫手拿.
神像很重的.
一起來一起搬走.
一起來跟著彈支派整個民族.
成為彈支派裡面的祭司.
二十七節到三十一節.
是今天讀的最後一個段落.
但人將米家所作的神像.
和祭司都帶到拉易.
見安居無慮的民.
就用刀殺了那民.
放火燒了那城.
並無人答救.
因為離西遁遠.
他們又與別人沒有來往.
城在平原.
那平原靠近的是伯利俠.
但人又在那裡修成居住.
照著他們始祖以色列的子.
丹的名字.
給那城起名叫丹.
原先那城名叫拉易.
但人就為自己設立那雕刻的像.
摩西的孫子.
格遜的兒子約拿丹.
和他的子孫作彈支派的祭司.
直到那地遭擄掠的日子.
神的殿在是羅多少日子.
丹人為自己設立米家所雕刻的像.
也在丹多少日子.
十八章記載著一個.
既是令人髮指.

$^{801}$但同時間也是很唏噓的一個場面.
這班丹的人.
他們要來窺探這個地方.
去到拉易這個地方住宅.
見到那個地方那些人手無寸鐵.
那些人又安居無慮.
於是就清殺那裡的人.
用火燒城.
然後就用他們自己的方法.
重新建立自己的丹這個城市.
建立了之後他們就有自己的地方.
不單止這樣.
丹支派也有自己敬拜的方式.
就按照之前住在米家裡的少年利未人.
在聖經裡他特別的再說.
本來身份尊榮的這班祭司的後裔.
摩西的孫.
格遜的兒子約拿丹.
和他的子孫.
現在成為了丹支派裡的這個祭司.
這個祭司是直到老弱之日.
神的殿在是羅有多久.
丹的人就為自己設立米家所雕刻的像.
在丹又有多久的日子.
這是一個荒謬的故事.
一個屬神的祭司.
去到走投無路.
到最後成為了幫他們敬拜.
一個以自己為首建立的神明.
米家的神像.
成為了他們心目中的耶和華上帝.
米家自己建立的這個家庭祭司的制度.
後來成為了丹支派的一個家庭制度.
當米家覺得有一個利未人成為他家的祭司的時候.
他知道就明白上帝怎樣賜福給他.
但他忘記了.
其實他之前偷了一個1100石赫勒.
他媽媽去咒詛他.
他沒有想過他做了這個.
是令他整個地方.

$^{841}$整個家庭被咒詛的一件事.
今天這段經文我講到這裡.
這裡讀上去你心裡會很不安樂.
一個被選召的上帝的子民.
成為了一班不堪入目的一班人.
我在讀的時候我從兩個角度去想.
第一從祭司的角度去想.
祭司的角度其實都是反映我們自己傳道和公.
我們今天作為神的公人.
我們的身份有一個什麼位置.
我們今天我們是不是有腰骨成為一個宣講上帝真理的人.
我在這裡勉勵自己.
我也是自我批判.
我會不會為了十個石赫勒而折腰呢?.
會不會為了更加好的路數.
作為旺角浸信會的祭司也好.
或作為另一個更加大堂會的祭司也好呢?.
表面上很吸引.
但是如果我當這個祭司或當這個目的是一份工作.
我真的可以很高無怨.
但是如果這個是上帝的召命.
如果這個是上帝給我們一個屬神的家.
我要上帝將我生命裡面.
我要服侍上帝成為我首要的時候.
其他的錢財所有的東西都已經不是我考慮的因素.
不過今天我很坦白說.
有些同道有時候都跟我說.
有時候他們在其他教會.
不是王浸.
有些其他教會服侍的同道說.
有時候是憤氣激氣.
接著就會說一句.
不過他自我安慰說.
不要緊啦.
我們都不是牧羊的了我們這班人.
那不是牧羊做什麼?.
那些是客人.
那些是客人我們要照顧的.
那我說你是牧者是什麼?.
他說Sorry,Anders不要說我們是牧者.

$^{881}$我們現在叫做Pastoral Service Provider.
牧羊服務提供者.
那班客人我們要照顧他們快樂.
令他們每一個星期有賓至如歸.
他們不開心就像聽顧客熱線一樣.
一聲吞下.
做牧者做到這樣很辛苦.
接著說為了兩餐不是這樣又如何?.
有時候會這樣.
除了牧者祭司這個角色之外.
其實祭司走到走投無路.
整個以色列群眾都有責任的.
當他們不按照律法來敬拜上帝.
當他們不是以上帝的心為心.
不是以真正的耶和華為他的耶和華上帝.
而是以他心目中的偶像成為了他的主.
到最後這班回眾都是一班客人.
這班回眾到最後只是來教會.
逢星期日周末各取所需.
拿自己心目中所拿的東西.
又或者過他們自己覺得理想中.
他自己心目中的信仰生活.
有本書的作者叫William William.
他講得很好.
今日這個牧羊關係是甚麼關係?.
是一個從世俗走出來的關係.
教會的目的不要把世俗觀念放進來.
我們要和這個世界貼地.
不過怎樣貼都好.
不是等同這個世界.
我們知道這個世界發生甚麼事.
但我們不是按照這個世界的軌道來走.
傳道人的角色不是成為一個故宮.
不是要提供牧羊服務.
不是成為一個服務者.
我們是服侍上帝.
同樣教會的群體.
我們是要成為一個以尊敬上帝為主的群體.
失去了這一切.
其實我們和米家神社.

$^{921}$米家自創有甚麼分別?.
沒有了律法.
沒有了牧者.
各人只是行自己眼中覺得為正的事.
今日這一段經文.
看上去我們覺得咬牙切齒.
但同時間都光照著我們.
今日我們怎樣看教會的生活.
怎樣看我們敬拜人生的生活.
求主繼續對我們講說話.
更新我們 調教我們.
\newpage



\section{士師記 19:1-30}
\label{sec:_4zbsWU6e2A}
\textbf{2023年10月1日崇拜直播｜陳明泉牧師｜「家醜說出口」系列:沒有上帝的麻木｜士師記19\_1-30}
\newline
\newline
連結: \href{https://youtube.com/watch?v=-4zbsWU6e2A}{\texttt{https://youtube.com/watch?v=-4zbsWU6e2A}} ~~~~ 語音日期: 2023-10-03
\newline
\newline
\hyperref[sec:jXXnWoB7EBk]{\small{< < < PREV SERMON < < <}}
~
\hyperref[sec:index_chronic]{\small{[返順時目]}}
~
\hyperref[sec:e5U96ffu0L0]{\small{> > > NEXT SERMON > > >}}
\newline
\newline
士師記 19:1-30
\newline
\begin{longtable}{cl}
\hline
\hline
章節 & 經文 (和合本修訂版)\\
\hline
19:1 & \begin{tabularx}{0.7\textwidth}{X} 當以色列中沒有王的時候，有一個利未人寄居以法蓮山區的邊界，他娶了一個猶大伯利恆的女子為妾。 \end{tabularx} \\ \\ \relax
19:2 & \begin{tabularx}{0.7\textwidth}{X} 這妾對丈夫生氣，離開丈夫，回到猶大伯利恆的父家，在那裡住了四個月。 \end{tabularx} \\ \\ \relax
19:3 & \begin{tabularx}{0.7\textwidth}{X} 她的丈夫起來，帶著一個僕人、兩匹驢跟著她去，要用好話勸她回來。女子就帶丈夫進到父親家裡。女子的父親看見了他，就歡歡喜喜地迎接他。 \end{tabularx} \\ \\ \relax
19:4 & \begin{tabularx}{0.7\textwidth}{X} 這岳父，就是女子的父親，留他住了三天。他們在那裡吃喝，住宿。 \end{tabularx} \\ \\ \relax
19:5 & \begin{tabularx}{0.7\textwidth}{X} 第四日，他們清早起來，利未人起身要走，女子的父親對女婿說：「先吃點東西，加添心力，然後你們才走。」 \end{tabularx} \\ \\ \relax
19:6 & \begin{tabularx}{0.7\textwidth}{X} 於是二人坐下，一同吃喝。女子的父親對那人說：「請你答應再住一夜，使你的心舒暢。」 \end{tabularx} \\ \\ \relax
19:7 & \begin{tabularx}{0.7\textwidth}{X} 那人起身要走，他岳父挽留他，他就留下，在那裡又住了一夜。 \end{tabularx} \\ \\ \relax
19:8 & \begin{tabularx}{0.7\textwidth}{X} 第五日，他清早起來要走，女子的父親說：「來，請加添心力，留到太陽偏西吧。」於是二人一同再吃。 \end{tabularx} \\ \\ \relax
19:9 & \begin{tabularx}{0.7\textwidth}{X} 那人同他的妾和僕人起身要走，但他岳父，就是女子的父親，對他說：「看哪，太陽下山，天快晚了，你們再住一夜吧。看哪，太陽偏西了，就在這裡住宿，使你的心舒暢，明天你們一早起來上路，回你的帳棚去。」 \end{tabularx} \\ \\ \relax
19:10 & \begin{tabularx}{0.7\textwidth}{X} 那人不願再住一夜，就備上兩匹驢，帶著他的妾起身走了，來到耶布斯的對面，耶布斯就是耶路撒冷。 \end{tabularx} \\ \\ \relax
19:11 & \begin{tabularx}{0.7\textwidth}{X} 將近耶布斯的時候，太陽快下山了，僕人對主人說：「來吧，我們進這耶布斯人的城，在這裡住宿。」 \end{tabularx} \\ \\ \relax
19:12 & \begin{tabularx}{0.7\textwidth}{X} 主人對他說：「我們不可進入外邦人的城，那不是以色列人的地方，我們越過這裡到基比亞去吧。」 \end{tabularx} \\ \\ \relax
19:13 & \begin{tabularx}{0.7\textwidth}{X} 他又對僕人說：「來，讓我們到基比亞或拉瑪的一個地方住宿。」 \end{tabularx} \\ \\ \relax
19:14 & \begin{tabularx}{0.7\textwidth}{X} 於是他們越過那裡往前走，將到便雅憫的基比亞的時候，太陽已經下山了。 \end{tabularx} \\ \\ \relax
19:15 & \begin{tabularx}{0.7\textwidth}{X} 他們進入基比亞要在那裡住宿。他來坐在城裡的廣場上，但沒有人接待他們到家裡住宿。 \end{tabularx} \\ \\ \relax
19:16 & \begin{tabularx}{0.7\textwidth}{X} 看哪，晚上有一個老人從田間做工回來。他是以法蓮山區的人，寄居在基比亞；那地方的人是便雅憫人。 \end{tabularx} \\ \\ \relax
19:17 & \begin{tabularx}{0.7\textwidth}{X} 老人舉目看見那過路的人在城裡的廣場上，就說：「你從哪裡來？要到哪裡去？」 \end{tabularx} \\ \\ \relax
19:18 & \begin{tabularx}{0.7\textwidth}{X} 他對他說：「我們從猶大的伯利恆過來，要到以法蓮山區的邊界去。我是那裡的人，去了猶大的伯利恆，現在要到耶和華的家去，卻沒有人接待我到他的家。 \end{tabularx} \\ \\ \relax
19:19 & \begin{tabularx}{0.7\textwidth}{X} 其實我有飼料草料可以餵驢，我和你的使女，以及與我們在一起的僕人都有餅有酒，甚麼都不缺。」 \end{tabularx} \\ \\ \relax
19:20 & \begin{tabularx}{0.7\textwidth}{X} 老人說：「願你平安！你所需用的我都會給你們，只是不可在廣場上過夜。」 \end{tabularx} \\ \\ \relax
19:21 & \begin{tabularx}{0.7\textwidth}{X} 於是老人領他到家裡，餵上驢。他們洗了腳，就吃喝起來。 \end{tabularx} \\ \\ \relax
19:22 & \begin{tabularx}{0.7\textwidth}{X} 他們心裡歡樂的時候，看哪，城中的無賴圍住房子，連連叩門，對老人，這家的主人說：「把那進你家的人帶出來，我們要與他交合。」 \end{tabularx} \\ \\ \relax
19:23 & \begin{tabularx}{0.7\textwidth}{X} 這家的主人出來對他們說：「弟兄們，不要做這樣的惡事。這人既然進了我的家，你們就不要做這樣可恥的事。 \end{tabularx} \\ \\ \relax
19:24 & \begin{tabularx}{0.7\textwidth}{X} 看哪，我有個女兒還是處女，還有這人的妾，我把她們領出來任由你們污辱她們，就照你們看為好的對待她們吧！但對這人你們不要做這樣可恥的事。」 \end{tabularx} \\ \\ \relax
19:25 & \begin{tabularx}{0.7\textwidth}{X} 那些人卻不肯聽從他。那人抓住他的妾，把她拉出去給他們。他們強姦了她，整夜凌辱她，直到早晨，天色快亮才放她走。 \end{tabularx} \\ \\ \relax
19:26 & \begin{tabularx}{0.7\textwidth}{X} 到了早晨，婦人回來，仆倒在留她主人住宿的那人的家門前，直到天亮。 \end{tabularx} \\ \\ \relax
19:27 & \begin{tabularx}{0.7\textwidth}{X} 早晨，她的主人起來開了門，出去要上路。看哪，那婦人，他的妾倒在屋子門前，雙手搭在門檻上。 \end{tabularx} \\ \\ \relax
19:28 & \begin{tabularx}{0.7\textwidth}{X} 他對婦人說：「起來，我們走吧！」婦人卻沒有回應。那人就將她馱在驢上，起身回自己的地方去了。 \end{tabularx} \\ \\ \relax
19:29 & \begin{tabularx}{0.7\textwidth}{X} 到了家裡，他拿刀，抓住他的妾，把她的屍身切成十二塊，分送到以色列全境。 \end{tabularx} \\ \\ \relax
19:30 & \begin{tabularx}{0.7\textwidth}{X} 凡看見的人都說：「自從以色列人離開埃及地上來，直到今日，像這樣的事還沒有發生過，也沒有見過。大家應當想一想，商討一下再說。」 \end{tabularx} \\ \\
[1ex]
\hline
\hline
\end{longtable}
$^{1}$願你繼續以聆聽聲道我們的敬拜.
今日我們選讀的經文是在舊約《士師記》第十九章.
一至九節和二十七至三十節.
分別在大家手上的聖經三百二十四和三百二十五頁.
請聽我讀出聖經的話語.
《士師記》第十九章.
當以色列終沒有王的時候.
有以法連山地那邊的一個利未人.
娶了一個猶大伯利恆的女子為妾.
妾行人離開丈夫.
回猶大伯利恆到了夫家.
在那裡住了四個月.
她丈夫起來帶著一個僕人兩匹驢去見她.
用好話勸她回來.
女子就引丈夫進入夫家.
他父見了那人便歡歡喜喜地迎接.
那人的岳父就是女子的父親.
將那人留下住了三天.
於是二人一同吃喝住宿.
到了第四天.
利未人清早起來要走.
女子的父親對女婿說.
請你吃點飯加添心力然後可以行路.
於是二人座下一同吃喝.
女子的父親對那人說.
請你再住一夜暢快你的心.
那人起來要走.
他岳父強留他.
他又住了一宿.
到第五天.
他清早起來要走.
女子的父親說.
請你吃點飯加添心力.
等到日頭偏西再走.
於是二人一同吃飯.
那人和他的妾和僕人起來要走.
他岳父就是女子的父親對她說.
看啊日頭偏西了請你再住一夜.
天快硬了可以在這裡住宿.
暢快你的心.

$^{41}$明天早起行回家去.
第27節.
早晨他的主人起來開了房門.
出去要行路.
不料那婦人撲倒在房門前.
兩手踏在門檻上.
就對婦人說.
起來我們走吧.
婦人卻不回答.
那人便將她托在路上.
起身回本處去了.
到了家裡.
用刀將妾的屍身切成十二塊.
使人拿著傳送以色列的四景.
凡看見的人都說.
從以色列人出埃及地直到今日.
這樣的事沒有行過也沒有見過.
現在應當思想.
大家相宜當怎樣辦理.
弟兄姊妹,平安.
平安.
我們繼續用士師記最後一個段落.
重新檢視一下我們自己的生命.
我們處於一種甚麼的境況.
有一個這樣的說法.
同一條魚說水是甚麼來的.
一來這條魚就不像人類.
可以解釋到水是一回甚麼事.
但更重要的是.
當他活在生存空間裡.
其實所有事都變得理所當然.
直至令到這條魚明白水是一回甚麼事.
就是將這條魚從水裡撈上來.
他不能呼吸.
他不能游動.
那種束縛,窒息.
魚就會知道甚麼叫水.
同樣的道理很多時候都放在我們福音信仰裡.
我們問一個福音或我們活在恩典當中的一個人.
上帝是一個怎樣的道理.

$^{81}$福音是一個怎樣的訊息.
我們每一個人活在恩典當中.
好人好者我們大概不是很體會.
直至到在我們生命裡面出現了一些狀況.
又或者很多的尾線離開了我們的時候.
我們就開始明白福音原來對我們生命是那麼重要.
今天我們用這個角度來看士師記第十九章.
跟之前的第十七章和第十八章的記載有些不同.
雖然同樣都是講一個利未人的故事.
很多經學家說是兩個利未人.
你當他一個也好或兩個也好.
在第十七章第十八章.
之前上個星期我們所讀的就有多次提及耶和華上帝.
又或者以上帝為名來實踐他們的生活.
不過去到第十九章.
一路讀上去我們完全看不到上帝介入在這個情況裡面.
剛才我們一頭一尾讀了這段聖經.
開頭的時候那個妾侍回到外家.
到尾的時候為何會被分成知解.
變為十二份呢.
究竟這件事發生了甚麼事呢.
一個沒有上帝的處境當中.
就發現了出現了一件這樣的事.
今天這件事讀上去真的很暗黑.
很多同工都不會特別去處理這段經文.
因為很多時候不知道怎樣去應用.
我們就帶著這段經文裡面.
雖然以色列民已經成為神國的子民.
但是他撇離了上帝.
離開了上帝之後.
他們整個民族走向一個甚麼的地步.
我們就帶著這個角度去思想今天這段經文.
我們先祈禱求主幫助我們同心祈禱.
生愛主我們求你在我們當中.
說領我們眾人.
讓我們今天能夠明白你的話語.
又讓我們心靈當中看到今天這段經文.
對每個世代作出的警惕.
又讓我們能夠真的安著你的心.
明白信仰耶穌基督對我們生命有多重要.

$^{121}$我們帶著恭敬的心在你面前去研讀.
求你帶領我們幫助我們.
垂聽我們的祈禱奉基督耶穌得勝的聖名祈求.
我們看第十九章.
在經文裡面一開始就很快很明確地記載著.
以色列沒有王的時候.
讓我們想起之前的循環.
沒有一個事司.
沒有人帶領整個民族的.
各人任意而為.
裡面就說到有一個人住在以法連山地的一個尼美人.
這是一個神職人員.
經文裡面特別說到他娶了一個妾.
在聖經裡面沒有特別禁止尼美人娶妾.
不過你可以明白一個神職人員.
原則上他的生命是交付給上帝.
他也是以服侍上帝為他生命裡面最重要的目標.
而當時的納妾.
很多人納妾除了繼後香燈之外.
更多是出於自私.
因為妾侍比家裡的僕人高一點.
但比家裡的正室低很多.
簡單來說就是一個妾侍.
就是一個廉價的勞工.
同時間滿足主人的性慾望.
你可以知道尼美人納了這個妾.
其實帶著一點點.
大家都會覺得有點奇怪.
不過經文裡面再繼續說.
第二節也有不同的說法.
妾行淫離開丈夫回白.
由大白離行到了夫家.
再納女住了四個月.
經文和合本翻譯是說妾行淫.
不過一邊讀上去你會覺得很奇怪.
行淫了然後老公為何不用律法將婦人致死呢.
可以用石頭扔死她.
所以和合修訂本也改了翻譯.
因為行淫這個字有兩個翻譯.
一個是說男女不倫的關係.

$^{161}$第二個處理是對丈夫生氣.
如果用對丈夫生氣.
即是發怒.
可能尼美人激怒了妾侍.
一怒之下她離開了回到夫家四個月.
讀上去上文下理似乎比較通順.
而且理解上也沒有製造更多困難.
我們用她生氣了丈夫.
向丈夫生氣.
所以回到外家.
她的外家在由大白離行住了四個月.
隔了四個月之後.
尼美人帶著一個僕人.
經文裡說帶著兩匹驢去見她.
用好話勸她回來.
你讀上去就知道.
如果行淫不會用好話勸她回來.
經文裡說尼美人帶著兩匹驢.
那匹驢不是讓她眼前的僕人坐.
而是希望用驢很尊貴地接妾侍回到家裡.
在她一直用好話勸她的時候.
經文很快就說到這個人去到岳父的家.
岳父見到女婿.
口水滴滴滴.
於是強行留她在裡面住宿.
希望留下接待她有三天的日子.
接待三天是當時猶太人的規矩.
遠道而來的遊客來到.
一般來說禮貌上是接待她三天.
三天後大家都歡唱.
而這個妾侍整件事都某程度是緘默的.
不過是這個妾侍引導尼美人和她的僕人.
去到岳父的家裡住了三天.
住了三天後經文記載.
岳父很喜歡尼美人.
一方面可能是一個自己的女婿.
他很盛情地去接待她.
另一方面其實尼美人這個身份.
都是帶著一份尊貴的.
所以在岳父一個尼美人的女婿.

$^{201}$他很想多留她一段時間.
一個服侍神的人來到他家裡留宿.
他覺得是一件很榮耀也很快樂的事.
所以留了第三天後再多留一天.
就是第四天了.
到了第四天想走的時候.
岳父就說都很晚了.
不如你多留一天.
太陽平西你出去趕路就很晚了.
於是就強留多一天.
到了第五天早上.
尼美人和僕人想走了.
他就說現在離開實在太倉促了.
再加上不如等到日間沒那麼曬.
當時早上或傍晚太陽沒那麼猛烈.
於是就再強留.
到了第五天黃昏的時候.
尼美人和僕人都忍不住要走了.
即使岳父怎樣強留她也好.
怎樣加添心力.
怎樣暢快住多一天.
尼美人都不會再聽.
於是他們就在日頭平西的時間.
即是傍晚.
尼美人帶著僕人和妾侍.
坐著兩匹驢離開這個地方.
大家記住剛才說離開的時間是傍晚的時間.
開始太陽落山的日子.
於是趕路的時候.
差不多回到家裡.
還沒到的地方已經進入天黑.
到了天黑不會再晚上趕路.
一來風沙大.
第二樣東西就是有很多盜賊.
很不安全.
所以他們就希望找一個地方.
來留宿一晚.
第二天早上再走路.
回到自己的地方.
在他們一路走的時候.

$^{241}$經文裡面記載著.
一路走到當時還不是以色列的首都.
那個地方叫做耶布斯.
經文括著告訴你.
那個地方叫做耶路撒冷.
這是第十節裡面的交代.
在耶路撒冷當時是屬於外邦人的地方.
有個僕人說不如進去住一下.
在耶布斯人的地方住宿一晚.
第二天早上再趕路.
但是利美的主人就說.
我們不可以進去.
因為耶布斯這個地方是外邦人住的.
他的考慮是什麼呢.
他的地方是屬於.
其他國家.
其他民族人住的地方.
所以在這裡.
寧願回到自己人.
變亞敏茲派集結.
集居的地方就是基比亞.
反正都是.
當然是回到自己人的地方.
來居住.
所以就沒有去到.
耶布斯這個外邦人居住的地方.
就進入了基比亞這個城市.
他心裡比較安全.
是變亞敏茲派.
也是自己人.
所居住的一個地方.
不過進入基比亞的時候.
其實太陽已經日落了.
去到那個地方裡面.
其實在街上.
在街上.
那些基比亞的變亞敏人.
或者自己的.
以色列人.
沒有人理會他們.

$^{281}$即使他們一直走著.
他們擺明是一個客人.
但是當地人就很冷漠.
人們就眼望望.
沒有特別來接待.
利未人和.
他的妾侍.
和那個僕人.
直到晚上.
在獨節那裡.
有一個老年人在田間裡工作.
回來了.
下班了.
原來他原本不是當地人.
他原本是住在以法連山地的人.
現在就住在基比亞.
在那個.
一個寄居在這個地方的老年人.
於是乎.
見到這個利未人的時候.
就接待他了.
經文裡面很快就會這樣說.
他們沒有地方居住.
還有他們有爐.
還有糧草.
不過今晚需要借宿一宵.
所以.
我和我的妾.
和我的僕人.
有餅有酒.
並不缺少什麼.
也就是說他不是來白吃的.
他只是需要借地方.
睡一晚.
不用來在街上.
或者在那裡睡.
他需要一個地方.
來借宿一宵.
豪情的.
像以前的龍門客棧.

$^{321}$有餅有酒.
我們什麼都不缺.
豪情的說了一番這樣的話.
於是這個老年人就說.
願你們平安.
你所要的我都會供給你的.
不過你不可以在街上過夜.
於是乎就將這幾個人.
接待到自己的家裡.
幫他們.
給他們一些爐.
有飼養的.
餵食他們.
接著幫他們洗腳.
和接待.
和他們一起吃喝.
這個老年人就.
帶著一份很恭謹的心.
來接待這位美人.
整個的家庭.
他們一邊吃飯.
一邊言談甚歡的時候.
大家很放鬆.
很歡暢的時間.
突然在家門外.
有人敲門.
啪啪啪.
有些人就在老年人的門外敲門.
為什麼敲門呢.
原來在基比亞.
有一些流氓.
經文裡面記載.
這些是匪徒.
匪徒以為我們是打家劫舍.
更好的說法.
就是這班人是流離浪蕩的流氓.
不務正業的.
就在這裡拍門.
啪啪啪.
接著就跟老人家說.

$^{361}$你把你家裡的男人.
帶出來.
我讀上去也有點尷尬.
我們要與他交合.
你聖經讀得很文雅.
不過.
我們知道那個狀況.
這一班是一班男人.
看到有一班客人.
去到老人家的家裡坐直.
他就要做一個.
同性戀的行為.
某程度上.
講得直白一點.
就是要雞奸他.
要將他交出來.
老人家就對他們說.
弟兄不可以這樣作惡.
這個人來到我家裡.
你們不可以做這樣的醜事.
第23節裡面這樣說.
24節繼續說.
這個老人家說.
我家裡有個女兒.
還是處女.
還有這個人的妾侍.
隨便你們.
怎樣處理他們.
但是這個男人.
不可以讓你們.
做這些醜惡的事.
同樣的事件.
或者這樣的情景.
如果你熟讀舊約聖經.
其實舊約裡面.
大概出現過一次這樣的狀況.
就是在創世記裡面.
第19章1-9節.
說到羅德接待兩個天使.
然後在.

$^{401}$索多瑪這個地方.
當羅德接待了.
兩個天使進入他家裡的時候.
索多瑪所有的人.
都要拍門.
由老到亂都要說.
這兩個天使是男的.
不要以為天使是女的.
是男的.
兩個天使來發生性行為.
整件事.
就好像一個翻版.
不過在.
舍斯記裡面.
記載的.
內容跟創世記有點不同.
在.
創世記裡面.
記載著.
說到上帝默默地保守.
羅德的家庭.
在一個連十個二人都沒有.
的索多瑪這個城市裡面.
上帝.
插手干預.
上帝保護羅德.
整個家庭.
這兩個是天使.
叫他們這群人.
去遊輪.
但在舍斯記的記載裡面.
上帝是執默的.
這群人也沒有想過.
呼求上帝.
為他們解決眼前.
他們面對的.
這麼大的危機.
反而.
這個老人家用了的方法.
就是家裡有個女兒.

$^{441}$還有這個人的妾侍.
讓這兩個女人.
成為這群基比亞的流氓的.
洩慾工具.
當然那群人聽到.
這樣的說話有什麼反應呢.
25節那些人不聽.
這個老人家所說的話.
拒絕了.
這樣的說法.
然後那個經文的最.
令你讀上去.
你會咬牙切齒的地方.
就是經文在.
電光火石之刻.
當拒絕.
在擾攘的時候.
這個利未人被稱為.
那人.
那人就將她的妾.
拉出去交給他們.
他們便與她交合.
終夜凌辱她.
直到天色快亮才放她去.
在經文裡面.
和合本就很短促地.
寫了一個字.
其實在原文裡面用的字眼.
那個字眼就是用.
利未人獻祭.
祭生的那個.
處理祭生.
因為他們其實既是祭祀的人.
也都是一個某程度的屠夫.
他們是處理那些生畜的.
用一個.
這樣的組合動詞.
就是拉.
抓住她.
然後拉出去.

$^{481}$整個這個動作.
是一個很年宵大打.
很純熟的動作.
不過一個利未人.
這個動作他做了什麼呢.
就是將剛才他在說.
要去到臥府裡面.
求回來的那個妾.
在電光火石之刻.
在她住的地方裡面.
抓住她.
可想而知那個婦人是掙扎的.
掙扎的時候拉住她.
然後就拉她出去.
推她出門.
然後關了門.
這個妾.
就被基比亞這一班.
中夜.
來去凌虐.
直到天色快亮的時候.
才放她離開.
你可以想像.
在這個情境.
是一個很.
這個晚上是一個漫漫的場景.
是一個痛苦的場景.
是一個.
好像野獸這樣的狀況之下.
將一個女子.
傷害.
徹底的.
來去凌辱她.
這是一個很痛苦的.
一個狀況.
到了第26節.
經文裡面再說.
天快亮了.
這個婦人回到她的主人住宿的門前.
就突倒在地.

$^{521}$直到天亮.
值得留意到第26節.
講到這個妾侍.
已經被完全的.
傷害.
身體已經完全的沒有力量.
倒在地.
經文裡面描述的這個利未人.
她改換了她的稱呼.
之前用那人.
或者利未人.
去到這裡就在說妾侍的主人.
你留意一下.
這裡有一點點地方.
待會值得我們再去深思.
不過當她.
完全沒有力量.
伏在地的時候.
她的主人.
好像沒有事發生一樣就打開了門.
然後經文再繼續.
來到她記載.
第27節.
她的主人起來.
打開了房門.
出去要行路的時候.
見到這個婦人.
倒在房的門前.
她的手.
搭在門欄上.
你想像.
整個的.
圖像.
這個婦人.
經歷過一晚被.
凌辱的那種痛苦.
她身體已經.
嚴重的傷害.
已經虛脫.
她連開門的力量都沒有.

$^{561}$手搭在門裡面.
但當一個.
如果你看到一個這樣的人.
如果你作為一個血肉之軀.
你心裡面會很痛苦.
你會覺得.
有一種.
一種完全不忍心.
不過說到這個.
利未人有什麼反應呢.
第28節.
對婦人說.
起來我們走吧.
這個.
淡淡言的一句.
回應.
好像這個利未人的心.
忽然間完全的.
好像關了門一樣.
完全的變得冷酷無情.
麻木的.
沒有任何反應.
半點眼淚都沒有.
就是.
婦人不回答他.
於是這個人就將.
他的妾侍.
就是.
拖在那個爐的上面.
起來就回到自己的家裡.
在聖經裡面.
用幾個字簡單地.
記載.
同樣你都呈現了一個.
被你看見的圖畫.
一個受害的婦人.
和一個.
原本是她的丈夫.
雖然妾侍也好.
這個丈夫.

$^{601}$竟然面對著.
在一個危急關頭裡面.
代替了她.
成為了被蹂躪的.
那個.
但是當他見到她的時候.
他就衝衝衝.
淡淡然的回到自己的地方.
去到第29節.
究竟.
這個婦人.
是之後.
被這樣蹂躪之後.
致死了.
還是回到家裡才死呢.
經文裡面沒有特別的來講.
所以我們沒有辦法.
來推敲究竟.
這個妾侍的死亡時間.
但去到第29節.
我們看到的就是.
第29節裡面講到這個利未人.
回到家裡的時候.
就將她的妾.
那個屍體.
切成12塊.
原文用的.
那個字眼.
用回她在聖殿.
或者回望裡面獻祭.
將那些祭僧.
經絡來順序的.
來解剖.
那些的生畜.
用回同樣的字.
拉.
抓住.
切割.
將她的妾侍.
身體分成12塊.

$^{641}$將她帶到.
傳送到以色列的死境.
整件事.
所以讀上去.
是一個很血腥的.
一個情境.
是一個很暴力的情境.
是一個很盲目不仁.
你讀上去.
你自己會很心痛的一段經文.
記載的.
很慘痛的一個歷史.
經文.
再繼續講.
當這個婦人的身體.
被分到死境的時候.
見到的人.
就會說.
由以色列出埃及地.
直到今天.
這些的事真的從來沒有.
行過也都沒有見過.
空前的.
講一個這麼殘暴的一件事.
竟然在他們當中來發生.
這一班人.
就繼續思想.
應該怎樣去辦理.
整個故事.
大概你掌握得到.
是一個很痛苦.
一個你完全估計不到.
為什麼會發生一個這樣的事.
以為兩夫妻吵架.
但到最後好像滾雪球一樣.
成為了一件.
這麼恐怖.
這麼兇殘的一個決案.
或者一件這麼令人髮指的.
一個強暴的事件.

$^{681}$不過這個.
經文裡面.
其實如果對應之前.
我們所讀的17章第18章.
多次提及耶和華.
去到第19章.
你見到通常.
如果在聖經裡面記載.
人在危急的時間.
會插手干預.
上帝會幫助.
甚至在舍斯記裡面說.
當人呼求上帝的時候.
上帝會介入.
在那班的.
以色列民當中.
拯救他們脫離這些兇惡.
不過在這裡.
整個過程.
打破了這個循環.
之後進入一個更加惡性的循環.
昔日當.
以色列民.
向上帝呼求的時候.
上帝插手.
但在這裡.
他們完全好像沒有上帝的狀況一樣.
他們只是用他們自己的方式.
來處理眼前的危機.
即或是這個利未人.
面對著眼前的狀況.
他都是用他自己昔日的技術.
用他自己擁有的能力.
抓住,拉出去.
這是他一直憲制的動作.
將這件事.
竟然成為了他處理危機的方法.
將這件事.
將上帝排他於.
他們的危難當中.

$^{721}$上帝也保持了整個的.
他們的危難當中.
整個的.
好像完全沉默的一樣.
在聖經裡面.
在這樣的狀況.
我們可以用來作一個應用.
當這個世界.
拒絕上帝.
排他上帝.
生命裡面很多的問題.
就會在一個.
被排他上帝的世代.
或者處境裡面.
出現.
在經文裡面講的有四個.
當這班人拒絕上帝的時候.
出現了一些狀況.
第一.
在一個沒有上帝的世代.
關係.
完全變得扭曲.
在整個的故事裡面.
都依然講有夫妻的那種愛.
有接待的那種愛.
是吧?.
對人的那種.
都有一些友善,有些人很冷漠.
但都有人來接待.
但當沒有上帝的時候.
在美好的關係.
面對著困難,面對著危機的時候.
一切都會變得扭曲.
特別是在講.
本來這個利未人.
是帶著那匹驢.
來接待.
勸回.
妾侍.
回到自己的地方.

$^{761}$可能他們的關係.
都是一個.
從屬的關係.
沒有一個明媛正娶.
之下的那種關係.
但我相信在這裡面,他們都有愛情的關係.
至少這個男人都願意.
付出一些代價.
希望讓這個妾侍回到自己的家裡.
不過在沒有上帝的處境裡面.
這些的愛情可能.
變得很短暫.
這些的愛情可能傾刻之間.
就變成了一個鐵石心腸的.
關係.
在一種這樣的關係的時候.
很多時候,女性都成為了一種犧牲品.
如果大家留意.
在一個沒有上帝的社會.
一個沒有真理的年代.
受害者.
永遠會成為.
更加大的受害者.
而在聖經的時代裡面.
一般來說,女性的地位都不及男性.
但是當人.
排除了上帝的時候.
女性的地位更加顯得.
更加好像犧牲品一樣.
女性成為了.
必然需要犧牲.
在這個狀況裡面.
在那些對話當中.
甚至乎.
在婦女關係裡面,老人家說自己還有一個女兒.
可以讓這班人去洩慾.
整件事.
都是完全.
顛覆了正常的人倫觀念.
沒有上帝的世代.

$^{801}$一切的關係.
是依然有這些關係.
不過這些關係到最終會被扭曲.
男與女.
親朋.
或者人與人之間.
人就變得很肉體.
按著人的感受.
需要來經營自己的關係.
經文裡面.
亦都值得對我們第二個反思.
在沒有上帝的時代.
任何的修和.
願意回復一些.
修復關係.
到最後是甚麼呢.
只是為了真正的.
衝突來作掩飾.
只是包不住火.
如果沒有了上帝的年代.
所有的修和.
到最後都是站立不住.
沒有了上帝的時代.
就是說要復和.
談何容易.
因為人沒有一個基礎.
令到人可以產生.
更加高尚的情操.
來原諒.
或者承認自己.
所犯過的錯.
這個男人.
最初去到岳父的家裡.
打算勸.
接事回來.
但在過程裡面.
經文記載.
其實這個男人沒有特別去說話.
說得最多對話的只是他和岳父.
岳父就說.

$^{841}$吃喝快樂,暢快你的心.
用快樂,用吃喝.
來遮掩他和.
這個接事之間.
可能.
那種不咬弦的.
或者吵架的這段關係.
整個過程裡面.
沒有針對問題,也沒有坦誠.
面對自己.
當然你說,那是今天.
輔導所說的,牧師.
你不要將今天的輔導套入經文.
是的,對的.
不過去到今天輔導.
又怎麼說呢?.
早幾天有機會.
回到神學院裡面上課.
和一些資深的老師.
談基督教的婚姻輔導.
我就問了一個問題.
問了我的老師.
老師,如果你除了.
那頂神學院老師的帽子.
你用一個坊間的輔導員.
這頂帽子來面對.
夫妻之間的那種爭執.
或者有一些.
的搖書,你會怎麼做.
令到那個.
夫妻之間可以產生.
搖書呢?.
然後這個神學院老師跟我說.
他說,其實.
坊間裡面說的輔導.
是不會說搖書.
成為了輔導的目標.
輔導的目標.
極其量是.
令到你會舒服一些.

$^{881}$給到你有一些.
方法處理眼前的衝突.
但是搖書.
已經不是輔導.
眼前坊間裡面所說的.
可以達至到的目標.
如果要說輔導.
到最後.
我們都是要訴諸上帝.
但是如果沒有上帝的年代.
搖書.
坦誠.
承認自己的問題.
這些一切都可以成為.
有其他東西來去遮掩.
我們只可以用一些方法.
令到自己舒服一些,過得好一些.
但是我們沒有辦法.
面對著生命最深處的傷害.
來用更高的情操.
用更高的神聖的態度.
來面對生命的傷害.
沒有上帝.
就沒有搖書.
沒有上帝的修和.
一切都是為衝突來掩飾.
第三方面.
沒有上帝的心靈.
是盲目不仁的心靈.
你會很奇怪的.
你讀這個第十九章的故事.
這是一個神跡人員.
一個利未人.
他服事上帝.
他服事一個.
或者成為了人和神之間的一個憲制.
一個很重要的橋樑.
他要明白上帝願意.
將福氣賜給人.
上帝那種憐憫.

$^{921}$上帝對人的那種歡懷.
是成為了憲制的一個很重要的基礎.
不過原來去到這裡.
說到利未人心裡面.
好像只有技術一樣.
抓住.
用他日常做的動作.
去面對眼前的危機.
他多麼辛苦.
希望勸退.
切事回來.
但到最後.
其實他對他都是盲目不仁.
面對著他的傷害.
他淡淡然說起來.
走吧.
根本這個切事都沒有力.
再去走面前的路.
沒有上帝的心靈.
其實都是一個盲目不仁的心靈.
盲目是一個什麼意思呢.
盲目我們今天.
我們看在經文裡面都有講過.
不過在輔導裡面.
在心理學裡面.
都有幾樣東西講盲目的.
第一.
盲目的人原來是不願意.
或者沒有辦法.
投入在你的生活當中.
如果我成為了一個盲目的人.
即是說.
我是有信仰.
但我的人生是抽離.
我投入不到.
不願意將自己投入.
在某一種的關係裡面.
盲目也代表著什麼呢.
盲目就是.
那個人不知道自己的感受.

$^{961}$不能夠講出自己.
當下裡面.
不論是憤怒.
悲哀.
或者是快樂.
基本上他沒有辦法.
去言說自己的感受.
又或者.
更加多的盲目的人.
心裡面積累大量的憤怒.
他成為了一個.
只是.
用很多的憤怒.
填積在他的心靈空間裡面.
用憤怒.
來過他自己的生活.
但他表面上是一個.
可以很平靜的人.
但心裡面是可以爐火中燒.
盲目就是一個這樣的心理狀態.
盲目的人.
也是和別人保持.
永遠保持距離.
是一個抽離的人.
你沒有辦法知道他心裡在想什麼.
你沒有辦法進入他的內心世界.
在他的口中裡面.
大部分都是一些.
超理性的一些說話.
和很多的.
技術字眼.
就是.
他生命裡面失去了對事物.
對人.
對生活的那些興趣.
你這樣講的時候你發覺.
有時候這個利未人.
和剛才心理學所講的.
都很近似.
不過更加近似的.

$^{1001}$今天這個世代很多的男人.
都是這樣.
我想不單止是男人.
今天這個都市上很多的人.
都是生活在盲目的.
心靈世界當中.
自己抽離.
心裡面不能夠沒有喜樂.
內心大量的苦毒.
負面情緒.
但表面上是很平靜.
但是.
他生命裡面就像一塊石板一樣.
沒有一個人.
與事可以進到.
這個心靈空間裡面.
這個就是盲目.
沒有上帝的世代.
就是一個盲目的世代.
人與人之間就帶著.
一份這樣的.
盲目來互相.
相處.
到最後只會積累更加多.
更加多的悲劇.
而這個.
亦都是第四點我要講的.
沒有上帝的處事.
就是將上帝.
排他於他們.
自己處理這個危機的處事.
一切.
只會製造更加多的仇恨.
為什麼.
這個利未人要將這個.
設事斬成十二份.
分給十二個支派.
目的只是一個.
他要將這件事.
挑起他.

$^{1041}$令到.
裡面這種仇恨的心.
當我們去讀的時候.
我們都很憎那些基比亞人.
不過當這種怒氣.
透過一個這麼血腥的.
一個身體.
來呈現的時候.
帶來的只是有更加大的.
那種仇恨的心.
其他那些人.
其實不應該涉及.
在當中或者要怎樣的公平處理.
是可以有.
有方法的.
不過這個利未人.
選取的就是.
用他自己心裡面.
用他的專業.
將他的設事肢解.
然後將他的殘肢.
分給.
血淋淋地告訴其他十二派.
十二個支派知道.
他們很想製造一個.
更加累積仇恨的.
一個情境.
弟兄姊妹.
今天這個是一個.
很負面的故事.
不過我們就用這個.
負面的故事來提醒我們.
這個世界很需要上帝.
我們很需要.
福音.
如果我們的心靈裡面.
又或者這個世代裡面.
沒有了那個救恩.
我們人生就是.
這樣.

$^{1081}$好像這個設施記第十九章.
裡面一樣.
男女都依然有.
男女的關係.
有淪捨之間的關係.
但到了面對著.
這樣的情景的時候.
一切都會變得.
好像血腥.
一切的仇恨都會變得累積.
所有的東西.
就走向一個負面的循環.
荷西亞書有講.
整段經文.
他用幾個字來講.
他講這個是敗壞之仇.
荷西亞才用.
這一段的歷史.
來講.
以法連心心的敗壞.
如在基比亞的日子一樣.
耶和華必.
紀念他們的罪孽.
追討他們的罪惡.
上帝在這一刻.
不代表上帝完全.
殺手不管.
有一天上帝要處理.
這一班人的罪惡.
同樣這個世代.
如果我們繼續走向一個沒有上帝的世代.
上帝有一天回來.
他會帶著.
公義的審判.
來面對.
求主繼續幫助我們.
當我們知道我們是認識上帝的人.
就讓這個負面的故事.
警惕著我們心靈.
警惕著我們行事為人.

$^{1121}$我們不可以沒有上帝.
我們就帶著這個故事.
作為我們的警惕.
警醒我們的內心.
讓我們不單止口講有上帝.
心靈我們都邀請上帝.
成為我們心裡面的.
我們的王.
帶領我們走面前人生的路.
願主繼續對我們講說話.
這就是我們說的說法.
\newpage



\section{約翰福音 3:1-15}
\label{sec:e5U96ffu0L0}
\textbf{2023年10月15日崇拜直播｜梁家麟牧師｜耶穌談重生｜約翰福音3\_1-15}
\newline
\newline
連結: \href{https://youtube.com/watch?v=e5U96ffu0L0}{\texttt{https://youtube.com/watch?v=e5U96ffu0L0}} ~~~~ 語音日期: 2023-10-17
\newline
\newline
\hyperref[sec:_4zbsWU6e2A]{\small{< < < PREV SERMON < < <}}
~
\hyperref[sec:index_chronic]{\small{[返順時目]}}
~
\hyperref[sec:jwnc_J9h4aM]{\small{> > > NEXT SERMON > > >}}
\newline
\newline
約翰福音 3:1-15
\newline
\begin{longtable}{cl}
\hline
\hline
章節 & 經文 (和合本修訂版)\\
\hline
3:1 & \begin{tabularx}{0.7\textwidth}{X} 有一個法利賽人，名叫尼哥德慕，是猶太人的官。 \end{tabularx} \\ \\ \relax
3:2 & \begin{tabularx}{0.7\textwidth}{X} 這人夜裡來見耶穌，對他說：「拉比，我們知道你是由神那裡來作老師的；因為你所行的神蹟，若沒有神同在，無人能行。」 \end{tabularx} \\ \\ \relax
3:3 & \begin{tabularx}{0.7\textwidth}{X} 耶穌回答他說：「我實實在在地告訴你，人若不重生，就不能見神的國。」 \end{tabularx} \\ \\ \relax
3:4 & \begin{tabularx}{0.7\textwidth}{X} 尼哥德慕對他說：「人已經老了，如何能重生呢？豈能再進母腹生出來嗎？」 \end{tabularx} \\ \\ \relax
3:5 & \begin{tabularx}{0.7\textwidth}{X} 耶穌回答：「我實實在在地告訴你，人若不是從水和聖靈生的，就不能進神的國。 \end{tabularx} \\ \\ \relax
3:6 & \begin{tabularx}{0.7\textwidth}{X} 從肉身生的就是肉身；從靈生的就是靈。 \end{tabularx} \\ \\ \relax
3:7 & \begin{tabularx}{0.7\textwidth}{X} 我說『你們必須重生』，你不要驚訝。 \end{tabularx} \\ \\ \relax
3:8 & \begin{tabularx}{0.7\textwidth}{X} 風隨著意思吹，你聽見風的聲音，卻不知道是從哪裡來，往哪裡去；凡從聖靈生的也是如此。」 \end{tabularx} \\ \\ \relax
3:9 & \begin{tabularx}{0.7\textwidth}{X} 尼哥德慕問他：「怎麼能有這些事呢？」 \end{tabularx} \\ \\ \relax
3:10 & \begin{tabularx}{0.7\textwidth}{X} 耶穌回答，對他說：「你是以色列人的老師，還不明白這些事嗎？ \end{tabularx} \\ \\ \relax
3:11 & \begin{tabularx}{0.7\textwidth}{X} 我實實在在地告訴你，我們所說的是我們知道的，我們所見證的是我們見過的，你們卻不領受我們的見證。 \end{tabularx} \\ \\ \relax
3:12 & \begin{tabularx}{0.7\textwidth}{X} 我對你們說地上的事，你們尚且不信，若對你們說天上的事，如何能信呢？ \end{tabularx} \\ \\ \relax
3:13 & \begin{tabularx}{0.7\textwidth}{X} 除了從天降下的人子，沒有人升過天。 \end{tabularx} \\ \\ \relax
3:14 & \begin{tabularx}{0.7\textwidth}{X} 摩西在曠野怎樣舉蛇，人子也必須照樣被舉起來， \end{tabularx} \\ \\ \relax
3:15 & \begin{tabularx}{0.7\textwidth}{X} 要使一切信他的人都得永生。 \end{tabularx} \\ \\
[1ex]
\hline
\hline
\end{longtable}
$^{1}$今日講到的經文是記載在約翰福音第三章一至十五節.
是耶穌向尼哥底姆講重生的事.
有一個法利賽人名叫尼哥底姆.
施猶太人的官.
這人夜裡來見耶穌說.
拉比 我們知道你是由神那裡來作師父的.
因為你所行的神蹟.
若沒有神同在 無人能行.
耶穌回答說.
我實實在在的告訴你.
人若不重生 就不能見神的王.
尼哥底姆說.
人已經老了 如何能重生呢.
豈能再進無福生出來呢.
耶穌說 我實實在在的告訴你.
人若不是從水和聖靈生的.
就不能進神的王.
從肉身生的 就是肉身.
從靈生的 就是靈.
我說 你們必須重生.
你們不要以為稀奇.
風隨著意思吹.
你聽見風的響聲.
卻不曉得從哪裡來 往哪裡去.
凡從聖靈生的 也是如此.
尼哥底姆問他說.
怎能有這事呢.
耶穌回答說.
你是以色列人的先生.
還不明白這事嗎.
我實實在在的告訴你.
我們所說的是我們知道的.
我們所見證的是我們見過的.
你們卻不領受我們的見證.
我對你們說地上的事.
你們尚且不信.
若說天上的事 如何能信呢.
除了從天降下.
仍舊在天的人子.
沒有人升過天.

$^{41}$摩西在曠野怎樣舉石.
人子也必照樣被舉起來.
叫一切順他的 道德永生.
提醒大家平安.
很多謝剛才志華姐妹的讀經.
讀得很清晰.
從約翰福音第三章開始.
就寄述了耶穌和幾個人的一連串對話.
包括了這裡的尼哥底姆.
四張的撒瑪利亞的婦人.
和一些外邦的官員.
五張不屍大辭旨旁邊的病人等等.
是一連串的對話.
很多人都相信尼哥底姆是一個人版.
不只是代表他自己.
是代表一群這樣的人.
代表在猶太教中間的那些有宗教和社會地位.
他擔任公職的猶太人的官.
同時他也是拉比 是律法師.
第十節耶穌說你是以色列人的先生.
所以他就是拉比.
第七章我們會看到他是工會成員.
尼哥底姆和他所代表的那一群人.
他們大概是比較中立 比較客觀.
不是帶著自己的利益保衛自己的小集團.
所以他們可以接納耶穌所做的事.
確實是有意思.
甚至肯定他正在做的是正面的.
是從神而來的事.
不過他們是否於是就成為了基督徒呢?.
這是另一回事.
但最少他們對耶穌有好感.
不是一棍打死全面否定耶穌.
不是的.
中間可能有人甚至想信耶穌.
不過因著現實利益的衝突.
他們不敢公開表白他們的信仰.
所以他們有時候需要別人給多些鼓勵幫助他們.
在《使徒行傳》第五章記載了一個叫加瑪烈的人.
這個人大概和尼哥底姆是差不多的人.

$^{81}$所以在耶穌時代早期教會.
也有一群猶太教所謂建制裡面的人.
但是對基督教對耶穌有好感.
耶穌基督的道理是適合所有人的.
凡有耳的就能夠聽得明白.
這個道理很有趣.
沒有讀過書的人.
沒有受過教育的人都能夠聽得明白.
擁有最高學問的人也能夠被折服.
問題是多數人不是完全持一個客觀中立的態度去聆聽耶穌的話.
他們總是帶著過去有的種種知識和經歷.
所凝固的觀點信念甚至是成見去批判.
選擇性的聆聽.
而他們對耶穌的話不是用一個完全開放的理性的抉擇.
他們很多時候是夾前很多社會性的利害關係的考量.
就是說我信耶穌會有多少社會代價.
我會不會被我身邊的人排擠.
我會不會升職加不到薪.
我會不會不能維持我既有的生活方式.
如果我的朋友知道我信耶穌的話.
他們去酒吧就不會找我.
他們去馬會就不會找我.
那我就很老套.
有很多生意接不到.
很多這些實際的利益的考量去決定他們要不要信.
所以越是有身份地位的人.
越是覺得自己是有的人.
總是擔心信耶穌會招致太多虧損.
就越不容易定義去信耶穌.
這是耶穌常常慨嘆財主難進天國的原因.
為什麼窮人才容易有福音可聽.
有錢人自己一大堆家當.
一大堆夾錢的東西.
總之麻麻煩煩.
你要信耶穌是難的.
尼哥底部就是一個這樣猶豫不定.
三心兩意的人.
在二章最後一章.
二章23-25節.
當耶穌在耶路撒冷過雨節的時候.

$^{121}$有許多人看見他所行的神蹟就信了他的名.
有人見到耶穌所做的事就說信.
不過耶穌卻不將自己交託他們.
因為祂知道萬人,又不用不著誰見證人怎樣.
因為祂知道人心裡所存的.
我們很多時候我們說想信.
不過我們就婆婆媽媽,以矣.
今天說信,明天就有反賭了.
什麼都可以.
所以尼哥底部大概是其中一個.
耶穌不會將自己交付給對方的對象.
不過尼哥底部卻是在那些有身份有地位的猶太人中.
最敢於公開表白對耶穌的好感.
甚至是支持的人.
他在第七章45-52節說到.
他在公會裡曾經就耶穌而有的爭議裡.
為耶穌說話,支持耶穌.
當然不是說耶穌完全對.
他只是開放一點.
如果他不對的話自然有天收他.
他這樣說的那些.
但最少是為他說話.
然後在耶穌死了之後.
他出錢去買香膏去糕耶穌的屍體.
這是第十九章38-42節.
所以他對耶穌也算有好感的.
一直暗地裡支持耶穌.
我們只是不確定他最終有沒有成為耶穌基督的門徒.
如果他曾經為耶穌做過很多事.
卻最終沒有成為基督徒的話.
你覺不覺得很可惜.
臨門脫腳.
這樣輸了.
在晚上.
尼哥底沒有來找耶穌.
有人估計他應該是被記的緣故.
擔心被人撞見.
所以才晚上偷偷來.
你可以這樣解釋.
不過也可以不用弄得這麼複雜.

$^{161}$因為猶太人的拉比.
總是習慣了晚上吃完飯才去吹水.
他們吹水的時間就是晚上吃完飯的時候.
現在吃飯的時間沒有我們這麼晚.
他們通常下午四五點就吃飯.
通常吃兩餐.
吃完晚飯就沒什麼好做.
就吹水.
以前大陸鄉下就是這樣.
晚上吃完飯就搬張椅子.
拿個大攜扇.
走去榕樹頭底下.
大家坐下圍著吹水.
不需要解釋.
他是偷偷摸摸的.
晚上吹水是很慣常的事.
耶穌基督和尼哥底姆展開一場宗教的對話.
談論的主題是重生.
不過你會看到.
這不是一個勢力均等的對話.
因為一切有內容的說話.
都是出自耶穌的口.
尼哥底姆只是扮演問問題的傻子的角色.
是嗎?是嗎?.
重生這個題目.
是耶穌親自主動提出來的.
並且得如其來提出這個主題.
尼哥底姆來找耶穌.
他首先道明來義.
並且他吹捧了耶穌一番.
他指出他相信.
耶穌是從上帝差來的人間的教師.
擁有上帝所賜的智慧和能力.
這是他如今.
如此堂德來拜訪耶穌的原因.
不過在他讚賞了耶穌.
還未提出他的問題.
他應該帶著問題想問.
還未提出他的問題之前.
耶穌就主動提出重生這個話題.

$^{201}$耶穌沒有給尼哥底姆提問的機會.
他對尼哥底姆的讚美支持都沒有回應.
他卻是扣著尼哥底姆說.
耶穌是由上帝那裡來作師父.
這句說話去作回應.
尼哥底姆確認耶穌是師父.
他來找耶穌的目的很明顯.
就是求道.
希望在耶穌身上有所學習.
耶穌既然是師父.
他就自認做徒弟.
徒弟找師父做什麼?.
學習.
對不對?.
你要讚尼哥底姆.
他應該是一個上了年紀的人.
否則他不能進入公會.
他應該有相當社會地位.
作為一個猶太教律法師.
他自己是一個拉比.
他是一個教師.
但他願意放下身段去找耶穌談道.
這真是一個虛心求道之心迫切的人.
他不覺得自己已經擁有完備的真理.
他繼續追尋.
並且耶穌比他年輕.
耶穌三十歲.
尼哥底姆一定是四十歲以上.
否則做不到長老.
不是長老.
猶太人四十歲以上叫老年人.
所以他一定過了四十.
所以一個年長的人對一個年輕的人.
並且是一個沒有出身的人.
耶穌就算不是滿腳牛屎.
也是加里尼基層出身的人.
沒有受過嚴格的經學訓練.
尼哥底姆願意虛心求教.
這真是很難得的事.
對猶太人來說.

$^{241}$求道不純粹是為了增長知識.
滿足你的求知慾.
更加是為了明白上帝的旨意.
從而得見並且進入上帝的國.
你記得勝經的說法.
敬畏耶懷智慧開端.
認識異性者便是聰明.
這個聖經真理.
所有拉比都知道.
所以你的認識和你對上帝的敬拜.
和他建立關係是連在一起的.
知識不只是知識.
知識和生命和我們的求道.
是有密切的關係的.
所以認識自聖者.
和自聖者建立終極的關係.
這是人生最高的追求.
耶穌突然提出重生這個主題.
不是無厘頭而是一針見血.
切中要害的.
他要告訴尼哥底姆.
你單單尋覓真理.
是無法讓你見到上帝的國的.
唯有重生的人才能見到上帝的國.
換言之尼哥底姆.
你現在去找耶穌.
不僅僅是想聽師父的一句話.
掌掌知識增加智慧.
你一定要預備好.
你的生命被徹底翻轉.
徹底改變.
第三節.
我實實在在的告訴你.
人若不重生.
就不能見上帝的國.
尼哥底姆說.
他從耶穌所行的神蹟裡面.
見到耶穌的屬天身份.
耶穌卻指出.
沒有人能夠憑肉眼.

$^{281}$就看到屬天的事物.
他們見不到上帝的國.
我實實在在的告訴你.
在尼哥底姆和耶穌的短短的交談裡.
耶穌三次說了這句短言.
證明耶穌是用非常嚴肅的態度.
去說一個非常重要的真理.
我實實在在.
可以很嚴謹的告訴你.
上帝的國.
指著一個將來全生永恆的國度.
耶穌曾經在馬達福音19:28說.
這是一個復興的時候.
一切都要更新的時候.
是一個萬物都要更新的世代.
能夠在其中存活的人.
都是更生了,重生的人.
重生這個字是elephant.
這個字可以解作從上而來.
born from above.
也可以解作再次出生.
born again.
耶穌這句話.
可以同時有這兩方面的含義.
不過重點一定不是born again.
而是born from above.
從上而來.
來源重要.
當你說再次的時候.
你只是說第二次.
你說第二次就即是說.
和第一次一定有些關係.
你會很自然的想起第二次.
這間餐廳我來的第二次.
旺鎮我第二次來了.
即是我第一次來過.
然後我來第二次.
兩次是同性質的.
所以我再做.
但是來個底無.

$^{321}$即是將重生完全看為again.
再次.
於是他立刻很困惑.
他說人已經老了.
如何能重生呢?.
豈能再盡母腹生出來呢?.
他將重生理解為再生一次.
並且和第一次生的時候一模一樣.
都是從母的子宮出來.
我這麼大塊.
我已經這麼老.
我如何能塞進母的肚子裡.
再生過來一次呢?.
是不是?.
耶穌沒有直接糾正.
來個底無的看法.
不過在回應他的時候.
再次強調.
關鍵不是在於再次.
是在於從上而來.
從上帝那裡來.
我實實在在的告訴你.
人若不是從水和聖靈生的.
就不能進上帝的國.
從肉生生的就是肉生.
從靈生的就是靈.
第五第六節.
在這裡耶穌強調重生的憑藉.
是水和聖靈.
並且將見到上帝的國.
在這裡就改為進上帝的國.
見上帝的國和進上帝的國.
是差不多意思.
從水和聖靈生.
多數學者都相信.
水是指日後基督教的洗禮.
而不是當時的猶太人.
或者施洗約翰.
所做的所謂悔改的禮儀.
約翰寫福音書的時候.

$^{361}$早期教會是非常注重水禮的.
視為做基督徒.
視為加入教會的起點.
耶穌在離世前的大使命.
也都吩咐門徒為眾人施洗.
所以水禮是新生命的開始.
至於水和聖靈兩者有什麼關係呢?.
其實都不容易確定.
前面一章33-34節.
施洗約翰曾經將水和聖靈連在一起說.
指出他只是用水替人施洗.
不過耶穌就用聖靈替人施洗.
所以耶穌見過他.
他玩鞋帶.
當然鞋帶都不配.
他用水施洗.
耶穌用聖靈施洗.
所以聖經的意思可能就是.
我們藉著參與水禮.
而標示我們和耶穌基督同死同埋葬同復活.
而我們都是在接受水禮的場合裡.
我們一起接受聖靈.
就像耶穌自己一樣.
不過耶穌說的話的重點.
就在聖靈而不是在水.
他強調從上而得的生命.
我們是藉聖靈所生.
然後他再也沒有提到水.
只集中說聖靈.
從靈生的就是靈.
從肉身生的就是肉身.
從靈生的就是靈.
人間的生命.
跟屬靈的生命.
兩碼子的事.
是不會混淆的.
有學者很簡單地解釋這句話.
There is no evolution from flesh to spirit.
這是一個很簡單的解說.
你不要期望從肉身可以轉變到靈.

$^{401}$兩者根本不搭配.
不相配.
肉身就肉身.
就不能變,不能改.
這裡指出有兩個出生.
首先我們是由母腹所生.
所以是擁有肉身的人.
接著我們有一個超自然的出生.
是藉聖靈而生的.
因為從聖靈生.
所以就是屬靈的生命.
強調從聖靈生的.
就是強調從上而來的生命.
上面,聖靈是指上帝.
不一定要說聖靈充滿.
聖靈的駛.
不一定有這方面的特別含義.
不是說一個特別的經歷的含義.
因為這裡說從聖靈生.
但是在14-15節.
耶穌的說法就是.
從聖子耶穌得生命.
他怎麼說?16節說.
摩西在曠野怎樣舉蛇.
人子也必照樣被舉起來.
結果就是信他的,也得永生.
是藉著耶穌基督得永生.
而16節,雖然不是我們今天讀的經文.
不過你是整本聖經最熟的一次聖經.
若不懂就要打屁股.
上帝愛世人,甚至將他讀生子,至急他門.
結果就是信他的,不致勿忘.
反得永生,也強調得永生.
不過就是由父上帝主導.
所以連續三節聖經.
前面說聖靈讓我們得永生,得生命.
接著說耶穌基督被舉起來,使我們得生命.
接著就說父上帝.
施展的救贖使我們得生命.
所以既是聖靈使人得生命.

$^{441}$又是聖子使人得生命.
又是聖父使人得生命.
賜人生命的是三一上帝.
所以就不需要特別說.
聖靈有什麼特別的含義.
有什麼特別的經驗,不一定的.
尼哥底姆不明白耶穌的意思.
耶穌接著就說.
我說你們必須重生.
你不要以為稀奇.
風隨著意思吹.
你們聽見風的響聲.
卻不曉得從哪裡來往哪裡去.
凡重生的也是如此.
耶穌看到尼哥底姆對重生這個道理.
感覺到難以消化.
就用比喻說.
就像我們無法掌握風的來處和去向.
不知道風從哪裡來.
但我們聽到風的聲音.
看到風所產生的影響.
你看見樹葉搖動.
一樣我們無法理解和掌握.
聖靈給人生生命.
這件事是怎麼發生.
但我們卻看到.
有生生命和沒有生生命.
兩者是有顯著的差異.
看到舊生命和新生命的差異.
你看不到聖靈.
但聖靈的作為是毫不含糊.
是清楚可見.
就像今天我們常常在教會.
在我們基督徒的生活裡.
常常說聖靈的工作.
你見到聖靈嗎.
我們沒有人見到.
曾經在聖經裡有人見到.
聖靈好像火一樣在你的頭上.
你的頭有沒有著過火.

$^{481}$我沒有試過.
對不起.
就算我有聖靈充滿.
我也沒有試過.
有火燄在我的頭上.
我沒有見過.
不過你問聖靈是否很虛.
對我們基督徒來說不是.
我切實經歷過.
他在我生命裡的改造.
帶領幫助.
尼哥底姆無法消化.
耶穌這個這麼嶄新的道理.
在第十節.
耶穌就略帶嘲諷的口吻說.
你是以色列的先生.
你是以色列人的先生.
還不明白這是什麼.
這麼淺的道理你都不明白.
你這個猶太拉比的位置.
是怎樣混回來的.
混吉地混回來.
不過耶穌沒有繼續.
跟尼哥底姆做理性的辯論.
他卻是從自己與眾不同的身份.
指出我是從天上來的那一位.
所以我知道天上的事情.
而我說有關天上的事情.
一定是正確的.
重生是從天上而來的新生命.
所以是天上的事.
不過他也在人間發生.
所以同樣是地上的事.
尼哥底姆是凡人.
沒有升過天.
所以不知道.
不明白天上的事.
也沒有什麼奇怪.
只是他必須相信耶穌所做的見證.
因為耶穌是唯一從天上來的那一位.

$^{521}$十三節.
從天降下.
凝究在天的人子.
沒有人升過天.
凝究在天這幾個字.
主要的古本都沒有.
所以我們不解釋這幾個字.
從天上降下就行了.
從天上降下的人子.
沒有人升過天.
耶穌基督是唯一的天外來客.
他不是曾經上過天.
看見天堂的奧秘.
你經常在坊間都會買到一些書.
有些人死了一陣子.
上過天堂門口然後又走下來.
那些是進去天堂做訪客.
他本來不是屬於天堂的.
耶穌本來是在天上.
而他去人間.
他不是偶然上去做訪客.
他是下來人間做訪客.
所以耶穌是唯一一個在天上.
看見,知道天上的事.
他是唯一有資格講論天上事情的人.
在這短短的說話裡.
耶穌透露了他兩方面的秘密.
第一,他的屬天性.
就是他的神性.
他是從天上來.
也是從父上帝那裡來的.
第二,他的拯救使命.
這是耶穌第一次在人前.
用民數記21章摩西舉蛇的故事宣告.
他要被釘在十字架上.
他要被舉起來.
好讓信他的人得永生.
11節.
我實實在在的告訴你.
我們所說的,我們知道的.

$^{561}$我們所見證的,我們見過的.
你們卻不領受我們的見證.
這裡用了「我們」.
我們所說.
我們所見證.
我們見過.
說明不僅僅是耶穌自己.
也包括他的門徒.
「我們」.
有學者相信.
這次聖經不僅僅是耶穌對尼哥底姆說的話.
也都是日後初期教會.
對猶太教公會.
或者有拉比背景的.
有律法訓練背景的人說的話.
我們在你們面前做的見證.
我們已經看完整段聖經.
整段聖經的中心教訓是第七節耶穌的說話.
你們必須重生.
這句話的主詞是眾數.
說明.
這不僅僅是耶穌對尼哥底姆的教導.
「你要重生」.
也都是對所有人的要求.
「你們要重生」.
耶穌指出.
人要進上帝的國.
首要的條件是從上帝而來的生命.
生命改變是門檻.
生命不改變.
只是增長知識是沒有用的.
正如我以前經常說的話.
好像說得很老套.
不過這是一個關鍵.
基督信仰永遠不是在你舊有的知識.
舊有的經驗之上做一些添加.
無論你的宗教知識是猶太教還是中國文化.
都不是.
都沒有用.
我常常說.

$^{601}$基督教信仰不可以是補充性.
一定是subversive(反覆性).
是顛覆性.
你不可以繼續維持你過去幾十年有的知識經驗.
然後加一個基督信仰下去.
錦上添花.
在你有的眾多的東西裡面.
加一個基督教.
No way.
沒有這回事.
基督教很霸道.
很兇的.
你要將基督形入你的生命裡面.
就麻煩你反轉了它.
你先拆掉你所有的東西.
捨棄拆毀是第一步.
一粒墨紙如果不掉在地上.
裡面死了仍然是一粒.
凡捨棄生命的就得著生命.
你以為你想保住生命.
你就限電.
你永遠得不到生命.
基督教不會跟你雞兔同籠.
不會在你既有的東西之上加一個信仰.
它不是一個supplement(補充).
你吃完三餐主餐.
再拍幾粒藥丸.
去補充補充.
你不會用藥丸做主菜.
你不會用飯吃.
你吃完所有飯菜.
然後加些營養劑.
加些補品.
基督教不是這樣的東西.
基督教不是supplement(補充).
基督教是subversion(反轉).
一個subversion(反轉).
是顛覆.
徹底的將你反轉顛覆.
所以你有幾基督徒.

$^{641}$這是他要問的問題.
How Christian you are(你是什麼樣的基督徒).
你有幾基督徒.
好像很古怪的說法.
你有幾被反轉.
你和過去有幾一刀兩斷.
你有幾狠的切斷你的過去.
就看你的信仰有幾真實.
我們舊的生命和新的生命.
不是一個連續體.
而是徹底斷裂彼此衝突.
好像保羅說成為基督徒之後.
我們同時有兩個我.
一個是舊人.
一個是新人.
歌羅西書三章九到十節.
不要企圖將舊人和新人握手言和.
二巴很難和好.
我們只能夠致死舊人.
努力擺脫舊人舊人的行為.
穿上新人.
生命更生.
生活才能更生.
重生是每一個基督徒都要確立的事實.
重點不在於再出生一次.
而是從上帝那裡得到屬靈的生命.
弟兄姊妹.
我不知道重生這個詞對你有幾多意思.
這個詞今天很常用.
用得很濫.
我上網搜尋.
原來大陸劇,韓劇,流行曲.
都有用這個名字做主題.
重生有大陸劇.
重生有韓劇.
未來這個星期看看劇.
不過我們說的重生沒什麼關係.
我們看穿越小說的時候.
最常用這個詞.
不知道你喜歡看網絡小說嗎.

$^{681}$凍結幾百萬字的小說.
總之說重生.
二十幾歲三十幾歲撞車.
死了馬上去唐朝.
做了唐朝人做了明朝人.
他們就叫重生.
重生意味著獲得一個新的生命.
從頭再來過再活一次.
這是重生的意思.
不過.
如同耶穌糾正尼哥底母的想法.
重生不是讓你重複塵世的生命.
在母腹中再誕生一次.
而是你要從上帝那裡.
獲得一個屬靈的生命.
你成為新造的人.
塵世的生命即或仍然存在.
對你再無主宰權.
亦都再無必然的延續性.
我們說哥林多後書的時候都很鄭重說過.
若有人在基督裡.
他就是新造的人.
舊事已過都變成新的了.
所以不要將以前的事混在一起說.
不要告訴我.
因為我爸爸卑鄙所以我卑鄙.
因為我媽媽軟弱所以我軟弱.
不要說因為我有DNA.
令我不能不肥.
不能不胃食.
不可以說我原生家庭一定可以.
讓我必須是這樣.
不是這樣沒有任何事再是注定.
今日我們的生命開了一個新的章節.
一個完全新的一頁.
過去所有事再無主宰權.
舊事已過都變成新的了.
我們有新的身份新的關係.
新的人生目標.
新的倫理責任.

$^{721}$過去一段很長的時間.
橫教會深受約翰.衛斯理.
乃至他所影響的.
中國教會的重要人物信尚哲.
和一班憤青報道家的影響.
總是將重生解說在一個特殊的經歷.
基督徒需要在某一個時刻.
經歷聖靈充滿.
讓自己的罪污得以完全潔淨.
然後脫胎換骨.
不更生.
覺得自己救我已死.
生命被翻轉這樣的經驗.
你有沒有聽過這些教導.
你沒有就證明你年輕.
在我十幾歲回教會的年代.
那些牧師來講道.
橫斑斑就問你重生了沒有.
即使是執事員都照樣問你.
你是牧師都照樣問你.
你顯然的說我應該重生了.
即是還沒有.
對他們來說.
重生是一個事件.
你要講得出.
我是2023年10月14日.
10點45分.
在旺鎮的正堂重生.
你講不出這個時間地點.
即是還沒有.
因為那個不是理念.
那個是一個事件一個經驗.
他們要求.
如果你還沒有這個經驗.
證明你信仰是不行的.
所以你要追求重生.
你沒有聽過這些理論.
老一輩都會聽過.
這個是以前很流行的說法.
重生是一個事件一個經驗.

$^{761}$怎樣追求呢.
他們的說法就是.
你請三天假.
然後拿一疊紙.
現在用電腦.
躲進房間.
最好買足夠的麵包糧水.
不要再出來.
躲進房間裡面.
然後將你由三歲.
即是有記憶開始.
你犯過的罪.
逐條逐條認.
逐條逐條寫下.
寫了十幾頁紙.
四十幾頁紙.
五十幾頁紙.
我真的犯罪累累.
然後你一路認一路害怕.
一路覺得.
我自己真的卑鄙無恥.
下流賤格.
為甚麼我是一個.
這麼衰格這麼衰格的人呢.
然後你切齒痛恨你自己.
然後你崩潰.
罪人赦免.
然後你經歷一陣聖靈的大光.
然後聖靈光照你.
然後就將你.
好像用一個鮑魚刷一樣.
然後回到內之外.
洗得乾乾淨淨.
你自己就變成一個脫胎換骨的人.
當你推出房門的時候.
理論上你的臉應該會閃光的.
人家瞎也見到你是.
已經重生了的人.
你有沒有聽過這些理論.
沒有嗎.

$^{801}$我們真的有generation gap(世代差別).
上一輩很流行說這些.
很流行.
所以一定講聖靈.
重生是一個經驗.
重生是一個特殊的經驗.
不過我必須要講.
重生是一個特殊的經驗.
這個說法.
我講得很盡.
完全沒有聖經根據.
完全.
重生這個觀念不是約翰獨有.
是遍及新約聖經的各處.
例如提多書三章五節.
比利前書一章三節二十三節.
約翰一書二章二十九節.
三章九節四章七節五章一節.
四節十八節.
你值得注意的是.
所有新約出現過有關重生的教導.
都只是講出.
我們被上帝更新了生命的事實.
沒有任何一處.
我再講沒有任何一處.
明事或暗事.
重生是一個特殊的事件或經歷.
沒有.
如果你找到你告訴我.
我就可以推翻我的說法.
約翰一書.
是提及基督徒由上帝所生.
最多的一卷書.
書中強調的是我們由上帝所生.
由上帝所生.
Born from above.
你是上帝的兒女.
你有上帝的生命.
而沒有說這是born again.
不是強調再次.

$^{841}$並且也說這是事實.
沒有任何含義說這是經驗.
耶穌和尼哥底姆談重生這個道理.
也沒有明事暗事.
重生是一個特殊經驗.
或必須藉著一個經驗才能獲得重生的事實.
只是指出這個進上帝國的人.
必要的條件.
並且耶穌用風隨自己的意思.
吹來比喻聖靈給我們的新生命.
也說明我們無法理解.
無法掌握這個事實是如何發生.
我們只能從效果上看到這個事實的存在.
如果重生是一個轟轟烈烈的經驗.
你就不會覺得難以判別它的存在.
保羅多次提到人接受水禮.
就是象徵得著一個新的生命.
所以受洗歸入基督.
與他同死同埋葬同復活.
保羅在提多書三章五節說.
他便救了我們.
並不是因我們自己所行的義.
乃是照他的憐憫.
藉著重生的洗和聖靈的更生.
他將重生連上水禮.
重生的洗.
就是象徵同基督同死同埋葬同復活的禮儀.
至於更生的工作是由聖靈做的.
不是在你接受浸禮那一刻就發生的.
不過我們在接受水禮的時候.
我們就得著重生這個事實.
或者我們在接受浸禮的時候.
我們確認我們已經生命改變的事實.
全本聖經都說明.
重生是一個事實不是一個事件.
重生意味著我們得著一個新的生命.
成為新造的人.
舊事已過都變成新的了.
這個不是一次的事件.
一次就讓我們有這些經驗.

$^{881}$就正如得舊是一個確定的事實.
你已經有了一個新的身份.
但是不等於我們馬上脫胎換骨.
變成無瑕疵的新人.
讓一書到去不斷提醒我們.
我們已經由上帝而生.
所以我們不可以再做跟這個新的身份.
新的生命不相稱的事.
不斷提醒你不要再做那些不相稱的事.
也表示你有可能.
你的生命改變但你的生活還沒變.
所以我們要努力使我們的行事為人.
跟我們所夢的因相稱.
保羅在哥倫西書三章九到十節.
勸導門徒脫去舊人舊人的行為.
也表示即或穿上了新人的我們.
仍然會受舊人舊人的行為影響.
不是一下子就變了.
我一直都困惑.
我說這樣的道理會不會太抽象.
不過事實上重新的道理就有點抽象.
讓我們確認我們已經是不一樣的人.
從肉眼看分不出來.
從理性分析解釋不到.
不過弟兄姊妹.
讓我們沉著氣.
放長雙眼去看長期的變化.
正如你掌握不到風.
但你看到風的影響.
我沒有任何自誇可言.
在你面前有什麼好自誇.
特別在我太太和女兒面前.
我怎麼自誇呢?我沒有自誇.
不過我真的要說.
我這個出身非常不好的人.
我是一個破碎家庭成長.
我從小沒有很好的教養.
我的脾氣很粗暴.
我真的不值得有一個幸福的家庭.
我不配有一個快樂的晚年.

$^{921}$那個不是一個血氣的我.
肉身的我.
就能夠走到今天.
我只能夠宣告.
我這麼奇怪的人.
可以擁有一個新的淑女身份.
淑女的生命.
這個不配合的.
不過我要說是千真萬確.
活了這把年紀.
當我回顧.
當我看到恩典處處.
當我看到我竟然可以.
活到這樣的地步.
有這樣的處境.
我不可以不說.
那個不是我做的.
是裡面的我幫我做的.
所以求主幫助我們.
我們在我們的具體生活裡.
看到裡面出來的能力.
裡面出來的抉擇.
然後你發覺.
竟然在不知不覺間.
你的氣質你的性格.
就慢慢被塑造被轉化.
然後你看到.
聖靈奇妙的作為.
不是我做.
是我裡面的那個做.
肅靈的事實.
不是我們三言兩語.
可以說出來.
不過如果你跟隨主.
你說你是基督徒一段長的時間.
你都沒有曾經經歷過.
你裡面的寶貝.
跟你外面不搭配的雅氣.
兩者的並存.
你都很糟糕.

$^{961}$無論如何.
今天第一步做的是.
記住我們擁有新的生命.
已經有新的生命.
新的身份.
所以一會兒如果.
你跟人爭位.
吃飯.
你想發脾氣.
想爆粗.
你提一提自己.
不可以爆了.
新的生命.
上帝轉化的生命.
不可以再做這樣的事.
上帝恩代我們每一個.
已經重生的弟兄姊妹.
\newpage



\section{士師記 21:1-15}
\label{sec:jwnc_J9h4aM}
\textbf{2023年10月29日崇拜直播｜陳冠成牧師｜「家醜說出口」系列:問題解決了嗎?｜士師記21\_1-15}
\newline
\newline
連結: \href{https://youtube.com/watch?v=jwnc-J9h4aM}{\texttt{https://youtube.com/watch?v=jwnc-J9h4aM}} ~~~~ 語音日期: 2023-10-30
\newline
\newline
\hyperref[sec:e5U96ffu0L0]{\small{< < < PREV SERMON < < <}}
~
\hyperref[sec:index_chronic]{\small{[返順時目]}}
~
\hyperref[sec:t62ZJvrM_ng]{\small{> > > NEXT SERMON > > >}}
\newline
\newline
士師記 21:1-15
\newline
\begin{longtable}{cl}
\hline
\hline
章節 & 經文 (和合本修訂版)\\
\hline
21:1 & \begin{tabularx}{0.7\textwidth}{X} 以色列人在米斯巴曾起誓說：「我們中誰都不把女兒嫁給便雅憫人。」 \end{tabularx} \\ \\ \relax
21:2 & \begin{tabularx}{0.7\textwidth}{X} 以色列人來到伯特利，坐在那裡直到晚上，在神面前放聲大哭， \end{tabularx} \\ \\ \relax
21:3 & \begin{tabularx}{0.7\textwidth}{X} 說：「耶和華－以色列的神啊，為何以色列中會發生這樣的事，使以色列今日缺了一個支派呢？」 \end{tabularx} \\ \\ \relax
21:4 & \begin{tabularx}{0.7\textwidth}{X} 次日，百姓清早起來，在那裡築了一座壇，獻燔祭和平安祭。 \end{tabularx} \\ \\ \relax
21:5 & \begin{tabularx}{0.7\textwidth}{X} 以色列人說：「以色列各支派中，誰沒有同會眾一起上到耶和華那裡呢？」因為以色列人曾起重誓說：「凡不上米斯巴到耶和華那裡的，必被處死。」 \end{tabularx} \\ \\ \relax
21:6 & \begin{tabularx}{0.7\textwidth}{X} 以色列人憐憫他們的弟兄便雅憫，說：「如今以色列中斷絕一個支派了。 \end{tabularx} \\ \\ \relax
21:7 & \begin{tabularx}{0.7\textwidth}{X} 我們既然向耶和華起誓說，必不把我們的女兒嫁給便雅憫人，現在我們該怎麼辦，使他們剩下的人可以娶妻呢？」 \end{tabularx} \\ \\ \relax
21:8 & \begin{tabularx}{0.7\textwidth}{X} 他們又說：「以色列支派中誰沒有上米斯巴到耶和華那裡呢？」看哪，基列的雅比沒有一人進營到會眾那裡， \end{tabularx} \\ \\ \relax
21:9 & \begin{tabularx}{0.7\textwidth}{X} 百姓被數點的時候，看哪，基列的雅比居民沒有一人在那裡。 \end{tabularx} \\ \\ \relax
21:10 & \begin{tabularx}{0.7\textwidth}{X} 會眾就派一萬二千名大勇士，吩咐他們說：「你們去用刀把基列的雅比居民連婦女帶孩子都殺了。 \end{tabularx} \\ \\ \relax
21:11 & \begin{tabularx}{0.7\textwidth}{X} 這是你們當做的事：要把所有男人和曾與男人同房共寢的女人全都殺了。」 \end{tabularx} \\ \\ \relax
21:12 & \begin{tabularx}{0.7\textwidth}{X} 他們在基列的雅比居民中，找到四百個未曾與男人同房共寢的處女，就帶她們到迦南地的示羅營裡。 \end{tabularx} \\ \\ \relax
21:13 & \begin{tabularx}{0.7\textwidth}{X} 全會眾派人到臨門巖的便雅憫人那裡，與他們講和。 \end{tabularx} \\ \\ \relax
21:14 & \begin{tabularx}{0.7\textwidth}{X} 當時便雅憫人回來了，以色列人就把所留下，基列的雅比活著的女子嫁給他們，可是還是不夠。 \end{tabularx} \\ \\ \relax
21:15 & \begin{tabularx}{0.7\textwidth}{X} 百姓憐憫便雅憫人，因為耶和華使以色列支派中有一個缺口。 \end{tabularx} \\ \\
[1ex]
\hline
\hline
\end{longtable}
$^{1}$我們一起來聽神給我們的說話.
這段說話是記載在.
《士師記》第21章1至15節.
請聽弟兄讀出.
以色列人在米斯巴曾起誓說.
我們都不將女兒給變雅敏人為妻.
以色列人來到伯特利.
坐在神面前直到晚上.
放聲痛哭說.
耶和華以色列的神.
為何以色列中有這樣缺了一支派的事呢?.
此日清早百姓起來.
在那裡築了一座壇.
獻繁祭和平安祭.
以色列人彼此問說.
以色列各支派中.
誰沒有和眾人上到耶和華面前來呢?.
先至以色列人起過大誓說.
凡不在米斯巴到耶和華面前來的.
必將他們致死.
第六節.
以色列人為他們的弟兄變雅敏後悔說.
如今以色列中絕了一個支派了.
我們既在耶和華面前起誓說.
必不將我們的女兒及變雅敏人為妻.
現在我們當怎樣辦使他們剩下的人有妻呢?.
又彼此問說.
以色列支派中誰沒有上米斯巴.
到耶和華面前來呢?.
他們就查出.
基列雅比沒有一人進營到會中那裡.
因為百姓被數的時候.
沒有一個基列雅比人在那裡.
第十節.
會中就把發一萬二千大勇士.
吩咐他們說.
你們去用刀將基列雅比人.
連婦女帶孩子都擊殺了.
所當行的就是這樣.
要將一切男子和已嫁的女子進行殺戮.

$^{41}$他們在基列雅比人中.
遇見了四百個未嫁的處女.
就帶到迦南地的侍羅營裡.
全會中打發人到臨門岩的變雅敏人那裡.
向他們說和睦的話.
當時變雅敏人回來了.
以色列人就把所存活.
基列雅比的女子給他們為妻.
還說不夠.
百姓為變雅敏人後悔.
因為耶和華使以色列人.
缺了一個支派.
聖經獨筆.
節目平安.
終於來到這個系列的結束一講.
上星期曾志牧師也和我們分享了上一章.
二十章.
你會看到其實.
以色列人原本他們只是打算做什麼.
來討伐住在變雅敏中間的基比亞人的匪徒.
原本只是打算這樣.
不過我們會看到.
變雅敏支派拒絕交出這些作惡的人.
結果就演變成.
以色列人和變雅敏支派的內訌內戰.
雖然上一章我們看到.
以色列人打贏了這場仗.
但是因為殺得過頭.
他們連變雅敏支派割成的人.
束併一切他們所遇見.
包括無辜的人全部.
通通都殺了.
並且他們放火燒毀了一切的誠意.
導致上一章最後就說到.
變雅敏人剩下多少個藍丁.
大家記不記得.
六百個.
厲害.
六百個藍丁.
即是說他們即將面臨一個滅族的危機.

$^{81}$所以緊接著的二十一章.
即是我們今天看的經文.
其實就記載了.
到底如何收拾這個爛攤子.
如何處理.
不過回到事件開頭.
我們回顧一點.
事件開頭是在十九章開始出現的.
你會記得.
最初是哪些人的問題.
我們說的是基比亞人.
住在變雅敏的基比亞人.
他們的淫行.
這是最初問題的開始.
由小撮人開始.
然後利未人亦有份參與這場惡行.
他將自己的妾侍推出去.
可以說是送死.
然後之後.
他有不盡不實的說話.
挑動到全部以色列人.
和變雅敏之派來對立.
再加上之後.
變雅敏人選擇護短.
不肯交出作惡的人.
導致一場內戰爆發.
最後因為以色列人.
本來是想懲治惡人.
現在變成了甚麼.
變成了濫殺無辜.
你會看到整件事.
我們說的就像滾雪球.
一路滾越大.
最初是一小撮人的問題或傷亡.
不正面處理.
不是靠上帝處理.
就會越滾越大.
越來越多罪惡發生.
越來越多傷亡出現在信仰群體內.
我們首先要看到這幅圖畫.

$^{121}$也是士師記作者.
想我們在後段看到的.
你會發覺以色列人.
他的敵人已經不是加拿大人.
是自己人.
不過無論如何.
我們說以色列人難辭其咎.
他都要為這場事件承擔責任.
正路他們應該怎樣做?.
如果按照上帝在聖經或律法的教導.
首先他們應該與上帝和好.
首先處理與上帝的關係.
因為很明顯整件事一定得罪了上帝.
亦傷害了上帝的百姓.
他們應該認罪悔改.
尋求上帝的寬恕接納.
這是第一件事應該做的.
第二件事當然是承擔所犯的罪.
所帶來的後果.
他有責任.
他需要作出補救.
很簡單就是將自己的女兒.
配給餘下的六百個辯雅文男丁.
幫他們繼後.
就可以重建辯雅文之派.
亦可以修補群體之間的破裂關係.
彼此復和.
正路是這樣.
不過剛才你看到的經文是否這樣?.
很明顯不是呈現這幅圖畫.
你見到沒有撥亂反正的出現.
反而是又再一次有殺戮的場面.
為何會這樣?.
其實《士司記》的作者不斷在.
最後幾章講述原因.
特別在這章結尾他講述了.
可能大家很熟悉.
我們說一讀《士司記》.
你一定會記得.
有句關鍵的詞語.

$^{161}$「那時以色列人會有王.
各人任意而行」.
記載了在21章的結尾.
或者更正確的解釋.
「各人任意而行」是甚麼意思?.
其實是說.
每一個人都照自己眼中看為對的事去做.
那怎會不亂?.
一定會亂的.
事實上你會看到.
《士司記》的作者很想告訴我們.
原來當上帝的子民.
他只看到自己對的事來做的時候.
他不再尊上帝為王.
不再聽從上帝的吩咐來生活的時候.
原來即使他們有信仰也好.
他們的道德光景可以很敗壞.
敗壞到甚麼地步?.
我們就看今天的經文.
你可以打開程序表的港道大綱.
印了一些相關的經文在當中.
以色列人其實真的有為這件事.
來後悔,表示難過.
還要不止一次,是三次.
印在當中的.
第二節,三節,第六節,第十五節.
他們有表示難過,甚至後悔.
不過你看到原因是甚麼?.
是以色列少了變雅文這個支派.
不知道大家聽起來會否覺得有點怪怪的.
覺得怪怪的嗎?.
類似甚麼情況呢?.
譬如家裡有個人很難賭的.
到處向人借錢,借彈,耳窿.
搞到神憎鬼厭.
經常有人上門追債,電話騷擾.
於是所有家人都怕了他.
要四散離開這家.
因為這個原因,所有家都四散了.
如果這個難賭的人真誠悔改.

$^{201}$他大概會怎樣說呢?.
他會說其實都是我不好.
是我難賭,到處向人借債.
令債主臨門.
現在我弄到所有家都四散了.
只剩下我一個.
他應該這樣反省自己的問題.
而不會只是說,我現在很慘.
家裡只剩下我一個,所有人都走了.
他不會這樣說,彷彿對自己.
一點責任都沒有.
彷彿說到自己不是罪魁禍首.
他不會這樣說的.
同樣以色列人都一樣.
如果以色列人當下真誠悔改.
他應該怎樣跟上帝說呢?.
他應該說,上帝.
是我們做得不好.
是我們煞得性起.
是我們摧毀了變亞敏茲派.
令到以色列現在變得不齊整.
起碼他應該這樣承擔自己的罪過.
而不會反過來只是問上帝.
神啊,為何現在以色列少了一個支派.
神啊,很慘.
為何你令到以色列失去了一個支派.
說到自己好像沒問題.
說到自己沒有任何責任.
你看到他三次悔改或所謂難過.
其實一字都沒有提及自己之前所犯下的罪行.
很清楚經文交代了.
很明顯你會看到.
百姓將原本自己作惡加害者的身份.
現在重新包裝轉成.
我是受害者.
我只是一個受害者.
這樣自然就會將自己要負的責任.
外化了,推走了.
將自己原本要承擔的罪咎.
慢慢地洗走.

$^{241}$所以你會看到他眼中.
再看不到自己有甚麼問題.
再看不到自己有甚麼責任.
他的後悔和難過.
不是因為他察覺到自己做錯了.
不是因為他察覺到自己傷害了人.
得罪了上帝.
所以認罪悔改.
他的後悔和難過.
其實只是一種自憐的心態.
我們少了一個支派.
現在多慘,不齊整了.
他覺得自己是受害者.
不是搞事的那個.
所以你看到他眼中.
不再看到別人的淒慘.
不再看到上帝的受虧損.
他只看到我們這個支派受虧損.
少了一個手足.
他也不再看到.
原來自己才有問題.
他只看到是其他人搞成這樣.
你會看到來到時思記的末段.
上帝的子民已經完全喪失了.
那個我們說「醒察己身」的能力.
那個承擔責任的能力.
不見了.
因為沒有了上帝.
中國人常有句說話.
「每日三醒五身」.
大家都聽過了.
其實意思就是每日.
「醒察三件事」.
我們看別人做事有沒有盡心竭力.
我們和朋友相處有沒有言而無信.
我們師長所傳授的道理.
我們有沒有認真的實踐.
說這三件事.
他提醒我們無論處事.
待人或者追求知識學問.

$^{281}$其實你都要時常反省自己.
如果有錯的你要去承擔.
你要去改正.
如果沒有的就當是一個面例.
不過.
說真的.
大家覺得現在這個價值.
有多少人還願意去擁抱呢?.
多不多?.
有人搖頭.
似乎你會感覺越來越少人.
擁抱這種價值觀.
相反.
埋怨推卸責任這種文化.
似乎越來越流行和普及.
確實我們心底裡都會明白.
找別人的錯處.
將責任推給別人.
其實一定會比舒服.
要面對自己的過失.
會安心一點.
所以當犯錯的時候.
確實有很多人第一時間的反應.
是會選擇.
不關我的事.
是別人的問題.
而不會去正視自己.
不會去反省己身.
從而承擔責任.
不過弟兄姊妹.
聖經提醒我們.
其實作為上帝的指引.
我們不可以這樣.
上帝期望我們.
全本聖經都經常說.
你要時刻審查.
你要敏銳聖靈對你的提醒.
聖經對你的督責.
上帝希望我們時常反省.
祂交給我們.

$^{321}$託付我們做的事情.
我們有沒有認真做好.
我們有沒有誠實無偽地.
去對待善待身邊的人.
又或者我們有沒有忠心地.
將聖經的教導實踐出來.
如果上帝的指引.
只是有例行的宗教活動.
但卻是欠奉了.
在日常生活裡面.
對信仰的反省.
和承擔的能力的話.
很有可能.
他們就會好像.
當時時事記的百姓一樣.
他們表面上仍然有宗教生活.
仍然敬拜上帝.
但實質上已經.
和上帝的背道而馳.
我們再看下去.
第四節這裡.
哭了一晚.
有後悔.
但第二天已經怎樣?.
百姓繼續向上帝行禮而喻.
留意他們有獻祭.
但獻什麼祭?.
正路應該是贖罪祭.
對不對?.
但他們現在獻的是.
凡祭和平安祭.
他們和上帝的關係.
彷彿已經雨過天青.
沒有事了.
所以可以如常地.
向上帝展示友好.
期望上帝月立.
繼續保佑他們順利和平安.
然後怎樣?.
繼續任意而行.

$^{361}$自己想想辦法.
怎樣幫助.
剩下的六百個南丁.
去找老婆.
整個圖畫就是這樣發生.
所以我們繼續看下去.
第五至九節這裡記載.
剛才也說過.
以色列人很清楚.
最簡單直接幫助.
變雅敏之派續後的方法.
就是將自己的女兒.
嫁給他們.
符合律法的要求.
亦可以讓他們自然地.
開枝散葉重建這個支派.
不過他們發覺.
執行起來時.
他們遇到一些困難.
是甚麼呢?.
經文這裡也交代了.
原來在他們攻打.
變雅敏之派之前.
他們曾經在米斯巴這個地方.
很衝動,很魯莽地.
發了兩個勢.
第一個就是.
凡是以色列支派的.
如果沒有上來米斯巴這個地方.
來表態.
表示要參與這場爭戰的話.
即是不同心.
如果凡是有這些人.
一律將他們致死.
這是第一個勢.
另一個就是.
他們為了要和變雅敏之派劃清界線.
當他們已經是仇人.
於是他們又起誓.
我們怎樣也不會將我們的女兒.

$^{401}$嫁給變雅敏人為妻.
他們發了這兩個勢.
所以現在.
他們如果要將自己的女兒.
嫁給變雅敏之派.
即是怎樣?.
即是違反誓約.
有後果嗎?.
有啊.
他們自己也說了.
會有咒詛可能淋到他們身上.
但如果不這樣做.
又怎樣?.
確實真的不知道怎樣.
幫變雅敏之派來繼後.
那怎樣辦呢?.
自己令自己落入一個兩難.
正如我們也說過.
應該是凝聚悔改.
求上帝寬恕.
推翻自己亂起的勢.
求上帝憐憫.
然後按照律法的吩咐.
將女兒嫁給他們續後.
但當然不會這樣.
因為已經沒有上帝了.
所以就任意而行.
怎樣任意而行?.
他們很厲害.
真的鑽空子鑽得很厲害.
想盡辦法不過是做壞事.
他們在兩個的誓約中間.
找到一個漏洞.
我們說是一條很邪惡的計謀.
是怎樣呢?.
一方面他們要鑽空子去找.
到底有沒有人沒有出現在米斯巴呢?.
如果有的話.
即是說他沒有參與這場起誓.
既然他沒有參與這場起誓的話.

$^{441}$即是說他們的女兒可以怎樣?.
可以嫁給變雅敏人.
但又沒有違反誓言.
這是他們第一個鑽到的漏洞.
另一方面他們也會想到.
既然如果有些人真的沒有出現在米斯巴.
那代表什麼呢?.
他們不同心.
這些不是自己人.
按照我們所起的誓.
我們要將他們致死.
既然如此.
他們也不可以干預我們怎樣處理他們的女兒.
毒不毒?.
其實也很狠毒.
你會見到他們沒有上帝的時候.
可以變得這麼邪惡.
於是他們覺得.
整件事已經解決了.
就立即經文記載.
以色列人急不及待去找代罪羔羊.
幫自己揹鑊.
他們找到一群叫做.
基列雅比人的.
原來真的沒有上到米斯巴.
來起誓.
於是隨即就打發一萬二千大能的勇士.
吩咐他們.
總之你去到當地.
見到有未出嫁的處女.
通通帶走.
至於其他人.
包括男人.
包括婦女.
包括小孩.
總之通通殺了他們.
結果.
經文裡很簡單地記載了.
他們去到基列雅比.
找到四百個未出嫁的處女.

$^{481}$就將他們帶走.
交給變雅敏人為妻.
其餘的基比雅人無一倖免.
慘遭殺害.
這件事再一次凸顯到.
當百姓心裡沒有上帝的時候.
他們的心態.
他們的行為.
可以變得有多麼扭曲.
多麼粗暴.
和令我們心寒.
這種邪惡的行為.
我即時想起一個電視劇集的人物.
以前有一套電視劇.
叫做《大時代》.
有沒有看過?.
就算你很年輕.
年輕的點頭就是看重播.
很明顯.
真的重播過.
應該也有幾次.
裡面有個奸角.
我還未說你知道是誰.
誰?.
裡面有個奸角很討人厭.
也有幾次《士司記》中描述的心態.
是誰?.
鄭少秋飾演的那個.
現在很當做的蟹.
丁蟹.
記不記得這個人?.
有沒有印象?.
丁蟹.
丁蟹有一個特色.
他每一次做錯事.
就算打死人也好.
他都會怎樣?.
他都會委過於人.
他永遠將自己的問題.
歸咎於別人身上.

$^{521}$你留意一下.
你有機會看回這套劇.
他一定說不是我的錯.
是對方的錯.
是他衰.
所以他活該.
你猜他真的.
不知道自己有問題嗎?.
其實你看起來他知道.
不過他永遠不想承擔這個責任.
他只想逃避問題.
於是他總是會怎樣?.
賴對方不好.
甚至拖對方下水.
幫他做替死鬼.
我之所以搞成這樣.
不是我有問題.
是對方有問題.
是社會的錯.
將自己的問題責任.
推得一乾二淨.
其實士司機百姓的情況.
都像是這樣的.
總是不肯面對.
其實誰是壞蛋.
自己才是壞蛋.
總是不想承擔責任.
還要反過來.
你會看見他會拖人下水.
強行誣陷別人做錯了.
誣陷別人.
替自己的罪來買單.
我們說這種.
很奇怪的心態.
很扭曲的心態.
總是透過做壞事.
來解決問題.
問題是否真的可以解決?.
可以嗎?.
你看到.

$^{561}$不可以.
沒有因為他們做完一場壞事.
殺了那麼多人.
問題就解決了.
事情未告一段落.
他們發覺死了.
剛才我們說.
抓了多少個基比亞的處女?.
400.
大家會加減數了.
600減400.
即是欠多少個?.
欠200.
怎樣做好?.
正路.
我們又會覺得.
應該是煞停自己.
想一想.
是否真的要繼續?.
又來?.
又去搶?.
但這不是士思記結尾讓我們看到的.
當人的心裡沒有上帝.
沒有王館的時候.
就會索性怎樣?.
一不做.
二不休.
繼續用邪惡的手段.
來解決自己的問題.
其實就像電視劇集.
很人性.
當沒有上帝的時候.
人性就是這樣黑暗.
你看到.
之後的經文.
雖然今天我們不會詳細說16節打後的經文.
留待大家在小組查經時.
再慢慢去研讀.
我會簡單說.
經文印了一點.

$^{601}$印了第19至22節.
簡單說這段經文是說.
以色列這次不只是普通百姓.
是長老.
長老建議.
他們去教唆那些.
還未找到老婆的變亞文人.
去搶回來.
你們去搶.
不是我去搶.
你們去搶.
教他們趁有節日的時候.
埋伏在葡萄園.
等那些侍羅女子.
到節日出來慶祝.
在那裡跳舞的時候.
趁他們沒有防備.
在葡萄園裡撲出來.
你們自己一人搶一個.
搶回去做老婆.
另一方面.
他知道一定會有人反抗.
被搶了女兒的家庭.
一定會出來反抗.
這班長老就幫他們.
擺平他們.
你放心.
我們會用我們的身份.
用我們的地位.
強行要他們接受.
這個迷而成趣的事實.
你放心.
你儘管去搶.
我幫你們善後.
整件事最後就是這樣.
他們覺得解決了.
弟兄姊妹可能我們會很奇怪.
不過很值得.
我們去想一件事.
雖然我們知道他們.

$^{641}$做的事真的很邪惡.
很令人不知.
但我們很值得去想.
到底人.
可不可以通過做壞事.
來換取好的結果呢?.
可不可以?.
我們可不可以透過做壞事.
來換取達至一個好的結果呢?.
我們正常覺得不可以.
但他們覺得可以.
他們就是這樣做.
他們很想不斷.
透過做一些壞事.
來達至他們心目中.
好的結果.
就是幫便雅敏派.
可以開枝散葉.
你會看到很奇怪.
他們通過兩個惡行.
希望換一個正面的結果.
就是幫便雅敏人計後.
經文告訴我們.
以色列人的處境是怎樣呢?.
他們現在每走一步.
都是為了解決.
之前所種下的問題.
但每走一步.
反而不能讓事情有好轉.
反而每走一步.
好像解決了之前的問題.
但同時製造更多的問題.
讓更多的人.
捲入這個我們說.
已經不是罪惡的循環.
而是罪惡的漩渦.
不單止平面捲入.
是越捲越入.
讓更多的人離開上帝.
更難搞的是.

$^{681}$當時人始終不明白.
到底在哪裡.
是誰出了問題.
弟兄姊妹.
你知不知道是誰出了問題?.
是誰出了問題?.
《士司記》的作者.
其實很想我們問.
問我們自己.
誰出了問題?.
結果事情越搞越窩窩.
我想我們也有聽過一個典故.
叫做周處除三害.
有沒有聽過?.
沒有人點頭.
可能這個典故.
現在不知道有沒有教.
不是最近上映的台灣電影.
你不要搞錯.
是一個我們以前的典故.
話說在西晉時期.
有個人叫周處.
周是誰的周?.
姓周的周.
處是四處的處.
他至少沒有人管教.
也不肯去讀書.
遊手好閒.
但他又長得高大.
力氣大.
所以恃強凌弱.
你當是技安吧.
所以當地人很討厭他.
也怕了他.
當時有三害.
在他們民間流傳.
一個是南山上的老虎.
經常會出廟傷害人畜.
另一個就是在一條橋.
長橋下面.

$^{721}$有一條蛟龍.
其實應該是鱷魚.
經常出來擾民.
嚇得人過橋很害怕.
於是大家就說.
南山的猛虎.
長橋的惡蛟.
即蛟龍的蛟.
還有惡霸周處.
合稱三害.
其中誰為何最深呢?.
周處.
有一天周處.
照樣在街上遊來浪蕩.
他看到百姓很悶不樂.
於是很好奇.
問一個老人家.
老人家.
其實今年大家收成也不錯.
挺好.
為何仍要愁眉苦臉?.
這個老人家就無奈地說.
三害不除.
大家怎會有安樂呢?.
其實周處第一次聽三害.
應該沒有人敢跟他說.
怕被他打.
於是一聽到三害.
是誰?是甚麼?.
這個老人家就告訴他.
南山的猛虎.
長橋的蛟龍.
再加上你周處.
就是三害.
周處聽到之後很震驚.
因為從來沒有人這樣跟他說.
告訴他自己.
他是最壞的.
他慨嘆自己竟然與.
猛虎和蛟龍並列三害.

$^{761}$還以為是三害之首.
於是他說了一句話.
既然大家都為了三害苦惱.
好,我就將他們消滅.
於是典故就是.
他拿刀拿箭殺了猛虎和蛟龍.
最後他如何收拾自己?.
是否一死以謝天下?.
不是.
典故記載他拜師苦讀.
修養自己的品格.
最後做了官.
成為一位中神.
徹底的改頭換面.
《鼎千尋》這個故事.
聽了多少次?.
好像聽過.
但他說了一件事是真的.
很多時最大的敵人.
真正的敵人是誰?.
是你自己.
士師記也是這樣跟我們說.
他們面對的最大敵人.
其實已經不再是周邊的家男人.
是他們自己.
而他們也不肯面對這個問題.
永遠將問題推給其他人.
甚至推給上帝.
如何解決?.
很難走出這個困局.
當越來越多人這樣想.
好像沒法挽回.
很難處理.
除非他們真的肯煞停自己.
願意謙卑下來.
承認自己的罪行.
願意認真處理.
首先和上帝的關係.
然後尋求上帝的寬恕.
才會慢慢見到出路.

$^{801}$才可以重回正軌.
再重新修復破壞的人際關係.
只可惜.
你會看到.
士師記呈現的畫面.
似乎不是這樣.
百姓好像一頭已經開行了.
煞不了的車.
我們說沒有煞車.
他已經駛到不歸路.
似乎沒有人可以勸服他.
讓他調頭.
不要再向前走.
撞車.
沒有人可以煞停他們.
所謂的「任意而行」.
到最後你會看到.
也印在這裡.
結尾.
「士師記交代」.
「於是便雅閩人照樣而行」.
「按著他們的數目」.
「從跳舞的女子中搶去為妻」.
「就回自己的地業去」.
「又重修成一居住」.
「當時以色列人離開那裡」.
「各歸本之派」.
「本宗族」.
「本地的業去」.
即是甚麼呢?.
即是代表.
百姓已經覺得問題解決了.
我的煩惱已經沒有了.
我可以繼續回到自己的地方.
如常的生活.
這是他們的想法.
但是作者補充了.
上帝的想法是甚麼.
我們可否讀一讀.
最後那25節.

$^{841}$預備起.
「但是以色列總會有過」.
「多人的禮儀」.
這是上帝的看法.
透過「士師記」的作者告訴我們.
換句話來說.
在上帝心目中.
問題根本一直沒有解決.
因為真正的問題沒有處理.
和上帝的關係.
由頭到尾都沒有擺平.
所以你看到.
即使這班百姓.
多麼絞盡腦汁.
用自己的方法.
來解決眼前的困難.
但是越搞問題越大.
上帝也選擇不作聲.
所謂的任由他們心硬.
但是你留意.
上帝不是袖手旁觀.
祂默默地部署.
祂的救贖計劃.
這是我們要知道.
也必須要謹記.
正如我之前說過.
讀「士師記」.
如果你只讀到這裡.
是會很灰心.
因為你見不到出路.
見不到曙光.
你會更加慨嘆.
為何信上帝的人.
都搞成這樣.
連未信的外邦人都不如.
確實這是「士師記」.
要讓我們知道的訊息.
重點是提醒我們.
前車可鑑.
不要撞下去.

$^{881}$有前人的例子看到.
不要重辦.
這是第一個.
「士師記」要告訴我們的訊息.
但這不是它唯一的訊息.
作者更加想我們明白的是.
即使是一片敗壞.
一片絕望.
當所有人都覺得.
不看好,不能拔的時候.
上帝說.
我雖然不作聲.
但不代表我不作供.
不代表我放棄了我的子民.
所以如果按正路來說.
純粹按邏輯.
以色列人自己選擇.
至今墮落到這個地步.
應該都玩完了.
放在現實裡很多關係的話.
應該玩完了.
如果是一個打工仔.
兩個決裂到這樣.
辭職或被解僱.
如果是父子.
登報脫離關係.
如果是夫婦.
可能分居離婚.
沒有盼望.
沒得搞.
但上帝說有得搞.
上帝說有得救.
因為他不放棄.
他堅持守施慈愛.
即使各人都任意而行.
不再尊重上帝.
但仍然不能阻撓上帝.
他仍然掌管萬物.
上帝定下的旨意.
沒有人可以推倒.

$^{921}$一定會成就.
所以聖經不是寫到.
在士師記就停了.
之後還有很多卷的經卷.
他們告訴我們.
上帝知道人的軟弱.
軟弱到不懂主動找上帝.
所以上帝調轉主動找你和我.
上帝知道我們欠他的實在太多太大.
做多少都還不了.
所以最後他透過他的兒子.
幫我們一次過還清.
上帝也知道我們是心硬.
很難搞.
但他不放棄.
他要處理的不是單單是我們的罪行.
他要處理的是我們的心.
他要將我們的心拆遠.
他要將我們變得柔和.
所以他一次又一次.
明明你都走了很遠.
他都盡量努力將你拉回來.
他不放手.
所以弟兄姊妹.
我們需要珍惜.
上帝這份不離不棄的愛.
我們不能像在史詩時代的百姓.
一而再再而三糟蹋辜負.
他的吩咐和期望.
所以求上帝幫我們.
今天我們所分享的.
多多少少有我們的影子.
我們會否不知不覺.
想透過做一些不好的事.
但希望達到正面的效果.
我們會否不知不覺.
想做好事.
但反而令更多人受傷害虧損.
我們能否煞停他?.
我們能否回轉.

$^{961}$回到上帝的身邊.
不要繼續這樣下去.
這是整個時事記給我們一個方向.
亦是上帝給我們的出路.
上帝說他不放棄我們.
他很想我們再一次認定.
他是我們唯一的王.
是我們唯一的主宰.
我們唯一要聽從的是他的話語.
我們唯一要信服的是他的大力.
讓我們一起祈禱.
這套上帝是我們藉著.
時事記這個系列的宣講.
最後確實讓我們看到.
人性那種黑暗.
甚至連上帝的子民.
他們的行為和心態都變得一團糟.
主啊!盼望這一切都成為我們今天的借鑒.
讓我們知道原來當人遠離上帝的時候.
不再尊重上帝你為王.
不再聽你的吩咐的時候.
真的可以一直一直崩壞下去.
縱然我們有多少良好意願.
但原來如果離開了上帝你的時候.
可能反而會讓更多的人受到傷害.
更多的人受到虧損.
包括上帝你都受到虧損.
主啊!大人.
我們知道這不是你所期望見到.
這亦不是我們終極要去走的.
要去達到的一個情況.
因為我們知道上帝啊!你不會袖手不顧.
你總是主動向我們手約,施慈愛.
一次又一次,多次多方.
從不同的方式,不同的人來挽回我們.
你很希望我們回轉.
重新來到你身邊.
認你為主,認你為王.
主啊!我們知道.
靠自己我們真的很多時候會走迷.

$^{1001}$我們不懂得煞停自己.
回頭看.
但我們知道聖靈啊!你可以幫我們.
求主啊!你藉著整個事司祭的提醒.
讓我們真的有一個柔軟的心.
不再心硬.
讓我們敏銳.
聖靈你的提醒.
讓我們願意放下.
我們自己一些不好的念頭.
回到上帝你的身邊.
主啊!我求你幫助我們每一個.
因為我們真的很需要你.
主啊!願你親自的.
藉著今天我們所看到的經文.
願你幫我們可以醒察.
讓我們有一個謙卑柔和的心.
來承認自己的軟弱,自己的幽暗.
讓我們有一個降服你的心.
願意回轉跟隨你.
讓我們有一個清潔的心.
讓我們再一次重視你的話語.
你的吩咐.
以此為我們生活言行的一個準則.
一個指引.
讓我們離開惡道.
重新歸回你的正路.
為此我們同心獻上祈禱.
奉耶穌基督你得勝的名字.
來到祈求.
Amen.
(影片完結).
\newpage



\section{撒母耳記下 6:16-23}
\label{sec:t62ZJvrM_ng}
\textbf{2023年11月5日崇拜直播｜陳明泉牧師｜唯獨你是王 - 牧者君王｜撒母耳記下6\_16-23}
\newline
\newline
連結: \href{https://youtube.com/watch?v=t62ZJvrM-ng}{\texttt{https://youtube.com/watch?v=t62ZJvrM-ng}} ~~~~ 語音日期: 2023-11-06
\newline
\newline
\hyperref[sec:jwnc_J9h4aM]{\small{< < < PREV SERMON < < <}}
~
\hyperref[sec:index_chronic]{\small{[返順時目]}}
~
\hyperref[sec:KCjIplGBEcc]{\small{> > > NEXT SERMON > > >}}
\newline
\newline
撒母耳記下 6:16-23
\newline
\begin{longtable}{cl}
\hline
\hline
章節 & 經文 (和合本修訂版)\\
\hline
6:16 & \begin{tabularx}{0.7\textwidth}{X} 耶和華的約櫃進大衛城的時候，掃羅的女兒米甲從窗戶裡往外觀看，見大衛王在耶和華面前踴躍跳舞，心裡就輕視他。 \end{tabularx} \\ \\ \relax
6:17 & \begin{tabularx}{0.7\textwidth}{X} 眾人將耶和華的約櫃請進去，安放在所預備的地方，就是大衛為它搭的帳幕中。大衛在耶和華面前獻燔祭和平安祭。 \end{tabularx} \\ \\ \relax
6:18 & \begin{tabularx}{0.7\textwidth}{X} 大衛獻完了燔祭和平安祭，就奉萬軍之耶和華的名祝福百姓， \end{tabularx} \\ \\ \relax
6:19 & \begin{tabularx}{0.7\textwidth}{X} 並且分給以色列眾人，所有的百姓，無論男女，每人一個餅，一個棗子餅，一個葡萄餅。眾人就各自回家去了。 \end{tabularx} \\ \\ \relax
6:20 & \begin{tabularx}{0.7\textwidth}{X} 大衛回去要為家裡的人祝福，掃羅的女兒米甲出來迎接他，說：「以色列王今日有好大的榮耀啊！他今日在臣僕的使女眼前露體，如同一個無賴赤身露體一樣。」 \end{tabularx} \\ \\ \relax
6:21 & \begin{tabularx}{0.7\textwidth}{X} 大衛對米甲說：「這是在耶和華面前的。耶和華已揀選我，在你父和你父的全家之上，立我作耶和華百姓以色列的君王，所以我在耶和華面前跳舞， \end{tabularx} \\ \\ \relax
6:22 & \begin{tabularx}{0.7\textwidth}{X} 我也必更加卑微，自己看為低賤。至於你所說的那些使女，她們反而尊重我。」 \end{tabularx} \\ \\ \relax
6:23 & \begin{tabularx}{0.7\textwidth}{X} 掃羅的女兒米甲，直到死的那日沒有孩子。 \end{tabularx} \\ \\
[1ex]
\hline
\hline
\end{longtable}
$^{1}$今日讀的經文是記載在《撒慕爾記》下第六章16-23節.
眾人將耶和華的藥櫃請進去,安放在所預備的地方,就是在大衛所搭的帳幕裡..
大衛在耶和華面前獻繁祭和平安祭,就奉萬君子耶和華的名給民祝福,並且分給以色列眾人,無論男女,每人一個餅,一塊肉,一個葡萄餅,眾人就各自回家去了..
大衛回家要看眷屬祝福,素羅的女兒米格出來迎接他,說:.
「以色列王今日在臣服的婢女眼前露體,如同一個輕賤人無恥露體一樣,有好大的榮耀喔!」.
大衛對米格說:「這是在耶和華面前,耶和華已揀選我,廢了你父和你父的全家,立我作耶和華民以色列的君,所以我必在耶和華面前跳舞,我也必更加卑微,自己看為輕賤.」.
大衛說:「你所說的那些婢女,她們禱告尊敬我.」.
大衛說:「素羅的女兒米格直到死日,沒有生養兒女.」.
聖經獨笠.
弟兄姊妹平安,去到十一月了,某程度上是二零二三年的倒數,也是我們用這幾個禮拜一步一步的思想,去到聖誕節就是我們的一個高峰..
相信之前的弟兄姊妹,我們一起來到,在過去幾個月一起讀士師記,你會看到整個民族每況愈下,見到他們一路一路進入一個令人不堪入目的處境..
不過去到某程度上一個最差的狀態,上帝沒有放棄這個民族,跟著下來就好像一個新的篇章一樣,上帝更新了整個民族,特別是上帝為這個民族預備君王..
在這裡值得注意的,不是純粹建立一個君王的制度,純粹有一個制度這麼簡單,更加重要的就是上帝要選照,要納若,帶領整個以色列民族的一個合適的君主,帶領整個民族走近上帝,尋求上帝的心意..
在《撒母耳記下》第五章裡面,上帝特別應許大衛,他這麼說,你必務讓我的民以色列作以色列的君..
如果你看回第五章裡面,他在說君主,統治者,成為以色列整個民族的牧者,我們想牧者很多時候就以為牧師傳道,不過在當代裡面,牧者這個字,其實不單止是在說驅寒問暖,令你很舒服,更加重要的就是帶領整個以色列走近上帝..
讓上帝的心意實踐在整個民族當中,如果一個人生命裡面,我們得著救恩,但是得不到上帝的指示,大概我們的人生都是漂泊流離,我們都不知道自己怎麼做,.
不過當我們生命裡面找到一個主宰,而且這個是一個致福的源頭,帶領我們一起走人生的路的時候,我們的生命是蒙福的..
當我們剛才一路唱詩歌的時候,說著上帝助爵為王,不單止是說某些政治或者某一種制度上助爵為王,上帝作王也在我們的內心當中..
接下來這個港獨系列,我們唯獨你是王,我們希望我們一路看著不同的君王,直到耶穌基督,究竟我們生命裡面怎樣以基督成為我們生命的主,成為我們生命的王..
接下來這幾次,幾密集,特別是說的是舊約經文,可能我們涉及的經文會較多,但是希望我們真的有一個歷史脈絡,來掌握整個上帝救恩的計劃,怎樣的計劃,以致我們更加懂得怎樣行事為人..
我們先祈禱,求上帝幫助我們,同心祈禱..
先父求你在我們當中繼續來對稅領,帶領我們眾人,讓我們能夠明白你的話語,讓我們生命當中不單止有宗教的禮儀,更加我們心靈裡面更加經歷上帝的恩典,和上帝要我們所行的路,讓我們快跑跟隨你..
我們恭敬將來領受你的話語時間交托,幫助孫講的講得清楚,幫助每位的靈聽鼎姐妹,他們聽得明白,獻上你禱告奉基督耶穌得勝的聖名,Amen..
剛才我們一起看過三毫爾記的記載,三毫爾記其實是記載著,特別是說大衛這個皇帝高立,和整個民族的恩爵這個王的建立,有些很不同的改變..
不過你值得留意的一樣東西,剛才我也有說過,這個王的角色和昔日的士司是有點不同的,士司主要的角色就是打仗,或者是做一些審判,評定是與非,這個是他主力要做的工作,王也是要做這些的..
不過去到大委身上,或者他成為了整個以色列的君王,不單只是說他是一個軍事專家,一個政治領袖,又或者是一個制度上建立的承接著權力的一個人,在聖經裡面說的這個君王,他也某程度上扮演了一些祭司的角色..
在聖經的記載裡面,他刻意要令到大衛的身份,不單只是說他是一個戰士,不是一個打仗的軍事人才,更加是說他對上帝的心意,和對整個以色列民帶起一個很重要的復興,將以色列民帶到上帝的角色..
所以我們剛才讀的,在撒母耳記下第六章裡面,記載著一個很特別的情況,就是大衛自己親身來迎接約衛.我們先看從第六節到第十五節,看回我們剛才所讀的經文的上下文..
在經文裡面說道,到了那銀的和場,因為牛失前提,烏薩就伸手扶著神的約衛,耶和華的怒氣向烏薩發作,神恩者冒犯在那裡擊打他,他就死在那裡,在神的約衛旁..
大衛因耶和華突然衝出撞死,烏薩就生氣,稱那地方為比列斯烏薩,直到今天..
那一天,大衛懼怕,耶和華說,耶和華的約衛怎可到我這裡來?於是大衛不願將耶和華的約衛接到大衛城他自己的地方,卻轉送到加特人俄別以東的家中..
耶和華的約衛停在加特人俄別以東家中三個月,耶和華似乎給俄別以東和他的全家..
有人高騷大衛說,耶和華因約衛的緣故似乎給俄別以東的家和一切屬他的,大衛就去歡欣喜地將神的約衛從俄別以東家接上來,到大衛城裡,抬耶和華約衛的人走了六步,大衛就獻牛和肥畜為祭..
大衛穿著細麻衣以佛德,在耶和華面前極力跳舞,這樣大衛和以色列全家歡呼催國,將耶和華的約衛接了上來..
我們看到這個六至十五節,就說到剛才讀的經文,大衛為什麼無端端跳舞呢?大衛為什麼會做一些獻祭的舉動呢?.
原來上文說到,因為之前打仗,整個以色列將一個最重要代表神童在的約衛,除了人沒有它之外,擄掠了這個約衛,去了非利士的地方..
大衛覺得這個約衛是非利士人的戰利品,而且拿到這個約衛,就像拿到偶像一樣,他們是這樣想的,.
不過他們沒有想過,這個約衛只不過是上帝同在的象徵物,上帝助爵為王,上帝統管萬有,.
當這個約衛拿到整個非利士民族的時候,非利士人就出現了一個很大的問題,就是整個民族生痔瘡,都很恐怖,.
個個都痛苦,叫苦連天,個個都生痔瘡,坐不定,不過非利士人就意識到這個約衛在以色列人當中非同凡響,.

$^{41}$於是在他打仗又打輸,上帝令到以色列民族打贏了他,然後非利士人就順勢把這個約衛交還給他們,.
他想著這個是不祥之物對他們來說,所以都歸還就OK了,當大衛拿回這個約衛的時候,第一次去拿這個約衛,.
他就犯下了一個嚴重的錯誤,為什麼呢?因為在經文裡面說,之前非利士人拿那個約衛,他用那些牛車來送,拿走這個約衛去他們自己的民族當中,.
到了大衛要拿回這個約衛的時候,第一次他想著拿這個約衛,他就用回非利士人混在偶像的方法,用一輛牛車來送這個約衛回到大衛城裡面,.
誰知道這個牛車走著走著的時候,去到蘭銀的禾場裡面,不是馬室前提,是牛室前提,然後牛就撲一撲,然後整個約衛就開始反了,.
在他快要反下來的時候,有一個人叫烏薩,應該不是祭司,而是大衛的將領,就立刻用手托著他,.
見到他其實是出於本能,也是出於責任感,扶著這個約衛,這個舉動如果對於一個押回押送者來說,這是很人之常情的事情,.
不過在上帝的律法當中,在民數記裡面,他特別說到這個約衛不是每個人都可以碰的,不是用手直接碰的,.
當這個烏薩是一個將士無端端碰了這個約衛的時候,聖經說道,上帝因為烏薩這個舉動覺得是被冒犯,於是怎樣做呢?.
就在那裡擊打,擊殺了這個烏薩這個人,於是這個人就死了在約衛的旁邊,在整件事裡面你會見到,整件事的錯誤在於什麼呢?.
大衛用了非利士人般偶像的方法來處理這個約衛,他想著上帝使他們得勝,他想著這件事是一個很簡單的運送的過程,好像叫高高VAN一樣,.
不過約衛不是可以用這樣的方法,於是他們就犯下了這個錯誤,當這個錯誤出現的時候,大衛就害怕,在第九節那裡大衛很懼怕耶和華,耶和華說約衛怎麼可以去我那裡,如果有一個擊殺的事情,於是害怕,.
第一件事就是放在別人家裡,事情就是這樣,於是就放在一個人,就是加特人,俄別以東的家裡暫放著,他害怕,不知道上帝的工作是怎樣,不過這個約衛去到俄別以東的人家裡,上帝卻賜福給這個家庭,.
具體的賜福我們不知道,因為經文的重點是說這個約衛不是不祥之物,而是那些人對上帝或者那種不謹慎,是上帝要給一個訊號給他們知道,上帝是興萬不得,上帝是有要求的,.
所以在這個過程裡面,見到上帝賜福給俄別以東的這個家庭,然後有人跟大衛說,上帝賜福給這個家庭,於是大衛就知道約衛其實是他的同在,代表著上帝的賜福,於是大衛就去接約衛重新去做運送的功夫,.
得了第一次教訓之後,第二次大衛做的事情就很不同了,大衛這次運送的時候,十三節,他們抬藥櫃,這次就不是用牛車,是用利未人,他派人親自用人來托整個藥櫃,.
藥櫃其實有兩支長竿,其實就代表著一些結證了的人,來抬著那個台竿,逐步逐步經歷上帝的同在,經文裡面講到這次大衛不單止學精,更加變得很謹慎,謹慎到一個地步,嚇你怕的,.
第十三節講到那些抬藥櫃的人每走六步,大衛就獻牛和肥畜為祭,都挺厲害的,走六步就獻一個祭,走六步就再獻一個祭,由俄別以東的家直至到一路步行,這樣獻祭,去到大衛城,究竟用多長時間,真的不知道,.
不過絕對不是一個短時間,用一個極大的誠意,來將這個藥櫃抬供到大衛城當中,除了找這些人來抬著這個藥櫃之外,經文裡面講大衛在接送的藥櫃有角色,.
經文裡面講大衛穿著細麻衣的伊弗德在耶和華面前極力跳舞,你會覺得很奇怪,首先講一下細麻布的伊弗德是什麼衣服呢?.
祭司穿的衣服,在伊弗德那裡有污淩土茗,就是問上帝的心意,那個是祭司的制服,大衛是君王,已經被高納了,成為了一國之君,他為什麼要穿那件祭司的制服呢?.
簡單來說,用一個比喻,一間公司的CEO,他不會穿那件清潔工的制服來工作的,不過大衛就用一個這樣的方式,他有份一起帶領以色列人去接藥櫃回到大衛城,.
第二件事,你會看到在耶和華面前極力跳舞,那個跳舞跟我們今天講的跳舞有些不同,同樣都是身體的律動,不過當時的跳舞是一個宗教制治儀式,而做這些宗教制治儀式,怎麼會是君王做的呢?.
甚至乎有些說法是連男人都不會做的,當時是找一些女士,找一些女孩子來跳舞的,這個宗教禮儀不是由一個高高在上的君王來做的,整件事你會覺得一反常態,一反大家覺得那個君王是一個很尊貴,走出來指點江山的做法,.
在這裡大衛穿上祭司的袍,在這裡實踐跳舞,實踐宗教禮儀,來催覺歡呼,把耶和華的藥櫃接到大衛城當中,所以在這裡你會看到整個上文,或者整件事是在講之前有前科可鑑,.
大衛就避免了再有出現上帝降和的日子,於是乎他用到他自己,即使是皇的身份也好,他都放下最大的身段,而且帶領整個以色列民一起來迎接這個藥櫃去到大衛城當中,經文繼續的來到他講,.
在第六章十四至十九節,大衛剛才讀過了,在耶和華面前跳舞,大衛和以色列全家歡呼催覺,把耶和華的藥櫃接了上來,第十六節,耶和華的藥櫃進到大衛城的時候,.
大衛獻完梵濟和平安祭,奉萬君之耶和華的名祝福百姓,並且分給以色列眾人,所有的百姓,無論男女,每人一個餅,一個肉,一個葡萄餅,各人就各自回家去了..
第十四至十九節,不單止是大衛迎接這個藥櫃,又跳舞,又穿上祭司的衣服,呈現一個祭司的形象,帶領以色列民人來迎接這個藥櫃,當這個藥櫃去到大衛城的時候,大衛做了什麼呢?.
大衛繼續獻祭,在耶和華面前獻梵濟和平安祭,然後奉耶和華的名祝福百姓,分餅,分肉,葡萄餅給眾人,在這裡你會看到大衛又再呈現一個什麼角色呢?.
祭司的角色,然後你讀的時候,如果你認真,或者你明白以色列歷史,你會知道先知君王,祭司是以色列整個民族裡面三個很重要的角色,是不是?.
但是我們覺得,我們很多時候想,這三個角色是分開不同的人的,是互相制衡的,甚至乎有一些角色是有一條界在這裡的,你不能做這些事情的,很多時候我們人是這麼想的,不過三寶義記他要打破我們一些想法,大衛作為一個君王,他也在實踐祭司的角色,他不單止在迎接的路裡面扮演祭司的角色,.
更加是在獻祭和祝福的過程裡面,將這個約季去到的時候,他也扮演祭司的角色,如果你記得的時候,大衛是整個以色列民族裡面第二個王,第一個王是誰呢?.
是蘇羅,蘇羅的失敗其中一個很重要的原因,就是有一次打仗,在三寶義記上第十三章,第八節至十二節,三寶義就命令蘇羅,當時的軍隊,大軍已經押警了,跟著那些外邦人要打蘇羅的軍隊了,蘇羅在等著薩姆爾獻祭,.
獻了祭之後,跟著就可以集合所有軍隊,就一起來抵抗這班外敵,這班外敵就已經去到密沒這個地方了,跟著在十三章第八節,三寶義記,蘇羅照著三寶義所定的日期等了七天了,但是三寶義都還沒有到吉格,百姓就離開蘇羅散去了,.
他正在吹雞,正在招募兵,準備打仗,等了七天都還沒有見到,百姓就回家了,軍心開始散緩了,於是蘇羅說,把繁祭和平安祭帶到我這裡來,蘇羅就獻上繁祭,他剛獻完繁祭,看著三寶義就到了,蘇羅出來迎接他,問他的安,三寶義說,你做了什麼?.
蘇羅說,因為我見百姓離開散去了,你又不照所定的日期來到,而且菲律士人已經在密沒集合了,我說,現在菲律士人已經下到吉格來攻擊我,可是我還沒有向耶和華禱告,所以我就勉強獻上繁祭,跟著之後,經文就說到,三寶義就說審判的話,上帝厭棄蘇羅作王,.
在這件事裡面,我們解釋一下,我們說,蘇羅為什麼被上帝厭棄作王?他做了一些他不應該做的事情,他為什麼無端端要獻祭呢?.
三寶義叫他等,由他來獻祭,他就見到有些人害怕,所以他就自己做了自己越界的繁祭,所以就令到蘇羅越走下坡了,就被上帝厭棄他作王了,如果你回看經文和剛才大衛所說的,兩個王都是獻祭,又是獻繁祭,蘇羅都沒有大衛那麼嚴重,.
大衛還要分餅,奉耶和華的名,祝福他的百姓,為什麼大衛可以,為什麼蘇羅不行?蘇羅做獻祭就會被耶和華譴責他厭棄他作王,但是大衛就可以繼續他的江山繼續得到上帝的祝福呢?.

$^{81}$關鍵是其實蘇羅和大衛同樣都是獻祭,蘇羅的獻祭只是出於他看這件事是一個宗教禮儀,他看這只是一個外在的儀式,他心裡面一心想著的只是打仗,他是以實用的角度來看當下這個獻祭的舉動,在他心目中他都等了七天,.
不過他看到有些事情已經比眼前的狀況更加危急,於是他就不再等了,自己來吧,自己勉強自己來做這個舉動,蘇羅在他心目中這個獻祭的舉動並不是一個發自內心,一種對上帝謹慎的舉動,這個獻祭流露了,反而輕忽了這件事,.
他只是覺得滿足了這個宗教禮儀,他就可以得到祝福,他就可以名正言順的招募軍隊來打仗,上帝是一個只是陶酉宗教外表,但是對他心裡面不是尋求的一個君王,相反來說,大衛他看到誣殺的事件,你猜他不謹慎嗎?.
看到有一個將領死了,碰一碰那個藥櫃就立即死了,對於這件事他不敢輕忽的,你想一下,如果他扮演祭司的角色,他不是很清楚上帝的旨意,隨時擊殺的是他自己,所以在這篇經文裡面,講到的那個獻祭,或者他帶頭做的這個祭司的角色,大衛不是輕忽的,.
或者反過來說,當大衛作為一國之君,他都帶頭成為一個謙卑,一個祭司的形象,來迎接上帝所重視的,代表上帝同在的藥櫃,其他人不敢不跟從的,.
我一個這樣的經歷,講到這裡,有時候在上或者在家裡面,每個人都有自己的地界,如果你做大的,你做長輩的,你去願意用最謙卑的方式去實踐信仰,其他人沒辦法不跟,甚至不敢不去跟,.
我記得有一次,我那時候是初出茅廬,去到機構裡面服侍,跟著查經,總幹事跟我們一起有些分享查經,進去的時候,椅子還沒放,進到一個課室的位置,跟我們有些圍圈的來分享,一般來說做這些圍圈的分享,大家就擔椅子,.
但大家手牽手不動,我其實也有做事,但我遲進去,我有一點遲到,已經不太尊重,跟著看到有些同工手牽手不動,圍圈,椅子就搭得很高,到最後是總幹事,其實是蔡元雲醫生,他老人家就擔椅子,逐張擔,坐吧,年輕人坐吧,.
那一刻我有什麼感覺呢?有一種折腹的感覺,你明白的,有時回不上,你也知道,你所尊崇的長輩忽然間幫你服侍你,跟著那種感覺,你死了,你為什麼會這樣,跟著立刻就立刻擔椅子,我想說的是什麼呢?.
一個有份量的,一個有身份的長輩,去做這些服侍的舉動,如果真心去跟的,那些人看到這件情景,他不會不為所動的,他一定會跟著去做的,在你家裡做長輩也好,在教會裡面做領袖也好,在你公司裡面做上司也好,如果你願意帶頭去做一個這樣的角色的時候,其他人一定跟的..
當你做了這樣的舉動的時候,你帶頭去讓所有人來尊重那個場合,尊重上帝,所有人你就可以帶來真正的秩序,相反來說,如果你有這麼多,但你不做,你抽離,你點人做呢?.
到所有人,每個人就互相點,斧頭作作,最小的那個就永遠做所有事情,在這裡面,大衛他扮演祭司的角色,他不是要自高,他不是要逾越那條界線,他不是說他很厲害,所以連祭司那一塊也混在一起,而是他帶著最謙卑,最謹慎的心,.
來接這個約會,在他們國民當中,他用最卑微的方法,甚至乎被人笑的,他跳舞,那些是最低下的人來做的宗教禮儀的先鋒部隊,他扮演那個角色,他在這裡就實踐這個最謙卑作王的身份..
我們繼續看經文,第20至23節,在這個段落裡面,最後的一部分,大衛要為家裡的人祝福,蘇羅的女兒米甲出來迎接他,說:以色列的黃金物,很大的榮耀呀,他今日在神僕的使女眼前露體,如同一個無賴赤身露體一樣..
大衛對米甲說:這是在耶和華面前的,耶和華已經揀選我在你的父和你的父家之上,納我作耶和華百姓以色列的君王,所以我在耶和華面前跳舞,我也是必更加卑微,自己捍衛低賤.至於你所說的那些使女,他們反而尊重我.蘇羅的女兒米甲直到死的那一日,沒有孩子..
經文裡面,最後的一段落是甚麼呢?當大衛帶領整個民族,所有人都去尊崇上帝,唯獨是他的枕邊人,這是蘇羅的女兒米甲,青梅竹馬的妻子,大衛成為了一個王..
米甲如何看待大衛這個舉動呢?他覺得他成何體統,用反華,很大的榮耀,其實是在用他做一些這樣的舉動,一個很卑微的舉動,作為一個王,你在跳舞,他說的那個字是很難聽的,.
在使女面前怒看,好像一個無奶赤身怒看一樣,那些話尖酸刻薄的,一方面他覺得他不合體統,第二樣就是藐視大衛他把自己的身段放到最低,用一個最卑微的方式在這班以色列人當中迎接上帝的約會..
在這裡有些字眼值得留意,在這段經文裡面,大衛說自己是耶和華納他成為以色列的君王,在聖經裡面和合本沒有刻意翻譯,不過其實在原文裡面有兩個字,.
一個字是納他做君王的君王的字,那個字叫管治者,就是管家的意思,是一個牧者的意思,幫上帝管理,他沒有終極的管治權,另外在聖經裡面講的君王的字,另外一個譯字叫做帝王,某程度上是一個大的霸王,.
管治者是指他沒有終極的管治權,他只不過是幫手做一個管家,是一個牧者,Malak,是指所有東西都是屬於我的,我是那個霸王,在聖經裡面用在大衛或大衛自稱,他是屬於上帝的Nagin,是一個管治者,.
但相反來說,掃羅看自己是Malak,當上帝要高納他,在聖經裡面掃羅將王的身份交付給他的時候,他看他所有東西都是屬於他自己的,他就是那個老闆,終極交代的那個是他,.
所以他將所有的責任放在自己身上,又或者將所有的權力放在他自己的身上,大衛在米格面前講這番話,說上帝納他作為以色列的君王,這個是他心底的確信,大衛和掃羅最不同的地方是,無論他的權力或人生是怎樣都好,.
這個國家到最後的主權是屬於耶和華上帝,不是屬於他自己的,但在掃羅眼中,他覺得他就是一國之君,所有東西都是屬於他,我們再用多一點經文你就知道掃羅在想什麼,.
撒母耳基上15章30節,當掃羅離棄上帝,上帝閹棄他作王,被撒母耳譴責的時候,15章30節怎麼說呢?.
掃羅說我有罪了,現在求你在我百姓的長老和以色列人面前尊重我,和我回去,我很敬拜耶和華你的神,你留意一下,當他做錯了事的時候,他第一件事求撒母耳的是什麼呢?.
給我面子,讓我在一群子民面前也有面子,以至我很敬拜耶和華,你的神不是他的神,掃羅在看他是一個帝王,所有東西都是屬於他,當他犯了罪的時候,.
撒母耳給我面子,讓我繼續維繫著這樣的身份,這樣的尊貴,相反來說,大衛在由第一刻開始,在做王的裡面,他知道這個國家不是屬於他,.
他是屬於上帝,上帝才是真正的王,他成為一個祭司,只是讓位,把真正的主權交回給上帝,他只不過是一個中間人,只不過是一個上帝使用的權力多一些的管家,但是最終的管事權不是他,.
在大委身上,雖然他的道德未必有掃羅這麼好的,不過他的心,他的價值觀,在他體會自己的身份,他看得很重,所以在聖經裡面說,上帝是上帝合乎上帝心意的一個王,.
今日這個觀念都值得放在我們的生命當中,今日我們人生裡面都有不同的位份,有不同的地界,我們會不會很強調,這件事是我的,我是老闆,我是屬於所有話事的,今日這個世代是喜歡多一些掌控感的,.
不過當我們看到大衛這個故事的時候,求主繼續對我們說話,今日我們所擁有一切其實不是屬於我的,上帝只是交付給我們,讓我們成為一個管家,.
我記得有一個很成功的商人曾經說過,你這麼多錢,你怎麼用呢?他說,這些錢不是屬於我的,這些錢只是上帝,不過上帝比我信得多,讓我去管多一些,但到最後我不是要自肥的,.
同樣的,弟兄姊妹,今日你所擁有一切,不要太過介意那些東西是不是屬於你,我們只不過是一個管家,我們只是被上帝信我們,託付我們去好好管理,在管理的過程裡面我們都有意識,不要覺得那些東西是屬於自己,.
在大衛的經歷裡面,這個王一直將以色列帶到新的一頁,另外一個重點就是我們可不可以謙卑下來,放下身段,不是越界的行為,只不過是在那一刻讓上帝作主,.
讓自己的位置放下,帶著一份謙卑,恭謹的心,謹慎地來服侍,謹慎地渡過你的生活,這也是大衛記載著這個令人摸不著頭腦,但也是很重要的經文所在,.
求主幫助我們,我們深信唯獨耶和華,唯獨基督是我們生命的王,我們只是一個管家,我們講不上有什麼很了不起,我們只是恭謹走面前人生的路..
\newpage



\section{以賽亞書 7:1-17}
\label{sec:KCjIplGBEcc}
\textbf{2023年11月12日崇拜直播｜陳冠成牧師｜唯獨你是王 - 以馬內利的兆頭｜以賽亞書7\_1-17}
\newline
\newline
連結: \href{https://youtube.com/watch?v=KCjIplGBEcc}{\texttt{https://youtube.com/watch?v=KCjIplGBEcc}} ~~~~ 語音日期: 2023-11-14
\newline
\newline
\hyperref[sec:t62ZJvrM_ng]{\small{< < < PREV SERMON < < <}}
~
\hyperref[sec:index_chronic]{\small{[返順時目]}}
~
\hyperref[sec:SCWtMBw63Gk]{\small{> > > NEXT SERMON > > >}}
\newline
\newline
以賽亞書 7:1-17
\newline
\begin{longtable}{cl}
\hline
\hline
章節 & 經文 (和合本修訂版)\\
\hline
7:1 & \begin{tabularx}{0.7\textwidth}{X} 烏西雅的孫子，約坦的兒子，猶大王亞哈斯在位的時候，亞蘭王利汛和利瑪利的兒子以色列王比加上來攻打耶路撒冷，卻不能攻取。 \end{tabularx} \\ \\ \relax
7:2 & \begin{tabularx}{0.7\textwidth}{X} 有人告訴大衛家說：「亞蘭與以法蓮已經結盟。」王的心和百姓的心就都顫動，好像林中的樹被風吹動一樣。 \end{tabularx} \\ \\ \relax
7:3 & \begin{tabularx}{0.7\textwidth}{X} 耶和華對以賽亞說：「你和你的兒子施亞‧雅述要出去，到上池的水溝盡頭，往漂布地的大路上，迎見亞哈斯， \end{tabularx} \\ \\ \relax
7:4 & \begin{tabularx}{0.7\textwidth}{X} 對他說：『你要謹慎，要鎮定，不要害怕，不要因利汛和亞蘭，以及利瑪利的兒子這兩個冒煙火把的頭所發的烈怒而心裡膽怯。 \end{tabularx} \\ \\ \relax
7:5 & \begin{tabularx}{0.7\textwidth}{X} 因為亞蘭、以法蓮，和利瑪利的兒子設惡謀要害你，說： \end{tabularx} \\ \\ \relax
7:6 & \begin{tabularx}{0.7\textwidth}{X} 我們要上去攻擊猶大，擾亂它，攻破它來歸我們，在其中立他比勒的兒子為王。 \end{tabularx} \\ \\ \relax
7:7 & \begin{tabularx}{0.7\textwidth}{X} 主耶和華如此說：這事必站立不住，也不得成就。 \end{tabularx} \\ \\ \relax
7:8 & \begin{tabularx}{0.7\textwidth}{X} 因為亞蘭的首都是大馬士革，大馬士革的領袖是利汛；六十五年之內，以法蓮必然國破族亡， \end{tabularx} \\ \\ \relax
7:9 & \begin{tabularx}{0.7\textwidth}{X} 以法蓮的首都是撒瑪利亞；撒瑪利亞的領袖是利瑪利的兒子。你們若是不信，必站立不穩。』」 \end{tabularx} \\ \\ \relax
7:10 & \begin{tabularx}{0.7\textwidth}{X} 耶和華又吩咐亞哈斯： \end{tabularx} \\ \\ \relax
7:11 & \begin{tabularx}{0.7\textwidth}{X} 「你向耶和華－你的神求一個預兆：在陰間的深淵，或往上的高處。」 \end{tabularx} \\ \\ \relax
7:12 & \begin{tabularx}{0.7\textwidth}{X} 但亞哈斯說：「我不求；我不試探耶和華。」 \end{tabularx} \\ \\ \relax
7:13 & \begin{tabularx}{0.7\textwidth}{X} 以賽亞說：「聽啊，大衛家！你們使人厭煩豈算小事，還要使我的神厭煩嗎？ \end{tabularx} \\ \\ \relax
7:14 & \begin{tabularx}{0.7\textwidth}{X} 因此，主自己要給你們一個預兆，看哪，必有童女懷孕生子，給他起名叫以馬內利。 \end{tabularx} \\ \\ \relax
7:15 & \begin{tabularx}{0.7\textwidth}{X} 到他曉得棄惡擇善的時候，他必吃乳酪與蜂蜜。 \end{tabularx} \\ \\ \relax
7:16 & \begin{tabularx}{0.7\textwidth}{X} 因為在這孩子還不曉得棄惡擇善之先，你所憎惡的那兩個王的土地必被撇棄。 \end{tabularx} \\ \\ \relax
7:17 & \begin{tabularx}{0.7\textwidth}{X} 耶和華必使亞述王臨到你和你的百姓，並你的父家，自從以法蓮脫離猶大的時候，未曾有過這樣的日子。 \end{tabularx} \\ \\
[1ex]
\hline
\hline
\end{longtable}
$^{1}$接下來是經訓.
今天的經文是《以賽亞書》第七章.
《舊約聖經》的九八一十三頁.
今天會有兩段.
同樣是第七章 一至七節.
和第十至十七節.
請聽我讀.
「烏西亞的孫子 約旦的兒子.
猶大王阿哈斯在位的時候.
亞蘭王利遜和利瑪利的兒子.
以色列王比加上來攻打耶路撒冷.
卻不能攻取.
有人告訴大衛家說.
亞蘭與以法連已經同盟.
王的心和百姓的心就都跳動.
好像林中的樹被風吹動一樣」.
第三節.
「耶和華對以賽亞說.
你和你的兒子斯亞亞述出去.
到上池的水溝頭.
在漂泊地的大路上.
去迎接阿哈斯.
對他說.
你要謹慎安靜.
不要因亞蘭王利遜和利瑪利的兒子.
這兩個冒煙的火把頭所發的烈怒害怕.
也不要心裡膽怯.
因為亞蘭和以法連.
並利瑪利的兒子.
設惡謀害你.
說我們可以上去攻擊猶大.
擾亂他.
攻破他.
在其中立他比立的兒子為王.
所以主耶和華如此說.
在這所謀的必立不住.
也不得成就」.
第十節.
「耶和華又曉遇阿哈斯說.
你向耶和華你的臣求一個兆頭.

$^{41}$或求顯在深處.
或求顯在高處.
阿哈斯說.
我不求.
我不試探耶和華.
以賽亞說.
大衛加亞.
你們當聽.
你們使人厭煩喜算小事.
還要使我的臣厭煩嗎?.
因此.
主自己要給你們一個兆頭.
必有同類懷孕生子.
給他起名叫以瑪奈利.
就是臣與我們同在的意思.
到他曉得.
棄惡擇善的時候.
他必吃奶油與蜂蜜.
因為在這孩子還不曉得.
棄惡擇善之先.
你所憎惡的那二王之地.
必自見棄.
耶和華必使亞述王.
攻擊你的日子.
臨到你和你的百姓.
並你的父家.
自從二百年離開猶大以來.
未曾有這樣的日子」.
聖經獨白.
(聖經獨白).
弟兄姊妹平安.
(平安).
剛才那一段經文.
你發覺它交代了.
我們很熟悉.
以瑪奈利.
臣與我們同在的典故.
不是出自《馬太福音》.
是來自哪裡?.
而是在書.

$^{81}$典故在那裡.
今天我們會看.
這段經文對當時的百姓.
到底帶來甚麼提醒.
又或者安慰.
稍為了解一下.
當時到底發生甚麼事.
我們看看相關的歷史背景.
在程序表.
或者你看這幅PowerPoint.
都會看到有一個地圖.
說回當時的一些情況.
麻煩給我那張圖出來.
對.
當時的亞蘭.
即是這裡寫著大馬士革.
或者今天敘利亞一帶.
與當時的以法連.
以法連即是甚麼呢?.
即是自從所羅門的王國分裂.
一分為二.
分了南北.
南國是猶大.
北國是以法連支派代表的以色列.
這兩個國.
為了對抗他們的共同敵人.
就是亞述.
看到很大塊橙色的那個嗎?.
你看到版圖超級巨大.
看到亞述帝國的入侵.
當時為了要聯手對抗亞述.
於是這兩個國家.
亞蘭和以法連.
就想遊說下面一些國家.
看到有摩亞,阿曼,猶大.
甚至乎以東.
他們很想拉攏這些國家.
組成一個反亞述的聯盟.
這些國家基本上都答應.
但唯獨猶大國不肯.

$^{121}$因為猶大王阿哈斯當時選擇的是相反.
他要靠攏的是亞述帝國.
所以拒絕與亞蘭和以法連結盟.
很明顯有後果.
不肯就範.
結果就觸怒了亞蘭和以法連.
於是他們就聯手進攻耶路撒冷.
就是我們今天所看到.
以賽亞書第七章的歷史背景.
並且他們不單進攻.
稍後會出現恐嚇.
他們擺明不只是攻打.
而是要取締.
要取締阿哈斯.
要散去他的王位.
然後安插另一個人.
成為亞蘭和以法連的內應.
簡單來說.
將猶大國變成一個傀儡國.
為了解除這場危機.
你看到阿哈斯做了兩個行動.
第一是隱約記載在第四節.
這裡說他去水溝頭.
很自然打仗.
如果怕被人圍城.
第一件事就要看什麼.
看水源.
他要確保假如敵軍真的殺到來.
甚至圍攻耶路撒冷的時候.
他要確保城內有足夠的食水供應.
另一方面亦要及早佈防.
因為如果敵人斷了水.
或是敵人在水溝下毒.
就沒辦法了.
圍住城內.
自己手下也沒人.
因為斷水.
即是遲早會被攻破.
所以第一件事他就要觀察水源.
另一方面.

$^{161}$在二次亞書沒有記載.
不過在列王記有記載.
他做了另一個舉動.
我們剛才說他靠攏亞述.
他去找亞述王幫忙.
他被兩個王夾擊.
於是找一個更強的王.
來幫忙解除這場危機.
經文印在《廣道大綱》裡.
列王記交代了.
阿哈斯差派使者去亞述王.
提甲拉比列色.
提甲拉比列色跟他說.
我是你的僕人.
我是你的兒子.
現在亞蘭王和以色列攻擊我.
求你來救我脫離他們的手.
阿哈斯將獵王說.
殿女和王公.
所有金銀都送到亞述王為禮物.
即是代表什麼?.
他現在臣服於亞述.
稱亞述為他的王.
簡單說.
很明顯.
這兩個舉動.
你會看到阿哈斯判斷眼前.
這是什麼危機.
我們今天術語是.
一個政治危機.
或者一個軍事戰爭危機.
我的國家被攻打.
甚至我的王位.
即將被人散奪.
他定性了這場是一個.
政治風波.
是一個軍事危機.
很自然.
他作為一國之君.
他會怎樣解決?.

$^{201}$透過軍事手段.
透過外交方法.
化解這場危機.
不過.
我們說阿哈斯除了是.
猶太國的君王.
他還有另一個身份是什麼?.
誰才是他真正的王?.
上帝才是.
他的身份除了是一國之君.
他更是上帝的子民.
上帝的僕人.
上帝才是他真正的主宰.
真正的君王.
換句話來說.
阿哈斯選擇用軍事外交.
來處理危機是一回事.
但上帝怎樣看待整件事.
這才是關鍵.
不可以解錯題.
解錯題.
你做多少都是白費心機.
什麼意思呢?.
另一段經文給我們很清楚.
同樣印在《廣道大綱》.
《歷代志下》第28章15節.
讀一讀真正的原因是什麼.
為什麼他讓兩個國家來攻打.
這裡說.
阿哈斯登基時二十歲.
在耶路撒冷作王十六年.
不像他祖大位.
行耶和華眼中看為正的士.
卻行以色列諸王的道.
又著做巴力的仗.
並且在恩怨子谷.
燒香.
用火焚燒他的女兒.
行耶和華在以色列人面前.
所趕逐的愛邦人那可憎的士.

$^{241}$並在幽壇上.
山崗上.
各青翠樹下獻祭燒香.
所以耶和華他的神.
將他交在亞蘭王手裡.
亞蘭王打敗他.
擄了許多的民.
帶到大馬士革去.
神又將他交在以色列王手裡.
以色列王向他大行殺戮.
透過這段經文.
我們開始看到真相.
表面上.
阿哈斯是因為拒絕.
與亞蘭和以法蓮結盟.
所以遭受攻擊.
這是表面的情況.
阿哈斯也是這樣判斷.
但實情是什麼?.
上帝透過他的話語.
告訴我們.
阿哈斯和百姓.
是因為行上帝眼中看為惡的士.
於是上帝透過亞蘭.
和以法蓮作為管教工具.
來刑罰教訓猶大.
換句話說.
阿哈斯做了這麼多.
但其實錯判了整個形勢.
例如今天考試.
考試第一件事.
對著題目最重要是什麼?.
解答題目.
解錯題目.
寫滿整頁紙都會怎樣?.
都是沒有用.
零分,一個交叉打下去.
同一個情況.
阿哈斯將一個屬靈的問題.
核心問題是他和上帝.

$^{281}$甚至乎猶大和上帝出了問題.
但他將它當作軍事外交問題處理.
即是解錯題.
又或者用另一個層面.
他只是處理一個表面問題.
我被人打.
但他沒有探究根源是什麼.
是我得罪了上帝.
所以我們說這叫做.
寫本續末.
就是阿哈斯的情況.
或許有時都是我們的情況.
例如在日常生活中.
特別你有小朋友.
可能他在幼稚園小學.
甚至中學階段.
總會離不開手機.
打機甚至上網.
當他有沉溺的時候.
我們當然很想幫他撕掉隱形手機.
但如果我們只針對.
子女表面的行為.
我們不去探究.
不去證實到底這些表面行為.
背後的原因是什麼.
他沉迷是否因為他學習太大壓力.
所以要找出路抒發.
他沉迷是否因為他很想在裡面.
找到社交,找到認同.
所以他不斷打機,不斷上網.
他沉迷是否因為他和家人很疏離.
沒有人關心他.
就靠手機,上網關心他.
如果我們不去探究這些問題.
反而只是處理他表面的問題.
我們就是所謂的「斷錯症」.
我們會用不同的方法.
去制止這些行為發生.
可能很多在座家長都試過.
首先勸他不要打.

$^{321}$吃飯,睡覺,總課.
勸不掉又怎樣?.
好像大家未試過.
未試過都知道會怎樣做.
就罵他.
甚至剷他.
通常都不收效.
可能剷一時不剷一世.
然後怎樣?.
然後開始有升級的行動.
手機.
拿來.
如果他那部是電腦.
你收不到.
放在這裡又怎樣?.
「蚊插數」.
斷電.
斷電都不行.
他還可以上網.
斷網.
其實都沒有用.
上網有WIFI.
所以最後.
可能短期內.
你會解決了.
阻止他上網的沉溺.
但長遠要找甚麼數?.
你會製造更多雙方的對立面.
他會越來越不滿你.
你又會越來越不滿他.
結果關係會破裂.
其實阿克斯的情況.
都在走這條路.
他可能最終真的解決到.
軍事問題,外交問題.
但和上帝的決裂會越來越深.
因為斷錯症.
所以這段經文.
短短幾節給我們看到.
原來擁有正確判斷問題的能力.

$^{361}$真的很重要.
在日常生活裡面.
你斷錯症.
其實是很大件事.
特別在危急,亂的時候.
突然有風波來臨的時候.
你更加需要一個.
正確判斷問題的能力.
來減少錯判形勢.
或做錯決定的機會.
正如今天.
我想很多時.
無論看新聞或上網.
很多時人們都會強調.
我們現在要對社會的發展.
要對國際的時局.
要有大局意識.
你可能聽不少這四個字.
要有大局意識.
你不能單單看到事物的表面.
不能單單這樣.
你更加需要了解.
事情背後的本質是甚麼.
真相到底是甚麼.
這樣你才可以為未來.
做一個正確的決定.
或做好風險的管理.
同樣基督徒都要有這個意識.
不過是屬靈上.
我們有沒有那種屬靈的大局觀?.
我們有沒有經常用上帝的目光.
來詮釋生活所遇到的人和事情?.
經文提醒我們要有這樣的意識.
我們懂不懂得.
把聖經的教導做一些整合.
做一些應用在生活不同層面.
不只是在教會.
在生活層面要把聖經活用出來.
你需要有反省和整合.
我們是否習慣用屬靈的目光.

$^{401}$來審時度勢.
又或我們是否敏銳.
上帝的大靈和提醒.
以致我們可以在不同處境.
無論對上帝,對人,對事.
你都可以有正確的決定.
你都可以做出合宜的回應.
不會解錯題目.
這段是經文給我們的提醒.
簡單來說,這種屬靈的大局觀.
回到一個最基本的操練.
就是我們和上帝的關係.
其實和祂是否緊密.
不只是知道祂是怎樣.
祂想我們怎樣.
我們有沒有跟隨.
我們要有這份意願.
我們才能慢慢建立屬靈的大局觀.
就好像內功修煉.
其實急不來.
經年累月來熬出來.
來培養出來.
但久而久之.
你會發覺你擁有一份.
獨立屬靈的洞察力.
久而久之你會發覺.
即使你多麼忙碌,多麼混亂.
其實你能夠靠聖靈的提醒.
靠聖經的教導.
你懂得擺正位置.
擺一個適當的位置.
你可以站穩.
意思是你會意識到.
無論天塌下來.
好像阿哈斯.
你第一時間想起.
你真正的身份是甚麼.
你是上帝的兒女.
你是上帝的子民.
你亦是祂的僕人.

$^{441}$你這樣想的時候.
你立即擺正位置.
我要回應真正服侍的.
到底是哪一個.
是這個世界.
還是上帝和祂的話語.
就算有危機.
有急難臨到的時候.
我們仍然能夠確信.
上帝沒有變.
上帝很穩穩地掌握著這個世界.
萬大事都由上帝擔當.
所以我們不會像.
經文所提到的.
阿哈斯或百姓.
過份害怕和憂慮.
不會因為這樣而亂尋出路.
我們懂得鎮定.
懂得耐心等候.
上帝對我們的指示和提醒.
不會像阿哈斯.
打到來.
唯有自己解決.
自己想辦法.
而看不到整個真正的大局.
其實是上帝管教.
第一件事回到上帝那裡.
所以你看到.
回到今日的經文.
你見到第二節面對這個場景.
經文形容得很好.
說他們心都跳動了.
因為敵人已經在附近駐紮.
準備進攻.
百姓現在的情況就像樹林的樹被風吹一樣.
用廣東話說.
有得震就沒得睡了.
不過你看到.
上帝很有恩典憐憫.
上帝主動透過先知.

$^{481}$來接觸阿哈斯王.
主動去找他.
為他提供安慰和指引.
我們繼續看下去.
第四節.
你見到.
上帝透過以賽亞先知.
提醒阿哈斯.
剛才也讀了.
要怎樣?.
四個字.
謹慎.
安靜.
不要讓你眼前的恐懼.
支配你自己的行動.
因而自亂陣腳.
在上帝眼中.
亞蘭和二法蓮的攻擊.
只不過是快要燒完的火把頭.
表面上你看下去.
很大煙.
你以為它準備燒得很大聲.
聲勢很大.
但上帝說.
不是,它很大煙.
其實是甚麼?.
就是快要燒完.
即是你有沒有燒東西.
試過炭燒完.
你倒水下去.
那時是否最大煙?.
就燒完.
但你如果錯估了.
你以為很大煙.
即是越來越厲害.
其實剛剛相反.
上帝說.
亞蘭和二法蓮.
只能攻擊猶大.
但不能打贏猶大.

$^{521}$為甚麼?.
我們試試用淑玲的角度去想.
為甚麼不會打贏?.
因為上帝的旨意.
剛才也說了.
祂只是透過亞蘭和二法蓮.
做些甚麼?.
是管教猶大.
並不是要透過他們.
消滅甚至取締大維加.
所以二法蓮和亞蘭.
他們想要取締猶大國.
即是和上帝對著幹.
和上帝作對.
上帝會否讓他們這樣做?.
自然不會.
所以上帝對阿哈斯說.
你放心.
他們不會得逞.
因為我不會讓他們得逞.
所以這裡很清楚.
上帝提醒阿哈斯.
無論你見到這兩個王.
他們對你說甚麼.
多麼威嚇你.
這不是重點.
根本事情不會按照.
他們的劇本發生.
重點是.
上帝這位天上的大君王.
現在應許了你甚麼?.
以及你是否認定.
上帝所說的話.
比敵人更有份量.
更值得你相信.
這才是重點.
我們聽到很多聲音.
有很多對你不利.
威嚇的聲音.
搞擾你的聲音.

$^{561}$但你聽到上帝對你應許過甚麼嗎?.
如果你聽到.
就會像上帝所說.
你能夠站穩.
你能夠安靜.
所以同樣.
上帝現在透過先知問阿哈斯.
你信誰多一點?.
信這兩個王所說的有份量.
還是信天上的大君王.
所應許的有份量?.
如果你聽得懂.
你有正確的選擇.
你就可以安穩.
你就能夠回轉.
專心依靠上帝保守.
不會再靠自己亂來.
也不會再去投靠亞述.
因為那不是你真正的王.
而為了激勵阿哈斯對上帝的信靠.
經文繼續說.
第十至十一節.
「又說有曉遇阿哈斯說.
你向耶和華你的神求一個兆頭.
或求顯在深處.
或求顯在高處」.
上帝的意思其實就是.
阿哈斯到底要怎樣才肯相信.
上帝不會被亞蘭和以法蓮的計謀得逞.
你要怎樣才相信?.
你要怎樣才相信.
上帝不會容許猶太被消滅?.
你儘管放膽向上帝求一個記號.
求一個印證.
就算天高或陰間深.
上帝說都可以.
只要你肯求.
只要你看見這個記號或印記.
你就會知道上帝對你的應許.
是一定會成就.

$^{601}$上帝對阿哈斯好不好?.
真的很好.
上帝用盡不同的方法.
去將不信任上帝.
想走開的阿哈斯.
怎樣都拉回來.
就好像上帝給了一張信心的支票.
給阿哈斯.
上帝簽了名或押.
不寫銀碼.
放膽填.
上帝說一定會兌現.
但是阿哈斯夠不夠信心?.
夠不夠膽填下去?.
還是已經填了心水不填?.
一定不會兌現.
彈票給上帝算了.
你看見阿哈斯有沒有領情?.
這句說話他講得很弔詭.
但我們知道他不領上帝的情.
他說我不求.
我不試探上帝.
表面看來是不是很勁虔?.
他又沒有說錯.
因為律法禁止人試探耶和華.
但現在的情況是甚麼?.
是上帝很清楚告訴你.
你向我求.
即是上帝要試驗阿哈斯的信心.
根本不存在試探上帝的可能性.
阿哈斯現在不肯祈求上帝反而是甚麼?.
擺明得罪上帝.
擺明抗旨.
但他又不夠膽明目張膽地說.
於是會怎樣?.
我們說所謂的扮勁虔.
我不求,我不試探耶和華.
希望可以這樣灑走上帝.
或是灑走二賽亞才知道.
阿哈斯這種情況有沒有讓你想起一類人?.

$^{641}$可能你在工作或不同場景都遇見過這類人.
這種人是怎樣?.
他很清楚自己做人很虛假.
譬如他明明要幫自己找著數.
就說我幫大家謀求福利.
我為大局.
但實際上是自己的好處被他放在口袋裡.
很虛假的.
他知道自己是這樣的人.
而你又知道他是一個很虛假的人.
然後他又知道你知道他是一個虛假的人.
知道我說甚麼嗎?.
知道我?.
非常好.
但他怎樣?.
他都繼續在你面前假扮虛假.
這種人是否很討人厭?.
很討人厭吧?.
他總是會有不同的說辭.
來替自己的醜惡掩飾.
總之底線是甚麼?.
大家知還知.
不要踢爆我.
不要拆穿我的戲.
我就會繼續演下去.
我會繼續虛假下去.
我們說人無恥又如何?.
可能真的會無敵.
你真的奈他不何.
阿哈斯也是這樣.
你奈他不何.
好像很難搞.
死都要賴皮.
明明不相信上帝都扮得很勁.
但我們不喜歡和這類人相處.
我們很怕和他們共事.
因為我們骨子裡都不喜歡虛假.
我們會覺得很反感.
甚至如這段經文所說.
很厭煩.

$^{681}$所以阿哈斯可以說是我們一個反面教材.
提醒我們不要落入這種虛假的試探.
特別當我們面對著.
我們的大君王上帝的時候.
與其我們假扮敬虔.
假扮相信他.
倒不如在上帝面前坦白承認.
我真的會有軟弱的時候.
我真的會有信心不足.
懷疑上帝的時候.
求上帝寬恕,接納.
重新建立我們對他的信靠.
這樣在上帝面前真誠否白.
反而可以蒙上帝月立.
所以這段經文提醒我們.
不要學阿哈斯.
企圖用虛假的敬虔.
來掩飾自己骨子裡其實是怎樣?.
其實是心硬,不相信.
上帝其實不會向你求任何兆頭.
任何恩基.
因為我根本沒有打算相信你.
我都沒有打算依靠你.
我已經決定了.
將自己和國家的前途.
交給了亞述.
壓在亞述王身上.
我不會改.
所以無論上帝怎樣說.
我都會站硬.
不過不是站在你那邊.
是站在亞述那邊.
大概這就是阿哈斯的心底話.
但又不想得罪上帝.
又怕上帝會震怒.
所以就扮下敬虔.
以為這樣可以胡混過關.
可以嗎?.
先知立即踢爆他.
先知立即斥責阿哈斯這種虛假敬虔.

$^{721}$連上帝我們說滿有恩典廉物.
不輕而發落的上帝都會覺得.
厭煩,很厭悶.
你不要再裝模作樣了.
已經踢爆了.
先知說.
阿哈斯既然你這麼心硬.
你就是不肯向上帝求一個兆頭.
上帝現在不理你.
上帝親自向猶太起一個兆頭.
上帝自己來.
顯明他對猶太的旨意.
是必定會成就.
你見到很有趣.
兩股好像不會動搖的力量.
一股是阿哈斯人的力量.
不會動搖.
另一股是上帝的旨意.
不會動搖.
誰較強?.
我們知道是上帝.
但我們骨子裡傻傻的.
會以為自己較強.
好像阿哈斯.
總是跟上帝對著幹.
這裡說.
大家讀出來好嗎?.
我想請大家.
十四至第十七節.
以賽亞書第七章.
預備起.
因此,主主既要求你每一個兆頭.
應有共對懷人心志.
及丹田明教以馬來歷.
到達曰德希臘七十四歲.
同求阿拉伊的大有與同文.
因為在這大小萬位的.
與德希國的親子之心.
你所身份的一萬五千萬.
國籍的必至見證器.

$^{761}$你共同之始.
一萬歲國光的地位.
及至立到的地位.
自從二十三年.
你開開大學.
你尋求的地位.
是真實的.
好,謝謝.
這裡有一段我們很熟悉的經文.
「必有同女懷孕生子」.
原文的意思其實.
不一定是一個神蹟.
即是性靈感人那種.
好像馬太夫所說.
原文的意思只是說.
一位已介乎適婚年齡的年輕女子.
即將要生一個孩子.
當然可以是性靈感人.
但亦可以是自然的分娩現象.
當然我們知道這個兆頭.
最終我們在馬太夫所見到.
是會完全應驗在耶穌基督身上.
這是必然的.
不過聖經學者亦提醒我們.
不要因為這樣而忽略了這個兆頭.
到底對當時的亞克斯王.
以及當時的猶大百姓.
到底的意義是什麼.
因為上帝這裡說得很清楚.
「這個孩子還未曉得器樂澤善之先」.
按照猶太人的標準.
大約要到十二至十三歲.
才有所謂器樂澤善的能力.
才能分辨到.
換句話說.
上帝這裡是預告.
你會見到將會有一個孩子出世.
而這個孩子在還未長到十二,三歲之前.
以法蓮及亞蘭早已被棄絕.
而歷史亦證實了.

$^{801}$亞蘭及以法蓮分別在兆頭發出後.
大約三年及十二年後.
被亞述帝國吞滅.
換言之.
上帝這裡說的是.
在不久的將來.
猶大不再需要面對.
亞蘭及以法蓮的攻擊.
因為在稍後時間.
這兩個國家不會再存在世上.
不過重點來了.
第十七節上帝這裡說什麼?.
沒錯,你現在所擔心的危機.
即將會解決.
不過取而代之的是.
會被亞蘭及以法蓮.
更可怕,更殘酷的亞述帝國.
來攻擊你.
而亞述對猶大的威脅.
上帝說是前所未有.
自從所羅門的王國一分為二後.
都未試過這麼嚴重.
所以.
以馬來尼.
上帝與我們同在.
這個名字.
包含了正反兩個意思.
一方面它象徵著.
上帝會與猶大同在.
所以他們不需要害怕.
鄰近這兩個小國的攻擊.
但同時也由於.
猶大在這次事件上拒絕.
信靠上帝.
反而走去倚靠.
一個更可怕的敵人.
期望以亞述來解決困難.
或者用屬靈的角度看.
就是以亞述王取代上帝這位君王.
所以稍後.

$^{841}$猶大會發現這位與他們同在的神.
將會使用他們所倚靠的亞述王.
來攻擊他們.
或者說教訓他們.
阿哈斯和猶大的心硬和不信.
沒錯,可能會解決眼前的危機.
但換來的是一個更大的恐懼.
所以,弟兄姊妹.
經文說得很清楚.
以馬來理的兆頭.
到底對當時的人是甚麼一回事.
當然亦成為我們今天.
要反省的事情.
首先.
神與我們同在.
對你來說是好消息.
還是壞消息?.
神與我們同在.
是好消息.
還是壞消息?.
不要這麼快回答.
因為不是是非題.
答對不一定有分.
答對之後做才有分.
所以這是一個反省問題.
是一個生活實踐的題目.
你看到阿哈斯.
他知不知道上帝與他同在嗎?.
知道了,已經說明了.
上帝與我同在,還不知道?.
一定知道.
但知道上帝與他同在.
對他的生命有沒有影響?.
有沒有影響?.
沒有.
他沒有因此而選擇.
上帝我會回轉,不靠阿術,信靠你.
沒有.
沒有回轉.
這正是我們要反省的地方.

$^{881}$知道上帝與你同在是不足夠的.
我們有沒有與上帝同行才是見真章.
無論我們倒背如流.
懂得寫原文也好.
你知道上帝與你同在不足夠.
你要因為知道上帝與你同在.
你如何回應上帝.
跟隨他,這才是信仰的真章.
很可惜,阿哈斯王.
他面對這個試驗.
他胖了.
我們呢?.
上帝說得很清楚.
要回轉,不可以再繼續心硬.
不可以再堅持不信上帝.
不聽先知的勸告.
阿哈斯的問題是.
就算他明知上帝介入.
他死也不改.
所以同樣的,姐妹們.
假如我們真的認定.
上帝與我們同在.
但我們怎樣也不肯讓上帝碰我們的生命.
動我們的生命.
不肯讓上帝剎停我們一些.
他不喜悅的事情或想法.
我們拒絕讓上帝在我們身上作王的話.
其實我們跟阿哈斯是一樣的.
以瑪麗莉神與我們同在這個消息.
大概不一定是平靜安穩的消息.
因為上帝的管教一定會想到.
上帝一定會迫使我們回到他身邊.
歸回他.
上帝不會對我們袖手旁觀.
當然是出於他對我們的愛.
和不離不棄.
而另一方面.
假如我們真的相信確認.
上帝與我們同在的話.
我們就要時刻想.

$^{921}$我們是否時刻與上帝同行.
無論環境如何.
際遇如何.
我們是否都竭力走近上帝.
相信他 信服他的旨意.
經文提醒.
單單知道上帝同在不足夠.
我們要操練如何配合上帝的旨意.
事實上你看阿哈斯今天的情況.
其實找到一些應用的方向.
假如上帝真的要求我們.
有些事情 有些人.
我們放下.
不要再過分操控.
過分的人為干預.
我們安靜.
等上帝工作.
可能是你的家人.
可能是你的工作夥伴.
甚至在教會裡的一些場景.
上帝說你放手.
看上帝工作.
你願不願意配合?.
可能真的會不習慣.
因為人多多少少要有掌控的慾望.
要你放手.
你第一個感覺就是不安全.
於是你很快又會惹上.
但我們是否真的相信.
上帝有最好的安排.
所以我寧願.
試試放手.
交回給上帝.
我們有沒有一個這樣的.
靈大局官認定我們只是.
一個暫託的管家?.
這是第一方面值得我們去想.
另一方面假如我們面對一些困境.
我們希望上帝有這麼快.
幫我們解決的話.

$^{961}$這段經歷也成為我們提醒.
假如上帝告訴我們.
我們面對的困難.
面對的危機會持續一段日子.
不會馬上解決.
未必會按我們的心意發展.
更加不會告訴我們.
要等多久.
但上帝應許.
你放心.
祂的恩典夠你用.
祂的保護也是周詳.
你願不願意靠上帝.
和這些逆境來共存?.
上帝在裡面.
你願不願意全心信服祂?.
不容易,弟兄姊妹.
因為我們今天講求時代.
講什麼?.
講快.
講即食.
昨天是什麼日子?.
昨天有很多日子.
可以是一生一世的夫婦節.
可以是.
如果紀念一次大戰.
是和平紀念日.
1918年那次的一次大戰.
又或者如果你比較熟悉國情.
昨天是光棍節.
鼓勵你購物.
淘寶,天貓.
什麼意思?.
不要等了.
快點買,便宜了,划算了.
不要錯過了.
與其你那部壞了.
慢慢修補,不如怎樣?.
買一部新的.
更快更便宜.

$^{1001}$社會不讓我們有耐性.
買東西是這樣的.
和人的關係也是這樣.
和上帝的關係會不會也是這樣?.
不習慣等.
不習慣等上帝的工作.
但你留意我們所相信的.
上帝聖經的形容是怎樣?.
不急躁,有耐性.
總是耐心等候我們回到.
祂也期望我們耐心等候.
上帝的工作.
所以弟兄姊妹,你有留意.
等候本來是我們.
屬靈生命的一部份.
你有機會在聖經裡.
搜尋「等候」這兩個字.
等候上帝的關鍵字.
你看看你找到多少個.
很多時候是等.
不可以急.
意味我們時時刻刻都要調教.
我們的生活節奏.
配合上帝這位君王的大力.
所以最後.
說了一些很灰的東西.
我們有些盼望才要.
最後我想和大家讀一段經文.
印在程序表最後.
這段經文不是在第七章.
是在第十二章.
其實第七章至第十二章.
是一個段落.
以蔡亞書來說.
作者說了很多比較悲觀的事.
他以一首詩歌.
一首盼望的詩歌作結束.
提醒我們無論這個世界變成怎樣.
甚至我們真的有軟弱.
有膽怯.

$^{1041}$甚至懷疑遠離上帝的時候.
但這段經文這首詩歌.
其實勉勵我們.
回到上帝那裡.
上帝想我們是這樣.
你重新信教上帝.
因為他才是你生命的王.
你認定他這樣的時候.
你就會有信心.
來與上帝同行.
因為一切仍然在掌管的裡面.
甚至你能夠讚美上帝.
就好像這首詩歌.
我們一起讀出來好不好?.
第十二章的一至二節.
預備起.
到那一日.
你必說:耶和華.
我要稱謝你.
因為你雖然向我發怒.
你的怒氣卻已消轉.
你又安慰我.
神賜我的拯救.
我的依靠.
並不懼怕.
因為主耶和華是我的力量.
是我的詩歌.
也成了我的拯救.
我們不單單宣讀.
我們現在唱出來好不好?.
你留意你的回應詩.
其實也是這一首的歌.
以賽亞書十二章第一至第二節.
我們一起唱吧.
請美心帶領我們.
\newpage



\section{約書亞記 14:6-14}
\label{sec:SCWtMBw63Gk}
\textbf{2023年11月19日長者主日崇拜直播｜羅志雄牧師｜老而彌堅的迦勒｜約書亞記14\_6-14}
\newline
\newline
連結: \href{https://youtube.com/watch?v=SCWtMBw63Gk}{\texttt{https://youtube.com/watch?v=SCWtMBw63Gk}} ~~~~ 語音日期: 2023-11-20
\newline
\newline
\hyperref[sec:KCjIplGBEcc]{\small{< < < PREV SERMON < < <}}
~
\hyperref[sec:index_chronic]{\small{[返順時目]}}
~
\hyperref[sec:_oDlWIvKGkE]{\small{> > > NEXT SERMON > > >}}
\newline
\newline
約書亞記 14:6-14
\newline
\begin{longtable}{cl}
\hline
\hline
章節 & 經文 (和合本修訂版)\\
\hline
14:6 & \begin{tabularx}{0.7\textwidth}{X} 猶大人來到吉甲，約書亞那裡，基尼洗族耶孚尼的兒子迦勒對約書亞說：「耶和華在加低斯‧巴尼亞指著我和你對神人摩西所說的話，你都知道。 \end{tabularx} \\ \\ \relax
14:7 & \begin{tabularx}{0.7\textwidth}{X} 耶和華的僕人摩西從加低斯‧巴尼亞差派我窺探這地的時候，我剛四十歲。我把心裡的話向他報告。 \end{tabularx} \\ \\ \relax
14:8 & \begin{tabularx}{0.7\textwidth}{X} 雖然同我上去的眾弟兄使百姓膽戰心驚，我仍然專心跟從耶和華－我的神。 \end{tabularx} \\ \\ \relax
14:9 & \begin{tabularx}{0.7\textwidth}{X} 那日，摩西起誓說：『你腳所踏之地必要歸你和你的子孫永遠為業，因為你專心跟從耶和華－我的神。』 \end{tabularx} \\ \\ \relax
14:10 & \begin{tabularx}{0.7\textwidth}{X} 現在，看哪，耶和華照他所說的使我活了這四十五年。當以色列人在曠野飄流的時候，耶和華曾對摩西說了這話。現在，看哪，我已經八十五歲了。 \end{tabularx} \\ \\ \relax
14:11 & \begin{tabularx}{0.7\textwidth}{X} 現今我還很健壯，像摩西差派我去的那天一樣；無論是戰爭，是出入，我現在的力量和那時的力量一樣。 \end{tabularx} \\ \\ \relax
14:12 & \begin{tabularx}{0.7\textwidth}{X} 請你將耶和華那日所說的這山區給我。那日你也曾聽說，這裡有亞衲族人，以及寬大堅固的城，或許耶和華會照他所說的與我同在，我就把他們趕出去。」 \end{tabularx} \\ \\ \relax
14:13 & \begin{tabularx}{0.7\textwidth}{X} 於是約書亞為耶孚尼的兒子迦勒祝福，把希伯崙給他為業。 \end{tabularx} \\ \\ \relax
14:14 & \begin{tabularx}{0.7\textwidth}{X} 所以希伯崙成了基尼洗族耶孚尼的兒子迦勒的產業，直到今日，因為他專心跟從耶和華－以色列的神。 \end{tabularx} \\ \\
[1ex]
\hline
\hline
\end{longtable}
$^{1}$今日要讀的聖經是約書亞記第十四章六至十四節.
聖經舊約第二百八十至二百八十一頁.
這段聖經是記載上帝的應許.
加勒德希伯倫遺孽.
約書亞記第十四章第六節.
那時猶大人來到吉甲見約書亞.
有基尼洗族耶夫尼的兒子加勒對約書亞說.
耶和華在迦底斯巴尼亞.
指著我與你對神人摩西所說的話.
你都知道了.
耶和華的僕人摩西從迦底斯巴尼亞.
打發我窺探這地.
那時我正四十歲.
我按著心意回報他.
然而和我上去的眾弟兄使百姓的心消化.
但我專心跟從耶和華我的神.
當日摩西起誓說.
你腳所踏之地.
定要歸你和你的子孫永遠為業.
因為你專心跟從耶和華我的神.
自從耶和華對摩西說這話的時候.
耶和華照他所應許的.
使我存活這四十五年.
期間以色列人在曠野行走.
看啊!現今我八十五歲了.
我還是強壯像摩西打發我去的那天一樣.
無論是征戰是出入.
我的力量那時如何.
現在還是如何.
求你將耶和華那日應許我的這山地給我.
那裡有阿納族人並寬大堅固的城.
你也曾聽見了.
或者耶和華照他所應許的與我同在.
我就把他們趕出去.
於是約書亞為耶夫尼的兒子迦勒祝福.
將希伯倫給他為業.
所以希伯倫作了基尼洗族.
耶夫尼的兒子迦勒的產業.
直到今日.
因為他專心跟從耶和華以色列的神.

$^{41}$剛才長者詩班唱詩時.
我想起去年的事.
不知道大家在長者主日崇拜時.
在哪個地方?.
在四樓.
阿東城,記憶猶新.
那個情景跟今天很不同.
我想是上帝的恩典.
在長者日.
香港今天是長者日.
我想每一位長者都能經歷上帝的恩典.
我們邀請大家一起有個禱告.
天父上帝懇求您憐憫與我們同行.
與我們同在.
我們經歷不同歲月.
但您仍然用您的恩典浮避.
在過程中,在我們的人生中.
有很多起起伏伏.
令我們有遺憾,有懊悔.
甚至有失敗的地方.
但上帝您仍然與我們同行.
今天我們藉著迦勒這段經歷.
幫助我們.
讓我們更有信心去投靠您.
願您帶領我們.
誰聽我們在您面前的禱告.
奉主耶穌基督的名.
阿們.
每當聽到長者.
會想起什麼呢?.
可能會想起您的年歲.
或者年歲的時候.
會與我們老人家掛勾.
不知大家呢.
問問大家.
你會不會想長壽?.
想長壽就舉手吧.
有人舉手.
最終沒有人舉手.
長壽意味著什麼呢?.

$^{81}$意味著老.
沒有舉手的人.
是不是要長壽不要老呢?.
我想這個框框.
也成為我們人生中.
其中一個很矛盾的狀況.
但其實我們很難打破這些框框.
要長壽不要老.
但原來人生的成長就是這樣.
原來一個發展心理學家艾力遜.
他提出人生長有八個階段.
由出生,嬰兒期,學前期,少年期.
青年期,成年期,中年期.
到晚年.
今天我們要講的.
關於我們人生的晚年.
是否老而離肩呢?.
還是我們覺得老到什麼都不想做呢?.
我們回看迦勒.
今天保齡姐妹為我們讀的經文.
迦勒在他八十五歲的時候.
回望他的人生.
但很奇怪.
為什麼會在《若書雅記》出現呢?.
我想是作者將這段經文放在《若書雅記》.
特別是分地的狀況.
其實真的反映出我們需要學習的一件事.
回看迦勒.
在這裡出現的時候.
其實就是一個晚年.
八十五歲的高齡.
他講這段說話.
由第六節的下半節.
一直到第十二節的時候.
其實我們用另一個角度來看的時候.
其實可能是迦勒一個獨白.
來向《若書雅記》講出他人生其中一個片段.
我們沒辦法知道迦勒出生是怎樣.
但如果我們重溫剛才所說的生命成長.
八個階梯.

$^{121}$我們可以推斷迦勒出生.
在他的人生裡.
大部分時間都是在埃及成長.
成長一段日子.
可能是青年期.
之後他就要跟隨整個以色列的族群.
離開埃及.
他就走進去曠野.
我相信這個就是他的壯年期.
但一直走的時候.
不知不覺.
甚至去到中年期.
當他進入迦南地的時候.
已經成為他的晚年.
迦勒回望他一生.
重整他的生命故事.
我們剛才已經看到.
這個就是迦勒的生命獨白.
但從迦勒的人生裡.
讓我們可以學習.
迦勒其實是一個認識自己的人.
怎樣可以看到呢?.
我們可以從第六節到第九節的時候.
可以看到這兩方面.
他能夠繼續仰望神.
他能夠知道神必定賜福給他.
在這段經文裡.
迦勒強調了兩次.
他說在迦底斯巴黎亞的事.
在這裡問問大家.
不知道大家人生裡.
有沒有經歷過迦底斯巴黎亞的事件呢?.
可能未必和約書亞或迦勒一樣.
但這件事件其實在民數記第十三章第十四章裡有記載.
到了迦勒回想的時候.
他沒有把很多細節重述.
但大家可以來看民數記十三十四章.
但對於迦勒來說.
當時這件事件歷歷在目.
以至當他和約書亞說的時候.

$^{161}$他有很多將它刪減.
用一兩句說話來表述他自己當時的狀況.
甚至在這些歷程裡.
這些事件裡.
他對自己的認識回望.
一直有所成長.
我很想和大家分享第七節.
第七節裡是這樣說的.
耶和華的僕人摩西從迦底斯巴黎亞打發我窺探者弟.
那時我正四十歲.
我安著心意回報他.
如果按照民數記裡的整個故事.
迦勒和其他十一個探子.
當時被摩西揀選.
他們應該是很出色的領袖.
以至屬於每一個支派裡的領袖.
迦勒也被選上.
可想而知.
當他回想的時候.
他和上帝的關係非常密切.
正如我們今天在教會服侍.
如果我們生命中有很多艱難.
但原來在艱難中面對上帝.
他和上帝的關係很好的時候.
我想這就成為他人生的閱歷.
他可以帶領著教會.
正如昔日迦勒那樣.
他帶領著宗族成為探子.
跟著摩西的說法.
成為探子走進迦勒的地方.
經文裡告訴我們.
看到當時的狀況.
所有探子,十二個人,包括迦勒在內.
他如實觀察到甚麼.
就告訴摩西.
但很明顯.
迦勒除了講到觀察實境之外.
他有一件事是講出他自己心底裡的說話.
他說「我按著心意回報他」.
新譯本講得比較清楚.

$^{201}$他說「我照著我心裡所想的.
向摩西回報」.
迦勒他有甚麼心所要說的呢?.
看到當時的實境.
阿納族是偉人.
當地城牆堅固寬大.
但他又看到.
是流瀨與物之地.
是上帝所應許他們一定能夠進去的.
他心裡這樣想的時候.
如果我們看民數記的時候很清晰.
迦勒就說在百姓面前.
在摩西面前安撫百姓.
我們立刻上去得拿地.
我們足能得勝.
原來迦勒所想的.
他心裡所想的.
要回報給摩西的.
原來那十一個人都沒有這樣想.
唯有迦勒.
他照著他心裡所想.
來向摩西安撫百姓.
如果大家留意民數記的時候.
就會知道.
說完之後.
所有人都聽到很多實情.
但他們選擇性地去聽.
城牆高大.
那裡有很多巨人.
還有七族的人.
不知道是甚麼人.
於是開始起哄.
整個過程.
在迦勒的生命裡歷歷在目.
所以在這個過程裡.
迦勒也形容.
當時那群人是怎樣呢?.
他說百姓的心都消化.
過程中.
那麼多一群人.

$^{241}$按照當時的估計.
有五,六十萬人.
全部在起哄.
當時的狀況.
可以說是失控.
但是迦勒.
因為他每一天.
跟上帝有很好的經歷.
他專注在上帝那裡.
他不是看眼前的環境.
他也不是去跟風.
那十個探子.
回報的時候.
按照民數記記載.
用噩信來回報.
即是說他們演繹的.
跟迦勒心想的不一樣.
原來每個人對實境.
對眼前所見的東西.
他看完之後.
有他自己內心的演繹.
是很看他跟上帝的關係.
今天迦勒.
給我們提醒.
當人心裡都消化的時候.
其實我們所專注的是什麼呢?.
是在人的身上.
還是在神的身上.
還是在一些環境裡.
我相信迦勒.
不是當時很快地說.
我們立刻上去德那地吧.
我們足能得勝.
因為他知道.
他有很多生命的經歷.
最震撼他的經歷.
就是過紅海.
這位摩西神所差派他.
能夠帶領著整個以色列的民族.
過紅海.

$^{281}$這件事我深信.
就是能夠成為迦勒.
他生命裡一件很震撼的事件.
再加上由摩亞地.
過約旦河的時候.
約書亞也是一樣.
在神裡面所展現的能力.
雖然是這樣.
結局不是一件大團圓.
所有人除了迦勒.
除了約書亞之外.
可以進入今天迦勒所說的這地.
其餘的人都死在曠野.
弟兄姊妹.
不知道你的迦底斯巴黎亞事件.
如何影響你的人生.
你的價值觀.
你和神的關係會是怎樣呢?.
當我去預備這段經文的時候.
我都要想一想.
人生有多少迦底斯巴黎亞事件呢?.
我相信未必是一件兩件.
可能有時我們很忙碌的生活.
我們都遺忘了.
原來有些重大的事件.
在我們人生的歷程裡.
成為我們的影響.
我們繼續看第十至十二節.
這裡迦勒轉向他生命的歷程.
以及他對前景的看法.
在他面對約書亞的時候.
他的年齡已經是85歲的高齡人士.
在座當中.
我相信很多都超過這個歲數.
又或者大家說.
我都未到這個年齡.
不需要想太多.
我只要想眼前就可以了.
但原來作為一個年長高壽的人.
他的生活仍然都要繼續.

$^{321}$有一位很出名的哲學家.
他曾經提過.
他說人生在老年裡我們怎樣過呢?.
原來是一個很艱難.
但透過我們怎樣過老年.
成為我們智慧的基礎.
而這位哲學家.
當他說這兩句話的時候.
在十九世紀的時候.
他說的時候.
其實他只有六十歲.
按照十九世紀當時的年歲.
平均年齡是四十五歲.
當我們今天回顧.
我們八十五歲或更長壽的時候.
二十一世紀.
有很多醫學的發達.
都令到我們人生.
去到這個年歲.
退休的年齡.
大家會做些什麼?.
可能會去享受.
休閒.
或者弄酸為樂.
其實現在說到.
照顧孫子都不容易.
退休生活去照顧孫子.
我聽過一些長者說.
比上班更慘.
時間表還比上班情況更緊湊.
如果有體力的.
你們會怎樣做呢?.
可能會有很多計劃.
學一些新事物.
又或者去見識很多.
但作為八十五歲高齡的家立.
他所選擇的是什麼呢?.
原來在第十節至十二節的時候.
他仍然記得.
神應許給家立.

$^{361}$他自己的應許之地.
所以當他這樣想.
他已經過了四十五年的日子.
如果再由迦底斯巴尼亞事件.
一直進入到迦南.
那裡已經有四十五年.
按照家立所說.
但中間有幾年的時間.
家立和他的支派.
會參加迦南地內的戰役.
如果不是上帝所保守的話.
他在戰役中出入.
縱使他強壯也好.
他在戰爭中.
很可能有機會戰死.
無關乎家立在第十一節說.
「我還是強壯.
像摩西打發我的那天一樣.
無論是征戰還是出入.
我的力量那時如何.
現在還是如何」.
當家立回望他的人生之餘.
他再與他今日生活的狀況.
作一個比較,一個檢視的時候.
他知道他這種狀況.
是上帝給他的恩典.
所以他去到約書亞的面前.
祈求的話.
就是因為有上帝的保守.
讓他能夠經歷生命的體驗.
因為他一直跟隨上帝.
縱然在曠野兜轉三十多年.
因為加底斯巴尼亞事件之後.
還要面對在迦南地.
流瀨與密之地的戰爭.
他經歷了很多風浪.
但他仍然相信自己有足夠的信心.
有足夠的把握.
可以看清楚自己的前景.
他人生裡去到八十五歲高齡.

$^{401}$他仍然是健壯.
就像當日摩西選召他一樣.
弟兄姊妹,我不知道今天的你.
如果你是年長者的時候.
你如何去看自己的生命.
可能有不同的年歲.
未必像八十五歲的迦勒.
但會不會覺得有很多.
現代人的病.
老人病,各樣的事情.
可能都會困擾我們.
但能否有新的眼光.
能否去看上帝給你的經驗.
給你的經歷.
仍然可以有這個生活去繼續.
大概在十月初.
我們一班年長者.
參加一些新的體驗.
健步足球.
不知道大家有沒有聽過.
沒有聽過吧.
應該是基督教青年會YMCA引入的.
其實這個Walking Football.
即是英文.
翻譯成中文叫健步足球.
其實是在英國開始的.
但來到香港.
發現原來很適合長者參與.
於是我知道就聯絡.
基督教青年會.
結果當日我們有三十多位長者.
參與健步足球.
其實在期間都知道.
我們最大年紀參加健步足球.
是多少歲?.
超過九十歲.
我跟職員說.
因為他們從來不會的.
他們由五十多歲開始.
五十多歲至六十歲.

$^{441}$最多是七十歲.
但我跟職員說.
我們有三個過了九十歲.
他們都很奇怪.
但你不要看他們不是坐著的.
他們真的下場.
經過教練來分享.
教他們怎樣打.
健步足球和一般足球的規矩不同.
只要健步是要走.
但對足球的接觸.
令到他們的身體機能更有活力.
這對一班長者是新的體驗.
新的嘗試.
甚至我問一班年長者.
他們說根本沒有打過足球.
其實足球是怎樣的我也不知道.
原來給他們一個新嘗試.
新體驗.
迦勒亦一樣.
他在九十五歲高齡的時候.
他回望他的人生.
他知道他還有一件事未完成.
他盼望他仍然健在的時候.
他要得到上帝曾經所應許的經歷.
就是這一塊所應許.
他腳踏之地.
是要歸迦勒和迦勒的子孫永遠為嚴.
因為迦勒是專心跟隨耶和華.
如果大家看摩西五經.
特別是民數記,申命記的時候.
專心跟隨耶和華.
其實是形容迦勒.
甚至約書亞.
而迦勒自己親自所說.
專心跟隨耶和華.
只是出現在剛才迦勒.
向約書亞說的讀八節.
然而同我上去的眾弟兄使百姓的心消化.
但我專心跟隨耶和華我的神.

$^{481}$這是唯一迦勒在他人生對自己的認識裡.
他再去檢視他現在原來.
他跟隨了這位他所相信的上帝.
幾十年的時候.
他認識自己.
他是專心跟隨耶和華.
縱然他等待了數十年.
他仍然未得到上帝所應許的.
和他子孫永遠遺孽.
但他不是豁出去.
他仍然有一個嘗試.
令他向約書亞求.
甚至在約旦河西分地裡.
原本約書亞要分給其餘九個半支派.
但約書亞在這裡記載的時候.
首先迦勒出場.
其實是要告訴我們什麼呢?.
我們有堅定的信念.
我們就要有行動去回應上帝.
給我們這份我們一定要得到的心志.
約書亞記載的裡面.
我們看到他沒有因為.
加底斯巴尼亞這件事件.
可能一般人.
甚至整個以色列的民族來看.
是成為一個咒詛.
因為是一件遺憾的事.
他不能夠進去攻取迦南地.
但對迦勒來說.
他沒有成為他的遺憾.
甚至沒有成為他的絆腳石.
反而在這一刻.
85歲的高靈.
他仍然健壯的時候.
是一個契機.
一個最好的機會.
甚至讓所有整個民族.
在分地的過程裡.
解除了這種咒詛.
當然這件事件.

$^{521}$令到上帝震怒.
是基於人對上帝的不信任.
基於人對神的應許.
不覺得要重視.
但對於迦勒來說.
他明白.
加底斯巴尼亞這件事件.
是因為人得罪了神.
侵奪了神的主權.
但當他回望.
現在他可以展望新的篇章.
兄弟姐妹.
在我們人生裡.
我們如何去過這經歷.
會不會有些事情.
很遺憾.
到了這個時刻.
我們很想.
我們曾經所願望的.
可以如願以償呢?.
迦勒本著對上帝的忠誠.
本著對上帝的認識.
和他對上帝的明白.
那種屬性.
他這一切.
也讓約書亞了解.
約書亞在聖經裡.
沒有表白和迦勒有任何對話.
因為當迦勒這樣表達的時候.
約書亞一定知道.
迦底斯巴尼亞事件.
所發生的來龍去脈.
所以到了第十二節的時候.
當迦勒去求.
要得山地的時候.
約書亞沒有推搪.
到了第十三節的時候.
約書亞就說.
經文裡記載.
約書亞為耶夫尼的兒子迦勒祝福.

$^{561}$反而約書亞是會祝福迦勒.
讓他在這個仍然是建狀的當下.
他能夠擺脫一些事件帶來的陰影.
他能夠繼續去成就這件事.
所以透過約書亞的祝福.
又加上迦勒他自己真正的努力.
而這個努力不是他自己.
因為他知道上帝的心意.
明白上帝給他的.
如果我們再看後一點.
約書亞記第十五章的時候.
有詳細的講述.
以致他的後代更加明白.
原來一生跟從神.
可以經歷到上帝的恩典.
不單單在約書亞記第十五章記錄.
迦勒再次上戰場的事件.
甚至在士師記再一次有相同的技術.
其實是要告訴我們.
今天每一個.
我們縱然未必有迦勒.
這種人生的閱歷.
但當我們回望的時候.
我們也未必有迦勒.
這種他有正面對人生的正向的看法.
可能我們會帶著很多遺憾.
以致有很多狀況.
我們會退縮.
特別有很多長者.
當我們探訪的時候.
他未必跟他們建立關係的時候.
他們未必能夠很主動甚至開放.
但是這段經文是要告訴我們.
鼓勵我們.
縱然我們的人生未必.
是一個很豐富.
是一個順風順利的人生的歷程.
但不要這樣.
我們會有失望和氣餒.
總括迦勒的一生.

$^{601}$在這裡提醒我們.
我們要擺脫我們生命裡面的一些遺憾.
我們要專注跟隨神.
不要讓這些困難再加深下去.
甚至要轉化我們生命成為有色彩.
甚至可以嘗試開展我們人生的精彩篇章.
因為當我們回望回顧我們的人生的時候.
我相信這就是智慧.
可以承傳下去.
正如摩西在他受膏年逸的裡面.
他看到人生縱然有很多灰暗.
有很多遺憾.
有很多70歲80歲那種歷程.
但原來都是勞苦受煩.
但是摩西沒有這樣放棄.
他仍然用他的生命歷程來鼓勵我們.
摩西說求你指教我們怎樣數算自己的日子.
好叫我們得著智慧的心.
弟兄姊妹.
我們人生的歷程.
在階段裡面.
可能有很多我們不如意的事情.
又或者甚至這些成為我們生命裡面的遺憾.
沒有辦法突破我們生命的框框.
但透過我們每一天的生活.
每一天與神建立好關係.
知道能否像迦勒那樣.
專心跟隨神.
過我們每一天的日子.
四十多年都可以過了.
又或者在大家的人生裡面.
已經經歷了很多歲月.
但約書亞的經歷.
約書亞的人生就是這樣鼓勵我們.
幫助我們.
仍然在前面的路上.
上帝會與我們同行.
上帝會用他的方法.
兼顧我們對他的信靠.
我們一起來禱告.

$^{641}$天父上帝我們感恩.
透過約書亞生命的經歷.
給我們有生命很好的反省.
我們要開展我們新的人生.
我們就要專心跟隨你.
我們對你有更多的認識.
因為認識我們自己.
有很多的限制.
有很多的軟弱.
但祈求你與我們同行.
與我們同在.
在我們未來的人生的路上.
無論是什麼年紀.
我們深信上帝會帶領我們.
讓我們不會帶著遺憾.
不會帶著後悔.
讓我們有正面的人生.
因為我們的生命有主.
有耶穌基督.
讓我們每一天都過著豐盛的生命.
奉主耶穌基督的名.
阿們.
\newpage



\section{約翰福音 3:1-15}
\label{sec:_oDlWIvKGkE}
\textbf{2023年11月26日崇拜直播｜梁家麟牧師｜你們要重生｜約翰福音3\_1-15}
\newline
\newline
連結: \href{https://youtube.com/watch?v=-oDlWIvKGkE}{\texttt{https://youtube.com/watch?v=-oDlWIvKGkE}} ~~~~ 語音日期: 2023-11-28
\newline
\newline
\hyperref[sec:SCWtMBw63Gk]{\small{< < < PREV SERMON < < <}}
~
\hyperref[sec:index_chronic]{\small{[返順時目]}}
~
\hyperref[sec:4NVCY6ZFpaA]{\small{> > > NEXT SERMON > > >}}
\newline
\newline
約翰福音 3:1-15
\newline
\begin{longtable}{cl}
\hline
\hline
章節 & 經文 (和合本修訂版)\\
\hline
3:1 & \begin{tabularx}{0.7\textwidth}{X} 有一個法利賽人，名叫尼哥德慕，是猶太人的官。 \end{tabularx} \\ \\ \relax
3:2 & \begin{tabularx}{0.7\textwidth}{X} 這人夜裡來見耶穌，對他說：「拉比，我們知道你是由神那裡來作老師的；因為你所行的神蹟，若沒有神同在，無人能行。」 \end{tabularx} \\ \\ \relax
3:3 & \begin{tabularx}{0.7\textwidth}{X} 耶穌回答他說：「我實實在在地告訴你，人若不重生，就不能見神的國。」 \end{tabularx} \\ \\ \relax
3:4 & \begin{tabularx}{0.7\textwidth}{X} 尼哥德慕對他說：「人已經老了，如何能重生呢？豈能再進母腹生出來嗎？」 \end{tabularx} \\ \\ \relax
3:5 & \begin{tabularx}{0.7\textwidth}{X} 耶穌回答：「我實實在在地告訴你，人若不是從水和聖靈生的，就不能進神的國。 \end{tabularx} \\ \\ \relax
3:6 & \begin{tabularx}{0.7\textwidth}{X} 從肉身生的就是肉身；從靈生的就是靈。 \end{tabularx} \\ \\ \relax
3:7 & \begin{tabularx}{0.7\textwidth}{X} 我說『你們必須重生』，你不要驚訝。 \end{tabularx} \\ \\ \relax
3:8 & \begin{tabularx}{0.7\textwidth}{X} 風隨著意思吹，你聽見風的聲音，卻不知道是從哪裡來，往哪裡去；凡從聖靈生的也是如此。」 \end{tabularx} \\ \\ \relax
3:9 & \begin{tabularx}{0.7\textwidth}{X} 尼哥德慕問他：「怎麼能有這些事呢？」 \end{tabularx} \\ \\ \relax
3:10 & \begin{tabularx}{0.7\textwidth}{X} 耶穌回答，對他說：「你是以色列人的老師，還不明白這些事嗎？ \end{tabularx} \\ \\ \relax
3:11 & \begin{tabularx}{0.7\textwidth}{X} 我實實在在地告訴你，我們所說的是我們知道的，我們所見證的是我們見過的，你們卻不領受我們的見證。 \end{tabularx} \\ \\ \relax
3:12 & \begin{tabularx}{0.7\textwidth}{X} 我對你們說地上的事，你們尚且不信，若對你們說天上的事，如何能信呢？ \end{tabularx} \\ \\ \relax
3:13 & \begin{tabularx}{0.7\textwidth}{X} 除了從天降下的人子，沒有人升過天。 \end{tabularx} \\ \\ \relax
3:14 & \begin{tabularx}{0.7\textwidth}{X} 摩西在曠野怎樣舉蛇，人子也必須照樣被舉起來， \end{tabularx} \\ \\ \relax
3:15 & \begin{tabularx}{0.7\textwidth}{X} 要使一切信他的人都得永生。 \end{tabularx} \\ \\
[1ex]
\hline
\hline
\end{longtable}
$^{1}$以下是讀經.
經文是約翰福音第三章.
新約聖經一百二十八頁.
請聽我讀.
三章一至十五節.
有一個法理菜人名叫彌哥底姆.
是猶太人的官.
這人夜裡來見耶穌.
說拉比.
我們知道你是由神那裡來作師父的.
因為你所行的神蹟.
若沒有神童在無人能行.
耶穌回答說.
我實實在在的告訴你.
人若不從生就不能見神的國.
彌哥底姆說.
人已經老了如何能從生呢?.
豈能再進母腹生出來麼?.
耶穌說.
我實實在在的告訴你.
人若不是從水和聖靈生的.
就不能進神的國.
從肉生生的就是肉生.
從靈生的就是靈.
我說你們必須從生.
你不要以為稀奇.
風隨著意思吹.
你聽見風的響聲.
卻不曉得從哪裡來往哪裡去.
凡從聖靈生的也是如此.
彌哥底姆問他說.
怎能有這事呢?.
耶穌回答說.
你是以色列人的先生.
還不明白這是麼?.
我實實在在的告訴你.
我們所說的是我們知道的.
我們所見證的是我們見過的.
你們卻不領受我們的見證.
我對你們說地上的事.

$^{41}$你們尚且不信.
若說天上的事.
如何能信呢?.
除了從天降下.
仍舊在天的人子.
沒有人升過天.
摩西在曠野怎樣舉蛇.
人子也必照樣被舉起來.
叫一切信他的都得永生.
各位姐妹早晨.
不要嚇一跳.
我不是忘記了上次的訊息.
上次的訊息.
做了很多解經的工作.
但在應用方面做得很少.
今天是上一篇道的續集.
我們全部時間都是用來做應用的.
將上個月所說的這篇道的內容.
應用在我們的生活中.
所以如果上個月沒有聽過這篇道.
可以回到我們的網頁.
聽那篇道.
找回我曾經說過的.
解過的經.
今天我們純粹做應用.
你們基本上已經唱出了.
我最後一大段要說的訊息.
今天的程序彼此配搭.
我不知道你們覺得做教徒容易不容易.
在我看來.
在今天的社會.
要做一個真實的教徒其實是非常難的.
問題不在於外面的社會和文化.
潮流和我們的信仰格格不入.
也不在於有很多人反對我們的信仰.
堅持信仰要付很重的社會成本.
不是.
真實的情況不是這樣.
外在的壓力.
或者會為實踐信仰帶來一定的困難.

$^{81}$但這些困難通常都能夠解決或忍受.
對信仰最大的挑戰.
在於我們如今太過自由.
我們失落了足以指導我們行為的外在權威.
古代的宗教.
我們中國人的宗教.
講求的是外在的禮儀.
藉著一套敬拜和獻祭的模式.
說得再俗一點.
藉著一隻金豬.
我們的罪孽就得到洩洩竭盛.
我們的困難就得到解決.
願望就得到達成.
一隻豬就能夠解決很多問題.
實踐敬拜和獻祭.
通常有固定的時間地點.
也有專業的宗教從業員幫忙.
他們會告訴我們.
按著什麼程序步驟才是正確的.
然後才可以達至預期的後果.
儀式和權威.
是我們與神明接觸和聯繫的外在憑藉.
這樣的宗教進路.
舊約的猶太人也是這樣.
中國的傳統也是這樣.
耶穌基督所鋪設的救贖之路卻不是這樣.
他宣告.
人能夠與上帝接通.
人能夠得到救贖的唯一出路.
是他自己,耶穌基督自己.
除他以外,別無拯救.
他是道路,真理,生命.
除了基督之外.
獻祭沒用.
律法沒用.
善行也沒用.
就像耶穌與郭井龐的婦人所說.
敬拜上帝不需要在這個山或那個山.
地域完全不要緊.
只要心靈和誠實就足夠.

$^{121}$耶穌說的話是不是很簡單?.
當然不是.
去這個山或那個山其實是容易的.
要確定自己時時都心靈誠實是很難的.
耶穌對尼哥底姆說.
人若不從生.
就沒有人能夠進上帝的國.
靠言讀和遵行律法可以嗎?.
不可以.
尼哥底姆如果僅僅視耶穌是一個拉比.
就期望從他那裡獲得一套智慧,啟迪.
他將一無所得.
他必須知道耶穌是從上帝那裡來的.
只有他知道上面的事.
尼哥底姆自己是屬地上的人.
永遠無法明白天上的事.
只能夠信任耶穌.
相信他,跟隨他.
然後才能夠進天國.
但是這個相信不可以是建基在尼哥底姆的自然本性得到.
尼哥底姆的生命需要被翻轉.
他要重生.
獲得從上帝來的新生命.
然後才能夠明白屬天的事.
然後才能夠進天國.
每個人都需要生命轉化.
然後才能夠得到拯救.
然後才能夠進天國.
用我們關上的術語.
每個人要跟上帝建立個人的關係.
有新的身份.
由上帝生.
是天父的兒女.
新的生命.
新的價值觀.
新的生活方式.
這是唯一的出路.
宗教儀式沒有用.
外在權威也沒有用.
不過.

$^{161}$話雖如此.
多年來.
基督教卻又取代了傳統的宗教.
繼續扮演外在權威和使行儀式的作用.
讓人可以跟隨教會.
獲得拯救的保證.
天主教的做法.
我們不需要多說.
基本上就是這樣的進路.
信教會.
信教會的權威.
望彌撒.
做各樣的聖禮儀式.
就是我們更正教.
基督教.
在一段很長的時間.
特別是保守的教會.
也是一直扮演著這樣外在權威的角色.
上個月.
10月17日.
說完日子.
全部說的是真話.
我在巴士上.
有一個我不認識的姐妹.
主動坐在我旁邊.
說要跟我聊聊天.
她說她看過我寫的很多書.
也在網上聽過我說的道.
算是我的粉絲.
她一個年輕的姐妹.
大概三四十歲.
她是屬於一間小群教會的.
那間教會其實離我們不遠的.
六十多人.
一間非常內聚的教會.
沒有牧者.
有幾個負責弟兄輪流講道.
教會仍然規定姐妹要蒙頭.
姐妹也不可以承擔教導的工作.
沒有我們教會的姐妹那麼厲害.

$^{201}$不過有趣的說.
那姐妹告訴我.
她的教會沒有受到太多社會運動.
和疫情的衝擊.
領導同工絕口不談政治.
不談社會問題.
弟兄姊妹也很少提問.
曾經有一個年輕弟兄.
不滿教會太內向.
曾經叫過說要走.
不過叫一下而已.
但是他就被其他信徒勸服留下.
因為他們教會那幾十人.
彼此之間的關係是很密切的.
密切的關係可以蓋過一切的分歧衝突.
他說即使在這兩年大移民潮裡面.
他的教會幾乎沒有受到影響.
只是有一個比較外圍的信徒移民.
其他的弟兄姊妹完全安守香港.
他們說上帝沒有指示他們離開.
他們就不會走的.
還有他們的教會.
這些弟兄姊妹很樂意生孩子.
到現在至少已經有兩個了.
至少那個姐妹和她教會的弟兄姊妹.
比我兒子對前景有信心多了.
我兒子死都不肯生.
對世界的前景沒有信心.
和這個姐妹聊完天之後.
我就在想.
在一個動盪多變的世代.
毀身在一個高權威的教會.
未嘗不是快樂的事.
最少信仰實踐會彼得簡單很多.
教會對信徒做了很多訓示規定.
教會設定了很多的規矩.
這些規矩也都成為了一個應許.
只要你跟著規矩而做.
你就會得蒙上帝月立.
你一定得蒙上帝的保守引導.

$^{241}$正不正?.
我在年輕的時候.
回到教會.
雖然沒有這間小群教會那麼具有權威.
不過都是差不多的.
教會為信徒定定了很多Do and Don't.
什麼可以做什麼不可以做.
不准喝酒.
不准看電影.
不准聽流行曲.
這些規矩.
很自然就將信徒和非信徒區分出來.
我在一般的同學中間是格格不入的.
他們玩的東西我都不玩.
那你多奇怪呢.
但是我是很清楚我的基督徒的身份.
盲的都看到了.
是不同的.
是不是?.
這叫做分別為勝.
只要你跟著這樣做.
你就對自己的得救有充分的把握.
沒有懷疑.
不用擔心.
教會還為信徒預備了一套又一套.
非常完備的信仰解說.
這些現成的解說.
就是對生活.
對世界的標準答案.
例如.
禱告必蒙英文.
如果不蒙英文.
就是你祈求不夠懇切.
你說不是啊.
我已經懇切到不得了了.
跪到膝蓋都出血了.
這就是妄求.
你所求的不符合上帝的心意.
遇到困難挫折.
有好結果的就叫做試煉.

$^{281}$沒有好結果的就叫做試探.
試煉來自上帝,試探來自撒旦.
這些解說.
能夠覆蓋生活的多數場景.
消解信徒多數的信仰疑難.
信徒不需要做個別的思考和尋索.
一切都是現成的.
是標準化的.
就像以前背的要你問答的教會神一樣.
不要怕,只要信.
時勢易.
今天的教會,特別是我們王浸.
已經不再扮演宗教權威的角色.
我們也很少為信徒設計一套嚴密的生活規範.
沒有訂定很多規矩.
也很少用高壓的方式.
向你們灌輸各種標準的信仰解說.
從某個角度來說.
你們比上一代的信徒更自由.
沒有太多外在的權威,生活規範.
你們可以看電影,遊行示威.
可以罵牧師,沒有問題的.
事實上,今天很多信徒的教會生活都是非常流動性的.
到處購物.
自己選擇合心意的教會.
這樣也很難建立對教會的長期關係和信任.
不要說權威.
牧者的權威形象也在失落中.
昔日,我說我年輕的時候.
我們對藤牧師,藤建輝牧師,鮑會員牧師.
總是恭恭敬敬的.
他們說的話,我們幾乎不會懷疑的.
今天,我們對牧者很難有這樣的信任.
他們在台上說的話.
for reference only.
僅供參考而已.
聽不聽由我決定.
說句坦白的話.
就算你們很信任我.
當我是權威.

$^{321}$我也不會這樣信任自己.
我不覺得自己是problem solver.
answer bearer.
是答案擁有者.
所以,教會不是權威.
牧者不是權威.
我們再沒有什麼外在的信仰權威.
你說沒有權威.
就靠自己.
我們成年了.
我們是自由人.
不過,自由的代價.
就是要完全自行承擔信仰的責任.
在很多生活場景.
自行做抉擇.
在一個多元化的社會.
沒有一個實踐信仰的固定模式.
沒有基督徒的樣板.
就是說,你不能照抄.
搭順風車.
你知道跟權威和跟風是兩件完全不同的事情.
跟權威.
如果發覺最後跟錯了.
你可以將責任.
罵權威.
你叫我投資什麼.
你叫我做什麼.
結果失敗了.
就罵你.
對不對?.
你怪梁牧師.
怪什麼叻叻.
師父教錯路.
是不是?.
你要罵他們.
但是.
如果你是盲目跟風.
跟潮流.
人做什麼你做什麼.
人去哪裡你去哪裡.

$^{361}$結果你去錯地方.
你埋怨你所跟隨的人.
對方可以說你發神經的.
誰叫你跟我走的?.
所以,不能賴的.
說那些什麼.
凡事都可行.
但不都有益處.
這句話其實是很艱難的.
很大挑戰.
那是什麼?.
有益處還是沒有益處?.
是不是?.
當然.
如果我們的抉擇.
當然是去餐廳.
選乾炒牛河還是星洲炒米.
是無相大雅的.
沒有所謂對與錯的.
選錯了也無所謂的.
是不是?.
不過如果我們在大時代.
常常要面對很多重大的抉擇.
這樣缺乏外在的權威.
缺乏清晰的指引.
就不是一件很愉快的事.
自由可以變成一個重擔.
在見到.
剛才說的那個姐妹的之後一天.
10月18日.
我說的全部都是真事.
我跟太太.
跟女兒.
就去日本玩.
我們在機場.
去一個很差的餐廳.
信用卡那些的貴賓室.
碰到另一對夫婦.
她是我之前服侍的教會的會友.
那間教會超過一千人的.

$^{401}$所以我沒有跟他們聊過天.
不算認識他們.
弟兄看到我.
還沒放下行李.
就撲到我身邊.
那時候我正在吃炒麵.
他跟我說.
他多年前聽過我在教會說的某兩篇道.
跟著他說這兩篇道.
怎樣對他帶來深刻的影響.
讓他在不明白上帝旨意的情況之下.
都不敢質疑上帝.
一邊說.
突然間就很激動.
就哭了起來.
一個中年男人哭.
我馬上關心一下他的情況.
詢問之下.
才知道他此行.
是陪他太太去英國.
被一家移民去英國.
兩個正在讀中學的兒子.
早已經送去英國不同的城市讀書.
他太太本來是個中學老師.
如今辭職去陪兩個兒子.
丈夫在香港還有些生意.
所以暫時會做太空人.
所以送了太太去.
坐一會之後就會回來.
跟我們一家去旅行.
那個心情是完全不同的.
我們帶著快樂的心情去.
他們是非常忐忑.
很多焦慮.
他們突然間遇到我.
就當是上帝給他們一個信心的證據.
弟兄馬上聯想起.
我幾年前說的港島信息.
這個在回憶裡面的信息.
今日竟然發揮作用.

$^{441}$我後來都不怕異相.
我沒試過在一個非常擠迫的地方.
總之坐到鼻子頂鼻子.
但我都抱著他們.
跟他們公開祈禱.
求上帝親自引領他們一家的前路.
這對夫婦突然間記得.
多年前我所說的信息.
我不覺得自己很厲害.
更加不相信自己是先知.
多年前已經預知他們今日的處境.
預先被教導他們.
限定他們不是這樣的人.
我相信人的記憶.
不是對過去客觀而全面的描述.
記憶是一個工具.
是過去對未來的引導.
我們在面對眼前的抉擇.
考驗的時候.
重新召喚那些瘋傳在腦海裡面的片段回憶.
是幫助我們能夠穩住情緒.
確定身份.
肯定方向是一個手段.
特別是當你面對生命嚴重的斷層的時候.
過去未來有斷層,斷裂的危機.
記憶就將過去和未來連在一起.
上帝曾經在很多年前已經提醒我.
應許我.
就執著那個「應許」.
就「bridge over」.
將過去和未來連在一起.
兄弟姐妹,我想說什麼?.
我們確實會在一個非常困難的大時代.
面對很多大小不一的抉擇.
去或留.
要不要在中學或小學階段送走自己的兒子.
要不要夫妻兩地分隔.
要不要推倒自己在香港的工作,專業,事業.
作為一個基督徒.
我們很難不將這些重大的抉擇.

$^{481}$綁上我們的基督信仰.
至少我們總是要問上帝的旨意何在?.
面前是否仍有祂的保守和帶領?.
在一個缺乏外在權威.
外在生活規範的環境裡面.
所有抉擇都純粹是個人的抉擇.
都完全由自己決定.
這真是一件非常不容易的事.
沒有外在的權威.
自己就是自己的權威.
或者你說沒問題.
我們是活躍的一代.
不喜歡聽別人說的.
不過我們總是會聽自己說的.
真相完全不是這樣.
我以前都跟你說過這些analogy(類似).
我們是很難聽自己的話.
特別一些不是出自我們心中真正的喜好.
只是理性上知道必須要做的事.
最爛的例子.
我們每個人都知道要節食減肥.
不過納節減肥又只有我.
減出來又不只有我.
這個爛笑話我講了很多次.
當然如果你的身體出了毛病.
醫生嚴令禁止你不准吃不健康的東西.
我們就會乖乖聽話.
你到我這個年紀就知道.
每一次見醫生都很害怕.
很敬畏的.
醫生是我很敬畏的人.
他說什麼我都聽.
驗完血或者什麼.
看報告那種成王成孔.
比以前去監務處拿成績表還要差.
很害怕的.
我也講過這個例子.
再頂尖的運動員.
都不可以沒有教練.
不可以說我只要自己盡力做好.

$^{521}$不用別人督促.
不可以的.
我跟自己說我已經竭盡所能.
已經跟皮力盡了.
不過有個教練他就跟你說.
做多十次.
信服自己的權威.
其實比信服外在的權威困難.
因為自己很難構成對自己真正的權威.
是這樣.
我們真的可以鎮壓自己心中的恐懼.
沮喪.
絕望.
仇恨.
這些負面的情緒.
我們真的可以攻克己身叫身服我.
有時候可以.
但不是時時都可以.
多數時間都不可以.
我這樣講你.
也講我自己.
為什麼聖經總是強調管教.
因為父母對子女的管教.
上帝對我們的管教.
才能使我們突破成長.
沒有管教就不會成長.
今天的教育生理學也證明這件事.
一個從小就完全放任的小孩.
沒有人將一套清晰的價值灌輸給他.
沒有人引導他.
甚至脅迫他.
去建立良好的生活習慣.
他日後的成長.
特別他要融入社會.
要socialize.
跟人建立關係就會困難重重.
講到這裡.
我們就明白.
基督信仰不可以僅僅是外來的添加.
而必須是生命的徹底改變.

$^{561}$我們必須要重生.
這是耶穌對尼哥底姆說的.
你們必須重生.
為什麼?.
因為我們處於一個急劇轉變.
前途未卜的時代.
我們面對很多影響深遠.
他需要快速決定.
沒有多少時間做考慮的抉擇.
我們處於各種憂患和危機中.
譬如說.
今天不僅僅是想移民的人.
需要嚴肅考慮.
確定你移民的理據.
就是你不想移民.
或者從來沒有想過要移民.
你都需要為自己和其他人.
提出不移民的理由.
幾年前在社會運動如火如荼的時候.
一個住在美國首都DC的家庭.
我認識他爸爸和他太太.
就是一對牧師師母.
他的女兒打電話給我.
是他的師母吩咐他打給我的.
跟我說.
安排好律師和其他各方面的佈局.
然後叫我全家立即遞旨移民去美國.
我對他這樣的性情非常感動.
也知道推卻他的性情.
就顯出我非常不近人情.
所以在推遲他的好意的時候.
我是有一點兒疑心疼的.
就是不知道怎麼說.
很感動也不知道怎麼說.
過去四十年.
不斷有人問我為什麼不移民.
最近三年.
問我的人除了出於關心之外.
也有表達懷疑甚至責備.
事到如今.

$^{601}$能夠走的卻不走.
如果不是心理變態.
那當然有特殊政治任務在身.
你信不信我是中聯辦派來滲透黃浸的人呢.
是不是.
所以不僅僅走的需要理由.
不走的也需要理由.
就好像從前結婚.
總是希望白頭到老的.
是不是.
但今天人人都最少來一次兩次的.
所以為什麼不來.
是需要思考和回答的問題.
我不知道你有沒有靜悄悄想過.
為什麼還沒離婚呢.
眾人以為美的事.
你為什麼不留意去作呢.
今天我們做什麼.
或者不做什麼.
都需要給理由.
轉變需要理由.
首相都需要理由.
移民需要理由.
不移民都需要理由.
離婚需要理由.
不離都需要理由.
並且不是一次過做思考和抉擇就算.
而是每個時候都要重新抉擇一次.
當太爵寺夫人跌倒的時候.
身邊的人問我為什麼不走.
大亞灣核電廠興建的時候.
身邊的人問我為什麼不走.
六四事件發生之後.
人們又問我為什麼不走.
社會運動發生之後.
人們又問我為什麼不走.
每個時刻都被人問.
都要自我詢問.
不抉擇都是一個抉擇.
正如過去身邊總有人跟我說.

$^{641}$你不表態就是一種表態.
不講政治就是政治.
無法選擇無法中立.
我們每個人都要做抉擇.
而抉擇總是反映我們的信仰.
我們的價值觀.
我不是說有信仰就能幫助我們做輕鬆的抉擇.
不過信仰能夠保證我們做最明智的抉擇.
我到最後只是說.
因為我們要做選擇.
所以我們必須要有牢固的信仰.
抉擇凸顯信仰的必須性和重要性.
在一個沒有客觀的對與錯.
沒有標準答案的時代.
我們都需要為自己做各種的決定.
給理由給價值判斷.
這些決定不僅僅關乎我們的現在.
更加影響我們包括我們的家人.
我們的現在我們的未來.
例如為了鋪排下一代的前景.
我們就要犧牲我們這一代的工作.
職業事業甚至專業.
為了關心下一代的前途.
我們就要放棄對上一代的照顧.
遺棄老人是香港非常嚴峻的問題.
資源有限.
我們總是在做一個燒心經.
一個零磨遊戲.
我們要為自己的決定負全責.
因為決策和行動所帶來的種種轉變和嚴重後果.
最終都是我們自己承受.
我只要講這個例子就一定很敏感.
我加拿大有一個很buddy的弟兄.
他媽媽在疫情的時候死了.
他沒有回來送殯.
我要講這個是他一生的傷痛.
那個傷口不知道什麼時候可以買得到.
就是這樣.
沒有客觀的對錯.
但主觀的對錯卻是無法避免的.

$^{681}$我們要預備好迎接各種的疑惑.
焦慮,不安,沮喪,失落的情緒.
所以要面對來自自己和身邊的人的挫折,憤怒,埋怨,內疚.
指責這些張力.
我們要問自己有沒有做對的決定.
今天沒有人能夠準確估算未來.
信仰也不能夠幫我們簡單估算未來.
我們確實是活在無知之中.
也因此我們活在種種的威脅之下.
疾病,不幸,痛苦,欺騙,惡意,背叛,墮落.
我們都深切經驗自身的有限.
我們無法保護自己.
所以都要咬緊牙根去面對這些困難和挑戰.
關鍵是你的信仰是獲得的.
是拿回來的.
還是你自己的.
是隨意借回來的.
還是真實的.
是內在的.
是你說說理論.
還是真實具有指導的能力.
能夠支持你承受各種生活的壓力和抉擇.
指導你向前行.
另外,在資源缺乏,生活壓力沉重.
人際關係變得緊張的日子.
我們其實不容易維持長時間的心理平衡.
亦很難藉著禮貌,教養,理性去控制.
就能夠鎮壓內心潛藏的膽怯,惡意,憤恨和敵意.
我們心中的黑暗面或魔鬼仔.
總是不自主地跳出來.
我們很容易變得心胸狹窄.
自私自利,埋怨別人,拒絕承擔.
我們會有自毀毀人,推倒一切重來的瘋狂念頭.
在風和日麗,太平無事的日子.
講信仰是喝杯茶,吃個包,吹吹水.
反正無傷大雅.
在波濤洶湧的日子.
一個理論,一套哲學.
是無法支撐我們的決策和行動.
我們必須有真正屬於我們自己的信仰.

$^{721}$不是借來抄襲別人的.
而這個信仰一定要連繫我們整個人.
我們必須要參與其中.
而不僅僅是相信或理解某個觀點,某個理論.
簡單說.
要麼你是真正信,徹底信.
要麼信不信也無所謂.
反正後果也無分別.
即是無作用.
半心半意的信.
守乎著泥.
往後看的.
不配盡上帝的國.
信一半是救不了你的.
簡單說.
如果你做了基督徒.
仍然過份體重,名利,富貴,得失成敗.
那你做基督徒實在太痛苦了.
太辛苦了.
信主,我和你坦白說.
明明是會為你帶來時間和金錢上的虧損.
當然會虧損.
除非坊間流行的成功神學是真理而不是歪理.
基督徒特別無上帝保守.
十個基督徒九個有錢.
十個基督徒九個長壽.
一個巴士翻車.
那些福利和健康的教誨就應許.
二百人死剩五個,五個就是基督徒.
不是基督徒的就全部死了.
不過.
如果不是這樣.
那就麻煩了.
如果非基督徒死,基督徒又死.
基督徒除了面對非基督徒一樣的困苦之外.
還要多加一個信仰危機.
質問上帝為什麼不保守我們免遭禍患.
人家只有一個身體問題.
我們還有多個信仰問題.
你這樣也很麻煩.

$^{761}$信耶穌帶給你更多的煩惱.
更麻煩.
正如我以前也舉過這個例子.
我有一個姐妹,媽媽死了.
她一直半年也不停問上帝為什麼要弄死我媽媽.
我說人人都有媽媽.
人人都會死媽媽.
我不知道基督徒死媽媽死得這麼麻煩.
我常常說對一些基督徒來說.
信仰帶給他們不是祝福,是父慮.
我在意大利一個城市港島.
有一個信主幾年的姐妹.
來跟我說她的故事.
她跟丈夫離婚之後.
才有人帶她信耶穌.
不過信耶穌沒有為她消去心中的寂寞和困苦.
她不斷追問身邊的基督徒.
也追問我.
為什麼上帝容許離婚在她生命發生.
為什麼上帝不使她的前夫回心轉意.
為什麼她多番祈禱也不蒙英文.
她不停問這些問題.
一天問,每個星期問.
令到她身邊的基督徒都被麻風地避開.
很麻煩.
每次問這些沒有營養價值的問題.
怎麼回答呢?.
不知道怎麼回答呢?.
終日埋怨.
在怨恨生活之餘.
還多了一個怨恨對象,就是上帝.
你看到嗎?.
他說他信了耶穌.
不過信仰完全沒有改變他的生活世界.
改變他的心思意念.
他的追求,他的盼望.
他只差未問.
上帝為什麼不給我六個號碼.
讓我種到六合彩.
然後我可以奉獻很多錢給教會.

$^{801}$可以藉著我的富貴去榮耀上帝.
我不知道我們中間有沒有人靜悄悄地問這個問題.
上帝為什麼不給我六個號碼.
我最近學習的一個說法是.
辭定永恆.
活在當下.
為了我們對永恆有真正十足的信心.
我們就可以擺脫.
對面前一年兩年十年的猜測和估算.
可以稍稍地.
壓低各種忐忑焦慮的心情.
也比較不受眼前的威脅和威嚇.
不確定所干擾.
專注活好今天.
活好當下.
我知道誰掌管明天.
我知道他握著我的手.
上帝不會成為我推卸自己責任的藉口.
上帝更加不是將我自己的抉擇合理化.
神聖化的一個神聖理由.
上帝卻是幫助我接納我自己的有限.
並且在各種有限無知之中.
仍然能夠比較從容淡定地向前行.
以求 以求 以求.
弟兄姊妹.
你被基督信仰徹底翻轉了沒有?.
你重生了沒有?.
你是否真是一個如假包換.
真真實實的基督徒?.
你的人生觀 你的價值觀 你的世界觀.
有沒有被基督信仰改變?.
舊約聖經有一個很有趣的說法.
叫做救自己.
你試試打這三個字.
很多 超過一百處經文.
或者叫做因他們的義.
救自己的生命.
以西結書十四章十四節二十節.
這些說法.
我就將因他們的義理解為因他們的信.

$^{841}$因信稱義 信就是義.
這亦是新約聖經裡面耶穌說的.
你的信救了你.
我們的信救了我們自己.
我們的信仰能否救我們自己.
進而救我們全家.
救我們身邊的人.
提摩太前書四章十六節.
弟兄姊妹.
今日如果我說的話少少沉重.
請你們原諒.
這確實是我這一兩年.
營養在心裡面最重擔的關懷.
我們的信仰是否真的.
我們的信仰是否看得見.
我希望我不在你們中間唱任何高調.
我和我的家庭都會和你們一起.
我們去面對香港種種的變化.
承擔旺朕面前的日子.
我們一起看看上帝怎樣做.
\newpage



\section{以賽亞書 36:1-12}
\label{sec:4NVCY6ZFpaA}
\textbf{2023年12月3日崇拜直播｜陳冠成牧師｜唯獨你是王 - 神必拯救我們｜以賽亞書36\_1-12}
\newline
\newline
連結: \href{https://youtube.com/watch?v=4NVCY6ZFpaA}{\texttt{https://youtube.com/watch?v=4NVCY6ZFpaA}} ~~~~ 語音日期: 2023-12-04
\newline
\newline
\hyperref[sec:_oDlWIvKGkE]{\small{< < < PREV SERMON < < <}}
~
\hyperref[sec:index_chronic]{\small{[返順時目]}}
~
\hyperref[sec:LudzCxftFrQ]{\small{> > > NEXT SERMON > > >}}
\newline
\newline
以賽亞書 36:1-12
\newline
\begin{longtable}{cl}
\hline
\hline
章節 & 經文 (和合本修訂版)\\
\hline
36:1 & \begin{tabularx}{0.7\textwidth}{X} 希西家王十四年，亞述王西拿基立上來攻擊猶大的一切堅固的城，將城攻取。 \end{tabularx} \\ \\ \relax
36:2 & \begin{tabularx}{0.7\textwidth}{X} 亞述王從拉吉差遣將軍率領大軍前往耶路撒冷，到希西家王那裡去。將軍站在上池的水溝旁，在往漂布地的大路上。 \end{tabularx} \\ \\ \relax
36:3 & \begin{tabularx}{0.7\textwidth}{X} 希勒家的兒子以利亞敬宮廷總管、舍伯那書記和亞薩的兒子約亞史官，出來見他。 \end{tabularx} \\ \\ \relax
36:4 & \begin{tabularx}{0.7\textwidth}{X} 將軍對他們說：「你們去告訴希西家，大王亞述王如此說：『你倚賴甚麼，讓你如此自信滿滿？ \end{tabularx} \\ \\ \relax
36:5 & \begin{tabularx}{0.7\textwidth}{X} 我說，你有打仗的計謀和能力，我看不過是空話。你到底倚靠誰，竟敢背叛我呢？ \end{tabularx} \\ \\ \relax
36:6 & \begin{tabularx}{0.7\textwidth}{X} 看哪，你所倚靠的埃及是那斷裂的葦杖，人若倚靠這杖，它就刺進他的手，穿透它。埃及王法老向所有倚靠他的人都是這樣。 \end{tabularx} \\ \\ \relax
36:7 & \begin{tabularx}{0.7\textwidth}{X} 你若對我說：我們倚靠耶和華－我們的神，希西家豈不是將神的丘壇和祭壇廢去，且吩咐猶大和耶路撒冷的人說：你們只當在這一個壇前敬拜嗎？ \end{tabularx} \\ \\ \relax
36:8 & \begin{tabularx}{0.7\textwidth}{X} 現在你與我主亞述王打賭，我給你兩千匹馬，看你能否派得出騎士來騎牠們。 \end{tabularx} \\ \\ \relax
36:9 & \begin{tabularx}{0.7\textwidth}{X} 若不然，怎能使我主臣僕中最小的一個軍官轉臉而逃呢？你難道要倚靠埃及的戰車和騎兵嗎？ \end{tabularx} \\ \\ \relax
36:10 & \begin{tabularx}{0.7\textwidth}{X} 現在我上來攻擊毀滅這地，豈不是出於耶和華嗎？耶和華吩咐我說，你上去攻擊這地，毀滅它吧！』」 \end{tabularx} \\ \\ \relax
36:11 & \begin{tabularx}{0.7\textwidth}{X} 以利亞敬、舍伯那、約亞對將軍說：「求你用亞蘭話對僕人說，因為我們聽得懂；不要用猶大話對我們說，免得傳到城牆上百姓的耳中。」 \end{tabularx} \\ \\ \relax
36:12 & \begin{tabularx}{0.7\textwidth}{X} 將軍說：「我主差遣我來，豈是單對你和你的主人說這些話嗎？不也是對這些坐在城牆上，要與你們一同吃自己糞、喝自己尿的人說的嗎？」 \end{tabularx} \\ \\
[1ex]
\hline
\hline
\end{longtable}
$^{1}$今日的經訓是記載在二賽亞書36章1至12節.
在大家手上的聖經舊約第957至958頁.
由你一同聆聽神的話語.
希西加王十四年.
亞述王西拿基納上來攻擊猶大的一切堅固性.
將成功取亞述王從拉吉差遣拉伯沙記率領大軍往耶路撒冷.
到希西加王那裡去.
他就站在上池的水溝旁.
在漂泊地的大路上.
於是希臘家的兒子迦載爾利亞徑.
並書記舍伯拿和亞薩的兒子史官約亞.
出來見拉伯沙記對他們說.
你們去告訴希西加說.
亞述大王如此說.
你所倚靠的有什麼可仗賴的呢.
你說有打仗的計謀和能力.
我看不過是虛話.
你到底倚靠誰才背叛我呢.
看啊 你所倚靠的埃及是那壓箱的偉將.
人若靠這將就必刺透他的手.
埃及王法老向一切倚靠他的人也是這樣.
你若對我說.
我們倚靠耶和華我們的神.
希西加豈不是將神的幽壇和祭壇廢去.
且對猶大和耶路撒冷的人說.
你們當在這壇前敬拜嗎.
現在你們當頭給我主亞述王.
我給你二千匹馬.
看你這一面騎馬的人夠不夠.
若不然怎能打敗我主神僕中最小的軍長呢.
你竟倚靠埃及的戰車馬兵嗎.
現在我上來攻擊毀滅這地.
豈沒有耶和華的意思嗎.
耶和華吩咐我說.
你上去攻擊毀滅這地吧.
以利亞經舍伯拿約亞對拉伯沙基說.
求你用亞蘭言語和僕人說話.
因為我們懂得.
不要用猶大言語和我們說話.
達到城上百姓的耳中.

$^{41}$拉伯沙基說.
我主差遣我來.
豈是單對你和你的主說這些話嗎.
不也是對這些坐在城上.
要與你們一同吃自己的糞.
學自己撩的人說嗎.
請姐妹平安.
繼續「唯獨你是王」宣講系列.
上一次我和大家分享過.
其實也是以上的書.
第七章.
如果大家還記得是猶大王希西嘉.
當他面對鄰國亞蘭和以法聯的聯合攻擊時.
他寧願依靠另一個更強大的勢力.
亞述帝國.
通過外交手段化解這場危機.
而不肯聽以賽亞先知的勸告.
回轉歸向上帝.
相信上帝的帶領.
於是大家還記得上一次我們說過.
上帝也有些激氣.
他親自對阿哈斯和猶大.
發出以馬來尼的兆頭.
表示阿哈斯現在倚靠亞述.
但亞述將來會成為一個更強悍的敵人.
並且對猶大作出更猛烈的攻擊.
而這個兆頭.
在以賽亞書三十六章.
即是我們剛才所讀的經文.
在阿哈斯的兒子.
希西嘉所面對的困境.
應驗了在他身上.
事實上如果比對.
以賽亞書第七章和第三十六章.
你會發覺這兩父子.
其實有很多相似的地方.
在他們的經歷上.
聖經學者也留意到.
原來他們同樣面對大軍的壓境.
事件的關鍵地點.

$^{81}$同樣一模一樣.
都是記載在上池的水溝旁.
在漂泊地的大路上.
一模一樣.
這兩位皇帝也試過.
用軍事,外交,聯盟的方式.
來化解危機.
當年的阿哈斯皇帝.
就是聯合亞述.
現在其實.
希西嘉也曾經試過.
聯合埃及來對抗亞述.
而這兩位皇帝也得到.
先知以賽亞的提醒.
勸他們不要倚靠人的力量.
而反而要倚靠上帝的幫助.
結果.
我們還記得上一次.
阿哈斯選擇了什麼?.
拒絕上帝.
上帝說什麼都沒有用.
我心意已決.
我要靠自己來拯救自己.
但時隔三十多年.
到今天的希西嘉.
他會不會重覆他父親的覆轍呢?.
還是他有不一樣的選擇?.
就是我們接下來要看的經文.
經文說.
希西嘉當時已經在位第十四年.
亞述的大軍向猶大發出攻擊.
經文亦記載.
當時大部分猶大的城邑.
原來已經落在亞述王的手上.
只剩下兩個.
就是耶路撒冷和拉吉這兩個城市.
大家看到程序表.
廣東大廣東也印了圖.
或者可以給出圖.
箭咀指著的兩個地方就是.

$^{121}$分別是藍色箭咀是耶路撒冷.
紅色箭咀是拉吉.
大家看到.
這兩個城市可以互相倚靠.
可以像一個倚閣一樣.
互相守望.
事實上這兩個城市未倒.
因為歷史記載他們有堅固的城牆.
是屬於易守難攻的城市.
換句話說.
只要拉吉守得住.
就會對進攻耶路撒冷的亞述軍構成威脅.
因為怕他們夾擊.
令他們腹背受敵.
反過來.
如果耶路撒冷站得穩.
也會成為現在正在圍攻拉吉的亞述大軍.
成為他們的一種隱患.
再加上當時還有另一個勢力.
埃及.
所以當時的亞述王要保留軍力.
以防埃及大軍突襲.
所以經文記載.
阿述王想到一個策略.
先包圍拉吉這個城市.
然後同時向耶路撒冷施壓.
迫他們投降.
我們說企圖不用打就贏.
不戰而勝.
於是他派出其中一位軍長拉伯沙基.
率領軍隊上到耶路撒冷.
準備向城內的人打一場心理戰.
大家可以看經文.
印在港島大綱裡.
我們先看第四至十節.
我做了一些筆記.
圈出很多與倚靠相關的詞語.
看到有很多嗎?.
用一個「呱呱」.
這是拉伯沙基打心理戰的主要內容.

$^{161}$重複提到倚靠.
他一開宗明義就說.
希西加所倚靠的.
根本不算什麼倚靠.
然後他就用這個論點.
將希西加或猶大所倚靠的對象.
逐一擊破.
第五節他首先提到.
希西加本人的實力.
他說有打仗的能力和謀略.
我看你不過是吹水.
事實證明大部份城市.
已經落在亞述王手中.
你還說自己有能力和謀略.
你到底倚靠誰來支持你.
去背叛亞述王和他作對.
你看你倚靠的埃及現在怎樣.
他只不過是一支壓傷了的蘆葦.
即是你一壓他就會怎樣.
其實會「啪」一聲就斷.
斷出來的位置還會刺傷你.
根本撐不住你.
拉伯沙基事實上也沒有吹水.
因為較早之前亞述大軍.
確實輕而易舉擊退了法路的襲擊.
嘲笑完希西加外交策略方面的失敗後.
拉伯沙基又開始針對另一件事.
就是希西加最弱的一環.
是他的軍力.
第八至九節.
這裡基本上好像挑機一樣.
希西加你可以和我的主人亞述王打賭.
我送你二千匹馬.
當然是戰馬.
能夠有戰鬥能力.
亦即是當時最先進的軍事武器.
我送你一個最先進的軍事武器.
看你能否湊夠人來騎戰馬.
換句話來說.
其實他在考問你們有沒有馬兵.

$^{201}$你們有沒有訓練人來騎戰馬.
我看死你沒有.
我送你二千匹戰馬也沒有用.
你也沒有人懂得騎.
所以你竟然夠膽依靠埃及來軍事援助你.
就算埃及為你提供戰車戰馬.
也沒有用.
你根本不可能打贏亞述王.
就連拉伯沙基一個小小的軍長.
希西加你也鬥不過.
他看死猶大這種弱小的國家.
根本沒有能力.
沒有軍力沒有武器.
來跟當時最強大的亞述帝國對抗.
不單止這樣.
心理戰還未打完.
最後他要攻擊猶大的信仰.
他試圖打擊希西加和百姓.
對上帝對耶和華的信靠.
經文繼續說.
你說你們依靠耶和華.
但你看看.
你連敬拜耶和華的優壇和祭壇.
不是全部被希西加拆走嗎?.
你還對百姓說.
不准在其他地方敬拜.
只可以在耶路撒冷敬拜.
第十節他再補刀.
你看看現在.
上帝派我們來攻擊耶路撒冷.
不是沒有上帝的心意嗎?.
上帝吩咐我們.
現在來攻擊毀滅這個地方.
聽起來有點亂.
馬利會發覺拉伯沙基.
很明顯我們知道.
當時打仗.
會有所謂的奸細.
和現在一樣.
很明顯亞述王派了臥底.

$^{241}$深入猶大的王宮.
來搜集他們的一舉一動.
他們似乎知道.
希西加曾經進行過宗教改革.
他們又知道.
先知亦曾預告過.
耶和華要把不聽話的猶大.
交給亞述.
來教訓他們.
於是拉伯沙基.
就把這些收集的情報.
混在一起.
說到希西加的宗教改革.
拆除了敬拜上帝的優譚.
現在就惹上帝生氣.
所以上帝派我們來攻擊.
來教訓你.
惹上帝生氣.
上帝不會再幫你.
也不會再成為猶大的倚靠.
你看到很厲害吧?.
是不是事實不要緊.
你看到某程度上.
拉伯沙基是.
委曲了.
曲解了事實.
扭曲了上帝的心意.
但重點是.
他成功製造了一番恐嚇的言論.
他已經成功削弱了.
猶大耶路撒冷的士氣.
影響他們對上帝的信靠.
你看到在場官員的反應.
你就知道這一點.
在場的三位官員.
他們知道拉伯沙基的意圖.
於是就怎樣呢?.
形勢敵人確實比自己強.
自己唯有低聲下氣.
來乞求拉伯沙基.

$^{281}$你改用亞蘭語.
即當時國際的言語.
你改用亞蘭語.
來跟我們對話.
不要再說猶大話了.
免得我們在城牆上的士兵聽到.
影響他們的軍心.
甚至將這番恐嚇的言語.
傳入城中百姓的耳中.
拉伯沙基知道自己的心理戰就臭了.
於是他就更加有恃無恐.
留意他接著的對話是甚麼.
我就是要耶路撒冷所有人.
都聽到我這番言論.
我就是要製造全城恐慌.
讓所有人對希西加.
甚至對耶和華都不信任.
覺得如果我們繼續.
依靠耶和華.
依靠希西加的話.
就一定不會有好下場.
這才是拉伯沙基真正的意圖.
正所謂我們常說.
攻心為上.
攻城其實是為下.
而為了徹底去粉碎.
猶大對抗亞述的念頭.
並且讓他們相信投降才是唯一出路.
拉伯沙基繼續說.
經文記載他還要大聲地.
用猶大話向耶路撒冷喊話.
我們再繼續看.
剛才未讀到的經文.
第十三至二十節.
拉伯沙基這次說的不是倚靠.
圈在經文中.
最主要是甚麼?.
拯救.
這次他用另一個說詞.
他強調拯救.

$^{321}$誰才能救你?.
你沒有東西可靠的.
再問一句.
誰才能救你?.
一方面他施壓.
對百姓說.
你們不要再被希西加欺騙.
誤導你們.
以為耶和華真的可以拯救你.
脫離埃及王的手中.
你看看周圍的國家.
有哪個國家的神.
能夠救他們自己的國家?.
能夠幫他們脫離亞述王的手?.
你看看.
一個都沒有.
他真的用數據.
用數據來支持他的論點.
周圍的國家都被亞述剷平了.
試問.
耶和華難道有例外?.
耶和華可以救你們脫離亞述王的手?.
拉伯沙基嘲笑.
你們所信的耶和華.
其實和列國的神一樣.
都不是亞述王的敵手.
上帝是沒有能力拯救猶大.
但同一時間.
他又伸出橄欖枝.
第十六到第十七節.
他又向這些百姓示好.
表示現在誰能救你們?.
只有亞述王一個.
是他才能救你們.
只要你們願意投降.
和亞述王友好.
那就會有安樂茶飯吃.
那就會沒事.
我們說其實.
拉伯沙基很明白.

$^{361}$打心理戰最重要的是兩件事.
什麼?.
威迫和利誘.
在我們剛才所讀的經文裡有齊.
他不斷製造壓力.
給所有城內的百姓.
造成一個假象.
以為這次真的會完蛋.
沒有出路.
當這些百姓面對高壓時.
覺得死路一條時.
他會怎樣?.
他突然給你一點點甜頭.
突然向你釋出善意.
告訴你.
有得救的.
令那些高壓的百姓.
突然有放鬆的感覺.
突然有釋懷的感覺.
在這時候就最容易動搖.
對上帝的信靠.
只要投靠亞述王.
似乎什麼事都沒有.
可以危機解決.
不過我們都很明白.
如果真的投靠亞述王.
其實就真的是噩夢的開始.
弟兄姊妹我們停一停.
其實你發不發覺.
拉伯沙基對耶路撒冷採取的心理戰手法.
和今天這個世界引誘我們基督徒.
不信靠上帝的手法其實都一樣.
都是威逼和利誘.
首先就是連珠炮發.
從不同的角度.
可能有人提醒你.
可能你在網上瀏覽很多資訊.
或者你現在身處的環境.
告訴你.
你這樣又不行.

$^{401}$你那樣又不行.
你看看你的工作.
好像不太好.
你的健康不用說了.
財富又不太好.
造成很多危機.
很多似是而非的訊息.
大量湧入.
然後就問你.
你還信得上帝?.
你還相信上帝可以幫到你?.
不是吧?.
首先動搖你對上帝的信靠.
覺得上帝在你身處一些環境.
都幫不到你.
然後就怎樣?.
就好像拉伯沙基那樣.
再給你一些甜頭.
很多時候都會這樣.
你會見到拉伯沙基一樣.
都是用一些真實的.
你現在多慘?.
現在上帝幫不到你.
你真的要繼續相信他?.
很多時候我們生活一樣.
會面對這個情況.
我們是否真的會動搖.
對上帝的信心?.
不知道大家還記不記得.
很久以前.
有一套電影.
湯姆漢斯影帝主演.
叫做《劫後重生》.
有印象嗎?.
《Cast Away》.
講述他是一個速遞員.
有一天要執勤.
坐一架飛機.
很不幸的.
飛機遇難墮海.

$^{441}$所有人都死了.
只剩下他一個.
漂流在荒島上.
最後他經過自己的努力求生.
最終脫險.
回去找回他的家人.
當時我們很小.
我們信主初初.
一班年輕人看完之後.
大家會討論在周會中.
問:如果你是湯姆漢斯.
你會怎樣做?.
信仰怎樣幫到你?.
年輕人會比較真實.
講一些天真一點的說話.
我當然會像他這樣.
看看有沒有水喝.
看看有沒有機會求救.
如果沒有.
算了.
不搞了.
沒有事可搞了.
有些人會問:你會不會祈禱?.
看看環境如何.
如果明知自己死定了.
為何要祈禱?.
不會的.
其實很實際.
當真實環境來到的時候.
當時問我們.
是否真的相信上帝的時候.
答案才是真的.
不是平日無風無浪的時候.
問我們的標準答案.
很多時候我們在環境的威嚇下.
在生命受威脅的時候.
在健康出問題的時候.
在財富有壓力的時候.
其實世界就在跟你說.
上帝可信嗎?.

$^{481}$你真的盡靠上帝?.
拉伯沙基對上帝的指紋是這樣.
這個世界對我們也是這樣.
目的就是要動搖我們對上帝的信靠.
另一方面.
又會向我們釋出善意.
這個世界.
其實除了上帝之外.
還有很多方法.
讓你生活過得好一點.
是嗎?.
很多方法.
何必勉強自己一定要跟隨上帝.
為什麼不試試轉投其他陣營.
依靠其他力量.
說不定有效呢?.
你還記不記得.
魔鬼撒旦也是這樣引誘人犯罪.
說不定.
這三個字很厲害.
他沒有叫你完全否定上帝.
他只是說.
說不定除了上帝以外.
還有其他拯救呢?.
還有其他方法解決你的困難呢?.
你試試模仿.
為什麼還要花這麼多時間投入教會.
追求信仰.
相信上帝.
弟兄姊妹這個世界.
其實真的會像拉伯沙基那樣.
用盡方法將我們拉走.
不再相信上帝是我們唯一的保障和拯救.
到底希西加.
他怎樣帶領白順去回應呢?.
我們繼續看.
我想大家一起讀一段經文.
我們揭開程序表旁邊的版本.
有兩段經文.
分別是三十六章的二十一至二十二節.

$^{521}$和三十七章的一至四節.
請大家幫我一起讀.
預備起.
(文字).
都撕裂衣裳.
來到希西加那裡.
將拉伯沙基的話告訴了他.
希西加王聽見.
就撕裂衣服.
披上麻布.
進了耶和華的殿.
使加載以利亞徑和書記舍伯蘭.
並祭司中的長老.
都披上麻布.
去見阿摩斯的兒子.
先知爾賽亞.
對他說.
希西加如此說.
今日是急難.
責罰.
凌辱的日子.
就如婦人將要生產頸鞋.
卻沒有力量生產.
或者耶和華.
你的神聽見拉伯沙基的話.
就是他主人亞述王.
打發他來辱罵永生神的話.
耶和華你的神聽見這話就發斥責.
故此.
謂謂為愚臣的民揚聲禱告.
好,謝謝.
你看到.
面對我們說的國難當前.
希西加他不像社會領袖那樣.
我們立即召開緊急戰略會議.
我們和一班官員商量.
看看如何反制對手.
反而你會見到.
他們兩個的行動.
顯出他們那種屬靈的智慧.

$^{561}$因為他正確判斷到.
擺在面前的.
不是國與國家的打仗.
是屬靈爭戰.
又或者.
是耶和華和亞述王的角力.
你留意他說什麼.
今天是急難責罰凌辱的日子.
他很清楚.
不只是對手凌辱我們.
裡面是上帝的管教.
很久以前上帝已經說了.
因為我們不聽上帝的吩咐.
是要靠攏亞述.
現在上帝就透過亞述來責罰我們.
但他的屬靈大局觀掌握得很清楚.
所以拉伯沙基問他.
耶和華真的能夠拯救耶路撒冷脫離亞述王的手嗎?.
其實沒有人能夠回答.
除了上帝之外.
沒有人能夠回答.
危機能否解除.
只有上帝才有答案.
就正如希西加他形容猶大的情況.
好像什麼.
好像婦人將要生產孩子.
卻沒有力量生產.
處於生死邊緣.
你面對生死邊緣.
但掌握生死的不是你.
因為你沒有能力.
我還記得我兒子出生前.
原本打算順產.
一直等.
陪著太太.
等啊等啊.
不對勁.
指數一跌.
醫院響機.
即是一堆人衝過來.

$^{601}$將我推開.
馬上開刀.
我連見面都不能見面.
被排拒於外.
就像現在這樣.
婦人將要生產.
但她沒有能力掌握這個過程.
她需要外力.
我需要上帝.
沒有人能夠幫助希西加.
沒有人能夠幫助猶大.
除了上帝.
他知道上帝才有能力起死回生.
因為他明白左右這場局勢.
其實不是亞瑟王.
是上帝.
所以你看到.
希西加做了幾個正確的判斷.
很值得我們學習.
首先他吩咐所有人.
什麼都不要說.
聽到什麼都不要回應.
無論敵人怎樣威嚇你.
無論敵人怎樣說什麼.
你一定要拒絕回應.
不要為敵人的挑釁.
再提供養分.
不要給助燃劑.
希西加他知道.
猶大最需要的不是.
如何與別人對嘴.
如何反駁別人的論據.
如何反擊敵人.
他們最需要的是.
專心回到上帝身上.
專心來到上帝面前.
尋求上帝的同在和指引.
所以你看到.
他有兩個動作.
他首先回到聖殿.

$^{641}$他要的是上帝的同在.
第二,他又派人去找先知爾賽亞.
他要的是上帝的指引.
他不會輕舉妄動.
這兩個態度很值得我們學習.
特別當我們面對別人的挑釁.
無論對你本人或對你的信仰.
有時最好的回應是.
不要立即回應.
先安靜.
特別當你很想以牙還牙.
很想罵對方.
與對方針鋒相對的時候.
這時候你選擇安靜.
仰望上帝反而更加有智慧.
只要我們拒絕為對方的挑釁提供養分.
其實就可以避免整件事失控越演越烈.
甚至明明動口變成動手.
所以與其在這些口試之爭裡找輸贏.
倒不如回到那位得勝的上帝身上.
尋求他的心意.
將這一切看為一場屬靈的操練.
事實上上帝正正在我們身邊.
有沒有真的榮耀他.
有沒有真的造就到人.
所以你看到.
希西加為何可以這麼冷靜.
大軍壓境死到臨頭.
為何他可以克服亞述帝國.
帶給他那種恐懼的威脅.
我們說是因為他有一個敬畏上帝的心.
他知道真正主宰的是上帝.
不是亞述王.
亞述王再惡再厲害.
都只不過是上帝手中的一支棍.
那支管教的杖.
你要怕那支棍還是拿著那支棍.
弟兄姊妹.
你要怕這個世界給你那種威嚇.
給你那種逆境.

$^{681}$你要懼怕這些多些.
還是敬畏那一位掌管逆境的上帝多些.
你應該懂得選擇.
正如希西加選對了一樣.
當我們懂得仰望敬畏上帝的時候.
其實你會否發覺.
做人不需要看人面色.
看人面色做人很辛苦.
我們都明白.
我以前曾經有個老師說得很好.
他說你出門口每天都要看天色.
看有沒有擔傘.
是否需要擔傘.
有沒有看上帝面色.
有沒有先看上帝面色去做人.
以致我們懂得在生活裡面.
做正確的選擇.
希西加他做了.
經文亦都期望我們先看上帝的面色.
定一定再想想.
上帝想我們怎樣回應.
而不是自己輕舉妄動.
即使有多大的困難.
我們要確信上帝依然可靠.
祂沒有倒下.
所以我們真的不需要獨自去面對.
人生很多困難.
今日最大的問題不是我們遇到困難.
遇到逆境.
是我們自己想一遍去面對.
是我們不願意和別人分享.
所以這段經文說得很好.
其實回教會真的有用.
找牧者傾真的有用.
希西加就是這樣.
他第一時間回聖殿就是上教會.
他去找仙子求問就是找牧者.
其實真的很有用.
當你遇到困難的時候.
你不要自己一個去想怎樣解決.

$^{721}$你找人來分享.
上帝會在那裡工作.
你回到神的家.
你找牧者傾.
找弟兄姊妹守望你.
一起去尋求上帝的幫助和指引.
相差很遠.
整個人會海闊天空.
因為我們確信上帝有出路.
上帝是我們唯一的拯救.
最後看最後的段落.
你見到上帝安慰希西加.
叫他幾件事很簡單.
不用怕.
亦不要回應亞述王對你的威脅.
甚麼都不要做.
專心看上帝出手.
上帝說我會為自己的名.
因為他辱罵我.
我會為自己的名出手對付他.
親自介入這個事件.
我不單止會讓亞述的大軍退兵.
我亦會讓亞述王死在刀下.
上帝說完了.
不難明白吧.
最難是甚麼.
最難是上帝沒有告訴你.
他何時出手.
可以是明天.
可以是一個星期.
甚至更久的時間.
總之上帝說你相信我.
我會出手.
你甚麼都不要做.
看我怎樣做.
你看到上帝要考驗希西加.
和百姓對他那份絕對的信心.
是否真的願意耐心等候上帝的工作.
弟兄姊妹.
你自己覺得自己是否一個有耐性的人.

$^{761}$做個簡單測試.
大家都有過馬路紅綠燈.
有試過按了紅綠燈按鈕.
按了一分鐘不轉.
兩分鐘不轉.
三分鐘不轉.
有試過嗎.
我有試過在曾容大廈對出的馬路.
間中會這樣.
這時你會怎樣做.
再按吧.
再按都沒有反應.
通常只有兩個答案.
衝呀.
沒有車就衝過去.
如果你不想違法.
你就找下一盞燈.
下一盞燈應該沒有壞.
下一盞紅綠燈.
你再讓一盞過就沒事.
我認識有個神學院老師.
可能因為教舊約的緣故.
他很守規矩.
一定不會衝燈.
如果那盞燈壞了.
就混路走.
走安全道.
怎樣都不犯法.
我們應該不是這樣.
我們是新約下的信徒.
我想我們沒有耐性.
去等它轉燈.
大部份人都會說服自己.
一定是那盞燈壞了.
是那盞燈有問題.
還是那盞燈有問題.
我趕時間.
沒辦法.
我要靠自己解決.
然後用自己的方法.

$^{801}$去過這條馬路.
不會有耐性繼續等下去.
我想我們都是這樣.
如果有這些情況.
待會你可以試試.
但我們對上帝的耐性.
是否都是這樣.
會否我們經常催迫上帝.
轉燈.
你還是紅燈.
那盞燈壞了嗎.
你還不轉綠燈給我過.
我等了很久.
我要自己衝過去.
我不等.
死人了.
等於是一個.
我們說你沒辦法不快的年代.
慢好像是一個罪.
慢好像被人鄙視.
我們覺得很有同感.
就是只等一件事.
那種感覺是很難受.
只等還要你什麼都不做.
你會覺得時間明明是一個小時.
你發覺好像等了一天.
我們不習慣等.
等候是我們最弱的一環.
基督徒也不例外.
所以上帝真的很好.
他會用他的方法.
你可能會覺得不舒服.
但他會逼你等.
你不等也沒用.
因為你沒辦法.
你唯一是依靠他.
他會用盡方法去調教.
適合我們的生命步伐.
來去馴服他的大靈.
所以有些時候他會幫你.

$^{841}$他會拯救你.
不過是按他的時間表.
不是按我們的意思.
上帝要熬練我們.
要我們學會.
什麼是信心等候的功課.
所以假如你真的.
好像現在經文所說的逆境.
當然我們真的很想盡快.
問題消失.
一切解決.
雨過天青.
但如果上帝的心意.
真的要我們多等一會.
多等一會.
這個困難會過的.
最終一定會解決.
並且我會與你同在.
我會一直在你身邊.
提供足夠的恩典和保護給你.
所以,靈姐妹.
我們願不願意等候上帝?.
我們願不願意.
上帝的旨意在我們身上成就.
我們同心祈禱.
天父亞正如.
我們踏入十二月期盼.
你親自的來到.
我們期望耶穌親自掌管這個世界.
掌管我們的生命.
上帝,我們確實知道你已經來了.
你在我們當中完成拯救.
你也應許凡信靠你的.
必然得救.
必然進入永生.
所以我們大膽求你.
幫助我們成就你的旨意.
讓我們即使處於困難的裡面.
我們知道上帝你沒有走開.
上帝你仍然掌管一切.

$^{881}$你的拯救功效仍然存在.
你仍然是那位絕對值得信靠的上帝.
求你幫助我們.
不要讓這個環境威嚇我們.
不要讓這個環境給我們的際遇.
給我們的困難,動搖,對上帝的倚靠.
願意我們真的面對這一切的時候.
我們首先安靜回到你那裡.
我們首先安靜回到神的家.
尋求弟兄姊妹的幫助和守望.
我們更加住在上帝的裡面.
尋求聖靈的安慰和指引.
以示我們的人生無論或順或逆.
上帝是你拯救的大能.
都在我們身上彰顯.
為此將我們每一位交託給你.
願上帝你與我們同在.
願你聽我們禱告.
奉基督耶穌你得勝的名字來祈求.
阿們.
\newpage



\section{路加福音 9:51-56}
\label{sec:LudzCxftFrQ}
\textbf{2023年12月17日崇拜直播｜梁家麟牧師｜耶穌與撒馬利亞人｜路加福音9\_51-56}
\newline
\newline
連結: \href{https://youtube.com/watch?v=LudzCxftFrQ}{\texttt{https://youtube.com/watch?v=LudzCxftFrQ}} ~~~~ 語音日期: 2023-12-19
\newline
\newline
\hyperref[sec:4NVCY6ZFpaA]{\small{< < < PREV SERMON < < <}}
~
\hyperref[sec:index_chronic]{\small{[返順時目]}}
~
\hyperref[sec:QUWqIbTXwuo]{\small{> > > NEXT SERMON > > >}}
\newline
\newline
路加福音 9:51-56
\newline
\begin{longtable}{cl}
\hline
\hline
章節 & 經文 (和合本修訂版)\\
\hline
9:51 & \begin{tabularx}{0.7\textwidth}{X} 耶穌被接上升的日子將到，他決定面向耶路撒冷走去。 \end{tabularx} \\ \\ \relax
9:52 & \begin{tabularx}{0.7\textwidth}{X} 他打發使者在他前頭走；他們進了撒瑪利亞的一個村莊，要為他作準備。 \end{tabularx} \\ \\ \relax
9:53 & \begin{tabularx}{0.7\textwidth}{X} 那裡的人不接待他，因為他面向著耶路撒冷去。 \end{tabularx} \\ \\ \relax
9:54 & \begin{tabularx}{0.7\textwidth}{X} 他的門徒雅各和約翰看見了，就說：「主啊！你要我們吩咐火從天上降下來，燒滅他們嗎？」 \end{tabularx} \\ \\ \relax
9:55 & \begin{tabularx}{0.7\textwidth}{X} 耶穌轉身責備兩個門徒。 \end{tabularx} \\ \\ \relax
9:56 & \begin{tabularx}{0.7\textwidth}{X} 於是他們就往別的村莊去了。 \end{tabularx} \\ \\
[1ex]
\hline
\hline
\end{longtable}
$^{1}$現在跟大家讀的經文是在新約聖經95頁.
路加福音9章51-56節.
聖經這樣說.
耶穌被接上升的日子將到.
他就定義向耶路撒冷去.
便打發使者在他前頭走.
他們到了撒瑪利亞的一個村莊.
要為他預備.
那裡的人不接待他.
因他面向耶路撒冷去.
他的門徒雅各·約翰看見了.
就說:主啊,你要我們.
吩咐火從天上降下來消滅他們.
像以利亞所作的魔.
耶穌轉身集被兩個門徒說:.
你們的心如何,你們並不知道.
人子來,不是要滅人的性命.
是要救人的性命.
說著就往別的村莊去了.
多謝梅主神.
我沒有離開約翰福音的書卷.
不過本來今天應該到約翰福音第四章.
講到耶穌在聚家井旁向婦人談道.
講到這裡不如我扯遠一點.
跟大家講一個更闊的題目.
講一下耶穌和撒瑪利亞人的關係.
這是一個背景.
講完之後我們下一次再講約翰福音第四章.
猶太民族是一個充滿仇恨和紛爭的群體.
無論在歷史還是現實.
幾乎沒有停止過有迫害,殺戮,戰爭,相伴隨.
最近牽動我們的心弦的是以巴衝突.
戰火繚搖,死了很多很多的人.
我去聖地超過十次.
每次進入耶路撒冷這個名為平安和平之城.
我都很深刻感受到那裡的居民.
是怎樣長期活在敵視,猜忌和恐懼之中.
最近幾年本來是好了.
不過我預計明年或後年.
戰爭止息有機會再去.

$^{41}$我都可以預計那種仇恨感會重現.
最重要的時候大概二十年前.
我記得去入救城.
每一個角落都有以色列幾個軍人駐紮.
用極其敵視的眼光去看著每一個來往的人.
而巴勒斯坦人那種冷漠,仇恨的眼光.
總之你去到就覺得是全身都不自在的感覺.
猶太人不僅與巴勒斯坦人互相衝突.
在歷史上他們都與基督徒和回教徒長期衝突.
甚至與種族的撒瑪利亞人相衝突.
撒瑪利亞人與猶太人都是阿帕拉罕的後裔.
都是以色列人.
在王國分裂之後.
撒瑪利亞人歸屬北國.
猶太人就屬南國.
其實為何叫撒瑪利亞人呢?.
因為北國曾經把國都叫做撒瑪利亞城.
所以他們就叫做撒瑪利亞人.
不過兩個具有相同祖宗和血緣的群體.
關係非常惡劣.
各種原因有歷史的也有現實的.
歷史的原因包括在主前八世紀.
不少撒瑪利亞人被擄到巴比倫和埃亞述.
他們與當地人通婚.
所以後裔在種族和宗教上都不再保持純潔.
自命血統純正的猶太人非常看不起他們.
這是歷史的原因.
而現實的原因是在耶穌的時代.
撒瑪利亞人是一般妒忌猶太人比較富有.
宗教和政治的地位也都比他們好.
無論如何在耶穌的時代.
猶太人和撒瑪利亞人曾經發生多次的衝突.
甚至有流血死亡的事件.
兩派勢成水火,壁壘分明.
約翰福音四章九節這樣說.
原來猶太人和撒瑪利亞人沒有來往.
猶太人很少和撒瑪利亞人有交往.
不會和他們一起吃飯.
連用一個器皿去喝一杯水都不會.
羅伽福音記載.

$^{81}$耶穌有一次和門徒從加利利前往耶路撒冷.
路過一個撒瑪利亞人的村莊.
但是因為他們是要向耶路撒冷去.
所以撒瑪利亞人都不用問他們是什麼人.
因為他們正前往耶路撒冷.
很明顯不是我們的人.
所以就不接待他們.
換言之,誰和耶路撒冷拉上關係.
就和撒瑪利亞人勢不兩立.
前面說過在耶穌的時代.
猶太人是被撒瑪利亞人富有的.
他們在羅馬帝國享受的政治待遇也被後者為優.
由於他們在政治和經濟地位上有優勢.
所以通常是他們歧視撒瑪利亞人.
而不是撒瑪利亞人歧視他們.
正如在香港.
一定是我們多於印巴裔人歧視我們.
他們不喜歡我們也可能.
不過我們佔主.
我們高位,我們優勢.
所以通常都是高位歧視低位.
低位的人又怎會有資格歧視高位的人呢?.
所以在那群自命,種族和宗教都是正統的猶太人眼中.
撒瑪利亞人是卑劣不堪的.
簡直淪為魔鬼的附庸.
耶穌基督在《約翰福音》8章48節.
就曾經被人指控.
我們說你是撒瑪利亞人並且是鬼附著的.
指責對方是撒瑪利亞人.
就成為了去誣衊,去詆毀的說話.
我說你是旺鎮的人,就立刻罵你.
最大的詆毀.
你可以想像那個標籤是什麼標籤.
並且說對方是撒瑪利亞人和被鬼附著.
差不多這麼嚴重.
你知道耶穌是很憎恨人.
很抗拒被人誣衊是魔鬼的使者.
耶穌曾經有一次這樣斷言.
這樣說的人是在褻瀆聖靈.
是永不得赦免的.

$^{121}$你可以想像耶穌的語氣有多沉重.
不過耶穌沒有為其他人誣衊他是撒瑪利亞人.
做任何辯護,沒有回應.
可能他根本不在乎這樣的稱呼.
他不覺得這是一個貶義.
說的人以為這是一個貶義.
耶穌自己不覺得.
耶穌基督如何看待這兩個族裔的矛盾呢.
在耶穌三年的事奉生涯中.
他很少跟撒瑪利亞人來往.
他曾經表示他定義集中拯救以色列的迷洋.
這是他自覺的使命.
在《約翰福音》四章二十二節.
他說因為救恩是從猶太人出來的.
所以傳福音的對象先是猶太人.
然後才是希臘人和其他外邦人.
這是上帝所定的救贖計劃.
不過在耶穌的事奉和教導中.
最少有三次是涉及撒瑪利亞人.
我們就這三次事件.
可以窺探到他對撒瑪利亞人的態度.
第一次就是我們剛才讀的聖經.
我剛才也有提及過.
他在某個猶太人和撒瑪利亞人.
集居的麻風病人村落.
治好了十個麻風病人.
你可以想像一下.
那個麻風病人的村落.
是有猶太人和撒瑪利亞人集居的.
為什麼呢?.
因為當你墮落成了麻風病人.
也沒有什麼好分辨的.
你也是人間最壞的人.
那些猶太人和撒瑪利亞人就不再分種族了.
就堆在一起.
他治好了十個麻風病人.
他們按照律法的條例去祭司驗明正身.
然後重獲正常人生活的資格之後.
而這個人,唯一的人就是撒瑪利亞人.
這個情況令耶穌非常感慨.

$^{161}$他嘆息說.
除了這個外族人.
再沒有別人回來歸榮耀與上帝嗎?.
這是路加福音17章11-19節的記載.
猶太人一向自命給撒瑪利亞人敬虔.
不過這一次,顯然是撒瑪利亞人更敬虔.
最後,耶穌宣佈對撒瑪利亞人的拯救.
這是第一件事件.
第二件事件.
是一次耶穌為了回答一個猶太律法師試探性的提問.
就說了一個很撒瑪利亞人的比喻.
這是大家很熟悉的.
路加福音10章25-37節.
他特別將一個祭司,一個利未人.
和一個沒有任何特殊的宗教背景的撒瑪利亞人.
做對照.
認明猶太人眼中最虔誠,最懂得摩西律法.
最遵守上帝的誡命的人.
在愛倫社的事上竟然比不上一個普通的撒瑪利亞人.
今天我們從耶穌的比喻學習愛倫社的道理.
我們很少注意這個比喻本身帶有強烈的挑釁性.
其實耶穌這一句對話不是一個溫柔的講道.
其實根本上是和吵架差不多的.
是針鋒相對的.
因為聖經講明了.
路加福音10章25-28節.
律法師是要試探耶穌.
不是一個平和的討論.
不過我們看不出針對性.
在當時期的人就一定看得出.
他們不可能聽不出耶穌說這番話的諷刺意味.
耶穌竟然將尊貴的祭司和諸狗不如的賤民放在一起.
並且說明那些祭司竟然在表現上不及那些賤民.
你可以想像他們怎樣覺得被冒犯.
第三個例子是.
第三個事例是.
我們下次會講的.
耶穌有一次在聚家城外的雅各井旁.
和一個撒瑪利亞婦人談刀.
這個故事其實也很有革命性.

$^{201}$耶穌打破了撒瑪利亞人.
特別是撒瑪利亞婦女往來.
不和他們共用一個不水的器皿的禁忌.
連撒瑪利亞婦人也覺得很意外.
在四章九節中.
你既是猶太人.
怎會向我撒瑪利亞婦人要水喝呢?.
耶穌不僅僅問她要水.
更主動和她談話.
由活水談到真正的敬拜.
藉著對敬拜地點的討論.
耶穌指出敬拜的地點根本不重要.
耶路撒冷沒有壟斷上帝聖地的資格.
只要心靈和誠實.
無論在那裡敬拜.
都獲得上帝的閱立.
如果地域也不重要.
族群的差異自然也不重要.
耶穌是不著痕跡的.
吻除了猶太人和撒瑪利亞人的族群差異.
指出他們有一位共同在天上的父.
結果撒瑪利亞人信了耶穌.
並且在他的傳福音下.
連帶同村的撒瑪利亞人都信了耶穌.
耶穌在他們中間住了兩天.
親自教導他們有關道理.
這是《約翰福音》4章39-42節所記載的.
耶穌對門徒說.
向撒瑪利亞人傳福音.
正是進行差我來者的旨意.
作成他的功.
耶穌為此高興快樂.
這是聖經唯一一次記述.
耶穌專門在撒瑪利亞群體裡工作.
也可以說是耶穌傳到生涯裡.
其中一個成功的經歷.
約翰在記述了這個故事之後.
再刻意提到耶穌在加利利本地不受人尊敬.
用他來對照他在撒瑪利亞人中間所得到的待遇.
三個有關撒瑪利亞人的故事.

$^{241}$不約而同都產生了和猶太人相對照的效果.
撒瑪利亞人的異.
對照猶太人的不異.
作為上帝嫡系的選民.
猶太人要謙恭自省.
不可以以自己的身份.
以自己的表現.
必蒙上帝月立自居.
從幾個科幻書的記述.
我們注意到的是.
聖經沒有刻意美化現實的撒瑪利亞人.
沒有將他們描述是敬虔正義的.
他們曾經拒絕接待耶穌和門徒.
令門徒非常憤怒.
那個痊癒了的麻瘋病人.
沒錯是比較敬虔的人.
不過罪加成的婦人.
肯定是一個道德敗壞的婦人.
由頭到尾都是好人.
不是的.
所以在現實裡面.
撒瑪利亞人有好有壞.
和其他族群一樣.
無論是屬於哪個族群.
都不會決定了個別的成員.
本人的好壞.
我們不可以單就人的血緣國籍.
而預先對他做道德的判斷.
在耶穌基督的比喻裡.
撒瑪利亞人最後往往成為了正派的角色.
顯示耶穌要刻意打破.
他對那班聽眾.
對撒瑪利亞人一貫的歧視.
從耶穌的比喻裡面.
我們可以確定.
耶穌自己對撒瑪利亞人.
是完全沒有任何的歧視.
沒有任何的解體.
我們剛才讀的聖經.
提到耶穌和門徒.

$^{281}$不被撒瑪利亞村莊接待的故事.
當發生這件事之後.
雅各和約翰就問耶穌.
主啊.
你要我們吩咐火從天上降下來.
消滅他們將以利亞所作的魔.
這裡所說的.
以利亞吩咐天火降下來.
消殺撒瑪利亞王赫阿哈謝.
一百個士兵的故事.
是記載在聶王記下第一章.
今天不詳細講這個故事.
我們不要因為門徒講的這番殘忍的話.
而感到詫異.
這是當時的猶太人.
對撒瑪利亞人典型的態度.
撒瑪利亞人應該被天誅地滅.
要維持一個族群歧視的態度.
總是要將對方典型化.
就是將整個族群所有成員.
都壓平.
看為一個的板.
然後再將他妖魔化.
就是將他們徹底非人化.
否則我們如何在自己的族群裡.
繼續傳遞繼續經營.
這種偏見和歧視呢.
我們的聖經記載耶穌這樣回答.
你們的心如何你們並不知道.
人子來不是要滅人的性命.
而是要救人的性命.
不過路加福音最古老的版本.
是沒有這段說話的.
這段說話有可能不是耶穌親自說的.
後來教會補充進去的.
或者是耶穌親口說過.
不過不是所有版本都將這段說話放進去.
路加本人最初沒有將這個傳統.
收納到他的作品.
後來才引補進去.

$^{321}$無論耶穌有沒有直接說過這番說話.
他肯定沒有同意門徒的建議.
而是和門徒神後轉就離去了.
這段說話令我有很多很多的感觸.
我拉了個道具.
這本書是我去年買的.
這本書叫做「無以神為名」.
Not in God's Name.
我看的是翻譯本.
我去年買.
幾個月前我有一次出門.
我在書架撿了它.
打算在路上看.
不過不是幾個月前.
一個月前.
不過一邊看一邊看.
我看到一種想作嘔的感覺.
這本書主要的評論是.
那些回教徒.
他們常常以聖戰為名.
去發動恐怖襲擊.
發動戰爭.
這是污衊上帝的名.
不可以用上帝的名義去發動戰爭.
這樣的論點.
我猜我們沒有人會反對.
不過剛剛好.
我看這本書的時候.
買的時候早點.
我早點看就沒事了.
我看的時候.
剛剛好就是二巴衝突的時候.
我在想.
如果我們用上帝的名義.
去發動恐怖襲擊.
這是一個妄稱耶和華上帝的名.
我們宣告.
這個地方是上帝給我們的.
於是我們就可以將一個.
二千多年.

$^{361}$住在這裡二千多年的地方的人趕走.
然後我們有大條道理.
理直氣壯地搬進去.
這是不是妄稱上帝的名呢?.
為什麼不是呢?.
這跟聖戰的論點有什麼不同呢?.
有什麼不同呢?.
你試試想想.
有什麼不同呢?.
確實.
如果我用讀歷史的人.
領土的主權的問題.
是一個太複雜的問題.
因為歷史從來都不是穩定.
不是一致的.
我們中國的版圖.
其實是隨著時代.
有大有小.
你要說中國版圖.
最好就是說元朝.
厲害了.
如果我們說.
元朝的版圖就是我們中國人應該有的版圖.
所以今天.
我們應該要維持元朝的時候.
所有的土地.
我們根本上就真的可以.
大條道理.
去問中亞.
問今天的俄羅斯.
甚至去到歐洲.
去爭取土地.
土地的問題是很複雜的問題.
今天中國跟周邊的好幾個國家都有領土的紛爭.
我不討論誰是誰非.
我只是很想簡單說.
事實上要論斷這個地方是屬於誰.
永遠都不是一個簡單可以說的東西.
巴勒斯坦地.
或者叫做菲利斯人地.

$^{401}$以色列人在那裡的時間其實一點都不長.
如果我們從.
14.
1473年.
約書亞帶領以色列進入迦南.
去到北國的滅亡.
在8世紀.
只是700年的時間.
就算數完南國滅亡.
也不到900年.
在歷史的長河裡.
900年其實是極短的時間.
為什麼這個地方.
之前又不是以色列人的.
之後又不是以色列人的.
為什麼必須是他們的呢.
這是一個很彎的問題.
我們唯一可以依據的就是信仰的理由.
上帝應許將這個土地歸屬於以色列.
所以就不理任何的情況.
不理要付多少性命的代價.
這個地方必須要屬於以色列人.
就是這樣.
但我不是說以色列人.
沒有權去爭取他們的土地.
就正如巴勒斯坦人.
有權去維護他們的故地.
我再說.
民族之間的衝突.
其實真的不容易簡單去論斷誰是誰非.
不過對我們.
作為第三者的基督徒來說.
不要簡單的.
特別用時代論的.
對我來說是相當爛的神學.
去爭辯我們有理由.
必須要支持以色列人.
去擴張他們的土地.
對我來說.
整本聖經.

$^{441}$必須要以耶穌基督作為解說.
所有舊約的應許.
是應驗在耶穌基督身上.
耶穌基督宣告他的國不在人間.
他沒有如猶太人所期待的.
成為一個政治的利益.
恢復一個實體的以色列國.
而我是拒絕任何宣稱.
舊約的應許.
舊約的預言.
可以在獨立耶穌之外.
要在耶穌基督之外.
另外實現的這些爭論.
我要說這是時代論的爛說.
這不是聖經的說法.
我的信仰.
永遠是以耶穌基督為中心.
如果不以耶穌基督來幫我們去了解.
去重溫.
去控制我們對舊約的詮釋.
今天我們不僅可以合理化聖戰.
我們還可以合理化種族滅絕.
並且當我們強調耶穌基督.
才是上帝自然在人間唯一的兌現.
唯一的計劃.
我們知道我們對猶太人最大的使命.
是傳福音.
正如我們對所有的民族.
包括回教徒.
我們的使命都是傳福音.
如果我們基督徒.
我們在政治上做了一個.
壁壘分明的抉擇.
我可以告訴大家.
我們就莫想再向十幾億的回教徒傳福音.
我們是和平之子.
我們在任何情況下.
都不會合理化,正義化一個戰爭,一個屠殺.
當然任何一方.
我不是說一方屠殺另一方.

$^{481}$是對還是錯.
我不是這個意思.
而是原則上我們拒絕.
合理化任何的戰爭和屠殺.
當然如果我們是當事人.
我們要打保衛戰的領計.
但作為第三者.
我不覺得我們在原則上可以這樣做.
我希望我做這些時評論.
我自己都有點不安.
不過這個確實是我非常非常大的困擾.
我真的要說我以前都跟你們說過.
基督徒是和平之子.
不要有太多鬥爭,抗爭的心態.
我以前都跟你們說過.
聖經的確有說過屬靈戰爭.
不過聖經說的屬靈戰爭和我們很多時候說的屬靈戰爭是完全兩碼子的事.
我們要披戴基督傳福的軍裝.
我們要做基督中庸的戰士.
不過你睜大眼睛看.
《伊忽所書》第六章.
說基督的軍裝.
即是那一段.
保羅是怎樣說的.
他有叫我們打嗎.
怎樣打.
你和撒旦怎樣打.
披戴軍裝就能夠抵擋魔鬼的詭計.
是叫我們抵擋.
不讓對方贏.
就不是你可以弄死撒旦.
你怎樣弄死撒旦.
你說話.
你弄得死嗎.
所以主禱文.
不叫我們遇見士譚.
救我們脫離兇惡.
只要我們不讓撒旦的計謀得逞.
就是贏了.
是Defensive(防衛).

$^{521}$唯一比較Offensive(攻擊)的.
就是在撒旦的營裡.
把人的靈魂搶救出來.
就是這樣.
所以有些戰爭的話語.
我是很不安的.
那些殺氣騰騰的萬世戰爭的說法.
其實我是很不安.
根本不是聖經的教導.
根本不是.
耳畔所說是很清楚的說.
我們如果今天要戰爭.
不是跟屬血氣的戰爭.
是跟天上屬靈界的力量戰爭.
跟天上屬靈界的力量戰爭.
就不是表示我們於是乎就要去問.
在人間誰是撒旦的化身.
誰是撒旦的爪牙.
誰是撒旦的走狗.
不是的.
然後到處把人標籤.
把一個族群標籤.
把一個我不喜歡的對象標籤.
把我的敵人標籤.
No No No No No.
聖經的教導倒轉頭來提醒我們.
縱然人間.
我們會面對很多的攻擊批評.
面對很多很多的反對.
我們可能會因此受虧損.
我們必須有屬靈的眼睛看到.
其實不是他.
不是那個攻擊我的對象群體.
而是背後.
撒旦.
才是我真正的敵人.
因此.
是這個這樣的發言認定.
叫我們不要將人間任何的對象妖魔化.
不將人妖魔化.

$^{561}$因為我知道.
真正迫害我的是撒旦.
不是眼前的人.
這個例子我講過不止一次.
在你中間.
在我剛剛信主的時候.
我媽媽非常反對我信主.
更加反對我做傳道.
因為我是家中唯一的大學生.
我都做了傳道.
家裡就不用指望可以翻身.
不會發達.
我媽媽難阻我.
難阻得很激烈.
我媽媽有沒有可能成為撒旦的工具呢.
當然有.
不過.
我會不會因此論斷我媽媽是撒旦呢.
如果我媽媽是撒旦.
她是不是我傳福音的對象呢.
她是不是我愛的對象呢.
你怎麼會愛撒旦呢.
你和撒旦有話談嗎.
有共同話題嗎.
不可能.
我媽媽不是撒旦.
開玩笑.
我媽媽是我愛的對象.
是我關心的對象.
是我傳福音的對象.
是不是.
她可以被撒旦利用.
但她不等於就是撒旦.
我仇恨撒旦可以.
但我不會仇恨我媽媽.
我不會.
正如.
當彼得難阻耶穌上十字架的時候.
耶穌立刻喝撒旦退我去吧.
但耶穌有沒有長久地認定彼得是撒旦呢.

$^{601}$怎麼會呢.
彼得是耶穌所愛的門徒.
是不是.
耶穌愛他一生.
對於那些迫害耶穌.
將耶穌釘上十字架的人.
耶穌的評論是什麼呢.
父啊.
赦免他們.
因為他們所做的.
他們不曉得.
他們不知道他們竟然成為了撒旦的爪牙.
是不是.
他們不知道.
起碼他們不知道.
所以我們可以憐憫他們.
起碼他們不知道.
所以他們情有可原.
起碼他們不知道.
他們還值得我們去拯救去挽回.
如果你將任何人.
任何族群標籤做了撒旦.
那你還有向他傳福音的動機嗎.
你和他還有任何可以談的動機嗎.
弟兄姊妹.
我們正在參與一場屬靈的戰爭.
是屬靈的戰爭.
這個戰爭真正發生的地點不在人間.
雖然在人間我們會碰到種種的遭遇.
不同的人和事.
有些對我們好.
有很多對我們不好.
有很多冒犯我們.
傷害我們.
辜負我們.
踐踏我們.
是的.
不過.
聖經是唱高調的.
不是梁家鑾唱高調.

$^{641}$是聖經唱高調.
以善成惡.
以愛還恨.
那些恨你的.
要愛他.
愛你的仇敵.
是不是高調.
是高調.
當然唱高調.
不過.
因為這是耶穌說的.
再高調我們都要硬吃.
是不是.
你可以不接納這個教導嗎.
我不可以.
不理我喜不喜歡.
我有衝地.
我認同不認同.
我也不可以不接納.
聖經的教導.
唯有這樣.
我們才能夠打破這個世界.
藉著仇恨建立的一重又一重的鎖鏈.
我們可以活得不一樣.
可以顛覆這個世界的秩序.
弟兄姊妹.
我真的有衝的期望.
我們活在人間.
我們真的周遭見到的.
都是對立,衝突,分裂,仇恨.
互相攻擊的現象.
我們的主就快在月底.
我們紀念祂降生.
這位是和平之子.
這位是和平之君.
在天上永要歸給上帝.
地上平安歸於祂所喜悅的人.
但願我們真的作為基督徒.
我們可以在人間帶來平安.
帶來喜悅.

$^{681}$努力壓低我們心中對人的仇恨.
對人的憤怒.
對人的不接納.
提醒自己.
我永遠要保持對人善意.
你們耳朵都起繭了.
永遠對人有good will.
一個人一天未死.
就不要對他蓋棺定論.
不要對任何人做本質性的論斷.
只要耶穌肯動一個人.
任何人都可以改變.
這是我們很傻的.
但我們生命深層的信念.
永遠不對任何人絕望.
耶穌可以.
耶穌能夠改變人.
為此我們就拗緊牙關.
希望我們用寬容的心.
我們活在人間.
傳播和平的福音.
我們遵循傳播與人和好的信息.
宣告人可以與上帝和好.
然後可以破除人間種種間斷的牆.
你記得這是以法所書的教導.
耶穌基督對撒馬利人的接納.
我希望是我們每一個人一個很好的提醒.
我不覺得我們在香港.
有很嚴重的種族歧視.
族群歧視的問題.
但完全沒有並不是真的.
華人的民族排他性不見得很少.
過去二百年國勢很衰弱.
我們一直被人歧視.
比別人歧視我們多.
我們在別人面前總有自卑感.
所以我們談不上歧視別人.
但對於比較弱小的族群.
例如印巴裔,菲律賓人或其他有色人種.
我們總有不同程度的歧視.

$^{721}$甚至對同族的大陸同胞.
我們仍然有很多不同程度的看法.
有個弟兄曾經告訴我一個故事.
他認識一個在南非出生的印度裔朋友.
對方經常來香港探望他.
他們有一次坐地鐵.
上車後他們看到地鐵有一個空位.
那位朋友就跟他說.
我們來做一個實驗.
如果你坐下我就沒有位子坐.
但如果我先坐下.
兩個人都有位子坐.
結果實驗是真的對的.
他坐下後旁邊的人就站起來了.
兩個人都可以坐下.
我曾經在馬來西亞某個城市.
住在一個教會長老的家.
他家是很大的豪宅.
他有幾個來自印尼的傭工.
好像我們那些非傭太傭.
那個女主人就跟我說.
她聘請外籍傭工的時候.
總是跟他們預先聲明.
不給他們任何假期.
聽清楚我說嗎?.
是任何假期.
一天都不給的.
每個星期有一天的例假是法定的.
他就用錢去替代.
不讓他們離開家門半步.
主要原因是怕他們.
一旦跟當地的馬來西亞男人接觸.
就會被勾引走掉.
我明白每個地方的國情不同.
所以我不可以簡單做任何的評論.
不過我聽到這樣的說法.
我心裡面總是有些不忍.
請你原諒我這種小男人.
總是有些心軟.
我想一年到頭都沒有假期.

$^{761}$也不可以離開家門半步.
你告訴我.
那跟坐牢有什麼不同呢?.
如果那個傭工在家裡.
不可以跟家主平輩論交.
成為朋友.
那傭工就連一個可以聊天的對象都沒有了.
你說現在可以用WhatsApp的.
你說多可怕.
多可怕.
很多年前.
其實我用了非傭幾十年了.
我記得都不止一次有人善意提醒我.
千萬不要給傭工星期天放假.
就算他們說要回教會也不好.
做這個建議的都是基督徒.
寧可給他們平日放假.
他們就不會在中環聚集成群.
被同鄉帶壞.
我的反應是.
我怎麼可以為一個已成年的人.
或者會被人帶壞.
而做任何防禦的公事呢?.
怎麼可能呢?.
我沒有什麼可以自誇的.
不過這件事我都可以很囂張的說.
我用了這麼多年非傭.
我一直堅持每天工作八小時.
所以我的非傭從來都不需要.
星期天像走難一樣一早疏散.
我今天早上出門的時候他都沒有起床.
或者最少都沒有開睡房.
他就在家裡嘻嘻嘻的賴著.
最少他不覺得這裡是他不可以在家的地方.
無論如何.
我們翻閱四卷的福音書.
耶穌沒有就族群歧視的問題.
做過任何正面的教導.
直接的教導.
我們也看不到他曾經公開批評猶太人歧視撒瑪利亞人的情況.

$^{801}$或者主張怎麼糾正他們的想法.
我們看到的卻是.
耶穌自己沒有以族群地域.
或者任何人間的元素.
而將人做先驗的判定.
卻是超越這些的分門別類.
然後接觸每一個個別的群體.
他是用新教而不是現教.
教導我們不要偏袒人.
好像彼得知道上帝拯救外邦人哥利留的時候.
他有的反應.
我真看出上帝不偏袒人.
《紹行傳》十章三十四節.
無論如何.
我估計這也是我們彼此的提醒.
今天我們不是說原則性的.
說了眾生平等也不一定很有效.
我希望我們開放我們自己.
進入每一個你能夠接觸到的人的心裡.
個別接觸和離遠評論是兩回事.
我有時候也會有很噁心的念頭.
當我在巴士上看到後面的人.
開著手機大喇叭播放YouTube.
我也有想掐死他的感覺.
正如我兒子也說過.
他最討厭看電影的時候後面的人說劇情.
後面的人說劇情.
他真的想用鞋子掐死後面的人的感覺.
你看到這些事情.
有時候心中也會有無名火氣.
也會想掐死他.
主要原因是我不認識他.
我媽媽生前說話非常大聲.
她從小也是張開嘴巴說話.
她的說話會為我們家人帶來很多尷尬.
跟她吃飯時張開嘴巴.
四桌子也聽到她說話.
有時候我也覺得很難受.
勸她小聲一點.
我媽媽總是瞪著我罵我.

$^{841}$但我不會因此討厭我媽媽.
因為我認識她.
我做院長就知道.
如果純粹是有同學犯事要處理.
要我公事公辦是很容易的.
我就是論事.
但如果我認識同學.
我進入過他的世界.
聽過他的故事.
我就心軟了.
我就會有很多囉嗦.
不想一下子就把他弄死.
不會的.
當你深入認識一個人.
你聽到他的故事.
你就不會簡單地將他標籤.
這是壞人 本質上的壞人.
這是魔鬼的化身.
我不會的.
當你認識他.
你對他就會看得立體一點.
不是表示你擁抱他所有不好的行為.
但最少你不會輕言.
把他弄死就算.
把他消滅就算.
你不會的.
我求主幫助我們.
認定每一個我們接觸到.
我們看到的都是有血有肉的生命.
正如每一個被恐怖襲擊.
殺死或俘虜的人.
每一個被炸彈所殺害的人.
個別都是有血有肉的人.
都有他的愛.
有他的憎.
有他的喜.
有他的惡.
真實的人.
而真實的人.
是不值得我們對他做簡單無段的論斷.

$^{881}$上帝幫助我們.
我們真的努力成為一個和平之子.
\newpage



\section{馬太福音 2:1-12}
\label{sec:QUWqIbTXwuo}
\textbf{2023年12月24日聖誕節主日崇拜直播｜陳冠成牧師｜唯獨你是王 - 耶穌帶來的驚與喜｜馬太福音2\_1-12}
\newline
\newline
連結: \href{https://youtube.com/watch?v=QUWqIbTXwuo}{\texttt{https://youtube.com/watch?v=QUWqIbTXwuo}} ~~~~ 語音日期: 2023-12-25
\newline
\newline
\hyperref[sec:LudzCxftFrQ]{\small{< < < PREV SERMON < < <}}
~
\hyperref[sec:index_chronic]{\small{[返順時目]}}
~
\hyperref[sec:iN3iqiIAXcs]{\small{> > > NEXT SERMON > > >}}
\newline
\newline
馬太福音 2:1-12
\newline
\begin{longtable}{cl}
\hline
\hline
章節 & 經文 (和合本修訂版)\\
\hline
2:1 & \begin{tabularx}{0.7\textwidth}{X} 在希律作王的時候，耶穌生在猶太的伯利恆。有幾個博學之士從東方來到耶路撒冷，說： \end{tabularx} \\ \\ \relax
2:2 & \begin{tabularx}{0.7\textwidth}{X} 「那生下來作猶太人之王的在哪裡？我們在東方看見他的星，特來拜他。」 \end{tabularx} \\ \\ \relax
2:3 & \begin{tabularx}{0.7\textwidth}{X} 希律王聽見了，就心裡不安；耶路撒冷全城的人也都不安。 \end{tabularx} \\ \\ \relax
2:4 & \begin{tabularx}{0.7\textwidth}{X} 他就召集了祭司長和民間的文士，問他們：「基督該生在哪裡？」 \end{tabularx} \\ \\ \relax
2:5 & \begin{tabularx}{0.7\textwidth}{X} 他們說：「在猶太的伯利恆。因為有先知記著： \end{tabularx} \\ \\ \relax
2:6 & \begin{tabularx}{0.7\textwidth}{X} 『猶大地的伯利恆啊，你在猶大諸城中並不是最小的；因為將來有一位統治者要從你那裡出來，牧養我以色列民。』」 \end{tabularx} \\ \\ \relax
2:7 & \begin{tabularx}{0.7\textwidth}{X} 於是，希律暗地裡召了博學之士來，查問那星是甚麼時候出現的， \end{tabularx} \\ \\ \relax
2:8 & \begin{tabularx}{0.7\textwidth}{X} 就派他們往伯利恆去，說：「你們去仔細尋訪那小孩子，找到了就來報信，我也好去拜他。」 \end{tabularx} \\ \\ \relax
2:9 & \begin{tabularx}{0.7\textwidth}{X} 他們聽了王的話就去了。忽然，在東方所看到的那顆星在前面引領他們，一直行到小孩子所在地方的上方就停住了。 \end{tabularx} \\ \\ \relax
2:10 & \begin{tabularx}{0.7\textwidth}{X} 他們看見那星，就非常歡喜； \end{tabularx} \\ \\ \relax
2:11 & \begin{tabularx}{0.7\textwidth}{X} 進了房子，看見小孩子和他母親馬利亞，就俯伏拜那小孩子，揭開寶盒，拿出黃金、乳香、沒藥，作為禮物獻給他。 \end{tabularx} \\ \\ \relax
2:12 & \begin{tabularx}{0.7\textwidth}{X} 因為在夢中得到主的指示，不要回去見希律，他們就從別的路回自己的家鄉去了。 \end{tabularx} \\ \\
[1ex]
\hline
\hline
\end{longtable}
$^{1}$今日經訓是在馬太福音第二章一至十二節.
是在新約聖經第二頁.
馬太福音第二章一至十二節.
請聽我讀出第二章一節.
「當希律王的時候.
耶穌身在猶太的伯利恆.
有幾個博士從東方來到耶路撒冷.
說:那身下來作猶太人之王的在哪裡?.
我們在東方看見他的聖.
特來拜他.
希律王聽見了就心裡不安.
耶路撒冷合成的人也都不安.
他就召齊了祭司長和民間的文士.
問他們說:基督當身在何處?.
他們回答說:就在猶太的伯利恆.
因為有先知記著.
說:猶太的伯利恆.
你在猶太諸城中並不是最少的.
因為將來有一位君王要從你那裡出來.
莫讓我以色列民.
當下希律暗暗地召了博士來.
誓問那星是什麼時候出現的.
就差他們往伯利恆去.
說:你們去仔細尋訪那小孩子.
尋到了就來報信.
我也好去拜他.
他們聽見王的話就去了.
在東方所看見的那星.
忽然在他們前頭行.
直行到小孩子的地方.
就在上頭停住了.
他們看見那星.
就大大的歡喜.
進了房子.
看見小孩子和他母親瑪利亞.
就俯伏拜那小孩子.
揭開寶盒拿黃金與香抹藥.
為禮物獻給他.
博士因為在夢中被主指示.
不要回去見希律.

$^{41}$就從別的路回本地去了.
聖經獨白.
弟姐們.
聖誕快樂.
我想剛才所聽到的經文.
東方的博士來訪尋耶穌基督的降生.
其實是我們很熟悉的聖經故事.
盼望我們今天透過這個.
「以數能詳」的故事.
我們再一次藉著上帝的靈的幫助.
我們對耶穌基督的降生.
有再一次新的體會和提醒.
話說剛才也說了.
耶穌基督降生在.
希律王統治巴勒斯坦地的年代.
先說說希律王是哪一位.
其實他有些少猶太人血統.
他本身是以東人.
大家記得以東是誰的後裔?.
以數.
雅各有兩個兒子.
大兒子是以數.
以數的後裔.
他有些少猶太人血統.
但他不是純正的猶太人.
他為何可以做到希律王這個位置.
其實就是當時羅馬帝國.
進攻巴勒斯坦地的時候.
他非常賣力替羅馬政府工作.
所以得到當時羅馬帝王的信任.
於是就冊封他.
成為一個分封王.
希律.
負責管理猶太地.
撒馬利亞和加利利一帶.
即是耶穌基督在新約時代活躍的幾個地方.
而希律的管治手法是怎樣呢?.
軟硬兼施.
一方面他會用強硬的手段.
來維持整個巴勒斯坦地的和平.

$^{81}$另一方面他也很懂得討百姓的歡心.
當社會不景氣的時候.
他真的會減稅.
來體恤當時百姓的辛勞.
希望減輕他們的負擔.
又或者當時如果有大饑荒的話.
歷史記載他曾經自己掏錢包.
購買穀物來鎮災.
再加上他本身也很喜歡.
做很多大型興建.
我們熟悉的新約時代.
耶穌基督的時代.
那個聖殿是很宏偉的.
其實就是當時的希律王.
他拍板要重建.
他的面積和輝煌程度.
其實比起所羅門時代的聖殿.
可以說是有過之而無不及.
不過.
希律作為一個管治者.
其實他也有一個陰暗面.
我們說他做人是沒有安全感的.
很多忌諱.
譬如他懷疑身邊有人.
威脅到他的王位和政權的時候.
他一定會想盡辦法將對方剷除.
到了晚年的時候.
這種疑心更加強烈.
甚至他懷疑他身邊的人.
歷史記載他懷疑他的妻子.
甚至兒子想對自己不利的時候.
他真的會殺了她.
所以當時的羅馬帝王.
就很刻薄地諷刺希律.
他說做希律的豬.
比做他的兒子好.
更加穩妥.
所以我們不難理解.
有了這個背景.
為什麼當希律王聽到.

$^{121}$耶穌基督降生的消息的時候.
聽到有一位新的猶太王誕生的時候.
他的心會這麼不安.
這幾位東方博士.
即是對當時的天文.
很有研究的星象學家.
他們從東方來到耶路撒冷.
很大機會他們是異教徒.
不過當猶太人昔日散居在東方的時候.
他們這班博士可能曾經接觸過.
當時的猶太人群體.
亦從他們身上得知.
舊約聖經曾經預言.
確實會有君王從猶大這個地方誕生.
並且他們都相信.
假如這位偉大的君王出生的時候.
天上一定會有一些異象.
例如有一些星獸的升起.
來揭示降生的消息.
現在正好.
他們看到那顆火象.
正正由猶大君王出生的星.
就從巴勒斯坦地上空的方向升起了.
於是他們按照常理.
君王應該在哪裡出生?.
皇宮.
正路一定在皇宮.
所以他們按照常識.
猶太人的君王應該在猶太人的皇宮出生.
於是就追星到耶路撒冷.
直接到當時希律的皇宮.
直接問希律.
那個生下來作猶太人的王在哪裡?.
我們在東方看見象徵著他出生的星.
現在升起了.
我們知道他應該誕生了.
所以我們專程遠道而來.
來朝拜他.
我們說這番博士就是.
崩口人的崩口碗.

$^{161}$對希律來說.
誰是猶太人的王?.
就是他.
我就是猶太人的王.
你們這樣說.
我不是真命天子.
換句話來說.
又會有人威脅到我的地位.
於是經文裡說.
希律很擔心他的王位又再受到威脅.
並且他這份不安.
影響到周邊的人.
你會發覺很奇怪.
為什麼希律的不安.
會令到耶路撒冷上下所有人都不安呢?.
其實不難理解.
譬如你上班.
你的老闆喜怒不容於色.
你今天看見他走來走去.
皺眉頭.
你就會知道.
麻煩了.
又不知道發生什麼事會否連累到我們.
事實上他們都是這樣想.
耶路撒冷上下的百姓.
他們很明白希律是一個疑心重的人.
現在希律聽見在我們猶太人中間.
會有人威脅到他的地位.
我們還有沒有好日子過?.
沒有了.
我們要打醒精神.
因為他們很熟悉希律一定會用盡不同的手段.
在我們中間抽出壞蛋.
然後將它消滅.
你看到希律這一種的不安.
很奇妙吧.
是因為耶穌基督的誕生.
而為了消除這份不安感.
當然他就要開始部署.
如何令耶穌基督消失在這個世界上.

$^{201}$他有很好的部署.
如果你細看這幾節經文.
首先他找一些很熟悉聖經的人.
就是召集當時猶太人的祭司,長和文士.
來到去.
你幫我查一下.
到底基督.
即是猶太人.
所盼望來到的那位救世主來塞亞.
他是在哪裡出生的?.
他要掌握到耶穌的出生地點.
然後祭司,長和文士按照舊約的尼加書.
先知裡面確實曾經預告過.
他說基督會出生在猶太的伯利恆.
因為上帝要在那裡興起一位君王.
牧養以色列民.
掌握了耶穌的出生地.
不單止,還不夠的.
不知道祂多少歲.
於是他又暗暗將這班博士.
召回來問他們.
你們見到這顆星升起的時候.
是甚麼時候?.
你們何時發現這顆星出現?.
他要進一步掌握.
耶穌已經出生了多久.
掌握了祂的年齡.
要將目標範圍收窄.
這樣夠不夠呢?.
都不夠.
因為在伯利恆.
不可能只有一個英孩出生.
於是他就吩咐這班博士.
你替我去尋訪這位小孩子.
你找到的時候就回來跟我說報信.
然後我去敬拜他.
所以你看到希律的部署是很周密.
不想太過張揚.
不想被人知道.
就假借東方博士的手.

$^{241}$來幫自己鎖定將要謀殺的目標.
然後就怎樣?.
派人去消滅耶穌.
從希律的反應我們見到.
某程度上反映現代人的真相.
你會否覺得身邊有不少人.
甚至你自己.
你害怕耶穌來到?.
你喜不喜歡耶穌來到?.
你會否覺得耶穌出現在你的生命裡.
其實是一個威脅.
會干擾到你的人生.
例如我們上一代.
很多人都不喜歡信耶穌.
因為一聽到我要信耶穌就要代表.
我失去了我的人生.
我要戒這.
不能做那.
我沒得作主.
因為我要聽耶穌說.
失去了自由.
所以不要搞我.
你還未說他聽到耶穌這三個字.
你還未叫他信.
他已經將耶穌彈開了.
不是所有人都喜歡耶穌的降生.
就好像希律一樣.
事實上確實是.
假如你一直是你自己人生裡.
作主的那個.
你一向是一個君王.
但現在耶穌說他才是.
但很可惜.
王位只有一個.
不可以兩個人同時坐在一起.
不可以你和耶穌基督同時坐在一張王座上.
換句話來說.
是第一個要讓位.
那是誰?.
是你還是耶穌?.

$^{281}$弟兄姊妹.
這不單是希律他要掙扎的位置.
其實也是經文要我們思考的.
當耶穌表明.
他要成為我們生命話事的一位的時候.
其實我們本能反應是怎樣?.
我們會不會都好像希律一樣.
其實心裡也有些不舒服.
有些不安.
拒絕被耶穌改變.
動搖你現在的生活狀況.
你現在做人的方式.
甚至為了消除你內心那種不安感.
你會用盡辦法和耶穌鬥力嗎?.
甚至你用很多方法試圖阻止.
耶穌基督不要出現在你的生命裡.
這是經文要我們自己思考的.
雖然我們很明白.
我們回教會那麼久.
我們慶祝過耶穌基督的降生那麼多次.
我們一定知道耶穌基督降生.
是要拯救我們.
是要成為我們生命的主.
生命的王.
我們需要將生命的主權交給祂.
那張王座是祂坐的.
我們很清楚.
由祂帶領掌管.
這才是生命之道.
我想沒有人會不清楚.
這個道理.
這個真理.
但現實裡.
你覺不覺得很多人.
其實是在走希律那套人生哲學.
覺不覺得?.
就是我凡事要自己親自部署.
親自指揮.
我要務求所有東西.
在我計算掌握之內.

$^{321}$我要確保事情.
盡可能按照我的意思來出現.
總之要做到什麼?.
好像商業管理一樣.
零風險.
沒有失誤的.
人生.
這樣才是贏家.
這樣的生命才是穩妥.
結果你發不發覺.
很多時候.
如果我們這樣想的時候.
所賺到的.
其實不是一個平安安穩的人生.
反而是一個長期被恐懼支配.
很多壓力.
很多焦慮的人生.
我不知道你身邊認不認識這樣的人.
我認識一些家長.
其實他很高壓.
因為子女從小到大.
他都擔心他的前程.
擔心自己.
由抑鬱出現身體很多症狀.
沒有平安.
問他整個過程是為什麼.
其實他說.
沒有擔心他們.
從小擔心到大.
擔心他出生.
擔心他長不大.
擔心他不能入讀好的學校.
擔心他升不上大學.
到他工作,結婚,生孩子.
都繼續擔心.
那怎麼會不生病呢?.
問他整個經歷.
因為我擔心.
所以我就要掌控.
因為我不安.

$^{361}$所以我就要把所有事情都抓得很緊.
我以為這樣會安心一點.
確實是有種掌控感的安全.
不過他發覺越掌控這個世界.
原來是失控.
不是所有事情都聽他講.
世界不是圍著他轉.
越掌控的時候.
發覺越失控.
試過想放手.
但是不行.
為什麼?.
因為放手的時候.
失去那份操控感.
更害怕.
更恐懼.
於是唯有再抓緊一點.
再抓緊一點.
再多一點焦慮.
就跌入了這個循環.
越抓緊越焦慮.
越焦慮越抓緊.
事實上弟兄姊妹.
今天這段經文給我們一個好好提醒.
到底生命.
是否真的我們可以掌握到?.
越想透過人為操控.
來獲得那份真正的平安.
原來是不行的.
原來是越沒有平安.
到底什麼才是安穩的生命?.
我想大家繼續看下去.
這個故事到底如何發展.
邀請大家一起讀.
印在程序表中間那段經文.
第十二至十五節和十九至二十節.
請大家一起讀.
預備起.
.
謝謝.

$^{401}$當我看到這幅圖畫的時候.
我想起自己小時候打機.
當對手武器很強大.
招招拿你命.
你根本無法跟對手鬥.
你可以做什麼?.
避開.
他多厲害都沒有用.
打不中我,我就什麼事都沒有.
其實這幅圖畫就是這樣.
我只要知道你多厲害.
但你拿不到我的命.
因為我全都避開了.
上帝幫我全都避開了.
事實上耶穌基督.
透過他的降生告訴我們.
生命的穩妥不是一點風險都沒有.
生命的穩妥是在於.
你確信上帝會保守帶領你.
這才是關鍵.
上帝容許手無寸鐵的耶穌.
當時他大概兩歲不夠.
加上他的父母都是平民百姓.
面對著希律那麼強大的追殺.
但上帝幫助他們化險為夷.
上帝一路給他指示.
教他如何避開.
並且你有沒有留意.
我們很熟悉博士.
送了幾個禮物給他.
黃金與香活藥.
一方面有他屬靈的象徵意義.
另一方面有他的實際用途.
當耶穌逃難到埃及的時候.
身無長物.
那怎麼辦?.
那些就是他的盤川.
變賣了就成為他足夠在埃及生活的費用.
所以你看到上帝全方位的保守.
供應他們各樣所需.

$^{441}$讓耶穌一家即使身處在風暴的裡面.
仍然安穩無恙.
弟兄姊妹.
很值得我們去反省.
不知道大家有沒有坐過山車?.
有嗎?.
喜歡坐嗎?.
有沒有人喜歡坐過山車?.
有嗎?.
很少的.
我自己也一般般.
因為我不喜歡那種刺激.
不舒服的感覺.
當然我會擔心機件出意外.
真的會沒命.
但人生難免會坐過山車一兩次.
特別是有小朋友.
我坐的時候一定會刺激到失神.
盡量閉上眼.
快點完.
我覺得這樣比較穩.
所以我很佩服那些玩過山車的人.
他可以這樣.
有沒有見過?.
有沒有試過?.
很開心地呼叫.
不是尖叫.
是呼叫.
很開心地呼叫.
放開雙手坐過山車.
很放膽地享受在驚險裡面的滿足.
那份樂趣.
其實人生有時真的可能會像坐過山車.
是有風險的.
重點是你夠不夠膽在風險裡面放手.
享受一下那種人生.
高高低低 上上落落.
其實裡面有樂趣.
當然我們明白.
其實坐過山車.

$^{481}$你和我只不過是乘客.
不是駕駛.
你有沒有試過駕駛過山車?.
坐在上面.
不是在控制室裡.
是駕駛不到的.
你越想掌控的時候.
其實你的挫敗感會越大.
因為無論你駕駛得多穩定.
那輛過山車是不會跟著你的意思走.
你不會改變它的軌跡.
不會改變它的速度.
你越想控制它.
你就會越不安越無力.
不過今天這段經文給我們好好提醒.
你不需要懂得怎樣駕駛.
你自己人生的過山車.
不過你要懂得.
誰可以駕駛到.
你要認識他.
你要奮身依靠他.
就好像耶穌基督那樣.
他不需要保護自己.
他靠上帝搞定.
他放心奮身.
依靠天父各樣周詳的供應.
這樣就行.
所以你看到其實.
英夏耶穌的生命.
不就像坐過山車一樣.
他沒有辦法掌握自己的生命.
但他確信他爸爸可以.
天父可以搞定.
所以他可以放心.
不怕身處危險的地方.
其實我們有沒有發現.
放手完全信靠天父的保守帶領.
其實就是耶穌基督來到世間上.
第一個向我們展示的人生的功課.
他第一個放手.

$^{521}$赤裸裸地面對挑戰.
面對追殺.
事實上如果你有留意到.
耶穌離世那一刻.
他也展示著放手交托的功課.
記得耶穌離世的時候說了一句什麼嗎?.
父啊.
我將我的靈魂交回給你.
你帶領.
所以弟兄姊妹.
其實今天這段經文很好的提醒.
就是我們慣不慣在我們人生裡.
特別是那些你抓得很緊的東西裡.
你放手交托給耶穌管理.
話說.
有一個很喜歡行山的人.
他有一次行夜山.
不小心迷失路.
天又黑風又大.
好像今天這麼冷.
他走著走著.
一個不小心滑了下山.
在漆黑一片的路上.
他當然是死撞.
結果他很幸運地撞到一條藤.
不過整個人就懸空在半空裡.
半天吊.
當時他很害怕.
於是當然大叫.
但很可惜當時已經天黑.
黑到漆黑一片.
根本沒有人會在那裡.
沒有人發現他.
他一邊叫一邊覺得不對勁.
最後他沒有辦法.
唯有祈禱.
上帝你救我.
就在這個時候.
他聽到一把很磁場的聲音.
他說.

$^{561}$放手吧.
放手就可以了.
他嚇一嚇.
他說你是誰?.
我就是你呼求的上帝.
你快點放手吧.
放手就沒事了.
結果這個行山人士.
他聽完之後更害怕.
有沒有搞錯?.
你看不到我現在什麼環境.
就死了.
你還叫我放手.
那不就死快點?.
我死都不放.
於是他抓著那條藤更加緊.
到後來.
搜索的人員發現他.
他依然緊握著那條藤.
不過活活地凍死了.
確實很可惜.
當時的搜索人員說.
他一定是太黑.
看不到自己身處的環境.
以為自己身處在半空.
下面就是懸崖.
所以死抓著那條藤不放.
結果凍死了.
事實上.
他只是離地十多呎.
其實有沒有想過.
我們一些生命範疇.
都好像剛才的故事一樣.
抓得太緊.
以為自己最清楚.
自己發生什麼事.
結果凍死了都不知道.
我們不肯將自己的主權.
交回給耶穌基督來大靈.
不過這段經文邀請我們.

$^{601}$或許是時候你放手了.
你放手.
將你自己人生的座位.
無論你坐了多少年.
是時候交回給耶穌去坐.
你放膽靠近耶穌.
讓祂接管你的人生.
因為祂才清楚你什麼狀況.
祂才知道你當下.
最需要的是什麼.
我們繼續看剩下的經文.
我們看到這幾位東方博士.
他們造訪耶穌的目的.
是為了敬拜那位.
生下來作猶太人君王的耶穌.
他們的目的很清楚.
不過我們回頭看.
這件事應該由誰來做?.
應該是耶路撒冷.
上上下下的人去做.
他們盼望已久的君王.
拯救主現在來到了.
誕生了.
不是應該歡喜雀躍嗎?.
但你看到上上下下.
無論是宗教領袖或百姓.
他們都沒有做他們當做的本份.
反而這個敬拜耶穌的責任.
落在了外邦人手上.
由他們來完成.
確實此情此景是值得思考的.
作為基督徒.
我們今天到底以什麼心態.
來回應耶穌基督的降生呢?.
會否像昔日耶路撒冷的百姓.
其實對耶穌的誕生.
已經沒有興奮雀躍.
又或是我們不知不覺.
好像用慶祝生日的態度.
來紀念耶穌基督降生呢?.

$^{641}$什麼是慶祝生日的態度呢?.
慶祝生日有什麼活動呢?.
沒什麼活動.
吃吃東西.
送送東西.
慶祝多久呢?.
一天而已.
第二天就沒事了.
回到現實.
什麼都沒有.
你不會因為今天收到一份禮物.
吃了一餐.
你可以延長到滿足興奮.
繼續延長到下星期或下月.
不會的.
甚至乎可能有些人.
都不喜歡慶祝生日.
有否感覺到?.
隨便吧.
又是這樣.
我們很多時候.
會否用慶祝生日的心態.
來跟耶穌慶祝生日呢?.
弟兄姊妹確實.
年復年地跟人慶祝生日.
其實有時我們真的會失去了.
那種驚喜.
那種興奮.
相比起我們慶祝耶穌生日.
我們會否覺得.
沒有什麼要回應.
唱唱歌.
大家熱鬧一下.
那就算了.
反正明年還會有.
分享一個.
跟聖誕有關的古老故事.
給大家聽.
話說以前有一對年輕夫婦.
他們很貧窮.

$^{681}$結婚後.
他們的生活都很清貧.
不過聖誕節臨近.
做太太的.
因為很愛她的丈夫.
一直很想送一份聖誕禮物.
來取悅他.
但真的沒有錢.
太窮.
於是怎樣呢?.
在聖誕的前夕.
他把自己最心愛.
最亮麗的那把長頭髮.
剪掉.
賣給當時的假髮店公司.
你知道以前.
沒有那麼多人做仿真髮.
真的會收集假髮.
然後賣.
是值錢的.
他把很漂亮的頭髮剪掉.
賣給假髮店.
來換取一些金錢.
也把換取的金錢.
買了一條錶鏈.
很希望把這條錶鏈.
掛在丈夫.
他有一隻很心愛的駝錶上.
這隻駝錶是丈夫的家傳之寶.
傳了兩代.
一直珍如重之.
他多窮都捨不得買.
回到家裡.
丈夫固然看到.
你為何剪了頭髮?.
你不是一直很心愛.
最漂亮,最引以自豪的東西嗎?.
為何你把你最心愛的頭髮剪掉.
來買一隻錶給我?.
丈夫看到很詫異.

$^{721}$因為原來.
他也把自己最珍愛的那隻錶.
賣了.
他把換取的錢.
買了一把.
鑲了寶石的梳.
為甚麼?.
因為太太有一把很美麗的頭髮.
是他珍而重之.
所以他也想買一把名貴的梳.
來襯托太太的這把瘦髮.
讓他每次梳的時候.
都感受到這份愛意.
他知道妻子一直很喜歡這把梳.
但根本沒有能力去買.
到最後.
這兩夫妻情深款款地.
看著對方相擁.
然後喜極而泣.
弟兄姊妹是否很傻?.
傻不傻?.
你會不會這樣做?.
死都不會吧?.
一對應該怎樣?.
很經濟的.
樹根應該怎樣?.
大家商量一下.
你買甚麼?我買甚麼?不要碰.
以免買回來不單不合用.
是得物無所用.
很掃興.
但其實.
會否這才是送聖誕禮物.
一個真正的意義?.
就是藉著我們為對方犧牲寫記.
來成就對方的需要.
我們可能忘記了聖誕的意義.
這份心意才是最珍貴.
事實上你看耶穌基督為何要渡誠育身?.
為何要來到這個世上?.

$^{761}$是否祂願意捨棄自己的尊貴.
放棄祂的榮華.
來到世間成為卑微的人.
甚至最後死在十字架上?.
耶穌基督的降生.
是否預示祂將要為我們犧牲寫記的生命?.
這份就是耶穌送給我們的聖誕禮物.
所以弟兄姊妹.
你收到了.
你如何回禮?.
你如何回應耶穌基督給你的禮物?.
你如何在你的生活中.
為祂犧牲付出?.
這幾位東方博士.
有些東西很值得我們學習.
你發覺他們為了敬拜耶穌基督.
他們千里迢迢艱辛走了相當遠的路程.
如果我們敬拜的生活.
有這份熱情,惡慕.
相信耶穌會很喜悅.
到底我們是否每個星期都渴望.
回到敬拜上帝.
很值得我們去想.
我們有沒有事前預備好?.
雖然我們不會帶禮物回來.
但我們的人應該是最好的禮物獻給耶穌.
讓祂去使用.
我們會否刻意不在星期日早上.
不在崇拜的時候.
安排一些其他的活動.
以致我們真的專注上帝優先.
特別我們知道在疫情後.
我們有網上崇拜.
我們方便了.
但這種方便會否反而削弱了.
我們對耶穌基督的尊崇.
我們今天晚了不來現場.
在家裡看雪.
我們會否不知不覺.
削弱了對耶穌基督的敬拜.

$^{801}$不願意為耶穌基督辛苦一點.
早起一點.
早一點回來預備.
不單單是有事奉崗位的人.
其實我們每一個都是敬拜的人.
我們每一個都應該有所預備.
來朝見耶穌.
不單單是聖誕這一天.
我們需要身體力行的敬拜耶穌.
而另一方面.
你看到當時耶穌很貧窮.
我們經常被聖誕節那幅浪漫的景象.
遮蓋了.
你看到馬槽的景象.
耶穌基督的頭頂有個光環.
周邊的人有個光環.
有美好的星星.
有羊,牧羊人.
所有博士齊齊.
但實際應該是很貧窮.
耶穌基督不可能.
有什麼可以回報.
給這幾位博士的厚禮.
但這樣就更加突顯出.
這幾位東方的博士.
他們怎樣無條件.
為耶穌基督獻上最好的敬拜態度.
其實提醒我們.
敬拜的焦點.
是向上帝付出.
而不是我今天有什麼得著.
當然.
我們不是說崇拜不期望有得著.
我們沒理由期望崇拜沒有得著.
崇拜一定會有得著.
因為得著來自上帝.
不過在今天我們說很講求消費主義.
個人主義底下.
其實我們有時會誤將有沒有得著.
成為崇拜的焦點.

$^{841}$你覺不覺得?.
我們會很在意.
今天的唱詩如何.
今天的訊息有沒有幫助到我.
我們不知不覺.
其實我們是會眾.
但變成了什麼?.
變成了觀眾.
覺不覺得?.
我們本來應該是一個敬拜者.
我們是一個獻禮的人.
但變成了一個觀眾.
什麼意思呢?.
做一個突擊測驗.
留意了.
如果作為一個敬拜者.
你要怎樣坐?.
你問怎樣坐?.
就這樣坐.
你和耶穌基督是有距離的.
祂是高高在上的君王.
我們是卑微.
你去崇拜這位大君王.
應該怎樣坐?.
這不是法理錯人主義.
是出自你那種有沒有敬畏的心.
如果你思考這件事.
每一次崇拜.
你應該知道怎樣才能夠坐.
如果你思考這件事.
每一次崇拜你應該知道.
怎樣才能夠正確坐.
有人曾經說得很好.
他說敬拜第一個姿態是前傾.
不是靠後.
因為你是敬拜那位大君王.
事實上弟兄姊妹.
我們確實很容易將焦點放在自己.
到底舒不舒適?.
好不好?.

$^{881}$有沒有記得以前崇拜.
我們常常坐很硬的長椅.
故意令你坐得不舒服.
讓你整個人坐直一點.
不能靠後.
靠後反而腰骨不舒服.
它會幫你建立一個敬虔的姿態.
來進入敬拜.
以致我們不會將崇拜的焦點.
放在事奉的人身上.
那次唱歌如果一般般.
阿威木講道又悶悶地.
怪怪地.
那怎樣好呢?.
那是不是不可以敬拜呢?.
這很值得我們去想.
其實當我們遇到崇拜.
人生百態.
有時電話響.
有時這樣那樣.
總會影響到你對上帝的專注.
但在這些時候更加考驗.
你是不是專注上帝.
當遇到這些.
你覺得不舒服的狀況的時候.
你第一時間是對準在那些人身上.
還是仍然對準在上帝身上.
這就是敬拜的真正意思.
敬拜不在於我們到底舒不舒服.
在於耶穌舒不舒服.
我們能不能夠先對準上帝.
然後為身邊的人來代禱.
這就是我們在敬拜裡面.
無論任何時候任何環境.
都可以獻情過給耶穌基督的態度.
所以弟兄姊妹其實.
敬拜的得著不在於事奉的表現.
並不是說不需要追求事奉的本意.
但敬拜的本意.
我們有沒有得著在於.

$^{921}$我們有沒有依靠上帝的靈.
有沒有依靠上帝的話語和恩典.
我們可以自由隨時來到耶穌基督的面前.
親近祂.
這本來就是敬拜的得著.
因為本來我們沒有這個機會.
是因為耶穌基督的降生.
我們才有這個福氣.
所以藉著今天.
我們去紀念耶穌基督降生的日子.
願意我們真的可以做一次審查.
我們怎樣去回應.
我們怎樣身體力行.
透過我們的付出.
透過我們的擺上.
不單單只在主日的裡面.
離開了禮堂.
之後的六天.
你怎樣在你的生活裡面.
藉著你的行為.
宣告耶穌基督是我們生命的主.
生命的王.
我們一起祈禱.
你昔日到城獄身.
來到這個世界.
為我們每一個預備救恩.
甚至當我們還不認識你的時候.
我們頂撞你的時候.
你已經親自成就這一切.
並且你在我們的生命裡.
親自來尋找我們.
讓我們有機會接觸你.
甚至認順你.
至今我們仍然活在這份恩典底下.
至今我們曾經立定心志.
要以你為王.
以你為主.
但在生活裡我們總有迷失.
我們會因為可能在生活裡的不順境.
我們很想去掌控祂.

$^{961}$漸漸我們將耶穌基督.
你本來應該坐在我們生命的王座裡.
我們將你趕走了.
我們自己坐回自己生命的寶座裡.
我們很想說事.
很想去作王.
但當我們越是這樣做的時候.
其實我們發覺我們越遠離你.
越失去那份平安.
耶穌我求你憐憫.
醫治我們每一個.
藉著你今天降生的日子.
我們再一次彼此提醒.
你要成為我們生命的主.
生命的王.
過往無論我們如何抓緊.
我們的人生.
不放也好.
盼望今天你的靈.
遊遠我們的心.
讓我們再一次歸向你.
放手完全倚靠你.
讓你完全掌管.
坐在我們生命的寶座上.
以致我們真的經歷到.
上帝你應許過給我們的豐盛.
天主上帝.
我們也求你幫助我們.
再一次重拾敬拜你的熱情.
主人我們確實有時真的會冷待了你.
我們確實有時真的沒有好好預備.
沒有珍惜每一次的主日.
來到你的面前.
我們忘記了.
這其實本來是恩典.
是我們不配.
上帝我求你.
幫我們身體力行的.
尊你為王.
尊你為大.

$^{1001}$以致我們能夠討你的喜悅.
無論在什麼情況之下.
我們首先都可以對準你.
以致我們知道.
如何回應你這份的心恩厚愛.
亦如何彼此守望.
互相幫助提醒.
我們要一起尊你為大.
為此我們同心的.
向耶穌基督你呼求.
願你幫助我們.
願你成全我們.
奉耶穌你得勝的名字來祈禱.
阿們.
謝謝.
\newpage



\section{詩篇 65:1-13}
\label{sec:iN3iqiIAXcs}
\textbf{2023年12月31日崇拜直播｜陳明泉牧師｜你以恩典作年歲的冠冕｜詩篇65\_1-13}
\newline
\newline
連結: \href{https://youtube.com/watch?v=iN3iqiIAXcs}{\texttt{https://youtube.com/watch?v=iN3iqiIAXcs}} ~~~~ 語音日期: 2024-01-01
\newline
\newline
\hyperref[sec:QUWqIbTXwuo]{\small{< < < PREV SERMON < < <}}
~
\hyperref[sec:index_chronic]{\small{[返順時目]}}
~
\hyperref[sec:2Yp8rtZS_PM]{\small{> > > NEXT SERMON > > >}}
\newline
\newline
詩篇 65:1-13
\newline
\begin{longtable}{cl}
\hline
\hline
章節 & 經文 (和合本修訂版)\\
\hline
65:1 & \begin{tabularx}{0.7\textwidth}{X} 神啊，在錫安，人都等候讚美你，也要向你還所許的願。 \end{tabularx} \\ \\ \relax
65:2 & \begin{tabularx}{0.7\textwidth}{X} 聽禱告的主啊，凡有血肉之軀的都要來就你。 \end{tabularx} \\ \\ \relax
65:3 & \begin{tabularx}{0.7\textwidth}{X} 罪孽勝了我；至於我們的過犯，你都要赦免。 \end{tabularx} \\ \\ \relax
65:4 & \begin{tabularx}{0.7\textwidth}{X} 你所揀選、使他親近你、住在你院中的，這人有福了！我們要因你居所、你聖殿的美福知足。 \end{tabularx} \\ \\ \relax
65:5 & \begin{tabularx}{0.7\textwidth}{X} 拯救我們的神啊，你必以威嚴秉公義應允我們；地極和海角遠方的人都倚靠你。 \end{tabularx} \\ \\ \relax
65:6 & \begin{tabularx}{0.7\textwidth}{X} 你既以大能束腰，就用力量安定諸山， \end{tabularx} \\ \\ \relax
65:7 & \begin{tabularx}{0.7\textwidth}{X} 使諸海的響聲和其中波浪的響聲，並萬民的喧嘩，都平靜了。 \end{tabularx} \\ \\ \relax
65:8 & \begin{tabularx}{0.7\textwidth}{X} 住在地極的人因你的神蹟懼怕，你使日出日落之地都歡呼。 \end{tabularx} \\ \\ \relax
65:9 & \begin{tabularx}{0.7\textwidth}{X} 你眷顧地，降雨使地大大肥沃。神的河滿了水；你這樣澆灌了地，好為人預備五穀。 \end{tabularx} \\ \\ \relax
65:10 & \begin{tabularx}{0.7\textwidth}{X} 你澆透地的犁溝，潤澤犁脊，降甘霖，使地鬆軟；其中生長的，蒙你賜福。 \end{tabularx} \\ \\ \relax
65:11 & \begin{tabularx}{0.7\textwidth}{X} 你以恩惠為年歲的冠冕，你的路徑都滴下油脂， \end{tabularx} \\ \\ \relax
65:12 & \begin{tabularx}{0.7\textwidth}{X} 滴在曠野的草場上。小山以歡樂束腰， \end{tabularx} \\ \\ \relax
65:13 & \begin{tabularx}{0.7\textwidth}{X} 草場以羊群為衣，谷中也長滿了五穀；這一切都歡呼歌唱。 \end{tabularx} \\ \\
[1ex]
\hline
\hline
\end{longtable}
$^{1}$今日頌讀的經文是詩篇65篇1至13節.
詩名是「讚頌主恩」.
今日是2023年最後幾日.
讓我們都藉著這首詩歌.
向我們的神獻上感恩 獻上頌讚.
詩篇65篇1至13節.
神啊 錫安的人 都等候讚美你.
所許的緣 也要向你償還.
聽禱告的主啊 凡有血氣的 都要來就你.
罪孽勝了我.
至於我們的個法 你都要赦免.
你所揀選使他親近你 住在你院中的.
這人 面為有福.
我們必因你居所 你聖殿的美福至足了.
拯救我們的神啊.
你必以威嚴秉公義 應允我們.
你本是一切地極和海上 遠處的人所倚靠的.
他既已大能蓄腰 就用力量安定諸山.
使諸海的響聲 和其中波浪的響聲.
定萬民的喧嘩 都平靜了.
住在地極的人 因你的神跡懼怕.
你使日出日落之地 都歡呼.
你眷戀地 降下透雨 使地大得肥美.
神的河滿了水 你這樣囂觀了地.
好為人 預備誤局.
你囂透地的泥溝 潤平泥脊.
降金林 使地軟和 其中發長的蒙尼似福.
你以恩典為連瑞的冠冕.
你的路徑都滴下滋油.
滴在曠野的草場上 小山以歡樂蓄腰.
草場以羊群為衣 谷中也長滿了誤局.
這一切都歡呼歌唱 聖經獨白.
弟兄姊妹平安.
平安.
2023年即將要過去.
今天我們用這篇詩篇.
來數算上帝的恩典.
在這篇詩篇裡面.
傳統很多的悉經學家都相信.
這篇詩篇是用在農作物收成之後.

$^{41}$到聖殿獻祭.
又或者在一年年中或開頭的時候.
用這首詩歌來歌頌上主的看顧和保守.
不過這篇詩篇裡面.
你讀上去.
你會覺得好像有一點點跳來跳去.
因為裡面其實是說三個不同地方的意象.
今天我們用這三個地方的意象.
以及裡面描述的聲音.
來講述一下究竟詩人想幫助我們.
用什麼樣的態度來回顧上帝的恩典.
用什麼樣的態度來敬拜上帝.
所以今天我們.
你好像看荷里活片一樣.
去完A地方之後.
我們去看B地方.
而這三個不同的意象.
就成為了我們今天詩篇裡面三個的段落.
到最後我們會總結一下.
究竟我們帶著什麼樣的心態去到上帝面前.
我們先祈禱.
求上帝幫助我們.
同心祈禱.
在過去的一年.
無論我們人生順與逆.
我們都認定你是我們生命的主.
在我們人生當中.
如果沒有了你的時候.
我們人生都不知道怎麼辦.
因為你的看顧.
你的保守.
你的話語.
成為我們生命裡面最重要的明燈.
當我們失落.
或者我們迷失的時候.
你會指教我們.
你會提醒我們.
今天我們一起來.
再一次重溫.
重讀你的話語.

$^{81}$求你再一次激勵我們.
提醒我們.
讓我們帶著合宜的心去到你的庚前.
恭敬將來的時間交託給你.
幫助每位聆聽的兄弟姐妹聽得明白.
宣講的講得清楚.
獻上我們衷心的禱告.
奉靠基督耶穌得勝的聖名祈求.
Amen.
今天我的聲音有一點.
比較厚一點.
大家忍耐一下.
我怕說著說著會走音.
我們一起看.
詩篇65篇.
第一節至第四節是第一個的意象和段落.
在這裡.
神啊!.
錫安的人都要等候讚美你.
所許的緣也要向你償還.
這是第一節.
第四節.
你所揀選使他親近你.
住在你院中.
這人便為有福.
我們必因你的居所.
你聖殿的美福.
知足了.
你見到第一節至第四節.
在說錫安.
和在說上帝的院舍.
上帝的居所.
和聖殿.
這一堆字眼其實在說同一個地方.
就是上帝的聖殿.
去到殿裡面去敬拜上帝.
詩人第一個要我們看到的意象.
他將鏡頭放在.
去到一個很宏偉的.
一個以色列人.

$^{121}$每年豐收感恩的時候.
去到這個聖殿當中.
敬拜讚美神的這個情景.
當你有豐收的日子裡面.
原則上大家會帶著歡欣雀躍的心.
興高采烈的.
甚至有時候會嘈吵.
去到這個聖殿裡面敬拜神.
不過在經文裡面.
詩人用一個我們不是很習慣的字眼.
去說敬拜.
去到這個聖殿當中的時候.
不是用嘈吵的方式.
在經文第一節裡面.
錫安的人都等候讚美你.
等候這個字是和合本的翻譯.
原文用的那個字眼.
是用沉默.
Dumiyah.
沉默這個字其實不只是在這裡出現的.
其實早一點的篇幅.
在詩篇第62篇第一節.
大家唱歌唱過的.
我的心默然無聲.
專等候神.
那個默然無聲.
默默無聲.
就是用回同一個字了.
Dumiyah.
詩人去到聖殿.
帶著一些的祭物.
帶著感恩的心.
去到聖殿當中.
他第一樣東西.
他要帶著的不是那種.
嘈圈巴閉.
不是興高采烈.
而是帶著一種肅靜.
安靜的心.
去到這個地方.

$^{161}$來敬拜神.
當然那個聖殿很宏偉.
令到人可能帶著一種莊嚴的心.
去到上帝的庚前.
不過.
經文裡面他再繼續說那個沉默.
就不是因為那個建築物有多麼宏偉.
又或者那個地方.
那個氣氛是很寧靜的.
令到大家不敢出聲.
然後在第二節到第三節就說到.
為什麼詩人會變得沉默呢?.
聽禱告的主.
凡有血氣的都要來咒你.
因為第三節裡面特別說.
罪孽性了我.
至於我們的過犯.
你都要赦免.
經文裡面他特別提到.
那個沉默.
那個敬拜者去到上帝的聖殿當中.
不出聲.
不是因為那個環境氣氛.
令到他有一種這樣的表達.
而是他想起自己.
他回想自己.
看到他自己的生命的真章.
看到自己生命的真相.
他發現.
他是充滿了罪的.
在這裡面說的那個罪孽的字.
其實不是只是說一件事.
或者那件罪行是什麼.
那個罪孽是說一個.
籠罩著他好像一個權勢一樣.
他在江湖裡面身不由己.
在他的生活當中.
或多或少.
他未必想的.
但是好像這個世界就有一種洪流.

$^{201}$有一種壓力.
令到他不知不覺.
又或者不情願之下.
就做了一些上帝不起悅的事情.
所以那個罪孽性的我.
是在說.
原來他活在一個世界裡面.
是敵黨神的.
是一個他自己不由自主.
甚至他自己想做好.
但是沒有能力.
以致被那種罪惡的權勢勝過.
那種面對著這個聖殿敬拜的時候.
他心裡面有一種慚愧.
有一種感受.
覺得很不配的那種經歷.
所以在這裡面使人.
去到上帝的聖殿當中.
他不出聲.
他安靜.
他那種安靜.
大致都是檢視自己生命裡面的狀態.
我不知道大家是否同意.
當然在聖經裡面.
很多時候都會說到.
我們要鼓失彈琴.
高聲地讚美上帝.
但是原來去到某一時的讚美.
去到某一些時間.
你發覺你不敢唱.
你未必是因為覺得戴得不好.
或者音樂等等的客觀條件.
而是當你在讚美的時候.
你想聽一下.
甚至你想靜一下.
去看看自己裡面發生了什麼事.
人大了.
我想我自己可能在座裡面.
有些男士都是這樣的.
人大了的時候.

$^{241}$那種沉默.
帶著的是一種.
恭敬的心去到上帝的庚前.
而這種沉默不是不合作.
不是和大家慶高采烈唱詩的時候.
在唱反調.
而是在說.
自己檢視到自己心靈的狀態.
你發現自己.
你要看自己的時候.
你想靜下來.
好好體會敬拜是一回什麼的事.
這種心態其實是一種很深度的經歷.
在新約裡面.
耶穌基督都用過這樣的比喻.
來去說.
敬拜人的那種狀態.
路加福音第十八章九至十四節.
耶穌就曾經用過聖殿的敬拜.
說兩班人的比喻.
來去說一個這樣的狀況.
耶穌在第十八章第九節.
耶穌向那將著自己是義人.
藐視別人.
就設一個比喻.
這個是說那些法利賽人的.
一眼兩面的.
說那些人.
覺得自己.
嘩.
很大聲.
很演戲的.
他就用這個比喻來去.
指他們的生命的盲點.
在這裡說.
有兩個人上殿去禱告.
一個是法利賽人.
一個是瑞利.
然後經文我省略了.
瑞利和法利賽人的禱告不同之處就是.

$^{281}$法利賽人在這裡數算自己.
做得多麼好.
一個星期他禁食兩次.
還有周濟窮人.
他又不停的去敬拜上帝.
還要說我不能侮辱某些人.
只是在敬拜神而已.
不過他好像為自己.
寫下一個在上帝面前.
稟報自己有多麼好.
這個是法利賽人.
至於那個瑞利.
耶穌特別說.
第十三節.
那瑞利遠遠的站著.
說的是距離.
連舉目望天也不敢.
只隨著空說.
神啊開恩可憐我這個罪人.
我告訴你們.
這人回家.
被那人倒算為異了.
因為凡至高的必降為卑.
至卑的必升為高.
在經文裡面.
其實耶穌基督要說的是一個敬拜.
去到聖殿當中.
去到一個神聖的地方.
去到敬拜上帝.
不過你帶著什麼心態.
那件事是很重要的.
不是一定要唱詩唱得很大聲.
或者祈禱很亮麗.
甚至數算自己做得很理想的地方.
那個禱告就蒙到月立了.
耶穌基督反而在這裡面說.
那個願意揭開自己的心.
看到自己生命的不時.
他以罪人自稱.
以罪孽壓過自己.

$^{321}$在生命裡面.
他發覺自己很輕忽的時候.
這種心靈的狀態.
他向上帝表達的時候.
耶穌說反而這個才稱為真正的義人.
宗教聚會崇拜不會盛化我們自己的生命.
我們來到這裡當中.
不是因為我們自覺自己做得很好.
而是我們發現我們自己的生命很需要上帝.
我們生命其實都是破破爛爛的生命.
有意無意都好.
這個世界好像一個籠罩著的權勢.
令到我們在生活當中.
被這個文化或者生命裡面.
過去有一些的積雜,遺憾.
以致我們沒有辦法好好地去踏前.
斯人他帶著一份這樣的心.
去到聖殿當中.
來敬拜上帝.
而這種敬拜是帶著默然無聲.
這種沉默也是敬拜的一種.
所以弟兄姊妹.
在這裡斯人的起點.
這裡去到每一年的感恩.
他是帶著檢視自己生命.
重新來到發現自己生命的真相.
去到上帝的殿中來作為敬拜.
這個場合.
這個教會的敬拜.
我們都是美輪美奐.
我們都很用心來預備的.
不過神聖的詩歌不會將我們聖化的.
我們說的話說得很漂亮.
都不會為我們自己的生命漂白的.
我們來到上帝的殿裡面.
我們只是將我們自己破爛不堪的生命.
呈現在上帝面前.
向上帝傾心吐意.
將我們自己生命的問題.
將我們自己生命的真相向上帝表明.

$^{361}$這個是敬拜的起點.
也是敬拜的第一步.
我看過一個在網上的故事.
那個故事是講到有一對兄弟.
這對兄弟每個星期都會去教會.
他們去教會的目的.
就是因為這兩個兄弟.
屬於一種掛名的信徒.
但是他們很勤力的.
雖然他們是掛名信徒.
但是他們希望藉著這些宗教聚會.
令他們自己給人感覺上是一個好的人.
但是離開了崇拜.
離開了禮拜堂之後.
其實他們的私德是很敗壞的.
總之壞事作盡.
但是每個禮拜他們都會回來.
而且都很慷慨.
願意給一點香油錢.
做一些善事.
來做一些奉獻.
來給教會.
這兩個兄弟其中一個弟弟因病去世了.
去世之後.
這個哥哥在安息禮拜的安排裡.
就跟主禮的牧師說.
牧師,在接下來的安息禮拜.
請你幫我一個忙.
在安息禮拜裡.
請你盡量把我的弟弟說成是一個聖人.
然後就說好好照顧他.
我希望你用你的講章.
整個事件聚焦在他身上.
唱好他.
讓他帶著聖人.
好人的身份離開這個世界.
這個牧師就答應了他.
當這個牧師答應了的時候.
這個哥哥.
有錢的哥哥就說.

$^{401}$當你答應了的時候.
教會的裝修包在我身上.
如果有人這樣說的話.
我都會答應.
接著這個牧師就預備了安息禮拜的第二天.
在第二天的時候.
在安息禮拜進行的時候.
這個牧師就一直說訊息.
一直說的時候,說到尾段.
就說死去的這個弟弟.
其實是一個壞事做盡的人.
他的私德甚為敗壞.
而且他是一個拋妻棄女.
一個很差的人.
壞事做盡.
一直數數數數.
當這個牧師一直說的時候.
那個哥哥臉都青了.
這樣都可以,答應了我,說好事.
唱好我.
接著下面其他那些人.
都沒有去過一個這麼真誠的喪禮.
說那個人的不是.
說完所有的東西.
一直數數數.
最後的時候.
這個牧師就加他一句話.
說死去的這個人,壞事做盡.
他在上帝面前.
上帝如果他有信心.
信主,他都是一個義人.
接著他再加一個柱腳.
還有相對於他的哥哥.
他是一個聖人.
明白我說什麼了嗎?.
他答應了他哥哥做的事.
他又說他是聖人.
答應了他哥哥所說的.
不過那個真相不是一個.
要漂白,唱好的一個真相.

$^{441}$同樣的,兄弟姐妹.
在這裡裡面.
私人他想我們.
真的去到上帝面前.
感恩敬拜.
你要帶著一份真誠.
今天坊間裡面所說的真誠.
很多時候就是在說我的感覺.
我的感覺.
那個就是我的看法為之真誠.
不過聖經裡面所說的真誠.
從來不是這些東西.
今天這個世代.
在說的那種真誠.
跟聖經裡面所說的真誠是完全兩碼子的事.
今天聖經裡面所說的真誠.
是要透視你的心靈空間.
看到你的罪才看到真誠.
看不到你自己的罪.
你不去檢視你自己生命破爛的地方.
所說的那些真誠.
其實都是沒有意思的.
那些是感覺來的.
求主幫助我們.
私人在這裡裡面去強調.
他自己默然無言.
他自己默然無聲.
他自己看到自己生命的真相.
那種沉默也是敬拜的一部分.
經文再繼續說.
這個私人除了沉默的敬拜.
他更加說到所許的願也要向你償還.
所許的願是什麼意思呢.
那個還願許願是什麼呢.
其實這裡是說.
其實這裡所說的許願.
或者當中所用的字眼.
和民數記裡面所說的贖罪祭.
那些字眼是一致的.
而這個許願不是純粹說.

$^{481}$上帝啊求你幫我.
來年順風順水.
我生意興隆.
或者我所做的事求你讓我順利.
這裡所許的願.
大致是說贖罪祭的許願.
求上帝赦免了他.
在未來的日子.
或者在他人生裡面.
得蒙上帝的月亮.
而那個美福.
在殿裡面享受的那種福氣.
或者在說.
那些美物.
在聖殿裡面.
聖殿的美福.
等等這些字眼.
其實所用的字眼.
其實是描述.
聖殿裡面的祭祀.
那個祭祀之後獻祭的那塊祭肉.
贖罪祭的那塊祭肉.
很多時候被稱為.
一塊美好的.
這些肉.
而這塊肉不會燒掉.
在聖殿當中.
祭祀之後那塊祭肉.
只能夠.
祭司可以在聖殿裡面.
來享用.
所以整件事是說.
帶著感恩.
也帶著上帝的祈求.
求上帝赦免他的罪孽.
他的許願.
也帶著一份.
求上帝赦免他的罪險.
免去他地上的債.
當上帝.

$^{521}$閱納了這個憲制的時候.
即使祭司享受.
這一塊.
那個祭肉.
那個私人憲制.
看到這個狀況.
他心裡面都感到安慰.
所以第一個.
上帝的敬拜.
其實是訴諸於指向心靈.
真實的那種.
你罪你自己生命的本相.
帶著這個態度.
來敬拜上帝.
這是第一個意象.
去到第五節至第八節.
就是第二個意象.
鏡頭一轉就轉到哪裡呢.
拯救我們的神.
你必以威嚴秉公義.
應允我們.
你本是一切地極和海上.
遠處的人所依靠的.
祂既以大能速邀.
就用力量安定諸山.
使諸海響聲和其中波浪的響聲.
並萬民的喧嘩都平靜了.
住在地極的人.
都因你的神跡懼怕.
你使日出日落之地.
都歡呼.
在第五節至第八節.
在第二個段落裡面.
鏡頭轉了去.
由聖殿轉到.
像一個懸崖望向大海.
一個很壯闊.
海浪聲拍岸的.
一個像.
一個像崖邊.

$^{561}$看著大海.
看著諸山的大自然.
的一種景地.
在這個景象裡面.
有兩個字眼需要大家注意.
第一個字眼就是「懼怕」.
「懼怕」的字.
除了在第八節.
其實在第五節.
「易」的字.
「必以威嚴」.
「威嚴」的字.
其實隱含著.
「可畏的事」.
上帝做了一些令人害怕的.
那種權柄.
「威嚴」就是這個意思.
所以一頭一尾.
是在說上帝帶你.
去看大自然.
去看到大自然的壯闊.
更加看到背後那位.
充滿力量.
令人望而生畏的.
上帝.
這個令人望而生畏的上帝.
當人.
見到他的時候.
重要的象聲詞就是.
「平靜」「安靜」.
剛才那裡說的那個沉默.
和這個安靜是兩個字.
不過都是在說.
一開始是看到自己的罪.
不說話.
而在這裡看到上帝的偉大.
看到上帝的能力.
所有的喧嘩.
靜止了.
不過不單止靜止.

$^{601}$到最後在日出日落之地.
眾民.
欣著這種敬畏的心.
又再歡呼起來.
整個的意象.
告訴你什麼呢.
這個珠山大海的意象.
其實讓你看到.
大自然的巨大.
和那種.
人力無法勝過.
大自然.
人很渺小的一種狀況.
當那些波濤洶湧.
這個象聲詞眼.
他就和那個萬民喧嘩.
是放在一起的.
很多人說大自然很安靜.
其實大自然一點都不安靜.
那些分貝吵得不得了.
那個.
拍案的啪啪聲.
詩人就用這些啪啪聲.
來等同於.
世界外族.
外邦萬民喧鬧的聲音.
嘈喧巴閉.
那些喧鬧是什麼聲音呢.
各自為政.
在人生裡面混亂.
為著他們.
所追求的東西叫囂.
又或者為著.
他們站在.
道德的高台裡面去審判.
這個世界的道德譴責的聲音.
五花八門.
不同的聲音.
在這個世界裡面充斥著.
詩人就用這些.

$^{641}$波濤洶湧的.
拍案的聲音.
來跟萬民的喧嘩.
都是一樣嘈喧巴閉.
但當上帝的出現.
上帝.
公義的實踐.
秉公義.
英運我們.
無論近處或者遠處.
地極或者.
在詩人所處的.
境地.
一齊見到上帝的作為.
見到大而可畏的上帝.
的時候.
所有人就安靜了.
就害怕了.
上帝的神蹟.
上帝的工作.
不是純粹像一個魔術師般的.
很可笑的.
過來看看.
新約裡面有人見到耶穌.
做這些神蹟.
就叫耶穌多給我一個神蹟.
耶穌說若拿的時代的神蹟.
如果你們都不承認.
這個世界沒有神蹟.
真正的神蹟是上帝的作為.
叫人悔改.
讓人生命得著改變.
當你看到上帝做事的時候.
你第一個反應應該是害怕.
如果你帶著一種.
興奮的心態.
大概是看表演.
你就失去了經文裡面.
詩人告訴你的.
你要肅然起敬的心態.

$^{681}$所以去到第二個景象.
詩人帶著感恩的心.
去到敬拜上帝面前.
那種敬拜.
你要帶著的.
就是一份敬畏.
你要害怕上帝.
今日我們的文化裡面.
我們喜歡將上帝描述成.
一個朋友.
當上帝.
或者英雲了你的時候.
你就跟他很友好.
當上帝有些地方不是按你心意的時候.
你會示意地.
批評,罵他.
不過在經文裡面說.
我們所信的上帝.
不是這樣的上帝.
不是你手裡的燈神.
也不是你的寶可夢.
他是.
在你生命裡面需要敬畏.
他是唯一敬畏的對象.
你要擺好.
那個焦點.
你的敬拜如果沒有了一份敬畏的心的時候.
大概.
你會將自己放大.
你會覺得自己的需要.
你會覺得自己是.
這個世界宇宙的中心.
但當你看到.
上帝的作為.
你將焦點放在上帝身上的時候.
你就會慢慢發現.
其實這個世界裡面一粒微塵都不如.
你只是看到自己的渺小.
敬拜.
是看到上帝的偉大.

$^{721}$看到人自己的渺小.
而這個上帝.
是一個.
不單止是說基督徒.
自己的上帝.
也是這個全世界的上帝.
上帝是一些未認識主的人的主.
不過那些人.
不知道而已.
他們不知道他們的生命有一個主宰.
他們有一個主宰.
不過在這裡講到.
日出日落之地都要歡呼.
住在地極的人.
都要因你的神蹟懼怕.
全部在講.
這個上帝是掌管宇宙萬物的上帝.
是一個要你.
驚懼一個的上帝.
今天我們對上帝的敬畏有多少呢.
我們今天帶著什麼心態.
來敬拜我們天上的這位主呢.
我們今天帶著什麼心態.
來敬拜我們天上的這位主呢.
我們今天帶著什麼心態.
來敬拜我們天上的這位主呢.
我記得我想這段經文的時候.
想起我以前年輕的時候.
回到團契.
我們唱一些詩歌.
剛才說的.
我的心默默無聲.
專等候神.
我們特意玩玩.
唱這首.
我的心默默無聲.
專等候神.
特意唱到很大聲.
我的心默默無聲.
唱到有點不順.

$^{761}$詩歌是要這樣唱的.
以前是這樣玩的.
然後就被導師.
馬上罵.
罵什麼呢.
他說Anders.
或者這麼多位年輕的弟妹.
你們沒想過你們平時唱的詩歌.
你們知不知道.
裡面在唱什麼.
你唱到這麼大聲.
你唱到朗朗上口.
是不是你心裡面.
所認同.
所相信的信念.
然後我們就開始.
有點害怕.
被團契導師.
板起面來質問.
然後他說.
其實不只是你們.
今天很多人也是.
我們唱詩歌.
將那些經文.
或者將那些內容唱到朗朗上口.
配以最好的音樂.
和最大的聲音.
你不是這樣想的.
你的心不是這樣想.
但你就朗朗上口.
當時代曲這樣唱.
你有沒有發覺有些問題呢.
這句話.
到今天還在影響我的內心.
今天我們.
每個星期.
每次聚會或者洗澡的時候.
我們在唱的那些歌.
除了好聽之外.
那個是不是.

$^{801}$我們心靈裡面所確信的信念呢.
那個是不是.
我們保持的那種態度呢.
可能我們未必做得到.
可能我們和詩歌裡面.
所講的內容我們還有點距離.
但至少我們的生命.
不是朝向那個方向.
我這樣講這番說話.
我不是要打殘大家.
大家不要誤會.
如果我說一些負面的話.
醒惡部就殺了我.
我的目的不是叫大家不唱歌.
而是我們每一次所唱的.
我們是不是帶著一份敬畏的心.
我們是不是尊重.
詩歌裡面所講的內容.
敬拜上帝.
如果失去了那份敬畏.
大概你只是感受良好的.
一種經歷.
我們經歷不到敬拜.
真正的意義.
詩人要將我們拉到大自然.
我們要看到.
大自然裡面.
創造自然背後.
那位滿有威嚴的上帝.
那個大而可畏.
令我們懼怕的上帝.
敬拜.
我們需要敬畏.
沒有了敬畏.
我們很難真正去敬拜.
創造天地萬物的主.
我們看第三個段落.
九至十三節.
你眷顧地降下透雨.
使地大得肥美.

$^{841}$臣的河滿了水.
你這樣囂冠了地.
好為人預備五穀.
你曉透地的泥溝.
潤平泥脊.
降金林.
使地軟和.
其中發長的蒙你賜福.
你以欣典為年歲的冠冕.
你的路徑都敵下之尤.
敵在曠野的草場上.
小山以歡樂束腰.
草場以羊群為衣.
谷中也掌滿了五穀.
這一切都歡呼歌唱.
這裡最後一個段落.
也是在說.
他來到聖殿敬拜的時候.
的核心.
為了感恩.
為了修成的充足.
後面所用的字眼.
是很美的字眼.
很有詩意.
你看到上帝降下雨水.
風調雨順.
上帝令到河滿了水.
囂冠了地.
而且.
風調雨順.
地又軟和.
土不是硬土可以耕作.
一切都有上帝的賜福在當中.
看到這一堆字眼.
都是說到.
修成滿有.
祝福的.
意象.
詩人第三個意象.
是一個豐收的意象.

$^{881}$而這個豐收的意象.
不只是說自己擁有很多.
豐收的意象.
多到一個地步.
是滿到寫的.
甚至乎不關事的地方.
上帝都賜福.
這個就是豐收.
豐到滿日.
多出來的那種意象.
所以看到降雨.
恩典,美物.
這些一切都是上帝的工作.
經文裡面.
他沒有特別說到.
我們從小讀書到大都知道.
農夫勤勞耕作.
我們經常都會說這些.
詩人從來沒有否定農夫的努力.
不過經文的重點不是說農夫有多努力.
所以他們有多農作物.
那個修成有多豐盛.
他是在說.
因為上帝的賜福.
因為上帝的供應.
在你的努力之下.
你會看到一些人力以外.
沒有辦法做到的.
那些福氣淋到.
他眼前的那種境況當中.
值得注意有幾個字眼.
這裡稍為讓大家知道.
第十一節裡面.
你以恩典為年歲的冠冕.
你的路徑都敵下之尤.
這句很多時候有些誤解的.
很多頂子妹.
首先說恩典的字.
其實那個恩典的字.
是用你的美麗.

$^{921}$上帝的你.
上帝好東西.
上帝將好東西給你.
和合本譯恩典.
不是譯錯的.
方向是對的.
不過譯不到的那種感覺是.
這是上帝給你的好東西.
你見到你會心花怒放.
那個就是恩典.
你的美麗的意思.
第二樣就是那個冠冕.
冠冕的字不是說希臘.
或者說新約裡面.
說菜刨裡面.
戴著的那個冠冕.
他這裡說的那個冠冕是說.
豐收滿到寫.
那種榮耀.
一個滿有收成.
一個農夫見到整車都是農作物.
那個是一個.
他自己的榮耀.
那個冠冕是說.
這個滿有收成.
那種的狀況.
那個路徑.
其實不是說那條路.
那個路徑是說.
那個.
座駕那架車.
運著糧食.
運著.
那些農作物的那架座駕.
那架車.
一路行的時候.
有東西掉出來.
有什麼掉出來.
我們以為滴下之油.
明明農作物會掉油出來.

$^{961}$肥lum lum 地滑死人.
其實原文.
不是強調那個油.
是強調.
滿到寫那些農作物.
一路行就一路掉一些出來.
一路多到一個這樣的地步.
那架車都載不完.
而且邊行就邊掉.
但是在說多到一個.
那麼豐收的那種狀況.
詩人用一個.
滿有豐收的.
那種意境.
說到生命真是豐富.
不單止這樣.
一路行一路掉一些.
農作物出來.
然後經文再繼續.
來到去說.
滴出來的.
滴到去曠野的草場上.
小山以歡樂蓄腰.
草場以羊群為衣.
谷中掌滿了五谷.
這一切都歡呼歌唱.
說到那個掉出來.
掉下去.
掉滿了那些.
其他曠野的田地.
沒有人在曠野裡面.
來種植的.
曠野是一些瘦土.
那些是化外之地.
基本上那些沒有人.
耕作的那些地方.
那些農作物.
那些的收成.
欣欣向榮的那種景象.
都充滿了.

$^{1001}$舊時很瘦弱的那些曠野之地.
而且那個地方.
掌滿了青草.
青草上就有滿佈.
所有羊群在那裡餵食.
嘩!你看到.
豐收.
豐收到一個地步.
超乎一般人的想像.
而且那個豐收不是自己的車.
而是整片的山頭.
他所見到眼前的景象.
全部都是欣欣向榮.
這一切的豐收.
就好像歡欣歌唱.
歌頌著上帝的榮美.
在一個這麼豐收的景象裡面.
詩人他想告訴我們.
我們看到.
生命裡面的豐富.
你要感恩.
生命裡面.
其實你的恩典是滿到寫的.
你的生活一無所缺.
今日我這樣講.
香港這個社會.
大家都擔心經濟.
又說.
旅遊又好像.
去外遊多過去香港.
又或者我們經常覺得.
我們自己.
就好像.
供不起樓等等.
有一種很大的恐怖感.
籠罩著你的內心.
你未經歷貧窮.
你先心裡面.
經歷那種貧窮的恐慌.
不過鄧英姐妹我想講.

$^{1041}$我們多麼害怕都好.
我們香港.
我們這個地方.
我們反而要講真話.
我們其實今日的生活.
是全世界頭的5\%.
你說差一點的.
頭的7\%.
全世界最富裕的頭的7\%.
那些人的生活.
就是我們今日香港人的生活.
我們怎樣窮都好.
我們都是一個.
豐足的那種生命.
社會上都依然有一些缺乏的人.
我們相對貧窮的.
是對的.
不過我們.
未去到那種.
極致的貧窮.
或者去到一個足襟見爪的.
那種處境.
又或者其實有時候.
富足和貧窮.
可能是在乎你自己.
怎樣看你自己所擁有.
有很多人在這個.
世界裡面是一個.
擁有很多但依然覺得.
自己擁有很少的一班人.
他們生活.
富足的.
但心靈是極之貧窮.
窮到一個地步.
害怕到一個地步.
經常都希望.
拿多一點錢.
拿多一點東西.
為自己製造更多的安全感.
這些本身就是富足的貧者.

$^{1081}$生活在富足當中.
但心靈是貧窮的.
但我亦都見到.
更加多的是.
其實不是很富裕.
但心靈裡面.
他能夠有那種富足.
他富足到一個地步.
他願意和別人來分享.
他願意.
看到上帝是一個傾福的上帝.
他的生命就流露了.
一種向人傾福的.
那種將福氣.
分享給別人的一種生命.
早幾日聽到一個見證.
我都聽到很亂.
話說教會裡面.
有一位姐妹.
現在在讀神學.
在神學院裡面派她去到.
教會裡面.
做一些輔導的實習.
她接完這個案件.
見完這個實習.
個案之後.
心裡面很沉重.
因為都是很痛苦的一個經歷.
一路想起這個案件.
就很沉重.
一路回家的時候.
想著這個案件.
令到自己可以放開一點.
但久久都不能夠釋懷.
於是她就去了一間便利店裡面.
她就想.
不如這樣吧.
今天這麼沉重.
買一支雪條來獎勵一下自己.
最冷的那幾天吃雪條.

$^{1121}$想買一支貴價的雪條.
寒天.
買一支貴價的雪條.
三十多元一支.
都很貴.
想獎勵一下自己.
看著雪條在選擇什麼口味.
打算買的時候.
接著.
便利店的收銀員就跟.
這個輔導員說.
「鼎姐妹,我今天請你吃,你選吧」.
便利店的收銀員說請她吃.
她都覺得,你為什麼無端請我吃.
推了幾次.
收銀員說我看你經常.
都在買那些平價.
一包包裝請人吃.
今天我請你吃.
於是.
當她選了口味之後.
就拿了這支.
請她吃.
三十多元的雪條.
這個姐妹還拍了一張照片給我看.
她說,真是值得回味.
我聽完這個故事之後.
我心裡面.
我第一個問題問.
這個收銀員.
其實他.
大概都是收最低工資.
說的那支雪條.
大概等於他一個小時.
已經達到了.
其實他用了.
他一個小時的工資.
請這個.
素未謀面,不認識的,只是熟口熟面.
這個人.

$^{1161}$知道他一直都是祝福人.
他就送一件東西給她.
這個人不是基督徒.
不是什麼的,不過.
他在一個.
很需要.
散心.
舒暢的時候.
這個人就像天使一樣.
這個收銀員.
其實他都不是一個很富裕的人.
但是他願意.
將他擁有的.
跟別人分享.
今天在我的生命裡面.
其實我們.
很多人都說2023年.
很多是不容易走的一年.
生活總是有.
一些困難.
有些不同的逆流.
不過.
多麼困難都好.
你總會見到.
恩典.
多麼.
不信心都好.
你會在這些困難當中.
上帝會默默.
透過不同的人來祝福你.
如果我們只是帶著.
我生命裡面.
很缺憾,上帝你欠了我.
我們敬拜不到.
因為我們常常覺得上帝欠我們.
但當我們轉換.
我們的焦點.
我們知道人生總是.
會有上有落.
有順有逆.

$^{1201}$但無論什麼情況都好.
上帝總是看著我.
我生命裡面我認定總有恩典的時候.
你會發現恩典處處.
多麼難走的路.
都可以見到有.
滴下滋柔的路徑.
多麼不容易.
多麼缺乏的日子.
你總會見到.
草場遍滿了羊群.
供應是充足的.
弟兄姊妹.
當我們心靈裡面.
有這一份感恩.
發現上帝供應的.
視野.
我們整個人生.
都很不同.
我們做人會幸福很多.
我們的敬拜.
也都到位很多.
如果我們能夠.
懂得感恩的時候.
我們所唱的詩歌.
就能夠口唱心和.
當我們懂得感恩的時候.
我們發覺.
生命裡面.
無論經歷低與高都好.
我們覺得我們的肩膀.
是很舒坦.
我們的骨可以挺直.
我們的眼目可以向上.
來去舉望.
我們人生可以不再.
頭低低地過我們自己的人生.
今日.
這一篇的詩篇六十五篇.
詩人帶我們看到三個景象.

$^{1241}$從聖殿的景象.
看到自己的罪.
敬拜是要看到自己生命的本相.
看到自己的罪.
明白自己.
體會自己.
我們的敬拜.
我們才帶著一份恭敬.
敬拜也都告訴我們.
我們需要敬畏.
我們敬畏的對象是向上.
那個是.
我們昭然的上主.
我們要帶著最大的敬畏.
來去面對他.
他不是一個輕忽.
我們不可以輕忽.
不可以隨意任意而為.
的一個敬拜的對象.
敬拜也都提醒我們.
帶著感恩的心.
無論日子過得怎樣也好.
帶著感恩的心.
你會看到上帝的作為.
在你人生當中.
當我們帶著這樣的經歷的時候.
帶著這個詩篇.
幫我們去默想的進程的時候.
你的人生你會覺得很開闊.
求主繼續的來致福我們.
今天晚上你帶著這個詩篇.
臨睡之前.
在現場.
也都祈願我們有更加精彩的.
2024年.
願主賜福你們中人.
主席.
\newpage



\section{創世記 18:16-33}
\label{sec:2Yp8rtZS_PM}
\textbf{2024年1月7日崇拜直播｜陳明泉牧師｜十個深化關係的禱告-為別人懇求｜創世記18\_16-33}
\newline
\newline
連結: \href{https://youtube.com/watch?v=2Yp8rtZS-PM}{\texttt{https://youtube.com/watch?v=2Yp8rtZS-PM}} ~~~~ 語音日期: 2024-01-08
\newline
\newline
\hyperref[sec:iN3iqiIAXcs]{\small{< < < PREV SERMON < < <}}
~
\hyperref[sec:index_chronic]{\small{[返順時目]}}
~
\hyperref[sec:H22XGgjIH6M]{\small{> > > NEXT SERMON > > >}}
\newline
\newline
創世記 18:16-33
\newline
\begin{longtable}{cl}
\hline
\hline
章節 & 經文 (和合本修訂版)\\
\hline
18:16 & \begin{tabularx}{0.7\textwidth}{X} 三人從那裡起程，面向所多瑪觀望，亞伯拉罕與他們同行，要送他們一程。 \end{tabularx} \\ \\ \relax
18:17 & \begin{tabularx}{0.7\textwidth}{X} 耶和華說：「我所要做的事豈可瞞著亞伯拉罕呢？ \end{tabularx} \\ \\ \relax
18:18 & \begin{tabularx}{0.7\textwidth}{X} 亞伯拉罕必要成為強大的國；地上的萬國都必因他得福。 \end{tabularx} \\ \\ \relax
18:19 & \begin{tabularx}{0.7\textwidth}{X} 我揀選他，為要叫他命令他的子孫和後代家屬遵行耶和華的道，秉公行義，使耶和華所應許亞伯拉罕的話都實現了。」 \end{tabularx} \\ \\ \relax
18:20 & \begin{tabularx}{0.7\textwidth}{X} 耶和華說：「所多瑪和蛾摩拉罪惡極其嚴重，控告他們的聲音很大。 \end{tabularx} \\ \\ \relax
18:21 & \begin{tabularx}{0.7\textwidth}{X} 我要下去察看他們所做的，是否真的像那達到我這裡的聲音一樣；如果不是，我也要知道。」 \end{tabularx} \\ \\ \relax
18:22 & \begin{tabularx}{0.7\textwidth}{X} 二人轉身離開那裡，往所多瑪去；但亞伯拉罕仍然站在耶和華面前。 \end{tabularx} \\ \\ \relax
18:23 & \begin{tabularx}{0.7\textwidth}{X} 亞伯拉罕近前來，說：「你真的要把義人和惡人一同剿滅嗎？ \end{tabularx} \\ \\ \relax
18:24 & \begin{tabularx}{0.7\textwidth}{X} 假若那城裡有五十個義人，你真的還要剿滅，不因城裡這五十個義人饒了那地方嗎？ \end{tabularx} \\ \\ \relax
18:25 & \begin{tabularx}{0.7\textwidth}{X} 你絕不會做這樣的事，把義人與惡人一同殺了，使義人與惡人一樣。你絕不會這樣！審判全地的主豈不做公平的事嗎？」 \end{tabularx} \\ \\ \relax
18:26 & \begin{tabularx}{0.7\textwidth}{X} 耶和華說：「我若在所多瑪城裡找到五十個義人，我就為他們的緣故饒恕那整個地方。」 \end{tabularx} \\ \\ \relax
18:27 & \begin{tabularx}{0.7\textwidth}{X} 亞伯拉罕回答說：「看哪，我雖只是塵土灰燼，還敢向主說話。 \end{tabularx} \\ \\ \relax
18:28 & \begin{tabularx}{0.7\textwidth}{X} 假若這五十個義人少了五個，你就因為少了五個而毀滅全城嗎？」他說：「我在那裡若找到四十五個，就不毀滅。」 \end{tabularx} \\ \\ \relax
18:29 & \begin{tabularx}{0.7\textwidth}{X} 亞伯拉罕又對他說：「假若在那裡找到四十個呢？」他說：「為這四十個的緣故，我也不做。」 \end{tabularx} \\ \\ \relax
18:30 & \begin{tabularx}{0.7\textwidth}{X} 亞伯拉罕說：「求主不要生氣，容我說，假若在那裡找到三十個呢？」他說：「我在那裡若找到三十個，我也不做。」 \end{tabularx} \\ \\ \relax
18:31 & \begin{tabularx}{0.7\textwidth}{X} 亞伯拉罕說：「看哪，我還敢向主說，假若在那裡找到二十個呢？」他說：「為這二十個的緣故，我也不毀滅。」 \end{tabularx} \\ \\ \relax
18:32 & \begin{tabularx}{0.7\textwidth}{X} 亞伯拉罕說：「求主不要生氣，我再說一次，假若在那裡找到十個呢？」他說：「為這十個的緣故，我也不毀滅。」 \end{tabularx} \\ \\ \relax
18:33 & \begin{tabularx}{0.7\textwidth}{X} 耶和華與亞伯拉罕說完了話就走了；亞伯拉罕也回到自己的地方去了。 \end{tabularx} \\ \\
[1ex]
\hline
\hline
\end{longtable}
$^{1}$(片段).
哪怕那些人跟我們素未謀面,不相識,但是我們要保守我們的心,讓我們的心不要變得盲目..
這種禱告是會拉近我們跟世界的距離,也拉近我們跟神之間的距離,求主繼續幫助我們!.
多謝您!.
\newpage



\section{尼希米記 1:1-11}
\label{sec:H22XGgjIH6M}
\textbf{2024年1月28日崇拜直播｜陳冠成牧師｜十個深化關係的禱告-以神國優先｜尼希米記1\_1-11}
\newline
\newline
連結: \href{https://youtube.com/watch?v=H22XGgjIH6M}{\texttt{https://youtube.com/watch?v=H22XGgjIH6M}} ~~~~ 語音日期: 2024-01-28
\newline
\newline
\hyperref[sec:2Yp8rtZS_PM]{\small{< < < PREV SERMON < < <}}
~
\hyperref[sec:index_chronic]{\small{[返順時目]}}
~
\hyperref[sec:_RD54XIpk84]{\small{> > > NEXT SERMON > > >}}
\newline
\newline
尼希米記 1:1-11
\newline
\begin{longtable}{cl}
\hline
\hline
章節 & 經文 (和合本修訂版)\\
\hline
1:1 & \begin{tabularx}{0.7\textwidth}{X} 哈迦利亞的兒子尼希米的言語如下：亞達薛西王二十年基斯流月，我在書珊城堡中。 \end{tabularx} \\ \\ \relax
1:2 & \begin{tabularx}{0.7\textwidth}{X} 那時，我有一個兄弟哈拿尼，同幾個人從猶大來。我問他們那些被擄歸回、剩下殘存的猶太人和耶路撒冷的情況。 \end{tabularx} \\ \\ \relax
1:3 & \begin{tabularx}{0.7\textwidth}{X} 他們對我說：「那些被擄歸回剩下的餘民在猶大省那裡遭大難，受凌辱；耶路撒冷的城牆被拆毀，城門被火焚燒。」 \end{tabularx} \\ \\ \relax
1:4 & \begin{tabularx}{0.7\textwidth}{X} 我聽見這話，就坐下哭泣，悲哀幾日，在天上的神面前禁食祈禱， \end{tabularx} \\ \\ \relax
1:5 & \begin{tabularx}{0.7\textwidth}{X} 說：「唉，耶和華－天上大而可畏的神，向愛你、守你誡命的人守約施慈愛的神啊， \end{tabularx} \\ \\ \relax
1:6 & \begin{tabularx}{0.7\textwidth}{X} 願你睜眼看，側耳聽你僕人今日晝夜在你面前，為你眾僕人以色列人的祈禱，承認我們以色列人向你所犯的罪；我與我父家都犯了罪。 \end{tabularx} \\ \\ \relax
1:7 & \begin{tabularx}{0.7\textwidth}{X} 我們向你所行的非常敗壞，沒有遵守你吩咐你僕人摩西的誡命、律例、典章。 \end{tabularx} \\ \\ \relax
1:8 & \begin{tabularx}{0.7\textwidth}{X} 求你記念所吩咐你僕人摩西的話，說：『你們若犯罪，我就把你們分散在萬民中； \end{tabularx} \\ \\ \relax
1:9 & \begin{tabularx}{0.7\textwidth}{X} 但你們若歸向我，謹守遵行我的誡命，你們被趕散的人雖在天涯，我也必從那裡將他們召集回來，帶到我所選擇立為我名居所的地方。』 \end{tabularx} \\ \\ \relax
1:10 & \begin{tabularx}{0.7\textwidth}{X} 他們是你的僕人和你的百姓，是你用大力和大能的手所救贖的。 \end{tabularx} \\ \\ \relax
1:11 & \begin{tabularx}{0.7\textwidth}{X} 唉，主啊，求你側耳聽你僕人的祈禱，聽喜愛敬畏你名眾僕人的祈禱，使你僕人今日亨通，在這人面前蒙恩。」我是王的酒政。 \end{tabularx} \\ \\
[1ex]
\hline
\hline
\end{longtable}
$^{1}$今天是經訓的時間.
我會選讀一段經文.
是舊約聖經的《尼希米記》.
在大家手上的聖經第591頁.
我們今天會讀出.
《利未記》第一章一至十一節.
哈加利亞的兒子.
尼希米的言語如下.
阿達是西王.
二十年基斯留約.
我在書山城的宮中.
那時有我一個弟兄哈拿尼.
同著幾個人從猶大來.
我問他們那些被擄歸回.
盛夏逃脫的猶大人.
和耶路撒冷的光景.
他們對我說.
那些被擄歸回盛夏的人.
在猶大省遭大能受凌辱.
並且耶路撒冷的城牆拆毀.
城門被火焚燒.
我聽見這話.
就坐下哭泣悲哀幾日.
在天上的神面前禁食,祈禱,說.
耶路說天上的神.
大而可畏的神啊.
你向愛你守你誡命的人.
守約,施慈愛.
願你睜眼看,側耳聽.
你僕人晝夜在你面前.
為你眾僕人以色列民的祈禱.
承認我們以色列人.
向你所犯的罪.
我與我父家都有罪了.
我們向你所行的甚是邪惡.
沒有遵守你藉著僕人摩西.
所吩咐的誡命,律例,典章.
求你紀念所吩咐.
你僕人摩西的話,說.
你們若犯罪.

$^{41}$我就把你們分散在萬民中.
但你們若歸向我.
謹守遵行我的誡命.
你們被趕散的人.
須在天涯.
我也必從那裡將他們招聚回來.
帶到我所選擇立為我命的居所.
這都是你的僕人,你的百姓.
就是你用大力和大能的手所救贖的.
主啊.
求你質以聽你僕人的祈禱.
和喜愛敬畏你名.
眾僕人的祈禱.
使你僕人現今坎通.
在王面前蒙恩.
我是作王酒政的.
有一段時間.
我們都沒有在講壇裡宣講尼希米基的信息.
如果你有留意.
尼希米基和他在聖經裡的上一卷書.
以斯拉記.
可以說是一冊.
一卷書分兩冊.
上下冊.
都是記載著當上帝的子民.
從老遠的波斯帝國回歸耶路撒冷之後.
到底發生了什麼事.
上一集的以斯拉記.
主要的信息就是.
怎樣幫助這群百姓重建聖殿.
更重要的就是.
重新定位律法書在百姓心目中的地位.
他們之所以亡國被擄.
正正是因為他們不遵守上帝的律法.
而下集的以斯拉記.
其實就是強調在重建聖.
剛才也說了.
城牆毀壞.
在重建城牆的過程裡.
怎樣帶領這群百姓.

$^{81}$向上帝重新立約.
強調約的關係.
透過這一切的過程來恢復.
不只是破壞了城牆.
是破壞了人和上帝的關係.
正如我們剛才所聽到的.
尼希米基的故事一開始就說到.
尼希米的弟兄哈拿尼和幾個同胞.
從猶太來到波斯的書山城的皇宮.
尼希米從他們的口中知道.
那些被迫歸回的以色列人.
在當時的猶太省.
其實當時的耶路撒冷或者猶太.
已經是在波斯帝國的統治下.
他知道猶太省正是遭患難受凌辱.
並且耶路撒冷的城牆遭到拆毀.
城門也被火焚燒.
聽到簡單這幾個字的描述.
你不難理解生活在耶路撒冷城的人.
當時的情況是怎樣.
沒有什麼好日子過.
因為失去了城門.
失去了城牆的保護.
意味著耶路撒冷就變成了一個無掩雞籠.
你要留意當時的猶太.
耶路撒冷已經不只是以色列人.
自從旺角被擄之後.
那裡充斥著所有的外邦人.
他們四處都是敵人.
失去了城牆的保護.
意味著他們四處的敵人.
可以自由進出耶路撒冷.
可以肆意騷擾他們.
可以搶他們的東西.
破壞他們所做的.
並且甚至殺了他們.
使得這群百姓.
即使回歸了應許之地.
但是生活仍然極其困苦.
每天都好像打仗戰亂一樣.

$^{121}$而另一方面從信仰的角度.
我們說耶路撒冷的城牆.
其實就像一條屬靈的界線.
一個屬靈的保護網.
將上帝的子民從世俗裡.
我們說分別出來.
保守他們不輕易被外邊的文化.
來去摻雜,污染.
也幫助他們不容易.
跟世俗的價值觀念來妥協.
可以專一敬拜上帝.
而當李希米聽到.
耶路撒冷和同胞的慘況的時候.
經文很清楚.
他記載了他的情緒反應.
坐下哭泣.
悲哀幾日.
你看到李希米他雖然.
人在波斯的書山城皇宮.
並且是當時的波斯王來工作.
但是他的心呢.
他的心掛在哪裡.
很明顯.
你會看到他仍然心是上帝的國.
神的家.
他知道聖城耶路撒冷.
其實就是代表著.
上帝國度的彰顯.
他即時這樣掛勾.
理論上.
上帝的子民現在已經回歸了.
應許之地.
應該安居樂業.
上帝應許幫助他們回歸.
但是為什麼仍然.
過著朝不保夕的生活.
城牆被毀壞.
並且百姓仍然受外敵欺壓.
其實就反映著.
上帝的國度仍然是荒涼.

$^{161}$上帝的榮耀竟然不能夠.
在屬於他名的耶路撒冷來彰顯.
李希米是從這個角度.
來理解整件事的本質.
再加上.
遠水不能夠近火.
我們說廣東話.
即是說李希米當時.
他雖然是波斯王的酒精.
擁有一定的身份地位甚至權力.
但是他可不可以隨便請假.
不可以的.
我們不可以用現代的角度來看.
鄉下有事.
老闆我遞張架子.
回鄉下幾天幫幫他們.
不可以的.
他不能夠馬上抽身.
回到耶路撒冷.
為同胞提供緊急的救援.
因為酒精的工作.
其實是一個很重要的事情.
他需要每天預先去嘗試.
王所飲用的酒.
國王飲之前他要先飲.
確保沒有人在裡面下毒.
一個這麼重要的職位.
其實不可能說走開就走開.
波斯王一定不會輕易的批准.
所以當我們了解了這些背景之後.
你會發覺當李希米聽到耶路撒冷的荒涼.
但是自己又無能為力的時候.
對一個很愛上帝.
很關心上帝國度的事情的人來說.
其實那份無力感.
怎會令到他不傷心.
不過李希米有趣的地方就是.
他不會因為自己這份無力感.
這份悲傷.
他病了自己,怨天尤人.

$^{201}$在他最難過,最無助的時候.
他懂得做一件事是什麼?.
他會在乎上帝.
他懂得去抓緊那位天上的神.
向祂禁食禱告.
很希望迫切尋求到上帝在這件事的介入.
事實上我們說.
李希米本身是一個很強調很重視禱告的人.
如果你有留意整本書只有13章.
但卻有9次收錄了他的祈禱.
今天我們看的第一章.
也可以說是在聖經裡很多的祈禱例子裡的一個典範.
裡面有很多的屬靈原則.
其實很值得我們借鏡.
應用在你和我的禱告生活裡.
首先我們會看到.
不知道大家有沒有留意地域的差距.
很多的差距.
李希米和耶路撒冷的百姓.
李希米身處在書山的王宮.
另一邊是遠方的耶路撒冷的百姓.
其實地域上已經距離很遠.
身份呢.
一個是當時政府的高官.
另一些只是普通平民的百姓.
甚至乎是一些難民.
他們的生活環境呢.
在摩斯城的書山王宮的酒鎮.
一定是生活富裕安穩.
耶路撒冷的百姓呢.
極之貧窮.
朝不保夕.
這麼大的差距是什麼.
將他們拉在一起.
到底是什麼.
禱告.
祈禱將他們拉在一起.
李希米的祈禱.
突破了雙方這些的界限.
讓他們現在都一起在上帝的面前.

$^{241}$他們準備經歷上帝的帶領.
成就上帝的旨意.
所謂經歷靈明的成長.
李希米這個情況讓我想起.
早兩個星期五.
不知道大家有沒有參與.
我們在這裡舉行了一個.
越南宣教士的分享.
我知道待續都有弟兄姊妹參加.
問問大家.
香港和越南為什麼會連成一線.
因為什麼緣故.
記得越南難民嗎.
越南難民記得我.
當時有很多難民因為.
當地的一些問題.
來了香港.
當時也有很多的教會.
傳道人 牧者為他們傳福音.
當他們回去之後.
他們就在當地興起教會.
隨之而後也慢慢有很多的宣教工作.
可以說其實.
福音的需要.
將這兩個遙遠的地方連在一起.
不過當時宣教士的分享.
你猜猜.
在香港有多少位越南宣教士.
有多少個.
只有一個.
就是分享那個.
並且他也快要退休了.
他很希望有接棒人.
當日我們為了越南的需要祈禱.
並且為了黃浸祈禱.
本來是毫不相關的兩個群體.
但你有沒有發覺.
其實這幾年我們走在一起.
疫情之前.
疫情 剛好去年.

$^{281}$疫情後.
黃浸也有支持越南的宣教工作.
神學發展的工作.
並且我們有兩次的探訪.
不知道大家有沒有留意.
是什麼東西令我們走在一起.
是福音的緣故.
是你和我的禱告.
禱告將兩群.
其實沒可能接觸到的人.
慢慢拉近.
事實上 弟兄姊妹.
當日我們禱告.
我們也承諾.
希望更加為這個地方.
來行上禱告.
並且懇求上帝.
會不會在我們當中.
興起更多的人.
來關心越南的科學工作.
我不知道將來我們會不會出產一個.
去越南的宣教士.
但我確信在未來的日子.
教會有更多的參與.
有更多在越南的服侍和短宣.
很希望大家來用禱告守望.
所以將來當有.
這一類活動的時候.
我盼望大家去留意.
用你的禱告.
將你和他們拉近.
一起去回應上帝.
在他的福音裡的呼召.
事實上 弟兄姊妹.
我想我們和尼喜米一樣.
某程度上.
我們都是人在江湖.
即使怎樣也好.
你一定有其他的工作身份.
就算你是家庭主婦也好.

$^{321}$你也有一些工作身份.
但我們真正的身份是什麼呢.
其實都是天國的子民.
這才是我們最真實的身份.
我們是耶穌基督的門徒.
到底上帝的事情.
我們能不能以此為優先.
我們的禱告生活.
能不能像尼喜米那樣.
將遠方或近處的一些需要.
你攬上身.
不要太過抽離.
你會不會將它都看為.
其實是屬於你的事情.
雖然地域很遠.
雖然雙方好像很有距離.
但如果你從上帝的角度來看.
其實都是你和我要關心的東西.
我們會不會甚至思考.
到底上帝讓我遇見這些事情的時候.
我可以怎樣進一步.
走近上帝這些被我們見到的異象.
去回應上帝給我們的呼召呢.
特別我知道每個禮拜.
我們的程序表.
如果你有留意.
最後都有一頁是代禱的.
有些是為遠方的地方代禱.
確實很有距離感.
你又不認識.
你又沒去過.
這期是印度.
你一想起你第一個反應是什麼.
都不太關我事.
因為我又不會去.
我又不懂.
我又不懂印度話.
印度餐我都不會吃.
所以我都放下了.
不會為它禱告.

$^{361}$但其實這個是雞和雞蛋的問題.
我們陌生所以不禱告.
但反過來其實.
我們不禱告就不是陌生.
雞和雞蛋的問題.
正因為這樣.
正因為我們不熟.
可能更加我們要留意.
我們更加需要去禱告.
如果我們當為這個是神家的事情.
你和我都需要參與.
無論多多少少遠遠近近.
我們都需要用信心的禱告.
來擺上的時候.
靖姐妹你信不信.
你的禱告能夠打破了.
這些你心裡面的界限.
你信不信上帝會透過你的禱告.
拉近這些和遠方的.
方便需要的距離.
讓我們明白.
他們的事是神家的事.
亦都和我們有關.
上帝期望我們.
有這一份代禱的承擔感.
特別在這個世代.
我們很需要明白.
就是他們的事情.
其實都反映著上帝的國度.
能不能夠彰顯.
如果我們從這樣的角度來看的時候.
其實這個都是我們的責任.
我們需要定期為他們守望.
事實上你會留意到.
在教會歷史你會看到.
很多上帝國度復興的工作.
是由什麼開始.
不是人的策劃.
是禱告.
如果你有留意.

$^{401}$或者你有看我們誠心所願.
上年我們做過一集的節目.
是說去年二月.
美國有一個大學叫做.
Asbury University.
亞斯不利大學.
出現了一場禱告復興的運動.
你想留意的話可以買廣告.
上年三月二十二.
誠心所願.
就是在說這個事情.
起初只是一群.
經常在教會聚會的弟兄姊妹.
來去敬拜.
但是結束的時候.
他們繼續敬拜讚美.
並且禱告認罪悔改.
分享見證.
甚至有些學生被聖靈的感動.
離開了那間教會.
完了聚會.
他再回去.
就是這樣形成了一種.
二十四小時不停的禱告敬拜.
甚至燃點了很多青少年.
他們願意來到上帝的面前.
去認罪悔改.
從CBN的新聞報導看到.
其實有成千上萬的人.
從各地湧去那間教會.
令到這個禱告復興運動.
不單止二十四小時.
持續了十七天.
十七天.
你在香港應該大概不會見到.
一個這樣拉到這麼長的禱告復興會.
所以.
阿斯匹利的學生他們形容.
其實這個一定是上帝的工作.
不是人為能夠策略出來的事情.

$^{441}$弟兄姊妹其實你看回教會的歷史也好.
聖經也好.
信徒的復興.
往往有兩個很重要的因素.
一個就是環境的動盪.
你看回《使徒行傳》也是這樣.
最初由一小撮.
看慕上帝的人開始.
從他們個人靈命的復興.
慢慢帶動整間教會群體.
到最後影響當地的社群.
出現一些更新的變化.
過程裡面有另一個很重要的元素.
不可以缺少的就是.
人向上帝一個認罪悔改的祈禱.
就好像我們今天這段經文這樣繼續看.
我想請大家去讀五至七節好不好.
看螢光幕或者看你的程序表.
第五至七節第二段經文.
預備起.
願你睜眼看則以聽.
你僕人就夜在你面前.
為你眾僕人以色列民的祈禱.
承認我們以色列人向你所犯的罪.
我和我家都有罪了.
我們向你所行的甚是邪惡.
沒有遵守你直人僕人摩西所吩咐的誡命律例典章.
好,謝謝.
李希米他開始他的祈禱.
不過你留意的內容是很特別的.
首先他不是求上帝.
上帝你要保護耶路撒冷的百姓.
你要幫助他們修復城牆.
讓他們不再受到敵人的攪擾.
他不是求上帝保護.
他也不是求上帝對付這些外敵.
他沒有譴責對手.
最壞是他們來搞擾我們的同胞.
求上帝教訓他們.
兩樣都不是.

$^{481}$而是什麼.
祈禱認罪.
為什麼李希米的祈禱.
首先選擇的是由認罪開始.
這個很值得我們去想.
或者很值得放在我們的生活裡.
用今天的術語就是.
他懂得從現象看到問題的本質是什麼.
他很明白耶路撒冷的光景.
為什麼會這麼荒涼.
真正的原因.
不是因為城牆不夠堅固.
失去了保護.
不是因為敵人麻煩.
勢利經常搞擾我們.
耶路撒冷的荒涼只有一個原因.
是百姓和上帝的關係出了問題.
他看到的本質是這樣.
因為明明已經按照上帝的應許.
回到耶路撒冷這個應許之地.
為什麼上帝的福氣沒有在耶路撒冷出現.
李希米他自己都一語道破.
在他的認罪祈禱裡.
是因為這班百姓回歸之後.
心有沒有改變.
沒有.
他仍舊好像他們的祖先那樣.
仿效他們的惡行.
沒有遵守上帝藉著摩西所吩咐的誡命律例典章.
其實這個是很有趣的.
他說的從摩西開始已經是這樣.
由他們出埃及到西納山立約.
這個都是上帝指紋的一直以來的問題.
根本問題.
去到哪裡都要處理這個問題.
到今天都要處理這個問題.
我們和上帝的關係如果出了問題的時候.
很多時候會反映在我們的現況裡.
要根治的不是那個事情.
而是那個關係.

$^{521}$是我們和上帝的關係.
李希米他看到這一點.
他明白因為人遠離上帝.
所以上帝按照他在五經裡面.
之前已經所講過的情況.
他是一定會管教的.
怎樣管教.
就是透過外敵來管教猶太.
很希望他們會回轉.
所以李希米他判斷出耶路撒冷的荒涼是一個屬靈的問題.
屬靈的問題自然要用屬靈的方法來解決.
所以他知道要重建的不是城牆.
也不是對付外敵.
他要重建的是首先破壞了的關係.
那個破壞了的人神關係.
這個才是首先要去修復.
所以我們看到李希米他很謙卑來到上帝的面前.
他承認自己和他的家族以及所有同胞.
以往一直所犯的事情.
不過值得注意的是.
李希米的禱告悔改.
他不是說上帝他們得罪了你.
他們很有問題.
他們甚是邪惡.
李希米不是這樣.
有沒有留意李希米用我們.
李希米是將他自己的身段放到最低.
願意和那些得罪上帝的人來看齊.
他一再強調是我們.
我們的父家.
我和我的父家都得罪了上帝.
事實上你會看到李希米將他自己的名字放在認罪名單裡.
他沒有和那些遠方耶路撒冷的百姓分割出來.
他很願意為他代討的對象來認同來看齊.
其實這件事在新約裡耶穌基督都完全反映出來.
大家有沒有留意.
耶穌基督來找施浸約翰的時候.
他要求接受施浸約翰的悔罪洗.
他有沒有罪?沒有.
但他卻願意盡諸般的義.

$^{561}$與罪人,與我們認同看齊.
耶穌被人覺得他是高高在上.
遙不可及.
他願意與人看齊這種禱告代討的姿態.
他願意與人禱告,與人看齊這種姿態.
其實很需要我們學習.
很需要我們放在我們的禱告生活裡.
唯有這樣才可以拉近人與人之間的聯繫.
我還記得早前同一位弟兄姊妹來祈禱.
她的生活很艱難.
她有很多疑惑.
甚至埋怨上帝.
她覺得她已經做得很好.
但上帝都沒有給她出路.
甚至更加新的艱難出現.
我都沒有答案.
我唯一能做的就是與她一起祈禱.
我與她一起說上帝我們真的受不了.
我們不知為何艱難一浪一浪地出現.
我們都是人.
我們不明白上帝的高深旨意.
但我們又知道上帝一定有原因.
求上帝就算你不解開那些難處.
你都給我們有力量去承載.
讓我們仍然維持住你.
讓我們在疑惑裡.
上帝維持住我們.
讓我們不要走開.
最終這個禱告完成了.
這位弟兄覺得舒服.
這位弟兄姊妹覺得舒服.
是因為不單止上帝會聽禱告.
是因為他們都明白.
有人沒有走開.
身邊有一個人和他一樣軟弱.
一樣有限制.
都是罪人.
都需要恩典憐憫.
所以其實有時代禱.
當我們用我們的時候.

$^{601}$弟兄姊妹就不知不覺拉近我們的關係.
我們不是抽離.
我們以一個高姿態的身份.
來幫對方.
而反過來在禱告裡.
我們都是一模一樣.
我們都需要上帝的拯救.
我們願意與對方一起去面對.
願意與對方一起去尋求上帝的旨意.
讓他知道.
讓我們彼此知道.
我們是同路人.
我們是不孤單.
我們一樣會經歷上帝的恩典.
確實.
在今日高舉個人主義的情況下.
其實基督徒都受影響.
我們有時都會將自己看成.
教會裡的一個獨立份子.
不知你覺不覺.
我們明明是一個身體.
但我經常覺得我們會看自己.
是一個獨立的個體.
那些事都不關我的事.
那個人我又不認識.
有什麼關係.
這種情況很吊詭.
所以我們不習慣將弟兄姊妹的一些甩漏過失.
我們會摟上身.
我們不習慣看為.
其實這些都是我們的事情.
但今天黎希敏的禱告提醒我們.
教會群體的軟弱.
其實是要彼此擔當.
我們要一起學習面對.
如何同心尋求上帝的憐憫和幫助.
正如新約所說我們是一個身體.
不同的功效.
但是一個身體.
你旁邊那個有事.

$^{641}$其實都影響到你.
影響到整個群體.
需要一起去守望.
而不是獨善其身.
事實上你會看到.
黎希敏的禱告正正就是反映這一點.
所以弟兄姊妹我們願不願意.
為身邊的人禱告.
哪怕那些事情.
你是不認識的.
你嘗試了解.
你透過禱告求上帝幫你.
和你待討的對象拉近距離.
以至上帝有進一步的帶領給你.
讓你可以多走一步.
我們再看下去.
第八至十節.
你看到黎希敏這裡的禱告.
他反映了一件很重要的事.
他很熟悉上帝的話語.
反映在他的禱告裡面.
他的禱告不是不著邊際.
也不是自由創作.
是很明顯他知道.
上帝對人的要求是什麼.
很明顯他清楚知道上帝的行事.
法則是怎樣.
他知道百姓只有一個出路.
就是回轉.
就是真誠向上帝悔改.
謹守進行上帝的誡命.
重新履行和上帝所納的約.
這是唯一的轉機.
唯有這樣才能再次經歷上帝的福氣.
所以其實你要看到.
黎希敏這一段的禱文.
雖然是很簡單.
不是很長.
如果我們讀出來.
一分鐘內就讀完.

$^{681}$但是卻是很清楚.
很有力.
是因為他見機在生命的應許.
所以他知道.
上帝一定會按照他的應許.
在上帝的時機裡.
兌現上帝自己的承諾.
所以你看到.
黎希敏最後那句說什麼.
第十節.
這些都是上帝你的僕人.
是上帝你的百姓.
就是你用大能大力.
的手所拯救出來.
他知道其實最關注的是上帝.
最肉緊的是上帝.
上帝做這麼多事.
都是希望管教他們.
讓他們回到上帝的身邊.
上帝有能力拯救他們.
並且上帝也願意拯救他們.
黎希敏確信這一點.
這也成為了他禱告的信心.
說了這麼久.
和大家做一點測試好不好.
你看看你旁邊那個.
你覺得他有沒有能力請你吃飯.
等下完崇拜.
有了.
看到大家很餓了.
第二個問題.
你覺得他.
等下完崇拜會不會請你吃飯.
願不願意.
為什麼你覺得不願意.
他不願意請你吃飯.
有人搖頭.
第一個問題大家都點頭.
第二個問題.
大家有人搖頭.

$^{721}$其實放在你的信仰裡很值得參考.
你覺得上帝是不是有能力拯救你.
是,但是上帝願不願意拯救你.
你會不會搖頭.
其實有人會搖頭.
有人會覺得我知道上帝可以.
不過他不會救我.
你只要少了一件事.
你的信心就會動搖.
接著你就會開始靠自己.
你一定是這樣.
但關鍵是你的禱告提醒我們.
他確信上帝既有能力拯救百姓.
也確信上帝願意拯救百姓.
有了這兩樣東西我們就叫信心.
就是信心.
雖然這件事還沒有實現.
但你正在捉緊上帝這兩個應許的時候.
你的禱告就是有力有信心.
不會不著邊際.
不會覺得是踩著空氣.
因為是建基於上帝說了的話.
上帝一定會承諾的.
只要你願意按照上帝的法則.
回到他的拯救道路.
就必定成就.
事實上這兩樣東西.
也成為了尼希米.
他代禱的信心.
我相信大家都認識.
早期有位教父叫奧古斯丁.
大家都認識.
但大家知不知道他年輕的時候.
其實生活很放蕩.
他甚至信奉異教.
他十九歲就有自己的私生子.
但這個人超級聰明.
但也相當反叛.
他怎樣回轉呢.
就是靠他媽媽的代禱.

$^{761}$他有一個很敬虔的母親叫Monica.
是一個虔誠的基督徒.
她看著兒子這樣但又沒有扶.
那怎麼辦呢.
禱告.
雖然一直都不見效.
但她深信.
她用膝蓋來托著兒子的禱告.
是可以的.
她知道上帝一定有能力幫助她的兒子.
上帝也願意.
於是她不斷堅持.
為自己的兒子守望.
終於有一天.
上帝在奧古斯丁三十一歲的時候.
引導了她悔改歸主.
並且我們知道.
她在教會歷史裡.
在早期教會成為了一位很偉大的神學家.
事實上.
頂智妹很多人能夠回轉歸向上帝.
是因為背後有人.
是有一個很偉大的神學家.
是因為背後有人.
持續為她祈禱.
很希望她能夠回到上帝那裡.
坦白說.
每一個人面對生命的軟弱.
甩漏甚至每個人墮落.
罪惡的裡面.
其實他們都想改變.
他們都想走出那個困局.
沒有人想過一個這樣的生活.
沒有人覺得自己天生.
喜歡為難.
喜歡污穢.
只不過有些人.
不知道怎樣改.
又或者有些人知道.
但他沒有這個能力自己走出困局.

$^{801}$當然也有些人.
很抗拒改變.
在這個時候.
他很需要身邊有人.
這個人.
既知道上帝是有能力幫他.
這個人也知道.
上帝是願意幫他.
願意向他守施慈愛.
並且這個人.
願意按照上帝的屬性和.
向他恆常.
為那些軟弱的人祈禱.
很希望上帝有一天.
讓他們可以回轉.
向他們守施慈愛.
弟兄姊妹.
到底為著上帝的緣故.
你願不願意成為.
一個這樣的代禱勇士.
一定有些人在你身邊.
這樣出現.
上帝讓你看到不是偶然.
你會不會動了慈心.
都將他擺在你的禱告生活裡面.
讓上帝.
都可以使用你.
成為他的守望.
好,看完最後一段經文.
最後這段經文.
11節這裡說.
主啊,求你則以聽你服人的祈禱.
和喜愛敬畏你命.
頌服人的祈禱.
使你的服人現今亨通在王面前蒙恩.
最後尼熙米.
他自己介紹自己的工作.
我是作王酒政.
之前也說過酒政這個工作.
其實很重要.

$^{841}$因為他要預先品嘗.
王每一餐所飲用的酒.
確保沒有人在裡面出古或下毒.
能夠當上這個重要職位.
你看到.
尼熙米有很多.
重要的條件.
首先他一定要深得王的信任.
還要留意他是一個外邦人.
對波斯人來說.
要取得王的信任.
一定有他個人之處.
他也必定是一個很睿智.
不會忽略細節的人.
再加上他一定會很熟悉.
官場的遊戲規則.
他知道宮廷鬥爭的險惡.
所以他也會懂得.
防備他的對手.
尼熙米他察覺到.
這一切職場裡面的磨練.
這一切累積的人脈.
其實也是上帝給他的裝備.
可以讓他在現在.
有機會回到耶路撒冷.
協助這班同胞.
重建城牆.
重建對上帝的關係.
特別身為酒精的他.
其實他每天都會見到王.
他最有機會接觸到波斯王.
向他提出這個請求.
尼熙米知道一定是他.
上帝選了他.
問題來了.
怎樣向王開口.
怎樣開口.
因為在.
以斯拉記.
尼熙米記的上一本書.

$^{881}$其實記載了亞達斯西王.
即是現在的王.
是因為誤信殘言.
以為耶路撒冷的百姓.
想造反想重建城牆.
所以下令.
不准興建城牆.
要停工.
所以如果.
尼熙米現在有了全盤計劃.
在他腦海裡.
怎樣幫這班百姓.
他也不敢貿然向王提出重建城牆.
我們說這麼敏感的事情.
為什麼.
用今天的術語.
提出這麼敏感的事情風險很大.
隨時會.
觸犯國安.
是會的.
所以尼熙米.
他沒有辦法.
他一定要上帝介入.
沒有上帝介入.
沒有天上的王搞定地上的王.
他就算多聰明想好了也沒有用.
他要求.
天上的王阻止.
任何地上的權勢.
阻礙上帝的國度彰顯.
這個也是.
尼熙米背後的意思.
他很確信.
如果上帝真的要差派我回到耶路撒冷的時候.
去事奉他.
這樣上帝一定會開路.
上帝一定有方法.
令到波斯王放我回去.
他知道時機掌握在上帝的手裡.
所以他行刺禱告.

$^{921}$等候這個時機出現.
當然這個時機.
最後在下一章出現.
今天不會在這裡交代.
有機會你可以自己看看這個時機是怎樣出現.
但讓姐妹們看到.
尼熙米他.
貴為波斯王的酒精.
其實他已經本身.
擁有一個很優越的身份.
富裕的生活.
用今天的術語他已經是怎樣.
上了岸.
等退休.
吃長糧.
為什麼他要冒這個風險.
為什麼他要冒這個隨時失去一切.
安逸的風險.
是因為他始終沒有忘記自己是誰.
他沒有忘記自己是上帝的子民.
他這個身份也驅使他願意放下他自己.
放下他擁有一些的安逸.
甚至榮華.
來回應上帝給他的託付.
結果他的人生下半場改寫了.
他之後也成為了猶大省的省長.
來成就上帝更大的事情.
帶領百姓進行一個靈明的復興.
一說起復興.
我想大家都可能有留意.
近來香港的市況都復興不了.
很低迷.
對很多人來說.
未來投資其實都很審慎.
持抱抱撲審慎的態度.
不會貿然出手.
可能你問可能你認識身邊有些人都會覺得.
市場都不知道什麼時候好回.
等好回的時候再想.
現在先不要動.

$^{961}$當然我不是和大家分析後市.
不是叫大家怎樣投資.
我是很想藉著這段經文.
其實我們一起反省.
我們屬靈上的投資.
我們會不會趁這個時候出手.
為上帝冒險.
放膽將你手上的恩賜.
你的才幹.
你的資源.
投資在上帝永恆的角度上.
這個一定賺.
不會輸的投資上.
事實上聖經告訴我們.
教會的歷史告訴我們.
每當人心動盪的時候.
福音的指數一定是怎樣.
不會跌的.
一定會升的.
人心動盪的時候福音指數一定會升.
因為人需要上帝.
人沒有了平安.
所以可能很多人.
都沒有想過.
都可能活了一段時間一把年紀.
都經歷過的時候.
都可以安逸了.
到現在竟然要經歷一些艱難的日子.
但是從屬靈的角度看.
這個反而是有利福音的擴展.
我們能不能設法.
讓身邊的人認識到福音.
知道耶穌.
才是他們唯一的拯救.
就是這段經文最後給我們.
一些提醒和祝福.
所以弟兄姊妹尼希米的禱告.
再一次喚醒我們.
我們身份是天國的子民.
我們需要像耶穌說的.

$^{1001}$先求他的國 先求他的義.
將上帝的家.
上帝的國度擺在優先的位置.
藉著你的禱告.
你真的可以修補.
別人一些破損的生命.
藉著你的禱告.
真的可以修補教會出現一些漏洞.
藉著你的禱告.
可以讓上帝的國度彰顯.
在我們的生活當中.
我們一起祈禱.
祈求上帝我們願意.
透過我們的禱告.
首先將我們擺到你的面前.
再一次提醒我們.
真正的身份是屬於你的子民.
是君尊的祭司.
是耶穌基督的門徒.
這個身份也定性了.
我們生活的焦點.
要以上帝你為首位.
先求你的國和你的義.
上帝盼望藉著今天的經文.
提醒我們.
當我們在生活中遇見.
別人給予我們的異象.
給予我們一些代禱的需要的時候.
讓我們都可以停下來.
讓我們都不要太快.
覺得很遙遠.
覺得好像不太關我自己的事.
讓我們確信.
這一切其實都是上帝你的美意.
讓我們願意去認識.
讓我們願意去開始.
點點滴滴地.
圍到這些需要的禱告.
當我們願意真誠.
來到上帝你的面前的時候.

$^{1041}$上帝眼睛幫助我們.
讓我們不單單走近你.
讓我們對這些需要.
更加去求問上帝.
怎樣去使用我們.
讓我們可以為上帝.
你做更多更大的事.
讓你的國度.
能夠更有效地.
在我們的生活裡面.
在地上彰顯.
上帝我們知道這是你的心意.
一直都沒有變.
只不過我們可能在生活裡面.
我們遺忘了.
求上帝你憐憫.
求上帝你寬恕.
藉著今日經文的提醒.
讓我們首先回到你裡面.
重新定位我們的身份.
以致我們生活的優先次序.
亦都願意對準你.
無論是近處,遠方的需要.
我們都用心去擺上.
求上帝你在我們這個群體裡面.
復興我們.
讓我們真的可以土地的喜悅.
祝福身邊更多的人.
我們這樣的禱告交託.
奉耶穌基督你的名字來祈求.
阿們.
\newpage



\section{歷代志上 4:9-10}
\label{sec:_RD54XIpk84}
\textbf{2024年2月4日崇拜直播｜陳明泉牧師｜十個深化關係的禱告-逆轉痛苦的禱告｜歷代志上4\_9-10}
\newline
\newline
連結: \href{https://youtube.com/watch?v=-RD54XIpk84}{\texttt{https://youtube.com/watch?v=-RD54XIpk84}} ~~~~ 語音日期: 2024-02-06
\newline
\newline
\hyperref[sec:H22XGgjIH6M]{\small{< < < PREV SERMON < < <}}
~
\hyperref[sec:index_chronic]{\small{[返順時目]}}
~
\hyperref[sec:xitIkhwfBC8]{\small{> > > NEXT SERMON > > >}}
\newline
\newline
歷代志上 4:9-10
\newline
\begin{longtable}{cl}
\hline
\hline
章節 & 經文 (和合本修訂版)\\
\hline
4:9 & \begin{tabularx}{0.7\textwidth}{X} 雅比斯比他眾兄弟更尊貴，他母親給他起名叫雅比斯，意思說：「我生他甚是痛苦。」 \end{tabularx} \\ \\ \relax
4:10 & \begin{tabularx}{0.7\textwidth}{X} 雅比斯求告以色列的神說：「甚願你賜福與我，擴張我的疆界，你的手常與我同在，保佑我不遭患難，不受艱苦。」神就應允他所求的。 \end{tabularx} \\ \\
[1ex]
\hline
\hline
\end{longtable}
$^{1}$今次讀的經文是記載於歷代至上第四章九至十節.
「亞比斯比他眾弟兄更尊貴.
他母親給他起名叫亞比斯.
意思說我生他甚是痛苦.
亞比斯求告以色列的神說.
甚願你賜福於我 擴張我的境界.
尚與我同在 保佑我不遭患難 不受艱苦.
神就應允他所求的」.
聖經讀筆.
弟子們平安.
我們今天繼續以十個禱告.
一路的這個系列來思考.
究竟我們今天在這個世代.
在我們的人生裡面.
如何實踐禱告的生活.
在大概公元二千年.
即是二十多年前.
出過一本書叫做《亞比斯的禱告》.
不知道大家有沒有看過.
很多弟兄姊妹都看過.
這本書都頗熱賣的.
一年之間就成為了.
在當時美國.
其實香港已經是銷售榜的第一.
排第一的了.
每個人都藉著這本書來學禱告.
這本書是有價值的.
因為它除了說.
用亞比斯的祈禱.
教我們如何去祈禱之外.
它也針對當時很多人.
已經禱告生活很不濟.
又或者跟神的關係很遙遠.
這本書某程度上也喚醒了很多人.
一起來成為禱告的佳母.
甚至很多人就組團.
成為了類似亞比斯的禱告團.
一起來跟隨亞比斯.
跟著他的格式來祈禱.
不過也有一些說法就是.

$^{41}$其實亞比斯的禱告.
如果你說說那個份量.
或者那個屬靈的深度.
其實沒有什麼內容.
兩句而已.
剛才我們已經讀完了.
所以如果你說將它成為了一個靈修.
或者禱告的佳母.
似乎在這裡的安排.
或者這種方式似乎有些缺陷.
所以有另外一批人就說.
其實你學禱告不是這樣學的.
甚至有時候有一種說法.
只是一兩句說話在聖經裡面.
應該要比較淡然一點去面對.
甚至有一種說法就是.
其實小人物的祈禱.
我們不需要特別的.
去將它放到這麼高的位置.
兩個說法都有其道理的.
不過今天我們讀這段經文.
我們就放在中間.
我們既不會過份覺得.
這個就是我們禱告的典範.
我們照樣跟著禱告就沒有死.
但同時間我們也不會完全輕看它.
因為當聖經記載著.
特別是在歷代志的記載當中.
有一個很特別的.
應該是一個家譜記載.
寧寧舍舍記載著.
這一兩節的經文.
對整個歷代志的理解.
對禱告的生活.
有一個很重要的訊息.
當我們今天去解清楚這段經文的時候.
我們就能夠將這段經文.
放在正確的位置當中.
希望我們今天一起去看的時候.
我們看到亮光.

$^{81}$今天要涉及的經文比較多.
希望大家我們用心去研讀.
也希望我們能夠在過程當中.
幫助我們.
讓我們的禱告生活更加有力量.
我們先祈禱求上帝幫助我們.
同心祈禱.
親愛的天父願你在我們當中.
帶領我們讓我們能夠明白你的話語.
求你親自臨隔在眾弟兄姊妹當中.
讓每一個聆聽的時候.
他們不單聽得明白.
更加是讓心靈可以敞開.
又讓宣講的講得清楚.
忠心的將你的話語.
向弟兄姊妹講明.
讓我們能夠明白你的時候.
我們更加明白.
怎樣用禱告來親近你.
更加明白你的心意.
恭敬將來的時間交托.
願你賜福.
天上禱告奉靠基督耶穌得勝的聖名.
Amen.
在我們看剛才那個歷代志裡面.
所記載的亞比斯禱告的時候.
有一些背景資料.
我要給大家知道的.
第一就是.
這兩句的經文.
其實竟然在歷代志的家譜裡面.
那個家譜其實很複雜.
甚至乎你這樣看.
會有些乏味.
我們看一看大綱裡面.
你會看到.
其實歷代志用九章的篇幅.
來寫家譜.
你一讀上去你自己覺得.
頭昏腦脹.

$^{121}$和合本經比較容易看.
為什麼呢?因為它會加多一點點的注釋.
加多一點點的意義.
誰生誰.
原文連那個生字都沒有.
只是人名.
但是當你看的時候.
你會發覺有些地方.
有一個.
pattern在裡面.
第一章的家譜.
由一章一節到第五十四節.
那個是.
創世記的家譜.
原則上它是記載著.
創世記裡面.
出現過的.
那些人的名字.
被記錄在第一章一節到第五十四節.
第二章一節到第五十四節.
第二個.
的家譜就是由第二章開始.
由第二章.
去到第.
四章的.
第二十三節.
這個第二個.
大段落其實是.
記載著猶大.
家譜.
亞國的.
兒子猶大的家譜.
中間夾雜著的.
就是大衛王朝.
大衛家的家譜.
三章一節到二十四節.
你看到.
亞比斯的祈禱.
夾雜著在猶大家譜裡面.
第四章九節到十節.

$^{161}$這個是.
亞比斯是.
猶大家裡面的.
這個家裡面的一份子.
是猶大支派的.
一個人來的.
好了,接著.
講完猶大.
接著到第四章.
二十四節到第二十六節.
記載其他.
以色列兩個半支派的家譜.
再加多一點點的事蹟.
西支派等等.
這兩個,不過這兩個支派.
是首先被魯巴比倫的.
所以在這裡面.
先講了這兩個.
第六章一節到八十一節.
就講利未濟斯的家譜.
接著.
剩下來的就是其他支派.
第七章一節.
八章四十節.
這兩個篇幅.
記載著其他.
以色列家的家譜.
其他支派.
最後,第九章一節到.
四十四節.
記載著猶大家.
被魯巴比倫之後.
再回到耶路撒冷記載著的.
那個家譜.
家譜就是.
有九章的經文.
夾雜少少的敘事.
但大部份都是人名.
記載著整個.
脈絡.

$^{201}$不過你剛才這樣讀你會看到有少少.
有少少東西你看到.
第一就是他講猶大的篇幅.
是不是講多一點.
他其實特顯就是在講猶大.
支派裡面.
十二支派裡面的大阿哥.
其實如果按照.
生育的次序.
猶大排第四.
他為甚麼會成為大阿哥呢.
歷代智去解釋給你聽.
他成為了.
君王之後.
在這裡他用家譜.
來呈現了這個狀況.
除此之外.
除了在講猶大.
整個家庭的家族.
他接著再講到猶大.
當中又再縮一縮窄.
就講到大衛和所羅門.
那個家譜.
特別講那個君王.
帝王的後代那條線.
因為猶大都幾個兒子.
不是每個兒子都是做皇帝.
猶大有五個兒子.
接著一直.
源遠流長去到大衛那時.
只是其中一個支派.
成為了這個家族裡面.
就是君王的支派.
其他的.
記載在其他的家譜.
剛才那個記載上.
所以在這裡你會見到.
猶大家.
和大衛家.
是一個重要的記載的主要內容.

$^{241}$其他的那些.
他都要按實如實.
來記載.
不過他的記載就變得好像很.
省略很快的.
不過在記載.
其他的家譜裡面.
其實整件事就講到一個.
好像整個以色列民族.
一個走向的.
一個民族歷史.
有皇朝.
但同時間後來一路一路.
急轉直下.
因為離棄上帝.
所以整個家庭.
或者整個民族就.
被擄被外邦.
直至到被擄之後.
再回到耶路撒冷.
歷代志的作者又再將這段歷史.
重新記載.
所以在這裡.
先拿到那個輪廓.
你要知道在這裡.
他講的是一個民族的歷史.
而在民族歷史.
講的人名裡面.
去到四章九節裡.
突然之間.
他就講到一個人.
猶大家裡面其中的.
一路源遠流長的一個後代.
這個人就叫做亞比斯.
四章九節裡面講.
亞比斯比他眾弟兄更尊貴.
他的母親給他起名.
叫亞比斯.
意思說我生他.
甚是痛苦.

$^{281}$在這裡很有意思.
在這裡裡面.
先解了這一句.
在這裡裡面.
當弘弘那麼多的名字.
出現成為家譜之後.
在這裡裡面講一個名字很獨特的.
叫做亞比斯.
這個名字再不是純粹.
記載一個人的名字.
誰生誰.
這個裡面他講多一點的背景.
給你聽.
原來亞比斯的出世.
他母親很辛苦.
他紀念他母親多過紀念亞比斯.
這裡裡面他特別講到.
因為他生他的時候特別痛苦.
不過你要留意.
坊間的書.
很多都會講.
因為他母親生得他那麼辛苦.
所以他的名字叫亞比斯.
亞比斯的名字就叫痛苦.
坊間有很多解釋是這樣講的.
不過.
我查了很多次.
原文其實不是這樣講的.
原文裡面.
痛苦那個字我寫給大家.
就叫做Azab.
即是A-Z-E-B.
是.
你沒有一個.
J字頭.
次序不同的.
而亞比斯就是J-A-B-Z.
是不是.
你見到兩個字是兩個不同的寫法.
痛苦是Azab.

$^{321}$亞比斯是Jabez.
我在找.
為甚麼有痛苦這個意思呢.
因為很多時候我們是望民生義.
因為母親說.
生他很辛苦.
所以兒子就叫辛苦.
我們很多時候是直線邏輯.
不過.
其實在原文裡面的字眼.
其實不是完全對等.
就將.
因為母親生他很辛苦.
所以就自動叫兒子做痛苦.
這樣的意思.
不過有很多聖經學者.
近來留意到.
其實將這個字Jabez這個字.
和痛苦這個字有兩個字母是.
調轉了的.
就是將那個Z.
和那個B字兩個調轉了.
然後再加上.
那個J字代表上帝.
這樣的意思.
有一些的.
聖經學家甚至說.
這個痛苦是.
母親改他的名字.
目的就是要將他的痛苦逆轉.
將那個B.
變為Z.
所以原本是痛苦.
生他出來很痛苦.
但因為這樣的時候.
母親更加有個寄願.
就是將這份痛苦逆轉.
加上那個J字.
你會見到的.
每一個聖經都有個J字.

$^{361}$其實就是一這個字.
一就是代表上帝.
上帝逆轉他的痛苦.
有這樣的文學去理解.
因為究竟.
Jabez究竟是一個什麼意思呢.
今天已經失傳了.
大概我們見到的.
另外出現的地方就是一個地名.
都是在歷代史上.
就是說一個地方名.
叫做Jabez.
其他的就沒有.
但有更加多的內容.
讓我們知道.
Jabez這個字.
究竟是具體說什麼意思.
不過如果按照剛才那個思路.
將那個B字.
和那個Z字.
反過來.
逆轉去說這個母親.
因為生產的時候很痛苦.
將這個兒子希望上帝逆轉去痛苦.
的這個心態.
似乎是說得通的.
這個是第一個我們需要留意的.
第二個需要留意的.
特別是坊間的書或者解經.
都是因為母親生得很痛苦.
很多時候就會順樂.
就改名叫痛苦.
所以Jabez一生就很痛苦.
然後所有死人的事都.
都在他身上.
來解釋給他.
這種解經的方法我們叫做Proof Text.
就是說我們先入為主.
我們有一個概念套入去經文裡面.
來理解.

$^{401}$不過.
經文不是想鋪陳一個這樣的意思.
告訴你.
在經文裡面從來都沒有繼續說.
Jabez很痛苦.
沒有說到他死人倒樓.
沒有說到好像《雨傘片》裡面.
屋漏更兼連夜雨等等.
沒有特別在說一些這樣的內容.
不過只是在說.
他的心願.
他要逆轉他生命裡面的.
那些痛苦.
在經文裡面反而.
強調的一樣東西.
如果孩子.
因為生得很痛苦的時候.
媽媽生他出來.
看他視為更加寶貴.
如果你用這個直觀來看.
你會選得明白很多.
坊間裡面的說法就是.
為什麼他這麼尊貴.
因為他懂得祈禱所以就很尊貴.
這個是那本書裡面所說的內容.
不過在這裡.
因為他生得很辛苦.
所以這個兒子很靈.
是不是這樣容易理解很多.
甚至乎在.
近來一些的《釋經》的說法.
都已經是用這個.
因為他生得很辛苦.
所以對這個兒子靈舍.
寵愛有加.
他看他相對於其他.
他出生的兄弟姊妹.
這個媽媽看他更加尊貴.
更加寶貴.
不過這個媽媽看他更加寶貴.

$^{441}$不代表他.
變得學壞.
反而是.
她的寄望就是將這個兒子.
將他生命的痛苦被逆轉過來.
用這個英文字.
Z和B.
在他的名字當中.
次序調換了.
讓他的生命活出一個.
不一樣的生命.
在這裡裡面.
其實那個名字代表著一生.
在歷代志裡面.
剛才我說過.
大部分的記載基本上只是一個人的名字.
連個生字都沒有.
所以讀上去很枯燥.
不過那個名字.
代表著一個人的一生.
一個人的一生某程度.
是受著他的名字.
就是一路帶著他.
人生要走未來的路.
就是.
大家每一個弟兄姊妹都有自己的名字.
你的名字是.
由父母對你的期望.
由我們父母裡面.
可能人生裡面面對的掙扎.
又或者在我們.
人生裡面.
如果你上一代是基督徒.
上帝給他一些特別的領袖.
按著這些領袖.
成為了你和我的名字.
我們的人生就被命名.
某程度上.
我們的人生就好像朝著我們的名字.
一起來走.

$^{481}$面前的路.
在聖經裡面.
我們解釋很多舊約的人物.
很多時候我們的起點是用什麼方法呢.
用他的名字來出發.
甚至乎有些人.
生命有一個轉捩點.
上帝要將這個人的名字.
給一個新的名字給他.
亞伯拉罕就是這樣.
亞伯蘭被重新命名為.
亞伯拉罕.
其實就是一個新一頁的開始.
一個新的人生.
一個新的命圖.
用名字來界定.
這個人的生活.
在經文裡面.
講到亞比斯這個名字.
被記載.
被改名.
然後有一個禱告.
在歷代志上.
特別講到猶太家.
好像是因為這個人的記載.
就好像發生了一些微妙的變化.
我們再看多一點.
你就會明白我說什麼.
我們看看歷代志.
上帝二章三至七節.
這個是在說猶太的家譜.
猶太家譜.
這個起點是怎麼記載的呢.
猶太的兒子是爾.
阿羅蘭.
和士拉.
這三人是嘉林女子拔書亞.
所生的.
猶太的長子爾在.
耶和華眼中看為惡.

$^{521}$耶和華就殺了他.
猶太的息父他馬.
為猶太生了法勒斯和謝拉.
猶太共有五個兒子.
猶太的兒父他馬.
給猶太生了法勒斯和謝拉.
猶太共有五個兒子.
法勒斯的兒子是希倫.
哈姆勒.
謝拉的兒子是森利爾.
探希萬甲國大拉共五人.
加米的兒子亞干.
這亞干在當滅的.
墓上犯了罪.
連累了以色列.
我們讀猶太二章.
第三至七節.
在這裡記載猶太的家譜.
開頭的時候怎麼記載呢.
他不是說猶太.
是眾子之首.
多麼英明神武.
在這裡記載猶太的家譜.
一記載就記載一些不光彩的歷史.
他講到猶太有五個兒子.
有三個是他正室所生的.
就是迦南女子拔書雅所生的.
大兒子叫爾.
第二兒子是俄男.
第三是謝拉.
爾因為行耶和華眼中.
看為惡的事.
所以上帝擊殺了他.
這裡沒有特別記載.
喪世紀有說.
他的二兒子俄男.
同樣都是犯.
耶和華眼中看為惡的事.
上帝都將他擊殺了.
所以猶太只剩下.

$^{561}$第三個兒子.
謝拉.
謝拉這個兒子.
原本想著因為他哥.
爾.
過世了.
他的遺孀他媽.
想著許配給拉仔.
原本給二兒子.
不想為他哥.
留後.
所以純粹.
和他發生性關係.
沒有和他留後.
以色列人的習俗.
所以上帝就擊殺了俄男.
而謝拉猶太害怕.
跟第三個兒子都是死了.
所以先保留著.
保留了一段日子的時候.
他媽想著婚期都到了.
兒子都長大了.
但是猶太好像無動於衷.
忽然間在《禪寶記》記載.
猶太見到一個女孩.
很漂亮.
原來那個就是他的媳婦.
不過化妝成為妓女一樣.
猶太就當她是妓女.
和她發生關係.
於是乎在聖經裡.
記載著.
猶太的息父他媽.
為猶太生了法律師和謝拉.
一生就生了一對媽出來.
所以在這件事裡.
其實都很不光彩.
很難看.
如果說作為記載著一個帝王的家譜.
一個長子的家譜.

$^{601}$起始點.
起始點都很低.
都是食物鏈底層.
相對於其他.
他的十一兄弟.
猶太其實.
他的起點.
說不上很有風範.
作為長子的.
那種風範的生命.
然後經文不單止.
記載這件事.
就是說猶太那五個兒子.
原配死剩一個.
後來在他的媳婦裡.
生多兩個.
然後再繼續說.
他下來的還有一些後代.
其中一個叫阿干.
阿干在當滅的墓上.
犯了罪.
阿干褻瀆上帝.
然後連累整個.
以色列人都被上帝.
來降禍.
等等.
幾節的經文.
其實你會看到猶太支派.
其實是一個不光彩的.
一個支派.
不過這個起點.
就是成為了在.
歷代至上開頭的時候.
記載的.
我們看回在最後的時候.
歷代至.
上帝四章.
二十一至二十三節.
就是加普最後記載猶太支派的結局.
或者記載著.

$^{641}$那一刻支派的.
段落結尾.
是怎樣結尾呢.
四章二十一節.
猶太的兒子是.
史拉.
史拉的兒子是利加的.
早爾馬利沙之.
早拉大和屬亞實比族.
即世麻布的.
國家.
還有約敬哥西巴人.
約亞斯薩拉.
就是在摩亞地掌權的.
又有亞叔.
比利恆.
這都是古時所記載的.
這些人都是瑤將.
是尼他應和基拉.
基底拉的居民.
與王同住.
為王做工.
你會看到.
結尾結束.
都有一點點成就.
一個挺好的結尾.
是一班.
挺好的結束.
成為了.
在第四章裡面.
被擄之前說的.
猶太家族.
已經成為了一個.
本來很差的起點.
到尾的時候.
是挺好的.
未至於君王.
因為君王在大衛記載.
但君王之後.
這班人都是不錯的中上產.

$^{681}$未至於整天都要偷騙拐騙.
來渡過日子的.
一班.
十惡不赦的人.
在這裡面.
中間記載著亞比斯的.
祈禱.
在近代.
很多的《釋經》裡面.
特別記載著.
亞比斯的禱告.
某程度上.
將整個家族.
在一個很不濟的處境當中.
他好像重新.
扭轉了整個家族的命運.
我們還有.
一樣東西需要留意.
其實在記載裡面.
剛才記不記得.
猶太有五個兒子.
大兒子,二兒子都死了.
剩下的就是.
和他媽所生的.
法勒斯和謝拉.
而法勒斯和謝拉.
就是大衛的後裔.
在那個媳婦所生的.
你值得留意的就是.
如果按著次序.
其實猶太的三個兒子.
應該要先記載的.
不過.
歷代史的作者不是這樣記載的.
第一,第二兒子死了.
然後就說和媳婦所生的.
按次序.
那兩個是小一點的.
不過就將這兩個帶出來.
就成為了君王之家.

$^{721}$然後到最後.
在這裡所說的.
一至二十三節就說到第三個兒子.
謝拉.
謝拉那裡的記載.
就在這裡裡面作為一個補充.
在次序上.
在那個時序上.
其實本身不是順序的.
不過在那個描述上.
聖經的作者.
很想說到.
這個家族被扭轉.
被改換了命運.
他刻意用亞比斯禱告.
來說.
他的家庭改變.
一個很簡單.
由一個罪惡的.
十惡不赦.
一個不理想的家庭.
改變為成為了一個.
上帝祝福的長子之家.
好了.
用這個角度來看.
第四章第十節你就會明白.
開始你會多一點理解.
亞比斯的祈禱.
亞比斯求告耶和.
以色列的神說.
保佑我的境界.
尚與我同在.
保佑我不遭患難.
不受艱苦.
神就應運他所求.
其實這段經文裡面.
兩個層面.
消極層面和積極層面.
消極面是說什麼呢.
在這裡和合本的譯法有一點.

$^{761}$令到我們未必掌握到.
他的那個.
消極面就是求上帝的手.
保護他.
保佑我不遭患難.
求上帝的手.
保佑他.
讓他不會.
沾染在.
惡事當中.
那個患難的字.
就是惡事的意思.
邪惡的事.
一些不好的.
道德上.
或者際遇上.
但這裡特別講的是道德的事.
而這個就針對.
猶大家族裡面所做的行為.
或說眼中看為惡的事.
亞比斯在這裡.
去到這一代.
他開始醒悟.
求上帝幫助他.
上帝的手淋在他身上.
遠離這些惡事.
以致到他們.
能夠不會受到痛苦.
因為做惡事.
面對痛苦.
在舊約裡面.
就有一個很重要的.
一個聖經的教導.
你行善的.
行正的事.
上帝就賜福給你.
行惡的.
你就會面對和患.
面對痛苦.
在這裡亞比斯的祈禱.

$^{801}$就是求上帝幫助他.
消極層面就是禁戒他.
遠離這些惡事.
讓他脫離痛苦.
積極的層面是甚麼呢.
就是求上帝賜福給他.
那個境界.
不是他的層次.
令他靈舍高.
而是令他的地圖的疆界更加闊.
所以在這裡.
很簡單.
亞比斯的祈禱就是.
上帝呀求你幫我.
不要做壞事.
我就不用受苦了.
求你賜福給我.
我想我的地更加多一點.
其實亞比斯的祈禱.
就是一個這麼簡單的祈禱.
一個這麼簡單的祈禱.
然後.
我們也不用輕看它.
不過在這個歷代志的技術當中.
一個.
第一個好像被清醒的後代當中.
忽然間本來是.
走向下坡的整個家族.
因為一個覺悟一個覺醒.
以致到整個家族的命運.
就被扭轉.
被改喚過來.
在聖經裡面.
其實.
它特別強調的.
在歷代志裡面.
其實歷代志下.
就是在說一個聖殿的興起.
一個聖殿的建造.
其實是幫助整個以色列民族.

$^{841}$扭轉他們自己的命運.
本來他們每一個在上帝面前.
都是行惡被上帝憎惡的.
不過.
因為有這個聖殿.
因為這個禱告.
他們就離惡.
然後生命得到發展.
成為一群新造的人.
成為一群可以重新被上帝喜悅的子民.
我們再繼續看.
這個我是借助高銘軒博士.
他的觀點.
他說如果要理解亞比斯的祈禱.
你沒辦法不看歷代志下.
七章第十四節.
這個是在說歷代志裡面.
很重要的鑰節.
這個是所羅門的祈禱.
在聖殿興建完之後.
所羅門是這樣說的.
這稱為我名下的子民.
若是自卑禱告尋求我的面.
轉離他們的惡行.
我必從天上垂聽.
赦免他們的罪.
醫治他們的地.
整段經文在說什麼呢.
其實.
所羅門教百姓祈禱.
這個殿怎麼用.
這個殿不是給你觀賞用的.
這個殿是用來成為.
你們禱告.
一起去求上帝的一個地方.
倘若你遇到天有閉塞.
或者.
生活裡面很多艱難.
很多艱苦.
面對著很多困難.

$^{881}$你第一件事要問一個問題.
你生命有沒有離棄上帝.
當你自己發現你有自省的時候.
你就去到這個聖殿當中.
向神謙卑祈禱.
藉著一份最真誠的禱告.
上帝會聽的.
會幫助你.
當你尋求上帝的時候.
找上帝.
轉離你們的惡行.
上帝就會必然垂聽你們的禱告.
就會赦免你們的罪.
醫治你們的地.
整個禱告的中心點就是.
悔改.
當你面對著生活裡面.
有很多困迫的時候.
你第一件事要問的不是.
上帝怎麼作弄我.
而是問你生命有沒有一些.
要悔改的地方.
所羅門用這個聖殿.
告訴人聖殿是這樣用的.
如果你建了一個這麼宏偉的建築.
你不是用來祈禱.
你不是去到上帝面前.
尋求上帝.
你誤用了它.
你浪費了聖殿.
這個聖殿最重要的就是幫助人.
找回上帝.
尋求上帝的面.
轉離你的惡行.
以至上帝就醫治你.
赦免你的罪.
讓你重新得到一個.
得福的生命.
在這裡說的.
轉離他的惡行.

$^{921}$和剛才亞比斯所說的.
遠離惡事.
其實是同一個意思.
不免有災難.
不是純粹說.
純粹被動式.
求天上的神.
看顧.
給我們一些好祝福.
在這裡更加帶著一份自省.
反省今天自己生命的光景.
在整個亞比斯祈禱裡.
最值得我們找到重要的養分.
亞比斯和他的祖先不同.
亞比斯的祖先.
一直好像循著某一條罪惡的軌跡.
一直像循著某一條罪惡的軌跡.
一直越走下坡.
但是在亞比斯的禱告.
或者由他開始.
他要作一些扭轉.
他要對自己的生命謹慎.
自省的生命.
尋求上帝的福氣.
尋求真正的發展.
在他這一代,在他自己生命當中.
其實亞比斯只不過是一個小人物.
在聖經裡面說不上是一個很偉大的人物.
他不是做了一些很偉大的事蹟.
他只是在他當下.
眼前有些什麼狀況.
他只是在他當下眼前有些什麼狀況.
他只是想扭轉改變.
他只是帶著一份很單純的心.
你可以說他是一個很自我為中心的人.
不過如果連自我都不知道.
有時候祈禱都不知道是些什麼.
亞比斯的祈禱.
是一個小人物的祈求.
他是用自己的關懷開始.

$^{961}$但是他意識到他們自己生命的罪.
求上帝幫助他.
讓他遠離惡事.
更加讓惡事離開之後.
他尋求一些新的可能.
尋求一些新的發展.
在這裡我想補充多一點.
你會覺得地界擴展.
覺得他很市儈.
又或者覺得他很.
求的東西都不是很高質量.
不過其實亞比斯.
要知道自己求些什麼.
我想再說.
今天很多弟兄姊妹祈禱.
說很多屬靈.
套語.
說一些很宏大很偉大的字眼.
甚至最偉大心靈.
所說的話就成為他們的祖國.
我不是說他們錯.
不過當我們覺得.
一種意義泛濫的社會.
說一些很好聽.
動聽高賢大智.
說些很厲害的話.
但自己是想怎樣他不知道.
近來我有時.
跟一些年輕人聊天.
他跟我說一些很偉大的理想.
他說好啊.
年輕人.
跟著我說.
其實你想上帝怎樣幫你.
跟著他再跟我說很偉大的理想.
說這個世界有多少黑暗.
所以他要怎樣為上帝焚燒幾生.
令到生命.
這個世界可以怎樣改變.
甚至要改變整個香港教會的生態.

$^{1001}$跟我說這些.
我說那你打算怎樣做呢.
都是說不到的.
我們今天.
大量空泛的字眼.
禱告裡面充斥著所有.
我們覺得好聽的.
話語.
但我們自己想怎樣呢.
我們說不到.
我們不知道.
亞比斯跟我們不同的地方是.
他知道自己想怎樣.
雖然你覺得.
他是小人物的祈求.
但小人物的關懷.
永遠是最具體的.
他知道.
他想擴展他的家族.
擴展他的地界.
由一個很實際的祈求.
開始慢慢來延伸.
改變了.
他整個家族的關係.
又或者.
只是一些小小的改變.
上帝未必一次過將.
很偉大的英雄之地.
不同的邊疆賜給他.
不過可能每一代又多一些.
每一代又多一些.
一些小小的改變.
積累而成.
世世代代經年累月.
他們的家族就得到改變.
今天我們很多弟兄姊妹.
當我們信靠上帝.
我們生命裡面未必.
尋求一些很偉大.
講述一些很厲害的理想.

$^{1041}$我們都是形形異異.
在生活當中.
為著生命不同的事我們掙扎.
搏鬥中.
但這些不同的需要放在上帝面前.
我們講我們心底.
真是很想祈求的事.
上帝是會聆聽.
不過重點是你要離棄罪惡.
你要有反省.
自己生命的.
很重要的這一環.
如果你只是尋求上帝的祝福.
不去反省自己的生命.
大概你都是誤讀了.
艾比斯的祈禱.
當我們反省.
自己的生命.
去到上帝面前.
尋求改變尋求上帝的面.
我們生命的素質改變.
我們的生命一些很踏實的.
祈願會一步一步.
的來到去實踐.
我講多一個例子.
這是一個外國的研究.
我之前很久之前和大家講過.
不過我覺得值得今天再和大家延伸.
有本書叫做.
Jokis Edwards.
這本書是講什麼呢.
這本書是講同代裡面有兩個人.
叫馬克尤克斯.
一個叫愛德華茲.
Jonathan Edwards.
這兩個人其實是來自兩個.
不同的價值觀.
Jonathan Edwards 是一個基督徒.
是一個基督教領袖.
是一個報導家.

$^{1081}$他是一個有上帝的人.
而馬克尤克斯.
就是一個無神論者.
他是否定上帝.
他甚至對Jonathan Edwards.
他們的跟隨者說.
我不會信你這個上帝.
講完這番說話之後.
有些社會學家就追蹤研究.
追蹤研究200年.
看看Jonathan Edwards.
和這個馬克尤克斯.
200年之後.
他的家庭發生了什麼事.
這本書就記載著.
200年之後.
逐代逐代來計算.
看看他家裡發生了什麼事.
200年之後出現了什麼狀況.
愛德華茲.
Jonathan Edwards 的家庭.
他們整個家庭的人口.
200年之後.
有1394人.
其中這1300多人裡面.
有100個是.
大學的教授.
14個是大學的校長.
70個是律師.
60個是法官.
60個醫生.
60個作家.
300個牧師和神學家.
3位議員和一位副總統.
這就是Jonathan Edwards.
在200年之後的一個有神家庭.
他們的家庭的狀況.
馬克尤克斯呢.
200年之後是怎麼樣呢.
馬克尤克斯的家庭總數.

$^{1121}$200年之後.
903人.
其中有310人.
是流氓.
130位是坐監.
13年以上.
有7個是殺人犯.
100個是.
沉溺於酒的酒徒.
60個小偷.
190位妓女.
20名商人.
其中20名裡面有10個.
因為在監倉裡面學會怎樣營商.
所以做偏門生意的.
這個就是馬克尤克斯.
家族.
在沉淪的家庭之下.
他們面對的狀況.
我想說些什麼呢.
我們今天.
我們覺得.
上帝要有一些.
很極大的改變在我們生命裡面.
可能會有的.
不過上帝在我們每一個人.
生命當中.
在我們這一代.
信主的這個人.
可能改變是很漸進很緩慢.
甚至乎不知不覺.
但當我們帶著一份.
很切實的祈求.
求上帝幫助我們.
遠離生命的罪.
尋求.
生命的發展.
尋求一個積極的改變.
可能未必在我們這一代.
有大的改變.

$^{1161}$但帶著這一種這樣的.
血脈.
源遠流長世世代代.
你的後代會.
因著你的信仰.
帶著福氣.
倘若你已經是基督徒.
我盼望你守住你的信仰.
不要輕易.
就.
向上帝投訴你生命際遇的不濟.
無論怎樣都好.
切實地.
揹實上帝.
你生命.
不會.
在上帝的恩手裡面白過的.
在上帝的恩手當中.
上帝會保守你.
在你這一代裡面依然緊守.
讓我們的後代.
讓我們接下來每一個世世代代.
都堅守.
耶穌基督是我們生命.
要追求的榜樣.
追求的對象.
唯有這樣我們才為之叫做.
世世代代的相傳.
如果今天我們從雅比斯裡面.
學禱告我們要學會.
怎樣凌駁.
行善怎樣去尋求改變.
更加重要.
是我們後一代的福祉.
我們一起來勉勵自首.
一起尋求上帝.
(音量調整中).
\newpage



\section{撒母耳記上 1:4-20}
\label{sec:xitIkhwfBC8}
\textbf{2024年2月18日崇拜直播｜吳真真牧師｜十個深化關係的禱告-愁苦中的禱告｜撒母耳記上1\_4-20}
\newline
\newline
連結: \href{https://youtube.com/watch?v=xitIkhwfBC8}{\texttt{https://youtube.com/watch?v=xitIkhwfBC8}} ~~~~ 語音日期: 2024-02-18
\newline
\newline
\hyperref[sec:_RD54XIpk84]{\small{< < < PREV SERMON < < <}}
~
\hyperref[sec:index_chronic]{\small{[返順時目]}}
~
\hyperref[sec:1GcKc4lpPUQ]{\small{> > > NEXT SERMON > > >}}
\newline
\newline
撒母耳記上 1:4-20
\newline
\begin{longtable}{cl}
\hline
\hline
章節 & 經文 (和合本修訂版)\\
\hline
1:4 & \begin{tabularx}{0.7\textwidth}{X} 每逢獻祭的日子，以利加拿把祭肉分給他的妻子毗尼拿和毗尼拿所生的兒女。 \end{tabularx} \\ \\ \relax
1:5 & \begin{tabularx}{0.7\textwidth}{X} 他給哈拿的卻是雙分，因為他愛哈拿。耶和華卻不使哈拿生育。 \end{tabularx} \\ \\ \relax
1:6 & \begin{tabularx}{0.7\textwidth}{X} 她的對頭毗尼拿因耶和華不使哈拿生育，就常常惹她發怒，要使她生氣。 \end{tabularx} \\ \\ \relax
1:7 & \begin{tabularx}{0.7\textwidth}{X} 年年都是如此。每當她上到耶和華殿的時候，毗尼拿就這樣惹她發怒，以致她哭泣不吃飯。 \end{tabularx} \\ \\ \relax
1:8 & \begin{tabularx}{0.7\textwidth}{X} 她丈夫以利加拿對她說：「哈拿，你為何哭泣？為何不吃飯？為何傷心難過呢？有我不比有十個兒子更好嗎？」 \end{tabularx} \\ \\ \relax
1:9 & \begin{tabularx}{0.7\textwidth}{X} 他們在示羅吃喝完了，哈拿就站起來。祭司以利坐在耶和華殿門框旁邊的位子上。 \end{tabularx} \\ \\ \relax
1:10 & \begin{tabularx}{0.7\textwidth}{X} 哈拿心裡愁苦，就痛痛哭泣，向耶和華祈禱。 \end{tabularx} \\ \\ \relax
1:11 & \begin{tabularx}{0.7\textwidth}{X} 她許願說：「萬軍之耶和華啊，你若垂顧你使女的苦情，眷念不忘你的使女，賜你的使女一個子嗣，我必使他終生歸給耶和華，不用剃刀剃他的頭。」 \end{tabularx} \\ \\ \relax
1:12 & \begin{tabularx}{0.7\textwidth}{X} 哈拿在耶和華面前不住地祈禱，以利注意她的嘴。 \end{tabularx} \\ \\ \relax
1:13 & \begin{tabularx}{0.7\textwidth}{X} 哈拿心中默禱，只動嘴唇，聽不到她的聲音，因此以利以為她喝醉了。 \end{tabularx} \\ \\ \relax
1:14 & \begin{tabularx}{0.7\textwidth}{X} 以利對她說：「你要醉到幾時呢？不要再喝酒了！」 \end{tabularx} \\ \\ \relax
1:15 & \begin{tabularx}{0.7\textwidth}{X} 哈拿回答說：「我主啊，不是這樣。我是心裡愁苦的婦人，清酒烈酒都沒有喝，只在耶和華面前傾心吐意。 \end{tabularx} \\ \\ \relax
1:16 & \begin{tabularx}{0.7\textwidth}{X} 不要將你的使女看作不正經的女子。我因極其難過和生氣，所以一直禱告到如今。」 \end{tabularx} \\ \\ \relax
1:17 & \begin{tabularx}{0.7\textwidth}{X} 以利回答說：「平平安安地回去吧。願以色列的神允准你向他所求的！」 \end{tabularx} \\ \\ \relax
1:18 & \begin{tabularx}{0.7\textwidth}{X} 哈拿說：「願你的婢女在你眼前蒙恩。」於是婦人上路，去吃飯，臉上不再帶愁容了。 \end{tabularx} \\ \\ \relax
1:19 & \begin{tabularx}{0.7\textwidth}{X} 他們清早起來，在耶和華面前敬拜，就回去，往拉瑪自己的家裡。以利加拿和妻子哈拿同房，耶和華顧念哈拿。 \end{tabularx} \\ \\ \relax
1:20 & \begin{tabularx}{0.7\textwidth}{X} 時候到了，哈拿懷孕生了一個兒子，給他起名叫撒母耳，說：「這是我從耶和華那裡求來的。」 \end{tabularx} \\ \\
[1ex]
\hline
\hline
\end{longtable}
$^{1}$以下我們要聆聽神給我們的說話.
今日要聆聽神的說話.
記載在撒慕爾記上的第一章四至二十節.
撒慕爾記上第一章第二節.
第四節.
「以彌加拿每逢憲制的日子.
將際欲婚給他的妻兒.
毗尼拿和毗尼拿所生的兒女.
給哈拿的卻是傷份.
因為他愛哈拿.
沒久耶和華不使哈拿生育.
毗尼拿見耶和華不使哈拿生育.
就作他的對頭.
大大激動他.
使他生氣.
每年上到耶和華殿的時候.
以彌加拿都已傷份給哈拿.
毗尼拿仍是激動他.
以致他哭泣不吃飯.
他丈夫以利加拿對他說.
哈拿啊!你為何哭泣不吃飯.
心裡愁悶呢?.
有我不給十個兒子還好嗎?.
他們在事牢吃喝完了.
哈拿就站起來.
祭司以利在耶和華殿的門框旁邊.
坐在自己的位上.
哈拿心裡愁苦.
就痛哭痛痛哭泣.
進入祈禱耶和華.
許願說:萬君知耶和華.
你若誰顧婢女的苦情.
眷念不忘婢女.
賜我一個兒子.
我必使他終身歸於耶和華.
不用剃頭都剃他的頭.
哈拿在耶和華面前不住的祈禱.
以利定睛看他的嘴.
原來哈拿心中脈搏.
只動嘴唇不出聲音.

$^{41}$因此以利以為他喝醉了.
以利對他說:.
你不要醉到何時呢?.
你不應該喝酒.
哈拿回答說:主啊!不是這樣.
我是心裡愁苦的婦人.
清酒濃酒都沒有喝.
但在耶和華面前傾心吐意.
不要將婢女看作不正經的女子.
我因被人激動受苦太多.
所以祈求到如今.
以利說:你可以平平安安地回去.
願以色列的神允許你.
向他所求的.
哈拿說:願婢女在你眼前蒙恩.
與時婦人走去吃飯.
而尚在不戴壽容了.
次日清早他們起來.
在耶和華面前敬拜.
就回拉瑪.
到了家裡.
以利加拿和妻哈拿同房.
耶和華顧念哈拿.
哈拿就懷孕.
一其滿足.
生了一個兒子.
給他起名叫薩姆爾.
說:這是我從耶和華那裡求來的.
聖經讀筆.
各位影子妹.
新年蒙恩身體健康.
今天是年初九.
所以在這裡恭祝大家.
不知道大家會不會每一年新年的時候.
都期望有個新開始.
我們都會有些習俗.
將一些壞事掃走.
迎來新的福氣.
不過大家不知道有沒有經歷過.
就是有些艱難的事情.

$^{81}$會跨年陪著你.
舉例.
12月31日大家回來聚餐的時候.
大家看不到我.
為什麼呢.
因為COVID-19在28日陪著我.
陪到1月3日.
我跨年都中了COVID-19.
這些都是小病.
但是我都會發覺.
有些時候.
我們有些很艱難的事情.
是跨著年陪著我們經過.
不是一年.
有時是兩年.
過完套年.
去到農年.
農年出年是蛇年.
都跨著陪著我們.
近年在黃泉牧會的時候.
我都和很多在艱難的弟兄姊妹.
一年一年這樣經過.
我真的有時會.
很希望新年的祝願.
是真真正正的發生.
我不知道你們有沒有想像過.
那些困難難處.
我連大年初一.
一起床的時候就沒有了.
想不想.
想吧.
我都很希望大年初一的時候.
大家的恭賀.
或者我所記掛的人.
他們都真的平安喜樂.
身體健康.
一家和睦.
豐豐足足.
我真的這樣期望.
不過事實上.

$^{121}$今天我們讀的這段經文.
是一個現實的反映.
在《撒慕爾記》上.
這一章裡面.
記載著一個婦女叫做哈娜.
我讀這段經文的時候.
最深的印象就是.
仇婦這個字.
張卓穆斯.
新年樓樓跟我說這個字.
但我真的很想看看.
如何在仇婦裡面.
有些轉化.
我們看看她如何在神裡面找到出路.
我們一起去看看.
第一段我們會看的經文.
就是在四至八節.
第十和第十五第十六節.
其實在《撒慕爾記》上.
這一章的聖經.
很重要的背景.
其實這是一個.
事司時代的背景.
大家記不記得我們去年.
讀了很多事司.
那時候是一個黑暗的時代.
那時候的人.
是隨己而行.
隨己而行是他們的代表.
的句子.
他們經常說他們沒有王.
其實自己就是自己的王.
他們隨自己而行.
還有一樣東西就是.
他們的宗教都是亂七八糟的.
不知道大家記不記得.
他們會.
當時的人滿天神佛.
外族人的偶像他們照拜.
然後他們會將.

$^{161}$耶和華神帶回家.
自己有祭司自己去供奉.
其實他們的宗教.
他們的道德.
全部都是很敗壞.
而在這個時候就有一個.
這樣的女子哈娜.
在《撒慕爾記》上的第一章.
其實在整本書的研究裡.
其實最主要就是記載.
撒慕爾的出世.
撒慕爾的出生.
是很特別的.
因為撒慕爾出生的記載.
是很獨特.
就算撒慕爾記載了一個.
很重要的人物.
大衛王.
聖經也沒有記載他的出世.
但是在這裡.
他在這裡很細緻地記載.
其實撒慕爾是怎樣來的呢.
原來是因為他有一個.
很不簡單的媽媽.
一個很非凡的媽媽.
因為他的出世.
是一個神蹟.
是源自他一位敬虔的媽媽.
在一個這麼混亂的時代.
在一個這麼黑暗的時代.
這個女子有什麼特別的地方.
值得我們去學習呢.
我們來看看.
第一章裡.
我們看到.
聖經裡面有很多.
短短的十幾節經文.
我們已經分別在不同的節數裡.
去看到.
很多關於仇婦的字.

$^{201}$譬如我們會看到.
第七節.
第八節.
哭泣仇門.
然後第十節.
哈拿心裡仇婦.
痛痛哭泣.
我們又看看第十五節.
第十六節.
他自己心裡仇婦.
又說仇婦太多.
我讀這段經文的時候.
不知道會不會讀到你也受邪.
我真的讀到我也受邪.
這裡的形容就是.
剛才我們在唱詩歌.
我們聽到詩歌說.
我們被恩典包圍.
但我覺得哈拿是被仇婦包圍.
在她的生命裡.
全部被仇婦包圍.
究竟為何她那麼仇婦.
原來我們看第五節.
第五節說了一句話.
就是哈拿有一樣東西.
第五節下.
無奈耶和華不使哈拿生育.
哈拿的仇婦是她生不了的.
聖經的直譯是.
耶和華關閉了她的子宮.
是耶和華令她沒有兒子生.
我想在這個時代.
大家覺得生不出就算了.
我現在做分錢輔導.
以前十對有八對說會生孩子.
九對說會生孩子.
我現在做分錢輔導.
十對有五對說不生.
這個時代變了.
但大家要知道那個時代.

$^{241}$傳宗接代很重要.
當時以色列人.
他們本來是一夫一妻的制度.
但如果妻子不能生育.
不能幫她傳宗接代的話.
她就會納妾.
她就會找一個他們的婢女.
或者妾侍.
幫她生下一代.
其實老公要跟人分.
是一個很大的羞恥.
哈娜的身份.
變得是一個很卑微的身份.
是一個很被人看低的身份.
很多時候我們都會經歷一樣事.
有很多不好的事發生在我們身上.
不是我們可以控制的.
哈娜不能生不是她可以控制的.
她不能生是鐵一般的事實.
是一個對她來說很殘酷的事實.
但是這件事是無法控制的.
所以我們有時仇苦仇煩.
是因為一些我們無法控制的事.
一些不好的事降臨在我們身上.
另外第六第七節.
我們看到當哈娜不能生.
伊麗嘉娜的丈夫娶了一個妾侍.
就是皮莉娜.
娶了皮莉娜之後她就生了.
我們看到第六第七節.
你看到這個人怎樣對待她.
怎樣對待哈娜.
聖經形容她是她的對頭.
形容她是她的競爭對手.
她還要做什麼.
她還要大大的去惹怒她.
激動她.
激動她而且不是一次.
特別是每次上去獻祭的時候.
都要激動她.

$^{281}$激動這個字是用來形容動物.
去捕食獵物的行為.
所以皮莉娜根本就當哈娜是獵物.
是要挑釁她.
是要拼湊她.
然後令她生氣.
生氣這個字不是惹笑.
生氣這個字的原文是說雷聲.
你會想像哈娜被皮莉娜的這些舉動激到瘋了.
激到情緒失控.
我們經常都以為.
在我們日常生活裡.
很多時候有些人我們不能避開.
有些人的確是在攻擊我.
有些人的確是因為不同的緣故在挑釁我們.
我們不能避開.
我們於是就產生情緒.
其實我們經常以為生氣很可怕.
哇 真是哈娜生氣.
哈娜生氣到爆炸.
咆哮.
但是你知不知道呢.
或者另外一樣就是我們又以為.
情緒低落的人只會一味的逃避.
抑鬱的人只會躲起來.
但並不是的.
你從哈娜這個這麼仇苦的人.
你會看到一樣東西就是.
其實很生氣很生氣.
只是她表面的情緒.
她底層被人挑到激氣的時候.
她底層是一種很傷心.
很悲哀和很無助的情緒.
抑鬱的人有時都會很憤怒.
是因為他們的內心很無助.
他們要將他們內心那種.
說不出的怒氣.
說不出的憂鬱去發洩.
所以有時我接觸一些社區的婦女.
她們會很生氣很無助的時候.

$^{321}$她們會打小孩.
不是她們狂躁.
其實可能她們是抑鬱.
她們是無助.
她們是內心傷痛到很無助.
我也聽過一些兄姊妹.
她們真的經歷過.
當她們經歷了很多艱難.
很多仇苦的時候.
她們連禱告都不懂.
我不知道怎麼禱告.
有沒有聽過?.
有沒有聽過人這樣分享?.
甚至我們看回第五節.
我們看到她們心裡悠悶.
她們心裡抑鬱.
你會發覺甚至沒有人可以幫她.
身邊的人幫不了她.
我們看看第五節.
聖經告訴我們.
第五節她的丈夫伊利加納很愛惜她.
他的祭物說他給她雙份.
有些翻譯不是形容是雙份.
是將上好的一份給她.
總之伊利加納很愛惜他太太.
不過他這份愛惜起不了作用.
甚至你可以想像.
可能他這些舉動.
這些對太太的愛惜.
更加挑起皮利娜對他的妒忌.
以至對他的攻擊越演越烈.
然後第八節.
伊利加納不單做到出面.
他給他太太雙份或上好的一份祭物.
他還開口說了一句.
有我不是比有十個兒子更好嗎.
他的意思是.
就算你有十個兒子.
他們對你的愛都不及我給你那麼多.
但對哈娜一點幫助都沒有.

$^{361}$因為你見不見他前面那句.
你為何哭 你為何不吃飯.
你為何心裡受煩.
你為何心裡受悶.
你為何那麼傷心.
這句話的意思是告訴我們.
你為何面帶愁容.
你的心為何不舒暢.
我對你那麼好.
為何你不舒暢.
這句話不是聖經說的.
是我說的.
男人很多時候想解決問題.
我給你雙份 你開心了.
但其實太太心裡最需要的那份.
是雙份的祭物嗎.
不是.
她心裡的哀愁和情緒還未解決.
所以她就哭了.
所以你就哀傷了.
所以她不想吃飯.
哈娜這份內心的受苦.
連最愛她的人.
她的丈夫都無法解決.
無法補救.
領子妹.
我相信我們每一個人.
就算你生命怎樣順風順水也好.
我們的生命都一定有些控制不了的事情會發生.
我們不想出現.
不想發生的事情淋到我們身上.
哈娜她想生 但無法生.
我們會經歷.
有些弟兄姊妹會病了.
一發現就是重病.
病極度治不好.
或者我們家人有病.
或者我們人與人之間的關係破裂.
家人與人之間的關係很疏離.
這些事情我們好像很無助.

$^{401}$很控制不了.
又或者我們好像哈娜一樣.
除了控制不了的事情發生之外.
就是身邊有些人.
你很想避開他.
但你無法避開.
他可能對你有批評.
對你有挑釁.
他可能看低你.
踩你.
有些人很知道怎樣激怒你.
他知道怎樣按著你的按鈕.
我們生命裡面就是有這樣的環境.
有這樣的光景.
你有沒有?.
新年還沒有走到.
我們看看哈娜怎樣做.
哈娜第十節告訴我們.
哈娜祈禱耶和華.
祈禱耶和華.
她沒有還擊.
她沒有像聖經裡面.
與她同樣遭遇的殺奈.
即是母子三.
她叫亞伯拉罕.
你與我在耶和華面前定奪.
現在下格欺負我.
取笑我 輕視我.
亞伯拉罕說好.
你喜歡怎樣對待她就怎樣對待她.
然後殺奈就虐待她.
但是哈娜沒有.
她沒有這樣做.
聖經告訴我們.
她祈禱耶和華.
哈娜是一個祈禱的人.
在短短一章裡面.
多次去形容她是一個祈禱的人.
撒慕爾記說祈禱.
只有兩個人.

$^{441}$一個是哈娜.
第二個是她的兒子撒慕爾.
他們兩個是祈禱的人.
哈娜沒有其他辦法.
她不知道自己還可以做些什麼.
事實上她的情緒告訴我們.
她真的無能為力.
她面對一個她無法避開的人.
她亦面對一個無法改變的問題.
她唯有祈禱.
她是怎樣祈禱的呢.
她的祈禱.
她的祈禱是心裡受苦.
痛痛哭泣.
她心裡受苦的字.
是精神上的受苦.
不是肉體上.
是精神上的受苦.
痛痛哭泣的原文是哭泣再哭泣.
所以她哭得很痛.
很痛心的痛.
這是罪事者對哈娜的形容.
哈娜又怎樣形容自己呢.
第十五節.
我們跳到第十五節.
第十五節她形容自己是心裡受苦的婦人.
在她的字白裡.
她也說自己受苦.
不過這個受苦與罪事者形容的受苦不同.
這個受苦的原文是源自一個農業的背景.
是說她背負了一個壓.
意思是她在表達她受的苦有多重.
有多艱難.
有多苛刻.
聖經說.
她去到神面前.
她說她傾心吐意.
傾心吐意是甚麼呢.
傾心吐意是傾倒你的生命.
傾倒你的靈魂.

$^{481}$甚至是你的所有.
你的情緒.
你的慾望.
哈娜在神面前就是這樣的狀態.
其實哈娜.
甚至後來你會見到.
以你以為她喝醉酒.
她一直脈搏.
一直在動.
她不斷地禱告.
她沉醉在禱告裡.
她向著上帝禱告的時候.
她傾盡她的生命.
傾盡她的靈魂.
她給了我們一個很好的榜樣.
她與神祈禱提醒了我們一件事.
我們與神祈禱有時很沙啞.
不知道你覺得嗎.
我們與神祈禱.
未必將最真實的自己呈現.
我們的祈禱有時都想過,度過.
不知道大家會不會.
但神知道你.
你不用遮掩,不用裝飾.
我不用計算.
上帝你會不會應許我這個呢.
我這樣求是不是網求呢.
我們可能有很多計算.
這樣求是不是會成真呢.
我們不用怕.
因為其實我們怎樣.
神很清楚.
你去到祂面前傾心道義.
祂是會閱立的.
昨天我們有一群組長聚集在一起.
我們一起看今年的主題書.
《選擇快樂是否》.
我們很榮幸.
我們請了這本書的作者歐醒華牧師.
來到我們當中.

$^{521}$他講到第一課的內容.
他說了一件事.
他說誠實.
我們第一課只學一件事就是誠實.
就是誠實.
誠實對自己.
誠實對我們身邊的人.
誠實於我們的關係.
誠實於上帝.
牧師說你們要分組.
三個三個一組.
然後每組裡面.
你找的人.
都要是你信任的人.
是真的.
如果你跟那個人的關係.
你不信任的話.
你沒有辦法誠實.
因為你心裡有很多懷疑.
你心裡有很多擔心.
你需要一個很有安全感.
你覺得信任你才可以傾心吐意.
有時我們代人祈禱.
或者我們祈禱都很功能性.
我們好像.
我是神的僕人.
我們要很規矩地祈禱.
不過我們不要忘記.
我們同樣是神的女兒.
神的兒子.
向爸爸禱告.
是可以帶著你的情緒.
帶著你的感情.
帶著你的痛痛哭泣.
去到神面前祈禱.
建立一份真真實實的關係.
神有時.
感受不到我們裡面的.
不是感受不到.
是祂知道的.

$^{561}$不過祂有時真的逼到我們埋牆角.
有時真的讓我們無位走.
去面對最真實的自己.
面對最脆弱的自己.
面對最不堪的自己.
我們才去向神赤露塘開.
我們就會經歷這份.
神在心底裡的改變.
我們就會知道怎樣祈禱.
有句說話說.
哪裡有痛苦.
神的愛就在哪裡.
然後我們看第十一節.
哈拿許言說.
萬君知耶和華.
你若垂聽你使女的苦情.
眷念不忘你的使女.
賜你的使女一個子嗣.
我必使她終身歸給耶和華.
不用剃刀剃她的頭.
哈拿向上帝陳明了她的苦情.
之後她向上帝納了一個願.
她對著的是萬君的耶和華.
她對神的稱謂.
祂是超乎地上一切能力的上帝.
祂是為她的子民.
可以提供無窮之源.
滿有能力的上帝.
這是哈拿對上帝的一份認信.
然後她很清楚.
很清楚自己的身份.
她是很卑微.
她只是求主的憐憫和顧念.
她知道自己是一個卑微的請求.
相對當時代的人.
他們是捉著上帝的手.
去成就他們的心願.
但哈拿懷著一個最卑微的心.
期望上帝聽她的祈禱.
接著她納了一個願.

$^{601}$如果上帝憐憫.
就賜她一個兒子.
我就使她終身歸給耶和華.
不用剃刀剃她的頭.
我們看到哈拿納了一個願.
她納了一個拿世爾人的願.
為她的兒子納了一個拿世爾人的願.
這個願最重要.
民數記三十章告訴我們.
納願最重要最重要的是.
我們一定要償還.
哈拿納了這個願是什麼願呢?.
哈拿納的願是.
她要將她的兒子獻給上帝.
以利加拿不是利未人.
不是憑著血.
所以撒母耳不是憑著血裔.
去做祭司的後代.
但撒母耳是納了一個拿世爾人.
代表了她不屬於任何的支派.
不屬於她母親的兒子.
就是說.
哈拿將會有一個兒子.
拿世爾人這個字.
是希伯來文語納撒勒的意思.
是分離的意思.
就是說撒母耳將要與她.
將自己和她的家人去分開.
和她的民族去分開.
以至她能夠全心全意.
完全奉獻給神.
致力於上帝的工作.
我不知道你們有沒有聽過類似的話.
如果我....
如果上帝給我....
我就....
我聽過.
如果上帝給我兒子讀了這間學校.
我就回教會了.
你們有沒有聽過這些話.

$^{641}$有吧.
但哈拿這個說話不是這樣的.
她不是信口開河的.
她不是和上帝做交易.
她不是捉著上帝的手去達成她的心願.
你想想一個這麼多年不育的婦女.
她要將她的兒子.
做拿世爾人.
要讓他以後離開自己.
和自己分開.
專心服侍上帝.
這是一個很艱難很艱難的決定.
你這麼多年沒有....
如果我們真的沒有兒子可以....
沒有兒子可以生.
如果我沒有兒子可以生.
如果我可以生.
我一定會當他是金波羅.
你一定會給他所有的.
去滿足他.
我還記得我兒子做手術.
他做心漏手術.
他之前是很憂鬱的.
因為他的心怎樣去泵血.
他都沒有辦法去供應.
所以他消耗很多能量.
心臟有事的人通常會瘦又會尿.
所以他說通常做完手術之後心情好.
家人.
他說通常這些小朋友之後就會變胖.
因為他之前沒有吸收.
瘦每晚.
他一做好手術.
心臟供應OK.
家人就一定會給他飽一飽.
給他飽一飽去吃.
所以通常都變胖.
因為覺得欠了他.
給他最好的.
因為撿回了一條命.

$^{681}$什麼都給他.
哈拿將他一直渴望想生的兒子.
將他給了上帝.
作為神職人員.
他就是要放棄了.
他因為有這個兒子.
所帶來的好處.
為什麼他可以這樣做.
為什麼可以這樣做.
是因為其實在這個禱告裡面.
他的內心徹底改變.
這個就是祈禱.
在他的禱告裡面.
在他和上帝相遇的裡面.
在他和上帝交流的裡面.
他的焦點.
他的願望.
轉變了.
他的焦點.
是在上帝那裡.
他關注的不是他自己.
他關注的是上帝自己.
他關注的是上帝的心意.
他關注的是上帝的國度.
他關注的是上帝.
想在這個地上的工作.
我很同意Tim Keller說.
Tim Keller就是我們去年.
一起查士斯記那本書的作者.
他這樣說.
他說在哈拿禱告之前.
他的兒子是他終極的目標.
而上帝就是他達成終極目標的一個媒介.
但是當他向上帝祈禱.
傾心討義的時候.
上帝才是他的終極目標.
而他的孩子就成為了.
服侍上帝的一個媒介.
原來當你坦誠地向上帝傾心討義的時候.
你的心意是會改變的.

$^{721}$然後我們看到第十二節.
到第十六節.
我們看到有一個人物出現.
就是以利.
以利是當時的祭司.
他看到哈拿祈禱.
他看到哈拿這樣脈討.
他以為他喝醉酒.
酒他說他.
你究竟要醉酒到什麼時候.
你不應該喝酒的.
哈拿就回答說.
主啊.
第十五節.
主啊 不是這樣.
我是心裡仇苦的婦人.
我清酒濃酒都沒有喝.
這也是一個拿西人應該有的規則.
但是在耶和華面前傾心討義.
不要把婢女看為不正經的女子.
我因被激動.
受苦太多.
所以祈求到如今.
因為他的祈禱.
一直不停地說.
祭司就以為他是一個醉酒的人.
我不知道當哈拿.
是一個之前覺得羞愧.
很屈辱的女子.
被當時的祭司這樣說她.
其實她是否再受委屈.
被他看為不正經的女子.
我們說二次傷害.
再令他委屈多一點.
但是你看看.
你看看哈拿.
他沒有再說.
你冤枉我了.
他沒有.
他很清晰很清醒地跟這位祭司說.

$^{761}$他來到這個殿的目的.
就是為了祈禱.
你想想一個受委屈.
現在被人加多一腳.
誣蔑他喝酒.
說他不正經.
我想受苦真的多了幾分.
如果是我.
可能已經破口大罵.
罵他.
你做什麼祭司.
但他竟然出奇地冷靜.
向這個祭司表達憂傷.
我深信.
是因為神.
在這個禱告裡.
已經醫治了他.
他不再覺得自己羞愧.
就算你的人.
在我面前說我.
我不覺得羞愧.
我很清楚地告訴你.
我經歷的是什麼.
大家要留意.
第十六節被人激動的字.
其實原文的翻譯.
更加貼切.
是極其難過和生氣.
他不是跟祭司.
以利去.
控訴說有人激我.
不是.
是說.
他用的字.
這個字在約伯記經常出現.
就是我.
其實有苦情.
我很哀情.
我有哀告的意思.
而且當你醉酒到什麼時候.

$^{801}$他告訴你.
我祈禱到如今.
我覺得很奇妙.
這個對答.
很清晰.
知道自己經歷的是什麼.
知道自己的情緒是什麼.
知道自己在做什麼.
他比以利清醒.
比祭司清醒.
而且我相信.
是因為上帝親自醫治了他.
好了.
祈願圖.
我們看看.
這位小女子.
她如何蒙神的眷顧和恩惠.
在第十七節說.
她說.
以利回答說.
平平安安回去吧.
願以色列的神允許你向祂祈求.
哈拿說.
願你的婢女在你眼前蒙恩.
於是婦人上路.
吃飯.
念上不再戴袖蓉了.
哈拿何時不再戴著袖蓉呢.
就是.
現在.
就是現在.
經過一輪解釋.
以利說奉神的名.
祝福你.
平平安安回去吧.
願你的祈求得到神的允許.
這是祭司的職份.
民數記六章告訴我們.
上帝吩咐摩西.
叫阿倫和他的子孫.

$^{841}$祝福他的子民.
讓耶和華的平安賜給他們.
要奉神的名去祝福他們.
賜福給他們.
祭司是沒有權.
也沒有能力去叫哈拿.
所求的成就.
不要想著.
找陳明泉牧師或曾珍牧師.
幫你祈禱.
你的事就一定會靈.
千萬不要.
不過以利做了.
最基本的事.
他作了神的吩咐.
祝福他的子民.
這是我自己很深的.
反省.
對於祭司以利的祝福.
因為.
我們做牧師.
有時為弟兄姊妹祈禱.
都會很小心.
因為有句說話很綑綁我們.
二人的禱告大有功效.
如果我幫你祈禱.
我祈禱的那句不成.
我是不是二人呢.
我怕我不是.
所以.
我自己一個很大的反省.
祈禱不是我想過.
度過你癌症第二期.
我就祈你痊癒.
第四期就祈什麼.
我們有時會陷入這個陷阱.
我不應該用.
我的方法.
用我的判斷去幫你祈禱.
然後.

$^{881}$我就這樣禱告.
我.
真的很願意.
其實我們每一個弟兄姊妹.
要為人代求的時候都是這樣.
就是願神.
按你所求的賜給你.
然後願神賜你平安.
因為我們.
絕對相信一件事.
要給的是上帝.
主權在上帝那裡.
祂必然有祂的應許.
和祂的心意.
你信不信一個耶和華.
祂是太陽.
是盾牌.
祂要賜下恩惠和榮耀.
祂未嘗留下福氣.
不給那些行動正直的人.
你信不信這個應許.
你信不信這個確定.
這個確據.
就是說.
我仍然是求.
不過上帝就有主權.
即或如不如我們所求的成就.
我也很渴望.
上帝.
對那些向祂祈求.
信靠祂的人.
所賜下的是最好的.
對祂是好處.
這個就是.
我心所願.
所以當哈拿.
聽了以利的祝福.
他知道他蒙恩的時候.
你要記住.
他不是第十九第二十節.

$^{921}$回去友濟山.
他才不戴壽容.
吃飯.
他是聽完.
以利的祝福.
他知道自己在神面前.
已經蒙恩.
未成就的一件事.
但是他已經.
不再戴壽容.
他不再吃飯.
不再不吃飯.
他回去吃飯.
是不是.
所以弟兄姊妹.
我們看到第十九第二十節.
當.
撒慕爾出世的時候.
他要見證.
他要見證的一件事.
就是他見證.
撒慕爾是他求回來的.
是上帝的神跡.
撒慕爾的名字的意思.
是神的名字是伊勒.
這個意思是.
Elohim.
就是伊勒.
這個意思就是說.
神就是神.
神真是神.
撒慕爾是神跡.
是我向神求回來的.
他要去見證上帝.
他要去歌頌上帝.
哈拿按著上帝聽完他的禱告.
給他成就之後.
他沒有忘記上帝.
他真的去還願.
這裡沒有記載.

$^{961}$我們日曆經文沒有說.
不過二十六至二十八節很清楚地說.
三歲之後.
三歲是斷奶的日子.
斷奶的日子之後.
他真的帶了撒慕爾去見義利.
然後告訴他.
在這裡祈禱的我.
我帶我祈求回來的兒子.
回來交給你.
因為我要將他奉獻給耶和華.
要為上帝工作.
二章一節.
其實很豐富.
二章一節到十節的哈拿之歌.
其實是一段很豐富的禱文.
不過我們今天不說了.
但是他第一節.
正正回應他之前的仇苦.
弟兄姊妹.
他怎樣說.
我們一起讀第二章第一節.
哈拿禱告說.
我的心因耶和華快樂.
我的國因耶和華高舉.
我的口向仇敵張開.
我的心因耶和華高舉.
我的口向仇敵張開.
我因耶和華的救恩歡恩.
他歡恩.
他快樂的源頭.
不是因為神給他的兒子.
這裡說不是.
是耶和華本身.
因為耶和華就是他的上帝.
他知道他的需要.
他將他救拔.
他打敗他的敵人.
將他從卑微的身份升高.
他不再被他的仇敵.

$^{1001}$不再被毗尼拿.
不再被這個世界的定義.
將他踩扁.
將他屈辱.
將他打垮.
其實不知道大家記不記得.
我在幾年前.
分享過一樣東西.
我懷著第二個兒子的時候.
我曾經經歷.
作小產.
睡在病床上一段日子.
我記得第二天.
我去照.
我之前那晚.
我是瘋狂地哭的.
因為我很怕沒有了他.
詳細.
不知道有沒有聽到.
很久之前.
我真的和上帝.
有一個這樣的祈禱.
我說上帝啊.
無論這個兒子.
他有什麼事都好.
我很想生他出來.
我很想見到他.
我希望他成長.
我很想他出生.
我會按你的心意.
多艱難.
我都要養育他.
就算他.
我沒有說什麼.
大家知道了.
上帝知道我想說什麼.
結果呢.
兒子心臟病.
他要做心臟手術.
經歷了一個很大的手術.

$^{1041}$我們陪他度過那些日子.
但是.
故事還沒完.
我的兒子還在成長.
但我的兒子.
他有特別的需要.
SEN.
教學的特別需要.
我不詳細說了.
其實在這段青春期的日子.
我很辛苦.
很辛苦.
我有時候都好像哭了.
上帝啊.
可不可以我一說他就馬上改變.
我們可以輕鬆自如地應付他的生活.
但是呢.
上帝給我預備了這段經文.
我在上帝面前.
很受苦地祈求.
上帝會聽.
他也會回應.
他會讓我.
很深很確切地知道.
在人生的旅程裡.
他沒有留下一樣好傳給我.
這是我很確實的確據.
過去一星期.
我有機會去到一個聚會.
那個聚會裡.
剛好有三個傳道人.
我和另外兩個傳道人.
其實大家不是很熟.
不過大家.
都有一個共通點.
我們都是服侍孩子和他們的家長.
除了這個共通點之外.
我們三個傳道人的孩子.
都有特別的需要.
都是可能有病.

$^{1081}$都是可能.
無論是情緒.
或者是腦.
或者是SEN.
都有很艱難的東西.
我們分享到.
我們帶著他們真的很艱難.
但是呢.
我記得大家有異口同聲地分享.
就是因為我們去做這些工作.
家長的工作的時候.
我們去和家長去談.
我們和他們一起分享.
我們去聆聽他們.
我們一起禱告的時候.
不知為何.
我們就聯繫到了.
我們很真實地.
和他們聯繫到.
去連結到.
我們可以.
很多的機會.
去印證到上帝的.
祝福和在我們身上.
怎樣一次又一次.
去幫我們經歷那些難關.
所以弟兄姊妹.
我們今天學習.
哈拿的榜樣.
不是.
有些秘訣班.
你們想像哈拿那樣.
禱告成功就要有.
一二三四步.
對不起我今天不是教這些.
如果你想學就不要找我.
我們不是想.
禱告如何.
成就速成班.
今天我們不是教這些.

$^{1121}$今天我們學到的是.
從哈拿身上我們學到.
我們要坦誠地面對自己的脆弱.
我們向上帝.
坦誠地打開我們的心扉.
我們向他.
祈禱.
傾心吐意.
而且他一定會醫治.
他一定會滿足我們.
他會幫我們.
對照.
教正.
幫我們的生命裡面.
經歷到原來只要有耶穌.
我們的生命就得到.
拯救.
我們的心就得到滿足.
他會幫我們由受苦.
變為喜樂.
所以我會祝願大家.
我用.
二章一節這樣祝願大家.
我祝願大家你的心.
因耶和華快樂.
你因耶和華的救恩.
新念.
\newpage



\section{約書亞記 4:1-24}
\label{sec:1GcKc4lpPUQ}
\textbf{2022年9月25日崇拜直播｜陳啟興牧師｜以石為記:這些石頭是什麼意思?｜約書亞記4\_1-24}
\newline
\newline
連結: \href{https://youtube.com/watch?v=1GcKc4lpPUQ}{\texttt{https://youtube.com/watch?v=1GcKc4lpPUQ}} ~~~~ 語音日期: 2024-03-13
\newline
\newline
\hyperref[sec:xitIkhwfBC8]{\small{< < < PREV SERMON < < <}}
~
\hyperref[sec:index_chronic]{\small{[返順時目]}}
~
\hyperref[sec:8EXqdqFDCIc]{\small{> > > NEXT SERMON > > >}}
\newline
\newline
約書亞記 4:1-24
\newline
\begin{longtable}{cl}
\hline
\hline
章節 & 經文 (和合本修訂版)\\
\hline
4:1 & \begin{tabularx}{0.7\textwidth}{X} 當全民都過了約旦河，耶和華對約書亞說： \end{tabularx} \\ \\ \relax
4:2 & \begin{tabularx}{0.7\textwidth}{X} 「你要從百姓中選出十二個人，每支派一人， \end{tabularx} \\ \\ \relax
4:3 & \begin{tabularx}{0.7\textwidth}{X} 吩咐他們說：『你們從這裡，從約旦河中祭司的腳穩穩站立的地方，取十二塊石頭，一起帶過去，放在你們今夜住宿的地方。』」 \end{tabularx} \\ \\ \relax
4:4 & \begin{tabularx}{0.7\textwidth}{X} 於是約書亞召集了他從以色列人中所選的十二個人，每支派一人。 \end{tabularx} \\ \\ \relax
4:5 & \begin{tabularx}{0.7\textwidth}{X} 約書亞對他們說：「你們要過去，到約旦河中，耶和華－你們神的約櫃前面，按以色列人支派的數目，每人各取一塊石頭扛在肩上。 \end{tabularx} \\ \\ \relax
4:6 & \begin{tabularx}{0.7\textwidth}{X} 這些石頭在你們中間將成為記號。日後，你們的子孫問你們說：『這些石頭對你們有甚麼意思呢？』 \end{tabularx} \\ \\ \relax
4:7 & \begin{tabularx}{0.7\textwidth}{X} 你們就對他們說：『這是因為約旦河的水在耶和華的約櫃前中斷；約櫃過約旦河的時候，約旦河的水就中斷了。這些石頭要作以色列人永遠的紀念。』」 \end{tabularx} \\ \\ \relax
4:8 & \begin{tabularx}{0.7\textwidth}{X} 以色列人就照約書亞所吩咐的做了。他們按以色列人支派的數目，從約旦河中取了十二塊石頭，正如耶和華所吩咐約書亞的。他們把石頭帶過去，到他們所住宿的地方，就放在那裡。 \end{tabularx} \\ \\ \relax
4:9 & \begin{tabularx}{0.7\textwidth}{X} 約書亞另外把十二塊石頭立在約旦河的中間，在抬約櫃祭司的腳站立的地方；直到今日，石頭還在那裡。 \end{tabularx} \\ \\ \relax
4:10 & \begin{tabularx}{0.7\textwidth}{X} 抬約櫃的祭司站在約旦河的中間，直到耶和華命令約書亞告訴百姓的一切事辦完為止，正如摩西所吩咐約書亞的一切話。於是，百姓急速過了河。 \end{tabularx} \\ \\ \relax
4:11 & \begin{tabularx}{0.7\textwidth}{X} 全體百姓都過了河之後，耶和華的約櫃和祭司才過去，到百姓的前面。 \end{tabularx} \\ \\ \relax
4:12 & \begin{tabularx}{0.7\textwidth}{X} 呂便人、迦得人、瑪拿西半支派的人都照摩西所吩咐他們的，帶著兵器在以色列人的前面過去。 \end{tabularx} \\ \\ \relax
4:13 & \begin{tabularx}{0.7\textwidth}{X} 約有四萬帶兵器的軍隊在耶和華面前過去，到耶利哥的平原，準備上陣。 \end{tabularx} \\ \\ \relax
4:14 & \begin{tabularx}{0.7\textwidth}{X} 在那日，耶和華使約書亞在以色列眾人眼前被尊為大。在他一生的年日中，百姓敬服他，像從前敬服摩西一樣。 \end{tabularx} \\ \\ \relax
4:15 & \begin{tabularx}{0.7\textwidth}{X} 耶和華對約書亞說： \end{tabularx} \\ \\ \relax
4:16 & \begin{tabularx}{0.7\textwidth}{X} 「你吩咐抬法櫃的祭司從約旦河上來。」 \end{tabularx} \\ \\ \relax
4:17 & \begin{tabularx}{0.7\textwidth}{X} 約書亞就吩咐祭司說：「你們從約旦河上來。」 \end{tabularx} \\ \\ \relax
4:18 & \begin{tabularx}{0.7\textwidth}{X} 抬耶和華約櫃的祭司從約旦河中上來，腳掌一落乾地，約旦河的水就流回原處，仍舊漲滿兩岸。 \end{tabularx} \\ \\ \relax
4:19 & \begin{tabularx}{0.7\textwidth}{X} 正月初十，百姓從約旦河上來，就在耶利哥東邊的吉甲安營。 \end{tabularx} \\ \\ \relax
4:20 & \begin{tabularx}{0.7\textwidth}{X} 約書亞把他們從約旦河取來的那十二塊石頭立在吉甲， \end{tabularx} \\ \\ \relax
4:21 & \begin{tabularx}{0.7\textwidth}{X} 對以色列人說：「日後，你們的子孫問他們的父親說：『這些石頭是甚麼意思呢？』 \end{tabularx} \\ \\ \relax
4:22 & \begin{tabularx}{0.7\textwidth}{X} 你們就讓你們的子孫知道，說：『以色列人曾走乾地過這約旦河。』 \end{tabularx} \\ \\ \relax
4:23 & \begin{tabularx}{0.7\textwidth}{X} 因為耶和華－你們的神在你們前面使約旦河的水乾了，直到你們過來，就如耶和華－你們的神從前在我們前面使紅海乾了，直到我們過來一樣， \end{tabularx} \\ \\ \relax
4:24 & \begin{tabularx}{0.7\textwidth}{X} 要使地上萬民都知道，耶和華的手大有能力，也要使你們天天敬畏耶和華－你們的神。」 \end{tabularx} \\ \\
[1ex]
\hline
\hline
\end{longtable}
$^{1}$願意我們繼續以聆聽聲道.
繼續我們的敬拜.
今日我們選讀的經文是舊約聖經.
舊約聖經.
大家在教會上的聖經是第266至267頁.
我們選讀的是舊約的約書亞記第四章.
約書亞記第四章第十九至二十四節.
請聽我讀出.
正月初十日.
百姓從約旦河裡上來.
就在吉格.
在耶利哥的東邊安營.
他們從約旦河中取來的那十二塊石頭.
約書亞就立在吉格.
對以色列人說.
日後你們的子孫問他們的父親說.
這些石頭是什麼意思.
你們就告訴他們說.
以色列人曾走乾地過這約旦河.
因為耶和華你們的神.
在你們前面使約旦河的水乾了.
等著你們過來.
就如耶和華你們的神.
從前在我們前面使紅海乾了.
等著我們過來一樣.
要使地上萬民都知道.
耶和華的手大有能力.
也是你們永遠敬畏耶和華你們的神.
今早很高興請到陳啟牧師為我們宣講.
他的題目是「以石為基」.
「這些石頭是什麼意思」.
有請陳牧師.
各位姐妹早安.
今天很感恩可以在大家當中.
繼續去一起學習神的話語.
但最重要的是我們傳一個謙卑的心.
我們去禱告仰望上帝.
親愛的主我們謙卑在你面前.
我們懇求聖靈你光照我們.
讓我們有一顆聆聽的耳朵.

$^{41}$開放的心靈.
讓我們的心是一個好土.
可以盛載你的話語.
結出很多的果實.
求主幫助.
我們不敢誇什麼.
我們只求你幫助我們.
聽你入耳.
去改變我們的生命.
我們仰望你.
奉耶穌基督的聖名.
阿們.
我小兒子叫做天成.
他胸口有一條疤痕.
有一條幾吋長的疤痕.
是他六歲的時候做了一個心臟大手術.
就留下了.
我相信這個疤痕.
這一生應該也會伴隨著他.
天成天生有一個心房中隔缺損.
我們叫做ASD.
其實這樣的情況是很普遍的.
就是左右心房中間的中隔穿了一個洞.
但是天成的問題就是.
他的洞就在中隔的上方.
這個就比較罕有.
通常在中間或者下方比較多.
他在上方就很接近一條靜脈.
就不可以用微創去處理.
就要做一個開胸的手術.
就是要剖開胸部.
做的時候要用bypass.
其實他做的時候六歲.
其實我也有想過這個念頭.
我可能會失去這個兒子.
這個六歲的人去面對這麼大的手術.
大神的恩典就保存了他的生命.
他運過了這一次的危機.
但是他的胸口就留了一條很長的疤痕.
對天成來說.

$^{81}$或者對我一家來說.
這個標記不是一個苦難的標記.
這個標記就是讓我們去紀念.
原來神真的是生命的主宰.
更加去敬畏他.
更加去尊崇他.
我們人身上可能多少也有疤痕.
我不知道.
如果你有疤痕.
這些疤痕對你來說.
有沒有重要性呢?.
還是你要設法將它.
用一些除疤膏去處理呢?.
對我來說最有意義的疤痕.
也不是天成身上的疤痕.
最重要最有意義的疤痕.
是耶穌手上的釘痕.
我想這些釘痕不僅對多瑪很重要.
對第一代的門徒很重要.
對今天我們每一個信耶穌的人都很重要.
這些釘痕就仿似耶穌復活的標記.
祂讓我們去紀念耶穌為我們的罪去受死.
但是祂勝過死亡.
祂勝過罪.
祂是我們的救主.
尊崇祂.
敬畏祂.
我們人生是需要有標記.
有時候是疤痕.
有時候是一些經歷.
這些標記幫我們去常常紀念神救贖的恩情.
神保守引領我們.
能夠過一個得勝人生的大能.
更加敬畏祂.
譬如在旺鎮.
每個月都會有聖餐.
聖餐是一個可見的記號.
一個標記.
一個聖禮.
是讓我們紀念主耶穌為我們受死.

$^{121}$祂成就了救恩.
幫助我們敬虔道逸.
等祂回來.
今天這段經文.
是講以色列人.
他過了四十年的曠野流浪.
他終於可以進入這個應許之地.
他們就立石為記.
去紀念上帝拯救的大能.
好像看得不太清楚.
可能我用錯了背景.
大家看不看得到.
其實也看得到.
如果不緊要.
看不到的話.
我就.
你照樣聽也可以.
其實這裡和上一章很有關係.
在上一章.
就是第三章的十四到十七節.
就講到以色列人過約旦河.
台約櫃的祭司.
他的腳一踏入約旦河的河水.
約旦河在很遠的地方.
那裡就是亞當城.
就立起城壘.
亞當城距離這些以色列人過約旦河.
就是耶利哥的對面.
有多遠呢?.
大約是三十公里.
不是很近.
三十公里以外的河水立起城壘.
這三十公里長的河水.
就一直向下流.
流入下面的死海.
就是阿拉巴海.
流完之後就露出河床的乾地.
以色列人就越過約櫃.
就過到約旦河的西岸.
這就是上一章.

$^{161}$尾段所講他們過河的情況.
你留意.
在上一章第三章最後那一節的最後那一句.
直到國民盡都過了約旦河.
和今天我們講四章一到二十四節的第一節的第一句.
很相似.
就是國民盡都過了約旦河.
幾乎一模一樣.
明顯就是聯繫上下文.
他們過河.
但是接著你看下去.
到了四章第十一節.
你發覺眾百姓盡都過了河.
就繼續這個故事的發展.
你會想.
四章一到四章十.
中間這一段.
發生什麼事呢?.
在寫作技巧上.
中間這一段.
我們會叫中段經文.
Interrupted Versus.
這是聖經的寫作技巧.
它的意思就是在中間加插進去.
這一段經文.
在四章一到四章十節.
不是過了河才發生的事情.
雖然第一節一開始說他們過了約旦河.
但是這一段是中段經文.
是加插進去的.
其實是在講過河前.
以及在過河所發生的事情.
所以你這樣看.
才可以明白四章一到十.
它在講什麼.
是作者特意加插進去.
很簡單.
揀選十二個人去搬石.
其實在三章十二節已經說了.
已經吩咐了.

$^{201}$所以在這裡是一個加插進去.
在這裡加插進去或者重複.
是有強調神吩咐約斯亞和以色列人聽命.
這樣的一個意味.
我解釋了一點背景.
開始進入經文的第一個段落.
就是第一到十節.
我們先看第一到三節.
他說國民進都過了約旦河.
耶和華就對約斯亞說.
你從民中要揀選十二個人.
每支派一人.
吩咐他們說.
你們從這裡從約旦河中.
祭司腳站定的地方.
取十二塊石頭帶過去.
放在你們今夜要住宿的地方.
很明顯這是上帝的吩咐.
就叫約斯亞找十二個人搬石.
十二個人.
每支派一個人.
是代表什麼呢?.
就是代表全部以色列人都參與其中.
就是整個群體.
當然根據民數記是六十萬人左右.
你不要管這個數字有多真確.
當是六十萬人.
就是能夠打仗的男人.
整個群眾可能隨時都有200萬人以上.
你不可以200萬人都下去搬石.
所以找十二個代表.
在哪裡搬呢?.
在祭司站立的地方搬石.
不是祭司厲害.
是祭司抬著的約櫃厲害.
或者更加清楚.
更加精準.
約櫃代表上帝的臨在.
所以其實也不是約櫃厲害.
是上帝很厲害.

$^{241}$神臨在其中約旦河水.
可以從中斷絕讓百姓過河.
所以他們要在代表神臨在的地方.
就是祭司腳掌所踏之地.
去撿石.
神如何吩咐約書亞呢?.
約書亞如何吩咐以色列人呢?.
是一個完全的聽命和跟從.
你看看第四到五節.
約書亞如何吩咐以色列人.
於是約書亞將他從以色列人中.
所預備的十二個人.
在三章十二節已經選擇了.
每個支派一人都照來了.
對他們說.
你們下約旦河中.
過到耶和華.
你們神的約櫃前頭.
按著以色列人十二支派的數目.
每人取一塊石剛在肩上.
很聽話的.
神如何吩咐約書亞呢?.
約書亞如何吩咐以色列人呢?.
那些石頭肯定不是石春.
肯定不是很小的石春.
也不會很大.
可以擔在肩上.
可能是幾十磅.
或是百多磅.
但肯定不會太大.
那些石頭有什麼意義呢?.
六到七節.
這些石頭在你們中間.
可以作為證據.
日後你們的子孫問你們說.
這些石頭是什麼意思呢?.
你們就對他們說.
這是因為約旦河的水.
在耶和華的約櫃前斷絕.
約櫃過約旦河的時候.

$^{281}$約旦河的水就斷絕.
這些石頭要作以色列人永遠的紀念.
「立石」有什麼意思呢?.
現在作者在這裡提供了一個.
相對簡單的解釋.
詳細解釋在後面.
22,23,24這三節.
「立石」是為了未來的人.
石頭不會說話.
但父母就會說話.
當後世子孫問.
為什麼要有石頭呢?.
父母就告訴他們.
告訴後世子子孫孫.
去紀念神令到河水中斷.
紀念神的拯救.
所以「立石」的第一個很明顯的意義.
就是紀念上帝的拯救.
上帝在你的生命裡面.
祂有什麼作為.
有什麼拯救呢?.
在這裡稍後最後的段落.
會多講一點「立石」的意義.
接著第八到第九節.
以色列人就照約書亞所吩咐的.
按照以色列人支派的數目.
從約旦河中取了十二塊石頭.
都遵耶和華所吩咐約書亞的行了.
他們很跟的.
上帝叫約書亞.
約書亞就叫以色列人.
你怎麼做呢?.
他們把石頭帶過去.
到他們所住宿的地方.
就放在那裡.
約書亞另把十二塊石頭.
立在約旦河中.
在台約會的祭司腳站立的地方.
直到今天.
那石頭還在那裡?.

$^{321}$很奇怪的.
究竟是立了一堆石頭.
還是立了兩堆石頭呢?.
上帝吩咐約書亞要立石.
約書亞就吩咐以色列人去立石.
以色列人就按吩咐而行.
福圖畫是很漂亮的.
他們經歷了神的拯救.
最重要就是按神的吩咐而行.
而不是自把自為.
但是第九節是很難解的.
在這一段.
最奇怪的那一節.
約書亞在約旦河的河床.
又另外立一堆石頭.
就是一堆在岸上.
一堆在河床裡面.
但是上帝好像沒有這麼吩咐.
他多立一堆在河裡面.
傳統的解法.
絕大部分的色經學者的解釋.
都是立了兩堆石頭.
一堆在岸上.
一堆在河床.
有幾個問題.
第一.
為什麼要另外立一堆石頭呢?.
起碼我暫時還沒有看到一個.
很令我很信服.
很合理的解釋.
要另外立一堆石頭.
在河床裡面.
第二.
你記得當百姓過河之後.
河水在30公里以外的亞當城.
湧過來.
湧過來的時候.
剛才說的石頭不是很大.
不是很大的石頭.
應該怎樣呢?.

$^{361}$沖散了.
30公里以外的河水.
沖下來這12塊石.
沖散了.
怎麼會直到今天呢?.
你記得經文所說.
直到今天.
就是約書亞的成書的時候.
它還在.
當然有學者的解釋.
約旦河真的不是很深.
約旦河的河床不是很深.
當水退或者乾旱的時候.
可能會看到一些石頭露出水面.
讓以色列人看到.
但是你沖下來.
正常都會全部沖散了.
我自己最不明白的地方就是.
神沒有吩咐約書亞.
另外立12塊石頭在河床.
這段經文你會發覺.
整個的信息鋪排.
很明顯的就是.
聽命.
神吩咐約書亞.
約書亞吩咐以色列人.
以色列人按吩咐而行.
但是.
在這裡.
約書亞好像沒有什麼聽命.
你看看第十節.
他說台約櫃的祭司.
站在約旦河中.
等耶和華曉遇約書亞吩咐百姓的事辦完了.
聽命.
這個信息不斷重複.
是照摩西所吩咐約書亞的一切話.
於是百姓急速過去了.
整個圖畫都是.
上帝怎麼說.

$^{401}$他怎麼做.
摩西以前怎麼說.
他就照做.
是按上帝的指示.
為什麼約書亞自作主張.
另外立一堆石頭呢?.
所以我想提供另外一個解釋.
這個解釋.
可能已經過了.
這個解釋.
其實我的理解.
約書亞沒有另外立一堆石頭.
在第九節的原文.
其實沒有"令"字.
另外這個字在原文沒有.
所以.
約書亞在祭司站立在河床的地方.
他選擇了十二塊石頭.
或者立了十二塊石頭.
然後就叫那十二個人.
崗台到岸上.
放在他們住宿的地方.
在第八節.
你留意八節最後有"那"類字.
就對應第九節最後"那"類字.
就是紅色的highlight.
這兩個"那"類字是指同一個地方.
你看下去就知道是吉甲.
所以如果你這麼看.
原來.
到了約書亞祭成書的時候.
人們還看到那堆石頭.
就是在吉甲.
在陸地上那堆石頭.
約書亞沒有在河床另外立一堆.
我不知道大家是否明白我所說.
他沒有另外立一堆.
他選擇了一堆石頭.
叫那十二個人搬到那裡.
那裡就是吉甲.

$^{441}$無論你怎麼去理解也好.
是傳統的解釋.
或者是我剛才提供的另外一個解釋.
其實也是給大家參考的.
我個人就傾向是另外一個解釋.
因為這個解釋就會很吻合整段的主旨.
聽命.
聽上帝說話.
不過怎樣也好.
無論一堆也好.
兩堆也好.
重點就是紀念上帝的拯救.
這就是立石的一個很重要的意義.
神拯救以色列人.
出埃及.
入迦南.
入迦南的最後障礙.
就是約旦河.
他們已經經過四十年的流浪.
到了約旦河東.
神為以色列人掃除了最後一個障礙.
讓他們過去得地.
所以他們要立石去紀念.
我不知道在你們的人生裡面.
有沒有經歷神為你們掃除障礙呢?.
有沒有經歷神為你們掃除困難呢?.
甚至拯救我們呢?.
好像我小時候所經歷一樣.
多數都有.
我們經歷過之後.
我們有沒有立我們的石呢?.
我們還是先過海就是神.
神救了我就OK了.
在經濟困難.
神給了我.
供應了我.
身體有病.
神醫好了我.
我們有沒有立我們的石.
真的去紀念呢?.

$^{481}$我記得.
說回我自己的一個經歷.
我自己讀神學的時候.
就是很多年前.
在我進入神學院之前.
我長子就出生了.
也很艱難的.
他在幾個月前出生.
然後我就讀書.
然後我在第一年年尾的時候.
我的外父.
就是曾真牧師的爸爸.
癌症就過身.
也是一個很大的衝擊.
隔了不久.
隔了兩個月左右.
我小兒子就出生了.
就是剛才做手術的那位.
那時候是我二年級的.
應該是剛開學.
他就出生了.
一個很悲哀的心情.
迎接一個新生命.
但是接著到了我三年級的時候.
就到了我媽媽過世.
這個對我有很大的打擊.
我們母子關係很好.
但是也要讀書的.
繼續讀.
其實我一直讀.
我媽媽在我三年級上學期就過世.
其實隔了幾個月.
我太太跟我說.
我好像很低落.
情緒很低落.
好像有一點抑鬱的感覺.
我媽媽過世對我有一點衝擊.
其實我在當時.
我是想過退學.
我想過不讀.

$^{521}$如果當時我不讀.
就應該今天不會站在這裡了.
弟兄姊妹.
這是我人生很大的經歷.
有四塊石頭.
當然有兩個就是我兒子.
有兩個就是我離世的親人.
四塊石頭.
這四塊石頭是做什麼的呢?.
紀念.
紀念上帝帶領我.
能夠完成那三年的神學課程.
對我來說.
能夠在神學院畢業.
其實是神蹟.
對我來說.
要帶著兩個很小的小朋友.
有兩個親人去離世.
弟兄姊妹.
不要讓你生命的這些經歷.
白白過去.
立石為記.
紀念上帝的拯救.
我們通常經歷了這些.
我們很多時候的反應是怎樣的呢?.
多數都是感恩的.
但是如果你留意.
你又聽以石為記的系列.
應該會說到.
譬如大衛對哥利亞.
應該會說到耶穌在曠野受試探.
在《養人荒》第八章.
有沒有扔行人的婦人.
你會發覺感恩的意味很少很少.
你看看這些經文.
這些石頭所代表的多數都不是說感恩.
感恩是起點.
弟兄姊妹.
我們千萬不要將上帝在我們生命中.
讓我們經歷這些事情.

$^{561}$是起點.
究竟上帝對你的拯救.
去掃除了你人生的這些障礙.
你是否更加敬畏祂呢?.
紀念祂對你的救贖.
不是單單多謝這麼簡單.
弟兄姊妹.
在我展覽會的過程中.
我發覺很多弟兄姊妹經歷了這些.
只有感恩.
千萬不要只有感恩.
感恩是起點.
在今天這段經文.
他們是紀念上帝對他們的拯救.
最後他們是敬畏上帝.
我們接著看下一個段落.
11-14.
在第11-13節.
他說:眾百姓盡都過了河.
過了中段經文之後.
開始故事的延續.
耶和華的約桂和祭司就在百姓面前過去.
裡面人.
迦德人.
馬來西亞半支派的人.
都照摩西的吩咐.
帶著兵器在以色列人前頭過去.
約有四萬人都準備打仗.
在耶和華面前過去到耶利哥的平原.
等候上陣.
摩西曾經跟裡面迦德.
馬來西亞半支派協議.
你們可以在約旦河東德地.
如果大家熟悉背景.
但是他們可以打仗的南丁.
就要過河幫九個半支派去攻取迦南地.
這兩個半支派就順手承諾.
有四萬人帶著兵器去到耶利哥的平原.
你知道一進入約旦河.
他們就在吉格安營.

$^{601}$夾在耶利哥的東邊.
就在耶利哥平原準備去攻取第一座城.
耶利哥.
當然大家知道結果就不用打仗.
圍著十三個圈.
城牆就塌了.
十二個人.
去台石.
是代表全體的以色列人.
都是經歷這個請求.
十二個支派.
現在全部都參與攻取迦南.
其實想說什麼呢?.
想強調他們是合一.
這群人是合一上帝的子民.
就像旺鎮一樣.
是合一.
是一個群體.
不是你有你做我有我做.
我們是有關係.
這是一個很重要對我們的提醒.
我們接著看第十四節.
當那日.
耶和華使約書亞在以色列眾人眼前專大.
在他平生的日子.
百姓竟為他將從前.
竟為摩西一樣.
這一節如果你熟悉上一章.
會令你聯想第三章第七節.
耶和華對約書亞說.
從今日起.
我必使你在以色列眾人眼前專大.
使他們知道我怎樣與摩西同在.
也必照樣與你同在.
這就是第三章第七節.
與第九章十四節.
其實是很相似的.
你看看.
我在紅色的高光那裡.
他令約書亞被尊為大.

$^{641}$這兩節的重點也是.
但是目的是什麼呢?.
就是要以色列人知道神與約書亞同在.
這一群人你聽約書亞的話.
就即是聽神話.
所以這裡有一個很重要的訊息.
就是聽話.
我們合一不是說我們關係很好.
在皇城彼此相愛.
我們合一做什麼呢?.
就是聽神話.
如果我們合一而不聽神話.
那就不如不要合一.
弟兄姊妹.
紀念上帝的拯救不是很純感性地紀念.
上帝很好.
對我真的很好.
我一禱告上上都去英文.
我們紀念上帝的拯救.
是要合一地聽命於上帝.
可以這麼說.
合一地跟從神.
你就會經歷上帝的眷顧.
我們經歷了上帝的眷顧之後.
我們的回應是什麼呢?.
更加合一地跟從上帝.
這個可以給我們整個群體一個檢視.
我們是不是很合一地跟從上帝呢?.
上帝的比旺朕有很多艱難.
我今天一回來.
看到天花有些問題.
都是一個困難.
很多的.
我們是不是很合一地跟從上帝呢?.
其實在整個第四章.
你應該很容易會發現.
有些字眼是重複重複再重複的.
祭司.
約衛.
和耶和華.

$^{681}$這些字眼是很頻密地出現.
祭司出現了八次.
約衛.
其實是一個季節.
不是一個約.
是一個季節.
季節再加上一個約.
有時候加上一個法.
或者法季.
季節出現了七次.
耶和華就是最頻密.
十四次.
所以.
很簡單.
二十四節的經文.
你會發覺這些字非常頻密地出現.
以色列人是跟著.
台,耶和華,約衛的祭司.
就可以了.
而耶和華的約衛就代表.
上帝現在臨在.
就是跟著神.
就可以了.
在今時今日.
香港地區.
旺鎮也好.
或者香港教會也好.
前景是怎樣的呢?.
我不知道大家心裡面有沒有這些思想.
究竟在香港前景是怎樣的呢?.
簡單來說.
合一.
聽命於神.
跟著神走.
就可以了.
前景可以很樂觀.
信徒也一樣.
你知道神的召命.
聽命而行就可以了.
其實很簡單.

$^{721}$我們的信仰不要搞得太複雜.
知道神想我怎樣.
憑信心走神想我怎樣走的路.
這樣就可以了.
不過我想問大家.
聽命跟從上帝難不難?.
有沒有人覺得很難?.
跟從上帝.
有沒有感恩?.
在這裡沒有一個覺得難的.
但我覺得很難.
有時候我覺得不容易.
我聽命跟從上帝.
我舉一個例子.
如果上帝在我年輕的時候.
呼召我獨身.
就是不結婚.
難不難?.
有時候對男人也不容易.
不結婚.
你個個像保羅,耶穌一樣.
不容易的.
一點也不容易.
當然了.
後來不幸碰到師母.
不是不幸.
就是碰到師母.
這樣就要結婚了.
一點也不容易.
有一個心理學家說過.
你每天做一兩件你應該做而不想做的事情.
你的生活會快樂一點.
你的壓力會少一點.
每天做一兩件你應該做而不想做的事情.
我自己有一個習慣.
我通常星期日就會在教會崇拜.
上午.
下午我多數都會跑步.
我有一個跑步的習慣.
這幾個月大家都知道.

$^{761}$天氣很熱.
很炎熱.
你想一想.
如果你是我.
坐在家的沙發.
省著冷氣.
喝一杯可樂.
看看昨晚的英超舒服一點.
還是在三十多度.
曝曬之下去跑步.
哪樣舒服一點呢?.
有時候我真的.
這樣的天氣.
居然曬得這麼厲害.
不如靠冷氣喝可樂比較好.
我應該去跑步.
但有時候我也覺得跑步也挺辛苦.
不過多數.
我自己多數.
我最後也是迫自己換件衣服.
先下去跑步.
幾乎每一次跑完步回來.
很舒暢.
吹一身汗.
很舒暢.
甚至有一點滿足感.
有一點成就感.
每天做一兩件你應該做.
而你很不想做的事情.
其實會讓你的生活快樂一點.
我將這句話應用在一個屬靈的層面.
每天做一兩件.
神想你做.
而你不太想做的事情.
你就會平安喜樂很多.
你試一試吧.
每天做一兩件.
你明知神想你做.
但你又不是太想做的事情.
你明知神想你要靈修.

$^{801}$硬是好像拉牛上樹一樣.
你明知神想你用這個鐘預備查聖經.
硬是好像不太行.
你試一試吧.
如果你可以.
你會平安喜樂舒暢很多.
為什麼呢?.
因為人生不是做我想做的事情.
人生就是做神想我做的事情.
人生很簡單.
永遠都是做神想我做的事情.
聽命而行.
這一段的重點.
合一.
聽命而行.
應該是一個最蒙福的人生.
我們再看最後的段落.
我們先看第十五到十八節.
「曉遇約書亞說.
你吩咐台法季的祭司從約旦河上來.
約書亞就吩咐祭司說.
你們從約旦河上來.
台耶和說約季的祭司從約旦河上來.
腳掌剛落旱地.
約旦河的水就流到源處.
營救將戈兩岸」.
你有沒有留意這個模式又重複出現?.
「神吩咐約書亞.
約書亞吩咐祭司.
祭司就按吩咐而行」.
吩咐而行.
這是跟你一到十節的模式.
又再一次出現.
神怎麼說.
他們就怎麼做.
「台著約季的祭司從約旦河上來」.
祭司從河水上來.
標誌什麼呢?.
全以色列的人已經成功踏入應許之地.
最後那幾個祭司都上了約旦河河西的陸地.

$^{841}$他們已經成功踏入了迦南地.
應許之地.
我想大家留意第十五節.
他用了「法季」這個字.
如果你熟悉這一章.
你看全章都是用「約季」.
是一個「季」加一個「約」字.
這個就是「季」加一個「法」字.
「法季」在約書亞記出現了一次.
就是這裡.
所以很奇怪的.
他很頻密地用了「約季」.
忽然轉向用了「法季」.
「法」這個字的基本意思就是證據或者見證.
神用「法季」這個字.
是用這個「季」來見證神對以色列人立約的忠誠.
這是見證.
神應許賜迦南地.
現在以色列人正式踏足迦南地.
「法季」是見證了神的救贖.
一個爸爸想幫兒子.
但如果這個爸爸有能力.
但他沒有愛心.
他就不會幫兒子.
但如果這個爸爸有愛心而沒有能力.
他又幫不到兒子.
愛心和能力要並存.
他才能幫兒子.
上帝是一樣的.
上帝是全能.
也有完全的愛.
所以他救贖了以色列人.
他救贖了我們.
他幫了我們.
以色列的歷史.
或者我們自己的生命經歷.
見證了這位滿有愛心.
滿有能力的神的救贖.
弟兄姊妹我們要很相信.
很相信.

$^{881}$上帝很愛我們.
上帝很有能力.
如果你這個相信缺一不可.
你相信上帝很有愛心而不夠能力.
他真的幫不了你.
你相信上帝很有能力而沒有愛心.
他不會幫你.
弟兄姊妹在任何的逆境裡.
我自己永遠都經常提醒.
就是神掌權.
神愛我.
在整個疫情裡的這兩三年.
神掌權.
神愛我.
當然不代表我不會中招.
但是神掌權.
神愛我.
在社運也好.
在戰亂也好.
整個世界很動盪也好.
神掌權.
神愛你.
你要堅持這個信念.
你一定會經歷到上帝的救贖.
我們看最後的一個段落.
就是講垃圾的意義.
在第十九節.
正月初十日百姓從約旦河上來.
就在吉甲在耶利哥的東邊安營.
正月就是猶太曆法的彌散月.
我們叫做.
就是在今天大約三月中到四月中左右.
正月初十.
如果你熟悉初一及二.
這個有意義.
初一及二的時候.
在哪一天是選擇逾越節的羔羊.
就是正月初十.
就是選擇逾越節的羔羊.
作者很明顯是想將過紅海出埃及.

$^{921}$和過約旦河進入迦南.
連結在一起.
這就是神應許的救贖.
接著他們從約旦河中取那十二塊石.
約書就立在吉甲.
對以色列人說.
日後你們的子孫問他們的父親說.
這些石頭是什麼意思?.
你們就告訴他們說.
以色列人曾走過乾旱.
乾地過著約旦河.
因為耶和華你們的神.
在你們前面使約旦河的水乾了.
等著你們過來.
就如耶和華你們的神.
從前在我們前面使紅海乾了.
等著我們過來一樣.
很明顯.
作者將過約旦河和過紅海結連.
出埃及和進入迦南結連.
立石的意義.
好像剛才第七節所說.
現在他再重複一次.
就是紀念上帝的救贖.
去到最後那一節.
第二十四節.
我很想大家一起讀好嗎?.
這一節是這一段最重要的節.
好,預備起.
(主持人說).
過約旦河不是純粹一個神蹟.
以色列人知道.
神很厲害.
過約旦河.
他們要立石為記在二十四節.
告訴我們這兩個目的.
第一,使地上萬民.
包括外邦人.
包括今天在旺朕以外的人.
他知道神救贖的大能.

$^{961}$第二,使你們.
就是以色列人.
信仰群體.
就是旺朕.
我們要永遠的去敬畏神.
這兩個很重要的向導.
立石是要信仰群體內.
和信仰群體外的人.
都知道神的大能.
以致我們可以敬畏祂.
弟兄姊妹,在你的生命裡面.
我不知道有什麼是很值得你紀念的.
可能是生日.
結婚紀念.
你得獎.
晉升等等.
這些很值得紀念.
這些不是錯的.
千萬不要.
這些不是錯的.
但我們要問.
這些能不能夠讓人.
那些未信的人.
知道神的大能呢?.
如果那些多數都不能夠.
讓未信的人知道神的大能.
是不是那麼值得紀念呢?.
能不能夠讓信的人更加敬畏神呢?.
永遠敬畏神.
我們紀念的那些.
如果做不到這兩個目的.
那些不是錯.
不過可以放輕一點.
反而我們要想一些.
什麼紀念能夠讓我們去敬畏神.
讓人知道上帝的大能.
最後就應用.
在疫情這兩三年.
我想大家想一想.
有什麼是值得你紀念的?.

$^{1001}$已經過了差不多三年了.
有沒有一些經歷值得你紀念的?.
你那十二塊石頭在哪裡呢?.
但更加重要的是.
紀念這些東西是為了什麼呢?.
為了自己?.
為了自己開心一點?.
雖然不是錯.
是不是那麼值得呢?.
我們紀念這些東西.
是不是為了讓人更加知道神的大能.
讓人更加去敬畏神呢?.
我們同心禱告.
親愛的主人.
在屬於你的子民.
以色列的群體裡面.
你施行大能的拯救.
其實我們也經歷了.
我們真以色列人.
我們經歷你的救贖.
我們不僅經歷你的救贖.
在我們人生裡面有很多的情況.
你真的因手能在.
讓我們更加去感激你的大能.
更加去敬畏你.
讓更加多未信的人去敬畏你.
主我們求你去施恩.
祝福旺振.
讓每一個人我們所紀念的.
是你自己大能的拯救.
讓我們敬畏你.
同心禱告.
奉耶穌基督的聖名.
阿們.
\newpage



\section{箴言 30:1-14}
\label{sec:8EXqdqFDCIc}
\textbf{2024年3月17日崇拜直播｜陳明泉牧師｜十個深化關係的禱告-智者的禱告｜箴言30\_1-14}
\newline
\newline
連結: \href{https://youtube.com/watch?v=8EXqdqFDCIc}{\texttt{https://youtube.com/watch?v=8EXqdqFDCIc}} ~~~~ 語音日期: 2024-03-18
\newline
\newline
\hyperref[sec:1GcKc4lpPUQ]{\small{< < < PREV SERMON < < <}}
~
\hyperref[sec:index_chronic]{\small{[返順時目]}}
~
\hyperref[sec:3GKCAsbLUbA]{\small{> > > NEXT SERMON > > >}}
\newline
\newline
箴言 30:1-14
\newline
\begin{longtable}{cl}
\hline
\hline
章節 & 經文 (和合本修訂版)\\
\hline
30:1 & \begin{tabularx}{0.7\textwidth}{X} 雅基的兒子、瑪撒人亞古珥的言語，是這人對以鐵和烏甲說的。 \end{tabularx} \\ \\ \relax
30:2 & \begin{tabularx}{0.7\textwidth}{X} 我比眾人更像畜牲，也沒有人的聰明。 \end{tabularx} \\ \\ \relax
30:3 & \begin{tabularx}{0.7\textwidth}{X} 我沒有學好智慧，也不認識至聖者。 \end{tabularx} \\ \\ \relax
30:4 & \begin{tabularx}{0.7\textwidth}{X} 誰升天又降下來？誰聚風在手掌中？誰包水在衣服裡？誰立定地的四極？他名叫甚麼？他兒子名叫甚麼？你知道嗎？ \end{tabularx} \\ \\ \relax
30:5 & \begin{tabularx}{0.7\textwidth}{X} 神的言語句句都是煉淨的，投靠他的，他便作他們的盾牌。 \end{tabularx} \\ \\ \relax
30:6 & \begin{tabularx}{0.7\textwidth}{X} 你不可加添他的言語，恐怕他責備你，你就顯為說謊的。 \end{tabularx} \\ \\ \relax
30:7 & \begin{tabularx}{0.7\textwidth}{X} 我求你兩件事，在我未死之先，不要拒絕我： \end{tabularx} \\ \\ \relax
30:8 & \begin{tabularx}{0.7\textwidth}{X} 求你使虛假和謊言遠離我，使我不貧窮也不富足，賜給我需用的飲食。 \end{tabularx} \\ \\ \relax
30:9 & \begin{tabularx}{0.7\textwidth}{X} 免得我飽足了，就不認你，說：「耶和華是誰呢？」又恐怕我貧窮就偷竊，以致褻瀆我神的名。 \end{tabularx} \\ \\ \relax
30:10 & \begin{tabularx}{0.7\textwidth}{X} 不要向主人讒害他的僕人，恐怕他詛咒你，你便算為有罪。 \end{tabularx} \\ \\ \relax
30:11 & \begin{tabularx}{0.7\textwidth}{X} 有一類人，詛咒父親，不給母親祝福。 \end{tabularx} \\ \\ \relax
30:12 & \begin{tabularx}{0.7\textwidth}{X} 有一類人，自以為純潔，卻沒有洗淨自己的污穢。 \end{tabularx} \\ \\ \relax
30:13 & \begin{tabularx}{0.7\textwidth}{X} 有一類人，眼目何其高傲，眼皮也是高舉。 \end{tabularx} \\ \\ \relax
30:14 & \begin{tabularx}{0.7\textwidth}{X} 有一類人，牙如劍，齒如刀，要吞滅地上的困苦人和世間的貧窮人。 \end{tabularx} \\ \\
[1ex]
\hline
\hline
\end{longtable}
$^{1}$今日經文是《箴言》三十章一到十四節.
經文都印在你們秩序表的內頁.
請聽我讀三十章一到十四節.
(朗讀).
經文讀筆.
(朗讀).
各位姐妹平安.
我們今日用《箴言》思想一下.
在《箴言》裡面的智者.
他是怎樣來祈禱的.
在我們進入經文之前.
我曾經遇過一個老牧者.
講開他的教會裝修的一個小片段.
我很值得和大家分享.
這個老牧者和我說.
教會有一段時間要擴堂.
要重新買一個新的地方.
將要入伙之際.
弟兄姊妹就組織祈禱會.
一齊來到衛教會.
將要搬到一個新的地方來祈禱.
他說其中一個很重要的土文.
或者弟兄姊妹的心願是祈求甚麼呢.
因為你去到一個新地方.
有很多的規例.
有防火又好.
或者有不同的政府上面的要求.
一班的弟兄姊妹.
他們就在祈求上帝神掩眼.
掩著那些官員的眼.
希望他們有一些不太合規格.
或者有些地方做得不盡善的.
他們已經想盡了.
但是可能還有漏了的.
忘記了的.
或者在漏洞裡面.
上帝就掩著那些官員的眼.
等於快快脆脆簽了紙.
讓他們可以快點入伙.
去到一個新的堂子裡面.

$^{41}$來敬拜上帝.
這個心其實是好的.
希望有一個新的地方.
但是這個老牧者跟我說.
如果我這樣教弟兄姊妹祈禱.
這些弟兄姊妹到最後.
她會經歷那種苦過.
很奇怪.
每個人的祈禱都希望快快脆脆.
一切順順利利.
這個老牧者說.
我跟他們的祈禱是相反的.
希望在未入伙之前.
所有問題要浮現出來.
他說我們做基督徒.
或者我們的祈禱.
如果我們的眼光.
只是解決當下眼前的問題.
其實我們沒有辦法.
傳遞智慧給弟兄姊妹.
我們要教會弟兄姊妹在祈禱的時候.
預示了人生.
或者將要發生什麼事.
在問題出現之前就杜絕了她.
那些官員不是跟你們這個入堂.
這個地方作對.
反而這班人是幫我們.
預示或者防止所有問題的出現.
接著這個牧者就說.
他力排眾議.
開頭的時候說了一個相反的祈禱.
令弟兄姊妹感覺上不思美意.
不過當他再去這樣的教導的時候.
這班弟兄姊妹明白.
人生不可以只是為眼前的問題.
搞定了.
就以為所有事情就迎刃而解.
我相信這個老牧者.
他在說的是一個智慧的人生.
我們的人生怎樣的祈禱.

$^{81}$不是純粹為眼前的需要.
或者我們眼前當下.
我們自己的生活.
智者和一般人的看法不同.
就是智者永遠看事物.
是一條直線.
有個方向 有個向度.
如果我們當下這個處境.
繼續發展下去的時候.
他能夠預示將來會有什麼情況出現.
在我們未出現問題的時候.
我們就能夠懂得去改喚心思.
避免在我們生命當中出現很多難處.
和很多我們不想經歷的問題.
今天如果用這個角度去看.
其實阿古爾的禱告.
都是一個這樣的禱告.
在大家看到的《箴言》裡面.
其實《箴言》裡面記載著.
討民或者祈禱的內容不是很多.
不過今天我們抽取這一段.
特別說的是智者阿古爾.
他怎樣教他兩個學生來祈禱.
不過這個經文不單止是說.
那個祈禱要怎樣祈禱.
或者要怎樣想.
整個第30篇.
在《箴言》裡面.
其實那個涵蓋面很闊.
由第一節到第33節.
其實是兩個大段落.
第一個段落其實是一至十四節.
也是我們今天所讀的.
而後半段由15節到第31節.
就是講到這個世界的荒謬.
接著到第32到33節.
再一次的來強化他的教導.
今天我們沒辦法看完33節.
不過讓大家知道有個框架.
回家再看看.

$^{121}$特別有趣的是.
在第15節到第31節.
講到一個數字箴言.
三四這樣寫出來.
但是他寫這個數字格式的狀況.
目的不是純粹玩文字.
其實他講到一個荒謬的世界.
每個時代都有不同的荒謬.
今天我們有我們當代的荒謬.
在箴言的世界.
阿古爾的世界也有一些荒謬.
你回去看看.
特別留意一下23節也很有意思.
在講醜惡的女子出嫁.
被女接續主母.
他看到的荒謬.
就是一個這樣的荒謬的世界.
一連串.
不過是一個荒謬的世界.
是一個他插不透.
在文青裡面.
一個智者如何教導他的學生祈禱呢.
今天我們就用這個角度.
來看1-14節.
我們先祈禱求上帝幫助我們.
天父上帝願你在我們當中帶領我們.
讓我們能夠明白你的話語.
願你的聖靈與我們眾弟兄姊妹同在.
讓每一位聆聽的我們聽得明白.
讓孫講的講得清楚.
忠心的將你話語來講解.
讓弟兄姊妹明白.
我們恭敬將我們的時間交託給你.
願你賜福.
獻上我們的禱告.
奉求基督耶穌得勝的聖名.
Amen.
好,我們一起看吧.
在《真言》第30篇裡面.
第一節裡面其實是一個挺有爭議性的一句.

$^{161}$雅基的兒子雅古爾的言語就是真言.
這人對耳鐵和烏鴿說.
這一句而已.
我們和學本讀得很順.
不過在原文裡面.
帶出來的翻譯其實有很多面向的.
有很多可能性的.
非常多.
不同的色經學家有不同的講法.
首先讓大家知道.
雅基是可以是一個人的名字.
雅基亦都是可以是一個美德的表現.
如果你明白舊約聖經裡面.
希伯來人講的那些名字.
名字背後可能是代表一個人.
或者一個名字背後其實是代表著一種美德.
有時候形容一個人.
他會用某一種人.
他會用某一種名來形容.
例如那個人不好.
他會說邪惡之子.
他用的名字.
你的名字叫邪惡.
有時候就是一堆人是這樣的.
和學本採取的翻譯.
就是將雅基成為一個人名.
雅基的兒子雅古爾.
雅古爾這個智者的爸爸就是雅基.
另一方面還有一點點地方值得大家留意.
在和學本裡面沒有特別注明.
在一些譯本裡面.
我在那個part point大家都會看到.
就是瑪沙地雅基的兒子.
瑪沙或者瑪薩一個譯音.
瑪沙或者瑪薩這個字眼.
可以解釋為一個地方.
這個地方其實是在講亞伯拉罕.
另外一個兒子以實瑪利的後代.
他們聚居的那個地方.
有一些說法就是.

$^{201}$瑪薩就是亞伯拉罕的另一個兒子.
以實瑪利的後代.
他們都有講智慧.
他們和亞伯拉罕的那個源遠流長的智慧不同.
不過在《箴言》裡面都編輯了他們一些很重要的講論.
所以就加了這個瑪薩或者瑪沙這個字.
但是和學本沒有譯的.
另一方面瑪薩也可以作為一個動詞.
就是在講負重擔或者上帝要講說話.
啟示神諭這個意思.
如果是這樣就不是純粹在講一個地方.
而是在講上帝藉著瑪基的兒子阿古爾講的這番說話.
我想我們現在下來的.
我們都是用一個這樣的角度.
未必一定要用外邦的.
以實瑪利的後代來講的.
因為這個瑪沙和後面下來所講的內容都是很一致的.
就是在講上帝要藉著他來說話.
為什麼這樣講呢.
接著後半節.
這人對以鐵和烏甲說.
這人對...說.
這個字呢.
這個格式.
其實在舊約聖經裡面很多時候就是在講先知說話.
他就會用這個句式和這個字的那音.
來講上帝有些很重要的內容.
要對著一眾的百姓子民.
或者對著人來宣講上帝的訊息.
如果你用一個這樣的角度來看.
其實如果大家同意.
瑪沙不是純粹講地.
是在講上帝要說話.
而上帝以阿古爾作為先知的角度.
來對著他的學生.
以鐵和烏甲來講以下的一番說話.
如果用一個這樣的層次來看.
《箴言》第30章裡面的教導內容.
就不是純粹講生活日常.
爸爸教兒子怎樣說.

$^{241}$這個層次就會被提升.
去到一個人生必須要學.
甚至每一個人都要虛心聆聽.
因為這是上帝對我們所說的那番說話.
我們用這個角度來研讀.
這樣下來的經文.
在他下來講的時候.
最特別的地方.
阿古爾講這番說話的引言.
他繼續用甚麼說話來講呢.
由第二節到第五節.
他就用了一連串的反文.
或者某程度上是一種反諷的反文.
來講他自己處於甚麼狀態.
第二節.
我比眾人更蠢笨.
也沒有人的聰明.
我沒有學好智慧.
也不認識自聖者.
誰升天又降下來.
誰罪風在掌握中.
誰包水在衣服裡.
誰立定地的四極.
他名叫甚麼.
他兒子叫甚麼.
你知道嗎.
神的言語句句都是練成的.
投靠他的.
他便作他們的盾牌.
你會見到在這裡.
一系列一連串的大量的問句.
是對自己作一個否定.
通常我們講話.
我們要人去聽.
或者告訴人.
這句話是上帝講的.
我們要用甚麼來印證呢.
我們會說我們有一些神跡歧視.
我很認識上帝.
我摸清上帝的心腸.

$^{281}$所以我現在要將這番說話告訴你.
這是人之常情的表達方法.
用上帝來增加自己的權威性.
不過阿古爾在這裡.
他作為教兩個徒弟.
以鐵和烏鴿.
他不是用這個角度.
來敦起一個很權威的款.
來講一個很重要的信息.
反而他將自己的位置踩到最低.
他在講我沒有人這麼聰明.
我比所有人都蠢.
沒有學好智慧.
又不懂得自勝上帝是誰.
誰升天又降下來.
誰罪風在掌握中.
誰包水在衣服.
這裡全部都是在講創造.
誰立定地的四極.
他的名字叫甚麼名字.
他的兒子是誰.
你知道嗎.
他在講他自己.
用一個好像無知.
或者甚麼對上帝都不認識的角度.
來表達這個一連串的提問.
或者一連串的講論.
不過你要明白.
很多時候.
真言所講的內容.
有些叫做反話.
他在講他不明白.
但是他描述的字眼.
在當代很多人來說.
都未必知道.
在講那個自勝者是誰.
特別注意第四節.
升天又降下來.
今天我們基督徒都未必明白.
不過他用的字眼就代表這件事.

$^{321}$道成肉身.
升天降下來.
在舊約的真言裡.
他就用道成肉身這個詞彙.
或者這個講法.
來講我不知道.
誰去創造.
誰是創造.
誰是聚攏的水.
誰設定四極.
他的名字叫甚麼.
他的兒子叫甚麼.
可能有些解經家以為是以色列.
但是如果我們基督徒看.
有些講法.
他的兒子是耶穌基督.
甚至推到極致.
可能這個阿古爾.
知道的東西比一般普羅大眾更加多.
不過他採取的就是一種.
訴諸於自己無知的一種狀態.
他要強調.
他所講的這番說話.
並不是來自他自己的經驗.
或者智慧有多高.
掌握了天地萬物主的心腸.
所以用一個居高臨下的姿態.
將下來的話語.
將他要教導或者勉勵的東西.
來告訴他的學生知道.
反而他用一個.
自我否定的角度.
他甚麼都不知道.
訴諸於自己是無知的一種.
一個角度.
來將上帝的權柄.
來表達出來.
所以第五節.
他一無所知.
他自己甚麼都不知道.

$^{361}$但總比我們知道的多很多.
但他所講的最終的權柄.
是放在上帝那裡.
上帝的言語每句都是唸剩的.
投靠他的.
他便作他們的盾牌.
然後他再勉勵他的學生說.
他的言語你不可以加添.
恐怕他責備你.
就顯為虛荒.
在這個第一個段落裡.
第一個值得我們去思考.
今天我們覺得智者是甚麼呢.
或者我們覺得.
用智慧角度來說.
我們要說的內容.
我們很多時候會為自己.
我有很多的經驗.
我讀了很多書.
或者我人生裡面.
有很深刻的體會.
這些不是錯的.
不過如果我們人生裡面只有這一面.
大概我們的人生.
我們就會被這些東西所規限.
我們對上帝的理解.
我們往往都是受著自己的經驗.
或者受著自己.
所知道的讀過的書.
來規限我們對上帝的認識.
上帝的超越.
和我們書本上.
或者我們經驗裡面所理解的上帝.
是一定有一個很大的距離.
在這裡面.
智者亞古爾說到.
他所認知的.
遠遠比真實的上帝.
有一個很大的距離.
他自己知道.

$^{401}$他窮他一生.
他都是一個無知者.
在上帝的面前.
成為一個無知者.
成為一個不知的人.
其實生命可以成為一個很謙卑的人.
求主今日幫助我們.
今日我們很多時候.
我們很喜歡.
或者我自己也是這樣.
我會用我自己的經驗.
用我讀過的書.
或者我涉獵過的一些東西.
我覺得這是世界的全部.
甚至我會站高一線.
很輕易就會給一些建議.
或者從一個全知者的角度.
來給不同的建議.
甚至有些頂子會說.
他有些困難.
他跟你去試的時候.
我們很容易會站高一線.
我們會覺得我們掌握了世界所有的問題.
我們深暗世界所有運作的道理.
我們用超然的角度.
來給一些建議給當時的人.
不過可能我們正正是.
有自己的聰明.
就成為了我們認識神的難處.
古爾他用一個無知者的角度.
之前有一個Apple的Steve Jobs說.
Stay foolish.
讓自己永遠保持愚蠢.
其實保持永遠都是最單純尋求的心.
用這種態度來找上帝.
我相信這會令到我們的經驗.
我們的背景不會破壞了我們.
以致我們不會成為一個武斷或驕傲的人.
很多人說或者聖經裡面.
新約裡面也說.

$^{441}$知識叫人自高自大.
所以我們不要知識了.
中國華人社會很多時候.
就會有一種反智的傾向.
不過這裡不是說一個這樣的狀態.
既然知識叫人自高自大.
我們成為一個沒有知識的人.
反而是知識叫我們謙卑.
我們掌握更多的知識.
告訴我們我們人生其實都是很無知.
我們甚至如何解釋是誰.
我們都沒有辦法去知道.
我們進入了一個題目.
但我們永遠都是窮一生.
來尋求那個答案.
尋求上帝 尋求智慧.
在座裡面可能有一些老師.
我不知道你同不同意.
如果學生上課的時候問他們.
有沒有明白 有沒有問題.
如果學生通常沉默 沒問題.
你會不會當他們全都懂了.
不會的.
越不問問題.
覺得沒問題的那些.
其實知道最大的問題.
是不是.
我小時候讀數學也是.
老師在黑板上寫出來的那些數.
我明白的 我看他寫出來.
是知道他在做什麼的.
老師問有沒有問題.
沒問題 你寫出來寫得很好.
這些字很漂亮 是不是.
但再自己做就不知道自己怎麼做了.
同樣道理.
人生裡面我們覺得有沒有問題.
經常都說沒問題.
行了 明白.
這其實是最大的問題.

$^{481}$反而你充滿了問題的時候.
你開始入題了.
你開始進入在那種思考的過程當中.
這個第一個段落.
阿古爾要顛覆一下.
整個智慧文學裡面.
說的是一個智商而下.
或者一個人掌握很多經驗.
教他的徒弟.
在這裡他教他的兩個徒弟.
第一件事他用一種最謙卑的角度.
來講上帝的話語不可以加添.
不要加鹽加醋.
接下來他就將那個態度成為了他的祈禱.
由第七節讀到第九節.
就是在說阿古爾的禱告.
我求你兩件事.
在我未死之先.
求你不要不恃給我.
求你使虛假和謊言遠離我.
使我也不貧窮也不富足.
似我需用的飲食.
恐怕我飽足不認你.
說耶和華是誰呢.
又恐怕我貧窮就偷竊.
以致褻瀆我神的名.
在第七到第九節.
其實說的是兩個很重要的禱告的內容.
第一個就是說.
讓虛假和虛謊遠離他自己.
這裡有兩個主動和被動去理解.
虛假和謊言離開自己.
第一件事就是求上帝給他力量.
成為一個真誠的人.
他成為一個不說謊的人.
在生活當中.
你苦心自問我是否一個誠實的人.
第一個反應是我當然誠實.
我們每個人都是這樣.
不過你再用顯微鏡看看自己的內心.

$^{521}$其實你生活中遍換了很多.
不論是黑色或白色的謊言.
你覺得有些事是無傷大雅的.
但有些地方你覺得萬不驚心.
又或者有些情況之下你情不得已.
甚至乎你自我形象.
或者你恐怕你恐懼.
你沒什麼安全感.
所以你要說一些和事實有距離的內容.
阿古爾在這裡求上帝.
第一件事就是要幫他自己成為一個誠實的人.
誠實的第一步最重要是什麼呢.
我們要承認自己是虛假.
承認,檢視自己.
有什麼地方製造了一些.
不論是白色或為自己保障自己的說話.
有時候我們可以選擇沉默.
但有時候我們更多的是.
用我們的嘴巴去說一些對自己有利.
或者令自己處於沒那麼大能處的一些話語.
以致我去保障自己.
令自己合法化去說自己是一個不說謊的人.
不過在經文裡阿古爾很謹慎.
他求上帝幫助他.
讓他不要成為一個說謊言的人.
不過第二個更值得我們注意.
說的是虛假和謊言.
不只是自己那麼簡單.
而是身處的世界.
他求上帝讓他遠離了這樣的狀況當中.
其實我們生活在一個謊言的世界.
我不知道大家同不同意.
其實我們今天日常接觸的媒體.
或者我們所看的資訊.
其實大部分都是謊言.
不過我們選擇的是什麼呢.
為了不要麻煩我相信.
又或者有些東西.
他呈現了最美麗的一面.
他沒有把暗黑的一面表達出來.

$^{561}$所以我們寧願相信.
我們會令自己舒服一點.
令自己感覺上好過一點.
我們覺得即使是謊言.
我們也會一聲吞下.
昨晚我看新聞.
我都很震驚.
看到新聞報告裡說的一個現況.
他說在國內有些很流行的.
預製了的食物.
就是梅菜扣肉.
有沒有看到.
他真的給你真正的豬肉.
不過他給的豬肉.
不是五花腩最好的那塊豬肉.
而是他給的那塊肉.
他騙你的.
他說是最好的五花腩.
其實不是.
他是給那些在頸上.
和身體接駁著的肥肉.
那些肥肉和五花腩最大的分別.
就是頸上的淋巴結會結起來.
所以那裡有很多毒素在裡面.
他也是肥膩膩的.
大家吃下去會淡淡的肉.
不過那些肉.
如果你長期吃.
對健康是有影響的.
他告訴你那塊肉是五花腩.
但實際上那塊是豬頸的淋巴結的肉.
那些人說.
在食物工場裡的那個人問他.
這些肉你都會做的嗎.
他說是 所以我不會吃的.
你吃啊.
他的客人吃啊.
但他自己不會吃的.
我們今天生活裡滿佈這些謊言.
我們生活在一個謊言.

$^{601}$有時候我們閉上眼睛就算了.
是這樣的.
我們就閉上眼睛.
相信了他.
忽視了那個真相.
阿古義在他的禱告裡.
求上帝幫助他.
遠離這些虛謊.
遠離這些謊言.
不過事實是一個很困難的處境.
我又看到一個這樣的研究.
有一個很奇怪的研究.
就是社會心理學.
1996年做過一個這樣的研究.
在外國做的.
那個研究的名字叫做.
Lying in everyday life.
就是說每天生活當中說謊話.
他用社會心理學的角度來看.
為什麼人們在生活裡會說謊話呢.
他就訪問了一系列的人.
在裡面得出了幾個結論.
為什麼人們會說謊話呢.
第一件事.
他說說謊話.
他們心裡覺得會帶來益處.
你說一點點謊言.
你會為自己帶來一些方便.
又或者.
你將你真實的形象.
和你口裡描述出來的.
說謊話的形象表達.
你覺得別人心目中對你的印象會好一點.
你說謊話.
用一個有利益的角度.
自利的角度.
所以就會說謊話.
這個我們經常都會明白.
但是第二件事更加重要就是.
他覺得說謊話很多時候.

$^{641}$不說謊話.
其實對於社會上忠誠的人.
永遠都是受害者.
忠誠是成為了一個吃虧的同義詞.
今天我們很多時候都是這樣想的.
第二件事.
他說到說謊話.
其實本小利大.
付出不是很多.
沒什麼代價的.
說一句錯的話.
但是可能帶來的利益是很多的.
本小利大.
所以人們就是這樣做.
第三.
每個人都是這樣說的.
這個世界都是一個謊言的世界.
每個人都是這樣說.
他也是這樣說.
有什麼問題呢.
所以整個世界都是這樣.
他自己就是這樣.
今天我不知道我們怎樣看我們身處的這個世界.
如果我們說一個荒謬的世界.
我們願意活在一個虛荒的裡面.
其實我們既是在罵他.
但同時自己又繼續享受在一種虛荒當中.
對阿古義來說是一個撥論.
他求上帝幫助他.
他看重的不是這個社會風氣.
不是說這個世界每個人都是這樣生活.
是這樣說謊話.
他就繼續和這個世界看齊.
他看重的是上帝的華語.
他看重上帝.
即使他用一個最謙卑的角度.
他不知道上帝是誰.
但他知道上帝的華語.
上帝的權柄.
遠超過他當下社會的情景.

$^{681}$他求上帝幫助他.
遠離虛假和謊言.
讓他在人生當中努力成為一個真誠的人.
又讓他身處的社會成為一個誠實的社會.
是一個忠誠的社會.
阿古義第二個的祈禱.
這個更加很有意思.
和這個世界的想法是相反的.
他求上帝不貧窮也不富足.
在這裡他特別用更多的篇幅來說.
賜給他日用的飲食.
在新約裡面耶穌基督教門徒的祈禱都是這樣.
我們日用的飲食今天賜給我們.
在阿古義裡面他用日用的飲食.
來求生活裡面的不需要多.
他不需要大量.
剛剛好去求上帝幫助他.
因為如果他有很多的時候.
他會覺得自己很有能力.
他就會驕傲.
令到他心裡面就會忘記上帝.
所以在經文裡面說.
恐怕我飽足就不認你說.
耶和華是誰呢.
他心裡面看重上帝.
但是原來過多的物質生活.
太過富足的生活.
是會看自己變得很膨脹.
看自己的需要.
看自己的實力.
和看自己的財富.
變得成為生命裡面最重要的偶像.
卻忽視了上帝.
不認上帝.
我們今天這個社會.
新年我們恭喜人發財.
不會恭喜你剛剛好.
我們不會這樣恭祝人家的.
不過如果阿古爾在當代出現.
他會祝你財富剛剛好.

$^{721}$不多不少.
在一個剛剛好的狀態.
讓人的心靈處於一種均衡.
不會因為過多.
他的日用飲食.
他說剛剛好對我們來說.
是貧窮線.
就是每天剛剛好的飲食.
和我們今天說的.
儲足半年的彈藥是兩回事.
不過在這裡調教我們.
我們生命裡面.
可能我們的財富為我們製造了煩惱.
我們很想拿多一點.
但是阿古爾說.
不要讓他變得很飽足.
以致他不認上帝.
阿古爾很實際.
他同樣都求.
你不要令到我很窮.
不要窮到一個地步連吃都不夠吃.
因為如果貧窮的時候.
第九節後半節.
又恐怕我貧窮就偷竊.
以致褻瀆我神的名.
如果他派貧窮.
走投無路.
他就走去偷又去搶.
又是褻瀆上帝.
他的心裡面.
他看財富.
看他生活的物質.
他祈求的就是剛剛好.
不多不少.
夠他日常生活所需.
我想這個也是我們今天.
作為基督徒.
我們生命裡面一個重要的提醒.
今天很多人.
我之前都說過.

$^{761}$我們以為財富多一點.
我的收入多一點.
或者我擁有的東西多一點.
我會令到我自己的安全感更加大.
不過其實有研究發現.
原來財富去到某個位置.
反而你的安全感比起你沒有.
還更加低.
你越擁有越多.
成王成孔.
經常都怕自己失去一些東西.
如果你一直都不介意.
你沒有什麼.
即使上上落落對你沒什麼影響.
有個這樣的弟兄說.
我都不買股票.
所以恆指上落對我沒什麼影響.
就是一個這樣的狀況.
我沒錢買.
對我來說.
那個市道旺與弱.
我都是剛剛夠好.
生活我就可以踏實安穩.
擁有太多反而煩惱更加多.
好 時間關係我們去到最後一個段落.
11至14節.
有一眾人周作父親不給母親祝福.
有一眾人自以為清潔卻沒有洗去自己的污穢.
有一眾人眼目何其高傲.
眼皮也是高舉.
有一眾人牙如劍 齒如刀.
要吞沒地上的困苦人和世間的窮乏人.
在這裡用的字眼要定義一下.
有一眾人.
一眾人這個字其實很奇怪.
你完全不知道在做什麼.
一眾聚就有了.
其實原文用的字是世代.
有一代人.
但是如果你用一代人.

$^{801}$你就以為那段11至14節是四代.
不同世代的人.
有一些比較貼近現在理解的.
其實有一種人.
不是一種人 是一種人.
有一類人.
不是一個半個.
也有一個集體.
不過在這個世界裡.
一樣是養四種人.
這裡就說一個這樣的道理.
這四種人是什麼呢.
其實是同出一轍的.
全部那種荒謬.
那種生活的形態.
其實就正正是那些人不認上帝.
藐視上帝.
看自己為很高的那些生活的表現.
第一種人就做父母不給母親祝福.
說他們不會實踐孝道的精神.
在猶太人裡面強調的孝道精神.
在這些極度自傲的人當中.
在看自己很重要的人.
他覺得父母是理應要供養他.
是要令他幫助他.
他沒有想過他要供養父母.
祝福母親.
在阿古爾的理解裡面.
這個就是.
這種人是一個自高自傲的一個人.
這裡稍微有一點題外話.
每一個民族.
都有他自己的民族文化特性.
中國人和猶太人在孝道上是很相似的.
中國人所說的孝道精神.
其實是家庭維繫的單位.
父慈子孝.
無論家人是怎樣都好.
我們都盡力用孝道的方式.
來孝敬他.

$^{841}$可能上一代有時候有些東西不懂得.
來幫助你.
或者天生天養.
不過我們是中國人.
我們就知道整個孝道本身就是.
不是因為他做了好事你孝順他.
而是正正的.
他的身份是你父親母親.
你就孝順他.
不過怎樣才算孝順.
這個可以商量的.
可以有調教.
不會說你不尊重父母.
或者當父母是朋友.
純粹當平輩看待.
在中國人的社會裡面.
是有這一面很重要的.
在猶太人裡面也是一樣的.
在十界裡面其中一界就是當孝順父母.
而整個的.
不論是中國人還是猶太人.
其實說的孝道精神.
其實是建構整個社會裡面的秩序的一個很重要的核心.
如果你沒有孝道精神.
其實你瓦解了家庭.
其實說的君臣父子社會制度.
在中文社會裡面是會瓦解的.
同樣道理.
在猶太人或者在舊約聖經裡面.
或者去到新約也是.
孝道精神是成為了敬拜上帝一個最基礎的訓練.
從小到大學會敬畏上帝.
也學會孝順父母.
孝順父母和敬畏上帝就是一個下而上的一種尊崇.
這種觀念其實是形成了他們敬拜的觀念.
如果沒有了這一點.
抽離了的時候.
在舊約裡面.
你不孝順父母.
某程度上就等於藐視上帝.

$^{881}$在阿古爾裡面.
我覺得這個是荒謬的.
如果對於一個口口聲聲說敬拜神的人.
但最後他不孝順自己的父母.
其實他的生活和信仰是撥亂來的.
是沒有意思的.
是荒謬的.
第二個.
自以為清潔卻沒有洗去自己的污穢.
這是一個生活的盲點.
那個潔淨和污穢是在道德上.
以為自己站在道德的高地.
原來他自己看不到自己身上的污穢.
耶穌基督去到新約.
說的那些法利賽人.
那些文士針對耶穌咬牙切齒要痛罵的.
通常就是一種這樣的生命狀態.
第三.
這個很直接了當.
眼目高傲.
眼皮也是高舉.
在這裡說的是白鴿眼.
看不起人.
永遠都是這樣看人的.
他的意思就是這樣.
我們生命裡有沒有覺得自己會比某些人優越呢.
覺得自己是一個優越分子.
即使你不是這麼優越.
你也覺得自己很優越.
這就是在《雅古意》裡面說的這個世代的荒謬.
第四種人.
不是說他的說話很靈麗.
不是說牙尖嘴利.
而是說牙齒和牙如劍.
齒如刀就是隨時吞噬.
準備要啄別人的東西.
啄那些人的東西呢.
就是地上困苦的窮乏人.
僅餘的東西.
他都在侮辱別人.

$^{921}$貪念.
或者想多些在別人身上.
獲取了別人僅餘的東西.
這就是以自我為中心的四種人.
這四種人其實都是同出一轍.
正如我剛才所說.
在自己很重要.
別人沒那麼重要.
以自我為中心.
在第一個段落.
用這四種人或者這四種的人.
來作為一個小結.
其實就是在說.
今天我們活在這個荒謬的世界.
我們會和這個世界的人同流合污.
一方面我在罵那些人.
但另一方面我正在成為那些人.
或者我謹慎.
我知道自己沒那麼了不起.
我只是尋求一個正在幫助我們的上帝.
在他的話語當中.
尋找我們人生應該要應有走出來的出路.
今天的心理學.
或者今天的輔導學很勵志.
我都是一個很陽光很正面的人.
不過我發覺原來過份勵志.
過份正面有一個危機.
原來過份勵志.
過份正面的結果.
會令人的內心變得很膨脹.
會覺得自己是全世界的中心.
覺得自己是很重要.
當我們面對這樣的狀況.
聖經是幫到我們.
聖經往往是將我們自己進行一些解構.
或者有一些否定.
這些否定不是說我們是地底泥.
上帝不要我們.
而是那種否定是要抗衡這個世界.
令到覺得自己很重要的觀念.

$^{961}$得到重新的解構.
加拉太書六章七節說到.
人若無有自要.
自己還以為有就是自欺.
今天這個世界可能有很多東西.
告訴你.
感覺你自己的感受.
你自己的觀點.
或者你自己覺得是全世界的中心.
很重要.
不過你沒有.
你說為有.
你自己騙自己.
羅馬書第十五章第三節.
我們看回耶穌是怎麼說的.
因為基督也不求自己的喜悅.
如經上所記.
肉麻你的肉麻都落在我身上.
這句話說什麼呢.
耶穌在他的全道生涯裡面.
他往往都是做一種自我否定的方式.
他不是否定他自己的神聖.
或者自己的任務.
而是他看自己當下的體會感受.
算得了什麼.
自我否定.
成為了一個屬靈裡面.
重新檢視自己的一個向導.
弟兄姊妹.
當今天我們活在這個世界裡面.
我覺得自己的感覺良好.
我覺得自己擁有的東西.
或者我覺得別人冒犯了我們.
問問自己.
你在上帝面前.
你不是真的想像中你自己這麼重要嗎.
我們生命裡面.
今天祈求的是為求上帝給我更多.
還是給我更真.
今天在我的生活當中.

$^{1001}$我祈求上帝給我更豐裕的物質.
還是我祈求上帝幫助我.
更加敬虔地度過我的人生.
阿古義這樣教導他兩個學生.
以鐵和烏鴿.
這樣去祈禱.
這樣去理解上帝.
盼望今天真言也這樣教導我們.
過我們人生.
踏實地過日子.
讓我們人生裡面.
不會像這個世界一樣繼續荒謬下去.
讓我們在上帝面前成為他所喜悅的子民.
(下集再見).
\newpage



\section{希伯來書 10:19-11:3}
\label{sec:3GKCAsbLUbA}
\textbf{2022年5月1日崇拜直播｜陳明泉牧師｜超越的基督:親近上帝之路｜希伯來書10\_19-11\_3}
\newline
\newline
連結: \href{https://youtube.com/watch?v=3GKCAsbLUbA}{\texttt{https://youtube.com/watch?v=3GKCAsbLUbA}} ~~~~ 語音日期: 2024-03-19
\newline
\newline
\hyperref[sec:8EXqdqFDCIc]{\small{< < < PREV SERMON < < <}}
~
\hyperref[sec:index_chronic]{\small{[返順時目]}}
~
\hyperref[sec:l8_t1yVcPHE]{\small{> > > NEXT SERMON > > >}}
\newline
\newline
希伯來書 10:19-11:3
\newline
\begin{longtable}{cl}
\hline
\hline
章節 & 經文 (和合本修訂版)\\
\hline
10:19 & \begin{tabularx}{0.7\textwidth}{X} 所以，弟兄們，既然我們靠著耶穌的血得以坦然進入至聖所， \end{tabularx} \\ \\ \relax
10:20 & \begin{tabularx}{0.7\textwidth}{X} 是藉著他給我們開了一條又新又活的路，從幔子經過，這幔子就是他的身體。 \end{tabularx} \\ \\ \relax
10:21 & \begin{tabularx}{0.7\textwidth}{X} 既然我們有一位偉大祭司治理神的家， \end{tabularx} \\ \\ \relax
10:22 & \begin{tabularx}{0.7\textwidth}{X} 那麼，我們該用誠心和充足的信心，同已蒙潔淨、無虧的良心，和清水洗淨了的身體來親近神。 \end{tabularx} \\ \\ \relax
10:23 & \begin{tabularx}{0.7\textwidth}{X} 我們要堅守所宣認的指望，毫不動搖，因為應許我們的那位是信實的。 \end{tabularx} \\ \\ \relax
10:24 & \begin{tabularx}{0.7\textwidth}{X} 我們要彼此相顧，激發愛心，勉勵行善； \end{tabularx} \\ \\ \relax
10:25 & \begin{tabularx}{0.7\textwidth}{X} 不可停止聚會，好像那些停止慣了的人，倒要彼此勸勉，既然知道那日子臨近，就更當如此。 \end{tabularx} \\ \\ \relax
10:26 & \begin{tabularx}{0.7\textwidth}{X} 如果我們領受真理的知識以後仍故意犯罪，就不再有贖罪的祭物， \end{tabularx} \\ \\ \relax
10:27 & \begin{tabularx}{0.7\textwidth}{X} 惟有戰戰兢兢等候審判和那將吞滅眾敵人的烈火了。 \end{tabularx} \\ \\ \relax
10:28 & \begin{tabularx}{0.7\textwidth}{X} 任何人干犯摩西的律法，憑兩個或三個證人，尚且必須處死，不得寬赦， \end{tabularx} \\ \\ \relax
10:29 & \begin{tabularx}{0.7\textwidth}{X} 更何況踐踏神兒子的人，他們將那使他成聖之約的血當作不潔淨，又褻慢施恩的聖靈的人，你們想，他不該受更嚴厲的懲罰嗎？ \end{tabularx} \\ \\ \relax
10:30 & \begin{tabularx}{0.7\textwidth}{X} 因為我們知道誰說：「伸冤在我，我必報應。」又說：「主要審判他的百姓。」 \end{tabularx} \\ \\ \relax
10:31 & \begin{tabularx}{0.7\textwidth}{X} 落在永生神的手裡真是可怕呀！ \end{tabularx} \\ \\ \relax
10:32 & \begin{tabularx}{0.7\textwidth}{X} 你們要追念往日；你們蒙了光照以後，忍受了許多痛苦的掙扎： \end{tabularx} \\ \\ \relax
10:33 & \begin{tabularx}{0.7\textwidth}{X} 一面在眾人面前公然被毀謗，遭患難；一面陪伴那些受這樣苦難的人。 \end{tabularx} \\ \\ \relax
10:34 & \begin{tabularx}{0.7\textwidth}{X} 你們同情那些遭監禁的人，也欣然忍受你們的家業被人搶去，因為你們知道自己有更美好更長存的家業。 \end{tabularx} \\ \\ \relax
10:35 & \begin{tabularx}{0.7\textwidth}{X} 所以，不可丟棄你們無懼的心，存這樣的心必得大賞賜。 \end{tabularx} \\ \\ \relax
10:36 & \begin{tabularx}{0.7\textwidth}{X} 你們必須忍耐，使你們行完了神的旨意，可以獲得所應許的。 \end{tabularx} \\ \\ \relax
10:37 & \begin{tabularx}{0.7\textwidth}{X} 因為「還有一點點時候，那要來的就來，必不遲延。 \end{tabularx} \\ \\ \relax
10:38 & \begin{tabularx}{0.7\textwidth}{X} 只是我的義人必因信得生；他若退縮，我心就不喜歡他。」 \end{tabularx} \\ \\ \relax
10:39 & \begin{tabularx}{0.7\textwidth}{X} 我們卻不是退縮以致沉淪的那等人，而是有信心以致得生命的人。 \end{tabularx} \\ \\ \relax
11:1 & \begin{tabularx}{0.7\textwidth}{X} 信就是對所盼望之事有把握，對未見之事有確據。 \end{tabularx} \\ \\ \relax
11:2 & \begin{tabularx}{0.7\textwidth}{X} 古人因著這信獲得了讚許。 \end{tabularx} \\ \\ \relax
11:3 & \begin{tabularx}{0.7\textwidth}{X} 因著信，我們知道這宇宙是藉神的話造成的。這樣，看得見的是從看不見的造出來的。 \end{tabularx} \\ \\
[1ex]
\hline
\hline
\end{longtable}
$^{1}$今日講到的經文是記載在希伯來書.
十章十九節至到十一章三節.
聖經這樣說.
弟兄們.
我們既因耶穌的血.
得以坦然進入至聖所.
是藉著他給我們開了一條.
又新又活的路.
從萬子經過.
這萬子就是他的身體.
又有一位大祭司治理神的家.
並我們心中天良的虧欠.
已經灑去.
身體用清水洗淨了.
就當存著誠心.
和充足的信心來到神面前.
也要堅守我們所承認的指望.
不致動搖.
因為那應許我們的是信實的.
又要彼此相顧.
激發愛心.
勉勵行善.
你們不可停止聚會.
可像那些停止慣了的人.
都要彼此勸勉.
既知道那日子臨近.
就更當如此.
因為我們得知真道以後.
若故意犯罪.
贖罪的祭就再沒有了.
唯有戰具等候審判.
和那消滅眾敵人的烈火.
人干犯摩西的律法.
憑兩三個見證人.
尚且不得連戌而死.
何況人踐踏神的兒子.
將那使他成聖滋藥的血.
當作平常.
又洩漫私恩的聖靈.
你們想他要受的刑罰.

$^{41}$該怎樣加重呢?.
因為我們知道誰說.
「申冤在我」.
「我必報應」.
又說.
主要審判他的百姓.
落在永生神的手裡.
真是可怕.
你們要追念往日.
蒙了光照以後.
所忍受大爭戰的各樣苦難.
一面被毀謗.
遭患難.
成了氣境.
一面陪伴那些受這樣苦難的人.
因為你們體恤了那些被困所的人.
並且你們的家業被人搶去.
也甘心忍受.
知道自己有更美長存的家業.
所以.
你們不可丟棄勇敢的心.
存這樣的心必得大上前.
你們必須忍耐.
使你們行完了神的旨意.
就可以得著所應許的.
因為還有一點點時候.
那要來的就來.
並不遲延.
只是二人必因順得生.
他若退後.
我心裡就不喜歡他.
我們卻不是退後入沉淪的那等人.
乃是有信心.
以致靈魂得救的人.
第十一章一節.
信就是所望知事的實體.
是未見知事的確具.
古人在這信上得了美好的證據.
我們因著信.
就知道諸世界是藉神話造成的.

$^{81}$這樣所看見的.
並不是從顯然之物造出來的.
弟兄姊妹平安.
今天見到教會坐得很滿座.
見到弟兄姊妹熟悉的面孔.
心裡一回來都想流眼淚.
不知道你是否這樣.
我是這樣.
今天能夠在早上回來.
一切都不是必然.
不過人有些很特別的.
惰性或者有些心態很有趣.
當我們.
不知道你回來之前帶著什麼心態.
很期待.
有些十五十六.
我聽過弟兄姊妹.
停了幾個月之後.
就真的很想回教會崇拜.
不單是回教會.
上學也是這樣.
上班也是這樣.
家裡困久了.
都想有實體這樣好一點.
但到真的有實體來到的時候.
心裡又想.
其實在家裡看Zoom也不錯.
看著方便很多.
不用這麼早起床.
也不用這麼忙.
人就很奇怪.
因為習慣了某種新舊的模式.
將會再轉回去.
總是有些東西卡住.
令自己好像捱不過.
或者很多時候沒辦法突破自己一步.
當有一個新的機遇來到的時候.
可能我們又已經習以為常.
有一些生命模式.
以致我們沒辦法繼續向前.

$^{121}$今天的讀經文.
剛才你很細心去聽的時候.
其實希伯來書的作者.
都是針對一班信徒.
面對生命的一種狀況.
特別是說到有一些弟兄姊妹.
或者他看到一些現象.
停止聚會已經停止慣了.
生命裡面可能進入了一種狀態.
好像上帝可有可無.
變得沒有上帝也習以為常.
面對一種這樣的狀態下.
希伯來書的作者怎樣去勉勵他們呢?.
希望今天我們藉著這一卷.
比較長的篇幅.
我們重新看看.
可能我們在舊有的模式裡面.
或者我們都習慣了某一種的生活形態.
如果我們想再一步.
重新去親近上帝.
我們可以怎樣做呢?.
我們可以帶著什麼樣的態度.
來對抗自己生命裡面.
習以為常的那種形態.
希望今天這篇講章可以幫助我們.
重新釋入我們熱切追求上帝的那種生命.
我們一起祈禱.
求上帝幫助我們.
求聖靈開啟我們的心靈.
我們一起去親近祈禱.
天父上帝願意在我們當中帶領我們.
讓我們能夠在希伯來書裡面.
我們拿到屬靈的養分.
我們生命裡面.
讓我們能夠可以繼續前進.
讓我們這一班弟兄姊妹.
我們生命被激勵.
讓我們帶著上帝的華語.
有力量走面前的路.
我們恭敬將來的時間交託給你.

$^{161}$幫助每位聆聽的弟兄姊妹.
他們每一個都聽得明白.
聽盡入心.
以致他們願意用生命來回應.
也幫助宣講的.
強而有力地將你的華語.
清楚地宣講.
讓人得著福音.
得著上帝的慰養.
讓我們每一個彼此.
為上帝努力.
我們恭敬將來的時間交託.
願你賜福.
獻上祈禱.
奉靠基督耶穌.
得勝的聖名祈求.
Amen.
希伯來書第十章第十九節開始.
你會看到語氣有一點改變.
一開始說的是弟兄們.
這個弟兄們.
很多時候我們用這些來做分段.
是一個很好的指引.
因為已經是丞相啟下了.
到了第十章第十九節某程度.
已經是進入應用和結論的部分.
整個希伯來書.
所以他不會再說一些比較抽象.
或者再說一些比較深奧層次的神學論述.
反而是將之前說了很大篇幅.
由第一章到第十章十八節.
那些關於基督耶穌是大祭司.
耶穌如何獻祭.
聖所是如何.
這些作為一個開始總結.
然後就說到弟兄姊妹.
如何做一個應用.
所以你看十章十九節裡.
他重複有些字眼.
這些字眼是說.

$^{201}$之前說過的神學內容.
例如說到耶穌的血.
大祭司.
之前我們說了幾次都是這樣.
進入聖所.
或者這裡譯為至聖所.
去到神的面前.
等等這些字眼.
他已經不再去仔細地去再.
告訴你.
解釋給你聽是一回甚麼事.
因為之前已經說了.
但將這一切的概念.
將這一切的神學觀念濃縮了.
他就說到.
我們為何今天可以敬拜神呢.
我們是因為耶穌的血.
所以可以坦然地進入至聖所.
之前說過至聖所是大祭司可以進入的.
現在因為拉著耶穌的衣服.
因為耶穌是大祭司.
以前一般百姓是不能進入的.
現在我們拉著耶穌的衣服.
我們可以進入到聖所裡面.
這個方式為何可以改變呢.
希伯來書的作者告訴我們.
是因為耶穌為我們開了一條.
又新又活的道路.
經過的萬子.
然後就是耶穌的身體.
然後就直接被大祭司治理神的家.
進入到聖所當中.
值得留意的就是.
在書裡面提到的又新又活的道路.
這條路.
提到新的道路.
我們當然想著有條舊路.
其實新不是對應舊.
新是對應舊時的沒有.
人能夠去到上帝的根前.

$^{241}$這麼近距離.
其實舊約裡面.
過去以色列人的先祖.
都沒有人做得到這件事.
最近是摩西了.
不過去到今天.
我們竟然每個人都能去到上帝的面前.
這條新的路.
這條讓人走的生命的活的道路.
因著基督耶穌的獻祭.
耶穌成為了大祭司.
以致我們可以進入到.
這條路進入至聖所去到上帝面前.
經文再特別的說.
我們心中的天良虧欠已經灑去.
天良弟兄就在那裡.
但心中的天良虧欠.
他的意思就是.
良心.
你舊時良心責備你的那種罪疚感.
現在就可以放下了.
我們今天每個人都是罪人.
每個人都是不配不理想.
不過因為耶穌的救贖.
我們這些不配不理想的人.
竟然可以去到上帝的面前.
希伯來書的作者更加呼籲我們.
將那些罪疚感灑去.
放下它.
耶穌已經洗乾淨了.
你的身體潔淨了.
他其實這樣說就是.
我們在上帝面前成為一個被潔淨的人.
哪怕我們以前是多麼的枯萎.
生命裡面多麼的千瘡百孔.
上帝藉著他的愛子耶穌.
令到我們個個潔淨了.
所以去到他的面前.
我們要帶著誠心和充足的信心.
去到上帝的面前.

$^{281}$在這裡再解釋多一點.
那個字就是洗去乾淨.
身體被潔淨.
心靈裡面的虧欠感被潔淨.
他用的字不是用洗潔精的洗.
他是用灑.
灑那個血灑去.
以致我們就得潔淨.
這個動詞為甚麼會這麼有趣.
用灑字的呢.
原來是舊時在舊約裡面.
那些祭司獻祭生.
那些祭生的血.
如果你留意回到利未記裡面.
他潔淨的那件事.
其實不是真的像我們今天所說的那種衛生.
找個刷刷,擦擦那些.
其實是在說那個祭生的血.
灑到那個人或那些祭物身上.
那樣東西就被潔淨了.
所以那個是一個被上帝閱立的那種方式.
灑的意思就是說.
耶穌的血就好像舊時祭司獻上祭生的那些血.
灑到那些不同的物件.
或是要獻祭的那些物件上.
按照今天的衛生條例那些是污糟的.
有血在那裡的.
但是在宗教來講.
那個是乾淨的.
同樣道理.
我們今天每一個.
心裡面的那個虧欠,罪疚感.
可以清除.
身體得到潔淨.
不是因為上帝做了一些什麼洗潔精那種功夫.
而是在神學,在建制那裡.
耶穌就將他的血灑在我們每一個身上.
所以我們每一個在上帝眼前成為一個被閱立,被分別出來的人.
我們帶著一個這樣的狀態.
就可以去到上帝的根前.

$^{321}$親身地去面對上帝.
面對上帝在舊約裡面是匪夷所思的.
罪人去到上帝面前.
只有一個結果.
就是被擊殺.
污糟的人在聖潔的上帝面前是無地自容.
不過今天.
因著基督耶穌的獻祭.
每一個人都可以去到上帝的根前.
接下來.
講到這個神學.
已經總結了.
之前講的東西已經重溫了.
你知道了.
然後怎樣去回應呢?.
由23節至到第25節.
其實就是一個用信望愛的框架.
既然人去到上帝面前.
人怎樣去回應上帝潔淨的這個恩典呢?.
他用信望愛來做一個框架.
輕輕地講.
第一個就是22節.
我們傳授誠心和充足的信心.
來到神面前.
在這裡你會看到.
誠心這個字其實譯得OK的.
譯得好的.
原文是在講真理.
真誠的那個心.
真實的意思.
傳授真心.
誠心其實我們都明白的.
即使我們多麼的污穢.
被上帝潔淨了.
我們都可以帶著一份真誠的心.
去到上帝面前.
不用裝模作樣.
不用裝敬虔.
不用刻意將自己提升了.
來到敬拜神.

$^{361}$你是怎樣.
就帶著你怎樣的款去到上帝面前.
希伯學書學者講.
在上帝面前甚麼都知道.
你的罪都已經被潔淨.
你就將你的生命原原本本.
放在上帝面前.
做回自己.
將你真實的你.
在上帝面前.
來到去敬拜.
第二個字.
充足的信心.
充足一.
一看這個字.
我們有種感覺就是.
數量上.
信心多與少的問題.
不過原本的意思不是說.
數量的充足.
大量的信心.
不是這個意思.
是說完備的信心.
即是說.
我知道無論我是怎樣都好.
我信得過上帝.
祂不會丟下我.
所以整個信心的框架.
意思就是.
我們每一個.
我們帶著我自己的真我.
去到上帝面前.
帶著一份確信.
帶著那份完備的信.
認定上帝.
無論我在甚麼狀態之下.
祂都不會撇棄我.
帶著一種這樣的心態.
走到上帝面前.
在牧羊的經歷裡.

$^{401}$我知道有些弟兄姊妹.
有時不敢回教會.
為甚麼呢?.
因為她看回自己.
看回自己生命裡面的紀錄.
自己知道.
覺得自己實在太差了.
一回到教會.
她覺得自己.
面對著那個講壇.
面對著那個十字架.
她就有一種很大的虧欠感.
這種虧欠感.
如果去到極致.
我們有時會採取一種.
逃避上帝的態度.
希伯來書的作者說.
不用這樣.
昔日你做錯了的事.
上帝知.
上帝已經潔淨了.
只要我們帶著一份真誠.
以及認定上帝不會放棄我.
我們就將我這個生命.
去到上帝跟前.
這個就是.
在這裡略略地先說了.
那種信心的框架.
來表達人對上帝的回應.
十章二十三節.
這個是在說盼望.
也要堅守我們所承認的指望.
不自搖動.
因為那應許我們的是信實的.
這個指望.
其實就是我們今天所說的盼望.
不過這個盼望的基礎是什麼呢.
這個盼望的基礎是在說上帝的信實.
上帝答應了人.
上帝在聖經裡面說.

$^{441}$上帝對人的應許.
我不撇下你成為孤兒.
我要與你永遠的同在.
無論人處於什麼狀態.
人約自己的罪.
神是公義的,信實的.
必要赦免我們的罪.
洗淨我們一切的不義.
等等這些.
上帝答應了我們的.
我們都是帶著這一份這樣的指望.
雖然眼前這一刻.
有時候覺得好像逆流而上.
但是我們深信.
我們帶著這樣的盼望.
去到上帝面前的時候.
我們堅守著這一種的狀態.
上帝是會閱立的.
除了盼望這個字是重點之外.
或者指望.
另一個字眼就是堅守.
堅守的意思就是.
不是一時一刻.
不是那一刻很興奮.
然後下一刻又再退.
人生常常都是這樣上上落落很累.
是不是.
他所說的堅守是.
我那一刻很興奮.
但是之後離開了這個會堂.
或者生活當中.
我都帶著.
不要說一直興奮.
但是帶著一份的確信.
在我的生活當中.
我保持著.
我自己的屬靈狀態是這樣.
所以.
不要將每一個星期.
好像提振一下.

$^{481}$催谷一下你那樣.
原則上應該是.
每一次崇拜之後.
你就帶著那份感恩.
和那種對上帝的信心.
你要保持著.
你的靈修其實就是.
幫你保持著你的屬靈生命.
以致到下星期我回來的時候.
我都是同樣的狀態.
不要讓你的屬靈狀況.
就好像不健康的人的血壓一樣.
有時高有時低.
血壓你都搞不定.
何況屬靈生命裡面的堅持.
保持在一種穩定的狀態.
堅守著你的那份盼望.
這是我們第二個對上帝的回應.
在這裡值得留意.
說信心和盼望.
在後來的篇幅會發揮更多.
所以在這裡.
希伯來書的作者.
輕輕地清停點水地說.
因為後來還有更多的篇幅.
是說如何堅守.
如何信靠.
就成為後面結論的部分.
一個很重要的回應上帝的方式.
所以在這裡.
他不是說很多的.
落筆去說這件事.
更加重要的是.
後來裡面所說.
他會將信心和盼望兩者.
混為一談.
很多時候我們其實說信望愛.
都不是完完全全可以.
清晰地來界分的.
兩者之間.

$^{521}$或者三者之間都有重疊的地方.
希伯來書的作者就在後面的篇幅.
將信心和盼望.
成為後面結論的部分.
這裡不多說.
我不再繼續用很多篇幅.
反而愛心在這裡.
就多說一點.
這裡裡面.
在第24節.
特別的來講.
又要彼此相顧.
激發愛心.
勉勵行善.
讀完25節.
你們不可停止聚會.
好像那些停止慣了的人.
不要彼此勸勉.
既知道那日子臨近.
就更當如此.
愛心如何去回應上帝呢?.
這個愛心.
他是在說.
不單止你對我好.
我對你好.
而是在說這個愛心.
就好像滾雪球一樣.
大家互相以愛相待的時候.
是會製造更加大的愛心.
來對待身邊的人.
而且這一份的愛心.
是有激勵性的.
他用的字眼很特別.
激發愛心.
這個激發這個字.
我們讀上去.
很正.
很有爆炸力.
不過在原文裡面的激字.
其實是很負面的字眼.

$^{561}$有些什麼字.
新約裡面如何用激字呢?.
在說吵架.
《士徒行傳》裡面說到.
保羅和巴拿巴吵架的時候.
大家奮然離開.
搞到吵架而不歡而散.
這個字就是激發這個字.
在舊約裡面.
七十事易本拿回來.
來譯舊約.
說的是耶和華的怒氣發作.
你可想而知.
這是一個很爆炸.
而且是一些很生氣的字眼.
負面的.
很多的用法.
普遍的用法.
是一個負面的.
不過來到這裡.
就不是負面的.
不過他就想說那種爆發.
吵架那種能量.
他用這個來成為愛心的表達方式.
你做這件事這麼好.
我又要做多點.
大家好像互相激勵.
以致愛心就有種爆炸性的果效.
讓身邊的人得著好處.
而且其他生命的人被激勵.
弟兄姊妹你今天回到教會.
你想一想.
我們會不會有時候.
我們偏向.
華人教會也好.
我們偏向有一種含蓄.
我們有時候會等待別人先做.
又或者我們有時候.
未必是錯的.
我們含蓄慣了.

$^{601}$我們很少做一個很主動.
去表達自己的愛心.
不過在希伯來書的作者裡面.
如果你要更新.
你要回應上帝.
耶穌基督帶你到上帝面前.
當你知道這個恩典這麼大的時候.
你就不要等了.
你的愛心就不要變得很被動.
或者變得一種很消極的信心.
那種愛心.
別人叫都好啦.
這樣想清想楚.
不是叫你不想清想楚.
只不過是不要常常處於一種被動的狀態.
處於一種很消極的態度.
來實踐愛心.
希伯來書的作者說到.
弟兄姊妹的愛應該是有一種爆炸性.
你越願意為上帝去做.
在希伯來書的作者裡面.
就是越能夠回應基督耶穌的這個恩典.
把你的愛心爆炸性地分享給身邊的人.
所以激發這個字是正面的.
不過我在大家的大綱裡面.
我寫了另外一個字眼.
一個負面的字眼.
就是消磨大家的愛心.
什麼叫消磨大家的愛心呢.
有時候我們做了一些事.
譬如我們很主動地想做一些愛心的行為.
不過今天社會有一些不好的風氣.
當我們做一些事的時候.
我們容易製造一些批評.
我們有時候對於人自發去做一些美事.
我們第一件事有時候去想.
有沒有這麼好死.
行不行的.
做得不好啊.
做得不好你不要做了.

$^{641}$不知道你會不會聽過這樣的字眼.
以致到本來我想有些心想做的時候.
接著就幾十萬句指點你.
或者批評你.
到最後只是做一點點事.
但是就製造了一個氛圍.
就好像付出最多的那些.
接受最多的批評.
那個就是消磨愛心.
我在這裡我也要幫我們教會的領袖.
一班的弟兄姊妹說一句話.
今天教會有時候.
有一些不是很好的習性.
做事的人永遠被針對.
就是有一些長期做事的人.
就長期被人針對.
就是我們將這個群體.
成為一個做事承擔的那些被批判的對象.
如果這樣做的時候.
到最後是什麼呢.
反正都做了被人罵.
不做不被人罵.
還可以罵人.
那我當然做罵人那個.
我當然不做承擔那個.
久而久之.
教會的關係就瓦解.
我講的消磨愛心就是這件事.
留意一下各位弟兄姊妹.
有沒有有時候我們很容易.
講一些話就批評了別人的努力.
有時候有些付出.
跟坊間那些最專業的沒得比的.
但是我們可不可以從一個激發愛心的角度.
來勉勵身邊的人.
多過我用一個最高的標準.
用坊間裡面最沒什麼人達到的層次.
來強加在身邊的人身上.
當我們成為一個這樣的人.
我可以這樣講.

$^{681}$我們群體就沒有前景了.
做事的你要格外來欣賞他的殷勤.
當他付出努力的時候.
希望你同樣被激發.
你同樣用愛心來找自己的位置.
來將這個愛心的行為來發大他.
希伯來書的作者提醒我們.
我們不要成為一個輕易批評別人努力的信徒.
一個弟兄姊妹.
反而我們的愛心是要互相激勵.
你做多一點我做更多.
讓整個愛心的餅越做越大.
這個群體才有那種衝勁和吸引力.
而不是互相消磨.
弟兄姊妹生命裡面擺上的一切美好.
經文裡面再繼續講.
除了愛心之外.
那個聚會.
剛才我也講過了.
那個聚會用的字眼就是.
停止慣了.
那個慣呢.
就是停一次.
然後停第二次.
很多時候都是這樣的.
一個習慣是怎樣建立而來的呢.
或者壞習慣是怎樣出現的呢.
連續21天.
有些人是這樣講過.
連續21天做一件事.
你就會建立一個習慣.
連續21天你不做某一件事.
你做了你不做.
你就去除了這個習慣.
如果在我們自己的信仰裡面.
連續21天我沒有靈修.
連續21天我沒有回教會.
沒有參與在教會裡面.
你的習慣性就會退減.
正面和負面都是這樣看的.

$^{721}$有時候.
這樣講.
剛剛我們有三個多月.
那個習慣就不是我們想要的習慣.
我們在疫情之下.
我們被迫著都要習慣.
不過當我們要抗衡這件事的時候.
我們就要重新學一個新的習慣.
重新建立自己的生活紀律一樣.
我們用一個很好的申命記律.
來幫自己重新建立一個生命的慣性.
所以頂尖我知道.
是逆流而上的.
在疫情之下我們停了實體這麼久.
然後再回來.
有很多事情很不習慣.
有很多事情覺得很奇怪.
不過.
在上帝的恩典當中.
我們能夠因為基督耶穌去親近神.
我們因為這個恩典.
我們願意改變自己.
待會兒會再講停止慣了的那件事.
不過在這裡首先我要總結.
今天我們第一個段落.
為什麼我們可以回來教會崇拜呢?.
為什麼我們可以.
我們覺得很順理成章.
為什麼我們要這樣做呢?.
不是因為我很偉大.
我為了上帝做一些事情.
不是覺得上帝很好.
我就做一些事情.
好像做一些功德一樣.
為自己做一些事情.
整個希伯來書不是在講這件事.
是在講.
因為耶穌基督幫我們開一條新的路.
以前人們都做不到.
現在將這個恩典.

$^{761}$廣泛地傾倒給你們每一個.
我們回應這個恩典最好的方式就是.
信忘愛.
所有事情都是因著基督耶穌.
以致我們有這個機會.
在舊約裡面.
人們想恨.
都恨不到我們今天那種形態.
但我們今天很多事情.
圍著很多很簡單的事情.
我們就放棄了.
我們就消磨了.
甚至有時候我們會採取一種.
不信任敵黨的態度.
聖經裡面提醒我們.
是因為基督耶穌的寶血.
讓祂為我們開一條新的.
又新又活的路.
我們可以坦然無懼.
去到上帝的致勝所.
去到上帝的根前.
我們帶著這一份確信.
我們再用信忘愛來調整我們.
我們的生命就會容易走得更遠.
我們繼續說.
十章二十六至第三十一節.
在二十六節裡面.
其實是回應剛才所說的.
停止聚會的再晉新版的問題.
二十六節.
因為我們得知真道以後.
若故意犯罪.
贖罪的祭就再沒有了.
唯有戰具等候審判.
和那消滅眾人的烈火.
人干犯摩西的律法.
憑兩三個人見證人.
尚且不得連戌而死.
何況踐踏神的遺址.
將那使他成聖之弱的血.

$^{801}$當作平常.
又洩漫私恩的聖靈.
你們想.
他要受的刑罰該怎麼加重呢.
因為我們知道.
誰說.
身冤在我.
我必道報應.
又說.
主要審判他的百姓.
落在永生神的手裡.
真是可怕的.
在第二十六至三十一節裡面.
作者就用一個比較負面.
反面的角度.
來講到.
你不要輕慢.
這條又新又活的道路.
既是一個因點.
這個銀的另外一面.
就是講到.
如果你忽略了這個救恩.
你輕忽了這件事.
你生命是落入一種很可怕的狀態.
經文特別提的字眼.
就是故意犯罪.
故意犯罪.
我們以為.
我們一讀.
我們會害怕.
因為我們想.
是不是我們生命裡面.
做了一些上帝不起悅的事.
做了一些不好的行為.
有時會做.
以致到叫做故意犯罪.
上文下理講的故意犯罪.
不是講我們做了一些什麼行為.
雖然上帝不會欣賞.
或者不起悅我們這樣做.

$^{841}$不過這個故意犯罪.
其實是有上文下理講的.
這個故意犯罪.
其實是在講敵黨的信仰.
敵黨基督信仰.
有一些不犯罪會的.
那個字眼.
停止罪會的停止.
Latel.
和後面贖罪制.
不再存留.
有個逆的字.
他同樣用的字根.
那個Latel的字.
是在講一種上帝.
對上帝的敵黨.
是一個進程.
這樣講好像很抽象.
解釋多一點.
有些弟兄姊妹.
可能生命裡面有些事.
在教會裡面有些生活不理想.
他停止了參與在教會的罪會.
停止了一段時間之後.
停了.
不過到最後.
他連認信他的信仰.
他都失去了.
以致到.
他好像一個屬靈.
一路走下坡.
很少人.
我還沒有聽過.
他停止了罪會.
但是他愛主的心.
更加的深.
如果愛主的心更深.
有時都是自己講.
他的信仰的投入程度.
毀身程度一定是越來越低.

$^{881}$因為你不返罪會.
很難毀身到在信仰裡面.
在希伯來書的作者裡面.
用同樣的字根.
講到一個程度上的分別.
不返罪會.
未至到故意犯罪.
不過不返罪會.
停止了罪會.
一路一路越走.
一路一路越走下坡.
到最後.
很多時就會成為一個.
敵黨信仰.
否定信仰的一個人.
我認識有.
有一些弟兄姊妹.
在他們年輕的日子.
在教會裡面.
返得很好.
很投入.
後來因為在教會裡面.
有一些人事.
和感情問題.
結果就停止了.
返教會.
我就說.
到最後你真的覺得很尷尬.
不知如何處理.
你不返這裡.
你還是返其他教會.
但有些返其他教會後.
覺得旺朕好像好一點.
或者形態和旺朕很不同.
返下ABCDEFG.
轉下轉下.
到最後就索性不返.
現在好些.
有得上網聽聽道.
偶爾聽聽道.

$^{921}$不過到最後.
一路一路將信仰離身.
有一次見到他.
在街裡見到弟兄姊妹.
我問他近來怎樣.
他帶著一種.
和你很有距離感的表達.
他覺得你不要和我聊天.
我其實現在某程度.
我都不是很相信.
你不要當這個是虛幻出來.
其實在現實生活裡面.
有很多基督徒.
很多生命.
其實面對著可能有些傷害.
他處理不到.
一路一路離開的時候.
到最後他走向一個棄教.
拒絕上帝的一種心靈狀態.
我們明白.
甚至乎我可以坦白講.
教會裡面真是.
有時我們人的群體.
都會充滿很多問題.
很多的荒謬.
甚至乎有時有些事.
真是很大的傷害.
我不會為這件事去護短.
實際上教會幾千年.
其實就是滿佈問題的一個群體.
不過上帝竟然選擇一個.
滿有問題的群體.
繼續延展他的福音工作.
我們因為這個群體而有問題.
而放棄了信仰.
那個代價實在太大了.
有一本書叫做.
靈魂倖存者.
Philip Yancey 寫的.
我很年輕就看這本書.

$^{961}$Philip Yancey 是一個什麼人呢.
是一個想事想到牛角尖的一個作家.
一個記者.
為什麼他會寫這本書呢.
他說自己是靈魂倖存者.
他說自己在成長的教會裡面.
他用一個字眼叫做toxic church.
有毒的教會生活.
他從小到大都被一些錯誤的教導.
很偏頗,很分割.
那種很法利賽式的成長背景下.
後來他自己才再回顧.
發現原來他在這樣的養分下長大.
他有一刻很厭惡,很掙扎.
他問為什麼要有什麼理由.
讓他自己繼續回教會.
他忽然間覺得自己在回教會.
為什麼每個星期都如是呢.
到最後他在這本書.
或者在他自己.
去到今天他都有一個結論.
今天人回教會不是因為教會.
或者投入在基督教信仰.
不是因為教會有多好.
他說的不是因為教會有多好.
所以我參與在當中.
而是我有多麼需要上帝.
在這個群體裡面.
問題荒謬.
其實每個人都承載著一部份.
但是我們共同最重要的核心.
不是這些問題.
我們共同需要的是.
我們每一個人都需要上帝.
在這個群體裡面.
我們就一個滿佈問題的群體當中.
一起仰望尋找上帝.
所以弟兄姊妹.
我特別向著網上的你.
不要因為教會的問題.

$^{1001}$不要因為你生命裡面.
曾經遇到的一些挫折.
有些嘴臉,有些傷害.
一步一步地抽離了自己.
甚至連信仰都放棄了.
希伯來書的作者說不要.
因為你會熟齡倒退一步一步.
去到最後你連信仰都會失去.
失去信仰的結果是甚麼.
他用了一些極端的字眼.
就是在說審判的火.
即是這種踐踏.
或者這種氣絕.
在舊約裡面說到氣絕.
連聽律法都難逃避死亡.
如果去到今天你氣絕基督耶穌的恩典.
離開了基督之恩.
人生命就走向一個沉淪.
希伯來書的作者某程度是嚇一嚇我們.
不過他說的話對我們是非常當頭棒喝.
提醒我們不要輕言放棄這個信仰.
我們投入在一個教會生活.
以及在我們未來的日子.
如何去保持這個信仰.
兩者是有關係的.
不要想著我不回教會.
我還可以很信上帝.
在希伯來書的作者裡面他不是這樣看的.
即使地上的教會滿佈問題.
不過我們共同的需要是我們需要上帝.
當然教會進步一路一路會更加理想.
不過不要將你的焦點放在人的問題上.
你會反而有力量.
菲利普·燕斯這樣提醒我們.
希伯來書的作者同樣這樣提醒我們.
我們再繼續看32節至到第36節.
你們要追念往日.
蒙了光照以後所忍受大爭戰的各樣苦難.
我在這裡跳了第35節.
所以你們不可丟棄勇敢的心.

$^{1041}$全這樣的心必得大賞賜.
你們必須忍耐.
使你們行完了神的旨意.
就可以得著所應許的.
32至到36節說的是什麼呢.
如果在他的敘述裡面.
剛才說了罵一罵你.
警告一下你.
去到這裡他又再勉勵你.
正面去看.
他如何勉勵呢.
他說其實希伯來書的讀者.
一班基督徒.
他本身的光景是一班什麼人呢.
你想起舊時.
追念往日.
蒙了光照的時候.
即是被上帝的光光照了他的心靈.
這班人原來是滿有力量的.
他們可以忍受大爭戰的各樣苦難.
很多的困難.
他都可以承受的.
被毀榜遭患難.
好像一場戲一樣被人看了.
又陪伴那些苦難的人.
連續被困所的人.
甚至乎有些家業.
那層樓那些家都被人搶去.
都可以甘心忍受.
在這裡面.
希伯來書的作者說.
他對著勉勵的那班頂子母是一班什麼人呢.
不是省油的燈來的.
他在說.
昔日你們這班人是一班什麼人來的.
你們是一班大有勇敢可以忍受各樣苦難的基督徒.
為什麼到了今時今刻你們放軟手腳.
沒有力量.
完全失去鬥志.
你們本質不是這樣的.

$^{1081}$希伯來書的作者說.
這班人你們要起來.
舊時的你.
你剛剛信主.
或者昔日的你.
是一個面對困難.
險阻.
兵來將擋.
水來土掩的一班勇敢的基督徒.
是一班忍受堅守的基督徒.
所以希伯來書的作者不是憑空捏造.
或者將一些難擔的重擔.
交給那班讀者身上.
告訴他們.
你學一下某個吧.
不是.
是這些事具體發生在這班人身上.
只不過這班人沒有了一副辦法.
同樣道理頂子妹.
我相信我們眼前一班的頂子妹.
鏡頭前的你.
我們每一個.
生命總有亮麗的時刻.
我們生命不是一直都是這樣的狀態.
我們有輝煌的日子.
那些故事是發生在我們身上.
我們可以挺起來.
我們可以為上帝的緣故.
我們可以勇敢地見證福音.
領人歸處.
讓人生命得著改變.
我們有牧羊的能力.
我們每一個都可以的.
可能因時制宇.
或者受到傷害.
我們就好像被打沉一樣.
但是如今.
我們要重新起來.
所以希伯來書的作者說.
你們不可以丟棄勇敢的心.

$^{1121}$不要將舊時的鬥志丟掉.
那樣東西在你的生命裡一直存留.
拿回那個鬥志回來.
將這份鬥志重新喚現在你的生命裡.
當你保持著它的時候.
到那一天.
耶穌回來的日子.
你會得著最大的賞賜.
忍耐吧.
有很多困難.
有很多荒謬.
一笑之之.
當笑話聽就算了.
不要因為這樣.
因為一句說話.
一些問題.
就引起了你的信仰.
當你忍耐的時候.
當你再拉遠一點去看.
那些時其實算不得什麼.
相對於永生.
相對於你生命的賞賜.
眼前的一句說話.
或者眼前做的一些很愚蠢.
很荒謬的事.
我比較粗俗一點.
不好意思.
算不得什麼.
比對一下那個價值.
我們知道我們追求這個.
我們就努力不懈.
來到去堅心.
來到去實踐.
好,我們時間關係.
我快速去看了.
37至11章一節.
他這裡.
他講到.
剛才說到那個信忘愛.
跟著他就再講.

$^{1161}$我們要去拿回那個竇字之後.
我們更加要用信心.
來到去回應上帝.
37節.
還有一點點的時候.
那要來的就要來了.
並不遲疑.
只是義人必因信得生.
他若退後.
我心裡就不喜歡他.
我們卻不是退後.
入沉淪的那等人.
乃是有信心.
以致靈魂得救的人.
就是這裡.
希伯來書再講.
我們是那班有信心.
不會等候沉淪的那班人.
我們將來會等候.
那個冠冕的一班.
上帝所視為義人的一班人.
基督將要來了.
我們就帶著這份信心.
等待他.
跟著他再講.
信是什麼意思呢?.
信就是所望之時的實體.
是未見之時的確具.
古人在這信上得了美好的證據.
我們因著信.
就知道世界是直神的話造成的.
這樣所看見的.
並不是從顯然之物造出來的.
在這裡說的是什麼呢?.
三個重點.
如果你是信.
你知道我不是那些沉淪的人.
我是一個有信心的人.
你做好這三件事.
不是你憑感覺.

$^{1201}$是真的按照聖經教導來做.
第一件事.
你的信心就是代表你的堅持.
你的信不是一時信一下.
一時又信下其他東西.
左右右擺.
我是持續地相信.
堅守著你的信仰.
無論在教會.
在生活.
在家庭.
在鄰舍.
無論什麼地方.
我都是帶著這一份確信.
來堅持走我信仰的路.
堅持是希伯來書裡面一個很重要的字眼.
失去了堅持.
等於失去你的信心.
所以在你的生活裡.
無論怎樣也好.
無論有多少困難.
你都要堅持.
堅持在這個信仰群體裡面.
堅持與上帝親近.
這就是一個很具體.
很落地.
告訴你.
你驗證自己是一個承受信心的人.
第二個.
信是預見.
我用這個預見.
不是聖經在那個處境裡面.
我嘗試給你容易一點理解.
因為他說的那句也不容易理解.
信就是所望之時的實底.
是未見之時的確具.
好像很抽象.
其實他說的意思是.
在你很多時候我們見到一些東西.
今天的科學就是你見到那個.

$^{1241}$驗證了.
那我就相信他.
有實據.
有證據證明的.
那我就相信.
不過在基督教信仰裡面.
有很多東西不是憑眼見可以見到的.
你即時的眼睛是看不到的.
不過你憑著信心的眼睛.
你可以知道將來的日子.
這一切會實現的.
不是你自己想像一個幻想的圖畫.
是按照上帝在聖經裡面的應許.
上帝怎樣說.
你就用信心的眼睛.
來預見將來是一個這樣的景象.
所以在你生命裡面.
這一種信心的預見就是.
你有一種超越眼前困境的能力.
如果你說你是一個有信心的人.
你會知道.
你會輕看眼前很多淋淋中的問題.
你不會被那些問題所打倒.
因為再多困難也好.
到最後那個美好的日子.
上帝預定的.
一定會來到.
路是彎彎曲曲.
但不代表我們去不到那條路.
帶著信心的眼睛.
你預見上帝為我們所預備的將來.
這就是信心的預見.
最後.
我們見到上帝.
我們知道這個世界.
其實是一種途徑認識神的途徑.
十一章第六節那裡說多一點.
人非有信就不能得神的喜悅.
因為到神面前來的人必須信有神.
且信他賞賜那尋求他的人.

$^{1281}$人能夠認識上帝.
能夠知道上帝.
那個途徑是憑著我的信.
今天說的那個信.
或者認識上帝.
今天坊間有很多東西在騙你的.
其中一樣東西就是感覺.
有些人說你感覺到上帝.
聖經裡面從來沒有說感覺到上帝.
不過你是不是信心.
你是不是憑信心去認上帝.
感覺上帝.
那個感覺其實是可以營造的.
可以製造一個空間.
神聖空間給你讓你感覺上帝.
如果大家有興趣.
我可以做一個神聖空間在這裡.
不過那個不是信仰的養分.
真正凝神的是.
在你生活跌跌撞撞很困難的日子.
在被人逼迫的日子.
在你生命如意的日子.
面對生與死的問題.
面對生命所有很多林林總總.
很多的困難.
我信.
我信上帝在這裡.
我信上帝與我同在.
我有這個信心的時候.
無論多麼困難.
多麼沒有感覺也好.
我知道.
我認識神.
我知道上帝在我生命裡.
與我這個破破爛爛的人.
一起來同行.
好,兄弟姐妹.
今天很長.
不過當你讀這段希伯來書.
當你認真去讀的時候.

$^{1321}$振奮人心.
希伯來書的作者面臨我們.
生命裡面有些東西.
可能令我們停止了.
生命裡面聚會.
或者激發愛心.
用信望愛來回應上帝.
今天我們重新拿起來.
拿起來不是說我們很偉大.
或者覺得自己很高尚.
只不過是因為基督耶穌為我們開了一條新的路.
這條路讓舊約的人恨都恨不到的.
一個可以親近上帝的路.
他為我們打開了這個新的路.
我們就按著他的旨意.
一起來實踐.
一起來用信望愛來回應上帝.
兄弟姐妹.
你預備好自己了沒有?.
你願意將你自己全心全意交給上帝嗎?.
求聖靈繼續對我們說話.
\newpage



\section{路加福音 18:9-14}
\label{sec:l8_t1yVcPHE}
\textbf{2024年3月24日崇拜直播｜余嘉敏姑娘｜十個深化關係的禱告-稅吏的禱告｜路加福音18\_9-14}
\newline
\newline
連結: \href{https://youtube.com/watch?v=l8-t1yVcPHE}{\texttt{https://youtube.com/watch?v=l8-t1yVcPHE}} ~~~~ 語音日期: 2024-03-24
\newline
\newline
\hyperref[sec:3GKCAsbLUbA]{\small{< < < PREV SERMON < < <}}
~
\hyperref[sec:index_chronic]{\small{[返順時目]}}
~
\hyperref[sec:1_gQDpS_bXo]{\small{> > > NEXT SERMON > > >}}
\newline
\newline
路加福音 18:9-14
\newline
\begin{longtable}{cl}
\hline
\hline
章節 & 經文 (和合本修訂版)\\
\hline
18:9 & \begin{tabularx}{0.7\textwidth}{X} 耶穌向那些自以為義而藐視別人的人講了這比喻： \end{tabularx} \\ \\ \relax
18:10 & \begin{tabularx}{0.7\textwidth}{X} 「有兩個人上聖殿去禱告，一個是法利賽人，一個是稅吏。 \end{tabularx} \\ \\ \relax
18:11 & \begin{tabularx}{0.7\textwidth}{X} 法利賽人獨自站著，自言自語地禱告說：『神啊，我感謝你，我不像別人勒索、不義、姦淫，也不像這個稅吏。 \end{tabularx} \\ \\ \relax
18:12 & \begin{tabularx}{0.7\textwidth}{X} 我每週禁食兩次，凡我所得的都獻上十分之一。』 \end{tabularx} \\ \\ \relax
18:13 & \begin{tabularx}{0.7\textwidth}{X} 那稅吏遠遠地站著，連舉目望天也不敢，只捶著胸，說：『神啊，開恩可憐我這個罪人！』 \end{tabularx} \\ \\ \relax
18:14 & \begin{tabularx}{0.7\textwidth}{X} 我告訴你們，這人回家去比那人倒算為義了。因為凡自高的，必降為卑；自甘卑微的，必升為高。」 \end{tabularx} \\ \\
[1ex]
\hline
\hline
\end{longtable}
$^{1}$弟兄姊妹,讓我們以聆聽聖道繼續敬拜.
今日我們選讀的經文是.
《新約聖經》路加福音第十八章九至十四節.
在教會聖經裡的第110頁.
或者在程序表裡也印製了出來.
請聽我讀出路加福音十八章第九節.
「耶穌向那些仗著自己是義人,藐視別人的」.
設一個比喻,說:.
「有兩個人上殿裡去禱告,一個是法利賽人,一個是瑞利.」.
「法利賽人站著自言自語地禱告說:.
也不仗這個瑞利,我是一個禮拜禁食兩次,.
凡我所得的都捐上十分之一.」.
「那瑞利遠遠地站著,連舉目望天也不敢,.
只垂著胸說:神啊,開恩可憐我這個罪人.」.
「我告訴你們,這人回家去,被那人倒算為異了,.
因為凡自高的必降為卑,自卑的必升為高.」.
各位弟兄姊妹早安.
早安.
我和老公拍拖時,他帶我去看一套電影.
是一套完全視覺沉浸式的朱古力電影.
不知道大家有沒有看過,叫做《朱古力掌門人》.
對我來說,大家認識我,應該知道我不喜歡朱古力.
所以對我來說是一套非常難忘的電影.
故事講述一個名叫Charlie的出身非常貧窮的男子.
他住在一個出產Wonka牌朱古力的工廠旁邊.
這個朱古力的牌子非常暢銷.
很多世界各地的小朋友都很喜歡這個朱古力.
Charlie每天都聞著濃郁的朱古力味.
夢想有一天可以進入這個朱古力工廠參觀.
有一天機會來了.
老闆Willie Wonka推出了一個Wonka的金獎活動.
把五張金獎券放進朱古力裡.
只有五張.
所以只有五個人和他的家人一個人可以參觀這個朱古力工廠.
這個活動可以想像.
會掀起很大的熱潮,搶購朱古力.
這個貧窮的Charlie很辛苦才有錢買了一排朱古力.
結果真的被他抽中了.
其餘四個分別是有家底,亦是有地位的小朋友.
一個是貪心的,一個是有公主病的.

$^{41}$一個是自以為是,看不起別人.
另一個是好好勝,心目中無人.
這四個小朋友加上Charlie一起參觀朱古力天堂.
我們一邊看電影就會看到.
他們因為貪心和自以為是,結果被淘汰.
一個一個逐出工廠.
唯有一個出身貧窮,謙卑良善的Charlie可以留下.
最後我們看到他更成為了這間朱古力天堂的承繼人.
這套電影令我聯想起今天耶穌所說的.
有關禱告的比喻.
主角提到有兩種人的代表.
分別是法利賽人和瑞利.
他們同一時間進入了聖殿祈禱.
但有兩種不同的結果.
我們一起從他們的禱告.
反思一下我們的屬靈生命的狀況.
以及學習如何深化我們與上帝的關係.
我們一起看禱告.
親愛的巴庫上帝.
多謝你留下寶貴的話語給我們.
教導我們如何做你喜愛的兒女.
求聖靈你光照我們的內心.
反照我們裡面的屬靈的光景.
教導我們如何謙卑地馴服在你的心意裡聆聽你的話語.
願意聽完之後我們願意行道.
願我們聽完這個訊息之後.
我們的生命被你更新改變.
求主你幫助我們.
祈禱你奉我主耶穌基督的聖名祈求.
阿們.
剛才提到禱告是一個比喻.
比喻通常是耶穌用來.
在新約時經常有的教導.
祂很喜歡用地上眾所周知的例子.
在真實的事或人做對比.
幫助我們理解一些比較深和複雜的屬靈真理.
而且會激發我們去多點想像力.
我們可以投入去那個場景.
想想我是誰.
如果我在那個境地.

$^{81}$我會怎樣去應用這個屬靈真理呢.
而這種教導其實很鼓勵我們.
主動去思考和反思.
更加深入應用這個真理.
耶穌的比喻通常都會說給特定的聽眾聽.
但這次究竟耶穌向著什麼人說呢.
而祂的反思又給我們什麼屬靈真理去想呢.
我們一起看看今天第九節的經文.
這裡說到「將著自己示意人」.
即是藐視別人設一個比喻.
「將著自己示意人」.
原文可以這樣說.
「信靠自己為義」.
耶穌的比喻對象是.
「信靠自己多於神,自以為義」.
認為自己做的事很合神心意.
而且更加藐視.
看不起別人的一些人.
當我們讀到這裡的時候.
我們很自然就會聯想起.
他一定是說我們下面會繼續說的法利賽人.
因為耶穌經常點名斥責他們.
說他們假冒為善.
好像粉色的墳墓一樣.
外面很漂亮.
但裡面都是裝滿死人的骨頭.
一樣的污穢.
但第九節我們看到.
其實他沒有點名指明那是法利賽人.
只是向他們說這個比喻.
但他反而是向一些自以為義的人.
說比喻給他們聽.
我想也可能包括今天我們在閱讀經文的我們.
任何一個自以為事.
自以為自己做得比別人好.
也看不起別人的人.
如果我們翻查原文聖經.
我們會看到這段短短的經文裡面.
有一個連接詞.
如果用中文去翻譯.

$^{121}$就是和或者也.
在英文就清楚看到.
King James Version.
開頭的時候會看到N字.
這裡說的是什麼意思呢?.
其實意味著這段經文.
和上文有一個連接的意思.
我們又想你們翻查聖經.
我們一起翻查前一點的第十七章.
第十七章二十至三十七節.
其實耶穌在教導門徒.
有關天國即將降臨.
人子會再回來的情況是怎樣的.
第十八章進入對耶穌說了兩個有關禱告的比喻.
第一個是第十八章的一至八節.
耶穌說了一個寡婦和不義的法官的比喻.
承接上面說天國的訊息.
耶穌想提醒一群信徒.
在等候主再來的時候.
你們要常常禱告,不可以灰心.
第二個比喻就是我們今天的經文.
就是九至十四節.
耶穌用了當時猶太人很熟悉的兩種代表人物.
教導信徒要持有什麼樣的態度.
去祈禱等候主再回來.
這裡我們先說第一個代表人物.
法利賽人,我們看第十節.
第十節說有兩個人一起去聖殿禱告.
一個是法利賽人,一個是瑞利.
在我們認知當中,法利賽人是一個怎樣的人.
他們是出名自以為是,都是有點壞.
假冒為善.
不過如果我們代入當時的社會.
原來當時的人會覺得法利賽人虔誠.
去嚴守律法的一些信徒.
所以他們的地位是很受尊崇.
只不過他們過度重視外在的一些律法規條.
忽略了內心的敬虔心.
而另一位上去祈禱的就是瑞利.
瑞利當時的社會裡的人很憎恨他們.

$^{161}$為什麼呢?因為他們效力羅馬政府.
他們會幫羅馬政府向猶太人收取稅.
當時羅馬政府是挺聰明的.
他們不直接向巴勒斯坦地的人民收取稅項.
他們會將收稅的專利權拍賣.
成功拍賣的稅利.
用我們今天的方式去理解.
其實這班是商人.
他們就承包了政府收稅的生意.
但羅馬政府就不理會他們收多少錢.
總之他們每年將特定的稅利金額交給羅馬政府.
你們可以想像.
如果你是瑞利人.
用我們基督徒的方式.
你會怎樣呢?.
你都會收多.
收多就由得代落代.
為自己去致富.
所以在猶太人眼中.
不受歡迎,罪大惡極的稅利.
另一個就是備受尊崇,嚴守律法的法利賽人.
頓時形成了強烈的對比.
這兩個人同時到聖殿祈禱.
但經文形容得很細緻.
你看他們兩個都有不同的禱告姿態.
首先第十一節說法利賽人是站著自言自語地禱告.
這個姿勢並不奇怪.
因為當時的猶太人通常都是站著.
雙手舉起望天禱告.
有趣的是這裡用了自言自語的字.
這個意思不是自言自語.
小小的吟吟地對自己說話.
而是原文直譯這個字.
就是自己對自己說話.
自己向自己禱告.
或者為自己的事禱告.
耶穌似乎刻意地諷刺那群法利賽人.
禱告的焦點好像放錯在自己身上.
你可能會說:不是的.
他們是刻意去聖殿找神.

$^{201}$更加是在聖殿中間.
舉手望天向上帝.
然後開口做了感恩的禱告.
這個焦點還不是對著上帝?.
不過我們留心再看一次他的禱文.
好像如數家珍一樣.
向上帝報告:我做了這樣的事.
不像其他人勒索,不義,姦淫.
更不像旁邊那個.
不像旁邊那個瑞利.
他在做什麼?.
他在靠自己的好行為.
他認為自己的好行為.
這樣的祈禱可以取悅上帝.
這樣可以和上帝再親近一點.
但第十二節他就說到.
他更加指出自己.
我每個禮拜禁食兩次.
凡我所得的都捐上十分之一.
意思是他不只做得對.
他還要做多.
其實在律法要求猶太人每年禁食一次.
就是那個贖罪日禁食.
他竟然是一個禮拜有兩次.
我不會計算,數字超出了不知多少倍.
而奉獻又有十分之一.
其實你可能會說.
不是呀,多好呀,做得對呀.
沒有錯的.
但問題在於他們的心態當中.
他的祈禱充滿自吹自擂.
他認為自己是正直的.
我是正義的,無可指責.
我不像其他人.
甚至貶低他旁邊那個瑞利.
覺得自己比別人優勝.
當人自以為是的時候.
眼角長在頭顱頂.
根本是目中無人.
眼睛又怎能放下來.

$^{241}$直視那塊鏡子.
那塊鏡子反照我們自己的本相.
在座的女士應該會知道.
出街的時候.
如果之前不照鏡.
頭髮亂了,眼線歪了.
你根本不會知道.
樣子醜了.
你去足全日街.
都可以自己騙自己.
夢在古裡,完全不知道.
說回朱古力掌門人.
結尾有一幕,我都很深刻.
就是拍攝那四個小朋友.
跟家人一起走.
被人趕出工廠.
這幕很有趣.
因為是一本故事書.
所以最後一幕.
Charlie和老闆.
坐著透明升降機.
從工廠慢慢走.
由低升到最高.
其中一個小朋友.
他看到這個景況.
很慘地走出工廠.
但他仍未吸取到教訓.
他繼續貪心地跟爸爸.
跟爸爸說.
「你買一部這樣的色」.
「我又要一部」.
爸爸一口拒絕,很生氣地罵他.
原來驕傲仍然蒙蔽著他的雙眼睛.
同樣,法利賽人的驕傲.
蒙蔽了他們的雙眼睛.
看不到他們內在的污穢.
承認不了自己很有限制.
仍然很需要神的憐憫和拯救.
他的禱告正正反映了.
他們屬靈的光景的欠缺.

$^{281}$他認為自己靠自己.
就可以得到夢神的喜悅.
甚至可以得到天堂的金獎.
相反,瑞利的禱告.
原文只有六個字.
簡而精.
已經可以反映到.
他內在的世界的謙卑.
我們一起看看第十三節.
瑞利的禱告.
與法利賽人的禱告相反.
他遠遠地站著.
似乎不敢靠近.
而且他的眼目.
連看向天都不敢.
這一連串的動作.
表明了瑞利覺得自己很不配.
他更一面掩著隨胸.
一面祈禱.
「神啊,開恩可憐我這個罪人」.
我們看到他的禱告,你覺得怎樣?.
我覺得他在表達極度深的懊悔之情.
他直認自己.
承認自己是一個罪人.
他亦沒有很多例子為自己辯護.
甚至沒有合理化所有的罪.
他只是很坦誠地走到上帝面前.
承認自己是一個罪人.
他懇求上帝開恩憐憫.
他的禱告正正反映了.
他在熟齡的光景中.
他很渴求神去寬恕他.
他知道他無法自救.
他救不了自己.
唯有神才能將他脫離罪的綑綁.
相比起法利賽人.
這個瑞利的禱告是真誠很多.
他的禱告表現出很謙卑.
很渴求神跟他親近的關係.
有一個朋友曾經跟我說過.

$^{321}$他服侍一個基督教機構.
出現了一個危機.
這個機構急求同工來開一個緊急會議.
希望大家可以集合大成.
尋找一些可以應對危機的方法.
在會議當中.
其中一個同工突然有感動.
開口跟大家說.
不如我們放下我們的智慧.
我們虛心地向神尋求憐憫和寬恕.
我覺得這個同工很勇敢.
我們試想像.
在公司開會.
突然有同事跟你說.
不如我們不要說那麼多.
我們一起祈禱.
我們一起認罪悔改.
別說認罪悔改.
隨便說我們先祈禱吧.
我們也未必夠膽.
或者我們想像教會群體.
可能會比較容易.
我們一起祈禱.
但我來教會工作了那麼多年.
甚少聽到弟兄姊妹.
或我們同工在開會期間.
說不如我們先為教會或自己.
做一個懊悔的謝罪的祈禱.
我覺得這是一個頗勇敢的行為.
同工的這麼無媚的邀請.
竟然眾人看為美.
於是當日上至總幹事.
下至一個新入職的同工.
他們一起為機構.
為自己的罪在上帝面前祈求赦免.
他們一面分享自己內裡的掙扎.
一面求神幫我們.
我們也不知道如何做決定.
他們更加分享到落淚.
這是一個非常美的畫面.

$^{361}$這個禱告不是突然有神奇功效.
令到他們祈禱後危機不用處理.
他們之後會繼續開會.
但同事之間的心被拉近.
最重要的是他們各人的眼.
由本來看著自己.
覺得可能是在用人的能力.
或人的需要上.
回轉對準上帝.
因為上帝才是領導這個機構.
甚至是我們每一個人的上帝.
我們回想我們.
我們很多時候都做這樣的禱告.
或是為自己或為人祈禱.
但我們很少做認罪悔改的禱告.
今日的經文提醒我們.
原來上帝看重我們謙卑認罪悔改的禱告.
但我們不要搞錯.
上帝不是要我們每一天都過著後悔.
內疚自責或自卑的生活.
因為自卑是一個會蠶食人生命的.
一種負向的心態.
對自己的價值和能力.
覺得自己不足夠甚至很沒信心.
自卑是無法帶來人的一種推動力.
不可以令我們向前.
自卑只會令我們越來越縮.
不敢面對人和面對上帝.
我回想受浸後的光景.
受浸時已經是屬靈的最高峰點.
覺得自己見證上帝.
但過了一段時間後.
我的屬靈生命就變得很低落.
我做了很多對不起上帝的事.
我覺得上帝不會喜悅我.
但我有回教會.
每星期都有回教會.
我坐在位置上.
向上帝求力寬恕我.
我又做錯了,又這樣那樣.

$^{401}$但我這份不是謙卑的生活.
我是內疚,我很自責,很自卑.
我經常想我回來做甚麼.
我回來面對上帝,我沒有面子見他.
於是這份自卑的心.
驅使我離開教會一段較長的時間.
直至我去了外國讀書.
重整我的生命後.
我才回到教會當中.
我相信自卑離開上帝.
不是神想要我們有的結果.
神看重我們是否有謙卑的心.
唯有謙卑是一個有建設性的正向心態.
真正謙卑的人.
不會高估自己的能力和價值.
亦不會像法利賽人那樣.
將自己凌駕在別人身上.
唯有謙卑的人才知道.
如何真誠面對自己.
願意將眼目從額頭放下來.
看自己的本相.
承認自己神是很有主權的.
亦承認自己很有限制.
我們只不過是一班蒙恩的罪人.
唯有這份謙卑的心.
才讓我們能夠謙卑地聽別人的意見.
虛心地信服神的指引和帶領.
弟兄姊妹.
請大家記得一個很重要的事.
神看重我們.
看重你們這份真誠,謙卑,禱告的心.
而不是叫你們自卑.
唯有謙卑的心.
才能夠將我們和上帝的關係.
越拉越近的這條很重要的鎖匙.
我們又看看最後一節第十四節.
耶穌親口宣告謙卑的重要性.
瑞利謙卑地承認自己的罪.
承認自己很不配.
最終反而被神稱為異.

$^{441}$我們看十四節.
我們一起讀吧.
十四節.
「我告訴你們,即梅加許,.
因為凡自高的必降為卑,.
自卑的必升為高.」.
這裡很有趣.
耶穌親口宣認瑞利.
在進入聖殿禱告後.
他回家後.
比起法利賽人倒算為異.
倒算為異是甚麼意思呢?.
意思是他本身不是異人.
但他靠著神的恩典.
神就算他為異人.
反而法利賽人.
在神的眼中.
他不及瑞利稱異.
回想剛才的場景.
我們在聖殿裡面.
法利賽人和瑞利的姿態.
畫面對比今天耶穌所說的局面.
簡直是極大的逆轉.
「凡自高的必降為卑,.
自卑的必升為高.」.
這句不僅出現在今天的經文.
其實在《馬太福音》23章.
《馬可福音》12章.
和《路加福音》前的位置.
其實都分別記載.
耶穌譴責那群文士和法利賽人.
說他們坐在摩西的位置上.
吩咐人嚴守律法.
但他們自己卻全部所做.
都是為了抬高自己.
所以耶穌譴責他們.
同樣說這句說話.
我們見到耶穌很重視.
這個天國的原則.
他務必要多次提醒我們.

$^{481}$謙卑尋求神憐憫的人.
必蒙拯救.
唯有謙卑的心.
才是將我們和神的關係.
越拉越近的那條鎖匙.
所以這就是今天.
上帝留給我們很重要的訊息.
我問問大家.
你內心像法利賽人.
還是你自己像一個瑞利人.
或者我們會這樣想.
我真的不夠膽在神面前.
自吹自擂地說自己多厲害.
我不是像大衛.
大衛真的有這樣向上帝禱告.
但我們會否都像法利賽人一樣.
我們很容易墮入.
自以為是的圈套.
以為我為了教會.
為了上帝做了很多事.
我比其他人更加屬靈.
更加優勝.
我們的禱告會否像法利賽人一樣.
弄錯了焦點.
對不起.
就像唸口往自言自語.
只著重外在的宗教儀式.
而忽略了內在自省的能力.
還是你會像瑞利一樣.
有一個謙卑知罪的心.
能夠坦然地向上帝面前.
承認自己的內在狀態.
認識到自己很軟弱.
有很多罪.
很謙卑地向神面前祈求憐憫和寬恕.
其實謙卑不只是說就有.
剛才我們說了一大段.
你也說是否我要學足.
像瑞利一樣說.
神啊 開恩憐憫我這個罪人.

$^{521}$這樣就代表謙卑.
大家也知道不可能吧.
真實的謙卑其實要操練和實踐.
有一個用了他一生去實踐和操練謙卑的.
勞運神父.
他是一個靈修神學的作者.
大家應該看過他的著作.
有看書的.
他曾經是哈佛大學和耶路大學的資深教授.
他曾經寫過.
他任教的時候.
他寫過一本書叫做.
向下的移動 基督的社稷之路.
當中他提到一個非常重要的.
很逆向的思維.
他提到神的兒子.
他選擇了成為拿撒拉人耶穌.
他走了一條非常謙卑.
向下的路.
去啟示他自己救贖計劃的奧秘.
他從本來掌權的.
走到毫無權力的境地.
他從本來很強壯的.
走到很軟弱.
而從很榮耀的.
走到被羞辱的一條路.
你看到耶穌越走越謙卑.
直至他走上十字架上.
為我們成就那份的救恩.
在我們世人.
向上移動就多了.
我們香港人都很喜歡向上移動.
好像法利賽人.
要有功 要果效.
要引人注目.
但這些從來都不是耶穌基督.
呼召我們成為基督徒的心意.
為了實踐謙卑.
盧雲神父在1985年.
決定離開哈佛大學.

$^{561}$告別向上游的試探.
他去了方舟的群體.
香港都有方舟之家.
1986年的時候.
他正式加入.
加拿大多倫多的黎明之家.
在一個機構裡.
服侍智障的群體.
直至他離開這個世界.
其實盧雲.
我不太認識他.
但從他的文字.
我都覺得他是一位書生.
他不是善於做家務.
他也不懂煮飯.
做這些體力勞動的工作.
但他在黎明之家.
第一個任務.
他要和一個嚴重智障.
殘障的Adam.
一起住.
負責照顧他.
幫他刷牙.
幫他洗澡.
幫他上廁所.
處理大小異變.
有時他因為疼痛不方便.
要經常抱他出入.
這些都是甚麼工作?.
很卑微吧?.
但盧雲.
在他的著作裡說.
他從Adam身上.
發現自己很多不足.
他從此學會更謙卑.
他更在Adam身上.
經歷了很深很深的神愛.
弟兄姊妹.
你們是否渴望像盧雲.
學懂謙卑.

$^{601}$經歷神那麼深的愛?.
我想邀請大家.
今天回家或在教會.
找一個空間.
在上帝面前安靜禱告.
首先求上帝光照你的內心.
承認真誠地面對自己.
反思並承認自己過往這段日子.
我的生命裡有沒有地方.
得罪了上帝.
然後求上帝幫助你.
如何成為謙卑的人.
謙卑地學習聆聽神的心意.
並且信服祂的心意和帶領.
第三,我相信大家需要一點時間消化.
歡迎上帝去管教自己.
尋找向下移動的服事機會.
服事一些我不擅長.
亦覺得很難的事和人.
因為我覺得.
令人感到可以學懂謙卑.
往往要將人經歷的.
自己引以為傲的事.
讓神去遙控和挑戰.
人才會發現自己很渺小.
我們很需要上帝的幫助.
神今天提醒我們.
在主再回來之前.
我們要常常禱告.
不要灰心.
祂更加喜悅我們這一夥.
謙卑真誠的禱告.
讓我們用以下的一首回應詩歌.
「讓我謙卑屈膝禱告」.
向神謙卑地尋求寬恕.
\newpage



\section{約翰一書 3:13-23}
\label{sec:1_gQDpS_bXo}
\textbf{2021年5月9日母親節崇拜直播｜陳嘉璐醫生｜愛在家中｜約翰一書3\_13-23}
\newline
\newline
連結: \href{https://youtube.com/watch?v=1-gQDpS_bXo}{\texttt{https://youtube.com/watch?v=1-gQDpS\_bXo}} ~~~~ 語音日期: 2024-03-27
\newline
\newline
\hyperref[sec:l8_t1yVcPHE]{\small{< < < PREV SERMON < < <}}
~
\hyperref[sec:index_chronic]{\small{[返順時目]}}
~
\hyperref[sec:7e2_k14U9kk]{\small{> > > NEXT SERMON > > >}}
\newline
\newline
約翰一書 3:13-23
\newline
\begin{longtable}{cl}
\hline
\hline
章節 & 經文 (和合本修訂版)\\
\hline
3:13 & \begin{tabularx}{0.7\textwidth}{X} 弟兄們，世人若恨你們，不要驚訝。 \end{tabularx} \\ \\ \relax
3:14 & \begin{tabularx}{0.7\textwidth}{X} 我們知道，我們已經出死入生了，因為我們愛弟兄。沒有愛心的，仍住在死中。 \end{tabularx} \\ \\ \relax
3:15 & \begin{tabularx}{0.7\textwidth}{X} 凡恨自己弟兄的，就是殺人的；你們知道，凡殺人的，沒有永生住在他裡面。 \end{tabularx} \\ \\ \relax
3:16 & \begin{tabularx}{0.7\textwidth}{X} 基督為我們捨命，我們從此就知道何為愛；我們也當為弟兄捨命。 \end{tabularx} \\ \\ \relax
3:17 & \begin{tabularx}{0.7\textwidth}{X} 凡有世上財物的，看見弟兄缺乏，卻關閉了惻隱的心，神的愛怎能住在他裡面呢？ \end{tabularx} \\ \\ \relax
3:18 & \begin{tabularx}{0.7\textwidth}{X} 孩子們哪，我們相愛，不要只在言語或舌頭上，總要以行為和真誠表現出來。 \end{tabularx} \\ \\ \relax
3:19 & \begin{tabularx}{0.7\textwidth}{X} 從這一點，我們會知道，我們是出於真理的，並且我們在神面前可以安心， \end{tabularx} \\ \\ \relax
3:20 & \begin{tabularx}{0.7\textwidth}{X} 即使我們的心責備自己，神比我們的心大，他知道一切。 \end{tabularx} \\ \\ \relax
3:21 & \begin{tabularx}{0.7\textwidth}{X} 親愛的，我們的心若不責備我們，在神面前就可以坦然無懼了。 \end{tabularx} \\ \\ \relax
3:22 & \begin{tabularx}{0.7\textwidth}{X} 我們一切所求的，就從他得著，因為我們遵守他的命令，行他所喜悅的事。 \end{tabularx} \\ \\ \relax
3:23 & \begin{tabularx}{0.7\textwidth}{X} 神的命令就是：我們要信他兒子耶穌基督的名，並且照他所賜給我們的命令彼此相愛。 \end{tabularx} \\ \\
[1ex]
\hline
\hline
\end{longtable}
$^{1}$大家可以回到程序表或教會聖經新約339-342.
約翰一書三章十三至二十三節.
弟兄們,世人若恨你們,不要以為稀奇.
我們因為愛弟兄,就曉得已經出死入生了.
沒有愛心的,仍住在死裡.
凡恨他弟兄的,就是殺人的.
你們曉得凡殺人的,沒有永生存在他裡面.
主為我們捨命,我們從此就知道何謂愛.
我們也當為弟兄捨命.
凡有世上財物的,看見弟兄窮乏.
卻塞住憐憫的心,愛神的心怎能存在他裡面呢?.
少先們,我們相愛,不要只在言語和舌頭上.
總要在行為和誠實上.
從此就知道我們是屬真理的.
並且我們的心在神面前可以安穩.
我們的心若責備我們,神比我們的心大.
一切事沒有不知道的.
親愛的弟兄們,我們的心若不責備我們.
就可以向神坦言無懼了.
並且我們一切所求的,就從他得著.
因為我們遵守他的命令,行他所喜悅的事.
神的命令就是叫我們信他兒子耶穌基督的名.
且照他所賜給我們的命令,彼此相愛.
聖經獨筆.
(祈禱).
願我們繼續用神的話言去敬拜神.
今天我們正在這裡,因為我們正在這裡有陳家洛醫生.
在我們這裡,為我們正在說神的話.
陳醫生可能有部分弟兄們都認識他.
他過往這一兩年都有試過在我們當中不同的場合.
和我們分享和說神的話.
陳醫生是香港註冊的精神科醫生.
但他也是長道衛理九龍塘的義務教師.
我們將時間交給陳醫生.
陳:親愛的弟兄姊妹你們好.
母親節快樂.
不單止做母親的快樂.
我們每一位有母親的,我們都是快樂.
其實我是說一節的.
其實是說一節的.

$^{41}$第十八節.
孩子們,我們相愛,不要只在言語或舌頭上.
總要在行為和誠實上.
其實我們神造,在家庭裡要和家人相處.
最容易相處就是一開口就說.
所以言語是非常容易表達.
不過要面對面說,要有眼神接觸.
不過有些人不是這樣,很間接地說.
不喜歡直接跟別人說.
你幫我說吧.
有一家人三兒子同桌吃飯.
母親問爸爸.
你問他,兒子叫他.
你問他喝不喝湯.
爸爸就問兒子.
你媽問你喝不喝湯.
兒子就告訴爸爸.
你告訴你老婆.
她兒子不喝湯.
這些,這樣行不行?.
完全沒有關係.
為什麼不直接說?.
我們可以面對面的,很親切.
不只說法,還有一些聲氣.
例如大小聲,高音.
沒有尾音,很膚淺.
有些拉長尾音.
例如,都說不去了.
說了很多次.
加重語氣.
行啦.
表現得很厭煩和討厭.
我有個小朋友就說.
陳醫生,我是當口頭禪.
我跟同學說.
有沒有搞錯?.
為什麼你這樣說?.
他們誤會,我不是覺得不對.
不過他們很生氣.
他們覺得我在說不好.

$^{81}$於是他們開了個戶口.
來罵我.
這樣的說話方法.
是否邀請人來罵你?.
所以如果你覺得有不好.
我們就反省一下.
停一停,想一想.
有些事要控制.
很多事我們都要過濾一下.
我最喜歡說的控制情緒.
就是在電影院.
有一齣電影很緊張.
很緊張的偵探片.
有個女士按捺不住.
突然起來.
究竟誰是兇手?.
後面有一把冷冷的男人聲說.
你現在不立刻坐下來.
你後面就有一個.
兩個都控制不了情緒.
因為有時我們控制不了情緒.
會罵人,對不對?.
說夠了沒有?.
你閉嘴.
就是頂撞她.
命令式,你閉嘴.
又罵人.
很討人厭.
尖叫.
真的語言傷害人很厲害.
但不只是說話.
不只是語言的洩透.
還有嘴唇都可以殺人.
我有個朋友說.
老闆收了我的東西.
很不高興.
接聲.
又將我的檔案丟到桌上.
然後鼻子就.
嘩.

$^{121}$嘴唇又可以殺人.
鼻息又可以殺人.
過濾一下.
如果我們過分會傷害人.
我們就不要做了.
有些不好的情緒.
叫做Delete.
內容是.
有個朋友說.
我去吃了些雪糕.
告訴我妹妹.
妹妹說.
你常說減肥.
拜託你不要吃雪糕.
那些雪糕卷.
是我媳婦送給我的.
她不知道.
我就很委屈.
人家送東西給我.
我怎可以不給呢.
我怎可以不吃呢.
她真的不明白.
而且我只吃一次.
又不是經常吃.
她也不了解.
她說話真的很不濟.
她不知道聽的人很辛苦.
又委屈.
但我說.
你不用聽那些令你刺傷的那一段.
你聽到妹妹對你的關懷.
她又想你不要這麼胖.
你聽到她的關懷.
然後你不要介意她說話這麼不濟.
那個寬容也在內.
所以懂得說.
也要懂得聽.
有聆聽的耳.
是積極的聆聽.
聽到好東西.

$^{161}$不好的東西.
就聽不到.
就叫做Selective Listening.
選擇性的聆聽.
不聽到.
但有些人真的傷害人.
有個媽媽說.
給那麼多錢你讀不到書.
浪費我的錢.
埋怨.
責任全都交給你.
這麼慘.
所以要傳遞正面的訊息.
欣賞人.
愛人.
你見到今天很好.
大家都說好聽的話.
如果很久沒見到某個家人.
你說很掛念你.
你好嘛.
很疼你.
祝福你.
多正面.
都是一句.
對不對.
你說掛念他.
你就打破隔膜.
你就表達你的慣念.
以及關懷.
人家說醫生是醫病.
家人是醫命.
說話說得對.
是很好聽的.
不需要說那麼責備.
罵人.
過濾一下.
好聽的就說.
不好聽的就不說.
這也是Selective.
Selective這個字很好用.

$^{201}$選擇性的說.
選擇性不好聽就收起來.
扔掉.
對不對.
好聽就說.
謝謝.
不好聽就說.
聽不到.
很厲害.
選擇性.
這些負面的想法.
想得自己頭都爆了.
非常憂慮.
就要停下來.
不要想了.
這些不好.
我剛才說負面的情緒.
全部都是情緒垃圾.
怎麼處理.
扔掉.
垃圾是不好的.
影響我們心境的平靜.
破壞我們的免疫系統.
所以要平靜.
就要扔掉.
有朋友說.
陳醫生多久扔一次垃圾.
一天一次好嗎.
我說好啊.
恭喜你.
你做了一整天的垃圾桶.
一有情緒.
馬上扔掉.
一有激氣.
不開心.
憂悶.
煩.
覺得很慘.
我們一點都不慘.
我們就轉一個正面的情緒.

$^{241}$馬上扔掉負面的情緒.
說些好聽的話.
你知道我這個身形.
當然要減肥.
吃什麼特別餐.
現在有一種新的奶粉.
很好.
我也想試試.
又要做什麼運動.
女兒很討厭.
媽媽.
你這麼煩.
你身材也很適中.
說得對.
距離好身材千分之遠.
不過適中.
什麼叫適中.
沒有標準.
你說適中就是適中.
有時我走出街頭可以看起來.
因為我的身材適中.
有自信.
所以一句話.
聖經說得合而.
好像金蘋果吊在銀網子裡.
然後人心憂悶.
屈而不伸.
一句良言.
洗心.
歡樂.
然後又說.
言語柔和.
洗腦消退.
言語暴戾.
觸動怒氣.
因為言語的力量是多大.
不單止可以毀壞人.
毀壞自己.
也可以建立人.
建立自己.

$^{281}$適當的用.
所以選擇性.
說些好聽的話.
不單止說方法.
說內容.
然後真的.
不只是嘴唇上這麼膚淺.
我們也要說.
深入一點.
因為在家裡.
我們教導我們的兒女.
母親教導的兒女.
是憑著主的教訓和警戒.
教導他們.
要小心.
因為現在資訊科技太發達.
馬可福音.
八章二十二至二十五節.
記錄一個.
醫好.
百菜大的盲人.
為什麼很特別呢.
因為是Two-Stage Healing.
要兩個步驟.
馬可福音八章二十二節.
他們來到百菜大.
有一個人帶盲人來求耶穌摸他.
耶穌拉著盲人的手.
帶他出村外.
在他眼睛上套唾沫.
為他按手問他.
你看見什麼.
他抬頭一看說.
我看見人.
他們好像樹木並且行走.
隨後耶穌又按手.
二十五節.
第二次.
按手在他眼睛上.
他定睛一看就復原了.

$^{321}$樣樣都看得清楚了.
有些解經家.
全部耶穌都是一次過醫人.
這個兩個階段.
因為不夠力.
分兩步.
好像打針要兩針.
可能還要三針.
不夠力才分段.
神來的.
怎會需要分段.
他不需要摸.
不需要弄唾沫.
他遙控.
說一句話.
你的信心夠了.
你都不是針對病.
記不記得.
白夫長的僕人已經好了.
用不用這麼複雜.
原來就是在這個.
資訊科技這麼發達的時代.
媒體發達.
衝擊之下.
神就有個啟示.
其實.
那個人第一次.
你看見什麼.
看見人好像樹木又行走.
你說什麼.
人又好像樹又行走.
不知道自己說什麼.
原來.
以為看見.
和清楚看見的實情.
是兩回事.
所以耶穌.
第二次再努力.
再按手.
那個人就看得清楚.

$^{361}$因為這個時代.
資訊可以又剪接.
又換.
又駁.
又加又減.
配音.
配樂.
改頭換面.
舊時的東西又翻做現在.
國家又調轉.
亂了.
偏頗的報道.
惡搞.
轉移視線.
都不知道信不信好.
所以.
我們更加要求神.
真的讓我們看清楚.
耶穌再按手.
耶穌.
按手就是帶著神的力量和智慧.
加上他自己定睛一看.
他定睛一看.
他努力看一看.
不只是神賜給我們力量.
我們自己也要做些事.
他努力一看.
就復原了.
樣樣都看得清楚.
在這個世代.
我們要教我們的兒女.
有辨別是非的智慧.
這個智慧.
這個律從小放在心裡.
就是神的話語.
只有一個.
不止現在這樣的惡搞.
有時在網上.
你跟不跟這個投資.
你信不信這個人是甚麼公安.

$^{401}$我說.
電光火石之間要作出決定.
這個智慧這個能力.
就是視乎你心中有沒有一個律.
這個最好的律就是神的話.
跟著神的聖經話語.
從小明白聖經.
這樣就行.
所以不怕我們求神給我們的智慧.
他就會醫治我們.
這個唯一的途徑是.
神再按手.
神帶著他的力量.
我們努力用神的眼光角度看.
由慢到模糊到清晰.
由演繹錯覺到正確的演繹.
是要得到神的力量.
我就弄了一句.
不少錢剩.
陳醫生這麼使喚.
不是.
因為有很多資訊.
你不喜歡看.
你知現在的AI.
你最喜歡看的.
他就努力傳給你.
你就會洗腦.
你說這個不合原則.
我不看.
但不行.
你不看就變成隱蔽.
於是你就少看.
少看.
不可以與世隔絕.
少看一點.
不少.
淺看.
這個不是錢.
淺看.
不要看得太深入.

$^{441}$不要受太大影響.
然後正.
就是正看.
看得正一點.
真一點.
闊一點.
遠一點.
高一點.
大一點.
你看到真相的成數就會高一點.
所以我們教導兒女.
是要有神的智慧.
求神給智慧.
然後心的律是勤讀聖經.
沒有一個指標.
是準確過神的話語自己.
然後現在的小朋友.
很高傲遠.
不勞而獲.
喜歡做網紅.
做YouTuber.
以前叫網紅.
KOL.
KOL不是不好.
以為任何人都做得到.
但我們要警告他.
KOL網紅一定要有幾點.
第一.
一定要有真材實料.
不是到處抄襲.
你要博學.
要有好的學識做基礎.
有這些學識經過事實檢查.
我剛才說.
你都不知道是不是真的.
亂說話.
害人無數.
不可以害人害己.
深信錯的東西.
經過事實檢查才好出街.

$^{481}$第三.
你要強健自己.
因為網上欺凌是很厲害.
不是人人稱讚你.
每個人都只是幾個手指頭.
有很多人惡毒地罵你.
就算你好好都可以.
所以他要有強健的心理.
頂得住網上欺凌.
不介意.
告訴他.
別人說什麼就是什麼.
不是的.
嘴巴在別人那裡.
你控制不了.
你堅定你的原則.
你的信仰.
你做對的事.
就行了.
負在暗中擦汗.
會幫助你.
會報答你.
不要介意這些欺凌.
因為有些人就倒下了.
很慘.
覺得自己很沒用.
第四.
要注重安全.
你去拍片.
為了要博like.
出位.
很危險的動作都做.
在懸崖上.
做一個怪招來自拍.
所以要很小心.
安全是非常重要.
不要亂來.
第五.
注重私隱.
不可以揭別人的私隱遺錄.

$^{521}$然後尊重別人的版權.
這些要教的.
因為現在的小朋友.
你不知道他們做了什麼.
很厲害.
他們的網球技術是非常好的.
他們做得很快.
就出街.
就糟糕了.
我們要強化他們的心靈.
然後就是.
言教.
第二.
一定是行教.
行教是很重要的.
中國人說.
上樑不正.
下樑歪.
你要榜樣影響他.
也有一句話.
其身正不令而行.
其身不正.
須令不行.
你經常叫子女不要打手機.
怎知道他爸爸機不離手.
你怎樣叫兒子停手呢?.
所以行為是很重要的.
我們的行教.
是顯出我們的好德行.
為他們做很多事.
但是.
也要有好效果.
你知道.
現在喜歡斷捨離.
收拾東西.
兒子就投訴.
陳醫生.
我媽媽經常幫我收拾東西.
一收拾完我的東西就不見了.
找不到.

$^{561}$要有效.
對不對?.
不是亂來的.
然後.
我們兩姐弟都外出吃飯.
回來.
弟弟一定有送樓.
我也有飯.
要公平.
你的行為.
要公道.
然後.
你關心他.
不斷打電話給他.
女兒說陳醫生.
媽媽不停打電話給他.
我在開會.
怎樣聽?.
不能聽.
最後一聽電話.
他就罵我.
為什麼不聽電話?.
你不可以惹人反感.
你不可以惹兒女的氣.
所以.
恰當.
好榜樣就是.
你又叫子女不要喝酒.
不要吸煙.
有個朋友說.
我媽媽煮飯都要喝威士忌.
廚房都要燒著.
你說怎樣?.
怎樣做新行教的榜樣?.
我有個朋友.
說舊時的事.
很刺激.
我的小學老師.
還說記得我.
經過30年.

$^{601}$因為他說.
我記得你.
你媽來問我拿兩分.
為什麼你媽拿兩分?.
他說我二三四年班.
都是第三名內.
誰知五年班.
就跌落五名.
只差兩分.
因為第四名都有獎.
第五名就沒有.
所以媽媽幫我爭取.
他說.
我都不知多慘.
我扯著他叫他不要.
我覺得老師給對分.
最後有沒有加分給你?.
沒有.
但老師記得我.
媽媽還說.
我一定要據理力爭.
這個你不是真你.
那個有少少偏頗.
幫了女兒.
不要緊.
顯示.
你又要欣賞你媽媽.
那麼肯為你出頭.
她真的很疼你.
不介意.
她都幫你.
你欣賞她.
你真的要更加愛她.
因為她這些行為.
以為令我很尷尬.
不是.
你又要欣賞她行為背後.
動機原來那麼愛你.
好.
小心行為的時候.

$^{641}$不要惹怒她的氣.
但是我們都要懂得看.
有時弟弟很煩.
又頂撞等等.
那又怎樣呢?.
忍不住.
罵他.
粗口.
打他一巴掌.
我告訴你.
我做了一句.
我說弟弟孩子最不可愛的時候.
就是他最需要愛的時候.
他知道自己不對.
他在發脾氣.
生氣.
你一巴掌打他.
就罵他.
慘了.
他更加淒涼.
但他預你打.
預你罵.
你不打不罵.
哎呀.
他就感覺你的寬容.
你的寬恕.
他明白你需要什麼.
寬容比懲罰更加感動.
更加可以改變人.
那一刻你控制不打不罵.
你就造就了他一部分.
因為愛.
因為被接納.
所建立出來的正性性格.
所以得到了解.
得到支持.
得到鼓勵.
寬恕.
改過的機會.
都是在我們家要做的.

$^{681}$有些人經常沒有正能量.
怎樣負能量呢?.
有位妹妹告訴我.
陳醫生我出生至現在.
二十多歲.
一天都聽到爸爸說.
爸爸就說.
唉.
明天又要上班了.
我覺得上班真的很辛苦.
他真的被壓迫.
他很無奈.
我真的很同情他.
但他為什麼還要上班.
他真的很慘.
害得我都害怕上班.
另外一個.
經常說.
唉.
快要被解僱了.
不過還好還沒解僱.
說來說去都還沒解僱.
說了二十多年都還沒解僱.
你天天說被解僱幹什麼.
害得全家被籠罩於一埋慘霧之中.
很怕失去支柱.
所以傳達的應該是正面的訊息.
行教是做出來的.
希望大家能夠接納.
而且有益於彼此.
有新教就是傳遞人生態度.
等等.
傳遞什麼人生態度呢?.
正面.
彼此相愛.
不要介意他們不好.
相愛是彼此的.
有回應的.
你做的事情對方會有回應.
不是單方面的.

$^{721}$所以相愛之中.
聖經說.
用荷西夏書.
教導是用慈誠哲獻人.
用慈誠哀索.
我要醫治他.
慈誠哀索不是整損人.
是有大令的.
哀索是拉著你.
讓你不能自由行走.
不要去哪裡.
或者對你有危險的地方.
是阻礙你的.
是約束你的.
子女最不喜歡父母管著.
我去哪裡他又要問.
和誰一起又要問.
什麼都要報告.
很討厭.
很煩.
我很不喜歡被他管.
其實不是.
我們喜歡被人管.
但當中我們得到愛.
我喜歡講弟子經.
弟子經有一句很好.
出必告.
返必面.
你出街一定要告訴父母.
你去哪裡.
去多久.
和誰一起.
回來吃飯嗎.
講完返必面才重要.
你一定要回來.
要給父母面子.
那時就要去老爺奶奶房.
叩頭.
老爺奶奶.
她的爺爺.

$^{761}$我回來了.
那我奶奶就要摸摸他的臉.
乖孫子你平安回來就好了.
原來你這樣講給家人聽.
他們放心.
他們知道你平安.
管有什麼不好.
他是用愛來規範你.
你是得到益處.
在原則.
應該在平安範圍內.
但現在的人藐視父母.
聖經說愚昧之子.
藐視父母親.
所以不要藐視.
這些慈情愛惻.
是欣賞的.
我們喜歡被這些規管.
是好的.
但我們要恰當的.
慈情愛惻.
不是去打他.
不是去罰他.
是很好的.
慈愛的性.
不會弄傷一隻牛.
只是令他得到益處.
又令他得到平安.
然後引導他應該走的路.
You raise me up.
就是這樣.
是幫助支持.
引領的.
好.
那.
但你怎樣才.
可以有這樣的力量呢.
有些人負面得太多.
最有力的就是.
聖靈的九種果子.

$^{801}$我喜歡.
聖靈的九種果子.
是仁愛喜樂和平.
忍耐恩慈.
良善信實.
溫柔節制.
這個裡頭.
是有素質素.
要有這些素質素.
就是充滿內在美.
但我在想原來那棵樹.
是用它最好的.
吸收.
最好的養分.
去送去果子那裡.
那它的果子才會長得.
大.
漂亮.
多汁.
好味道.
原來這個果子是用來繁衍宗族的.
它要繼續出其他樹.
延續後代的傳宗接代繁殖.
所以它是完全不計較的.
完全付出給了這些種子.
因為它明年又會再結果.
或者幾年後又會再結果.
是重複摘完又有.
生生不息.
因為它有陽光空氣.
有雨水的滋潤.
向下扎根.
向上生長.
原來人是仰望我們神的帶領.
神是栽種.
又令它生長.
所以不要緊.
這些果子是可以結出來.
結出來的就要用.
不要收起來.

$^{841}$因為遲些又有了.
連在基督裡面.
新鮮.
給人.
好味道.
予人有益.
有營養.
無私的.
不計較的付出.
不要緊.
因為源源不斷.
我們有生命力.
連在基督裡.
不是一次過的.
所以不要緊.
這個力量的傳源.
是聖靈的果子.
現在我們聖靈保衛師.
在我們這裡幫助我們.
彼此相愛.
是要彼此.
子女都要對父母好.
我們不惹兒女的氣.
總要按著聖經的教訓.
和篤澤.
好好地教育他們.
有個兒子說.
我知道了.
他說了很多次.
還在說.
我看一看他媽媽.
媽媽說.
他做.
我用不著說.
所以最好讓媽媽.
不要這麼嘮叨.
就是實踐她的教導.
子女都有責任.
不只是罵媽媽嘮叨.
你做好.

$^{881}$她用不著煞氣.
她喜歡省一把氣.
不知道多高興.
有些兒女.
很奇怪的.
有一個女朋友.
媽媽就來跟我說.
什麼什麼女朋友.
叫我幫她做功課.
我說做什麼功課.
她說原來讀設計的.
要出十五個版.
刺繡版.
叫我幫她做.
我說你就建一兩個給她.
教她.
她說不是.
她要我做完十五個.
她要我做完她的功課.
刺到我的手指都爛了.
我說這樣都可以.
她說不止.
她來我家吃飯.
下午打電話來.
指定菜單.
一定要很多芝士.
讓我去超級市場買.
哎呀.
我說不要緊.
哄哄兒子開心.
你也喜歡支持兒子.
她說不是.
她叫我做白痴.
我幫她做那麼多事.
她還叫我做白痴.
然後兒子給她看.
她跟我們拍的照片.
三個人.
我先生我和兒子.
那個女友剪了我的樣子.

$^{921}$我說不是這樣.
她說這麼暴力.
叫兒子想想.
這個也不知道是否好的人選.
還沒入門.
已經有那麼多要求.
那麼多奇怪的事.
嗯.
現在流行移民.
人們移民賣樓.
移移下移.
連媽媽也沒有房子住.
那媽媽怎麼辦.
媽媽不適應移民的生活.
不能讓她留在這裡.
入老人院.
這樣也行.
你看報紙有很多人.
逼媽媽賣樓.
賣完就分錢給她.
這樣也行.
聖經說.
如每人藐視母親.
自慰之子使父親歡樂.
如每人藐視母親.
令到媽媽非常不開心.
所以彼此相愛.
你愛她沒用.
愛她有用.
她愛你才行.
彼此相愛.
這是關係.
這是相對的.
要好的.
當然.
有時我們受別人影響.
有些不妥.
所以很重要的行.
就是為子女代求.
我們做不到.

$^{961}$我們不想那麼憎氣.
你就為她祈禱.
我也不知道她去了哪裡.
怎知道她在做什麼.
不要緊.
歌羅西書.
你看一章九節.
蘇羅也沒去過哥羅西.
他也不知道.
但他說不要緊.
我懇切為你們祈求.
不一定.
總之神知道.
我看不到.
但神知道.
求你們守潔心清.
遵行神的話.
令他們潔淨自己的行為.
神在管他們.
所以交給神是最好的方法.
行教.
是請神來教.
聖經還有一個關於家庭的事.
就是.
滿屋豐盈.
大家相憎.
不如有一塊餅乾.
大家相安.
這個平安是重要的.
在神子降生的時候.
天使天君唱歌.
在天上榮耀歸上帝.
在地上平安歸於他所喜悅的人.
平安是重要的.
耶穌復活之後.
周圍顯現的時候.
第一句話.
願你們平安.
平安是重要的.
在疫情當中有平安.

$^{1001}$在疾病,失落,迷茫,災難,徬徨,痛苦,死亡之中.
也是有平安.
是神自己的應許.
是給予祂所喜悅的人.
我們就是努力做神喜悅的人.
我們一定得到平安.
你看看這個平安.
是雙方面的.
你說你想平安.
對方不想平安.
也不行.
所以要努力.
平安一定有.
第一 物質方面,身體方面的平安.
沒有那麼多病.
要免疫力強一點.
身體強壯.
很重要的是睡眠.
睡眠足夠.
因為老人痔瘴症.
認知障礙症.
原因是睡眠不足.
少於八小時.
很多名人.
下輩子獎的人.
臨老後變成痔瘴症.
為什麼呢.
因為他有不眠不休的體力.
趁年輕的時候不睡.
不到你說了算.
腦部有一個排除垃圾的系統.
叫做Glymphatic system.
身體的血液裡的Glymphatic.
在我們的淋巴.
Glymphatic 不知道怎麼譯.
加上G.
就是把腦裡的腦細胞活動.
所產生出來的廢物.
加上Beta-amyloid.
如果存在的話.

$^{1041}$就令腦細胞萎縮.
令細胞多了.
於是就導致老人痔瘴症.
但怎樣才能令這個系統好呢.
這個排除廢物的系統.
怎樣才是最高效率呢.
原來不是全日都同樣速度開啟.
開動最強最快的時候.
是當你深層睡眠的時候.
Deep Sleep.
所以最好我們有足夠的睡眠.
多點深層睡眠.
希望八小時的睡眠.
就會清除垃圾.
清除得好.
腦部清醒.
大家年紀大的時候.
都會有更好的睡眠能力.
當然我現在不是講講座.
那些很重要.
身體是很重要.
現在的小朋友是不睡覺的.
他們要打機.
打到通宵達旦.
又對眼又壞.
什麼都壞了.
總之很不妥.
所以要告訴他們.
不要珍惜自己的身體.
不要自強.
平安是靠身體健康.
要懂得安全.
心境要平安.
Emotionally.
心境要平安.
就要做得好一點.
愛你周圍的人.
我有個朋友去探望媽媽.
都不想探望陳醫生.
為什麼.

$^{1081}$我媽一見到我就說.
你不是我鄉下的表妹.
然後就說.
為什麼你對我這麼好.
我是你的女兒.
我哪有這麼大的女兒.
真慘.
所以趁人還認得我們.
我們還可以表達愛.
令他感覺到合同的關係的時候.
我們就要珍惜.
努力做.
永遠不多.
你努力做.
時間是有限的.
所以你的情緒.
你做得夠多.
你就平安.
到某一天.
我們會特別平安.
我們做夠.
然後經濟上的平安.
我們年紀大一點.
穩定一點.
有些小朋友很刺激.
現在不是投資.
是網上投機.
很可怕.
倫敦金都出了.
現在買比特幣.
很可怕.
很不濟.
所以求神給他們智慧.
不要這樣賭.
要有他們不平安.
令我們也不平安.
他們以為很平安.
他以為很賺.
嚇得我們不得了.
然後社交上的平安.

$^{1121}$我們喜歡建立社交網絡.
我們教會有很多團契.
非常好.
人與人之間的關係.
是唯一令人少心臟病的方法.
有些研究.
很多研究都是這樣說.
所以家中有愛的時候.
平安是彼此感受得到.
我在這裡看「安」字很多.
有什麼「安」呢.
安慰,安心,安苦,安全,安靜.
安寧,安樂,安居,安頓,安儉.
安放,安家,安康,安舒,安壯.
原來我們都很喜歡平安.
耶穌說願我們平安.
很重要.
什麼境遇.
我們有神的時候就平安.
我們將這個平安帶回家.
但我們要先感覺到平安.
你才可以傳遞到平安的感覺.
所以人行教.
然後去新教.
新教是很重要的.
你有新的情緒.
有這樣的榜樣.
有這樣的氣質.
其實人不只這些.
聖經說你要做炎,做光.
炎我就在想.
厲害了.
炎是可以調和.
用炎調和.
我們做父母.
在家庭裡.
令家庭的關係和諧起來.
用炎調和.
大家都很舒服.
不要吵架.

$^{1161}$多點和氣.
多點處理衝突的技巧.
然後它令很多東西都好吃了.
又不會悶,不會寡.
現在說居家很悶.
我們做炎.
有新意.
多做新意.
好吃.
多煮新菜.
然後炎是防腐.
你知道我母親92歲.
我四年前和她吵架.
我說你一定要換智能手機.
她說不行.
有什麼用.
不懂玩,不懂用.
我說你不換.
你很快就會消失.
我送了一個給她.
她是用了.
她現在很開心.
因為她可以做很多事.
所以我們不可以發霉.
我們不可以糜爛.
我們要日久常新.
我們要追上時代.
不斷終身學習.
不可以只看很悶的電影.
我們也要學看Netflix等等.
不然就不可以.
兒女說我們不要惹怒她的電影.
你經常聽粵曲.
又被她罵.
我們可以帶耳機.
又送了耳機給她.
所以我們要追上時代.
這是好的榜樣.
讓她終身學習.
不然她說你不學我也不學.

$^{1201}$是不可以的.
是不斷學習.
不要發霉.
非常好.
我們又學了一些新名詞.
又叫網紅,又叫網友.
又叫「碌下機」.
上網瀏覽一些東西.
看網站.
自我感覺也良好.
第四是消毒.
鹽是消毒.
現在很流行.
現在喜歡防疫.
喜歡消毒.
但鹽最重要是改變.
鹽可以改變一些東西.
菇草無味.
不好的.
一小小就好了很多.
教養兒女.
我們是用教導的方式.
我們傳授給她的方式.
是資料的傳遞.
教她.
質她填鴨式學習.
目的就是formation.
我們想她們成為能達的人.
formation.
一個好人格的素組.
但不幸很多因素走偏離了.
變成deformation.
變成貓formation.
古怪.
生得怪了.
所以我們要努力.
放鹽.
放聖經教導.
用多些禱告.
這個叫做transformation.

$^{1241}$神可以改變.
改變transformation.
改變之後.
reformation.
更新.
做成一個新的人.
所以.
重要的是我們要做鹽.
然後我說.
做光.
聖經說.
你們就像光放在燈台上.
照亮一家人.
光很重要.
光一照.
東西有色彩.
看得清楚.
又不會悶.
光令人所有行為.
都是光明磊落.
坦坦蕩蕩.
不會藏起來.
不喜歡藏在黑暗的罪惡當中.
光明是好.
是乾淨的.
又照亮一家人.
因為你的智慧.
你的傳講.
你的好的榜樣.
你就令全家的人.
都有改變.
變得更好.
見得更陽光.
見得更喜樂.
當然第四.
光是希望.
因為.
黑夜將盡.
黎明又會到.
光是帶給人無量的希望.

$^{1281}$所以.
有一個.
余光宗寫了一本書.
他講武難日.
武難日就是他自己的生日.
以前母親要順產.
很辛苦.
他說.
父憂母難日也.
因為孩子的誕生.
意味著父親的憂心.
和母親懷胎十月的疼痛苦難.
所以他對母親.
寫下了生她的感恩之情.
他母親過世了.
他寫了一封信給母親.
寫給久別的母親.
他說.
慘了.
但我都不知道寄去哪.
天國什麼字頭.
不知道什麼字頭.
電話又不知道打去哪.
就算接通都不知道說什麼.
因為母親都不明白.
現在世代的事.
沒有什麼共同的話題.
但不要緊.
唯一不變的.
是我對母親永久的感恩.
所以如果有母親在.
我們更多的行動.
表達.
今天很多人會去喝茶.
送花.
送小禮物.
陪母親.
什麼都好.
什麼方法.
都是表達.

$^{1321}$你記得這件事.
你心中牽掛.
你肯來.
和她分享一些活動.
身體上來陪伴.
你有些東西給她.
大家可以享受這些.
你小小的禮物.
都是真心.
真誠.
是有意思的.
是真的給她.
什麼形式都好.
你的安排.
你的心思.
都是你的母親.
會感受得到.
但最重要跟她說.
就是感恩.
多謝我們的母親.
不只是母親節.
每天都是母親節.
我們的母親.
每天都是母親.
所以今天.
喜歡大家一起說.
但平時.
我們一樣愛我們的母親.
充滿感恩.
思想行為.
心理上.
都是很欣賞她.
為我們做這麼好的母親.
\newpage



\section{使徒行傳 2:37-47}
\label{sec:7e2_k14U9kk}
\textbf{2023年1月22日崇拜直播｜陳明泉牧師｜旺丁旺才的家｜使徒行傳2\_37-47}
\newline
\newline
連結: \href{https://youtube.com/watch?v=7e2-k14U9kk}{\texttt{https://youtube.com/watch?v=7e2-k14U9kk}} ~~~~ 語音日期: 2024-03-27
\newline
\newline
\hyperref[sec:1_gQDpS_bXo]{\small{< < < PREV SERMON < < <}}
~
\hyperref[sec:index_chronic]{\small{[返順時目]}}
~
\hyperref[sec:DVdI0qIQgiA]{\small{> > > NEXT SERMON > > >}}
\newline
\newline
使徒行傳 2:37-47
\newline
\begin{longtable}{cl}
\hline
\hline
章節 & 經文 (和合本修訂版)\\
\hline
2:37 & \begin{tabularx}{0.7\textwidth}{X} 眾人聽見這話，覺得扎心，就對彼得和其餘的使徒說：「諸位弟兄，我們該怎樣做呢？」 \end{tabularx} \\ \\ \relax
2:38 & \begin{tabularx}{0.7\textwidth}{X} 彼得對他們說：「你們各人要悔改，奉耶穌基督的名受洗，使你們的罪得赦免，就會領受所賜的聖靈。 \end{tabularx} \\ \\ \relax
2:39 & \begin{tabularx}{0.7\textwidth}{X} 因為這應許是給你們和你們的兒女，並一切在遠方的人，就是給所有主—我們的神所召來的人。」 \end{tabularx} \\ \\ \relax
2:40 & \begin{tabularx}{0.7\textwidth}{X} 彼得還用更多別的話作見證，勸勉他們說：「你們當救自己脫離這彎曲的世代。」 \end{tabularx} \\ \\ \relax
2:41 & \begin{tabularx}{0.7\textwidth}{X} 於是領受他話的人，都受了洗；那一天，門徒約添了三千人。 \end{tabularx} \\ \\ \relax
2:42 & \begin{tabularx}{0.7\textwidth}{X} 他們都專注於使徒的教導和彼此的團契，擘餅和祈禱。 \end{tabularx} \\ \\ \relax
2:43 & \begin{tabularx}{0.7\textwidth}{X} 眾人都心存敬畏；使徒們又行了許多奇事神蹟。 \end{tabularx} \\ \\ \relax
2:44 & \begin{tabularx}{0.7\textwidth}{X} 信的人都聚在一處，凡物公用， \end{tabularx} \\ \\ \relax
2:45 & \begin{tabularx}{0.7\textwidth}{X} 又賣了田產和家業，照每一個人所需要的分給他們。 \end{tabularx} \\ \\ \relax
2:46 & \begin{tabularx}{0.7\textwidth}{X} 他們天天同心合意恆切地在聖殿裡敬拜，且在家中擘餅，存著歡喜坦誠的心用飯， \end{tabularx} \\ \\ \relax
2:47 & \begin{tabularx}{0.7\textwidth}{X} 讚美神，得全體百姓的喜愛。主將得救的人天天加給他們。 \end{tabularx} \\ \\
[1ex]
\hline
\hline
\end{longtable}
$^{1}$我們以下要聽神給我們的說話.
祂的說話今天記載在.
《士道行傳》第二章的三十七節.
至到四十七節.
《士道行傳》第二章三十七節.
眾人聽見者說覺得擇心.
就對彼得和其餘的仕途說.
弟兄們我們當怎樣肯?.
彼得說:你們各人要悔改.
奉耶穌基督的名受浸.
叫你們的罪得赦.
就必領受所賜的聖靈.
因為這應許是給你們和你們的兒女.
並一切在遠方的人.
受事主勸勉我們說.
他們說:你們當救自己脫離這彎曲的世代.
於是領受他話的人就受了浸.
那一天門徒約添了三千人.
都恆心遵守仕途的教訓.
彼此交接,抹餅,祈禱.
眾人都懼怕.
仕途又行了許多奇事神蹟.
信的人都在一處.
繁物公用並且賣掉田產,家業.
照各人所需用的分給各人.
他們天天同心合意行節的在殿裡.
且在家中抹餅.
傳著歡喜誠實的心用飯.
讚美神,得眾民的喜愛.
主張得救的人天天加給他們.
聖經獨筆.
多謝各位.
主席.
弟兄姊妹平安.
新年夢福.
最重要是身體健康.
我們每一年都帶著積極的盼望.
迎接新的一年來臨.
今年我們教會有新裝修.
今天兒子一堂,大家很熱鬧.

$^{41}$我想起旺角浸信會.
旺角真是一個頗旺的地方.
所以我新的一年有一個很簡單的願望.
希望這個地方的群體旺丁旺財.
財不是很多錢的財.
是人才輩出的財.
希望我們這裡不單得救的人數加添.
而且更加有很多弟兄姊妹.
很多門徒可以興起.
所以我帶著這個心.
在聖經裡找一段很興旺的經文.
今天這段經文其實不是過年讀的.
不過在裡面的意象很識切.
我們今天藉著這段經文來彼此互勉.
除了用這段經文之外.
我今天也很想拋書包.
我中文不是很理想,不是很厲害.
不過剛才不論聽福音樂曲也好.
或者在街上看到.
在新年裡面有很多對聯,春聯.
那些字寫得很精煉.
幾個字就能寫出意境.
我再翻查一些舊的書.
看看唐詩.
很多唐代的詩.
其實都是春聯.
蘇東坡也有寫.
王羲之也有寫.
墨寶在新春的佳節裡.
只寫幾個字.
但把一些重要的訊息.
一些祝福跟家人說.
也跟身邊的鄰舍說.
今天我就希望用三對的對聯.
不是我作的,我抄回來的.
三對春聯.
希望將今天我們讀《仕途行傳》的這段經文.
可以有一些總結的方式來表達.
藉著這幾句或這一段的經文.
彼此祝福,彼此勉勵.

$^{81}$在新的一年我們更進一步.
我們一起去祈禱,求神幫助我們.
同心祈禱.
阿巴迪夫,我們謝謝你.
願你在我們當中帶領我們中弟兄姊妹.
我們深信上帝你一直看顧我們這個群體.
我們帶著感恩的心.
也帶著積極盼望.
來展望新的一年.
願主你繼續來對先恩.
與中弟兄姊妹同在.
有藉著你的話語給我們有激勵.
我們恭敬將來的時間交託給你.
願你賜福,願你人靈.
幫助宣講的說得清楚.
聆聽的每位弟兄姊妹聽得明白.
進到我們的內心.
給我們有力量實踐.
奉基督耶穌,得勝的聖名祈求.
Amen.
首先講一下這段經文.
今天選這段經文有一點背景.
這段經文是.
《史唐人傳》第二章.
話說耶穌復活升天.
完成了十字架的救贖大使命之後.
頒布門徒將福音棒傳給他們.
接著這班門徒一直等待.
等到五旬節來臨.
五旬節來臨的時候.
除了有神跡騎士聖靈.
很外顯地將能力挑灌在信徒和使徒當中.
更重要的是當這些使徒.
或者彼得.
特別是彼得在這裡是一個代表人物.
領受了上帝的聖靈之後.
他就很勇敢地對著百姓.
宣講上帝救恩的信息.
其實整個講論就在二章十四節開始.
一個很長的篇幅.

$^{121}$他所講的就是.
舊約所允許的聖靈.
好像一份禮物.
賜給他屬地的子民.
現在就允許了.
值得留意的是.
為什麼講聖靈這件事這麼重要呢?.
在舊約裡面.
聖靈其實是很少人被上帝揀選.
他的靈與他同在.
聖靈同在的通常是什麼人呢?.
或者有聖靈的能力.
那些領袖,君王,先知.
甚至乎有些藝術家.
上帝的靈都會與這一小撮人同在.
不過在上帝的計劃裡面.
上帝不是純粹將他的靈.
囂冠在小撮人身上.
藉著基督耶穌的救贖.
引悔改.
到了史唐人傳裡面這一局的講論.
就講到聖靈的能力.
聖靈來自上帝的新生命.
上帝的同在.
就要在這個時間裡面.
囂冠在眾百姓的身上.
比特就講一個很簡單的信息.
就是在講現在上帝的救恩.
上帝的聖靈現在來到了.
你們每一個都要領受.
你們每一個都能夠有這份恩典.
來領受上帝的靈與他們同在.
所以這是一個很簡單的信息.
當他講完這個很簡單的信息的時候.
在二章第37節.
就是剛才我們開始讀的那段經文.
那些百姓聽到比特講的這一番話.
眾人聽見這話就覺得責心.
就對比特和其餘的使徒說.
弟兄們,我們當怎樣行.

$^{161}$比特說,你們各人要悔改.
奉耶穌基督的名受贈.
叫你們的罪得赦.
就必領受所賜的聖靈.
因為這應許是給你們和你們的兒女.
並一切在遠方的人.
就是主我們神所召來的.
比特還用許多話作見證.
勸勉他們說.
你們當救自己脫離這彎曲的世代.
這裡這幾節我用一句春聯.
就是守經除咎嚴,普渡作新民.
為什麼要用這兩句話對戀來表達呢.
正如剛才所說.
這班百姓其實沒有想過.
或者一直等待一個君王來改變他們自己的生命.
不過上帝的計劃比他們這班以色列民.
所想所求更多.
走得更遠.
上帝不單止救贖這班人.
上帝藉著基督耶穌成為了他們的彌賽亞.
更加重要.
這個彌賽亞的工作.
不只是地上的生活這麼簡單.
彌賽亞的工作除了救贖他們的罪.
帶領這班以色列民之外.
這個更加是一個永恆的工作.
上帝將他們的靈與普天下的相信主耶穌的這一班人同在.
所以在這一班人裡面.
他們一直以來都在守著聖經.
他們對聖經所說的話一點都不陌生.
甚至彼得用《使徒行傳》所說的講論.
他除了在說基督耶穌的工作之外.
他不是憑空來捏造這些內容.
他是用《約珥書》和《舊約聖經》.
聖經裡面如何應許的那些說話.
就藉著耶穌基督彰顯了出來.
所以這一班人昔日.
他們守著那本聖經.
他們在讀的時候.

$^{201}$以為內宗有永生的道.
在等待.
但是他們沒有想過.
這個恩典這個應許已經實現了.
藉著基督耶穌.
聖靈就臨到這一班人身上.
改變他們.
他們只要願意接受基督耶穌.
奉基督耶穌的名去受浸.
和真正的悔改.
這一班人就得著永恆的福氣.
這裡有一句話.
或者在經文裡面有一個字眼很重要.
就是這一班人聽到這個信息的時候.
眾人聽見就覺得「擇心」.
「擇心」是什麼意思呢?.
「擇心」中文翻譯得很好.
和原文是一字.
「擇心」的意思就是說.
心裡面撩撩撩.
有刀子刺刀的那種感覺.
覺得不安樂.
心裡面有一種催逼感.
有一種心痛的感覺.
不過你看到彼得所說的內容.
其實目的不是說耶穌基督.
做很多很受苦的事情.
而是說一個榮耀的救贖.
聽到一個很榮耀的計劃.
原則上是雄心萬象的.
不過聽到一個雄心萬象.
一個這麼積極的福音信息.
但那些人的心裡面竟然覺得.
有刀子刺刀的那種滋味.
原因就是.
因為想起這個這麼榮耀的福音.
又想起自己生命的光景.
那個落差實在太大了.
如果我們每一個基督徒都知道.
上帝為我們預備了生命裡面一個豐厚的恩典.

$^{241}$但是我們是不是真的這麼配合呢?.
其實不是的.
這些福氣.
想起自己有時候自己行事為人.
自己所想的東西.
和這個福音其實有一段距離的.
但是現在說將這個應許淋在自己身上.
有一個大的應許.
又有一個美好的計劃.
但是自己的生命又不是這麼理想的時候.
這個落差就令到聽的百姓.
心裡面就覺得有一種責心的感受.
經文裡面.
當這些人覺得責心.
心裡面不安樂.
他就問彌德.
我們應該怎樣去做.
在這裡我再仔細的說多一點.
很多新來教會的朋友聽到福音訊息.
我們期望.
或者我們自己也是這樣.
期望聽一個很舒服的訊息.
聽一個令到我.
上帝很愛我.
令到我很寶貝.
那些很哄你.
很安慰的訊息.
我們覺得那個是福音.
這個不是錯的.
原來聽福音真正聆聽的人.
其實大部份的人不是越聽越安樂.
是越聽越不安樂.
覺得心裡面有些攪動.
為什麼呢?.
如果當你心裡面有這一份責心.
我恭喜你.
你用心去聽.
因為你知道.
你的生命處境和這個福音是有距離.
而這一種距離感.

$^{281}$這一種責心.
是令人尋問上帝究竟我可以怎樣做.
我可以怎樣回應一個這麼偉大.
這麼美好的福音訊息.
當彼得或者那些使徒聽到一班百姓.
問他們這個問題的時候.
彼得教他們一個很簡單的道理.
你們各人要悔改.
逢基督耶穌的名受針.
叫你們的罪得赦.
就必領受所賜的聖靈.
因為這應許是給你們和你們的兒女.
並一切在遠方的人.
就是主我們所召來的.
經文裡面其實很簡單.
怎樣去回應這種責心的體會呢?.
悔改.
將你生命的過去.
現在要做一個更新改變.
悔改這個字有時候我們說得很情感化.
悔改這個字很直接.
走那條路.
現在改換跑道.
因為知道舊時那條路這樣繼續走下去.
是不太通的.
走一條新的路.
走一條對自己生命有益處的.
對自己和對自己的後代有益處的一條路.
這個就是悔改.
當這班以色列民.
舊時他們的舊路是什麼呢?.
他們覺得守住字句.
看著聖經裡面.
跟著法利賽人的教導.
想著守住這些律法.
就會得到永恆的生命.
這個本身也是一種很敬虔的表現.
不過守住之後他們就不是守聖經了.
這班人就守住了法利賽人.
或者當代裡面.

$^{321}$所謂道德高地的這班人.
他們加了重重的枷鎖在真理當中.
又添加一點點.
本身他們添加的目的.
也不是出於自私的目的.
是基於敬虔的目的.
不過越加就令信仰越來越複雜.
令本來令人得生命的聖經的話語.
越來越多額外的詮釋.
在這些詮釋裡面又加了很多人為的傳統.
我們覺得那些傳統是很重要的.
有時候守那些傳統比守聖經更重要.
悔改的意思就是說.
放回正確的位置.
守回聖經最初的道路.
藉著耶穌基督的救贖.
重新找回自己生命要走的那條路.
悔改的意思就是說.
舊有的那條路不通的.
就轉換了.
走一條新的路.
當他們這樣做的時候.
這個悔改還要是有決心的.
藉著一個朕禮.
這個朕禮其實禮儀本身沒有結成的果效.
那個禮儀只不過是表達一種決心.
奉基督耶穌的名授朕.
就是說納入成為基督耶穌的一份子.
成為耶穌基督的門徒.
用這種決心來走這條人生新的路.
藉著這種生命.
來活現一個生命新的狀況.
藉著悔改受朕認罪.
罪就得到赦免.
領受所賜下的聖靈.
這個聖靈是一個活的.
上帝內在人的心裡.
人就得著這份盼望.
就將舊時的過去改變.
走一條新的路.

$^{361}$守經除舊掩.
普渡作新聞.
守著聖經的應許.
回到聖經那裡.
那些舊掩未必是指不好的行為.
那些掩當然是指大掩缸.
外面的世俗裡有很多問題.
所以彼得也說.
你要脫離這個世代裡很多的污穢.
班曲缽謀的世代.
要脫離這些舊掩.
不過第二個舊掩有另外一個意思.
就是有一些已完木掩.
有一些我們積累下來的.
我們覺得很重要的傳統.
因著聖靈光照之下.
我們就除去了這些想法.
我們走新的路.
我們不是反傳統.
不過有一些傳統.
或者有一些做法.
可能跟聖經有違背.
我們就要將它調和.
我們將它改頭換面.
或者基督化.
我們不需要一定要.
有一些民間習俗.
活上頭柱香.
我們沒有頭柱香.
不過年初一一起來敬拜神.
我們就將它基督化.
我們家裡的那些灰穿.
我們通常.
香港人有時候都很習慣.
恭喜發財.
我們香港人喜歡發財.
但不是錯的.
不過你會知道.
總是跟基督徒的生命追求有些不同.
那你就恭喜人家.

$^{401}$繼續生命步步高升.
對上帝或對人見人愛都可以.
你將它轉換了.
其實就是將一些舊有的.
積累下來的一些想法.
一些傳統.
有一些可能適合聖經.
你就保留.
有一些不適合聖經.
你就過濾.
重新調教它.
以至你成為一個.
活現一個生生命的基督徒.
帶著這一份信仰.
活現一個神國子民的新身份.
當你能夠這樣做的時候.
你會為世界帶來改變.
當這班人很聽話.
這班人問彼得怎樣做的時候.
這班人真的決意去悔改.
經文記載了.
二章41至第45節.
於是領受他的話的人就受了針.
那一天門徒約添了三千人.
都恆心遵守使徒的教訓.
彼此交折磨餅祈禱.
眾人都懼怕.
使徒又行了許多歧視神蹟.
信的人都在一處.
凡物公用.
並且賣了田產家業.
照各人所需用的分給各人.
經文裡面就講到.
當他們知道得救之道.
當他們知道很擇心.
需要悔改.
這班人就一齊去回應使徒的講論.
經文裡面記載了一個很誇張的場面.
當人聽到這個訊息的時候.
三千人願意受針.

$^{441}$如果真的做針禮都會針到手軟.
三千人.
但你可以想像.
不單是個人層面的悔改.
而是家家戶戶.
他們可能整個家族.
整個家庭.
整個社群.
這班人聽見彼得宣講的福音訊息.
全部人都願意去悔改.
藉著這個禮儀.
藉著針禮.
宣認他們是基督耶穌的門徒.
繼續在生活裡遵守使徒的教訓.
令他們繼續過新的生活.
很多頂展會.
一到針禮.
就以為是終點站.
其實那不是信仰的終點站.
那個針禮其實是一個中轉站.
這個中轉站就是宣認自己的關係.
和上帝是一個使徒之間的親密關係.
就正如結婚.
結婚在禮拜堂裡行禮.
那你沒理由交換戒指之後結婚.
然後愛情就結束.
是新的階段來開始.
更加親密的關係.
同樣道理在這裡針禮.
就帶來更加深入.
更加外顯.
與舊時有不同的新生命.
他們的新生命是如何表達出來呢?.
經文裡面就講到一個很震撼的場景.
就是那些人一起彼此交接.
交接的意思是團契.
我們聚會的意思.
一起聚會.
一起抹餅.
抹餅其實是吃的.

$^{481}$不過是分甘同味.
當然裡面有一些禮儀的元素.
藉著抹餅紀念基督耶穌的救贖.
祈禱.
藉著這些很簡單的生活禮儀.
他們就一齊來生活.
做這件事的時候.
剛才講到這三千人一齊來.
集體做一個很簡單的舉動.
但整個社會見到他們是怎樣呢?.
眾人都懼怕.
你覺得很奇怪.
做些事令人害怕.
那個懼怕的意思就是說.
這班人有上帝在.
開始帶著敬畏的心來看這班.
和當代社會裡面.
有一些不同的一種生命表現.
經文裡面再繼續的來講.
試圖繼續行許多的歧視神蹟.
做許多的醫治.
甚至乎趕鬼等等.
做許多外顯性的救贖工作.
令信的人一路一路的來到家庭.
信的人都在一處.
另一個很具體表現到.
這個令人懼怕的那種生命表彰就是.
人們將自己的東西繁蜜公用.
賣了自己的田產家業.
照各人所需的分給各人.
嘩!今天如果教會這樣做.
你真是嚇壞人了.
不過是嚇壞你的家人.
我這樣叫人.
叫人賣掉自己的東西.
我沒有所謂,我不是很多.
但在座下面的很多.
賣掉它.
人們都覺得這個教會.
是否有騙徒的手法.

$^{521}$不過整件事不是說.
那些試圖要威逼你.
要逼那些信徒來做這件事.
整件事是出於內心的自發.
知道因為他們地上的產業.
最終都不能永遠存留在自己那裡.
唯有上帝成為他們的產業.
他帶著感恩,帶著滿足.
帶著一種願意分享的心.
他才有價值觀的調教.
所以在這班人受浸了之後.
他們一路一路活現的整個信仰生命.
是一個價值觀重新調教的.
以前他們看這些東西是很重要的.
接下來他們就開始變得.
看那些地上的錢財不是最重要.
反而是能夠祝福身邊的人.
看到身邊有不同需要的人.
就讓那樣東西更加優先.
所以他們願意賣掉田產.
其實田產,家業是基於他們的心.
他們的價值觀被改變.
我想用第二句春聯來表達.
這班人以神為產業,愛侶度生涯.
這班人帶著基督耶穌的愛.
帶著一個新的價值觀念.
他們以上帝成為他們生命最豐富的產業.
他們有上帝,慢慢很多東西.
生命就有很多東西變得沒那麼重要.
以上帝一路一路擺在生命的首位的時候.
人生就有很多舊時為錢財.
為拿多一點的執著.
慢慢一路一路放下.
我自己成長.
我不是在旺鎮長大.
我也是在旺鎮長大的.
不過我是年紀比較大才在旺鎮長大.
我年輕時信主.
第一間回教會就是一間潮州人的教會.
我認識一些同輩.

$^{561}$他們每年新年都很開心.
為什麼呢?原來他們這一班人.
他們說他們上一代.
他們照樣賺錢.
舊時在香港,普遍都不是那麼豐富的日子.
這班鄉里就自己住在一起.
住在木屋區.
住在一些偷渡下來的地方.
有些是申請下來的.
總之在香港.
其實就是要隻身來到.
都是要靠鄉里互相的關懷.
有一些基督徒來到香港之後.
他們就組織這一班的同鄉.
一起來聚會.
這些同鄉可能.
舊時在他們自己的鄉下沒有信主.
來到香港就信了主.
後來大家轉為一起.
香港有些地方就叫基督教的村落.
現在就越來越少.
在九龍城舊時有一條村.
因為木屋燒了就重新建.
宣教士給錢的.
那條叫做博愛村.
有沒有聽過?.
有些弟兄姊妹聽過.
博愛村.
其實舊時有些木屋區.
後來因為火災.
教會給錢重新建.
然後就建了一些石屋.
在村落中間有一個禮拜堂.
小小的禮拜堂.
成為了他們所有年輕人.
所有平時聚集的一個很重要的場所.
我聽他們那些.
其實我不知道.
現在博愛村都拆了.
不過我聽他們每一次說.

$^{601}$我們以前在博愛村禮堂裡玩的.
有些說頑皮的.
在大禮堂前面一群弟兄姊妹聚集.
其實後台的背景有塊玻璃.
後面有個小房間.
我們在那裡做什麼?.
我們在那裡玩啤牌.
當牧師在那裡講道的時候.
後面有人在那裡玩啤牌.
你說多威風.
然後又試過一次打架.
我們在玻璃門後面打架.
我爆響後.
我們認識的一個很知名的一位牧者.
其實以前就在玻璃窗的玻璃窗後面打架.
現在已經做了知名的資深牧師.
打架兩個人.
兩個人都是打著打著就打了.
成為兩個重要的牧師.
究竟是誰呢?.
遲些會回來香港的.
有時候會去意大利.
那個呢.
那個他跟我說的.
他跟我說的.
不過雖然說打架或者玩.
但是這一班人有任何需要.
他們就互相的來幫忙.
缺乏的日子.
或者教養孩童.
那種很親密的關係.
無分彼此.
無分你我.
在一個這樣的群體當中.
在這個博愛村裡面.
這一班人就是這樣.
互相的成長.
今天我們旺角浸信會.
我們弟兄姊妹住在不同的地方.
我們分佈在港九新界.

$^{641}$我們要追回舊時博愛村那種生活.
是很難的.
不過在我們跟人相處裡面.
在弟兄姊妹之間的情誼.
我們的情可以拉近了我們地域上的距離.
我們彼此用愛相待.
我們看自己不要看得那麼重.
我們凡事都是看別人.
看我們的群體為先.
這種生命帶著一種很大的震撼力.
當你們彼此相愛的時候.
眾人就認得出你們就是上帝的門徒.
在這裡為什麼整個社會都懼怕他們.
怕不是因為覺得他們很恐怖.
而是覺得他們做了一些.
當代社會不會做的.
個個都是看自己的需要.
但是這群人卻可以無私地奉獻.
用上帝的心腸.
用那種價值觀.
新的價值觀來彼此相待.
這個群體就成為一個非常有吸引力的群體.
這三千人繼續地輪轉.
成為當代社會一個很重要的改變更新的現象.
經文還沒完的.
二章四十六到四十七節.
這群頂展媒繼續在教會.
在他們的信仰當中繼續來活現.
他們天天同心合意行徹地在奠禮.
且在家中抹餅.
傳著歡喜誠實的心用飯.
讚美神.
得眾民的喜愛.
主張得救的人天天加給他們.
在這裡值得注意的字眼就是天天.
天天都有新朋友.
天天都是同心合意.
這個教會真的很興旺.
三千多人你說很難管理.
原來不是很難管理.

$^{681}$只要有愛心.
有上帝的話語.
其實多少人都可以黏合在一起.
而且基數越大.
對於社會上的震撼力更加大.
這群人天天帶著單純.
願意分享.
帶著上帝的愛.
來跟別人分享生活的一切.
他們在殿裡祈禱.
在殿裡一起尋求上帝.
因為有上帝.
所以這群人放下了自己很多想法.
這群人就被聚焦起來.
聚合起來.
這群人在家中抹餅.
每一次吃粗茶淡飯也好.
或者像我們昨晚吃的年尾飯.
吃得很豐富也好.
無論吃什麼也好.
都是存著感恩的心.
誠實的心來用飯.
帶著歡喜.
無論吃什麼也好.
你吃得快樂的時候.
你的心快樂.
吃的每一樣東西都是這麼好吃.
他們這群弟兄姊妹.
他們的物質一定沒有我們今天這麼豐富.
他們只是帶著一份同心.
快樂.
這群人.
他們所吃的飯.
帶出一個充滿了人間快樂味道的飯菜.
藉著這種心懷念念.
他們又將一種感恩.
向上帝表明.
上帝為了這個群體.
他又每天都加增新的人給他們.
這段經文裡面說什麼呢?.

$^{721}$我想起一句對聯.
就是蘇東坡寫的.
向陽門第春常在.
即善人家慶有餘.
這句是在說什麼呢?.
向陽門第就是在說那扇門.
看到太陽.
經常看到太陽的家.
春天的生氣.
經常都在家裡.
其實這不是在說家庭坐向.
其實是在說家裡.
這群家人自己為內.
有沒有那種快樂.
正能量.
那種充斥在他們的內心當中.
如果這個家.
經常帶著笑聲.
很陽光的.
如果你家裡.
或者我們教會的群體.
打開門.
人家聽到的就是笑聲.
其他層樓聽到.
嘩!這裡怎麼笑得這麼開心.
都好奇想知道是什麼地方.
向陽門第春常在.
如果我們這個地方充滿了笑聲.
充滿了快樂.
剛才大家進來升降機的時候.
我弟兄姊妹自高奮勇.
我又幫忙恭喜一下.
真的抱著.
真的差很遠.
整個教會的溫度.
那種融入程度.
令你很想留在這裡.
來完之後.
接著再拜年.
你都會帶著這裡的快樂.

$^{761}$帶到世界當中.
那種生氣.
我們互相感染.
我們彼此帶著積極的心.
經文裡面再繼續說.
或者在春聯下半句.
回應經文.
即善人家,慶有餘.
經常做很多善事的家庭.
他們的歡慶會越來越多.
有餘.
我們經常都想快樂多一點.
做多點善事.
做多點真的落手落腳的善事.
真的去幫助身邊的人.
分享你生命裡面所擁有.
你覺得好的東西.
將你的生命一路一路擴展的時候.
這個群體就充滿了吸引力.
我希望我們旺角浸信會.
我們每一個都是笑逐顏開.
我們返教會返得很開心.
今天這個世界很奢侈.
不知道你們知不知道.
好像負面一點.
給我一點點兩句.
很多人說返教會不好.
牧師怎樣怎樣.
有這些有那些.
經常都很心煩意亂.
我們不要這樣.
返教會就要返得快樂.
很多事情.
其實很多煩惱.
可能是自己執著某一點.
覺得自己的想法最重要.
這樣就很難快樂.
只要我們反過來.
我們見到別人快樂.
願意將自己的快樂和別人分享.

$^{801}$有時可能自己吃虧一點.
但見到別人快樂很滿足.
你的心裡面有一種.
難以言喻的滿足感.
當我們每個人都走一步的時候.
這個群體就真的很陽光.
很有力量.
更加重要就是那種.
行善的能力.
不是別人逼的.
而是自發的.
當我們願意這樣去回應上帝.
又回應身邊不同人的需要的時候.
我們相信這個群體.
我們的快樂會繼續加添.
旺角浸信會.
如果帶著這樣的盼望.
或者這種心態.
相信我們真的成為一間旺丁旺財的教會.
今天的信息很簡單.
不過在整個信息裡面.
我們要認定.
上帝是一個賜福的上帝.
我們的信仰是由基督耶穌.
祂的生命換給我們.
祂將生命的豐盛.
藉著祂的十字架救贖.
現在已經給過你和我.
我們就帶著這份豐盛的恩典.
領受那份人力無法得到.
無法努力換取的聖靈的同在.
帶著這份活力.
帶著這份恩典.
我們開始有些改變.
我們不要看自己的需要.
或者看自己覺得重要的東西.
放到這麼高.
我們看上帝作為最重要.
調教我們的價值觀念.
調教我們的情緒.

$^{841}$我們笑逐顏開.
和身邊的人同行.
我相信我們帶著這種心態.
在未來的日子裡.
上帝會大大使用.
賜福給你和我.
求主繼續幫助我們.
我們要宣認.
旺國浸信會是一個.
「黃丁」「黃財」的教會.
願主賜福你們眾人.
\newpage



\section{ }
\label{sec:DVdI0qIQgiA}
\textbf{2024年3月29日 受苦節主餐崇拜直播｜仰望聖潔被殺羔羊}
\newline
\newline
連結: \href{https://youtube.com/watch?v=DVdI0qIQgiA}{\texttt{https://youtube.com/watch?v=DVdI0qIQgiA}} ~~~~ 語音日期: 2024-03-30
\newline
\newline
\hyperref[sec:7e2_k14U9kk]{\small{< < < PREV SERMON < < <}}
~
\hyperref[sec:index_chronic]{\small{[返順時目]}}
~
\hyperref[sec:g_Q6zBFvts0]{\small{> > > NEXT SERMON > > >}}
\newline
\newline
$^{1}$讓我們繼續定睛思想耶穌基督.
在十字架上所說的七句話.
我們看到即使耶穌基督.
無論是身體或心靈.
已經處於最痛苦,最脆弱的階段.
但他仍然沒有忘記世人的需要.
以及上帝所託付給他的使命.
當我們回望這七句話.
你會發現頭三句是耶穌表達.
對世人的一份關愛.
耶穌仍然記掛著那些陷害他.
有份釘他十架.
羞辱他,嘲笑他的人.
他們一再否定耶穌是上帝的兒子.
拒絕接受耶穌成為他們的救主.
但他們原來不知道.
他們將要面對永垷的結果.
將要永遠承受一個不認耶穌的苦.
但我們看到.
耶穌不單沒有以牙還牙.
沒有求上帝降禍來消滅他們.
相反,耶穌求天父赦免.
因為他們根本不知道自己正在得罪上帝.
因為這群人.
其實都是耶穌所愛,所要憐憫,拯救的對象.
另一方面,我們也看到.
耶穌紀念那一位和他同釘十架.
願意信靠他的那一個人.
耶穌不想他帶著一種遺憾.
帶著那種罪疚感離開世上.
於是耶穌對他說.
「我實在告訴你.
今天你要和我在樂園裡面」.
耶穌即時接納了他.
成為天國的成員.
賜給他永生的確居.
我們也看到耶穌.
他很記掛著那位和他即將分離的母親.
耶穌知道在這一刻很需要.
有一個人好像兒子一樣.

$^{41}$來和瑪利亞同行.
去照顧她,安慰她.
陪伴她過到這個哀傷的時刻.
我們看到這位被釘十架的耶穌.
他應驗了舊約聖經以賽亞裡所說的.
他不單只擔當了世人的罪孽.
他更加背負了每一個人的痛苦和憂患.
至於十架七弦裡面的第四句.
耶穌說「我的神,我的神,為甚麼離棄我?」.
以及第五句「我渴了」.
其實分別是來自舊約的詩篇.
二十二篇的一節.
和六十九篇的第二十一節.
我們會奇怪為甚麼耶穌在釘十架的裡面.
突然間唸起舊約聖經.
原來這兩篇的詩篇都是來自大衛.
大衛在他被仇敵圍剿.
在他被仇敵攻擊.
甚至乎癡笑他生命正在危在旦夕的時候.
大衛透過這兩篇的詩篇.
描述他的心路歷程.
在當中大衛說最令他難受的.
不單單是敵人對他的迫害.
而是上帝遲遲都不出手幫他,拯救他.
甚至大衛他自己覺得.
上帝彷彿掩面不看他.
離他而去.
我們見到大衛在受盡煎熬的時候.
在他苦無出路的時候.
他只有向他的上帝來呼喊.
他一方面來哀嘆.
一路問上帝為何要對自己不聞不問.
但同時他也哀求上帝.
上帝你要紀念我的苦況.
上帝你不會離棄我.
他也期望上帝會早日出手.
將他從那些苦難中拯救出來.
因此.
耶穌基督在十字架上.
唸起大衛這兩篇.

$^{81}$和受苦掙扎有關的詩篇.
一方面應驗了舊約聖經中.
指向耶穌的預言.
而更重要是表示.
耶穌不是那位高高在上.
離地生活的救主.
耶穌很願意親上世上的苦難.
甚至是最極端的痛苦.
不單被人離棄.
更被人出賣.
被最愛的人出賣.
也被他的同胞屈往.
被他們凌辱鞭打.
最後釘死在十架上.
耶穌不單只和大衛的痛苦來認同.
耶穌也要告訴我們每一個.
世人所經受的苦難.
耶穌是絕對理解的.
因為祂曾經切身處地感受過.
事實上.
不知道大家有沒有想過.
如果耶穌只需要.
替我們作為罪人來受死.
其實祂可以有很多相對輕散的方式.
祂不一定要用十字架.
死得那麼慘,那麼淒遠.
耶穌大可以像做烈士一樣.
轟轟烈烈.
痛痛快快風光地離開.
不需要受盡別人的煎熬和委屈.
為什麼耶穌要死得那麼痛苦?.
正正是因為.
祂不單只要承擔我們的罪.
祂更要背負我們所受的苦.
我們所經的憂患.
而這一切就是.
負上帝託付給他的使命.
所以.
耶穌基督比我們任何一個更了解.
其實苦難受苦.

$^{121}$到底是怎麼一回事?.
祂不是遠處的旁觀.
祂不是純粹用同情的心態.
來理解我們的苦難.
祂乃是親自踏上這條苦路.
徹身處地與世人一同去受苦.
並且耶穌祂知道.
即使在祂自己最痛苦.
最孤單無助的時刻.
甚至連負上帝都彷彿離棄了自己的時候.
耶穌仍然確信.
這不是祂的結局.
因此.
祂念的詩篇二十二篇和六十九篇的下半部.
我們會看到.
詩人不再停留在痛苦的掙扎裡.
而是向上帝發出一份稱讚.
一份讚美.
因為詩人確信.
上帝必定會垂聽他的禱告.
最終也會出手來拯救.
替他平反.
同樣耶穌基督在十字架的時候.
祂仍然確信.
負上帝不會離開他.
不會忽視在痛苦裡掙扎.
尋求他幫助的人.
就正如耶穌會即時回應.
那位釘十架的囚犯.
他會垂聽他的哀求.
賜給他永生的盼望一樣.
並且耶穌仍然確信.
縱使自己已經身處十架.
縱使好像一個惡人當道的世代.
但是一切仍然在負上帝掌握的手上.
上帝會藉著耶穌基督在十字架上的犧牲,捨命.
完成整個的救贖計劃.
所以我們看到.
來到剛才的十架第六燕.
耶穌是說「成了」.

$^{161}$耶穌不是說玩完了.
失敗了.
特別當我們細看.
約翰福音十九章二十八至三十節的時候.
其實你會發現「成了」這個詞語.
一共出現了三次.
第一次在二十八節的上.
耶穌知道各樣的事情已經完成了.
他們成就救恩.
近了.
第二次記載「成了」在二十八節的下.
和本譯做「應驗」.
其實也是解作「成了」成就的意思.
顯示耶穌那一刻說口渴.
其實並不是要解決他身體上的需要.
而是他刻意去成就.
舊約聖經指向他的應許.
至於第三次的「成了」.
在第三十節.
耶穌上過那一躁之後.
就說「成了」.
徹底完成上帝所託付給他的使命.
我們看到這三次的「成了」.
強調一件事情.
耶穌並不是像周邊的人所想.
他玩完了.
他失敗了.
相反.
「成了」是代表這位被釘十架的救主.
不單止沒有失敗.
他乃是默默.
在無人知曉的情況下.
完成了上帝的拯救.
弟兄姊妹.
耶穌基督對父上帝這份完全的信靠.
使他有力量.
來忍受各樣的痛苦.
各種的屈辱.
並且他堅持信服到最後一刻.
直到完成上帝所託付給他的工作.

$^{201}$耶穌在十字架上.
所呈現出來的力量.
其實也成為今天.
我們每一個的幫助和盼望.
試想像一下.
在這個世界上.
有沒有一份工作.
無論是職場上的工作.
或是在家中的工作.
是不需要伴隨任何的痛苦或屈辱.
即使那份工作會帶給你一些成就.
滿足或回報.
其實你在裡面也難免會遇到挫敗.
你會遇到委屈.
甚至乎所謂受苦的時候.
畢竟這個世界和身邊的人.
不是我們所想的那麼完美.
我們自己也不是.
在台灣裡邊.
有一個很出名的宣教士叫做馬佳.
他的英文名是George Leslie McKay.
這個人的名字在台灣裡邊.
可以說是無人不識.
因為在台灣有六間醫院.
都是以他命名.
馬佳雖然是加拿大人.
他也是長老會警察會的牧師.
在1871年他踏足了台灣.
展開了宣教的服務.
當時很多台灣人對外籍的宣教士.
仍然有很多的保留.
有很多的偏見.
譬如說會說他們.
是挖走那些死了的信徒的心臟.
或者眼珠來制約.
對他們有很多的謠言.
後來當清政府和法國.
來展開爆發大戰的時候.
馬佳也被波及.
他雖然是加拿大人.

$^{241}$但也因為他是外國人.
所以被很多台灣人.
都當成他和法國人是同一夥.
於是乎對他加以排擠迫害.
甚至破壞攻擊他所服侍的教會.
我們不難想像.
無故地遭受別人這樣迫害.
每一個人都會覺得委屈.
馬佳當然也深受屈辱.
但很奇怪.
在他晚年的時候.
他親自寫了一首詩送給台灣.
他對曾經深深傷害過自己的.
這片土地和裡面的人民.
表達了一份無限的關愛.
這首詩裡大概是這樣說.
我衷心所愛的台灣.
我全部的青春都獻給你.
我難捨難離的台灣.
我將我的人生都獻給你.
我一生的快樂其實就在這地方.
我盼望我人生的尾站.
都可以在這片地方.
在那些海浪拍岸的聲響裡.
在那些竹林搖曳的樹蔭裡.
找到我的歸宿.
為何馬佳可以這麼心繫.
這明明是令他深受屈辱痛苦的地方.
為何他仍然希望服侍當地的人.
到他的終老.
正正是因為.
耶穌基督賦予他的使命感.
遠遠大於他所受的屈辱和傷害.
所以他仍然繼續留在這個地方.
這個曾經傷害過他的地方.
完成上帝交託給他的事情.
所以今天我們看到.
以他命名的不單單只是醫院.
有很多也尊稱他為創辦人的機構,學校.
也在台灣四處林立.

$^{281}$我們看到信仰帶給他那種使命感.
讓他可以曾經在那些迫害過他的人身上.
彰顯耶穌基督的愛.
他視昔日的屈辱和傷害.
已經是過去式.
不足介意.
事實上弟兄姊妹.
我想沒有一份工作.
包括福音的工作.
是不用受苦的.
不需要面對屈辱.
就連我們的主耶穌.
他的福音工作也不例外.
他為了承受承擔上帝給他的使命.
他面對過各種的屈辱.
各種的迫害.
我們看到他的德性的關鍵.
就是他體重自己的使命.
過於他所承受的委屈痛苦.
使他有力量可以忍受十字架上的一切.
同樣.
今天我們說納治要跟隨耶穌基督.
其實也意味我們會和耶穌一樣.
經歷這個世界.
經歷世人對我們的不友善.
怎樣可以站立得穩.
直到完成上帝託付給我們的使命.
其實很視乎我們有沒有認定自己的身份.
我們有沒有認定自己是上帝的僕人.
認定他給我們的身份.
認定他託付給我們的工作.
從而我們可以輕看別人對我們的評價.
我們不再將我們生活的焦點.
困在別人的目光裡.
使我們很多時太急於討人的喜悅.
要認同別人的喜好.
忘記了基督徒真正的人生目標.
是追求上帝的喜悅.
當我們回望耶穌基督的一生的時候.
當我們回望他被釘十架的時候.

$^{321}$我們會發現耶穌基督.
沒錯,他走過一條艱辛的路.
但同時又是一條最有意義的路.
提醒我們每一個信徒的目標.
不是追求一個理想舒適的生活.
而是追求一個無論環境順逆.
我們都與主同行.
忠於所託付的一個人生.
雖然耶穌一生有很多坎坷.
不單被人誤解.
更被人鄙視,出賣,唾棄.
就連最親的門徒都離開他.
沒有人在耶穌釘十架的時候.
可以體諒他,明白他的處境.
有的只是那些積累,處心積慮.
要害死他的人.
種種負面的經歷.
其實耶穌絕對有理由自怨自艾.
但耶穌看這一切的痛楚.
卻是極具正面的意義和價值.
因為隨著他宣告「成了」.
這一切的苦難.
都通通轉化為他救恩的一部分.
亦都成為了每個面對痛苦的人的出路和盼望.
耶穌基督「成了」的宣告.
其實大大強化了我們在事奉上.
在人生上的抗逆,抗壓力量.
就正如在《腓立比書》一章十九節裡說.
因為我們蒙恩不單要得以信服基督.
並且要為他受苦.
顯示蒙恩和苦難是並存的.
一個蒙恩成為基督徒的人.
其實也代表他蒙恩為主受苦.
與耶穌基督的苦難有分.
所以,弟子們或許我們仍然覺得.
為主受苦好像是很遙遠的事情.
特別是今天一個講求舒適,成功的世代.
沒有人想經歷一個受苦的人生.
就連信徒可能也有這種想法.
總是想避開那些痛苦,那些苦難的經歷.

$^{361}$不過,我們的信仰由始至終離不開苦難這兩個字.
連我們信仰的最高領導耶穌基督.
他也是一位受苦的主.
他也通過苦難來成就上帝的吩咐.
建立教會.
所以作為教會一份子的我們.
自必然也會在我們的信仰裡.
在實踐福音的事情上.
經歷不同程度的苦.
我還記得有一位弟兄姊妹.
曾經分享她很希望她的信仰可以有些突破.
她求上帝給她一些痛苦的經歷.
我們當時還笑她.
覺得她奇怪.
怎會有人去求痛苦,求苦難.
不過當我們回想的時候.
其實這句話很熟悉.
我們時常說.
希望和耶穌基督同行.
希望跟隨他.
和他的生命有份.
其實背後不是提醒我們.
要預備隨時為主受苦.
事實上,每年我們去紀念耶穌.
舉行受苦之崇拜.
其實也是重複提醒我們.
為主受苦.
本來就是基督徒一個必定的配套.
耶穌基督通過受苦.
見證上帝的旨意在他身上怎樣彰顯.
同樣耶穌基督的門徒.
都必然要走這條蒙福的路.
事實上,弟兄姊妹.
通過受苦,達到靈性的成長.
跟耶穌和其他門徒聯合.
建立一份更緊密的關係.
這個對苦難這麼正面的思維.
其實不是今天才有的.
早在羅馬書第八章.
十六至十八節.

$^{401}$其實已經這樣勉勵我們.
我讀給大家聽.
斯圖普羅這樣說.
聖靈與我們的心同正.
我們是神的兒女.
既是兒女,就是後子.
就是神的後子和基督同作後子.
如果我們和他一同受苦.
也必和他同得榮耀.
我想,現在的苦楚.
比起將來要顯於我們的榮耀.
就不足介意了.
斯圖普羅勉勵我們.
重新去看,什麼是受苦.
他勉勵我們將痛苦看為.
一個與耶穌基督與上帝聯合的途徑.
他知道我們會因為愛神愛人的緣故.
而經受苦楚.
但他也勉勵我們.
要想到那一刻.
將來上帝極重無比的榮耀.
那一份屬天的福樂.
就輕看現今所受的痛苦.
因為保羅知道.
當我們受苦的時候.
耶穌基督不單止理解.
他也與我們同在.
與我們一起去受苦.
所以我們也可以像保羅一樣.
效法基督.
將自己在受苦的時候.
完全交託給上帝.
讓我們可以安穩在上帝的裡面.
並且確信最終.
我們會與耶穌基督同作後子.
一起去承受上帝的榮耀.
和永遠的福樂.
所以弟兄姊妹.
盼望今天晚上.
我們藉著耶穌十加七元.

$^{441}$這七句話.
再一次來互相激勵和提醒.
使我們的行事為人.
好像聖經所說.
與我們所蒙的恩相情.
願意以耶穌基督的心為心.
願意與耶穌受死的生命有份.
我們低頭一起去有個思想.
弟兄姊妹.
如果耶穌這一刻要你透過受苦.
以至更加能夠體會.
他在十字架上所經受的.
你願意嗎?.
又或者.
如果耶穌要透過你的淚水.
以至你能夠更加體會.
耶穌基督在十字架上所擔當的憂患.
你願意嗎?.
我們同心祈禱.
我們願意.
我們願意愛你.
是因為你先愛我們.
為我們犧牲.
耶穌啊.
我們願意犧牲.
是因為你先降悲.
為我們捨己.
我們願意捨己.
是因為耶穌基督.
你已經為我們受死.
我們願意捨你.
為我們受死.
所以主啊.
我們帶著這份恩典.
帶著這份感恩.
我們來到你的面前.
是我們不配.
但你卻願意給我們.
所以耶穌啊.
幫助我們願意.

$^{481}$以愛還愛.
我們願意為你受苦.
求你幫助成全我們.
可能我們會有掙扎.
但幫助我們仰望.
因為在十字架上的你.
以致我們能夠有更多的力量.
為上帝去堅持.
為上帝去多走一步.
耶穌說我們願意為你受苦.
是因為不單止你深愛我們.
更加是因為.
你願意.
在我們苦難的裡面.
和我們同在.
耶穌基督說我們願意.
在苦難的裡面與你有份.
不單止因為我們很愛你.
更加是因為.
你愛我們更深.
所以耶穌幫助我們每一個.
在仰望你的十字架.
仰望你的十架七賢的時候.
讓我們願意和你一樣.
背起你為我們預備的十架.
一生跟隨你.
讓我們在裡面.
同時會經歷到.
耶穌你給我們的盼望.
你給我們的平安.
使我們可以代表你.
去將你的平安和盼望.
播送給身邊有需要的人.
讓我們在痛苦的裡面.
我們同樣經歷你的醫治和保護.
使我們可以代表你.
成為別人的安慰和守望.
求聖靈你與我們同在.
求聖靈你幫助我們.
繼續定睛仰望耶穌你.

$^{521}$奉耶穌你的名去祈求.
阿們.
\newpage



\section{馬太福音 27:62-28:15}
\label{sec:g_Q6zBFvts0}
\textbf{2024年3月31日復活節崇拜直播｜陳明泉牧師｜無法封住的真相｜馬太福音27\_62-28\_15}
\newline
\newline
連結: \href{https://youtube.com/watch?v=g_Q6zBFvts0}{\texttt{https://youtube.com/watch?v=g\_Q6zBFvts0}} ~~~~ 語音日期: 2024-03-31
\newline
\newline
\hyperref[sec:DVdI0qIQgiA]{\small{< < < PREV SERMON < < <}}
~
\hyperref[sec:index_chronic]{\small{[返順時目]}}
~
\hyperref[sec:JykGqUs7644]{\small{> > > NEXT SERMON > > >}}
\newline
\newline
馬太福音 27:62-28:15
\newline
\begin{longtable}{cl}
\hline
\hline
章節 & 經文 (和合本修訂版)\\
\hline
27:62 & \begin{tabularx}{0.7\textwidth}{X} 次日，就是預備日的第二天，祭司長和法利賽人聚集來見彼拉多， \end{tabularx} \\ \\ \relax
27:63 & \begin{tabularx}{0.7\textwidth}{X} 說：「大人，我們記得那迷惑人的還活著的時候曾說：『三天後我要復活。』 \end{tabularx} \\ \\ \relax
27:64 & \begin{tabularx}{0.7\textwidth}{X} 因此，請吩咐人將墳墓把守妥當，直到第三天，恐怕他的門徒來把他偷了去，就告訴百姓說：『他從死人中復活了。』這樣的話，那後來的迷惑就比先前的更厲害了。」 \end{tabularx} \\ \\ \relax
27:65 & \begin{tabularx}{0.7\textwidth}{X} 彼拉多說：「你們有看守的兵，去吧！盡你們所能的把守妥當。」 \end{tabularx} \\ \\ \relax
27:66 & \begin{tabularx}{0.7\textwidth}{X} 他們就帶著看守的兵同去，封了石頭，將墳墓把守妥當。 \end{tabularx} \\ \\ \relax
28:1 & \begin{tabularx}{0.7\textwidth}{X} 安息日過後，七日的第一日，天快亮的時候，抹大拉的馬利亞和另一個馬利亞來看墳墓。 \end{tabularx} \\ \\ \relax
28:2 & \begin{tabularx}{0.7\textwidth}{X} 忽然，地大震動；因為有主的一個使者從天上下來，把石頭滾開，坐在上面。 \end{tabularx} \\ \\ \relax
28:3 & \begin{tabularx}{0.7\textwidth}{X} 他的相貌如同閃電，衣服潔白如雪。 \end{tabularx} \\ \\ \relax
28:4 & \begin{tabularx}{0.7\textwidth}{X} 看守的人嚇得渾身顫抖，甚至和死人一樣。 \end{tabularx} \\ \\ \relax
28:5 & \begin{tabularx}{0.7\textwidth}{X} 天使回應婦女說：「不要害怕！我知道你們是尋找那釘十字架的耶穌。 \end{tabularx} \\ \\ \relax
28:6 & \begin{tabularx}{0.7\textwidth}{X} 他不在這裡，照他所說的，他已經復活了。你們來！看看安放他的地方。 \end{tabularx} \\ \\ \relax
28:7 & \begin{tabularx}{0.7\textwidth}{X} 快去告訴他的門徒，說他已從死人中復活了，並且要比你們先到加利利去，在那裡你們會看見他。看哪！我已經告訴你們了。」 \end{tabularx} \\ \\ \relax
28:8 & \begin{tabularx}{0.7\textwidth}{X} 婦女們急忙離開墳墓，又害怕，又大為歡喜，跑去告訴他的門徒。 \end{tabularx} \\ \\ \relax
28:9 & \begin{tabularx}{0.7\textwidth}{X} 忽然，耶穌迎上她們，說：「平安！」她們就上前抱住他的腳拜他。 \end{tabularx} \\ \\ \relax
28:10 & \begin{tabularx}{0.7\textwidth}{X} 耶穌對她們說：「不要害怕！你們去告訴我的弟兄，叫他們往加利利去，在那裡會見到我。」 \end{tabularx} \\ \\ \relax
28:11 & \begin{tabularx}{0.7\textwidth}{X} 她們去的時候，看守的兵有幾個進城去，把所發生的事都報告祭司長。 \end{tabularx} \\ \\ \relax
28:12 & \begin{tabularx}{0.7\textwidth}{X} 祭司長和長老聚集商議，就拿許多銀錢給士兵， \end{tabularx} \\ \\ \relax
28:13 & \begin{tabularx}{0.7\textwidth}{X} 說：「你們要這樣說：『夜間我們睡覺的時候，他的門徒來把他偷去了。』 \end{tabularx} \\ \\ \relax
28:14 & \begin{tabularx}{0.7\textwidth}{X} 若是這話被總督聽見，有我們勸他，保你們無事。」 \end{tabularx} \\ \\ \relax
28:15 & \begin{tabularx}{0.7\textwidth}{X} 士兵收了銀錢，就照所囑咐他們的去做。這話就在猶太人中間流傳，直到今日。 \end{tabularx} \\ \\
[1ex]
\hline
\hline
\end{longtable}
$^{1}$陳鑑林議員.
「明日之事」第七章第六十二至二十八章十五節.
「次日,就是預備日的第二天,祭司長和法利賽人聚集來見彼拉多,說:大人,我們記得那誘惑人的還活著的時候曾說:三日後我要復活,因此請吩咐人將墳墓把守妥當..
直到第三日,恐怕他的門徒來把他偷了去,就告訴百姓說他從死裡復活了,這樣那後來的迷惑比先前的更厲害了..
彼拉多說:你們有看守的兵去吧,盡你們所能的把守妥當.他們就帶著看守的兵同去,封了石頭,將墳墓把守妥當..
「安息日將盡,七日的頭一日,天快亮的時候,沒大拉的瑪利亞和那個瑪利亞來看墳墓,忽然地帶震動,因為有主的使者從天上下來,把石頭滾開,坐在上面,他的樣貌如同閃電,衣服潔白如雪,.
看守的人就因他嚇得渾身亂顫,甚至和死人一樣.」天使對婦女說:「不要害怕,我知道你們是尋找那釘十字架的耶穌,他不在這裡..
照他所說的,已經復活了,你們來看安放在主的地方,快去告訴他的門徒說,他從死裡復活了,並且在你們以先往加利利去,在那裡你們要見他..
我已經告訴你們了.」婦女們就急忙離開墳墓,又害怕,又大大的歡喜,跑去要報給他的門徒.忽然,耶穌遇見他們說:「願你們平安.」他們就上前抱著他的腳拜他..
耶穌對他們說:「不要害怕,你們去告訴我的弟兄,叫他們往加利利去,在那裡必見我.」他們去的時候,看守的兵有幾個進城去,將所經歷的事都報給祭司長..
祭司長和長老聚集相宜,就拿許多銀錢給兵丁說:「你們要這樣說,夜間我們睡覺的時候,他的門徒來把他偷去了,上若這話被巡撫聽見,由我們勸他保你們沒事.」.
兵丁收了錢銀,就照所祝福他們的去行,這話就傳說在猶太人中間,直到今日,聖經讀不完,神賜福,常讀他話語,並且遵行的人..
弟兄姊妹平安,今天是慶祝基督耶穌復活的大日子,基督耶穌在十字架上經歷苦難,為我們眾人流血捨身,在祂死後的第三天,正式向世人宣告,祂已經戰勝死亡,祂已經復活..
不單止祂自己的復活,祂的生命,祂所走過的路,也成為了每一個信靠祂而來的一條要走的路,我們每一個基督徒都是跟隨基督耶穌的腳踵,我們都同樣得著一個永恆復活的生命..
基督的復活與拉撒路的復活不同,拉撒路的復活是說祂因為病患肉體死亡,祂是基督耶穌在祂身上作的神蹟,讓祂可以渡過由死而復活,就是身體重新復活的生命,.
而屈天,拉撒路,他也要再一次經歷肉體的死亡,不過基督耶穌的復活就不同,基督耶穌的復活是上帝讓他戰勝了死亡的權勢,死亡已經不再成為他的酬金,死亡成為了他的手下敗將,.
基督的復活就是帶進一個永恆的生命,將一種永生的盼望帶給這個世界,他自己走過這條路,也引領每一個信靠他的兒女走這條復活的道路..
復活不單止是肉體死亡得著永生,復活的能力也內在每一個願意信靠基督耶穌的人的生命裡,所以很多人因著復活的主,我們很多的困難可以得到改變,.
我們生命裡面很多的苦難可以被扭轉,生命裡面有很多昔日的對立關係,因著基督耶穌復活的生命,我們可以重新得著復活的一種生命,復活是基督教裡面最重要的信念,.
我們因著信靠耶穌,不單止是信他是一個道德偉人,更加相信他是一個死而復活,永生神的兒子,我們每一個因著復活,我們就得以改變每一個的生命,.
復活是我們核心信仰裡面最重要的核心之中的核心,保羅說如果有人說復活的事過去了,那個人就要被責備,甚至保羅也會覺得如果在說復活的事,不信耶穌復活的事,我們其實吃喝快樂都可以,因為我們人生都是現世過了就是這樣就算了,.
不過保羅也好,早期教會,直到我們今天,我們都是傳講一個復活的信息,我們的生命得以改變,我們信靠的那位上帝是一個復活的主,這個很簡單,我們理解了復活是一個很重要,對我們來說要傳揚的一個信息,.
那麼耶穌復活的這個消息是不是對於所有人來說都是很歡迎呢?今天我們所讀的這段經文我們看到,其實昔日和耶穌對著幹的那些對頭人,知道耶穌或者避免了耶穌復活這個信息,繼續的來傳廣,所以他們無所不用其技,.
但是在一個好像帶著不良動機,要壓制這個復活的信息傳廣的同時,上帝將壞事成為一個好像反面的證據,說著基督耶穌復活,今天我們用這段經文,比較長的篇幅,我們看看上帝怎樣去用這一班人,一班心裡不懷好意,不想復活的信息被傳開,.
但是竟然讓他們成為了一個佐證,讓復活的事,藉著他們的生命來傳揚開去,我們一起祈禱,求上帝幫助我們,.
天父上帝願你在我們當中帶領我們,讓我們繼續重溫你復活的大能,你的話語,你的信息,讓我們心靈裡面得著更大的確據,願你在此時此刻藉著你的話語,重新來幫我們認清我們所信的,.
求你在我們當中,你對我們內心說話,讓你復活的大能,佔據我們生命的全部,願你繼續帶領我們,祝福我們,我們下來聽受你話語的時間,願你幫助每位聆聽的弟兄姊妹聽得明白,宣講的講得清楚,恭敬將來的時間交託,願你賜福,獻上禱告奉基督耶穌得勝的聖名,Amen.
我們先看第一個段落,27章第62節至66節,在這裡講到耶穌基督受死,埋葬在墓穴,那個墓穴大概是一個財主的藥室,為他預備的一個地方安葬耶穌,在那天預備日的第二天,祭司長和快利賽人聚集,來見彼拉多,一般來說,祭司長和快利賽人不相往來,.
不過到了這一刻,他們就成為了一個聯盟,來找彼拉多,就是將耶穌宣告釘在十字架的那個官員,他們就希望彼拉多幫他們做一些配合,怎樣的配合呢?.
他說:「大人,我記得,那誘惑人還活著的時候曾說:「三日後我要復活.」因此請吩咐人將墳墓把守妥當,直到第三日,恐怕他的門徒來把他偷去,就告訴百姓說:「他從死裡復活了.」.
在第63,64節裡面,就這麼說,那個聯盟,快利賽,再加上祭司長,他們就求彼拉多給他們一些兵丁,把守著耶穌安葬的那個墳墓,那個墳墓跟我們中國人那些有些不同的,那個墳墓是一個山洞,山洞裡面還有一個大的幾噸重的石頭來堵塞住他的..
所以整件事基本上是一個已經很穩陣的了,不過這一幫法利賽人和祭司,他們都擔心如果那些門徒將耶穌的屍體偷去,那些人發現不到耶穌的屍體,就會將一個復活的信息傳遞出去..
他們心裡面想,如果耶穌的屍首不見了,就代表著復活的事就出現了.他們的心裡面不是想著耶穌會復活,只是他們將焦點放在耶穌的身體那裡,找不到那個身體就好像成為了一個佐證了..
所以他們連這個機會傳出去的信息都不容讓出現,於是就找兵丁來把守在那個墳墓當中,避免了後來要說要一個誘惑人的信息,誘惑人因為耶穌在他昔日裡面他曾經預告三天之後他要復活..
所以彼拉多就按照他們的心意,你們也有些看守的兵,就去將那裡把守妥當.於是整件事就說到拒絕耶穌復活,甚至他要確保所有東西穩穩當當,確保這件事他只能把守三天,過了第四天這件事就已經不成功了..
他就用這個士兵把守在耶穌的墳墓面前,確保耶穌的屍體不會被人偷走.事實上你說要偷走也不是那麼容易的,那個石頭是一個很重的,端以幾噸的一塊大的石頭,需要很多人才能推得動的..
現在再加上那些兵兵,有武器的,在那裡把守著,基本上整件事都是滴水不漏的了,基本上就沒有辦法將耶穌的屍體偷走.在這一幫人裡面,他們的思路是代表著什麼呢?.
一塊堵塞著墳墓的墓石頭,加上一幫士兵的把守,基本上這個已經是萬無一失,確保復活的信息不會傳遞出去的一條方程式..
在他們的思維,在他們很直線的想法,總之,使耶穌的屍體一直在那個墳墓裡面,那件事就不會被傳揚,他們視為誘惑人的信息就不會繼續傳遞出去..
為什麼那些文士,或者說那些祭司長和法利賽人,這麼緊張不容讓這個復活的信息繼續傳揚出去呢?在這裡沒有直接說的,不過在之前的篇幅裡面說到,如果耶穌基督的復活成為了真實,或者這件事流了出去的時候,.

$^{41}$對於當時佔領著論述的上風,或者那個話語權,對於上帝的話語權,法利賽人他們就完全沒有任何的地位,因為在耶穌昔日的傳講的信息裡面,某程度上都是針對法利賽人的講論的,.
不是針對法利賽人背後的律法,而是針對法利賽人他們那種傳遞,或者傳適律法的那種態度.當耶穌復活的時候,某程度上就代表著上帝為耶穌來辯護,以致顯出錯誤的就是法利賽人..
所以他們會用盡所有的方法,確保整件事,確保他們的話語權要籠罩在他們這個學派身上.當然對於祭司長來說,同樣道理,因為耶穌在潔淨聖殿裡面,某程度上都是針對當時的祭祀制度,.
也就是說當時的聖殿,就像敬拜神一樣,後來被耶穌視為賊窩,整個當時的敬拜系統,如果耶穌基督復活,可能隨時都會被瓦解..
所以對於這一班人來說,佔上這個話語的上風,確保耶穌不復活,確保耶穌復活的信息不被傳開,他們就可以繼續佔據著聖經話語權和解釋權,否則的話,他們就會失去位置,在猶太人當中,他們完全就沒有壓力了..
當他們做這些事情的時候,經文繼續說,上帝的工作並不是按照這些祭司長和法利賽人的心裡面來發展的..
接下來,在第28章開始這樣說:「安息日將盡七日的頭一日,天快亮的時候,末大拉的瑪利亞和那哥瑪利亞來看墳墓,忽然地帶震動,因為有主的使者從天上下來把石頭滾開,坐在上面..
他的相貌如同閃電衣服潔白如雪,看守的人因他嚇得渾身亂顫,甚至和死人一樣.」.
第五節:「天使對婦女說,不要害怕,我知道你們是尋找那釘十字架的耶穌,祂不在那裡,照祂所說,已經復活了..
你們來看安放主的地方,快去告訴祂的門徒,說祂從死裡復活了,並且在你們以先王加利尼去,在那裡你們又要見祂..
看,我已經告訴你們了.」由第一節到第七節,這裡說到,雖然那些法利賽人和祭司要把守著這個墳墓,.
但這個墳墓不是因為人的把守,耶穌基督復活的事件就不能夠來到它出現..
相反來說,因為上帝叫它復活,上帝的能力基本上超越了人,或者那些士兵把守的人的實力,.
上帝的能力與人的能力完全不足比較,上帝就在安息日那天就實踐了這個復活的大能..
經文裡面值得留意的是,在安息日將近,第七天的頭一天,相信是星期日,天快亮的時候,.
有兩個婦女特意要去看耶穌的墳墓..
值得留意,這兩個婦女當天還要包含耶穌的媽媽,甚至還有耶穌基督的門徒,.
特別是在耶穌基督受苦到被釘,甚至復活的日子裡面,.
一直陪著耶穌,與耶穌經歷著這一段苦路和復活,其實全部焦點既是放在耶穌身上,.
也是這一班婦女,特別是在說的莫大拉的瑪利亞和另外一個瑪利亞,他們目睹一切的..
上帝給予一個,又或者這兩個婦女,在耶穌面對種種的困難裡面,.
她們依然仰望著基督,在耶穌基督的受苦和復活的過程裡面,.
她們都是成為了第一個經歷復活的一些名不經傳的小角色,.
但是上帝卻是藉著她們,將這個信息由她們傳開..
不過在四個福音書裡面記載著婦女去看耶穌的墳墓的動機,.
幾個福音書的說法都有點不同,特別是在《瑪太福音》裡面,.
在這裡講到那些婦女去看耶穌,你有沒有想過為什麼她們去看耶穌的墳墓呢?.
她們大概都會知道,祭司長和法利賽人已經預備了兵丁在墳墓的洞穴門口,.
在墳墓的門口來把守,那個墳墓又找一塊大石頭來阻礙,.
所以那些婦女根本上沒有辦法進入墳墓裡面,再三的來高沒耶穌,.
或者望著耶穌的身體來緬懷一番..
在《瑪太福音》裡面的記載和其他的福音書記載裡面,.
她的說法有些不同,《瑪太福音》裡面在這個段落裡面特別強調,.
那個字就是漢字,去看,瑪太去導引讀者和當時第一代的信徒,.
帶領他們去看整個耶穌復活的事件,.
最初的時候這些婦女去看耶穌,大概如果那些法利賽人和祭司長都記得耶穌說過第三天將要復活的信息,.
同樣道理,真是清心的,真是期盼著基督復活的,可能就只有這兩個婦女,.
這兩個婦女大概都是帶著一份信心,當然整個運作是怎會不知道的,.
他們沒辦法看到神的能力怎樣,坊間裡面有不同的說法,.
在裡面那個果斯堡發出光等等,有很多人想像的神話在裡面,.

$^{81}$不過在經文裡面,她不是將這些描述放在《瑪太福音》裡面,.
《瑪太福音》只是說這兩個婦女相對於其他耶穌的門徒,.
他們去到墳墓裡面去重新,除了緬懷之外,更加重要的是,.
她記得耶穌說過,她要第三天復活,所以她就去到這個墳墓裡面,去看耶穌的墳墓,.
在經文裡面繼續說,當他們去看的時候,上帝的使者天使就顯現在他們當中,.
第六節說,她不在這裡,照她所說,已經復活了,你們來看看,.
安放主的地方,快去告訴她的門徒,說她從死裡復活,.
並且在你們以先往加利利去,在那裡你們要見她,可能我已經告訴你們了..
記不記得剛剛所說的,《瑪太福音》要導引我們去看東西,.
或者導引當代的那些兩個婦女,或者其他門徒去看,.
神的使者就將他的目光,從看墳墓的想法,就將焦點放在墓穴裡面,.
那塊石頭滾開了,叫她去看裡面耶穌不見了,耶穌不在墳墓裡面,.
在這個導引裡面,其實神的使者要告訴兩個婦女,.
基督耶穌已經復活了,你看不到祂的屍體了,.
你看祂所安放屍體的位置,不見了,祂不會在這裡,.
但是祂會在大家約定的地方,在加利利那裡,.
昔日耶穌和門徒約定在加利利那裡重新相見,.
如今祂告訴那些婦女,你去加利利那裡找耶穌,.
不是在這個墓穴裡面去找,.
以前我年輕的時候讀經,大家不理會四福音的記載,.
很多時候這兩個婦女就會成為批判的對象,.
因為他們覺得在空墳墓裡面找耶穌,當然找都找不到,.
然後就會覺得他們是一個沒有信心的人,.
但是馬太為他們平反,他們是最有信心的,.
他們去到這個地方,他們帶著耶穌基督的預言和應許,.
他們那裡很想見到復活的主,.
不過他們的焦點擺錯了,.
他們以為復活的主會在墳墓裡面,.
神的使者告訴他們,耶穌不是在那裡,.
耶穌在加利利約定的地方,.
所以神的使者就呼籲這兩個婦女去到加利利裡面等候耶穌,.
不單止這樣,也要將這個信息告訴耶穌的門徒,.
讓他們整個團隊和一些婦女去到加利利裡面集合,.
一起重新見證基督耶穌的復活..
在經文裡面,除了這些婦女見到神的使者,.
經文裡面也記載著剛才把守著墳墓裡面的那些冰釘,.
當這些冰釘見到神的使者,.
他的相貌好像閃電一樣,衣服潔白如雪,.
看守在那個墳墓裡面,原本看著死人,.
他們想著耶穌是一個死了的人,墳墓的人,.

$^{121}$忽然間變成了一個逆轉,.
經文說,像死人的那兩個冰釘,不是耶穌..
馬太丘是我們一脈的,原本活的人看著墓穴,.
看著死人,但是現在轉換了,死了的耶穌已經復活了,.
但是看守的人就好像死人一樣,.
因為經文裡面說,為什麼他們這麼害怕呢?.
因為他們害怕,害怕到好像死人一樣..
這兩個冰釘害怕什麼呢?.
或者這幾個冰釘看守著,.
大概是害怕一個超自然的現象,.
石頭竟然自動滾開,.
在他一直的職業生涯裡面都沒見過這樣的事..
另一方面,在他心靈裡面,他作為一個看守著墳墓的,.
他要負責看著這件事,.
現在這個被視為重犯,高度重視的犯人,.
他的屍體不見了,他們隨時都要面對大興問罪之師,.
他們自己都可能隨時都要死亡..
死亡可能就是他們沒辦法盡忠職守,.
在他的工作範圍裡面做好這件事,.
但同時他看到一個超自然的現象,.
耶穌復活了,神的使者顯現,.
是他一直始料不及的..
他的害怕帶著一份絕望,.
這幾個冰釘大概看到耶穌的復活,.
他心裡面害怕,.
但這種害怕卻不是驅使他走向上帝面前,.
只是令自己更加驚懼,更加無所適從..
在經文裡面同樣用害怕這個字來形容剛才說的顯現見到的婦女,.
那兩個瑪利亞,.
末大拉的瑪利亞和另一個瑪利亞,.
看到耶穌不在墓穴裡面,他們都很害怕,.
不過那個害怕是怎樣害怕呢?.
在第28章第8節裡面,.
婦女們急忙離開墳墓,又害怕,又大大地歡喜,.
跑去報給他的門徒..
在第8節裡面,說到那些婦女又是害怕,.
不過婦女的害怕和士兵的害怕是否不同呢?.
很不同,.
士兵的害怕是害怕自己快要死了,.
他們也不知道發生了什麼事,.

$^{161}$他們只知道將要面對一個重罰..
但是那些婦女的害怕是帶著快樂的,.
又驚又喜,.
因為耶穌基督昔日和他們應許的預言,.
竟然會出現,真的發生了,.
耶穌不在墳墓當中,.
昔日他們心裡面所痛心的,.
面對十字架上被鞭傷流血死身的,.
早兩天前也看到耶穌是一個歷歷在目,.
很痛苦的處境,.
這班婦女很心痛地陪伴著耶穌走上十字架,.
但是現在說他改喚了,.
他已經成為了一個復活的主,.
他心裡面有一份完全解釋不了,.
湧流出來的那份喜樂,.
他的害怕也代表著一種對神的敬畏,.
在他生命裡面未經歷過這樣的事,.
上帝將神跡中的神跡,.
呈現在這兩個婦女眼前,.
所以他們帶著一份快樂,.
一份歡喜,又驚又喜,.
就將這個訊息告訴耶穌基督的門徒知道,.
當他們一路走一路跑,.
跑回去,你看到兩個女人跑回去,.
她將這個訊息告訴門徒的時候,.
第九節,忽然耶穌遇見她們,.
說:願你們平安..
她們就上前抱著他的腳拜他..
第十節,耶穌對她們說:不要害怕,.
你們去告訴我的弟兄,.
叫他們往加里尼去,.
在那裡必見我..
第九第十節裡面,.
在他們這兩個婦女跑回加里尼,.
或者找門徒的途中,.
耶穌向他們首先顯現,.
耶穌第一樣跟他們說什麼呢?.
願你們平安..
一份經歷過苦難,.
再遇見主一份平安..

$^{201}$心裡面有一種,.
這幾句話充滿了份量..
很多時候我們今天說的平安,.
好像一個簡單的祝福語,.
但當你生命經歷風暴,.
當你生命裡面經歷了很多困難,.
在困難之後,.
你遇見基督,.
在你生命裡面經歷到上帝,.
平安這個字就變得很踏實,.
平安也代表著一種,.
完全逆轉的心靈狀態..
耶穌基督見到這幾個婦人,.
說回猶太人一直說的問安語,.
不過當說這個Shalom,.
願你們平安的時候,.
這些婦女就情不自禁,.
衝上去,走到耶穌的腳前,.
抱著耶穌的腳,像樹熊一樣,.
來去祈福敬拜耶穌..
耶穌對他們說,.
不要怕,不用怕,.
將這個信息告訴門徒知道,.
耶穌在這裡特別說,.
不是將消息告訴我的門徒知道,.
這裡特別改換了一個字,.
告訴我的兄弟,告訴我的弟兄知道,.
復活之後,.
連那個關係都重新被定義一樣,.
就是一個渙然一新的關係,.
一種渙然一新的心靈狀態,.
耶穌基督對他們說,.
叫婦女告訴這群兄弟,.
在加里尼,大家重新來相見..
復活的事件,.
改換了這群人對受苦的一些看法,.
復活的事件,.
改換了我們人生命的中心,.
在《瑪太福音》裡面,.
神的使者幫助這些婦女,.

$^{241}$將一個令人迷惑,.
或者不知道怎樣的一個焦點,.
慢慢放回在基督耶穌身上,.
由耶穌的空墳墓,.
帶領那些婦女看到真正的耶穌,.
按著耶穌的吩咐,.
讓這些婦女通知門徒和兄弟,.
在加里尼再一次見耶穌..
整個過程,.
瑪太就像一個導演一樣,.
他告訴你,你要看些什麼,.
他要告訴你,.
不要將你的焦點放在那些雜音身上,.
將你生命的焦點,.
將你的視線放在基督耶穌身上,.
由反證到真正的見面,.
整個的歷程,.
讓人經歷到一個復活的信仰,.
讓人經歷到耶穌基督復活的大能,.
在這裡我們值得再思考,.
今天我們的基督徒信仰,.
我們的信仰觀念,.
很多時候我們的焦點,.
很多時候會失焦,.
或者放錯了,.
很多弟兄姊妹,.
有時候他們會覺得生命很乏力,.
今天你打開Facebook,.
今天Facebook這段時間,.
最多出現什麼圖片?.
經常說全香港,.
好像所有店舖倒閉一樣,.
你同不同意?.
說到這個世界快要倒塌,.
甚至說到我們生命完全沒有盼望,.
你不要想這是媒體的工作,.
事實上今天我們進入了一種思潮,.
在一個叫後現代的思潮裡面,.
他講到這個世界完全沒有盼望,.
他講到一個很實在的狀況,.

$^{281}$他告訴你,.
你留意一下,.
你看看你現況,.
從你的現況裡面,.
你看看你有沒有盼望,.
後現代的其中一個特點,.
就是人失去了盼望,.
所以人的生活是短線的,.
按著你當下自己的感受是怎樣,.
你就用當下的感受來過你的生活,.
所以當你看到眼前的現象是這樣的時候,.
當你的盼望,.
當你人生的焦點放在這些現象裡面,.
你的人生是很容易失去了力量,.
失去盼望,.
所以今天這個世代裡面,.
他都帶著你看某些東西,.
帶著你看一些我們都經歷的是真正的現象,.
不過這種現象再推引下去,.
成為了一個絕望的人生觀念,.
看到一些失望的鏡頭,.
不停去看,.
以致我們活入失去盼望的一種經歷,.
另一種現象就是明天會更好的一種說法,.
這是什麼意思呢?.
我們不是很想今天的負面,.
有一個哲學家叫做希格爾,.
他講的是歷史是一個怎樣的進程,.
或者一個世界是怎樣走的呢?.
就是正反合,.
正和反,.
就是正面和負面,.
經過歷史的源流之後,.
大家就會有一個匯聚,.
就是一個叫做synthesis,.
在一個這樣的狀況之下,.
世界就會因為在不斷的正反合之後,.
就會朝向一個越來越至尊完善的一種情況,.
所以,.
希格爾的觀念看這個世界,.

$^{321}$好像很複雜,.
其實就是明天會更好,.
今天差一點,.
但因為後天就可以矯正昨天的錯誤,.
所以人類的歷史會不斷不斷進步,.
明天就會更好了,.
帶著一個比較樂觀的觀念來看未來的世界,.
不過,.
你活在地上,.
你就難免像後現代那些人說,.
是不是真的這樣?.
是不是明天會更好呢?.
有時候我們都不覺得,.
我們覺得世界好像走向一個越來越敗壞的狀態,.
甚至乎整個地球的氣溫,.
歷史裡面最高的氣溫,.
那十年,.
就是剛剛過去的那十年,.
整個地球就像火爐一樣,.
三月已經有三十多度的日子了,.
是不是?.
我們覺得,.
是不是真的明天會更好呢?.
基督教的信仰觀念,.
既不是明天會更好,.
也不是後現代所說的,.
因為我們的焦點,.
不是放在用現實當下的處境來定義,.
我們未來人生的走向..
基督教的觀念是什麼呢?.
基督教的觀念是說,.
整個世界,.
耶穌基督進入了苦難當中,.
這個苦難就是世界的現實,.
但是耶穌用一個不是越來越好的方式,.
而是用一個戰勝死亡,.
戰勝罪惡的那種死在十字架上的改喚,.
耶穌透過祂的死,.
更新了,改變了罪的權勢,.
那個盼望不是在於現實裡面,.

$^{361}$經濟環境,.
或者那些東西越來越好,.
我們就覺得心安理得,.
那個盼望是在於,.
因為我們知道這個世界的歷史,.
會由基督耶穌重新去改寫,.
即是說這個世界,.
有些東西越來越好像進入低點一樣,.
但是上帝就有能力,.
戰勝了一個很困難,.
苦難的狀態,.
突然之間,.
祂的復活就改喚了一切眼前的那種狀況..
基督徒的盼望,.
是放在基督耶穌的復活身上,.
我們欣著祂,.
即或眼前的那個狀況是怎樣,.
又如何?.
我們人生不是由外面的環境,.
來定義我們的生命,.
我們是由基督耶穌定義我們的生命..
所以弟兄姊妹,.
這裡面說得好像很複雜,.
不過我想說,.
這個復活,.
定睛在復活的主身上,.
其實告訴我們,.
我們的人生,.
即使你眼前有多少困難,.
即使你眼前很順境,.
那些都不足為道,.
不是代表著你接下來,.
你又寫包單,.
你未來的日子,.
有多麼好,.
或者會繼續差下去,.
反而我們生命的盼望,.
是由基督耶穌復活的生命所普捨,.
我們定睛在祂身上,.
我們的生命就可以得著改變..

$^{401}$我們經常說,.
有時信仰可以說得很抽象,.
但亦可以說得很具體,.
在我們接觸的弟兄姊妹當中,.
我們一班同工自己親眼見,.
我們見到有一位姊妹,.
剛剛回來教會的時候,.
她身體很軟弱,.
有時不笑,.
甚至有時說自己渾身骨痛,.
所以你見到她的笑容很低,.
剛剛回來教會,.
她自己又不太知道信仰是甚麼,.
不過當她信了主的時候,.
我們看著整件事,.
短短幾個星期,.
這位姊妹信了主的時候,.
她整個面容,.
整個人都不同了,.
我們覺得好像另一個人一樣,.
我們幾個同工在那裡,.
說她發生了甚麼事,.
為甚麼會這樣呢?.
以前有少許枯骨,.
有少許寒背,.
又沒有力量,.
但是當她信了主,.
她缺了自信主之後,.
隔了幾個星期,.
完全的人就渾然一新,.
我們都不知道發生了甚麼事,.
我再問羅牧師,.
羅記經常都笑著說,.
「哦,是的,信了主,.
信了主就很不同了,.
整個人,整個觀念都改變了.」.
我再跟這位姊妹聊天,.
「你真的很不同了.」.
她說「是的,耶穌在我心裡,.
好像覺得心靈定了很多,.

$^{441}$好像覺得整個人都不同了一樣.」.
弟兄姊妹,我不知道你今天的處境是怎樣,.
回憶你信主的時候,.
其實當你接受耶穌,.
你的生命就不同了,.
不過今天世界都有很多噪音拉走我們,.
我們又將眼目放在那些不太好的處境,.
又或者有一些,.
貿貿然在那裡叫囂打氣的說話,.
明天會更好,.
我們受著這些聲音,.
來將我們自己拉拉扯扯,.
但是真正令我們生命渾然一新的,.
不是外在的環境氣氛,.
單單仰望在這個復活的主身上,.
我們的生命就可以得著改變..
這班婦女經歷過耶穌的死亡,.
經歷祂的復活,.
又驚又喜,.
生命裡面脫胎換骨,.
成為一個新造的人..
我們看最後一個段落,.
11節至15節第28章,.
他們去的時候,.
看守的兵有幾個進城去,.
將所有經歷的事都報給祭司長,.
祭司長和長老聚集相議,.
就拿許多銀錢給冰冰說,.
你們要這樣說,.
夜間我們睡覺的時候,.
他的門徒來把他偷去了,.
倘若這話被巡撫聽見,.
有我們勸他保你們無事,.
冰冰授了銀錢,.
就照所祝福他們的去行,.
這話就傳說在猶太人中間,.
直到今天,.
在第11至15節,.
記載著一個荒謬中的荒謬,.
為什麼這麼荒謬呢?.

$^{481}$最初的時候,.
那些法利賽人和祭司長,.
他們想耶穌復活是一件什麼事?.
是一個誘惑,.
他們覺得耶穌復活是一個猶太教的異端,.
他們可能真心相信耶穌不會復活,.
不要把這件事傳出去就行了,.
把守著就安全了,.
但後來當他們知道,.
巡撫看到,.
或者那些冰冰將這件事告訴巡撫,.
告訴祭司長,.
耶穌復活了,.
看到那些使者的時候,.
這些人竟然,.
他繼續好像迷惑了自己一樣,.
他寧願相信自己的信念,.
也不相信眼前看到的那些冰冰所傳來的現實,.
更加重要的是什麼呢?.
他說到那些祭司長給錢收買那些冰冰,.
還有幫祂說好話,.
教祂說,.
在我們睡覺的時候,.
那些耶穌的門徒就將耶穌的身體偷去了,.
整件事是匪夷所思的,.
你想想,.
其實找他回來就是避免門徒去偷東西,.
但現在說到那些士兵睡覺,.
所以門徒去偷了,.
你覺得這班人真的做不了事,.
其實是有些教害的,.
其實你這樣說的時候,.
不過,.
這更加是說到,.
他們不相信,.
結果他們寧願捏造一些謊言,.
以一個謊言去掩蓋一個謊言,.
來證明他們一直的信念是對的,.
今天這個世界裡面很多人都是這樣的,.
寧願強辯,寧願否認,.

$^{521}$他們都不寧願相信耶穌基督復活的事,.
老實說,.
如果作為一個民眾去聽,.
復活的事是難以想像,難以相信的,.
是真的,.
不過冰冰所說的那些話,.
藏起來之後,.
大話再掩蓋大話,.
那些內容更加難以相信,.
更加荒謬,.
不過,.
這班人繼續選擇用這個方式,.
用謊言來掩蓋謊言,.
用謊言來捏造,.
以致他們繼續保持著昔日,.
他們自己話語權的權力,.
他們用不同的方式,.
來保住自己所說的話,.
證明自己所說的是對的,.
在經文裡面,.
很諷刺的就是,.
第十五節,.
冰冰受了這些銀錢,.
就照所吩咐他們的去行,.
這話就傳說在猶太人中間,.
直到今天,.
這句話如何理解,.
說什麼呢?.
就是說那些冰冰收了那些錢,.
就將那些大話說出來,.
不過,.
連冰冰說大話這件事,.
都是被傳說出來,.
簡單來說,.
你說秘密,.
你不要告訴別人,.
連這句話不要告訴別人,.
那句話都傳出去,.
告訴別人那句是是非,.
或者是大話,.

$^{561}$整個package,.
就將士兵說大話的整件事,.
傳說在當時的猶太民間當中,.
換句話來說,.
一銀兩面,.
這些人,.
證明自己的虛謊,.
證明自己所說的大話,.
某程度上也反證了復活事件的真實,.
上帝的工作真的很奇妙,.
這一班人原本要掩蓋著,.
不容讓復活的信息被傳開,.
但是到了最後,.
在第十一至十五節裡面說到,.
反而這件事更加被傳開,.
讓人對於復活的事件更加有興趣,.
更加想知道耶穌究竟是不是復活了,.
後來在哥林多前書裡面,.
再說到耶穌顯現給五百人看到,.
就是他的狀況,.
曾經有些人說,.
哇,耶穌復活是不是可信的,.
有些人曾經用過一些方法,.
特別是近代的心理學家,.
或者哲學家,.
會不會進入一種集體被騙的感覺呢?.
所以就傳出耶穌復活的信息,.
表面上很多證據都是成立的,.
不過有一件事是沒辦法推翻的,.
就是那些叫做group thinker,.
就是那些人抱持著同一個信念,.
看到和耶穌復活的幻象,.
可能大家心裡面投射出來的那些想法,.
很多時候是由一些人心裡面有很強烈的信念,.
所以有這樣的想法,.
但是這些畫像是不會出現在反對耶穌復活的人心靈裡面的,.
但是經文裡面記載著,.
那些看著耶穌,.
或者拒絕耶穌復活的信息的人,.
都見證著這樣的情景,.

$^{601}$耶穌復活是流言的不攻自破,.
今天我們去傳講這個信息,.
我們只是將我們所知,.
聖經所記載的東西,.
告訴我們身邊的人知道,.
可能今天很多人都選擇,.
寧願相信坊間裡面的一些說法,.
都不寧願相信更加可信的復活的事件,.
沒得逼的,.
不過作為基督徒,.
我們是很誠懇將這件事,.
將這個信息告訴身邊的人知道,.
有時候我不知道大家會不會,.
在講耶穌復活的事件,.
有時候說著說著會容易有些怒氣,.
會有辯論或者互教,.
我也有時候會這樣,.
因為我們覺得我們自己,.
我有一個責任,.
尤其是傳道,我是牧者,.
我有責任要跟別人解釋清楚,.
我要捍衛我的信仰,.
甚至我們要互教,.
將很多錯謀的東西,.
一一的去撇除,.
不過在講的時候,.
我發覺自己動了氣之外,.
我發覺越講,.
很多時候就按自己的智慧,.
人的想法去講,.
有一個前輩就跟我提一提,.
Anders,其實傳福音,.
我們不是做上帝的辯護律師,.
上帝不需要我們去辯護,.
我們的職責不是按著人的智慧,.
來辯解上帝復活的事件,.
真正將這個信仰帶出來的,.
是復活事件告訴我們的信息,.
我們只是如實將聖經呈現出來的面貌,.
告訴別人就足夠,.

$^{641}$不是我們去捍衛真理,.
而是真理保守我們,.
我再講多次,.
不是我們去捍衛真理,.
不是我們用我們覺得很有百般的智慧,.
很多的能力成為了上帝的辯護律師,.
我們怎樣做都做不到這個身份,.
是真理保守我們,.
在這個云云那麼多聲音的世代裡面,.
這個復活的真理在我們的內心,.
成為了我們生命的信息,.
成為了更新我們生命的內容,.
弟兄姊妹,.
今天講的整個片段,.
希望我們更加多些的確信,.
這個世代裡面很多人選擇相信方言,.
都不寧願選擇相信復活的主,.
但是對於我們每一個基督徒來說,.
我們親身經歷,.
我們明白,.
我們知道,.
復活的大能不單止成為了一個重要的教義,.
亦成為了我們生命重新詮釋,.
或者給我們面前人生有盼望的力量,.
我們就帶著基督耶穌這一份力量,.
走面前人生的路,.
復活的主與我們同在,.
聖經引導我們,.
請尋找本書的原文字幕.
謝謝大家.
\newpage



\section{路加福音 24:13-35}
\label{sec:JykGqUs7644}
\textbf{2024年4月14日崇拜直播｜陳明泉牧師｜十個深化關係的禱告-拿起餅祝謝的禱告｜路加福音24\_13-35}
\newline
\newline
連結: \href{https://youtube.com/watch?v=JykGqUs7644}{\texttt{https://youtube.com/watch?v=JykGqUs7644}} ~~~~ 語音日期: 2024-04-15
\newline
\newline
\hyperref[sec:g_Q6zBFvts0]{\small{< < < PREV SERMON < < <}}
~
\hyperref[sec:index_chronic]{\small{[返順時目]}}
~
\hyperref[sec:IOUGpe3_PlM]{\small{> > > NEXT SERMON > > >}}
\newline
\newline
路加福音 24:13-35
\newline
\begin{longtable}{cl}
\hline
\hline
章節 & 經文 (和合本修訂版)\\
\hline
24:13 & \begin{tabularx}{0.7\textwidth}{X} 同一天，門徒中有兩個人往一個村子去；這村子名叫以馬忤斯，離耶路撒冷約有二十五里。 \end{tabularx} \\ \\ \relax
24:14 & \begin{tabularx}{0.7\textwidth}{X} 他們彼此談論所發生的這一切事。 \end{tabularx} \\ \\ \relax
24:15 & \begin{tabularx}{0.7\textwidth}{X} 正交談議論的時候，耶穌親自走近他們，和他們同行， \end{tabularx} \\ \\ \relax
24:16 & \begin{tabularx}{0.7\textwidth}{X} 可是他們的眼睛模糊了，沒認出他。 \end{tabularx} \\ \\ \relax
24:17 & \begin{tabularx}{0.7\textwidth}{X} 耶穌對他們說：「你們一邊走一邊談，彼此談論的是甚麼事呢？」他們就站住，臉上帶著愁容。 \end{tabularx} \\ \\ \relax
24:18 & \begin{tabularx}{0.7\textwidth}{X} 兩人中有一個名叫革流巴的回答：「你是在耶路撒冷的旅客中，惟一還不知道這幾天在那裡發生了甚麼事的人嗎？」 \end{tabularx} \\ \\ \relax
24:19 & \begin{tabularx}{0.7\textwidth}{X} 耶穌對他們說：「甚麼事呢？」他們對他說：「就是拿撒勒人耶穌的事。他是個先知，在神和眾百姓面前，說話行事都大有能力。 \end{tabularx} \\ \\ \relax
24:20 & \begin{tabularx}{0.7\textwidth}{X} 祭司長們和我們的官長竟把他解去，定了死罪，釘在十字架上。 \end{tabularx} \\ \\ \relax
24:21 & \begin{tabularx}{0.7\textwidth}{X} 但我們素來所盼望要救贖以色列民的就是他。不但如此，這些事發生到現在已經三天了。 \end{tabularx} \\ \\ \relax
24:22 & \begin{tabularx}{0.7\textwidth}{X} 還有，我們中間的幾個婦女使我們驚奇：她們清早去了墳墓， \end{tabularx} \\ \\ \relax
24:23 & \begin{tabularx}{0.7\textwidth}{X} 不見他的身體，就回來告訴我們，說她們看見了天使顯現，說他活了。 \end{tabularx} \\ \\ \relax
24:24 & \begin{tabularx}{0.7\textwidth}{X} 又有我們的幾個人往墳墓那裡去，所發現的正如婦女們所說的，只是沒有看見他。」 \end{tabularx} \\ \\ \relax
24:25 & \begin{tabularx}{0.7\textwidth}{X} 耶穌對他們說：「無知的人哪，先知所說的一切話，你們的心信得太遲鈍了。 \end{tabularx} \\ \\ \relax
24:26 & \begin{tabularx}{0.7\textwidth}{X} 基督不是必須受這些苦難，然後進入他的榮耀嗎？」 \end{tabularx} \\ \\ \relax
24:27 & \begin{tabularx}{0.7\textwidth}{X} 於是，他從摩西和眾先知起，凡經上所指著自己的話都給他們作了解釋。 \end{tabularx} \\ \\ \relax
24:28 & \begin{tabularx}{0.7\textwidth}{X} 他們走近所要去的村子，耶穌好像還要往前走， \end{tabularx} \\ \\ \relax
24:29 & \begin{tabularx}{0.7\textwidth}{X} 他們卻強留他說：「時候晚了，天快黑了，請你同我們住下吧。」耶穌就進去，要同他們住下。 \end{tabularx} \\ \\ \relax
24:30 & \begin{tabularx}{0.7\textwidth}{X} 坐下來和他們用餐的時候，耶穌拿起餅來，祝福了，擘開，遞給他們。 \end{tabularx} \\ \\ \relax
24:31 & \begin{tabularx}{0.7\textwidth}{X} 他們的眼睛開了，這才認出他來。耶穌卻從他們眼前消失了。 \end{tabularx} \\ \\ \relax
24:32 & \begin{tabularx}{0.7\textwidth}{X} 他們彼此說：「在路上他和我們說話，給我們講解聖經的時候，我們的心在我們裡面豈不是火熱的嗎？」 \end{tabularx} \\ \\ \relax
24:33 & \begin{tabularx}{0.7\textwidth}{X} 於是他們立刻起身，回耶路撒冷去，看見十一個使徒和與他們正在一起的人聚集在一處， \end{tabularx} \\ \\ \relax
24:34 & \begin{tabularx}{0.7\textwidth}{X} 說：「主果然復活了，已經顯現給西門看了。」 \end{tabularx} \\ \\ \relax
24:35 & \begin{tabularx}{0.7\textwidth}{X} 於是，兩個人把路上所遇到，和耶穌擘餅的時候怎麼被他們認出來的事，都述說了一遍。 \end{tabularx} \\ \\
[1ex]
\hline
\hline
\end{longtable}
$^{1}$今日要讀的經文是.
路加福音第24章13至35節.
聖經新約第121至122頁.
這裡記載了兩個耶穌的門徒.
在二馬五師的路上的一些描述.
路加福音第24章第13節.
正當那日門徒中有兩個人往一個村子去.
這村子名叫二馬五師.
離耶路撒冷約有二十五里.
他們彼此談論所遇見的這一切事.
正談論相問的時候.
耶穌親自就近他們和他們同行.
只是他們的眼睛迷糊了.
不認識他.
耶穌對他們說.
你們走路彼此談論的是什麼事呢?.
他們就站著臉上戴著袖蓉.
二人中有一個名叫格留巴的回答說.
你做耶路撒冷作客.
還不知道這幾天在哪裡所出的事嗎?.
耶穌說.
什麼事呢?.
他們說.
就是拿撒勒人耶穌的事.
他是個先知.
在神和百姓面前.
說話行事都有大能.
祭司長和我們的官府.
竟把他戒去.
定了死罪.
釘在十字架上.
但我們素來所盼望.
要贖以色列民的就是他.
不但如此.
而且這是成就.
現在已經三天了.
再者.
我們中間有幾個婦女使我們驚奇.
她們清早到了墳墓那裡.
不見她的身體.

$^{41}$就回來告訴我們.
說看見了天使顯現.
說她活了.
又有我們的幾個人往墳墓那裡去.
所遇見的正如婦女們所說的.
只是沒有看見她.
耶穌對她們說.
無知的人啊.
先知所說的一切話.
你們的心信得太遲鈍了.
基督這樣受苦.
又進入她的榮耀.
豈不是應當的嗎?.
於是從媽西和眾先知起.
凡經上所指著自己的話.
都給她們講解明白了.
將近她們所去的村子.
耶穌好像還要往前行.
她們卻強留她.
說.
時候晚了.
日頭已經平息了.
請你和我們住下吧.
耶穌就進去.
要和她們住下.
到了坐直的時候.
耶穌拿起餅來.
祝謝了.
張開.
遞給她們.
她們的眼睛明亮了.
這才認出祂來.
忽然.
耶穌不見了.
她們彼此說.
在路上.
祂和我們說話.
給我們講解聖經的時候.
我們的心豈不是火熱的嗎?.
她們就立時起身.

$^{81}$回耶路撒冷去.
正遇見十一個使徒.
和她們的同人聚集在一處.
說.
主果然復活.
已經現給西門看了.
兩個人就把路上所遇見.
和抹餅的時候.
怎麼被她們認出來的事.
都述說了一遍.
頂展會平安.
今天我們一起來看盧家福音的這段記載.
這段記載是說耶穌基督復活之後.
耶穌復活.
某程度已經是戰勝.
不是某程度.
是完全戰勝死亡.
不過在門徒接受這個消息裡.
盧家福音用一個相當長的篇幅.
記載著跟隨耶穌的門徒.
他們面對著一些心路的掙扎.
這個掙扎甚至乎.
耶穌在他們面前也認不出來.
今天這段經文也比較長.
我們沒辦法完全逐個字去講.
不過這次主題是禱告.
十個拉近生命的禱告.
你會想.
哪裡有講耶穌的討文呢.
其實沒有什麼.
不過有一個很不著的地方.
就是耶穌的謝範祈禱.
這個謝範祈禱和一般的.
祝謝祈禱有些不同.
這個祝謝的祈禱.
令到整個門徒好像忽然間恍然大悟一樣.
今天我們就希望用這段的篇幅.
看看究竟這些門徒.
面對著耶穌基督復活的事跡.
他們心靈裡面的掙扎.

$^{121}$和耶穌跟他們同行的時候跟他們祈禱.
他們發生的一些心理歷程的改變.
這個歷程不只是門徒才有的.
我相信在我們生命當中也會有的.
希望我們今天用這段經文.
我們找到一些的亮光.
幫我們帶在我們的生活當中.
我們怎樣走我們生命的路.
我們先祈禱求上帝幫助.
父上帝願你在我們當中帶領我們.
願你的話語成為我們生命的激勵和亮光.
藉著今天我們一起來讀.
《路加福音》第24章.
我們求你親自對我們每一個內心說話.
幫助每位聆聽的弟兄姊妹他們聽得明白.
幫助孫說的說得清楚.
我們恭敬將我們的時間交託給你.
願你賜福 願你帶領.
獻上禱告奉靠基督耶穌得勝的聖名.
Amen.
我們一起看看這段經文.
在經文裡面.
第24章裡面.
特別先講到.
在耶穌和門徒一起進入一個村子.
這個地方叫做爾瑪五斯.
來耶路撒冷大概25里.
他們彼此談論所遇見的一切事.
談論的時候.
耶穌親自就近他們和他們同行.
只是他們的眼睛迷糊了.
不認識他.
在這裡面你會看到一個頗諷刺的地方.
耶穌和這班門徒.
這兩個門徒.
一起來同行.
這兩個門徒走的方向是一個怎樣的方向呢.
由耶路撒冷走去離開.
耶路撒冷進入一個.
現在已經是沒有辦法.

$^{161}$這些色經學家沒有辦法知道.
究竟是一個怎樣的地方.
這個地方叫做爾瑪五斯.
大概猜測.
今天有些說法是來耶路撒冷.
七公里的一個小鎮.
因為是小鎮所以沒有仔細地記載.
不過在這個路程裡面.
這兩個門徒帶著甚麼心情呢.
大概是帶著一份很低落的心情.
他們離開耶路撒冷.
本來耶穌帶著十二門徒.
或者其他的猶太人.
去到耶路撒冷.
去守這個逾越節.
不過今年的逾越節.
和昔日的很不同.
這個逾越節他們親眼目睹.
他們所信靠的那位彌賽亞耶穌基督.
釘身在十字架上.
當逾越節完了.
他們要回去了.
他們就離開耶路撒冷.
忽然之間好像這兩個門徒失去了生命的重心一樣.
帶著的一份失望和憂傷.
當他們離開耶路撒冷的時候.
耶穌和他們一起.
竟然他們的眼睛.
看不到東西.
迷糊了.
我不知道大家有沒有試過.
眼睛忽然之間蒙了.
就不是甚麼病態.
在聖經裡面說的眼睛迷糊.
通常是說那個信心.
或者我們叫做屬靈的核眼.
其實在路加福音裡面.
在第十一章都說到.
他說到如果我們的眼睛.
我們的眼睛就是身上的燈.

$^{201}$三十四節十一章.
他說如果你的眼睛瞭亮.
全身就光明了.
眼睛若分花.
全身就黑暗.
所以你們要省察.
恐怕你裡頭的光或者黑暗了.
若是你全身光明毫無黑暗.
就必全然光明.
如同燈的明光照亮你.
耶穌基督用這個眼睛的比喻.
或者說眼的視力.
很多時候就比喻我們的心靈.
我們對上帝的認識.
我們能不能夠發現上帝的足跡.
我們生命裡面有沒有智慧.
很多時候就用眼睛.
用你的視力.
來作為一個比喻.
去到這裡都是一樣的.
他講到兩個門徒.
耶穌和他們在一起.
他們竟然看不到他.
他們的眼睛好像看不到一樣.
耶穌是不是對這兩個門徒來說.
是一個陌生人呢.
在經文裡面.
他一直流露的就是說.
他知道了耶穌舊時的講論.
是鐵粉來的.
耶穌去那裡講論.
他就跟著去那裡.
而且耶穌所講的話.
是知道的.
甚至乎他們未必是一般的群眾一樣.
在聖經裡面.
因為形容這兩個人是耶穌基督的門徒.
跟隨過耶穌的.
又跟過他.
又聽過他.

$^{241}$他未是十二個門徒.
就是耶穌十二個門徒裡面.
那個核心的內圍.
但是怎樣呢.
他對耶穌不是完全毫無認識的.
不過在這個經文裡面.
就講到他們的眼睛迷糊.
不認識他了.
值得注意那個迷糊的字.
挺有意思的.
在路加福音裡面用的字眼.
和合本就譯他.
眼睛好像蒙了.
看不到東西.
但在原文裡面.
他用的字眼就是.
他們的眼睛好像被東西抓住了.
叫做抓住了.
我們不明白的了.
那個文化我們不會這樣講的.
看不到東西.
眼睛被東西抓住.
其實他所講的就是.
那種外在的力量.
令到這兩個門徒.
看不到耶穌.
就是跟隨過耶穌.
經歷過上帝的同在.
不過竟然耶穌和他們同行.
他看不到他.
是甚麼令到這兩個門徒.
眼睛被抓住呢.
接著我們看到.
經文就告訴你了.
他們就站在那裡念大仇容.
二人中有一個人名叫格留巴的回答說.
你在耶路撒冷作客還不知道.
這幾天在那裡所出的事嗎.
耶穌說.
甚麼事呢.

$^{281}$他們說.
就是那撒勒人耶穌的事.
他是先知.
在神和眾百姓面前.
說話行事都有大能.
祭司和我們的官府.
竟然將他解去定了死罪.
釘在十字架上.
但我們素來所盼望.
要屬以色列民的就是他.
不但如此.
而且這事成就已經三天了.
再者我們中間有幾個婦女使我們驚奇.
她們清早到了墳墓那裡.
不見他的身體.
就回來告訴我們.
說看見了天使顯現.
說他活了.
又有我們幾個人往墳墓那裡去.
所遇見的正如婦女們所說的.
只是沒有看見他.
這個經文裡面記載.
整個耶穌復活和婦女見到的事跡.
這兩個門徒其實很清楚.
不過另一方面.
耶穌和這群人同行.
耶穌沒有辦法.
這兩個門徒認出.
耶穌就某程度上就好像裝傻一樣.
什麼事啊.
發生什麼事啊.
其實發生在他身上的事.
但他透過一個好像問他們.
他們就將整個逾越節發生的事.
和耶穌復活的事.
這兩個門徒就清晰地表達出來.
不過說的時候.
這兩個門徒帶著什麼樣的狀態呢.
我提供給大家.
經文裡面就說到.

$^{321}$他們是站在那裡.
念上大雀壽容.
這兩個門徒.
走著走著的時候.
一說到耶穌發生的事.
走著回去.
站在那裡.
然後面就流露出一種極大的憂愁和痛苦.
他的痛苦是源自於.
這個是他們整個民族.
或者跟隨他的人.
一直以來的那種盼望.
盼望耶穌會為他們帶來一個新的狀況.
盼望耶穌成為一個宗教.
甚至乎政治的領袖.
將這一班.
生活在一種不容易狀況裡面的.
一班以色列人和這一班跟隨他的人.
都能夠有一種新的生活.
他們心目中所想的那種盼望就在於.
耶穌會推翻某個政權.
又或者至少.
他都成為一個宗教領袖.
不會繼續受到眼前的那種苦況.
不但他們期望的事情沒有實現.
而且耶穌還要死在十字架上.
耶穌在他們心目中成為了一個什麼形象呢.
成為了一個悲劇人物.
成為了一個受害者.
直著他的受害.
也告訴他們自己的心靈裡知道.
他們是一班絕望的人.
本來僅餘的希望.
現在連那個希望都沒有了.
死了.
所以這班人就帶著一份這樣的仇容.
不過這兩個門徒他再繼續說.
他說但是耶穌已經第三天死了.
然後那些婦女看到他復活.
又有天使顯現在他們當中.

$^{361}$現在耶穌的身體不見了.
其實整件事說到耶穌已經復活.
他們知道的.
他們也親身聽著婦女告訴他們.
看到耶穌的顯現.
和不見了耶穌的身體.
但是這一班的門徒.
他們都依然帶著一份愁雲慘霧.
離開耶路撒冷.
寧願選擇用絕望的方式來過生活.
他都不寧願相信婦女裡面所說的話.
如果你記得的話.
其實如果他們知道耶穌復活.
如果是跟隨他的人.
其實耶穌有一個約定.
就是在加利利再相見.
不過現在這些門徒.
他們選擇離開耶路撒冷.
選擇不會在加利利.
或者和11個門徒.
來到有再相遇的時間.
他們帶著一份絕望.
帶著在他們生命當中.
那種憂愁.
來離開耶路撒冷.
是什麼抓住他們的眼睛.
看不到耶穌呢?.
其實真正抓住的.
不是一些神秘的外力.
令他們的眼睛看不到耶穌.
是人的絕望.
是人寧願相信外面的世界.
那種絕望的那種悲觀的情緒.
都不相信基督耶穌的顯現.
和耶穌戰勝死亡的消息.
這裡我們稍微要思想一下.
我經常和弟兄姊妹分享.
可能我們今天的社會.
都瀰漫著一種很悲觀的情緒.
我不知道大家同不同意.

$^{401}$在我們的生活當中.
我經常說.
無論是你的媒體.
還是茶餘飯後.
我們香港人.
最影響我們的心情是什麼呢?.
通常是我們的經濟環境.
經濟好我們心情就好.
如果經濟不好.
或者我們賺錢艱難.
或者做生意的困難.
我們情緒就低落.
我們內心的感受.
我們對事物的看法.
很受著.
特別是香港這個社會.
很受著我們的經濟活動.
經濟繁榮度.
來決定我們的內心.
有時候我們基督徒也是這樣.
我們對前途有沒有信心?.
看經濟.
如果經濟好的.
我們信心就大一點.
經濟不好的.
我們的信心就會薄弱一點.
不過你留意一下.
如果我們久而久之.
我們繼續在這種文化土壤裡.
繼續去過生活.
其實我們和一個未信者.
是沒有什麼分別的.
我們的內心.
我們都是被那種外面的經濟環境.
來去左右了我們人生的喜樂程度.
又或者我們信仰裡面的踏實感.
我們受著經濟來去左右的時候.
我們的生命那份喜樂的基礎.
實在太薄弱.
甚至有時候我們活出來的生命.

$^{441}$和一個未信者.
都是沒有關係的.
沒有分別的.
求主幫助我們.
我們有時候可能都好像.
二馬五師路上的這兩個門徒一樣.
我們寧願相信.
外在世界.
即或我們生命裡面有一些.
基督耶穌的顯現.
在我們生命當中.
又或者.
在我們生命裡面.
知道上帝復活的這個訊息.
不過那個對我們來說.
覺得很離身.
覺得我們不是很貼地.
我們反而相信.
世界裡面給我們的訊息.
如果我們生命在這樣的狀況之下.
其實我們都某程度進入了.
一種叫做屬靈.
核眼的那種狀態.
我們的生命.
一路一路避著這個世界.
牽著它走.
曾經有一個靈修學家叫做.
Richard Foster.
早幾年來過香港.
屬靈生命禮讚.
富士德有個譯法.
來到香港.
他都明白香港某一些的狀況.
他曾經說過.
如果我們今天在這個世代裡面.
面對著幻變不定的世代裡面.
我們怎樣令到我們的生命.
有多些踏實感呢.
Richard Foster.
他提供一個答案就是.

$^{481}$你的教會.
或者在你的生命裡面.
要講另外一個的故事出來.
和這個世界某程度有少少離地.
因為界定我們人生命.
界定我們盼望.
就不是從這個世界的故事.
從這個世界的經歷.
來界定基督徒的生命力量.
反而是藉著基督耶穌的復活事件.
他甚至去到極致.
就是整個教會.
跟回教會年.
用那個週期.
來重新定義你每個星期.
裡面經歷的訊息.
來抗衡這個世界裡面.
對你的那份影響.
如果你能夠做到這樣.
你的生命就會多些力量.
而且你就不會受外面的影響那麼多.
當然這個是很極致的.
我們都不是那些禮儀教會.
我們都未必一定要.
好像他完全這樣說.
不過傅氏德說的.
亦都有一個很重要的道理.
我們今天的內心.
如果只是受外界的經濟環境影響.
我們受的只是我們的新聞.
訊息,媒體來訴訟我們內心的時候.
到最後我們難免會進入一種空虛.
和一種屬靈核硬的狀況.
我們的人生越走越無力.
我們的生命.
我們就沒有辦法經歷上帝而來的那份豐富.
不過這段經文.
不是叫我們去批判那兩個門徒.
說他們有多不濟.
或者讀的時候覺得我們有多無力.

$^{521}$經文裡面.
他特別說到耶穌基督面對著.
這兩個面對著屬靈,核硬.
或者有掙扎的人.
耶穌怎樣去幫他們呢.
是這段經文裡面一個很重要的重點.
我們繼續看.
耶穌怎樣去回應和幫他們.
25節.
耶穌對他們說.
無知的人.
先知所說的一切話.
你們的信心信得太遲鈍了.
基督這樣受害又進入.
他的榮耀豈不是應當的嗎.
於是從摩西和眾先知起.
凡經上指著自己的話.
都給他們講解明白了.
將近他們所去的村子.
耶穌好像還要往前行.
他們卻強留他說.
時間晚了.
日頭已經平息了.
請你和我往下住吧.
耶穌就進去.
要和他們住客.
耶穌就進去.
要和他們住客.
到了坐直的時候.
耶穌拿起餅來.
祝謝了.
張開了他們的眼睛.
他們的眼睛明亮了.
這才認出他來.
忽然耶穌不見了.
這段經文的記載.
是有很大的戲劇性.
為什麼戲劇性呢.
因為耶穌某程度上.
順著那兩個門徒的思路.

$^{561}$順著他們的想法.
一直追問究竟發生什麼事.
當門徒說完這些話之後.
耶穌突然間就好像.
從一個同行者.
變成了一個拉底的身份.
而且他第一件事.
對著那兩個門徒.
聽到他們這樣說之後.
耶穌說了一句不客氣的話.
無知的人才知道.
所說的一切話.
你們的信心.
你們的心信得太遲了.
這裡也很文雅的.
在翻譯.
你們無知的人.
不要以為耶穌很親暱地說.
傻子來的.
不是這個意思.
蠢人啊.
愚蠢的人啊.
蠢鬼啊.
就是這個意思.
這群人說得自己很優秀.
耶穌回應第一件事.
就是直斥其非.
頂他們.
站在拉底的角度.
說他們蠢.
然後蠢都不夠.
還要說他們遲緩遲鈍.
所以耶穌在這裡.
也挺直率的.
他不會跟他們客客氣氣.
聽到他們面帶仇容的時候.
那種看不到.
那種悲觀的情緒.
是因為他們看錯了.
不明白.

$^{601}$於是耶穌第一件事就糾正他們.
說他們是一群無知和遲鈍的人.
接著到了第26節.
是一個我們今天難以想像的一句話.
基督這樣受害又進入他的榮耀.
豈不是應當的嗎?.
我們對於這個世界裡的苦難.
我們今天面對著生命裡的困難.
我們第一件事會覺得這個世界是荒謬的.
在這個門徒的視角裡.
甚至面對著耶穌基督在十字架上受害.
被釘十架.
那種受害者的感受是很強烈的.
不過耶穌說.
耶穌這樣的受害.
基督這樣受害.
是應該的.
是對的嗎?.
這句話引起他們重新思考.
為什麼上帝的計劃會這樣.
接著耶穌就從創世記開始說.
摩西到眾善之.
基本上是整個舊約聖經.
和他溫習一次.
這段路真是一個溫書之旅.
耶穌由摩西到眾善之.
指著基督耶穌的召命和耶穌基督的工作.
都和這兩個門徒講解一次.
講到他們明白了.
但是一邊明白一邊點頭.
其實是不是真的明白呢?.
還沒有完全明白.
直到住在村莊裡.
看到耶穌還要繼續走.
這兩個門徒強留他.
和他們一起.
要耶穌在這裡過一晚夜.
再繼續走.
過夜之前大家吃東西.
吃一些很簡陋的麥餅.

$^{641}$這樣來做.
麥餅的時候.
才看到耶穌的祝謝禱告.
接著才重新發現.
眼前那個就是.
一直他們所仰望的基督耶穌.
是一個復活了的上帝.
在經文裡.
他用一個很戲劇性的描述.
他講到這一幫門徒.
其實他一直跟隨耶穌.
但是他從來沒有想過.
或者他都不明白聖經在說什麼.
他們只是帶著他們自己.
一直的想法.
來渡過他們的信仰.
他們的盼望只是建基於.
他們自己的期望上.
他們期望耶穌實踐了某一些功能.
或者他們心目中期盼的形象.
但是從來沒有想過.
基督的召命.
或者耶穌基督來到這個世界.
他就要用受苦的形態.
以致他要完成上帝的救贖工作.
由不明白到講解.
到開竅.
整個這段路.
就是這兩個門徒所經歷的.
在這裡我們都值得再去想一想.
我們有時候會不會都像兩個門徒一樣呢?.
我們期望耶穌基督.
成為了某一種形態.
我們生命裡期待的成功.
期待某一些生命裡我們所想得到的東西.
我們求祂.
但是到最後落空了.
又或者我們生活當中.
不僅得不到我們所期盼的東西.
因著信仰有時候.

$^{681}$我們甚至會失去了某些東西.
我們會帶著一份愁容.
帶著一份悲觀的情緒.
走面前人生的路.
最經典的例子.
有時候我們牧羊弟兄姊妹.
都有這樣的經歷.
有些弟兄姊妹期望.
可能一段感情.
會為她帶來新的開始.
或者期望一個新的工作崗位.
她期盼她可以在工作崗位裡見證上帝.
又或者她轉換一個新的處境.
她期盼耶穌基督.
她心裡有這樣的領袖.
她覺得那個領袖就是最終的絕對.
到最後原來她的領袖和現實有一段距離的時候.
她就帶著一份很悲觀的情緒.
甚至有一份憤怒.
有一份不甘心.
又或者用受害者來看基督耶穌.
也都看她自己的生命.
這種心態.
令到這兩個門徒進入了一種屬靈核硬的狀況.
但基督耶穌沒有嫌棄這兩個人.
進入一種這樣的狀況.
耶穌雖然直斥其非.
說他們是愚蠢和心靈遲鈍.
但耶穌和他們重新溫和書.
重新指著自己的召命.
將基督的那種由受害到復活的.
整個聖經指著他們生命的歷程.
重新講一次給他們聽.
講完一次給他們聽之後.
他們開始明白.
直到在那個祈禱裡.
在生活日常的謝飯祈禱.
祈完禱之後.
當這些門徒竹蔗.
撒開.

$^{721}$遞給他們的時候.
他們就想起昔日的那個夫子.
想起昔日五餅二魚.
想起逾越節.
那個為門徒洗腳的基督耶穌.
這個餅和這個禱告.
就好像成為了一個神聖的頓悟的那一刻.
讓他們重新看回自己的生命.
看回基督耶穌基督的工作.
整件事其實有一個思路.
耶穌直斥其非.
然後講解他的工作.
在講解的過程不是遙遠發功.
講解的過程其實耶穌和他們同行.
同住.
然後同住在一個地方裡.
再靠近的為他們竹蔗禱告.
眼睛就明亮了.
整件事講到.
那個禱告其實本身.
沒有一個神秘的功能.
那個抹餅祈禱不是說有什麼神秘.
不過那個抹餅的禱告.
卻是好像成為了之前.
耶穌和這兩個心靈合眼的人.
同行講解了的一些預備功夫.
到抹餅那一刻.
一下子清醒了.
開竅了.
讓他們的生命可以重新.
好像重新找回基督耶穌一樣.
這樣的一個歷程.
其實我們每一個信徒都會經歷的.
不過有時候我們不覺的.
我們營營役役的生命裡.
如果我們所有的東西都變得很.
生活日常.
所有的東西都是循規蹈矩.
這個是好的.
不過我們的人生如果所有東西都是.

$^{761}$都是循規蹈矩的.
大概我們人生.
活到八九十歲我們都是一樣的.
但我們的人生原來不是這樣走的.
我們的人生原來在我們很多規律背後.
會有一些類似火花的事件.
爆一爆.
而這些火花的事件.
就好像幫我們重新開竅一樣.
說說經歷.
大家可能會明白多一點.
我們每個星期三都有一個網台.
誠心所願.
我們訪問過很多不同的嘉賓.
近來很多嘉賓跟我說.
做這個誠心所願.
除了看得很好看之外.
他說他接受訪問的時候.
從一個被接受的人.
其實我一直在做的那些事.
我不知道有什麼這麼好訪問的.
他跟我說.
其實沒什麼特別.
我們每個人都有不同崗位.
大家努力拼勁地做自己想做的事.
每天都是這樣.
生活上我都是一個平凡的人.
不是什麼很值得跟別人分享.
不過在訪問裡面.
你給我們一些問題去想.
然後在一個談話的過程裡面.
說著說著.
我才發覺原來上帝真的有奇妙的帶領在我們生命當中.
如果不經歷過這個講解.
和回顧的歷程的時候.
我們的生命還是會繼續這樣去.
但如果有些回顧重新去看的時候.
我發覺我們看整件事都很不同.
還有看到上帝在我們生命裡面的這些足跡.
我們其實.

$^{801}$反過來說.
其實我們每個星期做.
我都不知道我們有沒有果效.
我都不知道有多少頂紫媒體.
有時候收視率會高一點.
有時候收視率會低一點.
但是我們都在如常運作.
不過在日常生活裡面.
突然之間有一些回顧.
或者有一些火花出現的時候.
就好像重新幫我們重新看回.
我們自己的人生.
我們自己的歷程.
藉著這些歷程.
我們重新認定上帝在我們生命當中的帶領.
其實我們人大了.
我們經不起人生裡面有很多驚天動地的改變.
我們都不會像年輕人一樣.
追求那種很火熱.
某一種令我們刻骨銘心的一些經歷.
因為我們知道人大了.
我們的心臟負荷不會來.
我們覺得沒有辦法來做.
而我們的生命也不是一定要這樣.
每一天都轟轟烈烈.
我們才經歷到上帝.
反而有時候在生活日常裡面重新回顧.
讓上帝罵一罵.
重新看回上帝的工作.
上帝在我們生命裡面的足跡.
重新整理的時候.
那些細水長流.
重新回顧.
反而幫我們更加有力量.
帶著一個insight.
我們就能夠有力量走面前的路.
耶穌基督在這裡的竹蔗禱告.
其實就好像畫龍點睛一樣.
其實最重要的是.
一直以來的那種陪伴,同行.

$^{841}$耶穌的講解,明白.
而這個祈禱就好像一個開竅的動作.
讓我們重新認定上帝的經歷.
同樣道理.
今天我們生活裡面有很多.
每一天都會藉飯祈禱.
早上起床你會祈禱.
臨睡前你會和上帝聊幾句話.
傾心吐意.
如果這些東西pattern(圖案)了成為生活日常.
有時候我們很容易變成某一種的習慣.
這些習慣是好的.
不過這些習慣如果只是令我變得恩秦.
我們的生命帶來不了改變.
唯有是那種習慣.
再加上生命裡面突然之間的頓悟.
重新明白上帝的感召.
或者上帝的工作的時候.
我們的人生才會有所成長.
兩個門徒就是這樣來見到耶穌基督.
在這樣的竹蔗祈禱之後.
耶穌就不見了.
當耶穌不見了的時候.
這兩個門徒是不是就立刻又再帶落呢?.
經文裡面再繼續說.
三十二節.
他們彼此說:在路上,他和我們說話,給我們講解聖經的時候,我們的心豈不是火熱嗎?.
他們就立時起身回耶路撒冷去.
正遇見十一個使徒,和他們同人聚集在一處.
說:主果然復活,已經遣給西門看了.
兩個人就把路上所遇見和抹餅的時候,怎樣被他們認出的事,都述說了一遍.
經文裡面,雖然說到耶穌不見了.
但是耶穌基督所說的.
聖經指著他所說的話.
仍然在他們的內心.
這種開竅的體會.
他們再想一想.
我們在聽的時候,我們的心不是很火熱嗎?.
現在耶穌在我們面前消失.
但是我們更加認定上帝的工作.

$^{881}$於是這兩個門徒有什麼反應呢?.
原本他們帶著愁容,悲傷離開耶路撒冷,進入二馬五世.
但是那時候已經日落平息了.
他們都急不及待.
由日落的時間趕回耶路撒冷.
要和十一個使徒重新來到會合.
要把他們所見所聞,所經歷的告訴門徒.
之前他們是愁雲慘霧.
但是現在他們帶著火熱的心.
去到耶路撒冷.
重新要把遇見基督的事,再述說一遍.
遇見上帝,找到上帝.
我們心靈不再黑眼.
其實就是一個火熱的經歷.
在那個經歷裡面,未必有很轟天動地的.
都是峰迴路轉的.
不過未至於當氣回腸.
很大場面,不是那樣的.
生活日常,但是明白聖經.
明白人生的時候.
心裡面就有這份火熱.
我近來呢.
我的弟兄姊妹說我有些不同了.
不是髮型不同了,不是戴了眼鏡不同了.
而是我讀這段經文的時候.
我發覺對我內心有一份很大的.
有一份很大的那種體會.
為什麼這樣說呢.
我在這裡回顧.
我忘記我在黃浸有多少年了.
19年了.
在這個群體裡面.
我在想,19年我有什麼做過呢.
好像每個星期都上班,又下班.
我做過什麼呢?.
我有什麼建立了呢?.
上帝,你呼召我.
究竟我的生命是怎樣的呢?.
不過呢.
當我見到很多弟兄姊妹.

$^{921}$聽到上帝的話語.
生命是忽然間有一些轉變.
見到生命軟弱的人.
他奉主名祈禱.
即使身體軟弱.
但心靈有平安,有力量.
見到有些弟兄姊妹.
可能在過去的人生裡面.
都很多東西卡住.
但是因為信仰的緣故被改變.
在小組當中,我進入那些小組裡面.
經歷過又疫情,又分組.
就是很亂的日子.
但到了將要分組.
大家將要重組了.
大家會覺得.
過去那段日子一起小組.
實在太過美滿.
有時候都是不堪入目的.
原來都成為了生命一個很重要的回憶.
這些一切.
我發現忽然間內心很踏實.
我發覺原來牧養.
信仰生命.
我們都不一定追求一些.
令人蕩氣回腸的一些.
一些很大場面的經歷.
反而是生命裡面.
逐漸逐漸地改變的時候.
反而經歷到那份信仰.
那份豐厚感和踏實感.
這段經文某程度是說.
耶穌基督復活.
顯現在門徒當中.
但很多人形形異異帶著悲觀.
經歷不到復活.
或者見不到耶穌.
但當我們重新看到上帝的工作.
看到教會的意義.
看到自己生命的路線.

$^{961}$走的一個走向.
重新找回自己.
走在這裡.
去到今天這個狀況.
知道未來下一站要怎樣走的時候.
心靈就變得很踏實.
弟兄姊妹.
不知道你信主有多久.
人生會不會經歷過迷茫.
有沒有經歷過失望.
又或者有時.
有些不堪入目的經歷.
甚至遺憾.
當我們經歷這些的時候.
我們想分開它.
我們想忘記.
不過.
生命的路其實不是往走的.
當你回顧看回.
你自己所走過的路.
看回基督耶穌那種受害.
進入荒謬世界.
用祂的死.
用祂的復活.
戰勝了生命裡面的荒謬的時候.
我們人生走過.
所謂的冤枉路.
其實一點都不冤枉.
反而是封口我們的生命.
讓我們生命有更加多的內容.
以致我們更加懂得.
怎樣去走人生的路.
這段經文特別是送給那些.
可能人生裡面.
有時走著走著都不知道怎樣走.
有時都會灰心.
失望的人.
希望我們今天.
將我們自己的生命.
代入在那兩個門徒當中.

$^{1001}$耶穌沒有放棄我們.
祂會罵我們.
但祂更加會與我們同行.
顯現自己.
當我們找到祂的時候.
祂生活日常的禱告.
就成為了我們開竅的一部份.
當我們找到耶穌.
我們就有力量.
火熱地走人生的路.
願主繼續幫助我們.
帶領我們.
\newpage



\section{列王記下 6:14-17}
\label{sec:IOUGpe3_PlM}
\textbf{2024年4月21日青少主日崇拜直播｜馬嘉言姑娘｜讓青少年親眼看見神｜列王記下6\_14-17}
\newline
\newline
連結: \href{https://youtube.com/watch?v=IOUGpe3-PlM}{\texttt{https://youtube.com/watch?v=IOUGpe3-PlM}} ~~~~ 語音日期: 2024-04-21
\newline
\newline
\hyperref[sec:JykGqUs7644]{\small{< < < PREV SERMON < < <}}
~
\hyperref[sec:index_chronic]{\small{[返順時目]}}
~
\hyperref[sec:K_X10pNASkM]{\small{> > > NEXT SERMON > > >}}
\newline
\newline
列王記下 6:14-17
\newline
\begin{longtable}{cl}
\hline
\hline
章節 & 經文 (和合本修訂版)\\
\hline
6:14 & \begin{tabularx}{0.7\textwidth}{X} 王就派遣車馬和大軍往那裡去，夜間他們到了，圍困那城。 \end{tabularx} \\ \\ \relax
6:15 & \begin{tabularx}{0.7\textwidth}{X} 神人的僕人清早起來出去，看哪，車馬軍兵圍困了城。僕人對神人說：「不好了！我主啊，我們該怎麼辦呢？」 \end{tabularx} \\ \\ \relax
6:16 & \begin{tabularx}{0.7\textwidth}{X} 神人說：「不要懼怕！因與我們同在的比與他們同在的更多。」 \end{tabularx} \\ \\ \relax
6:17 & \begin{tabularx}{0.7\textwidth}{X} 以利沙禱告說：「耶和華啊，求你開他的眼目，使他能看見。」耶和華開了這年輕人的眼目，他就看見了，看哪，滿山有火馬和火焰車圍繞以利沙。 \end{tabularx} \\ \\
[1ex]
\hline
\hline
\end{longtable}
$^{1}$今日讀的經文是在列王記下第六章.
在大家手上聖經的四百六十三頁.
列王記下第六章第十四 二十七節.
「王就打發車馬和大軍往那裡去.
夜間到了圍困那城.
神人的僕人清早起來出去.
看見車馬軍兵圍困了城.
僕人對神人說.
「愛在我主啊.
我們怎樣行才好呢?」.
神人說.
「不要懼怕.
與我們同在的比與他們同在的更多」.
以利沙圖告說.
「耶和華求你開這少年人的眼目.
使他能看見」.
耶和華開他的眼目.
他就看見滿山有火車火馬圍繞以利沙.
弟兄姊妹平安.
我是實習神學生.
我叫馬嘉賢.
大家也可以叫我馬姑娘.
或者清少….
清少部坐在哪裡?.
那裡.
少年部.
他們叫我卡媽.
或者大家可以叫我英文名Karen.
我主要負責少年部的侍工.
所以我很多時候都會在四樓的時間多一點.
平日就會在少崇那裡.
又或者星期六的時候.
就會參加中學星期日的大專團體等等.
今天藉著清少主日.
有機會上來六樓.
和大家一起聚會.
一起敬拜.
和可以和大家分享上帝的訊息.
我們開始.
講道之前.

$^{41}$不如我們一起祈禱.
親愛的天父 多謝你.
不論我們是幼年 青年 成年.
你對我們的愛都從來沒有改變.
上帝.
你一代又一代的門徒.
這一刻都同心聚集在你的十字架面前.
祈求聖靈運行在我們每一個人當中.
願孩子孫講清楚.
願弟兄姊妹聽得明白.
共主一命祈禱 阿們.
麻煩Viola.
這樣就讀了今天的經訓.
今天的經訓是記載在聶王紀下的六章十四到十七節.
其實是講述以利沙和他的門徒的一番說話.
正確來說 如果仔細一點.
其實是師父在回應徒弟一個求救的對話.
如果我們追溯一下.
究竟為什麼這段對話會發生的時候.
我們可以看回六章八節.
是在講述亞蘭王與與式的人爭戰.
亞蘭族是一個什麼族呢.
其實亞蘭族是羅亞的兒子.
閃所生的第五個兒子.
也就是說亞蘭族是羅亞的孫.
亞蘭族最初佔據的領土.
就在右法拉底河的東邊.
如果大家現在看回Power Point的話.
就會看到右邊用紅色圈圈住的位置.
當時亞蘭族有個首都.
這個首都就是大馬色.
由於在聶王紀的時代.
其實已經分成了南北兩國的時間.
所以當時的亞蘭王.
其實是與北國的約蘭王開戰.
亞蘭王的計劃.
其實他的軍事計劃.
就是主要想著在某幾個重要的據點安營.
然後與以色列軍進行突襲.
但是神人以利沙.

$^{81}$因為他是先知.
所以他就預先知道了亞蘭王的軍事行動.
第九節和第十節大家看回都會看到.
以色列軍因為有以利沙的提醒.
所以就成功防範了亞蘭軍的突襲.
接二連三的突襲都失敗.
令到亞蘭王心裡面驚異.
如果我們看回呂鎮忠的譯本的話.
我個人來說覺得譯得比較傳神一點.
他將亞蘭王的心情譯作了.
因了這事.
亞蘭王心裡十分激憤震盪.
亞蘭王將失敗歸咎於一個原因.
他認為當中有內奸.
所以亞蘭王就再次襲他的神僕.
然後他就很生氣地去問他的神僕.
我們當中到底誰出賣本王.
到底誰幫助以色列王.
在王大路的情況下.
其中有一個神僕回答亞蘭王.
我王你的僕人裡面其實是沒有內奸的.
我們行動被人洩漏了.
是因為以色列國當中有先知以利沙.
以利沙將王在房裡說的計劃.
預先跟以色列王通風報信.
然後他的神僕裡面再有人跟亞蘭王說.
這個以利沙我們知道他在哪裡.
他在多丹.
多丹這個地方不是經常在聖經裡面出現.
如果我們看回我給了一個地圖告訴大家在哪裡.
其實他就位於薩瑪利亞以北的17公里.
和在士劍大概22公里.
那是怎樣呢.
那就是你看到三個圈圍著的地方就是很近的距離.
那就是說多丹其實當時以利沙的位置.
就是跟薩瑪利亞非常近的.
當亞蘭王知道以利沙在多丹之後.
他就決定派發大軍前往多丹.
大家可以留意.
其實以利沙是一個手無寸鐵的先知.

$^{121}$最多他身邊有一群門徒跟著他圍繞著他.
但亞蘭王去捉拿這個手無寸鐵的先知和門徒的時候.
他不是派幾個簡單的精銳部隊去捉拿以利沙.
反而亞蘭王是派出了可以足以圍困整個多丹城的馬車和士兵.
從亞蘭王派出的這一隊兵力我們就可以看到.
其實亞蘭王是非常非常重視今次的軍事行動.
我們中國人也有一句說話.
射人要怎樣.
要先射馬.
所以亞蘭王就清楚知道.
如果他要成功擊敗以色列的話.
就必須首先去解決以利沙這個先知的問題.
所以亞蘭王就連日趕路.
趁著夜間圍困多丹城.
然後就去到我們剛才所讀的那段經文.
我們嘗試代入一下僕人在十四節裡面所想.
或者他那時候所身處的場景.
對於僕人來說.
第二天清早起床的時候.
眼前的畫面其實也很震撼.
前一晚.
我才跟我的朋友.
我的同伴.
我的師父.
安然生活.
安穩入睡.
但是一覺醒來之後.
所有的敵人.
都衝著師父而來.
整個城都被敵軍圍困著.
我是一介平民來的.
我只是想跟著師父學習.
為什麼突然間變得連命都難保.
插翼難飛.
所以僕人當時.
就很怕跟以利沙去求救.
他就說 死了 死了 死了.
死了 我的主啊.
我們現在怎麼辦.
我們現在應該要怎麼辦.

$^{161}$怎麼辦.
我相信我們每一個人.
其實在面對艱難的時候.
都一定會問這個問題.
一個How的問題.
不過成年人和年輕人.
去面對困難的態度.
有時候就有點不同了.
成年人去面對問題的時候.
很多時候就會傾向.
一個人想很多辦法.
這個問題我應該怎麼處理.
有A的方法 有B的方法.
年輕人他們都同樣想很多.
不過他們所想的.
很多時候圍繞著問題去想.
就是不停想那個問題.
然後將那個問題小視化大.
然後再去轉牛角尖.
就好像以利沙的僕人.
當他面對困境的時候.
他單單只看到自己.
被亞蘭軍圍困的困境.
僕人就是被困境的表面.
蒙蔽了自己的雙眼.
情況就好像.
面對一大塊欄杆.
顏色混在一起.
嘩 很雜亂 很混亂.
圖上面這個都真的很混亂.
如果年輕人看到這塊欄杆的時候.
他們第一個反應應該就會說.
嘩 很大塊.
嘩 很難拆.
嘩 很難搞.
他們缺乏了那個冷靜.
去梳理那塊欄杆的能力.
既然都缺乏了梳理的能力.
就更加不要想有沒有解決的方法.
所以結果面對一大塊欄杆的時候.

$^{201}$這麼大的問題.
年輕人就會好像僕人一樣說.
他就會跟別人求救.
其實告訴大家.
他不只是想跟別人求救.
他是想最好你幫他搞定.
幫他處理.
把那些鬍子撿走就最好了.
這個是他們心底裡的說話.
所以年輕人的看見就是.
看見欄杆的表面.
然後再找人求救.
希望有人幫自己去解決問題.
這一份看見.
聽起來好像都有點清潔.
有點膚淺.
不過.
這個就是青少年.
這個就是青少年人所擁有的特質.
有時候其實我自己都會想一下.
我都會想.
如果我年輕的時候.
有現在的智慧.
雖然現在都不是很多智慧.
那麼.
我應該會撞少一點板.
我應該不會.
頂撞我媽媽.
頂撞得這麼厲害.
我應該說話會比較圓滑.
不要這麼掘雷鎚.
不知道大家各位.
會不會有時候都好像我有這樣的回想呢.
可能我們都曾經年少輕狂.
可能我們都曾經狂妄自大.
我們都曾經跟這些年輕人一樣.
只看到事物的表面.
如果我們昔日.
都曾經有這樣犯過錯.
有這樣的特質的時候.

$^{241}$今天我們面對下一代.
如果他只可以看見表面的時候.
其實我們可不可以對他寬容一點呢.
我們可不可以不要直接批評他們.
甚至可以寬容到有時候讓他們有撞板的機會.
年輕人這些熟悉的看見.
在我角度.
因為我做青少年牧者.
我就覺得.
其實無傷大雅.
都不是真的那麼重要.
只要不是殺人放火.
去到無可挽救.
有時候我都會被他們撞板.
因為我自己會認為.
更加重要的是.
怎樣在他們.
看到他們有問題的時候.
或者只是有熟悉的看見的時候.
我自己可以成為那個橋樑.
去啟發他們有熟悉的看見.
不過在搭橋之前.
當然要先打好根基.
所以其實.
建橋之前就要有根基.
以利沙對僕人的回應.
就是展示出一個.
很重要的怎樣建立根基的關鍵.
16節以利沙是這樣回應的.
他說.
不要懼怕.
與我們同在的.
比與他們同在的更多.
坦白說.
其實以利沙這句話完全是答非所問的.
因為僕人是在問一個.
How?.
怎樣解決的問題.
但是以利沙.
沒有提供任何解決的方案.

$^{281}$如果我們多看一些.
列王記夏經文的時候.
回到第二章的時候.
其實打從以利沙升天之後.
聖經就記載了很多.
關於以利沙所行的神蹟.
關於以利沙很出名的神蹟.
就是他幫助書唸婦人的兒子復活.
他醫治了乃萬將軍的大麻瘋等等.
聖經記載這些神蹟.
其實是為了刻意要建立以利沙.
他是有承接以利亞先知的積分和權柄的.
所以.
以利沙的能力.
既然他已經承繼了積分和權柄.
他當然可以呼喚上帝的火車火馬.
一次過殲滅了亞蘭軍.
其實這個方法.
既可以即時解決當下的困難.
又可以回應了僕人的訴求.
絕對是一舉兩得的方法.
為什麼以利沙沒有即時實行這個神蹟呢?.
他首先要用十六節的經文去回應徒弟呢?.
因為以利沙在這個時刻.
不是以解決困局為優先.
以利沙在這個困境裡面.
不是以解決困局為優先.
他是利用這個困局去建立徒弟的屬靈生命.
不過.
有時候.
你不是說想建立年青人的屬靈生命.
是可以建立的.
我也知道.
你要年青人去聽我們說話的時候.
其實很多時候都需要以關係先行.
大部分的關係建立.
主要都是透過兩件事.
第一.
就是在同一個場景裡面共處.
第二就是在情感上有連結.

$^{321}$我們回看以利沙和僕人.
其實兩者同時都是面對著被敵人圍困的局面.
簡單來說.
師父和道弟.
兩代人.
就是在同一時間同一地點.
面對同一個困境.
當然我們都是經歷了兩代人.
可能不止兩代.
可能三代甚至四代在同一堂取得在一起.
往往因為我們面對同一件事的時候.
有不同的角度去詮釋.
所以.
幾代人之間就會很容易產生矛盾.
產生一些代溝.
我就找了一個圖片給大家看一下.
這個數目字就躺在地上.
左手邊的那個人就說.
啊 這是六字.
右手邊那個就說.
不是啊 這是九字.
其實.
如果兩個人都只是堅持幾件.
一個就一定說.
是六字.
一個就說.
不是啊 是九字.
其實他們永遠都不會看到.
這個數字是一個怎樣的數字.
或者不會看到為什麼對方會這樣說這個數目字.
所以如果兩個人堅持幾件.
不願意轉換立場的時候.
其實是一直不會有討論的結果的.
不過以莉莎和僕人去面對相同的處境的時候.
她沒有恃著自己的屬靈經驗比較多.
她沒有因為看到自己是看見上帝的同在.
因為看到火車火嘛.
她就責備年輕人.
就是看到表面困難.
反而以莉莎她是選擇和年輕人去同行.

$^{361}$她由左邊走到僕人的右邊.
她換位思考.
以年輕人的視角去看這個世界.
剛好剛剛就在這裡揭了程序表.
就看到上行者準備.
又要準備行山出隊.
好像去荔枝莊.
是嗎?.
我沒有看到那個地方.
總之我看到上行者剛剛準備行山出隊.
假設有一天.
我們少年部的團友.
就在說我們要和上行者一起行山.
都很難行的.
我知道他們的山都很挑戰性.
然後他們上行者就帶著我們一班少年部去行山.
期間如果有一個青年人跌進山洞裡.
發現自己受了傷.
然後這個年輕人就從山洞的底部.
向外面呼求救命.
有沒有人救我啊.
這裡很黑啊.
我的腳也動不了.
受了傷.
我被困住了.
我真的很害怕.
我受不了了.
快點找人來救我吧.
救命啊.
這個時候就像圖片一樣.
站在外面的山洞.
有一個弟兄姊妹就這樣回應.
她就說.
喂 年輕人.
有沒有事啊.
咦.
我聽到你回應我.
你有聲音的.
那就是沒什麼大礙了.
你說很黑.

$^{401}$你的腳也很痛.
你沒有即時的生命危險.
行了.
等一下.
我正在來.
我正在叫人.
正在召車來.
正在召直升機來.
你不用怕.
黑而已.
不會吃掉你的.
不用擔心.
我們正在找人.
很快就來救你了.
先撐著.
這個是其中一個人的回應.
另外一個人回應.
有點不同的.
他就選擇自己跳進了山洞裡面.
然後就跟受傷的年輕人說.
咦.
原來這個山洞.
真的很黑.
真的很擠迫.
真的很侷促.
跌下來那一刻.
你一定覺得很突然.
你一定很害怕.
你的腳是不是受了傷.
是不是動不了.
嗯.
是啊.
你也很擔心腳的傷勢.
不用擔心.
沒事的.
不要擔心將來腳會怎樣.
你問我為什麼要跳進你的山洞裡面.
因為.
我聽到你說受不了.
我聽到你說害怕.

$^{441}$我想陪著你.
我想陪著你一起等待救援人員.
一起去面對.
那你就不會那麼孤單.
你就不會那麼害怕.
兩個人總比一個人好.
弟兄姊妹.
兩個回應.
兩個回應其實都是跟年輕人說.
啊.
你在山洞裡面.
你不用害怕.
但是你自己聽到之後.
如果你是山洞裡面的年輕人.
你又會想聽到哪一個回應呢.
的確其實兩邊都好像跟以莉莎一樣說.
不要懼怕.
不過兩者的分別是.
後者在跟年輕人說不用害怕之後.
那位同行者願意首先放下身段.
逆地而處.
他願意陪伴年輕人一起走.
這份陪伴.
不只是坐在旁邊的那種陪伴.
而是願意站在年輕人角度裡面去換位思考.
因為用同一個角度去感受對方的感受的時候.
跟行者受傷.
就會覺得.
就算受傷他也有人跟自己連結.
所以我們經常說兩代人溝通.
可能很難.
都不太知道大家在想什麼.
我自己覺得那個關係建立.
的重要性或者關鍵就是陪伴.
陪伴一下年輕人.
以同理心去理解年輕人.
就好像以莉莎.
她知道的她明知道有火車火馬圍繞著自己.
但她不用擔心.
她也會知道.

$^{481}$以莉莎.
但她就沒有擔心.
甚至她可以換位思考地說出年輕人的恐懼.
不要懼怕這個說話是有很大安慰的力量的.
不過這份安慰的力量.
不是因為只是一個命令.
而是因為以莉莎她用陪伴.
用同理.
跟僕人一起去面對那份困境.
不要小看這些好像不是實在的力量.
這些陪伴這些同理心其實.
真的可以幫人去連結一起.
可以連結人與人之間的信任.
連結人與人之間的關係.
當有關係這個基礎的時候.
我們就可以引領年輕人從看見表面.
去看見上帝.
我們又再提一次.
十六字經文.
以莉莎說完不要懼怕之後.
她下一句就說.
與我們同在的比與他們同在的更多.
這句說話其實以莉莎是刻意使用了當下的處境.
去告訴門徒.
上帝在這個危難當中.
是和我們一起同在的.
上帝現在是在保護著我們.
我也會想.
其實為什麼以莉莎不用比較長的篇幅.
去解釋一下同在是什麼呢.
如果回到五經的話.
其實有很多經文都講解了什麼叫同在.
又或者我們真的可以去看多一點的時候.
其實她也可以說一下.
上帝與以莉亞同在.
然後又與她這個徒弟以莉莎同在.
她可以有很多經驗分享.
又或者以莉莎也可以訓斥她的門徒.
你明明看了我怎麼救奶萬的.
你明明看了我怎麼令書店兒子復活的.

$^{521}$你跟了我這麼久.
為什麼你還要這麼小信.
為什麼你在這些困境裡仍然要驚惶失措.
不過以莉莎應該不是一個很長氣的人.
我是很長氣的.
剛才我在裡面9點30分跟她說了十分鐘才出來.
以莉莎就不是很長氣的.
她也沒有責備.
她就跟門徒說了一句話.
她就說.
與我們同在的比與他們同在的更多.
這一句很簡單.
不過正正是因為這一句.
令到年輕人和自己和上帝的關係有一個連結.
他看到原來自己和上帝的關係是有份的.
所以其中一條或者第一條.
去引導年輕人看見的橋樑.
就是要描述上帝在年輕人生命上的工作.
令到年輕人可以和上帝有那份的連結.
描述上帝的工作這件事聽起來好像很大.
好像一個不容易處理的事.
可能甚至我們由創世記一直看到的啟示錄.
其實都不夠去描述上帝的屬性.
或者去描述上帝的大能.
如果我們真的單憑聖經的教導.
就可以令到年輕人和上帝連結的話.
我都想到一個辦法應該很簡單的.
我決定了下一季的主日學開始.
我們不要一季只看七至八章羅馬書.
或者《使徒行傳》看太少了.
我們應該要每一季教一本書.
之前有試過的.
每一季教一本《但以理書》.
教一篇詩篇.
一百五十篇.
真的很棒.
如果只要我計算過.
每一季去上教一本書的話.
讀完五點五年.
就可以看完整本聖經.

$^{561}$基本上和一個中學的課程是差不多的.
所以我都可以想到.
弄一個聖經的DSE給大家享受一下.
是不是真的單憑主日學.
單憑聖經的教導.
就可以幫助年輕人.
從熟世的眼光扭轉到熟齡的眼光呢.
當然按照聖經的教導.
一定是最穩陣的.
我知道的.
也不可以離開聖經.
不過如果只是硬邦邦.
將一些熟齡或者信仰的知識.
倒模一樣.
一式一樣.
用一些熟齡的金句.
金句很重要.
不過有時候有些熟齡的用語.
去教導和建立年輕人的信仰生命的話.
其實年輕人的信仰.
只會停留在這裡.
而不能進入他的心裡.
以五色珠為例.
如果我想和一個年輕人去傳福音.
如果我只是跟他說.
耶穌是創造了世界的.
然後你就有罪.
然後耶穌就為你死而復活.
你就成為一個新造的人.
這樣你就可以上天堂了.
我和一個年輕人這樣說的時候.
其實他就會.
哦.
好啊.
就是會有這樣的反應.
他其實看不到福音和自己的關係.
我們經常說福音是一件又真又活的東西.
福音又真又活.
那個活.
其實也有活潑的活在裡面.

$^{601}$所以如果我們想青少年和上帝連結的時候.
我覺得我們可以.
為青少年度身訂造他們的福音.
更準確來說.
這個度身訂造的福音.
我們當然不是說.
我們的福音要扭轉來說.
但是我覺得我們在敏銳上帝.
在年輕人身上工作的時候.
我們在抓緊他們面對的困難.
面對的處境.
利用一些教導的時刻.
英文就是teachable moment.
利用這些教導的時刻.
生活化地將信仰回應年輕人的困境.
以下這個例子.
我就有問過當事人可不可以說.
不過你們不會知道當事人是誰.
因為我不會告訴大家.
這張照片是什麼時候拍的呢.
就是去年冬天的時候.
我和幾位中學的導師.
帶著十幾個團友.
一起去了大美督附近BBQ.
如果大家有去過大美督BBQ的話.
其實那裡也比較多爛地.
很多爛地停車場.
有幾個小朋友.
也不是小朋友了.
他們都是中學生了.
他們就和我說.
Carmen我們燒完東西吃吧.
我想去後面踢球.
如果大家有去過四樓.
都知道我們的小朋友很喜歡踢球.
不知為什麼一把鏡子.
不斷地踢球也不爛.
真的很厲害.
繼續說下去.
結果我就說.

$^{641}$你們平平安安地去吧.
他們就平安地去.
愉快地玩.
驚慌地回來.
什麼事呢.
就是他們.
有一個其中一個人.
大力地一腳射.
射中別人的車.
就刮花了別人的車頭.
所以當時.
他們就像伊利沙的僕人一樣.
跑回來求救.
就說.
我在啊Carmen.
我現在怎麼辦.
當中你一言我一語.
有一個小朋友就說.
我們走吧.
我們馬上走人.
有一個人就說.
不關我的事.
他射的.
而且我不是射到刮花的位置.
他冤枉我.
我都沒有射到那個位置.
不是我刮花的.
最後導師繼續出手.
當然是跟別人賠錢.
難道真的跑路嗎.
很奇怪.
被人知道我射中.
真的搞不定.
但是.
問題就由導師解決了.
不過不代表這班年青人.
因為他射完這一球之後.
他的生命就會有改變.
所以我說要把握這個教育時刻.
我就看到這個教育時刻.

$^{681}$然後就抓住他們幾個.
就跟他們做解說.
做Debriefing.
Debriefing裡面.
其實我不是想要責備他們.
因為責備他們都沒有用.
都賠了錢.
錢都是要賠的.
又不是叫他們千萬不要再射球了.
不然遲早射爆四樓的鏡頭更危險.
其實我是要透過這個Debriefing的時間.
我是想去聽一下他們.
當時發生意外之後的心情.
然後聽完之後.
就去教導他們.
有錯就不是逃跑的.
有錯是要站在那裡承認的.
一起玩.
就要一起有份認錯.
一起要共同承擔這個懲罰.
在跟他們說的時間.
其實我都沒有引用很多聖經的經文.
我都沒有引經據典去說.
究竟彼得是怎麼犯錯的.
耶穌是怎麼寬恕他們的.
我沒有長篇大論的.
不過都說了很長的.
不過我就學校了以利沙.
其實當時就是抓住了.
年輕人面對的困境.
再用我們基督教的一些價值觀.
去建立年輕人的生命.
在他們的生命裡面去肯定.
其實上帝就是那位大領者.
上帝就是那位雕琢他們.
塑造他們生命成長的那位神.
我們這裡在座都有很多屬靈的前輩.
我想大部分都是我的屬靈前輩.
有時候我們其實都可以想一下.
我們真的有沒有去聆聽年輕人.

$^{721}$他們的一些困難.
他們的困難其實都很多.
學業 工作 原生家庭帶給你的.
不要浪費他們和我們分享的時間.
其實他們願意和我們分享.
是對我們一個很信任的態度.
裡面是包含了關係.
包含了他們對你 認定你.
是一個可以信靠 可以幫助他們的人.
所以在聽到這些困難的時候.
我們嘗試可以捉緊這些教育的時刻.
然後為年輕人去度身訂做他們的福音.
就好像我和他們做Debriefing.
其實我都沒有說耶穌拯救了他們.
因為我們知道賠錢不是耶穌賠的.
我只是告訴他們.
其實上帝不會想你無端端逃避你的責任.
我們知道上帝就是他們生命裡面的主.
剛才提到中學生踢球刮花了車的事件.
大家聽完之後都覺得很好笑.
就和他們做Debriefing.
Debriefing完之後我和他們做了些什麼呢.
和以莉莎一樣 或者和大家很多時候做的一樣.
都是和他們祈禱.
雖然我們比年輕人有更多的信仰經驗.
我們在信仰上是他們的前輩.
或者可以帶領他們說幾句話.
不過其實我們都只是一個人.
我們都是有人的限制.
一個人不可以只是用嘴頓或者用言語.
去令到另一個人有一些屬靈的經驗.
或者去經歷到屬靈的上帝.
一個人要經歷神其實很多時候是個人的經歷.
綜合了他個人的認知 個人的經驗 甚至是個人的詮釋.
其他前輩就算看到和解釋說.
上帝有很多的作為和解釋很多次的話.
都不代表他可以明白.
最多只是輔助.
正如大家不知道喜不喜歡看飲食節目.
看我的身形都知道我喜歡吃東西.

$^{761}$有時候我都會很喜歡在YouTube.
或者陪父母看飲食節目.
隔著螢幕看到賣相很吸引.
說得很便宜 現在又特價.
然後好像很好吃.
甚至主持人都會繪形繪聲說食物的味道 口感.
但其實我們看完不代表吃了.
這個節目最多就是吸引我去這間餐廳.
有些資訊告訴我這間餐廳的地址在哪裡.
唯有我親身去到那間餐廳的時候.
我聞到食物的氣味.
我真的品嘗到食物的味道的時候.
它才是一個完整的飲食體驗.
同樣地 今天的年輕人要看見上帝的時候.
都不是靠我們嘴頓.
都是要靠我們在禱告裡面.
為他屬靈 為他的生命去禱告.
這份禱告就是成為年輕人看見上帝的橋樑.
《下法爾利沙》的禱文說.
求耶和華開者少年人的眼目.
使他能看見.
為年輕人的生命懇切去祈禱.
祈求上帝去突破年輕人.
只能看到表面的這個屬靈盲點.
讓上帝親自介入或張開少年人的眼目.
使年輕人看見.
這份看見.
我們希望他們可以看到.
上帝一直以火車火馬圍繞著他們.
當年輕人看到這份屬靈的看見的時候.
他就可以定精於上帝那裡.
與上帝連結.
確立他們在上帝裡面兒女的身份.
希望他們可以一生跟隨上帝.
最後在講道的時候.
我就想跟大家分享一個年輕人的故事.
這個年輕人其實我小時候就認識他.
看著他在教會從小到大成長.
是一個信義代.
我看著這個年輕人.

$^{801}$從小到大都覺得他的生命.
他的信仰階段.
其實都頗混混惡惡的.
一直都只是活在別人的期望裡面.
這個混淆的狀況一直持續到大學.
不單止是毫無進步.
甚至是退步.
因為這個年輕人最後.
我看到他連大學都不可以畢業.
當時我看到這個年輕人.
就和以尼沙的僕人一樣.
看到環境的表面.
他就在跟我說.
這個人連大學學位都沒有.
那時候是沒有畢業證書.
到底怎樣可以在香港這個社會立足呢.
結果這個年輕人始終.
都是要硬著頭皮出來投身社會.
不過投身社會之後.
雖然他在信仰上混淆惡惡.
但他又有繼續回教會.
在教會裡面他開始去接觸一些牧者.
一些生命導師.
很記得這個年輕人.
繼續和我分享的時候.
他講過.
他說其實這些牧者.
我聽過他們講很多出色的講道.
這些生命導師每個星期小組的時候.
都會和他查經.
甚至很努力去預備經文.
不過再精彩的講道.
再好的查經預備.
都不及在我迷失的時候.
我只要打個電話.
我只要發個訊息.
我只要敲下牧者的門.
這些牧者和導師.
就會騰出空間開放家裡.
陪伴他.

$^{841}$聆聽他.
一邊聽他信.
一邊拆解這些問題.
再和他說一下上帝.
如何在困境中引領著掌權著.
不過其實這個年輕人.
每一次聽完之後.
他說最感動的地方.
是每一次那些人和他一起祈禱.
一起祈禱的時候.
他就會感動落淚.
這些牧者和導師.
在他祈禱承托著的時候.
就祈求上帝開他屬靈的黑眼.
讓他看見.
上帝是人命生命的主.
不是牧者不是導師.
就是上帝他自己本身.
一眨眼這個年輕人不吾倒業.
已經十年了.
有誰會想過.
曾經絕學的年輕人.
竟然最後會重返校園.
選擇再次讀書.
甚至修讀神學.
今天站在講壇上和大家分享上帝的道.
弟兄姊妹.
你的生命會不會都有屬於你的以利亞.
當日我們如何從我們的以利亞去被愛.
去被聆聽.
去被牧養.
去被祈禱的時候.
今天我們成為了以利沙.
我們如何可以指引我們寶貴的下一代.
同樣以關係作為漢業的基礎.
以牧養以禱告成為漢業的橋樑.
令到年輕人可以親眼看見上帝.
我們願意的時候一起有一個低頭期的時間.
上帝多謝你.
多謝你讓我們的生命當中.

$^{881}$有屬於我們的以利亞.
都有屬於我們的下一代.
多謝你愛年輕人.
即使他們只是有屬世的漢奸.
但你從來沒有放棄他們.
但願上帝你率領我們旺朕的年輕人工作.
使用我們每一個都以關係.
牧養.
禱告.
成為青少年人生命的信仰橋樑.
願上帝你親自打開青少年的眼目.
興起屬你青少年人.
作你的精兵.
以上禱告乃奉我救主耶穌基督的聖名自祈求.
阿們.
我們繼續回到家裡.
\newpage



\section{使徒行傳 7:51-60}
\label{sec:K_X10pNASkM}
\textbf{2024年5月5日崇拜直播｜陳明泉牧師｜十個深化關係的禱告-學禱告學做人｜使徒行傳7\_51-60}
\newline
\newline
連結: \href{https://youtube.com/watch?v=K-X10pNASkM}{\texttt{https://youtube.com/watch?v=K-X10pNASkM}} ~~~~ 語音日期: 2024-05-07
\newline
\newline
\hyperref[sec:IOUGpe3_PlM]{\small{< < < PREV SERMON < < <}}
~
\hyperref[sec:index_chronic]{\small{[返順時目]}}
~
\hyperref[sec:CNrc6_M4HZY]{\small{> > > NEXT SERMON > > >}}
\newline
\newline
使徒行傳 7:51-60
\newline
\begin{longtable}{cl}
\hline
\hline
章節 & 經文 (和合本修訂版)\\
\hline
7:51 & \begin{tabularx}{0.7\textwidth}{X} 「你們這硬著頸項，心與耳未受割禮的人哪，時常抗拒聖靈！你們的祖宗怎樣，你們也怎樣。 \end{tabularx} \\ \\ \relax
7:52 & \begin{tabularx}{0.7\textwidth}{X} 先知中有哪一個不是受你們祖宗的迫害呢？他們把預先宣告那義者要來的人殺了。如今你們成了那義者的出賣者和兇手了。 \end{tabularx} \\ \\ \relax
7:53 & \begin{tabularx}{0.7\textwidth}{X} 你們領受了天使所傳佈的律法，竟不遵守。」 \end{tabularx} \\ \\ \relax
7:54 & \begin{tabularx}{0.7\textwidth}{X} 眾人聽見這些話，心中極其惱怒，向司提反咬牙切齒。 \end{tabularx} \\ \\ \relax
7:55 & \begin{tabularx}{0.7\textwidth}{X} 但司提反滿有聖靈，定睛望天，看見神的榮耀，又看見耶穌站在神的右邊， \end{tabularx} \\ \\ \relax
7:56 & \begin{tabularx}{0.7\textwidth}{X} 就說：「我看見天開了，人子站在神的右邊。」 \end{tabularx} \\ \\ \relax
7:57 & \begin{tabularx}{0.7\textwidth}{X} 眾人大聲喊叫，摀著耳朵，齊心衝向他， \end{tabularx} \\ \\ \relax
7:58 & \begin{tabularx}{0.7\textwidth}{X} 把他推到城外，用石頭打他。作見證的人把他們的衣裳放在一個名叫掃羅的青年腳前。 \end{tabularx} \\ \\ \relax
7:59 & \begin{tabularx}{0.7\textwidth}{X} 他們正用石頭打司提反的時候，他呼求說：「主耶穌啊，求你接納我的靈魂！」 \end{tabularx} \\ \\ \relax
7:60 & \begin{tabularx}{0.7\textwidth}{X} 然後他跪下來，大聲喊著：「主啊，不要將這罪歸於他們！」說了這話，就長眠了。 \end{tabularx} \\ \\
[1ex]
\hline
\hline
\end{longtable}
$^{1}$《史獨行傳》第七章五十一至六十節.
主持人:你們這硬著頸項.
心如已未受割禮的人.
尚是抗拒聖靈.
你們的祖宗怎樣.
你們也怎樣.
哪一個先知不是你們祖宗迫不呢?.
他們也把預先傳說.
那二者要來的人殺了.
如今你們也把那二者賣了,殺了.
你們受了天使所傳的律法.
竟不遵守.
眾人聽見這話.
就極其牢牢.
向史提反拗牙切齒.
但史提反被聖靈充滿.
定睜望天,看見神的榮耀.
又看見耶穌站在神的右邊.
就說:我看見天開了.
人子站在神的右邊.
眾人大聲喊叫.
鼓著耳朵.
齊心湧上前去.
把他推到城外.
用石頭打他.
作見證的人.
把衣裳放在一個少年人.
名叫素羅的腳前.
他們正用石頭打的時候.
史提反呼籲主說.
求主耶穌接收我的靈魂.
又跪下大聲喊著說.
主啊,不要將這罪歸於他們.
說了這話就遂了.
素羅也喜悅,他被害.
丙姐妹,平安.
今日我們進入十個心法關係禱告的最後一講.
這一講,如果大家記得.
我們從舊約到新約.
今日是一個比較轟烈的場面.

$^{41}$去講禱告.
當然,今日所讀的經文.
處境上是一個比較轟烈的場景.
但我們應用今天的內容.
不一定要轟轟烈烈.
在生活日常裡.
要將聖經的教導應用在生活當中.
今日我們所講的內容.
我特別集中講關於聖靈充滿的主題.
如果我們講禱告的時候.
我們只是用一個平時.
我們怎樣著重我們的句語.
或者我們只是講一些生活性的東西.
其實我們就錯失了信仰裡面.
一個很寶貴的關於靈性的操練.
我們怎樣在靈裡面禱告.
又或者作為一個基督徒.
怎樣去實踐一個被聖靈充滿的生命.
今日我們以斯提反作為一個重要的例子.
事實上在《士堂行傳》裡面.
講聖靈充滿.
基本上講最多次是在斯提反身上.
在這麼轟烈的場景裡面.
其實在斯提反的生命當中.
不是純粹一個很轟烈的.
在前半部分我們稍後也會講論的時候.
我們也會涉及到的.
其實他的生命很多地方都值得我們學習.
他的禱告或者聖靈充滿為他帶來.
怎樣去面對一個很艱難的日子裡面.
我相信都成為我們今日很重要的養份.
在我們今日講這段經文之前.
我有一個這樣的例子.
這個是我信主二十多年前.
我大學的時候.
信主沒多久.
同學帶我去一些叫做特會.
不知道大家有沒有去過.
那些聖靈特會.
我們有些同學在大學裡面.

$^{81}$因為信主一起帶我去參與這些聚會.
這些聚會和我們一般崇拜很不同.
當時那些情景我還歷歷在目.
我記得那些同學帶我去.
說這個特會一定要去.
很棒的.
我說那些什麼不棒的.
只是聚會而已.
跟著去吧.
去到那時候那些詩歌很好聽.
整個鋪排都很流暢.
去到中間有一個環節.
我就開始有點不自然.
為什麼呢.
因為開始那段時間.
在台上領師和帶領的領子妹.
現在大家就可以在靈性裡面任意祈禱.
大家就在那裡祈禱.
接著就聽到很多.
以前我沒聽過的一些方言祈禱.
他們說用天使的言語.
我不懂的.
在那裡祈禱.
接著帶我去的同學.
都在那裡這樣.
我就拍一拍他.
你在做什麼.
你在祈禱什麼.
不要吵我.
我在那裡用方言祈禱.
我已經覺得很奇怪.
很神秘.
整件事跟我們平時參與教會的聚會很不同.
接著再去多一會.
音樂再澎湃一些.
接著去到最澎湃的音樂.
忽然間就靜了.
整個場面就沒有音樂.
只聽到人的祈禱聲.
我就看到我旁邊另外一個姐妹.

$^{121}$開始在那裡祈禱的時候.
她開始有點搖搖晃晃.
接著我問她.
她又笑又哭.
我不知道她在做什麼.
我看到有些人躺在地上.
有些人在跳舞.
有些人在嘔吐.
我幾天前.
腸胃炎.
看到嘔吐都害怕.
在那裡嘔吐.
發生什麼事.
什麼地方來的.
弄到這樣.
我再拍我旁邊那個帶我來的.
特別主張我去的姐妹.
你做什麼.
她說你不要搞我.
我的性命在充滿.
我跟她說.
我信主.
那時候說的是兩三個星期.
什麼.
為什麼信耶穌搞成這樣.
性靈充滿是什麼.
為什麼那個人又哭又笑.
是性靈充滿.
在那個場景裡面的人聚會.
變成一個這樣的聚會.
這個問題一直帶了我幾年.
到我神學的時候.
開始看書.
開始慢慢才有了解.
不過在這裡裡面.
可能很多人.
我們會覺得性靈充滿.
是一種這麼神秘.
或者這麼怪異的表現.
但是在聖經裡面.

$^{161}$原來從來都沒有這樣的情景.
是描述性靈充滿.
反而我們所讀的.
這段《仕途行傳》說的性靈充滿.
反而說到一個被性靈充滿的人.
可能是逆流而上的人.
希望今天這段經文幫到我們.
我們同心祈禱.
天父我們求您在我們當中.
教曉我們祈禱.
在我們生活裡面.
我們知道我們不是純粹單憑用理智來認識您.
我們在心靈誠實.
在我們的靈裡面.
我們需要與您交往.
但是我們怎樣去做.
才是合乎您的心意.
合乎聖經的教導呢.
求您就在我們當中裡面.
藉著您的話語啟示我們.
讓我們每一個心靈.
被您的話語激勵.
又讓我們能夠安著正確的認知.
以致我們不會走歪.
我們所走的道路.
我們恭敬將來今天我們所領受的經文.
和我們對您的認識交託給您.
求您賜福在我們當中.
幫助宣講的講得清楚.
幫助聆聽的每位電影節目聽得明白.
我們衷心的禱告.
求您垂聽.
獻上禱告奉基督耶穌得勝的聖名.
阿們.
我們一起看.
在我們講解今天的內容之前.
有少少幫大家補一補放.
我相信有些電影姐妹知道.
不過有些電影姐妹可能要重溫一下.
就是關於聖靈充滿和聖靈的洗.

$^{201}$或者聖靈的浸.
這個重要的意思.
其實每個基督徒信主是什麼原因呢.
有些人以為是我們做理性的選擇.
以為我們人有能力.
按著我們自己的理性.
好像超級市場買食物一樣.
喜歡選A.
我選對了選擇.
看中選擇月刊.
給我們做了明智的選擇.
我們很多時候以為是高舉理性的選擇.
不過整本聖經裡面啟示.
人能夠跟隨上帝.
就不是因為人有什麼力量.
人能夠選擇相信上帝.
其實一切都是聖靈的工作.
在聖經裡面說的是聖靈的浸.
聖靈的浸就是幫助我們.
重新活現一個新的生命.
按著聖靈的帶領.
讓我們能夠知道耶穌基督是我們生命裡唯一的救主.
能夠選擇上帝.
按上帝的心意來過我們的生活.
我們的決志.
我們的感動.
這些一切不是來自人的智慧或理性的結果.
而是來自聖靈的工作.
聖靈的浸其實就是一個這樣的心靈狀態.
其實如果在座裡面你已經決志相信基督耶穌.
你已經每一個已經經歷聖靈的浸了.
聖靈的浸不等於你做一個實體的水泥.
如果你做這個水泥之前.
你沒有聖靈的浸.
你沒有經歷重生.
你經歷明白信仰.
你根本做不到那個水泥.
那個水泥只是將你裡面所信呈現出來.
這個是聖靈的浸.
信主就是聖靈的浸.

$^{241}$在《士徒行傳》裡面因為第一代信徒.
聖靈的浸和聖靈充滿好像分開了.
好像一個叫做斷絕了的.
或者是discontinuity.
就是說不是延續性的.
所以它就分開來講.
但是去到我們今天.
我們其實所看的就是.
上帝已經將他的啟示.
上帝的工作在聖經裡面呈現了.
所以我們信主.
接著由聖靈的浸得著這個救恩.
我們已經進入了那個.
就是聖靈內住在我們的內心.
被聖靈洗乾淨那種的心靈狀態.
接著之後.
人信了主之後.
不等於你的信仰的終結.
那個只是你的起步點.
是不是.
就是你信了主之後.
你不會自動等上天堂的.
你的生命在聖經裡面講.
你要不斷的來到成長.
那個就叫做成聖.
而聖靈充滿就是一個很重要的.
一個.
你說的一個階段.
或者一個追求的目標.
幫助你的生命.
逐漸逐漸的.
一路變得成熟.
越來越像上帝.
像基督耶穌.
聖靈充滿在你的心靈當中.
指教你.
指引你.
所以在《以弗所書》裡面.
他將聖靈充滿.
某程度上成為了.

$^{281}$好像一個命令那樣.
來作為一個勸勉.
在《以弗所書》五章十七至十八節.
不要作糊塗人.
要明白主的旨意如何.
不要醉酒.
酒能使人放蕩.
乃要被聖靈充滿.
《以弗所書》的作者保羅在說什麼呢.
他說你的內心.
你要給什麼.
你的裡面.
要給什麼來充滿你呢.
不是給那些令人放蕩.
亂了性的酒來充滿你的祥福.
而是被聖靈充滿了你裡面的內心.
這種充滿會令你怎樣呢.
令你有智慧.
令你不會在這個世代裡面.
渾渾噩噩地作一個糊塗人.
反而在這個世代裡面.
多麼多衝擊都好.
因著聖靈充滿.
你能夠做一個明智.
明白上帝旨意的一個聰明人.
這個某程度上是一個.
很重要的一個表彰.
或者一個命令.
叫基督徒這樣來實踐.
聖經羅馬書裡面八章九節再繼續說.
如果神的靈住在你們的心裡面.
你們就不屬於肉體.
乃是屬聖靈.
人若果沒有基督的靈.
就不是屬於基督的.
經文裡面其實將肉體和屬於聖靈.
就好像二分了一樣.
他在說什麼呢.
其實每一個基督徒.
當你信了主之後.

$^{321}$你要說你是屬於上帝的.
你就時刻就要靠近.
或者追求聖靈充滿.
追求聖靈充滿在你的生命.
每一部份來作你的基督徒人生.
簡單來說.
你信了主之後.
你要繼續的來成長.
怎樣去成長呢.
邀請那個住在你內心裡面的聖靈.
充滿你生命的每一部份.
你的行事為人.
你的價值觀.
你的理性.
你的感受.
你的關係.
你的待人處事.
你的盼望.
等等一切所有的範疇.
你都要容讓上帝掌管一切.
這個過程不容易的.
你不要想著一個很簡單.
自動而言.
被動式自動就會充滿.
不是這樣的.
這種聖靈充滿的生命.
是需要有很多努力.
有很多的禱告.
很多的思考.
大概你明白了.
聖靈充滿某程度.
是上帝給我們一些重要的生命.
那種操練的取向和一個目標.
你說有沒有人100\%充滿呢.
沒有人敢說這件事.
只是我們今天比昨天滿一些.
一步一腳印來認識上帝.
一個誠信的過程.
這個聖靈充滿就是一個這樣的歷程.
正如我剛才所說的引子.

$^{361}$很多人以為聖靈充滿.
追求這種生命.
就會有很多神秘的經驗.
會說謊言.
用天使言語的方式來祈禱.
當然有些弟兄姊妹是懂的.
有這樣的能力去說這些.
一般人聽不明白.
或者聽下去你會覺得.
不是正常的禱告的方法來禱告.
我不否定有這些的可能性.
但是在聖經裡面沒有說到.
聖靈充滿的人就一定會說謊言.
沒有一個這樣對等的記號.
聖靈充滿的人也不是畫了一個等號.
等於那個人好像王大仙一樣.
走到每一步就醫好那個病.
說一句就搞定了.
說一句先知的話.
所有事情就按著他口中所說的.
他的領袖是怎樣呢.
全部跟他的心意來行.
聖靈充滿從來沒有一種.
某一種的生命的神跡.
其實一定要跟著這種方式來進行.
不過聖靈充滿的人卻是說到他的品格.
他對上帝的那種信服.
這個反而是一個很重要的表徵.
簡單來說.
有些人可能他有一些能力.
不過如果他生命裡面沒有那種信服.
或者按著上帝而活出的那種見證.
那種喜樂.
大概他未必真的有一種聖靈充滿的生命.
今天我們要專注的就是.
聖經裡面說到一個合而的.
聖靈充滿的人生是怎樣.
我們用《史提反》來看這個例子.
我沒有印給大家.
我讀給大家聽一段經文.

$^{401}$《使徒函傳》第六章第三節至第五節.
那段經文是說什麼呢.
就是在說《史提反》的出現.
由第六章第一節開始說起.
那時門徒曾多有說希利利話的猶太人.
向希伯來人發怨言.
因為在天天功級的上邊忽略了他們的寡婦.
十二使徒叫眾門徒來對他們說.
我們撇下神的道去管理飯食.
原是不合宜的.
所以弟兄們.
當從你們中間選出七個有好名聲.
被聖靈充滿.
智慧充足的人.
他們就派他們管理這事.
但我們要專心以禱告傳道為事.
大家都喜悅著說.
就揀選了《史提反》.
那時大有信心.
聖靈充滿的人.
又揀選了肥利,伯羅哥羅,尼加羅,提門,巴米那.
並進猶太教的安提安人尼戈拉.
叫他們站在使徒面前.
使徒禱告了就按守在他們頭上.
神的道興旺起來.
在耶路撒冷門徒數目加增的甚多.
也有許多祭司信從了這道.
在這六章裡面.
開頭所說的背景.
就是說那些跟隨耶穌的十二門徒.
就覺得一路信主的人越來越多.
教會要管理的事務實在太多了.
他不是純粹管理聚會.
是管理生活.
因為還有很多貧苦大眾.
大家一齊來抹餅.
一齊聚會.
一齊來生活.
所有的事務實在太多了.
那十二使徒根本搞不定.

$^{441}$於是說如果我把時間都放在做這些事.
他死定了.
他不用再去傳道.
不用再去禱告.
不用再去做上帝叫他們做的使命了.
於是他就把一些管理的職事.
交付給七個被聖靈充滿有智慧的人.
今天我們叫這些弟兄姊妹做執事.
其實原文來講.
只是找七個僕人.
一個願意承擔教會管理事務的.
一個有智慧被聖靈充滿的人.
值得留意的是.
這幾個人沒有特別說到他有什麼神職能力.
代表他們只是有智慧能力.
還有他們有心智.
還有就是在屬靈上.
他們是被聖靈充滿的.
所以按照這個規範.
選了史提反是其中一個.
他特別說到史提反是一個大有信心.
和被聖靈充滿的.
雖然找他做管理的職事.
但史提反不只是做管理飯食.
他管理飯食之外.
他也依然在做門徒要做的工作.
某程度上來說.
他比那十二使徒還要辛勞.
因為他又要做管理.
又要講道.
而講的道.
他講整個《士壇傳》裡最長的那篇道.
就是由史提反所講.
那個內容最尖銳的就是由史提反所講.
當我們知道這個背景的時候.
在經文裡面再繼續的來講.
史提反因為得著滿德恩惠能力.
在第八節那裡.
在民間裡面就行了很多的大歧視和神跡.
這裡是有講到聖靈充滿.

$^{481}$令到他有這些神跡歧視.
不過當他做完這些事之後.
當時有稱為利伯地拿會堂的幾個人.
並有古里內亞歷山大.
基里加亞細雅各處會堂的幾個人.
都起來和史提反辯論.
史提反以智慧和聖靈說話.
眾人敵擋不住.
這裡是六章第八節至第十節.
剛才講到管理教會的行政事務.
本來是史提反一個重要的身份.
一個角色.
不過他不單只是一個角色.
他的地界只是框著做這些事.
他要將他自己.
或者上帝呼召他.
他要做更加多.
他在他的管理之外.
還要去為著真理來抗辯.
很多人去做一些假見證.
或者誣衊跟隨耶穌的一幫門徒.
史提反就在那裡除了做一些行神職.
歧視和管理教會之外.
他也承擔起一個抗辯的角色.
當有幾個人誣衊他的時候.
史提反用智慧和聖靈來說話.
其他誣衊他的人完全都敵擋不住.
你看到那個聖靈的能力很厲害.
上帝用一個僕人.
加他的力量.
讓他能夠有一種這樣的生命.
來抵抗.
或者為上帝的正道.
來成為一個真正的僕人.
在這裡首先我們要留意的.
剛才說到那個聖靈充滿.
是上帝給予的一種能力.
一個目標.
但同時也是每一個信徒裡.
承聖的一個很重要的經歷.

$^{521}$在史提反的身上.
不是他禱告然後聖靈充滿.
而是他一直已經被聖靈充滿.
而在他面對不同困難的時候.
他繼續來禱告.
所以他的屬靈生命.
他不是呼天搶地.
去到某些關係有些困難的時候.
變身然後聖靈來.
然後他就充滿力量.
那些是日式的電影才會這樣做.
但真正的信仰運作不是這樣.
真正的運作是一個不停追求上帝.
主宰生命的僕人.
不斷的來祈求.
面對不同的挑戰.
求上帝加他力量.
他的日子如何.
他的力量又如何.
其實這個不是純粹在史提反身上.
每一個弟兄姊妹.
其實我們都可以實踐一個這樣的生命.
有些前輩說.
聖靈充滿只是某小撮人才有.
一些是個別事件.
不過其實你看看聖經.
其實在聖經裡說到.
聖靈充滿.
或者在史提反身上.
他更加說到.
是一個禱告.
以上帝為首的一個人.
他願意被上帝用的時候.
聖靈就會加力量給他.
讓他能夠好像天使那樣的能力.
有智慧來進行一種抗辯.
就是一種生命裡有上帝充滿的狀態.
有一種追求的目標.
這個真是一個上帝大用的僕人.
我們看到這部分.

$^{561}$我們會覺得很開心.
我們會覺得教會很有未來.
一個這麼有靈氣的人.
這麼有靈力的人.
來成為教會一個忠心的僕人.
不過到了第七章.
當史提反這樣抗辯的時候.
剛才說的誣衊他的人.
他們會不會這麼輕易放過他呢.
在第六章尾到第七章開始.
其實就說到那些剛才誣衊他的人.
不會這麼輕易放過他.
於是在六章十三節那裡說.
這些人就設下假見證說.
史提反這個人不住這樣來踐踏聖所和律法.
我們曾聽見他說.
說拿撒勒耶穌要毀壞這個地方.
也要改變摩西所交給我們的規條.
在公會裡面坐著的人都定睛看他.
見史提反的面貌.
但是見到史提反的面貌.
面對著這些質詢是怎樣呢.
他的面貌就好像天使的面貌一樣.
接著到第七章.
大祭司就說了.
這些事果然有嗎.
其實將史提反逼到一個牆角.
說他們是一幫什麼人呢.
就說耶穌說要毀壞聖殿這件事.
就無限誇大.
他只是說他字面的意思.
接著就在那裡咬文嚼字.
去針對他.
針鋒相對.
也都在說毀壞他們整個摩西的傳統.
哇就是這種罪行很大.
棄師滅祖.
對得起列祖列宗嗎.
於是就在會堂裡面.
透過大祭司的質問.

$^{601}$來質問史提反.
接著在第七章第二節開始.
到第五十節.
或者由一節到五十節.
因為篇幅關係我們沒辦法去說.
被質問之下.
史提反就說了一個驚世的一個信息.
這個信息是令他朝夜殺身之禍.
這個什麼信息呢.
史提反就由亞伯拉罕說起.
由上帝選召亞伯拉罕.
由一個人變成一個民族.
由一個民族變成一個國家.
由一個民族曾經被埃及欺壓.
摩西怎樣去帶領他們.
整個歷史事跡說出來.
不過你要留意.
在整個歷史事跡裡面.
重點放在哪裡呢.
史提反不是純粹數白欖.
如數家珍將整個以色列史.
編年史說出來.
他有個重點在整個講章裡面.
整個講章的重點就是說到.
由上帝選召的.
由一個人開始的民族.
成為一個活躍的民族.
是由第一天開始.
原本上帝選召他們.
但是這班人多次又多次地.
來抗拒上帝的帶領.
這一班是活躍.
橫亙抗拒聖靈的民族.
所以去到第51至第53節.
是這段講章裡面的一個總結句.
他講完整個歷史的目的.
他要講什麼呢.
51至53節.
你們這硬著頸汗.
心與意未受割禮的人.

$^{641}$嘗試抗拒聖靈.
你們的祖宗是怎樣.
你們都是這樣.
哪一個才知不是你們祖宗逼迫的.
他們也把預先傳說的乃而者.
要來的人殺了.
如今你們又把乃而者賣了殺了.
你們受了天使所傳的律法.
竟不遵守.
剛才說那些人在指控.
說耶穌要毀壞聖殿.
說這班使徒.
他們說要改了摩西的律法.
但是斯提反他講的這篇講論.
不單止是要回應.
更加要講到這班不理情由.
不理上帝帶領的這班活逆的百姓.
最大的活逆是他們.
而不是這班使徒.
他說的內容是很尖銳的.
甚至他在數歷史.
一邊數他都可以很輕手.
隱隱晦晦暗示都可以.
不過他不是這樣的.
一個聖靈充滿的斯提反很有勇氣.
你舊時的祖宗是怎樣.
你們現在都是這樣.
沒有改變過.
這是一個很辣的句語.
還要講到那些先知.
他們講的很守過的傳統.
那些先知.
那些舊時上帝所使用的天使.
或者上帝曉喻的那些僕人的說話.
他們覺得很遵守.
但全部都是給你.
舊時的人怎樣害死.
現在你們都是這樣將那些義者一一殺死.
當然他講到耶穌都是由這班活逆的百姓來殺.
現在甚至對於基督徒.

$^{681}$都是帶一種這樣的態度.
天使的律法所傳的.
他們到今天都不遵守.
這一句的說話.
這一種的態度.
我們不會看.
我們只是覺得他很有勇氣.
我們很多時候有些弟兄姊妹說.
有勇氣有吉時.
不過其實在聖經裡面講.
斯提反為何會講這番說話.
除了他個人有一種獨特.
那種夠魄夠勇氣之外.
其實是聖靈充滿.
一個被聖靈充滿的人.
是願意講真心話的.
是會講真誠的真相出來.
一個被聖靈充滿的人.
是敢於為上帝作見證.
是敢於面對罪惡.
呈現這個世界的荒謬.
一個敢於被聖靈充滿的人.
是願意用自己的生命來見證上帝.
在這裡裡面.
我們今天都要問一問自己.
各位弟兄姊妹.
如果我們今天要追求聖靈充滿.
或者我們要和上帝的關係緊密.
我們都要問自己.
你願不願意成為一個真誠的人.
願不願意講真心話.
不過今天我很坦白說.
今天我們為了維繫某些關係.
我們有時很怕和別人.
我們怕衝突.
我們不太喜歡講真誠的話.
有時會講真誠的話.
不過用什麼方式講呢.
用指桑罵槐.
用那些中國人有時講話不太好的方法.

$^{721}$那些就叫做.
講反話.
將一句正面的話反轉來說.
想著那個人會意會.
又或者講的東西就含沙射影.
就是那種方式.
表面上他好像在講真話一樣.
但實際上他沒有勇氣將他心底最真誠的話講出來.
如果我們經常給一些積雜在心裡.
我們經常講來講去.
我們都講不到真正的內容.
我們只是講了一大堆.
令到人憤怒.
但講不到內容的那些說話.
不過如果真正被聖靈充滿.
我們要有一份勇氣.
又或者我們要求上帝給我們有聖靈充滿的力量.
來將心底真誠的話表達出來.
當然有時講了一些說話.
如果講真相可能會帶來一些傷害.
不過你不講真相.
你不要想著沒有傷害.
只不過那個傷害是延後了.
或者是其他人來承受.
如果你真的愛那個人.
你真的想那個人好.
你需要講真誠的話.
當然你不是要攻擊他的方式來講.
你是為他好的角度來講.
帶著愛心來講誠實的話.
唯有這樣你才可以成為一個頂天立地.
成為一個可以由內而外都是一個真誠的人.
不過當士提反講完這番說話的時候.
那些人有什麼反應呢.
在經文裡面繼續的53節講完這番說話之後.
接著54節就講到那些人的反應.
眾人聽見這話就極其牢牢.
向士提反咬牙切齒.
但士提反被聖靈充滿.
再一次講聖靈充滿.

$^{761}$定睛望天看見神的榮耀.
又看見耶穌站在神的右邊就說.
我看見天開了.
人子站在神的右邊.
眾人大聲喊叫.
烏著耳朵齊心湧上前去.
把他推到城外用石頭打他.
作見證的人把衣裳放在一個少年人名叫素羅的腳前.
他們正用石頭打的時候.
士提反呼籲主說.
求主耶穌接收我的靈魂.
有跪下大聲喊著說.
主啊不要將這罪歸於他們.
說了這話就碎了.
素羅也喜悅他被害.
最後這個段落.
你會覺得和剛才的講論有很大的落差.
同樣都是被聖靈充滿的士提反.
他怎樣祈禱呢.
那班人對他咬牙切齒.
充滿了攻擊性.
他帶著真誠帶著上帝給他的勇氣.
帶著聖靈的能力.
講論了一個很尖銳不討好的訊息.
但是當他面對這些群眾將要將他殺害.
他將要快要死的那一刻.
他被聖靈充滿的時候.
他舉目望天.
他見到上帝的榮耀.
見到基督耶穌站在他的右邊.
接著他見到耶穌在神的右邊.
見到異象.
但是另一個鏡頭就拍著那些人咬牙切齒又要打他.
但是在那一刻你會滲透出一種.
士提反的心裡面是寧靜的.
在這樣的狀態之下.
他心裡面只有兩個願望.
求上帝接收他的靈魂.
這句話很熟面口.
其實在路加福音23章46節.

$^{801}$耶穌在十字架上.
同樣說了近似的說話.
父啊,我將我的靈魂交在你的手裡.
說了者說,氣就斷了.
耶穌在十字架上受苦的那句話.
十加七言,其中的一句話.
成為了士提反.
他學禱告,在聖靈充滿裡面所說的那一句話.
他將自己的生命交付給他見到異象的.
坐在寶座上的基督耶穌.
求基督耶穌接收他的靈魂.
第二個的祈願是什麼呢?.
當他見到那些人喜悅他被受害的時候.
六十節,跪在那裡大聲喊著.
主啊,不要將這個罪歸在他們的身上.
那些人是很討人厭.
那些人經世以來都是活逆上帝.
現在還要殘害使徒們徒的一班的白人的人.
但是在這一刻.
士提反心裡沒有任何的仇恨.
聖靈充滿的人.
心裡沒有任何的怨恨,寂怨.
他只是求上帝不要將這些罪歸他們的貪玩.
這句話又很熟面口.
路加福音23章34節.
父啊赦免他們.
因為他們所作的,他們不能要.
在極之困難,在極大的痛苦當中.
心靈被聖靈充滿.
任何的仇恨不佔據他的心靈空間.
聖靈充滿的人.
生命就是這樣的一種表現.
今天我們簡單地看了.
士提反整個生命的歷程.
用了第六章第七章.
其實是一個較長的篇幅.
但是對於一個很豐富的人生.
其實是一個很短的片段.
不過我們看到一個聖靈充滿的.
一個禱告的人.

$^{841}$他的生命不是特別有很多誇張的神蹟.
歧視,又或者有很多令人羨慕的能力.
他反而是一個禱告的人.
在多少個困難當中.
祈求上帝給他有力量,有智慧.
在他生活裡面.
面對有很多不同的為神所作的任務.
聖靈加他力量.
充滿他,賜他智慧.
他能夠做得好.
聖靈充滿的人.
在他生命裡面沒有很多的界線.
或者他不做,或者什麼.
他只是為上帝緣故.
他就願意豁出去.
在聖經裡面沒有嚴格的分工.
試圖去做這些.
那些被選召的人就做這些.
就嚴格,大家不做.
不是,而是選召更加多的僕人.
忠心,被聖靈充滿的人.
願意的來做.
承擔上帝天國的召命.
讓他們的生命呈現福音.
聖靈充滿的人.
亦都是一個敢說真話的人.
即或那句說話,不討好.
但他不會用指桑罵槐.
不會用戚戚然的態度.
來面對真相.
坦誠,真摯.
但不帶仇恨,將真相呈現出來.
聖靈充滿的人.
亦都是心靈裡面.
是一個豁達的人.
生命裡面做到最終.
他會見到上帝的榮耀.
即或面對有多少的艱難.
心裡不會懷恨.
這種聖靈的生命.

$^{881}$充滿生命是怎樣來的呢.
其實不是來自甚麼神秘的禱告.
很多教會歷史學家告訴我們.
其實當代的門徒.
當代第一代的信徒.
他們的生命是怎樣操練的.
其實耶穌怎樣說.
耶穌怎樣去禱告.
他們就跟著去禱告.
在中國人.
我們畫畫就會用這個字眼.
臨悟.
即是師父怎樣畫.
起了一個稿.
徒弟就跟著這個稿去畫.
所有生命的基礎.
所有生命最實在的.
都是由臨悟開始.
信仰都是一樣.
我們今天的祈禱是怎樣.
你喜歡祈甚麼就祈甚麼.
是嗎.
教會歷史.
在史提反的生命告訴我們.
他跟耶穌的祈禱.
跟耶穌怎樣祈禱.
他就怎樣去祈禱.
不過不是學那個語言.
是學耶穌的心法.
學耶穌的做人態度.
效法基督的意思是.
耶穌怎樣做.
他就跟著怎樣做.
按著他的祈禱.
按著耶穌禱告的行事為人.
所以史提反就是按著.
學耶穌的禱告.
跟著面對迫不得已的時候.
他都是用耶穌的方式.
來去回應眼前這些困難.

$^{921}$弟兄姊妹.
今天我們很長的篇幅.
我們沒有辦法完全去涉獵.
整個史提反的講論.
和他的生命事跡.
不過我希望.
在未來的日子裡.
我們如果追求聖靈充滿.
追求禱告的生活.
我們可不可以.
去一些基本功呢.
有些東西錯過了.
不要緊.
回頭再來.
不能太遲.
任何時間都可以開始.
我對於聖經裡面所講的禱文.
不是很清楚.
就重新看過.
由零開始建造起來.
重新建造起來.
重新建造起來去穩定它.
我的心底話是.
黃浸是一個.
很有聰明智慧的群體.
不過有時候.
我們會輕看一些基本功.
我們覺得我們懂了.
但其實我們不懂的.
我們錯失了那樣東西.
所以我們的信仰走得不深.
我們去到某些位置.
我們就像瓶頸一樣卡住了.
我們去到某些位置.
我們就像瓶頸一樣卡住了.
可以怎樣做呢.
回頭來.
由基本功開始.
由零開始.
逐步逐步來.

$^{961}$你會見到你的生命.
有革命性的改變.
不是甚麼神奇的工作.
只不過是基礎打得不好.
以致我們走得不夠遠.
走得不夠力量.
由主幫助我們.
最後一講.
十個改變生命的禱告.
就由聖靈充滿的禱告開始.
就由一個基本功的禱告開始.
希望我們這個系列.
鼓勵大家繼續努力.
朝向一個誠信的生活.
成為上帝喜悅的子民.
成為一個聖靈充滿的人.
肯定充滿的人.
\newpage



\section{詩篇申命記以賽亞書 22:9-10}
\label{sec:CNrc6_M4HZY}
\textbf{2024年5月12日崇拜直播｜區祥江博士｜神像母親養育我們｜詩篇22\_9-10;以賽亞書49\_15,66\_12-13;申命記32\_11}
\newline
\newline
連結: \href{https://youtube.com/watch?v=CNrc6-M4HZY}{\texttt{https://youtube.com/watch?v=CNrc6-M4HZY}} ~~~~ 語音日期: 2024-05-13
\newline
\newline
\hyperref[sec:K_X10pNASkM]{\small{< < < PREV SERMON < < <}}
~
\hyperref[sec:index_chronic]{\small{[返順時目]}}
~
\hyperref[sec:OGOXjFxluCg]{\small{> > > NEXT SERMON > > >}}
\newline
\newline
$^{1}$今次經文是記載於詩篇.
《以賽亞書》及《申命記》.
大家可以翻開程序表內頁.
詩篇二十二篇,九至十節.
「但你是叫我出母腹的.
我在母懷裡,你就使我有倚靠的心.
我自出母胎就被交在你手裡.
從我母親生我,你就是我的神」.
《以賽亞書》四十九章,十五節.
「婦人應能忘記她吃奶的嬰孩.
不連恤她所生的兒子.
即或有忘記的,我卻不忘記你」.
《以賽亞書》六十六章,十二節下至十三.
「你們要從中享受.
你們必蒙抱在肋旁,搖擺在膝上.
母親怎樣安慰兒子.
我就照樣安慰你們.
你們也因耶路撒冷得安慰」.
最後一節.
《申命記》三十二章,十一節.
「猶如鷹攪動巢窩.
在初鷹以上兩次擅展.
接取初鷹背在兩翼之上」.
很感恩今天由中國神學研究院.
副導科教授歐長江博士.
在中間為我們講訊息.
母親節的主題.
「神讓母親養育我們」.
請歐婆.
(歐哥) 第一件事是早安.
(記者) 早安.
其實講母親節的崇拜是有挑戰的.
因為每年都有母親節.
聖經裡的哪位母親你們沒聽過呢?.
瑪利亞的母親.
摩西的母親.
薩姆爾的母親.
想了想.
不如我找一些經文.
關於母親的跟大家分享.

$^{41}$不說特別一位母親.
不過母親在我們的成長裡.
其實扮演一個很重要的角色.
但我今天的角度是.
大家有沒有想過.
我們的神.
其實也像一個母親呢?.
當然我們通常稱呼神為天父.
好像有一個父親的形象.
但其實聖經裡有不少經文提到.
我們的神像母親一樣.
所以我今天選了幾段經文.
跟大家重溫一下.
其實我們的父神.
也有一個母親的特質.
母親的特質有幾個經常人會提到.
第一是關懷.
第二是養育.
以及保護.
這三樣都是母親一個很重要的角色.
我想透過今天這段經文.
講一下母親.
這裡有很多母親.
你們做母親的過程.
可能也經歷了我剛才的四段經文.
所經歷的事情.
有沒有想過我們的天父.
我們的父神.
其實也是在屬靈的方面.
同樣是在我們身上的呢?.
我自己其實這個星期是很開心的.
因為我做第二任窮工.
5月6日.
我第二個孫女出生.
因為我和我女兒同住.
所以我親手幫忙抱她.
3.03公斤.
其實很小.
我拍了照片.
她的腳掌放在我手掌上.

$^{81}$其實也是佔這麼多.
很小的.
但是抱著一個這樣的生命.
其實是很奇妙.
很感恩.
很開心.
所以也和大家分享這個喜悅.
因為見著自己的女兒做母親.
她也有餵奶.
有泵奶.
也有一些辛苦的時候.
因為她有一個兩歲大的孫女.
我又會不會增重.
所以她這段時間也有些挑戰.
我作為爸爸.
加上我太太.
其實這段時間也對她有些支援.
我們回到這段經文.
詩篇22篇.
九到十節.
這裡說.
但你是叫我出母胎的.
我再母懷裡.
你就使我倚靠.
有倚靠的心.
我想這是我們每個人生命的開始.
是你叫我出母胎.
有些聖經學者的翻譯.
是說到神好像一個接生婆.
見證著我們的出世.
而我們在母胎裡.
你就使我們有倚靠的心.
這是一個生命很奇妙的過程.
我們叫做懷胎十月.
一個生命在母胎裡.
其實是經歷一個完全無助的狀態.
她得到的生命資源.
都是透過與母胎的臍帶結連.
支取了養分.
原來我們出世就有一種.

$^{121}$作為人.
有一種很依靠的心.
這裡奇妙的地方.
第十節.
我自出母胎就被交在你手裡.
從我母親生我.
你就是我的神.
我們的神.
無論我們未出世之前.
我們出世的那一刻.
出世之後.
我們的母親就把我們交給神.
就是說我們整個人的生命.
無論你零歲到八九十歲.
整個生命都放在我們神的手裡.
當我思想這件事的時候.
其實心裡很感恩.
我想想.
小孩出世就是一個肉身的出生.
其實我們在屬靈裡.
都有經歷一個重生.
就是說信主的經歷.
都是一個新的生命的誕生.
大家想想信主的經歷.
其實是一個很奇妙的歷程.
跟這裡的描述相似.
在我們未信主之前.
其實神已經鋪排了很多人生經歷.
幫我們去到一個位置.
我們對神覺得不可以靠自己了.
要全然靠神.
所以把我們的生命交給神.
經歷了耶穌基督十字架的救贖.
有一個重生的經歷.
重生之後.
其實我們的人生怎麼走.
都一直在天父的手裡.
或者在天父好像母親一樣的手裡成長.
所以詩篇二十二篇九到十節.
可以說是一個肉身的媽媽的經歷.

$^{161}$但也是一個屬靈的經歷.
我們來看第二段經文.
《以賽亞書》四十九章十五節.
這裡說.
「婦人焉能忘記她吃奶的嬰孩.
不連恤她所生的兒子」.
即或有忘記的.
我卻不忘記你.
現在年輕的媽媽.
問問你們有沒有母乳餵哺.
那個感覺是怎樣.
現在母乳餵哺都是混合的.
餵人奶,奶粉.
有些也會這樣.
不知道大家有沒有想過.
其實我們母乳餵哺.
其中一個好處就是.
人奶有很多抗氧物質.
有豐富的營養.
對嬰孩的成長是很好的.
這是生理上的好處.
但其實更重要的好處.
是媽媽和嬰孩之間的情感連繫.
一個連繫.
是吧?.
我想.
這件事我會和太太比較.
「你做媽媽,我做爸爸」.
「為何父親節這麼容易找位置」.
「好像我的子女有些大細超」.
「好像多疼媽媽多於疼爸爸」.
有些不太憤氣.
「我做爸爸做得很好」.
但來到這些位置就無話可說.
生產的過程.
和母乳餵哺.
都是我們男人做不到的.
是吧?.
這兩件事.
就是媽媽和嬰孩.

$^{201}$最重要的情感連繫的位置.
無法代替.
我想一個媽媽.
如果母乳餵哺過自己的子女.
你和她有種很難.
因為有種身體的接觸.
或者她絕抱的時候.
其實你的身體會被觸碰.
有一種觸碰.
有個連繫在這裡.
所以這件事是相當寶貴的.
我都見證著.
現在我女兒都很崇尚這件事.
所以兩人都沒有乳餵.
見到她很辛苦.
又要泵奶.
但其實是很值得.
這個歷程對嬰孩的成長很好.
如果我們的神.
好像一個母親.
我拿了一個比喻.
比喻是.
其實我們吃什麼大呢?.
基督徒.
其實是聖經.
神的話.
我不知道大家有沒有回想.
回想就是我們剛剛信主.
或者在教會初信成長的時候.
其實喝慕靈奶的經歷是很深刻的.
有沒有覺得讀聖經很寶貴.
得到很多啟示.
很多神的話幫到我們成長.
初信的階段.
就等於媽媽在餵哺嬰孩.
很純粹在神的話裡.
支取到屬靈的餵養.
有沒有想過.
其實我們的神.
就像母親一樣.

$^{241}$透過她自己很寶貴的話.
育養了我們一段屬靈的生命.
所以今天一方面想.
歌頌母親的同時.
我們也想想我們的神.
其實祂是怎樣.
在我們屬靈生命的歷程裡面.
跟我們有很多相遇.
很多encounter(碰撞).
很多nurture(培育).
即是培育我們生命的地方.
是很值得我們感恩.
第三段經文.
又有少少不同.
我們一起來看.
第十二節.
後半節.
「你們必蒙抱在肋旁.
搖籠在室上」.
這是一幅甚麼圖畫.
甚麼圖畫?.
即不是嬰兒.
可能一歲兩歲.
最好玩的時候.
我們的嬰兒最好玩.
就是這些小孩子.
還肯讓你抱.
你可以抱著他.
或者在你的膝蓋.
彈彈.
這裡其實是說.
神給了我們甚麼呢?.
這個母親和嬰兒.
玩這些東西.
抱他.
或者在室上搖籠.
其實是說.
跟嬰兒有一種遊戲.
玩.
有一種.

$^{281}$跟他很輕鬆的接觸.
其實都是一個小孩子成長的.
要有的歷程.
然後十三節就說.
「母親怎樣安慰兒子?.
我就照樣安慰你們」.
其實母親一直以來.
都是我們一個很重要的.
傾訴對象.
是嗎?.
有沒有誰是你有心事.
有事不開心.
是找爸爸傾訴?.
還是找媽媽傾訴多些?.
其實媽媽是很擅長安慰的.
特別是肥肥的媽媽.
抱著你.
擁抱著你.
拍你的背.
你是很舒服.
很受到安慰的.
其實在我們的靈命成長裡.
我們有沒有經歷這些事?.
成長的時候很多東西玩.
我們信主之後.
可能學會事奉.
跟教會群體一起成長.
去營地.
或者人生裡面遇到的考試.
戀愛失敗.
成功.
或者出來工作.
其實母親都一直看著我們.
然後當我們遇到困難的時候.
她總是願意跟我們聊天.
安慰我們.
我回想起幾十年前.
我已經登陸了.
我年輕的時候.
都試過拍拖失戀.

$^{321}$其實媽媽都會問我.
關心我.
安慰我.
我想我們的信仰歷程裡.
我們給神安慰的時候.
其實真的很多.
好像母親.
母親怎樣安慰兒子.
我就照樣安慰你.
神自己說.
我就好像母親那樣去安慰你.
多寶貴.
很快的.
今天講的可能會短一點.
你們可以早點去喝茶.
新命記.
三十二章.
我們一起看看.
這裡形容一個.
如果前一點.
第十節的尾就這樣說.
或者我讀完第十節給大家聽.
這裡說:日和華遇見他在曠野.
荒涼野獸嘯叫之地.
就環繞他.
看顧他.
保護他.
如同保護眼中的同人.
然後就讀這段.
又如鷹攪動茶窩.
在初鷹耳上.
兩次善展.
接取初鷹.
背在兩翼之上.
其實這裡是說到.
神那種保護.
剛才說母親其中一個很重要的角色.
或者特質.
都是保護.
就像母雞保護自己的小雞.

$^{361}$不要被鷹襲擊一樣.
這裡的描述是挺特別的.
如果我要你們想想.
是一個人的什麼人生階段.
是像這裡的描述.
說鷹攪動茶窩.
本來小鳥在茶窩裡.
他就攪動.
其實是趕他出茶窩.
明白我的意思嗎?.
就像你躲在茶窩裡.
就太過安書了.
要離開你的安書區.
要學飛.
所以就把你撥出去.
接著呢.
初鷹就真的試試.
但可能起初都不太行.
母鷹就用兩翼護航.
然後就把他升上去.
其實他快要掉下來.
接一接他.
然後就把他推上去.
其實就是讓他去學飛.
是一幅很美麗的圖畫.
就像我們教小朋友踩兩輪單車.
扶著但都要放手.
如果跌倒也能抓住.
其實我們的父神.
在我們的成長裡.
其實就在做一樣.
幫我們去離開一些安書區.
當我們要去學飛的時候.
其實他是看著我們.
護航著我們.
如果快要掉下來的時候.
就打開雙翼.
然後就下來.
撥你上去.
讓你飛得更高更遠.

$^{401}$這個講一下現代的母親.
其實現在的母親.
我們都有些說法.
叫做直升機母親.
或者父母或者父媽.
什麼是直升機母親呢?.
就是經常盤旋在子女上.
他一有問題就衝下去救他.
而這種母親的說法.
其實就不太健康.
就是說你沒有讓他自己學.
去解決問題.
有事就投訴家長.
要怎樣.
其實對小孩的成長.
未必是一件好事.
這裡是相反的.
不是衝下去救他.
而是在巢裡撥他出去.
其實對於母親.
我們說人要離開父母.
在一個青少年成長的歷程裡.
我們的子女需要離開父母.
特別是離開母親.
因為父親一直都沒有怎樣理會.
所以不需要離開那隻小孩.
當然有些父親會這樣舉例.
母親肯放手.
因為媽媽很多時候很擔心.
很多憂慮.
很忌妒.
很多事情都為他安排好了.
但其實是需要讓他跌跌碰碰.
在座的媽媽.
神都這樣做.
神都好像母嬰一樣.
特意撥他出去.
又扶著他.
適當時候又幫他起飛.
但不是每次都去拯救他.

$^{441}$這是我們做母親需要學習的.
所以很快地和大家分享了這四段經文.
其實是講一個母親的歷程.
由你生孩子到你餵奶.
到她在兒童期你和她玩.
她不開心你安慰她.
到她青少年期.
其實你需要學她放手.
讓她自己成長.
我們見證每個母親都在做這件事.
是人生的歷程.
寶貴的地方是原來我們的父神.
我們的神都在經歷.
好像母親一樣.
和我們屬靈生命相處.
都是一個這樣的歷程.
所以有時候我們會覺得.
神為何不理會我們.
可能他已經覺得我們到了青少年期.
應該可以扔我們出去.
跌跌跌起飛.
要學習甚麼.
不得不說我們除了反思.
母親的角色之外.
我們懷著一個感恩的心.
如果大家記得.
有一本很出名的書是路運神父寫的.
路運神父寫這本書的時候.
他引用了一個荷蘭的名畫家的名作.
《狼子比喻》那幅畫.
畫家叫Rambrandt.
路運神父特別看到.
他看到爸爸擁抱著兒子.
其實他有兩雙不同的手.
一隻手粗獷一點.
一隻手溫柔一點.
其實他在說.
狼子比喻裡面的父親就等於父親.
其實他有威嚴.
有比人生的標準和規範的一面.

$^{481}$好像很強壯的那隻手.
他同樣有一隻很溫柔.
懂得呵護和愛我們的一面.
其實我們的父親.
又像爸爸又像媽媽.
不過我們慣常叫他做天父.
他也有作為一個媽媽的特性.
栽培著我們的成長.
神是按照自己的形象.
做男做女.
其實在我們的神裡面.
他有男有女.
有母性有父性.
所以我們基督徒的母親.
應該很感恩.
因為我們學做母親這件事情上.
原來我們的神.
其實也提供了很多經歷給我們.
去經歷一份母愛是一回事.
希望我們作為母親.
我們先給神的愛.
擁抱融化.
我們被愛之後.
我們就有更加大的心靈空間.
去愛我們的子女.
我們一起低頭祈禱.
因為你在聖經裡面.
留下了很多的說話痕跡.
讓我們看到.
你不單是那位公義超賢的神.
你也是那位臨近我們.
安慰我們,養育我們那位.
當我們思考母親.
作為母親的生命歷程的時候.
我們也很感恩.
你不單像我們在地上有母親.
而你自己也像母親一樣.
生我們,養育我們.
陪伴我們,安慰我們.
有時候你也會放手.

$^{521}$讓我們去經歷,讓我們去學習.
以致我們能夠成長.
在這個世界上,作為母親.
不論如何,我們也經歷到.
神給予我們自己的愛.
以致祂懂得怎樣愛我們的兒女.
在我們人生裡.
如果我們母親還在生.
我們把握機會去感恩,去淘寶.
所做的每一件事.
這個世界是你創造的.
人倫是你創造的.
在這些關係裡.
我們也經歷到你愛的奇妙偉大.
我們這樣禱告,奉賴你名及.
阿們.
\newpage



\section{創世記 28:10-22}
\label{sec:OGOXjFxluCg}
\textbf{2022年9月4日崇拜直播｜陳明泉牧師｜以石為記:伯特利的枕首石｜創世記28\_10-22}
\newline
\newline
連結: \href{https://youtube.com/watch?v=OGOXjFxluCg}{\texttt{https://youtube.com/watch?v=OGOXjFxluCg}} ~~~~ 語音日期: 2024-05-26
\newline
\newline
\hyperref[sec:CNrc6_M4HZY]{\small{< < < PREV SERMON < < <}}
~
\hyperref[sec:index_chronic]{\small{[返順時目]}}
~
\hyperref[sec:1HAPF8xhKJg]{\small{> > > NEXT SERMON > > >}}
\newline
\newline
創世記 28:10-22
\newline
\begin{longtable}{cl}
\hline
\hline
章節 & 經文 (和合本修訂版)\\
\hline
28:10 & \begin{tabularx}{0.7\textwidth}{X} 雅各離開別是巴，往哈蘭去。 \end{tabularx} \\ \\ \relax
28:11 & \begin{tabularx}{0.7\textwidth}{X} 到了一個地方，因為已經日落，就在那裡過夜。他拾起那地方的一塊石頭枕在頭下，就躺在那地方。 \end{tabularx} \\ \\ \relax
28:12 & \begin{tabularx}{0.7\textwidth}{X} 他做夢，看哪，一個梯子立在地上，梯子的頂端直伸到天；看哪，神的使者在梯子上，上去下來。 \end{tabularx} \\ \\ \relax
28:13 & \begin{tabularx}{0.7\textwidth}{X} 看哪，耶和華站在梯子上面，說：「我是耶和華—你祖父亞伯拉罕的神，以撒的神。你現在躺臥之地，我要將它賜給你和你的後裔。 \end{tabularx} \\ \\ \relax
28:14 & \begin{tabularx}{0.7\textwidth}{X} 你的後裔必像地上的塵沙，必向東西南北開展；地上萬族必因你和你的後裔得福。 \end{tabularx} \\ \\ \relax
28:15 & \begin{tabularx}{0.7\textwidth}{X} 看哪，我必與你同在，無論你往哪裡去，我必保佑你，領你歸回這地。我總不離棄你，直到我實現了對你所說的話。」 \end{tabularx} \\ \\ \relax
28:16 & \begin{tabularx}{0.7\textwidth}{X} 雅各睡醒了，說：「耶和華真的在這裡，我竟不知道！」 \end{tabularx} \\ \\ \relax
28:17 & \begin{tabularx}{0.7\textwidth}{X} 他就懼怕，說：「這地方何等可畏！這不是別的，是神的殿，是天的門。」 \end{tabularx} \\ \\ \relax
28:18 & \begin{tabularx}{0.7\textwidth}{X} 雅各清早起來，拿起枕在頭下的石頭，立作柱子，澆油在上面。 \end{tabularx} \\ \\ \relax
28:19 & \begin{tabularx}{0.7\textwidth}{X} 他給那地方起名叫伯特利；那地方原先名叫路斯。 \end{tabularx} \\ \\ \relax
28:20 & \begin{tabularx}{0.7\textwidth}{X} 雅各許願說：「神若與我同在，在我所行的路上保佑我，給我食物吃，衣服穿， \end{tabularx} \\ \\ \relax
28:21 & \begin{tabularx}{0.7\textwidth}{X} 使我平平安安回到我父親的家，我就必以耶和華為我的神。 \end{tabularx} \\ \\ \relax
28:22 & \begin{tabularx}{0.7\textwidth}{X} 我所立為柱子的這塊石頭必作神的殿；凡你所賜給我的，我必將十分之一獻給你。」 \end{tabularx} \\ \\
[1ex]
\hline
\hline
\end{longtable}
$^{1}$《創世記》第28章十至二十二節.
陳明荃牧師為我們宣講信息.
《以石為記》.
《伯特利的錦首石》.
足夠法定人數.
請丙姐妹打開今日的經文.
記載在《創世記》.
在我們唐用聖經第34頁.
聖經的第一卷書《創世記》.
請聽我讀出.
第十節.
(主持人讀書).
經文讀筆.
(主持人讀書).
弟兄姊妹平安.
我們由今個月開始.
連續十次.
我們會開展一個新的講道系列.
《以石為記》.
這個主題是講什麼呢?.
我們想做一些新意思.
過往我們去講道.
都會就著一個神學主題.
或者有一個方向去組織.
或者有一個書卷去講道.
今次這個系列我們有少許不同.
我們這個系列.
我們會用聖經裡面的一種象徵物.
今次是用石頭來貫穿.
由舊約到新約.
你又會說牧師石頭都好講.
到處都是.
不過在聖經裡面.
當我們認真查考的時候.
到處都是的石頭.
原來是一塊小小的石頭.
或者是一塊大的石頭.
原來背負著聖經裡面很多的人物.
有很多人對神的關係.
又或者藉著一塊石頭.

$^{41}$來講道人對上帝.
或者人與人之間有很多恩怨情仇.
藉著石頭來講上帝的故事.
所以我們接下來這個系列.
與其說用石頭.
其實更加講到上帝與人之間的關係.
也藉著一塊平平無奇的石頭.
讓我們重新體會.
上帝由舊約到新約.
甚至乎到現在.
原來上帝都在我們生命裡面.
一草一蜜來對我們講說話.
今日我們由《創世記》開始看.
在舊約裡面.
以色列人的先祖.
雅各又名以色列.
看上帝在這個人身上.
由一個落魄的走投無路的人.
到一路一路發展成為以色列國.
一段的經歷.
剛才大家聽聖經的時候都看到石頭的字.
不斷重複.
希望我們看這一段經文.
我們除了領會一些道基督教訓之外.
更加重要的是體會上帝在我們生命裡面.
對於我們每一個平平無奇的人.
上帝有什麼工作.
我們一起祈禱.
求上帝幫助我們.
讓我們同心領受祂的話語.
阿爸天父願你在我們當中帶領我們.
幫助眾弟兄姊妹.
讓我們能夠在聆聽你話語的時候.
你親自對我們講說話.
我們恭敬將今天我們所讀的經文.
一段關於雅各的生命.
你如何領受你的恩典.
直著石頭展現的這段經文.
我們祈求仰望你.
願你同樣對我們每一個.

$^{81}$你的訊息帶到我們心靈當中.
改變我們對你很多.
有時候我們不是對你很理解的地方.
幫助我們.
讓我們更加認清你是我們生命的主.
求你在我們當中帶領我們.
幫助每位聆聽的弟兄姊妹聽得明白.
宣講的講得清楚.
獻上禱告奉靠基督耶穌得勝的聖名.
Amen.
我們一起看看.
在《創世記》裡面.
其實講到.
在第十二章開始講到.
其實我們經常講以色列人.
以色列人這個民族是怎樣來的呢?.
他們就一個一個那些先祖.
那些民族的領袖.
一個又一個地去描述出來.
由第十二章《創世記》已經開始這樣講.
先是講亞伯拉罕.
然後就講以色列.
其實講以色列的篇幅不是很長的.
講完以色列之後.
就講雅各了.
雅各其實佔在《創世記》裡面.
都有相當的篇幅.
雅各這個人.
我們一般想.
一個領袖.
或者一個某程度是.
開機的一個人.
一定是一些很偉大.
很叻.
甚至乎.
道德.
很偉大的一班人.
我們很多時候都會這樣想.
不過在聖經裡面就確實呈現了.
雅各作為整個以色列民族裡面.

$^{121}$真是開枝散葉的那個.
不是一個這樣的人.
在經文裡面就馬上來記載.
在第28章裡面.
就先記載著.
剛才我們所讀的第十節裡面.
就講到.
雅各出了別時巴.
向哈蘭走去.
無端端講雅各這個人要離開別時巴.
要去哈蘭.
去到那個地方.
上文下理原來就交代了.
剛才講過.
雅各不是一些很偉大的人.
他有一個孿生哥哥叫做爾素.
兩兄弟.
在舊時的當代社會裡面.
他的傳統就是說.
大兒子.
長子.
就會承受雙倍的家產.
承受雙倍的家業.
小兒子當然少一點.
除了承受雙倍家產之外.
更加是成為家裡面的當家.
一家之主.
你不要忘記.
雅各和爾素是孿生的.
他們只是輸一隻腳位.
在這裡活出大半生的人生裡面.
那個爾素是甚麼呢.
喜歡玩的.
喜歡去打獵.
喜歡到處去.
你猜對於頭家是不是很有魄力去管理呢.
在聖經裡面說不是那些人.
相反來說.
雅各比較沉靜安穩.
和母親之間的關係比較好.

$^{161}$因為經常躲在家裡面.
但是對於家裡面的那種管理.
似乎心是在家裡面多一點.
但是去到某一刻.
他父親爾薩就開始患病.
年老患病.
於是雅各竟然做了一件事.
他就要扮他哥哥.
整個身上都是毛.
他就扮他哥哥.
整個身上都是毛.
他父親就病得瘟瘟瘟瘟.
摸摸他.
咦 這麼多毛.
就為他祝福.
就藉著這個過程.
就騙了那個長子的名份.
其實很簡單.
這樣祈禱就對了.
將家產過了.
今天我們又說要寫平安紙.
當代沒有的.
祈禱.
祝福了.
那個過程就是了.
雅各就是帶著一個欺騙.
就拿了這個長子的祈禱.
拿了他父親爾薩的祝福.
當他的媽媽生了哥哥.
二叔知道這件事怎麼辦呢.
生氣得飛起來.
於是就要殺他.
所以要追殺他的時候.
雅各為了逃避哥哥的追殺.
所以就由自己的別士巴.
就投靠他的舅父哈蘭.
在哈蘭這個地方裡面.
就是希望再遠一點.
其實不是巴蘭.
是哈蘭的再遠一點.

$^{201}$不過去到哈蘭這個地方.
由別士巴到哈蘭整個的路程.
大概有多遠呢.
我拉過一些尺.
做過一些研究.
是885公里.
885公里.
885公里是一個什麼的概念呢.
我拉過一條直線距離.
由香港去到越南的河內.
步行.
可想而知.
雅各為了要逃避.
他媽媽生了哥二叔的追殺.
和要投靠他的舅父.
他是步行從香港走到越南的河內.
走了相當長的日子.
你會想像得到.
孤軍作戰.
在一個這樣的狀態之下.
他就帶著一份落魄.
雖然可以拿了長子的祝福.
但是他享受不到長子的那個名分.
帶來的好處.
他要離開自己的父家.
要去投靠他的舅父.
在經文裡面.
就是一路一路的來到去這樣做.
其實在《創世記》27章41至44節.
那裡都有一點記載著那個心.
就是說因為他很怕瓜哥會殺他.
所以他就逃到哈蘭.
投靠李百嘉.
他媽媽的哥哥拉班那裡.
和他住一些日子.
希望他哥哥消了氣之後.
他回去就拿回長子的名分.
那麼一路一路走的時候.
去到一個地方.
這個地方後來都寫了叫做路思.

$^{241}$路思是一個不知名的城市.
一個不知名的.
你當是礦業也好.
普通的地方也好.
路思的意思是甚麼呢.
行人.
即是眾行人的地方.
那麼你大概可以想像.
雅各一路逃離.
一路步行.
走了都有相當的日子.
接著就落魄地去到一個.
名不經傳的一個小城市.
或者一個小地方.
因為在經文裡面講到.
這個人就很累.
雅各累到一個地步.
太陽一落.
他就要在那裡住宿.
睡的時候.
你會睡甚麼呢.
你不要想著真的有客棧.
有地方給他staycation(宿舍).
沒有staycation可言的.
即是在經文裡面講到.
他是投靠一個地方.
植地而睡.
接著就拾起那地方一塊的石頭.
接著就再拿來躺臥睡了.
這個是第十一節.
可想而知其實一個逃避住.
親生哥哥來追殺的一個人.
他是帶著一種恐懼.
帶著一種心裡面一種孤單.
來鄉別井.
更加面對著路上有很多的驚惶.
用落魄來形容他.
大概都不會講錯的了.
在一個步行.
長足這麼長途的路程裡面.

$^{281}$身體疲倦.
心靈困乏.
接著就植地而睡.
拿著一塊石頭.
就隨意地用一個地方.
靠著繼續的來趕路.
這塊石頭代表著的就是.
同樣好像表達到.
雅閣的那個心態.
為了甚麼呢.
人生為了甚麼呢.
想著拿到長子的名分.
想著得到心裡面很想得到的東西.
但卻是帶來的就是擔驚受怕.
帶來的就是家裡面眾叛親離.
甚至乎自己親生骨肉的哥哥.
都覺得要追殺他.
心裡面帶著一份很沉重的落魄.
和一種疲乏.
這塊石頭.
亦都是今次裡面所講的.
以石為記的.
這塊枕頭的石.
成為一個這樣的象徵.
當他一直睡著的時候.
在經文裡面第十二節到十七節.
就在雅閣的夢境裡面.
見到一個很奇異的現象.
他見到甚麼呢.
他在第十二節裡面講.
夢見有一條梯子.
立在地上.
梯子的頭頂頂著天.
有神的使者在梯子上上去下來.
在經文裡面.
其實都是帶著當代的世界觀.
他們當代其實有很多平衡的.
不同的那些民族裡面.
我們叫的那些叫做民間故事.
個個都是這樣講的.

$^{321}$天和人間.
是有一條很長的樓梯來貫穿著的.
其實中國人都是這樣講的.
那些天梯.
其實我們舊時那些人.
我們的文化.
很多的文化都是這樣.
貫穿人與天是有一條路的.
不過因為它是高的.
所以要上樓梯去做.
在雅閣裡面發夢都是見到一條這樣的樓梯.
不過這條樓梯裡面見到.
有神的使者.
那些天使.
或者上帝差遣的僕役.
在樓梯上就很忙碌地上上落落.
有一種上去下來的一種這樣的現象.
而上帝.
耶和華上帝就在梯子以上.
當然有些講法是在祂的旁邊.
不過更加多的支持就是.
上帝是在天上.
就是一個很淡定.
不是像那些使者勤勤勞勞.
不過一個很高超然的上帝.
就對著雅閣在夢裡面.
就講了一句這樣的說話.
我是耶和華李祖亞伯拉罕的神.
也是以實的神.
我要將你現在所妥協的地方.
賜給你和你的後裔.
你的後裔必像地上的塵沙那樣多.
必向東西南北擴展.
地上萬族必因你和你的後裔得福.
我也願你同在.
你無論往那裡去.
我必保佑你.
令你歸回這地.
總不離棄你.
直到我成全了向你所應許的.

$^{361}$在那個夢裡面.
剛才講到.
上帝答應了.
或者上帝將一個承諾.
向雅閣表達出來.
這個承諾是甚麼呢.
有幾樣東西.
第一.
告訴他.
他的後裔如塵沙那樣多.
第二樣東西.
他所訓的這個地方.
將來就是成為他家業之地.
第三.
他要向東西南北擴展.
不是訓那個位置.
那個位置給他為地.
比香港更慘.
他要告訴他.
他由這個地方.
中心點那樣.
不斷地來擴展.
成為一個有疆土的一個國.
一個民族.
地上的萬族.
都要因雅閣和他的後裔得福.
講到不單止是他自己.
而是藉著他這個民族.
上帝賜福給其他外邦的人.
接著.
上帝再繼續講.
你要回來這裡.
你要歸回這地.
上帝更加應許.
總不離棄他.
你會在這個夢境.
或者在這裡聽的時候.
你會想像得到.
很宏大.
很厲害.

$^{401}$講的那些東西.
不過你對比雅閣當下的處境.
他只不過是一個孤身一人.
現在要逃難.
睡覺的地方都沒有.
找一塊石頭.
來當作枕頭來睡覺.
在什麼都沒有.
這個地方都不是他的.
然後說這個地方賜給他為業.
多荒謬.
是不是真的可以的.
這樣都可以.
再加上.
他面對著的就是.
很怕他的親生哥哥.
對他的追殺.
隨時性命不保.
面對著這麼多的困難.
上帝這個應許是不是真的.
好像和他現下當下的處境.
成為一個很強烈的對比.
經文都是他跟.
接著再講.
這個夢.
上帝就是這樣來到.
去曉如之後.
十六節.
雅閣睡醒了.
就說.
耶和華真的在這裡.
我竟然不知道.
然後就懼怕說.
這地方何等好為.
這不是別的.
乃是神的殿.
也是天的門.
在這裡.
當雅閣睡醒.
他驚醒.

$^{441}$或者那一刻.
他醒過來的時候.
這個什麼地方.
這個不是杏仁.
還杏仁的地方.
這個地方是連接著.
人間與天上的一道門.
一個連斬首.
即是一個落魄的人.
來到這個地方.
都不知道是巧合還是什麼.
但對他來說.
這個地方絕對不是一個簡單.
普普通通的.
一個可以休息.
或者讓他歇下來的一個地方.
而是一個能夠貫穿人與神.
是一個上帝顯現的殿.
是天的門.
在經文裡面.
對於一個落魄者.
對於一個欺騙長子名分的雅閣.
竟然上帝在這裡.
來呈現給他知道.
上帝與他同在.
上帝會作出一些的應許.
值得注意的是.
在這段經文上帝所曉喻裡面.
他所說的內容.
其實在聖經裡面.
之前記載過的.
在《創世記》裡面.
是記載過的.
記載過什麼呢.
原來上帝說的同樣的內容.
在《創世記》第十三章.
十四至十七節.
其實這個是說.
上帝應許亞伯拉罕.
同樣的約.

$^{481}$那個應許.
現在在雅閣的夢境裡面.
重新說一次.
所以為什麼雅閣會.
發到這個夢呢.
在《創世記》的記載裡面.
這個不是偶然.
不是日有所思夜有所夢的.
一個情景.
這個是說.
上帝答應了亞伯拉罕.
答應了他的祖先.
雅閣瓜爺亞伯拉罕.
答應了他的事情.
去到如今在雅閣這一代.
來實踐出來.
昔日上帝藉著亞伯拉罕.
藉著他成為萬國的祝福.
如今他就要將這個應許.
兌現在這個.
生命有些不濟.
鴉鴉片片的.
這個小一點的小媽.
但卻成為長子的.
這個雅閣的身上.
在這裡我們再繼續下去看.
我們就有些地方.
有些再思想的空間.
在第十八節我們再看.
經文裡面說到.
當雅閣知道這個地方.
這麼可畏.
這麼神聖.
於是十八節就說了.
他立刻起了身.
清早起來.
剛才說的睡覺的枕頭石.
他就拔了出來.
然後就將他.
將所枕的石頭.

$^{521}$十八節.
立作為柱子.
你以為是一塊磚頭.
一個睡覺都不會沾不掉的石頭.
硬穿自己的後腦.
平平的.
一塊石.
然後他就將那塊.
原本很代表著落魄孤單的這塊石.
然後他就找多幾塊石.
一起來砌砌.
砌了一條柱出來.
有一個很不起眼的.
不過聖經是用這個字的.
立作柱子.
那個立作那個字.
那個立不是隨便用的.
那個立那個用字.
其實是用在宗教場所.
在聖殿才用這個字.
我們有時候說起興建.
興建有上下一代的.
大興土木那個字.
我們不會說興建一間木屋.
不會這樣的.
我們不會興建一塊Lego.
不會興建一塊積木.
但是在這裡.
雅各好像用一些睡覺的石.
好像砌積木一樣砌了一條柱出來.
但是用的字眼就是興建.
雅各在這一堆的石柱裡面.
撩上一些油.
然後在這堆石上.
他就說這個地方叫做伯特利.
是上帝的殿.
一個落魄的人.
連枕首的地方都沒有.
但是因為上帝的顯現.
他的心態改換了.

$^{561}$他發現這裡.
是上帝與人相遇的地方.
他身邊僅餘的只有一塊枕首的石.
拿幾塊石一起堆砌成為一條柱.
原本這個地方叫杏仁路斯.
現在改為伯特利.
這個名字是雅各自己改的.
其他人不是這樣叫的.
不過是雅各改這個地方是神的殿.
然後他就藉著這條柱.
藉著腰柔的.
很簡單的很原始的祭祀儀式.
他就向上帝許願.
神如果與我同在.
在我所行的路上保佑我.
又給我食物吃.
又給我衣服穿.
使我平平安安回到我父親的家.
我就必以耶和華為我的神.
我所立的柱子也必作神的殿.
凡你所賜給我的.
我必將十分之一獻給你.
這裡看到.
剛才說的雀草.
但這裡卻是說到.
他的心態改變了.
他祈願上帝與他同在.
他祈願上帝保佑他.
他求的東西和剛才耶和華上帝給他的那樣東西.
完全是兩個距離.
雅各在這裡求什麼呢.
有食物.
有衣服穿.
回到家裡.
我們大部分人都滿足到.
很基礎的需要.
很基本的.
他能夠有這些基礎的需要滿足到.
然後他就說.
我會以這裡成為神的殿.

$^{601}$還要將十分之一.
很尊貴的.
奉獻給上帝使用.
在這裡.
剛才我們也對比過.
他所求的東西和上帝應許的東西.
其實是一段距離.
但是聖經裡面不是否定了他眼前這些需要.
不過在一個落魄的人心中.
你叫他講一些很宏大.
藉著一個人.
來致福萬國.
又或者如陳沙那麼多後裔.
那些東西很離身的.
很不落地的.
眼前他所求的.
只不過是眼前的這些需要.
用這些石頭.
來作出一個許願.
在這裡我們稍微有一些思想的空間.
我們今天很多頂展會.
我們看聖經也好.
或者我們去教導.
或者去和別人分享聖經.
我們很多時候.
很多時候最後的結束.
或者結論.
我們是教別人做一個道德的人.
這個不是錯的.
我們不會反道德的.
在聖經裡面也有講很多道德倫理.
教你怎樣來作一個頂天立地的君子.
宗教也導人向善.
基督教不會導人向惡.
一定是導人向善的.
不過在聖經裡面.
和云云的宗教.
那個層次有些不同.
除了我們教別人導人向善之外.
聖經更加講到.

$^{641}$其實我們生命的問題.
不是不懂得什麼叫善.
而是在我們掙扎的人生裡面.
我們不知道怎樣去選擇自己的路.
在聖經裡面就引導著.
每一個尋求上帝的人.
可能我們自己都帶著自己.
一些不太理想的過去.
又或者帶著一些.
帶回教會或者認識信仰.
可能心都不是很純正.
我們不是很學問式地追求信仰.
大多數都是什麼呢?.
眼前有些需要.
求上帝滿足我.
以致我一直認識這個信仰.
講回我自己.
我最初那時候信主.
除了覺得扮東西.
我云云宗教看了這麼多單張.
儒釋道和基督教來做一些比較.
到最後我最靈智就選了基督教.
以前我是大學的時候.
扮學者是真的這樣想的.
是這樣講的,不是這樣想的.
是這樣講的.
但心底裡是什麼呢?.
那時候我認識了一個基督徒女朋友.
自己未信主.
我妹妹跟我說.
哥,你認識了一個基督徒女朋友.
你死了,她上天堂.
你就下地獄.
你沒有辦法在一起.
那一刻她說一說,我真的害怕.
人生第一次真的認真拍拖.
心裡這樣都可以?.
天地如此分隔.
直教生死相許.
一直講,心裡很混亂.

$^{681}$帶著的是什麼呢?.
帶著的就是不想跟當時基督徒女朋友分開.
一個不是很純正信耶穌的心.
但就要證明自己是很理性的.
所以就看了不同的儒釋度.
做一些比較.
我知道也有一些弟兄.
我不怕揭穿.
有些弟兄回來教會,為了什麼?.
有些笑,暗笑.
認識女孩子.
有些真的是這樣的.
以前我們見到很多.
當然我們有教會.
回來很久的.
我們當然有左右護法.
護著那些想回來認識女孩子的人士.
保護著我們的姐妹.
也好好招呼那位純粹認識女孩子的朋友.
但我也認識很多.
最初都是帶著這個心態.
後來成為一個很忠心的門徒.
成為一個很好的基督徒.
每個人都是認識上帝的途徑.
有些人的起始點.
未必是我們今天心目中這麼純正.
或者帶著很多自己心裡的罪性.
帶著自己很多的議題.
來到教會群體.
來到在上帝面前.
但是上帝不是一個論色的上帝.
在聖經裡面更加呈現一個圖畫.
就是雅各一個.
我們讀上去,怎麼搞的?.
這樣也可以?.
但上帝竟然對著這個人.
在他一個很落魄的路途當中.
上帝親自的來顯現.
將一個由他亞伯拉罕以撒的應許.
竟然承諾在一個.

$^{721}$我們覺得道德不是很理想的一個人身上.
更加重要的就是在這個生命裡面.
很落魄,很困難的日子.
上帝將一個更加宏大的應許.
放在這個人身上.
在雅各的眼中裡面.
他求的其實都不是.
上帝也不是說上帝.
曉喻給他知道那個應許的藥.
他都不是求那些東西.
他只是求有得吃,有得穿.
可以回到父家.
他都是求這一樣東西.
不過這個卻是成為了雅各.
認識上帝一個的起始點.
由一塊枕頭之石.
他興起上來.
要將它成為一個石柱.
來獻祭找回生命裡面.
可以扶持的一個上帝.
在經文裡面第28章就說到雅各.
就帶著一份這樣的心.
繼續的走到他舅父拉班那裡.
來經歷他人生未來的旅程.
當他投靠了他的舅父.
投靠拉班.
去到那個地方裡面.
大概他又精於計算.
其實人家也是這樣計算他.
舅父也是這樣計算他.
到最後在他人生裡面經歷了.
這段日子有20年的時間去投靠.
一走就走了20年.
在這20年裡面經歷過被舅父騙.
但是他又可以成家立室.
舅父兩個女兒就被他娶了.
不漂亮的先娶.
漂亮的遲一點.
多做幾年工.
等等一直一直去經歷的時候.

$^{761}$本來又一個孤家寡人.
後來就成為家兒子那.
一群兒女.
12個支派.
12個兒子.
就在這一段鴉鴨片片.
幾滄桑的日子.
竟然就捱出一頭家出來.
當他捱出一頭家的時候.
他想著就要回去.
重建找回他的父家.
要回去自己所長大的地方.
一來他想拿回長子的名分.
但第二件事.
他心裡面那個是他的根據地.
他要回自己的家.
經文裡面一直發展的時候.
就說到在第34章和35章.
特別是35章.
就說到雅各回到伯特利這個地方.
本身雅各為什麼會去到伯特利呢.
原因不是特意要去做什麼儀式.
或者回去尋根.
在他一頭家回去和以素他的親生大佬相遇.
重新復和.
經文裡面說到雅各很害怕.
頂硬上 禮物先走.
最後才自己去.
到最後他兩兄弟復和.
這裡有機會要說回這段經文.
不過我們今天的重點放在伯特利這個地方.
當兄弟復和解決了他二十多年的心結之後.
但是在他路上裡面.
一直回到自己根據地的地方的時候.
他雅各其中的女兒.
被外邦人事件姦污了.
於是乎見到雅各的女兒這麼漂亮就說要提親.
但是雅各其他的兒子就很生氣.
於是乎要做什麼呢.
就要設計來將事件和他身邊的人殺死.

$^{801}$於是乎就騙他們.
就說我們以色列那些承受應許的人.
個個都行國禮.
就說你娶我妹妹都可以.
全部男丁 整個民族.
個個都行國禮就差不多了.
那班人真是傻傻地.
整個群族的男士全部都行了國禮.
做完手術了 大家明白了.
做完手術之後 大家都剩下半條人命.
然後在這個時間 雅各兩個兒子.
利味和西面就殺了那裡的男丁.
當殺了男丁之後就惹下大禍.
雅各說死了.
我們怎麼及別人外邦那些人多勢眾.
所以就逃走.
當他想著很擔憂的時候.
上帝就再一次對雅各說.
三十五章第一節.
神對雅各說 起來 上伯特利去.
住在那裡 要在那裡築一座壇 給神.
就是你逃避你哥哥耳掃的時候.
向你顯現的那一位.
上帝叫雅各去伯特利.
於是第二節.
雅各就對他家人並一切與他同在的人說.
你們要除掉你們中間的外邦神.
也要自結 更換衣裳.
我們要起來 上伯特利去.
在那裡 我要築一座壇給神.
就是在我遭難的日子.
應運我的禱告.
在我行的路上保佑我的那一位.
他們就把外邦人的神像.
和他們耳朵上的環子.
那些戴耳環的環子.
就交給雅各.
雅各就藏在事件那裡的橡樹底下.
他們便起行前往.
神在那周圍成邑的人都甚驚懼.

$^{841}$就不追趕雅各和眾子了.
於是雅各和一切與他同在的人.
到了迦南地的路司就是伯特利.
他就在那裡築了一座壇.
就蓋了那地方 起名叫伊勒伯特利.
在這裡第七節 頭七節.
你就看到整件事了.
他回到伯特利幹什麼呢.
上帝叫他的.
他都是著草 帶著逃亡.
不過這次的逃亡去到伯特利.
和上一次很不同.
上一次是孤軍作戰.
孤家寡人.
現在他是帶著一群子女.
整個家族這樣回到去.
上一次在他心中裡面.
耶和華在這裡了.
不過這次他已經教定他的子女.
放下那些外邦神明.
代表著一種對神的效忠.
這次再回去 他更加知道.
那個是有了他昔日.
在二十多年前.
在伯特利經歷的路上.
保有的這個經歷.
如今他帶著那個記憶.
帶著那份經歷.
來到他帶著感恩的心.
去到這個伯特利的地方.
經文很簡單 很巧妙地說.
那些人不是追殺雅各和他整個家族嗎.
經文裡面簡單幾個字就說.
上帝令到他們那些追殺他的人很害怕.
就不再追捕他們了.
簡簡單單這樣.
不過去到雅各.
在他的生命裡面.
他是帶著一份嚴肅.
帶著一份敬畏.

$^{881}$重新去到伯特利這個地方.
在他生命裡面.
他和昔日的雅各不同.
他帶著很害怕和哥哥之間的那種追捕.
但是這次來到這個地方.
他和哥哥已經復和了.
他也不知道哪一個是耶和華上帝.
很模糊.
亞伯拉罕怎麼說 以撒怎麼說.
他就怎麼相信.
但他這次是第一身經歷.
這個真實的信仰.
他願意丟棄了那些外邦的神明.
那些偶像.
他的名字也被改變了.
本來叫雅各 叫做「維住」.
如今他的名字叫做「以色列」.
就是上帝給他的一個新的名字.
也是跟著他整個家族.
整個國家的一個名字.
他有十二個兒子.
承受這個豐盛的產業.
最後他講到他的獻祭.
也是不同的.
昔日就是隨便用一些油來燒上去.
那些燈油等等那些.
這次他可以在石柱上再敲點一些油.
加上一些酒.
比之前已經進步了.
升了級.
但這個依然是一條石柱.
是一個代表著.
不過生命裡面.
跟他最初的原點已經不同的一條石柱.
如今是帶著上帝的恩典.
走在面前的路的一條石柱.
整個今天所講的.
對照著這個伯特利的石柱的故事.
或者這一段經歷.
對我們呈現一個甚麼的訊息呢?.

$^{921}$正如我剛才所講.
每個人都有信仰的起點.
我們有些的起點.
可能在身邊的人覺得.
你有沒有搞錯?.
這樣的?.
這個起點在哪裡?.
最重要的不是我們永遠都在這個起點當中.
每個人的信仰.
可能信耶穌的心.
跟隨耶穌.
可能最初的時候心不是很順.
但我們信得過上帝.
會帶領我們一步一步.
我們的終點.
或者我們生命裡面.
是由這個起點繼續來到前行.
當我們掌握到這件事的時候.
人生我們可以輕散一些.
坦白講.
有時候我們會覺得.
那個人回來都不知道為了甚麼.
我們心裡面就招待.
我們的樣子就招待他.
心裡面就抗拒他.
是嗎?.
不過今天這段經文提醒我們.
每個人都有信仰的起點.
我們學習以上帝的心腸.
真心地來接待那些.
你大概知道他是什麼狀態的那些人.
不用這麼快來批判他.
反而是我們用真理.
用上帝的恩典.
引導他走面前人生的路.
經文裡面再繼續提醒我們知道.
上帝對人是不離不棄.
令到人生命可以改變的原因.
不是因為人有什麼道德的力量.
令到人生命裡面可以改變.

$^{961}$是因為有一個納若斯慈愛的上帝.
上帝對人昔日他怎樣答應了亞伯拉罕二撒.
他如今都兌現了這個承諾.
在雅各和在他未來的後裔當中.
今天我們人生裡面可能都經歷到.
林林總總有一些落魄困難不濟的時間.
但是我們不要忘記.
生命裡面有一個看顧我們的主.
他答應過我們.
他會對我們不離棄.
會保守我們.
我們就帶著這一份很單純的信.
在他的面前繼續尋求出路.
當我們生命裡面有這份確信的時候.
外面再艱難也好.
再覺得我們很難受的一些際遇也好.
上帝會為我們開路.
當我們回頭看的時候.
我們就見到恩典處處.
生命裡面就好像雅各所行過.
上帝為他預備恩典的神蹟一樣.
讓他經歷到上帝而來的公認和豐盛.
求主幫助我們.
讓我們作為這個系列的開始.
用一塊枕頭石成為一個敬拜上帝的祭壇的石.
用這塊石開始來更加明白上帝的心意.
讓我們能夠在我們生命跌跌撞撞的人生裡面.
經歷上帝與我們同在.
經歷上帝的施恩.
我們同心祈禱.
天父上帝願你在我們當中.
繼續對我們每一個講說話.
願你的話語在我們的心靈裡面.
繼續遙憾.
改變我們的心思意念.
讓我們所作的被你閱納.
又讓我們生命更加發現上帝的足跡.
求你繼續來施恩與頌頂展會同在.
我們多謝你.
垂聽祈禱.

$^{1001}$奉靠基督耶穌得勝的聖名.
\newpage



\section{腓立比書 1:1-11}
\label{sec:1HAPF8xhKJg}
\textbf{2024年5月26日崇拜直播｜陳冠成牧師｜靈性·關係·成長－健康的屬靈關係｜腓立比書1\_1-11}
\newline
\newline
連結: \href{https://youtube.com/watch?v=1HAPF8xhKJg}{\texttt{https://youtube.com/watch?v=1HAPF8xhKJg}} ~~~~ 語音日期: 2024-05-27
\newline
\newline
\hyperref[sec:OGOXjFxluCg]{\small{< < < PREV SERMON < < <}}
~
\hyperref[sec:index_chronic]{\small{[返順時目]}}
~
\hyperref[sec:OamF1FcoO7k]{\small{> > > NEXT SERMON > > >}}
\newline
\newline
腓立比書 1:1-11
\newline
\begin{longtable}{cl}
\hline
\hline
章節 & 經文 (和合本修訂版)\\
\hline
1:1 & \begin{tabularx}{0.7\textwidth}{X} 基督耶穌的僕人保羅和提摩太寫信給住腓立比、在基督耶穌裡的眾聖徒，以及諸位監督和執事。 \end{tabularx} \\ \\ \relax
1:2 & \begin{tabularx}{0.7\textwidth}{X} 願恩惠、平安從我們的父神和主耶穌基督歸給你們！ \end{tabularx} \\ \\ \relax
1:3 & \begin{tabularx}{0.7\textwidth}{X} 我每逢想念你們，就感謝我的神， \end{tabularx} \\ \\ \relax
1:4 & \begin{tabularx}{0.7\textwidth}{X} 每逢為你們眾人祈求的時候，總是歡歡喜喜地祈求， \end{tabularx} \\ \\ \relax
1:5 & \begin{tabularx}{0.7\textwidth}{X} 因為從第一天直到如今，你們都同心合意興旺福音。 \end{tabularx} \\ \\ \relax
1:6 & \begin{tabularx}{0.7\textwidth}{X} 我深信，那在你們心裡動了美好工作的，到了耶穌基督的日子必完成這工作。 \end{tabularx} \\ \\ \relax
1:7 & \begin{tabularx}{0.7\textwidth}{X} 我為你們眾人有這樣的想法原是應當的，因為你們常在我心裡；無論我是在捆鎖中，在辯明並證實福音的時候，你們都與我一同蒙恩。 \end{tabularx} \\ \\ \relax
1:8 & \begin{tabularx}{0.7\textwidth}{X} 我以基督耶穌的心腸切切想念你們眾人，這是神可以為我作證的。 \end{tabularx} \\ \\ \relax
1:9 & \begin{tabularx}{0.7\textwidth}{X} 我所禱告的就是：要你們的愛心，在知識和各樣見識上，不斷增長， \end{tabularx} \\ \\ \relax
1:10 & \begin{tabularx}{0.7\textwidth}{X} 使你們能分辨是非，在基督的日子作真誠無可指責的人， \end{tabularx} \\ \\ \relax
1:11 & \begin{tabularx}{0.7\textwidth}{X} 更靠著耶穌基督結滿仁義的果子，歸榮耀稱讚給神。 \end{tabularx} \\ \\
[1ex]
\hline
\hline
\end{longtable}
$^{1}$經文:《新約聖經》第273頁.
請聽我讀出.
基督耶穌的僕人保羅和提摩泰.
寫信給梵柱 肥納比.
在基督耶穌裡的眾聖徒.
和諸位監督 諸位執事.
願因為平安 從神我們的父.
並主耶穌基督歸於你們.
我每逢想念你們.
每逢為你們眾人祈求的時候.
嘗試歡歡喜喜地祈求.
因為從頭一天直到如今.
你們是同心合意的慶王復任.
我深信那在你們心裡動了善功的.
必成全者功.
直到耶穌基督的日子.
我為你們眾人也有這樣的意念.
原是應當的.
因你們常在我心裡.
無論我是在困所之中.
是變命 證實福音的時候.
你們都與我一同得恩.
我體會基督耶穌的心腸.
切切地想念你們眾人.
這是神可以給我作見證的.
我所禱告的.
就是要你們的愛心.
在知識和各樣的見識上多而有多.
使你們能分別是非.
作誠實無過的人.
直到基督的日子.
並靠著耶穌基督.
結滿了仁義的果子.
叫榮耀稱讚歸於神.
願神保守.
祝福常讀他話語.
亦都進行的弟兄姊妹.
弟兄姊妹平安.
在來近五月尾.
到七月其實.

$^{41}$同工都會用肥辣比書.
宣講一個新的系列.
亦都引了在程序表裡.
講到大綱.
叫做「靈性關係成長」.
某程度上亦都呼籲教會.
今年主題是「凝聚動力」.
事實上為何我們會選擇這本書呢?.
肥辣比書.
首先如果你有留意.
保羅在新約書信裡面.
共有十三封.
其中有八封保羅的書信.
保羅都用使徒.
或是奉主的名作使徒.
來稱呼自己.
用一個很高的姿態.
來強調自己.
是行使著作為使徒的權柄.
不過在剛才所讀的肥辣比書.
你會發現.
保羅稱呼自己是.
「瀑人」.
基督耶穌的「瀑人」.
你會看到他一下子.
將自己的身段.
「賴範帝」.
和收信的人.
基本上是同一個水平.
顯示其實保羅寫肥辣比書的意圖.
不是像其他書信那樣.
在教會提出教訓.
指正甚至獨澤.
事實上你會看到.
剛才都隱隱地發現.
肥辣比書其實是一封.
很有溫情,很窩心.
很有溫暖的書信.
可以說是保羅這麼多封書信裡.
寫得最親切,最開心的一封.

$^{81}$你看下去就會發覺.
整本書其實充滿了很多.
保羅對肥辣比書的稱讚和肯定.
一句責備的話.
基本上你不會在肥辣比書裡找到.
事實上很多聖經學者.
都這樣形容保羅和肥辣比信徒的關係.
好像一種亦私亦友的關係.
所以都將肥辣比書.
稱為一封友誼的書信.
相比起保羅所建立的其他教會.
譬如我們認識的哥林多教會.
你會發現.
保羅和哥林多教會的信徒的關係.
相對沒有那麼好.
裡面存在很多問題.
甚至很多的嚴規.
需要保羅透過書信來澄清.
去交代一些誤會.
但你會看到.
保羅和肥辣比信徒的關係.
其實可以說是非常好.
書信裡面佈滿了很多喜樂.
同心,一起事奉的字句.
你會看到.
透過這些內容反映.
保羅和他們之間.
有一份很深厚的友誼.
不只是大家信耶穌.
不是大家去服侍.
他們是有一份主內的友情.
但很特別的是.
如果你細心發現.
保羅和肥辣比教會的信徒.
卻是最少來得多.
他們不是經常在一起.
一方面其實在當時宣教的緣故.
保羅不可能停留在他所建立的教會太久.
他很快建立了教會.
短時間內就要去下一站.

$^{121}$繼續去傳福音.
建立另一間教會.
另一方面保羅寫肥辣比書的時候.
其實聖經學者都相信.
他當時是一個被囚或被軟禁的情況.
身處在羅馬.
基本上他不可能去肥辣比當地.
去埋身牧養這些弟兄姊妹.
但我們又見到這種.
秘密上的阻隔.
無阻他們雙方建立一種.
好像你中有我.
我中有你一份這麼親密的關係.
到底秘訣是什麼呢?.
這就是我們希望在接下來的宣講.
跟大家一起分享.
到底我們除了一起聚會之外.
我們怎樣才可以增潤.
姐妹之間的友誼.
怎樣才可以在主裡.
建立一份持續又健康成長的.
一份屬靈關係呢?.
在進入講道之前.
進入經文之前.
我們先有個祈禱好嗎?.
我們一起求聖靈幫助我們.
預備我們的心.
我們同心祈禱.
天父上帝我們恭敬.
將我們每一個的心.
交在你的恩手.
願聖靈你內處在我們當中.
預備好我們的新心靈.
讓我們藉著你的話語.
得到啟迪.
讓我們的生命得到更新.
讓我們可以更認識.
耶穌基督你的心意.
走近你.
亦與弟兄姊妹彼此更加靠近.

$^{161}$願你幫助我們.
求你隨聽我們的禱告.
奉基督耶穌你的名字來祈求.
阿們.
我們看回經文.
第三至四節保羅怎樣說呢?.
他說.
每逢他上了菲勒比教會.
每逢為他們祈求的時候.
嘗試歡歡喜喜地祈求.
不知道你有沒有試過這種體會.
你一想起.
心目中某個人.
你就會很感恩.
你每次為他祈禱.
想起他就會笑眉眉.
有沒有這個經歷?.
你想起他.
就會很開心.
笑眉眉.
甜在心頭.
有沒有?.
有時好像想不到.
有沒有去過一些婚禮.
當一對新人去行禮的時候.
父母在台下.
甜在心頭.
有沒有見過這些場景?.
捕捉過?.
我見過不少.
因為可能某程度上.
父母看到自己的兒女.
生成了.
成家立室了.
終於可以見到他們.
有自己一頭家.
所以甜在心頭.
笑眉眉這種感覺.
你會想到.
當你想起一個人.

$^{201}$笑眉眉的時候.
一定會想起什麼?.
你跟他有些美好的回憶.
你跟他有一些共同的經歷.
你一想起這個人.
就想起那些開心愉快的畫面.
所以你才會像保羅一樣.
很喜樂,很感恩.
相反.
如果你想起那個人.
就想起他那些不長進的事.
他那些壞事.
怎樣激你的時候.
你會不會甜在心頭?.
大概不會.
你會怎樣?.
氣在心頭.
想起就極為氣.
不要想了,不要看了.
所以你會見到.
保羅對於肥納比教會.
之所以情有獨鍾.
我們說其實有兩方面的原因.
都在經文裡面說了.
首先.
保羅說是因為.
肥納比教會.
從頭一天直到如今.
都是同心合意.
慶王復任.
我們細微這句說話.
什麼是從頭一天直到如今?.
從頭一天直到如今.
如果你看回《使徒行傳》.
他有交代到.
保羅當日創立肥納比教會的時候.
《使徒行傳》第十六章.
當時你猜猜只有多少信徒?.
大概十隻手指數不清.
只有呂底雅一家.

$^{241}$和禁卒一家信了.
就只有這麼多.
但是十多年後.
保羅寫肥納比書的時候.
教會已經變成怎樣?.
有沒有細心看第一節說什麼?.
我們在PowerPoint可以看第一節.
這裡說什麼?.
寫信給繁殖在肥納比.
在基督耶穌裡面的眾聖徒.
諸位監督,諸位執事.
就是說這十年左右.
肥納比教會已經由.
少於十個人變成怎樣?.
變成怎樣?.
變成一間很有規模.
很有體系的教會.
還有一群很有承擔.
願意毀身愛主的信徒.
不同了,不同當日了.
也就是說,雖然過去十幾年.
過去十幾年.
保羅未必有機會親身去.
當地提供牧養指導的工作.
但是肥納比教會長不長進?.
很長進.
他們很懂得自動波.
將保羅當日的教導實踐出來.
他們很自發追求屬靈的成長.
來和心合二,興旺福音.
所以保羅為何這麼喜歡這間教會.
而另一方面.
保羅也透過肥納比教會.
這種增長,這種成長.
他看到背後其實是誰的工作?.
正如他自己所說.
他看到的是耶穌基督的工作.
保羅第六節說.
他知道其實是耶穌基督.
從頭一天開始直到如今.

$^{281}$在肥納比信徒的生命裡動了善功.
誰首先殺種?.
是保羅.
耶穌基督透過保羅.
向當地的信徒殺種.
然後這些種子.
在耶穌基督的恩典下.
在他的帶領下漸漸成長.
到現在保羅說我們見到成果.
教會越來越壯大.
更加多人信主.
保羅不單看過去.
看到現在.
他還在展望將來會是怎樣.
他知道耶穌基督的工作.
不會停在這裡.
我們常說耶穌基督是創始成眾.
保羅因為過去他的殺種.
看到現在肥納比教會的修成.
他也在前瞻將來.
耶穌基督要繼續帶領.
耶穌基督會繼續幫助他們.
完成在地上的福音工作.
直到耶穌基督的日子.
直到耶穌再來的日子.
保羅看到這個長遠的景象.
所以保羅時常為肥納比的信徒.
來歡喜,感恩,禱告.
並不是沒有原因.
是因為肥納比教會.
保羅和耶穌基督是怎樣的?.
三者是同心的.
是連在一條線.
福音的使命使他們.
和耶穌基督連在一起.
並且孕育出一份深厚的主內友誼.
讓我們看到真正健康的索領關係.
是和福事,福音離不開的.
一定是離不開的.
所以為什麼教會.

$^{321}$很多時候鼓勵弟兄姊妹.
不只是個別參與事奉.
不只是參與崇拜.
而是要怎樣?.
你要有一個相交的群體.
你要入伍.
你要和你的組員一起去參與事奉.
實踐福音的工作.
因為唯有這樣.
才會孕育出保羅和肥納比的信徒.
他們這份更健康的索領關係.
我還記得以前.
在自己的舞會裡.
最初信主.
信主後便開始回團契.
很快成為團契職員.
由康樂的職位.
慢慢做副團長,團長.
然後再去蒙召.
我發覺那幾年是一個.
很快成長的經歷.
一起合作事奉.
一起追求成長.
不單在教會裡.
也有對外服侍,科學工作.
或者聯誼,露營,捉蜜魚.
很多很好的回憶畫面.
當然也有些面對困難.
一起拆解人與人之間的糾紛經歷.
一些挫敗經歷.
但就是這樣.
種種的經歷加起來.
我們就連在一起.
我們就一起進入耶穌基督裡面.
所以到現在過了二十年.
即使我們不是經常見面.
但我們也會定期在WhatsApp.
或者電話裡來往.
或者每年我們都會在家庭裡相聚.
一起說一些信仰的近況.

$^{361}$大家事奉成怎樣.
保持著這份主內友情.
所以弟兄姊妹其實某程度上.
是保羅透過他與菲臘比信徒.
這份友誼.
其實也在問我們自己.
我們與弟兄姊妹之間的友情.
其實見機在什麼地方?.
我們在小組,團契裡.
我們有沒有像保羅與菲臘比信徒.
有些共同追求的信仰目標?.
有沒有一些活動.
可以一起實踐福音?.
並且我們時常為這些目標.
來禱告,感恩甚至喜樂.
如果你看整卷的菲臘比書.
你會發覺其實很多時候.
出現一個同工,同心.
一起參與這些字眼.
反覆出現在書信裡.
保羅他想提醒我們.
一份健康的屬靈關係.
其實我們一定要有.
共同追求,努力的目標.
不可以離開這件事.
否則我們未必一定是一個信仰團體.
與其他坊間很多團體都是一樣.
我們之所以分別為性.
是否因為我們有共同的方向使命.
有共同信仰的追求.
我們有沒有這方面的努力目標.
讓上帝的旨意.
讓耶穌基督大使命的吩咐.
成為弟兄姊妹之間.
共同關注,參與.
和努力的一個方向.
我們有沒有讓這些共同目標.
成為我們每次聚會.
其中一個主要分享和喜樂的源頭.
通過一起釋放.

$^{401}$一起參與福音的工作.
讓我們可以連於基督.
成為不單單只是團友.
不單單只是組員.
你和他其實應該是戰友.
是隊友.
因為我們要一起承擔.
耶穌基督大使命的吩咐.
換句話來說.
保羅和肥立比教會.
這種緊密的互動.
其實是測試我們和弟兄姊妹的關係距離.
首先我們每個人都一定贊同.
我們要有屬靈夥伴.
我想沒有人不需要屬靈夥伴.
但在現實生活裡.
我們不妨想一想.
你試試心裡數一數.
有多少弟兄姊妹.
你覺得真的可以稱得上是屬靈的隊友.
真是隊友.
真是戰友.
換句話來說.
你們一定有這些經歷.
一起打過福音的仗.
一起共同經歷艱難.
亦一起給予面子成長.
我們有沒有這些屬靈戰友?.
這是經文問我們的.
其次就是.
假如我們真的想不起來的時候.
我們會否主動邀請對方成為你的屬靈夥伴.
一起去事奉.
一起去追求成長.
並且定期彼此聚在一起禱告.
又或者當你看見對方成長的時候.
你不單止心裡為他高興.
你更加用口去肯定他.
鼓勵他.
以致他可以繼續在信仰裡追求.

$^{441}$相反又或者.
如果真的在事奉的路上.
我們跌低的時候.
我們能否彼此扶持.
彼此守望.
繼續前行呢?.
這一切正正是保羅.
他很想透過他和菲勒比書.
菲勒比信徒的友情.
來啟發我們.
而另一方面我們再看下去.
第七至第八節.
保羅對菲勒比教會情有獨鍾的第二個原因.
是菲勒比教會一直支持.
保羅在外面宣教的工作.
第七節保羅說.
因為你們常在我的心裡.
無論我在困所之中.
在辨明證實福音的時候.
你們都與我一同得恩.
得恩這個字其實原文可以解.
好像團契.
一起有份的意思.
即使保羅不在他們身邊.
他們都是一條心.
他們都是彼此的扶持.
當保羅在羅馬被囚被軟禁的時候.
菲勒比教會主動派人來守望他.
稍後的章節你會看到.
菲勒比教會更加派專人.
送上他們的一些葵贈.
帶給保羅.
並且那個人是留下支援保羅的.
你看到菲勒比教會可以說是保羅宣教工作.
一個我們說好像支援團隊.
一個後勤支援的部隊.
給保羅很多的鼓勵.
很多的支持.
事實上我們說要建立一份緊密的關係.
要維繫它.

$^{481}$我們常說有五個A.
印在大家的港島大綱裡.
你會發現這五個A.
都出現在保羅和菲勒比教會的關係裡.
第一個是affection.
即是說一段好的關係.
不只是工作.
大家是有情感連繫的.
我不知道你和你身邊的弟兄姊妹.
多不多情感交流.
不只是說事情.
也會說自己的一些感受.
很重要的.
因為我們是有情感的.
我們是需要和人建立一份關係.
我們被造本來就是一份關係.
第二個A就是.
emulation.
互相欣賞.
我們的努力是可以得到彼此的肯定.
是可以得到彼此的認同.
我們是可以互相感染.
很重要的.
可以互相豐富我們的生命.
你有沒有發覺.
你和你身邊的弟兄姊妹是不同的.
每個人都有獨特性.
上帝不會創造兩個一模一樣的人.
上帝之所以擺不同的人在我們身邊.
是希望我們可以互相感染.
互相激勵.
第三個A就是advantage.
一段關係裡.
大家都可以互相貢獻.
我們未必有很宏大的成就.
但是你會發覺你在小組裡.
你在團契裡.
你會找到成就感.
你會發揮到你的亮點.
你會發揮到你的價值.

$^{521}$這件事很重要.
正如保羅和菲律賓教會.
他們可以互相欣賞.
互相發揮.
互相激勵對方.
第四個A.
當對方真的有需要的時候.
你會在.
可能你未必馬上趕到現場.
但是你可能一個電話.
即時一個WhatsApp.
表示你會支持對方.
最後那個A稍後就會說.
我們看到其實.
雖然保羅和菲律賓教會的關係.
剛才我們都說了.
是受到地域的限制.
但是他們仍然是連繫著.
是因為耶穌基督在他們中間.
耶穌基督和他的科林使命.
打破了地域限制.
讓他們可以彼此連於基督.
從而雙方連繫著一起.
並且他們雙方都知道.
當我有需要的時候.
對方是會支持我.
當我有需要的時候.
對方是會守護我.
並且對方會記掛著我.
保羅他很體會到.
菲律賓教會對自己這份不離不棄的一份情義.
即使在遠方.
仍然無時無刻支援自己.
支援在困所中的自己.
所以你看到保羅.
雖然他有很多限制.
甚至乎他有很多張力.
但他同時也會出一份很堅韌的生命力.
我記得我看過一篇文章.
他說.

$^{561}$我們都記得可能在二次大戰裡.
很多猶太人被困在集中營.
事件過後.
曾經有人訪問那些.
劫後餘生的災民.
他說你在集中營那幾年.
你是靠著什麼來支持過?.
於是那個人就說.
其實是很絕望的.
不過他每一次想起.
在遠方.
他的家人.
他的太太.
仍然愛他.
仍然記掛著他.
不知為何他仍然有韌力支持下去.
他表示他事實上.
其實我不知道他們是否還在生還.
我只相信.
他們是會愛我.
他們是會記掛我.
就是這樣點點滴滴.
他就撐過了在集中營.
最絕望的那幾年.
這份想像.
這份愛意.
超越了他的艱難.
超越了他的限制.
同樣我們看到.
保羅他之所以能夠.
跨越那個被囚的逆境.
展現出那種堅韌的生命力.
是因為他懂得怎樣?.
他懂得重新塑造.
他正在看的那幅圖畫.
我們說的重新畫面.
他懂得重新看自己的遭遇.
為自己被困的逆境.
製造一個新的意義.
在困難限制的裡面.

$^{601}$保羅不只是看到自己多麼慘.
自己如何被監禁.
他同時看到.
耶穌給他的可能性.
他不能夠出去傳福音.
但是怎樣?.
人們反過來.
是冒名來探望保羅.
為何這個人為了福音.
會如此奮身.
周邊所有的人.
無論猶太人或外邦人.
都主動到他遠近的地方.
來找保羅.
了解他發生了甚麼事.
亦就是因為這樣.
人們主動去找他的時候.
變得反而有很多傳福音的機會.
大把跟人辯明福音的機會.
你會看到保羅在逆境裡.
他懂得重新再想像自己的遭遇.
他不是單純為了一件壞事.
而是看到上帝的帶領.
另一方面.
他在這個逆境.
在缺乏的裡面.
他看到怎樣?.
在遠方.
有人愛他.
有人守望他.
就是這樣.
雖然他現在面對的困難.
未解決到.
但他內心的抗疫.
抗壓的能力是怎樣?.
提升了.
他可以超越被囚的限制.
他可以繼續為上帝事奉.
他可以繼續關心身邊的人.
祝福他們.

$^{641}$不受自己的困難限制.
我們說這是一種.
突破困難超越的生命.
事實上.
姐妹姐妹假如我們在生活裡.
亦都可以慢慢培養超越的目光.
其實有時當我們面對逆境的時候.
不一定要想像得那麼可怕.
我們不一定要.
那些令我們逆境一定是壞事.
一定是不幸.
如果從信仰上的目光來看.
我們可以效法保羅.
你試試換一個看東西的框框.
你試試從新去另一個角度看.
其實困難.
就是我們更加連結上帝.
更加連繫他人一個機會.
你想想.
最困難的時候.
不是困難未解決.
最困難的時候.
你發覺你孤單.
但當你找到同伴.
你發現上帝仍然在懷抱你的時候.
你的盼望就出來了.
你就能夠活出那份.
好像保羅一樣.
我們說的屬靈的韌力.
事實上你會見到.
保羅這種堅韌.
這種抗逆力.
不是純粹來自他個人意志堅強那麼簡單.
更加重要的是.
他這份抗逆力是來自他的信仰.
是建基於他和耶穌的緊密關係.
是建基於他和弟兄姊妹.
那個美好的人際關係.
他知道.
即使現在落入一個困難的裡面.

$^{681}$上帝沒有走開.
上帝容許這個試煉臨到他.
即是說上帝也會和他同在.
他也知道周邊遠方.
仍然有人記掛他.
為他禱告與他同行.
所以雖然是困難.
雖然壓力很大.
但他得到相應的承受能力.
得到同樣的韌力.
所以你看到保羅的心境仍然健康.
他的心靈依然乾淨.
沒有很多的苦毒.
沒有很多的苦澀在他的生命裡面.
他不會影響到他的科研工作.
不會對他的生命構成很多嚴重的威脅.
我還記得以前認識一位弟兄姊妹.
她的婚姻出現了一些問題和配偶.
其實積落了很多年.
當我聽她說的時候.
我發覺真的有很多問題.
不是那麼容易解決.
更重要的是.
她的另一半沒有什麼意欲去解決.
於是弟兄姊妹該怎麼做?.
很難一次過處理這麼多的問題.
當見到一個絕望的弟兄姊妹.
來到你面前的時候.
第一件事該怎麼做?.
很多時候未必是幫她解決問題.
而是你幫她重新去塑造.
看回她現在的經歷.
你幫她有些盼望感.
我嘗試跟那位弟兄姊妹說.
確實是很艱難.
好像沒有什麼希望.
但我鼓勵她相信.
我們今天的對話.
不就是證明上帝仍然在這裡.
因為上帝仍然在這裡.

$^{721}$我們才有機會仍然為了你的婚姻.
去想想辦法.
來對話.
嘗試這樣去重新幫她.
用另一個角度去看回她的遭遇.
讓她本來是無望的感覺.
變回有些盼望.
因為上帝仍然在這裡.
身邊仍然有人支持她.
讓她看到.
本來覺得自己沒有辦法了.
但不是.
只是未找到對的方法.
讓她看到.
她的努力好像白費.
但我說不是.
其實你只是要等上帝給你的時機.
預備好對方.
亦預備好你的心.
就是這樣.
建立她那種韌性.
建立她那種抗疫力.
雖然問題仍然未解決.
但幫她連結上帝.
鞏固她和上帝的關係.
底氣就出來了.
韌力就出來了.
慢慢地.
去熬過最難的日子.
後來她和她的配偶關係.
都慢慢變好.
所以弟兄姊妹.
當我們看到身邊的人.
有困難的時候.
我們怎樣去支援她.
當然如果能夠馬上解決問題.
是絕對好.
但我們也明白現實裡.
很多問題.
不是一時三刻可以處理到.

$^{761}$我們更加需要.
和那些問題共存的一種.
韌性抗疫力.
我們需要幫助你身邊.
有困難的弟兄姊妹.
去換一個框框.
看同一件事.
你要幫助她.
對焦在上帝那裡.
幫助她連結和上帝的關係.
事實上你會看到.
耶穌在釘十架之前.
祂有沒有處理好.
人類所有的問題.
沒有.
仍然有很多人世間的問題.
疾病的問題.
人與人相處的問題.
還沒有解決.
耶穌就釘十架.
耶穌唯一解決的是什麼問題.
就是我們和上帝的關係.
解決了這件事.
換句話來說.
我們人生遇到什麼問題.
出路的源頭是什麼.
你先回到上帝那裡.
你恢復和上帝的正常關係.
所有出路就會慢慢見到.
所有的盼望也會慢慢出來.
你會確信無論現在將來.
上帝仍然在.
並且幫助你.
在看到困難的同時.
你也看到那位超越困難的上帝.
與你同在.
所以當我們見到.
身邊的弟兄姊妹有困難的時候.
你很想幫助他們的時候.
你嘗試幫助他們學習將逆境.

$^{801}$轉化為連結上帝的一個機會.
你嘗試幫助他們明白.
其實有時人生遇到逆境.
可能是上帝想我們.
更加抓緊他們的方法.
從而增加我們身邊弟兄姊妹.
那種屬靈的韌力.
而另一方面你會看到.
第九至十一節.
保羅他為了菲勒比教會的愛心.
來禱告.
原因是我們實踐愛心的時候.
是很需要有智慧.
愛是一定有實踐的部分.
而實踐的時候.
我們需要從上帝而來的.
知識和智慧.
以致我們實踐愛的時候.
好像保羅所說.
首先能夠分辨是非.
不會對人過份地產生溺愛.
縱容.
又或者也不會被人濫用了.
我們的愛心.
另一方面就是.
我們實踐愛的時候.
也需要真誠不虛假.
當我們實踐愛心的時候.
不代表我們完全沒有了自己.
你也要誠心去面對.
自己裡面的反應.
可能有些人你會覺得.
很難去愛.
有時不能強行來.
你需要回到上帝那裡.
重新想起耶穌基督對你的愛.
你才有能力去愛身邊的人.
不能強行靠自己.
來輸出你的愛心.
因為我們愛.

$^{841}$只有神才愛我們.
第三方面.
保羅之所以求上帝.
讓他們有更多.
實踐愛心的智慧和知識.
是因為他們要為上帝.
結更多有好品格的果子.
可以榮耀上帝.
事實上保羅.
你會看到他沒有具體地說.
到底是什麼原因.
以至他為了肥勒比教會的信徒.
求更多有關愛心的知識和智慧.
大概不是因為他們.
在這方面無知.
而是我們看到.
即使肥勒比教會.
在保羅心目中.
已經有很多美好的表現.
即使保羅非常喜愛他們.
但是保羅不滿足於.
肥勒比教會現在的情況.
保羅覺得怎樣?.
你們還有進步的空間.
你們要繼續成長.
所以他求上帝.
幫助這班頂子妹.
繼續將他們的潛能發揮出來.
他求肥勒比教會的愛心.
可以不斷地成長.
更加純正.
更加真誠.
來表達對上帝對人的關愛.
為上帝結更多的果子.
可以榮耀上帝.
這裡亦提醒我們.
一份健康的基督徒關係.
其實要有第五個A.
就是我們和頂子妹互相負責.
意思是.

$^{881}$我們不單止自己追求成長.
我們也有責任.
你要鼓勵你身邊的人成長.
你要向對方這方面問責.
我們要放膽有勇氣.
邀請別人一起成長.
不要任由他.
我最近看過一篇文章.
很有趣.
它說基督徒要怎樣做.
才能保證你的靈命不成長.
我們做了什麼.
會保證我們的靈命不能成長呢?.
留意一下.
它有三方面.
我認同和大家分享.
第一就是.
如果我們不刻意營造.
我們的屬靈生命.
我們就不會成長.
意思是什麼呢?.
當我們覺得信仰生命的成長.
只是憑感覺.
而不是刻意去經營的話.
其實大概我們的信仰歷程.
一定會緩慢甚至倒退.
如果我們認為.
看聖經,祈禱.
甚至是事奉這些基本功.
不是我們優先的選項的話.
大概我們和上帝和頂子妹的關係.
可能會是不難不易.
確實.
屬靈的事實上.
我們說不進就怎樣?.
就會倒退.
不像現在.
很多時候大家會說.
我的利息很高.
你把錢放在這裡.

$^{921}$都會自動有利息.
但屬靈上沒有所謂被動收入.
你不會把屬靈的狀況.
放在這裡十年後.
它會自己增值.
它只會怎樣?.
我們常說.
「食老本」.
想著儲起很多屬靈的經歷.
那些本錢.
然後不動它.
不去增值.
結果就慢慢蠶食你的老本.
所以頂子妹其實.
我們是需要刻意經營.
我們的屬靈生命.
否則我們會慢慢退步.
落後.
和上帝和頂子妹疏遠.
第二方面就是.
我們怎樣才.
容易令自己的靈命不成長呢?.
就是你常常和合適的人在一起.
你只願意和合適的人在一起.
你就一定不會成長.
事實上上帝.
改變我們的生命方式.
很多時候是怎樣?.
在你身邊擺一個.
總是會掩蓋你的人.
有沒有發覺?.
那個可能是你的配偶.
覺不覺得?.
有時我常常勉勵自己.
我老婆要求我成長的時候.
有時我會不斷勉勵自己.
這就是上帝.
透過她幫助我.
是真的.
上帝會刻意將一些.

$^{961}$和你不同甚至相反的.
頂子妹放在你身邊.
那又怎樣?.
就開闊你.
拉闊你的屬靈視野.
以致你不單單.
和同口徑.
合口味的人相處.
你也要學懂.
如何和你不同.
甚至你不是很投契的人.
去學懂相處.
在那裡學習怎樣.
透過愛上帝.
關愛對方.
上帝也會因為這些情況.
你願意降服的時候.
讓你的愛心變得更成熟.
更像他.
我還記得有一次.
和之前沒有的一位.
頂子妹分享.
分享一些在教會生活.
人事上的事.
她最後說了一句話.
你也看得很化.
我不是說看得化.
是看得闊.
很不同吧?.
看得化是一個好笑的事.
但看得闊.
因為上帝.
給不同的人.
有很多不同的彈性.
同樣頂子妹.
我們也要看得闊一點.
在我們和頂子妹相交的時候.
你要努力發掘她的優點.
也可能透過她一些.
你不太舒服的提醒.

$^{1001}$看到自己的盲點.
或者看到自己.
可以成長的地方.
最後一樣.
我們要留意的就是.
我們有時會太過相信.
自己現在已經很OK.
你覺得你的紓困生命.
是否已經很OK?頂子妹.
應該沒有人敢舉手吧?.
有沒有人覺得自己的生命.
已經很OK.
不需要長進?.
事實上.
我們有時真的會.
滿足於現在的紓困狀況.
我們覺得現在很好.
差不多了.
已經可以了.
平平淡淡地過完之後的生活.
但我總是相信.
中國人有句說話.
活到老,學到老.
真的.
我好像也說過這句說話.
我說上帝是很環保的.
如果一個人.
無論他年紀如何.
當上帝說.
你學了所有.
在地上的人生功課的時候.
上帝就會怎樣?.
上帝就會接走你.
省下一些資源.
所以某程度上.
我們今天仍然在這裡.
仍然活著,代表什麼?.
上帝說你有進步的空間.
你有突破的空間.
一定是有些功課.

$^{1041}$上帝要你學會.
無論是和他的關係.
或者和自己.
和身邊的頂子妹的關係.
你是還有成長進步的可能.
所以頂子妹.
不要滿足於你現在的狀況.
你要為你所愛的人.
為你所愛的上帝.
你所愛的人來成長.
也要為那些愛你的人.
為那位愛你的上帝.
為愛你的家人.
頂子妹來成長.
不要辜負上帝和頂子妹.
對你的期望.
最後.
保羅認定.
菲律賓教會是一間.
很有前景的教會.
所以他時常為到.
菲律賓教會的成長.
來禱告.
期望他們長進.
同樣望針.
要成為一間有前景的教會.
我們也需要靠主不斷成長.
並且我們需要同心的.
來彼此禱告守望.
一起參與教會的工作.
讓教會不單在質方面成長.
在數量方面也能成長.
培養出更健康緊密的屬靈關係.
所以.
兩個的呼籲.
我知道這幾年.
有很多的頂子妹.
透過不同的途徑加入了望針.
我們也知道.
當中有些是沒有小組生活的.

$^{1081}$我勉勵你.
不要只崇拜.
慢慢投入在望針.
這個大家庭裡.
你去參與小組.
開始參與事奉.
因為成長關係的建立.
一定在一個群體裡.
不是單打獨鬥.
第二方面.
假如你已經在群體的一部分.
可能已經有你的小組.
有你的團契.
你也要追求鼓勵.
和你身邊的頂子妹一起成長.
不要只顧各個.
我們需要像保羅和肥立比教會.
將我們的生命連於基督.
連於科林的使命.
我們才能發揮到.
培育到那份健康緊密的屬靈關係.
我們一起祈禱吧.
因著我願你的靈與我們同在.
讓我們再次藉著肥立比書.
去檢視我們和耶穌基督的關係.
和頂子妹的關係.
讓我們透過保羅的提醒.
以及他和肥立比教會.
頂子妹的互動.
讓我們再次去重新去檢視.
到底我們和頂子妹.
和上帝的距離.
在我們的關係裡.
我們有沒有讓福音成為.
我們最基本最重要的元素.
以致我們透過一起去事奉.
我們首先連於基督.
一起去緊靠你.
以致因為這份主內的友情.
我們的關係可以變得更加健康.

$^{1121}$更加緊密.
上帝幫助我們.
我們知道確實生命裡有很多拉扯.
很多困難.
但在這些困難裡.
讓我們看到上帝你的同在.
你可以幫助我們.
突破這些困難.
轉化成為可以親近你.
靠近頂子妹的途徑.
願意上帝你幫助我們.
堂讓我們身邊仍然沒有加入小組團契的頂子妹.
願聖靈領你預備他們的心.
讓他們可以進一步走近教會.
走近我們這個群體.
也都讓小組團契裡每一位頂子妹.
我們能夠更加的同心.
一起在福音的事上有份.
藉此我們不單單成為團友,組員.
我們更加在主理裡成為戰友,成為隊友.
一起為耶穌基督打拼.
一起在福音的事上彼此同心守望.
我們今晚的禱告交託.
奉基督耶穌的名字來祈求.
Amen.
\newpage



\section{希伯來書 12:1-3}
\label{sec:OamF1FcoO7k}
\textbf{2020年4月10日 受苦節崇拜直播｜希伯來書12\_1-3}
\newline
\newline
連結: \href{https://youtube.com/watch?v=OamF1FcoO7k}{\texttt{https://youtube.com/watch?v=OamF1FcoO7k}} ~~~~ 語音日期: 2024-06-06
\newline
\newline
\hyperref[sec:1HAPF8xhKJg]{\small{< < < PREV SERMON < < <}}
~
\hyperref[sec:index_chronic]{\small{[返順時目]}}
~
\hyperref[sec:9ogZwUemHN4]{\small{> > > NEXT SERMON > > >}}
\newline
\newline
希伯來書 12:1-3
\newline
\begin{longtable}{cl}
\hline
\hline
章節 & 經文 (和合本修訂版)\\
\hline
12:1 & \begin{tabularx}{0.7\textwidth}{X} 所以，既然我們有這許多見證人如同雲彩圍繞著我們，就該卸下各樣重擔和緊緊纏累的罪，以堅忍的心奔那擺在我們前頭的路程， \end{tabularx} \\ \\ \relax
12:2 & \begin{tabularx}{0.7\textwidth}{X} 仰望我們信心的創始成終者耶穌，他因那擺在前面的喜樂，輕看羞辱，忍受了十字架的苦難，如今已坐在神寶座的右邊。 \end{tabularx} \\ \\ \relax
12:3 & \begin{tabularx}{0.7\textwidth}{X} 你們要仔細想想這位忍受了罪人如此頂撞的耶穌，你們就不致心灰意懶了。 \end{tabularx} \\ \\
[1ex]
\hline
\hline
\end{longtable}
$^{1}$在希伯來書十二章一至三節.
裡面是這樣記載.
我們既有這許多的見證人.
如同運彩圍著我們.
就當放下各樣的重擔.
脫去容易纏累我們的罪.
全心忍耐.
奔那擺在我們前頭的路程.
仰望為我們信心創始成眾的耶穌.
他因擺在那前面的喜樂.
就輕看羞辱.
忍受了十字架的苦難.
便坐在神寶座的右邊.
那些受罪人這樣頂撞的.
你們要思想.
免得疲倦灰心.
接著下來有陳明荃牧師.
透過這段經文來勉勵我們.
請陳明荃牧師.
在這段經文裡面.
提及到受苦者的苦難.
和受苦者的痛苦.
請陳明荃牧師.
在這段經文裡面.
提及到受苦者的苦難.
和受苦者的痛苦.
請陳明荃牧師.
在這段經文裡面.
提及到受苦者的苦難.
和受苦者的痛苦.
請陳明荃牧師.
在這段經文裡面.
提及到受苦者的苦難.
和受苦者的痛苦.
請陳明荃牧師.
在這段經文裡面.
提及到受苦者的苦難.
和受苦者的痛苦.
請陳明荃牧師.
在這段經文裡面.

$^{41}$提及到受苦者的苦難.
和受苦者的痛苦.
請陳明荃牧師.
在這段經文裡面.
提及到受苦者的苦難.
和受苦者的痛苦.
請陳明荃牧師.
在這段經文裡面.
提及到受苦者的苦難.
和受苦者的痛苦.
請陳明荃牧師.
在這段經文裡面.
提及到受苦者的苦難.
和受苦者的痛苦.
請陳明荃牧師.
在這段經文裡面.
提及到受苦者的苦難.
和受苦者的痛苦.
請陳明荃牧師.
在這段經文裡面.
提及到受苦者的苦難.
和受苦者的痛苦.
請陳明荃牧師.
在這段經文裡面.
提及到受苦者的苦難.
和受苦者的痛苦.
請陳明荃牧師.
在這段經文裡面.
提及到受苦者的苦難.
和受苦者的痛苦.
請陳明荃牧師.
在這段經文裡面.
提及到受苦者的苦難.
和受苦者的痛苦.
請陳明荃牧師.
在這段經文裡面.
提及到受苦者的苦難.
和受苦者的痛苦.
請陳明荃牧師.
在這段經文裡面.

$^{81}$提及到受苦者的苦難.
和受苦者的痛苦.
請陳明荃牧師.
在這段經文裡面.
提及到受苦者的苦難.
和受苦者的痛苦.
請陳明荃牧師.
在這段經文裡面.
提及到受苦者的苦難.
和受苦者的痛苦.
請陳明荃牧師.
在這段經文裡面.
提及到受苦者的苦難.
和受苦者的痛苦.
請陳明荃牧師.
在這段經文裡面.
提及到受苦者的苦難.
和受苦者的痛苦.
請陳明荃牧師.
在這段經文裡面.
提及到受苦者的苦難.
和受苦者的痛苦.
請陳明荃牧師.
在這段經文裡面.
提及到受苦者的苦難.
和受苦者的痛苦.
請陳明荃牧師.
在這段經文裡面.
提及到受苦者的苦難.
和受苦者的痛苦.
請陳明荃牧師.
在這段經文裡面.
提及到受苦者的苦難.
和受苦者的痛苦.
請陳明荃牧師.
在這段經文裡面.
提及到受苦者的苦難.
和受苦者的痛苦.
請陳明荃牧師.
在這段經文裡面.

$^{121}$提及到受苦者的苦難.
和受苦者的痛苦.
請陳明荃牧師.
在這段經文裡面.
提及到受苦者的苦難.
和受苦者的痛苦.
請陳明荃牧師.
在這段經文裡面.
提及到受苦者的苦難.
和受苦者的痛苦.
請陳明荃牧師.
在這段經文裡面.
提及到受苦者的苦難.
和受苦者的痛苦.
請陳明荃牧師.
在這段經文裡面.
提及到受苦者的苦難.
和受苦者的痛苦.
請陳明荃牧師.
在這段經文裡面.
提及到受苦者的苦難.
和受苦者的痛苦.
請陳明荃牧師.
在這段經文裡面.
提及到受苦者的苦難.
和受苦者的痛苦.
請陳明荃牧師.
在這段經文裡面.
提及到受苦者的苦難.
和受苦者的痛苦.
請陳明荃牧師.
在這段經文裡面.
提及到受苦者的苦難.
和受苦者的痛苦.
請陳明荃牧師.
在這段經文裡面.
提及到受苦者的苦難.
和受苦者的痛苦.
請陳明荃牧師.
在這段經文裡面.

$^{161}$提及到受苦者的苦難.
和受苦者的痛苦.
請陳明荃牧師.
在這段經文裡面.
提及到受苦者的苦難.
和受苦者的痛苦.
請陳明荃牧師.
在這段經文裡面.
提及到受苦者的苦難.
和受苦者的痛苦.
請陳明荃牧師.
在這段經文裡面.
提及到受苦者的苦難.
和受苦者的痛苦.
請陳明荃牧師.
在這段經文裡面.
提及到受苦者的苦難.
和受苦者的痛苦.
請陳明荃牧師.
在這段經文裡面.
提及到受苦者的苦難.
和受苦者的痛苦.
請陳明荃牧師.
在這段經文裡面.
提及到受苦者的苦難.
和受苦者的痛苦.
請陳明荃牧師.
在這段經文裡面.
提及到受苦者的苦難.
和受苦者的痛苦.
請陳明荃牧師.
在這段經文裡面.
提及到受苦者的苦難.
和受苦者的痛苦.
請陳明荃牧師.
在這段經文裡面.
提及到受苦者的苦難.
和受苦者的痛苦.
請陳明荃牧師.
在這段經文裡面.

$^{201}$提及到受苦者的苦難.
和受苦者的痛苦.
請陳明荃牧師.
在這段經文裡面.
提及到受苦者的苦難.
和受苦者的痛苦.
請陳明荃牧師.
在這段經文裡面.
提及到受苦者的苦難.
和受苦者的痛苦.
請陳明荃牧師.
在這段經文裡面.
提及到受苦者的苦難.
和受苦者的痛苦.
請陳明荃牧師.
在這段經文裡面.
提及到受苦者的苦難.
和受苦者的痛苦.
請陳明荃牧師.
在這段經文裡面.
提及到受苦者的苦難.
和受苦者的痛苦.
請陳明荃牧師.
在這段經文裡面.
提及到受苦者的苦難.
和受苦者的痛苦.
請陳明荃牧師.
在這段經文裡面.
提及到受苦者的苦難.
和受苦者的痛苦.
請陳明荃牧師.
在這段經文裡面.
提及到受苦者的苦難.
和受苦者的痛苦.
請陳明荃牧師.
在這段經文裡面.
提及到受苦者的苦難.
和受苦者的痛苦.
請陳明荃牧師.
在這段經文裡面.

$^{241}$提及到受苦者的苦難.
和受苦者的痛苦.
請陳明荃牧師.
在這段經文裡面.
提及到受苦者的苦難.
和受苦者的痛苦.
請陳明荃牧師.
在這段經文裡面.
提及到受苦者的苦難.
和受苦者的痛苦.
請陳明荃牧師.
在這段經文裡面.
提及到受苦者的苦難.
和受苦者的痛苦.
請陳明荃牧師.
在這段經文裡面.
提及到受苦者的苦難.
和受苦者的痛苦.
請陳明荃牧師.
在這段經文裡面.
提及到受苦者的苦難.
和受苦者的痛苦.
請陳明荃牧師.
在這段經文裡面.
提及到受苦者的苦難.
和受苦者的痛苦.
請陳明荃牧師.
在這段經文裡面.
提及到受苦者的苦難.
和受苦者的痛苦.
請陳明荃牧師.
在這段經文裡面.
提及到受苦者的苦難.
和受苦者的痛苦.
請陳明荃牧師.
在這段經文裡面.
提及到受苦者的苦難.
和受苦者的痛苦.
請陳明荃牧師.
在這段經文裡面.

$^{281}$提及到受苦者的苦難.
和受苦者的痛苦.
請陳明荃牧師.
在這段經文裡面.
提及到受苦者的苦難.
和受苦者的痛苦.
請陳明荃牧師.
在這段經文裡面.
提及到受苦者的苦難.
和受苦者的痛苦.
請陳明荃牧師.
在這段經文裡面.
提及到受苦者的苦難.
和受苦者的痛苦.
請陳明荃牧師.
在這段經文裡面.
提及到受苦者的苦難.
和受苦者的痛苦.
請陳明荃牧師.
在這段經文裡面.
提及到受苦者的苦難.
和受苦者的痛苦.
請陳明荃牧師.
在這段經文裡面.
提及到受苦者的苦難.
和受苦者的痛苦.
請陳明荃牧師.
在這段經文裡面.
提及到受苦者的苦難.
和受苦者的痛苦.
請陳明荃牧師.
在這段經文裡面.
提及到受苦者的痛苦.
請陳明荃牧師.
在這段經文裡面.
提及到受苦者的苦難.
和受苦者的痛苦.
請陳明荃牧師.
在這段經文裡面.
提及到受苦者的苦難.

$^{321}$請陳明荃牧師.
在這段經文裡面.
提及到受苦者的苦難.
請陳明荃牧師.
在這段經文裡面.
提及到受苦者的苦難.
請陳明荃牧師.
在這段經文裡面.
提及到受苦者的苦難.
請陳明荃牧師.
在這段經文裡面.
提及到受苦者的苦難.
請陳明荃牧師.
在這段經文裡面.
提及到受苦者的苦難.
請陳明荃牧師.
在這段經文裡面.
提及到受苦者的苦難.
請陳明荃牧師.
在這段經文裡面.
提及到受苦者的苦難.
請陳明荃牧師.
在這段經文裡面.
提及到受苦者的苦難.
請陳明荃牧師.
在這段經文裡面.
提及到受苦者的苦難.
請陳明荃牧師.
在這段經文裡面.
提及到受苦者的苦難.
請陳明荃牧師.
在這段經文裡面.
提及到受苦者的苦難.
請陳明荃牧師.
在這段經文裡面.
提及到受苦者的苦難.
請陳明荃牧師.
在這段經文裡面.
提及到受苦者的苦難.
請陳明荃牧師.

$^{361}$在這段經文裡面.
提及到受苦者的苦難.
請陳明荃牧師.
在這段經文裡面.
提及到受苦者的苦難.
請陳明荃牧師.
在這段經文裡面.
提及到受苦者的苦難.
請陳明荃牧師.
在這段經文裡面.
提及到受苦者的苦難.
請陳明荃牧師.
在這段經文裡面.
提及到受苦者的苦難.
請陳明荃牧師.
在這段經文裡面.
提及到受苦者的苦難.
請陳明荃牧師.
在這段經文裡面.
提及到受苦者的苦難.
請陳明荃牧師.
在這段經文裡面.
提及到受苦者的苦難.
請陳明荃牧師.
在這段經文裡面.
提及到受苦者的苦難.
請陳明荃牧師.
在這段經文裡面.
提及到受苦者的苦難.
請陳明荃牧師.
在這段經文裡面.
提及到受苦者的苦難.
請陳明荃牧師.
在這段經文裡面.
提及到受苦者的苦難.
請陳明荃牧師.
其實我們今天.
我們所說的罪.
大部分其實我們都知道的.
我們自己知道.

$^{401}$自己的軟弱.
我們知道哪些是應該要做.
哪些是不應該做.
不過罪就是有一種.
那種魔力.
或者那種.
令到我們跌倒的那種.
核制的能力.
我們很想去做的事.
我們沒有辦法去做.
我們想做的就一路一路.
好像引人入勝那樣.
令到我們進入罪的圈套裡面.
聖經說.
如果我們一意孤行.
我們明知固犯.
我們不靠基督的時候.
我們就好像將耶穌.
重新釘在十字架裡面.
將地上罪的.
帶來羞辱再一次.
加諸在耶穌的身上.
我們深硬的表現.
到最後我們成為一個.
不孤行的心.
這個信仰再沒有辦法在我們生命裡面.
發揮果效.
基督的救恩.
我們悔改的心.
再不能夠大到在我們上帝的根前.
曾經認識過一些弟兄姊妹.
我預備這個時候.
我也想起一些.
曾經在教會裡面.
參與過的聚會.
甚至乎一段很長的日子.
因為生命的罪.
他覺得沒有面目.
在這位上帝.
或者因為這些罪.

$^{441}$核製著他的心靈.
他自己提不起勁.
再回到教會當中.
他覺得愧對身邊的弟兄姊妹.
更加愧對.
在十字架上的耶穌基督.
就因為.
這種態度.
令到他越走越遠.
可能在教會裡面.
很熱心服侍的.
就因為這些罪的緣故.
他跌倒.
他不尋求神.
於是他離開上帝.
越走越遠.
離開這個救恩.
就好像希伯來書的作者說.
這些信徒就隨流失去.
弟兄姊妹.
我們生命裡面.
我們必須要正視那些.
令到我們離棄神的重擔.
以及那種.
纏累我們生命的罪.
如果我們不是.
在上帝面前再乞求.
那個恩典.
到最後我們都會成為沉淪的人.
我們會離開上帝.
會越走越遠.
希伯來書的作者.
第一提醒我們.
要將這些重擔和罪.
除下.
讓我們可以輕生.
跑這條信仰的路.
在第二節裡面.
再繼續說.
同樣繼續出現忍受這個字.

$^{481}$仰望各位.
為我們信心創始成中的耶穌.
他因擺在.
前面的喜樂.
就興翰羞辱.
忍受了十字架的.
患難.
便坐在神寶座的右邊.
在這裡面特別提出.
一個很重要的字眼.
就是創始成中.
創始成中成為.
耶穌基督的一個重要的形容詞.
什麼意思呢.
創始和成中.
是分開解釋的.
創始的意思是.
耶穌成為一個.
開啟新道路的先鋒.
他走過的路.
成為我們接下來.
每一個跟隨他的人.
好像一個開路先鋒一樣.
將面前的經濟.
和疾禮.
破除了這些障礙.
他首先自己走過這條路.
跟隨的人.
按著他的腳踏.
一步一步走這條信仰的路.
創始的意思就是.
他成為我們生命的榜樣.
成為我們生命裡面的先鋒.
成中的意思.
其實是說.
更加好的字眼.
就是用成全這個意思.
耶穌不單止.
為我們作開路的先鋒.
耶穌也將.

$^{521}$我們的生命變得完全.
在我們生命最軟弱的地方.
他的恩典.
會陪我們一起離開同行.
所以我們信仰的路.
不單止是說耶穌.
走在我們前面.
耶穌也走在我們旁邊.
他為我們打氣.
他扶持我們.
扶起我們.
讓我們走這條路的時候.
不會只覺得自己孤身一人.
耶穌更加.
走在我們的後面.
耶穌的恩典.
令我們不要再繼續走.
自己昔日未信主的回頭路.
他在後面.
來推動我們.
讓我們繼續按著.
上帝的心意奔走這條.
屬天的道路.
創始成中的耶穌.
就是這樣來帶領我們.
經文再繼續.
來講.
創始成中.
成全我們救恩的基督耶穌.
他怎樣看十字架.
這個苦難.
可能我們今天.
去思想十字架.
或者思想耶穌受難的時候.
我們也有很多的悲情.
我們心裡面.
也有一種很沉鬱的體會.
我相信.
每一個愛耶穌的人.
每一個真是很堅持.

$^{561}$很明白耶穌心腸的人.
我們為著耶穌的受苦.
我們心裡面.
也有這種體會.
不過在經文裡面.
他講到耶穌上十架並不是.
單純的痛苦苦澀.
或者覺得自己很可憐的角度.
來走這條十架的道路.
在經文裡面.
他特別講.
耶穌因為擺在前面的.
那個喜樂.
他就輕看那些羞辱.
忍受了十字架的苦難.
便坐在神.
寶座的右邊.
如果只是想.
耶穌只是單純.
聚焦十字架的痛苦.
那條路是很難走的.
不過經文裡面.
他特別強調.
耶穌聚焦的.
是十字架苦難之後.
救贖的喜樂.
耶穌看見的.
不是單純在眼前的.
那種苦況.
他的視野超越了.
眼前一切的苦難.
他的視野.
是在講將來.
屬天的那份.
為上帝.
為這個世界預備救恩.
而得到的那份喜樂.
雖然耶穌是面對.
困難,面對十字架.
那種痛苦.

$^{601}$不過他的心靈裡面.
希伯來書的作者講到.
他方言前的喜樂.
大於眼前的痛苦.
以至到眼前.
所有冰冰對他的靈慾.
或者在十字架上面.
面對所有的艱難.
他都可以一一輕視.
輕看這些羞辱.
讓他有力量.
走面前的路.
希伯來書的作者.
他為我們提出.
一個很重要的錦囊.
如果我們今天在我們的生活裡.
我們只是覺得基督徒的生命.
只是剩下苦難的時候.
我們的人生.
我們會覺得.
我們很被迫地走我們自己信仰的路.
我們覺得做每一件事.
都是很艱難,很痛苦.
因為我們只是.
看到受苦的盡頭.
都是繼續受苦.
不過基督教信仰告訴我們.
苦難不是.
人生命的終結.
上帝的福樂.
才是我們生命的總結.
將來我們在得贖的日子裡.
我們可以勝過.
這些苦難.
我們眼前要看的.
除了眼前這些不同的.
困難,阻隔之外.
我們的目光視野.
我們更加要看到.
從神而來的那份福樂.

$^{641}$我們看到這些福樂的時候.
我們就可以知道.
今時今刻.
或者眼前這些困難.
只是短暫.
上帝還有更長的.
美好的福氣.
上帝更加大的喜樂.
屬天的兒女.
經文裡面再繼續說.
因為耶穌.
看見眼前的那種福樂.
他能夠.
忍受十字架的苦難.
在這裡我們值得.
繼續思考.
什麼叫忍受十字架的苦難呢?.
很多時候我們去說.
這一段耶穌忍受十字架的時候.
第一件事我們想起的.
就是耶穌在皮肉上面.
面對的艱辛.
我們知道.
羅馬.
那些兵兵用的.
那些鞭.
或者那些酷刑.
全部都是非人化的.
全部都是要為人帶來很多的折磨.
那個折磨是.
肉身的那種折磨.
那些鞭有勾的打下去.
就要令到皮開肉爛.
那個十字架.
其實令到一個人死亡.
其實有很多的方法.
但是這個十字架卻是在竹柳地.
在這個葛國他山上.
將一個犯人.
在一個高山上.

$^{681}$還要懸掛一個十字架.
來張開他的死狀.
其實是一種.
很大的.
心靈和肉體上的折磨.
十字架的苦難.
第一件事就是.
身心的那種折磨.
不過對於耶穌基督來說.
他不單止面對這一種.
身心那種折磨的苦難.
在耶穌作為神的兒子.
這一處.
那個十字架的苦難.
更加深一層的就是.
這個滿有榮耀.
這個.
掌管萬有.
神的兒子.
在地上裡面.
他滿有能力來施行神跡.
其事.
不過在這一刻裡面.
他成為一個.
手無寸鐵.
他沒有辦法.
來改變眼前.
這些困難的.
一個被釘的犯人.
他沒有犯過任何的罪.
不過在這一刻裡面.
神的兒子.
滿有大能的兒子.
竟然成為一個最軟弱的姿態.
被釘在.
十字架上.
當耶穌在.
十字架上被釘的時候.
身邊有很多的.
百姓或者有些人.

$^{721}$在下面來叫喊.
如果這個是.
神的兒子.
為何會進去行神跡.
他這一刻不能夠救自己.
如果是你.
或者我是耶穌.
我這一刻裡面.
我就會變走所有的苦難.
馬上就用我.
最強人的姿態.
去扭轉所有的苦難.
這個神跡夠偉大了.
不過耶穌.
沒有在十字架上.
施行任何的神跡.
一個滿有大能的.
神的兒子.
收起他自己的榮光.
以神的能力.
為同等來強奪.
他用最軟弱的姿態.
不用神跡.
甚至面對著.
下面每一個村民.
百姓.
說一些調侃的說話.
他默然.
來承受.
這些一切冷嘲熱諷.
作為神的兒子.
還有第三方面.
這個.
帶著.
上帝拯救計劃的.
神的兒子.
他一直以來都是將上帝的福音.
將神的使命.
他用他自己的愛.
去感染更新改變.

$^{761}$所有跟隨他.
或在這個世界裡所有的人.
不過.
在.
他釘十字架的時候.
下面的.
每一個.
昔日曾經.
慾勇他這些群眾.
他們高聲情願釋放.
重犯巴拉伯.
甚至面對著.
他自己那種.
血肉痛苦.
下面的人.
繼續帶著.
猙獰的面貌.
將耶穌移得快快處死.
耶穌沒有做過.
任何傷害他們的事.
但這群人就.
揮拳.
很高氣.
那種.
帶著很大的情緒.
來將耶穌置之死地.
面對這群.
頂撞的人.
耶穌默然地來忍受.
他沒有差遣.
天使天君.
直接將這些憎恨的人一一殲滅.
他只是.
承受十字架.
為世人.
所流的血.
更加重要的是.
面對一群這麼憎恨的人.
一群這樣的罪人.
耶穌在十字架上說.

$^{801}$「父啊.
請赦免他們.
因為他們所作的.
他們不知道」.
面對苦害他們的人.
耶穌不單沒有.
帶著報復的心.
來對待他們.
他更加在十字架上.
為他們每一個代求.
祈求天父的憐憫.
祈求天父的赦免.
臨到他們每一個當中.
最後忍受的是什麼呢?.
忍受就是.
上帝的旨意.
好像很荒謬一樣.
在這個罪惡滿盈的世界裡.
神的愛竟然這麼大.
神的愛竟然.
讓自己的.
獨生子耶穌走上.
這條十字架的苦路.
當耶穌.
在馬利園.
祈禱之後.
他要信服上帝的時候.
他就信服至死.
在十字架上.
他用最軟弱的方式.
用最信服的方式.
來表達神的愛.
他張開自己的手.
用最軟弱的方式.
向這個世界宣告.
上帝最堅定的愛.
最荒謬的方法.
竟然用上帝自己的生命.
上帝自己的生命.
換取地上每一個人的生命.

$^{841}$耶穌.
忍受這一切的痛楚.
滿有權能.
神的兒子.
他將他的生命.
傾刻流出他的.
寶血.
來挽回我們地上的每一個.
可能弟姐妹你會說.
他是神來的.
他能夠這樣忍受.
我們是血肉之軀.
我們如何忍受這樣的信仰迫迫.
或者我們信仰的實踐.
沒錯.
我們只不過是血肉之軀.
耶穌.
他的層次.
或者他的生命.
跟我們真的無法比.
不過希伯來書的作者說.
當我們想到耶穌基督的生命.
是我們的先鋒的時候.
其實我們今天.
我們所走過的路.
對他來說輕省了很多.
在第十二章第四節裡面.
他講到我們今天.
我們征戰都沒有流血的地步.
為什麼我們面對著.
眼前的挑戰.
耶穌這麼艱難.
他都沒有放棄.
為什麼我們今天這麼輕易.
就放棄了我們自己的信仰.
寧願將自己的生命放逐.
寧願將自己隨流失去.
我們都不堅持繼續我們這個信仰.
耶穌成為我們的先鋒.
他走過比我們的路更加艱辛.

$^{881}$我們只是走一條比他容易的路.
他的愛.
為了我們生命每一個.
他為了我們的緣故.
他將最難的.
最荒謬的.
傾倒在他自己的身上.
反而.
他勸勉每一個在地上跟隨他的人.
信仰的路是不容易走.
但是我們可以放心.
耶穌基督的生命.
是我們每一個人的生命.
我們可以放心.
耶穌基督已經勝過世界.
我們都可以靠著他的恩典.
勝過眼前一切的困難.
不要為了眼前的困難.
我們就成為我們.
前類我們自己不能夠繼續走信仰的路的.
這些絆腳石.
我們看看這位基督耶穌.
他為我們開路.
他為我們披荊斬棘.
他也為我們預備.
我們應該要走的路.
十二章第三節.
和合本這樣說.
那忍受罪人這樣鼎壯的.
你們要思想.
免得疲倦灰心.
在這裡我們讀的時候.
我們很多時候.
我們很多時候將主和客的位置.
混亂了以至我們誤會了和合本.
或者原本聖經的意思.
當我們讀的時候.
那忍受罪人鼎壯的.
我們很多時候讀的時候我們會覺得.
我們將我們自己被人鼎壯的那種經歷.

$^{921}$那種經歷.
那些不好受就套到耶穌身上.
意思大概是我們都被人鼎壯.
耶穌也被人鼎壯.
所以我們的焦點就放在.
我們被人鼎壯的經驗裡面.
不過更加多的譯本.
特別是最近和合修定本.
它也這樣譯的.
經文是這樣記載的.
你們要仔細想想這位.
忍受了罪人如此鼎壯的耶穌.
你們就不自心灰意冷了.
經文的作者他想告訴我們.
我們今天眼睛的注視.
不是放在.
多少人鼎壯我.
而是將我的鼎壯和.
耶穌的被鼎壯成為一個對等.
將我們自己的經驗.
讀了在聖經裡面.
希伯來書的作者.
他要調教我們的焦點.
今天我們看的不是.
有多少人鼎壯我.
是看世界上有多少人.
鼎壯的那位耶穌.
這個.
剛才說的忍受很多艱難.
忍受十字架上.
很多的苦難的.
那位主.
他的一生其實都是面對著.
一個逆流而上.
將上帝福音.
宣揚出去的那種生命.
當我們看到.
他這樣走的時候.
同樣我們.
看見那位耶穌.

$^{961}$他走過的路.
我們心裡知道.
當我們比較下去的時候.
蚊鼻和牛鼻.
我們所經歷的.
算不得什麼.
我們.
和耶穌所走過的路.
最難的其實.
由他已經走過.
我們只是跟著.
他的腳步一步一步地走.
今天或者我們.
很安穩的社會裡面.
我們只是.
在我們的生活裡面.
堅持多一些.
放下那些生命裡面.
面對我們的罪.
當我們能夠.
將自己的焦點.
從我們自己覺得逆流順受.
很多我們覺得.
很不好的經歷放大了.
覺得耶穌都在經歷這些.
我們轉換.
我們的眼光.
我們看見原來.
生活裡面所有的艱難.
耶穌都經歷過.
我們看他怎樣走.
我們就重新調教我們今天.
走這條信仰的路.
這條路是輕散的.
這條路是需要.
我們堅持.
但絕對沒有耶穌堅持的程度.
那麼高.
經文更加提醒我們.
頂撞.

$^{1001}$耶穌的罪人.
不是外面身邊的那些人.
頂撞耶穌的罪人其實都是我們有份.
在我們人生裡面.
我們都是曾經這樣來頂撞他.
或者我們今天.
在我們今天的信仰實踐裡面.
我們都是成為一個.
頂撞耶穌.
不明白他的心懷.
甚至乎.
安著自己的心意.
來走自己信仰的路的人.
聖經勉勵我們.
我們不是成為一個這樣.
和上帝越走越遠的.
反而我們要跟隨基督的腳蹤.
他跑他的路.
我們跟著他跑.
這種忍耐.
不是帶著絕望.
來繼續承受信仰的代價.
這種忍耐.
是有曙光.
因為基督在.
他自己受苦的過程裡面.
他知道面前將會有.
曙光在他生命當中.
這一份的喜樂.
令他興翰眼前的羞辱.
同樣我們每一個.
我們都要.
看著眼前.
我們生命都會繼續有曙光.
黑暗的盡頭.
就是上帝恩典的來臨.
我們就繼續跑.
這條信仰的路.
這種忍耐.
是會有盡頭.

$^{1041}$這個忍耐.
是會帶領我們生命.
有一種豁然開朗的靈性的建立.
這種忍耐.
是會讓我們成為.
一個更加成熟的人.
以致我們的信仰.
可以不單止自己去跑.
我更加可以成為.
我們後繼每一個人.
跑的一個先鋒.
更加重要的是.
因為這條路不是絕望的路.
所以我們.
沒有任何的資格.
放棄這個信仰.
我們更加沒有任何的資格.
來放棄.
這個.
付上沉重代價.
愛我們的.
基督耶穌.
無論我們經歷什麼處境.
我們仰望祂.
我們看著祂.
來到做人.
看著祂昔日怎樣走過的路.
成為我們今天.
效法的路.
弟兄姊妹.
今天希伯來書.
給我們一個很重要的提醒.
可能我們.
和當時希伯來書的讀者一樣.
可能我們.
雙手已經垂下來.
我們的腳.
變得酸軟.
甚至乎在聖經裡面.
講到那些人跑信仰的路.

$^{1081}$已經好像脫離了教一樣.
不可行.
但是當我們思念.
這位十字架上.
為我們犧牲的主.
我們就有力量將我們.
發酸的.
垂下來的手提起來.
將我們.
沒有力量的雙腿.
重新站立起來.
當我們覺得很多外來的挑戰.
令我們沒辦法去走的時候.
我們仰望這個主的時候.
我們.
有力量奔走.
面前的路.
十字架.
不單止是令到我們.
提升我們對上帝.
那種肅然起敬.
十字架更加不是.
純粹令我們流下.
覺得耶穌很痛苦的眼淚.
十字架更加是勉勵我們.
在這個充滿逆流.
充滿挑戰.
充滿試探引誘的世界裡面.
我們基督徒.
可以挺起腰骨.
仰望基督耶穌.
奔走前面信仰的路.
十架的基督.
繼續對我們每一個講說話.
願十字架的耶穌.
祂的愛.
時刻激勵我們.
我們同心祈禱.
基督耶穌.
竟然你為我們流血.

$^{1121}$捨身在十架上.
當我們都不知道.
自己是這麼寶貴的時候.
你看我們是最珍貴的.
你用你的寶血.
潔淨我們.
天父啊.
我們心靈裡面.
或者我們在家裡面的弟兄姊妹.
如果我們知道我們自己聖靈裡面有罪.
我們向你.
乞求你的憐憫.
求你寬恕.
求你潔淨.
求你赦免.
可能在我們生活當中.
可能有很多東西阻礙了我們奔跑.
這條信仰的路.
我們想放棄.
對你已經不冷不熱.
求你的愛喚醒我們的心靈.
讓我們重新思考.
你在十架裡面的愛的時候.
讓我們人生沒有力量.
再一次將我們每一個.
我們的心思意念交託給你.
幫助我們.
賜給我們有力量.
賜給我們生命裡面.
有你的愛.
口口充滿我們.
再一次多謝耶穌.
令愛我們.
愛我們到底.
奉基督耶穌.
得勝的聖名祈求.
\newpage



\section{腓立比書 1:27-2:11}
\label{sec:9ogZwUemHN4}
\textbf{2024年6月16日父親節崇拜直播｜陳冠成牧師｜靈性·關係·成長－與身份相稱的生活｜腓立比書1\_27-2\_11}
\newline
\newline
連結: \href{https://youtube.com/watch?v=9ogZwUemHN4}{\texttt{https://youtube.com/watch?v=9ogZwUemHN4}} ~~~~ 語音日期: 2024-06-16
\newline
\newline
\hyperref[sec:OamF1FcoO7k]{\small{< < < PREV SERMON < < <}}
~
\hyperref[sec:index_chronic]{\small{[返順時目]}}
~
\hyperref[sec:0Rt7bY88Mpk]{\small{> > > NEXT SERMON > > >}}
\newline
\newline
腓立比書 1:27-2:11
\newline
\begin{longtable}{cl}
\hline
\hline
章節 & 經文 (和合本修訂版)\\
\hline
1:27 & \begin{tabularx}{0.7\textwidth}{X} 最重要的是：你們行事為人要與基督的福音相稱，這樣，無論我來見你們，或不在你們那裡，都可以聽到你們的景況，知道你們同有一個心志，站立得穩，為福音的信仰齊心努力， \end{tabularx} \\ \\ \relax
1:28 & \begin{tabularx}{0.7\textwidth}{X} 絲毫不怕敵人的威脅；以此證明他們會沉淪，你們會得救，這是出於神。 \end{tabularx} \\ \\ \relax
1:29 & \begin{tabularx}{0.7\textwidth}{X} 因為你們蒙恩，不但得以信服基督，而且要為他受苦。 \end{tabularx} \\ \\ \relax
1:30 & \begin{tabularx}{0.7\textwidth}{X} 你們的爭戰，就與你們曾在我身上見過、現在所聽到的是一樣的。 \end{tabularx} \\ \\ \relax
2:1 & \begin{tabularx}{0.7\textwidth}{X} 所以，在基督裡若有任何勸勉，若有任何愛心的安慰，若有任何聖靈的團契，若有任何慈悲憐憫， \end{tabularx} \\ \\ \relax
2:2 & \begin{tabularx}{0.7\textwidth}{X} 你們就要意志相同，愛心相同，有一致的心思，一致的想法，使我的喜樂得以滿足。 \end{tabularx} \\ \\ \relax
2:3 & \begin{tabularx}{0.7\textwidth}{X} 凡事不可自私自利，不可貪圖虛榮；只要心存謙卑，各人看別人比自己強。 \end{tabularx} \\ \\ \relax
2:4 & \begin{tabularx}{0.7\textwidth}{X} 各人不要單顧自己的事，也要顧別人的事。 \end{tabularx} \\ \\ \relax
2:5 & \begin{tabularx}{0.7\textwidth}{X} 你們當以基督耶穌的心為心： \end{tabularx} \\ \\ \relax
2:6 & \begin{tabularx}{0.7\textwidth}{X} 他本有神的形像，卻不堅持自己與神同等； \end{tabularx} \\ \\ \relax
2:7 & \begin{tabularx}{0.7\textwidth}{X} 反倒虛己，取了奴僕的形像，成為人的樣式；既有人的樣子， \end{tabularx} \\ \\ \relax
2:8 & \begin{tabularx}{0.7\textwidth}{X} 就謙卑自己，存心順服，以至於死，且死在十字架上。 \end{tabularx} \\ \\ \relax
2:9 & \begin{tabularx}{0.7\textwidth}{X} 所以神把他升為至高，又賜給他超乎萬名之上的名， \end{tabularx} \\ \\ \relax
2:10 & \begin{tabularx}{0.7\textwidth}{X} 使一切在天上的、地上的和地底下的，因耶穌的名，眾膝都要跪下， \end{tabularx} \\ \\ \relax
2:11 & \begin{tabularx}{0.7\textwidth}{X} 眾口都要宣認：耶穌基督是主，歸榮耀給父神。 \end{tabularx} \\ \\
[1ex]
\hline
\hline
\end{longtable}
$^{1}$領子妹今天的經訓是記載在.
肥立彼書第一章二十七節.
至到第二章十一節.
在新約聖經二百七十四頁.
只要你們行事為人.
與基督的福音相稱.
叫我或來見你們.
或不在你們那裡.
可以聽見你們的警方.
知道你們同有一個心志.
站立得穩.
為所信的福音齊心努力.
凡事不怕敵人的驚嚇.
這是證明他們沉淪.
你們得救.
都事出於神.
因為你們蒙恩.
不但得以信服基督.
並要為他受苦.
你們的征戰.
就與你們在我身上從前所看見.
現在所聽見的一樣.
所以在基督裡.
若有什麼勸勉.
愛心有什麼安慰.
聖靈有什麼交通.
心中有什麼慈悲憐憫.
你們就要以念相同.
愛心相同.
有一樣的心思.
有一樣的意念.
使我的喜樂可以滿足.
凡事不可結黨.
不可貪圖虛浮的榮耀.
只要全心謙卑.
各人看別人比自己強.
各人不要單顧自己的事.
也要顧別人的事.
你們當以基督耶穌的心為心.
他本有神的形象.

$^{41}$不以自己與神同等為強奪的.
反倒虛己.
取了奴僕的形象.
成為人的樣式.
既有人的樣子.
就自己卑微.
全心信服.
以至於死.
且死在十字架上.
所以神將他升為至高.
又賜給他那超乎萬名之上的名.
叫一切在天上的.
地上的.
和地底下的.
因耶穌的名.
無不屈實.
無不口稱耶穌基督為主.
聖經讀筆.
願神賜福常讀他話語.
並且進行的人.
弟兄姊妹平安.
我們今天繼續.
在腓立比書的宣講.
上一次如果大家記得.
我們講到一章的25至26節.
保羅他很渴望.
再一次回到菲律賓教會.
但是他自己對於.
能否離開當時被囚的羅馬監獄.
仍然是一個未知之數.
其實他沒有什麼把握.
能夠離開回到菲律賓.
所以你會看到.
在剛剛所宣讀的經文第27節開始.
他要求菲律賓的信徒.
無論彼此能否重逢.
他們都要行事為人.
與基督耶穌的福音相稱.
換句話來說.
保羅這裡很清楚地說.

$^{81}$菲律賓信徒的靈命成長.
是不是和保羅有關係.
沒有必然的關係.
他們的成長不在於.
保羅能否回到他們的中間.
而是單單取決於.
他們本身是否有成長的意願.
保羅向菲律賓的弟兄姊妹.
提供一個很簡單直接的建議.
努力使自己的生活.
與你的身份相稱.
既然已經接受了耶穌.
信了主 認耶穌為主.
那就要怎樣.
活出一個與基督的福音相稱的人生.
以福音主導你的人生.
堅守耶穌基督所託付給你的使命.
保羅這裡說得很清楚.
特別在第27節這裡說.
行事為人.
其實原文有一個含義.
就是做一個公民.
過一個公民的生活.
意思是其實菲律賓的信徒.
其實他們本身都擁有一個.
我們說的羅馬帝國殖民地的身份.
殖民地公民的身份.
他們未信主之前已經擁有公民的身份.
並且一直以活出羅馬公民的身份.
來感到自豪.
以效忠羅馬帝國來為榮.
不過現在身份不同了.
多了一個身份.
他們同時也已經接受了.
以耶穌基督為主.
天國公民的身份.
行事為人.
自然就像保羅所說的.
以基督的福音為依歸.
到底要效忠哪一個身份呢?.

$^{121}$兩個身份.
一個是羅馬帝國公民的身份.
另一個是耶穌基督門徒.
天國子民的身份.
很多時候會惹來一些矛盾.
甚至一些衝突.
最簡單的就是.
如果你記得.
菲律賓教會是怎樣成立的.
當保羅搭乘中歐洲大陸.
向他們傳福音.
你會回顧在《史唐傳》.
當保羅在菲律賓這個地方.
這個外邦人的地方.
行使天國公民的身份的時候.
向當地人傳福音的時候.
討不討好?.
其實招頭難押.
你會看到.
當然有人信主.
但同時也遭到當地的市民.
甚至微信者的反對.
因為保羅是搞亂的.
你傳說一些羅馬人.
不能夠接受.
甚至不應該遵守的規矩.
你是在破壞我們社會的秩序.
於是你在《史唐傳》第16章.
你會看到他們對保羅.
就是追打他.
誣告他.
甚至最後將他收入監牢.
事實上這種.
我們說外界對基督徒.
基督教不友善的態度.
是不是只發生在保羅身上?.
很明顯不是.
你會看到.
其實也必定會發生在.
保羅所建立的肥納比教會的身上.

$^{161}$因為當他們要為新信仰的時候.
也必然會像保羅那樣.
面對攻擊.
面對排擠甚至迫不得已.
所以保羅就在之後的28至30節.
這樣勉勵他們.
不要害怕.
敵黨福音的勢力.
首先原因是.
教會不是因為做錯了事.
不是因為做了壞事.
所以被人嫌棄甚至被人威嚇.
相反是因為什麼原因?.
是因為教會.
做得對.
所以才遭受反對.
做得對堅持信仰.
所以才遭受威嚇.
保羅表示教會遭受苦難.
表面上是因為敵人的反對聲音.
但實際上他講得很清楚.
是出於誰?.
是出於上帝.
你選擇了蒙恩不救的道路.
就必然會有反對的聲音.
當你堅持走一條正確的道路.
很多時候遇到的不一定是討好的事情.
相反是一些令你不舒服的際遇.
就像保羅所說.
為主受苦.
其實本來就是上帝救恩的一部分.
正如保羅所說.
我們蒙恩不但得以信服基督.
並且為主受苦.
他帶出了一個很重要.
但我們可能經常不為意的蜀靈意義.
蒙恩和受苦是並存的.
過一個蒙恩的基督徒生活.
某程度上也等於過一個為主受苦的生活.
我們蒙恩信主成為基督徒.

$^{201}$也同時蒙恩為主面對艱苦.
與耶穌基督的苦難有份.
事實上,你想清楚保羅這句話.
就像一個當頭棒喝釘醒我們一樣.
我們的信仰其實對苦難不陌生.
因為我們的信仰從來沒有離開過受苦.
就連我們信仰的最高主人耶穌基督.
他本身也是一位受苦的主.
他透過受苦成就了上帝的拯救計劃.
透過受苦承傳了教會的建立.
所以作為他的門徒,作為教會的一份子.
我們自然也會走回耶穌的路.
在你和我的信仰歷程裡.
特別當你為福音堅持去作見證的時候.
你一定會吃到不同程度的苦.
所以為什麼每年教會要去紀念受苦節.
不只是紀念耶穌受苦.
同時也要提醒我們要受苦.
每一年都提醒我們一次.
為主受苦是你和我作為基督徒的必然配套.
通過受苦,耶穌基督成為了世人的恩惠.
同樣我們作為他的門徒也會怎樣.
上帝會讓你吃一些苦頭.
透過你為主受苦來成為別人的恩惠.
我們也會走回耶穌基督所走過的道路.
無論你願不願意.
最近有一套電影可能大家也有機會看過.
叫做《九龍城寨圍城》.
這套電影不單叫好也叫座.
也因為這個緣故令昔日活躍在城寨的人事物再次成為熱話.
其中包括了一位當年積極參與宣教工作的宣教士.
有港版《德蘭修女之稱》的潘寧卓.
Jacqueline Puddinger.
她在1966年已經來到香港.
當時她只有22歲.
當時畢業於英國皇家音樂學院的她.
她買了一張船票.
這張船票很特別.
一來要最便宜.
二來可以去最多的城市.

$^{241}$讓她可以到處遊歷海外宣教.
她經過了很多不同的地方.
當她到達香港維多利亞港的時候.
她覺得這就是她要找新的地點.
於是她就留下.
結果一留就留了58年.
到現在她還在香港.
潘寧卓走到當時我們說的龍蛇混雜.
稱為三不管.
三不管是什麼意思呢.
大家知不知道.
香港政府不敢管.
英國政府不想管.
中國政府因為還沒回歸.
所以不能管.
就是三不管的意思.
她走到當時稱為三不管的九龍城寨.
她發現有很多無數的吸毒者.
從城外跑到城內.
所以當時人口很密集.
有三萬多人.
潘寧卓對這些被社會遺棄的邊緣人很有負擔.
特別是那些十多歲就已經開始吸毒的年輕人.
又或者那些去到五十多歲.
仍然要靠賣淫來滿足毒人的忌諱.
她很心痛.
城內的年輕人或者妓女.
他們擺脫不了惡性循環.
甚至有很多都誤入歧途加入了黑幫.
最初因為她沒有警察會的支持.
所以她唯有自籌自己的生活費.
於是一邊教英文教音樂.
讓她賺到的金錢來自資.
在城寨裡開設一間青年中心.
幫助當地人去解除毒癮.
相反很多人對她充滿懷疑.
為甚麼?.
因為一個外國人貿然跑來一個中國人香港人的地方.
自然很多人以為她是英國派來的間諜.
被她誤入歧途.

$^{281}$甚至她家門窗或家居曾經被毒販潑上糞便.
後來因為參加她的中心的人越來越多.
於是也惹來黑幫的注意.
除了派人去威嚇她恐嚇她之外.
也派人去跟蹤她.
但很奇妙的是.
那些跟蹤她的年輕人反而信了主.
還將多年的毒癮解除了.
類似的事件一再發生.
所以也令當時一些黑幫的頭目.
對她也開始注意起來.
還要約見她.
很奇妙的.
但那些黑幫中人竟然說.
只要幫會派來的人.
如果能夠幫她解毒.
他們都可以脫離社團.
我不追究.
是不是很奇妙?.
慢慢地.
所以上帝就使用潘寧卓.
變成一個全時間的服侍者.
來關顧城寨裡面的貧窮人.
那些癮君子和妓女.
二十多年來.
讓城寨裡很多人的生命變成不一樣.
後來街坊跟她說.
要不然.
你就是英國派來的間諜.
所以你才對我們這麼好.
有頸滾的.
要不然.
就像你所說.
耶穌是真的.
你是真心愛我們.
雖然我們知道九龍城寨.
回歸前已經拆了.
消失了.
但潘寧卓在一次接受訪問時.
也說.

$^{321}$其實香港周圍還有很多無形的城寨.
還有很多需要被關顧的邊緣人.
他曾經在一次訪問裡.
去指責一些有錢的人.
特別是基督徒.
為何只追求自己的幸福.
卻漠視了身邊的貧窮人的需要.
他有個座右銘.
他的任務就是在於.
每一分鐘每一刻.
我是否向所遇見的每一個人.
顯出上帝的恩賜呢.
我們看到潘寧卓用自己的生命.
見證了為主吃苦受苦.
使自己和別人蒙福.
蒙恩惠.
是一件真正的事情.
是一件很有意義的事情.
事實上頂梓梅.
可能一說到吃苦受苦.
我們總覺得好像很遙遠.
很離身.
特別是今天.
我們講求一個很崇尚個人主義.
成功主義.
或者很需要舒適生活享樂的世代.
有時不要說為主受苦.
會不會為主受氣.
我們都忍受不了.
接受不了.
我還記得之前有個經歷.
一早起床.
上完教會前.
那天早上起床很開心.
沒什麼人.
去到一間茶餐廳.
很多位置.
自然你會選擇什麼位置坐.
選擇卡位坐.
舒舒服服.

$^{361}$沒人.
但是.
(咳嗽).
對不起.
事件就很特別.
事件可能.
他們有個規矩就是.
卡位留給兩個人以上的人坐.
所以就說.
坐這邊吧.
先生不要坐那邊.
我就有點.
明明有很多位置.
為什麼不讓我坐呢.
沒人.
算了.
不要緊.
順服.
坐去.
另一個那些.
那些叫搭台位.
雖然有人已經在搭台.
但是走到那裡.
另一個侍應正在收東西.
他一收的時候.
我走過去.
他一收一轉身.
他那碗沙嗲牛肉麵的汁.
就啪.
濕透了.
我還沒吃沙嗲牛肉麵.
都已經聞到了.
氣不氣啊.
生氣不生氣啊.
生氣啊.
真的很生氣.
我有點想爆出來.
有一點想.
但是覺得.
不是很好.

$^{401}$忍住.
因為我這個身份.
因為我有個基督徒.
忍住.
他已經殺了.
對不起.
馬上拿杯水給我.
拿些紙巾給我.
意思是自己清洗.
我幫不了你什麼.
除了說對不起之外.
除了遞上紙巾和一杯清水之外.
其實沒有了.
他已經盡了他諸官的義.
其他的就是我的操練.
我的操練.
我可以有個決定.
就是我罵他.
我罵他也有權.
不要搞錯.
做事那麼不小心.
我決定吞下肚子.
為主的緣故.
我想有時這些都是日常的操練.
如果我們連為主受苦.
受氣.
有時也覺得不爽.
其實更加難進一步.
我們說為主吃苦.
是不是.
很難進一步去面對別人對信仰的挑戰.
所以.
三十節這裡說.
其實保羅他形容我們信仰本身就是一場什麼.
爭戰.
他用一個爭戰的字來提升我們的警醒度.
是要爭戰.
不一定是你真的和別人真刀真槍地對抗.
而是說.
不只是和別人.

$^{441}$是你內心也會有很多的爭戰.
你很想去動對方.
用口咀嚼他的時候.
你能不能忍住.
試試為主受苦.
為主受氣.
是不是.
事實上保羅他說.
無論他之前在菲臘比教會或是羅馬經歷.
他都當為他戰役的一部分.
並且他說這一場的爭戰.
也會以受苦的方式.
出現在他所建立的菲臘比教會.
他門徒的身上.
所以保羅他用自己的戰鬥經驗.
用自己爭戰的經歷.
來教導菲臘比的信徒.
你怎樣去守住你這條信仰的戰線.
首先.
保羅說得很清楚.
你對外一定要站穩.
你要做好為主受苦的心理準備.
你預料到了.
預料受苦就好很多.
當受苦來到的時候.
你沒有那麼反感.
預料到了.
並且你知道你受苦.
是出於上帝.
你見證耶穌.
而受苦.
這個是上帝的安排.
既然知道是上帝.
看著的時候.
相排的時候.
你的抗苦抗壓能力自然會高一點.
你不再去太過介意.
或者害怕別人對你的威嚇.
另一方面.
保羅也說到了.

$^{481}$不單單對外要站穩.
更加是對內你要同心.
特別面對著外界.
一個不友善的處境.
教會更加要對內團結一致.
事實上也是真的.
很多時候我們都會理解.
破壞教會合一的往往是怎樣.
是內部問題.
不是外面給我們的攻擊.
我想大家讀一讀.
之後的經文好不好.
第二章的一至四節.
節序表也有的.
講到大綱.
或者你看螢光幕.
預備起.
(法文).
好,謝謝.
你會看到其實.
勸勉,安慰,交通,相交.
慈悲,憐憫.
其實這些都是基督徒.
在耶穌的裡面.
在聖靈的裡面.
一些普遍共同的經歷.
不單單是菲律賓教會的弟兄姊妹.
我們也會經歷過.
所以保羅這裡說.
若有.
其實他不是假設.
他不是說.
不是很確定.
到底菲律賓的弟兄姊妹.
他們有沒有這些共同的信仰經驗.
他也不是假設.
他們現在沒有了.
他們沒有這些經驗.
期望他們將來會擁有.
保羅不是這個意思.

$^{521}$保羅這裡其實是強調.
他們已經擁有.
這些共同美好的信仰經歷.
有的了.
因為他們本身在保羅的口中.
他們已經是一個在基督的裡面.
在聖靈的裡面.
互相的勸勉,安慰,交通.
憐憫的一個群體.
既然如此.
保羅就勸勉他們.
在這些美好的基礎上.
要進步.
不要只停留在這裡.
你們要更加繼續去成長.
培育出保羅在之後所說.
在意念上.
在愛心上.
在心思上.
你們都要邁向相同合一.
這樣的話.
保羅即使無論能否來到菲律賓教會.
他說他的喜樂.
都可以更加滿足.
更加安慰.
留意保羅這裡所說.
相同.
他當然不是說.
所有人的愛心,思想都要或一.
不是或一.
保羅不是不容許.
在教會裡面有不同的意見和看法.
好像那些工廠那樣.
啤膠出來.
啤那些產品出來怎樣.
每件都一樣.
一模一樣.
或一不是這種意思.
保羅是說信徒的心思意念.
應該朝向一個相同的方向和目標.

$^{561}$來到一同以耶穌基督的吩咐.
以聖靈的帶領作為一個會合點.
是說邁向同一個方向.
這種合一.
而實際到底要怎樣做.
保羅提出了一些很具體的建議.
首先他提到凡事不可結黨.
不可貪圖虛浮的榮耀.
結黨其實我想我們在社會上.
不難去體會或理解.
結黨最簡單的意思就是.
只會和那些對自己有利益.
給自己人好處的人結交.
組隊.
所謂人多勢眾.
所以你會看到在社會上.
有不少人會用結黨的方式.
我們說集氣.
來增加自己的影響力.
或增加自己有更多的話語權.
這就是結黨.
在這個原則下.
屬於黨內的自己人自然會怎樣.
大家同聲同氣.
大家互相去保衛.
但在黨以外的.
還會怎樣?還會是自己人嗎?.
就不是了.
那些不是自己人的就自然不用理會.
冷落他.
甚至排擠打壓他.
這樣還能不能夠合一?.
不能夠了.
這樣是破壞合一.
另一方面保羅也提到.
不要貪圖虛浮的榮耀.
首先其實.
虛浮的意思是空心.
即是那些空心的榮耀.
保羅說如果人不斷去爭取.

$^{601}$這些空心的榮耀.
其實是在破壞合一.
例如是甚麼呢?.
可能過於炫耀自己給人一些優越的條件.
又或者不斷追求別人對你的吹捧.
只能接受別人讚你.
漸漸你會怎樣?.
你只會和那些喜歡你的人.
和那些like你的人組隊.
又或者明知自己是有錯的.
但因為面子.
因為不能失去的緣故.
所以怎樣?.
就死撐.
堅持自己的尊嚴.
堅持己見.
甚至怎樣?.
甚至和別人對著幹.
結果老羞成怒.
保羅提醒我們.
要杜絕這些結黨貪圖虛浮榮耀的行為.
因為當一個群體.
越來越多人以自我為中心.
以自己的利益行先.
將自己當成為主人的時候.
這種心態行為就必然會怎樣?.
破壞教會群體的合一.
所以保羅說一定要杜絕.
一定要留心.
而另一方面.
保羅也提醒.
一個信仰群體要一起成長.
不能只顧著自己的需要.
信仰從來不只是我們和上帝的關係.
信仰某程度上已經包含了.
我們和弟兄姊妹的關係.
換句話來說.
保羅提醒我們要有一個大局的意識.
不能只看自己.
我們要顧及彼此.

$^{641}$所以他就要求信徒要存謙卑的心.
要看別人比自己強.
不單顧自己的事.
也要顧別人的事.
首先謙卑其實是說什麼呢?.
是一種正確的自我評估.
其實我們明白.
無論我們有多少才幹.
有多少能力.
是出於誰?.
是上帝給予的.
是上帝的恩賜恩典.
不是靠我們賺回來的.
我們的本相在上帝面前是渺小的.
是不配和不完全.
如果我們將這種謙卑的做人態度.
應用在一個群體的人際關係裡.
就是無論我們有多強大.
多有才幹.
總會有些情況.
我們是不夠的.
總會有些情況我們是軟弱.
我們是不濟的.
又或者說你多厲害.
人生總是有限制.
總會有我們自己覺得無能的時候.
並且我們會看到.
即使對方確實有些地方不及我們自己.
但是怎樣?.
有另一些範疇.
是我們不及對方.
是別人給我優勝.
總是有些亮點上帝給他.
沒有賜給我們.
從而我們可以由衷地去欣賞對方這些美善和重要.
所以保羅所說看別人比自己強.
其實有另一個盡心的意義.
就是我們能否看到別人在群體裡的價值和重要性.
看別人比自己重要.
事實上我們會很明白.

$^{681}$上帝創造每一個人.
其實都有他的獨特性和重要性.
你看看你旁邊的那個.
不單單樣子和你不同.
和你的氣質,才幹所有東西都不同.
上帝就透過他.
透過你的不同和差異來怎樣?.
互相祝福和合作.
否則上帝做我們全部一樣就行了.
為何要做我們不同?.
就是要我們學習和你不同的人相處.
來學懂如何互相接納和欣賞.
並且彼此的配搭和服務.
假如你試想像.
我們總是沒辦法去欣賞對方.
總是常常看別人不順眼.
覺得自己總是表現好一點.
嫌對方總是做得不好.
甚至覺得他一事無成.
這樣其實還有沒有得配搭?.
很難的.
你本身已經厭惡了他.
你怎樣和他再有美好配搭呢?.
更難說你和他同心合意事奉上帝.
所以弟兄姊妹.
我自己很相信.
上帝安排一些和我們不同的人.
在我們身邊.
甚至你覺得差勁的人在你身邊.
不是想我們去找錯處.
找錯處很容易的.
你發不發覺?.
太容易了.
我記得有個例子很深刻.
有個老師在黑板上寫著.
一加一等於幾?.
二加二呢?.
三加三等於六.
四加四呢?.
他寫了一個九字.

$^{721}$於是所有家長都說.
老師錯了第四條.
是錯的.
四加四等於八.
為什麼寫了九呢?.
老師說其實我故意的.
為什麼大家只看到我錯的東西.
看不到我答對了三題.
找錯處是我們天生的強項.
有沒有試過和子女對功課?.
我們最差的是.
你做錯了.
欠了一個零.
欠了一個點.
我們最厲害就是抓錯處.
但我們有沒有在抓錯處之前.
原來你做對了這麼多.
沒有的.
我們不理會他做對了的東西.
我們只理會他做錯了的東西.
所以上帝將那些你覺得差勁的人.
放在你身邊.
其實不是要你抓錯處.
太容易了.
這麼容易上帝不用你做.
上帝是想你學會.
如寶樂所說的.
欠卑.
看別人比自己重要.
學會這個功課.
所以不妨操練一下.
如寶樂所說的.
不要只顧自己的需要.
不要只看到自己眼睛的良目.
不要只看到別人的錯.
你也要顧別人的事情.
與其我們嫌棄對方做得不好.
不正經.
不斷埋怨投訴.
不如學寶樂所說的.

$^{761}$換一下心態.
對方不是做不到.
是怎樣?.
未做到.
已經很不同了.
你看死他做不到.
還是你覺得仍然有盼望.
他未做到.
當你覺得他未做到.
你會怎樣?.
我給他多點時間.
我給他多點耐性.
幫他成長.
同時也求上帝給我們智慧耐性.
我們怎樣鼓勵引導他達標.
又或者.
可能對方的強項優勢.
真的不在那些地方.
那又怎樣?.
你發掘其他他的優點.
你發掘他的亮點.
一定有的.
因為他是上帝所做的.
跟你一樣.
你有你的生命亮點獨特之處.
他也有.
不過你未發掘到.
你試試發掘出來.
並且你提升他.
你鼓勵他這些亮點.
讓他的亮麗.
也能夠發揮出來.
這樣的話.
我們就可以彼此祝福.
真的邁向保羅所說的.
那種同心合一.
而最後.
我很想再看餘下的經文.
第五至十一節.
保羅從上述的教導.

$^{801}$轉向了耶穌基督的身上.
我們看到.
他勉勵菲勒比的信徒.
要在謙卑合一的事情上.
以基督耶穌為榜樣.
效法耶穌基督的心志.
大家可否再讀一次.
雖然我們已經很熟悉.
我請大家再讀一次.
第二章的五至十一節.
預備起.
「耶穌基督的心」.
「靠有心的形象」.
「體現自己與神同等的愛情」.
「但我會記取祂的形象」.
「成為人類的姿態」.
「以祂的名字」.
「就自己的氣味」.
「全心信我」.
「以至至死」.
「且死在神之家上」.
「所以神將祂的性格」.
「傳授給有持不怕」.
「受挑撥萬民之上的民」.
「願你一切在天上的地上的」.
「和我們在地下的」.
「因耶穌的名」.
「無法不實」.
「無法考證耶穌基督」.
「為主」.
「使我們要歸於神上」.
好,謝謝.
我們不妨去看看.
耶穌基督.
如何一路一路降卑自己.
英文說.
耶穌基督本身雖然是上帝.
他擁有與上帝相同的尊榮.
甚至乎權柄.
但卻沒有「抓緊不放」.

$^{841}$耶穌沒有因為擔心自己失去身份地位榮譽權柄.
他抓住了.
沒有抓住他應有應得的利益.
相反他選擇一路一路降低自己.
如何降低?.
試想像一下.
耶穌由最尊榮的一格.
一降就降到人.
然後還要人之間.
不是上等人.
富貴人是怎樣?.
奴僕.
奴僕還不夠.
還要被人誣告.
最後變成罪人.
還不夠.
要用當時最羞辱的方法.
釘十字架而離世.
試想像中間的距離.
有多少重?.
我們降格了多少級別?.
但你會看到.
有沒有人逼他?.
沒有人逼耶穌這樣做.
亦不是因為他們三位一體.
上帝,聖靈,聖父和聖子.
三個猜包剪鎚.
耶穌不好意思.
你猜輸了.
所以你下去吧.
你下去拯救世人.
你下去釘十字架.
不是這樣.
而是為了拯救.
成全我們的生命.
耶穌是自願降低虛己.
一路向下.
向每一個人和罪人提齊.
耶穌基督為了我們的好處.
不是為了祂的好處.

$^{881}$祂甘願捨己.
捨棄自己的權利,自由和生命.
我們看到耶穌基督.
祂親身的示範.
什麼是謙卑,捨己的極致.
祂用身體力行.
來見證上帝行事為人的方式.
就是沒有死亡.
就沒有生命.
沒有死透就不能重生.
沒有捨棄就沒有得著.
所以中國人說得好.
捨得.
要捨棄才能得著.
沒有降卑就沒有提升.
耶穌基督如果不願意降卑.
我們就不能被提升.
你見到耶穌基督.
祂甘願成為人世間最卑微的那一位.
為的是要配合上帝的旨意.
上帝最終也透過耶穌.
成就了祂的拯救計劃.
並且祂和耶穌平反.
祂提升耶穌成為最卑微的那一格.
升上宇宙最高最尊榮的那一位.
並且也成就了我們和上帝和好.
和人和好的好消息.
保羅也以耶穌基督.
這種捨棄謙卑的心腸.
來叫我們成為效法的榜樣.
勸勉我們要活出一個.
以捨棄服侍的生命為依歸.
要全心謙卑.
努力地和弟兄姊妹.
和你身邊的人追求同心合一.
而我也很想.
趁著今天我們說是父親節.
透過這段經文.
來勉勵每一位家庭成員.
我們如何去效法基督的心腸.

$^{921}$以耶穌基督的心為心.
如何去學習放下那些不必要的執著.
來一心追求和睦.
追求合一.
我們可以想像一下.
如果耶穌基督不願意放下自己的權益.
不願意放下自己的面子.
能否成就耶穌復活的工作.
不可能.
同樣地.
很多時候家人之間.
很多時候會有很多的摩擦.
不能夠復活.
其實不也是因為我們有很多東西要扛著不放嗎.
我們有很多東西扛得太緊.
看自己太過重要.
有時我們可能是太過看重自己的得失需要.
太過執著於面子.
各人都看自己為家庭的主人.
以致我們總是想怎樣.
一有什麼不舒服.
或者對方和你不同的時候.
我們的焦點就會去證明是我對你錯.
又或者即使自己有錯也好.
我們都要拉對方下水.
有沒有聽過這些爭拗.
我是有做得不好.
你很好嗎.
你比我還要壞.
然後攬炒.
我們很多時候執意在.
要自己贏自己說了算.
忘記了耶穌才是我們一家之主.
覺不覺得.
雖然我們頭腦上知道.
但我們走出來.
我們總是自己做主.
所以最後我們可能很多時候.
在家人的爭拗裡我們爭拗贏了.
但是.

$^{961}$嚴規就是這樣產生.
關係就是這樣慢慢破壞或者疏遠.
結果贏的人開心嗎.
不開心.
他贏了但他發現所有人都輸了.
這樣贏有什麼意思.
根本都沒有贏過.
我記得有一本書叫做.
活出群體的美好.
好像是校園出的.
我看過一個例子.
好像跟大家分享過.
不過我想再說一說.
很深刻.
那本書說到一個例子.
如果我們對著一班加拿大的印第安小朋友.
說一說.
譬如他問一條問題.
最先舉手答到的就可以得到一份禮物.
你猜他們會怎樣反應.
譬如我問大家.
今天是什麼節.
舉手有獎品.
大家就會怎樣.
馬上就像聖誕節玩問答遊戲.
馬上舉手拿獎品.
鬥快嘛.
香港人當然是這樣.
獎品給一份.
當然是鬥快拿.
他說他們不會鬥快舉手.
他們會聚在一起.
一起去想答案.
談好之後.
然後一起答出答案.
因為他們不能接受.
只有我一個贏.
其他人都輸.
他們沒辦法接受這種情況.
事實上很有意思.

$^{1001}$如果家裡都是這樣.
多麼美好.
家庭一定會遇到困難.
遇到問題.
重點是我們能否在這些困難裡.
轉化成祝福.
我們可否慢慢有商有量.
可否體諒協調.
彼此之間有差異.
步伐不同.
以至最後我們都可以.
Win Win 雙贏.
比起只是我贏.
你輸.
其實一家人有時齊上齊落.
可能更加有意思.
所以弟兄姊妹.
其實無論我們是做人配偶.
做人父母或子女.
我們都需要留心.
有時我們會否過於.
抓住一些你覺得很重要的權益.
死都不放.
導致影響破壞了.
你與家人的關係.
破壞了你對上帝的寫記.
那種信服.
事實上.
家裡是我們回應上帝召命的地方.
上帝呼召你在家裡.
來服侍你的家人.
服侍上帝.
見證上帝.
換句話來說.
裡面一定會通過寫記.
受苦.
這些過程.
才會成就上帝的心意.
要服侍即是甚麼意思.
你要放下某些執著.

$^{1041}$權益自由.
以達至別人蒙福.
上帝得到榮耀.
所以假如我們與家人有吵架的時候.
有埋怨對方的時候.
甚至很想數落對方的時候.
我們可以停一停.
想一想今天的經文.
看一看那位在十字架上的耶穌.
我們回想起來.
我們不是經常得罪耶穌嗎.
有沒有誰沒有得罪耶穌.
無端端坐在這裡.
每個人都得罪過.
我們也試過對耶穌不公道.
我們現在答應了耶穌的事.
有沒有試過走數.
有了不用說了.
大家一定都很悶.
我們經常都會這樣.
但耶穌有沒有每次都在計算.
耶穌總是饒恕我們.
總是接答我們.
所以當我們看到耶穌的時候.
我們又何苦要和你所愛的家人.
太過斤斤計較.
抓住他不放.
另一方面你看到耶穌.
不單單是寬恕我們.
他看好我們.
他總是會用我們.
一次一次給我們機會.
不單單幫我們改過.
給我們機會發揮成長.
所以其實作為一家人.
我們能不能夠像保羅所說.
看別人比自己強比自己重要.
我們不要再經常都.
憂慮著你家人那些未改的甩鈕.
不要只找錯處.

$^{1081}$你找找他的亮點.
你稱讚他鼓勵他.
有時可能這樣.
反而增加大家對家裡的歸屬感.
和投入感.
試想想你回家時.
常常被人盯著.
你會不會想回去呢.
不想了.
但你每次回去.
總有人稱讚你.
不是虛假那種.
真心找到你的亮點.
那你就終於回家了.
所以弟兄姊妹盼望我們.
藉著今天這段經文.
我們能夠在家庭裡活出.
和我們身份相稱的生活.
我們以耶穌基督的心腸.
來善待我們的家人.
我們去努力發揮.
來發掘.
欣賞你家人那些亮麗之處.
一起去追求和睦合一.
讓上帝在我們家裡得到榮耀.
也讓我們家人可以一起成長.
一起成為蒙福的上帝兒女.
我們同心圖告.
願你的恩典來到我們當中.
來到教會當中.
來到我們家庭當中.
來到弟兄姊妹每一個人的身上.
讓我們再次明白.
我們蒙受恩惠.
不單是因為我們認信耶穌.
我們更加要蒙恩與.
耶穌基督的受苦有份.
因為我們蒙恩.
是要進到不同的崗位裡.
回應上帝給我們的呼召.

$^{1121}$無論在教會裡.
在家裡.
我們都要學習成為.
這幾路復的心腸.
來服侍耶穌基督.
託付給我們每一個人.
在過程裡.
我們需要像耶穌基督那樣.
謙卑,降卑.
要犧牲我們一些可能的意願.
我們一些面子.
甚至我們一些需要.
為的是要讓上帝得到榮耀.
為的是要讓我們身邊的人.
能夠蒙福.
能夠從我們身上體會到.
耶穌基督的愛.
過程裡我們很需要.
耶穌幫助我們.
因為我們靠自己.
我們真的做不到.
求主讓我們能夠時刻仰望你.
特別當我們同人有厭愧的時候.
我們求主讓我們.
看到你在十字架上的犧牲.
看到你在十字架上.
怎樣沒有和我們計較的時候.
我們也放下一些無謂的執著.
我們願意和我們身邊的人.
竭力和好,追求合一.
為的是要見證耶穌你.
為的是要通過我們一些.
可能吃苦的事情.
來互相祝福.
一起進到耶穌裡面.
求主你藉著這段經文.
賜福給我們每一位.
賜福給教會.
賜福給每一個家庭.
奉基督耶穌你的名來祈求.

$^{1161}$阿們.
\newpage



\section{約翰福音 4:5-15}
\label{sec:0Rt7bY88Mpk}
\textbf{2024年6月9日崇拜直播｜梁家麟牧師｜為有源頭活水來｜約翰福音4\_5-15}
\newline
\newline
連結: \href{https://youtube.com/watch?v=0Rt7bY88Mpk}{\texttt{https://youtube.com/watch?v=0Rt7bY88Mpk}} ~~~~ 語音日期: 2024-06-16
\newline
\newline
\hyperref[sec:9ogZwUemHN4]{\small{< < < PREV SERMON < < <}}
~
\hyperref[sec:index_chronic]{\small{[返順時目]}}
~
\hyperref[sec:_sNrPBKmvwc]{\small{> > > NEXT SERMON > > >}}
\newline
\newline
約翰福音 4:5-15
\newline
\begin{longtable}{cl}
\hline
\hline
章節 & 經文 (和合本修訂版)\\
\hline
4:5 & \begin{tabularx}{0.7\textwidth}{X} 於是到了撒瑪利亞的一座城，名叫敘加，靠近雅各給他兒子約瑟的那塊地。 \end{tabularx} \\ \\ \relax
4:6 & \begin{tabularx}{0.7\textwidth}{X} 雅各井就在那裡；耶穌因旅途疲乏，坐在井旁。那時約是正午。 \end{tabularx} \\ \\ \relax
4:7 & \begin{tabularx}{0.7\textwidth}{X} 有一個撒瑪利亞婦人來打水。耶穌對她說：「請給我水喝。」 \end{tabularx} \\ \\ \relax
4:8 & \begin{tabularx}{0.7\textwidth}{X} 因為那時門徒進城買食物去了。 \end{tabularx} \\ \\ \relax
4:9 & \begin{tabularx}{0.7\textwidth}{X} 撒瑪利亞婦人對他說：「你是猶太人，怎麼向我一個撒瑪利亞女人要水喝呢？」因為猶太人和撒瑪利亞人沒有來往。 \end{tabularx} \\ \\ \relax
4:10 & \begin{tabularx}{0.7\textwidth}{X} 耶穌回答她說：「你若知道神的恩賜，和對你說『請給我水喝』的是誰，你早就會求他，他也早就會給了你活水。」 \end{tabularx} \\ \\ \relax
4:11 & \begin{tabularx}{0.7\textwidth}{X} 婦人對耶穌說：「先生，你沒有打水的器具，井又深，哪裡去取活水呢？ \end{tabularx} \\ \\ \relax
4:12 & \begin{tabularx}{0.7\textwidth}{X} 我們的祖宗雅各把這井留給我們，他自己和兒女以及牲畜都喝這井裡的水，難道你比他還大嗎？」 \end{tabularx} \\ \\ \relax
4:13 & \begin{tabularx}{0.7\textwidth}{X} 耶穌回答，對她說：「凡喝這水的，還要再渴； \end{tabularx} \\ \\ \relax
4:14 & \begin{tabularx}{0.7\textwidth}{X} 誰喝我所賜的水，就永遠不渴。我所賜的水要在他裡面成為泉源，直湧到永生。」 \end{tabularx} \\ \\ \relax
4:15 & \begin{tabularx}{0.7\textwidth}{X} 婦人對他說：「先生，請把這水賜給我，使我不渴，也不用到這裡來打水。」 \end{tabularx} \\ \\
[1ex]
\hline
\hline
\end{longtable}
$^{1}$今日經訓於《新約聖經》第129頁,約翰福音第4章5-15.
請聽我讀出.
於是到了撒瑪利亞的一座城,名叫聚迦,靠近雅各給他兒子約瑟的那塊地,在那裡有雅各井,耶穌因走路困倦,就坐在井旁..
那時約有誤證,有一個撒瑪利亞的婦人來打水,耶穌對她說:「請你給我水喝.」那時門徒進城買食物去了..
撒瑪利亞的婦人對他說:「你既是猶太人,怎麼向我一個撒瑪利亞人要水喝呢?」原來猶太人和撒瑪利亞人沒有來往..
耶穌回答說:「你若知道神的恩賜,我對你說給我水喝的是誰,你必早求他,他也必早給了你活水.」.
婦人說:「先生,我有打水的器皿,井有心,你往哪裡得活水呢?」.
「我們的祖宗雅各將這井留給我們,他自己和兒子,並生畜也都喝這井裡的水,難道你比他還大嗎?」.
耶穌回答說:「凡喝這水的,還要再喝.人若喝我所賜的水,就永遠不喝.我所賜的水,要在他裡頭成為泉源,直湧到永生.」.
婦人說:「先生,請把這水賜給我,叫我不喝,也不用來這麼遠打水.」.
各位觀眾,早安!.
看聖經是很長的,我們讀經只是讀了其中一部分,因為如果讀完整張聖經就太過冗長了..
耶穌在雅各井旁和婦人談道的故事,我們都很熟悉的,這是一個內容很豐富的故事,可以從不同的角度去詮釋去理解..
例如,有人就耶穌一句話,建立起他們的崇拜神學,這樣算不算是小題大做呢?可以相確的..
不過,今天我們只是專注去看,這是一個談道的故事,用我們慣常的說法,耶穌是向一個陌生的婦人傳福音..
雖然聖經明說耶穌這三年半的公開傳道生涯,主要做的就是醫病趕鬼,傳天國的福音,.
不過在福音書裡面,耶穌主動和個人傳福音的技術其實不多的..
就像雅各井旁的案例,技術如此的詳細,將兩個人的多個對白都詳盡記錄的,就絕無僅有..
我看不到有另外的例子..
我們因此可以理解,這應該是耶穌傳道的一個示範..
如果我們說我們要效法耶穌,學耶穌傳福音,約翰福音第四章應該是一個常見的範本..
有個題外話,當耶穌和那個婦人談道的時候,門徒去了找食物,所以現場只有耶穌和那個婦人,.
他們說了些什麼門徒理論上是不知道的..
由於門徒不認識那個撒瑪利亞婦人,所以由那個婦人事後將她整大段的對白告訴門徒的可能性其實不是很高..
所以門徒之所以知情,約翰之所以能夠將整個故事記述下來,極有可能是耶穌傳告給他們的..
耶穌在故事發生之後,將他和婦人談道的詳細過程告訴門徒..
而他之所以告訴門徒這個故事的細節,極有可能那個目的是教育性的..
希望門徒藉著這樣來學習傳福音,.
這樣也強化了我們將整個故事理解為耶穌傳福音的示範的理據..
必須說,耶穌和撒瑪利亞婦人傳福音是一個非常精彩的傳福音案例,.
說它是一個完美的案例也不過分,中間有很多值得我們學習的地方..
首先我們注意到,耶穌和撒瑪利亞婦人彼此是完全沒有交集的..
兩人不認識對方,連生活圈子,人際關係網絡都沒有重疊之處..
第九節說,原來猶太人和撒瑪利亞人沒有來往..
這句話的原意可能是說,原來猶太人是不會用撒瑪利亞人用過的餐具..
這是呼應著前面耶穌要撒瑪利亞婦人用她的裝水器皿給她喝水的說法..
所以當耶穌開口向婦人求助,要求她拿水喝,.
是令到那個婦人感覺到詫異的..
那個婦人說,既然你是猶太人,你怎麼會向我這個撒瑪利亞婦人要水喝呢?.
所以不要以為猶太人和撒瑪利亞人同源同宗,即使不是同文化,都是近文化,傳福音很容易..

$^{41}$他們彼此之間互相敵視,老死不相往來的立場,比我們向那些異文化的人傳福音還要難..
就好像今天以色列人和巴勒斯坦人傳福音,容易嗎?.
俄羅斯人和烏克蘭人傳福音,容易嗎?.
又或者熱刺的球迷寧願輸球給曼城,都不願意給同事叔弟阿仙奴贏球,是嗎?.
這個原則大家都明白的,我寧願自己攬炒,我都不會讓你贏,這是很簡單的事..
在耶穌時代,撒瑪利亞已經不再是一個獨立的政治實體,是被羅馬帝國歸聘到猶太的領域..
不過猶太人和在這個地方居住的撒瑪利亞人仍然彼此視為世仇,老死不相往來..
為什麼耶穌會和撒瑪利亞婦人遇上呢?.
事緣耶穌和門徒要從南面的猶太地前往北部的加利利,中間就會經過撒瑪利亞這個地區..
當然,猶太人還有另外的路可以選擇,他們可以避開撒瑪利亞,沿著約旦河谷走就可以了,不過這樣走的路會長一些..
無論如何,耶穌只是路過和撒瑪利亞婦人偶遇..
中午時分,他們來到雅各井旁,耶穌覺得累,就坐下來休息,門徒就入城購買食物做午餐..
然後這個婦人就出現,她在中午烈日當空時分獨自來到雅各井取水..
這個不是當時期的婦人,通常都是女人去取水的,即是慣常取水的時間..
在那個時期,婦人通常在早上或傍晚約結伴來取水,順便聊天..
這個是打牙交,講八卦的時間..
就好像菲律賓傭工早上送小朋友上學或購物時,都是他們相約在商場聊天..
這個婦人可能是因為名聲不佳,被人排擠,所以才選擇用這麼古怪的時間來取水..
所以整個場景只有耶穌和婦人兩個..
聖經記載,是耶穌主動開口跟那個婦人說話,請她幫忙取水..
耶穌的口渴可能是真實的情況,但也不排除耶穌是找藉口跟那個婦人搭訕..
在前面,約翰福音第三章記載了耶穌和尼哥底姆的詳細對話..
兩章聖經,三章四章記載了兩段對話,不過說話的對象完全不同..
尼哥底姆是正統猶太人,受過高深的教育,受人尊重的律法師,生活聖潔的宗教領袖..
但婦人卻是一個被視為是雜種混雜的撒瑪利亞人,無知,處於社會的底層,生活敗壞,遭人鄙視..
不過耶穌對這兩個極度不同的人同樣感興趣,也認定他們都是傳福音的對象..
是的,不論種族,性別,教育程度,社會地位,也不管對方的宗教或道德的景況是好是壞,.
人人都是耶穌傳道的對象,人人都需要耶穌..
在福音面前,不存在任何篩選的機制,沒有人是罪惡到敗壞到一個地步,不值得將福音傳給他..
反倒耶穌宣稱,有病的人才需要醫生,他們是罪人悔改..
到最後那些覺得自己很輕香,覺得自己很正,覺得自己很棒的人,他們還是耶穌所選擇的對象之外..
無論如何,福音要破除我們對任何人,任何先入為主的成見,偏見,我們自己無法避免會對人有看法的,.
只不過我們學習以耶穌的眼光來對待身邊每一個人..
耶穌請婦人幫忙取水,以此作為藉口,開始了兩個人的對話..
在婦人質疑為何會問一個撒瑪利亞婦人取水之後,耶穌巧妙地利用「問人拿水飲」這個話題,.
然後將主客調轉,主動邀請婦人「問耶穌拿水飲」..
第十節這樣說:「你若知道上帝的恩賜和對你說話說給我水喝的是誰,你必早求他,他也必早給了你活水.」.
耶穌含蓄地指出他的特殊身份和他能夠給人的特殊禮物..
你若知道我是一個沒有化妝的姜濤,你一定會問我拿簽名,還會問我拿下一場演唱會的季婚票,你知道我是誰,你就一定會問我拿東西..
至於這個特殊的禮物,上帝的恩賜是什麼呢?.

$^{81}$耶穌就順應「水」這個話題,指出他能夠送出的是活水..
什麼叫活水呢?婦人直覺理解的是活動的水,流動的水,也就是河流..
於是她質疑耶穌如何找活水給她..
一方面,她連打水的器具都沒有,.
二方面,如果耶穌自己有水,他怎麼會問婦人拿水喝呢?.
你自己有水嗎?.
第三,這裡方圓數十里都沒有河流,如果有河流,他們的祖先就不需要大費周章打出這個井來,沒有河流又怎麼會有活水呢?.
耶穌糾正婦人的觀念,活水不是指著流動的河水,而是指著賜人生命的水..
所以活是賜生命的意思,Living是賜生命,不是指著活動..
能夠一次解決人乾渴的問題..
13,14節耶穌說,「凡喝者水,還要再喝.人若喝我所賜的水,就永遠不喝.我所賜的水,要在他裡頭成為泉源,直湧到永生.」.
活水這個觀念在舊約已經存在,通常是指著上帝的供應和幫助..
今天我們不詳細糾纏那些自圓自圓的解釋了..
簡單我們看到的是,耶穌從那個婦人的感覺需要,引導她明白她的真實需要,她的real need..
婦人一時之間思想轉不過來,仍然理解水和口渴是物質上,生理上的東西,於是很開心地立即回答,.
「好啊,如果有這樣喝了一次就不需要再喝的水,我就不需要再花氣力來這個井打水了,對不對?」.
耶穌沒有直接糾正婦人理解的錯誤,祂是即時轉移了話題,「你去叫你丈夫也道謝你來,你叫你老公來這裡.」.
婦人對耶穌的吩咐感覺到很噁心,即時回答她,「我沒有丈夫.」.
婦人這樣說,耶穌就立即回答,「你說你沒有丈夫是沒有錯的,你已經有過五個丈夫,你現在有的並不是你的丈夫,你這話是真理,十七,十八節.」.
毫無預警的,耶穌一下就揭露了婦人的私生活極度混亂的情況,她曾經和很多人結婚或同居,如今和她在一起的都不是一個合法的丈夫..
值得注意的是,哪怕是講到婦人生活的問題,耶穌在這裡的措辭還是非常謹慎的,祂避免做任何直接的道德判斷,沒有斥責婦人道德敗壞,也沒有用上奸淫,通奸,同居等字句,祂只是說如今和婦人在一起的不是祂的丈夫..
事實上,從始至終,你看所有的對白台詞,耶穌都沒有對那個婦人所說的任何否定的判斷,當婦人說她沒有丈夫的時候,耶穌的回應是,「你說得對,你沒有說錯.」.
耶穌也沒有對那個婦人的行為和生活做直接的否定.耶穌說這番話的目的不在於批評,而是有兩個目的..
第一,耶穌亮了一手,展示了一些理智,讓婦人知道,耶穌其實是明白她的底蘊的,「我知道你的,我認識你.」.
第二,讓婦人不要再在那些空洞的話語上兜來兜去,直接面對個人生命的真相,直接正視自己的問題和需要..
我們很多時候和別人談道,你都會知道,對方不停地將來一些客觀的問題,理性的問題,跟你說來說去,都言不及義..
如果信仰不是坐落在人生命深處的課題,就永遠都掃不著樣處..
不過,無論如何,從耶穌那裡我們學到的是,在和別人談道的時候,盡量避免對別人作出直接的否定和批評,.
無論是對方的觀點,想法,還是對方曾經犯過的可能的錯誤或者罪惡..
你否定對方的觀點,總是會激起個人的情緒,造成不必要的口舌之爭,即使你贏了辯論,也輸了人心,是不是?.
現在不是辯論比賽,不要駁斥,不要否定..
你直指對方犯罪,更加會惹來反感,「我就是犯罪,關你什麼事?憑什麼你來指指點點?」.
對不對?維持彼此間的同理心,讓人明白,你努力明白他,理解他的想法和感受,他才願意繼續和你談下去,是不是?.
婦人被耶穌揭穿個人的隱私,非常震驚..
她第一個反應是這樣說的,19節:「先生,我看出你是個先知.」.
這句話表明什麼?婦人承認耶穌所說的話是符合她的事實,不然她應該這樣說:「先生,我看出你是個假先知.」.
她沒有說她是假先知,就證明了她說的是對的,她沒有否認,是不是?.
她們從未謀面,而耶穌竟然知道她的故事,這個知識只能夠理解是來自上帝的啟示,是上帝通水讓她知道,不然怎麼知道呢?.
所以耶穌是一個先知.從這句話,你看到這個婦人對耶穌的態度尊敬了很多,.
對於一個知道自己生命真相的人,你最好客氣一點,因為他拿著你的底牌..

$^{121}$不過,婦人不願意在這個陌生人面前多說自己私生活的問題,於是她馬上就轉移了話題,從個人問題轉到神學的問題,.
就是瑟瑪利亞人和猶太人有一個重大的差異.20節,我們的祖宗在這個山上禮拜,你們倒說應該禮拜的地方是耶路撒冷..
大概在主前400年,瑟瑪利亞人在基里森山建立了他們自己的聖殿,雖然這一座聖殿在主前200年被毀,.
不過瑟瑪利亞人仍然習慣在基里森山上敬拜上帝,他們拒絕承認耶路撒冷這個宗教中心的地位..
耶穌知道婦人不想說個人私德的問題,不過耶穌也沒有咄咄逼人,繼續執著於婦人有多少個老公這個話題不放,.
畢竟她無意做審判官,對人定罪也不是她人間的目的.再說,婦人將話題轉到宗教方面,也符合耶穌的期望..
於是耶穌接過婦人的提問,鄭重地說了這段比較長的說話,21到24節,「婦人,你當信我.」.
什麼叫你當信我呢?就是跟耶穌說認真話,他就會這樣說,我實實在在告訴你們,一樣的意思..
「事後將到,你們拜父也不在這山上,也不在耶路撒冷,你們所拜的你們不知道,我們所拜的我們知道,因為救恩是從猶太人出來的..
事後將到,如今就事了.真正拜父的要用心靈和誠實拜他,因為父要這樣的人拜他,上帝是個靈,所以拜他的必須用心靈和誠實拜他.」.
這段說話傳遞了三個信息,第一,儘管從過去到現在為止,猶太人和撒瑪利亞人都在爭論應該在哪裡敬拜上帝,各執一詞..
不過耶穌說,宣告在不久的將來,這個問題不再是一個有意義的問題..
耶穌是兩次說道「事後將到」,21節,23節..
什麼叫「事後將到」?什麼時候?就是來塞亞降臨的時候,也就是耶穌的時候..
從耶穌施行代屬以後,敬拜的地點就不再是一個問題,因為連聖殿都不會再有價值了..
唯一蒙上帝閱立的敬拜,是在耶穌基督裡面,地點完全不重要,聖殿也不重要..
耶穌強調,我們拜的不是一般的神明,而是我們的父..
我們之所以稱上帝為父,我們之所以成為天父的兒女,是藉著耶穌基督而得的身份和關係,是不是?.
所以,在基督裡,我們就拜父,不再是拜一個神明上帝,是我們的父..
所以,有耶穌就足夠,地點完全不相干..
第二點,在這一段信息裡面傳遞的是,猶太人和撒瑪利亞人對敬拜地點的爭論,.
沒有誰對誰錯,誰優誰劣的分別..
不過,耶穌卻是指出,救贖之路是出自猶太人,不是撒瑪利亞人,因為耶穌基督是在猶太人那邊..
第三點,耶穌宣告在往後的日子,就是在他施行了救贖之後,蒙上帝閱立的敬拜應該是怎麼做的?.
經典不再重要,憲制不再重要,什麼禮儀都無關痛癢,依靠什麼來敬拜才是關鍵性..
對很多人,特別是古代的人來說,在哪裡拜神,怎麼拜神是一個重要的問題..
錯廟,拜錯神,是嚴重失誤的..
不過,這樣的問題其實是將神明部落化了,不同的神明各據山頭,.
上帝管這裡,撒旦管另外一處,所以你進錯廟,你進了撒旦的廟,進了撒旦的地頭,.
那麼上帝就收不到風了,聽不到你說話了..
這些說法是將上帝局限變成一個部落神明,.
在耶穌眼中,去拜神根本不是重要問題,.
連怎麼拜神都不是重要的,關鍵是用什麼態度,什麼境況來拜神..
耶穌一語道破這個關鍵,上帝是一個靈,所以拜祂的,必須用心靈和誠實去拜祂..
一直以來,有不少學者主張,心靈和誠實最好翻譯成為聖靈和真理,.
他們在說,真理不是指一般的教義或者理論,而是指道成肉身,.
充充滿滿的有恩典,有真理的耶穌基督,耶穌是真理..
所以在新約之後,真正的敬拜要藉著聖靈和基督而進行,.
唯有奉基督的名義,靠聖靈引導的敬拜才是真正的敬拜,才能蒙上帝月立..
你們應該聽過這些說法,不過我是完全不同意這樣的詮釋,.

$^{161}$這個詮釋是荒謬的..
耶穌基督為什麼要和一個撒瑪利人說敬拜神學呢?.
毛澄澄在這裡說什麼敬拜神學呢?.
他如今要做的是帶引這個婦人有生命突破和改變,.
所以什麼以基督中心,聖靈中心的敬拜,.
通通言不及義,和上下文根本不相干..
耶穌要和這個婦人說的是,在哪裡敬拜上帝不是問題,.
真正的關鍵是,我們必須用我們自己真實的本相,.
用我們的心靈和本來是一個靈的上帝接觸,聯繫,.
地點不重要,形式不重要,真誠真心才重要,.
in spirit and in truth指的是人內在的生命,而不是什麼三位一體,.
spirit這裡指的是人的靈,不是聖靈..
耶穌這樣的主張其實是貫穿在整個福音書裡面,.
和敬拜上帝有關的所有,我們說的是所有的教導,.
特別是在登山寶訓裡面,獻什麼禮物給上帝不重要,.
你不要帶著仇恨和負面情緒去見上帝才重要,.
祈禱和禁食的形式不重要,真誠真實面對那個在你還沒開口,.
祈禱之前就已經認識你的本相的上帝才重要..
我相信耶穌在這裡說的是心靈和誠實,.
祂要求婦人回歸自己的本心,正視自己的現實,.
包括那個很不理想的現實..
Need for complete sincerity and complete reality in our approach to God,.
真實的本相,誠實的面對自己,.
這個才是我們見到上帝的第一步,.
你永遠不可以裝模作樣見上帝,.
你永遠不可以化妝見上帝,.
當然我不是反對姐妹化妝,.
你永遠是以我們最真實的面貌,.
因為反正上帝都會將你打回原形,.
我們永遠是在我們的生命的原點,.
你會見到你的渴求,你的需要,.
你的脆弱,你的失敗,你的賞忘..
從講到怎樣才是真正的敬拜,.
婦人馬上聯想到無論是猶太人還是撒瑪利亞人都等候的彌賽亞,.
在25節她說,我知道彌賽亞就是那個稱為基督的要來,.
他來了必將一切的事都告訴我們,.
她這樣說,耶穌馬上就應她,.
這和你說話的就是他,26節,.
耶穌是彌賽亞,耶穌是基督,.
這就是福音的核心,.

$^{201}$從水到活水,從丈夫到敬拜,.
兜了一個大圈,.
終於來到耶穌是基督這個核心的訊息,.
基督教不是說我們去找耶穌,.
因為耶穌基督有答案有出路,.
我們是說我們來找耶穌,.
因為耶穌是答案和出路,.
我們不是藉著耶穌去找解決我們問題的方法,.
我們是要確認耶穌基督就是我們的方法,.
我們的出路..
有趣的是,耶穌和這個婦人談道的技術在這裡就結束了,.
沒有提到婦人對耶穌自認為基督的說話有什麼回應,.
這裡沒有直接說她因此信了耶穌,.
或者門徒在這個時候剛剛買了食物回來,.
婦人看到有這麼多人在,就趁勢就終止了對話,.
聖經說她匆匆離去,連打水的器皿都沒有帶走,.
不過她不是帶著恐懼和羞愧而去,.
她卻是走到城裡,走到人群中,.
將她所遭遇的耶穌介紹給她們,.
29節說:你們來看,.
有一個人將我素來所行的一切事都給我說出來了,.
莫非這就是基督嗎?.
我遇見了一個人,他很可能是基督,.
因為他竟然知道我過去所做的一切事,.
所以不如你們來看一看他吧..
最後就帶了一些人出城外井旁去看耶穌..
我們可以相信耶穌這一次的傳福音是成功的,.
因為在38,39節有記載,.
那一副人不僅是相信了耶穌,.
他還帶了很多人去信耶穌,.
然後他們還邀請了耶穌去他們的城裡住了兩天,.
然後耶穌在那裡再帶了更多人去信耶穌,.
然後那些信了耶穌的人就責怪那副人,.
我不是聽你的見證而相信,.
是我自己聽到,我自己相信的..
所以很明顯,這個婦人信了耶穌應該沒有什麼走彈的,.
並且不僅僅她信了耶穌,.
她還立刻成為見證者,.
帶了很多人同時信耶穌..
這是第一個我們可以相信,.

$^{241}$這個傳福音的事件是成功的原因,.
第二個我們相信它是成功的原因,.
是因為耶穌自己對這次傳福音很滿意,.
當門徒帶了一些食物回來請耶穌吃的時候,.
耶穌說不吃了,我飽了,.
他說你沒有食物,你怎麼會飽呢?.
耶穌怎麼回答?.
他說我的食物就是完成了猜我來的父上的功,.
我能夠做到父上要我做的事,.
我自己很開心,開心到飽了,.
不想再吃東西了..
耶穌一定很滿意他這次傳福音的經歷,.
覺得自己很開心,.
做到父上對他的委託和期望,.
很明顯這次一定是一個成功的傳福音案例,.
否則耶穌不會這麼開心的..
一次成功的傳福音案例,.
我們可以簡單做一些總結,.
雖然有很多東西可以說,.
不過我們唐崇拜不能處理很多東西,.
從婦人在人前所做的見證,.
我們看到她對耶穌的說話最印象深刻的,.
不是說活水,不是敬拜地點的討論,.
而是什麼?.
耶穌知道她的過去,.
28節,38節,.
我再重複說,.
她將我的真相底牌都揭穿了..
這個真相底牌,.
是這個婦人以前都不想面對的,.
被人揭發自己過去的醜事,.
特別是一直迴避,.
不想別人知道,.
不想別人討論的隱密事,.
這個婦人理論上應該很不開心的,.
應該覺得很尷尬,.
很羞恥,.
很受傷害的,.
不過這個婦人竟然沒有負面的感受,.
反倒是得著釋放,.

$^{281}$到處告訴別人,.
有人知道我的過去,.
並且用這個來為耶穌做見證,.
為什麼?.
真的很有趣,.
我相信婦人過去過著一個不道德的生活,.
也活在一個遭到其他人異樣眼光的環境,.
周圍的人都好像在批評她做錯了,.
她是一個罪人,.
為此她選擇逃避人群,.
盡量不和別人交往,.
在逃避別人的同時,.
她其實是在逃避自己,.
她不願意以自己的真面目和別人交往,.
沒有以自己真實的面貌和別人建立關係,.
唯獨這一次和耶穌相遇,.
一個不用她招供,.
就已經知道她所有真相的人,.
竟然完全沒有帶著批判否定的眼光,.
這令她感覺到很舒服,.
很自在,.
這種自在的和別人相處,.
和自己相處的經驗,.
是她從來沒有過的,.
無論如何,.
從經文我們看到,.
耶穌和婦人談到的整個過程,.
一直以一個很友好的氣氛進行,.
耶穌指出一種溫暖,.
關懷,喜愛,.
興趣和尊重的正面態度,.
沒有讓婦人感覺到被人評價,.
沒有被人指點點的威脅感,.
並且耶穌沒有糾纏於婦人所做過的各種壞事之中,.
只是輕輕讓她覺察到,.
這個過去是她如今的起點,.
她可以從這個起點鋪開一個新的旅程,.
一個正在發展,.
正在形成中,.
on becoming,.

$^{321}$becoming的一個過程,.
沒有被過去定格,.
沒有受過去約束,.
弟兄姊妹,.
我們傳耶穌基督的福音,.
我們不是要展示這個世界如何罪孽深重,.
遠離上帝,.
而是要講上帝來我們中間尋找我們,.
和我們建立關係,.
讓我們的生命得以突破和成長,.
有不一樣的可能性,.
我們不是帶著批判和否定的眼光去接觸人群,.
而是只能夠,.
也必須要展現我們的愛,.
寬恕和接納,.
讓人解除防範,.
得以正視自己的本相,.
然後用心靈,.
用誠實,.
來和做他愛他的上帝建立關係,.
我們對這個世界只有一個信息,.
只有一個稱為福音的,.
就是上帝接納人有新的可能..
\newpage



\section{腓立比書 2:12-18}
\label{sec:_sNrPBKmvwc}
\textbf{2024年6月23日崇拜直播｜羅志雄牧師｜靈性·關係·成長－活潑的生命｜腓立比書2\_12-18}
\newline
\newline
連結: \href{https://youtube.com/watch?v=_sNrPBKmvwc}{\texttt{https://youtube.com/watch?v=\_sNrPBKmvwc}} ~~~~ 語音日期: 2024-06-24
\newline
\newline
\hyperref[sec:0Rt7bY88Mpk]{\small{< < < PREV SERMON < < <}}
~
\hyperref[sec:index_chronic]{\small{[返順時目]}}
~
\hyperref[sec:l1Zxe_Q_Nro]{\small{> > > NEXT SERMON > > >}}
\newline
\newline
腓立比書 2:12-18
\newline
\begin{longtable}{cl}
\hline
\hline
章節 & 經文 (和合本修訂版)\\
\hline
2:12 & \begin{tabularx}{0.7\textwidth}{X} 我親愛的，這樣看來，你們向來是順服的，不但我在你們那裡，就是我現在不在你們那裡的時候更是順服的，就當恐懼戰兢完成你們自己得救的事； \end{tabularx} \\ \\ \relax
2:13 & \begin{tabularx}{0.7\textwidth}{X} 因為是神在你們心裡運行，使你們又立志又實行，為要成就他的美意。 \end{tabularx} \\ \\ \relax
2:14 & \begin{tabularx}{0.7\textwidth}{X} 你們無論做甚麼事，都不要發怨言起爭論， \end{tabularx} \\ \\ \relax
2:15 & \begin{tabularx}{0.7\textwidth}{X} 好使你們無可指責，誠實無偽，在這彎曲悖謬的世代作神無瑕疵的兒女。你們在這世代中要像明光照耀， \end{tabularx} \\ \\ \relax
2:16 & \begin{tabularx}{0.7\textwidth}{X} 將生命的道顯明出來，使我在基督的日子得以誇耀我沒有白跑，也沒有徒勞。 \end{tabularx} \\ \\ \relax
2:17 & \begin{tabularx}{0.7\textwidth}{X} 我以你們的信心為供獻的祭物，我若被澆獻在其上也是喜樂，並且與你們眾人一同喜樂。 \end{tabularx} \\ \\ \relax
2:18 & \begin{tabularx}{0.7\textwidth}{X} 你們也要照樣喜樂，並且與我一同喜樂。 \end{tabularx} \\ \\
[1ex]
\hline
\hline
\end{longtable}
$^{1}$(教會禮儀)今天經訓是記載在新約聖經.
肥臘比書第二章十二至十八節.
愛尼教會的聖經二百七十四頁至二百七十五頁.
新約肥臘比書二章十二節.
本來這樣對肥臘比信徒說.
這樣看來我親愛的弟兄.
你們既是常馴服的,不但我在你們那裡.
就算我如今不在你們那裡更是馴服的.
就當恐懼戰驚,作成你們得救的功夫.
因為你們立志行事.
到時神在你們心裡運行.
為要成就他的美意.
凡所行的都不要發怨言,起爭論.
使你們無可指責,誠實無畏.
在這彎曲薄貌的世代.
作神無瑕疵的兒女.
你們顯在這世代中.
好將明光照耀.
將生命的道表明出來.
叫我在基督的日子.
好誇我沒有胸袍也沒有屠奴.
我以你們的信心為貢獻的祭物.
我若被囂墊在旗上也是喜樂.
並且與你們眾人一同喜樂.
你們也要照樣喜樂.
並且與我一同喜樂.
.
各位靈姐妹主內平安.
今日會延續這個系列.
靈性關係成長關於肥立比書.
是第四講.
我們一齊開始之前有個禱告.
求你的靈與我們每一個同在.
讓我們能夠在上帝你自己的聖靈引導裡.
我們的心懷意念都蒙你保守.
我們可以聽到上帝你自己對我們親自的說話.
讓我們的生命能夠在這個世代見證你.
求你與我們同在.
我們來獻上祈禱奉主耶穌.
阿們.

$^{41}$在這裡先問大家一個問題.
你問我們現在身處在一個什麼世代.
我相信可能是一個惶恐.
混亂甚至迷失的世代.
我們剛才聽到弟兄和我們讀的這段聖經.
再對應經文裡肥立比教會的時候.
我們就會去想一下.
我們怎樣才能夠活出一個活潑基督徒的生命呢?.
我們看回這段經文.
剛才所讀的.
是承接著上一段第五節至十一節的經文.
這裡是說到以耶穌基督的心為心.
來到這裡的時候.
保羅就以轉接.
就說這樣看來.
開始要告訴我們.
耶穌基督是以謙卑馴服的生命.
願意降生為人的形象的樣式.
雖然祂本有神的形象.
但是祂釘死在十字架上.
是要信服神的旨意.
成就神救贖的心意.
不過在神的應許裡.
十字架我相信不是一個終局.
而是通往榮耀的一個里程上.
所以保羅來到勸勉信徒.
學習耶穌基督信服的生命.
是要以喜樂地.
活潑地在地上過我們的生活.
直到正如保羅所說的.
在基督耶穌裡面的日子.
我們來看一看.
這一段經文第十二至十三節.
是一個怎樣信服的生命呢?.
這樣看來,我親愛的弟兄,.
你們既是常信服的,.
不但我在你們那裡,.
就是我如今不在你們那裡,.
更是信服的..
就當恐懼戰驚,.

$^{81}$造成你們得救的功夫,.
因為你們立志行事,.
都事在神你們心裡運行,.
為要成就他的美意..
其實,.
在肥勒比教會已經有信服的基礎.
不知道大家知不知道.
究竟肥勒比的信服基礎在哪裡?.
或者我們簡單地重溫一下.
肥勒比的信徒.
他們信服在基督的權下.
因為保羅一開始就說.
他們被稱為基督裡的聖徒.
第二就是.
肥勒比的信徒信服福音.
因為保羅在肥勒比成立教會的時候.
一群當時還沒有認識耶穌的人.
還沒有認識福音的人.
他們能夠接受福音.
但是在那個時候開始.
他們就努力地傳福音.
所以保羅說他們是同心合意.
興旺福音.
而最後.
這群信徒.
他們都是信服在保羅的教導裡面.
在很多的教導上.
其中有一樣就是.
這群信徒很願意支持保羅的事工.
無論他去到哪裡都是這樣.
所以在教後的篇幅書信裡都會有提到.
說到這個信服的前提.
保羅勉勵他們.
就是如今我不在你們那裡更是信服.
就當恐懼戰經.
造成你們得救的功夫.
因為你們立志行事.
都使神在你們心裡運行.
為要成就他的美意.
在我預備這段經文的時候.

$^{121}$我都看過不同的聖經版本的翻譯.
在聖經當代譯本修訂版.
他就將這個更是信服.
就當恐懼戰經造成你們得救的功夫.
是譯為要戰戰兢兢地.
活出得救後應有的生活.
而在聖經信息版.
Eugene Peterson.
The Message Bible裡面.
他就將它這樣譯為.
你們生活在有求必應的信服當中.
繼續努力吧.
加倍努力.
當中信服.
如果在Eugene Peterson裡面.
原本只是Obedience.
但是他在前面就加了一個字.
就叫做Responsive Obedience.
剛才讀的時候就變成了有求必應的信服.
究竟是指什麼呢.
這裡原來就是說.
隨時準備回應信服的一種態度.
弟兄姊妹.
不知道有沒有發覺.
我們每個信徒都會向神作出回應的信服.
不知道大家有沒有這樣做過.
第一種他們參與教會服事的時候.
就會在教會.
看到有什麼需要的地方.
他們會去協助幫助.
甚至可能不只是做一些很簡單的服事.
甚至會擔起一些事奉的崗位.
可能做某一部部長或部員.
去教主學.
詩班.
唱詩.
甚至如果覺得自己音樂好.
有恩賜的話.
可能就會做指揮.
或者做詩琴.

$^{161}$但為了在事奉中有更好的果效.
可能會進修不同的課程.
甚至會去神學院進修一些神學的知識.
來幫助他們在事奉上.
可以更加明白上帝的心意.
又或者去參與學一些其他的課程.
因為他們覺得其他的知識.
可以幫助他們在事奉上.
可以更有發揮.
甚至會成為教會領袖.
發揮他的恩賜才幹.
這些回應福音的服事.
不知道大家怎麼看.
只是在做.
事奉式的服事.
我想上帝也是欣賞的.
不過原來有另一種.
要回應福音的服事.
不知道大家能做到多少呢?.
這就是內心.
心靈深處一種靈性的服事.
一種的服事.
它是要經歷與神的關係.
因為當你立志行事.
加上神藉著聖靈.
在你的心裡運行的時候.
我相信這個服事會更加持久.
更能夠成就神在信徒身上的美意.
而這種信服是一種內在的信服.
一種必然的信服.
所以當我們要向神作出.
有求必應的服事回應的時候.
就是剛才所說的兩種.
一種就是外顯.
外在的.
或者我們常說的doing.
這一種信服.
當然在事奉裡面.
我們回應神的召命.
但是另外剛才所說的.

$^{201}$內在靈性的信服.
如果我們沒有這種靈性.
服事的這種操練.
去回應上帝的召命的時候.
我相信很可能會有些偏差.
又或者在事奉的時候.
往往就會成為.
用人的心意來服事.
多於由聖靈來帶領大家.
以致到如果我們加上有.
又doing又being的時候.
這一種信服.
我相信教會整體來說.
就會有健康的成長.
兄弟姐妹.
我們當然會知道.
按照神在我們心裡面.
神的靈在行事為人的時候.
我們就能夠活出耶穌基督的樣式.
但是很多時候.
我們真的未必能夠做到.
所以保羅他的經驗.
來跟我們分享.
他說立志為善由得我.
此事行出來由不得我.
我不知道大家會不會有這種狀況.
因為很多時候我們都容易被罪.
來引誘去試探我們.
以致我們服事教會的時候.
就以為完全是信服在神的裡面.
保羅夢耶穌親自向他顯現的時候.
他要成為外邦人的使徒.
向外邦宣講福音.
不過在這之前.
他有幾年的時間.
操練自己的靈性.
與神建立關係.
然後才開始他的服事.
所以,靈姐妹.
我們常常都要保持一個.

$^{241}$「戰經」謹慎.
在聖靈保守的底下.
我們才能夠活出一個.
與我們蒙召的恩相稱的.
一個基督徒的生活.
一個的生命.
所以我們要常常讀聖經.
提醒我們不要靠自己.
以為我們自己立志.
就可以單單靠自己.
相反我們更加需要依靠.
在我們心裡運行神給我們的立志.
以致有神的保守.
我們才不會偏離正確福音使命的方向.
藉著信服主耶穌基督.
來過聖潔的生活.
如果這兩種信服配合起來.
我相信在神的事宮裡.
就能夠完成了他美好的旨意.
我相信對史托德牧師John Stock.
這位牧者大家都很熟悉.
他曾經有一個分享.
他說他每天有不同的服事.
有很多的宣講.
甚至是教學.
他分享的是他每天都要以雙膝來服事.
雙膝的意思是他用膝頭來服事.
來達成他這種信服神.
他怎樣服事呢?.
他說這個才能成為他的動力.
在他任何的情況下.
他都會透過禱告,讀經,退修.
和常常在神的面前更新他自己的生命.
內裡的生命.
以至他能夠過一個信服的生活.
有一個分享是在上星期前.
我重新看到他.
在我以前在教會團契成長的弟兄.
因為他和太太移居了新加坡.
因為他的太太是馬來西亞的華人.

$^{281}$所以就去了那邊生活.
最近就透過Facebook聯絡他.
他回來香港探親.
一起出來喝下茶,喝下高茶的時候.
除了說一下我們這幾年的生活之外.
大家都不約而同地說教會的狀況.
原來他和太太都在教會裡來事奉.
他自己也是神學院的一個校董.
一直分享香港和新加坡教會的狀況的時候.
發覺不約而同地說.
大家都很拼命地做.
不停地做,正如我剛才所說.
但做完之後.
每個人都會burn out(失去活力).
甚至現在的狀況.
香港和新加坡的弟兄姊妹.
年齡層好像是差不多.
年青的一輩的狀況都不太理想.
但再說下去的時候.
為什麼我們的生命.
在教會裡一直成長.
會去到這樣的狀況呢?.
我們一直在分享.
原來我們內裡的靈性的狀況.
我們的being(存在)出了偏差.
出了問題.
於是他就繼續分享.
他說現在我們一有人信耶穌的時候.
就立刻開始門徒訓練.
差不多是門徒特訓.
因為他用的英文叫intentional discipleship.
不知道怎樣是intention.
原來他參照了美國教會的門徒訓練.
他將我們日常生活.
我們日常生活相信有七個層面.
我們的家庭,財務,成長過程.
人生的成長.
可能包括了宗教信仰.
甚至我們的健康狀況.
無論家庭的健康.

$^{321}$或者其他的健康.
可能都包括在教會上.
社交,人際關係,靈性,工作.
在這七個範疇裡.
他就用一個叫特訓式的門徒訓練.
即是intentional discipleship.
來幫助弟兄姊妹成長.
我想這讓我們更加清晰明白.
當然我們知道任何一些訓練工具.
它只是一種工具.
但我們是否認真對待屬靈的狀況.
如果我們只是去做.
又或者沒有檢視到屬靈的狀況裡.
究竟有多少軟弱.
有多少缺點.
我們是否願意謙卑下來.
信服在神的面前.
然後才能讓我們在聖靈裡.
在屬靈的生命裡成長.
這一種門徒訓練.
應該是將我們日常生活裡的訓練.
加在聖經的真理裡.
我想這不是我們今天發明的.
如果大家打開福音書裡.
去看耶穌如何訓練門徒的時候.
其實這幾個範疇.
耶穌也是不斷地.
重重複複地和一些門徒實習.
我相信.
如果是這樣將神的話.
將真理 將生命之道.
融入在我們生命裡.
我相信我們整全的生命.
無論在靈性 在關係 在成長上.
都會有更精彩.
更土神喜悅的一面.
我們來看第十四至第十六節.
這裡也是承了上面第十二至第十三節的吩咐.
保羅是要信徒將得救.
以兒女的身份實實在在地.

$^{361}$在生活中展現出來.
我們在座當中.
相信我們都會承認我們是神的兒女.
究竟我們會有什麼樣的生活態度呢?.
在這裡很想邀請弟兄姊妹來讀第十四至第十六節.
我們一起來讀第十四至第十六節.
凡所行的.
(念念不忘).
好 麻煩弟兄姊妹.
我曾經聽聞某一間教會.
因為一些重大的內部事務而產生不和.
彼此發怨而起爭論.
這些狀況不斷擴大.
牧者當然會收到很多這些訊息.
事前會做很多的調解.
有很多的勸勉.
甚至會找幾個核心的成員.
來做一些教導和勸勉.
但都沒有辦法收拾這個局面.
其中一方就說要提告法庭.
這更加違反聖經的教導.
如果有任何的爭執.
我們要在教內處理.
但另一方又怎樣呢?.
他說好啊 奉陪.
不過呢.
這些重費就要由教會支付.
不知道大家聽到這些狀況的時候.
有什麼想法呢?.
我相信實在會失去基督徒生命的見證.
弟兄姊妹 再問問大家.
最近有沒有發怨言?.
有沒有起爭論呢?.
有沒有?.
好靜啊.
其實我都會抱怨.
有發怨言 有爭論的時候.
不過當我們停一停的時候.
我們會想一想整件事情發生的經過.
我們是否會需要發怨言起爭論呢?.

$^{401}$又會問一下為什麼呢?.
當我們這樣做的時候.
很可能我們從來都沒有想過.
後果會是怎樣的.
原來發怨言起爭論.
在保羅的書信裡是孖公仔出現的.
在哥林多書信裡.
保羅就這樣警告教會的弟兄姊妹.
你們也不要發怨言.
像他們有發怨言.
就被滅命的所滅.
在這裡我們剛才所讀的經文第十四至第十五節.
保羅意有所指.
因為他正在用發怨言起爭論.
來警惕菲律賓的信徒.
發怨言起爭論.
如果有明白的話.
這裡其實是引致出埃及記.
如果大家知道當時摩西帶領以色列人出埃及的時候.
發生了什麼事.
這個例子就來勸勉菲律賓的信徒.
要引以為鑒.
另外保羅也要求信徒要無可指責.
不知道會不會覺得保羅的要求很高.
他叫價實在太高.
標準很離地.
但是我們真的會知道.
無論有多努力.
在任何情況上.
你所做的事可能做到很完美.
做到很頂級.
自己覺得是.
但你猜猜一定會不會有人喜歡或欣賞呢.
要無可指責.
相信要超越那種完全.
可能是一個完美主義.
不過當我們認定我們的身份的時候.
我們是蒙恩德的人.
應該不會再在此事而非.
或者在罪惡世界裡.

$^{441}$這些彎曲薄茂世代裡.
這樣的權勢之下.
受到影響.
因為我們是屬神的.
因著他的應許.
接受了耶穌基督的福音的時候.
我們就和耶穌和基督有結連.
所以我們被稱為神的兒女.
既然我們是神裡分別出來的人.
雖然我們還沒有到完全的地步.
但在神的恩典裡.
在他的引領之下.
我們是朝著這個完全的方向走.
所以保羅就這樣.
很特別很鄭重地提醒哥林多信徒.
在屬靈的路上的時候.
應該邁向神應許的完全.
而必須避免怨言和爭論.
這就是神兒女第一個品格的特質.
耶穌基督教導門徒要彼此相愛.
因此眾人就認出你們是我的門徒.
彼此相愛.
我們都很熟悉.
亦是門徒的標記.
但有沒有想過耶穌教導門徒.
又命令門徒彼此相愛的時候.
是什麼狀況呢?.
耶穌當時是一個最不友善的情況.
是受到不能避免的逼迫.
很多艱難裡.
向門徒提出的要求和命令.
所以我們今天.
縱然同樣面對這些艱難.
困境或前路迷茫的時候.
保羅提醒我們.
不要彼此埋怨.
不要彼此爭論.
以至破壞彼此合一.
彼此相愛.
彼此同心.

$^{481}$反而我們更要謙卑地.
來到神的面前.
同心仰望救贖我們這位主.
仰賴祂的幫助.
我們才能達到成為.
神給我們的兒女的身份.
第二個神兒女的品格.
就是誠實無偽.
保羅用這個字非常獨特.
可以譯為純真無邪.
是指品格的本質.
是不會有雜質.
如果你是鑽石.
應該是不會有雜質.
不會有其他成份.
參雜在裡面.
所以誠實無偽.
保羅要提醒.
信徒的內心.
不應該參雜任何邪惡.
任何詭詐的地方.
要誠實無偽.
像光明的子女.
很多次聖經裡說.
神是光在他毫無黑暗.
所以.
神是真理的源頭.
我們也是祂的兒女.
我們要將神的真理.
福音反映在我們的生命當中.
在充滿著黑暗.
扭曲.
惶恐迷惑的世代裡.
來彰顯神的生命之道.
最後來到經文裡的部分.
雖然短短有兩句.
但其實充滿著.
保羅服侍的心境.
就是他最關切的生命取代.
如何才能真正地.

$^{521}$精彩又謙卑的生命事奉.
第十七節.
保羅這樣說.
我以你們的信心為貢獻的祭物.
我若被曉奠在其上.
也是喜樂.
並且與你們眾人一同喜樂.
你們也要照樣喜樂.
並且與我一同喜樂.
人的生命.
如何才能覺得精彩.
覺得喜樂.
覺得活潑.
我不知道大家如何想.
如何看.
可能不同人有不同的想法.
也有不同的取態.
但對於保羅來說.
生命是有限制.
甚至軟弱.
有很多困苦.
雖然如此.
但保羅仍然以見證耶穌.
成為他終生的職志.
一個照明.
所以保羅親自建立.
菲律賓教會的時候.
其實他的處境.
實在都是一樣.
可能與我們今時今日一樣.
在一個對福音不友善的環境裡.
建立菲律賓教會.
但保羅神給他的恩典.
他能夠看到信徒可以成長.
整個信仰群體.
可以在靈性上長進.
彼此的關係.
會有美好的見證.
這會令保羅心滿意足.
願意為他們獻上祭物.

$^{561}$所以保羅說.
以你們的信心為貢獻的祭物.
而自己若被囂墮在堤上.
也是喜樂.
保羅的意思是.
菲律賓教會的信心.
是一個主要的祭物.
而它只是一個.
以陪伴的形式.
是一杯酒.
是一杯電酒.
即是說它是一個.
祭物以外的附屬物.
它只是繞在這些祭物之上.
就算是很短暫也好.
這個時間出現在上面.
就會消失.
但保羅來說.
他也極其滿足和喜樂.
無怪乎保羅在這兩節經文裡.
表達出.
菲律賓信徒關於信心的時候.
他就說有極大的喜樂.
這個喜樂.
如果大家有留意.
數到有多少次?.
應該有四次吧?.
有聖經學者說.
兩次是屬於敘述性的描述.
保羅自己喜樂的態度和狀態.
另外兩次.
保羅對菲律賓信徒的命令.
無論是哪一個.
是敘事還是命令.
保羅同樣以兩個相關動詞來表達.
其中一個是指出.
他個人的喜樂.
另一個是指出.
這個喜樂必須要共同擁有.
和一同分享.

$^{601}$所以.
如果大家留意經文的時候.
他的狀況.
其實對於保羅來說.
就是我和你.
你和我.
一同的喜樂.
當然我們明白.
我們在座當中.
我們的喜樂很多時候都是我們個人的事.
你分享給別人聽.
或者在你的小組.
在你的團契裡.
他會不會跟你一樣.
有同樣的領受.
或者有同樣的經歷.
同樣有同樣的喜樂呢?.
我相信可能會有.
但是否完全一致的喜樂呢?.
但保羅在這裡提醒.
既然在一個信仰群體的時候.
有一個這麼密切的關係.
所以真正的喜樂.
就要在群體裡.
互相分享出來.
我們會看到.
保羅和菲勒比信徒的關係.
是非常親密的.
他們的情誼.
已經超出了弟兄姊妹之間的情誼.
所以保羅的喜樂.
一定有他的影響力.
所以能夠帶動教會的信徒.
在面對逼迫.
不友善.
在困難的情況下.
都要鼓勵菲勒比的信徒.
能夠產生喜樂的滿足.
弟兄姊妹.
這些是不是真的很難?.

$^{641}$真的挺難的.
難怪保羅在第二章較早的部分.
他說要鼓勵菲勒比的信徒.
不要只顧自己的事.
要顧別人的事.
弟兄姊妹.
你現在會不會遭遇到.
或者正在遭遇一些困難.
令你很憂心很掛慮.
內心的喜樂可能已經褪色.
甚至毫無喜樂.
經文在這裡.
很想鼓勵你.
重新想想.
你經歷了耶穌基督的福音.
得到這份救贖的時候.
你去再想一想.
我們都知道.
患難困苦確實是常數.
亦都是我們無法避免.
但這本經文是鼓勵.
鼓勵這些弟兄姊妹.
你要再一次回到你的信仰群體.
重新跟他們結連.
在基督的福音裡.
有一個轉換心懷意念的機會.
重拾從上帝而來給你的喜樂.
在結束之前.
我會講一個小小的.
關於喜樂生命的感染力的故事.
話說有一位兒哥.
他的弟弟回到我們教會.
回到旺角浸信會.
回來十多年了.
這位弟弟.
我們稱他為弟兄.
當他回到教會的時候.
我知道他身體行動不方便.
甚至工作都沒有能力.
但在他仍在工作之前.

$^{681}$我聽說.
當爸爸有事的時候.
他辭了工作.
專心照顧爸爸.
直至爸爸離開這個世界.
他家人看到這位弟兄.
覺得真的需要照顧他.
因為他工作能力有限.
各方面都需要特別照顧.
他當時住在旺角附近的地方.
很近旺角浸信會.
所以是一次偶然的機會.
我相信是上帝帶領他回來.
在這裡聚集.
不知道大家會不會認得他.
這位弟兄已經安息主懷.
他的狀況就是.
他起下的生命.
其實是感染了這位義哥.
他見到教會幫助這位弟兄.
在他最艱難的時刻.
會照顧他.
甚至在他病患的情況下.
要做頸椎手術.
甚至之後中風.
要坐輪椅.
要去老人院.
甚至當他做頸椎手術的時候.
上帝恩典給他.
由旺角這個地方搬到.
遠離旺角區.
去到沙田大衛的屋村住.
那時他還可以回來教會.
但之後當他頸椎手術之後.
沒有辦法回來.
我們去探望他.
說他住的屋很好.
但其實他不開心.
他怎樣分享呢?.
他說其實我回不來教會.

$^{721}$回不來旺鎮.
和弟兄姊妹一起聚集.
回來團契.
但他義哥到最後.
和我們分享的時候.
他說就是他弟弟.
看到教會的弟兄姊妹.
教會的傳道人.
不斷去探望他.
而他這位弟弟.
在他這樣的人生的經歷裡.
他毫無怨言.
反而他覺得.
他不能回來教會.
他的生命好像少了一些.
這位義哥他分享的時候.
他說他原先的太太和兒子.
不斷叫他回教會.
他說他都不會回來.
但他見到他這個弟弟.
他的生命是很艱難的裡面.
但他仍然生命裡面.
都是綻放著.
上帝耶穌基督的福音.
這份喜樂的恩典的時候.
到最後.
他說他願意跟他的太太.
跟他的兒子回教會.
雖然他不是回我們旺角浸信會.
但原來這份喜樂.
生命的感染力.
確實是會感染到其他人.
在總結的時候.
不知道大家怎樣看.
今日肥立比書.
第二章十二至十八節.
短短七節的經文.
這裡其實提醒我們.
作為信徒.
作為基督的門徒.

$^{761}$我們怎樣過一個活潑的生命呢?.
經文提醒我們.
要將信服的生命.
守住神兒女身份.
這個召命.
和喜樂的生命.
三者結合起來.
縱然要面對一個.
惶恐失迷的世代.
甚至在不同的人生觀.
價值觀的衝擊下.
我們能否活出.
一個不一樣的生命.
同時又守住生命之道.
見證基督.
過活潑生命.
等候基督再回來的日子呢?.
我們一起來圖稿.
天父上帝我們這樣衝心的.
去讚頌.
將榮耀歸給你.
因為基督的福音.
拯救了我們每一個.
讓我們被稱為神的兒女.
求聖靈這一刻.
在我們每一個人的心裡動功.
我們能夠經歷上帝.
你自己的恩典.
我們學習信服.
學習保羅提醒我們.
我們生命靈女.
要信服在你自己的裡面.
在這個彎曲薄幕的世代裡.
來見證你.
因為你已經賜下你自己.
喜樂生命給我們每一個.
主求你來幫助我們.
我們生命是有軟弱.
是有不足.
但是主願你親自來調教我們的生命.

$^{801}$讓我們朝向上帝.
讓我們要走的道路.
我們要靠召命來經歷你.
求主來帶領我們.
我們獻上祈禱.
奉主耶穌基督的名.
\newpage



\section{腓立比書 2:19-30}
\label{sec:l1Zxe_Q_Nro}
\textbf{2024年6月30日崇拜直播｜陳啟興牧師｜靈性·關係·成長－工作不離三兄弟｜腓立比書2\_19-30}
\newline
\newline
連結: \href{https://youtube.com/watch?v=l1Zxe_Q-Nro}{\texttt{https://youtube.com/watch?v=l1Zxe\_Q-Nro}} ~~~~ 語音日期: 2024-06-30
\newline
\newline
\hyperref[sec:_sNrPBKmvwc]{\small{< < < PREV SERMON < < <}}
~
\hyperref[sec:index_chronic]{\small{[返順時目]}}
~
\hyperref[sec:8t3Xeh3MCK8]{\small{> > > NEXT SERMON > > >}}
\newline
\newline
腓立比書 2:19-30
\newline
\begin{longtable}{cl}
\hline
\hline
章節 & 經文 (和合本修訂版)\\
\hline
2:19 & \begin{tabularx}{0.7\textwidth}{X} 我靠主耶穌希望很快能差提摩太去見你們，好讓我知道你們的事而心裡得著安慰。 \end{tabularx} \\ \\ \relax
2:20 & \begin{tabularx}{0.7\textwidth}{X} 因為我沒有別人與我同心，真正關懷你們的事。 \end{tabularx} \\ \\ \relax
2:21 & \begin{tabularx}{0.7\textwidth}{X} 其他的人都求自己的事，並不求耶穌基督的事。 \end{tabularx} \\ \\ \relax
2:22 & \begin{tabularx}{0.7\textwidth}{X} 但你們知道提摩太是經得起考驗的，他與我為了福音一同服侍，待我像兒子待父親一樣。 \end{tabularx} \\ \\ \relax
2:23 & \begin{tabularx}{0.7\textwidth}{X} 所以，我一看出我的事怎樣了結，我希望立刻差他去， \end{tabularx} \\ \\ \relax
2:24 & \begin{tabularx}{0.7\textwidth}{X} 但我靠著主自信我不久也會去。 \end{tabularx} \\ \\ \relax
2:25 & \begin{tabularx}{0.7\textwidth}{X} 然而，我想必須差以巴弗提到你們那裡去。他是我的弟兄、同工和戰友，是你們差遣來供應我需要的。 \end{tabularx} \\ \\ \relax
2:26 & \begin{tabularx}{0.7\textwidth}{X} 他很想念你們眾人，並且極其難過，因為你們聽見他病了。 \end{tabularx} \\ \\ \relax
2:27 & \begin{tabularx}{0.7\textwidth}{X} 他真的生病了，幾乎要死。然而神憐憫他，不但憐憫他，也憐憫我，免得我憂上加憂。 \end{tabularx} \\ \\ \relax
2:28 & \begin{tabularx}{0.7\textwidth}{X} 所以，我更要盡快送他回去，好讓你們再見到他而喜樂，我也可以減少憂愁。 \end{tabularx} \\ \\ \relax
2:29 & \begin{tabularx}{0.7\textwidth}{X} 故此，你們要在主裡歡歡喜喜地接待他，而且要尊重這樣的人， \end{tabularx} \\ \\ \relax
2:30 & \begin{tabularx}{0.7\textwidth}{X} 因他為做基督的工作不顧性命，幾乎至死，為要補足你們供應我不夠的地方。 \end{tabularx} \\ \\
[1ex]
\hline
\hline
\end{longtable}
$^{1}$今天的經訓是記載在.
肥立比書二章十九至三十節.
在各位的聖經新約二百七十五頁.
聖經這樣說.
「我靠主耶穌指望快打發提摩太去見你們.
叫我知道你們的事.
心裡就得著安慰.
因為我沒有別人與我同心.
實在掛念你們的事.
別人都求自己的事.
並不求耶穌基督的事.
但你們知道提摩太的名正.
他興旺福音與我同流.
待我像兒子待父親一樣.
所以我一看出我的事要怎樣了結.
就盼望立刻打發他去.
但我靠著主自命自信.
我也必快去.
然而我想必須打發.
以巴伐提到你們那裡去.
他是我的兄弟.
與我一同作供.
一同當兵.
是你們所差遣的.
也是供給我所需用的.
他很想念你們眾人.
並且極其難過.
因為你們聽見他病了.
他實在是病了.
幾乎要死.
然而神連恤他.
不但連吻他.
也連恤我.
免得我憂上加憂.
所以我越發急速打發他去.
叫你們再見他.
就可以喜樂.
我也可以少些憂愁.
故此你們要在主裡.
歡歡樂樂地接待他.

$^{41}$而且要尊重這樣的人.
因為他作基督的功夫.
幾乎致死.
不顧性命.
要補足你們供給我的.
不及之處.
聖經獨筆.
好,弟兄姊妹早安.
看見樣子有點尷尬.
我們先將自己交託給上帝.
我們同心禱告.
親愛的主我們求你的靈與我們同在.
得著我們.
讓我們的心.
因你的聖經更加歸向你.
更加跟從你.
將我們整個人生交託給你.
同心禱告奉耶穌基督的聖名.
阿們.
有一間教會需要去抗唐.
唐主任帶領整個抗唐的侍工.
他很緊張這些侍工.
因為要買新的地方.
他就先去收奉獻.
他幾乎是強迫每一個弟兄姊妹.
都要定額奉獻.
有一個名單.
你不奉獻在名單上.
沒有這樣的金額.
你就不愛教會.
不屬靈了.
然後又強迫教會裡面很多的專業人士.
就是建築師.
設計師.
很多專業人士.
全部都要用大量的.
工餘的時間.
去協助抗唐的侍工.
弟兄姊妹真的被他抓得很辛苦.
下班.

$^{81}$還要有很長時間的加班.
最終就搞定了.
買到的地方.
開了一個獻堂禮.
很多會友都沒有出席.
這個侍工雖然成功.
但是你一看回頭.
我的形容就是.
屍橫遍野.
很多弟兄姊妹都受傷了.
沒有人再想跟這個唐主任合作.
弟兄姊妹.
結果重要.
還是關係重要呢?.
在一間教會裡面.
我們的心就是.
當然最好兩者都兼備.
一個侍工.
結果是重要.
但是關係同樣都是重要.
結果是一時的.
但是關係是長遠的.
一個事奉團隊.
如果.
它沒有這麼重要的元素.
其實是很難的.
這個元素就是.
同心.
同行的關係.
在一間教會裡面.
如果沒有一個同心同行的關係.
你的侍工.
其實到最後可能也有很多虧損.
我今天一回教會.
看到陳明全牧師.
他不止一次跟我說.
辛苦你了.
其實我心想.
港獨是不是最辛苦?.
應該也不是.

$^{121}$教主的學是不是很辛苦?.
應該也不是.
在教會裡面什麼最辛苦?.
就是關係破裂.
要處理一些破裂關係.
這個是最辛苦的.
在你的事奉路上.
我不知道你有沒有一些.
跟你同心同行的人.
我們叫同行戰友.
如果你沒有.
就不要氣餒.
快點去找.
孤零零.
孤單作戰.
從來不是上帝的設計.
保羅提摩泰和爾巴忽提.
就是作弓不離三兄弟.
是同行的戰友.
這段經文所顯現的是一幅.
很美麗的事奉圖畫.
他們有一些的事奉心態和品格.
是很值得我們學習的.
很值得我們操練的.
說一點背景.
在保羅的書信裡面.
當提到人物交代.
誰是大信人.
通常放在哪裡呢?.
都是書信的最後.
很少放在書信的中間.
好像這裡這樣.
二章尾就是.
在《肥勒比書》中間的部分.
為什麼會這樣呢?.
保羅有什麼特別.
要將這兩位弟兄放在中間呢?.
如果你以後看《肥勒比書》.
其實一章二十七節開始.
保羅提出一個問題.

$^{161}$就是要同心合意.
興旺福音.
同心合意是很重要的.
大家不能夠同心合一的原因.
保羅一直說.
因為我們自我中心.
自我膨脹.
要同心合一.
就要先謙卑自己.
所以在二章五節他說了.
你們要以耶穌基督的心為心.
耶穌基督的心腸是怎樣的呢?.
二章五到十一很清楚.
他是謙卑.
捨己.
是我們一個典範.
你要效法耶穌這個謙卑.
才能夠去同心合一.
所以.
在二章尾.
提出提摩太和以巴弗提.
這兩位同工作為一個例子.
是我們怎樣效法耶穌謙卑一個具體的實例.
所以思路一直沿著.
保羅在一章尾一開始說的.
他中間就加了這兩位同工.
作為一個例子.
當時保羅是在羅馬坐牢.
提摩太我們估計應該也在羅馬陪伴保羅.
而以巴弗提就是在亞西亞西岸.
在馬其頓.
肥納比教會的一個.
就是差派去羅馬支援保羅的一位弟兄.
帶了一些鬼陣.
想去支援和服侍保羅.
背景就是這樣.
看兩個段落.
第一個段落是十九到二十四.
就是提摩太和保羅.
接著二十五到三十.

$^{201}$就是講以巴弗提和保羅.
看看這三個人.
我用了作公不離三兄弟.
應該也是抄以前的一套電影.
叫作《作死不離三兄弟》.
不是講那套電影.
是在講這三位主內的同工.
他們的關係有些不同.
關係親疏有別.
靈命也未必一樣.
但是他們的心智是很相同.
我們先看十九到二十四這個段落.
「我靠主耶穌指望快打發提摩太去見你們.
叫我知道你們的事.
心裡就得著安慰」.
知道你們的事是指甚麼事呢?.
很可能是指剛才我講的一章二十七節所講.
菲律賓信徒很同心去慶王服音.
很同心去傳福音.
所以保羅的心很安慰.
我一知道你們這樣.
我就會很安慰.
為甚麼?很簡單.
保羅很著重傳福音.
他看重耶穌基督的事情.
然後他說.
「因為我沒有別人與我同心.
實在掛念你們的事.
別人求自己的事.
並不求耶穌基督的事」.
他用了「因為」開始.
在第二十節.
是解釋保羅為甚麼想猜猜提摩太.
他不是猜猜其他人.
為甚麼我特別要選提摩太呢?.
有以下幾個原因.
在二十節的上半節.
他與保羅同心.
在二十節的下半節.
他們是真真正正去關懷菲律賓信徒.

$^{241}$在二十一節.
他們是追求耶穌基督的事.
有這幾個原因.
所以保羅就決定猜猜提摩太.
如果大家熟悉提摩太這位聖經人物.
提摩太有很多限制.
他身體軟弱.
他經常患病.
他的性格比較柔弱膽怯.
所以在提摩太後書.
保羅提醒他.
神給我們的不是膽怯的心.
是剛強仁愛自制的心.
所以提摩太無論身體性格.
都有不少的限制.
可能在座的很多狀況.
都比提摩太優越.
是理想的.
即使如此.
但提摩太仍然可以有很美好的事奉.
所以弟兄姊妹不要被自己的限制嚇怕.
人有很多限制.
可能現在是工作的薄煞期.
很多限制.
可能你剛生了小朋友.
很多限制.
可能你的家人在患重病.
甚至你自己有重病.
很多限制.
但我鼓勵大家不要被這些限制嚇怕.
因為上帝的恩典是超越任何的限制.
如果你能夠拿到上帝的恩典.
靠主.
你就能夠做上帝要你做的事奉.
不是做你自己想做的事奉.
別人是求自己的事.
他求甚麼事呢?.
很可能是指第二章三四節所說的.
貪圖虛榮.
只顧自己.

$^{281}$不顧別人的人.
這些人表面上很可能會很熱衷於事奉.
可能做很多事情.
但他們求甚麼呢?.
求自己的榮耀.
在教會也有.
我猜每間教會都有.
我不太認識黃浸民.
不太認識黃浸的弟兄姊妹.
在香港的母會裡面.
我看到也不少.
我們的事奉帶著很多私慾.
是為了自己.
別人肯定我.
別人去讚賞我.
讓牧師去稱讚我.
在人面前得到榮耀.
但提摩泰是怎樣呢?.
提摩泰是求耶穌基督的事.
這個很重要.
他關心耶穌所關心.
他是為耶穌做事.
我們留意.
我們有時做很多事.
我們的動機是甚麼呢?.
我們是否以完成耶穌基督託付的使命.
是我們的目標呢?.
如果是這樣.
我們真的以耶穌基督的心為心.
其中傳福音是無可推諉的.
傳福音一定是耶穌很關心的事.
因為你看看第22節.
「但你們知道提摩泰的名正」.
「他興旺福音與我同路」.
「待我像兒子待父親一樣」.
我們是否.
只是以耶穌所掛念的事情.
是我們的掛念呢?.
一想這個問題.
我就很慚愧.

$^{321}$我經常掛念其他事情多於掛念耶穌的事情.
這個我很坦白去承認.
有時可能.
今晚歐國盃.
究竟英格蘭是怎樣的呢?.
今晚看不看呢?.
又這麼晚.
很記掛這些.
弟兄姊妹.
我們會否經常記掛.
那些不是耶穌記掛的事情.
甚至超越了耶穌記掛.
所以我們做信徒又不是很複雜.
我經常都是這一句.
我關心耶穌所關心的事情.
我就不關心耶穌不關心的事情.
盡量是這樣.
譬如我今早和太太吃早餐.
有ABCDEFG很多選擇.
我估計耶穌不太關懷我選A還是B.
我選了A之後.
究竟是奶茶還是咖啡.
我相信耶穌不太關懷.
所以我沒有禱告.
耶穌是奶茶還是咖啡呢?.
A還是C呢?.
我沒有.
我有吃飯.
我有吃之前禱告.
所以有些事情不關心.
但是你疑不疑問.
耶穌很關心.
你做甚麼工作.
耶穌很關心.
你和誰拍拖結婚.
耶穌很關心.
影響你的一生.
你怎樣運用你的時間.
運用你的金錢.
運用你的健康.

$^{361}$耶穌很關心.
我們是否去關懷耶穌所關懷的.
特別是如果這裡的經文.
這個向度就是全服音.
有人估計.
也根據較新的統計.
香港現在大約有四十多萬的基督徒.
這個應該是相似的.
因為有很多已經移民了.
只剩下四十多萬.
香港目前如果根據最新的政府統計.
七百五十萬人口.
很多人走了.
很多人也南下.
來了香港.
在香港如果你這樣看.
有多少人還未信主.
起碼七百萬.
七百萬人還未信主.
我們不需要經常想我們去哪裡宣教.
我們去哪裡短宣.
其實香港已經有極大的和長.
在你的職場裡面有多少人還未信.
很多.
我現在服務的金融界.
在香港有二十多萬人.
現在我們服務的是二百多人.
在我們的小組裡面.
很多人是未得之民.
你的左鄰右里.
你的鄰居有多少人還未信.
你的家人也可能有很多人還未信主.
問題是我們是不是很關心耶穌基督的事情.
很著意很積極向這些人傳福音呢.
在福音派的人眼中.
有一個異端.
叫做耶和華見證人會.
這個是我們看的.
根據耶正.
他在2023年有一個叫做全球的總報告.

$^{401}$他每年都出的.
應該是2022的9月到2023的8月.
大約是這一年.
叫做2023的報告.
他們有862萬5千位傳道員.
當然傳道員不是傳道人.
傳道員就是會傳.
就是去分享耶正的福音的人.
他們全年用了17.9億小時.
去傳他們的福音.
你算一算.
就是平均每人.
每一個傳道員.
每一個耶正的傳道員.
每天他們傳他們的福音有34分鐘.
也不是很多.
34分鐘.
但是你想一想.
自己.
你每天用多少時間去傳福音.
可能已經不太記得.
這個月有沒有用多少時間.
上個月.
在2024年用了多少時間傳福音.
可能有些人加起來也未必有34分鐘.
提摩太和保羅都是非常關注耶穌所關心的事情.
其中傳福音.
這個很重要.
所以保羅教導我們.
在這段經文.
我們要學習他和提摩太.
去求耶穌基督的事.
同心去傳福音.
提摩太的事奉心智是有印證的.
我們一路看下去.
有兩方面的印證.
第一就是在第22節的上半節.
在肥納比信徒的印證.
但是你們肥納比信徒.
知道提摩太的名證.

$^{441}$他興旺福音.
與我同路.
他們知道提摩太經得起考驗.
為了福音很努力.
第二個印證就是保羅自己的印證.
他們為福音一同去服侍.
提摩太待保羅好像兒子待父親一樣.
保羅就跟著那兩節.
保羅說他一知道自己受審訊的結果.
如何了結呢.
他能否放監.
能否離開羅馬呢.
他就打算差遣提摩太去到肥納比.
保羅自己也相信很快他能夠去到肥納比.
可以看到保羅很記掛他所牧養的羊.
他跟肥納比信徒的關係比較親切.
你看肥納比書.
你會了解到.
很有情誼的保羅和肥納比教會.
跟你看提摩太.
跟你看哥林多前書後書.
整個感覺都很不同.
我想大家留意一下.
在這個段落的最開頭和最結尾.
十九節.
「我靠主耶穌」.
第廿四節.
「我靠著主」.
這兩組的詞語.
首尾呼應.
保羅的心態是.
做得完這些事.
是要靠主.
不是靠自己.
如果我們.
剛才說我們要關心耶穌所關心的事.
其實也不是靠自己去關心耶穌所關心的事.
說得白一點.
就是我們靠著耶穌.
幫我們去關心耶穌所關心的事.

$^{481}$永遠都不可以離開耶穌基督.
保羅和提摩太的榜樣是甚麼呢?.
我們同心求耶穌的事.
尤其是傳福音.
弟兄姊妹.
在你們的事奉人生裡面.
有沒有這樣的同行戰友呢?.
跟你同心同行.
是求耶穌基督的事.
如果你有.
我很感恩.
如果你沒有.
千萬不要氣餒.
繼續找.
一定有的.
如果你用心去找.
一定有.
我們看第二個段落.
二十五到三十.
在第二十五節.
「然而我尚必須打發以巴弗提到你們那裡去.
他是我的弟兄.
與我一同作供.
一同當兵」.
「一同當兵」就是戰友.
這個字就是戰友的意思.
「是你們所差遣的.
也是供給我所需用的」.
保羅稱呼以巴弗提是弟兄.
同供.
和戰友.
弟兄和同供.
如果你看保羅的書信.
是保羅對其他事奉拍檔.
一個很慣常用的稱呼.
但是戰友就比較少用.
這個比較少.
保羅這樣稱呼以巴弗提.
顯示他們之間.
除了有親口的情誼.

$^{521}$他們有共同的目標.
戰友就是我們有共同的目標.
一同為這個目標奮鬥.
所以在教會.
我經常覺得.
同供.
教務同供.
長輯.
事奉人員.
我們不應該只是有公事公辦的關係.
這些是工作關係.
在公司就很普遍了.
在教會就不太好.
我自己覺得.
在教會裡面.
或者我們現在在機構裡面.
同供我們.
不是單單公事公辦的關係.
應該是同行戰友.
有一個這樣的.
我不知道你在旺站.
有沒有人可以稱呼.
你是我的弟兄.
你是我的同供.
這個就容易了.
你是我的戰友.
我們真的為著一些目標.
我們去奮鬥.
這是上帝給我們的目標.
保羅再次強調同心事奉的重要性.
他和提摩泰很同心.
他和爾巴弗提一樣.
很同心.
是弟兄同供戰友.
所以如果我們用保羅第二章的用語.
如果我們不是意念相同.
愛心相同.
我們很容易是甚麼呢?.
就會起爭論.
輕則口角.

$^{561}$關係很差.
重則是甚麼呢?.
會分裂的.
香港在過往這幾年.
大家都知道.
有很多人是拉隊離開教會.
因為不同的可能政見.
其實有人觀察北美也是這樣.
北美的華人教會出現很多新堂.
最主要的原因就是分裂.
原來這個分裂比主動去植堂而建立的新堂數量還要多.
這就是看那個地方的一些特色.
我們千萬不要這樣.
接著保羅說.
他(即是爾巴弗提).
很想念你們眾人.
並且極其難過.
因為你們聽見他病了.
大家知道爾巴弗提.
由馬其頓的菲勒比.
去到羅馬供應.
他有兩個任務.
第一個任務就是帶一些鬼陣.
去供應保羅.
第二個就是去到那裡服侍保羅.
代表菲勒比教會去服侍保羅.
其實估計由馬其頓的菲勒比.
去到羅馬大約最快也要四個星期.
當時沒有飛機.
要陸路去.
周居勞頓可能爾巴弗提就病了.
但是很奇怪的.
為何爾巴弗提會難過呢?.
爾巴弗提病了.
他不是因為病了.
他很不舒服.
他難過.
他是因為自己的病.
而令到別人.
別人就是菲勒比的信徒.

$^{601}$很憂傷而難過.
你想想.
他自己病了.
反而因為那些很記掛他的人.
很憂傷.
他就很難過.
這種是我為別人.
為他者的一個心腸.
這也是耶穌基督的心腸.
保羅怎麼看這件事呢?.
保羅我覺得在這裡.
他對自己比較嚴.
他對人比較寬恕.
他不把自己的事奉標準強加在其他同工身上.
我保羅三次宣教.
沒事的.
現在保羅已經是三次宣教之後.
他就在羅馬坐牢了.
你爾巴弗提一來到就病了.
如果保羅用自己的標準.
他肯定是比不合格的.
就是爾巴弗提.
但是保羅沒有.
保羅是為爾巴弗提著想的.
爾巴弗提應該來到去服侍他.
反而病了.
反過來要保羅服侍爾巴弗提.
保羅的著眼不是你也幫不了我.
幫不了忙.
我是使徒.
我的身份多厲害.
我比你大很多.
爾巴弗提.
坦白說在《新約聖經》.
也不是一個很著名的牧者.
保羅很重視這位弟兄.
很重視這位同工.
保羅也是一個為他者的事奉態度.
我記得當我在武會事奉的時候.
有一次我去台灣台中短孫.

$^{641}$教會在台中一個叫霧峰的地方.
它叫做和睦中心的教會.
教會有一句使命宣言.
這個使命宣言真的很貫徹.
落實在教會裡面.
他們教會整個事奉的理念和精神.
他們的事奉宣言就是.
從別人的需要看見自己的責任.
這句就是他們的使命宣言.
他們不是真的說的.
因為那段時間我就跟他們到處走.
他們去教會服侍小朋友.
提供免費的功課輔導.
這個很多教會都做.
他們服侍長者.
又幫他們量血壓.
長者有空來教會坐可以聊天.
提供地方給他們喝茶.
沒問題.
他們服侍大學生.
那些大學生畢業之後.
因為那個地方不是一個很發達的城市.
在台中.
畢業之後就出城打工了.
他們九年在我去的時間.
之前的九年.
只有一個大學生.
願意留在這間教會服侍.
不是做童工.
是留在那裡.
你想想.
九年你服侍這些人.
回報率都很低.
如果我們從金融界經常說這些術語.
繼續做.
有一次我和他們去醫院.
打算去探病.
去探什麼呢.
那些重症.
其實每個人都沒有機會復原.

$^{681}$那些是瀕死的病人.
他們彈結他.
唱歌給他們聽.
希望他們能夠表達信耶穌.
就動一下手指.
動一下眼皮.
這些就肯定沒有回報的.
就是這些的事奉.
他們不問回報的.
他看到別人的需要.
就盡上自己的責任.
他們的服侍就是很著重傳福音.
讓人去信耶穌.
所以這一間令我很感動.
我今天也很記得.
也有一點塑造.
接著在舞會的事奉.
為他者.
一個為他者的教會.
一個為他者的信徒.
接著保羅很想盡快將以巴弗提送回菲律賓.
他有幾個目的.
在第27,28節.
我不讀了.
你待會兒自己看經文.
保羅可以減少一些擔憂.
可以鬆一口氣.
如果把以巴弗提送回去.
另外就是.
菲律賓信徒因為重遇以巴弗提.
會有喜樂.
第29節.
菲律賓信徒要好好接待和尊重以巴弗提.
為什麼保羅要菲律賓教會尊重以巴弗提呢?.
你想一想.
以巴弗提不是有兩個任務的嗎?.
你帶禮物,帶些饋贈給保羅.
他應該做得到的.
但是服侍保羅這件事是否達標呢?.
不達標.

$^{721}$你去到病了.
要保羅服侍你.
如果我們用香港今日的公司.
可能以巴弗提回到菲律賓教會.
辦事不力.
有負所託.
不達標.
可能會責備以巴弗提.
所以保羅叮囑菲律賓教會.
你們要在主裡歡歡樂樂地接待他.
而且要尊重這樣的人.
為什麼以巴弗提值得尊重呢?.
保羅也不想菲律賓教會.
你不要責備他.
反而他是值得尊重的.
我們再看最後一節.
因為他作基督的功夫幾乎至死.
不顧聖命.
要補足你們供給我的不及之處.
當然不及之處不是指.
菲律賓教會所做的不足夠.
就是帶不夠錢給保羅.
以巴弗提要自己用錢包去填補.
不是這個意思.
不及之處是說.
一件原本已經很美好的事情.
他沒有了這部分就不完整.
以巴弗提所做的.
就是讓一件原本已經很好的事情.
得以圓滿結束.
能夠得以完成.
但重點是.
以巴弗提真的有值得尊重的地方.
因為他做基督的工作不顧聖命.
我不知道大家的視頻是怎樣.
香港也不多做到不顧聖命.
我們只會被牧師抓到沒有聖命.
不好意思.
這是說陳牧師的風格.
不顧這個字在新約只出現了一次.

$^{761}$這個比較特別.
這個字當時是一個在賭場裡面的用語.
意思就是一個人將賭注放下去之後.
你拿不回來.
你壓了注.
你拿不回來.
不顧聖命就是你放下去了.
你拿不回來.
就是有些壯士一去不返.
以生命做賭注.
這樣的意味.
這種是全然擺上的服事.
我全力去服事.
在中國的教會裡面.
如果喜歡看傳記.
也有很多這些人.
其中有一個很有影響力的牧者.
宋尚哲是其中一個.
宋尚哲牧師臨死之前說什麼呢?.
他說我的為人不及王明道.
我的講道不及魏拓星.
我的交際不及繼志文.
這三個都是很出名的中國教會的牧者.
但是主耶穌知道.
我將一切都擺上.
所以宋尚哲的人生也影響了很多人.
他是不顧聖命的去擺上.
我又想起一個.
應該是在09年過世的一個牧師.
這個大家未必很熟悉.
這個是學園傳道會.
是韓國學園傳道會的創辦人.
以前的會長.
他叫金俊坤牧師.
金俊坤牧師的事奉.
真的是不顧聖命.
我在哪裡知道呢?.
就是在學傳的創辦人Bill Bright.
就是伯立德牧師的傳記裡面.
就看到金俊坤牧師的生命見證.

$^{801}$因為伯立德和金俊坤是非常緊密的同行戰友.
這兩個是經常交換茶經心得.
怎樣去禱告.
怎樣去開拓時宮.
怎樣去傳福音.
是一個很親密的同行戰友.
伯立德的傳記裡面.
記載了一段金俊坤.
他還沒有加入學傳之前的經歷.
當時在韓國.
在一個下雨的日子.
當時是很多很多年前.
當時的共產黨游擊隊.
突襲金俊坤所住的村落.
見人就殺.
金俊坤的太太和他爸爸當場被殺.
金俊坤被打到變成體淪相.
奄奄一息暈了.
可能人們以為他死了.
就走了.
當金俊坤清醒之後.
就帶著他的小女兒.
去到山區避開屠殺.
他是這個村莊裡面僅餘的兩個倖存者.
但是金俊坤是一位很認真的基督徒.
他記得勝經教他要愛自己的仇敵.
好像我上次講道.
好像講過這一段.
要愛你的仇敵.
要為你的仇敵.
即是迫害你的人去禱告.
他如何面對這些殺他太太.
殺他爸爸的敵人呢?.
神的靈感動他.
他要他向游擊隊的領袖.
親口告訴他.
我愛你.
耶穌也愛你.
耶穌也會饒恕你.
將神的大愛.

$^{841}$將福音傳給他.
我不知道.
如果給你.
第一.
給我我肯定做不到.
這個剛剛是沒有多久之前.
殺妻殺父仇人.
還有你去找這個游擊隊的領袖.
隨時連命都沒有.
真的不顧性命.
但是神的靈感動他.
他真的照這個感動.
去下山找這個游擊隊的領袖.
結果.
那個共產黨首領.
就跟金俊坤談了一段時間之後.
他信得住.
一起跪下禱告.
在很短時間之內.
很多.
即是這一隊人.
很多的這些共產黨員.
就歸順了耶穌.
金俊坤就幫他們建立了一間教會.
直到他自己.
後來金俊坤就去到美國.
深造.
然後就停了.
沒有再服侍這間教會.
所以白納德在他的傳記回顧的時候.
他說.
他跟金俊坤不斷的很多配搭.
同心同行.
分享讀經心得.
怎樣傳福音.
怎樣瞞訓.
怎樣跟各地教會去配搭組織.
共同去完成主所託付的大使命.
弟兄姊妹.
在聖經我們看到.

$^{881}$作公不離三兄弟.
保羅提摩泰爾巴伐提.
其實現代也有的.
白納德跟金俊坤也是這樣.
金俊坤是一個不顧性命的人.
我不知道你有沒有為著某一些的目標.
某一些工作.
某一些事工.
是不顧性命幾乎致死的呢.
其實我們香港每個人都有.
不過多數都是為工作.
我們為工作可以不顧性命.
犧牲健康.
捱更底夜.
有些壞一點的.
就是為英超.
通宵達旦.
三點鐘都看那些.
就是不顧健康.
弟兄姊妹.
我們是不是關心耶穌所關心呢?.
你有沒有曾經因為想服侍主.
而不顧性命呢?.
保羅提摩泰和爾巴伐提三個人.
他們的靈命不同.
恩賜不同.
甚至親疏都有別.
保羅跟提摩泰的關係是比較深的.
就是親如父子.
跟爾巴伐提就是同工關係比較多.
雖然是這麼多不同.
但一樣可以這樣.
去事奉配搭同心為主.
我想總結一下.
他們這個同心有什麼很值得我們.
去效法的一些心態呢?.
你效法到這個心態.
你就不會經常為了侍工的結果.
那些可以放輕一點.
關係才重要.

$^{921}$提摩泰追求耶穌基督的事.
他關心耶穌的事.
尤其關心傳福音.
爾巴伐提做基督的工作.
不顧性命.
完全是為他者的心腸.
所以如果你的人生能夠.
只關心耶穌所關心的事情.
你的人生多些是為他者.
放下自我.
這樣你就有些作功不利.
三兄弟的經歷.
保羅就跟他們同心同行.
是同行戰友.
保羅一直在支援這兩位弟兄.
教會呢.
我經常覺得.
教會不怕沒錢.
教會也不怕人很少.
教會更加不會怕地方很小.
這些坦白說不太重要.
教會最怕的是什麼呢?你問我.
最怕的就是不同心.
不同心會引起很多紛爭.
你有很多的心力資源.
去搞這些不同心的事情.
在事奉的路上.
有同心為主的人.
做你的同行戰友.
是非常重要的.
如何找到同行戰友呢?.
我鼓勵.
在教會裡面.
不要做旁觀者.
你要委身投入事奉.
你投入事奉.
你就有同心同行的經歷.
你就有機會找到同行的戰友.
你永遠是孤零零.
一個回來十點鐘崇拜.

$^{961}$崇拜完就馬上搭電梯下去走.
你永遠不會有.
同心同行的戰友.
願意神裝備我們.
都有保羅提摩太.
以巴弗提.
這種事奉的心腸.
又為我們預備同心的戰友.
如果你想想.
如果整個旺鎮.
大部分的人.
都有同心同行.
我們關心耶穌所關心的.
一個為他者的心腸.
旺鎮的未來.
真的你完全不用擔心.
我們同心禱告.
親愛的主.
你是我們的元帥.
我們靠著你.
才能夠彼此聯絡.
成為同心同行的戰友.
願意你祝福旺鎮.
成為一個.
可以常常都是為他者.
為旺鎮以外的人.
去存在的教會.
更加幫助弟兄姊妹.
讓我們多些關心.
你所關心的事情.
我們的人生很想去.
完成你所託付給我們的任務.
衷心去完成照命.
討你的喜悅.
求你加力.
求你祝福.
仰望你奉耶穌基督的聖名.
阿們.
\newpage



\section{腓立比書 3:1-16}
\label{sec:8t3Xeh3MCK8}
\textbf{2024年7月7日主餐崇拜直播｜陳明泉牧師｜靈性·關係·成長－甚麼偷走了我的快樂｜腓立比書3\_1-16}
\newline
\newline
連結: \href{https://youtube.com/watch?v=8t3Xeh3MCK8}{\texttt{https://youtube.com/watch?v=8t3Xeh3MCK8}} ~~~~ 語音日期: 2024-07-10
\newline
\newline
\hyperref[sec:l1Zxe_Q_Nro]{\small{< < < PREV SERMON < < <}}
~
\hyperref[sec:index_chronic]{\small{[返順時目]}}
~
\hyperref[sec:BxDFSB8mRZw]{\small{> > > NEXT SERMON > > >}}
\newline
\newline
腓立比書 3:1-16
\newline
\begin{longtable}{cl}
\hline
\hline
章節 & 經文 (和合本修訂版)\\
\hline
3:1 & \begin{tabularx}{0.7\textwidth}{X} 末了，我的弟兄們，你們要靠主喜樂。我把這些話再寫給你們，對我並不困難，對你們卻是妥當的。 \end{tabularx} \\ \\ \relax
3:2 & \begin{tabularx}{0.7\textwidth}{X} 應當防備犬類，防備作惡的，防備妄自行割的。 \end{tabularx} \\ \\ \relax
3:3 & \begin{tabularx}{0.7\textwidth}{X} 因為真受割禮的，就是我們這藉著神的靈敬拜、以基督耶穌為誇耀、不依靠肉體的。 \end{tabularx} \\ \\ \relax
3:4 & \begin{tabularx}{0.7\textwidth}{X} 其實，我也可以靠肉體；若是別人以為他可以依靠肉體，我更可以。 \end{tabularx} \\ \\ \relax
3:5 & \begin{tabularx}{0.7\textwidth}{X} 我出生後第八天受割禮；我是以色列族、便雅憫支派的人，是希伯來人所生的希伯來人。就律法說，我是法利賽人； \end{tabularx} \\ \\ \relax
3:6 & \begin{tabularx}{0.7\textwidth}{X} 就熱心說，我是迫害教會的；就律法上的義說，我是無可指責的。 \end{tabularx} \\ \\ \relax
3:7 & \begin{tabularx}{0.7\textwidth}{X} 只是我先前以為對我是有益的，我現在因基督的緣故而當作是有損的。 \end{tabularx} \\ \\ \relax
3:8 & \begin{tabularx}{0.7\textwidth}{X} 不但如此，我已把萬事當作是有損的，因我以認識我主基督耶穌為至寶。我為他已經丟棄萬事，看作糞土，為要贏得基督， \end{tabularx} \\ \\ \relax
3:9 & \begin{tabularx}{0.7\textwidth}{X} 並且得以在他裡面，不是有自己因律法而得的義，而是有信基督的義，就是基於信，從神而來的義， \end{tabularx} \\ \\ \relax
3:10 & \begin{tabularx}{0.7\textwidth}{X} 使我認識基督，知道他復活的大能，並且知道和他一同受苦，效法他的死， \end{tabularx} \\ \\ \relax
3:11 & \begin{tabularx}{0.7\textwidth}{X} 或許我也得以從死人中復活。 \end{tabularx} \\ \\ \relax
3:12 & \begin{tabularx}{0.7\textwidth}{X} 這不是說我已經得著了，已經完全了；而是竭力追求，或許可以得著基督耶穌所要我得著的。 \end{tabularx} \\ \\ \relax
3:13 & \begin{tabularx}{0.7\textwidth}{X} 弟兄們，我不是以為自己已經得著了；我只有一件事，就是忘記背後，努力面前的， \end{tabularx} \\ \\ \relax
3:14 & \begin{tabularx}{0.7\textwidth}{X} 向著標竿直跑，要得神在基督耶穌裡從上面召我來得的獎賞。 \end{tabularx} \\ \\ \relax
3:15 & \begin{tabularx}{0.7\textwidth}{X} 所以，我們中間凡是成熟的人，總要存這樣的心；若在甚麼事上存別樣的心，神也會把這些事指示你們。 \end{tabularx} \\ \\ \relax
3:16 & \begin{tabularx}{0.7\textwidth}{X} 然而，我們達到甚麼地步，就當照這個地步行。 \end{tabularx} \\ \\
[1ex]
\hline
\hline
\end{longtable}
$^{1}$《腓立比書》第3章1至16節.
第一節.
弟兄們,我還有話說.
你們要靠主喜樂.
我把這話再捨給你們.
於我並不為難.
於你們卻是妥當.
應當防備犬類.
防備作惡的.
防備妄自行各的.
因為真受割禮的.
乃是我們這以神的靈敬拜.
在基督耶穌裡誇口.
不靠著肉體的.
其實,我也可以靠肉體.
若是別人想他可以靠肉體.
我更可以靠著了.
我第八天受割禮.
我是以色列族.
變雅敏支派的人.
是希伯來人所生的希伯來人.
就律法說,我是法理賽人.
就熱心說,我是迫不教會的.
就律法上的二說,我是無可指責的.
只是,我先前以為與我有益的.
我現在因基督都當作有損的.
不但如此,我也將萬事當作有損的.
因我以認識我主基督耶穌為至寶.
我為他已經丟棄萬事,看作糞土,為要得著基督.
並且得意在他裡面.
不是有自己因律法而得的義.
乃是有信基督的義.
就是因信神而來的義.
是我認識基督,曉得他復活的大能.
並且曉得和他一同受苦,效法他的死.
或許我也得以從死裡復活.
這不是說我已經得著了,已經完全了.
我乃是竭力追求,或許可以得著基督耶穌所得著我的.
弟兄們,我不是以為自己已經得著了.
我只有一件事.

$^{41}$就是忘記背後努力面前的向著標竿直跑.
要得神在基督耶穌裡從上面照我來得的獎賞.
所以我們中間凡是完全人,總要存這樣的心.
若是在什麼事上存別樣的心.
神也必以此指示你們.
然而我們到了什麼地步,就當照著什麼地步行.
聖經讀完.
弟兄姊妹平安.
平安.
我們今天一起來用肥勒比書第三章.
我們一起來思考快樂,喜樂這個主題.
在肥勒比書裡,喜樂是一個重要的字眼.
出現了很多次.
在我們之前幾次的講道裡.
我們都有提及過這個主題.
以及如何實踐喜樂的一種態度和生命.
今天更聚焦到保羅.
他用自己作為一個例子.
如何實踐或經歷到從內而外的喜樂生活.
不過在我們進入經文之前.
大家問問自己.
你覺得自己是快樂的人嗎?.
你覺得自己正在經歷從信仰而來的快樂嗎?.
很多時候弟兄姊妹會說.
我都經歷了,為什麼呢?.
因為很多恩典,數算.
這是真的.
但有些弟兄姊妹會轉過頭來.
不是的,我有經歷過這些痛苦.
有這些病患.
人生有這些遺憾.
有時候我們見機.
或是看自己是否快樂.
我們用際遇來看自己的人生.
不過今天我們看這段經文.
正正就是保羅即使經歷了一個不太理想的際遇.
他都能夠活出一個喜樂的生命.
究竟他的秘訣是什麼呢?.
又或者我們和保羅都是有同一個信仰.
為什麼有時候基督徒信著信著.

$^{81}$沒有了那份喜樂.
但保羅可以保持住呢?.
希望今天我們在這段經文裡找到答案.
我們先祈禱求上帝幫助我們.
天父上帝願你在我們當中帶領我們.
讓我們能夠在你的話語裡.
我們讀得到我們生命的亮光.
指教我們,調教我們.
讓我們在你的話語裡.
我們可能有時候會想錯了.
或者我們有偏差.
求你的話語成為我們生命的準繩.
讓我們能夠聚焦在你的身上.
基督耶穌成為我們生命追求的目標.
讓我們在一瞬間我們明白你的話語的時候.
在生活裡實踐.
求你繼續來說領我們.
幫助每一位聆聽的弟姐妹聽得明白.
幫助孫講劇講得清楚.
堅守我們禱告奉靠基督耶穌得勝的聖名.
Amen.
我們一起看吧.
首先我們看第一個段落.
第三章第一至第三節.
在這裡說.
弟兄們,我還有話說,你們要靠主喜樂.
我把這話再寫給你們.
於我並不為難,於你們卻是妥當.
應當防備犬類,防備作惡的,防備妄自行覺的.
因為真受割禮的乃是我們這以神的靈敬拜.
在基督耶穌裡誇口,不靠著肉體的.
在這裡保羅就很斬釘截鐵地說.
他自己的生命是喜樂的.
他要將一個訊息告訴弟兄姊妹.
值得留意的是.
如何才算是靠主喜樂呢?.
在這裡我們並不清楚知道.
他不是說靠主喜樂有什麼步驟.
其實他在說靠主喜樂.
在經文裡留意一下.

$^{121}$其實他有當時在經文處境裡.
針對某些人.
而那些人令基督徒無法實踐喜樂的生活.
所以在這裡他有一種相對.
第一就是靠主得到喜樂.
第二個字眼就是靠肉體.
靠肉體是什麼人呢?.
他在經文裡有三個字眼.
重複又重複地告訴你.
你要謹慎防備.
防備這個字眼其實就是.
用廣東話來說就是看著.
盯著.
不是看著他來做.
那種看著就是要避開的.
類似說有個水桶掉下來.
看著那個水桶.
不是看著水桶你的頭撞過去.
被他打下來的.
看著就遠離.
這個防備就是這個意思.
看著一個很嚴厲的字眼.
就是在說什麼人呢?.
就是那些犬類.
作惡和妄自行惑.
其實三個字眼都是說同一批人.
就是當時在信仰裡.
有一些猶太主義者.
他們覺得信了耶穌都不夠.
他們還要加上舊約的一些律法.
加上那些宗教實踐.
這樣才能夠得救.
保羅說.
你要靠耶穌基督而得到喜樂.
而不是靠這些外在的宗教行為.
如果你靠這些外在宗教行為.
你的生命就達不到一種狀態.
就是那個喜樂和妥當的字眼.
在第一節那裡.
於我並不為難.

$^{161}$於你們卻是妥當.
什麼意思呢?.
妥當和合本詞好像OK的意思.
其實不是.
妥當的字眼本來就是安全保護的意思.
安全感.
你看著不要跟隨那些人.
你靠著基督.
你的生命才有一份安全感.
一份來自上帝的那種心靈的守衛.
你不看著上帝.
你不是靠上帝.
你靠著那些外在的宗教行為.
到最後你的心靈就會覺得還得還失.
之前我在課堂裡跟弟兄姊妹說過一次這樣的經歷.
話說我太太和我去到商場裡逛街.
在我進了廁所的時候.
有一個小妹妹問我太太.
問她什麼問題呢?.
就問她.
太太你是不是基督徒?.
類似是在說信仰的實踐.
她說我不明白一段聖經.
在《哥林多前書》裡說保羅說守的主餐不按道理.
按理守的主餐就是犯罪.
她就問我太太.
按理的道理是什麼意思?.
有時候忽然間考聖經問題.
雖然是師母.
但腦子有時候也有機會空轉.
在她在想的時候我就從廁所出來.
我太太就指指點點.
你問她吧.
她沒有說我是什麼身份.
你慢慢問她吧.
她就問我同樣的問題.
我說順民學裡就說到.
那些人分門別類.
她說不是的 哥哥.
她叫我哥哥也懂得做.

$^{201}$不是的 哥哥.
那段按理的意思是不按猶太人的律法來守主餐.
為什麼你會看到這件事.
她就立即揭開一大疊.
揭開了舊約.
聖經裡說逾越節晚餐怎麼守.
然後她就立即說.
聖經保羅說人們不守.
不按逾越節的規矩來守餐.
所以你們很多基督徒的生命一般般.
因為保羅也在罵你.
然後我聽完之後.
我這個哥哥就說.
你揭開這麼大疊.
但保羅順民夏里也不是在說那件事.
上面的順民夏里你不跟.
你揭開一大疊.
你就斷章取義來跟.
似乎不太妥當 妹妹.
我就說回去.
但在那裡似乎大家.
她好像說了也不明白.
然後我再說.
我們又不是猶太人.
為什麼要跟猶太人的規矩.
來守這個逾越節晚餐呢.
聖經裡說這個主餐和逾越節晚餐.
不是同樣對等的東西.
然後沒有聲音.
然後她知道搞不定這個.
打不定大佬.
於是找她的大姐來.
然後很快就立刻有個大姐走出來.
又說同樣的話.
不過找另外一個人來說.
然後我太太就說.
你罵她不贏的.
她是牧師.
然後我們兩個.
我說我很欣賞你們.

$^{241}$不過你們信錯了東西.
不是這樣看聖經的.
有機會來旺角浸信會.
這個好像是一個經歷.
不過其實這裡正正是說.
保羅時代裡面.
他們面對的處境.
在恩信清二的教義之外.
還要額外加添.
多一些宗教細節.
多一些猶太人的禮儀.
他們這群人沒有得救的把握.
沒有來自信仰裡面.
對他們的那份應許的信心.
反而在他們信之外.
再要額外加一些宗教行為.
他們心裡面才得救.
不過加這些外在行為.
可以這麼說.
是一個滾雪球的效應.
做了A之後你會覺得更加不夠.
你會更加再做B.
久而久之就製造了.
大量的宗教實踐.
在信仰裡面.
最純粹的那一面他不再說.
反而就疊加了很多.
那些宗教的外在行為.
禮儀等等那些東西.
在這裡特別說到.
那些人特別再去做那個國禮.
那些外邦人其實不需要再去做這個國禮.
但是保羅很針對性的說.
這些外邦基督徒.
反而要走回頭.
信猶太人的猶太主義的禮儀.
再行這個國禮.
想著這樣才成為神國子民.
保羅針對的這班人拿不到喜樂的原因.
就是他們疊加了很多.

$^{281}$在信仰以外的困難.
防備的意思就是.
在信仰實踐裡面.
會不會在我們人生當中.
加了很多那些人做的.
或者我們覺得是好的宗教實踐.
但是越做的時候.
越強調這些宗教外在規範.
製造更加多的禁忌.
只是令到你的心靈更加被束縛.
令到你自己覺得自己還得還實.
保羅在這裡說.
真正的敬拜不是在乎那些外在的規範.
你跟著什麼條件.
才叫做敬拜.
反而是用心靈誠實.
來敬拜我們的主.
所以在這裡他再繼續說.
敬拜神是靠上帝.
在基督裡面誇口.
直著聖靈.
聖靈以至到我們每一個都達到上帝的根前.
就不再是透過人為的力量.
透過我們自己外在宗教的行為.
來幫我們自己經歷這個信仰.
在這裡雖然在說.
過去保羅的時代.
有很多猶太主義的禮儀.
我們今天大概未必會這樣.
除非有一些所謂的異端.
他們可能會加了這樣的狀況.
不過在基督徒的實踐.
在我們的信仰.
我們即或.
我們今天很多時候.
我們自己都是純正的信仰.
但有意無意之間.
在我們的信仰實踐裡面.
我們可能在我們的生活當中.
加添了其他信仰以外的東西.

$^{321}$而將這些其他以外的東西.
成為了一個很重要的信仰元素.
例如我們要做某一些信仰的功夫.
做某一些功德.
我才覺得我得到上帝而來的喜樂.
又或者我們很擔心我們會做錯.
舉多一個例子就明白了.
有些弟兄姊妹.
家裡遇到家人有喪事,有白事.
很多弟兄姊妹問牧師.
應該要怎樣做.
牧師解釋了整個.
我們說我們黃浸.
沒有什麼繁文縟節.
很簡單的.
以紀念和回想.
紀念家人和安慰為主.
所以你說有一些神聖效力的禮儀.
我們沒有的.
沒有做這些事.
不過很多弟兄姊妹擔心.
他們覺得要做多一些事.
才能表達對先人的尊敬.
又或者不會在信仰裡面.
有一些抵觸.
做了一些傷害神人之間感情的事.
我經常說.
安慰不用擔心.
我們每一個基督徒.
安息禮拜其實不是為了.
以逝去的先人來做.
是為了地上的人.
我們教會一起努力.
就是為了這樣來安慰你.
陪你度過這段時間.
不用擔心.
有些事做了有什麼衝突.
違反一些忌諱.
不過很多弟兄姊妹心裡很混亂.
因為我們覺得我們要做些事.

$^{361}$我們才能平息上帝的憤怒.
在這裡保羅第一個提醒我們.
我們的信仰.
不要為我們製造太多.
不必要的繁文縟節.
在當時的猶太社會.
猶太人的禮儀.
成為了這些絆腳石.
去到今天.
在我們的信仰裡面.
在我們的文化當中.
我們都有很多這些絆腳石.
回歸最基本的.
因信基督.
我們心裡得平安.
這份安全感.
就是因為來自我們和神.
有一個第一身的交往.
單單這件事就足夠我們.
經歷一份喜樂的生活.
所以你的安全感.
你信仰裡面的實在.
不是透過外在行為來建立的.
是透過和神之間.
一份緊密的關係.
靠基督就是一個這樣的意思.
相信他.
你仍然仰望他.
你就可以得到這份從神而來的幫助.
經文裡面再繼續說.
第四節至第八節.
這是第二個段落.
這個段落是保羅用他自己的經歷.
來說.
我讀了一段經文.
我再引入多一點的內容給大家聽.
其實我也可以靠肉體.
若是別人想他可以靠肉體.
我更可以靠著了.
我第八天受割禮.

$^{401}$我是以色列族變眼敏之派的人.
是希伯來人所生的希伯來人.
就律法說我是法理賽人.
就以心說我是迫不及待的.
就律法上的易說.
我是無可指責的.
只是我先前以為與我有益.
我現在因基督都當作有損的.
不但如此.
我也將萬事當作有損的.
因為我以認識我主耶穌基督為至寶.
我為他已經凋棄萬事.
看作糞土.
為要得著基督.
在第四至第八節裡.
保羅就引入自己的經歷.
和自己的價值觀.
來作為一個展述.
要如何靠主起來.
在這裡給多一點色經上下文給大家知道.
在這裡所寫的整個保羅的筆調.
他雖然講自己.
他講自己過去好像很高.
如果引以為傲可以誇口.
我就是這樣的人.
一路講.
但一路講的時候.
由一個引以為傲覺得很.
某程度在當時人們說很厲害的一個履歷.
講這裡.
他越講就講到自己越來越低.
由高至低來描述自己.
甚至他將自己過去成為糞土.
過去不屑一顧.
覺得是廢物.
看舊時他覺得這麼高.
個個覺得很亮麗的背景.
他如今看為糞土一樣.
由高至低的寫作手法.
原來是在回應第二章裡面.

$^{441}$耶穌基督的描述.
你當以基督耶穌的心為心.
然後就講到耶穌坐在高天之上.
然後一路降低為人.
去到最謙卑的方式.
而且死在十字架上.
然後死而復生.
掌管萬有.
毫無不可.
口稱耶穌基督為主.
使榮耀歸於父神.
講的是耶穌的榜樣是由高至低.
再由低回到最高.
是一個V型反彈.
同樣保羅在這裡.
他用耶穌的生命經歷.
他用自己來回應福音.
耶穌的生命是這樣的經歷.
他都是用他自己的生命.
他開始明白耶穌的心腸.
耶穌基督的經歷.
由高至低.
甚至過去的日子.
他覺得很了不起.
很亮麗.
他到現在覺得不屑一顧.
在經文裡面有幾節出現一個字眼.
當作和漢作.
Hagelmai.
這是希臘文的字眼.
漢作和當作是同一個字.
剛才叫弟兄姊妹要防備.
那個又是漢.
不過今次漢作或當作.
不是看別人.
是看自己.
或者看自己回顧自己的價值觀念.
他覺得很重要的事.
因為相信基督.
現在他覺得那些不但幫不到他.

$^{481}$而且那些在信仰裡面是有損的.
因著基督.
顛覆了他過去裡面覺得重要和不重要.
以前逼迫教會的人.
如今就是極力去建立教會.
以前以他作為一個希伯來人.
當中的希伯來人的法理錯人為榮.
如今他就是講到.
那些實踐這些律法的人.
不屑一顧.
他用一個顛覆價值觀的觀念.
來看自己的生命.
有時我們人生不快樂.
或者信了主.
有時都經歷不到那種喜樂.
很大程度上.
我們信了主.
但我們沒有更新我們對人與事的價值觀念.
我們依然帶回舊時未信主的觀念.
放回我們信仰當中.
我們心裡昔日拜的偶像.
信了主後.
如果我們的價值觀不顛覆.
我們只是將偶像的標誌撕走.
放了耶穌的頭像.
我們覺得這樣就是信耶穌.
但這幫不到我們由內而外經歷上帝的豐富.
價值觀需要重新顛倒我們.
我們看甚麼為之有價值.
看甚麼為之不屑一顧.
當我們信了主的時候.
我們不要再帶著過往.
我們沾沾自喜.
覺得引以為傲的經歷.
成為我們生命中經歷不到上帝喜樂的半腳石.
我有這樣的經驗.
大家都知道.
小弟讀書不出色.
有一次家裡搬屋的時候.
我發了瘋去找.

$^{521}$家人問我找甚麼.
找了這麼久.
我想找我之前讀書不好的會考成績單.
家人問我為何要找這些.
家人覺得我有點無聊.
我說不是的.
我想找出來講見證.
家人問我講甚麼見證.
我只有一科合格.
回考只有一分.
十四分入中六.
我第一次考回考只有一分.
其他全部都是拿個U.
這麼噁心的東西還到處拿.
你不覺得很醜嗎.
其實真的挺醜的.
當時我記得拿了這張成績單後.
我第一樣反應是zar 皺了它.
我不敢撕掉它.
我沒有那種氣概.
我只是zar 皺了它.
zar 完一塊.
最後都是摺在袋子裡.
然後放在家裡的抽屜.
但為何我要拿呢.
我發覺這是我人生的經歷.
我正正是因為拿這一分.
上帝給我有恩典可以重新讀.
讀書不好不代表我人生.
永遠無法再行前.
這代表這一分.
回考這麼失敗的成績表.
比較我今天.
我仍然覺得自己還可以讀書.
我覺得是一個很大的落差.
一個反照.
我很想用這一張一分的成績單.
來勉勵身邊很多人.
所以我很努力去找.
不過對很多人來說.

$^{561}$那是不光彩的經歷.
但對我來說.
我過了它.
我覺得那張成績單是一個恩典的標記.
對我來說.
那不是什麼醜怪的經歷.
反而是我相信上帝.
因為我經歷信仰.
我的人生可以重新來過.
有蛻變的一種經歷.
價值觀改變.
過往覺得很棒的東西.
可能在保羅眼中視為糞土.
過往你覺得不堪入目的經歷.
因著上帝的更新.
那個成為了你講見證.
成為你生命裡面一個很重要的一頁.
我們求主幫助我們.
如果我生命裡面經歷到這一份價值觀.
或者我生命裡面優先次序.
重要不重要的排序重新去看的時候.
我們會發覺我們人生豁然開朗.
相反來說.
你只是藉著信仰來強化了過往.
你活在這個世界的價值觀念.
你沒辦法在信仰裡面拿到養分.
因為你仍然繼續帶著這個世界的角度來過你的生命.
保羅說過去所有的事對我來說一切都是有損的.
我現在因基督耶穌當作有損.
我又將萬事當作有損.
因我以認識我主耶穌基督為至寶.
我為祂已經凋棄萬事.
看作糞土為要得著基督.
得著基督就是要將你過往的價值觀捨棄.
丟下.
重新經歷一個從神而來.
讓你懂得想對的方法來過你的生活.
我們看第三個段落.
九節到十一節.
並且得意在他裡面.

$^{601}$不是有自己因律法而得的義.
乃是有信基督的義.
就是因信神而來的義.
使我認識基督.
曉得他復活的大能.
並且曉得和他一同受苦.
效法他的死.
或者我也得以從死裡復活.
在這段小段落裡.
首先要解釋一下第十一節.
或者我也得以從死裡復活.
看上去跟你上文下理好像很奇怪.
剛才說到因信基督.
完全倚靠基督.
他就得到死而復生的生命態度.
但為何第十一節說到無伶兩可.
或者得呢.
其實不同的色經學家有不同的說法.
不過在這裡.
最多色經學家同意的就是.
或者我也得以從死裡復活.
是一個很謙卑的說法.
他不是否定因信清義.
或者信的人在基督裡死了.
也必復活這個教義.
只不過是他在信仰實踐當中.
他在生命的復活.
或者將來的日子.
和基督一同經歷那種豐富.
他不視為老奉.
他反而是主要保守我的信心.
持守我的生命到底.
有一天我也帶著這份信心離開這個世界.
也會因著上帝的應許.
重新經歷一個復活的生命.
所以保羅在這裡.
他是篤信上帝的工作.
但他不信的.
或者他不敢篤定的.
就是自己有多麼厲害.

$^{641}$自己的信心有多麼堅強.
面對逼迫的日子.
有時也是逆流而上的日子.
他只是帶著一份最大的謙卑.
來講主若許可.
上帝的美意是這樣的時候.
我也可以經歷一種從死裡復活.
當你知道這是一個很謙卑的說法的時候.
然後再回到前面的說法.
說的是一個因律法得而.
和因信基督而來的義.
是兩種的義.
舊時的他.
他是一個猶太人.
一個法利賽人.
強調的是固守所有的律法.
做好所有的律法的規矩和規定.
按照律法來說.
頂呱呱是好事.
不過就是靠人力做好所有的事.
他反而失落了.
沒辦法經歷上帝而來的豐富.
但因信基督的義.
那個義是因為.
單單是我們每一個相信.
因信的正義.
不是我們說我們有多麼的了不起.
我們有多麼好的能力.
只不過是一份單純的信靠.
以致他有一份這樣的能力.
保羅更加相信.
這個因信的正義.
不是說一個將來式.
將來他復活得生的.
一個很抽象的教義.
他更加說到.
這種復活的能力.
在今天他已經經歷得到了.
這種復活的能力就是.
面對著生活有很多的困難.

$^{681}$有很多的信仰迫不得.
但他都有能力去面對.
今天很多時候我們信仰裡面.
我們沒辦法經歷快樂.
原因是我們經常覺得自己要做些事.
我們覺得人力是做到很多.
改變到很多的事.
我們用自己的方法來操盤.
我們靠我們自己.
靠我們的智慧.
靠我們自己的能力.
我們覺得那件事會成功.
到最後.
當結果和我們的努力不相配的時候.
我們會失望.
我們會質疑.
我們會灰心.
保羅在這裡裡面.
他知道他的能力感.
在他的信仰實踐裡面.
能夠面對這種逆流.
能力感不是來自.
自己有多麼強的堅強意志.
而是來自上帝復活的大能.
來自上帝在他生命當中的介入.
將這種復活的能力.
讓他的內心改喚一生.
這種體會自己無能.
但經歷了上帝的大能的心態.
成為了他實踐信仰一個很重要的.
一個里程碑.
一個幫助他走面前的路.
很多時候今天坊間裡面.
有些正向心理學.
說一個能力感.
我做到的.
我自己做得到的.
這個都是正面的.
好的.
不過人越大.

$^{721}$你更加明白.
人再多努力也好.
我們人生沒有辦法做盡所有的能力.
改變所有的問題.
更加多.
其實都是讓乃上帝的恩典.
人生裡面.
我們體會自己的無能為力.
我們更加發現上帝復活的大能.
這個就是成為了我們生命裡面.
可以圍繞上帝一個很重要的經歷.
我經常想起.
有一段早幾年的經歷.
其實都是幫助.
人難免有低潮.
我想起那段經歷.
我心裡面就有力量.
大家記得我媽媽有一段時間.
在患癌症離開之前.
有一段都是很.
很貧破.
很憂心的日子.
我記得那一年.
2020年剛好是新冠肺炎的時候.
家人一起出去吃完一餐.
父親節飯之後.
媽媽第二天就開始發燒.
身體的情況急轉直下.
我沒有特別跟大家說.
其實我媽媽進了醫院之後.
第二天.
進了私家醫院.
醫生說你要有心理準備.
然後他說肺的呼吸量很低.
馬上又要戴一個很大的氧氣罩.
用到最高濃度的氧氣.
來把氣泵放進去.
讓他可以維持生命.
醫生給我的建議就是.
你們要有心理準備.

$^{761}$第二件事他跟我說的就是.
現在的設備在私家醫院是很貴的.
你們有沒有想過.
媽媽是急.
但是有沒有想過.
如果你的經濟沒有辦法實踐的時候.
可能你要轉去公立醫院.
這個很實在.
當我聽到這個消息的時候.
擔心.
半夜的時候醫生跟我說這番話.
第一擔心媽媽的狀況.
第二擔心我有沒有辦法找數.
醫療費是怎樣.
半夜又不知道怎樣去做.
祈禱 上帝我不懂得做.
不懂得怎樣去做.
可不可以幫我.
躲在浸會醫院的角落.
在車裡面忍泣和哭.
半夜三更覺得特別淒清.
不過一直在祈禱的時候.
其實不會有天上的聲音.
你放心.
只不過我知道這件事不是靠自己.
不是靠人的能力.
仰望上帝.
讓媽媽在醫院裡面.
放心地去醫治.
一個星期之後.
媽媽的氧氣罩就拔掉了.
醫生說她的肺奇蹟地恢復得非常好.
然後一路一路.
身體有一個比較好的狀態.
走完她人生最後兩個多月的日子.
雖然媽媽都離去.
回到天家.
但這段經歷對我來說很重要.
上帝沒有即時讓媽媽回到天家.
沒有讓我們變得很遺憾.

$^{801}$吃完飯 媽媽就離開了.
反而是在COVID的日子.
我們家庭 幾兄妹 所有家人.
有足夠的時間陪伴媽媽.
很多時間聊天.
很多時間讓她接受我們的關懷.
同時也有一些共處的日子.
現在回想這兩個多月.
是上帝一種能力在我們生命裡面呈現.
雖然從結果中媽媽都會離開.
但在心靈當中.
我知道上帝特別給了我們一段時間.
跟家人相聚.
是一個很寶貴的一段時間.
這段經歷.
當我人生有低潮的時候.
上帝是這樣幫我的.
所以就不需要擔心.
人生裡總會難免有一些.
你自己處理不到的難題.
回想上帝那種人力無法控制.
無法操控的狀況.
你第一生經歷上帝的幫助.
你知道上帝願你同在.
你就從這些經歷 回憶當中.
支取你的力量.
認定篤信上帝願你走過人生高高低低的經歷.
生命裡面可能都會經歷生老病死離別.
但當我們知道生老病死這種離別.
不是無情的.
好像關燈一樣.
當中有很多恩典來對偏執而來的時候.
心裡面就踏實安穩.
心裡面就有一份確信.
你知道這個世界不是你自己孤身作戰.
不是靠你自己的力量來面對一切.
保羅的能力感來自他一直傳福音.
雖然經歷逼迫.
雖然經歷患難.
但他知道他生命是上帝所看著.

$^{841}$上帝與他同在.
用這些經歷構建他生命那份能力感.
最後一個段落.
十二至十六節.
其實這幾節可以足以講一篇講章.
不過今天我快講.
我們整個脈絡去看.
「這不是說我已經得著了,已經完全了.
我乃是竭力追求.
或者可以得著.
基督耶穌可以得著我的.
弟兄們,我不是以為自己已經得著了.
我只有一件事.
就是忘記背後努力面前.
向著標竿直跑.
要得神在基督耶穌裡.
從上面召我來得賞賜.
所以我們中間.
凡是完全人總要存這樣的心.
若在什麼事上存別樣的心.
神也必以此指示你們.
然而我們到了什麼地步.
就照著什麼地步行」.
在這裡某程度上.
保羅用一個比較勵志.
用一個很前進很動態的方式.
講到他的信仰實踐.
他生命裡有喜樂的原因.
他認定.
他說自己不是已經得著了.
這是一個很謙卑的說法.
他信了耶穌.
他得著救恩.
但在他心中.
他不是走到信仰的終點.
他只不過是繼續朝著上帝為他所預備的終點.
繼續努力.
所以那個標竿才是人生的終點.
他當下站在這裡.
他在教育書上經歷了二十多年的傳道日子.

$^{881}$他在基督教界是大佬.
但他也不是用一個強人之態.
去講他自己的信仰生命.
他不覺得自己已經是一個最高典範.
他只說.
我今天站在的不是一個終點.
前面有一個標竿在等著我.
我就踏實朝著那個標竿.
來走面前這個信仰勁跑的路.
在這裡保羅特別提了幾個字.
第一就是要忘記背後.
忘記什麼呢?.
今天我們讀這段經文.
我們說忘記背後.
忘記一些遺憾.
忘記一些不開心.
不過剛才你看到那個經文的脈絡.
可能更加多的忘記.
是在說昔日令他沾沾自喜.
他覺得理想的生活狀態.
那些昔日可以誇口的經歷.
保羅說我忘記它.
與其說忘記.
倒不如不要用那個來比對今天.
昔日多麼彪炳也好.
人生都是向前行.
走向前面的標竿.
過去不要成為今天的難所.
我們可能有些成熟一點的.
年紀的靈子們有時會慨嘆.
可能今天的日子相對於年輕.
可能已經世界或自己.
很多東西已經改變.
有時我們有份唏噓的感覺.
在保羅的信仰實踐裡沒有唏噓.
他知道身體會軟弱.
人會老.
但他的心靈是可以繼續.
還有很多力量可以向前奔跑.
不過奔跑的方式不同.

$^{921}$但不代表今天的你停止了進步.
忘記背後那些令你沾沾自喜的經歷.
活在當下我還可以努力的地方.
用不同的方式朝著標竿去跑.
無論什麼年紀.
無論什麼狀態.
我們生命都可以有進步.
所以保羅提醒我們.
他也提醒自己.
忘記背後 努力面前 向著標竿直跑.
第二件事他更要說的是.
跑步的競賽者的對手是誰.
這是跑步競技場的用語.
向著標竿直跑 跑馬拉松.
不過真正要比較的不是看旁邊的跑得如何.
而是和自己比較.
如果和別人比較.
刻意和身邊的人較勁.
你會很失落.
你會很大壓力.
你亦不覺得自己做到心目中想要朝著的標竿.
因為你的眼不是看著標竿.
而是看著旁邊的.
你越跑就越覺得費勁.
和自己的過去做一些比較.
昨日的你可能是這樣.
明天的你會更進一步.
那一步未必很大.
一小步一進步都是好的.
用你自己的過去.
成為你自己生命的對手.
用上帝給我們的經歷.
用我們的成熟度來承托我們每一步的步伐.
這是保羅給我們最重要的經歷.
保羅看掌相是甚麼呢?.
在當時來說.
馬拉松的掌相是甚麼呢?.
對今天來說.
覺得是貴官.
找些植物積黃冠.

$^{961}$但不是用金做的.
某程度上你會覺得官面有點膽怯.
會借的.
不過重點不是拿那個官.
重點是你進步.
你朝著那個標竿一切的努力.
最終的掌相其實不是最重要.
在過程中你所得到的才是最重要.
今天這段經文很豐富.
不過今天有甚麼偷走了我們的快樂呢?.
可能我們將安全感放在一些外在行為裡.
我們經歷不到那份喜樂.
將你的安全感放在基督裡.
這是第一個今天這段經文的提醒.
第二個.
你的生命的喜樂.
可能延續你以前世界的價值觀念.
所以你心裡依然無法體會信仰的寶貴.
經歷不到那種徹頭徹尾的豐盛.
重新顛覆一下你自己.
今天你信主的原因.
你看甚麼為重要.
重新調教你的價值觀念.
你會發現生命很多苦惱.
來自錯誤的價值觀念.
求主幫助我們.
第二個是自我調教我們的價值觀.
第三我們的能力感.
我們的能力感不是基於我們自己有多大的能力.
而是我們知道我們人越來越成熟的時候.
我們知道我們的人生不是靠我們肉體的力量.
來克服眼前所有困難.
上帝成為我們生命裡面最大的能力.
上帝介入在我們生命當中.
我們人生不是自負盈虧的方式來過我們的生活.
上帝介入在我們生命.
數算主因.
你會發覺原來上帝一直在幫助我們.
最後我們的信仰要有目的感.
我們要有目標.

$^{1001}$無論甚麼年紀,甚麼狀態.
我們總是可以比昨日更加進步,更加成熟.
今日這篇信息成為一個很重要為我們打氣的信息.
為了盼望我們整個教會.
我們一眾弟兄姊妹在基督裡面.
我們得著真正的喜樂.
\newpage



\section{約翰福音 5:18-47}
\label{sec:BxDFSB8mRZw}
\textbf{2024年7月21日崇拜直播｜梁家麟牧師｜與父上帝平等的耶穌｜約翰福音5\_18-47}
\newline
\newline
連結: \href{https://youtube.com/watch?v=BxDFSB8mRZw}{\texttt{https://youtube.com/watch?v=BxDFSB8mRZw}} ~~~~ 語音日期: 2024-07-23
\newline
\newline
\hyperref[sec:8t3Xeh3MCK8]{\small{< < < PREV SERMON < < <}}
~
\hyperref[sec:index_chronic]{\small{[返順時目]}}
~
\hyperref[sec:rzRD6n7AW0w]{\small{> > > NEXT SERMON > > >}}
\newline
\newline
約翰福音 5:18-47
\newline
\begin{longtable}{cl}
\hline
\hline
章節 & 經文 (和合本修訂版)\\
\hline
5:18 & \begin{tabularx}{0.7\textwidth}{X} 為了這緣故，猶太人越發想要殺他，因為他不但犯了安息日，而且稱神為他的父，把自己和神看為同等。 \end{tabularx} \\ \\ \relax
5:19 & \begin{tabularx}{0.7\textwidth}{X} 於是耶穌回答，對他們說：「我實實在在地告訴你們，子憑著自己不能做甚麼，惟有看見父所做的，他才做；父所做的事，子也照樣做。 \end{tabularx} \\ \\ \relax
5:20 & \begin{tabularx}{0.7\textwidth}{X} 父愛子，將自己所做的一切事指示給他看，還要將比這更大的事給他看，使你們驚訝。 \end{tabularx} \\ \\ \relax
5:21 & \begin{tabularx}{0.7\textwidth}{X} 父怎樣叫死人復活，賜他們生命，子也照樣隨自己的意願賜人生命。 \end{tabularx} \\ \\ \relax
5:22 & \begin{tabularx}{0.7\textwidth}{X} 父不審判任何人，而是把審判的事全交給子， \end{tabularx} \\ \\ \relax
5:23 & \begin{tabularx}{0.7\textwidth}{X} 為要使人都尊敬子，如同尊敬父一樣。不尊敬子的，就是不尊敬差子來的父。 \end{tabularx} \\ \\ \relax
5:24 & \begin{tabularx}{0.7\textwidth}{X} 「我實實在在地告訴你們，那聽我話又信差我來那位的，就有永生，不至於被定罪，而是已經出死入生了。 \end{tabularx} \\ \\ \relax
5:25 & \begin{tabularx}{0.7\textwidth}{X} 我實實在在地告訴你們，時候將到，現在就是了，死人要聽見神兒子的聲音，聽見的人就要活了。 \end{tabularx} \\ \\ \relax
5:26 & \begin{tabularx}{0.7\textwidth}{X} 因為父怎樣自己裡面有生命，也照樣賜給他兒子自己裡面有生命， \end{tabularx} \\ \\ \relax
5:27 & \begin{tabularx}{0.7\textwidth}{X} 並且賜給他施行審判的權柄，因為他是人子。 \end{tabularx} \\ \\ \relax
5:28 & \begin{tabularx}{0.7\textwidth}{X} 你們不要對這事感到驚訝，因為時候將到，凡在墳墓裡的，都要聽見他的聲音， \end{tabularx} \\ \\ \relax
5:29 & \begin{tabularx}{0.7\textwidth}{X} 並且要出來：行善的，復活得生命；作惡的，復活被定罪。 \end{tabularx} \\ \\ \relax
5:30 & \begin{tabularx}{0.7\textwidth}{X} 「我憑著自己不能做甚麼。我怎麼聽見就怎麼審判，而我的審判是公平的，因為我不尋求自己的意願，只尋求差我來那位的旨意。」 \end{tabularx} \\ \\ \relax
5:31 & \begin{tabularx}{0.7\textwidth}{X} 「我若為自己作見證，我的見證就不真。 \end{tabularx} \\ \\ \relax
5:32 & \begin{tabularx}{0.7\textwidth}{X} 另有一位為我作見證，我也知道他為我作的見證是真的。 \end{tabularx} \\ \\ \relax
5:33 & \begin{tabularx}{0.7\textwidth}{X} 你們曾差人到約翰那裡，他為真理作過見證。 \end{tabularx} \\ \\ \relax
5:34 & \begin{tabularx}{0.7\textwidth}{X} 其實，我所受的見證不是從人來的；然而，我說這些話是為了使你們得救。 \end{tabularx} \\ \\ \relax
5:35 & \begin{tabularx}{0.7\textwidth}{X} 約翰是點亮的明燈，你們情願因他的光歡欣一時。 \end{tabularx} \\ \\ \relax
5:36 & \begin{tabularx}{0.7\textwidth}{X} 但我有比約翰更大的見證：父交給我去完成的工作，就是我正在做的，為我作證是父差遣了我。 \end{tabularx} \\ \\ \relax
5:37 & \begin{tabularx}{0.7\textwidth}{X} 那差我來的父也為我作了見證。你們從來沒有聽見他的聲音，也沒有看見他的形像。 \end{tabularx} \\ \\ \relax
5:38 & \begin{tabularx}{0.7\textwidth}{X} 你們並沒有他的道存在心裡，因為你們不信他所差來的那一位。 \end{tabularx} \\ \\ \relax
5:39 & \begin{tabularx}{0.7\textwidth}{X} 你們查考聖經，因你們以為其中有永生；而這經正是為我作見證的。 \end{tabularx} \\ \\ \relax
5:40 & \begin{tabularx}{0.7\textwidth}{X} 然而，你們不肯到我這裡來得生命。 \end{tabularx} \\ \\ \relax
5:41 & \begin{tabularx}{0.7\textwidth}{X} 「我不接受從人來的榮耀， \end{tabularx} \\ \\ \relax
5:42 & \begin{tabularx}{0.7\textwidth}{X} 但我知道，你們沒有愛神的心。 \end{tabularx} \\ \\ \relax
5:43 & \begin{tabularx}{0.7\textwidth}{X} 我奉我父的名來了，你們並不接納我；若有別人奉自己的名來，你們倒會接納他。 \end{tabularx} \\ \\ \relax
5:44 & \begin{tabularx}{0.7\textwidth}{X} 你們互相受榮耀，卻不尋求從獨一神來的榮耀，怎能信我呢？ \end{tabularx} \\ \\ \relax
5:45 & \begin{tabularx}{0.7\textwidth}{X} 不要以為我會在父面前告你們；有一位告你們的，就是你們所仰望的摩西。 \end{tabularx} \\ \\ \relax
5:46 & \begin{tabularx}{0.7\textwidth}{X} 如果你們信摩西，也會信我，因為他寫過關於我的事。 \end{tabularx} \\ \\ \relax
5:47 & \begin{tabularx}{0.7\textwidth}{X} 你們若不信他的書，怎能信我的話呢？」 \end{tabularx} \\ \\
[1ex]
\hline
\hline
\end{longtable}
$^{1}$今日讀經是在約翰福音五章十九至三十節.
在新約聖經一百三十一至一百三十二頁.
請聽我讀出第十八節.
「所以猶太人越發想要殺他.
因他不但犯了安息日.
並且稱神為他的父.
將自己和神當作平等.
耶穌對他們說:.
「我實實在在的告訴你們.
只憑著自己不能作甚麼.
唯有看見父所作的,子才能作.
父所作的事,子也照樣作.
父愛子將自己所作的一切事.
只給他看,還要將比這更大的事.
只給他看,叫你們稀奇.
父怎樣叫死人起來.
使他們活著.
子也照樣隨自己的意思使人活著.
父不審判甚麼人.
乃將審判的事傳教於子.
叫人都尊敬子.
如同尊敬父一樣.
不尊敬子的,就是不尊敬差子來的父.
我實實在在的告訴你們.
那聽我說:.
「有送差我來者的,就有永生.
不至於定罪,是已經出死入心了.
我實實在在的告訴你們.
時候將到,現在就是了.
死人要聽見神兒子的聲音.
聽見的人就要活了.
因為父怎樣在自己有生命.
就賜給他兒子也照樣在自己有生命.
並且因為他是人子.
就賜給他行審判的權柄.
你們不要把這事看作稀奇.
時候要到,凡在墳墓裡的.
都要聽見他的聲音,就出來.
行善的,復活得新.
作惡的,復活定罪.

$^{41}$我憑著自己不能作甚麼.
我怎麼聽見,就怎麼審判.
我的審判也是公平的.
因為我不求自己的意思.
只求那差我來者的意思」.
聖經獨白.
(主持人:聽見會早晨).
(主持人:少少說明今天講的道).
(主持人:第一要講的就是今天篇道會有多悶).
(主持人:因為講教義的東西是很悶的).
(主持人:請大家先有心理準備,有受苦心志).
(主持人:這是第一樣東西).
(主持人:第二樣東西就是).
(主持人:如果心水清一點).
(主持人:你會見到我其實skip了一段經文).
(主持人:五章,其實第一段的經文).
(主持人:是耶穌在畢氏大廁子醫好那個).
(主持人:據說好像是癱瘓的人).
(主持人:那段經文我以前講過的).
(主持人:並且講得也算不錯的).
(主持人:又將它寫出來出了書的).
(主持人:叫做改變人生命的耶穌).
(主持人:如果有人看過的話).
(主持人:或者沒有的話,再去找這本書來看).
(主持人:因為我自己寫完之後).
(主持人:我自己也沒有辦法).
(主持人:我想了很久也無法突破我以前寫的觀點).
(主持人:我跟陳牧師問).
(主持人:我將編出的講章再讀一次好不好玩).
(主持人:他說不是太好).
(主持人:所以我一跳就跳了它).
(主持人:所以就沒有了那一段).
(主持人:不過如果大家要refer那段經文).
(主持人:很容易找得到的).
(主持人:改變人生命的耶穌,那一段).
(主持人:於是就來到這一段).
(主持人:五章十八節).
(主持人:其實整段經文是很長的).
(主持人:你看到我讀是讀三十字).
(主持人:但是那裡是漏了一大段).

$^{81}$(主持人:我們會一看就看到四十七節).
(主持人:就是由十八節一路跳到四十七節).
(主持人:是很長的一段經文).
(主持人:本來我可以分兩段來講).
(主持人:不過這麼悶的東西一次講完就好了).
(主持人:長痛不如短痛).
(主持人:所以就短短地講完它).
(主持人:因為太複雜太長).
(主持人:所以我就很罕有地用powerpoint).
(主持人:麻煩幾位弟兄同工幫我們搞powerpoint).
(主持人:先聲明一下).
(主持人:這個powerpoint是我太太幫我做的).
(主持人:所以要給她一個credit).
(主持人:好,這一個大段).
(主持人:五章十八節一路到四十七節).
(主持人:記載了耶穌一篇非常長篇的講道).
(主持人:也是在約翰福音裡面).
(主持人:繼重生和活水之後).
(主持人:耶穌第三個長篇的教導).
(主持人:表面上看這段經文沒什麼特別).
(主持人:耶穌說話有點重重複複).
(主持人:又沒什麼比較醒目的詞彙表達).
(主持人:諸如我是生命的量).
(主持人:我是世上的光).
(主持人:所以這段經文).
(主持人:我們很多時候自己看).
(主持人:都會看完就算了).
(主持人:是不是?).
(主持人:看就看過了).
(主持人:不過沒什麼深的印象留下).
(主持人:但是其實這段經文).
(主持人:是耶穌基督非常重要的講論).
(主持人:跟其他耶穌在約翰福音的講論比較).
(主持人:他的重要性只是有過之而無不及).
(主持人:為什麼這段教導這麼重要).
(主持人:因為在這裡).
(主持人:耶穌是鄭重的有系統的全面的).
(主持人:講述他的神聖身份).
(主持人:他跟父的地位和關係).
(主持人:他被差遣到人間).

$^{121}$(主持人:要完成彌賽亞的使命).
(主持人:簡單說).
(主持人:耶穌基督在這裡做了一個).
(主持人:我是上帝的神聖宣告).
(主持人:雖然你揭破整個約翰福音).
(主持人:乃至連同婦女福音四卷福音書).
(主持人:在聖經裡面).
(主持人:耶穌在人間從來沒有直接說過).
(主持人:我是上帝).
(主持人:沒有的).
(主持人:你找不到一句說話).
(主持人:耶穌說我是上帝).
(主持人:不過他不喜歡這麼噁心).
(主持人:這麼明顯地說我是上帝).
(主持人:但他卻宣告).
(主持人:他跟你們信的上帝是同等的).
(主持人:我跟父上帝同等).
(主持人:這個說話的意思).
(主持人:其實等於說他是上帝).
(主持人:對不對).
(主持人:我跟他同等).
(主持人:他這麼大我這麼大).
(主持人:他是怎樣我是怎樣).
(主持人:所以這是一個極其霸道的宣告).
(主持人:我跟父平等).
(主持人:我來自上帝).
(主持人:父上帝的旨意就是我的旨意).
(主持人:我的旨意就是父上帝的旨意).
(主持人:他是這麼說的).
(主持人:必須要說).
(主持人:耶穌這樣做宣告是非常不容易的).
(主持人:如果他是突然之間).
(主持人:好像外星人那樣降落在加利利).
(主持人:那些人不知道他是誰).
(主持人:他這麼說).
(主持人:就是裝神騙鬼).
(主持人:還可以騙到一些如夫如婦).
(主持人:但加利利是耶穌出生成長生活的地方).
(主持人:這裡不少人是認識他的過去).
(主持人:為什麼他突然宣告他是上帝呢).

$^{161}$(主持人:就好像樹根今天突然在我們中間說).
(主持人:我告訴你一個秘密).
(主持人:我是上帝的化身).
(主持人:你們會有什麼反應呢).
(主持人:玩弄嗎?要玩到這麼大嗎?).
(主持人:要玩到這麼盡嗎?).
(主持人:英雄關鍵亦常人).
(主持人:我不知道你的背景底細).
(主持人:就被你騙而已).
(主持人:離譜成這樣).
(主持人:你真的好像在冒犯我的智商).
(主持人:約翰福音四章四十四節).
(主持人:耶穌自己慨嘆).
先知在本地是沒有人尊敬的.
這一點也不奇怪.
在旺朕多年弟兄姊妹姐妹.
也是沒有人尊敬的.
你們要有心理準備.
大家都知道大家是誰.
尊什麼敬呢?.
但你們要知道.
要宣告自己得到上帝的啟示.
宣告自己被上帝委任做先知.
擁有神聖的使命.
比宣告自己是上帝要合理得多.
容易被人接受得多.
雖然我和你很熟悉.
我知道你平凡的過去.
我也無法排除.
上帝昨晚可能向你報夢.
給你一個特殊的知識.
特殊的能力.
讓你水鬼升城王.
一步登天.
一招得志.
是不是可以?.
是可以的.
但如果你說你是上帝.
一直都是.
向來是.

$^{201}$我就連查證也不想.
不信就當你傻了.
我信你就是我傻了.
每個人都會這樣想的.
所以當耶穌在眾人面前.
鄭重宣告他的神聖身份.
那些敬奉上帝的猶太人就想殺他.
因為他的宣告是狂妄而具洩毒性.
五章十八節是這樣說的.
所以猶太人越發想要殺他.
因他不但犯了安息日.
並且稱上帝為他德父.
將自己和上帝當作平等.
所以耶穌這個宣告是什麼意思?.
如果猶太人想殺他.
就證明他其實是在做一個.
我是上帝的宣告.
我和父平等.
他這麼大,我這麼大.
犯安息日是小事.
自稱和上帝平等.
是嚴重的褻瀆神明罪.
這篇長篇的講論.
可以分為三個信息段落.
第一是耶穌和父上帝的關係.
第二是耶穌擁有審判世人的權柄.
第三是司一若翰和舊約聖經.
是耶穌身份的證明.
先講第一段.
耶穌和父上帝的密切關係.
我再講.
耶穌在人間從來沒有直接說自己是上帝.
但他卻多次強調.
猶太人所拜的上帝是他的天父.
而他和天父本質相同,地位平等.
耶穌首先講在十九節.
我實實在在告訴你們.
只憑著自己不能作什麼.
唯有看見父所作的,只才能作.
父所作的事,只也照樣作.

$^{241}$耶穌說他沒有自作主張做自己喜歡的事.
他所做的都是天父已經做的.
都是他看到天父做,然後自己做的.
看著眼,這句話好像在說.
耶穌沒有自己的主見,甚至沒有自己的意志.
他只是copycat(複製).
照抄天父所做的.
他只是工具,是打手.
天父要他做什麼,他就做什麼.
好像說得他很低,其實不是.
事實上,耶穌強調的是什麼?.
是他和天父的一致性.
和相依性.
兩者完全一致.
就是說,你不要企圖區分.
耶穌和天父.
不要用天父的旨意來挑戰耶穌的主張.
喂,耶穌你這樣做,和上帝的教導不同.
你違反上帝的心意.
耶穌說,他不僅在事實上沒有違反天父的旨意.
更加是原則上不可能違反天父的旨意.
因為耶穌的旨意就是天父的旨意.
兩者是完全相同的.
為什麼原則上不可以有差異?.
因為天父和耶穌本性相同.
所以當你看到耶穌的想法和做法.
就知道這個同時就是天父的想法和做法.
耶穌的教導就是上帝的教導.
因此是終極性的真理.
20-21節.
父愛子將自己所作的一切事子給他看.
問要將比這更大的事子給他看.
叫你們不稀奇.
父怎樣叫死人起來,使他們活著.
只是照樣隨自己的意思使人活著.
從來沒有人見過上帝.
只有父懷裡的獨生子將祂表明出來.
耶穌基督是隱藏的父上帝的外顯.
我以前經常說.
天父如果是一個隱藏的上帝.

$^{281}$耶穌就是那個顯明的上帝.
耶穌所做的就是父上帝的啟示.
這是天父在隱藏中所做.
而被耶穌看見.
然後耶穌做給我們看.
耶穌告訴我們.
你們如今看到我的教導和作為.
覺得很難受,很差異.
我告訴你,好戲在口頭.
耶穌還有更激烈的作為.
父上帝的旨意會有更爆炸性的披露.
這個更激烈的事是什麼呢?.
就是賜生命.
父上帝是賜人生命的上帝.
耶穌也賜人新的生命.
叫人得著永生.
在24到25節.
耶穌為他能賜人生命做進一步的解釋.
24到25節.
我實實在在的告訴你們.
聽我說,有信猜我來者的就有永生.
不至於定罪.
是已經出死入生了.
我實實在在的告訴你們.
是由將到,現在就是了.
先要聽見上帝兒子的聲音.
聽見的人就要活了.
信耶穌得永生.
凡是信耶穌的,遵行祂的話的人.
就不致滅亡,反得永生.
耶穌為什麼能賜人生命?.
耶穌說,原因是26節.
因為父怎樣再自己有生命.
就賜給他兒子.
也照樣再自己有生命.
再自己有生命是什麼意思呢?.
一方面可以解作上帝的自有永有.
Self existence.
也可以解作上帝是生命的源頭.
就好像一章四節說.

$^{321}$生命在他裡頭.
而我更喜歡的解釋是.
生命的主權是在父上帝和耶穌的手裡.
這裡說的是主權.
由祂說了算.
耶穌讓人生,人就能夠生.
耶穌要人死,人就不能不死.
復活在我,生命也在我.
意思就是,我說就是.
我說誰復活,誰就復活.
耶穌強調的是祂的主權.
不是祂的能力這麼簡單.
是祂的權柄.
耶穌在23節又說.
叫人都尊敬子,如同尊敬父一樣.
不尊敬子的,就不尊敬差子來的父.
耶穌和天父有同等的尊榮.
我們怎樣尊榮天父上帝.
就要怎樣尊榮耶穌基督.
不尊榮耶穌的,就等於不尊榮天父.
我們沒有可以只是尊榮天父.
不需要尊榮子的選擇,沒有.
總結這個段落.
耶穌強調祂和父上帝緊密相連.
性格相同,地位平等.
耶穌所做的就是天父所做的.
耶穌的想法就是天父的旨意.
耶穌和天父有同等的尊榮.
人們對待耶穌的態度.
一如他們對待天父的態度.
天父擁有生命的主權.
耶穌同樣擁有生命的主權.
如果創造是由父上帝主導.
將生命帶來人間.
救贖就由耶穌完成.
耶穌賜人永遠的生命.
耶穌和天父上帝之間的關係.
是父和子的關係.
兩者由愛去連繫.
二十節,耶穌簡單一句.

$^{361}$父愛子導出了父子關係.
一個最大的奧秘.
今天我們沒有時間去講這部分.
好,我們來到第二個段落.
第二個段落強調的.
耶穌所說的是.
耶穌擁有審判世人的權柄.
耶穌賜人生命.
主要說的是末日的時候.
這裡指著末日的審判.
二十二節,耶穌說.
父不審判什麼人.
乃將審判的事全交於子.
二十七節,並且因為他是人子.
就賜給他行審判的權柄.
父上帝不負責審判.
耶穌基督全權負責審判.
他是末日的審判主.
三十節,耶穌又說.
我憑著自己不能做什麼.
我怎樣聽見就怎樣審判.
我的審判也是公平的.
因為我不求自己的意思.
只求那差我來者的意思.
因為天父的心意和耶穌完全相同.
子怎樣審判就等於父子的成就.
這裡沒有上訴機制.
你不可以說耶穌的審判是初審.
如果你不服判決.
可以上訴去父那裡.
耶穌的審判是終極性的.
因為他跟父的旨意完全一樣.
為什麼父上帝會將審判權完全交給耶穌呢?.
真正的原因我們不能知道.
沒有人可以參透上帝的心意.
不過我們可以做一個合理的推斷.
就是在末日的時候.
我們在審判台前能不能過關.
最重要的關鍵.
是在於我們跟耶穌的關係.

$^{401}$無疑我們最後接受審判.
基督的審判是我們在人間的行為.
善有善報,惡有惡報.
耶穌在29節也說.
行善的復活得生,作惡的復活定罪.
保羅在哥林多後書5章10節也說.
因為我們眾人必要在基督台前顯露出來.
叫惡人,惡人按著本身所行的.
或善或惡受報.
我們將來是按著我們的行為去接受審判.
不過聖經也告訴我們.
人間沒有一個二人.
沒有人能夠靠行為滿足上帝的要求.
所以其實審判的結果是不用審.
不審也知道.
人人都過不了關.
人人都被定罪.
這裡沒有懸念,沒有變數.
但是上帝卻為人類預備了一個救法.
一個唯一能夠逆天改命的途徑.
就是信耶穌.
被祂的寶血遮蓋得稱為義.
如此,人人都是罪人,是沒有變數的.
但是是不是相信耶穌就是唯一的變數.
在24,25節,耶穌在這裡很清楚地告訴我們.
我實實在在地告訴你們.
聽我說,有信猜我來者的就有永生.
不至於定罪,是已經出死入生了.
我實實在在地告訴你們.
時候將到,現在就是了.
死人要聽見上帝兒子的聲音.
聽見的人就要活了.
耶穌在這裡連續用兩次最鄭重的語氣.
我實實在在地告訴你們.
指出這個很重要的事實.
唯有信耶穌,相信祂彌賽亞的身份.
又遵行祂命令的人.
就不會被定罪,可以出死入生.
祂說,凡是聽到上帝兒子聲音的人都能得到生命.
這句話是什麼意思呢?.

$^{441}$我可以想像,大概就在末日.
無數人,所有人都站在審判台前.
基本上都是不用審的.
人人按照行為都死定的.
但耶穌走出來,就讀出一連串的名字.
庶金,加侖,馬上就能夠執到.
你明不明白?.
當祂讀出我們的名字的時候.
就意味我們過關.
不致滅亡,反得永生.
所以末日其實真的不用怎麼審的.
難道還要拿這些證據.
全知全能的上帝跟你玩這個遊戲.
你還有什麼好辯護的?.
我可以解釋,我其實是不小心的.
不覺意的.
說這些話都是多餘的.
沒有營養價值的,是不是?.
所以,只是代耶穌下命.
唯有信耶穌的人,才能得永生.
人在審判台前能不能夠過關.
關鍵端在於和耶穌的關係.
耶穌讓你過,你就過.
耶穌不讓你過,就沒有生路.
既然是這樣.
父上自然將審判的權柄完全交付耶穌.
因為反正是耶穌說了算.
基督徒,我們知道.
我們之所以獲得永生.
永遠不是審判的結果.
而是出於上帝的恩典.
我們得救是本乎因,也因著信.
不是出於行為,免得有人自誇.
如果你和我能夠最終過關.
一定不是因為我們行善.
是大善人,所以能過關.
是因為我們信耶穌.
耶穌認識我們.
你記得他說過.
我們在人間認耶穌.

$^{481}$將來,他就認我們.
他就說他認識我.
他認識我,是我能夠過關的唯一關鍵.
我們背的《紹浩信經》說.
耶穌基督將來必從那裡.
那裡是天上降臨.
審判,活人,死人.
這是我們的信仰.
好,我們來到最後一段.
我們剛才讀經沒有讀.
第三個訊息是.
耶穌神聖身份的見證.
耶穌說了上述的說話.
這真是一個驚世大奧秘.
自然讓跟隨他的門徒非常驚嚇.
讓那些本來反對他的人很憤怒.
宣稱自己擁有神聖.
宣稱自己與上帝平等.
這多麼的褻瀆.
就是那些跟隨耶穌的人.
不免都會半信半疑.
夫子,你是否信口雌黃.
你是否吃錯藥,胡言亂語.
你說這些天馬行空的話.
究竟有什麼證據呢.
是的,耶穌需要為他這個驚人的宣告提供證據.
但是,耶穌是上帝這個宣稱.
在人間還可以找到什麼證據呢.
人如何證明上帝呢.
當然是無法證明.
只有上帝能夠證明自己.
所以,耶穌的宣稱最穩當的證明是什麼呢.
就是天父上帝走出來證明.
只要父上帝說耶穌是上帝.
就過關了.
就像耶穌在受洗的時候.
天父公開宣告,這是我的愛子.
不過,不是每個人都聽到上帝的宣告.
有人當他是打雷而已.
不是聽得很清楚.

$^{521}$並且,我不在受洗的現場,我怎麼知道呢.
除非上帝每天都出來宣告幾次.
不過,如何證明那個聲音.
不是我們耳鳴.
是真的,父上帝的聲音也是一個問題.
所以,37節,耶穌這麼說.
差我來的父也為我作見證.
你們從來沒有聽過他的聲音.
也沒有看見他的形象.
這是他的遺憾.
除了父上帝親自作證之外.
耶穌提出另一個有效的證據.
就是由耶穌將要成就的救恩來證明.
36節,他說.
但我有比約翰更大的見證.
因為父交給我要成就的事.
我所作的事,這便見證.
我是父所差來的.
當人看到耶穌的使命成就.
就不可以不相信他的話是真的.
特別將來的末日.
當我們都自願被迫走到耶穌的審判台前.
見到,哎呀,原來站在審判台上的我真的是他.
你想不承認也不行,對不對.
你想不承認他是上帝也不可能.
28節,他說.
你不要把這事當作稀奇.
時候要到,凡在墳墓裡的.
都要聽見他的聲音就出來.
耶穌在約翰福音12章32節也說過.
我若從地上被舉起來.
就要吸引萬人來歸我.
所以,最終我們也會看到耶穌所說的是真的.
不過,風水佬騙你十年八年.
耶穌說的這個證據.
就不止十年八年.
從耶穌說的時候.
到末日,現在也還沒到,真可憐.
就要到二千多年後.
你覺不覺得很奇怪.

$^{561}$我給你一個證據.
不過證據是二千多年後才看到的.
說還沒說,就多餘了,對嗎.
所以,第二個可能的證據.
耶穌說,就看我做的事.
最終的結果你們就知道了.
現在不會知道了.
所以第二個證據也是很奇怪的.
耶穌於是只能夠提出兩個比較次要的證據.
天父上帝的見證是第一個證據.
他的事最終成就是第二個證據.
第三,第四個證據是什麼呢.
第三個證據是.
施洗約翰為他做見證.
32到35節.
另一位給我作見證.
我知道他給我作的見證是真的.
你們曾猜人到約翰那裡.
他為真理作過見證.
其實我所受的見證不是從人來的.
然而我說者之話.
我要叫你們得救.
約翰是點著的明燈.
你們情願暫時喜歡他的光.
我們知道施洗約翰在約翰的福音裡.
佔了很重要的位置.
前面有兩段經文提及施洗約翰.
約翰福音一章15節.
19到37節.
然後二章22到30節.
福音提到有關約翰的故事.
主要不是說他的生平.
而是轉引他為耶穌做見證的說話.
施洗約翰是當時多數人公認的先知.
也是猶太社會在三百多年來.
唯一出現的先知.
耶穌曾經稱施洗約翰是舊約時代最後.
也是最偉大的先知.
六個福音七章26到28節.
施洗約翰治國最重要的任務.

$^{601}$不是施洗不是叫人悔改.
而是預備主的道.
修籍主的路.
是為耶穌做見證.
不過有關施洗約翰的見證.
我們之前說過三次.
三篇道我們說過.
今天不重複了.
第四個耶穌提出有關神聖身份的證據.
就是聖經.
就是律法書和先知書.
就是我們所說的舊約聖經.
39節你們查考聖經.
因你們以為內宗有永生.
給我作見證的就是這經.
摩西所說的話.
都是指向耶穌的.
猶太人如果信摩西.
就應該信耶穌.
46,47節.
耶穌和日後的使徒都宣稱.
整本舊約聖經都是為耶穌做見證的.
先知有關末後的日子.
有關彌賽亞的所有預言和應許.
都是應驗在耶穌身上的.
也是為預告耶穌的身份和使命而說的.
是嗎?.
給我作見證的就是這經.
給耶穌作見證的就是舊約聖經.
這句話我以前也說過.
對外邦基督徒.
包括你和我在內是非常重要的.
沒有人要先成為猶太人再成為基督徒.
我不是猶太人.
我也不是猶太教徒.
我是信耶穌的.
猶太經典對我有什麼權威?.
我為什麼要讀律法書和先知書?.
唯一的原因.
耶穌宣告給他作見證的是這本聖經.

$^{641}$律法書和先知書.
對比耶穌所訂立的新約.
被稱為舊約.
對我來說.
沒有什麼唐乃馨或先知的獨立地位.
沒有.
獨立一本經典.
我會看它和《可蘭經》一樣.
因為是宗教經典.
但是.
在我眼中.
只有依附耶穌基督並他釘十字架的新約的舊約.
你聽明白我的意思吧.
相對於新約它叫舊約.
沒有舊約何來新約?.
舊約從來不被獨立看待.
必須要做基督論的詮釋.
我們要在裡面看到耶穌基督.
因為耶穌清楚指出.
給我作見證的就是這經.
回到耶穌神聖的證據這個題目.
如果父上做見證.
耶穌的工作果效做見證.
即時被人看見.
施一約翰的見證和聖經的見證.
就是最能夠拿出來的證據.
畢竟對當時的猶太人來說.
施一約翰是他們多數人都公認的先知.
而律法書有權威更加無庸多言.
但這兩個證據能不能夠簡單清楚證明耶穌的神聖身份呢?.
還是有很多相確餘地.
我們知道,猶太律法師和法理賽人.
不一定都信施一約翰.
群眾信施一約翰.
但不是所有猶太教的領袖都信施一約翰.
我們看六合福音七章三十節就知道.
而將律法書和先知書.
所有應許都指向耶穌.
更加不是所有人都認同的觀點.
所以根根究底.

$^{681}$耶穌的神聖身份仍然是難以證明的.
耶穌自己也承認這是無法得到客觀證明的一件事.
祂指出最終都是信心的問題.
因為信仰是不會有百分百客觀證明這件事.
另外更加重要的是信仰態度的問題.
就是我們是否將宗教信仰變成一個自編自導自演的遊戲.
耶穌說在42到47節.
但我知道你們心裡沒有上帝的愛.
我奉我父的名來.
你們並不接待我.
若別人奉自己的名來你們也要接待他.
你們互相受榮耀.
卻不求從獨一之上帝來的榮耀.
怎能信我呢?.
哪有想我在父面前要告你們?.
可以告你們的就是你們所仰賴的摩西.
你們如果信摩西也必信我.
因為他書上有指著我寫的話.
若不信他的書怎能信我的話呢?.
這句話有很豐富的含義.
信的本質是什麼?.
因為時間所限今天就不說了.
有機會再和你們說有關信仰態度的課題.
我們對今天的經文做一個總結.
你們很辛苦了.
引我講了這麼久.
就是這三點.
月前前面說的.
耶穌從來沒有直接說他是上帝.
不過他是很清晰的宣稱.
他和父上帝同質同權.
思想行為完全相同.
對當時的猶太人而言.
這跟耶穌說他是上帝是毫無分別的.
所以是最具洩毒性的話.
耶穌用間接的方法宣告他是上帝.
這讓所有要思考.
如何對待他的人毫無退路.
很多人欣賞耶穌的教導.
打算視他為道基督教師.

$^{721}$再不是偉大的先知.
偉大的宗教家.
但是說他是上帝.
我們就必無選擇.
或者不信.
從而認定他要麼就是騙子.
要麼就是傻子.
要麼你就相信.
你認順我的主.
我的上帝.
耶穌和天父上帝平等.
耶穌是上帝.
在耶穌仍然以人的身份活在人間的時候.
這個宣告確實是難以讓人接納的.
不要說那些反對他的猶太教領袖.
跟隨他的門徒.
都不見得能夠即時理解和接受.
不過當耶穌成就十字架的救恩.
復活升天之後.
要接納這個事實就容易很多.
不過今天我要說的是.
我們是否承認耶穌是上帝.
成了我們真實的基督徒的分水嶺和試金石.
我們很多人都會相信耶穌.
我們都會相信耶穌是神聖.
不過作為中國人的我們.
如果我們仍然承襲著一些傳統中國人的宗教信仰觀念.
其實我們不是每個人都想玩這麼大.
很多研究中國宗教特質的人都說過.
學者都說過.
中國人相信滿天神佛.
但他們通常選的神明都不會選很大的.
很大的神明是沒有群眾基礎的.
在長洲那些神明階級最高的是關公.
關公廟在山頂上沒有人拜.
拜得最多的是北帝廟.
是觀音廟.
因為中國人很現實.
拜一個這麼崇高的幹什麼.
我又不是要他什麼.

$^{761}$我只是要他給我六個號碼.
他能夠給我六個號碼就夠了.
拜一個低低的容易說.
太崇高很難受.
他講完我沒什麼可以討價還價.
中國人不用拜這麼高.
不怕官最怕管.
我搞得定我的鎮的鎮長就夠了.
我不用去到習近平的地步.
那不關我的事.
中國人拜祖先不用很厲害的神明.
只要他能夠保佑子孫就夠了.
我不需要拜到這麼盡.
因為拜到盡其實是將我們自己逼入牆角.
我們別無退路.
如果耶穌是一個次等神明.
你渾水摸魚.
說幾句就可以過關.
全知全能的上帝其實很難受.
兄弟姐妹.
我希望我說的這些是廢話.
對你完全沒有意義.
你根本沒這樣想過.
是真的.
我們今天信的一個創天造地.
知道我的過去現在未來的上帝.
活在一個這樣的上帝面前.
你和我可以怎麼做.
應該怎麼做.
我們是否真的由衷地認定.
這位上帝是最高的上帝.
是一個對我有百分百的主宰和管轄權的上帝.
我的主,我的上帝.
\newpage



\section{腓立比書 4:1-9}
\label{sec:rzRD6n7AW0w}
\textbf{2024年7月28日崇拜直播｜陳冠成牧師｜靈性·關係·成長－靈性安穩與探索｜腓立比書4\_1-9}
\newline
\newline
連結: \href{https://youtube.com/watch?v=rzRD6n7AW0w}{\texttt{https://youtube.com/watch?v=rzRD6n7AW0w}} ~~~~ 語音日期: 2024-07-29
\newline
\newline
\hyperref[sec:BxDFSB8mRZw]{\small{< < < PREV SERMON < < <}}
~
\hyperref[sec:index_chronic]{\small{[返順時目]}}
~
\hyperref[sec:OxtvYBdmNrI]{\small{> > > NEXT SERMON > > >}}
\newline
\newline
腓立比書 4:1-9
\newline
\begin{longtable}{cl}
\hline
\hline
章節 & 經文 (和合本修訂版)\\
\hline
4:1 & \begin{tabularx}{0.7\textwidth}{X} 我所親愛、所想念的弟兄們，你們就是我的喜樂，我的冠冕。我親愛的，你們應當靠主站立得穩。 \end{tabularx} \\ \\ \relax
4:2 & \begin{tabularx}{0.7\textwidth}{X} 我勸友阿蝶和循都基要在主裡同心。 \end{tabularx} \\ \\ \relax
4:3 & \begin{tabularx}{0.7\textwidth}{X} 我也求你這真實同負一軛的，要幫助這兩個女人，因為她們在福音上曾與我、革利免和我其餘的同工一同勞苦，他們的名字都在生命冊上。 \end{tabularx} \\ \\ \relax
4:4 & \begin{tabularx}{0.7\textwidth}{X} 你們要靠主常常喜樂。我再說，你們要喜樂。 \end{tabularx} \\ \\ \relax
4:5 & \begin{tabularx}{0.7\textwidth}{X} 要讓眾人知道你們謙讓的心。主已經近了。 \end{tabularx} \\ \\ \relax
4:6 & \begin{tabularx}{0.7\textwidth}{X} 應當一無掛慮，只要凡事藉著禱告、祈求和感謝，將你們所要的告訴神。 \end{tabularx} \\ \\ \relax
4:7 & \begin{tabularx}{0.7\textwidth}{X} 神所賜那超越人所能了解的平安，必在基督耶穌裡，保守你們的心懷意念。 \end{tabularx} \\ \\ \relax
4:8 & \begin{tabularx}{0.7\textwidth}{X} 末了，弟兄們，凡是真實的、凡是可敬的、凡是公義的、凡是清潔的、凡是可愛的、凡是有美名的，若有甚麼德行，若有甚麼稱讚，你們都要留意。 \end{tabularx} \\ \\ \relax
4:9 & \begin{tabularx}{0.7\textwidth}{X} 你們從我所學習的，所領受的，所聽見的，所看見的事，你們都要繼續去做，賜平安的神就必與你們同在。 \end{tabularx} \\ \\
[1ex]
\hline
\hline
\end{longtable}
$^{1}$今天「經憤」記載在腓立比書四章 一字九字.
是保羅對腓立比教會的勸導.
我所親愛 所想念的弟兄們.
你們就是我的喜樂 我的冠冕.
我親愛的弟兄 你們應當靠主 站立得穩.
我勸有阿蝶和秦都基 要在主裡同心.
我也求你 借真實同父一呃的 幫助這兩個女人.
因為她們在福音上 曾與我一同勞苦.
患有隔離免 並其餘和我一同作功的.
她們的名字 都在生命冊上.
你們要靠主常常喜樂 我再說 你們要喜樂.
當叫眾人知道你們 謙揚的心.
主已經近了 應當一無掛淚.
將你們所要的 告訴神.
神所賜出人意外的平安 必在基督耶穌裡.
保守你們的心懷義念.
弟兄們 我還有未盡的話.
凡是真實的 可敬的 公義的 清潔的.
可愛的 有美名的 若有什麼得行.
若有什麼稱讚 這些事你們都要思念.
你們在我身上所學習的 所領受的.
所聽見的 所看見的.
這些事你們都要去行.
賜平安的神 就必與你們同在.
弟兄姊妹 平安.
來到腓立比書最後的一講.
如果你記得上一次 其實 波木跟我們提到第四章的一開頭.
裡面講到保羅勉勵菲律賓的信徒 要靠主 站立得穩.
一方面提醒他們 在一個不容易見證信仰的世代.
要像下法保羅一樣 堅守他們天國公民的身份.
與此同時 這句說話帶出下文 就是保羅最後的勉勵.
他希望幫助菲律賓教會的信徒兩方面.
包括如何處理在教會內部信徒之間的糾紛.
同時也正視基督徒與世界的關係到底如何.
幫助教會的信徒 無論在教會裡面或外面的生活.
都可以靠主耶穌來站穩.
經文提到幾個名字.
你會發覺 雖然保羅一直對菲律賓教會讚不絕口.
但不代表裡面是一些問題都沒有.
畢竟我們都明白有人的地方會有什麼問題.

$^{41}$有人事的問題.
保羅提到首先有阿蝶 秦都基兩位的姊妹.
雖然我們不清楚他們的背景 因為我們不是第一身的收信人.
但從經文中可以看到他們可能是菲律賓教會開創時.
初期發展時他們已經信了主.
並且他們很熱心地事奉.
因為保羅提到兩個人都在福音的事上.
和保羅一起童工.
並且一起的勞苦.
這兩個人大概都是菲律賓教會一些主要成員.
保羅我們不難想像他從爾巴弗提的口中.
知道這兩位姊妹近來關係出現了一些問題.
不過你見到事情應該未是鬧大.
因為保羅沒有長篇大論去講他們的情況.
大概還未影響到菲律賓整體教會的關係.
所以保羅沒有嚴厲地譴責這兩位姊妹.
不過他很清晰正中點名.
他指明這兩位姊妹要在主的同心.
勸他們盡早化解彼此之間的矛盾.
保羅提到另一個人的名字.
這裡提到一位真實和福音的信徒.
同樣我們不知道他的身份.
不過很明顯保羅拜託的一定是教會.
一些有熟齡分量的人.
熟齡程度相當成熟的人.
並且應該善於處理人際關係.
所以保羅特意拜託他幫這兩位姊妹和好.
經文又提到第三個名字叫做甲利敏.
Kellerman這個名字.
其實在羅馬帝國當時是一個很平凡普遍的名字.
不過很明顯他在教會有沒有好的明星.
因為保羅說他的名字都寫在生命冊上.
保羅就借用甲利敏這個名字來表示.
其實有蝶同秦都基和甲利敏一樣.
他不是壞人.
這兩個姊妹他們都有好的明星.
他們的名字同樣寫在生命冊上.
有好的素齡質素.
所以換句話來說這兩位姊妹的爭拗.
是不是好人和壞人之間的矛盾.

$^{81}$不是,是不是壞人和壞人的爭拗.
更加不是,而是怎樣.
是兩個好人的爭拗.
是不是很有趣.
你有沒有發現教會有時都是這樣的.
不是壞人和壞人,或者好人和壞人之間的問題.
是兩個都愛主的人的問題.
明明兩個人都愛主,熱心事奉.
他們的名字保羅都說寫在生命冊上.
但卻有可能因為性格,意見,處事的手法不同.
就會有嚴窮甚至不和.
嚴重的話可以影響到一個教會群體造成分裂矛盾.
窒礙教會的成長.
在這個時候就很需要有人在身邊提醒他們.
有人在身邊勸他們協助他們化解彼此之間的矛盾.
我們看到保羅其實他做了一個非常好的示範.
首先保羅保持中立.
你會見到他沒有偏袒任何一方.
經文原本是說他勸有阿諜,亦都勸秦都基.
兩個都勸,沒有說誰對誰錯.
保羅提醒他們兩個都要在主裡面去同心.
表明首先這是不是他們兩個人的私事.
不是,在主裡面即是說這件事是甚麼.
跟上帝有關的,跟耶穌有關,跟教會有關是公事.
這兩個人的不和是很大機會會有機會令到教會受虧損.
令到耶穌基督的身體受虧損.
事實上我們不難理解.
可能我們回教會來臨都會經歷過類似的情況.
我還記得我以前初信的時候.
我在團契裡面那時候大概信了主三四年左右.
當時我是團長.
當時那個團契即稱團契30歲左右.
你知道這些時間很容易會有男女拍拖.
拍拖的時候亦都有機會不和的時候會怎樣.
造成兩班人的對立.
你很難說哪一班人對錯有時感情的問題.
但自然很喜歡人就會在裡面分對錯.
雙方都有自己的認同者.
你站在中間怎樣做?.
不容易,有人提醒我一定要中立.

$^{121}$因為這不是純粹他們的感情糾紛.
這是在主流的公事.
你要持平,聽到雙方的意見.
並且讓他們看到,不是純粹感情的問題.
這是教會相交的事情.
要處理得很好,不只是解決你們的感情糾紛.
而是要怎樣竭力在主裡面去合一,去包容.
確實,弟兄姊妹,正如保羅所說.
他勸誘蝶,秦都基要回到主爺秀的裡面.
因為他們的關係不是靠人的意志或情商.
來達成和好.
而是要從屬靈層面來化解彼此之間的分歧.
唯有雙方都順服在耶穌基督的主權裡.
就像之前的腓立比書一直說的.
要以耶穌基督的心為心.
為了福音的緣故,為了上帝的緣故.
願意先放下自己.
願意謙卑來到上帝的面前.
看別人比自己重要.
看到對方其實和自己一樣都是上帝所做.
都是上帝所愛的兒女.
大家其實都是服侍同一位主.
唯有這樣在主的裡面.
人與人之間的分歧才有機會化解.
重新達致合一和同心.
事實上,弟兄姊妹,其實我們很明白.
我們屬靈群體真正的敵人.
第一個是誰?.
是誰?.
耶穌說是撒旦,是魔鬼.
但是在現實生活.
我們會否發覺有時基督徒花了很多力氣.
處理人事的問題?.
覺不覺?都覺?.
甚至有時會視弟兄姊妹為自己的對頭人.
甚至視她為敵人.
說她確實是魔鬼撒旦常用的伎倆.
來破壞基督身體的合一.
她很想轉移我們的視線焦點.
讓我們不知不覺將精力,人生耗費在不必要的爭拗裡.

$^{161}$無法履行上帝所託付給我們的使命.
我相信教會裡沒有人全心想做壞人,想傷害別人.
在座任務當中,你是立心回教會害人做壞人.
舉手給我看.
不會吧?.
如果有的話,待會再聊.
我們沒有立心回來害人,壞人.
我們每一個人都一定帶著良好的想法,語言,背景.
希望如何?我是為你好,我想教會好.
我們都是以這個出發點來和人相交.
不過,聖經裡面說我們的本相又是怎樣?.
我們是一群蒙恩的罪人,走在一起,越近又怎樣?.
越有機會互相刺傷對方.
我們都避免不了,因為我們真的不是完美.
我們看東西總會有看不完的地方.
我們有時看錯,看漏,有時可能是聽錯,聽漏.
又或者我們本來是想做一件好事,但時機掌握不了.
或者我們的用詞與氣不恰當,導致對方接收不了.
人與人之間的厭隙有時不是這樣產生.
是兩個好人之間的矛盾.
在這個時候,保羅提醒我們,真的要回到耶穌基督裡面.
你要將你的淑女生命安穩在耶穌裡面.
讓上帝監察我們的心思意念.
讓上帝潔淨我們的心思意念.
以致我們真的可以再一次以上帝的目光.
以耶穌基督由和謙卑的心腸.
回到我們的人際關係裡面.
使我們的言行真的可以榮耀上帝造就身邊的人.
這是保羅剛才提醒我們.
要再主流去面對人際之間的糾紛.
另一方面,你會留意到保羅在規勸尤亞諜和秦都基的同事.
他沒有忘記肯定他們的尾線事情.
他加許這兩位姐妹其實都是忠心愛主.
直至強調,直至強化他們和好的動力和意願.
除此以外,保羅可能都考慮到.
始終我們常說當局者迷.
兩個在吵架的人未必有動力.
未必有足夠的智慧和意願主動和好.
所以保羅就拜託另一位真實同父恩丫的同工.
來協助劃船,促進他們姐妹之間的和好.

$^{201}$保羅向我們示範了一個我們說勸說的智慧.
一方面你要有一個柔和的心,謙柔的心.
不要只是看到別人的問題.
不要只是追著別人做得不好的地方.
而忘記了對方其實有些事情值得欣賞,值得肯定.
因為對方和我們一樣都是上帝所做的.
有上帝的尾線,有上帝給他的獨特性.
我們在提醒,我們在督責的同時.
不要忘記去發掘欣賞,肯定對方這些尾線.
特別是他一些良好的動機和品格.
特別不知道大家有沒有留意.
我們勸說的時候,經常說「不過」這兩個字.
有沒有留意?我們和別人對話的時候.
可能我們先讚好,你看上去都不錯,這件事都做得不錯.
不過,這兩個字又如何?.
將你之前的欣賞肯定化解了.
所以我看過有人說得好.
他說其實和別人溝通,特別在你勸說的時候.
化解恩怨的時候,你少些說「but」.
多些說「and」.
意思是什麼呢?其實你都做得很好.
只要你再做這件事,就更好.
如果你肯這樣留心這件事,就會更好.
豈不是更加能夠造就人,更加能夠幫助人改進呢?.
所以真的,很多時候我們很容易說一大堆讚美的說話.
肯定的說話,但突然間說「不過」.
大家都會知道,那個才是你的重點.
但很多時候這句說話,就是將你之前的欣賞肯定前功盡廢.
好像倒了進海一樣,對方聽不入耳.
另一方面,你會看到,在勸說的過程裡面.
找一個對的人是很重要的.
找一個懂得去滑船,懂得去勸教的人.
是非常之重要的.
有沒有見過一些不太成熟的人.
他越勸教又如何?越壞事.
我們說「拿著火水去救火」.
拿著火水去救火,即是怎樣?.
越潑就越燒得含含聲.
很多時候因為是勸教的人,或者說篤責的人.
他自己都沒有在主的裡面.

$^{241}$他都沒有首先去疏離他自己的情緒.
可能他對當事人都有點不滿.
在勸說的時候,一併爆發出來.
熱至到火上加油.
越勸說,反而越勒著更多的人.
擴大更多的矛盾對立面.
相反,保留這裡說.
你拜託一位在靈裡面成熟的長輩.
一方面他會保持中立.
另一方面,他不會偏袒任何一方.
他會聽取雙方的意見.
並且他會堅持在主裡面的大原則.
他確實說的話不是要拆毀人.
不是要增添人的煩惱.
而是要造就人,幫人帶到耶穌基督的面前.
發揮穿針人線的作用.
所以丁姐妹其實單單保留這兩句說話.
其實提醒我們.
我們避免不了人際關係的問題.
你跟我回教會.
本身一個很重要的課題就是怎樣?.
相教.
我們信了耶穌之後.
不是孤零零自己躲在房間裡讀聖經就搞定.
我們一定要面對人,避免不了.
面對人的過程裡面.
越陷入陷阱,就有機會沒有風險被傷害.
或者反過來是你傷害人.
所以其實我們都需要多多少少學會怎樣去面對處理人際關係.
所以你會發覺今年教會都開了一些.
與人際關係相關的主一學或者說門訓營會.
很想幫助大家來到我們怎樣在上帝的裡面.
在耶穌基督裡面有智慧去認識自己認識別人.
又或者進一步幫助自己和別人.
在那些分歧裡面怎樣處理得合意.
怎樣可以榮神益人建立一個健康的人際關係.
屬靈關係.
好,我們再看下去.
我們再看餘下四至七節.
我想請大家讀一讀,好嗎?.

$^{281}$我們有一段很熟悉的經文.
《腓立比書》四至七節,預備起.
你們要教主上上基督,我在說,你們要理內.
當要眾人知道,你們必須要穿上主義的衣服.
主義的衣服,應當人好化了.
主要是萬事節的高高,以外就不安靜.
將你們所要的基督教會.
告訴神,神所賜出人意外的平安.
必在基督耶穌裡保守你們的心懷意念.
你會發覺這段經文我們熟悉得不得了.
裡面很多都是我們很喜愛的金句.
保羅他說,面對教會,無論在外面有威脅.
內部亦都有一些隱患的時候.
他提醒了幾個很重要的原則.
要求信徒來執行.
首先他命令丁姐妹要常常喜樂.
不單單說一次,要說兩次.
要靠主常常喜樂.
不過,實際如何去遵行這個命令.
人有時,是否真的可以時時刻刻.
都維持做一個歡喜快樂的心情.
可不可以?.
可不可以時時刻刻都維持一個歡喜快樂的情緒.
做不到的,現況.
到底保羅這段,我們說.
看上去每個字都明,靠主喜樂,常常喜樂.
但實際如何做得到呢?.
話說有個小朋友跟媽媽說.
媽媽,你暑假了,你買部遊戲機給我.
我很恨,你買給我,我就有喜樂.
於是父母就用聖經跟他說.
兒子,你的喜樂與有沒有好事.
臨到你的身上是沒有關係的.
你的喜樂與有沒有這部遊戲機.
亦都不是必然的.
你的喜樂與你的期望達不達到.
亦都沒有一定程度的關係.
喜樂是你個人本身的責任.
這句說話相當有意思.
喜樂是我們個人的責任.

$^{321}$喜樂是上帝要求我們培育出一種做人的態度.
偶爾我們可能會聽到身邊的人說.
等我找到這份好功,我就會有喜樂.
等我找到我所摯愛的,與他長相廝守.
我就會有喜樂.
等我賺夠,退休,就會有喜樂.
等我捱大了子女出生,就有喜樂.
等我這個病治好,我就可以有喜樂.
有沒有聽過?.
可能有時我們自己都這樣想.
不過你會發覺上述的例子.
最後即使他們的心願目標都達成.
是不是一定就會常常喜樂?.
卻又不是.
是不是他們不夠努力?.
不是.
是不是他們很希望達成目標的意念不夠執著?.
通常也不是.
而是很多時候人會將喜樂的次序調轉了.
他們將喜樂視為一種結果,視為一種副產品.
等我得到某樣東西滿足的時候,我才可以喜樂.
保羅這裡卻是要改變我們這種想法.
他提醒我們,首先是要靠主來喜樂.
他不是呼籲我們要維持一種喜樂的情緒.
靠自己的意志.
他要提醒我們.
你要確信凡事有上帝的掌管.
並且認定上帝是美善的,上帝是會愛我們.
所以我們可以憑信心宣告.
沒有任何事情,沒有任何權勢.
可以破壞上帝的美善旨意在我們身上彰顯.
因為上帝看著我們.
所以我們可以上上喜樂.
就是這個意思.
沒有特別原因,不是因為我某些願望達成.
不是因為我跨過某個逆境.
甚至不是我在困難裡找到一些人生意義.
全部都不是.
單單是因為有上帝.
是因為一切上帝包底,上帝搞定.

$^{361}$所以你可以安心靠近上帝的喜樂.
換句話說,保羅強調上上喜樂.
和我們環境的信義沒有必然關係.
上上喜樂反而反映你和上帝的關係.
之所以能夠上上喜樂是因為什麼原因.
就是因為你和上帝常常結連.
就是這麼簡單.
你常常靠近上帝.
你知道無論你經歷什麼上帝都看著你.
在你旁邊的時候.
你便能夠寬容,你便能夠喜樂.
因為你知道不是靠你自己來控制.
不是靠你自己掌握.
是靠著上帝靠他帶領你.
你會發現.
在這個世界裡面.
沒有任何事情能夠奪去上帝給你的那份淑天的盼望.
你要確信自己是活在上帝的恩典,保守和帶領的裡面.
亦因為這樣,我們可以在人的前面.
無論別人對我們怎樣都好.
可以像保羅所說.
展現出那份謙和的特質.
我們對事物,對自己的際遇.
不會怨天尤人,不會走極端.
亦不會過於抓住自己的權益來不放手.
我們更加不會對人過份的偏激.
反而我們願意去相信別人.
我們願意去對人寬容.
即使我們面對那些對我們不好的人.
我們亦不會以牙還牙.
這就是保羅所說.
因為喜樂而有那種謙和的主事態度.
簡單來說.
保羅這裡說.
喜樂就是依靠上帝你坦然的面對自己.
面對自己的際遇.
面對你身邊的事物.
一種積極正面的做人態度.
這是保羅第一個要求弟兄姊妹培養的.
第二方面.

$^{401}$又是我們很熟悉的聖經禁句.
應當一無掛律.
會不會對這個命令覺得很困難.
我聽過很多弟兄姊妹覺得.
明白,明白為甚麼.
不過做不到.
人怎可能沒有憂慮.
憂慮本來就是一個.
我們與生俱來的情緒.
很難叫人不憂慮.
說回幾年前COVID你有沒有憂慮.
過了,我忘記了.
你回想起幾年前.
COVID一開始最憂慮的是甚麼.
一層層地來的.
沒有口罩,沒有搓手液.
沒有廁紙,用品.
這件事害怕,那件事害怕.
這是第一層.
第二層呢.
開始解決了這些.
糟了,小朋友怎樣上學.
安排不到,跟你一起.
家居,那怎樣.
打大交.
或者家裡有老人家.
要看醫生覆診,又怎樣.
很多憂慮再遠一點.
經濟開始下滑落王.
有些人的工作不保.
很多商舖結業.
甚至教會怎樣.
教會其實.
不是說王宅.
普世教會其實都在COVID裡.
受到很大的影響.
實體都回不了.
怎樣去牧養,怎樣去聚會.
一堆一堆的憂慮.
怎樣去不憂慮.

$^{441}$說近一點.
上星期有沒有人坐飛機.
有沒有人要.
剛剛遇到星期五.
星期五,Window,單機.
有沒有人出事.
或者你身邊的人.
如果你看新聞的時候.
你會發現很難不憂慮.
因為不只是.
飛機出發不了.
然後所有他的行程.
酒店,甚至要轉機.
通通怎麼辦.
你只是賠機票給我沒用.
我都一堆憂慮要處理.
憂慮就好像變了骨牌效應.
很多旅客.
不單單行程.
心情都大受影響.
事實上在生活裡面.
我們一定有一些東西.
是掌握不了,預計不了.
因為這樣我們一定會有憂慮.
不過.
保羅這裡給了一個方法.
我們不是要讓憂慮消失.
是不要讓它困住你的人生.
保羅這裡說.
凡事藉著禱告祈求.
並且帶著一個感恩的心.
將你一切的需要.
告訴上帝.
保羅.
這個教導很簡單.
原因是.
他說主已經近了.
意思其實.
有兩方面,一方面是說.
空間上的接近就是.

$^{481}$時間上的接近.
就是耶穌基督很快回來.
另一個就是空間上的接近.
就是其實耶穌很近.
在你旁邊.
你有什麼顧慮,有什麼需要.
你隨時告訴他.
他是知道的.
他願意去聆聽.
所以之所以要隨時.
來禱告教託.
不是因為耶穌基督忘記我們.
善忘,所以我們要常常.
他日餓夜餓,餓到他.
幫我們搞定為止.
教託其實是.
培育我們.
時常連結著耶穌基督.
那個關係.
提醒我們其實耶穌沒有忘記我們.
他每時每刻都眷顧我們.
保守我們.
當然有些時候.
他不一定按我們的心意來成就.
不一定立即解決.
我們遇到的困難.
有時他甚至.
會容許一些困難落在你身上.
不離開.
他只是要你學會.
你搞不定,你要放手.
你要挨笨到耶穌.
讓他幫你.
你要學習信心等候的功課.
因為他允許.
一定能搞定.
正如保羅所說.
當我們願意禱告.
祈求教託的時候.
沒有允許即時解決我們的問題.

$^{521}$卻是允許.
他會賜下足夠的.
平安,出人意外的平安.
保守我們在耶穌基督裡的心懷意念.
正如主耶穌基督.
他臨離世前.
他也跟門徒說,在這個世界上.
一定會有什麼?.
苦難,耶穌沒有違反商品說明條例.
你一定遇到.
但同時耶穌說什麼?.
他說他搞定了.
他沒有說我們會搞定.
他說他會搞定.
你相信他.
他已經勝過世界.
你安穩在裡面.
一定會有平安.
所以就算環境多惡劣.
多嚴峻.
保羅提醒我們.
教託就是我們確信.
神應認長管.
如果我們不確信上帝長管.
我們教託給他做什麼?.
最後你自己拿回來做.
教託就意味著.
你一定要放手.
信上帝,不是信自己.
並且確信耶穌.
很愛我們.
有需要的時候.
一定向我們施恩典幫助.
所以.
我們之所以可以凡事.
帶著感恩的心來到上帝面前祈求.
是因為我們確信凡事.
上帝都有長管.
凡事上帝都有足夠的供應.
和保護.

$^{561}$事情的好壞不再是.
我們感恩的焦點.
上帝本人.
才是我們感恩的焦點.
換句話來說.
感恩其實也好像起落一樣.
是一個信心的宣告.
宣認上帝在我們身上.
的主權和大能.
沒有任何事情能夠將上帝.
在我們身上的慈愛拿走.
上帝依然是萬有的根基.
所以感恩其實不是被動的.
不是你等到有一些好的事情.
發生然後我們出現的行為.
感恩反過來是我們凡事.
主動去連結上帝.
一種信心的表達.
我還記得.
曾經聽過一個馬丁路德的故事.
有一次他自己.
在宗教改革的過程中.
面臨很大的危機和逼迫.
極其的沮喪和很多的憂慮.
他太太見到他.
就特意穿上一件.
喪服.
即是全身黑色的服飾.
出現在他面前.
馬丁路德就很驚訝.
老婆你為什麼穿成這樣.
誰過世了.
馬丁路德這麼愛主.
當然很生氣.
怎麼搞的上帝死了.
上帝怎麼可能死了.
他太太就笑著這樣說.
上帝都沒有死.
既然是這樣你為什麼這麼愁.
上帝都還在.

$^{601}$上帝仍然掌管他一定幫助你渡過難關.
你又何必.
自己抱著太多憂慮在身上.
馬丁路德於是.
謊言大悟.
他從上帝以來那種信心和喜樂.
其實也提醒我們.
陰無掛淚的關鍵.
不是我們身處的環境.
有沒有改善.
而是因為我們時刻.
和上帝結連.
事實上兄姐妹.
我們每一個人在世上都有自己.
或者說上帝.
要你面對的難關.
但是那些平安.
喜樂的人.
他們有一個秘訣.
怎樣轉化.
有沒有玩過拍照.
手機也好.
單反也好.
你拍照一定要做一個動作.
對焦.
每遇到一件事.
你就先對焦.
對準上帝.
看清楚一件事.
然後才按掣.
你拍出來的東西會怎樣.
就會清晰.
就會美好.
首先連結對準上帝的時候.
上帝還在.
沒有離開.
你就可以坦然面對你的逆境.
因為是上帝和你一起去面對.
並且上帝幫我們去面對.
事實上平安和喜樂.

$^{641}$就好像一對孖公仔.
平安就是我們知道上帝掌管著.
即使在風浪裡面.
我們可以安穩.
而喜樂就是我們享受.
其實上帝在風浪裡面.
和我們同在.
我們可以在風浪裡面.
前行.
繼續去認識探索這個世界.
在裡面事奉榮耀上帝.
再看最後的經文.
儘管.
剛才也說過.
保羅或是腓立比教會.
他們身處的環境.
對基督宗教是不友善的.
充滿很多的敵意.
甚至會有一些不合理.
不公道的事情發生在他們的生活裡面.
但保羅卻是提醒.
腓立比的信徒.
仍然是天父所創造.
所會有的世界.
上帝仍然.
賜福給世人.
人世間其實仍然.
有很多美善的事物.
是值得信徒.
來欣賞.
甚至乎學習.
所以保羅在第八節裡說.
弟兄們,我還有未盡的話.
凡是那些真實的,可敬的,公義的.
清潔的,可愛的.
有美名的,若有什麼得行.
若有什麼稱讚.
這些事,我們都要去思想.
但我想大家都有看新聞的習慣.
有沒有發覺其實.

$^{681}$這幾年我們看新聞的時候.
多了很多的感受.
可能特別是負面的感受.
為什麼我們身處的社會搞成這樣呢?.
為什麼世界搞成這樣呢?.
立立亂.
為什麼好像越來越糟糕.
身邊的人做事為什麼會做成這樣呢?.
以前不是這樣的.
為什麼人的道德操守好像越來越差呢?.
不知道你有沒有這些同感?.
最近看過一篇文章.
他說香港人其實.
有時對事,對人的態度.
其實負面多了,否定多了.
這篇文章說到.
沒有貶義,純粹是筆者的體會.
他說香港人有五個缺點.
其中一個不知道你認不認同.
他說很多的挑剔.
很多的吹毛求疵.
很多人喜歡在一些很少的事上找錯處.
只要發現對方跟自己的做法不同.
就會批評.
反對.
平時處事對人很mean.
即是很苛刻.
談的多,讚的少.
所以香港人都習慣性批評.
因為這樣的緣故.
大家都不敢表達自己.
因為一表達就被人公審.
所以變得很被動.
丁姐妹,無論你認不認同這種說法.
其實有一樣東西我們否定不了.
幾個圖都走不掉.
我們起落不了其中一個原因.
可能是因為我們對這個世界.
對身邊的人積累了很多的不滿.
積累了很多的不滿和紅素.

$^{721}$事實上我們不難去理解.
如果我們對這個世界對世人.
真的積累很多的負面看法的時候.
你很難去愛他.
你憎一個人你怎樣去愛他.
很難的.
如果我們積存了太多負面批評的看法的時候.
我們很難去代表耶穌祝福他.
我們更加很難在他們面前.
告訴他們耶穌愛他們.
所以保羅勉勵我們.
我們要思考身邊周遭的事物.
我們學習不要過分的極端.
我們要正面積極地認識探索這個世界.
這個天父所創造的世界.
來接觸身邊的人.
發掘他們.
欣賞他們.
其實有長處有微善的地方.
今早有沒有看新聞?.
或者昨晚有沒有看奧運?.
有啊.
我們香港的劍擊女子重劍的選手.
江維文.
拿了金牌.
如果有機會看他的歷史.
上一屆飲恨東京奧運.
但是怎樣堅毅地.
去跨過壓力.
跨過自己身體的創傷.
和退役的掙扎.
明明決賽對著東道主法國的選手.
輸了1比6.
怎樣咬牙關.
追到12平手.
到最後加時的時候.
一劍就定江山.
香港人這種毅力.
很值得我們欣賞.
他信不信主我們不知道.

$^{761}$但是這種美善的做人態度.
豈不是.
很值得我們去欣賞.
去認同.
甚至去學習.
事實上弟兄姊妹.
不要因為遇到這個世界上.
可能一些惡人惡事.
而懷疑上帝在這個世界的管理.
不要動搖上帝.
看顧我們的信心.
不要否定這個世界.
仍然有一些.
美善的地方.
更加不要因為你見到.
壞人一時得逞.
公義未能彰顯.
你就抱怨上帝.
覺得上帝不公道.
甚至乎你自己都不知不覺遠離上帝.
不再持守上帝的吩咐.
將自己.
由上帝的兒女變成上帝的仇敵.
所以保羅.
最後第九節.
無論這個世界變成怎樣.
無論身邊的人怎樣對你.
我們都要說好耶穌的故事.
我們都要跟保羅一樣.
你將你的目光轉頭上帝.
我們可以在保羅身上學習.
他所持守.
他怎樣去考慮.
他如何去承認.
他如何去持守.
他如何去效法基督.
我們如何去學習保羅效法基督.
保羅的意思其實是.
身處在一個混亂充滿風險的環境世代裡.
甚麼才是我們生活最大的保障.

$^{801}$就是你奮身相信耶穌.
沒有其他選擇.
怎樣奮身相信耶穌.
就是你捉緊他的吩咐.
你遵守他的吩咐.
為他作見證.
完成他交託給我們的使命.
保羅說上帝會將他的平安賜給我們.
他會與我們同在.
而有上帝同在的生命.
就一定是最安穩的人生.
所以弟姐妹其實.
喜樂平安不在於.
我們的生活是否一帆風順.
是否無風無浪.
在於我們有沒有捉緊上帝創造我們的目的和意義.
我們的生活其實.
好像一艘船.
有時遇到太巨大的風浪.
它真的需要怎樣.
它真的要停在避風港.
讓那裡獲得安穩.
可以遮風擋雨.
但一艘船可否永遠停留在避風港.
又不可以.
因為這不是它被造的目的.
船被造就是要出海.
就是要經歷風險.
冒險.
甚至有時會有損傷.
撞傷刮花.
然後又回到避風港.
回到船塢裡面維修.
維修完之後.
再出發.
去探索這個世界.
探索這個生命.
就是這樣不斷循環累積.
一些真實的歷練.
人就會漸漸長大.

$^{841}$其實我們被造的目的都是這樣.
我們要和世界結連.
要和別人建立關係.
幫助自己.
和別人在裡面更加認識.
上帝的旨意.
這就是我們的寧靜成長.
在過程裡面.
會有損傷的.
會有風險.
甚至乎你會卻步.
但耶穌允許.
祂和你同在.
祂會和你同行.
並且祂會有.
足夠的保護平安.
讓我們最後.
經歷一個得勝的人生.
到達終點.
好盼望這個系列.
讓我們讓孫逸宣說.
能夠幫助弟姊妹.
我們的靈命真的可以安穩在耶穌裡面.
一方面經歷祂的保護.
一方面經歷祂的醫治.
但同時.
我們也相信耶穌會堅持祂.
我們出去探索這個世界.
和你身邊的人.
建立一個健康的關係.
我們同心討論.
恩主啊,我們多謝你.
讓我們真的無論在什麼環境裡面.
讓我們都可以.
依靠耶穌你.
在主理裡面找到出路.
找到安息.
找到醫治.
我們再可以繼續出發.
去履行上帝所託付給我們的使命.

$^{881}$走入人群裡面.
去為你作見證.
與身邊的人的結連.
彼此的祝福.
上帝我們知道過程不容易.
正如你走進人間的時候.
你也會遇到敵黨.
我們也知道耶穌你已經為我們.
完成了這趟人生的之旅.
成為我們美好的示範.
也成為我們跟隨的一個榜樣.
以致我們真的可以.
跟著你來行的時候.
我們就知道.
我們可以走不全程.
求主你幫助我們.
幫助每一位弟兄姊妹.
無論我們在什麼光景裡面.
我們可能與人有厭曲.
或者與人關係有阻滯.
又或者我們在這個世界裡面.
遭受過傷害.
求主你也醫治.
讓我們有重新出發的勇氣.
去為你作見證.
我們這樣的禱告交託.
奉耶穌基督你的命來祈求.
\newpage



\section{耶利米書 20:7-18}
\label{sec:OxtvYBdmNrI}
\textbf{2021年6月13日崇拜直播｜陳冠成牧師｜從聖經看喜怒哀愁 - 告別憂愁｜耶利米書20\_7-18}
\newline
\newline
連結: \href{https://youtube.com/watch?v=OxtvYBdmNrI}{\texttt{https://youtube.com/watch?v=OxtvYBdmNrI}} ~~~~ 語音日期: 2024-08-06
\newline
\newline
\hyperref[sec:rzRD6n7AW0w]{\small{< < < PREV SERMON < < <}}
~
\hyperref[sec:index_chronic]{\small{[返順時目]}}
~
\hyperref[sec:Hufazqfe3oE]{\small{> > > NEXT SERMON > > >}}
\newline
\newline
耶利米書 20:7-18
\newline
\begin{longtable}{cl}
\hline
\hline
章節 & 經文 (和合本修訂版)\\
\hline
20:7 & \begin{tabularx}{0.7\textwidth}{X} 耶和華啊，你欺哄了我，我也被你欺哄了。你比我強，並且得勝。我終日成為笑柄，人人都戲弄我。 \end{tabularx} \\ \\ \relax
20:8 & \begin{tabularx}{0.7\textwidth}{X} 我每逢講話的時候，就哀嘆，我喊叫：「有暴力和毀滅！」因為耶和華的話終日成了我的凌辱和譏刺。 \end{tabularx} \\ \\ \relax
20:9 & \begin{tabularx}{0.7\textwidth}{X} 我若說：「我不再提耶和華，也不再奉他的名講論」，我心裡便覺得似乎有燒著的火悶在我骨中，我忍受不住，不能自禁。 \end{tabularx} \\ \\ \relax
20:10 & \begin{tabularx}{0.7\textwidth}{X} 我聽見許多的毀謗，四圍都是驚嚇；連我知己朋友都看著我跌倒：「告他吧，我們要告他！或者他被引誘，我們就能勝他，在他身上報仇。」 \end{tabularx} \\ \\ \relax
20:11 & \begin{tabularx}{0.7\textwidth}{X} 然而，耶和華與我同在，好像可怕的勇士。因此，迫害我的都絆跌，不能得勝；他們大大蒙羞，由於行事沒有智慧，必永遠受那不能忘懷的羞辱。 \end{tabularx} \\ \\ \relax
20:12 & \begin{tabularx}{0.7\textwidth}{X} 考驗義人、察看人肺腑心腸的萬軍之耶和華啊，求你使我得見你在他們身上報仇，因我已將我的案件向你稟明了。 \end{tabularx} \\ \\ \relax
20:13 & \begin{tabularx}{0.7\textwidth}{X} 你們要向耶和華唱歌！要讚美耶和華！因他救了窮人的性命脫離惡人的手。 \end{tabularx} \\ \\ \relax
20:14 & \begin{tabularx}{0.7\textwidth}{X} 願我出生的那日受詛咒！願我母親生我的那天不蒙福！ \end{tabularx} \\ \\ \relax
20:15 & \begin{tabularx}{0.7\textwidth}{X} 報信給我父親說「你得了兒子」，使我父親甚歡喜的，願那人受詛咒。 \end{tabularx} \\ \\ \relax
20:16 & \begin{tabularx}{0.7\textwidth}{X} 願那人像耶和華所傾覆而不憐惜的城鎮；願他早晨聽見哀聲，中午聽見吶喊； \end{tabularx} \\ \\ \relax
20:17 & \begin{tabularx}{0.7\textwidth}{X} 因他沒有在我未出胎就把我殺了，以致我母親成為我的墳墓，她卻一直懷著胎。 \end{tabularx} \\ \\ \relax
20:18 & \begin{tabularx}{0.7\textwidth}{X} 我為何出胎見勞碌愁苦，在羞愧中度盡我的年日呢？ \end{tabularx} \\ \\
[1ex]
\hline
\hline
\end{longtable}
$^{1}$讓我們繼續以聆聽聖道去敬拜我們的主.
首先我們由邱小英姐妹為我們讀出.
耶利米書第20章7到18節.
接著陳貫成牧師會為我們宣講.
他的題目是.
從聖經看喜 勞 哀 樂 之告別憂愁.
(教育).
今日頌道的經文是在.
耶利米書第20章7至18節.
記載在舊約聖經的1053至54頁.
耶利米書20章7節.
耶和華啊!你曾勸導我.
我也聽了你的勸導.
你給我有力量且勝了我.
我終日成為笑話.
人人都戲弄我.
我每逢趕弄的時候就發出哀聲.
我喊叫說:有強暴和毀滅.
因為耶和華的話終日成了我的靈欲.
既持.
我若說我不再提耶和華.
也不奉他的名講論.
我便心裡覺得似乎有燒著的火.
閉塞在我骨中.
我就含忍不住.
不能自禁.
我聽見了許多人的殘暴.
四圍都是驚嚇.
就算我知己的朋友也都窺探我.
怨我跌倒.
說:告他吧!我們也要告他.
或者他被引誘.
我們就能勝他.
在他身上報仇.
然而耶和華與我同在.
好像甚可怕的勇士.
因此逼迫我的必都半跌.
不能得勝.
我們必大大蒙羞.
就是受永不忘記的羞辱.

$^{41}$因為他們行事沒有智慧.
試驗二人.
察看人肺腑深長的萬軍之耶和華.
求你容我見你在他們身上報仇.
因我將我的案件向你稟明了.
你們要向耶和華唱歌.
讚美耶和華.
因祂救了窮人的性命脫離惡人的手.
願我生的那日受咒助.
願我母親產我的那日不蒙福.
給我父親報信說:你得了兒子.
使我父親甚歡喜的.
願那人受咒助.
願那人將耶和華所傾福而不後悔的誠邑.
願他早晨呈見愛聖.
享悟呈見立憾.
因他在我未出胎的時候不殺我.
使我母親成了我的墳墓.
胎兒就是常重大.
我為何出胎見奴隸受苦.
使我的連日因羞愧消滅呢?.
弟兄姊妹平安.
來到這個系列最後一講了.
今天我們講「壽」字.
「壽」字怎樣寫呢?.
「心」上面加個「秋」.
是「迎」聲字.
可能和秋天也有關係.
秋天令我們想起陣陣的秋風.
捲起那些樹葉.
或者草木開始凋零了.
很容易令人觸景深情.
想起一些比較暗淡或淒涼的畫面.
所以在「秋」字上加個「心」.
就是用來形容人的一種很鬱悶.
或者很傷感的情懷.
中國人也有很多這些作品.
可能大家也有聽過.
宋代有位詩人申氣擇.
你不認識他.

$^{81}$可能你也會聽過他其中一首著名的作品.
《少年不識「壽」之味》》.
愛上層樓 愛上層樓.
為父身馳 強說「壽」.
而今識盡「壽」之味.
欲說還憂 欲說還憂.
卻道天良好過秋.
申氣擇描寫自己年輕時入世未深.
其實不太真正體會到什麼是「壽」.
雖然他的詩詞在年輕時也用了「壽」字.
但可能到他人到中年.
對「壽」的理解比識入豐富了.
覺得以前說自己很「壽」.
其實也不算什麼.
比起從前.
其實自己不太認識什麼是「壽苦」之味.
現在真的開始明白了.
一方面可能因為「壽」多了.
另一方面是「壽」深了.
有些「壽水」甚至不知怎樣形容.
不知怎樣說.
或者不方便跟人說.
說來也沒用.
對事情沒有幫助 解決不了.
在欲言又止之間又止好怎樣?.
好像他這首詩最後一句.
卻道天良好過秋.
即是說很涼爽的秋天.
意思是今天說的話打完槍.
今天天氣不錯.
來表達他的「壽情」.
事實上中國人有很多以「壽苦」「憂愁」為題材的文學作品.
例如明天是什麼節?.
即使不能爬龍舟.
你也記得吃粽子.
明天是端午節.
紀念誰?.
屈原.
不知會否令你想起他也很「壽」.
他是春秋戰國時期的楚國士大夫.

$^{121}$他有一首著名的作品叫做「離疏」.
有否聽過?.
「離」最簡單 字面理解.
「離別」.
「疏」其實也是解「愁」.
「憂愁」.
「離疏」可以說是他試圖失意滿肚牢疏那段時間寫出來的.
我們大概對屈原有些背景理解.
他有一腔愛國的熱誠.
深受那時的楚懷王重用.
不過後來怎樣?.
通常都是這樣.
小人當道.
他被奸臣挑撥離間.
他雖然有一個很好的自覺想法.
不過無用無之地.
明明向皇帝陳明「國之利弊」.
卻不但不被人採納.
還要被人排斥.
接著皇帝開始對他不信任.
最後他被迫流放.
我們也記得他的情況.
更仇的是什麼?.
更仇的是他看著楚國衰敗.
最後被秦國消滅.
自己卻無法阻止.
內心有一種揮之不去的鬱悶.
或者那種怨憤.
慢慢就變成毒.
原來是毒的.
有毒的.
最終我們知道他選擇了.
投身了密羅江.
結束了自己的生命.
用死來告別內心的仇恨.
不過遠在地球的另一邊.
有一個比阿屈原早幾百年出道的以色列人.
就是我們今天看經文所說的耶利米.
耶利米先知.
他和屈原像不像?.

$^{161}$其實都像.
拼在一起.
兩個人都同病相憐.
兩個人可以說是.
面對著受到想死的處境.
首先他們都是服侍君王.
但很不幸忠言逆耳.
君王又不喜歡聽他們說.
其次就是他們雖然不受歡迎.
甚至有人去害他們.
攻擊他們.
但他們仍然熱愛他們所身處的國家.
總是肆意不捨.
苦口婆心地勸他們回頭.
第三就是他們多愁善感的類別.
都是憂國憂民.
所以他們有個共通點.
就是喜歡寫東西.
喜歡寫文學作品.
抒發心中的情懷.
剛才我們說屈原就有黎蘇.
耶利米有什麼作品?.
最簡單是我們今天看耶利米書.
還有一首出名的《愛歌》.
這些作品都是深深影響後世.
所以我們運用我們的想像力.
如果耶利米真的可以跨越時空.
和屈原見面.
大家聊天的時候.
你猜會怎樣?.
這兩個人相遇的時候.
大概交換卡片.
問問做姓何.
然後就開始大吐苦水.
大家可能就喝一杯.
不要太愁.
或者說一句.
我們真是相逢恨晚.
來我們再喝一杯.
希望可以借酒消愁.

$^{201}$不過你會看到.
雖然這兩個人的遭遇都相似.
但當他們面對.
他們仇苦的態度的時候.
有些不同.
是不是?.
弟兄姊妹.
所以我想.
不如趁著中秋節.
不,端午節就快來了.
中秋節未來.
下次再說.
我們一邊吃總應節.
但又一邊學習一下.
耶利米如何告別憂愁.
我們一起有個禱告.
我們交托.
聖靈願你與我們同在.
幫助我們.
讓我們真的能夠.
從上帝你的話語裡得到亮光.
得到生命的力量和方向.
讓我們可以在聽到的裡邊.
我們能夠被上帝感動.
讓我們真的能夠在生活裡.
從自己的教導裡活現出來.
多謝你聽我們禱告.
奉耶穌你得勝的名字.
來祈求.
阿們.
先看回耶利米.
為什麼他這麼仇.
我們要有一點背景.
因為他要向當時一群.
自認生活安穩.
覺得上帝很保守他的百姓.
宣告一個不中聽的訊息.
你們快點悔改回轉.
否則上帝就會審判我們.
上帝就會透過巴比倫.

$^{241}$來懲罰管教我們.
不中聽嗎?.
不難理解.
用今天的角度.
他在說危害國家安全的訊息.
當百姓都覺得.
我們明明生活得很好.
為什麼你發放這樣的訊息呢?.
於是全國上下反對他.
笑他 孤立他.
甚至你會看到.
稍後將他又拉又鎖.
逼他閉嘴.
不要再發放我們說的.
顛覆國家的言論.
但是受就受在.
上帝叫他堅持下去.
即使如此.
你堅持下去.
我們看一看.
耶利米先知蒙召的經過.
第一章十九至二十節.
應該在你們的程序表.
講到大綱內.
在這裡說.
上帝和耶利米先知說.
「看!我今日使你成為堅城鐵柱銅牆.
與傳帝和猶大的君王.
首領祭司.
並地上的眾民反對.
他們要攻擊你.
卻不能勝你.
因為我願你同在.
要拯救你.
這是耶和華說的」.
上帝不單止叫他堅持下去.
還允許了.
「我會與你同在.
使你變得堅固」.
有多堅固?.

$^{281}$如果上帝要使你變成堅城鐵柱銅牆.
你要面對的就不簡單了.
否則上帝也不需要給你強勁的保護.
上帝說會讓他變得堅固.
足夠抵抗全國上下的攻擊.
確保沒有人能夠打敗耶利米.
厲害嗎?.
厲害.
但不代表不能打傷耶利米.
明白嗎?.
打不倒他.
但不代表他不會受傷.
不會痛.
不代表耶利米的心靈不會痛苦.
不會疲乏.
正如我們在現實生活中都明白.
上帝確實會保守我們.
但不代表那些令你苦澀的際遇.
就會馬上變得甜.
對不對?.
所以我們看到耶利米每一天都要面對很多掙扎.
他要繼續宣講.
但沒有人聽.
甚至反對他.
他每天都活在矛盾中.
所以說真的挺受苦.
第七節我們看到.
剛才也讀過.
「耶和華,你曾勸導我,我也聽了你的勸導,你比我有能力,且勝了我,我終日成為笑話,人人都戲弄我,我每逢講論的時候就發出哀聲,.
我喊叫說:有強暴和毀滅,因為耶和華的話終日成了我的凌虐基詞.」.
第七節這裡和學本也做勸導.
其實也很客氣.
其實原本的意思是「欺騙如弄」.
意思是耶利米是在向上帝表達他的不滿.
上帝沒錯,我知道你是大能的上帝.
你的說話一定是對的.
我明白,理性上我明白.
但感受上他覺得什麼?.
上帝是覺得你在欺控我.
上帝是在誤導我.

$^{321}$我當日被你「講服」了.
「講贏」了.
所以去向百姓傳講審判的訊息.
但現在你看看上帝.
我每次向群眾大聲疾呼.
有強暴和毀滅要來.
他們當我是傻子.
他們不停地取笑我.
為什麼?.
因為我一直講,講了這麼久.
什麼事都沒有發生.
他們仍然很安全.
上帝沒有出手懲罰他們.
所有事都安然無恙.
所以百姓漸漸習慣了.
耶利米所講的審判訊息.
當它是什麼?.
是狼奶料.
每逢耶利米這樣講.
這個人又發瘋了.
又講狼奶料了.
於是耶利米心中有個鬱結.
就是上帝讓我傳講你的說話.
但你的說話反而令我成為.
所有人的笑柄.
糾不糾結?.
其實真的很糾結.
我明明只不過是忠心上帝.
難道這樣也有錯?.
為什麼所有人都要說我不對.
所有人都要針對我?.
很難受的.
其實耶利米知道真的很難受.
所以他有想過.
他也好像現代人.
工作這麼難受,怎麼辦?.
不做了,辭職吧.
他有想過.
第九節他有想過.
我不想再和上帝發聲.

$^{361}$甚至我投其所好.
和別人說平安,合心意的話就算了.
這樣大家又開心.
我又可以告別我內心的受煩.
但上帝批不批?.
上帝不批.
耶利米表示.
若他不宣講上帝的話語.
內心的痛苦就像有火.
在骨頭裡燃燒.
一直燒著你,但又出不來.
那種你又抓不到他.
不知道怎樣處理好那種痛苦.
即是不宣講.
不面對仇敵.
更加辛苦.
那慘了,不能走,那又怎樣?.
唯有繼續留在熱廚房工作.
但情況又一天一天沒有改善.
不會放月下吧?.
第十節這裡說了.
很多人誹謗我,告我.
甚至.
昔日和我拍肩膀.
昔日在Facebook裡friend我.
現在又怎樣?.
現在unfriend我了.
不單止,還搜集我的證據.
跟蹤我,監察我.
想戳我背,弄死我.
我們不會去到.
耶利米這麼艱難.
甚至人生受威脅的情況.
我相信.
但我們多多少少總會遇到.
進退失據.
左右為難的局面.
無論如何都好像沒有出路.
例如我們有時會問上帝.
「上帝,我真的很不喜歡現在的工作環境」.

$^{401}$「我真的很不喜歡每天面對很難受的人」.
「但你又不讓我走」.
「你又要我繼續留下來面對」.
並且情況越來越差.
上帝,我真的覺得我做什麼都沒有力.
很難撐下去.
對方不單止不領情.
還要伸我一腳.
上帝真的很仇.
很難受.
我每天雖然在服侍你.
但實際我每天都被人揍一身.
仇不仇?.
真的挺仇.
等姐妹當面對這種仇苦.
沒有出路的時候.
聖經又讓我們看到有什麼出路.
你見到剛才一開始我們說.
耶利米和屈原的際遇很相似.
但他卻沒有像屈原那樣.
選擇一聲「咚」.
在以色列就可以投去哪裡?.
沒有密羅岡就可以投去約旦河.
或者找加利利海.
他沒有這樣選擇.
投進約旦河或加利利海了結自己.
他選擇投身在剩下的最後防線.
剩下的最後信仰防線是什麼?.
他自己說「我還有上帝」.
「我還有上帝在我身邊」.
這就是他最後的防線.
第十一節他說「然而」.
背後意思是.
就算我人生一團糟也好.
我怎樣仇也好.
所有人都不喜歡我.
甚至攻擊我也好.
但我仍然確信耶和華你與我同在.
並且這不是他一句信仰口號.
他真的會運用這句宣告.

$^{441}$在他的處境裡.
他知道當耶和華與他同在時.
他懂得怎樣?.
有些事他會交給上帝.
他不會所有事都全部.
全部負責自己.
他懂得將自己困難的處境.
交給上帝處理.
包括第十二節.
上帝我搞不定這班仇敵.
甚至我被他們攻擊.
所以上帝我還是將他們.
交給你來對付.
一來我沒有能力面對他們.
二來上帝你也沒有吩咐我.
去對付我的仇敵.
上帝你只是吩咐我.
要繼續跟他們宣講.
勸他們回頭.
這是第一方面.
他將他自己搞不定的事.
上帝沒有吩咐他要做的事.
交給上帝.
第二就是第十三節.
他將自己的命.
交回上帝手中.
因為他知道.
當全國要攻擊他的時候.
你怎樣防都防不了.
他唯有.
他不會再自救.
他知道拯救是上帝的工作.
所以交回給上帝.
上帝我將我的命放回你手中.
我們看到耶利米有他面對.
仇苦或燒仇的秘訣.
不難理解.
其實就是將自己搞不定的難題.
交回上帝.
讓上帝按他的旨意.

$^{481}$按他的時間表來幫助你.
但要留意一點.
重點不是.
這樣我就可以脫難了.
從此可以舒服地生活.
重點不是這樣.
重點是幫自己更加專注服侍上帝.
教託是為了讓自己更專注服侍上帝.
事實上你看到.
耶利米才知道.
他將自己搞不定的事情.
交回上帝處理.
背後有一個意思.
上帝我不再逃避.
我願意留在廚房.
繼續履行上帝託付.
交給我要去處理的任務.
就是我要繼續.
在這群不聽話的百姓面前宣講.
這些是上帝吩咐我去處理的.
我要繼續堅持.
但有些事情我真的搞不定.
上帝也沒有吩咐我去處理.
我交回給上帝.
上帝你看著我.
必要時你出手幫我.
救我.
這就是他真真正正體會到.
耶和華與他同在.
然後應用在他的生活處境裡.
不難理解的弟兄姊妹.
新約我們也熟悉.
這個情況.
你還記得耶穌基督在.
赫西瑪利園的禱告嗎?.
當耶穌要面對.
十架苦杯的時候.
你留意看回福音書.
他一樣是極度的憂愁.
憂愁到要死.

$^{521}$於是他又怎樣?.
他也只有禱告這一招.
持續的禱告.
讓耶穌能夠保持警醒.
一次不夠就怎樣?.
兩次吧.
兩次不夠就三次.
經文沒有繼續形容.
如果還有時間.
可能還有第四第五次.
目的是否要爭取贏上帝?.
是否要打敗上帝.
讓他不用喝苦杯?.
他有這樣表達他的情感.
但你知道真正的目的是甚麼?.
真正的目的.
不是為自己求勝利.
不用釘十架.
耶穌持續向上帝禱告的目的.
是為了求自己打輸.
是為了求敗.
被誰打敗?.
上帝你打敗我.
然後我可以降服你的旨意.
真正的意思是這樣.
我們很多時候禱告交託.
就是要爭取贏上帝.
我要贏.
上帝你拿走我一些不舒服的東西.
但耶穌的祈禱是.
上帝你打敗我.
讓我可以繼續降服你.
照你的意思.
不照我的意思.
所以弟兄姊妹.
你會看到耶利米斯先知其實都是這樣.
透過他不斷的禱告交託.
讓自己持續獲得.
可以和仇虎共處的力量.
一次不夠兩次.

$^{561}$兩次不夠就三次.
你看到聖本耶利米斯.
有很多來來去去.
才知道他和上帝禱告交託的對話.
雖然情況沒有改善.
敵人仍然在.
困難沒有解決.
但很奇妙的就是.
這樣不斷禱告交託來來回回.
上帝的僕人就沒有離開.
繼續留在火場裡.
他可以有力量咬緊牙關.
直到他完成了.
上帝託付給他的任務.
就是這樣過渡了.
所以真的弟兄姊妹.
當面對那些令我們感到痛苦.
受煩的人或事的時候.
其實我們很自然.
第一個想法是.
上帝你快點幫我拿走他.
讓我可以脫苦海.
我們很想上帝聽我們的呼求.
將問題解決.
讓我們離開火場.
但很多時候我們都明白.
當我們越是這樣站立的時候.
有時會怎樣?.
反而越失望.
越痛苦.
越更加愁.
因為上帝有自己的時間表.
不會每次都按照我們的意願.
來行事.
我曾經聽過.
有個家長這樣分享.
他說我的子女.
開始步入青春期.
反叛了.
不聽我說了.

$^{601}$所以我每次都祈禱.
每天都祈禱.
求上帝幫助我的子女.
能夠聽教聽話.
當然他越期越愁.
越失望.
子女繼續頂撞他.
繼續不聽他說.
不過在這個時候.
他開始有些事情轉變了.
他開始發覺.
我不斷地求上帝.
令我的子女聽教聽話.
那我自己呢?.
我自己也是上帝的兒女.
那我是否也是聽教聽話呢?.
我有沒有自己帶頭做好榜樣.
成為一個聽教聽話的上帝兒女呢?.
就是這個轉念.
他開始明白.
有些事情是要放手的.
有些事情不要太執著.
因為上帝沒有吩咐他去擺平.
所以他開始選擇.
不再執意要子女聽教聽話.
反過來他開始嘗試放下身段.
我去聽下子女怎麼說.
聽下他們的心聲.
事情沒有改變.
問題沒有解決.
反叛的情況仍然繼續.
沒有馬上消失.
但他卻不需要帶著太多的愁水.
陪伴子女過渡了青春期.
所以頂子妹耶利米提醒我們.
人生有些事情.
明明上帝要你放手.
要你交托.
我們不要千方百計.
要將事情拿回來.

$^{641}$自己去處理.
自己去擺平.
相反.
如果有些事情.
上帝拜託你去處理.
你就不要拖延.
也不要迴避.
你想辦法堅持下去.
想辦法做好上帝交託你去處理的事情.
就像耶利米先知那樣.
這樣做不代表問題會馬上解決.
但是在一個仇苦的情況裡.
我們卻可以像耶利米那樣.
向上帝歌頌讚美.
就像第十三節那樣.
偶爾當我們懂得去交托.
將自己從苦毒釋放出來的時候.
你就可以向上帝歌頌讚美.
再看下去.
第十四至十八節.
這裡確實看到耶利米的情況不穩定.
剛剛第十三節才向上帝交托.
歌頌讚美完了.
轉眼下一節又怎樣?.
又去咒詛自己.
一個處於受損的人.
狀況不一定很穩定.
他咒詛自己的生日.
甚至咒詛跟父親報喜訊的人.
來發洩自己現在的狀況.
上帝我生不如死的狀況.
當然我們可能會問.
這樣會否得罪創造的上帝?.
會否不屬靈?.
畢竟你也是先知.
不過我們要留意的是.
無論我們有多不習慣.
這種信仰表達的方式.
在經文裡你有沒有看到.
上帝因為這樣責備耶利米?.

$^{681}$你為何要咒詛自己?.
不可以這樣.
做人要開心一點.
上帝沒有這樣回應.
所以當我們看到.
耶利米先知咒詛自己的時候.
我們要問的不是這樣合不合意.
而是應該問一個問題.
為何上帝會將.
耶利米最脆弱的一面.
保留在聖經.
讓我們今天看得到.
為何上帝要保留這一幕.
讓你和我看得到?.
這才是我們要問的問題.
弟兄姊妹,如果你還記得.
先知祂稱呼上帝是什麼?.
萬君之耶和華.
是給祂有能力,大能的上帝.
背後代表著先知祂確信.
上帝我是不需要因為要保護.
你的面子或保護上帝你的情感.
而讓自己虛假起來.
我不需要這樣做.
因為上帝是大能的上帝.
我也不需要用一些不相干的言語.
來隱藏自己內心的切膚之痛.
因為我知道上帝是你撐得住的.
所以你看到耶利米.
祂知道如果我繼續在上帝面前.
抑壓自己的仇苦.
其實是反而怎樣?.
太過有禮貌.
禮貌地將上帝推開了.
好像我們和別人互動.
「你沒事吧?」.
「明明有事的,不過你禮貌地」.
「OK啦」.
將人推開了一樣.
我們也懂得用這些技倆.

$^{721}$但我們也會用在上帝身上.
耶利米才知道如果他再在上帝面前.
隱瞞他這種切膚之痛的話.
其實是和上帝太過見外.
就好像我們在祈禱或崇拜的裡面.
無論你的內心有多麼鬱結愁悶.
總之我一埋眼就習慣了讚美感恩.
不論如何也說「Hallelujah」.
將所有正面的我懂得的詞彙都數出來的時候.
這樣我們是否真的和上帝交流?.
是否真的敬拜祂?.
其實可能不是.
就算我們禱告完.
其實和禱告前也沒有分別.
所以頂枝梅我們看到.
耶利米才知道他沒有虛假美化了他的人生.
他承認受苦,孤單,焦慮.
甚至絕望是真實的.
是他生命裡面的一部分.
他沒有去隱藏.
也沒有在上帝面前逃避.
他卻選擇在上帝面前承認自己是脆弱的.
上帝我真的受到不想做人.
我很希望我從來沒有來過這個世界.
我們也聽過.
甚至我們自己也想過.
我們明白耶利米才知道.
他只是想在上帝面前盛夏,怨夏.
反正他又不是真的想要他自盡.
其實又何妨?.
當然.
向上帝吐完苦水,發完牢騷.
問題仍然存在.
但很奇妙的是.
當我們這樣做的時候.
我們就會在上帝裡面.
支取到一些繼續活下去,撐下去的力量.
我相信我們也體會過.
有時有些事情是透過說出來,申訴出來.
事情沒有解決.

$^{761}$但卻有活下去的能量.
所以廣東話有一句說話很有趣.
叫做「吊頸都要透下氣」.
吊頸都要透下氣.
我和兒子一起溫習.
正在考這個成分試.
昨天和他操練數學.
操練到頭昏腦脹.
問他可否休息一下看電視.
我說「也好」.
「吊頸都要透下氣」.
「什麼來的?吊頸都要透下氣?」.
「吊頸都要透下氣」.
突然有否想過.
為什麼吊頸都要透下氣?.
吊頸明明是想怎樣?.
想死.
但透下氣呢?.
想生.
是想生還是想死?.
其實很有趣的黑色幽默.
表面看似乎是想死,求解脫.
實際上仍然很想求生.
仍然很想求活.
透一透氣為了繼續吊.
吊,但又不是想死.
所以懂得說吊頸都要透下氣的人.
他會明白人生一定是苦樂參半.
兩件事甚至可以同時間發生.
只不過不按比例出現在人生一些場景裡.
懂得吊頸都要透下氣的人.
他會在他很愁的時候.
他會找辦法怎樣?.
先休息一下.
做些令自己開心一點的事.
讓自己透下氣.
然後再問自己.
我想怎樣繼續好好地生活下去.
我曾經聽過一位太太的傾訴.
她說她的婚姻關係一般般.

$^{801}$和丈夫的關係很差.
她覺得很愁.
覺得丈夫經常對她不好.
她愁得想死.
她說如果可以,求上帝把她帶走.
不是把她丈夫帶走,是把自己帶走.
我沒有否定她的愁遂.
我沒有跟她說.
不要這樣想,基督徒應該開心一點.
沒事的,讓我和你丈夫說說就行了.
我沒有這樣說打氣的話.
我只是跟她說.
你真的可以和上帝這樣傾訴.
因為上帝是大能的上帝.
就算你不說,他也知道你在想什麼.
所以就算你這樣和他訴苦.
說自己真的愁得不想做人.
你放心,上帝定得順.
真的定得順.
當然,上帝應不應允你又是另一回事.
這是上帝的主權.
關鍵是只要你不要堅持.
你現在的想法是絕對正確.
你願意仍然開放自己.
隨時讓上帝糾正你,更新你.
使你可以繼續有力量生活下去.
這樣就行了.
所以每次聽太太分享的時候.
我都給她機會先嘆一口.
先怨一口,甚至先哭一口.
處理了她情感需要的時候.
就再慢慢幫她拆解婚姻裡的相處問題.
不難理解的時候,弟兄姊妹這個道理.
很多年前有套電影.
可能大家都有看過.
叫做《Inside Out》.
有沒有看過?.
香港是做什麼的?.
是不是玩轉老朋友?.
玩轉老朋友.

$^{841}$音樂其實很好.
裡面有個角色叫做阿冰冰.
好像一隻象,但又有幾片鬚.
他有一次很失意.
他身邊有兩個朋友.
一個叫做快樂的樂.
一個叫做仇.
就是我們今天的主題.
優仇的仇.
樂想盡辦法來哄他.
擠他,扮鬼臉.
叫他玩遊戲.
一直跟他說.
不要這麼不開心.
不用這麼仇.
結果有沒有效?.
沒有效.
樂沒辦法來哄冰冰.
相反的仇又如何?.
仇其實沒什麼做.
他只是走過去.
運用一點我們今天輔導技巧的同理心.
就是聽聽他說.
讓冰冰伸一下,掃一下苦水.
給他空間哭一場.
哭完之後又如何?.
冰冰自己說.
「OK的,我沒事了」.
就是這樣.
於是樂很不明白.
他問仇.
「你是怎樣弄好他?」.
然後仇說.
「不知道,我看他傷心」.
「就聽聽他說」.
就是這樣.
所以其實我們也很羨慕樂觀的性格.
但現實讓我們看到.
單靠樂觀的性格沒辦法.
應對人生種種複雜問題.

$^{881}$有時真的需要大哭一場.
滲一下,怨一下.
將自己的情緒適當地表達出來.
反而可以幫我們正面地消愁.
漸漸從人生的低谷裡走出來.
事實上,我們沒有天生要做小丑.
沒必要逼自己時時刻刻都笑面迎人.
明明很愁的時候.
仍然勉強自己要討好全世界.
就好像今天這段經文提醒我們.
上帝真的很愛我們.
對不對?.
但上帝又可能會安排一些.
永遠不喜歡你的人在你身邊出現.
對不對?.
無論你怎樣做.
他都可能看你不順眼.
對不對?.
沒有原因.
可能你很努力.
但他仍然不喜歡你.
就好像耶利米那樣.
所以這段經文也提醒我們.
不要太在意你要討世上每一個人的喜悅.
不要太在意他對你好還是不好.
當然,聖經教導我們要盡力與人和睦.
但卻不需要我們對別人.
對我們的看法過度耿耿於懷.
事實上,我們很明白.
順得歌情又如何?.
失數意.
順得人情又如何?.
失上帝的意.
越是想千方百計討好全世界.
有時就越會增添自己的壽水.
所以,耶利米給了我們一個示範.
就是你專一的投靠上帝.
首先專心做好上帝已經交託給你辦的事情.
另一方面,你覺得真的做不來的事.
就是怎樣?.

$^{921}$透過你的禱告.
行切的禱告.
交回給上帝.
等他幫你去處理.
最後,和大家讀一段經文.
耶利米愛歌第三章32至33節.
我們重讀,好嗎?.
預備起.
「全天水使人有仇,何能要小他?.
天官的前代大連民,.
一一壓定德金三師的仇父,使人有仇」.
好,謝謝.
很真實寫實的一段經文.
上帝容許憂愁出現在你的生活裡.
已經說了.
但是,他仍然會用他的慈愛憐憫來幫助我們.
因為他不甘心我們永遠活在憂愁痛苦的裡面.
他會幫助我們.
所以,就好像今天這段經文.
當我們感到失意,憂愁的時候.
你隨時來到上帝的面前交托.
你謹記上帝是我們的父.
不要收起自己,孤立自己.
當你真的受不了的時候.
放膽向上帝滲,哭,甚至怨.
他受得起,頂得住.
你亦放心.
他會醫治我們.
他會保守我們.
繼續堅持走在上帝喜悅的道路裡.
我們同心禱告.
上帝藉著今天這段經文.
幫助我們再一次.
投靠在你的懷裡.
讓我們明白當上帝與我們同在的時候.
我們總得將人生一些我們無法擺平的事情.
放手交給上帝.
亦讓我們有更專注,更多的忍耐.
去忠心做好上帝在我們人生裡.
已經安排,定義要我們辦妥的事情.

$^{961}$讓我們在禱告交托過程裡.
我們不逃避,不離開.
我們能夠繼續堅持.
來土上帝的喜悅.
直到完成上帝的吩咐.
上帝亦讓我們能夠在你面前赤露闖開.
讓我們願意將我們最脆弱的一面.
向上帝陳明.
事實上上帝你知道我們的心思意念.
你明白我們.
但你亦同時願意希望我們透過分享.
來與你聯繫.
讓我們在過程裡.
可以得到醫治,得到幫助.
讓我們可以更有力量的.
繼續為上帝你活下去.
特別我將每一位弟兄姊妹都交托給你.
或許我們人生都會面對.
面對著不同程度的難處.
一些困局.
我們沒有辦法離開.
但亦沒有辦法解決.
我們在裡面經歷著受苦.
上帝你知道.
你容許我們在這個狀況.
但你亦知道我們不會永遠都是這樣.
讓我們看見從你而來的盼望.
讓我們看到你怎樣幫助我們.
讓我們總是知道你與我們同在.
以致我們真的可以.
能夠將我們的壽水.
按著上帝你的旨意.
擺在一個適當的位置.
讓我們可以繼續有力量去事奉你.
與逆境同行.
多謝上帝你.
讓我們有寶貴的話語.
願你聽我們在你面前的禱告.
奉基督耶穌你特性的命旨來祈求.
阿們.

\newpage



\section{民數記 12:1-15}
\label{sec:Hufazqfe3oE}
\textbf{2024年8月4日崇拜直播｜陳冠成牧師｜正視你的情緒－嫉妒惹的禍｜民數記12\_1-15}
\newline
\newline
連結: \href{https://youtube.com/watch?v=Hufazqfe3oE}{\texttt{https://youtube.com/watch?v=Hufazqfe3oE}} ~~~~ 語音日期: 2024-08-06
\newline
\newline
\hyperref[sec:OxtvYBdmNrI]{\small{< < < PREV SERMON < < <}}
~
\hyperref[sec:index_chronic]{\small{[返順時目]}}
~
\hyperref[sec:0SAEU456WxE]{\small{> > > NEXT SERMON > > >}}
\newline
\newline
民數記 12:1-15
\newline
\begin{longtable}{cl}
\hline
\hline
章節 & 經文 (和合本修訂版)\\
\hline
12:1 & \begin{tabularx}{0.7\textwidth}{X} 摩西娶了古實女子為妻。米利暗和亞倫因他娶了古實女子就批評他， \end{tabularx} \\ \\ \relax
12:2 & \begin{tabularx}{0.7\textwidth}{X} 他們說：「難道耶和華只與摩西說話嗎？他不也與我們說話嗎？」耶和華聽見了。 \end{tabularx} \\ \\ \relax
12:3 & \begin{tabularx}{0.7\textwidth}{X} 摩西為人極其謙和，勝過地面上的任何人。 \end{tabularx} \\ \\ \relax
12:4 & \begin{tabularx}{0.7\textwidth}{X} 忽然，耶和華對摩西、亞倫和米利暗說：「你們三個人都出來，到會幕這裡。」他們三個人就出來了。 \end{tabularx} \\ \\ \relax
12:5 & \begin{tabularx}{0.7\textwidth}{X} 耶和華在雲柱中降臨，停在會幕門口，叫亞倫和米利暗。二人就出來， \end{tabularx} \\ \\ \relax
12:6 & \begin{tabularx}{0.7\textwidth}{X} 耶和華說：「你們要聽我的話：你們中間若有先知，我－耶和華必在異象中向他顯現，在夢中與他說話； \end{tabularx} \\ \\ \relax
12:7 & \begin{tabularx}{0.7\textwidth}{X} 但我的僕人摩西不是這樣，他在我全家是盡忠的。 \end{tabularx} \\ \\ \relax
12:8 & \begin{tabularx}{0.7\textwidth}{X} 我與他面對面說話，清清楚楚，不用謎語，他甚至看見我的形像。你們為何批評我的僕人摩西而不懼怕呢？」 \end{tabularx} \\ \\ \relax
12:9 & \begin{tabularx}{0.7\textwidth}{X} 耶和華向他們怒氣發作，就離開了。 \end{tabularx} \\ \\ \relax
12:10 & \begin{tabularx}{0.7\textwidth}{X} 當雲彩從帳幕上離開時，看哪，米利暗長了痲瘋，像雪那麼白。亞倫轉向米利暗，看哪，她長了痲瘋。 \end{tabularx} \\ \\ \relax
12:11 & \begin{tabularx}{0.7\textwidth}{X} 亞倫對摩西說：「我主啊，求你不要因我們愚昧，因我們犯罪，就將這罪加在我們身上。 \end{tabularx} \\ \\ \relax
12:12 & \begin{tabularx}{0.7\textwidth}{X} 求你不要使她像那一出母腹、肉已侵蝕了一半的死胎。」 \end{tabularx} \\ \\ \relax
12:13 & \begin{tabularx}{0.7\textwidth}{X} 於是摩西哀求耶和華說：「神啊，求你醫治她！」 \end{tabularx} \\ \\ \relax
12:14 & \begin{tabularx}{0.7\textwidth}{X} 耶和華對摩西說：「她父親若吐唾沫在她臉上，她豈不蒙羞七天嗎？現在要把她隔離在營外七天，然後才領她回來。」 \end{tabularx} \\ \\ \relax
12:15 & \begin{tabularx}{0.7\textwidth}{X} 於是米利暗被隔離在營外七天；百姓沒有起程，直等到米利暗回來。 \end{tabularx} \\ \\
[1ex]
\hline
\hline
\end{longtable}
$^{1}$今天所讀經課是在民數記第十二章.
(文字).
他們三人就出來了.
(文字).
趙阿倫,黃米利安就出來了.
(文字).
弟兄姊妹平安.
來到八至十月.
我們會開始一個新的講壇系列.
大家在滴注表可能都看到.
「正是你的情緒」.
我們希望透過.
舊約聖經一些人物的生命故事.
來解構他們在這些事件中的情感反應.
以及上帝如何介入在裡面.
從而培育我們可以從不同的角度.
敏銳我們自己.
或者身邊的人的情緒和感受.
幫助我們和上帝.
和我們身邊的人建立一個更加健康的穩妥關係.
今天首先跟大家講.
剛才所讀的一段經文.
文蘇記的十二章.
有關於我們一個.
嫉妒的情緒反應.
經文剛才一開始已經說了.
米利暗阿倫因為摩西娶了一位古實的女子.
就向他毀謗.
首先我們要了解這三個人的關係.
大家可能都清楚.
米利暗是大家姐.
阿倫是哥哥.
而摩西是弟弟.
為什麼會因為這件事來毀謗摩西呢?.
首先我們要了解.
裡面提到古實女子的身份到底是什麼.
基本上有兩個可能的解釋.
第一個.
古實就是當時埃及南部一個地方.
換句話來說摩西除了他妻子.

$^{41}$他的元配.
米顛人西波拉之外.
又另娶了一位外邦女子.
於是就惹來哥哥阿倫和姐姐米利暗說她的壞話.
這是其中一個可能解釋.
另一個更加合理的.
可能是參考哈巴谷書第三章七節.
聖經學者透過這次經文.
他們理解到其實.
古實也就是米顛同一個地方.
古實就是米顛的意思.
換句話來說.
古實女子所說的就是摩西唯一的太太.
她的元配米利暗.
事實上你會看到.
為什麼這麼奇怪.
明明摩西很久以前已經娶了這個太太.
但摩西什麼時候娶了米利暗.
在他打死了人在埃及.
逃走的時候逃亡在米顛的曠野.
就是遇到他這位妻子米利暗.
說的是《出埃及記》第二章.
但現在來到《民數記》十二章.
為什麼隔了這麼久才拿這件事來去誹謗摩西.
我們就試試解構裡面到底發生什麼事.
首先我們先了解.
原來在《出埃及記》第十八章.
他說摩西之前曾經打發他的太太和兩個兒子.
你們先回到米顛.
到他們出埃及以色列人.
及至到達曠野的時候.
摩西的岳父葉蒂羅才帶回西波拉.
就是摩西的太太.
和她的兩個兒子來到西乃山和摩西會合.
而在那個時候.
西波拉也正式成為以色列人的一份子.
連帶他的外家也得到摩西的重用.
首先我們有印象.
如果你有看《出埃及記》.
摩西不是很苦惱怎樣去管理這麼大群人.

$^{81}$於是他的岳父葉蒂羅就為他出謀獻策.
你將他們分成什麼.
千夫長,百夫長,五十夫長,十夫長.
這樣就容易管理.
不用一次過一對幾百萬人.
省下很多功夫.
建議他這樣去管理以色列人.
摩西也接納了.
而另一方面我們又看到.
來到近一點.
《民數記》第十章.
摩西看中了他的內卿.
即是米利安的哥哥.
一個叫荷巴的人.
他看中了這個人有在曠野紮營的經驗.
於是他又邀請了內卿荷巴.
你和我們一起同行.
和以色列人同行.
我答應我會厚待你.
你看到摩西對他妻子外家的重視.
使得他和他自己的原生家庭.
和米利安姐姐和哥哥亞倫的關係出現了變化.
你要留意.
事實上米利安在《出埃及記》中.
他不單是摩西的姐姐.
還要是救命恩人.
摩西出生的時候.
埃及要殺死所有男嬰.
他的姐姐米利安出了一個辦法.
把箱子裡還在保存的摩西.
放進河裡一直盯著.
結果摩西的嬰兒.
就在河裡一直留在法老的女兒身上.
於是米利安就想收養男嬰.
於是就這樣救回摩西的一命.
另一方面你也知道.
在過紅海的事件.
米利安和摩西並肩作戰.
《出埃及記》第十五章.
米利安作了一首詩歌.

$^{121}$帶動一群婦女擊鼓.
鼓舞以色列人的士氣.
不過現在摩西對外家的重視.
似乎被米利安覺得有什麼想法.
「咦!弟弟摩西你好像現在.
對家人反而沒有對外人那麼親密」.
「弟弟你好像站在古實女子身上.
其實她是外人來的.
你站在她那邊.
站在她外家那邊.
好像多過站在自己人那邊」.
有沒有似曾相識的場面?.
有沒有看粵語長片?.
總有些情況父母出來.
「如果不是父母含辛茹苦養大你」.
下一句是什麼?.
「你有今天」.
「現在娶了老婆又怎樣?」.
「忘記自己家了」.
「忘記自己家人那份恩情了」.
「弟弟如果不是當日姐姐救你一命.
你會有今天做到以色列人的領袖」.
有點酸酸的感覺.
這是其中一個.
為什麼他們會借摩西娶米利安這件事.
來毀謗摩西.
而另一個原因.
其實也和上文.
今天我們沒有看.
但你可以回去看看.
《民數記》的十一章其實有很大的關係.
裡面記載到當時百姓在曠野都漂流了一陣子.
他們發覺每天都吃什麼?.
每天都吃馬欖.
吃到口裡有點怪怪的.
於是又想起埃及.
我們當日日子多好.
又有肉吃.
拿些蒜炒一下.
他覺得很香.

$^{161}$覺得不開心.
於是埋怨摩西.
摩西覺得受不了.
開始想劈炮不幹.
上帝就出手幫他解決百姓沒有肉吃的問題.
裡面記載了在秩序表說到大綱.
《民數記》十一章十六至十七節.
這裡簡單說.
上帝為了解決摩西的仇苦.
他就叫摩西.
你在你當中長老那裡找七十個人.
到時我會和你說話.
上帝會和摩西說話.
亦都將降在摩西身上的靈.
分給這七十個人.
為了幫摩西分擔.
他在大靈上的辛苦.
那個重擔.
但是有沒有找米利安和亞倫?.
沒有.
沒有找他們兩個.
所以透過這兩件事件.
我們看到.
如果西波拉的出現.
是令到米利安覺得自己在家庭的地位有所動搖.
令到他事奉的感覺不良好.
摩西和七十個長老.
被上帝的靈感動.
這一份屬靈經歷.
亦都同樣挑戰著.
米利安的屬靈地位.
事實上.
她不單單只是摩西的姐姐.
救回摩西的命.
她本身是什麼?.
是以色列人的女善知.
上帝都會找她說話.
她都是上帝其中一個的代言人.
但是為什麼?.
上帝只選擇摩西.

$^{201}$只感動那七十個長老.
這次沒有如我.
我都是女善知.
為什麼上帝有一點大勢超?.
為什麼上帝親疏有別?.
所以你留意到.
她真正投訴的原因是什麼?.
為什麼上帝只和摩西說話.
不和我們和亞倫說話?.
這個才是更加核心的問題.
米利安覺得.
我好像不再受重視了.
上帝不再選擇我.
不再用我.
這也是她誹謗摩西的原因.
覺得自己的地位.
似乎也金飛色比.
不再受上帝或摩西的重用.
我們說事奉的舞台.
不再屬於自己.
事奉的大光燈.
似乎不再照在自己身上.
於是對上帝對摩西.
產生了不滿.
相對亞倫的地位.
有沒有感受影響?.
亞倫的地位好一點.
為什麼?.
她還是以色列人.
唯一的大祭司.
只有她一個.
她的地位仍然崇高尊貴.
上帝事實上也曾經.
起碼有一次.
直接和亞倫說話.
但是米利安沒有.
所以你會看到.
沒錯,經文雖然這裡和本易.
亞倫也有份和姐姐一起針對摩西.
但是毀謗這個字的源源.

$^{241}$是一個陰性的單數.
即是說誰才是主犯?.
米利安.
他是整件事的主謀.
亞倫只不過是.
附和的.
共犯來的,同謀.
所以也解釋到.
為什麼上帝之後.
責備亞倫.
沒有米利安那麼重仇.
來到這裡.
所以我們看到.
米利安之所以攻擊摩西的動機.
其實是出於什麼?.
是否出於內心的正義感?.
不是.
其實是出於什麼?.
嫉妒.
上帝只選擇摩西.
不選擇我.
廣東話吃不到葡萄.
就一定是酸的.
就一定的.
事實上你會留意到.
嫉妒很多時候也發生在.
我們人際關係裡面.
很多的矛盾.
很多的衝突.
也都是源於嫉妒引起.
事實上你看到.
聖經裡面記載.
第一宗的謀殺案.
也是嫉妒闖出來的禍.
看到嗎?.
該忍殺也罷.
原因是什麼?.
上帝為什麼看中弟弟的共物.
看不中我的共物.
於是乎.

$^{281}$接納不到上帝這個決定.
就妒忌弟弟.
並且將他殺死.
接下來你會看到.
聖經所揭示的.
有不少.
和嫉妒有關的故事.
都是發生在哪裡?.
在我們家裡.
賽世記記載了很多.
最經典就是.
弱勢的哥哥們.
因為不看得起弱勢.
為什麼爸爸,哥哥對你這麼好?.
只有你有這件彩衣.
我們沒有.
於是乎就對摩西.
產生了這份厭惡感.
弟弟來的.
親弟弟.
但卻是一個很親密的仇人.
是不是很吊詭?.
在你家裡最親密的那個.
有時不知為何.
會成為你的仇人.
事實上.
嫉妒這件事.
如何產生呢?.
有些心理研究告訴我們.
其實嫉妒的出現.
通常一定有四個條件.
你符合這四個條件.
基本上你就會有嫉妒.
聽聽吧.
第一個就是.
你嫉妒的對象.
他的社會地位.
應該和你差不多.
什麼意思呢?.
你會不會嫉妒.

$^{321}$英女皇的皇室生活?.
會不會?.
會不會嫉妒.
現在英王查理斯.
為何你每天都.
有這麼奢華的生活?.
會不會?.
不會的.
因為你和他的地位是怎樣?.
差太遠了.
除非你也是皇室貴族.
你才有可能嫉妒.
和你身份相似的人.
這是第一點.
第二點就是說.
通常說的.
嫉妒的對象.
他的工作或生活的場景.
和你是相近.
甚至乎一樣的.
又舉個例子.
除非你是從事拍戲的.
甚至乎你自己.
也是擔正主角的.
否則大多數.
你不會嫉妒.
王子華一套爆紅的電影.
你不會嫉妒.
除非你是另一個人的粉絲.
你不會嫉妒.
眼紅他的成就.
但是如果兩個同事.
在同一個辦公室.
容易嫉妒嗎?.
就容易了.
大家的工作場景很相似.
會因為工作表現.
或上司對他的評價.
而不妥對方.
嫉妒對方.

$^{361}$這是第二個條件.
第三個就是.
無論那件事.
是因什麼事情嫉妒也好.
那件事通常都是.
嫉妒的人.
他沒有能力自己得到.
或者做得到.
舉米利暗的例子.
是誰選擇摩西的?.
上帝嘛.
上帝不侵米利暗玩.
米利暗有什麼福份?.
死了也沒用.
上帝不侵米利暗玩.
不選擇你.
你無論如何也不會靠你的努力.
爭取到.
第三個條件.
最後一個條件就是.
嫉妒的人通常會覺得.
對方你現在擁有的.
其實不是你應得的.
簡單來說.
你以為你很厲害嗎?.
只不過你幸運.
上帝選擇你.
不選擇我.
是不是有點似曾相識?.
你符合這幾個條件.
很大機會你會有.
吃不到葡萄是酸的質度感覺.
我還記得在我初信成長.
在母會的時候.
有一次.
被牧者邀請.
在崇拜裡面.
去講成長見證.
很短的.
五分鐘十分鐘左右.

$^{401}$講一下你怎樣信主.
信了主之後.
怎樣在教會裡面.
團契事奉.
這樣去成長.
過了不久.
有位弟兄姊妹打電話給我.
她說.
我不值得你.
那次這樣.
原因是什麼呢?.
其實她和我.
差不多時間初信.
她也很努力在教會成長.
她也熱心事奉.
但問題是什麼?.
為什麼找我不找她?.
我又沒得答.
我沒得答.
因為不是我爭回來.
我沒得答.
但我很感恩.
我意識到當她肯打電話給我.
講這件事的背後原因是什麼?.
她消化到.
她消化到才找我.
她消化不到.
如果還沒有.
還卡住她的質度感覺.
她不能找我來表白.
我感恩她可以轉化到.
她將裡面的質度化為.
其實是欣賞.
也沒有否定.
她自己在上帝面前的努力.
她最後決定轉化這份感覺.
成為自己繼續努力的原動力.
我為她感恩.
因為我拆解不到.
我只能夠聽.

$^{441}$願意坦誠地分享.
也很感恩上帝在我們中間.
化解了這個嚴規.
事實上之後我們繼續合作.
各位在團契裡.
在全科任務裡繼續合作.
並且這位弟兄姊妹.
她之後做了牧者.
事實上弟兄姊妹.
你會看到質度的感覺爽不爽?.
我問大家.
你有沒有試過質度別人?.
爽不爽?.
其實很特別的.
質度的一種感覺是不爽的.
但是你有沒有留意.
其實質度引發出來的後果.
不一定是壞的.
有時可以是好的.
如果我們身邊有人比我們厲害.
勝過我們.
比我們優秀.
其實我們有兩個方法去面對.
我們有兩個方法去消解.
彼此之間的距離.
第一個是什麼?.
將他拉下來.
抹黑他.
唱衰他.
就好像米利和亞倫那樣.
我誹謗他.
你厲害很厲害嗎?.
我唱衰你,你拉你下來.
那你就變成和我同一輩子.
這是其中一個.
另一個是如何消除彼此之間的差距.
就是我認同你的優秀.
你的強項.
不過我也有我的亮點.
我提升我自己.

$^{481}$我努力我自己.
這是第二個方法.
請問你會選擇哪種方式.
來處理我們一些質度的感受?.
事實上最近大家一定有看奧運.
你有沒有發覺這兩種做法.
其實經常出現在頒獎的情況裡.
金牌永遠只有多少個?.
一個.
只有一個.
但是拿金牌和銀牌的兩位選手.
或者銅牌的兩位選手.
其實是否真的差很遠?.
不是.
你看游泳就知道.
金牌銀牌銅牌.
差多少秒?.
零點不知道多少秒.
其實大家一樣的水平.
不過金牌只有一個.
但很多時候你會見到.
拿不到金牌.
有些人就會替他.
替不了葡萄.
他很希望在網上.
把拿了金牌的人拉下來.
覺得不公道.
甚至裁判,黑哨.
這樣出來.
你會見到有人選擇.
質度如何回應質度.
將對手拉下來.
毀謗他和他身邊的人.
而另一個呢?.
另一個就是相反.
你會見到有些選手.
即使他拿不到金牌.
拿不到銅牌.
甚至沒有獎牌.
但是怎樣?.

$^{521}$他沒有否定自己的努力.
因著這個緣故.
我繼續做好自己就可以了.
拿不到金牌.
沒有被揀選.
不代表他比別人差.
事實上弟兄姊妹.
很值得我們去思考.
有時面對質度這種感覺.
避免不了.
不過與其我們去葡萄別人.
倒不如怎樣?.
中國人講四個字.
見賢思齊.
你去提升改善你自己.
正面處理你內裡的一些質度.
這樣就更加健康.
更加蒙上帝月亮.
這是經文提醒我們.
我們繼續看事件如何發展.
第六至八節你會看到.
其實美利安不單止毀謗摩西.
其實背後是挑戰上帝的安排.
既然是挑戰上帝.
上帝自然會出手.
你看見上帝的出現.
不是覺得對待你們不好.
所以補償他們.
令他們有與摩西同等待遇.
上帝沒有這樣做.
上帝反而很清楚.
一出場就說.
他有權用不同的方式.
與不同的人說話.
上帝很清楚在第六節說.
他與先知的交流方式.
基本上是怎樣?.
透過異象.
讓人感受到.
甚至見到他在這裡.

$^{561}$他的同在.
又或者透過異夢.
讓人聽見上帝所說的話.
這兩個啟示方法.
基本上上帝說適用於一般代言人.
不過上帝之後說重點.
他與摩西的說話方式就不同.
完全不一樣.
第七節你會看到.
上帝突顯摩西在他面前的獨特性.
雖然摩西也是先知.
但上帝刻意稱他為.
我的僕人.
你留意這個稱謂在舊約裡不多.
或者應該這樣說.
大部分稱呼我的僕人.
主要落在摩西身上.
有少部分落在大衛.
再少部分落在約書亞.
摩西佔了很大的比重.
上帝不單稱這是我的僕人.
還要稱讚他.
他在上帝的全家.
在以色列的帶領面前.
是最忠心,最忠誠.
因為這個原因.
上帝給他一個很特殊的身份.
就是他與摩西溝通的方式.
與一般的先知不同.
上帝說.
我與摩西是面對面明講.
不用猜謎.
不用透過異象和異夢.
首先一方面反映.
摩西真的有這個實際需要.
如果留意摩西帶領的.
以色列人數有多少.
可能幾十萬.
如果看文述記錄.
打仗的南丁也有六十萬.

$^{601}$你當成兩倍,三倍.
過百萬人.
摩西要下達一個指令.
可不可以每次都猜謎語.
不可以.
指令一定要很清晰.
因為要傳達下去.
確保以色列人明白上帝的話語.
所以上帝按這個實際需要.
不需要你猜測.
不需要你擔憂.
我直接告訴你.
要怎樣做.
然後你直接告訴以色列人.
他們要怎樣做.
就符合上帝的心意.
這是實際需要.
所以為什麼面對面要明講.
不需要猜謎語.
另一方面.
上帝直接向摩西啟示.
直接跟他說話.
其實是反映兩個人的關係.
關係是非比尋常.
遠勝於他和其他先知的關係.
上帝表示獨有摩西.
看到他的形象.
這裡記載在.
《出埃及記》三十二.
三十三,三十四章.
不說了.
上帝說無論在視覺上.
聽覺上.
摩西對神的認識.
是遠勝拋離一切的先知.
你會看到.
顯示摩西的身份檔次.
是高不止一個檔.
更重要是強調.
上帝有權這樣做.

$^{641}$上帝有權這樣揀選對待摩西.
所以上帝說完之後.
責備米利暗和亞倫.
為什麼你們不識死.
敢挑戰我的決定.
我的安排.
你誹謗我揀選的僕人摩西.
其實是挑戰上帝的安排.
和權威.
不怕上帝.
這是上帝的表示.
你看到上帝對米利暗.
對摩西的回應.
引起我們一個很重要的反省.
弟兄姊妹.
上帝對人可不可以親疏有別.
可不可以.
我們聽到親疏有別.
覺得很不舒服.
你看到米利暗和亞倫的答案.
上帝不可以對人親疏有別.
你對摩西是這樣.
你對我和亞倫都一定要這樣.
這樣就為之公平.
因為這樣這種心態.
很多人特別不信上帝的人都會有.
他們覺得上帝的絕對主權這麼霸道.
很不公道.
於是對上帝這個絕對主權提出質疑.
不過弟兄姊妹.
我們要想深一層.
公平到底是否上帝做事的首先原則.
公平是否上帝行事.
首先一個最重要的原則.
你想想.
我舉個例子.
明明是我們世人犯了罪.
但誰幫我們埋單.
誰?.
主耶穌.

$^{681}$公不公平?.
不公平.
我們犯罪當然是我們自己埋單.
這樣就是公平.
為何要一個無罪的耶穌.
替我們埋單贖罪受死.
不是因為公平.
為何上帝堅持這樣做.
是因為有些事凌駕在公平之上.
是恩典.
是上帝的慈愛.
是祂的主權.
上帝給我們恩典,慈愛.
原來本來是不公平的事情.
因為我們不配領受.
如果公平的話.
上帝又如何?.
如羅馬書所說.
人人都有一死.
就是公平.
所以弟兄姊妹.
從這個角度我們會很明白.
可能在生活中有些事情.
你覺得很棒,很好.
你很希望擁有得到.
和別人一樣.
但上帝偏偏沒有應允你.
沒有批准.
是出於甚麼?.
一方面出於祂的主權.
出於祂有智慧的安排.
但不代表祂不愛我們.
正如摩西被上帝任命.
你覺得他的工作很好做嗎?.
容易做嗎?.
交換好嗎?.
沒有人搶頭.
大家都知道好不好做.
不好做.
但米利安和亞倫只是羨慕他.

$^{721}$看到他的光環.
看不到背後有多艱難.
不是他們可以承擔.
所以上帝沒有應允他們的訴求.
不一定是偏心大細超.
是出於愛,出於保護.
只不過人當下.
可能特別是在他的血氣,疾度裡.
他真的沒有足夠的智慧,耐性.
去體會,理解上帝這份苦心.
這份愛.
以為上帝真的對他們不好.
事實上,弟兄姊妹.
這段經文提醒我們.
要謹守一句說話.
無論你遇到什麼警況.
無論上帝在你身上的安排.
你信不信詩篇裡面說.
上帝本為善.
這句說話沒有任何背景.
沒有逆境才是上帝本為善.
順境才是上帝本為善.
是上帝的本質是美善.
換句話來說,他對你的安排.
他說是好,有用,有意思.
可能我們不明白.
他對我們的慈愛是永遠長存.
我們能否因為這樣的緣故.
總是帶著一份敬畏上帝的心.
來和他同行.
我們尊重上帝在我們身上.
或在別人身上一些主權的安排.
因為上帝愛我們.
我們尊重他.
因為上帝有智慧,我們重視他.
因為上帝有主權,我們願意接受.
弟兄姊妹,這句話很重要.
因為上帝雖然是慈愛上帝.
但他亦是興邁不得的上帝.
繼續看下去,你就會明白.

$^{761}$9至15節,上帝的烈怒和刑罰.
馬上臨到,說完上帝就馬上走.
隨之而來的就是刑罰.
你會看到米利暗和阿倫.
沒有機會辯駁.
不能招架.
毫無預警下,大麻風就馬上出現在米利暗身上.
上帝興邁不得.
阿倫他為何一應一應認得出.
因為他大祭司,見慣,檢驗慣大麻風.
他一看就知道闖了大禍,大件事.
你留意,上帝沒有馬上擊殺這兩個人.
其實已經留有恩典.
留有一個代求認罪的機會.
於是阿倫很明白,回轉的機制.
沒有第一條路,就是認罪悔改.
在摩西面前.
表示他和米利暗得罪上帝只是出於無知.
我們不是真的有心頂撞上帝.
我們是小朋友,不懂得死.
希望上帝可以釋怒.
經文強調,阿倫並非向上帝直接祈求.
而是向摩西求你醫治米利暗.
你見到阿倫稱呼摩西為「我主」.
換句話說,他在承認上帝給摩西的尊貴身份.
給阿倫和米利暗更尊貴.
同時,阿倫這句說話也是收回他之前對摩西的誹謗言論.
他重新承認上帝是有權在摩西身上宣示主權.
因為這樣,阿倫沒有進一步受到刑罰.
至於米利暗,我們看到.
最後他固然得到上帝的憐憫,甚至醫治.
不過,要承擔後果.
你看到第十三至十五節.
其實是按照當時的文化背景.
如果女兒得罪了父親或女兒犯罪.
父親會在她的臉上套唾沫作為懲罰.
令她蒙羞七天.
照樣,米利暗現在得罪父親上帝.
也應當被關在刑外,視他為不潔淨.
讓他在眾人面前蒙羞七天.

$^{801}$你看到米利暗的嫉妒.
他對上帝的埋怨換來是甚麼?.
是自己的羞辱,是一個公開的羞辱.
這七天你會象徵著上帝不看他的眼.
掩面不看他,不跟他說話.
並且其他人也是這樣.
其他人也不跟他說話.
米利暗被公開的羞辱.
不過,你會看到上帝已經重慶發落.
因為七天而已.
他的名字並沒有從以色列人的名單中被剪除.
當然,你會見到米利暗關在刑外這七天.
以色列人也沒有前行.
給我們一個很重要的提醒.
意味著原來屬靈領袖一時犯錯.
並不是純粹自己的事.
屬靈領袖控制不了自己的情緒.
控制不了自己的質度.
有可能影響到整個群體的屬靈進展.
屬靈成長.
所以,最後我想和大家回頭看摩西.
我們看看摩西如何顧全大局.
這種的事奉態度.
我們看第二節的尾和第三節.
這裡耶和華聽見.
即是他聽到米利暗和亞倫對摩西的攻擊.
對上帝的埋怨.
耶和華聽到.
然後作者的註腳.
他形容摩西在這件事.
他展現出一種什麼的屬靈質素.
摩西為人極其謙和.
勝過世上眾人.
換句話來說.
如果其他人可能會怎樣.
其他僕人領袖可能會.
怎樣也要認對方兩句.
怎樣也不值得.
或者說找證據證明自己不是.
證明自己清白.

$^{841}$甚至拉攏自己的粉絲.
來支持自己.
有可能是這樣.
但這裡說摩西不是這樣.
摩西很明白發生在眼前的.
不單單只是家人的爭拗.
你想想.
剛才我們一開始說了.
米利暗是什麼.
女先知.
亞倫呢.
大祭司.
摩西呢.
以色列人的領袖.
我們說以色列三巨頭.
大家如果吵架會怎樣.
是不是家事.
就不是了.
你想像到以色列三個巨頭.
他們有沒有自己的粉絲.
怎樣也有.
多多少少也有.
如果這三個人真的反目的時候.
群體會怎樣.
撕裂.
一定.
教會我們有時也見到.
後果一定不敢設想.
摩西很明白他面對的.
不單單是自己的家事.
是神家的事.
他沒有能力擺平.
怎樣擺平.
一個是女先知.
一個是大祭司.
一個是姐姐.
一個是哥哥.
很難擺平.
有理說不清.
他知道唯有上帝才有能力擺平.

$^{881}$所以你看到他.
他沒有即時開口反駁去還擊.
他而是安靜地等.
他不做任何回應.
他等什麼.
等他的主人出世.
所以你見到黑二經文描述怎樣.
面對著米利暗和亞倫的夾擊.
第一個有反應的不是摩西是誰.
是上帝.
既然上帝都聽到了.
做僕人的摩西又怎樣.
當然是懂得站位.
摩西知道上帝是知道的.
上帝聽到.
這是一個很重要的提醒.
沒有一樣東西上帝是不知道的.
並且上帝不會置之不理.
做僕人既然明白這件事.
就懂得自己站好位置.
他不會自己衝出來.
不會自己頂不住頸.
和那些誹謗他的人.
去頂嘴和他們拼命.
他而是懂得克制自己.
而是懂得修煉自己.
耐心等上帝替他來作主.
這就是摩西極其謙和.
勝過世上所有的人.
這個事奉心態.
謙和通常的翻譯是.
那些很貧窮.
很受壓抑的情況.
受逼迫的情況.
用來形容那些.
總是依靠上帝.
以上帝為優先.
以上帝為樂的人.
因為他們確信上帝會聆聽他們的苦情.
並且知道上帝是願意.

$^{921}$適當時候出手相助.
不會忘記他們.
坦白說.
其實被家人這樣惡意重傷.
傷不傷?.
很不舒服.
很不開心.
但你見到摩西的謙和.
令他可以勒住自己.
不會以牙還牙.
謙和也驅使摩西.
不會因為私事.
而忘記了他的公事.
他忍受上帝給他的任命.
謹守他先知的職份.
他願意替米利暗代求.
摩西這種謙和.
我們說.
好像心理學今天強調.
擁有一個良好的.
自我區分的能力.
什麼意思呢?.
即是說摩西既會明白自己心路歷程.
他可以和人的情緒連繫.
更加和上帝結連.
但是他又不會因為別人.
一些負面情緒.
一些激勵的情緒.
影響.
否定了自己.
做一個正確的判斷思考.
他有這種能力.
因為他知道上帝還在.
上帝聽見.
上帝會替他作主.
換句話來說.
摩西這個謙和.
不是純粹說他一個美好的品格.
謙和是反映他和上帝的美好關係.
他懂得以上帝優先.

$^{961}$他知道上帝在.
他會為上帝顧全大局.
無法冷戰.
他將自己一些需要.
看為次要.
並且他相信上帝會樂意.
為他這位忠心的僕人.
挺身而出.
親自為他平反.
事實上頂梓梅.
謙和就像一件你看不到的避彈衣.
你穿著它會怎樣?.
你可能會中箭.
可能會中彈.
穿著避彈衣來做什麼?.
就是你預計有機會避不開.
我們常說別人的嘴在別人身上.
你怎樣去控制?.
總有機會遇到無理的惡待.
甚至乎攻擊.
但是謙和就像一件避彈衣.
幫你抵擋.
抵禦別人的攻擊批評.
有時中彈難免你會不開心.
但你不會重傷.
不會沒命.
不會得罪上帝.
謙和可以保守我們.
不會一味聚焦在別人的言論.
被人牽著鼻子走.
謙和是幫我們取主動.
對準上帝.
我們跟著上帝走.
信服在祂的主權.
將我們生命的焦點.
再一次放在上帝喜悅的事情上.
所以盼望今天藉著這段經文.
作為這個系列的宣講.
讓我們好好面對.
我們自己的心路歷程.

$^{1001}$求上帝監察我們的心思意念.
我們有個簡單的祈禱.
祈求上帝願你住在我們心裡.
監察我們的心思意念.
或許我們像經文所說.
被別人的說話令我們很不舒服.
令我們很想反擊.
我們求主你保守我們.
讓我們因著上帝你知道的緣故.
我們勒住自己.
勒住自己的衝動.
勒住自己的舌頭.
我們為了上帝的緣故.
我們願意去顧全大局.
為上帝你作見證.
又或者反過來.
我們真的很想用我們的口.
去說一些不造就人的說話.
可能是出於覺得對我們不公道.
或許甚至如經文所說.
我們真的可能妒忌別人的一些成就.
為何上帝不選擇我.
卻選擇祂.
求主你幫助我們.
化解我們這些心結.
讓我們知道其實上帝你就是愛我們那一位.
上帝有時你沒有愛我們的心意.
而來成就我們.
是因為出於你的主權.
出於你的愛.
出於你奇妙自秘的安排.
我們知道上帝你對我們本為善.
你的慈愛亦對我們永遠長存.
求主幫助我們.
閱立我們在你面前的禱告.
奉耶穌基督你的名字.
你也可以祈求.
阿們.
\newpage



\section{約翰福音 6:1-14}
\label{sec:0SAEU456WxE}
\textbf{2024年8月11日崇拜直播｜梁家麟牧師｜吃飽了不餓｜約翰福音6\_1-14}
\newline
\newline
連結: \href{https://youtube.com/watch?v=0SAEU456WxE}{\texttt{https://youtube.com/watch?v=0SAEU456WxE}} ~~~~ 語音日期: 2024-08-19
\newline
\newline
\hyperref[sec:Hufazqfe3oE]{\small{< < < PREV SERMON < < <}}
~
\hyperref[sec:index_chronic]{\small{[返順時目]}}
~
\hyperref[sec:Ef4j6FO5fXg]{\small{> > > NEXT SERMON > > >}}
\newline
\newline
約翰福音 6:1-14
\newline
\begin{longtable}{cl}
\hline
\hline
章節 & 經文 (和合本修訂版)\\
\hline
6:1 & \begin{tabularx}{0.7\textwidth}{X} 這些事以後，耶穌渡過加利利海，就是提比哩亞海。 \end{tabularx} \\ \\ \relax
6:2 & \begin{tabularx}{0.7\textwidth}{X} 有一大群人因為看見他在病人身上所行的神蹟，就跟隨他。 \end{tabularx} \\ \\ \relax
6:3 & \begin{tabularx}{0.7\textwidth}{X} 耶穌上了山，和門徒一同坐在那裡。 \end{tabularx} \\ \\ \relax
6:4 & \begin{tabularx}{0.7\textwidth}{X} 那時猶太人的逾越節近了。 \end{tabularx} \\ \\ \relax
6:5 & \begin{tabularx}{0.7\textwidth}{X} 耶穌舉目看見一大群人來，就對腓力說：「我們到哪裡去買餅給這些人吃呢？」 \end{tabularx} \\ \\ \relax
6:6 & \begin{tabularx}{0.7\textwidth}{X} 他說這話是要考驗腓力，他自己原知道要怎樣做。 \end{tabularx} \\ \\ \relax
6:7 & \begin{tabularx}{0.7\textwidth}{X} 腓力回答他：「就是兩百個銀幣的餅也不夠給他們每人吃一點點。」 \end{tabularx} \\ \\ \relax
6:8 & \begin{tabularx}{0.7\textwidth}{X} 有一個門徒，就是西門‧彼得的弟弟安得烈，對耶穌說： \end{tabularx} \\ \\ \relax
6:9 & \begin{tabularx}{0.7\textwidth}{X} 「這裡有一個孩子，帶著五個大麥餅和兩條魚，但是分給這麼多人還算甚麼呢？」 \end{tabularx} \\ \\ \relax
6:10 & \begin{tabularx}{0.7\textwidth}{X} 耶穌說：「你們叫大家坐下。」那地方的草多，人們就坐下，男人的數目約有五千。 \end{tabularx} \\ \\ \relax
6:11 & \begin{tabularx}{0.7\textwidth}{X} 耶穌拿起餅來，祝謝了，就分給坐著的人，也同樣分了魚，都照他們所要的來分。 \end{tabularx} \\ \\ \relax
6:12 & \begin{tabularx}{0.7\textwidth}{X} 他們吃飽後，耶穌對門徒說：「把剩下的碎屑收拾起來，免得糟蹋了。」 \end{tabularx} \\ \\ \relax
6:13 & \begin{tabularx}{0.7\textwidth}{X} 他們就把那五個大麥餅的碎屑，就是大家吃剩的，收拾起來，裝滿了十二個籃子。 \end{tabularx} \\ \\ \relax
6:14 & \begin{tabularx}{0.7\textwidth}{X} 人們看見耶穌所行的神蹟，就說：「這真是那要到世上來的先知！」 \end{tabularx} \\ \\
[1ex]
\hline
\hline
\end{longtable}
$^{1}$今日的經訓是記載在《約翰福音》六章一至十四節.
在新約聖經一百三十三頁.
這節聖經是記載五餅而儀的神蹟.
《約翰福音》六章一節.
「這是以後耶穌到過加利利海.
就是提比利亞海.
有許多人因為看見他在病人身上所行的神蹟.
就跟隨他.
耶穌上了山.
和門徒一同坐在那裡.
那時猶太人的逾越節近了.
耶穌舉目看見許多人來.
就對肥力說.
我們從那裡買餅.
叫這些人吃呢.
他說者說是要試驗肥力.
他自己原是知道要怎樣行.
肥力回答說.
就是二十兩銀子的餅.
叫他們各人吃一點也是不夠的.
有一個門徒.
就是西門彼得的兄弟安德烈.
對耶穌說.
在這裡有一個孩童.
戴著五個大麥餅.
兩條羽.
只是分給這許多人.
還算什麼呢.
耶穌說.
你們叫眾人坐下.
原來那地方的數很多.
眾人就坐下.
數目約有五千.
耶穌拿起餅來捉著了.
就分給那坐著的人.
分魚也是這樣.
都隨著他們所要的.
他們吃飽了.
耶穌對門徒說.
把靜下的靈髓收拾起來.

$^{41}$免得有糟蹋的.
他們便將那五個大麥餅的靈髓.
就是眾人吃了靜下的.
收拾起來.
裝滿了十二個籃子.
眾人看見耶穌所行的神蹟.
就說.
這真是那要到這世間來的先知.
鄧炳智,早安.
今天我都有點擔心.
今天講的道.
不容易有共鳴.
上一篇道比較複雜.
因為講的是神學性的東西.
大家可能在理解上有些困難.
今天講的東西很簡單.
淺到不得了.
大家一定明白.
不過.
我猜六十歲以下.
或者家境比較豐裕的.
都不知道我講什麼.
大家聽聽便知道了.
這是五餅二魚的神蹟.
我很熟悉的.
耶穌基督在人間所行.
最偉大的其中一件神蹟.
四卷福音書都有記載這個故事.
不過在細節上.
是有一點點不同的.
三卷福利福音.
都有交代故事發生的背景.
這是希律王斬了施洗約翰的頭.
沒多久才發生的.
路加特別指出.
有人將耶穌傳道的故事.
告訴希律王.
並且說坊間有傳言.
耶穌就是另一位舊約先知的復活的顯現.
希律王於是說.

$^{81}$我要見耶穌.
意思大概就是要抓他.
一不做二不休.
能夠殺死施洗約翰.
就不差連耶穌也殺了.
任何先知的出現.
都對政權有威脅的.
一定要防範於未然.
滿途將他們被通緝的消息.
告訴耶穌.
他們便立即靜悄悄地離開.
坐船去到白菜大.
不過他們出走的消息.
被群眾發現.
很多人在岸上.
見到耶穌的船.
不知道耶穌去哪了.
他們就在岸上.
繞在岸邊追求耶穌的船.
他們一直追求耶穌.
船只能走得很慢.
追得到的.
並且他們一直走.
一直高聲說.
耶穌來了.
耶穌去哪了.
於是吸引其他沿岸漁村的人.
加入追趕的行列.
於是最後超過五千人去追求耶穌.
由於群眾突然決定.
被吸引加入追趕的行列.
他們沒有任何的預備.
所以他們沒有帶食物.
他們聽完耶穌講道之後.
就發覺大家都沒有東西吃.
耶穌最後施行了五餅二魚的神蹟.
神蹟發生之後.
群眾和門徒都開始對耶穌的身份有不同的討論.
他們先知他是救藥人物的復活等等.
而耶穌基督也第一次向門徒披露.

$^{121}$他即將被殺和復活的使命.
這是《路加福音》九章十八到二十七節所記述的.
今天我們專注看.
約翰福音所記述五餅二魚的故事.
在約翰的北下.
這個五餅二魚的故事有兩個獨特的地方.
第一,這裡有提名的兩個門徒.
一個是肥力,一個是安德烈.
都是平日出鏡率很低的兩個人.
我們慣常在耶穌的故事裡面看到的是.
彼得,雅各和約翰,對吧?.
彼得是最勇於發言,敢於任事的一個.
肥力和安德烈除了在被召的故事提過他們的名字之外.
紀卷,福音書幾乎都沒有提過他們的名字.
第二個特點是.
四卷福音書都有記載五餅二魚的神蹟.
但是只有約翰提到.
捐出自己的食物的是一個小孩.
其他書卷都是說.
門徒就這麼說.
我們這裡有五個餅,兩條魚.
沒有說那個來源.
所以你可以看到.
約翰是特別在這個故事裡面提到一些小人物.
就是平時不出鏡的門徒.
提到一個小朋友.
約翰沒有直接提到,才是約翰被殺.
耶穌被希律盯著.
所以要和門徒離開這個背景.
只是說耶穌他們坐船去到提比利亞海.
然後又指出群眾是因著耶穌有治病的能力.
而跑來跟隨他.
無論如何,耶穌和門徒在某個地方登岸.
爬上一個山崗,然後在那裡教訓眾人.
和符雷夫音的技術不同.
不是門徒首先發現群眾有缺糧的問題.
而是耶穌自己先提出這個問題.
不是門徒束手無策.
被迫前來尋求耶穌的幫助.
而是耶穌主動用這個困局.

$^{161}$來考驗門徒的信心.
第五第六節是這樣說.
耶穌舉目看見許多人來就對肥力說.
我們從哪裡買餅叫這些人吃呢?.
他說者說,是要試驗肥力.
他自己原知道要怎麼行.
每個意外,每個突發的處境.
每個危機困局.
都可以是一個信心的試驗.
這個不僅僅是測試.
我們對幕下的警方的一個研判.
我們所能將來的資源.
所能設想的解決方案.
更加是測試我們心中所持定的信念.
以及我們對這些信念的堅定堅固的程度.
很多人說信上帝.
很多人說會相信上帝幫助.
不過他們常常憂心忡忡.
那個徬徨焦慮的樣子.
跟那些不信上帝的人其實不是有太大差異.
他們一向宣稱上帝是幫助他們的上帝.
不過處於這個境況.
他們不免就動搖信念.
一個弔詭的事實是.
最需要信心的時候.
卻是最容易失去信心.
就好像我們在街上.
你不用截的士的時候.
你就總是見到有空的士在你身邊.
你要截的時候就沒有了.
不知為何.
要截的時候就連打鑼都找不到的士.
所以每當遭遇困難的境況.
我們除了急於去尋求各種解決方法.
包括切切向上帝尋求幫助之外.
我們都需要騰出我們的心靈空間.
問上帝要我落在這樣的境地.
有祂的心意嗎?.
祂期望我有什麼回應和表現呢?.
祂是否要藉著這件事來測試我的信心呢?.

$^{201}$祂是否要藉著這樣的事.
來對我有新的啟悟.
讓我學習成長呢?.
聖經說.
耶穌說者說.
是要試驗肥力.
他自己原知道要怎麼行.
耶穌對問題是有答案的.
對困境是有解決的出路的.
不過他就是不喜歡.
不肯即時出手解決.
即使我們祈禱去求問他.
他都不肯即時回覆.
原因是他要試驗我們.
如果求什麼就立刻有什麼.
如果要求幫助就立刻得到幫助.
我們還有試驗信心的空間嗎?.
我們還有成長的需要嗎?.
我們都成為了那些所謂屬靈上的巨嬰.
Giant Baby.
不夠錢嗎?.
問上帝吧.
想不通?.
問上帝吧.
做不通?.
問上帝吧.
上帝是王大仙.
有求必應.
上帝一切問題的Easy way out.
簡易出路.
有上帝.
我就不用負任何責任.
不用過負責任的人生.
弟兄姊妹.
我必須坦白地說.
在外觀上.
基督徒和非基督徒.
其實差異很少.
少到幾乎分不出來.
或者唯一.

$^{241}$我比較容易見到的就是.
基督徒吃飯前會借飯.
基督徒說信上帝.
非基督徒就說信命運.
信自己.
當遇上危機的時候.
他們焦慮.
我們一樣焦慮.
他們怨天尤人.
我們也怨聖在道.
我們在困難面前的那個騰雞樣.
和非基督徒其實毫無分別.
唯一不同的是.
我們還會埋怨上帝.
在別人有的經濟危機.
政治危機.
社會危機之餘.
我們全都有.
我們多過信仰危機.
非基督徒病急亂投醫.
基督徒一樣會四出張羅.
答案出路.
除了信上帝之外.
他們其他什麼古靈精怪的偏方都信.
他們信不過自己.
非基督徒信不過自己.
我們也不是很信得過上帝.
我們說的上帝.
和他們說的命運其實沒什麼分別.
大家都是等運到.
我在網上看到一段有趣的信息.
是一些對話.
有一個基督徒就這樣提問.
他說他被詐騙集團騙走不生即出.
他問可不可以問教會尋求接濟幫助.
於是一大堆人去回應他.
不過回應的人很奇怪.
主要是批評這個提問的人.
說他沒智慧.
你怎麼可以被人騙走所有的錢.

$^{281}$基督徒不是有屬靈眼光.
屬靈智慧嗎.
怎麼會被人騙這麼蠢.
我覺得問題有奇怪.
答案更奇怪.
我根本不覺得讀聖經.
包括你讀《箴言》.
就能夠使你有避免被人欺騙的智慧.
我不相信.
聖經很清楚說.
今世之子比光明之子要聰明.
你以為你讀了聖經就避免被人騙.
讀了智慧書就不會被人騙.
我不相信這一套.
我更加關心的是.
不是基督徒被人騙的問題.
而是基督徒為什麼要找快錢的心態.
妄圖不需要付勞力.
就可以獲百倍收穫的秘笈.
沒智慧不是問題.
貪心只是問題.
我不知道你同不同意.
我是一個完全不懂投資的人.
我知道我們中間有很多財經界的精英.
你們可能很鄙視我講以下的話.
我這個愛行者我這樣說.
如果有人想賺快錢.
我寧願他去買六合彩.
包括對基督徒.
我也是這樣說.
買六合彩.
因為很少人會買六合彩.
每次都會搞大包圍.
買六合彩搞到傾家蕩產.
我沒聽過.
我只知道買六合彩多數是.
坐在那些投注站外面的老伯.
他們只是買十多二十元.
消磨時間.
但你跟人去做投資.

$^{321}$對方允許你每年40\%賺錢.
於是你就押整副身家下去.
我覺得風險大過你.
賺幾十元去買六合彩.
你同不同意嗎?.
而不道德的程度.
我看不到有分別.
你不要告訴我這個叫投資.
那些叫賭博.
我這些不懂投資的人.
我這樣說.
為什麼你會期望.
跟人去做一件事.
每年賺40\%.
回到我們說的聖經.
這套聖經給了我們一個角度.
讓基督徒可以活得跟非基督徒不一樣.
什麼不同.
就是我們視所有人生際遇.
不管是信逆成敗.
都是上帝對我們信仰的測試考驗.
是這樣的.
我們所有的事件.
所有的行動.
都是信仰事件.
信仰行動.
都是信服或是搏逆的選擇.
我被公司解僱.
我家裡立即出現財政壓力.
我個人的自我形象都不免受損.
我落在一個灰心沮喪的地步.
聖靈此刻提醒我.
或者是上帝考驗我.
能否在逆境仍然對祂有信心的測試.
即使背後沒有像約伯記那樣.
上帝跟撒旦談判的故事.
上帝藉著逆境來考驗我.
讓我在其中有學習成長.
仍然是很可能的.
很可能的.

$^{361}$我上個月跟一對夫婦吃飯聊天.
太太半年前得了一個怪病.
全身沒有力氣.
有一段時間連洗澡的基本功能都沒有.
要丈夫抱出抱入.
她跟公司取了三個月病假.
丈夫笑著說.
這是考驗她在結婚之前.
在結婚的時候所做過的承諾.
是否真實的時候.
那張無論健康疾病都愛你.
這張名牌照.
她說現在要兌現.
是嗎?是的.
現在大家風騷.
那張空頭支票只是說說而已.
有一天可能要兌現.
我常常說基督徒不僅在傳福音.
和參與教會服事的時候.
才算履行信仰的任務.
我們的婚姻和家庭是信仰任務.
我們在職場對待同事和顧客.
是信仰任務.
我們起居作息,待人處事.
通通是信仰任務.
是這一點使得我們活得跟非基督徒有所不同.
亦是這樣的想法.
使得我們將生活和信仰緊密的結連.
我們回到今天的經文.
面對這個突然的測驗.
肥力的答卷是怎樣的呢?.
他的回應是:無法解決.
一萬個不能.
正如肥力回答說.
就算二十兩銀子的餅.
叫他們各人吃一點也是不夠的.
需要太大.
遠超出我們的能力.
所以無法解決.
一開始就舉手投降.

$^{401}$另一個門徒插嘴.
就是西門彼得的兄弟安德烈.
他說:就他所知.
有一個小朋友拿著五個大麥餅兩條魚.
所以他們不是好像肥力說的一無所有.
有,不過所有的是微不足道.
六章八到九節說.
只是分級者許多人還算什麼呢?.
你要注意.
餅是大麥餅.
意思是大麥做的餅.
Bali bread.
不是很大的麥餅.
大字不只是餅很大.
魚也不是大魚.
五個餅和兩條魚.
是剛剛好一個小朋友的飯量.
大概這個小朋友在出門之前.
家裡為他預備了便當.
一個人份的食物.
怎麼能夠餵飽.
男生就五千人.
加上女生就過萬人.
怎麼夠呢?.
馬太福音記載.
還有門徒提出建議.
立即遣散群眾.
讓他們各自去不同地方找食物.
門徒耶穌只管嘴領聚會.
不管飯吃的.
就像我們這樣.
你們下午出去找食物.
問教會拿飯吃是沒有的.
我們不會有飯吃供應的.
不管一萬幾千人能不能夠.
在這個曠野附近找到食物.
不過遣散他們就解決了問題.
對不對?.
最少問題不在我眼前發生.
五千人吃飯的問題.

$^{441}$一萬人吃飯的問題.
當然非常嚴重.
問題就在眼前.
門徒也試過尋找各種解決問題的答案.
耶穌是否滿意他們提出的答案.
他們能否通過耶穌的信心測試.
我們不知道.
不過總之就是他們試過做不到.
現在測試完畢.
當然是耶穌要出手.
耶穌真的出手了.
耶穌第一件事就是控制現場的秩序.
他吩咐門徒將群眾安頓好.
分成一堆堆坐在草地上.
六章十節.
馬拉和路加都提到.
群眾被分成五十人或一百人一堆.
我們知道一大批饑民蜂擁爭奪食物是很危險的事.
搞不好可以打架甚至踩死人的.
讓他們分批坐下然後有秩序地分派食物.
不讓爭先恐後.
可以減少意外的發生.
這是一個很周詳的安排.
然後耶穌就拿起安德烈從一個小朋友所拿來的五個餅和兩條魚.
向天舉起, 竹蔗抹開.
然後交給門徒分送給坐在地上的群眾.
我們不知道那個做法是怎樣.
是耶穌在抹餅的時候一次過抹成超過二萬五千個或五萬個餅.
如果小朋友要吃五個餅才飽.
他一次抹成幾萬個餅.
抹了幾千條魚.
抹開再給門徒, 遞給群眾.
群眾再抹開再分食的時候.
一邊抹一邊變多.
還是耶穌一次過變完.
我們不知道整個故事的細節.
不過我自己傾向相信.
不是耶穌一次過抹完幾萬個餅.
而是他們一邊抹一邊傳開的時候.
餅一邊變多一邊變大.

$^{481}$所以每個人不僅僅是神蹟的見證者.
得益者.
其實每個人都是其中的參與者.
實際上, 耶穌刻意讓門徒參與在神蹟裡面.
耶穌大可以從無到有.
自行將實物變出來.
又或者像他受試探的時候.
魔鬼所給的建議.
也可以將石頭變成食物.
不過耶穌決定以門徒找回來的五個餅.
和兩條魚做基礎.
以少變多.
這顯示門徒雖然自覺無能無用.
所能夠有的貢獻微不足道.
不過在耶穌手中還是有價值有作用的.
門徒質疑自己所有和所做的.
還算什麼呢?.
耶穌要求你們給他們吃飽.
我們不足夠, 我們算不得什麼.
但在耶穌的幫助下.
他產生意想不到的後果.
這都是我們常常經歷的事實.
全能的上帝說有就有, 命立就立.
願可以自行成就他的旨意.
他要全世界的人信耶穌.
這都是但此間可以成就的事.
不過上帝卻邀請我們和他同工.
他使用我們在他的國度裡.
上帝的大能可以繞過我們所有的作為.
反正我們的資源和能力都很有限.
盡力而為都不見得有多大用處.
不過上帝定義要在我們簡陋的工作裡.
彰顯他的大能作為.
將水變酒, 化腐朽為神奇.
在人的工作裡看到上帝的工作.
上帝需要若拿的宣告服事.
才能使得你們未成的人悔改嗎?.
以若拿的心意具體行動表現.
他的參與其實是幫倒忙的.
不過上帝就在若拿的宣講服事裡.

$^{521}$而不是繞過他的服事.
來彰顯他奇妙的救贖作為.
同樣的, 黃國俊信會必須有你和我.
才能使這教會運作下去嗎?.
我們有什麼人是不可或缺的呢?.
沒有, 沒有.
故事的結局我們都知道.
不過這個結局是很好玩的結局.
耶穌用五餅二魚餵飽了在場所有人.
我要說, 只要是男人就五千.
連同婦女和小孩, 一定超過一萬人.
十二十三節說:他們吃飽了.
耶穌對門徒說:把剩下的零碎收拾起來.
免得有糟蹋的.
他們便將那五個大麥餅的零碎.
就算眾人吃了剩下的收拾起來.
也將十二個男子裝滿了.
故事加上這樣的結束.
用現代的觀點說, 是否宣揚環保?.
當然環保很現代.
古代的觀念叫「識食識福」.
就像我媽媽一粒飯掉在地上也會撿起來.
食物要珍惜.
是不是這樣呢?.
鼓勵我們吃剩打包, 不要浪費食物.
有這個意思的.
因為耶穌說過:把剩下的零碎收拾起來.
免得有糟蹋的.
就是這個意思.
不過, 我覺得這不是聖經將結尾放在這裡的最主要原因.
第一,剩下的食物其實不是很多.
十二個籃子, 有多大的籃子我不知道.
不過手提籃有多大?.
一萬人吃剩十二籃.
即是一千人吃剩一籃, 很多嗎?.
一籃不是最多十幾十個餅嗎?.
一千人吃剩幾十個餅, 很多嗎?.
坦白說, 耶穌不做這個安排.
你真的相信在場的人會將餅放在地上?.
私下拿都拿了, 放自己的袋子裡, 對不對?.

$^{561}$每個人都帶走.
其實, 你說會怕糟蹋食物, 我不是覺得這是一個很合理的理由.
坦白說, 那時候食物那麼珍貴.
怎麼會被他放在地上?.
每個人都帶走, 對不對?.
我相信整個結尾, 最主要的意思是.
表明耶穌變出來的食物是夠多的.
不是一萬人吃一萬個餅.
每人吃一個, 擋一下肚子, 即是不要餓壞, 即是不會餓暈.
是二思二事.
而是數量大到可以讓所有人都放開懷抱吃.
盡情吃, 任意吃.
吃到最後是人人都塞不下去.
而還剩下十二籃的食物.
十一節有一句說話常常被人忽略.
都隨著他們所要的.
NIV是譯成As much as they want.
翻成我們喜歡的說法就是All they can eat, 對不對?.
即是任吃.
六章二十六節, 即是耶穌comment那班人.
他們來找我, 是想吃餅得飽.
重點不是吃餅, 是得飽.
Worth it.
五餅二魚的神蹟, 其實沒有什麼特別的食物.
餅和魚是當時最普遍, 最低級, 最日常的食物.
最簡單.
所以不是說你吃了一次, 包心刺肚.
很深刻印象.
食物沒有任何印象.
群眾事後主要的回憶.
不是耶穌請我們吃了鮑魚燕窩, 魚翅黑松露.
而是讓我們吃了一餐飽飯.
吃餅得飽.
餅不稀罕.
得飽才稀罕.
我再說.
一開始說.
60歲以下.
即是家裡很富裕的, 不會知道我說什麼.
你們每餐都可以吃.

$^{601}$每餐吃不同的美食.
從來只是擔心吃得太多, 有三高要減肥.
從來沒擔心過吃不飽.
從來不追求吃飽.
所以你們從來都不會明白.
吃飽了的幸福感.
我太太在這裡, 我也這樣說.
她永遠只吃完一餐美食之後就內疚一次.
糟了, 吃多了, 又重一公斤.
我要說.
在食物貴乏, 價錢昂貴的時代.
很多人其實很少機會吃飽飯.
我小時候一家五兄弟姐妹.
不要說菜, 根本沒什麼菜.
一碟蒸水蛋, 七個口分, 根本做不到.
我只知道每餐煮四元米.
一人趕快吃兩三碗.
連飯焦都沒有.
搶的, 我真的沒印象.
我一餐飯吃得飽.
我只知道.
我經常餓的樣子.
我後來在酒樓工作.
最大的好處就是吃得飽.
那時在酒樓很厲害.
每餐吃飯, 用鴨腿湯飯盤去盛飯.
一盛盛三碗.
三碗鴨腿湯飯盤的飯.
有關飢餓的小故事, 我很多.
和你分享少少的兩個.
有一次我大概十二歲那年.
要做手術, 進了東華醫院.
有個親戚晚上來探病.
媽媽已經回家.
親戚送來一包六個橙.
我平時在家幾乎沒機會吃水果.
有一個橙都是五分分.
五兄弟姐妹分, 一人吃兩三塊.
我那晚忍不住.
吃完一個又一個.

$^{641}$最後六個橙一次吃完.
媽媽第二天知道.
罵了我幾年, 不是幾次, 是幾年.
另一個故事, 我太太可以作證.
我進了大學之後.
有一次在飯堂吃完晚飯.
一個客飯, 再問善後婆.
添多一次飯, 兩大碟飯.
吃了足足兩碟.
飯後去宿舍小賣部.
買了一條生命麵包.
藍白色那種最差的.
本來打算留下來做早餐.
誰知晚上在宿舍讀書.
吃一塊又一塊.
最後一口氣吃了一條生命麵包.
又被宿舍笑得我面黃.
我年輕時期那種飢餓的記憶.
使幾十年後的我.
仍然認定吃飽是很重要的.
吃不飽是很討厭的事.
吃餅吃飽, 吃飯也吃飽.
擺盤擺得再美麗的日本餐,法國餐.
吃不飽就令我很反感.
這與我兒女的想法是剛剛相反的.
他們要求吃得精緻, 但怕吃得多.
所以他們都不喜歡自助餐.
根本沒有我以前那種自助餐的回憶.
那些小便宜的自助餐.
吃揚州炒飯也吃幾碟那種.
那時我在加拿大溫哥華讀書.
有一間餐廳叫Uncle Willy.
十元加幣不用就可以任吃炸雞.
那時常常去, 每一盤炸雞都是垃圾.
總之就拼命塞進肚子裡.
你們不會明白的.
正因為古代社會很少機會被人放開肚皮進食.
所以五餅二魚的神蹟就是保證吃飽的一次很難得的經驗.
都隨著他們所要的吃餅得飽, 有魚有成.
這是生命豐盛的一次表徵.

$^{681}$我相信這是十二男子的真正含義.
證明給大家看.
吃到大家頂著肚子, 上了頸, 你明白嗎?.
整個故事就講完了.
想想其中有什麼含義.
五餅二魚的神蹟當然是很具震撼性.
一個小孩子的份量的食物.
竟然可以數千倍增加到餵飽超過五千人.
甚至上萬人之上.
這是多麼不可思議的事.
用我們今天的我們來看.
食物最多是關乎錢.
有錢就可以買到食物.
所以二十兩才能買到一些食物分給人.
要吃飽的可能要一百兩太貴.
不過在古代社會, 食物是關乎生存.
那個重要性是遠遠大於今天的想像.
如果食物不是那麼重要.
其實我們中國人以前的事, 大家都知道.
我們做事叫做「搵食」.
「搵食」是頭號大事.
聖經也說「搵食」真的很重要.
耶穌的那句話.
「人活著不是單靠食物」.
正好從側面說明食物的重要性.
如果食物不重要.
耶穌就不會將食物和上帝的話作對比.
為什麼不說「人活著不是單靠洗澡」.
而是單靠其他?.
不如單靠食物.
因為食物對人來說是頭號大事.
主禱文也是求上帝供應我們一用的食物.
「民以食為天」.
古代人幾乎用一生全部的時間精力去填飽肚腹.
解決溫飽, 佔據了人的關懷心思.
吃餅得飽, 豐衣足食.
衣食無憂, 這是生活最高境界.
耶穌這個神蹟所變出來的東西.
正是人最渴想的東西.
是這樣, 五餅二魚的神蹟在古代社會.

$^{721}$真的有和現代社會不一樣的意義.
所有的民變,暴動,社會革命.
除了有粗糙的民族主義情緒之外.
更多是出於現實的需要.
民不聊生, 活不下去.
就不怕死, 乾起耳.
災荒, 特別是饑荒.
常常是民變產生的背景.
簡單來說, 仍然是為了謀生.
你想想, 如果有一個人.
可以憑空變出大量的食物.
永遠不受缺糧,斷糧.
跟著他就最安全.
一定有得食, 一定食得飽.
約翰福音第六章, 接著就說.
其後有很多人跟隨耶穌.
跟著耶穌謀生.
因為有耶穌就有得食.
耶穌明白地質問他們.
你們找我並不是因見了神蹟.
乃是因吃餅得飽.
六章二十六節.
吃餅得飽比見神蹟更重要.
當然了.
耶穌因為這個神蹟.
而擁有大量的跟隨者.
他們因著這個神蹟而相信.
耶穌是從天上差來的先知.
就是真是那要道世間來的先知十四節.
你知道他說的是哪個先知嗎?.
是以利亞, 為什麼?.
原因有兩個.
第一, 舊約說以利亞沒有在人間死亡.
他是被火車,火馬接到天上去的.
所以嚴格來說, 以利亞還沒有死.
他最少沒有在人間死過.
所以很多人相信.
他被上帝接到天上.
他也會從天上再差下來.
所以他們猜耶穌是以利亞.

$^{761}$這是第一個原因.
第二個原因.
以利亞曾經在撒勒伯.
一個寡婦家裡.
施行過一個神蹟.
將食物以少變多.
雖然這個其實和五餅二魚比較.
是一個太蚊型的神蹟.
只是餵飽寡婦和兒子.
餵飽三個人.
跟耶穌餵飽一萬人.
根本無法比較.
不過總是性質類似.
將食物以少變多.
除了猜測耶穌是以利亞之外.
更加有人想推動耶穌出來做政治領袖.
帶領以色列人推翻羅馬帝國的統治.
成為政治上的理賽亞.
你想想.
如果耶穌搞革命不用籌募糧食.
帶多少兵都沒有問題.
這支軍隊一定打贏.
對我們今天的我們來說.
餅和魚的吸引力重要性.
一定沒有古人那麼多.
這個神蹟的意義在什麼地方.
可以發揮的地方很多.
剛才我一直說這個故事的時候.
耶穌說了很多話.
例如耶穌要求我們首先承擔.
生活中所遭遇的各種問題.
我們視生活每個遭遇是信仰的考驗.
或者耶穌是我們生活基本需要的供應者.
耶穌是將我們的生命和事奉點石成金的人.
我們做一點點.
然後就將水變酒.
這都可以是我們在這個故事裡面反省到的.
不過我覺得最重要的是十四字的一句話.
眾人看見耶穌所行的神蹟就說.
「謝真是那要到世間來的先知」.

$^{801}$聖經裡面的群眾.
因著看到耶穌所行的神蹟.
而對他有新的看法.
耶穌以他的服侍讓人認識他.
聖經作者之所以記述耶穌各種的事蹟.
目的就是要讓我們更真確的認識他.
正如約翰在約翰福音20章31節說.
「但記這些事,要叫你們信耶穌是基督,是上帝的兒子,並且叫你們信了他,就可以因他的名得生命」.
所以我們在聖經裡面看到耶穌各個的故事.
包括在這裡的五餅二魚的神蹟.
我們在我們個人的經歷.
在教會群體的經歷裡面.
我們遇過耶穌.
我要說每一個的遭遇.
每一個的經歷.
每一個的故事.
都逼我們不斷再問.
再問.
耶穌是誰.
我們這位主是誰.
對我.
耶穌是誰.
這是我們不停問的問題.
我的主.
我的上帝.
\newpage



\section{撒母耳記上 10:9-13}
\label{sec:Ef4j6FO5fXg}
\textbf{2024年8月18日崇拜直播｜陳明泉牧師｜正視你的情緒－焦慮的糾纏｜撒母耳記上10\_9-13;19\_18-24}
\newline
\newline
連結: \href{https://youtube.com/watch?v=Ef4j6FO5fXg}{\texttt{https://youtube.com/watch?v=Ef4j6FO5fXg}} ~~~~ 語音日期: 2024-08-19
\newline
\newline
\hyperref[sec:0SAEU456WxE]{\small{< < < PREV SERMON < < <}}
~
\hyperref[sec:index_chronic]{\small{[返順時目]}}
~
\hyperref[sec:cerRyXhHqiE]{\small{> > > NEXT SERMON > > >}}
\newline
\newline
撒母耳記上 10:9-13
\newline
\begin{longtable}{cl}
\hline
\hline
章節 & 經文 (和合本修訂版)\\
\hline
10:9 & \begin{tabularx}{0.7\textwidth}{X} 掃羅轉身離開撒母耳，神就改變他，賜給他另一顆心。當日這一切徵兆都應驗了。 \end{tabularx} \\ \\ \relax
10:10 & \begin{tabularx}{0.7\textwidth}{X} 他們來到那座山，看哪，有一隊先知遇見掃羅。神的靈大大感動他，他就在先知中受感說話。 \end{tabularx} \\ \\ \relax
10:11 & \begin{tabularx}{0.7\textwidth}{X} 所有先前認識掃羅的人看見了，看哪，他和先知一同受感說話，百姓就彼此說：「基士的兒子遇見了甚麼呢？掃羅也在先知中嗎？」 \end{tabularx} \\ \\ \relax
10:12 & \begin{tabularx}{0.7\textwidth}{X} 那地方有一個人說：「這些人的父親是誰呢？」因此就有一句俗語說：「掃羅也在先知中嗎？」 \end{tabularx} \\ \\ \relax
10:13 & \begin{tabularx}{0.7\textwidth}{X} 掃羅受感說完了話，就上丘壇去了。 \end{tabularx} \\ \\
[1ex]
\hline
\hline
\end{longtable}
$^{1}$接下來是經訓.
今天的經文記載在.
撒母耳記上的第十章和第十九章.
經文都印在你們的秩序表內頁.
你們可以聽我讀.
第十章的第九節.
「蘇羅轉身離別撒母耳.
神就賜他一個身心.
當日這一切兆頭都應驗了.
蘇羅到了那山.
有一班仙子遇見他.
神的靈大大感動他.
他就在仙子中受感說話.
蘇萊認識蘇羅的.
看見他和仙子一同受感說話.
就彼此說.
「既是的兒子,遇見甚麼了.
蘇羅也列在仙子中嗎?」.
那地方有一個人說.
「這些人的父親是誰呢?」.
此後有句俗語說.
「蘇羅也列在仙子中嗎?.
蘇羅受感說話已不.
就上幽壇去了」.
然後是第十九章.
第十九章的第十八節.
「大衛逃避來到拉瑪見撒母耳.
將蘇羅向他所行的事.
述說了一遍.
他和撒母耳就往拿約去居住.
有人告訴蘇羅.
說大衛在拉瑪的拿約.
蘇羅打發人去捉拿大衛.
去的人見到一班仙子都受感說話.
撒母耳站在其中監管他們.
打發去的人也受神的靈感動說話.
有人將這事告訴蘇羅.
他又打發人去.
他們也受感說話.
蘇羅第三次打發人去.

$^{41}$他們也受感說話.
然後蘇羅自己往拉瑪去.
到了西姑的大井問人說.
撒母耳和大衛在那裡呢?.
有人說在拉瑪的拿約.
他就往拉瑪的拿約去.
神的靈也感動他.
一面走一面說話.
直到拉瑪的拿約.
他就脫了衣服.
在撒母耳面前受感說話.
一晝一夜怒體躺臥.
因此有句俗語說.
蘇羅也列在先知中魔.
聖經讀畢.
弟兄姊妹平安.
我們今天繼續正視情緒的系列.
不知道你有沒有.
有時候人生裡面.
面對有焦慮的一些狀況.
焦慮其實在這個定義裡面很闊.
醫學上有定義焦慮的.
甚至乎可以視為情緒病的一種.
不過人的焦慮.
不等於你有焦慮.
就等於你有情緒病.
其實每個人在一個動盪.
或者是一些轉變的環境當中.
我們都會出現一種焦慮.
我回想起我最焦慮的一次.
就是學車的時候.
不知道大家有沒有經歷過.
第一次坐在車上.
然後我還要學棍波.
之後我現在是在開自動波的.
中間有些故事我會再說.
不過第一次坐在棍波車上.
然後覺得感覺很良好.
但是開車的時候.
一踏實那種焦慮感就上了.

$^{81}$然後那個師傅說.
進球了,然後踩煞車,然後去吧.
他說得很容易.
但是原來很難.
然後我的手好像打結一樣.
都不知道自己在做什麼.
然後一路走的時候.
車子就不停地死火.
所以第一次學車的時候.
我發覺我自己焦慮是爆燈的.
然後完結之後.
師傅就學了大概九個字.
然後師傅就嘆了一口氣.
我從來沒見過男士像你這麼笨手笨腳的.
然後我自己都頭大.
我跟師傅說.
很久都沒有人這樣說過我.
謝謝你.
人生一定有些焦慮的.
有些焦慮你可以一笑置之.
但是有些焦慮如果在你的生活.
在你的家庭.
在你的工作裡面.
你不去面對它的時候.
焦慮的那種情感.
或者你自己想繼續做一些事情.
來克服焦慮的時候.
可能會越踩越遠.
甚至乎可能會令到自己進入一種更加大的混亂當中.
今天我們讀的這段經文.
其實我們讀的篇幅會更加多.
不過我們想聚焦在一個人物身上.
就是以色列國的第一個王 蘇羅.
蘇羅王本身是高人一等.
在一個真的某程度上是一個君主的一種形象.
不過在他生命裡面產生大量的焦慮感.
而這些焦慮感他沒有好好面對的時候.
就越走越下.
成為了上帝厭棄的一個君王.
今天我們從蘇羅的身上.

$^{121}$我們看看有沒有一些功課值得我們學習.
又幫助我們檢視一下我們的內心.
究竟我們今天處於一種什麼樣的狀況.
我們先祈禱求上帝幫助.
上主願你幫助我們.
讓我們能夠在你的話語當中.
我們生命得著激勵和調教.
願你在此時此刻.
你對我們眾弟兄姊妹講說話.
對我們的內心講說話.
你的話語進到我們的心扉當中.
更新改變也都督責我們.
讓我們行在你的旨意裡面.
願你在此時此刻.
你幫助每位聆聽的弟兄姊妹聽得明白.
幫助孫講的講得清楚.
我們謙卑苦伏在你面前.
求你教導我們.
獻上禱告奉靠基督耶穌得勝的聖名.
Amen.
我們看蘇羅有一個很重要的字眼.
就是聖靈與他同在.
在撒慕爾記上.
基本上由第九章開始.
到第二十七章.
一直講到蘇羅和上帝的靈之間的互動.
在頭半段的時候.
聖靈成為了蘇羅王的一個很重要的印記.
告訴整個以色列民族.
他是上帝與他同在的一個選立的君王.
所以剛才我們讀的兩段經文裡面.
第一段經文很重要.
因為那段經文講到.
聖靈感動他.
給他一個新的心.
撒慕爾記上十章九至十三節裡面講到.
離開了撒慕爾.
其實他受到撒慕爾的高納.
上帝給他一個新的心.
一個新造的人.

$^{161}$然後之前蘇羅要找那個驢仔.
一切的兆頭就應驗了.
經文裡面特別記載一個很小的.
但是重要的篇幅.
就是講到蘇羅去到那個地方.
那個山裡面有一群先知.
見到蘇羅.
神的靈大大感動他.
他就在那些先知當中受感說話.
蘇羅認識蘇羅的.
看到他和先知一同受感說話就彼此說.
幾時的兒子遇見了什麼.
蘇羅也列在先知中的嗎.
那地方又有人說.
這些人的父親是誰呢.
此後有句俗語說.
蘇羅也列在先知中的嗎.
蘇羅受感說話已不.
就上幽壇去了.
在經文裡面你多次見到.
神的靈感動他.
然後上帝的靈感動他.
令他受感說話.
受感說話是什麼意思呢.
有些不同的聖經都出現過這個字.
有些聖經說說如言.
不過受感說話.
其實特別指向上帝的靈.
入到某一些人的身上.
特別是在先知身上.
在先知被上帝感動的靈.
然後他就說出上帝的訊息.
不是每個人都有的.
這不是一個普遍有的一種經歷.
所以是上帝特別選定的某些人.
當時的先知.
而在這裡蘇羅不是先知.
不過蘇羅是上帝選立的君王.
是選選的君王.
他被高納了.

$^{201}$上帝用他的靈囂冠在他身上.
以至他就像先知一樣.
說出一些語言受感說話.
這個受感說話的內容究竟說什麼呢.
經文沒有特別記載.
我們不推敲的.
不過可以肯定的.
或者到了今天現代.
有些近似是說謊言.
某一些宗派很強調.
某一些頂子會領受了.
有種能力能夠說謊言.
不過這不是一個普遍的現象.
有些說法就是說.
可能他好像先知那樣.
安著上帝的感動.
說到一些平時以他的水平.
或者以他平時的生活.
不會說出的那些話語.
這個就是受感說話.
在第十章那裡.
上帝的靈與他同在.
然後就印證了.
要告訴整個以色列民聽.
他就是一個上帝所選立的君王.
其中那句話.
蘇羅也列在先知中的嗎.
在第一個段落裡面.
是一種詫異.
原來這個男人.
那時候大概三十歲.
這個男人竟然.
以前我見到他高人一等.
原來想不到上帝.
讓他成為先知.
選立了他成為一個與眾不同的人.
這是詫異句.
不過這句話.
後來成為俗語之後.
開始在蘇羅身上這句話.

$^{241}$就成為了貶義詞.
因為在蘇羅的身上.
他活出不了那種聖靈同在.
或者嬰幼的那種.
一個王的那種風采.
反而是他開始.
生命裡面那些焦慮感.
或者人的性.
那個人性.
在他內心裡面不斷地糾纏.
以致到他一路.
十級以下一路走向下坡.
去到第十九章.
經文再用這一句話.
不過他就帶著一些負面的技術.
我們再看.
大衛逃避來見拉瑪見三寶爾.
將蘇羅向他所行的事述說一遍.
他和三寶爾往那約去居住.
有人告訴蘇羅.
說大衛在拉瑪的那約.
蘇羅打發人去捉拿大衛.
去的人見一班先知都受感說話.
三寶爾站在其中監管他們.
打發去的人也受神的靈感動說話.
首先在這裡說到.
蘇羅派人去找大衛.
去捉拿大衛.
打發那些人去捉大衛的時候.
見到有一班先知受感說話.
講預言.
又見到薩姆爾在監管他們.
監管的意思是看著.
薩姆爾和這一班先知同在.
接著蘇羅打發去的那些人.
聽到這些受感說話.
那些打發的人又經歷了.
上帝同在受到感動而講那些預言.
21節有人將這事告訴蘇羅.
他又再打發第二批人去.

$^{281}$他們都受感說話.
蘇羅第三次打發人去.
他們也受感說話.
你看到這班人不是特別什麼先知.
不過這班人是蘇羅打發的.
去捉拿大衛的士兵和手下.
見到先知受感說話.
然後他們就跟著去.
最後蘇羅自己往拉瑪去.
去到西姑的大井問人說.
薩姆爾和大衛在哪裡呢?.
有人說在拉瑪的拿約.
他就往拉瑪的拿約去.
神的靈也感動他.
一面走一面說話.
直到拉瑪的拿約.
他就脫了衣服.
在薩姆爾面前受感說話.
一晝一夜露體躺臥.
因此有句俗語說.
蘇羅也列在先知中的嗎?.
你看到嗎?聖經作者用這一句俗語.
蘇羅也位列先知中的嗎?.
這句說話.
蘇羅因為要捉拿大衛.
要將他置諸死地.
因為覺得這是威脅他王位的一個人.
用盡不同的方法.
打發了三班人去.
那三班人都受感說話.
然後到蘇羅自己去找.
甚至他一路走的時候.
被那班打發的手下更加嚴重.
一路走就一路說.
而且受感的程度更大.
去到什麼狀況呢?.
在薩姆爾面前受感說話.
而且脫了衣服.
赤裸裸地一晝一夜.
整天來受感說話.

$^{321}$你看到後面的描述受感說話.
其實說到蘇羅已經是一個不體面的狀況.
來講上帝的靈感動.
他所說的預言.
但這絕對不是一種追求敬虔.
或不是要上帝的靈.
在他身上彰顯上帝要揀選的受告者.
反而在這裡有一種羞辱的感覺.
一個曾經被上帝選召的人.
如今都被上帝的靈感動.
不過要令他當眾或他所差演的人.
全部都出醜.
聖經的作者他用這樣的記載.
他首先告訴我們.
有些神秘的經歷.
有很多弟兄姊妹都會追求.
算不得什麼.
那個神秘經歷.
不代表上帝與你同在.
即是上帝與你同在.
未必代表你生命真的走向上帝.
要你走的方向當中.
那個神秘經歷.
並不是表達出.
你是一個敬虔追求上帝的人.
即是殺害或追捕大衛的人.
都會受這樣說話.
不過更重要的是.
在一個上帝同在被揀選的人.
如今卻被上帝所嫌棄.
究竟發生什麼事呢.
在蘇羅的生平中.
有什麼出現過.
以致上帝由揀選到最後要厭棄他呢.
我們今天抽取幾個重要.
蘇羅生命的轉捩點.
來看看究竟發生什麼事.
第一個段落我們要看的.
就是在第十四章.
在經文中提到.

$^{361}$蘇羅照著撒慕爾所定的日期.
等了七天.
這個是在說.
撒慕爾原本要為蘇羅去打仗.
蘇羅約定撒慕爾.
一起來準備獻祭.
接著就打仗了.
蘇羅照著撒慕爾所定的日期.
等了七天.
撒慕爾都還沒到吉格.
百姓都已經開始離開蘇羅散去了.
蘇羅說.
將繁祭和平安祭帶到我這裡.
蘇羅就獻上繁祭.
剛獻上繁祭.
撒慕爾就到了.
蘇羅出去迎接他.
要問他的好.
撒慕爾說.
你做的是什麼事呢?.
蘇羅說.
因為我見百姓離開我散去.
你也不照所定的日期來到.
而且非利人聚集在脈絡.
所以我心裡說.
恐怕我沒有禱告耶和華.
非利人下到吉格攻擊我.
我就勉強獻上繁祭.
在第十五章裡面.
蘇羅第一次經歷一個考驗.
這個考驗就是.
約定了他們和撒慕爾.
約定要去打仗.
獻完繁祭之後.
一個很重要的儀式.
尋問上帝的旨意.
然後準備迎戰非利人.
非利人已經大軍壓境.
等了七天.
撒慕爾還沒到.

$^{401}$另一方面.
他焦聚的軍隊.
那些百姓已經齊備.
但是打仗的人就知道.
打仗最重要的是士氣.
如果你叫士兵純粹在等.
等一個獻祭的.
他們視為宗教.
但是他們覺得.
動刀動槍的更加實質.
於是等了七天都沒見.
蘇羅心裡就焦急.
於是他自己起來.
自己上去獻祭.
值得注意.
獻祭不是每個人都可以獻.
在當時裡面.
指定了祭司才可以獻祭.
如果其他人所獻的.
就叫做繁火.
一個不潔之身來到獻祭的時候.
會招惹上帝的和煥.
更加重要的是.
他的位置是打仗的.
他要做他自己的工作.
祭司或者先知.
按著上帝的感動.
既是先知說話.
祭司就去負責獻祭.
如果沒有祭司的獻祭.
他們不能夠來到去做.
某程度這一段經文裡面記載著.
上帝要考一考蘇羅的牌.
他夠不夠耐性.
他夠不夠信服.
來等待撒母耳.
為他們獻祭.
不過這次考牌是胖子的.
蘇羅見到.
那個情勢已經開始急迫.

$^{441}$他特別說到.
百姓開始離開我散去.
說到整個軍隊.
他焦聚的百姓軍心開始動搖.
他見到那些人的狀況.
心裡面來的焦急.
更加重要的就是.
他很冠冕堂皇的.
第十二節裡面說.
恐怕我沒有禱告耶和華.
非理士人下刀吉甲攻擊我.
他說到他怕自己不夠敬虔.
不夠禱告.
這個是蘇羅很多時候的用語.
不過他心裡面是否看重禱告呢.
他反而看重那些百姓離開他而去.
再加上一句.
再加上他就要將這個.
某程度是那種過失.
推在撒母耳身上.
你不照所定的日期來到.
簡單來說就是.
百姓的心又散緩了.
你又所愛弱.
我怕我不夠祈禱.
所以我就勉強來獻梵濟.
不過表面上這件事好像解決了這個問題.
但蘇羅裡面內心的焦慮.
令到他要做些事.
逼著他覺得.
怎樣都好.
硬著頭皮來做.
結果他就.
這件事就惹來上帝對他的那種嫌棄.
甚至後來撒母耳也因為這經歷.
來譴責蘇羅.
他不是信服上帝的君王.
所以第一件事我們要知道的焦慮感.
焦慮感通常會有一種.
伴隨出來的生命態度是什麼呢.

$^{481}$自己很想做些事.
去補救.
但那個補救可能未必是你的位置.
或者你需要去做的.
有時候越做.
可能好心就會幫到忙.
有時候又或者我們.
很急迫的狀況之下.
勉強自己去做一些.
逾越了某些界線的那些位置.
做完之後又會將那個過份.
委過於人.
推卸在身邊的某一個.
本來要去做的那些人的身上.
我不知道在我們生活當中多不多.
在工作職場裡面.
很多時候都會有的.
我見到那個人沒做了.
那個客人又在催那些東西在做.
我唯有我自己做了進去.
來做.
我試過一次差點出事.
話說就是我們以前.
我第一份工作是做電力工程的.
那些文件.
有一些的.
通過某一些的圖織等等.
簽名那個一定要是註冊的工程師.
來去簽的.
大家明白了.
不過那時候急迫.
那個客人又催.
我想說不如我幫他簽.
我打個電話問了他.
OK.
我想我自己代簽的.
在我想下筆去簽的時候.
我姐姐說你做什麼.
你不要越界.
犯法的.

$^{521}$還有你做的事情一定惹來大禍的.
在下筆差點要去做的事情的時候.
被人提一提才收起那枝筆.
很多時候我們人生都會這樣.
我們去到某一些的位置.
我們覺得我有責任.
或者我自己裡面.
肯定是想做些事情.
來做補救.
但是如果我們越界的時候.
可能到最後帶來的結果.
可能被你保持住.
那個結果更加嚴重.
焦慮在這裡面.
令到蘇羅勉強獻上梵濟.
逾越了那個祭司的那條界線.
不單止是這個.
然後經文裡面.
再繼續的來記載.
上帝說蘇羅你是受告者.
你繼續去打仗.
於是上帝就吩咐蘇羅.
去將阿瑪利人進行的殲滅.
這個某程度是說不留生口的.
是一個很極致的舉動.
在撒慕爾記上15章第9節裡面.
蘇羅和百姓聽到上帝的吩咐.
但是他們怎樣做呢?.
蘇羅和百姓卻憐惜阿革.
阿革就是阿瑪利的王.
也愛惜上好的牛羊牛犢羊羔.
並一切美物不肯滅絕.
犯下賤修約的盡都殺了.
即是在說優生的.
靚的留下.
渣的就殺了.
結果就惹來上帝的憤怒.
上帝的吩咐.
他只是聽了一部份.
他沒有做真正上帝要他執行的任務.

$^{561}$然後去到20-21節.
蘇羅對撒慕爾說.
我實在聽從了耶和華的命令.
行了耶和華所差遣我行的路.
擒了阿瑪利王阿革來.
滅盡了阿瑪利人.
百姓卻在所當滅的物當中.
取了最好的牛羊.
要再結合獻於耶和華.
你的神.
這件事是甚麼呢?.
就是說到當他們蘇羅和百姓.
不遵守上帝的命令之後.
撒慕爾就代表上帝去訓示蘇羅.
蘇羅的答辯是甚麼呢?.
我實在聽了耶和華的命令.
我自己有聽上帝的命令.
行上帝差遣我的路.
也擒了阿瑪利人的阿革.
不過上帝吩咐是連王都要清除.
但他說擒了他回來.
淡化了他眼前的錯誤.
然後滅盡了阿瑪利人.
那些錯誤是誰呢?.
不過百姓卻在當滅的物中.
取了最好的牛羊.
不過百姓拿了牛羊是有原因的.
因為他要獻祭給耶和華.
你的神.
整件事是甚麼呢?.
他不用負責任.
他覺得他好的.
或者他自己面對上帝對他的吩咐.
他作為王其實有最終的責任.
那些百姓走在怎樣的路上呢?.
其實他要負上這個責任.
他沒有辦法將責任推給他所領導的百姓.
但在這裡他選擇性地為自己去答辯.
而整件事他所說的.
就是我和你分得很清楚.

$^{601}$他說到我做了我自己的本份.
我做了我的事.
我實在是聽了上帝的吩咐.
不過那些人做不到.
或者他們選擇性聽從.
就將原本他自己的愛惜.
或者他自己的東西.
就輕輕卸下給身邊的人.
到最後那些百姓這樣做是有原因的.
他既是卸下了罪給他.
但同時也在說那些百姓都是為了甚麼呢?.
為了獻祭而已.
他都是想將美好的牛羊獻給耶和華.
你的神.
在這裡說到那個字眼很可圈可點.
因為耶和華的神不只是撒母耳的神.
那個也是蘇羅選召他的神.
但他在這裡就畫清楚我和你.
在這裡整件事他就將自己抽離了出來.
他的焦慮令他做錯了事.
但這個焦慮感也令他自己成為一個自我保護的機制.
他以自己的想法為中心.
以自己成為了最終他要保障的.
那個很重要的最後底線.
所以在這裡你看到蘇羅.
他的焦慮一直推使他.
令他變得越來越自我為中心.
所有事情都是自辯.
所有事情都是為了自己做的好處.
他不肯承認自己在這裡是有錯誤的.
然後到蘇羅被撒母耳一路責備的時候.
到二十四節就開始軟化了.
一路罵的時候.
蘇羅就開始有些改變了.
口風上開始改變了.
二十四節蘇羅對撒母耳說.
我有罪了.
我因懼怕百姓.
聽從他們的話就違背了耶和華的命令和你的言語.
現在求你赦免我的罪.

$^{641}$和我回去.
我很敬拜耶和華.
三十節蘇羅說.
我有罪了.
雖然如此.
求你在我百姓和長老和以色列人面前抬舉我.
和我回去.
我很敬拜耶和華.
你的神.
在聖經作者裡面.
很簡單的幾個字.
你會看到.
在上帝直著撒母耳.
來到罵蘇羅的時候.
他開始知道自己有罪了.
不過這個有罪是不是真的知道.
或者為了自己心悅誠服.
來承認自己的活逆上帝呢?.
在他裡面的對話流露了不是.
其實他最擔心的焦慮是什麼呢?.
他擔心的焦慮就是那些百姓不聽他的話.
他作為一國之君.
但到最後原來他最怕的就是那些百姓不擁戴他.
所以他寧願順著那些百姓的說法.
那些百姓的心意.
以至到他違背耶和華的命令和言語.
你看到他那種焦慮感.
他以別人的看法來定義自己.
即是說他最初的時候知道上帝選照他.
但去到某一些地方.
人生一路走的時候.
他在意別人的眼光.
多於在意上帝如何看待他眼前的任命.
我們有時候不知道會不會面對著這種狀況.
看別人的眼光.
我們嘗試討好身邊的人.
或者那個人所做的事是錯的.
我們怕他不開心.
所以我令自己做很多事情來討好他.
淡化了我們眼前應該做的事或立場.

$^{681}$卻犧牲了上帝給我們的生命的一些命令.
我曾經和大家說過一個經歷.
我在大學時期.
我曾經有一個信主不久後.
有一個信仰很大的衝擊.
就是教我一個很仰慕的老師.
他是一個無神論者.
這個老師每句話上課.
我都很留心聽清楚.
前排第一個.
有一次他知道我信了主.
然後他就責我.
他說基督徒很迂腐.
整本聖經我只覺得最有價值去信的.
就是耶穌.
他很佩服他.
不過聖經其他人一視不顧.
說保羅跟人口水尾沒用的.
我聽到這番說話剛初信.
我的大學老師說這件事的時候.
我心裡覺得我剛剛在讀保羅書信.
剛剛沒多久.
然後他說保羅跟他說是不屑一提的.
這個人.
那時候我口震.
面對著他這樣說.
然後我竟然說出一句.
我為了討好他.
其實我覺得保羅好像沒什麼實力.
說完之後我心裡面.
見完這個教授之後.
心裡面極大不安.
不是說保羅這句話.
而是我好像要討好一個我很仰慕的老師.
以致我正在讀的真理.
正在讀的聖經.
我予以一個負面或者否定的態度.
我們在生活裡面.
我們會不會有時候都是這樣呢?.
我們為了要討好.

$^{721}$以致我犧牲了.
甚至我覺得這是合法地.
令我們違背上帝命令的一個表達方法.
素來這裡就是這樣.
很在意百姓怎麼看他.
他的焦慮感其實都是來自百姓.
以致他經常怎麼介意別人怎麼看他.
他就活在別人的眼目當中.
他不能自己.
去到三十節.
他求殺無疑.
求殺無疑是什麼呢?.
你跟我一起回去吧.
還要在百姓和長老面前抬舉他.
Give me face.
給我一個面子.
讓我很敬拜耶和華.
你的神.
經文裡面說到.
其實到最後.
他和上帝好像沒有關係一樣.
他要人抬舉他.
令他有一個超然的身份.
但是他沒有想過.
整件事他一路一路離開上帝.
漸行漸遠.
耶和華.
我的神變了耶和華你的神.
我們再看最後一個段落.
這個段落.
這個焦慮更加令到.
素羅成為了一個.
好像野獸一樣的人.
十八章第六節至到第九節.
大衛打死了那非理士人.
和眾人回來的時候.
婦女們從以色列各城裡出來.
歡歡喜喜.
打鼓擊盤.
歌唱跳舞迎接素羅王.

$^{761}$眾婦女舞蹈唱和說.
素羅殺死千千.
大衛殺死萬萬.
素羅甚發怒.
不起悅者話就說.
將萬萬歸給大衛.
千千歸我.
只剩下王衛沒有給他了.
從這日起.
素羅就怒視大衛.
你看到.
這個大衛出現.
大衛打贏大衛的疾度.
其實是出現在.
素羅自己的焦慮之後.
是素羅自己裡面有問題.
然後再見到.
有一個好像比他更優秀的人出現的時候.
他就產生一種疾度之情.
而且將這個.
本來和他無仇無怨.
又用音樂治療.
當他發瘋的時候.
彈琴.
令他可以平靜心的一個年輕人.
他竟然以他為敵.
特別是那些婦女.
追星的.
以前都有的.
現在追MIRROR.
以前追那些將士.
其實他不是說素羅廢.
他只是說素羅殺的是千千.
大衛殺死更加多.
他想說的是讚美大衛.
多於貶抑素羅.
不過素羅聽到這句說話就不高興.
就覺得.
現在每個人都讚美他.
他殺死萬萬.

$^{801}$我殺千千.
給他一個王.
這些是發晦氣的說話.
在這裡面.
就因為這些婦女所說.
亦都是別人的眼光怎樣看他.
他就進入了一種焦慮.
將自己和別人來比較.
別人比較別人一回事.
你控制不了別人的.
將自己和別人來比較.
你自己感覺越來越差的時候.
那個就是將自己放在一個不理想的境地.
結果他製造了這種比較之後.
第九節說到素羅.
就勞事大衛.
以他為敵.
然後這一種憤怒.
一路發酵.
甚至要將大衛殺害.
去到第22章6-8.
素羅在基比亞的拉瑪座在垂絲柳樹下.
手裡拿著槍.
眾神僕士納在左右.
素羅聽見大衛和跟隨他的人在何處.
就對左右士納的神僕說.
「比亞敏人啊,你們要聽我的話.
耶西的兒子能將田地和葡萄園賜給你們國人嗎?.
能納你們作千夫長百夫長嗎?.
你們竟都結黨害我.
我的兒子和耶西的兒子結盟的時候.
沒有人告訴我.
我的兒子挑唆我的臣子謀害我.
就如今天的光景.
也沒有人告訴我.
為我憂慮」.
去到這裡.
你看讀到他的心態嗎?.
他將自己變成一個孤島.
全世界都成為了結黨.

$^{841}$來傷害自己的敵人一樣.
去到這個境地.
素羅生命完全沒有安全感.
見到所有人都變成了敵人.
他最親的兒女都成為了他的敵人.
其實不是他的兒女以他為敵.
而是他心裡覺得身邊的人都結黨謀害他.
沒有人為他分憂.
所有人都謀著他的王位.
每一個人都好像在傷害他一樣.
他自己面對一種孤單.
既是對外界有一種憤怒.
但同時對自己最核心的那些關係.
他都覺得很沒有安全感.
焦慮如果沒有正面處理的時候.
我們身邊所有人都成為敵人.
身邊所有人都成為了危害自己安全的那些關係網絡.
覺得自己不由自主.
其實說到素羅.
如果你要數,還有很多片段.
不過我們去到這裡.
讀上去你都覺得好像很暗黑.
很負面.
不過聖經的作者他不是純粹將一個人.
那種困難描述到極致.
然後就令到我們變得很痛苦.
其實在撒慕爾記上或者去到新約.
其實上帝都為這個.
他所選擇的君王提供出路.
不過素羅沒有聽入耳.
出路是什麼呢?.
在第十五章二十二至二十三節.
這個既是一個很重要的責備.
但同時間也是一個對素羅很重要的出路和提醒.
耶和華喜悅凡濟和平安濟.
喜悅人聽從他的話.
聽命聖於憲制.
信從聖於羔羊的自由.
搏逆的罪與行邪術的罪相等.
頑梗的罪與拜虛神和偶像的罪相同.

$^{881}$你既厭棄耶和華的命令.
耶和華也厭棄你作王.
這句話好像一個宣判一樣.
厭棄他作王.
但是這句話裡面其實有幾樣東西.
值得提醒當時的素羅.
也提醒我們今天每一個.
第一樣東西說到.
當你生命裡面面對著一種.
開始你覺得自己開始有點發癲.
你的焦慮感爆燈.
第一樣東西你要想.
你人生裡面你為了什麼來努力.
你聽從的是誰.
所以裡面特別說到.
聽命聖於憲制.
在上帝面前真的重新去聆聽.
在你很混亂的時候.
我鼓勵眾弟兄姊妹.
你真的要停下那些密密手.
你要不停補位.
或者很焦慮而做的那些工作.
放下它.
靜下來.
聽一聽你心靈的聲音.
聽一聽上帝的聲音.
我在這裡也要否白.
我自己有一點點的經歷.
我發覺有一段時間.
我發覺自己發了瘟.
什麼意思呢?.
發覺很忙碌.
很多位要補.
然後有很大量的工作.
覺得眼前有很多事情要去做.
然後就不停地做.
做到一個地步.
去到面對著一個境地.
就是覺得.
我自己在工作.

$^{921}$我覺得身邊的人.
好像既幫不了.
但同時間也覺得.
好像覺得自己孤軍作戰.
情緒不好.
然後看到所有事情都覺得很不入眼.
一直做的時候.
都經歷了一段很長的時間.
但有一晚.
自己靜下來祈禱的時候.
如果我再繼續下去.
我真的隨時會發瘋.
隨時會變成掃羅.
於是.
忙碌的事情都放下.
真的跪下祈禱.
苦服.
然後心裡有很多做的.
很多擔憂的事情就哭出來.
那些情緒都出來.
不過去到最後.
我只是做一個禱告.
我說上帝.
我不可以靠自己.
求你幫助我.
求你赦免我那種專橫.
或者覺得自己可以做到很多事的那種驕傲.
那種自義.
甚至覺得自己比別人更勤力.
更努力的心態.
那段時間開始.
心裡開始靜.
那種暴風開始平靜下來.
我看身邊的人.
看身邊很多弟兄姊妹.
其實很多人都很努力.
不過我只是看到.
將焦點放在自己那裡.
以為這個世界只有自己努力.
放下這些成見.

$^{961}$放下這些焦慮之後.
發覺豁然開朗.
身邊所有人都變得跟你一起同工一樣.
弟兄姊妹.
我不知道在你的生活裡.
有沒有經歷過這樣的焦慮.
在你的家.
在你的職場.
在你與人之間的關係.
你想自己去操盤.
你想自己去做一些事來補救.
你覺得很乏力.
你覺得所有人都不按著你的心意而行.
停下來.
安靜地聽一聽.
聽一聽上帝.
聽一聽你內心.
你自己裡面的呼聲.
聽完之後.
上帝不會只叫你停.
當你靜下來.
你會知道上帝有一些任命.
放你某個位置.
其實祂想你信服.
上帝要你在某個位置裡面.
做上帝要你做的事.
篩選,信服,甘心.
在某個崗位裡面.
繼續努力.
信服的意思就是.
不按自己的議程.
按自己的心意來過自己.
或者將自己的焦慮.
成為動力帶你去做.
你自己想做的事.
而是按著上帝的議程.
來做上帝要我們去做.
素羅的事業.
如果他知道自己好好做好一個王.
不要做了祭司那份.

$^{1001}$可能他會變得.
帶領的時候會少很多困難.
又或者上帝厭棄他做王.
不要叫他不斷去殺害.
或者去糟蹋大衛.
以至他不會走向自己最後的絕路.
生命裡面能夠信服.
比起我們做其他.
更加風高為即更加重要.
所以聽命勝於獻祭.
順從勝於羔羊的自由.
獻多少祭做多少彪炳的業績.
及不上一個很簡單的信服.
在這裡上帝是這樣.
紮臂也是指示素羅要走的方向.
最後素羅大量的焦慮感.
來自人怎樣看待他.
不過人生如果將焦點放在.
人家怎樣看待你.
你會很累.
你會覺得生命要滿足所有人.
對自己的期望.
你覺得是不可能.
但你整天放自己在這樣的境地.
你會迷失.
你會憤怒.
你的情緒就會受著.
人家一句說話.
一個在Facebook上說的話.
你就可以令到你摧毀你的人生.
我們的生命不是由這些說話定義.
上帝才是我們生命的主宰.
他怎樣定義我們.
他怎樣看待我們.
我們的人生以上帝為中心.
我們才真正找到人生的出路.
今天很簡單.
這一段經文.
我們看了素羅的失敗.
不過在他的失敗當中.

$^{1041}$我們都看見.
其實上帝不是想放棄他.
厭棄他的作王.
不等於要放棄一個人.
不過似乎他的焦慮.
令到他越走越差.
唯有我們真是願意停下來聽.
信服上帝.
學習以上帝為中心.
我們的生命才可以抗衡這種.
令到我們發癲的那種焦慮.
願主繼續幫助我們.
帶領我們走面前的路.
\newpage



\section{詩篇 119:105-112}
\label{sec:cerRyXhHqiE}
\textbf{2024年8月25日崇拜直播｜張培良傳道｜活在上帝的光中｜詩篇119\_105-112}
\newline
\newline
連結: \href{https://youtube.com/watch?v=cerRyXhHqiE}{\texttt{https://youtube.com/watch?v=cerRyXhHqiE}} ~~~~ 語音日期: 2024-08-26
\newline
\newline
\hyperref[sec:Ef4j6FO5fXg]{\small{< < < PREV SERMON < < <}}
~
\hyperref[sec:index_chronic]{\small{[返順時目]}}
~
\hyperref[sec:lCXbCQNDTgA]{\small{> > > NEXT SERMON > > >}}
\newline
\newline
詩篇 119:105-112
\newline
\begin{longtable}{cl}
\hline
\hline
章節 & 經文 (和合本修訂版)\\
\hline
119:105 & \begin{tabularx}{0.7\textwidth}{X} 你的話是我腳前的燈，是我路上的光。 \end{tabularx} \\ \\ \relax
119:106 & \begin{tabularx}{0.7\textwidth}{X} 你公義的典章，我曾起誓遵守，我必按著誓言而行。 \end{tabularx} \\ \\ \relax
119:107 & \begin{tabularx}{0.7\textwidth}{X} 我極其痛苦；耶和華啊，求你照你的話將我救活！ \end{tabularx} \\ \\ \relax
119:108 & \begin{tabularx}{0.7\textwidth}{X} 耶和華啊，求你悅納我口中的讚美為甘心祭，又將你的典章教導我！ \end{tabularx} \\ \\ \relax
119:109 & \begin{tabularx}{0.7\textwidth}{X} 我的性命常在我手掌中，我卻不忘記你的律法。 \end{tabularx} \\ \\ \relax
119:110 & \begin{tabularx}{0.7\textwidth}{X} 惡人為我設下羅網，我卻沒有偏離你的訓詞。 \end{tabularx} \\ \\ \relax
119:111 & \begin{tabularx}{0.7\textwidth}{X} 我以你的法度為永遠的產業，因這是我心中所喜愛的。 \end{tabularx} \\ \\ \relax
119:112 & \begin{tabularx}{0.7\textwidth}{X} 我的心傾向你的律例，謹守到底，直到永遠。 \end{tabularx} \\ \\
[1ex]
\hline
\hline
\end{longtable}
$^{1}$而是經訓的時間.
今日我們選讀的經文是舊約的詩篇.
舊約詩篇第119篇.
以下我們會選讀105至112節.
在你們舊約聖經第805頁.
經文也印製在秩序表裡面.
請聽我讀出詩篇119篇105節.
你的話是我腳前的燈 是我路上的光.
你公義的典章 我曾起誓遵守.
我必按誓而行 我甚是受苦.
耶和華求你照你的話將我救援.
耶和華求你閱納我口中的讚美為貢物.
又將你的典章教訓我.
我的性命尚在危險之中.
我卻不忘記你的律法.
惡人為我設下網羅.
我卻沒有偏離你的訓詞.
我以你的法道為永遠的產業.
因這是我心中所喜愛的.
我的心專向你的律例永遠遵行.
一直到底 聖經獨筆.
第一節目平安.
先介紹一下自己.
我是張培良 全道.
來自澳門百甲巢浸信會.
今天很開心代表百甲巢浸信會.
今天不只是我來的.
剛剛大家敬拜隊當中.
也有一班我們年輕的弟兄姊妹.
所以希望今天帶一些活力過來給大家.
這次是我第二次來到武會這裡講道.
來到這裡講道.
常常帶著一種緊張但又興奮的心情.
因為可以看看上帝在武會當中的工作.
到底是怎樣來開啟開展.
今天來到當中.
我希望抱著學習的心態.
但同時帶著聖經的說話來彼此勉勵.
正如當我們進入教會.
我不知道大家有沒有這樣的感覺.

$^{41}$就算你是什麼心情也好.
你進來教會總是有一種被光照的感覺.
可能今天星期天有點單調.
今天早上外出跟兒子,老公,老婆吵架.
進來後聖師播放戲你不敢再跟他們吵架.
通常會先唱詩歌.
唱完詩歌通常不能再吵架.
因為心情已經被上帝充滿調整.
正如我們聽完敬拜.
多謝獻奏的戲樂小組.
很優美的旋律.
聽完之後整個心情都很平和,很平安.
很多時候我們很多人.
不知道在有當中有沒有第一次來教會.
或者第一次進入教會.
或者你不是第一次.
但你記得你第一次進入教會.
都經歷過這種感覺.
我們很希望這種上帝的平安.
上帝的光照能夠常常在我們生命當中發生.
所以今天我揀選了這個主題.
就是「活在上帝的光中」這個主題.
很希望藉著今天我們所讀的詩篇經文.
去彼此勉勵.
如何能夠活在上帝的光當中.
並且這份光不是只有一刻.
不是只是星期天回來聽那十幾二十分鐘.
而是一生都可以帶著我們.
我不知道大家想起光.
你想起一個什麼畫面.
我想光是有很多狀態的.
例如我今天一出來.
走過來這個旺鎮.
我覺得今天的太陽太強了.
最近太陽太強了.
有些警報指外線光很高.
這是其中一種光的狀態.
但光不是只有一種狀態.
晚上你以前燙著一支火柴.
點著的光.

$^{81}$或者回到比較遠舊的年代.
我們用一些蠟燭.
油燈點著.
在晚上當中燙著燙著的那種.
都是光.
所以當我們活在上帝的光中的時候.
我不知道你立刻想到的是什麼畫面.
可能你想到很熱很溫暖的畫面.
但我想今天我們所讀的經文.
可能未必是這種畫面.
今天我們所讀的經文.
有一節大家很熟悉的.
上帝你的畫是我們腳前的燈.
路上的光.
這個畫面到底是否說.
我們回到日間.
太陽很曬.
我想不是.
這個畫面比較像在晚上.
黑漆漆.
然後古時的人想回家.
看不到路怎麼辦.
只好拿起一盞油燈.
或者火把之類的東西.
然後我們現代人很難感受.
因為晚上不會拿油燈.
最多拿手機.
但我知道舞會.
也有上行者.
團計會行山.
我不知道這裡有沒有上行者.
有沒有行過夜山.
或者晚上去行山.
你看不到路.
然後你會打開你的電筒.
或者拿手機燈亮了也好.
然後會怎樣.
你就會發覺前面那一格的路.
就很光猛.
但遠一點會怎樣.

$^{121}$有些人會覺得.
越光還會更光.
但你可以想像到畫面就是這樣.
當我們這些精文化.
在你腳前的燈路上的光的時候.
其實不是在說光猛猛.
是在說在黑暗當中.
你拿著一盞油燈.
眨下眨下.
你看不到最遠的地方.
但起碼.
你看到你眼前的那一步.
可以怎樣走好.
不要在眼前的石頭.
卡到你.
我想這段經文一開始.
給了我們這樣的背景.
看到原來在上帝光中.
有一種狀態就是在黑暗當中.
上帝一步一步.
帶我們走前面的路.
這首詩歌.
既然我剛剛說到.
他是在黑暗當中.
去照明我們.
黑暗我們也要理解一下.
也要描述一下黑暗到底是什麼.
因為很多人會將.
我這段時間真的很黑.
跟黑暗會有些搗亂.
我想兩種有些不太一樣.
如果說很黑的話.
有什麼很黑呢.
我本來今天想來舞會講到.
都很帥氣.
怎知道.
昨天來到這裡.
突然眼睛有毛巾還是什麼原因.
我老婆說你的眼睛腫了.
大家不要突然看我的眼睛.

$^{161}$有點尷尬.
我昨晚已經.
滴眼藥水.
今天想帥一點.
但不要緊.
待會直播時不要太近.
這樣也很黑.
如果早前.
要論黑仔.
不知道大家記不記得.
幾個星期前有一個軟件壞了.
香港應該也受影響.
一個軟件壞了.
全球的飛機也受影響.
如果你在機場.
你也會說為什麼我不黑仔.
我想像有些黑仔的畫面.
如果今天告訴你.
你手機裡的資料突然離開教會.
就沒有了.
舉例.
你馬上想我去了哪間教會.
這些位置都很黑.
但是.
這些黑和黑暗有點不同.
黑暗是什麼呢.
我想.
這裡.
我應該要下去.
黑暗是什麼呢.
可能大家不太清楚.
但我希望吃住黑暗.
讓你覺得黑暗看不清楚.
黑暗是什麼呢.
我想黑暗是一件事.
我去描述.
在人生裡的黑暗.
有時候是突如其來.
你本身今天可能健健康康.
今天在教會.

$^{201}$彼此打招呼很開心.
但突然之間.
可能是身體的疾病.
可能是突如其來的悲劇.
一個傷痛.
你突然之間人生可以一蹶不振.
可以突然之間.
覺得自己的人生是永不翻新.
大家記不記得死亡當中.
有一個叫做很薩瑪利亞人的故事.
我們經常都記得那個.
很薩瑪利亞人.
但有沒有人很特別印象深刻.
那個人.
有沒有想過那個人其實在上一秒.
還很開心.
敬拜完上帝.
可能很開心.
帶著東西走了.
誰知道去到山路的時候.
突然之間被人打劫.
打到半死.
然後他在地上什麼都做不到.
他只能等候經過的人去救他.
我想真正的黑暗.
有點像是這樣.
黑暗可能是.
城市人特別有一種.
很長久地都覺得自己的生命.
很疲乏.
每天都覺得很累.
甚至有些恐怖一點.
我們這個暑假.
很多人都會暑假去旅行.
去完旅行回來.
通常人們說充電充電.
不知為什麼越充都越累.
去旅行都很累.
行程為什麼排得這麼密.
每個景點都要去.

$^{241}$所以去完旅行回來.
是不是真的讓你的生命.
重置得力呢?都未必一定是.
常常這份疲倦.
都是我們生命當中纏繞我們.
我聽過一個.
弟兄姊妹分享裡面.
他說那種累是什麼呢?.
就像肩膀一直都堅擔著很多.
其實不應該是我擔的擔子.
我們生命裡面.
承擔著很多人對於.
我們那種錯誤的期待.
而我們又覺得我們應該要滿足.
這些期待.
所以每一天都覺得很疲倦.
每一天都覺得很無力.
每一天都說.
我要一些me time.
但那些me time是否真的幫到我們呢?.
黑暗亦都可能有.
生命很多的問題.
甚至可能那個問題都未發生.
但我們已經在恐懼.
我們已經在憂慮.
我們已經在搞到頭破血流.
未發生但我們已經害怕.
這是什麼?.
這是很多心理疾病.
憂慮,抑鬱.
根據世界衛生組織.
計算.
其實未來在人類當中.
最大的殺手在不是癌症.
是精神疾病.
是抑鬱疾病.
不是好意思這樣說.
不過根據統計.
在城市人裡面.
三個就有一個.

$^{281}$這裡有多少個呢?.
可能我是第一個.
所以那樣東西.
黑暗不是離我們很遠.
黑暗不是那麼一件.
是不可探知的事.
隨便找一個弟兄姊妹來.
她都可以分享給你聽.
她生命裡面經歷什麼黑暗和痛苦.
但今天.
我們來到上帝的殿當中.
上帝告訴我們一件什麼事呢?.
在聖經當中.
我們見到很多人在黑暗的裡面.
聖經當中都有很多無助.
曾經無助的人.
但他們卻是和上帝相遇.
這個相遇的過程.
令到他們的生命見到那個光.
感受到那個光.
然後他們能夠活在這個光當中.
但今天不停止在那裡.
我最後想和大家說.
你當是彩蛋也好.
今天上帝不是叫你活在光中.
今天上帝要你成為光.
嘩!你從來沒有想過.
活在光中都已經夠開心.
我還可以成為光.
到底是怎樣的呢?.
帶大家看看聖經的人物.
我不知道電視機能否看清楚.
這幅圖.
這幅圖我大概十年前.
看到了.
我很喜歡,一直留著.
後來找出來.
裡面說到.
你有沒有看到聖經裡面很多人物.
那些我們覺得很出名.

$^{321}$或者我們很認識他們.
我們看聖經就是看這些人物.
亞伯拉罕,摩西,約瑟,諾亞等等的人物.
嘩!說起他們就覺得很厲害.
上帝很使用他們.
如果你熟讀聖經的時候.
你知道這些所有人物背後都有他們的黑歷史.
他們都有他們人生的黑暗時間.
他們有他們人生無力軟弱的時候.
就等同於你和我一樣.
我們試一下想一想.
第一個,我今天說三個.
因為太多個了,說不出.
亞伯拉罕,大家看到他有個「朵」.
這個「朵」叫甚麼?.
叫甚麼?大膽一點.
信心之父.
說信心之父我們都要有信心.
不要叫信心之父.
好像對他沒甚麼信心.
信心之父的亞伯拉罕.
這個信心之父.
不是隨便一個人可以拿到信心之父這個「朵」的.
你和我不會拿到這個「朵」的.
他做了甚麼呢?.
你發覺在上帝對他的生命.
他要離開他的本族.
父家,去到上帝應許之地的時候.
過程中間發生了甚麼事?.
他因為自己的膽怯.
他差一點,還差一點沒有.
差一點賣了自己老婆幾次.
真的,你看聖經.
他真的賣了自己老婆.
賣了她,拿回一條命也好.
到最後上帝出手.
救回她.
是這個老婆,沒得轉.
這個信心之父.
也是後來上帝應許.

$^{361}$他要有一個上帝應許的兒子.
他也很想要這個兒子.
但到最後.
他用他自己的方法要了這個兒子.
大家記得創世記所講.
找個妾侍阿環過來.
這些聖經經常出現的片段.
到最後.
他自己搞出來的兒子.
搞到家嘈屋閉.
這個家差不多散了.
起碼先散了一半.
先散了一半.
這個就是信心之父.
摩西這個.
是世經當中最偉大的人物.
帶領以色列人離開埃及的領袖.
他過往是一個.
皇宮裡面的貴族.
但當他年輕的時候.
他以氣用事殺了人.
到最後怎樣呢.
他要逃跑離開埃及.
去到曠野放羊.
放了幾十年羊.
去到他八十歲的時候.
上帝再來到呼召.
這位摩西.
在座有沒有人八十歲.
可以舉起手嗎.
或多於八十歲.
有沒有.
看到了.
八十歲以上.
果然沒有的會很年輕.
所以很有勇氣.
承認自己八十歲以上.
我想跟八十歲以上或就八十歲的第一知名者說.
你不要以為退休.
八十歲才來了.

$^{401}$摩西就是這樣.
八十歲的時候.
上帝要來.
給你一個好工作.
帶你以色列人離開埃及.
摩西有什麼反應.
好啊 這麼好.
快點吧 我做吧.
當然不是.
摩西怎樣.
他擔驚受怕.
耍手棍 扣我住.
扣我便是 請不要找我.
一個最偉大的領袖.
說自己不會說話.
請不要找我.
用盡一切方法.
找你哥哥幫忙也不能走.
到最後他真的肯.
用了很多神蹟騎士.
真的帶領了以色列人離開埃及.
結果是什麼.
那班帶著的子民天天都嫌棄他.
天天都說.
為什麼帶我離開埃及.
我去埃及吃火鍋不好嗎.
聖經真的這樣說.
我回去要吃肉鍋.
我回去為什麼不吃火鍋.
如果你是摩西真的會吐血.
再說最後一位.
大家一定很熟悉的.
那個彩衣少年.
藥室.
他是掌上明珠.
我不知道男士這樣形容對不對.
總之就是被爸爸捧在手心上.
但是卻在人生當中.
他被兄長嫉妒.
到最後賣他珠仔.

$^{441}$賣到埃及.
到最後做了奴隸.
不止還要被人冤枉.
坐牢.
他很努力.
坐牢都很努力.
然後幫助了很多人.
最後幫助的人離開監獄.
忘記了他.
我想如果這些事.
今天發生在我們.
《時但一位》的節目身上.
我們都足夠黑暗.
我們都應該可以一蹶不振.
我們都應該有理由可以怨天尤人.
可以說上帝你是不是作弄我.
上帝你是不是玩我.
但正正就是這些人物.
他們因著遇見上帝.
上帝的說話.
上帝的光.
照了他們的生命.
令到他們的生命軌跡.
不是走向黑暗的那一邊.
而是走向上帝引領.
那些斬下斬下的光的那一面.
正如我們今天所讀的.
這一篇詩篇.
詩人第一節就說了.
105節就說了一個真理.
他說我們的生命要活在上帝的光中.
唯有這樣我們才能夠在黑暗和黑暗.
以及這些危險的裡面.
我們能夠放膽的繼續向前.
我想我們今天生命都很需要這一份的光.
或者在進入聖戰之前.
不如我們有少少感受一下這個光的時間.
不如大家可以閉上眼.
我們一起去感受一下.
既然是詩篇我們應該有些詩意.

$^{481}$我們閉上眼.
我想你可以想像一下.
你前面都是黑暗的.
一片的黑暗.
無力.
你不知道向哪個方向走.
但同樣在這個時候.
我邀請你在你的腦海.
或者在你心裡.
你想起一節上帝對你所說的.
聖經,經文.
或者一句說話都好.
在黑暗當中.
你想起這句說話.
然後你發現這一句說話.
慢慢慢慢在黑暗當中.
變得越來越光.
甚至這份光.
讓你的眼前一亮.
慢慢慢慢讓你沉浸在其中.
讓你一起祈禱.
神在你的說話就是我們生命的光.
沒有你的說話.
我們的生命當中.
我們永遠不知道我們何時會陷入黑暗.
但唯有當你的說話進入我們生命當中.
你不單帶給我們溫暖.
亦帶給我們光芒.
照亮我們人生眼前的這一步.
讓我們繼續有勇氣向你踏前.
繼續向原理的方向前行.
願你的話語繼續指引我們.
幫助我們.
但願以下的時間.
讓我們的心被你打開.
祈禱感恩奉主名求.
聖經的說話來幫助我們.
今天我們所讀的這篇詩篇.
他一開始就描述神的說話.
就是那個光.

$^{521}$他在這裡告訴我們.
他用了很多不同的形容去描述上帝的說話.
例如他講了律例.
典章 法道.
上帝的訓詞等等.
但其實都是講同一句說話.
就是上帝的說話.
上帝的說話就是生命的光.
光照著這個黑暗的前路.
這個黑暗的人生.
但在這個段落當中.
我們如何去深入一點去理解.
上帝的說話成為生命的光.
在這篇詩篇中.
有一些部分想跟大家分享.
很有意思的.
因為詩篇119篇.
大家可能都聽過.
其實是一首字母詩.
就是希伯來人用他們的希伯來字母.
去寫成這首詩篇.
所以有很多節.
但很容易記住.
因為希伯來人有22個字母.
然後每一個字母寫了8節.
所以你看到聖經每一個段落都是8節.
所以8乘22就是多少節.
176節.
大家的計算很快.
所以你就知道119篇的詩篇.
總共有176節.
每個字母有22個字母.
每個字母8節.
而我們今天所讀的這個段落.
那個正是第14個字母.
是一個N字開頭.
很有趣的.
和英文的字母第14個都是N字.
所以有些類近.
是一個N字.

$^{561}$不過希伯來人叫做Lum.
Lum這個字.
其實在這個段落裡.
我想和大家分享兩個字.
在這個N字開頭裡.
有兩個很特別的字.
我相信是詩人特別用這兩個字.
放在這段經文當中.
去提醒我們到底生命的光.
到底上帝的說話.
是怎樣幫助我們呢.
第一個字和大家分享.
這個字的原文是Nadabah.
這個字.
其實是說Free Will Offering.
或者是一種自由意志.
願意甘心的拜賞的意思.
經文說到.
是我嘴唇上的甘心的拜賞.
如果是用這個和修版.
他就會說是我口中的讚美.
成為甘心祭.
甘心祭的意思是.
我是願意自由意志選擇去拜賞.
這是第一個字.
出現了自由意志這個字.
第二個字.
這個字其實我們很熟悉.
第109節說我的性命.
第二個字叫做Nephes.
希伯來人Nephes這個字.
我經常聽到這個字.
我就想起我們放假很想看的.
那個叫什麼.
叫Netflix.
我們喜歡看的叫Netflix.
不過沒有任何文獻.
說到他其實是抄兩個字.
不過Nephes這個字.
其實是說生命.

$^{601}$性命或者是說靈魂.
在聖經當中有不同的翻譯.
但是在這裡.
這個Nephes這個字很有趣.
你一定看過.
因為創世記的第一章.
上帝創造所有生命的時候.
所有有生命生氣的活物.
都是用這個字.
Nephes.
創造人類的時候.
都是用這個字.
叫做Nephes這個字.
但是後來當上帝創造了人類之後.
大家記不記得創世記第二章.
第七字裡面.
上帝做了一個動作.
他將一口氣吹進人生命當中.
然後聖經說什麼.
他就成為了一個有靈的活人.
Nephes.
Haya這個字.
所以如果用我理解.
你怎麼理解這段經文.
你和我都是有生命的人.
我們有呼吸.
Nephes這個字在猶太人當中.
其實有另一個意思.
是解喉嚨的意思.
和呼吸道很有關係.
所以聖經裡面說.
我的呼吸.
我的性命都是說什麼呢.
就是在說.
我有氣頭.
就是在說Nephes這個字.
但是今天.
當聖經在說Nephes.
或者Nephes這個字.
你怎麼讀都好.

$^{641}$他在說我們原來有性命的意思.
但是原來今天.
當你生命裡面有性命.
你有自由意志.
都好.
原來都還不夠.
施人告訴我們.
原來一個有性命的人.
一個有生命氣息的人.
原來如果沒有了上帝的說話去指引的話.
和創世記所說.
你和其他的生物.
如果沒有意思.
俗一點說.
你和那隻貓.
你和那隻狗.
你和那隻獅子.
你和那隻老虎.
你和那隻大象.
這些創造的生物是沒有分別的.
唯有什麼呢.
當上帝的說話.
當上帝的靈氣.
當上帝的呼吸.
進入到你生命的時候.
你才成為一個有靈的活人.
其實作者某程度上是在用這個畫面.
第一次你回想一下.
回想一下自己的生命.
如果今天沒有上帝的說話.
沒有上帝的真理進入你人生當中.
其實我們就是形形異異地過你一生.
這個世界怎樣活我就怎樣活.
人家怎樣做我就怎樣做.
其實有時候我們真的.
用今天的說法.
那隻狗過得比我們好.
我相信大家應該明白這件事.
寵物過得比我們好.
我們比其他的Netflix.

$^{681}$我們還過得差.
但這不是上帝的心意.
上帝告訴我們.
原來上帝要讓我們活在他的光中.
透過什麼呢.
透過上帝的說話.
透過上帝的導.
讓Netflix有上帝的真理.
有上帝的靈在裡面的時候.
我們的生命就能夠在不同的境況當中.
都活出上帝那種光芒.
所以倒過來是什麼意思呢.
如果你生命裡面沒有上帝的話.
你變成什麼.
你不是變成你的寵物.
不要搞錯.
你變成跟動物沒有分別.
所以我們要跟自己說.
我不要活成一隻動物.
雖然我的狗過得很好.
我的貓也過得很開心.
但我不是要成為牠.
牠只是我的很開心的伴.
牠是我的心靈安慰的一種方式.
但牠不是我要過的生活.
上帝要透過他的說話.
幫助我們人生.
指引我們.
這段經文很有意思的就是.
剛剛我說了一個net devoir.
這個字解釋自由意志的意思.
其實自由是我們今天人類裡面.
我想我們很需要的.
特別是我們教會做很多年輕人工作.
很多年輕人說我要自由.
你不要逼我選.
做父母的應該很明白.
你不要逼我選.
可能到最後還是選一個.
但我還是要自己選.

$^{721}$我要自由意志.
但最有意思的是.
斯人竟然將自由意志.
free will這個字放在什麼段落.
放在上帝的律法.
放在上帝的說話.
放在上帝的訓詞的段落當中.
你覺不覺得很衝突.
我們心裡面想的那種自由.
是無拘無束.
沒有規條.
我喜歡回來教會穿什麼就穿什麼.
做什麼就做什麼.
我人生想怎樣就怎樣.
這個才叫自由.
但斯人告訴我們.
原來真正有靈的人.
一個活在上帝光中的人.
不是這樣的.
活在上帝光中的人.
有上帝的說話.
有上帝的指引去光照他.
是可以過一個.
可能是自律.
可能是有上帝的原則.
可能有上帝的指引的生命.
這個是斯人所說.
然後斯人在這裡當中.
繼續告訴我們.
這種自由意志.
為我們的生命帶來了什麼.
就是帶來了後來我們越來越喜愛.
上帝的律法.
這種衝擊我會這樣理解.
就是斯人覺得.
在他的理解裡面.
上帝的律例.
上帝的說話.
他的誡命.
不是他生命的束縛.

$^{761}$反而成為了他人生黑暗當中的.
指引.
這個就是斯人想表達的意思.
而之後.
另一個意思就是說.
斯人認為上帝的說話.
就是他生命黑暗當中的火把.
那個油燈.
那個腳前的燈.
路上的光.
我不知道今天.
你想一想.
你是否這樣去看上帝的說話.
還是你覺得上帝的說話.
是你人生當中的束縛呢.
如果我們.
認為上帝的說話是我們人生的束縛的話.
我相信你很難經歷到.
活在上帝的光中.
這個意境.
因為上帝的光.
和上帝的說話不是分開的.
兩件事.
是同一起發生的.
所以今天我們很蒙福地.
我們回來教會的群體裡面.
我們能夠不斷去聆聽上帝的說話.
一起去經歷上帝的說話.
一起學習實踐上帝的說話.
讓我們能夠經歷.
活在光中的這個境況.
但是這個境況.
是不是等於之後很容易呢.
不是的,在1806節裡面.
的時候,詩人說.
他起誓和上帝.
立誓,我要遵行你的指引.
你的教導.
但是接著他說什麼呢?.
1807就說,我甚是受苦.

$^{801}$原來今天活在.
上帝的光中.
活在上帝的說話裡面的時候.
你不一定是順風順水.
你每次都很順利,很順境.
走出去,那一塊雲就剛剛遮住你.
上面你就不怕曬.
走出去買東西喝,最後那一罐就被你買到.
所有家裡的事情.
都很平安很順利.
是的,有些人可能是這樣.
但是我想更大的境況是.
當你去實踐.
當你說我要來遵行上帝的說話的時候.
你遇到的困難挑戰.
比以前更加大.
你遇到困難的挑戰.
比以前可能更加入心.
但正正是這種境況.
讓你更深的感受到.
上帝的光.
到底是什麼.
當受苦的境況裡面.
斯人他怎樣描述呢?.
在詩篇裡面繼續說,斯人就按著.
他對上帝的認識來祈禱.
基督徒可以怎樣做呢?.
在困難當中就祈禱.
他怎樣祈禱呢?.
他按著對上帝的認識來祈禱.
從一開始他說上帝的說話.
是那麼有能力,是我們的光的時候.
所以他說什麼呢?.
他說上帝的說話必定可以將我.
在黑暗當中救活過來.
必定可以幫助我.
在這個性命危難當中.
拯救我的生命過來.
斯人這句說話有什麼重要的提醒呢?.
弟兄姊妹,我想告訴你.

$^{841}$上帝的說話不是.
頭腦的知識.
上帝的說話不是我們.
認識了多一些美善的概念.
上帝的說話也不是.
哲學的思維.
上帝的說話.
是真實的能力.
是真實賜給我們.
生命盼望.
那種的幫助.
108節斯人.
我想107節和108節中間.
已經經歷了一段時間.
雖然是寫詩,但已經經歷了一段時間.
107節祈願禱之後.
108節斯人經歷了神的恩惠.
然後就去到我們說的那個字.
斯人用什麼?.
他用自由意志獻上的讚美.
來獻情給上帝.
自由意志的讚美.
Free will offering.
什麼意思呢?.
就是那個讚美.
剛剛我們唱詩不是被迫的.
不是機械性.
不是回來教會只是一種習慣.
而是我們用了.
帶著自願.
帶著喜樂.
帶著生命的感恩.
來獻上讚美.
如果今天.
你回來教會要預備什麼.
就是要預備.
生命當中.
用你的自由意志.
用你的自願.
用你的自由意志.

$^{881}$我今天選擇星期日的早上.
我不是去喝茶.
我不是和朋友去喝啤酒.
我自由意志.
我選擇來到這裡敬拜上帝.
這是一個何其美的選擇.
這是一個何其蒙福的行動.
斯人正正就是一種回應的時候.
同時間他繼續求上帝.
將他的說話教導他.
豐富他的人生.
去走過每一個黑暗.
每一個受苦的境況.
109節到110節.
我們會有一個.
非常非常非常.
深刻的.
109節到110節.
你會發覺這個.
這個軀體這個性命.
這個作者.
他一直陷在一些危險當中.
這個危險指的是什麼呢.
是不是有人用刀.
把他抗衡住呢.
我想可以有很多的.
我們剛剛也說了很多黑暗.
但這裡他講的可能特別是指.
有很多惡人.
有很多邪惡.
他們設下了一些陷阱.
讓你走到偏離上帝的道路上.
今天在我們生活當中.
在我們生命裡面.
原來有很多上帝.
不喜悅的東西.
魔鬼撒旦或惡人所設下的網羅.
是要讓你離開.
上帝的道路.
因為害怕,因為恐懼.

$^{921}$因為無望等等原因.
最後離開上帝要你走的路.
魔鬼撒旦.
最終的目的就是這樣.
但我們看到作者的回應是什麼呢.
作者很堅定地.
來講.
我不忘記你的律法.
我不偏離你的訓詞.
這兩句的意思在說.
作者認定了上帝.
相信上帝的說話.
會成為他一生的保守.
會走過人生.
每一個的年頭.
甚至到最後.
要成為他生命要繼續傳承下去的產業.
剛剛在座的一些舉手.
我有八十多歲.
或是我們不同年紀也好.
我相信在你的生命當中.
你也經歷過上帝光照我們.
帶領我們人生走過一個又一個的難關挑戰的時刻.
所以今天我們坐在這裡敬拜上帝.
甚至我們邀請人回到教會.
想將上帝的恩典.
上帝的救贖.
能夠傳承下去.
成為產業.
當使人遵守上帝的說話的時候.
神的光就來到一度.
引領他的生命繼續向前.
他就活在上帝的光中.
如果用新約的說法.
堅定在生命裡選擇上帝的話.
其實就是你堅定選擇耶穌基督.
因為耶穌基督就是上帝說話的顯現.
我們選擇耶穌基督.
就是活在光中的秘訣.
剛剛我已經說了.

$^{961}$你認定了耶穌基督.
你選擇了耶穌基督.
不代表你的人生一帆風順.
無風無浪.
而是當你經歷著這些黑暗.
和你身邊的人那些微信者一樣.
你發現上帝的光就在當中對你說話.
就在當中指引你走前面的那一步.
讓你繼續有勇氣.
繼續有信心可以繼續為主去努力.
但是剛剛說了.
不單止這麼簡單.
來到新約的時候.
耶穌對著那些跟從他的門徒說什麼.
你們是世上的光.
光要照在人的面前.
叫他們看見你的好行為.
就將榮耀歸給天上的父.
這一節的經文我們可能都很熟悉.
但是弟兄姊妹.
如果我們只是將重點放在好行為上.
你可能是抓錯重點了.
因為我想很坦白的跟大家說.
好行為是不會令你的生命歸榮耀給父神的.
起碼不是直接關係.
因為未信主的人他也可以生命有好行為.
所以這段經文說.
藉著你們的好行為.
將榮耀歸給天父背後的光.
成為光的重點是什麼呢.
我想就像我們剛剛所說的.
那個光不是你的光.
那個光是上帝的光.
是神說話的光.
是由你生命裡面認定上帝的說話.
去選擇耶穌基督而來的這份光.
如果基督徒失去了這份對上帝的堅持的話.
好行為是不會帶來真正的光.
好行為很多時候帶來的可能只是驕傲.
可能帶來的只是覺得.

$^{1001}$你看看我做得多好.
你看看我是一個多麼虔誠的基督徒.
正如我們剛才所看到的文字範圍長一樣.
所以我想說.
我們基督教也會說很多善行.
但基督教的善行.
和其他宗教最大的分別在於什麼呢.
在於一切都不是出於我們自己本身.
我們所做的一切的事奉.
我們所做的一切的好行為.
一切的擺上.
是因為神的光.
是因為神的救恩在我們生命當中開啟.
而我們作出自由,意志,感恩的回應而開始.
所以讓上帝在我們身上的恩典.
成為一個很美麗的產業傳承下去.
這就是基督徒的事奉.
所以正正因為這樣.
基督徒所做的善行.
有時候會帶著很多人的限制.
你不要以為我們做的好事一定全部都很好.
也有很多一面兩體的.
但我們不會恐懼而不去做.
因為我們知道基督徒所做的.
不只是停留在人的限制.
不只是停留在你和我的力量的有限.
而是他在分享神的愛.
分享神的榮耀的一個行動.
所以我們可以很有信心.
因為上帝會包底.
因為上帝會帶領我們.
所以我會說基督徒的善行.
拒絕驕傲,否定一切的個人主義.
記不記得耶穌的名字叫什麼.
叫以馬來尼.
神與我們同在的意思.
耶穌都說得這麼白了.
我和你同在.
你還怎麼可以對人說.
不關上帝的事,是我做的.

$^{1041}$我們剛剛所看到的聖經人物.
我想給大家看看他們的生命的光在哪裡.
透過聖經裡面看到.
亞伯拉罕,我剛剛說了.
他的生命有很多黑歷史.
但亞伯拉罕的光.
不只是單單說他帶領整個民族.
整個家族去到這個應許之地.
不是他獻上耳朵殺這些都很厲害.
不是他成為了後來的信心之父等等.
我覺得亞伯拉罕的光.
如果你看聖經裡面有很多片段.
我其中覺得有一個片段.
就是當他和他的侄羅德.
去到約旦河的時候.
因為兩個家族太大了.
大家互相攻擊.
在這個時候亞伯拉罕選擇什麼呢.
他選擇不爭執.
讓出選擇權給羅德選擇上哪.
然後羅德就選擇了最漂亮的平原.
最漂亮的地方約旦河東.
到最後變成了索多瑪和俄摩拉.
然後亞伯拉罕就自己選擇.
去到比較山地一點的地方.
他生命裡面選擇不去爭執.
不去爭取自己想要的東西.
亞伯拉罕的光可能是創世記後來描述.
亞伯拉罕為了索多瑪和俄摩拉向神講價.
因為羅德住在那裡.
他不想上帝滅那個城.
所以亞伯拉罕好像在街市和上帝買菜一樣.
便宜一點可以嗎.
再便宜多兩塊可以嗎.
再少一個可以嗎.
十個可以嗎.
你自己回去看那段經文.
真的好像和上帝在街市買菜一樣.
到最後好像應該可以.
應該有十個的.

$^{1081}$好了差不多了.
收工了.
如果不是上帝去改變亞伯拉罕的生命.
我相信他不會作出這樣的行動.
他生命裡面發出了上帝的光.
摩西的光在哪裡呢.
我想摩西的光很多人看到的.
帶領以色列人離開埃及的十災.
我們今天看電影都要看那些.
來到那個地方.
到埃及進入英雄之地.
各樣的神蹟奇事.
這個可能都是摩西的光.
但是我想分享一個很多時候人會忽略的片段.
就是有一次當那群以色列人在埋怨上帝.
在搞事.
上帝很憤怒的時候.
摩西又上去山上跟上帝談話.
上帝一氣之下說什麼.
上帝說我真的很憤怒.
我不要理了.
我滅了那些以色列人.
我就是興起你摩西一族.
曾經有一段記載.
然後摩西的反應是什麼呢.
摩西說上帝.
上帝你是不是傻了.
你怎麼可以這樣做.
你滅了以色列人的話.
那些埃及人就笑我們.
你帶我們出來曠野就是為了滅我們.
那為什麼不是埃及滅我們呢.
摩西就力爭為那群以色列人祈求上帝不要發出他的異怒.
但是大家不要忘記.
那群以色列人一直頂著摩西.
頂心慘.
下山又要被他們說.
又要再假倫.
我經常在想.
如果那一刻摩西都沒有人看到.

$^{1121}$他自己上山而已.
如果那一刻摩西有一點點自私的心態.
也好.
那群人真的說不上去.
我跟兒子說容易一點.
我跟自己家族說容易一點.
只要有這麼一絲的自私經過他的腦海.
我想整個《出埃及記》的故事就要改寫.
就不再不是十二之派.
我不知道會繼續怎麼寫.
但是摩西沒有這樣做.
每次我讀到這裡.
我覺得摩西在這個地方真的很值得敬佩.
他也有黑暗.
但是很值得敬佩.
他不是看到他自己.
他看到的是上帝的計劃.
到最後弱失的光.
不是他成為了埃及的宰相.
一人之下萬人之上.
不是他很厲害的政績.
七年的旱災他都能夠處理.
不是這些東西.
弱失的光.
我想更多的是創世記最後的記載.
當弱失見到那些賣他哥哥的時候.
大家記不記得那些哥哥怕得要死.
你是弱失嗎?.
糟糕了.
但是弱失的反應是什麼?.
當然弱失也有一點小心思.
搞一點東西.
撿一些東西.
又抓一個弟弟.
但是到最後創世記50章的時候.
聖經來記載.
弱失寬恕那些哥哥.
他用安慰的說話來跟哥哥說一句很重要.
我覺得是整個創世記最重要的一句話.
你們的意念你們的行動是帶著惡意.

$^{1161}$但是神的旨意是好的.
弟兄姊妹如果不是一個有上帝光照的人.
是說不出這句話.
你經歷了人生的黑暗之後.
你跟隔壁的人說.
或者你跟那些得罪你的人說.
跟那些令你很不開心的人說.
你的行動.
這句話我能說出來.
你的行動是帶著惡意的.
這句話我們每個人都能說出來.
下一句我說不出來.
神的旨意是好的.
如果不是上帝的光光照在弱失的生命當中.
他見到上帝一步一步帶領他.
我覺得他說不出這句話.
但是今天我們見到這些生命.
他們不單止經歷了上帝的光.
更重要的是他們的生命都成為了我們的光.
成為了那個時代的光.
我知道在香港或者在這個城市裡面的處境.
都有很多不容易和黑暗.
我們澳門未必完全明白.
但我們經常看TVB.
我們經常看你們的電視.
所以有些地方我們會了解.
但是我想.
是的,黑暗的,真的有黑暗的.
真的有些地方我們都無力.
但是上帝提醒我們.
不單止要認定他.
認定上帝的說話繼續向前.
並且不要忘記在這些黑暗當中.
你和我要成為那個光.
帶給身邊的人.
見到上帝的榮耀.
見到上帝的恩典.
見到上帝的盼望.
無論你什麼年紀都好.
你老到80歲.

$^{1201}$你老到120歲都好.
上帝都可以繼續有美好的旨意在我們生命當中發生.
我帶著這一句話.
希望我們能夠激勵彼此.
能夠活在上帝的光中.
讓上帝繼續帶領著我們眾弟兄姊妹.
在這個時代當中.
把握機會能夠做到什麼.
我們就努力去做.
我們一起祈禱.
因主幫助我們.
我們是很軟弱.
我們是很無力的一群.
不過你的話語卻是光照我們.
能夠照亮我們的生命.
讓我們見到有盼望.
不單止讓我們見到有盼望.
而是你願意使用我們的生命.
成為上帝你自己光的一部份.
所以當今天我們一起來敬拜.
當我們口說出上帝的話語的時候.
我們的生命去回應那些黑暗的時候.
上帝讓我們有機會.
有份去成為光一樣.
去光照這個世界.
今天你不是要我們每一個弟兄姊妹.
好像聖經那麼偉大的人.
一樣要做到很大的事.
但可能你卻要邀請我們.
在生命很微小的.
在家庭,工作,城市中.
那些微小的位置裡.
願意為上帝你持守.
願意你行出上帝你的話語.
求願你帶領我們眾弟兄姊妹.
在這個時代當中.
繼續活在光中.
並且成為光.
給我們信心.
帶領我們.

$^{1241}$祈禱感恩.
奉主耶穌基督的名字祈求.
Amen.
\newpage



\section{創世記 16:1-16}
\label{sec:lCXbCQNDTgA}
\textbf{2024年9月1日崇拜直播｜吳真真牧師｜正視你的情緒－受屈者的出路｜創世記16\_1-16}
\newline
\newline
連結: \href{https://youtube.com/watch?v=lCXbCQNDTgA}{\texttt{https://youtube.com/watch?v=lCXbCQNDTgA}} ~~~~ 語音日期: 2024-09-02
\newline
\newline
\hyperref[sec:cerRyXhHqiE]{\small{< < < PREV SERMON < < <}}
~
\hyperref[sec:index_chronic]{\small{[返順時目]}}
~
\hyperref[sec:uj5FB161Gvc]{\small{> > > NEXT SERMON > > >}}
\newline
\newline
創世記 16:1-16
\newline
\begin{longtable}{cl}
\hline
\hline
章節 & 經文 (和合本修訂版)\\
\hline
16:1 & \begin{tabularx}{0.7\textwidth}{X} 亞伯蘭的妻子撒萊沒有為他生孩子。撒萊有一個婢女，是埃及人，名叫夏甲。 \end{tabularx} \\ \\ \relax
16:2 & \begin{tabularx}{0.7\textwidth}{X} 撒萊對亞伯蘭說：「看哪，耶和華使我不能生育。你來和我的婢女同房，也許我可以從她得孩子 。」亞伯蘭聽從了撒萊的話。 \end{tabularx} \\ \\ \relax
16:3 & \begin{tabularx}{0.7\textwidth}{X} 於是亞伯蘭的妻子撒萊把她的婢女，埃及人夏甲，給了丈夫為妾；那時亞伯蘭在迦南已經住了十年。 \end{tabularx} \\ \\ \relax
16:4 & \begin{tabularx}{0.7\textwidth}{X} 亞伯蘭與夏甲同房，夏甲就懷了孕。她看見自己有孕，就輕視她的女主人。 \end{tabularx} \\ \\ \relax
16:5 & \begin{tabularx}{0.7\textwidth}{X} 撒萊對亞伯蘭說：「我因你受了委屈。我把我的婢女放在你懷中，她見自己懷了孕，就輕視我。願耶和華在你我之間判斷。」 \end{tabularx} \\ \\ \relax
16:6 & \begin{tabularx}{0.7\textwidth}{X} 亞伯蘭對撒萊說：「看哪，婢女在你手裡，你可以照你看為好的對待她。」於是，撒萊虐待她，她就從撒萊面前逃走了。 \end{tabularx} \\ \\ \relax
16:7 & \begin{tabularx}{0.7\textwidth}{X} 耶和華的使者在曠野的水泉旁，在書珥路上的水泉旁遇見夏甲， \end{tabularx} \\ \\ \relax
16:8 & \begin{tabularx}{0.7\textwidth}{X} 對她說：「撒萊的婢女夏甲，你從哪裡來？要到哪裡去？」她說：「我從我的女主人撒萊面前逃出來。」 \end{tabularx} \\ \\ \relax
16:9 & \begin{tabularx}{0.7\textwidth}{X} 耶和華的使者對她說：「你要回到你的女主人那裡，屈服在她手下。」 \end{tabularx} \\ \\ \relax
16:10 & \begin{tabularx}{0.7\textwidth}{X} 耶和華的使者對她說：「我必使你的後裔極其繁多，多到不可勝數。」 \end{tabularx} \\ \\ \relax
16:11 & \begin{tabularx}{0.7\textwidth}{X} 耶和華的使者又對她說：「看哪，你已懷孕，要生一個兒子。你要給他起名叫以實瑪利，因為耶和華聽見了你的苦楚。 \end{tabularx} \\ \\ \relax
16:12 & \begin{tabularx}{0.7\textwidth}{X} 他為人必像野驢。他的手要攻打人，人的手也要攻打他；他必常與他的眾弟兄作對。」 \end{tabularx} \\ \\ \relax
16:13 & \begin{tabularx}{0.7\textwidth}{X} 夏甲就稱那向她說話的耶和華為「你是看見的神」，因為她說：「他看見了我之後，我還能在這裡看見他嗎？」 \end{tabularx} \\ \\ \relax
16:14 & \begin{tabularx}{0.7\textwidth}{X} 所以這井名叫庇耳‧拉海‧萊，看哪，它位於加低斯和巴列的中間。 \end{tabularx} \\ \\ \relax
16:15 & \begin{tabularx}{0.7\textwidth}{X} 後來夏甲為亞伯蘭生了一個兒子；亞伯蘭給夏甲生的兒子起名叫以實瑪利。 \end{tabularx} \\ \\ \relax
16:16 & \begin{tabularx}{0.7\textwidth}{X} 夏甲為亞伯蘭生以實瑪利的時候，亞伯蘭年八十六歲。 \end{tabularx} \\ \\
[1ex]
\hline
\hline
\end{longtable}
$^{1}$主持人:今日的經訓是在《創世記》第十六章.
在教會的聖經舊約第十六頁.
請聽我讀出.
亞伯蘭的妻子撒萊不給他生兒女.
撒萊有一個侍女名叫夏甲,是埃及人.
撒萊對亞伯蘭說:耶和華使我不能生育.
求你和我的侍女同房,或者我可以因她得孩子.
亞伯蘭聽從了撒萊的話.
於是亞伯蘭的妻子撒萊便將侍女埃及人夏甲給了丈夫為妾.
那時亞伯蘭在迦南已經住了十年.
亞伯蘭與夏甲同房,夏甲就懷了孕.
她見自己有孕,就小看她的主母.
撒萊對亞伯蘭說:我因你受屈.
我將我的侍女放在你懷中.
她見自己有了孕,就小看我.
願耶和華在你我中間判斷.
亞伯蘭對撒萊說:侍女在你手下.
你可以隨意待她.
撒萊輔待她,她就從撒萊面前逃走了.
耶和華的使者在曠野樞爾路上的水泉旁遇見她.
對她說:撒萊的侍女夏甲,你從那裡來?要往那裡去?.
夏甲說:我從我的主母撒萊面前逃出來.
耶和華的使者對她說:你回到你主母那裡,伏在她手下.
又說:我必使你的後裔極其繁多,甚至不可生數.
並說:你如今懷孕,要生一個兒子.
可以給他起名叫爾實瑪利.
因為耶和華聽見了你的苦情.
她為人必像野狼,她的手要攻打人,人的手也要攻打她.
她必住在眾弟兄的東邊.
夏甲就稱對她說話的耶和華為看顧人的神.
因而說:在這裡我也看見了看顧我的魔.
所以這名叫比爾拉海內者者正在迦底斯和巴列中間.
後來夏甲給亞伯蘭生了一個兒子.
亞伯蘭給他起名叫爾實瑪利.
夏甲給亞伯蘭生爾實瑪利的時候,亞伯蘭年八十六歲.
(自動翻譯).
各位弟兄姊妹平安.
今天我們繼續《正是你的情緒》系列.
我們繼續用舊約的故事去認識神賜給我們的禮物.
就是我們的情緒.

$^{41}$因為有著這份禮物和情緒.
讓我們看到自己的需要.
可以與人結連.
今天我們要看的情緒是受屈.
那種委屈,受屈辱的感覺.
其實在日常生活中,我們很難避免這種感覺.
當我們覺得自己有時不得到公平的對待.
或者我們不被人尊重.
或者我們被忽視,被輕視.
甚至在一些生活情況中.
我們會被欽著,被侮辱.
以致我們都會有些委屈的感覺.
這種感覺不是在某些地方才出現.
反而有時在家庭,工作,朋友之間,教會中.
我們幾乎都會遇到這些問題.
這種感受經常會出現.
不要小看這種感覺.
因為這種感覺對自己,對人的破壞都很大.
今天我們會看一家庭.
這家庭都面對一些他們覺得受屈的事件.
我們一起去看看.
當我們感覺到受屈時的情緒會是怎樣.
當我們面對這種情緒時.
原來我們會有一些應對機制.
那些又是甚麼機制呢?.
我們又看看神怎樣讓這種情緒讓我們成長.
讓我們在生命中突破.
讓我們一起去祈禱.
透過你讓我們與人相處.
透過我們的情緒.
透過你自己聖經裡寶貴的應許.
你信實的話.
你讓我們在每一個情緒裡找到出路.
我們可以認識自己.
我們可以與你遇見.
上帝一救,就求你在這一刻.
我們今天認真去聽你的道.
我們希望聖靈引導我們明白你的心意.
作為一個祭獻給你.
求你去閱立,求你去引導,求你去輔助.

$^{81}$安慰我們.
因為有很多時候我們都會有時.
會在這些感受當中.
願意你自己親自成為那個出路.
我們這樣禱告交託.
奉主耶穌基督的聖名自祈求.
阿們.
剛才還讀的經文是《創世記》第十六章.
大家單單看這段經文.
就會發覺這段經文是在說一個問題.
是一個什麼問題呢?.
簡單來說就是一個家庭裡.
沒有兒子.
生不了後代這個問題.
你覺得可能在香港這個社會.
現在這個世代,這不是什麼大事.
不過其實在神救贖的歷史裡.
這段經文是佔了一個很重要的位置.
自從《創世記》第十二章第一節開始.
神將亞伯蘭從他的本族呼召出來.
他要離開他的家鄉.
離開他的文化.
甚至離開當時影響他多元的宗教.
而是去學習,單單專一去跟從耶和華.
神應許他要成為大國.
說他的後裔要像海裡的沙.
天上的星星那麼多.
不過十年之後.
亞伯蘭的家仍然沒有所出.
撒萊的年紀越來越大.
已經七十五歲了.
他可以生育的機會也越來越低.
他們的信心再一次受考驗.
不過這個考驗最終帶來的就是.
家無寧日的局面.
在這段落裡有兩個很重要的人物.
兩位女主角,分別是撒萊和夏婕.
他們都有一個共通點.
就是他們蒙受了屈辱.
讓我們一起去看看他們兩個如何蒙受屈辱.

$^{121}$我們讀第一至第四節.
「亞伯蘭的妻子撒萊不欲她生兒女.
撒萊有一個侍女名叫夏婕.
是埃及人,撒萊對亞伯蘭說.
耶和華使我不能生育.
求你和我的侍女同房.
或者我可以因她得孩子.
亞伯蘭聽從了撒萊的話.
於是亞伯蘭的妻子撒萊將侍女埃及人夏婕.
給了丈夫為妾.
那時亞伯蘭在迦南已經住了十年.
亞伯蘭與夏婕同房.
夏婕就懷了孕.
她見自己就小看她的主母」.
我們很清楚知道.
連撒萊都很清楚知道.
能不能生育不是人能控制的.
她和丈夫在第二節說.
「耶和華使我不能生育」.
「生育」這個字直譯是.
耶和華關閉了.
耶和華上鎖了.
對生產這件事上鎖了.
以致她不能生育.
所以她想說罪不在她.
問題不在她身上.
就算怎樣努力.
都不可以靠她自己的能力做到.
因為我們很清楚知道.
能否生育是上帝的掌權.
不過在當時的社會裡.
婦女不能為家族傳宗接代.
是一個極大的羞恥.
是一個被人看低.
看為沒有尊嚴的一件事.
所以的確撒萊的第一個屈辱.
就是社會的眼光.
她確實承受了不能生育的悲痛.
這是她第一個受屈.
而第二個受屈.

$^{161}$就是因著她想解決第一個受屈而帶來的.
既然她不能為她的家族繼後香燈.
這個污名令她被人藐視.
被人輕看.
她就嘗試去找一個出路.
她就向她的丈夫獻計.
她說「求你和我的侍女同房.
或許我可以因她得孩子」.
「得孩子」這個字在原文來說.
是被建立.
意思是撒萊期望.
兩夫婦可以因著有她的兒子.
就可以傳宗接代.
這個家就變得完整和完全.
所以她很想擺脫這個屈辱.
很想用自己的方法去解決這件事.
於是她就向丈夫獻計.
第一節告訴我們.
她就找她的侍女下鴿.
來服侍這位主人的侍女下鴿.
去和她的丈夫同房.
根據當時的法例和習俗.
撒萊是下鴿的主人.
她把自己的侍女當作是亞伯蘭做妾侍.
即是下鴿可以代表撒萊去生子女.
而下鴿所出的子女.
就會歸在撒萊的名下.
就成為撒萊的兒女.
這個家庭就因此而完整.
她的丈夫亞伯蘭很信服.
他沒有反對.
他信服妻子的建議.
這次的亞伯蘭很特別.
他沒有像第十五節一樣.
當上帝說他有後代.
去接續繼承他的家業的時候.
他就問上帝.
我現在這麼久都沒有兒子.
究竟會不會是我的僕人.
我的管家已變已借.

$^{201}$是我的後裔呢?.
他有問上帝.
但這次他沒有.
而且這次他沒有考慮到.
上帝從他的本家本族.
那些文化社會的習俗.
去呼召他出來.
去跟從上帝設立一夫一妻的制度.
去採納妻子的建議.
去享齊人之福.
這個建議.
令他們的家帶來了很多麻煩.
結果夏甲真是懷孕.
不過這個懷孕.
令夏甲心裡生起了一些驕傲.
他抱著金菠蘿.
他覺得自己可以母憑子貴.
身心驕傲.
然後他就不尊重他的主人薩萊.
本來薩萊是想.
因著夏甲為他生兒子.
他的生命被建立.
從卑微,從被藐視去建立.
但反而因著這個緣故.
他再被輕視,被詛咒.
他本來在家裡.
他是家主的夫人地位.
都被剝奪.
所以薩萊其實是帶著兩個屈辱.
另外我們又看看夏甲的屈辱.
我們留意第六節.
第六節說.
亞伯蘭對薩萊說.
「使女在你手下,你可以隨意待她」.
薩萊苦待她.
她就從薩萊面前逃走了.
其實除了薩萊之外.
這個故事的另一個女主角夏甲.
她同樣受到屈辱.
夏甲本來的身份是薩萊的侍女.

$^{241}$極有可能她在第十二章.
亞伯蘭因為饑荒而逃到埃及的時候.
她由那時候開始成為薩萊的僕婢.
作為僕婢是沒有尊嚴.
也沒有選擇.
她沒有為自己的婚姻作主的話事權.
她出身汗味.
當薩萊命令她要和亞伯蘭同房的時候.
她是沒有選擇.
不過很特別.
這也是人的需要.
當人感到自己汗味委屈的時候.
心裡往往有一個傾向上.
向上提升.
離開底層的慾望.
夏甲也因為被揀選代主人生子.
而覺得自己的地位被提升了.
她於是就作出了一些蔑視她主人的行為.
不過她和薩萊一樣.
當她以為她可以無憑只貴脫離屈辱的時候.
反而最終她的下場也是蒙受屈辱.
為什麼呢?.
剛才那段經文告訴我們.
薩萊在亞伯蘭的准許下.
可以按照自己的心意對付這個令她受屈的人.
所以結果薩萊選擇了輔代夏甲.
去壓迫她,以至對她造成身心靈的虐待.
結果夏甲就受盡屈辱.
而另一方面,當她受盡屈辱的時候.
她不和薩萊一樣.
薩萊可以向她的丈夫申訴.
但她委屈的是與她同房的主人亞伯蘭.
對薩萊這些舉動視若無睹.
她的痛苦沒有被看見.
也令她受屈辱.
所以我們看到兩個受屈的人.
在這個情況下.
他們的心靈有很多痛苦.
我們看看當我們有這些受屈的痛苦的時候.
會有什麼機制和反應呢?.

$^{281}$讓我們一起讀第五,第六節.
薩萊對亞伯蘭說.
「我因你受屈,我將我的屎女放在你懷中」.
「她見自己有了人,就小看我」.
「怨耶和華在你我中間判斷」.
亞伯蘭對薩萊說.
「屎女在你手下,你可以隨意待她」.
「薩萊輔待她,她就從薩萊面前逃走了」.
當薩萊承受了夏革輕視她的屈辱.
她就找她的丈夫亞伯蘭尋求協助.
她希望她的丈夫主持公道.
而且她還說她希望上帝出手為她申冤.
不過,有沒有留意到薩萊將這個委屈的原因.
這個受屈的原因.
完全加諸在丈夫的身上.
第一句就說「我因你受屈」.
意思是,最壞的是你.
如果不是你的話,我就不用這麼委屈.
但大家有沒有想過.
當初是誰叫亞伯蘭和夏革同房的?.
有些經學者細心去研讀這段經文時.
他們會發覺這段經文的描述.
和創世記聽夏娃的話.
是分辨事惡樹的果子描述.
在結構上,在詞組的運用上極之相似.
聖經說「夏娃拿果子給丈夫,丈夫聽從了妻子」.
而原文「薩萊」亦是「拿」同一個字.
「拿夏革給丈夫,丈夫聽從了妻子」.
兩個的情況有一個共通點.
就是妻子的建議和方法.
都讓夫婦二人一起離開了神的心意和計劃.
而且「受屈」這個字.
原文是指極度強暴和殘忍對待的感受.
「薩萊」有強烈被迫害的感覺.
同樣,這個字亦有意思.
就是她得到不公平的對待.
所以她想找耶和華主持公道.
要讓丈夫知道她受的傷害完全是因為她.
而沒有提及自己的責任.
但其實當初她希望為丈夫立後而被建立.

$^{321}$這個心意她沒有提出.
她沒有看到.
所以她訴諸她的委屈.
覺得自己無辜.
我會形容這是受害者的心態.
為自己受害.
而忽略了自己有責任的部分.
原來委屈也會有陷阱.
就是我們會將傷害完全掩蓋自己.
去看自己的責任.
我們會躲藏在委屈的背後.
外面是生氣,憤怒,嫉妒.
而忘記了委屈背後也有自己的責任.
而亞伯蘭在撒萊的控訴下.
他最終決定逃避一家之主的責任.
他說:侍女就在你手下.
你可以隨便待她.
他將夏革滅時撒萊所引起的衝突處理.
拱手相讓給撒萊.
於是就讓受辱的陷阱繼續沉下去.
受屈時我們會有這個反應.
我們會漠視,覺得自己很委屈.
陷入陷阱裡.
甚至會作出一些攻擊性的報復.
撒萊得到這個令牌後.
他就苦待夏革.
目的有兩個.
第一是希望夏革可以分清莊瀚.
警告她不要再做出輕視我的行為.
做出侮辱我的行為.
更加可能是他想脫去這個眼中釘.
將她趕走.
令自己不用再受被蔑視的痛苦.
的確當人受到委屈.
得到不公平的對待的時候.
有時我們就認為自己有權制裁令我們受害的人.
這也是我們要小心的一個惡性循環不斷延伸的結果.
我們看到第六節.
夏革最後如何回應呢?.
夏革最後從撒萊面前逃走.

$^{361}$夏革可能因為自己本來成為亞伯蘭的妾侍而提升.
感到驕傲.
對她的主人不尊重.
結果換來被撒萊欺負.
比她之前做侍女的時候的地位更差.
結果她就決定逃離這個令她痛苦的地方.
聖經說她逃走.
同樣夏革亦可能因為被虐待的痛苦.
她也看不到自己的問題.
是她先輕視她的主人.
不過同樣在這個時候.
很混亂的她也有一個報復的心.
既然你利用我去建立自己.
你也得不到你想要的.
我就逃走,跟你一拍兩散.
其實我們真的會看到.
當人透過情緒去面對神.
去與神相遇的時候.
其實是可以帶來靈命的成長和後果.
如果我們躲在情緒的後面的時候.
用受傷去做擋箭牌.
妨礙了真正處理問題的動力的時候.
我們往往會越走越深.
在情緒裡更加不能自拔.
我們發覺我們用人的方法去實現神的心意.
結果就帶來一男子的問題.
自己也帶來受苦和受屈的感覺.
用人的方法沒有正視自己的傷痛之處.
卻帶來了憤怒,嫉妒和仇怨.
帶來一個不可收拾的殘局.
我們就看看下面跟著七至十六節.
上帝的介入如何將這個屈辱轉化成為祝福.
讓我們一起讀第七至第八節.
下甲這裡說到兩個婦女分別在不同的屈辱下.
下甲因為受盡痛苦,因為撒萊苦待她.
鏡頭一轉已經去到第七節.
已經去到曠野裡面的一個水泉旁邊.
你會看到在這個時候上帝的使者.
耶和華的使者出現.
耶和華的使者出現跟她相遇.

$^{401}$雖然這個叫相遇.
但大家不要被這個相遇以為是偶遇.
因為「圓文一」這個字是指尋找她.
上帝定義要尋找下甲.
要幫助這個深受屈辱,痛苦的婦人.
但你會發覺很有趣.
上帝尋找她的地方是樞爾路上的一個水泉旁.
樞爾路上的水泉旁.
就是在迦南,即是她原先的地方.
通往埃及的必經之路.
下甲在那個時候定義要返回.
我想她很想返回自己的家鄉.
我們估計下甲在迦南無依無靠.
她是一個婢女.
於是乎樂方兒徒的情況下.
如果說到這個地方.
我們估計她大概已經走了一百公里的路程.
一百公里的路程是一個孕婦一兩天未必做得到的事情.
但如果上帝是想下甲.
即是返回她的主母那裡.
我們剛才讀經文知道.
她要返回她的主母那裡.
為何到這個時候她才出手攔她呢?.
我相信因為這個時候是最佳的時機.
最好的時機.
因為很多時候當我們在混亂的情緒下.
不能理性思考的時候.
實在需要一些冷靜的時間.
才能認清環境.
聽到上帝的心意和計劃.
上帝的介入正是要下甲.
正視自己在委屈之下.
她內心的傷痛之處.
讓我們看看他們當中是怎樣對話的.
首先,她稱呼下甲為「殺來的侍女下甲」.
也不是直接叫她下甲.
也不是叫她「亞伯蘭的妻」.
「亞伯蘭的妾」.
她是要告訴她.
「你叫做殺來的侍女下甲」.

$^{441}$這個稱呼對下甲來說一定很驚訝.
因為為何這個人知道我的身份呢?.
她一定不是普通人.
她會是神差來的侍者.
所以接下來這個對話的展開.
令下甲有更多的心去聆聽.
她覺得這個人將要對我說的話有份量.
然後她問了她甚麼問題?.
她問她的問題是.
「你從哪裡來,要往哪裡去?」.
我想侍者既然知道她是誰.
她當然知道下甲要從迦南逃出來.
也要去埃及.
但為何她會這樣問呢?.
她其實是想問下甲為何會弄成這樣?.
你將來會有甚麼打算?.
我們看到下甲只回答一半.
她回答了上半部分.
就是我在我的主母殺來面前逃跑出來.
她回答了為何會弄成這樣.
很多時候,當我們覺得自己受傷害時.
我們就會停下來.
我們會注目過去.
而未必有勇氣或想到如何面對將來.
正正因為這樣.
上帝的侍者介入.
就是要幫她思考.
事情既然發生了.
你面對這件事,究竟要怎樣解決呢?.
她提醒她.
你要認清你作為殺來侍女的身份.
你應該做你俠嶼的事.
他們的相遇令下甲聽到她對她的稱呼.
喚起她自己原本的身份.
原來我就是殺來的婢女.
讓她反省到.
她之所以要逃跑出來.
是因為她僭越了主人的地位.
所帶來的惡果.
然後耶和華給她一個出路.

$^{481}$她說:你回到你的主母那裡.
伏在她的手下.
意思是你按著你這個身份.
按著殺來侍女這個身份.
去順服神的旨意去做就夠了.
其實我覺得很殘忍.
這個方法很難想像.
耶和華竟然吩咐她.
做一個如此不近人情的吩咐.
好像將她再次送養入虎口一樣.
不過你要留意一件事.
原來第六節的「虎代」和第九節的「伏」字.
這兩個字的字根.
都是同一個希伯來文的字根.
裡面都是帶著痛苦的.
神要下格面對受屈辱的痛苦的方法.
不是從痛苦逃走.
而是要帶著這個痛苦回去.
要自己謙卑自己.
自己將自己降卑.
繼續伏在殺來的權柄下.
不過我們從下面的經文.
我們就會知道.
上次的「虎」是被動的.
它是被人虎代.
但今次的「虎」是上帝吩咐她.
要主動回去伏在她的權柄下.
是有不同的.
我們會看回第十到第十二節.
耶和華繼續去應許她.
「你的後裔極其繁多,甚至不可生數」.
這是一個神對她回去伏下.
在殺來權柄下的一個應許.
祂要她生一個兒子.
這個兒子叫做「以實瑪利」.
因為耶和華聽見你的苦情.
當她活在困苦裡.
她一定會想自己和胎兒.
如何可以保得住.
我的命和兒子的將來會是怎樣.

$^{521}$她覺得沒有保證.
但神應許她.
「你會生孩子.
而且你還要改他的名字做以實瑪利」.
而以實瑪利這個名字.
是預言她的生命不再受核制.
不像她媽媽一樣要遺老.
使者告訴她.
這個應許是因著神聽了她的苦情.
的一個見證.
所以這件事就是神主動尋求夏革.
夏革向神發出了她的愛情.
她的傷痛,她的絕望.
神是聽見的.
當她回到伏在撒萊手下的時候.
當她每一次去叫她兒子的名字的時候.
她就會記得.
神聽她的苦情的經歷.
上帝和上次不同.
祂不會像亞伯蘭那樣.
對她的受苦視而不見.
袖手旁觀.
而是上帝要賜福給她.
要讓她和亞伯蘭,撒萊一樣.
她的後裔要多不勝數.
當夏革聽到神對她的應許的時候.
她就對神說.
神是看顧人的神.
她稱呼神是看顧我苦況的上帝.
夏革因為這個痛苦和神相遇.
最後夏革為她和神的使者相遇的井.
改了一個名字.
這個名字就是一個譯音.
就是「比爾拉海來」.
這個井或水泉的名字.
譯出來就是活著.
而且看見我的那位神的井.
她在這裡和上帝相遇.
認識了上帝.
知道祂是信實,慈愛,滿有恩典.

$^{561}$而且樂於隨時幫助她的上帝.
我們看見神給她應許.
給她很大的祝福.
她有很多後代.
結語第十五,第十六節告訴她.
她果然生了二十萬.
有數之不盡的後代.
不過我相信最大的祝福.
是她遇上看顧她的上帝.
其實我再說.
受屈受辱可以有很多情緒表現.
我們可以被憤怒遮蓋.
我們可以被疾度遮蓋.
我們會內心沉鬱,悲傷.
甚至會帶來仇恨.
但其實受屈有一個很核心的情緒.
就是我們那種受傷,絕望,無助的情緒.
但當夏革與神相遇.
上帝就將她的受屈受辱轉化.
她向上帝展示她生命的痛處.
她讓上帝的愛湧流在生命的痛.
讓她感受到這些痛不是她生命的全部.
她有被神聽見.
她有被神看顧.
她反而被神的慈愛包圍.
她不再尋求報復.
她不再被傷痛糾纏她的生命.
反而她勇敢面對這些傷痛.
但有主的看顧她不再怕.
有主的醫治她不再受傷.
弟兄姊妹,你相信上帝會聽你的苦情嗎?.
你相信上帝會看顧你,照顧你嗎?.
在我生命裡有一位恩師.
他在我的屬靈生命裡有很大的重要的教導和榜樣.
他本身是一位中學生物科老師.
在教學生涯裡教授過不同的學校.
他告訴我每一次神差派他去每間學校的時候.
他很清楚神給他的獨特使命.
不過他在學校裡都會有很多階級.
都會有很多不同的困難.

$^{601}$他亦都嘗試過經歷很多很受委屈.
很令他痛苦的地方.
他覺得有時對待,他很用心去教學.
很用心去付出,但都得不到公平的對待.
我有問他為何他不離開.
他說我多年來的生涯裡.
我沒有一次因為我覺得受委屈,痛苦而離開.
原因是因為神沒有叫我走.
反而他堅持在這些痛苦裡.
他更多給上帝的愛,保護,恩典.
去帶領他克服一個又一個的難關.
讓他最終可以完成使命.
他就會帶著這些新的使命去到新的工場.
其實在我們的環境裡.
到處我們都可能會經歷到這些.
不被尊重,不公平對待,受委屈的感覺.
但我們真的要學習撒萊一樣.
學習下甲一樣.
當他面對神的時候.
神就告訴他你要往那裡去.
那個勇氣是神才能給予我們.
詩篇103篇說.
「耶和華施行公義,為一切受委屈的人伸冤.
上帝為我們伸冤的方法很獨特.
但很肯定的是.
他不會點一條黑路給我們走.
他的愛是很豐富地包圍我們」.
我想起耶穌基督.
他自己絕對是受屈者的代表.
他在這個世上所受的屈辱.
可以告訴我們更多.
但他是用無條件的愛去回應.
我們每一個人的生命裡都有耶穌.
我深信在他的裡面.
我們會讓他的慈愛包圍.
我們會得意志.
我們的生命會有盼望.
讓我們一起祈禱.
撒萊天父,我們滿心感恩.
因為在你的愛裡.

$^{641}$我們不再懼怕.
我們不需要為著我們所受的傷害.
而苦苦地沉屈和痛苦.
我們可以經歷痛.
但我們可以因為你.
去捍衛這個不只是苦.
而是當中有甘甜.
有上帝的同在.
願以神讓我們在你的裡面.
在你的愛裡面找到出路.
願你跟我們每一個人說話.
提醒我們.
你與我們同在.
多謝你聽我們祈禱.
禱告交託.
奉主耶穌基督得勝.
命之祈求.
阿們.
\newpage



\section{腓立比書 4:1-9}
\label{sec:uj5FB161Gvc}
\textbf{2024年2月11日新春崇拜直播｜梁家麟牧師｜一個叮嚀三個祝願｜腓立比書4\_1-9}
\newline
\newline
連結: \href{https://youtube.com/watch?v=uj5FB161Gvc}{\texttt{https://youtube.com/watch?v=uj5FB161Gvc}} ~~~~ 語音日期: 2024-09-16
\newline
\newline
\hyperref[sec:lCXbCQNDTgA]{\small{< < < PREV SERMON < < <}}
~
\hyperref[sec:index_chronic]{\small{[返順時目]}}
~
\hyperref[sec:QTbAXdDuZOw]{\small{> > > NEXT SERMON > > >}}
\newline
\newline
腓立比書 4:1-9
\newline
\begin{longtable}{cl}
\hline
\hline
章節 & 經文 (和合本修訂版)\\
\hline
4:1 & \begin{tabularx}{0.7\textwidth}{X} 我所親愛、所想念的弟兄們，你們就是我的喜樂，我的冠冕。我親愛的，你們應當靠主站立得穩。 \end{tabularx} \\ \\ \relax
4:2 & \begin{tabularx}{0.7\textwidth}{X} 我勸友阿蝶和循都基要在主裡同心。 \end{tabularx} \\ \\ \relax
4:3 & \begin{tabularx}{0.7\textwidth}{X} 我也求你這真實同負一軛的，要幫助這兩個女人，因為她們在福音上曾與我、革利免和我其餘的同工一同勞苦，他們的名字都在生命冊上。 \end{tabularx} \\ \\ \relax
4:4 & \begin{tabularx}{0.7\textwidth}{X} 你們要靠主常常喜樂。我再說，你們要喜樂。 \end{tabularx} \\ \\ \relax
4:5 & \begin{tabularx}{0.7\textwidth}{X} 要讓眾人知道你們謙讓的心。主已經近了。 \end{tabularx} \\ \\ \relax
4:6 & \begin{tabularx}{0.7\textwidth}{X} 應當一無掛慮，只要凡事藉著禱告、祈求和感謝，將你們所要的告訴神。 \end{tabularx} \\ \\ \relax
4:7 & \begin{tabularx}{0.7\textwidth}{X} 神所賜那超越人所能了解的平安，必在基督耶穌裡，保守你們的心懷意念。 \end{tabularx} \\ \\ \relax
4:8 & \begin{tabularx}{0.7\textwidth}{X} 末了，弟兄們，凡是真實的、凡是可敬的、凡是公義的、凡是清潔的、凡是可愛的、凡是有美名的，若有甚麼德行，若有甚麼稱讚，你們都要留意。 \end{tabularx} \\ \\ \relax
4:9 & \begin{tabularx}{0.7\textwidth}{X} 你們從我所學習的，所領受的，所聽見的，所看見的事，你們都要繼續去做，賜平安的神就必與你們同在。 \end{tabularx} \\ \\
[1ex]
\hline
\hline
\end{longtable}
$^{1}$今日經文是肥立比書第四章一到九節.
經文都印在你們的秩序表裡面.
請聽我讀.
四章一至九節.
我所親愛所想念的弟兄們.
你們就是我的喜樂.
我的冠冕.
我親愛的弟兄.
你們應當靠主站立得穩.
我勸友阿諜和秦都基.
要在主裡同心.
我也求你這真實同父一額的.
幫助這兩個女人.
因為她們在福音上.
曾與我一同勞苦.
還有甲理冕.
並其餘和我一同作供的.
她們的名字都在生命冊上.
你們要靠主常常喜樂.
我再說你們要喜樂.
當叫眾人知道你們謙讓的心.
主已經近了.
應當一無掛慮.
只要凡事藉著禱告.
將你們所要的告訴神.
神所賜出人意外的平安.
必在基督耶穌裡.
保守你們的心懷意念.
弟兄們.
我還有未盡的話.
凡是真實的.
可敬的.
公義的.
清潔的.
可愛的.
有美名的.
若有什麼德行.
若有什麼稱讚.
這些事你們都要思念.
你們在我身上所學習的.

$^{41}$所領受的.
所聽見的.
所看見的.
這些事你們都要去行.
賜平安的神.
就必與你們同在.
今天年初二我們特別高興的.
有我們熟悉的梁家倫牧師.
為我們宣講的.
題目是.
「一個叮嚀 三個祝願」.
麻煩梁牧師.
各位頂智妹新年快樂.
身體健康.
家庭幸福.
剛才說的祝願.
有些是祝願.
身體健康不完全受我們控制.
不過其他的.
不僅僅是一個祝願.
是一個命令.
是一個信仰的命令.
這是我們的任務.
我們要讓我們自己.
和身邊的人.
覺得幸福.
我們要讓我們自己.
和身邊的人常常喜樂.
從年三十晚到年初一.
大家應該都吃得差不多了.
新春講道.
一定不可以heavy dose.
不可以太重劑量.
要清淡一點.
所以只能夠說一些小品的東西.
太難消化.
或者太沉重的.
不是一個很適合時機的東西.
所以我和大家選了.
在《肥納比書》最後一章.

$^{81}$一段聖經裡面.
有些散散發發.
不過同樣是很重要的教導.
想從那裡.
幫我們一起去思考.
很簡單地說.
在這章聖經裡面.
保羅在做一些人事的安排.
特別是對某些人做一些提醒.
亦都說應該怎樣處理.
這些人事的風波.
《肥納比書》寫作的背景.
我們今天不詳細說.
我們知道的就是.
保羅正在坐牢.
他在所謂的house arrest.
即是在家中軟禁.
很多人可能說一點題外話.
問我為甚麼.
經常說軟禁這些東西呢.
如果你去過意大利 去過羅馬.
你會有一個名勝.
就是彼得和保羅的監獄.
很有趣的.
你可以去參觀.
彼得和保羅的監獄.
其實就是一個地下室.
一個小區.
一個chamber.
最初我們去的時候.
就覺得很奇怪.
監獄之為監獄.
羅馬是羅馬帝國的首都.
為甚麼監獄只有一個房間.
少成這樣.
酒店最少也有十多個房間.
這麼小 為甚麼是監獄呢.
我要告訴大家.
在古代.
就不同今天.

$^{121}$第一古代資源不夠.
第二古代人比我們現在聰明.
政府不會建一個大監獄養著人.
建監獄是很貴的.
又要給他吃又要給他住.
又要找人看守住.
其實是很貴很貴的事情.
以前的政府沒有這麼蠢.
所以基本上不會給人坐無期徒刑.
你拿來搞.
虧了錢.
砍了個頭算了.
總之就是即捕即審即殺.
這是最簡單的做法.
真的如果要坐牢.
就是你在家裡坐.
就是禁止你出門.
有時候會有兵兵在門口看著.
那些重犯.
如果不是重犯.
就間中來家訪一下.
總之就是在自己家裡困著.
就不會給你一個監獄.
哇 建一個很貴的監獄.
哪有這麼多錢.
保羅在那個時候是坐牢.
他寫的監獄書信是在家裡寫的.
就是House of Rats.
在家裡被軟禁.
他在一個完善不理想的狀態.
然後他寫信給菲律賓的教會.
這間教會我們在看第四章.
你會看到裡面有一些人事的問題.
人事的風波.
簡單的說法就是.
在一個不算很理想的環境裡.
我最近在準備一個系列的講道.
是幫年底有一間教會預備的.
Summer Camp.
他們不是叫Summer Camp.

$^{161}$叫讀經營.
我給的題目是一個不完美的信仰實踐.
刺激我有這個題目的想法是.
在今年一月.
我幫的意大利環球學院有一個神學營.
請了一些學者來討論AI時代的信仰實踐.
在那個營裡有談到有關AI時代.
AI世界的問題.
和教會相關的就是.
很多人會選擇去AI的教會.
或者去找一個AI的牧者.
現在牧者這麼短缺.
找AI的牧者就最好了.
你可以想像一個AI的牧者可以有多好.
他可以和我相處幾次之後.
就已經摸了我的思路.
他已經可以知道我有什麼想法.
我最近的情緒.
我最近的感受.
他就可以量身訂做.
為我去預備一個講章.
你說不是很新.
哪有新的.
全球擺上網的講章.
可以有幾千萬甚至上億篇.
我沒有誇張.
香港一個星期產出2000篇.
一年也有10萬篇.
10年就有100萬篇.
很誇張的.
所以有幾千萬篇.
我天下文章一大抄.
左抄右抄.
左剪右剪.
一定適合用.
哪有這麼多新意.
不用這麼多新意.
所以AI時代有個AI的牧者.
你說多好.
他完全符合我的心意.

$^{201}$想我想想的事.
說我想要知道的事.
這是不是一個很正的安排.
於是很多台下的討論就是.
將來會不會有一間教會.
找一個AI的牧者.
那個AI的牧者會搶走所有的客人.
所有的影子妹.
包括我們王浸.
都全部一起拉隊過主.
去了那間教會.
我們談出來的結論是不會的.
因為如果真的有這樣的AI時代.
有這樣的AI的牧者.
那就不如一人有一個教會.
你明不明.
我怎會要走在一起.
就算我和我老婆都不是完全.
心心相印到任何事都一樣的看法.
我弄一個AI的牧者.
他弄一個AI的牧者.
對不對.
不止AI的牧者.
我作為一個牧者.
我不是AI的牧者.
不如我弄一間AI教會.
你明不明.
我度身訂做一班AI的執事.
完全合我的.
一班AI的回眾.
我講什麼他都笑.
會傻笑的.
你明不明.
我一個AI的牧者.
我一個正常的牧者.
我找一班AI的回友.
我就有一間我的教會.
完全good perfect.
合我的口味.
你明不明我說什麼.

$^{241}$一人一個教會.
不需要走在一起.
然後我們又說.
會不會有AI伴侶的出現呢.
度身訂做一個老婆.
一個老公.
完全合我的.
完全和我的想法一樣.
你說多完美呢.
多完美呢.
我要說這樣的情況是會發生的.
非常會發生的.
正如今天很多人不想結婚.
或者結婚不想生孩子.
很多人聲稱寧願養狗都不養子女.
都是一個這樣的選擇.
養子女和養狗最大的分別在哪裡.
養狗比較under control一點.
比較不會誤譯我.
所以沒那麼勞氣.
我去旅行.
雖然有時候都會想念.
不過我放牠在寵物酒店.
都是容易安排的.
你敢不敢放你兒子去寵物酒店呢.
你不行的.
你發覺養寵物.
我可以控制得好一點.
我可以一切在我的掌握.
在我的感受.
在我的意志的指導之下.
將來我們可以有一個AI的伴侶.
完全按我的想法來做.
按我的想法來回應.
你說多完美呢.
我可以配合牠不同的樣子.
還可以換screen換樣子.
你說多好呢.
不過我在想.
假如我們將來真的可以有一個AI的伴侶.

$^{281}$那個算不算是愛情呢.
我這麼保守這麼傳統的人.
我就說不是.
最多是愛自己.
是self love.
我只愛我自己.
我不是真愛一個對象.
如果我可以度身訂做一個.
完全合我的口味的人.
他的意志就是我的意志.
這樣的人不是愛.
讓我很老套地說.
我們去背誦《臨前十三章》.
如果不用恆久忍耐又有恩賜.
不用凡事包容.
凡事盼望.
凡事相信.
凡事忍耐.
那個不是愛.
如果有一個關係.
你不用《臨前十三章》的這些要求.
那個不是真愛.
你只愛自己.
你對自己不會有自誇張狂.
作害羞恥的問題.
永遠是對身邊的人才需要.
我想說的是.
我們是否真的可以有一個完美的環境.
有一個完美的對象.
讓我們去實踐信仰.
讓我們去愛.
去關心.
去付出.
我只能夠說的是.
今天就是我們實踐信仰最好的時間.
你看到你身邊的這一班人.
這班岩岩沉沉.
不是盡都如你期望的人.
就是我們實踐信仰的最好可能.
唯一的可能.

$^{321}$追求尋找一個最理想最完美的對象.
等待一個最完美的時機.
讓我們去實踐信仰.
讓我們去愛.
讓我們去付出.
那你省一點吧.
你省一點吧.
你等到主再來都等不到.
都等不到.
是的.
真實的世界.
就不會是完美的世界.
真實的人.
就不會是盡都如我們所期望的人.
是不是.
不可以.
不過.
這個就是我們實踐信仰的場景.
上帝今天呼召你和我.
在這個世界.
在這個世代裡面.
去踐行祂的使命.
去兌現祂的教訓.
我們別無選擇.
別無其他的可能.
不要構想.
想像一個完美的環境.
期待一個完美的AI伴侶.
AI牧者.
AI教會.
AI信徒.
沒有這樣的事.
在腓立比書四章一到九節.
這一段聖經.
第一部分.
首先要處理一些人事的問題.
他提到兩個姐妹的名字.
有爹和秦都基.
兩位姐妹.
他沒有說很多他們的事.

$^{361}$他只是輕輕提到.
他們大概有爭拗.
有衝突.
有不開心.
並且不開心的衝突.
甚至已經影響整個教會.
不僅僅是他們私下的一些紛爭.
是什麼事.
保羅沒有詳細說.
不過他勸導教會的負責同工.
幫助這兩個姐妹解決問題.
很明顯已經上升到教會要處理的層次.
不僅僅是他們私下擺平的問題.
這兩個姐妹不知道什麼背景.
不過如果保羅離開菲律賓教會這麼久.
他仍然認識這兩個姐妹.
又聽到這兩個姐妹的問題.
很明顯這兩個姐妹一定是資深的姐妹.
不會是剛剛信主.
剛剛來的.
剛剛來的保羅都不會認識.
或者如果真是一個青年團契的兩個人吵架.
沒理由需要勞動保羅去處理.
一定是會影響比較大的.
走一步路都會震地.
那些人才會構成問題.
才會構成保羅要親手處理.
所以兩個姐妹一定是很資深的姐妹.
並且保羅還提到這兩個姐妹.
曾經在福音上跟他一起路苦.
更加說明他們不僅僅是聚會的姐妹.
這兩個應該是在教會裡面有事奉有參與.
並且很可能是曾經出錢出力.
就是教會的支柱的人物.
兩個教會資深的姐妹.
人事有衝突.
這不是一個很理想的情況.
不是一個很理想的情況.
不過這個情況我們大家都慣常知道.
可能你在旺鎮一段比較長的時間.

$^{401}$你被中壞了.
我們教會其中一個特徵就是我們關係很和諧.
我們很少派系衝突.
也沒有人事鬥爭.
所以你會覺得這是理所當然的情況.
但我要說的其實不是真的.
教會最難搞的是人事問題.
最複雜的是人事問題.
人事問題解決了.
其他所有的問題都容易解決.
家和萬事興這句中國的老土說話是很對的.
總之教會一旦有人事衝突.
就會投人生氣.
招頭爛額.
不是開玩笑的.
保羅很簡單要求教會的領袖.
去幫助這兩位姐妹去解決問題.
不過他反映給我看的.
仍然是我以前經常掛在口邊的一句話.
保羅仍然是希望用一個積極,善意,正面的態度.
來看待,處理教會的人事衝突.
他沒有具體去說這兩個人的問題.
沒有具體說保羅用什麼指示去處理他們的衝突.
沒有 他只是叫教會的領袖去幫他們解決問題.
讓他們能夠恢復同心.
但是保羅最少是提醒教會的同工.
這兩個姐妹都是好人.
這兩個姐妹都是對教會,對福音工作有幫助.
是正面的人物.
維持一個正面,樂觀的態度.
我要說是非常非常非常重要.
特別是在今天我們所處的這個時勢.
我不知道大家覺得活在我們這個世代有什麼感受.
我常常的感覺就是.
很多人給我一個很灰,很負面,很否定.
甚至很hostile的感覺.
和家人吃飯.
其中提到最近政府說要推行的垃圾徵費的問題.
一片嘮叨聲.
我很好奇.

$^{441}$我問是否因為徵費要平均徵100多200元.
是對很多人很大的負擔.
對那些基層的人當然是有負擔.
政府需要幫助那些中門或最基層的人去解決問題.
但徵費本身的原則是很難反對的.
你們羨慕你們追求去的國家.
台灣一個星期收三次不同的垃圾.
你們知道的.
你們在家要嚴格分類.
他們不是每間房子都很大.
不過都是要儲存垃圾.
定期送下樓.
在日本你們喜歡的地方.
一樣是做這件事很久了.
在美加就不用說了.
澳紐全部都是要垃圾分類並且要徵費.
原則上為甚麼要反對呢.
接著就說這個政府非常非常之無能.
你就是看他分類之後最後都是堆填的.
你就知道這個政府卑鄙無恥.
下流賤格.
所以他做的所有事都是錯.
他官司不好是一件事.
年輕人.
環保是關乎你們下一代的利益.
你明白我說甚麼嗎.
我的感覺就是很堅定.
總之甚麼都是先否定.
否定甚麼呢.
就可以不停轉移.
A不是就B.
B不是就C.
總之立場先堅定了.
否定然後找理由去否定.
A理由不行就B理由.
B理由不行就C理由.
我不是說我要用批判的眼光.
沒有的.
我只是感覺我們大家好像在一個很深的困局裡面.
那個困局先不要說是外在的困局.

$^{481}$我們自己在一個很強的情緒的困局裡.
總之就是很煩躁.
樣樣都看不慣.
並且是真的很憤世嫉俗.
我要說當我們經常有這樣的感受.
不要說對世界有甚麼問題.
沒有甚麼影響.
是對自己.
我要說.
你很難說得到保羅接著做的三個吩咐.
很難的.
所以對人對上帝常懷good will.
這個真是我經常掛在口邊的說話.
但它是真的很重要.
常懷善意.
善意 善意 善意.
不是在於我們要去拆甚麼人的鞋.
No No No.
是讓我自己可以拜得正.
我自己能夠實踐到信仰.
我常說.
終日忿忿不平.
終日怨氣滔天.
你很難活一個基督徒的樣子出來.
然後在這段聖經裡面.
在勸勉部分.
保羅做了三個很老土的要求.
簡單但重要的要求.
這三個要求可以變成揮春.
我便想起三個揮春.
所以便和大家說這三個要求.
第一個要求.
保羅提到常懷一個喜樂的態度.
常常喜樂.
四節到五節.
你們要靠主常常喜樂.
我再說你們要喜樂.
當叫眾人知道你們謙讓的心.
主已經勤了.
其實前面保羅在三章一節.

$^{521}$已經說過要靠主喜樂.
喜樂是菲勒比書中心主題的經文.
今天不和大家說太多喜樂的東西.
喜樂是一個積極的人生態度.
是一個信心的宣告.
保羅特別提到一句.
我覺得值得我們停住思考.
當叫眾人知道你們謙讓的心.
謙讓的心.
這是一個喜樂的要求的關鍵.
謙讓降低一些要求.
降低一些期望.
減少一些堅持.
太有品味的人.
是很難喜樂的.
我一吃你那一口蘿蔔糕.
不是糕來的.
蘿蔔剁得這麼碎.
太有品味了.
我這些什麼都吃下去都覺得很好.
一個過口都過不了很多時間.
所以我都未感覺到是什麼感覺.
就吃下去了.
總之有一種飽的感覺.
有一種滿足的感覺.
對我來說.
我吃東西容易快樂很多.
比你快樂.
太有品味.
太有要求的人.
是很難快樂的.
你同意嗎.
完美主義的人是永遠不會快樂的.
是永遠不會快樂的.
謙讓的心.
不僅僅是一個對身邊的人.
我讓一讓他位.
對任何事物.
一個讓步.
一個這樣的態度.

$^{561}$可以退後一點.
不要這麼堅持.
我不是沒有品味的人.
我不會監自己.
明明是很難吃.
都要說自己覺得很好吃.
但最少.
學我媽媽從小的教導.
好好醜醜吃一餐.
不是時必要精雕細琢的東西才吃的.
不一定要弄到最好的.
我才肯吞下去.
最少我的幅度大很多.
我快樂的機會又大很多.
沒有人叫你垃圾吃下去都覺得快樂.
你真的太變態了.
你幾乎每樣都當是垃圾.
你就沒有什麼快樂的了.
我保證的.
你沒有什麼可以快樂的了.
總之你的垃圾範圍越小.
你的快樂幅度就越大.
總之就是這樣.
當叫眾人知道你們謙讓的心.
我們的喜樂態度.
是反映我們在生活上的要求.
對人對己.
對這個世界的要求.
我覺得是很好的一個提醒.
我是不是一個謙讓的人.
還是一個事事都批評.
事事都覺得被冒犯.
事事都覺得這個世界對我不友善.
如果我是一個這樣的人.
我先不說其他.
最少就是你很難快樂.
你生氣的機會一定多過.
壓倒性多過你快樂的態度.
時間對不對.
然後他提到主已經近了.

$^{601}$這是解說我們為什麼要有一個謙讓的態度.
簡單來說.
我們的末日意識.
一定要幫我們去減少執著.
放下怨恨.
增加喜樂.
有很多基督徒說末日.
越說就越殺氣騰騰.
越說就越怨氣滔天.
看看你們這班罪人.
這班罪孽深重的人.
耶穌回來怎麼弄死你們.
看你們怎樣.
一歲收尾.
欣豬撈尾幾年.
上刀山下油鍋.
看著你們怎麼死.
上上好像彼得和耶穌說.
叫他降天火.
去懲罰我們看不慣的那些人.
如果我們的末日意識.
是讓我們上上咬牙切齒.
對這個世界充滿怨恨和咒詛.
這是一個扭曲錯誤的末日態度.
保羅提主已經近了.
這個字的真正意思是主是近的.
主是近的.
其實是說我們的主很近我們.
是我們很近我們.
當我們知道這個世界很多東西會成為過去.
人間眼前的事會很快結束.
我們就知道有些執著其實不需要.
有些心理糾纏其實不需要.
我們沒有什麼是需要堅持到最後一滴血.
需要執著到盡.
不用的.
贏贏輸輸其實很小事.
其實很小事.
很小事.
我們的主.

$^{641}$我們最期望的那個chapter.
很快會展開.
會鋪開.
今天我們就簡單一點.
隨意一點.
我們不需要太糾纏.
不需要爭最後一滴血.
寸土必爭.
不用這樣.
基督徒可以隨意自在.
就可以謙讓.
就可以上懷起樂.
上起樂第一個要求.
第二個要求是不要憂慮.
是不是很老套.
不過輝聰也是寫這些東西的.
不要憂慮.
第六第七節.
應當一無掛慮.
將你們所要的告訴上帝.
上帝所賜出人以外的平安.
必在基督耶穌裡.
保守你們的心懷意念.
交托.
這個是不要憂慮的關鍵.
我們為什麼可以不憂慮.
是因為我們學會交托.
學會交托有兩個認定.
第一個認定.
當然我們都知道的就是.
上帝一定行.
上帝行.
是不是.
上帝行.
不過.
確定上帝行的另一個前提是.
我知道我自己不行.
什麼時候你真的心悅誠服.
知道自己經營不起你的生活世界.
經營不起你的各種的關係.

$^{681}$各種的責任.
你的交托.
是真真實實的交托.
而不會.
還要拿著.
拿著.
說交托.
但是又瞻前顧後.
又苟且不憂.
為什麼我們可以一無掛慮.
就是耶穌在登山寶訓所說的.
我們知道我們掛慮沒有用.
一無掛慮不是信心的表現.
或者不是知道自己信心堅強的表現.
而是知道自己軟弱的表現.
耶穌是很真實很殘酷的跟你說.
你想得怎樣呢.
你掛慮有用嗎.
你的思慮可以使你受數增加一刻嗎.
有用嗎.
想 想什麼想.
你那麼少.
你當然想不通.
你以為你想有用嗎.
沒有用.
既然沒有用就收手吧.
對不對.
想不通就不要想.
然後真是.
由衷的.
真真實實的交托.
交托不是簡單的hand over.
而是真正的surrender.
是不是.
你來 我不行.
我不行.
學習一無掛慮.
將我們所有的東西.
我需要的.
我想望的.

$^{721}$告訴我的主.
好了.
我不是說禱告後的效果.
效果就是.
祂會供應 祂會幫助 祂會解決.
我們就會心想事成萬事勝意.
到頭祂說的就是.
你交托的效果就是.
你可以維持心中的淡定平安.
上帝會賜出人以外的平安.
去控制著.
佔據著你的心懷意念.
是不是.
平安會長駐在你裡面.
有那種淡定.
安穩 實在的感覺.
我都說了給主聽.
是不是.
祂都知道我的需要了.
祂怎麼供應我不知道.
總之祂在這裡.
我可以定點.
我們上上祈禱.
上帝賜我們平安喜樂.
喜樂的前提是謙讓.
而平安的前提是交托.
是不是.
有謙讓就容易喜樂.
真實的交托就能夠平安.
最後一點.
凡事學習.
第三個要求.
弟兄們我還有未盡的話.
凡是真實的可敬的.
公義的 清潔的 可愛的 有美名的.
有什麼德行 有什麼稱讚.
這些事你們都要思念.
你們在我身上所學習.
所領受 所聽見 所看見的.
這些事你們都要去行.

$^{761}$賜平安的上帝.
就必與你們同在.
這段經文對八第九節.
我以前跟大家說過.
我們盡量去學習.
美善的事.
眼中多些看到身邊的人和事的美善.
然後我們就努力逼自己去學習.
不要看到別人太多的缺點.
看他的強項.
看他正面的好處.
然後就要求自己見然思齊.
見然思齊.
要求自己學習.
所有人你都覺得他差.
你就沒有東西可以學習.
你已經所向無敵.
還有什麼可以學習呢.
永遠看到自己不行.
身邊的人有好處 有強項.
永遠看到.
這個世界有很多值得我們學習.
值得我們欣賞.
讓我們感動.
讓我們羨慕的地方.
然後就要求自己.
可以跟著改變自己.
前兩天.
我回建律師學院.
講早會.
今次的早會很淪盡.
因為在早會沒多久之前.
我們有兩個學院的小事件.
一個就是我以前的同事.
黃錫賢老師.
可能有人聽過他的名字.
很突然被主接去.
然後第二個就令學院更加震動的.
就是我們有一個同學.
在教會事奉中.

$^{801}$突然之間就死了.
二十多歲的女兒.
法醫現在說可能是急性主動脈爆裂.
可能是有些隱病.
大家不知道.
當然這個就造成整個學院.
很大的困擾和震動.
所以以前要講的訊息.
準備好的訊息都全部不能講.
臨時之前一晚.
本來我說不如不要講.
縮骨的.
反正我也不認識那個同學.
我也沒有教過他.
不過總之撿出來就講.
我跟同學講一些安慰鼓勵的話.
最後我在早會上提醒眾人.
基督徒我們有兩個最關鍵的信仰任務.
第一個是悔改更新.
第二個是實踐使命.
實踐使命今天不講了.
悔改更新是很重要的.
真的.
把握現在.
香港的竹馬說.
有錯就要認.
打就站定.
總之真的覺得自己生命有某些瑕疵.
立刻改.
不要拖了.
是不是.
有錯就認.
在上帝面前盡快認.
然後要求自己可以改變可以突破.
是不是.
動自己.
悔改更新是我們基督徒的其中一個主要信仰任務.
本來這個任務是伴隨你一生.
是不可以被其他事物取代的.
本來講過什麼.

$^{841}$你傳福音給別人.
你都可能自己會被棄絕的.
是不是.
所以你要努力的恭克己身.
叫身服我.
我們要努力完善自己的生命.
兌現信仰的要求在自己裡面.
努力使自己成為更好樣的基督徒.
努力讓自己的生活表現和信仰的要求吻合一點.
沒有人敢說做得到.
但吻合一點.
你總可以找到一個今天你努力實踐的議程.
要做的事.
然後你就會發覺.
信仰是紮紮實實.
不是不斷重覆.
說那些慣常的屬靈Dragon.
而是真的對我自己有要求.
謙虛 溫柔 忍耐.
是要求來的.
是真實的要求.
我知道我的死穴.
我知道我的問題.
我努力在對付中.
我努力在改變中.
凡事學習.
不僅僅是要我們在知識和經驗上盡心.
更加重要的是.
努力更新自己.
繼續更新自己.
這個信仰任務是伴隨我們一生.
永遠都未完成.
永遠是今時今日.
你可以做 你可以做 你可以做.
你努力在做中.
三個很簡單的信仰要求.
上懷喜樂.
不要憂慮.
凡事學習.
求主藉著一段這麼簡單的聖經.

$^{881}$這麼簡單常識性的要求.
提醒我們今年.
我們努力.
努力在主面前做好.
我們所在的世界不一定是一個理想的世界.
不一定是一個理想的世界.
香港的經濟環境不好.
香港的政治佈局也不理想.
我們面對很多風風雨雨.
連踢球也有很多事是非非.
是很多很多的麻煩.
很多的問題.
我只是說.
一個不完美的世界.
是上帝交付給你和我.
去實踐信仰.
去見證基督的世界.
讓我們.
就在這個不完美的世界.
是我們唯一的給予.
我們唯一的吸定.
整定.
我們在這裡.
我立志在這裡預見我的主.
我立志在這裡去榮耀我的主.
我立志在這裡讓基督的名.
被人不僅僅聽到.
是看到.
上帝與我們大家同在.
大家同意就這樣.
\newpage



\section{約書亞記 1:1-9}
\label{sec:QTbAXdDuZOw}
\textbf{2024年9月15日崇拜直播｜陳冠成牧師｜正視你的情緒－從神而來的勇氣｜約書亞記1\_1-9}
\newline
\newline
連結: \href{https://youtube.com/watch?v=QTbAXdDuZOw}{\texttt{https://youtube.com/watch?v=QTbAXdDuZOw}} ~~~~ 語音日期: 2024-09-16
\newline
\newline
\hyperref[sec:uj5FB161Gvc]{\small{< < < PREV SERMON < < <}}
~
\hyperref[sec:index_chronic]{\small{[返順時目]}}
~
\hyperref[sec:QHvBwcg9HEw]{\small{> > > NEXT SERMON > > >}}
\newline
\newline
約書亞記 1:1-9
\newline
\begin{longtable}{cl}
\hline
\hline
章節 & 經文 (和合本修訂版)\\
\hline
1:1 & \begin{tabularx}{0.7\textwidth}{X} 耶和華的僕人摩西死了以後，耶和華對摩西的助手嫩的兒子約書亞說： \end{tabularx} \\ \\ \relax
1:2 & \begin{tabularx}{0.7\textwidth}{X} 「我的僕人摩西死了。現在你要起來，和眾百姓過這約旦河，往我所要賜給以色列人的地去。 \end{tabularx} \\ \\ \relax
1:3 & \begin{tabularx}{0.7\textwidth}{X} 凡你們腳掌所踏之地，我都照我所應許摩西的話賜給你們了。 \end{tabularx} \\ \\ \relax
1:4 & \begin{tabularx}{0.7\textwidth}{X} 從曠野和這黎巴嫩，直到大河，就是幼發拉底河，赫人的全地，又到大海日落的方向，都要作你們的疆土。 \end{tabularx} \\ \\ \relax
1:5 & \begin{tabularx}{0.7\textwidth}{X} 你一生的日子，必無人能在你面前站立得住。我怎樣與摩西同在，也必照樣與你同在；我必不撇下你，也不丟棄你。 \end{tabularx} \\ \\ \relax
1:6 & \begin{tabularx}{0.7\textwidth}{X} 你當剛強壯膽，因為你必使這百姓承受那地為業，就是我向他們列祖起誓要給他們的地。 \end{tabularx} \\ \\ \relax
1:7 & \begin{tabularx}{0.7\textwidth}{X} 只要剛強，大大壯膽，謹守遵行我僕人摩西所吩咐你的一切律法，不可偏離左右，使你無論往哪裡去，都可以順利。 \end{tabularx} \\ \\ \relax
1:8 & \begin{tabularx}{0.7\textwidth}{X} 這律法書不可離開你的口，總要晝夜思想，好使你謹守遵行這書上所寫的一切話。如此，你的道路就可以亨通，凡事順利。 \end{tabularx} \\ \\ \relax
1:9 & \begin{tabularx}{0.7\textwidth}{X} 我豈沒有吩咐你嗎？你當剛強壯膽，不要懼怕，也不要驚惶，因為你無論往哪裡去，耶和華你的神必與你同在。」 \end{tabularx} \\ \\
[1ex]
\hline
\hline
\end{longtable}
$^{1}$我們以下要聆聽神的說話.
今天我們要聽聽經文是記載在約書亞記第一章的一至九節.
約書亞記第一章.
耶和華的僕人摩西死了以後.
耶和華曉遇摩西的幫手.
亂的兒子約書亞說.
我的僕人摩西死了.
現在你要起來和眾百姓過這約旦河.
和我所有次級以色列人的地去.
凡你們腳掌所達之地.
我都照著我所應許摩西的話次級你們了.
從曠野或者黎巴嫩直到巴拉大河.
下人的全地有到大海日落之處.
都要作你們的警戒.
你平生的日子.
必無一人能在你面前站立的住.
我怎樣與摩西同在.
也必照樣與你同在.
我必不撇下你.
也不丟棄你.
你當剛強壯膽.
因為你必使這百姓承受那地為業.
就是我向他們列祖.
起誓應許設給你們的地.
只要剛強大大壯膽.
謹守遵行我僕人摩西所吩咐你的一切律法.
不可偏離左右.
使你無論往哪裡去都可以順利.
這律法書不可離開你的口.
總要就耶思上.
好使你們遵守遵行這書上所寫的一切話.
如此你的道路就可以亨通.
凡事順利.
我希末有吩咐你摩.
你當剛強壯膽.
不要懼怕.
也不要驚惶.
因為你無論往哪裡去.
耶華華你的神必與你同在.
聖經獨筆.

$^{41}$各位姐妹平安.
本來這星期梁家倫院長來到港島.
但因為他還在外地.
趕不及回港.
所以在講壇內有少少調配.
今天我們先回到.
「正視情緒」系列的宣講.
今天我們會透過剛才所讀的經文約書亞記.
第一章來到去.
我們說人生難免會遇到一些.
你害怕或懼怕的事情.
我們怎樣來到去透過信仰.
向上帝支取足夠的勇氣.
來面對呢?.
正正是今天想和大家分享.
如果你剛才留意.
經文一開始交代了時空背景.
第一句就說.
耶和華的僕人摩西.
怎樣?.
死了.
如果你留意.
耶和華的僕人.
在聖經內不是一個普遍的稱呼.
不同我們今天說.
我們每個人都是上帝的僕人.
不是一個簡單的稱呼.
其實這個字原文在舊約聖經.
出現了23次.
舊約這麼厚.
只有23次.
哪些人拿得最多呢?.
摩西.
摩西拿了18次.
剩下的有5次.
2次給約書亞.
2次給大衛.
1次是用來負面形容上帝的百姓.
即是諷刺他們.
你可以想像而知.

$^{81}$上帝不是那麼輕易.
稱呼他的子民.
做耶和華的僕人.
約書亞你猜猜他何時才拿到這個稱呼?.
如果你有留意.
剛才我們讀了一段經文.
一章一節.
說當時約書亞的稱呼是甚麼?.
摩西的助手.
摩西的幫手.
他何時才拿到這個稱銜呢?.
經文說到大綱印了.
約書亞的結尾.
24章第29節.
到他離世的時候.
他就拿到這個稱銜.
被稱為耶和華的僕人.
換句話說.
如果根據聖經學者所說.
約書亞大概是80歲開始.
接手摩西的工作.
帶領以色列人進入迦南地.
到約書亞的結尾.
他離世的時候.
記載了他當時是110歲.
他那一刻被稱為耶和華的僕人.
上帝花了多少年在他身上?.
這個訓練.
起碼30年.
30年來到他輔助約書亞.
由摩西的幫手.
慢慢晉升為耶和華的僕人.
我們要看到上帝給予一個人的栽培.
事實上這30年容易嗎?.
不容易.
雖然在這30年裡.
摩西確實離開了大眾的視線.
離世了.
但他在百姓心目中.
仍然是一個偉大的領袖.

$^{121}$你翻查一下.
就算到新約都還在講哪個律法?.
摩西的律法.
你就知道他的影響力多厲害.
摩西可以說是整個以色列民族.
他的歷史地位.
基本上是不會忽視的.
沒有人會忽視他.
他不會容易動搖.
事實上就算你只看約書亞記.
摩西的名字.
即使離世了.
也出現了50多次.
包括提及他的名字.
或提及他的教導.
你可以想像在百姓心目中.
摩西的聲望影響力.
已經深入人心.
百姓自然對他的繼任人.
一定會有期望.
很自然的想法.
至少期望約書亞會怎樣?.
作為繼任人.
起碼都維持現況.
起碼都延續摩西帶隊的效果.
人很自然這樣期望.
沒理由期望他變差.
甚至最好是怎樣?.
超越摩西的往績.
這樣就更加好了.
在這個情況下.
作為繼任人.
作為摩西的助手.
約書亞會怎樣想?.
有沒有壓力?.
他帶隊有沒有壓力?.
其實一定有的.
不難理解.
一定會承受不少壓力.
在座有沒有人喜歡看英超?.

$^{161}$有沒有?.
其實我不看的.
不過曾仔牧師也看.
英超.
你覺得英超最難帶隊的是哪一隊?.
有人說英超的曼聯.
是英格蘭足球壇.
第二難做的工作.
做他的領隊.
是第二份最難的工作.
第一份是帶領英格蘭國家隊.
為什麼?.
因為很久沒拿過世界盃.
很難帶.
而且球星太多.
英超的領隊是一份很艱難的工作.
因為我們知道.
在曼聯的領隊是一份很艱難的工作.
因為曼聯有一個很傳奇的主帥.
費甲遜.
他在他帶領下.
為球會開創了最輝煌的成績.
但當費甲遜2013年退休的時候.
球迷自然怎樣?.
也期望他的繼任人.
可以延續他的佳績.
所以大家一定記得當年.
誰接班做費甲遜帶領曼聯.
記不記得?.
不記得?.
費甲遜親自欽點莫耶斯.
有沒有印象?.
有看英超應該有印象.
莫耶斯來做曼聯的新領隊.
如果大家有記得.
就算我沒看英超我也知道.
莫耶斯有個綽號是什麼?.
當他被費甲遜欽點的時候.
就變成了操身瘟.
有沒有印象?.

$^{201}$大大橫額.
你就是那個真命天子.
是你的了.
選了你做接任人.
做繼任人.
你就是那個被選上的.
但我們知道.
莫耶斯接棒.
不單止接棒.
還接了什麼棒?.
接了火棒.
燒了自己.
球迷對他寄予厚望.
但可惜事與願違.
那一屆曼聯在賽事裡四大皆空.
什麼獎都沒有.
從天堂一下子掉到地底.
更加失去了.
來屆參加歐聯的資格.
其實很正常.
因為當年費甲遜.
大家不記得了.
他最初帶隊.
曼聯也不是那麼好.
也要熬了幾年才有成績.
但沒辦法.
莫耶斯有一個高峰接手.
突然下滑.
自然他突然變成千夫所指.
最後他也逃不過被炒友的結果.
其實他上任很短.
一年也不夠.
他就由The Chosen One.
不知道之後球迷撇起他.
寫了一句什麼橫額.
The Wrong One.
慘不慘?.
你是真命天子.
一年不夠.
Sorry 我們選錯你.

$^{241}$你不是.
很慘吧.
但自此之後.
其實你有留意.
每一屆曼聯新上任的領隊.
他們帶隊成績一定會被人拿來比較.
就是和最厲害的費甲遜比較.
和他往績比較.
有沒有壓力?.
其實一定有.
一定有一種無形壓力.
繼任人是不容易做的.
所以你看看約書亞.
他是不是The Chosen One?.
他是不是?.
其實他是.
他是被選擇.
你不要以為在約書亞的記才任命.
在申命記上帝選了他.
一早已經是上帝任命他.
成為摩西的接班人.
並且不是空降.
你留意一下抗疫那四十年.
約書亞一直在摩西身邊做助教.
協助摩西帶領管理百姓.
約書亞也親自帶隊.
替摩西打贏過不少仗.
所以其實約書亞.
他不是新人.
他不是初出茅廬.
他是久經訓練.
但在這一刻要自己擔正去領隊.
很明顯其實也有壓力.
也有憂慮.
所以你看到為何上帝會主動向他顯現.
特別你留意剛才讀經到結尾第九節.
上帝叫他什麼?.
不要害怕.
不要驚惶.
大家覺得約書亞害怕什麼?.

$^{281}$首先不知道你有沒有試過.
被指定繼任人.
有嗎?.
其實本身就是一種壓力.
什麼意思?.
譬如我指著你.
是你的了.
是你,不要忘記你.
就是你.
選了你.
It's you.
即是代表什麼?.
你不能避.
不能躲.
以前你可以站在後面.
但現在你站在最前.
你不能卸給你旁邊的人.
你只可以怎樣?.
你面對這個崗位.
你硬頂上.
所以本身你被指定為一個繼任人.
你本身就會有壓力.
一定有.
以前你是執行指令.
但現在你發私號令.
你變成了最高的負責人.
但如果做得不好.
就會變成最高的文責人.
你就會變成像莫耶斯.
所以一定會有壓力.
特別假如你的前任成績是好的.
是優秀的話.
繼任人都會給自己壓力.
特別可能有時看新聞都知道.
如果一些家族生意的承繼人.
上一代做得很好.
現在生意交到自己手上.
自己又怎樣?.
你會拼命去做.
因為你不想有辱家門.

$^{321}$所以繼任人很容易會對自己.
產生不合理的要求或期望.
給自己很多壓力.
甚至產生恐懼感.
另一方面你有留意.
大家覺得礦業是否一個很理想的生活環境?.
應該不是吧?.
但你有否留意.
約書亞和百姓住在那裡多久?.
40年了.
甚麼都習慣了.
不好的,難受的.
總有辦法去適應和克服.
所以他們一定很明白.
礦業的生存之道是怎樣.
所以礦業對他們來說.
是一個安書區.
住慣了40年.
但現在呢?.
現在你要離開一個安書區.
進入迦南.
不是短期旅行.
不是一轉日本.
兩個星期後回來.
是在那裡住一輩子.
你將要踏入一個.
不行.
將要離開一個你熟悉的地區.
踏入一個全新.
有很多未知的生活環境.
你一定會面對陌生感.
你一定會面對很多無法掌握.
甚至控制的狀況.
所以你一定會害怕.
再說我們會知道.
當時在迦南是否沒有人住?.
不是.
不是沒有人的地方.
是已經有很多不同的民族.
佔據了很多年.

$^{361}$有很多地膽.
有很多地方的勢力在裡面.
有很多甚至是敵人.
比以色列人更強大.
所以.
這確實是上帝應許的留賴與密之地.
但同時又很吊詭的是甚麼?.
是敵人盤踞的地方.
進入迦南一定不是一步到位的事情.
你會預見到是一種長期的爭戰.
過程裡面因為爭戰.
所以也有機會有損傷.
所以.
進入迦南確實是上帝應許的一件好事.
但約書亞身為領袖.
我們說任重道遠.
他難免會比別人有更多的壓力.
更多的驚惶.
我想正如大家.
你有沒有試過承接一些新的崗位?.
有沒有試過升職?.
由被人管理到管理人?.
可能也有.
有沒有試過升級做父母?.
有,大家眨眨眼.
有沒有試過在教會.
由一個普通的基督徒.
變成團契或小組的職員團長?.
有嗎?也有.
甚至乎成為教會的領袖.
蒙召去讀神學.
裝備自己成為全職的牧者.
通通這些都是進入一個新的崗位.
人生新的階段.
你要由一個被帶領的人.
成為一個帶領者.
是不是好事?.
大家都靜下來.
部長是不是好事?.
領袖是不是好事?.

$^{401}$應該是吧,我逼你說.
沒得說不是吧.
固然是上帝所賜的美事.
是值得興奮和喜悅.
但同時你總是夾雜著什麼感覺?.
憂慮或害怕的感覺.
因為你怕什麼?.
你怕可能做得不好.
做得不夠.
別人會怎樣看我?.
我能否升任?會質疑自己.
我會否辜負上帝給我的期望.
辜負眾人的期望?.
我們會有很多這些.
我們說的負面的狀況.
於是面對這些狀況.
我們通常有三個情況.
大家有留意嗎?.
你試試檢視一下.
我們說三個F.
第一個就是Fight.
Fight.
即是怎樣?.
你去對抗.
你拼命去回應.
第二個呢.
多於L.
Flight.
即是什麼?.
你避開.
你不去回應.
逃避.
不想碰他.
第三個F呢.
你Freeze了.
凍結了.
僵硬了.
即是你不是迴避.
又不是面對.
你不懂得反應.

$^{441}$我們說所謂的「當機」.
弟兄姊妹.
如果用這個角度來看.
其實約書亞的處境.
其實我們一定會面對.
我們怎樣可以.
從信仰裡面去支取到.
上帝給我們的勇氣和力量.
來坦然面對新的挑戰.
裡面所帶來的壓力.
甚至恐懼呢?.
我們一起看看上帝怎樣去回應.
或者怎樣引導約書亞.
我們去到第二節.
上帝顯現.
跟約書亞說.
摩西已經離世了.
現在你要起來.
和眾百姓過這約旦河.
往我所要刺給以色列人的地方去.
當一個人.
特別在事奉裡面.
經歷懼怕.
上帝第一件事是怎樣?.
讓他穩定下來.
要一個人穩定.
最重要是任命和方向.
你看見上帝就是做這件事.
上帝親自任命約書亞.
呼召他.
亦給了一個很清楚的.
job nature.
job description.
很清楚的方向給他.
你要帶領百姓入迦南.
其實是提醒約書亞甚麼?.
你不是接手摩西的工作.
你是回應上帝對你的呼召.
這就是事奉.
跟我們做一般工作不同.

$^{481}$任何事奉.
首要是回應上帝給你的呼召.
不是別人對你的邀請.
雖然上帝會透過人.
對你發出邀請.
但核心是來自上帝.
對你的呼召和任命.
雖然約書亞一直是摩西一手大大.
他的行事為人的方式.
亦是建基於摩西的教導.
但上帝在這裡清楚表示.
你跟摩西是不同的.
你們不是相同.
約書亞你可以有自己的思考方式.
和帶領的風格.
你不需要成為另一個摩西.
你亦不需要模仿摩西所做的一舉一動.
變成了他的複製品.
變成了摩西的2.0.
重點是.
實質捉緊上帝給你的呼召和應許.
清楚上帝給你要走的路.
這樣你就不會容易.
在事奉的路上迷失或懼怕.
並且謹記上帝說.
我不會丟下你.
我不會丟下回應我呼召的人.
我必與你同在保守你.
正如我對摩西一樣.
第三節特別說.
凡你們腳掌所踏之處.
我都照著我所應許摩西的話.
賜給你們了.
換句話說.
這裡表示.
上帝給我們的福氣.
其實遠遠大於我們所想像.
有時我們經歷上帝的帶領太少.
經歷他的福氣太少.
這段經文.

$^{521}$告訴我們答案.
不是上帝不願意給.
上帝很樂意給.
而是我們夠不夠膽去拿.
我們夠不夠膽.
有沒有勇氣去領取.
上帝的應許.
上帝的福氣.
我們有沒有信心走近他.
簡單來說.
來回應他對我們的呼召.
他對我們的任命.
事實上這裡給我們一個重要訊息.
弟兄姊妹.
上帝的呼召.
上帝的應許.
包括他的同在.
就是你與我事奉的根基.
任何事有迷失的時候.
你回到這三個東西.
你就站穩了.
你知道不是出於人的緣故.
而是出於上帝.
你就可以在事奉的風浪裡.
你不會倒下.
你可以懂得重新定位.
因為上帝給你一個方向.
給你一個肯定.
你可以專注在他給你的使命裡.
你這樣就會慢慢有力.
走起來.
所以弟兄姊妹.
假如你都在事奉路上.
遇到瓶頸.
遇到挫折.
我盼望這段經文.
首先引導我們.
不是亂解決問題.
是回到上帝那裡.
你記住是上帝任命你.

$^{561}$上帝放你在這個位置.
他不會丟下你.
你首先回應他對你的呼召.
他交給你的使命.
讓你可以穩藏在.
上帝這個避難所被封港的裡面.
你經歷他的醫治和安慰.
以至你可以重拾.
從上帝而來那份信心和勇氣.
重新踏上那條事奉之路.
我還記得.
我在教會師班事奉.
大家猜猜有多少年.
我2012年開始.
接替洪姑娘的工作.
我翻查我的記錄.
現在已經十三年了.
你猜我有沒有疑惑?.
有很多人問我.
你有沒有讀過音樂.
你有沒有彈琴.
我說我通通沒有.
我經常懷疑.
為何上帝要放我在這個位置.
我不是做音樂出身.
我不是靠唱歌賺飯錢.
沒錯.
我之前在舞會.
我有參加過師班事奉.
我也在崇拜令的師.
但你知道嗎?.
做指揮是兩回事.
我沒有學過.
對一個沒有學過指揮的人.
他每個星期出來指揮.
你覺得是很痛苦的事.
偶爾我都會困惑.
我會卻步.
我會所謂懷疑人生.
你這樣安排真好.

$^{601}$沒有人了嗎?.
有一次我和一位弟兄姊妹.
分享這個掙扎.
她說上帝不會選錯人.
你去吧,繼續吧.
有點被困的感覺.
但很重要.
她說上帝不會選錯人.
不是我供奉他.
是上帝選的.
如果你給我,我不會選.
因為害怕.
這句說話叮囑我.
是上帝呼召我參與他的工作.
所以主角是誰?.
主角是他,不是我.
我只是一個配合他的器皿.
由始至終不是因為教會缺人手.
上帝沒有人使用.
所以找我.
更不是因為我有什麼能耐.
可以擔任這個崗位.
他選我.
他呼召我.
是因為他主動走近我.
他想和我建立關係.
上帝都期望.
我透過我的行動.
努力去回應這份關係.
回應他的呼召.
在這個關係裡面.
一路去經歷.
他的同在的應許.
和他奇妙的供應.
所以多年來.
我體會到.
其實是誰在指揮我?.
是上帝在指揮我.
透過我一個牧者的身份.
在每次練習.

$^{641}$每次獻狀.
向上帝分派的力量.
圍繞給每一個人.
就是這樣奇妙的.
年復年.
向上帝支取勇氣和信心.
一步一步走近上帝.
認識他.
亦都活出上帝在我生命裡面.
一些獨特的旨意.
頂尖的這個才是事奉.
不是KPI.
我們常常覺得是業績.
是成績.
可能這些都是重要.
但不是最重要.
上帝不是沒有人用.
上帝甚至自己處理好都可以.
為何要找你找我?.
是因為關係.
唯有透過事奉.
我們才走近他.
亦都知道他就在我們身旁.
所以頂尖姐妹.
真是唯有我們.
投入在事奉裡面.
回應上帝給你的呼召和應許.
你才可以近身認識到上帝是怎樣.
你才在上帝裡面.
體會到你真實的價值是甚麼.
你真正的身份是甚麼.
當然在事奉裡面.
是會有壓力.
是會有阻力.
因為事奉的本質是甚麼.
你委身聽上帝講.
不聽自己講.
一定有壓力.
一定有阻滯.
你委身聽上帝講.

$^{681}$不聽隔離的講.
怎會沒有人際方面的不和諧呢?.
是一定會有的.
但這一切的掙扎.
這一切的疑慮.
其實是讓我們發現自己的本相.
你看到自己的不足.
看到自己的限制.
看到自己的脆弱.
於是怎樣?.
你便快點回到上帝那裡.
你快點抓緊它.
連結上帝.
你便立即經歷到他的應許.
經歷到他的安慰.
甚至乎是醫治.
你亦都可以從他的壓力.
從他的驚惶裡面.
你可以放出來.
你可以被治好.
其實就是這樣.
點點滴滴強化了我們和上帝的關係.
事實上我們和上帝的關係.
大家有沒有玩過攀石或攀山?.
或者見過人去攀石場玩?.
如果他是初學者.
下面一定有教練.
還有一條甚麼繩?.
有條救生繩.
是誰去爬?.
是否教練去爬?.
不是吧?.
教練看著學員去爬.
是吧?.
我們的關係其實就是這樣.
和上帝的關係.
是吧?.
因為有教練保護著.
看著.
所以做學員的可以怎樣?.

$^{721}$你放膽去爬吧.
放膽去爬吧.
一路向上爬吧.
你不需要擔心.
失足或掉了繩.
或者說失足或掉手.
是吧?.
因為你知道教練.
他拿著一條安全繩.
是吧?.
他隨時可以救你.
亦都一定救到你.
同樣上帝.
其實他有辦法獨自完聖所有工作.
是吧?.
為何很多事情.
他仍然邀請我們去參與?.
他是想鼓勵我們.
相信他.
相信他的話語.
你放膽去建立教會.
放膽去探索不同的福音機遇.
在過程中會否有失手.
會否有出錯?.
當然會有.
我們只是軟弱的塵土.
但是否在上帝計算中?.
亦都當然是.
如果我們認定上帝是掌管的話.
掌管背後的意思是甚麼?.
如果我們有機會失手.
但他都能處理.
所以他有絕對的能力保護和拯救我們.
並且即使我們行錯了.
他會帶我們回正路.
就像他對約書亞所說.
就像他對約書亞的英雄所說.
所以最後我們看六至九節.
我想大家一起讀一讀.
印在秩序表或看螢光幕.

$^{761}$好,六至九節,預備起.
(念念不忘).
謹守遵行我僕人摩西所吩咐你的一切律法.
不可偏離左右.
使你無論往哪裡去都可以順利.
這律法書不可離開你的頭.
總要就耶穌上.
保持你謹守遵行這史上所示的一切話.
如此你的道路就可以行動.
凡事順利.
我豈沒有吩咐你呢.
你剛剛強壯膽.
不要懼怕也不要驚慌.
因為你無論往哪裡去.
也和你的神必與你同在.
我常說重要的事說多少次.
你找到多少個剛強壯膽.
也是三次.
很明顯剛強壯膽是段落一個焦點.
一方面固然是因為他們即將要面對征戰.
他們準備承受迦南地為業.
亦表示他們要和迦南人展開戰役.
固然需要上帝給予他們勇氣.
剛強壯膽面對迦南人和他們的征戰.
但更重要的是.
他們打贏了所有仗.
安定居住在迦南之後.
還需不需要剛強壯膽.
特別在生活遇到隱有.
試探的時候.
是否需要勇敢進行上帝一切的吩咐.
是否需要?.
這就是重點.
假如純粹搬到迦南住.
安居在那裡.
就是上帝指明的終極目標.
其實最快捷的方法是甚麼?.
其實不用征戰.
甚至不用漂流曠野.
上帝應該怎樣?.

$^{801}$我們經常說曠野直送.
拯救了你出埃及.
直接送到迦南.
又或者上帝先用他的大能.
親自在迦南進行大清潔.
掃走所有不喜悅的人.
所有外敵.
掃乾淨後.
交鎖匙給摩西或約書亞.
交吉.
你們搬進去住.
這就是完成目標.
是最簡單直接.
一定不會出錯.
你拜託人去工作.
一定出錯.
聖經會讓我們知道.
以色列人沒有因為搬進迦南住.
從此安居樂業.
沒有.
反而你去到約書亞.
記下一本書.
士師記說甚麼?.
全部進去住.
已經沒有打仗.
但因為沒有上帝的話語.
更大件事.
就正如.
我們不會因為每天.
回到旺朕的禮堂.
這個建築物.
每天回到.
搬進來住也好.
不會自動安居樂業.
是就好.
明天一定迫爆.
不會.
上帝說得很清楚.
人真正的保障是甚麼?.
是聽他的話.

$^{841}$真正的保障不是住在哪裡.
不是住在迦南.
而是持續遵守他的吩咐.
特別第七至八節.
這裡說得很清楚.
只要剛強大大壯膽.
謹守遵行我僕人摩西所吩咐.
你一切律法.
不可偏離左右.
使你無論往哪裡去.
不限於迦南地.
都可以順利.
這律法書不可離開你的口.
總要晝夜思想好.
使你謹守遵行.
這書上所寫的一切話.
如此你的道路.
亦都沒有說明在哪裡.
就可以亨通凡事順利.
其實你會留意整個五經.
特別由《出埃及記》到《申命記》.
甚至乎《約書亞記》.
摩西帶領以色列人出埃及.
擺脫了他們適有為奴的生活.
顯明了上帝拯救的決心.
不過出埃及只是拯救過程中.
其中一站.
百姓的忠誡目的.
不單止要進入迦南.
住在迦南.
承受應許之地.
更重要的是什麼?.
安穩在上帝的應許.
這才是真正的安息.
如果我們用希伯來書所說.
就像第八節所說.
「這律法書不可離開你的口.
總要晝夜思想.
好使你謹守遵行.
這書上所寫的一切吩咐.

$^{881}$就是你的道路可以亨通凡事順利」.
其實這裡讓我們看到.
以色列人進入迦南之後.
他們仍然有很多.
不同的信仰挑戰和衝擊.
會引誘他們離開上帝.
離開亨通的道路.
不再以上帝的吩咐為生活依歸.
最簡單就是.
他們以前在抗疫吃什麼?.
馬欖.
進入迦南.
吃好了.
住好了.
但是要自己耕種.
自己畜牧.
問題來了.
還守不守安息日?.
星期日做什麼好?.
安息日做什麼好?.
謀生好?.
還是聽上帝的話好?.
再來一個問題.
敬拜迦南的偶像.
他保我們風調雨順.
有好的收成和收入.
上帝好像沒有說這番話.
跟誰好?.
他們會面對的問題.
我們也會面對.
考驗上帝的子民.
對上帝的忠誠.
今天也一樣.
有不同的活動.
來考驗我們對上帝的忠誠.
我們對崇拜,小組,團契.
對事奉的重視.
尊重程度是怎樣?.
會不會也被周邊的事物吸走了?.
這也是我們的考驗.

$^{921}$會不會漸漸偏離上帝給我們的話語?.
聖經裡強調.
我們存活的核心.
是要順靠上帝口裡所說的一切話.
舊約是這樣說.
耶穌也是這樣說.
上帝所說的一切話.
包括聖經的教導.
包括他的呼召.
包括他的應許.
等於是上帝提醒我們.
我們要相信他所說.
多於這個世界說給我們知道的事.
我們要相信上帝多於自己.
所以為此.
上帝叮囑我們.
無論去到哪裡.
都要以他吩咐的聖經話語.
來回應這個世界給我們的挑戰.
或者人有.
這樣我們才可以安穩.
在上帝喜悅的道路裡.
經歷到他的公認帶領.
過一個得勝的生活.
上帝要求我們.
晝夜思想他的說話.
晝夜思想其實可圈可點.
以我為例.
我現在這個身形.
我經常晝夜思想.
可以怎樣變回二十年前的樣子.
可不可以只是晝夜思想就能搞定?.
不可以的.
不睡覺也不行.
晝夜思想代表什麼?.
代表兩個向度.
你所做的事是否對你有益.
第二是.
你怎樣做好事.
使自己有益.

$^{961}$套上去.
聖經裡面就是.
世界鼓吹的價值.
我們所做的事.
是否合乎聖經的教導.
是否合乎上帝的心意.
有沒有榮臣益人.
我們要晝夜思想.
而正面來看.
無論在什麼環境.
什麼情況.
我們能否靈巧.
以不同的方式.
竭力.
直接間接.
將上帝的吩咐.
落實實踐出來.
並且.
上帝對約書亞說.
也意味著對整個群體所說.
不單單是教會的領袖.
而是整個教會群體.
一起用盡方法.
搞盡腦汁.
怎樣去將上帝的話語.
不單單只在我們這一代活出來.
並且在每一個世代.
都重視上帝的話語.
所以弟兄姊妹.
其實.
今日經文講得很清楚.
每一個蒙拯救的人.
不可能總是漂流在曠野.
如果你有讀過.
啟示錄的特別.
你會發覺.
啟示錄.
你沒有發現.
講得救這兩個字.
他assume.

$^{1001}$他的讀者都已經得救.
但是他十三次提到.
要得勝.
他很強調耶穌基督已經得勝.
並且他勉勵我們.
都要靠著耶穌來得勝.
當耶穌基督再來的時候.
誰會和他有份.
莫稱為忠心良善.
獲賞賜榮耀冠冕.
啟示錄說.
就是那些勇敢.
堅持主導.
靠主得勝的門徒.
所以上帝不單單.
要求我們靠他得救.
上帝要求我們要靠他得勝.
而約書亞記.
記錄著上帝的子民.
如何從得救.
到得勝這條道路.
亦勉勵我們.
能否今天捉緊上帝給我們的呼召.
應許和話語.
勇敢地過得勝的人生.
只要我們願意貫徹.
上帝的吩咐.
上帝應許無論你去到那裡.
無論你落入什麼境況.
祂和我們同在.
祂會幫我們勝過一切的試煉試探.
事實上新約.
耶穌也是這樣說.
你實踐祂給你的大使命.
你就會經歷到.
上帝對你的同在.
所以我相信在座我們.
絕大部分.
都已經是信主得救的弟兄姊妹.
那麼你們的屬靈的靈情.

$^{1041}$現在走到哪裡了?.
在曠野.
去了迦南.
還是已經安穩在上帝的裡面.
讓我們勇敢去回應.
上帝今天透過這段經文.
給我們的呼召.
我們一起祈禱.
天主上帝我們多謝你.
因為是你親自揀選.
呼召任命我們.
讓我們不單單成為你的兒女.
更加成為君尊的祭司.
但願我們都好像.
今天經文所提醒.
是你呼召我們.
是你親自揀選.
讓我們在教會裡.
可以親近你,事奉你.
你很想藉此讓我們認識你.
讓我們見識你的大能.
讓我們透過事奉來榮耀你.
來和你建立這份親密的關係.
上帝在過程裡我們會有膽怯.
因為我們很多時候.
沒有將目光對焦在你身上.
我們對焦在自己的成敗.
別人的目光.
或者這個世界的價值當中.
以致我們真的會有壓力.
我們會氣餒甚至膽怯驚惶.
求主你幫助我們.
讓我們首先作為你得救的兒女.
我們認定要靠主你過一個得勝的人生.
讓我們再次安穩.
在耶穌基督的理論裡.
有你成為我們的避難所,避風港.
讓我們經歷你的醫治,經歷你的安慰.
可以再重新得力.
認定你讓我們走的方向.

$^{1081}$讓我們勇敢踏出一步.
為主你作見證.
我們這樣禱告交託.
奉耶穌基督你的名來祈求.
阿們.
\newpage



\section{希伯來書 4:1-11}
\label{sec:QHvBwcg9HEw}
\textbf{2022年3月13日崇拜直播｜羅志雄牧師｜超越的基督:進入安息｜希伯來書4\_1-11}
\newline
\newline
連結: \href{https://youtube.com/watch?v=QHvBwcg9HEw}{\texttt{https://youtube.com/watch?v=QHvBwcg9HEw}} ~~~~ 語音日期: 2024-09-19
\newline
\newline
\hyperref[sec:QTbAXdDuZOw]{\small{< < < PREV SERMON < < <}}
~
\hyperref[sec:index_chronic]{\small{[返順時目]}}
~
\hyperref[sec:U7nTwzqoXV8]{\small{> > > NEXT SERMON > > >}}
\newline
\newline
希伯來書 4:1-11
\newline
\begin{longtable}{cl}
\hline
\hline
章節 & 經文 (和合本修訂版)\\
\hline
4:1 & \begin{tabularx}{0.7\textwidth}{X} 所以，既然進入他安息的應許依舊存在，我們就該存畏懼的心，免得我們 中間有人似乎沒有得到安息。 \end{tabularx} \\ \\ \relax
4:2 & \begin{tabularx}{0.7\textwidth}{X} 因為的確有福音傳給我們像傳給他們一樣；只是所聽見的道對他們無益，因為他們沒有以信心與所聽見的道配合。 \end{tabularx} \\ \\ \relax
4:3 & \begin{tabularx}{0.7\textwidth}{X} 但我們已經信的人進入安息，正如神所說：「我在怒中起誓：他們斷不可進入我的安息！」其實造物之工，從創世以來已經完成了。 \end{tabularx} \\ \\ \relax
4:4 & \begin{tabularx}{0.7\textwidth}{X} 論到第七日，有一處說：「到第七日，神就歇了他一切工作。」 \end{tabularx} \\ \\ \relax
4:5 & \begin{tabularx}{0.7\textwidth}{X} 又有一處說：「他們斷不可進入我的安息！」 \end{tabularx} \\ \\ \relax
4:6 & \begin{tabularx}{0.7\textwidth}{X} 既有這安息保留著讓一些人進入，那些先前聽見福音的人，因不信從而不得進去， \end{tabularx} \\ \\ \relax
4:7 & \begin{tabularx}{0.7\textwidth}{X} 所以神多年後藉著大衛的書，又定了一天—「今日」，如以上所引的說：「今日，你們若聽他的話，就不可硬著心。」 \end{tabularx} \\ \\ \relax
4:8 & \begin{tabularx}{0.7\textwidth}{X} 若是約書亞已使他們享了安息，後來神就不會再提別的日子了。 \end{tabularx} \\ \\ \relax
4:9 & \begin{tabularx}{0.7\textwidth}{X} 這樣看來，另有一安息日的安息為神的子民保留著。 \end{tabularx} \\ \\ \relax
4:10 & \begin{tabularx}{0.7\textwidth}{X} 因為那些進入安息的，也是歇了自己的工作，正如神歇了他的工作一樣。 \end{tabularx} \\ \\ \relax
4:11 & \begin{tabularx}{0.7\textwidth}{X} 所以，我們務必竭力進入那安息，免得有人學了不順從而跌倒了。 \end{tabularx} \\ \\
[1ex]
\hline
\hline
\end{longtable}
$^{1}$領子妹讓我們一同去聆聽一段聖經的說話.
在希伯來書四章一至十一節.
經文裡是這樣說.
我們既蒙留下有進入他安息的應許.
就當畏懼.
免得我們中間或有人似乎是趕不上了.
因為有福音傳給我們.
將傳給他們一樣.
只是所聽見的道與他們無益.
因為他們沒有信心與所聽見的道調和.
但我們已經相信的人.
得以進入那安息.
正如神所說.
我在路中起誓說.
他們斷不可進入我的安息.
其實造物之公.
從創世以來已經成全了.
論到第七天.
有一處說.
到第七天神就揭了他一切的公.
又有一處說.
他們斷不可進入我的安息.
既有必進安息的人.
那先前聽見福音的.
因為不順從不得進去.
所以過了多年.
就在大衛的書上又限定一日.
如以上所引述的說.
你們今日若聽他的話.
就不可硬著心.
若是約書也已叫他們享了安息.
後來神就不再提別的日子了.
這樣看來.
必令有一安息日的安息.
為神的子民傳流.
因為那進入安息的.
乃是揭了自己的公.
正如神揭了他的公一樣.
所以我們務必竭力進入那安息.
免得有人學那不順從的樣子跌倒了.

$^{41}$今天早上有羅茲雄牧師.
在我們當中宣講希伯來書的主題.
第二講.
講題是進入安息日.
(安息).
等子妹平安.
在我們開始之前.
我們一起有一個禱告.
讓我們有生命氣息.
可以聆聽你自己的道.
願上帝你打開我們屬靈的耳朵.
我們能夠聽到上帝你自己的教導.
亦都願講的.
都是經歷上帝你自己的洗滌.
心靈裡面能夠湧發出上帝你自己的恩典.
一齊的分享上帝你自己的說話.
我們獻上禱告.
奉主耶穌基督的名.
阿們.
上個主日.
陳牧師已經講到有關.
基督是神親自曉喻.
並且是超越天使和先知的信息.
作者講到教導的時候.
當時候要提醒信徒.
雖然他們面對很大的挑戰.
靈性低沉.
加上生活上面面對很多困難.
呼籲他們仍然持守所相信的道.
不要垂留失去.
又勸勉信徒.
不要走回頭路.
離開所信的道.
今天這個信息.
我們從另一個層面.
看看作者.
來勉勵我們一群信徒.
如何獲得神真正的安息.
中文聖經會譯為「安息」.
英文會用「rest」.

$^{81}$不知道大家有沒有想過.
休息是否等於不做任何事.
即是我們躺平.
舒舒服服.
這個時候面對疫情仍然嚴峻.
很多時候我們都會在家工作.
甚至任何時候.
都會上網處理一些工作.
又或者通電話.
這樣是否變得完全沒有休息呢?.
弟兄姊妹.
究竟甚麼是休息?.
做事和不做事.
是否當中不能夠有休息有安息呢?.
舊約有一個重要的人物.
摩西是神中心的僕人.
他可以與神面對面傾談.
在以色列裡.
其實他的地位也相當崇高.
不過舊約多方多次曉喻.
神的兒子耶穌.
也是同樣盡忠來服侍上帝.
在希伯來書的第三章.
第一至第六節.
作者是這樣講述.
基督身為神的兒子.
無論他在地位和事奉上.
都是比神的僕人摩西更加超越.
摩西雖然做了很多很有影響.
和有心軟的一些事情.
摩西就帶領了以色列人離開這個圍爐之家的埃及.
想帶他們脫離被牢獄多年的一個地方.
但最終他不能夠帶領這群百姓.
真正的失去他們的奴婦.
我們今天這個主題進入安息.
聖經裡面也有教導.
我們要安息.
特別是我們回看《創世記》的時候.
神在六天做工來創造天地.
到第七天他就休息.

$^{121}$細心了解一下.
原來他們要守安息日.
就是因為摩西領受了神所頒布的十誡.
告訴他們其中的誡命就是要守安息日.
什麼工作都不可以做.
這個就成為猶太人世世代代遵守安息日.
如果大家在疫情之前有去過巴勒斯坦的以色列這個地方.
我們去旅遊.
又或者參加一些聖地遊.
如果遇到以色列地在安息日的時候.
可能就會有一番的體會.
而希伯來書的作者.
現在在第四章同樣是講到安息的課題.
他是講到守安息.
或者守安息日有什麼不同呢?.
作者就是這一段經文來表達.
安息.
真正的安息是什麼.
我們怎樣得到安息.
神所應許給我們的安息又是什麼呢?.
我們再進入去看這一段聖經.
在巴勒斯第四章第一節至第三節的上半節.
我們既蒙留下有進入他安息的應許.
就當畏懼.
免得我們中間或有人似乎是趕不上了.
因為有福音傳給我們.
像傳給他們一樣.
只是所聽見的道與他們無益.
因為他們沒有信心.
與所聽見的道調和.
但我們已經相信的人.
得以進入安息.
在這裡第一節就已經講明.
我們這一群當時候的信徒.
是已經有一個安息應許.
在中文聖經裡面沒有譯為所以.
這個連接詞.
如果是這樣看的時候.
其實第一節就是連上去第三章十九節那裡.
這樣看來他們不能進入安息.

$^{161}$是因為不信的緣故.
由第三章第七節.
一直到第四章十三節.
這個大段落.
作者就以耶穌基督是超越摩西.
之後就帶出摩西所做的其中一個最重要的任務.
就是將以色列人帶來埃及.
但是他們究竟有沒有進入安息之地呢?.
這個也是申命記有提.
迦南就是一個安息的地.
所以在這裡作者是要給出警告和勉勵.
頂子妹短短的這三節經文.
其實提醒我們.
很重要的字眼就是「就當畏懼」.
希伯來書的作者提出這個畏懼的警告.
就是因為在以色列人他們經歷了很大.
和一個很震撼令他們失敗的一個結局.
以致他們不能夠進入安息.
就是在第十三章十九節所提到.
因為他們是被神拒絕.
不能夠進入神應許的迦南地.
一個可以安歇,長居久安的地方.
而這個記載.
相信在民數記裡面有很詳述的記載.
他們在迦底斯這個地方.
將近可以進入迦南地的時候.
摩西領受了神的吩咐.
派探子做先頭部隊.
當然他們回來的結果是非常好的.
並且都帶了很豐富的農產品來回報.
他說:實在是流瀨與物之地.
神是沒有說錯的.
我們可以進去.
但當中大部分的人都有異議,不同意.
他們說:沒錯,這是流瀨與物之地.
但是人高大,城牆堅固.
如果我們進去這個地方.
就是等於送死.
當時聖經是這樣記載神法路.
因為他們叛逆試探神.

$^{201}$於是他們就在曠野漂流四十年.
因為他們的不順從.
他們就倒閉在曠野.
這個重要的失敗經歷.
就成為以色列人.
整個舊約的先知詩裡面.
不斷地重提.
成為他們一個重要的記憶和警惕.
弟兄姊妹.
究竟不能夠進入神的安息是什麼原因呢?.
在第三節的時候.
神說:我在路中起誓.
他們斷不可進入我的安息.
如果我們可以稍為看回第三章的時候.
其實這就是詩篇95篇11節的引文.
在這裡.
第七節到十一節的時候.
其實對照詩篇同一段三至十一節的時候.
其實經文幾乎是一樣的.
只不過在詩篇第八篇.
95篇第八節的時候.
他就是這樣描述.
你們不可硬心.
像當日在米利巴.
就是在曠野的瑪薩.
原來瑪薩,米利巴.
這兩個字翻譯出來.
就是鬥爭和試探.
原來他們在整個曠野的旅程裡面.
不斷地與神抗爭.
甚至試探上帝.
所以希伯來書的作者.
就在這裡有一個結論.
縱然詩篇也有這樣提法.
他們就是一個背叛,得罪神的人.
他們不是一次.
不是半次.
而是屢屢試探神.
蔑視神.
於是他們就變得心硬.

$^{241}$最後這樣的演變下去.
就成為一個滾雪球的效應.
結果他們不相信上帝.
真的可以帶領以色列民.
進入迦南神的安息.
最終他們是自絕於安息地的門外.
又或者從另一個角度看.
一個正面的.
就是要進入應許的安息.
我們就是要有信心.
正如經文裡面提到的.
那些人沒有信心.
於是提醒希伯來書的讀者.
做一些正面的事情.
就是要有信心.
與所聽見的道調和.
丙姐妹.
或者我們在這裡停一停.
去想想這一三節裡面.
給我們有什麼應用.
經文裡面.
都會提到.
人很多時候未必生活在一個安恕的環境裡面.
就好像舊約以色列人一樣.
他們經歷沙漠曠野.
他們有沒有什麼.
來幫助他們行過呢?.
沒有食物,沒有水.
但上帝眷顧他們.
他們有從上帝而來的食物.
有從盤石而出的水.
甚至當他們遇上危難的時候.
很艱難.
甚至會受到敵人挑戰的時候.
他的情況會怎樣呢?.
一定會很危險.
怎樣可以得到休息呢?.
我相信同樣希伯來書的讀者.
又或者包括我們今天一眾弟兄姊妹.
其實我們會不會就是在一趟.

$^{281}$屬靈的曠野旅途上呢?.
我們能不能夠經歷到.
神賜給我們的安息.
如果是這樣的話.
我們就要信服.
與上帝的道來調和.
第二.
那些以色列人在曠野那裡.
不能夠進入安息.
因為神其實沒有拒絕他們.
而是他們的心眼.
惹怒上帝去試探神.
所以以色列人這一段黑骨銘心的歷史.
就成為當時候希伯來書的讀者.
一個反面的教材.
要時刻懼怕不要報他們的後塵.
從另一個角度來看.
說到關於應許.
這個也是由第四章希伯來書開始.
第一次出現的.
如果我們一直往後來看.
希伯來書的時候.
總共出現了有18次這麼多.
這個相信是希伯來書裡面的一個關鍵詞.
它是指出神的應許是有內容.
不是說上帝應許一點.
就這樣說應許而是有內容.
也有具體的行動.
甚至也描述了這位有應許的神.
他的承諾.
對相信他的人是有承諾.
希伯來書的作者.
就說出第一個很重要的應許.
就是安息的應許.
這個就是表明激勵信徒.
神賜福給聽道又行道.
屬於真正了解福音.
了解神的道的人.
說到應許.
有些聖經統計大概有八千個這麼多.

$^{321}$當時我一直在教會年輕的時候.
也聽到很多的牧師.
說應該有上二萬個應許.
我想一想.
會不會是成本的聖經大部分都是有神的應許呢?.
我們現在生活在一個非常的世代.
又有疫情.
世界的秩序又不穩定.
政治上,軍事上,經濟上.
甚至突如其來在歐洲又爆發了戰爭.
我們的心靈真的能不能夠休息呢?.
會不會情緒因著這些我們都喜樂不定呢?.
聖經原來在這裡勸勉我們.
祂會給我們安慰.
在這裡提醒我們.
我們要將視線轉移在神的應許上.
其實應許.
我相信每一位弟兄姊妹都很想知道.
究竟這個應許何時出現.
何時兌現呢?.
有時候我們去想的時候.
就會覺得很遙遠的距離.
但原來希伯來書的作者.
在往後一直講的時候.
都有一個放心.
當然我們今天去看希伯來書的時候.
我們就更加看到.
希伯來書裡面對應許的提醒.
這一段就是在希伯來書的十一章第一節.
他說:信就是所望知事的實底.
是未見知事的確具.
這個我相信.
就是希伯來書的作者對信徒的一個鼓勵.
甚至在整個初期教會的時候.
使徒彼得也有這樣的提醒.
他說:主所應許的尚未成就.
有人以為他是擔言.
其實不是擔言.
乃是寬容你們.
不願有一人沉淪.

$^{361}$乃願人人都得救.
原來應許在上帝的裡面.
何時成就.
我們每個人.
每個相信主的人.
我們都知道.
是有上帝的時間.
上帝是出於愛.
他的應許.
就是要更多的人可以進入上帝的道裡面.
耶穌也有他的應許.
耶穌也有很多的應許.
但是最擇心的應許就是說.
復活在我.
生命也在我.
信我的人雖然死了.
也必復活.
凡活著信我的人.
必永遠不死.
於是主耶穌就說.
你信這話嗎?.
另外一位的使徒保羅.
他也在他的書信裡.
常常提到關於應許.
他說:親愛的弟兄.
我們既有這等應許.
就當潔淨自己.
除去身體,靈魂一切的污穢.
敬畏神.
得以成聖.
我相信希伯來書的作者.
就是看到.
我們要信所聽見的道.
來調和.
我相信就是.
要在我們的生命裡面.
來應用.
接下來我們就看看.
第三節下部分.
到第十節.

$^{401}$這裡希伯來書的作者.
就是重申安息日應許的定義.
對我們明白.
真正安息日.
其實是十分重要.
我們看看這段經文.
正如神所說.
我在路中起誓說.
他們斷不可進入我的安息.
其實眾物之宮.
從創世以來已經成全了.
論到第七日.
有一處說.
到第七日神就揭了他一切的宮.
又有一處說.
他們斷不可進入我的安息.
既有必進安息的人.
那先前聽見福音的.
因為不順從.
不得進去.
所以過了多年.
就在大衛的書上.
又限定一日.
如以上所引的說.
你們今日若聽他的話.
就不可硬著心.
若是若書也叫他們享了安息.
後來神就不再提別的日子了.
這樣看來.
必有另一安息日的安息.
為神的子民傳流.
這段經文很長.
當中亦有一些引文.
我們很想在這裡.
作一個簡單的分析.
來講解兩種的安息.
傳息的框架.
以致我們能夠有一個.
比較清晰的思路.
去了解希伯來書的作者.

$^{441}$為何要有這段引文.
我相信.
如果再看的時候.
由第三節下.
到第五節.
其實這就是一個舊約引文.
作為一個基礎.
去藏息.
關於安息.
和安息日的安息的看法.
希伯來書第四章和學本.
翻譯為安息.
在英文裡是rest.
在這裡第四章出現了.
九次這麼多.
及apostas.
這是一個關鍵字.
安息在希臘文裡.
有兩個意思.
第一個是停止工作.
或一種活動的狀態.
所以是一種安歇.
歇了的狀況出現.
正如在經文第四節所說.
到第七日.
神歇了他一切的工作.
其實都是用了希臘文的字.
歇了.
另外一個意思.
具體是表達到.
一個安息之所.
即是一個地方.
人可以在這裡休息.
和居住的地方.
當然這個安息.
涵蓋了這兩種意思.
所以希伯來書作者.
引用詩篇95篇.
很明顯地告訴我們.
這是一個地域性.

$^{481}$甚至整段的技術.
其實是一個歷史.
但歷史並不是沒有意義.
正如以色列人常常.
都將這件事掛在口邊.
甚至經文裡有提到.
以色列人出埃及幾百年後.
去到大衛的時候.
都引述這段經文.
詩篇95篇第七至第十一節.
我想知道這是一個.
非常重要對猶太人.
關於安息的觀念.
他們亦常常盼望.
到最後安息的時候.
他們能夠確實地.
真正得到神所允許.
為他們所預備的安息之所.
經文第九節.
「這樣看來,必另有一安息日的安息,為神的子民傳流.」.
安息日的安息.
希臘文Sabbathesimus.
原來只在這裡出現一次.
有些解經家說.
希伯來書的作者.
刻意地創造了這個字.
這個字的意思是什麼呢?.
是守安息,或者慶祝安息.
它不單單指一個安息日.
而是指在安息日裡.
關於有關安息,節慶的一種景況.
是一種守安息,慶祝安息.
圍繞著安息日去慶祝的時候.
表達喜樂的心境和一種狀態.
在那天的時候.
所有的信徒一起讚美,稱頌上帝.
安息日的安息.
就是希伯來書的作者對舊約的傳息.
他將安息.
猶太人的觀念安息.

$^{521}$與安息日的觀念連接起來.
所以希伯來書的作者呼籲讀者進入這種安息.
或者我們再仔細看看第三節,下半節到第五節.
其實這是一個很重要的.
作者要解釋.
安息和安息日的安息.
這個重要的引文框架.
當然了.
希伯來書的作者文學造詣很高深.
所以當他做教導的時候.
他引了這篇詩,九十五篇來做他的釋經.
這個技巧非常高超.
但重要的是什麼呢?.
重要的是讓人明白神的真理.
神真正的安息.
當在這個框架裡面.
其實另外也是提醒我們.
我們需要有一個謙卑受教的心.
以致我們真的有這個歷史給我們的教訓.
我們才能夠體會到.
上帝所應許的安息是什麼.
所以一直下去,經過第七節到第十節.
甚至到第九第十節的時候.
其實這就是作者對希伯來讀者一個最高潮的解釋.
關於人找到安息日的安息.
然後我們應該怎樣做呢?.
當我們清楚安息.
引入安息日的安息.
一個重要的脈絡的時候.
我們就會清楚明白.
在這裡第八節.
聖經作者這樣引出了另一個事實.
「若是約書亞已經叫他們安息了.
後來神就不再提別的日子」.
我相信希伯來書的讀者.
除了認識摩西之外.
約書亞他們不會不認識.
約書亞就是承接著摩西帶領以色列人進入迦南.
去到神所應許的神的安息.
沒錯,如果我們這樣回顧這段以色列人的歷史的時候.

$^{561}$他們確實是進入了神所說的迦南地.
是神的安息.
但要問,弟兄姊妹.
他們究竟有沒有得到真正的安息呢?.
大家都知道約書亞所帶領的以色列人進入迦南地.
是第二代的以色列人.
即是上一輩在摩西帶領的時候.
他們在曠野已經倒閉了.
當希伯來書的作者引用第八節的時候.
就不是在詩篇裡面所說的.
但這只不過是地域上的安息.
而這個狀態的安息.
在約書亞之後.
可以說是短暫.
又可以說是反反覆覆.
最終整個以色列人.
在迦南神所賜給他們的神安息的地方.
他們是否真正得到安息呢?.
這就是經文裡面.
作者很想提出來.
讓當時的猶太基督徒去明白.
甚至說我們今天都需要了解去明白.
另外一個更加有趣的.
就是我們要留意一下.
作者提出提及約書亞這個名字.
約書亞是希伯來文.
如果翻譯為希臘文.
其實跟耶穌這個希臘文是一模一樣.
我相信作者在這裡.
既然約書亞做不到的事.
他雖然帶領百姓進入迦南地.
這個地域上的安息.
真正上帝所賜的安息.
他們還沒有得到.
那誰可以承傳這樣去做呢?.
希伯來書的作者.
就是用這種表述.
他說:沒錯.
在詩篇95篇裡所說的.
百姓倒閉在曠野.

$^{601}$約書亞時期.
以色列人真真正正進入迦南.
這個很實際的地方.
但他們依然得不到安息.
唯有耶穌基督.
這一位超越摩西的耶穌.
這一位基督.
才能夠.
亦都是唯一的能夠帶我們進入神的安息.
另一方面在這一段亦提到.
當希伯來書的作者引用詩篇95篇11節時.
我在路中起誓說.
他們斷不可進入我的安息.
其實我們看回第三章的時候.
同樣都在此出現過.
本身在這一段第四章亦都出現了兩次.
其實是否暗示上帝這一位的神.
不單單是要對曠野世代的百姓定罪.
還是說信徒我們今日都要定我們的罪呢?.
我相信這不是作者的意思.
他就是要提醒我們.
無論是當代的基督徒.
跟隨耶穌的人.
又或者我們之後很多世代.
今日的我們.
我們都要來找到福音.
找到所聽見的道.
神的話.
很真實地來跟道調和.
然後才能夠有一個真正的安息之所.
在這裡.
經文裡面說.
神稱這就是我的安息.
這就是神自己所享受的安息.
亦都是同時允許我們.
凡是跟隨他.
相信他.
行他的道.
這一班的信徒可以同享.
或者在這裡我們有一個簡單的小結.

$^{641}$經文講到安息的看法.
提到迦南地的安息.
這是一個歷史性的.
是一個地域性的.
亦都是一個暫時的安息.
第二個就是安息日的安息.
弟兄姊妹.
不知道你現在有沒有安息呢?.
我指的安息.
就是你連迦南地這個安息.
你有沒有得到呢?.
或者你自你相信耶穌基督.
接受耶穌基督成為你生命的救主的時候.
你聽過耶穌的道.
你有沒有覺得.
此時此地.
此刻.
你有迦南地式的這種安息.
縱然短暫.
是一個暫時性.
但是.
在這裡思想一下.
有沒有這種的安息呢?.
假如你沒有的時候.
我相信上帝的說話.
神的道.
就會引導你.
可以進入去.
最低限度.
在這個屬靈的曠野路裡面.
你去尋索.
可能會跌跌撞撞.
可能會好像昔日以色列人那樣.
有些是背叛.
有些是試探.
但是上帝的愛.
會讓你.
只要你真的謙卑下來.
不會好像所有的以色列人一樣.
只有迦納,約書亞.

$^{681}$他們就本著這種相信的時候.
最低限度.
他們都能夠去到迦南地安息之所.
或者弟兄姊妹.
你在哪一個狀態.
會不會又很期盼著.
安息日的安息呢?.
這一種是真正.
一個屬於上帝.
他領受了上帝的恩典.
而他很盼望.
正如在希伯來書第十一章裡面所說的盼望.
雖然現在仍然感覺到是勞苦.
有很多的愁煩.
又或者面對現在生活上的種種困難.
以致到.
看聖經.
靈性上和上帝有很大的掙扎.
有一種張力.
我何時才能脫離這種勞苦愁煩.
去到上帝給我們安息日的安息呢?.
我相信這一段經文.
很好的來激勵我們.
提醒我們.
對於猶太的基督徒.
在當時希伯來書所寫的時候.
我相信這是一個.
嶄新和具有顛覆性的.
一種的詮釋.
原來由安息日安息.
引入安息日的安息.
這一種的視覺.
當我們繼續看下去的時候.
其實經文就是這樣來鼓勵我們.
我們怎樣能夠進入.
鼓勵我們進入安息.
這就是第十節到第十一節.
經文是這樣記述.
因為那進入安息的乃是揭了自己的躬.
正如神揭了他的躬一樣.

$^{721}$所以我們務必竭力進入那安息.
免得有人學那不順從的樣子跌倒了.
在第十節裡面.
如果仔細詮釋這段聖經.
其實有不同的詮釋的說法.
但我簡單的來說的時候.
其實作者要我們.
當分辨到這兩種安息.
對這兩種安息的新的理解.
新的視覺理解的時候.
其實作者就是在鼓勵我們.
要踏實進入安息日的安息.
作者在這裡鼓勵.
鼓勵信徒要竭力進入這種安息.
日的安息.
所以我們務必竭力進入那安息.
免得有人學那不順從的樣子跌倒了.
丙姐妹,竭力其實就是告訴我們.
我們要盡一切的努力,一切的能力.
我們手上所有的一些資源.
關於尋求上帝的心意.
去明白上帝的道.
這一種方式.
否則的話,我們只有我們人的性情.
只有人的那種想法.
我們就成為那些不信的人一樣.
他們終有一天就倒閉在曠野裡.
不要說要進入安息.
就算是進入到迦南地.
這個都不是那個的安息.
所以第十一節.
雖然我們看到都是用一般希臘文的字.
但是當由傳息了第九節之後.
這個安息就是那安息.
就是安息日的安息.
丙姐妹,究竟我們要怎樣做呢?.
希伯來書的作者沒有就這樣停在這裡.
他就說你們自己想方法吧.
你們想用什麼方法就盡你們所能去做吧.
但又不是.

$^{761}$希伯來書繼續的告訴我們.
唯一一個實踐的方法.
和這個要持續的一個進程.
我們要有一個迫切的態度去追尋.
這個就是什麼呢?.
這個就是神的道.
這個就是我們整個信仰重要的核心.
一個重要的價值.
我們看這一段經文十二到十三節.
神的道是活潑的,是有功效的.
比一切兩刃的劍更快.
甚至魂與靈,骨折與骨碎.
都能刺入,否開.
連心中的思念和主義都能辨明.
並且被造的沒有一樣在他面前不顯然的.
原來萬物在那與我們有關係的主眼前.
都是赤露闖開的.
我再讀還有一些其他的日本的譯文.
有新漢語譯本和環球聖經譯本.
他就說都是赤露闖開的.
原來還有一句.
我們都必須向他較仗.
原來作者當他寫希臘文這一段的時候.
十二十三節.
是只有一句的句子.
開頭的時候他就用道.
到結尾,我讀了我們都必須向他較仗.
較仗這個字,其實原文也是道.
就是說這一段經文.
十二十三節,一頭一尾.
用道來彼此呼應.
我相信這個就是希伯來書作者所寫的希伯來書.
上次陳牧師已經提過是一個很優美.
用他很出色的希臘文來表達.
在其他的章節同樣可以看到.
無論是怎樣,這一節的經文其實總結了一個方法.
我們要進入去安息日的安息.
我們就只有神的道.
神的道是唯一.
但是作者在這裡運用了一個意象.

$^{801}$一個意象來講.
神的道是活潑的.
是有功效的.
甚至是說到是劍.
一把兩刃的利劍.
當我一路看這段經文的時候.
我腦海裡也浮現出有關於劍.
不知道大家有沒有想到.
保羅在這份所書裡同樣提出關於劍.
如果有記得的話.
我們要怎樣在世間.
我們的身體怎樣能夠抵擋魔鬼呢?.
保羅提醒原來我們要穿戴神所賜的全副軍裝.
其中一樣就是聖靈的寶劍.
聖靈的寶劍是什麼呢?.
保羅就是說神的道.
沒錯.
這兩位在新約書卷裡的作者.
同樣都以劍來寓意神的道.
這筆所書裡保羅提的全副軍裝.
我們拿著劍對抗屬靈的惡魔.
但如果我們再想一想.
那把劍其實也可以做人體解剖.
人體解剖我們會看看.
如果我們有病.
可能做一些深層的檢查.
我們普通的醫生斷症.
可能都未必看清楚.
甚至可能會切一些細胞來檢查.
看看有什麼最好的治療方法.
這個就是一些手術刀.
說起人體解剖的時候.
我有一個經歷.
有一個弟兄.
他安息禮拜.
原來他的安息禮拜.
就是在一間醫學院裡做.
原來這位弟兄在他在世人生的時候.
他承諾在他離開這個世界的時候.
他整個身體整個軀體.

$^{841}$要來做一些有益的事.
要來給一班醫生解剖.
當我去到那間醫學院門外.
去到這個部門.
他有兩句話.
今日吾軀無言教.
他朝良醫有心人.
原來當我們去幫助我們去剖析的時候.
可能只是解剖身體.
但原來這一段的經文.
如果我們再深層一點去看的時候.
我們從另一個角度.
因為經文說到.
骨折與骨髓.
靈與魂.
神的道.
一把鋒利無比又很快的一把刀.
原來是刺入人內心最深層的深處裡面.
相信這個就是屬靈的解剖.
兄弟姐妹.
我自己也去想一下.
我有沒有做過這樣的靈性手術.
讓上帝來解剖我自己呢.
以致神的說話.
告訴我們.
連我們的心思意念.
我們的主義.
神也一一知道.
一一很清楚.
就算我們想了沒有說出來.
上帝也知道.
甚至我們要想下一步的時候.
我們也沒有想出來.
上帝也知道.
神也知道.
神的道就是這樣去了解.
去理解我們每一個人.
將我們心層最不會讓人知道的地方.
最污穢,最黑暗,最邪惡的.
上帝也可以去剖開.

$^{881}$甚至連我們自己也沒有想過.
原來我們自己在上帝的面前.
也是污穢不堪.
甚至比頂撞上帝更嚴重.
我想這就是這段經文.
來提醒我們.
神的道.
生命之道.
是大有功效的.
所以.
兄弟姐妹.
當我們看到這段經文.
我們在探索的時候.
究竟在我們這個曠野屬靈之旅裡.
我們又願不願意.
被上帝用他的道.
鋒利的這把刀.
來刺透我們呢?.
這個比保羅所說的.
聖靈的利劍,神的道.
更加的深入.
那裡所提及的.
我相信就是一個對外在的環境.
我們拿著這把劍.
但我們今時今日是不會穿著.
上帝給我們這副軍裝.
如果這樣的話.
走出來一來會怪怪的.
第二是怎樣呢?.
我們怎樣去走呢?.
原來這是一個隱喻.
就是說神的道.
是幫助我們在日常生活上.
在日常生活裡.
我們藉著神的道.
給我們有從神而來的智慧.
我們做一些判斷.
做一些事情.
我們都需要按著神的教導.
有些更宏大的.

$^{921}$我們更加需要神的道幫助我們.
就算是一些很簡單.
我們認為很簡單,很瑣碎.
我們日常所經歷的,所要的.
其實這些都需要上帝的道.
來調和我們的生命.
幫助我們.
至於這個靈性上的解剖.
其實就是要對付我們最深層.
有一些驕傲.
其實我們是不自覺的.
有一些惰性.
其實我們都是不自覺的.
弟兄姊妹.
我也有聽過.
弟兄姊妹會分享的.
也會埋怨.
特別是回教會很久的弟兄姊妹.
她說:我越回教會越久.
回了幾十年.
總是心裡沒有那份安穩.
沒有那份喜樂.
沒有那種踏實.
沒有那種平安.
甚至說安息.
有沒有問一下是什麼原因呢?.
又或者有些弟兄姊妹.
我不管了.
除了崇拜之外.
我只是回來聽道.
不過我真的有聽.
但一邊聽就頭暈.
但究竟她的道.
聽到的道.
有沒有在她的生命裡融和出來呢?.
又或者有一些弟兄姊妹.
說我回來教會.
我只是事奉.
我會不停做.
但是當她宦官不斷做的時候.

$^{961}$有些弟兄姊妹就離開了.
甚至在事奉的過程裡.
她們都會產生一些問題.
可能因為小事.
又會起了一些爭執.
甚至又聽過.
她說我會聽道.
不過不是為我自己聽.
上帝你的道.
不要像一把手術刀那樣.
切入我這裡.
我是為誰聽的.
因為我聽到她.
她又得罪人又得罪我.
這篇道不是跟我說的.
就算跟我說了.
我也要叫她悔改.
向我道歉.
這是不是我們聽道的態度呢?.
又或者.
我們聽到很多的見證.
有些弟兄姊妹.
當他們信了之後.
原來她分享的時候.
她說我知道回來教會是很好的.
但其實我信的那一套更好.
所以她回來的時候是怎樣呢?.
她就懷著一種挑戰的心態.
牧師傳道人說的道.
我就要挑戰他.
甚至好像去踢館那樣.
弟兄姊妹.
不知道你有沒有經歷過這些呢?.
如果是的話.
有這樣的情況的時候.
很想和大家分享.
其實希伯來書的作者.
在這兩節的經文裡.
其實就是要提醒我們.
他活潑的道.

$^{1001}$有功效的道.
就是一把利劍.
去刺開我們.
是一個靈性的一種解剖.
你願意嗎?.
到最後有一個總結.
在今天.
這篇希伯來書第四章.
到第十三節.
作者很想我們能夠去明白.
就是我們要進入去安息的時候.
我們要對神的道認真.
我們願不願意讓神的道.
勝在我們的生命呢?.
我們是否以信心信服.
當然.
我們一個人做.
是好的.
但是我們能不能夠一個人.
奔跑到這條屬靈的道路呢?.
其實我們是有需要弟兄姊妹.
一起彼此相勸.
進入去安息日的安息.
這條道路上.
迦南地安息這個模式.
是不是弟兄姊妹今天.
理所向往的呢?.
抑或我曾經有.
只不過是一個歷史的記憶.
又或者只是短暫.
在某一個地方.
我在A教會.
來到B教會或C教會呢?.
還是說.
無論你去到哪一個地方.
哪一間教會.
你都會領受到最低限度.
是迦南地這種模式的安息.
又或者你不能夠.
離開這個第一種的安息的模式.

$^{1041}$我相信你很難進到去安息日的安息.
因為你沒有將神的道.
載在你的心裡.
但希伯來書的作者.
縱然有警告.
但他亦有勉勵.
有很豐富的提醒.
正如耶穌基督.
都是這樣來邀請我們.
耶穌其實早就已經邀請我們.
「凡勞苦擔重擔的人.
可以到我這裡來.
我就使你得安息.
我心裡柔和謙卑.
你們當乎我的額.
學我的樣式.
這樣你們心裡就必得享安息」.
耶穌就是我們每一個.
在地上避風港灣.
你有沒有將你生命的鬧.
放在這個港灣.
耶穌的裡面呢?.
以致你的人生不會失去方向.
不會隨流失去.
任由這個時代的洪流.
來淹沒我們.
我們一起來禱告.
讓你的道進入我們的生命裡面.
因為你的道就是生命之道.
有活潑,有功效.
來將我們靈命裡面.
最黑暗,最敗壞的.
你都會來看清.
讓我們看清我們的本相.
以致我們需要的.
就是上帝的道來潔淨.
幫助我們回到上帝.
要我們去的安息日的安息.
獻上禱告奉主耶穌基督的名.
阿們.

\newpage



\section{撒母耳記下 12:15-23}
\label{sec:U7nTwzqoXV8}
\textbf{2024年9月22日崇拜直播｜陳明泉牧師｜正視你的情緒－為悲傷劃界線｜撒母耳記下12\_15-23}
\newline
\newline
連結: \href{https://youtube.com/watch?v=U7nTwzqoXV8}{\texttt{https://youtube.com/watch?v=U7nTwzqoXV8}} ~~~~ 語音日期: 2024-09-23
\newline
\newline
\hyperref[sec:QHvBwcg9HEw]{\small{< < < PREV SERMON < < <}}
~
\hyperref[sec:index_chronic]{\small{[返順時目]}}
~
\hyperref[sec:JCVpdjhfj4w]{\small{> > > NEXT SERMON > > >}}
\newline
\newline
撒母耳記下 12:15-23
\newline
\begin{longtable}{cl}
\hline
\hline
章節 & 經文 (和合本修訂版)\\
\hline
12:15 & \begin{tabularx}{0.7\textwidth}{X} 拿單就回家去了。耶和華擊打烏利亞的妻子為大衛生的孩子，他就得了重病。 \end{tabularx} \\ \\ \relax
12:16 & \begin{tabularx}{0.7\textwidth}{X} 大衛為這孩子懇求神。大衛刻苦禁食，到裡面去，躺在地上過夜。 \end{tabularx} \\ \\ \relax
12:17 & \begin{tabularx}{0.7\textwidth}{X} 他家中的老臣來到他旁邊，要把他從地上扶起來，他卻不肯，也不同他們吃飯。 \end{tabularx} \\ \\ \relax
12:18 & \begin{tabularx}{0.7\textwidth}{X} 到第七日，孩子死了。大衛的臣僕不敢告訴他孩子死了，因他們說：「看哪，孩子還活著的時候，我們勸他，他尚且不聽我們的話，我們怎麼能告訴他孩子死了，讓他做出不好的事呢？」 \end{tabularx} \\ \\ \relax
12:19 & \begin{tabularx}{0.7\textwidth}{X} 大衛見臣僕彼此低聲說話，就知道孩子死了。他問臣僕說：「孩子死了嗎？」他們說：「死了。」 \end{tabularx} \\ \\ \relax
12:20 & \begin{tabularx}{0.7\textwidth}{X} 大衛就從地上起來，沐浴，抹膏，換了衣服，進耶和華的殿敬拜。然後他回宮，吩咐人為他擺飯，他就吃了。 \end{tabularx} \\ \\ \relax
12:21 & \begin{tabularx}{0.7\textwidth}{X} 臣僕對他說：「你所做的是甚麼事呢？孩子活著的時候，你為他禁食哭泣；孩子死了，你卻起來吃飯。」 \end{tabularx} \\ \\ \relax
12:22 & \begin{tabularx}{0.7\textwidth}{X} 大衛說：「孩子還活著，我禁食哭泣，因為我想，或許耶和華憐憫我，會讓孩子活下來。 \end{tabularx} \\ \\ \relax
12:23 & \begin{tabularx}{0.7\textwidth}{X} 現在孩子死了，我何必禁食呢？我能使他回來嗎？我必往他那裡去，他卻不能回到我這裡來。」 \end{tabularx} \\ \\
[1ex]
\hline
\hline
\end{longtable}
$^{1}$也是中讀經訓的時間.
今天所中讀的經文.
是薩姆爾記下的十五節至二十三節.
經文記載在程序表的內頁.
也可以翻開聖經.
唐用的聖經第三百九十二頁.
經文記載了大衛和拔士巴.
所得的兒子病逝的一段經文.
「那丹就回家去了.
耶和華擊打烏利亞妻.
及大衛所生的孩子.
使他得病重病.
所以大衛為這孩子懇求.
神而且禁食進入內室.
中夜躺在地上.
他家中的老臣來到他旁邊.
要把他從地上扶起來.
他卻不肯起來.
也不同他們吃飯.
到第七日孩子死了.
大衛的神僕不敢告訴他孩子死了.
因他們說孩子還活著的時候.
我們勸他.
他尚且不肯聽我們的話.
若告訴他孩子死了.
豈不更加憂傷?.
大衛見神僕彼此低聲說話.
就知道孩子死了.
問神僕說:孩子死了嗎?.
他們說:死了.
大衛就從地上起來.
沐浴,抹膏,換了衣裳.
進耶和華的殿敬拜.
然後回宮.
吩咐人擺飯.
他便吃了.
神僕問他說:.
你所行的是什麼意思?.
孩子活著的時候.
你禁食哭泣.

$^{41}$孩子死了.
你倒起來吃飯?.
大衛說:孩子還活著.
我禁食哭泣.
因為我想.
或者耶和華連恤我.
使孩子不死也未可知.
孩子死了.
我何必禁食?.
我豈能使他返回呢?.
我必往他那裡去.
他卻不能回我這裡來.
然神賜福常讀句話語.
並且進行的人.
弟兄姊妹平安.
我們今天繼續.
正視情緒系列的一系列講道.
很多弟兄姊妹曾經問過.
牧師日你講這個.
或者我們整個同工團隊.
他所講的正視情緒系列.
用意是什麼呢?.
其實情緒是我們生命的一部分.
我們每一個人生活在地上.
其實我們或多或少都會經歷一些事情.
牽動了我們的情緒.
而這個系列的目的就是.
藉著情緒成為了我們屬靈生命裡.
一個檢視我們自己處於什麼狀態.
又或者在我們的情緒背後.
可能我們有一些價值觀.
有一些對上帝的看法.
甚或者可能有一些情緒.
成為了我們信仰裡的樽頸位.
如果當我們去正視.
我們生命裡出現的情緒.
其實某程度上都會幫助我們個人.
生命裡有成長.
所以我們由聖經裡.
看不同的人物.

$^{81}$看他們如何面對自己.
不論是喜怒哀驚.
或者有不同的情緒.
我們很想進入他們的內心世界.
看他們走出來藉著上帝.
和上帝祈禱也好.
或者生命裡有一些表現.
可不可以為他們在生命這種處境中.
走出來.
今天我們看的大衛.
剛才讀的這段經文.
某程度上都是大衛生命裡的一個轉捩點.
之前的轉捩點就是大衛由一個牧羊仔.
變成了一國之君.
這段經文就說到.
當他已經掃羅追殺他的王已經死了.
他已經進入一個權力穩固的地步.
但是卻在他生命裡.
犯上了上帝所憎惡的罪行.
以致到.
接下來在撒母耳記下.
由第十一章開始一直到最後.
其實都是滿佈著大衛的眼淚和大衛.
生命裡有很多的艱難.
今天我們就從這個轉捩點來看.
大衛如何面對生命裡有一些不光彩的經歷.
我們更加在他身上.
大衛其實是很複雜的讀他.
有時候會成為我們的榜樣.
有時候會成為我們的借鑒.
但有時候也是從他身上.
我們看到上帝如何去琢磨一個人的生命.
我們今天就從一個上帝如何琢磨一個人.
或者如何回應上帝的角度.
來讀這段經文.
我們一起祈禱求聖靈幫助我們.
親愛的天父,求你在我們當中帶領我們.
讓我們能夠借著你的話語.
讓我們生命得著亮光和激勵.
願你在此時此刻.

$^{121}$您臨在我們眾弟姐妹心靈當中.
幫助孫說的說得清楚.
幫助聆聽的每一個人.
都能夠聽得明白.
我們俯伏在你面前.
恭敬祈求你的話語臨在.
獻上禱告奉扣基督耶穌得勝的聖名.
Amen.
在這裡我們先有一些背景資料.
剛才提到大衛面對著他的兒子.
就是從烏尼亞的妻.
給大衛所生的孩子.
上帝令他得到重病.
我們首先看看背景.
究竟這個孩子是怎樣出生的.
我們看回前一章.
撒母耳記下第十一章第二至第五節.
在這裡提到大衛的權力已經穩定.
在第二節提到.
有一天太陽平西.
大衛從床上起來.
在皇宮平頂遊行的時候.
見到有一個婦人沐浴.
容貌甚美.
大衛就警察打聽那個女人是誰.
有人說她是以年的女子.
黑人烏尼亞的妻拔士巴.
大衛就警察去將婦人接來.
那時候他的月經才得潔淨.
他來了大衛與他同房.
他就回家去了.
第五節提到.
於是他懷了孕.
打發人去告訴大衛說.
我懷了孕.
幾節經文已經交代了.
其實大衛就跟他的部屬.
他的下屬.
烏尼亞的妻子.
就有一些不倫的關係.

$^{161}$這個不倫關係.
有後續的.
就是在說這一次的不倫關係.
就令到拔士巴懷了新妓.
拔士巴就將這件事告訴了大衛.
都很肯定這個懷孕的孩子.
一定不是黑人烏尼亞的.
為什麼呢?.
然後經文繼續說.
他說因為烏尼亞上了戰場打仗.
所以他基本上沒有回過家.
老公出了去打仗.
無端端太太懷了孕.
就一定是有姦情的.
於是大衛第一件事要做的就是.
他很想掩飾這件事.
於是就由戰場裡面徵召他的部屬.
烏尼亞就回來.
跟他太太團聚.
想藉這個機會.
可以掩飾了他這種姦情.
誰知道這個烏尼亞就忠心耿耿.
在前線打仗的時候.
大衛叫他回來.
烏尼亞死都不進去.
怎麼都不回家.
他說的就是說.
當我的兄弟在戰場裡面打仗的時候.
我怎麼可以回家享受家庭的溫馨.
我作為戰士.
我不是這樣打仗的.
忠心耿耿.
當大衛聽到這個忠心耿耿的戰士.
不肯按照他的想法.
回家的時候.
於是在第十一章十四至十五節.
大衛就做了第二個舉動.
次日早晨大衛寫信給約翰.
這是大衛某程度上的心腹.
交烏尼亞隨手帶去.

$^{201}$信內寫著說.
要派烏尼亞前進.
到陣勢極險之處.
他們便退後.
使他被殺.
簡單來說.
他就說.
既然他不肯回家.
於是大衛就叫約翰和一隊新的軍隊.
就上到前線.
當到上到前線的時候.
就叫烏尼亞去衝.
接著他們一幫約翰大衛其他軍隊就送.
就退後.
烏尼亞一個怎麼抵擋其他的敵軍呢.
於是就用敵軍借刀殺人.
借阿滿人的手.
將這個忠心耿耿的部屬.
在戰場上殺害了他.
當烏尼亞戰死沙場的時候.
大衛還沒做完他要做的事.
十一章第二十六至二十七節.
烏尼亞的妻婦聽見丈夫烏尼亞死了.
就為他哀哭.
哀哭的日子過了.
大衛差人將他接到自己的宮裡.
他就作了大衛的妻.
給大衛生了一個兒子.
但大衛所行的這些事.
耶和華甚不喜悅.
簡單來說.
十一章二十六至二十七節就說到.
當烏尼亞死了.
大衛就借這個機會.
就用回舊時的猶太人的文化.
就是禮遇.
禮遇孤兒寡婦.
其實只是寡婦.
沒有孤兒.
其他人不知道整件事.

$^{241}$接了妻子去到皇宮.
禮遇她.
在其他人眼中這是一個很偉大的舉動.
大衛真是愛錫下屬了.
但實際上只不過是藉著這件事.
淫淫妻子殺人丈夫.
於是這件事.
以為只有他自己知道.
整件事好像天衣無縫.
不過在經文裡記載著.
洞悉萬事的上帝.
耶和華上帝看到整件事的發生.
所以二十七節裡面說到.
耶和華甚不喜悅.
上帝極之不喜悅.
大衛做了一件這樣犯罪的事.
接著當大衛想著這件事可以瞞天過海.
想著這件事沒有人知道的時候.
上帝就差遣拿丹仙芝.
到了第十二章裡面.
藉著拿丹仙芝來宣告大衛的罪行.
大衛所做的一切事.
上帝都放在眼內.
到了第十二章第九至十節.
拿丹仙芝很嚴厲.
宣告耶和華的命令.
接著第九節說.
你為什麼藐視耶和華的命令.
行他眼中看為惡的事呢?.
你借雁門人的刀殺赫人烏尼亞.
又取了他的妻為妻.
你記未笑我.
取了赫人烏尼亞的妻為妻.
所以刀劍必永不離開你的家.
所以整個段落你見到了.
這是說大衛已經解除了所有問題之後.
在安逸的狀況之下.
接著犯下一個離棄耶和華命令的罪.
淫人妻子殺人丈夫.
而且用一個瞞天過海.

$^{281}$用一個掩飾的方式.
來遮蓋自己的罪行.
更加要把光環放在自己身上.
令所有人覺得大衛真偉大.
真是一個很愛錫下屬的部下的王.
真是太好了.
不過暗地裡.
生命裡上帝神不喜悅.
因為所有最卑污的事.
發生在他身上.
當上帝指出大衛的問題.
指出大衛的那種不是之後.
上帝不是就這樣就收科.
上帝就開始在第十二章裡.
就要施行他對大衛的審判.
在大衛被上帝審判的內容裡有幾樣東西.
第一,上帝說你暗中做的這些事.
將來明明要報應在你身上.
你拿了別人的妻子.
成為你自己的妻子.
將來你的兒女都會.
你種什麼果.
你種什麼因.
後來就會出現什麼果子.
種什麼你就收什麼.
還有,因為你製造了一個殺人的罪.
所以刀劍永遠都不能離開你整個家庭.
這個判刑都很深的,是不是?.
不過,喃喃先知再繼續說.
上帝已經原諒了他.
死罪可免,活罪難饒.
他眼前的生活.
他接下來面對的就是.
上帝施行公義的審判.
要在這個家庭裡.
他萬不以有罪為無罪.
上帝都會報應在這個家庭當中.
不過,上帝最嚴重的離棄他.
這件事已經赦免了他.
他的生命跟之前的王素羅很不同.

$^{321}$大衛可以繼續在皇位裡面.
又或者在他的管治當中.
繼續領受上帝的恩典.
不過,昔日他在地上所犯的罪.
上帝會按照公義來審判他.
好了,你知道這個背景之後.
接著這個孩子.
剛才說到跟拔士巴所生的孩子.
上帝就擊打他.
上帝就要讓他得到重病.
當他得到重病的時候.
大衛如何去面對這件事呢?.
十二章第十六節至第十七節.
所以大衛為這孩子懇求神.
而且禁食,進入內室.
中夜躺在地上.
他家中的老神來到他旁邊.
要把他從地上扶起來.
他卻不肯起來.
也不跟他們吃飯.
原來在這裡記載著.
這個孩子患了重病的時候.
大衛也不是當這個孩子.
只是云云這麼多的兒子裡面.
其中一個沒什麼所謂.
他為了這個孩子患了重病.
他帶著一份很大的悲傷.
甚至他做了一個很極致的舉動.
就是用禁食的方式.
來為這個孩子祈禱.
值得留意的是.
他的禁食不是只不吃一餐就算了.
連續禁食了六天.
不吃飯.
在內室裡面躺在地上.
俯伏在地上來禱告.
希望上帝來幫助他.
其實在撒慕爾記下.
記載著大衛的禁食.
或者聖經裡面說的禁食.

$^{361}$也有好幾次的篇幅.
第一次記載著大衛的禁食是什麼時候呢?.
就是提拔他王素羅死.
和約拿丹死的時候.
他們被部屬砍了頭.
或者知道他的死訊的時候.
大衛就撕裂衣服.
撒慕爾記下一章十一至十三節.
說到撕下衣服.
然後就哀鳴英雄何故倒下來.
然後就為了素羅和約拿丹.
大衛就禁食禱告.
不吃飯.
來表達他對素羅的尊敬也好.
或者那份恩情.
又或者和這個好朋友約拿丹.
為了他的死來悼念.
他就禁食禱告.
除了這一次之外.
然後隔了不久的篇幅.
到了第三章.
撒慕爾記下第三章三十五至三十六節.
在素羅舊的部下裡面.
有一個大將軍叫做厄尼爾.
厄尼爾知道素羅已經死了.
於是大衛的權力又未穩固.
厄尼爾就帶著他自己的部下.
他的軍隊.
就投誠大衛當中.
當他投誠的時候.
大衛自己原本的部屬約拿丹就看不過眼.
於是就將厄尼爾殺害了.
殺了他之後.
當大衛知道了.
約拿丹派人殺了厄尼爾.
他圍著投誠的這個部將.
他又禁食禱告.
他心裡咬牙切齒.
為什麼要殺害這個投誠的人.
但同時又圍著選兵結帳的這件事.

$^{401}$他又禁食禱告.
然後他到第三次禁食禱告.
就是圍著這個孩子.
其實大衛家裡還有很多事情發生.
他的兒子暗戀.
暗戀他的女兒哈瑪.
然後就被亞沙隆殺害.
又沒有禁食.
不過這次他就圍著他的兒子.
撕裂衣服.
哀傷禱告.
後來他某程度上很疼愛的一個兒子.
最聰明的.
不過謀朝篡位的兒子亞沙隆.
後來也死了.
他又圍著亞沙隆的死.
他又哀痛撕裂衣服來哀傷.
整個三寶衣機下.
其實記載著大衛有好幾次禁食.
和極致哀傷的表達.
不過這次這個孩子.
去到拔士班所生的第一個孩子.
他面對重病的禁食.
和之前的不同.
之前的禁食全部都是那些人已經死了.
那種禁食是一種哀悼.
但在這裡.
這個孩子還沒死.
這個孩子只不過是面對一個很重的病患.
大衛就圍著這個孩子.
懇切的來禁食祈禱.
希望上帝可以收回成命.
希望上帝可以連戳這個孩子.
犯罪的是大人.
這個小孩是沒事的.
希望上帝藉著大衛這種禁食.
求神可以回心轉意.
讓這個孩子得到救贖.
某程度上這個孩子得到救贖.
也代表著上帝會不會收回成命.

$^{441}$之前一些很嚴厲的審判.
說刀劍離不開他的家.
又或者他所犯過的家人罪.
所有事情上帝可以一筆勾銷.
這種禁食也帶著.
會不會祈求上帝可以輕手.
可以有憐憫淋到在他的身上.
不過到了第七天的時候.
這個禁食就停止了.
到了經文裡面再繼續的.
來到說十八節到二十節.
到第七天.
孩子死了.
大衛的神僕不敢告訴他孩子死了.
因他們說孩子還活著的時候.
我們勸他.
他上者不肯聽我們的話.
若告訴他孩子死了.
豈不更加憂傷.
大衛見神僕彼此低聲說話.
就知道孩子死了.
問神僕說孩子死了嗎.
他們說死了.
大衛就從地上起來.
沐浴抹膏.
換了衣裳.
進夜無畫的殿敬拜.
然後回宮.
吩咐人擺飯.
他便吃了.
在這裡說到頭六天.
大衛面對著孩子患重病的時候.
他禁食的.
他乞求上帝.
希望上帝可以收回成命.
但當第七天.
當知道孩子死了.
不是直接知道.
而是見到神僕吱吱喳喳.
低聲說話.

$^{481}$很醒目的大衛知道.
一定心知不妙.
於是他便停止禁食.
沐浴抹膏.
換上衣服.
本來禁食的人是最頹的.
他停止禁食之後.
梳洗.
令自己的身體比較體面.
然後去到上帝的殿敬拜.
回到宮殿吃飯.
整件事會匪夷所思.
因為一般來說.
以前的禁食.
是為了死去的人禁食.
素羅也好.
厄尼爾也好.
若拿丹也好.
那些人死了.
會很慘很傷痛.
然後撕裂衣服.
在地上禁食.
不過今次的表達很不同.
今次神僕以為.
頭六天知道兒子死了.
那種哀傷更大.
在很高的哀傷中.
再升級.
再飆升.
加添那種仇苦.
不過大衛沒有這樣做.
反而他停止了禁食.
在這裡.
借這個地方.
也要順帶和弟兄姊妹.
說一下禁食的意義.
禁食除了是一種.
懇求上帝,尋求上帝的憐憫之外.
禁食有時候.
也代表著一種.

$^{521}$對自己身體的操練.
一個令自己更加專注.
祈求上帝一種禱告的表現.
有些人以為.
這段經文記載著大衛禁食.
上帝宣判之後禁食.
這個禁食不吃東西.
是代表著大衛做一些事情.
來懲罰自己.
不吃東西.
來糟蹋自己的身體.
令上帝的怒氣平息.
不過不是的.
禁食從來不是一個.
自我懲罰的方法.
當你生命裡面遇到一些難處.
遇到一些困難.
你不是用糟蹋自己的方法.
來面對自己生命的憂傷.
或者自己做錯很後悔的事情.
禁食只不過是.
令自己可以專注地去祈求上帝.
禁食也代表著一種懇求.
清心.
連飯也不吃.
專注在上帝面前.
祈求上帝的赦免和臨在.
禁食也代表著一種悲傷的表現.
面對著生命一些很重要的人.
或者一些關係失去的時候.
又或者國破家亡的日子.
藉著這種禁食.
刻苦己心.
令自己可以在上帝面前.
好好大徹大悟地去禱告.
不知道大家有沒有經歷過.
生命裡面遇到一些這樣艱難.
或者令到自己禁食.
代表著心靈裡面.
一些很大的悲傷的禁食.

$^{561}$我也分享我自己的一些經歷.
其實在我自己家人.
這幾年家人很多相繼過身.
心裡面真的有很大的悲傷.
曾經有一段日子.
弟兄姊妹說牧師你瘦了.
一來我正在做168.
那時候斷食.
生童.
第二件事也是有一些.
我不介意說.
就是希望自己可以專注在禱告裡面.
不吃東西來禱告.
不過我想說.
其實在禁食裡面.
真的專注去祈禱.
最初那一刻是真的.
反而專注不到的.
因為當你很想祈禱的時候.
你的肚子不停在那裡嘰嘰咕咕的.
你的身體告訴你.
其實你很需要.
你的身體開始累.
需要飲食.
但經歷過最肚餓的日子.
再進入禱告的時候.
你發覺那種功克己心.
那種專注在上帝面前祈禱的時候.
你發覺食物其實對你生命.
不是一些很重要的東西.
反而祈禱才是最重要的一種經歷.
我特別說一個經驗.
就是在我家人離開的時候.
有一段時間.
去做一個禁食祈禱.
在我自己最悲傷的時候.
穿著禁食.
那種專注祈禱.
很多時候就帶來了.
我生命裡面有一些安慰.

$^{601}$在禁食裡面上帝提醒.
生命過去裡面.
和家人的美好的關係.
那些不會成為.
今日生命裡面這些遺憾.
好好帶著一些.
舊時生命裡面美好.
可能覺得很短暫.
但這些美好的.
很短暫的時光.
其實都是我生命的一部份.
就單單聽到這句說話.
我開始有力量.
我也因為在禱告裡面.
跟著我說上帝.
我知道我不再需要.
再怎樣功克己心.
我就好好活在當下.
我就吃東西了.
我就開始回復自己正常禱告.
平時生活作息的禱告的生活.
有時候.
悲傷如果我們繼續鑽牛角尖.
或者生命裡面有一些不好的際遇.
你繼續鑽下去的時候.
你會越鑽越深.
生命裡面會面對這些艱難.
會越來越嚴重.
有些弟兄姊妹可能經歷生命裡面.
有一些悲痛的際遇.
她沒有走出來.
沒有停止它.
反而繼續給油下去.
繼續鑽下去.
沒有走向禱告裡面.
反而走向自己感懷身世.
或者覺得自己生命裡面.
有很多悲慘的際遇.
沒有辦法從一個悲傷裡面走出來.
如果我們繼續這樣走下去的時候.

$^{641}$不單止我們沒有辦法.
走出眼前的困境.
更加重要的是在我們的信仰靈性裡面.
我們會走向一個.
泥足深陷的地步.
大衛的生命不是一個.
大衛犯罪的面對這件事.
然後大衛面對上帝的審判.
要被擊殺的這個遺址.
他這些經歷絕對不是值得我們學校的.
不過他經歷了這些不濟的事件之後.
他值得我們去學校.
或者從他負面經歷裡面.
看到這個打不死的人的生命裡面.
有一樣東西最重要的就是.
他在這些逆境當中他抓住上帝.
他更加不容讓他自己負面情緒.
繼續被誘惑下去繼續去踩.
反而去到一個極致的悲傷當中.
他就停止了.
他就不容讓自己再繼續在.
悲傷的泥漿裡面繼續去掙扎.
反而他停止了這個禁食.
回到上帝的面前祈禱.
然後繼續他正常的生活.
他的老臣子.
他身邊的服侍他的人.
不明白大衛.
21至23節大衛就圍著他這種禁食.
這種表現來做一些解說.
神僕問他說.
你所行的是什麼意思呢?.
孩子活著的時候你禁食哭泣.
孩子死了你倒起來吃飯.
大衛說.
孩子還活著我禁食哭泣.
因為我想.
或者耶和華連說我.
使孩子不死也未可知.
孩子死了我何必禁食呢?.

$^{681}$我豈能使他返回呢?.
我必往他那裡去.
他卻不能回我這裡來.
在這裡其實大衛很清楚.
他告訴神僕.
既然孩子已經死了.
這已經是上帝定義的事實.
原本他禁食的時候.
他希望挽回上帝的手.
用盡他的方法.
用盡他的力量.
來求上帝讓這個審判.
或者上帝可以輕手.
讓這個孩子不用死.
所以他說到.
孩子的死未知.
不知道的.
還是一個未知之數.
盡他最後的能力.
希望可以挽回這個現狀.
但當孩子死了.
所有事情已經塵埃落定.
所有事情已經定數了.
他就知道.
已經是不能逆轉的事實.
既然不能逆轉.
他就順應了.
於是他就不再禁食.
他再說一句很唏噓的話.
我豈能使他返回呢?.
他沒有辦法將他起死回生.
反而他講到.
我必往他那裡去.
他更加明白.
他自己有一天好像會像孩子一樣.
他的肉體的身體會死亡.
在這個禁食裡面.
在這種悲傷當中.
某程度大衛開始明白上帝的心腸.
明白上帝的旨意.

$^{721}$亦都體會到自己罪孽的問題.
這些一切都已經成為上帝的定案.
面對上帝的這個定案.
他唯有的就是.
很好的來順應上帝的安排.
順著上帝這個審判.
繼續走面前人生的路.
頂子妹.
我們人生裡面.
難免會面對這些悲傷的時刻.
很多時我們面對這些悲傷.
面對人生這些我們覺得很痛苦的際遇.
我們經常有種想法.
如果人生可以坐時光機回到過去.
如果這件事可以重新來過.
我就不會是這樣.
我們有很多的回想.
很多的遺憾.
很想回到過去.
又或者不做這些決定.
又或者可以珍惜回到昔日.
可以相處的日子.
林林種種.
我們覺得如果人可以回到過去.
我們就好了.
不過人生永遠都要向前走.
時間是不會回到過去.
所有事都有上帝的命定.
所有事都是在上帝的掌管當中.
今日我們插不透.
不明白.
甚至為著過去我們做過的一些事.
我們覺得很遺憾.
不過既然這些事情已經出現了.
我們就要為我們自己的悲傷.
為我們自己的遺憾來指洩.
不要再讓自己繼續泥足深陷下去.
反而是我們令到自己.
不要再在這種悲傷裡面繼續打滾.
我再多說一遍.

$^{761}$承接著我剛才所說的家人過世的日子.
我講過很多次.
不過我覺得這段是對我生命很重要的經歷.
當我媽媽過世的時候.
我心裡面很悲傷.
在那種悲傷裡面.
我第一件事.
我在問上帝.
上帝,我這個家人對我這麼好.
大家都這麼疼愛他.
我都很努力事奉.
為什麼你不放在地上.
讓他在地上的日子可以多留一些.
讓我所愛的家人.
可以讓我和他多一點相處.
讓我有多一點時間和他接觸.
經歷上帝的恩典.
我和上帝說.
我媽媽都還沒有上栽培班.
我都還沒有跟她說成長百貨.
整個信仰我都還沒有透徹地告訴她.
她就離開了.
我生命裡面為了這件事很傷痛.
醫生甚至在媽媽過世之前.
說到因為她的癌細胞已經擴散.
擴散到什麼地步呢.
擴散到她的心臟裡面有一個叫心包.
不知道大家有沒有聽過.
心臟裡面有一個叫心包.
是另外一個器官.
心包裡面都有癌細胞.
所以她的身體和心臟.
就會被那些癌細胞擠壓住.
她呼吸很辛苦.
她的心臟每跳一下都會有很大的壓力.
我心想我祈求上帝.
你有神蹟將這些癌細胞清除.
求上帝苦苦哀求.
希望可以將媽媽挽回過來.
不接受醫生說.

$^{801}$她在香港的日子很短促.
可能隨時都會離開.
我心裡面大量的悲傷和憤恨.
希望上帝可以幫我們.
有很多的祈禱.
直至到媽媽最後都會過世.
在她過世的時候.
我再回到剛才禁食的祈禱裡面.
我發覺其實很多事情.
人生命我們沒有辦法強求.
如果我繼續去問.
以前如果不是看醫生.
或者用這些治療.
可能媽媽不會是這樣.
問很多這些假設性問題.
但到最後其實沒有答案.
帶來的只是內心裡面更加多的悲傷.
唯有接受眼前這個現實.
接受人生命裡面是短暫.
但也認定在她在世的日子.
我們都是頂天立地.
我們真心地去愛我所愛的人.
我珍惜這一段關係.
在短暫的日子裡面.
我們盡力了.
生命在人盡力之後.
沒有辦法逆轉.
面對的將來就是接受.
好好地去過自己的生活.
接受自己生命的悲傷.
停止自己不要再在悲傷裡面泥足深陷.
這是在大衛的身上我們看到的功課.
弟兄姊妹,今天講這段經文有點沉重.
不過我想講.
真的,今天我們要正視情緒.
我們的情緒其實是壓在我們心靈裡面.
或者有很多的祈禱.
或者我們生命裡面有些際遇.
我們解釋不了,處理不了.
我們不妨今天聽完這篇信息之後.

$^{841}$我們再去回想.
正視我們生命這點的悲傷.
不是叫我們再找出來.
然後再令自己自憐一番.
而是正視它.
不否定這一段的經歷.
不否定我們自己有這一份的情感.
但在上帝的恩典當中.
讓我們這些負面的情緒.
這些情感可以重新畫了一條界線.
我們向上帝祈求.
這些過去就成為過去.
亦都求上帝給我們有力量.
正視我們生命裡面上帝為我們所預備的將來.
唯有是這樣.
我們生命的瓶頸才可以被突破.
我們生命這個信仰才能夠幫助我們.
去檢視我們生命內心深處裡面.
一直未處理好的這些負面情感.
求主繼續幫助我們.
藉著大衛這一段不光彩的經歷.
這一段大衛的轉捩點.
在他生命當中.
我們讀得到一些的亮光.
讓我們生命繼續走向光明.
願主繼續彌陀賜福我們眾人.
\newpage



\section{約翰福音 6:24}
\label{sec:JCVpdjhfj4w}
\textbf{2024年10月6日崇拜直播｜梁家麟牧師｜生命的糧｜約翰福音6\_24 - 59}
\newline
\newline
連結: \href{https://youtube.com/watch?v=JCVpdjhfj4w}{\texttt{https://youtube.com/watch?v=JCVpdjhfj4w}} ~~~~ 語音日期: 2024-10-11
\newline
\newline
\hyperref[sec:U7nTwzqoXV8]{\small{< < < PREV SERMON < < <}}
~
\hyperref[sec:index_chronic]{\small{[返順時目]}}
~
\hyperref[sec:hd2ElWc7xho]{\small{> > > NEXT SERMON > > >}}
\newline
\newline
約翰福音 6:24
\newline
\begin{longtable}{cl}
\hline
\hline
章節 & 經文 (和合本修訂版)\\
\hline
6:24 & \begin{tabularx}{0.7\textwidth}{X} 這時眾人見耶穌和門徒都不在那裡，就上了船，往迦百農去找耶穌。 \end{tabularx} \\ \\
[1ex]
\hline
\hline
\end{longtable}
$^{1}$今日講到的經文是約翰福音六章二十四至五十九節.
不過我們只選讀其中的三段.
就是二十四至二十七節.
三十二至三十五節.
還有四十八至五十九節.
六章二十四節.
眾人見耶穌和門徒都不在那裡.
就上鳥說:我加伯隆去找耶穌.
既在海那邊找著鳥.
就對他說:拉比,是幾時到這裡來的?.
耶穌回答說:我實實在在的告訴你們.
你們找我並不是因見了神徵.
乃是因吃餅得飽.
不要為那必壞的食物勞力.
要為那全到永生的食物勞力.
就是人子要賜給你們的.
因為人子是父神所印證的.
三十二節.
耶穌說:我實實在在的告訴你們.
那從天上來的糧不是摩西賜給你們的.
乃是我父將天上來的真糧賜給你們.
因為神的糧就是那從天上降下來賜生命給世界的.
他們說:主啊,上將這糧賜給我們.
耶穌說:我就是生命的糧.
到我這裡來的.
必定不餓.
信我的,永遠不渴.
第四十八節.
我就是生命的糧.
你們的祖宗在曠野吃過晚了.
還是死了.
這是從天上降下來的糧.
叫人吃了就不死.
我是從天上降下來生命的糧.
人若吃這糧就必永遠活著.
我所要賜的糧.
就是我的肉.
為世人之生命所賜的.
因此猶太人彼此爭論說.
這個人怎能把他的肉給我們吃呢.

$^{41}$耶穌說:我實實在在的告訴你們.
你們若不吃人子的肉.
不喝人子的血.
就沒有生命在你們裡面.
吃我肉.
喝我血的人.
就有永生.
在末日我要叫他復活.
我的肉真是可吃的.
我的血真是可喝的.
吃我肉.
喝我血的人.
常在我裡面.
我也常在他裡面.
永活的父怎樣差我來.
我又因父活著.
照樣,吃我肉的人.
也要因我活著.
這就是從天上降下來的糧.
吃這糧的人.
就永遠活著.
不像你們的祖宗.
吃過晚了還是死了.
這些話是耶穌在加伯隆回堂裡.
教訓人說的.
很長一段聖經.
內容也有少許鬧彆扭.
今天篇道大家要有一點仇父心智.
緊接著五餅二魚的神蹟.
耶穌就接續食物這個主題.
同門道,講道.
祂才是人要期盼的食物.
真是超級食物.
故事的起始.
在五餅二魚這個神蹟發生之後.
很多人對耶穌崇拜的程度.
達到顛峰的程度.
四出去找耶穌.
總是要跟著祂.
一個能夠變出大量食物出來的宗教領袖.

$^{81}$一個能保證你們吃得飽.
衣食無憂的人.
你不跟他還跟誰呢?.
有奶就是糧.
耶穌是誰不重要.
祂的教訓是什麼都不重要.
跟著耶穌有飯吃才重要.
這些人從發生五餅二魚神蹟的.
加里尼湖的東岸.
開始四出去找耶穌.
終於他們在安息日.
在西岸的加伯隆的猶太會堂裡面.
就找到耶穌了.
對於這班狂熱的粉絲.
耶穌並沒有感到興奮.
這班群眾對他的身份和使命認識非常有限.
還帶著不少的偏差.
這很明顯是因誤會而結合.
所以即使他們如今很熱情地跟著耶穌.
沒多久他們就發現真相不是這樣.
因著了解他們就會分開.
耶穌很感慨地說.
在26到27節.
我實實在在告訴你們.
你們找我並不是因見了神蹟.
乃是因吃餅得飽.
必要為那必壞的食物努力.
要為那全度永生的食物努力.
就是人子要賜給你們的.
因為人子是負上帝所印證的.
猶太人士.
神蹟是一個符號.
一個sign.
這個符號是指著某些屬靈的實體.
屬靈的真相.
群眾追求的不是想認識背後的真相.
而只是想看到那個神蹟.
就是變出來的那個餅.
耶穌看到他們滿頭大汗.
到處東奔西跑去找他.

$^{121}$耶穌是憐憫他們的.
不過他們這樣.
只是真的名乎其實.
是胃口不馳.
餵兩餐什麼都肯吃.
耶穌就勸勉他們.
不如用這樣謀生的動力.
來追求有永恆價值的東西.
是的.
在自然界.
在一個資源比較貴乏的社會.
謀生是頭號大事.
消耗人最多時間精力.
生存最大的動力就是為了生存.
不計代價.
不擇手段.
當你覺得你自己陷到survival這個層次.
你就什麼都可以了.
什麼都可以了.
聖經也說.
要汗流滿面.
然後才可以得戶口.
如果我們用謀生.
相同的決心和動力去做一件事.
追求一樣東西.
幾乎無往而不利.
群眾只是為求餅而來.
但耶穌能夠給他們的不只是餅.
還有更珍貴更值得追求的東西.
餅吃了.
飽得一時還會肚餓.
還要再吃.
但耶穌卻能夠給他們永生的食物.
不僅僅是維持有限生命的口糧.
一天兩天.
而是能夠讓他們永遠活著的食物.
群眾當然不明白耶穌所說的話.
他們先問.
28節我們當行什麼才算作上帝的功呢?.
他們誤會耶穌的說法.

$^{161}$以為是耶穌說他們做得不對.
然後才不莫上帝所賜永恆的食物.
因此他們就問.
怎樣的穩食才是正確的方向和態度.
才能得到上帝所賞賜永恆的食物呢?.
耶穌的回答很乾脆.
29節.
信上帝所差來的.
就是做上帝的功.
意思就是.
信我吧.
信我.
就能夠得到永恆的食物.
不是做.
不是做什麼.
是信.
群眾再問.
30到31節.
信你?.
我們為什麼要信你?.
你說你有永恆的食物.
就是天糧供應我們.
我們不知道什麼叫天糧.
但我們知道從前祖宗.
在曠野吃過摩西所供應的嗎哪.
嗎哪從天而降.
大概就是天糧.
是不是?.
那你是摩西嗎?.
你有摩西的本事嗎?.
他們就問這個問題.
耶穌的回答.
在32到33節.
我實實在在告訴你們.
那從天上來的糧不是摩西賜給你們的.
那是我父將天上來的真糧賜給你們.
因為上帝的糧.
就算那從天上降下來.
賜生命給世界的.
我將它翻譯成這樣.

$^{201}$摩西.
摩西不是很大.
意思就是.
耶穌自己覺得大過摩西很多.
嗎哪也不是由摩西變出來的.
是天父上帝供應給以色列百姓的.
就像我今天說的天糧.
也是天父上帝親自賜給你們.
群眾聽到這裡.
就要求耶穌將天糧賜給他們.
明白不明白不重要.
有得吃最重要.
他們要求耶穌給他們天糧.
耶穌於是就做了一個偉大的宣告.
我就是生命的糧.
到我這裡來的必定不餓.
信我的永遠不渴.
至使我對你們說過.
你們已經看見我還是不信.
凡父所賜給我的人.
必到我這裡來.
到我這裡來的.
我總不丟棄他.
因為我從天上降下來.
不是要按自己的意思行.
乃是要按那差我來者的意思行.
差我來者的意思.
就是他所賜給我的.
叫我一個也不失落.
在末日卻叫他復活.
因為我父的意思.
是叫一切見紙面而信的人得永生.
並且在末日我要叫他復活.
35到40節.
耶穌不僅僅將天糧賜給群眾.
耶穌自己就是天糧.
凡是信祂的人.
就是屬耶穌的.
天父上帝就應允.
讓信耶穌的人在死之後.

$^{241}$可以復活.
得著永遠的生命.
所以信耶穌的人.
就永遠不餓不渴.
得著永恆的飽足.
我們見到前面第四章.
耶穌與雅各井旁的婦人談道的時候.
曾經說過祂自己是活水.
是生命的活水.
凡喝了祂賜的活水.
就永遠不渴.
這裡耶穌又提.
祂是生命的糧.
耶穌在水井旁用水的比喻.
在說到食物的時候.
又用食物的比喻.
耶穌要強調的很簡單.
祂是我們的供應者.
祂使我們得著永恆的飽足.
祂是我們終極追求和滿足的那一個.
在這裡我們可以稍稍欣賞一下.
耶穌的談道方法.
你可以看到耶穌的傳福音.
是真是見縫插針.
善於把握一切能夠引入福音的機會.
在我們這群庸俗的小老百姓面前.
耶穌從來不扮清高.
不說什麼玄妙高深的道理.
祂總是用人最基本.
最真實的追求來作為談道的話題.
甚至用祂來比喻福音.
比喻自己.
你要喝水.
耶穌就是活水.
你要吃餅.
耶穌就是生命的糧.
對我們香港小市民來說.
叉燒飯,菠蘿包.
是我們日常追求的食物.
耶穌在這裡就會說.

$^{281}$我就是叉燒飯.
我就是菠蘿包.
吃我的叉燒飯.
永遠飽肚.
吃我的菠蘿包.
終結滿足.
耶穌沒有覺得.
把自己扯上叉燒飯,菠蘿包.
是貶低自己.
破壞形象.
關鍵是他沒有看不起叉燒飯,菠蘿包.
沒有看不起小老百姓日常生活的追求.
耶穌不是微俗.
不要覺得祂是微俗.
祂是迎合民眾的低品味.
耶穌是肯定百姓日常的生活.
日常的追求.
你以為很清高嗎?.
你以為很高檔嗎?.
耶穌不是這樣看的.
耶穌可以和民事,法理,賽人.
做高層次的論道.
又可以坐在鴨寮街,洋屋裡.
和街坊阿伯,阿婆聊天.
福音在高峰論壇有效.
福音在街市魚檔同樣有效.
講到這裡.
你覺不覺得我們的主很迷人.
我看到我的主的談道方法.
真的覺得很過癮.
會堂是猶太人聚集的地方.
在場的猶太人.
當聽到耶穌說.
祂是從天上降下來的娘.
就議論紛紛.
不以為然.
什麼什麼從天上降下來.
吹就吹.
41到42他說.
我們認識這個人.

$^{321}$他不是阿墨水的兒子嗎?.
我們連他的家底都知道.
都起了.
還裝什麼神聖呢?.
加伯隆是耶穌傳道生涯.
慣常住宅的地方.
有不少耶穌的親戚.
同鄉在這裡.
他們認識約瑟.
認識瑪利亞.
根本不用做什麼網絡人肉搜尋.
我都知道你的底.
耶穌沒有興趣.
就他的神聖來源.
身份去做答辯.
確實這個超過了.
人的認知範圍.
你不要大家議論.
若不是猜我來的附給引人.
就有人能到我這裡來.
到我這裡來的.
在末日我要叫他復活.
在先知書上寫著說.
你們都要蒙上帝的教訓.
凡聽見父子教訓要學習的.
就到我這裡來.
這不是說有人看見過父.
唯獨從上帝來的.
他看見過父.
我實實在在的告訴你們.
信的人有永生.
43到47節.
只有耶穌來自天上.
認識父上帝.
所有人都不認識父上帝.
自然也無法確認.
耶穌的神聖身份.
耶穌承認除非是.
父上帝特別的挑選和感動.
沒有人可以認出.

$^{361}$耶穌的神聖身份來.
也沒有人會自然的.
接受耶穌的教訓.
這是神跡而獲得的信仰.
理性的討論是很有限的.
不會有效果的.
我們知道信耶穌.
在關鍵的部分.
是無法用理性有需的.
只是要靠信心的跳躍.
Lib of faith.
所以每個人的信主.
都是聖人的工作.
都是神蹟.
耶穌沒有興趣糾纏.
在他自己的神聖身份那裡.
卻堅持要.
賜下生命的糧這個主題.
在48到51節.
他會再說.
我就是生命的糧.
祖宗在曠野吃過嗎哪或死鳥.
這是從天上降下來的糧.
叫人吃鳥就不死.
我就是從天上降下來生命的糧.
人若吃這糧就必永遠活著.
我所要賜的糧.
就是我的肉.
是為世人的生命所賜的.
耶穌邀請人去吃生命的糧.
這個就是他的肉.
耶穌的說話很聳人聽聞.
令在場的人非常不安.
他們驚呼.
這個人怎能將他的肉給我們吃呢?.
耶穌很奇怪.
他沒有理會猶太人的反應.
反而他還要多加幾錢.
很奇怪.
變本加厲.

$^{401}$再三重申他的肉是可以吃的.
但必須給人吃.
唯有吃了他的肉才有永生.
53到58節.
我實實在在的告訴你們.
你們若不吃人子的肉.
不喝人子的血.
就會有生命在你們裡面.
吃我肉喝我血的人就有永生.
在每天我要叫他復活.
我肉真是可吃的.
我血真是可喝的.
吃我肉喝我血的人常在我裡面.
我常在他裡面.
永護上帝怎樣餐我來.
我又因父活著.
照樣吃我肉的人也要因我活著.
正如從天上降下來的量.
吃這量的人就永遠活著.
不像你們的祖宗吃過嗎哪.
還是死了.
我不知道你們聽這段話.
有什麼感覺.
耶穌不僅沒有打算.
消去猶太人的疑慮.
反而是加強.
並且在說吃他的肉的同時.
再加上喝他的血.
吃人肉喝人血.
血腥非常恐怖.
我們知道.
舊約律法是嚴禁吃動物的血.
所以耶穌叫人喝他的血.
這是跟律法的禁令對著幹的.
在當時期.
是可以撕裂衣裳的.
是非常古怪的說話.
我們解經解到這裡.
我們做一些分析.
有不少學者主張.

$^{441}$約翰福音六章.
22到58節這一段經文.
不是耶穌說的.
甚至不是約翰寫的.
是後來的信徒加插進去的.
他們的理據有兩個.
第一.
耶穌在這段話多次說到末日的情況.
他們認為耶穌應該在比較後的階段才會說到末日.
或者耶穌根本沒有這個末日的想法.
是後來的信徒才有的.
這是第一個他們的主張.
第二個他們的理據是.
耶穌說到他的肉可以吃.
基本上只能夠指向.
聖餐神學的變則說.
而這個變則說的理論.
是在耶穌死了之後.
早期教會才慢慢發展出來的理論.
不可能出自耶穌的口.
不少學者是這樣解釋的.
要說這段聖經不是.
仕途約翰一口氣寫成的作品.
甚至不是約翰寫的.
而是後人偽作的.
如果說這段說話不是耶穌的原話.
耶穌根本沒有這樣的想法.
我猜這也不是我們能夠選擇的主張.
作為福音信仰者.
這兩個說法我們都拒絕的.
我們認為耶穌當然可以有這樣的末日觀點.
他在世界末日重臨人間.
叫信他的人復活得著永生.
耶穌不僅僅是一個這樣的觀點.
耶穌是知道這是事實.
這是事實而不是觀點.
所以耶穌何時都知道有這樣的事實.
他在末日會再回來.
會使得我們得永生.
這不是一個主張.

$^{481}$這是他預告的事實.
所以不是後世信徒倒讚出來的觀點.
至於說耶穌的肉可以吃的說法.
其實絕大多數的聖經學者.
都將它聯繫到聖餐神學的變質說的說法.
跟大家上一堂很快捷的神學課.
基督教對聖餐有三個主要的主張.
第一個是變質說.
Transubstantiation.
這是天主教,東正教.
以及聖公會的部分.
更正教的觀點.
認為餅和杯在主禮人築謝之後.
就已經變成了耶穌基督的身體和寶血.
雖然外形仍然是餅是杯.
但實際上不是餅不是杯.
不要被你的眼睛騙了.
它已經實際上變成了耶穌的身體和寶血.
這是第一個主張.
變質說的主張.
第二個主張是同質說.
Transubstantiation.
就是餅和杯在祝謝之後.
同時兼具餅和杯.
和耶穌的身體和寶血.
這兩重特性.
神人異性.
就好像耶穌具有神人異性一樣.
但它們是雙重的特性.
又有神性又有物質性.
馬丁路德主張的臨在說.
Real Presence.
耶穌在我們築謝之後.
臨在餅和杯.
所以我們在領的時候.
同時兼領耶穌的身體和寶血.
這是同質說的說法.
加爾盟約路德的觀點加以修訂.
改為一個肅靈的臨在.
Spiritual Presence.

$^{521}$不是討論耶穌是否客觀地坐在餅和杯上.
讓我們一起飲喝.
不說這件事.
只要你認定你領的是耶穌的身體和寶血.
你就得著耶穌基督身體和寶血的果效.
就像我們中國傳統儒家說.
濟愚在,濟臣愚臣在.
你當有就行了.
當你認定了身體和寶血.
你就得著身體和寶血的果效.
這就是所謂的肅靈的臨在.
這是第二個主張.
第三個主張是紀念說.
這是由宗教改革家慈雲李.
首先提出的主張.
以後就成為多數更正教的主張.
餅和杯在築戒之後.
沒有性質上的變化.
我們仍然是領餅和杯.
只不過我們將餅和杯作為一個象徵.
藉以紀念耶穌的受死和犧牲.
是不是很複雜?.
大概就是這三個主張.
這是最快速的一堂聖餐神學.
主張變質說的人.
多數認為約翰福音第六章.
耶穌這段說話是指著聖餐而說的.
例如我現在在意大利服侍.
主要是溫州教會背景的人.
溫州教會普遍主張變質說.
相信餅和酒在築戒之後.
已經變成了耶穌的身體和寶血.
所以信徒領的不再是餅和杯.
而是耶穌的身體和寶血.
溫州教會變質說的主張.
主要的聖經根據就是這裡.
約翰福音第六章.
所以耶穌這段說話.
是在說聖餐神學的變質說.
我們旺角浸信會.

$^{561}$和幾乎所有浸信會的人.
都不會相信變質說.
所以我們不會用聖餐神學來解說.
約翰福音第六章耶穌的說話.
你有不同的看法.
你私下跟陳牧師說.
看他怎樣處置你.
我們是不相信這一套的.
回到這段聖經的解釋.
我既不相信這段說話.
是後世主張變質說的人.
就是生安白造.
冒充耶穌的話.
我同時也不相信.
耶穌會在這裡無端端說聖餐神學.
在加伯倫的會堂裡.
和小老百姓這段偶然的對話.
怎會牽扯到一百年後.
複雜的聖餐神學裡.
耶穌連聖餐都未設立.
怎會先說聖餐神學呢.
除非耶穌和當時的人對話.
純粹是借題發揮.
預先超過一百年後的人.
才明白的道理.
而不是要讓當時的人.
明白理解祂說什麼.
確實有很多門徒和群眾.
是聽不明白耶穌說的話.
不過這是他們資質有限.
理解能力有限.
不是耶穌根本在說火星話.
從來沒有打算讓他們明白.
就正如我今天在你們中間說到.
我是希望你們每一句都聽得明白.
當然我這是苛求你們.
每一句都和你們發生關係.
我是為今天而說的.
我不是想在這裡留下.
劉伯溫的燒餅哥式的寓言.

$^{601}$二百年後的人才明白.
我今天在說什麼.
原來梁家倫在預告.
地球人大量移居火星的情況.
解決地球人和火星人.
混分帶離的神學問題.
我現在先處理一個.
二百年後的神學問題.
怎麼可能.
我不相信這段話和聖餐有關.
日後教會在爭論.
聖餐神學的問題的時候.
為了搜羅聖經根據來支持他們的觀點.
宣稱《約翰福音》第六章是講論聖餐神學.
並且提倡辯則說.
這是將官理大冤枉耶穌.
講他們的觀點.
不僅僅耶穌不會在這個時候.
做聖餐的教導.
約翰也不可能在他的福音書裡.
講到聖餐或者聖餐神學.
如果你看《約翰福音》.
和《馬太福音》這些符類福音比較.
你會發覺很有趣.
《約翰福音》並沒有記載耶穌的受洗.
並沒有記載耶穌設立最後晚餐.
前面第一章我們看過.
約翰三次提到施洗約翰.
提到施洗約翰在約旦河邊替人施洗.
也提到施洗約翰宣告耶穌是上帝的羔羊.
不過就是沒有提到施洗約翰為耶穌施洗.
沒有提到這件事件.
同樣的《約翰福音》有記載耶穌的受苦死亡.
不過約翰寧願花很長的篇幅.
去記載耶穌對門徒的最後教訓.
也沒有提到最後晚餐的設立.
即是聖餐的事沒有提到.
有不少學者因此認為約翰有反聖禮的傾向.
至少他不是很注重.
以後教會兩大聖禮的基督論的根據.

$^{641}$我們怎麼可能會相信.
這位約翰會在第六章.
無情白事會安插一段.
不支持變則說的聖餐實行的講論.
此外我們要知道的是.
耶穌在這裡說的話.
和《福音》記載日後他設立聖餐的時候.
所說的話是有很大的差異.
在這裡耶穌說的是.
我的肉真的可以吃.
我的血真的可以喝.
吃我的肉喝我的血的人常在我裡面.
我也常在他們裡面.
在第56節.
耶穌赤裸裸地強調.
吃肉喝血.
但在設立聖餐的時候.
耶穌用的詞是非常委婉的.
他說身體.
body,soul,mind.
而不是肉.
不是flesh.
不是soul,mind.
並且他只是比較迂迴地說.
這是我的身體為你們捨的.
是捨身為我們.
而不是直接說.
你們吃我的肉喝我的血.
在設立聖餐的時候沒有這樣說.
我們很明白.
將這一段耶穌的話解釋為對聖餐的討論.
可以消去耶穌在鼓吹吃人肉的尷尬.
但是對我來說.
為什麼我們要用字面去解釋耶穌的肉和血呢?.
如果不用字面解釋.
就沒有吃人肉的困擾.
從上文下理看.
耶穌明顯是在找話題去說福音.
而他剛剛好用了群眾吃餅吃飽的訴求.
他順著群眾要求吃餅的心態.

$^{681}$指出他才能夠真正滿足他們的食物.
就像他在亞國井旁對婦人說.
他能夠賜下活水一樣.
耶穌實際上不是水.
實際上不是麵包.
只不過他比水和麵包更加能夠滿足人的需要.
在亞國井旁.
耶穌不是曾經有做過一個食物的比喻嗎?.
他說他以他完成父上帝的旨意.
等於他吃了一餐豐富的大餐.
我的食物就是遵行差我來者的旨意.
作成他的功.
耶穌用食物來做比喻的例子是很多的.
生命的糧根本不是一個特例.
耶穌在講道的時候.
常常用了大量象徵性的詞句.
例如在約翰福音.
耶穌說了七個「我是」.
包括了這裡的「我是生命的糧」.
之後他又說「我是世上的光」.
「我是羊的門」.
「我是好牧人」.
「我是復活和生命」.
復活在我,生命也在我.
「我是道路真理生命」.
以及「我是真葡萄樹」.
七個「我是」.
從來沒有人會用字面去理解耶穌的比喻.
說耶穌真是羊欄的門.
耶穌真是一個發光體.
今天教會我們常常會用吃.
用食物來做象徵.
例如今天的培靈會是一場屬靈的盛宴.
我們吃得飽.
沒有人會誤會.
培靈會是大食會.
我們最後要問.
耶穌說吃.
祂的肉和祂的血.
作為一個象徵性的語句.

$^{721}$指的是什麼意思呢?.
我們簡單將耶穌前後兩個說話對照一下.
54節和40節.
「吃我肉喝我血人就有永生」.
「在末日我要叫他復活」.
「叫一切見子而信的人得永生」.
「並且在末日我要叫他復活」.
兩個說話.
基本上後面的意思完全一樣.
結構也一樣.
因此很明顯.
「吃我肉喝我血」.
就是「見子而信」.
「吃我肉喝我血人就有永生」.
就是藉著「信」來到主的面前.
跟祂結連.
住在祂裡面.
也讓祂住在我們裡面.
所謂「共同入住」.
就是你住在我這裡.
我住在你這裡.
更重要的是.
當耶穌說祂是生命的糧的時候.
祂指的是祂要犧牲祂自己的生命.
讓人得生命.
接受祂的代屬替死.
也從祂那裡得著新的生命.
就是能夠展現到永恆的生命.
耶穌說祂是生命的糧.
吃祂喝祂就得到永生.
這句話有效性不僅僅限於聖餐的時候.
是任何人任何時候都有效的.
我們在領聖餐的時候吃喝耶穌.
在沒有領聖餐的時候.
我們仍然靠著耶穌得生命.
關鍵的是我們要信祂.
四十節說.
因為我父的意思是.
叫一切見子而信的人得永生.
四十七節又說.

$^{761}$信的人有永生.
人是因為信耶穌而得永生.
不是因為領聖餐而得永生.
重點是相信不是聖餐.
所以千萬不要宣揚領聖餐得永生.
不領聖餐就注定滅亡.
這樣的歪理.
我們教會沒有這樣的歪理.
還有耶穌不是說.
吃了聖餐就算你相信變質雪.
已經變了耶穌基督的身體和寶血.
你就可以有永生的效果.
因為聖餐不同妖精吃了唐僧的肉.
唐僧的肉是長生不老藥.
我們知道的.
吃了就是青春常駐.
與天同壽.
對不起.
聖餐就沒有這樣的滋補能力.
耶穌說的是.
差我來者的意思.
就是祂所賜給我的.
叫我一個也不失落.
在末日卻叫祂復活.
因為我父的意思是.
叫一切見之而信的人得永生.
並且在末日我要叫祂復活.
39到40節.
44節也說到我者來.
在末日我要叫祂復活.
你看到嗎.
是耶穌在末日使我們從死裡復活.
不是聖餐使我們永遠不死.
是祂使我們死了也能復活.
耶穌將祂自己.
與出埃及事件中的嗎哪作比較.
58節說.
不像你們的祖宗吃過嗎哪還是死了.
耶穌的意思不是說.
以色列人吃嗎哪死了.

$^{801}$嗎哪有毒.
毒死百姓.
耶穌只是說嗎哪沒有為人帶來永生.
吃嗎哪不是直接讓人致死.
同樣吃聖餐也不能讓人直接得永生.
嗎哪和聖餐都不是重點.
關鍵在於耶穌.
耶穌讓人得生命.
信耶穌才得永生.
幾節說這麼久.
你們忍這麼久也算很厲害.
最後跟大家說一個.
我在讀經文時的感受.
耶穌在這裡用了很激烈的言辭.
吃肉喝血的比喻.
說的時候又說得實牙實齒.
非常認真.
你簡直覺得他咬牙切齒.
我肉真的可以吃過.
我血真的可以喝過.
吃我肉喝我血的人就有永生.
聽到你真的想死.
耶穌說這些話時完全不理.
聽這些話的人的心理反應.
甚至會被誤導.
耶穌為什麼要這樣說呢?.
我一直在問耶穌為什麼要說這些話呢?.
為什麼要說這麼激烈的話呢?.
我唯一能夠解釋的是.
耶穌在這個時候.
他自己的心情都是很沉重很激動.
雖然他不是真的要將自己的肉和血分給我們吃.
割自己的肉給我們吃.
不是.
不過耶穌確實是真實地準備好要捨棄他自己的一切.
包括他的生命來拯救我們.
是不是?.
這是基督版的賣肉養孤儀.
這是基督版的割肉餵嬰.
耶穌已經預備好為我們這群人做出這樣的犧牲.

$^{841}$不過可能他有些煩躁.
對著我們這群傻乎乎的人.
他已經準備好要為我們犧牲.
而我們唯一要做的就是相信他.
接受這個白白的救恩.
就是這麼簡單.
相信和接受是我們唯一要做的事.
就好像在最後晚餐前.
耶穌為門徒洗腳.
彼得冷謙卑地說.
他不配接受耶穌的服侍.
耶穌就很嚴肅地回應.
「我若不洗你,你就和我無份」.
耶穌正正是要將他的肉割出來給我們吃.
不吃他的肉.
我們就沒辦法得著永生.
耶穌的失是我們的得.
耶穌的死是我們的生.
耶穌用很沉重的語句.
來講明這個沉重的事實.
「那知他為我們的過犯受害.
為我們的罪孽壓傷.
因他受的應罰,我們得平安.
因他受的鞭傷,我們得醫治」.
我只有一個形容.
就是我們親愛的主耶穌.
我好愛好愛的耶穌.
\newpage



\section{耶利米書 15:15-21}
\label{sec:hd2ElWc7xho}
\textbf{2024年10月27日崇拜直播｜陳明泉牧師｜正視你的情緒－當你感到自憐時......｜耶利米書15\_15-21}
\newline
\newline
連結: \href{https://youtube.com/watch?v=hd2ElWc7xho}{\texttt{https://youtube.com/watch?v=hd2ElWc7xho}} ~~~~ 語音日期: 2024-10-27
\newline
\newline
\hyperref[sec:JCVpdjhfj4w]{\small{< < < PREV SERMON < < <}}
~
\hyperref[sec:index_chronic]{\small{[返順時目]}}
~
\hyperref[sec:limrMw42910]{\small{> > > NEXT SERMON > > >}}
\newline
\newline
耶利米書 15:15-21
\newline
\begin{longtable}{cl}
\hline
\hline
章節 & 經文 (和合本修訂版)\\
\hline
15:15 & \begin{tabularx}{0.7\textwidth}{X} 耶和華啊，你是知道的；求你記念我，眷顧我，向迫害我的人為我報仇；不要把我取去，因你不輕易發怒，要知道我為你的緣故受了凌辱。 \end{tabularx} \\ \\ \relax
15:16 & \begin{tabularx}{0.7\textwidth}{X} 耶和華－萬軍之神啊，我得著你的話就把它們吃了，你的話是我心中的歡喜快樂；因我是稱為你名下的人。 \end{tabularx} \\ \\ \relax
15:17 & \begin{tabularx}{0.7\textwidth}{X} 我並未坐在享樂人的會中歡樂；因你的手，我就獨自靜坐，你使我滿心憤慨。 \end{tabularx} \\ \\ \relax
15:18 & \begin{tabularx}{0.7\textwidth}{X} 我的痛苦為何長久不止呢？我的傷痕為何無法可醫，不能痊癒呢？難道你以詭詐待我，像流乾的河道嗎？ \end{tabularx} \\ \\ \relax
15:19 & \begin{tabularx}{0.7\textwidth}{X} 所以耶和華如此說：「你若回轉，我就使你歸回，站在我面前。你若能將寶物和無用之物分別出來，你就可以當作我的口。他們必歸向你，你卻不可歸向他們。 \end{tabularx} \\ \\ \relax
15:20 & \begin{tabularx}{0.7\textwidth}{X} 我必使你向這百姓成為堅固的銅牆。他們必攻擊你，卻不能勝過你；因我與你同在，要拯救你，搭救你。這是耶和華說的。 \end{tabularx} \\ \\ \relax
15:21 & \begin{tabularx}{0.7\textwidth}{X} 我必搭救你脫離惡人的手，救贖你脫離殘暴之人的手。」 \end{tabularx} \\ \\
[1ex]
\hline
\hline
\end{longtable}
$^{1}$今日經訓係記載於耶利米書15章15節至21節.
《舊約聖經》1046至1047頁.
耶利米書15章15節至21節.
耶和說:「你是知道的,求你紀念我,眷顧我,向逼迫我的人為我報仇,.
不要向他們忍怒取我的名,要知道我為你的緣故受了凌辱.」.
耶和說:「萬君之臣啊,我得著你的言語就當食物吃了,.
你的言語使我心中的歡喜快樂,因我是稱為你名下的人,.
我沒有坐在燕洛人的會中也沒有歡樂..
我因你的感動獨自靜坐,因你使我滿心憤恨,.
我的痛苦為何長久不止呢?我的傷痕為何無法醫治不能痊癒呢?.
難道你待我有鬼詐像流竿一河道嗎?」.
耶和說:「如此說,你若歸回,我就將你再帶來,使你站在我面前,.
你若將寶貴的和下賤的分別出來,你就可以當作我的口,.
他們必歸向你,你卻不可歸向他們,.
我必使你向這百姓成為堅固的銅牆,.
他們必攻擊你,卻不能勝你,因我與你同在,要拯救你,搭救你.」.
這是耶和說的,「我必搭救你脫離惡人的手,救贖你脫離強暴人的手.」.
聖經獨筆.
弟兄姊妹平安!.
今日我們是最後一講,正視情緒系列.
今日講的情緒是比較複雜的情緒,.
這個情緒叫做自憐.
自憐是一個混合很多個人的看法.
對環境,對自己和情緒的複雜表現.
我不知道大家有沒有經歷過這種自憐的狀態.
簡單來說,自憐就是一種覺得自己很可憐的感受.
這種感受不是純粹有一種愛自己的表現.
反而是一種憤怒的表現.
很多時候我們的憤怒是向外表達出來的.
那些叫做Anger Out.
我們憤怒就被人知道這是一個訊息.
在做一些冒犯我,令我不喜歡的事情.
不過自憐卻是Anger In.
就是那些憤怒是沒有辦法有對象去接受的.
你作為妻子,你發怒的時候,你丈夫知道發生什麼事.
就懂得去調教你的反應,或者調教你的生活態度.
但是Anger In就是身邊的人忽視你,不理你.
他憤怒的過程主要是憤怒自己.
甚至乎他連憤怒的對象是誰都不知道.
在這種狀態下,往往那種情緒表現會將自己抽離.

$^{41}$有時候會看外面的挑戰實在太大.
自己的能力實在太渺小.
自己的堅持和外在的不友善.
成為了很強大的落差.
那你聽到那種情緒很複雜.
不過在人的成長裡面,或多或少都會有這樣的經歷.
在我自己的成長裡面,我記得十五,十六歲.
我曾經經歷過一個這樣的階段.
話說那時候讀書不是很厲害.
我記得我中三的時候有個老師跟我說.
陳詠全你是可以讀書的.
然後我覺得原來我在老師眼中是一個可以讀書的人.
是一個可以透過努力拿到好成績的學生.
於是我真的很用心開始想讀好書.
不過一直讀的時候我發覺事與願違.
我很想做好的科目,完全做不到.
而且一落千丈,尤其是計算,那些數字很差的.
當一條不懂,然後整個學期一直累積下去的時候.
就覺得那樣東西實在太過龐大.
自己沒有辦法再去追.
結果是做些什麼呢?.
以前我也是一個比較活潑開朗的年輕人.
但是經歷了這種自憐之後,我開始自己收起來.
以前我很喜歡跟別人一起吃午飯.
但當我自己開始覺得自己有點廢的感覺.
覺得自己不行的時候,我開始抽離.
開始自己在吃飯,然後一直演變為自己回家吃飯.
再演變就是在午餐吃完午餐之後.
有時有個多小時吃飯,然後就不再上學.
然後下半週就逃學.
經常說自己生病,說自己又不舒服.
一直到最後的階段,到中午放榜.
成績出來了,是一個很差的會考成績.
會考只有一分,大家知道的.
我收到這張成績單,有一段日子.
我是找個大被蓋著頭.
不想見任何人,覺得辜負了家人.
辜負了身邊所有覺得自己是可以的人.
甚至乎,我問自己,我在做些什麼.
覺得自己是全世界最可憐的人.

$^{81}$不過,其實這一段經歷.
相對於今天坊間很多病態,我不是太長.
可能經過一段日子之後,慢慢重新起來就可以爬起來.
不過我自己去看回自己.
為什麼這麼快呢?.
因為我性格某程度底都不是至於很悲觀.
或者是一些很執著的人.
所以有些調教會容易一點.
但是有些朋友,有些弟兄姊妹,性格比較執著一點.
執著某一些信念.
他的信念和環境的落差越大.
他面對著那種感受的強烈就會越來越深.
今天我們去看的這段經文,耶利米就是一個這樣的人.
耶利米先知,他不是一個面面俱圓,說話很好聽的先知.
甚至乎,因著他所說的話語,他面對著很多的艱難.
今天很多人說,或者外國有一些的港元就說.
自憐是一種毒素.
如果面對著這件事的時候.
有時候就要處理好像罪的方式來處理.
他要跟他斷絕關係.
不過是不是一兩句鼓勵的話.
又或者有一些很正面的信息.
就可以處理好你內裡的那種自憐的感受呢?.
今天的輔導學家或者心理學家,他說的不是.
反而可能自憐是你成長生命裡面一個很重要的契機.
覺得自己可憐,可能就是正正進入了一個自我更新的境地.
在聖經裡面,我們今天看的耶利米書.
耶利米也進入了一種這樣的狀態.
聖經很直率地說他的心路掙扎和歷程.
今天第十五章就記載著他跟上帝的那種質問.
甚至乎他在說他自己那種很痛苦的情緒.
上帝怎樣去幫助他,幫他繼續走面前的路呢?.
希望今天我們從這一段經文.
我們會長一點的篇幅去看.
希望在耶利米生命當中,我們找到一些生命亮光.
我們先祈禱求上帝幫助.
阿巴天父求你在我們當中帶領我們.
讓我們在你的話語當中,找到我們生命要走的路.
幫助我們藉著你的話語,藉著昔日的先賢.
他怎樣經過信仰的路.

$^{121}$這些的心路歷程都能夠成為我們今天取得亮光的一些經歷.
幫助我們每一個,讓靈聽的弟兄姊妹每一個.
不單止聽得明白,也能解開我們的心結.
又幫助宣講的講得清楚.
將你的話語,忠心地向弟兄姊妹宣講.
讓我們一同領受,我們彼此得益.
獻上禱告奉求基督耶穌得勝的聖名,阿們.
我們一起看這一段經文.
耶利米先知,剛才說過了,是一個什麼先知呢?.
耶利米先知其實是在整個以色列國.
整個猶太國滅國之前的一個很重要的先知.
他所說的信息其實很簡單.
四一字,在中文來說,審判和悔改.
在當時的國家來說.
整個以色列國離棄上帝到了一個地步.
已經使得上帝不可以不再管教.
或者上帝不可以再啞忍了.
所以上帝就派遣他的僕人.
來將一個重要的審判信息告訴當時整個猶太國.
整個以色列民族知道.
他要說的就是,他在說這個國家將要被擄.
這個國家已經是沒得救了,在現在這個階段裡面.
要經歷過一個好像重生一段很受苦的過程.
他們再歸回的時候,他們才可以重新振作.
所以這個審判的信息是一個不好聽的信息.
最重要的是除了在說上帝的審判之外.
也要說這一幫人需要悔改.
上帝的審判源自這一幫人死牛一面頸.
或者滿足眼前的燕落,或者眼前的生活.
忽略了上帝,或者敬拜外邦的偶像.
他們整個國家,由上至下,他們全部都需要悔改.
所以耶利米很單純地說這個信息.
在當時的社會裡面流行說一個信息.
就是說什麼信息呢?就是平安的信息.
是說外族來到不用怕.
一切上帝會幫我們化解,平安平安.
一切都能搞定的.
不過耶利米就跟這些當時的假先知.
被稱為假先知的人.
耶利米很清楚地說.

$^{161}$說的平安其實是沒有平安的.
你所說的搞定,其實沒有搞定.
如果你們生命不徹底悔改的時候.
一切都沒有平安.
所以當時的主流社會在說這樣的信息的時候.
耶利米就跟他們開倒車.
可以想像的就是.
一個經常說一些不中聽的說話的先知.
在社會上是不受歡迎的.
再加上耶利米被上帝呼召的時候.
上帝對他特別嚴格.
對他嚴格到一個什麼地步呢?.
叫他終生不娶不嫁.
就是說他不會有婚姻的.
他要過著最簡樸的生活.
在說這個社會上被人嫌棄的那種生活.
耶利米就要經歷一個這樣的生活.
他是孤軍作戰.
他沒有任何的群體去支持他.
他就是極力一個人.
來面對著說一些不好聽的信息.
也面對著所有人以他為首.
所以當時他長達幾十年的先知生涯裡面.
他面對著的就是.
自上的君王到在下的平民百姓.
見到他都覺得他這個人真的很怪.
這個人實在太過離譜了.
這個人說話實在太不中聽了.
每一個人都排擠他.
說他進入一種叫做非社會的狀態.
或者你說反社會也好.
非社會也好.
就是說他的生活.
他沒有任何的群體.
他唯一支撐著他的只是上帝的話.
是上帝他自己的.
到了第十五章的時候.
耶利米面對著一種這樣的狀況.
十五章十五節.
耶和華你是知道的.

$^{201}$求你紀念我 眷顧我.
向逼迫我的人為我報仇.
不要向他們忍辱取我的命.
要知道我為你的緣故受了凌辱.
在這裡面他講到他的狀態.
他的心靈狀態是什麼呢?.
他求上帝眷顧他.
他其實覺得自己很慘.
他的慘是上帝知道的.
他甚至乎覺得所有人都以他為敵.
他到了最後他沒有辦法來掃尊.
他那種憤怒.
他又不可以自己來報仇.
唯有他求上帝幫他來報仇.
不過眼前的生活.
或者這個禱告.
我想他在經文裡面.
或者整個書《耶利米哀歌》裡面.
這一類型的禱告不止一次.
是多次.
可想而知一個力排眾矣.
逆流而上的先知.
面對著心靈的掙扎有多大.
在經文裡面他特別提到.
你不要忍辱取我的命.
不要恨他們忍辱取我的命.
簡單來說就是.
上帝你不要容忍他們繼續.
然後你就要剝奪我的性命.
即是侮辱我.
你濫用了我.
來縱容這一班逼迫他的人.
所以上帝你又不輕易發怒.
寫有豐盛的慈愛.
你就對著那些人有豐盛的慈愛.
對我你就覺得這麼嚴苛.
其實耶利米在這裡是發脾氣的.
他在這裡甚至乎他求上帝幫他報仇.
所有的一切.
所有的凌辱.

$^{241}$原因是什麼呢.
為了你啊上帝.
我為了你受了凌辱.
你看到其實他講的這一句話.
或者他求上帝幫助他.
不是純粹一種情感的反應.
耶利米他講的這番話.
其實背後就是來自他的呼召.
在耶利米書一章六節.
這個是他年輕時期.
年輕的日子的時候.
上帝呼召耶利米才知道.
耶利米和上帝的對話是怎樣呢.
一章六節八節.
我就說主耶和華.
我不知怎樣說.
因為我是年幼的.
耶和華對我說.
你不要說我是年幼的.
因為我差遣你到誰那裡去.
你都要去.
我吩咐你說什麼話你都要說.
第八節這裡最重要了.
你不要懼怕他們.
因為我與你同在.
要拯救你.
這是耶和華說的.
在他的呼召裡面.
耶利米說十幾歲或者年輕的時候.
我只是年輕的.
我怎樣跟世界逆流而上.
耶利米這樣發出呼聲.
上帝就說你不用怕.
你不要怕他們.
我叫你說什麼就說什麼.
最重要的就是我與你同在.
我一定要拯救你.
上帝一定會拯救耶利米.
不過在耶利米當下的感受是什麼呢.
上帝都不幫他.

$^{281}$如果幫他.
他不需要受這些的凌辱.
他不需要受這些的逼迫.
在耶利米當下的感受當中.
如果是上帝幫他的.
他會幫他拔走所有的困難.
幫他所說的話.
當他振臂一呼的時候.
所有人就悔改悔改悔改.
令到他不再面對這些的艱難.
又或者這麼說.
耶利米也不是省油的燈.
他未至於很單純地覺得.
他說的話可以令人立即信服.
但至少在這些逼迫的過程中.
上帝要令他好過一點.
拯救他.
讓他不會面對這些的艱難.
經文裡面除了說這一句話之後.
然後由第十六節到第十八節.
耶利米進一步說他的心靈狀態.
十六節.
耶和華萬軍之神.
我得著你的言語就當食物吃了.
你的言語使我心中歡喜快樂.
因我是稱為你名下的人.
我沒有坐在燕洛人的會中.
也沒有歡樂.
我因你的感動獨自靜坐.
因你使我滿心憤恨.
我的痛苦.
為何長久不止.
我的傷痕為何無法醫治不能痊癒.
難道你待我有鬼咋.
像流肝的河道嗎.
在十六節到十八節再一次表白他的痛苦.
他這個痛苦.
他其實在說的內容.
其實對照上帝之前應許他的內容.
甚至乎他面對的處境.

$^{321}$是上帝叫他這樣做.
帶來的更大的苦果.
他說道.
你的言語是我心中的快樂.
上帝的話語很甘甜.
聽到上帝的話語.
心很有方向感.
很有力量.
他是這樣想的.
不過在他的生活當中.
毫無快樂.
當人們去聚餐的時候.
當人人在一起喝酒作樂.
吃東西很開心.
他自己坐在一旁.
他自己是沒有快樂的.
他對上帝完全的聽從.
上帝要他怎樣做他就怎樣做.
但這種聽從.
不是為他帶來生命中更多的快樂或滿足感.
帶來的只是失去了快樂的情緒.
然後他再繼續說.
我是被稱為你名下的人.
而你不是說我是屬於你的才知道.
我就是跟耶穌說.
上帝僅僅扣連的被稱為這個人.
當時的社會每個人都知道.
甚至乎他自己都是以上帝的名義.
來講上帝的話語.
他被稱為上帝名下的人.
在舊約裡面說.
屬神的人會怎樣.
凡事都凡他所作的盡都順利.
思想上帝的話語會比蜜柑甜.
比風防下敵的蜜柑甜.
生命得著他就得著喜樂.
舊時舊約是這樣說的.
屬上帝的人會得到一些屬靈的好處.
非物質的好處都得到一點吧.
不過他說沒有.

$^{361}$什麼都沒有.
他只是獨自靜坐滿心憤恨.
面對著所有人.
這個對他的那種敵黨.
一個被稱為屬神的人.
得不到尊榮.
只是成為了社會上被唾棄的人.
他成為一個孤單抽離的一個人.
再繼續講的是第十八節.
這是一句最尖銳的話.
他說上帝你是不是騙我.
我的痛苦為什麼長久不止息.
我的傷痕無法醫治不能痊癒.
難道你待我又詭詐.
這句話本來也是很客氣.
有點客觀式.
原文是說上帝你騙我.
難道你騙我嗎.
騙他什麼呢.
如果看回上帝對他的二章十三節.
你不是輸.
上帝說自己是活水傳元.
當人倚靠上帝的時候.
生命是源源不息.
就好像井旁婦人見到耶穌.
耶穌說她是活水.
飲於基督的時候.
生命就永留不息.
你就會有那種生命力.
上帝是這樣應許耶利米的.
不過耶利米當下覺得什麼.
他覺得他是流乾的河道.
生命很乾燥.
人生這樣走下去.
他也在問是否要繼續這樣走.
他說上帝說謊.
應許過他的事沒有兌現.
他面對的只是艱苦困難.
頂子妹不知道如果你是耶利米會怎樣.
如果面對上帝叫你做一些事.

$^{401}$你一直想著做的時候.
會有好的結果出現.
但一直做的時候你發覺.
越做越低.
越做越沒有果效.
越做越沒力.
越發覺身邊的人以你為敵.
你覺得自己的生命.
都不知道為了什麼.
你的人生就好像對著什麼來努力呢.
上帝有這個感召在你身上.
但當你一直實踐的時候.
你體會不到實踐上帝召命的.
那份踏實感和滿足感.
生命可能面對很困難.
我想我們有時候都未必.
去到耶利米這個境地.
但有時候我們都會問問自己.
在生活當中.
蒙祿為了什麼呢.
我們為了人生.
形形意意地擺上.
但到最後事與願違.
期望祝福.
卻換來別人對自己的咒詛.
為了什麼呢.
楊木屋牧師寫這段經文寫得很好.
他的注釋寫得非常好.
他說如果他是耶利米.
大概有三條道路可以選擇.
三條路.
第一條路.
如果是耶利米可以選擇的.
第一就是否定過去.
什麼叫否定過去呢.
就是說年輕的時候.
上帝呼召耶利米.
叫他成為一個先知.
耶利米如果面對這些艱難.
可以將這一段經歷抹走.

$^{441}$將他一路走過的路.
撇下.
對過去採取一種完全否定的態度.
只是一整晚.
或者他自己收錯了.
所以面對這些艱苦.
是我自己收錯了.
我誤解了.
或者我傻了.
年少氣盛.
所以將昔日的感召.
將昔日覺得令他很真摯.
很初心的回應.
他覺得是一個自棄的表現.
這些過去變得沒有意義.
沒有價值.
將這些一切的經歷.
從他自己腦袋的回憶當中.
全部刪走.
否定過去.
可以成為今天很多人的一些經歷.
面對著一些基督徒.
我曾經說過.
有些基督徒.
他知道在聖經讀的價值.
和這個世界的價值.
有很大的落差.
我跟大家說過.
有個律師.
他覺得他所面對的罪惡世界.
根本這個信仰就無法實踐.
到最後他放棄了信仰.
帶著憤怒.
來面對基督徒.
他說.
如果這個神是真的.
為什麼這個世界這麼多不公義.
所以他選擇否定.
這些過去在成長階段裡面.
所領受的召命.

$^{481}$這個是第一個.
楊木菊說.
耶利米大可以這樣做.
第二個可以去得再盡一點.
否定神.
這個世界根本就不會有神.
這個世界既然這麼不公義.
他徹底地將上帝排除在他的生命當中.
索性和這個世界一起的.
來到走齊.
跟著這個世界的步伐來走.
按著這個世界的世界觀念.
來為自己打平.
上帝只不過是前人所說的話.
這個世界沒有上帝.
否定神.
將上帝否定在他的生命當中.
如果你覺得這個說法很抽象.
你有機會去看一套電影.
或者一本書叫做《沉默》.
有些氣教者採取的就是這種方式.
踐踏耶穌的那個象.
氣教.
跟著過回他們自己的生活.
不過真的可以這麼容易.
放棄了.
你否定了上帝.
上帝就不在了.
在生命裡面.
這一班氣教者.
心靈裡面隱隱作痛.
當他自己一路走人生的路的時候.
根本沒有辦法去否定這件事.
楊木屋說.
最後一條路.
也是耶利米最終記載在聖經裡面所說的路.
就是跟上帝對質.
跟上帝爭辯.
所以你看到整個耶利米書.
耶利米哀歌.

$^{521}$我們讀的時候.
我們不可以用系統神學.
這樣格格地讀.
因為它是對話.
大部分是跟上帝爭辯.
上帝你這樣答應我.
但為什麼沒有這樣實現.
上帝為什麼你叫我做這些事.
我要面對這些痛苦.
如果我是吸毒.
我是做一些壞事.
我面對這些受苦.
心甘命瀝.
但我現在做對事.
按著你的旨意而行.
但卻面對著一個這麼大的困難.
迫不可避.
為了什麼.
為什麼我的生命要面對這樣.
其實耶利米整個書卷.
那種哀歌的形式.
就是用一種跟上帝對質.
這種表達.
來呈現他的痛苦.
和面對他生命的經歷.
我說楊木谷寫得很好的原因就是.
你猜耶利米.
他的爭辯是為了什麼.
很多人在說.
剛才在說.
那些人去向人說上帝的道.
其他人是不能勝過他.
沒有人可以救他.
耶利米到了這一刻.
到了最痛苦的時刻.
他和上帝爭辯.
他是不是為了要贏上帝呢.
楊木谷的主色說不是.
他和上帝爭辯是為了輸.
他輸的原因是.

$^{561}$上帝你提升我的層次.
你令我心悅誠服.
跟你照明.
所以對質的目的.
不是代表著一種不敬虔的表現.
又或者對不敬的態度.
我們中國人學道.
我們經常說對老師不能挑戰.
不能對質.
但是到了這一刻.
到了心靈裡最沉重的狀態.
那種對質.
目的不是要爭辯贏.
你覺得你高於一個層次的對象.
不是和神爭一日之長短.
而是要令自己心悅誠服.
要令自己咬輸.
咬輸在這個爭辯當中.
藉著這個爭辯處於失敗.
讓耶利米進入另一個層次.
進入一種生命的轉化.
頂梓梅我不知道你在人生裡有沒有和上帝爭辯.
你為了什麼和上帝爭辯.
為了贏還是為了輸呢?.
今天我們心裡的信念是什麼呢?.
有時候很悲壯地說.
今天很多人我們都是活在一個.
我們有信仰但沒有信念的世代.
我們生命裡究竟有什麼是我們一直堅持著的呢?.
自憐的感覺反而是值得寶貴的原因.
就是因為你堅持而面對一些困難.
如果我們生命裡所有東西都是很順.
或者所有東西都是按著自己的心意去做.
沒有掙扎,可能我們生命裡沒有真正堅持的位置.
痛苦往往是因為生命裡有一些很重要的價值.
有一些信念在堅持著.
正正就是這種信念和環境有很大的落差.
以致我們面對這種張力.
心靈裡有憤怒,有不甘,有被愛等等複雜的情緒.
上帝面對著耶利米這一連串的對質.

$^{601}$在聖經裡,或者整本聖經裡你會覺得很有趣.
我們經常覺得對質,思辨是一些很不禮貌.
很不客氣的舉動.
甚至上帝你騙我,或者這些是很冒犯上帝的句語.
不過在聖經裡,整本聖經從來沒有提到人對上帝的對質成為了上帝批判的對象.
反而整本聖經裡最批判的就是.
你表面是很順服,但裡面是和上帝抗拒的.
你離棄上帝,那種雙面的人反而是被聖經最批判的對象.
真誠找上帝對質的,反而上帝採取兼容和和這些人進行對話.
在第十九至二十一節就是這一段經文最後的出路.
「耶和華如此說:你若歸回,我就將你再帶來,使你站在我面前.
你若將寶貴的和下賤的分辨出來,你就可以當作我的口.
他們必歸向你,你卻不可歸向他們.
我必使你向這百姓成為堅固的銅牆.
他們必攻擊你,卻不能勝你.
因我與你同在,要拯救你,答救你,這是耶和華說的.
我必答救你脫離惡人的手,救贖你脫離強暴人的手」.
耶和華上帝如何回答耶利米的爭辯呢?.
其實上帝的答案不是什麼很大的新意.
大部分尤其是從第二十至二十一節.
基本上是copy and paste.
copy and paste是什麼呢?.
剛才讀過給大家聽.
耶利米書第一章六至八節.
我再讀給大家聽.
「耶和華,我不知怎麼說,因為我是年幼的,耶和華對我說:.
你不要說你是年幼的,因我差遣你到誰那裡去,你都要去.
我吩咐你說什麼話,你都要說.
你不要懼怕他們,因我與你同在,要拯救你,這是耶和華說的」.
一章十八至十九節,再繼續的來到他講.
我再讀給大家聽.
「汗啊,我今日使你成為堅城鐵柱銅牆,與全地和猶大的君王,首領,祭司並地上的眾民反對.
他們要攻擊你,卻不能勝你,因為我與你同在,要拯救你,這是耶和華說的」.
上帝將他自己昔日和耶利米所講的這番說話.
某程度上複製和複製重要字眼,放在這裡,你都可以回應耶利米.
不過在這裡所講的內容裡,他加了一個很重要的字眼.
「歸回」.
耶和華如此說,你若歸回.
「歸回」這個字是一字多而有不同的譯法.
接下來,再帶來是同一個字.

$^{641}$接著到十九節最後面.
你若將寶貴和下賤的分別出來,你就可以當作我的口.
他們必歸向你,你卻不可歸向他們.
在同樣的字眼裡,同樣的呼召.
不過面對著耶利米那些嘰嘰咕咕的那種自憐,那種有憤怒,痛苦的感受.
上帝的回覆是「歸回」,或者是悔改.
「去回」你的照明當中.
很奇怪,這是一個過去的語言.
整件事是在講上帝對耶利米的出路.
不是在講你計劃一下前面要怎麼走,感受一下當下你的感覺.
不是,是回到你的初心.
昔日上帝如何呼召他,他如何領受.
回到這些日子裡,重新站立起昔日上帝如何呼召他.
他如何回應,用這個呼召來應對當下生命裡的這些嘰嘰咕咕.
是真的,其實耶利米是一個衷心的僕人.
是一個痛苦悲慘的僕人.
到他所講的話實現了,整個以色列國真的被擄巴比倫了.
他所講的上帝的審判來到了.
到最後他依然實踐著他所講的照明.
耶利米到最後他選擇當別人的外邦王當他禮遇他.
你回來我給你做個宰相.
你這樣幫我潰敗了整個以色列的軍心.
你簡直是我們國家的恩人,巴比倫禮遇他的.
不過他說我不會跟你在一起.
我效忠的對象是上帝.
耶利米到最後他還是回到耶路撒冷.
跟那些剩餘的聖民在那個頹垣敗瓦的首都裡面一同渡過餘生.
他回到那個照明那裡.
回到那天昔日年輕的時候他覺得自己沒有可能.
但是上帝呼召他找回那個昔日的初心.
這個初心正正是守護著他的生命.
不會落入一種自憐或是萬劫不復的境地.
這個忠心的僕人在生命很多困難當中.
有時候會異化了,改變了.
好像成為了怪獸一樣.
人也會有的.
我頂持著某些信念,我繼續的來到一行.
但面對著太多的困難,太多的不理解或是排斥.
人的性格會扭曲的.
不過當你知道自己扭曲的時候.

$^{681}$找回自己,又對照今天.
為什麼今天成為了怪獸,成為一個這麼奇怪的人.
上帝叫他歸回,當你歸回到這個照明的時候.
你就可以平衡,你再看回你自己的生命.
不要再繼續落入一種這樣的境地.
上帝說你要成為一個同場.
成為一個沒有人可以攻擊你.
不能夠勝利的一個很堅強的人.
是真的,這種的美德,這種的生命.
就是在力排眾議,在逆流而上裡面孕育出來.
同場不會是自動起一幅銅牆.
然後人們受不了你,就等於同場.
而是在很多的艱難,很多的掙扎當中.
依然站立得穩.
依然生命裡面你都堅持著你要所行的信念.
就是因為這樣你的信念才有價值.
就是因為這樣實踐信仰.
所以你的生命才成為一個同場.
同場不是在順風順水,眾人愛戴之下建立出來的.
而是由很多困難,很多排斥之下.
那種的硬度是由這樣的方式在逆境裡面磨練出來.
回到那個呼召裡面.
其實上帝已經說了.
上帝一直保守他的心,上帝一直與他同在.
丁姐妹,我不知道今天聽完這篇訊息.
這一段對話,你心裡面有什麼體會.
我自己都問自己,我的初心是什麼呢?.
我的奉獻,我們的呼召,我們的人生.
我們為了什麼呢?.
當我在讀的時候,我心裡面覺得很溫暖.
我找回自己.
有時候人生去到某些時候會容易迷失.
有一些令你很羨慕的東西會令你走差.
但是回歸你的照明當中的時候.
你會發現你的生命可以不再一樣.
又或者如果我的生命歸化在這個世界裡面.
這個世界所同化,我們就失去了這個信仰.
失去了呼召對這個世界真真有詞的影響力.
丁姐妹,在你的生命當中,你有堅持嗎?.
你的信念是什麼呢?.

$^{721}$今天耶利米蘇第十五章.
喚醒我們自己.
可能我們有時候會覺得自己可憐.
覺得自憐的這種經歷.
可能就是你自己成長的契機.
在自憐當中,可能找回你自己生命裡面.
有一些價值,有一些信念是你所堅持.
找回這個信念來好好地過日子.
今天這個世代,需要有信念的人實踐真理.
希望你和我都成為一個有信念的人.
\newpage



\section{詩篇使徒行傳 2:1-9}
\label{sec:limrMw42910}
\textbf{2024年11月3日崇拜直播｜余嘉敏傳道｜詩篇裡的基督－受膏僕人｜詩篇2\_1-9;使徒行傳4\_23-31}
\newline
\newline
連結: \href{https://youtube.com/watch?v=limrMw42910}{\texttt{https://youtube.com/watch?v=limrMw42910}} ~~~~ 語音日期: 2024-11-04
\newline
\newline
\hyperref[sec:hd2ElWc7xho]{\small{< < < PREV SERMON < < <}}
~
\hyperref[sec:index_chronic]{\small{[返順時目]}}
~
\hyperref[sec:uf_k9SIAd7s]{\small{> > > NEXT SERMON > > >}}
\newline
\newline
$^{1}$以下要讀出的聖經有兩段.
第一段是在舊約詩篇第二篇.
第一至第九節.
在大家手上的聖經677頁.
第二段在《仕途行傳》第四章23至31節.
《新約》《仕途行傳》在新約166至167頁.
詩篇第二篇第一節.
「愛邦為甚麼爭鬧.
萬民為甚麼謀算虛妄的事.
世上的君王一齊起來.
神宰一同相議.
要敵黨耶和華並他的受告者.
說:我們要掙開他們的裙袍.
脫去他們的繩索.
那坐在天上的必發笑.
主必癡笑他們.
那時他要在奴中責備他們.
在列奴中驚嚇他們.
說:我已經立我的君.
在石安我的聖山上了.
受告者說:我要傳聖旨.
耶和華曾對我說:.
你是我的兒子,我今日生你.
你求我,我就將列國賜你為基業.
將地極賜你為前產.
你必用鐵掌打破他們.
你必將他們如同搖杖的瓦器摔碎」.
第二段是新約《使徒行傳》.
第四章二十三至三十一節.
在新約聖經一百六十六頁.
二十三節.
「二人既被釋放,就到會遊那裡去.
把子子上和長老所說的話都告訴他們.
他們聽見了,就都同心合意地高聲向神說.
主啊,你是造天地,海和其中萬物的.
你曾藉著聖靈託你僕人.
我們祖宗大位的後恕.
外邦為甚麼爭鬧.
萬民為甚麼謀算虛妄的事.
世上的君王一齊起來.

$^{41}$神在也聚集,要敵黨主.
並主的受高者.
希律和本雕彼拉多.
外邦人和以色列民.
果然在這城裡聚集.
要攻打你所高的聖僕耶穌.
成就你守和你兒子所預定必有的事.
他們恐嚇我們.
現在求主鑒察.
一面叫你僕人大放膽量講你的道.
一面伸出你的手來醫治疾病.
並且使神跡其事.
因著你聖僕耶穌的名行出來.
31 字.
禱告完了,聚會的地方震動.
他們就都被聖靈充滿.
放膽講論神的道.
.
弟兄姊妹,早安.
今天時間過得很快.
已經是十一月的第一個星期.
很快我們要紀念耶穌基督.
受苦,不是受苦.
廣新的大日子.
我們開始一個新的廣道系列.
叫做《詩篇裡的基督》.
當中我們會有四個星期.
以舊約的詩篇和新約的經文表述.
讓弟兄姊妹更加立體.
去認識我們的耶穌基督.
今天我會嘗試.
用詩篇的第二篇.
和《使徒行傳》的第四篇.
去反思耶穌基督.
是一個受高的僕人.
這個身份和我們有什麼關係呢?.
我們先一起有個禱告.
親愛的天父.
也留下你的大使命.
讓我們在你的福音上有份.

$^{81}$求聖靈你高滅我的口.
讓我能夠說得清楚.
也讓我所說的話都蒙你喜悅.
也聖靈你預備我們的心.
去被你的道去感動,引導.
多謝你聽我們禱告.
奉主耶穌基督的聖名祈求.
阿們.
幾年前我去過法國的梵爾賽宮.
欣賞了很多博物館.
黃色裡的英明神武.
(先不要太快).
(對不起).
英明神武的人像.
還有很多描述了.
軍王海船龜打仗的輝煌場面.
畫不是像現在的畫廊.
是很高的,看起來很宏偉.
其中一張大家應該都認識.
就是《拿破崙穿越阿爾卑斯山》的油畫.
當中描述拿破崙第一次帶領軍隊.
穿越阿爾卑斯山進入意大利的場景.
是不是很英明神武?.
當時我聽到引導的聲音.
原來這油畫跟現實有點不相符.
首先我們知道拿破崙的身高是怎樣的?.
這張照片很英明神武,很高大衛猛.
而且他原來不是騎馬,是騎驢.
所以有一個比較寫實的畫家.
創作了另外這張油畫.
相比之下是否截然不同?.
其實弄得這麼英明神武.
無非是想營造一個令人景仰.
令人可以臣服他的君王形象.
想清楚一點,是真的.
一個授權被加冕成為君王的領袖.
當然有既定的形象.
如果像這幅畫般,騎著驢仔進城.
被外族攻打,被人笑,被人敵擋.
相信人民看見這幅畫.

$^{121}$會否想去效忠他?.
應該不太想.
《聖經》也有提及.
原來被授權,加冕成為君王.
祭司和先知的.
他們都被稱為授高者.
意思是被神揀選.
選立執行神的任務.
猶太人一提到授高者.
就一定會想起這篇詩篇.
就是詩篇第二章.
當中提及授高者.
並非像剛才那幅畫.
騎馬的景象.
不是萬人景仰,萬邦臣服.
究竟這授高者是誰?.
我們一起看詩篇第二篇.
看經文前,我想交代背景.
詩篇第二篇是一篇公認的皇室詩.
描述上帝對以色列君王的應許.
通常在君王加冕禮時.
他們會採用這篇詩歌.
這篇詩歌亦會在當時的會堂.
用來敬拜和教導時中讀.
有學者認為這作是.
大衛時代至所羅門時代的作品.
但我們看新約,《仕途行傳》第四章.
使徒保羅,不,使徒彼得在祈禱中.
引述了這篇詩篇第二篇的一至二節.
他說:「神藉著聖靈.
託他僕人大衛的口所講述」.
所以今天我們嘗試用這篇.
彼得的這番說話.
貫穿舊約大衛的詩篇.
去理解今天的經文.
這篇詩篇,我一邊看覺得很有趣.
好像在看一部短劇.
當中有幾個角色在對話.
第一至第三節是作者反問叛逆者的說話.
第四至六節是耶和華的反應.

$^{161}$第七至九節是受高的王所引述的說話.
我邀請大家代入短劇當中.
嘗試投入角色,讓大家更深入經文.
我們一起看第一至第三節.
我們看到一個景象.
耶和華和受高者處於被人攻擊的局面.
敵人是誰?.
外邦軍隊,萬民,世上的君王和神宰.
他們要聯合在同一個陣營.
去反抗神和受高者.
第三節說,叛逆者在喧鬧.
我們要掙開他們的束縛,脫他們的繩索.
束縛的意思是「阿」.
大家可能熟悉經文.
「阿」是把他們困在繩索上的意思.
繩索是約束.
意思是說,這群叛逆者很想獨立自主.
很想擺脫神的命令和神的律法.
不想再受束縛.
神和受高者好像處於劣勢.
外面的陣營好像很強大.
但很有趣的是.
第一位出場的就是第一身作者大衛.
他竟然用諷刺的口吻反問這群人.
「外邦為何爭鬥?萬民為何圖謀虛妄的事?」.
他認為他們所做的事是虛妄的,是徒然的.
他沒辦法相信,哇,堂堂大國.
為何可以做一些注定失敗的事呢?.
第四至六節,鏡頭一轉.
因為鏡頭拍到天上.
我們就看到天上的就是神,耶和華.
祂很冷靜.
祂一點也不擔心列國他們的計劃.
也沒有被外邦的喧鬧影響.
第四節,神必譏笑,主必癡笑他們.
甚至說要烈怒去責備他們.
驚嚇他們,宣告.
「我已經立我的軍在石安我的聖山上」.
留意,神所高立的軍王.
神原來早就命定了.

$^{201}$儘管這群叛逆者聯手攻擊和攻打神.
神原來計劃已經做好了.
神已經高立了祂的聖山.
在石安山上.
石安山是甚麼地方呢?.
原來石安山有大衛的城.
在舊約《撒母耳記》第五章.
大衛被索羅追殺.
索羅死後,在希伯倫的地方.
在耶和華面前和眾支派立約.
大衛被告為以色列王.
第六至七節,大衛王去耶路撒冷時.
攻打耶布斯人.
最後,大衛王攻取了石安的堡壘.
石安的城是大衛的城.
從歷史角度去看.
神的確在石安山上.
高立了大衛作為以色列的王.
建立了他的城.
但是否一帆風順呢?.
也不是,我們知道大衛一生都在戰爭中.
愛邦不斷攻打他.
大家想起拋棄的民族.
非利士人,摩訶人,索巴王,雅蘭人等等.
一個戰士接一個.
打仗好玩嗎?.
不好玩,很艱難.
但是大衛所經歷的.
每一次他打仗的時候.
有愛邦爭奪他的時候.
耶和華都率領他.
讓他能夠抵抗戰勝那些愛邦人.
沒錯,愛邦的喧鬧.
敵黨其實早已是計劃當中.
原來有一個比他們勢頭更強.
更加超然的上帝在他們的上面.
掌管整個勢力.
無論他們做什麼.
這群人都無法抵擋神和受高者.
神的計劃最終都會得勝.

$^{241}$這就是第二篇詩篇告訴我們一個很重要的訊息.
所以天上的神見到這樣的狀況.
你們這群人在喧鬧.
神真的沒辦法,只能發出譏笑.
彷彿在跟這群叛逆者說.
你們知不知道你們在抵擋的受高者是誰?.
這一切都在我計劃當中.
你們只會白費心機.
究竟這個受高者是誰呢?.
我們看第七至第九節.
這裡提到神高立的君王終於出場.
他宣告神的聖旨.
他就說:耶和華曾對我說:.
你是我的兒子,我今天生你.
聖旨宣告了一個很重要的事.
就是王的身份.
他是神的兒子.
其實在舊約我們看神和人立約.
是父子之間很親密的盟約.
神是父親.
以色列成為了神的兒子.
父親在舊約中會有自己的產業和福氣.
在臨終時祝福他的孩子.
賜給他的兒子.
所以我們從歷史角度看.
神和大衛王納了一個永遠的盟約.
神賜列國給大衛成為產業.
並且把地極的所有地方都賜給他成為田產.
第九節說:受高者必定用鐵杖打破他們.
那些外邦的叛逆者.
如同搖杖的瓦器擊碎.
其實這是一個政治用語.
近東有一個傳統.
就是把瓦器上寫上芙蓉國的名字.
當鐵杖打破瓦器時.
代表宗主國對芙蓉國有絕對的主權.
但如果我們看「希伯來文」這個字.
打破這個字.
其實也有另一個意思.
就是管教和牧養.

$^{281}$即是告訴我們.
原來受高者有能力.
有主權去牧養和管教外邦的叛逆者.
想表達甚麼呢?.
想表達受高者是超然的.
面對外邦的敵黨時.
受高者必然會得勝.
正如我一開始所說.
這篇詩篇是公認的黃室詩.
而且會堂也經常中毒.
所以猶太人應該很熟悉這篇詩篇.
但回看歷史.
他們其實心知肚明.
大衛沒有幫助他們復國.
這個盼望好像幻滅了.
猶太人的拉比更加認為.
這篇詩篇所指向的.
將來會來到受高者彌賽亞.
他將會登基.
幫他們擊打.
打敗外族.
以色列就會怎樣?.
以色列國復國.
然後得享神永遠的福氣.
其實那時新約的時代.
耶穌在人民面前.
見到很多人.
傳福音.
很多都被人看到.
耶穌是那位彌賽亞.
神的兒子.
我們回看剛才說的第七節.
「你是我的兒子,我今日生你」.
是不是有點熟口熟面?.
原來在施洗約翰.
為耶穌施洗的時候.
天上有聲音說.
「這是我的愛子,我所喜悅的」.
他很清楚地宣告了一件事.
耶穌就是神的兒子.

$^{321}$而且早在耶穌出生的時候.
天使對瑪利亞說.
「你將會懷孕生子」.
「要給他起名叫耶穌」.
「他將要為大」.
「稱為至高者的兒子」.
甚至在登山變相的時候.
其實也有說過.
也有清晰地說.
耶穌是誰.
耶穌是神的兒子.
而且他會在亞國家作王.
直到永遠.
你可能會說.
那些叛逆者.
那些祭司長.
他們都沒有看過新約.
很後期才寫的.
他們沒有看過.
但願意跟從耶穌的基督徒.
他們有沒有看過?.
他們都沒有看過.
他們聽聞耶穌是一個怎樣的人.
聽聽耶穌所說的話.
然後他們就相信.
只不過那些叛逆者.
那些法利賽人祭司長.
他們期待著.
彌賽亞有既定的身份和魅力.
他們期待受高的君王.
能夠帶領他們.
威脅國家的政治君王.
相反.
耶穌.
看見祂騎著路仔進城.
然後又是無形無格的君王.
即使祂做了多少神蹟.
甚或祂所做的事.
是多能夠成就.
他們很熟悉的詩篇.

$^{361}$裡面所說的彌賽亞的特徵和預言.
他們都不願意去接受和相信.
耶穌就是受高者彌賽亞.
不單是這樣.
還要去攻擊他們.
攻擊耶穌.
要將他們釘在十字架上.
在《使徒行傳》記載.
耶穌復活之後.
這群黑眼的叛逆者.
仍然在逼迫那群耶穌的跟從者.
我們轉個鏡頭.
看那群跟從耶穌的人.
他們和那群叛逆者不同.
他們看見受高者就是耶穌基督.
讓他們有信心.
更加放膽地傳揚福音.
究竟他們看了些什麼.
我們一起去看.
如果你們有聖經.
可以去看《使徒行傳》.
新約的《使徒行傳》.
在《使徒行傳》裡.
講述這群門徒領受了聖靈之後.
他們就像效法耶穌一樣.
四處傳道.
在第三至第四章.
講述使徒彼得和約翰.
奉耶穌的名.
去醫治了一個.
由戰武福騎腿的四十歲人.
他又向百姓講述.
耶穌在死裡復活的福音.
令那群叛逆者聽完.
哇! 神上善掃飆升.
於是就去捉拿他們.
去拘留所拘留他們.
而且威嚇他們.
不可以再奉耶穌的名.
去傳揚福音.

$^{401}$第23節我們看到.
他們兩個最終被釋放.
於是他們回到會友.
可能是回到教會.
就和這群耶穌的跟從者.
講述那群祭司長長老所說的話.
告訴他們.
我們試試代入.
如果你是當時的會友或使徒.
你會有什麼想做?.
有什麼反應?.
你們會不會商議一下.
下一步怎麼做?.
究竟是逃走還是分散?.
還是談談有什麼方法.
傳得低調一點.
不要太高調.
很害怕吧?.
一定很害怕.
他們的威嚇是針對我們的生命.
是生命威脅.
他們剛剛跟從耶穌.
剛剛才被那群叛逆者.
迫到釘上十字架.
縱然我們知道會得勝.
但是他們在十字架上受到刑罰.
其實都很可怕.
他們的威嚇是很真真實實的.
甚至稍後我們在經文裡.
我們看到第十二章.
這群反叛徒是因為傳耶穌的道而殉道.
害怕嗎?.
不害怕就假了.
不過經文看到這群跟從耶穌的人.
他們沒有逃走.
他們也沒有逃避.
第二十四節我們看到.
他們竟然聽了之後.
是同心合意高聲向神祈禱.
禱告一開始.

$^{441}$他們就宣認.
認定神的身份.
他說:這就是創天造地萬有的上帝.
然後他們更加認定詩篇裡.
第二篇所說的受高者.
就是彌賽亞就是耶穌基督.
他說:神藉著聖靈去啟示.
在至少一千年前.
即是出生一千年前.
大衛估計登基的時候.
應該是在主前一千四十五年左右.
所以當時的使徒就說.
這個預言是至少一千年前.
原來藉著祂們的祖先大衛的口中說過.
而且他們更加肯定.
這篇詩篇是說耶穌的原因.
就是耶穌在升天之前.
曾經和這群門徒.
留下大使命給他們的時候.
跟他們說詩篇所寫的.
凡指著我的話.
都必須應驗.
不單止這樣.
他們親眼看到耶穌.
釘十字架.
在十字架上被擊打.
這些片段和畫面.
其實跟詩篇裡說的受高者.
有很多相似的地方.
所以他們看到的.
就是神的計劃.
第25至26節說到.
使徒彼得記起.
並且引用詩篇第二篇的經文.
我完全看到經文的時候.
我感受到他的興奮.
他彷彿在說.
詩篇裡說的預言.
現在終於都發生應驗了.
耶穌基督在世的時候.

$^{481}$大家一起抵擋他.
如今我們跟從的門徒.
同樣都受到這些迫迫和打擊.
然而彼得這群使徒.
他們不是看見自己被迫迫.
很害怕.
而是他們看到上帝的恩典.
上帝的計劃.
他們看到一切被迫迫的背後.
有神在掌權.
因此他們很興奮和驚歎.
引用詩篇第二篇的經文.
第27節的禱告繼續說.
希律,本·迪拉多,愛邦人和以色列民.
果然聚集在城裡.
攻打受告聖母耶穌.
他們回想起這群叛逆者.
就像詩篇第二章一至二節所說.
原來劇本早已寫好.
這群受告者攻打耶穌.
因為耶穌君亡.
但這個受苦必然要發生.
原來劇本亦寫明.
耶穌甘願受苦被人唾棄.
甘願取奴僕的樣式.
走上十字架.
成就了救恩.
成為了我們眾人的彌賽亞.
這個君王就不像.
我們剛才說的那部《拿破崙》的複名話.
孤大衛猛騎著駿馬進城.
去作世界的王.
而是很謙卑.
作萬人的僕人耶穌.
這完全顛覆了我們.
所有人對君王的形象.
彼得和約翰看見上帝的計劃.
他更加肯定.
今天他自己和那群跟從耶穌的人.
所受的迫擊.

$^{521}$是必定會發生的.
這些全部都是為了迫擊耶穌.
才會發生的事.
於是他們放膽向神祈禱.
第29至31節.
第29節他們說.
他們恐嚇我們.
現在我們求主鑒賊.
一面叫你的僕人放膽.
放大膽量去講道.
一面伸出你的手.
去醫治疾病.
並且使神跡騎士.
因著你的聖僕耶穌的名字.
能夠走出來.
他們祈禱完後.
聚會的地方就震動了.
地震一樣.
之後他們被聖靈充滿.
放膽去講神的道.
不要忘記.
其實耶穌離開了這群門徒之後.
這是門徒第一次經歷迫擊.
當他們越大膽去講道.
大膽去傳揚福音的時候.
你們會看到迫擊將會怎樣?.
會升溫.
甚至有生命的危險.
但這只是開始.
但聖靈很好.
祂藉著這個開始之先.
讓他們看見這詩篇第二篇所預言的.
一一都應驗在耶穌的身上.
而這份看見也讓他們堅信.
耶穌基督就是真正的彌賽亞.
他們堅信一切的迫擊是必然發生的.
上帝的計劃是一定得勝.
而他們竟然有份參與並見證.
這個千年以前已經預定好的救贖計劃.
難怪他們會這麼激動.

$^{561}$這麼振奮.
充滿信心.
上帝一定會得勝的.
現在他們受的苦楚.
原來是有意義的.
這就成為了仕途.
這群跟從者的動力.
他們要繼續傳揚耶穌.
見證耶穌.
更加放膽去傳揚耶穌的道.
因為唯有祂才是拯救萬民的君王.
由我這裡有個小小的總結.
世界的敵黨耶穌是必然會發生的事.
跟從耶穌的人同樣都會受到逼迫.
但是受高者必然得勝.
跟從他的也必然得勝.
弟兄姊妹.
今天我們做基督徒比較輕散.
我們沒有像昔日的門徒一樣.
會受到很多的逼迫.
我們可以選擇.
信了主之後.
我們可以選擇過得安逸.
我上個月去了英國旅行.
我參加了我朋友的一個主日崇拜.
剛好那天是福音主日.
他邀請了一個外來的港元來.
他講了福音的報道訊息.
那個港元牧師.
他在那天就分享了.
原來跟從耶穌.
是要預備好作主的門徒.
走上受苦逼迫的責勞.
他邀請了未信主的人.
就決志.
但是相信了主的.
願意跟從耶穌的人.
他就問我們.
你願不願意受苦為主走這條責勞.
你猜當天有沒有人舉手決志?.

$^{601}$我現在問你們.
你們都心大心細.
很可惜沒有人舉手決志信主.
甚至我自己也有些猶豫.
我回去反覆思量.
問題出在哪裡呢?.
我想現實沒有人喜歡受苦.
沒有人喜歡被人逼迫.
或者走上要死的路.
如果我們單單看著這些苦.
而我們看不到盼望.
肯定沒有人會這樣做.
這條責勞不要搞了.
問題是基督教的信仰.
不只是單單講受苦.
如果我們忽略了.
神最終必得救的福音.
我們只會停留在灰心失望的當中.
我們沒有辦法向前行去見證上帝.
我想跟大家看一條短片.
這條短片講述耶穌基督.
基督教自從耶穌降生.
被傳到碧白.
甚至是釘上十字架之後.
因為這條片是2015年製作的.
直到二千多年前.
這段歷史.
究竟基督教的發展是怎樣.
我想大家留意.
白色的部分是基督教.
其他顏色是其他宗教.
我們一起看這條片.
基督教的信仰的確受到很多迫迫.
甚至有些階段.
白色的部分收窄了.
好像快要失去宗教信仰.
但我們最終看到.
歷史告訴我們.
白色的部分不斷擴展.
甚至到2015年.

$^{641}$基督教差不多到達地極.
其實跟從耶穌.
這個謙卑的受苦的僕人.
的確會受到很多迫迫.
不太受歡迎.
但聖經告訴我們.
耶穌基督是忠必得勝.
你見到嗎?.
你相信嗎?.
你的生活裡有沒有一些外在的環境.
都有很多勢力.
勢頭很強大的外力.
讓你無法宣認耶穌基督是王.
無法去傳揚耶穌基督的道.
是不是你的工作裡.
有太多同事是非基督徒.
他們的眼光,嘴臉.
甚至上司都要擺放風水陣.
老闆都不喜歡.
我怎能夠抵擋他們呢?.
我怎能夠反抗呢?.
還是你的家人,朋友圈.
他們個個都不喜歡基督教.
我講了那麼多.
我怎能夠抵抗他們呢?.
我怎能夠制衡這些力量呢?.
我怎能夠去傳揚福音呢?.
其實真的很艱難.
很艱難.
我自己都覺得很艱難.
我們可能會想.
停在我相信了就行了.
耶穌基督已經為我們.
這個僕人已經為我們成就了救恩.
我們只要去相信祂.
認定祂是我們的救主.
我們就能夠救了.
不用做那麼多事.
那就夠了.
不過使徒告訴我們.

$^{681}$祂沒有停在這裡.
祂沒有停在自己個人得救的層面.
而是去效法受高的僕人耶穌基督.
放下自己的權利.
去傳揚神的福音.
好,叫到這個好的訊息.
這個福音能夠傳到萬國萬邦.
只要我們願意.
我和你都一定有份.
參與這個早在創世之先.
已經預定好的救贖計劃.
去見證神最終必定會得勝的大能.
最後我想用一首詩歌.
作回應詩去作結束.
我都想但願望浸這個群體.
可以同心合意.
向神同唱.
主我願意.
謝謝.
\newpage



\section{約翰福音 6:60-7:13}
\label{sec:uf_k9SIAd7s}
\textbf{2024年11月17日崇拜直播｜梁家麟牧師｜被人厭棄的耶穌｜約翰福音6\_60-7\_13}
\newline
\newline
連結: \href{https://youtube.com/watch?v=uf_k9SIAd7s}{\texttt{https://youtube.com/watch?v=uf\_k9SIAd7s}} ~~~~ 語音日期: 2024-11-20
\newline
\newline
\hyperref[sec:limrMw42910]{\small{< < < PREV SERMON < < <}}
~
\hyperref[sec:index_chronic]{\small{[返順時目]}}
~
\hyperref[sec:230hZbWKaH0]{\small{> > > NEXT SERMON > > >}}
\newline
\newline
約翰福音 6:60-7:13
\newline
\begin{longtable}{cl}
\hline
\hline
章節 & 經文 (和合本修訂版)\\
\hline
6:60 & \begin{tabularx}{0.7\textwidth}{X} 他的門徒中有好些人聽見了，就說：「這話很難，誰聽得進呢？」 \end{tabularx} \\ \\ \relax
6:61 & \begin{tabularx}{0.7\textwidth}{X} 耶穌心裡知道門徒為這話私下議論，就對他們說：「這話成了你們的絆腳石嗎？ \end{tabularx} \\ \\ \relax
6:62 & \begin{tabularx}{0.7\textwidth}{X} 如果你們看見人子升到他原來所在之處，會怎麼樣呢？ \end{tabularx} \\ \\ \relax
6:63 & \begin{tabularx}{0.7\textwidth}{X} 聖靈賜人生命，肉體毫無用處。我對你們所說的話就是靈，就是生命。 \end{tabularx} \\ \\ \relax
6:64 & \begin{tabularx}{0.7\textwidth}{X} 可是你們中間有些人不信。」耶穌起初就知道哪些人不信他，哪一個要出賣他。 \end{tabularx} \\ \\ \relax
6:65 & \begin{tabularx}{0.7\textwidth}{X} 於是耶穌說：「所以，我對你們說過，若不是蒙我父的恩賜，沒有人能到我這裡來。」 \end{tabularx} \\ \\ \relax
6:66 & \begin{tabularx}{0.7\textwidth}{X} 從此，他門徒中有很多退卻了，不再和他同行。 \end{tabularx} \\ \\ \relax
6:67 & \begin{tabularx}{0.7\textwidth}{X} 耶穌就對那十二使徒說：「你們也要離開嗎？」 \end{tabularx} \\ \\ \relax
6:68 & \begin{tabularx}{0.7\textwidth}{X} 西門‧彼得回答他：「主啊，你有永生之道，我們還跟從誰呢？ \end{tabularx} \\ \\ \relax
6:69 & \begin{tabularx}{0.7\textwidth}{X} 我們已經信了，又知道你是神的聖者。」 \end{tabularx} \\ \\ \relax
6:70 & \begin{tabularx}{0.7\textwidth}{X} 耶穌回答他們：「我不是揀選了你們十二個嗎？但你們中間有一個是魔鬼。」 \end{tabularx} \\ \\ \relax
6:71 & \begin{tabularx}{0.7\textwidth}{X} 耶穌這話是指著要出賣他的加略人西門的兒子猶大說的；他本是十二使徒裡的一個。 \end{tabularx} \\ \\ \relax
7:1 & \begin{tabularx}{0.7\textwidth}{X} 這些事以後，耶穌周遊加利利，不願在猶太往來，因為猶太人想要殺他。 \end{tabularx} \\ \\ \relax
7:2 & \begin{tabularx}{0.7\textwidth}{X} 這時猶太人的住棚節近了。 \end{tabularx} \\ \\ \relax
7:3 & \begin{tabularx}{0.7\textwidth}{X} 耶穌的兄弟們對他說：「你離開這裡上猶太去吧，好讓你的門徒也看見你所做的事。 \end{tabularx} \\ \\ \relax
7:4 & \begin{tabularx}{0.7\textwidth}{X} 因為人要揚名，沒有在隱祕的地方行事的，如果你要做這些事，該把自己顯明給世人看。」 \end{tabularx} \\ \\ \relax
7:5 & \begin{tabularx}{0.7\textwidth}{X} 原來連他的兄弟們也不信他。 \end{tabularx} \\ \\ \relax
7:6 & \begin{tabularx}{0.7\textwidth}{X} 於是耶穌對他們說：「我的時機還沒有到，你們的時機卻隨時都有。 \end{tabularx} \\ \\ \relax
7:7 & \begin{tabularx}{0.7\textwidth}{X} 世人不會恨你們，卻是恨我，因為我指證他們的行為是惡的。 \end{tabularx} \\ \\ \relax
7:8 & \begin{tabularx}{0.7\textwidth}{X} 你們上去過節吧！我現在不上去過這節，因為我的時機還沒有成熟。」 \end{tabularx} \\ \\ \relax
7:9 & \begin{tabularx}{0.7\textwidth}{X} 耶穌說了這些話，仍然留在加利利。 \end{tabularx} \\ \\ \relax
7:10 & \begin{tabularx}{0.7\textwidth}{X} 但他的兄弟們上去過節以後，他也上去，不是公開去，卻似乎是祕密地去的。 \end{tabularx} \\ \\ \relax
7:11 & \begin{tabularx}{0.7\textwidth}{X} 節期間，猶太人尋找耶穌，說：「他在哪裡？」 \end{tabularx} \\ \\ \relax
7:12 & \begin{tabularx}{0.7\textwidth}{X} 人群中有許多人對他議論紛紛，另有的說：「他是好人。」有的說：「不，他是迷惑群眾的。」 \end{tabularx} \\ \\ \relax
7:13 & \begin{tabularx}{0.7\textwidth}{X} 可是沒有人公開談論他，因為他們怕猶太人。 \end{tabularx} \\ \\
[1ex]
\hline
\hline
\end{longtable}
$^{1}$以下是讀經.
今日的經文是.
約翰福音六章六十節到七章十三節.
經文已印在你們的秩序表內.
麻煩你們打開.
我們以啟認的方式讀.
首先由我讀.
然後由會眾讀下一節.
約翰福音六章六十節.
請聽我讀第一節.
他定門途中有好些人聽見了.
就說:這話甚難,誰能聽呢?.
約翰福音六十節到七章十三節.
就說:這話甚難,誰能聽呢?.
倘或你們看見人子.
升到他原來所在之處.
怎麼樣呢?.
約翰福音六十節到七章十三節.
如啟示無疑的.
當你們所說的話.
就像是真實的命.
只是你們中間有不信的人.
耶穌從起頭就知道.
誰不信他.
誰要賣他.
要做人子.
所以當你們說過.
若有真信和不可悔的因子.
會有人能到往者的路.
從此.
他門途中多有退去的.
不再和他同行.
(耶穌說:你們不選擇門徒).
西門彼得回答說.
主啊!你有永生之道.
我們還歸從誰呢?.
(我們已經信了).
(就知道你是神的使徒者).
耶穌說:我不是選擇了.
你們十二個門徒嗎?.

$^{41}$但你們中間有一個是魔鬼.
(耶穌說:是主啊!在西門彼得的前面).
(祂把祂的十二個門徒的一切).
(向後來要賣給耶穌).
這事以後.
耶穌在加里里遊行.
不願在猶太遊行.
因為猶太人想要殺祂.
(當時猶太人的住棚節).
耶穌的弟兄對祂說.
「你離開的住棚節」.
「你離開這裡上猶太去吧」.
「叫你的門徒也看見你所行的事」.
(人們要以為人們是聖).
(沒有了主的行事之力).
(你們要把行事之事).
(就當作主的言論和世人說).
因為連他的弟兄說者說.
是因為不信他.
(耶穌說:我對你們的行事).
(我的使徒們沒有任何).
(理會的事後嘗試說不定).
世人不能恨你們.
卻是恨我.
因為我指正他們所作的是善惡的.
(你們要上猶太節).
(我們要再上猶太節).
(因為你們要嘗試).
(和我辦理任何的事).
耶穌說了這話.
「仍舊住在加里里」.
(但當你們上猶太以後).
(大家上猶太節).
(不信你的行事).
(才會嘗試說不定).
正在節期.
猶太人尋找耶穌.
說:他在哪裡?.
(所有人的家庭紛紛亂端).
(有些事發生無人).

$^{81}$(有些事不見發生無人).
只是沒有人明明的講論他.
因為怕猶太人.
聖經獨筆.
頂師傅早安.
當耶穌說完.
他是生命的糧.
他的肉可以吃.
他的血可以喝.
吃他的肉喝他的血的人才得到永生.
這個這麼噁心.
這麼奇怪的道理之後.
立即招來極大負面的反應.
不僅猶太教的領袖要殺他.
連一直追隨他的門徒.
群體也產生不少的震動.
有些人因此就離開了他.
還有連耶穌的家人.
就是他的弟弟.
也表達對他的不信任.
可以看到一石激起千重浪.
耶穌有些處於四面楚歌.
處處不討好的困境裡面.
人對耶穌的態度.
當然取決於耶穌說了什麼.
做了什麼.
如果耶穌說了他們喜歡聽的話.
例如他們會得到上帝賜福.
他們會心想事成萬事勝意.
他們當然樂意接受.
如果耶穌做了他們期待的事.
例如顯神蹟.
醫病趕鬼.
特別是有飯吃.
五餅二魚的神蹟.
他們當然會為之發狂.
蜂擁而至.
貼身跟隨.
但如果耶穌說了一些.
害人聽聞.

$^{121}$離經叛道的話.
顛覆了他們一貫的認知和價值觀.
又或者.
耶穌刻意挑起.
跟猶太教當局.
或者羅馬政權一些對立的事端.
多數小老百姓.
就會雞咁腳就走.
正如在這一次.
在五餅二魚的神蹟之後.
耶穌明明是賺取了大量的粉絲.
很多人到處去找他.
要聽他的道.
要藉著他.
一再吃餅得飽.
從追隨者的數量來說.
這是耶穌全道生涯的一個高峰.
不過耶穌不知道.
真是視何居心.
竟然就說吃肉喝血這樣的道理.
造成群眾極大的不安.
這是名副其實的反高潮.
用今天的說法.
耶穌是在趕客的.
這樣的情況就好像美國總統候選人.
在接受電視訪問之後.
因為表現不好.
說了一些不對的東西.
就造成搖擺州份.
大量選票流失一樣.
吃肉喝血的說話.
固然是非常惹人反感的.
即使在今天.
吃人肉還是令人作嘔的事.
對於嚴禁吃血的猶太人來說.
耶穌呼籲人喝他的血.
這是跟律法直接對上的.
真是明挑律法的機.
不過對猶太人的領袖來說.
最反感的也不是吃肉喝血.

$^{161}$而是耶穌宣稱.
上帝是他的天父.
他隱隱將自己和上帝的地位.
看為是平等.
又說他是從天上來的.
只有他明白天父的心意.
這些說話是非常褻瀆神明的.
大逆不道.
欲與天公事比高.
這是極其狂妄的言論.
人間怎能容納如此的狂徒呢?.
所以他們真是想除耶穌而後快.
剛才我們讀了一段很長的聖經.
記載了三批人對耶穌的狂言反應.
第一批是廣泛的追隨者.
聖經稱他們為門徒.
不過跟接下來要說的十二門徒有分別.
不過他們既被稱為門徒.
顯然他們跟隨耶穌也有一段時間.
不是偶然來一次.
他們人數很多.
也是最有社會輿論的代表性.
他們對耶穌的說話第一個反應是甚麼呢?.
「這話甚難,誰能聽呢?」.
六章二十字.
「這話甚難」不是很難明白.
而是很難接受.
很難忍受.
是否?太具冒犯性.
整句話的意思翻成廣東話就是.
「耶穌說甚麼?語無倫次」.
不知道他發了甚麼神經.
「誰能聽呢?」.
就是受不了他.
基於群眾一貫都是比較含蓄膽怯.
他們沒有將對耶穌的不滿宣之於口.
只是在心底裡嘰嘰噤噤.
跟身邊的人嘔吐嘔吐.
不過耶穌是讀懂人心中的想法.
他沒有迴避.

$^{201}$假裝不知道.
避而不談.
他專挑黃蜂鬥.
將人心中的芥蒂.
將它明說.
將它挑出來.
簡單說就是要將衝突擴大.
耶穌直接問他們.
「這話是叫你們厭棄嗎?」.
這句話有兩個意思.
第一個是簡單自貶的意思.
「你們不喜歡聽我這番說話?」.
第二個是引伸的意思.
「這樣說你們已經受不了嗎?」.
「我還有更激烈的事」.
「你們聽了不是要去死?」.
耶穌接著說.
在6章22-24節.
「討厭你們看見人子」.
「升到他原來所在之處」.
「怎樣呢?」.
「既然活著的是靈」.
「肉體是無益的」.
「我對你們所說的話就是靈」.
「就是生命」.
「只你們中間有不信的人」.
耶穌再次強調.
祂是從天上來的.
說的是天外之音.
耶穌是說天籟.
這些天籟就是從天上來的聲音.
是你們這班凡夫俗子.
有限智慧的人.
無法消化的.
你們跟耶穌根本是不同的級別.
一個檔次.
他們對耶穌的說話.
唯一正確的態度.
不是企圖去明白.
你這麼小,明白什麼?.

$^{241}$更加不是要加上你們那些.
無謂的喜惡愛憎.
你喜不喜歡?.
什麼事?.
有什麼關係?.
你們只能做的是.
無條件地相信和接受.
你們沒有資格去理解.
沒有資格去表達你們個人的喜惡愛憎.
你只能夠全盤接受.
叫人活著的是靈.
肉體是無益的.
耶穌的意思不是叫人用靈異.
不用字面去理解祂的話.
也不是那些什麼三原論者.
說要用靈,不是用魂去理解耶穌.
而是乾脆指出.
作為凡夫俗子的他們.
你們無論如何都不會明白我所說的話.
靈異不行.
用靈和魂分開都不行.
除非是聖靈.
這裡說的靈是聖靈.
親自將你翻轉.
將你重生了.
你才明白我說的話.
叫人活著的是靈.
指的是聖靈.
耶穌說的是靈,是生命.
不是肉體的層次.
人能不能夠理解.
是不是忠義.
通通不相干.
請自量一下.
你們有能力明白嗎.
你們有資格說不忠義嗎.
在肉體商店.
你只能買到一塊肉路.
你怎會聽到屬靈的話.
要聽屬靈的話.

$^{281}$只能來旺站.
很明顯.
耶穌沒有打算安撫那些追隨者心中的不滿情緒.
反倒是惹是鬥非.
直接批評他們.
進一步激發他們的反感.
耶穌有說.
六章十五節.
所以我對你們說過.
若不是無我父的恩賜.
沒有人能到我這裡來.
相信耶穌本身就是一個神跡.
沒有神跡你信不了耶穌.
需要上帝親自作為.
只有那些被父上帝揀選.
用我們佛教的說法.
有慧眼有慧根的人.
才能夠明白耶穌的話.
才有資格跟隨耶穌.
耶穌不安慰不解釋.
還說那些人沒有慧根.
沒有資格去跟隨.
你說說話說到這個地步.
那群群眾那群門徒還有什麼選擇.
當然是走.
不走還有什麼可以做.
我都跟你分享過.
曾經有一次我在另一間教會.
我講完一篇我自己覺得很有型.
很有動見的一篇講道之後.
有個弟兄.
即是完了崇拜之後就來找我.
跟我說.
他說牧師.
我不知道我是否應該轉會.
剛才你講的那段道.
我一句都聽不明白.
我當時的反應是想.
找一面牆撞頭過去.
我立刻安慰鼓勵那個弟兄.

$^{321}$問題不是出在你身上.
是我剛才發台瘟.
亂說話.
所以不關你的事.
如果在那個時候我沒有自責.
我卻是充滿自意地說.
你沒有慧根慧眼.
你怎麼會聽得懂我說的話呢.
我的道根本不是為你預備的.
那你說他有什麼反應.
他不走?.
你有沒有搞錯?.
講到這個地步當然是走.
66節說.
從此他們途中.
多有退去的.
不再和他同行.
本來有幾千甚至近萬粉絲.
一下子絕大部分都跑了.
剩下小貓兩隻.
真的可以說是眾叛親離.
這裡講的是在加利利發生的故事.
耶穌一時間被很多人追隨.
很快被群眾凋棄.
我們知道同樣的故事.
發生過在耶路撒冷.
在耶穌進入聖城的時候.
是不是曾經有一大堆粉絲追隨他.
鋪路.
弄一個人造紅地毯給他.
賀山娜.
馮主名來的王是當受讚美的.
很快他們就加入了要釘死他的行列.
耶穌常常的故事都是大上大落的.
一時之間很受歡迎.
一時之間神憎鬼厭.
就好像有很多人在中學階段.
很積極很熱心參加教會活動.
大學畢業之後出來工作.
沒多久就慢慢離開教會.

$^{361}$因為人生有更多值得他們追求的東西.
大批人離開之後.
耶穌身邊只剩下孤零零.
小貓幾隻.
那十二門徒就顯得很突出.
耶穌就轉頭問他們.
你們是不是也想走.
十二門徒之首的彼得.
總是最勇於發言.
馬上就向耶穌表態.
六十八,六十九節.
主啊,你有永生之道.
我們還歸從誰呢.
我們已經信了.
又知道你是上帝的聖者.
在門徒中間.
彼得對耶穌的認識算是最多的.
他是第一個說出耶穌是基督的人.
他對耶穌的忠心也不需要懷疑.
彼得不一定明白耶穌說吃肉喝血的含義.
但就捉住前面耶穌說.
祂的話就是生命這一句.
然後就宣告.
你有永生之道.
對耶穌的話.
不用執著你不明白的細節.
我讀聖經的時候也不是句句明白.
捉住其中的要旨才是正理.
同樣,我也沒辦法明白.
耶穌在我生命裡的布局的每一個細節.
為什麼這樣搞,為什麼這樣發生.
你要我每一個細節都能解釋.
我解釋不了.
所以我自己也解釋不了.
你就不要經常找你的故事來問我為什麼.
我更加解釋不了.
很多東西我當然不明白.
我怎麼知道祂為什麼.
我又不是祂的發言人.
我真的不知道.

$^{401}$我又不祈求神,我怎麼會知道.
不明白就不明白.
不過,我們只需要確定.
跟隨祂的大方向就夠了.
彼得稱耶穌為上帝的聖者.
強調祂是彌賽亞的身份.
耶穌有神聖的來源,神聖的身份.
我們知道,如果你讀聖經的話.
這個神聖身份首先是撒旦說出來的.
你看馬福音一章二十四字就知道了.
然後才被門徒認出來.
彼得的認信帶給耶穌一點點的安慰.
畢竟跟其他追隨者不同.
甚至試圖是耶穌逐個逐個選出來的.
不是他們自告奮勇跟隨的.
所以,如果他們放棄相信.
不僅僅是他們自己失敗.
也是耶穌的失敗.
因為耶穌手握他們.
選成這樣,都失敗了.
你說慘不慘呢?.
不過,縱使彼得說了一些很肯定性的說話.
很鼓勵性的說話.
耶穌還是要潑冷水.
他不客氣這樣說.
即使在他親自選的十二個試圖中.
還是有一個失敗者.
有一個落敗者.
有一個會走的.
「我不是選了你們十二個門徒嗎?.
但你們中間有一個是魔鬼」.
魔鬼是一個比較籠統的翻譯.
原意是誣告者.
即是反對耶穌的人.
約翰在這裡立刻加插一句說話做解釋.
做一句旁白來解釋一下.
《習錄》71節.
「耶穌說,這話是指著加略人西門的兒子猶大說的.
他本是十二個門徒裡的一個.
後來要賣耶穌的」.

$^{441}$前一年,約翰有說過.
在《習錄》64節.
耶穌從起頭就知道誰不信他.
誰要賣他.
耶穌對所有跟隨他的人.
都是一清二楚,心知肚明.
他知道誰是真心.
並且長久地跟隨他.
誰是真正的門徒.
誰是一時頭腦發熱.
跟隨他,但又逃脫.
誰只是因誤會而結合.
因了解而分開.
誰甚至會在最後關頭捅耶穌一刀.
耶穌知道真相.
所以他不會因為很多人一時間跟隨他而很高興.
很多人又離開他,就很失望.
耶穌要的是真實的門徒.
就是那些毫不保留,相信他.
不離不棄,跟隨他的人.
我們在看到耶穌說了吃肉喝血的教訓之後.
猶太教的領袖就要置耶穌於死地.
大量的群眾離棄耶穌.
只有十二仕徒仍然堅定跟隨他.
約翰接著再交代.
耶穌身邊一個群體對他的反應.
那就是他的親人,他的弟弟.
七章一節說.
耶穌知道猶太教的領袖想要殺死他.
就盡量避免去南部的猶大耶路撒冷.
耶路撒冷是猶太教的總部.
他只留在北邊的加利利區.
因為那是公會邊場莫及的地方.
加利利不是耶路撒冷的權勢可以去到的地方.
比較安全.
不過作為一個守宗教規條的猶太人.
他總是需要在那三大節期前往耶路撒冷聖殿守節.
第二節就說.
「住彭節到了」.
也就是說又要去耶路撒冷守節的時候.

$^{481}$耶穌的弟弟,這不是說一個人.
因為後面說的是眾數,不是單數.
他們既不相信耶穌的神聖身份和地位.
又對耶穌經常說的那些奇奇怪怪的.
就是經世海俗的說話不以為然.
他們總是覺得耶穌是一個不安分.
經常做白日夢的人.
認定他那些怪異的言論和行動.
會為家庭帶來麻煩,是駝衰家的.
特別是你專門去和那些宗教權貴為敵.
和那些社會領袖為敵,有沒有好結果?.
你經常和那個姓李的去吵架.
大哥,我們會死的,對不對?.
對小老百姓的他們,這是極其不利的,對不對?.
耶穌的弟弟看見耶穌這次不打算去耶路撒冷.
哎呀,怎麼這麼陰春?.
就諷刺他說.
你離開這裡,上猶太去吧.
叫你的門徒也看見你所行的事.
人要顯揚明星,沒有在暗處行事的.
你如果行這事,就當將自己顯明給世人看.
這段話的意思是.
喂,耶穌,為什麼你不去耶路撒冷?.
連你的門徒也去了.
你很應該在那裡再顯明神蹟給你的門徒看.
還有,如果你好高騖遠,想博出名.
那你也應該去耶路撒冷這個更大的平台.
不要留在加利利這個窮鄉僻壤,是不是?.
就是怎樣?你經常說中國呼音化.
你在香港說這句話有什麼營養價值?.
你去人民大會堂門口派單張.
你去新華門派,派中南海.
擒賊先擒王,整個政治局信了耶穌.
就是中國呼音化.
在香港說這些無謂多餘的東西.
簡單說,要認什麼就認什麼.
請你走遠一點,去耶路撒冷吧.
不要在這些地方癢癢癢癢.
耶穌當然知道他的弟弟對他的不信任和不屑.
就這樣回應,七章六到八節.

$^{521}$我的時候還會有道,你們的時候常是方便的.
世人不能恨你們,卻是恨我.
因為我指正他們所作的事是惡的.
你們上去過節吧,我現在不上去過節.
因為我的時候還會有滿.
耶穌知道自己在承擔什麼使命.
當他揭露人的罪惡,指斥人的謬誤的時候.
他肯定會招來很多人的恨惡.
他們要置耶穌於死地.
耶穌的使命是玩命的.
耶穌知道自己如果上耶路撒冷將會面對什麼風險.
耶穌不是不知道,不是不小心抓到地雷.
他明知山有虎,他知道他在做什麼.
要承受什麼後果.
他不是因無知犯錯,耶穌不怕死亡.
他只是為迎向死亡而離人間.
但上帝有他預定的時間.
耶穌說他預定被殺害的時間還沒到.
所以他不想在這個時候冒這個生命的危險.
耶穌對他的弟弟說.
你們不用用激將法,不用跟我說這些蛙虎的話.
我正在做什麼,我要做什麼,我心中都有數.
這一次我不打算去耶路撒冷首節,你們自己去吧.
有趣的是,不知出於什麼原因.
耶穌最後決定去耶路撒冷首節.
不過這次是比較低調的做法.
不是帶著門徒,撫平引雷,而是自己偷偷去.
不過耶穌的行蹤還是被人發現了.
這裡招來很多人的議論.
這個後續的故事今天就不再說了.
我們已經看完整段經文.
在這裡我們可以問.
大批跟隨者離棄耶穌.
耶穌的親人也嫌棄耶穌.
原因是什麼?.
是因為他們不明白耶穌的話?.
是因為他們對耶穌吃肉喝血的教訓.
覺得反感?.
都有可能.
不過我不認為這兩件事是主要的原因.

$^{561}$事實上跟隨耶穌的門徒.
不明白耶穌說的話是太普遍出現的情況.
耶穌說的道理其實是很難受的.
他對律法的嶄新解釋.
耶穌說的很多比喻.
都不是那些知識,見識都淺薄的漁夫,牧人.
才能夠明白的,是不是?.
不要說他們.
就是今天受過高等教育的我們.
都不敢說對耶穌說的話.
都能夠聽得明白的,是不是?.
所以即使十二使徒.
他們到了耶穌上十字架的時候.
甚至是耶穌復活之後.
還是不明白耶穌降世拯救罪人的道理.
還沒有完全明白.
你看《約翰福音》二十章九字就知道了.
所以就令耶穌覺得氣盡.
「蓋歎,你們還不醒悟,還不明白嗎?.
你們的心還是如幻嗎?」.
這是耶穌經常很生氣的說話.
無知的人才會所說的一切話.
你們的心信得太遲鈍了,是不是?.
這是耶穌復活之後.
在二百五十路上對那兩個門徒說的話.
所以不明白有什麼奇怪的?.
到了最後關頭他們都不明白,是不是?.
所以不明白耶穌是不理想的情況.
但不一定要逃跑的.
至於他們說會被耶穌說的那些很可人聽聞的言論嚇倒.
甚至認定耶穌說的那些離經叛道.
褻瀆神明的話,這也可能的.
不過有這樣的想法的人通常都是高階人士.
那些宗教社會的領袖.
他們才會介意他們要捍衛真道.
誰說的是不是離經叛道.
他們會很介意的,是不是?.
民事法理才人就最介意你說的話有沒有離經叛道.
不過小老百姓對於什麼叫正統,什麼叫異端.
大概都是蒙查查而已.

$^{601}$蒙查查的人怎麼會很介意.
哎呀,你說的是異端.
在我看來,離棄耶穌的人主要是跟風.
過去他們跟隨耶穌,原因是跟風.
如今他們離棄耶穌,原因都是跟風.
如果很多人跟隨,人人跟隨.
他們很自然就會加入跟隨者的行列.
在香港,我們看到有條路.
我都要立刻先排好,然後叫太太去前面問排什麼.
先排好,對不對?.
當然是沒有執書,走前面,是不是?.
跟耶穌,反正有免費午餐吃.
吃得飽,不跟就笨,對不對?.
不過如果多數人離棄耶穌.
特別是那些德高望重的社會領袖.
知名人士都離開耶穌.
他們都不會想跟,要跟耶穌撇清關係,是不是?.
你說跟風,哎呀,是不是叫人云亦云.
缺乏主見,沒有個人立場呢?.
不過跟風一定不是一個無意識的行為.
其中有很重要的現實考量.
如果多數人認為好是對的.
我跟著說好,說對.
即使最後證明我錯了.
我要承擔的後果都很有限,對不對?.
但如果眾人認為是不對,不好.
我仍然堅持跟隨.
在未曾揭曉最終是好是壞.
是對是錯的答案之前.
今天我做爬壓水的事情.
已經要讓我付高昂的社會成本.
Social cost,是不是?.
未來贏不贏還是未知之數.
今天會輸是百分百的.
你說我如果是聰明人.
我是識相的人,我會怎麼選擇?.
跟大衛對抗,跟潮流對抗.
需要的不是一般的信心和勇氣.
很多基督徒大學生跟我說.
他們不敢在迎新營拒絕叫那些噁心的口號.

$^{641}$是因為怕犯眾怒.
他們在課堂上不敢暴露自己基督徒的身份.
不敢就同性戀,毒品合法化.
或者其他倫理的議題表達性經的立場.
不敢挑戰,自由,包容.
這些在今天被奉為至高無上的價值.
連大學生知識分子都缺乏道德勇氣.
一般的小老百姓.
我們能夠對他們有什麼要求?.
耶穌呼籲人要進閘門走閘路.
你們要進閘門.
因為引導滅亡,那門是寬的,路是大的.
進去的人也多,引導永生.
那門是窄的,路是小的.
找著的人也少.
《麥提夫音》七章十三十四節.
閘門主要指的不是門很窄.
所以要逼進去很辛苦.
胖的人逼不進去.
閘路也主要不是指路很窄.
所以很難走很遠,要踩出街.
窄是指不受歡迎.
選擇的人不多.
不是多數人自覺理所當然的選擇.
人人都穿長,我穿短.
人人都穿紅,我穿白.
奇怪嗎?難受嗎?.
很難受,對不對?.
我要說,這道聖經所描述的情況.
正好應合在我們所在的今天.
弟兄姊妹.
基督教在十九世紀到二十世紀初.
曾經是最風光.
可以說是基督教歷史的prime time.
曾經是全球佔最具優勢的宗教.
人數最多,對世界影響力最大.
基督教文化被等同為西方文化.
也是民主,科學,現代化進步的同義詞.
信仰基督教的人是受尊重的.
在社會是享有很多紅利的.

$^{681}$不過隨著時代的轉變.
基督教在西方世界越來越被邊緣化.
在二十一世紀,2000年.
美國回教會的人仍然佔人口的70\%.
二十年後,2020年已經跌到50\%以下.
千禧世代,即是說80年代出生的人.
更加只有36\%,三分一.
歐洲的情況更嚴重,不說這些數字了.
在2022年,英國的統計.
數前十年,2012到2022年.
十年間,有二千多間教堂.
因為沒有人參加聚會而被抓.
然後出賣,二千多間.
今天,回教會過信仰生活的人.
不是再多人做的事.
在香港更加不用說.
另外,聖經的教導和今天西方主流社會的教導格格不入.
特別是對待異教徒的態度.
對待同性戀,對待各種性倫理的課題.
我們顯得極其不包容.
對個人的選擇不夠尊重.
我要說的就好像昔日那些門徒跟著耶穌的時候.
為耶穌說的話感到羞愧,覺得丟架.
耶穌說的令我覺得丟架.
今天,對不起.
我們也有可能將聖經的真理說出來有丟架的感覺.
我未必說我丟架.
但我都要說,其實很多時候.
在說的時候都非常困難.
不要說對非基督徒,我對基督徒都有困難.
我做牧師,你不要懷疑我的信仰立場.
不過我要說,很困難,很困難,很困難.
困難到不懂得形容那種困難.
面對兩個基督徒要求離婚.
如果你說其中一方有婚外情.
那我立刻就捉到誰是石堅,插死他.
你搞婚外情就是犯奸淫.
但如果他們只是說他們性格不合.
有時候你真的不知道可以說什麼.
真的不知道可以說什麼.

$^{721}$我都跟你說過這個例子,說過不止一次.
不過我每次說的都是真實發生的生活例子.
所以說來說去都是那些例子.
有個弟兄來跟我說.
牧師,教教我不是說人是罪人.
人是罪人,因上上犯錯,得罪上帝誤會他的心意.
我說是,教教是不是常常說.
人只要認錯,願意向上帝悔改.
上帝就給他再次來過的機會.
我說是,然後他就說.
牧師,我現在承認十年前我娶這個太太.
是錯誤理解上帝的心意.
我完全做錯了,現在我認錯,我要悔改.
你給不給我悔改.
還是你要我堅持一錯再錯,一錯到底.
他跟我說這些,我怎麼回答他.
如果不從上帝配合這些屬靈的說法.
純粹說我們兩個是兩情相悅,然後我們走在一起.
現在我不止不悅,我還恨他.
為什麼要逼我走在一起,要我無期徒刑.
你解釋給我聽,我怎麼堅持他們要走在一起.
我見到他我就火困,那你怎麼辦.
我見到他我就情緒低落,甚至有心理病.
那你怎麼辦.
我當然知道什麼叫堅持.
但對著一個具體的對象個案.
我真的有時候有太多的不忍.
特別是今天的婚姻,兩方面對對方的依賴都很少.
生活生理的依賴都很少.
所以基本上就是兩情相悅,是感情.
如果沒有感情,你要他們勉強走在一起.
其實我又不知道怎麼說.
至於性別的問題,同性婚姻的問題就更大.
在香港問題不大,在加拿大問題已經是無法救了.
小學,人類四種性別.
四種性別,我是男人,我知道自己是男人,這是第一種.
第二種,我是男人,我變性變成女人.
第三種,我是男人,我覺得自己是女人.
所以你知道今天在美國加州還要問.
要問你想我怎麼叫你,He 還是She.

$^{761}$你不要亂問,不要說謊,你不要見到他好像男人,你就叫他He.
他可以立刻告你,就正如我現在進入女廁.
你們不要吵一聲,我告你不是你告我.
是你傷害我,你冒犯我.
我覺得自己是女人,我就有權進入女廁.
這是我主觀的認定,這是第三種.
第四種就是自由轉換,我早上是男人,下午是女人.
不行,不需要固定的.
你告訴我,整個社會,學校從小學就教這東西.
你教我,我可以說什麼?.
現在不准讓父母知道,子女要變性.
你告訴我,我們可以說什麼?.
你覺不覺得我們的信仰現在和整個社會對著幹.
和整個社會主流對著幹.
是的,你喜歡也好,你不喜歡也好.
這都是沒辦法的事.
我不知道你有沒有因為父母沒有受過什麼教育.
在公開場合胡說八道.
你就有些很丟臉,很羞恥的感覺,覺得他們死蠢.
今天你會不會對耶穌的教導,對聖經的教導.
都有很丟臉,很羞恥的感覺,我不知道你會不會.
你敢不敢表達我是基督徒的立場.
我的聖經告訴我,有這樣的倫理原則.
如果信耶穌不再是一件流行,受歡迎.
甚至和主流社會對著幹的事.
我們還會不會堅持信耶穌的決定.
我們會不會還高舉聖經的立場.
耶穌看到很多離祂而去的追隨者.
再轉向孤零零剩下的幾十人.
包括十二使徒在內.
問你們都要走嗎?.
意思是你們走不走?.
要不要跟著其他人走?.
今天同樣的問題也向我們發出.
你們也要去嗎?.
但願我們都和應彼得的反應.
主啊,你有永生之道.
我們還跟隨誰呢?.
我們已經信了.
以後知道,你是上帝的聖者.

\newpage



\section{希伯來書詩篇馬太福音 95:6-11}
\label{sec:230hZbWKaH0}
\textbf{2024年11月24日崇拜直播｜陳冠成牧師｜詩篇裡的基督－安息之主｜詩篇95\_6-11;希伯來書4\_1-2;馬太福音11\_28-30}
\newline
\newline
連結: \href{https://youtube.com/watch?v=230hZbWKaH0}{\texttt{https://youtube.com/watch?v=230hZbWKaH0}} ~~~~ 語音日期: 2024-11-24
\newline
\newline
\hyperref[sec:uf_k9SIAd7s]{\small{< < < PREV SERMON < < <}}
~
\hyperref[sec:index_chronic]{\small{[返順時目]}}
~
\hyperref[sec:a62m7d20RCU]{\small{> > > NEXT SERMON > > >}}
\newline
\newline
$^{1}$讓我們繼續以聆聽聖道敬拜.
今天的經訓是記載在詩篇95篇6至11節.
希伯來書4章1至2節.
馬太福音11章28至30節.
詩篇95章95篇1至11節.
來呀!我們要向耶和華歌唱.
向拯救我們的盆石歡呼.
因耶和華為大神為大王超乎萬神之上.
地的深處在他手中.
山的高峰也屬他.
海洋屬他是他造的.
罕地也是他所造成的.
來呀!我們要屈身敬拜.
在做我們的耶和華面前跪下.
因為他是我們的神.
我們是他草場的羊.
是他手下的民.
唯願你們今天聽他的話.
你們不可硬著心.
像當日在米利巴.
就是在曠野的瑪薩.
那時你們的祖宗視我探我.
四十年之久.
我繁衍那世代.
說:這是心裡迷糊的百姓.
竟不曉得我的作為.
所以我在路中起誓.
說:他們斷不可進入我的安息.
希伯來書四章一至二節.
我們既夢留下有進入他安息的英雄.
就當畏懼.
免得我們中間或有人似乎是講不上了.
因為有福音傳給我們.
像傳給他們一樣.
只是所聽見的道與他們無益.
因為他們沒有信心與所聽見的道調和.
萬太福音十一章二十八至三十節.
煩惱負擔重擔的人.
可以到我這裡來.
我就使你們得安息.

$^{41}$我心裡柔和謙卑.
你們當乎我的呃.
學我的養息.
這樣你們心裡就必得享安息.
因為我的呃是容易的.
我的擔子是輕散的.
願神的話光照我們的生命.
請姐妹平安.
平安.
我們繼續《詩篇裡的基督》這系列宣講.
如果你有留意.
剛才三段經文都與一個主題有關.
安息.
不知道一聽到安息這兩個字.
會令你聯想起什麼?.
很多人都會.
很自然的.
一聽到這兩個字就會聯想起.
死亡.
或者離世.
譬如可能如果你有在社交平台裡.
活躍開的.
無論Facebook或者IG.
當可能你見到有人.
發佈一些新故的消息的時候.
可能因為年紀大.
或者是疾病.
或者戰亂.
或者意外.
甚至乎可能有時.
他心愛的寵物離開世界的時候.
就會有人在他的消息下面留言.
三個英文字是什麼?.
大家很厲害.
又玩這個媒體.
RIP.
英文是Rest in Peace.
中文就是.
願死者安息.
又或者我們基督教的喪禮.

$^{81}$我們不會叫喪禮.
我們會叫什麼?.
大家都很有體會.
我們會叫安息禮.
我們不會叫喪禮.
我們會說死去的人.
我們稱呼他是怎樣?.
安息主懷.
或者說息奴歸主.
我們看到其實安息.
給我們一個很深的印象.
或者前設就是.
針對那些已經離開世界的人.
不過如果我們看回聖經.
特別今天的經文.
你其實見到有很多.
關於安息方面的教導.
是針對活人的.
是針對那些仍然在世上的人.
他強調人應該在地上.
仍然在存活的時候.
就可以經驗到.
享受到上帝的安息.
所以聖經裡面所說的安息.
不只是說RIP.
不只是願死者安息.
更多是涉及什麼?.
願生者安息.
在生的人.
所以很盼望.
藉著我們剛才所宣讀的幾段經文.
詩篇,希伯來書和馬達福音.
幫助我們明白.
並且可以進入.
上帝安息的應許.
我們先看第一段.
詩篇的95篇.
剛才讀了6至11節.
其實它全部的.
我們剛才1至5節是未讀的.

$^{121}$你會發覺.
如果我們看完整個95篇.
程序表已經印了.
講到大綱那裡.
其實是涵括了一些.
敬拜或崇拜的內涵.
11節裡面其實可以分為三個部分.
第一個部分由來也開始.
1至5節.
裡面呼籲百姓如何去敬拜.
要用一個感恩的心.
要用一個歡欣的心.
來去敬拜上帝.
因為上帝就是我們的拯救盤石.
祂又是那位偉大至高的君王.
不單止這樣.
祂還是那位獨一的真神.
超越超乎一切偶像.
那些虛假的偶像.
祂又是那位創造的主.
這裡一連串說.
天地萬物都是祂首創造.
都是屬乎上帝.
所以詩人呼籲.
我們當見到上帝這份偉大和超越的時候.
很自然應該是讚歎.
並且從心底裡.
向上帝發出那份歡呼歌頌.
這是第一部分.
至於第二個部分.
你會看到第二個「來」出現.
在第六節.
這部分比較短.
六至七節的上半部分.
他在說呼籲百姓.
這次不是歡呼了.
不是歌頌.
是屈身跪下.
表達一種很謙恭的態度.
來到上帝面前.

$^{161}$他提醒百姓.
其實在我們歌頌,歡呼的時候.
在我們情感處於高漲激昂的時候.
不要過了龍.
不要興奮到忘記了自己的身份.
第七節說到我們的身份.
上帝是我們的神.
這句話其實很重要.
他已經分了高低.
上帝是我們的神.
即是強調上帝在我們身上擁有主權.
上帝固然喜悅我們.
來到祂面前親近祂,敬拜祂.
但不代表我們和上帝是平起平坐.
我們是一種從屬的關係.
正如詩人繼續說.
我們是祂草場的羊.
是祂手下的物.
看到從屬高低的分別嗎?.
這兩句話其實是代表我們要承認.
我們必須依靠上帝的供應和保護.
我們必須順從祂的安排和帶領.
我們才可以存活.
不過詩人似乎也很明白.
因為罪的緣故.
百姓通常會怎樣呢?.
口裡是承認的,點頭的.
但回到現實生活裡.
離開了敬拜的場所之後.
很多時候只想要上帝的供應.
只想要祂的保護.
不過祂的安排和帶領就不太好.
我不太想遵守.
所以你看到第三部分的氣氛很不同嗎?.
七章後半節至十一節.
第七節後半節至十一節.
整個氣氛都轉變了.
開始變得嚴肅,嚴厲.
詩人開始告誡讀者.
「惟願你們今天要聽上帝話」.

$^{201}$「不可以硬著心」.
很明顯整個氣氛不同了.
為了強化教導的果效.
他用了一個例子.
以色列人非常熟悉的一段慘痛的歷史教訓.
作為一個借鑒.
這段歷史其實就在民數記十四章裡.
當日探子事件.
以至上帝大怒.
讓他們有一群人只能留在曠野.
然後再去迦南.
經文說有兩個地方.
米利巴是爭鬧的意思.
瑪薩是代表試探的意思.
反映以色列人出埃及後.
在這個地方試探上帝.
爭論有沒有水學.
最後上帝透過摩西擊打盆石.
解決了他們喝水的問題.
不過你會看到.
在《五經》記載.
這群人在曠野.
雖然經歷了上帝的拯救.
一次又一次見識到.
上帝對他們的供應,保守和帶領.
但卻沒有因為這樣而累積到.
對上帝的信服.
那種完全信靠他的態度.
反而是怎樣呢?.
好像經文說.
經常試探上帝.
懷疑上帝是否真的對我這麼好.
你叫我出來.
還要去迦南.
是否真的這麼好.
好像埃及那樣.
經常試探上帝.
質疑上帝的說話的可信性或可靠性.
最後我們當然知道.
有一部分人.

$^{241}$他們背棄了上帝.
他們很想回到埃及.
過他們遺留的生活.
結果就像詩篇所說.
上帝發裂腦.
上帝說我厭煩.
夠了.
我受夠了.
上帝說厭煩那一代的百姓.
不再寬容.
並且在路中起誓.
他們斷不可以進入我的安息.
上帝很斷言.
這句話是甚麼意思呢?.
他們不能回到起點.
不能回到埃及.
能不能去到終點?.
又不能去到終點.
不能進入迦南.
不能去到上帝應去的安息地方.
即是怎樣呢?.
不上又不下.
上帝說.
滯留在這裡.
滯留四十年.
直到你們離開世界.
所以這裡說.
他們不能進入上帝的安息.
因為不聽上帝的吩咐.
不上不下.
結果滯留.
不能進入應許的安息之地.
事實上弟兄姊妹.
給我們看到一個圖像.
昔日百姓雖然出了埃及.
但進不了迦南安息地.
給我們反省的機會.
原來不是所有前來敬拜上帝的人.
都一定得到上帝安息的應許.
為何得不到?.

$^{281}$經文裡說得很清楚.
因為總是硬著心.
不願意聽從上帝的話語.
就等同於今天.
不少人雖然信耶穌.
以耶穌為主.
但最後卻得不到耶穌的安息應許.
是因為什麼?.
是因為骨子裡仍然是自己扯得很緊.
靠自己來判斷.
靠自己生活.
或只是選擇聽從上帝一些.
他覺得合聽的吩咐.
他不肯將我們的生命主權.
完全交給上帝.
我們呢?.
我們今天都是前來敬拜上帝的人.
我們在這段經文光照下.
我們都要負心自問.
我們能否進入上帝的安息?.
不是等離世的時候才問.
是在世的時候.
就要不斷問自己這個問題.
我們崇拜大約兩個小時.
我們在裡面.
在這個時空.
無論唱詩或各樣的項目.
我們都會視耶穌為主.
我們都會雅門點頭.
我們會口裡承認.
祂是我們的王.
我們要信服祂.
離開了這個教會.
剩下的166個小時.
我們下次回來崇拜的時候.
我們是否仍然繼續願意.
順從上帝的吩咐和大令?.
詩篇正在問我們這件事.
不只是在禮堂這個空間.
不只是在當時的會幕或聖殿.

$^{321}$離開了.
回到生活.
那個長時間裡面.
我們是否仍然敬拜上帝.
尊重祂的話語.
承認祂在我們身上的主權?.
弟兄姊妹.
經文這裡說.
進入上帝安息是什麼?.
是今天.
今天就要把握.
今天就要聽從上帝的吩咐.
換句話來說.
我們這些已經蒙上帝拯救.
有信仰生活.
有聖經教導的人.
我們已經出了埃及.
我們也正在走進迦南.
但是我們有沒有真的.
聽從上帝的吩咐來生活?.
還是不知不覺.
我們重演昔日第一代以色列人的歷史.
在曠野裡徘徊滯留.
總是無法掌握到.
享受上帝給我們的安息.
事實上.
不單只是詩篇作者.
他提出的問題.
你會發覺到了新約.
希伯來書還在問這個問題.
到現在.
作者仍然在問.
聖經仍然在問.
我們進入上帝安息這個問題.
事實上你會看到.
希伯來書是參考至詩篇九十五篇.
希伯來書第四章一至二節.
或者說再前一點第三章.
在聖靈的引導裡.
希伯來書作者引用了詩篇九十五篇.

$^{361}$作為第一代教會信徒的借鑒.
上帝昔日對以色列人的安息的應許.
作者說今天仍然是有效.
對新約的信徒仍然有效.
並且惠及的範圍更廣.
不再單單是給以色列人.
而是給每一個凡應信耶穌為主的基督徒.
包括我們.
不過作者鄭重提醒.
要兌現上帝安息的應許是有條件的.
說得很清楚是有條件的.
首先他說人要用一個畏懼的態度.
來回應上帝安息的應許.
什麼是畏懼呢?.
我們聽到好像很負面.
畏懼不是害怕.
不是畏首畏尾的那種狀況.
或者用另一個字我們會比較明白.
敬畏,小心,謹慎.
這個態度令我們面對上帝的話語.
不敢抬攬.
面對上帝的話語我們不敢輕慢.
並且我們知道.
如果我們持續心硬.
不肯聽上帝的吩咐.
是有後果的.
昔日的百姓有後果.
今天也會有.
甚至可能是後果更嚴重.
如何體驗一份敬畏,小心,謹慎的態度呢?.
早幾年大家都體會到.
你覺不覺得?.
疫情的時候.
你是最敬畏,謹慎地過生活的.
有沒有留意?.
你因為很擔心感染.
影響到你身邊的人.
不能上班.
甚至面對生命的威脅.
於是沒有人會心存僥倖.

$^{401}$說我不戴口罩出街.
現在因為過了.
但在那段高峰時間.
沒有人會想著沒事的.
不戴口罩也沒事的.
照出街也沒事的.
每一個人都帶著謹慎,敬畏的心.
做足所有防護措施.
盡可能減低風險.
同樣.
雖然上帝的安息.
這個應許已經留給每一個新約的信徒.
但是.
作者亦警告.
沒有以敬畏態度來回應的人.
是要承擔風險的.
有可能.
好像作者所說.
趕不上這個應許.
趕不上的意思是.
你去不到目的地.
就像昔日以色列人一樣.
去不到迦南應許地.
被排拒於門外.
就像昔日抗疫的百姓.
最終被神拒諸於應許的大門之外.
而另一方面.
作者亦勉勵我們.
不可以只有心.
只有心不足夠.
只有敬畏,畏懼的態度不足夠.
還要有行動.
連同行動來配合上帝這個安息的應許.
作者在第二節說.
要用信心和所聽見的道來調和.
調和的意思是.
混合它.
打成一片.
或者所謂的「連合起來」.
簡單來說.

$^{441}$單單知道上帝給我們安息的應許.
甚至我們知道前人失敗.
拿不到這個應許.
這樣不足以我們自己.
可以進入上帝的應許.
我們必須要在行為和態度上.
作出合宜的調節和回應才可以.
或者套用雅各書的說法.
我們會明白多些.
進入上帝的安息.
不可以單單靠聽到來得到.
安息不是單靠聽回來的.
要運用信心,行道,實踐真理.
才可以拿到這個應許.
周貞瑜.
在座有沒有人定期做運動?.
有沒有?.
大家都有,有人點頭.
我們知道.
做運動可以幫助自己保持健康.
我們知道這個道理.
我們都知道.
身邊有很多成功的例子.
透過做運動.
來籤體或身體健康.
但我們不會因為單單知道這兩個資訊.
你會認為自己已經成功保持健康.
不會的.
不過很有趣.
在現實生活中.
我們真的有這些錯覺.
在不同的事情中.
我們總會有種感覺.
我聽了很多次.
我明白了.
我很認同.
我還想過怎樣去做.
這樣是否等於做了?.
但我們會覺得是做了.
我們想過,知道,聽過,認同了.

$^{481}$有時會覺得是做了.
你看我現在健不健康?.
我以前很健康.
有機會給你看我結婚前的照片.
瘦過現在30磅.
你想像不到吧.
黃震是幾號的恩典.
把我養得這麼肥.
你想想.
我20年前很健康.
我一直都想過想變回健康.
我也很明白怎樣可以變健康.
不過.
我可否透過我單說單想.
單想像.
晝夜思想也好.
看很多Youtube減肥片也好.
我就變了健康?.
不會吧.
你看我現在就知道.
我一直不做運動.
不改變我的飲食生活習慣.
我是在騙自己.
給自己一個感覺.
哄自己以為變了健康.
其實根本沒有變健康.
有保持肥胖.
同樣希伯來書說的安息.
是一種持續聽道行道.
與神同行同工的一種屬靈狀況.
超越了詩篇裡說的.
不是去迦南地.
今天我們不需要去迦南地.
我們要維持那種持續聽上帝吩咐.
與神同工同行.
是一種靈情.
是一種很積極.
很動態的安息.
正如作者說.
進入安息是怎樣.

$^{521}$是潔力.
有沒有覺得很違和.
潔力但進入安息.
安息不是不用做嗎?.
錯.
很多弟兄姊妹有時做得很辛苦.
他們會說一句.
我等上帝接我回去.
等我不用做.
我就安息.
是也不錯.
不過我們大部份人都不是.
如果是的話.
上帝應該我們一信祂.
就立刻讓你安息.
這樣就很好.
但不是.
很明顯安息是要你仍然在世.
甚至在勞苦工作裡.
在那裡經歷.
所以上帝不接你回去住.
是有原因的.
你應該好好享受.
在苦中.
在辛勞裡.
你是會找到安息的.
正如這裡說.
作者繼續說.
我想大家讀.
我用了新漢語譯本.
大家會明白一點.
第四章十至十一節.
程序表也有.
或者螢光幕也有.
請大家同讀.
預備起.
(上帝說話).
好.
你看到希伯來書的作者.
他用上帝的創造進一步解釋.

$^{561}$什麼是進入上帝的安息.
我們大部份人都明白.
上帝的創造的記載.
創世記二章一至三節.
我們知道上帝用了六天.
來做齊天地萬物.
到第七天.
上帝就如經文所說.
(上帝說話).
不過我們要想.
上帝安息是否因為.
連續六天工作太辛苦.
太累.
所以第七天.
無論如何都要放假.
無論如何都要讓自己休息.
休息充電.
是否這個意思.
絕對不是.
這是人的需要.
上帝是不需要的.
永不疲乏的.
聖經說.
上帝的安息.
事實上不代表.
他靜止.
什麼都不做.
正如耶穌所說.
(上帝說話).
上帝沒有一刻停止.
對這個世界.
對我們的管教.
安息其實原本的意思.
是欣賞.
是滿足.
上帝用了六天.
做齊萬物.
包括他對萬物的心意.
萬物要如何運作.
才能配合上帝的旨意和計劃.

$^{601}$上帝說.
完成了.
做好了.
所以第七天.
他卸了自己的工.
停了創造的工作.
是因為什麼?.
沒有什麼需要加減.
完成了.
完聖所有東西.
沒有什麼需要加減.
畫上一個完美的句號.
所以上帝在第七天.
安息是說.
他欣賞自己的完美創造工作.
欣賞當萬物受造物.
包括我們.
按照他的旨意運行發揮的時候.
一切按照上帝的秩序進展的時候.
上帝感到心滿意足.
這就是上帝安息的意思.
我們如何能享受上帝的安息?.
考考大家.
就是我們一直說.
你聽他的話.
而不是翹起雙手.
什麼都不做.
事實上.
當萬物包括我們.
什麼時候按照上帝的旨意運作.
什麼時候我們順服上帝的吩咐.
聽從他的話語.
活在他創造的計劃.
他的旨意和秩序的裡面.
我們就和神一樣進入上帝的安息.
我們就可以和上帝一樣.
經驗到真正心滿意足.
相反.
如果我們像昔日第一代百姓.
我們不肯被上帝帶領.

$^{641}$硬要自己作主的時候.
我們就是離開了上帝創造的秩序和法則.
就不會安息.
無論你再努力.
再做多些事情.
憑自己去死.
即使你有這樣做.
你很想透過這樣去經歷安穩的生活.
也沒有辦法.
因為你離開了上帝創造的旨意.
你不會找到真正的安穩.
這是聖經所說.
所以要得到上帝安息的應許.
第一件事.
就像經文所說.
先停手.
協助你自己的功.
停止再依靠自己去忙亂的生活.
亂走.
第二,轉向.
你回到上帝身上.
你要竭力去攻克自己.
不要配合上帝的帶領.
不要竭力回到上帝身上的旨意.
第三,這不可少.
你要持續和耶穌基督連繫.
你要和祂結連.
讓祂帶領你.
不是你自己去做.
讓祂幫你活出上帝在你身上的旨意.
最後一段經文.
大家很熟悉的.
馬來福音第十一章28至30節.
我想大家再讀一次.
慢慢讀.
慢慢去接受耶穌想說什麼.
預備起.
到我這裡來.
我就使你們得安息.
我心裡柔和謙卑.

$^{681}$你們當乎我的厄.
學我的弱色.
這樣你們心裡就必得享安息.
因為我的厄是容易的.
我的擔子是輕散的.
耶穌很清楚.
說了兩次.
祂願我們在祂裡面得到安息.
祂願意把這個應許給我們.
首先耶穌這番話.
是針對當時宗教領袖.
給百姓一份的重擔.
令百姓以為.
我們要遵守很多他們所訂立的條文規矩.
這樣才可以換取上帝的眷顧和幫助.
假如在生活上我們遇到困難.
感受不到上帝的幫助的時候.
就會以為一定是我守律法守得不好.
一定是我守得不夠.
於是就被要求.
你要再守好一點.
你要再守多一點律法.
於是宗教反而變成了無法令人安息的枷鎖.
這是耶穌說這番話.
針對當時的情況.
在這裡耶穌卻是指出.
得到安息的關鍵.
不是我們的信仰知識再多一點.
不是我們有再多的宗教行為.
而是耶穌說得很清楚.
「道他那裡來」.
到耶穌那裡.
這位安息之主告訴我們.
真正的安息取決於我們和他的關係.
是一份關係.
是一份信任的關係.
是一份願意被他帶領的關係.
耶穌的應許.
當我們願意接受他的邀請.
和他連結在一起.

$^{721}$耶穌說就可以得到安息.
實際是什麼意思呢?.
耶穌可能擔心他的聽眾會不明白.
所以他用了當時群眾熟悉的事物.
來解釋這番話.
他用了牛的額.
古代耕田.
要拖牛去爬田.
通常有一個用木造的額.
圈著牛.
讓牠順服主帶領.
來爬田.
耶穌的意思其實是.
反映古代猶太人耕作的場景.
通常不會用一頭牛.
用兩頭牛.
兩頭牛一起扶額來爬田.
並且用一頭年長.
有經驗的牛.
帶著另一頭年輕.
沒什麼經驗的牛來做.
你開始明白耶穌想說什麼.
耶穌比喻自己是哪一頭牛.
就是那一頭有經驗.
年長的牛.
帶著我們這些沒什麼經驗.
年輕的牛.
來學習跟隨上帝.
一起扶上帝的額.
一起完成上帝的旨意.
所以真正其實.
誰做得多一點.
耶穌說是他.
其實是他做得多一點.
他出力最多.
誰才是真正完成上帝的工作.
也是他.
我們其實只是配合他.
向他學習.
如何事奉上帝.

$^{761}$如何愛人.
這種態度.
事實上耶穌.
也是一位創造的主.
他很清楚.
什麼才是適合我們.
什麼額才能讓我們輕生.
他同樣是那位救贖的主.
即使我們扶額時.
可能會亂來.
我們會走迷甚至行錯.
但是他會怎樣?.
他有能力帶我們回去.
他有能力拯救我們.
所以無論是工作.
無論是休息.
無論是停下或前進.
我們可以放心被耶穌帶領.
我們可以進入祂的裡面.
不需要像以前的百姓.
總是懷疑祂.
總是跟祂角力對抗.
可以預期.
如果我們真的願意.
接受耶穌這份扶額的邀請.
你的生活真的會變得輕生.
你做同一樣的事.
但你會變得輕生.
因為有耶穌帶領你回到正軌.
回到上帝的創造裡.
你可以進入上帝的安息.
直到永遠.
所以等於姐妹這裡說.
跟耶穌和父的額.
不是點頭後各有各做.
而是我們真的決心.
配合耶穌的步伐.
而不是再自己亂為主奔跑.
不過今天不容易.
因為今天我們一定是說快的.

$^{801}$什麼都要立即做.
慢就像一種罪.
但我們確實.
要像耶穌這裡說.
你要從容跟耶穌的步伐.
不知道大家認不認識.
中國一位很重要的基督徒.
叫做王明道.
他是上世紀在中國.
其中一位很重要的牧者.
影響了很多基督徒.
曾經有位外國學者跟他交談過.
他們談論起.
與主同行這個課題.
王明道問對方.
你是如何與上帝同行的?.
對方說了很多屬靈的操練.
譬如祈禱,讀經,靈修,默想.
沒想到王明道突然間打斷他.
說錯了.
與主同行是要用步行的速度.
來與主同行.
他不是說你做什麼.
而是說你是否配合耶穌的步伐.
這句話很有意思.
值得我們細味.
步行的速度.
不是快,也不是慢.
給我們更多心靈的空間.
看見耶穌在做什麼.
你可以細味敏銳到.
耶穌的步速是怎樣.
以至你有這個空間.
調教你自己配合耶穌的步速.
可能是加快,可能是減慢.
可能是停下來.
我還記得上星期三.
我在這裡帶星期三晚的默想祈禱會.
帶著的時候.
大家知道那天有颱風.

$^{841}$我當時還未到八時.
在升降機上跟阿邦傳道說.
打不成的.
誰知帶到八時多.
當大家默想的時候.
有風聲飄上來說.
等下過了八日.
我馬上嚇了一跳.
很緊張.
到底我是否要加快腳步.
快點默想完.
大家快點散.
但我又想那段經文很好.
我就是帶著平靜風浪的那段經文.
我想耶穌還在船上.
還未完.
雖然有風浪,但耶穌沒有離開.
站在那裡.
我相信這是你和我也需要學的功課.
如何配合在風浪裡.
發現耶穌還在.
留意祂的步速.
祂不離開時,你不需要急於掙扎.
祂走的時候,你就要緊緊跟隨.
耶穌邀請我們同行.
同父恩厄.
但太多時候我們心急.
我不知道你是否一個心急的人.
你總是想超前耶穌.
耶穌不動手時.
反而不能讓你平靜,更加混亂.
為何還不出手.
我要比耶穌快.
我要跑在耶穌前面.
甚至跑著跑著.
耶穌也不知道去了哪裡.
你自己一個人亂跑.
你還以為是為了耶穌跑.
我們錯過了很多和耶穌相交.
認識祂,認識自己的機會.

$^{881}$事實上弟兄姊妹.
今天耶穌很想我們.
看看祂的走法.
看看祂如何負壓.
然後跟隨祂學習.
學習如何愛上帝.
學習如何服侍人.
讓我們在事奉中.
在生活中少一點急躁.
少一點焦慮.
多一點像耶穌般的從容.
弟兄姊妹盼望.
好像這段經文.
今天我們把握.
上帝給我們安息的應許.
我們接受耶穌給我們的管理.
接受祂給我們的壓.
一起和祂進入上帝的安息.
我們同心圖告.
帖幅說我們多謝你.
你給我們有安息的應許.
不是要等到我們離開這個世上.
而是現在.
當我們活著的時候.
我們就可以依靠耶穌.
和耶穌一同去負壓.
我們經歷這種的安息.
這一份的安穩.
不是因為我們不需要面對難處.
不是因為我們不需要工作.
不是因為我們需要面對種種的限制.
而是因為有耶穌在.
而是因為耶穌已經拯救我們.
並且你願意帶領我們.
幫助我們一同去完成.
上帝在我們生命裡的旨意.
以致我們知道原來人生無論順逆.
無論是輕省或是辛苦.
其實都在上帝的旨意裡.
在上帝的管理裡.

$^{921}$這一切不能難阻上帝的供應.
你的豐盛,你的保守.
以及你對我們的愛.
上帝盼望當我們.
不只是在禮堂裡說阿們.
當我們離開.
我們回到工作,生活裡的百多個小時.
我們繼續對這段經文說阿們.
我們用我們的行動.
調教我們的步伐.
敏銳你在我們生命裡的安排和旨意.
以致我們的生命得到成長.
更加謹靠耶穌.
我們這樣的禱告.
奉耶穌你的名來祈求.
阿們.
(影片完結).
\newpage



\section{希伯來書詩篇 102:3-28}
\label{sec:a62m7d20RCU}
\textbf{2024年12月1日主餐崇拜直播｜陳明泉牧師｜詩篇裡的基督－永活神子｜詩篇102\_3-28;希伯來書1\_10-12}
\newline
\newline
連結: \href{https://youtube.com/watch?v=a62m7d20RCU}{\texttt{https://youtube.com/watch?v=a62m7d20RCU}} ~~~~ 語音日期: 2024-12-02
\newline
\newline
\hyperref[sec:230hZbWKaH0]{\small{< < < PREV SERMON < < <}}
~
\hyperref[sec:index_chronic]{\small{[返順時目]}}
~
\hyperref[sec:_MkLayWrDNI]{\small{> > > NEXT SERMON > > >}}
\newline
\newline
$^{1}$今天還要讀的有兩段經文.
第一段是在舊約詩篇102篇23至28節.
第二段是在新約希伯來書第一章10至15節.
請聽我讀出第一段是在舊約第779頁詩篇102篇.
他使我的力量中度衰弱 使我的年日短少.
我說 我的神啊 必要使我終年去世.
你的年數世世無窮.
你起初立了地的根基 天也是你所做的.
天地都要滅沒 你卻要長存.
天地都要如外衣漸漸舊了.
你要將天地如女衣更換 天地就都改變了.
唯有你永不改變 你的年數沒有窮盡.
你僕人的子孫要長存 他們的後裔要建立在你面前.
第二段經文是新約希伯來書第一章10至12節.
詩人引用剛才所讀的詩篇.
又說 主啊 你起初立了地的根基 天也是你所做的.
天地都要滅沒 你卻要長存.
天地都要像衣服漸漸舊了.
你要將天地捲起來 像一件外衣.
天地就都改變了 唯有你永不改變 你的年數沒有窮盡.
弟兄姊妹平安.
我們繼續用詩篇思想基督耶穌.
在12月 特別是將要預備自己慶祝基督耶穌降生的日子.
我們對於基督耶穌的認識有多少呢?.
我們對祂的身份 對祂的生命.
或者祂所為人類所作的作為 我們有了解有多深呢?.
今天我們希望藉著剛才的詩篇102篇和希伯來書.
來我們更深入的理解基督耶穌對我們的工作.
在我們進入經文之前 有一個這樣的機會.
跟我的同輩聊天 中年男子.
比我大幾天的 同年.
我記得他跟我說了一番話 很深刻的.
那時聊天 可能我們都很近似.
所以人生到中年都有一些打拼.
去到有些位置回顧自己的成長.
弟兄跟我說 他覺得人生現在都算是這樣了.
他想著自己踏入40歲就會死了.
他跟我說 為什麼你會這樣想呢?.
背後說在他很年輕的時候.
他的爸爸大概40歲就去世了.

$^{41}$所以他覺得人生的大限 或者人生裡面.
盡的了 在他從小到大的想法裡面.
覺得40歲就是人生男士的一個極限.
所以他覺得在他很小的日子裡.
他已經很努力 很打拼.
他覺得人生既然這麼短暫.
他就要用盡他的一生.
來爭取他自己想要的東西.
所以他很發奮圖強讀書 很努力.
他已經結婚生兒養女了.
他覺得自己已經比他爸爸那個年紀.
已經拿多了東西.
所以他有一種滿足感.
但是我問他你快不快樂?.
他就沒有回答我.
他沒有特別回答我這一條問題.
他只是笑笑笑.
同樣的問題 問完他之後 我也問一下我自己.
我也發覺 其實在我自己成長裡面.
我爸爸在我四十幾歲的時候.
我爸爸四十七歲去世.
原來我也有跟他有類似的想法.
我覺得四十多歲是人生男士的一個極限.
接著我現在賺了.
我已經被我爸爸長壽了.
我已經入了五了 小弟 大家知不知道.
真的入了五了 不止 是吧.
像登陸 是吧.
去到中年這些日子就開始想.
我生命裡面 我快不快樂呢?.
我覺得怎樣呢?.
相對於我成長裡面我見到的爸爸.
我自己跟他有什麼不同呢?.
當我一直想這個問題的時候.
我想起今天的詩篇和希伯來書所講.
事實上有一些外國的研究.
2020年有一個心理學的發表.
他講到人生最不開心的時間.
原來就是在四十五至五十五歲.
這十年是人生的U型谷底.

$^{81}$他講到人生的快樂.
一出生最開心 一直下降.
谷底就是四十五至五十五歲.
接著又再慢慢升上來.
他用一個人生的週期歲數.
來看人生的快樂程度.
不過他的解釋或研究裡面.
他有些註腳.
這個U型未必是一定的.
這個U型有些人會短一點.
有些人未必到谷底.
有些人可能是因著信仰.
或者找到生命意義.
這個U型的谷底就下降得沒那麼深.
又或者在人生裡面是V型.
很短一站就上去了.
不知道我們今天生命狀況是怎樣.
可能那個是你.
可能是你的家人.
又或者是你的未來.
希望今天這篇信息幫到我們.
我們先祈禱求主幫助.
天父上帝願你幫助我們.
讓我們在你話語當中.
我們既體會我們人生的週期.
我們藉著我們生命所經歷的歷程.
我們更加明白基督耶穌的心腸.
幫助我們.
讓我們今天我們一起來讀你話語的時候.
給我們內心有激勵有反省.
也讓我們帶著你的話語.
好好地過我們的生命.
我們恭敬將此時此刻交託給你.
幫助孫講的講得清楚.
幫助聆聽的每一個靈姐妹聽得明白.
願我們被你聖靈所繼續來瑞領我們.
獻上禱告奉求基督耶穌得勝的聖靈祈求.
Amen.
我們先看今天主要看的詩篇.
詩篇102篇.

$^{121}$詩篇102篇大概28節.
我沒辦法因為時間所限.
我沒辦法將整個詩篇跟大家解釋.
不過給大家有一些概覽.
整個詩篇分為三個段落.
第一個段落就是詩篇第102篇的1至11節.
102篇1至11節說什麼呢.
其實是詩人向上帝禱告哀求.
因為這篇詩篇在猶太人傳統中叫作懺悔詩.
不是說他做了什麼罪行令他要懺悔.
而是面對著那種苦難.
面對著他的國家.
大概是猶太人被擄之後的作品.
面對自己的民族國家.
被外邦所滅.
他們的同胞被擄巴比倫.
又或者舊時以錫安城耶路撒冷為榮的聖殿.
整個城牆被毀.
面對著身邊的那種狀況.
詩人就用這首詩篇來向上帝發出哀鳴.
說的是一個哀歌.
面對著一種這樣的狀況.
他向上帝祈禱.
經文中特別說到.
詩人不單是說一個抽象的國婆家亡.
在第2節.
102篇第2節.
他說我在急難的日子求你.
向我側耳不要向我掩面.
我呼求的日子求你快快應允我.
在這段中說到他自己第一身.
面對著這些苦難.
面對著這些很大的痛苦.
接下來他在數.
數算他自己面對的一些困難.
他面對飢餓.
他面對飢餓到什麼地步呢.
他說到用一些爐灰來當飯吃.
就是沒東西吃.
用一些灰瓶來放進口裡.

$^{161}$類似上一代.
或者再上一代人說日本人打仗的時候.
吃樹皮吃樹根.
吃泥土來充飢.
這個就是詩人裡面第一個.
面對著國家的那種痛苦的時候.
他所經歷的那種苦況.
他更加說到他自己的身體.
他的骨頭好像被火燒一樣.
他的心受傷害.
如草枯乾.
甚至忘記吃飯.
因為他哀嚎的時候.
他的肉緊貼骨頭.
你可想而知.
他那裡說什麼呢.
肉緊貼骨頭的意思就是說他瘦.
瘦到被捕就是皮包骨.
又不想吃東西.
心裡面哀痛.
同時間面對著國破家亡.
又經歷戰火.
面對有敵人的時候.
又被敵人中日的來咒罵.
所以由第二節開始.
到第十一節.
整個段落.
其實是訴苦的.
跟上帝說有多大的那種苦難.
當他面對這種痛苦的時候.
他不是停在自怨自艾的地步.
然後到第二個段落.
由第十三節開始.
讀到第二十二節.
這個就是詩人裡面.
詩篇的第二個段落.
在這個段落裡面.
他就將焦點.
由自己的痛苦轉念.
轉到哪裡呢.

$^{201}$轉到在上帝當中.
唯有你 耶和華傳到永遠.
你可以紀念的名也傳到萬代.
你必起來憐憫錫安.
因為現在可憐他的時候.
日期已經到了.
你的僕人原來喜悅他的石頭.
可憐他的塵土.
列國要敬畏耶和華的名.
世上諸王要敬畏你的榮耀.
在這裡一邊祈求一邊數船.
他將他剛才想到的苦況.
去到第二個段落.
就成為了一個很大的改變.
他的筆鋒就轉移.
去描述上帝的永恆.
和上帝對以色列民的那些的應許.
所以在經文裡面.
你會看到.
雖然面對自己眼前國破家亡.
但是他思想上帝.
上帝答應了這個國家民族.
甚至乎答應了他自己.
他相信上帝永恆的能力.
是會奉庇這個國家.
奉庇他自己.
所以去到第二個段落.
他將上帝的永恆.
一連串地數算出來.
也說到上帝有永恆的恩典.
淋到他自己的國家民族當中.
這個就是整個經文裡面.
某程度上就轉念了.
講到上帝的永恆.
你看到剛才整個經文裡面.
頭的篇幅是講自己很痛苦.
第二個篇幅就講到上帝的永恆.
在中間的這一部分.
就成為了第三部分的詩篇.
23節至第28節.

$^{241}$也是剛才我們所讀的那些內容.
當他面對自己很苦.
見到上帝的永恆.
他的祈願是什麼呢?.
我們再看看這幾節經文.
他使我的力量終度衰弱.
使我的年日縮短.
我說,我的神啊.
不要使我終年去世.
你的年歲西西無窮.
你起初立了地的根基.
天也是你所做的.
天地都要滅沒.
你卻要長存.
天地都要如外衣漸漸舊了.
你要將天地如女衣更換.
天地就都改變了.
唯有你永不改變.
你的年歲沒有窮盡.
你僕人的子孫要長存.
他們的後裔要建立在你面前.
在這一段落裡面.
他在講什麼呢?.
或者他最終的祈禱是什麼呢?.
面對人生命的軟弱,痛苦.
和永恆的主.
他的祈願其實是.
你覺得他不是很高尚.
他只有一個很單純的祈願.
就是生存.
不要死.
還有他求壽.
他用的字眼特別講到.
上帝令他在這些日子裡面.
終度衰弱.
在他中年的日子.
甚至乎他講到.
面對著這種國破家亡.
那種苦難.
他的年日短少.

$^{281}$某程度他真的覺得.
上帝是一個.
他的旨意.
要令他們這一代人.
面對一個這樣的苦難.
但是即使是這樣.
他求上帝不要使我終年去世.
你的年歲世世無窮.
很簡單.
私人其實就是只是想.
希望自己在一個這麼困難的時間裡面.
上帝.
求你讓我在地上.
多一些的日子.
你是一個永恆的主.
而我是一個地上.
短促生命的一個.
塵土裡面生活的一個人.
你可不可以開恩.
讓我這個短促的生命.
可以長久一點.
不要在終年.
在生命裡面最鼎盛的時間.
反而就要離開了人世.
他求上帝讓他長久一點.
然後他再繼續講.
他又用很多的篇幅.
講到上帝的永恆.
他用什麼來表達呢.
他講到天與地.
天與地其實就是我們自然.
造物界裡面.
我們覺得永恆不變的定律.
如果我們有科學家知道.
有一些恆久不變的定論.
就在創造的天地萬物來改變.
私人訴諸於創世所見到的創造的萬物.
然後他講到這個創造的世界的萬物.
是來自那位創造的主.
他立了地的根基.

$^{321}$天也是他的手所造的.
不過這些都是被造之物.
在上帝的手中裡面.
天與地有一天都要改變.
所以第26節.
天地都要滅沒.
理卻要長存.
天地都要如外衣漸漸舊了.
你要將天地如女衣更換.
天地就都改變了.
唯有你永不改變.
你的年歲沒有窮盡.
在這裡他就用天地的改變.
和上帝的永存.
上帝的永恆來做一些對比.
天與地對於人世間來說.
那個都是很長久的.
不過天與地的長久.
都及不上.
都是受造之物.
在上帝的手中裡面.
天與地就好像一件衣服一樣.
脫了它.
會捲起.
然後就換另一件.
無論外衣也好.
女衣也好.
內外外外.
天地就好像衣服一樣.
會有更換改變的一天.
他用這個比喻說到天地都會改變.
有一天都會消失於永恆當中.
甚至天地不是永恆的存在.
不過對於上帝來說.
上帝就是永恆的主.
將天地要改變.
上帝的年歲卻沒有窮盡.
他用這個比喻來強調.
上帝那種超然.
超越萬物的能力.

$^{361}$以致到詩人繼續訴諸.
他求上帝.
你的僕人子孫要長存.
他們的後裔要建立在你面前.
這就是詩篇裡面最後的結語.
詩人求上帝給他長久的日子.
他也求上帝給他生命裡.
有更加長久的日子之外.
希望他的子子孫孫.
可以延續在這個世界當中.
不會淹沒在歷史的洪流裡.
這就是詩篇102篇裡所呈現的內容.
整個詩篇裡所講的內容.
簡單來說就是.
特別是在苦難中.
這個詩人見到很多困苦.
他自己第一生經歷.
人生很苦.
但既然人生這麼苦.
他不是氣派.
他不是說上帝讓我上天堂.
他不是這樣.
雖然人生很苦.
但他很想繼續在地上日子長久一點.
他希望他的後代.
可以繼續地延綿不絕.
在掌管歷史.
掌管天地永恆的主面前.
他祈求上帝.
他求上帝讓他日子長久一點.
他的後裔世世代代延續下去.
所以其實詩篇102篇.
他的內容或祈願的想法.
不是一些很深奧或高言大智的內容.
只不過是一個很痛苦的人.
他在生活中掙扎求存.
他是一個倖存者.
他更加期望的是.
在他短促的人生當中.
上帝可以給他多一點時間.

$^{401}$讓他可以見到他的後代.
也讓他的後代可以繼續地繁衍.
在上帝的永恆當中.
上帝永恆的主答應過.
所以他就是這樣來祈求.
當我們明白這件事的時候.
我們稍為有時間去想一想.
我們今天人生在祈求什麼呢?.
在我們自己經歷每一年.
一年差不多要過去了.
將要踏入新的一年.
你的心裡在想什麼呢?.
有些弟兄姊妹可能會覺得.
人生實在太苦了.
所以其實在世的日子未必需要很長.
反而要的就是要精彩地過人生.
我又聽過一些說法.
人生既然這麼苦.
人生為什麼要這麼努力呢?.
人生既然這麼苦.
反正死亡也要臨近在我們的生命當中.
為什麼我們要這麼努力.
去爭取我們覺得地上重要的東西呢?.
不過說這件事.
他就棄牌,然後就順應自然.
大概像中國人所說的道家思想.
類似都是這樣的.
以有涯除無涯,桌以.
人生裡既然有這麼多書要讀.
人生裡面對這麼多困難.
為什麼要逆流而上呢?.
不如順應自然.
甚至乎輕輕地過人生.
這樣人生就不錯了.
有一些說法是.
人生既然有這麼多困難.
不如順應自然.
不需要力爭上游.
不過今天我們看到的詩篇102篇.
似乎作者.

$^{441}$他給我們的詩篇有另外一個圖畫.
生命裡即使是短促.
但我們也要好好爭取.
努力地活下去.
也希望我們生命可以.
祝福到我們自己的後代.
我們希望看到世世代代.
繼續地來生存.
我們沒辦法掌控自己的壽命有多長.
但是在一些未知.
在人生裡不知道的狀況下.
我們繼續打拼.
繼續為自己的生存.
為自己生命的福祉.
繼續祈求和努力.
剛才說過.
跟我同年紀的弟兄.
說人生的一個問題.
我問他覺得開心嗎?.
他沒有回答我這個問題.
他只說.
我人生現在眼前.
我要為我的家庭兒女.
即使自己不開心也好.
或者自己有很多壓力也好.
我都要好好地讓他們生活.
生活得快樂.
他們的快樂會令我心裡有滿足.
人生我們沒辦法決定自己的壽命長與短.
但是今天面對生命.
有時候我們覺得困難.
不理想.
我們採取放棄的態度.
或者我們依然積極.
我們祈求上帝.
讓我們生命可以繼續延續下去.
也將福氣延續在我們的後代當中.
這是詩篇102篇裡的訊息.
當我們掌握這個訊息的時候.
你就會覺得很有趣.

$^{481}$希伯來書引用詩篇102篇.
剛才我們所讀的.
特別是第24至27節.
剛才我們所讀的段落.
是人生意義.
祈求上帝的福氣.
希伯來書引用這段.
不是說這件事.
我們再看希伯來書.
希伯來書的第一章.
10至12節.
他這樣引用的.
主人你起初納了地的根基.
天也是你所做的.
天地都會消滅你卻長存.
天地都會像衣服漸漸舊了.
你要將天地捲起來像捲一件外衣.
天地像衣服都會改變.
你卻永不改變.
你的年歲沒有窮盡.
在這裡裡.
為什麼希伯來書的作者會引用這一段呢.
我沒有印給大家.
我讀給大家聽.
其實上文下你講什麼的.
希伯來書.
希伯來書的作者.
在一章第五節.
是整傳的段落.
他在講天使與基督耶穌來做比較.
所有的天使神從來對那一個說.
你是我的愛子.
我今日生你.
又指著那一個說.
我要作他的父.
他要作我的子.
然後他就列舉不同舊約的詩篇.
舊約裡面的經文.
來突顯耶穌基督的身份.
在第八節那裡.

$^{521}$論到子卻說.
神啊你的寶座是永永遠遠的.
你的國權是正直的.
你喜愛公義.
恨惡罪惡.
所以神就是你的神.
用喜樂由高你.
勝過高你的同伴.
在第八節那裡.
一章八節.
他特別強調就是.
聖子坐在寶座當中.
他就是那位掌管萬物的君王.
他的寶座是永永遠遠的.
然後他就引用剛才所讀的那段經文.
值得留意他的引用.
原本的段落.
就是由剛才詩篇102篇那裡.
他是在說.
他是求壽的.
那個詩人是求不要終年去世.
希伯來書的作者沒有寫那段.
他中間就插入了.
就說到上帝的永永遠遠.
所以嚴格上來說.
這段經文其實他不是解釋.
不是解釋102篇裡面的內容的.
他只是借助.
102篇詩篇裡面的內容.
來凸顯耶穌基督的那種永恆.
那種的能力.
我們都知道的.
聖誕節慶祝的就是耶穌生日.
是不是.
永恆的主.
為什麼會有在地上出生的日子呢.
如果他在地上是受造之物.
他會有一天終結的.
是不是.
不過希伯來書的作者說到.

$^{561}$上帝的兒子一早已經存在了.
不過道的存在.
接著他插入歷史的就是.
道成了肉身.
而這個肉身的耶穌基督.
他是一個永恆的主.
他是世世代代.
他繼續.
顏面不絕地守望.
愛眷.
以及將恩典淋在他的子民當中.
所以當我們.
如果用很偏頗的理解.
就是受造的這個人.
他沒辦法有永恆的.
但是當他其實.
是先存於道成肉身.
接著他加入歷史裡面.
接著他的生命其實一直都存在.
只不過.
是道成肉身的時候.
我們親身看到神的形象.
不過上帝不是由歷史才開始.
耶穌不是由歷史的時候才開始.
耶穌是自有永有.
由起初到末了.
他就是那位不變的神.
所以希伯來書的作者.
就是用這一段經文.
來說到耶穌基督.
是一個這樣超然的主.
人會改變.
但是他永不改變.
希伯來書的讀者.
其實可能都在面對.
詩篇102篇裡面.
那種心路歷程.
當然沒那麼大困難.
沒去到死的.
不過那些人借.

$^{601}$那些人信仰隨流失去.
又或者生命裡面困苦流離.
希伯來書的作者告訴他.
你信的是誰?.
那個是永恆的主.
所以這個耶穌.
是納根基的上帝.
天地是由耶穌基督所做的.
天地都會改變.
耶穌卻是會長存.
天地都會好像衣服那樣漸漸變舊.
耶穌基督會將天地捲起來.
像捲一件外衣.
天地都好像衣服那樣改變.
唯有基督耶穌卻永不改變.
耶穌基督是沒有窮盡的年歲.
當我們這樣讀.
我們就明白了.
我們信靠著那個.
我們在祈禱的那個.
就是一個永恆的主.
他帶領我們走人生的道路.
我們知道了.
我們有什麼回應呢?.
其實這一句.
其實是在回應希伯來書.
最終他要給我們的出路.
希伯來書十三章第七至九節.
我們一起看看這幾節的經文.
從前引導你們全神的道.
給你們的人.
你們要紀念他們.
效法他們的信心.
回顧他們為人的結局.
耶穌基督昨日.
今日一直到永遠是一樣的.
你們不要效法種種怪異的教訓勾引了他.
因為人的心.
靠恩典德堅固才是好的.
並不是靠飲食.

$^{641}$那在飲食上用心的.
從來沒有得到益處.
這個經文他在說什麼呢?.
其實這個是結論部分.
整個希伯來書.
也是用剛才永恆主的神學觀念.
他告訴基督徒.
我們今日這個信仰是有價值的.
那個價值不是在於你.
只是現世活得舒服一點.
或者令你覺得幸福一點.
這個他不否認的.
不過如果我們將我們信仰的重點放在.
我們今生眼前的信仰生活.
我們的禍與福眼前的那些際遇.
如果我們只是帶著這種信仰觀念.
那種比較表面浮淺的信仰.
去到人生某些階段是不夠用的.
如果我們對基督耶穌.
對我們的信仰.
對我們的人生體會.
跟我們對基督耶穌的認識不配合.
有一天我們去到某些位置.
我們會覺得沒有意思.
我這樣說很複雜.
我再舉個例子就明白.
我都在黃浸有大概19,20年的事奉日子.
我認識有些弟兄姊妹.
在年輕的時候見到他回來教會.
很熱心.
但到他成長的時候.
到他進入中年.
一個一個開始離開信仰離開教會.
他心底裡說我都會信的.
我都會問他你們信不信.
他說他都會信.
但我說你會不會投入在信仰群體.
你會不會繼續回教會.
他們就灑手伶頭.
我問他為何會有這樣的狀況.

$^{681}$他說 講得直白一點.
他說有時他覺得教會很低能.
大家這樣說你會覺得.
為何會這樣說.
他說有時我們的形態.
還是想用中學生的形態.
來牧養一個成年人.
他不是批評我.
不過他講整個有時的形態.
我們經常都很想去中學生的生活.
他說我人大過了.
我追求的不是這些.
不是玩幾個遊戲.
我就心裡可以有啟發.
我很想知道究竟聖經對我生命有甚麼重要.
但很多時我們的信仰.
我們就徘徊在很浮淺的地步.
偶爾有些人情味是開心.
但面對生命的困難承托不起.
當然我對於他所講的判斷.
我有些不同意.
我覺得不是完全由他偏頗的角度去講.
不過有一樣東西都值得我們去問.
我們今日的信仰生活.
會不會經歷一種浮淺的關係.
我們對上帝的認識.
我們都徘徊在初信.
在我年青時我認識很深的時候.
我們就停止了.
我們沒有辦法繼續向前.
我們沒有辦法.
我們自己個人的成長經歷加深.
但我們的信仰的理解我們沒有理解深入.
以致我們出現了這種落差.
我們對信仰的理解.
承托不了我們走面前信仰人生的路.
有神學家講過.
今日信仰最大的敵人是什麼呢?.
不是信仰本身.
而是浮淺的認知.

$^{721}$如果我們人生徘徊在浮淺.
這個信仰既沒有解釋能力.
也沒有辦法帶我們在生命很深度的困苦.
或者面對人生的壓力當中.
能夠帶來新的出路.
唯有重新找回信仰的根底.
更加深入進入信仰.
更加明白基督耶穌的生命.
我們才有力量走面前人生的路.
希伯來書的作者.
其實他都是用這樣的角度.
整本希伯來書都是講這件事.
他講為什麼信徒會垂流失去呢?.
為什麼教會有什麼問題?.
為什麼外面有多兇猛?.
而是我們沒有下力去竭力追求.
帶領我們永恆不變的主.
當我們認識祂不夠深.
當我們跟祂的關係不夠近的時候.
我們就缺乏力量去應對眼前所有的挑戰和困難.
希伯來書的作者在這裡的結論.
其實上帝一直不變.
變的只是這個世界.
正正就是因為這個世界不停改變.
人生命裡面際遇不斷的來更換.
時局我們沒有辦法控制.
但是我們必須要拿著一個昨日今日一直到永遠.
是一樣的上帝.
世界上有五花八門的道理時刻來影響我們.
但是我們不要將我們的眼目放在五光十色.
五花八門的新理論.
或者這個世界裡面很多的政局.
又或者媒體裡面渲染的不同的信息.
我們反而要將焦點放在這個永活的上帝當中.
唯有這樣我們才有力量走面前人生的路.
在希伯來書裡面也提醒.
或者在問我們一個問題.
今日我們人生在世.
我們是努力生存還是活出真正的生命.
這個問題不容易回答.

$^{761}$因為兩件事可能是重疊的.
我們既是生存也在尋找生命.
很難說完全二分來看這個題目.
不過希伯來書的作者挑戰我們.
你不妨看看在教會裡面我們有一些先賢前輩.
看看他們怎樣憑信心來過生活.
看他們的人生過得怎樣.
我們沒有辦法為自己.
我們怎會有水晶球知道我們自己人生怎樣走.
不過看看他們憑著信心.
那些先賢前輩他們的人生是怎樣.
可能他們的人生會長一些.
可能有些人生會短一些.
但他們帶著信心而活的時候.
他們的生命是精彩的.
他們的生命是亮麗的.
我們教木同工除了近來經常看那套電影之外.
其實我們教木同工也見過很多.
見過很多生死.
最令我們欣慰的是.
在為他們送別的日子裡面.
見到遺體的面容是寬容.
帶著少少的笑容.
這些不是化妝可以化妝得到的.
家人對離世者的愛顧.
生命裡面留下一點點美好的故事.
我見到這些小弟兄小姐妹.
我心裡面求上帝.
如果有一天我像他們這樣就好了.
如果有上帝我們生命可以.
如果沒有上帝我們垂留失去.
可能我們活在這個世界的困苦流離當中.
弟兄姊妹今天你怎樣選擇?.
求主繼續幫助我們.
\newpage



\section{詩篇馬太福音 118:19-26}
\label{sec:_MkLayWrDNI}
\textbf{2024年12月8日崇拜直播｜陳冠成牧師｜詩篇裡的基督－房角石頭｜詩篇118\_19-26;馬太福音20\_33-46}
\newline
\newline
連結: \href{https://youtube.com/watch?v=_MkLayWrDNI}{\texttt{https://youtube.com/watch?v=\_MkLayWrDNI}} ~~~~ 語音日期: 2024-12-09
\newline
\newline
\hyperref[sec:a62m7d20RCU]{\small{< < < PREV SERMON < < <}}
~
\hyperref[sec:index_chronic]{\small{[返順時目]}}
~
\hyperref[sec:jH3ZJLIEHsY]{\small{> > > NEXT SERMON > > >}}
\newline
\newline
$^{1}$今日重讀的經文有兩段.
分別是《馬太福音》21章33至46節.
和詩篇118篇22至25節.
請聽我讀出.
《馬太福音》21章33至46節.
你們再聽一個比喻.
有個家主栽了一個葡萄園.
周圍圈上泥巴.
裡面掛了一個壓酒池.
蓋了一座樓.
租給員戶就往外國去了.
收果子的時候近了.
就打發僕人到員戶那裡去收果子.
員戶拿著僕人打了一個.
殺了一個.
用石頭打死一個.
員主又打發別的僕人去.
比先前更多.
員戶還是照樣待他們.
後來打發他的兒子到他們那裡去.
意思說他們必尊敬我的兒子.
不料員戶看見他兒子就彼此說.
這是承受產業的泥巴.
我們殺他佔他的產業.
他們就拿著他推出葡萄園外殺了.
員主來的時候要怎樣處置這些員戶呢?.
他們說要下毒手除滅那些惡人.
將葡萄園另租給那按著時候交果子的員戶.
耶穌說.
經上寫著.
將人所棄的石頭已作了房閣的頭塊石頭.
這是主所作的.
在我們眼中看為稀奇.
這經你們沒有念過嗎?.
所以我告訴你們.
神的國必從你們奪去.
次級那能結果子的百姓.
誰掉在這石頭上必要奪誰.
這石頭掉在誰的身上.
就要把誰砸得稀爛.

$^{41}$祭司長和法利賽人聽見他的比喻.
就看出他是指著他們說的.
他們想要捉拿他.
只是怕眾人.
因為眾人以他為先知.
詩篇118篇22至25節.
將人所棄的石頭已成了房閣的頭塊石頭.
這是耶和華所作的.
在我們眼中看為稀奇.
這是耶和華所定的日子.
我們在其中要高興歡喜.
耶和華求你拯救.
耶和華求你使我們亨通.
聖經獨筆.
請姐妹平安.
今天已經是張臨奇的第二個主日.
換句話說.
距離聖誕節的日子再近一步.
按照教會的傳統.
張臨奇有兩個重要的意義.
或兩層的意義.
都是與預備有關.
首先是預備慶祝耶穌基督.
第一次降生來到世上.
這一天已經發生了.
至於另一個預備.
是迎接耶穌基督第二次榮耀的日子.
這一天還未發生.
顧名思義.
我們說預備是指.
不是一個被動.
不是一個只在等待的狀況.
其實預備是提醒我們.
在剛才說的兩個日子.
兩個與耶穌有關的來臨日子中間的時間.
即是我們現在身處的時間.
我們如何去審查自己的生命.
我們有沒有用行動來順從配合.
耶穌基督的吩咐呢?.
我們有沒有在生活中.

$^{81}$持續轉向上帝.
不斷為祂結出果子呢?.
其實就好像我們剛才唱的三首歌.
我們一方面慶賀耶穌基督的降生.
亦檢視我們自己有沒有成為世上的光和炎.
來預備耶穌基督再來的日子.
不過,知道歸知道.
我們看聖經.
它告訴我們.
不是所有上帝的子民.
都會懷著這種警醒預備的心情.
來回應耶穌的道理.
事實上我們看聖經.
我們很明白.
兩千多年前.
當耶穌來到的時候.
當時猶太人的靈性.
我們說是處於一個低谷.
他們的信仰生活.
開始形式化,僵化了.
道德亦日漸敗壞.
表面上每一個人都口稱.
我們很期盼上帝光照我們.
期盼那位救主來到.
但實際上你又會看到.
在他們的骨子裡已經習慣了.
活在黑暗裡.
習慣了以自己的心思來過活.
即使他們有宗教的生活.
甚至乎即使明知.
耶穌就是上帝坐差來的拯救者.
不過他們都會否定耶穌.
甚至乎處處和耶穌作對.
千方百計要將耶穌除滅.
其實就好像我們剛才所看到.
《馬太福音》那段經文.
這個我們說是空樂園戶的比喻.
我想在座很多人都聽過.
正正是耶穌進入耶路撒冷之後.
祂面對當時的宗教領袖和群眾.

$^{121}$質疑祂的權柄到底是從何而來.
於是耶穌就說了這個比喻.
如果你看這個比喻.
其實有祂舊約的背景.
其實是直根在《以賽亞書》第五章.
它裡面說到.
上帝就是葡萄園的園主.
對葡萄園.
其實就是在說以色列.
對葡萄園悉心的栽種建立他們.
然後交給園戶即是以色列人.
每一代的以色列人來打理葡萄園.
並且按時候.
上帝會差派他的僕人.
即是歷世歷代的眾先知門.
向他的百姓來收取果子.
提醒百姓是有責任向上帝交賬的.
因為納了藥.
不過以色列人卻拒絕.
進行配合上帝的吩咐.
並且好像那個比喻所說.
一再惡待上帝所差派來的僕人.
不過即使是這樣.
你會看到上帝是怎樣.
上帝沒有以惡還惡.
上帝仍然是繼續一再寬容.
去打發更多的先知.
來到境界的百姓.
很希望他們可以.
真是有一天會回心轉意.
不過比喻這裡繼續說.
他們照舊糟蹋上帝的愛心.
繼續苦待上帝所差派的先知.
到後來比喻裡的家主.
他就決定.
既然先知和僕人都搞不定.
我派我自己的兒子親自出馬.
家主一心去想.
我兒子就是我的繼承人.
和我一樣擁有相同的權柄.

$^{161}$見到我兒子.
代表見到我本人.
他們怎樣壞也好.
怎樣壞也好.
怎樣不給我差去的僕人面子也好.
怎樣都應該尊敬我兒子.
怎樣都應該尊敬我.
我們看到家主這種既往不咎.
仍然對兇惡的元戶.
深存那份盼望的態度.
現實裡你應該見不到.
注無僅有.
但卻正正反映上帝對子民的寬容.
的真實寫照.
上帝真的這樣對待他的百姓.
但沒想到.
百姓不是這樣回應上帝.
比喻這裡說.
元主沒有因此而迴轉.
反而是動起惡念.
不如我們趁機殺了元主的兒子.
侵吞他的家產.
原來按照當時的習俗.
假如作為地主的.
他連續四年都無法向他的租戶.
收取租金.
某程度上他就要被迫放棄.
他的產業.
他要放棄那塊地.
收不到.
就好像被人逆權侵佔.
另外如果土地的承繼人.
長期的失蹤甚至死亡.
最早霸佔土地的人.
就自動可以享有那塊地的業權.
換句話來說.
比喻裡面這些兇惡元戶.
大概以為家主離世了.
所以現在他的兒子.
以承繼人的方式.

$^{201}$向我們問責.
要求我們交出果子.
我們只要殺了他.
就沒有王館.
我們就可以將葡萄園據為己有.
於是就像比喻所說.
他們一心下狠手.
將家主的兒子拿著.
推出葡萄園外殺了.
什麼意思呢?.
為什麼要推出葡萄園外?.
不只是不要弄髒葡萄園這麼簡單.
推出葡萄園外.
將家主的兒子殺了.
即是將命案裝扮成.
不關我們事.
他在外面被人殺了.
是其他人做的.
以掩飾自己這個邪惡的行為.
耶穌暗示當時他面對的宗教領袖.
稍後也會像比喻的兇惡緣故那樣.
對付他們.
並且借助外邦人的手.
將上帝的兒子殺害.
用一個邪惡的手法.
來保住他們在宗教上的地位和權柄.
但是耶穌在比喻的結尾.
斬釘截鐵說清楚.
他們不會逍遙法外.
不會得逞.
於是耶穌就問一個問題.
當那個玄主來到的時候.
要怎樣處理那群玄戶.
他問了所有的聽眾.
我們看到猶太人的宗教領袖.
在群眾眾目睽睽的情況下.
他們一定要說出一個正確的答案.
他們說那些玄戶這麼壞.
當然要狠狠地消滅那群惡人.
並且將葡萄園交給那些.

$^{241}$按時候結果指的好的玄戶.
你看到耶穌也很厲害.
耶穌是間接逼那些要對付他.
要謀害他的人.
你自己定下自己的罪.
你自己預告自己的結局.
來承認自己所作所為.
其實是極之邪惡.
一點都不配得.
所以我們要求神帝再寬容.
再饒恕.
所以某程度上.
剛才這個比喻.
空惡玄戶的比喻.
其實也揭示了人性最幽暗醜陋的一面.
就是當人明知自己做得不好.
明知自己理虧.
但卻拒絕回轉.
拒絕悔改的時候.
真的可以變得很邪惡.
不單止.
他會對那些前來善意提醒的人.
感到煩厭.
感到憎惡.
還會到一個地步.
你明明是好意提醒他回轉.
他會滅你聲.
他會將你趕走.
甚至乎像這個比喻.
他要讓你消失在地上.
玄戶不單止沒有盡到自己的本份.
為家主結出好的果子.
反過來將所有前來勸戒的人.
通通解決.
就等同於耶穌時代猶太人的宗教領袖.
不單止拒絕自己.
在信仰裡出了問題.
反過來推卸責任.
說誰搞亂黨.
耶穌你搞亂黨.

$^{281}$你破壞我們一向行之有效的規矩.
於是乎屈往.
問題在耶穌那裡.
不在我們那裡.
我們要設法將它消滅.
等同於今天.
可能你也聽過.
見過一些人做事的手法.
只要我壓制消滅了那些.
指責我有問題的人.
那我又怎樣?.
我便沒有問題.
手法同出一轍.
弟兄姊妹.
耶穌這個比喻.
我想給我們每一個人.
有一個思考的機會.
好像一條哲學問題.
到底人可否透過不斷否認.
遮掩自己的過失.
最終會得到一個好的下場和結果呢?.
可不可以呢?.
在兇惡的緣戶.
或敵黨耶穌的人.
他們認為可以的.
惡行是可以結出善果的.
他們真的這樣想.
所以才這樣做.
他們覺得惡行真的可以帶來好的下場.
但聖經清楚告訴我們.
這是錯謬的想法.
這種想法只會讓人.
越踩越深.
陷入罪的深淵.
但值得注意的是一件事.
比喻當中的緣戶.
即以色列人.
其實他們本身的條件好不好?.
他們本身是得天獨厚.
上帝揀選的子民.

$^{321}$他們有一個很好的開始.
他們也曾答應上帝和上帝立約.
他們會好好打理上帝的國度.
按時為上帝結出良善的果子.
但為什麼最後仍然遭上帝的懲罰.
被上帝嫌棄.
這就值得我們思考.
亦讓我們了解一個情況.
原來最初我們的信仰條件.
無論多麼優厚.
不保證我們最終有份於上帝的國.
關鍵在於我們信主之後的表現.
我們有沒有配合上帝的心意.
我們有沒有活出耶穌基督的吩咐.
確實我們有份於上帝的國.
我們說完全出於上帝的恩典.
但如果我們只恃著上帝給我們恩典.
我們就放任沉醉在自我的世界當中.
某程度上我們是在揮霍上帝的恩典.
我們是在糟蹋上帝的愛心寬容.
我們是在漠視上帝一再的境界.
上帝的救恩.
如果這樣持續下去.
其實我們可能跟比喻當中的緣故.
跟耶穌時代的法理錯人或文士.
結局可能都會相似.
就正如施浸若汗.
當日他呼籲人預備好自己的生命.
迎接耶穌基督的來到.
我們每一個人都要時刻去審查.
主動審查自己的言行.
是否與所蒙的恩來相稱.
這樣做會驅使我們的心靈.
由遠起來轉向上帝.
而這種轉向也帶動我們.
不會只停留在做一個星期日的信徒.
他會迫使我們回到生活中.
在每一個場景中.
我們都積極尋求上帝的心意.
在我們一些生活細節中.

$^{361}$我們都努力配合耶穌基督的吩咐.
來建立自己和別人的生命.
所以當耶穌說完這個比喻時.
他引用了一段舊約經文.
就是我們剛才讀的第二段詩篇經文.
耶穌引用當時猶太人非常熟悉的.
詩篇一八一十八篇.
他沒有全部引用.
他只引用了二十二至二十三節.
「將人所棄的石頭.
以作了房閣的頭塊石頭.
這是主所作的.
在我們眼中看為稀奇」.
「將人」其實是說.
當時面對敵黨的宗教領袖.
他稱呼他們「將人」是怎樣?.
這班人是受過訓練的.
這班人是得天獨厚的.
對聖經滾瓜爛熟.
對律法更是倒背如流.
他們應該有慧眼.
一眼就看出耶穌是誰.
應該看到耶穌是上帝所差派來的拯救者.
但是因為心硬.
不是知識的問題.
而是心的問題.
他們遮蓋了本來應該有的.
那種屬靈的洞察力.
認為耶穌只是一個普通碎料的人.
耶穌只是一塊普通無用的石頭.
所以就嫌棄.
這就是詩篇裡所說.
「將人所棄的石頭」的意思.
但上帝很特別.
上帝是人棄我取.
上帝很奇妙地.
你們嫌棄的那位.
你們覺得不起眼的那個.
我正正是揀選了他.
成為房閣裡的頭塊石頭.

$^{401}$什麼意思呢?.
頭塊石頭其實.
在古代建築.
我們知道不是像今天那樣.
會打樁的.
那怎樣去訂立根基呢?.
就是靠一塊所謂的房閣石頭.
在安放根基的時候.
猶太人會在一座建築物.
最外圍的角落.
就會放一塊基石.
就是經文所說的房閣石.
它不單止是整座建築物裡.
最大,最堅固,最精細的一塊石頭.
其實它同時也規範了.
之後那座建築物要怎樣建築.
換句話來說.
這塊房閣石一定在那角落之後.
所有之後加上去的石頭都要怎樣呢?.
都要跟它拼在一起.
都要跟它對準.
你不可以自己起歪了.
或者說斜了.
你要跟它對準.
這個房閣石.
其實就像我們所說.
一個準繩一樣.
如果你有做從事裝修.
你可能見過.
現在很先進的.
有一個像雷射的工具.
一放在地上.
左右兩邊直角三邊.
射出綠色線.
好像雷射一樣.
就好像這樣.
你要對準.
你要依著這塊房閣石來建造.
換句話來說.
房閣石制定了整座建築物之後的發展方向.

$^{441}$而耶穌就表示.
祂就是安放在以色列那塊房閣石.
百姓所有的信仰工程.
都一定要經過耶穌檢驗.
那些不合乎聖經真理的言行思想.
最後耶穌都會怎樣呢?.
揭出來.
甚至推倒它.
但是猶太的宗教領袖.
卻拒絕耶穌建造以色列的房閣石.
他們想怎樣呢?.
就是要在耶穌以外.
他要另立一個信仰的根基.
結果.
耶穌這塊房閣石就成為了反對者的.
都是石.
不過是半腳石.
變成了擊跌他們的石.
他們不遵從耶穌的吩咐.
他們想撞倒耶穌.
踢走耶穌.
但反而怎樣呢?.
一踢到耶穌身上.
自己被擊跌.
自己跌傷.
所以耶穌說.
誰掉在這塊石頭上.
必要跌碎.
這塊石頭掉在誰的身上.
就要把誰砸得稀爛.
即是說所有敵黨耶穌的人.
最終都要面臨審判和定罪.
並且上帝的國.
要從他們身上收回.
賜給誰呢?.
就是賜給我們.
賜給教會.
賜給新約的信徒.
由他們繼續代表上帝.
向外結果治.

$^{481}$換句話來說.
今天這段經文很關我們事.
既然我們在新約下.
領受了耶穌基督的使命.
成為新的元戶.
要為上帝的按時結果治的時候.
我們需要問責嗎?.
我們是需要.
那個空岳元戶的比喻.
其實也與我們有關.
只不過我們今天可以有選擇.
我們不一定要成為空岳元戶的結局.
你可以有不一樣的結局.
就是你願意配合耶穌這塊防角石.
來按時為上帝這位元主結果治.
耶穌一定會檢驗.
我們今天每一個人的信仰工程.
到底是否符合祂的標準.
有沒有結出美善的果子.
並且耶穌說祂是那塊最堅硬的防角石.
你不要想著動祂.
動不了.
耶穌不會因為人的罪.
軟弱人的惰性.
而強行降低祂的標準來遷就人.
但耶穌會怎樣?.
耶穌會幫我們達標.
祂會給我們恩典.
祂會加力給我們.
祂會與我們同行.
一起去完成上帝的吩咐.
所以弟兄姊妹.
經文這裡很清楚.
我們都明白.
唯有依靠著耶穌的吩咐.
人生才會建立得穩妥.
否則一定會面對建築物傾斜.
崩壞甚至倒塌的危機.
所以耶穌會定期.
用我們身處的環境.

$^{521}$來提醒我們.
來搞動或改變我們的生命.
看看我們有沒有緊靠祂的防角石.
事實上大家身處香港這幾年.
你覺得容易嗎?.
這個問題不用回答.
大家都反映在你面前.
難道很容易才走過幾年的疫情.
以為會好起來嗎?.
誰知這幾年經濟還是一樣.
香港人最關注樓市股市.
好像有一點點變化.
但沒有什麼盼望.
再加上明年.
誰會上場?.
特朗普上場.
大家擔心會否加關稅.
貿易戰影響香港前景.
很多負面的事物.
於是在這些狀況下.
人一定會將防角石.
移到自己身上.
要想如何賣樓.
轉移資產.
避風險.
早幾年很多移民.
企圖保住所擁有的東西.
企圖找到前景的安逸.
當然有些人覺得.
先看清楚.
按兵不動.
觀望一下.
可能會好起來.
也不奇怪.
坦白說.
我們沒有預言的能力.
我們手上沒有未來的水晶球.
沒有人可以通過自己的籌算.
百分百找到保障.
這是事實.

$^{561}$但這事實.
到底是迫使我們.
更加自己活生生地活生生命.
還是像今日經文提醒.
看看有沒有靠著耶穌的防角石.
你依靠的到底是什麼?.
事實上作為上帝的管家.
我們要為未來計劃籌算.
但同時我們也不能忘記.
誰才是真正生活的保障.
如果我們認定.
耶穌就是那個防角石.
祂最穩妥.
祂最堅固可靠的時候.
就要靠著祂.
一切籌算都要問過祂.
都要讓祂告訴你.
應該怎樣做.
事實上.
聖經告訴我們.
信靠耶穌的人.
不一定事事順利.
什麼都贏.
每一個決定一定是贏.
但信靠耶穌的人.
聖經應許你有平安.
你有那份安穩.
你會經歷到上帝的恩典.
就像剛才我們讀那段詩篇最後說.
你會亢通.
亢通不是說我們做人生勝利組.
亢通是你在上帝的道路裡.
你會走得通順.
你會問心無愧.
因為你知道.
你正在進行上帝的旨意.
因為你知道你願意讓耶穌帶領你.
調教你人生的方向和步伐.
所以弟兄姊妹.
耶穌呼籲我們要單純地.

$^{601}$專一依靠祂.
效法祂.
當然意思不是說.
我們人生之後不需要再籌算.
而是相反.
你在籌算的同時.
你要認定上帝的安排才是最穩妥.
所以你可以放心靠下去.
並且你可以在那個處境裡.
學習分別為勝.
我想我們身邊的人.
很多未信主的人.
有著我們剛才所說.
同樣對前景的憂慮.
你又憂慮.
他又憂慮.
但是我們有什麼分別?.
信耶穌的人在這裡有什麼與別不同.
那個就是聖經經常所說的分別為勝.
當人依靠自己.
你卻是相信耶穌的時候.
那個就成為你的見證.
成為你的安穩.
也可以將這份安穩.
傳遞給你身邊的人.
另一方面.
專心跟隨.
依靠耶穌.
並不是說.
你放棄你所有人生的理想.
你變成一個沒有大志的人.
而是你做任何事都好.
什麼決定都好.
你都先考慮.
能不能榮耀上帝.
能不能討上帝的喜悅.
以這個為一個目標.
專一的跟隨耶穌.
也不是說我們以後要遠離世俗.
我們過一個隱居的生活.

$^{641}$反過來.
當你進入世界.
面對那些不公義的事.
面對上帝不起悅的事.
甚至乎像耶穌那樣.
面對那些惡人,惡事的時候.
你會不會像耶穌那樣.
效法他.
仍然堅持上帝看為對的價值.
事實上.
當我們看回剛才所讀的詩篇.
118篇.
它有一個背景.
就像程序表裡面.
港澳大綱印了在那裡.
我們看多一點.
10至14節.
我想大家讀一讀.
這段詩篇的10至14節.
了解當時詩人的背景.
請大家幫我讀.
預備起.
(國民黨為義和和.
國家的貴婦和貴妃.
必消滅貪污.
貪污難饒我.
我乃君王.
我靠耶和華的力.
必消滅貪污.
貪污雖無功.
雖無能.
但我將燒荊棘的火.
必被你消滅.
我靠耶和華的力.
必消滅貪污.
你推我.
要叫我跌倒.
但耶和華幫助了我.
耶和華是我的力量.
是我的事多.

$^{681}$阿也成了我的前家.
好,謝謝.
最後那句.
我們待會回應詩都會唱.
你會看到其實這段.
這幾次的經文.
是一個很清晰的對比.
是不是?.
有很多人要消滅詩人.
但是詩人一再強調.
上帝會幫助我.
上帝會幫我反過來對付他們.
是不是?.
讓我們看到信仰.
有它敵對的一面.
信仰不只是做一個好人.
信仰是說你堅持.
上帝看為對的事.
而在這個情況下.
確實有機會得罪.
那些敵擋上帝的人.
令到你自己也很不安.
所以看回這幾節的詩篇.
一方面反映了詩人那種.
四面受敵的情況.
是不是?.
身邊每一個人.
他堅持走上帝的價值觀.
所以攻擊他.
討厭他.
這就是耶穌基督當時面對的困境.
沒有一個人看好耶穌.
沒有一個人尊重他.
彷彿全世界都像詩篇所說.
要他跌倒.
要他失敗.
甚至乎哭著要消滅他.
在這個情況下.
弟兄姊妹你會怎樣選擇?.
容易的路是立即放棄妥協.

$^{721}$真的.
立即埋隊.
靠邊站.
這樣就解除威脅.
這是很實際的事情.
堅持反而是困難.
堅持.
在一片反對聲音裡.
堅持上帝看為對的事.
就會好像詩人所面對.
耶穌所面對的一樣.
但為什麼耶穌仍然堅持.
上帝看為對的事.
是因為他知道.
最終上帝會得勝.
最終上帝會讓他得勝.
正如詩人所說.
那些為上帝堅持.
倚靠上帝的人.
神會替他們平反.
並且去懲治.
對付敵黨上帝的人.
耶穌和詩人一樣.
都知道最後得勝的是.
站在上帝那邊的人.
堅持上帝公義的人.
所以他認定.
父神就是他的力量.
他的詩歌也成了他的拯救.
這句宣告成為了他.
繼續有勇氣,有力量.
堅持下去的一個憑據.
弟兄姊妹.
耶穌之所以能夠.
坦然面對反對敵黨的聲音.
堅持上帝看為義的事.
絕不讓步.
正正是一種極致的信仰表達.
事實上.
耶穌這種對真理的持守.

$^{761}$我想也是我們面向世界.
需要效法的.
確實你看回.
世上有不少價值觀念.
其實直接或間接挑戰.
衝擊我們的信仰.
這方面我想不需要多說.
大家都很明白.
他會游說我們.
妥協一點.
開放一點.
包容多一點.
不要太堅持.
不要太執著保守.
例如圍繞西方世界.
這幾年關於性別的議題.
關於濫用藥物或毒品的問題.
甚至在這些國家裡越演越烈.
因為這些問題.
被包裝為一種進步的價值觀.
將性別的混淆.
毒品的濫用.
包裝為道德進步.
人生的覺醒.
這樣很亮麗.
於是你不跟進是落伍.
你反對呢?.
反對就是歧視.
你很難跟他對抗.
因為高舉自由.
高舉自我.
其實這股勢力一直都在.
確實基督徒要接觸這個世界.
要處理不同的人際關係.
很自然我們會顧及所謂人情世故別人的感受.
聖經告訴我們.
若善能行.
總要與人和睦.
但聖經也同時提醒.
在生命裡總會有些事情.

$^{801}$你要為主堅定堅守.
正如我上次提到.
近代中國基督徒.
很重要的領袖王明道.
我今天想再說多一點.
關於他的事跡.
王明道是一位.
神學立場極之保守的人.
他是一位作風很硬朗的教會領袖.
有學者怎樣形容他呢?.
現代版的施浸若寒.
你就知道他多麼雷厲風行.
在信仰方面很硬朗.
他面對罪惡.
面對否定真理的人.
他絕不手軟.
他不會畏懼.
毫不妥協.
所以王明道這種.
這麼硬朗的純靈生命.
的確造就了上世紀.
無數中國基督徒的靈命.
不過同樣.
他最被人批評得多.
最被人討厭的.
亦是他這種過份硬朗.
不可以兼容別人.
很容易得罪人.
這種想法和態度.
是不是很有趣?.
儘管如此.
絲毫不影響.
王明道對中國近代基督徒的影響力.
很多學者都相信.
他的講道基本上.
大部分內容都不是很學術.
他本身不是一個很卓越學術的學者.
但他卻是一個.
最受人敬重的其中一位基督徒.
一位的屬靈領袖.

$^{841}$在他的講道裡.
基本上道德的議題.
佔了絕大部分.
他會指導當時.
那些沒什麼學識的信徒.
如何具體地活出.
美好的屬靈品格.
很多人亦因為.
我見到這個人的品格.
我見到他的言行.
他真的說得出做得到.
因為他這種對信仰的執著和堅持.
就是這樣相信耶穌.
事實上,弟兄姊妹.
我想生活裡.
我們當然不想被別人討厭.
有沒有人想不斷被別人討厭?.
沒有的.
我們也不想被別人批評.
但如果我們只顧著.
過一個受人愛戴.
受人歡迎的基督徒生活.
這樣我們的信仰生命.
遲早會不自由.
遲早會越來越困難.
假如我們只對周邊的人.
監貌辨識.
看人面色做人.
我們因為這樣而要吞下.
要說的真理.
一聲吞下肚子不說.
我們要忍住不為耶穌做見證.
如果一切只是為了迎合別人.
或這個世界的口味而活.
就好像我們唱第三首歌.
做不到光和炎.
反過來.
世界會把我們重重捆綁.
所以.
經文勉勵我們要效法耶穌.

$^{881}$要有勇氣堅守上帝看為對的事物.
當然我們每一個都渴望.
被人認同,被人肯定.
我們沒有人想得罪別人.
沒有人想被人討厭.
但剛才的詩篇.
或者耶穌基督面對.
一些反對聲音的時候.
他無可避免.
或者我們無可避免.
為了上帝的緣故.
為了真理的緣故.
我們有需要.
被這個世界.
被你身邊的人.
來討厭一下.
我們有這種覺悟.
要培育這份勇氣.
確實我們不容易去改變.
那些批評,討厭耶穌的人.
但我們亦不應該.
因為這樣而影響.
我們跟隨耶穌的決心和信心.
雖然按照社會規則.
遊戲法則.
那些單純為上帝堅持的人.
好像是比較吃虧.
好像很不安.
因為很多人可能會不滿意.
跟這個世界所說的.
那種成功,美好的生活條件.
距離好像越來越遠.
但如果按照耶穌的生命之道.
祂說你持守上帝看為對的事的時候.
你是最終不會吃虧的.
反過來你會得到耶穌基督的看顧.
祂願意去投靠那些.
凡信靠祂的人.
所以弟兄姊妹盼望.
上帝藉著這段經文給我們勇氣.

$^{921}$讓我們能夠在你的生活處境裡.
為上帝去堅持.
依靠耶穌基督這塊防國盤石.
來為耶穌,為神.
而出見證.
讓我們同心祈禱.
天父上帝願意我們藉著這段經文.
好好來預備.
不單單慶祝耶穌基督的降生.
也去迎接耶穌基督再來的日子.
讓我們到時可以贏得耶穌的稱許.
得著你稱我們為忠心良善的僕人.
獲得上帝你榮耀冠冕的賞賜.
為了這份盼望,這份期待.
我們在當下堅持.
上帝你捍衛對的事情.
謹守上帝你託付給我們每一個的崗位.
哪怕在裡面可能真的會遇到挑戰,衝擊.
甚至不認同的聲音.
想我們妥協的聲音.
求主你讓我們不單單有那份智慧.
柔和謙卑的心腸.
也有從上帝你而來的勇敢.
為上帝你堅持.
願以耶穌基督.
你的生命成為我們效法的榜樣.
讓我們依靠著你這塊堅固的防國盤石.
來建造我們的生命.
來籌算我們的人生.
以致我們在動盪的世代.
我們可以經歷到上帝你所賜的應許.
經歷到上帝你給我們各樣的恩典.
保守和盼望.
我們這樣的禱告.
是奉耶穌基督你的名來祈求.
阿們.
\newpage



\section{約翰福音 2:1-12}
\label{sec:jH3ZJLIEHsY}
\textbf{2023年2月19日崇拜直播｜梁家麟牧師｜以水變酒的神蹟｜約翰福音2\_1-12}
\newline
\newline
連結: \href{https://youtube.com/watch?v=jH3ZJLIEHsY}{\texttt{https://youtube.com/watch?v=jH3ZJLIEHsY}} ~~~~ 語音日期: 2024-12-12
\newline
\newline
\hyperref[sec:_MkLayWrDNI]{\small{< < < PREV SERMON < < <}}
~
\hyperref[sec:index_chronic]{\small{[返順時目]}}
~
\hyperref[sec:KLXa_sv_cYI]{\small{> > > NEXT SERMON > > >}}
\newline
\newline
約翰福音 2:1-12
\newline
\begin{longtable}{cl}
\hline
\hline
章節 & 經文 (和合本修訂版)\\
\hline
2:1 & \begin{tabularx}{0.7\textwidth}{X} 第三日，在加利利的迦拿有一個婚宴，耶穌的母親在那裡。 \end{tabularx} \\ \\ \relax
2:2 & \begin{tabularx}{0.7\textwidth}{X} 耶穌和他的門徒也被請去赴宴。 \end{tabularx} \\ \\ \relax
2:3 & \begin{tabularx}{0.7\textwidth}{X} 酒用完了，耶穌的母親對他說：「他們沒有酒了。」 \end{tabularx} \\ \\ \relax
2:4 & \begin{tabularx}{0.7\textwidth}{X} 耶穌說：「母親，我與你何干呢？我的時候還沒有到。」 \end{tabularx} \\ \\ \relax
2:5 & \begin{tabularx}{0.7\textwidth}{X} 他母親對用人說：「他告訴你們甚麼，你們就做吧。」 \end{tabularx} \\ \\ \relax
2:6 & \begin{tabularx}{0.7\textwidth}{X} 照猶太人潔淨禮的規矩，有六口石缸擺在那裡，每口可以盛兩三桶水。 \end{tabularx} \\ \\ \relax
2:7 & \begin{tabularx}{0.7\textwidth}{X} 耶穌對用人說：「把缸倒滿水。」他們就倒滿了，直到缸口。 \end{tabularx} \\ \\ \relax
2:8 & \begin{tabularx}{0.7\textwidth}{X} 耶穌又說：「現在舀出來，送給宴會總管。」他們就送了去。 \end{tabularx} \\ \\ \relax
2:9 & \begin{tabularx}{0.7\textwidth}{X} 宴會總管嘗了那水變的酒，並不知道是哪裡來的，只有舀水的用人知道。於是宴會總管叫新郎來， \end{tabularx} \\ \\ \relax
2:10 & \begin{tabularx}{0.7\textwidth}{X} 對他說：「人家都是先擺上好酒，等客人喝夠了才擺上次的，你倒把好酒留到現在！」 \end{tabularx} \\ \\ \relax
2:11 & \begin{tabularx}{0.7\textwidth}{X} 這是耶穌所行的第一個神蹟，是在加利利的迦拿行的，顯出了他的榮耀來，他的門徒就信他了。 \end{tabularx} \\ \\ \relax
2:12 & \begin{tabularx}{0.7\textwidth}{X} 這事以後，耶穌與他的母親、兄弟和門徒都下迦百農去，在那裡住了不多幾天。 \end{tabularx} \\ \\
[1ex]
\hline
\hline
\end{longtable}
$^{1}$今日的經文是在約翰福音第二章一至十二節.
是在新約一百二十七二.
請聽我讀出.
第三日,在加利利的迦那,有取親的筵席.
耶穌的母親在那裡?.
耶穌和他的門徒也被請去筵席.
走運進了,耶穌的母親對他說:他們沒有走了.
耶穌說:母親,我與你有什麼相關?我的時候還沒有到.
他母親對傭人說:他告訴你們什麼,你們就作什麼.
照猶太人潔淨的規矩,有六口石缸擺在那裡,每口可以盛兩三桶水.
耶穌對傭人說:把缸倒滿了水,他們就倒滿了,直到缸口.
耶穌就說:現在可以舀出來,送給管賢直的.
他們就送了去,管賢直的嘗了那水變的酒,並不知道是那裡來的.
只有要水的傭人知道,管賢直的便叫新郎來.
對他說:人都是先擺上好酒,等客喝足了,才擺上次的.
你都是把好酒留到如今,這是耶穌所行的頭一件神蹟.
是在加里里的加拿行的,顯出他榮耀來,他的門徒就信他了.
這事以後,耶穌與他的母親,弟兄和門徒都下加班去,在那裡住了不多日子.
各位姐妹早晨.
我們教會經過了六十年,我們現在慶祝六十週年,一個重要的傳希的日子.
不過對我自己來說,我一直思想的是,教會最基本的是什麼?.
信徒最需要的是什麼?.
剛才鎮琴也說,一數就有四十年信仰,我猜我們大家可能就算沒有六十年,也有相當的時日.
我們也要問,在這幾十年過去,對我們信仰來說,有什麼是我們抓得最緊,或是必須抓緊的?.
不要因為時間越久,我們的組織架構複雜了,人際關係複雜了.
我們上主學,越上越高深,讀了很多延伸課程,讀了很多複雜的理論.
我們失去了最基本的東西,什麼是我們信仰最必須抓得緊.
我們教會最基本的東西,我們必須抓得緊,這是一個真切的問題.
今天我想跟大家用這一段,其實不是很拉筋的經文,去思考這樣的課題.
加拿大婚宴,耶穌以水變酒的故事,大家很熟悉,熟悉得不得了.
今天還是新春的日子,用這段經文來說是很合理的.
婚宴是一個喜慶的日子,我不知道你們知不知道,聖經有關喜慶的日子,沒有多少個.
聖經不是喜歡記載喜慶的,因為猶太人很奇怪,他們新年是大贖罪日.
他們是悔罪的,不是像我們這樣大鑼大鼓的慶祝,聖經也沒有慶祝的日子.
加拿位於拿撒勒以北約14公里的地方,同屬加利利省.
前面第一張耶穌所呼召的門徒拿旦葉就是加拿人.
我們看約翰福音21章第二節就知道.
加拿不是一個大地方,今天我們去聖地一定會去這個地方.
也是一個小村鎮,非常小,不過相對於耶穌的家鄉拿撒勒.
他還算熱鬧一點,先進一點,多兩個舖.

$^{41}$所以當拿旦葉這個加拿人第一次聽到耶穌出身拿撒勒.
就說拿撒勒還能出什麼好東西呢.
這個明顯就是一百步笑五十步,自己也是蚊帳的村鎮.
不過拿撒勒還小一點,所以看不起拿撒勒.
就像我們旺角人看不起深水Po 人一樣.
雖然旺角也不是很好的地方,總之好像比深水Po 好一點.
可能深水Po 人不同意.
無論如何,約翰在頭兩章記載了五天的故事.
用兩章聖經記載五天的故事,這是第三天的故事.
有一個婚宴在加拿大舉行,根據猶太人米茲娜的決定.
婚禮通常在星期三舉行.
聖經沒有說一對新人是誰,只是說耶穌的母親瑪利亞在那裡幫忙.
耶穌和他的呼召門徒也有出席婚宴.
在十二節說到耶穌的弟兄,耶穌的弟弟也有參加.
這說明應該是合乎統請的.
耶穌的家人,從瑪利亞開始率領所有人都有參加婚宴.
大概新婚夫婦和耶穌一家的關係是很密切的.
不然就不僅僅請他們來喝,全家來喝.
還要求瑪利亞下手幫忙.
當時沒有什麼酒樓食肆,通常是在家裡舉行.
自己買酒買食物,邀請親朋好友來幫忙做菜.
中國傳統農村也是用這種形式舉行婚宴.
新界的圍村也會這樣.
猶太人的婚宴可以綿延幾天,甚至一個星期.
視乎主人家的財力有多大.
延期大概進行了一段時間.
瑪利亞突然跑到耶穌的庚前對他說.
他們沒有酒了,第三節.
婚宴的主人家竟然沒有預備足夠數量的酒.
這真是一個很尷尬和掃興的情況.
為什麼會出現這樣的情況呢?.
是因為客人多了? 主人家不夠錢買酒?.
或者是有人喝多了? 喝大了?.
也有可能,聖經沒有說明,我們不知道實際的原因.
總之無論如何,酒不夠了.
當時農村地區市場不發達.
有錢也不能馬上上街找酒舖買酒.
向鄰居籌酒是可以的.
不過,大概同村的人都離開參加婚宴.
一旦向他們開口借酒.

$^{81}$就會將尷尬的情況凍開.
日後就會成為整條村的笑柄.
你知道農村社會是一個小圈子.
什麼出洋相的小事情都可以張揚.
並且可以說很久.
因為沒有新的話題,可以說一年兩年都可以.
被人笑到臉都黃了.
猶太人和中國人一樣.
很重視一個所謂的恥感文化.
Shame Culture,很講面子.
是不可以丟架的.
所以主人家大概不會想這樣做.
不想向鄰居借酒.
馬尼亞無計可施,只好來找耶穌.
從馬尼亞這麼早就發現這個問題.
你可以看到他在現席不是扮演一個幫忙的小角色.
他是後台的總負責人.
前面我們假設主人家是他的至親.
這也是一個證據.
馬尼亞發現問題之後.
沒有直接跟主人家說.
然後將問題扔給主人家.
你們想吧.
而是將問題抱在自己身上.
希望靜悄悄替主人家解決問題.
免得他們出醜.
也可以看到他和主人家的關係是非比尋常的.
還有我們看到馬尼亞吩咐傭人.
無論耶穌吩咐他們做什麼.
他們都要聽命.
傭人真的聽著馬尼亞吩咐.
耶穌說什麼他們就照做.
也可以看到馬尼亞的地位.
在那個場合的地位.
馬尼亞為什麼要跟耶穌說這個情況呢.
當然他自己在一個無計可施的處境之下.
找自己的兒子說一下.
但他找兒子說什麼呢.
他是來訴苦.
還是他期待他的兒子會出手幫忙解決問題呢.

$^{121}$如果按照常理推測.
我們應該相信馬尼亞只是找耶穌來訴苦.
畢竟耶穌在過去的日子.
從來沒有行過神蹟.
一件都沒有行過.
馬尼亞沒理由突然間.
會相信耶穌會行過神蹟出來.
所以他來找耶穌.
應該只是訴苦吧.
但是如果我們仔細看整個故事的發展.
我們就會相信.
馬尼亞是真的認定耶穌能夠解決問題.
然後來找耶穌幫忙解決問題.
他相信耶穌能夠行神蹟.
即使耶穌沒有直接答應去幫忙這個家庭.
更加沒有說他會行什麼什麼神蹟.
馬尼亞已經自行吩咐那些傭人.
無論耶穌要他們做什麼.
他們都要聽命.
所以耶穌不僅僅是他個人的輔導員.
是問題的解決者.
不僅僅是輔導員.
是問題的解決者.
馬尼亞把主人家的問題告訴耶穌.
耶穌的回應是什麼?.
一句很冒失,很噁心的話.
母親,我與你有什麼相干?.
我的時候還沒有到第四節.
獲本,亦做母親.
不過他也不敢大包大攬.
也要指出母親這個字的原文是婦人.
我們很多聖經學者都想幫助耶穌兜里兜路.
包裝得讓耶穌不會那麼不孝順.
不會那麼忤逆.
把這段經文解釋和改了基本意思.
很多聖經學者都說.
耶穌稱母親為婦人不是一個很特別的情況.
不代表他不尊重母親.
不過在這裡.
你看上雲下理.

$^{161}$我們有百分之一百相信耶穌不太尊重母親.
整段說話是非常不客氣.
耶穌總是在宣示他比瑪利亞有更高地位的時候.
就稱瑪利亞為婦人.
不稱她為母親.
另外一次在十字架上.
他把瑪利亞托付給約翰的時候.
他就稱她為婦人.
不是母親.
我與你有什麼相干?.
聖經是這樣寫的.
不過很多學者當然是繞來繞去.
例如你看啟導本的注釋.
為什麼不讓我來做決定呢?.
不過如果這樣解釋.
當然是對母親不敬.
如果解釋成這樣.
你跟著耶穌說那句.
我的時候還沒有到.
完全不相干.
為什麼讓我決定.
但我的時候還沒有到.
你讓我決定但我不想決定.
我還未當值.
我已經下班了.
你告訴我兩個字怎麼接起來呢?.
所以很多人想幫聖經去分析.
不過其實聖經是這樣分析的.
根據上下文.
這句話真正的意思是.
你為什麼要拿這些來煩我?.
為什麼要拿這件事來煩我?.
所以卡爾遜也說.
說這樣的說話.
我與你有什麼相干.
其實是想將他自己.
跟這件事撇清關係.
推開他.
與我無關.
不是我的問題.

$^{201}$這不是我的事.
我沒有需要做這件事.
所以不是說讓我來做.
不是這樣的意思.
耶穌是故意.
用一個不客氣的口吻去說話.
因為在這裡他要撇清.
他與瑪利亞的母子關係.
他當然是瑪利亞的兒子.
但他也不僅僅是瑪利亞的兒子.
如果他僅僅是瑪利亞的兒子.
僅僅是一個凡人.
他也不會有能力去解決.
瑪利亞所提出的問題.
耶穌在這裡申明了他的神聖地位.
他是至高上帝的獨生子.
他是上帝.
當瑪利亞期待耶穌施行神蹟的時候.
他不可能用母親的身份.
去吩咐耶穌怎麼做.
他只能謙卑地以一個凡人的身份.
來就近上帝.
祈求上帝施行恩典.
在這裡耶穌不再稱瑪利亞為母親.
只是稱她為婦人.
就是這個原因.
弟兄姊妹.
人來到上帝的面前.
必須保持極大的敬畏和謙卑.
這是我們必須注意的態度.
不管你和我跟隨了耶穌多少年.
四十年,五十年.
甚至更長的時間.
我都不可以自認跟他很親切.
再沒有隔膜.
你想找我幫你的兒子.
想辦法去某間名校讀書.
我可以跟你說沒問題.
我和這間學校的校長很熟.
我開口叫他幫忙.

$^{241}$他一定會賣個人情給我.
但如果你想請我為你的病祈禱.
我就不可以說沒問題.
我和上帝很熟.
我開口叫他醫治.
他通常都不會托我手肘.
他不太好意思托我手肘.
我敢這樣說嗎?.
我怕神經.
你和我和上帝都不熟到這個地步.
上帝不會也不需要賣人情給我們任何一個人.
他不欠我們任何東西.
沒有任何東西是本然應得的.
都是他的恩典.
我們只能謙卑地趴在這裡向他祈求.
再沒有其他的方法.
耶穌不客氣地告訴瑪利亞.
他的時候還沒到.
這不是他應該施行神蹟的時間.
他不應該拿這樣的事去麻煩他.
換言之.
耶穌知道瑪利亞不僅僅是來找他套苦水.
而是真的期待他出手幫助.
並且是用非凡的能力.
去解決眼前這個現實的問題.
瑪利亞聽到耶穌的回應.
他並沒有覺得被兒子冒犯了.
他很清楚兒子的性格.
他是一個口硬心軟的人.
批評他兩句.
然後就會出手相救.
耶穌也沒有答應出手.
瑪利亞就轉頭跟用人說.
他吩咐你們做什麼.
你們就跟著做吧.
耶穌還沒有開口答允之先.
耶穌還沒有說他會怎麼做.
瑪利亞已經有信心.
他跟耶穌說的話會得到應運.
並且他準備好.

$^{281}$沒有條件配合耶穌所說的話.
耶穌要求我們做什麼.
我們就跟著做什麼.
耶穌果然真的口硬心軟.
開口吩咐用人.
把放在門口潔淨用的六口石缸都注滿水.
然後再拿出來.
交給前台分庫.
管賢直的送到每一台上.
用人不知道發生什麼事.
只跟著他所說的就這樣做.
我們不知道.
為什麼耶穌不直接將水.
由河或水井吸出來.
交給管賢直.
他說要倒入石缸.
然後再拿出來這麼麻煩.
有人說.
耶穌在這裡做一個神蹟.
就好像在做一個魔術.
他要吩咐人.
將那六個口缸的石缸注滿水.
倒到滿天出來.
看一看,只是水.
裡面沒有酒糟.
沒有任何酒溝.
要顯明全部只是水.
證明神蹟如假包換.
不過故事往後的發展.
我們看到耶穌根本不想張揚神蹟.
他現在不在舞台上玩這個遊戲.
所以他根本不需要做這些證明.
水而已,水而已,水而已.
耶穌根本不想你知道.
根本不想人知道有這個神蹟發生.
所以為什麼耶穌要做這個手續.
我們不知道為什麼.
一切的行動都是用極低調的形式進行.
耶穌沒有說咒語.
沒有做出一些稀奇古怪的手勢.

$^{321}$甚至沒有直接去摸水缸和水.
他只是吩咐用人將水注入缸.
然後再拿出來.
這些都是用人做的,不關他的事.
甚至動人做的時候.
也不知道自己在做什麼.
他們更加不知道一個神蹟只是藉著他們的手而發生.
總而言之.
除了耶穌和瑪利亞之外.
沒有人知道發生了神蹟.
在過程中不知道.
事後,我猜是一段很長的時間.
也沒有人知道.
但是神蹟卻是明明白白地發生.
當用人將水拿出來之後.
給管筵席的人.
管筵席的人,即前台監督.
嘗了一下.
就將筵席主人新郎叫來.
跟他說,人都是先放上好酒.
等客喝足了才放上次的.
你倒一把好酒留到如今.
第十節.
管筵席的人道出了當時宴客的習慣.
在賓客味蕾還敏銳的時候.
先上好酒.
等他們喝到差不多.
口腔都麻木了之後.
再差的酒他們也分不出來了.
總之是倒進去,不是倒進去.
這是省錢的方法.
但是這一場筵席都是反其道而行.
先上不是最好的酒.
然後才上最好的酒.
這是整個故事最後一句話.
有一個毫不知情的管筵席.
跟一個更加不知情的新郎說的.
整個對白說明.
耶穌變出來的不僅僅是酒.
而是上好的酒.

$^{361}$再沒有耶穌的話.
也都再沒有瑪利亞的話.
神蹟就是這樣發生了.
聖經的結語是.
這是耶穌所行的頭一件神蹟.
是在加里利的迦那進行的.
顯出他的榮耀來.
他的門徒就信他了.
這句話其實是廢話.
不是很準確.
對不起,我說聖經不準確.
好像很古怪.
我要說的是.
不要說當時甚至之後.
絕大多數門徒.
是不知道在筵席當下.
發生了這個神蹟.
這個神蹟在四卷呼音書裡.
單單若罕有記載.
馬太,馬可,盧迦都沒有記過這個神蹟.
如果我們相信呼音書是由門徒自己.
或者由他的直傳弟子.
門徒的直傳弟子寫出來的.
為什麼不記載.
大概是因為門徒根本不知道有這件事.
才不記載.
否則怎麼會錯過耶穌的處女作.
第一件神蹟.
第一件神蹟不記載?.
沒理由的.
很難想像怎麼會不記第一件神蹟.
第十件不記載.
記錄一件,無所謂.
第一件不會不記載的.
問題是門徒.
這裡沒有說有多少個.
不過大部分有參加這個婚宴.
他們為什麼不知道這件神蹟呢?.
即使他們沒有看到耶穌做婚婦的動作.
之後耶穌應該在跟他們相處的兩年多.

$^{401}$應該跟他們說一下.
沒有,有理由相信的.
他們根本不知道有這場神蹟發生過.
根本不知道主人家發生過這個危機.
整件事極低調進行.
為什麼約翰不能記載呢?.
我們的推論是.
因為約翰日後是負上照顧瑪利亞的責任.
瑪利亞在他晚年一直跟約翰共處.
可能在這個時候.
瑪利亞年老,當然是說以前的事.
想當年,說著說著就說回這個故事.
這個故事只有約翰知道.
只有約翰有記述.
其他三卷福音書都沒有記述這個故事.
這是一個非常特別的神蹟.
它不僅僅是耶穌施行的第一個神蹟.
也是耶穌所施行的眾多神蹟中.
同類型的唯一一個.
前無古人,後無來者.
耶穌日後施行過很多神蹟.
例如洗疾病得全愈.
洗死人復活.
在水面行走.
五餅二魚餵飽五千人.
很多更加轟轟烈烈的神蹟.
不過跟水變酒不同.
因為多數的神蹟都是增加數量.
延長它的效用.
五餅二魚餵飽很多人.
將一個本來應該受命終結的人延長他的壽命.
但是這個神蹟是改變事物的性質.
水是水,酒是酒.
水和酒是兩種完全不同的東西.
雖然樣子都是白色.
其實那時候的酒不是白色的,是紅色的.
都是液體.
不過就算外表沒有很多不同.
你做過化學實驗就知道.
是兩種不同的化學結構.

$^{441}$怎麼會是一樣的東西呢?.
猶太人不是喝米酒,而是喝葡萄酒.
所以顏色真的不同.
所以耶穌是改變了事物的性質.
使它由一樣東西變成完全不同的另一樣東西.
弟兄姊妹.
耶穌將水變酒的故事.
我們通常第一個很簡單的形容是.
我們生命的改變.
耶穌改變我們的生命.
這個確實是一個最自然的對這個故事的形容.
我常常說,是這樣的.
對我們汪津來說.
這六十年來,我們有很多很多特殊的經歷.
我們教會,都搬過幾次.
我們傳道人又換過很多個.
我們弟兄姊妹,也送走一代.
不過,如果你問我們最值得珍惜.
或者是什麼東西使我們的教會能夠一直延續到如今.
一定不是我們背後那麼厲害的屏幕.
不是我們教堂的各樣的設施.
而是我們這裡有很多生命改變的故事.
你和我為什麼回來教會.
一直回來教會.
不僅僅是因為這裡有班朋友.
可以大家敘舊一下.
而是我們確實經歷到生命改變.
我們信耶穌.
我們的生命真的不一樣.
我們過去有性格暴躁.
因著信了耶穌的緣故,變得溫柔了.
我們本來有比較消極悲觀的.
信了耶穌之後,四萬人笑口常開.
我們本來不接納自己,很自卑.
信了耶穌之後,腰部也挺直了.
我們本來夫妻關係很糟糕.
家庭有很多危機.
因為信了耶穌的緣故,而改善了.
沒有耶穌基督的水.
就不會變成今天有永恆生命的酒.

$^{481}$就好像加拿大水變成酒的神蹟一樣.
一切其實是靜悄悄發生的.
這裡沒有人說咒語.
沒有人做特別的手勢.
我們只在這裡如常的敬拜.
如常的團契.
學習聖經,一起祈禱.
一起彼此服侍.
但是神蹟就這樣不知不覺的發生了.
如果不是在浸禮,或者其他場合.
聽頂子妹,例如在睡晚除夕.
聽頂子妹講她的見證故事.
有時候我們自己也不知道發生了這些神蹟.
這裡人頭湧湧.
我們如常吃喝,參加婚宴的.
他們根本不知道有神蹟發生.
只是覺得上一輪的酒沒有現在這一輪的好喝.
有時候還會埋怨管筵席的不按規矩上酒.
次序混亂.
也有像傭人或者管筵席的.
我們在不同的崗位擔當服侍的工作.
努力做好份內的角色.
倒酒的倒酒,布菜的布菜.
人家叫我做什麼就做什麼.
但卻不知道其中有什麼巧妙.
有什麼背後的故事發生了.
現在我們如常參與運作這間教會.
有人現在教主日學.
有人在做兒童工作.
有人在這裡領唱.
有人做其他音樂的服侍.
有人在外面做接待.
你覺得你做的事很有價值嗎?很重要嗎?.
其實也不是很重要.
我們只是做自己的一小部分.
做好我們的角色也很重要嗎?.
也不是很重要.
我們知道不是很重要.
布羅在這份書中提到四重或五重最重要的角色.
他所賜的有仕途,有先知,有傳福音,有牧師和教師.

$^{521}$他沒有提到這五重角色各自有什麼重要的使命.
各自的使命是不重要的.
整體整個教會的使命只有一個.
就是成全聖徒,各盡其職,建立基督的身體.
就算你扮演好主席的角色又如何呢?.
我們是在參與一個大的使命.
是讓弟子們的生命得到改變.
教會得到成長.
這才是我們整體的一個目標.
我們黃浸只有一個共同目標.
不是各自每個部門自己劃定自己的目標.
做好自己份內的事就算了.
不,共同參與,共同承擔一個目標.
但這個目標又不是我們能做到的.
是不是很矛盾?.
是的,布羅有這樣說.
你栽種了,你囂觀了.
但唯有上帝叫生命成長.
叫生命成長與你與我無關.
不是你與我做到的.
所以布羅接著說.
可見栽種算不得甚麼.
囂觀也算不得甚麼.
一切只在乎那個叫生命成長的上帝.
我不是叫我們小看我們在演的角色.
我們只是知道我們在做好自己的角色.
而我們看不見背後有一隻無形的手.
我們的主在混淆為壓.
並且靜悄悄地一路使生命得到改變.
你們未必知道.
剛才你們帶唱的那首詩.
某些歌詞在動一個弟兄,一個姐妹.
你不需要知道.
你需要刻意做一些戲劇性的表述去動它嗎?.
不需要.
聖靈在我們中間做默默的工作.
神蹟在我們中間發生.
而那個神蹟不是我們看得見的.
不是我們看得見的.
是的.

$^{561}$我們努力教導我們的主一學.
做一些純秀的工作.
我看到你們做兒童工作.
我就覺得很煩.
我最怕鬥小朋友.
不過你們做到都頗猛靜.
都頗有挫折感的時候.
可能你們不知道.
你們正在訓練一個將來的藤近輝.
將來的葛培莉.
沒有人知道.
我們只是做我們自己微小的一部分.
沒有人發現我們中間有神蹟.
是一點都不奇怪.
旺角山東街能夠出什麼好東西呢?.
這裡不扔油,不賣白粉.
還有什麼可以做呢?.
有時候我們都不覺得.
我們自己有什麼特別.
我們不是一間大型的教會.
我們又不是幾千人,上萬人的教會.
我們沒有豐厚的資源.
事實上你們不良禽擇木而棲.
去那些大教會獲得最好的服侍.
我都不敢說.
我都不明白.
不過我常常說的是.
香港有超過七成的教會是基層教會.
我沒有嚴格的統計.
超過七成以上的教會領袖.
如今的教會領袖.
是來自基層教會.
我教會的神學生.
多數是來自小教會,中小的教會.
大教會養著很多後座會友.
只有你們才需要在這裡落手落腳.
我們在這裡不經意地塑造了本土.
不經意地帶領人成長.
不經意地看到有人向上帝呼召毀身.
無論如何.

$^{601}$能夠施行神蹟的不是我們.
是在我們中間坐直的耶穌基督.
只有祂能夠將水變酒.
化腐朽為神奇.
我們的責任是邀請耶穌.
在每一個人身上施行神蹟.
聖經多次記載.
上帝因為人的迴轉.
聖經用後悔來對這個字去說.
就是上帝改變了祂本來要懲罰人的決定.
迦納的故事也告訴我們.
人的祈求可以改變耶穌基督的時間.
本來祂說祂的時間還沒有到.
但是祂最終提早施行了神蹟.
迦納故事裡面出現的眾多人物.
除了耶穌之外.
最突出的應該是瑪利亞.
這個真是一個非凡的女性.
雖然她負責將耶穌生出來.
將耶穌撫養成人.
但她在聖經裡佔的篇幅不多.
在約翰福音這裡是第一次出場.
因為約翰福音沒有記載耶穌出生的故事.
所以第一次說瑪利亞在約翰福音就是這裡.
除了耶穌出生的故事之外.
迦納婚宴是她極少數.
有一點點位置.
就算不是第一主角.
也好像是第一女配角.
就是女主角的角色.
全世界所有人裡面.
瑪利亞是最有機會看到耶穌平凡一面的.
看耶穌餵奶,換尿片.
沒有任何神秘感或神聖感.
如果拿丹葉可以說.
拿撒勒還能出什麼好東西.
如果連耶穌的兄弟在一段很長的時間都不相信他.
瑪利亞應該是最不相信耶穌是上帝.
英雄慣見逆常人.
才會在家鄉不受歡迎.

$^{641}$但瑪利亞卻不是這樣.
或者她從天使領報得知要懷孕生子的一刻開始.
她就知道這個兒子是極不尋常的.
他是上帝特別的挑選.
將來要去拯救整個以色列民族.
雖然耶穌在很長的一段時間.
30年的時間沒有行過一個神蹟.
有時候會說一些特別的話.
當失老那一次.
說過一些很噁心的話.
但你最多覺得他是五逆.
或者覺得他好像是天造英才.
但你不會覺得他是神人.
並且在約瑟死後.
我們相信耶穌是擔當照顧弟妹的角色.
甚至繼承父業,在做木匠.
瑪利亞一直認定這個兒子是不尋常的.
她沒有見過他行神蹟.
但相信他能夠行神蹟.
這是她將主人家的問題帶到耶穌跟前來的原因.
我常說瑪利亞是一個非常不聰明的人.
她沒有讀過書,沒有受過教育.
她對彌賽亞所知不多.
耶穌說過一些稀奇古怪的話.
她也一句都聽不明白.
聖經經常說.
耶穌說了一句話.
瑪利亞將耶穌傳在心裡反覆思想.
但沒有說她想到什麼.
她的轉數這麼低,B.
她想不到的,只是想而已.
負責想,但不負責想到.
不過,這個婦女是一個敬畏上帝.
最懂得的就是敬畏上帝.
在天使領部的故事裡.
她說了一句非常重要的話.
我是主的侍女.
情願照你的話.
成就在我身上.
而當她見到以利沙白.

$^{681}$更加說了一句.
抗古力金的名言.
我心尊主為大.
我寧以上帝我的救主為樂.
這兩句話.
充分顯示她用謙卑,信服,敬畏.
去對待上帝.
這是她的性格.
在加拿芬言的故事裡.
瑪利亞這兩個性格也充分彰顯出來.
不需要明白太多.
只需要謙卑和敬畏.
就不會做錯事.
一個謙卑和敬畏的人.
是永遠不會走錯路的.
我不需要明白.
主人的話,我怎會句句明白呢.
我明白得完,我就是主人.
但他叫我做,我就跟著做.
在教會裡.
最需要的就是像瑪利亞.
這樣謙卑和敬畏的人.
我相信我們往生.
有這樣的性格的人是不少的.
我們不需要懂得太多.
不一定站在台前指手畫腳.
但我們謙卑,信服,敬畏上帝.
默默關心身邊人的需要.
為別人的缺乏,可能有的危機去召急.
並且將人的需要帶到耶穌跟前.
我們不知道耶穌會做些什麼.
但我們認定耶穌會出手幫助.
所以我們只管聽從,吩咐.
無論耶穌叫什麼,我們都跟著做.
明的做,不明的也做.
瑪利亞是一個非常出色的代求者.
弟兄姊妹,你們知不知道.
我們中間有很多人渴望.
渴望將他們自己的水變成酒.
你們願不願意看到他們的生命得到改變.

$^{721}$是的,我們本來都是平凡人.
過著平凡的生活.
上班,下班,跟太太說些沒營養的話.
打麻將,每個星期用幾十元去買綠合彩.
買個夢想,養大子女.
希望退休後可以過一些無憂無慮的生活.
僅此而已.
這是香港人一般生活的軌跡.
多數香港人是期望這樣的生活.
我也都期望.
不過,我們信了耶穌之後.
我們願意離開這個慣常的軌跡.
我們願意經過考驗磨練被上帝改變.
我以前都跟你們說過.
我很認定我們每個人.
都是走著一個英雄之路.
由水變酒.
所有成為門徒之路.
都是一個英雄之旅.
一個信心英雄的故事.
我們教會很注重門徒分裂.
今年的主題是門徒不斷.
門婚就是自己生命得到改變.
人看到我們生命改變.
也都願意跟隨改變.
我們由水變酒.
也都吸引人由水變酒.
最重要的是.
我們確認耶穌基督的主權.
沒有兩個生命成長的故事是一樣的.
沒有任何生命成長的故事.
是奉旨的,是理所當然的.
都是耶穌的特殊恩典.
過去的成功不等於今天的成功.
我們只能憑信心繼續向耶穌呼求.
仰賴祂,繼續出手相助.
以前有自強的故事.
我們今日期待有更多不同名字的故事.
我們可以相信的是.
我們旺針所擺上的酒會是越來越好的.

$^{761}$過去六十年我們曾經擺上了不錯的酒.
未來我們能夠擺上的一定是更好更多的.
當然我們必須確認的是.
不管我們中間出現過多少神蹟.
就算我們旺針能夠發展到二千人三千人.
我們都不可以在此突出,吹捧我們自己.
而要藉著我們在此所見到的神蹟.
顯出耶穌基督的榮耀.
我們在此見到耶穌基督的榮耀.
\newpage



\section{詩篇使徒行傳 110:1-7}
\label{sec:KLXa_sv_cYI}
\textbf{2024年12月15日崇拜直播｜楊珊珊傳道｜詩篇裡的基督－永遠為祭司的君王｜詩篇110\_1-7;使徒行傳2\_29-40}
\newline
\newline
連結: \href{https://youtube.com/watch?v=KLXa_sv_cYI}{\texttt{https://youtube.com/watch?v=KLXa\_sv\_cYI}} ~~~~ 語音日期: 2024-12-15
\newline
\newline
\hyperref[sec:jH3ZJLIEHsY]{\small{< < < PREV SERMON < < <}}
~
\hyperref[sec:index_chronic]{\small{[返順時目]}}
~
\hyperref[sec:75p0crrJvGI]{\small{> > > NEXT SERMON > > >}}
\newline
\newline
$^{1}$弟兄姊妹,以下是經訓的時間.
今日所讀的經文記載在詩篇110篇1至7節.
另外,另一段是記載在《仕途行傳》第二章29至40節.
聽我讀出.
對我說:你坐在我的右邊,讓我使你受敵作你的腳凳.
耀華必使你從石安伸出能力的將來.
你要在你受敵中掌權.
當你掌權的日子,你的民要以聖潔的裝飾為衣.
以聖潔為裝飾,甘心犧牲自己.
你的民多如清晨的甘露.
耀華起了誓,決不後悔說.
你是照著墨基使得的等次,永遠為祭司.
在你右邊的主,當他發怒的日子,必打傷獵王.
他要在列邦中刑罰惡人,屍首就遍滿國柱.
他要在許多國中打破仇敵的頭.
他要喝路旁的河水,因此別抬起頭來.
弟兄們,先祖大衛的事,我可以明明的對你們說.
他死了,也葬埋了.
並且他的墳墓,直到今日,還在我們這裡.
大衛既是先知,又曉得神曾向他起誓.
要從他的後裔中,立一位坐在他的寶座上.
就預先看明這事,講論基督復活說.
他的靈魂不撇在陰間,他的肉身也不見扭歪.
這耶穌,神已經叫他復活了.
我們都為這事作見證.
他既被神的右手高舉,又重複受了那所應許的聖靈.
就把你們所看見,所聽見的囂觀下來.
大衛並沒有升到天上,但自己說.
主對我主說:.
「你坐在我的右邊,讓我使你受敵作你的腳凳」.
故此,以色列全家當確實的知道.
你們釘在十字架上的這位耶穌.
神已經立他為主,為基督了.
眾人聽見這話,覺得紮心.
就對彼得和其餘的使徒說.
弟兄們,我們當怎樣行?.
彼得說:你們各人要悔改.
奉耶穌基督的名受浸.
叫你們的罪得赦.
就必領受所賜的聖靈.

$^{41}$因為這應許是給你們和你們的兒女.
並一切在遠方的人.
就是主我們神所召來的.
彼得還用許多話作見證.
勸勉他們說:你們當救自己脫離這彎曲的世代.
聖經讀筆.
願神賜福常讀他話語並且遵行的人.
弟兄姊妹平安.
可能大家比較少見我.
我自我介紹,我是楊珊珊傳道.
大家親切一點可以叫我珊珊.
我加入旺鎮已經快四個月了.
通常我會在四樓的少年崇拜.
所以未能好好在大堂和弟兄姊妹交流.
希望將來有機會和大家深入認識.
每一個月的主餐崇拜.
我們這班中學生和大專生都會上來和大家一起敬拜.
希望大家平時在教會見到我們教會的少年人.
都主動多些和他們打聲招呼.
在講到之前,讓我們一起祈禱.
並願意差派你的愛子基督耶穌來擔當我們的罪.
釘在十字架上.
受死並且復活.
救贖我們每一個.
又賜下聖靈給我們.
求主開啟我們的耳朵.
緩化我們的心.
幫助孩子能夠清楚明白地講解.
願你的話語得以成就.
奉主名求.
阿們.
弟兄姊妹,你們有沒有試過和人傳福音或參加短宣?.
我見到有些人點頭.
剛好我們有一隊八人的短宣小隊.
在12月4至9日去了越南.
容讓我和大家分享我的見聞.
原來在越南進行福音的工作是非常困難.
因為那裡沒有像我們香港那樣有宗教的自由.
首先原來教會是很難向政府機構.
政府部門申請註冊成為一個合法的機構.

$^{81}$而且申請的過程往往需要十多年.
這實在對教會或是牧者來說太費事.
亦太虛耗人的耐性.
不過原來福音的工作往往都有聖靈的同工.
有位當地牧者跟我們說.
他非常感恩.
因為他的教會只用了五年時間成為一間政府認可的教會.
其次原來基督徒不可以在公開的場合傳福音.
否則會被執法部門截查或是被擋.
不過特別的是原來近年來越南希望提升自己的國際地位.
所以他逐漸開放給一些外國機構.
進行福音的盛宴.
有牧者跟我們分享.
雖然每次大型的福音盛宴有很多人信主.
不過在越南現在的基督徒人口分佈很少.
只有1\%.
所以當地的教會根本不夠人手跟進和栽培.
所以流失率很高.
另外我們還認識了一些同工.
他們非常有心智去山區和不同少數族裔的人傳福音.
但很可惜他們往往都是面對被捕的風險.
由於在共產制度的教育底下.
信基督教的人經常都覺得是不良分子.
甚至被人說是信羊教.
基督徒經常受到嘲笑或輕視.
試想一下.
在這樣的敵對環境之中.
你還要跟一開始對你有誤解的人傳福音.
是否很困難?.
試想一下.
你去見一份工作.
知道新的公司的同事很不濟.
對你很難回答.
上司又很刻薄.
對你很苛刻的話.
你會不會想去那裡工作呢?.
當然不會.
更何況是有基督徒跟這些人分享生命見證.
跟他們宣講神的道.
不過原來在越南有這班同工.

$^{121}$無論是上山下鄉還是在城市裡面.
他們面對人手不足的困境.
跟不同單位的周旋.
他們都很堅持,很忠心地.
在上帝交給他們的使命裡面.
他們很努力去完善.
或者我們會問.
究竟他們是怎樣可以做得到呢?.
又或者我們更加需要問的是.
我們所信的耶穌基督究竟有多強大.
竟然能夠讓人在艱難的情況下.
仍然選擇宣講他的事蹟.
或許這就是我們基督徒的特色.
早在二千年前.
初期教會的使徒.
在面對嚴重逼迫的時候.
仍然勇敢宣告.
因為他們深信耶穌基督就是他們生命的主.
四面受敵.
甚至面臨流亡.
福音的工作仍然不斷向外發展.
從耶路撒冷開展擴展到地極.
這就是使徒行傳的故事.
使徒行傳第二章記載了.
彼得在耶路撒冷領受了聖靈之後.
第一次的講道.
剛好今天也是我第一次在旺鎮的大堂講道.
不過當時的會眾應該沒有大家那麼友善.
我們不妨進入一下經文.
如果你有聖經也可以翻開來看.
使徒行傳二章四節提到.
當聖靈在五長節降臨在這群使徒身上.
這群使徒就按照聖靈所賦予他們的語言能力.
他們能夠用不同的外語.
與來自五湖四海的人聊天.
並且還能夠有足夠的能力.
與人分享神的偉大作為.
如果我們繼續看第九至十一節.
我們會看到有十多種民族.
這意味著他們突如其來可以說出十多種不同的語言.

$^{161}$弟兄姊妹厲害嗎?.
試想像我們在六樓的禮堂上.
有十多種不同的語言.
簡直就像聯合國一樣.
但不要以為大家聽了就立即相信.
他們的第一個反應是怎樣?.
他們覺得很驚奇,很疑惑.
互相在問:什麼鬼?搞什麼?.
What's going on?.
有些人甚至嘲笑他們.
說他們喝醉酒,是發酒癲.
於是彼得在講道開始前回應他們說.
其實現在才九時,怎會喝醉?.
不過有一件事你們應該很明白.
亦要留心聽我所說的話.
於是彼得引用了幾段舊約經文證明.
「施行大能,歧視和神跡的拿撒勒人耶穌.
正是以色列民一直期待的救主彌賽亞」.
首先彼得先引用「若以先知」的語言.
來講述聖靈的囂觀.
亦將「若以書」二章中的「以後」改為「在末後的日子」.
強調在上帝救贖計劃的歷史中.
已經進入了一個新的時代.
這是聖靈帶領下教會的時代.
正因為聖靈的降臨.
不同的使徒就說了不同的語言.
這個情況好像應驗了「若以書」二章所說的.
「以後我要將我的靈囂觀凡有血氣」.
「你們的兒女要說預言」.
「你們的年老人要做義夢」.
「少年人要見異象」.
「在這些日子我要將我的靈囂觀我的僕人和侍女」.
我記得三年前的暑假.
在疫情期間,當我仍然在讀神學的時候.
我的母院,見道神學院就做了一件很瘋狂的事.
我們在長洲的校園舉辦了一個有三百多人的大專四日三夜營.
當時我仍然是一年級生.
我負責擔任營會的助理和小組組長.
其中有兩個片段令我印象很深刻.
第一個片段就是.

$^{201}$我記得有一天,我在幕後拿著帆盒.
同一時間收拾物資的時候.
有一個白色短頭髮的女生,也是大專生.
很尷尬地走過來問我一個問題.
她說:「珊珊,為什麼你讀神學?」.
於是我就一邊塞著飯進她的口.
一邊回答她.
其實當時我只是想快點回答完她就完結了.
於是她聽完之後就說了一句「哦」.
然後就走了.
之後我沒有再特別留意這個女生.
因為她不是小組的組員.
第二個片段令我很深刻.
就是在營會最後一晚.
我的神學院老師竟然很勇敢地.
向這班大專生發出呼召.
問他們有沒有人願意踏上全時間事奉的路.
成為傳道人或宣教士.
其實當時我聽到.
我也覺得應該沒有很多人走出去.
因為他們都在讀學位.
還未讀完子學位又怎會想其他東西呢?.
不過很感恩.
當晚有超過30位大專生.
走上台回應神的呼召.
令我很驚訝的是.
那位白色短髮的女生.
也走出去回應神的呼召.
之後老師又發出了第二個邀請.
她說:「如果你有見到朋友或弟兄姊妹.
或認識的走出去回應呼召.
而你也希望支持他們.
願意跟他們同行.
不如你也走出來台前祈禱吧!」.
那一刻我感到很糾結.
因為我覺得神也在邀請我.
陪伴這班年輕人.
為神的國努力.
當我終於願意走出去台前的時候.
我發覺原來已經有百多人.

$^{241}$衝出台前為願意為神擺相的年輕人祈禱.
我發覺原來我連那位女生的衣服尾巴也碰不到.
我覺得那一刻自己實在太低估神的能力.
原來即使在這個世代如此艱難.
有些環境你覺得很惡劣.
神從來沒有離開過香港.
也從來沒有離開你和我.
祂還呼召了一班比我和你更年輕的人.
開闊他們的眼界.
讓他們在一個黑暗的世代裡.
看見神給他們的遺像.
弟兄姊妹.
聖靈的工作從未停止.
有心的人其實還有很多.
所以我們不是少數.
不知道你今天願不願意相信並且看見聖靈.
同樣降臨在我們的皇寢當中呢?.
使徒行傳的故事並未停在聖靈降臨的那一刻.
在耶穌復活之後.
使徒行傳第一章第六節.
這班使徒就向復活了的耶穌問了一個問題.
祂說:主啊!你是否在這個時候要復興以色列國?.
換句話來說.
其實在他們眼中.
耶穌只是那位帶領猶太人脫離羅馬帝國政權迫害的王.
所以彼得這篇講到.
其實是要更正他們對彌賽亞君王的理解.
彼得之後其實引用了以色列民一個最偉大的君王大衛的詩篇.
來教導這班弟兄.
究竟耶穌基督是一位怎樣的君王呢?.
於是彼得就用了詩篇第十六篇和第110篇來講解.
根據以色列的傳統.
詩篇第110篇被視為皇室的詩篇.
所以它比第十六篇更能夠凸顯耶穌基督作為君王的形象.
所以我們今天會集中看110篇.
這首詩通常會在國王登基時獻唱.
我小時候也有看TVB的宮廷劇.
通常我會想像皇帝登基.
鏡頭一轉就是紫禁城.
朝廷的文武官員和妃嬪都出來迎接皇帝.

$^{281}$背景音樂都是很雄壯激昂.
我們想像就是這個感覺.
彼得就引用詩篇第110篇.
向那些很熟悉舊約的猶太人.
講述一個新的君王開始.
他開始接受上帝賦予他的皇權.
這是標誌著一個新的皇朝和新時代的誕生.
這首詩其實講述了耶和華和大衛的主之間的對話.
即是天父和作為彌賽亞的兒子耶穌基督之間的對話.
如果將這首詩分為兩個部分.
前面頭三節主要描述了耶穌作為偉大君王的形象.
而後面四至七節指出耶穌是我們永恆的大祭司.
我們首先一起看看第一個部分.
在前面的頭三節.
天父向基督賦予了三個應許.
第一節:你的仇敵會被打敗.
第二節:使錫安伸出能力的權杖.
即象徵著基督的國度將會擴張.
第三節:你的人民將樂意投身於戰鬥.
即天父會賜一個軍隊給基督.
因為耶穌願意被釘十字架受死.
隨後復活而升天.
所以祂現在已經坐在父神的右邊.
右邊象徵著榮譽,威嚴和權力.
這位祭司君王的形象.
代表著耶穌擁有完全屬天的尊榮和力量.
就像菲勒比書2:9-11所說.
因此神高舉了他.
他指的是耶穌基督.
賜祂的超乎萬名之上的名.
使天上,地上,地底下一切.
全部都屈膝敬拜在耶穌的名字面前.
全都口裡宣認耶穌基督是主.
將榮耀歸給父神.
在詩篇110篇第一節.
有一個很特別的字就是「腳凳」.
「腳凳」其實象徵著一個完全的勝利.
試想想.
你將你的腳放在敵人的頸上.
敵人怎會不死呢?.

$^{321}$早幾年在美國有一個新聞.
有一個執法人員.
他的腳跪在黑人的頸上.
其實一定會死的.
在這裡我們看到有兩個視角.
在耶和華的眼中.
耶穌為世人所做的事情.
使祂能夠勝過罪惡的權勢.
成為君王.
但是在世人的眼中.
祂是嘲笑和貶低基督.
值得注意的是.
其實剛才在《使徒行傳》2:36.
我們都聽到.
杜經元幫我們讀出.
彼得沒有說羅馬兵丁或外邦人.
那些人是釘死耶穌的人.
而是面對著這群熟悉耶和華律法的人說.
是你們釘死耶穌基督.
弟兄姊妹.
今天我們是否仍然在罪中作樂.
又或者在跟隨基督成聖的路上.
停滯不前呢?.
耶穌基督已經為我們釘在十字架上.
讓我們不要再留戀傷害自己.
和他人的關係或事物.
站在耶穌那邊.
總比做祂的敵人更好.
因為祂是我們祭司君王.
祂是那位德性的君王.
並且祂也應許賜聖靈給我們.
我們求聖靈幫助我們.
如果我們走錯了路.
不要猶豫.
轉身走回城裡便可.
這就是悔改.
讓我們繼續看詩篇第二部分.
因為詩篇的重點是放在第四節.
在聖經裡.
除了耶穌基督.

$^{361}$只有一位同時擁有君王和祭司身份的人.
那就是默基洗德.
默基洗德這個名字.
在聖經裡有三個出處.
第一個出處是在《創世記》第十四章.
提到他是一個撒冷的王.
撒冷事實上是耶路撒冷的地方.
他也是那裡至高神的祭司.
當時他還為未改名的亞伯拉罕.
即是亞伯蘭去祝福.
第二個出處是詩篇110篇.
第三個出處最多的是在新約的希伯來書.
七章第二節裡.
默基洗德這個名字.
解釋為公義的王.
撒冷的王意思是和平的王.
平安的主.
平安的王的意思.
他無父無母無家族無出生日子.
無生命的終結.
聽起來和神的耶穌有點相似.
他永遠作祭司.
沒有人可以繼承他的地位.
這個祭司的等級.
是不同利尾支派亞倫祭司的等級.
因為要成為默基洗德這個高等的祭司.
是不可以根據肉體上的條例.
即是他不是指他是哪個支派所出.
而是根據他有一個不扭生命的能力.
基督作為祭司君王.
就像默基洗德一樣.
他不單是猶太人的祭司.
因為他是至高神的祭司.
所以他不受地上任何一個國家或民族的限制.
他不特別屬於某一個群體.
所以耶穌基督的拯救和祝福.
是會臨到各個不同國家的人身上.
彼得之所以引用詩篇110篇.
就是希望我們可以認定.
耶穌基督就是那位偉大的君王.

$^{401}$和永遠至高神的大祭司.
這首詩的開頭是這樣表達的.
「主對我主說」.
值得注意的是.
第一個主指的是耶和華.
而第二個主就是我主的意思.
聽起來好像很接近.
其實讀音有少許不同.
他們都是同等級別.
不過是指兩位不同的主.
耶和華指的是.
以色列文的獨一真神.
而我主指的是耶穌基督.
在猶太人的傳統裡.
主和我主的稱呼只限於耶和華.
但彼得和使徒.
用我主來稱呼了復活的耶穌基督.
這代表著耶穌和神有同等的地位.
正正因為復活的基督.
是永活的祭司君王.
祂成為了我們回到神的出口.
丙姐妹.
耶穌對你來說究竟是誰呢?.
你的信仰是第一手信仰還是第二手信仰?.
是因為其他人信而你相信?.
還是你每時每刻都堅定.
認定耶穌基督就是你個人的救主呢?.
我其實是在準備入大學的時候.
回教會的.
當時我記得我參加的青年團契.
是在高峰期.
所以我們有30至40人.
但經過十多年後.
有些人因為拍拖,工作,生活忙碌.
人際衝突,家庭移民或有牧者的變遷.
各種原因.
漸漸地大家都慢慢脫離了團契.
或是大家都覺得疲憊,忙碌,懶惰.
對信仰變得懶賴.
團契的聚會次數由每周一次.

$^{441}$變成每月一次.
然後每季一次.
其實當我們稱耶穌為「我主」的時候.
我們不可以口頭上說「是的,祂是」.
「Yes」.
但我們的內心和行為卻說「不」.
我們對耶穌說「No」.
耶穌曾經說過.
「不是每一個對我說主呀主呀的人,我都認識」.
「能夠進入天國的,唯有進行我天父旨意的人,才可以進去」.
如果今天是你第一次回來教會.
又或者因為朋友邀請你先來.
我鼓勵你多認識耶穌基督.
因為祂真的很堅.
不信的話,你可以問問你的朋友.
我希望弟兄姊妹.
我們要記得.
我們所信的耶穌是永遠的祭司君王.
祂是帶給我們公義和平安的祭司.
亦都是擁有尊榮和能力的君王.
如果你在信仰的路上遇到困難.
願你記得.
世上沒有困難比基督更大.
在物質豐富的香港.
其實我們有時候都很難願意放下自己做主.
我們有時候可能是前途做了我們的主.
戀愛做了我們的主.
父母做了我們的主.
孩子做了我們的主.
但我們又願不願意選擇.
看見我們一切所擁有的.
都是我們主耶穌基督所賜予的呢?.
耶穌基督才是我們的主.
是我們所關注,關心的人和事物的主宰.
我明白有時候跟隨神是需要付出代價.
我真的絕對明白.
坦白說,我也試過感到疲倦,灰心.
我甚至曾經離開教會一段日子.
不過,當我清楚明白我所信的是誰.
耶穌不單只是一個歷史的偉人.

$^{481}$更不是一位有求必應的王大仙.
他是至高神的兒子.
永遠的祭司君王.
然而,這位永遠的祭司君王.
卻願意為了你和我.
書專降矣.
有一首歌叫《Love was when》.
中文譯作《愛就是那時》.
其中有一句歌詞我很喜歡.
它是這樣寫的.
「愛就是神化身為人之時.
受困於時空之中.
沒有名,沒有地位」.
「愛就是神化身為人之時.
受困於時空之中.
沒有名,沒有地位」.
耶穌曾經嘗過人間的疾苦.
他明白我們作為人的限制.
即使我們人釘他上十字架上.
他仍然為我們代求.
他說「負我赦免他們一切所做的.
因為他們所做的他們不知道」.
當彼得提到這位耶穌就是基督的時候.
那些會眾非常擇心.
原文擇心的意思就是被人捅了一刀.
又或者你的良心好像受到傷害.
如果用廣東話翻譯.
我會選擇良心發現.
心裡覺得很不舒服.
因為這班人找到生命的主.
他們願意向彼得和其他使徒詢問.
「弟兄,我們應該怎樣做?」.
這班會眾是從曾經釘耶穌十字架的人.
轉變為願意放下手上的錘.
跟隨耶穌基督.
最終加入使徒的行列.
傳頌耶穌的故事.
經文指出他們有四個反應.
第一,悔改.
第二,奉耶穌基督的名受洗.

$^{521}$第三,聚得赦免.
第四,領受聖靈.
頭兩個回應就是人願意宣揚耶穌是主的表現.
彼得挑戰他們的生活一定要有改變.
悔改這個詞是路加非常喜歡用的動詞.
它不單是說頭腦上的改變.
因為當時希羅的《智學》是說頭腦要怎樣變.
要怎樣認識智慧.
它不是這麼簡單.
它是說心意上的改變.
更加是說放棄一切罪惡的行為.
聖經說這種改變會讓我們回到上帝.
一旦放下罪惡.
我們就不要向後望.
當我們眼睛還要回望,眷戀那些罪惡的時候.
我們要記得用力看著耶穌.
當我們願意相信耶穌的時候.
神賜予我們之後的兩個應許.
當我們願意悔改,轉向神的時候.
我們的罪惡就會得赦免.
並且可以領受聖靈.
這裡提到的應許不單只是給猶太人或他們的兒女.
神的應許也是賜給那些在遠方的人.
當我們繼續讀《仕途行傳》的時候.
就會看到很多跟外邦人傳福音的事蹟.
有些人這樣說.
仕途行傳有多少章.
對了,我看到有些口型.
28章.
其實我們是《仕途行傳》的第29章.
29章怎樣寫呢.
其實很看我們怎樣跟主耶穌相遇.
耶穌應許我們.
他說天上地下一切的權柄都賜給了他.
所以我們要去使萬民作他的門徒.
奉父子聖靈的名幫他們施浸.
並且教導他們遵守他吩咐我們一切的.
最重要的是耶穌會經常跟我們在一起.
因為他答應賜下聖靈.
聖靈會在我們軟弱的時候.

$^{561}$幫助我們直到世界的終結.
弟兄姊妹,願今天的經文.
能夠幫助我們更加認識.
耶穌基督這位永活祭司的君王.
所以即使無論我們現在.
或將來所面對的困難.
怎樣也好.
要記得耶穌比困難更加大.
不要害怕.
灰心的時候不要喪志.
聖靈會一直與我們同在.
願主賜福給大家.
\newpage



\section{路加福音彌迦書 5:2-5}
\label{sec:75p0crrJvGI}
\textbf{2024年12月22日聖誕節主日暨嬰孩奉獻崇拜直播｜陳明泉牧師｜萬世之君誕於小城｜彌迦書5\_2-5;路加福音2\_1-20}
\newline
\newline
連結: \href{https://youtube.com/watch?v=75p0crrJvGI}{\texttt{https://youtube.com/watch?v=75p0crrJvGI}} ~~~~ 語音日期: 2024-12-22
\newline
\newline
\hyperref[sec:KLXa_sv_cYI]{\small{< < < PREV SERMON < < <}}
~
\hyperref[sec:index_chronic]{\small{[返順時目]}}
~
\hyperref[sec:NqwhnMUzemI]{\small{> > > NEXT SERMON > > >}}
\newline
\newline
$^{1}$《聖經》第五章二至五節上及二十二節中的第一段:《尼加書》第五章二至五節上及二十二節中的第二段:《馬太福音》二章一至六節.
《尼加書》第五章第二節.
「伯利恆,爾弗他,你在猶大諸城中為少,將來必有一位從你那裡出來,在以色列中為我作掌權的.他的根源從耿古,從太初就有,耶和華必將以色列人交付敵人,直等那新產的婦人生下子來..
乃是掌權者,其餘的弟兄必歸到以色列人那裡.他必起來,以靠耶和華的大能,並耶和華他臣之名的威嚴,募養他的羊群.他們要安然居住,因為他必日見尊大,直到地極.這位必作我們的平安.」.
第二段經文是在《馬太福音》二章一至六節,在新約第二頁.
請聽我讀出第二章一至六節.
「當希律王的時候,耶穌身在猶太的伯利恆,有幾個博士從東方來到耶路撒冷,說:那生下來作猶太人之王的在那裡.我們在東方看見他的聖,特來拜他..
希律王聽見了,就心裡不安.耶路撒冷合城的人也都不安.他就召齊了祭司長和民間的民事,問他們說:基督當身在何處?他們回答說:在猶太的伯利恆,因為有先知記著說:猶大地的伯利恆啊,你在猶大諸城中並不是最少的,.
因為將來有一位君王要從你那裡出來,莫讓我意識滅盟.」.
各位弟妹平安!.
我們今天帶著歡喜,快樂,感恩的心來到教會當中,一起慶賀那位為我們降生的救主耶穌基督..
我們一起來用歌聲,用很美的佈置,用很熱烈的氣氛,來到這個很熱鬧的地方裡面,我們要唱頌這位拯救我們的主..
不過不知道大家心裡面有沒有問一個問題,為什麼這位為我們降生的孩子耶穌,在他沒有做任何事情的時候,全天下甚至歷代歷代,到今天,或者在12月這個日子裡面,.
大家都會用不同的語言,用不同的方式來慶賀他的出生..
這個孩子其實原則上在他出生的時候,沒有任何事情是為我們做得到出來的,只不過是因為上帝藉著這個兒子展開了一個救贖的計劃,.
但是他的出生,在當時的社會裡面,都沒有做任何的事情,都沒有任何的果效可以看到,但是所有人都為著他的出生來慶賀..
然至到今天,我們每一個基督徒,或者每一個主內大英姐妹,或者朋友,我們都用不同的方式,一起來慶祝..
但是為什麼這個嬰孩為我們帶來這麼重要的改變,我們要慶祝呢?.
今天我們看的這一段經文,某程度上都在說一些很重要的,一些的應許實現..
我相信,單單耶穌基督的降生,都顛覆了整個世界,或者我們人類歷史裡面,很多我們一直保持的價值觀念..
當這個世界覺得人性有很多困難的時候,耶穌基督的降生,就好像扭轉了這個世界很多對人與人之間的看法,.
又或者國與國之間產生很多的嚴規,甚至打仗,耶穌基督的降生,好像顛覆了扭轉了這個世界不同國家的命脈..
有一個心理學家更加說,他說人裡面心靈的平安,因著有一個救主,因著有福音來到,.
人對自己的看法,人對身邊的人的看法,甚至人對世界的看法,就自然改變了..
希望我們今天從這一段經文,由耶穌基督出生,到耶穌基督出生之前的七百多年的一段記載,.
我們看看,究竟這個應許在一個怎樣的狀況下實現出來的..
到七百多年之後,這個應許實現了,那些人怎樣去回應,好不好?.
我們用這段經文希望成為我們生命裡面一個很重要的亮光,我們先祈禱,求上帝幫助..
天父上帝願你在我們當中帶領我們,讓我們能夠明白你的話語,.
在你歷久不衰的真理當中,讓我們拿到生命的養分..
願主你此刻在我們當中,賜福給每一位聆聽的弟兄姊妹,讓她們不單聽得明白,.
也將你的話語藏寄在心中,實踐在生活裡面,又讓孫講的講得清楚,.
將你的話語中心的樣子向弟兄姊妹去展明,我們恭敬將來的時間交託給你,願你賜福..
獻上禱告奉求基督耶穌得勝的聖名,Amen..
我們先看看,今天我們主要看兩段經文,第一段就是舊約的尼加書,.
尼加書這一段的經文,他有一些背景的,尼加書在一個什麼狀況之下寫呢?.
大概成書的日子,或者尼加先知是在耶穌基督出生之前的七百多年,.
大概主前750至710年的這段時間,他出生..
所以尼加先知是不知道上帝未來具體的計劃,他更加不知道,他看不到耶穌基督是誰,.
在耶穌到誠育身之前,他就講出整個書卷的內容,.

$^{41}$而尼加先知身處的是一個什麼時代背景呢?.
在他身處的年代裡面,其實整個以色列民族已經分了南與北兩個國家,一分為二..
類似一個國家的興衰,有時候都會因著權力,或者因著人的那種慾念,.
甚至乎在以色列民族裡面,因著他們敬拜的神明的觀念不同,.
以致他們分開了南北兩國..
在當時來說,在大概公元前735年,有一個外邦叫做亞述帝國..
這個亞述帝國其實在歷史裡面,我們在不同的典注裡面都會看到的,.
是一個軍事國家,不斷的吞併,由敘利亞到不同的地方裡面,.
包括以色列的北國,簡稱叫做以色列國,當時來說已經南北分開了,.
南國叫做猶大,北國叫做以色列..
在亞述不斷擴展的時候,他已經開始侵佔,甚至乎已經佔領了當時的北國以色列..
可想而知,其實整個以色列民族都面對著一個外邦壓境的局面..
過往的日子裡面,以色列這個民族是挺有錢的,他們都生於安逸當中,.
以致他們在安逸的環境裡面,他們就忘記了上帝,又或者他們失去了以前他們先前的那種警醒,.
要聽上帝的勸勉,要怎樣來治理這個國家..
因為他們有錢了,他們興奮了,甚至乎追隨了一些外邦的神明,.
以致他們的國家一直由富有逐漸走下坡,直至到外邦來壓境,.
以致北國就開始覆亡..
所以在這樣的背景下,尼加書五章一節裡面,他特別的來講,.
他說成群的民,現在你要聚集成隊,因為仇敵圍攻我們,要用仗擊打以色列審判者的臉..
在這裡裡面,其實他在講,其實整個第四章上面都一直在講,.
就是外邦已經壓境了,現在那些神國的子民,整群的女子,.
他在講整個民族的人,他們要興起來,集結成軍隊,.
以前他們是賺錢的,生活安逸的,現在他們要警醒了,要開始來了,.
因為外邦的仇敵已經圍攻著整個民族,不單止是這樣,.
他不單止是打仗,他更加是羞辱以色列的審判者,.
審判者這個字,其實舊約裡面講的士司記,那個士司是同一個字,.
就是在講以色列那些領袖被外邦人侮辱,.
其實是在講一個國家,面對著一種很大困難的處境之下,.
寫成的這個尼加書,更加重要的就是,去到第五章第二節開始,.
其實去到一路五章一路打後,全部都是在講一些未來的預言,.
面對著國家犯境,又面對著他們自己國家裡面那些人,.
可能就不公義的實踐信仰,或者離棄上帝,.
在第五章二到第五節,特別我們今天要看的那段經文,.
他就作出在未來裡面的一些上帝對直著這個群體所作的一些預言,.
裡面所講的預言其實你一路看下去,你用心看你會覺得非常奇怪,.
我們再看第五章第二節,.
在第五章二節裡面就這樣講,.
伯利恆的爾法他,你在猶大諸城中為少,.
將來必有一位從你那裡出來,在以色列中為我作掌權的,.

$^{81}$他的根源從耿古到有,從耿古從太初就有,.
我們先看第二節在這裡頭半部分,.
他講到伯利恆的爾法他,在猶大諸城中為少,.
將來必有一位從你那裡出來,在以色列中為我作掌權,.
在這裡一定要有一些補充,大家才明白那句話的奇怪地方,.
首先在當時整個以色列民族,如果想著有一個君王的出生,.
我們通常想都會在那些大城市或者出生很好的,.
或者出生很好的人一定在那些大城市,.
在大家焦點底下的城市裡面就會興起一個孩子,.
被高度栽培,以至成為明日的君王,.
一般人都是這樣想的,.
不過在這段經文裡面講到未來的君王出現的地方就是伯利恆這個城市,.
伯利恆有多大呢?其實他的人口等於旺角兩棟大廈,.
一千多人的人口,為什麼這樣說呢?.
他說那個字諸城,諸城以為很多的城市,.
其實是在說小城市,一千人人口裡面的村落仔,.
在那些村落仔一千人打下的,他還是叫小的城市,.
可想而知他的人口其實是一個微不足道的,.
不是舉足輕重的,當時影響經濟或者很多生活上面,.
或者政治裡面一個很重要角色的一個城市,.
但是在這個經文裡面他就很直率的在講伯利恆將來會出一個君王,.
不過伯利恆這個地方真的是福地,.
為什麼這樣說呢?因為伯利恆本身這個字,.
他的名字伯利恆其實在希伯來文的意思就是等於麵包屋的意思,.
就是說他的物產是豐裕的,有很多糧食在當中的,.
所以雖然這個城市小小的,但是他們會有多少的出產,.
更加重要的就是在整個以色列民族的歷史裡面,.
有一個很重要的偉人就在那裡出生了,.
在撒慕爾記上十七章十二節,那裡就記載著整個以色列民族,.
第二代的君王大衛就在那裡出現了,.
大衛是猶太地伯利恆的爾法他耶西的兒子,.
耶西有八個兒子,在這個記載裡面我們知道了,.
整個以色列民族在講大衛王朝,大衛的掌權,.
整個以色列當下裡面的民族的君王,.
就是先祖先王,他都在那裡出生的,.
不過雖然這個地方出了第一代君王,.
但是他的兒子,大衛小小的代代,.
他們已經搬進耶路撒冷,已經住在大城那裡,.
他不會回到鄉下,再回到自己的根源地裡面,.
又再出現一個新型的君王出來的,.

$^{121}$所以在這裡面其實很顛覆,.
當時人們想的,會有一個新的大衛,.
回到以前的鄉下,現在還是窮乏,.
不富庶的一個小城市,.
來出一個新的君王,.
另外重要的就是,其實他要回顧伯利恆,.
因為他講大衛,大衛的出生,.
大衛是一個小的城市裡面,.
還以為耶西有八個兒子,.
他是個小子,.
所以他是小城市裡面一個最微不足道的小子,.
尼加書他要講到,.
在將來重新要作王的,.
就要找一個小子,.
找一個將所有肩負起和平的責任,.
就會放在一個眼前裡面微不足道的一個孩子,.
來將所有上帝的應許實現在他的身上,.
所以在這裡其實就顛覆了第一個我們的想法,.
在這個應許裡面講到,將來的君王,.
不是由一個很顯赫的王宮來出現,.
而是一個窮鄉僻壤,.
我們再繼續看,.
在經文裡面再繼續講的,.
這裡就更加令我們不容易理解了,.
因為他的時序其實都頗令你頗混亂的,.
為什麼這樣講呢?.
經文裡再繼續講,.
在第二節下半部,.
將來必有一位從你那裡出來,.
在以色列中作掌權的,.
他的根源從耿古從太初就有,.
你留意一下他的時序,.
將來就是未來式,.
但是在講那個作掌權的年輕人,.
這個將來式,.
但是那個年輕人是什麼人呢?.
他的根源從耿古從太初就有,.
就是說是一個創世未有的,.
一個已經存在的掌權者,.
那時候對我們人類來說,.

$^{161}$將來有一個從耿古已經有的人來掌管,.
是什麼概念來的?.
我們以為將來來的就是一個年輕人,.
他也會長大,.
然後成為掌權者,.
是尼加的後人,.
幾百年之後的一個很重要的人被興起,.
不過他講到這個人不是後者,.
而是前者,.
耿古和太初用的字就是創世以前的,.
未有世界,未有人類歷史的時候,.
那個人已經存在了,.
所以你會覺得很奇怪,.
整個時序將來有一個舊的掌權者,.
來帶領整個以色列,.
然後去到第三節,.
然後續節的,.
每一句都是令人難以明白和費解的,.
第三節,.
耶和華必將以色列人交付敵人,.
直等到那生產的婦人生下來,.
那時掌權者其餘的弟兄,.
必歸到以色列人那裡..
講第三節這一句的經文,.
更加對於當時聽的人來說,.
一句不中聽的話,.
記不記得剛才尼加在講,.
面對著國家有危難的時候,.
外族入侵的時間,.
來到寫成的,.
如果我是一個振奮人心的先知,.
我就會說,.
這個人興起來,.
然後就幫我們打仗,.
然後我們就所向披靡,.
就戰無不勝,.
我們以色列就光復了等等,.
類似一個民族精神的一個信息,.
不過尼加就不是這樣說,.
尼加是在說,.

$^{201}$上帝會將以色列人,.
其實整個民族交付給敵人,.
簡單來說,.
整個民族是要亡國,.
是要敵人外邦,.
基本上是要進佔他們,.
他們要被擄,.
被外邦地,.
其實在第四章裡面,.
都有出現巴比倫的,.
他還沒講到這麼遠,.
近來的亞述,.
對於以色列北國,.
都一定已經被擄,.
甚至他們已經戰敗給敵人手上,.
所以這個不是什麼,.
那個應許你覺得,.
其實是說一些很令人整個民族洩氣的話,.
不過,.
他在說這個不是以色列人的結局,.
到那個生產的婦人生下子來,.
其實是婦人生兒子,.
不過,.
如果我們不明白的,.
他在說一個嬰孩將要降生的日子,.
那個掌權者其餘的弟兄,.
被擄去到外邦地的那些人,.
都要重新歸到以色列人那裡,.
所以整件事是在說,.
他這一句話,.
似乎不是很中庭,.
但是在說一個更遠的將來,.
就是那些四散的那些以色列人,.
或者亡國的那些,.
離開了本土的那些弟兄,.
有屈天,.
當這個孩子降生的時候,.
他們在外邦所有人,.
又會匯聚回到以色列人當中,.
這個是第二個他們難以理解,.

$^{241}$甚至聽上去覺得是一個很難的應許,.
接著第四節,.
他必起來,.
依靠耶和華的大能,.
並耶和華他神之名的威嚴,.
牧養他的羊群,.
他們要安然居住,.
因為他必日見尊大,.
直到地極,.
這位必作我們的平安,.
整個句式第四節那裡,.
到最後講到的這個掌權者,.
一個生出來的明日的君王,.
他會依靠上帝的大能,.
和上帝的威嚴,.
牧養他的羊群,.
我們想著牧養他的羊群,.
想著這個君王就是很慈祥的牧者,.
用一個很溫柔的方式去帶領,.
這個不是錯的,.
不過我們不要純粹對號入座,.
一想就是牧養羊群,.
就純粹去想,.
當時的人就一定覺得是一個溫文爾雅的牧師,.
來帶領他整個民族,.
他這裡所講的牧養那個字,.
其實是一個政治字眼,.
因為當時那些國邦,.
那些國家叫自己的皇帝,.
都是叫自己的皇帝做牧者,.
所以那個牧者其實是一個帶領的領袖,.
而牧養羊群其實是在講管治他們整個的民族,.
這個是當時裡面最普遍的用法,.
他們就想這個新星,.
這個新的君王興起來,.
會做他們的軍事或者國家元首領袖,.
會令到所有四散的以色列人,.
都要重新歸在他的民族當中,.
跟著他要令到整個以色列民安然居住,.
即是在講他們有好日子過了,.

$^{281}$而且這個君王會日見尊大直到地極,.
他的影響力偏及萬邦,.
跟著第五節那裡說,.
這位必作我們的平安,.
同樣那個平安的字和牧養的字是相呼應的,.
那個平安的字不是在講很舒服,.
內心裡面的那種體會,.
那個平安是一個政治用語,.
是在講四境不再打仗,.
人在這個地方裡面可以安居樂業,.
而且為什麼不用打仗呢?.
因為他自己夠強力,.
這個君王帶領整個以色列民族,.
那些人面對著一個這麼強盛的國家,.
沒有人夠膽打他,.
所以沒有人夠膽犯警,.
這個平安就是來自一種政治,.
當時的人這樣想,.
來作為一個整個理解這個的應許,.
梅加書講的這個應許,.
或者在這裡裡面,.
他對於一個外族壓境的那種狀況,.
其實他作出一句這樣的一段上帝的工作,.
某程度要針對當時那些外邦人壓境,.
以至他在上帝的領袖當中,.
他希望,或者不是自己他,.
而是整個民族都希望有一個真正強盛的軍事領袖,.
一個真正的彌賽亞,.
來到政治裡面,.
令到他們一切都可以享受到平安,.
這個是在耶穌降生前七百多年,.
然後一直後來他們都是這樣想的一種觀念,.
而在以色列歷史裡面都曾經出現過,.
類似有一些軍事強人,.
馬加比,他們希望用一個軍事的形式,.
來為整個民族帶來平安,.
甚至興起一個新的君王,.
帶領整個以色列民族,.
不過上帝的計劃和人的想像和人的想法,.
甚至是在講尼加書,.

$^{321}$講了這個應許之後,.
當時那些人怎麼看呢?.
是完全有一個很大的距離,.
是去到新約馬太福音第二章一至六節,.
他就講到這個,.
或者整個福音書,.
講到這個君王降生,.
但這個君王的降生,.
在七百多年之後,.
帶來的整個改變,.
已經不是和他們之前最初那種想法是一致,.
我們再看馬太福音第二章,.
第一節至到第六節,.
當希律王的時候,.
耶穌生在猶太的伯利恆,.
有幾個博士從東方來到耶路撒冷說,.
那生下來作猶太之王的在哪裡呢?.
我們在東方看見他的聖,.
特來拜他,.
希律王聽見了,.
就心裡不安,.
耶路撒冷合成的人也都不安,.
他就召齊了祭司長和民間的民事,.
問他們說,.
基督當生在何處?.
他們回答說,.
在猶太的伯利恆,.
因為有先知記著說,.
猶大地的伯利恆,.
你在猶大諸城中並不是最少的,.
因為將來有一位君王,.
要從你那裡出來,.
目讓我以色列民..
這裡是截錄了一部分,.
在這裡說到耶穌,.
按著上帝的應許,.
七百多年之後,.
在一個不起眼的麵包城,.
經過七百多年,.
雖然有這個應許,.

$^{361}$但是以色列人又亡國,.
又生活形形異異,.
某程度已經忘記了,.
又或者不想到是這樣實現,.
七百多年之後,.
在伯利恆這個小城裡面,.
再一次做了一個動作,.
以至到耶穌出生在那裡,.
耶穌自己本身,.
他的家人,.
他的媽媽,.
其實是拿撒勒人,.
更加名不經傳的一個小地方,.
因為當時要收稅,.
所以所有人按著人口普查,.
要回到自己的鄉下,.
來數點人數,.
而耶穌的家族,.
他們就屬於伯利恆的,.
因為大衛的後裔,.
所以他們就回到伯利恆那裡數點人數,.
目的是收稅,.
就在這個時間裡面,.
耶穌就在伯利恆一個客店裡面,.
都沒有地方住,.
但是在這樣的安排之下,.
耶穌就出生在伯利恆這個小城裡面,.
按著約瑟或者瑪利亞,.
他們自己本身成長的地方,.
伯利恆不是他們成長的地方,.
但是按著加普,.
他們是,.
所以他們就去到這裡,.
去到這裡的時候,.
然後就有東方的博士,.
或者智者來到耶路撒冷,.
要找一個新的猶太王,.
誰知道他們做這件事,.
就冒犯了當時的欺律,.
就是那個時候掌權的欺律,.

$^{401}$一聽到這樣的信息,.
就問那些文士,.
發來錯人,.
那個生來的基督,.
猶太人的王在哪裡出現?.
然後,.
依書直說,.
說回尼加書那段經文,.
就在猶太地,.
然後就找那幾個東方博士,.
就去找,.
其實希律很聰明的,.
等那幾個博士去找,.
幫他去找,.
找到之後,.
然後希律就可以一次過,.
將所有新生的王剿滅,.
以免他自己的統治權受到威脅,.
所以,.
當希律王聽到,.
耶穌將要降生,.
或者耶穌降生在伯利恆,.
經文裡面第一樣東西說,.
希律王聽見了,.
就心裡不安,.
耶路撒冷合成的人也都不安,.
這個跟剛才那段經文,.
一個很大的諷刺,.
耶穌的降生,.
理應是帶來世界的平安,.
但是,.
當耶穌降生的消息,.
被那些當時的希律,.
或者耶路撒冷的人聽見,.
他們視為一個不安的信息,.
對於希律來說,.
他的不安感,.
不是因為來自耶穌,.
耶穌是一個嬰孩,.
耶穌連事情都沒做,.

$^{441}$怎麼會為他帶來不安呢?.
真正令他不安的就是,.
他心裡怕他的權力被人瓜分,.
他怕自己在歷史的洪流當中,.
已經淹沒了,.
當他失去權力的時候,.
也都代表著他失去自己的存在,.
所以在希律裡面那種不安感,.
其實不是來自上帝,.
是來自他自己的,.
來自他自己那種,.
懼怕失去權力的那種恐懼,.
令他不安,.
而因為這個希律王的不安,.
他主力去管治的那些猶太人,.
特別是居住在耶路撒冷的人,.
更加不安,.
因為那個在上者心裡面,.
經常都戚戚言的時候,.
不知道會不會又出現一些很奇怪的,.
又或者當時為了權力鬥爭,.
而帶來的一些政策,.
又或者令到他們所有的百姓,.
不知道又要面對一個新的,.
一些什麼的制度,.
所以合成的耶路撒冷的人,.
都是不安,.
耶路撒冷的人的不安,.
都不是來自耶穌的,.
最主要的不安就是,.
他們將人生所有的寄望,.
放在或者他們的生命,.
他們的人生,.
他們的去向,.
放在希律的那個喜護身上,.
所以當耶穌出現的時候,.
耶穌出世的時候,.
所有人都面對著一種很不安的那份情緒,.
不過都是那句話,.
耶穌只不過是一個嬰孩,.

$^{481}$他連神蹟都還沒有行,.
不過就單單他的出現,.
他的出世,.
一個萬世的君王,.
在一個小城裡面出現,.
就帶來一個世界裡面,.
翻天覆地的改變,.
整個經文或者去到新約,.
怎樣將這個鷹許實現呢?.
上帝的計劃是一個怎樣的計劃呢?.
整個計劃是說到,.
上帝要將這個掌權的萬有之君,.
用最卑微的方式來實現,.
當我們這個世界,.
我們經常覺得想,.
所有一些背景最彪炳的,.
或者最受到良好教育,.
甚至乎他的家裡最好是有錢,.
他的上一代是名門望族,.
出來那個就一定是明日之星,.
明日的君王,.
不過上帝的計劃卻是藉著,.
在一個木匠的家庭裡面,.
藉著一個很單純的少女瑪利亞,.
她單單這樣依靠上帝,.
耶穌就在這個家庭裡面,.
來到她出生,.
在這裡的成長,.
她的家庭不是決定了,.
她成為一個怎樣的人,.
真正影響她的,.
只是上帝的計劃,.
真正影響她的是上帝的話語,.
和上帝的鷹許,.
無論這個孩子,.
在一個卑微的家庭,.
或者在未來的日子裡面,.
用神蹟騎士來實踐救贖的工作,.
她卑微的家庭,.
絕不影響到耶穌基督實踐的救贖的大能,.

$^{521}$又或者這樣說,.
上帝看得起世界上最卑微的,.
最不起眼的小人物,.
藉著這些小人物,.
來成就這個救恩,.
我們今天,.
不知道你會怎樣看待自己的成長,.
弟兄姊妹,.
會不會我們都力拼,.
我們希望自己的背景,.
或者我們為我們的下一代,.
製造一個更加優良的背景,.
製造一個更加理想的氛圍,.
讓我們的下一代可以更加順著這個勢頭來成長,.
甚至我們去到一個地步,.
有時候這個世界有一種說法,.
你的家庭都是靠父幹的,.
如果你的生命裡面,.
你的父親能夠為你賺取第一桶金,.
有很多很豐富的背景,.
你的人生就很昌達,.
我們看這件事甚至好像一個宿命論一樣,.
爸爸有錢,.
家中有中產,.
我們就注定是一個人生的勝利者,.
相反來說,.
如果家庭是一個很基層的,.
我們就沒有辦法經歷這個生命的那種,.
那麼亮麗的一個經歷,.
不過整個福音是顛覆了我們這些價值觀念,.
上帝他的兒子耶穌基督,.
這個將要管治萬邦的君王,.
卻是選擇生在一個卑微的家庭裡面,.
上帝也顛覆了這個世界上的權力觀念,.
這個世界的權力觀念就是講到,.
誰的手臂夠硬朗,.
就可以帶來裡面他所帶領的民族的那種平安,.
平安感是建基於自己夠硬朗,.
夠強盛,.
不過上帝計劃卻是選擇用最軟弱的方式,.

$^{561}$來成就上帝最偉大的救恩,.
當這個世界的價值觀是講到,.
你要變得強盛,.
其實不是講你的強盛要此世寧人,.
或者用你的力量來壓過身邊的人,.
真正強大的就是甘於自己成為一個軟弱的人,.
甘於將自己放在一個最謙卑的位置,.
耶穌基督道成肉身,.
以奴僕的形式來成就這個救恩,.
讓這個救恩用最低調,.
最卑微,.
最不起眼的方式來成就出來,.
耶穌基督的救贖,.
改換了這個世界對人,.
對權力的看法,.
伯利恆這個地方是一個小城市,.
但是在歷史當中去到今天,.
剛才我們唱的詩歌都在講,.
伯利恆卻成為了一個偉大的城市,.
不是因為他物產豐富,.
而是上帝選擇在一個很小的城市裡,.
藉著他的兒子在這個地方誕生,.
這個地方就變成一個不一樣的城市,.
如果突然之間有去過一些歷史旅遊的景點,.
有時候你會覺得有一個狀況挺有趣的,.
我去過一個地方,.
這個地方就是某位作家用過的寫字台,.
很多人觀察,.
放了一些寫字台的設定,.
這個作家以前就是這樣寫文的,.
作了很多句子,.
一邊看的時候我覺得挺有趣,.
其實我看的都是一張寫字台,.
其實這張寫字台放在哪裡都是一樣的,.
這張寫字台不會塑造一個作家出來的,.
而是因為這個作家令這張寫字台提升了,.
人生不是靠那些外在物質,.
有這些東西來塑造一些好的人,.
或者成為一個最佳的客觀條件,.
不過這個世界就講了一個這樣的價值觀念給你,.

$^{601}$你要跑步,.
你有一雙漂亮的跑步鞋,.
你就跑得更加快,.
你要有一隻漂亮的錶,.
你就彰顯你自己是一個有名望,.
有品味的人,.
或者你穿的衣服,.
我們透過這些外在的東西來告訴別人,.
我是一個很有價值的人等等,.
不過聖經卻是重新要顛覆我們,.
真正有價值的不是那些外在的物件,.
不是那個地方,.
而是在上帝裡面,.
我重新怎麼看回我自己,.
真正的平安是在說我們在主裡面,.
合而體待自己的生命,.
我怎麼看自己的生命呢?.
真正的平安是來自我怎麼學懂尊重身邊,.
跟我一樣成長的這個人,.
真正的平安是在說這個世界,.
即使多麼困難也好,.
都是上帝所掌管的世界,.
我們因著上帝,.
我們能夠在這個世界裡面,.
平安地度日,.
我們不需要拿起我們的罩頭,.
我們也不需要用很多外在的物件,.
來標榜自己的價值,.
我們只是單單因著這個上帝,.
他的遺旨都是這樣來改變世界,.
我們就跟著耶穌這個小嬰孩,.
跟著這個應許的成就,.
我們生命慢慢得著改變,.
兄弟姐妹,聖誕節是一個充滿了喜樂,.
也充滿了很多消費,.
充滿了很多佔據你心靈的一個季節,.
就在這個時節裡面我們重新想一想,.
我們生命的價值在什麼上面呢?.
在那些外在的禮物,.
或者我們很期望想得到的那些物品身上,.

$^{641}$或者是在我們所信求的一個以軟弱方式,.
以卑微的方式成就救恩的基督耶穌身上,.
求主繼續來帶領我們,.
當我們這樣想的時候,.
我們生命會踏實很多,有力量很多的..
\newpage



\section{約翰福音 7:14-31}
\label{sec:NqwhnMUzemI}
\textbf{2024年12月29日崇拜直播｜梁家麟牧師｜令人不安的耶穌｜約翰福音7\_14-31}
\newline
\newline
連結: \href{https://youtube.com/watch?v=NqwhnMUzemI}{\texttt{https://youtube.com/watch?v=NqwhnMUzemI}} ~~~~ 語音日期: 2024-12-30
\newline
\newline
\hyperref[sec:75p0crrJvGI]{\small{< < < PREV SERMON < < <}}
~
\hyperref[sec:index_chronic]{\small{[返順時目]}}
~
\hyperref[sec:H2Yyyhsu_Qc]{\small{> > > NEXT SERMON > > >}}
\newline
\newline
約翰福音 7:14-31
\newline
\begin{longtable}{cl}
\hline
\hline
章節 & 經文 (和合本修訂版)\\
\hline
7:14 & \begin{tabularx}{0.7\textwidth}{X} 節期已過了一半，耶穌上聖殿去教導人。 \end{tabularx} \\ \\ \relax
7:15 & \begin{tabularx}{0.7\textwidth}{X} 猶太人驚訝地說：「這個人沒有學過，怎麼那樣熟悉經典呢？」 \end{tabularx} \\ \\ \relax
7:16 & \begin{tabularx}{0.7\textwidth}{X} 於是耶穌回答他們，說：「我的教導不是我自己的，而是差我來那位的。 \end{tabularx} \\ \\ \relax
7:17 & \begin{tabularx}{0.7\textwidth}{X} 人若立志要遵行神的旨意，就會知道這教導究竟是出於神，還是我憑著自己說的。 \end{tabularx} \\ \\ \relax
7:18 & \begin{tabularx}{0.7\textwidth}{X} 憑著自己說的人是尋求自己的榮耀；但那尋求差他來那位的榮耀的人，他是真誠的，在他心裡沒有不義。 \end{tabularx} \\ \\ \relax
7:19 & \begin{tabularx}{0.7\textwidth}{X} 摩西不是傳了律法給你們嗎？你們卻沒有一個人守律法。為甚麼想要殺我呢？」 \end{tabularx} \\ \\ \relax
7:20 & \begin{tabularx}{0.7\textwidth}{X} 眾人回答：「你是被鬼附了！誰想要殺你呢？」 \end{tabularx} \\ \\ \relax
7:21 & \begin{tabularx}{0.7\textwidth}{X} 耶穌回答，對他們說：「我做了一件事，你們都驚訝。 \end{tabularx} \\ \\ \relax
7:22 & \begin{tabularx}{0.7\textwidth}{X} 摩西傳割禮給你們（其實割禮不是從摩西開始，而是從列祖開始的），你們就在安息日給人行割禮。 \end{tabularx} \\ \\ \relax
7:23 & \begin{tabularx}{0.7\textwidth}{X} 人若在安息日受割禮，是為了不違背摩西的律法，我在安息日使一個人痊癒了，你們就向我發怒嗎？ \end{tabularx} \\ \\ \relax
7:24 & \begin{tabularx}{0.7\textwidth}{X} 不要憑外表斷定是非，總要按公平斷定是非。」 \end{tabularx} \\ \\ \relax
7:25 & \begin{tabularx}{0.7\textwidth}{X} 於是耶路撒冷人中有的說：「這個人不是他們想要殺的嗎？ \end{tabularx} \\ \\ \relax
7:26 & \begin{tabularx}{0.7\textwidth}{X} 你看，他還公開講道，他們也不對他說甚麼。難道官長真的認為這是基督嗎？ \end{tabularx} \\ \\ \relax
7:27 & \begin{tabularx}{0.7\textwidth}{X} 然而，我們知道這個人從哪裡來；可是基督來的時候，沒有人知道他從哪裡來。」 \end{tabularx} \\ \\ \relax
7:28 & \begin{tabularx}{0.7\textwidth}{X} 那時，耶穌在聖殿裡教導人，喊著說：「你們認識我，也知道我從哪裡來；我並不是憑著自己來的。但差我來的那位是真實的，你們不認識他。 \end{tabularx} \\ \\ \relax
7:29 & \begin{tabularx}{0.7\textwidth}{X} 我卻認識他，因為我從他那裡來，是他差遣了我。」 \end{tabularx} \\ \\ \relax
7:30 & \begin{tabularx}{0.7\textwidth}{X} 於是他們想要捉拿耶穌，只是沒有人下手，因為他的時候還沒有到。 \end{tabularx} \\ \\ \relax
7:31 & \begin{tabularx}{0.7\textwidth}{X} 但人群中有好些人信他，他們說：「基督來的時候，他所行的神蹟難道會比這人行的更多嗎？」 \end{tabularx} \\ \\
[1ex]
\hline
\hline
\end{longtable}
$^{1}$今日經訓記載於約翰福音第七章第十四至三十一節.
到了節期,耶穌上殿裡去教訓人.
猶太人就輕啟說:.
「這個人沒有學過怎麼明白書呢?」.
耶穌說:.
「我的教訓不是我自己的,乃是那差我來者的」.
「人若立志遵著他的旨意行,就必曉得這教訓或是出於神,或是我憑著自己說的」.
「人憑著自己說,是求自己的榮耀,唯有求那差他來者的榮耀,這人是真的,在他心裡沒有不義」.
「摩西豈不是全律法給你們嗎?」.
「你們卻沒有一個人守律法,為什麼想要殺我呢?」.
眾人回答說:.
「你是被鬼附著了,誰想要殺你?」.
耶穌說:.
「我作了一件事,你們都以為稀奇」.
「摩西全國禮給你們,其實不是從摩西起的,乃是從先祖起的」.
「因此,你們也在安息日給人行國禮」.
「人若在安息日行國禮,免得違背摩西的律法」.
「我在安息日,叫一個人傳言好了,你們就向我生氣嗎?」.
「不可按外貌斷定是非,總是按公平斷定是非」.
耶路撒冷人中的柔的說:.
「這不是他們想要殺的人嗎?」.
「你看他還明明的說道,他們也不向他說什麼」.
「難道官長真知道這是基督嗎?」.
「然而,我們知道這個人從哪裡來」.
「只是基督來的時候,沒有人知道他從哪裡來」.
那時耶穌在殿裡教訓人大聲說:.
「你們也知道我,也知道我從哪裡來」.
「我來並不是由於自己,但猜我來的是真的」.
「你們不認識他,我卻認識他」.
「因為我是從他來的,他也是猜了我來」.
「他們就想要抓拿耶穌,只是沒有人下手」.
「因為他的時候還沒有到」.
「但眾人中間有好些信他的說」.
「基督來的時候,他所作的神蹟」.
「豈能比這人所行的更多麼驚憤獨不」.
弟兄姊妹,預祝你們新一年新年快樂.
生活豐豐富富.
我們繼續《約翰福音》漫長的旅程.
上一次說到在駐彭捷.
耶穌本來不打算去耶路撒冷守捷.

$^{41}$因為恐怕猶太教當局要殺害他.
不過後來他仍然低調去了.
不過像耶穌這樣的風頭人物.
去到哪裡都惹人注目.
要低調也很難.
這一段記載了耶路撒冷所引起的很多爭議.
耶穌真是一個搞亂天下的人.
比保羅的搞亂天下還要早很多.
這一趟有幾個風波.
第一波的爭議出現在耶穌和一班猶太人中間.
爭論的焦點是耶穌說到的權威.
耶穌雖然說要低調.
不過他竟然在聖殿裡站出來教訓人.
或者耶穌覺得與其讓身邊的人各自去猜測.
估計耶穌是誰,他說什麼,言滿天非.
真真假假假假真真.
不如他親自出來講道,講他自己的觀點.
聽眾同意不同意是一件事.
最少可以消除一些不必要的謠傳誤會.
按著約翰福音的記載.
這是耶穌第一次在聖殿裡講道.
之前他曾經在聖城附近行過神蹟.
但沒有提到他講道.
他的聽眾在聖殿裡.
多數都是從各地來守節的猶太人.
很多人第一次聽耶穌講道.
就感覺到很震撼.
耶穌在加里尼山頭野嶺講道講過很多次.
在聖殿,耶路撒冷的聖殿是第一次.
第一次聽他講道的人很震撼.
聽眾中有人稀奇說.
耶穌沒有受過經學訓練.
他從哪裡得著律法的知識呢?.
他們稀奇的不僅僅是.
耶穌具備講經的律法知識.
更加稀奇的是.
耶穌講的東西.
跟他們過去一直聽的觀點很不同.
他有特別的動見.
並且耶穌的講道有種特別的魅力.

$^{81}$總之你聽他講話.
不僅聽到他的觀點.
你只會覺得很有權威.
為什麼這個人講話這麼有權威呢?.
很想去跟隨他.
他們問耶穌的觀點.
是從哪一個拉比學派.
Mabinic School來的呢?.
為什麼從來沒有聽過呢?.
那時候的猶太人解經有幾個不同的學派.
他們叫拉比學派.
就好像我們的武當,阿彌,空洞.
有很多派的.
但是耶穌是哪一派呢?.
他們馬上問這個問題.
我們知道耶穌的講道是多麼奇怪.
聽他講道的人都有種稀奇.
《馬太福音》記載.
耶穌在登山寶訓之後.
聽眾的反應就是.
眾人都稀奇他的教訓.
因為其他教訓他們.
像有權柄的人.
不像他們的文士.
為什麼這麼有權威呢?.
耶穌在聖殿裡回答這些人.
我的教訓不是我自己的.
乃是那差我來者的.
人若立志遵著他的旨而行.
就必曉得這教訓.
是或是諸於上帝.
或是我憑著自己說的.
第十六十七節.
和傳統中國人一樣.
猶太人是一個很崇尚古人的民族.
不喜歡講某個觀點.
某個解說是自己的.
我們和古代人不同.
我們抄回來都當是自己的原創.
古代人明明是自己原創.

$^{121}$都當是抄回來的.
要假裝抄回來的.
所以我們看古代文件.
有很多「偽XX」的文件.
「書刀XX」.
即是說是孔子寫的.
但其實不是孔子寫的.
是某個人寫了.
但不喜歡出自己的名字.
就說是孔子寫的.
我們今天很難想像這樣的情況.
不過古代是有的.
因為你想自己的作品流傳於世.
你說梁家鑾寫的沒有人看.
你說王明道寫的沒有人看.
於是就寫了王明道.
即是我的筆名.
古代人喜歡說來說去.
都說甚麼子曰,語雲.
孔子怎麼說,論語怎麼說.
耶穌在這裡宣稱.
祂的傳息不是祂自創的.
強調也不是來自任何人間的拉比傳統.
而是從上帝來的.
律法由上帝頒佈.
耶穌的傳息如果是由上帝而來.
就具有終極權威.
是嗎?不可能有錯誤的.
接著耶穌說了一句很重要的話.
祂是一個真正敬畏上帝.
願意遵祂的道而行的人.
也就是說.
你聽到是準備毀身於真理.
有那個誠信的人.
就一定能夠分辨出.
耶穌說的教訓到底是出自祂自己.
到底是出於上帝還是出於人意.
我們知道.
講信仰,講神學,講宗教的人.
有兩種.

$^{161}$一種純粹是吹水活動.
將神學變成思想或概念的遊戲.
天馬行空為辯論而辯論.
所討論的與個人信仰和生活毫無關係.
與教會和福音使命更不相干.
另一種目的是.
要更準確地明白上帝的生意.
然後付諸實踐.
要更悉切,更有效地完成福音使命.
後者就是有誠信的人.
正如我經常說.
在大學中教系讀基督教.
與在神學院讀基督教是兩碼子的事.
宗教學是將宗教讀為社會現象,文化現象.
去吹水,講理論,思想.
天馬行空.
說說不准.
不需要有信仰.
但我們要求的是一個有委身心智.
真正相信你所說的.
願意委身在其中.
很多年前.
在國內有一班研究基督教學者.
說了很多漢語神學.
不知道大家有沒有聽過.
港台也有一班基督教學者加入唱和.
宣稱他們做的是文化神學,公眾神學.
有別於教會的教會神學.
教會神學是狹隘的.
他們的公眾神學才能夠出台面.
在我看來.
這些所謂神學不折不扣就是吹水活動.
漢語神學的領軍人物.
叫劉曉峰.
理論天馬行空.
生活極其美濫的人.
結離三次.
又結又離.
生活非常混帳.
這叫文化基督徒.

$^{201}$我從一開始就看不起這些人.
恥與他們為伍.
聖經說.
「立志遵上帝的旨行」.
這是真信仰和假信仰的分別所在.
真神學和假神學的分別所在.
上帝不可以作為一個被研究,被分析的對象.
讓你去猜測,想像.
必須是你敬畏祂,願意聽祂說.
是憑信心實踐的行動.
你才可以認識祂.
上帝唯獨是你在信仰中.
I am believing.
我正在信的時候你才認識到祂.
再說尋索上帝旨的人.
不可能是一個置身事外的局外者.
沒有人可以用一個旁觀者.
一個bystander的身份去明白上帝的旨意.
必須要立志成為上帝旨意的一部分.
這樣才能明白上帝的旨意.
上帝對未來的旺鎮有什麼旨意?.
上帝對未來的香港有什麼旨意?.
你吹牛是沒有意義的.
唯有願意投身其中.
祈禱,行動,承擔.
然後我們祈禱.
願上帝的旨意成就在我們中間.
我們願意委身其中.
一個願意參與委身在旺鎮.
塑造旺鎮未來的人.
就可以說旺鎮的未來是怎樣.
我是否一個極極的人?.
你聽懂我說話嗎?.
置身事外指指點點.
慳跌慳跌.
你不會知道.
因為我講了一句.
聽起來有些離經半道的話.
我們教務同工.
常常在講台上.

$^{241}$跟你們說上帝是誰.
上帝是怎樣怎樣怎樣.
但真相是.
沒有人可以告訴你上帝是怎樣的.
連聖經都做不到.
其他人或聖經.
可以跟你介紹有關上帝的資料.
但能讓你認識上帝是誰的.
只有上帝和你自己.
當你和上帝面對面的時候.
你就認識從前風聞的那個他.
現在親眼看到他.
上帝,原來是你.
同樣的,甚麼是上帝旨意?.
當你親耳聽到他的聲音.
聽到他對你的教訓和猜險.
你也聽到你自己.
自願或不自願的說.
我在這裡,僕人敬聽.
你就遇上了上帝的旨意.
回到聖經.
耶穌突然將話題一轉.
就指出,講道的時候.
專門講個人想法的人.
目的是要抬高自己的地位.
追求自己的榮耀.
唯有只是講上帝的信息的人.
要求的是上帝的榮耀.
而不是自己的榮耀.
是全心高舉真理的人.
心裡無不疑.
十八節,「人是真的,在他心裡沒有不疑」.
耶穌說的當然是祂自己.
很多人從實用主義或功能主義的角度去看真理.
問,這個真理對我有甚麼好處?.
高舉真理為我帶來甚麼好處?.
如果只有我發現真理.
我就可以壟斷真理,抬高自己,貶低別人.
耶穌教導我們.
追求真理的目的是預備好.

$^{281}$躲在其中,全心全情地服從真理.
預備好所有的代價去堅持真理.
像孔子說的,「朝問道,直死可疑」.
耶穌轉過來質問猶太人.
十九節,「摩西希不是傳律法給你們嗎?」.
「你們卻沒有一個人遵守律法,為甚麼要殺我?」.
耶穌不是經驗上知道眼前所有人都守不住律法.
而是原則上認定沒有人能夠守律法.
像保羅對律法的講論.
律法只能寫人之罪,羅馬書三章的說法.
在律法的標準對照之下.
每個猶太人都是不義的人.
耶穌問,不義的人竟然想殺一個沒有不義的他.
這是為了甚麼呢?.
惡人殺善人,這樣有天理嗎?.
跟耶穌對話的猶太人立即否認.
二十節,「你是被鬼附著了,誰想要殺你?」.
耶穌揭露了他們心底的想法.
不過由於他們還未付諸行動.
當然就此口否認.
未做就即是沒有.
他們倒過來指控耶穌,你是被鬼附的.
指責對方被鬼附是最常對人攻擊的手段.
完全不用證據.
耶穌常常遭遇這樣的指控,好幾次.
耶穌列舉了他們反對他的理由.
並且指出這個理由的荒謬.
耶穌之前來的耶路撒冷.
在西羅亞遲止.
叫一個躺在玉子的三十八年的坦子希尼行走.
我們沒有解讀這段經文.
因為我以前講過這段經文.
你應該記得,西羅亞遲止.
治好三十八年躺在那裡的病人.
那天正好是安息日.
猶太人在那一次就指控耶穌破壞守安息日的規定.
為了對照容.
耶穌提到割禮二十二節.
摩西傳割禮給你們.
其實不是摩西起的,而是從祖先起的.

$^{321}$因此你們也在安息日,給人行割禮.
割禮的規定是早於摩西的律法.
在亞伯拉罕的時候就已經確定了.
所以其實不是摩西創造的.
不過摩西將割禮納入在律法裡面.
成了律法的一部分.
這出在《十二章》第四十四節和後面的文章.
律法規定,一個男孩出生後第八天就要守割禮.
如果那個男孩是在安息日出生.
那麼他就在第二個安息日守割禮.
你明白我說什麼嗎?.
第八天嘛,第八天就是七天後.
如果他在安息日出生.
那麼律法怎麼規定都不可以禁止人在安息日出生.
生孩子什麼時候出生,怎麼輪到你開刀.
總之安息日出生了,然後安息日就要守割禮.
耶穌於是問,治病是破壞安息日規定.
但是猶太人在安息日為人施行割禮.
那是不是破壞安息日呢?.
猶太人的說法就是.
割禮的規定是超越安息日的禁令.
為什麼行割禮比治病更加緊要更加緊急呢?.
行割禮可以在安息日進行.
但是治病不可以呢?.
生命的價值不是最高的嗎?.
救命不是最緊急最重要的事嗎?.
不是,他都癱瘓了38年都不差一天.
不急於一定要安息日做.
那是呀,那守割禮有什麼緊急性呢?.
退後一天,星期一做行不行呢?.
是吧,不行的,怎麼會不行呢?.
那怎麼說呢?猶太人的做法都是說不通的,是不是?.
耶穌於是做了這樣的評論,在24節.
不按外貌斷定是非,還要按公平斷定是非.
不要糾纏於某些外在行為能做不能做.
總要將這些外在行為變成一堆真假對錯的判斷標準.
例如在中國教會,特別是華中地區農村的教會.
仍然有規定婦女要蒙頭,並且不可以吃血的爭論.
很多人將蒙頭不蒙頭,吃血不吃血.
看為正統和異端的分別所在.

$^{361}$從前保守的教會,包括我的建築學院.
都有很多行為的規定,很多do and don't.
不准看電影,不准聽流行音樂,不准潮流服飾.
為了便於執行,我們總是會畫一些紅線出來.
定出什麼是正確,不正確,屬靈不屬靈.
女生穿的長裙要長到膝蓋下面一吋.
那一吋就決定你屬靈不屬靈.
你差了0.1,那就不屬靈了.
總之就是這樣,定了這個標準.
以前中學也規定男生頭髮不長到耳朵.
不准碰到耳朵,碰到耳朵就不屬靈了.
總之就是這樣,定了這樣的規矩.
今天我們教會對這些標準應該不會再堅持了.
以前是會堅持的.
18世紀的流行音樂是聖樂.
21世紀的流行音樂就是俗樂.
為什麼三個世紀前是聖樂,現在是俗樂呢?.
是沒得解釋的.
總之現在我們沒有這個規定.
什麼電影,什麼音樂是好是壞.
要區分對待,不是全盤否定.
這是第一波的爭議.
第二波的爭議是在耶路撒冷的猶太人中間興起的.
他們在爭論耶穌是不是基督.
25-27節,「這不是他們想要殺的人嗎?」.
「你看他們說明明的講道,他們也不向他說什麼」.
「難道官長真知道這是基督嗎?」.
「然而我們知道這個人從哪裡來」.
「自從基督來的時候,每個人都知道他從哪裡來」.
他們首先猶太人猶太人仍然大肆在聖殿講道.
沒有被猶太教的管理層去干預他們.
他們問猶太教的領袖中庸.
他講道是不是意味他們認定了,肯定了耶穌的基督身份呢?.
如果連領袖都肯定耶穌的基督身份.
原則上他們不應該反對.
他們要信奉宗教領袖.
但是他們心中有一個極有前.
舊約提到彌賽亞的時候.
總是說他是來無終,去無影的.
但是這個耶穌,他們是知道他的出身的.

$^{401}$是出身於不會出事的那撒勒.
所以耶穌怎麼可能是基督呢?.
宗教領袖說他是基督.
他們都說我們沒辦法接受這個事實.
你知道在古代社會.
家庭,社會,家族出身是很決定性的因素.
龍生龍,鳳生鳳.
老鼠的兒子會打動.
畢竟社會資源缺乏.
多數教育或職業訓練都是由家庭提供的.
家庭為你提供開展事業的一切.
所需要的土地,財政資源,人脈關係.
全部都是家族供應給你的.
個人奮鬥成長空間很少.
什麼白手興家,萬中無一.
所以每個人都靠父幹.
也因此被限制了職業的選擇和發展的可能性.
中國歷史常說農民革命,官逼民反,揭竿起義.
但是你看這麼多個朝代.
真正農民出身的開國領袖只有明朝的朱元璋.
沒有了.
農民暴動,農民革命最後都是被貴族揀了便宜.
你不是貴族出身,省一點吧.
你怎麼可能坐莊做皇帝呢?.
前面猶太人詫異耶穌為什麼會解釋律法.
因為如果不是出身文士和法理賽人世家.
一般人連律法書的書皮都摸不到.
都不識字,怎麼有機會去學習《釋經》講道呢?.
那時候還沒有什麼網上神學課程,省一點吧.
因此當群眾知道耶穌出身自北部的加利利的拿撒勒的時候.
雖然未必馬上說他是賤民階層.
不過也不相信那裡會跑出一個怎樣厲害的人物.
什麼天眾英才,無私自通.
省一點吧,沒有這樣的事.
耶穌知道群眾對他的出身的議論和質疑.
就在聖殿講道的時候做出一個公開的回覆.
28-29節.
耶穌有兩本護照.
一本是天上的,一本是地上的.
群眾執著的是耶穌人間的出身.

$^{441}$他的直觀是拿撒勒人.
但耶穌卻強調他的神聖出身.
他是從父上帝所差來的.
要傳遞的是父上帝的旨意.
這番話在猶太人的耳中是非常具褻瀆性的.
為什麼有人會宣稱自己是從上帝所差來的呢?.
如果我真的相信耶穌講的這番話.
我就要趴在他面前.
古代皇帝派遣那些傳聖旨的天差來.
那些宦官叫天差.
當天差宣告聖旨到的時候.
人們要怎麼做?.
就趴在那裡.
如果你說你真的從天父那裡來.
我就要趴在你面前.
有沒有搞錯?.
在古代假傳聖旨的人要被殺頭.
所以耶穌講了這番話.
如果我們不在他面前趴下.
我們就要殺他了.
當時已經有人想要捉拿耶穌.
不過由於在場的人.
沒辦法達成一致的意見.
有人反對他.
也有人支持他.
所以沒辦法立即做那個逮捕的行動.
引起社會的騷動.
支持耶穌是基督的人.
他們推理的理由是.
耶穌行了這麼多神蹟.
如果他不是基督.
而是有真基督來.
大概那個真基督.
也行不了耶穌這麼多神蹟的.
31節.
這段經文我們解讀到這裡.
其實整張聖經都是在講這幾個爭論.
有四波的爭論.
我們今天講了兩個.
你可以看到.

$^{481}$耶穌的出現.
為當時的人帶來很多的困擾.
耶穌是誰?.
他是彌賽亞嗎?.
如果說他是.
他有很多地方不是很符合.
當時的人對彌賽亞的期望.
例如他出生在加利利.
例如他出生在卑賤.
他沒有家型美容.
如果耶穌出生在猶太的顯貴家庭.
或者是荊學世家的文士和法利塞人.
就會減少很多人對他的懷疑.
因為耶穌不是彌賽亞.
但是他有很多言語行為.
顯出他真的不同範品的特性.
例如他講到的內容不僅僅很.
很深刻.
更加常常顯出.
有很多不尋常的肅靈權病.
耶穌經常這樣說.
你們聽過古人這樣說.
「主使我告訴你們」.
哇.
耶穌的權威好像比所有古人都大.
比古時的拉比都大.
這些權病是出生在窮鄉僻壤的.
無名之輩.
很難裝空作勢.
更加不要說耶穌行了很多神蹟騎士.
治好癱子.
五餅二魚餵飽了幾千人.
這也不是那些遊走江湖的.
招搖撞騙的神棍能夠做到的.
你可以用什麼樣的掩眼法.
可以讓一萬人吃飽呢.
不可能的.
所以說耶穌是彌賽亞.
他不完全像.
說耶穌不是彌賽亞.

$^{521}$他又有點像.
你自覺知道.
作為耶穌的開路先鋒的斯一約翰.
在他坐牢等死之前.
他也托門徒問耶穌一個問題.
這是斯一約翰埋藏在心底多年的困擾.
馬特伯恩十一章三節說.
那將要來的是你嗎.
還是我們要等候別人呢.
連斯一約翰都十五十六.
不確定耶穌是不是彌賽亞.
他明明應該知道.
他自小就應該知道.
馬尼亞的身份會比彌賽亞更年輕.
應該斯一約翰的父母告訴他.
但是他看到耶穌的樣子.
整個樣子.
有點像梁家倫.
喜歡吃喝的那些.
都不太熟悉.
應該很熟悉的.
不吃人間煙火.
不會吃雞腿.
不會喝紅酒.
做這些事.
整個樣子都不像彌賽亞.
斯一約翰都十五十六.
很忐忑.
是嗎?不是嗎?.
所以耶穌並未符合斯一約翰心中所有彌賽亞的標準.
我們知道耶穌是否彌賽亞.
這個問題是沒有中立的可能.
每個人都要做決定.
要表態.
如果不是彌賽亞就要趴在他面前.
要敬拜他跟隨他.
如果不是彌賽亞.
他就是最大的宗教騙子.
褻瀆神明.
由於人對耶穌產生對立分歧的意見.

$^{561}$耶穌的出現就為猶太人的社會帶來矛盾.
衝突.
混亂甚至分裂.
從社會影響來說.
造成混亂和分裂是必然的壞事.
作為被羅馬軍隊監控的亡國之民.
任何風吹草動都會帶來統治者的暴力鎮壓.
為人帶來傷害和苦難.
從社會和諧穩定壓倒一切的政治考慮之下.
耶穌是為社會帶來危害的人物.
最好是將他除掉.
就算不從社會的高層次去考慮.
純粹從個人的角度思考.
耶穌的出現同樣是為我們帶來不安甚至威脅.
當我們認識一件新的事物的時候.
我們總是喜歡將它歸聘到原來的思想架構裡.
分類歸黨.
這樣就可以減少那件新事物的外來性.
不致構成對我們原有的想法的威脅.
你跟我說一件事.
我會說「就是這樣」.
「這樣」的時候我就會將你的案例.
將它拼在我既有的想法裡.
如果你講的東西是完全無法分類歸黨的話.
你就搞亂黨.
成為了麻煩友.
你會破壞我原有的想法.
破壞我原來的生活秩序.
碰到這樣的情況.
我們通常只有兩個選擇.
第一是徹底拒絕外來的新事物.
第二是準備好將自己原有的思想架構重新整理.
推倒重來.
然後去接納這個新的事物.
我以前都跟你們說過.
有兩種信仰的態度.
你可以將信仰看為你的補充劑.
或者信仰看為你的顛覆者.
信仰如果只是一個補充品.
你的生活是如常不變的.

$^{601}$你可以維持你原有生活的一切平衡.
加上一點信仰.
就像你以前如果你拜觀音.
拜黃大仙.
他完全不干涉你繼續打老婆.
繼續打,無所謂的.
繼續賭錢,繼續做壞事.
無所謂的.
他不干涉你的生活.
但是如果信仰是顛覆者.
是subversive(反覆)的.
你的生活就準備好要被全盤翻轉.
所以,兄弟姐妹.
如果你信耶穌.
你覺得這個耶穌是左頭左勢.
你不相信.
他經常在你生命裡.
閃閃耀耀,嘰嘰噤噤.
這個上帝就是一個左頭左勢.
比你的外母還難受.
比你太太還可怕的那個.
如果你信的信仰.
是比你的外母更陰春.
大概你的信仰多數是偶像.
不是真的信仰.
當你覺得這個信仰.
是威脅到你既有的生活的時候.
就是你自己準備好.
被拆毀,毀壞,傾覆,建立和栽植的時候.
兄弟姐妹.
耶穌一定是你無法將祂安置好.
祂是無法定義你的.
祂不僅僅是你的拯救者和幫助者.
祂不僅僅帶給你安慰和賜福.
不僅僅做一些你能夠明白並且同意的事情.
不僅僅做一些符合人情常理.
在你計劃之內.
讓你有足夠的心理準備去迎接的安排.
如果上帝是上帝.
祂一定會帶給你很多驚喜.

$^{641}$更多的驚嚇.
會多次擾亂你既有的安排.
為你加插很多的意外.
甚至會挑起你的憤怒,憤慨,埋怨.
讓你產生這個那個的信仰危機.
聖經多次預告.
約伯記三十七章五節.
恆大事我們不能夠測透.
約伯記十一章七到八節.
你考察就能測透上帝嗎?.
你豈能盡情測透全能者嗎?.
他的智慧高於天.
你還能作甚麼?.
心於心間.
你還能知道甚麼?.
這樣的知識奇妙.
是我所不能測的.
至高是我所不能及的.
《詩篇》一百三十九篇第六節.
《儀賽亞書》五十五章八到九節.
我們更加熟悉.
「說:我的意念非同你們的意念.
我的道路非同你們的道路.
天怎樣高過地.
照樣我的道路高過你們的道路.
我的意念高過你們的意念」.
保羅在解預定論.
解到自己一頭糊塗一頭煙之後.
就做了一個讚美.
《羅馬書》十一章三十三到三十四節.
心在上帝的豐富智慧和知識.
他的判斷何其難測.
他的終積何其難尋.
誰知道說主的心.
誰作過他的謬視呢?.
丁姐妹.
問一問你自己.
你覺得上帝最近那一次跟你說話.
你覺得信仰為你帶來干擾,帶來麻煩.
最近那一次是何時?.

$^{681}$你面對上帝心中作難的時候.
最近是何時?.
你有多少次感覺到.
你不明白上帝在你生命裡做了甚麼.
為何他會這樣教我.
你最近那一次是何時?.
我曾經說過信主幾十年.
我的經驗總是.
每當上帝剝奪我某些東西的時候.
就是他要跟我說特別的話.
我的靈命可以有特別的突破的時候.
同樣每一次當我感覺到強烈的不安,不解,不認同的時候.
通常就是迫我直接面對上帝.
亦面對我自己.
從而帶來新的關係,新的認識的時候.
一個太平靜,平順的信仰.
是很難有甚麼突破的.
我求主向你們每一個人養念.
讓你們真實地去認識他.
我求主識到的.
不要過份.
識到的去動你們.
讓你們在焦慮,不安裡面跟他面對面.
你們會見到他,聽到他.
一個真真實實的上帝.
會翻轉,改變我們的生命.
直到我們完全符合他的心意.
願他的旨意成就在我們每一個人之上.
謝謝大家收看 再見.
\newpage



\section{申命記 1:19-46}
\label{sec:H2Yyyhsu_Qc}
\textbf{2025年1月5日主餐崇拜直播｜陳明泉牧師｜朝向應許之地－要懂得選擇｜申命記1\_19-46}
\newline
\newline
連結: \href{https://youtube.com/watch?v=H2Yyyhsu_Qc}{\texttt{https://youtube.com/watch?v=H2Yyyhsu\_Qc}} ~~~~ 語音日期: 2025-01-06
\newline
\newline
\hyperref[sec:NqwhnMUzemI]{\small{< < < PREV SERMON < < <}}
~
\hyperref[sec:index_chronic]{\small{[返順時目]}}
~
\hyperref[sec:YgvzlIV1PPE]{\small{> > > NEXT SERMON > > >}}
\newline
\newline
申命記 1:19-46
\newline
\begin{longtable}{cl}
\hline
\hline
章節 & 經文 (和合本修訂版)\\
\hline
1:19 & \begin{tabularx}{0.7\textwidth}{X} 「我們照著耶和華－我們神所吩咐的，從何烈山起行，經過你們所看見那一切大而可怕的曠野，往亞摩利人的山區去，到了加低斯‧巴尼亞。 \end{tabularx} \\ \\ \relax
1:20 & \begin{tabularx}{0.7\textwidth}{X} 我對你們說：『你們已經到了耶和華－我們神所賜給我們的亞摩利人之山區。 \end{tabularx} \\ \\ \relax
1:21 & \begin{tabularx}{0.7\textwidth}{X} 看，耶和華－你的神已將那地擺在你面前，你要照耶和華－你列祖的神所說的，上去得那地為業。不要懼怕，也不要驚惶。』 \end{tabularx} \\ \\ \relax
1:22 & \begin{tabularx}{0.7\textwidth}{X} 你們都來到我這裡，說：『讓我們先派人去，為我們窺探那地，把我們上去該走的路線和該進的城鎮回報我們。』 \end{tabularx} \\ \\ \relax
1:23 & \begin{tabularx}{0.7\textwidth}{X} 這話我看為美，就從你們中間選取十二個人，每支派一人。 \end{tabularx} \\ \\ \relax
1:24 & \begin{tabularx}{0.7\textwidth}{X} 於是他們起身上山區去，到以實各谷，窺探那地。 \end{tabularx} \\ \\ \relax
1:25 & \begin{tabularx}{0.7\textwidth}{X} 他們的手帶著那地的一些果子，下到我們這裡，回報我們說：『耶和華－我們的神所賜給我們的是美地。』 \end{tabularx} \\ \\ \relax
1:26 & \begin{tabularx}{0.7\textwidth}{X} 「你們卻不肯上去，竟違背了耶和華---你們神的指示， \end{tabularx} \\ \\ \relax
1:27 & \begin{tabularx}{0.7\textwidth}{X} 在帳棚內發怨言說：『耶和華因為恨我們，所以將我們從埃及地領出來，要把我們交在亞摩利人的手中，除滅我們。 \end{tabularx} \\ \\ \relax
1:28 & \begin{tabularx}{0.7\textwidth}{X} 我們上哪裡去呢？我們的弟兄使我們膽戰心驚，說那裡的百姓比我們又大又高 ，那裡的城鎮又大，城牆又堅固，如天一樣高，並且我們在那裡看見亞衲族人。』 \end{tabularx} \\ \\ \relax
1:29 & \begin{tabularx}{0.7\textwidth}{X} 我就對你們說：『不要驚惶，也不要怕他們。 \end{tabularx} \\ \\ \relax
1:30 & \begin{tabularx}{0.7\textwidth}{X} 在你們前面行的耶和華－你們的神必為你們爭戰，正如他在埃及，在你們眼前為你們所做的一樣； \end{tabularx} \\ \\ \relax
1:31 & \begin{tabularx}{0.7\textwidth}{X} 並且你們在曠野所行的一切路上，也看見了耶和華---你們的神背著你們，如同人背自己的兒子一樣，直到你們來到這地方。』 \end{tabularx} \\ \\ \relax
1:32 & \begin{tabularx}{0.7\textwidth}{X} 你們在這事上卻不信耶和華---你們的神。 \end{tabularx} \\ \\ \relax
1:33 & \begin{tabularx}{0.7\textwidth}{X} 他一路行在你們前面，為你們尋找安營的地方；他夜間在火中，日間在雲中，指示你們當走的路。」 \end{tabularx} \\ \\ \relax
1:34 & \begin{tabularx}{0.7\textwidth}{X} 「耶和華聽見你們的怨言，就發怒，起誓說： \end{tabularx} \\ \\ \relax
1:35 & \begin{tabularx}{0.7\textwidth}{X} 『這邪惡世代的人，一個也不得看見我起誓要賜給你們列祖的美地； \end{tabularx} \\ \\ \relax
1:36 & \begin{tabularx}{0.7\textwidth}{X} 惟有耶孚尼的兒子迦勒必得看見，並且我要將他所踏過的地賜給他和他的子孫，因為他專心跟從我。』 \end{tabularx} \\ \\ \relax
1:37 & \begin{tabularx}{0.7\textwidth}{X} 耶和華也因你們的緣故向我發怒，說：『你也不得進入那地。 \end{tabularx} \\ \\ \relax
1:38 & \begin{tabularx}{0.7\textwidth}{X} 那侍候你，嫩的兒子約書亞必得進入那地。你要勉勵他，因為他要使以色列承受那地為業。 \end{tabularx} \\ \\ \relax
1:39 & \begin{tabularx}{0.7\textwidth}{X} 你們的孩子，你們說要成為擄物的，就是今日尚不知善惡的兒女，必進入那地。我要將那地賜給他們，他們必得為業。 \end{tabularx} \\ \\ \relax
1:40 & \begin{tabularx}{0.7\textwidth}{X} 至於你們，要轉回，從紅海的路往曠野去。』 \end{tabularx} \\ \\ \relax
1:41 & \begin{tabularx}{0.7\textwidth}{X} 「你們回答我說：『我們得罪了耶和華！現在我們願遵照耶和華---我們神一切所吩咐的上去爭戰。』於是你們各人帶著兵器，以為很容易就能上到山區去。 \end{tabularx} \\ \\ \relax
1:42 & \begin{tabularx}{0.7\textwidth}{X} 耶和華對我說：『你對他們說：不要上去，也不要爭戰，因我不在你們中間，恐怕你們在仇敵面前被擊敗。』 \end{tabularx} \\ \\ \relax
1:43 & \begin{tabularx}{0.7\textwidth}{X} 我就告訴了你們，你們卻不聽從，竟違背耶和華的指示，擅自上到山區去。 \end{tabularx} \\ \\ \relax
1:44 & \begin{tabularx}{0.7\textwidth}{X} 住在那山區的亞摩利人像蜂群一樣出來迎擊你們，追趕你們，在西珥擊敗你們，直到何珥瑪。 \end{tabularx} \\ \\ \relax
1:45 & \begin{tabularx}{0.7\textwidth}{X} 你們就回來，在耶和華面前哭泣；耶和華卻不聽你們的聲音，也不向你們側耳。 \end{tabularx} \\ \\ \relax
1:46 & \begin{tabularx}{0.7\textwidth}{X} 你們照著所停留的日子，在加低斯停留了許多日子。」 \end{tabularx} \\ \\
[1ex]
\hline
\hline
\end{longtable}
$^{1}$《申命記》第19-46節.
我照著耶和華我們神所吩咐的.
從荷列山起行.
經過你們所看見的大而可怕的曠野.
往阿摩利人的山地去.
到了迦底斯巴尼亞.
我對你們說.
你們已經到了耶和華我們神所設給我們的阿摩利人之山地.
看啊!耶和華你的神已將那地擺在你面前.
你要照耶和華你列祖的神所說的.
上去得那地為業.
不要懼怕也不要驚慌.
你們就近來對我說.
我們要先打發人去為我們窺探那地.
將我們上去該走何道必進何城都回報我們.
這話我以為美.
就從你們中間選了十二個人.
每支派一人.
於是他們起身上山地去到以實國谷窺探那地.
他們手裡拿著那地的果子下來.
到我們那裡回報說.
耶和華我們的神所設給我們的是美地.
你們卻不肯上去.
竟違背了耶和華你們神的命令.
在帳篷內發怨言說.
耶和華恩恨我們.
所以將我們從埃及地領出來.
要皆在阿摩利人手中除滅我們.
我們上那裡去呢?.
我們的弟兄使我們的心消化.
說那地的民比我們又高又大.
城邑又廣大又堅固.
高得頂天.
並且我們在那裡看見阿納族的人.
我就對你們說.
不要驚慌也不要怕他們.
在你們前面行的耶和華你們的神.
必為你們征戰.
正如祂在埃及地和曠野.
在你們眼前所行的一樣.

$^{41}$你們在曠野所行的路上.
也曾見耶和華你們的神撫養你們.
如同人撫養兒子一般.
直等到你們來到這地.
你們在這世上卻不信耶和華你們的神.
祂在路上在你們前面行.
為你們找安寧的地方.
夜間在火柱裡日間在雲柱裡.
只是你們所行的路.
耶和華聽見你們這話就發怒.
起誓說.
這惡世代的人.
連一個也不得見我起誓.
應許賜給你們列祖的美地.
惟有耶夫尼的兒子迦勒必得看見.
並且我要將他所踏過的地.
賜給他和他的子孫.
因為他專心跟從我.
耶和華為你的緣故也向我發怒.
說:你必不得進入那地.
事後你的孫子約書亞.
他必得進入那地.
你要勉勵他.
因為他要使以色列人承受那地為業.
並且你們的婦人孩子.
就是你們所說必被擄掠的.
和今日不知善惡的兒女.
必進入那地.
我要將那地賜給他們.
他們必得為業.
至於你們要轉回.
從紅海的路往曠野去.
乃是你們回答我說.
我們得罪了耶和華.
情願照耶和華我們神一切所吩咐的.
上去征戰.
於是你們各人帶著兵器增仙上山地去了.
耶和華吩咐我說.
你對他們說不要上去.
也不要征戰.

$^{81}$因為我不在你們中間.
恐怕你們被受敵殺敗了.
我就告訴了你們.
你們卻不聽從.
竟違背耶和華的命令擅自上山地去.
住那山地的阿摩利人.
就出來攻擊你們追趕你們.
如蜂擁一般.
再西耳殺退你們.
直到荷爾蒙.
你們便回來在耶和華面前哭號.
耶和華卻不聽你們的聲音.
也不向你們質疑.
於是你們在迦底斯住了許多日子.
弟兄姊妹平安.
平安.
預祝大家在新的一年裡面.
我們經歷上帝的恩典更深.
在2025年我們過得更加精彩.
更加有力量.
在進入新的一年裡面.
我們開始一個新的講道主題.
朝向應許之地.
我們這次這個系列要講半年.
我們用這半年時間.
一起來看申命記.
看看整個以色列民族.
如何經歷上帝的帶領.
他們生命裡面經歷高高低低.
他們如何透過上帝的引導.
逐步逐步踏進應許之地當中.
很多人去解釋應許之地.
或者申命記的整個段落.
有不同的解釋方法.
有些解釋方法是用歷史文化.
了解當時的人面對的處境.
以至在裡面看到歷史的教訓.
有些是用新約的角度.
來看上帝如何引導以色列人.
走在面前的路.

$^{121}$接下來我們用的方式.
大概將剛才兩個夾雜.
亦會加入一個靈性的角度.
我們將自己.
我們將我們每一個.
比喻為整個以色列民族.
上帝如何將我們這一班的回眾.
和我們當代的這一班人.
如何引導我們進入靈性的階段.
在我們的人生裡面.
我們面對著一些甚麼阻礙.
以致我們沒有辦法.
根據上帝吩咐我們做的事.
我們能夠一一實踐.
一邊說你便會覺得沒那麼抽象.
不過我們最重要的是.
我們從靈性角度.
來代入在以色列人當中.
當我們去講這個書卷的時候.
大家不妨想像一下.
我們就是當代裡面的以色列.
將要進入迦南地的一班回眾.
我們生命裡面需要拿甚麼教導.
我們真的可以踏進生命裡面.
有一些目標和上帝吩咐我們做的事.
我們能夠有力量來實踐.
有一個這樣的經歷.
我在神學院裡面.
我接受幾年的讀書過程裡面.
我其實主修一條stream.
我就是主修基督教教育.
我記得那個老師是我已故的老師.
黃錫賢博士.
因為我很喜歡上這一科.
很喜歡教導.
每一次做功課就寫好教案.
劃盡心思.
寫好不同的目標.
很有創意地去想每個課堂.
但是每一次功課的成績不是那麼好.

$^{161}$氣氛很不順.
於是我就去找黃老師.
我說為甚麼我那麼努力.
做到甚麼都做不到呢.
你看我識經識得多好.
我又有情感.
我知情義行的目標.
基本上我都做得到.
我就拿著那份功課和老師來談.
老師說我看到你很有心機.
解經識經的過程都很仔細做得很好.
情感也很有dynamic很有投放.
不過你覺得參與在課堂裡面的學生.
弟兄姊妹.
他們會不會帶來行為的改變呢.
我說做好這兩件事其實是足夠的.
他說你的人生是不是這樣.
是不是聽到一篇很蕩氣迴腸的道.
你就會改變呢.
是不是聽到一些很感人至深的話.
你就會有行動來回應呢.
其實你寫的教案也反映了你對牧羊的理念.
他說今個世代的人.
或者在我們這一代裡面.
最缺乏的不是知識也不是情感.
最缺乏的就是改換.
將這些東西轉化成為行動的意志.
當我們沒有了這份意志.
或者沒有了勇氣去實踐.
這些剛才所明白的真理.
或者情感打動你的真理.
到了最後沒有了這一步的時候.
你是無法化行動的.
今日我們看的申命記.
第一章已經在說一個這樣的主題.
如果我們從靈性角度去看.
這班以色列人.
他們到了最後.
用什麼來幫他們作出選擇呢.
他們到了最後能不能夠立起心智.

$^{201}$來進入迦南地呢.
今日這一篇第一個講章.
或者去到第一個段落裡面.
摩西就在說一個這樣的主題.
希望今日這一篇的訊息.
能夠幫助到我們.
能夠在新的一年展開新的展望.
我們先祈禱求主幫助.
父上大人願你在我們當中帶領我們.
讓我們能夠明白你的話語.
在我們生命當中.
我們知道透過自己的力量.
是無法達成任何你為我們所預備的吩咐.
求你在我們當中.
讓我們不單止聽得明白.
也激發我們的鬥志.
在未來的日子裡面.
我們要實踐你給我們的照明.
求你的話語.
你對我們的愛激勵我們.
幫助我們.
恭敬將來的時間交託.
幫助孫講得清楚.
聆聽的每一個聽得明白.
獻上禱告奉教基督耶穌得勝的聖名.
Amen.
我們看看新面記第一章.
我們看看第一節到第二節.
這裡我讀了兩節經文給大家聽.
以下所記的是摩西在約旦河東的曠野.
疏忽對面的阿拉巴就是巴蘭.
陀忽拉班哈西路.
在撒哈中間.
向以色列眾人所說的話.
從荷列山經過西伊到迦底斯巴尼亞.
有十一天的路程.
在新面記第一章一至二節裡面.
其實是很簡單很明快的.
是摩西一篇的講章.
申命記這個字.

$^{241}$有一些的說法.
就是如果你用希伯來文第一句.
第一個字來成為書卷的名字.
其實就是說話的意思.
不過這個說話是說.
重新將以色列人經歷了四十年的經歷.
以及在上帝裡面.
摩西舊時藉著律法的教導.
這班以色列人怎樣走呢.
將這些事跡重新說過一次.
摩西五經裡面特別對著百姓的講章就是申命記.
所以去到今天我們去翻譯.
或者我們去理解申命記.
其實不是說一些新意.
或者一些新的東西.
而是說將他們過去上一代的經驗.
重新去累積.
告訴他們.
上一代走過這些.
某程度上不是很理想的路.
他要勉勵新的一代.
你們要重新.
要引以為鑒.
以至你們生命可以前進.
生命就是重新述說律法一次.
這樣的意思.
所以其實說的不是甚麼新意.
不過這個是他們過去裡面實踐得不好的.
我給了一個地圖給大家看一下.
說這番說話我們看一下下一頁.
我們看的地圖.
看不到那個圈圈.
好,看一下地圖.
我們看到字比較小一點.
我想圈給大家看.
這個地方.
迦底斯.
這個就是摩西說這一番.
申命記第一章.
一至二節裡面.

$^{281}$講這個講章裡面的地點.
剛才講到第一節裡面.
講到有不同的地方.
那個地方的起點在哪裡呢?.
就在最南這裡.
西乃山.
逆河列山.
其實河列山和西乃山是同一個地方.
由最南一直走到這個地方.
本來如果按一般人腳的路程.
第二節裡面講的.
十一天就走到這裡.
不過整個以色列民族.
由這裡走到這裡.
而迦南地就是這裡.
迦南地去到這裡其實已經是臨門一腳.
原本就將要進入.
但是他們因為生命裡面的失敗.
以致進入不到.
這個十一天的路程.
成為了以色列人.
留落在這個地方裡面.
四十年的地方.
如果你看一些比較立體的地圖.
你更加知道.
這裡原地中海.
這些是一些很富庶.
很肥沃的地方.
這些你見到.
巴蘭曠野.
西乃曠野.
舒爾曠野.
全部的伊坦曠野.
全部是尋曠野.
見到全部都是一些荒蕪.
曠野之地.
內陸地方是辛苦的.
本來有一個流奶與蜜之地.
有一個等待他們要進入的美地.
他們進入不到.

$^{321}$他們反而要在曠野裡面.
由十一天的路程.
變成了四十年的光景.
在這個申命記裡面.
這裡再講到.
現在摩西已經去到.
剛才所講的迦底斯的地方.
四十年前來到這裡進不來.
然後他們就留落在這裡.
四十年後日子到了.
摩西再帶領以色列民族.
在這裡重新再講一次.
迦底斯這個地方.
告訴他們日子已經到了.
他們將要進入迦南美地.
所以下來的第一章.
就講到他們進入迦南美地之前.
首先要回顧在這個地方裡面發生的事.
所以你要留意.
你讀申命記的時候.
有些記載或講論.
不是完全順序.
因為他要抽取的.
就是按著地方.
就著那個地方的特性.
他就要將那個訊息告訴你.
而第十九節開始.
剛才我們所聽的那段經文.
講到以色列人.
去到迦底斯裡面失敗的經驗.
因為這個篇幅很長.
我分了幾個段落.
我們從這幾個段落裡面.
你會看到.
以色列民他們正在經歷.
生命裡面.
逐步逐步朝向離棄上帝.
或者失敗的那些道路.
第一個段落就是.
十九節到第二十五節.

$^{361}$我們照著耶和華我們神所吩咐的.
從荷列山起行.
經過你們所看見.
大而可怕的曠野.
往阿摩利人的山地去.
到了迦底斯巴尼亞.
剛才講的那個地圖.
我對你們說.
你們已經到了耶和華我們神.
所賜給我們的阿摩利人之山地.
看啊.
耶和華你的神已將那地擺在你面前.
你要照耶和華你列祖的神.
所說的上去得那地為業.
不要懼怕也不要驚慌.
Sorry啊.
你們就都近我來.
你們要先打發人去.
為我們窺探那地.
將我們上去該走的河道.
必進河城.
都回報我們.
這話我以為你.
就從你們中間選了十二個人.
每支派一人.
於是他們起身上山地去.
到以實國谷窺探那地.
他們手裡拿著那地的果子.
下來到我們那裡回報說.
耶和華我們的神所賜給我們的是美地.
這個第一個段落.
在這個段落裡面.
他講到其實.
以色列人去到迦底斯.
四十年前.
去到迦底斯這個地方.
其實他們某程度已經臨門一腳.
舊時是在說圍爐的地方.
在埃及地裡面.
他們沒有自己的地方.

$^{401}$沒有自己的身份.
到如今.
經歷神蹟之後.
摩西帶他們離開埃及.
然後走過一個很難走的曠野.
在這裡特別.
摩西在說的就是大而可怕的曠野.
那個地方.
如果你去過沙漠地曬一曬.
你都覺得很辛苦.
是吧?.
在說幾百萬人.
又老又亂.
而且沒有自己固定居所.
其實是一個極之痛苦的歷程.
不過.
當他們經歷過這樣的曠野之後.
上帝說已經將美地.
放在你們面前.
擺出來了.
你看到的就是曠野.
和美地之間有一個很強烈的對比.
原則上.
面對著這麼大的對比.
人的本性就是.
我們就要去那裡了.
不過.
在經文裡面說到.
摩西第一件事就是說.
你們不要懼怕也不要驚慌.
為什麼他們害怕呢?.
如果你退一萬步.
或者你再去回一回帶.
你會發覺.
如果你是那些以色列民你都會害怕的.
為什麼呢?.
在說的是之前的日子.
其實他們只不過是家裡的奴隸.
他們是做體力工夫的.
他們其實是沒有自主的成份的.

$^{441}$甚至很多都是那些婦孺百姓.
他們處於一個光景.
離開了.
埃及已經是一件很不得了的事情.
但是現在他們在說.
他們這班以前做奴隸的人.
現在有自己的屬地.
而且他們那個地方雖然很富庶.
但那個是陌生的地方.
那個陌生的地方.
他們要進入的時候還要跟那些人說.
上帝讓我們來拿.
Claim這個地方是屬於我們的.
其實這件事也是挺奇怪的.
再加上.
其實他們做這個過程.
只要被說這個地方是屬於我的.
就自動.
原則上是要打仗的.
是一班婦孺百姓.
以前是一班奴隸.
進入一個地方要打仗.
你想想這班人.
其實面對著一個.
他們完全沒有想像過的處境.
所以這件事對於他們來說.
上帝真的給了他們一個很艱鉅的任務.
他們的驚慌.
他們的恐懼.
其實是人之常情.
甚至乎是可以理解的.
對於一種這樣的恐懼.
或者一些不理解的狀況.
他們這班人在這一刻都很理性的.
他們在說什麼呢.
不如這樣吧.
我們進去就派一些探子進去.
窺探一下.
了解一下那個美地.
究竟是什麼的走道.

$^{481}$或者我們怎樣去拿那個地方.
我們都知道那個地方是怎樣的.
他們是要畫一些地圖的.
如果進去就潰敗成軍.
所以他們就派了12個探子進去.
去窺探那個地方.
所以整件事是非常理性.
也是非常合理的.
在經文裡面特別說到.
摩西聽到他們這樣的建議.
這句話是在說摩西.
或者說我以為美.
簡單來說就是.
摩西都非常同意他們這個想法.
所以就在他們當中選了12個人.
每支派派一人.
就是12個支派.
12個人.
起來就進去那個山地裡面.
以實覺谷.
窺探迦南地是一個什麼樣的地方.
簡單來說.
這12個探子進去做要做什麼呢.
他不是要看看那裡有多美.
那裡的地有多麼富庶.
他進去真的要說.
他們接下來要進攻那個地方的時候.
他們要拿到路線圖.
他們要知道如何進入迦南地.
這個是他們最重要的目標.
好了這12個探子回來了.
他們帶著什麼呢.
在聖經裡面說到.
帶了那地的果子.
流奶與物之地.
帶了一些水果回來.
很棒啊.
有很肥美的出產.
帶了這些東西來.
然後他們的結論就是說.

$^{521}$耶和華我們的神所賜給我們是美地.
你說這個探子.
他們做的事情.
在經文裡面說.
他們是不是稱職呢.
一半.
因為那個美地一早上帝已經說是美地.
他們離遠在那個礦業裡面看過去.
陸游游一遍都一定知道是美地.
美地這個是結論.
他們最重要的就是要知道那個地方.
探聽他們的虛實.
以及知道路線.
以致他們知道要怎麼做.
不過第一個段落裡面說到這個理性.
本來是幫他們實踐的理性.
就變為了一個辯解.
反而一個好的結論就帶來了負面的果效.
接著26至到第33節.
你們卻不肯上去.
竟違背了耶和華你們神的命令.
在帳篷內發怨言說.
耶和華因為恨我們.
所以將我們從埃及地領出來.
要交在阿摩利人的手裡面.
除滅我們.
我們上哪裡去呢.
我們的弟兄使我們的心消化.
說那地的民比我們又大又高.
城邑又廣大又堅固高得頂天.
並且我們在那裡看見阿納族的人.
我就對你們說不要驚恐.
也不要怕他們.
在你們前面行的耶和華.
你們的神必為你們爭戰.
正如他在埃及和曠野.
在你們眼前所行的一樣.
你們在曠野所行的路上.
也曾聽見耶和華你們的神撫養你們.
如同撫養兒子一般.

$^{561}$直等到你們來到我這地方.
你們在這事上卻不信耶和華你們的神.
他在路上在你們前面行.
為你們找安寧的地方.
夜間在火柱裡.
日間在雲柱裡指示你們當行的路.
26至33節.
這裡說的是什麼呢.
剛才說的那些探子窺探那個地方.
不過回來之後得出的結論.
這是一個美地的時候.
那些以色列人帶來的.
不是一個正面的回應.
他們帶來的就是他們不肯去.
違背上帝的命令.
在帳篷裡發怨言.
你想想多有趣.
剛才說的所有的條件.
所有的東西.
全部都是很正面的.
是這樣去的.
那個是美地.
上帝把這個地方賜給我們.
不過將這個好消息.
某程度上不知道怎麼做的好消息.
帶回來的時候.
那些人的恐懼感就發作了.
那些以色列人覺得.
不可以這樣.
我們不去了.
或者他們要拒絕上帝.
叫他們進入迦南地的命令.
他們甚至把探子帶來的結果.
和上帝的命令視為一個什麼呢?.
他們覺得是上帝要恨他們.
要將他們交在阿摩利人當中被除滅.
猶太的色經學家.
他用一個字眼.
可能在玩那些希伯來人.
很喜歡玩wordplay的字眼.

$^{601}$他在說那個恨字.
那個恨字叫做Sinah.
原文讀法叫Sinah.
而這個Sinah這個字.
和西乃山他們從這裡十一日路程的起點.
西乃山的西乃.
基本上差一點點.
一個是Sinah.
就是西乃.
恨戶Sinah就是Sinah.
讀音是差不多的.
不過就由西乃到恨戶.
差一隻字.
就是這隻字.
就好像那些以色列人對上帝的那種反應.
上帝在西乃山裡頒佈律法.
上帝給他們命令要進入迦南地.
這個一切都是美好的安排.
不過在以色列文當中.
他看這些上帝美好的計劃是Sinah.
是上帝憎我們.
是上帝要將我們交給那些外邦人.
要滅絕我們.
所以整個將正面的東西負面來解讀.
以致他們心裡面的那種負面情緒.
就不斷的來醞釀.
他們甚至乎用不同的方法.
來合理化他們的拒絕上帝.
他們說當時迦南地那些阿納族的人是巨人.
誠邑又廣大又堅固高到上天.
我們見到阿納族那些人好像蚱蜢一樣.
誓死無生.
上帝將我們說要交到美地.
其實上帝將我們交給那些阿摩尼人.
將我們送死.
其實這個段落.
他就在說那些以色列人.
就帶著一個這樣的心態.
來面對上帝的美意.
摩西在這裡裡面.

$^{641}$他要補充的.
他說你們說上帝憎恨你們.
你想想上帝怎樣對你們.
上帝一直為你們來去爭戰.
你們在埃及在曠野.
上帝都是走在你們的前面.
上帝帶領你們就好像小蚊子一樣.
無微不至.
雲處火處指引你們的道路.
在你們缺乏當中補足你們.
你們本來是一群什麼都沒有的人.
現在將要進入一個上帝預備的迦南美地.
你們竟然將這些美好的事視為負面的經歷.
似乎覺得上帝要糟蹋他們.
恨他們.
所以摩西在這裡就說到.
他要為上帝去辯解.
同時間在這裡裡面說到.
這群以色列人是一群不信耶和華你們的神的人.
在這裡裡面我們稍微要停一停想一想.
記不記得剛才說過那些以色列人.
是一群恐懼的人.
他們的處境面對著一個新的境地.
其實那種不安那種恐懼.
其實是人之常情的.
不過這群以色列人.
他們沒有在上帝面前呈現他們真正心靈的狀況.
如果他們直接說上帝我們很害怕.
進入一個這樣的地方.
你給我們一些確據給我們.
給我們一些信心.
我們就容易進入一點.
不過他們不是這樣看的.
他們第一樣看的就是.
他們的恐懼就將自己的情緒扭曲.
扭曲到一個地步.
就將正面的東西曲解成為負面的經歷.
以致他們和上帝的關係就越走越遠.
成為了一個不信的人.
在經文裡面特別更加要說.

$^{681}$進入迦南地不是純粹像你吃牛河.
或是吃乾炒肉片河的A餐或是B餐的選擇.
而是在說上帝的命令.
上帝的心意上帝的旨意.
要叫他們進入去.
他們怎樣都要進去.
這個既是一個命令也都是一個應許.
不過在過程裡面.
人就會覺得看到上帝辣的一面.
覺得是一個命令就帶著負面的情感.
來看這些一切美好的事.
產生了一種極之混亂扭曲的負面情緒.
我不知道頂展門你會否有時面對這種狀況.
在我們人生裡面上帝有很多美意.
有很多美善的經歷.
或者上帝美好的聖兆.
淋到我們的生命當中.
但我們很多時看這些東西.
是覺得上帝作弄我,糟蹋我.
或者我們用不同的方法.
用我們的理性來辯解.
拒絕上帝對我們每一個人生裡面美好的照明.
我年輕的時候聽過一個老牧者說.
這是一個培靈粉兄的牧者.
他說過一個令我印象深刻的例子.
他說他在教會裡面每年都講培靈會.
他見到台下有一個年紀比較大的成熟的人.
在培靈會裡面決志來奉獻.
在培靈會完結後.
這個牧者和年紀比較大的來談.
你都相當大了.
你打算將來的路是怎樣走.
但這個年長的人不停哭.
不停哭的時候.
他說每年都決志.
我都決意要跟隨著.
但每一年我帶來的都是遺憾.
這個培靈粉兄的牧者問.
為什麼你會這樣做呢?.
他問那個在哭的弟兄.

$^{721}$他說在我十幾歲的時候.
我聽培靈會已經知道我的生命要交給主.
但我的理性或我自己的經歷告訴我.
讀多幾年書成熟一點的時候.
我才會踏上.
讀了幾年書大學畢業的時候.
他又聽到上帝的召命的時候.
他覺得都是時間了.
不過會不會做多幾年事.
人生成熟一點的時候.
他才要將自己奉獻出來成為一個傳道人呢?.
於是又隔了幾年.
隔了幾年之後事業又回復正常.
家裡已經開始有兒有女.
呼召又再次來到他的生命當中.
他就說家庭的責任都是我召命的一部分.
會不會令我的兒女長大了之後.
我就去回應上帝的呼召呢?.
如是者人生就幾十年.
到他老了.
兒子都長大了.
他都進入了年紀很老邁的階段.
身體的魄力已經開始衰殘.
當日那個培靈奮興的訊息.
再一次叫他來呼召的時候.
他心裡很大的感動.
但是他慨嘆.
人生裡面錯過了最有力量的青春.
他將他的青春放在他覺得很重要的事業.
金錢.
甚至乎他覺得要美滿的家庭等等一切.
他將這些視為生命最重要.
就放棄了眼前這個召命.
生命裡面.
因為恐懼呼召帶來的改變.
以至到他要選擇次好的.
又或者拒絕上帝美好的呼召.
剛才所說的這個弟兄.
和以色列的人民有點不同.
這個弟兄都是帶著正面.

$^{761}$來看上帝的命令和召命.
但這裡的以色列人帶著負面.
來解讀上帝的旨意.
覺得上帝恨他們.
以至到他們的生命完全扭曲.
在進入迦南地的事上.
他們就敵擋上帝.
上帝面對著以色列人有什麼反應呢?.
34節.
耶和華聽見你們的話就發怒起誓說.
這惡世代的人連一個也不得見我.
豈世應許賜給你們列祖的美地.
惟有耶夫尼的兒子迦勒必得看見.
並且我要將他所踏過的地.
賜給他和他的子孫.
因為他專心跟從我.
耶和華為你的緣故也向我發怒說.
你必不得進入那地.
事後你亂的兒子約書也.
他必得進入那地.
你要勉勵他.
因為他要使以色列人承受那地違約.
並且你們的婦人孩子.
就像你們所說的.
必被擄掠的.
和今日不知善惡的兒女.
必進入那地.
我要將那地賜給他們.
他們必得違約.
至於你們要轉回.
從紅海的路往曠野去.
上帝面對著那群以色列人的那種.
那種不順.
那種負面的解讀.
上帝有什麼反應呢?.
耶和華發怒.
上帝和人之間是有互動的.
人的敵黨.
人的那種負面詮釋上帝的美意.
招惹上帝的憤怒.

$^{801}$上帝稱呼那一代那些人.
那些以色列眼前的那些.
由埃及領出來的那第一代.
他說那一代是惡世代的人.
那種惡不是在說.
他們道德上的惡.
而是在說那種完全扭曲.
將上帝美意.
或者心靈裡無法承受.
上帝美好的那一代人.
心生怨言的人.
成為了上帝眼中的惡世代.
所以上帝很清晰說.
這一代的人.
不能夠見到應許之地.
其實是一個很大的審判.
是一個很大的懲罰.
因為其實剛才那個地圖你看到.
多走幾步.
幾天的路程就到了.
但是就因為他們的不信.
因為他們那種扭曲的那種負面情感.
使他們無法進一步進入應許之地.
換來的是什麼呢?.
上帝第一次用你們要轉回.
轉回這個字很多時候說是悔改.
不過這個字今天說不是.
走回曠野的路.
由臨門一腳將要進入迦南美地.
現在你回頭去曠野.
他們原本在民數記裡說.
他們甚至想回到埃及.
不過上帝說你回到曠野.
繼續經歷那種歷練.
不是回到那裡.
是回到40年.
人生的不信.
惹來上帝的憤怒.
在我們人生裡.
我們有沒有經歷過兜兜轉轉.

$^{841}$我們會不會有時候.
我們執意而行.
又或者拒絕上帝.
負面地看上帝的美意.
以致我們被上帝很深刻地.
來管教我們的生命.
時間關係我們再看下一個段落.
41-46節.
這是一個更加令人嘆息的段落.
那時你們回答我說.
我們得罪了耶和華.
情願耶和華我們神一切所吩咐的上去征戰.
於是你們各人帶著兵器.
爭先上山地去.
耶和華吩咐我說.
你們對他們說不要上去也不要征戰.
因為我不在你們中間.
恐怕你們被仇的.
殺敗了.
我就告訴了你們.
你們卻不聽從.
竟違背了耶和華的命令擅自上山地去了.
住那山地的阿摩利人就出來攻擊你們.
追趕你們.
如蜂擁一般.
再篩以殺退你們.
直到荷爾瑪.
你們便回來在耶和華面前哭號.
耶和華卻不聽你們的聲音.
也不向你們質疑.
於是你們在迦底斯住了許多日子.
這是某程度上.
以色列人在第一個講章中面對的結果.
他們的結果是甚麼呢.
當上帝很憤怒叫他們回到曠野.
講得直白一點.
以色列人就知死.
知死的.
然後做一些更加極致.
更加敵黨的事.

$^{881}$現在我們上去打仗了.
知道要回到曠野更慘.
於是憑著匹夫之勇.
現在就去打.
現在就去打阿摩利人了.
於是就拿兵器爭先.
這裡面講排第一.
大家要盡快爭.
我們要進去迦南地.
爭先恐後地上山去打仗.
之前怎麼踢都不動.
但到現在知道有後果了.
就準備去了.
上帝怎樣跟他們說呢.
上帝直著摩西說.
叫他們不要去.
不要打.
不用打.
我不在你們中間.
你們一定會打輸.
上帝不在這群以色列人當中.
上帝的靈不跟他們同在.
如果你硬要上去.
一定會被仇敵殺敗.
這群以色列人聽到摩西的這番說話.
有甚麼反應呢.
繼續去.
他們卻不聽從.
竟違背耶和華的命令.
擅自上山地去.
他們再一次不聽上帝的吩咐.
再一次的來犧牲.
於是阿摩利人就出來攻擊他們.
追趕他們.
然後就殺敗他們.
他們在上帝面前哭號.
經文講到上帝好像很無情.
上帝不聽你們的聲音.
不向你們的質疑.
其實真正不聽的是那些以色列人.

$^{921}$上帝已經將結局告訴他們.
不過那些以色列民.
視而不見.
聽而不聞.
繼續硬心過他們自己的生活.
結果他們就流落在曠野裡.
新年開始用一個不太老套的經文.
好像很負面.
不過這個負面其實是一個很重要的歷史.
或者對以色列人是一個很重要的一課.
我們的人生會不會有時覺得.
失去了那個時機.
然後我們就要做一些事.
強行做一些事來做補救的功夫.
或者要強行憑血氣.
憑著那種心裡面的衝動.
來完成.
當我們很想去做的時候.
上帝說我不跟你同在.
你不用強行去.
不要這樣去.
上帝不是要挫敗那些以色列人.
只不過是要那些以色列人學懂這個功課.
上帝的說話要聆聽.
他也要搞清楚他們自己心靈的狀況.
當我們去讀以色列人的歷史.
如果用回靈性的角度.
會不會有時上帝也有一些吩咐叫我們去做.
我們去拒絕.
然後到最後我們經歷過一些功課之後.
我們就強行逼自己去做.
不是重新重整過自己.
聽下上帝想我們去怎樣.
而不是盲目去實踐某一些的照明.
今日的經文似乎是一個很負面的經歷.
但對於我們今天來說.
我們可以正面的去看.
我們正面的去看.
今天有沒有一些東西難了我們進入一個迦南美地.
上帝有一些吩咐叫我們去做的時候.

$^{961}$我們有沒有真是聽清楚.
或者我們有沒有那種鬥志來實踐.
抑或是我們心裡面太多顧慮.
我們有太多的恐懼.
以致我們沒有辦法踏前.
有這些顧慮.
有這些恐懼是不要緊的.
最重要的是你願不願意將你真實的一面交給上帝.
不要將你的恐懼.
或者將你的憂慮壓低了.
然後就負面詮釋上帝的美意.
我可以說.
今日很多基督徒其實都落入一種這樣的狀況.
用一個直率的說.
有時候有些弟兄姊妹無端端生氣.
生氣教會.
生氣牧者.
我們都不知道發生什麼事.
但再談.
其實都不是我們做了什麼.
而是可能在生命某些階段裡面.
遇到一些恐懼.
在成長裡面遇到一些關口.
心靈很糾結.
很混淆.
於是他有一種負面情緒.
他都不知道怎樣面對的時候.
當有一個牧者.
當有教會.
有一些說話呈現出來的時候.
就好像刺了他一樣.
其實那個教會的信息.
或者牧者.
其實都不是特意要針對的.
正正是這些說話可能對自己有一種威脅感.
以至產生了很大量的負面情緒.
我們被這些負面情緒籠罩的時候.
我們生命就失去了方寸.
我們不知道怎樣走.
我們只知道每一天帶著負面的感受.

$^{1001}$來過我們自己的生活.
如果以色列人去正視他生命的恐懼.
祈求上帝幫他疏解.
不再將這些負面情緒累積的時候.
我相信迦南地就成為了他下一步要進入的目標.
而不是再回到曠野當中.
同樣我們的人生都是.
我們正視我們心靈裡面的真實狀況.
將我們裡面的那種很真實的感受.
情緒恐懼來呈現出來的時候.
弟兄姊妹,教會,牧者是很願意陪你一起來同行.
如果我們捱起爭頭.
我們壓低自己的那種感受.
帶來的就是生命裡面更加大的混亂.
經文裡面更加講到.
這班以色列人其實進入迦南地很需要勇氣.
這份勇氣是他們過去裡面沒有經歷過.
同樣我們今天的信仰都需要這一份勇氣.
一份和這個世界和自己過去不同的那份勇氣.
失去了這份勇氣.
你就沒有辦法有力地實踐信仰.
我很喜歡C.S. Lewis講的這一句.
他說勇氣是什麼呢?.
他說勇氣是每一個美德.
不是其中一個的美德.
而是在最關鍵因素裡面實踐.
所有美德裡面最重要的那一環.
簡單來說就是如果你沒有勇氣.
你生命裡面沒有辦法有力量去實踐生命其他所有的美德.
如果你要愛,你沒有勇氣你愛不到.
你要忍耐,你沒有勇氣你沒有辦法忍耐.
你要服事,你沒有勇氣你沒有辦法服事.
林林總總,勇氣成為了我們這個世代裡面最重要的一種美德.
昔日,以色列人失去了這份勇氣.
他用他的理性來包裝自己的怯懦.
以致他們沒有辦法進入新的一年.
弟兄姊妹,我們每年都為自己立志.
我們有沒有勇氣來好好實踐.
讓上帝的美意在我們生命當中來成就.
上帝進入迦南地.

$^{1041}$你不要覺得剛才說過不是A餐和B餐的選擇.
而是一個命令和應許.
上帝應許要把迦南地賜給你和我.
上帝有一些美好的旨意成就在你和我身上.
你怎樣避開也好,上帝都會怎樣管教你,帶你進入去.
你想痛苦地進入,還是榮耀地進入呢?.
今日這段經文提醒我們.
我們正視心靈真實的狀況.
有勇氣地面對我們的人生.
迦南地就在我們前面,我們一齊進發,我們祈禱.
求主幫助我們,讓我們藉著申命記第一篇的訊息.
提醒我們,我們生命不要再在曠野裡兜兜轉轉.
幫我們去正視心靈裡的負面恐懼.
幫我們能夠好好在這個世代裡.
活出屬神而屢的那種風榮.
天父上帝,在新的一年裡我們立志.
我們更加求你的聖靈,剛強我們的內心.
耶和華上帝在我們前頭引領.
讓我們的心靈能夠有耶穌基督的同在.
踏實地走我們人生.
求你與我們每一個弟兄姊妹同在.
引領我們,說領我們走面前人生的路.
願上禱告奉教基督耶穌得勝的聖名.
Amen.
\newpage



\section{約翰福音 7:32-52}
\label{sec:YgvzlIV1PPE}
\textbf{2025年1月12日崇拜直播｜梁家麟牧師｜令人不安的耶穌(二)｜約翰福音7\_32-52}
\newline
\newline
連結: \href{https://youtube.com/watch?v=YgvzlIV1PPE}{\texttt{https://youtube.com/watch?v=YgvzlIV1PPE}} ~~~~ 語音日期: 2025-01-12
\newline
\newline
\hyperref[sec:H2Yyyhsu_Qc]{\small{< < < PREV SERMON < < <}}
~
\hyperref[sec:index_chronic]{\small{[返順時目]}}
~
\hyperref[sec:uiuY8tue_2Q]{\small{> > > NEXT SERMON > > >}}
\newline
\newline
約翰福音 7:32-52
\newline
\begin{longtable}{cl}
\hline
\hline
章節 & 經文 (和合本修訂版)\\
\hline
7:32 & \begin{tabularx}{0.7\textwidth}{X} 法利賽人聽見群眾對耶穌這樣議論紛紛，祭司長和法利賽人就打發聖殿警衛去捉拿他。 \end{tabularx} \\ \\ \relax
7:33 & \begin{tabularx}{0.7\textwidth}{X} 於是耶穌說：「我跟你們在一起的時候不會太久了，我要回到那差我來的那裡去。 \end{tabularx} \\ \\ \relax
7:34 & \begin{tabularx}{0.7\textwidth}{X} 你們要找我，卻找不到；我所在的地方，你們不能去。」 \end{tabularx} \\ \\ \relax
7:35 & \begin{tabularx}{0.7\textwidth}{X} 於是猶太人彼此問：「這人要往哪裡去，使我們找不到他呢？難道他要往散居在希臘的猶太人那裡去教導希臘人嗎？ \end{tabularx} \\ \\ \relax
7:36 & \begin{tabularx}{0.7\textwidth}{X} 他說『你們要找我，卻找不到；我所在的地方，你們不能去』，這話是甚麼意思呢？」 \end{tabularx} \\ \\ \relax
7:37 & \begin{tabularx}{0.7\textwidth}{X} 節期的最後一天，就是最隆重的一天，耶穌站著，喊著說：「人若渴了，到我這裡來喝！ \end{tabularx} \\ \\ \relax
7:38 & \begin{tabularx}{0.7\textwidth}{X} 信我的人，就如經上所說：『從他腹中將流出活水的江河來。』」 \end{tabularx} \\ \\ \relax
7:39 & \begin{tabularx}{0.7\textwidth}{X} 耶穌這話是指信他的人要受聖靈說的；那時還沒有賜下聖靈，因為耶穌還沒有得到榮耀。 \end{tabularx} \\ \\ \relax
7:40 & \begin{tabularx}{0.7\textwidth}{X} 眾人聽見這些話，有的說：「這真是那先知。」 \end{tabularx} \\ \\ \relax
7:41 & \begin{tabularx}{0.7\textwidth}{X} 另有的說：「這是基督。」但也有的說：「難道基督是出自加利利嗎？ \end{tabularx} \\ \\ \relax
7:42 & \begin{tabularx}{0.7\textwidth}{X} 經上不是說『基督是大衛的後裔，出自大衛的本鄉伯利恆』嗎？」 \end{tabularx} \\ \\ \relax
7:43 & \begin{tabularx}{0.7\textwidth}{X} 於是眾人因耶穌而分裂了。 \end{tabularx} \\ \\ \relax
7:44 & \begin{tabularx}{0.7\textwidth}{X} 其中有人要捉拿他，只是沒有人下手。 \end{tabularx} \\ \\ \relax
7:45 & \begin{tabularx}{0.7\textwidth}{X} 警衛們回到祭司長和法利賽人那裡。他們對警衛說：「你們為甚麼沒有帶他來呢？」 \end{tabularx} \\ \\ \relax
7:46 & \begin{tabularx}{0.7\textwidth}{X} 警衛回答：「從來沒有像他這樣說話的！」 \end{tabularx} \\ \\ \relax
7:47 & \begin{tabularx}{0.7\textwidth}{X} 於是法利賽人說：「你們也受了迷惑嗎？ \end{tabularx} \\ \\ \relax
7:48 & \begin{tabularx}{0.7\textwidth}{X} 難道官長或法利賽人中有信他的嗎？ \end{tabularx} \\ \\ \relax
7:49 & \begin{tabularx}{0.7\textwidth}{X} 但這些不明白律法的眾人是被詛咒的！」 \end{tabularx} \\ \\ \relax
7:50 & \begin{tabularx}{0.7\textwidth}{X} 其中有尼哥德慕，就是從前去見過耶穌的，對他們說： \end{tabularx} \\ \\ \relax
7:51 & \begin{tabularx}{0.7\textwidth}{X} 「不先聽本人的口供，查明他所做的事，難道我們的律法還定他的罪嗎？」 \end{tabularx} \\ \\ \relax
7:52 & \begin{tabularx}{0.7\textwidth}{X} 他們回答他說：「你也是出自加利利嗎？你去查考就知道，加利利是不出先知的。」 \end{tabularx} \\ \\
[1ex]
\hline
\hline
\end{longtable}
$^{1}$今日要讀的經份是記載在.
約翰福音第七章三十二至五十二節.
聖經新約第一百三十六至一百三十七頁.
約翰福音第七章第三十二節.
法利賽人聽見眾人為耶穌這樣紛紛議論.
祭司長和法利賽人就打發差役去捉拿他.
於是耶穌說.
我還有不多的時候和你們同在.
以後就回到差我來的那裡去.
你們要找我卻找不著.
我所在的地方你們不能到.
猶太人就彼此對問說.
這人要往哪裡去叫我們找不著呢.
難道他要往散著希利利中的猶太人那裡去.
教訓希利利散著的人嗎.
他說你們要找我卻找不著.
我所在的地方你們不能到.
這話是什麼意思呢.
節期的末日就是最大之日.
耶穌站著高聲說.
人若渴了可以到我這裡來渴.
信我的人就如經上所說.
從他腹中要流出活水的江河來.
耶穌說是指著信他之人要受聖靈說的.
那時還沒有賜下聖靈來.
因為耶穌尚未得著榮耀.
眾人聽見者說.
有的說:這真是那先知.
有的說:這是基督.
但也有的說:基督豈是從加利利出來的嗎.
經上豈不是說:基督是大衛的後裔.
從大衛本鄉伯利恆出來的嗎.
於是眾人因著耶穌起了紛爭.
其中有人要捉拿他.
只是無人下手.
差役回到祭司長和法理賽人那裡.
他們對差役說:你們為什麼沒有帶他來呢.
差役回答說:.
從來沒有像他這樣說話的.
法理賽人說:你們也受了迷惑嗎.

$^{41}$官長說:法理賽人豈有信他呢.
但這些不明白律法的百姓是被咒了的.
內宗有尼波底姆.
就是從前去見耶穌的.
對他們說:不先聽本人的口供.
不知道他所作的事.
難道我們的律法還定他的罪嗎.
他們回答說:你也是出於加利利嗎.
你且去查考.
就可知道加利利沒有出過先知.
我們接著繼續說.
耶穌在主彭節.
在耶路撒冷聖典裡.
引起眾多不安騷亂的故事.
前面已經說過.
耶穌與不同的人群接二連三.
發生很多衝突和爭議.
合共有四波.
上次我們說了兩波.
今天我們說其餘的兩波.
第三波的爭議是緊接第二波爭議之後.
同樣是在聖殿裡的爭議.
當猶太教的最高決策層.
祭司長和民事聽到耶穌說的話.
引起各種爭議不安之後.
他們決定採取行動.
派出衙差(聖殿衛隊)去捉拿耶穌.
那些衙差去到聖殿的時候.
卻聽到耶穌向在場的群眾宣稱.
我還要不多的時候你們同在.
以後就回到差我來的那裡去.
你們要找我找不著.
我所在的地方你們不能去.
33到34節.
耶穌不僅宣告他從天父那裡來.
他也宣告他要回到天父那裡.
我所在的地方你們不能到.
我知道這句話的意思很簡單.
當耶穌從死裡復活之後.
就回到天上坐在聖父的右邊.

$^{81}$這個階段是沒有人能夠跟他去的.
即使是他的十二使徒.
也只能夠在人間.
在陰間等候.
直到耶穌基督第二次回來的時候.
然後才可以跟耶穌重聚.
群眾不明白耶穌說.
人們無法尋著他的意思.
以為他準備去外地宣教.
這是在主彭節的最後一天.
節期的第七天.
也是整個節期的高潮所在.
耶穌在聖殿裡面說了一篇很重要的信息.
人若渴了.
可以到我這裡來.
信我的人就如經上所說.
從他腹中要流出活水的江河來.
37,38節.
這句話耶穌曾經在雅各井旁.
跟那個婦人說過.
你記得吧.
喝她賜的水就流出活水江河.
主彭節起源是紀念.
以色列民離開埃及之後.
在曠野漂流.
住在象棚40年之久.
期間一直有上帝的供應.
由於主彭節定在我們的西曆.
即主曆每年的9月,10月.
這是降秋雨的時候.
即秋耕開始的時候.
因此節期逐漸就變成跟農業拉上關係.
祈求上帝降雨.
滋潤大地.
好使來年豐收.
耶穌在這個節期宣稱.
祂是水的供應者.
是很應節的.
一看為耶穌這番說話做了一個註解.
從腹中流出活水江河.

$^{121}$是指接受聖靈.
我們必須說.
這是日後教會對耶穌的話的理解.
耶穌是說人要接受聖靈.
不過我自己很保留.
我有理由懷疑.
耶穌怎會在這個時候.
在聖殿對著一群還未跟隨的人.
無端端說要接受聖靈.
連耶穌都未接受.
接受什麼聖靈呢?.
對不對?.
耶穌應該在離開人間前不久.
對門徒預告.
他會賜下另一個保衛師.
來代替他的位置.
那時候才說到聖靈.
在他復活升天前.
他再次提到聖靈賜下.
所以沒什麼理由.
耶穌是對著一群還未信的人.
在聖殿裡無端端說到聖靈.
這純粹是作者約翰.
在日後的教會的詮釋.
事後孔明的解釋.
耶穌公開呼籲人相信他.
並且說凡是信他的人.
都會得到真正並完全的滿足.
耶穌這番說話引起現場極大的爭議.
再次出現兩個對立的陣營.
有支持耶穌的.
相信他就是先知以利亞的復活.
也有更相信他是基督的.
不過也有反對耶穌是基督的.
他們提出的一個反對的理由是.
按照舊約先知的預言.
基督應該是出身自大衛的本鄉伯利恆.
而不是加利利.
他們不知道耶穌童年.
曾經有一個很大的轉折.

$^{161}$所以由南部的伯利恆.
遷到北部的加利利.
現在我們來到第四波的爭議.
這是在猶太教決策層內部出現的.
奉派前去捉拿耶穌的衙差.
因為聖殿現場出現激烈的爭論.
群情洶湧.
不敢貿然行動.
只能空手而回.
他們回去見祭司長.
祭司長就問他們.
為什麼沒有把耶穌帶來.
他們就回覆說.
46節從來沒有像他這樣說話.
他們在現場聽了耶穌的講道.
覺得有些猶豫.
不確定是不是要抓這個人呢.
或者他被耶穌的氣場說服了.
太厲害了.
猶太人聽到這班衙差這樣說.
就很生氣.
法利賽人說.
凡是有些知識的.
有些見地的官長和法利賽人.
都不會相信.
只有那些不明白律法的百姓.
才會受這個耶穌的迷惑.
這個耶穌是受咒詛的.
他們中間一個法利賽人.
尼哥底姆.
就是從前曾經見過耶穌的.
他對耶穌有好感.
就為耶穌說話.
51節.
你不先聽本人的口供.
不知道他所做的是什麼.
難道我們的律法就能定他的罪嗎.
你怎樣也要讓他答辯一下.
不要妄下結論.
不過反對耶穌的人.

$^{201}$就只堅持一句.
一點.
52節.
你也是出於加利利摩.
你且去查考.
你便可知道.
加利利沒有出過先知.
意思就是.
就耶穌在加利利出生的事實.
就可以將他判死刑.
不需要其他證據.
不需要酒情罪.
讓他答辯.
不用.
必須要注意.
這裡說的先知.
英文是有定貫詞的.
希臘文也有定貫詞的.
definite article(定義文).
英文是the prophet(神父).
即是那個先知.
意思就是彌賽亞.
舊約從來沒有經文預言.
彌賽亞會在加利利出生.
但不是說.
加利利真的從來沒有出現過先知.
先知約拿和先知拿洪.
都是出自加利利人.
是加利利人.
祭司長和法利賽人.
沒理由這樣也搞錯.
這樣也搞錯就很糟糕.
這是基本常識.
無論如何.
不是加利利沒有出過先知.
只是加利利沒有出過彌賽亞這個先知.
所以.
耶穌不可能是基督.
他今天被人誤會他是基督.
扮基督樣.

$^{241}$就是妖言惑眾.
罪無可恕.
對祭司長和法利賽人而言.
耶穌不應該安息日治病.
耶穌不應該說語重不同.
譁眾取寵的話.
造成社會騷亂.
耶穌不應該在加利利出生.
這些情節都構成他們拒絕耶穌的充分理由.
至於耶穌說了什麼.
行了什麼神蹟.
都變得不重要.
他們心中有對上帝一個很牢固的看法.
上帝必須是這樣.
必須是這樣.
不可以有其他樣.
不可以離開他們對上帝的認知.
凡是和他們的期望.
他們所想的不同的.
通通不可能是上帝.
簡單地說.
他們已經知道上帝是怎樣.
真實的上帝是不可以逾越他們的知識.
他們的上帝知識大於真實的上帝.
上帝不可以有別的樣貌出現.
我們總是對上帝有很多預設的想法.
上帝應該是這樣.
上帝應該做這樣的事.
上帝應該這樣善待我.
上帝應該擊殺我所憎恨的人.
祭司長這些宗教權威.
自以為很了解上帝.
很知道彌賽亞.
但他們執著對彌賽亞出身的詮釋.
拒絕眼前真真實實的彌賽亞.
我們對上帝既有的想法.
竟然大於眼前的上帝.
大於上帝自己.
耶穌,你未符合我所理解的彌賽亞的標準.
對不起上帝,我要拒絕相信你.

$^{281}$因為你不是我心目中所要的上帝.
看國內的YouTube上親節目.
不知道你有沒有看過.
很好笑的.
一個女士開出條件.
我要的對象四十歲以下未婚.
身高一百八十,年薪一百萬以上.
市區有一套已經付清房價的心房住宅.
房要寫我的名字.
開的車要八十萬以上.
每月要給女士兩萬元零用錢.
一張很長很長的清單.
唯有符合這堆條件的才是我理想的結婚對象.
我想這個女士是希望找一個心目中的米飯班主.
多於想找一個廝守終身的伴侶.
找對象這樣找都有的.
同樣很多人不是想找上帝.
只是想找他們心目中的上帝.
或者和黃大仙和聖誕老人差不多.
問題是我們都未見過上帝不認識他.
怎可以定出上帝的樣子呢.
上帝一定要符合我們心中的條件.
是不是有些狂妄呢.
和大家講一部戲.
我上個月看了一部戲叫做《教宗選戰》.
《Cave》有沒有人看過.
大陸也有秘密會議.
一部對我來說都挺震撼的戲.
那部戲講天主教的教廷.
藉著書記團或者紅衣主教團.
College of Cardinals去選任新教宗的故事.
整部戲帶出的訊息是一個.
確定是信仰的最大敵人.
也是教會合一的最大敵人.
Certainty is the enemy of faith and church unity.
Certainty 確定是信仰的敵人.
是教會合一的敵人.
這個訊息非常有深度.
也非常切合事宜.
首先我要講.

$^{321}$真實的信仰總是無法百分百確定的.
信仰總是伴隨著懷疑甚至是不信.
所以Believe and Unbelieve是雙胞胎.
在什麼都確定的情況下.
你是不需要講信心的.
你要不要信一加一等於二.
你要不要相信孫中山是人.
不用信的.
因為是事實.
又沒有信.
像保羅說的.
我們的基督教是在乎盼望.
只是所見的盼望不是盼望.
誰說盼望是他所見的呢.
但我們若盼望所不見的.
就必須忍耐等候.
看到了還有沒有盼望.
你盼望聽梁家倫講到.
有沒有盼望.
你現在盼望他快點下台.
因為他還沒有下台.
顛著的事有沒有盼望.
你是未到的事才是盼望.
講信心一定是不確定.
值得懷疑.
有理由不信的元素.
並且有趣的是.
往往是和現實越不相符.
越和外在的證據相衝突的宣告.
才是越需要我們有信心.
信心的大與小.
不是我們自己講.
我們自己的信念的強度.
很堅定很強.
或我所做的宣稱的力度.
很大聲叫.
是決定的.
不是.
而是取決於我們所講的事.
和客觀的現實有多大距離.

$^{361}$當親眼看到復活的主.
站在我們面前的時候.
我們講我們信他.
不需要什麼信心.
你見到了.
但當你目睹耶穌掛在木頭上的時候.
我們講我們仍然信他.
跟隨他.
哇!這個信心就很大.
耶穌對多瑪是這麼說的.
你一看見了我才信.
沒有看見就信的就有福了.
未看到的時候信.
是比看到才信重要.
信心的大與小.
和我們持守的信念有多強固.
沒什麼關係.
否則精神病患者.
就是最有信心的人.
我總是認定身邊有人在謀害我.
不需要你怎麼解釋.
我都認定有人在謀害我.
你拿什麼證據.
我都不同意有人在謀害我.
這樣是不是很有信心.
這不是信心.
這是精神病.
這是妄想症.
信心和自己信念的強弱無關.
所以聖經也不是要求我們要有很大的信心.
耶穌說你的信心只要有芥菜種這麼大就夠了.
一點點就夠了.
信心的大與小.
和個人信念的強弱無關.
是從外在的環境.
其他人的態度有關.
外在的環境越是和我們的宣告有差距.
大量的證據要迫我們懷疑.
甚至放棄我們的信仰.
或者我們身邊多數人.

$^{401}$都不是這樣選擇的.
我們自己孤身獨行.
自己做這個選擇.
這樣就需要大信心了.
不是在於你裡面有多強.
而是在於和外在的環境有多少差距.
正如我以前說過.
當人人都一窩蜂跟耶穌的時候.
跟隨耶穌是不用信心的.
你隨眾而已.
跟風而已.
當人人都捨他而去的時候.
我對他的不離不棄宣告.
你有永生之道.
我們還跟隨誰呢.
這樣就需要很大的信心了.
什麼叫信心呢.
就是你看到眼前的環境.
跟你所說的完全不一致.
雖然無花果樹不發黃.
葡萄樹不結果.
但我最後仍然一夜嘩嘩歡呼.
因我的救主喜樂.
這個需要信心嗎.
這個需要極大的信心.
是不是.
醫生宣告我沒得救.
我仍然深信靠主我可以得醫治.
這個需要很大的信心.
信心是跟外在的環境有多大落差有關.
不是在於你自己的信念有多強固有關.
所以在這樣的情況下.
我們永遠是在一個充滿懷疑.
叫我們不信的環境裡相信.
我們在懷疑中相信.
在不確定中相信.
其中叫我們敢於去做.
跟現實相對立的抉擇.
不在於你夠頑固.
夠堅強.

$^{441}$不在於你的信念夠強.
而端在於你對耶穌我們信仰的對象的信任.
我再說一次.
Prestige這個字是指信任.
Trust.
我信任他.
不是指你持守一套教條.
一套主義一套理論.
是信任.
對你跟隨的對象的信任.
信仰是信任耶穌基督.
接受他為我們成就的贖罪請求.
有期盼他按應許會再接我們上天堂.
做基督徒是信耶穌.
信任耶穌.
耶穌呼籲人來接受福音.
祂對人的邀請是什麼.
信我.
耶穌要求我們信祂.
我們很熟悉的經文.
復活在我生命也在我.
信我的人.
雖然死了也必復活.
凡活著信我的人必永遠不死.
你試試數數信我這個詞.
在耶穌口中出現了多少次.
你自己回去做這個功課.
信我.
信任一個對象.
不是在於你充分了解他.
而是你放棄對對方和自己充分掌握.
我去外國.
坐別人的車.
我信任他會載我去目的地.
我根本不用問他走什麼路線.
反正我也不知道.
專心欣賞沿途風光.
除了信仰總是夾雜著懷疑.
不可能百分百確定之外.
我們對我們所信的對象.

$^{481}$也不可能有百分百的認識和掌握.
我們信的活著的上帝.
不是一套宗教理論.
所以我們永遠不能說.
我們已經充分掌握完全明白.
你可以完全明白一個理論.
但你無法明白一個活著的上帝.
這個上帝是一個主題.
不是一個物體.
你只能安靜在他面前.
讓他告訴你他的想法.
他的做法.
確定是信仰的敵人.
應該說在宗教信仰上.
持百分百確定態度的人.
極可能是上帝的敵人.
一個人既然自以為.
已經充分認識上帝.
掌握了他的心意.
摸清他的行事規則.
他還需要繼續認識上帝.
尋求他的旨意.
期望上帝向他顯現.
向他說話嗎?.
不需要了.
God is fully known.
都知道了.
我用什麼再問他?.
他自己可以扮演上帝的角色.
Pray God.
宣告上帝是怎樣怎樣怎樣.
替天行道.
他自己就是上帝的代言人.
代理者.
他的宗教系統已經不需要上帝.
也容不下真實的上帝.
我們在上帝的面前.
必須要發現自身的罪惡.
承認自己是罪人.
繼續尋求上帝的赦免和幫助.

$^{521}$明白自覺.
只有病人才需要醫生.
耶穌是來召聚人悔改.
我們如果不讓他洗腳.
就與他無分.
同樣的.
我們在上帝面前.
也要承認自己的無知愚昧.
否則我們很難接受十字架的道理.
這個十字架的道理.
是要廢掉人的聰明.
我們很熟悉的.
《箴言三章》的教導.
你要專心仰賴耶和華.
不可以依靠自己的聰明.
在一切所行的事上都要認定他.
他必指引你的路.
不要只為有智慧.
要敬畏耶和華.
遠離惡事.
這個說法我以前說過.
其實很簡單.
耶穌要我們變回小孩子.
並且說.
唯有小孩子才能變上天國.
我講過.
小孩子不是天真無邪的意思.
小孩子.
應該翻成廣東話是「靚仔」.
在上帝面前千萬不要裝大腦.
不要二文二武.
唯有你承認自己是「靚仔」.
你然後才可以見到他.
聽到他.
我從來不在他面前指手畫腳.
說我懂什麼就懂什麼.
閒一點.
上帝是獨行其事的上帝.
上帝的意念非同人的意念.
上帝的旨意是人永遠無法察透的.

$^{561}$所以在信仰的道路上.
總是有驚嚇有驚喜.
不斷有新的啟迪和發現.
Faith is adventure.
信仰是一個力鉗.
我們一生都在學習中.
也一直有成長的空間.
我今年剛剛順主五十週年.
1975年決志.
這五十年.
你告訴我.
我《約翰福音》看過多少次.
超過五十次.
一定超過五十次.
我事奉將近四十年.
一年講一百篇.
一定不誇張.
一個星期最少兩篇.
週間都要講.
我講了四千篇.
如果假裝老練.
左執右執.
可以湊夠料都可以這樣一餐.
我太太可以作證.
我每篇道最少用五十到七十個小時去預備.
不要當自己懂.
我常說.
如果你還是用紙本的聖經.
如果這本聖經硬到花了.
就丟掉或放著.
找一本新的.
當然手機聖經很簡單.
不要反錯看你以前看過的東西.
每次認認真真去問上帝.
今天你就《約翰福音》三章十六節.
對我有什麼說話.
你對我們教會有什麼要求.
有什麼提醒.
每早晨都是新的.
不要讓你過去舊有.

$^{601}$你以為全都懂的知識.
將你困住.
以致你聽不到他跟你講新的話.
我常說.
老練很多時候是信仰的敵人.
祈禱.
我和你分享過不止一次.
幾年前的學習.
就是盡量祈禱祈得多一點.
不要太信口漏.
一開口就可以創世記祈禱啟示錄.
多少經文可以拋出來.
不要祈禱那些.
每句說話認真一點.
想清楚說.
就算說得多一點.
你覺得做牧師祈禱那麼差.
不要緊.
不用專業的.
因為我面對一個真實的上帝.
在他面前說話很難不口吃的.
很難不口吃的.
學會打回原形.
回到最基本的.
確信也是教會合一最大的障礙.
一個宣稱掌握了聖經通盤真理.
掌握了上帝所設定的屬靈規律的人.
總會宣稱他的認知是唯一壟斷性排他.
此次一家別無分店.
所有不同於我的理解和選擇的都是錯的.
都是不可接受的.
持有和我的想法不一樣的人.
顯然不僅是錯.
他還是別有用心.
志在混亂真理.
混亂教會.
獨斷的人只是期望有追隨者.
不能期望有合作夥伴.
由於真理是唯一的.
不可能有兩個真理.

$^{641}$那些有異於我的想法的人.
一定就是異端.
這是很多時候教會的衝突所在.
我們將意見變成真理.
真理和意見混亂了.
無論如何.
我們對上帝對他所創造的萬事萬物.
常存恭謹敬畏的心.
承認自己的無知.
承認自己的微諸.
向主講.
主下請說.
僕人敬聽.
我們除了對上帝這樣之外.
我們對身邊所遭遇的人和事.
也都應該是這樣.
《教宗選戰》這套電影最末段的情節.
給我們一個驚嚇的就是.
那個被選上為新教宗的人.
竟然是天生的雙性人.
你記得嗎?.
如果你看過這部電影的話.
他肯定有男性的生殖器官.
不然他不可能在修院住了這麼多年.
集體生活.
不過他的身體裡竟然有卵巢.
有子宮.
這個就是不開刀都不知道的.
所以後期才發現的.
你要注意.
他不是做了變性手術.
由男變女.
那就是我篡改上帝的創造.
他沒有.
他天生就是這樣.
所以.
他是沒有信仰或倫理責任的.
不是我搞到自己變成這樣.
我本來就是這樣.
問題是.

$^{681}$我們能不能承認生命是複雜的.
多過我們原來黑白分明的認知.
男的.
女的.
帶有女性器官的男性.
你能不能夠忍受這樣的事實.
什麼叫做男性.
生命太複雜.
或者說.
每一個生命都是獨特的.
也因此合成的世界是複雜的.
我專心去聽一個身邊的小兄弟.
講他的故事.
不要太快講兩句.
他講兩句我就馬上將他歸檔.
他大概是我認識的.
Cat A.
Section 2.
那種人.
馬上就對他做了足診斷.
不要.
每一個生命故事都是獨特的.
都蘊含上帝特別的旨意.
也因此都是讓我們眼界開闊的機會.
還有.
我們承認生命裡有很多辦法解說的奧秘.
不要隨意套用既有的屬靈八股的解說.
不要用人無知的言語.
使上帝的旨意暗昧不明.
是的.
我們所做的很多神學解釋.
很多時候不是了解解明上帝的旨意.
是混亂上帝的旨意.
就是越解釋越困惑.
越解釋越偏離上帝的旨意.
約伯三個朋友不是唯一的例子.
我很喜歡保羅的這段經文.
不會解釋的.
因為時間所限.
這本書我已經得著.

$^{721}$已經完全了.
我也是竭力追求.
或者可以得著基督耶穌.
所以得著我的.
等一會.
我不是以為自己已經得著了.
我只是一件事.
就算忘記背後努力面前.
你向著標桿直跑.
要得上帝在基督耶穌裡.
從上面召我來得的掌上.
所以我們中間凡事完全人.
總要存在於心.
若在什麼事上存不在於心.
上帝必以此指示你們.
而我們到了什麼地步.
就當照著什麼地步行.
總結說.
上位恭謹謙卑的心.
不要以為自己已經掌握一切真理.
我是屬靈人.
所以我已經參透萬事.
永遠對上帝持開放的態度.
對面前的日子持開放的態度.
上帝仍然會和我說話.
藉著我很熟悉的經文和詩歌.
上帝今天仍然有話要和我說.
我仍然有進步的空間.
我都和你說過.
我最快樂的一件事就是.
靈修的時候讀一些很熟悉的經文.
聽一首很熟悉的詩歌.
心中仍然有一些觸動.
觸動不是在於.
又有新見解.
不是這麼多的.
而是你覺得上帝還沒有放過我.
他還和我說話.
你覺得很過癮嗎.
上帝仍然肯和我說話.

$^{761}$這個多棒.
他仍然相信我受教.
仍然相信我有進步的可能.
是不是.
弟兄姊妹.
上帝最近動你是什麼時候.
你聽到他和你說話是什麼時候.
過去成功的事奉經驗.
不可以保證今天仍然成功.
所以保羅說.
忘記背後光輝的歲月.
專注眼前的任務和挑戰.
小心區分真理和真理的應用.
真理不變.
那個implications and applications.
什麼時候都可以有修訂的可能性.
我不是以為自己已經得著了.
已經完全了.
我們仍然走在路上.
對我們的信仰.
唯一的可靠.
是在於我們眼前.
這個昨日今日.
直到永遠都不變的耶穌基督.
\newpage



\section{約拿書 4:1-11}
\label{sec:uiuY8tue_2Q}
\textbf{2020年11月29日 崇拜直播｜約拿書4\_1-11}
\newline
\newline
連結: \href{https://youtube.com/watch?v=uiuY8tue-2Q}{\texttt{https://youtube.com/watch?v=uiuY8tue-2Q}} ~~~~ 語音日期: 2025-01-17
\newline
\newline
\hyperref[sec:YgvzlIV1PPE]{\small{< < < PREV SERMON < < <}}
~
\hyperref[sec:index_chronic]{\small{[返順時目]}}
~
\hyperref[sec:KYvJDDo3N1U]{\small{> > > NEXT SERMON > > >}}
\newline
\newline
約拿書 4:1-11
\newline
\begin{longtable}{cl}
\hline
\hline
章節 & 經文 (和合本修訂版)\\
\hline
4:1 & \begin{tabularx}{0.7\textwidth}{X} 這事令約拿大大不悅，甚至發怒。 \end{tabularx} \\ \\ \relax
4:2 & \begin{tabularx}{0.7\textwidth}{X} 他就向耶和華禱告，說：「耶和華啊，這不就是我仍在本國的時候所說的嗎？我知道你是有恩惠，有憐憫的神，不輕易發怒，有豐盛的慈愛，並且會改變心意，不降那災難。我就是因為這樣，才急速逃往他施去的呀！ \end{tabularx} \\ \\ \relax
4:3 & \begin{tabularx}{0.7\textwidth}{X} 耶和華啊，現在求你取走我的性命吧！因為我死了比活著更好。」 \end{tabularx} \\ \\ \relax
4:4 & \begin{tabularx}{0.7\textwidth}{X} 耶和華說：「你這樣發怒，對嗎？」 \end{tabularx} \\ \\ \relax
4:5 & \begin{tabularx}{0.7\textwidth}{X} 約拿出城，坐在城的東邊，在那裡為自己搭了一座棚。他坐在棚子的蔭下，要看看城裡會發生甚麼事。 \end{tabularx} \\ \\ \relax
4:6 & \begin{tabularx}{0.7\textwidth}{X} 耶和華神安排了一棵蓖麻，使它生長高過約拿，影子遮蓋他的頭，使他免受苦難；約拿因這棵蓖麻大大歡喜。 \end{tabularx} \\ \\ \relax
4:7 & \begin{tabularx}{0.7\textwidth}{X} 次日黎明，神卻安排一條蟲來咬這蓖麻，以致枯乾。 \end{tabularx} \\ \\ \relax
4:8 & \begin{tabularx}{0.7\textwidth}{X} 太陽出來的時候，神安排炎熱的東風，太陽曝曬約拿的頭，使他發昏，他就為自己求死，說：「我死了比活著更好！」 \end{tabularx} \\ \\ \relax
4:9 & \begin{tabularx}{0.7\textwidth}{X} 神對約拿說：「你因這棵蓖麻這樣發怒，對嗎？」他說：「我發怒以至於死，都是對的！」 \end{tabularx} \\ \\ \relax
4:10 & \begin{tabularx}{0.7\textwidth}{X} 耶和華說：「這棵蓖麻你沒有為它操勞，也不是你使它長大的；它一夜生長，一夜枯死，你尚且愛惜； \end{tabularx} \\ \\ \relax
4:11 & \begin{tabularx}{0.7\textwidth}{X} 何況這尼尼微大城，其中不能分辨左右手的就有十二萬多人，還有許多牲畜，我豈能不愛惜呢？」 \end{tabularx} \\ \\
[1ex]
\hline
\hline
\end{longtable}
$^{1}$今日選讀的經文是《約拿書》第四章.
在《舊約聖經》一千二百五十九頁.
《約拿書》第四章.
「這是若拿大大不完.
且甚乏勞.
就禱告耶和華說.
耶和華啊.
我在本國的時候.
豈不是這樣說嗎.
我知道你是有恩典.
有憐憫的神.
不輕易發怒.
有豐盛的慈愛.
並且後悔不降所恕的災.
所以我急速逃往他思去.
耶和華啊.
現在求你取我的命吧.
因為我死了.
比活著還好」.
耶和華說.
「你這樣發怒合乎理嗎.
於是若拿出城.
坐在城的東邊.
在那裡為自己搭了一座棚.
坐在棚的任下.
要看看那城究竟如何」.
耶和華神安排一夥悲罵.
使其發生高過若拿.
影而遮蓋他的頭.
救他脫離苦楚.
若拿因這夥悲罵大大喜樂.
此日黎明.
神卻安排一條蟲子咬這悲罵.
以致呼告.
日頭出來的時候.
神安排炎熱的東風.
日頭捕殺若拿的頭.
使他發昏.
他就為自己求死.
說.

$^{41}$「我死了比活著還好」.
神對若拿說.
「你因這夥悲罵發怒合乎理嗎」.
他說.
「我發怒以致於死.
都合乎理」.
耶和華說.
「這悲罵不是你栽種的.
也不是你培養的.
一夜發生.
一夜乾死.
你尚且外色.
何況這歷歷未大成.
其中不能分辨左手右手的.
有十二萬多人.
並有許多生畜.
我豈能不外色呢」.
各位弟兄姊妹.
今日我們很高興.
有我們以天水為勳堂.
天恩勳勳堂的.
邱潔成傳道.
在我們當中正講上帝的說話.
今天的講題是「服唔服」.
我們有請邱傳道.
各位母堂的弟兄姊妹平安.
(弟兄姊妹:平安).
支持無論在祈禱,人力,物力方面.
都一直支持勳堂.
不知道大家還記不記得.
我們舉辦過一個「稀房」.
早前經過呼籲後.
有許多物資.
床,櫃等等.
都有很多.
最後弟兄姊妹.
都親力來到我們當中.
佈置和油油.
做了許多工作.
下星期.

$^{81}$應該說這個星期.
就會開始有第一個住戶.
我們會和他們簽約.
接著亦會有第二個.
希望這些我們同心合力.
一起可以將有需要的人幫助到.
亦有機會將福音和他們分享.
另外我們亦舉辦了一個「功課扶桃班」.
隨時山長水遠.
一星期兩次.
來到我們當中.
服侍一班有需要的中學生.
中一至中三.
他們那份無條件的付出.
付出都不計較.
不僅不收錢.
還要買糖果和食物.
去鼓勵他們.
希望透過這些服侍.
讓他們體會到神的愛.
亦讓我們有機會和他們傳福音.
這些等等的事.
我們亦會繼續在天水圍.
努力做福音的工作.
今日我和大家分享約拿書.
相信約拿書在大家當中.
都不會是一本很陌生的書卷.
當中有很多神蹟.
亦有很多段落.
大家都會認識的.
在短短的四章裡.
我們看到一些神蹟.
是很特別的.
譬如約拿在漁福裡面三天.
不死的.
經過很多歷史考察.
究竟是怎樣的呢.
真的很神蹟.
約拿去到利美尾城.
說了一句說話.

$^{121}$再過四十天.
利美尾就要毀滅了.
然後全城上下的人都悔改.
是不是很厲害.
我們舉辦一個報道會.
很多準備的.
又要請嘉賓.
又要想想誰講完.
又有很多報紙.
可能不是很多人相信.
但約拿就這樣說了一句.
全城上下就有很多人相信了.
即是他相信了神.
然後又有一個.
剛才我們也讀到一個神蹟.
就是有一個悲罵.
一日就長成.
一聲一日就謝了.
這些我們看到了很多神蹟.
究竟我們看約拿書.
是不是很羨慕.
我們可以經歷到這些神蹟.
很體會到神出手的作為呢.
還是約拿書有一個更加重要的核心訊息.
跟讀者去講呢.
約拿書是用一種.
諷刺文學的手法去寫成的.
在四章的約拿書裡面.
我們看到約拿跟神之間.
有一種很緊密的互動關係.
今天我們就嘗試去看一看.
這些互動關係.
神究竟跟約拿想講些什麼呢.
而約拿對上帝的反應.
上帝又如何去忍耐.
如何去教導他呢.
這個約拿可以說.
神經常要他向左走.
他就要向右走.
無論如何都不肯聽神的話.

$^{161}$在這些對話和互動裡面.
這本書究竟是發出什麼訊息呢.
神究竟想跟我們講些什麼呢.
我們就逐項看下去吧.
我們看到.
我也是第一次用.
黑色字的就是神的反應.
或者神的主動.
藍色字的就是約拿的反應.
我們開始一看.
在第一章第一節.
耶和華臨到阿米泰的兒子約拿說.
開始書卷第一個訊息.
你起來往利美尚大城去.
向其中的居民呼喊.
因為他們的惡達到我面前.
所以神就吩咐約拿去做事.
叫那些人去悔改.
作為一個先知神的代言人.
他的責任就是聽了神的話.
然後就將神要說的話向人說.
不過約拿不單止沒有聽神的話.
他的選擇就是調頭走.
我們看到.
大家看不看得清楚.
我試試看.
按不到.
這裡.
總之約拿那裡.
本來是這個位置.
神就叫他去利美尚那裡.
然後他就調頭走去他斯.
在當時的地方.
他斯可以說是天涯海角.
可以去到最遠的地方.
而我們看這個距離.
離利美尚和他斯是三倍距離.
而且要去他斯經海路.
是一個很困難.
當時很凶險的旅程.

$^{201}$但是約拿仍然選擇不去.
而是去一個很危險的地方.
所以我們看到約拿那份決心.
反抗神.
或者不聽神的決心是多麼強.
我們繼續看下去.
經文說.
他選擇了去一個遠遠的地方.
與其說他是逃避神.
倒不如說他逃避責任.
因為作為一個先子.
應該說神要說的話.
不過他不肯.
約拿不願意去利美尚的原因.
一方面他覺得這是他們敵對的國家.
所以他不想去.
第二就是利美尚是一個外邦人.
外邦的城市.
他認為上帝的救恩.
他們不值得得到上帝的救恩.
所以在這件事上.
神未曾去審判利美尚的人.
約拿已經作出審判.
認為他們亦不需要給機會悔改.
所以最好的做法就是.
不去.
讓他們聽到審判訊息.
都沒有機會.
所以約拿給他們一個什麼機會呢.
給他們一個零的機會.
令他們繼續犯罪.
然後神就會毀滅他們.
以往我小時候都聽過一些長者.
或者一些基督徒長者.
因為經歷了戰爭的痛苦.
他們一方面不想和日本人交往.
亦有些基督徒說.
我不會和日本人傳福音.
上一代.
他們那種心情.

$^{241}$或者他們的痛苦經歷.
我們未必完全體會到.
因為我們不是當事人.
我們未必完全明白.
不過在我們心目之中.
有沒有一些人.
我們認為他們不值得神去拯救呢.
可能這些人曾經和我們有過節.
或者因為他們太清楚他們的為人.
覺得他們太差.
甚至他們做了.
你知道他們做了很多傷害人的事.
所以在我們心目中.
覺得他們是無法拯救.
神是否應該拯救他們.
他們是否配得神的恩典呢.
很自然我們就不會願意去接近他們.
亦不想向他們傳福音.
如果真是這樣.
如果我們真的這樣想.
我們是否已經為神做了一個裁決呢.
經文繼續這樣說.
約拿逃到他斯上了船.
上了船之後.
經歷了大風浪.
經文說.
然而夜禍使海中起大風.
海有狂風大浪.
甚至船幾乎破壞.
約拿選擇去到艙底.
然後說不理會.
但因為艙底太顛覆.
船也差不多要沉.
約拿對船員說了些話.
他對他們說.
你們將我抬起來.
拋在海中.
海就平靜了.
我知道你們遭這大風浪.
是因我的緣故.

$^{281}$約拿很清楚.
為何無端端有狂風大浪.
約拿心知肚明.
這些狂風大浪是神出手.
才會出現.
不過約拿怎樣呢.
約拿繼續選擇和神鬥氣.
寧願不聽神的命令.
竟然叫船員將他拋下海.
約拿那份要死的決心.
或不肯聽神命令的決心.
我們看到是多麼強烈.
故事的發展是怎樣呢.
很神奇.
神安排了一條大魚.
將約拿吞下肚子.
吞下肚子三天.
約拿在這段安靜的時間.
甚至可以說是經歷死亡的時間.
經歷生死的時間.
他就作出了一個祈禱.
他說: 諸水環繞我.
幾乎掩沒我.
深淵圍住我.
海水纏繞我的頭.
我下到山根.
地的門將我永遠關住.
在經文裡面有提及.
他已經很接近死亡.
在心海裡面.
他講到: 我下到山根.
地的門將我永遠關住.
在希伯來傳統裡面.
地的門將我關住.
這個意思就像是陰間的意思.
在希伯來人的觀念.
就像是陰間 接近死亡.
當約拿真真正正面對死亡的一刻.
他也害怕.
所以在他禱告的其中一句說話.

$^{321}$他跟神說:.
我所許的願.
我必償還.
這句說話暗示了他願意聽神說.
去到尼尼美城替神發出審判的訊息.
弟兄姊妹.
我不知道你們有沒有經歷過困難.
但如果困難可以令我們更明白自己的限制.
看到自己的不是.
明白到自己不是時時都可以自作主張.
學習再次誠心去倚靠神.
經歷困難.
就可以說我們在人生.
在神那裡上了一個寶貴的一課.
約拿是否真正上到這個寶貴的一課呢?.
在我們的生命 在我們的人生經歷裡面.
在我們跟神那種關係互動裡面.
我們有沒有經歷過神一些很深刻的教導.
以致我們能夠學到放下自己.
不再有不必要的堅持呢?.
經文我們繼續看下去.
耶和華吩咐雨.
然後把約拿套出來.
套出來後.
神第二次跟約拿說.
「起來往尼尼微城去.
向其中的居民宣告我所吩咐的話」.
因為約拿自己答應過神.
答應過神要說一些宣告的訊息.
審判的訊息.
所以約拿就去.
去到尼尼微城.
說了一句很短的說話.
就是「再過四十天,尼尼微城就要傾覆了」.
這句說話原來在希伯來文裡面.
只有五個字.
我不知道你有沒有跟人傳過福音.
我最快手跟人傳了福音.
已經是很神跡.
都要一小時.

$^{361}$說了很多東西.
然後他才肯信.
一次就完成.
已經很快了.
但約拿只說了五個字.
說了五個字就說完了.
其實尼尼微城是一個很大的城.
經文是這樣說.
「約拿便照著耶和華的話起來往尼尼微去.
在尼尼微城極大,有三天的路程」.
即是起碼要走三天,繞一圈.
這樣才走完這個城.
但約拿走到那裡.
走了一天.
然後就說了一句說話.
我們看到約拿是多麼有心.
是多麼有誠意.
所以我覺得約拿是說完了.
我答應了你,阿神.
我已經說了.
就說完了.
說完了就算了.
不過就出現了一個很戲劇性的結果.
竟然全城上下.
由黃到一些卑微的市民.
全城都信,歸信耶和華.
他們也作出了一些行動.
表達了他們真心悔改.
於是神察按他們的行為.
見他們離開罪惡.
他就後悔.
即是說神不再用想著去毀滅他們.
不把所說的災禍降於他們.
我們看到人相信神.
應該是很開心的.
你沒有看到人相信神.
沒有感覺,這是第一.
看到人相信神,OK.
或者看到人相信神.
有沒有搞錯?.

$^{401}$或者會發怒.
約拿看到一群人悔改.
大大不悅.
且甚發怒.
並且求死.
不開心.
發怒.
和要求死.
約拿才知道.
有什麼大理由去發神脾氣呢?.
他的原因是.
因為神有恩典,有憐憫.
不輕易發怒.
有豐盛的慈愛.
最後不把災禍淋到離離未成.
約拿不久前經歷了神在生死的邊緣.
神拯救他.
在餘福那裡.
他知道他快要死了.
所以他去求神.
神也拯救了他.
他也正正經歷了神的恩典.
憐憫.
神沒有向他發怒.
也經歷了神的慈愛.
但是當約拿看到.
離離未人去經歷神的恩典.
憐憫.
也不發怒.
又看到那份慈愛的時候.
約拿很生氣.
很生氣.
生氣得想死.
因為在約拿心目中.
只有他或屬於他的人.
神的子民才可以有這份恩典,憐憫,慈愛.
而那些離離未人.
作惡滔天的人.
豈有此理.
神怎可以對他們這樣.

$^{441}$所以他說.
對神說.
不如你拿了我的命來算了.
為什麼他這樣說呢.
因為約拿看到他最不想看到的結果.
正如他所料.
離離未人.
他們有機會悔改.
他們悔改之後.
他們也知道神是那份.
有恩典,有憐憫.
亦想拯救人的神.
想給機會人的神.
如果正如他想.
他初初逃去他施的時候.
那些人沒有機會聽到的時候.
神就會審判他.
所以他看到最不想看到的結果.
他就生神氣.
生得想死.
不過我們很.
這種出現是很奇怪的事情.
先知不是看到人悔改應該開心嗎.
看到神願意赦免一些.
歸順悔改的人.
應該開心.
不過.
我們知道.
約拿是先有立場.
後有判斷.
因為他事先有既定的立場.
所以他已經判斷了.
不應該這樣發生.
不想這樣發生.
離未未人和約拿之間.
站在一個對立的位置.
所以在約拿的心目中.
早已判定離未未人應該要死.
我們會否對我們不喜歡的人.
傳福音呢.

$^{481}$或者當知道一些我們心目中不喜歡的人.
信了耶穌.
反而會覺得.
神啊,為什麼你會救他.
我們是否因為有了立場.
造成了我們人與人之間.
那份永遠的牆呢.
抑或是相反.
是否因為我們看到神的赦罪.
而能夠化解我們人與人之間那份牆呢.
在第四章的第四節.
神問約拿.
為什麼你發脾氣.
你發脾氣是對的嗎.
武彩神到離未未城外.
搭了一座棚.
要飛了.
搭了一座棚之後.
慢慢看著城內的情況.
約拿為什麼要這樣做呢.
他這樣做有什麼作用呢.
如果約拿真的不滿意神寬恕那群離未未人.
以他的性格.
他便不看了.
走了,算了.
上帝想這樣就這樣吧.
但是他仍然選擇留在那裡.
看著離未未城內所發生的情形.
很明顯他也曾經經歷過.
有些人悔改之後會回頭.
大家都見過.
所以他很想看見離未未人.
如果他們再次犯罪.
他便可以大聲跟神說.
他們就是一群叛逆神的人.
上帝你應該去審判他們.
不過當然這些事情沒有發生.
我們會否有時所謂關心一些人.
收一下資料,八卦.
關心一些我們覺得心目中.

$^{521}$不太看得過眼的人.
這些所謂的關心.
究竟是真的想關心他們,幫助他們.
還是想知道他們何時失敗.
如果知道的時候便會說.
是否呢?一早就說了.
他們怎會這麼好.
約拿作為一個先知.
他的本份應該是去守望那群人.
希望他們離開罪惡,重回正道.
作為一個基督徒.
會否都很想自己認為所謂的壞人.
會繼續好.
還是想他沒有好結果呢?.
我們有沒有為他祈禱.
為他守望他們.
我們會否有約拿這種心態.
還是像保羅所說.
教導我們那種心態.
你們當以基督耶穌的心為心.
我們知道耶穌基督.
是一次又一次給我們的機會.
就像彼得一樣.
一次又一次給彼得的機會.
經文繼續說.
耶和華安排了一棵悲麻.
使其發生高過約拿.
影而遮蓋他的頭.
救他脫離苦楚.
因為他坐在那裡.
在那個地方很熱.
神安排了一棵樹.
可以借陰.
約拿很開心.
約拿因這棵悲麻大大喜樂.
為何約拿會因為有遮陰的樹.
有少許甜頭而開心呢?.
約拿一個那麼堅強的人.
一個不怕辛苦的人.
被太陽曬一下又如何呢?.

$^{561}$為何會這樣又會開心呢?.
可能他的感覺不是因為有棵樹.
遮陰而開心.
反而是因為他認為神看著他.
神仍然很疼他.
神仍然不想他那麼微小的事.
神仍然願意幫助他.
他覺得神仍然是一個重要的人.
有時可能做了不好的事.
你的上司或尊敬的人.
跟你說幾句好說話.
解一下困.
你都會特別開心.
因為你會覺得他願意幫助你.
我在他心目中仍然有一份位置.
故事的發展就是這樣.
次日黎明.
神卻安排一條蟲咬樹.
令你枯乾.
還不止.
日頭出來的時候.
神安排炎熱的東風.
即是再加碼.
日頭曝曬若拿的頭.
使他發昏.
所以我們看到.
神不是做些事去哄若拿.
相反神是要若拿明白神的心意.
所以他安排這棵樹.
很戲劇化的生和死.
讓若拿自己去明白自己的矛盾.
表面上他想的很合理.
但其實他想的很自私.
神不斷用一些實物的教材.
去勸導他.
使若拿明白.
希望他能夠回轉.
希望若拿能夠悔改.
我們看到風大浪的經歷.
跌下海也是一條魚的經歷.

$^{601}$他曾經說過一句話.
全城上下那麼多人都能歸順神.
一棵飛馬樹的生死.
一次又一次去教導若拿.
希望他消除自己一些的成見.
但若拿去到這個時刻.
他就說.
他就為自己求死說.
我死了比活著還好.
第二次問若拿.
神對若拿說.
你因這棵飛馬法路對不對?.
你就是因為這棵樹.
有得遮陰.
沒得遮陰.
就對我法路.
就發我脾氣.
對不對?.
但若拿說.
我法路已至於死.
怎樣?都合乎你.
我這樣法路是對的.
是正確的.
我不知道你有沒有對神這樣說話.
當然這一段說話我們會看到.
好像很戲劇性.
去反映了一些人內心.
那種的心額.
那種的自以為是.
但神聲好氣跟他說.
這棵飛馬不是你栽種的.
也不是你使它長大的.
一夜成長.
一夜死去.
你尚且外適.
即是這棵樹都與你無關.
你沒有努力過.
不是你栽種.
你若對它沒有感情.
你也會那麼珍惜.

$^{641}$其實我們說了.
那份珍惜不是真的珍惜那棵飛馬樹.
是珍惜背後.
即是上帝對它的看法.
若拿很明顯.
關注的不是那棵樹.
關注的是他自己.
關注的是神跟他那種的關係.
作為先知他很明白.
這一棵樹的一夜生死.
是神的作為.
所以神這樣對著他.
用一種方式插進他的內心.
那種的想法.
所以他更加憤怒.
更加在神面前求死.
神繼續對若拿說.
若拿你不認同去拯救尼尼微城.
尼尼微人.
若拿不是不知道.
上帝想拯救這些人.
若拿是不認同.
若拿也不是不知道,不理解.
是不同意和不肯接受神的做法.
神既然看到尼尼微人能夠悔改,回轉.
神當然會愛惜這些生命.
給他們回轉的機會.
因為神不是憎恨某個人或某些人.
神是憎恨罪惡.
神不是憎恨人.
而是憎恨罪惡.
不過人有時將憎恨人和憎恨罪惡.
混在一起,分不開.
所以當一個人.
可能是一個罪人.
或是一個我們跟他有牙齒印的得罪過的人.
就算他悔改了.
在我們心目中,可能我們仍然不喜歡這個人.
甚至放下不了那些仇恨.
可能有一個得罪過的人.

$^{681}$他很真心跟我道歉之後.
我們是否真的能夠冰釋前嫌呢?.
在整卷的約拿書作者.
用了利美尼全城的人悔改.
和若拿那份堅持執著.
作了一個很強烈的對比.
若拿去到一個罪惡滔天的地方.
說了五個字.
希伯來文五個字.
全城的人就悔改了.
但我們看到在四張約拿書裡.
上帝用了種種的方法勸他.
上帝親自跟他說話.
亦有很多經歷讓他經歷過.
但若拿就是死牛一面頸.
無法談,死撐,寸步不讓.
最難悔改的不是那些罪惡滔天的人.
反而是一個認識上帝的先知.
大家覺得是否很諷刺呢?.
很多讀者看完約拿書之後.
看不到結果.
若拿究竟會怎樣呢?.
他最後會怎樣呢?.
他會聽神的話嗎?.
還是他自行決定呢?.
我覺得這本書寫下來.
作者不是想我們去猜測結果.
最重要的是當我們看完約拿書之後.
如果我們有同樣的經歷.
我們會怎樣去選擇呢?.
我們繼續選擇自己的看法.
堅持自己的判斷.
抑或我們願意回轉悔改.
當然神不會赦免一些仍在罪中生活的人.
但當人願意認罪,離開罪惡.
神亦會給人重新的機會.
我們又怎樣呢?.
我們都是神給我們機會.
我們才能重新過一個.
合乎上帝生活的基督徒.

$^{721}$但我們會否有太強烈的個人判斷.
就算神已經用多種方法提醒我們.
都不聽神的話.
這種偏執在聖經中形容為心硬.
如果看過聖經,你都知道甚麼人最心硬.
究竟我們行事為人的方法.
是用我們自己的尺去量度.
抑或學習用神的尺去量度呢?.
我們願不願意放下自己的執著.
放下自己的看法或成見.
去信服上帝的教導呢?.
讓我們一起低頭祈禱.
天賦上帝.
我們認為自己是認識上帝的人.
亦領受上帝的恩典,憐憫,赦罪.
以致到今天我們願意去跟從上帝.
但與此同時,當我們跟人相處的時候.
可能我們仍然有很多自己的看法,堅持或執著.
這些是否必要或正確呢?.
求上帝在我們日常生活提醒我們的時候.
讓我們能夠放下自己一些很難放下的想法.
求上帝改變我們.
我們這樣祈禱.
主席.
\newpage



\section{申命記 3:18-28}
\label{sec:KYvJDDo3N1U}
\textbf{2025年1月26日崇拜直播｜鄭國邦傳道｜朝向應許之地－寬容與堅拒｜申命記3\_18-28}
\newline
\newline
連結: \href{https://youtube.com/watch?v=KYvJDDo3N1U}{\texttt{https://youtube.com/watch?v=KYvJDDo3N1U}} ~~~~ 語音日期: 2025-01-26
\newline
\newline
\hyperref[sec:uiuY8tue_2Q]{\small{< < < PREV SERMON < < <}}
~
\hyperref[sec:index_chronic]{\small{[返順時目]}}
~
\hyperref[sec:cuIe7N8imyc]{\small{> > > NEXT SERMON > > >}}
\newline
\newline
申命記 3:18-28
\newline
\begin{longtable}{cl}
\hline
\hline
章節 & 經文 (和合本修訂版)\\
\hline
3:18 & \begin{tabularx}{0.7\textwidth}{X} 「那時，我吩咐你們說：『耶和華－你們的神已將這地賜給你們為業；你們所有的勇士都要帶著兵器，在你們的弟兄以色列人前面過去。 \end{tabularx} \\ \\ \relax
3:19 & \begin{tabularx}{0.7\textwidth}{X} 但你們的妻子、孩子、牲畜，可以住在我所賜給你們的各城裡，我知道你們有許多牲畜。 \end{tabularx} \\ \\ \relax
3:20 & \begin{tabularx}{0.7\textwidth}{X} 等到耶和華讓你們的弟兄像你們一樣，得享太平，他們在約旦河另一邊，也得了耶和華－你們的神所賜給他們的地，你們各人才可以回到我所賜給你們為業之地。』 \end{tabularx} \\ \\ \relax
3:21 & \begin{tabularx}{0.7\textwidth}{X} 那時，我吩咐約書亞說：『你親眼看見了耶和華－你們的神向這兩個王一切所做的事，耶和華也必向你所要去的各國照樣做。 \end{tabularx} \\ \\ \relax
3:22 & \begin{tabularx}{0.7\textwidth}{X} 不要怕他們，因為那為你們爭戰的是耶和華－你們的神。』」 \end{tabularx} \\ \\ \relax
3:23 & \begin{tabularx}{0.7\textwidth}{X} 「那時，我懇求耶和華說： \end{tabularx} \\ \\ \relax
3:24 & \begin{tabularx}{0.7\textwidth}{X} 『主耶和華啊，你已開始將你的偉大和你大能的手顯給你僕人看。在天上，在地下，有甚麼神明能像你行事，像你有大能的作為呢？ \end{tabularx} \\ \\ \relax
3:25 & \begin{tabularx}{0.7\textwidth}{X} 求你讓我過去，看約旦河另一邊的美地，就是那佳美的山區和黎巴嫩。』 \end{tabularx} \\ \\ \relax
3:26 & \begin{tabularx}{0.7\textwidth}{X} 但耶和華因你們的緣故向我發怒，不應允我。耶和華對我說：『你夠了吧！不要再向我提這事。 \end{tabularx} \\ \\ \relax
3:27 & \begin{tabularx}{0.7\textwidth}{X} 你上毗斯迦山頂去，向東、西、南、北舉目，用你的眼睛觀看，因為你必不能過這約旦河。 \end{tabularx} \\ \\ \relax
3:28 & \begin{tabularx}{0.7\textwidth}{X} 你卻要吩咐約書亞，勉勵他，使他壯膽，因為他必在這百姓前面過去，使他們承受你所要觀看之地。』 \end{tabularx} \\ \\
[1ex]
\hline
\hline
\end{longtable}
$^{1}$弟兄姊妹平安.
平安.
我這樣能夠在這裡跟大家分享神話.
因為快要新年了.
很快今天已經是年廿七了.
明天就是年廿八洗垃圾的日子.
所以我相信下個星期和這個星期.
大家一定很忙碌.
正在準備新年的事.
所以在這裡先跟大家拜個早年.
因為下個星期應該是年初二.
對,應該是年初五.
已經是新年的時候了.
所以祝大家.
新年裡面能夠蒙神的因.
能夠身體健康.
每天都能夠靠著主因去更新.
那麼我想剛剛說了.
可能這段時間大家也忙碌.
正在準備新年的事.
早幾天如果大家經過彌敦道.
都看到匯豐銀行也排隊了.
每年新聞都會報道.
就是很多百姓,市民排隊唱新人紙.
可能你未必會唱新人紙.
可能你也有很多事情要預備的.
是不是?.
剛剛說了.
明天就是要洗垃圾的日子.
可能正在丟新的東西.
添置新的東西.
或者可能要買一些揮春.
或者利是封.
即使你信了主也會有這些習俗的.
是不是?.
你覺得這份是我們的文化.
就是我們的新年.
即使信了主.
我們也有一些祝福的揮春等等.
那麼.

$^{41}$除了新年有很多事情要預備之外.
我想大家都知道很多不同的習俗.
新年裡面有很多事情要留意的.
除了要準備很多東西之外.
我們有些日子也會很留意的.
就是大家.
我有一次在搭地鐵的時候.
聽到旁邊有兩位女士聊天.
我覺得她們的內容.
是很重溫以前一直聽的一些內容.
她們說什麼呢?.
就是說.
我們鄉下的傳統是這樣的.
初一開始就不可以掃地.
不知道大家有沒有這種傳統.
不可以掃地的.
因為如果初一年大清頭.
年初一,二,三的時候掃地.
掃走運氣和財氣.
然後那個人就問.
什麼時候可以掃地呢?.
初五.
初五才可以倒垃圾的.
然後她說不僅如此.
然後大年初一和初二也不可以洗頭的.
大家的姐妹小時候有沒有經歷過這些呢?.
我也經歷過.
不可以洗頭的.
洗是不好的.
要等到初三才洗頭的.
其他的日子大家都知道了.
初三就是赤口.
如果他們有其他的宗教或者信仰.
就要去廟或者其他地方.
然後在那裡聽的時候.
哇.
這麼多的嗎?.
以前可能我只聽過年廿八洗垃圾.
賣冷賣到年三十晚的那些.
這些都是我們習俗去做的一些事情.

$^{81}$那個人就不斷去說.
然後另一個人就說.
喂!不要說這麼快.
再說一次.
我要記一下這些時間是什麼時候等等.
那一刻可能因為我也在預備今天的講道.
一邊聽的時候.
哇.
他就在申命記住.
他就重提這些過往的條例或者這些事情.
告訴一個可能不太認識.
不太知道的人.
你怎麼做才可以有這份福氣.
有這些平安.
這些才會是好的.
當然.
我覺得那個人只是再說一次讓他知道.
但是聖經裡面的申命記.
我想大家都知道不是的.
不是一回事.
不是只是說.
新年的時候才做這些事情或者等等.
當我們看回申命記的時候.
會發覺有相似的地方.
好像重提一些事情.
但是也有不同的地方就是.
我們知道申命記是他們在一個處境裡面.
準備進入應許之地.
當講到申命記的時候.
我想大家都會知道在過往的25年開始之後.
我們也開始了申命記的講道系列.
兩位陳牧師也已經講過朝向應許之地系列了.
到了第三講的這個裡面.
大家可能也會知道.
朝向應許之地的脈絡和結構.
按照申命記的結構我們會知道.
這一次裡面是摩西在申命記三次的講論裡面.
第一次接近最後的時候.
之前兩位陳牧師也講述了.
他們複述了很多關於以色列人.

$^{121}$在出埃及到曠野準備踏入應許之地的一些經歷.
在裡面成為一些提醒.
當今天的回顧以前的事情.
來到最終段落的時候.
摩西也以他進不了迦南這件事情.
作為一個轉接.
以致帶下去告訴以色列民.
你要聽上帝的話.
要聽他的吩咐.
所以當我們再一次看回申命記這個技術的時候.
就好像剛剛所說.
不同的地方.
不只是將一些歷史.
過往的舊事.
再重提一次.
而是大家再退後一步想起來的時候.
申命記本身是一個很獨特的文體.
它記載著不只是大事回顧等等一些事情.
而是.
你會看到.
他在裡面選擇了一些很重要的片段.
再複述給他.
以致成為一個很重要的基礎.
再向前踏一步的一種敘述.
我記得當我初信的時候.
剛剛讀聖經.
就是大家打開申命記.
不知道第幾章第幾節.
那時候我找了很久.
申命記啊.
哪裡有生.
因為我腦子裡面是想著生命.
就是生命的申命記.
我腦子裡面沒有申命記.
找不到的.
然後導師就說.
不是那個生.
不是生活的生.
是重生的生.
因為申命記是要重新律法.

$^{161}$再說一次.
當我去理解的時候.
我都以為.
申命記究竟再複述的意義在哪裡呢?.
為什麼又要再複述一次過往的事情呢?.
正如剛剛所說.
既然編輯這個申命記.
或者摩西.
或者那個編者.
是有心去選擇這些舊事的時候.
他一定不只是將這些事情.
搬字過紙地放回申命記.
以前發生過這些事情.
搬回來提起.
而是當他編輯這個申命記的時候.
一定有他的心思意念在裡面.
以至當他揀選這些事情過來的時候.
他一定會想串連一個結局.
讓他們知道一些提醒.
所以當我們知道.
我們自己提起一些舊事的時候.
通常是什麼時候呢?.
要麼就是吵架的時候.
吵架的時候.
你以前是這樣的.
以前你就是這樣對我.
以前這樣的事情.
為什麼要複述一些舊事呢?.
就是讓自己有權力告訴你.
你一直都是這樣的嘛.
一直都是這樣做事的.
這樣的情況.
那麼.
當然也有一些另一邊廂的.
什麼時候你會提起一些舊事呢?.
就是通常一群舊同學.
舊同事.
說回以前的日子多好啊.
以前的日子多好.
或者是說回你自己一些.

$^{201}$很厲害的歷史.
很威風的歷史.
那麼你提起以前吧.
哇!以前我就是這樣.
打滾了一段時間.
終於有這些成就了.
可能你又在另一種情況之下.
提起以前的事情.
當然也有一些比較中性的.
就好像剛剛.
可能二四年十二月三十一日.
那麼你打開電視就是那些.
啊.
舊日的足跡啊.
或者是那些什麼大事回顧啊.
或者等等.
就是複述這些事情的過去.
又或者呢.
如果你有參加我們的感恩的分享會.
就是可能你什麼時候會複述以前的事情呢?.
你就會覺得.
啊.
以前的這些經歷.
都是上帝很多的恩典.
妹妹當你再去想這些以前的事情的時候.
為的是什麼呢?.
而申命記.
在這裡再一次提起以前的事情.
目的又是為了什麼呢?.
其實當我們看回申命記整個脈絡的時候.
你會知道.
去提起以前的事情.
重點不是再在以前那裡.
不只是數你以前有多好.
或者數你以前有多壞.
或者說回以前的事情.
哦.
完了.
就是這樣.
他再重提這些事情呢.

$^{241}$不只是為了過去.
是為了將來.
是為了後一輩的人.
或者之後直到我們看聖經的人.
能夠有一個基礎.
知道怎麼在這些過去裡面.
再踏一步向前走.
所以你看到呢.
申命記裡面選的材料呢.
不是什麼都放進去的.
就是有什麼大事啊.
或者等等.
就是呢.
如果讓大家去回顧.
就是黃浸的歷史也是一樣的.
哇.
黃浸發生過這麼多事情.
怎麼可以讓你重新一次走回那條樓梯.
經歷回60年的回顧.
那你也看到不是所有事情都放在上面的.
就是可能有時候.
爆水管啊.
或者什麼都不會寫在裡面的.
是有些重要的人.
或者重要的事情才會放在那面牆上.
然後去提醒我們.
同樣的.
如果你其實將申命記.
比對整個民數記.
利未記或者出埃及記.
倒數進去的時候.
摩西不是每樣事情都提到的.
至少可能有幾件.
其實本身也很厲害的事情.
你看民數記.
記得哪個巴蘭先知嗎?.
就是騎著驢也會說話的那個.
然後又會就坐以色列那些.
沒有提到的.
又或者其實有些事情.

$^{281}$你會看回在申命記裡面.
其實個個都沒有提到的.
在申命記裡面加插了進去.
甚至有些內容.
以前的經文.
沒有很深入地提醒.
但是在申命記裡面再轉深一點.
再深化它.
你就會看到.
申命記.
不只是真的複述歷史這麼簡單.
而是很想藉著歷史.
建構一個新的里程.
讓新的這一群以色列人.
一方面不要重蹈覆轍.
重提舊時其中一個方向就是.
不要做以前的傻事.
或者得罪上帝的事情.
另一方面.
就是喚醒他們一些記憶.
當然這些記憶對新一代的以色列人來說.
可能是他們二十歲前.
如果你看回經文.
就是二十歲前那一群才可以.
過了約旦河去回應他們.
但是.
新的一群.
他們小時候的事情.
以至再去重提的時候.
不要重蹈覆轍之餘.
就是喚醒他們這個集體的意識.
讓他們想回自己的根源是什麼.
這一群人.
跟著以前那一輩出埃及.
來到曠野經歷四十年.
來到約旦河前.
他們會問一個問題.
我是誰?.
我在這裡做什麼?.
我將來會是怎樣?.

$^{321}$可能他們的疑惑都在這裡.
再一次摩西藉著這些故事.
這些假時.
再去提醒他們這個根源.
以至幫助他們.
藉著知道上一代怎樣去做.
而他們都可以怎樣跟著去做.
或者想到.
找到.
究竟現在這一刻.
我所做的是什麼呢?.
我要做什麼在前面的日子裡面.
去回應上帝的帶領呢?.
當然了.
當我們去想的時候.
我們今天可能也會問這個問題.
在我們的生活裡面.
無論年輕的時候.
或者到年長了.
到你有些階段轉變了.
你也會想究竟我是誰呢?.
我在這裡做什麼呢?.
我接下來要做什麼.
去回應我現在的狀況呢?.
可能我們會用很多方法去想.
可能你會去做一些性格分析.
或者是摸索一下前面的路的一些幫助.
去見一下輔導等等.
但是對我們今天最難理解.
要代入在申命記裡面的時候.
很難的一個位置就是.
上一代怎樣幫我.
去找到這個答案呢?.
因為很多時候我們現在自我意識比較強一點.
就是自己的個體的時候.
很多時候我們去想這個問題的時候.
我就是我.
我不用任何人去幫我界定的.
我要找我自己的價值.
為什麼要看這些東西.

$^{361}$我自己去摸索.
藉著我的知識.
藉著我做的事情.
我去找回我自己.
就會忘記了上一輩帶給我們的提醒和幫助.
但是其實當我們讀入申命記的背景裡面的時候.
當時比現在更加強的一個觀念就是.
那個群體.
那個宗族帶給他的身份認同和意識.
以前那個意識很強的時候.
因為最重要的一件事情就是.
你是一群什麼人.
或者你是什麼背景的人.
很影響了他們怎樣對待你和去看你.
其實也不是離我們很遠的.
有時候我們在生活裡面.
我是哪個鄉下的.
或者我是什麼的時候.
你突然之間會有共鳴感的.
或者會說你出自哪間學校的.
或者等等的時候.
其實這個意識也有的.
只不過我們有沒有去想過.
原來上帝藉著這個身份意識.
也提醒了我們下一步.
究竟要去哪裡呢?.
在當時來說.
對他們來說.
上一輩在做什麼呢?.
對他們來說很重要的.
因為呢.
當他們一直聽那些故事.
爸爸怎麼做.
那些小叔伯又怎麼經歷上帝.
上一輩怎麼去跟從上帝的時候.
他們要在這個基礎上.
再去發掘.
再去承襲一下.
究竟他們過去在做什麼呢?.
以至在這個基礎裡面再去想.

$^{401}$啊.
那我又可以藉著他們在做的事情.
怎麼建構下去呢?.
怎麼樣可以幫助我.
再去發掘這個前路裡面的時候呢?.
去找到成為幫助.
就是說其實.
今天我們也會有時候有這樣的狀況.
是不是?.
就是有一些.
特別是食肆.
就是說三傳的秘方.
三代祖傳秘方.
又或者是幾代家傳之寶.
傳下來的東西.
其實這些寶貴是很重要的.
因為你會知道.
這些傳得下來的東西.
是經得起風浪的.
如果我們貿貿然就是說.
你這個是傳統來的.
不用管了.
你就會失卻了這個家傳之寶.
這一個可能他們經歷裡面.
是很重要的事情.
當然.
銀子兩邊都是.
有時候可能就會覺得.
這些會不會是歷史包袱呢?.
背負了這件事情.
會不會阻礙了我們呢?.
一會兒會再說.
那麼下去申命記裡面了.
其實到了第三章.
這裡提到了十一個那氏.
就是那些年.
啊.
以前是怎麼樣的.
以前你們是這樣.
那氏是這樣做了這些事情.

$^{441}$導出了一些.
反映他們究竟是一群什麼人的選擇.
那時候.
你面對這些事情就是這麼選擇的了.
那時候你面對這些事情.
上帝就是這麼幫助你們的了.
那時候在這種處境之下.
上帝就是這麼吩咐你的了.
這十一個那氏.
建構了他們一些共同的經歷.
幫助他們去想想.
啊.
究竟當日他們是怎麼經歷神的呢?.
其實當你看回那十一個那氏的時候.
不是每一件事情都輝煌的.
每一件事情都說那氏你很乖啊.
跟從神啊.
那氏你做得好啊.
這樣子去做.
如果你回頭看那十一個那氏.
其實有些好的地方.
也有不好的地方.
成為借鏡.
在數這些那氏的時候.
其實就是為了建構這一刻的現在.
當你一直細數的時候.
一到三章不斷提很多的那氏.
到第四章就轉接了.
現在我就要將上帝的命令.
教導你以色列人.
摩西當時是這麼說的.
一個轉接由十一個那氏轉到現在.
就是要提醒上一代經歷的事情.
現在這一刻.
怎麼幫助這一輩去前行呢?.
當我們回想我們自己.
啊.
細數那些那氏的時候.
你的心態是什麼呢?.
剛才說了.

$^{481}$除了可能在那些背景之下.
有時候呢.
當你去跟下一輩說那氏的時候.
為的是什麼呢?.
跟申命記那種那氏.
有沒有什麼不同呢?.
咦?.
當我們下一輩了.
就是第二代了.
經常聽到上一代說那氏的時候.
你又在想什麼呢?.
這些我們都一起經歷過的.
你也做過上一代.
可能也做過下一代.
不斷循環.
妹妹當你去說這些歷史的時候.
你的出發點.
為的是什麼呢?.
上一代的時候.
哇.
我由名不經傳到現在.
家傳戶曉.
做了一些很偉大的事情.
That's all.
說完了.
咦?.
你想表達什麼呢?.
被人稱讚你?.
被人說你厲害?.
還是不是呢?.
可能你分享的不只是你的輝煌.
甚至像剛才說的十一個那氏.
有很多軟弱的經歷.
怎樣在跌跌碰碰裡面.
很多血肉的見證.
都跟下一代分享.
不只是說那些輝煌的地方.
當你是下一代了.
聽那些那氏的時候.
你在想什麼呢?.

$^{521}$唉.
又是那三幅被子.
我知道你當年是這樣.
我知道你當年經歷這些事情.
唉.
行了.
都聽過了.
你怎樣聽那些那氏呢?.
你只是聽表面那一層.
他做了什麼厲害或者不好的事情.
有沒有在這些那氏的故事裡面.
聽到弟兄姊妹背後.
那個上帝是誰呢?.
聽不聽到背後.
他想表達的上帝.
怎樣帶領著他經歷.
以至成為你今天的祝福呢?.
近這幾十年.
其實也不止十年了.
教會界裡面也不斷經常說.
教會稱王不接.
就好像.
如果你們去一些聯合聚會.
或者什麼的.
哇.
每個人都五十歲以上.
或者等等.
所以我這些一出去聯合的.
不只是聯會.
免得大家一讀就以為是浸聯會.
很多聯合的地方.
幾個中排.
我們這些年紀一出去.
哇.
你們是哪裡人呢?.
這麼年輕的.
或者什麼的狀況.
所以一直到這些場合.
都很有負擔.
就是怎樣可以連繫上一代.

$^{561}$和下一代之間的那種狀況呢?.
經常說稱王不接.
說好像下一輩.
又好像不太能回應上一代所說的話.
上一代又好像不太能聽下一代說的話.
究竟裡面的狀況.
是因為什麼呢?.
因為很多事情好像只是硬推.
下面要硬吃.
又或者覺得.
下一輩就覺得舊的不適用.
或者等等.
究竟怎樣可以更加適切.
讓兩代之間一起去同行呢?.
摩西就在這個轉折點裡面.
去提醒以色列人.
以致當我們能夠去跨代同行的時候.
我們聽的心態.
都要去轉變.
說的心態.
都要去轉變.
能夠祝福這個群體.
所以過往在三中.
三中就是在說浸信會,宣導會和播道會.
有一些活動都很想做這些聯繫.
我參與其中的時候.
真的發覺不容易.
當一群不同的年紀坐在一起.
想學習傾聽.
都不容易.
但是最重要的那一刻.
我感受到就是.
至少這一群大英姐妹.
這一群牧者.
願意還走在一起.
去試試去聽.
試試去說.
試試去想.
怎樣可以互相成為祝福.
所以盼望這一個.

$^{601}$都成為我們去塑造我們群體的一種很重要的提醒.
而且我們知道.
當《申命記》讀到這個階段的時候.
剛才說了.
摩西是要引入四章一節.
頭那裡說.
當你了解這些那時的時候.
現在我就要將上帝的律例,教導.
放在你們的群體裡面.
再一次去提醒.
這個教導.
如果沒有這個根源的時候.
只會變成什麼呢?.
世界級的黃金定律.
總之放到哪裡都可以.
媽媽是女人.
太陽由東邊升起.
總之你做好事就是一個好人.
這些定律真的放在哪裡都可以.
什麼群體都可以.
這些教導是沒有根的.
但是為什麼摩西要再重提一次背景就是.
這些就是他的根源.
而接下來四章的教導就要直根在這些根源裡面.
教導以色列人.
你們過往是這樣.
以致我們能夠有這種特質.
去承擔著這些上帝的教導的時候.
這個教導就來得更加貼切.
更加 tailor-made.
更加是為這個群體去準備的一種教導.
以致當他成集的時候.
就能夠想起他的成長經歷.
以致再踏前一步.
怎樣去走.
所以當想到成全的時候.
我們都要想.
究竟無論旺盡的DNA是怎樣.
你的生命成長的歷程是怎樣.
都要去想.

$^{641}$以致去承繼著上帝的教導的時候.
能夠更加適切在你的生命裡面.
當然當說到一成全的時候.
其實有時候很容易走向兩個極端的.
之前因為一直去裝備自己的時候.
跟過一個機構叫做好爸爸中心.
就是教弟兄們怎樣去做一個好爸爸.
當然.
沒有教他們.
其實真的很多事情是在內在的經歷等等.
但很多時候我們會發現一個模式.
就是上一代爸爸.
其實是絕對和很影響.
你怎樣做下一代的爸爸.
就是爺爺怎樣去對待那個爸爸.
是很重要地影響那個爸爸.
怎樣對待兒子或者女兒.
這個成集很容易走向兩個極端.
就是以前爺爺這樣做.
我打死都不這樣做的.
我永遠很討厭爺爺這樣做.
我一定不會做的.
就會走另一個極端.
就是就算有多好或者有多差.
都完全摒棄爺爺的形象.
又或者另一個極端.
我以前真的很疼愛我爸爸的.
我很喜歡我爸爸.
所以將所有事情無論好和壞.
照 copy and paste(複製和貼上去).
這樣放在自己的子女裡面.
而那個好爸爸課程就是要教一些爸爸.
你要分哪些要承傳下去.
就要秉持.
抓緊住.
哪些你看到爸爸的狀況不好.
你要 deline(刪減).
就是要不再這樣承傳下去.
或者將那種地方要斷捨.
要切割.

$^{681}$以致將這些事情.
去承傳下去.
成為祝福.
所以今天當我們在兩代裡面的時候.
我們要拿捏這個平衡點.
以致將這一份家傳之寶.
能夠承傳下去我們的下一代.
那麼繼續下去.
當摩西講完這個那時.
到了三章裡面的時候提到.
那一群以色列人.
他們開始在約旦河東分地了.
這個技術也很特別的.
因為其實當新一代的以色列人.
其實他不僅僅是剛剛說的.
照板煮碗.
所有事情都承傳上一代做的事情下去.
其實這一群新的以色列人.
他有機會繼續藉著過去的結果子.
去探索新的可能.
其中一個新的可能.
被摩西寄述下來.
就是三章裡面說的.
他們在探索這個可能的時候.
在第三章第十二節開始說.
他就開始跟摩西說.
那時候我們佔領了這些地方.
不如我們在約旦河東.
有兩個半支派都留在那裡.
你想想當我們讀這個概念的時候.
因為我們可能拿捏了整個舊約的時候.
就覺得這群人這麼快就分地了.
但其實如果你跳進去.
當時來說.
是一個很嶄新的概念.
一向這十二個支派.
整群以色列人想著就是一起過河的.
一起過到河裡面再分地的.
新一輩不同想法了.
喂.

$^{721}$我有一個新想法.
摩西.
可不可以這樣承傳呢?.
我想試試.
不如我們既然家產這麼多.
或者這麼多的產業.
可不可以讓我們在約旦河東也劃一些地.
反正剛剛打完兩個王.
有些地在那裡.
我們留在那裡吧.
如果你看申命記裡面.
其實是很平序地說這件事情而已.
如果你有時間回到民數記裡面.
去看三十二章的時候.
你會看到故事沒有這麼平.
沒有這麼單調.
因為當時當這群新一輩這樣提出的時候.
摩西很驚訝.
哇.
怎麼可以呢?.
怎麼可能你們會留在這裡呢?.
他們在裡面甚至是複述.
你們怎麼可以像以前你們的祖先.
你們的先祖一樣.
那十二個支派的探子那樣.
想背叛我們.
想陷害我們.
陷我們於不義.
所以就罵那兩個半支派.
所以.
其實他們當時對質了一段時間.
但是呢.
來到申命記這裡.
整件事情的格調完全不同了.
當你一直看第三章.
一開始說佔完的地方.
再分地的時候.
摩西就說.
我們恩爵臣的帶領得到這些地方.
我現在就吩咐你們了.

$^{761}$這些地方就分給你們兩個半支派.
去使用了.
但是你們一定要答應我們.
要帶你們的男丁軍隊.
幫我們進入應許地.
接著打贏了.
那你就可以回到你自己河東那邊.
約旦河前那個地方.
回到你家裡.
去得著你們的地方了.
整個鋪排很順理成章.
好像都安排好了.
還有納入了在這個技術裡面.
是上帝的吩咐一樣.
為什麼摩西來到申命記.
會有這樣的轉變呢?.
當然剛才說了.
當我們一讀進去的時候.
很快就看了結果.
這群人不好啊.
最終被人滅的不就是這兩個半支派.
接著就出事了.
很快我們就會讀進了結局.
但是.
要先hold一hold.
不要這麼快.
不要這麼快將這些你已有的想法放進去.
當你按回整個脈絡的時候.
你就會想.
不是的.
這個舉動其實.
摩西是有澄清的.
如果沒有澄清.
上帝沒有說什麼的時候.
上帝已經擊殺了他們了.
已經出事了.
他們不用分那兩個半了.
你不肯過河.
可能已經死在河東了.
你會看到整個技術上.

$^{801}$上帝沒有任何舉動.
沒有任何阻攔.
沒有任何意見.
容許他.
寬容他.
給他.
在那裡留下.
總之他答應.
帶著裝備過河.
跟我們攻打就可以了.
在這裡我們會看到.
摩西也認同了和接納了這一群新一輩的意見.
以致讓他們知道.
原來只要一起去找回核心價值.
就是一起去打仗.
一起在裡面得到的時候.
無論結局如何.
OK.
你知道最重要的地方如何就可以了.
為什麼要hold住結論.
不要這麼快跳到.
你們就是被人滅的那個了.
你想想.
當你看回這個敘事.
再加上之後的敘事.
絕鬱的地方就在這裡.
為什麼摩西要將這兩件事放在一起呢?.
本身就這麼說.
這兩個半支派分地.
其實是一個很普通的敘事.
即使在民數記好像爭吵了很久.
到了申命記.
沒有任何東西.
但是拍下去.
摩西說.
神啊.
既然如此.
可不可以讓我有一個突破一點的想法呢?.
我知道.
你是一個大能的神.

$^{841}$你看到摩西的祈禱是很謙卑的.
是一個很好的祈禱的template(圖案).
先說上帝你做了這麼多大事.
這麼偉大.
這麼有能力.
來吧.
不如讓我過河.
看看櫻水地.
一次看看也好.
這是我家的.
看看櫻水地.
既然如此.
不如試試改變一下原本的想法.
讓我試試吧.
當摩西想試試突破神之前的想法的時候.
神立刻很斷然拒絕了他.
這一個你會看到是編者很特別的安排.
因為本身摩西不能夠過河的技術.
其實比分地更早.
就是舊的其他的技術.
當你去想的時候.
可能你會覺得很奇怪.
摩西一件歸一件.
大家在說公事.
在說整個意識列文.
你為什麼加插了使徒在這裡呢?.
大家會不會覺得很奇怪?.
就是說以前的事.
然後轉接下去就是.
說要上帝如何教訓他.
為什麼生命的技術者摩西.
會放一個自己的圖文在這裡呢?.
他想表達什麼呢?.
其實當你看回整個申命記.
摩西做了很多復筆的.
他其實隱藏了很多句.
在一章,三章,四章和最後三十二章.
都說了很多次.
是因為你們這群意識列文.
你們的緣故導致我無法進入迦南地.

$^{881}$這就是你會看到在整個申命記裡面.
很多的復筆.
當然.
如果你看回整個技術裡面的時候.
你看回第一章的一章三十四至三十七節.
其實已經提過這個內容.
大概的內容.
我很快地讀給大家聽.
就是說什麼呢?.
耶和華聽見你們這番話就發怒.
氣勢就說.
這個惡世代的人.
連一個都不能夠見我所氣勢.
應許賜給你們列祖的美地.
唯有耶夫利的兒子迦納就看得見了.
並且我要將他所踏過的地方.
賜給他和他的子孫.
因為他專心跟從我.
三十七節說.
耶和華為你的緣故也向我發怒說.
必你必不得進入那地.
事後你的亂的兒子約書亞.
他就能夠進入那地了.
你要勉勵他.
因為他要使意識列人.
承受那個地違約.
這個技術就是在說十二個探子回來的時候.
只有兩個報好消息.
十個就說不要去了或者等等.
當時呢.
所以神就說.
只讓約書亞和迦納進入應許地.
但其實呢.
當你回到民數記十四章.
二十六至三十二節.
其實當時神沒有那麼明顯說不讓摩西進入.
他只是在整個技術就是說.
我要氣勢你們不能夠進入居住的地方.
只有迦納和亂的兒子約書亞能夠進去.
這就是民數記的技術.

$^{921}$但當在申命記裡面的時候.
摩西就再加一句.
就說因為你們的緣故.
所以我進不去.
但如果你再抽絲剝繭.
再找回整個脈絡的時候.
究竟摩西什麼時候知道自己進不去呢?.
在當時民數記他們十二個探子回來的時候.
其實摩西還不知道.
應該還不知道自己進不去.
為什麼在申命記裡面他再加一句就是說.
因為你們的緣故.
所以原來我也不行.
可能當時摩西還以為自己可以.
不只是是不是神說漏了.
因為你看到整個敘述.
他還在幫那群以色列人求.
不要啊.
因為神說我要興起你.
滅掉他們.
但摩西說不要啊.
其實他還在為那群以色列求.
摩西究竟什麼時候知道自己進不去迦納呢?.
就在民數記二十章.
就是米利巴水事件.
大家如果熟悉背景就是說.
因為在二十章十節就說.
摩西和亞倫就召集那些會眾來到盆石前面.
因為他們要水嘛.
喝水沒有水喝.
摩西就跟他們說.
聽著啊.
你們這群活逆的人.
我要叫盆石流出水來給你們.
接著摩西就舉起手.
用那個杖打了盆石兩下.
就有水流出來了.
會眾跟他們的生畜.
都能夠有水喝了.
但是耶和華就對摩西和亞倫說.

$^{961}$因為你們不相信我.
沒有在以色列人眼前尊我為聖.
所以你們不能夠領這群會眾進入我所賜給你們的地方.
民數記二十章才寫明.
摩西不能進入迦納地.
繼續在民數記二十七章.
他就再一次摩西領受上帝的話.
摩西當時領受上帝的話是什麼呢?.
第二十七章十二節就說.
你要上這個山去.
看看我賜給以色列人的地.
看了之後.
你就要歸到你的祖先那裡.
跟你哥哥亞倫一樣.
因為你在尋的曠野.
當會眾爭鬧的時候.
違背了我的命令.
在取水的地方.
沒有在會眾面前尊我為聖.
這兩個技術全部都比以前早.
當以色列人打西宏,打巴山.
兩個半支派分地.
這些技術都比以前早.
所以回到申命記裡面.
為什麼會這麼奇怪呢?.
分地.
說OK了.
突然加插了摩西.
可不可以讓我過河啊?.
不要這樣啊.
這樣的情況.
讓我過去看看.
當我們理解到摩西的心的時候.
可能我們又很快跳過去結論了.
因為摩西是最接近神的男人.
所以上帝對他要求最高.
所以一點錯都不可以.
你要尊爺和為聖.
一點錯都不可以.
我不是說這個教導不對.

$^{1001}$也是有他的可取地方.
但是聽人會怎麼樣呢?.
哇.
那我還是不要跟神那麼近.
我還是跟神遠一點.
就是不知道大家.
以前的下令會,培令會.
有沒有弟兄跟你說.
我不是很想祈禱.
但時一期.
神太近我.
我還是跟他有一個距離.
結論就是因為摩西跟神很近.
所以他不能夠進入迦蘭地.
這個結論好像是不希望.
讓你不想再走一步.
多於讓你想再走一步.
所以我相信這個不會是當時的結論.
不會是因為想讓他腳步.
就不要走了.
我們不如留在這裡吧.
所以當我們看回的時候.
如果這個結論只是讓你腳步.
不是他的心意.
更何況其實整個生命的最後.
對摩西的認同是很高的.
他說沒有一個先知好像摩西一樣.
跟神是面對面所認識的.
可能就留一個問題給大家再想想.
為什麼要將這件事放在中樞位.
說完那時.
去到現在的時候.
就記述摩西這個祈禱.
這個詩禱.
成為一個很重要的中間位.
而且你要知道.
摩西在生命裡面不斷重複.
是因為你們的緣故.
如果你看回剛才民數記不是這麼說的.
是因為你不尊爺和為.

$^{1041}$為什麼摩西現在又會說.
是因為你們的緣故.
他進不去呢?.
當你繼續去想想裡面的內容的時候.
你留意一下一些細節.
當你看回第三章二十六節裡面的時候.
聖經裡面說什麼呢?.
第三章二十六節.
「但耶和華因你們的緣故.
向我發怒,不應允我.
對我說:罷了,你不要向我再提這事.」.
這個罷了.
可能我們一聽的時候覺得很絕情.
上帝連摩西再說一次.
這個詩禱都不要說了.
夠了,停了.
但當你再看回這個罷了的時候.
如果你看其他譯本.
或者和修本.
或者新漢語譯本等等.
他說你夠了.
摩西夠了.
在這裡夠了.
這個夠了.
不只是說那種嚇著他.
停了.
你說夠了.
原來這個夠了的用字.
是和一章六節.
二章三節.
說這群以色列人.
在曠野的地方.
在山上的地方.
夠了的字是一樣的.
就是說這些日子都夠了.
這些狀況都夠了.
你要轉變.
當時一章六節說.
你們在山上住的日子夠了.
要轉起行.

$^{1081}$去到上帝要帶他們的地方.
二章三節.
說你們在這個山上.
要住要行的日子都夠了.
現在要轉向北去行.
這個夠了不只是說.
你停了.
不讓再說了.
沒有後來了.
然後就沒有那種後來的夠了.
而那個夠了是要轉的.
摩西當聽到上帝說夠了.
不要再提這件事的時候.
是他在這件事裡面.
停留的時間夠了.
你轉去山上.
去看看上帝的應許.
要怎樣成就.
摩西一直在困惑的.
就是這種方式.
才是成就應許的象徵.
或者是表達.
神說不要再停留在這種狀況.
轉變過來.
你看看.
這一群新的以色列人.
將要進入這個地方.
而你看見這個應許的地方的時候.
已經足夠.
已經知道上帝的應許.
將要成就.
當上帝說夠的時候.
就是要摩西.
你要變.
你在你糾結的事裡面夠了.
不要再停留在這個階段.
不要再停留在這種.
你覺得要看見去踏足.
切身處地走過去.
你才能夠覺得滿足的.

$^{1121}$這種心態已經夠了.
用你的信心去看.
上帝覺得在你的生命裡面.
已經足夠了.
當我們看前面的事的時候.
可能只是覺得.
這件是摩西的遺憾.
一邊看.
哇.
一個這麼忠心的僕人.
一個這麼好的.
只是一點點錯.
這樣上帝就不給很大的遺憾.
甚至你會覺得.
這個遺憾.
是摩西你做錯了一點.
但不至於是這樣吧.
當你會想的時候.
哇.
他在擔當這個百姓.
再加上這個遺憾.
他會覺得.
哇.
那怎麼走下去呢?.
上帝就是想在這一刻跟他說.
夠了.
轉變下去.
將你的期盼去轉化.
成為力量繼續走下去.
如果他知道.
這個新的開始.
原來是在這種新的轉化裡面.
摩西也轉變了.
能夠繼續仰望上帝的大靈.
當你看回這段的技術的時候.
我不知道大家會不會覺得.
一開始的時候覺得.
好像上帝有兩把尺似的.
為什麼對著那兩個半支牌.
又讓他試.

$^{1161}$又讓他這樣去做.
對著摩西.
又好像堅決地.
不讓他這麼做呢?.
為什麼不願意改變他的遺憾.
好像那兩個半支牌.
讓他去看看吧.
究竟想教導以色列人什麼呢?.
其實這個狀況是很吊詭的.
正如今天的主題.
還有你會看到.
今天的大綱.
掛著很多字.
那些字眼.
你一邊看的時候.
覺得好像很相對.
一邊對著他寬容.
但是另一邊又對著他堅決.
有些東西應該承傳.
有些東西應該斷捨.
有些東西應該改變.
有些東西不應該改變.
好像本身兩樣東西是兩邊的.
但是最吊詭的地方就是.
原來這些所謂的極端.
都是並存在我們的生命裡面.
而上帝的教導也是一樣.
有時候好像很矛盾.
但是又並存.
原來這種不是非此則彼的信仰.
不是這樣做就一定是好的.
你只要全部人做完這件事.
就一定是好的.
有時候又不是的.
為什麼會這樣呢?.
原來就是摩西藉著這個技術.
很想教導這些新的以色列人.
藉著這種張力.
這種矛盾.
認識什麼是真正的上帝.

$^{1201}$當我們去教導認識上帝的時候.
跟一般的教育不同.
就好像我們在查經等等的時候.
教導上帝.
對上帝認識.
有時候教育或者教導.
很多時候重視的就是.
我們如何掌握那個技能.
掌握那個理論.
掌握這種科學或者人民的心理的時候.
就變成了一條道路.
讓我們有一種力量或者權力.
去管治我們正在學習的事情.
知道的事情.
知識上知道的事情.
經常有一句話就是.
知識就是力量.
我知道了多一點.
我學會了多一點.
知道這件事情多一點.
我就可以改變了它.
認識上帝對你來說.
會不會是這樣呢?.
認識上帝多一點.
我就可以知道.
如何在神裡面得到更多福氣.
認識上帝.
我就可以在祂身上.
得到祂的賜福.
幫助更多.
這個認識的心態本身.
也是有一個問題在這裡.
因為上帝不是一條方程式.
不是說做這件事.
上帝一定賜福我.
所以以色列人就聰明了.
做這些.
我過河就行了.
主權不是在學習那個人.
如果我們相信上帝是一個有慧格.

$^{1241}$有性情.
是互動的.
跟我們互動的神.
不是只是一條.
世界的原則.
一條方程式的時候.
祂有祂的意識和回應.
以色列民的方式.
祂的應許也不是草引花.
你草得越多獎賞越大那一種.
祂的應許是跟祂的生命有一種互動.
你的生命跟上帝的生命.
而產生的互動.
而摩西藉著這個技術.
就是很想以色列民知道一件事情.
你認識上帝.
為的是敬畏上帝.
你認識上帝越多.
你敬畏的心就會越大.
而敬畏最需要的知道就是.
所有的大權.
都是上帝的歌徒.
所以在這個故事裡面.
你會看到好像.
神給又不給的時候.
因為好像看起來很對立.
但其實主權.
回到都是在上帝的大權裡面.
祂是按著祂的大圖畫去建立.
和醫治這個群體.
我們深信.
祂不是破壞和傷害.
每一個屬神的人.
和屬神的群體.
所以認識神.
讓我們能夠有這份敬畏的心.
知道上帝是賜予的那個.
我們是蒙神愛.
而讓我們能夠踏在這段路裡面.
跟從神的人.

$^{1281}$即使有時候.
可能我們也會想.
好像摩西那樣.
很不公平而已.
摩西你真的是值得擁有的.
甚至可能加多一句.
你看看上帝多麼對你多麼不好.
這樣對你.
給你這些事.
摩西不給你面子的.
或者調回自己.
喂上帝.
我值得擁有的.
即是英文有句話.
You deserve it.
當然兩看.
你值得還是你活該.
即是究竟你值得擁有.
我真的做了這麼多.
我值得的.
為什麼神不給我.
我做到今天.
我覺得也OK.
為什麼你也不讓我經歷.
改變我的遺憾.
上帝.
我以前這些經歷.
可不可以抹走它.
或者改變它.
為什麼讓我仍然停留在這個遺憾裡面呢?.
這個停留.
如果只是讓你停留在遺憾.
也不是上帝的轉化.
如果你只是視這件遺憾.
就是反映了上帝對你的一種認識.
或者你對上帝的一種認識.
也是局限著你和上帝的關係.
而今天你和上帝之間.
有沒有讓你覺得是遺憾.
或者好像過不了的事卡住了.

$^{1321}$停留了.
不能踏向前.
令你覺得.
一是兩個方面.
哎.
我本身可以這樣.
是上帝不好.
上帝你不好.
不讓我這樣前行.
如果你讓我,我就走.
你不讓我,我就不走.
還是好像摩西一樣.
過往的事是這樣.
轉化.
令你繼續前行.
去move on(前進).
去繼續向下一步.
放下你的顧慮.
接受你的遺憾.
去改變前面的路.
歷史已經過去了.
我知道你很努力想改變歷史.
我知道你很想努力.
去改變以前有的遺憾.
但是.
放下它.
夠了.
是時候踏上這個新的路途.
按著上帝的心意.
過去已經定了.
今天你每一個決定.
都在取決你明天.
如何繼續經歷神.
當你知道.
好像摩西身上.
他不能夠進入迦南這個地.
不代表他已經對神灰心.
死心.
他仍然靠著神去轉化.
祝福下一輩.

$^{1361}$讓他們在這些困苦掙扎裡面.
去走前面的路.
當你今天處身在你的生活裡面.
每逢回想以前.
你會覺得.
啊 上帝你的大靈好.
還有啊 耿耿於懷.
有些事情.
上帝為什麼你不去改變.
夠了.
轉變過來.
讓這一切成為上帝和你的足跡.
再去展望下一個階段.
對準你對上帝的認識.
為什麼有這些過往.
為什麼有這一刻的你.
再展望你的應許地.
是什麼呢.
到最後有時候.
其實我們真的在這些弔詭狀況裡面.
很大的困難.
有一個圖文.
是成為我自己同樣在這些張力或者矛盾裡面.
很大的幫助.
是一個大家可能認識的寧靜圖文.
一個神學家尼布爾在1933年裡面.
他三句的圖文就是說.
請賜我寧靜.
去接受我沒辦法改變的事情.
請賜我勇氣.
去改變.
我能夠改變的事情.
請賜我智慧.
以分辨兩者的不同.
在圖文裡面好像常常說.
改變不改變.
你怎麼知道呢.
沒有人知道.
唯獨是你回到上帝求智慧.
去分辨究竟這件事.

$^{1401}$臨的心意.
上帝你怎麼成為祝福呢.
盼望我們能夠從過往的事裡面去學習.
立新的志向.
砥礪前行.
我們一起祈禱.
親愛的主祈求你帶領著我們每一個.
幫我們能夠在每一件事裡面.
看到你的心意.
有些事情我們可能看似.
是上帝你缺席.
離開.
甚至是一大遺憾.
但我們深信.
正如摩西的身上所經歷.
這是一個中樞或轉捩.
我們要放下這一切.
在你的懷抱.
你的同行裡面.
繼續前行.
以致看見上帝你.
繼續讓我們朝向這個英雄地裡面的時候.
我們如何緊緊抓著你.
繼續前行.
我們祈禱.
是奉主耶穌基督的名字而求.
阿們.
謝謝大家.
\newpage



\section{申命記 4:1-14}
\label{sec:cuIe7N8imyc}
\textbf{2025年2月2日主餐崇拜直播｜陳明泉牧師｜朝向應許之地－安居樂業的願景｜申命記4\_1-14}
\newline
\newline
連結: \href{https://youtube.com/watch?v=cuIe7N8imyc}{\texttt{https://youtube.com/watch?v=cuIe7N8imyc}} ~~~~ 語音日期: 2025-02-02
\newline
\newline
\hyperref[sec:KYvJDDo3N1U]{\small{< < < PREV SERMON < < <}}
~
\hyperref[sec:index_chronic]{\small{[返順時目]}}
~
\hyperref[sec:uDlv_6lnHKk]{\small{> > > NEXT SERMON > > >}}
\newline
\newline
申命記 4:1-14
\newline
\begin{longtable}{cl}
\hline
\hline
章節 & 經文 (和合本修訂版)\\
\hline
4:1 & \begin{tabularx}{0.7\textwidth}{X} 「現在，以色列啊，聽我所教導你們的律例典章，要遵行，好使你們存活，得以進入耶和華－你們列祖之神所賜給你們的地，承受為業。 \end{tabularx} \\ \\ \relax
4:2 & \begin{tabularx}{0.7\textwidth}{X} 我吩咐你們的話，你們不可加添，也不可刪減，好叫你們遵守耶和華－你們神的命令，就是我所吩咐你們的。 \end{tabularx} \\ \\ \relax
4:3 & \begin{tabularx}{0.7\textwidth}{X} 你們已親眼看見耶和華因巴力‧毗珥所做的。凡隨從巴力‧毗珥的人，耶和華－你的神都從你中間除滅了。 \end{tabularx} \\ \\ \relax
4:4 & \begin{tabularx}{0.7\textwidth}{X} 只有你們這緊緊跟隨耶和華－你們神的人，今日全都存活。 \end{tabularx} \\ \\ \relax
4:5 & \begin{tabularx}{0.7\textwidth}{X} 看，我照著耶和華－我的神所吩咐我的，將律例和典章教導你們，使你們在所要進去得為業的地上遵行。 \end{tabularx} \\ \\ \relax
4:6 & \begin{tabularx}{0.7\textwidth}{X} 你們要謹守遵行；這就是你們在萬民眼前的智慧和聰明。他們聽見這一切律例，必說：『這大國的人真是有智慧，有聰明！』 \end{tabularx} \\ \\ \relax
4:7 & \begin{tabularx}{0.7\textwidth}{X} 哪一大國有神明與他們相近，像耶和華－我們的神在我們求告他的時候與我們相近呢？ \end{tabularx} \\ \\ \relax
4:8 & \begin{tabularx}{0.7\textwidth}{X} 哪一大國有這樣公義的律例典章，像我今日在你們面前所頒佈的這一切律法呢？ \end{tabularx} \\ \\ \relax
4:9 & \begin{tabularx}{0.7\textwidth}{X} 「但你要謹慎，殷勤保守你的心靈，免得忘記你親眼所看見的事，又免得在你一生的年日這些事離開你的心，總要把它們傳給你的子子孫孫。 \end{tabularx} \\ \\ \relax
4:10 & \begin{tabularx}{0.7\textwidth}{X} 你在何烈山站在耶和華－你神面前的那日，耶和華對我說：『你為我召集百姓，我要叫他們聽見我的話，使他們活在世上的日子，可以學習敬畏我，又可以教導他們的兒女。』 \end{tabularx} \\ \\ \relax
4:11 & \begin{tabularx}{0.7\textwidth}{X} 那時，你們近前來，站在山下；山上有火燃燒，直沖天頂，並有黑暗、密雲、幽暗。 \end{tabularx} \\ \\ \relax
4:12 & \begin{tabularx}{0.7\textwidth}{X} 耶和華從火焰中對你們說話，你們聽見說話的聲音，只有聲音，卻沒有看見形像。 \end{tabularx} \\ \\ \relax
4:13 & \begin{tabularx}{0.7\textwidth}{X} 他將所吩咐你們當守的約指示你們，就是十條誡命，並將誡命寫在兩塊石版上。 \end{tabularx} \\ \\ \relax
4:14 & \begin{tabularx}{0.7\textwidth}{X} 那時，耶和華吩咐我將律例典章教導你們，使你們在所要過去得為業的地上遵行。」 \end{tabularx} \\ \\
[1ex]
\hline
\hline
\end{longtable}
$^{1}$以下跟大家讀出的經文是在舊約《新命記》第四章一至十四節.
是在大家手上的聖經舊約的218至219頁.
聖經這樣說:.
「以色列人啊,現在我所教訓你們的律例殿章.
你們要聽從遵行,好叫你們存活.
得以進入耶和華你們列祖之神所賜給你們的地.
承受遺孽.
所吩咐你們的話,你們不可加添,也不可刪減.
好叫你們遵守我所吩咐的.
就是耶和華你們神的命令.
耶和華因巴力比爾的事所行的,你們親眼看見了.
凡隨從巴力比爾的人,耶和華你們的神都從你們中間除滅了.
唯有你們專靠耶和華你們神的人,今日全都存活.
我照著耶和華我神所吩咐的,將律例殿章教訓你們.
使你們在所有進去得遺孽的地上遵行.
所以你們要謹守遵行.
這就是你們在萬民眼前的智慧聰明.
他們聽見這一切律例必說.
這大國的人真是有智慧有聰明.
哪一大國的人有神與他們相近.
像耶和華我們的神在我們求告他的時候與我們相近呢?.
由哪一大國有這樣公義的律例,殿章.
像我今日在你們面前所陳明的這一切律法呢?.
你只要謹慎,殷勤保守你的心力.
免得忘記你親眼所看見的事.
又免得你一生只是離開你的心.
總要傳給你的子子孫孫.
你在可列山站在耶和華你神面前的那一天.
耶和華對我說:你為我焦聚百姓.
我要叫他們聽見我的話.
使他們存活在世的日子.
可以學習敬畏我.
又可以教訓兒女這樣行.
那時你們近前來站在山下.
山上有火焰沖天.
並有昏黑,密雲,幽暗.
耶和華從火焰中對你們說話.
你們只聽見聲音,卻沒有看見形象.
他將所吩咐你們當守的約指示你們.
就是十條戒.

$^{41}$並將這戒寫在兩塊石板上.
十四節.
那時耶和華又吩咐我將律例,殿章教訓你們.
使你們在所有過去得遺孽的地上遵行.
丁姐妹新年快樂.
在新的一年,恭祝大家身體健康.
生命更上一層樓.
今日我們讀這段新命記.
是關於祝福及重要願景的經文.
在我們理解這段經文前.
有一個片段.
早幾天我在澳洲回來.
在澳洲回來時.
飛機機師通常會有一段廣播.
將要降落.
降落前,機師講完所有資料後.
用英文說「恭喜發財」.
然後整個飛機的人都在笑.
有一個虛構片段.
小朋友聽到「恭喜發財」.
小朋友問爸爸.
爸爸,為何平時你跟我說「恭喜」.
你沒有說「恭喜我發財」.
小朋友問是有道理的.
我們街上廣東人或香港人.
一見到人,新年打招呼就是「恭喜發財」.
但不會見爸爸和兒子「恭喜發財」.
一個爸爸聽到這句說話.
爸爸很有智慧地回答.
兒子,其實你人生需要的不是發財.
你人生需要的是智慧.
所以你要讀書聰明伶俐.
繼續有好的成長.
這就是爸爸對你的祝願.
是一個這樣的小片段.
不過你會見到.
其實一般在街上的祝福語.
和有關係緊密的祝福語.
兩個內容具體是有不同的.
有關係的祝福語會杜毀一些.

$^{81}$和知道對象的需要.
一般來說,你不知道對方.
就會說一些比較行.
或是世界價值所追求的祝福語.
你用這個角度來看.
今天在《申命記》第四章.
你就會明白摩西所說的這番說話.
是一個度身訂做.
為新一代的以色列子民.
來送上的一個祝福和一個糞便.
希望今天我們都拿了這段經文.
由《申命記》第四章一至十四節.
我們理解一下將要進入.
去得地為業.
將要安居樂業的這班以色列人.
摩西用什麼樣的勸勉來祝福他們.
我們先祈禱求上帝幫助.
上主願你幫助我們.
讓我們在新春的日子.
我們思想上帝你賜下的福氣.
又讓我們今天.
我們一同頂展會來聚集.
我們不單歡慶相聚.
也一同領受你的話語.
求你賜我們心靈悟性.
一起來明白你對我們所說的話.
也賜福給孫子說的.
讓孩子所說的話忠於聖經.
按著真理來頂展會領受.
讓我們在生活裡一起實踐.
恭敬將來的時間交託.
願你賜福.
獻上禱告奉求基督耶穌得勝的聖名.
Amen.
今天我們看的第四章一至十四節.
我們分為三個段落.
第一個段落是第一節至第四節.
這是整個申命記裡面.
一個很重要的神學觀念.
就叫做新典神學.

$^{121}$我先讀一段經文來聽.
以色列人.
現在我所教訓你們的律例典章.
你們要聽從進行.
好叫你們存活.
得以進入耶和華你們列祖之神.
所賜給你們的地.
承受為業.
所吩咐你們的.
說你們不可加添也不可刪減.
好叫你們遵守我所吩咐的.
就使耶和華你們神的命令.
耶和華因巴力比爾的事所行的.
你們親眼看見了.
凡隨從巴力比爾的人.
耶和華你們的神.
都從你們中間除滅了.
唯有你們專靠耶和華你們神的人.
今日全都存活.
在這裡其實是在說一個因果關係.
這個因果關係就是貫穿了舊約裡面.
由申命記到列王記裡面.
這幾個篇幅裡面.
幾個聖經裡面.
一個很重要的神學觀念.
就叫做新典神學.
新典神學在說什麼呢.
就是說在申命記裡面.
摩西將一個上帝與子民關係的賜福和和幻.
這個生命裡面的因與果.
就表達了出來.
其實這個因與果.
建基於一個盟約的邏輯.
上帝與他的子民立約.
上帝對於立約的子民.
他是有要求的.
同樣立約的子民.
因為與上帝有一個這樣的立約關係.
上帝就將一個應許之地.
應許的福氣.

$^{161}$贈予.
送予給以色列民.
所以這是一個盟約的關係.
不過說盟約或者說約.
很多時候我們會有些誤解.
我們想那個約.
就好像我們今天所做的房地產買賣一樣.
我們想著那張合約.
我們覺得就是買賣雙方.
一個買一個賣.
然後大家履行利益的責任.
大家出個價.
然後互相交換.
可能做商品交易.
或者買樓買地.
一買一買的關係.
我們想這個就是一個約的關係.
這個是叫做contract.
就是說一些商業的合約.
不過上帝與人之間的約的關係.
就不是這樣的一買一買的關係.
在聖經裡面說到.
上帝與人的契約.
盟約就叫做covenant.
這個covenant強調的是.
上帝主動的.
來跟他的子民.
來做一些恩典的表達出來.
這一種的約背後有一種神聖和永恆性的.
一般的商業合約.
有一個時限的.
會有停的.
就是買賣交易停了就停了.
你買一層樓.
你不會再跟上帝主還有很多的交往.
那個利益就是那次的交易就停止了.
完了之後.
但是去到covenant.
這個關係是一個傳到永遠.
和神聖的一份的關係.

$^{201}$而且那個約的建立.
是有一種動力在裡面的.
這種的動力就不是跟一般商業合約.
所說的利益成為了動力.
愛成為了covenant.
上帝神與人之間立約的那種動力.
上帝因為愛人.
所以用這個約的方式.
既是規限住自己.
也都是允許將他所給的福.
給以色列人.
同時間人就履行在這個契約裡面.
對神的那份的效忠和那種信實.
所以整個的covenant.
神人之間的那個契約.
強調的就是一份的忠誠關係.
如果你覺得說這個都好像很抽象.
其實在聖經裡面.
說神人之間的契約之外.
有三個很重要的字眼.
都在人世間裡面的契約.
都會出現的.
就是離開,聯合,合一.
很熟悉.
這個就是創世記裡面.
人要離開父母.
與妻子聯合.
二人成為一體.
這裡的用字.
其實在舊約裡面說的.
這幾個用字就是一個人與人之間的.
那種盟誓.
一生一世的.
那種盟約的納誓.
這個放在婚姻裡面.
同樣這個也成為了神人之間的類比.
神與人之間.
上帝無條件地將福氣賜給那班以色列民.
同樣祂期望.
納了約之後的那些以色列民.

$^{241}$生命裡面好像納了約之後沒有東西.
不是的.
祂期望他為了回應這個舊屬的恩典.
回應了這個這麼豐盛恩惠的生命.
所以人的行為.
人的生命都隨著這個約.
來塑造了他的行為表現.
而在聖經裡面所說的.
特別是在說摩西這一段的經文.
所說的律法.
就是在說立約之後的子民的生命.
是一個豐盛的生命.
配得上上帝與你立約的福氣.
你也要回應上帝的時候.
你如何去表達.
你的生命是與上帝立約的呢?.
律法就是一個很重要的藍圖.
告訴那些以色列民.
告訴每一個神國子民.
安著上帝的心意好好地去生活.
以致他能夠經歷上帝賢尼的豐盛.
在剛才所讀的經文裡面.
還有一個很重要的字眼.
其實是不停提醒你的.
他告訴那些子民.
今天教訓你的律例.
將你們要聽從進行.
好叫你們存活.
在第四節裡面.
唯有你們專靠耶和華你們的神.
今天全都存活.
那個存活的意思.
不單止是生存.
還要活得到.
能夠屹立在當時進入一個新的地方.
充滿了異教的文化.
人與人之間有不同的利益私利.
不過不同的部落.
不同的文化.
會隨著歷史一一改變.

$^{281}$在摩西裡面.
你見到有一些拜巴力的.
到最後他們被消滅.
但你們這班.
以前在埃及當奴隸的人.
現在你們可以經過曠野四十年.
而且還可以開枝散葉.
還有第二代.
還有一個越來越好的願景.
因為上帝在我們當中.
上帝與他的子民立約.
所以那個存活是說.
因為守約的結果.
而帶來的一些美好的果效.
更加重要的就是.
同時間也告訴那些以色列人.
這個與上帝進行立約的遵守誡命.
這種生命表現.
是一個生與死的決定.
你不是想著錦上添花.
有些事情可以聊聊.
然後你就會覺得.
過得好一點.
而是說.
是一個好.
是一個生存和滅亡的選擇.
是一個生命與死亡的.
一個很重要的選擇.
摩西將這一段的神學觀念.
將這個上帝立約的邏輯.
告訴那些以色列人.
將這樣的福與禍.
放在他們的面前.
這一代的以色列子民.
他要怎樣選擇呢.
就將這個球交給他.
他不能強迫他們一定要去做.
因為律法從來都是由心底裡.
來認同的.
但是當你去遵從.

$^{321}$你就會經歷到那份豐盛和美滿.
當你拒絕.
或者你去搏弈的時候.
你面對的就是.
好像那些外邦.
巴力的那些外邦一樣.
就得不到那種永恆的祝福和生命.
當我們明白了.
這個神學觀念.
這個新典神學的觀念的時候.
接下來.
摩西就會將一個.
進入去迦林地裡面一個美好的願景.
告訴那些以色列人知道.
由第五節到第八節.
摩西就繼續用第一個願景.
來告訴他們.
不過這個願景很有趣.
這個願景進入迦林地.
不是說那些以色列人進去的時候.
你就會衣食無憂.
或者很有富有.
他不是將這個東西放在他們面前.
摩西用什麼來說一個.
真正的願景給那些.
昔日為奴隸的一群以色列人呢.
我們聽聽第五至第八節.
我照著耶和華我神.
所吩咐的律例.
典章教訓你們.
使你們在所要進去.
得遺孽的地上遵行.
所以你們要謹守遵行.
這就是你們在萬民眼前的智慧.
聰明.
他們聽見這一切律例必說.
這大國的人真是有智慧有聰明.
那一個大國的人有神與他們相近.
將耶和華我們的神.
在我們求告他的時候.

$^{361}$與我們相近呢.
又那有一個大國.
有這樣的公義律例典章.
將我們今天在你面前.
所陳明的這一切律法呢.
在這裡你見到.
有一個很重要的.
類似一個反問句.
但是又重複出現的.
就是有哪一個大國.
說的是一個國家的強盛.
有哪一個國家這麼強盛到一個地步.
被其他外邦人稱讚.
今天我們說的.
那個大國的觀念.
曾經有一句話這樣說的.
一個國家或者一個地方的強大.
看什麼呢.
看它的建築物.
如果那些建築物很宏偉的.
而且那個設計很美輪美奐的.
那個呢.
那個有一定的實力.
那個地方那個國家.
就是一個強大的國家.
但是有一點文化學家說.
一個地方的建築.
令到一個地方變得大.
或者變得強大.
但是真正令到一個民族.
變得偉大的.
是只有文化和那些人民生命的素質.
所以在這裡面.
摩西所說的那個大國.
哪有一個大的國家.
它不是說有一些宏偉的建築.
或者有一些很強的軍事實力.
它說的是那個大那種偉大.
體現於幾樣東西的.
第一.

$^{401}$就是說到那些人民行事.
或者在萬民當中.
裡面有智慧.
有聰明.
那些的國家民族.
能夠成為一個強大.
就能夠顯現出.
在眾人在混亂的世界裡面.
他們的生命的智慧聰明.
是比其他國家來的高.
我們一想這些就想IQ.
想他們的那些創意.
可能都是的.
不過這裡摩西所說的那個智慧聰明.
就不是說人世間裡面.
說的那些IQ.
或者那些創造能力.
這個是說.
他們能夠承接著上帝給他們的那種關係.
以致他們整個文化.
他們真的開竅.
他們真的明白上帝.
明白這個世界裡面的走向.
而能夠掌握和上帝之間有一個緊密的關係.
不跟這個世界的潮流去走.
在《箴言》裡面說到.
驚威而不是智慧的開端.
在說那個智慧和敬畏上帝.
永遠都是畫一個等號在這裡.
而這個智慧聰明.
往往就是跟其他的民族.
其他的邦國比下去的時候.
就會有一些很不同的那種看法.
簡單來說.
神國的子民.
想事物.
做決定.
是會成熟一點.
他的成熟不是代表著.
因為這個民族那些DNA.

$^{441}$那些人的IQ零碩高.
而是因為上帝教他們.
上帝幫他們.
以致他們真的每一個人都能夠開竅.
如果盲目地遵守律法.
純粹說是這麼做的.
律法是這麼叫的.
這樣就很糾扯.
或者帶著被迫的心態.
來實踐這些的律法.
是得不到智慧的.
但是在聖經裡面.
你要留意.
上帝說的律法.
從來不是純粹說一條規條的.
他會解釋給你聽.
他也說到上帝的創造心意.
上帝的旨意是一個怎樣的旨意.
所以背後是會有解釋.
有因果關係.
上帝在當中也會指教.
令到人會繼續思想.
所以在這裡說的智慧聰明.
是說上帝不單止將一個律法裡面.
所說的界線.
停止了.
不要這樣做.
或者跟這個世界分別出來.
跟這個世界的價值觀有所不同.
上帝更加將他的價值觀念.
上帝將一些他的想法.
人與人在地上裡面.
美善的道德旨意.
他會指教給那些以色列民.
那些以色列民一路實踐.
一路思考的時候.
他生命裡面有所提升.
所以在這裡.
第一個他體現的.
就是說到他的大國.

$^{481}$就是說到這班人有智慧有聰明.
第二件事他繼續說.
這班國民跟其他世界的神明是不同的.
其他的神明.
或者拜巴力的那些.
當時社會的那些.
他們敬拜他們的神.
最重要的動力是來自利益.
不過是出產的.
他拜的那些巴力.
或者說的舊時裡面.
去敬拜的那些偶像.
往往都是因為那些偶像.
有一種形態就是賜給你有豐收.
賜給你生活無憂.
又或者有一些世界很想追求的某些價值觀念.
就放了成為了人的圖騰或者偶像.
人就去敬拜這些偶像.
來滿足心靈或者現實裡面.
沒辦法達至的一些理想.
不過在聖經裡面說到神人之間.
如果說到是一個納若的關係.
神與人之間.
就是在說一種互動的關係.
上帝跟人是不停的來互動的.
這種的納若不是靜態.
這種的納若.
這種的實踐.
是整個民族和每一個人生裡面.
一路實踐一路去思考的時候.
一路會豐富起來.
一路明白上帝的旨意.
所以這個神跟那些子民是很相近的.
是很近距離.
而不是很有距離感.
你找東西來賄賂他.
或者用一些東西來爭取更多的利益.
第三個說到那些律例典章.
那些是實踐起來.
到最後是能夠有公義.

$^{521}$是令到每一個人都心悅誠服的.
曾經有一個AO.
就是政府的那些官員.
聊天.
他說你不要想我們做官員好好做.
你想一條法例.
想到頭都爆了.
想到頭髮都掉了.
因為要做一條.
能夠滿足社會各界.
都能夠覺得是合宜的.
公義的.
在過程裡面實在太難.
一隻字.
如果表達得不清楚.
都會引起很多的誤會.
或者就會引起了世界社會上.
都有很多的混亂.
所以他們說那條法例.
那些律例是怎樣去執行呢.
是一個非常之大能度的工作.
我不知道.
我們想官員只是出去解說去說.
但背後裡面的努力是沒有人知道的.
用這裡.
在這裡所說的.
是外邦說那些以色列人.
因為他們實踐了一個很公義的律例典章.
這個公義的典章.
是超乎人的智慧可以想像的.
也都杜絕了.
也都看盡了人性的那種本質.
以致到這些律例典章實踐起來的時候.
每一個子民都能夠心悅誠服.
再加上上帝是掌管萬物的主.
他有超然的威權.
那個的看法一定是超越地上人間裡面的作品.
所以在這裡裡面.
那些外邦見到以色列人.
就從聰明智慧.

$^{561}$上帝與他們親近.
以公義來描述整個以色列人.
因為守律法.
因為他專心遵循上帝的教導.
以致到在其他外邦中.
他們就可以鶴立雞群.
被人稱為一個強大有力量的民族.
在這裡我們稍為思想一下.
我們今天很多時候.
我們想.
我們讀聖經.
我們參與教會的聚會.
我們會覺得是不是真的聽道理.
是不是真的這樣.
不過如果人世間的道理是這樣聽.
你會聽到厭煩.
但如果那個道理是來自上帝.
當我們第一手去看上帝的話語.
我們真的去思考上帝給我們的教導的時候.
你再深入去想多一層.
也都去再和你自己的人性作一些對比的時候.
你會發覺上帝所說的話是非常有睿智.
以及超越了人世間所有的宗教哲學.
或者其他人所說的那些覺得很高尚的.
很高深的道理.
上帝的道理.
上帝的律法既是淺白.
但同時間影響卻是深遠.
在我們每一天.
我們說有讀經計劃.
我們去讀聖經的時候.
如果我們帶著有一個目標.
這樣都是好的.
不過如果只是說.
好像有些大家做一些壯舉.
逼自己去做一些事.
你是體會不到讀聖經的樂趣.
但如果我們每一次讀聖經.
我們就好像尋寶一樣.
淘寶吧.

$^{601}$淘寶大家會容易明白一點.
現在大家沒有人去尋寶.
淘寶會多一點.
你淘寶一樣.
淘出一些好東西.
淘出一些金.
你就會發覺.
這些東西不是人世間的產物.
是唯有上帝來作對我們生命的提醒.
所以弟兄姊妹.
即日那些以色列民將要進入迦南地.
摩西一而再再而三告訴他們.
就是因為講到上帝的律法.
當你熟讀了它的時候.
你的生命.
你會跟別人是很不同的.
你們的生命會比別人.
會有一些超然一點.
會超越一點.
如果你的生命裡面.
只是在人世間的那些生活當中.
來生存的話.
你沒有辦法過一個超越世界.
也能過一個有智慧的生活.
有一個很慚愧的經歷.
話說有一次我跟子女.
幾年前一起去到一個商場拍車.
就遇到一個中學的同學.
那個中學同學沒有信主.
中學畢業之後.
他就過自己的生活.
當時他見到我.
熟練了.
給我多一小時拍車的時間.
便便宜了我.
我的子女問.
爸爸為什麼大家都在同一間中學.
為什麼大家的形態好像很不同呢.
我不是自傲.
我只是說.

$^{641}$如果爸爸沒有認識主.
可能我也是差不多.
我也未必能夠.
生命裡面有不同的體會.
我在那裡跟子女說.
大家受的教育一樣.
但因為有上帝.
生命我們走在有福氣那裡.
而不會成為一個煙不離手.
或者生命裡面有很多的惡習.
或者生命有很多價值觀.
很混亂的一個人.
就因為爸爸認識了信仰.
有上帝與我們同在.
所以整個生命就改變了.
我要說的不是.
我們要做一些自誇.
我們每一個其實本質都是一樣的.
如果按照這班以色列人.
他們的本質本身都是一班奴隸.
如果在埃及地裡面.
為什麼這班奴隸會成為一個大國呢.
不是因為這班人有多努力.
DNA有多強盛.
單單是因為上帝與他們立約.
上帝教導他們.
幫助他們.
數十年之間.
讓他們成為一個散門的奴隸群體.
成為一個幫國.
成為一個承受應許.
神國指紋的一個代表.
所以在這裡面.
摩西將這個願景.
告訴新一代的以色列人.
叫他們學校上帝的律法.
遵行緊守.
以至他們的生命成為萬國羨慕的國家.
我們看最後一個段落.
第九至十四節.

$^{681}$你們要謹慎.
因勤保守你們的心靈.
免得忘記你親眼所看見的事.
又免得你們一生.
這是離開你的心.
總要傳給你的子子孫孫.
你在荷列山站.
在耶和華女神面前的那一天.
耶和華對我說.
你為我焦罪百姓.
我要叫他們聽見我的話.
使他們存活在世的日子.
可以學習敬畏我.
又可以教訓兒女者讓行.
那時你們近前來站在山下.
山上有火焰沖天.
並有昏黑,默暈,幽暗.
耶和華從火焰中對你們說話.
你們只聽見聲音.
卻沒有看見形象.
祂將所吩咐你們當手的約.
指示你們就是十條界.
並將這界寫在兩塊石板上.
那時耶和華又吩咐我.
將律例典章教訓你們.
使你們在所有過去得違約的地上進行.
最後一個段落是說什麼呢?.
這個願景的實踐.
它訴諸於昔日在西乃山裡立約的片段.
最重要的字眼.
就是放在中間的學習敬畏.
學習敬畏首先要告訴以色列新一代.
這是整個歷史民族的經驗.
他們的上一代因為進入曠野時.
他們真正敬畏耶和華上帝.
所以到最後他們進入不了應許之地.
他們整個民族都不斷學習中.
要繼續學習如何敬畏耶和華上帝.
所以訴諸的第一點是.
你們新一代歷歷在目.

$^{721}$所以你要看著昔日的情景.
也要將這段民族的記憶.
印在他們整個這一代.
也要世世代代地提醒他們.
不要忘記上帝.
所以他提醒他們.
謹慎因勤保守你們的心靈.
這種謹慎的態度.
為了什麼呢?.
就是免得忘記你親眼所見的事.
免得你們一生.
這時離開你們的心.
要令到學習敬畏.
第一件事就是不要忘記舊時的功課.
舊時走錯的路.
不要再走下去了.
那些曾經令到無法承受應許.
或者惹怒上帝的一些經歷.
不要再走那條錯路了.
而且還要不單止自己知道.
還要世世代代提醒你們的子女子子孫孫.
讓他們走在正確的路上.
錯誤的走過去就算了.
不過那個算的意思就是.
不要再耿耿於懷.
要成為一個真正的功課.
以致他們謹守遵行.
所以在這裡.
他叫他們第一.
不要令到整個民族失憶.
第二個重要的字眼.
就是後來他所說的.
他們去到河獵山的時候.
去做上帝的約.
約上帝頒佈的方式其實是挺可怕的.
他說到約的頒佈是直著摩西一個人說的.
在山上拿著兩個法板.
然後就將十戒頒佈給以色列民.
但是頒佈十戒的時候.
上帝在山裡面是隱密之處.

$^{761}$但是上帝就用一個火焰沖天.
昏黑,默暗,幽暗.
上帝在火焰中對以色列民說話.
你想起這個場面.
是很令人可怕的.
是一個大自然的力量.
令到所有人都變得驚懼.
這種驚懼.
在以色列民覺得恐懼的狀態下.
上帝就頒佈了這個律法.
這個驚懼的用字.
其實用在敬畏的字是同一個字.
在聖經裡面說.
不過那個態度是不同的.
如果那些子民永遠都是想著.
上帝是一個很恐怖的上帝.
很驚懼的上帝.
人就會敬而原之.
就會與上帝拉遠了關係.
因為人的本性一見到驚懼的東西.
就會害怕就會逃跑.
那個驚懼是保護自己的.
不過上帝說的敬畏不是叫你.
因為怕他所以就遠離他.
反而是在一種驚懼.
在上帝的威榮之下.
學習怎樣去親近上帝.
上帝有大而可畏的一面.
但上帝也有自福恩慈慈愛的一面.
我們真正能夠和上帝親近.
就是藉著我們明白上帝的旨意.
明白上帝與子民立約.
所以在這裡.
摩西用的字就是當首的約指示你們.
他在說上帝有這麼恐怖的一面.
不過到了最後.
真正的關係建立是說.
興邁不得的上帝.
都是很樂於和人建立約的關係.
人就帶著一份親近上帝.

$^{801}$或者上帝其實是親近人.
主動和人建立關係的上帝.
我們就帶著這一份敬愛.
和上帝好好相處.
這一份敬畏是一個很大的藝術.
今天很多弟兄姊妹怕神.
因為怕自己做得行差踏錯.
所以經常會步步為營.
本來這種謹慎是好的.
但如果基於那個觀念.
只是覺得上帝很恐怖.
怕自己隨時落地獄.
或者受到刑罰.
這種的靈性是不健康的.
我們知道上帝興邁不得.
但同時也知道上帝樂於賜福.
將祂的恩典賜給我們的神.
我們只管坦然無懼.
進到上帝的面前.
領受祂的恩惠.
也在生命裡繼續堅持.
我們不想上帝不開心.
而不是怕上帝發怒來懲戒我們.
這種敬畏的態度.
不單是新一代的以色列民要學習.
經文摩西更加說.
你要一直傳給你的子子孫孫.
學習敬畏.
學習有一種對神有敬畏.
對人有尊重.
你就能夠在地上好好地過活.
經文大致說到這裡.
不是很複雜.
說的是一個新典神學.
用這個框架就很容易理解.
我們今天如何應用這段經文.
對我們有什麼重要幫助呢.
我想起有兩點.
第一點.
今天作為基督徒.

$^{841}$其實我們可以停棄空堂.
或者我們所信的信仰.
你要有百分之百的信心.
有時候實踐信仰並不容易.
但你知道日子有功.
你就會在不容易的日子裡.
你的生命會變得很不同.
甚至因為你專心依靠耶和華上帝.
你明白聖經的教導.
你的改變是會由內而外改變.
日子久了的時候.
你就會看到這種生命的功力.
而這種呈現.
跟這個世界截然不同的表現.
當你信主信得久的時候.
你再比對一下.
未信主的跟你一起長大的朋友.
你不是去批判他.
不過你會看到.
如果沒有上帝.
可能你的生命跟他們.
那些價值觀念.
或者那些處事不尊重人.
或者對神沒有敬畏的心.
你就會呈現出來.
但當你生命裡.
認真地實踐這個信仰.
日子久遠.
你就會看到你生命截然不同.
甚至你的生命的本質.
是會令人很羨慕的.
今天很多人說.
傳福音最有效的方法是什麼?.
除了能夠延傳之外.
其實就是在關係裡傳福音.
關係除了大家很親密之外.
大家之間.
弟兄姊妹很有真誠相處.
弟兄姊妹之間.
有一份很濃烈的熱情.

$^{881}$對上帝有一種敬畏的心.
一起唱詩.
一起渡過生命的高高低低.
這種生命本身.
都令這個世界非常羨慕.
今天如果我們真的要好好實踐這個信仰.
其實聖經裡有教導我們.
我們每個星期都在做這些事.
不過即使變成果效.
好像吃了止痛藥.
一吃就馬上改變.
未必是.
不過是一個長年累月.
日子有功.
當你隔了一段時間再回望的時候.
你就會發現自己不同了.
由主幫助我們.
我們努力實踐這些基本功.
日子久遠的時候.
你會看到有很多的改變.
第二件事.
這個改變.
你不要只放在你自己這一代.
或者自己在我們今天處境當中.
我們的信仰需要傳承.
我們的信仰需要世世代代傳遞下去.
不單滿足於我們當下這一代.
我們這一班人生命的滿足.
我們更加寄望的是我們的子子孫孫.
他們的生命同樣都得到上帝的應許.
盼望見到你的兒女.
在他們的同輩當中.
他們成為一個decent的人.
成為一個真是.
上帝喜悅的一個年輕人.
一個上帝喜悅的孩子.
而不是成為一個生骨大頭菜.
好好地管教他們.
世世代代地.
從小到大已經將上帝的價值觀.

$^{921}$將上帝的美善教導他們.
而不要將外面社會的職責放進去.
如果你將這個世界的觀念放進去的時候.
到最後其實孩子還是沒有辦法經歷上帝的豐盛.
求主幫助我們.
當我們今天在新春的日子.
我寄望我們旺角浸信會.
我們的世世代代.
我們下一代又興起.
全部都是上帝所喜悅的子民.
更加願我們這一代.
我們的生命.
活出良藝的生命.
活出一個令人羨慕的生命.
求主繼續來帶領我們.
祝福我們.
(今天視頻就拍到這裡啦 拜拜~).
\newpage



\section{約翰福音 }
\label{sec:uDlv_6lnHKk}
\textbf{2025年2月9日崇拜直播｜梁家麟牧師｜我也不定你的罪｜約翰福音 八 1-11}
\newline
\newline
連結: \href{https://youtube.com/watch?v=uDlv-6lnHKk}{\texttt{https://youtube.com/watch?v=uDlv-6lnHKk}} ~~~~ 語音日期: 2025-02-09
\newline
\newline
\hyperref[sec:cuIe7N8imyc]{\small{< < < PREV SERMON < < <}}
~
\hyperref[sec:index_chronic]{\small{[返順時目]}}
~
\hyperref[sec:oex63GVBJWc]{\small{> > > NEXT SERMON > > >}}
\newline
\newline
$^{1}$今日經訓是在約翰福音8章1至11節.
新約聖經137頁至138頁.
「於是各人都回家去了.
耶穌卻往橄欖山上去.
清早又回到殿裡,眾百姓都到他那裡去.
他就坐下教訓他們.
文士和法利賽人帶著一個行淫時被拿的婦人來.
叫她站在當中,就對耶穌說.
「夫子,這婦人是正行淫之時被拿的.
摩西在律法上吩咐我們.
把這樣的婦人用石頭打死.
你說該把她怎麼樣辦呢?」.
他們說者說:乃試探耶穌.
要得著告她的把柄.
耶穌卻彎著腰,用指頭在地上畫字.
他們還是不住的問她.
耶穌就直起腰來對他們說.
「你們中間誰是沒有罪的.
誰就可以先拿石頭打她?」.
於是又彎著腰,用指頭在地上畫字.
他們聽見者說.
就從老到少一個一個的都出去了.
只剩下耶穌一人.
還有那婦人仍然站在當中.
耶穌就直起腰來對她說.
「婦人,那些人在哪裡呢?.
沒有人定你的罪嗎?」.
她說:主啊,沒有.
耶穌說:我也不定你的罪.
去吧,從此不要再犯罪了.
《大紀元》字幕組.
弟兄姊妹早安.
早安.
今天看的故事是我們很熟悉的故事.
我自己也說過很多次.
還有一篇講章更加出版了.
就是《彎腰的耶穌》.
有興趣的弟兄姊妹自己拿來看看.
今天我嘗試用一個新的角度.
來重讀這一段經文.

$^{41}$首先要指出.
主人崇拜很少講這件事.
不過這是真事,重要的事.
《約翰福音》八章一到十一節.
這一段經文.
在最早期的古經版本.
通通都沒有出現過.
卻是在二世紀之後.
才被收入《約翰福音》裡面.
所以學者幾乎一致的意見都是.
這個故事應該不是出自《約翰福音》的手筆.
而是來自其他傳統.
什麼叫其他傳統呢?.
例如一些非正典的福音書.
Non-canonical gospel.
你知道早期教會.
流傳的不僅僅是我們今天那四卷福音書.
有很多其他的福音書.
然後後來才被拼入去.
《約翰福音》七章五十二節後面.
變成第八章的第一段.
絕大多數的學者.
都不會認為這段故事.
是《約翰福音》的原來部分.
所以過半的專業學者.
乾脆在解經的時候.
就跳了這一段,根本不解它.
因為根本不當這一段是《約翰福音》的一部分.
你看波曼.
甚至是福音派的Leon Morris.
基本上都完全skip了這一段.
就不處理它.
不過學術研究歸學術研究.
我們在教會裡面.
就有完全不同的態度.
在教會裡面我們尊奉聖經.
正典的權威.
我們今天擁有的這一本聖經.
都是我們要依從的真理.
我們不會拿著今天的聖經.

$^{81}$然後就說這一段要守,那一段不守.
如果這樣就很大亂了.
太危險了.
所以我們還是認定.
不管這一段聖經是不是出自.
《約翰福音》的手筆.
這個故事還是《約翰福音》.
不過分割的一部分.
必須要申明的是.
這個故事即使不是.
約翰寄述.
不代表它不是真實發生的.
因為所有正典的福音書.
都不可能將耶穌.
33年所發生的.
所有事跡都收錄下來.
漏記,缺記的部分.
是多到不得了的.
約翰自己都承認.
耶穌在門徒面前.
另外行了許多神蹟.
會有記載在這書上.
但記者之時.
要叫你們信耶穌是基督.
是上帝的兒子.
並且叫你們信了他.
就可以因他的名得生命.
又說耶穌所行的神蹟.
還有許多.
若是一一都寫出來.
我想所寫的書.
就是世界也容不下了.
世界也容不下了.
就是有點誇張的.
不過你想一想都覺得很恐怖.
耶穌在公眾事奉生涯中.
平均假如.
假如每天跟門徒只說兩小時.
24小時說兩小時.
不算多吧.

$^{121}$那三年半就1200天.
1200乘以2就2400小時.
我不知道這2400小時.
有多少是重複的內容.
不過如果真的要將每句說話都記載下來.
我們今天拿著這本聖經都很可怕.
不知道該怎麼辦.
所以.
我相信這個故事.
應該不是約翰自己寫的.
不屬於最早期版本的約翰福音的一部分.
但他收錄在我們今天的正典中.
是有聖靈的心意的.
我認定這個故事是真實發生的.
是耶穌事奉生涯中其中一個的片段.
故事.
如果我們黏著約翰福音七章解決.
是發生在聖殿裡.
耶穌教訓人的時候.
如果我們相信這個故事跟之前的約翰技術是連在一起的.
那就要假設這個故事同樣是發生在主彭節之後.
我們已經說過在主彭節期間.
耶穌跟猶太教當局有幾次的衝突.
後兩次的衝突是在主彭節的最後一天發生的.
這樣.
耶穌在主彭節結束之後沒有立刻離開.
祂是在橄欖山多留一晚.
第二天就是節後的一天.
祂再去到聖殿教訓人.
這些聽眾或者是平日.
慣常來聖殿的人.
他們多數是住在耶路撒冷附近的人.
來聖殿是很方便的.
所以不只是大年大節才來.
每天都來.
又或者像耶穌那樣.
在節期結束之後不是立刻離開.
多留兩天才離開的那些人.
就在耶穌教訓人的當下.
有一群文士和法利賽人.

$^{161}$帶著一個在行淫時當場被捕獲的婦人.
抓了他們來.
然後質問耶穌要怎樣處置這個婦人.
講一點題外話.
這段聖經不是出自約翰的手筆.
有很多證據.
包括了這裡有記載文士和法利賽人.
經文你會覺得很熟悉.
《寬恩書》這個說法在馬太盧家也很常見.
不過約翰只有一次.
就是這裡.
之前之後從來沒有文士和法利賽人的技術.
然後叫耶穌做夫子.
在約翰福音十萬零一次只有這裡.
前後都不是這樣叫的.
但是其他的《寬恩書》是這樣叫的.
所以我們可以說明.
這段經文不應該是約翰的手筆.
因為用字都不一樣.
無論如何.
他們把這個婦人抓到耶穌面前.
問耶穌要怎樣處置.
我們知道聖殿是一個神聖的地方.
本來容不下罪人眾.
這裡也不是一個審判場所.
聖殿不是高法院.
不過聖殿卻是討論律法.
要怎樣理解和執行的地方.
用今天的說法.
聖殿是一個神學論壇.
很多律法師會帶著他們的弟子來到聖殿.
公開討論律法要怎樣理解和怎樣應用.
在這個時候那些法利賽人.
那些文士把婦人抓來.
都是本著同樣的心態.
他們把婦人帶來這裡.
不是問耶穌要怎樣審她.
耶穌又不是法官.
耶穌沒有權審.
他們把她變成一個案例.

$^{201}$問耶穌我們可以怎樣理解律法.
怎樣應用在這個人身上.
要耶穌參與這個律法應用的討論.
不是公審.
這裡不是嚴格意義的犯人或被告.
這是一個辯論擂台.
那個婦人是一個被觀察的標本.
一個被公開議論的案例的主角.
4到5節.
夫子,這婦人是正行淫之時被拿的.
摩西在律法上吩咐我們.
把這婦人用石頭打死.
你說該把她怎樣呢?.
問事他們先表達兩件事.
第一是事實.
這個婦人是行淫當下被現場捕獲的.
所以罪狀是坐實了.
就沒有誤會,以為.
現場沒有辯解的餘地.
第二是法律的規定.
他們把律法的明文規定說出來.
並且搬出摩西的權威.
我們參考利未記十章十節.
新聞紀二十二章二十二節.
律法是這樣規定的.
這個婦人要被治死.
在這兩個基礎上.
一個是事實,一個是法律觀點.
然後他們問耶穌的意見.
又要說一點題外話.
根據猶太的律法.
犯奸淫行淫是要處死的.
但根據羅馬的法律.
行淫是小事.
羅馬人對道德倫理的觀點.
跟今天歐美差不多.
有什麼所謂.
一夜情,你好我好,大家好.
所以羅馬的法律.
行淫是不處死的.

$^{241}$是沒有死刑的.
而我們又知道.
當時猶太人亡國了.
工會有權審案.
但工會是沒有權判死刑的.
所以耶穌也不可以.
工會直接判他死刑.
要捉他去彼拉多.
所以那班文士和法理人.
其實他們都沒有能力處死這個婦人.
他們可以起哄群眾用石頭扔死她.
這是可以的.
因為群眾是無人性.
捉不到的.
只是暴民.
但這些文士和法理人是有頭有面的人.
如果他們說處死.
他們直接跟羅馬帝國對著幹.
他們都要背大鑊.
所以他們都沒有能力處死這個婦人.
只是用石頭扔死她.
不是他們主張要整死她.
而耶穌說不主張整死她.
耶穌不跟他們意見.
他們不敢這樣說.
你可以說是設定了耶穌.
整蠱耶穌.
現在是試探耶穌.
看他們怎樣看這件事.
不是他們真的覺得可以處死這個婦人.
這些題外話.
不過希望你明白.
這些事情.
一個證據確鑿的犯罪者.
一條清清楚楚規定犯奸淫者當至死的律法.
你說耶穌還有什麼遊刃的空間呢?.
很有趣的是.
文氏和法利賽人.
專門將這個案例.
拿來詢問耶穌.

$^{281}$明顯的是.
他們知道耶穌.
會有跟摩西的律法不一樣的觀點.
如果耶穌的回應是.
這個婦人犯了奸淫.
那就整死她吧.
那就沒戲了.
因為文氏和法利賽人.
不是要思考怎樣處置那個婦人.
那個婦人根本不是他們關注的重點.
他們想藉著這個婦人的案例.
去找耶穌犯錯的把柄.
聖經說的是他們要抓耶穌的把柄.
處置耶穌才是他們要處理的事.
是不是?.
聖經說.
他們說者患乃試探耶穌.
要得就告他的把柄.
所以如果耶穌的答案跟他們一樣.
那就是沒有把柄可以抓.
整個行動就是白費的.
是不是?.
文氏和法利賽人.
在未詢問耶穌之前.
已經猜到耶穌應該會不同意處死這個婦人.
為什麼他們會有這樣的猜想?.
第一.
耶穌對律法的解釋常常跟傳統不一樣.
有很多新觀點.
第二.
更加重要的是.
他們知道耶穌慈悲為懷.
對人有強烈的憐憫的心.
不忍將犯罪者處以極刑.
他們利用耶穌善心的弱點.
來對他進行攻擊.
這個我以前講道也講過.
敵人通常是最了解自己的人.
反倒朋友和學生不一定明白.
與耶穌為敵的法利賽人文氏.

$^{321}$都知道耶穌對人慈悲為懷.
不輕易傷害生命.
但是耶穌的門徒.
就不一定知道.
他們竟然因為不被接待的事.
就求耶穌降天火.
燒死全村的撒瓦里亞人.
你記得這個故事吧.
路加福音9:54.
雅各和約翰這兩個主角.
是耶穌最貼身的弟子.
竟然說出這些喪心病狂的話.
無怪耶穌憤怒地罵雅各和約翰.
你們的心如何?你們並不知道.
主意不是他們不明白耶穌的心.
不明白耶穌的心已經夠死了.
而是他們不知道自己的心.
他們不知道自己在滿口義正詞義.
替天行道的說話背後.
他們的心是多麼惡毒污穢.
我們經常說聖戰.
說要替天行道.
要斬草除魔.
但其實在正義的話背後.
我們的心其實是最惡劣最污穢.
你們的心如何?你們自己不知道.
文士和法利賽人很了解耶穌.
用他的善良作為攻擊他的軟肋.
耶穌的門徒卻不一定了解耶穌.
你們是不是很有趣?.
面對文士和法利賽人的質詢.
耶穌沒有馬上回應他們.
他做了一個很奇特的動作.
彎腰在地上寫字.
耶穌在地上寫了什麼字我們不知道.
有人說耶穌這個寫字的舉動.
純粹是想退後.
退後.
耶穌不想參與整件事.
要置身事外.

$^{361}$耶穌不喜歡這些討論.
所以做一些你覺得奇奇特特.
很無理的事.
有人估計耶穌是在想如何回應文士的提問.
如何為那個婦人的困境解困.
有人估計耶穌可能是藉著寫字的行動.
來平息心中的怒火.
這些人竟然以他的善良作為攻擊點.
就像我們經常在國內看到的碰瓷.
一個老人家假裝跌在地上.
有人去扶他.
馬上有人圍過來.
你推他跌,你弄傷他.
於是要賠一大筆錢.
這是一個非常卑鄙的行徑.
做這些事的後果是.
日後就沒有人敢再行善.
見到真的有需要跌在地上的人都不敢扶.
我怎知道不是有一個陷阱.
這是很卑鄙的.
踐踏利用別人的善良.
這是最卑鄙無恥的做法.
為了對付耶穌.
他們就踐踏他的善良.
踐踏他對人的尊重.
還有他們把婦人帶到聖殿來示眾.
完全不理會婦人當下的尊嚴感受.
也是一個極其殘忍的行為.
所以耶穌如果真的很憤怒.
去平息自己心中的怒火也是可以想像的.
文士再三追問.
一定要耶穌表態.
最後耶穌就發話.
但不是直接回應文士的提問.
這個婦人究竟有沒有犯罪.
應該按照律法怎樣處分.
耶穌是轉向在場的群眾說.
你們中間誰是沒有罪的.
誰就可以先拿石頭打他.
第七節.

$^{401}$耶穌不是在說應不應該用石頭扔死這個婦人.
耶穌只是建議一個扔死她的程序.
讓沒有犯罪的人先扔.
有犯罪的人等一等.
遲些再扔.
耶穌也不是說.
如果一個人是自覺有罪的就沒有權懲罰人.
不是.
耶穌不是這樣說.
第一, 律法沒有這樣的規定.
耶穌不可以擅自增加篡改律法的條文.
第二, 這個主張在實踐上是沒法執行的.
如果人人都有罪.
人人不能執法,司法.
這個世界就大亂了.
警察有罪, 所以警察不可以抓人.
法官有罪, 法官不可以判人罪.
這個世界還有天理嗎?.
這個世界就大亂了.
有罪的人不可以去懲罰人.
這是荒謬的.
第三, 其實這也是明目張膽.
挑戰整個舊約精神.
舊約是怎麼規定的?.
律法規定.
一個祭司在為百姓獻上挽回祭之前.
要先為自己獻祭.
贖自己的罪.
這個律法規定是意味著什麼?.
祭司有沒有罪?.
有罪的.
祭司一定要知道自己有罪.
沒有罪的祭司.
所以先結證自己.
然後才為百姓行結證儀式.
所以在希伯來書就說.
只有耶穌一個, 不用先結證自己.
就可以為百姓獻罪.
所以, 祭司要先赦免自己的罪.
然後才可以赦免別人的罪.

$^{441}$這是整個舊約精神.
是講明了.
人人有罪, 無人例外.
所以, 祭司不是高人一帝的無罪者.
是先除了自己的罪.
然後協助除別人的罪.
這是舊約的精神.
人人都有罪, 人人都要除罪.
所以, 如果耶穌說.
無罪的才可以除別人的罪.
這是直接與整個舊約相衝突.
耶穌只是說.
讓自覺無罪的先扔石頭.
問題在於在場.
沒有人自覺無罪.
所以沒有人撿起第一塊石頭.
沒有第一塊, 就沒有第二塊.
對不對?.
僵持了一段時間.
不知多久.
大家覺得很尷尬.
你眼望我.
大家都覺得氣氛很有壓迫性.
於是大家覺得沒有人, 沒有氣體.
於是一個一個閃了.
一個一個離場.
當所有人離開後.
耶穌轉向那個婦人說.
「我也不定你的罪, 去吧!.
從此不要再犯罪了!」.
「我也不定你的罪」.
注意, 耶穌不是先定罪.
然後赦罪.
他是先赦罪, 然後定罪.
在宣告赦罪之前.
耶穌沒有開口說過.
任何責備那個婦人的話.
祂是首先宣告赦罪.
不過, 這個赦罪的宣告.
已經包含了對罪的判定.

$^{481}$沒有罪, 何來需要寬赦?.
宣告寬赦, 即是假設你有罪.
我不用你還錢.
即是意味著你欠了我的錢, 對不對?.
否則我無端端跟你說.
我不用你還錢, 是嗎?.
所以是先赦罪, 然後定罪.
如果耶穌說.
沒有罪者, 先扔石頭.
是為那個婦人解困.
這是表現出祂的機智.
這裡「我也不定你的罪」.
就流露出祂對犯罪者憐憫的愛心.
耶穌不是否認他犯罪.
不是無視他犯罪.
而是給他悔改重生的機會.
你是有罪的, 但我現在不定你的罪.
不過, 悔改的機會不是永遠有的.
改變是必須的.
好好做人, 以後不要再犯罪.
耶穌是對人的在乎著緊.
這是憐憫.
憐憫是憐憫軟弱的人.
不是人的軟弱.
是軟弱的人, 人是重點.
憐憫是認定他所犯的罪.
抹殺犯罪者的人格和個性.
婦人雖然是罪犯.
但婦人仍然是真實的人.
有尊嚴, 有價值.
憐憫是看出人的可能性.
所犯的罪是後天的, 是偶然的.
不是本質性的.
不是他本質邪惡, 本性難移.
所以他有改變的可能性.
給第二次機會是值得的.
有用的.
我們經解讀到這裡.
和大家做一些思考.
拉得比較闊的思考.

$^{521}$一開始我們說.
這個故事不是最早版本的.
約翰福音的一部分.
而是在二世紀之後.
才被拼入約翰福音.
我們於是要問.
初期教會為什麼要將這個故事.
放入正典裡面呢?.
有學者做了一個很好的建議.
我相信這是一個很合理的解釋.
二世紀的教會.
面臨著嚴酷的政治迫害.
有不少信徒殉道.
不過更多信徒叛教.
不用說的.
在生死關頭.
叛教的人一定壓倒性多於殉道的人.
物以罕為貴.
如果人人都殉道.
殉道就像吃生菜一樣容易.
那就不用歌頌殉道.
有沒有歌頌呢?.
就是因為殉道少,殉道難.
殉道不是很多人.
所以才要歌頌.
甚至將殉道士放為聖人.
正正是因為迫害很嚴峻.
叛教的人很多.
早期教會很自然的.
就會大力推崇信仰中正的人.
鄙視軟弱跌倒的人.
確實我們知道.
叛教和不叛教的人.
人間後果是天與地的懸殊差異.
你說吧.
說一句「我就在耶穌」.
我就放你了.
說一句「我就在耶穌」.
我就可以保住身家性命財產.
可以維持我家人的安全.

$^{561}$保住家庭的完整.
保住我的社會地位.
拒絕說「我就在耶穌」.
拒絕叛教的人.
就要家破人亡.
人間的名利富貴盡失.
這個很悲慘.
對不對?.
好了,迫害不是永遠綿延的.
是定期發生的.
幾年之後又停一停.
迫害暫時終止了.
教會還要面對.
是否讓那些叛教的人回來教會.
他們說「我道歉」.
說一句「對不起」.
就等於沒事發生.
旁邊那些堅持不叛教.
不過還沒死去.
還沒死去.
那些家破人亡的人.
你跟我排排坐在旺站嗎?.
有沒有搞錯?.
我輸光光,你什麼都沒輸.
你說一句「對不起」就跟我完全一樣嗎?.
我怎麼能忍受一句「對不起」.
就可以解決你叛教的問題呢?.
對不對?.
日後教會還怎麼鼓勵信徒堅守信仰.
為主背十字架嗎?.
耶穌是否很清楚說過.
「凡在人前認我的,我在天上的父面前也必認他」.
「凡在人面前不認我的,我在天上的父面前也必認他」.
我們知道在三世紀羅馬帝國兩次最嚴重的宗教迫害之後.
教會都出現大分裂.
一次是所謂Norvigianism.
洛伐天主義.
另一次是多立陶主義.
分裂的原因都是因為.
那些叛教的人要回教會.

$^{601}$教會要怎麼看待那些曾經叛教的人.
這是一個大課題.
正正是因為政治形勢很嚴峻.
教會普遍的氣氛.
傾向對信徒提出高端的要求.
藐視那些軟弱跌倒的人.
這樣很容易在一個群體裡面產生屬靈精英主義.
Spiritual elitism的氛圍思想.
推崇少數能夠攀登最高境界的人.
輕視那些無法達標的多數.
高舉意志.
Real power.
鄙視感情,情緒.
在軍隊裡面情緒是最沒用的.
撐著,很劣就去死吧.
你怕是沒有地位的.
尊崇強者.
排斥弱者.
成王敗寇.
歌頌成功者.
厭棄失敗者.
對失敗者缺乏憐憫關懷的心.
視他們的失敗為純粹至今墮落.
他們要背負所有責任.
是罪有應得,不值得同情.
鑒於二世紀教會普遍存在著.
屬靈精英主義的思維.
有心的人,通常是有權力的.
可能某一間教會的主教.
刻意將耶穌的故事收入約翰福音.
以對信徒的想法作出調整和修正.
如果主教正要處理犯罪的人.
如何執行紀律處分.
他想對犯罪者做寬大處理.
就援引這個流傳的耶穌故事.
將它收入約翰福音.
我們知道.
福音書記載耶穌的故事.
通常是兩件事.
第一是耶穌行過的神蹟.

$^{641}$治病趕鬼.
第二是耶穌的教導.
登山寶訓.
但我要說.
在這個故事裡.
約翰福音8章兩件事都不是.
又不是教導,又不是神蹟.
首先,耶穌沒有行過甚麼神蹟.
他沒有用大能拯救行人被捉的婦女.
一個神蹟都沒有行過.
第二,耶穌也沒有講過甚麼教訓.
沒有,沒有講過教訓.
他之前就在聖殿教訓人.
直到耶穌將婦人拉到來為止.
然後耶穌就一直沒有說話.
一直彎腰寫字.
不肯說一句說話.
對不對?.
他後來被迫要說話.
他只是做了一個提議.
他不是教導.
誰沒有犯罪才扔.
他只是做了一個這樣的提議.
他不是說可以扔,不可以扔.
他沒有這樣說過.
只是做了一個提議.
所以不是原則性的教導.
是扔石頭的次序的建議.
這個故事又沒有神蹟,又沒有教導.
要傳遞給我們甚麼訊息?.
唯一的訊息.
只是流露了耶穌對罪人的關愛憐憫.
唯一的訊息.
就是耶穌憐憫這個犯罪的人.
這是耶穌的性格.
是耶穌的心.
他是這樣看待罪人的.
如果我們要在這個故事裡有所學習.
就是這一點.
這是唯一要學的功課.

$^{681}$他是至關要緊的功課.
好像前面說.
耶穌不是包容罪惡.
是憐憫罪人.
不是視有罪為無罪.
而是容讓罪人悔改.
給人重新做人的機會.
要跟今天的自由主義者.
把奸淫搶劫吸毒.
通通非刑事化.
打劫950美元以下的財物就不算犯罪.
吸食毒品不算犯罪.
完全不同.
罪是罪.
罪是上帝定的標準.
這是不可以改變的.
聖經說上帝萬不以有罪為無罪.
這是原則.
但上帝卻憐憫犯罪的人.
給他們迴轉的機會.
耶穌確定人犯罪.
耶穌也有權定人的罪.
不過如果定了罪.
人就死了.
再沒有翻身的可能性.
所以耶穌決定放人一馬.
給人悔改.
重新開始的機會.
耶穌來到人間施行拯救.
目的是為了赦罪.
不是為了定罪.
這是約翰早已經很清楚說的.
約翰福音三章十六十七節.
前面那節.
你閉上眼睛也會背.
上帝愛世人.
甚至將他獨生子賜給他們.
以一切信仰大帝.
不致滅亡.
反得永生.

$^{721}$因為上帝差他的兒子降世.
不是要定世人的罪.
而是要叫世人因他得救.
這是我們向這個世界所宣告的訊息.
定罪不是福音.
赦罪才是福音.
我要用最簡單的說話.
去說什麼叫基督教的福音.
就是耶穌這句話.
我也不定你的罪.
這就是福音.
我們有罪.
不過耶穌說.
我們的主說.
他不定我的罪.
他赦免我的罪.
這就是福音.
福音就是這樣的訊息.
犯罪者受懲罰.
這是理所當然的事.
我們說罪有應得.
反倒罪得赦免是不合理的.
不是人應得的.
不應得的就叫恩典.
福音是恩典.
十字架是恩典.
耶穌的救贖百分百是恩典.
基督是上帝恩典的見證者.
教會是恩典的見證者.
正義凜然譴責罪惡.
不是教會的本業.
會出恩典的見證.
處處流露出對人的包容節合.
對罪人的軟弱跌倒.
施以憐憫關愛.
這才是教會在人間的使命.
弟兄姊妹.
我們怎樣看同性戀者.
怎樣看吸毒者.
怎樣看自暴自棄的邊緣青少年.

$^{761}$我們怎樣去看殘暴的哈馬斯.
殘暴的以色列鷹派政府.
是切齒痛恨.
呼求上帝盡早降下天火.
消滅他們.
我不知道你們看這些新聞的時候.
會不會心中有這樣的想法.
燒死這些人.
因為我們對世人.
陷在罪中不能自拔.
感覺到傷痛.
對人的愚昧無知.
有更大的同情關愛.
我們對身邊罪人的態度.
像耶穌多些.
像那些自覺正義.
自覺是真理化身的.
法利賽人多些.
這裡再總結一次.
整個故事.
沒有提及耶穌任何重要的教訓.
做過什麼神蹟.
只是流露出耶穌對罪人的關愛憐憫.
豁現了耶穌的肺腑心腸.
我們在這個故事要學習的.
就是體會耶穌的心.
體貼耶穌的心.
我們要以耶穌的心為心.
黃國贈信會弟兄姊妹字幕.
我們要努力追求成長.
人人都要竭力與所蒙的恩相稱.
但我們教會卻拒絕.
屬靈精英主義的思想.
我們沒有任何互相競賽.
大家鬥屬靈.
而是像保羅所說的.
向下看齊我們教會.
有智慧的照顧那些沒有智慧的.
有能力的照顧那些沒有能力的.
我們律己以嚴.

$^{801}$待人以寬.
永遠對人包容接納.
營造一個讓人可以在這裡自由成長.
可以跌倒.
可以軟弱.
但是又重新爬起來.
再嘗試的空間.
我們的信念是.
在恩典裡面永遠有可能.
在恩典裡面我們拒絕對人絕望.
畢竟我們每個人都經歷過.
被耶穌放一馬.
兩馬三馬.
很多次放一馬的經歷.
我們知道.
這是我們今天能夠坐在這裡的原因.
蒙恩者學習成為私恩的人.
\newpage



\section{申命記 4:15-31}
\label{sec:oex63GVBJWc}
\textbf{2025年2月16日崇拜直播｜陳冠成牧師｜朝向應許之地－禁戒雕刻偶像｜申命記4\_15-31}
\newline
\newline
連結: \href{https://youtube.com/watch?v=oex63GVBJWc}{\texttt{https://youtube.com/watch?v=oex63GVBJWc}} ~~~~ 語音日期: 2025-02-16
\newline
\newline
\hyperref[sec:uDlv_6lnHKk]{\small{< < < PREV SERMON < < <}}
~
\hyperref[sec:index_chronic]{\small{[返順時目]}}
~
\hyperref[sec:bBtvGdEATA0]{\small{> > > NEXT SERMON > > >}}
\newline
\newline
申命記 4:15-31
\newline
\begin{longtable}{cl}
\hline
\hline
章節 & 經文 (和合本修訂版)\\
\hline
4:15 & \begin{tabularx}{0.7\textwidth}{X} 「所以，你們為自己的緣故要分外謹慎；因為耶和華在何烈山，從火中對你們說話的那日，你們沒有看見任何形像。 \end{tabularx} \\ \\ \relax
4:16 & \begin{tabularx}{0.7\textwidth}{X} 惟恐你們的行為敗壞，為自己雕刻任何形狀的偶像，無論是男像或女像， \end{tabularx} \\ \\ \relax
4:17 & \begin{tabularx}{0.7\textwidth}{X} 或地上任何走獸的像，或任何飛在空中有翅膀的鳥的像， \end{tabularx} \\ \\ \relax
4:18 & \begin{tabularx}{0.7\textwidth}{X} 或地上任何爬行動物的像，或地底下任何水中魚的像。 \end{tabularx} \\ \\ \relax
4:19 & \begin{tabularx}{0.7\textwidth}{X} 又恐怕你向天舉目，看見耶和華－你的神為天下萬民所擺列的日月星辰，就是天上的萬象，就被誘惑去敬拜它們，事奉它們。 \end{tabularx} \\ \\ \relax
4:20 & \begin{tabularx}{0.7\textwidth}{X} 耶和華將你們從埃及帶領出來，脫離鐵爐，是要你們成為他產業的子民，像今日一樣。 \end{tabularx} \\ \\ \relax
4:21 & \begin{tabularx}{0.7\textwidth}{X} 「耶和華又因你們的緣故向我發怒，起誓不容我過約旦河，不讓我進入耶和華－你神所賜你為業的那美地。 \end{tabularx} \\ \\ \relax
4:22 & \begin{tabularx}{0.7\textwidth}{X} 我只好死在這地，不能過約旦河；但你們必過去得那美地。 \end{tabularx} \\ \\ \relax
4:23 & \begin{tabularx}{0.7\textwidth}{X} 你們要謹慎，免得忘記耶和華－你們的神與你們所立的約，為自己雕刻任何形狀的偶像，就是耶和華－你神所禁止的， \end{tabularx} \\ \\ \relax
4:24 & \begin{tabularx}{0.7\textwidth}{X} 因為耶和華－你的神是吞滅的火，是忌邪的神。 \end{tabularx} \\ \\ \relax
4:25 & \begin{tabularx}{0.7\textwidth}{X} 「你們在那地住久了，生子生孫，若行為敗壞，為自己雕刻任何形狀的偶像，行耶和華－你神眼中看為惡的事，惹他發怒， \end{tabularx} \\ \\ \relax
4:26 & \begin{tabularx}{0.7\textwidth}{X} 我今日呼天喚地向你們作見證，你們在過約旦河得為業的地上必迅速滅亡！你們在那地的日子必不長久，必全然滅絕。 \end{tabularx} \\ \\ \relax
4:27 & \begin{tabularx}{0.7\textwidth}{X} 耶和華必將你們分散在萬民中；在耶和華領你們到的列國中，你們剩下的人丁稀少。 \end{tabularx} \\ \\ \relax
4:28 & \begin{tabularx}{0.7\textwidth}{X} 在那裡，你們必事奉人手所造的神明，它們是木頭，是石頭，不能看，不能聽，不能吃，不能聞。 \end{tabularx} \\ \\ \relax
4:29 & \begin{tabularx}{0.7\textwidth}{X} 你們在那裡必尋求耶和華－你的神。你若盡心盡性尋求他，就必尋見。 \end{tabularx} \\ \\ \relax
4:30 & \begin{tabularx}{0.7\textwidth}{X} 日後你在患難中，當這一切的事臨到你，你必歸回耶和華－你的神，聽從他的話。 \end{tabularx} \\ \\ \relax
4:31 & \begin{tabularx}{0.7\textwidth}{X} 耶和華－你的神是有憐憫的神，他不撇下你，不滅絕你，也不忘記他起誓與你列祖所立的約。 \end{tabularx} \\ \\
[1ex]
\hline
\hline
\end{longtable}
$^{1}$《新命記》第四章 第15至20節.
第23至31節.
在我們手上的聖經219頁.
也印在了情在彼的裡面.
《新命記》第四章15節.
所以你們要憤愛謹慎.
因為耶和華在何列山.
從火中對你們說話的那日.
你們沒有看見什麼形象.
唯恐你們敗壞自己.
挑剋偶像.
彷彿什麼男像女像.
或地上走秀的像.
或空中飛鳥的像.
或地上爬蜜的像.
或地底下水中魚的像.
見耶和華你的神.
為天下萬民所擺立的日月星.
就是天上的萬象.
自己便被勾引.
敬拜事奉他.
耶和華將你們從埃及領出來.
脫離鐵路.
要特作自己產業的指紋.
像今日一樣.
《新命記》四章二十三節.
你們要謹慎.
免得忘記耶和華你們神.
與你們所立的約.
為自己雕刻偶像.
就是耶和華你神所禁止你.
作的偶像.
因為耶和華你的神.
乃是烈火.
是忌邪的神.
你們在那地駐久了.
生子生孫.
就雕刻偶像.
彷彿什麼形象.
敗壞自己.

$^{41}$恆耶和華你神.
眼中看為惡的事.
惹他發怒.
我今日呼天喚地.
向你們作見證.
你們必在個約.
但何得遺孽的地上.
速速滅盡.
你們不能在那地上長久.
必進行滅除滅.
耶和華必使你們.
分散在萬民中.
在他所領你們到的萬國裡.
你們剩下的人數稀少.
在那裡.
你們必事奉人手所造的神.
就是用木石造成.
不能看.
不能聽.
不能吃.
不能聞的神.
但你們在那裡.
必尋求耶和華你的神.
你盡心盡性尋求他的時候.
就必尋見.
日後你遭遇一切患難的時候.
你必歸回耶和華你的神.
聽從他的話.
耶和華你神.
源是有憐憫的神.
他總不撇下你.
不滅絕你.
也不忘記.
他起世與你滅祖所立的約.
弟兄姊妹平安.
再繼續我們.
《新命記》這個系列的宣講.
剛才那段經文其實你都會發覺.
就好像今天我們這個講題.
很清楚了.

$^{81}$是講與偶像有關的事情.
雕刻偶像的事情.
事實上.
來到《新命記》這個段落.
你會發覺摩西開始訴諸於.
以色列民和上帝所立的約.
透過約裡面的內容.
提醒第二代的以色列人.
他們在進入迦南的時候.
要謹記一些事情.
包括雕刻偶像的事情.
摩西提醒以色列人.
好像經文所說.
當日上帝在河列山.
與百姓立約的時候.
百姓其實只聽到的是什麼?.
聲音.
他們只聽見神的聲音.
其實沒有看見過.
上帝的形象到底是怎樣.
上帝從烈火裡面對百姓說話.
將祂的律例典章教訓百姓.
要求他們進入迦南的時候.
得地為業的時候.
就在當地遵守上帝的教訓.
藉此在萬民裡面.
顯出與眾不同的智慧.
為上帝作見證.
由此可見.
我們在這裡看到.
以色列人的信仰.
與外邦人有一個很明顯的分別.
就是他們敬拜的上帝.
看不看得見?.
其實是看不見的.
他們是敬拜一個.
看不見形象的上帝.
他們不像當時外邦人.
他們膜拜偶像.
一定要有一個實物.

$^{121}$一定要製造一個.
看得見的神像來敬拜.
換句話來說.
以色列人的上帝之所以與眾不同.
不是因為上帝耶和華.
本身看上去有什麼獨特.
因為看不見.
不是因為耶和華本身身光頸靚.
高大衛猛.
上帝說我沒有顯露過.
我的形象給人看.
以色列人之所以與眾不同.
是什麼呢?.
是因為他們自己.
透過遵守上帝的律法.
而與眾不同.
所以不是上帝看上去.
有一個形象與別不同.
是上帝要求他的子民.
分別為性.
當他們進入迦南的時候.
就要看起來不一樣.
意思是什麼呢?.
上帝說得很清楚.
就是遵守上帝的吩咐.
遵守上帝的吩咐的時候.
你就在萬民中分辨出來.
有不一樣的社會秩序.
不一樣的道德標準.
不一樣的敬拜原則.
這個我們經常說.
以色列人如何在萬民中.
有分別的意思.
換句話來說.
假如上帝的子民.
真的抵受不了誘惑.
他仿效迦南人.
膜拜偶像.
膜拜一個看得見的神像.
這種習性.

$^{161}$很自然他就會失去什麼呢?.
失去了他的獨特性.
你想像一下.
大家在迦南一樣生活.
一樣耕種.
他有拜他的神像.
你也有耶和華的神像.
其實沒有分別.
大家都是用同一種敬拜的方式.
換句話來說.
假如以色列人.
抵受不了雕刻偶像的試探.
他們就不能夠在萬民中分辨出來.
亦因為這樣.
他們也破壞了和上帝索納的約.
因為上帝要求他們.
要與別不同.
在萬民中作見證.
既然破壞上帝的約.
自然上帝會履行約的吩咐.
上帝會管教.
甚至是原例管教.
收回他對子民所應許的一切.
所以你會看到摩西.
為什麼要在這裡鄭重提醒新一代的百姓.
不要禁戒雕刻偶像.
特別當以色列人進入迦南定居之後.
你會發覺他們的生活方式會大大改變.
過去三十多四十年.
他們在抗疫是做什麼的?.
怎樣謀生?.
考考大家.
雖然我們還未讀到初一及二.
他們是靠上帝降嗎哪.
不用耕種.
每天按上帝吩咐.
每天撿一份就行了.
四十年都是這樣.
但現在進入迦南又如何?.
你會發覺上帝稍後就會停止降嗎哪.

$^{201}$他們要很務實地耕種.
他們要轉型用農耕種植的方式生活.
自然對天氣對氣候的需求一定會大增.
也很容易受迦南當地的宗教引誘.
敬拜供奉當地人以為能夠掌管天氣氣候的偶像假神.
事實上這裡說偶像敬拜是什麼回事呢?.
是說人試圖用一個實物.
用一個物件來理解或代替一些比較抽象的觀念或事物.
舉個例子.
有沒有人用iphone?.
iphone有個形象是什麼?.
當初推出的時候.
就是簡約.
以前的手機有很多按鈕.
很複雜的.
現在一顆圓形就能搞定.
如果純粹說iphone的概念很簡約.
很簡樸的.
那你會怎樣?.
你不會理解的.
直到實物出現.
你就會知道什麼是簡約.
一顆圓形就能搞定.
這個就是所謂將比較難理解的抽象理念.
去具體化形象化視覺化.
你見到就會明白.
如果用當時的迦南文化.
迦南人會普遍會膜拜一個女神.
叫做阿瑟拉或阿斯他祿.
但他們不會在意阿斯他祿是怎麼來的.
不會理會他的由來.
最重要的是他的功能是什麼?.
偶像敬拜是強調功能.
阿斯他祿或阿瑟拉是掌管生殖和生產的.
於是迦南人如何將這個概念具體化?.
幫自己明白.
為阿斯他祿雕造一個擁有女性生殖器官的實體神像.
並且他看到這個神像就知道.
他是負責掌管生產和生育的.
就將這個概念具體化了.

$^{241}$看得見.
每天他看到這個阿斯他祿的神像.
他就會想起.
只要我敬拜供奉他.
他就會保佑我農作物可以生產.
保佑我人丁興旺.
事實上偶像敬拜很強調的就是這種功能性.
所以摩西告誡百姓.
不可以這樣做.
不可以奉行當地人的偶像敬拜.
為自己雕造偶像.
第十七節開始說.
不可以將你眼所見的想像幻化成為一個實體敬拜的木像.
不可以雕造.
無論是男的像也好,女的像也好.
飛鳥,走獸,爬蟲,魚也好.
不可以雕個偶像來供奉他們.
也不可以看天空.
看到星星,月亮,太陽.
這些比較難去雕造.
不可以去敬拜天像.
也不可以.
有些聖經學者注意到.
剛才說的偶像排列.
不可以雕刻人類.
動物,飛鳥,魚類,日月,星辰.
其實你會發覺和創世記第一章上帝創造的秩序是怎樣.
基本上是相反的.
反過來.
似乎摩西是想說出一件事.
假如人去敬拜自己雕刻出來的偶像.
其實就是違反了上帝創造的秩序.
你搞亂上帝的秩序.
當人搞亂上帝的秩序.
他自己的生活也會怎樣.
會亂七八糟.
會失去秩序.
出錯.
就是這麼簡單.
所以上帝為什麼禁止以色列人去雕刻偶像.

$^{281}$就算你說.
我現在雕一個偶像.
我叫他做事.
我敬拜他.
我沒有敬拜別神.
我真的敬拜上帝.
不過我有一個對象.
有一個實物.
可以嗎.
不可以.
也不可以.
以色列人就是因為這樣.
曾經敗壞自己.
幾乎搞到滅亡.
事實上你看文素上一代的以色列人就是這樣.
當時摩西和上帝在山上四十晝夜.
遲遲都未下來.
山下的百姓又開始怎樣.
開始害怕.
摩西不回來是否已經有什麼不測.
上帝又不出聲.
不如這樣.
阿倫就說我們做一個像.
我們收集一些金器出來.
我們打造一隻牛犢.
我們叫他做上帝.
耶和華.
這個就是耶和華.
他帶領我們.
他也是在敬拜耶和華.
但是卻是惹來上帝的震怒.
並且刑罰他們.
換句話來說.
即使你雕刻一個偶像.
叫他做上帝.
耶和華.
我供奉他.
都不可以.
原因是摩西十五節說得很清楚.
上帝和我們立約的時候.

$^{321}$並沒有顯現.
他實際是什麼模樣.
他是通過說話.
讓我們認識他的旨意和要求.
我們是這樣和上帝建立關係.
不是通過形象.
是通過上帝的聲音和說話.
換句話來說.
沒有人可以憑自己肉眼.
建構想像出上帝的形象.
任何人間想像出來的偶像.
就算你叫他做耶和華上帝也好.
都只會是一個不完全矮化了.
甚至扭曲了的上帝.
代表你對上帝的完美.
是一種虧欠.
造成一種褻瀆.
很清楚.
我們不難理解.
不過.
知歸知.
人總是明知不可為.
而為之.
為什麼呢?.
因為偶像的敬拜.
確實有一種很強大的吸引力.
就是他會給人一種安全感.
掌控的感覺.
又再舉個例子.
我們說迦南人除了敬拜阿斯塔諾這個女神之外.
都會敬拜她的丈夫.
一個男性的神.
叫做巴力.
你熟悉的巴力.
你見到在舊約.
經常處理敬拜巴力的問題.
迦南人知道自己沒有能力掌控大自然.
他沒有這個能力.
但可以怎樣呢?.
我可以雕一個偶像出來.

$^{361}$雕一個實體的木像.
我叫他做巴力.
我賦予巴力他的能力和權力.
巴力你就是掌管世間的風風雨雨.
降雨.
風調雨順.
就是巴力你去控制.
然後我就透過膜拜供奉巴力你.
心理上我會怎樣呢?.
我會安穩一點.
因為我可以控制巴力.
我可以得到掌握風雨的感覺.
這樣實在一點.
看天色做人是很難的.
因為你控制不到.
你不知道何時下雨.
對你的龍更修成是很有影響.
現在我有個巴力就行了.
我敬拜他.
我控制到他的時候.
我就確保我的生活安穩.
個人自然會踏實很多.
事實上不只是以色列人有這個情況.
中國人民間宗教.
你會發覺其實都是提供類似效果.
用什麼方式呢?.
例如可能.
抬一隻燒豬去.
抬一隻燒肉去.
不夠要抬什麼呢?.
抬一隻燒豬去.
那代表什麼呢?.
我再神針一點.
用賄賂交易的方式.
來討好他所相信的神明.
來換取祝福.
又或者不用這招.
我給人家加倍努力敬虔.
通常是過年那幾天.
就會一早去某間廟排隊.

$^{401}$比上班逼地鐵還厲害.
用盡渾身解數.
撞開逼開.
旁邊的人是誰呢?.
那一刻你會當他是敵人.
對手.
撞開之後.
我就可以插到頭炷香.
總之用我的誠心.
打動我所相信的神明.
但背後其實都是一種操控的方法.
透過他很想操控他心裡所想的神明.
來換取自己的需要.
事實上.
所有雕刻偶像的敬拜.
背後都有慾望操控的元素.
通過掌控自己建構出來的偶像.
滿足自己物質或心靈的需要.
這樣你生活就多一份掌控的感覺.
多一份安全感.
如果從這個角度來看.
你和我今天打死都不會自己.
做一個偶像出來敬拜.
都不敢.
我們有沒有雕刻偶像的問題?.
其實我們都會有.
我們仍然要面對偶像敬拜這方面的事態.
正如已故的提摩泰凱勒牧師提摩太摩泰凱勒牧師曾經說過.
他說人類本身就是不斷製造偶像.
追求自己的假神.
現在是一間工廠.
每個人心裡多多少少都有這種心態.
因為很想滿足自己的需要.
在他的著作《山寨版上帝》.
你可能也有看過.
他說到什麼是偶像呢?.
凡是你覺得比上帝更加重要.
你很想.
那樣東西很吸引你.
並且你很想從他身上得到的滿足.

$^{441}$比上帝的滿足更加大.
假如這個事物一旦失去.
你就會混亂.
你會忐忑.
你會失去做人的分寸和方向.
很可能.
那就是你追求的偶像和假神.
我曾經看過一篇文章.
《誰是你心中的現代偶像》.
當然不是說明星和歌星.
是從信仰的角度去看.
這篇文章提到幾年前美國機構有一個調查.
訪問很多教會的牧者.
他們想了解一下.
影響基督徒最深遠的現代偶像是什麼呢?.
你猜猜是什麼呢?.
給你五秒時間想想.
影響基督徒最深遠的偶像是什麼呢?.
你一定猜有錢吧?.
開估吧.
第一位,六成七.
是舒適的生活.
你看看是不是.
第二位,也有五成六.
就是我剛才說的掌控的感覺.
安全的感覺.
第三位才是金錢,也有五成五.
第四位,可能你沒想過.
原來我們爭取別人的認同.
很容易都會變成一種偶像的敬拜.
你很想得到別人的認同.
很想得到別人的稱讚.
一過了龍.
就會變成你偶像的敬拜.
原來這些都成為了基督徒潛藏的偶像.
影響作為一個信徒的價值觀念.
其實很值得我們注意.
舉個例子.
旺鎮是不是一間很蒙福的教會?.
你不敢說不是吧?.

$^{481}$你坐在這個地方已經代表你很蒙福.
如果你有空看一看旺鎮的歷史.
你就知道上帝很祝福我們.
雖然我們一直從小的地方搬到大的地方.
不同的變遷.
但是上帝一直都讓我們越來越壯大.
到今天我們有一個很舒適.
很安穩.
屬於自己敬拜的地方.
聚會的場所.
這件事好不好?.
當然好.
是來自上帝的.
但是如果這種安逸舒適的教會生活.
不知不覺成為了教會的潛藏偶像.
我們只顧著追求.
總之我們要坐得舒適.
總之要招呼周到.
就會變成什麼?.
就會歪了.
這不是上帝要我們追求的目標.
我們就變成離開上帝的使命.
追求了我們自己心裡的偶像.
事實上當我們越是追求安逸舒適的生活的時候.
逐漸我們減退的就是什麼熱誠?.
我們犧牲社稷.
服侍他人的熱誠就會減退.
必然的.
基督徒一旦過度追求生活的舒適安逸.
受虧損的一定是上帝.
受影響的一定是上帝的吩咐.
所以.
今日這段經文提醒我們.
需要懂得分別什麼是我們心裡潛藏的偶像.
正如耶穌所說.
你的心在哪裡?.
下一句.
你的財寶就在哪裡?.
很實際.
當你花最多的時間.

$^{521}$心思努力.
甚至金錢.
來到效力的對象.
往往一定是你心裡的摯愛.
當然聖經所說的應該是誰?.
是上帝.
是耶穌.
但聖經也同時說.
那個對象很可能是其他的東西.
取代了耶穌基督最重要的地位.
那是什麼?.
香港人最喜歡的是什麼?.
打完風就想起上班.
你不會見全世界這麼奇怪.
打完工就想起車子到哪.
排隊排隊.
結果排了一個小時都沒等到班車.
上班很特別.
上班背後代表什麼?.
是一種安穩掌控的生活.
每天都做重複的事.
然後有安穩的收入.
安穩的機會讓你發揮貢獻能力.
其實是很大的操控和慾望的需要.
如果這東西超越了上帝.
他就會變成你的偶像.
又或者手機.
不用說了.
我們香港人或基督徒滾手機的時間.
很可能比看聖經的時間多.
又或者有些人很著重子女的前途.
子女的學業.
每逢考試測驗.
就會怎樣?.
我要幫他請假.
我要幫他向教會請假.
向主日學請假.
那他會不會變成偶像呢?.
很值得我們去想.
又或者香港人購物,旅遊,消費.

$^{561}$不妨檢視一下.
會不會這些都成為我們的偶像.
正如我自己.
我很喜歡高達.
我很喜歡模型.
這是我的嗜好.
這是我的興趣.
但如果我每天都砌高達.
每天都模型.
又不奉獻.
又不寫講章.
又不崇拜.
那已經不是嗜好.
不是興趣.
是你的偶像.
即使我沒有拿著高達不斷擺放.
那也是我的偶像.
所以弟兄姊妹.
當我們不妨檢視一下.
你每天的時間分配.
你的生活消費模式的時候.
你慢慢會辨識到.
什麼是潛藏在你心裡的偶像.
而另一個很有效的方法.
就是你試一下觀察一下.
你是怎樣應對那些落空了的期望.
那些上帝沒有應運你的禱告.
譬如說.
假如你很想成就一件事.
但最終沒有成功.
你當然會不開心.
你也會失望.
人之常情.
但是不會影響你和上帝的關係.
你繼續好好按照聖經吩咐.
為上帝而活.
表示這件事情.
雖然你很想追求.
但它沒有成為你生命最重要的主人.
相反.

$^{601}$你很努力去達成一件事.
或者一個願望.
你又做到最盡.
又禱告.
很想心願達成.
但是事與願違.
於是你開始埋怨你身邊的人.
甚至埋怨上帝.
你開始憤怒.
又或者感到驚慌絕望.
這一刻留意.
你心裡真正效力服務的偶像.
可能就是現形.
你不知不覺已經在服務他.
表面上好像是你控制他.
但最終原來他反過來控制你.
所以,靈姐妹.
求主幫助我們.
我們要懂得去分辨.
哪些是我們心裡的偶像.
並且我們求主幫助我們.
要除去這些假神.
使我們可以重新將人生的焦點.
放回上帝身上.
正如第二十節裡說.
耶和華將你們從埃及領出來.
脫離了鐵路.
要特作自己產業的子民.
將今天一樣.
鐵路這裡是說以色列人.
昔日在埃及圍爐被轄制的生活.
就好像在鐵路裡.
被火燒的鐵.
每天都煎熬.
很慘,不得自由.
不過好消息是.
上帝已經施行了拯救.
使得以色列民已經不再需要.
被這些偶像來轄制.
這些生活中他們自己.

$^{641}$營造出來的偶像所管轄.
他們現在只需要專一服侍上帝.
過往他要服侍無數的偶像.
因為無數的偶像.
他覺得要服侍他.
才能給他平安安穩.
上帝說不用了.
不用了,你尊重我.
你只需要專一效忠服侍上帝.
脫離所有的偶像.
那你就可以經歷真正人生的福氣.
事實上,這就是上帝一直提醒.
百姓要謹記的事.
而另一個.
上帝要求他們.
斷絕所有雕刻的偶像.
取決於上帝一個的屬性.
二十四節這裡說.
上帝表示自己是.
烈火,還有.
忌邪的神.
上帝訴諸於他的屬性.
來到警戒以色列人.
要敬畏他.
千萬不要試探上帝.
偶像的敬拜不是開玩笑的.
百姓若敢這樣做.
一定是得罪上帝.
等於引火自焚.
烈火也反映了上帝的.
那種的聖潔和公義.
他不容許有偶像假神.
佔據他子民的心.
他一定會將他憎惡的事情消滅.
因為上帝說自己是.
忌邪的神.
留意忌邪這個字.
黃學本翻譯為忌邪.
並不代表上帝很忌諱.
很骯髒的東西.

$^{681}$不像廣東人和中國人.
很忌諱這些邪的東西.
不要碰它,不要碰它.
不是那些忌諱的意思.
英盟約這個字翻譯為.
Jealous.
妒忌,嫉妒.
聽起來好像不太好.
我們不喜歡嫉妒人.
為什麼要妒忌別人比你厲害.
比你幸福.
不好的玻璃心.
其實嫉妒好像不好.
這種情感反應不是很好.
不過在一段親密的關係裡.
嫉妒卻是發揮了保護的作用.
因為嫉妒就代表排他.
在一段關係裡不容許有第三者的存在.
意思是什麼呢.
假如一段精忠的愛情關係.
被發現有小三的時候.
你是配偶,你會怎麼做.
假如那個想出去搞小三的.
男或女都好.
他問.
我可不可以偶爾出去結識一下異性.
發展一下情,可以嗎.
你放心,我保證.
你永遠是我心裡的摯愛.
你在我心中是第一位的.
不會動搖的.
他這樣說你會有什麼反應.
你會不會說.
你偶然逢場俗興.
無傷大雅.
最重要的是你心裡還有我.
最緊要的是你心裡還有我.
這樣就行了.
行不行.
當然不行.

$^{721}$你是他,你會怎麼做.
你聽到你的配偶想出去搞小三的時候.
你當然像預場片那樣.
扭他的耳朵.
斬釘切鐵地扭下去.
你不用指望.
你試一下都不行.
想一下都不好.
你試一下出去搞小三.
我一定怎樣.
我和你拼了,我不會放過你.
你看得多了,預場片.
為了維護這段婚姻關係.
你一定會嚴厲訓斥對方.
你禁止任何會破壞這段關係的行為.
弟兄姊妹,這就是聖經裡說的.
嫉妒,忌邪的意思.
上帝說他是忌邪.
嫉妒的上帝代表他和以色列民的關係排他.
不容許他和上帝子民之間多了一個偶像.
多了一個第三者.
霸佔了子民的心.
百姓不能夠在信仰裡捏花野草.
搞三搞四.
不過以色列人的問題是什麼呢.
他們不是完全不相信上帝.
而是什麼呢.
聖經裡經常說.
信仰裡的不忠.
對上帝不夠專一忠誠.
又敬拜事奉上帝.
但同時又敬拜偶像.
左右逢源.
反映當時他們要面對的風氣是多神的世界.
多神的主義.
事實上多神主義在古代是一個很共通的文化.
沒有人會只相信一個神這麼吃虧.
他們覺得是吃虧的.
當然是多一點.
就可以多拜多一點神明.

$^{761}$自然有更多機會獲得祝福.
如果這個不行.
那怎麼辦.
不要緊.
還有千千萬萬的假神.
我可以去找他們幫忙.
這個不行.
還有下一個.
這個也不行.
還有另一個.
再不行我自己建構一個出來.
總有一個滿足到我的需要.
事實上以色列人進入到迦南.
他們也沒有辦法完全擺脫這種思潮.
想像一下一個很現實的場景.
剛才也說過.
以色列人在曠野是靠嗎哪來維生.
其實他們經歷了三十多四十年.
應該不懂得種植.
沒有耕種的經驗.
他們進入到迦南定居的時候.
看到當地人很自然就會.
看到他們可能有收成.
他們有種植.
就會過兩天來看看有什麼心得.
讓我也有些收成.
我又不懂得種植.
剛學種植.
你一定會去問他們.
這個時候迦南人自然會說什麼.
他們不會說我們是怎樣下肥料.
怎樣殺種.
一定是說我們拜巴力.
我們拜阿斯塔諾.
你處理好這兩個神明.
討好他們的歡心.
他們會保佑你豐收.
一定是這樣去解說.
不過當他們這樣說的時候.
以色列人第一刻一定會警醒.

$^{801}$不可以.
上帝說過.
只可以敬拜他.
不可以刁黑偶像.
這樣上帝會生氣.
會得罪他.
但是對一個多神主義的迦南人來說.
有沒有問題.
沒有問題.
沒有人叫你放棄上帝.
他仍然是你的摯愛.
不過你認識多一個.
拜多一個偶像.
又有什麼所謂.
多個人保佑你不就更好.
其實就是這樣慢慢地滲.
動搖以色列人對上帝的忠誠.
好像一個包著糖衣的毒藥.
慢慢地滲.
今天偶像敬拜也是這樣.
沒事的.
沒有人叫你不回教會.
沒有人叫你不用聖經.
不過.
謀生也很重要.
考試也很重要.
那些東西都不可以放.
慢慢地上帝就不是你的第一.
所以弟兄姊妹.
上帝要求他的子民.
對他的忠誠度是百分百.
不可以打折扣.
上帝的話.
因為一打折扣又怎樣.
這份關係就會變質.
就會走樣.
連帶上帝的話語都開始不被尊重.
失去了在百姓身上那種影響力.
和神的關係自然也會隨流失去.
結果就好像27至28節那裡說.

$^{841}$我想大家讀一讀.
27弟兄姊妹也有印了.
最後那段經文.
只是讀27至28節.
好,預備起.
(教會禮儀).
好,謝謝.
結果你看到了.
上帝一定會出手.
上帝雖然包容了他們很久.
但最終也會出手.
重罰這群在信仰上不忠的百姓.
將他們從被擄.
從被擄到真的拜偶像的地方.
求仁得仁.
現在你這麼想拜偶像.
現在你日對夜對.
日日去供奉對著那些.
用木頭做,用石頭做.
但是不能看見,不能聽見.
不能吃,不能聞的偶像.
你事奉他吧.
你這麼想他們的時候.
留意這裡,上帝說什麼.
他們會人數變得稀少.
這個和上帝的心意完全相反.
上帝對亞伯拉罕的應許是.
你們人數會加增.
但是當百姓不活在.
上帝的盟約裡的時候.
上帝說我收回,人數稀少.
不過,這並不是以色列人的忠告.
上帝表示他們可以有不一樣的選擇.
仍然有轉機.
他們可以選擇另一條路.
就是跟著29至31節說.
你悔改,你脫離你的惡行.
轉回上帝裡面.
盡心盡性尋求上帝的吩咐.
聽從他的話.

$^{881}$這樣上帝都會怎樣.
上帝都會回心轉意.
將你們在被擄之地救出來.
因為上帝自己說.
他是那位有憐憫的神.
他不會丟下.
不會撇下,不會滅絕.
這一班的子民.
因為他從來沒有忘記過.
他向列祖所立的誓.
所以,弟兄姊妹.
看起來好像沒有矛盾.
上帝一時又愛他們.
一時又罰到他們死.
但這就是上帝的愛.
上帝對子民的愛是認真負責任.
就是會這樣.
既有嚴謹的要求.
但同時也有恩典和憐憫.
我常常想起.
我們跟小朋友相處.
我記得我兒子還在小學的時候.
我常常要跟隨他的功課和考試.
你猜我做牧師會生氣嗎.
當然會生氣,只是你看不到.
其實很生氣.
你叫他做就做.
溫習一下.
當然不會盡力.
開始會生氣.
失去了掌控的感覺.
開始不安心,就會生氣.
所以我說很容易會變成偶像.
但不是這個意思.
當我心情緊張的時候.
有時會動氣,有時會用重言語.
事後會後悔,跟他道歉.
所以有時平心靜氣.
我會用另一個方式跟他相處.
其實我以前問他.

$^{921}$其實你想不想爸爸這樣罵你.
罵得你這麼狠.
其實我自己也不開心.
有時也想算了.
不要破壞關係和諧.
不如我不要理會你的事.
你猜他說好不好.
他說好啊.
你爸爸理會我,不要理我.
我說不會,他沒有.
他說這樣也不行.
你不提醒我,我搞不定.
你要提醒我,不過不要這麼兇.
做爸爸很難做.
弟兄姊妹,這就是上帝對我們的.
當小朋友在欲緊這段關係.
他知道這段關係是真的.
愛他的人,不只是純粹疼他.
是會負責任,會提醒他,會管教他.
但也有恩典憐憫.
管教是不開心的.
但他不會放棄這份關係.
因為他知道爸爸疼他.
同樣以色列人和上帝的關係也是這樣.
上帝對他們的愛是認真負責任.
當他們行差踏錯,甚至不忠的時候.
上帝不會不理會.
如果上帝這樣也不理會.
就代表他不愛這群子民.
上帝一定基於他的愛.
他會管教,甚至嚴厲,不留情面的管教.
目的是想彰顯他那份敬愛.
愛之深責之切,我們不難理解這個道理.
而當百姓願意悔改.
尋求他,聽從他的心意的時候.
上帝都會施恩憐憫.
鬆開他的手,讓子民得到拯救.
事實上,弟兄姊妹.
今日對經文強調.
我們和以色列人一樣.

$^{961}$我們享有獨一無二的身份.
亦有一個使命.
上帝期望我們不是在人世間.
靠自己的努力尋求生活幸福.
而是在上帝的旨意裡尋求豐盛.
你活出上帝的心意.
那才是討上帝喜悅的事情.
上帝也自然會供應.
祝福你世上所需要.
不是靠你自己去餵.
而是靠上帝的供應.
他說過一定會供應救濟.
我們要做的是過一個不一樣的生活.
如何在這個世界裡.
透過進行聖經的教導.
過一個分裂為勝的生活.
上帝期望我們.
好像以色列人一樣.
我們要讓人看得見.
別人看不見上帝.
但透過我們.
會見到上帝是怎樣.
見到上帝對這個世界的心意是怎樣.
也見到一群對上帝忠心的人.
他們應該有什麼行為.
正因為這樣.
聖經常常提醒我們.
好像今天的講題說.
你要防範那些在你生命裡出現的偶像.
你要拒絕那些影響你.
專一愛上帝的假神.
而堅持這一切屬靈教導的.
就是上帝對我們不離不棄的愛.
所以弟兄姊妹.
求主加給我們力量和勇氣.
幫助我們勝過正面對偶像的試探.
無論如何.
你都堅持追求活出聖經的真理.
以這成為你愛上帝.
最實際的證據.

$^{1001}$我們求心禱告.
因為你已經將我們從罪的綑綁裡.
從偶像的綑綁裡.
其實你已經拯救了我們.
然而在你再來的日子之前.
我們仍然在地上生活的時候.
難免我們會面對很多人有試探.
很多事情本來其實都是好的.
但是在罪人的眼中.
慢慢地那些就會變成.
不知不覺變成我們的偶像.
我們追求他們.
我們來服侍他們.
過於我們追求服侍上帝的心.
而我們不知不覺被那些事情吃勢.
我們遠離上帝你.
忘記了你的吩咐.
然而上帝你很好.
因為你愛我們.
你對我們的愛負責任.
你會親自提醒我們.
讓我們知道原來由始至終.
我們是靠上帝你的恩典.
來守護我們和你的關係.
其實不是靠我們自己.
正因為這樣我們知道這段關係.
最終一定穩妥.
因為上帝你傳言的掌管.
保守和提醒.
特別今天你再次藉著你的話語.
藉著聖靈提醒我們.
要禁戒偶像的事情.
要專一地來信靠你的時候.
主讓我們有勇氣.
有從你而來的智慧.
我們懂得去分辨.
我們懂得去剔除.
這些纏繞我們的偶像.
我們求主求你的令.
也幫助我們.

$^{1041}$讓我們真的有這份堅持.
有這份力量.
來抵抗偶像的試探.
讓我們的心百分百歸向你.
榮耀你.
特別我們將每一位.
弟兄姊妹交在你的恩手.
求主引你.
你知道我們每一個的心思意念.
我們的狀況是怎樣.
我們正面對的是什麼挑戰試探.
願你剔除我們裡面的污穢.
讓我們真的成為聖結.
合乎主用.
為你作見證.
我們這樣禱告交託.
奉基督耶穌你得勝的命.
來祈求.
阿們.
\newpage



\section{申命記 5:1-6:9}
\label{sec:bBtvGdEATA0}
\textbf{2025年2月23日青少主日崇拜直播｜楊珊珊傳道｜朝向應許之地－十誡與大誡命｜申命記5\_1-6\_9}
\newline
\newline
連結: \href{https://youtube.com/watch?v=bBtvGdEATA0}{\texttt{https://youtube.com/watch?v=bBtvGdEATA0}} ~~~~ 語音日期: 2025-02-23
\newline
\newline
\hyperref[sec:oex63GVBJWc]{\small{< < < PREV SERMON < < <}}
~
\hyperref[sec:index_chronic]{\small{[返順時目]}}
~
\hyperref[sec:C0mf0rNYQN4]{\small{> > > NEXT SERMON > > >}}
\newline
\newline
申命記 5:1-6:9
\newline
\begin{longtable}{cl}
\hline
\hline
章節 & 經文 (和合本修訂版)\\
\hline
5:1 & \begin{tabularx}{0.7\textwidth}{X} 摩西召集以色列眾人，對他們說：「以色列啊，要聽我今日在你們耳中所吩咐的律例典章，要學習，謹守遵行。 \end{tabularx} \\ \\ \relax
5:2 & \begin{tabularx}{0.7\textwidth}{X} 耶和華－我們的神在何烈山與我們立約。 \end{tabularx} \\ \\ \relax
5:3 & \begin{tabularx}{0.7\textwidth}{X} 這約耶和華不是與我們列祖立的，而是與我們，就是今日在這裡還活著的人立的。 \end{tabularx} \\ \\ \relax
5:4 & \begin{tabularx}{0.7\textwidth}{X} 耶和華在山上，從火中，面對面與你們說話。 \end{tabularx} \\ \\ \relax
5:5 & \begin{tabularx}{0.7\textwidth}{X} 那時我站在耶和華和你們之間，要將耶和華的話傳給你們，因為你們懼怕那火，沒有上山。他說： \end{tabularx} \\ \\ \relax
5:6 & \begin{tabularx}{0.7\textwidth}{X} 「『我是耶和華－你的神，曾將你從埃及地為奴之家領出來。 \end{tabularx} \\ \\ \relax
5:7 & \begin{tabularx}{0.7\textwidth}{X} 「『除了我以外，你不可有別的神。 \end{tabularx} \\ \\ \relax
5:8 & \begin{tabularx}{0.7\textwidth}{X} 「『不可為自己雕刻偶像，也不可做甚麼形像，彷彿上天、下地和地底下水中的百物。 \end{tabularx} \\ \\ \relax
5:9 & \begin{tabularx}{0.7\textwidth}{X} 不可跪拜那些像，也不可事奉它們，因為我耶和華－你的神是忌邪的神。恨我的，我必懲罰他們的罪，自父及子，直到三、四代； \end{tabularx} \\ \\ \relax
5:10 & \begin{tabularx}{0.7\textwidth}{X} 愛我、守我誡命的，我必向他們施慈愛，直到千代。 \end{tabularx} \\ \\ \relax
5:11 & \begin{tabularx}{0.7\textwidth}{X} 「『不可妄稱耶和華－你神的名，因為妄稱耶和華名的，耶和華必不以他為無罪。 \end{tabularx} \\ \\ \relax
5:12 & \begin{tabularx}{0.7\textwidth}{X} 「『當守安息日為聖日，正如耶和華－你神所吩咐的。 \end{tabularx} \\ \\ \relax
5:13 & \begin{tabularx}{0.7\textwidth}{X} 六日要勞碌做你一切的工， \end{tabularx} \\ \\ \relax
5:14 & \begin{tabularx}{0.7\textwidth}{X} 但第七日是向耶和華－你的神當守的安息日。這一日，你和你的兒女、僕婢、牛、驢、牲畜，以及你城裡寄居的客旅，都不可做任何的工，使你的僕婢可以和你一樣休息。 \end{tabularx} \\ \\ \relax
5:15 & \begin{tabularx}{0.7\textwidth}{X} 你要記念你在埃及地作過奴僕，耶和華－你的神用大能的手和伸出來的膀臂領你從那裡出來。因此，耶和華－你的神吩咐你守安息日。 \end{tabularx} \\ \\ \relax
5:16 & \begin{tabularx}{0.7\textwidth}{X} 「『當孝敬父母，正如耶和華－你神所吩咐的，使你得福，並使你的日子在耶和華－你神所賜給你的地上得以長久。 \end{tabularx} \\ \\ \relax
5:17 & \begin{tabularx}{0.7\textwidth}{X} 「『不可殺人。 \end{tabularx} \\ \\ \relax
5:18 & \begin{tabularx}{0.7\textwidth}{X} 「『不可姦淫。 \end{tabularx} \\ \\ \relax
5:19 & \begin{tabularx}{0.7\textwidth}{X} 「『不可偷盜。 \end{tabularx} \\ \\ \relax
5:20 & \begin{tabularx}{0.7\textwidth}{X} 「『不可做假見證陷害你的鄰舍。 \end{tabularx} \\ \\ \relax
5:21 & \begin{tabularx}{0.7\textwidth}{X} 「『不可貪戀你鄰舍的妻子；也不可貪圖你鄰舍的房屋、田地、僕婢、牛驢，以及他一切所有的。』 \end{tabularx} \\ \\ \relax
5:22 & \begin{tabularx}{0.7\textwidth}{X} 「這些話是耶和華在山上，從火焰、密雲、幽暗中，大聲吩咐你們全會眾的，再沒有加添別的話了。他把這些話寫在兩塊石版上，交給我。 \end{tabularx} \\ \\ \relax
5:23 & \begin{tabularx}{0.7\textwidth}{X} 山被火焰燒著，你們聽見從黑暗中發出的聲音，那時，你們各支派的領袖和長老都挨近我。 \end{tabularx} \\ \\ \relax
5:24 & \begin{tabularx}{0.7\textwidth}{X} 你們說：『看哪，耶和華－我們的神將他的榮耀和他的偉大顯給我們看，我們也聽見他從火中發出的聲音。今日我們看到神與人說話，人還活著。 \end{tabularx} \\ \\ \relax
5:25 & \begin{tabularx}{0.7\textwidth}{X} 現在這大火將要吞滅我們，我們何必死呢？若再聽見耶和華我們神的聲音，我們就必死。 \end{tabularx} \\ \\ \relax
5:26 & \begin{tabularx}{0.7\textwidth}{X} 凡血肉之軀，有誰像我們一樣，聽見了永生神從火中講話的聲音還能活著呢？ \end{tabularx} \\ \\ \relax
5:27 & \begin{tabularx}{0.7\textwidth}{X} 求你近前去，聽耶和華－我們神所要說的一切話，將耶和華－我們神對你說的話都傳給我們，我們就聽從遵行。』 \end{tabularx} \\ \\ \relax
5:28 & \begin{tabularx}{0.7\textwidth}{X} 「你們對我說的話，耶和華都聽見了。耶和華對我說：『這百姓對你說的話，我聽見了；他們所說的都對。 \end{tabularx} \\ \\ \relax
5:29 & \begin{tabularx}{0.7\textwidth}{X} 惟願他們存這樣的心敬畏我，常遵守我一切的誡命，使他們和他們的子孫永遠得福。 \end{tabularx} \\ \\ \relax
5:30 & \begin{tabularx}{0.7\textwidth}{X} 你去對他們說：你們回帳棚去吧！ \end{tabularx} \\ \\ \relax
5:31 & \begin{tabularx}{0.7\textwidth}{X} 至於你，可以站在我這裡，我要將一切誡命、律例、典章傳給你。你要教導他們，使他們在我賜他們為業的地上遵行。』 \end{tabularx} \\ \\ \relax
5:32 & \begin{tabularx}{0.7\textwidth}{X} 所以，你們要照耶和華－你們神所吩咐的謹守遵行，不可偏離左右。 \end{tabularx} \\ \\ \relax
5:33 & \begin{tabularx}{0.7\textwidth}{X} 你們要走耶和華－你們的神所吩咐的一切道路，使你們可以存活得福，並使你們的日子在所要承受的地上得以長久。」 \end{tabularx} \\ \\ \relax
6:1 & \begin{tabularx}{0.7\textwidth}{X} 「這是耶和華－你們的神所吩咐要教導你們的誡命、律例、典章，叫你們在所要過去得為業的地上遵行， \end{tabularx} \\ \\ \relax
6:2 & \begin{tabularx}{0.7\textwidth}{X} 好叫你和你的子孫在一生的日子都敬畏耶和華－你的神，謹守他的一切律例、誡命，就是我所吩咐你的，使你的日子得以長久。 \end{tabularx} \\ \\ \relax
6:3 & \begin{tabularx}{0.7\textwidth}{X} 以色列啊，你要聽，要謹守遵行，使你可以在那流奶與蜜之地得福，人數極其增多，正如耶和華－你列祖的神所應許你的。 \end{tabularx} \\ \\ \relax
6:4 & \begin{tabularx}{0.7\textwidth}{X} 「以色列啊，你要聽！耶和華－我們的神是獨一的主。 \end{tabularx} \\ \\ \relax
6:5 & \begin{tabularx}{0.7\textwidth}{X} 你要盡心、盡性、盡力愛耶和華－你的神。 \end{tabularx} \\ \\ \relax
6:6 & \begin{tabularx}{0.7\textwidth}{X} 我今日吩咐你的這些話都要記在心上， \end{tabularx} \\ \\ \relax
6:7 & \begin{tabularx}{0.7\textwidth}{X} 也要殷勤教導你的兒女。無論你坐在家裡，走在路上，躺下，起來，都要吟誦。 \end{tabularx} \\ \\ \relax
6:8 & \begin{tabularx}{0.7\textwidth}{X} 要繫在手上作記號，戴在額上作經匣； \end{tabularx} \\ \\ \relax
6:9 & \begin{tabularx}{0.7\textwidth}{X} 又要寫在你房屋的門框上和你的城門上。 \end{tabularx} \\ \\
[1ex]
\hline
\hline
\end{longtable}
$^{1}$今日經訓是在《新命記》五章一至二十一節.
舊約二百二十至二百二十一頁.
以及在《新命記》六章一至九節.
舊約聖經二百二十二頁.
請聽我讀出.
《新命記》五章一至二十一節.
摩西將以色列眾人召了來.
對他們說:以色列人啊.
我今日曉遇你們的六禮典章.
你們要聽.
可以學習.
謹守遵行.
耶和華我們的神在河的山與我們立約.
這約不是與我們列祖納的.
乃是與我們今日在這裡存活之人納的.
耶和華在山上.
從火中面對面與你們說話.
那時我站在耶和華你們中間.
要將耶和華的話傳給你們.
因為你們懼怕那火.
沒有上山.
說:我是耶和華你的神.
曾將你從埃及地遺牢之家領出來.
除了我以外.
你不可有別的神.
不可為自己雕刻偶像.
也不可作甚麼形象.
彷彿上天,下地和地底下.
水中的白物.
不可跪拜那石像.
也不可事奉它.
因為我耶和華你的神.
是既邪的神.
恨我的.
我必追討他的罪.
至父及子直到三,四代.
愛我.
守我誡命的.
我必向他們發慈愛.
直到千代.

$^{41}$不可罔稱耶和華你神的名.
因為罔稱耶和華名的.
耶和華必不以他為無罪.
當照耶和華你神所吩咐的.
守安息日為聖日.
六日要勞碌作你一切的工.
但第七日是向耶和華.
你神當守的安息日.
這一日你和你的兒女.
瀑皮,牛奴,新畜.
並在你城裡寄居的客女.
無論何工都不可作.
使你的瀑皮可以和你一樣安息.
你也要紀念你在埃及地.
作過勞瀑.
耶和華你神用大能的手.
和新出來的傍臂.
將你從那裡領出來.
因此耶和華你的神.
吩咐你守安息日.
當照耶和華你神所吩咐的.
孝敬父母使你得福.
並使你的日子.
在耶和華你神所賜你的地上.
得以長久.
不可殺人.
不可姦淫.
不可偷渡.
不可作假見證陷害人.
不可貪戀人的妻子.
也不可貪圖人的房屋.
田地,瀑皮,牛奴.
並他一切所有的.
然後在新命紀六章一至九節.
在舊約聖經二百二十二頁.
請聽我讀出.
第六章一至九節.
這是耶和華你們神.
所吩咐教訓你們的誡命.
律例典章.

$^{81}$使你們在所有過去.
得遺孽的地上進行.
好叫你和你子子孫孫.
一世敬畏耶和華你的神.
謹守他的一切律例誡命.
就是我所吩咐你的.
使你的日子得以長久.
以色列啊.
你要聽.
要謹守進行.
使你可以在流乃與物之地.
得以享福.
人數極其增多.
正如耶和華你列祖的神.
所應許你的.
以色列啊.
你要聽.
耶和華我們神是獨一的主.
你要重心盡性盡力.
愛耶和華你的神.
我今日所吩咐你的話.
都要記在心上.
也要殷勤教訓你的兒女.
無論你在在家裡.
行在路上.
躺下.
起來.
都要談論.
也要繫在手上為記號.
戴在額上為經文.
又要寫在你房屋的門框上.
並你的城門上.
聖經獨白.
(聖經).
弟兄姊妹平安.
平安.
不知道你們有沒有試過.
與一些非基督徒朋友聊天的時候.
他們會說我們做教徒很不自由.
因為有很多規條要守.

$^{121}$例如不建議與未信者拍拖.
不要有婚前性行為.
不說髒話.
不買六合彩金多寶.
可能你會覺得這些規矩.
人們都覺得是綁手綁腳.
想約我們去行山喝茶.
又拒絕說星期日要回教會崇拜.
弟兄姊妹.
我不知道你有沒有試過.
星期日回教會很掙扎.
累,懶,家人有事.
地鐵服務受阻.
心情低落不想見人.
上班,上學.
生病了要考試.
要去送機.
要去旅行等等.
以上都是我曾經掙扎過的原因.
如果你有共鳴的話.
可以看看我微笑點頭.
起碼都讓我知道我們是不孤單的.
沒錯.
掙扎的原因可能有很多個.
不過回來的原因.
一個就夠.
對於我來說.
就是很想面對面遇見神.
我很想和大家一起.
很真誠地去敬拜我們的上主.
還有經歷祂給我們的恩典和大能.
不知道你唯一回來教會的原因.
又是甚麼呢?.
今天無論是我們的少年人.
還是大人.
還有網上的朋友.
我盼望你今天可以找到.
這個問題的答案.
在我們進入申命記這段很長的經文之前.
讓我們一起祈禱吧.

$^{161}$昔在,今在,永在的耶和華.
我們的上帝.
求你親自向我們每一個人.
啟示你自己.
讓我們能夠面對面認識你.
在社會,學校,工作環境.
家庭裡面.
甚至在教會裡面.
看見你.
並且教導我們用真理和愛.
去引導我們怎樣彰顯你的榮耀.
幫助我們能夠聆聽.
以你的話語為我們生命的中心.
盡心,盡性,盡義去愛你.
並且一生與你同行.
奉主明求.
阿們.
今天我們要討論的十誡.
至今仍然被很多人視為.
一個很明智和良好的道德標準.
即使沒有信仰的人.
都應該會覺得應該要尊重和遵守.
例如不可殺人.
殺人放火就金腰帶.
是不對的.
不可姦淫.
我記得小時候看TVB的劇集.
都說姦夫淫婦要浸豬籠.
不屬於你的.
我們不拿.
不要偷竊.
不要貪心搶別人的錢.
說真相是很難.
但至少不要張開眼睛說謊.
年輕人記得.
做人要真誠.
還有每一個父母都應該很喜歡這條.
當孝敬父母.
但十誡其實不是一套純粹的道德標準.
要記住這些誡命出現.

$^{201}$其實都是因為有一個特定的故事背景.
這個背景就是耶和華救贖以色列人逃離埃及的事件.
《申命記》五章二節提到頒布十誡的場景就是.
耶和華我們的上帝在荷列山與我們一起納藥.
十誡是一個盟約.
十誡是一份合同.
十誡是一個承諾.
大家有沒有去過婚禮.
十誡很比一對新人在婚禮面前讀一個盟誓.
他們會說.
從今以後無論順境逆境.
富裕貧窮.
疾病健康.
都盡從上主的聖法.
愛顧你.
珍惜你.
直到終生.
這是我嚴肅的誓言.
十誡正是耶和華對以色列人所立的三文海誓.
用來維繫雙方關係的條款和內容.
熟悉摩西五經的朋友都知道.
以色列人總共要遵守613條條款.
而十誡就是在這613條條款中.
最重要的那一頭.
所以他就放在首位.
就好像幫我們和耶和華這段關係.
定了立場和生命的位置.
十誡正正就是反映.
上帝期望這群出埃及的以色列人.
擁有一個怎樣的身份.
和他們的生命應該怎樣被上帝塗造.
以前這群以色列民.
不知道怎樣和神建立一個正確的關係.
因為他們生活的劇本和角色設定.
就是塑造為了去事奉法老和埃及的神明.
但現在耶和華要將他們從奴役中釋放.
並且帶領他們去到曠野.
一個渺無人煙的地方.
他們再沒有自己的土地.
亦都沒有社會身份.

$^{241}$更加重要的是他們不再是奴隸.
耶和華正在幫他們建立一個全新的身份.
所以十誡的誡命或律法.
就是神用來重新塑造.
這群以色列民的身份角色.
和他們的社會倫理關係.
重建他們的三觀.
編寫他們的人生故事.
雖然有613條規矩.
但其實上帝的律法不是要我們做到完美.
又或者跟著這些規矩.
我們便會過到一個得蒙救贖的人生.
在猶太教或基督教的傳統裡.
這些律法其實都是為我們好的.
因為當我們選擇活出我們的生活時.
就是要顯出我們是上帝所揀選的子民.
簡單來說.
其實遵守律法不是為了討好上帝.
而是為了展示上帝的形象給全世界看.
我們之前可能被其他人和事物奴役.
不過我們現在可以選擇.
可以選擇活得更加自由.
不需要再被人矮化.
被非人的形象束縛.
我們可以重新擁有上帝的形象.
弟兄姊妹.
你現在的生命處於一個什麼樣的光景呢?.
是只是在生命裡的奴役.
還是察覺到神已經將你從綑綁中釋放出來呢?.
我們不要再回頭看著埃及.
上帝救了我們出來.
我們和上帝是有一份愛的關係.
並且他今天要和我們納藥.
你準備好了嗎?.
不過我猜其實心水清的弟兄姊妹.
你會問.
其實申命記是摩西的臨別贈言.
這番說話明明是和第二代的以色列人說的.
他們很多都在曠野裡出生.
根本沒有經歷過在埃及被奴役的洗禮.

$^{281}$就好像我們這班年輕人.
他們應該從小都上教會.
應該至少活在上帝的盟約下.
應該都沒有什麼問題吧?.
沒錯.
他們是有很豐富的物質.
家長會給他們很多東西學.
坦白說我來了旺鎮半年.
我發覺我們的年輕人真的很厲害.
因為我去過四間教會實習.
沒有一間教會是有青少年的敬拜隊.
只不過我們的年輕人.
仍然都要面對著學校的要求.
朋輩的壓力.
父母的期望.
受著一些社會.
其實有些我跟你根本都不認同的影響.
我很想問在座的年輕人.
又或者是成年人.
你有沒有曾經感受過迷失.
徬徨.
又或者有一份說不出來的無力感呢?.
去年有一個專門做年輕人工作的機構.
「突破」.
它是一個調查發現.
它說現在的年輕人的痛苦.
不再是物質上的缺乏.
而是與其他人的比較.
青少年從小就活在父母的期望當中.
心裡面不其然都會有一把尺.
又或者有一個很嚴厲的法官.
所以在這裡我很想呼籲各位家長和大人.
特別如果你是有年資回來教會的基督徒的話.
我們更加要紮穩我們的馬步.
因為我們是年輕人的第一本聖經.
如果你敬畏上帝的話.
你就每日活出見證.
我求神給我們一班人.
有智慧地幫助年輕人.
知道自己是與這位創天造地的上帝有關係.

$^{321}$而不是用這個社會的生存法則.
去定義他們的身份和價值.
這位耶和華和我們和他們.
今天已經簽了盟約.
就好像以前某某電訊公司那樣.
那份合約刪除也刪不掉.
讓我們彼此鼓勵建立真實的關係.
雖然不容易.
是需要時間.
是需要試.
但是我們這樣才可以顯示一個真實的人生.
告訴年輕人.
可以在這個彎曲薄茂的世代.
可以真實地與上帝和弟兄姊妹.
面對面.
邊境兼併兼.
能夠達得穩.
成為神無瑕疵的兒女.
並且可以在黑暗的地方作炎作光.
新命記五章的一至五節.
是講納若的雙方.
而第二部分.
六至二十一節.
就是十誡的內容.
十誡的意義.
不單只是一組法律.
更是講我們與上帝之間.
有一個深厚的關係的核心.
當我們聽到這一句.
神說:我是你們的主.
這不是單單是十誡的開端.
而是一個強而有力的宣告.
提醒我們這些誡命的發言人是誰.
上帝以祂的權威.
向我們發出這些誡命.
因為祂就是那個將我們從埃及地.
從奴役中拯救出來的上帝.
新命記四章四十節.
摩西這樣說.
摩西今日將神的律例誡命曉如你.

$^{361}$你要遵守使你的子孫可以得福.
並使你的日子.
在耶和華所賜的地上得以長久.
十誡其實是為已經得救的人制定.
不是用來救贖人.
上帝是透過律法重新塑造.
和重建以色列人共同的身份.
故事和價值觀.
前面四個誡命.
講述以色列人和耶和華之間的關係.
而之後的誡命規範了以色列人.
應該怎樣對待其他人.
每一條誡命都講述我們要尊敬耶和華.
同時也要尊敬同樣有神形象的他者.
由於時間有限.
所以我今日想重點講四條誡命.
第一條誡命說.
除了我以外.
你不可有別的神.
這在五章第七節寫了.
因為除了耶和華之外.
我們沒有另一個神救我們出埃及.
事實上如果任何將上帝的關係.
降到第二位的東西.
都是另一個神.
你可能問.
即使我將耶和華放在第一位.
我將其他神.
例如巴力或太陽神放在第二位可以嗎?.
Plan A 不可以.
就用Plan B.
我們香港人很流行.
斷乎不可.
這裡只有耶和華.
這個才是我們立約的對象.
在這裡.
神進一步用第二條誡命.
去解釋我們和祂的關係.
祂說不可為自己雕刻偶像.
也不可做什麼形象.

$^{401}$就連天空,地下或水中的白物都不例外.
這些偶像全部都是創造物.
也包括你和我人類.
雖然我們擁有上帝的形象.
但我們不是上帝.
弟兄姊妹.
跟隨基督絕對不是靠自己完成的事情.
我不是叫大家不努力.
而是提醒大家.
我們不可以將自己的能力成為唯一的希望.
不要將自己變成自己的偶像.
唯有基督的恩典和憐憫.
才是我們真正可以仰望的.
再者.
我們不要再跪拜其他A4.
並且事奉他們.
更重要的是.
我們不可以將眼前的人取代上帝.
特別是我們這群教導人的人.
千萬不要將人帶到我們面前.
因為這是一件很危險的事情.
我們應該將人帶到上帝面前.
我們也不需要為上帝製造偶像.
因為上帝不需要我們這樣做.
就讓我們單單相信上帝.
他是一位永活的真神.
在過去.
上帝怎樣拯救以色列人出埃及.
並且帶領他們進入迦南地.
今天.
他仍然會釋放屬於他被生活壓制的人.
曾經我在一間基層教會實習.
那裡的弟兄姊妹.
很多都是透過教會的食物銀行得到幫助而信耶穌.
並且有些牧者很貼身地陪伴他們.
所以他們都感受到神的供應和看顧.
坦白說.
其實他們信主之後的生活.
都是很拮据,充滿挑戰.
但這群弟兄姊妹.

$^{441}$他們願意奉獻時間和勞動力.
例如.
有些婦女.
她們會很努力學做盆菜.
之後教會一群大專生.
都會邀請區內獨居的長者.
或是街坊一起回教會享受盆菜宴.
我記得當時我的督導人.
跟我說了一句話.
他說:我們教會沒有錢.
什麼都缺.
不過靠著耶穌.
有一無所缺.
其實我們想深一層.
自己的能力.
生活的保障.
工作,家人.
錢財生愛物.
都是上帝所賜.
故此這些東西.
不應該成為我們的偶像.
主宰和定義我們的人生.
接下來.
我想說第四條誡命.
守安息日.
香港人最難守安息日.
因為香港人真的很喜歡上班.
舉世聞名.
我上網找到Google說.
香港是全世界最過勞.
最過度勞累的城市.
差點看不到.
不知在在有多少人.
真的可以下班呢?.
經文說.
「六日要勞碌做一切的工.
但第七日無論何工都不可作」.
如果我們對比出埃及的二十章版本.
申命記其實是給了另一個解釋.
這裡說.

$^{481}$「因為上帝大能的手.
和伸出來的蚌臂.
將你從埃及遺老之地領出來.
所以你要記得守安息日」.
做奴隸有沒有得安息?.
有沒有休息?.
當然沒有.
不知大家有沒有聽過.
其實在香港.
都是一個人口販賣的中轉站.
什麼叫奴隸呢?.
就是通過欺騙,強迫又或者暴力的手段.
從而讓他人失去了自由的狀況.
不知大家有沒有聽過.
其實現在在我們香港社會.
有一些東南亞的婦女.
她被中介公司又或者他們的僱主沒收了護照.
一個星期做足七天.
無法休息.
而有一些人口販賣機構.
就專門去救這些婦女.
上帝就是要提醒這一班的以色列人.
以前你作為奴隸的身份.
你七天都沒得選擇.
要被奴役.
但現在你和上帝有了這份關係.
你的身份就不同了.
你要和你以前在埃及的日子分別出來.
在你可以控制的範圍.
即是你自己,你的子女.
奴僕,牛奴,畜牲以及客女.
都應該與你一同享受安息.
因為你要記得.
被人奴役的日子真的很慘.
所以己所不欲,勿施於人.
你現在已經不是奴隸了.
因為你已經離開了埃及.
上帝照顧著你.
你的生活再不是一個「做」字.
Relax,上帝在大廳.

$^{521}$第五條是和做子女的我們說.
當孝敬父母.
父母聽到是不是很棒呢?.
最好每個星期都是青少主日.
昨天有位媽媽跟我說.
「珊珊,可否幫我教我的年輕人?」.
我聽到也不知道該說什麼.
因為我也覺得很汗顏.
不如大家一起努力吧.
我們一起回到《申命記》的五章第九節.
讀給大家聽這段經文.
「恨我的,我必追討他的罪,至父及子,直到三,四代.
愛我,守我誡命的,我必向他們發慈愛,直到千代」.
回到第二條誡命.
其實我們看到父母有一個很沉重的責任.
如果一個人對神的關係有錯誤的解讀.
他們的錯誤會禍及子孫.
做父母,做家長的你.
是有責任去愛神,守神的誡命.
但如果你本身也不敬畏神.
你就很難做到一個好的榜樣.
讓你的子女見證你信仰的真實.
正如剛才我跟大家說的一句話.
是我以前的牧者說的.
「我們就是未信者的第一本聖經.
如果我跟你不相信聖經是期間限定的話.
今天我們又如何用生命去翻譯上帝的話語給下一代聽呢?.
同樣地,做子女的也有一份責任.
這個責任就是孝敬父母」.
剛才我讀了《生命的事》四章四十字給大家聽.
「摩西今日將神的律例誡命曉完你.
你要遵守,使你和子孫可以得福.
並且使你的日子在耶和華你神所賜的地上得而長久」.
如何才能達至這個效果呢?.
就是當父母負責任教育子女有關上帝的盟約.
而子女另一邊也有一顆願意的心去聽父母的教導.
於是留下來的就不是一份心意.
而是和神的盟約.
這個誡命告訴我們一個盟約群體需要有兩層的關係.
頭四條的誡命就是描述.

$^{561}$如果我們和神的關係不正確.
就無法和人建立一個正確的關係.
而最後五條的誡命就是專門處理人與人之間的關係.
即是當我們對神的關係正確的時候.
我們也要依賴人與人有一個正確的關係.
第五條的誡命.
當孝敬父母正正處於這兩個極端的中間.
特別是針對家庭的情況.
這是和最親密的人際關係的領域.
亦是一個群體核心.
當以色列人聽完摩西講十誡之後.
其實他們正面對一個處境.
就是他們要用一個什麼樣的態度去實踐這個盟約.
讓我們看看今日經文的第三部分.
我邀請大家讀出《申命記》的六章四至九節.
即是以色列那裡開始.
準備.
一,二,三,起.
以色列啊,你要聽.
耶和華,我們神是獨一的主.
你要盡心盡性盡力愛耶和華,你的神.
我今日所吩咐你的話都要記在心上.
也要殷勤教訓你的兒女.
無論你坐在家裡,行在路上,躺下,起來都要談論.
也要繫在手上為記號,戴在額上為經文.
又要寫在你旁屋的門框上,並在聖門上.
多謝大家.
幾千年以來.
每日早晚,猶太人都會讀出一個很出名的土文.
叫做Shema,你們要聽.
用來表達他們對上帝的虔誠.
聽不只是指於耳朵.
更加代表著注意和專心.
譬如在《創世記》的二十九章三十三節裡.
尼亞將其第二個兒子命名為西面.
西面的希伯來文就是Simon.
好像和Shema,Simon.
意思是耶和華因為聽見尼亞的失寵.
聽見,所以他又賜了這個兒子給她.
Shema其實包含了聽,留意和對所聞之事作出一個回應.

$^{601}$當我們看詩篇時.
我們經常見到詩人向上帝求助.
他說:Shema,上帝Shema.
我用聲音呼籲時求你去誰聽.
並求你憐憫我,英雲我.
同樣地,當我們向上帝說Shema時.
上帝也要人Shema他.
在這六章的四節中.
其實是Shema,Shema,說兩次.
意思是叫他們聽清楚並且要服從這個誡命.
上帝說:你們要聽.
因為我是你們唯一的神.
因此以色列人要遵守他的誡命.
這就是愛上帝的內容.
不過怎樣愛?.
經文告訴我們.
你要盡心盡性盡力去維繫這份關係.
不如讓我們逐一看看.
這三個盡是什麼意思.
盡心不是感受.
而是代表我們的心思意念.
即是你的腦在想什麼.
你的腳就會跟著去做.
什麼意思呢?.
即是你心思思想念你的男女朋友.
你就會打電話去找他吃飯.
意思是盡心的時候.
用盡心思去愛上帝,找上帝.
盡性,英文的翻譯是soul,即是靈魂.
但其實看希伯來原文的意思.
即是生命或性命.
盡性的意思是用你整條命去愛上帝.
第三,盡力.
大家有什麼力量可以盡呢?.
所以我可以想想.
我們香港人可以有財力,心力,體力.
這裡也包括我們的資源.
例如我們的知識,人脈.
用盡我們所有力量去愛上帝.
當然,說愛主又談何容易呢?.

$^{641}$我記得有一次.
在我事奉到很灰心.
對前路很疑惑的時候.
我坐在其中一間實習教會的教堂.
我望著眼前的十字架.
這十字架跟我們往往的十字架有點不同.
這十字架上面有耶穌.
耶穌釘在木製的十字架雕刻上.
我去到台前.
我默想耶穌的表情.
祂戴著一個很磁場的袖容.
很磁,但又很磁場.
雙手微微張開.
其實我發現原來這裡有一個祝福世人的意思.
我想起耶穌的受苦是帶著祝福的.
祂從來都不離開祂的召命.
祂一直都為殘破不堪的世界.
或每況愈下的處境帶來祝福.
這也提醒我蒙召的初心.
讓教會跟隨耶穌成為社會不同人士.
包括邊緣群體傷痕累累的人的避難所.
讓我們每一個人都可以在教會與主相遇.
一眾旺角浸信會的弟兄姊妹.
你要聽,耶和華是我們獨一的主.
我們要盡心盡性盡力愛耶和華.
這一種盡心超越我們自身的感受.
跟隨神難免會有難受的時刻.
但感受不代表一切.
更不代表神不愛你.
抓緊我們是上帝女兒的身份.
第二,這種盡性就是要擺上被上帝塗灶.
有時候有可能要被破碎.
但相信上帝必定會重建.
最後,這種盡力可能代表我們要把自己覺得對的東西或利益放下.
有時候會不忿,但公道在上帝那裡.
但願我們一起盡心盡性盡力愛上帝.
因著上帝先愛我們.
無論何境況都要站主撲.
願主常見,理我眾生.
中心.

\newpage



\section{申命記 7:1-11}
\label{sec:C0mf0rNYQN4}
\textbf{2025年3月2日主餐崇拜直播｜陳明泉牧師｜朝向應許之地－專一｜申命記7\_1-11}
\newline
\newline
連結: \href{https://youtube.com/watch?v=C0mf0rNYQN4}{\texttt{https://youtube.com/watch?v=C0mf0rNYQN4}} ~~~~ 語音日期: 2025-03-03
\newline
\newline
\hyperref[sec:bBtvGdEATA0]{\small{< < < PREV SERMON < < <}}
~
\hyperref[sec:index_chronic]{\small{[返順時目]}}
~
\hyperref[sec:mfCFFhs9u_c]{\small{> > > NEXT SERMON > > >}}
\newline
\newline
申命記 7:1-11
\newline
\begin{longtable}{cl}
\hline
\hline
章節 & 經文 (和合本修訂版)\\
\hline
7:1 & \begin{tabularx}{0.7\textwidth}{X} 「耶和華－你的神領你進入你要得為業之地，從你面前趕出許多國家，就是比你更強大的七個國家：赫人、革迦撒人、亞摩利人、迦南人、比利洗人、希未人、耶布斯人。 \end{tabularx} \\ \\ \relax
7:2 & \begin{tabularx}{0.7\textwidth}{X} 當耶和華－你的神把他們交給你，你擊殺他們的時候，你要完全消滅他們，不可與他們立約，也不可憐惜他們。 \end{tabularx} \\ \\ \relax
7:3 & \begin{tabularx}{0.7\textwidth}{X} 不可與他們結親；不可將你的女兒嫁給他的兒子，也不可叫你的兒子娶他的女兒。 \end{tabularx} \\ \\ \relax
7:4 & \begin{tabularx}{0.7\textwidth}{X} 因為他必使你的兒女離棄我，去事奉別神，以致耶和華的怒氣向你們發作，迅速將你除滅。 \end{tabularx} \\ \\ \relax
7:5 & \begin{tabularx}{0.7\textwidth}{X} 你們卻要這樣處置他們：拆毀他們的祭壇，打碎他們的柱像，砍斷他們的亞舍拉，用火焚燒他們雕刻的偶像。 \end{tabularx} \\ \\ \relax
7:6 & \begin{tabularx}{0.7\textwidth}{X} 「因為你是屬於耶和華－你神神聖的子民；耶和華－你的神從地面上的萬民中揀選你，作自己寶貴的子民。 \end{tabularx} \\ \\ \relax
7:7 & \begin{tabularx}{0.7\textwidth}{X} 耶和華專愛你們，揀選你們，並非因你們人數比任何民族多，其實你們的人數在各民族中是最少的。 \end{tabularx} \\ \\ \relax
7:8 & \begin{tabularx}{0.7\textwidth}{X} 因為耶和華愛你們，又因要遵守他向你們列祖所起的誓，耶和華就用大能的手領你們出來，救贖你脫離為奴之家，脫離埃及王法老的手。 \end{tabularx} \\ \\ \relax
7:9 & \begin{tabularx}{0.7\textwidth}{X} 所以，你知道耶和華－你的神，他是神，是信實的神。他向愛他、守他誡命的人守約施慈愛，直到千代； \end{tabularx} \\ \\ \relax
7:10 & \begin{tabularx}{0.7\textwidth}{X} 向恨他的人，他必當面報應，消滅他們。凡恨他的，他必當面報應，絕不遲延。 \end{tabularx} \\ \\ \relax
7:11 & \begin{tabularx}{0.7\textwidth}{X} 所以，你要謹守我今日所吩咐你的誡命、律例、典章，遵行它們。」 \end{tabularx} \\ \\
[1ex]
\hline
\hline
\end{longtable}
$^{1}$願意繼續我們以聆聽聖道.
繼續我們的敬拜.
接下來我們會讀出.
今日我們的經訓.
在舊約聖經新命期.
舊約聖經新命期第七章一至十一節.
在大家手頭上教會的聖經是第二百二十三頁.
請聽我讀出新命期第七章.
耶和華神領進入耀德日頁之地.
從你面前趕出許多國民.
就是黑人,格加薩人,阿摩利人,迦南人,.
比利西人,希美人,耶布斯人.
共七國的民.
都比你強大.
耶和華神將他們交給你擊殺.
那時你要把他們滅絕,淨盡.
不可與他們納藥.
也不可憐憫他們.
也不可與他們結親.
不可將你的兒女嫁他們的兒子.
也不可叫你的兒子娶他們的兒女.
因為祂必使你的兒子轉你不跟從主.
去事奉別臣.
以至耶和華的怒氣向你們發作.
就速速地將你們滅絕.
你們卻要這樣待他們.
拆毀他們的祭壇.
打碎他們的柱像.
撼下他們的木牛.
用火焚燒他們雕刻的牛像.
因為你歸耶和華.
你神為聖潔的民.
耶和華你神從地上的萬民中揀選你.
特作自己的子民.
耶和華尊愛你們.
揀選你們.
並非因你們的人數多於別民.
原來你們的人數在萬民中是最少的.
只因耶和華愛你們.
又因要守祂向你們列祖所起的誓.

$^{41}$就用大能的手領你們出來.
從遺牢之間救贖你們.
脫離埃及王法老的手.
所以你要知道.
耶和華你的神.
祂是神是信實的神.
向愛祂守祂誡命的人.
守約思慈愛直到千代.
向恨祂的人當面報應他們.
將他們滅絕.
凡恨祂的人必報應他們.
決不辭言.
所以你要謹守遵行.
我今日所吩咐你的誡命律例典章.
請大家多多支持.
多謝.
請大家多多支持.
多謝.
頂子妹平安.
平安.
繼續來到去研讀新命記.
講到一個民族.
本來是在埃及裡面圍爐.
沒有自己的產業.
自己成為人家的產業.
但是上帝卻是有恩典臨到.
在這個以色列民族當中.
經歷過四十年流落曠野.
如今他們將要進入一個新的新天新地.
一個迦南地上帝應許給他們的地方.
在進入這個地方之前.
摩西就對這個年輕的世代.
作出很重要的說話.
今日我們讀到第七章.
你讀上去的時候你會發覺.
裡面都幾辣的.
因為裡面對以色列的要求.
或者在未來的日子裡面.
他們要怎樣去選擇.
怎樣去行.

$^{81}$是一個很靠近.
但是也要很多的決心和恆心.
來實踐一番的豢面.
在我們看經文之前.
我曾經看過這樣的預道故事.
這個是來自一個希臘的預道故事.
在那個故事裡面說.
有一個國王叫做Midas.
這個國王因為善待了他的導師.
叫做西勒洛斯.
因為善待這個人.
天上的神明見到他.
哇 這麼好.
於是就給他一個獎勵.
給他什麼獎勵呢.
他賜給剛才那個王叫做Midas.
有一種能力.
就是點石成金的金手指的能力.
他的手指只要指向什麼物件.
或者接觸某一個物件的時候.
他就會令到那樣物件.
由一個普通的形態變為金.
所以他的手指就叫做金手指.
不是他手指用金做.
而是他點石成金.
石頭可以變金.
樹可以變金.
哇 他發達了 是不是.
所有東西.
總之他手指碰過的.
全部都會變成金.
他就成為了這個世界裡面最富有的人.
不過他擁有世界上所有的黃金.
但是他卻是很多東西失去.
其中一個片段就是.
他見到一樣東西.
他最喜歡吃雞腿.
他想吃的時候一拿.
就變成了金雞腿.
能不能吃呢.

$^{121}$吃不到.
不過這個故事最後最令到人黯然的地方就是.
他的女兒要出嫁了.
當他女兒要出嫁的時候.
他擁抱自己的女兒.
女兒就成為了一個金人.
在他生命裡面.
他的起居飲食.
他最重要的親情.
因為黃金.
因為他有這種特別能力.
所有東西都變成了冷冰冰的黃金.
失去了世界上所有對人生命裡面最重要的東西.
這個預渡故事是在說.
我們人生裡面.
人生有很闊的.
擁有不同的東西.
但是我們很多時候.
為了一些次要.
或者為了別人覺得好的東西.
我們放棄了生命裡面.
一些最重要的部分.
今天我們用這個角度來看.
摩西去勉勵一班的以色列.
告訴這個新的世代.
不要因為迦南地裡面.
你覺得那樣東西的好.
或者其他邦國裡面覺得好的地方.
你就盲忠忠.
跟著這些人來追求.
又或者那些人.
所經歷的一切.
未必是整個以色列民族要追求的一切.
我們希望今天用這個角度.
來重新思考.
《申命記》第七章裡面所說的內容.
我們先祈禱求聖靈幫助.
上主願你幫助我們.
讓我們在你的話語當中.
讀到生命的亮光.

$^{161}$賜我們智慧.
讓我們能夠明辨.
你對我們所說的每句說話.
賜我們信心.
讓我們帶著信心來聆聽.
又讓我們有恆心.
在實踐當中.
我們銘記你對我們的教導.
幫助每一位弟兄姊妹聽得明白.
幫助孫講的講得清楚.
我們恭敬將來的時間交託.
願你賜福.
獻上禱告奉靠基督耶穌得勝的聖名.
Amen.
我們一起看看.
在《申命記》裡面第七章.
我們先看第六節至到第八節.
這個是整個段落裡面.
一個很重要的神學內涵.
因為你歸耶和華.
你神為聖潔的民.
耶和華你神從地上的萬民中揀選你.
特作自己的子民.
耶和華尊愛你們揀選你們.
並非因你們的人數多於別民.
原來你們的人數在萬民中是最少的.
只因耶和華愛你們.
只因要守他向你們列祖所起的誓.
就用大能的手領你們出來.
從遺牢之家救贖你們.
脫離埃及法老王的手.
在這裡面你會看到.
在這裡面他特別強調一個很重要的觀念.
就是揀選.
上帝揀以色列民.
我們在聖經裡面經常會讀到揀選這個字.
揀選其實具體的意思是甚麼呢.
就是在說上帝去選擇我們.
上帝揀我們.
主權就放在上帝那裡.

$^{201}$我們今天很多弟兄姊妹.
有時我們對基督教的信仰.
我們有種這樣的體會.
我們覺得是我們揀上帝.
我們揀神來相信.
我們比較世界上不同的宗教.
然後作出一個理性的決定.
這種的想法是說.
十七十八世紀啟蒙時期的人.
對自由,平等,選擇的鼓吹.
以致整個世界都是這樣想的.
強調選擇的力量.
不過在聖經裡面說到.
人能夠得著這個救恩.
人能夠在救恩上帝的任備當中.
主動權或決定性的因素不是在人的選擇那裡.
而是在上帝的選擇.
所以揀選這個意思.
其實是說到摩西對著那些以色列人說.
是上帝揀你們的.
但是以色列人有什麼優越之處.
好像選美那樣.
是不是有些個人之處讓上帝來揀呢.
摩西說不是.
萬民當中.
人數最少的.
人數最少的意思不是純粹是人數.
而是整個民族.
沒有什麼可以值得上帝覺得你正來揀選.
你的人數少代表你的實力其實很弱.
你的人數少代表你沒有什麼很深厚的文化觀念.
你的人數少代表其實這個民族是不堪一擊.
是一個處於萬族萬邦裡面.
很容易被淹沒在歷史當中的一個民族.
所以在這裡上帝揀以色列.
不是以色列有什麼個人之處.
不是說這個民族有很多優越的地方.
單單只是上帝的主權.
上帝說要揀你.
以致這個民族承接著上帝給他不同的福氣.

$^{241}$經文裡面再繼續說.
剛才我說過了.
你們並非你們的人數多於別人.
原來你們的人數在萬民中是最少的.
並非和原來這個巨色提醒他們.
他們不是很厲害.
不是想像中那麼有實力那麼厲害.
他們什麼都不是.
只不過是一切出於上帝的恩典.
我們可不可以捉到路呢.
可不可以知道上帝憑什麼來揀人呢.
在聖經裡面說沒有的.
不知道的.
去到新約裡面.
保羅也說過.
他也不明白上帝為什麼要揀選世上愚拙的.
是那些不體面的.
在世界上被遺棄的人.
上帝偏偏就揀這些人成為他的子民.
究竟上帝裡面.
他那把尺是怎麼量度的呢.
沒有人知道.
我們今天來到教會當中.
我們今天成為基督徒.
不是因為我們自己很厲害.
我讀很多書.
我能夠擁有不同的智慧能力.
單單只是上帝的主權.
上帝將你放在這個群體裡面.
上帝將你放在信仰.
在基督裡面.
以致我們就可以得著這個恩典.
經文裡面再繼續說.
他們什麼都不是.
他們是怎樣稱呼這些以色列民呢.
他用幾個字眼.
他說第一你是聖潔的民.
第二他是特選上帝的子民.
上帝尊愛你們.
你會見到上帝在雲雲的民族裡面.

$^{281}$揀選其實可以有很多選擇.
他揀了以色列人.
但他揀了之後不是當他們是地底泥.
而是放在一個很重要.
放在一個很高的層次裡面.
上帝這樣來看他.
首先看聖潔的民.
我們一說聖潔.
我們第一樣想的.
我們就想道德觀念.
這個不是錯的.
不過聖潔在舊約聖經.
特別在聖經申命記裡面.
說的不是純粹道德觀念這麼簡單.
聖潔其實是說一個和神之間的關係.
所以你見到.
第六節說到.
因為你歸耶和華你神為聖潔的民.
其實那個歸耶和華為聖.
其實是一個整全的字眼.
是說到上帝在雲雲這麼多的民族裡面.
拉了你出來.
成為他所和他建立關係的一個民族.
所以在這裡面.
那個歸耶和華為聖.
是說一種特別緊密的關係.
也是一種上帝藉著這個民族來敬拜他.
將主權在地上萬族萬邦裡面.
他就好像成為一個.
凸顯了萬族萬邦裡面.
一個特別不同的身份位置.
去到新約.
基督徒就是萬民當中揀選我們.
我們就成為了上帝特別愛眷的一班子民.
所以聖潔第一樣東西是說的一種關係.
第二樣東西.
特別說到的那個字眼就是.
特作自己的子民.
這裡應該是在第六節尾那裡.
這裡和合本譯不是譯錯.

$^{321}$不過譯下去就不夠到位.
因為原文用那個特作自己的子民.
他不是純粹特別這樣的意思.
和合修訂本和修本去用那個.
就是作自己寶貴的子民.
有些英文譯本是說到.
上帝的私人珍藏.
所以在這裡摩西在說.
這班以色列民看在上帝的眼中.
他們是上帝的珍藏.
上帝當他們是寶的.
如果你用這個角度去看.
那個聖潔的意思.
就不是純粹說.
上帝純粹建立一些關係.
上帝要將這班人視為.
一班最重要.
上帝心上最記掛的一班子民.
你和我今天我們去到新約.
在基督耶穌裡.
我們都是在基督耶穌當中.
被視為寶貴.
不是因為我們有多麼厲害.
有多麼好.
單單是上帝的主權.
上帝不單止救你.
而且將你放在一個很重要的位置.
放在一個上帝心裡面.
做每一個舉動.
他都是要將福氣臨到你和我身上.
臨到那些上帝的子民當中.
所以在這裡面.
其實他的用語是貫穿了.
整個簡選的意思.
經文裡面再繼續說.
然後到第三部分.
就是在說聖潔.
就是後面.
既然上帝是這樣看你.
是這樣看那班子民.

$^{361}$這班人不是任意輕率地過生命.
這班人就要和這個世界分別出來.
所以我們經常讀.
分別為聖.
分別出來的原因.
就是基於上帝的簡選.
基於上帝對那班子民的厚愛.
以至這班人和其他世俗很不同.
他們要過一個和其他世俗不同的生命.
就是單單用一個很獨特的關係.
來和上帝建立一個很親密的交往.
如果這樣說也很抽象.
我們用婚姻來做比喻.
你就知道.
一男一女在上帝面前.
共證婚盟.
他身邊站在那裡的男士.
就和所有的男人不同.
這個男士和他立約的丈夫.
就是這個妻子裡面.
用獨一的愛.
用完全的效忠.
完全的愛.
學習這樣來愛他.
同樣道理.
這個丈夫就要用獨一的愛.
世界上其他女孩子.
再漂亮也好.
比自己的老婆漂亮也好.
他都不會入眼.
不會有任何風吹草動.
他只是帶著單單愛自己妻子的態度.
來經歷他未來的人生.
這種婚盟的態度.
就是很多時候聖經裡面的比喻.
以色列民和上帝之間的關係.
一生一世.
一種完全效忠.
一種分別.
從這個世界出來.

$^{401}$你對丈夫和你對其他男人是不同的.
就是一種這樣的關係.
其實在這裡.
剛才一直說這個揀選.
其實你會一直推演下去.
你就知道.
完全基於的是一種什麼的精神呢?.
就是上帝與以色列民立約.
所以摩西在這裡說.
只因耶和華愛你們.
又因要守他向你們列祖所起的誓.
就用大能的手領你們出來.
從遺牢之家救贖你們脫離埃及王法路的手.
在這裡說的那個誓.
其實就是上帝與亞伯拉罕.
跟著亞伯拉罕以撒雅各所納的這個約.
上帝和以色列民就是一個約定.
有一個covenant.
是一個一生一世完全效忠的一種關係.
弟兄姊妹.
今日我們信主的人.
我們生命其實都是和上帝有這份立約的關係.
可能你會覺得你人生裡面會有跌跌撞撞.
會有一些跌倒的時間.
又或者有一些行為上不理想.
但是那個行為或者你生命裡面不理想的地方.
上帝因為基於一個約裡面.
他會愛你到底.
會拯救你到底.
我們就帶著這一份的信心.
來調教我們自己的生命.
我們人生裡面不會納了約之後.
就成為一個完美的人.
我們納了約之後.
我們人生都不斷的轉化.
我們不斷的進步.
以致到我們生命裡面所承受的恩典.
和上帝和我們納的約.
慢慢互相相稱.
新約保羅說.

$^{441}$願你蒙照的恩相稱.
就是在說這個意思.
不是你很完美.
然後就和上帝結合關係.
上帝說你是他的聖潔子民.
然後你人生就不斷不斷的進步.
逐步逐步將自己變得神聖.
是一個歷程.
是一個更新的過程.
以致到我們和上帝.
給我們這個恩典越走越近.
用這個角度你就知道後面的那兩個.
其實所說的那兩段是一些應用.
屬靈意義放在我們今天裡面的.
我們首先就再看第一至第五節.
耶和華你神領你進入要得為業之地.
從你面前講出許多國民.
就是黑人,格拉薩人,阿摩利人,迦南人,比利西人,.
希美人,耶布斯共七國的民.
都比你強大.
耶和華你神將他們的家族擊殺.
那時你要把他們滅絕正盡.
不可與他們立約.
也不可連戌他們.
不可與他們結親.
不可將你的女兒嫁給他們的兒子.
也不可叫你的兒子娶他們的女兒.
因為他必使你兒子主尼不跟從主去事奉別臣.
以致耶和華的怒氣向你們發作.
就速速地將你們滅絕.
你們卻要這樣待他們.
拆毀他們的祭壇.
打碎他們的柱杖.
砍下他們的木偶.
用火焚燒他們雕刻的偶像.
在這裡我們第一件事.
我們看到最多出現的就是要滅絕.
然後後面就是要砍掉那些偶像.
拆毀他們的祭壇.
以及火燒他們那些雕刻的偶像.

$^{481}$我們覺得純粹是一個偶像做專業出來.
不過你再用剛才那個觀念去看.
在這裡那些外在行為目的就是要除去.
整個以色列民族將要進入迦南地的時候.
不要受著迦南那七個族群的文化去沾染.
或者令到以色列民族迷失.
以致他們就離棄上帝.
記不記得剛才所說的.
在迦南地裡面.
人數最少的就是這班以色列民族.
他比起進入迦南地裡面其他的族群.
人數都要少.
一個小小的民族.
你覺得你真的可以去用你自己的文化.
去同化人家或者去感染人家.
其實不是的.
在這裡摩西很清楚告訴那些以色列民族.
不可以的.
不要給自己有後路.
要去到一個盡的階段就是滅絕正盡.
在文化角度裡面.
不要受這些文化影響你自己和你的兒女.
否則的話你們就會沾染世俗.
你們就會逐漸逐漸離棄上帝.
去到今天我們知道那個歷史處境.
今天我們不會用這些字眼來過我們的生活.
又或者我們不會照字面.
那些不是基督徒我們要殺了他.
不是這樣想的.
所以我們要用屬靈的角度去想.
這個就是上帝要叫我們在生活裡面.
有一些令我們不能夠專注在上帝.
不能夠好好跟從上帝的這些絆腳石.
我們要用減法的方式去剔走他.
以致我們真的減低了生命裡面.
令到我們分心或者令到我們沒辦法專注的.
有一些的事和物.
經文裡面再繼續說.
這班的以色列人不可以跟他們去立約.
不可以跟他們有婚姻關係.

$^{521}$又做很多永恆的立約關係.
不是說他反對國際貿易那樣東西.
那個立約的意思就是一種完全效忠的意思.
以色列民只有一個效忠的對象就是上帝.
在他們進入去迦南地的時候.
他不可以將這個立約的對象溝淡了他.
他是要真的專注在上帝當中.
所以這一種的捨去法.
這一種的減法.
目的就是要令到那些以色列民專注在上帝那裡.
有一些的靈修學家.
我特別喜歡看的古倫神父.
他講幾樣東西.
今天對於我們都市人來講.
可以應用這一種的方法.
這一種的減法來幫我們.
第一種就叫做對話.
跟你裡面有些東西會令到你很想追求.
你就跟你心裡面的聲音.
為什麼你那麼喜歡那樣東西.
為什麼那樣東西對你那麼重要.
講得白一點.
有時候事奉.
有時候我們很介意別人怎麼看自己.
事奉是為了上帝的.
但為什麼別人的聲音.
別人所講的話對我那麼著意.
這些就可以進入對話.
很老實講.
每個人都會經歷的.
但我們就不管的.
我們繼續我們自己的生活.
靈修裡面提醒我們.
有時候我們都是做同樣的宗教行為.
但我們帶著的可能是.
帶著異心的來做.
我們需要正視.
我們需要對話.
所以第一樣東西就是.
將你生命裡面你呈現了的.

$^{561}$或者你耿耿於懷的東西拿出來.
你自己就拿來祈禱.
跟自己來對話.
第二樣東西就是進入一種默想.
默想你的情緒.
你的情感.
為什麼某些東西會令你發瘋失控.
在你生命當中.
為什麼有些東西會容易開啟你的制.
這些情感是來自什麼原因.
你在上帝當中.
你著重的東西.
還有那些人對你.
或者你的生活裡面.
是不是挑戰著你自己的存在價值.
這些情緒.
是可以成為了你對話.
或者令你自己可以更加明白自己.
以至你可以減低了你自己分心的狀況.
反而是你清楚自己的時候.
你就可以一路一路.
在明白自己的過程當中.
更加追求上帝.
第三個就是那些叫做負面駁斥法.
這個是耶穌用的.
魔鬼在曠野裡面試探耶穌.
耶穌不是斥責魔鬼的.
撒旦退後邊去吧.
你不要以為耶穌才用.
其實生命裡面我們有很多.
令到我們會分心.
我們當然不會在街上罵的.
但是你心裡面.
你就趕走這些念頭.
這些東西不是屬於上帝的.
用這種抗衡的方式.
來避免了有一些.
令到自己一路一路分心.
或者勾引你脫離上帝旨意的.
這些感受.

$^{601}$這些情緒.
或者這些想法.
古倫神父會用這個.
另外還有一個就是.
將那些令到你容易.
容易引起你分心的東西.
拔走它.
不要放在你的眼前.
我們廣東話說的叫.
斬手指避沙蟲.
這個是有效的.
其實有一些試探我們是要逃避的.
你不要想著透過人的能力可以做.
弟兄來說.
見到一些衣著暴露的.
你真的要掩眼的.
你真的要離開的.
你不要想著看一會兒.
無所謂.
第一秒你不是犯罪.
第二秒你都差不多了.
看多第三秒你就來了.
將自己抽離.
以致到避免自己進入那種狀態.
又或者今天.
有一些操練.
未必去到那麼的.
但是有時候我們生活.
會被手機佔據我的所有空間.
將手機拿遠一點.
我去過一次.
早一輪去過一個基督教的晚宴.
我覺得很冒犯的.
進去的時候.
為什麼呢.
進去的時候.
那個總幹事.
他說今晚我們大家一起來吃一餐飯.
一起分享異象.
所以請大家幫我做一件事.

$^{641}$將你的手機關掉.
就放在桌子上.
那個轉盤.
吃飯的那個轉盤.
然後就轉一轉.
所以你眼前那個.
就不是你自己的手機.
你不會聽別人的手機.
就放上去.
然後說.
搞錯啊.
我這麼大個兒子.
我去很多那些基督教的宴會.
都沒有試過這樣的做法.
我覺得很冒犯.
很冒犯.
嘰嘰咕咕.
然後那個總幹事說.
就因為沒有手機.
所以你和身邊的人結緣.
聊聊天.
他說這件事我開始明白.
其實很多時候.
我們出去見人.
其實我們有部手機.
掩飾了我們的尷尬.
我們不知道說什麼好.
沒有什麼話題.
就拿部手機出來.
本來沒有WhatsApp.
硬要查有沒有WhatsApp.
本來Facebook都沒有update.
都還要看一看.
目的就是不想看人.
又不想聊天.
就想營運自己的小小空間.
過了一個多小時就算了.
不過那一晚我開始覺得很冒犯.
但正正就是因為做了這樣的動作.
我發覺我認識了很多.

$^{681}$基督教裡面很多.
很志同道合的一些童工.
有一個過程裡面.
很深入的來去交流.
而且當晚裡面所說的.
那個意象分享.
我收得很清晰.
我發覺我離開.
到轉盤拿回在我面前.
吃完飯就走了.
我發覺挺好的.
如果王浸每次開座都是這樣.
挺好的.
求主幫助我們.
有時候很多事情.
未必是很大的罪.
但會令我們分心.
會令我們沒辦法專注在上帝那裡.
我們有時候會用這些方法.
來掩飾自己的尷尬.
又或者令自己感覺上.
好像容易一點去kill time.
但實際上我們是浪費時間.
多過是專注在上帝那裡.
所以在這裡面.
對上帝專注要用減法.
第二個由第九節至第十一節.
這個是在說家法.
所以你要知道耶和華.
你的神他是神.
是信實的神.
向愛他守他戒命的人守約.
思前愛直到千代.
向恨他的人當面報應他們.
將他們滅絕.
凡恨他的人必報應他們.
決不辭而.
所以你要謹守進行我今日所吩咐你的戒命律例典章.
在這個最後一個段落裡面.
你會發覺很強調律例典章戒命.

$^{721}$我們覺得是這些規條.
不過你留意一下.
他這裡面說了一個正反的結果.
如果對於那些守戒命.
守約的人.
上帝會將慈愛恩典.
一直放在那些以色列民身上.
所以這個顏面不絕的結果.
是去到千代.
一千代都能夠領受這個守戒命.
愛他.
實踐他照命的人.
你就有一個這麼長遠的效果.
相反來說.
你不理他.
你恨他.
你帶來的就是即時的死亡.
即時的滅絕.
一個時間裡面一個長與短.
一個禍與福的因果關係.
摩西就將這件事放在那些以色列民面前.
來到他做.
你要做選擇.
你要講顏面千代的祝福.
一或是講即時被上帝滅絕.
恨惡你.
令你得不到福氣的一種生命.
你自己選擇.
讓你選擇.
跟著他講.
所以你今天要謹守進行我吩咐你的戒命律例典章.
這個吩咐其實是一個禍與福的選擇.
之後有機會會再仔細講這個禍與福的神學觀念.
不過在這裡面.
摩西已經將權衡輕重.
告訴了那些以色列民知道.
不過更加重要的就是.
既然是這個時候.
那個遵守律例典章.
就是一個很重要的關鍵因素.

$^{761}$他可不可以得到這個永恆的福氣.
我們一講想起這些想法.
我們覺得遵守律例典章是一些很艱鉅.
很違逆我們的人性.
很痛苦.
甚至乎是很規條主義的那種想法.
不過在這裡面講.
其實不是一些很複雜的東西.
其實是一些很簡易的實踐.
在你生命當中.
你尊爺和華偉大.
你記起他的話語.
你值得留意.
摩西那時候講的.
沒有今天那麼厚的那本聖經.
他舊約.
啟示都未完全啟示.
摩西所講的那些以色列文.
就背和記.
每日記一些.
每日預備多一些.
就是銘記在心靈當中.
這些其實是成為生活的一部份.
當你將所有這些的誡命律例.
成為你生活一部份的時候.
你根本就不當是一回事.
就是用一個不是很貼切的例子.
我們今天從小到大都學.
我們要成為一個奉公守法的良好公民.
奉公守法良好的公民.
我們不會做每一件事都拿著那本基本法來看.
很多東西其實是生活的常識.
你所有東西其實是你跟著.
你其實進入了他的生命當中.
那樣東西你照著行.
就是了.
你說做每一樣東西都揭聖經.
你不會的.
因為那樣東西.
如果你平時讀經有功的時候.

$^{801}$那樣東西已經內化在你的心靈當中.
你自己會.
你裡面有個律來指教你.
指引你.
讓你走在一條對的路上.
那個重要性.
其實不是說這本聖經很難行.
而是我們願不願意將聖經放在生命當中.
早前要趕神學院一份功課.
我看了一份很有趣的文獻.
他講到一個人的轉化.
他做了一個質性的研究.
做一些研究人可以改變的.
這個題目很有趣.
他講到原來人的改變來自主要兩樣東西.
第一樣東西.
人的改變來自生命裡面有一些突然其來的突變.
那些可能是對你生命有威脅的.
那些是令你很痛苦.
或者是很突然一聲.
好像忽然間沒有了那種經驗.
這種突變會令人的價值觀忽然間扭轉.
不過他說在研究裡面講到.
人真的因為這些突變而改變.
或者真是人生裡面無端端出現這些突變.
不是那麼經常.
如果天天都是突變.
那個就不是突變了.
有些死人倒樓你不會天天都死人倒樓的.
這些出現的事件.
在你人生裡面大概佔兩成.
兩成的事.
這些突變會令你有些改變.
他說六至七成的改變來自.
每天以如木掩.
很低調.
很少量.
逐步逐步的那些改變.
每天是很少的.
可能零點零零零一.

$^{841}$但日子有功.
日復日年復年.
然後你再回頭一看.
嘩 這種那麼少的改變.
原來就帶來一個.
生命裡面一個很巨大的影響.
我一讀我就明白了.
今天講的那些律例典章.
其實很多都是一些很低調.
很細微.
不是一些很誇張的.
激變的那些道理.
生命裡面只是累積一點點.
一點點.
每天讀經.
讀一點點.
你讀Facebook的時間.
一定會比你讀經多.
你看其他文字.
一定會比你看每天一張聖經的文字多.
但我們有時候會選擇看那些.
放棄每天讀一點點.
每天讀一點點.
讓你生命改變一點點.
年復年日復日.
帶來生命徹底的改變.
弟兄姊妹.
在這裡講的那個加法.
不是加很多的重擔.
不是加一個很難扶的重握.
在你們的生命當中.
每天只是加一點點.
很少.
你絕對能力可以應付.
但經歷過日子有功.
你的生命你會發覺不再一樣.
今天我們看這段經文.
申命記第七章.
最重要的神學觀念講到.
上帝與我們納藥.

$^{881}$這個納藥不是廉價的藥.
上帝看我們寶貴.
不是純粹祂看.
然後我們繼續過一個世俗的生活.
每天學習.
讓上帝的神聖改變我們.
我們知道我們是一個.
與上帝納藥的子民.
我們就好好活出那種.
聖潔被上帝喜悅的那種生命.
活出你的豐盛.
活出那種基督徒生命應有的那種風采.
這種風采未必是要經歷一些很大的事件.
逐少逐少.
日積月累.
你會發覺你的生命不再一樣.
在生命當中有很多東西.
會難了你.
成為你半腳zoung .
拿走它.
不要放在你的眼前.
逃避這些試探.
讓我們緊緊跟隨基督耶穌的腳zoung .
這個是申命記第七章.
給我們專一的教導.
願主幫助我們.
\newpage



\section{ }
\label{sec:mfCFFhs9u_c}
\textbf{2025年3月6日預苦期首日塗灰禮晚禱崇拜直播}
\newline
\newline
連結: \href{https://youtube.com/watch?v=mfCFFhs9u-c}{\texttt{https://youtube.com/watch?v=mfCFFhs9u-c}} ~~~~ 語音日期: 2025-03-06
\newline
\newline
\hyperref[sec:C0mf0rNYQN4]{\small{< < < PREV SERMON < < <}}
~
\hyperref[sec:index_chronic]{\small{[返順時目]}}
~
\hyperref[sec:crDuuHRkfCw]{\small{> > > NEXT SERMON > > >}}
\newline
\newline
$^{1}$弟兄姊妹,今天我們用這段大家可能都很熟悉.
我們讀過很多次,甚至有時候我們都默想過很多次的一段經文.
我們重新思考一個題目,生命的真相.
這段經文雖然對我們來說是很熟悉.
但當你仔細去看的時候,你會發覺.
其實這段經文所出現的情景.
如果相對於耶穌行的神蹟.
或者在門徒所經歷的,其實這是一個很生活化的片段.
在這裡講到耶穌在格尼撒勒河邊那裡.
有很多人擠湧著祂,要聽耶穌的道.
為了不用被人群擠湧著.
耶穌就上了彼得的船那裡.
稍微離開岸邊,離遠的來將天國的信息.
將耶穌基督所要作的教導,就向群眾說.
說完之後,對彼得就簡單作了一個很直率.
甚至是一個好像常識一樣的,告訴西門兄.
將船開到水深那裡,在那裡盼望打魚.
然後西門彼得就說,我們整晚努力了很久.
都沒有打著任何的漁獲.
但是既然夫子你這樣說,我們就下網吧.
於是他就出去了水深的地方,下了網,圈著很多的魚.
在一個他們意想不到的情景之下.
竟然網裡面就再滿了漁獲.
多到一個地步,就是那艘船都差不多要沉下去了.
他還叫其他的船過來幫忙,整個這次的打魚.
足足再滿了兩艘船,整艘船誇張到一個地步.
差不多要沉到海裡一樣.
這是一個很生活化,是一個豐收的情景.
不過這個情景卻是為彼得帶來生命一個很重要的轉捩點.
因著這個情景,如果一般人來說.
見到一個這麼巨大的漁獲.
心裡面想著以後就要抓住耶穌.
每一次出海打魚一定要抓住他.
他就可以永遠都是帶著這麼豐收的漁獲去過他的生活.
不過這裡就說到彼得不是這樣做.
彼得第一樣看到的是上帝的神聖和自己的庇護.
所以在這裡他跟耶穌說主啊離開我,我是個罪人.
其他人看到這麼多的漁獲.
同樣都有一種很詫異的經歷.
於是雅各,約翰都是這樣對著耶穌.

$^{41}$帶著一份敬畏,帶著心裡面有一份覺得很不配的感受.
耶穌勉勵他不用怕,從今之後你要得人了.
整個故事其實你看到是一個很生活化,一個豐收的片段.
有什麼生命真相留在這裡呢?.
第一樣東西其實在這麼生活化的片段裡面.
其實他告訴我們可能我們的生命跟西門彼得很近似.
我們可能人生裡面作過很多的努力.
但是我們的人生充滿了很多的限制.
我們有時候會期望一分耕耘就一分的收穫.
我們作了很多的努力.
我們期望我們自己的努力可以換取到相應的效果.
不過當我們人越大,人生越來越豐富的時候.
我們會越來越發現有時候世界的法則不是我們心裡面的想像是一致的.
甚至乎可能很多時候會覺得徒勞無功.
生命裡面付出了很多的努力.
但是永遠跟心中所期望所得到的很多時候都不是成為正比.
又或者有時候我們會覺得這個世界很不公平.
有些人可能沒怎麼付出努力的.
但是反而他得到多一些.
自己負心自問,自己牢牢碌碌.
但是得到的,或者好像彼得那樣.
整夜努力卻打不著什麼.
人生第一個真相是.
原來我們的生命不是像我們預期一樣.
我們的生命充滿了限制.
我們期望用自己的努力可以換取更加多.
不過單靠自己人的努力.
卻是沒辦法達到我們心中所期望的效果.
當罪進入這個世界的時候.
上帝跟阿當說過.
他說當你犯了罪的時候.
土地已經不再為你效力.
男人要汗流浹背.
才地裡面勉強有出產.
人就好像逆流而上一樣.
大自然對我們來說不是想像中那麼友善.
有時候可能我們面對著很多的艱難.
我們逆流而上.
我們人生,我們所期望的.
我們的努力擺上.

$^{81}$但可能得不到我們預期的收穫.
這是生命真相的第一點.
第二點,在經文裡.
其實當耶穌基督叫彼得開往水深之處.
去打魚的時候.
其實如果你再仔細去看.
耶穌的本行是做木匠的.
一個木匠叫一個打魚的人.
去開船到水深之處.
其實是一個常識.
如果對於一個經常出海.
甚至以捕魚為生的漁夫.
你講一些常識給他聽.
其實是一份很大的冒犯.
不過在這裡.
彼得反而說.
夫子既然你這樣說.
我就樂望打魚.
你會覺得為何彼得這麼容易就範呢?.
為何彼得願意聽耶穌基督.
去作出一個近乎是常識的建議呢?.
除了耶穌帶著權能.
對著群眾講道之外.
其實在盧卡福音裡也埋下了伏線.
在第四章裡曾經記載過.
38至39節.
他講到彼得的外母曾經患了一個很嚴重的熱病.
然後耶穌去到他的外母當中.
很快地將他外母身上的病治好了.
於是他的外母就馬上來侍候耶穌.
彼得第一身見到耶穌醫病的大能.
他也覺得這個夫子講的話語.
他所講的天國道理.
和一般的文士,法利賽人不同.
很多人要擁擠著他.
要聽他所講的道理.
他就帶著一份某程度.
因為是夫子所講的.
既然你說到開口要做這件事.
我就按著你的心意去做.

$^{121}$可以想像的是.
其實彼得從來都沒有想過.
耶穌這個直率到不得了的.
這麼直接的命令.
卻是換取了他生命裡.
都沒有經歷過這麼豐收餘禍的經歷.
在這裡講生命的真相第二點.
其實人和彼得一樣.
有時候我們的視野很狹窄.
通常我們判斷事物.
我們去看這件事如何發生.
我們如何做決定.
往往都是根據自己的經驗去做決定.
有時候可能身邊的人.
講一些我們覺得都是常識的東西.
我們都知道.
但我們心裡總是覺得.
我們的直覺告訴我們.
或者我的經驗告訴我們.
這件事是不可行的.
看有什麼人跟我們說.
如果有些德高望重的人.
可能我們都會勉強去做.
但如果有些人.
我們覺得我們自己所擁有的知識比他更多.
可能我們就會嗤之以鼻.
有時候我們甚至乎.
我們會用一種鄙視的態度.
去面對那些給我們建議的人.
不過在這裡.
彼得的視野就好像我們的視野一樣.
我們的人生其實往往.
我們都是聚焦在我們自己所經歷的事.
我們將我們的經驗奉為.
我們生命裡最高的標準.
但我們可能忽略了.
我們人生裡有很多東西.
可能未必是我們經驗裡經歷過.
但卻是對我們生命很重要.
生命的第三個真相.

$^{161}$也是對彼得來說.
這是一個很簡單的捕魚經驗.
但對他來說是一個生命的轉捩點.
當看到這麼多的漁獲.
拿了上來的時候.
彼得他第一句對耶穌說.
「我是個罪人,請你離開我」.
你會覺得這是一個很奇怪的一句話.
事實上在經文的上下文裡.
沒有記載過彼得做過任何犯罪的行為.
在這裡也沒有留到.
彼得半點對耶穌的不敬.
所以看到這個漁獲.
看到一個這麼豐收的景象.
如果對於一般人來說.
是一個高興到不得了.
不得了的一個經歷.
不過在這裡.
正正是一個這麼大的恩典.
卻是讓彼得心裡.
感受到一份他前所未有的自覺.
他是一個罪人.
在我們看歷史.
或者在我們自己看電視劇也好.
要自認是一個罪人.
很多時候我們在電視劇也會看到.
叫人去承認這個罪.
大部分原因是因為人們怕刑罰.
甚至有時候會嚴刑逼供.
在強迫之下.
為了避免疲弱受苦.
所以人在這些威逼之下.
他就承認自己有罪.
又或者承認自己做錯了這些事.
可能他真的有做.
又或者可能他沒有做.
屈打成招.
以致他承認.
用口裡說我是一個罪人.
不過在這裡.

$^{201}$完全不是那個圖畫.
耶穌用恩典.
用他的豐盛.
來令人訴諸自己的良知.
人在上帝豐饒的恩典當中.
重新有一份自覺感.
原來人在世界上實在太渺小了.
相對於上帝一個偉大的恩典.
相對於上帝這麼偉大.
這麼豐足的供應.
再看回自己.
去檢視自己生命的時候.
發現原來自己的生命是這麼疲乏.
發現自己的生命原來是一個.
這麼不堪的一個人.
大概彼得在這裡.
他發現生命的真相是.
在上帝豐盛的恩典當中.
他看見自己的不堪.
看見自己最污穢的一面.
不過這一份污穢.
卻不是令他立刻要撲向耶穌.
求耶穌饒恕他.
這一份覺得自己很不配.
很污穢的心態.
他第一件事就是求耶穌.
求耶穌離開他.
他想和上帝製造一個安全的距離.
面對著一個豐盛恩典的主.
面對著一個沒有能力的耶穌.
人第一件事在這裡.
看見上帝想敬而遠之.
在這裡可能也是我們生命的寫照.
有時候我們不是覺得基督教不好.
而是因為我們覺得基督教實在太好.
我們再看回自己的生命.
我們覺得自己有時候不堪入目.
有時候我們會逃避上帝.
有時候我們想和上帝製造一些安全距離.
因為他的好.

$^{241}$實在令我們心裡有一種覺得.
完全不配的那份體會.
有很多弟兄姊妹曾經和我說.
他們年輕的時候信主.
但每一次唱詩歌的時候.
聽到詩歌很雄壯的音樂.
又或者見到講道裡面.
耶穌基督所講的內容.
一當聽到這些說話的時候.
再影照自己生命的真相.
他們第一件事很想避開.
因為他們覺得自己是污穢的人.
面對聖潔的主.
他們覺得完全不配來接觸這個信仰.
甚或有些弟兄姊妹可能生活在很多的艱難.
很多的沉溺行為.
甚至很多不堪在口中說出來.
一些很扭曲的關係.
面對著這些林林種種.
我們覺得自己生命是不配.
我們很想和上帝建立一段最安全的距離.
其實故事可以是這樣來發展的.
不過經文裡面講到.
耶穌的目的不是要拒絕人.
也不是要令到這個世界所有的人覺得自己很難堪.
聲稱自己是一個罪人.
並不是要令到自己找一個山洞.
來逃避榮耀或者聖潔的上帝.
上帝要人從內心裡面訴諸於良知.
看到自己生命的真相.
原來祂很想幫助人.
脫下自己生命的假面具.
不再扭曲地過人生.
所以在這裡耶穌基督特意和西門說.
不要怕,從今以後你要得人了.
有些譯法是不要怕,從今以後你要得人,如得如了.
耶穌將這個豐盛的漁獲.
比喻為在未來的日子裡.
你不需要怕.
你將來的生命就像一個得著漁獲的漁夫一樣.

$^{281}$不過未來的日子.
他所得到的就不是這些魚.
而是得到人的生命.
藉著他的生命見證.
藉著他回應耶穌基督的呼籲.
他的生命就可以參與在耶穌基督的團隊裡.
一起將福音改變這個世界.
生命第四個真相是.
人雖然是被污.
或者人生命裡有很多遺憾.
有很多不濟.
但從來上帝沒有嫌棄我們.
相反來說.
上帝想用恩典來看到我們自己生命的真相.
當我們自己知道自己是一個怎樣的人.
我們不是停留了.
去定性了自己永遠都是這麼庇護.
反而是藉著自己的真相.
我們願意面對上帝.
耶穌基督在這裡呼籲彼得.
要跟隨他.
彼得,約翰和雅各撇下了這個慾望.
全心全意地跟隨耶穌.
他走這個全道的生涯.
成為基督耶穌最重要的生命裡的門徒.
弟兄姊妹.
不知道你今天讀這段經文.
你有沒有同感呢?.
我們會不會有時候都會覺得人生短暫.
形形異異.
很多我們期望得到的東西.
其實沒有辦法得到.
我們自覺生命充滿了很多限制.
有時候當我們再仔細去看的時候.
我們也會發覺自己的視野很狹隘.
我們只關心我們自己覺得重要的東西.
又或者以我們自己的經驗.
來判斷所有事物的對與錯.
我們只是用自己的角度去看這個世界.
不過在這個經文裡.

$^{321}$或者在這個經歷裡.
耶穌最想看到的是.
不只是純粹看到你的限制.
或者你的狹隘視野.
更加要訴諸你的良心.
看到你生命的本相.
我們每一個人都犯了罪.
規缺了神的榮耀.
每一個人都會有這些過去.
但上帝從來沒有嫌棄,沒有撇棄我們.
反而他誠意邀請你們每一個跟隨他.
放下昔日的遺憾.
背起十字架.
面向上帝所走過的路.
我們一步一腳印地跟隨他.
走未來人生的道路.
今天晚上的陶奎禮.
除了告訴我們.
我們生命本是塵土.
仍要歸於塵土.
我們的生命本上其實沒有什麼地方很了不起.
不過正正是這一份對自己生命的自覺.
反而是為我們帶來一個新的出路.
因為我們知道.
即或我們本於塵土.
歸於塵土的生命.
但上帝願意將一個榮耀的即時.
附託給你和我.
我們既帶著上帝的赦免.
亦都帶著上帝的召命.
繼續走未來人生的道路.
一瞬間.
藉著這個灰燼.
塗上在你的額頭或手背上的十字架.
提醒你.
你的生命不再靠自己而過.
而是靠著我們所養賴.
我們所信靠的基督耶穌.
走未來人生的道路.
主繼續來激勵我們.

$^{361}$幫助我們.
愛我們到底的耶穌基督.
願意為我們而死.
要我們活著的不再是我.
而是耶穌你活在我們裡面.
我們可以靠你去與父神.
與自己.
與其他人和好.
願聖靈幫助我們.
在這四十天裡面.
初戀我們能夠親近基督更多.
學掌基督更多.
天父啊.
我們祈求你.
我們尋找你.
我們向你叩門.
願天父給我們勇氣.
智慧.
叫我們愛人愛神.
願基督為朋友寫命.
和犧牲的愛.
時常在我們心裡.
以致我們.
激勵我們.
願意去實踐基督的愛更加多.
奉主耶穌基督得勝明志祈求.
阿們.
讓我們藉著.
哥林多後書第六章.
一至十節.
一同去領受.
上帝託付給我們的使命.
我們與神同供的.
也勸你們不可徒受他的恩惠.
因為他說.
在月立的時候.
我應運了利.
在拯救的時候.
我答救了利.
在祈求的時候.

$^{401}$我答救了利.
看啊.
現在正是月立的時候.
現在正是答救的日子.
我們凡事都不叫人有妨礙.
免得這積分被人毀謗.
反倒在各樣的事上.
表明自己是神的用人.
就如在許多的患難.
忍耐.
貧乏.
困苦.
變打.
監禁.
擾亂.
勤勞.
警醒.
不食.
廉潔.
知識.
行忍.
恩慈.
聖靈的感化.
無偽的愛心.
真實的道理.
神的大能.
人而的摒棄.
在左在右.
榮耀.
羞辱.
惡名.
美名.
似乎是有人的.
卻是誠實的.
似乎不為人所知.
卻是人所共知的.
似乎要死.
卻是活著的.
似乎受認罰.
受責罰.

$^{441}$卻是不自喪命的.
似乎憂愁.
卻是常常快樂的.
似乎貧窮.
卻是叫許多人付足的.
似乎一無所有.
卻是樣樣都有的.
讓我們同心低頭禱告.
主耶穌.
藉著剛才的經文.
提醒我們每一個.
我們所夢的因.
全是因為上帝.
你捨身來成就出來.
讓我們知道.
我們作為門徒.
我們要活出一個.
與蒙恩相稱的生活.
因為我們知道.
我們的主.
也是一位受苦的主.
通過苦難來成就救恩.
換句話來說.
作為你的門徒.
求主幫助我們.
願意為主你去受苦.
讓我們也通過.
在生活裡面.
面對不同的困難.
不同的逆境.
我們可以在當中.
為主你作見證.
以致我們在這些艱難的裡面.
在這些逆境的裡面.
我們可以和耶穌基督有份.
並且我們知道.
我們也可以靠著.
耶穌基督的幫助.
你的帶領.
我們可以跨過這些難處.

$^{481}$我們知道耶穌.
你能夠將這一切都轉化成.
對我們的幫助.
對我們的成長和祝福.
以致我們的生命.
獲得轉化.
我們能夠在生活裡面.
因著主你的緣故.
有更多的韌力.
有更多的生命力.
無論順境,逆境.
無論得失,成敗.
我們都要堅守.
上帝的吩咐.
耶穌基督你所負的使命.
以致無論我們得事不得事.
我們總要叫.
耶穌你的名照常顯大.
並且我們確信.
在裡面.
耶穌說你要與我們同在.
你要透過我們的言行.
讓世人明白.
讓世人經歷何為善,何為美.
並且確信.
沒有任何事情.
能夠難阻耶穌基督.
你對我們每一個的愛.
所以為此.
我懇求上帝幫助我們每一個.
求聖靈你.
內注在我們的生命裡面.
掌管我們的一切.
讓我們整個生命.
都能夠被你使用.
讓我們無論在人生的高高低低.
起起跌跌的裡面.
我們都要活出一個.
與蒙恩相稱的生活.
我們不單止與耶穌你同苦同死.

$^{521}$也會經歷與耶穌你同生同活.
我們為此懇求上帝你垂聽.
我們在你面前的禱告.
奉耶穌基督你得勝的名.
來祈求.
Amen.
弟兄姊妹我們一同站立.
同心領受上主而來的祝福.
同心祈禱.
願上帝你幫助我們.
在地上的日子.
讓我們能夠剛強壯膽.
靠賴你的恩典.
我們生命活得有意義.
更加我們要找到生命的方向.
撇除我們生命裡面的庇護.
讓我們重新立志在你面前.
立志要跟隨你一生.
走你要我們所行的道路.
深願我主耶穌基督救贖恩惠.
天父上帝的慈愛.
聖靈的感動激勵.
上上願你們每一個同在.
從今時直到永永遠遠.
Amen.
\newpage



\section{申命記 8:1-10}
\label{sec:crDuuHRkfCw}
\textbf{2025年3月9日崇拜直播｜陳冠成牧師｜朝向應許之地－謹記遵守盟約｜申命記8\_1-10}
\newline
\newline
連結: \href{https://youtube.com/watch?v=crDuuHRkfCw}{\texttt{https://youtube.com/watch?v=crDuuHRkfCw}} ~~~~ 語音日期: 2025-03-09
\newline
\newline
\hyperref[sec:mfCFFhs9u_c]{\small{< < < PREV SERMON < < <}}
~
\hyperref[sec:index_chronic]{\small{[返順時目]}}
~
\hyperref[sec:P0a0tJAFKd4]{\small{> > > NEXT SERMON > > >}}
\newline
\newline
申命記 8:1-10
\newline
\begin{longtable}{cl}
\hline
\hline
章節 & 經文 (和合本修訂版)\\
\hline
8:1 & \begin{tabularx}{0.7\textwidth}{X} 「我今日所吩咐你的一切誡命，你們要謹守遵行，好使你們存活，人數增多，可以進去得耶和華向你們列祖起誓應許的那地。 \end{tabularx} \\ \\ \relax
8:2 & \begin{tabularx}{0.7\textwidth}{X} 你要記得，這四十年耶和華---你的神在曠野一路引導你，是要磨煉你，考驗你，為要知道你的心如何，是否願意遵守他的誡命。 \end{tabularx} \\ \\ \relax
8:3 & \begin{tabularx}{0.7\textwidth}{X} 他磨煉你，任你飢餓，將你和你列祖所不認識的嗎哪賜給你吃，使你知道，人活著，不是單靠食物，乃是靠耶和華口裡所出的一切話。 \end{tabularx} \\ \\ \relax
8:4 & \begin{tabularx}{0.7\textwidth}{X} 這四十年，你身上的衣服沒有穿破，你的腳也沒有腫。 \end{tabularx} \\ \\ \relax
8:5 & \begin{tabularx}{0.7\textwidth}{X} 你心裡要知道，耶和華---你的神管教你，像人管教兒女一樣。 \end{tabularx} \\ \\ \relax
8:6 & \begin{tabularx}{0.7\textwidth}{X} 你要謹守耶和華---你神的誡命，遵行他的道，敬畏他。 \end{tabularx} \\ \\ \relax
8:7 & \begin{tabularx}{0.7\textwidth}{X} 「耶和華---你的神必領你進入美地，那地有河流，有泉源和深淵的水從谷中和山上流出。 \end{tabularx} \\ \\ \relax
8:8 & \begin{tabularx}{0.7\textwidth}{X} 那地有小麥、大麥、葡萄樹、無花果樹、石榴樹，那地也有橄欖油和蜂蜜。 \end{tabularx} \\ \\ \relax
8:9 & \begin{tabularx}{0.7\textwidth}{X} 那地沒有缺乏，你在那裡有食物吃，一無所缺；那地的石頭是鐵，山中可以挖銅。 \end{tabularx} \\ \\ \relax
8:10 & \begin{tabularx}{0.7\textwidth}{X} 你吃得飽足，要稱頌耶和華---你的神，因為他將那美地賜給你。」 \end{tabularx} \\ \\
[1ex]
\hline
\hline
\end{longtable}
$^{1}$今日經訓是《新命記》八章一至十節.
是記載在回眾的座位前面的聖經九曰部分.
224頁至225頁.
聖經這樣記載.
我今日所吩咐你一切戒明.
你們要謹守遵行,好叫你們存活,人數增多.
且進去得耶和華向你們列祖起誓應許的那地.
你也要紀念耶和華你的神,再抗夜引導你這四十年.
是要苦練你,試驗你,要知道你心內如何.
肯守他的戒明不懇,他苦練你,任你饑餓.
將你和你列祖所不認識的嗎哪賜給你吃.
使你知道人活著不是單靠食物.
乃是靠耶和華口裡所出的一切話.
這四十年你的衣服沒有穿破.
你的腳也沒有腫.
你心裡思想,耶和華你神管教你.
好像人管教兒子一樣.
你要跟守耶和華你神的戒明.
進行他的道,敬畏他.
因為耶和華你神領你進入美地.
那地有河,有泉,有源.
從山谷中流出水來.
那地有小麥,大麥,葡萄樹,無花果樹.
石榴樹,橄欖樹,黃麥.
你在那地不缺食物,一無所缺.
那地的石頭是鐵.
山內可以挖銅,你吃得飽足.
就要稱頌耶和華你的神.
因為他將那美地賜給你了.
弟兄姊妹平安.
剛才閻明弟兄所讀的這張申命記.
一共有二十節.
如果留意到聖經裡的記載.
剛才讀了上半部分.
一至第十節.
裡面提到提醒百姓要紀念.
上帝昔日在曠野四十年的作為.
並且以敬畏的態度.
去承受迦南地的福氣.
至於下半部分,稍後我們會說.

$^{41}$十一至二十節.
裡面就不同了.
裡面提醒百姓不要忘記.
摩西警誡百姓不要因為.
生活漸漸富足而忘本.
仍然要記得上帝.
事實上紀念可以說是這張聖經.
貫穿了主題.
貫穿了段落.
紀念是聖經常常提到的.
常常出現的題目.
當然不只是回憶.
上帝以前做了什麼.
其實紀念有一個很重要的意思.
就是你要重新檢視.
你和上帝的關係到底是怎樣.
你和上帝的距離到底是近還是遠.
紀念是說.
我們和上帝立過什麼約.
並且在這份約的內容是什麼.
譬如說上帝應許了他要做什麼.
另一方面.
上帝又要求我們怎樣守約.
並且說明.
如果你遵守的時候.
上帝會應許什麼.
反過來你不遵守的時候.
違約的時候.
上帝又會怎樣做.
紀念驅使我們.
更加認識上帝是怎樣的神.
你要更加委身忠心於上帝.
並且願意遵從他所說的話.
讓上帝的話語塑造你日常生活.
那個態度.
那個優先次序和價值觀念.
所以我們說.
聖經裡面說的紀念.
一定不是靜態.
不只是腦袋的運轉.

$^{81}$紀念是透過你回顧.
上帝在你生命裡面的作為.
你一定有相應的行動.
來配合上帝的吩咐.
正如今天我們所看的.
申命記第八章.
你會留意程序表都印了所有經文.
我們說有三個段落.
每一個段落你會發現.
第一節,第六節和第十一節.
紀念一定會生出一個效果.
就是你要謹守.
遵行上帝的誡命.
也就是你要謹守.
你和上帝所納過的約.
也就是今天的講題.
我們逐個段落去看.
首先第一至到第五節.
摩西提醒百姓.
要紀念在曠野.
上帝如何引導他們這四十年.
這段日子我們之前都說過.
是因為上一代.
我們說第一代的以色列人.
他們不聽上帝的吩咐.
背叛上帝.
所以遭受懲罰的結果.
不過摩西現在和新一代的以色列人.
回顧這段慘痛的歷史.
他要求他們.
用另一個角度去看這四十年.
在那位滿有智慧和恩典的上帝底下.
其實這四十年的意義已經轉化了.
從上一代試探神是一件惡事.
現在上帝逆轉了.
對新一代來說.
這是一個靈命成長的功課.
上帝很想透過這四十年.
雖然不容易.
經文裡也說了.

$^{121}$是一個苦練.
是一個試驗.
但卻可以煉證以色列人的生命.
引導他們在曠野.
這個我們說很困乏的歲月裡.
他們會知道三件事.
經文裡也有三個方框.
他要知道三件事.
在這四十年.
第一個就是知道.
到底是否忠於上帝的吩咐.
廣東話說患難就是見真情.
這是很實際的.
身處在曠野是一個什麼地方.
你一定要想像一下.
我們今天沒有什麼機會經歷.
什麼都沒有.
很缺乏的地方.
再加上你不知道前路是怎樣.
不是短期.
不是四個鐘四十天.
是四十年.
在這個情況裡.
人裡的真面目.
什麼都會顯現出來.
特別是裡面的負面情緒.
會有什麼呢?.
憂慮.
會有擔心.
你會有很多埋怨.
甚至怕.
怕到一個地步你會憤怒.
你這一切的負面情緒.
都會在上帝面前表露無遺.
上帝看得一清二楚.
但正好也是一個測試.
測試在這些情緒裡.
你的靈命是怎樣.
無論外面的環境是怎樣.
就算多惡劣也好.

$^{161}$你內心可能真的有很多不安不滿也好.
你能不能管理到它.
情緒固然是正常的.
負面情緒也是一些人的機制.
上帝給我們一些反應.
但是你怎樣按著上帝的應許.
按著上帝的帶領.
你可以剋性你裡面的情緒.
將它合而表現出來.
不會過了農.
不會因為你肚餓就責罵上帝.
不會因為你有很多埋怨.
就將其他人完全摒棄.
來忽略他們.
事實上弟兄姊妹.
曠野正好是考驗.
一個人在缺乏的裡面.
在充滿負面的裡面.
他仍然能不能夠因為上帝指明的身份.
按著上帝的話語和應許.
他仍然可以好好跟自己的情緒相處.
跟身邊的人相處.
跟上帝相處.
仍然堅信上帝的帶領是好.
上帝這一刻的安排仍然是好.
我仍然要遵守上帝的吩咐.
上帝就是透過抗疫這四十年.
我們說是一個長期的考試.
來了解百姓的心想什麼.
你怎樣想上帝.
你怎樣想其他人.
你是不是對上帝忠心.
當然在這個考驗裡面.
百姓也會意識到.
原來我都很難處理.
為什麼上帝要用四十年?.
因為我很難處理.
我的靈命是很橫緊.
不然四十天就能處理.
四十年也反映人是很難處理.

$^{201}$所以上帝才用長期的磨練.
來讓他們成長.
所以這四十年.
不單是上帝要試驗百姓的內心.
其實也是給機會百姓苦心自問.
我會否合格?.
我的靈命能否過上帝的標準?.
這是第一個要知道的東西.
在疫情裡面.
你要知道你跟上帝的關係是怎樣.
你要知道你的心怎樣想上帝.
怎樣想你周邊的群體.
第二個要知道的是什麼?.
第三節這裡說.
我靠什麼來活著?.
經文這裡說.
上帝苦練你,任你飢餓.
將你和你列祖所不認識的嗎哪賜給你.
使你知道.
這句我想就算你沒看過《申命記》.
你也知道出自哪裡.
耶穌基督也引用過.
在他面對試探的時候.
人活著不是單靠食物.
而是靠耶和華口裡所說的一切.
聖經學者估計.
當時以色列人離開埃及的時候.
人數可能有二百萬.
甚至三百萬人.
這麼多人在曠野.
而且是長期的生活.
很自然就會面對什麼困難?.
缺水缺糧.
你要留意.
曠野是你想謀生也沒有機會.
種不到東西的.
什麼都沒有.
在這樣的情況下.
飢餓一定是每一個人.
每一天.

$^{241}$每一餐都要面對.
但又如何?.
你沒辦法解決.
你費盡窮荒之力也解決不了.
因為根本沒有東西給你做食物.
你沒能力去解決.
但上帝就按著他的應許.
我會照顧你的.
我吩咐天降嗎哪.
解決百姓的需要.
摩西在這裡特別強調.
以色列人的先祖.
包括他們第二代的以色列人.
他們不知道嗎哪是什麼.
不認識.
不認識不單止.
不知道嗎哪是什麼成份組成.
不認識包括.
我不知道如何製造嗎哪.
我每天都在用.
但我沒能力自己製造.
複製嗎哪出來.
嗎哪是完全100\%.
我們說是神製造.
百分百是上帝生產製造.
完全要依賴上帝的供應.
並且百姓不單止這方面依靠上帝.
用嗎哪的時候.
上帝說有使用守則.
有使用守則.
你要跟著上帝的吩咐來享用.
按照《出埃及記》記載.
嗎哪是每一個百姓.
每家每戶.
你要清早.
日頭還未曬的時候.
就要清早出來.
按照個人的飯量.
要出外拾取.
晚一點是不行的.

$^{281}$就好像有時你去酒店.
有些酒店包早餐.
說明6-9點.
晚了怎麼辦?.
沒有了.
收了自己處理.
嗎哪也是一樣.
你要清早去拾取.
否則太陽一出來發熱.
嗎哪很奇妙的.
溶化消失了.
不見了.
你想找也找不到.
所以嗎哪是每天的神蹟.
你拾取到的嗎哪.
也不能留過隔夜.
不要想著拾取多一點.
等明天吃.
明天休息一天.
不行.
上帝命定了.
第二天生蟲發臭.
你想留也留不了.
另一方面.
嗎哪也不是每天那麼簡單.
是每個星期的神蹟.
為何這樣說呢?.
上帝命定.
安息人不可以拾嗎哪.
因為不可以動工.
於是上帝又吩咐.
你們要在安息日前一天.
拾取.
安息日前一天.
拾取雙份.
拾取雙份後.
預早煮好.
等安息日第二天.
你可以拾取來吃.
上帝保證.

$^{321}$這種隔夜的嗎哪.
是不會發臭的.
沒有問題.
你看到嗎哪.
其實有一個很重要的功課.
就是迫使以色列人怎樣呢?.
你不能積谷防飢.
你不可以想著儲起來.
遲些享用.
不可以.
嗎哪是迫使.
這群百姓要完全相信.
上帝每天供應是足夠的.
每天供應足夠的.
不要想第二天.
想這一天.
換句話來說.
嗎哪表面上的功能是什麼呢?.
填飽我們的肚腹.
實際上.
好像聖經所說.
引導百姓這四十年來.
你要學會.
每天靠上帝所說的法則來生活.
不是靠嗎哪這麼簡單.
是靠上帝說.
怎樣去享用嗎哪.
使他們知道.
不是嗎哪本身使他們活著.
而是那位供應嗎哪的上帝.
使他們活著.
並且是因為他們.
遵行上帝的話語.
使他們活著.
所以對百姓來說.
這就是他們長期的操練.
而上帝對他們說過什麼呢?.
不單單是嗎哪怎樣享用.
也包括上帝在跟他們納藥的時候.
上帝說過什麼呢?.

$^{361}$上帝說好了什麼呢?.
簡單來說就是.
上帝承諾.
我會帶領保守我的子民.
一定會這樣做.
上帝承諾了.
而在這個藥裡面.
百姓也承諾了.
我們會百分百聽上帝吩咐.
上帝說什麼我們照做.
於是乎在抗疫的裡面.
在飢餓的裡面.
上帝遵守了他的承諾.
降嗎哪滿足百姓的需要.
百姓也遵守了他的承諾.
按照上帝的吩咐來享用嗎哪.
雙方都遵守了承諾.
這兩樣東西加起來.
就是以色列人這四十年.
存活的關鍵.
讓百姓真實體會到.
好像這段經文所說.
到底我們靠什麼來活著.
事實上,弟兄姊妹.
今天我們不需要去拾取嗎哪.
基本上我們都是上帝.
透過我們的工作.
來供應我們生活所需.
但是我們在日常生活.
有沒有一些熟能操練.
可以幫我們每一天.
都像色列人那樣.
明白我們是靠什麼來活著.
問問大家.
你今早是怎樣起床的?.
鬧鐘?.
手機?.
還有沒有?.
家人叫?.
自然醒?.

$^{401}$自然醒.
自然醒都不錯.
剛才我都會用鬧鐘.
用手機.
偶爾家人都會叫.
但有時這三樣都不湊巧.
有試過嗎?.
睡過籠.
誰叫醒你?.
上帝拍醒我.
剛才有人說自然醒.
有沒有想過自然是不會醒?.
為什麼你必然是自然醒?.
為什麼不可以是自然不醒?.
沒有上帝的批准.
任憑你手機鬧鐘.
家人自然醒.
通通都不湊巧.
你醒來是因為上帝准你醒.
沒有上帝的准許.
那些我們覺得自然.
理所當然不過的事.
是不會發生.
所以你今天不用吃嗎哪.
但你每天起床想起什麼?.
原來我起身活著.
是上帝讓我醒.
是上帝讓你醒.
這一下很動搖.
這一下主宰你過一天的言行.
基本上主宰了.
原來我今天醒.
代表上帝仍然有事給你做.
仍然有任務交給你.
仍然有功課給你學.
這樣你整個人就不同了.
你知道你是為主而活.
不是我去打工.
不是我對著那些很難受的人.
你是為上帝作見證.

$^{441}$事實上上帝已經指明.
如果連最基本最簡單.
睡醒起床.
都是賴於上帝的供應.
更何況說到生命的意義.
人生的目的.
這些重大更複雜的事.
你不能不相信上帝.
依靠他.
聽從他的吩咐來生活.
這是第二個你要知道的事.
第三個.
三個要知道和合本益造思想.
在第四至第五節裡.
我們要知道上帝和我們有什麼關係.
確實在曠野這四十年.
是不容易過渡的.
是很難捱的.
但摩西讓他們在困境裡.
你不要只看困境.
你看看.
剛才Jenna大師也有.
你看看你的衣服.
有沒有爛掉.
沒有.
你看看你的腳.
有沒有腫脹.
沒有.
應該發生的事.
為什麼不發生.
是應該會爛的.
是應該會腫脹.
為什麼不發生.
很明顯不是因為我們幸運.
不是幸運.
更加不是我們的本事.
是出於上帝的供應和保守看顧.
沒錯.
上帝說到明是要管教我們.
但上帝不是把我們當犯人教.

$^{481}$不是把這些以色列人當犯人.
如果把我們當犯人教會怎樣.
一不聽話一犯錯.
第一件事就是.
斷水斷糧.
第二件事.
把你的衣服和鞋子拿走.
這就是把我們當犯人教.
摩西提醒.
上帝對百姓就像父親管教兒子.
為了要讓上帝的兒女成長.
上帝是有嚴格要求的.
但同時也有恩典和憐憫.
他會體恤百姓的軟弱.
明白可能一時三刻.
有些事真的未達標.
所以上帝願意花時間慢慢教.
弟兄姊妹.
我想你在生活中.
也遇過一些不容易相處的人.
甚至有時是你的親人.
偶爾頂撞你.
不聽你的話.
氣得你死就死.
有時要和他們相處.
四十天也不容易.
更何況是四十年.
你願意嗎?.
你不想的.
又或者你去到書室.
基督教書室也好.
一般的書局也好.
你是父母來的.
你見到有兩本書放在你面前.
一本就是四十天.
可以改變你子女的生命.
旁邊還有一本.
四十年才改變到.
你買哪一本?.
或者相反.

$^{521}$你覺得哪一本會好買?.
不用說了.
沒有人說的.
你買四十年的那本?.
不會的.
不會有人傻得出一本書.
叫做四十年才改變生命.
為什麼?.
四十年後我不知道去哪了.
沒有人會有這樣的耐性.
現代人就是要快.
哪有可能要四十年.
我才看到對方改變.
浪費力氣.
人不會有這樣的耐性.
就算是地上的父母.
也未必有這麼長情.
但很奇怪.
聖經告訴我們.
上帝就是這樣.
上帝是用四十年.
真的很長情.
引導一班不聽話的人.
慢慢去成長.
為什麼上帝要用四十年.
來管教百姓?.
是不是因為上帝太過嘮叨?.
說多了也不能表達.
所以說多了也不能表達.
所以說了四十年.
還是上帝真的很有耐性.
祂對這班子民很有愛.
才肯放四十年下去?.
到底是上帝沒有能力管教.
不懂得怎麼教.
還是他管教的對象.
真的很翹口參手.
很難教.
教極也不改.
於是上帝才這麼包容.

$^{561}$接納他們.
到底是上帝硬要強迫百姓.
接受他那套說辭.
還是這關乎生死禍福的真理?.
上帝想你存活.
不想你滅亡.
所以你一直不懂.
上帝也要你學懂為止.
否則你根本不可能走下去.
弟兄姊妹.
摩西很希望百姓知道.
神是因為當他們是兒女一般.
才會管教他們.
否則放棄了.
並且摩西很想他們知道.
上帝這位父親對兒女的管教一定是好.
一定是美善有益.
所以應該要樂而接受上帝的管教.
有得被上帝管教.
代表什麼?.
還有得救.
還有盼望.
上帝利用撤走他們的管教.
就大件事了.
上帝掩面不顧的時候.
你不要以為很過癮.
那一刻是最慘的.
換句話來說.
摩西提醒.
我們要尊重上帝作為父親的管教和帶領.
要敬畏他.
否則你進入迦南也沒有用.
摩西吩咐他們.
要用一個敬畏的心.
到日後進入迦南的生活做好準備.
我們再看第二個段落.
六至十節.
你會發現這個段落.
重複出現一個字叫地.
那地,美地.

$^{601}$意思是他們稍後要進入迦南應去之地.
與曠野相比.
差天拱地.
剛才也讀過了.
與曠野的生活.
來到迦南地簡直是人間樂土.
摩西表示那裡以前沒有的.
現在什麼都有了.
水源充足.
農作物豐富.
礦產取之不盡.
樣樣都不會缺乏.
什麼都有.
很多無比的地方.
不過在說這一切之前.
摩西劈頭說.
你要承受這一切美好的事物.
有一個先決條件.
你要敬畏神.
如果不敬畏神.
這一切反而會成為你的禍患.
什麼意思呢?.
我們說摩西說迦南地這麼棒.
這麼吸引一份禮物.
照道理比在你身上.
你是不是搶著要?.
應該是吧?.
有沒有人會撒手寧頭?.
真的沒有?.
上一代就撒手寧頭.
上一代的以色列人.
也是面對著這麼棒的地.
但他們沒有選擇敬畏上帝.
而是懼怕他的仇敵.
他懼怕.
就好像剛才蔣介石那一番.
那麼巧.
懼怕那些巨人.
只有迦納和約書亞看到.
不只是看到巨人.

$^{641}$還看到上帝.
上一代的以色列人.
因為懼怕仇敵的強大.
所以對於上帝.
給他們這麼棒的禮物.
你不要給我.
我不要,撒手兼寧頭.
我死都不去.
結果當然是死在曠野.
換句話來說.
其實你會看到上帝所賜的福地.
這份禮物確實是好得無比.
但這份禮物原來也是一個挑戰.
是考驗現在.
上一代又給你看了.
現在新一代百姓.
你到底仍然是畏懼.
那些敵人的強大.
還是你這次懂得敬畏神呢?.
同一個處境.
但經過四十年之後.
會不會有更新,不同的看法.
這次這群百姓願意敬畏上帝.
聽從上帝的吩咐.
勇敢地邁向應許之地.
事實上,兄姐們.
敬畏這個功課.
也是他們這四十年.
在曠野透過吃嗎哪.
每天都操練的.
剛才我們也說過.
任憑你怎樣努力.
嗎哪是無法多勞多得的.
你想撿多點,沒用.
留不住隔夜發臭.
你無法儲備.
也留不住過夜.
上帝命定了.
除了安息日之外.
每天只能取當日的份量.

$^{681}$當然這不是鼓勵百姓.
做夠三十六便算了.
其他人可以躺平.
而是一種操練敬畏上帝的生活態度.
你努力做好上帝叫你每天要做的事.
當盡你每天的份.
同時你也相信.
仰賴上帝每天的應許.
供應是足夠的.
這就是敬畏.
不單是在曠野.
操練這四十年.
都是為了之後的生活.
想遠一點.
將來你進到迦南.
你不再需要撿嗎哪了.
上帝會改用種植農耕的方式.
繼續供應你.
但是.
敬畏的態度.
你依然要保持住.
摩西提醒百姓.
第十節這裡說.
當你進到那塊美地.
吃得飽足.
亦要稱頌耶和華你的神.
因他將那美地賜給地了.
稱頌這裡其實有另一個意思.
就是你感恩.
即是說你將.
之後生活一切那種風雨.
你將這個功勞歸回給上帝.
因為是上帝的緣故.
上帝將這塊美地賜給你.
你必須帶著一個敬畏感恩的心.
來享用上帝給你的豐饒.
免得落入一種自高自大的心態.
慢慢習慣了以為.
我之所以有這麼好的環境.
是因為我拼搏的成果.

$^{721}$是因為我的努力.
我的本事.
為什麼這樣說.
繼續看下去你會明白.
我們剛才未讀到十一至二十節.
現在讀完它.
請大家一起讀螢光幕.
看看程序表.
十一至二十節預備起.
(念念大德 以德望基耶和華你的神).
(不受大人大命 天上和地下).
(就使我為人所感悔的氣力).
(苦難未吃得飽足).
(建造美好福堂屋居住).
(你的牛羊家多 你的金銀增滿).
(你的所有的儲蓄都加增).
(你就心高氣大).
(望基耶和華你的神).
(就是將你從埃及地).
(遺留之家領出來的).
(引你經過那大而可畏的曠野).
(那裡有火蛇 蠍子).
(乾旱無水之地).
(他曾為你使水).
(從堅硬的盆石中流出來).
(又在曠野將你到).
(所不認識的嗎哪次級領遺).
(是要苦練你 試驗你).
(叫你終究享福).
(滿下你心裡說).
(這才才是我的力量).
(我能力得來的).
(你要紀念耶和華你的神).
(因為得才才的力量).
(是他給你的).
(為要堅定他向你列祖).
(你誓所立的約).
(像今日一樣).
(望基耶和華你的神).
(隨從必神).

$^{761}$(事奉敬拜).
(你們必然滅亡).
(這是我今日警戒你們的).
(耶和華在你面前).
(怎樣使列國的人滅亡).
(你們也必照樣滅亡).
(因為你不聽從).
(耶和華你們神的話).
摩西這裡預先告訴百姓.
他們稍後進入迦南的生活.
是會漸漸富足起來.
為什麼?.
因為上帝說了.
上帝應許他們這樣做的時候.
他們會守承諾.
將福氣賜給他們.
但是上帝同時.
也將這個考驗給了他們.
如果百姓能夠繼續.
將他們在曠野那四十年.
操練到敬畏感恩的態度.
來承受將來迦南地的福氣的話.
上帝會持續遵守他的承諾.
繼續賜福給他們.
上帝會首約供應百姓各樣所需.
甚至豐足有餘.
但相反,經文這裡說得很清楚.
若百姓因為漸漸富足.
就變得心高氣傲的話.
以為一切的財富.
是自己的能力,努力的成果.
忘記了上帝才是他們存活的根本.
即是忘記了他們在曠野所學成的功課.
不再聽從上帝的吩咐的話.
即是違背上帝所立的約.
你猜上帝會怎樣做?.
上帝會否首約?.
上帝是會首約的.
上帝是信實的,不會不首約.
不過不是賜福那部分.

$^{801}$而是降禍那部分.
上帝說我照首約.
但是降禍那一部分.
20節這裡說得很清楚.
耶和華在你們面前怎樣使列國的民滅亡.
你們也必照樣滅亡.
因為你們不聽從耶和華你們神的話.
說得很清楚,不需要釋經.
儘管以色列人是上帝的子民.
但無情說.
當一再不聽從上帝的話的時候.
其實也是將自己帶離了.
他和上帝所立的約.
這樣也難免上帝的刑罰管教.
甚至走向滅亡之路.
當然我們知道在以色列人的歷史裡.
上帝不會一下子就這樣.
上帝真的有很多恩典.
恩威並施.
就算往後的日子.
他的國家沒有了.
但上帝的計劃沒有變.
上帝的國度依然要彰顯.
而在這裡摩西也語重心長地.
預先告誡他們.
我們真的不要忘本.
不可以忘記上帝.
不可以因為你變得很厲害.
就輕慢了上帝的吩咐.
因為那是你存活的根本.
是你和上帝立了的約.
弟兄姊妹,事實上.
聖經的講解大概來到這個部分.
我想有一些反省應用.
可以在這裡彼此提醒.
首先,我們說上帝讓我們去享用各樣的美物.
不代表你擁有.
在我們今天活在資本主義的底下.
我們很強調.
資本財產就屬於我私人擁有.

$^{841}$我決定怎樣用,因為是我的.
但聖經卻表明.
一切的資本財產是屬於上帝.
這裡說,十八節.
就連賺取這些財富的力量.
上帝說都是他給的.
上帝要收回的是輕而易舉.
上帝讓我們享用一切的財物.
甚至乎可能是豐厚.
沒錯,我們是用家.
但不是用家至上,而是用家大.
經文這裡清楚說.
擁有權的主權由始至終在上帝手上.
譬如證據.
你戶口儲存了五百萬,一千萬.
你用什麼心態去看.
你認為這是你努力賺取儲存的成果.
所以我有權任意去使用.
還是你認為這是屬於上帝的資源.
上帝先放在我手上.
我是託管的.
我受託好好運用這些資源.
兩者的心態完全不同.
出來的結果也截然不同.
取決於我們對上帝的態度.
我們對上帝的公認.
有沒有一個敬畏感恩的心.
知道我們雖然是用家.
但同時我們是管家,不是主人.
我們會否按照聖經的教導.
善用上帝給我們賜這些資源.
讓上帝得到榮耀.
讓人因著上帝得到造就和祝福.
這是第一點我們值得去思考.
第二點經文給我們應用的向度是.
富足是否等於幸福.
中國人的文化總覺得.
我越富裕,越充足.
就是等於幸福.
越多的財富物質.

$^{881}$就代表我的生活越滿足,越幸福.
所以這個想法也令到.
很多人會將生活的不滿.
覺得自己不幸福.
歸咎於什麼?.
我的物質不夠,我缺乏.
我賺得不夠,我未夠富足.
但是,《新命記》清楚告訴我們.
很特別,無論貧窮缺乏的日子.
也可以是百姓最近上帝的日子.
抓住上帝不放的日子.
因為貧窮和缺乏.
你就知道沒有東西可以靠.
你想發力也沒法力.
你一定要抓緊上帝.
因為只有他供應你.
反而這個時刻,你因為缺乏.
你抓得上帝最緊.
你最信服,信靠他.
你的靈性反而是好.
相反,物質風搖.
可能看似是幸運.
但是如果根據經文所說.
也可以是禍患.
是你信仰一個隱患.
導致我們會忘記了上帝.
離棄了他的吩咐.
所以,弟兄姊妹,貧窮或者富足.
是好是壞.
其實是視乎你用什麼心態去應對.
正如《箴言》的作者.
他很明白這一點.
他也寫下他對貧窮和富足的領袖.
我們一起讀,好嗎?.
最後那節經文印在這裡.
《箴言》30章的八節下至第九節.
請大家一起讀,預備起.
使我也得貧窮也得富足.
使我需要的飲食.
恐怕我飽足不見.

$^{921}$你少,給我八十歲命.
又恐怕我貧窮就糾纏.
以至始終我身未平.
好,其實我想.
粗略去估計.
世人向他們所信的神.
無論是上帝也好.
祈求的內容.
大概你會發現.
多數都是求富足.
很少人會去求自己貧窮.
因為人總覺得.
我要富有了才會滿足.
所以很少人會像.
《箴言》作者的祈求.
既不貧窮也不會富足.
但這正正是他作為智慧人.
對生命的本相的理解.
確實,問世間有多少人.
能夠忍受極度貧窮.
而不會對上帝發怨言.
又有多少人能夠忍受極大的缺乏.
仍然堅持上帝的真理.
不會背棄上帝的吩咐.
很少.
《箴言》作者很體會到.
人真的很渺小,很脆弱.
人生也充滿了這些試探.
另一方面,又有多少人.
能夠在大富大貴的生活中.
仍然懂得謙卑下來.
依靠上帝.
清醒地知道一切來自上帝.
並且他不是擁有者,只是管家.
又有多少人在生活物質風搖的時候.
總是保持一個清醒敬畏的心.
仍然以上帝的吩咐為優先.
《箴言》作者亦知道.
他沒有能力勝過富足的試探.
所以與其這樣,他寧願放棄富有.

$^{961}$他寧願與上帝保持緊密的關係.
既不貧窮,也不富足.
只是求上帝的恩典.
剛剛夠用的禱告.
就是《箴言》作者.
所追求的人生智慧.
他既承認自己是軟弱的.
很差勁,很容易失腳.
但亦尊重上帝的吩咐.
才是我生命的本.
正如我們今天看申命記.
他勉勵我們.
無論你的日子如何.
無論是順是逆.
無論貧或富.
或如耶穌所說.
我們會否先求上帝的國和耳.
在上帝眼中,貧窮富足.
其實也可以成為祂祝福我們的方法.
但原則是.
上帝的供應是足夠.
上帝是信實的.
上帝會堅守這份約.
與此同時.
他亦希望與他立約的子民.
包括你和我.
無論在甚麼環境.
都堅守與上帝所立約.
認定上帝才是我們人生唯一的保障.
一個有上帝同在.
與神同行的生命.
才是真正滿足豐盛的人生.
我們同心祈禱.
多謝天賦你給我們.
這段寶貴的話語.
亦讓我們明白生命的本相.
我們活著是在於你.
在於你的應許.
在於你的拯救.
在於你對我們的忍耐和恩典.

$^{1001}$並且是在於你的話語.
所以我們才能夠存活.
然而上帝是我們的確實.
我們是脆弱的.
我們是塵土.
我們很多時候都不明白這個真理.
我們以為活著是要靠我們自己.
我們以為要靠我們自己努力.
來賺取,來積存,來打拼.
殊不知上帝可能正在走一些.
與我們心意相反的方向.
但我們知道上帝很有耐性.
正如你對你的百姓,你的人都很有耐性.
你願意慢慢教我們.
總是給我們機會.
很想我們迴轉.
明白原來人生無論順逆貧或富.
都離不開上帝你的愛.
上帝你的法則.
才是我們存活的根本.
盼望上帝幫我們重新認定這一點.
當我們紀念上帝你的救恩.
紀念上帝你在我們身上的作為的時候.
又或者特別當我們每天醒來的時候.
我們知道原來是因為上帝你的恩典.
我們才可以活著.
既然上帝你讓我們仍然活著的時候.
我們就要按照聖經的吩咐進行你的話語.
因為這是上帝你做我們的旨意.
你拯救我們的目的.
以致我們可以在地上為你作見證.
讓上帝你的國度臨到地上.
盼望上帝你幫助我們.
無論我們面對的僵境是貧或富.
順或逆.
讓我們知道上帝你的恩典救勇.
讓我們在難處裡.
同時看到上帝你的時候.
我們就生活有力.
有勇氣.

$^{1041}$有決心.
來跟隨上帝你而行.
我們這樣的禱告交託.
奉主耶穌基督你得勝的命來祈求.
阿們.
(下集待續).
(音樂播放).
\newpage



\section{申命記 }
\label{sec:P0a0tJAFKd4}
\textbf{2025年3月16日崇拜直播｜羅志雄牧師｜朝向應許之地－恩典中的悔改與更新｜申命記918-29}
\newline
\newline
連結: \href{https://youtube.com/watch?v=P0a0tJAFKd4}{\texttt{https://youtube.com/watch?v=P0a0tJAFKd4}} ~~~~ 語音日期: 2025-03-16
\newline
\newline
\hyperref[sec:crDuuHRkfCw]{\small{< < < PREV SERMON < < <}}
~
\hyperref[sec:index_chronic]{\small{[返順時目]}}
~
\hyperref[sec:LKA6QvIWId8]{\small{> > > NEXT SERMON > > >}}
\newline
\newline
$^{1}$經文:舊約新命記,9章18-20,25-29,226-227.
一群以色列人得罪了神,為自己鑄造偶像,摩西如何為他們祈禱?.
因你們所犯的一切罪,行了耶和華眼中看為惡的事,惹他發怒,我就像從前伏伏在耶和華面前四十晝夜,沒有吃飯也沒有喝水..
我因耶和華向你們大發烈怒,要滅絕你們,就甚害怕..
但那次耶和華又應允了我,耶和華也向阿倫甚是發怒,要滅絕他..
那時我又為阿倫祈禱.我因耶和華說要滅絕你們,就在耶和華面前照舊伏伏四十晝夜..
我祈禱耶和華說:主耶和華,求你不要滅絕你的百姓,他們是你的產業,是你用大力求救贖的,用大能從埃及領出來的..
求你紀念你的僕人亞伯拉罕以撒雅各,不要想念這百姓的橫行邪惡罪過,免得你令我們出來的那地之人說:.
耶和華因為不能將這百姓領進他應許之地,又因恨他們,所以領他們出去,要在曠野殺他們..
其實他們是你的百姓,你的產業,是你用大能和新出來的,傍被領出來的..
弟兄姊妹主內平安!.
上個主日,Ray牧師宣講了關於謹記遵守文約,提出摩西要警告提醒新一代以色列人進入迦南應許之地,.
他們那裡吃得飽足,享美物時,也提到我們要紀念耶和華,.
要堅守耶和華的約,不要因為有飽足美食而得意忘形,以致心高氣傲,忘記神..
摩西警告他們,如果不聽從耶和華神的話,就必滅亡..
今日剛剛陳玉峰弟兄和我們讀出的程序表,這兩段經文,大家都聽到,.
弟兄姊妹有沒有聽到呢?.
去到第九章的時候,第一次是這樣說,.
以色列你當聽,王震弟兄姊妹你當聽,.
究竟聽些什麼呢?不如聽聽歷史好不好?.
我先說我自己的歷史,可能都是一些臭屎,.
話說我大概六七歲的時候,.
當時你知道六七歲,你知道我什麼年紀,.
你會知道那些環境是怎樣的,.
我爸爸在一個鄉村,應該是八鄉的地方,.
但他不是耕田,他在開山寨廠,.
做藤椅,你知道做藤椅,做工,有時候也會有些錢,.
我不知道為什麼,無端端去偷錢,.
偷爸爸的錢,衣袋,褲袋,甚至他賴在桌上的地方,.
但每次偷多少呢?其實那年代都是偷很少錢,.
不知道大家有沒有見過那些五元紙,.
啡色那些五元紙,偷偷埋埋,其實都是很少錢,.
當時在鄉村做什麼呢?那時候的小朋友,.
當然要群埋他們,要有得食,玩,有得食,.
但被爸爸發現了,你知道發現了什麼事,.
當然是藤條炆豬肉,打完之後都會偷,.
偷偷地偷,有一天,爸爸和媽媽就走來說,.
你下期要轉校,其實那時候已經轉了一兩次校,.
繼續轉校,轉去青山新墟的地方,.
在白鄉去新墟都很遠,不是每天上學放學,.

$^{41}$而是去那裡寄宿,寄宿在那裡,年輕人兩兄弟,.
去那裡讀書,沒辦法,爸爸媽媽說要這樣做,.
在那裡,那間又不是宗教學校,.
但有聖經堂上,聽了很多,.
偷錢的罪,如果是這些蒼霸,越早揭,當然越好,.
今天這一段聖經告訴我們,其實都很嚴重,.
弟兄姊妹,不知道你們的人生裡面有沒有一些歷史蒼霸呢?.
仍然都深深藏在你們的內心,.
但久不久就若隱若現,但你又很想包得密密實實,.
但有時候又不想提,一提起就有些東西撩起,.
但原來當揭開的時候,其實那個狀況是怎樣?.
就會很腐爛,甚至會很臭,會流些膿出來,.
但這樣算了,不要搞那麼多,如果再理他,.
都會很大件事,就當沒事發生,.
不過聖經裡面告訴我們,如果就這樣蓋起來,不理他,.
當沒事發生,上帝會發怒,會有所行動,.
今天剛才讀的經文,在第九章裡面,.
不知道大家有沒有開口讀出第九章呢?.
剛剛過去的禮拜三,我們誠心所願的嘉賓告訴我們要讀聖經讀出來,.
我真的回去試,預備這段經文讀第九章,.
我今天沒有時間讀,大家留意一下,.
經文裡面說的話,我一邊讀都很害怕,為甚麼呢?.
上帝說要滅絕,硬著頸看,橫梗,各樣東西,.
當我們一邊看第十八節,剛才我們在經文十八節開始,.
但如果再看前幾節的時候,耶和華說甚麼呢?.
看到那些百姓犯了很多罪,硬著頸看,.
但神又很奇怪,在那裡他跟摩西說甚麼呢?.
他跟摩西說,我看到這班百姓是硬著頸看的百姓,.
你且遊著我,我要滅絕他們,將他們的命從天下屠沒,.
當上帝跟摩西這樣說的時候,摩西清楚回顧這段歷史,.
大家一直看下去都知道,是金牛毒的事件,.
他說的細節是非常清晰,是向著第二代的以色列人來說的,.
在這裡我又想起一件臭事,小朋友,我都有一個兒子,.
小學階段讀書,父母都會罰小朋友,.
不知道大家怎樣罰呢?我也是一樣,.
但我一次罰的時候,覺得罰完,但兒子仍然不聽話,.
他做錯事,但仍然不聽話,罰了也在搗亂,.
於是那次真的很生氣,我立刻罰了他拉開大門,.
罰出去站著,我從來沒有這樣做,罰小朋友站在大門口,.
我又把大門打開,兒子當然哭了,哭的時候當然驚動媽媽,.

$^{81}$媽媽走出來問什麼事,我心想,你由得我,.
如果我一早告訴他,你猜我會不會罰他站在大門口?.
大家可以想想,經文告訴我們,上帝神就是這樣,.
其實就是要滅絕那些以色列人,沒得救了,.
但當摩西聽到神這樣說,他就說你由得我去做,.
殺光他們,他就知道這件事是非常之嚴重的事,.
所以摩西說,因你們所犯的一切罪,.
行了耶和華眼中看為惡的事,惹他發怒,.
摩西就像從前俯伏在耶和華面前四十晝夜,.
沒有吃飯,也沒有喝水,我因耶和華向你們大發烈怒,.
要滅絕你們,就甚害怕,但那次耶和華又應運了我,.
耶和華也向阿倫甚是發怒,要滅絕他,那時我又為耶和華禱告,.
在最關鍵的時刻,摩西非常之害怕,.
俯伏在神的面前,結果耶和華應運了他,.
當然,如果大家留意第九章,那些罪證確實存在,.
就算是禱告也好,後果也會負上沉重的代價,.
當時百姓會受到上帝的審判,.
如果有留意經文,有些平衡的經文就說,當時確實會有死人,.
死了一批百姓,被上帝,被神擊殺,滅絕,.
如果不是摩西來禱告,正如詩篇裡說,.
所以他說要滅絕他們,若非有他所揀選的摩西站在當中,.
使他的怒氣轉燒,恐怕神就滅絕他們,.
全部都死了,變成連第二代的以色列人都不可能進到迦南地..
弟兄姊妹,這就是神的恩典,祂有憐憫,.
否則後果就非常難以想像,至於亞倫又如何呢?.
亞倫作為大祭司,耶和華向亞倫甚是法勞,.
即是說層次高了,甚是法勞,他是極度非常之法嬲,.
是對著亞倫的法勞,是要滅絕的,不過摩西仍然嘗試為亞倫禱告,.
否則他也會喪命,如果我們一直看第九章的時候,.
金牛族的事件確實是有難有臭的瘡疤,.
這一段的歷史,如果大家繼續想下去,.
它的破壞力非常之大,亞倫作為大祭司是當時候的領袖,.
他沒有憤戚去阻止百姓要做出這件事,.
他竟然不知發生什麼事,上身發瘋,竟然去做,.
跟他們說,你拿光金銀首飾去做這個偶像,.
所以神就是這樣要在他的身上去擊殺他,.
摩西在申命記回顧這件事的時候,.
是要我們去想一想,今天有很多的教會,.
大教會也好,人數多也好,人數少也好的堂會,.
如果我們當中的牧者或者領袖出了問題,不能夠見證神,.

$^{121}$好像亞倫那樣,他沒有去糾正會眾所做錯的事,.
犯罪的事,是神憎惡的事的時候,.
他不單止是這樣,他還要做一份的時候,.
我相信上帝今天仍然會擊殺他,.
當我看到這裡,我在教會裡服事,.
如果不是戰戰兢兢,還有很多我沒有被揭發,.
是被大家揭發的時候,我想都很難想像,.
弟兄姊妹,揭開歷史的瘡疤,.
摩西是不是要在第二代的以色列人面前,.
來羞辱他們的先祖,他們的上一代呢?.
不知道大家,當看第九章,是不是有這種想法?.
為什麼摩西這麼差?.
是不是要將這些事件揭出來呢?.
摩西這樣做,聖經也是這樣記載,.
如果我們看回聖經其他的章節,.
一些人的生命裡面,他們這些瘡疤仍然記載在聖經裡面,.
他的目的是什麼呢?.
目的就是告訴我們,人在很多人生的歷史裡面,.
實在有很多失敗的地方,.
我們不能夠抹去的,.
你睡了覺,或者你做一些你覺得很開心的事,.
其實都抹不去,.
但聖經又不是叫我們去遺忘它,.
而是在我們的人生,在我們的成長裡面,.
給我們一個警惕,.
如果不悔改,不回轉,.
不離開罪,就沒有更新的機會..
第二就是儘管以色列人犯了很重的罪,.
摩西仍然很迫切,.
苦服在神的面前,.
來禁食禱告,.
在人生,在他們的歷程裡面,最黑暗的時候,.
上帝,這位施恩慈愛的神,.
仍然開一條出路,摩西知道,.
摩西就要仰望神,.
求神來寬恕他們..
第三,揭開以色列人這個歷史的瘡疤,.
我相信最重點就是,.
其實是要讓新一代的以色列人,.
能夠進入迦南地,.

$^{161}$不要重蹈覆轍,.
腐敗這種罪惡..
如果我們再看二十一節的時候,.
剛才我們沒有讀的經文,.
摩西去回顧的時候,.
看到神的憐憫,.
他馬上有行動,.
摩西在代求之後,.
他的行動,我相信是非常重要..
我把那叫你們犯罪所鑄的牛毒,.
用火焚燒,.
又倒碎磨得很細,.
以至細如灰塵,.
我就把這灰塵殺在從山上流下來的溪水中..
摩西來到禁食禱告,.
但他知道禁食禱告,.
他一定要這樣做,.
但上帝是怎樣,.
是否要滅絕,.
這個摩西是管不到,.
但他要做的,.
就是要把金牛族,.
這個腫瘤,.
是要徹底的清除,.
他要把他磨到成粉末,.
然後要由水沖走,.
大家有沒有印象,.
金牛毒是怎樣做的?.
如果你們留意,.
這道聖經第九章沒有說,.
原來當時候他們離開埃及的時候,.
神都很好,.
你叫埃及人,.
其實都不是埃及,.
因為當時埃及發生很嚴重的事件,.
就叫他們快點走,.
於是把埃及人的金器,銀器,.
所有東西,.
以色列人都拿走,.
原來金牛族就是用這些物料來製造,.

$^{201}$摩西要將它完全的粉碎,.
完全的剷除,.
就要告訴新一代的以色列人知道,.
不可以再有這些物料在他們的身上,.
如果有留意的話,.
自這件事之後,.
以色列人他們的身上,.
是不會有這些金飾,.
一些銀器,.
帶在身上的,.
弟兄姊妹,.
不如我們停一停想一想,.
想一想大家如果是有的,.
一些罪的蒼巴,.
又未向上帝揭開,.
其實神正正在等候著,.
你去到上帝的面前,.
闖開你自己內心裡很深處的一些罪,.
很想向上帝來認,.
但是認罪的時候,.
有沒有將這些,.
來到徹底的,.
來到粉碎,.
見到連微塵都見不到呢?.
其實上帝要更新你,.
祂在等著你,.
最後一段的經文,.
第26至29節,.
其實這一段就是摩西的土文,.
就是接著第20節那裡,.
這段的土文,.
經文是這樣說,.
我因耶和華說,.
要滅絕你們,.
就在耶和華面前照舊,.
苦服四十晝夜,.
我禱告耶和華說,.
主耶和華,.
求你不要滅絕你的百姓,.
他們是你的產業,.

$^{241}$是你用大力救贖的,.
用大能從埃及領出來的,.
求你紀念你的僕人,.
亞伯拉罕,.
以撒,.
雅各,.
不要想念這百姓的環境,.
邪惡,.
罪過,.
免得你領我們出來的那地之人說,.
耶和華說,.
因為不能將這百姓領進他所應許之地,.
又恨他們,.
所以領他們出去,.
要再曠野殺他們,.
其實他們是你的百姓,.
你的產業,.
是你用大能和新出來的傍臂,.
領出來的..
摩西實實在在地知道,.
神是有恩典有憐憫,.
因為他曾經經歷過,.
也見證著神對他的信實,.
所以他向新一代以色列人回顧這一段,.
他們先輩的歷史蒼巴的時候,.
他也再一次複述他的禱告,.
以色列人得罪耶和華,.
當時正值摩西領受上帝給予以色列人的約,.
正正山下以色列人,.
違反了最重要的,.
不可有別的神,.
也不可有偶像,.
所以摩西在這段禱告再一次重溫的時候,.
他要求神不要滅絕這一班百姓,.
沒錯,他們是罪有應得,.
但是他求神不要想念這些百姓的環境,.
邪惡,罪過,.
不要想念,.
神啊,你不如不要看他們這些蒼巴,.
反而要看上帝你自己,.

$^{281}$他要求神去看上帝自己的屬性,.
不知道大兄姐妹,我們祈禱,.
我們經常犯罪,.
不是經常犯罪,有意無意犯罪,.
那天突然上帝感動你,.
你會去祈禱認罪,.
但祈禱認罪會不會好像摩西的祈禱一樣,.
很徹底,看的不是看自己的罪,.
是看上帝屬性的罪,.
上帝的屬性,.
這裡有三方面,.
我們如何找到找住神的信實和應許,.
第一就是神的救贖,.
摩西在這裡說的禱告,.
就是讓上帝,.
神啊,你看看以色列這個民族,.
這一批的百姓,.
是你大能的神從埃及救贖出來的,.
他是你的產業,.
你不要滅絕他們,.
第二個就是,.
神啊,你要向這一班以色列人的列祖,.
曾經有應許有立約,.
神啊,求你紀念他們的先祖,.
你曾經和他們立約的,.
是要說,.
亞伯拉罕以撒雅各的後裔,.
是要人數眾多,.
你要實現,.
如果你滅絕了,.
你的應許就不會出現了,.
第三個理由,.
他向神來禱告,.
就是神的命和神的榮耀,.
如果以色列人被滅的時候,.
當時的列國,.
一定會誤解,.
神啊,你真是無能,.
以致神的命被褻瀆,.
甚至失去了以色列人見證神榮耀的機會,.

$^{321}$所以摩西在這篇討論中,.
向第二代百姓鄭重宣告,.
他們以色列先祖的身份,.
一定要記住,.
他們是神救贖的產業,.
神不會棄絕自己親手所建立的關係,.
摩西告訴新一代的以色列人知道,.
這表明了神是信實到底,.
不改變的神,.
所以讓我們確實抓住神的屬性,.
神是守約並且思慈愛的神,.
在我們往後的日子裡,.
我們實在不應該再找到任何藉口,.
再次重複褫瀆神,.
羞辱神的命,.
當然大家可以說,.
會不會有更多的深入了解這段禱告呢?.
如果要用很深的神學理解,.
我相信在這個時刻,.
不是這個做法,.
其實這段聖經,.
這段禱文,.
摩西的教導提醒,.
是要我們每一個,.
用單純的信靠的信心,.
和持久謙卑的態度,.
來到神私恩座面前禱告,.
奧古斯丁,.
我想大家都聽過,.
是早期一位很出名的教父,.
如果大家有看過他的人生歷史,.
我相信都會覺得很臭,.
他小時候經歷認識神的人,.
但一路在他成長裡,.
他離開了上帝,.
他撥逆,.
他做了很多在他人生裡,.
不單是信仰迷失,.
甚至是撥逆神,.
人生搞到很多千瘡百孔的污穢,.

$^{361}$但他母親就是這樣,.
不離不棄,.
雖然看到自己的兒子奧古斯丁,.
人生不堪入目的各樣的事,.
但他就是這樣,.
用單純信靠神的心,.
持續地不斷地,.
為他的兒子禱告,.
他兒子奧古斯丁,.
其實都是常常犯罪,.
如果他母親不斷地,.
看到兒子一件件所犯的事,.
來禱告,.
我相信他一定心灰意冷,.
但神讓奧古斯丁的母親看到,.
神那種憐憫,.
神那種信實,.
他一直禱告,.
十多年後,.
奧古斯丁就遇見他的恩師,.
他信仰生命大逆轉,.
悔改回到上帝那裡,.
他有一本著名的著作,.
懺悔錄,.
如果弟兄姊妹有看的時候,.
我相信不單單是他生命的寫照,.
甚至可能會反照我們的生命,.
我們的生命其實都會有很多難,.
很多千瘡百孔的地方,.
但原來神有能力去赦免我們,.
縱然我們會滅絕,.
但神的心意不是要滅絕我們,.
是要我們繼續向著神給我們的應許之地前行,.
弟兄姊妹,.
今天神就是藉著耶穌基督的救贖,.
我們就有他寶貴兒子身份的名份,.
所以從歷史我們看到教訓,.
就是不要重複以色列在曠野裡的失敗,.
我們就要學習在歷史的瘡疤裡脫身而出,.
我們的更新不是在於我們是不是完美無罪,.

$^{401}$而是在於我們要持續悔改,.
找著基督給我們的恩典,.
基督救贖的恩典,.
而更加要聚焦在神的信實榮耀和他的應許,.
我們要再次出發,朝著神應許的未來,.
我又說回我的歷史的臭屎,.
偷完錢,結果是怎樣呢?.
原來越早脫離偷錢,.
因為除了打之外,我爸都會訓斥我,.
他說你現在小時候只會偷錢,.
長大後會偷什麼?.
長大後當然會偷金,.
所以有理無理就去了寄宿學校,.
但原來在學校,剛才也有分享過,.
是可以讀到聖經故事的,.
當小學畢業升到中學的時候,.
在天主教學校讀,.
但我又不是很習慣禮儀,.
雖然進去是這樣的,.
當時升到中學,.
爸爸在山寨廠轉型,.
轉做車皮夾手袋的工場,.
但爸爸仍然和他的行家有聯絡,.
有一次,.
原來我聽過世叔和我爸說,.
不如帶你兒子回教會吧,.
那時候回教會也很特別,.
在某一次的星期日,.
帶我們去喝茶,.
平時去喝茶是去元朗,.
因為元朗是白鄉,.
那次他載我去上水喝茶,.
喝完茶後載我去教會,.
就是那間粉嶺教會,.
很奇怪,.
當我這樣想的時候,.
原來上帝是不會放過人生命的那種蒼巴,.
他沒有去擊殺他,.
反而要帶他去教會,.
讓他在教會裡聚會讀聖經,.

$^{441}$當我在教會讀聖經,.
我又想起小學讀的聖經故事,.
有多適合,.
於是就留下在那裡,.
在教會裡成長,.
當然不會每個星期我爸爸載我回教會,.
喝完茶,.
要你自己去做,.
原來上帝就是這麼奇妙,.
透過我們人生裡的軟弱,.
一些不為人知的蒼巴,.
上帝就會感動其他人,.
他會帶這個人去上帝那裡,.
我去預備經文的時候,.
去想起原來我回教會,.
到現在數一數都半個世紀,.
我會看神是信實,.
他沒有從我在偷錢的污點裡不理我,.
反而用不同的方法來接觸,.
讓我去經歷上帝是一個怎樣的上帝,.
今天,弟兄姊妹,.
你生命裡的一些蒼巴,.
罪的痕跡,.
你沒有願意,.
或者和弟兄姊妹分享,.
但是上帝有計劃,.
有他的計劃能夠轉化,.
甚至將你生命裡的金牛毒,.
這些蒼巴,.
他很希望來到磨成灰燼,.
而轉移到上帝那裡,.
就讓你的生命來去改變,.
去見到神的信實,.
我們一起來低頭禱告,.
我們確確實實,.
人生裡有很多很多一些罪的蒼巴,.
我們甚至不敢碰,不敢提,不敢碰,.
大主啊,你仍然在我們的生命裡,.
很想我們親自的,.
將這些揭示出來,揭開出來,.

$^{481}$去到你的面前,.
求你來幫助我們,.
讓我們承認我們的軟弱,.
承認我們的過犯,.
主啊,我們在罪中,.
你必定的會滅絕我們,.
但是在你的恩典,.
在你的慈愛裡面,.
主啊,你很想我們來到有禱告,.
有認罪的時候,.
我們就要轉念,.
看到上帝你自己的信實,.
和你自己救贖,.
和你自己榮耀的名,.
你是施恩憐憫,.
大有慈愛的主,.
願你來赦免我們,.
讓我們重新出發,.
能夠行走上帝你為我們預備的,.
一個美好的道路,.
誰聽我們的禱告,.
奉主耶穌基督的名,.
阿們..
\newpage



\section{約翰福音 21:20-23}
\label{sec:LKA6QvIWId8}
\textbf{2021年4月25日崇拜直播｜陳明泉牧師｜與你何干?｜約翰福音21\_20-23}
\newline
\newline
連結: \href{https://youtube.com/watch?v=LKA6QvIWId8}{\texttt{https://youtube.com/watch?v=LKA6QvIWId8}} ~~~~ 語音日期: 2025-03-18
\newline
\newline
\hyperref[sec:P0a0tJAFKd4]{\small{< < < PREV SERMON < < <}}
~
\hyperref[sec:index_chronic]{\small{[返順時目]}}
~
\hyperref[sec:aQhmLPtYC40]{\small{> > > NEXT SERMON > > >}}
\newline
\newline
約翰福音 21:20-23
\newline
\begin{longtable}{cl}
\hline
\hline
章節 & 經文 (和合本修訂版)\\
\hline
21:20 & \begin{tabularx}{0.7\textwidth}{X} 彼得轉過身來，看見耶穌所愛的那門徒跟著，就是在晚餐時靠著耶穌胸膛說「主啊，出賣你的是誰」的那門徒。 \end{tabularx} \\ \\ \relax
21:21 & \begin{tabularx}{0.7\textwidth}{X} 彼得看見他，就問耶穌：「主啊，這個人怎樣呢？」 \end{tabularx} \\ \\ \relax
21:22 & \begin{tabularx}{0.7\textwidth}{X} 耶穌對他說：「假如我要他等到我來的時候還在，跟你有甚麼關係呢？你跟從我吧！」 \end{tabularx} \\ \\ \relax
21:23 & \begin{tabularx}{0.7\textwidth}{X} 於是這話在弟兄中間流傳，說那門徒不死。其實，耶穌不是說他不死，而是對彼得說：「假如我要他等到我來的時候還在，跟你有甚麼關係呢？」 \end{tabularx} \\ \\
[1ex]
\hline
\hline
\end{longtable}
$^{1}$讓我們繼續敬拜這位主.
以上帝的說話幫助我們.
繼續敬拜.
經份是約翰福音21章20至23節.
在大家的教會聖經是160頁.
約翰福音第21章20至23節.
彼得轉過來.
看見耶穌所愛的那門徒跟著.
就是在晚飯的時候.
靠著耶穌胸膛說.
主啊,賣你的是誰?的那門徒.
彼得看見他,就問耶穌說.
主啊,這人將來如何?.
耶穌對他說.
我若要他等到我來的時候,與你何干?.
你跟從我吧.
於是者說,存在弟兄中間.
說那門徒不死.
其實,耶穌不是說他不死,乃是說.
我若要他等到我來的時候,與你何干?.
今天宣講的題目是「與你何干?」.
由旺角浸信會主任牧師陳明全牧師宣講.
我們把時間交給陳牧師.
弟妹平安.
我們今天來到教會.
一起來敬拜,一起來聆聽他的話語.
有一個虛構的故事是這樣說的.
Robert仔剛剛談戀愛.
Robert仔夢寐以求有一個夢幻的約會.
跟他剛剛進入愛河的女朋友.
很想哄她開心.
於是Robert仔存了很多錢.
就去香港某一間很知名的酒店.
跟他女朋友共進晚餐.
但存夠錢的時候,他發現一件事.
其實他長這麼大,都沒有去過這麼華貴的酒店吃晚餐.
當他想起這件事的時候,又怕自己失禮.
存了這麼多錢.
於是Robert仔問他的朋友.
去到這麼華貴的酒店要怎樣做才不失禮人.

$^{41}$又令女朋友開心呢?.
這個朋友就給他一個辦法.
他說你不用怕.
總之去到吃飯的時候.
旁邊的客人做什麼,你就跟著做什麼就行了.
他點什麼餐,侍應都不會理你.
你做什麼,就跟著旁邊的客人做就行了.
Robert仔說是很巧.
於是他按照這個朋友的建議.
那天晚上就收拾得很整齊.
西裝不挺地跟女朋友一起來共進晚餐.
旁邊的客人點A餐.
應該不會點A餐,茶餐廳都會點A餐.
總之點餐,聽到就知道.
他點了這個餐,我就跟著點這個.
要什麼餐湯.
旁邊的客人點了什麼湯,我就點什麼湯就行了.
然後整件事就很平順.
餐具怎麼放,他就看著.
餐具有分次序的,他就看著別人這樣做.
女朋友看到Robert仔.
挺懂事的,餐桌禮儀一切都很有風度.
覺得是很難忘的一晚.
這一晚Robert仔吃飯的過程都覺得很棒.
結帳,旁邊的客人也結帳.
就想著這次過關了.
Robert仔很開心.
結帳之後離開.
這間酒店有一條很漂亮的螺旋樓梯.
通常婚禮,看到公主下來的樓梯.
結帳就一直下去.
旁邊的客人又跟著女朋友.
一直很雍容華貴地走這條螺旋樓梯.
誰知道旁邊的男人就滑了一腳.
就這樣蹬了幾級.
下樓時摸了一下屁股.
Robert仔看到整件事.
死了,難道這件事也要學嗎?.
結果是怎樣呢?大家想一想.
搞這個爛gag,是一個道讚出來的故事.

$^{81}$不過今天我們人生很多時候都是這樣.
我們人生不知道自己的路要怎樣走.
或者有時候我們跟隨基督耶穌.
信了主之後.
我們生命的路應該要怎樣走呢?.
我們通常最近的參考就是看旁邊的那個.
看大我信主的基督徒的生命是怎樣.
或者看我的祖長的榜樣是怎樣.
我們跟著他來做.
信仰裡面一定要有這些響道.
我們經常都鼓勵的.
不過原來如果我們純粹想複製了.
我身邊大我信主的響道.
或者看著他來過我們信仰生活的時候.
很可能我們就有類似Robert仔的感覺.
究竟他會不會都要這樣走自己的路.
今天這段經文是講到彼得和耶穌之間.
一個很親密的對話.
在這段經文裡面.
他特別交代了耶穌基督對這個門徒的呼召.
剛才我們一起讀的那段經文我們知道.
在約翰福音第21章裡面.
耶穌特別講到他親自呼籲這個門徒.
他未來要走一條怎樣的路.
其他人怎樣走和他有什麼關係呢?.
耶穌要親自呼召這個門徒.
走他自己所走的路.
今天這篇經文內容可能涵蓋範圍.
比剛才我們看的那些經文更闊.
特別是我們針對在第21章約翰福音裡面.
所講的內容更闊.
希望今天我們仔細看這段經文的時候.
今天我們可以看見旁邊的那個之外.
也看見上帝在我們生命裡面怎樣引我們走前的路.
我們一起祈禱吧.
求聖靈幫助我們讓我們明白這段經文.
我們同心祈禱.
阿爸天父願你在我們當中.
順領我們.
求你對我們說話.

$^{121}$藉著今天這段經文.
你寶貴的教導.
讓我們藏在心裡面.
幫助孫說的說得清楚.
幫助聆聽的每位弟兄姊妹聽得明白.
我們恭敬.
虛心地領受你的話語.
願你的聖靈親自與我們同在.
我們多謝主隨天祈禱.
奉基督耶穌得勝的聖名祈求.
Amen.
其實如果你仔細去看約翰福音.
你會發現的.
其實在第20章.
第30至31節.
其實某程度上來說.
約翰福音已經完結了.
特別是說到耶穌基督復活.
而且顯現了給門徒知道.
在月頭的時候和大家說.
多瑪的懷疑.
到耶穌向多瑪顯現.
其實耶穌的復活已經是整個約翰福音裡面的高潮.
而且好像一個結尾一樣.
所以你讀上去你會知道.
約翰也是這樣結尾的.
在20章第30至31節.
他這樣寫的.
耶穌在門徒面前另外行了許多神蹟.
沒有記載在書上.
但記者則是要叫你們信耶穌是基督.
是神的兒子.
並且叫你們信了他.
就可以因他的名得生命.
那你見到了.
是一些總結的說話.
原則上約翰福音可以停在這裡.
不過在福音書再流傳的時間裡.
大概是耶穌主後70至90年這一段時間裡.
福音書完成了一個大浪的初稿.

$^{161}$但面對主後70至90年這一段日子.
其實耶穌所愛的門徒已經殉道.
為了信仰的緣故.
已經相傳說他倒吊在十字架上.
又有些傳言說到.
這個唯一存留在世上的使徒.
寫了約翰福音的使徒.
他不會死的.
有很多的流傳.
所以很多的色經學家.
大家都同意.
這個21章是對於整個約翰福音裡裡面.
一個補充的筆.
或者是一個後記.
或者是今天電影裡面說的彩蛋.
不是為了貪過癮.
或者還有些東西想補充一下.
而是要針對當時教會的一些流傳.
以為說使徒約翰不會死.
又特別紀念到殉道.
舊時的門徒的大哥彼得.
所以他特別記載了第21章的內容.
這個內容特別的除了要澄清.
約翰不是不會死.
只不過是給其他門徒.
存留在世上會長久一些的日子.
他更加要說到在聖靈降臨之前.
和耶穌復活之間.
耶穌顯現給門徒的時候.
他有些特別呼籲給一班門徒.
大概你用這個框架來看.
其實是一個補充的內文.
內容也是使徒約翰所寫的.
不過他會帶著一些焦點.
他會說到彼得如何被呼召.
也會說到人的生命究竟要怎樣.
來繼續為主努力.
當我們知道這個時候.
經文第20章1至14節就設定了一個場景.
它也是一個很完整的技術.

$^{201}$這個技術是說到耶穌第三次向門徒顯現.
在復活之後第三次.
聖靈有未降臨.
在第三次的顯現裡面.
在經文裡面說到一班門徒集結.
一起在加利利海或者提比利亞湖.
在那裡等待.
在那裡吃東西.
在那裡聚集.
一班門徒就下海去打魚.
彼得為首的一班漁夫.
打了一整晚魚.
也沒有任何的漁獲.
於是在失望之際.
就聆聽到離岸有聲音傳過來了.
喂!小伙子!吃了早餐了嗎?.
嘩!聲音這麼有趣.
還沒吃早餐的就在船右邊下網吧.
那裡有很多魚.
這一班漁夫.
不過打了一整晚都沒有魚.
在失望或者無癮之際.
他不知道那聲音來自哪裡.
不過覺得挺熟悉.
又這麼說.
於是在船的右邊就下網了.
一下網就撈了很多的魚獲.
153條魚.
153條沒什麼特別的意思.
不要過分詮釋.
不止是153條很多的魚獲.
就拉上來.
當拉上來的時候.
使徒約翰就知道.
那個說的聲音是耶穌來的.
於是彼得就赤著自己的身子.
脫下衣服馬上衝下海.
游到耶穌那裡.
一路游的時候.
原來在岸上.

$^{241}$耶穌已經在等待門徒.
耶穌已經預備了炭火.
預備了幾個餅.
然後叫門徒拿幾條魚來.
一路燒魚,吃餅.
由一個早餐約會.
來展開第21章的內容.
在吃完早餐之後.
也是我們今天要講的主要主題.
第15至17節.
他們吃過了早飯.
耶穌對西門彼得說.
約翰的兒子西門.
你愛我比這些更深嗎?.
彼得說:主啊!是的!你知道我愛你.
耶穌對他說:你餵養我的小羊.
耶穌第二次又對他說.
約翰的兒子西門.
你愛我嗎?.
彼得說:主啊!是的!你知道我愛你.
耶穌說:你牧養我的羊.
第三次他說.
約翰的兒子西門.
你愛我嗎?.
彼得因為耶穌第三次對他說.
你愛我嗎?.
就憂愁對耶穌說.
主啊!你是無所不知的.
你知道我愛你.
耶穌說:你餵養我的羊.
在第15至17節.
當門徒一起吃完早餐.
耶穌將焦點放在彼得的對話當中.
約翰特別記載.
大概不用想其他門徒是否聽著這段話.
都在一起.
不過焦點要記述的不是閒雜的對話.
耶穌在這裡特別的焦點放在.
耶穌與彼得之間的那份情.
耶穌在這段對話裡三次問彼得.

$^{281}$你愛我嗎?.
或者你愛我比這些更深.
這些代表什麼呢?.
第一次問彼得.
彼得愛這些東西比其他更深.
這些有什麼意涵呢?.
在這裡其實很模糊.
不過離不開這三個說法.
第一個.
彼得愛這些打魚工具.
或者打魚的生活.
愛這件事比愛耶穌更深.
這是第一個說法.
所以如果說這個說法.
就會很批判彼得在河邊回到加利利打魚.
他覺得松草固葉撇棄了耶穌.
這是第一個詮釋.
第二個詮釋就是.
這些不是指那些物件.
而是耶穌要彼得與這班門徒做比較.
你愛我比這些門徒更深.
這是對象.
這個解法是指.
耶穌既然要彼得做大佬.
那麼你做大佬會不會愛主比其他門徒更深呢?.
這是第二個說法.
不過這兩個說法並不是無懈可擊.
因為當你看到那些魚具.
其實耶穌並不是批判這班門徒在打魚.
事實上這班門徒也要吃的.
他們自己下去打魚並不等於松草固葉.
只不過他們基於生活生存.
他們也要生活.
所以他們就去海裡打魚.
第二個.
其實耶穌在後面說.
每個門徒都要走自己的路.
他不會回到一個狀態.
就是叫彼得與其他門徒之間比較.
誰愛主更深.

$^{321}$這不是耶穌一貫對門徒的教導.
反而是勉勵門徒盡心盡性盡力愛主.
每個門徒都盡自己的力量.
而這個盡力不是人與人之間的互相較勁和比較.
所以第三個可能會比較好.
第三個意思是.
這些是在說世上所有的一切.
彼得愛世上的一切多一些.
還是愛主來的多一些.
這些其實是一個虛範的詞語.
不過在這裡訴諸於問彼得心裡.
究竟填滿在彼得內心當中.
最重要的是什麼.
說愛有時候我們對神的愛.
如果在聖經裡讀我們會覺得很習慣.
不過你抽身來說.
你看到兩個大男人.
或者耶穌成為一個男人.
問你愛我媽.
你想想這個場景其實挺奇怪的.
我們甚至乎中國人的文化裡.
我們不會私逃關係.
或者神與人之間的關係.
我們不會用愛來表達.
不過在聖經裡面說到人與神之間.
或者神人之間的關係和維繫.
就不是我們文化裡可以完全媲美的.
甚至說到神對人的愛.
就是一種犧牲自己獨生兒子的愛.
人亦希望用他對生命的一些很重要的對象.
那種擺上犧牲.
來表達人對神的那種尊崇和深切的愛.
這種愛不是純粹一些愛在規條.
或者訓練出來的.
而是一個發自內心來將那種心意表達.
所以你說愛其實是一種怎樣的關係呢.
愛其實就是那樣東西.
那個對象會不會成為你生命的全部.
你愛上帝.
上帝就會成為你生命的全部.

$^{361}$當有人說上帝的褻瀆.
或者污衊上帝的名.
你會很生氣的.
因為當你生命的全部被人詆毀的時候.
你會有一種反抗.
但當你所愛的對象憂傷的時候.
你都會和他同步來憂傷.
神人關係.
或者耶穌要問彼得就是訴諸一種這樣的關係.
耶穌自己.
或者上帝在彼得心目中裡面.
你有沒有一種這樣的關係呢.
你愛我比你生命的一切會不會更加重要.
上帝的愛會不會佔據你整個心懷的意念呢.
第一次彼得回答你知道我愛你.
第二次彼得回答你知道我愛你.
同樣的問題問兩次.
大概他都是同樣的答案.
可能都是真誠的.
但到第三次再問的時候.
有些解經瓜說是不是要回應彼得三次不認主.
用這裡有三次那裡有三次來對應.
如果人生這麼簡單就好了.
問三次就可以解決生命裡面很大的罪疚.
約翰福音特別說.
第三次的問是令彼得變得憂愁.
不是憂愁是來自他自己過去.
或者那三次不認主對應那個內容.
而是這個夫子不滿意彼得自己提供的答案.
事實上你看彼得回答.
太太通常都會問丈夫.
很少老公問太太.
老婆問老公你愛我嗎.
老公通常怎麼回答.
當然愛.
尤其是睡覺的時候才問.
臨睡前問老公都想睡.
你愛不愛我.
愛.
是不是真的愛我.

$^{401}$是.
愛不愛我.
不要這麼吵.
通常男人是這樣的.
不耐煩或者有種感覺.
同樣問題問三次.
不過如果你細心去想.
大概你知道你所說的答案.
是不能夠滿足了發問的人的心.
但彼得的憂愁是覺得.
耶穌問三次是不是代表著他的答案.
不是一個完美的答案.
不過在經文裡說到彼得.
事實上你回顧他整個跟隨耶穌的歷史.
如果在人之常情.
你愛我比這些更深.
很直觀的答法.
你會用一些東西來證明自己.
我很愛你.
我做了什麼.
我為你怎樣去做.
這樣來證明自己愛神.
我們有時候也會這樣.
向上帝祈禱.
不過在這裡彼得不是用這樣的答案.
因為事實上在彼得的經歷裡.
即使他以前怎樣去跟隨耶穌.
但他經歷過一個很重要的失敗.
就是曾經否認主.
在他的人生裡.
他還有什麼業績可以誇.
還有什麼誇口他覺得.
他說他很愛主.
他根本沒有辦法證明自己.
或者在他的經歷裡.
告訴耶穌他是一個很深愛主的人.
不過他不抽空一件事.
他只是整體的來講.
耶穌你知道我真的很愛你.
耶穌你知道.

$^{441}$我知道你在我生命裡.
一個很重要的.
是我的全部.
他用一個整體的來講.
將自己的心來表達出來.
所以即使去到第三次.
耶穌去問彼得的時候.
他說耶穌你是無所不知的.
你知道我愛你.
彼得大概沒有辦法訴諸於自己的經驗.
再去證明自己是一個很完美的門徒.
他生命裡大概都充滿了很多的遺憾.
在這個使徒關係當中.
在這個救主,恩主面前.
他其實有很多事情.
不是能夠站立的.
在這裡我稍微思想一下.
弟兄姊妹.
如果耶穌問你.
你愛我比姐姐更深嗎?.
你愛我嗎?.
你會怎樣回答呢?.
在我預備聖經的時候.
我也想像如果耶穌問我這件事.
大概我也是沒有辦法.
像琅琅上口說到我自己.
證明自己多麼愛上帝.
生命裡有大量的遺憾.
即使作為一個牧者,牧師.
我都不是一個完美的人.
生命裡有很多做錯了的事.
有很多令上帝甚至可以責備的事.
有很多事情力有不逮.
有心無力.
如果按照我自己的本相.
如果上帝這樣問我.
我說不上.
我是一個完美的基督徒.
可能我看我的童宮更完美.
我看我自己破破爛爛的.

$^{481}$我大概和彼得一樣.
上帝,我沒有什麼可誇張的.
但我知道,你很明白我.
我的心是真誠的.
你無所不知.
我整體而言.
上帝,我是很愛你的.
今天這段經文.
似乎為我們很多弟兄姊妹.
帶來一些.
要多想一點的去問自己.
今天我們當然.
除了我自己之外.
如果真誠地面對自己.
其實我們每一個都為很多失敗.
很多不理想的地方.
甚至我們力數加增.
怎樣數都好.
數不上是一個完美的基督徒.
不過,我們的心是真誠的.
我們知道上帝知道我們的人.
不是真的有心.
或不是很願意走在罪惡離棄上帝的道路裡.
不過我們人生有時真的不知道為何會這樣走.
上帝,你知道,我愛你.
當彼得回答耶穌之後.
三次裡面,耶穌都是用同樣的話來回應他.
我們知道,其實耶穌不是不滿意他的答案.
耶穌在答第一次.
你愛我媽.
然後他就說,你餵養我的羊.
第二次,同樣都是.
你餵養我的羊.
第三次,約翰兒子西門,你愛我媽.
耶穌說,你餵養我的羊.
牧羊,餵養,小羊或羊.
大概都是那些字眼.
約翰都是用參雜的方式來用.
你不用特別細分.
因為在約翰的寫作特色裡.

$^{521}$包括這個愛字,他夾雜著來用.
大概寫東西不用寫得那麼參.
用不同的字眼來表達同樣的意思.
不過在這裡,如果耶穌對這個門徒所答的答案不滿意.
大概他給的答案不會是一個呼召.
但在這裡,耶穌給的答案是一個呼召.
你牧羊,你餵養上帝或耶穌基督的羊.
當彼得回答完,你知道我愛你.
你就餵養我的羊.
知道自己對神的愛.
上帝就叫他化成一種行動.
餵養他的羊.
這個說話對於彼得來說是很深的意思.
如果你有機會回去再看.
在彼得前書第五章第二節.
彼得也是用同樣的意思來勉勵教會的長老.
彼得前書第五章第二節裡面說要牧養群羊.
按著主的心來牧養.
不是出於勉強是甘心的.
不是出於貪財而是樂意的.
是要付錢,付力,付心機來牧養這些人.
這些小羊你不要克制他.
而是你要成為群羊的榜樣.
彼得在他將要殉道之前寫彼得前書.
其實也是一直在回應耶穌基督.
在這段經文裡面對他的呼召.
你愛我,你就好好牧養這些小羊.
幫助他們,甘心樂意地帶領他們.
讓他們成長.
如果你要表達你的愛.
就在牧養這些羊身上.
來將對上帝的愛表現出來.
今天我們說需要很多愛.
我們的動力,第一件事我們要明白.
不是那個人很可愛,所以你很愛他.
今天我們對身邊的社會弟兄姊妹要愛.
其實是基於我們對神的愛.
我們不完美,但我們知道.
既然我們帶著破破爛爛的生命.
去到上帝面前表達.

$^{561}$上帝我真的很愛你.
上帝就要叫我們.
將這種愛要慰給身邊的小羊.
那些可能比你遲一些認識主的朋友.
讓他們得著福音的好處.
這是第一點我們要知道.
彼得要重建彼得的.
需要重建彼得的愛心.
呼召,跟從.
第一個最重要的起始點就是愛.
今天我們所做的.
是基於愛上帝.
還是基於愛自己呢?.
我們想自己感覺良好.
還是我們想呼召上帝.
感到滿足.
這是今天第一點.
我們從這段經文看到的.
在第十八至十九節.
這是第二點.
我想跟大家仔細看.
當耶穌呼召彼得之後.
耶穌繼續引申.
再說:我實實在在地告訴你.
你年少的時候自己束上帶子.
隨意往來.
但年老的時候你要伸出手來.
別人要把你束上帶你到不願意去的地方.
十九節.
耶穌說者說是指著彼得要怎樣死榮耀神.
說了者說就對他說.
你跟從我吧.
在這裡.
耶穌繼續說.
要將餵養羊之後.
他的人生要怎樣呢?.
某程度來說.
耶穌說了一個寓言.
與其說寓言.
其實是在說上帝的計劃在彼得身上.

$^{601}$彼得身上的計劃是怎樣呢?.
年輕的時候.
隨意往來.
束上帶子.
就是說令自己行動方便一點.
進入一種戰鬥.
或者一種行動很靈巧的.
隨意往來的一種狀態.
不過到了年老的時候.
那條帶子不是束上.
令自己更加靈活.
那條帶子是變成束上他的手.
讓他去到他不願意去的地方.
束上這個字在後面的解法.
有些人的說法就是.
戴上十字架.
手環.
束上就是同樣用這個字.
所以在這裡.
耶穌很隱含地說.
到了他年老的日子.
他將要像束上十字架一樣.
走上一條被約束.
甚至是殉道的道路.
約翰怕大家不明白.
怕解錯.
或者耶穌說的這個.
好像很含糊的一個寓言.
約翰第十九節補充說.
其實這句話是說.
彼得將來要死.
直著殉道.
來榮耀上帝.
不過耶穌對著彼得.
其實說這番話.
彼得大概不是很清楚.
自己將要這樣走這條路.
不過他稍微將一個.
上帝的命定在這個.
大門徒大弟子身上.

$^{641}$他講了出來.
束上帶子.
隨意往來.
到年老的時候.
被人束上.
帶到不願意去的地方.
年輕的時候束上帶子.
代表著他自由奔放.
他覺得自己有很多的可能性.
在信仰裡面.
他有很多的發揮.
不過這個不是等於.
有這個東西就有那個結局.
這兩者沒有因果關係.
不過後者是說到.
他將來的命定.
在基督裡面的命定.
就是說到.
他將來要以殉道.
來榮耀上帝.
在整個對彼得的呼召裡面.
其實都是在說.
他正在走的路.
等於耶穌正在走的路.
耶穌如何走上十字架.
走上去一個榮耀的道路.
同樣作為他的門徒.
都要跟隨他的腳踏.
來到這樣走.
不過弟兄姊妹.
你不要想著.
每個人都是這樣的.
在聖經裡面說.
以死來到榮耀神.
不是每個人都有這份資格.
或者不是每個人都是一條.
這樣的路來做的.
不過彼得在這裡.
特別接受耶穌.
很清晰的計劃告訴他.

$^{681}$將來的日子他要受苦的.
他要面對他很不願意承受的那件事.
不過整個過程.
整個人生.
是以基督耶穌的命途.
來定義他今天.
這個人.
他要未來所走的每一個步履的時候.
都是以耶穌基督走過的步履.
來跟隨他.
耶穌怎樣走.
他就跟著怎樣走.
有些人走的方式.
和耶穌一模一樣.
有些人可能跟得沒那麼貼近.
但都不離開.
在耶穌的足印裡面.
來走面前的路.
今天我們.
我不知道你思考你自己的人生.
你會用什麼來計劃你自己的人生.
不過.
作為一個認真.
尋索聖經的基督徒.
大概我們都要不停問上帝.
上帝你想我生命怎樣走.
不過今天我們不習慣這樣思考.
我們通常習慣思考什麼呢?.
思考哪裡有機遇.
哪裡感覺良好.
哪裡少一點.
少癮在工作裡.
哪裡找個好一點的環境.
哪裡賺得多一點錢等等.
我們很多時候用這些角度.
來思考我們自己人生的將來.
在這段經文裡面.
如果跟隨主.
耶穌說你跟從我吧.
大概祂想調教我們的焦點.

$^{721}$用耶穌想我們怎樣走的路.
來走面前人生的路.
以耶穌想我們怎樣做的.
我們就去怎樣做.
更加重要的是.
與其說剛才我說的.
耶穌預言給我們的將來.
與其說預言.
倒不如說這是命定.
上帝定了你是這樣的.
你的人生不是你任意想怎樣就怎樣.
你的人生.
其實上帝已經為你鋪陳一些軌跡.
其實你都是跟著上帝的軌跡.
來走在我們的人生道路當中.
昔日有些不起眼的經歷.
到最後.
原來這些經歷會成為你人生裡面.
一個祝福或是你要奮鬥的目標.
這些一切都是成為我們人生的.
很重要的找回我們人生的足跡.
或是未來所走的每一步.
這些是印證.
曾經有一個小小負面的例子.
有一個傳道童公.
我們自己也會說自己的夢召經歷.
有一個傳道童公.
我問他為何他會讀神學.
為何他會做一個牧者.
這個童公說.
我的夢召很特別.
其實那次我在培靈賢經大會裡面.
聽完培靈賢經大會的講道之後.
我決志要跟隨主.
為主服侍.
我說我也是這樣.
他說不過我跟你看的東西有點不同.
我問他看到什麼.
當天我不太記得那個訊息說什麼.
但很記得講完在台上呼召的義氣風法.

$^{761}$我已經有點無本動地說這些話.
他說然後怎樣.
你看見義氣風法.
我覺得將來上帝要我像台上的牧者.
號令武林會.
然後一叫就振臂一呼.
所有人就要悔改歸主.
迎合上帝的計劃旨意.
就要這樣來毀身.
我聽完之後就覺得.
打了冷顫.
有些人這樣說話也有.
我說你真的這樣相信嗎.
他說是的.
當我閉上眼的時候.
我就覺得站在台上穿西裝的那個就是我.
不過我也很不客氣.
我說那個是呼召.
還是幻想出來你自己感覺良好的一個想像.
究竟是意象還是想像.
是你自己想做.
還是上帝叫你去做.
今天很多朋友很多弟兄姊妹.
我們自己想做的事.
我們找上帝過橋.
但裡面的內容其實大部分.
都是心裡面自己想出來的.
我們今天如果說真的呼召跟隨.
我們要放下我們心裡面很多不必要的想像.
不要覺得自己太過重要.
覺得自己很重要的人.
很難面對上帝的呼召.
因為到最後他也是以自我為中心.
來走自己想走的路.
求主幫助我們.
在這裡上帝命定了我們人生怎樣走.
你尋索一下你自己人生的足跡.
可能有些群體上帝一直放在你身邊.
那些人可能是你需要向他傳福音的對象.
有些經歷是你才獨有的.

$^{801}$可能這個經歷上帝要透過你生命的掙扎.
甚至乎痛苦.
你怎樣走出來.
用這些生命的經歷來幫助.
可能在未來的日子裡.
有相同經歷的人走過人生的路.
上帝在我們人生裡有很多的命定.
我們這些的命定.
我們知道了我們接受了它的時候.
我們可以放下那些不必要的幻想.
我們接受上帝的命定.
我們更加知道我們的人生是交付在一個友情.
和上帝為我們預備最好的上帝的手裡.
命定和宿命不同.
宿命是說一盤數.
或者是社會上的力量驅使你.
沒有情的.
不會說哪些對你好.
宿命論就是說那些.
但是命定是說我們的生命是交付在上帝的手裡.
他按著他的慈愛.
引領你.
帶領你走人生的路.
你將你人生.
放心地交付給這個主.
他不會給我們一條冤枉路.
給我們一條錯路走.
最後一個段落.
20至第23節.
彼得轉過來看見耶穌所愛的冷門徒跟著.
就在晚飯的時候靠著耶穌胸膛說.
主啊,賣你的是誰?的冷門徒.
彼得看見他就問耶穌說.
主啊,這人將來如何?.
彼得對他說.
我若要他等到我來的時候.
與你何干?.
你跟從我吧.
於是者說,傳到弟兄中間.
說冷門徒不死.

$^{841}$其實耶穌不是說他不死.
乃是說.
我若要他等到我來的時候.
與你何干?.
與你何干是一個很重要的字眼.
你覺得文巢巢不知道在說什麼呢?.
回到廣東話你就知道了.
關你什麼事?.
約翰將來如何?.
關彼得在這一刻跟隨主.
關他什麼事?.
其實耶穌就是這樣回答.
彼得的問題.
彼得的問題是為什麼無端端要這樣呢?.
大概還有一些劣根.
或者還有一些想法.
如果我將來要這樣榮耀主.
那我身邊的兄弟會怎樣?.
他們會不會很好的?.
那我豈不是很糟糕?.
可能是這樣的,是不是?.
又或者.
上帝有呼召臨到你身上.
上帝要你做一個怎樣的人.
未必一定要讀神學.
可能在你的工作崗位裡繼續榮耀神.
那他說.
我跟他的大團契.
他走的路這麼容易.
我的路這麼難走.
我這樣走不就很糟糕?.
上帝很不公道.
我們會忿忿不平.
經常都會問這些問題.
耶穌回答.
關你什麼事?.
他們的際遇是怎樣?.
跟你如何領受呼召.
兩者有什麼關係?.
不過我們今天很習慣.

$^{881}$將別人的際遇跟自己的際遇比較.
見到別人很順.
你自己就感覺不好.
但見到別人比你差.
你又覺得沾沾自喜.
我們人生不停地比較.
不過.
人之常情.
我們經常都比較.
我們的社會也是.
或者我們生活的空間.
經常都用比較來凸顯自己.
去到什麼位置.
不過.
去到耶穌的呼召.
去到人生如何跟隨主.
我就不可以用比較這種形式.
去知道自己的位置.
我們也不需要看別人的際遇.
可能跟隨主的人.
有些人走在這條路.
他會比較順.
你不要羨慕他的順.
你羨慕的那個是救恩.
有些人是信了主之後.
飛黃騰達.
在職場裡面發光發熱.
有些人可能信了主.
他面對的很多困難.
如果你只是看他的際遇.
你自己不會走的.
我只是單單看.
上帝想我怎樣.
我生命裡面.
很多事情上帝已經為我們預備了.
我們有一條軌跡.
我們跟著這條軌跡.
來走我們自己的人生.
不要作那麼多的比較.
人與人.

$^{921}$比較就比死人.
彼得.
他的.
這句說法背後.
就帶著.
如果其他門徒比他更差.
他可能會墊高自己.
或者他自己更轟烈一點.
他覺得自己英雄感強一點.
不過耶穌說.
那些不是你需要考慮的因素.
單單就看著.
主要怎樣.
要他走的路.
那些.
跟他沒有關係.
最後.
當耶穌這樣回答的時候.
彼得沒有特別的被記載.
他有什麼回應.
不過在他一生知道.
他是跟隨主.
當彼得接受耶穌的責備.
那些事情與你無關.
他醒一醒.
他知道他人生要走自己的路.
甘心接受上帝為他預備的人生的這段路.
各位弟兄姊妹.
我們每一天.
都要在這段時間.
去看一看.
耶穌的回應.
各位弟兄姊妹.
我們每個人都有不同際遇.
有些人可能在地上.
存留的日子.
會長久一點.
高壽地用歲月.
來見影神.
大概仕途若寒就是一個這樣的身份.

$^{961}$不過有些.
有些門徒.
可能在很短暫的人生.
匆匆的一生.
他們就用他人生.
來見證主.
甚至用他的死來榮耀神.
每個人.
所走的路是不同的.
你和我.
將來要走的命途.
都不一樣.
我們唯一共同的就是.
我們都是以耶穌基督成為我們的主.
我們就.
靠賴.
拉著主的衫尾.
跟隨他的腳蹤.
走上帝為我們預備.
走走的人生的道路.
當我們甘心的時候.
我們生命裡就沒有那麼多.
那些雜質.
或者有很多不忿氣.
在我們生命裡出現.
我們更加知道.
我們生命是很獨特的.
我們就用我自己獨特的生命.
見證主耶穌.
今天這段經文.
好像很闊.
說的話很涵豈.
不過希望今天這段經文.
或者今天離開這個禮堂.
你問自己三個問題.
你愛耶穌.
比謝西更深嗎?.
第二條問題.
你甘不甘心.
接受耶穌在你生命裡的命定?.

$^{1001}$你願不願意接受你生命的限制?.
第三.
你知不知道你自己生命裡.
有些獨特的路上帝要你走.
那條獨特的路.
沒有人可以取代你的.
你願不願意走這條獨特的路?.
求主繼續幫助我們.
讓祂的話語成為我們的激勵.
成為我們的激勵.
\newpage



\section{申命記 14:1-29}
\label{sec:aQhmLPtYC40}
\textbf{2025年3月23日崇拜直播｜陳明泉牧師｜從生活細節,區分特殊身份｜申命記14\_1-29}
\newline
\newline
連結: \href{https://youtube.com/watch?v=aQhmLPtYC40}{\texttt{https://youtube.com/watch?v=aQhmLPtYC40}} ~~~~ 語音日期: 2025-03-23
\newline
\newline
\hyperref[sec:LKA6QvIWId8]{\small{< < < PREV SERMON < < <}}
~
\hyperref[sec:index_chronic]{\small{[返順時目]}}
~
\hyperref[sec:P1VP0bkSNZc]{\small{> > > NEXT SERMON > > >}}
\newline
\newline
申命記 14:1-29
\newline
\begin{longtable}{cl}
\hline
\hline
章節 & 經文 (和合本修訂版)\\
\hline
14:1 & \begin{tabularx}{0.7\textwidth}{X} 「你們是耶和華---你們神的兒女。不可為了死人割劃自己，也不可使額上光禿； \end{tabularx} \\ \\ \relax
14:2 & \begin{tabularx}{0.7\textwidth}{X} 因為你是屬於耶和華－你神神聖的子民，耶和華從地面上的萬民中揀選了你，作自己寶貴的子民。」 \end{tabularx} \\ \\ \relax
14:3 & \begin{tabularx}{0.7\textwidth}{X} 「凡可憎的物， 你都不可吃。 \end{tabularx} \\ \\ \relax
14:4 & \begin{tabularx}{0.7\textwidth}{X} 可吃的牲畜是：牛、綿羊、山羊、 \end{tabularx} \\ \\ \relax
14:5 & \begin{tabularx}{0.7\textwidth}{X} 鹿、羚、麃子、野山羊、瞪羚、羚羊、山綿羊。 \end{tabularx} \\ \\ \relax
14:6 & \begin{tabularx}{0.7\textwidth}{X} 凡蹄分兩瓣，分趾蹄而又反芻食物的牲畜，你們都可以吃。 \end{tabularx} \\ \\ \relax
14:7 & \begin{tabularx}{0.7\textwidth}{X} 但那反芻或分蹄之中不可吃的是：駱駝、兔子、石獾，雖然反芻卻不分蹄，對你們是不潔淨的； \end{tabularx} \\ \\ \relax
14:8 & \begin{tabularx}{0.7\textwidth}{X} 豬，雖然分蹄卻不反芻，對你們也是不潔淨的。牠們的肉，你們一點都不可吃；牠們的屍體，你們也不可摸。 \end{tabularx} \\ \\ \relax
14:9 & \begin{tabularx}{0.7\textwidth}{X} 「水中可吃的是這些：凡有鰭有鱗的都可以吃； \end{tabularx} \\ \\ \relax
14:10 & \begin{tabularx}{0.7\textwidth}{X} 凡無鰭無鱗的都不可吃，對你們是不潔淨的。 \end{tabularx} \\ \\ \relax
14:11 & \begin{tabularx}{0.7\textwidth}{X} 「凡潔淨的鳥，你們都可以吃。 \end{tabularx} \\ \\ \relax
14:12 & \begin{tabularx}{0.7\textwidth}{X} 不可吃的是：鵰、狗頭鵰、紅頭鵰、 \end{tabularx} \\ \\ \relax
14:13 & \begin{tabularx}{0.7\textwidth}{X} 鸇、小鷹、鷂鷹的類群， \end{tabularx} \\ \\ \relax
14:14 & \begin{tabularx}{0.7\textwidth}{X} 各種烏鴉的類群、 \end{tabularx} \\ \\ \relax
14:15 & \begin{tabularx}{0.7\textwidth}{X} 鴕鳥、夜鷹、魚鷹、鷹的類群、 \end{tabularx} \\ \\ \relax
14:16 & \begin{tabularx}{0.7\textwidth}{X} 鴞鳥、貓頭鷹、角鴟、 \end{tabularx} \\ \\ \relax
14:17 & \begin{tabularx}{0.7\textwidth}{X} 鵜鶘、禿鵰、鸕鶿、 \end{tabularx} \\ \\ \relax
14:18 & \begin{tabularx}{0.7\textwidth}{X} 鸛、鷺鷥的類群、戴鵀與蝙蝠。 \end{tabularx} \\ \\ \relax
14:19 & \begin{tabularx}{0.7\textwidth}{X} 凡有翅膀卻爬行的群聚動物對你們是不潔淨的，都不可吃。 \end{tabularx} \\ \\ \relax
14:20 & \begin{tabularx}{0.7\textwidth}{X} 凡潔淨的鳥，你們都可以吃。 \end{tabularx} \\ \\ \relax
14:21 & \begin{tabularx}{0.7\textwidth}{X} 「凡自然死去的動物，你們都不可吃，可以給城裡寄居的人吃，或賣給外人，因為你是屬於耶和華－你神神聖的子民。「不可用母山羊的奶來煮牠的小山羊。」 \end{tabularx} \\ \\ \relax
14:22 & \begin{tabularx}{0.7\textwidth}{X} 「每年，你務必從你播種的一切收成，田地所出產的，取十分之一獻上。 \end{tabularx} \\ \\ \relax
14:23 & \begin{tabularx}{0.7\textwidth}{X} 要在耶和華－你神面前，就是他選擇那裡作為他名居所的地方，吃你所獻十分之一的五穀、新酒和新的油，以及牛群羊群中頭生的，好讓你天天學習敬畏耶和華－你的神。 \end{tabularx} \\ \\ \relax
14:24 & \begin{tabularx}{0.7\textwidth}{X} 當耶和華－你的神賜福給你的時候，耶和華－你神選擇立他名的地方若離你太遠，路途太長，使你不能把這東西帶到那裡去， \end{tabularx} \\ \\ \relax
14:25 & \begin{tabularx}{0.7\textwidth}{X} 你可以把它換成銀子，把銀子包起來，拿在手中，往耶和華－你神所選擇的地方去。 \end{tabularx} \\ \\ \relax
14:26 & \begin{tabularx}{0.7\textwidth}{X} 在那裡，你可以隨心所欲用銀子或買牛羊，或買清酒烈酒，或買任何你心所想的。你和你的全家要在耶和華－你神面前吃喝歡樂。 \end{tabularx} \\ \\ \relax
14:27 & \begin{tabularx}{0.7\textwidth}{X} 「住在你城裡的利未人，你不可離棄他，因為他在你那裡沒有分得產業。 \end{tabularx} \\ \\ \relax
14:28 & \begin{tabularx}{0.7\textwidth}{X} 每三年的最後一年，你要把那一年收成的十分之一取出來，積存在你的城中； \end{tabularx} \\ \\ \relax
14:29 & \begin{tabularx}{0.7\textwidth}{X} 那沒有與你一起分得產業的利未人，和城裡的寄居者，以及孤兒寡婦，都可以前來，吃得飽足，好讓耶和華－你的神在你手裡所做的一切事上賜福給你。」 \end{tabularx} \\ \\
[1ex]
\hline
\hline
\end{longtable}
$^{1}$今日宣讀的經訓是記載在舊約《新命記》14章1至21節.
「你們是耶和華你們神的兒女.
不可為死人用刀劃身.
也不可將額上剃光.
因為你歸耶和華你神為聖傑的民.
耶和華從地上的萬民中.
揀選你特作自己的子民.
凡可憎的物都不可吃.
可吃的牲畜就是牛.
綿羊 山羊 鹵.
羚羊 豹子 野山羊.
尾鹿 黃羊 青羊.
凡分體成為兩範.
又倒著的走獸.
你們都可以吃.
但那些倒著或分體之中不可吃的.
乃是駱駝 兔子 沙蕃.
因為是倒著不分體.
就與你們不結證.
豬因為是分體卻不倒著.
就與你們不結證.
這些瘦的肉你們不可吃.
死的也不可摸.
水中可吃的乃是這些.
凡有刺有鱗的都可以吃.
凡無刺無鱗的都不可吃.
是與你們不結證.
凡結證的鳥你們都可以吃.
不可吃的乃是鵰.
狗頭鵰 熊頭鵰.
箭 小鷹 妖鷹與其鱗.
烏鴉與其鱗.
鴕鳥 夜鷹 魚鷹.
鷹與其鱗.
囂鳥 貓頭鷹.
角雌 堤壺.
禿鵰 驢鰭.
冠 老獅與其鱗.
大任與蝙蝠.
凡有刺旁爬行的物.

$^{41}$是與你們不結證.
都不可吃.
凡結證的鳥你們都可以吃.
凡致死的你們都不可吃.
可以給你城裡寄居的吃.
或賣與外人吃.
因為你是歸耶和華里神.
為聖潔的民.
不可用山羊羔母的奶煮山羊羔.
各位聽眾 平安.
你會很奇怪 今天是我在講道.
今天講完 駱建源博士昨晚九點多打電話給我.
他病得很厲害.
咳得很厲害.
其實他在星期六 昨天幫我們黃昏崇拜.
其實他也是在支持.
看完醫生 照完肺之後.
聞了那些氣.
聞完之後 還是沒有好轉.
昨晚還差了.
接著他就沒有辦法了.
於是乎九點多就預備這篇講章.
這篇講章其實我是在少崇今天會講的.
但是對大人或成年人來說 可能要調教一下.
另外剛才的經文 讀經的阿嬸.
馬上要查剛才你讀的字.
大家平常都不會這樣讀的.
今天你就學到了.
今天是講一段很冷門的經文.
但是今天這段很冷門的經文.
其實背後有一些很重要的意義給我們.
希望大家今天看的篇幅會再長一點.
由第一節讀到第29節.
不過我們精簡.
我們剛才讀經 聆聽到21節.
今天這段經文 其實最主要你看到.
吃什麼不吃什麼.
其實飲食本身是一種很重要的文化.
這種文化觀念構成了不同的地方.
或者不同民族裡面背後的精神.

$^{81}$一個很簡單的例子.
如果你去吃薄餅.
你去到意大利你告訴別人加菠蘿.
你知不知道會有什麼效果.
你們有沒有試過.
如果你加菠蘿.
意大利人覺得你對他是一種冒犯.
因為意大利的薄餅是不會有菠蘿的.
美式的薄餅才有菠蘿.
如果你說要加菠蘿 他真的跟你拼了.
去年我去意大利.
還有你怎樣吃薄餅 人家都很講究的.
我們平時吃薄餅用刀叉.
很骨子的 手不髒的.
我用刀叉來切薄餅.
放進口裡的時候.
那個店長 老闆.
他看到我一下子拍了我肩膀.
意大利人都這樣的.
然後我不知道他在做什麼.
我說 delicious.
他再拍多一下.
然後他說薄餅不能用刀叉吃.
是用手拿著吃的.
所以他拍我.
然後他想餵我 用手拿著.
叫我拿著薄餅吃那個pizza.
我覺得這些很有趣.
用刀叉吃不用這麼嚴格.
和用手吃都是一樣.
不過對於不同的民族.
對於不同文化的人來說.
其實你怎樣去處理食物.
食物的方式是怎樣.
其實是一種民族身份.
今天如果你用這個角度來看.
摩西去教導一班以色列人.
在講什麼動物可以吃.
什麼動物不能吃.
背後就有一種民族精神在這裡.

$^{121}$而這個也是預備那班以色列新的一代.
他們將要進入到加拿大.
他們怎樣來過他們的生活.
在這些生活小細節裡面.
怎樣去表達他們是一班屬神的子民.
今天我們用這29節.
我們蓋藍色來了解一下.
他們在文化上怎樣去適應成為一個神國的子民.
我們很需要上帝幫助我們.
我們先祈禱 求聖靈幫助我們明白經文.
天父上帝願你在我們當中帶領.
幫助孫講的講得清楚.
願你的話語能夠進入弟兄姊妹的心靈當中.
求你也幫助眾弟兄姊妹.
在我們每一個聆聽的時候.
我們帶著謙卑.
我們繼續來仰望你.
願你指教我們今天我們的行事為人.
求你的聖靈與我們眾人同在.
帶領我們下來領受你話語的時間.
獻上我們的禱告奉求基督耶穌得勝的聖名.
Amen.
今天我們看的第十四章.
我大概不是很平均的篇幅.
我分為三個段落.
第一個段落是從第十四章第一節.
主要是這一節來看.
第十四章第一節裡面這樣說.
你們是耶和華你們神的兒女.
不可為死人用刀劃身.
也不可將額上剃光.
在這裡你會看到.
什麼叫劃身呢.
其實劃身今天我們不用這個字.
我們用紋身.
其實是同一個意思.
就是在身體裡面.
將顏色打進去的身體.
現在要用的機器.
以前是用刀.

$^{161}$將一些顏料.
割傷了你的皮膚.
然後將那些顏料.
植入你的身體當中.
這個與一般的印水紙不同.
這種植入皮膚裡面的顏料.
這些顏色.
相對來說是一個很長久很永恆的顏色.
在這裡為什麼摩西特別要說這一段.
或者用這一節來說不可劃身呢.
除了他們將入去迦南地之外.
可能會受著迦南的.
有一些當時的文化風氣影響.
更加重要其實有這個劃身.
或者紋身的傳統.
主要來自埃及.
在古埃及裡面.
當時有很多人有一些考古發現.
有些木乃伊身上.
是有這些紋身這些記號的.
通常這些記號是來自幾樣東西的.
他們對神明的敬拜.
有那些鳥仔.
有些獅子.
又或者有一些人的神像.
曾經有一些說法說.
會不會是紋身代表著那些低下階層.
然後那些人就用這種方式.
古埃及的時候.
用這些來保佑自己.
不過後來在古埃及的考古發現.
原來是貴族有錢人.
才做得到這個劃身的記號.
除了這些神明之外.
也有一些人會用一種劃身.
來將為了那些已逝去的人的亡靈.
既是一種哀悼.
記住的家裡的人.
或者一些很重要的人.
同時間也好像是.

$^{201}$祈求他們的亡靈.
會依附在他們的生命當中.
在他們的身體裡面.
保佑他 保護他.
以致他的生命得著平安.
所以你從這個角度來看.
這種的紋身或者這種的劃身.
背後是有宗教.
有祭祀.
甚至乎祈求一種神秘的能力.
來保護紋身的那個人.
在這裡裡面.
他們以色列人將要進入迦南地.
他們進入的時候.
他們不是文化真空的.
他們不會什麼都沒有的.
他們都有自己的傳統習俗.
經過幾百年在埃及地作為奴的那些人.
那種的生活.
即使脫離了他們.
但是生活習慣.
或者在那些以色列百姓當中.
或多或少都會將舊時古埃及的那些文化.
帶到他們的生活當中.
而進入迦南又進入一個新的文化.
如果在這裡裡面.
摩西不幫他們清一清.
或者有一種意識.
清晰知道他們做的這件事.
不是上帝的吩咐.
或者這件事不合乎上帝的心意.
他們將來就會很迷失.
所以在這裡裡面第一個節數第一節.
他就說到你們既然要脫離埃及.
你們要進入迦南地.
就不可以為了死去的人.
劃自己的身體.
甚至不是純粹為死去的人.
其實劃身這個已經成為了禁令.
就不可以藉著這些顏料.

$^{241}$在身體裡面去整.
以致到好像有一種神明的庇佑.
上帝的庇佑不是用這種方式去表達.
反而是用會幕來講上帝與人同在.
不是用身體割傷這種方式去做.
另外一個字眼.
也不可使額上光禿.
額上光禿是和合本的譯法.
有些譯本寫額上那個字.
其實是在講兩眉之間.
就是說眼眉不要剃眼眉.
不過你也覺得我們人很醜.
怎麼會剃眼眉呢.
其實這個也是那些埃及的傳統.
如果你有機會看過.
以前有上下年紀的那些看.
以前我們看過一套戲《奔雨》.
尤伯連納.
尤伯連納是光頭的.
千萬不要吸煙那個.
記不記得.
如果你記得你跟我差不多年紀.
就是光頭的.
其實舊時的古埃及那些皇室.
貴族人員.
他們會將身體的毛髮刮乾淨.
他們很乾淨很整齊.
所以你看到那些.
不論是那些木乃伊的神像.
或者那些皇的照片.
所有身上面上的毛髮都剃乾淨.
這種剃乾淨.
既是代表著一種乾淨.
同時間也代表著一種貴族走出來.
很有光鮮的一種表達.
如果你記得的.
你看《創世記》.
摩西將要為法老解夢.
在監倉之前他就做乾淨.
剃掉這些毛.

$^{281}$就是做這件事.
剃乾淨.
剃掉所有毛.
代表著他就歸入去.
好像埃及的那種文化傳統裡面.
不過在這裡面.
摩西提醒那些以色列人.
你不要將舊時埃及的那些文化.
劃新剃乾淨那些毛.
帶進去迦南地裡面.
你要做回以色列人.
你生命裡面.
你是神的兒女.
你就要用神兒女的方式.
來處理你的兒童.
其實在第一節裡面.
摩西在這裡面特別說.
你們是耶和華你們神的兒女.
這句話不只是在這裡面說的.
在第二節裡面再三的強調.
因為你是屬於耶和華.
你神聖的子民.
耶和華從地面上的萬民中.
揀選了你.
作自己寶貴的子民.
之前說過了.
寶貴的子民.
在我上一次講到有說過.
上帝的私人珍藏.
上帝隆而重之.
最寶貴的那班子民.
就是他們了.
摩西在這裡面再三強調.
這班以色列人.
不同舊時在埃及地裡面.
圍爐的這班人.
現在是上帝視為重要的.
而這班被上帝視為重要的人.
在你的言行舉止.
在你的外表裡面.

$^{321}$你們要與眾不同的.
這種的表達不是標旗立意.
而是說這班人是屬於耶和華上帝.
他們的生命是效忠耶和華上帝.
藉著他們的外表.
藉著他們的生活細節.
他們就呈現一種.
與世俗很不同的一種形態.
第二節剛才強調是神聖的子民.
去到第21節.
關於吃東西的段落.
其實都是再一次來講.
21節的後半節就講到.
因為你是屬於耶和華.
你神 神聖的子民.
21節 一頭一尾.
他是這樣包著.
來叮囑 勉勵整個以色列民族.
你們是和一般的其他子民.
或者其他民族是不同.
你們就要好好地凸顯自己.
今天這一段經文.
特別是第一節講關於劃新的經文.
有些弟兄姊妹就會覺得.
牧師會不會純粹是咎藥.
因為後來在講吃不吃的那些東西.
今天我們基督徒也吃.
包括吃豬.
那塊豬腩肉我昨晚才吃完.
沒理由的.
那些只是應用在以色列人那裡.
未必我們是新約信徒.
我們是跟耶穌的.
會不會凌駕於不需要跟隨這些.
泛民族節身上呢.
很多人看這些經文.
就覺得咎藥歸咎藥.
不關自己的事.
所以他就覺得那樣東西.
讀了當作一個資訊.

$^{361}$知道就算了.
不過如果我們這麼輕忽.
來處理咎藥.
我們有很多中間重要的元素.
我們錯失了.
講回剛才那個劃新的例子.
有些年輕的弟兄姊妹問牧師.
我可不可以去做民生.
其實他不是問你可不可以.
如果你說不可以.
他也照樣去做.
很多時候年輕人是這樣.
其實他是想覺得你.
好像是祝福他.
支持他祝福他這個決定.
對我來說.
我覺得我不會贊成.
不會鼓勵弟兄姊妹去做這件事.
即使你紋一個耶穌像在那裡.
弄一個十字架在那裡.
我都覺得不是這個聖經.
對我們的教導.
首先在這裡裡面.
他講到那個民生背後.
這種的形態是有異教.
舊時那些古埃及的那種.
迷信的色彩.
含在當中的.
而這種的觀念.
去到今天社會.
同樣都是一樣有這種觀念.
今日很多坊間裡面.
雖然五花八門.
但是我們對於我們中國人社會.
那個文化.
那個民生的文化記號本身.
都是在講一些的外邦神明.
又或者有一些.
他們自己的神秘的經歷.
如果你真的見過這些.

$^{401}$你心裡面會打一個突.
我以前在突破.
有一個同事.
他以前是黑社會.
那次去到海灘.
游泳.
他一打開衣服.
嘩 你是基督徒.
整個觀音紋在背後.
然後我說很厲害.
後面那個公仔.
他說是呀.
我年輕時很頑皮.
就紋了.
很貴的這樣說.
我說你信了之後怎樣.
撕不掉.
沒有辦法.
現在都做不到.
所以整個就是紋著觀音在背後.
就一起唱詩祈禱.
在海灘裡面打噴嚏.
就是這樣.
很多時候今天背後的民生.
都是有這種這樣的.
文化意涵在.
我再問你以前為什麼會做這樣的舉動.
他說我們以前劈油.
或者我們出來走.
保佑.
戴個像都不夠.
這些隨身攜帶.
你就跟著.
所以那種同在感很強.
以前我們是這樣想的.
不過我信了耶穌之後.
那個同在知道不是用這種方式.
其實我以前不懂得想才這樣做.
我要說什麼呢.
弟兄姊妹.

$^{441}$那個民生背後.
可能去到今天沒有那種神秘果效.
但是在文化裡面.
我們要區分.
我們不要做一些叛島人的工作.
即是說你紋一個十字架在那裡.
聖經從來都沒有叫我們.
將十字架紋在身上.
就令到我們變得很敬虔.
或者這件事提醒自己.
真正的話語不是紋在身上.
真正的話語要藏在心裡.
如果我們真正要去跟隨基督.
我們不是從肉體.
從顏料裡面.
令到我們與眾不同.
反而是我們的心靈裡面.
有上帝的話語.
以致到我們的行事為人.
藉著這些話語.
來端正我們的人生.
所以第一點.
在這裡摩西.
雖然篇幅不是很長.
但是他就告訴那些以色列人.
你不要帶回埃及的文化.
進入迦南地裡面.
整個文化.
整個以色列文.
他們是上帝寶貴的子民.
是屬神的子民.
好了我們看第二個段落.
因為字不適同.
我怕獻醜.
我不讀了.
剛才阿沈讀了.
不過在這一段經文裡面.
由第二節到第21節.
不斷又不斷的來到重複.
說可以吃.

$^{481}$有些東西是不可以吃.
我們先看看第三節.
到第八節那裡.
就是說有些可憎的肉.
然後他就說到.
說完可憎的肉之後.
他就說到哪些肉是可以吃的.
很具體說到.
一些是普遍在家裡.
或者放牧養的.
甚至有些野生的.
這些肉是可以吃的.
牛,綿羊,山羊,鹿,羚羊,豹子.
是不是讀豹.
我忘記了.
等等.
這一類的動物.
是分題分成兩反.
有杜絕的.
分題一分為二.
而且可以杜絕的.
這些在上帝眼中.
視為潔淨.
他們就可以吃了.
但是如果那些是杜絕.
或者分題之中不可吃的.
就是有分題或者杜絕.
那些就視為不潔.
那些視為不潔的物.
你摸都不可以吃.
更加不可以吃.
在這裡很多弟兄姊妹說.
那些潔淨不潔淨.
是從什麼角度去看的呢.
曾經有一些色經學家說.
會不會是從衛生角度呢.
豬有豬兔蟲.
其實牛也有寄生蟲.
如果純粹用.
當時我們用衛生角度去看.

$^{521}$有些是解釋不了的.
還有一些就是有文化習俗.
這個也是對的.
有些文化是會吃這些動物.
有些文化是不會吃的.
不過有些地方也是解不通的.
到最後其實色經學家也是說.
潔淨與不潔淨的動物.
到最後訴諸於什麼呢.
只是上帝主觀或者上帝的旨意.
上帝說了這東西可以吃就可以吃.
這東西不可以吃就不可以吃.
這個最後是訴諸於上帝.
對動物吃得動物的那種看法.
所以我們今天就不能完全.
去爭辯的或者我們完全沒辦法.
有一個通則來到去.
有條定例來到去講的.
摩西在這裡講到賭雀.
和昏蹄的.
昏為兩反的那些才可以吃.
其他那些都視為不潔.
就不可以吃了.
還有一樣東西就是.
那些不潔的被視為吃落肚不潔的那些動物.
不等於上帝很憎那隻豬.
不是上帝在講可憎的意思.
不是在講上帝創造出來又很憎他.
不是在講那種情感上的憎惡.
而是吃落肚裡面的那種憎惡.
我這樣講好像很複雜.
你用文化角度你就知道了.
早幾十年前或者舊時國內.
有些人沒東西吃或者喜歡吃.
吃狗肉.
有些人喜歡的.
有些人點頭.
你沒吃過.
有些人說喜歡吃狗肉的.
但今天這個世代我們不會吃的.

$^{561}$我們覺得是朋友.
如果你在西方社會.
你說吃狗肉人家會打你.
就是一種這樣.
那個憎惡不是那種動物可憎.
而是吃這個行為令到人家很憎惡.
這樣你都這樣做.
又或者那次有機會去到日本.
都是很現代的社會.
那次去吃.
我差點吃了就是吃馬肉.
怎麼會吃馬肉.
我們覺得這麼噁心.
但我見到隔壁桌子吃到碎碎聲.
我們覺得不是那種動物骯髒.
而是吃這個行為令人帶來憎惡.
你用這個角度去看.
那種潔淨與不潔淨.
不是看那種動物本身.
而是看那個吃.
或者你怎樣對待牠的行為.
來界定牠是否潔淨與不潔淨.
和可憎與不可憎.
如果你用這個角度去看.
接下來所說的那些魚.
在水中有刺有鱗的就可以吃.
沒有鱗沒有刺的就不可吃.
我們今天吃的有些魚.
鱗鱗鱔應該沒有鱗沒有刺的.
那就不可以吃.
今天我們吃盤龍鱔.
我們中國人.
但在舊約在摩西裡面說到.
這些食物對他們來說是可憎惡之物.
你們不可以吃.
到最後就是吃那些雀鳥.
雀鳥他們都會吃的.
這裡說了有不同的動物是不可以吃的.
有鵪鶉就可以吃.
上帝用東風吹的鵪鶉.

$^{601}$在以色列人在曠野裡面吃這些雀鳥.
來誘肉吃.
令他們有些補充一些蛋白質.
在可以和不可以吃的過程裡面.
潔淨與不潔淨.
其實就是從這些生活的細節裡面.
藉著他們吃的食物.
以致反映了他們自己生命的人生.
西方有一句諺語說.
You are what you eat.
你就是你所吃的.
這句說話不是純粹說你喜歡吃什麼.
你就是一個什麼人.
而是背後有一種文化.
我們背後的文化.
端正著左右著我們的口味.
我們對於事物的判決.
即使多麼飢餓我們都不會吃.
這些就是這種文化觀念.
塑造了我們整個文化的人格.
摩西在這裡面再三的強調.
這些不潔之物或者這些只可以吃的東西.
藉著他們所吃的食物.
來界定他們.
他們究竟是從什麼角度來看.
其他外邦或者其他的文化.
可能看到有食糧.
可以大量養殖就去吃了.
但對於當時的以色列人來說.
摩西就覺得這東西不能吃.
就不可以這樣做.
值得留意還有一句.
很多人都視為孤句的.
就是在第21節的後半節.
我再讀一讀21節給你聽.
凡致死的你們都不可吃.
不可以給你城裡寄居的吃.
或賣與外人吃.
因為你是龜爺或你神為聖傑的民.
跟著後面這句.

$^{641}$不可用山羊羔母的奶煮山羊羔.
這句說話特別抽出來.
有些和合本會抽了.
獨立了出來.
什麼意思呢.
為什麼會這樣寫呢.
不同的色經學家有不同的說法.
最主要來自兩個最重要的看法.
第一就是在舊約裡面.
特別是在說以色列條例.
如果你要宰山羊羔或者沒有山羊.
是可以吃的.
不過你屠宰的過程裡面.
你不可以一次過把牠的母子.
烹調了來吃.
他覺得這東西背後就是一種不人道.
你好像滅了他整個家一樣.
背後會將一種人愛的精神推展出去.
所以你宰一隻山羊羔.
你不可以把牠的母子也宰了.
一起去做獻祭或者去吃.
不可以這樣做.
隨便一隻.
這是第一個有這樣的條例.
第二個更加承接著這個說法.
不過他更加用一種憐憫的角度.
那些山羊羔即是那些小隻的BB山羊.
牠吃的是什麼.
喝媽媽的奶.
即是沒有山羊的奶.
如果那個奶是用來餵牠.
但是當人們去烹調的時候.
用餵牠的奶來煮牠.
你不覺得那種感覺.
為什麼會這麼殘忍.
你吃牠但你還要用.
他自己用來喝的奶來烹調.
讓他更加好吃等等.
有些說法就是為什麼會這樣做.
原來在迦南地有些文化.

$^{681}$就是說原來這樣的煮法.
原來是會促進了當時.
對於農業社會裡面.
豐收生養眾多.
原來這樣煮.
這個本身是一個.
類似祭祀的一種烹調過程.
而且這樣煮的東西吃是特別好吃的.
在迦南有這樣的一個文化.
我們從聖經的角度.
其實是很一致的.
肉是可以吃的.
不過我們要尊重生命.
背後就是帶著一份這樣的人愛的心.
來處理食物.
對於一群以色列神國子民來說.
他們不單是吃.
他們不只是純粹想滿足他的肚腹.
那樣東西好不好吃.
或者那樣東西是不是上帝起源.
在烹調的過程裡面.
他是不是帶著上帝的憐憫.
來看這個蒼生之物.
所以在這裡.
那種吃與不吃.
背後就有一個很深厚的.
上帝對人的觀念.
這個是第二個段落.
我們看第三個段落.
時間關係我們要很快.
22節讀到第29節.
我們先看22-23節.
每年你務必從你播種的一切收成.
田地所出產的取十分之一獻上.
要在耶和華你神面前.
就是他選擇那裡作為他名居所的地方.
吃你所獻十分之一的五穀.
新酒和新的油.
以及牛群羊群中投生的.
好讓你天天學習敬畏耶和華你的神.

$^{721}$在這裡第22-23節就講到一個奉獻.
以色列人那種的宗教制治.
或者對於食物之後.
對上帝有什麼表達呢.
在這裡他就強調.
每年你要將十分之一的出產.
獻在上帝的面前.
要在耶和華面前獻祭.
獻了這個祭的祭物.
他特別講到擺上十分之一之後.
其實那些的食物.
獻祭者是可以吃的.
大概這個是一個平安祭.
一個感恩祭.
你獻上然後放在上帝面前.
做了儀式之後.
一家人一起來領受那些上好的牛羊.
那個十分之一其實就強調.
是一種帶著感恩的心.
摩西在這裡講的那個十一為什麼這麼重要呢.
就是說原來這種獻祭的過程.
就是讓你天天學習敬畏耶和華你的神.
奉獻原來是一種學習.
這個學習不是說我給東西.
我給東西上帝.
或者我交換一些條件.
原來這個學習是在生活當中.
當你生命裡面有豐收有收成.
你藉著你擺上這些十分之一.
來感恩學習如何敬畏上帝.
今天我們很多時候會講敬畏神.
我們唸聖經也是這樣唸.
不過那種敬畏從來不是純粹在口裡.
又或者是觀念上.
在摩西的教導當中.
這是一種行動.
這種敬畏其實就實踐了.
在他們每年的生活當中.
所以這個第23節就講到.
那個十一奉獻是要學習對上帝的敬畏.

$^{761}$我們繼續看第24至26節.
當耶和華你的神賜福給你的時候.
耶和華你神選擇納他名的地方.
若離你太遠路途太長.
使你不能把這東西帶到那裡去.
你可以把它換成銀紙.
把銀紙包起來拿在手中.
往耶和華你神所選擇的地方去.
在那裡你可以隨心所欲用銀紙.
或買牛羊或買清酒烈酒.
或買任何你心所想的.
你和你的全家要在耶和華你神面前吃喝歡樂.
第二個段落是講什麼呢.
如果那個路途和獻祭的地方.
上帝所指定獻祭的地方很遙遠.
真的帶不到.
你可以奉獻金錢.
去到奉獻金錢的時候.
那個金錢去到你獻祭的地方.
你就可以買牛羊.
你喜歡的牲畜.
既是獻祭.
獻完祭之後就可以一家人去吃喝歡樂.
這個吃喝歡樂.
其實吃的肉上帝不會吃掉你的肉.
而是上帝藉著你這樣的來擺上.
讓你的家庭都能夠經歷到.
那種分享這份上好的肉的那種歡愉.
有些色經學家說到.
你見到剛才說食物不可摸不可吃.
到這裡你隨意買隨意吃.
那些字眼你想一下好像很熟.
面孔從哪裡來呢.
創世記.
伊甸園裡面.
上帝說伊甸園裡面.
園中所有樹上的果子你都可以吃.
唯獨是這一棵不能吃.
就是分辨善惡樹.
吃了你會死.

$^{801}$同樣這個意象都放在食物和分享的神學觀念裡面.
在你的吃喝當中.
是有一條界線.
有些東西你能吃.
有些東西不能吃.
而能吃的目的上帝就是要把豐盛快樂帶給你.
將這份喜悅和你家人分享.
所以這個隨意吃.
隨心所吃.
在上帝的規矩裡面.
在上帝的規範當中.
這群人是可以吃喝快樂.
而且是上帝的心意.
我們再看最後一個段落.
27-29節.
除了自己家裡快樂之外.
27-29節.
他要將你的快樂和別人分享.
住在你城裡的利未人.
你不可離棄他.
因為他在你那裡沒有分得產業.
每三年的最後一年.
你要把那一年收成的十分之一取出來.
積存在你的城中.
那沒有與你一起分得產業的利未人.
和城裡的寄居者.
以及孤兒寡婦.
都可以前來吃得飽足.
好讓耶和華你的神.
在你手裡所作的一切事上.
賜福給你.
在這裡上帝要將焦點放在.
這些能夠有能力有出產的人.
十一奉獻每年之外.
每三年還要做多個額外的十一.
額外的十一奉獻是為了什麼人呢?.
原來在分地裡利未人是沒有產業的.
這一群在回望或者在當時負責.
祭祀的一群祭師裡面.
這一群人以上帝耶和華為他的產業.

$^{841}$所以對於其他以色列的其他那些支派.
他們就要每三年就要分享一次.
他所有的來讓這些利未人.
都可以得享這一份上帝而來的恩典.
除了沒有產業的利未人之外.
上帝特別的來看顧那些.
在他們城裡寄居的那些無家可歸的那些人.
那些窮苦大眾.
又或者是那些孤兒寡婦.
那些孤兒寡婦沒有產業.
上帝不是看他們是局外人.
上帝的恩典是遮蓋他們.
他要鼓勵那些以色列人.
同樣都要有這一份的心腸.
除了自己家裡吃喝快樂.
他年年都可以享受.
但是每三年一次就要額外來照顧.
在城裡面這些人.
所以整個律法.
或者今天我們看的這些生活細節.
很瑣碎或者背後的精神.
其實是在說上帝對世界的關愛.
對生命的那種尊重.
文化上的那條界線.
他怎樣敬拜神明.
他怎樣經歷上帝的同在.
他怎樣來制治.
在民生裡面就看到了.
還有他怎樣分享食物.
以前只是自己一家人享用.
但是對於上帝國的子民來說.
那些孤苦流離的人.
那些沒有產業的利尾人.
那些孤兒寡婦.
上帝的眼目都要看到這些人.
他都希望那些以色列人.
來照顧他們.
供應他們所需.
今天這段經文看完了.
好像很瑣碎一樣.

$^{881}$我想今天對我們基督徒的意義是.
我們今天會不會在這個世代裡面.
能夠分別為聖呢.
那個神聖從來不是說.
你經常說一些很深奧的道理.
道德是重要的.
不過不是純粹口說道德.
在我們的生活細節裡面.
我們都活現了一些.
跟這個世俗有不同的地方.
我近來讀這段經文的時候.
我自己有一個祈禱.
我有一點擔心.
因為這個世代裡面.
很多時候覺得傳福音的方法就是.
令基督徒變得很不基督徒.
令到教會變得很不像教會.
很多朋友覺得.
令到教會變得好像康樂中心一樣.
令到你的團契.
好像在一個不覺得是宗教聚會一樣.
你的行事為人.
你不像基督徒.
這樣就不會令到別人排斥你.
這樣你就能夠有效傳福音.
跟這個世界有一些對話.
或者你讓別人覺得.
你不是跟這個世界脫節.
我覺得這個很重要.
不過當我們人生裡面.
所有的東西.
都跟這個世界沒有分別的時候.
我們基督徒的獨特之處在哪裡呢?.
當我們經常都是用這種.
覺得越不像基督徒的方式.
來活在這個世界.
這個不就是真正可以傳到福音嗎?.
還是因為我越來越不像.
久而久之我被這個世界同化了.
或者我用最冠冕堂皇的方式.

$^{921}$我是傳福音.
但到最後我們可能在這個信仰裡面失落了.
被這個世俗文化淹沒了.
令到我們站不穩陣腳.
求主幫助我們.
今天我們看申命記第十四章.
我們多些文化的意識.
幫我們有明確的判斷.
上帝的同在不是在身上.
上帝的同在在我們的心靈當中.
我們要活現和這個世界很不一樣.
又令到這個世界羨慕的生活方式.
我們活在一個恩典當中.
我們就成為一個祝福別人的人.
求主繼續幫助我們.
\newpage



\section{馬太福音 4:24-5:3}
\label{sec:P1VP0bkSNZc}
\textbf{2025年3月30日崇拜直播｜張永信牧師｜虛心蒙福｜馬太福音4\_24-5\_3}
\newline
\newline
連結: \href{https://youtube.com/watch?v=P1VP0bkSNZc}{\texttt{https://youtube.com/watch?v=P1VP0bkSNZc}} ~~~~ 語音日期: 2025-03-30
\newline
\newline
\hyperref[sec:aQhmLPtYC40]{\small{< < < PREV SERMON < < <}}
~
\hyperref[sec:index_chronic]{\small{[返順時目]}}
~
\hyperref[sec:I13jDOQMkLQ]{\small{> > > NEXT SERMON > > >}}
\newline
\newline
馬太福音 4:24-5:3
\newline
\begin{longtable}{cl}
\hline
\hline
章節 & 經文 (和合本修訂版)\\
\hline
4:24 & \begin{tabularx}{0.7\textwidth}{X} 他的名聲傳遍了敘利亞。那裡的人把一切病人，就是有各樣疾病和疼痛的、被鬼附的、癲癇的、癱瘓的，都帶了來，耶穌就治好了他們。 \end{tabularx} \\ \\ \relax
4:25 & \begin{tabularx}{0.7\textwidth}{X} 當時，有一大群人從加利利、低加坡里、耶路撒冷、猶太、約旦河的東邊，來跟從他。 \end{tabularx} \\ \\ \relax
5:1 & \begin{tabularx}{0.7\textwidth}{X} 耶穌看見這一群人，就上了山，坐下後，門徒到他跟前來， \end{tabularx} \\ \\ \relax
5:2 & \begin{tabularx}{0.7\textwidth}{X} 他開口教導他們說： \end{tabularx} \\ \\ \relax
5:3 & \begin{tabularx}{0.7\textwidth}{X} 「心靈貧窮的人有福了！因為天國是他們的。 \end{tabularx} \\ \\
[1ex]
\hline
\hline
\end{longtable}
$^{1}$接著是讀經.
今天的經文是《馬太福音》第四章二十四節到第五章第三節.
經文都印在你們的程序表內頁.
請聽我讀.
《馬太福音》四章二十四節.
「他的名聲就傳遍了敘利亞」.
「那類的人把一切害病的」.
「就是害各樣疾病,各樣疼痛的和被鬼附的」.
「癲癇的,癱瘓的都帶了來」.
「耶穌就治好了他們」.
「當下有許多人從加利利」.
「底加波利,耶路撒冷,猶太」.
「約旦河外來跟著他」.
第五章.
「看見這許多的人就上了山」.
「基於座下門徒到他跟前來」.
「他就開口教訓他們說」.
「虛心的人有福了」.
「因為天國是他們的」.
《聖經》讀筆.
跟著是宣講的時間.
今天的講員是張詠順牧師.
容我介紹一下.
張牧師本身是.
以前在建道神學院和恩福神學院的老師.
現在已經退休了.
不過他退而不休.
還是很有心的去栽培新一代的牧者.
是很願意為神所用的牧者.
我們請張詠順牧師.
很感動.
崇拜的領唱.
人歌合一是最高境界.
今天我看見你們是敬拜隊.
是人歌合一.
好.
跟著你的守靈隊.
換來換去你有沒有搞錯.
我真的很擔心你拿錯東西.
真的很厲害.

$^{41}$好啊.
很感動啊.
謝謝.
你們訓練有素.
這幾排講到.
我心裡面常常有一種千言萬語.
不知如何說.
因為我想講整個偉大.
我想講我們的卑微.
不知如何說.
我今天嘗試又講一下.
其實來黃浸我是舊地重遊.
為什麼呢.
因為去年年尾.
我帶一班恩福學生來黃浸取經.
先去見到神學院長洲取經.
取完之後.
我就跟人傾.
來哪間教會取經好呢.
有些人堅持來黃浸.
證明你們教會有實力.
跟著陳牧師就brief我們.
陳牧師他居公自位帶我們走又介紹.
那班學生回到教會.
回到神學院debriefing.
我讚不絕口.
這樣.
所以我今天是有少少舊地重遊.
好,我們回到我們所讀的經文.
我試一下powerpoint行不行.
我們先看看經文的背景.
我講到宏觀和微觀.
宏觀是講背景.
微觀是講色機.
講到這兩個point.
我們先看宏觀.
你看背景.
背景其實是四章的尾.
23節開始.
是一個summary.

$^{81}$一個概述.
耶穌走遍加利利.
在各會堂裡教訓人.
傳天歌的福音.
然後醫治百姓的病症.
你看到排列的次序.
耶穌先教訓.
teaching.
傳天歌的福音.
preaching.
然後才醫治.
次序是影現了耶穌的優先.
然後明星傳遍了敘利亞.
好有趣.
是馬太說傳遍敘利亞.
所以學者說.
馬太福音是寫於敘利亞.
安提沃.
這個很大件事.
我不想花太多時間去講.
那裡的人不一切害病的.
就像害各種疾病.
作者就講鬼婦.
癲癇.
癱瘓等等.
各樣奇難雜症.
都來找耶穌.
耶穌就自好探.
耶穌不會說我是傳福音.
我是教導.
你不要煩我.
耶穌不是這樣的.
他動了慈心.
見到人們像羊毛牧人一樣.
又醫他們又餵飽五千人.
耶穌是一個這樣的人.
我學耶穌.
大家知道我舊神學.
牧師又講到.
有時就.

$^{121}$這個當然是我的優先次序.
我早上經常去一間叫大XX吃早餐.
大家明白嗎.
每天都去那間.
因為我有個老人卡.
有三元雞.
他們還在推銷.
吃飯的時候.
早上有個老人家搶到.
什麼都沒有.
拿杯水在那裡.
不知道在幹什麼.
我動了慈心.
我不會說我講福音.
不理你.
我不是這樣的人.
我有一天終於動了慈心.
神教我.
你買個早餐給他吃.
好.
於是我就鼓起勇氣.
走去他前面.
老太不如我買個早餐給你吃.
你猜他怎麼回答我.
猜猜來.
互動一下.
他不要.
第一我下去吃的時候.
他自己在吃早餐.
你明白我的意思嗎.
做基督徒不只是講傳福音.
你有同情心.
有愛心.
耶穌就是這樣.
他傳到.
接著講登山.
就是報訓.
你明白嗎.
他見到人需要.
他都依.

$^{161}$來者不拒.
當下有許多人從加利利.
底加玻利.
耶路撒冷.
猶太 猶太.
哇.
萬人空巷跟耶穌.
好多人.
第五章.
耶穌看見很多人.
耶穌眼見.
萬頭轉動.
人頭湧湧.
我做主任牧師.
我做了很久主任牧師.
我見到人數多.
我都很開心.
但有時我又提醒自己.
虛火上升.
小心點.
我不是說你們.
不要對外入在.
好.
於是.
耶穌又上了山.
上山了.
繼而坐下.
坐下.
又上山又坐下.
門徒都跟前來.
所以登山補佛.
是為門徒而設.
他就開口.
教訓.
談說虛心的人.
一幅了.
耶穌又去到基本教導.
耶穌不是因為人多.
就充分了頭腦.
這班人可能是烏合之眾.

$^{201}$一打擊就潰不成軍.
耶穌說不要搞.
現在來一個門訓.
門徒.
先打好根基.
這是最重要.
所以其實登山寶訓.
主要是門訓.
整個馬太福音都是說門訓.
馬太福音被譽稱為.
The Teaching Gospel.
是講耶穌偉大的拉比.
你看這裡怎麼說.
他說他開口教訓他們.
你教訓人當然開口.
你不開口怎麼教訓.
這樣就變成耶穌偉大的拉比.
他好像拉比一樣坐下開口說.
你研究馬太福音.
你會覺得偉大到無論.
山上.
海邊.
會堂.
聖殿.
曠野.
在船上.
在屋裡.
你明白嗎?.
耶穌差不多每次都教訓人.
我又不是太德.
我通常都是廣東.
廣談就教訓平時.
我都懶得教訓人.
還有耶穌說.
基本上每樣東西都可以成為教材.
百合花.
飛鳥.
麻雀.
葡萄樹.
橄欖樹.

$^{241}$吃飯的餅和杯.
這是我的身體.
這是我的血.
你看到嗎?.
任何東西都可以成為耶穌教材.
真是厲害.
學者說耶穌的教導.
八成是poetic.
詩體.
詩歌題材.
你看看.
虛心的人.
看不看得到平衡句子?.
甚至主通文都是平衡.
天下得富 圓圓的尊 圓圓圓.
你看不看到?.
教導每一句都是poetic.
嘩!大哥.
我都想教書.
每一句都是.
詩作對.
做不到.
作者說.
我聽到耶穌是超級拉美.
很厲害.
所以太毒華.
好了.
我說完背景.
現在我們論經文.
虛心的人有福利.
這是八福的第一福.
肯定是重中之重.
你明白我的意思嗎?.
我想今天先說這個福.
因為八福之中的排頭位.
一定是最重要.
每一個基督徒必須要明白.
作耶穌的門徒.
必須具備一個很重要的心理質素.
虛心的人有福利.

$^{281}$原文是poor in spirit.
靈女貧窮.
心靈貧窮.
或者靈魂貧窮.
poor in spirit.
Blessed are the poor.
大家看不到.
是要種數.
每一個跟隨的人.
都要首先記起一件事.
排頭位.
就是心靈貧窮.
所以和合修定本.
心靈貧窮的人有福利.
原文是這樣表達.
虛心也可以.
心靈貧窮就是虛心.
虛心求教也可以.
不過是有個意義.
原文就是心靈貧窮.
或者靈魂貧窮.
或者心靈貧窮.
為什麼耶穌會這樣說.
學者可能是針對性.
聖經上的說話都是針對性.
因為當時有些人在那裡.
有個背景的位置.
耶穌突然這樣說.
我們法利賽人.
他們經常都覺得自己很棒.
很自滿自足.
在安書區裡.
覺得很舒服.
我們法利賽人守了律法.
很棒.
所以當時法利賽人有個格言.
有兩個人上天堂.
其中一個一定是法利賽人.
覺得自己很棒.
在安書區裡很厲害.

$^{321}$那些人是罪人.
那些人是瑞利.
我們聖潔無憾.
所以這句話.
耶穌說外表不等於所有.
心靈貧窮很重要.
是針對性說.
我和太太.
太太也是牧師.
兩個人在談一個問題.
究竟安書區有什麼問題.
安書區一個最大的問題就是麻木.
你的思維平白.
人在安書區裡一個最大的問題.
我們有一點習以為常.
腦袋不再醒覺.
學者說.
基本上科學家說.
我們的腦子只用了十分之一.
有沒有聽過.
安書區裡十分之一都沒有.
心靈平白.
安書區是最大的問題.
所以心靈貧窮很重要.
這兩天出了一個教會普查.
有沒有看.
看完之後很up .
所有的數字都下跌.
有些跌得很淒慘.
安書區.
各位聽眾.
我們需要心靈貧窮.
我們需要醒覺.
心靈不要再停擺.
很危險的.
殺到埋身都不知道.
你看地震.
災禍忽然臨到.
耶穌早就說了.
他警告我們.

$^{361}$拜託了.
平時做好心理準備.
做主角末道的心理質素很重要.
不是表面為為樂樂.
不是這樣的.
因為生活充滿考驗.
人生不如的事十之八九.
來到這裡.
我想介紹兩位哲學家.
第一個叫Socrates.
第二個叫Zenos.
我先介紹第一個.
因為這兩個人和心靈貧窮很有關係.
蘇格拉底是古希臘的第一位哲學家.
把滿天神佛的宙斯那些東西.
拉到地上去.
講人.
他第一個哲學家講人論.
講人的問題.
他的名句就是.
學生拍拉圖就說他教學很喜歡問人問題.
不斷問.
問到最後那個人說.
阿Sir我不懂了.
他說那你懂了.
你不懂就是懂.
所以在蘇格拉底的教導裡面.
人歸乎自知.
知己知彼百千百姓.
The first thing we have to know.
is self understanding.
自我認識很重要.
你讀到輔導學.
輔導學第一關就是教你自我認識.
究竟我是誰.
他說當你開始認識自己的時候.
第一樣你要知道的就是.
You have to know that you don't know.
Then you know.
You don't know then you know.

$^{401}$這樣的說法有些震撼性.
你是大學哲學家.
為什麼你說我們認識自己.
原來我是不懂的呢.
他說你不懂就開始是智慧的開端.
有時候我教書有講這件事.
學生說.
聖經上說經威神是智慧的開端.
那怎麼拆解.
當你知道不懂的時候.
你就經威神.
那就是.
所以耶穌說心靈貧窮是重要.
為什麼貧窮是重要.
因為貧窮轉移以求神.
為什麼財主教不信耶穌.
因為財主勇才自重.
所以財主轉移耶穌.
比下台的執行更難.
因為他勇才自重.
貧窮反而OK.
因為山窮水盡.
義無路.
林林不平.
以求神.
所以耶穌說虛心的人.
他想說的是.
請我們了解我們自己.
其實我們不懂很多東西.
請你信這位造物主.
Then you begin.
Type.
你開始.
明白.
開始知道.
開始有以求神的智慧.
第二位.
這個人.
比蘇格拉底帥.
跟他差不多同期.

$^{441}$這個人是斯多阿學派的祖師爺.
他高大衛猛.
帥氣高富帥.
學生很多人擁護他.
但他教學很多東西.
講的東西都不太肯定.
終於有一個學生.
他說阿Sir你這麼厲害.
這麼講述語.
這麼帥這麼威猛.
為什麼你講的東西不肯定.
他說你跟我來.
他就放了一盤沙在桌上.
畫了兩個圓圈.
一個大一個小.
小代表你.
大代表我.
圈內就是我們認識的東西.
我認識的東西比你多.
不過圈外不認識的東西.
我又比你知道多.
所以我知道其實我不知道的東西.
很多所以我說話不肯定.
我不知道比我知道的更多.
所以.
虛心的人.
令你憑虛的人.
有福了.
很多東西我們不知道.
我最近參加了一個.
天文學會.
開始研究天文.
這東西我不懂.
跟著跟著.
不好意思我還要下載.
Earth is one of 3.2 trillion.
少.
少就是一字後面12個零.
3.2 trillion planets.
in our galaxy.

$^{481}$the sun is one of 200 billions.
一個billion.
一萬億.
stars in the milky way.
and the milky way.
is just one of 2 trillion galaxies.
in the observable universe.
最慘就是observable.
那個字.
已經測到宇宙.
地球我們人類已經測到無論.
何況observe不到.
你就想想創世的神創造天地.
由無變為有.
神說有光就有光.
整個宇宙.
我跟學生說.
神用六日創造宇宙.
目的是想告訴以色列人.
為什麼要守安息日.
因為第七天休息.
神是全能的.
六日就搞定.
你明白我的意思嗎.
你一想你會發覺神的全能.
真的不知道你怎樣去形容.
整個宇宙.
這樣呈現在眼前.
最近我收到最新的料.
是說我們在黑洞裡面.
我回去想.
會不會黑洞又黑洞.
你一想到.
聖經說.
神的榮耀.
Heavenly care.
Heavens.
舊約的人得到神的啟示.
知道原來天外有天.
宇宙是universes.

$^{521}$你可以想像.
一個創造主.
這樣創造一個大到無論.
你無法去認識宇宙的時候.
你想想他多偉大.
偉大到我不知道怎樣形容.
所以我不知道怎樣說.
有一次有個學生問我.
牧師.
我們所認識的神.
和我們不認識的神.
怎樣比就百分比.
一比無限大.
我只能用保羅句說.
心在神的智慧何其難測.
他的足跡何其難尋.
我想到這些的時候.
我覺得自己渺小到無力.
我神我要敬畏你.
因為我這麼偉大的神.
我沒甚麼可以選擇.
我不知道能不能說到.
神何其大.
我接著上網問.
這個宇宙有沒有邊緣呢.
有沒有boundary呢.
答案.
宇宙在空間上的無限性.
只能體現於無數有限的.
其具體空間.
好像很難.
而不能全面的.
好像很難.
而不能全面的.
出來這些具體空間.
宇宙總是有有限.
又無限.
又有邊緣又無邊緣.
甚麼來的.
網上答案.

$^{561}$你問宇宙有沒有邊界.
他沒辦法告訴你.
現在哈佛最好望鏡.
望到你observe universe.
找不到邊緣.
我不懂理解.
你們可以採取多少人.
你明不明.
恩福大樓都很大.
有沒有邊界.
前一陣子上星期.
有個培靈會.
城鎮有沒有邊界.
有哪個地方沒有邊界.
我沒辦法理解.
甚麼叫沒有邊界.
還是observe universe.
我還是等朱天術說聖德榮耀.
我講不到.
言有主進而義無窮.
我都不知道怎麼說.
人到底是什麼.
你看看創世記.
看一看聖經.
第一章.
神創造天地.
第23節有晚上有早晨.
是第五日.
神說地要生出.
門密來各從其內.
生出.
昆蟲野獸各從其內.
是就正性了.
於是在第六日.
記住.
神說地要生出門密來各從其內.
生出.
昆蟲野獸各從其內.
神說地要生出.
門密來各從其內.

$^{601}$昆蟲.
是他們管理海裡的魚.
大家很熟.
人之所以與其他.
第六日所創造的.
其他動物.
唯一的不同是我們有神的形象.
如果不是其實沒什麼分別.
記住這一點.
我們之所以與其他動物有分別.
因為我們有神的形象.
我們可以與神連繫.
當我們不連繫的時候.
有一點點我們與其他東西.
沒什麼分別.
所以你明白為什麼這個世界這麼亂.
為什麼會這樣.
打仗都打了很多次.
打完第一次的大戰.
又打第二次的大戰.
愛因斯坦和羅斯福說.
第三次世界大戰是核子戰爭.
所以他說發明核子彈.
是人類最大的災難.
德國哲學家.
黑格爾說了一句話.
人類從歷史學到功課.
只有一個就是我們學不到功課.
打完又打.
We don't learn.
這個我們學到的唯一的功課.
就是希格.
我真的感覺到.
我們沒有神的形象.
我們失去了不同神的聯想.
我們就是這樣.
你看看經文怎麼說.
讀第一句.
我們一起讀.
你這蟲.

$^{641}$不用了.
那個是蟲.
那你和我是什麼.
我是一條可憐蟲.
我離開了神.
我是一條可憐蟲.
我病到.
七五八六.
不知道怎麼說.
兩年前我有一個Liver App.
去年我有一個戰線癌擴散了.
我現在又不死了.
都很神奇.
我懶得說見證.
沒有了神.
我真是一條可憐蟲.
但他說不用怕.
因為耶和華幫我們.
你明白嗎?.
人類和神失聯我們是蟲.
因為我們是第六天做.
我們是和其他動物一起做的.
你明白嗎?.
你不要以為自己天下無敵.
你以為吧.
不是的.
我們只不過是黃土地裡的小流危.
我經常說我是黃土地裡的流危.
不過神看得起我.
神和我聯想.
所以我今天可以活著.
可以站港台.
我都很吵架.
據我老婆說.
我不是很好.
你明白嗎?.
我和神失聯的時候.
我真的不是很好.
我只是蟲.
我和其他動物.

$^{681}$蟲的法則.
沒有分別.
昨天早上我和一個姊妹聊天.
她們都被人欺凌.
昨天我和我女兒聊天.
她們來我家的時候.
我又說這個問題.
都有欺凌.
欺凌那些是因為家長教他們欺凌.
那些家長不是基層.
是高深教育.
教青他們都是蟲的法則.
弱肉強食.
死命.
接著那位姊妹.
不知道該怎麼辦.
我也不知道怎麼辦.
你想想蟲能化心.
就是動物.
何事何悲.
我想說.
因為這樣.
祈禱是很強大的.
當你依靠這位動物主.
那位禱告.
是超乎你所求所想.
你想想.
如果動物主是這麼偉大.
是在於他來說.
沒有難性的事.
在人不能在他.
凡事都能.
我們不明白什麼叫凡事.
你明白嗎.
我們不明白何謂凡事.
有人問我牧師.
他老公不信耶穌怎麼辦.
我為他祈禱.
他說可以.
不是祈禱可以.

$^{721}$是他可以.
我見過太多這類個案.
因為他.
我最近在學.
祈禱功課.
我還未說太多.
因為我還未好.
我祈禱都.
有一次我入醫院.
說說祈禱.
將軍澳醫院.
男性病房.
一間房八個男人.
我因為 liver epithelial.
肝有個瘡.
像雞蛋般大.
隔壁那個因為神石.
對面那個因為大XX.
吃豬骨飯.
吃到豬骨.
那個老人家就在出院.
但又在打又在吊鹽水.
醫生說不出去打.
隔壁那些人.
跟我聊天.
喂你第一次來嗎.
我說哦.
大哥你想我幾多次.
他說我來過很多次.
哦失覺失覺你是我前輩.
他神石.
做他的餐館.
食無定時.
跟我們聊天.
聊完之後我動了時辰.
很慘那班男人.
在房間裡胡說.
我說不如我們一起祈禱.
他們說好啊.
很奇怪.

$^{761}$個個肯祈禱.
病就肯祈禱.
祈禱之後.
兩天之後隔壁那些人出院.
打電話.
前天有個傳統人跟我們祈禱.
我祈禱真的很靈.
我今日出院了.
他說祈禱靈.
不是祈禱靈是他靈.
我去殺你種.
原來住醫院也可以殺種.
因為如果我不住醫院.
我永遠不會遇到那班病友.
原來人生.
大小事情順機應適.
都有主教安排.
如果你看得通的話.
在乎你怎樣看.
所以當你有病有冤.
你不要問為什麼是我.
你問下這件事神安排有什麼意義.
意義.
等一會.
不要問為什麼是我.
為什麼你.
因為神恩代你看看你.
考驗你.
不如這樣說.
耶穌說.
是因你們信心.
少說鬼說不出來.
我實在告訴你們.
你們信心好像.
芥菜種那麼小.
就是對這座山說.
你從這邊.
落到那邊.
他也必拿去.
並且你們也門有.

$^{801}$一件不能作的事.
真的嗎.
祈禱可以延生.
填海.
不是祈禱.
是他.
你想想這個宇宙那麼大.
宇宙中的宇宙他都做得出.
有什麼會難到他.
當然要符合他的意義.
好.
看看祈禱的經文.
到那一日.
你們什麼也.
不問我了.
也就不問我了.
我實在告訴你們.
你們若向父求什麼.
他必因你的名.
賜給.
因我的名賜給你們.
向來你們沒有.
奉我的名求什麼.
如今你們求就必得著.
得著.
叫你們的喜樂可以滿足.
信不信.
真的全體父母.
如何了解那個神.
求就必得著.
叫你們的喜樂.
可以滿足.
神想我們成為一個快樂.
幸福健康的人.
他呀.
你回去慢慢想.
想了一通了解.
我都在想這個問題.
我想我.
是不是時間夠了.

$^{841}$我停一停.
我還有些話要說.
下一個午堂再說.
有沒有好笑.
等一下.
我能夠說的就這麼多.
我怕到有限.
但我肯定知道這個創造主.
偉大到無論.
他不是全能,他是超能.
看你認識他有多少.
我們祈禱.
最終我們在你面前.
我們離開你.
我們不過是蟲.
我們距離死亡.
只有一線之差.
地震告訴我們.
大自然的威力是何等大.
是無可擋.
何況你的威力.
主呀,幫我們.
讓我們望尋頂.
才能夠參透萬事.
讓我們有熟悉的眼光.
看到天外有天.
這個宇宙.
是你的作為.
以致我們能夠全然全心.
成為一個虛心的人.
全然去投靠你.
以致我們能夠.
活出基督來.
我們祈禱奉耶穌的名.
多謝各位.
(完).
\newpage



\section{申命記 15:1-18}
\label{sec:I13jDOQMkLQ}
\textbf{2025年4月6日主餐崇拜直播｜陳明泉牧師｜朝向應許之地－顧念別人｜申命記15\_1-18}
\newline
\newline
連結: \href{https://youtube.com/watch?v=I13jDOQMkLQ}{\texttt{https://youtube.com/watch?v=I13jDOQMkLQ}} ~~~~ 語音日期: 2025-04-07
\newline
\newline
\hyperref[sec:P1VP0bkSNZc]{\small{< < < PREV SERMON < < <}}
~
\hyperref[sec:index_chronic]{\small{[返順時目]}}
~
\hyperref[sec:CrZ2VP_EO18]{\small{> > > NEXT SERMON > > >}}
\newline
\newline
申命記 15:1-18
\newline
\begin{longtable}{cl}
\hline
\hline
章節 & 經文 (和合本修訂版)\\
\hline
15:1 & \begin{tabularx}{0.7\textwidth}{X} 「每七年的最後一年，你要施行豁免。 \end{tabularx} \\ \\ \relax
15:2 & \begin{tabularx}{0.7\textwidth}{X} 豁免的方式是這樣：凡債主要把手裡所借給鄰舍的全豁免，不可向鄰舍和弟兄追討，因為耶和華的豁免已經宣告了。 \end{tabularx} \\ \\ \relax
15:3 & \begin{tabularx}{0.7\textwidth}{X} 你可以向外邦人追討；但你弟兄欠你的，無論是甚麼，你都要放手豁免。 \end{tabularx} \\ \\ \relax
15:4 & \begin{tabularx}{0.7\textwidth}{X} 其實，在你中間不會有貧窮人；因為在耶和華－你神所賜你為業的地上，耶和華必大大賜福給你。 \end{tabularx} \\ \\ \relax
15:5 & \begin{tabularx}{0.7\textwidth}{X} 只要你留心聽從耶和華－你神的話，謹守遵行我今日所吩咐你這一切的命令， \end{tabularx} \\ \\ \relax
15:6 & \begin{tabularx}{0.7\textwidth}{X} 因為耶和華－你的神會照他所應許你的賜福給你，你必借給許多國家，卻不需要去借貸；你要管轄許多國家，它們卻不能管轄你。 \end{tabularx} \\ \\ \relax
15:7 & \begin{tabularx}{0.7\textwidth}{X} 「在耶和華－你神所賜給你的地上，任何一座城裡，你弟兄中若有一個貧窮人，你不可硬著心，袖手不幫助你貧窮的弟兄。 \end{tabularx} \\ \\ \relax
15:8 & \begin{tabularx}{0.7\textwidth}{X} 你總要伸手幫助他，照他所缺乏的借給他，補他的不足。 \end{tabularx} \\ \\ \relax
15:9 & \begin{tabularx}{0.7\textwidth}{X} 你要謹慎，不可心起惡念，說：『第七年的豁免年快到了』，你就冷眼看你貧窮的弟兄，甚麼都不給他。他若為你的緣故求告耶和華，你就有罪了。 \end{tabularx} \\ \\ \relax
15:10 & \begin{tabularx}{0.7\textwidth}{X} 你要慷慨解囊，給他的時候不要心疼，因為耶和華－你的神必為這事，在你一切的工作上和你手所做的一切賜福給你。 \end{tabularx} \\ \\ \relax
15:11 & \begin{tabularx}{0.7\textwidth}{X} 因為地上的貧窮人永遠不會斷絕，所以我吩咐你說：『總要伸手幫助你地上困苦貧窮的弟兄。』」 \end{tabularx} \\ \\ \relax
15:12 & \begin{tabularx}{0.7\textwidth}{X} 「你弟兄中，若有一個希伯來男人或希伯來女人賣給你，已服事你六年，到了第七年就要讓他自由離開你。 \end{tabularx} \\ \\ \relax
15:13 & \begin{tabularx}{0.7\textwidth}{X} 你讓他自由離開的時候，不可讓他空手而去， \end{tabularx} \\ \\ \relax
15:14 & \begin{tabularx}{0.7\textwidth}{X} 要從你的羊群、禾場、壓酒池中取一些，慷慨地送給他；耶和華－你的神怎樣賜福給你，你也要照樣給他。 \end{tabularx} \\ \\ \relax
15:15 & \begin{tabularx}{0.7\textwidth}{X} 要記得你在埃及地作過奴僕，耶和華－你的神救贖了你。為此，我今日將這事吩咐你。 \end{tabularx} \\ \\ \relax
15:16 & \begin{tabularx}{0.7\textwidth}{X} 他若對你說：『我不願意離開你』，因為他愛你和你的家，並且他在你那裡很好， \end{tabularx} \\ \\ \relax
15:17 & \begin{tabularx}{0.7\textwidth}{X} 你要拿錐子在門上穿透他的耳朵，他就永遠成為你的奴僕了。你待婢女也要這樣。 \end{tabularx} \\ \\ \relax
15:18 & \begin{tabularx}{0.7\textwidth}{X} 你讓他從你那裡自由離開的時候，不要看作困難，因為他已服事你六年，相當於雇工雙倍的工錢。這樣，耶和華－你的神必在你所做的一切事上賜福給你。」 \end{tabularx} \\ \\
[1ex]
\hline
\hline
\end{longtable}
$^{1}$我們今天要聽神的話是記載在.
《申命記》第十五章一至十八節.
《舊約聖經》二百三十三至二百三十四節.
《申命記》第十五章第一節.
「每逢七年末一年.
你要施行豁免.
豁免的定例乃是這樣.
凡債主要把他所借給鄰舍的豁免了.
不可向鄰舍和弟兄追討.
因為耶和華的豁免已經宣告了.
若借給外邦人.
你可以向他追討.
但借給你弟兄.
無論是什麼.
你要鬆手豁免了.
你若留意聽從耶和華.
你神的話.
謹守遵行.
我今日所吩咐你一切的命令.
就必在你們中間沒有窮人了.
在耶和華你神所賜你為業的地上.
耶和華必大大賜福於你.
因為耶和華你的神.
必照他所應許你的賜福於你.
你必借給許多國民.
卻不似向他們借貸.
你必管轄許多國民.
他們卻不能管轄你.
在耶和華你神所賜你的地上.
無論那一座城裡.
你弟兄中若有一個窮人.
你不可忍著心.
撥著手不幫補你貧乏的弟兄.
總要向他鬆開手.
照他所缺乏的借給他.
補他的不足.
你要謹慎.
不可以心裡起惡念.
說.
「捨七年的豁免年快到了.

$^{41}$你便可.
烏眼向你窮乏的弟兄.
甚麼都不給他.
以至他因你求告耶和華.
罪必歸於你了」.
你總要給他.
給他的時候.
心裡不可受煩.
因耶和華你的神.
必在你這一切所行的.
並你手裡所辦的事上.
賜福於你.
原來那地上的窮人永不斷絕.
所以我吩咐你說.
總要向你地上.
困乏窮乏的弟兄鬆開手.
第十二節.
你弟兄中若有一個希伯來男人.
或希伯來女人被賣給你.
服侍你六年.
到第七年就要任他自由出去.
你任他自由的時候.
不可視得空手而去.
要從你羊群.
和場.
走廊中多多給他.
耶和華你的神怎樣賜福於你.
你也要照樣給他.
要紀念他.
紀念你在埃及地作過奴僕.
耶和華你的神將救贖.
因此我今日吩咐你這事.
他若對你說.
「我不願意離開你」.
是因他愛你和你的家.
且因在你那裡很好.
你就要拿錐子將他的耳朵.
在門上刺透.
他便永為你的奴僕了.
你代婢女也要這樣.

$^{81}$你任他自由的時候.
不可以難事.
因他服侍你六年.
交給故宮的公家.
多加一倍了.
耶和華你的神.
就必在你所作的一切事上.
祝福你.
聖經獨筆.
(聖經獨筆).
 各位姐妹平安.
(各位姐妹平安).
 我們今天一齊讀新命記.
在我們讀這段經文之前.
近來大家看新聞.
會發現全世界轟動的一個消息.
就是美國要全世界納關稅.
我不是很經濟地分析.
不過我對這段新聞很有興趣的.
就是一個鏡頭.
美國總統說關稅的原因就是.
他經常說一個字.
就是全世界欠了他.
他看美國和外國的貿易逆差.
就是外國人在美國騙了他們的錢.
所以他們就要收回這個對等關稅.
以致他們被欺騙.
花出去的錢就可以藉著這個關稅.
討回一個公道.
再加上他說這個已經是一個很仁慈的關稅.
所以就說30\%多.
去到最高的可能是幾十\%.
還有一些花邊新聞就是在說.
有一個島只有我住的.
那裡也收關稅.
但是有些國家.
他們視為敵對的國家.
也沒有收關稅.
其實整件事是很有趣的.
不過我很想特別針對的就是.

$^{121}$這個總統他裡面經常說.
就是全世界欠了他.
他要追回這些錢.
如果你在新聞裡面每天看著.
這一類的論述或者這種語調.
其實你的心靈會有時候很不平衡的.
你會經常覺得身邊所有人都欠了你.
欠了你.
如果以緣木厭進入你的內心當中.
有時候你的價值觀都會受到扭曲.
心理學裡面有一種說法.
就叫做索債心理.
索債心理其實就是覺得別人欠了自己.
在心理學家裡面說到這個索債心理.
有外顯和內在的.
外顯就會經常拐口唇邊的.
就經常說還錢.
你欠了我的.
內斂的,內隱形的那些.
就會整個心裡面默默的計算.
就是在你的心裡面.
你又做了一些事情虧欠了我.
就是找一本小喜簿.
記錄了所有別人對自己的那種虧欠.
你不要輕視這種心態.
如果我們久而久之.
我們不面對這種價值觀念.
不論是世界的政局也好.
或者個人的內心心理也好.
如果經常覺得有一種別人欠自己的時候.
你做人是覺得不開心.
甚至你會經常覺得.
對於身邊的人.
就會處於一種根根計較的狀況裡面.
為什麼今天我要用這個例子呢?.
因為今天我們所看的那段經文.
正正就是要面對人性裡面的那種根根計較.
今天特別講的兩個重點.
一個就是豁免年.
就是說借貸的債主和債制之間的關係.

$^{161}$第二個就是勞瀑關係.
其實也是承接著剛才講的豁免年.
不過一個是欠錢.
而後半段就是沒有錢還.
就要在別人家裡做奴隸.
究竟聖經裡面摩西是一個.
那時候還沒有很多貧富懸殊.
或者在經濟狀況和今天沒有那麼複雜的時候.
他已經講定一個上帝的心意.
這一種可能對我們今天來說.
我們要明白上帝的心腸.
以致我們轉化使用在我們的生命當中.
究竟這段經文講了什麼呢?.
我們先祈禱.
求聖靈指教我們.
讓我們能夠明白祂的話語.
同心祈禱.
天父上帝願你在我們當中帶領.
賜福給我們眾弟姐妹.
願你的話語成為我們生命的幫助.
求你在我們當中.
你的聖靈指教.
讓每一位弟兄姊妹讀這段聖經的時候.
我們不僅讀得明白.
更加是讓你的心腸.
成為我們行事為人的準則.
又求你幫助宣講的講得清楚.
將你的話語如實地帶給弟兄姊妹.
讓我們在這個世代裡面.
成為你所喜悅的子民.
恭敬將來的時間交託.
願你賜福.
獻上我們禱告奉靠基督耶穌得勝的聖名.
Amen.
好 我們一起看《新命紀》.
在我們講這段經文之前.
有一點補充的.
由《新命紀》第十二章開始.
其實這是一系列的講章.
去到第十五節.

$^{201}$或者由第十二章.
十三章 十四章 十五章.
甚至第十六章.
其實講的很多就是由一些生活細節.
生活裡面出現的處境.
由下而上.
將這些生活處境.
再去推翻上帝的那種心意.
所以你會看到.
其實你一直讀.
在《新命紀》裡面讀的時候.
他會將生活裡面不同的處境.
他就會講出來.
不過摩西他不是純粹講給你聽.
你怎麼做.
他會透過這些生活場景.
去教導那些子民.
怎樣去面對.
而面對的那些原則.
就是來自上帝頒佈的律法和吩咐.
所以在這裡.
其實是很完備的.
他舉例在生活裡面.
面對這些狀況.
舉一反三的.
來體會上帝.
律法的心意和那種心腸.
我們看到第十五章這裡.
在第十五章裡面.
第十一節他特別講到一個很重要的字眼.
原來那地上的窮人永不斷絕.
所以我吩咐你說.
總要向你地上困苦.
窮乏的人的弟兄.
送開手.
在這裡裡面.
十一節那裡.
我特別將這句放上來的.
原因就是.
上帝講了這個現況.

$^{241}$世界上貧窮是沒辦法完全杜絕的.
貧窮是一個相對的概念.
你有錢一點.
那個沒有你這麼有錢.
相對來說他就叫做相對貧窮了.
當然也有一些的.
就有一些叫做絕對貧窮.
可能就是沒有三將.
基本需要都不能夠滿足的.
不過在聖經裡面.
我們幫的.
或者面對著那個社會.
窮乏或者貧窮人.
在現有的制度.
或者在上帝的眼光裡面.
你沒辦法.
完全杜絕了這一種的窮人.
多麼努力都好.
總會有一些生活條件.
比你們弱的一些的人.
面對著這種這樣的狀況.
摩西吩咐那些的以色列民.
你要對著這些困苦窮乏的弟兄.
送開手.
他特別留意到.
他講到的就是.
他有分遠和近的.
外邦人的對待方式.
跟你以色列的那些弟兄的對待方式是不同的.
這裡特別講的是對於你自己的弟兄姊妹.
你自己同邦國的人.
對於同邦國以色列民族的人.
這些窮乏的人.
你一定要鬆開手的.
和合本譯這個鬆開手的譯得挺好的.
它的原文的意思就是張開雙臂.
張開雙臂就拿那個意象就是.
你就要不要跟跟計較.
你就要真的用你所擁有的.
來幫助他們.

$^{281}$在這裡面一個很重要的處境就是講到.
世界上所有貧窮人.
無處不在.
所以你面對著這些.
跟你有關係的這些同族.
同民族的人.
跟你血統關係的.
或者你的鄰舍你的弟兄.
你不可以暫不知.
你不可以忽視他.
反而你要對他們就鬆手了.
如果用這個原則下去.
然後你再看.
從第一節讀到第二節.
這個就講到豁免年的做法.
每逢七年的末一年.
你要施行豁免.
豁免的定例乃是這樣.
凡債主要把所借給鄰舍的豁免了.
不可向鄰舍和弟兄追討.
因為耶和華的豁免年已經宣告了.
其實這個豁免年很簡單的.
那個人欠你債.
最多還債期是六年.
第七年你就免了他.
你就免他債了.
無論欠多少都好.
這個債不可以終身的帶下去.
最多只是帶六年.
到了第七年.
這個就是在講七年的末一年.
就是第七年.
你就一定要施行豁免.
不再追討了.
在這裡背後就是在講.
一個人的那種.
可能偶爾面對著生活有些困難.
或者天災人禍.
又或者家裡有一些重要成員過世.
他們可能要投靠一些弟兄.

$^{321}$或者投靠他的鄰舍.
當他投靠的時候.
他幫你做工.
他沒錢還.
於是就去打工.
或者他幫你做事.
或者你借了一些錢給他.
一路一路退回來還給你的時候.
最多的年期.
他只能夠還六年.
你就收他六年的錢.
就不可以再收.
或者他終身.
他還不了就叫他的子女還.
沒有的.
只是最多就是用六年.
經文裡面用這個字眼.
或者講這句話.
他講到的規格是很高的.
首先就講到豁免的定例.
這個定例是誰定的呢.
不是約定俗成.
不是社會裡面定的.
在這裡裡面就講到.
是耶和華的豁免年已經宣告了.
在講的這個定例.
是由上帝自己親自頒佈的.
這個宣告了這個字.
在聖經裡面.
最早出現的就是講到.
上帝創造宇宙萬物.
神說.
神說那個字就是神宣告了.
天地宇宙萬物.
用這種權柄.
來創造出來.
同樣這個豁免年.
這個位置.
是在講上帝帶著一個這樣的權柄.
來講這群以色列民必須要執行.

$^{361}$所以這種的豁免.
或者對於窮人的那種鬆手.
這個是上帝的命令.
這個是上帝的應許.
你不可以說.
我心情好一點就去這樣做.
心情沒那麼好就不去這樣做.
這條界線是不能移動的.
因為這個是上帝所定下的命令.
所以那些以色列民.
或者摩西呈現出來的.
要做一個這樣的豁免年.
上帝定下了就是這樣的了.
如果你.
你是精明的那些債主就說.
那不要緊.
那我就分期付款.
就是你欠我的錢分六年還.
六年最後.
就是最後那一天.
12月31日.
這樣就剛剛好還清.
這樣大家就最平衡了.
是不是?.
就是最公道.
不過呢.
現實有時候不是你說.
不像我們今天這樣供樓的.
我們今天有月薪的.
我們很固定的那種人工.
但是在摩西的年代.
如果你看一個遊牧民族.
其實他們的出產.
他們的財產.
就視乎那個天氣.
或者生活裡面那種.
是不是風調雨順等等.
要看天造人的.
他們未必是可以固定的來還錢的.
正正就是因為一種這樣的那種.

$^{401}$有時候生活不是很穩定的那種狀態.
上帝要求那些債主.
即是說如果到最後他沒辦法.
在六年裡面清還所有的那種欠款.
第七年.
你都要算的了.
他要求的就是債主.
要有一種胸襟.
要有憐憫的心.
來對待那些窮乏的.
或者沒能力的那種的人.
除了一個硬梆梆.
一個宣告出來的一個命令之外.
由第三節到第六節.
摩西就將那個焦點.
又轉了另一個角度來說.
若借給外邦人.
你可以向他追討.
但借給你弟兄.
無論是什麼.
你要鬆手豁免了.
你若留意聽從.
耶和華你神的話.
謹守遵行我今日所吩咐你的一切命令.
就必在你們中間沒有窮人了.
在耶和華你神.
所賜給你遺孽的地上.
耶和華必大大賜福允你.
因為耶和華你的神.
必照他所應許的.
應許你的賜福允你.
你必借給許多國民.
卻不至向他們借貸.
你必管轄許多國民.
他們卻不能管轄你.
在第三至第六節裡面.
摩西就用一個.
上帝應許的角度來說.
豁免或借貸的方式.
如果你記得剛才說到.

$^{441}$地上總會有窮人.
這是全世界的普遍現象.
但如果以色列民自己對弟兄.
對自己的國民.
你願意鬆手的時候.
在經文裡面特別說到.
你們中間就沒有窮人了.
即是說如果我們自己.
自己人互相照顧.
我們就能夠互相得益.
以至以色列民族自己內部.
就可以成為一個沒有窮人的民族.
外邦人可能繼續有.
你借給他們可以繼續超過七年.
不過在內部裡面.
因為你互相照顧.
因為你互相體恤.
自己民族裡面人的需要.
所以這種貧窮的現象.
如果大家都樂意這樣做的時候.
這個民族會變成一個富庶的民族.
成為一個大大自福.
得著福氣的民族.
值得注意的就是.
在這裡特別將上帝所說的福氣.
用一種很具體的方式.
去告訴以色列民智.
上帝是一個應許之地.
將一個地賜給以色列人.
同時上帝不是只給了你進去.
在他們生活當中.
上帝會繼續將福氣.
托著他們這個民族.
一直將福氣賜給他們.
這是從天而下的那種福氣.
不過這個福氣不是純粹從天而落地.
有時候也是人為操作.
如果人與人之間.
沒有一種根根計較的觀念.
弟兄姊妹之間.

$^{481}$或者民族自己之間.
都可以成為一個互相祝福的途徑.
所以在這裡他說到.
他用一個應許的角度.
既是這種願意幫助人的.
就會得到上帝的賜福.
同時在人間的制度裡面.
因為人與人之間互相幫忙.
那種貧窮就減低.
他未必叫他關懷到全世界.
未去到很大愛.
不過至少在他們自己的民族裡面.
他們就能夠經歷到那種小康.
或者那種沒有貧窮人的願景.
所以在這裡.
最後第六節.
上帝會將應許的福賜給你.
你就可以將來有餘.
你有能力借給別人.
那你就可以借給那些外邦人.
而那些外邦人到最後.
就成為你所管轄的對象.
又不會反過來.
因為你們經常不肯借.
然後就外邦人借錢給你們自己窮乏的人.
到外邦人管轄你自己.
這個就是講到那個.
你當是與大相關.
但也講到一個很現實的現象.
自己的弟兄姊妹.
你一定要幫助他.
你不可以當他是外人.
否則他的需要就找了外邦人的時候.
就由外邦人來管轄自己的弟兄.
所以你就會令到你們自己.
處於一種困難的狀況.
所以這個應許其實挺複雜的.
不過你讀上去就知道.
其實就是要願意.
躺開你的心去幫助弟兄姊妹之間.

$^{521}$或者在他的民族裡面互相幫助的時候.
你就能夠令到整個民族都可以富強.
好了,到了第七節到第九節.
除了講到剛才應許和宣告之外.
接著他就要針對人的人性.
第七節.
「在耶和華你神所賜給你的地上.
無論那一座城裡.
你弟兄若有一個窮人.
你不可忍著心.
撥著手不幫補你窮乏的弟兄.
總要他鬆開手.
照他所缺乏的謝給他.
補他的不足.
你要謹慎不可心裡起惡念.
說:第七年的豁免年快到了.
你便惡眼看你窮乏的弟兄.
什麼都不給他.
以致他因你求高耶和華.
罪便歸於你了」.
這裡特別一個重要的字眼.
就是起惡念.
起什麼惡念呢?.
只是借貸而已.
你借了給他已經很好了.
現在叫他還錢而已.
一借一還.
這也是一個很公道的行為.
不過這裡他特別講到債主.
他怎樣對待債主呢?.
第七年快到了.
準備要豁免了.
在第五年半已經開始眼紅他了.
他還不了.
還錢還錢.
現在還是第五年半而已.
第六年還沒過.
快點揍他.
這是第一個惡念.
就是用盡你的權力去捏住別人.

$^{561}$撥著手.
即是說殺手不管.
不管了.
他有需要你不管.
債主在你眼中就是還錢還錢還錢.
S兩棟.
看到他就是這樣的角度.
第二件事就是.
可能到了第七年差不多到的時候.
你就刻意惡待債主.
刻意的.
甚至是剝削他.
希望在豁免年來到之前.
你就拿到你盡量獲取你最大的.
你借出去的那些錢或者那些利益.
倘若債主是沒有錢還.
他在你的家裡成為了你的奴僕.
那所有事情就叫他做.
所有最重擔的.
最下欄的功夫.
或者不停在做.
連睡覺都不讓他睡覺.
當他一個機器來使用的時候.
這就是惡念的意思.
你行使盡你作為債主的權柄.
剝削那些欠你債的那些人.
有一本書.
很出名的一本小小的詩集.
作者是一個外國人叫Martin Buber.
那本書叫做I vowed.
即是我與你.
這本書好像用詩的方式來寫.
這本書裡面說的是什麼呢.
他說原來全世界人與人之間的關係.
就是一個我與你的關係.
很平常而已.
我與你的基礎在於.
我看你和我是平等的.
即是說我的收入比你多一點.
或者我是債主你是債人.

$^{601}$但是在人與人之間的交往裡面.
我會當你是一個人來看待.
你和我之間是有來有往是平等的.
雖然經濟或者你的生活背景有些不同.
但是不會構成了.
就因為我有錢一點.
我就高人一等我看不起你.
我與你這本書裡面說到.
如果一個平等尊重的關係.
無論你處於什麼狀況.
你對人就是一種平等的關係.
相反來說.
如果我看那個人是比我低下的.
我就不是純粹一個我與你的關係.
而是一個我與他的關係.
那個他是死物的他.
it.
I it的關係.
這個I it的關係就是說.
我就以我為中心.
你要滿足我的需要.
到最後是說將其他人看不起.
不尊重他就將他物化了.
即是視為死物來看待身邊的人.
這本書的作者他說到.
這種我與你或者我與他.
不單只是說一個人對身邊的人的那種關係.
當我與你是這樣平等的時候.
同樣你也會對我是一個平等的那種看法.
這個是很對等的.
但當我物化別人.
我看別人是一個不尊重人的工具.
一個工具人這樣看待身邊的人.
其實同樣的思維都是別人這樣看待我.
人與人之間的關係就變成了一個工具.
效益的層次.
用這本書的角度來看你就明白了.
如果債主他覺得自己是至高一等.
那個債主是無權無勢的.
你不尊重他.

$^{641}$當他是一個搖錢樹.
是一個物化了的工具人.
這樣來看待這個奴隸的時候.
你不尊重他.
到最後其實在他的心靈當中.
他也不會尊重你.
而當他不尊重你的時候.
他是求高惹禍.
說這種對人不尊重的罪.
就會加諸在你的身上.
所以在這裡面.
這個起惡念.
摩西他特別說到.
其實去到最核心就是.
你對人尊不尊重.
那個人不是欠了你錢.
但不代表他欠了你錢.
你就不尊重他.
我相信黃枕頂子會不會.
不過也要說的.
可能我在公司裡面.
我們是主管.
我們下面有一些人幫我工作的.
在制度上我是分配工作.
我是主管.
但不是在本質上.
因為我是主管.
我就看不起我的下屬.
我都當他是一個平等的人來看待.
又或者家裡幫你工作的.
我身上的印傭姐姐.
不同的姐姐.
她來你家裡打工.
她不是你的奴隸.
你就要尊重她.
這種就是對人的那種.
你有沒有一種這樣的心腸.
如果你對人不尊重.
用你的經濟條件來壓榨.
或者剝削他的時候.

$^{681}$他求高惹禍說.
你的心裡面就陷在罪孽當中.
你的心裡面因為你看不起別人.
而導致你生命裡面陷在罪孽當中.
所以在經文裡面.
第七至第九節.
你就要檢視.
究竟你的心是一個怎樣的人.
去到第十至第十一節.
也是豁免年的最後一個小段落.
你總要給他.
給他的時候心裡不可受煩.
因耶和華你的神.
必在你這一切所行的.
並你手裡所辦的事上賜福願你.
原來那地上的窮人永不斷絕.
所以我吩咐你說.
我要向你地上困苦窮乏的弟兄鬆開手.
剛才說到.
你不要心裡面帶有惡念.
看不起別人.
去到這裡.
你就要學習繼續用鬆開手這個字.
其實鬆開手這個字.
不斷又不斷重複很多次.
不過這個鬆開手不斷重複的意思.
除了是在說語色之外.
或者那種勸告之外.
我自己去體會的就是.
這個鬆開手也成為了我們今天生活的形態.
你怎樣對身邊的那些欠債的人.
其實跟你平時怎樣對人那種人際關係.
是一致的.
你不會說我對債務很寬容.
但我平時對人很敏感的.
不是這樣的.
是這一種由生活細節所有的東西推開.
由你對身邊的人那種尊重.
那種鬆開手.
即使沒有欠與還的那種關係.

$^{721}$但平時對人你也有一種慷慨.
那種願意鬆開手的心腸.
這些平時訓練了你.
你有這種視野.
以至到你就可以將這種視野.
即使別人欠你債的時候.
你也用這種對等的方式.
對人尊重的方式來實踐.
在豁免年這麼重要的吩咐上.
我們需要練習的.
我們需要對身邊.
善待身邊的人.
我們不要經常覺得有一種虧欠.
或者我幫助人那種高高在上.
在經文裡面說到.
借東西給別人.
或者你對人身邊.
很願意鬆開手的這種生活形態.
你不要覺得很煩.
不要經常覺得很混亂.
反而是在你做的事情上.
如果你能夠帶著一份這樣廣闊的胸襟.
上帝一定賜福給你的.
摩西其實就是用這種角度.
來勉勵那些以色列子民.
今天我們有時候會斤斤計較.
為什麼呢?.
我們怕沒有.
我們怕撥了它.
就被人家豪華了.
我們又或者怕借錢不還.
令我們有一種感覺.
不事未宜.
也會有的.
可能世界上也會有些人.
利用了人的愛心.
來製造了這樣的狀況.
不過我們的生活裡面.
我們真的有時候要想一想.
我們會不會因為有很多失敗的經驗.

$^{761}$就令到我們的心變得很麻木.
對於身邊的需要.
變得很麻木不仁呢?.
這裡要很小心說.
如果有電話打給你.
忽然間說那個人是你某個親戚.
聲稱你都要查證一下.
現在很多詐騙的.
那種善心.
那種願意鬆開手.
不要成為了一種.
不受完全查證的.
那種亂來的輕率.
我們是要小心查證.
他真的知道那個人有需要.
你認識他的.
他當面的來問你.
你真的不可以硬著心的.
但是如果有個電話來.
你真的要查證.
現在太多電騙了.
不過正正有很多.
今天有很多欺騙的事件出現.
我們索性硬著心.
所有的需要我們都完全視而不見.
完全不聽不聞.
這不是聖經對我們的教導.
反而是我們要學習.
在這些困難當中.
我們會不會願意鬆開手.
躺開我們的心.
我看過有些人物傳記.
不說名字.
不過大家會知道.
在一個坊間裡面.
有一個上海的梟雄.
他臨終的時候.
梟雄可能有很多朋友.
也有很多牙齒印.
在他臨死之前.

$^{801}$在遺囑裡面加了一個註腳.
他就說在他死了之後.
就叫他的後人在甲萬裡面.
把甲萬裡面所有的借據.
全部燒掉.
然後有些人說.
你借了很多錢出去.
應該全部燒掉.
梟雄就說.
如果他有心還的.
沒有借據他都會還給你.
所以這個借據是沒有意思的.
但是如果那些人沒有心還.
他的後人強行去追那些錢.
就會有機會朝野殺身之禍.
又或者因為這樣去追討的時候.
可能為他們帶來更加多的困難.
這些是梟雄的價值觀.
不過有時候都是.
我們心裡面有沒有一份的胸襟.
我們會不會經常捏著.
還錢還錢.
我將人成為了S兩棟.
或者是我對人願意的.
有一份的尊重.
甚至乎我經常操練自己.
我們的心是一個慷慨就義的人.
我們繼續看.
由第十二節到第十八節.
是最後一個段落.
李大興中藥有一個希伯來男人.
說希伯來女人被賣給你.
服侍你六年.
到第七年就要任他自由出去.
你任他自由的時候.
不可使他空手而去.
要從你的羊群和牆走渣之中.
多多地給他.
耶和華你的神怎樣賜福願你.
你也要照樣給他.

$^{841}$要紀念你在阿吉地作過奴僕.
耶和華你的神將你救贖.
因此我今日吩咐你這件事.
他若對你說.
我不願意離開你是因他愛你和你的家.
且因你那裡很好.
你就要拿錐子.
將他的耳朵在門上刺透.
他便永為你的奴僕了.
你代婢女也要這樣.
你任他自由的時候不可為難事.
因他服侍你六年.
較比故宮的功架多加一倍了.
耶和華你的神.
就必在你所作的一切事上賜福願你.
第十二至第十八節.
其實是一個再進一步的仔細引申.
剛才的豁免年.
不過這裡是說人力資源.
他特別強調的就是.
希伯來男人或希伯來女人.
即是那些是弟兄.
那些是姐妹.
是同族的人.
如果因為他沒有能力.
他要出賣自己的奴隸.
他便將自己賣給你.
以前的奴隸.
埃及地做奴隸.
以色列人賣給埃及人是終身的.
是一個終身奴隸.
不過在這裡.
摩西班牟這麼重要的定例.
就是說他賣.
是賣六年給你的.
是六年的奴隸.
在你的家裡居住做工作.
第七年你就要放他走.
任他自由.
當他走的時候.

$^{881}$他便自由離開.
就這樣.
經文第十四節說到.
你要將羊群和長走渣之中.
多多給他.
即是說他做了那六年苦工.
他不會只是做.
你要帶著感恩之情.
來有一些的饋贈.
給他來迎向自由.
如果他自由了.
什麼都沒有.
到最後他便.
由A家自由了.
便去B家做奴隸.
繼續.
整個制度就是希望.
整個國幫能夠邁向小康.
經歷過最困難的日子.
有人幫助之後.
他便可以重新起來.
所以在這裡.
那個自由不是純粹給他最大的恩典.
而是你要祝福他的自由.
你要給一些你的羊群和長走渣.
裡面的農作物給他.
你這樣做的時候.
上帝會賜福給你的.
上帝怎樣去.
昔日怎樣幫助你.
你就帶著這一份感恩的心.
來幫助身邊的人.
摩西再繼續說.
你們以前在埃及也做過奴僕.
你做過這樣的辛苦.
你就不要將以前你所經歷的苦.
套用在今天你自己做了債主的時候.
你就這樣去壓削那些奴隸.
反而你自己經歷過這些痛苦的時候.
上帝怎樣幫助你.

$^{921}$你就幫助身邊的弟兄.
或者這些弟女姐妹.
所以在這裡.
其實他要訴諸於的.
不是純粹說那些外在的規條.
他這裡更加內心的那種動力.
你也做過這種苦況.
你也做過奴僕.
上帝怎樣幫助你.
你帶著這一份感恩的心.
來去善待那些在你家打了六年工的奴僕.
不過經文裡面再繼續說到.
有另外一個.
有些奴僕原來覺得自己未必想自由.
經文裡面再繼續說.
十六節如果他跟你說他不願意離開.
不想走.
是因為你對他的愛.
你是一個好的僱主.
是一個好的債主.
令到他不願意離開.
因為他很愛你.
結果是怎樣呢.
要做一個記號.
姐妹會做的.
弟兄也有做.
就是在門上拿個鎚.
就在門的耳珠位置.
一下子打下去有個洞.
穿耳.
現在穿耳是為了戴耳環漂亮.
以前穿耳代表著.
我永遠在這個家裡面.
我是屬於這個家的.
我感恩這個恩主.
來去善待自己.
在困難當中.
以至他終身在這個家庭裡面.
男的也是這樣.
婢女也是這樣.

$^{961}$你要用這個方式來做一個記號.
當你善待人.
人家會死心塌地.
願意為你.
這個舉動.
用鎚子釘耳洞.
好像是一個很危險的記號.
也是一個.
你看起來也是很血腥.
很血淋淋的一個代價.
但是一個這樣的舉動.
那些人都願意去做.
可想而知.
你善待人.
人家就會很願意用一生.
和用一個這麼危險.
這麼血淋淋的記號.
來代表著跟隨這個家主.
終身來服侍他.
今天我們願不願意善待人.
關鍵在於.
我們的心有沒有感恩之心.
如果我們覺得.
所有事情都是靠我自己努力得來.
我承受著.
那件事是我的.
沒有想過.
上帝有很多福氣在我們生命當中.
我們就沒有辦法.
有這種心腸.
來對待身邊的人.
我想起一個故事.
這是我的長輩告訴我的.
是我爸爸的一個故事.
我爸爸其實是.
以前在大陸偷渡下來的.
他整個身上.
只有三個膠袋和一條泳褲.
他跟我說.
為什麼有三個膠袋呢.

$^{1001}$因為他是游泳下來的.
就是從以前的保安游泳.
游去流浮山.
那些膠袋就吹脹了他.
讓自己有一點浮力.
游泳的時候不用這麼多力.
就放在泳褲上.
他就是這樣來的.
就來到香港.
後來就做製衣.
他以前是不會做的.
由零開始跟別人學師做製衣.
然後慢慢成為了.
一個山寨廠的老闆.
然後就開始收徒弟.
我聽我的長輩.
我爸爸以前也有跟我說.
其實我們都是拆手空拳下來.
只有一條游泳褲.
所以什麼都沒有.
就養大你們了.
然後就成家.
他特別說到.
以前在工廠裡面.
有些哥哥在我那間山寨廠做.
原來那些哥哥是什麼人呢.
原來是一些遠房親戚.
就是那些兒子讀不了書.
有些在長洲住的.
讀不了書.
就一群結隊.
出來社會做事又要給保贖金.
很難的.
以前出來香港社會.
於是就跑來.
去我爸爸做的山寨廠裡面做學師.
所以我爸爸就從零開始教他.
怎樣裁衣服,車衣服,下欄功夫.
什麼都做了.
沒有薪金的.

$^{1041}$或者後來他一直成熟一點.
就有一點點收入給他.
最初的時候包吃包住.
在工廠裡睡覺.
小時候我跟他們一起睡.
就是這樣來到生活.
有一天我記得那個哥哥.
他說我要走了.
他叫我爸爸的師父陳生.
他說陳生我學滿師了.
我想出去開檔.
如果是你.
以前有些訓練熟手技工.
又正直壯年.
他忽然間要走.
都會很不小心的.
我記得我家人跟我說.
我爸爸就說.
對的.
你長大了都要出去闖一些天下.
我爸爸就出了一份比較厚的薪金給他.
就介紹了以前給我們做的客人.
就介紹給他.
他說你可以做.
我介紹客人給你.
不過你不要跟我競爭.
但是他祝福他.
他自己繼續開廠.
自己開一片天地.
我爸爸都過了身三十多年.
媽媽都過了身幾年.
每年過節過節.
這些哥哥都會問候我.
會一次說這些經歷.
他自己的眼睛都會有眼淚.
我每一次聽的時候.
我會覺得很自豪.
我家原來不是一個.
斤斤計較的家庭長大.
後來聽我家人說.

$^{1081}$我爸爸都是穿著游水褲來到香港.
他怎樣經歷這個香港發展的浪.
同樣他都想祝福他的徒弟.
經歷這個浪.
他有自己的事業.
能夠自己可以養妻活兒.
今日在我們的生命當中.
我們會不會很多時候有斤斤計較.
如果我們今日能夠學習鬆開手.
未必是對待身邊的寨子.
但是如果我們由生活.
對待人接物開始.
我們能夠鬆開手.
放開我們的胸襟.
尊重他.
善待身邊那些經濟條件.
美利沒那麼好的人.
上帝會賜福的.
上帝的福是怎樣我不知道.
但是我可以肯定.
願意施予的人.
心靈是闊綽.
是綽綽有餘.
生命不會覺得很多人謙虛.
或者生命是不會貴乏.
反而經常斤斤計較的人.
生命就會落入一種困乏的境地.
他怎樣經歷.
怎樣大魚大肉.
他都無法經歷到豐盛的生命.
今日摩西用這個豁免年.
用這個不要令到那些奴僕空手而去.
調教我們今日的價值觀念.
求主繼續藉著這段經文.
在這個世代有很多困難.
經濟不景的年代裡.
上帝繼續對我們說話.
\newpage



\section{申命記 16:1-17}
\label{sec:CrZ2VP_EO18}
\textbf{2025年4月13日崇拜直播｜陳冠成牧師｜朝向應許之地－循環生活中記念上帝｜申命記16\_1-17}
\newline
\newline
連結: \href{https://youtube.com/watch?v=CrZ2VP_EO18}{\texttt{https://youtube.com/watch?v=CrZ2VP\_EO18}} ~~~~ 語音日期: 2025-04-13
\newline
\newline
\hyperref[sec:I13jDOQMkLQ]{\small{< < < PREV SERMON < < <}}
~
\hyperref[sec:index_chronic]{\small{[返順時目]}}
~
\hyperref[sec:96uE5WjMIVg]{\small{> > > NEXT SERMON > > >}}
\newline
\newline
申命記 16:1-17
\newline
\begin{longtable}{cl}
\hline
\hline
章節 & 經文 (和合本修訂版)\\
\hline
16:1 & \begin{tabularx}{0.7\textwidth}{X} 「你要守亞筆月，向耶和華－你的神守逾越節，因為在亞筆月，耶和華－你的神在夜間領你出埃及。 \end{tabularx} \\ \\ \relax
16:2 & \begin{tabularx}{0.7\textwidth}{X} 你當在那裡，耶和華選擇作為他名居所的地方，從羊群牛群中，將逾越節的祭牲獻給耶和華－你的神。 \end{tabularx} \\ \\ \relax
16:3 & \begin{tabularx}{0.7\textwidth}{X} 這祭牲不可和有酵的東西一起吃。因為你曾匆忙離開埃及地，你要吃無酵餅，就是困苦餅七日，好讓你一生的年日記得你從埃及地出來的那一日。 \end{tabularx} \\ \\ \relax
16:4 & \begin{tabularx}{0.7\textwidth}{X} 在你全境內，七日不可見到酵母。第一日晚上所獻的肉，一點也不可留到早晨。 \end{tabularx} \\ \\ \relax
16:5 & \begin{tabularx}{0.7\textwidth}{X} 你不可在耶和華－你神所賜的各城中，任何一座城裡，獻逾越節的祭， \end{tabularx} \\ \\ \relax
16:6 & \begin{tabularx}{0.7\textwidth}{X} 只可在那裡，耶和華－你神選擇作為他名居所的地方，在晚上日落的時候，就是你出埃及的時候，獻逾越節的祭。 \end{tabularx} \\ \\ \relax
16:7 & \begin{tabularx}{0.7\textwidth}{X} 你要在耶和華－你神所選擇的地方把肉烤來吃，次日早晨就回到你的帳棚去。 \end{tabularx} \\ \\ \relax
16:8 & \begin{tabularx}{0.7\textwidth}{X} 你要吃無酵餅六日，第七日要向耶和華－你的神守嚴肅會，不可做工。」 \end{tabularx} \\ \\ \relax
16:9 & \begin{tabularx}{0.7\textwidth}{X} 「你要計算七個七日：從你用鐮刀開始收割莊稼時算起，一共七個七日。 \end{tabularx} \\ \\ \relax
16:10 & \begin{tabularx}{0.7\textwidth}{X} 你要向耶和華－你的神守七七節，按照耶和華－你神所賜你的福，獻上你手裡的甘心祭。 \end{tabularx} \\ \\ \relax
16:11 & \begin{tabularx}{0.7\textwidth}{X} 你和你的兒女、僕婢，以及住在你城裡的利未人、在你中間寄居的和孤兒寡婦，都要在那裡，耶和華－你神選擇作為他名居所的地方，在耶和華－你神面前歡樂。 \end{tabularx} \\ \\ \relax
16:12 & \begin{tabularx}{0.7\textwidth}{X} 你要記得你在埃及作過奴僕，也要謹守遵行這些律例。」 \end{tabularx} \\ \\ \relax
16:13 & \begin{tabularx}{0.7\textwidth}{X} 「你收藏了禾場和壓酒池的出產以後，就要守住棚節七日。 \end{tabularx} \\ \\ \relax
16:14 & \begin{tabularx}{0.7\textwidth}{X} 在節期中，你和你的兒女、僕婢，以及住在你城裡的利未人、寄居的和孤兒寡婦，都要歡樂。 \end{tabularx} \\ \\ \relax
16:15 & \begin{tabularx}{0.7\textwidth}{X} 在耶和華所選擇的地方，你要向耶和華－你的神守節七日，因為耶和華－你的神要在你一切的收成上和你手裡所做的一切賜福給你，你就非常歡樂。 \end{tabularx} \\ \\ \relax
16:16 & \begin{tabularx}{0.7\textwidth}{X} 「你所有的男丁要在除酵節、七七節、住棚節，一年三次，在耶和華－你神所選擇的地方朝見他，不可空手朝見耶和華。 \end{tabularx} \\ \\ \relax
16:17 & \begin{tabularx}{0.7\textwidth}{X} 各人要按自己手中的能力，照耶和華－你神所賜你的福，奉獻禮物。」 \end{tabularx} \\ \\
[1ex]
\hline
\hline
\end{longtable}
$^{1}$今日要讀的聖經是新命紀第十六章一至四節和九至十七節.
在聖經舊約第二百三十四至二百三十五頁.
這裡記錄了以色列人三個節日.
逾越節,七七收穫節和駐彭節他們所遵守的一些條例.
請聽聖經的記載.
新命紀十六章第一節.
你要注意亞畢月向耶和華你的神守逾越節.
因為耶和華你的神在亞畢月夜間領你出埃及.
你當在耶和華所選擇要納為他命的居所.
從牛群羊群中將逾越節的祭生獻給耶和華你的神.
你吃這祭生不可吃有教的餅.
七日之內要吃無教餅就是困苦餅.
你本是急忙出了埃及地.
要叫你一生一世紀念你從埃及地出來的日子.
在你四境之內七日不可見面教.
頭一日晚上所獻的肉一點不可留到早晨.
第九節.
你要計算七七日從你開廉收穫和嫁時算起.
共計七七日.
你要照耶和華你神所賜你的福.
守女拿著金心祭獻在耶和華你的神面前.
守七七節.
你和你兒女瀑被並住在你城裡的利未人.
以及在你們中間寄居的如孤兒寡婦.
都要在耶和華你神所選擇納為他命的居所.
在耶和華你的神面前歡樂.
你也要紀念你在埃及作過奴僕.
你要謹守遵行這些律例.
你把和祥的菊酒榨的酒收藏以後.
就要守住彭節七日.
守節的時候.
你和你兒女瀑被並住在你城裡的利未人.
以及寄居的孤兒寡婦都要歡樂.
在耶和華所選擇的地方.
你當向耶和華神守節七日.
因為耶和華你神在你一切的土產上.
和你兒女所辦的事上要賜福於你.
你就非常的歡樂.
你一切的藍丁要在除膠節.
七七節,駐彭節一連三次.

$^{41}$在耶和華你神所選擇的地方朝見他.
卻不可空手朝見.
各人要按自己的力量.
照耶和華你神所賜的福分奉獻禮物.
靈姐妹平安.
剛才寶靈姐妹其實在讀經前已經幫我們開了頭.
今天我們會講關於以色列人節旗的一個訊息.
首先節旗在原文裡的意思是約會.
以色列人在特定的日子和地點和上帝相會.
不只是見面那麼簡單.
而是在過程中重溫上帝對整個民族的作為.
並且互相提醒他們和上帝有立約的關係.
他們需要守約遵守上帝的吩咐.
而在以色列人這麼悠長的歷史裡.
其實有三個我們說要回到耶路撒冷朝聖的日子.
所謂的三大節旗.
就是剛才申命記裡見到的.
逾越節,七七節和駐彭節.
如果你有留意剛才經文的細節.
你會發現出現了除教,收割和收藏.
我們說這些是甚麼用語?.
是農耕的用語.
暗示這三個節旗背後和古代的農業社會有密切的關係.
而在申命記之前的初一及二.
是這樣記載這三個節旗.
大家可以翻開程序表.
或者看powerpoint.
初一及二這裡說.
同樣是一年三次要守上帝的節旗.
你要守除教節.
「照我所吩咐你的,在亞不月裡所定的日子,吃無教餅七日」.
「誰也不可空手朝見我,因為你是這月出了埃及」.
另外又要守收割節.
「所收的,是你田間所種,勞碌得來初熟之物」.
最後要在年底收藏.
守甚麼?收藏節.
「一切的藍丁都要一年三次朝見主耶和華」.
其實聖經學者都認為.
除教節,收割節和收藏節.
本身是因應著循環的農耕生活週期.

$^{81}$每年都是這樣.
除教定時定候.
接著就到收割.
收割完就收藏.
是這樣循環的農耕生活衍生出來的三個節日.
再看剛才的表.
隨著以色列人的信仰歷史的發展.
有些信仰元素加入了農耕節期.
例如在表的最左邊.
最初是收割大麥的日子.
即是今天的三,四月.
對猶太人來說就是亞不月.
剛才看申命記就看到亞不月.
原本是除教節在埃及記.
但後來到之後的利未記,民數記甚至申命記.
多了一個名字.
叫逾越節.
再過兩個月.
五,六月開始收割小麥和其他農作物.
原本是收割節日子.
但來到民數記已經變成七七節.
或之後有個名字叫五秦節.
稍後我們會講解.
到年底九至十月.
基本上全年的收割都完成了.
那就是收藏節.
但來到民數記.
你會看到為何變成住棚節呢?.
很明顯是多了一些信仰元素.
你會看到其實.
上帝將一些信仰元素.
放在了原本的農耕節期.
令到以色列人守節期的目的.
由原本我們只是慶祝農作物有收成.
但現在提升了一個層次.
就是紀念上帝對整個民族的救恩.
所以每年守節期.
變相重新確立上帝子民的獨特身份.
他們不是普通農民過節.
而是上帝的子民.

$^{121}$以他們的獨特身份守上帝的節日.
提醒整個民族.
他們都有責任遵守上帝所立的約.
讓以色列人在慶祝農耕節期時.
可以從外邦人身上.
我們常說分別為勝.
有分別,大家都是那時候過節.
但以色列人過節有另一重重要意義.
這就是聖經強調分別為勝的意思.
有不同的.
怎樣有不同?.
我們逐個節期來看.
首先是逾越節.
我知道按照教會讀經計劃.
你已經讀完馬可福音.
如果你有留意馬可福音記載.
特別十四章裡說到.
除教節的第一天.
就是載逾越節羔羊那一天.
你會看到很特別.
兩個節日緊接相連.
甚至可以當成一個節期.
除教節裡說的教其實是甚麼意思呢?.
如果你今天在家裡做麵包.
你會很明白.
你要將麵教加在麵團.
令到它發酵.
發酵後就產生空氣膨脹.
吃下去會鬆軟一點.
如果沒有發酵膨脹.
吃下去會硬硬的.
很辛苦.
所以麵教令到麵團製成麵餅後.
就會有口感香氣容易咀嚼.
而在古代世界裡.
他們怎樣做教呢?.
其中一個方法就是.
將大麥和水混合.
然後等它發酵.
發酵成為細小的粉末.

$^{161}$你當做像今天的酵母菌.
類似這些東西.
並且將它儲起來備用.
但你知道.
其實今天的酵母菌也好.
星期的麵教也好.
可不可以長期擺放呢?.
又不可以.
是會有期限的.
也有過期的.
過期就不只是不能用那麼簡單.
是會產生毒素的.
所以在每年農耕社會裡.
他們的三,四月春季.
正值像現在這樣.
回暖潮濕.
東西容易發霉的時候.
除教節就在那些時間開始.
也就是剛好.
又是他們收割大麥開始的時間.
農民會在除教節裡.
除教七天.
將那些積存的教.
全部不要.
斷絕它.
所謂的斷教.
所以變相只能夠.
吃那些無教餅七天.
過了七天之後.
他們就可以用剛收割的大麥.
來重新培育新的麵教.
這是除教節的由來.
至於為什麼.
逾越節和除教節會扯上關係.
剛才我們看新面記很清楚.
除教節本身是亞畢月.
即是以色列人的正月來舉行.
而亞畢月剛好又是什麼日子呢?.
經文說了.
上帝領以色列人出埃及的月份.

$^{201}$同一個月份.
你看看《廣道大綱》裡.
新面記第十六章一至四節說了.
亞畢月他們要守逾越節.
因為上帝是那個月領他們出埃及.
上帝吩咐當日以色列人.
各家各戶要在正月十四的黃昏.
事先預備好那隻沒有殘血的羔羊.
要宰殺牠獻祭.
並且取一些羔羊的血.
將各家各戶塗在門框上.
到了晚上.
當上帝巡遊全國埃及地的時候.
要擊殺所有投生的時候.
那些人可以沒事?.
就是門框上有羔羊血記號.
上帝看到這個記號.
他就會越過所謂的「過節」.
這個就是逾越的意思.
越過不會擊殺的長子.
這個就是逾越節的由來.
現在摩西吩咐百姓.
不只是當日要守逾越節.
而是之後每年的亞畢月都要守逾越節.
吃逾越節的羔羊.
紀念上帝在那一個月.
將以色列人從埃及圍爐之地領出來.
我們很明白這個典故.
特別是上個月已經看完《出埃及記》.
在守節期的時候.
百姓不可以吃油餃餅.
要連續吃無餃餅七天.
目的是紀念當他們出埃及的時候.
走得很倉促,很蒼茫.
或者像經文說的很苦困.
如果大家記得《出埃及記》第十二章.
走難的情景是怎樣的.
以色列人在正月十四黃昏.
到黃昏之後.
各家各戶就要開始做一連串的行動.

$^{241}$宰殺逾越節的羔羊.
然後要把血塗在門框上.
然後又要急急忙忙.
將宰殺的羔羊烤煮.
然後又要馬上吃掉.
因為那些製肉不能留到過夜.
到了半夜.
上帝擊殺所有埃及長子之後.
法老就連夜急召摩西和阿倫.
說你帶走所有的以色列人.
馬上走.
那裡有多少人?.
可能有過百萬人.
立即離開埃及.
甚至連埃及人自己都趕走他們.
快點走.
當他們是瘟神送他們走.
於是所有百姓都急忙到一個地步.
就像經文說的.
拿著未發酵的生面.
馬上慌忙離開.
試想像一下.
大家都有去旅行的經驗.
如果一晚之內有這麼多行程.
辛苦嗎?.
當然辛苦.
你會罵死導遊.
突然說走就走.
你會覺得很匆忙.
更何況當時以色列人是逃難.
是逃亡.
所以你不難明白.
當日是多麼困苦.
你看到新命紀第十六章三節.
無教餅又有一個什麼名字?.
困苦餅.
它有一個註腳阿古屋.
因為以色列人當日.
本來是很急忙離開埃及.
所以無教餅七天.

$^{281}$不再是籠作的原因.
不是純粹將舊教送走.
迎接新教.
而是有一個新的屬靈意義.
就是要讓世世代代.
都紀念這段舊恩歷史.
雖然出埃及時非常困苦.
但上帝仍然保守住這群百姓.
拯救他們脫離這一切困苦.
至於逾越節對新約教會.
對我們基督徒.
其實我們都不陌生.
我們都明白.
逾越節那隻被宰的羔羊.
最終是指向誰?.
耶穌基督,我們都很明白.
所以我們每年.
特別是這段時間.
我們最清楚.
之前的大齋期.
和即將到來的壽苦節.
我們都紀念耶穌基督.
這隻被宰的逾越節羔羊.
祂為世人捨命.
將我們從罪惡的困苦中.
拯救出來.
我們很清楚.
逾越節本身對我們有什麼意義.
比較陌生的是之後那兩個.
七七節和鑄棚節.
而七七節是什麼意思呢?.
經文說九至十一節.
是開廉收割的日子.
即是他們開始收割大麥.
剛才我們說三四月.
要計算之後的七個七日.
也就是逾越節之後的五十日.
這就是七七節.
所以剛才我們說.
七七節有另一個名字.

$^{321}$就是五秦節.
五秦節我們比較清楚.
我們新約教會.
很多教會都會慶祝五秦節.
聖靈降臨.
在短時間內.
那些大麥差不多收割完.
現在就開始陸續收割小麥.
和其他農產品.
譬如葡萄.
或者各種夏天的果實.
譬如橄欖.
所以七七節的前身.
是一個我們說收割.
慶祝農作收成的日子.
但是隨著以色列人.
那個救恩歷史的發展.
守節的原因.
也同時有所改變.
根據猶太拉比的傳統推算.
他們說逾越節之後的第五十天.
其實剛好是以色列人離開埃及.
去到西乃曠野那一段日子.
期間發生了哪兩件大事?.
考考大家.
你們看了出埃及了.
去到西乃山.
最重要的兩件事是什麼?.
就是上帝和以色列人納藥.
納了藥之後還有什麼?.
頒佈律法.
你看到.
這就是以色列人最大的收穫.
最大的收割.
使得收割節原本只是.
慶祝農作收成的日子.
現在提升到七七節.
慶祝他們一個屬靈的豐收.
變成一個節期.
因為在這段時間.

$^{361}$以色列人收穫了一個新的身份.
是和上帝納藥的子民.
並且他們獲得上帝頒佈他們的律法.
所以.
以色列人每年守七七節.
或者所謂五旬節.
其實是一個對西乃之約的重溫.
紀念上帝和他的律法.
怎樣成為了以色列人的生活衣櫃.
提醒他們.
之所以有農作物的豐收收成.
原因不是純粹是因為氣候環境使然.
也不是人努力耕種.
是因為上帝在律法所說.
他守約,思,持,愛.
眷顧他的百姓.
透過農作業的產品.
來賜福,供應他們.
與此同時.
也是因為百姓願意遵守上帝的約.
遵行律法的吩咐.
所以他們透過農作物的收成.
經歷了上帝所賜的福.
整件事不同了.
因為屬靈的意義.
加進了他們循環不息的農耕週期.
變成了日復日的生活.
每年都是這樣.
但是有了信仰元素就不同了.
因為知道在營營役役的裡面.
背後是上帝的祝福和供應.
值得一提的是.
猶太人在每年過七七節的時候.
他們會重讀一卷舊約書卷.
叫做《路得記》.
為何讀《路得記》呢?.
一方面如果有看過《路得記》.
你會記得.
《路得記》出現了很多農作物收割的場景.
與七七節本身收割的背景相似.

$^{401}$另一方面.
《路得記》很強調一個主題.
就是要以恩慈去對待.
身邊有需要的弱勢社群.
這就是《申命記》強調的主題.
第十一節說的是甚麼呢?.
提醒以色列人.
在七七節的裡面.
不只是與家人.
不只是與你的仆彼.
而是要與城裡的利未人.
和寄居的.
和孤兒寡婦.
來分享上帝給他們的修成.
要一切的來歡樂.
這些人因為不同的原因.
他們沒有自己的土地.
換句話說.
他們沒有辦法通過耕種.
來獲取生活的基本需要.
所以上帝就吩咐他的子民.
在勞碌為自己收割的日子.
不要忘記了律法的吩咐.
不要忘記了自己是上帝子民的身份.
要顧念周邊的貧窮人.
或者寄居者.
一些基本生活的需要.
盡你的所能.
這裡說的是獻上金心祭.
上帝如何祝福你.
上帝賜給你什麼.
你就盡你的力量.
為上帝擺上.
去祝福身邊的人.
透過獻上金心祭.
很自然就供應到祭司利未人.
他們在飲食上的需要.
並且當他們收割的時候.
上帝透過律法吩咐.
你不要割盡.

$^{441}$你不要收盡.
要留一些給寄居的孤兒寡婦.
讓他們自己去拾取享用.
特別我印了另一段申命記.
教後的二十四章.
你會看到九至十二節.
這裡提醒以色列人.
在田間收割和價的時候.
如果你真的留下了一些.
忘記拿的.
你不要快去拿.
那一群不是你的了.
那一群是留給寄居的孤兒寡婦.
他們自己去拾取.
上帝透過這樣去供應他們的需要.
又或者當你打橄欖的時候.
你留下了一些.
你不要再打橄欖.
或者你拾葡萄的時候.
你不要再拾剩下的.
刻意要你留下.
讓寄居的孤兒寡婦.
享受上帝給他們的收穫.
所以這就是七七節.
當時一個很重要的意義.
對我們來說.
對現代的信徒來說.
又有什麼樹齡意義呢?.
還記得在新約裡面.
什麼時候聖靈降臨嗎?.
五重節在小堂轉.
那一天厲害了.
那一天怎麼樣?.
彼得說到三千人信主.
收割了三千個新的果子.
多大收穫嗎?.
這就是教會.
初期教會最大的收穫.
豐收.
其實也間接提醒.

$^{481}$每一個基督徒.
我們都是上帝收割回來的果子.
當我們在勞碌.
每天營營役役.
上班賺錢的時候.
不要忘記我們真正的身份是什麼.
我們是和上帝納入的子民.
我們是上帝的兒女.
耶穌基督的門徒.
在日常生活裡.
我們要遵行上帝的吩咐.
實踐耶穌的教導.
去為上帝收果子.
收莊稼.
這才是我們的身份.
可能我們在公司裡上班.
甚至是家庭主婦.
我們有自己的工作.
我們在那裡收割.
我們所需要的.
獲取我們所需要的.
但不要忘記.
我們有另一個身份.
我們要替上帝傳福音.
結果子.
讓別人無論在生活上.
或是在靈性上.
他們都得到需要.
這才是我們真正的收穫.
也是收割節.
七七節.
或是五秦節.
給我們今天的意義.
最後一個.
主棒節.
這裡說.
不如大家讀一讀.
第十三至十五節.
印在程序表裡.
十六章的十三至十五節.

$^{521}$預備起.
以馬和上帝保護.
受上帝的收穫以後.
就要收主棒節出入.
收節的時候.
讓我們以熱淚為伍.
和氣.
接觸到我們成為的敵人.
以及敵對的如虎而化虎.
都要歡樂.
在我們和下所選擇的地方.
讓我們安穩和下的精神.
收節七天.
因為和下的精神.
在我們一切的表現上.
和我們所累積的思想上.
要似乎與你.
有非常的歡樂.
好,謝謝.
其實這裡說的是.
和場的菊已經收齊了.
酒渣釀的酒收藏好了.
全年的收割工作已經完成.
所有的農產品,農作物.
也都收好在倉庫裡.
即是說準備慶祝過年.
這個就是收割節.
本身的節日背景.
為什麼之後會變成駐棚節呢?.
當然,這也與以色列人.
舊恩歷史的發展有關.
首先我們很明白.
駐棚節是為了紀念什麼?.
以色列人什麼時候駐象棚?.
就是他們離開埃及後.
在曠野漂流的時候.
搭棚居住.
並且棚是不會永遠停留在那個地方.
它要跟著上帝走.
回望何時移動,棚何時拆.

$^{561}$回望何時停留,它們又再搭建.
所以那只是一個短暫的居所.
這就是駐棚.
所以駐棚只是紀念.
昔日以色列人離開埃及後.
在曠野駐棚居住.
那四十年的漂流歲月.
在那個時候.
百姓對上帝是怎樣的?.
你見到很難搞的.
上帝的關係有時好,有時不好.
但上帝沒有因為百姓的過犯罪惡.
而棄絕了他們.
在這四十年駐棚的生活裡.
仍然對他們手藥師慈愛.
讓他們存活.
很耐心地引導他們.
跨過了這四十年駐棚的歲月.
最終去到迦南美地.
而即使之後的百姓.
當他們進入了迦南生活.
在那個時候就不用駐棚了.
因為開始擁有自己所建造的房屋.
但上帝仍然吩咐他們.
每年在駐棚節.
在上帝揀選的地方.
搭建一個簡陋的臨時帳棚.
全家來居住首節七天.
意思是什麼?.
不是我們今天像都市人一樣.
住了私樓之後.
趁假期就一家大小出去露營.
不是這個意思.
駐棚其實是一個很實際.
一個實體的信仰教育場地.
為什麼這樣說呢?.
就是讓所有以色列人.
他們可以集體紀念.
原來我們這個民族.
是有過一段漂流曠野.

$^{601}$居無定所的日子.
從而反省這段日子.
為什麼我們可以熬得過呢?.
很艱難的.
是因為什麼緣故?.
不是因為我們生得夠銀利.
是因為上帝.
是因為上帝憐憫.
守約思慈愛.
我們才可以在駐棚的日子.
活過來.
即使之後.
我們有了自己的田地.
有了自己的土產.
有了自己的瓦遮頭.
生活開始變得安穩.
所謂的有保障.
但是生命的本質.
有沒有分別?.
駐棚節就是這個提醒.
即使你現在多富裕.
你的本質也是駐棚的.
換句話來說.
我們和昔日在曠野漂流的.
以色列人沒有分別.
我們之所以存活是因為.
上帝公認保守.
不是因為我們勤力拼搏夠聰明.
所以駐棚節.
是提醒每一代的以色列人.
在我們很努力.
在快要過冬之前.
我們積存了足夠的糧食.
積存了很多的財富.
我們生活開始富裕.
覺得生命似乎有了保障的時候.
駐棚節就像一棍打下來.
你不要搞錯.
上帝才是你真正的保障.
所以弟兄姊妹.

$^{641}$駐棚節雖然我們今天.
不需要去守這個節日.
但是它給了我們一個很好的提醒.
我們今天會不會因為不同的憂慮.
我們過分為自己.
收藏,積存了一些不必要的財寶.
我們希望在那些地上的財富裡面.
找到一個堅固的保障.
我們會不會希望他們.
給我們真正的平安.
駐棚節提醒我們一個很重要的事實.
這個世界不是我們永恆的居所.
換句話來說.
駐棚節其實是教導我們.
你要帶著一個客裡的態度.
你要帶著一種永恆的意識.
在地上生活.
你要明白上帝現在給你的.
你收藏中的,你積累中的.
是短暫的.
那些東西不能帶去永恆.
那些東西有一天一定會過去.
會扭壞.
所以我們就懂得將地上短暫的資源.
你要透過分享,透過施予.
這些善行.
來化成天上永恆的財寶.
這是駐棚節的意思.
事實上駐棚節給我們的屬靈意義.
我想剛好對應今天我們所身處的環境.
特別近期大家都知道關稅戰.
全球經濟危機.
你有沒有擔心你的錢包少了錢?.
看一看你的全接強積金.
會不會突然間好像少了一個零.
所以你會看到坊間很多資訊教你.
我們現在手上的資金要怎樣轉移,避險.
希望可以保住.
你一直積累,收藏的財富.
那些你覺得人生的保障.

$^{681}$弟兄姊妹,這會不會是我們所憂慮的?.
但是駐棚節卻告訴我們.
我們真正的保障是上帝和祂的話語.
那才是我們永遠的保障.
真正的平安,耶穌說.
你見機在祂的裡面.
你藏在祂的裡面.
最後一個段落.
第十六節至第十七節.
這裡還你一切的南丁.
要在除教節,七七節,駐棚節.
一連三次.
在耶和華萬里神所選擇的地方.
朝見他.
卻不可空手朝見.
那個人要按自己的力量.
照耶和華萬里神所指的方法.
奉獻禮物.
你留意.
十六節這裡摩西吩咐以色列人.
每連三次.
要在上帝所選擇的地方.
過節,首節期.
我們說重要的事情要說多少次?.
至少三次吧.
如果你有留意.
上帝這句話說了六次.
在一至十七節裡面.
剛才最後一次就在十六節.
你發覺,十五節說一次.
然後呢.
在十一節就說一次.
另外有兩次沒有引述出來.
六七節就說了.
第二節就說了.
類似的吩咐說了六次.
你可想而知有多重要.
我們知道以色列人.
當他們一旦進入迦南生活的時候.
這十二個支派就會開始.

$^{721}$分散來居住.
他們會分散到.
整個巴勒斯坦.
整個迦南地.
不同的地方.
居住.
但在首節期的時候.
百姓卻不可以.
在自己居住的地方.
貪我方便.
便利我自己.
在自己的城市裡.
自己設立一個首節.
敬拜獻祭的場所.
可不可以?.
上帝說不可以.
不可以貪方便.
上帝要求散居在各地的以色列人.
每年三次.
要專程回到上帝指定的地方.
即是回望的地方.
或者說.
聖殿所在地耶路撒冷.
你一定要在那裡.
守你的節日.
朝見他.
獻上禮物.
這樣就麻煩了.
你住得近就好.
你住在猶大地就好.
可能隔幾條街就到了.
你住在迦尼尼.
你就知道.
你會有很多不便.
上帝的吩咐.
對散居的百姓.
會造成不便.
特別是住得遠離.
回望和耶路撒冷的以色列人.
你想想.

$^{761}$一年三次.
都很頻密.
五月節朝聖完結後.
隔兩個月又來了.
五旬節七一節又要去朝聖.
再隔四個月.
又要去住彭寺.
並且上帝說.
不只是你要帶自己路上的物資.
不是帶衣服,鞋子,食物,用品.
上帝說你不可以空手朝見他.
即是你要帶什麼?.
你要帶祭身去.
你要帶禮物去.
雖然說是南丁.
即是滿了十二歲以上的男孩.
但很多時候不只是男人去.
留意耶穌受好節期的時候.
一家大小.
起碼媽媽,瑪利亞都有去.
因為你小的十二歲去了.
媽媽也會去.
一群男人去不懂照顧自己.
婦女有時也會跟著去.
所以每次朝聖.
也不容易.
可能沒有我們今天家庭旅行那麼複雜.
又要訂機票,訂酒店.
當小朋友請假等等.
但絕對不是簡單的事情.
每次朝聖.
百姓都要花時間,心思,努力去預備.
不過是很麻煩的.
但正好是考驗.
百姓是否對上帝信服.
你對上帝有沒有心?.
對上帝有沒有愛?.
就在那裡見真章.
事實上不容易.
有很多外面的人.

$^{801}$會叫你不要這樣做.
試想像一下.
以色列人進入迦南生活之後.
他們周邊的迦南人.
他們的獻祭不用那麼複雜.
他們的守節期不用那麼麻煩.
沒有那麼多規矩.
可以很隨心.
你照自己的想法.
你在山上搭個壇.
來獻祭,守節也可以.
你在山下找棵樹.
搭個壇,守節,獻祭也可以.
非常方便,簡單.
你有心就可以了.
弟兄姊妹.
你會看到當以色列人進入迦南之後.
會被各式各樣的偶像.
敬拜的風俗影響.
引誘他們慢慢偏離上帝的吩咐.
這些看來很簡單,很方便.
很吸引人的敬拜方式,過節的方式.
卻不是上帝所喜悅,所揀選.
上帝之所以要求以色列人.
每年三次專程回到.
他所指定的地方來守這些節日.
是要考驗百姓對他的忠誠.
面對著不同的引誘,挑戰的時候.
我們能不能夠堅持.
來討上帝的喜悅.
順從上帝的吩咐.
揀選,走他所揀選的道路.
以致我們即使和外邦人一起生活.
我們仍然可以竭力在他們中間.
分別出來,讓他們看得到.
從而讓他們有機會認識.
我們所相信的上帝.
所以,弟兄姊妹.
駐彭節給我們一個提醒.
或者說節期給我們一個提醒.

$^{841}$特別是為了下星期.
我們也有新約教會要守的節日.
我們也有壽苦節,復活節.
我們準備好了嗎?.
特別是上帝的吩咐.
不可以空手朝見他.
換句話來說,這是一個挑戰.
回來也不夠.
來對地方也不夠.
不可以空手.
但上帝說得好.
上帝說:你一定有能力.
因為經文說了什麼?.
「按上帝賜給你的福」.
按你自己的能力.
即是說上帝不是要難為我們.
你是做得到.
是否願意.
是可以的.
弟兄姊妹,事實上你會發覺.
也因為這個緣故踏入了三月.
從屠夫禮開始.
每個星期三,剛才達哥也有說.
我們會有不同的主題.
幫助弟兄姊妹進入壽苦,復活節.
這段節期.
讓我們真的可以身體力行地.
實踐耶穌基督的吩咐.
預備好如何和耶穌基督同生同死.
也讓身邊的人透過我們言行.
可以經歷到上帝所賜的福氣和豐盛.
所以盼望藉著今天.
雖然舊日的節期離我們很遠.
但裡面的屬靈意義.
卻是色切每一代的基督徒.
願意我們帶著一個預備的心.
帶著一個朝見上帝的心.
不單止在節期.
甚至每一個主日.
我們都尊崇敬拜祂.

$^{881}$向祂獻上禮物.
我們同心禱告.
七愛因主願你自己的話語.
成為我們生命的亮光.
引導我們好好在每一個崇拜.
每一次節期裡朝見你.
帶著你給我們的力量.
帶著你所賜給我們的各樣事物.
我們預備好來當作禮物獻上給你.
事實上上帝讓我們知道.
你喜悅的是我們將我們全人.
當作活祭獻上來敬拜你.
服侍你.
所以為著今天我們所領受的經文.
能夠成為我們今天生活的提醒.
讓我們知道我們在每日循環生活裡.
輪營日夜的工作裡.
其實上帝你揀選拯救了我們.
給我們一個新的身份.
我們真正就是上帝的子民.
我們是耶穌基督的門徒.
在生活裡我們有更加重要的職份.
就是替上帝來收莊稼,結果子.
以致我們可以被上帝使用.
讓上帝你自己各樣的福氣.
你的恩典藉著我們的言行.
能夠成為我們身邊.
特別是未信的人.
祂們的祝福和幫助.
為此我求上帝.
藉著今天我們所領受的經文.
來帶領我們.
讓我們重新去檢視我們和你的關係.
讓我們明白.
我們在世上只是一個暫住的居所.
我們所仰望的是永恆的天家.
讓我們能夠有從你而來那份屬天的智慧.
將地上我們所擁有的,所積存的.
轉化成屬天的財寶.
能夠成為你所喜悅的活祭.

$^{921}$為此我們同心獻上禱告.
願上帝你垂聽.
奉基督耶穌你得勝的名.
來祈求.
阿們.
其他事項.
\newpage



\section{ }
\label{sec:96uE5WjMIVg}
\textbf{2025年4月18日受苦節崇拜直播}
\newline
\newline
連結: \href{https://youtube.com/watch?v=96uE5WjMIVg}{\texttt{https://youtube.com/watch?v=96uE5WjMIVg}} ~~~~ 語音日期: 2025-04-19
\newline
\newline
\hyperref[sec:CrZ2VP_EO18]{\small{< < < PREV SERMON < < <}}
~
\hyperref[sec:index_chronic]{\small{[返順時目]}}
~
\hyperref[sec:9gw1hxqzLeo]{\small{> > > NEXT SERMON > > >}}
\newline
\newline
$^{1}$為甚麼會這樣?.
耶穌明明是無罪的.
祂曾經醫治我.
在我身上趕走七隻鬼.
祂給了我新的生命.
為甚麼?為甚麼祂要承受這一切?.
我們都不明白.
現在只能陪伴祂到最後一刻.
我記得那天.
我用珍貴的香膏來塗抹祂的時候.
有人說我浪費.
但主就說.
「隨它吧,它是為我安葬之日預備的」.
難道今天就是那天.
我會被祂承受嗎?.
當年天使向我顯現.
祂說「你必懷孕生子」.
「他必稱為至高者神的兒子」.
今天我的兒子飽受折磨.
我還要去做我自己的計劃.
如果….
如果耶穌真的是神的兒子.
為甚麼祂不從十字架下來?.
為甚麼祂要承受這一切?.
或許我們現在看不清楚.
但神必定有祂的旨意.
你們還記得那天嗎?.
當主來到伯大利的時候.
我坐在祂的腳旁聽道.
當時我的心裡充滿著平安.
或許現在我們也需要想起.
祂所說的一切說話.
是的,天使曾經宣告.
在神沒有難成的事.
雖然眼前的一切令我心如刀割.
但神從來不會失信於人.
我們應該如何面對這麼殘酷的現實?.
我們所愛的人正在受苦.
或許我們現在無法理解.
但我們必須相信神是掌管一切.

$^{41}$主必定有祂美好的旨意.
在祂最痛苦的時候.
祂選擇赦免.
這是我的兒子.
是至高者的兒子.
或許我們現在所經歷的一切.
是有更深遠,更榮耀的目的.
主,雖然我不明白.
但我願意相信你的應許必定成就.
主,你曾經說過.
我的香膏是為你安葬之日預備的.
現在我終於明白了.
求你幫助我們繼續相信你的話.
無論如何,我們都要繼續跟隨耶穌.
就算前面的路充滿未知和黑暗.
我們就彼此扶持.
一起等待神應許的那一天.
你答應我們,你答應我們.
《聖經》.
當耶穌基督受難的時候.
被釘十架的時候.
《聖經》裡記載除了使徒約翰之外.
就是剛才那幾位婦女.
釘十架可以說是當時.
一種最可恥羞辱的懲罰.
所以沒有什麼人願意.
跟那些被釘十架的人拉上關係.
但是這群婦女.
她們不介意別人的目光.
她們堅持陪伴她們所愛的主.
走到最後的一刻.
這群婦女她們所表現出來的勇氣.
你會看到.
跟那些四散逃亡的門徒.
形成一種很鮮明的對比.
後者我們看到.
因為很害怕被牽連.
於是乎都紛紛離開了主耶穌.
事實上我們知道.
這群婦女其實不是沒有風險.

$^{81}$她們都有可能被牽連.
被當成是耶穌的同黨.
將她們一併的來去刑罰.
之所以能夠毫無懼怕地跟隨耶穌.
是因為那一份對主的愛.
激發在她們內心.
一種很強大的勇氣.
她們決定選擇留在現場支持耶穌.
她們所展現出來的一份力量.
可以說成為我們今天每一個.
好像打開了一道窗戶.
希望我們每一個可以透過她們看到.
當正面對苦難掙扎的人.
其實可以怎樣做.
可以怎樣依靠耶穌.
跟耶穌同行.
雖然遭受到創傷的是耶穌本人.
但你會看到耶穌所受的苦難.
同樣也成為這群婦女.
內心對她們造成了創傷.
她們覺得很委屈.
很不甘心.
明明巡撫已經查明了.
耶穌是沒有罪.
很應該釋放耶穌.
但耶穌反而是無辜枉死.
她們也質疑上帝.
為何神要這樣.
來折磨苦待我們所愛的耶穌.
甚至我的兒子.
難道這樣真的可以將罪人拯救出來.
她們也懷疑耶穌的能力.
耶穌既然有本事將我身上的污穢趕出.
為何你不運用這份能力.
來拯救自己.
你可以從十字架上下來.
為何你不這樣做.
她們也很無奈,很無助.
主啊.
你平日多麼愛護我們.

$^{121}$幫助我們.
但這一刻.
我們除了站在這裡.
我們除了為你哭泣.
我們甚麼也做不到.
我們看到耶穌基督的十字架.
和每一個跟隨祂的人.
好像變成了一種命運共同體.
衝擊著當時每一個跟隨耶穌的人.
祂的信仰.
耶穌不是本來滿有榮耀能力嗎.
為何現在卻受盡屈辱.
祂自己也變成了無能為力.
耶穌你不是說過要帶領世人走向美善嗎.
為何這個世界到處仍然有很多惡人,惡事,四處的張狂.
耶穌你明明已經為我們的生活帶來了很多盼望.
為何這一刻你又突然叫我們陷入失望甚至絕望的裡面.
耶穌你所說的.
不是福音,不是好訊息嗎.
為何這個福音反而會叫我們痛苦,叫我們悲傷.
弟兄姊妹,我想這一切都是歷世歷代的基督徒.
包括你和我總會面對到的信仰衝擊.
我們看到這一班很愛主的婦女.
她們卻選擇毫不忌諱地正視自己這種信仰的掙扎.
她們甚至願意分享她們內心的矛盾.
雖然始終她們無法在理性上,認知上.
她們明白為何耶穌要受苦.
但當她們在分享內心痛苦的過程中.
她們開始回憶耶穌對她們說過甚麼.
她們開始想起和耶穌的互動的片段.
就是這樣,你會看到她們漸漸從耶穌和自己的創傷中.
開始找到一些支撐力.
找到支撐自己繼續走下去的那種韌力.
沒錯,現實是很殘酷,很痛苦.
我們失去了耶穌,失去了摯愛.
我們這一刻是無比的心痛.
但我們仍然選擇這種悲慘.
不是耶穌,也不是我們的結局.
確實,縱使我們會有難過,會悲傷.
但我們仍然有跟隨著耶穌的理由.

$^{161}$因為這個世界還是上帝掌管著.
包括世上的苦難,仍然在上帝掌握之中.
而這個世界會像耶穌所說.
最終會朝向上帝的美善旨意前進.
所以我們仍然有理由為耶穌好好活下去.
並且我們知道,原來在苦難的裡面.
我們不是唯一,我們不是孤單.
有一班好姐妹,我們可以互相扶持.
我們可以依傍,我們可以一起等候.
上帝的應許來實現.
確實我們會看到,聖經沒有對每一個人.
他所經歷的苦難,提供一個既定的標準答案.
但很奇妙,上帝的真理.
卻是會邀請人來跟那位苦難的主同行.
在苦難的路上一起探索,裡面的意義是甚麼.
而這班婦女,她之所以沒有被內心的痛苦打倒.
是因為她們懂得運用這種屬靈資源和力量.
她們將她們痛苦的經歷慢慢轉化.
當成一趟旅程,一趟屬靈的旅程.
當她們一路去憶述,耶穌曾經對我說過甚麼.
說了甚麼應許.
當她們一路去回想,耶穌怎樣幫過我.
怎樣跟我互動.
她們就一路跟那位受苦的主保持聯繫.
在傾訴自己的想法,自己的感受的時候.
她們也同時願意去聆聽真理.
她們讓真理對她們的內心發言.
引導她們可以怎樣做.
事實上弟兄姊妹面對苦難的問題.
如果單靠我們在情感上的一些抒發.
其實很可能會令當事人慢慢陷入一種自憐或自瀆的心態.
事實上在苦難當中.
我們更加需要的是倚靠上帝的真理來光照我們.
更新我們裡面面對苦難的那種略想法,那種態度.
清教徒有一位領袖,他叫巴赫斯特.
他在他有一本著作,叫做信靠信服的裡面.
為人生的苦難提供了一些很重要的神學反省.
幫助我們怎樣用上帝的話語,怎樣用真理.
來轉化苦難裡面的情感.
事實上巴赫斯特的人生都充滿苦難.

$^{201}$在他晚年的時候,他受到很多的逼迫.
他所有的財產都被沒收.
更加雪上加霜的是他的太太.
以及他一班好朋友都在那段時間相繼離世.
你可以想像到在他的晚年有多麼的坎坷,多麼的辜負.
他之所以讚寫信靠信服這本書.
原先是很想為自己在這個這麼苦難的階段.
整理一下自己,很希望為自己寫一些靈修材料.
來過渡人生最痛苦的時刻.
巴赫斯特他表示,他說對上帝那種真實的信靠信服.
並非要否定我們每一個人的內心感受.
因為這一切的情感原本都是來自上帝.
給人一種信號,提醒人現在落入一種什麼狀況.
對人都是有益的.
而當面對苦難逆境的時候.
這些情感就好像信號一樣提醒自己.
令到我們的身體和心靈會感應到相應的一種痛苦.
但是巴赫斯特他提醒我們.
我們作為上帝的兒女,我們是上帝的管家.
我們都有責任好好去管理這一切的情感.
包括在苦難裡面那些憤怒,那些哀傷和恐懼的情緒.
避免我們的情緒陷入失控,完全失去秩序.
釀成一種不良的後果.
甚至產生一些怨天尤人,埋怨上帝的想法和行為.
他提醒我們,每一個信徒在苦難的時候.
要認定上帝沒有變.
他的主權,他的公義,他的智慧和他的慈愛.
仍然是有效的.
從而幫助每一個人,可以在苦難裡懂得回到上帝那裡.
仍然保持他的理性和意志.
來去順應和配合上帝當下在他身上的安排.
既接觸自己的內心和情感.
但同時又用上帝的真理來去管理,駕馭自己這一切的感受.
為的是要成就上帝的旨意.
這一切都是這位屬靈前輩巴克斯特.
留給我們一些很重要的信仰遺產.
確實在座當中,沒有一個人喜歡痛苦苦難臨到.
但是身為上帝的兒女.
其實我們又應該最清楚明白.
苦難本身是你和我生命的一部分.

$^{241}$不可以分割.
如果你看聖經,其實聖經一早說了幾件事情很清楚.
是關於苦難.
聖經說當罪進入世界的時候.
苦難也都一定來到.
所以在亞當之後,每一個人.
他們都需要通過母親懷胎生產所受的苦難.
才會來到.
聖經也告訴我們,在這個罪惡的世界裡.
我們每一個人都帶著被傷害或加害人.
這兩個角色來生活每一天.
所以某程度上來說.
我們都有份製造參與這一切的苦難.
但是聖經也教導我們要感恩.
因為上帝的兒子耶穌基督來到這個世上.
他為了我們的罪受死.
他擔負了我們一切的罪債.
藉著耶穌基督的苦難.
他成就了上帝的救恩.
將我們從人生裡最大的永遠的痛苦.
就是和上帝隔絕這個苦況.
徹底的拯救出來.
並且很吊詭的是.
聖經說我們作為基督徒.
我們信了主的一群人.
我們夢照原來是要為主受苦.
和耶穌基督的苦來有份.
耶穌邀請我們一同去經歷他的苦難.
通過苦難.
不單單只你和我的生命.
得到練正,得到成長.
也通過我們的苦難.
來幫助,祝福那些陷入困苦的人.
聖經最後也說到.
地上這一切的苦難.
其實都會過去.
而在苦難的盡頭.
耶穌已經為我們預備了永恆的福樂.
所以,聖經去幫助我們.
在苦難的裡面建立那種屬天的思維.

$^{281}$來轉化地上這一切的受苦.
我們固然仍然悲嘆.
為什麼人世間會有這麼多的疾苦.
他已經將我們拯救.
脫離了地獄永遠的苦.
我們一方面仍然去重視.
人在苦難當中是有真實的情感.
但是我們又不會被這一切的情感氾濫.
完全的失控.
甚至高舉情緒.
多於上帝的旨意和他的應許.
我們也會明白.
苦難確實可能會奪走了我們很珍愛的東西.
包括生命.
但是聖經應許.
沒有任何的事情.
包括苦難.
可以奪去上帝給我們的福樂.
祂給我們的平安和盼望.
所以基督徒不會將他今生所虧損的.
所損失的.
來到他無限放大.
當成好像失去了上帝般痛苦.
當成好像失去了救恩般痛苦.
並且我們懂得和那些活在苦難中的人.
來到同行.
我們會為他們難過的同時.
我們更加為他們的靈魂仍然失望.
仍然未信主來到著急.
因為這個我們知道對他們來說其實是真正的苦難.
我們會努力把握剛恩.
來為他們信主得救.
來代禱.
甚至有實際的行動.
所以弟兄姊妹為這一切的緣故.
我們和這班婦女一樣.
我們有足夠的理由.
在這個苦難世界裡面.
你保持一種積極的生活態度.
就好像昔日這班與耶穌同行的婦人一樣.

$^{321}$你仍然努力為主而活.
沒有將自己完全陷入苦難的情緒當中.
耶穌基督曾經說過.
若有人要跟從我.
就當寫記背起他的十字架來跟從我.
因為凡要救自己生命的.
必喪掉生命.
凡為我喪掉生命的.
必得著生命.
你看到對那些逃跑了的門徒來說.
要愛那一位平時四處傳道.
醫治趕鬼.
一切都很成功很順境的耶穌的時候.
他們很樂意.
非常的嚮往.
但是這一刻.
仍然要他們繼續愛.
去跟隨這一位已經被釘十架的主.
覺得玩完的主.
這一位失敗受苦的主的時候.
對門徒來說.
這一刻的包袱太沉重.
太難受.
內心強烈的不安和恐懼.
令他們卻步.
他們不敢繼續走近耶穌.
並且跟耶穌要保持距離.
劃清界線.
門徒這個表現其實也讓我們很明白.
原來今天我們只知道.
我們要愛的對象是耶穌.
不夠支撐你的信仰人生.
苦難原來迫使所有的基督徒.
要重新檢視.
我和耶穌的關係到底是怎樣.
我愛耶穌的態度和力量.
到底夠不夠.
會不會像昔日的門徒一樣脆弱.
一碰就散.
還是像這一班婦女一樣堅韌.

$^{361}$怎樣都繼續跟隨著耶穌.
怎樣都愛耶穌到底.
我們看到這班婦女的堅持提醒了我們.
愛耶穌不是一種靜止的心靈狀態.
愛耶穌是一個持續的互動.
當中的關鍵是對主付出.
事實上我們要問.
為什麼這班婦女仍然願意在生死關頭冒險付出愛.
來愛耶穌和祂同行.
答案我們知道.
聖經也有說.
是因為耶穌先愛他們.
是因為耶穌首先毫無保留付出了祂的生命.
不單單是耶穌釘十字架犧牲寫明果國的時候.
而是在平時.
祂和這班婦女在日常的互動.
平日耶穌怎樣透過祂的宣講.
透過祂的教訓.
怎樣透過祂的醫治服侍.
來將祂的愛.
包括祂對生命的熱誠.
那份愛人愛神的態度.
就是這樣點點滴滴.
傳遞到這班婦女的身上.
來牧養豐富她們的人生.
同樣今日.
耶穌基督不單單為我們犧牲寫明.
耶穌也要將祂的生命力.
傳遞給我們賜給我們.
包括你和我在面對困難苦難的時候.
怎樣去跨越的力量.
耶穌說這一切都要賜給你和我.
使我們的生命真真實實邁向完整和豐盛.
並且耶穌也要透過苦難.
來引導我們成長變得更加堅毅.
祂也要透過我們這種生命.
和你身邊可能只是面對困苦的人.
來同行提升他們的生命能力.
聽姐妹玩具說來說.
耶穌將他的苦難轉化為我們的祝福.

$^{401}$堅固了我們和他的關係.
也提升了我們每一個對生命的熱誠和態度.
和大家分享一個生命的故事.
有一個人他叫做山姆伯恩斯.
就是你現在看到的這一位.
在他大概兩歲的時候.
他被確診患有一種罕有的病叫做早衰症.
衰老的衰.
意思是這個病會使一個人的身體的老化速度異常驚人.
平均的年齡只有13歲.
所以山姆他從小就知道.
自己會比正常人所面對的生活.
不單單是更加短暫.
更加無力.
也會更加困難很多的障礙.
但很特別他沒有因為這樣而怨天尤人.
他很樂天知命地繼續生活.
雖然最後山姆他只能夠活到17歲.
但在他離世之前的一個月.
他舉行了一次的演講.
還拍了一條影片.
他分享到自己即使是疾病纏身.
為什麼仍然可以樂觀面對他人生的三個秘訣.
第一個就是不要聚焦在你做不到的事情上.
而是關注任何時刻你仍然可以做到什麼.
因為人生很苦短.
很多美好的事物瞬間就會即逝.
你要好好把握.
山姆他表示即使他患有早衰症.
但原來他大部分的時間不是花在這個病上.
他做一些與這個病無關痛癢的事.
包括完成他人生理想.
包括音樂,包括打鼓.
這是他第一個秘訣.
而第二個就是每天讓那些你很想跟他共處的人.
環繞在你身邊.
尤其是那些會帶給你正面能量的人.
帶給你盼望正向思維的人.
你跟他們在一起.
山姆雖然患病.

$^{441}$但他擁有一個很幸福美滿的家庭.
他也有一班要好的朋友.
他們大家會時常相聚.
一起做一些喜歡的事.
有需要的時候亦會互相守望.
良好的家庭,良好的關係和同伴.
幫山姆他能夠跨過他的病患.
不會過份擔憂.
第三個原則是說.
你要持續向前看.
你要期待有更好的未來.
山姆他曾經承認.
也有過一段自怨自艾的時間.
但因為他很喜歡看卡通片.
他很深受和老迪士尼這個人所說的話影響而改變.
我永遠都會讓自己有期待的事物.
讓我有繼續奮力向前的動力.
使我的生命更加充實.
這也成為了山姆的座右銘.
從此他對每一天也充滿了盼望.
心境也變得越來越正面.
事實上弟兄姊妹山姆的生命故事.
可以說是應對苦難一種人生的智慧.
提醒我們與其在地上追求一個沒有痛楚的人生.
或者你對著一些困難,苦難左閃右避.
倒不如你去培育一份能夠承載著苦難的堅韌力.
畢竟你和我其實都會明白.
苦難它是會不同時間用不同的方式.
它不會問過你.
它就會闖入你的世界.
闖入你的人生裡面.
打亂了你一切的生活節奏和秩序.
所以懂得倚靠耶穌.
跨越面前的困難.
將你人生一些的痛苦.
視為你整全生命.
不可以缺少的一部份的時候.
這一個才是我們面對苦難的出路.
事實上主耶穌在踏上苦路之前.
他曾經跟門徒這樣說.

$^{481}$他說在世上我們是會有苦難.
但是我們可以放心.
因為他已經勝了世界.
所以每一年的復活節.
每一年的受苦節.
我們不單單只在紀念耶穌基督為我們犧牲.
其實我們更加要深入去思考.
我們怎樣和這位受苦的主來同行.
無論是我們個人.
或者這個群體所面對的艱難苦難.
其實他們首先都通過上帝允許.
才發生出現在我們的人生裡面.
在裡面我們可以怎樣去效法耶穌.
我們怎樣展現那份對上帝的愛和信服.
讓我們的生活仍然可以緊貼著耶穌.
變得更加的豐富和整全.
弟兄姊妹.
耶穌又或者今天我們看到.
跟隨耶穌這一班的婦女.
他們透過他們的苦難.
告訴我們甚麼才是豐盛的生命.
甚麼才是完整的生命.
不是生活裡面毫無掙扎或者苦楚.
而是擁有更加多.
從上帝而來的那種適應力.
那種堅韌.
就在苦難裡面.
你不要花時間去想.
你做不到甚麼.
而是像這班婦女一樣.
專注在當下我仍然可以為耶穌做到甚麼.
你會發現在稍後.
上帝會使用他們所擺上.
他們成為了第一批見證耶穌.
傳揚主耶穌復活信息的侍女.
這個是他們想都想不到.
另一方面在苦難裡面.
我們也會學習這一班的婦女.
來邀請你身邊的弟兄姊妹和你結伴同行.
特別是那些素齡成熟的弟兄姊妹.

$^{521}$你們在路上可以互相照應.
這樣你們就會有更多的力量.
來面對苦難帶來的衝擊和挑戰.
最後我們要認定.
上帝永遠是同在和掌管.
在苦難裡面我們要學會主動尋求上帝的介入和幫助.
並且我們要認信.
耶穌已經為我們預備好新天新地.
有一天祂會來接我們到祂那裡去.
我們就帶著上帝給我們的這份應許和盼望.
好好地來回每一天.
\newpage



\section{羅馬書 8:31}
\label{sec:9gw1hxqzLeo}
\textbf{2025年4月20日復活節崇拜直播｜陳明泉牧師｜勝利的體會｜羅馬書 8\_31 -39}
\newline
\newline
連結: \href{https://youtube.com/watch?v=9gw1hxqzLeo}{\texttt{https://youtube.com/watch?v=9gw1hxqzLeo}} ~~~~ 語音日期: 2025-04-20
\newline
\newline
\hyperref[sec:96uE5WjMIVg]{\small{< < < PREV SERMON < < <}}
~
\hyperref[sec:index_chronic]{\small{[返順時目]}}
~
\hyperref[sec:eeqdwBmS648]{\small{> > > NEXT SERMON > > >}}
\newline
\newline
羅馬書 8:31
\newline
\begin{longtable}{cl}
\hline
\hline
章節 & 經文 (和合本修訂版)\\
\hline
8:31 & \begin{tabularx}{0.7\textwidth}{X} 既是這樣，我們對這些事還要怎麼說呢？神若幫助我們，誰能抵擋我們呢？ \end{tabularx} \\ \\
[1ex]
\hline
\hline
\end{longtable}
$^{1}$接著是經訓時間.
今天我們選讀的經文是新約聖經羅馬書.
新約聖經羅馬書第八章三十一至三十九節.
聖經是這樣說.
「既是這樣,還有什麼說的呢?.
神若幫助我們,誰能抵擋我們?.
神既不愛惜自己的兒子,為我們眾人捨了.
豈不也把萬物和他一同白白地賜給我們嗎?.
誰能控告神所揀選的人?.
有神稱他們為異了.
誰能定他們的罪?.
有基督耶穌已經死了.
而且從死裡復活,現今在神的右邊也替我們祈求.
誰能使我們與基督的愛隔絕呢?.
難道是患難嗎?是困苦嗎?.
是迫不得已嗎?是飢餓嗎?.
是赤身怒體嗎?是危險嗎?是刀劍嗎?.
如經上所記,我們為你的緣故終日被殺.
人看我們如張載的羊.
然而,靠著愛我們的主.
在這一切的事上已經得勝有餘了.
因為我深信,無論是死,是生,是天事,是掌權的,是有能的.
是現在的事,是將來的事,是高處的,是低處的,是別的受造之物.
都不能叫我們與神的愛隔絕.
這愛是在我們的主基督耶穌裡的.
頂子會平安.
可否再用一種復活的心態來回應.
頂子會平安.
在二千多年前,基督耶穌從死裡復活.
耶穌第一樣對著門徒說:願你們平安.
我相信門徒回應聽到耶穌問的這句話.
大概他們在幾天裡經歷過山車一樣的情感.
但是看到復活的主,他們的心裡滿懷確信.
在二千多年來,基督教裡有很多的詩歌描述基督耶穌的受苦.
但是有更多的詩歌描述因著基督耶穌的復活.
他是得著榮耀的萬王之王,萬主之主.
每一個信靠他的頂子們,我們的生命都可以變得良麗.
從此我們的生命不再活在黑暗當中.
因為信靠這個復活的主.
我們人生的價值觀念,我們看善物.

$^{41}$又或者我們生命裡怎麼去普捨我們的人生.
通通因著基督耶穌完全改喚過來.
所以今天是很值得慶賀的日子.
今天也是我們心靈裡永遠要銘記的一個復活節.
今天跟大家去看的這段經文.
剛才Pelton幫我們讀的時候.
我們看到裡面寫到一個很豪邁,豪情壯語.
洋洋灑灑將整個福音信息呈現給大家.
不過這不單止是說一種聚得勝面的心靈體驗.
更加說到一個基督徒作為一個信靠主的人.
可以挺起胸膛,面對著世界上有很多的凶猛.
他都能夠有力量去面對.
保羅用勝利這個字眼來描述.
所有基督徒的生命都是一個德勝有餘的生命.
不知說起勝利,德勝有餘.
你怎樣看自己的人生.
你相不相信自己的人生是一個勝利的人生.
你會不會覺得自己是人生的勝利組呢?.
還是你覺得人生都是一個受害者.
或者我是被外面的社會,這個世界.
或者人的際遇來管轄.
有時候想發圍,但卻是無力.
不過在聖經裡卻是說到.
每一個基督徒是一個德勝者.
靠著基督耶穌我們德勝有餘.
這個德勝的體會.
不是純粹是說取得了什麼外在的得益或者成就.
更加是一種心靈的狀態.
我記得中學的時候有一位男孩子.
你當是我吧.
中學的時候有一個男孩子.
會戰運動會的時候.
他通常跑步,運動,全部都拿不到獎.
但每一次運動會的時候.
這個男孩子和他的同學.
一群人都一定會穿著某一個牌子的球鞋.
一個Tick那個.
知道為什麼嗎?.
他們說拍照漂亮,有型.
更加重要是什麼呢?.

$^{81}$說真的,每年的校運會.
未必拿到獎牌.
但穿著這對有Tick的.
知道那個牌子的意思是什麼嗎?.
勝利的意思,Nike,Nikkei.
就是希臘文的意思.
穿著一對勝利的球鞋.
心態上已經贏了.
贏不到比賽,但贏到心態.
贏到一張漂亮的照片,你說多好呢?.
很多人想著勝利的感覺是一種很主觀的.
我們靠外在的物件來定義自己.
不過,其實在勝利的體會裡面.
其實有幾個層次.
第一個層次是.
我們真的贏了別人和我競賽的對手.
有一個對手和我們走向同一個目標.
但到最後我比他更快.
或者比他更強.
或者比他更高.
這就是贏了別人的勝利.
這是在場上的一種比喻.
當然,輸贏有時候.
可以是正面的較量.
也可以是很骯髒的.
在運動場上你去比賽.
你贏了別人就不光彩了.
所以有時候這種勝利包含著輸與贏.
有時候是光彩的.
但也有時候可能是卑鬱.
這是第一個層次.
第二個層次,其實勝利是不是一定要有一個對手呢?.
後來有些人說.
其實勝利從來都不是純粹看對方.
是看自己的.
是要贏自己.
所以人裡面自己訂立了一些目標.
訂立了一些方向.
我完成了.
我就會有一種正面的成就感.

$^{121}$這也是一種勝利的感覺.
這是第二個.
不是再跟別人比較.
是跟昔日的自己.
以及看自己的目標來做一些衡量.
當然,這個目標是按著你自己的實力來定下的.
但第三個層次,勝利就是一種自我的超越.
就是說,過往自己做不到的.
但我要將自己提升一個很困難的層次.
令到自己可以超越自我的限制.
那個極限.
當我去到那個境地的時候.
我達到某個時數.
或者做到某一樣東西.
這是過往我的經驗裡面做不到的.
這個就是超越自我的那種勝利的體會.
第一種就是對人比較.
第二,第三種就是跟自己比較.
不過到了最終.
或者聖經裡面所強調的.
是第四種.
這個不是一種比較的勝利.
而是分享基督耶穌的生命.
是一種靈性的轉化.
因著基督耶穌死而復活的那種生命.
他的德性.
他經歷生命的路.
我們就將我們的生命.
跟他一起走過這一段路.
藉著我們參與在基督耶穌的生命當中.
我們都能夠體會到.
耶穌基督戰勝死亡那份喜悅.
即或人生處於什麼制度.
什麼處境.
又或者我們的生命處於什麼狀況.
但是如果我們跟上帝.
跟這位我們所信的主基督耶穌.
緊緊扣連的時候.
我們就能夠體會到.
什麼叫做勝利的那種體會.

$^{161}$今天那段經文.
其實都是說這件事的.
我已經濃縮了大概說這件事.
不過具體怎樣去實踐到.
或者讓我們去明白.
這種勝利的體會呢.
今天希望用兩個段落.
我們重新去理解一下.
究竟上帝所說的勝利.
是一回什麼事.
我們先祈禱求上帝幫助.
天父上帝願你幫助我們.
讓我們能夠在你的話語當中.
我們讀到生命的亮光.
我們很需要你與我們同在.
我們每一個虛心在你面前.
領受你對我們的教導和勸勉.
幫助我們眾弟兄姊妹.
每一個聆聽的時候.
對著我們的內心說話.
讓我們聽得明白.
又在生命當中呈現你的話語.
又讓宣講的說得清楚.
願你的話語能夠清晰地帶給弟兄姊妹.
讓我們一同造就.
我們的生命被你建立和激勵我們.
我們恭敬將來領受你的話語的時間交託.
願你賜福.
獻上禱告奉求基督耶穌得勝的聖名.
Amen.
我們一起看看羅馬書.
羅馬書其實是一個不容易看的書卷.
我簡單來說.
我做一個簡單的總結給大家知道.
羅馬書第一章至第七章.
這是保羅第一生的經歷.
他這個神學的陳述.
或者說的是上帝的救恩.
上帝將人稱為義.
整個神學觀念.

$^{201}$他不是基於他自己是一個律法師.
加瑪列門下.
或者按照律法來說.
而是按照他自己生命中.
他對律法的體會和人生的體會.
所以在第一章至第七章.
其實是說他對人性的一份掙扎.
然後經歷過掙扎之後.
他明白基督的心腸是怎樣.
上帝的救贖.
藉著耶穌的受苦和復活.
就叫人不再被罪所控訴.
所以第一至第七章.
某程度上就是建立這個理論.
用他第一生的經驗.
再加上教義的思考.
就濃縮成為第一至第七章的總結.
而去到第八章.
某程度上就開始進入一種.
更加清晰的立論的部分.
在第八章第一節和第二節裡面.
是這樣說的.
如今那些在基督耶穌裡的就不定罪了.
因為似生命聖明的律.
在基督耶穌裡釋放了我.
使我脫離罪和死的律了.
這兩句經文說的是甚麼呢.
他說到人因為信主.
人就在基督耶穌裡面.
有兩個很重要的核心.
第一就是我們不被定罪.
不被定罪是甚麼意思呢.
罪的結果就是死.
在這裡面說到我們不被定罪.
就是當我們信靠主的時候.
我們生命就不會進入死亡隔絕的境地.
所以不被定罪的第一件事說到.
我們不需要承受罪的刑罰.
當信主耶穌.
我們就不需要面對罪的那種痛苦.

$^{241}$第二個就不單止是說當下舊時的罪.
被免去或者被赦免.
不被定罪.
而是接下來說後半段的生命的律.
舊時每一個世上的人都活在罪的捆綁當中.
有些事你不想做的.
但偏偏去做.
又或者有些事你想做的.
又沒有能力去做.
保羅說這個真是痛苦.
人就被這種罪的律.
合制住心靈.
整天都是唯一自己.
良心又責備自己.
又充滿了很多的遺憾.
不過因為信耶穌.
因著基督耶穌的得勝福音.
人亦都可以不再活在舊有的這個律裡面.
我們可以用一個聖靈賜給我們新的生命的律.
我們每一個能夠可以按著上帝的帶領.
我們有能力可以不犯罪.
我們見到心裡面很期盼做的美好的事.
我們有力量來實踐.
這些力量是怎樣來的呢.
在經文裡面就說到.
那種力量是來自叫耶穌死而復活的聖靈.
那個力量是賦予在地上的每一個基督徒生命當中.
所以我們生命裡面.
每一個基督徒其實都有一種生命力.
在第八章第十一節裡面這樣說.
然而叫耶穌基督從死裡復活者的靈若住在你們心裡.
那叫基督耶穌從死裡復活的也必藉著住在你們心裡的聖靈.
使你們必死的身體又活過來.
八章十一節裡面說到.
上帝的靈聖靈激活了我們每一個.
這種活力令到我們人生裡面充滿了期盼.
以前如果你未信主.
所有事情都很悲觀.
所有事情你都無能為力.
但當你信了主之後.

$^{281}$未必一下子可以變身.
馬上就很有力量.
不過你知道你的能力感會強了.
你知道你有信心去面對生命所有的艱難.
因為這份能力不是自己提出來的自信.
不是訓練回來的.
不是你的自我形象或者你想通了.
而是有一份由神而來.
聖靈內在你的生命當中的力量.
綻放在你生命當中.
所以弟兄姊妹.
你有時候要值得驗證一下.
你問一問自己.
你生命裡面你信了主.
不論長短日子.
你覺得你自己生命有沒有這份活力呢.
你覺得你自己信了主之後都依然垂頭喪氣.
還是充滿了期盼 充滿了可能.
心裡面你帶著無比的積極和正面.
來迎接生命的挑戰.
即或生活有很多艱難.
但當你信了主的時候.
你依然擁有無比的力量.
你就可以按著上帝給你的感動.
來跟這個世界的律.
來有一些很不同的生命表現.
我記得以前有位我認識的宣教士.
他的媽媽 宣教士的媽媽.
都是一個很敬虔的老姐妹.
我還記得她.
其實一個老人家她不是很懂得字.
不是很懂得像年輕人一樣.
世界這麼闊.
不過她是一個很忠心的人.
雖然她身體很軟弱.
去到八十多九十歲.
身體很軟弱.
不過她說她有空.
她就拿著一張單張在街上派.
可能你會覺得這個很老套.

$^{321}$或者這個不合時宜的一種方式.
不過對於一個老人家.
對他來說.
我做到的就是這些.
除了袋套之外.
我能夠送東西給人.
送張單張給人.
你只要肯做就行了.
但對我們今天來說.
我們覺得不想做.
我們覺得不行.
甚至乎颳風下雨.
或者有一些東西.
我們經常都有一種被迫的感覺.
但這位老人家很有趣.
他無論風雨.
或者身體多麼軟弱也好.
除非他真的差到不得了.
基本上他每個星期.
他就走遍不同的地區.
買一疊單張在街上派給人.
人們會笑他.
或者不理他.
忽略他.
但對他來說.
上帝的靈在他生命當中.
果效如何.
問他能夠派多少.
多不多反應.
他說未必收.
那你會否很感到失落.
他說做了上帝要我做的事.
我就喜樂.
還問這麼多的效果.
我只知道.
我作為一個九十多歲的老人家.
我做到的事.
上帝叫我做.
我盡力去做.
我做到.

$^{361}$我的心裡面.
我就覺得很有力量.
而且我覺得很開心.
弟兄姊妹今天不知道你的生命裡面.
會否有這種力量.
保羅說當聖靈.
內在每個基督徒的生命當中.
你就會有這種呈現.
這種活力呈現出來.
當我們知道這個框架的時候.
其實整個框架裡面.
他講了兩件事.
第一就是在法庭上.
上帝辯護這個人.
這個人無罪釋放.
所以這是一個不定罪的福音.
另一方面.
上帝確確保這個.
舊時被人唾棄.
看不起.
甚至乎在當代裡面.
有很多基督徒有身份危機.
他都不知道自己想怎樣.
或者被那些同族人排他的時候.
上帝說.
直著聖經裡面說.
他做到他自己生命逆流而上的能力.
即或這個世界唾棄他.
但他都能夠有逆流而上的那種活力.
這個律就是在第八章裡面.
由第一節至到第三十節.
籠統地就是一個這樣的框架.
接下來他再繼續說.
由三十一節到三十九節.
某程度就是這個段落裡面一個小節.
三十一節至到三十四節.
既是這樣.
還有什麼說呢.
神若幫助我們.
誰能抵擋我們呢.

$^{401}$神既不愛惜自己的兒子.
為我們眾人寫了.
豈不把萬物和他一同白白地賜給我們嗎.
誰能控告神所揀選的人呢.
有神稱他為二了.
誰能定我們的罪呢.
有耶穌基督已經死了.
而且從死裡復活.
現今在神的右邊也替我們祈求.
幾節經文.
但是說了一些很豐富的神學資源.
讓我們每一個人拿來放在生活當中.
第一件事他說到.
上帝要幫助的我們每一個人.
有誰能抵擋我們呢.
抵擋的意思.
當然對信仰或對信靠上帝的人.
有一些可能不懷善意的態度.
那種抵擋未必一定是外在的.
有時候那些抵擋可能是我們內心.
今天很多基督徒很多頂子妹.
心靈經常打架.
做完一件事之後自己罵自己.
你不要覺得這是敬虔表現.
有時候做完一件事之後.
我們叫做發來發去.
做完一件事走了一步.
然後退回三步.
撞來撞去.
心靈就是有一種患得患失.
不平安的那種抵擋.
那種抵擋可能是外在.
又或者是你自己內裡.
在這段經文裡面保羅說.
如果有上帝幫助我們.
還有誰能抵擋呢.
《離直氣壯》說到我們的生命.
因為上帝稱我們為義.
我們就不再變得患得患失.
然後經文再繼續說.

$^{441}$上帝這個福音是一個慷慨的福音.
是一個施予的福音.
上帝成為了我們的供應.
是恩典的供應者.
祂藉著祂的愛子耶穌.
為我們捨棄生命.
這個捨棄生命然後再復活.
然後讓每一個信靠的人.
每一個人都一同得著這個豐盛的恩惠.
大家知不知道.
我們不單止是身體復活.
在將來的日子不單止是純粹說.
我們不用死這麼被動.
最底線是這件事.
但是我們還有高的那一線.
我們生命裡面會和基督耶穌.
分享祂的榮耀.
所以在今世裡面.
我們自己要淡定.
我們的人生怎樣走都好.
我們終極有一個冠冕等著你們每一個.
不過我們今天會不會帶著.
這一份這樣的期盼過我們的生活.
當我們生命裡面有這一份期盼的時候.
我們就朝著這個目標來進發.
我們的生命就變得不再一樣.
有一個這樣的生活小例子.
你當是一個小例子.
和這裡這麼驚悚.
很大落差的.
話說我神學院有一個同學.
他以前在他的事奉崗位裡面.
他教那些小學生的主日學.
那些小學生個個都很頑皮.
他第一天去到.
他差不多全散了.
他都覺得很痛苦.
他看完這段經文之後.
或者看完聖經之後.
他有一種感覺.

$^{481}$與其整天批評他.
說他做得不好.
不如稱讚一下他.
有些小朋友經常都很搞砸.
於是就吵吵鬧鬧.
然後我同學就靈機一觸.
就說你好乖 你坐下.
明明吵吵鬧鬧很搞砸的.
他就說他乖.
小朋友都呆了.
我這麼搞砸 你說我乖.
他說是 我看到你的心.
你好乖.
其實是回應他.
不過這個小朋友真的態度改變了.
過了幾年之後.
神學院的這個牧者告訴我.
當你說他是乖的時候.
他就朝著乖的方向去走.
他的生命是朝著好的方向去進發.
我們生命都是.
我們經常自己打自己.
我們經常敵擋自己.
我們經常說自己不好.
說自己很差.
我們差極有上帝救贖我們.
再差也好.
上帝的兒子已經為我們寫命.
有公義的官面為我們傳流.
如果當我們知道我們的生命.
是寶貴到一個地步.
是耶穌為我們寫命的時候.
我們就沒什麼條件再踩自己.
覺得自己不好.
上帝是我們恩典的供應者.
上帝都保護我們的心靈.
保護我們.
然後他再繼續說33節.
誰可以控告神所揀選的人.
沒有人有這樣的位置.

$^{521}$包括你自己都不可以控告自己.
質疑上帝對你的揀選.
上帝已經稱你為義.
沒有人可以定你的罪.
有基督已經死了.
而且從死裡復活.
現今在神的右邊也是我們祈求.
當我們知道生命裡.
我們在不理想的狀態裡.
但是上帝救贖我們.
去幫助我們.
我們就沒條件放棄自己的生命.
沒什麼條件.
也不要再成為控訴自己.
保定自己的一個人.
弟兄姊妹你知嗎?.
其實很多人來到教會當中.
很多都本身是基督徒.
不過可能生命裡.
每個人都有自己的故事.
每個人都有自己的遺憾.
因著這些過去.
這些可能不堪入目.
或者不濟的過去.
他覺得和上帝離開很遠.
他覺得沒有眼目見這位上帝.
不過保羅在這裡說.
當上帝要揀選這個人.
哪個人可以定他的罪.
哪個人可以控告他.
我們就只管坦然無懼.
去到上帝的面前.
更加重要的是.
上帝從來沒有放棄我們.
那個死而復活的上帝.
現今繼續活著.
這個活著的基督.
繼續為我們每一個人來代求.
如果你記得基督的代求.
他是為彼得在跌倒之前.

$^{561}$他說我已經為你代求.
你回頭過來.
你就要堅固你的弟兄.
耶穌的禱告.
是要守望我們的生命.
當基督都是這樣看我們的時候.
我們就帶著正面.
樂觀積極的心.
來接受這個恩典.
放過自己.
讓自己繼續活在恩典當中.
有力量走面前人生的路.
所以在這裡第一個.
說到上帝對我們生命是一個保障.
是我們生命的拯救.
他也是我們生命的代討者.
我們不要貿然.
不停在質疑自己.
也不讓外面的質疑.
來否定自己.
我們鐵錚錚地站立在地上.
有上帝的恩典與我們同在.
第二個段落.
由35節到39節.
誰能使我們與基督的愛隔絕呢?.
難道是患難嗎?.
是困苦嗎?.
是迫伐嗎? 是飢餓嗎?.
是赤身奴隸嗎? 是危險嗎? 是刀劍嗎?.
如經上所記.
我們為你的緣故終日被殺.
人看我們如張載的羊.
然而靠著愛我們的主.
在這一切的事上已經得勝有餘了.
因為我深信無論是死是生.
是天使是掌權的.
是有能的 是現在的 是將來的.
是高處的 是低處的.
是別的就做之物.
都不能叫我們與神的愛隔絕.

$^{601}$這愛是在我們主耶穌基督裡.
35至39節更加豪情壯語.
不過這裡的豪情壯語.
特別將轉向說到.
我們每一個是基督的愛對象.
這個世界沒有一個人可以與基督的愛隔絕.
生命裡患難 困苦 迫伐 飢餓 危險 刀劍.
都不能叫我們與基督的愛隔絕.
有時候 弟兄姊妹 你們知道嗎?.
人生有一些逆境際遇.
有時候我們會用這些際遇來定義自己.
是不是上帝還愛我.
當一個人面對著很多的難處.
生離死別 很多的難關的時候.
我們第一件事就會質疑.
上帝為什麼我要經歷這麼多的苦難.
這些刀劍 患難 困苦留在我身上.
是不是你撇棄我.
《胡若聖經》裡很多時候的哀歌都是說這件事.
我們用際遇來定義我們生命是不是被上帝所愛.
不過這裡面 剛才所說的一羅列 一系列的負面遭遇.
其實保羅自己第一身經歷過.
所以第一章至第七章是說他第一身的體會.
他傳道的經歷 面對著他恩著信耶穌.
面對昔日他擁有優厚位置的身份.
傾刻之間好像什麼都沒有了.
他沒有這些東西是不是代表上帝不愛他呢?.
保羅說這些負面的經歷.
不會叫上帝的愛與他隔絕.
因為愛不是等於好景 順景.
反而在逆境當中我們都能夠經歷上帝的愛.
又或者這麼說 無論順與逆都好.
順景或逆境其實只不過是一個處境.
但上帝的愛在順境可以愛你.
在逆境當中上帝都可以愛你.
你的人生無論你病患 疾苦 孤單 所有事情.
你都不會與上帝的愛隔絕.
保羅深信靠著愛上帝.
在這些一切的事上我們就得勝有餘.
得勝有餘這個字 剛才有Nike的字.

$^{641}$他再加多一個字就是Hooper.
Hooper就是說強烈或征服的意思.
Hooper L'Aquilo 得勝有餘.
用廣東話來說就是贏到開巷.
你的生命贏到開巷不單止是僅勝或慘勝.
不是悲壯地贏人生的仗.
是贏到很多很從容不迫地打贏這場仗.
打贏這場仗的核心是因為有基督的愛在我們生命當中.
我們就要問基督的愛代表什麼意思呢?.
基督的愛首先第一件事不是.
愛不是一種回憶.
基督的愛不是純粹用一種回憶.
那種浪漫來表達.
這句說話是一個神學家John Piper說的.
耶穌的愛不是愛的記憶.
愛的記憶是什麼意思呢?.
即是老婆經常質問丈夫.
我結婚和你結婚二十年了.
為什麼你不像以前拍拖那樣對我.
為什麼沒有初戀的那種感覺.
初戀的那種感覺就是愛的記憶.
但你結婚幾十年了.
我們還求初戀的感覺.
有時候也挺難為丈夫的.
基督的愛從來不是一種回到過去.
回看舊時的懷念.
你年輕的時候剛剛信主.
很甜蜜的那種愛.
那不是一個過去式.
基督的愛是一個現在進行式.
你當下也在經歷基督的愛.
不過愛的方式是不同的.
你父母小時候愛你的方式.
是會驅寒問暖抱你.
會餵你吃貓貓.
但當你長大成人.
用這個方式來對待你的時候.
你會覺得是對你一種羞辱.
在不同階段.
基督的愛就有不同的呈現.

$^{681}$但保羅許郭卻是說到.
基督的愛是一個現在進行式.
不是追回舊時那種感覺.
這種愛是一種怎樣的愛呢?.
是一種代求紀念.
上帝紀念你的愛.
所以耶穌基督不是一刻幫你代求.
而是持續不斷紀念你的生命.
在上帝的右邊.
復活之後都繼續提名.
點名為你和我每一個.
紀念我們生命不同處境的主.
基督紀念我們.
這種愛持續不斷在我們生命當中.
基督的愛第二個層次.
就是立約的愛.
立約的愛就是說到.
在他的救恩,歷史.
由昔日怎樣對以色列人.
到今天怎樣對我們這班外邦的信徒.
上帝答應過你.
他對你那份不離不棄.
和你立約的態度.
他今天都會手約.
私前愛在我們每一個弟兄姊妹心靈當中.
第三方面.
基督的愛更加說到.
當我們見到基督耶穌的受苦.
我們知道我們今天的受苦.
不是說上帝很變態地.
外加一些慘痛經歷在我們生命當中.
基督的愛是說到.
我們生命的罪令到他要承擔這些苦.
以致到當我們面對苦難的時候.
我們都知道.
原來上帝是一起和我們走過這條路.
我們某程度在苦難的裡面.
我們明白基督耶穌的心腸.
具體說一個例子就是.
可能很多弟兄姊妹做了父母的時候.

$^{721}$面對著生活.
自一個家的艱難.
就開始明白爸爸媽媽那種偉大.
當你經歷過痛苦,難過.
你就明白對你好的人那種心腸.
基督的愛在這裡.
對保羅來說是很濃烈.
他說到耶穌的受苦.
和他自己的受苦是相關.
當他一直受苦的時候.
他覺得他正在分擔.
正在分享基督耶穌的使命.
他覺得這個偉大的證明和他有關.
他就帶著這份感恩的心.
繼續的來向前.
弟兄姊妹.
今天這段經文.
我們大致都明白掌握了.
在今天這個世代.
不只是香港.
很多人都垂頭喪氣.
很多人都被很多問題打倒.
我們很多人.
未必去到最窮困,最艱難的境地.
但是我們的感受.
我們經常都好像落入一種情緒的低谷.
這些情緒都是正常的.
我們都不需要否定這些情緒.
不過保羅在基督耶穌.
從死裡復活的教義裡面.
他將一個生命的價值觀告訴我們.
我們生命除了可以自怨自艾.
或者還是覺得自己活在痛苦裡面.
上帝為我們開啟一條新的出路.
當我們看著基督耶穌的生命.
由最低點到最榮耀的高升.
同樣這也是我們生命的寫照.
我們今天生命無論落入什麼境地.
上帝說這是我揀選的人.
我們沒有條件再去自我質疑.

$^{761}$又或者被別人質疑自己.
你只管挺起胸膛走你信仰的路.
有基督的愛在我們生命當中.
我們的生命有重要性.
我們在上帝眼中不是微塵一粒.
而是祂用祂的重價救贖我們.
我們明白基督的愛.
我們知道眼前這些苦難.
是我們生命的一部分.
我們更加有份於耶穌的救恩當中.
我們一起來努力.
放開我們的心懷.
我們放鬆我們緊皺的眉頭.
提起發酸的手.
肩撞起我們那雙軟弱無力的腳.
我們就奔向面前.
充滿了復活盼望的生命.
求主繼續幫助我們.
\newpage



\section{創世記 22:1-14}
\label{sec:eeqdwBmS648}
\textbf{2025年4月27日崇拜直播｜林嘉鴻牧師｜生命的最重要｜創世記22\_1-14}
\newline
\newline
連結: \href{https://youtube.com/watch?v=eeqdwBmS648}{\texttt{https://youtube.com/watch?v=eeqdwBmS648}} ~~~~ 語音日期: 2025-04-27
\newline
\newline
\hyperref[sec:9gw1hxqzLeo]{\small{< < < PREV SERMON < < <}}
~
\hyperref[sec:index_chronic]{\small{[返順時目]}}
~
\hyperref[sec:u7h0Nna94Mk]{\small{> > > NEXT SERMON > > >}}
\newline
\newline
創世記 22:1-14
\newline
\begin{longtable}{cl}
\hline
\hline
章節 & 經文 (和合本修訂版)\\
\hline
22:1 & \begin{tabularx}{0.7\textwidth}{X} 這些事以後，神考驗亞伯拉罕，對他說：「亞伯拉罕！」他說：「我在這裡。」 \end{tabularx} \\ \\ \relax
22:2 & \begin{tabularx}{0.7\textwidth}{X} 神說：「你要帶你的兒子，就是你所愛的獨子以撒，往摩利亞地去，在我指示你的一座山上，把他獻為燔祭。」 \end{tabularx} \\ \\ \relax
22:3 & \begin{tabularx}{0.7\textwidth}{X} 亞伯拉罕清早起來，預備了驢，帶著跟他一起的兩個僕人和他兒子以撒，劈好了燔祭的柴，就起身往神指示他的地方去了。 \end{tabularx} \\ \\ \relax
22:4 & \begin{tabularx}{0.7\textwidth}{X} 到了第三日，亞伯拉罕舉目遙望那地方。 \end{tabularx} \\ \\ \relax
22:5 & \begin{tabularx}{0.7\textwidth}{X} 亞伯拉罕對他的僕人說：「你們和驢留在這裡，我和孩子要去那裡敬拜，然後回到你們這裡來。」 \end{tabularx} \\ \\ \relax
22:6 & \begin{tabularx}{0.7\textwidth}{X} 亞伯拉罕把燔祭的柴放在他兒子以撒身上，自己手裡拿著火與刀；於是二人同行。 \end{tabularx} \\ \\ \relax
22:7 & \begin{tabularx}{0.7\textwidth}{X} 以撒對他父親亞伯拉罕說：「我父啊！」亞伯拉罕說：「我兒，我在這裡。」以撒說：「看哪，火與柴都有了，但燔祭的羔羊在哪裡呢？」 \end{tabularx} \\ \\ \relax
22:8 & \begin{tabularx}{0.7\textwidth}{X} 亞伯拉罕說：「我兒，神必自己預備燔祭的羔羊。」於是二人同行。 \end{tabularx} \\ \\ \relax
22:9 & \begin{tabularx}{0.7\textwidth}{X} 他們到了神指示他的地方，亞伯拉罕在那裡築壇，把柴擺好，綁了他兒子以撒，放在壇的柴上。 \end{tabularx} \\ \\ \relax
22:10 & \begin{tabularx}{0.7\textwidth}{X} 亞伯拉罕就伸手拿刀，要殺他的兒子。 \end{tabularx} \\ \\ \relax
22:11 & \begin{tabularx}{0.7\textwidth}{X} 耶和華的使者從天上呼喚他說：「亞伯拉罕！亞伯拉罕！」他說：「我在這裡。」 \end{tabularx} \\ \\ \relax
22:12 & \begin{tabularx}{0.7\textwidth}{X} 天使說：「不可在這孩子身上下手！一點也不可傷害他！現在我知道你是敬畏神的人了，因為你沒有把你的兒子，就是你的獨子，留下不給我。」 \end{tabularx} \\ \\ \relax
22:13 & \begin{tabularx}{0.7\textwidth}{X} 亞伯拉罕舉目觀看，看哪，一隻公綿羊兩角纏在灌木叢中。亞伯拉罕就去牽了那隻公綿羊，獻為燔祭，代替他的兒子。 \end{tabularx} \\ \\ \relax
22:14 & \begin{tabularx}{0.7\textwidth}{X} 亞伯拉罕給那地方起名叫「耶和華以勒」 。直到今日人還說：「在耶和華的山上必有預備。」 \end{tabularx} \\ \\
[1ex]
\hline
\hline
\end{longtable}
$^{1}$今日選道的經訓是記在舊約《創世記》二十二章一至十四節.
經文是記載亞伯拉罕獻以殺的事跡.
「這些事以後.
臣要試驗亞伯拉罕就呼叫他說.
亞伯拉罕他說.
我在這裡神說.
你帶著你的兒子就是你獨生的兒子.
你所愛的兒子往摩利亞地去.
在我所要指示你的山上把他獻為反制.
亞伯拉罕清早起來.
被上路帶著兩個僕人和他兒子而殺.
也劈好了反制的柴.
就起身往神所指示他的地方去了.
到了第三日.
亞伯拉罕舉目遠遠地看見那地方.
亞伯拉罕對他的僕人說.
你們和奴在此等候.
我與童子往那裡去拜一拜.
就會到你們這裡來.
亞伯拉罕把反制的柴放在他兒子以殺身上.
自己手裡拿著火魚刀.
於是二人同行.
以殺對他父親亞伯拉罕說.
父親啊.
亞伯拉罕說.
我兒我在這裡.
以殺說.
請看火魚柴都有了.
但反制的羊羔在哪裡呢.
亞伯拉罕說.
我兒神必自己預備作反制的羊羔.
於是二人同行.
他們到了神所指示的地方.
亞伯拉罕在那裡築壇.
把柴擺好.
捆綁他的兒子以殺.
放在壇的柴上.
亞伯拉罕就伸手拿刀.
要殺他的兒子.
耶和華的使者從天上呼叫他說.

$^{41}$亞伯拉罕 亞伯拉罕.
他說.
我在這裡.
天使說.
你不可在這童子身上下手.
一點不可害他.
現在我知道你是敬畏神的了.
因為你沒有將你的兒子.
就是你獨生的兒子.
留下不給我.
不料有一隻公羊.
兩個扣在稠密的小樹中.
亞伯拉罕就取了那隻公羊來.
獻為反制.
代替他的兒子.
亞伯拉罕給那地方起名叫.
耶和華爾納.
直到今天.
人還算.
在耶和華的山上必有預備.
.
今天很開心有一個很熟很熟的.
廣州朋友在我們當中.
林家洪牧師.
相信在座的弟兄姊妹很多都認識他.
他也很熟的廣州人.
林家洪牧師現在是荃灣浸信會.
東涌福音堂的主任牧師.
我一來到.
他就叫我叫他做Rambo牧師.
我把時間交給Rambo牧師.
弟兄姊妹平安.
平安.
很感恩回到家中.
多謝岑牧師的邀請.
可以回到家中的感覺.
很熟悉.
見到很多阿姨.
見到很多很熟悉的叔叔.
覺得很親切.

$^{81}$回到家中的感覺.
好像走了船回來家中的感覺.
我不知道走了船是怎樣.
好像很久沒有回來.
回來後很熟悉.
感到很感恩.
見到弟兄姊妹.
大家都很精神.
每個人都說Rambo你又肥了.
又老了.
是的,事實是又肥又老了.
都離開了一段很長的時間.
回來後跟大家分享整個話語.
真的很感恩.
我們從一個禱告開始.
求主開我們的眼睛.
讓我們看到主你的奇妙和信實.
求你開我們的耳朵.
聽到你的教導.
聽到神對我們的說話.
求你開我們的心.
讓我們的信心被你建立.
願意信服於你.
緊緊跟隨你.
我們禱告奉耶穌基督的名字.
禱告,阿們.
今天跟大家分享一段經文.
大家都很熟悉.
如果回教會久了一定聽過.
我們看亞伯拉罕願意獻兒子以撒.
這個經歷其實真的不簡單.
所以歷代信徒都稱呼亞伯拉罕為.
信心之父.
這個名字真的不是浪得虛名.
當我們看到聖經的時候.
原來信心之父其實經歷了很多很多事情.
亞伯拉罕的名字叫做亞伯蘭.
亞伯拉罕之前的名字叫做亞伯蘭.
是上帝為他取名叫做亞伯拉罕.
我們看到亞伯蘭的生命.

$^{121}$其實不斷被上帝建立.
被神教導.
他也經歷過很多的試驗.
當中生命經歷過很多的甜酸苦辣.
他對神的認識其實是不斷增加.
他對神的信心也不斷被建立.
正正因為他這份信心.
神喜悅亞伯蘭.
他創世記十五章第六節.
說亞伯蘭信耶和華.
耶和華就以此為他的兒.
今天我們看到亞伯蘭對上帝這份信心.
這個行動反映他對上帝是完全的信服.
上帝為什麼有這麼的吩咐.
真的很難理解.
原來神一個人有時候也會對他作出一些考驗.
這次的事件原來是神對亞伯拉罕.
做了一次很重要的考驗.
一個DSE(Decisive Exam).
第一部分是說什麼呢.
第一節和第二節.
他說神對亞伯拉罕做了一個很重要的試驗.
他說「這些事以後」.
神要試驗亞伯拉罕.
經文一開始提到「這些事以後」.
我們很容易忽略了這幾個字.
「這些事以後」其實之前發生了很多事.
很多事情發生.
我們看回《創世記》第十二章開始.
我們看到亞伯拉罕其實經歷了很多事情.
神一直教導亞伯拉罕.
一直建立亞伯拉罕這份信心.
在第十二章一開始的時候.
他看到耶和華口口賜福亞伯蘭.
當時亞伯蘭已經七十五歲.
已經很年長.
神要他領受很大的祝福.
他要離開本地本族父家.
去神所指示的地方.
他要成為大國.

$^{161}$要成為別人的祝福.
神要應許亞伯拉罕.
要應許他的土地賜給他的後裔.
在第十三章他這樣說.
神再次應許亞伯蘭.
他說他的後裔極其繁多.
如同地上的塵沙那麼多.
但是有一件事很困擾亞伯蘭.
是什麼事呢?.
他一直沒有兒子.
沒有兒子.
今時今日很多年輕夫婦都不會生小孩.
養隻貓養隻狗就可以了.
但是當時古代的世界.
兒子是非常重要.
代表著家族的繼承.
是上帝那種賜福.
亞伯蘭當時已經一把年紀.
七十幾歲了.
他仍然沒有兒子.
代表著沒有人承繼他的財產.
所以他到了第十五章.
亞伯蘭好像有點失望.
跟神說.
神啊!神啊!.
我既無子,你還賜我什麼呢?.
他有一點點投訴.
但其實也不是.
他可能跟神說.
你還賜我什麼呢?.
於是亞伯蘭的妻子.
撒拉.
當時叫撒奈.
嘗試用人的方法.
為亞伯蘭生一個兒子.
大家知不知道?.
他安排他的侍女夏格.
做他丈夫的妾侍.
代替撒拉生一個兒子.
最後夏格真的生了一個兒子.

$^{201}$叫做以實瑪利.
以實瑪利就是妾侍所生的兒子.
雖然亞伯蘭生了這個兒子.
但這個兒子並不是神所應許的.
後來這個家庭搞得一塌糊塗.
記載主母撒拉不喜歡夏格.
趕他們母子倆出門.
於是亞伯蘭就預備了一些水.
預備了一些糧食.
送走以實瑪利.
送走夏格.
其實他去了哪裡?.
他去了曠野.
真的隨時會死.
從此亞伯蘭再也見不到這個兒子.
第十七章.
耶和華向亞伯蘭顯現.
向他立約.
他說這個名字不要再叫做亞伯蘭.
要叫做亞伯拉罕.
因為神已納他成為多國之父.
多國之父.
所以亞伯拉罕的希伯來文名字叫做亞伯拉罕.
阿伯是爸爸的意思.
多國的爸爸.
國道是從理而出.
君王要從理而納.
神對他說.
撒拉要生一個兒子.
必有百姓的君王從撒拉而出.
撒拉所出的才是神納的約.
唯有撒拉所生的才是耶和華所應許.
不單止兒子這麼簡單.
是神所應許.
而且是耶和華和亞伯拉罕所納的約.
終於亞伯拉罕到了人生一百歲的時候.
他就生了這個以撒.
這個兒子以撒就是唯一的兒子.
神的應許等了很多很多年.
經過種種的辛酸.

$^{241}$應許終於成就了.
亞伯拉罕和他的妻子撒拉終於老來得子.
生了一個萬千寵愛在一身的以撒.
相信他們一定非常痛惜這個兒子.
是他們的心肝寶貝.
而且這個兒子也是代表神的應許.
代表神和他納的這份約.
所以短短幾個字.
「這些事以後」.
大家要看很長的經文.
由十二章開始一路看.
原來發生這麼多的事情.
他的兒子以撒瑪利已經走了.
現在他和撒拉只是年紀很大.
只剩下這個兒子.
這個以撒的兒子.
他是亞伯拉罕唯一的希望.
但今天的經文創世的二十二章.
亞伯拉罕他要面對人生一個很大的考驗.
相信是他人生最大的掙扎.
神吩咐亞伯拉罕.
要將他的命根.
將他心愛的兒子以撒獻上.
亞伯拉罕要作出一個非常困難的選擇.
他要怎樣面對.
要獻上他的命.
神要亞伯拉罕這樣做.
原來神有清楚的目的.
就是第一節這樣說.
神要試驗亞伯拉罕.
神要考驗他的信心.
是否完全信服上帝.
這一次亞伯拉罕是一個最大的考驗.
當我們被神使用的時候.
有時我們都會被試驗.
神都會有時試驗我們的信心.
我們能否忠於上帝.
以信心去仰賴我們的主.
我們要被上帝不斷建立.
我再說一遍.

$^{281}$上年我都經歷了一個很大的挑戰.
我上年八月尾的時候.
肺部感染了嚴重的肺炎.
因為兩年前我換了腎.
其實看不出來.
換了腎.
上年八月份感染了肺炎.
入了醫院住了一個星期之後.
情況越來越差.
入了深切治療部.
一直守望.
其實當時我都很擔心.
我真的怕出不回醫院.
但我九月的時候.
上年九月我安立了牧師.
我跟醫生說我九月要安立牧師.
那時是九月.
他說這樣嗎?你可否請假?.
不知如何請假.
跟醫生說不如你為我祈禱.
醫生說好啊,我為你祈禱.
很感恩.
在醫院裡感受到上帝對我的考驗.
真是考驗.
雖然過程很辛苦.
但我明白什麼叫信心.
原來信心不是我信就行.
原來是要將自己放在上帝身上.
由上帝帶領你.
雖然沒有力氣.
但我要知道神一定會帶領我.
帶領我走過困難的日子.
亞伯拉罕同樣經歷這份信心.
這份考驗.
雖然他不明白.
但他願意踏上.
所以第二節.
他就說.
神說.
你帶著你的兒子.

$^{321}$就是你獨身的兒子.
你所愛的以撒.
往摩尼亞地去.
在我所要指示的地.
山上裡面.
把他獻為梵濟.
大家知不知道什麼叫梵濟?.
梵濟是要將一個滯腎完全焚燒.
是毫無保留的.
將它燒盡.
神要吩咐亞伯拉罕.
要親手殺死.
這個非常難得的獨身兒子.
這個以撒.
將他完全燒盡.
我以前聽完這個故事.
覺得很難理解.
為什麼神要吩咐亞伯拉罕.
做這些事情.
這個要求是完全超越常理的.
是毫無道理的.
但我們的神.
原來是有他的目的.
亞伯拉罕當時.
可能以為上帝真的要他這樣做.
這個很容易理解.
因為當時亞伯拉罕在迦南地.
迦南地的文化.
是真的會將人獻祭.
當時當地的神明.
譬如巴力.
會將人獻祭.
這是古代人的理解.
他們會將自己生產的部分.
譬如農作物.
譬如一些動物.
將小孩獻為繁祭.
但唯獨以色列的神.
唯獨上帝.
他是獨一的上帝.

$^{361}$他多次向以色列人表明.
他是反對將人獻祭的.
這是耶和華所珍惡的.
與當時的外族的異教.
有很大的對比.
在經文當中我們看到.
耶和華絕對不需要.
以人獻為繁祭.
最後預備了一個供養.
代替亞伯拉罕的兒子.
作為繁祭.
所以我相信當時的亞伯拉罕.
很大可能真的以為.
上帝要他這樣做.
真的要將以色列人獻為繁祭.
要將他的兒子獻上.
他不知道這次是一個考驗.
他不知道.
他真的以為神要將.
他最心愛的心肝寶貝.
獻上去繁祭.
神三次跟亞伯拉罕.
說你要帶著你的兒子.
就是你獨身的兒子.
你所愛的以色列.
上帝連續三次.
去描述這個兒子.
若亞伯拉罕真的要想清楚.
這次行動是很大的犧牲.
是神對他的應許.
他願不願意信服上帝的吩咐.
當我們知道亞伯拉罕.
之前所經歷的事情.
就會深深體會這個以色列.
對他來說是非常重要.
但這次神要求亞伯拉罕.
帶著這個兒子.
不是去野餐.
而是將他的兒子獻上.
獻為繁祭.

$^{401}$這一刻亞伯拉罕面對這麼大的挑戰.
他怎樣去回應呢?.
所以故事第二部分.
就是第三節至第十節.
就是亞伯拉罕的信服.
亞伯拉罕願意信服上帝.
經文第三節.
不如大家幫我讀.
經文第三節.
準備,一二三.
亞伯拉罕請祖彼來.
比上帝帶著兩個人.
和萬子而上.
加上好料,繁祭的餐.
就你身我身所致使.
帶來的地方,去你的宿舍.
亞伯拉罕完全沒有跟神討價還價.
亞伯拉罕沒有跟上帝有什麼爭辯.
他一早起來.
就跟神說.
他願意清早起來.
就預備好路.
預備伯人和兒子以撒.
一同前往這個地方.
他沒有跟神說.
不如我多獻些牛羊給你.
他沒有.
他也可以跟神說.
以前你應許過我.
以後由以撒所生的.
從以撒生的才稱為你的後裔.
這是你的約.
如果獻為繁祭,這個約就斷了.
這個約就破壞了.
他可以跟神說很多理由.
但他沒有.
他清早起來就預備路.
預備好柴.
就去神所指示的地方.
如果細心去看.

$^{441}$去摩利亞山這個地其實很遠.
經文說要三天以上.
要三天.
到了第三天.
所以那個地方.
真的要走三天以上的路程.
這幾天.
伯拉罕的心情.
我想是相當痛苦.
可能神要他想清楚.
你是否真的毫無保留.
將以撒獻上給我.
他是不是真的要.
反悔.
神要這三天內想清楚.
這三天.
伯拉罕可能想起很多.
跟兒子一起的片段.
那種心酸很難理解.
之後伯拉罕對一輩人說.
你們在這裡等.
我和同志往那裡去拜一拜.
就回到你們當中.
之後伯拉罕就將繁殖的柴.
放在兒子的身上.
自己手裡拿著火和刀.
兩父子一起走上這條殺路.
這裡提到的同志.
其實不是我們K1 K2的小朋友.
原來同志在這裡已經十多歲.
如果猶太人的傳統.
30歲.
覺得應該是30歲.
無論如何.
他都是一個青年人.
他有氣有力.
所以可以背起柴.
這兩父子一起走這個痛苦的山路.
這一刻.
這兩父子二人同行.

$^{481}$父子一起走上這條痛苦的道路當中.
兩父子在路上說什麼.
這段經文.
以撒的說話不多.
但他有一個問題.
第七節.
他問.
以撒對他的父親伯拉罕說.
他說父親.
就好像我們現在叫爸爸.
爸爸.
他說.
火和柴都有.
但繁祭的羊羔在哪裡.
他可能自小都和爸爸一起獻祭.
都知道獻祭是怎樣.
有羊羔.
這次有火有柴.
沒有羊羔.
可能心裡面都有點不妙.
可能爸爸想.
你就是羊羔.
伯拉罕說.
兒子.
神必自己預備了繁祭的羊羔.
接著二人同行.
經文再一次提到.
二人同行.
兩父子一起走上信心的道路當中.
上一代和下一代.
一起學習信心之路.
信服的道路當中.
一起呈獻給上帝.
是很美麗的圖畫.
這裡給我們一些反省.
我們怎樣和下一代相處.
上一代和下一代.
都走上這條事奉的道路.
走上這條十字架的道路.
你怎樣和他一起走.

$^{521}$於是二人同行.
走上信心的道路當中.
伯拉罕深信.
神必親自預備.
神會親自預備.
我們很難明白.
當時伯拉罕在想什麼.
很可能希伯來書的作者這樣說.
他認為兒子將會復活.
可能伯拉罕真的相信.
兒子會復活.
即使死了也會復活.
也不一定.
但無論如何.
這只是他的想法.
他的期望.
他回問伯拉罕.
我真的不想兒子死.
他真的不想.
但他們到了神所指示的地方.
伯拉罕在那裡竹壇.
把柴放好.
捆綁兒子以撒.
放在壇上的柴.
我看到以撒.
他很信服.
剛才說到.
他不是一個小孩子.
他是一個青年人.
甚至二三十歲的人.
你不要說小孩子.
如果現在十多歲.
你想追也追不到.
如果以撒要反抗.
一定沒問題.
追不到他.
但以撒沒有這樣做.
以撒沒有說一聲.
他很靜悄悄.
他真的信服.

$^{561}$爸爸所說的.
要將自己獻上.
你問以撒.
他一定說.
我不想死.
我真的不想死.
但他沒有反對.
他沒有逃避.
他仍然很信服爸爸.
躺在柴上.
他那時怎會知道.
他知道.
即使父子知道有神的預備.
但這個行動要很大勇氣.
經文沒有詳細交代以撒的反應.
我看到父子一樣信服.
很可能亞伯拉罕.
小時候看到爸爸的獻祭.
所以知道這個獻祭.
真的將他燒盡.
他完全信服.
一同拜上.
故事高潮第十節.
亞伯拉罕伸手拿著刀.
要殺他的兒子.
真的要下手.
要殺他的兒子.
耶和華出手.
神化解了這次危機.
第三部分.
第十一至第十四節.
我看到神的介入.
神的介入和預備.
第十一節.
耶和華使者從天上呼叫他.
亞伯拉罕.
我在這裡.
天使說.
你不可在這童子身上下手.
一點不可害他.

$^{601}$現在我知道你是敬畏神的鳥.
因為你沒有將你的兒子.
即使是獨生的兒子.
留下不給我.
亞伯拉罕.
看到一隻公綿羊.
兩個潛在棺木中.
牽著羊.
代替兒子以撒.
將牠們獻為飯祭.
代替兒子獻上給上帝.
亞伯拉罕鬆了一口氣.
Hallelujah.
神藉著一隻公羊.
化解了這次危機.
兒子有了一條命.
爸爸也有了兒子.
表面上神吩咐亞伯拉罕.
要獻上自己的獨生兒子.
但我看不是.
實際上獻上的是亞伯拉罕.
他自己.
做父母會明白多些.
做父母會明白.
如果看到自己的子女病重.
在醫院當中.
其實父母也很心痛.
父母在醫院看到自己的子女.
如果病重的時候.
他寧願躺在醫院的那個是自己.
亞伯拉罕.
他要親手結束自己兒子的生命.
雖然他知道耶和華必預備.
但他心裡也很難過.
很可能寧願死的是他自己.
但他完完全全信服上帝當中.
甚至這個吩咐.
對他來說是不可理喻.
完全不合道理.
但他也願意去信服.

$^{641}$亞伯拉罕對上帝是完完全全信服的.
因為他的生命是敬畏神的.
他沒有將自己的命根.
他的兒子留下不給上帝.
他沒有.
耶和華這位創造天地的主.
在他生命當中永遠是最重要的.
他願意將生命的一切獻給上帝.
願意遵行神的命令和吩咐.
即使你那一刻覺得這個不合道理.
但他也願意去信服.
這個就是敬畏.
現在我們怎樣看我們的信仰呢?.
是否有平安和喜樂就足夠呢?.
這是好的.
但是否足夠呢?.
我們很少說敬畏上帝.
認為神只是很慈愛.
很有恩典.
是的,神很慈愛,很有恩典.
但我們要看到上帝的屬性.
上帝是一個可敬可畏的上帝.
我們對上帝有沒有一份敬畏的心呢?.
我們回來教會敬拜.
我們有沒有一份敬畏的心.
去敬拜我們的上帝呢?.
耶和華透過這一隻公羊.
代替了亞伯拉罕的獻祭.
是神為他的預備.
亞伯拉罕沒有失去這個獨生的兒子.
以撒也沒有因此失去自己的生命.
神預備了這一隻公羊代替了以撒.
亞伯拉罕的心情.
相信由谷底回到上來.
失而復得.
所以第十三至十四節.
看到神預備.
耶和華必預備.
於是亞伯拉罕就給地方起名叫.
耶和華以勒.

$^{681}$你看這個細字其實也有一個解釋.
就是耶和華必預備的意思.
耶和華以勒這個字.
原來還有另外一個意思.
就是耶和華必看見.
耶和華必看見.
他會看到.
耶和華必看見.
耶和華必定會預備.
我們有時看到.
其實不單止看到.
而是神會注意到.
神會看到.
神會為我們預備一切.
神會負責一切.
一開始神要試驗亞伯拉罕.
最終他能夠成功通過這個考驗.
神當然不是真正想要他的兒子以撒的生命.
神真正想要亞伯拉罕獻上的.
就是他對上主的心.
他對神的敬畏.
對神的信服.
那種倚靠.
弟子們.
同樣神是喜悅我們對他的信心.
喜悅我們那種生命.
願意跟隨上帝一生去倚靠神.
敬畏上帝.
這是神所喜悅的.
亞伯拉罕對神的敬畏.
成為他心中一個很重要的根基.
我們看到亞伯拉罕沒有去.
死去惟著神所賜的祝福.
他沒有.
即使萬過不捨.
他仍然沒有惟著這個獨生兒子.
反而他惟著的是什麼呢.
就是賜祝福的上帝.
他惟著的是賜祝福的上帝.
主耶穌亦都這樣教導門徒.

$^{721}$你們要先求他的國和他的兒.
這些東西都要加給你們了.
我們願不願意跟隨主耶穌的吩咐.
信不信他這個應許.
我們會不會只是要神的祝福.
我們只是要惟著自己的命根.
卻沒有惟著這位賜祝福的上帝.
什麼東西才是我們生命中最重要.
是我們生命中的耳塞.
是我們生命中的中心.
其實神豈不是已經預備了.
一個最寶貴的羊羔.
正正是阿巴拉罕所說的話.
第八節.
他說神必自己預備作繁祭的羊羔.
神必自己預備.
神愛我們到一個情況.
他甚至將他的獨生兒子.
耶穌基督賜給我們.
神親自成為這個羊羔.
他將自己獻上.
上星期是壽苦節.
是復活節.
就是紀念主耶穌基督的犧牲.
在基督教的傳統.
阿巴拉罕獻耳塞的事件.
正正是在預表.
父上帝獻上他的獨生兒子.
主耶穌基督.
天父已經為世人預備.
作繁祭的羊羔.
藉著耶穌基督十字架的犧牲.
成為一隻被殺的羔羊.
主耶穌在十字架上的刑罰.
代替了我們罪所受一切的刑罰.
神去預備一切.
神去犧牲自己.
去救贖我們.
流出他的寶血.
去寫命.

$^{761}$他要成為這個羊羔.
要獻的是我們.
我們犯罪的是我們.
我們要死的是我們.
但耶穌他甘願成為這隻羊羔.
去代替我們.
去成為這個受苦的基督.
所以保羅在羅馬書第八章三二節所說.
神既不愛惜自己的兒子.
為我們眾人卸了.
豈不把萬物.
和他同伯伯的賜給我們.
神為世人的緣故.
願意將他獨生的兒子獻上為祭.
成為世人的救贖.
這是神對世人最大的預備.
所以第八節說.
神必自己預備了犯罪的羊羔.
我們每一個作為主耶穌的門徒.
我們需要將.
上帝放在我們心中最重要的位置當中.
學習以信心去仰望我們的主.
去回應我們的神.
神愛我們.
願意將他獨生的兒子賜給我們.
我們如何去回應上帝 跟隨上帝呢.
神是我們的幫手.
還是我們的心靈安慰.
還是我們願意將一生跟隨這位神.
成為我們生命中最重要的位置.
我們願不願意信服神的帶領.
放下自己所認為重要的東西.
其實這重要的東西.
神是知道的.
神知道我們最需要什麼.
神知道.
甚至我們願意為上帝放下很多的一切.
願意納至將上帝成為我們生命中最重要.
最核心的位置.
我們願不願意將自己獻為犯罪.

$^{801}$所以保羅在羅馬書十二章一節說.
所以弟兄們.
我以神的慈悲勸你們.
將身體獻上.
當作活祭.
是聖潔的.
是神所喜悅的.
你們如此事奉.
乃是理所當然的.
所以保羅提醒我們要獻為活祭.
這個活祭就是好好地事奉神.
好好地跟隨上帝.
這是神所喜悅的.
我們如此事奉.
其實是理所當然.
你的一切需要.
其實神已經看見.
神已經看到.
神已經注意到.
上帝已經預備.
耶和華已立.
神已經看到我們的需要.
面對前面的道路.
可能有很多的困難.
很多的艱辛.
或者未來經濟方面.
可能面對一些困境.
世界很混亂.
但你願不願意跟隨上帝.
將你的手拖著上帝的手.
握著上帝.
這位賜祝福的上帝.
而不是握著自己.
認為重要的東西.
在任何的緊急情況當中.
無論順境.
無論逆境.
我們都願意敬畏神.
忠於神.
將神成為我們生命中.

$^{841}$最重要的東西.
成為我們最重要的上帝.
放在我們生命當中.
跟隨祂.
你願不願意.
我向你禱告.
天父上帝多謝你的恩典.
多謝你救贖我們.
我們雖然有很多的罪.
但你都願意.
當人犯了罪的時候.
你為我們死.
神你的愛就在此顯明了.
主啊求你幫助我們.
能夠緊緊的跟隨你.
這條道路裡面很多的崎嶇.
很多的隱憂.
很多的挑戰.
大神啊.
我們真的要將你.
放在我們生命中最重要的位置.
有時候我們做不到.
求你寬恕.
但我們要願意.
去跟隨你.
求聖靈幫助我們.
緊緊的去跟隨你.
緊緊的去依靠你.
抓緊你的恩手.
我們透過教堂.
奉耶穌基督的名字而祈求.
阿們.
好.
\newpage



\section{ }
\label{sec:u7h0Nna94Mk}
\textbf{2025年4月27日主日講道｜胡康偉牧師:天國的筵席與和平之子}
\newline
\newline
連結: \href{https://youtube.com/watch?v=u7h0Nna94Mk}{\texttt{https://youtube.com/watch?v=u7h0Nna94Mk}} ~~~~ 語音日期: 2025-04-28
\newline
\newline
\hyperref[sec:eeqdwBmS648]{\small{< < < PREV SERMON < < <}}
~
\hyperref[sec:index_chronic]{\small{[返順時目]}}
~
\hyperref[sec:Wk72Zd8NFhw]{\small{> > > NEXT SERMON > > >}}
\newline
\newline
$^{1}$.
早安,鄧英展會.
很開心,因為回到教會,這是我的家.
我自己計算日子的話,我1996年離開香港.
但我相信耶穌還要早一點,我是70年代.
我們有一個80年代,很多人來.
我是70年代,所以比較老一點.
每次回來都覺得很親切.
我覺得今天是一個很特別的日子.
不是偶然的日子.
你知道嗎?.
有些人坐的位置永遠是定的.
你坐我的位置,就好像犯了罪一樣.
我坐你的位置,教會有些人永遠坐我的位置.
但今天你坐在我旁邊,我覺得是非常特別.
你可否跟我說,我等了一個星期,就是要坐在你旁邊.
不只是坐在你旁邊,我要送福給你.
我很喜歡勵志.
我覺得太多消極的事發生在我們周圍.
有很多戰爭,政治,甚至人與人之間.
有很多問題發生.
但除了這些問題之外,我覺得有很多勵志的事發生.
有一間連鎖店在美國,叫做Dennis.
如果你在美國,很多人都會來Dennis吃東西.
因為便宜.
最近他推出了一個套餐,叫做NVIDIA Breakfast Bites.
4.95元不用早餐,在美國來說是很便宜的.
因為我們面對的通脹是很厲害的.
為什麼叫NVIDIA呢?.
我想大家都應該聽過NVIDIA的一個科創.
王仁勳剛剛去了中國大陸,我不知道他回來了沒有.
王仁勳在台灣長大,移居到美國讀書.
今天王仁勳應該是全世界其中一個最有錢的人.
如果你買了他的股票,你應該很發財.
Dennis為什麼要去做NVIDIA的套餐呢?.
就是紀念王仁勳.
王仁勳很有錢,但他去美國的時候不是那麼有錢.
而且他是讀書的.
所以他在一間餐廳裡面打工,就是在Dennis裡面打工.
他幫人洗碗,還有做侍應.

$^{41}$所以下次你去元記,在樓下元記吃東西.
千萬不要看低服侍你的那個.
他可能將來會很發達的.
所以你可能將來要靠他.
因為這樣的緣故.
所以Dennis就用他的.
Inspired by a legend built for hunger.
靈感源於自傳為飢餓而生.
飢餓,其實我上次也有講飢餓的.
我稍後也會再講.
飢餓,我們現在吃得太飽了.
你有沒有經歷過餓和飢餓.
我想我們在座當中應該沒有經歷過飢餓.
或者有一點點餓.
飢餓,有時候我們喜歡講希望.
但飢餓,我比較喜歡飢餓.
這個故事是一個所謂的勵志故事.
讓我們知道我們每一個人都有機會為神.
為神可以改變的命運.
這是講我自己.
我是一個棄嬰.
我沒有不平凡的恩典.
什麼叫不平凡的恩典.
我們每一個人都有恩典.
我們每一個人都會講.
但不平凡就是說.
當人聽到你的恩典的時候.
哇,神或者某件事這麼厲害.
可以令人一擲.
我出生的時候.
因為家人窮.
所以我媽媽出生之後.
我在上環長大.
在東華醫院出生之後.
我們剛剛跟朱長老聊天.
原來很多人都經歷過.
被媽媽吃很多白花油.
什麼藥要墮胎.
因為當時大家的家庭環境很貧窮.
我就是這樣.

$^{81}$但我媽媽沒辦法把我墮胎.
我出生的時候.
我媽媽當時已經45歲.
但我爸爸當時已經70歲.
厲害嗎?.
我爸爸厲害.
男人你行的.
I speak to you.
不要說你不行.
葉麗儀.
唱上海灘的葉麗儀.
她出生的時候.
她爸爸是80歲.
生她的.
所以我們有很多消極的思想.
還有我們有很多框框.
我們覺得某個年齡就不做這件事.
或者有一個環境.
你就不可以超越這個.
是可以的.
我媽媽把我出生之後.
結果我出生了.
因為太窮.
上星期沈牧師說.
有十兄弟姐妹舉手.
記不記得.
其實我沒有舉手.
因為我不是十兄弟姐妹.
我是十六兄弟姐妹.
所以我爸爸很厲害.
我是最小的那個.
我爸爸有三個老婆.
我是最小的那個.
因為家貧.
60年代的時候.
香港很貧窮.
所以把我扔了在上環東華醫院.
今天還在的醫院.
我就住在芽菜街.
可能你沒聽過.

$^{121}$正式叫栗打街.
為什麼叫芽菜街呢.
因為當時我們整條街都是摘芽菜.
摘芽菜然後賣給酒樓賺錢.
我在我們那條芽菜街的上面.
如果你知道.
這是很古老的故事.
就是拿打素醫院.
最原來第一間是拿打素醫院.
旁邊有一個叫合一堂.
今天還在的合一堂.
我就扔在醫院門口.
因為我大姐姐.
當時比我大18年.
所以18年她知道.
知道有弟弟出生.
看到媽媽一個人走回家.
弟弟呢?太窮了.
我們養不起.
所以我姐姐說.
不行,我一定要弟弟.
所以就追問媽媽.
弟弟在哪裡.
我媽媽就說扔在醫院門口.
所以我姐姐就想不通.
就跑去醫院撿我回來.
撿我回來.
因為大家都不想養我.
所以我姐姐就說.
沒有人要你.
所以沒有人願意跟你改名.
你的名字都是我改的.
但我從來沒有生我的媽媽的氣.
沒有生我的家人的氣.
因為我知道當時我出生.
會影響整個家人的生活.
其他的弟弟可能沒得吃.
一個棄嬰.
你看到我今天.
在敘利亞.

$^{161}$我們有一個叫希望之家的中心.
在去年12月.
這個政變.
敘利亞改變的時候.
就是那一天.
就是一個棄嬰.
我們那裡有一個.
叫希望之家的中心.
這對夫婦.
是一個回教徒夫婦.
人會影響人.
當你越積極的時候.
隔壁的人會很積極.
很自然.
但當你消極的時候.
所以如果你有depression.
千萬不要找一個depression的人.
你一定要找一個開心的人.
那個開心的人才能幫到你.
不要想著.
大家同病相憐.
越說你就越慘.
所以你要找一個開心的.
跟你聊天.
人會影響人.
有很多地方我們不能傳福音.
你不能直接跟他說耶穌.
但我們的生命.
可以影響其他人的生命.
雖然是回教徒.
雖然不信耶穌.
但他看到我們怎樣愛他的小孩.
怎樣教他的小孩.
結果父母聚會他也來.
就是那一天我們有聚會.
散會之後.
這對回教的夫婦.
走過一個垃圾站.
走過這個垃圾站的時候.
聽到一個聲音.

$^{201}$是從哪裡來的.
在膠袋裡.
找到一個氣嬰.
手裡有條帶子.
最後找到他的媽媽.
那個帶子.
就是那天他出生.
他發展於年敏.
他有四個小孩.
在敘利亞.
在黎巴嫩.
在很多人的家庭裡.
他們很喜歡吃餅.
不是喜歡吃餅.
因為這是他們唯一的生存.
一個大餅.
如果晚上有一個大餅.
一家人可以吃點東西.
可以飽來睡覺.
但有很多家庭.
連一個大餅都沒有.
所以很多晚上.
他們的小孩要經歷肚子餓.
這樣去睡覺.
他有四個小孩.
環境是很不好的.
但是發展於年敏.
看到一個氣嬰.
把他帶回去.
我們要找他的父母.
他找到他的父母.
打電話給他媽媽.
他媽媽跟他說.
你不要還給我.
如果你還給我.
我會再一次扔掉.
他就收留了他.
當然我們會對這個小孩有幫助.
這就是當日的氣嬰.
為什麼叫做天恩呢.

$^{241}$因為當我發現這件事的時候.
很受感動.
因為同理心.
我也是氣嬰.
所以看到這樣的氣嬰.
我就寫了幾篇幾個文章.
這個故事.
我鼓勵我們有一個這樣的氣嬰.
我們在美國教會有一個台灣的姐妹.
看了這個故事之後.
她說我一直哭著為這個氣嬰祈禱.
我也是.
一個氣嬰.
到今天.
我想我都不錯.
不是很好.
但是一個被拋棄的.
所以我對這個氣嬰有個期望.
我相信這個氣嬰.
可能比我更加好.
她雖然生長在一個回教國家裡面.
她將來會帶很多人信耶穌.
她也會信耶穌.
所以我會一直為這個氣嬰祈禱.
我會一直跟著她的腳踵去幫助她.
就在一個這樣的垃圾堆裡面.
最近大家都聽到.
敘利亞和黎巴嫩的戰爭很厲害.
這個是敘利亞.
我剛剛去完.
這些垃圾堆是很多的.
大概像港台那麼大.
還要大多一點.
一個家庭.
真的是一家人.
就在中間裡面開餐.
因為要撿垃圾.
總之要找一些能夠生存的地方開餐.
她就在這樣的地方裡面.
找到這個氣嬰的膠袋裡面.

$^{281}$這個就是她的媽媽.
我們台灣的姊妹.
很感動.
給了她一封利是.
寫了一些字.
天恩這個名字.
當然不是阿拉伯的名字.
這個姊妹特別.
她叫天恩.
所以我們也告訴她.
她的名字叫天恩.
她將來叫什麼名字都好.
她其實叫另外一個名字.
很美麗的故事.
雖然我們不知道她將來會怎樣.
人間有情.
這個世界裡面.
我們有很多消極的事.
在我們周圍裡面發生.
但也有很多積極的事發生.
我們需要更積極地去面對我們的人生.
這是我長大的六十年代.
我在上環長大.
如果大家在網頁裡面可以看到.
這是我師母.
我們也是這樣長大.
師母當時住在雞寮火速.
她住在觀塘附近.
所以她去觀塘.
當時香港有很多天台學校.
她在天台學校裡面讀書.
然後她在慈雲山長大.
當時香港有很多這樣的事.
我太太有接受一些幫助.
有人領養她.
她每年收到五元美金.
有些收到更多的錢.
有些禮物給她.
但這些人從來不會來香港.
整間學校都有這些被收養的.

$^{321}$但偶然才有一對美國夫婦.
走來探望這個.
這個我領養了很多年.
我會帶他們出去玩.
但我太太沒有這個幸福.
只能拿到五元.
我記得我小時候也是這樣.
在上環長大.
現在還有一個叫合一堂.
就是舊的拿達數醫院旁邊.
很古老的一間教會叫合一堂.
逢聖誕節,復活節.
有什麼節日的話.
我就會走上去那裡.
因為有禮物牌,奶粉牌,米牌.
我們會走上去.
我不是信耶穌.
但因為他們派.
我們就拿.
我想這些故事.
在我們50年代,60年代.
是很多的.
為什麼這些人會這樣對我們.
這個也是我的童年.
最近有一部電影《九龍城寨》.
這個是我的所謂三合會.
14K,和聖和.
我不是藝人.
所以沒有聽過這些東西.
我也經歷過神打.
我未信耶穌之前也學過神打.
除了黑社會之外.
很多壞事也做過.
透過巴木斯,姐姐,太太.
將我帶信耶穌.
這個是七聖五代.
這是我.
變了樣子.
變了形.
這是福音教會.

$^{361}$在西營門的福音教會.
另外那個不是.
是我老婆.
當年的我太太.
我們人生裡.
神會將我們放在環境裡.
令我們很不舒服.
誰想做個棄嬰.
誰想遇到困難.
沒有人想的.
但就算你是一個很現代化的社會.
很發達,很有錢的環境裡.
我們也不可以離開有困難.
或者離開一些不舒服的.
可能被人打完胸口又打背脊.
又被人攻擊.
一定會有這樣的事發生.
當我們有這樣的事發生的時候.
聖經裡面有這樣說.
你幫我讀.
你要專心養賴耶和華.
不可以靠自己的聰明.
在你一切所行的事上.
都要認定他.
他必指引你的路.
你要專心養賴耶和華.
這個很重要.
我們周圍裡面太多事物發生.
不是每一件事都是神的旨意.
不是每一件事都是對你有益處.
所以我們要專心養賴耶和華.
我做錯了很多錯事.
我一生裡面.
我不是謙卑地說這句話.
我真的做錯了很多事.
我得罪了很多人.
但神的恩典.
將一些錯事.
明明好像看見是錯的.
或者得罪了.

$^{401}$無法反究.
但神可以去改變.
我想每一個人在裡面.
做基督徒的一段日子裡面.
都經歷過.
專心養賴耶和華.
他可以將錯改為對的事.
但你要專心養賴耶和華.
沒有一個人有權並為自己作一個結局.
或者一個定局.
我們可能會看醫生.
可能科技很發達.
以前都不知道自己有癌症.
現在很快就知道.
我們周圍的人都有癌症.
或者特別是美國的公司.
沒有人情味.
你早上上班.
突然有封信給你.
你放下.
甚麼都沒有.
下午你收拾東西離開.
或者你工作.
或者突然收到醫生報告.
你有癌症.
你沒有了.
那個醫生報告是否一個最後的定局.
誰有這個權並.
能夠作最後的決定權.
所以你要專心養賴耶和華.
Who have the final say?.
是耶穌基督.
我今天講天國的筵席和和平之子.
在這之前.
因為上星期是復活節.
其實我對復活節有很多講章.
因為我覺得很興奮.
當然復活節我們很多時候強調.
當然是耶穌基督為我們的死.
但耶穌基督的死帶來很多積極的事.

$^{441}$我們有救恩.
我們有因為耶穌基督的緣故.
你和我的生命會改變.
耶穌基督我們要有一個好的事發生的話.
很多時候可能在前面.
要發生一件沒有的事.
耶穌基督的復活.
他怎樣可以復活.
他要死才能夠復活.
他怎樣能夠死.
他受很痛苦.
被釘死在十字架上.
但因為他前面所受到的事.
後面就帶來一個結果.
耶穌當復活那一天.
那一早在《約翰福音》裡說.
他去到門徒那裡.
當時門徒很懼怕.
因為怕受逼迫.
把所有人關在一個房間裡.
聖經裡說他們關在一個房子裡.
然後突然間耶穌基督出現在他們面前.
大家都很懼怕.
為什麼很懼怕.
因為耶穌沒有敲門.
沒有爬窗.
耶穌那時候有一個超越自然的身體.
穿過了牆壁和門或窗.
去到他們的中間.
他們看到耶穌的時候.
不知道他是誰.
但很難想像.
這些人跟隨了耶穌已經三年多.
怎麼會連耶穌的聲音都認不出來.
但聖經裡形容的是.
這些人認不出耶穌.
耶穌要讓他們知道.
我是耶穌.
但耶穌用什麼方法讓他們知道.
聖經裡說.

$^{481}$祂伸出了祂的手.
和把他的肋旁打開.
讓他們認識.
我就是耶穌.
當時祂說了一句話.
願你們平安.
這句話很重要.
在很短的時間裡.
耶穌說了兩次.
願你們平安.
今天這個世界裡.
沒有什麼東西能夠跟平安比較.
我們需要平安.
世人需要平安.
所以耶穌讓他們看.
我覺得這件事很奇怪.
像今天這樣.
如果我讓他們看我的手就很簡單.
但如果讓他們看我的肋旁.
就算我穿今天的衣服.
也要打開衣鈕.
好像有點不自然.
如果是以前穿的袍.
可能會有更大的動作.
為什麼耶穌要這樣做.
耶穌要讓他們知道.
我是真真正正的耶穌.
我有傷口.
我有一個疤痕在裡面.
疤痕.
一個傷口.
我在黑社會的時候.
我也有疤痕.
因為打架.
那是為自己而受傷的疤痕.
但我們在主裡面.
有沒有疤痕.
我想要細想一下.
我想每一個人.
你信了一段時間.

$^{521}$你也會有傷口.
所謂的疤痕.
你說疤痕是什麼.
從那天開始.
從耶穌基督復活.
把他的傷口給門徒看的時候.
傷痕就成為教會的標誌.
好像一個商標一樣.
商標很重要.
你和我都需要一個商標.
在這個世界裡面.
商標是很值錢的.
如果你有一個好的品牌.
好的名字.
我們有些名牌.
為什麼那麼值錢.
就是因為有名字.
名字就慢慢把它建立出來.
耶穌基督那時候.
那個標誌.
我們在世界各地裡面.
你看到那個大M.
大家都知道.
這個是麥當勞.
你上到麥當勞的時候.
你都不用想其他的東西.
你知道你進去麥當勞吃什麼.
吃漢堡包.
然後喝可樂.
基本上你都知道.
因為這是一個深刻印象的東西.
你和我.
在別人面前.
或者在教會.
我們以教會.
或者以個人.
我們有沒有一個商標.
有沒有一個.
讓別人認定.
好像耶穌基督.

$^{561}$讓他們看到.
從今天開始.
傷痕.
就是我和你的商標.
然後耶穌更加認定他們.
這個很重要.
然後聖經說.
祂吹一口氣給他們.
接受聖靈.
這個時候.
耶穌確認他們.
我們都要被耶穌所認定.
但是我們都要有一個商標.
清楚知道我們在哪裡.
耶穌基督就是這樣.
願你們平安.
帶給這個世界.
今天想和你們說一下宴客.
因為我去過世界各地.
有時候不同的地方.
其實有很多尷尬的事情.
因為不熟悉.
有些人情.
或者每一個國家.
甚至吃飯的時間都不同.
我記得早年我去大陸的時候.
那時候還很小.
我說的是六十年代.
七十年代.
或者八十年代初的時候.
因為大聖經給他們.
和教會.
農村教會就很客氣.
叫我們坐.
我們一坐.
他們就叫我們去哪裡坐.
我們很信服.
就去那裡坐.
但其實很不應該.
因為有先後次序和長幼之別.

$^{601}$所以聖經裡面說.
可不可以幫我讀七和八節.
耶穌見所請的客.
揀擇首位.
就用比喻對他們說.
你被人請去.
父婚姻的筵席.
不要坐在首位上.
恐怕有比你尊貴的客.
被他請來.
所以我們做客人.
要有禮貌.
有時候你要坐在後面.
等到有人請你去.
你才坐在尊位的位置.
這就是聖經所教導我們.
在一個飲宴裡面.
一個筵席裡面.
我們應該是這樣.
請你們的人前來對你說.
讓座給這一位吧.
你就羞羞愧愧地.
退到末位上去了.
你被請的時候.
就去坐在末位上.
好叫那請你的人來對你說.
朋友請上座.
那時你在同席的人面前.
就有光彩了.
OK.
再繼續下去.
一個筵席.
我們現在說筵席.
筵席通常有不同的崗位.
有客人.
有主人.
其實我覺得.
應該還有第三類人.
就是招呼客人的人.
當然請吃飯的.

$^{641}$很明顯是主人.
我們把這個比喻放在教會裡.
有些人.
特別在北美服侍.
是很辛苦的.
因為北美的人可能太富裕.
所以叫他們回教會.
要來來來.
很久沒回教會了.
你們來吧.
這些人就算30年來教會.
他們是什麼位置呢.
他們都是客人.
每個人來坐都要被人服侍.
當然還有很多其他的東西出現.
你在教會裡.
你的角色究竟是主人還是客人.
或者你是招待人的人.
一個大的筵席.
我們說聖經裡的那個.
如果請一兩個人吃飯.
就沒所謂.
幾個人自己煮都可以.
但在聖經裡說.
超過10個人.
20個人.
要做很多預備的工作.
所以耶穌基督在這個比喻裡說.
他預備了要請這些人.
要來的.
所以聖經裡說.
同籍的人聽見這話.
對耶穌說.
在神的國裡吃飯的人有福.
就是被服侍的人真的有福.
在這個大筵席裡請了很多人.
到坐席的時候.
就打發了僕人去請這些人.
請來的每一個人都齊備.
耶穌基督已經預備了這個很大的筵席.

$^{681}$預備我們每一個人來.
這些僕人聽到主人的命令的時候.
就去請那些人.
第21節請讀.
那僕人回來,把這扇刀告訴了主人.
家主就動怒,對僕人說.
快出去到城裡大街小巷.
領那貧窮的,殘廢的,乞眼的,瘸腿的來.
僕人說:主啊,你所吩咐的已經辦了.
還有空位.
這個筵席裡本來是預備了.
已經有請,請,選去預備這些人.
但這些人都不來.
聖經裡的比喻就說到.
舊恩本來是給猶太人.
主人再一次說.
你出去路上的,在尼巴的.
勉強他們來坐滿這個房子.
我告訴你們.
先前所請的人沒有一個得償我們筵席.
其實我們每一個人都會有一個特別的位置.
在這個筵席裡.
只是我們願意來參加這個筵席.
客人就邀請了.
結果各方面都準備好.
然後有很多貧窮的,殘疾的,乞眼的,瘸腿的.
這些代表了猶太人中的瑞利,妓女等等.
然後有第二批的人.
這批人在尼巴或路上的人.
代表了你和我.
所以在《士徒行傳》13章46節.
「自因你們猶太人棄絕者道.
我們就轉向外邦人去」.
這就是你和我得到的舊恩.
通常我們用筵席來說這個比喻.
但聖經裡在同一個章節裡.
是說到另一個筵席.
這個筵席和前面的筵席是很不一樣的.
第一個筵席只是很短.
12節到14節.

$^{721}$第二個筵席15節到24節.
所以說了很多細節.
怎樣邀請人來.
但是另外一個筵席就很簡單.
他說「不要請好朋友.
恐怕他們也請你們報答」.
不要請好朋友.
我覺得耶穌.
你是不是有點不懂得人成世故.
為什麼不要請好朋友.
我們也想請好朋友.
你待會想不想請我吃飯.
不是很想.
我們也想請好朋友.
我們請好朋友是怎樣的.
你請我一餐.
我請你一餐.
不會說一個人請.
十餐也有一個人付錢.
我們以前有的.
巴木斯每餐都付錢.
你們沒有這樣的福分.
因為我們小時候.
我們通常都是這樣.
禮尚往來應該是人情.
但聖經裡耶穌說.
不要請我們的好朋友.
什麼意思.
老實說.
我一年只來講一次.
你又沒有我講一次.
父親教會也是這樣.
或者你請我吃飯.
其實對我們人生沒什麼改變.
我們叫做.
Entertainment at the church.
我們大家去娛樂對方.
當然我們有很多好處.
我們彼此多點吃飯.
多點了解.

$^{761}$多點溝通.
但整體來說.
我們是好好好.
好在我們裡面.
對外面的人來說是沒好處的.
這是另類的現值.
有沒有人吃過這些東西.
這是在北京拍的.
不知為何有這麼多蠍子.
到現在還是.
有些烏鴉.
現在在中國或者泰國.
你也會吃到這些.
這是另類的現值.
不是你吃到的東西.
窮人沒有能力報答你.
要等到二人復活的時候.
所以你請我吃飯.
就算我請不到你吃飯.
可能你得到報答的東西.
也很少.
因為我們起碼認識.
有些感情.
所以起碼感情會回答你.
雖然沒有請你吃飯.
也是你付錢.
但當一個人請你吃飯.
他不能夠報答你的時候.
你才可以得到.
怎樣使人和睦.
使人和睦的人有福.
因為他必稱為神的兒子.
現值和使人和睦.
這是我最近在孟加拉.
剛去拍了照.
幾個孟加拉人和我聊天.
和平很重要.
平安很重要.
平安是我們怎樣創造出來的.
在以前.

$^{801}$我記得在某個時候.
有一個總統應該是卡特.
開始說這句和平對話.
和平對話的意思是.
不是在吵架的時候.
勸架.
叫做和平對話.
就是在吵架之前.
當你聞到味道.
知道有些事情發生.
可能引起一個衝突.
在之前的時間.
你去創造一個機會.
或者進入那個情況下.
去停止那個衝突.
要這樣的話.
你是需要聖靈的恩高.
絕對是需要聖靈的恩高.
你的聰明是有限的.
我記得有一年我在巴基斯坦.
因為我很喜歡拍照.
在一個墓地裡.
星期五.
每個人都在敬拜.
我走進去.
拿著相機去拍他們.
不知道我是誰.
我只是撞進去.
最後有人問我.
我放了一個腳架在裡面.
照著他們.
他們問我.
喂,你做什麼?.
你是記者嗎?.
我說不是.
我是遊客.
遊客就會拍照.
他們問我想做什麼.
我說你拍照來做什麼.
我說我從來沒見過這些東西.

$^{841}$我們家人很少見到這些東西.
中國沒有這些東西.
所以我想拍照給家人看.
我開始去談.
有一個負責叫做E-MOM.
他說E-MOM想和你談.
我說好啊,沒所謂.
我就和E-MOM談.
通常回教徒都會.
我個人的經驗.
都會問我同一個問題.
What is your faith?.
你的信仰是什麼?.
你看看.
一個中國人和一個巴基斯坦回教國家裡面.
他一定是有一個帽子戴在頭上.
你中國人嗎?.
共產黨來的.
你中國人嗎?.
拜佛的.
無神論的.
當他這樣問我的時候.
我沒有和他傳福音.
但你問我,我就要回答你.
我說我是一個基督徒.
我是一個牧師.
因為他的期待是這樣.
但我給他的期待是另一個期待.
就開始去聊天.
有很多事情去談.
除了我們宗教之外.
你還有很多人性.
你是爸爸,我是爸爸.
回教徒很保護自己.
因為他怕別人說你是回教徒.
你就一定是恐怖分子.
我說不是,基督徒也有不好的.
你是爸爸,我是爸爸.
你疼愛你的子女,我疼愛我的子女.
一樣的.

$^{881}$你想你的子女讀好的書.
我也是想我的子女讀好的書.
所以很多人性的事情.
我開始和他溝通.
溝通之後,他請我進去.
他問我有沒有時間.
進去我們的內室.
那個內室就是他們散會之後.
一些長老聚集的地方.
我跟著說.
我可以去你想去的地方.
我進去了內室.
他又請我坐.
我當然不知道回教的規矩.
他請我坐高位.
我坐在這裡.
我是遊客.
所以繼續把我的影片.
把我的提示板動起來.
拍攝整個東西.
他說遲些我們會有一個.
更高級的人.
另外一個很快就來到.
我要把我的位置移給他.
因為他比我更高級,更加honor.
我們開始說人性的事情.
最後,整個旅程我覺得很奇妙.
他請我星期五有時間過來.
可惜我星期五要走.
否則我應該在回教史裡講道.
接著我被邀請到另一個.
因為我們是一個比較新的民族.
去一個新的國家.
我被邀請到一個真正的peace talk.
在基督徒和回教徒之間.
因為他們有太多的爭執.
太多的殺戮.
所以他們的長老想大家去談.
當時我也被邀請.
老師,你可不可以來.

$^{921}$在一個peace talk裡.
你和我要成為和平之子.
peace talk.
要和平的話.
我想,為什麼我們會有爭執.
因為你對我對.
我對你錯,所以才和你有爭執.
但我們要創造和平的話.
我們就要把我們的對放下.
你可不可以把你的對.
跟別人說對不起.
當你願意這樣做的時候.
你就可以創造和平.
聖經裡說我們是世上的炎和世上的光.
可不可以讀經文.
你們是世上的炎,你們是世上的光.
你的光也當這樣照在人前.
讓他們看見你們的好行為.
便將榮耀歸你們在天上的府.
我們要做炎做光.
一定要有一個行為走出來.
我們行為是怎樣去對待別人.
除了我們在教會裡.
我們在教會裡.
彼此相愛,我不需要再強調這一點.
但我希望能夠分享到.
這個炎夕是什麼意思.
我們都是蒙召成為和好的橋樑.
所以我們叫橋樑.
要挺身而出,履行神聖的呼召.
成為和平之子.
以憐憫,寬容,尊重為基礎.
為世界帶來和平.
這是我自己多年來從開始的教會.
因為伯木斯一直對中國很有負擔.
所以我去了美國多年,去了外國多年.
我的負擔是沒有離開過中國.
所以我對中國有個期望.
或者對世界的華人有個期望.
無論你走到新加坡,台灣.

$^{961}$走到世界每個角落裡.
你可以看到華人是曾經享受過這個炎夕的人.
我是,我的太太是.
可能在座當中都很多曾經享受過這個炎夕的人.
我們是享受的人.
今天華人世界可以為這個世界做些什麼.
我們已經脫離了我們以前的貧窮.
如果我們要做回以前宣教士所做的事.
我們是綽綽有餘.
我們要擺設一個現實.
我在很多地方,特別是貧窮國家.
很難說奉獻,但我很喜歡說奉獻.
因為我覺得這是一個真理.
當我們有奉獻的時候.
你可以看到神在一個人身上的改變.
我一生就是這樣.
我的氣嬰在一個這樣的環境裡長大.
沒有讀書,直到現在.
你知道我很滿足嗎?.
我很開心.
因為我幫助了很多沒有飯吃的人.
我看到他們吃飯的時候.
看到他們吃三文治的時候.
我很滿足,很開心.
所以我希望將這份喜悅.
跟我認識的人分享.
這張照片是我在中國的農村裡廣道之後.
有一個應該是婆婆.
她跟我說,你不要走好不好.
你等我,我回家拿些東西來.
我要奉獻.
她就拿了這些東西.
後來我們拿來香港賣了.
六千多元港幣.
很珍貴,下面的戒指.
那些戒指不是新的戒指.
這些戒指或者這些飾物.
應該是幾十年前結婚的時候.
有的戒指,有的飾物.
她收藏了很多年.

$^{1001}$但在這個時間,她要將它奉獻.
我記得我在巴基斯坦一個貧民庫裡.
你要講真理,弟兄姊妹.
我不會怕在貧民區裡.
我可以在這裡講同一個道理.
貧民區是一樣的.
我挑戰你們擺設現值.
我也挑戰他們擺設現值.
現值不是你有多少錢.
我跟他們分享的時候.
甚至我直接跟他們說.
你巴基斯坦人跟我們中國人一樣.
你見到白人的時候.
你就會有一個美元的尺寸.
白人,美國來是很有錢的.
歐洲來更加有錢.
千萬不要這樣.
聖經是給所有的人.
貧窮和富裕.
所以我們的真理不是因為在某個國家去改變.
而是真理.
我覺得姐妹是很寶貴的.
這個媽媽拿了巴基斯坦的錢給我.
她說這個不是給僑娘的.
是給你自己用的.
我拿了,很難收的.
很難收這些奉獻.
但我知道我應該收.
我收了這些奉獻.
我只有一元美金.
但神預立了這一元美金.
不是你奉獻多少.
這是黎梵特的故事.
可不可以幫我讀這個.
先訴情由的,似乎有理.
但論寫來到,就拆出實情.
我們都被傳媒所渲染到.
我們看一些東西都覺得.
戰爭很恐怖.
這個地方不可以去,那個地方不可以去.

$^{1041}$所以結果我想搞很多短訊都沒有人去.
人們答應了,就算在美國也好,在台灣也好.
打仗你還去?.
結果都沒有人跟我去.
但是不是真的這麼亂,這麼慘呢?.
我不否定傳媒所說的東西.
但傳媒所說的不是全部的故事.
我們還有很多真實的東西.
如果我們要知道真實的東西.
我們要跟當地的人才能說出當地的故事.
所以我來之前,我們去了敘利亞和黎巴嫩.
其實打仗的時候我們都有去.
我們沒有停過在那邊的工作.
這幾個地方,如果在傳媒裡面你都會聽到.
在北邊的黎波里市,Tripoli.
最近靠近敘利亞.
黎巴嫩裡面有很多巴勒斯坦的難民營.
另外一個叫貝卡谷的地方.
還有一個叫巴勒貝克.
這幾個地方裡面都是以色列攻擊的主要對象.
因為我們在貝卡谷,第三個地方.
貝卡谷這個地方.
所以我們成年人裡面.
真的,今時今日就算打仗.
WhatsApp都可以收到.
有線路.
打開一條供,我們都很害怕.
電話,沒事.
然後,砰.
沒事,那邊離鄉.
那個地方爆炸,不是爆炸我們這裡.
周圍裡面,作為童工.
我們沒有任何東西可以說.
我們只能說,你離開吧.
離開,找個安全的地方.
因為我們不能夠做到生命.
因為生與死.
他們都跟我說,我們不離開.
他們無論是當地人.
他們都有歐洲護照.

$^{1081}$甚至我們當地的童工.
在打仗最厲害的時候.
我們童工去了埃及宣教.
回來的時候,他們拍了一張照片.
降落飛機的時候,在炸機場旁邊的時候.
經歷了很多這樣的事件.
童工就說,我們不可以走.
因為我們服侍他們這麼多年.
如果我們走,他們會怎樣看我們.
平安你就留在這裡.
不平安你就走.
所以這個時候,我們更加需要留在這個地方服侍他們.
這就是和平之子.
大家看新聞裡面,黎巴嫩的南部.
因為接壤著以色列.
被炸得最厲害.
這個地方離開以色列,南部.
聖經裡面說的推羅和西頓.
就是以色列的南部.
所以我們在聽這兩個地方.
都是我們很熟習的地方.
最近這次比較鬆一點.
我和太太開車去到最南.
差不多我們看到以色列這麼近的地方.
然後去到最北靠近最底下的地方.
所以看到的.
可能和你們所看到的新聞.
我覺得有一點不一樣.
人仍然是這樣去生活.
不是每個人都可以走.
可能,當然有逃難.
比如南部去北部.
所以我們中心,在黎巴嫩的中心.
我們要開放,所謂做shelter.
當時很多頂尖會感動.
對我們的奉獻.
所以我們去重建一些房屋.
和將我們的中心開門.
將一些被炸爛了屋的人.
或者沒有地方要找.

$^{1121}$安全的地方來住在我們中心裡面.
你看到這個是推羅南部.
大概只離以色列50公里.
很多都是這樣.
但是兩旁的建築物又沒有.
那些人繼續住在這裡.
你不要想著沒有人住.
這邊我覺得.
這個只是我其中一個.
你看到中間有個洞.
有兩個洞在這裡,三個洞在裡面.
根本上跟以前的戰爭不一樣.
一個炸彈穿進這個地方.
它只是炸這間屋.
沒有毀壞建築物.
那些人在樓上樓下照樣生活.
我剛剛去的時候.
那時候已經停戰.
所以比較.
有很多生活上的不方便.
我們本來是用GPS.
用Google Drive去看去哪裡.
所有這些東西都被以色列封鎖了.
Google那些全部沒有了.
只是認住那條路.
知道怎麼走.
有很多不方便.
這是另外一個地方.
其實這個地方是最重要的.
巴勒貝克.
這個地方曾經被羅馬統治過的神廟.
這是很多恐怖分子的地方.
因為很多人都不願意跟我去.
所以只有一些.
我想白人不怕死.
有些南非人,有些澳洲人.
這是我們幾個星期前去的.
去探我們安提加教會.
他們會有難民營.
去服侍他們,為他們祈禱.

$^{1161}$然後我們去了敘利亞.
因為我們忘記帶香港護照.
所以我們拿著美國護照進去.
以前我們都不敢.
因為美國和敘利亞這些都不敢.
但因為我們忘記了.
所以我們硬著頭皮.
神廟很奇妙.
所以你要試.
你不要說不行.
你願意踏出一步.
不是你所想像的.
我和我太太拿著美國護照.
進去敘利亞.
他們讓我進去.
但現在的政權對基督徒是很不好的.
以前的政權是保護少數民族.
因為他們都是少數民族.
但現在的政權是ISIS的後裔.
即是主要的人.
所以對基督徒和少數民族是殺戮的.
所以我們在那裡的童工和基督徒.
我們和他們聚會的時候.
他們表現到他們真的恐懼.
因為有些回教徒來.
為什麼你還不悔改.
真的當街這樣和他們說.
但因為我們的童工和我們說同樣的話.
我們不可以走.
我們的童工我從來沒有聽過他們哭.
過去這麼多年.
我覺得他們都很勇敢.
但這次是唯一一次.
當政變頭幾個月.
12月1月2月的時候.
我打電話給我的童工.
WhatsApp給他們的時候.
我的童工哭.
怎麼辦?.
十幾年了.

$^{1201}$我們堅持了十幾年留在這裡.
就是因為我們要愛我們的國家.
但為什麼要繼續這樣.
不知道前途是如何.
但敘利亞教會沒有放棄.
它不可以在外國做一個筵席.
但敘利亞教會繼續在國內為主義而言.
堅持它的信心.
繼續擺設他們的筵席.
這是我們的教會.
這是我住的地方.
夫婦的小組.
大家一起.
和你和我一樣.
過著這樣的教會生活.
彩虹在眼.
我看到很多.
我覺得這是神蹟.
為什麼穆斯林掩著頭.
走進我們的教會.
聽我們講道.
甚至拿一本聖經來看.
這對我來說絕對不是稀奇.
我每次去每個地方.
我都看到有這樣的人.
神還在做工.
神沒有停止.
彩虹可以再現.
這是我們青年團裡.
因為他很受到幫助.
他影響了父母.
把父母帶來.
和我們一起慶祝聖誕節.
或者去其他聚會.
這個.
十四歲十五歲.
兩姐妹.
對於我們來說.
十四歲十五歲.
應該有一定的見識.

$^{1241}$但因為她生長在一個非常傳統的回教家庭裡.
她的父母不容許她走出外面.
她大概兩個月前來到我們中心裡.
對於她來說.
是一個新的世界.
因為她從來沒有接觸外面.
她每天的生活就在難民營裡.
從來沒有拿起一支筆.
十四歲十五歲.
不知道寫自己的名字.
只知道自己叫什麼名字.
兩個月之後.
她懂得寫自己的名字.
突然覺得我存在.
你知道這些人的希望就是.
沒有了.
十五六歲之後.
就被父母找人嫁了.
她一生人就是生孩子.
就是這樣.
但我覺得很滿足.
我能夠教她寫一個字.
讓她認識一些東西.
這個六歲的小孩子.
大概也是兩個月前發生的.
很喜歡來我們中心裡.
我們有一百二十個小孩子.
每天來.
我們給他們一些教育.
然後教他們祈禱.
給他們一些聖經故事.
我們有影片給他們看聖經故事.
去教育他們.
每個星期有兩餐給他們吃.
這個很可愛.
已經來了兩年的時間.
突然她說.
我要回家跟先生說.
先生覺得很奇怪.
她一向都很喜歡來上課.

$^{1281}$為什麼要回家.
不斷跟先生說我要回家.
休息了之後就OK了.
休息之後說我要回家.
後來才發覺.
她昨晚沒有吃飯.
所以很餓.
她要回去.
她想找一些東西吃.
我們中心經常都有預備食物.
這個不是一個很特別的例子.
有很多時候都有小孩子.
是這樣的.
我們敘利亞中心也是.
他就是這樣.
所以我們預備一些食物.
隨時可以給小孩子吃.
對於我來說.
我是很滿足的.
我不會覺得這是一個負擔.
要為你做什麼.
又要為你籌錢.
我不會.
因為我很滿足.
因為我是一個棄嬰.
看到他們能夠吃飽.
我很開心.
我兒子不在.
他很多年沒回來.
吃到肥大.
他回來.
因為很多年沒吃過香港的東西.
我很喜歡.
你想不到的東西.
平時不吃的東西.
吉野家,大家樂.
他掛念的東西.
現在下午要去小雨天.
吃這些東西.
很開心.

$^{1321}$因為他的理由是沒有港幣.
我都要給錢.
但我看到他吃得很開心.
我不會因為我付錢而不開心.
因為這是我兒子.
他吃得開心.
我就開心.
我想你對你的子女.
或者父母,親人.
你都會很開心.
你能夠請父母去吃飯.
你會很開心.
但這個筵席.
是否為了我們的家人.
孝道當然需要.
但今天的挑戰.
耶穌基督給我的挑戰是.
我們要為一些我們不認識的人.
可能我們的鄰居.
可能他們不是因為沒有飯吃.
心靈很空虛.
他需要你為他擺設一個筵席.
去歡迎他.
讓他感覺到.
回到教會的時候.
這是不是我們做主人.
我討論過我們每一個人都是主人.
每一個人都很火熱去愛主.
別人進來的時候.
不覺得有客人在.
除了真正的客人.
這次第一次來.
你會覺得.
為什麼每個人都好像有歸屬感.
我們做一個主人.
做你喜歡做的事.
這是我上次有分享過.
一開始的時候.
你說黃仁勳.
飢餓.

$^{1361}$Legend built by hunger.
我們的傳奇就是因為我們有飢餓去追求.
一個台灣的窮學生.
去到美國讀書.
今天成為最有錢的人.
為什麼.
如果你沒有追求.
你沒有渴慕.
你不會有這樣的成績.
弟兄姊妹.
今天我不管你是在教會服侍.
或者你在外面工作.
你有工作或服侍你的生意.
這是不是你最喜歡的事.
我告訴你.
如果你是你最喜歡的事.
你會很開心.
你一定會有成就.
所有你看到的傳奇.
都是因為飢餓.
因為有個渴慕.
或者我們基督徒說得比較熟悉.
或者我們願意為這個意象而飢餓.
我不是將意象.
只是領受意象.
無論什麼時候發生.
我不能等.
這個意象要成就.
因為這個意象只是一個起步.
後來會更加多.
Stay hungry.
Stay foolish.
保持你的愚蠢.
當你的意象有目標的時候.
你看每一個傳記裡面.
都被人說是蠢的.
我上次說Steve Jobs.
蘋果的創辦人.
為什麼蘋果全世界的app都不同.
他的ISO.

$^{1401}$其他都不一樣.
為什麼他要這樣做.
因為他有個理想.
我覺得我做的事會比你好.
為什麼黃仁勳能夠做到這樣的事.
他有個理想.
他有個Hungry.
他不只是這樣.
他熱愛.
然後他沒有放棄.
他才能夠有這樣的成功.
弟兄姊妹.
你要成功.
為主而成功.
這些人.
我不需要再說.
一定經歷過很多的犧牲.
很多晚上不能夠睡覺.
無論去研究.
做所有的事.
可能睡覺的時間比你和我都要少.
他才能夠成功.
很多的犧牲.
甚至很多的攻擊.
復活節的真正意義.
耶穌的犧牲.
耶穌的釘痕等等.
在他身上.
求主幫助我們.
讓神跟我們每一個人說話.
當我們結束.
乌?.
(笑).
\newpage



\section{申命記 24:1-22}
\label{sec:Wk72Zd8NFhw}
\textbf{2025年5月4日主餐崇拜直播｜陳明泉牧師｜朝向應許之地－情理兼備的國度｜申命記 24\_1-22}
\newline
\newline
連結: \href{https://youtube.com/watch?v=Wk72Zd8NFhw}{\texttt{https://youtube.com/watch?v=Wk72Zd8NFhw}} ~~~~ 語音日期: 2025-05-04
\newline
\newline
\hyperref[sec:u7h0Nna94Mk]{\small{< < < PREV SERMON < < <}}
~
\hyperref[sec:index_chronic]{\small{[返順時目]}}
~
\hyperref[sec:eFUsT6910pQ]{\small{> > > NEXT SERMON > > >}}
\newline
\newline
申命記 24:1-22
\newline
\begin{longtable}{cl}
\hline
\hline
章節 & 經文 (和合本修訂版)\\
\hline
24:1 & \begin{tabularx}{0.7\textwidth}{X} 「人若娶妻，作了她的丈夫，發現她有不合宜的事不喜歡她，而寫休書交在她手中，打發她離開夫家， \end{tabularx} \\ \\ \relax
24:2 & \begin{tabularx}{0.7\textwidth}{X} 婦人若離開夫家以後，去嫁別人， \end{tabularx} \\ \\ \relax
24:3 & \begin{tabularx}{0.7\textwidth}{X} 後夫若恨惡她，寫休書交在她手中，打發她離開夫家，又或者娶她為妻的後夫死了， \end{tabularx} \\ \\ \relax
24:4 & \begin{tabularx}{0.7\textwidth}{X} 那休她的前夫就不可在婦人玷污之後再娶她為妻，因為這是耶和華所憎惡的。不可使耶和華－你神所賜為業之地蒙受玷污。 \end{tabularx} \\ \\ \relax
24:5 & \begin{tabularx}{0.7\textwidth}{X} 「人若娶了新娘，不可從軍出征，也不可派他辦理任何事情。他可以在家清閒一年，使他所娶的妻快活。 \end{tabularx} \\ \\ \relax
24:6 & \begin{tabularx}{0.7\textwidth}{X} 「不可拿人的石磨或上面的磨石作抵押，因為這是拿人的命作抵押。 \end{tabularx} \\ \\ \relax
24:7 & \begin{tabularx}{0.7\textwidth}{X} 「若發現有人綁架以色列人中的一個弟兄，把他當奴隸對待，或把他賣了，那綁架人的就必處死。這樣，你就把惡從你中間除掉。 \end{tabularx} \\ \\ \relax
24:8 & \begin{tabularx}{0.7\textwidth}{X} 「關於痲瘋的災病，你們要謹慎，照利未家的祭司一切所指教你們的留心遵行。我怎樣吩咐他們，你們要照樣遵行。 \end{tabularx} \\ \\ \relax
24:9 & \begin{tabularx}{0.7\textwidth}{X} 要記得，在你們出埃及後的路途中，耶和華－你的神向米利暗所做的事。 \end{tabularx} \\ \\ \relax
24:10 & \begin{tabularx}{0.7\textwidth}{X} 「你借給鄰舍，無論是甚麼，不可進他家拿抵押品。 \end{tabularx} \\ \\ \relax
24:11 & \begin{tabularx}{0.7\textwidth}{X} 要站在外面，等那借貸的人把抵押品拿出來交給你。 \end{tabularx} \\ \\ \relax
24:12 & \begin{tabularx}{0.7\textwidth}{X} 他若是困苦的人，你不可用他的抵押品蓋著睡覺。 \end{tabularx} \\ \\ \relax
24:13 & \begin{tabularx}{0.7\textwidth}{X} 日落的時候，總要把抵押品還給他，讓他用那件外衣蓋著睡覺，他就為你祝福。這在耶和華－你的神面前就是你的義行了。 \end{tabularx} \\ \\ \relax
24:14 & \begin{tabularx}{0.7\textwidth}{X} 「困苦貧窮的雇工，無論是你的弟兄，或是住在你境內，在你城裡寄居的，你都不可欺負他。 \end{tabularx} \\ \\ \relax
24:15 & \begin{tabularx}{0.7\textwidth}{X} 要當日給他工錢，不可等到日落，因為他困苦，需要靠工錢過活，免得他因你的緣故求告耶和華，罪就歸於你了。 \end{tabularx} \\ \\ \relax
24:16 & \begin{tabularx}{0.7\textwidth}{X} 「不可因兒子處死父親，也不可因父親處死兒子；各人要因自己的罪被處死。 \end{tabularx} \\ \\ \relax
24:17 & \begin{tabularx}{0.7\textwidth}{X} 「不可對寄居的和孤兒屈枉正直，也不可拿寡婦的衣服作抵押。 \end{tabularx} \\ \\ \relax
24:18 & \begin{tabularx}{0.7\textwidth}{X} 要記得你曾在埃及作過奴僕，耶和華－你的神從那裡救贖了你，所以我吩咐你遵行這事。 \end{tabularx} \\ \\ \relax
24:19 & \begin{tabularx}{0.7\textwidth}{X} 「你在田間收割莊稼，若忘了一捆在田間，就不要再回去拿，要留給寄居的、孤兒和寡婦；好讓耶和華－你的神在你手裡所做的一切，賜福給你。 \end{tabularx} \\ \\ \relax
24:20 & \begin{tabularx}{0.7\textwidth}{X} 你打了橄欖樹，枝上剩下的不可再打，要留給寄居的、孤兒和寡婦。 \end{tabularx} \\ \\ \relax
24:21 & \begin{tabularx}{0.7\textwidth}{X} 你摘葡萄園的葡萄，掉落的不可拾取，要留給寄居的、孤兒和寡婦。 \end{tabularx} \\ \\ \relax
24:22 & \begin{tabularx}{0.7\textwidth}{X} 你要記得你曾在埃及地作過奴僕，所以我吩咐你遵行這事。 \end{tabularx} \\ \\
[1ex]
\hline
\hline
\end{longtable}
$^{1}$《申命記》第24章1至22節.
《舊約聖經》第243頁.
《申命記》第243頁.
「人若娶妻而後見她有什麼不合理的事,不起悅她,就可以寫休書交在她手中,打發她離開夫家,婦人離開夫家而後可以去嫁別人..
後夫若恨惡她,寫休書交在她手中,打發她離開夫家,或是娶她為妻的後夫死了,打發她去的前夫不可在婦人沾污之後再娶她為妻..
因為這是耶和華所憎惡的,不可使耶和華你臣所賜為業之地被沾污了.新娶妻之人不可從君出征,也不可託他辦理什麼公事,可以在家清閒一年,使他所娶的妻快活..
不可拿人的全盤磨石或上磨石作當頭,因為這是拿人的命作當頭.若遇見人拐帶以色列中的一個弟兄,當奴才待他,或賣了他那拐帶人的就必自死,這樣便將那惡從你們中間除掉..
在丹麻風的災病上,你們要謹慎,照祭司利未人一切所指教你們的留意遵行.我怎樣吩咐他們,你們要怎樣遵行?當記念出埃及後,在路上,耶和華你臣向米利安所行的事..
你這級鄰舍不軀是什麼,不可進他家拿他的當頭,要站在外面,等那向你借貸的人把當頭拿出來交給你.他若是窮人,你不可留他的當頭過夜.日落的時候,總要把當頭還他,使他用那件衣服蓋著睡覺.他就為你祝福..
這在耶和華你臣面前就是你的二了,困苦窮乏的故宮,無論是你的弟兄,或是在你城裡寄居的,你不可欺負他.要當日給他功架,不可等到日落.因為他窮苦,把心放在功架上,恐怕他因你求告耶和華,罪便歸你了..
不可因子殺父,也不可因父殺子,凡被殺的都為本身的罪.你不可向寄居的和孤兒屈往靜靜,也不可拿寡婦的衣裳作當頭,要紀念你在埃及作過奴僕.耶和華你的臣從哪裡將你救贖,所以我吩咐你者讓行..
你在田間收割莊稼,若忘下一菌,不可回去再取.要留給寄居的孤兒和孤兒寡婦..
這樣,耶和華你臣必在你手裡所辦的一切事上賜福於你.你打橄欖樹,枝上剩下的不可再打.要留給寄居的與孤兒寡婦.你摘葡萄園的葡萄,所剩下的不可再摘.要留給寄居的與孤兒寡婦.你也要紀念你在埃及地作過奴僕,所以我吩咐你者讓行..
願臣賜福上讀去華語有遵行的事..
請姐妹平安.我們今天開始,我們又再返回之前的系列,朝向英雄之地,幫大家重溫一下,我們隔了一個四月,有時候會忘記了,究竟整個系列我們要講些什麼呢?.
我們藉著新命記,摩西如何教導新一代的那班以色列子民,告訴他們,他們將要進入加拿美地,但不是純粹拿個地方這麼簡單,那個地方既是上帝幫他們打仗,說的英雄之地是上帝賜給這班以色列民,但最重要的就是要建立一個願景,幫助這班子民百姓去到一個地方裡面,.
他要抗衡當時在迦南也好,或者在世俗文化裡面,有一些不是按著上帝心意的那些行為,那些做法,上帝要藉著這班以色列民,成為了這個萬國裡面的一個代表,甚至藉著這班子民去到迦南地裡面,重新建立一個合乎上帝心意的國度..
所以我們一直讀的不同的誡命,或者十誡也好,或者很多不同的內容,有一些是由上而下,由一些理念開始,然後就應用在生活當中來實踐,然後在新命記的後半段的時候,就將那些生活情境就呈現出來,將這些生活情境可能加上有不同的誡命的一些想法,可能是彼此交織的..
藉著這個在處境裡面重新思考上帝的心意,所以你看到新命記,他有一些地方會講很多概念,理念層面的篇幅,那個是頭半部分,新命記裡面摩西對以色列民所講的,但是後半部分就會講一些生活例子,.
在這些生活場景裡面,怎樣去想上帝由下而上,來改變跟外邦人的生活有所不同..
今天我們看的新命記第24章,也是一個很典型的講論的系列,不過你會看到那些生活處境是很零散的,一兩句就完結了,所以今天我嘗試將它分為三個大段落,一個講離婚,一個講家庭友善,家庭友善是今天的terms,聖經舊時不是用這個字眼的,.
不過我們都用這個怎樣去對待別人的家庭,怎樣去對待社群,最後怎樣去對待社會上最窮乏的孤兒寡婦,藉著這三個段落,我們去反思一下今天我們作為一個上帝的子民,作為一個基督徒,我們怎樣經歷我們的信仰生活,我們又怎樣去實踐,我們這個時間求主幫助我們,我們先祈禱,同心祈禱..
上主願你幫助我們,你的話語帶著能力,你的話語也帶領我們行在你的道路,行在你的美善當中,我們恭敬將來的時間交給你,幫助孫講的講得清楚,讓你的話語清晰地帶給弟兄姊妹,又讓每一個聆聽的弟兄姊妹,我們每一個聽得明白,我們不單止知道我們想怎樣行,.
更加在我們生活裡面,內化你的話語在我們的心靈當中,在我們的行事為人,按著你的心意,一起來建立上帝你所喜悅的國度,我們祈求你,與我們每一個同在,恭敬將來的時間交託,願你賜福,獻上禱告奉基督耶穌,得勝的成名,Amen..
我們一起看看,我們先看第一節至第四節,這個是關於離婚的處理,人若娶妻作了,她的丈夫發現她有不合宜的事不喜歡她,而寫憂書交在她手中,打發她離開夫家,夫家,婦人若離開夫家,而後去嫁別人,後夫若恨惡她,寫憂書交在她手中,打發她離開夫家,又或者娶她為妻之後,後夫死了..
那憂她的前夫就不可再婦人玷污之後,再娶她為妻,因為這是耶和華所憎惡的,不可使耶和華你臣所賜為業之地,蒙受玷污..
在這裡有幾個字眼需要定義一下,你才明白裡面所講的內容..
首先第一個字眼就是在第一節講的不合宜的事,不合宜的事是什麼意思呢?其實原文裡面用的那個希伯來文,是赤裸的意思..
不過這個赤裸未必是裸露全身,可能露多點手臂,對著另一個男士,身體暴露出來的程度,未去到奸淫的狀況,但是已經有些不合當時文化裡面的體統..
而通常女性這樣做的一定是跟男士可能有一些隱情,在逢象裡面譯的這一句,他說到夫妻之間的妻子有隱情,這個字其實是很含糊的,究竟去到什麼程度呢?.
沒有清晰界定,不過那個程度一定未去到申命記第22章裡面講的奸淫的行為,又未去到這個地步..
但是開始夫妻之間的感情,另外有一些外遇,有一些變質,所以這個就是不合宜的事..
當有不合宜,即是夫妻之間的感情變質,男方不再喜歡她,即是破壞了婚姻的忠誠度,於是乎他就可以寫休書來交到婦人手裡面..
休書是什麼呢?休書其實是一個官方正式文件,就是說跟這個婦人已經脫離夫妻之間的關係,這個中國人也有,而在當時來說這個休書是重要的..
如果這個婦人因為有這個隱情,她有這份休書之後,她就可以嫁給之前有隱情的男士,她就可以休妻休書,然後再另外嫁娶..
如果沒有這個休書,她就犯奸淫了,她就會去到申命記第22章裡面講的是男女之間的那種不倫關係,那個是可以接受死刑的..
所以這份休書對於這個婦人來說是一個重要文件,代表著她可以再嫁,跟另外她有隱情的男士建立一個婚姻關係..
另外一方面還有一個字眼大家需要留意的,去到第四節那裡,休她的前夫就不可再婦人癲烏之後,再娶她為妻..
那個癲烏跟後面第四節裡面那個承受遺孽之地蒙受癲烏,雖然和合本是用同一個字眼,但是原文是用兩個字來的..
首先講一下第一個第四節裡面講的那個字眼,那個是宗教上面,就是說這個女人已經被休了,跟之前的前夫已經沒有關係了..

$^{41}$這個女人對於那個前夫來說已經斷絕了任何夫妻之間的關係,如果她又再嫁給前夫,這個關係就是癲烏了..
但是這個女士嫁給另一個男人就不是這個女人髒,是相對於那段婚姻而言,她的前夫而言,這個是一個癲烏的關係..
但是她嫁給另一個男人,那個男人沒有犯姦淫,而是休了之後再這樣,就不是屬於癲烏了..
所以這個是一個宗教禮儀上面的那種癲烏,是相對於那個前夫而言的..
第二個那個癲烏,那個程度就不同了,那個是犯罪來的,那個是惹惹禍禍的憤怒,所以那個某程度是要接受懲罰的..
你會看到一個純粹是一個宗教上面,那個婦人再嫁給,那個婦人對前夫來說已經離婚了,不要再拖泥帶水,不要再回去,去前夫那裡..
因為這段關係已經癲烏了,已經不可以回頭,你這樣理解就可以了,不能回頭,如果繼續再回頭,這件事就是犯罪的事..
這一段經文,第一節到第四節是在說什麼呢?.
如果你說上帝在這個角度裡面,他要講一個藉著這個離婚的制度,他在講一個什麼道理呢?.
其實到最後他要保護的,他是保護那個婦人..
在新約裡面,有些法利賽人問耶穌,摩西都容許離婚,所以沒什麼所謂的..
其實這段經文的意思不是在說鼓勵人離婚,很多人誤讀了以為是鼓勵人離婚..
其實這段經文是在說,婚姻和離婚是一個極之大的一個很嚴重的決定,不是一個輕率的決定..
也不能回頭的,如果你回頭的話,你會製造更多的悲瘡,.
以及你會製造更多的一種兒戲感,令到上帝會有憤怒在這個家庭裡面,特別是在男士那裡..
而這個優先或者製造這個制度就是去保護這個婦人,可能有一些的忍情,又沒有去到姦淫這個地步..
而當這個丈夫覺得這個婦人已經有異心了,不能夠繼續走下去的時候,.
他這個優先就成為了這個婦人的保障,讓她可以重新來過一段新的婚姻關係..
但是如果沒有這個優先或者沒有這個制度的時候,這個女人就好像人球一樣被人踢來踢去..
特別是在當時一個父系的社會,男人為主導的社會裡面,一個離婚的婦人或者一個失婚的婦人,.
她就失去了社會上的位置,她就成為了社會上最悲慘或者沒有人保護的一個最弱勢的群體..
但是如果有這個制度的時候,摩西既是強調婚姻大事,結婚已經是一個很重要的事了,離婚是比結婚更加重要..
如果你做這件事又輕率,那個婦人又再回來重婚舊時的前妻,就會搞到整件事變得非常複雜,.
也會令到那個女士就好像一件貨物一樣,當球一樣被人踢來踢去..
所以在這裡裡面,整個的段落,她不是鼓勵人離婚,她只是在說,倘若真的有這件事發生的時候,都要處理得恰當..
我想這個也是我們今天教會裡面,我們都需要正視的..
今天在坊間裡面,或者在這個社會裡面,可能都會面對著有些弟兄姊妹有離婚這個決定..
我們不要說她是不是前夫,因為新約又有另外一種說法,但是至少有一樣東西我們知道,.
我們在一個家庭裡面,我們要一個離婚的家庭,一個離異的家庭,最脆弱的那些人是哪些人呢?.
教會的恩典要覆蓋在這些人的身上,究竟誰是受害者,誰是加害者呢?.
在一個家庭裡面,很難完全說得清,但是面對著這些離婚,離異的家庭,.
我們真的要保護那些最脆弱的,不論是男還是女,我們藉著恩典要防備他們..
也希望他們的婚姻關係不要繼續延續著很多的悲滄在他們的生命當中..
所以今天我們去讀這段經文的時候,不是我們鼓勵別人去作離婚的決定,.
我們覺得結婚是很重要,但是我們就是因為要維繫著婚姻的關係,.
所以我們叫別人弟兄姊妹的離婚就不可以變得很輕率..
倘若真的有離婚的處境的時候,我們要保護最弱勢的那一方..
這個就是第一個我們需要知道的段落..
接下來由第五節,我們一直讀到第十六節,就是我第二個段落..
這個我定義為一個叫做家庭友善的一系列的條例..

$^{81}$我們先看第五節,都是繼續講婚姻的..
新娶妻之人不可從軍出征,也不可託他辦理什麼公事,.
可以在家清閒一年使他所娶的妻快活..
在這裡講第五節,在講那個男士,以前有些以色列打仗要進入迦南地,.
他們其實那些都是百姓子民,他們不是一個僱傭兵..
所以當有戰士要打仗的時候,那些男士就會徵召入伍成為軍隊..
但是有些人是例外的,就是剛剛結婚的男士就不用上戰場打仗..
可以想像的,如果剛剛一結婚就上戰場打仗,如果他戰死沙場,.
那麼那個剛剛新婚的婦人就成為了寡婦了..
所以在這裡首先他要至少有一年的時間,那個新婚的夫婦,特別是男士就不用上戰場打仗..
他可以放心地跟他太太好好共度Honeymoon,Honey year,.
就是用這個新婚的一年來好好相處,讓他的妻子得到新婚的那種快樂,.
甚至可以傳宗接代,這個是一個很人道的一種想法..
在律法裡面就規定了要保護這些新婚的家庭..
好了,接著去到第六節開始,除了新婚之外,對於那些窮乏的家庭又怎麼做呢?.
不可拿人的全盤磨石,或是上磨石,當作當頭,.
因為這是那人的命,來作當頭..
第六節那裡說什麼呢?我們今天很抽象,我們不知道..
以前那個上磨石,磨是家裡主要的糧食工具,.
你摘了那些和菇,或者那些墨水,你不是就這樣煮來吃的,.
你要磨了它成為粉,然後再開漿,然後你才可以做到那些餅來沖飢..
而沒有了這個磨,是做不到那樣東西的..
有些人家裡窮到一個地步,要典當抵押一些值錢的東西來過生活,.
當作沒什麼好當了,唯有家裡家徒四壁,.
最後剩下的就是這個石磨和那個上磨石,就是用來吃飯的工具..
這個在以色列人當中,你不可以拿別人的糧食工具來作抵押,.
來借錢給他,或者幫助他們分些糧食給他..
因為他說如果你沒有了這個石磨,就等於你拿別人的命一樣,.
拿別人的命來當頭,就是說產生不了食物,.
以致家裡整個家庭就餓死了..
在當時來說,如果面對一個家裡連這些當頭的都沒有,.
就完全借給他,送給他那些食物,就不會再要他抵押,.
又要拿他的東西來典當,然後借錢給他..
反而是見到這些最窮乏的人,無條件地送東西給他,.
以致他們可以渡過最飢餓的日子..
所以在第六節,在第十節同樣都是講當頭的,除了吃飯之外,.
第十節那裡,你借給鄰舍不虛是什麼,不可進他的家拿他的當頭,.
要站在外面,等那向你借貸的人把當頭拿出來給你..
他若是窮人,你不可留他的當頭過夜,日落的時候,.
總要把當頭還他,使他用那件衣服蓋著睡覺,.

$^{121}$他就為你祝福,這在耶和華你神面前,就是你的餘了..
去到第十節到第十三節,繼續講當頭這東西,抵押,抵押品,.
他這裡講的抵押品,你怎樣拿抵押品?.
你不可以借錢給他,然後進他的家拿一件最值錢的,.
你不可以這樣,你真的要抵押,.
你找些東西來借錢的時候拿信用,.
你要站在窮人家門口,等窮人家裡拿一件值錢的東西來做抵押,.
不是你說什麼就什麼,不是你想怎麼樣就怎麼樣,.
所以在這裡借錢給人做抵押,你都要帶著一份愛心,.
帶著一份恩典來做,第二件事就是,.
那個當頭你不可以拿過夜,為什麼拿過夜呢?.
其實舊時的人的生活是日新的,.
每天餐飲餐食餐餐清,.
基本上如果他肯做事,他有工作機會,.
即日收到一些薪金,他就即日可以還,.
所以你不可以拿人家的當頭,.
然後你要長時間保持住他,來核製債,.
以致他沒辦法還給你,你經常都要威脅他,.
你不可以這樣,總之你借他抵押,.
當下一日,一日你就還給他,你知道他有信用的,.
他又即日還錢給你,如果還不到的,你都要疏手給他,.
所以這是第二點,就是你的當頭,你抵押品不可以過夜,.
你不可以用這個抵押品來核製那個人,.
第三,除了剛才說的那個上磨石不可以做當頭抵押品之外,.
衣服都是不行的,衣服對我們來說是好,.
衣櫃很多件衣服,我認識一個姐妹,.
她說家裡很頭痛,她斷捨離,因為她有八個衣櫃衣服全部擠滿了,.
她說很頭痛,怎樣去處理這八個衣櫃的衣服,.
今天我們富足到一個地步,煩惱是為了衣服多了出來怎樣處理,.
但舊時那件衣服,舊時的人一兩件,.
當時如果在巴勒斯坦,或者在當地,你知道沙漠地區的衣服,.
白天就很熱,晚上就很冷,一收太陽,.
那件衣服白天用來遮住身體,晚上就是保暖,.
如果你拿了人家的衣服作為抵押品,.
晚上他就沒有衣服來遮蓋他的身體,就不能取暖,.
他就處於一個很痛苦,熬凍,甚至會凍死的一個狀態,.
所以這裡說抵押品不准拿人家的衣服,.
不准拿人家的上磨石,不准進人家家裡拿,.
還要不准過夜,對於借貸的人來說,.
都頗純粹要求借貸的人多過還錢的人,.

$^{161}$不過這樣的規限,其實就是他要這班借貸有能力的人,.
你要眷顧那些沒能力的人,你要體會他們的難處,.
不是純粹你有能力就去克制基層人士,那些沒條件或者窮乏的人,.
你怎樣都不能成為一個霸主來克制那些人,.
所以這個第二點是重要的,就是那個當頭,.
去到第七節,窮的人就是什麼都會做的,.
若遇見有人拐帶以色列中一個弟兄,.
當奴才代他,或是賣了他,那拐帶人的就必自死,.
這樣便將那惡從你們中間除掉,.
去到第七節是講到一個比較嚴重,.
你見到摩西講這個律法,他講到那個結果是要進行死刑的,.
因為是拐帶人口,將人成為貨物,.
和剛才那個當頭都近似的,因為好像那些物件不行,.
那個上磨石不行,那件衣服不行,.
那我就捉走那個人的女兒,或者偷走那個人的女兒,.
那個子女來我家做奴隸,用他的勞力,.
或者買賣他,用這個方式來還錢,.
或者用這個方式來令自己逃離,.
在聖經裡面講,不行,不可以,這個人的命是很重要的,.
所以不可以將這些人口作為販賣,.
也都不可以拐帶人家,如果做這些事的時候,.
那個人就一定被自死,所以摩西在這裡,.
你見到循序漸進,一路講,那些對人家的家庭,.
對於窮乏的家庭,他有一個保護網在那裡,.
令到那些最窮乏的人,他不會很擔心,.
去到第八節,這裡就有些奇怪的,.
因為他跟上面所講的那個窮乏,是有少少不咬弦的,.
不過你要看他背後的精神,第八節,.
在大麻封的災病上,你們要謹慎,.
照祭司尼美人一切所指教你們的留意進行,.
我怎樣吩咐他們,你們也要怎樣進行,.
當紀念出埃及後,在路上耶和華,.
你臣向米利暗所行的事,第八節講的是什麼呢?.
其實就講到那個大麻封,這個大麻封就不是說一定是世紀絕症的,.
患者有些皮膚病,會傳染的那些皮膚病,.
輕則可能隔離幾天就沒事了,有一些可能終身的,.
他就籠統成為一個叫做大麻封的災病,.
當有這些人患上這些病的時候,.
要照尼美人的吩咐,最重要是做些什麼呢?.
或者耶和華上帝向米利暗所行的是什麼呢?.

$^{201}$就是那個典故,就是民數記裡面講到,.
因為米利暗不喜歡,他不喜歡大麻封,.
不喜歡摩西娶了一個外邦女子作為妻子,.
於是琪琪姑姑和阿倫,結果上帝就降災病在米利暗身上,.
米利暗身上就患了大麻封,全身都白了,.
而耶和華上帝怎樣對米利暗呢?就叫他在營外隔離七天,.
其實最重要,講的是耶和華上帝在米利暗所做的就是,.
有這個災病,上帝要將他隔離七天,.
用這個例子,又或者可能是第一個有皮膚災病的例子,.
來成為了這班以色列人要想的,.
就是實踐這個家庭友善的一個條例,.
如果你用上文下理去理解,.
剛才講對窮乏人,對於災病的人,.
你一定要這樣去對待,.
特別是一個密集的社群,一個家庭裡面,.
你不可以因為他有皮膚病,你要愛惜他,你要照顧他,.
你就沒有了隔離衛生的條件,.
到最後你都要聽祭司,利未人的吩咐,.
將這個人隔離,避免將這個災病在社群裡繼續傳開,.
這個是一個愛心原則,這個隔離,.
不過你都知道這個隔離不是治病的方法,.
只不過是令到病不要繼續傳染,.
真正的醫治要去到基督耶穌那裡才做得到,.
去到神的恩典才能夠做,.
不過這裡就是要照顧家庭,不要全部野開,就是這個原因,.
第十四,第十五節,困苦窮乏的故宮,.
無論是你的弟兄或者在城裡寄居的,你不可欺負他,.
要當日給他功架,不可等到日落,因為他窮苦,把心放在功架上,.
恐怕他因你求高夜活,罪便歸你了,.
十四,十五節,這裡說你不可以欺負那些為你打工的那些人,.
那些打工通常是甚麼人呢?.
就是你收割的時候,你用日薪來聘請那些寄居的人,.
或者窮苦人來做一些收割的工作,.
特別是在雨季來到之前,如果一落雨,那些農作物就會被水淹沒,就會浪費了,.
所以很短暫就要快速地一有收成,就要拉大隊,找工人來做收割,.
在這個收割工人裡面是計日薪的,你不可以欺負他,.
他所說的欺負他,欺負他是甚麼呢?原來他說到就是出糧的時候,.
你怎樣出糧給他,就界定了你是不是欺負他,.
你不要等到真的做完日落的時候,到最後一刻,摘了,你才出糧給他,.
在這裡裡面,如果他摘了最後一粒,他在做的時候,就在詛咒僱主,.

$^{241}$你最壞,你撕到死,他就這樣心去擺一個工架,老闆不出糧給他,.
他就將這些求高惹和,罪就歸到僱主那裡,這是一個勸勉式,.
簡單來說就是,出糧不用等到最後一刻才出,這個對我們今天都,.
有些弟兄姊妹,我想是值得去留意,在座裡面可能有些弟兄姊妹,.
在家裡有請其他外傭姐姐,我認識有一個,每個月有31號,.
是11點59分才出糧,合乎法例,我不知道你有沒有聽過,.
不要虧蝕給他,你猜姐姐心裡面怎樣想,姐姐一直做的時候,.
到11點58分還未出糧給他,心裡就很亂,預早出糧給他,.
其實我有機會去聽,很多姐姐儲錢都是寄回鄉下,.
有時候可能他們都會背著一些借貸,他來香港工作的時候,.
又要給按金又要給訓練費,其實又要借了一些錢,.
以至到他才可以來到香港打工,其實每個月除了自己使用那些之外,.
那些就是要還債,然後才寄給家裡,你會見到其實在香港的生活一點都不容易,.
但如果作為僱主的我們,很斤斤計較,我們只是做法例上的標準,.
我們將條線拉到最底,不是不行的,不過在經文裡面說到,.
如果那個姐姐在祈禱裡面去求高夜和,罪就歸到你身上,.
這個好像很負面,但他亦都提醒我們作為僱主,.
或者我們有經濟能力的人,我們會不會願意鬆綁,.
或者我們有多些恩慈對待在我們身邊的這些為我們打工的工作的這些人,.
最後就是十六節,不可因子殺父,也不可因父殺子,.
凡被殺的都為本身的罪,這個就是說父親的罪不會由兒子來承擔,.
父親的債亦都不可以由兒子來承擔,相反兒子有罪就不可以追討他父親,.
誰犯罪就追討誰,這個就是一個很均真的做法,.
由剛才所讀的,一路由五節到第十六節,這十幾節的經文裡面,.
其實說一個家庭不是對你的家庭友善,而是對你身邊的窮乏人,他們都有家庭,.
他們都有面對不同的艱難,甚至乎他們的條件沒有我們今天富裕的人那麼好,.
我們要多走一步去想他們,我們要想他們家裡要供養不同的人,.
又或者他們足襟見肘,我們作為基督徒,或者上帝指紋我們的思考方式,.
就不是每一樣東西都是做到只是律法准許,而是我們要有多一點恩典的空間,.
來跟這些社會上基層困難的人相處,我們看最後一個段落,.
十七節到第二十二節,.
你不可向寄居的和孤兒屈往正直,也不可拿寡婦的衣裳作當頭,.
要紀念你在埃及作過奴僕,耶和華你的神從哪裡將你救贖,所以我吩咐你這樣行,.
你在田間收割莊稼,若忘了一菌,不可回去再取,要留給寄居的與孤兒寡婦,.
這樣耶和華你的神必在你手裡所辦的一切事上賜福願你,.
你打橄欖樹,枝上盛下的不可再打,要留給寄居的與孤兒寡婦,.
你摘葡萄園的葡萄,所盛下的不可再摘,要留給寄居的與孤兒寡婦,.
你也要紀念你在埃及地作過奴僕,所以我吩咐你這樣行..
在這裡某程度上是一個結論的經文,在這個結論裡面多次提及,.
你要向寄居的孤兒寡婦成為了這個段落的一個重要對象,.

$^{281}$而如何去對待寄居的與孤兒寡婦呢?.
你要想起你在埃及的時候作過奴僕,具體怎麼做呢?.
就是有田有地的那些以色列人,你收割的時候,收割盡量,.
你當然是豐收,盡量收多一點的果子和農作物,.
不過你收割的時候你忘記了有一群,你不要拿了,.
不記得拿回也不要帶入你的倉庫裡面,你就留在田裡面,.
留在田裡面就是給那些孤兒寡婦可以自己摘來吃,.
同樣的就是那些葡萄,橄欖,剩下的你未必100\%摘完,.
你怎樣都會剩下一些,剩下的那些你就是留給那些孤兒寡婦自己摘來吃,.
在社會裡面就是另外一個的說法,田不可割盡四角,.
就是說你田收割的那些農作物,你不要割盡了,.
你拿了所有的農作物,你怎樣都要留一些在這裡,.
留一些的目的就是要給社會上有些人沒有田沒有地的,.
你給他們,他們可以摘來吃,在他們生活裡面困乏的,.
社會上都有一些田地可以供他們來養活他們,.
就是一個這麼簡單的道理,.
但是當你做這麼簡單的道理,上帝說他會賜福給願意做這些事的人,.
如果你不願意去做的,你回想起你以前在埃及地你是怎樣的人,.
上帝沒有說會降和的,不過你想一下,你問心,.
你以前做過埃及的奴隸,你怎樣過貧窮的日子,.
你要怎樣去待身邊的人,今天這段經文都比較長的,22節,.
但是你看到有一個中心點,我很喜歡有一個外國的舊約學者叫Christopher Wright,.
他定義這一章24章,他講到整個律法其實就是要建立一個關愛的社會,.
做每一樣事不要去到盡,反而是在律法裡面,.
他想藉著這些律法去教導那些以色列人,不要做每一樣事去到盡,.
我去看這段經文的時候,我心裡面有很大感受,.
我家裡曾經富裕過,也被人騙過錢,所以當我們窮乏的時候,.
那一點點的恩惠對我們銘記於心,那一點點的恩惠令我永不忘記,.
我講一個例子給大家聽,在我爸爸去世的時候,家裡不懂得搞白事,.
很徬徨,媽媽和我,在那段時間裡面,我們到處找人,.
因為我們家裡那時候還未信主,就找那些藍和老,十幾二十萬搞一次,.
我們已經孤兒寡婦,怎會拿到錢來做這些白事呢?.
我媽媽就說不要找這些,不知道是不是不停地花錢,.
第一件事就是找我姑婆以前回過的教會,阿加魯街神教會,找那個幹事,.
然後我媽媽說,姑婆已經死了很久了,無端端找教會問人家怎樣做白事,會不會很唐突?.
我媽媽說沒辦法,沒錢,都要問,我們真的親身去到教會,去問那個行政幹事,我們應該怎樣做?.
他們真的幫我們,去介紹我們,教我們怎樣做,然後我們告訴他們,我們沒有錢,.
然後那個幹事說不用怕,很多其實不用用很多錢,然後我媽媽再說,.
我們不是用基督教儀式,不用怕,我們介紹的公司,什麼宗教都做的,你找他吧,.
很快很便宜就處理了爸爸的後事,這件事我那時候十六歲,.

$^{321}$一直走過,每一次走過神教會,想起這段經歷,我們有很多不容易的日子,.
但是上帝有憐憫,在一些窮乏的人當中,你願意有一些恩典在他們生命裡面,對他們影響很大,.
我講這些說話,聲淚俱下,不要介意,令我哭了,我們今天所做的對於窮乏人是很有意思的,.
我們今天其實未至於,我們未至於足襟見爪,特別是這個社會現在,.
經常說經濟不好,我們勒緊了錢包,我們只想自己,但是上帝的心腸不是想我們這樣,.
我鼓勵大家,在我們不說很豐足,但是我們有餘的日子裡面,我們都能夠願意,願意擺上,願意供給,.
我們不以那些窮乏的人,不覺得他們不在社會存在,我們考慮每一樣東西,都要考慮他們的需要,.
這個就是上帝國道的心腸,律法到最後他要教這些以色列人是什麼呢?.
他不是要死守那些律法,因為上帝降禍,令他們不要做一些惡患上帝的事,.
這個律法是教導那些以色列人有恩典,有憐憫,有愛心地建立這個社群的關係,.
今天我們讀申命記第24章,我希望大家,我們有多些憐憫的心,.
我們的心變得柔軟,體會社會上困難的人,我們唯有這樣才能成為上帝喜悅的子民,.
願主席著申命記第24章,祝福我們,幫助我們..
(今天視頻就拍到這裡啦,下次見).
\newpage



\section{約翰三書使徒行傳 16:31-34}
\label{sec:eFUsT6910pQ}
\textbf{2025年5月11日母親節主日崇拜直播｜吳真真牧師｜喜樂家庭｜使徒行傳16\_31-34;約翰三書1\_4}
\newline
\newline
連結: \href{https://youtube.com/watch?v=eFUsT6910pQ}{\texttt{https://youtube.com/watch?v=eFUsT6910pQ}} ~~~~ 語音日期: 2025-05-11
\newline
\newline
\hyperref[sec:Wk72Zd8NFhw]{\small{< < < PREV SERMON < < <}}
~
\hyperref[sec:index_chronic]{\small{[返順時目]}}
~
\hyperref[sec:6xido1UHEiY]{\small{> > > NEXT SERMON > > >}}
\newline
\newline
$^{1}$今日經訓是在兩段經文.
第一段是《仕途行傳》16章31至34節.
以及《約翰三書》第四節.
我首先讀出是《仕途行傳》16章31至34節.
在新約聖經186頁.
請聽我讀出第31節.
他們說:當送主耶穌.
你和你一家都必得救.
他們就把主的道.
講給他和他全家的人聽.
當夜就在那時候.
禁足把他們帶去.
洗他們的傷.
他和屬乎他的人立時都受了浸.
於是禁足令他們上自己家裡去.
給他們擺上飯.
他和全家因為信了.
神都肯喜樂.
我們接著在《約翰三書》第四節.
在新約聖經345頁.
第四節.
我聽見我的兒女們按真理而行.
我的喜樂就沒有比這個大的.
聖經獨筆.
願神的話語祝福大家.
(自我介紹).
各位兄弟姐妹各位朋友平安.
我知道今天有些新朋友來到我們當中.
我知道他們手上會有一個.
我們送給他們的小飾物.
如果有新朋友請你揮揮手好不好?.
Hello.
那邊有幾位.
是的.
這裡也有.
我們歡迎你們.
歡迎你們.
是啊.
很開心見到大家.
今天是母親節.

$^{41}$今天是母親節福音主人.
我想知道在座有哪位是媽媽的?.
(笑聲).
怎麼了?.
是.
好.
我在這裡衷心恭祝大家母親節快樂.
加上母親節喜樂.
因為今天的主題是喜樂家庭.
在節日裡.
我們很多時候都會祝人快樂.
母親節快樂,新年快樂.
我們很少祝人喜樂.
是不是?.
好像不知道是不是因為我們中國人,廣東人.
祝人喜樂好像和喜樂樂.
就是上上落落的諧音是一樣的.
好像你祝人母親節快樂.
好像說到你人生古事大喜大樂.
這樣不是很好.
我們很少這樣祝人喜樂.
但是我今天很想告訴你.
喜樂的意義在哪裡?.
我很想祝願在我們當中的媽媽.
她們的生命得到真真正正那種內心出來的喜樂.
不過我想說一件事.
其實人生有喜樂是不是常態?.
應該是吧?.
我覺得作為一個媽媽.
她一定喜樂.
為什麼呢?.
就是喜喜樂樂,上上落落的那個.
因為養兒一百歲.
長幼多少?.
九十九.
無可避免地我們都會被我們的子女.
我也是媽媽.
我們都會被我們的子女.
他們的狀況牽涉到我們.
以致我們的情緒都會上上落落.

$^{81}$子女好,子女順利.
我們又開心.
子女一般般.
經歷挫折,困難.
我們又會很難,很擔心.
所以我覺得媽媽.
或者人生有喜樂是很正常.
第一就是媽媽為了投家.
心情一定常常都喜喜樂樂.
第二就是家庭生命的週期.
本身都會有喜樂.
大家不知道看過這個表沒有?.
這個叫婚姻滿意度的表.
什麼叫婚姻滿意度呢?.
我一說到喜樂就想起這個表.
原來隨著人生.
你們家庭的週期不同的時候.
原來夫婦之間對於婚姻滿意的.
滿足的狀態是會有起跌的.
不知道大家看不看到.
其實最低谷是什麼時候?.
就是小朋友在六歲至十三歲.
大概是青春期.
其實是青春期,小學.
為什麼那時候特別辛苦呢?.
第一就是我們很多壓力.
工作通常在那段時間.
是比較事業的一個比較重要的時刻.
另外小朋友在學業.
還要步入青春期.
又要面對我們自己更年期.
又要面對上面有人要照顧.
所以我們有很多很多的壓力.
所以生命的家庭週期.
都會影響到我們生命的起起跌跌.
然後我們都會想到.
家庭是否可以像社會學或研究所說.
成為一個發揮到正面作用.
可以彼此去支持扶持.
情緒上可以去支援的一個地方呢?.

$^{121}$但現在這個世界告訴我們一件事.
就是現在的家庭挑戰越來越多.
家庭功能失效也越來越嚴重.
我們看一些數據.
譬如去年.
政府統計處的資料告訴我們.
有4萬4千6百對人註冊結婚.
但同時那一年就有1萬8千9百多宗的離婚個案.
即是說其實2.4對結婚的人裡面.
有一對是會離婚的.
並且其實最近有一個.
關於女性幸福指數的調查.
訪問了1千3百多個女性.
問她們覺得單身女性幸福還是已婚女性幸福呢?.
超過57\%的女性覺得單身幸福.
所以你會想像到.
其實家庭究竟能否發揮到功效呢?.
是否能夠做到互相扶持.
能夠有一個喜樂的家庭呢?.
在座的每一位.
你都無可避免地成為家庭的一份子.
你可能是兒子,女兒,孫子,祖父母.
甚至我們透過血緣,婚姻或領養.
我們都可以成為一個家庭的一份子.
但你的家呢?.
你的家是一個可以互相扶持.
充滿喜樂的家.
還是一個感到佛力沮喪的家庭呢?.
讓我們今天從聖經裡探討一下.
究竟我們如何可以經歷一個喜樂的家庭.
讓我們一起祈禱吧.
因為你讓我們有家.
縱然家裡有很多挑戰,困難.
但剛才我們唱的那首詩歌.
我們在這個家裡.
因為有你的恩典臨在.
其實可以有很多的祝福在裡面.
求你幫助我們去認識.
究竟一個家裡有些很關鍵的東西.
我們如何去掌握.

$^{161}$我們如何去實踐.
求主你在當中引導我們.
觸碰我們的內心.
讓我們真的有勇氣去面對.
無論我們的家是否能夠扶持我們.
或有時是能夠給予我們一些負累.
或讓我們受傷的地方.
但我們很願意主你的愛和恩典臨到.
我們這樣禱告,交託,奉主明求.
阿們.
今天我們會從一個禁卒.
什麼叫禁卒呢?.
即是監獄裡的守衛.
家庭的改變.
來看喜樂.
如何充滿喜樂的家庭呢?.
我們先看看背景.
其實在《使徒行傳》第16至30節提到.
保羅和西拉是兩個傳道人.
他們第一次到一個叫做肥納比的地方去傳道.
他們去到的時候遇到一個侍女.
那個侍女是什麼人呢?.
那個侍女原來有鬼附在她的身上.
當有鬼附在她身上的時候.
她就可以說預言.
就像我們中國人的文米.
鬼附在她身上.
她就能說出一些很靈驗的說話.
操控她的主人.
就可以靠著這個侍女去賺錢.
這個侍女一見到保羅和西拉.
就一直跟著她.
大家可以看到圖片.
一直跟著她.
而且一直說.
這些是至高神的僕人.
要對你們傳講救人的道.
一直在後面不停地說.
但是她說的內容是對的.
但是其實她是一個騷擾.

$^{201}$她的動機是不純正.
而且會令人很疑惑.
究竟那個神是誰的神呢?.
因為侍女也會說一些說話.
也有能力.
會令人混淆.
誤會他們是一黨.
對保羅和西拉來說是一個滋擾.
於是保羅就趕走了她身上的鬼.
糟糕了.
那個侍女從此就再沒有說預言的能力.
她的老闆當然是大罵.
一隻生金蛋的雞.
不可以生金蛋.
他們就不可以透過去賺錢.
於是就將憤怒發洩在西拉和保羅身上.
抓了他們出來.
誣衊他們.
帶他們去見裁判官和群眾.
誣衊他們.
去擾亂社會的秩序.
她就煽動那些官長和群眾.
用暴力對待他們.
裁判官就脫下他們的衣服.
下令用棍打.
聖經形容不是普通打.
是打到他們重重地打.
然後就下命令將他們收監.
這個就是背景.
在這個時候第24節.
就是我們今天說的主角.
就是禁卒出場.
本來她就安安分分.
聽了上頭的話.
於是就將他們收監.
將保羅和西拉收監.
而且用一個腳鎖.
在腳上鎖了一個木狗.
什麼叫木狗呢.
就是剛才在圖上.

$^{241}$就是一個有通窿的.
一個用木製的刑具.
對於羅馬人來說.
甚至有一些羅馬的刑具.
將兩個洞鎖得很開.
分得很開.
以至他們的腳要鎖得很大.
才可以插進去.
令他們很痛很痛.
讓他們受苦.
而保羅和西拉就被這些木狗鎖住.
進去內監裡面.
監是最裡面的.
因為要給那些嚴重重犯.
不要讓他們逃脫.
所以就困在裡面.
困在裡面.
還被鎖住.
應該安然無恙.
但是沒想到.
到了半夜的時候.
保羅和西拉竟然不是在身陰.
很痛苦.
很痛.
他們在做什麼呢.
他們在禱告.
還要怎樣呢.
禱告神救我們.
他們還要唱詩讚美上帝.
在這個時候監牢的地基就搖動了.
監門全部打開.
鎖住犯人的鎖鏈都鬆開了.
這個噤出精神.
他看到所有監門打開的時候.
他第一時間就想.
糟了,走犯了.
所有犯人都走了.
他意識到自己大禍臨頭.
在這個時候他很慌張,很害怕.
他在這個情況下.

$^{281}$他決定要怎樣呢.
拔劍,要自盡.
在這個驚惶失措的噤出.
看不到保羅,看不到西拉.
以為他們走了.
不過上帝的侍者保羅和西拉.
看到他.
在這個電光火石之間.
他就在監獄裡叫出來.
不要傷害自己,我們還在.
就將他從絕望的生死邊緣裡拔出來.
這個時候金卒醒了.
他停止了這個魯莽的行動.
因為已經是半夜了.
周圍的環境和氣氛都很暗.
金卒就叫人拿起火把.
令到監獄的黑暗可以驅散.
他衝進監獄裡.
在西拉和保羅面前.
很顫抖,很害怕.
伏在他們面前.
因為他相信他們一定是神聖的侍者.
然後他將他們領出來.
開口問他.
兩位先生.
我應該做什麼才可以得救呢?.
我們今天講的經文就是由這個問題開展.
兩位先生.
我應該做什麼才可以得救呢?.
其實金卒問這個問題有什麼意義呢?.
我想有兩個原因.
第一原因就是走反.
會令到他有殺身之禍.
事實上在《仕途行傳》第十二章.
另一位神的僕人彼得.
他也曾經被監禁.
後來被天使解去他的鎖鏈.
並且令他離開監獄.
結果看守他的金卒要被處決.
原來他有機會失去性命.

$^{321}$所以在這個時候.
金卒可能關心自己現在面前.
最切身的問題和危機.
去問這兩位先生.
究竟他可以如何脫離眼前的危機呢?.
但另外一個他可能更加重要要問的問題.
是可以參考之前.
我們剛才第十七節出現過的經文.
記不記得那個被鬼附的侍女.
跟著保羅和西拉.
一直在後面他們說什麼?.
他們如何稱呼她?.
她說她是至高神的僕人.
她要去傳講救人的道.
她擾攘了保羅和西拉一段日子.
保羅和西拉才幫她趕走鬼.
這些傳聞和傳說.
加上他們進來監獄.
他們為何要被捕.
我相信金卒應該聽聞過這兩人的來歷.
這兩人的背景.
在這個時候.
他看見眼前這麼神奇的奇蹟.
當他們唱贊美詩的時候.
監牢也震動.
他們的鎖鏈自己解開.
不是拿鑰匙來解開.
他就開始問.
這兩人是否真的至高神的僕人?.
這兩人可以講救人的道理.
我不是問他嗎?.
所以他問他這條問題.
有兩個很重要的意義.
其實弟兄姊妹.
在我們人生裡.
一樣的.
一樣會有意想不到的衝擊發生.
可能是自己患重病.
或者我們家人患重病.
或者我們想要的目標,理想.

$^{361}$受到挫折,失敗.
或者我們有些很需要我們照顧的孩子.
SEN的孩子,患重病的孩子.
我們一直要去照料.
或者我們工作很不如意.
或者我們被人陷害,被人騙.
或者我們會有意外.
或者甚至我們身邊的人會突然離開.
你們有沒有試過當這些衝擊來的時候.
你都慌躁大亂?.
你有沒有試過覺得很絕望?.
好像沒有出路.
好像沒有黑暗.
好像很黑暗.
我們費盡努力去找辦法,找出路.
但是和這個禁足一樣.
我們可能去問的,去找的.
未必是我們最需要的.
未必是我們人生裡最重要的事.
可能我們會找人醫病.
可能我會找專家去幫我教育孩子.
可能我會找一份好工作去解決當前的問題.
但事實上我們最關注的事.
未必是最重要,最核心,最關鍵的事.
但當我們去找對的人,問對的問題.
我們的生命就會有機會轉化.
生命就會豁然開朗.
其實大家在這麼多年的日子裡.
不是第一次聽到耶穌.
不是第一次聽到寶貴的福音.
不過很多時候我們聽完之後就會忘記.
或者忽略了它.
但是當生命裡面遇到一些很奇妙的事.
當你遇到一些衝擊,一些生命的驟變.
我們是很值得去思考人生的問題.
就好像禁卒一樣.
而且當你問的時候.
禁卒見到監門打開.
其實福音的門也為你打開.
我們一起看看禁卒如何回應.

$^{401}$禁卒回應之前.
首先要保羅回應.
我們會看到.
當保羅一直引導禁卒.
而且不單止禁卒.
還有他一家的時候.
他們整個家庭就成為充滿喜樂的家庭.
我們首先看看禁卒問保羅如何回應.
我們看到喜樂家庭的秘訣.
就是保羅所回應的.
「當信主耶穌」.
禁卒問我要做些什麼才可以得救.
保羅和西拉的回應是.
「當信主耶穌」.
相信主耶穌不是你要做什麼.
是你要信耶穌.
這是人生得著拯救的唯一途徑.
然後他說「你和你一家都必得救」.
其實余文哲也是.
「你必得救,你的一家也是這樣」.
這是很清晰的表達.
就是說家人不是因為禁卒得救.
所以全家都得救.
不是這樣的.
這裡說的是.
當他們像禁卒一樣.
聽到這個福音的時候.
和禁卒一樣.
相信耶穌的時候.
他們都會得救.
因為凡信耶穌的人就可以得救.
於是禁卒聽了這個回應之後.
他就帶同他的家人.
一起去聽保羅和西拉告訴他們.
信主耶穌的真理教訓是什麼.
其實他說要信主耶穌.
其實得救的核心就在於耶穌基督.
祂是神.
但因為愛世人的緣故.
祂願意渡誠肉身來到這個世界.

$^{441}$為了拯救我們.
為什麼要拯救?.
其實我們每一個人都是罪人.
我們每一個都是自我中心的人.
去武將做我們的上帝.
愛我們的上帝.
是我們的天父的上帝放在眼內.
我們去做一些違反祂心意的事情.
聖經說我們都如羊走米.
各人偏行幾路.
於是人的罪就讓我們同神隔開了.
罪在這個世界裡.
是讓我們人與人之間有很大的隔膜.
家庭裡為什麼不可以和諧?.
人與人之間為什麼不可以和諧?.
因為罪在我們當中.
我們會彼此去傷害.
心靈裡也得不到真正的安寧.
就是因為這個緣故.
就是因為罪令我們同神分開了.
也令我們在人世間裡有很多煩惱.
很多憂慮.
耶穌基督於是在二千年前.
祂來到這個世界.
祂真是有血有肉地活了三十三年.
最後祂被誣衊成為處死的犯人.
本來我們要為自己的罪去承擔應得的懲罰.
但是耶穌基督卻帶著無罪之身.
信服天父的吩咐.
甘願犧牲自己.
釘在十架上.
代我們受罰.
最後,他死的時候.
在十架上說了一句話.
「成了」.
這個意思就是一分一毫的罪債.
都幫我們還清.
他的寶血洗淨了我們的罪.
而且他死後第三天復活.
戰勝了死亡.

$^{481}$只要我們信他.
就是那位愛我們.
創造天地.
上帝賜給我們這份寶貴的禮物的時候.
我們就可以跟他恢復一個寶貴的關係.
所以這條就是一條永生的道路.
得救的道路.
耶穌說:我來是要叫人得生命.
並且得的更豐盛.
其實只要我們承認自己是罪人.
我們得罪了人.
更加嚴重的是我們得罪了神.
我們願意悔改.
不再過自我中心的生活.
我們相信耶穌.
將我們生命的主權守衛.
交回給祂.
這樣我們就得救.
其實當他的家人聽到.
這個寶貴的福音.
一個以基督為核心的福音.
他們相信了耶穌.
接受了福音.
而且他帶保羅和西拉出去監獄.
外面的一個井婆.
他做什麼呢?.
其實是為他們預備受寢.
幫他們施浸.
但是通常禁足這些粗獷的人.
他們很多時候都是退伍的軍人.
都是很冷酷的.
但是當他們相信耶穌的時候.
他們不同了.
他們竟然留意到這個犯人.
但是現在要為他施浸的人.
他身上所受的傷.
他很細心地為他們清洗傷口.
而且他還帶保羅和西拉到他家.
為他們送上飯菜.
對待他們如上賓一樣地招呼.

$^{521}$聖經告訴我們.
他們全家人都因為相信耶穌.
就有大大的喜樂.
滿心的歡喜.
是的,相信耶穌.
得而重生.
自然就有喜樂.
這個喜樂不是我們所說的快樂.
大家知不知道.
英文的happy是快樂的意思.
但是happy的字的hap.
意思是機會.
是運氣.
你開心嗎?.
其實完全取決於你正在遇的際遇和環境.
例如老闆加薪水.
你開心嗎?.
老闆炒魷魚.
你開心嗎?.
你不開心吧.
所以其實happy就是.
因著周圍的情況而有的快樂反應.
但是我要祝母親節快樂.
為何要加喜樂呢?.
因為喜樂這個字.
在聖經裡.
因為信主帶來的喜樂.
耶穌說沒有人可以奪去.
是發自內心的.
不是被環境去掌控的.
因為我們喜樂可以超越環境.
因為我們心裡有盼望.
因為耶穌已經擔當我們的罪債.
耶穌戰勝了死亡.
我們和耶穌同主恢復了愛的關係.
所以我們靠著這位耶穌.
藉著信.
我們心裡就充滿了盼望.
藉著信.
我們有永生.

$^{561}$藉著信.
我們信得過這麼愛我們.
連命都給我們的耶穌.
祂的大能大力.
可以帶領我們過一個豐盛的生命.
所以這份喜樂.
是從心底裡出來的.
是不受環境影響的.
你看看保羅和西拉.
他們受盡奴役之苦.
被人打得全身受傷.
但是他們毫不畏懼.
難道有天使幫他們擋住棍子嗎?.
不是的.
你看到他們是要去洗傷口的.
甚至他回顧.
他在《哥林多前書》回顧.
他自己經歷這些事的時候.
他都說.
這些是基督僕人所經歷的苦楚.
但是他們竟然在監獄裡.
他們可以禱告.
他們可以讚美上帝.
他們甚至可以和身邊的犯人一起去分享.
在監獄裡好像開一個報道會一樣.
因為他們就是經歷神的實在.
監門為他們打開.
足以證明他們的清白.
所以這完全是超越環境.
超越苦難,縱有痛苦,仍然有盼望的喜樂.
有牧者形容.
他說超越環境的喜樂.
就好像媽生孩子一樣.
經歷陣痛.
雖然是痛.
但是兒子很快要出生了.
孕婦雖然經歷痛楚.
但是都會期望孩子出生所帶來的喜樂.
因著這份盼望而喜樂.
所以信主的人的盼望是由基督而來的.

$^{601}$但事實上世界上真的有很多挑戰.
我們的確會經歷很多情緒的起伏.
這是人之常情.
彼得前書一章五至七節說.
「你們因信夢神能力保守的人.
必能得著所預備到末世要顯現的救恩.
因此你們是大有喜樂.
但如今在百般的試煉中暫時憂愁.
叫你們的信心既被試煉.
就比那被火試驗仍然能壞的金錢更顯寶貴.
可以在基督耶穌顯現的時候得著稱讚榮耀尊貴」.
彼得前書很清楚地說.
信耶穌不是沒有煩惱的.
不是沒有挑戰.
不是沒有苦難.
但他說得很清楚.
我們確實會有百般的試煉.
我們暫時會憂愁.
但這些試煉並不是枉然白費.
因為在試煉當中.
大能的神會保守我們.
祂對我們如金紙更加珍貴.
我們的生命會加以成長.
讓我們的信心得以試煉.
更加有能力去面對艱難的日子.
而且我們已經有救恩在我們生命中.
在耶穌再回來.
我們再去見祂的時候.
我們會得著稱讚.
這不是很值得喜樂的事嗎?.
另外最後一點.
其實這個喜樂不只是信主那一刻有.
而是我們想家庭喜樂.
或是我們想生命喜樂.
我們一定要按照真理而行.
我們才會喜樂.
相信主耶穌並不是單單理性同意.
而是積極地信靠或依靠耶穌.
約翰三書四節說.
「我聽見我的兒女們按真理而行.

$^{641}$我的喜樂就沒有比這個大」.
這段經文是耶穌所愛的門徒.
在年老的時候所寫的.
他作為教會的長老.
他說他最大的喜樂是什麼?.
或者作為父母最大的喜樂是什麼?.
沒有比這件事更喜樂.
就是他們的兒女.
或者屬靈的兒女.
是按照真理而行.
信服神的帶領.
謙卑地去生活.
因為只有這些人才是真正真正.
信靠神的人.
信靠天父代表了悔改.
你的價值觀已經不像世人一樣.
以前的你可能追求金錢,愛情.
或其他東西來滿足自己.
讓自己有安全感.
又或者你追求名利,權勢.
可以換取掌控別人.
帶來的滿足.
但是信靠神的人.
自然就很想明白.
究竟神的心意是怎樣.
按照祂的心意生活.
在生活裡一定會帶來轉變.
以至生命會自然產生果效.
結出一些聖靈的果子.
所以我很想鼓勵.
如果在我們當中作為父母.
我很希望你們最大的喜樂.
不是你的兒女有什麼成就.
而是他們能夠按照神的心意去行.
在他們成長的階段.
你就成為他們的榜樣.
不要讓家裡的基督是我的主的牌.
變成一個裝飾品.
而是整個家庭堅持去實踐的信念.
讓孩子學校理.

$^{681}$因為你學校基督.
如果你家裡只有你一個信耶穌.
我想告訴你.
你更加要倚靠這位上帝.
因為家裡除了你.
沒有人能夠將耶穌介紹給他們認識.
你就是喜樂家庭的起始點.
所以靠你了.
在我的經歷裡.
我很認同一句說話.
他說:神在你心裡很大的時候.
你的問題就會很小.
但如果神在你心裡很小的時候.
你的問題就會很大.
我相信同樣可以放在我們家裡.
神在你家裡很大.
你家裡的問題就會變得很小.
神在你家裡很小.
你家裡的問題就會變得很大.
其實說起喜樂的家庭.
家庭裡喜喜樂樂.
我家裡就是一個很典型的例子.
有沒有人知道我還沒有做傳道人.
還沒有在教會服侍之前.
我是做什麼的?.
做生意.
有沒有人聽過?.
其實我是一個太子女.
但在做太子女之前.
我不是含著金鎖匙出生.
其實小時候我家裡很窮.
一家七口住在窩居裡.
我還經常聽媽媽和爸爸聊天.
偷聽.
你知道小朋友最厲害的是什麼嗎?.
不是你正正經經跟他們說話.
他們不會聽.
但他們會聽你含含沉沉的小聲說話.
所以你跟小朋友說話.
你自己私下說就可以了.

$^{721}$我偷聽到我爸爸媽媽.
他們經常在籌募金錢.
要問誰要借錢.
去哪裡擲米.
但經過一輪努力.
爸爸在我中學階段開了一間小工廠.
然後我們一家人.
我中二時已經經常回工廠工作.
主隔時工作至通宵.
因為聖誕貨期我們開紙品廠.
做卡通盒.
然後我們家裡生活得到大大改善.
不止在香港開工廠.
在大陸開了幾間工廠.
亦做了一些貿易生意.
漸漸我們由窩居搬到私人屋苑.
然後去到西貢的三層高別墅.
我們家裡亦試過有很多架車.
這是我家庭裡由低處到高處.
經濟上.
但在1998年金融風暴來到時.
我們家裡面對很大的挑戰.
因為我們被欠債超過七位數字.
以致我們無能力供樓.
家裡的西貢別墅變成了銀主盤.
要抵押給銀行還債.
其實經濟上下滑得很厲害.
但很多人都說不要緊.
窮不要緊.
大家齊心協力一起捱過便可以.
不過很可惜的是.
在這個時候我們家裡面.
爸爸因為他在大陸.
當時很多男人都在大陸包二奶.
所以我媽媽情緒非常低落.
抑鬱.
她多麼堅強都試過企圖自殺兩次.
所以家裡面真的很不濟.
我有很多方法去幫他們.
但都沒有辦法.

$^{761}$但作為神的兒女的我.
我曾經很生爸爸的氣.
我真的生他氣不想見他.
直至有一個姐妹問我一個問題.
她問我是否覺得我比爸爸更配得上帝的拯救.
你覺得你配一點.
他沒有那麼配.
當她這樣問的時候.
我心裡面很痛.
我覺得其實我都不配.
我都是有罪的人.
我可以得救只是神的恩典.
不是我做什麼.
很奇妙.
就在那個時候我定義.
按真理而行.
堅定地關愛我的爸爸.
常常為他信主祈禱.
很多年之後.
真的很多年了.
爸爸終於了倦之患.
他回到香港.
不過當時他已經患了兩個重病.
一個是抗藥性的肺癆.
一個是癌症的末期.
我記得他出入醫院的時候.
我還要托著肚子.
很不容易.
雖然如此.
但我很感恩.
因為當時他認識了耶穌.
他接受了耶穌做他的救主.
他得到最寶貴的東西.
就是信靠耶穌.
我看著他去倚靠耶穌.
經歷怎樣痛的痛楚.
他都向神祈禱.
一一捱過.
我很感恩.
但我心裡面當時也很忐忑.

$^{801}$我的家是我信了耶穌.
我爸爸信了耶穌.
但我家其他人還未信.
我心裡祈禱.
我希望我家其他人.
都可以接受這個寶貴的福氣.
特別是我媽媽.
因為爸爸傷得最深的是媽媽.
就算他信了耶穌.
他爸爸信了耶穌.
他要離開這個世界.
他們一生就能離開.
但媽媽和他之間的關係.
那種傷害仍然未能解決.
我一路祈禱.
一路祈禱求主憐憫.
有一天醫院打來說.
爸爸不行了.
於是我們全家人.
趕去見他最後一面.
那一天是母親節之後.
我們去見爸爸最後一面.
不是,是母親節之前那天.
星期六.
我們去見他最後一面.
當我們一家人去到床邊.
爸爸曾經的腦癌.
去到腦部.
他意志不清.
但那天他真的很清醒.
他用最弱的氣力.
看著我媽媽.
示意她要說話.
然後我把媽媽拉到爸爸的肩前.
讓他們牽著手.
我爸爸說了一句話.
他說對不起.
這是我沒有想過的.
因為我爸爸以前和我爺爺一樣.
全部都有兩個老婆.

$^{841}$他不覺得有什麼問題.
但他竟然說對不起.
而且我怕媽媽聽不清楚.
要他再說對不起.
媽媽一時之間.
不知道如何回應.
在這個時候.
我就問媽媽.
爸爸跟你說對不起.
你原不原諒他?.
媽媽已經忍不住哭了.
她什麼都說不出.
她只是不停點頭.
然後她說原諒.
最奇妙的是.
她跟爸爸說.
我原諒你.
我要跟你信耶穌好不好?.
爸爸點頭.
接著.
爸爸安詳地離開.
然後我弟弟也信了耶穌.
我看到一個家庭.
其實可以有很多苦難.
很多挑戰.
我爸爸選擇了信靠耶穌.
願意悔改.
就成為了他人生最大的轉捩點.
由他以前沒有錢.
破產了.
沒有家.
人家也覺得他丟下家不好.
他也沒有健康.
到他後來被媽媽原諒.
心裡有勇氣.
平安地面對死亡.
有勇氣向傷害過的人說對不起.
他的喪禮有幾百人來送他.
因為他的見證.
真的幫了很多人認識耶穌.

$^{881}$我媽媽也因為這樣而有很大的轉變.
我媽媽本來是一個很安靜的人.
她是一個很害羞的人.
很內向的人,明顯是一個哀人.
加上爸爸這樣.
她很怕見親戚.
她覺得是很醜的事.
自己被丈夫拋棄.
但因為她信了耶穌.
我覺得她過了人生最快樂的階段.
後來她在教會很積極地探訪.
很積極地服侍人.
而且她還有勇氣去分享見證.
告訴別人神在我們家奇妙的作為.
所以我看到一個本來可以很破碎的家.
因為耶穌的救恩.
因為耶穌來到我們當中.
以致我們轉化成為一個喜樂的家庭.
今天媽媽雖然患了中度的認知障礙.
在生活中帶來了很多不便.
在我們照顧上帶來了很多難處.
但我們很深信.
我們家仍然是一個喜樂的家庭.
因為那份從神而來的喜樂.
是沒有人可以拿走的.
所以我覺得要有喜樂的家庭.
必須回到我們天父,爸爸的家.
回到這個永恆愛的家.
喜樂就充滿我們的心.
我在這裡很想挑戰大家.
我向三類人去挑戰你們.
第一,如果你是家中唯一信主.
或者你是家中少數信主.
你千萬要繼續去倚靠上帝.
因為你的家人可以透過你去認識耶穌.
將你的家成為基督的家.
轉化你們的家.
第二,我挑戰那些信了主的弟兄姊妹.
我鼓勵你們.
好像約翰三書第四節所說.

$^{921}$但願我們每一個都按著真理去行.
讓基督真的成為我們家的主.
將我們的家交託給他.
最後,我挑戰一些.
你之前未決自信耶穌.
未將你的生命交給耶穌的朋友.
我很鼓勵你們.
我家就是一個很活的見證.
當有基督的救恩臨到你們家的時候.
你們的家可以轉化.
我鼓勵你們.
就像按著真理一樣.
一聽到這麼好的消息.
這麼好的福音的時候.
你打開你的心.
告訴耶穌.
我願意相信你.
讓我們一起祈禱.
七位天父.
我們為在座的每一位弟兄姊妹.
我們交託給你.
願聖靈在我們心裡說話.
我願意今天來的新朋友.
或者來了很久.
但他們未真正將自己的生命交託給主的朋友.
我求你.
聖靈引導他們.
開口向你說.
主啊.
我願意相信你.
我願意成為你的兒女.
我承認我有罪.
求你寬恕我.
我願意悔改.
一生交給你管理.
依靠你.
按真理去行.
過一個喜樂的人生.
我又願意在當中.
信靠你的.

$^{961}$我們堅持為我們家人的信主.
我們努力這樣去禱告.
將他們帶到你的庚前.
讓一家人同在主的愛裡.
經歷你的恩典.
得享這份喜樂.
又願我們做父母的.
求主你幫助我們.
在你裡面得能力.
有你的智慧.
有你的愛心.
帶領我們的孩子.
成為一個按真理而行的人.
主啊我們將每一位交託.
恭敬仰望.
奉主耶穌基督的聖名自祈求.
阿們.
\newpage



\section{使徒行傳 11:1-18}
\label{sec:6xido1UHEiY}
\textbf{2025年5月18日崇拜直播｜鄭家恩牧師｜摻埋我玩｜使徒行傳11\_1-18}
\newline
\newline
連結: \href{https://youtube.com/watch?v=6xido1UHEiY}{\texttt{https://youtube.com/watch?v=6xido1UHEiY}} ~~~~ 語音日期: 2025-05-18
\newline
\newline
\hyperref[sec:eFUsT6910pQ]{\small{< < < PREV SERMON < < <}}
~
\hyperref[sec:index_chronic]{\small{[返順時目]}}
~
\hyperref[sec:wyPYrchtH6g]{\small{> > > NEXT SERMON > > >}}
\newline
\newline
使徒行傳 11:1-18
\newline
\begin{longtable}{cl}
\hline
\hline
章節 & 經文 (和合本修訂版)\\
\hline
11:1 & \begin{tabularx}{0.7\textwidth}{X} 使徒和在猶太的眾弟兄聽到外邦人也領受了神的道。 \end{tabularx} \\ \\ \relax
11:2 & \begin{tabularx}{0.7\textwidth}{X} 等到彼得上了耶路撒冷，那些奉割禮的信徒和他爭辯， \end{tabularx} \\ \\ \relax
11:3 & \begin{tabularx}{0.7\textwidth}{X} 說：「你竟進入未受割禮之人當中，和他們一同吃飯！」 \end{tabularx} \\ \\ \relax
11:4 & \begin{tabularx}{0.7\textwidth}{X} 彼得就開始把這事逐一向他們解釋，說： \end{tabularx} \\ \\ \relax
11:5 & \begin{tabularx}{0.7\textwidth}{X} 「我在約帕城裡禱告的時候，魂遊象外，看見異象，有一塊好像大布的東西降下，四角吊著從天縋下，直來到我跟前。 \end{tabularx} \\ \\ \relax
11:6 & \begin{tabularx}{0.7\textwidth}{X} 我定睛觀看，見內中有地上四腳的牲畜、野獸、爬蟲和天上的飛鳥。 \end{tabularx} \\ \\ \relax
11:7 & \begin{tabularx}{0.7\textwidth}{X} 我還聽見有聲音對我說：『彼得，起來！宰了吃。』 \end{tabularx} \\ \\ \relax
11:8 & \begin{tabularx}{0.7\textwidth}{X} 我說：『主啊，絕對不可！凡污俗或不潔淨的東西從來沒有進過我的口。』 \end{tabularx} \\ \\ \relax
11:9 & \begin{tabularx}{0.7\textwidth}{X} 第二次，有聲音從天上回答：『神所潔淨的，你不可當作污俗的。』 \end{tabularx} \\ \\ \relax
11:10 & \begin{tabularx}{0.7\textwidth}{X} 這樣一連三次，然後一切就都收回天上去了。 \end{tabularx} \\ \\ \relax
11:11 & \begin{tabularx}{0.7\textwidth}{X} 正當那時，有三個從凱撒利亞差來見我的人，站在我們所住的屋子門前。 \end{tabularx} \\ \\ \relax
11:12 & \begin{tabularx}{0.7\textwidth}{X} 聖靈吩咐我和他們同去，不要疑惑，還有這六位弟兄也跟我一起去，我們進了那人的家。 \end{tabularx} \\ \\ \relax
11:13 & \begin{tabularx}{0.7\textwidth}{X} 那人就告訴我們，他如何看見一位天使站在他家裡，說：『你派人往約帕去，請那稱為彼得的西門來， \end{tabularx} \\ \\ \relax
11:14 & \begin{tabularx}{0.7\textwidth}{X} 他有話要告訴你，因這些話你和你的全家都可以得救。』 \end{tabularx} \\ \\ \relax
11:15 & \begin{tabularx}{0.7\textwidth}{X} 我一開始講話，聖靈就降在他們身上，正像當初降在我們身上一樣。 \end{tabularx} \\ \\ \relax
11:16 & \begin{tabularx}{0.7\textwidth}{X} 我就想起主的話如何說：『約翰用水施洗，但你們要在聖靈裡受洗。』 \end{tabularx} \\ \\ \relax
11:17 & \begin{tabularx}{0.7\textwidth}{X} 既然神給他們恩賜，像在我們信主耶穌基督的時候給了我們一樣，我是誰，能攔阻神嗎？」 \end{tabularx} \\ \\ \relax
11:18 & \begin{tabularx}{0.7\textwidth}{X} 眾人聽見這些話，就不說話了，只歸榮耀給神，說：「這樣看來，神也賜恩給外邦人，使他們悔改得生命了。」 \end{tabularx} \\ \\
[1ex]
\hline
\hline
\end{longtable}
$^{1}$以下是以聖道敬拜主的時間.
今日的經訓是記載在新約的《士徒行傳》11章1至18節.
在唐庸的聖經新約部分177頁.
「士徒和在猶太的眾弟兄聽說外邦人也領受了神的道,.
及至彼得上了耶路撒冷,那些奉國禮的門徒和他爭辯說:.
「你進入未受國禮之人的家和他們一同吃飯了?」彼得就開口把這事挨斥給他們講解說:.
「我在約伯城裡禱告的時候,.
「雲遊象外,看見異象,有一物降下,好像一塊大布,繫著四角從天墜下,.
我且聽見有聲音向我說:彼得,起來宰了吃!.
我說:主啊,這是不可的,凡俗而不潔淨的物從來沒有入過我的口!.
第二次,有聲音從天上說:神所潔淨的,你不可當作俗物!.
這樣一連三次,就都收回天上去了,正當那時,.
有三個人站在我們所住的房門前,是從凱撒利亞差來見我的,.
聖靈吩咐我和他們同去,不要疑惑,同著我去的,.
還有這六位弟兄,我們都進了那人的家,那人就告訴我們,.
他如何看見一位天使站在他屋內說:你打發人,往約帕去,請那稱呼彼得的西門來,.
他有話告訴你,可以叫你和你的全家得救..
我一開講,聖靈便降在他們身上,正當初降在我們身上一樣..
我就想起主的話說:約翰是用水施浸,但你們要受聖靈的浸,.
神既然給他們恩賜,就像我們信主耶穌基督的時候給了我們一樣..
我是誰能阻擋神?眾人聽見這話,就不言語了,只歸榮耀於神,說:.
這樣看來,神也次恩給外邦人,叫他們悔改,得生命了..
敬文讀筆,願神賜福常讀他話語,並且遵行的人..
現在是正道的時間,今天我們很高興有見道神學院鄭家恩牧師在我們當中正講神的話..
鄭牧師是見道神學院的入學主任,跟我們很多位牧師和工人都很相熟..
我們把時間交給鄭牧師..
各位等姐妹,大家早安..
今次來到大家當中,數一數應該是第二次..
第一次是什麼時候呢?其實是有一次準備接任入學主任..
當時學院舉辦了一個培練獻身會,請了梁院長在當中講道..
我既然準備接任,大概是五月左右,我來到這裡第一次摸上門..
所以這次算是熟悉一點路,上來大家當中..
剛才坐在一起聆聽師班會,我們一起獻給神,我心裡非常滾動..
為什麼呢?這首詩歌是我當年按木的時候選擇的一首詩歌獻給神,我願一生事奉神..
今天再一次在我來到大家當中宣道的時候,能夠神直在大家,直在孩子自己,.
用師班員的口再一次唱出來,對我心裡再一次宣告,我願一生獻給神..
今天跟大家講的經文,其實很合切這首詩歌..
這首經文講什麼呢?大家剛才朗讀了一次,其實是一個故事..
我想將它再口述一次..
大家聽過我們口頭上有句話叫做「重要的事講三次」,.

$^{41}$其實這首經文在聖經裡只是使圖行轉,就是講了三次..
為什麼同一件事講三次呢?我們數一數,除了剛剛大家讀的位置,.
前面就是整個故事的開頭,中間,結尾..
到這裡,是彼得被質問的時候才再重述一次故事,這是第二次..
然後講了第三次是什麼時候呢?就是這裡之後幾章,第十五章就是耶路撒冷會議..
大家又不妥,就來挑戰為什麼外邦人信主應該要國禮,應該要守法律,又開一次會..
在我這裡,彼得再講是第三次.前前後後我幫他計算..
加起來有多少節數?是72節..
那我用一些比較才知道72有多厲害..
五旬節有印象吧?聖靈降臨,是教會開展的大事..
五旬節整件事的記載,包括講道,有多少節?40節,40而已..
另外一件事是很重要的,將福音從耶路撒冷拓展開去,那次是史提反被殺,.
那些人受害受迫,史提反也被殺了,石頭也被扔死了,那些人就四散,福音就由那時候開始,從耶路撒冷傳開去..
那這件事很重要吧?那史提反也是講整篇道,那加起來有多少節?都是60節..
五旬節40節,史提反受害60節,剛剛這個前前後後重要的事,講三次,加起來有多少節?.
真的很重要,今天我就在這個很重要的故事中,選擇中間這個位置,.
講三次,我講第二次,在裡面跟大家一起思考.經文裡面一開始就說是由耶路撒冷來的弟兄姊妹,.
就不是很妥,不妥什麼事呢?就是彼得,彼得你不是吧?你是教會領袖,厲害吧?教會領袖,你走去違反律法,去外邦人家裡跟他們相交..
當時裡面用的字眼是第二節,彼得上耶路撒冷,那些奉國禮的門徒和他爭辯..
查一查什麼叫爭辯,爭辯就不是很客氣地你一句我一句.原來這個字也用在一個場景,在尤大書說,.
是天使長和魔鬼在爭辯,就用這個字.那你撞了他們兩個,天使長和魔鬼說話,不是很客氣,請坐,喝杯茶先..
應該是很激烈.你想像當時那些在耶路撒冷的信徒,不是吧?主任牧師,彼得,當時我們也說是頭,.
你這樣跟外邦人相交,進入未受國禮的人家裡,跟他們吃飯,住在一起,不行,大家就爭辯..
然後第四節起,彼得就開始將整個故事說一次.當時第五節說,他在約柏一個地方,在做什麼呢?.
這個字其實用在其他經文是說「鷗嘴」,即是驚訝.他為什麼這麼驚訝,這麼鷗嘴呢?.
他這裡不是用中文發白日夢,什麼意思?「玩偶像,我是鷗嘴,驚訝驚奇」什麼意思?.
鷗嘴是因為他看見一件事,哇,厲害了,這個異象就是一塊布,那塊大布有四個角,有條繩子扯著,這樣降下來..
那塊布裡面是什麼呢?厲害了,他一眼就知道了,從小到大,猶太人不會吃的東西,一眼就認出來了..
例如龍蝦,鮑魚,撻沙,這些全部都是舊約裡面說的無刺無鱗的海鮮,凡是海鮮要吃,就要吃合乎大多數的樣子..
有鰭,大家賣魚給鰭慘生,有鰭的,還有鱗,要刮鱗的那些.撻沙的鱗都可以,一點點都不覺得存在,但是沒有鰭,不值得..
除了龍蝦,貝殼這些更加不行,沒有鰭,沒有鱗.另外,如果是走獸,又要反戳,我們人是不懂的,如果吐就試過了,吐完就吐了,不會再來一次的,明白我的意思嗎?.
如果是生畜要吃,就要反戳,反戳後,還有手指,有兩隻,要分開的,不分開的,有個膊在那裡,又不能吃的那些..
很多這些要求,但給得至少受訓,一眼就認出,不能吃,不能吃,不能吃.我記得有一次去到內地,入鄉隨俗,有時候會驚訝,有些肉放在面前,就吃,吃著吃著,跟我說,這些是我們這裡很流行吃的兔子,算不算是野味?.
我去到當地的機場,有手信,包好可以拿走的,我就看到他們怎樣把兔子弄好,拉直,你想像一下,然後四肢就這樣屈起來,一條像大香腸一樣,給你手提回來香港..
回到主題,當時我是四體不勤,五穀不分的讀書人,不下田,不太會煮飯,不太會分這些奇怪的肉,小心不要亂吃..
但給得不是這樣,給得至少受訓,很認得,所以他看著塊布,我的食物,馬上說什麼?他說,咦?他首先聽到聲音,叫他給得起來,起來吃了它,他就開口說,煮一下不行,煮一下,我從小到大,這些俗而不潔的東西,我從來都沒有吃過,不要吧..
透過給得至親口說,他說這番你來我往,吃吧,不行,我從來都沒有吃過,吃吧,不行,我從來都沒有吃過,來回了多少次?第十節說,三次,又是重要的事,說三次..
說到三次之後,有人敲門,在他家樓下,不是家,他借住的一個地方,到處去的一個地方,在他家樓下有人敲門,原來有某個地方,凱撒尼亞這個地方,外邦人的一個軍長,叫伯夫長,意大利營,那些人是外邦人,去做外邦人的官長..
一個叫哥尼劉的人,猜派自己一個兵,和自己的親人,請一個地方的西門的彼得來,當他敲門,聖靈說不要再和他猜乒乓球了,你來我往,猜了三次..
聖靈就出聲,第十二節說,不要疑惑,同著我去,聖靈吩咐他,不要不要疑惑,同著他一起去,當時在這間屋其他基督徒,有六位一起去,去哥尼劉家..

$^{81}$我特別提第十二節的不要疑惑,這個不要疑惑是聖靈跟他這樣說,以至他就複述出來,不要疑惑的節奏,就是剛才我說天使長和魔鬼在爭辯的字眼,你不要再爭辯了,不要再說不行不行了,去吧..
彼得就複述聖靈當時的提示,不要再爭辯了,你去吧.他就起來去,第十三節他再複述,去到那個人家,就是哥尼劉,他就問哥尼劉,什麼事?聖靈都這樣跟我說話,叫我來這裡,來這裡做什麼?.
哥尼劉又說話,其實早在這一天之前,數上去四天,如果今天是第四天,就是前面第一天,那天我正在祈禱,下午三點,很虔誠,經文都這樣稱讚他,是一個虔誠人,愛邦人信主,不過沒有受割禮..
一個虔誠人,下午三點,正在祈禱的時候,有一個天使走出來,跟哥尼劉這個愛邦人說,你平日的祈禱,你平日做的周祭,見到貧苦的人,有些奉獻給他們的周祭,你平日的祈禱,你平日的周祭,神知道,已經達到神面前,現在你派人去,請一個叫西門的彼得來,他會將得救的道告訴你..
一五一十,哥尼劉就告訴彼得,彼得聽了,哦,原來是這樣,真是神厲害,將福音傳給愛邦人,好,我現在就告訴你,耶穌是誰,得救的道是怎樣怎樣..
開報道會,開報道,彼得就在那裡說,還沒說完,報道會還沒說完,就發生什麼事,有印象嗎?聖靈降臨,聖靈降臨的現象是這班哥尼劉,哥尼劉當時在等,由第一天三點收到這個異象,他就帶人去,往還回來,去到第四天,所以他就叫了一班親朋戚友,自己密友,來來來,我見到異象,即將有,不知道是誰,.
即將有不知道什麼事情發生,所以當時有班報道會的人,有哥尼劉和一家他自己的好朋友,就來了,發生的事情就是,還沒說完報道會,聖靈就降臨,你記得五重節降臨的時候,怎樣顯明出來的?.
聖靈如火,有個東西,不知道怎樣,聖靈如火一樣降在他們頭上,以及他們的口出什麼呢?謊言,謊言不是今天我聽不懂的那些,是你都聽得懂的,因為在五重節裡面說到,當時萬邦來朝,很多人在不同的地方,來朝聖,去到耶路撒冷,.
其他可能是客家話,潮州話,上海話,英文,法文,德文,總之就是在說一些人聽得懂的,只不過是我英文不叻,聽不懂,我德文不叻,聽不懂的意思,所以當時他們那一流的一班朋友,一班親人,所說的就是這些客家話,潮州話的意思,可能說來幹什麼?說什麼內容?.
經文說他們說出來是讚美神,彼得就說,你們現在的樣子,你們現在聖靈降臨的樣子,就好像我們當日彼得,雅各,馬太一,大丘,這些十一個門徒,當時經歷的五重節,都是這樣,就是說這些客家話,潮州話,這些謊言,用來讚美神,那就是像什麼呢?像一個chop印,如果當日神用這個chop印,chop,.
你們蒙聖靈降臨,今天這班外邦人,都是同樣一個chop印,沒有變過,一模一樣,不用擔心,不用爭辯,一模一樣,都是這個chop印,他們是蒙神越立,揀選的人,就是這樣,彼得不再說那些,今天有感動的就來到台前,或者在座位舉手,信耶穌,不用了,你們已經是蒙神揀選了,那柯克就為他們洗禮,.
彼得將整個故事都說了,重要的事說三次,為什麼這麼重要,路加福音,路加自己,寫這個《士道行傳》,說三次這麼重要,我就在裡面看到一個很重要的真理,今天講題是什麼呢?.
我想今天跟大家說,「參埋我玩」,先算一下自己,大家分不清人生月曆,說過這個話嗎?說過,小時候跟別人玩的時候,走過去參埋我玩,有沒有說過?.
下街,以前公屋有條走廊,在走廊奔跑,在街頭跟別人石屎做的乒乓球台,走過去參埋我玩,很多這些叫參埋我玩,後來就變了,變了一些比較個人化的玩具,打機舖那些,我人生只進過一次,報案,.
原來是這樣的,將一塊硬幣,當時一塊硬幣,放在遊戲機的螢幕旁邊,放下來叫做「跟下一步我」,變了不用說參埋你玩,自己玩,後來變了網絡上自己打機,我們是不是很多日子都沒有說過參埋我玩?.
今天跟大家說的是參埋我玩,參埋誰玩呢?是彼得,我看這個故事是一個人故意要彼得去做這件事,其實是故意到很誇張,你數算一下整件事的結果是要外邦人信主,要外邦人哥利留信主,其實誰已經是第一個去到哥利留面前,叫他去猜險人去找彼得?.
誰已經去到?天使,天使都來到了,天使都來到了,為什麼不是天使馬上去跟哥利留說,你當信主耶穌,你和你一家必得救,你記得保羅是怎麼歸主的,在前一點,這段經文前面,保羅的歸主就是耶穌馬上來到他面前,.
保羅,訴羅,訴羅,你為什麼逼迫我?耶穌直接去找保羅,保羅就直接信了,現在天使去到哥利留身上,前面,三點鐘下午祈禱,救吧,你馬上說就搞定了,為什麼這麼麻煩?.
警察去到一個叫裕阿柏的地方,叫彼得過來,他會說一些你聽不懂的東西,他會說一些你聽不懂的東西,侵買彼得,其實不需要這麼麻煩,就是這樣,叫彼得來,你記得我們被造或創世記一章,.
經文裡說過我們是按著神的形象樣式去做,然後吩咐我們做,管理海裡的魚,天空的飛鳥,替他生養眾多管理大地,這是神交給我們管理的職份,其實你想想,想想,要不要交給我們?.
整個天地是神創造的,說有就有,命立就立,有的,有的,齊齊齊,神眼中,好,好,好,但今時今日我們,他交給我們之後,我們就搞到這個世界亂糟糟,亂糟糟,今天好像內地有句說話,青山綠水就是金山銀山,.
意思就是搞旅遊業,保護好環境就搞旅遊業,就是說現在我們很容易不保護環境,很容易破壞環境,神將這個天地交給我們,我們搞到會有人際間的紛爭,國家之間的打仗,我們有經濟上的欺壓,我們有權力就會言語暴力別人,.
我們有體力就會打人,我們管理這個大地,神交給我們,我們就搞到亂糟糟,神知不知道呢?一早就知道我們這麼差勁,但為什麼要交給我們?.
大家有兒女有姨孫,你會不會放手讓他們學走路?很危險,會跌的,會撞到頭,你讓他們學嗎?我記得小時候有個畫面,很可愛,我現在想起很可愛,當時應該是初小,很小,.
我對上一個哥哥,我爸爸,有一次趁他們兩個出街買菜,我跟我弟弟,我弟弟更小,我們就很想拖地,當時沒有那些比較智能的,壓下去,然後就噴水出來,那些叫洗地桶,沒有這些東西,.
要拿兩隻手,我跟弟弟,在地拖棍裡面,結果我們兩個就闖禍了,地拖棍一直都不乾,一滴水,我們兩個就知道闖禍了,那怎麼辦?就趕快假裝睡覺,我們兩個馬上去假裝睡覺,然後聽到我爸爸,我哥哥回來的時候,他們說什麼?火災啊!.
我們兩個當然很小,不夠力去做,但我們很想去做,大個多一點,當然要開始做,要拖地,我想說的是,有時候我們自己的管理,又不夠力,又不懂得做,就以為自己做得到,我爸爸同時也沒有說收起我,不讓我去拖地,一定是教我,讓我的手多一點力,一定要做,.
交東西給別人做,是不是容易?不容易,近年流行親子農莊,有聽過嗎?給小朋友親自去餵動物,給你一條蘿蔔條,紅蘿蔔條,你自己去餵,什麼羊仔,草泥馬,我記得早前,看手機應用程式,看看人們的近況,.
看到一條片,是我牧羊長大的一個姐妹,她已經是小朋友的媽媽,我按來看看,原來是去了這些動物農莊,親子農莊,一條片不是太長,但反映她要蹲下,然後讓兒子在她的保護前,鼓勵兒子拿紅蘿蔔條穿過鐵絲網給羊仔吃,.
我看這條短片,當然,很厲害,媽媽回答,你很厲害,很勇敢,在這個往來的影片,看到小朋友開始不敢,媽媽鼓勵他,我特別感動的是,整個母親是否站著?.
想像一下,是怎樣的?蹲著,整個情蹲著,我就在這裡腦補整件事,她沒有只叫兒子餵一次,你可以想像,整個親子農莊餵小動物吃東西,媽媽每次蹲下,又要鼓勵,不用怕,可以的,其實整件事,你想想,我的羊,需不需要一條這樣的金筍仔?.
我的金筍仔可以說是羊的零食,只是塞牙縫,真的餵飽羊,不是小朋友可以做得來的,為什麼媽媽要這樣做呢?媽媽想羊飽?一定不是,媽媽想小朋友成長,想他成長,但媽媽要付出不少的事情,蹲下起身,蹲下起身,餵的是小朋友能夠有膽量,成長多一點..
跟大家說一個很觸動的見證,我過去還沒有做神學院的入學主任,我在一間教會做青少年事工的牧者,主力親自貼身牧養大學生和初職生,而他們就牧養中學生,我每一年都會搞一些短孫,短孫其實同時間除了祝福當地,也是一些信徒領袖的訓練,訓練他們..
記得有一年,我去道當地,其中一環節是去少數民族上山探訪,我去到就坐下來,我通常分好工,你去到就負責去問人,關心人,你負責大詩歌,最後你負責大祈禱,這些都不想我自己做,我想分給大學生做..
然後有一年他們回來,就跟我說,我們去到同一個地方,同一條村,再一次進去,裡面有一個媽媽抱著一個小孩,抱著一個小孩,三個月左右,跟我說,這個是去年,當時我還沒有按牧,鄭傳道當時為我祈禱,.
當時我是一個貧物,貧繁流產的婦人,當時其中一個祈禱事項就是想安胎,當時我還沒有小孩,一年後,每年暑假去,整整一年,一年後,這兩三個小孩就出生了,告訴我們,這個就是去年你們為我代禱,然後有的出生,安然來到這個世界的小孩..
我抱著這個小孩,心裡覺得,神啊,我當時祈禱了一次,你聽我,聽我祈禱,而我有份參與,重點是參與,然後我一年除了抱著這個小孩,我們就去了一個地方,又是要行山才能去到山頭的家庭,裡面的爸爸媽媽坐下來問他有沒有什麼代禱,媽媽馬上哭了,因為我們有個兒子..
他兩歲了,現在兩歲了,也不懂得走路,因為出生的時候遇到一些生產過程,有毛病,有腦癱,腦部有些事情,兩歲了也不懂得走路,這兩夫婦很為這個小孩憂心,今時今日,你知道有鼓勵大家在內地生兩胎,生三胎,當時是沒有的,沒有的,只是一胎,所以他們為了這個小孩的人生路很憂愁..
很擔心,我們為他們祈禱,現在很想小孩能夠走路,我們就祈禱,一年後,我們又再去這個位置,他們告訴我們,這個小孩,你看,能走路,當然腦部還未準備好,繼續持續,然後他們再說,你知道現在我們,他們準備偷偷懷孕多一個,希望將來自己兩老離開,.
還有個兄弟姐妹可以照顧這位哥哥,我們準備懷孕,現在有一個,在肚子裡,是女兒,掙扎很久,要不要打掉她,就放了她,不是的,剛剛變了,因為原來我們第一個小孩有問題,有些身體的問題,結果當時還未說兩胎,三胎的日子,都已經比你多懷一胎,.

$^{121}$有時候我們對祈禱的想法,就是加油,要儲印花,加油給神,神你不夠力,我幫你,這些民間宗教的想法,不是我們天父聽祈禱的情況,為什麼要我們祈禱呢,其實所謂要你祈禱,我們試一下從親近我玩的角度,我剛才跟大家說,我所經歷的這些,很開心,親近我玩,我有份見證,我有份見到,.
比起大家聽我福術,遠了一點,聽我福術,我的見證又遠了一點,跟大家說上個禮拜,上個禮拜我們學院搞一個培練獻身會,我聽台上做師班那個,一下台,我們就有陪談員,因為叫人獻身來台前,那個師班員,完了整個聚會,在樓下一樓有陪談區,那個師班員走過來跟我說,.
很感動,我就是師班員,我站在這個位置,我見到有人走過來,我很感動,然後那天我負責為這班蒙字的人做祈禱的那個是我們的老師,他上台說,哭著說,緊咽說,通常人哭我們都會感動,眼泛淚光,很感動,.
事後,剛剛這個禮拜三,我們同工會,教員會開會,請他分享一下,你都一個準備榮休的牧者,榮休,見過很多世面,我真的問,為什麼當天你哭著祈禱,他說,真的不行,真的不可以不感動,他說,人們如潮湧走上來,已經感動這些師班員見到一個個,神為他的名,樓下七千個未向巴黎屈臣,.
不止,這個準備站上台,其實已經坐在前面,事奉人員席,我見到有一個人走上來,是穿著校服,穿著校服,現在DSE都完了,看來他不是中六,又猜,可能中四五補習完了,然後就來這個培靈慶生會,又走出來,.
真的不行,見到一個穿著校服的走出來,就很感動,當日來台前的有六十位,因為他們要填紙,我負責收回應紙,我只有六十個,我做了什麼呢?我不是台上講道的那個,但我為這些畫面很觸動,甚至我見證到我有份參與,這個師班員,你說他有沒有做,他做了些什麼,唱詩歌,.
這個做最後結束祈禱,也不是講員,他只是結束祈禱,為這些人祝福祈禱,他也很觸動,就是侵買我,弟兄姊妹,我們是不是沒有了,我們不行,天父是不是沒有了我們參與,我們不行,一定不是,反而我們謙卑些,有我們,先論論盡盡,甩甩漏漏,你教你的小朋友拖地,教他們洗碗,有沒有試過打爛碗,.
你教他們切菜,切肉,有沒有試過弄傷手,自己也有,但你也要讓他們試,因為對他們來說是成長,天父絕對不需要我們草引花,不需要我們一起搞才算是事情搞得成,天父說有就有,命立就立,所以只有一個緣由,侵買我們玩,.
在這個情況,在經文裡,神侵買彼得玩,你看看彼得為了要彼得在這件事上有份,神做了什麼,我們一起來數,首先他要去到彼得下午準備吃飯,正午的時候,他在約柏這個地方,在天台上祈禱的時候,他又看到一個異象,特意給異象看,.
然後還要往還了多少次,他在爭辯三次,不行不行不行,三次,聖靈要催促他,不要爭辯,不要再爭辯了,下樓跟他們同往,去吧,他們叫你做什麼,會告訴你,你去,已經第二劑了,聖靈出手,.
到第三劑了,彼得去到現場,明明聽到高麗樓說見證,他說四天前,聖靈,天使親自來跟我說,去到這個位置,你還不知道聖靈已經做了什麼,他都繼續開報道會,開了報道會,聖靈都等不住了,不用你講了,我自己來,終於自己來了,給聖靈充滿你們,講方言,.
神為了讓彼得在這個當中有份,他前前後後,加上高麗樓見異象,其實做了四件事,我還是那句,最簡單,其實神親自叫高麗樓,你信耶穌就夠了,他要彼得插下來,一起,他很想彼得能夠直彼得的口,他自己見證,親自講出來,.
神越立外邦人,神的福音不止於自己本地,本家,其實這個早在《小學行傳》第一章第八節已經講了,福音要從耶路撒冷去到撒瑪利亞,再傳到他們,根本早就講了給你們聽,不過他們繼續在耶路撒冷,.
所以從一個這樣的角度看,史提反被逼迫,以至教會四散,其實神都使用,使用來讓福音傳開,不過他們還是迷迷糊糊的,還是自己人的,.
讓彼得你來,親自去見證,加上不止他,因為他回來,彼得這樣做,其他那些人不滿意,就是那些人一直守住,信耶穌就需要由國禮,由摩西律法,所以彼得這次自己一個,還要連同六位弟兄一起做見證人,.
你有沒有印象,要作見證,要有兩三個人的口才能定準,給你一個雙倍,六個,加上你彼得七個,一起去做見證,證明神越立外邦人,搞這麼多事情,天父就是想侵埋我們玩,侵埋彼得玩,.
我來之前在網上看了大家的程序表,很好,5月31日有一個短訊去到越南,然後讓大家都可以侵埋你們玩,我們中國人說什麼未吃五月種寒衣未入籠,現在未到中秋節,但應該都可以執拾一下寒衣,原來大家都可以參與寒衣收集的行列當中,.
我們都知道大家很關心社區的教會,所以我看到大家的相片,不是在影樓拍的,是在旺角街頭群眾拍的,其實是一個很能夠彰顯和表現旺鎮是一間社區當中的教會,.
讓福音傳開去,是不是沒有我們不行呢?你想像一下,在約伯的影片中提及過,神親要和約伯聊天,在山上,你沒見過,沒人見過的小鹿,媽媽,生小孩的時候,誰會看顧他們的?.
大到你以為我的河馬,鱷魚,你以為你可以捉住牠們嗎?你做不到的,耶穌在天父當時就說了很多東西給約伯聽,其中一個說得我真的很動心,就是一些孤苦的,或者困難,危難生小孩,會有難處的,.
他說你能顧到嗎?你顧不到的,我知道,弟兄姊妹,很多患難,包括在香港社區,源自於越南這個地方,你想像他們是不是沒有我們一分錢,天父會不看顧他們?不會的,不會的,大神讓我們有份,.
在這教會讓大家知道的一些情況,可能是世界代討的事,可能是具體的物質上,金錢上的一些奉獻,以至教會可能邀請你參與一些事奉,一次過的,去街頭報道,或者是比較長期的委身的某些崗位,想像記住今天的講題是什麼,我們一起讀,.
「侵買我橫」,願大家記住我們是和天父一起參與整個國度的擴展,福音值我們的雙手,值我們的口,甚至值我們的笑容,去傳揚開去,都是因為神很想我們在同宮當中有經歷,有成長,我們一起去祈禱,.
因主耶穌我們感恩你,感恩你能夠在萬人當中,揀選我們相信你,亦將全福音的責任交給我們,其實是真的,你自己一個已經可以讓事情成就,你邀請我們參與,讓我們見證你的作為,見證福音廣傳而我們心得喜樂,見證到我們做的微小的丁點,而可以見到我們有份,主多謝你,.
願主你激勵我們一班頂尖在旺站能夠彼此見證,我們有這份力,我們有這份錢,我們就參與在當中,願我們在這個國度有份,願我們在他日天家的時候,蒙你稱讚你,我們是良善有中心的僕人,我現在禱告奉主耶穌得安眠,阿們..
\newpage



\section{申命記 26:16-19}
\label{sec:wyPYrchtH6g}
\textbf{2025年5月25日崇拜直播｜鄭國邦傳道｜朝向應許之地－自由與順服｜申命記26\_16-19}
\newline
\newline
連結: \href{https://youtube.com/watch?v=wyPYrchtH6g}{\texttt{https://youtube.com/watch?v=wyPYrchtH6g}} ~~~~ 語音日期: 2025-05-25
\newline
\newline
\hyperref[sec:6xido1UHEiY]{\small{< < < PREV SERMON < < <}}
~
\hyperref[sec:index_chronic]{\small{[返順時目]}}
~
\hyperref[sec:dxgI8DcWOBk]{\small{> > > NEXT SERMON > > >}}
\newline
\newline
申命記 26:16-19
\newline
\begin{longtable}{cl}
\hline
\hline
章節 & 經文 (和合本修訂版)\\
\hline
26:16 & \begin{tabularx}{0.7\textwidth}{X} 「耶和華－你的神今日吩咐你遵行這些律例典章，所以你要盡心盡性謹守遵行。 \end{tabularx} \\ \\ \relax
26:17 & \begin{tabularx}{0.7\textwidth}{X} 你今日宣認耶和華為你的神，承諾遵行他的道，謹守他的律例、誡命、典章，聽從他的話。 \end{tabularx} \\ \\ \relax
26:18 & \begin{tabularx}{0.7\textwidth}{X} 耶和華今日照他所應許你的，也認你為他寶貴的子民，叫你謹守他的一切誡命， \end{tabularx} \\ \\ \relax
26:19 & \begin{tabularx}{0.7\textwidth}{X} 要使你得稱讚、美名、尊榮，超乎他所造的萬國之上，並且照他所應許的，使你歸耶和華－你的神為神聖的子民。」 \end{tabularx} \\ \\
[1ex]
\hline
\hline
\end{longtable}
$^{1}$多謝守宗代表我們會眾作出剛才的獻祝.
接下來願意我們繼續以聆聽聖道.
繼續我們的敬拜.
今天我們會選讀一段經文.
是舊約聖經新命記.
舊約聖經新命記的二十六章十六到十九節.
請聽我讀出.
耶和華,你的神今日吩咐你行這些律例典章.
所以你要盡心盡性謹守遵行.
你今日應耶和華為你的神.
應許遵行他的道.
謹守他的律例,誡命,典章.
聽從他的話.
耶和華今日照他所應許你的.
也應你為他的子民.
使你謹守他的一切誡命.
有使你得稱讚美名尊榮.
超乎他所做的萬民至上.
並照他所應許的.
使你歸耶和華你神為聖傑的民.
各位兄弟姐妹平安.
能夠再次分享上帝的說話.
並且跟他遵行和實踐主義的教導.
今天已經是五月二十五日.
大家在二零二五年已經習慣了.
寫日曆的還不會寫二四年.
已經過了這麼多個月.
已經一段時間了.
大家除了日復日的工作和生活之外.
有沒有想過二五年.
有沒有什麼很特別的事情要感恩呢?.
畢竟都差不多過了一半了.
這一年裡面你經歷了什麼呢?.
當然了,當你看回教會的時候.
我們今年的年題就是聖經導引.
這個過程其實就像一個旅程.
跟大家每天的經歷一樣.
你們可能每一年有些目標.
去做到一些方向.
教會也一樣.

$^{41}$在二五年的時候也訂立了這個主題.
我們現在差不多走到一半了.
大家有沒有機會想想.
教會這半年又發生了什麼事呢?.
就好像上個星期.
大家如果有參與的話.
就有去到讀經日營.
當中一起早上有安靜靈修的時間.
下午也有高sir的分享.
今年如果你參加小組查經的時候.
可能也會發覺進度有一點不同了.
如果有弟兄姊妹.
有小組跟隨教會的方向的時候.
除了可能在知識上增多了之外.
我們更加著重的是.
在聖經裡面如何引導我們的生命.
有些不同呢?.
有些分別和應用是怎麼樣呢?.
每天的讀經計劃.
有些弟兄姊妹也一同參加了.
如果一起讀的時候.
可能你也會留意教會的Facebook專頁.
裡面每天也有一些短短的文章.
講一下今天讀的那段經文裡面的一些體會.
我們牧者也會一同去有一些文章分享.
甚至乎年頭.
其實如果大家有留意我們誠心所願.
有幾集也有訪問一些.
可能是夏達華的聖經世界.
以及一些不同的聖經機構.
裡面也是按照這個方向.
想讓大家有更多的資訊和補充.
甚或乎來到講壇裡面.
我們講道系列.
今天大家也知道主題.
就是《朝聖應許之地》這個主題.
這個是我們自家牧者的講道系列.
裡面也是想講這個方向.
就正如在《申命記》裡面.
我們所運用的經文.

$^{81}$也是在講摩西再一次重新的律法和解明一樣.
很希望藉著這一切的重新.
讓那一群要進入應許之地的新一輩以色列人.
知道他們今天憑著什麼.
來到今天的這個光景.
也是憑著什麼應許.
可以進到面前這個應許地.
另外也重新再宣告.
上帝的律法帶著上帝的約的關係.
藉著很多的律法的解明重點.
再去講述的時候.
其實是在塑造這群以色列人的身份.
以致他們知道今天在做什麼.
所以當今天.
我們以教會年題.
就是《聖經導引》為題的時候.
其實我們牧者們.
大英姐妹們也用了不少時間.
想讓大家得著這件事.
對,《聖經》對我來說是些什麼呢?.
教導我們的是什麼呢?.
而且用了很多時間去研讀.
律法背後究竟它的意義是什麼呢?.
我們不會說《聖經導引》.
總之今天《聖經》說什麼.
你自己跟著做就行了.
而是大家會用很多的心機和時間去研讀.
這些律法背後其實想告訴我們什麼道理.
以致今天我們可以應用呢?.
但是當你真的按照教會年題去想的時候.
不知道你會不會開始想.
問題就來了.
究竟今天對我們來說呢?.
摩西那時候就有這些律法.
他們五經裡面的教導.
很多的律法去教導他們.
今天對我們來說的律例典章是什麼呢?.
是不是就是那些我們經常說.
你做基督徒.
應該做這些不應該做這些.

$^{121}$全部的東西是不是只是說那些呢?.
另一方面.
就說朝著應許之地.
大家好像都很乖.
又沒有問我們.
你說朝著應許之地.
什麼應許之地呢?.
我現在就坐在這裡.
我又要有房子就住在房子裡.
我又不是要移民.
朝著一個什麼應許之地呢?.
我們在哪裡出來.
要去哪裡呢?.
這個方向其實是想教導什麼呢?.
可能這個也是大家可能會想的東西.
而我自己一直去研讀生命的時候.
都會想同樣的問題.
如果不是呢.
只是讀了別人的歷史.
啊,是喔.
以色列人怎麼樣信服神啊.
很乖啊.
帶著他.
雖然都不是很乖.
就是看著他們怎麼樣.
好像跟著一步一步去到應許之地.
但是我們今天又怎麼去想這件事呢?.
信服是在說什麼呢?.
信服只是在說.
啊,我聽到這些東西跟著去做就可以了.
還是可能是指究竟這些東西是教導.
怎麼在你和我每天的生命裡面.
每一個選擇裡面.
都在產生變化.
當然了.
這個變化可能你也會說.
我聽了這些教導.
我應不應用在我生活裡面是我的自由.
你沒人逼我.
你沒人說因為我是基督徒.

$^{161}$道德綁架我.
我一定要這樣做.
這個跟不跟是我的自由.
所以就更加讓我想到今天主題.
就是自由和信服究竟是什麼呢?.
究竟在神的應許.
神的律法面前.
我們是不是自由的?.
是不是可以選擇跟和選擇不跟呢?.
在神的律例典章裡面.
我們自然就信服.
那就好像沒有自由.
我們有沒有選擇呢?.
盼望這個成為我們思考的方向.
我們先一起祈禱.
親愛的主,求你給我們智慧.
以致當我們每星期每一天.
去研讀你的話語的時候.
讓你的話語成為我們腳前的燈.
路上的光.
讓我們在每天裡面.
知道怎樣前行.
以致我們能夠知道.
我們縱然可能一邊聽上帝的教導的時候.
我們心底裡也有些想法.
但我們知道這個過程.
就是我們去成長.
去跟從主理的每一個機遇.
我們祈禱是奉主耶穌基督的名字而求.
阿們.
剛才讀經的時候.
大家讀了26章的16至19節.
其實16章一開始讀的時候.
他說今天上帝吩咐你所進行的這些律例典章.
你就要盡你全部的心思.
全部的性命.
如果你看回和學本.
就是盡心盡性去謹守進行.
16節這裡就要提醒百姓.
去行這些律例典章的事.

$^{201}$其實這是一個山括弧.
那開括弧在哪裡呢?.
開括弧其實在12章1節.
我讀出來.
你有機會可以揭開聖經.
看回申命記12章1節.
裡面說什麼呢?.
你們存活於世紀日子.
在耶和華.
你們列祖的神.
所賜你們遺孽的地上.
要謹守進行的律例典章.
乃是這樣.
這樣就一直說.
從12章開始說.
說到26章.
到最後.
所以16節就說.
這些就是上帝吩咐你進行的.
誡命和律例.
你就要盡心盡性去進行.
所以原來12章1節.
到26章.
裡面都包含了.
所有的律例.
當中從整個角度.
互相呼應.
讓我們知道這是告一段落.
而中間就重申了很多不同的律法.
如果你記得過往這幾個月.
牧者在講壇上都說了很多.
這些律例就是這樣去跟的.
當時來說為什麼要跟呢?.
就是因為有這樣的原則.
那應用在今天.
又是代表著什麼的意義呢?.
過往這段時間.
就說了很多這些內容.
隨後呢.
繼續你看下去.

$^{241}$第17 18節.
就講述了.
這些律例典章背後.
表達了為什麼他們要和應.
這些上帝的誡命呢?.
就是17 18節表達了.
以色列這個民族.
和神彼此毀身的承諾.
這個承諾就是一個文約.
至於在當中最重要的.
就是在說耶和華神.
就是以色列這個國家的宗主國.
然後以色列就是耶和華的民.
就是那個從屬國.
為什麼用宗主國和從屬國呢?.
大家記不記得.
陳木斯一開始的時候.
講到一件事就是.
我們這個約的關係呢.
不是在說今天我們簽的合約.
不是那些合約精神.
就是不是說對等的.
一買一賣.
大家雙方的那種約的關係.
這種約的關係呢.
更加像以前的時代那些.
一個強國去打另一些國家.
接著打贏了.
我不殲滅你了.
但是我要和你建立一個約的關係.
我就是宗主國.
你就是那個從屬國.
那我呢.
和你簽了之後呢.
我就不會亂殲滅你.
或者亂對你的.
我會對你好好的.
那個從屬國呢.
也都會回應那個宗主國.
就說我會奉獻給你.

$^{281}$我會聽你的說話.
那我就隸屬於你了.
那我就會回應你了.
所以這個的盟約的關係呢.
更加像這一種了.
更加像這一種的.
宗主國和從屬國的關係.
而這份信念呢.
其實對以色列民來講呢.
不陌生的.
他們一直都知道這個關係.
因為呢.
由一開始的時候.
上帝聽到.
以色列這個民族的呼求.
在埃及做奴隸的時候.
他們的呼求上帝聽到.
所以就釋放了他們.
從奴隸當中解脫出來.
得到自由的這個盼望.
並且呢.
不單止得到自由.
還進入這個應許之地呢.
成為一個新的國家.
一個帶著上帝同在的國家.
所以他們很明白一件事.
不是無緣無故的.
突然之間.
他的國家很好.
哎呀 神就選你了.
我選你成為我的國家.
那你接著就成為從屬國了.
跟著我的律法去做呢.
你就可以換到好處了.
而他們心明的一件事就是.
他們從奴隸一個地方.
從艾輔.
從被逼北釋放出來.
得自由.
進入這個應許地.

$^{321}$而這一份呢.
就是被拯救的恩典.
也是我們一切的基礎.
而當我們呢.
其實你有留意到.
這個段落的時候呢.
雖然短短三節而已.
三節裡面呢.
你會看到有一樣東西很重要.
不斷提醒的.
就是今天這個字.
一開始十六節裡面就是說.
今天上帝吩咐你所進行的一切律例典章.
接著十七節就說.
今天你去宣認上帝是你的神.
第十八節.
今天上主上帝也向你宣認你是祂的子民.
三次的今天呢.
都是連繫著這個短短的段落.
甚至呢.
我們可以看到.
意思就是呢.
剛才說了.
這個段落由十二章到這裡.
你看回生命的一至十二章呢.
是在說重新一些的歷史.
一些以色列民在曠野和跟隨主的歷史.
到十二章打後至二十六章呢.
就是在說將來了.
他們去到這個應許地.
要跟著這些的律例去生活.
過去他經歷過這些.
將來要遵守這些的東西.
今天呢.
今天你知不知道自己和神是什麼關係.
而今天你知不知道自己和神是什麼關係呢.
是絕對左右了和決定了.
你將來會不會信服在神的帶領下.
將來我為什麼要跟上帝你說這些話.
就是因為今天我決心要這樣去做.

$^{361}$因為過往的這些經歷.
我今天知道前面要怎麼去走.
事實上呢.
以色列民是這樣經歷的.
他們經歷出埃及.
代表著上帝怎樣去顧念人.
也都回應人在痛苦裡面的呼求.
拯救受苦者脫離困苦.
今天呢.
今天我們都一樣.
舊恩的歷史不只是停留在他們出埃及的裡面.
今天也藉著耶穌基督的救贖.
去成就在我們身上.
令我們不是那些只是.
哦 是啊.
我覺得信耶穌其實都挺好的.
信耶穌好像在好人卡可以打個卡.
基督徒是好人來的.
所以我信耶穌在卡裡加難好處.
不是這樣的.
而是從罪裡面我們深感到.
我們在這些困苦.
這些不能夠自己擺脫的裡面.
釋放出來的自由.
因而呢.
能夠跟大君王結上這個關係.
所以其實我們跟以色列民呢.
有相似的地方.
但是為什麼剛才說.
如果沒有今天你很清楚這個應信.
和上帝是什麼關係.
會影響了將來你怎樣跟從神呢.
因為從經文裡面我們會看到.
人很容易呢.
走著走著又走回自己的路.
聽著聽著的時候.
又會只是想選擇那些自己覺得合聽的話去聽.
正如17節如果你看的時候.
他說什麼呢.
你今天宣認上帝是你的神的時候.

$^{401}$你就要遵守他的教導.
遵守這個原文本身是說走.
走在上帝的道路裡面.
謹守他的律例典章.
戒鳴並且聽他的聲音.
原來不只是口裡說.
今天我要信耶穌.
我要認上帝.
是你的腳和你的耳.
都要在今天去聽上帝的話.
去走在上帝的聲音裡面.
當然了.
以前那個年代.
明白的.
那種人的獨立性和個人主義.
還不算是很強烈的時候.
信從其實本身是基本.
就像我們上個年代聽父母說的.
聽長輩說的.
聽老師說的.
或者很久以前那些.
根本這個不需要疑惑的.
是啊.
他們的話就應該要聽.
但是今天當然主體比較重要.
個人化的年代.
要講信從有時候是相對比較危險的.
這個年代講信服很危險的.
當然不難明白.
因為當資訊很多的時候.
我們都會看到那些邪教.
或者那些有企圖的領袖.
怎麼經常說你們要聽.
要信服才可以.
用這種權柄去支配人.
所以很多時候今天.
當我們面對著打開聖經.
講信服和我們領受著的自由的時候.
這兩樣東西好像很對立.
為什麼對立呢?.

$^{441}$彷彿就好像如果說.
我要有個人的那種選擇.
我個人有一種自由的時候.
我不需要信服任何的東西去約束我自己.
我自己就會選擇到我自己覺得對.
我覺得自己做到一些好的事情.
那就是最好了.
這個就是今天自由對我們的教導.
一旦如果我表達.
我信服在任何的教導或者權威下面.
就好像失去了我的自由.
那我就不自由了.
我就要聽你所說的.
所以不難明白.
今天我們傳福音的時候.
和一些未信主的朋友就會說.
你這群基督徒都很沒有自由.
好像想跟你打牌又說不行.
想跟你去開心時光.
去談天說地開心一下.
去喝酒又說不行.
星期日去行山又說不行.
其實你們都很沒有自由.
然後又說要奉獻.
十一啊.
所以一聽就覺得好像很沒有自由.
好像要跟進要聽.
所以我們都明白.
為什麼社會的人或者周圍的人會這樣看.
也不難明白開始有些基督徒會轉化.
開始整合一下.
其實我也有自己的選擇.
那我就會選擇什麼?.
我做一個casual的基督徒.
不知道大家有沒有聽過這些詞.
那時候在大學裡面第一次聽.
有些教授說我是一個casual的基督徒.
我不是不信主.
我信主的.
但我不會跟你們那一套.

$^{481}$我知道你們信的是好.
聽你們的道理都很好.
我就不會像你們那樣.
就覺得這樣就是通往自由的道路.
我有真理.
但我有我自己的選擇.
我就會覺得這樣就最好了.
其實如果仍然覺得信服就等同失去自由的話.
其實是反映你還沒真正想通什麼是真正的信服.
什麼是真正的自由.
將這個問題放在心裡.
什麼是真正的自由?.
什麼是真正的信服?.
在你心裡沉澱一下再想一下.
我們看回今天的經文.
當你看回今天的經文.
16-19節很簡單的一個段落.
第16節說什麼呢?.
摩西他首先強調一件事.
你們聽完這些吩咐律例典章的時候.
你們要盡心盡性去跟從這個教導.
新漢語譯本裡面說什麼呢?.
他說要盡你全部的心思.
全部的性命去遵守這個教導.
意思是什麼呢?.
是思想和身體.
你的心思和你的生命就是你所想的和你所做的.
不只是說我思想上認同你.
但我身體不動.
OK的你說得不錯.
但我繼續.
又不是那種我做出來.
但其實我腦子裡不認同你.
我聽你說我才這樣做.
裡面也不是說分裂了這兩樣東西.
是怎樣盡心盡性.
盡你全部的心思和全部的性命去走這件事.
而是兩者都要盡.
不容易的真的不容易.
當然那個不容易.

$^{521}$不單只是說不容易去做到.
第二方面那個不容易.
是別人都不知道你有沒有盡.
因為那個盡.
真的只有你自己才知道是不是盡.
因為有很多很簡單的比喻.
可能有些大英姐妹看到.
葛邦你好盡力啊.
然後我心想你都不知道我吃了兩個午餐.
吃了兩個小時午餐.
你見到我很盡力其實我可能很輕鬆.
反過來可能你說.
葛邦這個活動你盡力找人.
你找到三四個就.
對我來說已經很盡力了.
找三四個已經很好了.
已經很盡力邀請你了.
那個盡.
別人不知道.
你有多盡.
其實是你自己和上帝才知道.
甚或乎如果你看回新約裡面.
那個使徒行傳的經文.
有一對夫婦欺控聖靈.
他就賣了財產然後收了一部分.
然後去到使徒面前說.
所有來的了.
事實上使徒都不知道是不是盡的.
唯獨是聖靈的感動.
彼得問是不是所有.
他說是呀是所有的了.
他也認了是所有的了.
然後就撲倒了.
就因為他有沒有盡力.
他們自己知道和上帝知道.
所以這個盡.
我認為不是只是凸顯你擺上多少.
或者多努力讓人看見那一種.
而是從心裡面.
有沒有那種徹徹底底的去遵行.

$^{561}$所以當你理解什麼是真正盡心.
盡力的時候.
你就會明白.
信服其實是真正的自由.
為什麼這樣說呢.
當你盡心盡力去強調我們去信服上帝的時候.
這種是在說我們裡面.
是很喜歡上帝的律例.
我很喜歡去選擇去進行神的教導.
這個是我們心靈裡面.
深處的一種渴望.
事實上如果沒有自由意志的信服.
不是真信服.
如果你也沒有選擇.
我逼你.
你跟著我.
不然你就死定了.
這種信服不是一種從心裡面的真.
真正的信服是要有那種.
我自由的選擇.
願意有這種自由意志之下.
我仍然選擇跟從你.
這種才是真正的信服.
沒有自由的信服.
不是真正的信服.
反過來.
沒有信服的自由.
其實不是真正的自由.
這份自由不是真正領略自由的寶貴.
底牌這個自由其實是自我.
我喜歡做什麼就做什麼.
或者我自己的選擇.
我們都知道.
就好像愛神一樣.
或者愛其他人一樣.
愛不是一種衝動一時之間的事.
愛是意味著隨時準備為你愛的人.
盡力獻上你的生命.
所以最深刻的表現之一.
就是將我們的意願.

$^{601}$跟我們所愛的人的意願.
連合在一起.
你想要的東西.
就是我想要的東西.
而你想知道的東西.
就是我想知道的東西.
我要這東西.
是因為你要它.
正如你要它的時候.
我也覺得在它身上得到這份好處.
當你能夠有這份愛.
和這份心底裡的愛慕和渴慕的時候.
你就能夠連合到.
原來信服這東西.
是自由.
是你的選擇.
也因為信服是一種彼此的答應.
其實不只是彼此的約定這麼簡單.
其實真正的自由也是一樣.
不只是單方面的.
你這樣做就行了.
好吧,你這樣說我就這樣做.
自由不只是說很自由自在.
我將自身以外的其他人.
或者其他他者.
所有都不管.
那我就最自由了.
不用理會其他人.
總之我做所有事情就行了.
將自身以外的事.
所有都不管.
不是說自由.
我在我覺得正確的路上走.
就是我的自由.
你在你覺得正確的路上走.
就是你的自由.
但是你說你要我走你的路.
那我就覺得不是自由了.
如果你仍然停留在這種裡面.
都是和這些關係割裂分開.

$^{641}$彼此的約定.
其實也是信服.
也是自由.
就好像我們和上帝.
這個關係.
說完這個可能很理論了.
簡單一點.
拿一些生活例子來做一個簡單的想法.
雙方而已.
不是完全的.
拿旅行來做一個簡單的例子.
大家喜歡自由行還是跟團?.
當然是很自由的分兩派的.
當然喜歡自由行.
喜歡自己去哪裡就去哪裡.
選擇時間或者去哪裡都可以.
但是有些人覺得跟團也不錯.
你帶我去所有地方.
有些人覺得自由行當然是好.
因為自由行通常不是計劃的那個.
就是你去的那個.
通常交給一個人.
然後交給他就想好了.
我們自由行.
跟著去吧.
如果你真的自由行.
你是負責想的那個.
和跟足時間表走的那個.
其實你不是很自由.
你跟著家裡的人.
去吧去吧.
夠鐘了夠鐘了.
要趕這班車了.
那你自然就要跟著他們去.
帶著他們去.
然後去到那些景點.
你又要做很多資料搜集的.
有沒有試過.
打算大條道理去到那裡.
一個館去參觀.

$^{681}$怎知道原來星期三休館.
而你的schedule本身就是安排星期三去那裡.
然後通常就說.
你不看的嗎?.
你不知道的嗎?.
那我拿到的資料就是這樣.
自由行就是這樣.
那反過來.
可能對年輕一輩.
就會覺得你們跟團當然不好.
早上九點多就要起床了.
然後定時定點要去這些地方.
但是最近我有一個新的體會.
是什麼呢?.
就是因為早前三月的時候.
我要帶隊去馬來西亞.
帶一些教務同工.
去當地一些教會去參觀.
和看看他們的服侍.
和一些年輕人的工作.
有一個很深刻的團友去分享.
說什麼呢?.
其實原來跟著你們去都挺好的.
又不用安排車.
又不用安排景點.
然後這些東西顧慮放下之後.
他真的可以在教會那班弟兄姊妹的時間.
很自由地和他們相處去體會這種人情.
世味當中的那種教導.
不用那個朋友想.
什麼時候要走.
什麼時候要跟下一個行程.
什麼時候要去做這件事.
他反而在這種跟團裡面.
得到真正的自由.
其實正如剛才所說.
跟團和自由行不一定是相對的.
但是最慘的是哪些人呢?.
就是跟團的時候.
有機會說有沒有搞錯.

$^{721}$這行程是這樣的.
又這樣或者什麼.
都自由行了.
又說為什麼去這樣的地方.
自由行而已.
都不用去這樣的地方.
和我去旺角差不多而已.
我聽過的.
安排完行程.
跟我在香港去旺角一樣.
然後就馬上說有沒有搞錯.
帶你來之後說這樣的東西.
最慘就是這種.
自由行又嘰哩咕嚕.
跟團又嘰哩咕嚕.
最終因為他忘記了旅行的目的.
其實自由行也好.
跟團也好.
旅行的目的都是想開開心心.
去體會一下當地的行程.
想知道一下當地的生活是怎樣.
一些風土人情.
其實沒有相對哪樣東西對.
哪樣東西錯.
但如果這種心態的時候.
就失去了旅行本身的意義.
正如回到這段經文裡面.
咦.
當我們在跟從著律法.
在那種自由和信服的那裡.
嘰哩咕嚕.
好像跟了你我就沒有自由了.
好像自由自在.
沒有了那種律法的.
那種東西好啊.
其實最終律法背後的意思是什麼呢.
你和上帝的關係.
和這些律法的關係是什麼呢.
我們一定要細思.
以致當我們知道這個自由的基礎是什麼的時候.

$^{761}$就是因為我們和上帝之間這個愛的關係.
被救贖的關係成為基礎.
以致我們能夠在愛裡面去信服上帝的教導.
當然另一邊廂.
當你看回整個經文的時候.
難道你覺得.
以色列人信服了之後.
他比其他的人自由了嗎.
你看看.
我們這群以色列人跟從上帝的律例典章之後.
我們多好多自由.
其實你看起來可能不是.
其實對當時來說.
以色列只是一個很小的國家.
比起周圍的強國.
那些很強的國.
很差的.
軍力又弱.
都不用計算那些GDP.
就是國民生產總值.
小到無論.
真的是一個很小的國家.
難道你們跟了上帝的律法就好一點嗎.
就自由自在嗎.
其實你可以反過來想.
如果你不跟上帝的律法.
難道你又覺得迦南人自由自在嗎.
他們很開心.
如果真的有機會.
看看當時的人情或者歷史.
迦南人拜偶像的習俗.
比律法裡面的跟從還要可怕.
有機會再說.
就是考古上的那些.
所以整個的時候.
如果你將兩件事放在一種相對的時候.
很容易就失去了本身律法的好處.
其實本章這個段落.
最後的時候.
也提醒我們一樣東西.

$^{801}$遵行律法本身.
不只是目的.
我做這件事就是要達到這個目的.
去換來上帝的保障.
換來上帝的一些好處.
其實最終來說.
律法真的出於這份愛神.
出於愛我選擇.
放棄我這些選擇.
選擇去跟從上帝.
也在裡面告訴我們一樣東西.
原來當我們在這些律例.
將這些上帝聖經的教導裡面.
進行的時候.
原來還有一種宣教的目的.
當你看回第十八節裡面.
和十九節裡面說什麼呢.
當你願意宣認上帝是你的神.
上帝也宣認你是祂的子民.
祂就會高舉你.
過於祂所做的萬過.
使你成為聖潔的子民.
歸於上帝你的神當中.
原來當我們願意.
以這個律法的基礎.
去建立這個關係.
和在裡面願意跟從的時候.
神的名就在我們中間被高舉.
為什麼你會做這些事.
為什麼你會這樣選擇.
因為我愛我的神.
我的神想我這樣做.
我會這樣做.
以至在當中的時候.
神的名就在萬民當中.
讓人去認識受尊崇.
所以原來律法典章的作用.
除了幫助以色列人.
找回他的身份.
和知道他的使命之外.

$^{841}$並且在當中成為一種很重要的見證.
當神的子民.
忘記了上帝的恩惠.
和忘記了這個身份的時候.
就失去了那份愛神的心.
會怎樣呢.
就會扭曲了律法.
通常只有兩個結果.
一是完全走向極端的律法主義.
就想著跟足那些來做就行了.
上帝你喜歡我做這些事.
你不要理我.
總之我做給你看.
又或者另一個極端.
就是反律法主義.
為什麼要跟這些.
我自由來的.
我不做這些.
這些傾向.
是因為他失去了神的那個恩典.
失去了我們作見證的那種動力.
有一個靈修學者叫做傅氏德.
他說了一件事.
真正的馴服.
是一定要運用自由.
而馴服的結果.
是為了得到更多的自由.
就是要從那些非要不可.
照我的意思去做的那些想法.
那些執著釋放出來.
於是你就更加深空廣闊.
對人和對身邊的事.
那份體諒就會多了.
有更多的自由去愛身邊的人.
神既然將我們從罪惡和不義.
那種自我自私的世界裡釋放出來.
我們要用這份真正的自由.
去愛身邊的人.
去回應上帝的愛.
而這份馴服.

$^{881}$也是運用在基督裡面的自由.
去選擇那些好像體力很不自由.
又不能做這件事.
又不能做那件事.
自願放棄這些權利.
自動地將自己的身份.
放在一種獨人的身份.
來服侍上帝.
來服侍教會.
而不是那個心裡面.
就是心不甘情不願.
無可奈何.
你這樣說我這樣做.
也不是懼怕刑罰.
也不是出於只是怕任何的事.
而是出於愛神.
而作出的自我約束.
當你有這個心和見證的時候.
就成為了今天這個世代.
所謂純相各種自由的.
一種逆向的見證.
為什麼你會甘願選擇這一種生活呢?.
是因為我享受真正的自由.
到最後.
在人類歷史上.
我們很清楚知道一件事.
沒有任何的人.
或者任何人的行為.
彼得上我們的耶穌基督.
在十字架上.
這個最自由.
最馴服的選擇.
是自己選擇走上這個十架的路.
為我們犧牲.
讓我們帶著這份心去回應上帝的愛.
我們一起祈禱吧.
神愛主,求你去豐富我們的心.
使我們的心得到真正的自由.
並且完全信服在主你的大靈.
即使今天我們內心有很多疑惑.

$^{921}$神愛主,求你一路去光照.
繼續讓你的說話.
成為我們生命每一個選擇的燈和光.
去照亮我們.
以致我們能夠成為.
國道裡的子民的見證.
能夠將你的名在萬國裡被高舉.
我們祈禱.
是奉主耶穌基督的名字而求.
阿們.
謝謝大家.
\newpage



\section{申命記 30:11-20}
\label{sec:dxgI8DcWOBk}
\textbf{2025年6月1日主餐崇拜直播｜陳冠成牧師｜朝向應許之地－務要揀選生命｜申命記30\_11-20}
\newline
\newline
連結: \href{https://youtube.com/watch?v=dxgI8DcWOBk}{\texttt{https://youtube.com/watch?v=dxgI8DcWOBk}} ~~~~ 語音日期: 2025-06-01
\newline
\newline
\hyperref[sec:wyPYrchtH6g]{\small{< < < PREV SERMON < < <}}
~
\hyperref[sec:index_chronic]{\small{[返順時目]}}
~
\hyperref[sec:Buo5z3C6Kgc]{\small{> > > NEXT SERMON > > >}}
\newline
\newline
申命記 30:11-20
\newline
\begin{longtable}{cl}
\hline
\hline
章節 & 經文 (和合本修訂版)\\
\hline
30:11 & \begin{tabularx}{0.7\textwidth}{X} 「我今日所吩咐你的誡命，對你並不困難，也不太遠； \end{tabularx} \\ \\ \relax
30:12 & \begin{tabularx}{0.7\textwidth}{X} 不是在天上，使你說：『誰為我們上天去取來給我們，使我們聽了可以遵行呢？』 \end{tabularx} \\ \\ \relax
30:13 & \begin{tabularx}{0.7\textwidth}{X} 也不是在海的那邊，使你說：『誰為我們渡海到另一邊，去取來給我們，使我們聽了可以遵行呢？』 \end{tabularx} \\ \\ \relax
30:14 & \begin{tabularx}{0.7\textwidth}{X} 因這話離你很近，就在你口中，在你心裡，使你可以遵行。 \end{tabularx} \\ \\ \relax
30:15 & \begin{tabularx}{0.7\textwidth}{X} 「看，我今日將生死禍福擺在你面前。 \end{tabularx} \\ \\ \relax
30:16 & \begin{tabularx}{0.7\textwidth}{X} 我今日所吩咐你的，就是要愛耶和華－你的神，遵行他的道，謹守他的誡命、律例、典章，使你可以存活，增多，而且耶和華－你的神必在你所要進去得為業的地上賜福給你。 \end{tabularx} \\ \\ \relax
30:17 & \begin{tabularx}{0.7\textwidth}{X} 倘若你的心偏離，不肯聽從，卻被引誘去敬拜別神，事奉它們， \end{tabularx} \\ \\ \relax
30:18 & \begin{tabularx}{0.7\textwidth}{X} 我今日向你們申明，你們必定滅亡；在你過約旦河進去得為業的地上，你的日子必不長久。 \end{tabularx} \\ \\ \relax
30:19 & \begin{tabularx}{0.7\textwidth}{X} 我今日呼天喚地向你作見證：我已經將生與死，祝福與詛咒，擺在你面前。所以你要揀選生命，好使你和你的後裔都得存活。 \end{tabularx} \\ \\ \relax
30:20 & \begin{tabularx}{0.7\textwidth}{X} 要愛耶和華－你的神，聽從他的話，緊緊跟隨他，因為他是你的生命，必使你的日子得以長久，可以在耶和華向你列祖亞伯拉罕、以撒、雅各起誓要給他們的地上居住。」 \end{tabularx} \\ \\
[1ex]
\hline
\hline
\end{longtable}
$^{1}$今日的經分是記載在申命記三十章十一至二十節.
是在你們的聖經舊約二百五十一至二百五十二頁.
聖經那裡這樣記載.
十一節.
我今日所吩咐你的誡命.
不是你難行的也不是離你遠的.
也不是在天上使你說.
誰替我們上天取下來.
使我們聽見可以進行呢?.
也不是在海外使你說.
誰替我們過海取了來.
使我們聽見可以進行呢?.
這話卻等你甚勤.
我就在你口中.
在你心裡使你可以進行.
看那我今日將生與死與福.
死與和陳明在你眼前.
吩咐你愛耶和華你的神.
遵行他的道.
遵守他的誡命律例典章.
使你可以存活.
人數增多.
耶和華你的神就必在你所要進去得為業的地上.
賜福於你.
倘若你心裡偏離不肯聽從.
卻被勾引去敬拜事奉別神.
我今日明明告訴你們.
你們必要滅亡.
在你薄弱彈劾.
進去得為業的地上.
你的日子必不長久.
我今日呼天喚地.
向你作見證.
我將生死與福.
陳明在你面前.
所以你要選擇一個生命.
使你和你的後裔都得存活.
且愛耶和華你的神.
聽從他的誡命律例.
專靠他.

$^{41}$因為他是你的生命.
你的日子長久也在乎他.
這樣你就可以在.
耶和華向你列祖亞伯拉罕與撒雅各起誓.
應許所賜的地上居住.
請聽到畢.
弟兄姊妹平安.
剛才所讀的這段申命記.
其實連同上一章第29章.
其實是摩西他離開世界前.
我們說最後一次對以色列人的宣講.
你聽到他告誡將要進入迦南地的新一代百姓.
要他們在剛才你聽到的兩條道路裡.
一條是關於生命關於蒙福.
另一條是關於死亡遭禍患.
摩西要他們在當中做一個二選一的抉擇.
我們不妨先想像一下.
摩西宣講的時候.
他帶著什麼情懷.
帶著什麼語氣.
很多時候中國人說人之將死.
是如何?.
其言也善.
一定是語重心長.
一方面他想到這班百姓.
終於可以進入迦南.
固然很值得興奮喜悅.
但摩西同時也對這班百姓多多少少充滿憂心.
為什麼這樣說呢?.
因為畢竟摩西跟他們太熟了.
每天都對著.
其實這幾十年.
摩西很了解他們的屬靈狀況是如何.
他們跟上帝的關係到底是如何.
這四十多年來.
以色列人雖然經歷了上帝的倚靠.
經歷了嗎哪.
經歷了各種的公認.
但同時他們也對上帝有很多投訴.
有很多不相信.

$^{81}$甚至一些叛逆.
雖然他們已經跟上帝建立了盟約的關係.
也從上帝領受了戒命律例典章.
但我們在《五經》中看到.
百姓始終不能貫徹實踐上帝一切的吩咐.
所以摩西在他臨離開世界之前.
再一次鄭重地勸告.
他帶領牧養了幾十年的百姓.
一定要認認真真地看待上帝所說的話.
摩西在這段經文中強調了三次.
「今日」.
你看看程柱理表的廣東大綱.
已經寫明了.
很清楚的三個段落.
三個段落都用「今日」.
其實是一種迫切的感覺.
迫使百姓不能避開.
上帝的話語不能迴避.
上帝的吩咐不可以側側膊.
不能說偷懶.
有空才做.
不要等到明天.
不要等到將來.
你現在就要回應.
面對上帝的話語.
人唯一能夠做的就是.
做正確選擇.
不能迴避.
我們逐個段落來看.
首先第十一至十四節.
摩西在這裡說.
我今日所吩咐你的誡命.
不是難行的.
也不是離你遠的.
意思是什麼呢?.
不難行的其實是不難理解的.
也不難執行.
上帝也說離你不遠.
即是說你很容易聽到.
你很容易聽取到上帝的話語.

$^{121}$不需要像經文形容.
摩西說的.
不需要找人上天下海.
好像我們以前看電視.
或者看小說西遊記.
要找團人攀山涉水.
長途拔峽去取經.
上帝說不需要的.
你不需要這麼長途拔峽.
你才能將律法帶來給回族.
上帝是怎樣的?.
上帝是啟示的.
上帝主動將他的誡命律例典章.
啟示給百姓知道.
你不知道上帝的心意是怎樣.
這是我們的限制.
所以上帝主動告訴我們.
他的心意是怎樣.
怎樣可以得到生命.
怎樣可以得到拯救.
這是基督教和其他宗教語別不同的地方.
由上而下地啟示.
讓百姓可以聽到.
可以進行.
所以摩西的表示是.
上帝的吩咐對百姓來說.
不是離身的.
不是遙遠的.
是距離很近.
百姓可以隨時接觸到.
親近到.
也就是說他可以聽到.
他可以讀到.
他也可以背上帝的話語.
就像經文說的.
放在心上.
隨時隨地在你生活中.
就拿出來實踐.
換句話來說.
大家懂得聽摩西這句話嗎?.

$^{161}$背後其實是說.
原來認識和進行.
誡命上帝的吩咐.
是每一個上帝的子民.
包括你和我.
無可推諉的責任.
不能說我不想.
不能說我不願意.
假如有個以色列人說.
我真的沒有時間.
我要謀生.
我沒有辦法.
我不是很聰明.
我真的沒有那麼多的能力.
去認識和遵守上帝的話語的話.
試圖想迴避這個責任.
你猜摩西會怎麼說?.
摩西如果跟從這段經文.
他就會斬釘截鐵地說.
不可能的.
上帝是有恩典.
有憐憫的神.
他不會訂立一些要求.
不會訂立一些吩咐.
是他的百姓做不到的.
上帝不會這樣.
上帝說得出的東西.
是一定做得到的.
百姓是一定有能力去辦妥.
之所以沒有辦法遵行.
配合上帝的吩咐.
其實不是人欠缺能力.
不是欠缺時間.
不是欠缺方法.
而是缺少了什麼?.
對上帝的話語的尊重和信服.
或者說得白一點.
對上帝的愛不夠.
愛自己多過上帝.
事實上如果有留意.

$^{201}$我們之前在申命記第八章裡.
也說過.
上帝透過嗎哪的功課.
用了四十多年.
讓百姓知道.
人活著是靠什麼.
記不記得那句金句.
耶穌基督也用過.
人活著是靠.
神口裡所說的一切話.
嗎哪表面上的功能.
是填飽肚子.
但實際是引導在抗疫的百姓.
這四十年來.
你盡照上帝的吩咐.
來拾取享用嗎哪.
讓他們知道.
其實不是嗎哪本身讓他們存活.
是那位公認嗎哪的上帝.
以及他的吩咐.
令百姓存活.
所以對百姓來說.
上帝所說的話.
自然是包括什麼?.
他和百姓立約的內容.
對百姓來說.
就是遵行上帝的吩咐.
所以在抗疫裡看到.
上帝一方面履行了他的承諾.
他帶領照顧這群百姓.
降下嗎哪.
同一時間.
百姓也履行了他在約裡的承諾.
他們願意遵照上帝的吩咐.
來享用嗎哪.
兩件事加起來.
就構成了以色列人存活的關鍵.
少一件事也不行.
讓百姓真實體驗到.
原來我們是靠上帝和他的話語.

$^{241}$來過生活.
從而懂得尊重和信服上帝.
所以弟兄姊妹.
今天我們當然不會像抗疫的以色列人.
我們要透過拾取嗎哪.
來做一個屬靈操練.
但是什麼是我們今天要培養的.
一個屬靈生命.
以致我們懂得重視和尊重上帝.
如果對應摩西這裡的吩咐.
就是在說我們有沒有一個習慣.
去認識上帝的話語.
你有沒有培養一個讀經的習慣.
讓上帝的話語來塑造你和我的生命.
將聖經的價值觀.
將裡面所說所吩咐的言行.
在你的生活裡實踐.
我聽過德國人有一句慣用語.
叫做你用什麼來做糧食.
你就會變成一個怎樣的人.
意思是什麼呢?用廣東話.
你用什麼你就會變成那樣的人.
舉個例子我以前初出來工作.
我跟一個上司.
他是一個韓國人.
知道韓國人喜歡吃什麼嗎?.
是劍鰭.
你不要這麼地道.
韓國人很喜歡吃生蒜頭.
我老闆特別喜歡吃.
每一餐都會有生蒜頭.
就這樣放進口裡.
另外他很喜歡吃煙.
他那時候應該四十多歲.
如果根據這句話.
他吃什麼就變什麼.
你猜他變了什麼?.
他每次走過.
他不需要開口.
我就會發覺有顆蒜頭在我旁邊飄過.

$^{281}$有支點著的煙在我旁邊經過.
他就是那支煙.
他就是那顆蒜頭.
所以你有沒有吃蒜頭.
我一定知道.
你不用開口.
我走過就聞到.
吃了蒜頭.
準確率我想達八成以上.
我現在還記得那股味道.
換句話來說.
你吃什麼.
你就會變成一個怎樣的人.
其實也適用在其他生活層面.
最簡單來說.
人如果缺乏一些優質均衡的正常飲食.
你總是喜歡吃麥當勞,薯片.
各種煎炸的食物.
這些無益有害的東西的雜糧.
遲早你的身體會出事.
同樣我認為作為基督徒.
如果你沒有經常用聖經上帝的話語.
餵養你的生命.
反而你用其他的東西.
你用媒體.
你煲劇.
你用不同的資訊.
很多扭曲了,雜亂的資訊.
霸佔了你的心靈.
其實說真的.
你很難孕育出上帝的喜悅生命.
再白一點說.
上帝想怎樣.
你知不知道?.
聖經說什麼.
我們知不知道?.
我們可能真的不知道.
所以弟兄姊妹.
經文第一個提醒我們.
我們有沒有認真看待上帝的話語.

$^{321}$在日常生活裡面.
你有沒有好好抽時間.
去研讀你手上的聖經.
可能我們因為種種原因.
你沒有參加教會讀經計劃.
但你沒有理由拒絕讀經的責任.
摩西在這裡說了.
無可推諉.
你用你的方式.
將上帝的話語放在心裡.
弟兄姊妹.
摩西首先要我們培養一個.
重視和信服上帝的態度.
我們繼續看下去.
我們看第二個今日.
第二個今日在十五節到十八節.
摩西這裡很清楚.
擺了兩條路出來.
第一條就是步向生命.
一個祝福的道路.
一個蒙福的道路.
走在這條道路的百姓.
他會不斷用行動來建立和上帝的關係.
經文這裡說了三件事.
三個動詞.
第一個愛上帝.
不可能沒有愛.
你讀經不可能沒有愛上帝的態度.
否則你很難讀下去.
第二個進行.
不只是看.
還要進行上帝的吩咐.
還要緊守他的誡命.
什麼意思呢?.
即是你遇到挑戰的時候.
你很想讓自己的意見先行.
放下上帝的心意的時候.
你緊守上帝的心意.
仍然抓緊上帝給你的吩咐.
應難而上.

$^{361}$如果你這樣做的時候.
摩西就說了.
上帝會賜福給百姓.
又是有三件事.
首先你最底層的那一層.
存活.
然後呢.
你可以開始有自己的後裔.
人口眾多.
不只是你能夠有好的生養.
是你也能夠養活.
即是你開始安居樂業.
並且你能夠在迦南德地為業.
即是說.
你可以勝過你周邊的外敵.
不被攪擾.
你可以可持續發展這個民族.
這是第一條路.
第二條路.
反過來.
是一個步向死亡.
是一個會遭禍患的道路.
走在這個道路的百姓.
不單止心會慢慢偏離上帝.
不聽上帝的吩咐.
更加會被勾引去敬拜事奉別神.
對我們今天來說.
偶像我們之前也說過.
不只是假神有專像.
那個霸佔了你心目中的守衛.
但不是上帝.
其實那個已經是偶像.
可以是你的家人.
可以是你的子女.
可以是你任何的慾望.
摩西話.
走在這條道路的百姓.
有一個很嚴重的警告.
就是會滅亡.
就算過了約旦河.

$^{401}$進入了迦南.
日子也不會長久.
事實上這種情況是最難處理的.
為什麼呢?.
因為他不是完全不相信耶和華.
他覺得自己還是耶和華的子民.
有憲制.
有敬拜的生活.
不過有些時候.
我左右逢源.
他又好像是耶和華那邊.
但又好像是其他偶像那邊.
我們說有些魚.
在鹹淡水交界.
兩邊都可以生活.
鹹水也可以.
淡水也可以.
上帝的話語.
我又可以守著.
迦南的風俗.
世界社會的世俗的想法.
我又可以跟著.
兩邊都領了好處.
這樣就最好了.
世俗裡也有很多教我們這樣想.
不過摩西話死路一條.
不要想著左右逢源.
上帝是那位嫉妒忌邪的上帝.
你只能選擇他.
不可以選擇他以外的上行.
來敬拜.
等於姐妹摩西展示了這兩條道路.
其實我們會怎樣選擇?.
你又會引導你的下一代.
選擇哪一條路?.
其實你看到這段經文.
不需要任何複雜的釋經.
每個人都聽得明白.
什麼叫生死.
什麼叫禍福.

$^{441}$每個人都懂得怎樣選擇.
因為摩西在宣講的時候.
他已經假設了.
任何頭腦稍為清醒的人.
是不會選錯的.
不會把頭放在死亡,遭禍的那條路上.
不過現實也告訴摩西.
真的有人選錯.
因為人的本質是軟弱的.
確實有很多以色列人.
他見過不自覺地.
明明在走一條存活無福的道路.
但走著走著就走到旁邊那條路.
走到滅亡,遭禍的道路.
如果我們有機會問那些昔日.
被上帝懲罰死在江南.
死在曠野.
不能進入江南的那些以色列人.
你問問他們.
為什麼他們那麼傻?.
為什麼他們那天選了這條死亡.
遭禍患的道路來走?.
你覺得他們會怎麼回答?.
大概我想他們可能會說.
其實我也不知道.
我不知道為什麼會這樣.
我記不起.
但就是這樣發生了.
但有一件事我可以很肯定.
我一開始不是這樣想的.
我沒有這樣打算過.
相信沒有一個以色列人.
他會說這條路是死路.
是滅亡之路.
不過我試試是否真的.
我走過去試試.
我挑戰一下上帝說的話是否真的.
我想沒有人會那麼大膽.
這樣去挑戰上帝.
同樣我想沒有一對準新人.

$^{481}$會在一個婚禮當中.
抱著早晚都會分手.
我試試進入婚姻.
然後發這個婚姻的誓言.
沒有人會那麼傻.
也沒有人會期待.
在上帝的婚禮上.
期望有一天.
就被弟兄姊妹說了.
期望有一天.
就會跌倒或是半倒人.
沒有人會這樣想去事奉.
也不會有父母.
在生小朋友之前會想.
我將來都沒有時間去理會他.
我將來的關係都不好.
你不會帶著一個奇怪的期望.
而進入一個生活裡.
但是人有時就是這樣.
或者像保羅所說.
立志行善是由得我.
不過行出來卻身不由己.
我本意是想聽上帝的話語.
我本意是想將生命托付給上帝.
好好在我不同的人生範疇.
會出上帝的旨意.
經驗他所有的福氣.
但是走著走著.
就好像第一代以色列人一樣.
我不知不覺因為.
我貪戀自己的需要.
我比較自我.
我覺得我要將生命拿回來控制.
於是我就將我的生命根基.
建築在上帝以外.
建築靠攏在那些很短暫.
那些有限.
經不起考驗的事物裡.
我將我的生命慢慢遠離上帝.
也遠離了上帝所賜的生命.

$^{521}$我還是存活.
不是因為我行得對.
是因為上帝憐憫.
是因為上帝仍然有恩典給我們.
希望我們回頭.
上帝想我們選擇生命.
這是摩西一直強調.
到底什麼是生命呢?.
第二十次摩西說.
上帝不就是你的生命嗎?.
好像很奇怪.
好像很哲學.
什麼是生命呢?.
上帝就是你的生命.
你選擇上帝就是選擇生命.
意思是當我們選擇.
聽上帝的吩咐.
遵守上帝的誡命.
走上上帝的道路的時候.
你會洞悉什麼是生命的本相.
因為你連結上帝創造你的主.
你會明白上帝創造你的心意.
到底是什麼.
你開始意識到.
你要怎樣活一個上帝喜悅的人生.
其實很簡單的道理.
你不要插錯插蘇.
你將你連結生命的插蘇.
插回上帝身上.
你接駁上帝.
通上帝就行了.
你持續連結於.
那位賜給你生命的上帝.
連結於上帝給你的生命之道.
那就是你生命的根本.
所有事情就會水到渠成.
所以摩西在這裡說.
其實選擇生命不是選擇一次.
同樣你決志其實不是缺一次.
是你之後每一天都決志.

$^{561}$由面向世界.
依靠世界.
轉向依靠上帝.
依靠耶穌.
是每一天每一次.
每一種抉擇的時候.
你都要選擇上帝.
選擇生命.
就是你不斷回應上帝給你的吩咐.
選擇走上帝的道路.
特別當你落入了試探.
當你的人生混亂.
很多事情搞不定的時候.
你被罪核制的時候.
你問一問你自己.
上帝想不想你這樣.
你自己想不想這樣.
然後你再問一問自己.
你選擇的是一條什麼路.
會不會像摩西這裡所說.
走一條遭禍死亡的道路.
然後你選擇要不要回頭.
要不要回轉.
選擇一條上帝喜悅的道路.
弟兄姊妹.
摩西就是這樣勸勉我們.
你今天就做選擇.
你今天就要做決定.
現在就是你改變人生的機會.
上帝透過這段話語.
提醒面對我們.
你又會怎樣選擇呢?.
或者我們看看摩西最後.
想我們怎樣選擇.
我們讀19至20節好不好?.
最後一個段落了.
預備起.
(上帝說話).
你的日子長久也在乎他.
這樣你就可以在耶和華.

$^{601}$向你列祖亞伯拉罕.
以撒瓦各喜誓.
應許所賜的地長居住.
你看到摩西擺出兩條道路.
好像可以選擇.
但其實是不能選擇.
因為40年來.
他見太多人選錯了.
救都救不回.
於是他就直接.
怎樣?.
他幫我們選了.
摩西說你要選擇生命.
你要選擇上帝.
因為上帝就是你的生命.
什麼是選擇上帝呢?.
就是你信服上帝的吩咐.
過於信服你自己的心意.
你單單愛上帝本身.
多於愛上帝給你的好處.
這就是選擇上帝.
其中一個最具體的例子.
我們也很熟悉.
阿巴拉罕獻以撒.
我們知道上帝賜給阿巴拉罕一切福氣.
所有的應許.
全盤的計劃.
都要透過阿巴拉罕的兒子.
以撒來成就.
換句話來說.
阿巴拉罕失去了以撒.
就等於失去什麼?.
失去所有上帝答應他的福氣.
好處等等.
所以.
這個獻以撒的吩咐.
對阿巴拉罕來說.
一定是相當糾結和掙扎.
事實上.
他真的可以和上帝上訴.

$^{641}$爭論.
上帝.
我等了二十五年.
你拿他.
現在才賜以撒給我.
你為什麼要收回?.
你違約.
他有權這樣和上帝爭論.
我不明白你想怎樣?.
上帝,你解釋清楚.
沒有.
阿巴拉罕.
他沒有這樣做.
即使是這樣.
阿巴拉罕仍然願意獻上以撒.
讓我們看到幾點.
第一個就是他很相信.
誰是生命的主宰.
人最終是回到上帝身上.
上帝有權賜生命.
也有權收回.
阿巴拉罕很清楚.
即使是以撒.
上帝賜的.
上帝有權收回.
第二點就是.
他很明白.
雖然上帝有安排.
有吩咐.
我真的不明白.
我真的不懂.
因為我是人.
我理解不到上帝那種超越的想法.
上帝也沒有和我解釋.
他也沒有必要和我解釋.
因為他是上帝.
但我相信上帝一定有他的理由.
只是我不明白.
我可以選擇相信.
他仍然愛我.

$^{681}$事實上他由始至終都要賜福給我.
他不會騙我.
也從來沒有騙過我.
他沒有虧欠過我.
所以.
阿巴拉罕這份信心讓他確信.
如果人生有些事情只能夠二選一.
他選擇什麼?.
要不選擇上帝.
要不選擇上帝給你的好處.
兩樣都好.
選擇什麼?.
選擇最好的那樣.
上帝本人.
次好的那樣.
上帝的好處.
都可以不要.
他懂不懂得選擇?.
容易不容易選擇?.
不容易選擇.
所以你明白.
信心之父不是浪得虛名.
你不要看他創世記好像不太好.
他有很多事情.
你和我都未達到他的目標.
所以,弟妹們.
你會看到阿巴拉罕這種堅持選擇上帝.
堅持追求上帝的旨意成就.
將上帝的心意放在第一位.
這個人生態度.
其實也在挑戰我們.
因為選擇上帝.
選擇生命其實意味.
我們要為神的緣故.
你可能真的要放下一些東西.
放下你正在抓緊的東西.
你正在追求的事物.
那些可能是你很想要的財富.
可能是你的事業.
你的學業.

$^{721}$你的愛情.
你的家庭.
甚至你的健康.
這些事情本身不是壞事.
這些事情.
上帝都可以用它們成為祝福你的途徑.
但是我們要明白.
它們不是上帝.
你一旦做錯了.
你就走在另一條路.
如果有需要的時候.
上帝會幫我們放下.
我們說是次要的福氣.
迫使我們專一跟隨他.
得到最好的福氣.
就是上帝他自己.
所以摩西.
你看到他這裡.
呼天喚地.
他呼喚天地為他宣講.
作見證.
他很期望百姓真的可以認真聽他.
這段語重心長的勸勉.
你選擇上帝吧.
你選擇生命吧.
摩西對百姓.
我們說對他們的數年生命.
那種執著.
那種迫切和欲緊.
毋庸置疑.
但是我們也看到.
當百姓進入迦南.
過了幾代之後又怎樣?.
又固態復萌.
又再出現很多不信.
背叛上帝的表現.
讓我們看到.
其實摩西也有限制.
摩西再欲緊.
再奮身.

$^{761}$再努力.
他也得以接受一個限制.
就是什麼呢?.
他沒有辦法代替百姓.
來馴服上帝.
他不能做完百姓那份.
你有沒有能代你身邊的人.
成為一個好的基督徒?.
沒有的.
你怎樣努力都好.
你只可以勸他.
你不可以為他祈禱.
你不可以做完他那份.
就正如你作為家長.
在獄中.
有時也不可以.
代你的子女生性做人.
說白一點.
他考試不好.
你也不能代他去考.
雖然我知道.
聽過很多新聞.
家長會找人代子女去考試.
代他們去做功課.
但是你可以代多少次?.
終歸人要成長.
就要面對承擔自己的責任.
包括面對他的信仰.
同樣我們無論多慾緊.
多著急都好.
我們只能勸勉.
守望我們的家人.
希望他們信主.
我們不能代他們信耶穌.
我們不能代他們.
這是我們的無奈.
是我們的無能為力.
同時也是摩西的無能為力.
迫使我們如何?.
你會實上帝.

$^{801}$會實上帝.
期望上帝最後幫你包底一切.
所以等姐妹如此看來.
到底什麼支撐著摩西.
她繼續事奉.
面對可能事奉沒有果效.
仍然灰心放棄.
仍然一次又一次努力.
去作出勸勉.
是因為她看到.
上帝對以色列人.
那份不離不棄的愛.
她仍然看到上帝展示出那份憐憫.
那份忍耐和盼望.
事實上摩西知道.
上帝已經預告了.
百姓將來會仍然叛逆上帝.
甚至去到一個地步.
會被嚴厲管教.
趕出應許之地.
被擄到外邦.
但是上帝也留有一手.
當百姓願意謙卑.
回轉重新去聽上帝的吩咐的時候.
上帝又會再一次去施恩憐憫.
將他們從被擄之地焦聚回來.
上帝這份對人的恩慈和期待.
總是給犯錯的人.
一次悔改的機會.
一次又一次的機會.
就成為了摩西繼續事奉的動力.
然後摩西也深深體會到.
事奉的焦點是什麼.
不在於她對百姓的宣講和勸告.
或者不在於她的事奉.
是否立竿見影.
是否即時有美好果效.
而是無論得時不得時.
她都選擇聽上帝吩咐.
我回應的是上帝.

$^{841}$即使百姓聽了她的宣講.
不知道會有什麼反應.
無論將來他們身處什麼光景.
摩西首要的是.
我做好上帝的本分.
我謹守了上帝給我的崗位.
我也確信上帝會為我自己.
為百姓的軟弱限制.
或者甩撈.
上帝通通都會包底.
最終上帝會帶領她的子民.
向著她的旨意來成就.
所以弟兄姊妹摩西.
今天將她這篇信息.
也成為了我們的勉勵.
她想我們選擇生命之道.
她也勉勵我們一定要這樣選擇.
並且活出來.
弟兄姊妹盼望我們在這一刻.
透過禱告.
我們回應上帝今天對我們的提醒.
我們同心禱告.
因主無論我們身處什麼境況.
無論我們正在走哪一條道路.
上帝並意讓我們對準你.
回到你的生命之道.
正如摩西藉著她自己的宣講.
勸勉我們每一個.
唯一要選擇的就是選擇上帝.
唯一要選擇的就是.
遵行上帝的吩咐.
上帝不是一刷的選擇.
而是我們一生要跟隨.
應信的功課.
上帝幫助我們.
我們其實也會像以色列人一樣.
我們有軟弱.
我們會走迷.
甚至得罪你.
我們覺得世界其他的法則.

$^{881}$好像更好.
給我們更多即時的好處.
有時我們會選擇.
但我們不知不覺其實.
遠離上帝.
甚至得罪上帝.
盼望主因你藉著今天這段經文.
重新讓我們愛慕你的話語.
讓我們知道你的話語是可親近.
不難理解,不難遵守.
因為上帝你是那位有恩憐,有憐憫的神.
讓我們養成中讀你的話語.
那一份的操練.
那一份的態度.
也讓我們努力實踐.
特別是做人生選擇的時候.
幫助我們能夠揀選上帝.
揀選生命.
走上帝的道路.
我們獻上禱告.
奉基督耶穌你的名.
來祈求.
阿們.
(影片完結).
\newpage



\section{申命記 31:1-8}
\label{sec:Buo5z3C6Kgc}
\textbf{2025年6月8日崇拜直播｜陳明泉牧師｜朝向應許之地－你們不孤單｜申命記31\_1-8}
\newline
\newline
連結: \href{https://youtube.com/watch?v=Buo5z3C6Kgc}{\texttt{https://youtube.com/watch?v=Buo5z3C6Kgc}} ~~~~ 語音日期: 2025-06-08
\newline
\newline
\hyperref[sec:dxgI8DcWOBk]{\small{< < < PREV SERMON < < <}}
~
\hyperref[sec:index_chronic]{\small{[返順時目]}}
~
\hyperref[sec:f__YU0k5QKA]{\small{> > > NEXT SERMON > > >}}
\newline
\newline
申命記 31:1-8
\newline
\begin{longtable}{cl}
\hline
\hline
章節 & 經文 (和合本修訂版)\\
\hline
31:1 & \begin{tabularx}{0.7\textwidth}{X} 摩西去把這些話吩咐以色列眾人， \end{tabularx} \\ \\ \relax
31:2 & \begin{tabularx}{0.7\textwidth}{X} 對他們說：「我已經一百二十歲了，現在不能照常出入。耶和華曾對我說：『你不得過這約旦河。』 \end{tabularx} \\ \\ \relax
31:3 & \begin{tabularx}{0.7\textwidth}{X} 耶和華－你的神必在你面前過河，把這些國從你面前除滅，你就必得他們的地。約書亞要在你面前過河，是照耶和華所吩咐的。 \end{tabularx} \\ \\ \relax
31:4 & \begin{tabularx}{0.7\textwidth}{X} 耶和華必對待他們，如同從前待他所除滅的亞摩利人的王西宏與噩，以及他們的國一樣。 \end{tabularx} \\ \\ \relax
31:5 & \begin{tabularx}{0.7\textwidth}{X} 耶和華必將他們交在你們面前，你們要照我所吩咐的一切命令待他們。 \end{tabularx} \\ \\ \relax
31:6 & \begin{tabularx}{0.7\textwidth}{X} 你們當剛強壯膽，不要害怕，也不要畏懼他們，因為耶和華－你的神必與你同去；他必不撇下你，也不丟棄你。」 \end{tabularx} \\ \\ \relax
31:7 & \begin{tabularx}{0.7\textwidth}{X} 摩西召了約書亞來，在以色列眾人眼前對他說：「你當剛強壯膽！因為你要和這百姓一同進入耶和華向他們列祖起誓要給他們的地，你也要使他們承受那地為業。 \end{tabularx} \\ \\ \relax
31:8 & \begin{tabularx}{0.7\textwidth}{X} 耶和華必在你前面行，他必親自與你同在，必不撇下你，也不丟棄你。你不要懼怕，也不要驚惶。」 \end{tabularx} \\ \\
[1ex]
\hline
\hline
\end{longtable}
$^{1}$今天的經文是在《新命記》第31章.
我將會讀出第一至第八節.
在《舊約聖經》252頁.
《新命記》31章1至8節.
摩西去告訴以色列眾人說.
「我現在一百二十歲了,不能照常出入」.
耶和華也曾對我說.
「你必不得過這約旦河」.
耶和華「你們的神必引導你們過去.
將這些國民在你們面前滅絕.
你們就得他們的地」.
「約書亞必引導你們過去」.
正如耶和華所說的.
耶和華必待他們.
如同從前待他所滅絕的.
阿摩利,二王,西弘與我.
以及他們的國一樣.
耶和華必將他們交給你們.
你們要照我所吩咐的一切命令待他們.
你們當剛強壯膽.
不要害怕.
也不要畏懼他們.
因為耶和華「你的神和你同去.
他必不撇下你,也不挑起你」.
摩西照了約書亞來.
在以色列眾人眼前對他說.
「你當剛強壯膽.
因為你要和這百姓一同進入耶和華.
向他們列祖.
起誓應許所賜之地.
你們要使他們承受那地為業」.
耶和華必在你面前行.
他必與你同在.
必不撇下你.
必不挑起你.
不要懼怕.
也不要驚怕.
弟兄姊妹平安.
我們今天繼續用《新命記》第31章.
來作一個基本上這個系列裡.

$^{41}$差不多最後兩次的總結.
在我們看這段經文之前.
近來有這樣的經驗.
踏入六月份有很多學校的畢業禮.
早幾個星期前.
飛機去到意大利看神學院的畢業禮.
這幾個星期就會有不同的學校.
教會合作的學校也好.
自己的女兒也畢業了.
一邊看或者參與不同的畢業禮.
我發現有一個很重要的體會.
對於我們這些在旁的參與者.
我看到在台上所說的.
特別是不同學校的師長所說的話.
他們是充滿了很多語重心長的說話.
如果對著一些學生.
在這六年或者很多年的相處裡.
這班師長對學生也有相當的掌握.
通常他們給的那些勸勉.
都是度身訂做的.
有些學生頑皮一點的就笑笑笑.
有很多那些頑皮的片段.
但是如果懂得聽的也知道.
很多學生特別是大學的.
他們會完全的體會和領受.
有時候也會勾起他們的眼淚.
我看到也有神學院的老師.
來分享他們畢業的感言.
是不是有些聲音.
有些歌聲.
畢業感言裡面.
特別對神學院的師長.
對於將要去到事奉的一班牧者.
他們做一個按手猜拳.
和講語重心長的說話.
那裡是很深情的.
對於我們來說.
我們覺得那些都是一般聖經的教導.
但是對於當時的人.
或者領受接受話語的人來說.

$^{81}$那些說話就成為了他們生命裡.
一些很重要要提醒的重要的美德.
你用這個角度來看.
申命記第31章到34章.
其實某程度上就像一個畢業禮的場合一樣.
不過畢業的不是以色列人.
而是摩西他自己.
摩西120歲了.
將要被上帝接走.
他沒辦法和這班以色列人進入加拿美地.
不過在他將要和這班以色列人.
這班百姓離開.
經歷過40年共處的日子.
裡面有喜怒哀樂.
有很激心的.
有很多眼淚的.
去到這一刻臨別依依了.
他一定不會再講一些很表面.
或者很華麗的措辭.
他一定要按著那些以色列人.
生命裡面最需要學習的地方.
來將最後的信息告訴他們.
今天我們看的申命記第31章.
是這個講論裡面一個起始的部分.
這個部分摩西記是針對著以色列人生命的困難.
他們一直以來面對著的盲點.
同時間他都要將那個權力.
帶領的那個權兵.
將這個權兵公開地告訴那班以色列人知道.
也藉此確立約書亞未來帶領的位置.
究竟今天這段經文要講什麼信息呢?.
也對我們今天生命有什麼重要呢?.
希望今天我們短短地去思考.
31章我們會闊一點的看這段經文.
希望從中我們拿到一些養分.
我們先祈禱求主幫助.
(主教).
願你的話語率領我們走人生的道路.
我們恭敬將來我們將每一個心教託給你.
求你的話語進到我們的心靈當中.

$^{121}$指教我們帶領我們.
讓我們不會在偏行幾路的處境裡.
幫助孫講的講得清楚.
願你的話語中心向弟兄姊妹表達出來.
讓我們每一個一同領受一同學習.
在我們生命當中繼續以你成為我們生命的元帥.
恭敬將來的時間教託.
願你賜福.
獻上禱告奉基督耶穌得勝的聖名.
Amen.
我們一起看.
申命記第31章1至8節是第一個段落.
在這個段落裡面.
第一二節那裡.
莫西將他的處境再一次告訴以色列人.
莫西他把這些話吩咐以色列眾人.
對他們說我已經120歲了.
現在不能照常出入.
也和我曾對我說.
你不得過這約旦河.
在這裡其實很清晰的說.
莫西已經120歲了.
除了他在說他自己年老之外.
你不要覺得他因為年老.
他體魄不夠魄力.
不能照常出入.
純粹在說他身體狀態.
莫西他繼續說.
他說不是身體狀態的問題.
因為在聖經裡面記載.
莫西到了120歲.
眼目沒有分花.
身體強健.
他死的時候還要爬山.
上山來被上帝接去.
所以不是魄力的問題.
這是上帝命定的問題.
在他生命當中.
他是以色列人偉大的領袖.
但是他生命也有一些污點.

$^{161}$因為擊打盤石.
就是將上帝的怒氣.
成為了自己的怒氣.
以致他沒辦法進入了迦南美地.
這是上帝命定.
在這個命定下.
莫西他特別的.
他將要回到天父那裡.
他就要將一些很重要.
以色列人面對著一個最大最大最大的問題.
要指出來.
他要告訴他們不要再重蹈覆轍.
在第三節到第八節.
其實裡面的渲染就是.
以色列人最大的盲點是什麼呢.
我讀的版本有一點點不同.
大家聽我讀.
耶和華你的神必在你面前過河.
把這些國從你們面前除滅.
你就必得他們的地.
若孫雅要在你面前過河.
是照耶和華所吩咐的.
耶和華必對待他們.
如同從前待他.
所除滅的阿摩利人的皇,西,弘,與,鄂.
以及他們的國一樣.
耶和華必將他們交在你面前.
你們要照我所吩咐的一切命令待他們.
你們當剛強壯膽不要害怕.
也不要畏懼他們.
因為耶和華你的神必與你同去.
他必不撇下你也不丟棄你.
莫西照料若孫雅來在以色列眾人眼前對他說.
你當剛強壯膽.
因為你要和你者百姓一同進入耶和華.
向他們列祖起誓要給他們的地.
你也要使他們承受那地為業.
耶和華必在你前面行.
他必親自願你同在.
必不撇下你也不丟棄你.

$^{201}$你不要懼怕也不要驚惶.
這裡跟大家讀的和合本.
近似的,相約的.
基本上翻譯的意思是不變的.
不過在這裡重複重複.
當剛強壯膽不要害怕.
我們經常都會引用這一句.
來勉勵人.
不過這個剛強壯膽和害怕.
在那個句式之前.
剛才讀的那段經文.
和合本沒有譯到.
或者他用了其他字眼.
就是耶和華必在你面前.
上帝的面.
在你的面前那句說話.
在原文裡面出現很多次.
就是說.
在說剛強壯膽.
在說那些以色列人不需要害怕.
不要怕.
為什麼呢?.
因為上帝在你面前去行.
上帝也都猜遣了約書亞.
在這班以色列人面前來行.
所以這班以色列人是跟隊.
有開路先鋒上帝.
也都有大元帥約書亞.
在你們面前來行.
有什麼危難.
有什麼上帝幫你擋.
約書亞幫你擋.
這一班以色列人.
他們就要帶著一份無比的勇氣.
來跟著新一代的首領.
進入迦南地.
為什麼說這是最大最大.
以色列人要學習的一個地方呢?.
如果你記得.
或者這段時間.

$^{241}$跟教會的讀經計劃.
在民數記裡面說到.
那些以色列人看到那些外邦人害怕.
看到那些人比自己高頭一截.
即使看到迦南地是流奶軟蜜之地.
那裡肥沃豐富.
有所有很多物產.
很棒的.
不過當他們見到當地的迦南人.
或者在這裡說的阿摩利人.
甚至乎那些其他的迦南外族.
他們害怕.
人的恐懼就佔領人的內心.
這種恐懼感.
就對整個民族進行一個極大的扭曲.
這種扭曲令到那些以色列人說.
我們現在不出埃及了.
上帝害死我們.
現在我們要立首領.
要回到埃及裡面.
又或者這種扭曲恐懼的心.
令到人產生那種權力鬥爭.
這個權力他們要興起新的領袖.
不再跟摩西.
不再跟阿倫.
他們要另立一個合乎他們自己的想法.
或者令到他們舒服的領袖.
令到他們可以安然地.
覺得心裡面有多些踏實感.
來回到舊時的道路.
摩西在這裡裡面.
他在說這個講章.
他第一樣東西.
他直指的就是以色列人.
整個民族的核心.
恐懼,懼怕.
而這種懼怕驅使到人不跟隨上帝.
人生裡面有時候.
我們都可能好像以色列人一樣.
我們有很多事情害怕.

$^{281}$而這種懼怕.
我們看威嚇我們的東西.
比起上帝更加大.
我們害怕那些東西.
令到我們的心靈變得扭曲.
有時候甚至乎這種不是外顯的恐懼.
在內心裡面會影響我們的人際關係.
又或者看自己的那種體會.
早前有一個這樣的弟兄姊妹.
聊天.
一個對話.
他就在一個工作了相當日子的一個機構裡面.
在公司裡面工作.
做著做著的時候.
可能見到其他那些.
因為他做了相當日子.
他身邊有一些下屬.
指指尋尋見到他.
見到這些說話.
這些類似指指尋尋.
他心裡面很不安.
然後在公司裡面.
公司委派一些更加艱巨的任務交給他.
因為他很久了.
所以上面有一些挑戰.
下屬就指指尋尋.
經常看著他在那裡說話.
他心裡面就感覺很不安樂.
但是他又不是很想去澄清.
又不是很想去抗辯.
於是這些說話.
這些人的指指尋尋.
就影響了他的心態.
然後他就不停在質問自己.
我這麼久在這間公司裡面.
為什麼一點份量都沒有呢?.
我做每一樣東西.
為什麼人們會有這麼多的說話呢?.
又或者在這個公司裡面.
為什麼上頭不是很體恤他的困難.

$^{321}$還要不斷不斷地交一些艱巨的任務給他.
然後浮現腦裡面有很多的感受.
然後到最後他就覺得.
這個群體或者這間公司不是屬於他的.
他想走.
然後就自暴自棄.
甚至覺得這間公司只有自己做事.
一邊談的時候.
其實他真正的恐懼是覺得.
別人怎麼看他.
他真正覺得的那種體會是覺得.
自己孤軍作戰.
但是其實可能是客觀的環境.
他還沒有查證過.
究竟這些人指指尋尋.
是不是可能在讚他.
或者這個老闆真的好.
不知道.
不過在我們人生裡面.
有時候我們看到這些.
觀察到這些東西.
我們就會內化了.
甚至那種不安全感.
令到自己失去了那份.
在對面前有些挑戰.
或者有很多面對不同的困難的時候.
我們的心裡面就產生一種孤獨感.
在聖經裡面其實先知都曾經出現過.
不同的人物都經歷過這種心路.
聖經給我們的出路是.
你要剛強壯膽.
如果你越害怕的東西.
你更加要面對它.
有一句這樣的說話.
曾經有一個名人說過.
我覺得這句很受用.
他說面對懼怕最好的方法是什麼.
就是你面對它.
面對到你做到成功為止.
你一直這樣做的時候.

$^{361}$你就會克服到這個懼怕.
這樣說好像很抽象.
這是我親身經歷.
我小時候.
我家人就回大陸去游泳.
那時候年輕的時候去大陸.
沒有水泡的.
不懂得游泳.
我爸爸就把我放到那些叫沖邊.
就是那些小餐廳小河那些鄉下的.
他說沒有水泡.
但舊時鄉下有些用木造的浮木.
就綁了身.
我覺得很安全.
但如果你試過.
你只能夠浮木綁一個身.
和一個水泡綁一個身.
那種體會是不同的.
浮力是不同的.
我爸爸把我放下去.
讓我下去浮在沖邊.
拿著木頭在那裡浮浮浮的時候.
突然之間就有一艘小艇撞.
一直駛過來.
我在沖邊那裡.
很深的.
在沖邊那裡.
那艘艇一下撞到我的頭上.
然後我的人就沉下水了.
這樣滾來滾去.
因為一般來說.
有塊木頭會浮著.
撞一撞就會浮.
但因為我的人被艇撞.
人在艇底那裡.
你能不能想像到那種恐怖.
我一浮的時候.
那艘艇頂著我.
那種感覺就是不停地喝水.
我家人就馬上從河邊.

$^{401}$拉我上來.
把我拉上來.
然後我沒有死.
被淹的那種感覺.
滿口都是水.
濁到很恐怖.
然後有大量的時間.
晚上做噩夢.
就想起這個情景.
有時候偶爾狀態不好.
都會做這個噩夢.
想起小時候的情景.
然後很害怕.
但看到別人.
一直長大的時候.
看到別人游泳.
很自由自在.
覺得很想學.
但又很害怕.
然後到最後是怎樣呢.
那時候我記得.
我找回我爸爸.
爸爸你陪我游泳.
他就想著他綁著我.
有什麼會被淹.
怎料我爸爸自己游泳.
把我放在泳池裡.
然後我想起那個情景.
我年輕的時候.
我還沒有讀聖經.
不過我第一件事.
既然你害怕的事情.
你就克服它.
於是自己浸在水裡.
感受一下被水過頭的感覺.
然後浸過頭之後.
那個人覺得浮.
然後就游了.
前行了.
不怕的.

$^{441}$到地的.
我夠高的.
於是慢慢自學.
就學會游泳.
我要說什麼呢.
我不是說我自己有多厲害.
人生裡有些懼怕的經驗.
如果你懼怕經驗.
或者失敗經驗.
令你沒辦法前行的時候.
你要靠上帝.
在經文裡說到.
上帝走在你面前.
如果上帝叫你做的事情.
你心裡懼怕的時候.
你要知道.
你不是要靠上帝.
更加重要的就是.
面對著你覺得害怕的事情.
你不要避開它.
你面對它.
你越面對它.
做到成功為止的時候.
你的恐懼.
你的喪膽就會克服.
在人世間裡.
未必一定有挑戰.
有時候可能在我們關係裡.
說一句對不起.
又或者.
在我們的生命當中.
我們接受某一些任務.
這些一切.
那種恐懼.
或者那種心裡的不安樂.
我們是需要靠上帝來帶領我們.
在第一個段落裡.
摩西說你不用怕.
上帝走在你面前.
上帝藉著他的僕人約書亞.

$^{481}$走在你面前.
你就看著那些人的腳zoung .
看著上帝的作為.
來做.
上帝叫你去做.
上帝吩咐你去做的事.
他必然成就.
上帝不會代替你.
完成你眼前的任務.
但上帝會允許你.
他會與你同在.
他會為你製造適合的條件.
但最終人要負起.
提升勇氣的責任.
承擔它.
面對它.
你就可以進入迦南美地.
也完成了上帝給你的託付.
這是第一個段落.
所以他對以色列人.
第一件事他要他們克服的.
就是勇氣的問題.
第二個段落.
是由第九節至第十三節.
我再讀一次那段經文給大家聽.
用我的譯本讀給大家聽.
摩西寫下這律法.
交給台耶和華約櫃的利未人祭司.
和以色列的眾長老.
摩西吩咐他們說.
每逢七年的最後一年.
就是定期的豁免年.
在駐彭捷的時候.
當以色列眾人來到.
耶和華理神所揀選的地方.
朝見他的時候.
你要在以色列人.
以色列眾人面前.
念這律法給他們聽.
要召集百姓.

$^{521}$男人,女人,孩子和你成旅居居的.
叫他們都得以聽見.
好學習敬畏耶和華你們的神.
謹守遵行這律法的一切話.
你們的兒女就是那未曾認識的.
也可以聽.
學習敬畏耶和華你們的神.
你們一生的日子.
在你們過躍但何得為業的地上.
都要這樣做.
第二個段落.
同樣都出現面前的字.
不過今次的面前.
就不是在說耶和華在你面前行.
第二個段落.
是在第十一節說.
用一個節期的方式.
上帝的吩咐.
上帝的律法.
要在以色列人面前宣讀.
同樣都用這個work play.
即是玩字眼.
但今次的面前.
是說上帝的話語.
放在以色列人面前.
值得細節去看.
在第九節.
特別強調摩西寫下這律法.
交給祭司和以色列長老.
每逢七年最後一年.
就是豁免年或是欺年.
在駐彭節的時候.
在百姓面前宣讀這些律法.
在這裡.
這幾節.
為整個以色列民族帶來一個.
嶄新的局面.
為什麼這樣說呢.
過往的日子.
在他們在曠野昔日的日子.

$^{561}$那個律法的吩咐是口傳的.
上帝吩咐摩西說什麼.
摩西就copy and paste.
或者加入一點點.
但他還是將上帝的傳話.
傳達給以色列人.
不過在這裡.
就將上帝的話語.
成為了點注.
成為了律法筆錄下來.
記下來.
而且要重複來中讀.
所以由一個口述的傳統.
慢慢變成了一個.
既是口述教導.
同時間也是有一個文字的傳統.
藉著這些文字.
令到整個民族不會民族失憶.
要記得.
還有一個週期.
一個節期性重新來中讀.
放在百姓面前的中讀.
為了什麼呢.
學習敬畏.
重複又重複.
不斷又不斷地去做.
目的就是要訓練一代又一代.
世世代代.
來幫助百姓.
記得起律法.
一代過去下一代也要記得.
在摩西大令以色列人的慘痛經歷裡面.
那四十年的百姓.
那個第一代.
他們經歷過這麼沉痛的經驗.
他不會任由他們過去.
這些沉痛經驗.
就成為了一個很重要的案例.
或者後世要參考的.
一些歷史根基.

$^{601}$上帝怎麼吩咐.
這些人怎麼活躍.
到最後這些人怎麼迴轉.
後來上帝怎麼去帶領.
一代又一代.
來帶領他們.
幫助他們.
以致這一班人.
可以世世代代遵守上帝的吩咐.
這個世代慢慢開始.
變得很輕忽上帝的話語.
今天有時候有些弟兄姊妹.
會將基督教成為了什麼呢.
成為了一個很溫馨.
令人很舒服的一個宗教.
我們講很多愛心.
講很多同在感.
或者有很多溫情.
但是可能信仰裡面.
其實沒有什麼內涵.
沒有什麼硬朗的內容.
弟兄姊妹.
當你真心去讀聖經的時候.
上帝不是純粹安慰你.
有時候你讀聖經.
他會罵你.
如果人生.
我只喜歡找一些我聽得舒服的.
話語來滋潤我的心靈.
我用不著讀聖經.
坊間你去讀那些.
心靈雞湯式的書.
大把是這樣的.
不過真正生命.
對我們生命有重要的道.
他總是既安慰.
同時也是獨澤.
他既是引導.
同時也是滋潤.
是很平衡的一個生命的教導.

$^{641}$我們今天讀聖經.
最終要學的.
就是一種學習敬畏上帝的一種生命態度.
如果我們讀聖經.
只是到最後純粹為了感覺良好.
這個就忽略了.
昔日由摩西.
到今天這個我們整個基督教傳統裡面.
最重要的.
核心的理念.
求主幫助我們.
在我們現在這段時間.
有些弟兄姊妹說.
牧師讀聖經.
摩西五經有些東西很重複.
有些枯燥.
有些悶.
我們都盡量調動一下.
令他不悶.
不過有些東西真的要硬吃.
如果你因為悶而不看聖經.
那你其實.
都是以自我為中心.
來讀經.
我們變成屬靈的偏食.
看一些峰迴路轉.
劇味甚重的經文.
但是對於一些比較表面.
你不仔細看.
是看不到的.
那些內容.
你全部忽略它.
你的信仰.
是會有偏食.
學習敬畏.
在上帝在百姓面前.
去宣讀經文.
這是摩西第二個.
對整個以色列文最重要的渲染.
時間關係我們看第三個段落.

$^{681}$十六節至到第二十二節.
耶和華對摩西說.
看啊.
你必和你的列祖同遂.
這百姓要起來.
在他們所要去的地上.
在那地的人中.
隨從外邦的神明行任.
離棄我.
違背我與他們所立的約.
那時我的怒氣.
必向他們發作.
我必離棄他們.
轉念不顧他們.
以致他們被吞滅.
並有許多禍患.
災難臨到他們.
在那日人必說.
這些禍患臨到我.
豈不是因為我的神.
不在我中間嗎.
在那日因人偏向別神.
所行的一切惡事.
我必定轉念不顧.
現在你們要寫下這首歌.
教導以色列人.
放在他們口中.
使這首歌成為我指責.
以色列人的見證.
我們跳到二十二節.
當日摩西就寫了一首歌.
教導以色列人.
整個十六節到二十二節.
第三個出現的「面」那個字.
就是耶和華掩面不看.
第十六節.
第十八節.
在這裡講到.
那些以色列人.
這是上帝發的寓言.

$^{721}$以色列人進入了迦南地.
進入了安日.
在那裡竟然和外邦神明行淫.
即是離棄上帝.
敬拜那些偶像.
違背上帝與他們所立的約.
所以上帝發怒.
離棄他們.
他會轉念不顧他們.
他會轉念不顧上帝的轉念.
這是一個寓言.
他講到那些以色列人.
進入安日的時候.
已經好像完成了重要的任務.
然後他們就會離棄上帝.
你值得留意的.
那個轉念的次序.
不是上帝一進到迦南地.
就轉念不顧他們.
而是在講那些百姓.
因為安日.
享受.
以致他們和那些神明.
外邦神明行淫.
第二步.
第三步就是離棄上帝.
所以先以色列人離棄上帝.
然後惹動上帝的憤怒.
以致上帝離棄他們.
轉念不顧他們.
所以上帝的轉念.
是基於這一班以色列民.
上帝寓言.
他們將會敗壞.
所以他們就.
會經歷不到.
上帝的同在.
究竟上帝的不同在.
是否任由他們.
讓他們自生自滅.

$^{761}$上帝的離棄代表著什麼呢?.
我們看回.
如果你留意看到歷史.
到今天以色列人都還在.
在這個世界裡存在.
我特別翻查.
在他們以色列人將要進入.
迦南地的最重要.
最大的民族.
或者幫國.
亞馬力人.
這個亞馬力人在歷史裡消失了.
那個舊時最強悍的外邦.
非利士人今天都消失了.
昔日以色列人的宿敵.
罪惡的外邦.
在歷史裡淹沒了.
相反來說.
這些以色列人經常都跳來跳去.
又敬拜神.
又參拜外邦神明.
這樣遊走之下.
上帝說要離棄他們.
但最終.
以色列這個民族.
或者聖經裡所啟示的這群人.
他們沒有在歷史裡淹沒下來.
關鍵是原來上帝的離棄目的.
離棄不是最終的終點.
離棄其實是在第十七節.
當那些人經歷被上帝離棄的時候.
那些人在那天裡說.
那些禍患淋到我.
豈不是因為我的神不在我中間.
他們在上帝的離棄之下.
就發現原來沒有上帝那麼可憐.
這種離棄是令他們勾起這群人的反思.
為什麼他們會經歷這些禍患.
本來是很幸福的日子.
因為上帝的不同在.

$^{801}$他們就經歷禍患.
而這些禍患就迫使他們.
重新去想想他們和神之間的關係.
所以上帝的離棄就好像.
我們有時候帶小孩的爸爸媽媽.
你這麼頑皮 我不要你了.
我們說 語言暴力.
上帝經常說這些語言暴力.
不過真正的不是上帝真的不要他們.
而是說你不願意去聽.
那你經歷一下你自己的後果吧.
我不要你的意思是代表著.
我就是因為痛心而來.
你想經歷一下你自己覺得好的日子.
讓你去經驗一下.
承擔一下你自己的後果.
著急一下.
原來過去裡面一直活在恩典當中.
所以在聖經裡面.
希伯來書有時候會用另一個字眼.
說的就是管教.
上帝的離棄就等於上帝對人的管教.
如果去到極致.
羅馬書裡面說.
如果上帝真的放棄的話.
他任憑你.
但上帝對以色列人.
或者對我們每一個信神的人.
上帝說的離棄.
只不過是一個暫時的管教.
暫時的管教你會有眼淚.
你會經歷困難.
但到最後你回到上帝的時候.
你就會經歷一份悔改的生命.
你會重新走回生命的正軌當中.
所以弟兄姊妹.
上帝藉著摩西說一個這樣的預言.
他叫摩西說這個預言的時候.
他要將這些內容成為一首歌.
摩西之歌成為《親命記》第32章裡面.

$^{841}$一個主要的內容.
這個摩西之歌.
一般來說摩西之歌.
或者聽聖經裡面的歌.
是歌頌上帝.
或者是說人多麼厲害.
以色列如何打勝仗.
不過這首歌.
世世代代唱的.
就是說到人們會失敗.
這些人會因為順境.
就離棄上帝.
藉著這首歌.
上帝說.
就好像成為這班人的見證一樣.
你世世代代唱.
你就比對一下這首.
在入迦南地之前已經作好的這首歌.
比對一下這個生命.
看看是不是你們正在走.
在一個這樣離棄上帝的軌跡裡面.
這首摩西之歌.
就指導甚至指出這些以色列人的不是.
今天這段經文不是很複雜.
整段經文用上帝的面.
或者用面的字.
來玩語帶相關.
第一節八節說到上帝的面.
叫這班以色列人.
克服恐懼.
勇往直前.
進入迦南地.
第二個面.
用的就是說.
上帝的吩咐.
要成為典章.
成為律法.
世世代代來作教導.
在這些眾人面前去宣讀上帝的律法.
就是為了要學習敬畏.

$^{881}$第三個.
人的離棄.
會令到上帝轉念.
人面對著上帝轉念.
而面對的惡果.
你不用驚奇.
不過在這些困難當中.
重新思考.
反思一下.
我們的生命是不是離上帝甚遠.
以至我們正在經歷眼前的困難.
今天這段經文對我有什麼體會.
或者對大家有什麼體會呢.
我覺得今天.
我們每個人都有為自己的生命目標去努力.
不過我們很少問.
上帝在我們的群體.
這個集體裡面有什麼群體的目標.
曾經有些人說.
教會的目的.
是要幫助每一個獨立的弟妹找到個人目標.
我不否定.
我覺得這個是重要的.
不過除了教會幫助你找到你自己人生個人目標之外.
我們這一班人的聚集.
我們會不會都有一個上帝給我們群體的目標.
如果我們說我們人生裡面要進入.
旺角浸信會的迦南地.
下一步要怎樣做.
我們有沒有這樣的急事.
因為有大量的恐懼.
沒有錢 經濟不好.
我家裡又沒搞定.
我自己又千瘡百孔.
這些很多的事情.
到最後我們就繼續停滯不前.
可以這樣繼續蹉跎歲月.
不過上帝如果將我們每一個人招聚成為.
旺角浸信會這個群體.
我想我們不只是純粹單純來到這裡.

$^{921}$就是滿足我們.
要教會滿足我個人.
或者建立我個人自己的志業.
除了個人之外.
群體繼續怎樣去走.
我們必須要虛心求上帝去帶領.
我想起.
想這段經文的時候.
我想起一個片段.
大家記得我們的執事會主席阿叔.
我經常都說.
不過每一次想起這些情景.
我都會想起他.
我記得他臨終的時候.
我去到醫院去探望他.
我捉住他的手.
阿叔跟我說.
他說教會會如是.
旺角浸會將會有一個很大的發展.
他最三的叮囑除了發展之外.
他說要溝通多一點.
這個阿叔經常都說.
我們的執事溝通多一點.
努力去溝通.
建立好弟姐妹之間的關係.
那些是我們的弟弟妹妹.
是我家人.
教會向前.
我想阿叔說的這番話.
不是一些很隔去的騷擾.
雲雲美德裡面.
要拿一個出來說一下.
打發我這個牧師.
這個新手的主任牧師.
而是他知道.
我們這個群體一定要學習.
在無論什麼處境裡面.
我們要坦誠.
我們要溝通.
我們要建立深度的關係.

$^{961}$而這份關係.
不是叫我們只在這裡輕輕按摩.
這份關係是說.
如是教會會有大發展.
我們要繼續來到勇往直前.
在發展的過程.
可能有很多溝通.
有很多誤解.
但是如果沒有溝通.
大家各自各做.
vision就變了division.
但是如果大家互相溝通.
大家一起朝著某一個應許之地.
奮發向前.
那個就是未來.
黃國浸信會要走的路.
弟姐妹.
當你讀這個申命記的時候.
希望我們每一個都.
拿到這班人.
他們將要前進的方向.
同樣我們讀申命記.
我們都要找到黃國浸信會.
要前進的方向.
上帝在面前等我們.
你跟嗎?.
\newpage



\section{約書亞記 1:1-9}
\label{sec:f__YU0k5QKA}
\textbf{2025年6月15日父親節主日崇拜直播｜劉佩婷博士｜50歲的剛強壯膽｜約書亞記1\_1-9}
\newline
\newline
連結: \href{https://youtube.com/watch?v=f-_YU0k5QKA}{\texttt{https://youtube.com/watch?v=f-\_YU0k5QKA}} ~~~~ 語音日期: 2025-06-15
\newline
\newline
\hyperref[sec:Buo5z3C6Kgc]{\small{< < < PREV SERMON < < <}}
~
\hyperref[sec:index_chronic]{\small{[返順時目]}}
~
\hyperref[sec:EfFf_ZeLMAM]{\small{> > > NEXT SERMON > > >}}
\newline
\newline
約書亞記 1:1-9
\newline
\begin{longtable}{cl}
\hline
\hline
章節 & 經文 (和合本修訂版)\\
\hline
1:1 & \begin{tabularx}{0.7\textwidth}{X} 耶和華的僕人摩西死了以後，耶和華對摩西的助手嫩的兒子約書亞說： \end{tabularx} \\ \\ \relax
1:2 & \begin{tabularx}{0.7\textwidth}{X} 「我的僕人摩西死了。現在你要起來，和眾百姓過這約旦河，往我所要賜給以色列人的地去。 \end{tabularx} \\ \\ \relax
1:3 & \begin{tabularx}{0.7\textwidth}{X} 凡你們腳掌所踏之地，我都照我所應許摩西的話賜給你們了。 \end{tabularx} \\ \\ \relax
1:4 & \begin{tabularx}{0.7\textwidth}{X} 從曠野和這黎巴嫩，直到大河，就是幼發拉底河，赫人的全地，又到大海日落的方向，都要作你們的疆土。 \end{tabularx} \\ \\ \relax
1:5 & \begin{tabularx}{0.7\textwidth}{X} 你一生的日子，必無人能在你面前站立得住。我怎樣與摩西同在，也必照樣與你同在；我必不撇下你，也不丟棄你。 \end{tabularx} \\ \\ \relax
1:6 & \begin{tabularx}{0.7\textwidth}{X} 你當剛強壯膽，因為你必使這百姓承受那地為業，就是我向他們列祖起誓要給他們的地。 \end{tabularx} \\ \\ \relax
1:7 & \begin{tabularx}{0.7\textwidth}{X} 只要剛強，大大壯膽，謹守遵行我僕人摩西所吩咐你的一切律法，不可偏離左右，使你無論往哪裡去，都可以順利。 \end{tabularx} \\ \\ \relax
1:8 & \begin{tabularx}{0.7\textwidth}{X} 這律法書不可離開你的口，總要晝夜思想，好使你謹守遵行這書上所寫的一切話。如此，你的道路就可以亨通，凡事順利。 \end{tabularx} \\ \\ \relax
1:9 & \begin{tabularx}{0.7\textwidth}{X} 我豈沒有吩咐你嗎？你當剛強壯膽，不要懼怕，也不要驚惶，因為你無論往哪裡去，耶和華你的神必與你同在。」 \end{tabularx} \\ \\
[1ex]
\hline
\hline
\end{longtable}
$^{1}$今日是父親節.
我在此恭祝在座所有父親.
父親節快樂.
公公爺爺有雙份的快樂.
今早重讀的經文是《若書雅記》一章一至九節.
經文記載在《聖經》舊約部分263頁.
《若書雅記》一章一至九節.
耶和華的僕人摩西死了以後.
耶和華曉遇摩西的幫手.
亂的兒子若書雅說.
我的僕人摩西死了.
現在你要起來.
和眾百姓過這約旦河.
往我所要次級以色列人的地去.
凡你們腳掌所踏之地.
我都照著我所應許摩西的話.
次級你們了.
從曠野和者黎巴嫩.
直到幽法拉底大河.
客人的傳遞.
又到大海日落之處.
都要作你們的警戒.
你平生的日子.
必無一人能在你面前站立得住.
我怎樣與摩西同在.
也必照樣與你同在.
我必不撇下你也不丟棄你.
你當剛強壯膽.
因為你必使這百姓承受那地遺孽.
就是我向他們列祖.
豈世應許次級他們的地.
只要剛強大大壯膽.
謹守遵行我僕人摩西所吩咐你的一切律法.
不可偏離左右.
使你無論往哪裡去.
都可以順利.
這律法書不可離開你的口.
總以晝夜思想.
好施你謹守遵行.
這書上所寫的一切話.

$^{41}$如此你的道路就可以亨通.
凡事順利.
我豈沒有吩咐你嗎?.
你當剛強壯膽.
不要懼怕.
也不要驚慌.
因為你無論往哪裡去.
耶和華你的神必與你同在.
聖經獨白.
各位弟兄姊妹早晨.
讓我今天一個很特別的主人.
可以跟大家分享.
父親節.
無論你是父親還是不是父親.
或者不會是父親.
怎樣也好.
祝大家今天可以和你的家人.
快快樂樂地度過.
今天我選的題目.
《若書雅記》.
《五十歲的剛強壯膽》.
我想看回剖開.
先說說題目.
剛強壯膽四個字.
剛才弟兄幫你讀出.
重複了三次.
在短短第一章的一至九節裡.
其實在第一章還有第四次.
剛強壯膽.
不過那是以色列人對若書雅說.
而上帝親自對若書雅說.
就是三次.
剛強壯膽.
只要剛強大大壯膽.
你要剛強壯膽.
不要懼怕.
為甚麼要說三次這麼多呢?.
為甚麼上帝要對若書雅說這麼重要.
這麼短的文章裡面出現了三次.
我猜這應該是中心主題.

$^{81}$若書雅是否很弱.
很懦弱.
很怕事.
很膽小.
所以要說三次這麼多.
我們待會看看是否這樣.
為甚麼又要叫五十歲呢?.
我引用李思敬博士.
他在他的《舊約識經講道》系列裡面.
提及到關於若書雅的.
他說若書雅寫在.
若書雅當時踏入人生最後的三分一.
當時摩西死了.
一個新的時代.
新的階段的開始.
根據聖經記載.
當時摩西死的時候.
若書雅剛開始的時候.
若書雅大約是八十歲.
用今天的計算方法.
參考李思敬博士.
他說若書雅寫給.
準備進入人生最後階段的若書雅.
所以如果用今天的數字去計算.
我們就叫它做人生的中期.
當時的八十歲.
計算今天大約是五十歲.
六十歲這個階段.
一個階段.
視乎人的壽命.
現在我們香港人的壽命長了很多.
大約在這個期間.
簡單來說就是.
若書雅記載在人生中期的若書雅.
人生中期.
我將若書雅的人生大約分成三部分.
中間就是我們今天要說的若書雅的位置.
聖經記載她有110歲.
大約將它分成三個階段.
這個階段究竟發生什麼事呢?.

$^{121}$我們一想到人生最後三分之一.
這一句話我們就會有很多聯想.
我們會想起坐輪椅.
睡在床上.
要靠人照顧.
沒什麼用.
或者好一點積極一點的.
功成身退.
或者是都可以享受一下.
人生下半場.
通常我們都會這樣想的.
我們看看若書雅的人生.
她在頭四十年其實是一個奴隸.
當時她在埃及受很多艱苦.
她在這四十年裡.
跟著摩西出紅海的時候.
她親眼見過超自然的實在.
而紅海也是在她面前分開的.
這個神蹟.
第二個四十年就很精彩了.
在聖經裡若書雅第一次出現.
其實就是被摩西吩咐她.
叫她預備找一班人去打仗.
當時她剛剛離開埃及沒多久.
就在曠野.
就阿瑪利人要來打以色列人.
摩西就叫若書雅.
為什麼叫她呢?.
因為她是以法蓮之派的代表.
就叫她找一班人.
其實是志願軍來的.
記住他們一班人都是奴隸.
不會接受過任何軍事訓練.
也沒有任何的裝備.
只是臨時拉夫.
準備之下就去.
所以她只是一班志願軍.
就去打仗.
當然在聖經裡記載.
我們都知道若書雅很有名的.

$^{161}$她是十二探子之一.
去窺探迦南地.
在很多的色經學者說.
若書雅是一個很有勇猛.
很有信心的將士.
一個兵士.
既然她的人生都這麼豐富.
這麼精彩.
也是摩西的一個重要的幫手.
為什麼上帝還要叫她剛強壯膽呢?.
是不是若書雅有些事情很害怕.
很擔心.
很恐懼.
我們不知道的呢?.
她那些是外表的勇猛.
還是裡面還有些什麼呢?.
我今天就很大膽.
就用我自己的本行.
用心理學去猜測.
究竟一個人在人生中期.
他會遇到什麼機遇和挑戰.
用今天的圖像去理解.
如果我將二十年的人生.
為一個階段的話.
那介乎四十至六十歲.
就是一個人生的中期.
人生中期.
如果一班中年人坐在一起.
最多說的話是什麼呢?.
就是這句.
你退休了嗎?.
你什麼時候退休?.
你退休之後有什麼做?.
或者你什麼時候退休?.
有什麼打算?.
這個不只是香港的.
其實在外國也一樣.
人生中期.
很常見的一個社會話題.
或者社交話題.

$^{201}$就是退休.
可能很多人都在想.
我甚至都被邀請去講一些講座.
教弟兄姊妹決定退不退休.
其實我都不太懂得做這個講座.
這個真是一個很不容易的題目.
其實退休是一個法律上的定義.
究竟我們人的心態.
是另一回事.
退休之後有什麼好做呢?.
旅行?.
照顧孫子?.
其實我試過問一些弟兄姊妹.
退休後有什麼打算.
或者有什麼人生意義目標.
其實很多都回答不了.
很多都講不了.
很多都不知道.
接下來的時間是怎樣過.
我試一下從一個傳統的心理學角度.
去認識一下人生的中期.
其實這個人生階段是很多人.
由以前到現在.
都是很有興趣去研究的.
人生中期用這個角度去看.
就是剛剛人生的巔峰.
出名的心理學家容格·卡爾·雍.
他說人生的中期.
就是一個afternoon.
午後的時光.
中午之後的時光.
他的意思就是叫人好好享受.
剩餘的陽光和溫暖.
也有一些心理學家說.
這是人生的活力期.
是前所未有的那麼多.
那麼好的機會.
因為人生的上半場.
累積了很多豐富的經驗.
可以承擔一些獨特的角色和責任.

$^{241}$可以為下一代做貢獻.
這是充滿愛心的表現.
也有人說.
人生好像一場球賽.
上半場就累積經驗,財富.
人際關係網絡.
下半場先來個中場休息.
調整策略.
然後繼續在人生下半場.
追求有意義的人生.
人生的中場.
如果我們用大約這個階段.
四十多歲,五十多歲.
這些階段.
有些人會趁中場做些什麼呢?.
如果大家打球.
都知道一個球賽.
你中場休息.
回到休息室.
你不會罵自己.
真是蠢.
錯失了一個球.
為什麼接不到那個球.
為什麼打得這麼差.
你會不會用中場休息去罵自己呢?.
你不會的.
因為你會快點.
有什麼打得不好.
快點補充體力.
快點定下來.
快點再去想想.
之前做得不好.
怎樣調整策略.
以至爭取下半場.
作為一個難得的機會.
繼續拼勁下去.
繼續把握機會.
然後追求一些有意義.
有價值的生活.
不過很可惜.

$^{281}$不是每個人都有能力去做到.
今天說到人生的中期.
最常見一個很主流的觀點.
就是中年危機.
這個差不多是.
每個人說到中年都會問的.
是不是真的有中年危機呢?.
在學術上.
不是每個人都會經歷中年危機.
不過在生活上.
的確有很多事情.
在中年階段的發生.
引發起我們想一些很重要的問題.
那些引發的事件.
包括病患.
無論是身邊的家人還是自己.
或者突然之間過世.
或者是突然之間失去工作.
突然之間失去一些很重要的關係等等.
於是就會問.
人生價值.
我還在做什麼?.
形形異異.
幾十年現三十年.
我還在做什麼?.
人生價值如何?.
我是誰?.
在傳統的心理學上.
女性的中年危機.
一定會提及空巢期.
因為當時的主流觀點.
大部分的女性都會是媽媽.
子女長大之後.
空巢期就很影響到女性的精神健康.
簡單來說.
中年危機就是被人去.
有一個所謂危機.
就是突然之間發現.
人生很短暫.
原來不好好把握機會.

$^{321}$就會錯過很多難得的機會.
在心理輔導角度.
我們常見一些個案.
多數都是這樣.
在某某年突然之間做一些很重大的決定.
我認識幾個家庭.
男士突然之間察覺.
可能覺得自己體力開始弱.
有大肚腩.
有脫髮.
突然發現自己不再年輕.
人生原來沒有什麼機會.
於是毅然買了一層樓.
要做老闆.
追求他年輕時候的夢想.
當年他旁邊的太太發狂.
於是引發很大的家庭危機.
亦有些.
今天其實無論男性女性都會.
有很多在輔導裡面的個案.
覺得如果我不再為人生做一些決定.
我不想下半生再在糟糕的婚姻裡面被折磨.
於是毅然結束婚姻.
結束一些關係.
中年裡面最大的標誌就是遺憾.
很多事情覺得自己要做的沒有做到.
當然做錯的事情亦不少.
於是很盡力去做很多補償.
所以中年一個很大的標誌.
很多時候都會為自己想.
覺得以前只是為他人而活.
開始為自己想.
面對疫情之後.
2020年之後.
心理學家提出一件事.
因為結婚的人少.
生子女的人又少.
中年已經不是像我們剛才傳統所說.
未必有這麼多人有空巢.
未必有這麼多人有經歷這些家庭關係的轉變.

$^{361}$或者我們未來的時代.
就是中年時間長了.
因為人的壽命長了.
即是這段中年期.
可能隨時比想像中的20年更長.
於是有心理學家提出.
這可能是人生的關鍵時期.
關鍵時期的意思就是.
在這個階段可能要為人生做一些很重要的決定.
這個決定在於.
要平衡和顧及很多兩樣的處境.
中年其實是處於世代之間.
一個人兼了很多角色.
可能是父母.
亦都是別人的子女.
是別人的上司.
亦都是別人的下屬.
不知不覺之間.
可能在行業裡面已經是你的長輩.
已經是你資深的角色.
你已經可能要被要求去擔當很多的職位.
除了你本身的工作.
工會,教會.
甚至可能是你樓下.
你作大廈的管理委員會.
總之很多地方都因為你的資歷心,經驗心.
很想你幫忙做很多事.
責任亦都前所未有那麼多.
人的壽命長了.
自己的父母,配偶的父母.
甚至你好朋友的父母,家人.
照顧責任大了很多.
我認識一個人.
他童年有些不幸.
他用剛才的說法.
頭二十年.
他經歷父母雙親都離開.
因病,因意外都過世.
所以他很年輕.
跟他的弟弟兩姐弟.

$^{401}$相依為命長大.
他當然很重要.
他信了耶穌回教會.
所以當他第二個二十年.
即二十至四十歲.
他就用他的生命去服侍很多有困難的家庭.
甚至到海外去幫助貧困的小朋友.
他被呼召回香港讀神學.
他想著上帝召他去牧會.
或者去教會服侍.
誰知原來不是.
在一個機會之下.
讓他在人生的中年去做縣務.
做縣務來說.
他年紀比較資深.
到中年才去做縣務.
雖然已經是中年.
不過在縣務裡面.
他是一個新丁,新人.
去做縣務的時候.
我都認識這個人.
我跟他聊天的時候.
我就說.
其實你有沒有留意.
你在中年裡面得到的機會.
不是偶然.
是上帝在你的生命裡面.
讓很多事情發生,經歷.
以致你的人生變得圓滿或者完整.
不是經過以前.
經過很大的失去.
失去父母.
那種很大的哀傷.
以致今天你不會承擔到一個.
這麼重要,很稱職的使命.
縣務.
去到人生的中期.
一個關鍵的時期.
我們可能要有平衡很多的得與失.
當然亦都會面對很多的挑戰.

$^{441}$身體的確已經和以前不同.
白頭髮,老花,肚腩.
這些的確是.
越來越不像年輕時.
做得這麼快.
記憶力這麼強.
但是我們依然.
都是要為自己的人生做一些決定.
去到中年時期.
所追求的個人成長.
是什麼成長呢?.
我引用美國一位學者.
他說的一句話.
我將它翻譯成中文.
個人成長不是學習新的知識.
而是要忘記舊有的限制.
有一位大學教授.
他在他的行業裡已經很資深.
很多人因為他的名字.
跟他做研究.
跟他寫論文.
跟他學習.
他都退下來.
因為他已經夠年紀要退休.
記者訪問他的時候.
已經接近70歲.
但他依然都去外國上課讀書.
那班記者去訪問他.
問他教授.
你都是一個很資深的學者.
都很有影響力.
在那個專業裡面.
已經是一個資深的學者.
你為什麼還要去外國上課呢?.
這位退休教授說.
不是的.
他坐下來上課.
教授的老師.
當然年輕很多.
年輕過他很多.

$^{481}$他說老師說的話.
我從來沒有想過.
這位退休教授說.
我沒有想過這個題目.
原來可以這樣想.
原來可以這樣去學.
或者這樣去教同學.
原來同學是明白的.
這位退休教授.
他說的話.
深深影響記者.
因為沒有想過.
他已經這麼資深.
他應該去教人.
應該很多人跟著他.
但原來不是.
他去學.
學什麼不是知識.
原來他去看到.
另一個景象.
很多人以為人到中年.
快點去上課.
我教輔導的.
不少人中年來讀輔導.
都是想幫自己.
或者想盡心學一些東西.
這當然是好事.
我是輔導的老師.
我都很鼓勵.
不過去到中年.
我們的個人成長.
不是要讀多個證書.
上多個工作坊.
不是要學一些自己不懂的東西.
或者沒有讀過的東西.
不是靠背誦.
不是靠記憶.
不是靠考試.
中年的個人成長.
是要忘記以前.

$^{521}$你的想法方式.
你的做事方式.
英文我們叫unlearning.
這是非常艱難.
去到中年.
要放下.
很難.
過去人生的經歷.
組成了今天的自己.
要去忘記.
要去放下.
有很大的掙扎.
可能因為這樣.
上帝給我們一個教導.
在二創書.
《耶和華》.
講了一句話.
不要記念從前的事.
不要思想古時的事.
我要做一件新的事情.
當我想到中年的時候.
我常常都去想一個課題.
就是記不記得以前的事好呢?.
經常說放下放下.
記不記得好呢?.
記住和忘記.
又如何去學習呢?.
研究二塞那書的學者.
他說記憶可以帶來希望.
例如以色列人.
記住上帝當日如何.
讓紅海分開拯救他們.
記住那個因為很重要.
但是記憶可以取代希望.
當我們什麼都不做.
坐在那裡等.
等上帝這次又會叫紅海分開.
我等.
我坐在那裡等.
等我小時候的奇蹟再次出現.

$^{561}$我等.
我等上帝以前呼召我做信徒的時候.
他在我家走的奇蹟行事.
要再發生.
我就坐在那裡等.
但是記憶可以取代希望.
所以這位色經學者.
他說故意忘記過去.
可能就是要將所有注意力.
集中在今天上帝要做新的事情上.
今天上帝要做新的事情.
讓我們要忘記過去.
因為今天上帝要做新的事情.
不是我們想到的事情.
因為我們不懂.
我們也沒有經歷過.
這位是德國的一位社會心理學家.
費諾姆.
他寫了一本書叫做.
To Have or To Be.
中文翻譯做.
以佔有為導向的人生.
以存在為導向的人生.
我們從小到大生活在一個社會裡.
將人變成一個物件.
要有某些表現,功能,標準.
人成功的標準可能是.
有房子,有子女,有太太,有事業.
我記得有一位求助者也說過.
他什麼都有.
還有很多憂慮.
人到中年就是這樣.
但這是社會一直給予我們的標準.
量度,衡量一個人的標準.
於是到了這些階段就想.
究竟我是誰呢?.
究竟我活到今天的意義在哪裡呢?.
我們回看約書亞.
約書亞在很多聖經學者的研究裡.
說他是一個很勇猛,很成功的人.

$^{601}$我同意.
但我覺得他最成功的位置在哪裡呢?.
我覺得他最成功的.
就是在他人生裡.
做了至少三個很重要的決定.
第一.
剛才我提到.
他在以色列人離開埃及到曠野之後.
打的第一場仗.
打亞馬利人.
如果根據聖經的記載.
當時摩西叫他.
你找一群人去打仗.
然後聖經說約書亞就去做.
這是他很重要的決定.
摩西吩咐他就去做.
就是因為他這個決定.
帶來一個很重要的影響.
就經歷了上帝的保護和能力.
第二,他人生裡很重要的決定.
就是在民數記裡記載.
剛才也提及過.
他是其中一個探子去窺探迦南地.
他做了一個決定.
就是向以色列人說.
那塊地是極美之地.
是不是?.
可能是上帝賜給我們的.
不要背叛耶和華.
不要害怕.
耶和華會與我們同在.
這是他人生裡一個很重要的決定.
他勇於說他看到的東西.
他相信的東西.
這是他人生裡一個很重要的決定.
如果大家看聖經.
我們都知道.
十二探子.
是一個我們常常讀聖經.
都會很熟悉的.

$^{641}$一個很重要,很關鍵的決定.
第三個決定.
我覺得是約書亞人生裡.
做一個很關鍵的決定.
就是在約書亞記開始的時候.
他回應上帝的呼召.
就剛強壯膽.
繼續承接摩西的工作.
帶領以色列人進入迦南地.
他這個決定.
引來什麼重要的影響?.
就是他親眼看到.
上帝如何帶領以色列人過約旦河.
過程中.
他面對上帝如何使用一個忌諱.
在當中令整件事發生.
第二就是.
他剛強壯膽地做了這個決定.
跟從耶和華之後.
在他人生最後的三分一.
他看到上帝如何將耶利哥城.
交給以色列人.
人生的成功.
不是我們擁有什麼.
人生的成功.
而是在乎你生命中做了什麼重要的決定.
而這些決定不是帶來財富.
而是帶來下一代的福利.
我們說到下一代.
會想到一個英文字.
Legacy.
Legacy這個字是解為遺產.
或者我們用一個比較簡單的說法.
當你人不在的時候.
你留下了什麼影響.
未必是給你的子女.
而是給下一代.
社會上也好.
家庭也好.
教會也好.

$^{681}$一個群體也好.
Legacy這個字.
遺產.
或者我們說遺留下來的影響.
究竟是什麼呢?.
我再說多一點.
遺產或者Legacy這個字.
本身是很男性的.
女性我們很少用Legacy.
因為女性和下一代的關係.
很多時候都透過生殖能力.
所以那種情感關係.
很多時候已經存在.
所以遺產遺留下來的影響.
這個字本身是很男性的.
約書亞遺留下來給後代的影響.
在哪裡看到呢?.
約書亞記最後二十四章.
已屬能詳的二十四章十四節.
「至於我和我家.
我們必定事奉耶和華」.
約書亞在約書亞記.
耶和華親自跟他說了三次.
「當剛強壯膽」.
害怕嗎?.
當然害怕.
「要再一次重新立志信服上主」.
無論上主帶你去哪裡.
無論要你付任何代價.
你都決意跟從.
當然害怕.
預備我便害怕到死.
都中年了.
讓我放假也好.
人生上半場.
我信耶穌信了四五十年.
什麼崗位都做過.
執事會主席都做過.
唱詩教主也做過.
什麼都做過.

$^{721}$現在準備退休.
讓我放假也好.
讓我做其他事也好.
當然想.
人生中期.
心理學家的研究.
有時說得很對.
掙扎位在哪裡呢?.
就是自我專注.
你再想幫自己想多一點.
還是其他專注.
你幫別人想多一點.
別人不只是子女.
你可能沒有子女.
你可能沒有其他很重要的人.
但其他可能是一些.
你影響力可以做到的人.
那個不是你有什麼本事.
而是你的信心會不會影響到身邊的人.
李子敬博士說約書亞記.
他說一點.
我是非常同意的.
約書亞的遺產.
遺留下來.
對以色列人的影響.
是什麼呢?.
就是他對上帝的信心.
我相信一群跟隨他的人.
不是看見他很健碩.
很厲害.
打人也死.
射槍也中.
不是那些.
而是他的決意.
決心.
和那種信心.
我們有沒有想過.
其實我們跟從上帝的人.
我們從來都不是獨行俠.
其實在我們的生命裡.

$^{761}$我們一直都是跟上帝一起.
上帝都是跟我們一起.
在無論人生什麼經歷.
我們看回之前這個.
可能在我們人生上半場.
可能有很多不同的經歷.
但這些經歷不是目的.
頭二十年.
頭四十年.
這些經驗.
無論多成功.
多燦爛.
多光輝.
或多坎坷.
多不幸.
都是過程.
不是你人生的目標.
而我們必須經過這些過程.
才能為下一個階段作準備.
我用這個圖.
約書亞頭四十年是奴隸.
中間四十年.
在曠野流浪四十年.
他這些經歷.
就造就了他.
人生最後的三分一.
可以見到上帝.
那個很真實的同在.
所以在聖經裡記載.
約書亞的出現.
經常帶著一個標誌.
就是他是摩西的幫手.
在聖經出現了很多次.
包括約書亞記第一章.
但在約書亞記最後.
在他一百一十歲.
預備過生.
人生最後階段的時候.
聖經記載.
上主的僕人約書亞.

$^{801}$沒有以前.
怎會有現在呢?.
沒有以前.
怎會有現在.
那個完美完整的人生呢?.
其實如果用這個圖表.
你看摩西的一生.
摩西頭四十年在埃及.
王子有很高級的教育.
中間四十年在米田.
學了什麼?.
學了如何在曠野生活.
然後在人生最後的階段.
最後的四十年.
去做以色列人的領袖.
在曠野生活.
沒有以前.
怎會有現在?.
人生的中期.
今天的標題五十歲.
只是代表著人生的中期.
我們沒有人知道.
我們是不是已經在人生最後三分一.
我們沒有人知道明天會如何.
今天可能就是我們人生的關鍵時期.
為我們的人生做決定.
再重新立志信服上主.
會不會今天.
其實上帝有很特別的想法.
我們不知道.
上帝的意念.
永遠都高於我們的意念.
所以在約書亞記.
最後的時候.
約書亞記已經年紀大的時候.
他向以色列人說.
大大壯膽.
各位.
今天我們看約書亞記第一章.
當剛強壯膽.

$^{841}$這是上帝給我們機會.
再一次為我們自己的人生做決定.
不是要去補償過去.
而是這個決定.
留下給下一代一些影響.
今天我講這個題目.
legacy留下給下一代的影響.
送給每一位做父親的.
當我們要去問一些作為子女.
你記得你爸爸的什麼呢?.
很多時候很多人的答案都會是.
他經常都不在家.
他很少說話.
我也不知道他在想什麼.
爸爸在傳統下來.
遺留下來.
給子女的影響力.
其實是什麼呢?.
有些訪問過一些朋友.
一些弟兄姊妹.
他們說爸爸給他們影響.
就是不要隨便信人.
這個世界很危險.
很多威脅.
或者是.
總之要多想自己.
等等這些.
今天我們做父親的.
我們想留下什麼給下一代.
不是財產.
一想到這個就很害怕.
當我們很害怕的時候.
就背若書雅記第一章.
剛強壯膽.
怕不是問題.
而是我們背若書雅記.
然後就好好為自己去祈禱.
我們一起祈禱.
.
多謝你.

$^{881}$多謝你在我們人生裡面.
做奇妙的工作.
多謝主.
我們以前風雲有你.
今天要親眼見你.
求主你繼續在我們生命裡面.
做奇妙的工作.
求主你繼續開我們的眼睛.
在一個我們很迷惘.
很混亂的人生階段.
讓我們再一次立志跟從你.
讓我們今天為我們人生.
一個很關鍵的時期.
做一個很重要的決定.
我們再次將身體獻上.
將我們的人生獻上.
當作活祭.
主我們這樣做.
你必然喜悅.
亦都是應當的.
求主你繼續使用我們.
我們不知道明日會如何.
亦都不知道會有什麼.
預見的事情.
我們亦都知道.
人生真的很短暫.
我們求主繼續閱立.
我們在你面前的祈禱.
我們當願意再一次.
將我們的生命.
交給你的時候.
主你讓我們有平安的心.
讓我們有智慧.
讓我們去立志.
多去為下一代去思考.
去著想.
亦都願意我們所做的.
所說的.
所做的決定.
都是合乎上帝你的心意.

$^{921}$求主你繼續引領我們.
在座每一位弟兄姊妹.
在我們的人生裡面.
在關鍵的時候.
做重要的決定.
以致讓更多人可以得到祝福.
求主幫助.
我們同心的祈禱.
奉耶穌名求.
阿們.
我們下次見吧.
\newpage



\section{約翰福音 8:12-30}
\label{sec:EfFf_ZeLMAM}
\textbf{2025年6月22日崇拜直播｜梁家麟牧師｜耶穌與天父的關係｜約翰福音8\_12-30}
\newline
\newline
連結: \href{https://youtube.com/watch?v=EfFf_ZeLMAM}{\texttt{https://youtube.com/watch?v=EfFf\_ZeLMAM}} ~~~~ 語音日期: 2025-06-22
\newline
\newline
\hyperref[sec:f__YU0k5QKA]{\small{< < < PREV SERMON < < <}}
~
\hyperref[sec:index_chronic]{\small{[返順時目]}}
~
\hyperref[sec:12V0_gKaK_A]{\small{> > > NEXT SERMON > > >}}
\newline
\newline
約翰福音 8:12-30
\newline
\begin{longtable}{cl}
\hline
\hline
章節 & 經文 (和合本修訂版)\\
\hline
8:12 & \begin{tabularx}{0.7\textwidth}{X} 耶穌又對眾人說：「我就是世界的光。跟從我的，必不在黑暗裡走，卻要得著生命的光。」 \end{tabularx} \\ \\ \relax
8:13 & \begin{tabularx}{0.7\textwidth}{X} 法利賽人對他說：「你是為自己作見證，你的見證不真。」 \end{tabularx} \\ \\ \relax
8:14 & \begin{tabularx}{0.7\textwidth}{X} 耶穌回答他們，對他們說：「即使我為自己作見證，我的見證還是真的，因為我知道我從哪裡來，到哪裡去。你們卻不知道我從哪裡來，到哪裡去。 \end{tabularx} \\ \\ \relax
8:15 & \begin{tabularx}{0.7\textwidth}{X} 你們是以人的標準來判斷人，我不判斷任何人。 \end{tabularx} \\ \\ \relax
8:16 & \begin{tabularx}{0.7\textwidth}{X} 即使我判斷人，我的判斷也是真確的，因為不是我獨自在判斷，而是差我來的父與我一同判斷。 \end{tabularx} \\ \\ \relax
8:17 & \begin{tabularx}{0.7\textwidth}{X} 你們的律法也記著說：『兩個人的見證才算為真』。 \end{tabularx} \\ \\ \relax
8:18 & \begin{tabularx}{0.7\textwidth}{X} 我是為自己作見證，還有差我來的父也為我作見證。」 \end{tabularx} \\ \\ \relax
8:19 & \begin{tabularx}{0.7\textwidth}{X} 於是他們問他：「你的父在哪裡？」耶穌回答：「你們不認識我，也不認識我的父；若是認識我，也會認識我的父。」 \end{tabularx} \\ \\ \relax
8:20 & \begin{tabularx}{0.7\textwidth}{X} 這些話是耶穌在聖殿的銀庫房裡教導人的時候說的。當時沒有人捉拿他，因為他的時候還沒有到。 \end{tabularx} \\ \\ \relax
8:21 & \begin{tabularx}{0.7\textwidth}{X} 於是耶穌又對他們說：「我去了，你們會找我，而你們會死在自己的罪中；我所去的地方，你們不能去。」 \end{tabularx} \\ \\ \relax
8:22 & \begin{tabularx}{0.7\textwidth}{X} 猶太人說：「他說『我所去的地方，你們不能去』，難道他要自殺嗎？」 \end{tabularx} \\ \\ \relax
8:23 & \begin{tabularx}{0.7\textwidth}{X} 耶穌對他們說：「你們是從下面來的，我是從上面來的；你們是屬這世界的，我不是屬這世界的。 \end{tabularx} \\ \\ \relax
8:24 & \begin{tabularx}{0.7\textwidth}{X} 所以我對你們說，你們會死在自己的罪中，你們若不信我就是那位，就會死在自己的罪中。」 \end{tabularx} \\ \\ \relax
8:25 & \begin{tabularx}{0.7\textwidth}{X} 他們就問他：「你到底是誰？」耶穌對他們說：「我從起初就告訴你們了。 \end{tabularx} \\ \\ \relax
8:26 & \begin{tabularx}{0.7\textwidth}{X} 我有許多事要講論你們，判斷你們；但差我來的那位是真實的，我從他那裡所聽見的，就告訴世人。」 \end{tabularx} \\ \\ \relax
8:27 & \begin{tabularx}{0.7\textwidth}{X} 他們不明白耶穌是對他們講父的事。 \end{tabularx} \\ \\ \relax
8:28 & \begin{tabularx}{0.7\textwidth}{X} 所以耶穌說：「你們舉起人子以後就會知道我就是那位了，並且知道我沒有一件事是憑著自己做的。我說這些話是照著父所教導我的。 \end{tabularx} \\ \\ \relax
8:29 & \begin{tabularx}{0.7\textwidth}{X} 差我來的那位與我同在；他沒有撇下我獨自一人，因為我一直行他所喜悅的事。」 \end{tabularx} \\ \\ \relax
8:30 & \begin{tabularx}{0.7\textwidth}{X} 耶穌說這些話的時候，有許多人信了他。 \end{tabularx} \\ \\
[1ex]
\hline
\hline
\end{longtable}
$^{1}$《新約聖經》第138頁.
耶穌有對眾人說:.
「我是世界的光,.
跟從我的,就不再黑暗裡走,.
必要得著生命的光.」.
法利賽人對他說:.
「你是為自己作見證,.
你的見證不真.」.
耶穌說:.
「我雖然為自己作見證,.
我的見證還是真的,.
因為我知道我從哪裡來,往哪裡去,.
你們卻不知道我從哪裡來,往哪裡去,.
你們是以外貌判斷人,.
我卻不判斷人,.
我就是判斷人,我的判斷也是真的,.
因為不是我獨自在這裡,.
還有差我來的父與我同在.」.
「你們的律法上也記著說,.
兩個人的見證是真的,.
我是為自己作見證,.
還有差我來的父也是為我作見證.」.
他們就問他說:.
「你的父在哪裡?」.
耶穌回答說:.
「你們不認識我,也不認識我的父,.
若是認識我,也就認識我的父.」.
這些話是耶穌在店裡的庫房.
教訓人士所說的,.
「也沒有人拿他,.
因為他的時候還沒有到.」.
耶穌又對他們說:.
「我要去了,你們要找我,.
並且你們要死在罪中,.
我所去的地方你們不能到.」.
猶太人說:.
「我所去的地方你們不能到,.
難道他要自盡嗎?」.
耶穌對他們說:.
「你們是從下頭來的,.

$^{41}$我是從上頭來的,.
你們是屬這世界的,.
我不是屬這世界的,.
所以我對你們說,.
你們要死在罪中,.
你們若不信我是基督,.
必要死在罪中.」.
他們就問他說:.
「你是誰?」.
耶穌對他們說:.
「就是我從起初所告訴你們的,.
我有許多事講論你們,.
判斷你們,.
但那差我來的是真的,.
我在他那裡所聽見的,.
我就傳給世人.」.
他們不明白耶穌是指.
著父說的,.
所以耶穌說:.
「你們舉起人質以後,.
必知道我是基督,.
並且知道我沒有一件事.
是憑著自己作的,.
我說這些話乃是.
照著父所教訓我的,.
那差我來的是與我同在,.
他沒有撇下我獨自在這裡,.
因為我當作他所喜悅的事.」.
耶穌說這話的時候.
就有許多人信他..
(主持人: 請問您是否在聽話?).
(主持人: 您是否在聽話?).
各位節目早安!.
幾個月沒有和大家見面了.
我們繼續漫長的約翰福音之旅.
我們今天這一段聖經.
是比較神學性的.
所以在應用上比較單薄.
要大家多些忍耐.
去聽一些理論.

$^{81}$這段聖經記載了.
耶穌和法利賽人.
兩個神學的爭論.
都是圍繞耶穌基督的身份和使命.
爭論發生的地點.
在耶路撒冷聖殿庫房旁邊.
如果你還記得幾個月前.
我講上一次的那一段.
多數解經家都相信八章一到十一節.
耶穌去寬恕在行淫被捉的婦人.
那段不是最早版本的約翰福音的一部分.
是後來加上去的.
八章十二節這裡.
其實應該是連接著七章五十二節.
就是跳過了前面那一段.
七章五十二節那裡.
我們記得七章是宗教領袖對耶穌的神聖身份.
有很多的質疑爭論.
這一段是延續七章的爭論.
確實這是古往今來一個非常重要的爭論.
耶穌是誰?.
祂是不是來創丫呢?.
祂是不是上帝呢?.
如果祂是上帝.
我們要怎樣去面對祂?.
我們要怎樣去經營我們自己的人生?.
這是一個大災民的問題.
這也是約翰在寫這一卷福音書的時候.
所辭的主題主旨.
正如我們其實這段話.
印過不止一次.
我向你保證.
未來在講約翰福音的時候.
我會繼續去讀這段話.
二十章三十一節.
約翰說.
但記這些事.
要叫你們信耶穌是基督.
是上帝的兒子.
並且叫你們信了祂.

$^{121}$就可以因祂的名得生命.
整卷約翰福音目的是想做這件事.
叫你信耶穌.
信耶穌.
信祂是上帝的獨生子.
信祂可以為你帶來生命.
七章五十二節之前的爭論.
主要是關乎出身自拿撒勒.
這個窮鄉僻壤的地方的耶穌.
可不可能是來創丫的問題.
這裡有兩個爭論.
就是以耶穌所說的兩句話.
作為導火線.
第一個爭論是八章十二到二十節.
針對的是耶穌一個自稱.
他自稱是世界的光.
八章十二節.
耶穌對眾人說.
我是世界的光.
跟從我的就不再黑暗裡走.
必要得著生命的光.
耶穌又對眾人說.
意思是連接前面七章.
即是二十二節之前的爭論.
所以這不是孤立的一句話.
之前一大堆的爭論.
耶穌說我是世界的光.
在整卷的約翰福音.
耶穌做了七次「我是」的宣告.
第一次在六章三十五節.
耶穌宣稱我是生命的糧.
我就是生命的糧.
我來的必定不餓.
所以我必定不渴.
這裡說我是世界的光是第二次.
之後還有第三次.
羊的門.
我是羊的門.
第四次我是好牧人.
第五次我是復活和生命.

$^{161}$復活在我.
生命都在我.
第六次我是道路真理生命.
第七次我是真葡萄樹.
對我來說.
耶穌的七個「我是」是很重要的.
我是在足足五十年前.
1975年11月18日.
在大球場聽葛培里報道大會.
講報道會.
我在那裡決志.
那篇葛培里講的訊息.
就是耶穌的七個「我是」.
我不知道有沒有人聽過.
在那裡.
葛培里講的那裡.
他不停翻炒.
我之後幾十年後.
他再來香港講道.
也是講這篇.
聽回一樣的訊息.
他來來去去只有幾篇訊息.
耶穌的七個「我是」是很多人聽過的訊息.
我是在聽這篇訊息的時候信主.
耶穌的七個「我是」在《約翰福音》.
這是第二個「我是」.
我是世界的光.
在《約翰福音》.
耶穌是世界的光.
其實是一個很早的說法.
比生命的良其實要早.
在第一章的序言.
我們記得.
約翰怎麼說.
生命在他裡頭.
生命就是人的光.
光照在黑暗裡.
黑暗卻不照受光.
真光照亮一切生在世上的人.
聖經宣稱耶穌是光是很早的.

$^{201}$第一章就已經宣稱了.
不過當然在這裡.
講的是耶穌的自稱.
耶穌自我宣稱.
祂是世界的光.
耶穌是光.
光是相對於黑暗來說的.
光的作用是照明.
就是要驅就黑暗.
在大白天.
你不需要點燈的.
反正點了都沒用.
沒有黑暗就不需要光明.
光是因應黑暗而存在的.
耶穌是世界的光.
是為了這個黑暗的世界而來的.
世界是黑暗的.
所以需要光.
光是真光照亮一切生在世上的人.
耶穌在這裡宣稱.
我是世界的光.
跟從我的.
就不在黑暗裡走.
必要得著生命的光.
耶穌在黑暗中為人帶來光明.
凡跟隨祂的.
就行走在光明之中.
黑暗和光明是相對的.
沒有對黑暗的經歷恐懼.
就很難有面對光明的喜悅.
香港是一個福地.
我們已經很久沒有停電的經驗.
很多地方經常停電.
我在羅馬也經常停電.
我們對黑暗的經驗變得很少.
很少很少.
我家裡到現在還保存著陽燭和電筒.
不過很久沒有用過.
不知道電筒還能不能用.
因為沒有這樣的需要.

$^{241}$對我們來說.
光是主要的.
黑暗是要營造出來的.
晚上關燈之外.
還要拉上窗簾.
才有黑暗的環境.
不過黑暗的環境.
你就會閉上眼睛睡覺.
所以你沒有知覺.
香港跟大多數的大城市一樣.
都是不夜天.
晚上燈光如白晝.
就算我在長洲生活.
長洲晚上路上都很光.
鄉郊地方都很光.
基本上不用電筒.
所以我們很難看到滿天星斗的星空.
很難的.
因為光源不足.
基本上不能看到.
我上個月跟太太的女兒.
在西班牙行朝聖路.
我們去巴塞隆拿的時候.
就住了一間非常潮.
很貴的.
很潮的酒店.
整間房間沒有傳統的燈光.
只是在大門口和床邊.
各有一個屏幕.
所有燈光都是在屏幕按來按去.
我們住三人房.
靠窗的沙發拉出來.
變成單人床.
我睡單人床.
我太太跟女兒睡雙人床.
所以我旁邊是沒有燈光的.
我半夜起床去廁所.
我們住酒店.
通常為了安全起見.
通常都會開了廁所燈不熄.

$^{281}$要不然黑漆漆都不知道怎麼辦.
誰知她是很現代化的.
那盞燈一開了她自己就熄了.
自己熄了.
整間房七黑一片.
我起床一直摸著牆.
希望可以找到燈光.
一直摸到廁所都沒有燈光.
又沒有窗.
那種徬徨的感覺.
摸了十分鐘都找不到燈光.
幾乎有點不方便.
最後要叫醒女兒.
叫她在床頭幫我摸燈開給我.
這是很少有黑暗中的經驗.
唯有我們在黑暗中.
有那種恐懼無助.
有這種體會.
你才會歡迎光明.
歡迎耶穌這個世界的光.
無論是猶太人或希臘人.
都有光明代表真理.
代表正確的道路.
代表智慧的說法.
Inner Light.
裡面有光.
光照.
真理,智慧都不是來自人間.
是來自天上.
舊約聖經也有說.
是我的亮光.
是我的拯救.
因為在你裡面有生命的源頭.
在你的光中我不必得見光.
這些大家很熟悉詩篇的說法.
當耶穌宣稱他是世界的光.
為人類帶來光明,真理和生命的時候.
其實他在說什麼?.
他在宣告他的神聖來源.
和神聖身份.

$^{321}$他有一個Divine Origin.
他的神聖身份.
耶穌為世界帶來光明.
世界原來是黑暗的.
所以耶穌的光不會來自世界.
從哪裡來?.
當然是從天上來.
所以耶穌在這裡同時強調.
他和上帝特殊的關係.
他稱上帝為父.
他是由父差來人間的.
在場反對耶穌的那群法利賽人.
當然明白耶穌這個宣稱的含義.
他們先不駁斥耶穌說的話的內容.
卻簡單批評.
耶穌這個宣告是自說自話.
自己說的.
沒有任何實質的證據.
他們說.
你是為自己作見證.
你的見證當作是真的.
這句話翻譯成廣東話.
就是說得就說.
什麼證據都沒有.
耶穌承認他的宣告.
是無法拿到具體證據.
不過不等於他自己亂說.
他指出問題不在於他不想拿到證據.
而是他根本不會有這樣的證據.
第一.
耶穌說的事實超過了.
在場所有人的經驗.
十四節.
我雖然為自己作見證.
我的見證當作是真的.
因我知道.
我從哪裡來往哪裡去.
你們卻不知道.
我從哪裡來往哪裡去.
耶穌是從天上來的.

$^{361}$所有人.
你和我都只能活在人間.
從來沒有去過天上.
又怎能證明.
耶穌的神聖來源呢.
下蟲不可遇,遇冰.
你沒有可能見過的東西.
又怎能證明呢.
一個剛剛登陸月球的宇航員.
即是太空人.
跟我說月球是怎樣怎樣的.
我只有兩個選擇.
一是相信他,一是不相信他.
我無法雙重檢查.
無法檢查.
檢查甚麼.
我都沒有去過.
亦都去不到.
我只能相信他.
或者不相信他.
第二.
除了我們.
耶穌說的話超出了我們的經驗之外.
第二.
耶穌說的話.
亦都超出了我們的能力.
所能夠判斷.
甚至.
你們是以外貌判斷人.
我腳不判斷人.
凡夫熟足子.
只能夠用感官去認知物理的世界.
只能夠憑肉眼去觀察事物的外貌.
耶穌的神聖來源.
卻是超越了眼見才為真這個層次.
要藉肉眼去判斷.
就肯定無法證明.
或者否證耶穌的宣告.
是嗎.
是那個評估的工具不對.

$^{401}$就無法做有效的評估.
你拿一把彈弓尺.
去量度這間屋的長度高度.
你量度到.
要你去量地球和月球之間的距離.
那就省一點.
你拿一把尺子.
怎樣量都量不到出來.
是嗎.
二十世紀實證論.
Positivism.
同樣指出.
宗教命題.
例如上帝存在.
上帝是慈愛的.
都是無法用人間的經驗.
來加以證明.
to verify.
或者否證.
to falsify.
是嗎.
因此他們推斷說.
宗教命題是沒有認知意義的.
沒有quantitative meaning.
是呀.
你想找實驗室的儀器.
去證明上帝存在.
判斷靈界事物的真偽.
你就省一點.
是嗎.
實證論者.
以此去拒絕宗教信仰的合理性.
我們卻認為.
他們只是揭露了.
人的理性.
人的知識的局限性.
是嗎.
不是上帝是否存在.
而是你無法用你的方法.
來證明上帝存在.

$^{441}$是嗎.
一個科學家.
拿著一支試管問.
你可以在實驗室.
證明上帝存在嗎.
我搖搖頭.
證明不到的.
科學家就說.
說明上帝根本不存在.
宗教信仰是不科學的.
我就說.
這只是證明.
科學方法本身有局限性.
與上帝存不存在.
是沒有關係的.
耶穌強調.
祂不會像法利賽人那樣.
只看到事物的外表.
憑肉眼判斷.
往下的結論.
祂對事物的理解.
特別是對屬靈事物的講論.
都是超越了人間的層次.
耶穌擁有神聖的身份和能力.
耶穌來自天上.
由父上帝差遣來到人間.
祂是唯一有資格去講論.
有關上帝和天上的事情的.
一章十六節.
我們很熟悉的經文.
從來沒有人見過上帝.
只有父懷裡的獨生子.
將祂表明出來.
當說到判斷.
法庭上常用的動詞的時候.
耶穌補充指出.
就是判斷人.
我的判斷也是真的.
因為不是我獨自在這裡.
猜我來的.

$^{481}$父與我同在.
你們的律法上也記輿說.
兩個人的見證是真的.
我是為自己作見證.
猜我來的父也是為我作見證.
古代社會我們知道沒有CCTV.
沒有科學鑑證.
所以要搜集證據.
重組案情是非常困難的.
他們大多數只能夠依據.
在場當事人的.
或者目擊證人的證供.
是說的.
陳明全殺死了一個弟兄.
只能夠在場的.
有人說是,有人說不是.
沒辦法可以重演原來的案情.
所以全部靠人證,口供.
當然人可以作假證供.
所以必須多幾個目擊證人的證供.
互相參照.
根據律法規定.
例如新命紀十七章六節.
要憑兩三個人的口作見證.
將當死的人致死.
不可憑一個人的口作見證.
將他致死.
越是重要的案件.
越不可以一個人的口作見證.
需要多幾個證明.
否則很容易會判錯證.
耶穌自稱來自天上.
為人間提供光明.
祂所做的宣告其實是一個口作見證.
這不可信.
不過耶穌強調.
祂的神聖來源,神聖身份.
不是祂一個人說了算.
是有父上為祂撐腰.
做祂的見證.

$^{521}$可惜的是.
世人不認識祂也不認識天父.
他們憑肉眼和有限的宗教經驗.
也無法明白上帝的旨意.
十九節.
他們問祂說你的父在那裡.
耶穌回答說.
你們不認識我也不認識我的父.
若是認識我也就認識我的父.
耶穌說祂的宣告有事實證據.
只不過法理賽人無法認識這些證據.
耶穌說祂的宣告有見證者.
只不過凡夫俗子的他們看不到.
這是第一個.
源自耶穌宣稱祂是世界的光而產生的爭論.
第二個爭論.
我說過這段經文有兩個爭論.
第二個爭論源自耶穌的另一段說話.
二十一節.
耶穌對他們說.
我要去了.
你們要找我.
並且你們要死在罪中.
我所去的地方你們不能到.
前面.
耶穌說到祂的神聖來源.
祂從天上來.
祂被天父差離人間.
這裡就說到耶穌的去處.
祂之後會去哪裡.
祂宣告祂會很快離開人間.
回到天父那裡.
祂說祂要去的地方.
暫時這群人是去不了的.
耶穌說祂去的地方.
他們不能去.
在場圍觀耶穌和法理賽人辯論的猶太人.
很多人馬上誤會.
以為耶穌打算自盡.
我去的地方你們不能去.

$^{561}$是因為黃泉路上只能獨自行走.
耶穌解釋.
祂不是自盡.
祂要回到祂原來的出處.
就是天家.
耶穌強調祂和世人不同.
祂不是屬於這個世界.
祂是來自天上.
二十三節.
你們是從下頭來的.
我是從上頭來的.
你們是屬於這個世界的.
我不是屬於這個世界的.
耶穌也說到.
在場猶太人將來的去處.
就是他們命運的歸宿.
耶穌從天上來.
世人活在地上.
耶穌來人間的目的是要接世人回天上與他同住.
不過前提是.
他們要先信耶穌.
唯有接受耶穌是彌賽亞救世主的.
才能得到上天堂的福氣.
信耶穌有永生.
不信耶穌的就沒有永生.
眼前這群猶太人大多數不信耶穌.
所以耶穌斷言.
他們終極的命運是死在罪中.
耶穌要去的地方就是天堂.
他們永遠都去不到.
二十四節.
所以我對你們說你們要死在罪中.
你們若不信我是基督.
必要死在罪中.
死在罪中這句話.
死在罪中.
是很重要的一句話.
很多人誤會.
基督徒主張信耶穌有獎就是上天堂.
不信耶穌就有罰就是落地獄.

$^{601}$有人因此說基督教太霸道.
要逼人信耶穌.
如果不信耶穌就罰他們落地獄.
永不超生.
我們知道這個說法是錯的.
是完全錯的.
真相是人落地獄不是因為不信耶穌而受罰.
而是因為他們自己所犯的罪.
是因為犯罪.
世人都犯了罪.
所以罪該萬死.
注定要落地獄.
所以他們是死在罪中.
Die in their sins.
信耶穌是讓人得著免於死亡逃出生天的機會.
所以人不是因為不信耶穌而被罰落地獄.
而是因為不信而錯過了得蒙拯救的機會.
錯過了這個機會.
就好像我患癌.
我拒絕看醫生結果病死.
我不是因為拒絕看醫生就被醫生懲罰.
被醫生整死我.
我是死在自己的癌症裡.
醫生只是醫人.
醫生沒有整死人.
醫生只是拯救罪人.
沒有害人落地獄.
我再重申.
沒有人是因為不信耶穌而被罰落地獄.
這不是聖經的說法.
你是因為犯罪.
所以你注定滅亡.
信耶穌是給我們一個機會.
可以免於滅亡.
這是聖經給我們很清楚的提醒.
猶太人或有興趣知道耶穌的來處和去處.
他們一定最關心自己的將來.
自己死後的去處.
人不為己天誅地滅.
每個人都最關心自己的前景.

$^{641}$他們聽到耶穌對他們未來的論斷就大衛緊張.
馬上追問耶穌你是誰.
25節.
這個問題不僅僅是問耶穌的神聖來源.
身份的問題.
更加要問他們自己命運的問題.
耶穌在25,26節跟他們說.
我從始初告訴你們的.
我有許多事講論你們.
判斷你們.
但猜我來的是真的.
我在他那裡所聽見.
我就傳給世人.
耶穌說.
祂早就告訴他們.
祂是世界的光.
祂是被天父猜來的彌賽亞.
為了拯救他們免受罪的刑罰.
一旦說到天父和天上的事.
在場很多人都聽不明白.
這完全超過了他們的認知範圍.
耶穌只能這樣總結.
在28,29節.
你們舉起人子以後.
必知道我是基督.
並且知道我沒有一件事是憑著自己作的.
我說這些話乃是照著父所教訓我的.
猜我來的是與我同在.
他沒有撇下我獨自在這裡.
因為我常作他所喜悅的事.
唯有在耶穌被釘在十字架完成救贖工程之後.
然後人才真正明白祂的身份和使命.
認識耶穌的神聖地位.
認識耶穌和天父的緊密關係.
才開始有人信祂.
耶穌在這次在聖殿聲嘶力竭地講道.
還是有果效的.
在三十節聖經中說.
有一些人信了祂.
這段聖經講完了.

$^{681}$是不是很辛苦呢?.
只是一段很悶的聖經.
講完這段聖經之後.
我們做一點信仰思考.
還是悶的.
我們已經讀到約翰福音第八章.
整本書超過三分之一.
在福音書中.
我們看到耶穌剛出道不久.
就已經顯出祂非凡的智慧和能力.
祂能夠醫病趕鬼.
能夠行很多驚心動魄的神蹟.
祂的講道非同一般的文士.
有不屬於人間.
應該是來自天上的智慧.
所以你看到一個這麼厲害的耶穌.
你推論祂是非凡的人.
推論祂是先知.
這就有很鐵一般的事實證據.
但耶穌其實不滿意.
人對祂有這樣的理解.
祂希望我們看到祂的神聖.
耶穌具有先知的職能.
但耶穌不是先知.
就算你叫祂做最大的先知也不是.
施洗約翰才是最大的先知.
耶穌比施洗約翰大很多.
施洗約翰幫祂解鞋帶都不配.
所以耶穌不是先知.
耶穌是誰?.
必須說這是超越了人間所有的理解.
耶穌是人?.
還是耶穌不僅僅是人.
祂是神.
祂具有人所不能擁有的神性.
對猶太人來說.
他們是信奉獨一上帝的民族.
這與希臘羅馬社會.
可以接受很多神明是不同的.
羅馬人多神教.

$^{721}$所以不介意多拜一個神.
耶穌是神?.
不要緊,多拜一個.
反正已經滿天神佛.
也不差多一個.
中國人也一樣.
多拜一個關公.
多拜一個孔明.
是沒有關係的.
但猶太人不是.
耶和華是獨一上帝.
如果耶穌基督宣稱自己具有神性.
就要處理.
祂與獨一上帝之間的關係問題.
如果處理不好.
很難想像會有人相信耶穌.
跟隨耶穌.
所以耶穌在這裡.
要處理祂與獨一上帝之間的關係問題.
祂將耶和華上帝稱為父.
強調祂與上帝有父與子的關係.
這個關係是獨一無二的.
不像舊約.
舊約也會稱獸為物.
稱人類為上帝的眾子.
對上帝.
耶穌不是一般的泛稱兒子.
而是獨特的臣子.
這個關係是本然的.
不是偶然的.
不是某個時候.
被上帝認為是乾兒子.
不是.
是永恆的.
這個說法.
很難讓人明白和接納.
我們再看28-29節.
耶穌這麼說.
你們舉起人子以後.
必知道我是基督.

$^{761}$並且知道我每一件事.
是憑著自己作的.
我說這些話.
乃是照著父所教訓我的.
那差我來的是與我同在.
他沒有撇下我獨自在這裡.
因為我常作他所喜悅的事.
連同前面.
耶穌在第五章十九節說過的一句話.
我實實在在的告訴你們.
只憑著自己不能作什麼.
唯有看見父所作的.
只才能作.
父所作的事.
只也照樣作.
你試試將這兩段經文加起來.
很多人直觀的理解是什麼.
直觀的理解是.
耶穌好像是自己否定.
自己貶低自己的獨立人格.
他要宣稱.
他其實沒有自己想法.
他只是傳聖統而已.
他是打手.
他是一字不漏的.
去重塑天父的說話.
一絲不苟的遵行天父的意志.
一心一意做天父喜悅的事.
在教會歷史裡.
有不少人因此就推論.
耶穌其實是次等神明.
Secondary God.
父的出意志.
耶穌是執行者.
耶穌沒有自己的意思.
耶穌沒有說自己的話.
耶穌沒有做自己的事.
所以耶穌不可以和天父平起平坐.
父是差遣者.
只是執行者.

$^{801}$父是定義者.
只是完成者.
差遣者和定義者一定比執行者要大.
是不是.
弟兄姊妹.
如果我們不小心看.
就會錯誤理解耶穌說這番話.
耶穌說這兩段話其實不是假裝謙卑.
宣稱祂只是父上帝的馬仔.
祂沒有自己的意志和觀點.
只是按天父心意辦事.
恰恰相反.
耶穌這段話.
是超級無敵的不謙卑.
極其招析.
極其誇張.
非常狂妄.
耶穌說這段話其實是宣稱.
祂和天父是完全沒有差異的.
天父的想法就是祂的想法.
天父的心意就是祂的心意.
你們如今在人間只是看到耶穌的教導和行動.
但是你們在耶穌身上看到的.
百分之百就是天父的想法和行動.
沒有一滴滴例外.
耶穌指出.
你們看到的不僅僅是我的做法.
我的想法.
也百分百是天父的想法和做法.
所以請你不要質疑我的教導和行為.
是否符合上帝的心意.
不存在這個質疑的空間.
因為我如今的說法.
完全是天父的想法.
我如今的行動.
就是符合上帝最喜悅的行動.
不要試圖離間上帝和耶穌.
區分上帝不可以說他們.
是單數的.
總之天父和耶穌的想法和做法.

$^{841}$耶穌有祂的想法.
天父會不會有另外的想法呢?.
不會有.
他們之間毫無間隙.
一而二而一.
耶穌說這些話.
其實和祂其他的說話連在一起.
你就知道那種震撼性.
就像十章十節.
耶穌宣告什麼?.
我與父原為一.
我和父是同一個.
耶穌曾經以不同的形式表達過這個真理.
十四章十節.
我在父裡面.
父在我裡面.
十二章四十四,四十五節.
信我的不是信我.
是信那個差我來的.
人看見我.
就是看見那差我來的.
你看見我就看見天父.
我就是天父的化身的顯現.
《約翰一書》亦有這樣說.
凡不認子就沒有父.
認子的連父也有.
在《易書》亦有說過.
凡越過基督教訓不常守著的就沒有上帝.
常守著教訓的就有父又有子.
你聽耶穌的話就有父也有子.
認子就等於認父.
不需要分開認.
遵守基督的教訓就等於百分百遵從天父的旨意.
零例外.
弟兄姊妹.
我們無法完全明白三位一體的奧秘.
我們無法知道父子聖靈之間的關係和角色.
但我們要清楚耶穌迫切要傳講給我們的訊息.
不要試圖區分父和子.
不要試圖在耶穌的想法和做法以外.

$^{881}$去識別天父的想法和做法.
耶穌的想法完全就是天父的想法.
耶穌的做法完全就是天父的做法.
從來沒有人見過上帝.
道成肉身的耶穌.
就是我們唯一見到的上帝.
耶穌是上帝在人間的彰顯.
是上帝自己來到人間.
活在人間.
所以我經常說.
我們是基督徒.
只相信耶穌基督.
我們強調三位一體.
因為三位一體是保衛耶穌的神聖.
我們不用說那麼多.
三一平衡.
專一高舉耶穌.
歌頌耶穌.
傳揚耶穌.
跟隨耶穌.
活出耶穌.
全都足夠.
耶穌.
祂所說的就是父上說的.
你見到耶穌就是見到父上.
我與父原為人.
我說過不拉奸的比喻.
我在人間.
我不用問那麼多.
我將來會是怎樣.
我見到我的老婆.
就見到我的將來.
我老婆就是我最有可能的將來.
耶穌基督就是我能夠認識.
百分百完全的上帝.
有祂.
我就有我的今天.
有我的將來.
All and all in Jesus.
一切在耶穌裡.

$^{921}$耶穌是我的一切.
\newpage



\section{申命記 32:45-52}
\label{sec:12V0_gKaK_A}
\textbf{2025年6月29日崇拜直播｜楊珊珊傳道｜朝向應許之地－話別與離別｜申命記32\_45-52;34\_1-12}
\newline
\newline
連結: \href{https://youtube.com/watch?v=12V0_gKaK_A}{\texttt{https://youtube.com/watch?v=12V0\_gKaK\_A}} ~~~~ 語音日期: 2025-06-29
\newline
\newline
\hyperref[sec:EfFf_ZeLMAM]{\small{< < < PREV SERMON < < <}}
~
\hyperref[sec:index_chronic]{\small{[返順時目]}}
~
\hyperref[sec:cD5y395bElI]{\small{> > > NEXT SERMON > > >}}
\newline
\newline
申命記 32:45-52
\newline
\begin{longtable}{cl}
\hline
\hline
章節 & 經文 (和合本修訂版)\\
\hline
32:45 & \begin{tabularx}{0.7\textwidth}{X} 摩西向以色列眾人吟誦完了這一切話， \end{tabularx} \\ \\ \relax
32:46 & \begin{tabularx}{0.7\textwidth}{X} 對他們說：「我今日以這一切話警戒你們，你們都要記在心中，要吩咐你們的子孫謹守遵行這律法上一切的話。 \end{tabularx} \\ \\ \relax
32:47 & \begin{tabularx}{0.7\textwidth}{X} 因為這不是與你們無關的空話，而是你們的生命；因遵行這話，你們的日子必在你們過約旦河得為業的地上得以長久。」 \end{tabularx} \\ \\ \relax
32:48 & \begin{tabularx}{0.7\textwidth}{X} 就在那日，耶和華吩咐摩西說： \end{tabularx} \\ \\ \relax
32:49 & \begin{tabularx}{0.7\textwidth}{X} 「你上摩押地的亞巴琳山脈，到面對耶利哥的尼波山去，看我所要賜給以色列人為業的迦南地。 \end{tabularx} \\ \\ \relax
32:50 & \begin{tabularx}{0.7\textwidth}{X} 你必死在你所登的山上，歸到你祖先那裡，像你哥哥亞倫死在何珥山上，歸到他祖先那裡一樣。 \end{tabularx} \\ \\ \relax
32:51 & \begin{tabularx}{0.7\textwidth}{X} 因為你們在以色列人中得罪了我，在尋的曠野，加低斯的米利巴水那裡，在以色列人中沒有尊我為聖。 \end{tabularx} \\ \\ \relax
32:52 & \begin{tabularx}{0.7\textwidth}{X} 我所賜給以色列人的地，你只可從對面觀看，卻不得進到那裡去。」 \end{tabularx} \\ \\
[1ex]
\hline
\hline
\end{longtable}
$^{1}$教會法案 - 二讀(恢復辯論).
《聖經》第十二章第五十二節.
「莫世向以色列眾人說完了這一切的話,又說:.
『我今日所警教你們的,你們都要放在心上,.
要吩咐你們的子孫,謹守遵行這律法上的話,.
因為這不是虛空與你們無關的事,乃是你們的生命..
在你們割約但可要得遺孽的地上,別因這事日子得以長久.』.
當日耶和華吩咐摩西說:.
『你上這阿巴林山中的尼泊山去,.
在摩押地與耶利哥相對,.
你必死在你所登的山上,歸你列祖去,.
像你哥哥亞倫死在荷爾山上,歸他的列祖一樣..
因為你們在尋的曠野,.
加第四的米利巴水,.
在以色列人中沒有尊我為聖,得罪了我..
我所賜給以色列人的地,.
新命記第34章,第一節.
摩西從摩押平原登尼泊山,.
上了那與耶利哥相對的皮斯加山頂..
耶和華被巴基列傳遞直到達,.
拿忽他里傳遞以法列嗎哪西的地,.
猶帶傳遞直到西海南地,.
和中著乘耶利哥的平原直到索爾,.
都只給他看..
耶和華對他說:.
「這就是我向亞伯拉罕,以撒,瓦割,.
起誓應許之地,說:.
『我必將這地賜給你的後裔.』.
現在我使你眼睛看見了,.
你卻不得過到那裡去.』.
於是,耶和華的僕人摩西死在摩押地,.
正如耶和華所說的,.
耶和華將他埋葬在摩押地,.
巴皮爾對面的谷中,.
只是到今天沒有人知道他的墳墓..
摩西死的時候年一百二十歲,.
眼目沒有昏花,精神沒有衰敗..
以色列人在摩押平原為摩西哀哭了三十日,.
為摩西居喪哀哭的日子就滿了..
亂的兒子約書亞,.

$^{41}$因為摩西曾按手在他頭上,.
他就被智慧的靈充滿,.
以色列人便聽從他,.
照著耶和華吩咐摩西的行了..
以後,以色列中再沒有興起先知像摩西的,.
他是耶和華面對面所認識的,.
耶和華打發他,.
在埃及地向法老和他的一切神僕,.
並他的傳遞,行各樣神蹟奇事,.
又在以色列中人眼前顯大能的手,.
行一切大而可畏的事..
.
弟兄姊妹平安,.
善上的朋友平安,.
不經不覺我們旺朕的講壇已經用了半年的時間講《申命記》,.
今天我們來到這本書的終結篇,.
《申命記》可以說是摩西和以色列人最後的講道,.
而整本書的重點就是上帝準備帶這群以色列人進入迦南應許之地,.
而整個過程建立在他和以色列民之間的盟約關係上..
第32章尾段記載了摩西的離別和華別,.
因為耶和華預言摩西會死於尼泊爾山上,.
所以摩西留下了一番對世世代代以色列人的遺言..
而到了第34章,就是摩西人生的終結篇,.
經文提及耶和華神在摩西臨終前,.
叫他上尼泊爾山,.
在山頂由北望到南,.
俯瞰整個應許之地,.
看完這幅全景後,.
經文說耶和華的僕人摩西照著耶和華所說,.
死在摩亞地..
聖經沒有記載摩西是怎樣死的,.
但強調了一個很特別的事情,.
就是耶和華親自埋葬了摩西在摩亞地的一個山谷裡,.
到今天都沒有人知道他的墳墓在哪裡..
最後,申命記34章的10至12節,.
是一段後人對摩西生平的評價,.
他們說以後在以色列中再沒有興起一位先知,.
好像摩西一樣,.
他就是耶和華面對面所認識的..
我不知道大家對摩西的離世有什麼感受,.

$^{81}$通常在安息禮,.
我想大家都參加過安息禮,.
我們會聽到親友講述去世的親人的生平,.
這個環節通常叫述史或懷緬..
如果我們為摩西寫一篇述史的話,.
我就參考了《出埃及記》,《民數記》和《申命記》,.
以下是我一篇很簡單寫摩西的述史..
摩西組織希伯來人,屬利美茲派,.
家中排行最細,.
尚有哥哥亞倫和姐姐米利安..
摩西出生在以色列人在埃及受壓迫的年代,.
當時埃及的法老採取很嚴厲的政策,.
不單止逼迫以色列人從事繁重的苦工,.
更加下令希伯來的接生婆殺死所有新生男嬰,.
以控制以色列人口的增長..
然而,希伯來的接生婆秉承對神的敬畏和良知,.
拒絕執行這個殘酷的命令,.
令到以色列人的人口依然持續增長..
面對這個情況,法老進一步頒布命令,.
他說凡希伯來人生下來的男嬰都必須扔到尼羅河..
摩西正正就是在這段黑暗的時代出生,.
當摩西三個月大的時候,.
因為男嬰的身份難以再隱藏,.
所以他媽媽就將他藏在一個防水的籃子裡,.
放在河邊..
因著神的保守和帶領,.
法老的女兒發現了這個嬰孩,.
並且心裡很憐憫他,.
於是就將他帶回皇宮撫養,.
取名為摩西,意思是從水裡拉上來..
摩西在埃及的皇宮長大,.
接受當時最頂尖的教育,.
他學會了埃及人各樣的智慧,.
《使徒行傳》七章二十二節說,.
他可以說話行事都大有能力,.
然而縱使他是一個埃及王子,.
摩西沒有忘記自己是希伯來人的身份..
成年之後,摩西因為一次跟同胞出頭,.
打死了一個埃及人,.
結果被法老追殺,.

$^{121}$被迫逃亡到米甸..
在米甸,摩西邂逅當地祭司的女兒,.
西婆拉,二人相識相愛,.
婚後育有兩名兒子..
摩西以為自己的人生這樣下去,.
不過其實神對他的呼召剛剛才開始..
當摩西八十歲的時候,.
他在荷列山放羊,.
就在這個很普通的日子裡,.
他第一次見到一個大意象,.
就是消不盡的荊棘..
在這個意象裡,.
耶和華親自向摩西顯現,.
當時的摩西很害怕,很疑惑,.
還覺得自己完全不夠資格去完成神的使命..
他在神面前自嘲,.
他說我是誰,為什麼會是我?.
他甚至連怎樣介紹神給法老和以色列人,.
都不知道怎樣做,也不懂怎樣說,.
覺得自己完全勝任不了..
不過神給了摩西一個很重要的提醒,.
他說不在乎你是誰,摩西,.
而是在乎我,這位自有永有的耶和華神,.
是會與你同在..
神還告訴摩西,.
他沒有忘記和以色列人祖宗所納的藥,.
也看到以色列人在埃及受苦的情景,.
他要拯救他們的子民從埃及出來,.
帶他們去到應去之地..
為了讓摩西更確信這個使命,.
神給了他三個神跡作為引證..
不過摩西依然很猶疑,.
還說自己是絕口笨舌,.
沒有能力和法老或百姓說話..
說到這裡,神聽到也生氣,.
說「好吧,你叫你哥哥阿倫幫你說吧.」.
摩西這次終於沒話好說,.
乖乖地回到埃及開始了他的使命..
事實上帶領以色列人離開埃及絕不是一件容易的事,.
摩西一方面要面對法老和埃及的術士和他們大鬥法,.

$^{161}$另一方面卻要面對百姓的埋怨和不信任..
因著法老的心硬經過十災之後,.
摩西終於成功帶領以色列人出埃及,.
更親眼見證了神跡紅海分開,.
讓百姓安全渡過紅海..
回想起摩西的人生,.
他曾經在四十歲的時候,.
因為幫自己的同胞出頭,.
被自己的同胞嫌棄,.
說「誰要立你做我們的首領和審判官?」.
但是到了八十歲的時候,.
上帝卻藉著消不盡的荊棘,.
呼召他做以色列人的首領和救贖者,.
從埃及地到紅海,再到曠野,.
摩西在四十年裡走了無數的神跡騎士,.
很忠心地完成了神給他的使命..
在《申命記》四章,摩西提到耶和華神在荷列山和以色列民立約,.
這些律法都是為了我們的好處,.
我們選擇這樣生活,.
是要顯出我們是上帝子民的身份..
摩西到了一百二十歲的時候,.
他在臨死之前寫好了律法書和詩歌,.
提醒新一代的以色列民要將神的話放在心裡,.
他在《申命記》三十二章四十六節警告百姓,.
他說要將我一切的話都放在心上..
在希伯來文的文化裡,.
心就是思想的根源,.
意思就是說你心想什麼,.
你就會行動出來,.
所以如果你忘記了一些很重要的東西,.
別人就會說你沒有心臟在..
摩西也提醒百姓,.
單單你們眼前這一代人相信神是不夠的,.
因為上一代信不代表下一代會信,.
父母可以為子女代求,.
但不可以代替子女去信神,.
所以每一代的人都需要重新跟神續約..
弟兄姊妹,你們還記不記得之前跟神簽了這份合約?.
而今天你們續約了沒有?.
當我們跟上帝有了盟約的關係,.

$^{201}$我們就要用心裝著這些律法上一切的話,.
還記不記得十誡和大誡明,.
不是為了為難我們,.
或者幫忙,.
而是幫助我們活出上帝子民的身份,.
不再是埃及的奴隸,.
又或者是曠野的遊牧民族..
這些律例典章對其他國家來說,.
以色列人是很怪的,.
因為律例這個字,.
原文的意思是跟神有一個關係,.
他們只是拜耶和華這個神,.
典章這個字的字源,.
是跟公義是同一個字根,.
說的是人與人之間有一個合宜的關係,.
其他國家的當權者,.
他們會以他們的利益作法律的出發點,.
但是上帝的典章,.
是會考慮和保障卑微的人,.
窮人,無權無勢的人,.
要所有屬於他的子民,.
都要愛倫社,如同自己..
摩西在《申命記》32章47節,.
他釘附以色列民,.
他說:喂,這些律法不是虛空的,.
不是跟你們沒有關係的,.
千萬不要當神的話語沒有實際性,.
否則你的命真的凍過水了..
因為摩西早在《申命記》的8章2至3節,.
已經說了,.
你要記得,這40年來,.
耶和華,你的上帝在曠野一路引導你,.
是要磨練你,考驗你,.
要知道你的心是怎樣,.
是否願意遵守他的誡命,.
神磨練你,任你飢餓,.
將你和列祖所不認識的嗎哪,.
即是那些不知是甚麼的,.
賜給你吃,.
讓你知道人活著..

$^{241}$大英姐妹,是甚麼?.
不是單靠食物,.
那是靠耶和華口中所說的一切話..
他們吃過嗎哪,.
要用心記住耶和華,.
緊緊跟隨他,.
因為神就是你的生命,.
必使你們在應許之地日的日子穩打穩陣,.
長長久久..
就是這個摩西跟我們說完心的大日子,.
耶和華就叫他做了三個動作,.
首先我想說頭兩個動作,.
第一就是上尼泊爾山,.
應該是48節,沒錯,.
給大家一個提示,.
查考聖經的時候,.
其中一樣很重要的事情,.
就是要留意一些Landscape,.
即是地方的樣貌,.
意思是我們要留意那個地方的名字,.
地形如何形容那個地方,.
還有如何形容那個地,.
都是很值得我們深思的,.
譬如我們作為香港人,.
其實我們很少需要向另一個香港人講解旺角是一個什麼地方,.
因為大家都知道旺角,.
最旺的那隻角,.
是香港著名旅遊和購物的區域,.
交通非常方便,.
所以大家回來教會都很方便,.
香港的中心地帶,.
不需要再說,.
但如果我特意向旺鎮的弟兄姊妹介紹旺角這個地方,.
一定是因為有些很特別的東西想提,.
又或者有些很重要的地方想大家知道,.
例如很重要的事情就是,.
如果大家有看我們的「繼往開來」,.
裡面應該有說教會搬過幾次的地點,.
為什麼會搬去那裡,.
又或者那條街有什麼店舖,.

$^{281}$背後的意思就是想我們看到上帝在我們不同年代,.
他的心意,.
恩典的供應和神的帶領和保守,.
所以大家可以回去看看這本精美的「繼往開來」,.
弟兄姊妹,我們試想像一下,.
摩西走了在曠野40年,.
你覺得他這40年看到的曠野是什麼樣貌,.
我不是說天上的雲柱火柱,.
而是說曠野的風景,.
你有沒有想過曠野的風景是怎樣的,.
曠野的意思就是什麼都沒有,.
很荒蕪,也很危險,.
明明只需要11天的路程,.
因為以色列民不斷抱怨,.
摩西這位領袖足足帶領他們走了40年,.
看了40年的人在野,.
而就在這一天,.
上帝讓摩西看到一個新的風景,.
不知道有沒有尼泊爾山的風景,.
一個是他夢寐以求的風景,.
這是我一個朋友在以色列拍攝的,.
當時還沒有打仗,.
這是一個家的風景,.
不需要再漂泊的風景,.
尼泊爾山這個地方,.
就是坐落在皮斯加山麥大的山頂,.
由這個位置,摩西可以想像,.
他對眼看到的景象,.
就是由左手邊的北邊,.
幾年一路掃過去,.
未到中間之前,.
距離100英里的丹,.
然後又看到西北邊的以法連和馬來西,.
再看到中間的猶大泉地,.
還看到西海,.
繼續向右邊看就是南地,.
耶利哥,.
最右手邊就是索爾,.
由左到右,.
這樣看下去,.

$^{321}$這是一個鳥瞰式的全景,.
是神和摩西分享應許之地的視覺,.
這個地貌,.
就是代表著以色列人夢寐以求的家,.
亦都是代表神的應許,.
很真實,.
不過,.
神給了摩西第三個指令,.
就是要他死在這裡,.
為什麼摩西要死在尼泊爾山,.
到底發生了什麼事?.
大家看我們的程序表,.
32章的51節,.
提到一件很關鍵的往事,.
就是米利巴水的事件,.
民數記20章記載了.
事情發生在以色列出埃及後的第40年,.
當時回縱無水飲,.
於是聚集攻擊摩西和阿倫,.
其實那群百姓不是第一次集體投訴摩西和阿倫,.
之前在民數記的16章,.
可拉事件或百姓金牛毒事件,.
他們有爭吵,.
他們對摩西和阿倫的態度,.
就像現在年輕人所說的「靈專」,.
即是一點尊重也沒有,.
現在摩西再一次面對一個靈專的處境,.
神吩咐摩西拿起一支杖,.
焦醉回縱,.
對盤石說話,.
這樣盤石就出水,.
百姓就有水飲,.
不過摩西沒有完全照著神的吩咐去做,.
第一,她沒有對盤石說話,.
反而對以色列人說話,.
她稱他們為背叛的人,.
呼籲他們聽摩西自己說話,.
第二,摩西舉起支杖,.
用力擊打盤石兩下,.
心水清的弟兄姊妹就會聽得出,.

$^{361}$耶和華的吩咐和摩西的行動有些差異,.
在這裡我要強調,.
上帝沒有責備摩西,.
她發怨言或是她生氣,.
因為上帝也是這樣罵以色列民,.
很叛逆,摩西是神的代言人,.
說他們很麻煩,.
背叛人的人也沒有大問題,.
不過摩西發怒,.
這口怒氣她沒有好好消化,.
所以成為了她的燃料,.
令她心煩氣燥,.
結果被燃料驅使她選擇做一個很魯莽,.
甚至是僭越上帝的行為,.
沒有照神的意思去做,.
不相信或偏偏不向盤石說話,.
反而用自己的方法擊打盤石,.
這個行動會令百姓誤以為石有水出,.
水出是因為摩西之仗勁,.
而不是因為神的話語,.
可能摩西那刻很想用自己的方法,.
威風給百姓看,宣洩她的怒氣,.
不過正正因為這個偏差,.
如果擊打盤石和向盤石說話,.
都能使水流出,.
即是說人意和神意是沒有分別,.
我們還需要將神放在眼裡嗎?.
我們還相信神說話的大能嗎?.
我們有沒有將神放在我們的心上?.
如果沒有分別,.
又如何稱之為尊神為聖呢?.
摩西在這件事上,.
其實很完美地示範了她自己的遺言,.
就是不將耶和華的話語放在心裡,.
結果就自食其果,.
神是興曼不得的..
弟兄姊妹,聖經說,.
我們的人的心比萬物都詭詐,.
壞到極處,.
有誰能夠看通看透呢?.

$^{401}$我知道我們當中有很多很厲害的弟兄姊妹,.
甚至有些弟兄姊妹很有能力,.
不過有時候我不知道大家有沒有同感,.
就是越有能力,反而越難依靠神..
我今天很想和大家分享我作為傳道人的掙扎,.
現在我和姑姐一起住,.
雖然我姑姐不是信徒,.
但其實她每個星期都要回教會,.
其實她都回了幾年教會,.
她過去十幾年,.
她天天都看靈修資料,.
最近這半年她更向教會借了一本聖經來看,.
時不時我都回到教會看她翻看聖經,.
甚至主動問我一些問題,.
姑姐還很疼我,.
特別收留了我這個傳道人照顧我的生活,.
不過我姑姐是拜神,.
所以她還不是信徒..
我記得上星期有一晚,.
我下班回家很不開心,.
覺得自己很辛苦,.
就和姑姐說,.
姑姐我很不開心,.
誰知我姑姐竟然和我說一句,.
這樣是因為你不夠依靠神..
當時我心裡簡直是嘰哩咕嚕,.
還反問姑姐,.
我已經很努力了,.
這樣都不算是依靠神?.
誰知姑姐再說一句,.
你就是不依靠神..
當時我那一刻覺得,.
我姑姐不明白我,.
裡面有很多反感,.
亦不想再和她討論神學,.
不過當我靜下來,.
開始反思,.
其實我的好努力是甚麼意思?.
是否靠自己去靠神?.
但是依靠神是否需要努力?.

$^{441}$到底是要努力還是放手?.
有些事是否神叫我先做?.
還是我走得比較快,比較實?.
走得快過實,.
所以我這麼累?.
這些問題在我心裡不斷打轉,.
直到最後我只能在禱告裡面說一句,.
主啊,我很小信,.
求你幫助我的不信..
主啊,我很小信,.
求你幫助我的不信..
我盼望我們旺鎮的弟兄姊妹,.
不要像我上次高估自己那樣,.
俗語也說,.
差之毫厘,謀之千里..
有時候有些代價是我們可以承受,.
不過不是每次都可以,.
就好像摩西這次,.
因為他的不信,他得罪神,.
沒有在以色列人中尊耶和華為聖,.
結果神下了命令,.
不容許他親自踏足應許之地..
其實我覺得神也合情合理,.
因為如果領袖做得不好,.
小的就會學壞手勢,.
不能親自踏足迦南這個地方..
可能你會覺得很可惜,.
因為摩西只差一步,.
就可以踏足迦南這個地方..
如果他當初在米利巴水事件,.
面對班半逆的人,.
他稍微忍一時風平浪靜,.
退一步海闊天空,.
結局可能就會完全不同..
弟兄姊妹,我們有沒有想過,.
其實更多時候,.
上帝不是要我們單單忍耐環境和人,.
而是要我們學懂等候上帝..
接著我就會有另一個問題,.
要等到何時?要等到多久?.

$^{481}$坦白說,這是神的主權,我不知道..
不過我相信,.
在我們每一個人的生命當中,.
或多或少都曾經經歷過上帝,.
描繪我們的祈禱..
而同時我們都知道,.
從來人都是不可信的,.
唯有神才是完全信實的..
所以我很感恩,.
因為我們和神之間的合約,.
另一邊簽名的是上帝祂自己,.
只有神的屬性才是完全信實的..
其實我的見證還未完,有下一部分..
剛剛在星期三,我去了一個報道會,.
講一個報道信息..
我發現原來上帝利用了那個信息,.
即是用了一個繪本的故事,.
我發覺上帝透過繪本告訴我,.
祂很愛我,和祂肯定接納我..
就好像我們剛剛唱的那首詩歌,.
神會以真理和愛引導我們..
我知道神看見我,祂接納我,.
我就夠了..
於是我記得當晚報道會結束後,.
校長請了我們一班事奉人員喝糖水,.
我問校長可否多買一個糖水,.
請我姑姑喝..
於是我帶了那瓶糖水回家,.
我留了一張紙給我姑姑說,.
我今天去營會工作,講福音,.
所以校長請我和家人喝..
要倚靠神,神感動人,請我們喝糖水..
第二天,我姑姑很甜蜜,.
她回覆了一張紙,.
她說:「多謝.」.
送給黃浸穎弟兄姊妹,.
她說:「要有信心,神會幫你.」.
當年在西瀨山,.
上帝親自在雲中降臨,.
站在摩西面前,.

$^{521}$向她宣告祂的名,.
亦即祂的屬性..
出埃及記34章6節,.
祂說:「耶和華,耶和華,.
有憐憫,有恩惠的上帝,.
不輕易發怒,.
且有豐盛的慈愛和信實.」.
即使摩西不完美,.
不能夠「Finish well」,.
但其實上帝已經好恩待摩西..
你想想,摩西是在120歲.
才在地上的老虎釋放出來,.
完成了她在地上所有的責任..
即使來到我們今天科學昌明的時代,.
120歲絕對是一個長壽的標誌,.
而摩西在120歲那一年,.
她的身體依然健壯,.
雙眼視力沒有分花..
有些英文的譯本還說,.
摩西的視力是銳利的,.
她有一個很鮮明的眼睛,.
否則她又怎會上到.
比斯加沙漠的尼泊爾山的山頂,.
還要有足夠的眼力.
可以親眼看到.
英雄之地遼闊的景象..
摩西也不用擔心,.
她自己是黃絲街頭,.
又或者好像西藏的天葬,.
被禿鷹吃了她的身體,.
因為耶和華親自為摩西埋葬,.
我相信在聖經的記載裡,.
只有摩西一個人是由上帝親手埋葬,.
這份恩典是多麼之獨特..
至於摩西最掛心,.
又最激心的百姓,.
神早已安排了亂的兒子,.
約書亞擔任摩西領導的角色,.
帶領這群在曠野長大的以色列人,.
最終成功進入迦南地,.

$^{561}$實現神的應許..
由此可見,即使人失信,.
神仍然可信..
上帝面對面認識她的僕人摩西,.
祂很愛錫摩西,.
不單對她負責到底,.
即使摩西在米利伯水事件中失焦,.
也無間神和人對她的評價..
在最後申命記34章10至12節,.
聖經是這樣評論摩西,.
「以後在以色列中再沒有興起一位先知.
像摩西一樣,.
她是耶和華面對面所認識的.」.
耶和華差遣她在埃及地.
向法老和他一切的神僕,.
以及他的傳弟行一切的神蹟和奇事..
摩西又在以色列眾人眼前.
行一切大能的事和一切大而可畏的事..
最後我盼望摩西一生中的高高低低,.
都能夠與每一位你們的生命對話,.
無論她的得失成敗都能夠成為我們的借鏡..
生命道路的選擇從來都不是一次性,.
而是一場不斷回應和選擇的旅程..
摩西在120歲的時候安息處懷,.
「我不知道弟兄姊妹你有多長命,.
因為我也不知道自己有多長命,.
只有上帝知道.」.
有一句經文是送給今年考完DSE的同學,.
我覺得這句話也很適合送給在座的每一位,.
就是詩篇90篇12字,.
「求上主教導我們,明白人生有限,.
使我們做有智慧的人..
願上帝教導我們,愛惜光陰,.
活出智慧的生命,榮耀祂的名,.
奉主耶穌基督得勝的名祈求.
阿們.
\newpage



\section{馬可福音 6:1-13}
\label{sec:cD5y395bElI}
\textbf{2025年7月6日崇拜直播｜陳明泉牧師｜耶穌給門徒的十課－學懂被拒絕｜馬可福音6\_1-13}
\newline
\newline
連結: \href{https://youtube.com/watch?v=cD5y395bElI}{\texttt{https://youtube.com/watch?v=cD5y395bElI}} ~~~~ 語音日期: 2025-07-06
\newline
\newline
\hyperref[sec:12V0_gKaK_A]{\small{< < < PREV SERMON < < <}}
~
\hyperref[sec:index_chronic]{\small{[返順時目]}}
~
\hyperref[sec:OaFcCMZ21q4]{\small{> > > NEXT SERMON > > >}}
\newline
\newline
馬可福音 6:1-13
\newline
\begin{longtable}{cl}
\hline
\hline
章節 & 經文 (和合本修訂版)\\
\hline
6:1 & \begin{tabularx}{0.7\textwidth}{X} 耶穌離開那裡，來到自己的家鄉；門徒也跟從他。 \end{tabularx} \\ \\ \relax
6:2 & \begin{tabularx}{0.7\textwidth}{X} 到了安息日，他在會堂裡教導人。眾人聽見，就很驚奇，說：「這人哪來這本事呢？所賜給他的是甚麼智慧？他手所做的是何等的異能呢？ \end{tabularx} \\ \\ \relax
6:3 & \begin{tabularx}{0.7\textwidth}{X} 這不是那木匠嗎？不是馬利亞的兒子雅各、約西、猶大、西門的長兄嗎？他姊妹們不也是在我們這裡嗎？」他們就厭棄他。 \end{tabularx} \\ \\ \relax
6:4 & \begin{tabularx}{0.7\textwidth}{X} 耶穌對他們說：「先知除了在本鄉、本族和自己的家之外，沒有不被尊敬的。」 \end{tabularx} \\ \\ \relax
6:5 & \begin{tabularx}{0.7\textwidth}{X} 耶穌在那裡不能行甚麼異能，不過為幾個病人按手，治好他們。 \end{tabularx} \\ \\ \relax
6:6 & \begin{tabularx}{0.7\textwidth}{X} 他也詫異他們不信。 \end{tabularx} \\ \\ \relax
6:7 & \begin{tabularx}{0.7\textwidth}{X} 他叫了十二個使徒來，差遣他們兩個兩個地出去，也賜給他們權柄制伏污靈， \end{tabularx} \\ \\ \relax
6:8 & \begin{tabularx}{0.7\textwidth}{X} 並且吩咐他們：途中不要帶食物和行囊，腰袋裡也不要帶錢，除了手杖以外，甚麼都不要帶； \end{tabularx} \\ \\ \relax
6:9 & \begin{tabularx}{0.7\textwidth}{X} 只要穿鞋子，也不要穿兩件內衣。 \end{tabularx} \\ \\ \relax
6:10 & \begin{tabularx}{0.7\textwidth}{X} 他又對他們說：「你們無論到何處，進哪家，就住在哪裡，直到離開那地方。 \end{tabularx} \\ \\ \relax
6:11 & \begin{tabularx}{0.7\textwidth}{X} 若有甚麼地方的人不接待你們，不聽你們，你們離開那裡的時候，要跺掉你們腳上的塵土，證明他們的不是。」 \end{tabularx} \\ \\ \relax
6:12 & \begin{tabularx}{0.7\textwidth}{X} 使徒就出去傳道，叫人悔改， \end{tabularx} \\ \\ \relax
6:13 & \begin{tabularx}{0.7\textwidth}{X} 又趕出許多鬼，用油抹了許多病人，治好他們。 \end{tabularx} \\ \\
[1ex]
\hline
\hline
\end{longtable}
$^{1}$我們要聽神給我們的說話.
我們要聆聽的是.
馬可福音第六章一至十三節.
馬可福音第六章第一節.
耶穌離開那裡,來到自己的家鄉,門徒也跟從他.
到你安息日,他在會堂裡教訓人,眾人聽見就甚稀奇,說:.
「這人從哪裡有這些事呢?所賜給他的是什麼智慧?他所作的是何等的藝能呢?.
這不是那木匠嗎?不是瑪利亞的兒子雅各約西猶大西門的長兄嗎?.
他們每門不也是在我們這裡嗎?他們就嫌棄他.」.
耶穌對他們說:.
「大凡先知,除了本地親屬,管家之外,沒有不被人尊敬的.」.
耶穌就在那裡不得行什麼藝能,不過按守在幾個病人身上治好他們.
他也詫異他們不信,就往周圍鄉村教訓人去了.
第七節.
耶穌叫到十二個門徒來,差遣他們兩個兩個的出去.
也賜給他們權柄,制服,烏龜,並且囑咐他們.
行路的時候不要帶食物和口袋,腰袋裡也不要帶錢.
除了拐杖以外什麼都不要帶,只要穿鞋,也不要穿兩件襪子.
又對他們說:.
「你們無論到哪兒都何處,盡要人的家就住在那裡,直到離開那地方.
那處的人不接待你們,不聽你們.
你們離開那裡的時候,就把腳上的塵土奪下去,對他們作見證.」.
門徒就出去傳道,叫人悔改,又趕出許多的鬼.
用油抹了許多病人,治好他們.
聖經獨筆.
弟兄姊妹平安.
我們接下來開始一個新的講道系列.
我們用了半年的時間一起去看摩西和整個以色列民族埃及.
在《申命記》裡面的記載.
但是到了下半年我們又轉換,一起讀新約.
在這個系列裡面,我們就叫這個叫耶穌給門徒的十課.
其實耶穌給門徒的課是不止十次的.
不過我們濃縮在《馬可福音》裡面,由第六章讀到第十章.
在這裡的門徒訓練,其實就記載著耶穌帶著門徒.
或者他差遣門徒去做很多不同的任務和使命.
在過程裡面會看到門徒自己的生命狀況.
當然在整個傳道的歷程裡面.
耶穌也會將很多很重要的一些事奉的心路或者經歷.
透過門徒的實踐,耶穌也會預定地告訴他們.
今天我們上的第一課,我們第一個開始看的.

$^{41}$就在第六章,由第一節讀到第十三節.
有些弟兄姊妹問,穆斯林為什麼不由第一章開始讀起呢?.
那也是真的,一來我們篇幅所限,我們沒有辦法,整個《馬可福音》全部看完.
更加重要的是,我們這次特別要看這一段.
門徒開始落手落腳,由第一章到第五章.
其實過去裡面,門徒主要跟隨或者被耶穌呼召.
真的落手落腳去做,不是很多.
但是在這裡開始,門徒真的被差遣出去實踐.
就正如你學車一樣,如果你說學車不停學理論科學.
學take multi,那些選擇題,你不覺得自己在學車.
唯有開著車在路面裡面,你怎樣去駛那台tie ,怎樣踩油門,怎樣踩煞車.
你有那種真實感覺,你就真的正式覺得自己在學車.
現在這裡都是一班新手的門徒,真的學車了.
耶穌教他們,第一個要出去的,要做什麼,要怎樣學呢?.
學生命什麼的經歷呢?今天要特別要他們學的就是學識被拒絕.
我們先祈禱求上帝幫我們,同心祈禱.
天父,願你的話語成為我們的幫助.
帶領我們,讓我們能夠在你的話語當中,讀得到生命的亮光.
又讓一世間的孫講,願孫講的講得清楚明白.
又聆聽的弟兄姊妹,讓他們每一個都能夠聽得入心.
在我們生活當中實踐開去,求你與我們同在.
我們恭敬將來的時間交託,願你賜福.
獻上我們的禱告,奉靠基督耶穌,得勝的聖名祈求,Amen.
在我們讀的今天這段經文之前,有一些生命經歷.
需要勾起你,以至幫你讀這段聖經的,就是被拒絕的經驗.
其實我們從小到大,都在經歷一些被拒絕的經歷.
從你小時候開始,可能你上學的時候,你跟其他同學玩.
同學說你很奇怪,就拒絕了你.
其實那一刻你已經是學識被拒絕.
如果你的心在那一刻消化不了.
覺得這種被拒絕就佔據你心靈所有空間.
你就很不喜歡上學.
又或者你心裡面就製造了一些憤怒.
看到同學這麼壞,不跟我玩,於是小息的時候打他.
這些可能就是用報仇的方式去面對.
到你中學階段成長的時候.
開始看到一個女孩子,或者看到一個男孩子.
你想示愛,寫一張情信給他,表達你的愛意.
但人家說,對不起,我們只是做普通朋友比較好.
做同學比較好,也是被拒絕的.

$^{81}$如果你沒有風度,那些男孩子,我也試過沒有風度.
然後就說你很厲害嗎?你以為你很厲害嗎?.
就會唱衰人家,或者本來你上一刻很喜歡他.
因為人家拒絕你,然後你就說人家很差.
又或者你被拒絕的時候感受不好.
你就覺得自己是一個很差的人.
於是就自暴自棄.
到你成年人,到工作崗位裡面.
可能你要對你的客戶,你預備了很多的時間.
預備了一個Tender,但到最後失落了,被拒絕.
林林總總,細緻到個人,到後來的國家大事.
其實這個世界未必一定按著我們的心意來運作.
懂得怎樣去面對被拒絕.
是成為了人類世界裡面一個很重要的態度.
如果你用這個角度來看今天,我們看馬可福音第六章.
你會讀到有一點點端倪.
這段經文我分為三個部分,我們先看第一個段落.
由第一節讀到第六節.
經文裡面是這樣記載的.
耶穌離開那裡來到自己的家鄉.
門徒也跟隨他.
祂安息日,祂在會堂裡教訓人.
眾人聽見就稀奇說.
這人從哪裡有這些事呢?.
所賜給他的是什麼智慧?.
他所作的是何等的異能呢?.
這不是木匠嗎?.
不是瑪利亞的兒子,雅各約西猶大西門的掌興嗎?.
他們會門也不在我們這裡嗎?.
他們就嫌棄他.
耶穌對他們說.
大凡先知,除了本地親屬,本家之外,沒有不被人尊敬的.
耶穌就在那裡不得行什麼異能.
不過按守在幾個病人身上治好他們.
他也詫異他們不信.
就往周圍鄉村教訓人去了.
其實這個段落跟之前一至五章的段落的記載.
是一個很強烈的反差.
如果你有仔細研讀或看過《馬可福音》一至五章.
基本上耶穌去到每個地方都是萬人迷.

$^{121}$因為他行到的神蹟異能.
他大能的趕鬼.
又或者幫助很多病人.
那些世紀絕症或沒辦法治療的病.
根據耶穌的權能和能力.
很多人都得到醫治.
更重要的是耶穌不單止是醫病趕鬼.
在會堂裡面作教導的時候.
他的權柄跟法利賽人是完全不能媲美的.
因為法利賽人說的是醫書直說.
就照這樣說.
但耶穌的權柄是來自上天.
而且他自己知道真正上帝的心意是怎樣.
他所講論或所講的道充滿了能力.
他所做的充滿了神蹟奇事.
所以他去到各城各鄉宣講的時候.
一定是一個萬人迷.
不過這些風頭或耶穌的能力.
當他回到自己的鄉下.
他帶著門徒一起回到自己的鄉下的時候.
整件事就完全不同了.
值得留意的是.
這次去自己的家鄉.
耶穌刻意帶著門徒.
門徒一直跟著他.
耶穌某程度上示範一下.
他回到自己的家鄉是怎樣被對待.
去到安息日的時候.
當他教訓人的時候.
不是耶穌質疑了那些人.
而是那些人聽見耶穌在會堂裡面講道的時候.
第一件事掃尊於他們對這一家大小.
對耶穌自己的成長作一些質疑.
他們說這個人是從哪裡來的.
為什麼會做這些事.
他就是舊時的那個木匠.
他整個家我都認識.
瑪利亞亞國約西猶大西門的掌興.
就是這個耶穌.
他整群妹妹都是住在我們那裡.

$^{161}$現在這個人有什麼厲害.
學了什麼智慧在這裡特別講.
所以耶穌帶著這一群門徒.
去到自己的家鄉的時候.
看見的是自己從小看到大的那些鄉親父老.
是對耶穌的質疑.
旁邊三叔婆就說.
你為什麼這麼厲害這麼厲害.
說這番話不是肯定他.
而是否定他.
更加重要就是.
出去遊了一個圈子.
學了什麼新的.
掌握到什麼智慧.
才能做到這些的異能.
不過當他們說一個這樣的能力.
其實他們的心裡面是不相信的.
因為這個從小看到大的.
他是什麼底蘊.
他都看得通通透透.
所以這一群的鄉親父老.
這一群耶穌一起長大.
或者看著耶穌長大.
甚至乎可能有些是耶穌的後輩.
他們不覺得耶穌有什麼特別的能力.
在第六章第三節裡面.
特別說一個字.
他們就厭棄他.
有些譯本就寫.
他們就被叛賭.
被叛賭這個字.
讀上去都好像有點抽象.
原文用的那個字就叫.
Scandalism.
這個字和英文字.
Scandal是同一個字.
英文的字就由這裡轉化而來.
Scandal是什麼意思呢.
是一些醜聞.
叛賭的醜聞.

$^{201}$就是說這一群人.
他某程度看到耶穌.
作這個神蹟異能的這種能力.
這一群人得出的結果.
不是要將榮耀歸給上帝.
反而是覺得.
好像當是一些八卦新聞.
或者是一些叛賭人的消息.
同樣一件神蹟.
在其他的鄉村裡面.
彰顯上帝的能力.
不過回到耶穌自己的鄉下.
這些異能或者醫病趕鬼的能力.
就成為了一些人洩洩細語.
在裡面就像是在說一些是是非非.
去質疑耶穌的身份.
所以在這裡不是覺得耶穌很討厭.
而是覺得耶穌所做的事情不是來自神.
或者他們會很詫異.
究竟是什麼搞成這樣.
這麼誇張.
用不著這樣.
這些就是經文裡面的語氣.
當耶穌看到這一群.
從小看到他長大.
或者一起生活的.
在自己的同鄉裡面.
耶穌都很詫異.
耶穌也都說了一句話.
他對於這一種的抗拒.
那種的程度他覺得很詫異.
但是那些鄉村父老.
來拒絕他這種的態度.
耶穌不覺得是很大問題.
因為他特別說.
那些先知去到自己本地.
親屬本家.
是不會受到尊重的.
反而去到外邦其他地方.
先知就得到尊重了.

$^{241}$簡單來說.
耶穌其實都認定.
他這個回到自己鄉下的旅程.
他都是得不到那些人的那種尊崇.
愛戴或者歸榮要給神.
但是他沒有想過.
這些人那種的心.
冷淡到一個地步.
完全好像拒絕上帝一樣.
所以在第五節那裡.
特別記載著.
這句好像很難解的.
他說到耶穌就在那裡.
不得行什麼義能.
耶穌因為那些人的不信.
或者被判到.
所以耶穌就行不到義能.
但是耶穌的能力是不是.
因為那些人信他才行得到呢.
信是不是耶穌行神蹟的條件呢.
不過後面他告訴你.
不是這個意思.
不過按手在幾個人.
幾個病人身上.
就治好他們.
簡單來說.
耶穌都依然做到神蹟.
治到病.
趕到鬼.
不過最終的結論.
不是上帝得著榮耀.
所以那個義能.
那個意思.
他這裡說的Miracle.
就是那個神蹟.
這個義能有兩部分.
第一就是做到一些的醫治.
超自然上帝的介入.
而後半部就是將榮耀歸給上帝.
耶穌只是做到上半部.

$^{281}$在這一群不信.
在這一群拒絕相信的人當中.
耶穌只能夠做到一些.
其他醫生做不到的事.
不過這一群人卻沒辦法將榮耀歸給神.
反而帶來的就是大量的謠言.
大量叛道人的說話.
以致到耶穌就沒辦法行義能.
所以這一節義能的意思.
就是這樣來理解.
整個經文其實是一個.
這個段落講到耶穌不是很討好.
也不是很爽的經驗.
反過來想想你就明白.
作為一個師父.
帶著一群門徒.
如果我呢.
我自己帶著我們年輕人出去.
通常那些長輩一定在後輩面前.
展示肌肉.
一定要示範一下我怎樣做.
我記得以前在旺鎮裡.
我們試過做一個.
都很瘋狂的做法.
就是那時候流行的free hug.
在街上自由擁抱.
那些外國人有些文化.
大家想著互相接納.
但我們不會做這樣的自由擁抱.
於是我們就改了它叫做free prayer.
自由祈禱.
我不知道有沒有跟大家說過.
我們在星期二.
崇拜完了之後.
跟一群年輕人上街.
在旺角街頭舉個牌.
free prayer.
看到有一個人就說.
我隨時為你祈禱.
多有意思多好玩.

$^{321}$其實不過這個舉動不是我建議的.
是我的團友建議的.
然後我說.
好啊好建議啊去吧.
然後他說不是啊牧師.
你也一起去吧.
心裡想.
不是吧.
我幾十歲還要搞這樣的東西.
但是沒辦法.
他這樣叫到.
頂硬上.
好吧要去就去吧.
下去的時候.
一群人站在中橋地下.
撿了一堆.
沒有人敢上去舉牌.
到最後你自己也頂硬上的.
是不是.
於是就自己舉牌.
然後找人來祈禱.
然後就覺得好好示範.
但到最後其實那些人走.
看到牌就不是走到我那裡.
我是趕客的.
那些人就去那些年輕人那裡.
怎樣也好.
怎樣也找了一兩個.
到他們說有傳福音.
然後叫我去.
就傳福音.
你作為一個長輩.
或者一個師父.
你目的就會在門徒面前.
要展示肌肉.
無論怎樣也好.
你都要頂硬上.
是不是.
不過耶穌在這裡.
祂自己知道.

$^{361}$先知在自己本族本鄉.
是不被認同的.
耶穌特別將這個經驗.
用祂自己一個被拒絕的經驗.
祂帶一些門徒.
看到這樣的一個狀況.
所以這個經驗.
耶穌被拒絕.
或者被人嫌棄的經驗.
就成為了耶穌帶門徒.
一個很重要的教材.
去見到人家怎樣拒絕耶穌.
同樣.
祂就用這個經驗.
來拆遣門徒.
所以接下來的.
耶穌在某程度.
祂自己失敗的經驗裡.
或者一個不太討好的經驗裡.
第七節.
耶穌就除職了.
接著第七節.
耶穌叫了十二個門徒來.
拆遣他們兩個兩個地出去.
也賜給他們權柄制服污鬼.
並且囑咐他們.
行路的時候.
不要帶食物和口袋.
腰袋裡也不要帶錢.
除了拐杖以外.
什麼都不帶.
只要穿鞋.
也不要穿兩件褂子.
在第七至第九節.
在耶穌被拒絕之後.
當然不是隨即.
在耶穌的家鄉拆遣門徒.
耶穌因為去到自己的家鄉.
已經沒有辦法.
做到任何異能.

$^{401}$於是耶穌就去了其他的鄉村.
去了別處.
去了周圍的鄉村.
耶穌也除了詫異他們不信之外.
也沒有說將來末日的日子.
這群人就慘了.
有些鄉村是這樣的.
不過去到自己的家鄉.
耶穌沒有說一些咒詛.
或者很強烈的話.
只是默默淡淡言地離開.
他自己的家鄉.
然後他就拆遣他的門徒.
兩個兩個出去.
去做回耶穌所做的事.
值得留意的是這次耶穌的拆遣.
他有些要求給門徒.
他有些什麼要求呢.
他就說先給他們有權柄制服污鬼.
所以那個門徒是有來自耶穌.
來自上天的能力.
以致他們都能夠像耶穌一樣.
可以實踐醫病趕鬼的那種能力.
不過在這次裡面.
他不單止給他們權柄.
他更加要告訴他們.
不要帶什麼.
基本上你留意一下.
我們上行者行山.
裝備都沒有那麼少.
他那裡的裝備是極之少的.
他說什麼呢.
走路的時候不准帶食物.
不准帶袋.
兩手差不多都要搖晃.
然後.
醫護袋裡也不要帶錢.
你想想.
沒有錢.
沒有食物.

$^{441}$連袋子都沒有.
那就是說.
好像真的很輕散的輕裝上陣.
唯一可以帶的是什麼呢.
除了拐杖以外.
什麼都不要帶.
只要穿鞋.
也不要穿兩件褂子.
唯一一樣東西要帶的就是.
穿一雙鞋子.
基本裝備.
然後.
衣服都不要帶多件.
不要帶兩件.
帶一件.
唯有真的要叮囑他們帶的就是帶拐杖.
和合本易拐杖.
以為你的腳有點麻痺才會用的那一支.
其實那支是叫手杖.
那支手杖有什麼用呢.
以前的牧羊人都拿著手杖.
就是如果遇到猛獸.
用來防身的.
又或者.
後來這支手杖.
就代表著一個權柄的象徵.
但是這裡.
其實他的權柄.
不用一支手杖.
不是舉起就能治病.
耶穌就自動將權柄賜給他.
他不需要道具來實踐他的權柄.
所以簡單來說.
那支手杖只不過是防身之用.
手上有東西拿著.
萬一有山賊的時候.
都是防衛保護自己.
等等.
耶穌在這裡.
他要這班門徒完全的輕裝上陣.

$^{481}$剛才示範過如何被拒絕.
接著進去.
夫子就要吩咐這些門徒.
兩個人一個小組.
接著什麼都不帶.
不設防的方式.
用最低度的那種態度.
去到入城裡面.
來實踐耶穌之前所做的趕鬼醫病.
甚至乎有一些傳福音.
傳講天國福音的一些使命.
就是這樣被拆遣出去.
你可以想像的.
其實耶穌在這裡將整個的拆遣.
或者門徒裡面的任務.
他壓縮到一個最基本的單位.
這班門徒有什麼可以依靠呢.
除了是傳福音的人接待他之外.
因為他沒有食物,沒有袋子.
他不可以一個人接觸了之後很好.
接著拿一大量的麵包放在袋子裡.
他沒有額外的東西來擋住額外的食物.
他只有當下裡面.
他就要放下身段.
去到不同的人當中.
來帶著上帝的權柄來幫助他.
這班門徒其實是某程度.
放到自己的身份.
是一個最謙卑的位置.
甚至連基本需要都沒有辦法滿足.
就是一個這樣的方式.
一個這樣的旅程.
來到他們拆遣出去.
值得注意的就是.
只要穿鞋第九節.
也不要穿兩件褂子.
那件外衣.
以前在古代人來說除了是備體之外.
最重要的就是.
多一件衣服其實是保暖用的.

$^{521}$因為晚上的巴勒斯坦地.
或者耶穌的居住地.
晚上會冷的.
沒有太陽之後.
所以有一件衣服就是用來遮住身體保暖.
只是穿下午那件衣服.
一件衣服.
多一件都不行.
那就是說.
晚上的時候.
你一定要找到地方來居住.
否則的話你就快要冷死了.
哇.
你說耶穌給這個任務給門徒.
都很艱鉅.
不過耶穌就是這樣來訓練他們.
賜他們權柄.
但是用最輕散的方式來做.
我們再繼續看.
最後我就會說一下.
究竟說什麼.
整段經文.
第十到第十三節.
第三個段落.
又對他們說.
你們無論到何處.
盡了人的家.
就住在那裡.
直到離開那地方.
何處的人不接待你們.
不聽你們.
你們離開那裡的時候.
就把腳上的塵土堵下去.
對他們作見證.
耶穌就出去傳道.
叫人.
耶穌門徒出去傳道.
叫人悔改.
又趕出許多的鬼.
用油抹掉許多的病人.

$^{561}$治好他們.
第三個段落.
是耶穌再一次.
他在說.
這班門徒出去的時候.
如果被人不接待.
被拒絕的時候.
他要用什麼方式.
來面對.
首先第十節說到.
你要住在別人的家.
你要放下身段.
你擁有的權柄.
醫病都好.
不是你呼風喚雨.
我奉主的命.
現在要徵用你的家來住.
不是用這個態度的.
他只是隨走隨傳.
帶著最大的誠意.
帶著上帝的權柄和話語.
問有沒有這樣的需要.
如果別人接待他.
他就入屋了.
別人不接待他.
他就去另一家.
就是這麼簡單的一個做法.
而且在那裡吃.
在那裡住.
直到離開.
你不可以拿多東西.
你只是當下有什麼需要.
你就吃什麼.
講一句題外話.
這句經文.
以前我在讀神學的時候.
我的老師經常提醒我們.
就是說被人接待.
放上什麼就吃什麼.
他講什麼意思呢.

$^{601}$他說其實你們傳道人.
有時候也有弟兄姊妹請你吃飯.
未必一定吃好東西.
雖然有些弟兄姊妹.
通常都很愛傳道人.
都會請人吃好東西.
但是有些弟兄姊妹家裡.
真的比較基層的.
或者他們真的沒有辦法.
他們放上什麼.
你就吃什麼.
我記得這句話.
經常都在我腦裡.
我記得有一次.
我們和弟兄姊妹上山區.
他們上去的時候.
山區最窮.
所以他們的孩子都不夠營養.
他們吃的就是用那些粟米.
山上的粟米很瘦.
他們煮一些水燙燙的粥.
他們說我們就吃這些.
大家舀一碗在這裡.
舀在這裡的時候.
就有一堆蒼蠅飛下來.
你們吃不吃.
放上什麼就吃什麼.
第一個吃.
樣子還不可以有一點點嫌棄.
哇 正啊.
我也很會演戲.
如果你讓我選.
我不會吃的.
但是他們說 哇 正.
然後他們說.
我們預備了一些最好的東西給大家吃.
什麼呢.
就是鳳爪.
雞腳.
獨立包裝.

$^{641}$用膠袋包著.
每一隻都真空的.
這些就是我們山上的蛋白質.
唯一的蛋白質.
然後我撕開那裡.
有一陣酸酸的醋味.
醃進去.
鳳爪的指甲都很尖.
一拉出來的時候.
它不像我們平時吃的酒樓的鳳爪.
它不發開的.
很緊緻的.
所以一咬下去的時候.
根本咬不到肉的.
然後一邊吃一邊吃的時候.
那個長老說.
牧師啊 其實這個是我們山區裡面最好的蛋白質.
因為我們山區裡面什麼都沒有.
我看到他頂著姊妹選了兩樣東西就扔掉了.
我說不行.
吃掉它.
那些肉要吃掉.
這些是山區裡面最好的東西.
逐一的肉吃掉.
放上什麼吃什麼.
未必是有豐富的東西的.
傳道的人就是要這樣去實踐.
我想大致耶穌叫門徒去到別人的家裡面.
不是叫別人去找一些大富大貴的人家來吃好東西的.
只不過是接待你的.
你就當下進去就吃了.
不是選擇或者要做什麼.
別人擺什麼給你招呼你.
你就吃那些東西.
所以在這裡接待就是這樣去面對.
第十一節就說到.
何處的人不接待你們不聽你們.
你們離開那裡的時候.
就把腳上的塵土躲下去.
對他們作見證.

$^{681}$這個是什麼呢.
這是一些文化的表達.
有一個典故是這麼說的.
具體是什麼呢.
我看不到一些舊約的經典來講的.
都是看過一些旁敲側擊的說的.
話說就是躲下去的塵土.
就是那些猶太人和外邦人談不攏.
談不攏的時候.
表達就是在那些外邦人面前.
猶太人就會踩踩腳.
踩踩腳就把鞋裡面的灰塵踩走了.
這個動作就代表著你和我之間就互不相干.
就是一個這樣的做法.
不是故意弄髒人家的地毯或者什麼.
只不過是踩下那個塵土.
就是說我剛才跟你所講的.
我作出一些的邀請.
你不接受的時候.
我就跟你互不相干.
這個的習俗或者這種的做法.
是舊時猶太人對外邦人的那種做法.
但是耶穌就將這個做法放在.
那些不接待猶太門徒的那些人的屋口門口裡面.
就是那些家庭不接待你.
不聽你講了.
你就在家門口裡面踩踩腳.
將那些塵土留下.
那你就離開了.
在經文裡面講到那個腳下塵土.
躲下去為他們作見證.
那個見證是什麼呢.
就是不聽不講.
但是那個見證有多麼強力呢.
到處都是灰塵.
你是認不出來的.
而且那個見證只是為當下的門徒.
心裡面覺得他不接待我了.
就是一個這樣的見證.
你不跟舊約裡面.

$^{721}$在出埃及記裡面.
在門徽裡面畫一些血.
然後天使經過.
有些記認就不去.
這個塵土不是什麼記認.
也沒有辦法可以留下很長久.
有傷害性或者有些指責性.
很清晰分辨到.
因為有些塵在這裡.
間間戶戶都站有塵的.
所以那個見證.
不是給天使看到.
也不是給基督人看到.
只是為當下的門徒.
那種感受而言.
所以簡單來說.
人家拒絕門徒.
門徒最多只是做這件事.
然後就要離開找第二個家庭.
不過在這個安排裡面.
那個結果是怎麼樣呢.
第十三節告訴我們.
門徒回來是成功的.
趕出許多的鬼.
用油抹了許多的病人.
治好他們.
耶穌這個差遣.
是帶來門徒一個很好的經驗.
不過在這個差遣裡面.
其實他也告訴了門徒.
他要有心理準備.
人家未必一定完全接待他.
整段經文.
如果我們從一節到第十三節.
如果是一個整體來看.
耶穌由自己被拒絕的經驗.
到他差遣門徒.
而且之後教門徒.
怎樣去面對拒絕.
這個就一氣呵成.

$^{761}$究竟這段經文.
耶穌給了什麼錦囊給門徒.
來面對這些拒絕.
或者傳道實踐的經驗呢.
我用四個F來代表.
第一個.
耶穌用自己失敗的經驗.
來呼召門徒.
其實耶穌可以早點呼召門徒.
在一至五章.
耶穌做了一些事情之後.
在耶穌做事情最巔峰的時候.
最厲害的時候.
他就說門徒,看著了.
看呼自己做了多少事情.
神國的權柄有多大.
所以你好像我這樣做.
你就做到事情了.
其實耶穌可以把時間放早一點.
不過耶穌刻意.
要門徒看到他自己被拒絕的經驗.
他要塑造門徒的期望.
他要調教他們的期望.
在他的事奉道路.
或者在他人生裡面.
你要預計被人拒絕.
即使他們所跟的夫子都被人拒絕.
他們作為徒弟.
有什麼可能比先生厲害呢.
所以拒絕是人生的常態.
在這裡同樣對我們都重要.
我們今天覺得自己面皮很薄.
我們經常覺得我們做每一件事.
成功是我們基本一定會出現的結果.
所以我們面對著不成功的經歷.
失敗的經歷或者被拒絕的經歷.
我們就覺得自己好像被人比下去.
大夫心智問鼎姐妹.
其實你人生裡面做多少件事你真的成功.
或者真的很多的學財聲音.

$^{801}$很多的鎂光燈放在你身上.
根本上不是的.
我們面對大部分.
是大部分都是被拒絕的經驗.
只要我們重新去調教我們自己的期望.
我們的人生就不會過得這麼難.
相反來說我們將每一樣東西.
我都將自己的人生.
我做每一樣東西都要成功的時候.
我們的人生實在太大壓力了.
我也覺得我們沒可能做得到.
以致到如果我們經常都是用一個這樣的定位.
做每一樣東西都要成功.
做一樣東西都是很棒的.
我們的人生就好像寸步難行.
小弟是讀市場學.
簡單來說.
每一年都會給一個市場導師你做.
都是由銷售員做起的.
一聽到銷售員就.
以前我是第一份工是暑期工.
你們最討厭的事.
以前我做的就是打電話.
幫某間電訊公司問人買不買電話機.
或者有電話的就買多部電話機.
沒有電話機的就叫他買多部電話.
就是這麼簡單.
在別人的手裡每個月拿多二三百元.
打電話就是找那些客戶.
經常被人吸線.
開始做了幾天我就會覺得.
我這麼大個兒子我做銷售員也好.
我出去銷售我也很少這麼大的失敗經歷.
現在每天被人吸線被人罵.
不知道做什麼的.
然後就想辭職.
我的同學跟我一起去的都已經辭職了.
只剩下自己.
我已經覺得自己的毅力很強.
但是當同學走的時候.

$^{841}$我就想遞辭職信給我老闆.
我老闆其實不是一個基督徒.
不過他說了一番話.
他說一點點失敗你都面對不了.
你怎麼做人.
然後我說不是一點點失敗.
我整個暑假.
每天我都面對著.
還有說成功率只有100個電話裡只有一個.
有99次失敗.
我真的接受不了.
然後他說你每個同學都接受不了.
但是我覺得你可以.
你憑什麼覺得我可以.
因為最遲走的那個是你.
到現在你都還沒走.
你現在才跟我辭職.
不過我真的聽入心.
我覺得真的.
為什麼每個人都放棄.
為什麼自己不繼續呢.
我是不是純粹為了成功.
還是在過程裡.
我其實生命裡掌握到一些東西.
然後這個supervisor他再繼續說.
只要你將失敗當成你生命的一個很普遍的一部分.
別人給你一線就不覺得是什麼.
如果每給你一次線.
你都覺得很受傷.
每一天至少打1000個電話.
你有900多次來傷害你.
你做人都很坎坷.
但是你為的就是為那一兩次.
結束了收到單.
那你就做得到.
那你就為了得著的來興奮.
因為失敗是平常事.
成功就變成一個二數.
所以你就要好好珍惜成功的機會.
他反思我.

$^{881}$我現在是製造一個成功經驗給你.
而不是製造你失敗.
果然是sales的.
他說完這句話之後.
接著那兩天.
我一次過買了17部call機.
做了top sales.
不過人生裡要做什麼呢.
我想說在你人生裡都是一樣.
面對著人生的失敗.
你不要覺得是一件很大的事.
習慣它.
面對它.
其實人生裡基本上是這樣.
正正因為失敗是生活的一部分.
所以成功就很值得慶賀.
如果你的成功.
你做的所有事都是成功.
你慶祝什麼呢.
你的成功.
我們要調教我們的expectation.
同樣傳福音都是一樣.
傳福音有些人禮不虛發.
讀完三幅.
你以為讀完說幾句金口.
別人就可以嗎.
其實不是的.
傳福音大部分時間.
都是覺得別人客氣地拒絕你.
因為你說的是一個很冒犯性.
人生觀念的一個訊息.
你說他在拜的.
在信的東西是錯的.
只有耶穌是對的.
多大的一種否定.
所以那個拒絕是正路的.
但是不要因為這種拒絕.
你就放棄.
第二個F.
在這段經文裡.

$^{921}$耶穌刻意要令門徒無可依靠.
即或他有很多權柄.
即或他來自上天.
有神給他的那種趕鬼的能力.
他其實更加要依靠上帝.
什麼都不帶.
只是帶著最大的誠意.
去到各家各戶面前.
問人有沒有什麼需要.
放下身段讓人接待自己.
這是一個信靠的表現.
信靠上帝.
而不是信靠那些接待自己的人.
這些村民有機會拒絕他.
有些人會招待他.
但是真正的依靠是來自上帝.
生命裡我們做的每一件事.
我們知道我們做的事有什麼價值.
我們就依靠上帝而做.
今天我們最不理想的狀態就是.
成功就覺得是上帝.
失敗就不是來自上帝.
但是在耶穌的講論裡說到.
失敗成功都是上帝給我們的經歷.
我們只管坦然無懼.
靠著上帝為我們所預備的.
在這個門徒課程裡繼續的經歷.
信靠神.
在做每一件事之前.
好好禱告上帝.
可能面對拒絕的.
靠上帝面對這些負面感受.
成功的.
因為知道是靠上帝.
所以功勞不是在我這裡.
而是在上帝那裡.
上帝給我有機會看到他的作為.
第三個了.
For Go.
第三個F裡面強調.

$^{961}$面對人的拒絕.
你不要特別去執意要贏他.
不要強求他.
拒絕就拒絕.
躲一層土在別人的門前.
你沒得逼的.
你沒得逼別人.
或者罵爆別人.
又或者用惡言惡語.
去罵那些拒絕你的人.
這個世界本身每個人都不明白真理.
在這個世界裡每個人都是任意而行.
按人自己的價值.
所以你強求別人.
在你說幾句話就完全改變.
除了神跡之外都沒什麼會出現.
你不要面對別人的拒絕.
覺得他無可救藥.
或者強求他.
用辯論的方式去爭取贏他.
你只用最低度最溫柔的方法.
算了,沒辦法了.
我就走了.
不要死裡難了.
不要強求.
唯有面對有風度的方式.
來面對別人的拒絕.
你的人生才可以繼續.
讓別人來搭前.
否則你就繼續原地踏步.
你繼續強求.
你都沒辦法繼續前行.
最後就是前進.
失敗就是失敗.
不理想就是不理想.
繼續前行.
你要知道如果你不前行.
你的結果是什麼呢?.
你只有一件襪子.
你是沒有食物的.

$^{1001}$你是沒有口袋的.
你不前行.
你跟著人家沒有退路可言.
我想說,弟兄姊妹.
我們經常覺得.
退一步海闊天空.
耶穌在這裡說沒有.
不是海闊天空.
是你凍死在巴勒斯坦地.
在困難裡面.
繼續前行.
這個不是一個無奈或者悲壯的決定.
這個是生存之道.
生命繼續來搭前.
是一個最佳面對著.
被人拒絕的生命最好的道路.
今天這段經文不是很複雜.
但很有意思.
我讀這段經文.
一想起耶穌被拒絕.
但他怎樣教導門徒.
怎樣去面對,學習.
在他未教門徒.
怎樣得到什麼果效之前.
他就教門徒.
怎樣去面對生命裡面.
有一些無理取鬧.
或者拒絕他的人.
這個不單止是對當代門徒.
一個很重要的功課.
也是對我們今天每一位基督徒.
對我們能夠在這個世上裡面.
和人和平共處.
尊重別人.
一個很重要的功課.
耶穌還有很多堂課教我們的.
我們拭目以待.
接下來的馬學福音.
看看耶穌怎樣教門徒.
也教導我們.

$^{1041}$我們同心祈禱.
天父願你在我們當中.
繼續對我們講說話.
率領我們前行.
引導我們.
讓我們走人生的道路.
我們活得有意義.
而且有你同在.
我們心靈安穩踏實.
多謝你.
求你繼續來賜福給眾靈姐妹.
獻上禱告奉基督耶穌得勝的聖名.
Amen.
阿們.
\newpage



\section{馬可福音 6:14-29}
\label{sec:OaFcCMZ21q4}
\textbf{2025年7月13日崇拜直播｜鄭國邦傳道｜耶穌給門徒的十課－付代價｜馬可福音6\_14-29}
\newline
\newline
連結: \href{https://youtube.com/watch?v=OaFcCMZ21q4}{\texttt{https://youtube.com/watch?v=OaFcCMZ21q4}} ~~~~ 語音日期: 2025-07-13
\newline
\newline
\hyperref[sec:cD5y395bElI]{\small{< < < PREV SERMON < < <}}
~
\hyperref[sec:index_chronic]{\small{[返順時目]}}
~
\hyperref[sec:SbKqVXHzrXE]{\small{> > > NEXT SERMON > > >}}
\newline
\newline
馬可福音 6:14-29
\newline
\begin{longtable}{cl}
\hline
\hline
章節 & 經文 (和合本修訂版)\\
\hline
6:14 & \begin{tabularx}{0.7\textwidth}{X} 耶穌的名聲傳開了，希律王也聽見。有人說：「施洗的約翰從死人中復活了，因此才有這些異能在他裡面運行。」 \end{tabularx} \\ \\ \relax
6:15 & \begin{tabularx}{0.7\textwidth}{X} 但別人說：「他是以利亞。」又有人說：「是先知，正如先知中的一位。」 \end{tabularx} \\ \\ \relax
6:16 & \begin{tabularx}{0.7\textwidth}{X} 希律聽見卻說：「是我所斬的約翰，他復活了。」 \end{tabularx} \\ \\ \relax
6:17 & \begin{tabularx}{0.7\textwidth}{X} 原來，希律為他兄弟腓力的妻子希羅底的緣故，派人去抓了約翰，把他綁了在監獄裡，因為希律已經娶了那婦人。 \end{tabularx} \\ \\ \relax
6:18 & \begin{tabularx}{0.7\textwidth}{X} 約翰曾對希律說：「你佔有你兄弟的妻子是不合法的。」 \end{tabularx} \\ \\ \relax
6:19 & \begin{tabularx}{0.7\textwidth}{X} 於是希羅底懷恨他，想要殺他，只是不能。 \end{tabularx} \\ \\ \relax
6:20 & \begin{tabularx}{0.7\textwidth}{X} 因為希律怕約翰，知道他是義人，是聖人，所以就保護他，雖然聽了他的講論十分困惑，仍然樂意聽他。 \end{tabularx} \\ \\ \relax
6:21 & \begin{tabularx}{0.7\textwidth}{X} 有一天，恰巧是希律的生日，希律擺設宴席，請了大臣、千夫長和加利利的領袖。 \end{tabularx} \\ \\ \relax
6:22 & \begin{tabularx}{0.7\textwidth}{X} 他的女兒希羅底進來跳舞，使希律和同席的人都很高興。王就對女孩說：「無論你要甚麼，向我求，我都會給你」； \end{tabularx} \\ \\ \relax
6:23 & \begin{tabularx}{0.7\textwidth}{X} 又對她多次起誓說：「無論你向我求甚麼，就是我國家的一半，我也會給你。」 \end{tabularx} \\ \\ \relax
6:24 & \begin{tabularx}{0.7\textwidth}{X} 她就出去對她母親說：「我該求甚麼呢？」她母親說：「施洗約翰的頭。」 \end{tabularx} \\ \\ \relax
6:25 & \begin{tabularx}{0.7\textwidth}{X} 她就急忙進去見王，求他說：「我願王立刻把施洗約翰的頭放在盤子裡給我。」 \end{tabularx} \\ \\ \relax
6:26 & \begin{tabularx}{0.7\textwidth}{X} 王就很憂愁，然而因他所發的誓，又因同席的人，不願食言， \end{tabularx} \\ \\ \relax
6:27 & \begin{tabularx}{0.7\textwidth}{X} 就立刻派一個衛兵，吩咐拿約翰的頭來。衛兵就去，在監獄裡斬了約翰， \end{tabularx} \\ \\ \relax
6:28 & \begin{tabularx}{0.7\textwidth}{X} 把頭放在盤子裡，拿來給那女孩，她就給她母親。 \end{tabularx} \\ \\ \relax
6:29 & \begin{tabularx}{0.7\textwidth}{X} 約翰的門徒聽到了，就來把他的屍體領去，放在墳墓裡。 \end{tabularx} \\ \\
[1ex]
\hline
\hline
\end{longtable}
$^{1}$今天選讀的經訓是記載馬可福音6章14-29節新約聖經第55頁.
是講述施震約翰的死 是他付代價的死.
耶穌的名聲傳揚出來 希律王聽見了就說.
施震迫約翰從死裡復活了 所以這些義能由他裡面發出來.
但別人說:是以利亞!又有人說:是先知正像先知中的一位.
希律聽見卻說:是我所斬的約翰 他復活了.
先是希律為他兄弟肥力的妻子希羅底的緣故 差人去拿住約翰所在艱雷.
因為希律已經娶了那婦人 約翰曾對希律說:.
你娶你兄弟的妻子是不合理的.
於是希羅底懷恨他 想要殺他 只是不能.
因為希律知道約翰是義人 是聖人 所以敬畏他 保護他.
聽他講論 受多召著行 並且樂意聽他.
有一天 恰巧是希律的生日 希律擺設筵席.
請了大臣和千夫長 並加莉莉作首領的.
希羅底的女兒進來跳舞 使希律和同席的人都歡喜.
皇就對女子說:你隨意向我求什麼 我必給你.
又對她喜誓說:隨你向我求什麼.
就算我國的一半 我也必給你.
她就出去對她母親說:我可以求什麼呢.
她母親說:施朕約翰的頭.
她就急忙進去見皇 求他說:.
我願皇立時把施朕約翰的頭 放在盤子裡給我.
皇就甚憂愁 因他所起的誓 又因同席的人就不肯推辭.
隨即差一個護衛兵 吩咐拿約翰的頭來.
護衛兵就去 在監裡斬了約翰 把頭放在盤子裡.
拿來給女子 女子就給他母親.
約翰的門徒聽見了 就來把他的屍首拿去 葬在墳墓裡.
(大型節目平安).
我們今天平安回到當中.
當我預備訊息的時候 我也很怕會有極端天氣.
令大家不能回來當中一起敬拜.
但今天天氣也很好 回來當中一起相聚.
把這份榮耀頌讚歸給天父.
我們這個星期承接上星期陳牧師開始了下半年的港道系列.
就是當耶穌差派他的門徒開始傳道的時候.
開始準備他們和實踐門徒訓練.
如何在繼續承接耶穌的工作中準備他們.
以致讓他們在新的心靈有一個充分的預備開始.
上星期陳牧師講的內容是第六章開始.
一開始的時候 講耶穌在家鄉傳道 但受鄙視.

$^{41}$人們說「我都知道你是誰」 「你的家人我都知道是誰」.
所以在家鄉所傳福音的工作不太順暢.
同時 耶穌也差派了十二個門徒出去傳道.
去講福音的訊息 去醫治 去趕鬼.
同樣也教導他們不要為自己預備太多.
因為一方面要準備被人接待.
第二方面就像上星期的講題一樣.
也要學習被拒絕.
因著福音的緣故 因著上帝的緣故.
可能也會被拒絕.
這就是上星期的內容.
而當他們出去傳道後 結果是怎樣呢?.
就是今天六章十四節前一點.
六章十二節說什麼呢?.
門徒就出去傳道 叫人悔改.
又趕了很多鬼出去.
用油抹了很多的病人 醫好了他們.
這就是接著上星期的內容.
發現這些門徒出去傳道的時候很有果效.
真的奉主的名去做了很多不同的事.
好像很成功地經歷了回來.
事實上應該第十三節就直接跳到第三十節.
如果你有聖經可以打開看.
三十節就是說什麼呢?.
使徒聚集在耶穌那裡.
將一切所做的事 所傳的道都告訴他.
因為這個段落就是直接說.
當他出去之後 做了很多神蹟 趕鬼.
然後就回到耶穌那裡述職.
告訴他 發生了這些事.
於是十二節到三十節.
其實不是很應該說今天說的內容.
施浸若汗 被人砍這個結局.
應該說什麼呢?.
就是說這段時間門徒出去.
做了什麼特別的事呢?.
有什麼神蹟很值得記述下來的?.
又或者有什麼事在裡面說了.
這群門徒經歷了什麼呢?.
又或者可能不說門徒了.

$^{81}$至少也說一下耶穌.
耶穌你把十二個門徒踹了出去.
那你不用做嗎?.
你只是叫了人出去就算了.
那你自己那段時間做了什麼呢?.
可能其實如果我們正路地想.
如果一個故事去講述的時候.
十二到三十節應該是說這些內容.
但似乎為什麼突然間說了另一些的東西呢?.
加插下去的意思是想說什麼呢?.
相信這個不是不小心的.
馬可寫馬可福音的時候.
不會把一些東西堆在一起.
就像我們小學那樣.
把一些有關的東西放在一起.
然後就交功課了.
馬可一定不是這樣.
所以他放這個段落下來.
一定有些東西想表達.
當你看回整個內容那種感覺的時候.
馬可好像想表達一件事.
就是說當這十二個門徒出去傳福音.
我們要小心.
不要以為已經是得勝了.
已經很順利了.
或者一切都好像童話故事一樣.
王子和公主就美滿生活下去.
很開心了.
當他講回這個段落的時候.
好像要告訴你.
是,門徒出去傳福音.
好像是一個小階段的勝利.
但同時那個反對上帝國道的訊息.
那個勢力沒有減少.
那份力量仍然在醞釀著.
那你以為門徒奉主的名出去的時候.
好像很順利嗎?.
是不是門徒好像成為了耶穌的開路者?.
不是,原來當他在做這些工作的時候.
黑暗的勢力仍然在醞釀著.

$^{121}$仍然在當中有一份張力.
同樣在這個並排裡面.
在信景裡面原來逆勢在醞釀著.
但當然在逆境裡面.
其實信景仍然在當中前行.
所以當這個加插下去的時候.
成為了一個很重要的補充.
不知道小時候你們有沒有看過.
男兒當入樽,足球小將那些.
明明起手射過三分球.
幾秒的時間.
可能要不進去,要不不進去.
但突然之間編輯就說.
突然飛回三年前.
明明那幾秒的時間又可以喊了幾集.
事實上都相似.
明明傳福音或者報喜訊回來.
還沒有跟耶穌述職.
為什麼中間要加插這麼大的段落.
用這麼大的篇幅去補充一些重要的資料.
因為這個真的很重要.
當補充這個資料之後.
當這班門徒回來跟耶穌述職.
又是另一個高峰.
就是五餅二魚的神蹟.
下個星期Raymond就會說.
五餅二魚,然後耶穌做的一些事情.
當這個高低起伏的時候.
正是馬可很想讓我們讀者.
當時的讀者去感受到.
原來傳福音不是只是求一味順暢.
又或者一味要受苦等等.
在這種高高低低的張力裡面.
正是我們每一個跟從主的門徒要經歷的.
當焦點回到今天的那段經文.
六章十四節至二十九節裡面.
一開始是因為門徒去傳福音.
然後以至耶穌的名聲.
傳揚出去做一個開始.
承接著門徒所做的工作.

$^{161}$當時的人怎樣看待耶穌這個人呢?.
在裡面他們有很多不同的聲音.
在裡面我們會看到.
唯獨是希律王很鐵定覺得.
耶穌的出現是因為施浸約翰.
很明顯他是有一個很堅持的.
同樣也當時解答了.
耶穌真正的開路者的先鋒.
施浸約翰怎樣壯烈犧牲.
所以這個插曲就像剛才所說.
明明門徒出去傳福音很開心.
一群興高采烈覺得有小成功的門徒.
好像突然被馬可淋一盆冷水.
突然說這麼開心的事.
突然說施浸約翰的結局.
由耶穌在家鄉不受愛戴.
到傳福音好像有些成果.
突然間又降落到一個.
好像很多的危機.
很多的殺害等等的事.
正是馬可正在建構剛才所說的.
這份張力.
不只是求順利.
也不只是在逆境裡面沒有出路.
有時候我們以為傳福音.
去做一些建立人工作.
或者生命的工作.
那個人一定某程度上要受愛戴.
或者很多人覺得你也不錯.
做這麼多好事.
無論是不是因為耶穌的名.
你做很多善事一定會受人愛戴.
但其實當你想這些表面成功的時候.
可能在現實裡面.
也有很多暗裡的攻擊.
會慢慢醞釀而且越來越尖銳.
因為這個勢力在這一切背後.
充滿著很多的苦毒和剛硬.
我之前在教會牧養之前.
我是在浸會大學服事.

$^{201}$浸會大學你也知道.
浸會大學已經是基督教名頭的一間大學.
浸會雖然有些同學真的不知道.
不知道為什麼叫浸會大學.
一個信仰背景的大學.
或者在校園裡面去接觸學生.
跟他們一起去玩.
一起去建立生命.
跟他們做輔導.
同行.
可能你也會覺得同學應該也有一點尊敬.
或者覺得你們做的事情挺好.
但有幾個片段我很深刻.
就是同親去迎新的時候.
在學校的背景裡面.
除了你們這些基督徒.
這些傳道牧者.
在做福音的工作和建立生命的工作.
但其實同時間還有很多派對.
不同的人在這裡.
有些可能是學生會.
有些是不同的興趣小組.
同樣的.
當在九月開學的時候.
不知道有沒有一些升大學的.
當你九月的時候.
你一去註冊日.
就好像大家要搶人一樣.
搶什麼人呢?.
就是要大家參與你的活動.
和參與當中一些的迎新活動.
去到你的團契的迎新活動.
可能就會撞上學生會.
撞上其他興趣小組的時間.
他就去不了.
我記得有一次是什麼呢?.
在campus裡面正在做迎新.
有一班同學聽完註冊信息出來.
然後有一個另一邊的.
當然不跟大家說是什麼背景.

$^{241}$總之有一個學生就很大聲說.
你不要去那些邪教活動.
來我們這些.
那些很恐怖的.
我心想邪教?.
浸信會大學.
我們是浸信會裡面的一個department.
你這麼大聲說我們是邪教活動.
但他不知道自己在說什麼.
我相信那一刻.
我都不會覺得他很冒犯.
我只是覺得很可憐.
你知道自己在說什麼嗎?.
那你會不會跟大家說.
你不要進來這間邪教大學讀書.
這樣吧.
所以你會更加明白.
當你正在做這些的時候.
原來真的有很多人在誤解.
不知道你在做什麼.
以為你在辦教會的時候.
會有很大的張力.
所以當你在做form的活動的時候.
不要以為每個人都很愛戴你.
或者很明白你在做什麼.
另一邊廂當然不只是大學的背景.
有時候你聽一些遠目的分享也是一樣.
可能你來做很多關懷生命.
在臨終等等.
跟生命同行的時候.
可能你以為大家一定會方便一下.
或者會幫助你.
不是的.
可能有些醫護會覺得.
你不要妨礙著.
你不要弄得很麻煩.
所以原來今天.
當我們在做很多.
可能很成功做到一些生命的工作.
或者傳科音的時候.

$^{281}$其實背後原來也有著這些危機.
或者這些張力在裡面.
跟這個世界原來有這種拉扯.
如果我們只是沉醉在.
做到這些成功.
或者這些事的時候.
沉醉在這些成功的得著裡面的時候.
如果意識不到原來也有這些阻力.
或者這些困難的時候.
當這些挑戰成為一種強勢.
就好像希律王接下來要做的事情一樣.
我明明在做這些這麼有意義的事情.
為什麼你要這樣說我.
為什麼我做這些這麼好的事情.
你要逼迫我.
如果沒有這個準備的時候.
很容易陷入一種失望.
或者甚至失敗.
覺得我還是不做這些事情了.
而當繼續下去這個段落的時候.
雖然耶穌絕對不是這次故事的中心人物.
但是我們可以看到.
馬可藉著這個故事的細節.
在暗示將來耶穌也會面對這些的迫害.
這些的審判.
這些的定罪.
暗示將來耶穌也會面對著這個結局.
同樣當耶穌也面對這個結局.
他這一班背後的門徒.
也同樣將要面對這些事情.
正所謂學生不能夠高過先生.
耶穌也要受迫害的時候.
你覺得門徒可以不用受這些迫害嗎?.
甚至他們也要預備這個殉道的可能.
當看回這個希律王.
我們會知道這個希律王.
如果你看回歷史上.
真的會找到的.
不是傳說不是故事.
這個希律王全名叫做安提柏.

$^{321}$希律安提柏.
是主前四年至到主後三十九年.
統治著加利利和比利亞這個地方.
是大希律王的兒子.
在裡面其實他不是當時整個以色列地的王.
他只是一個分封王.
當大希律過世之後.
他只是處理著這幾個小地方的一個分封王.
當希律知道耶穌的事情的時候.
我不知道大家是否感受到他的那種肯定.
你這樣說不是.
他這樣說不是.
總之我知道耶穌是誰.
是怎樣來的.
他心裡有數.
從經文裡面所說.
其實我們會知道.
他很多人的身邊.
都對耶穌的行神跡.
傳福音有不同的看法.
並且將這個看法.
有些人就將他和施浸約翰的死.
再一次連繫上來.
當你看第十四節.
就說耶穌的名傳揚出去.
希律王聽見了就說.
本身有些譯本不是說就說的.
是說希律王聽見了.
就有人說施浸約翰是從死裡復活.
所以這些異能又許出.
然後十五節就說.
另一些人就說是以利亞.
又有人說是先知.
但是十六節.
如果用那種譯本的時候.
就更加凸顯到.
希律王堅定的聽見.
就說這是我所斬的約翰.
他復活了.
意思是什麼呢?.

$^{361}$有人說是施浸約翰從死裡復活.
有人說是以利亞.
有人說是先知.
但是希律全部都不理.
或者很肯定.
只知道一個原因.
就是這是他所斬的約翰.
為什麼會這樣想呢?.
因為可以看得出.
當時其實耶穌在做的事情.
對當時的人來說.
有很多不同的看法.
所以馬可更加需要.
用這個篇幅來解釋.
究竟耶穌的名聲.
傳揚出去是為了什麼呢?.
原來當時有極多的誤解.
和極多的雜質在裡面.
連當時的人.
一些領袖都不太知道.
究竟耶穌雖然名聲傳開了.
但是是因什麼緣故呢?.
當繼續看下去的時候.
你會看到.
希律王聽到耶穌的名聲.
但是他不知道是什麼內容.
他也沒有特別說.
聽到他行神蹟.
聽到他的醫治.
等等的事情.
好像很模糊地知道.
耶穌在做的.
是一些很大的事情.
令很多人回想.
而那個名聲的詞.
不只是說名字.
耶穌的名字很響亮.
耶穌的名字或者什麼.
耶穌的名聲其實包含了很多東西.
他的為人.

$^{401}$他說的話.
他在做的事情.
也包括了奉他的名.
所做的事情的人.
不是只是說耶穌的名字這麼簡單.
是包含了這個組織.
裡面所發生的一切事情.
而希律很堅定地覺得.
這一切的名聲背後.
蘊含著的只是一個原因.
就是當日.
他斬了的施浸若寒的原因.
他不理其他事情.
他不理是不是以利亞.
不理是不是以前的尼采.
或者一直在說的.
先知所說的事情.
他堅定地覺得.
這跟昔日他做過這件事情的.
施浸若寒有關.
因為在他心底裡.
他很清楚這件事情.
如果真的有一個神在的時候.
這個位置一定是他不會放過我的地方.
這個位置一定是他知道.
他人生裡面最大得罪上帝的罪行.
所以當有一個說.
奉神命來的人的時候.
他立刻就想起.
是的 當年做了這件事情.
一定是這個人要來追討.
或者要在當中.
去回應他做了的這件事.
很明顯就是要顯示出.
希律當時正在受良心責備.
很害怕.
正如這個世界.
我也不明白那個學生.
為什麼會說我們是邪教.
但一定是有些東西.

$^{441}$隱藏在他裡面.
收藏在裡面.
當有一群基督徒.
或者一群奉主命.
奉神命來的人的時候.
就立刻想起.
一定是這樣.
一定是為了那件事而來.
所以在裡面.
他就有這個肯定的話.
第十六節你看到他就說.
他卻說是我所斬的約翰.
他清楚知道自己做了的事情.
現在他復活了.
當時復活這個訊息.
不是普遍的訊息.
但是在裡面.
他能夠肯定.
這個約翰復活了.
究竟約翰做了什麼.
最後段落就告訴我們.
十七至二十節就告訴我們.
希律王和施浸約翰的關係.
很微妙的.
你會看到.
在當中他好像很不喜歡他.
但又收監困住他.
但又好像很想聽他說.
關於上帝的道等等的事情.
但在這個關係裡面.
又加插了第十九節裡面.
希羅底對約翰的那種仇視的態度.
你就會看到在這種醞釀裡面.
不祥的事情就即將發生了.
而當中提到希羅底.
其實你會想起.
救了一個很重要的人.
當你想到約翰的時候.
就想到以利亞.
這個希羅底很明顯.

$^{481}$就是你會想到那個耶駒別.
而這個希律王.
就會更加想到當年的阿哈.
你就想一個黑心的王赫奧.
一個沒用的皇帝.
加一個至死忠心的神僕.
很明顯這個組合.
就是在約翰的身上出現了.
這個故事的背後.
其實也挺錯綜複雜的.
如果你看回馬可說的話.
就比較簡單一點.
總之就是你娶了你兄弟的老婆.
所以不合神的心意.
但當你找回一點背景的時候.
特別是你看回約瑟夫的猶太故事.
成為一個旁證.
你去想一想.
就更加明白到那種.
恐怖,張力和複雜性在哪裡.
其實在哪裡呢?.
就是原來希羅底.
這裡聖經馬可也有說.
第一任的丈夫其實是肥力.
然後你就想.
難道希律王娶了他兄弟的太太.
但那個兄弟還沒死.
如果他說他死了.
承接下來就說.
但他還沒死.
而那個希律王.
其實也不是第一時間.
娶希羅底做太太.
原來你看回歷史.
希律王本身是娶了阿拉伯王的一個女兒.
去做老婆的.
但後來當他娶了阿拉伯王的女兒.
做老婆之後.
他回去羅馬.
探望他這個同父異母的兄弟.

$^{521}$就是肥力.
當他探望的時候.
竟然愛上了他這個同父異母的兄弟的太太.
接著這個太太同樣也不是野小.
你會看到她除了黑心之外.
其實她也有很多想法.
很多野心.
兩人一拍二合.
然後這個希律王安提帕.
就將希羅底帶回加利利.
就是搶了別人的老婆.
同樣搶了別人的老婆之後.
就休了自己的老婆.
就是將阿拉伯王的女兒休了.
將她送回阿拉伯.
正式娶希羅底為妻.
希羅底離棄了丈夫.
離婚.
然後就和希律王同住.
阿拉伯王的女兒也不甘受辱.
回到阿拉伯之後就發展了一場戰爭.
以致當時希律王要和阿拉伯去打仗.
打仗打到一段時間.
希律王損失慘重.
接著他就向羅馬皇帝求救.
要羅馬皇帝下一個命令.
壓住了阿拉伯王.
他才免於滅國.
所以是很錯綜複雜的背景.
馬可當然不特別去講詳細.
總之他知道這兩個人所做的事.
是多麼亂 當中多麼錯.
而希律王和希羅底的關係.
其實當中為什麼會牽涉到斯贈約翰呢?.
如果他們只是一群皇帝.
很荒淫很亂.
就由他們去就算了.
但問題是希律王又要管住加利利.
管住加利利的時候.
又為了討好猶太人.

$^{561}$他又說自己是一個猶太人.
很相信你們那一套.
讓我成為猶太人.
又自稱是一個猶太人.
你又要自稱.
當然又要有義務.
去遵守上帝的律法.
所以斯贈約翰自然就要糾正他.
出來告訴他 你不可以這樣.
在馬太科學裡面.
不是只說了一次.
其實約翰是多次跟希律王說過這個問題.
所以在糾正裡面.
他就是要讓他們知道.
你既然又說要成為我們猶太地的分封王.
又要自稱是猶太人的時候.
你不要忘記你自己在做的事.
你要悔改.
你要在當中去改變過來.
但當然結局你會看到.
他們完全相反.
一方面就抓了約翰.
一方面當然不希望這個出名的先知.
繼續在公開場合去批評他.
但另一方面你又會看到.
在聖經裡面說他是想保護他.
為什麼要保護他呢?.
你想想你最後又斬了人.
一開始又抓了人.
為什麼要保護他呢?.
原來你會看到.
他除了保護他.
其實是他怕.
怕一件事就是如果他.
在公開的場合繼續這樣說的時候.
無論希羅弟.
無論其他的人.
將他殺死.
一殺死就大件事了.
因為很多人知道他是先知.

$^{601}$所以他就要將他收在監獄裡.
覺得困住他.
好像有一個地方保護著他.
然後繼續讓他在當中.
又同時間去聽上帝的話.
其實你會看到希羅王.
心裡面很亂.
又好像不太想聽他說.
又好像在裡面不想繼續說.
但當然你會看到希律.
不像希羅弟那麼決絕.
剛才說黑心的皇后.
希律王不是想置約翰於死地的決定.
相反他經常會走回監獄裡.
去聽約翰說道.
心裡面其實又很害怕.
但又很想知道神的心意是什麼.
所以常常處於猶豫不定的狀況.
同樣在馬太福音.
在另一個福音書裡也有說到.
希律王不殺約翰的另一個原因.
就是因為他害怕百姓.
因為他是一個先知.
所以這種三心兩意.
是很符合這個皇的特質.
也刻意突出了他這種張力.
很想跟從神.
但現實又不想.
生活上又不想跟從.
但又知道道是有他的意思.
馬可突出了這個狀況.
很想告訴我們.
如果像希律王那樣.
只是聽道.
又很想兼容世俗.
又很想聽聖人約翰說的話.
是不可以的.
原來在當時的屬靈真善.
百業化的時候.
每個人都要有他的立場.

$^{641}$即使你是王也好.
那種猶豫不決.
不能讓你面對屬靈的挑戰.
故事繼續發展下去.
就像希律那樣.
想著將這種張力.
不理就算了.
總之我又有聽道.
保護一下他就算了.
事實上希羅底沒有讓他有好教訓.
故事繼續發展下去.
有一天他的女兒在跳舞的時候.
向希律王說.
跳完舞就說獎勵你.
然後就問媽媽.
媽媽就說要施浸約翰的頭.
你會看到這個女子.
都是很重要的一個角色.
完全將剛才那種張力突破了.
你既然在膠著.
總之你太太和女兒幫你決定了.
幫你決定了你的立場.
毫無疑問.
當時的希律王.
都不知道在半醉半醒的狀態下.
做了這個決定.
就答應了他的女兒的要求.
當然一方面告訴她.
你不要以為王很厲害.
你看看一個王.
你以為至高無上.
個個都尊敬.
雖然是一個分封王.
但原來這個王.
都不是至高無上.
都不是沒有人可以阻止到他.
原來在當時的榮辱文化.
和當時人的罪性裡面.
即使你去到王.
你不要以為自己所向無敵.

$^{681}$不要以為你成為王的時候.
什麼都可以做.
名義上是.
但你看看這個技術裡面的希律王.
多麼的被動.
多麼的沒用.
你想想一個王.
他現在竟然在這種張力下.
屈服了下來.
希律王當然你會看到.
他說就說.
你跳完舞我獎你.
將我半個國給你都可以.
吹噓的.
他是分封王.
你以為他真的可以給一個王.
他也是在一個大的王裡.
他可以給多少人.
所以你看到這些承諾.
沒有什麼份量.
但當他要賜下這個獎勵的時候.
他已經說了.
他以為自己是一個王.
我有什麼做不到.
他就是做一件事.
令他也很疑惑.
他竟然什麼都不要.
要施浸約翰的頭.
當時他也很糾結.
但礙於剛才所說.
當時的榮辱文化.
當時人的罪性.
他也只好屈服.
被動地做了這個決定.
這個章節.
如果你想到那種張力和那種狀況.
其實你很容易可以跳下去.
在獵王技上.
剛才說了.
一個沒用的皇帝.

$^{721}$一個黑心的皇后.
和一個先知.
一個忠心的僕人的張力.
獵王技上21章.
如果你有聖經或者寫下來.
回去可以看一下.
是發生了一件什麼事呢.
就是當時阿哈王.
他在以色列地作王.
有一天他作王的時候.
看到旁邊有一個葡萄園.
是一個屬於拿伯的葡萄園.
他很正路.
一個皇宮旁邊有一個葡萄園.
跟他說.
你可不可以把你的地賣給我.
我會給你工價.
你給我一些葡萄.
我屬於他的直銷供應商.
你有葡萄就直接運給我.
我又會給你錢.
做這個交易.
誰知道拿伯就跟他說.
雖然我們很近.
但是因為我敬畏神.
我不可以把我先人.
即是祖先留給我們.
就這樣賣給你.
所以你想想一個王.
想著這樣做.
做不到.
那你自己就再想想.
無論是交易上.
甚至強硬去做.
這個阿哈又不是.
什麼都沒有做.
就回到皇宮.
回到皇宮裡面.
就吃不安坐不下.
就被他老婆.

$^{761}$這個.
等一下.
爺洗別.
沒錯.
突然間卡住了.
爺洗別發現了.
爺洗別就看到他就說.
為什麼你吃不安坐不下.
然後阿哈就說.
因為我向這個拿伯要地.
他不給我.
他說不賣地給我.
所以我很不開心.
然後連皇后爺洗別.
都跟阿哈說.
你是智利以色列的一國之君.
你想辦法處理他.
你現在不要不開心.
總之我會幫你搞定他.
那怎樣搞定這個地.
總之有一天.
他就找兩個匪類.
就是找兩個強悍的人.
就邀請拿伯去一個聚會裡面.
在那個聚會裡面.
突然間那兩個壞人起來就說.
這個拿伯蝕族臣.
然後又說要推翻那個王.
然後拿伯就被人拉了出去.
用石頭扔死了.
所以那個地就變成沒有了主人.
阿哈同時間就可以順利地吞併了他的地方.
你會看到這個王真的很沒用.
他只是想要那塊地.
然後吃不安坐不下.
就要太太想辦法幫他處理.
當時這件事不要以為沒有人知道.
上帝一定知道.
所以就興起了當時的先知.
以利亞就來到跟阿哈說.

$^{801}$在裡面他就說.
耶和華如此說.
第十九節說.
你殺了人又得了他的產業嗎.
所以既然你做了這件事.
狗在哪裡舔拿伯的血.
也在哪裡舔你的血.
將這個懲罰審判到阿哈.
怎知道阿哈在裡面就知道.
糟了 我得罪了神.
你猜他有沒有回轉.
一個這樣的王.
他有回轉.
有跟神說我得罪了你.
我不對等等的事.
在第十五節裡面有一個小小的段落.
去講述阿哈究竟是一個怎樣的王.
二十五節說什麼呢.
從來沒有像阿哈的因他自賣.
恨耶和華眼中看為惡的事.
受了皇后耶駛別的聳動.
在裡面你會看到這個王是壞的.
但是他壞的地方就是.
他完全是被人操控.
在裡面被耶駛別影響了他.
然後他就撕裂衣服又禁食等等的事.
你會看到這個王.
好像剛才說的希律一樣.
又會聽神的話.
但是又會被他的皇后聳動.
影響了又會做了一些壞的事.
在裡面讓我們更加看到.
這一切的事.
想表明這個施浸若寒.
就是那個為彌賽亞修植路的那個以利亞.
當這個以利亞來到的時候.
又不同以前舊約那個.
以利亞最後就安然無事.
就面對著這個耶駛別的攻擊等等的事.
最後當然以利亞看不到他的終局.

$^{841}$但是也知道他也不會有一個好的結局.
但是在耶穌時期的施浸若寒的結局.
不能夠免於死刑.
以致在裡面我們會看到.
甚或乎施浸若寒的死因.
是多麼無稽多麼無奈.
當施浸若寒的生命結束了.
他的門徒知道的時候.
就將他的屍首領去埋葬了.
就沒有再出現.
就沒有下文了.
你要看到若寒他好像當時的以利亞一樣.
他講出人的問題.
但是因著人的罪性.
人的那些無謂的榮辱感.
遭受殘忍的對待.
這樣看來原來地上的這個權勢.
沒有因為人去傳福音.
或者去做這些事好像有改變.
甚或乎反過來.
人在這些權勢面前很無力.
甚至好像看起來很失敗.
而且從施浸若寒的結局.
我們都看到.
不是像以利亞一樣有石破天驚的結局.
所以馬可更加在這個敘述裡面.
凸顯了當我們要去傳福音.
去按著上帝心意的時候.
去做的時候.
這些敵黨的權勢仍然存在.
神的角度看起來好像很渺小,很無助.
但是我們知道若寒既然是為了彌賽亞所預備的以利亞的時候.
而耶穌就是那個彌賽亞.
耶穌自己或者他跟從者.
縱然是彌賽亞.
但是同樣都要面對著這些權勢的壓迫.
但是若寒和耶穌的分別在於什麼呢?.
若寒的門徒很好.
把老師的身體埋葬.
但是耶穌的門徒就全走了.

$^{881}$當釘完十字架就全走了.
另外更大的分別是什麼呢?.
若寒最終死在墳墓當中.
但是耶穌沒有被墳墓所囚禁.
三日後升天.
這一份就是當時對門徒來說最大的力量和盼望.
在當我仰望這個力量和盼望之前.
他們要去傳福音的時候.
一定要好像剛才一樣.
我們整全明白到.
跟從耶穌.
作為耶穌工作的時候.
要付的代價.
這個代價不是一定是很正路那種.
我出來為神說一些話.
去傳福音.
然後你迫不及待我.
我已經預料到.
有些的迫害.
有些的困擾.
不是你所想的那麼正路.
就好像若寒一樣.
原來我死是因為那個女孩跳完舞就斬了我.
我死得多麼冤枉.
原來這種張力是必然的.
這種生命.
原來在神的手中的時候.
是會經歷這一切的事.
所以馬可福音14-29節.
就是要補充這個插曲.
不是無關痛癢或者寫錯東西進去.
而是要預表著耶穌和門徒.
將來都同樣好像若寒一樣.
要面對著恩著這個上帝的工作.
面對著這些受難.
這些挑戰.
其實正正就是我們的使命.
和我們要付上的代價.
當我們繼續去看的時候.
我們會看到.

$^{921}$這一班跟從主的門徒.
做完這些事的時候.
他們要深刻的知道一件事.
就是這個代價將來要去付上.
為什麼我們在裡面可以看得出呢.
因為我們知道耶穌不是只是要那些.
聽到知道有神蹟.
或者因為有很多好事而走過來的.
這些人作門徒.
這些人和希律沒有什麼分別.
而真的作門徒的人.
是要遠超過這個只是淺層的.
很基本的表現.
如果要做門徒的話.
就好像這裡所說.
我們一定要跟從耶穌到底.
背上這個十字架.
去承受作門徒的後果.
而施浸若寒也不只是一個先鋒.
更加是作門徒的榜樣.
如何付代價走上.
同樣地.
當我們知道這個日子.
這個可能將要去面對的時候.
我們更加要預備.
這個付代價的心智是什麼.
當我們今天認識了上帝.
是很美好的.
可能是很多的祝福.
在很多的恩典裡面.
但是如果我們今天的基督教.
或者信主的路.
傾向著很多的神蹟.
很多的憤興.
很多的德性.
很多的這些好處的時候.
不是不好的.
也好的.
一開始的時候是很吸引人的.
又或者是很有效讓人家.

$^{961}$信主多好.
信耶穌是很好的事.
但是單單靠這件事.
不能夠應付我們生活裡面.
很多的挑戰.
回到真實的層面裡面的時候.
我們的人性.
我們的問題.
都沒有解決.
結果會是怎樣呢?.
如果不是我們在裡面.
只是高舉著這種德性.
或者什麼的時候.
我們就會開始發現.
你不開心嗎?.
你不屬靈.
你不德性嗎?.
你不夠好.
一定是你有這些問題.
甚或乎可能就會更加多的生活規範.
就是覺得.
你總之信耶穌.
或者等等.
就OK的了.
就會限制了很多生命力.
最終可能很多的信徒.
就好像經歷不了神話.
可能這個信仰都不是真實的.
那就走了.
又或者這個群體.
只是不斷靠很多新的人進來.
去壯大這個群體.
而這些的成功.
只是表面的.
將一些失敗.
或者這些軟弱的張力.
蓋在一邊.
但是馬可.
很明顯就是要糾正這個觀念.
跟從耶穌.

$^{1001}$一開始的時候.
門徒很開心.
做到很多事情.
但是當你繼續跟從耶穌的時候.
要好像施浸若寒.
好像耶穌一樣.
要活在這種張力下.
經常說這種張力.
在說的是什麼呢?.
就是那個代價.
你不要以為付出了什麼.
做了什麼.
就是這樣.
那種張力是令你困在裡面.
覺得好像做到一些事情.
又好像做不到.
好像表面成功.
但是又很失敗.
但當我以為很失敗的時候.
原來又是成功的.
這種就是信仰張力的代價.
以至在裡面.
你要知道你靠的是誰的大能.
以至我們今天.
能夠為耶穌獻上什麼呢?.
獻上這份心智.
如何跟從真理去生活.
當這個信息.
一直在我心中去想的時候.
我都很困惑.
今天我在傳道.
在傳福音.
活在信仰裡面.
在付出什麼代價呢?.
昨晚在開小組之前.
在預備的時候.
在路上走在街上.
一直在想.
如何分享這個內容呢?.
在兩街走的時候.

$^{1041}$覺得人來人往.
大家都是生活著.
大家都有不同的壓力.
或者有不同的事情.
究竟今天當我們說.
要付上代價.
生活在這種張力裡面的時候.
我應該要怎麼做?.
到底要怎麼做呢?.
也許當你落在這個世界裡面.
你知道的.
這個世界看的東西.
可能是很膚淺的.
更加是現在這個年代.
電話那些.
看的東西就是這樣.
甚至現在的AI都騙人.
有些什麼穿越時空的片.
都是很真實的.
當我們今天.
要去傳福音.
報喜訊.
去領人歸主的時候.
其實可能很多人會誤解我們.
你做這些是為了什麼呢?.
有什麼利益?.
有什麼好處?.
或者等等的事情.
其實我們可能已經.
被人標籤了.
或者被人說了一些話.
這個會不會是我們.
那種張力的代價呢?.
別人這樣說你的時候.
你會說不是啊.
還是仍然默默去做呢?.
今天世界的荒誕.
大家看新聞都會看到很多.
或者有一天.
可能我只是說了.

$^{1081}$我覺得我的立場.
因著信仰.
因著聖經.
因著耶穌教的時候.
就會遭受一切的對待.
有沒有這些準備呢?.
有沒有預備這一切呢?.
又或者.
可能你覺得.
我已經為上帝付代價了.
做很多事情.
用很多時間和心力.
是好的.
但是我昨天去走的時候.
想起一件事.
想起過往有一個.
服侍的片段.
最後一個分享就是.
當我有一天.
跟一群小朋友.
全部都是四五八歲.
一起去有一個活動.
跟他們分享一些情緒.
這個就是開心.
這個就是傷心.
這個就是生氣.
這個就是害怕.
四個很簡單的情緒.
跟幾個小朋友分享.
你有沒有試過開心.
有沒有試過不開心.
有沒有試過生氣.
我問有沒有試過生氣.
他們很乖.
沒有啊.
沒有生氣過.
我說不要緊.
你們沒有生氣過.
有沒有見過別人生氣.
有一個就說.

$^{1121}$有啊.
我見到媽媽生氣.
因為什麼呢?.
因為我做功課的時候.
字寫得不漂亮.
媽媽很生氣.
我說是嗎?.
生氣你會怎樣做?.
跟他們一起分享.
另一個.
我也見過生氣.
我見過媽媽生氣.
我說又是因為什麼事?.
因為爸爸偷了媽媽的錢.
媽媽沒有錢生活.
很生氣.
那一刻我呆了.
當然因為那個背景.
是一個無家者的家庭.
今天我們以為自己.
為神付了很多代價.
在做很多事.
但原來對有些家庭來說.
他本身的生活就是這樣.
當你以為為主.
去小賜旅行.
為主做了這些事.
一個家庭.
一個家庭裡面.
就是已經是這樣.
你以為付了很大的代價嗎?.
今天的代價是什麼?.
我們為主做了什麼?.
當你以為自己付了這個的時候.
其實上帝怎樣去看?.
同樣.
當我們要去付上這一切的時候.
是正常的.
這種張力是當你.
為主做這些事.

$^{1161}$未必有一個好的結果.
但這個代價.
上帝知道.
耶穌也知道.
以致今天.
當我們付上這個前行的時候.
求主加力給我們.
我們一起祈禱吧.
親愛的主耶穌.
求你帶領我們.
當我們願意去跟從你的時候.
我們知道.
這是美好的前路.
這段路一定是上帝你所祝福.
但我們知道.
踏上這條路.
我們要付代價.
這個代價不是我們為你.
放下了什麼.
這個代價是我們生命裡的張力.
我們以為前行嗎?.
但原來是停滯.
以為我們不能為主做什麼的時候.
卻是主賜下大能.
讓我們前行.
所以求主繼續引領我們.
讓我們在你的同在.
你的帶領下.
願意傾出一切.
去跟從主.
以致我們今天.
回到生活裡的時候.
我們知道如何為主你做.
為主你作見證.
為主你得著更多的靈魂.
求主你帶領我們.
我們祈禱.
奉主耶穌基督的名字而求.
阿們.
\newpage



\section{馬可福音 6:30-44}
\label{sec:SbKqVXHzrXE}
\textbf{2025年7月20日崇拜直播｜陳冠成牧師｜耶穌給門徒的十課－給眾人吃飽｜馬可福音6\_30-44}
\newline
\newline
連結: \href{https://youtube.com/watch?v=SbKqVXHzrXE}{\texttt{https://youtube.com/watch?v=SbKqVXHzrXE}} ~~~~ 語音日期: 2025-07-21
\newline
\newline
\hyperref[sec:OaFcCMZ21q4]{\small{< < < PREV SERMON < < <}}
~
\hyperref[sec:index_chronic]{\small{[返順時目]}}
~
\hyperref[sec:gyB32xPImr0]{\small{> > > NEXT SERMON > > >}}
\newline
\newline
馬可福音 6:30-44
\newline
\begin{longtable}{cl}
\hline
\hline
章節 & 經文 (和合本修訂版)\\
\hline
6:30 & \begin{tabularx}{0.7\textwidth}{X} 使徒們聚集到耶穌那裡，把一切所做的事、所傳的道全告訴他。 \end{tabularx} \\ \\ \relax
6:31 & \begin{tabularx}{0.7\textwidth}{X} 他就說：「你們來，同我私下到荒野的地方去歇一歇。」這是因為來往的人多，他們連吃飯的時間也沒有。 \end{tabularx} \\ \\ \relax
6:32 & \begin{tabularx}{0.7\textwidth}{X} 他們就坐船，私下往荒野的地方去。 \end{tabularx} \\ \\ \relax
6:33 & \begin{tabularx}{0.7\textwidth}{X} 眾人看見他們走了，有許多認識他們的，就從各城步行，一同跑到那裡，比他們先趕到了。 \end{tabularx} \\ \\ \relax
6:34 & \begin{tabularx}{0.7\textwidth}{X} 耶穌出來，見有一大群的人，就憐憫他們，因為他們如同羊沒有牧人一般，於是開始教導他們許多事。 \end{tabularx} \\ \\ \relax
6:35 & \begin{tabularx}{0.7\textwidth}{X} 天已經很晚，門徒進前來，說：「這地方偏僻，而且天已經很晚了， \end{tabularx} \\ \\ \relax
6:36 & \begin{tabularx}{0.7\textwidth}{X} 請叫眾人散去，他們好往四面的鄉鎮村莊去，自己買些東西吃。」 \end{tabularx} \\ \\ \relax
6:37 & \begin{tabularx}{0.7\textwidth}{X} 耶穌回答他們說：「你們給他們吃吧！」門徒對他說：「我們要拿兩百個銀幣去買餅給他們吃嗎？」 \end{tabularx} \\ \\ \relax
6:38 & \begin{tabularx}{0.7\textwidth}{X} 耶穌說：「你們有多少餅？去看看。」他們知道後就說：「有五個，還有兩條魚。」 \end{tabularx} \\ \\ \relax
6:39 & \begin{tabularx}{0.7\textwidth}{X} 耶穌吩咐他們，叫眾人一組一組地坐在青草地上。 \end{tabularx} \\ \\ \relax
6:40 & \begin{tabularx}{0.7\textwidth}{X} 眾人就一群一群地坐下，有一百的，有五十的。 \end{tabularx} \\ \\ \relax
6:41 & \begin{tabularx}{0.7\textwidth}{X} 耶穌拿著這五個餅和兩條魚，望著天祝福，擘開餅，遞給門徒，擺在眾人面前，也把那兩條魚分給眾人。 \end{tabularx} \\ \\ \relax
6:42 & \begin{tabularx}{0.7\textwidth}{X} 他們都吃，並且吃飽了。 \end{tabularx} \\ \\ \relax
6:43 & \begin{tabularx}{0.7\textwidth}{X} 門徒把餅和魚的碎屑收拾起來，裝滿了十二個籃子。 \end{tabularx} \\ \\ \relax
6:44 & \begin{tabularx}{0.7\textwidth}{X} 吃餅的男人共有五千。 \end{tabularx} \\ \\
[1ex]
\hline
\hline
\end{longtable}
$^{1}$今日經訓記載在《馬可福音》.
六章三十至四十字.
有關耶穌給五千人吃飽的神蹟.
請聽神的話.
「使徒聚集到耶穌那裡.
將一切所作的事.
所傳的道傳告訴他.
他就說.
『你們來.
給我暗暗的到曠野地方去掀一掀.
這是因為來往的人多.
他們連吃飯也沒有功夫.
他們就坐船.
暗暗的往曠野地方去.
眾人看見他們去.
有許多認識他們的.
就從各城步行.
一同跑到那裡.
被他們先趕到了.
耶穌出來.
見有許多的人.
就憐憫他們.
因為他們如同羊沒有牧人一般.
於是開口教訓他們許多道理.
天已經晚了.
門徒遵前來說.
『這是野地.
天已經晚了.
請叫眾人散開.
他們好往四面鄉村裡去.
自己買什麼吃』.
耶穌回答說.
『你們給他們吃飽.
他們就吃飽』.
門徒說.
『我們可以去買.
二十兩銀子的餅給他們吃嗎?』.
耶穌說.
『你們有多少餅可以去看看?』.
他們知道了.

$^{41}$就說『五個餅,兩條魚』.
耶穌吩咐他們.
叫眾人一幫一幫的坐在青草地上.
眾人就一排一排的坐下.
有一百一排的.
有五十一排的.
耶穌來著這五個餅,兩條魚.
望著天祝福.
抹開餅.
第二級門徒擺在眾人面前.
也擺那兩條魚分給眾人.
他們都吃.
並且吃飽了.
門徒就把碎餅碎魚收拾起來.
裝滿了十二個籃子.
吃餅的男人共有五千.
聖經獨筆.
弟兄姊妹平安.
五餅二魚這個故事.
相信大家都知道.
是耶穌基督所行的眾多神蹟裡.
唯一一個同時收錄在四卷的福音書裡.
反映這個故事對每一個福音書作者來說.
或對初期教會來說.
是一個極大的意義.
事實上你會留意到.
福音書作者有意引導.
每一個看這個故事的人.
他們很不其然地將.
耶穌在曠野餵飽五千人的神蹟.
聯想到上帝昔日.
同樣在曠野透過赤下嗎哪.
來餵養群眾這一幕的景象.
顯示出耶穌基督就是上帝.
所差派來的那一位大牧人.
要將生命的糧賜給每一位百姓.
至於五餅二魚裡同時記載著.
耶穌如何使用那五個餅.
拿起餅內,捉藉料,.
打開遞給門徒.

$^{81}$亦很自然地讓每一個讀者聯想到.
在初期教會裡.
他們每一次手煮餐的時候的景象.
所以你會看到.
五餅二魚這個神蹟對初期教會來說.
特別是那些昔日沒有親眼見過耶穌基督.
未曾和祂一起生活過.
經歷過的信徒來說.
其實別具意義.
提醒他們.
耶穌就是上帝所差派來.
那位世間的終極好牧人.
又是那位的救主彌賽亞.
並且將來祂要是那位天國賢直的主人.
幫助每一位信徒.
藉著這個盼望.
他們可以有力量面對當時生活中的種種挑戰.
或者在信仰中所帶來的衝擊.
而對今天的我們而言.
五餅二魚這個故事.
確實我們已經聽過很多遍.
耳熟能詳可以說.
但是上帝的話語.
我相信今天仍然要啟迪我們.
所以我們一起去看.
今天馬可他怎樣記載五餅二魚這個故事.
在他的福音書裡.
你會看到有一個背景.
就是耶穌基督在差派.
十二個門徒.
兩個兩個出去傳道趕鬼.
然後他們聚集了.
一起回到耶穌基督的面前.
回報他們所工作的一切.
耶穌也很體恤.
這一班門徒都辛苦了一整天.
非常疲倦.
經文裡也記載他們連飯都未吃.
所以他就吩咐門徒我們一起上船.
退到曠野去.

$^{121}$打算去揭一下.
不過耶穌這個計劃也落空了.
事實上他和門徒.
無論是傳道治病趕鬼的那一刻.
大眾都有目共睹.
令到可以說耶穌和他十二個門徒.
人氣都飆升.
他們的一舉一動.
很自然也受到當時普羅大眾的注意.
群眾看到耶穌他們上船離開了.
於是他們就像現在追星一樣.
他們就從陸路來追蹤著耶穌.
我們說通常粉絲行動迅速.
一定是被追捧的明星.
往往早到一步現場.
你看到當耶穌在船上出來的時候.
他就看到原來廣大的群眾.
已經在曠野恭候他們.
事實上耶穌本來是打算退休的.
他可以有充足的理由我們上船.
我們去另一個地方避開這班群眾.
不過經文記載.
耶穌卻是憐憫這班他所見到的百姓.
馬科音很特別.
他給其餘福音書記載了一些細節.
就是耶穌選擇了憐憫.
原因是因為耶穌見到這班群眾.
如同羊沒有牧人一般.
這一句是馬科音裡面獨有的.
事實上如同羊沒有牧人.
不單止是舊約所提到的主題.
其實翻開舊約聖經.
我們經常會見到.
特別在猶太傳統裡.
以色列的領袖.
如摩西.
經常被譽為一位牧羊人.
帶領著以色列人.
來認識和遵守上帝的吩咐.
另一方面在先知書裡.

$^{161}$我們也會見到.
上帝透過先知多次責備.
那些以色列的領袖.
他們失職.
不能像一位好的牧羊人.
做好帶領牧羊百姓的本分.
甚至他們會叛倒百姓.
令他們遠離甚至得罪上帝.
而到了新約.
我們看到情況似乎沒有什麼改善.
就像耶穌現在見到.
這班廣大群眾如同羊沒有牧人一樣.
他看到百姓得不到屬靈的餵養.
於是他動了持心.
他決定放下他原先和門徒的退休計劃.
他優先照顧面前羊群的需要.
彰顯他作為終極大牧人的職份.
至於為什麼耶穌優先不是給他們餅吃.
而是要開口教訓群眾許多道理.
是因為耶穌明白.
這班群眾確實像羊一樣.
他們不能分辨到生命的真正需要是什麼.
我回想起我小時候.
不知道大家有沒有去過荔園.
以前荔園除了玩機動遊戲外.
還可以餵羊.
小時候知道餵羊要買飼料.
買完草和蘿蔔之後餵完.
沒有了怎麼辦呢.
以前很窮很沒錢.
剛好發現之前我去的攤位遊戲.
有些訂街磚的.
訂中了有些贈品.
有些香口膠.
剛好袋子裡有些綠色的物體.
就是綠戰香口膠.
不知道大家有沒有吃過.
現在不知道還有沒有賣.
看到又是綠色的.
不如試試餵.

$^{201}$看他吃不吃.
誰知道羊真的照吃.
你給牠什麼牠不會理會.
到底是什麼顏色.
牠就會照吃.
而我試了一塊.
得的時候就拿第二條.
又試試餵.
身邊其他的羊.
看到有東西吃.
就會一窩蜂湧過來.
牠們不管什麼都好.
總之你閉上牠們的口.
就照吃.
當然小時候頑皮.
後來就知道不應該這樣做.
所以純粹是一個分享.
你會看到其實羊.
牠真的不懂得分什麼是好.
什麼是壞.
什麼是對牠身體有益.
羊也很盲目的跟風.
牠看到那邊的羊好像有好處.
有東西吃.
牠就不管三七二十一.
牠就會湧過去.
同樣群眾當時.
是因為看到或者聽聞到.
耶穌大有能力.
無論是治病,趕鬼.
耶穌都好像很厲害.
於是他們就一窩蜂.
盲目的追捧耶穌.
他們期望可以在耶穌身上.
得到自己想要的好處.
這也是耶穌憐憫他們的.
另一個原因.
因為耶穌很明白他們搞錯了.
生命的優先次序.
到底是什麼.

$^{241}$他們只是關注眼前.
譬如治病,趕鬼.
這些屬於我們說.
次要,短暫的人生需要.
耶穌卻是期望這群白色.
知道其實他們真正需要的是什麼.
就是認識,接受耶穌基督的福音.
得以進入上帝的國度裡.
所以你看到.
馬考特特別強調.
耶穌是率先用上帝的話語.
來慰養這群白色.
來顯示他作為那位好牧人所要做的.
事實上耶穌所做的.
也在呼應舊約.
特別我們剛剛之前說完的.
系列申命記裡面所說.
「人活著不是單靠食物.
乃是靠神口裡所出的一切」.
我們看到.
雖然耶穌和門徒.
他們本身已經很疲倦.
甚至很餓.
但是耶穌認為.
迫切教導這群白色真理的需要.
遠遠凌駕於自己和門徒的肉體需要.
所以即使多疲倦.
耶穌仍然很努力滿足他們.
這個屬靈的需要.
不過可能時間不早了.
經文形容當天晚的時候.
門徒似乎也忍不住.
於是開聲建議耶穌.
夫子時候不早了.
我們也累了.
我們也餓了.
飯也還沒吃.
這裡又是野地.
不如我們趕快叫群眾散開.
讓他們可以去附近的村莊休息.

$^{281}$他們自己買東西吃.
門徒這一問.
也引發了接下來五餅二魚的神蹟的發生.
事實上你看到整件事很奇妙.
原本一個只屬於耶穌加上十二門徒.
十三人的退休會.
突然變成一個超大型的陪靈會.
大家的體力已經透支.
其實很需要休息.
所以門徒這個建議.
我們說是合情合理的.
耶穌也認同.
是時候要放飯了.
不過耶穌很奇怪.
耶穌突然考一考門徒.
要求門徒順便幫群眾安排食物.
門徒可能已經很餓了.
我想都很混亂.
可能你試過肚子餓.
想不到辦法.
有點生氣.
門徒也不例外.
他們馬上提出質疑.
他們反問耶穌.
難道耶穌要我們去買二十兩銀子.
相當於當時一般人二百日薪金.
要我們用這麼多錢去買餅給群眾吃嗎.
你看到他們很快速地.
用他們的理性思維.
把耶穌餵養眾人的吩咐放在一起.
瞬間就看到自己的缺乏.
簡單來說.
他們的反問是.
耶穌我們沒有這麼多錢.
他們的結論是.
耶穌你的要求我們做不到.
根本不合理.
不過其實.
稍稍停下來.
經文給了一個很有趣.

$^{321}$很重要的提醒我們.
原來我們看為合理的建議.
或者分析.
原來不一定是合乎上帝的旨意.
事實上.
信仰的本質.
原來不是純粹遵守一些.
我們認為合理的吩咐.
信仰既然講求我們對上帝的信心.
對祂的完全信靠.
其實仍然會涉及一些.
你覺得不合理的東西.
你有沒有想過.
譬如聖經裡的吩咐.
耶穌說你要愛受敵.
你要愛那些逼迫你的人.
為他們禱告.
其實是不合理的.
一點都不合理.
特別是你放在今天的社會標準裡.
你覺得有沒有搞錯.
想反抬.
不過信仰正正有這種吊詭.
你覺得合理的東西.
耶穌說不一定是這樣.
反過來你覺得不應該這樣做的時候.
耶穌說你試試去做.
正正反映上帝.
其實是一位超越的神.
代表上帝的旨意.
其實遠超於人的盤算和想法.
因為只有上帝才知道.
到底要怎樣做.
才對我們的生命有益處.
所以你會發覺其實不單止是門徒.
上帝有時也會刻意.
在你意想不到的時候.
安排一些很奇妙的吩咐.
要你去遵守.
來考驗一下你.

$^{361}$信不信得過他.
就正如耶穌要這班門徒明白.
確實要你們去餵飽五千人.
並不見得是合情合理的事.
不過耶穌說是一件應該做的事.
耶穌認為憐憫這班群眾.
被他們吃飽.
是當下要優先處理.
並且他要求門徒.
你來配合我.
你來跟我一起做.
所以你看到.
耶穌沒有再糾纏.
他沒有再從門徒爭論.
到底要去哪裡賺這麼多錢買餅.
耶穌直接問門徒.
你先看看你手上有多少食物.
於是門徒很快就回答.
只有五個餅和兩條魚.
就是這樣.
耶穌接著會向他們示範.
怎樣運用手上僅有的資源.
來滿足龐大的需要.
從而改變門徒一種思考模式.
原本他們覺得我們不足夠.
我們只有五餅二魚.
不可能滿足群眾的需要.
或者耶穌的吩咐.
但現在耶穌要透過這個神蹟.
來改變他們的思維.
雖然只有這麼少.
只有五個餅和兩條魚.
但既然是耶穌的吩咐.
就必定足夠有餘.
他要讓門徒明白.
當我們領受上帝給我們的吩咐.
給我們的吩咐任務的時候.
我們不要只看到困難.
不要只看到問題.
而看不到其實那位呼召.

$^{401}$差派我們的主一直和我們同在.
他會不斷的供應我們.
同樣耶穌也會考問我們.
弟兄姊妹.
當我們面對上帝託付給我們的事情的時候.
我們會不會因為覺得眼前的需要太大.
太過乏力.
又或者你很想氣餒.
很擔心的時候.
你會不會想一想.
其實耶穌還在你身邊.
其實耶穌不會要你孤單作戰.
他一定會有夠用的安排和帶領.
來和你同行.
事實上這段經文也提醒我們.
上帝不會要我們從無變有.
因為這是他只能夠做到的工作.
你留意上帝行事的方式.
很多時候是怎樣.
是會先用你手上已有的資源.
他已經安排了給你的資源.
來去成就他的旨意.
正如耶穌沒有要求門徒.
你們去找五千個餅來.
你們變五千個餅.
耶穌沒有這樣做.
耶穌只是很簡單的跟他們說.
你看看你手上還有什麼.
所以弟兄姊妹.
其實這段經文雖然我們已熟能詳.
但其實有很多很重要的提醒.
就是我們不要限制了.
上帝使用我們的方法.
我們不要以為我們這麼微小.
我們怎能夠替上帝來成就大事.
事實上在耶穌的眼中.
只要是我們全心全意的擺上.
哪怕我們看為微不足道也好.
上帝也有辦法使用.
上帝也可以用我們.

$^{441}$來成為別人的祝福.
成為我們成長的養分.
事實上你會看到.
聖經裡面一再提醒一個道理.
在實踐天國使命的過程裡.
需求太大或資源太少.
從來不是耶穌要顧慮的事情.
正如我們也很明白.
他沒有難成的事.
我們繼續看.
看看耶穌怎樣去成就整個的神蹟.
我們說39至43節.
經文裡面記載了三件事.
耶穌吩咐門徒安排群眾坐下.
然後叫他們分派餅和魚.
之後收拾餘下的魚和餅.
你會看到耶穌通過這三件事.
這三個吩咐.
其實是讓門徒上了一課成長課.
好像提供一場門徒訓練.
我們逐個來看.
首先39至40節.
耶穌吩咐門徒安排群眾.
50個也好,100個也好.
分組坐在青草地上.
你留意.
經文其實一直提及.
那地方是曠野.
突然間會有一大片青草.
可以容納到五千人.
門徒當時也沒有航拍機.
可以在高空看下來.
看看有什麼地方適合安排這麼多人.
井井有條地坐下.
他們沒有.
他們很簡單.
耶穌既然這樣說.
我們就憑信心走出那一步.
試試排位.
試試吩咐群眾.

$^{481}$如何井井有條地坐下.
你看到他們按照信心.
信服耶穌的吩咐.
開始安排座位的時候.
結果就成事了.
經文裡記載.
這麼多群眾.
最後就井井有條.
或五十或一百坐在青草地上.
很明顯整件事.
是耶穌早就知道有這個需要.
亦是祂事先為門徒.
預備了這個地方.
事實上.
在許多時候.
當我們順著耶穌.
順著聖經的帶領去做的時候.
你會發覺很多事情.
其實上帝在祂吩咐你之前.
祂預備好.
在你未行動之前.
可能會有很多憂慮.
你會看不清前路.
畢竟我們很有限.
我們沒有一個.
可以預測未來的水晶球.
可以知道往後我們所做的事情.
最終會如何發展.
但正正這段經文提醒我們.
原來在我們行動之前.
上帝已經安排妥當.
重點是我們能否憑信心.
踏出那一步.
放心走每一步.
在行的過程裡.
我們一直敏銳著.
耶穌基督的帶領.
配合祂的作為.
這是經文給我們一個.
很重要的提醒.

$^{521}$事實上每年到了七月八月.
我都會想起我兒子.
以前如何找到一間幼稚園.
到今天我又想起這段經文.
很想和大家再一次分享.
我如何經歷這個信心.
和我太太.
其實我們以前住在荃灣.
有一次在大埔看到一間.
很適合我兒子發展的幼稚園.
但你知道不單止要報名.
還要舉家搬過去.
如果要去報這間幼稚園.
經過禱告.
經過面試後.
覺得都平安.
校長本身是基督徒.
有基督教的背景.
於是去遞表申請.
要一步一步走.
過了幾個星期後.
他又給了一封信.
我們接受你申請.
他說我們沒有全日制課程.
只剩下半日或下半日.
他亦沒有告訴你.
到底是上午班還是下午班.
要何時知道呢?.
當日你來到就會抽獎.
到時告訴你.
你去不去好呢?.
做父母.
大部分特別是雙職父母.
你會期望幫小朋友.
找到一間全日制幼稚園.
方便他學習之餘.
亦可能方便接送.
如果退役求其次.
不行的話,就上午班.
但這間幼稚園沒有告訴你.

$^{561}$你是讀上午還是下午.
我們就很混亂.
萬一抽到下午班.
還去不去呢?.
我和太太都很相信.
既然我們沿途.
認為是上帝帶領.
我們知道是上帝帶領.
你走下去吧.
你唯有走到最後.
有了結果.
你才會知道.
到底上帝最終.
要為你預備什麼.
於是我們就繼續憑信心.
當日就去到抽獎.
當日去到有三十多個家長.
校長就宣佈.
好消息.
現在有四個全日制的位置.
放出來.
因為有四個舊生退學.
很開心.
有四個位置很開心.
校長就問.
有哪個家長想轉讀全日制?.
全部都舉手.
在場有三十多個家長.
校長很公平.
就拿出A4箱.
把洞戒掉.
假設裡面有三十個乒乓球.
編號1至30.
總之抽到1,2,3,4.
就有得讀.
很公平.
我們大概抽了二十多個.
二十四至二十五.
第一個一加抽.
已經抽了1.

$^{601}$慘了.
已經沒有機會了.
只剩下三個機會.
如是者抽抽抽.
第二和第四又出來了.
剩下3.
到我的時候.
我後面還有六七個人.
即是只有六分一機會.
或者七分一機會.
我突然有壓力.
雖然我是基督徒.
我還是傳道人.
我都怕我會否不夠守潔心清.
抽不中呢?.
很沒面子.
很大壓力.
抽中抽不中好像是關我事.
但我那一刻抽了.
都要抽.
放手進去.
很奇妙.
那一刻.
你當我好像有聲音.
我抽旁邊那一隻.
拿出來.
嘩!三個.
很厲害.
第一件事就是對著老婆說.
三個很厲害.
但下一刻就覺得.
上帝你很厲害.
上帝你知道我們的擔憂是什麼.
上帝你的預備也是很奇妙.
第二刻我是有些懼怕.
或者說有些敬畏.
因為我體會到上帝一直帶著.
你不去到最後一步.
原來你是不會知道上帝.
給你的路是怎樣走.

$^{641}$事實上那一次經歷.
為什麼到現在都這麼深刻.
過了十幾年.
不是因為我兒子讀了.
一間心儀的基督教幼稚園.
而是我知道上帝在這裡.
我知道無論將來是怎樣.
上帝就露這一手給我看.
我在這裡看著你.
無論怎樣也好.
無論他將來讀哪間成績是怎樣.
我在這裡看著你們.
這比他到底讀了哪間學校重要.
那是次要.
最重要是我經歷到.
上帝很真實的同在和帶領.
所以弟兄姊妹.
其實這段經文五餅二魚.
也是提醒我們.
很多時候上帝的吩咐.
上帝的帶領.
你真的要憑信心走一步.
一路走下去.
你才會知道.
原來上帝是多麼的厲害.
我們再看下去.
41至42節.
耶穌在這裡說.
拿起餅來和魚.
竹蔗了,抹開第七門徒.
擺在群眾面前.
經文如果你有留意.
當時餵飽的人.
藍丁也有五千人.
假如你是門徒.
當你拿著耶穌所抹開的餅和魚.
分派的時候.
你也會有些擔心.
你可能也會.
假如我們禮堂裡有五百個會眾.

$^{681}$我們也會限住.
每人第一次拿一條魚.
一個餅,不要拿太多.
等派完一輪.
每個人都有的時候.
如果有剩的.
你才會再派.
你也不會任由群眾.
按自己心意隨便拿.
因為你怕最後不夠.
正如大家也有去過.
飲宴的印象.
有一味可能有魚翅湯.
以前很流行.
從舀魚翅湯.
你就會知道侍應是熟手還是生手.
熟手的通常會分得很均勻.
為什麼會分得很均勻.
因為每一碗只能舀七成.
舀七成後.
最後發現餘下的魚翅很大.
他才舀第二碗.
這樣比較好看.
這些就是熟手.
不熟手的會怎樣.
第一,二,三,四,五,六碗.
就會舀滿了.
最後會怎樣呢.
最後不夠,只剩下半碗.
我就見過有個侍應可能生手.
到時怎麼結尾呢.
大碗沒有魚翅了.
他只能在第一,二,三,四碗舀一點.
你就會見到很難看.
事實上我們也會這樣.
當你不肯定夠不夠派的時候.
在當時的場景.
你一定會保守一點.
不會讓他們亂舀.
但在這裡很特別.

$^{721}$這裡的經文說了.
擺在群眾面前.
即是怎樣.
任你們舀.
沒有限制你們.
你們就按自己的需要.
和吃飯的份量.
來領取餅和魚.
結果經文記載.
所有的群眾都拿到餅和魚.
並且吃飽了.
作者要顯示.
耶穌就是那個生命量.
祂本身就是那位的生命量.
祂的供應絕對的充足.
無窮盡.
遠超過門徒.
遠超過群眾所能夠想像.
所以也能給予門徒.
或我們一個很重要的提醒.
我想我們日常工作.
也會有一個習慣.
我們會在事前計劃好.
我們一定會用一個最穩妥的方案.
按自己現有的能力資源.
來匹配我們能夠做多少事情.
盡可能留有一點餘地.
有一點額外來防範.
來避免意外出現.
這也是我們行之有效.
一個好好做事.
甚至是事奉的方法.
但是會有它的盲點.
或有它的不足.
我們有時過於依賴自己的籌算.
我們會貶低了上帝的超越.
限制了上帝的豐盛.
我們不夠膽去做.
這段經文很提醒.
有時我們太精於計算.

$^{761}$太多理性.
太多擔心.
太多相信自己的判斷的時候.
我們不能夠建立對上帝真正的信求.
我們也很難體會到.
原來耶穌是多麼樂意.
將祂的豐盛來跟人分享.
我想就好像這段經文所說.
唯有你真的相信耶穌會包底.
唯有你相信祂是供應的源頭.
你放膽全心全意擺上的時候.
我們才會看到神蹟的發生.
事實上.
真正清楚我們前路.
我們人生的藍圖的是耶穌.
不是我們.
耶穌其實知道.
在我們身上發生什麼事.
我們需要怎樣走.
在什麼階段.
需要什麼改變.
什麼裝備.
耶穌通通都知道一清二楚.
所以你會看到.
耶穌不單單對門徒.
對我們也會定期有些妥善.
不過可能你已想得到的成長課.
安排給我們.
就好像將會放暑假.
家長會為不同的小朋友.
安排不同的成長課程.
因為你清楚知道.
他在什麼階段.
需要有什麼發展.
有什麼改善.
同樣耶穌也是這樣.
耶穌也會定期.
祂作為我們的主.
作為我們的父.
祂也會安排定期的成長課給我們.

$^{801}$所以這段經文很好.
我想特別提醒.
有小朋友的家長.
當你安排了不同的成長課程.
給小朋友進行的時候.
你期望他聽你吩咐.
信服的時候.
亦提醒我們自己本身.
我們是天父的兒女.
我們願不願意信服上帝.
安排給我們的成長課.
我們願不願意配合.
上帝在我們成長.
在我們生命裡的不同階段.
給我們一些歷練和學習.
這是經文要我們思考的.
最後43節.
耶穌吩咐門徒.
把剩下的食物收拾.
你會看到最後.
他們收拾了有12個男子的食物.
為什麼耶穌在這一刻.
要門徒來做一個收拾呢?.
當然重點一定不是純粹.
環保不要浪費食物.
而是門徒一定要親身體驗到.
原來我們剩下的資源.
比我們原先放上的還要多.
而這個體驗.
在他們在派餅.
在分魚的時候.
非常忙的時候.
他們是不會察覺到的.
什麼時候才開始察覺.
就是當他們真的.
完成整件事.
他們開始收拾的時候.
他們就會發覺.
原來我們得到的.
比付出的還要多.

$^{841}$事實上,兄弟姐妹.
門徒這幾個成長的經歷.
我們今天如何可以內化.
成為我們生活的養分呢?.
又或者成為我們.
恆常事奉上帝的一種態度.
其實就好像耶穌這三個吩咐.
安排坐下.
分餅和收拾零碎.
意思是什麼呢?.
其實安排坐下.
好像是叫我們活於當下.
你現在就聽耶穌的吩咐.
接受耶穌的邀請.
你會發覺其實在你行動之前.
原來耶穌已經為你.
預備好你所需要.
甚至超過你想像.
至於分派食物.
就是說你如何.
一路憑信心展望將來.
你要確信其實耶穌.
不單止他差派你.
吩咐你.
他也會在背後.
不斷的供應,保守你.
所以你要對未來有信心.
因為耶穌掌管著.
他會不斷供應你的需要.
引導你成長.
至於收拾零碎.
其實就是我們要定期回顧.
數算上帝的恩典和大靈.
我們要在忙碌的生活裡.
你真的要找機會定時坐下.
數算一下.
到底過往的日子.
上帝如何帶領你走過呢?.
信仰有時有一個盲點.
就是我們一味只顧著做.

$^{881}$忘記了上帝.
在我們一些生命片段裡.
他如何介入.
如何慢慢改變我們.
當我們經常有這樣的回顧時.
慢慢你就會培養一種信仰態度.
就是你知道上帝.
定期要裝備改變你.
因為我們最終.
要成為他合容的器皿.
所以頂子妹你會看到.
伍丙兒這個神蹟.
其實不單止提供一個門訓.
也不單止給我們一個概念.
他很真實地要你下場參與.
就好像耶穌要門徒下場參與一樣.
事實上有人這樣形容教會.
就好像一場足球比賽.
就是在場上參與的人.
往往只有少數.
只有十多個.
但大多數的人.
這不是一個健康的景象.
如果在教會裡.
因為我們都知道.
一個健康的教會.
其實應該是下場參與比賽的人.
多於在台上看東西的人.
所以耶穌很期望.
我們藉著這段經文的提醒.
你回應祂給你的呼召.
你回應祂給你的邀請.
來實踐天國的使命.
最後說一個例子.
話說有一位.
他很擅長走鋼線的高手.
他有一次安排了一場表演.
有在場的觀眾.
他在兩座山峰裡.
拉了一條鋼索.

$^{921}$他要表演.
由山的一端走到另一端.
於是他開始.
他沒有平衡桿.
單靠雙手慢慢移步.
由一個山頭走到另一個山頭.
他成功做了創舉.
於是在場的人.
這位高手覺得意猶未盡.
他決定加碼.
這次我要綁著雙手在後面.
減低我的平衡力.
我再做一次.
大家信不信我可以.
成功從這座山走到另一座山.
大家因為很想知道結果.
所以就說我們相信你.
你一定做得到.
於是這位走鋼線高手.
果然又真的成功.
又再做出一次創舉.
綁著雙手都可以跨過那條鋼線.
群眾很興奮.
給了很多讚好.
表演還未完.
這位高手決定來一個終極示範.
他在群眾中找了一個小孩子.
他對大家說.
其實這是我的兒子.
我現在要將他放在我的肩膀上.
像騎駁馬一樣.
然後我仍然將雙手綁在後面.
我這樣要嘗試走過去這條鋼線.
你們信不信我可以做得到呢?.
群眾說你這麼厲害.
我們一定相信你.
於是這位鋼線高手再問.
你們真的相信我?.
這樣都可以?.
很難的.

$^{961}$群眾說你一定行的.
沒問題的.
於是這位走鋼線高手最後問.
既然大家這麼相信我.
我將我的兒子放下.
換成你的兒子.
有哪個兒子呢?.
大家都知道.
即時在場所有人鴉雀無聲.
沒有人再敢說我相信你.
弟兄姊妹.
其實有時我們面對信仰.
面對主耶穌.
可能我們都是這樣.
我們所謂相信他.
會不會只是掛在嘴邊.
一旦真的要下場參與.
真的要埋身寫記的時候.
我們又會不會退縮?.
所以弟兄姊妹.
其實今天這段經文提醒我們.
唯有我們相信耶穌.
奇蹟才會發生.
唯有我們真的願意和耶穌同工.
我們的生命才會成長.
信仰永遠不能只是掛在嘴邊.
一定要下場參與.
所以盼望今天的故事.
五餅二魚的神蹟提醒我們.
當耶穌要求我們去服侍.
餵養我們身邊的人的時候.
我們不要想都不想就說不行.
我們不要想都不想就拒絕.
你嘗試一下.
好像今天這段經文一樣.
看看你現有的資源.
耶穌不會無緣無故來到呼召你.
祂一定會為你預備好你所需要.
你嘗試將你的目光.
移回到耶穌.

$^{1001}$那位能夠成就萬事的主身上.
並且你確信祂不會撇下你.
祂一定有妥善的安排.
你可以將你的擔心掛慮交託給祂.
從而配合祂的帶領和安排.
就好像今天這段經文所分享.
盼望耶穌基督的神蹟.
不單單幫助昔日的門徒.
也能夠成為我們成長的一課.
請阿細倫帶領我們唱會認詩.
\newpage



\section{約翰福音 8:31-59}
\label{sec:gyB32xPImr0}
\textbf{2025年7月27日崇拜直播｜梁家麟牧師｜放棄耶穌的兩大理由｜約翰福音8\_31-59}
\newline
\newline
連結: \href{https://youtube.com/watch?v=gyB32xPImr0}{\texttt{https://youtube.com/watch?v=gyB32xPImr0}} ~~~~ 語音日期: 2025-07-27
\newline
\newline
\hyperref[sec:SbKqVXHzrXE]{\small{< < < PREV SERMON < < <}}
~
\hyperref[sec:index_chronic]{\small{[返順時目]}}
~
\hyperref[sec:woNoWOUM52k]{\small{> > > NEXT SERMON > > >}}
\newline
\newline
約翰福音 8:31-59
\newline
\begin{longtable}{cl}
\hline
\hline
章節 & 經文 (和合本修訂版)\\
\hline
8:31 & \begin{tabularx}{0.7\textwidth}{X} 耶穌對信他的猶太人說：「你們若繼續遵守我的道，就真是我的門徒了。 \end{tabularx} \\ \\ \relax
8:32 & \begin{tabularx}{0.7\textwidth}{X} 你們將認識真理，真理會使你們自由。」 \end{tabularx} \\ \\ \relax
8:33 & \begin{tabularx}{0.7\textwidth}{X} 他們回答他：「我們是亞伯拉罕的後裔，從來沒有作過誰的奴隸，你怎麼說『會使你們自由』呢？」 \end{tabularx} \\ \\ \relax
8:34 & \begin{tabularx}{0.7\textwidth}{X} 耶穌回答他們：「我實實在在地告訴你們，所有犯罪的人就是罪的奴隸。 \end{tabularx} \\ \\ \relax
8:35 & \begin{tabularx}{0.7\textwidth}{X} 奴隸不能永遠住在家裡；兒子才永遠住在家裡。 \end{tabularx} \\ \\ \relax
8:36 & \begin{tabularx}{0.7\textwidth}{X} 所以，神的兒子若使你們自由，你們就真正自由了。」 \end{tabularx} \\ \\ \relax
8:37 & \begin{tabularx}{0.7\textwidth}{X} 「我知道，你們是亞伯拉罕的後裔，你們卻想要殺我，因為你們心裡容不下我的道。 \end{tabularx} \\ \\ \relax
8:38 & \begin{tabularx}{0.7\textwidth}{X} 我所說的是在我父那裡看見的；你們所做的是在你們的父那裡聽到的。」 \end{tabularx} \\ \\ \relax
8:39 & \begin{tabularx}{0.7\textwidth}{X} 他們回答耶穌：「我們的父是亞伯拉罕。」耶穌對他們說：「你們若是亞伯拉罕的兒女，就會做亞伯拉罕所做的事。 \end{tabularx} \\ \\ \relax
8:40 & \begin{tabularx}{0.7\textwidth}{X} 我把在神那裡所聽見的真理告訴了你們，現在你們卻想要殺我；亞伯拉罕沒有做過這樣的事。 \end{tabularx} \\ \\ \relax
8:41 & \begin{tabularx}{0.7\textwidth}{X} 你們是做你們父的工作。」他們就對他說：「我們不是從淫亂生的，我們只有一位父，就是神。」 \end{tabularx} \\ \\ \relax
8:42 & \begin{tabularx}{0.7\textwidth}{X} 耶穌對他們說：「假如神是你們的父，你們會愛我，因為我本是出於神，也是從神而來，我不是憑著自己來，而是他差我來的。 \end{tabularx} \\ \\ \relax
8:43 & \begin{tabularx}{0.7\textwidth}{X} 你們為甚麼不明白我的話呢？無非是你們聽不進我的道。 \end{tabularx} \\ \\ \relax
8:44 & \begin{tabularx}{0.7\textwidth}{X} 你們是出於你們的父魔鬼，你們寧願隨著你們父的慾念而行。他從起初就是殺人的，不守真理，因他心裡沒有真理。他說謊是出於自己的本性，因他本來是說謊的，也是說謊者之父。 \end{tabularx} \\ \\ \relax
8:45 & \begin{tabularx}{0.7\textwidth}{X} 但是，因為我講真理，你們就不信我。 \end{tabularx} \\ \\ \relax
8:46 & \begin{tabularx}{0.7\textwidth}{X} 你們中間誰能指證我有罪呢？既然我講真理，你們為甚麼不信我呢？ \end{tabularx} \\ \\ \relax
8:47 & \begin{tabularx}{0.7\textwidth}{X} 出於神的，必聽神的話；你們不聽，因為你們不是出於神。」 \end{tabularx} \\ \\ \relax
8:48 & \begin{tabularx}{0.7\textwidth}{X} 猶太人回答他：「我們說你是撒瑪利亞人，並且是被鬼附的，這話不是很對嗎？」 \end{tabularx} \\ \\ \relax
8:49 & \begin{tabularx}{0.7\textwidth}{X} 耶穌回答：「我沒有被鬼附；我尊敬我的父，你們卻不尊敬我。 \end{tabularx} \\ \\ \relax
8:50 & \begin{tabularx}{0.7\textwidth}{X} 我不尋求自己的榮耀，但有一位為我尋求榮耀，判斷是非。 \end{tabularx} \\ \\ \relax
8:51 & \begin{tabularx}{0.7\textwidth}{X} 我實實在在地告訴你們，人若遵守我的道，就永遠不經歷死亡。」 \end{tabularx} \\ \\ \relax
8:52 & \begin{tabularx}{0.7\textwidth}{X} 於是猶太人對他說：「現在我們知道你是被鬼附了。亞伯拉罕死了，眾先知也死了，你還說：『人若遵守我的道，就永遠不經歷死亡。』 \end{tabularx} \\ \\ \relax
8:53 & \begin{tabularx}{0.7\textwidth}{X} 難道你比我們的祖宗亞伯拉罕還大嗎？他死了，眾先知也死了，你把自己當作甚麼人呢？」 \end{tabularx} \\ \\ \relax
8:54 & \begin{tabularx}{0.7\textwidth}{X} 耶穌回答：「我若榮耀自己，我的榮耀就算不了甚麼；榮耀我的是我的父，就是你們所說的你們的神。 \end{tabularx} \\ \\ \relax
8:55 & \begin{tabularx}{0.7\textwidth}{X} 你們不認識他，我卻認識他。我若說不認識他，我就是說謊的，像你們一樣；但我認識他，也遵守他的道。 \end{tabularx} \\ \\ \relax
8:56 & \begin{tabularx}{0.7\textwidth}{X} 你們的祖宗亞伯拉罕歡歡喜喜地仰望我的日子，他看見了，就快樂。」 \end{tabularx} \\ \\ \relax
8:57 & \begin{tabularx}{0.7\textwidth}{X} 猶太人就對他說：「你還沒有五十歲，難道見過亞伯拉罕嗎？」 \end{tabularx} \\ \\ \relax
8:58 & \begin{tabularx}{0.7\textwidth}{X} 耶穌對他們說：「我實實在在地告訴你們，還沒有亞伯拉罕我就存在了。」 \end{tabularx} \\ \\ \relax
8:59 & \begin{tabularx}{0.7\textwidth}{X} 於是他們拿石頭要打他，耶穌卻躲開，走出了聖殿。 \end{tabularx} \\ \\
[1ex]
\hline
\hline
\end{longtable}
$^{1}$今日要讀的聖經是.
約翰福音第八章第三十一至五十九節.
聖經新約第一百三十八至一百四十頁.
約翰福音第八章第一節.
「耶穌對信他的猶太人說:.
「你們若常常遵守我的道,就真是我的門徒..
你們必曉得真理,真理必叫你們得以自由.」.
他們回答說:.
「我們是亞伯拉罕的後裔,從來沒有作過誰的奴僕..
你怎麼說你們必得以自由呢?」.
耶穌回答說:.
「我實實在在的告訴你們,.
所有犯罪的,就是罪的奴僕..
奴僕不能永遠住在家裡,.
兒子是永遠住在家裡..
所以,天父的兒子若叫你們自由,.
你們就真自由了..
我知道你們是亞伯拉罕的子孫,.
你們卻想要殺我,.
因為你們心裡容不下我的道..
我所說的是在我腹那裡看見的,.
你們所行的是在你們的腹那裡聽見的.」.
他們說:.
「我們的腹就是亞伯拉罕.」.
耶穌說:.
「你們若是亞伯拉罕的兒子,.
就必行亞伯拉罕所行的事..
我將在腹那裡所聽見的真理告訴了你們,.
現在你們卻想要殺我,.
這不是亞伯拉罕所行的事..
你們是行你們腹所行的事.」.
他們說:.
「我們不是從淫亂生的,.
我們只有一位腹,就是神.」.
耶穌說:.
「倘若神是你們的腹,.
你們就必愛我,.
因為我本是出於神,.
也是從神而來,.
並不是由著自己來,.

$^{41}$乃是祂差我來..
你們為何不明白我的話呢?.
無非是因為你們不能聽我的道,.
你們是出於你們的腹魔鬼,.
你們腹的私慾,.
你們偏要行..
祂從起初是殺人的,.
不守真理,.
因祂心裡沒有真理,.
祂說謊是出於自己,.
因祂本來是說謊的,.
也是說謊之人的腹..
我將真理告訴你們,.
你們就因此不信我,.
你們中間誰能指證我有罪呢?.
我既然將真理告訴你們,.
為何不信我呢?.
我將出於神的,.
別聽神的話,.
你們不聽,.
因為你們不是出於神..
猶太人回答說,.
我們說你是撒瑪利亞人,.
並且是鬼附著的,.
這說豈不正對嗎?.
耶穌說,.
我不是鬼附著的,.
我尊敬我的父,.
你們倒輕說我,.
我不求自己的榮耀,.
有一位為我求榮耀,.
定是非的..
我實實在在的告訴你們,.
人若遵守我的道,.
就永遠不見死..
猶太人對他說,.
現在我們知道你是鬼附著的,.
亞伯拉罕死了,.
眾先知也死了,.
你還說人若遵守我的道,.

$^{81}$就永遠不常死嗎?.
難道你比我們的祖宗亞伯拉罕還大嗎?.
他死了,.
眾先知也死了,.
你將自己當作什麼人呢?.
耶穌回答說,.
我若榮耀自己,.
我的榮耀就算不得什麼,.
榮耀我的乃是我的父,.
就是你們所說是你們的神..
你們未曾認識他,.
我卻認識他..
我若說不認識他,.
我就是說謊的,.
像你們一樣,.
但我認識他,.
也遵守他的道..
你們的祖宗亞伯拉罕,.
歡歡喜喜地仰望我的日子,.
饑餓了就快樂..
猶太人說,.
你還沒有五十歲,.
若見過亞伯拉罕呢?.
豈見過亞伯拉罕呢?.
耶穌說,.
我實實在在地告訴你們,.
還沒有亞伯拉罕,.
就有了我..
於是他們拿石頭要打他,.
耶穌卻躲藏從殿裡出去了..
寶玲,姐妹讀得很好,.
讀聖經讀得很動聽..
我們仍然看耶穌在聖殿裡跟群眾對話,.
有趣的是,.
這段聖經,.
作者約翰特別在31節提到,.
耶穌對信他的猶太人說,.
耶穌說話的對象,.
不是跟他徹底對立,.
完全否定他的身份和言論的那些法利賽人和法師,.

$^{121}$而是對他有好感,.
願意坐下來聽他說話的猶太人..
這些人在某個程度上是相信他的,.
不過信什麼呢?.
他們大概相信耶穌是先知,.
這是最低消費的..
耶穌所傳的道,.
有屬天的智慧,.
值得聆聽學習,.
所以他們才站在他身邊,.
假裝聽眾,.
不過他們對耶穌的信,.
未曾達到耶穌的期望,.
耶穌是一個人信他,.
但不是將他當作律法師或先知,.
而是確認他是上帝的兒子,.
從上帝而來,.
所宣講的完全是父上帝的心意,.
即是那一位,.
但耶穌沒有明講,.
耶穌是期望人知道他是上帝,.
人對耶穌的信,.
和耶穌期望他們的信,.
總是有一個懸殊的距離,.
那麼耶穌怎麼宣?.
先收課,.
諒解他們的理解程度有限,.
容讓他們先慢慢跟著,.
再調教對耶穌的認識,.
這是我們多數人比較溫和的做法,.
不過耶穌在人間的時間不多,.
所以他不是這樣做,.
所以耶穌總是用激烈的言辭,.
要求要強迫群眾正視他的真實身份,.
要求他們在這個身份面前表態,.
你信還是不信,.
如果是信他,.
就要信他是上帝的獨生子,.
否則就是完全不信,.
耶穌寧願將這些因誤會而結合的跟隨者趕走,.

$^{161}$要嗎?徹底信,.
要嗎?就拉倒,.
割斷關係,.
或者會破壞群眾對他的好感,.
甚至由愛變恨,.
但是他在所不惜,.
歷代歷世都有不同人對耶穌有不同的理解,.
印度的甘地就稱耶穌為人類的導師,.
不只是印度人,是整個人類的導師,.
他特別喜歡耶穌登山寶訓的教導,.
但是甘地或者這班人,.
不因此就成為真實的基督徒,.
也不會因著他們信耶穌是人類導師,.
就可以和那個永活上帝建立和好的關係,.
你想想,.
耶穌基督渡成肉身降世為人,.
經歷33年的人生禍難,.
最終為我們釘在十字架上,.
難道就僅僅是想我們終結他的教訓,.
覺得他很有智慧,.
真的可以做人類的導師嗎?.
不是的,耶穌的教訓不僅僅有智慧,.
啟發你的思考,.
很啟發,.
更加是拯救性的真理,.
是一個救贖的真理,.
要幫助人改變,.
耶穌渴望人知道他的身份,.
唯有你知道他的真身份,.
你才能準確領悟他的道理,.
接受他的救恩,.
不過在現階段,.
他又不方便直接說他是上帝,.
甚至亮出他的神聖能力,.
例如顯大能,.
秒殺羅馬帝國的軍隊,.
這樣行不行呢?.
或者吩咐亞伯拉罕,.
以撒,亞國,摩西,大衛,.
所有偉人,.

$^{201}$通通四而復生,.
站在他後面顯現,.
即刻所有人都趴下,.
不過如果耶穌這樣做,.
他還需要上十字架嗎?.
還需要為人類代贖犧牲嗎?.
他打死了羅馬軍隊,.
誰拉他上十字架?.
所以耶穌只能吞吞吐吐,.
半明半忍,.
不斷作出各種側面提示,.
但又不能直接說出標準答案,.
你可以想像耶穌的鬱結非常大,.
真的很鬱悶,.
在今天我們所讀的一大段對話中,.
耶穌牽扯到兩個話題,.
第一是自由,.
第二是亞伯拉罕,.
先說自由,.
耶穌對那些對他有正面好感的猶太人說,.
31到32節,.
你們若常常遵守我的道,.
真是我的門徒,.
你們必曉得真理,.
真理必叫你們得以自由..
耶穌原本是祝福他們要聽耶穌的話,.
然後他就引入了自由這個觀念,.
指出唯有他所說的真理,.
才讓他們得到真正的自由,.
耶穌這番說話有沒有產生讓人聽從他的說服力?.
沒有,.
剛剛好弄巧反拙,.
還造成新的難阻和困惑,.
讓人得著自由,.
這句說話太過offensive,.
太過冒犯性,.
誰才要自由?.
就是那些原本不自由的人,.
唯有失去自由的人才需要自由,.
耶穌這樣說,.

$^{241}$就是暗示在場的聽眾都是不自由的人,.
他們馬上就回應,.
33節,.
我們是亞伯拉罕的後裔,.
從來沒有做過誰的老僕,.
你怎麼說你們必得自由?.
很明顯,.
耶穌口中的自由和群眾所理解的完全是兩樣的,.
群眾所理解的就是政治上或者社會上的自由,.
就是沒有成為奴隸,.
沒有失去行動上的自主性,.
沒有失去自由就是自由人,.
現在耶穌說的是屬令上,.
人因犯罪而成為罪的老僕,.
活在罪中無法自拔,.
沒有不犯罪的自由,.
沒有和上帝建立正面聯繫的自由,.
耶穌指出一個很重要的事實,.
罪的一個本質是合制性,.
不管犯什麼罪,.
都會使我們被拘束,.
被迫去做一些事情,不能不做,.
我們以為我們得著吸煙的權利,.
得著吸大麻的自由,.
但是想一想,.
吸大麻的人自由,.
還是不吸大麻的人自由?.
將人從罪中釋放出來,.
是耶穌在人間最主要的使命,.
耶穌在出道的時候,.
他在納薩勒的會堂拿起以賽亞書就宣讀,.
《羅格福音》四章十八到十九節,.
主的靈在我身上,.
因為他用羔告我,.
叫我傳福音給貧窮的人,.
妻我報告被擄得釋放,.
乞訝得看見,.
叫那受壓制的得自由,.
報告上帝閱立人的欺憐,.
然後再宣告,.

$^{281}$以賽亞的話,.
今天兌現在他們眼前,.
被擄得釋放,.
乞訝得看見,.
叫那受壓制的得自由,.
耶穌使我們由罪的盧伯變成自由的人,.
這也是耶穌在這裡回答群眾的話的意思,.
三十四到三十六節,.
我實實在在的告訴你們,.
所有犯罪的就是罪的盧伯,.
盧伯不能永遠住在家裡,.
兒子是永遠住在家裡,.
所以天父的兒子若叫你們自由,.
你們就真自由了,.
耶穌在這裡關心的主要是人和天父的關係,.
自從犯罪之後,.
人和上帝的關係就嚴重破損了,.
人失去了親近上帝的資格,.
喪失了永生的資格,.
耶穌來人間主要就是修補這個破損了的關係,.
讓人得以重回天父的懷抱,.
住在天父的家裡,.
猶太人聽到耶穌這番話就很不高興,.
他們向來自認是上帝的選民,.
和上帝有特殊又密切的關係,.
他們不覺得他們和上帝的關係是斷裂了,.
他們反駁說,.
耶穌說的父上帝也是他們向來的上帝,.
他們一直和上帝有聯繫,.
一直在事奉他,.
這是一節,.
耶穌他是斷然指出他們和上帝沒有關係,.
他列出一個證明,.
這班人反對他,.
拒絕他的教訓,.
甚至想殺害他,.
這樣的人怎麼可能是依從上帝的心意呢?.
42節,.
耶穌說,.
倘若上帝是你們的父,.

$^{321}$你們就必愛我,.
因為我本來是出於上帝,.
有時從上帝而來,.
並不是由著自己來,.
那時他差我來..
47節,.
也說,.
出於上帝的必聽上帝的話,.
你們不聽,.
因為你們不是出於上帝..
很明顯,.
耶穌仍然是順著他一貫的堅持,.
我們上一次也說過,.
耶穌堅持他的教導,.
完全是出於上帝,.
耶穌的心意就是父上帝的心意,.
耶穌和父上帝毫無差別,.
所以愛上帝的,.
就必愛耶穌,.
並且聽他的話,.
拒絕耶穌的,.
就是拒絕父上帝,.
和耶穌敵對的,.
就是和父上帝敵對..
這也呼應了,.
約翰在一章十八節所說,.
從來沒有人見過上帝,.
只有父懷裡的獨生子,.
將祂表明出來,.
耶穌是上帝在人間的唯一代理,.
壟斷了我們對上帝的接觸和認知,.
祂是我們唯一能夠認識的上帝,.
和耶穌敵對,.
甚至想殺害耶穌的人,.
一定不屬於上帝,.
只能夠屬於撒旦,.
所以這些猶太人,.
耶穌說,他們其實是撒旦的兒女,.
44節,你們是出於你們的父魔鬼,.
你們父的使辱,.

$^{361}$你們偏要行,.
他從起初是殺人的不守真理,.
因為他心裡沒有真理,.
他說謊是出於自己,.
因為他本來是要說謊的,.
也是說謊之人的父..
我們說了第一部分,.
有關自由的課題,.
然後我們來到說第二個的題目,.
是亞伯拉罕,.
第二個的theme,.
在耶穌和猶太人的對話裡面,.
除了牽扯到自由這個觀念之外,.
也將亞伯拉罕這個以色列人的祖宗,.
也拉了下來,.
事緣猶太人為了證明自己是沒有受過奴役的自由人,.
就宣稱他們是亞伯拉罕的子孫,.
33節,.
耶穌就這樣回答,.
39到40節,.
你們若是亞伯拉罕的兒子,.
就必行亞伯拉罕所行的事,.
我將在上帝那裡所聽見的真理告訴你們,.
現在你們卻想要殺我,.
這不是亞伯拉罕所行的事..
如果他們真是亞伯拉罕的子孫,.
又怎會和耶穌對著幹,.
甚至想加害耶穌呢?.
在51節,.
耶穌又提到,.
我實實在在的告訴你們,.
人若遵守我的道,.
就永遠不死..
對此猶太人的反應是彈起,.
你說謊,.
不要說耶穌沒辦法讓人長生不死,.
耶穌自己也不可能長生不死,.
他們的常識推理是,.
連他們最偉大的祖宗亞伯拉罕都死了,.
所有過去曾經存在的先知都死了,.

$^{401}$耶穌說讓人不死,.
豈不是他比亞伯拉罕,.
比眾先知還要偉大嗎?.
如果有人跟你兜售長生不老藥,.
你會不會相信呢?.
耶穌這樣回答,.
56節,.
你們的祖宗亞伯拉罕,.
歡歡喜喜仰望我的日子,.
既看見了,就快樂..
耶穌沒有明講,.
但耶穌的意思其實很清楚,.
喂,亞伯拉罕不過是我的馬仔,.
他見到我的時候,.
大氣都不敢喘,.
你們如果是他的後裔,.
又怎可以這樣對我不敬呢?.
亞伯拉罕仰望耶穌的日子,.
什麼叫耶穌的日子?.
當然是耶穌第二次從天降臨,.
結束地上一切,審判萬民的時候,.
亞伯拉罕雖然是舊約的人物,.
卻仍然一直盼望主在內的日子,.
因為舊約所有應許的落實應驗,.
以色列民的最終得贖,.
都要在耶穌的日子才完全兌現,.
亞伯拉罕在猶太人心中,.
是極其崇高的人物,.
但對耶穌來說,.
不過是一個凡人,.
受造物和創造主無可比性,.
亞伯拉罕見到耶穌只能夠.
「服服在地叩拜」,.
耶穌不會特別尊敬他,.
猶太人越發不能接受耶穌的話,.
他們質疑說,.
在57節,你都未到50歲,.
你怎會見過亞伯拉罕呢?.
耶穌就這樣回答,.
我實實在在的告訴你們,.

$^{441}$還沒有亞伯拉罕就有了我,.
他比亞伯拉罕更早存在,.
他是自由永有的,.
猶太人認定耶穌說的是潛望的話,.
就拿石頭打他,.
不過好彩耶穌就躲避了,.
離開聖殿,.
這段聖經我們都看完了,.
在福音書裡面,.
我們常常看到耶穌和人,.
牛頭不對馬嘴的對話,.
耶穌說的話,.
他們完全無法理解,.
甚至他們只認定耶穌說的是潛望的話,.
聘請議論,問題不出在群眾那裡,.
而是出在耶穌身上,.
因為耶穌根本不是在說人話,.
而是在說神話,.
完全超出猶太人的理解,.
甚至想像的範圍,.
假如我和你說,我見過秦始皇,.
我和孔子吃過一餐飯,.
你們有什麼反應呢?.
大概你當我語無倫次,.
胡說八道,.
所以我們看到耶穌在人間的教導,.
是多次失效,.
被人拒絕,.
被人當發神經說謊,.
甚至說他褻瀆神明,.
口出狂言,.
他們不會覺得奇怪,.
那些認定耶穌是凡人的,.
不可能聽得懂耶穌說的話,.
民事和法理錯人拒絕他,.
群眾誤解他,.
連跟隨他的門徒,.
基本上到他死之前,.
都還沒有真正了解他,.
其實耶穌自己也知道,.

$^{481}$他曾經有一次慨嘆,.
在《馬特福音》11章6節,.
當他知道連屍體約翰都懷疑他的時候,.
他就這樣慨嘆過,.
「凡不因我跌倒的,就有伏料.」.
他知道很多人在他面前會被絆倒,.
我們看到一位非常寂寞的主,.
不被理解,.
不受尊重,.
沒辦法跟人有效的溝通,.
看完這一段經文,.
我們做一些信仰思考,.
這一段經文從一開始說,.
耶穌是對信他的猶太人說的,.
31節,.
聽道的猶太人,.
本來是對耶穌有一定的信心,.
不過在耶穌做完這一次宣講之後,.
人不是由不信變成信,.
由小信變成大信,.
而是剛剛相反,.
最初有一點點信心,.
變成徹底敵對,.
最後要用石頭扔死耶穌,.
不止拒絕耶穌的信息,.
連講完也要殺,.
講到竟然帶來這樣的反效果,.
是不是很奇怪?.
在耶穌三年半的公開事奉生涯裡面,.
我們確實看到有很多本來有一點信心,.
半信半疑,.
又信又不信的人,.
跟著在耶穌身邊,.
在耶穌傳道的高峰期,.
圍在他身邊的有幾千人以上,.
特別是他行了五餅二魚這麼偉大的神蹟之後,.
蜂擁而來跟隨的人很多很多,.
講這些人對耶穌沒有好感嗎?.
不可能,.
講他們對耶穌完全沒有信心嗎?.

$^{521}$不可能,.
耶穌是先知,.
耶穌能夠行神蹟而能,.
是他們一致公認的判斷,.
在後期更加有不少人,.
是相信耶穌是舊約所允許的彌賽亞,.
用來復興以色列的國家,.
他們從來都是用政治的角度去理解,.
彌賽亞的使命,.
所以他們要擁納耶穌做王,.
群眾不是對耶穌徹底沒有信心,.
但他們對耶穌的想法,.
跟耶穌的真相,.
跟耶穌期待他的認識,.
是有很懸殊的距離,.
正因為他們對耶穌的身份和使命是一知半解,.
不是完全錯,.
他是片面的,.
對和錯是混在一起,.
未曾窺見耶穌的真貌,.
所以他們某個階段,.
對耶穌熱情的敬拜跟從,.
在另一個時段又會棄他而去,.
甚至想殺死他,.
他們一時對耶穌滿懷希望,.
一時又徹底失望,.
他們的情緒高低起伏,.
非常波動,.
他們對耶穌失望,.
當然耶穌也對他們失望,.
不要說外圍群眾,.
連耶穌手挽的,.
親自揀選的門徒團隊,.
只有十二人這麼多,.
也有一個人對耶穌徹底失望,.
最後選擇出賣背叛耶穌,.
我們不會假設,.
交流人從一開始就不信耶穌,.
不過他後來真的失望,.
耶穌不是他所期望的宗教領袖,.

$^{561}$朝夕跟耶穌共對的門徒,.
還有放棄耶穌的人,.
一般群眾信心不夠牢固,.
搖擺不定,.
是可以想像的,.
人為何會改變對耶穌的想法呢?.
他們見異思遷的原因何在呢?.
當然,.
總理由就是他們個人的期望,.
跟現實的耶穌有落差,.
他們對耶穌失望,.
他們不喜歡耶穌的教導,.
不過仔細看,.
從這段聖經,.
我們可以找到兩個比較具體的理由,.
就是剛才我們說的兩個主題,.
第一,.
第一個理由,.
他們不肯承認自己犯罪和徹底無能的本相,.
這也牽涉到他們的驕傲,.
他們的自豪,.
猶太人原來是喜歡聽耶穌的講道,.
耶穌幫他們長知識,.
但是耶穌很多說話又為他們帶來威脅,.
耶穌的說話是一個flag,.
因為要強迫他們接受一些他們不想承認的真相,.
例如,.
耶穌說他們是罪的奴僕,.
不是自由的人,.
這就傷害了他們的自我形象,.
侵犯了他們的尊嚴,.
赦罪的前提是要認罪,.
治病的前提是你自覺需要醫生,.
他們不想承認自己有罪,.
就沒辦法接受耶穌的赦罪,.
他們拒絕承認自己是罪的奴僕的事實,.
就沒辦法得到耶穌的釋放,.
真理真的會讓人得自由,.
但真理首先會照出人的本相,.
拆穿人艱苦經營的自我保護機制,.

$^{601}$我不知道你有沒有這樣的情況,.
你年輕就不會的,.
年紀大一點點就像我這樣就會了,.
每次看醫生,每次做身體檢查,.
我都覺得有壓力,.
因為要被迫正視一些我們不喜歡聽見的數據,.
我的血糖又高了,.
我的膽固醇又超標,.
我的BMI指數不合格,.
你喜歡這樣的調查結果嗎?.
我當然不喜歡,.
所以我每次驗血之前都會戒甜品一個星期,.
正如我以前說的那個笑話,.
先找一個不熟悉的牙醫,.
洗牙後才見自己的牙醫,.
你說這個叫慰疾忌醫嗎?.
這是人之常情,.
有個笑話我最近聽的,.
有個姐妹走來跟牧師說,.
牧師,我要認罪,.
因為我覺得自己很漂亮,.
經常顧著照鏡,.
就沒心機讀聖經,.
你說這是不是犯了驕傲罪呢?.
牧師就很嚴肅回答,.
這肯定是罪,大罪,.
不過這不是驕傲的罪,.
是說謊和自我欺騙的罪,.
港台的信息冒犯或傷害了人,.
這是常常發生的事,.
聖經對人的罪性真相的診斷,.
不是人喜歡聽的,.
我要說,就算不說負面的評斷,.
不指責人犯罪,.
只說正面的信息,.
高舉聖經的教導,.
都同樣會帶來威脅,.
我上星期在荷蘭主講一個夏令營,.
星期日在聯合崇拜,.
中文和荷語全部混在一起,.

$^{641}$幾百人,.
我說一個很簡單的題目,.
聖經對婚姻和家庭的教導,.
我要告訴你,.
我講完之後,引起了很大的騷動,.
我講正面,全部講正面,.
引起了很大的騷動,.
竟然有幾個分別來找我,.
是年輕的荷語人,.
他們只是聽翻譯,.
跟我說很震動,很好聽,.
但我又看到,.
我一邊講的時候,.
太太已經在她旁邊的老公耳邊不斷說,.
你都沒有這樣做,.
你明不明白,.
為什麼你不是這樣,.
為什麼你這麼壞,.
然後再有幾個人來跟我說,.
他以前怎樣受過他父母的傷害,.
現實的他和聖經的理想有嚴重的差距,.
有個年輕信徒,.
跟他母親說,.
他母親也是姊妹,.
都在那個營裡,.
就這樣跟他母親說,.
你知不知道你過去經常傷害了我,.
聖經的理想就對照了現實的不理想,.
耶穌問在場的猶太人,.
我將真理告訴你們,.
你們就因此不信我,.
你們中間誰能指證我有罪呢,.
我既然將真理告訴你們,.
為什麼不信我呢,.
45到46節,.
耶穌只是講真理,.
沒有責備他們,.
只是宣告他們可以得到自由,.
不先說他們是罪的奴僕,.
不過也已經構成對他們的冒犯,.

$^{681}$人必須在上帝面前承認自己有罪,.
真誠面對自己的罪過和無能,.
我們是罪人,.
我們無力自救,.
我們才需要真理的釋放,.
我們需要耶穌,.
什麼人最終會放棄信仰耶穌,.
拒絕接受自己徹底無能無有的人,.
拒絕承認自己最大惡極,.
最無可恕的人,.
我要說,.
信主這麼多年,.
其實不停被迫問真相,.
有時候我也想輕輕地跟耶穌說,.
我都讓得夠多,.
還迫你?.
還要這麼盡力嗎?.
是不是?.
給我一點面子,.
你不用拆毀我,否定我到這個地步,.
我不知道你會不會有這樣的反應,.
什麼人想守住自己尊嚴的底線,.
想守住自己認為還可以拿得住的資本,.
你最終會被耶穌冒犯,.
這是第一個原因,.
那些信耶穌一點點的人,.
最終拒絕耶穌,.
因為耶穌迫他們,.
接受他們不想再知道的真相,.
第二,.
第二個原因,.
他們認為自己在耶穌之外,.
尚有可以依次的地方,.
some other thing, reliable,.
不是完全helpless,.
完全無助的,.
留有一手,留有退路,.
不喜歡corner自己,.
不喜歡陷入一個dead end的狀態,.
猶太人是很驕傲的民族,.

$^{721}$他們相信自己是亞伯拉罕的子孫,.
上帝的選民,.
比世界上任何民族都要優秀,.
上帝對他們零眼相看,.
儘管上帝有時候會責備他們,.
不過這是愛之深責之切,.
最終上帝會拯救他們,.
真正有詞說,.
在41節,.
我們不是從淫亂生的,.
我們血統純正,.
我們只有一個父,.
就是上帝,.
他們是天父的兒女,.
保羅說,.
以色列民族,.
即使如今拒絕耶穌,.
最終都會全家得救,.
都是這樣的信念,.
他們不會輸光,.
所以當耶穌說他們不自由的時候,.
他們馬上回應,.
我們是亞伯拉罕的子孫,.
不是任何人的老僕,.
我們知道,.
他們說他們從來沒有做過任何人的老僕,.
這是錯的,.
他們其實是亂說的,.
是不是?.
他們最少在埃及圍奴了400年,.
這是鐵一般的事實,.
再說,.
有一個學者,.
D.A. Carson,.
都說得很好,.
說得很諷刺的,.
他說其實,.
幾乎當世所有的強權,.
無論是埃及,亞述,巴比倫,希臘,敘利亞,羅馬,.
以色列一個都沒有漏,.

$^{761}$都做過他們的老僕,.
以色列做老僕的歷史很長,.
經驗豐富,.
並且很多主子,.
很多不同的主子,.
所以就算是現在,.
耶穌跟他們說話的時候,.
他們都是羅馬帝國的亡國奴,.
怎麼會沒有做過人家的老僕,.
他們說話,.
是不是?.
當然,.
他們的說法可能是,.
他們沒有在宗教上放棄耶和華信仰,.
沒有去跟隨那些統治者的信仰,.
他們仍然相信亞伯拉罕,.
以色列和亞國的上帝,.
在宗教上他們是獨立的,.
猶太人以祖宗亞伯拉罕為榮,.
這個是令人敬佩的事,.
不過他們以為祖宗亞伯拉罕,.
他們的選民身份和神聖的血緣,.
是他們可以依仗的東西,.
這個在耶穌面前就有些麻煩,.
耶穌會不會像他們那樣尊重亞伯拉罕,.
耶穌會不會賣面子給亞伯拉罕,.
抱歉,不會,.
正如我前面說的,.
耶穌只是將亞伯拉罕看為小弟弟,.
是不是?.
怎麼會特別尊重他呢?.
耶穌告訴他們,.
亞伯拉罕對上帝都只能夠俯首聽命,.
如果他們是亞伯拉罕的子孫,.
怎麼可以不聽耶穌的話呢?.
正如我以前多次說,.
我們必須要被打回原形,.
用一個赤裸裸的小孩子的身份,.
才可以近見上帝,.
我們什麼時候都不可以在上帝面前自我介紹,.

$^{801}$我是梁家倫牧師,.
是見道神學院的榮譽院長,.
我是梁家倫醫生,.
我哥哥是梁朝偉,.
我爸爸是前特首梁振英,.
你省一點吧,.
在上帝面前,.
所有身份,職銜,地位,關係,.
通通一文不值,.
連提都不要提,.
多餘的,.
我們在耶穌面前沒有任何人間的依恥,.
只能夠全然降服,.
俯首聽命,.
不要倚仗任何地位和關係,.
也不要倚仗自己的聰明智慧,.
好像曾賢說的,.
你要專心仰賴耶和華,.
不可以倚靠自己的聰明,.
在一切所行的事上,.
都要認定祂,.
祂必指引你的路,.
什麼人最終會拒絕耶穌,.
放棄相信,.
就是他認為自己還有其他路,.
他以為自己有計劃B,.
他以為自己可以留一手,.
多過後門,.
是這樣的,.
耶穌不是他唯一,.
耶穌不是他最大終極的要求,.
覺得自己還有卡可以出,.
有本錢可以賴,.
但願我們每個人都老老實實的認罪,.
專心一意的跟隨耶穌,.
是這樣我們才可以一根到底,.
至死終生..
\newpage



\section{馬可福音 6:45-56}
\label{sec:woNoWOUM52k}
\textbf{2025年8月3日崇拜直播｜吳真真牧師｜耶穌給門徒的十課－祂是誰?｜馬可福音6\_45-56}
\newline
\newline
連結: \href{https://youtube.com/watch?v=woNoWOUM52k}{\texttt{https://youtube.com/watch?v=woNoWOUM52k}} ~~~~ 語音日期: 2025-08-04
\newline
\newline
\hyperref[sec:gyB32xPImr0]{\small{< < < PREV SERMON < < <}}
~
\hyperref[sec:index_chronic]{\small{[返順時目]}}
~
\hyperref[sec:qauwj9FGRZc]{\small{> > > NEXT SERMON > > >}}
\newline
\newline
馬可福音 6:45-56
\newline
\begin{longtable}{cl}
\hline
\hline
章節 & 經文 (和合本修訂版)\\
\hline
6:45 & \begin{tabularx}{0.7\textwidth}{X} 耶穌隨即催門徒上船，先渡到對岸，到伯賽大去，等他叫眾人散去。 \end{tabularx} \\ \\ \relax
6:46 & \begin{tabularx}{0.7\textwidth}{X} 他辭別了他們，就往山上去禱告。 \end{tabularx} \\ \\ \relax
6:47 & \begin{tabularx}{0.7\textwidth}{X} 到了晚上，船在海中，耶穌獨自在岸上。 \end{tabularx} \\ \\ \relax
6:48 & \begin{tabularx}{0.7\textwidth}{X} 他看見門徒因風不順，搖櫓很苦。天快亮的時候，他在海面上走，往他們那裡去，想要超過他們。 \end{tabularx} \\ \\ \relax
6:49 & \begin{tabularx}{0.7\textwidth}{X} 但門徒看見他在海面上走，以為是鬼怪，就喊叫起來； \end{tabularx} \\ \\ \relax
6:50 & \begin{tabularx}{0.7\textwidth}{X} 因為他們都看見了他，甚為驚慌。耶穌連忙對他們說：「放心！是我，不要怕！」 \end{tabularx} \\ \\ \relax
6:51 & \begin{tabularx}{0.7\textwidth}{X} 於是他到他們那裡，一上船，風就停了；他們心裡十分驚奇。 \end{tabularx} \\ \\ \relax
6:52 & \begin{tabularx}{0.7\textwidth}{X} 這是因為他們不明白那分餅的事，心裡還是愚頑。 \end{tabularx} \\ \\ \relax
6:53 & \begin{tabularx}{0.7\textwidth}{X} 他們渡過了海，在革尼撒勒靠岸，泊了船， \end{tabularx} \\ \\ \relax
6:54 & \begin{tabularx}{0.7\textwidth}{X} 他們一下來，眾人立刻認出是耶穌， \end{tabularx} \\ \\ \relax
6:55 & \begin{tabularx}{0.7\textwidth}{X} 就跑遍那整個地區，聽到他在哪裡，就把有病的人用褥子抬到哪裡。 \end{tabularx} \\ \\ \relax
6:56 & \begin{tabularx}{0.7\textwidth}{X} 耶穌所到的地方，或村中、或城裡、或鄉間，他們都把病人放在街市上，求耶穌讓他們摸一摸他的衣裳繸子，摸著的人就都好了。 \end{tabularx} \\ \\
[1ex]
\hline
\hline
\end{longtable}
$^{1}$耶穌隨即催門徒上船,先到那邊白菜大墟,等他叫眾人散開..
他既辭別了他們,就往山上去禱告,到了晚上,船在海中,耶穌獨自在岸上..
看見門徒因風不順,搖擾甚苦.耶律若有四更天,就在海面上走,往他們那裡去,意思要走過他們去..
但門徒看見他在海面上走,以為是鬼怪,就喊叫起來,因為他們都看見了他,且甚驚慌.耶穌連忙對他們說:「你們放心,是我,不要怕.」.
於是到他們那裡上了船,風就止住了.他們心裡十分驚奇,這是因為他們不明白那分餅的是心裡還是如何..
基督過去來到隔離撒勒地方,就靠了岸,一下船,眾人認得是耶穌,就跑遍那一帶地方,聽見他在何處,便將有病的人用玉紙抬到那裡..
凡耶穌所到的地方,或村中,或城裡,或鄉間,他們都將病人放在街市上.求耶穌只用他們摸他的衣裳睡紙,凡摸著的人就都好了..
經訓讀筆.
各位弟兄姊妹平安..
我中學的時候,團契的導師邀請我們,問我們心目中耶穌是誰..
他不單是問,還要我們用畫作表達,我們畫一幅畫,去表達耶穌是誰..
我想問你們有沒有試過讓肆虐有這樣的活動?有沒有?有,有,有,有,沒有.我想隨著年紀大一輩便會做這些事.年輕的未做過,我挑戰一下,如果有中學或大專團契的導師都可以做..
我很記得當時我畫了一幅畫,那幅畫是耶穌牽著我.因為我很深印象,那時候有首歌叫做「耶穌是我的知心友」..
在我當時的年紀,我認識的耶穌,我覺得耶穌是我的好朋友.我有不開心,有困難,我會向祂傾訴,祂就是我的好朋友..
我還記得在我身邊的人,有些人畫十字架,有些人畫心,有些人畫耶穌,但畫出滿身的光芒..
我在預備這篇講座時,我問自己,幾十年後的今天,如果要我再畫這幅畫,耶穌是誰?我會畫些什麼?我會表達些什麼?.
我相信在我們信仰的歷程當中,我們不斷有機會反思耶穌究竟於我,祂是誰?.
雖然我這幾個月不在,但我仍然有看每一篇關於耶穌教門徒的十堂課的講道..
我記得陳牧師說,如何選取從六章開始?就是因為門徒要出去工作,要行動,而不是只聽..
所以就開始了從第六章開始的十堂課.上一次,Ray牧師在十號風球下,大家有沒有聽到,有沒有看到?.
就是他講到五餅而如這個神蹟,這是最偉大的神蹟.這個神蹟在六章三十節一開始時,是講述門徒出門傳道後,回來回報給耶穌,究竟他們傳了什麼道,做了什麼事..
但當時耶穌沒有對他們的任何活動做出任何評價,只是說我們去曠野歇一歇..
但因著耶穌對這麼多群眾,萬多二萬群眾的憐憫,他就開始去講道,慰養他們,最後就有五餅而如這個神蹟的發生..
完了之後,馬可福音的作者沒有講到門徒和群眾有什麼反應,因為如果我們再讀下去,我們會發覺原來五餅而如這個神蹟和耶穌理海的神蹟是很緊扣的..
最後,馬可形容門徒是很愚頑,他們對耶穌的認識仍然是很膚淺..
縱然門徒和耶穌是日夕相對,但他們仍然對耶穌是誰,他們原來不是真的認識..
所以我們今天很想一起讀這段經文,去探討一下究竟耶穌想告訴他們他是誰..
當他們要行動,當他們要傳道,當他們要醫病,趕鬼,去傳揚主的時候,他們究竟知不知道他所傳揚的耶穌是誰?.
讓我們一起去祈禱,天父求你指教我們,幫助我們,讓我們去認識你,並不單單是頭腦的認識,並不單單是坐著去認識,.
而是我們用我們的行動,讓我們的生命的投放,讓我們對你的交託,以至我們對你有更深的認識..
求你開我們的心,引導我們,幫助我們明白你所要我們去明白的..
多謝你聽我們祈禱,禱告交託,奉主明求.阿們..
今天我們會有三個段落,我們會說到耶穌有三個特點,我們很想去知道的..
分別是,耶穌是一個禱告的主,耶穌是一個拯救我們的神,和耶穌是一個服事的主..
我們一開始第45節,我們會看到,原來完了五餅如魚這個偉大的神蹟之後,沒有任何的debelief,沒有任何的檢討,沒有..
耶穌很快地做了一個動作,就是隨即催門徒上船,然後叫他們要渡過加利利海的東邊,去到白菜大這個地方,在地圖上,見到是白菜大這個地方..
然後,與此同時,耶穌自己親自去解散面前的一群信眾,一群群眾..
我心裡一直在想兩個問題,第一,為什麼耶穌要這麼急,馬上催門徒上船,去一個東北面的白菜大的地方,這是第一..
第二,通常耶穌會吩咐門徒給他們食物,但這次,祂自己親自去解散這群群眾,究竟為什麼呢?.
我們在約翰福音第六章,給了我們一個參考,因為在約翰福音第六章,祂有說到一件很重要的事情,就是當群眾聽完耶穌的道,領受完五餅二魚之後,他們有一個想法,.

$^{41}$他們先知,而且他們很想逼他們做王,原來耶穌知道群眾的意願,想逼他們做王,.
但耶穌很清楚知道,他們來到這個世界,不是要做政治上的領袖,只是群眾想他們,.
但祂擔心門徒也受到群眾的壓力所影響,所以打亂了他們做門徒訓練的計劃,於是祂就以極快的速度,叫群眾散開,和叫門徒離開..
第四十六節告訴我們,當他們離開門徒之後,耶穌再一次像馬可福音一章三十五節所說,.
祂在繁忙的事奉中,去到一個安靜的地方,獨自安靜地親近神,向神禱告..
我覺得很有趣,我們禱告,不知道我們禱告最迫切是什麼時候,有時候我們禱告最迫切的時候,.
就是我們可能不懂,很困難,很失意,很想神幫助,我們會祈禱..
但耶穌在這個時刻祈禱,是在祂最成功,實現了最偉大的神蹟之後,祂要祈禱..
祂在祈禱什麼呢?我們沒有仔細地說,不過我很相信耶穌的祈禱,是因著剛剛發生的事..
祂要面對的就是活生生在祂面前的那一萬多人,那群群眾,那群以色列民,對祂的期望,.
希望祂能夠做王的期望,但是又和祂來這個世界,上帝給祂的使命有不同,.
祂的內心會有些掙扎,祂一定是很想尋求神的心意,明白天父的旨意如何,給祂力量繼續向前走..
我覺得很好的提醒,作為神子,作為耶穌基督,祂仍然是切切地尋求神的旨意..
其實我們每個人都會祈禱,我們在教會裡開會時會祈禱,我們崇拜會祈禱,我們團祈會祈禱,.
我們自己個人都會向神祈禱,但是我們反思一下,我們的祈禱不是求神去成就我的旨意,.
而是我們真的有一個謙卑的心去求神的旨意在我身上成就,這是很重要的..
前播道神學院的院長,現任愛丁堡基督教國際神學院的院長楊永祥博士,.
他分享一件事,他說從小到大,他從來沒有想過自己不結婚,.
所以當他二十多歲,三十多歲的時候,他看著身邊的朋友一個一個結婚,他很不服,.
他問上帝為什麼他到現在都沒有拖拍,為什麼他沒有給他一個婚姻,.
不過後來在一個靈會的聽道裡面,講完提到拉薩路的死,.
他聽到彷彿耶穌在跟他說一句話,當他看到瑪利亞不明白耶穌要在拉薩路的身上行更大的事的時候,.
好像在跟他說一句話,因為他們覺得拉薩路死了,這件事不好,.
但是他們看不到耶穌的心意卻是要這件不好的事,要彰顯上帝要做更大的事..
所以楊院長他就清楚明白,眼前的事好像捍衛不好,好像不合我的心意,.
但是他知道上帝有更好的心意,更大的事要在我們身上,他就因言順服,.
不是要我的旨意神去成就,而是神的旨意他願意去配合..
所以我們從耶穌的身上知道,他是一個禱告的主,而且他是願意上帝的旨意臨到他身上的主..
然後我們看到第47節,這個段落告訴我們耶穌很想揭示一個很重要的訊息,.
他就是拯救我們的神,我們看看這段經文是怎樣說的,.
他說耶穌在禱告完結後,一個人去到岸邊,他去到岸邊的時候看到一個情況,.
這個情況就是門徒要上的一堂課,他們遇到一個警方,那個警方說風不順,.
他看見他們風不順,你要留意,這個風不順和《馬學福音》四章同樣在海上,.
耶穌去施行神蹟是很不同的,當時四章的時候,你們看見一班門徒嗎?.
耶穌在船上,但是他們說刮起大風浪,然後他們說船滿了水,然後他們就跟耶穌說我們快要沒命了,.
他們去到一個危險的情況,四章的時候他們在一個危險的情況,耶穌在船上和他們同在,.
但是今天這個警方不同,他說風不順,那個意思不是代表他們有危險,而是代表我們人生總會有一些逆境,.
面對這些逆境,這班門徒他們有努力的,他們怎樣努力?聖經說他們搖腦甚苦,.
搖腦甚苦的意思是他們撐著槳,試圖用槳,用他們的努力去駕駛那艘船,在當中的確被風折磨得很厲害,.
是會很辛苦的,是會很痛苦,是要很出力的,耶穌見到他們在逆境當中,見到他們辛苦,因為這個逆境的時候他們辛苦,.

$^{81}$然後在四更天,四更天的意思是羅馬的一個時間的概念,他們將黃昏六點到第二天的六點,這十二個小時分成四更,.
每三個小時一更,四更天是凌晨三點到早上六點,在這個天黑或快天亮的時候,耶穌採取了一個行動,.
他走在水面上,有個岸上,走在水面上,聖經說一個很重要的事情就是耶穌要超越他們,.
其實是的,我們面對辛苦的時候,聖經說耶穌是看顧的,耶穌是看顧的,他看顧的時候是有一個過程,我們都需要努力,我們都需要辛苦,.
但他看顧不代表他不照顧我們,他不出手是因為他要考驗我們是否信得過他,有人說神答應你的禱告就是因為你信得過神,.
但如果神不答應你的禱告,其實是神信得過你,是神信得過你,要考驗你,大家一起看看漢見這個字,在48章耶穌漢見他們,.
在逆風遙老甚苦,但我們看到第49,50節我們再看到漢見這兩個字,這次漢見這兩個字是誰漢見?是門徒,.
不過門徒的漢見就以為他是鬼,門徒的漢見其實是認不清楚耶穌,門徒的漢見是有限,但耶穌卻是漢見我們,.
所以即使有時在我們的生活中我們看不到上帝,好像神不在,但我們都要學習信靠,信得過耶穌看顧,我們要學習的是等候他出手,.
就好像我們以色列的祖先亞伯拉罕,他75歲離開哈蘭時神應許成為大國,但到二十多年後他的兒子以撒才出生,.
又或是我們看到摩西帶領以色列人出埃及,要帶他們去應許之地,但結果在四十年的曠野生活中,他自己都未嘗入到去,.
其實我們真的很需要耶穌的看顧,但當我們看不到耶穌看顧時,我們真的需要耐心等候,耶穌叫我們等候不是因為他不懂,他不能夠,.
你看看馬可福音,馬可形容的耶穌,就像剛才所說,他催門號,催他們隨即上船,耶穌的動作可以很快,很多形容詞形容耶穌可以很快,.
但耶穌如果不快,你還未看見他,是因為他在考驗我們的信心,叫我們要學習等候..
在這段經文中,耶穌很想展現一個很重要的身份,他很想表明他是上帝,他是神,他是基督的身份,為何我這樣說呢?.
第一,他在海面上行,這是上帝顯現的一個表徵,是上帝顯現的一個情況,在舊約裡面,約八記第九章第八節,.
顯然上帝是獨自鋪張蒼天,步行在海浪之上,創世記第一章第二節,神的靈是運行在水面上,所以耶穌在水面上,.
他要告訴他們,我的神聖,我就是上帝,然後第二,我剛才說過,他在水面上有一個目標,目標是他要超越他們,其實超越是很有趣的,.
為什麼這樣說呢?其實耶穌要幫助他們,為什麼不像上次那樣上了一艘船,叫風停止就可以了?.
但這次是故意要讓他們看到,要超越在他們面前,讓他們看到..
其實這個走過他們的走過,和七十事日本舊約裡面常常說的經過,即是出埃及記所說的經過是同一個字,.
經過是出埃及記十九章裡面,耶和華這樣說,我要顯出我一切的恩慈在你面前經過,宣告我的名,我要恩待誰就恩待誰,要憐憫誰就憐憫誰,.
耶和華在摩西面前經過所用的字眼,就和馬可這個字是一模一樣的,他就是要讓這班門徒意識到,.
其實在他們面前超越他們的這一位,就是滿有慈愛與他們同行的那位上帝,.
就好像摩西和意識的人要啟程出埃及之前,他們很恐懼,但神保證他們在他們面前經過..
然後第五十節,當耶穌看到他們很害怕的時候,他跟他們說你們放心,示我,.
你不要小看這個「示我」這個字,因為「示我」這個字是神的自稱,在出埃及的三章十四節,神在燃燒的荊棘裡面對摩西說,.
我是自由永有的,這個就是「示我」,他跟這班門徒說,我是自由永有的上帝,我和你們同在..
所以凡此種種,我們看到在短短的經文裡,耶穌一而再,再而三地告訴他們,.
我是神,我是神,我是大自然的主宰,我是拯救你們的那個..
不過,耶穌很用心良苦,但門徒卻看不到神的榮耀,四十九節告訴我們,他們以為耶穌是鬼怪,嚇得大叫,然後魂不附體..
你可能會說,是的,黑漆漆的,前面突然有句說話,真的很害怕..
是的,但他們的害怕是正常的事,不過這件事是他們認錯了耶穌的一件事,他們認錯了耶穌..
我們很多時候都會認錯神,有人說這個世界有兩個神,一個是創造我們的神,一個是我們創造出來的神..
我們應該相信創造我們的神,但我們很多時候會依靠我們自己創造出來的神,.
我們創造出來的神可能是金錢,可能是權力,錯信了一些權力,或者我們信自己,甚至是我們營造出來的基督教信仰的信念,但那些是錯的信念也不一定..
除了我們會認錯,認了我們自己創造出來的神做神之外,我們也有時候會錯誤地理解我們的上帝..
有一個猶太人的拉比,因為她的兒子患了超速衰老症,在她十四歲那年,因為她很快就衰老,所以十四歲那年她就已經去世了,因為她老了,她很不明白..
於是她就寫了一本書,叫做《When Bad Things Happen to Good People》,我想有些人也看過這本書,.
她在裡面流露了一件事,原來一般人在苦難中的心態是覺得神不是萬能,神的創造是不完美的,.

$^{121}$但這就是錯誤的觀念,用了我們的想法,用了我們的態度去想我們的上帝,我們錯誤地認識上帝,我們也要留意在我們的生活中,我們有沒有錯誤地解讀了我們的上帝呢?.
但神沒有責怪他們,神沒有罵他們,反而上帝,即耶穌基督是安慰他們,祂說你們放心,你們放心,不要怕..
這句話令我想起《耳塞書》第41章第4及第13節,.
「誰幸做成就者事,從起初宣召歷代,就是我耶和華,我是首先的也與沒後的同在.」.
這裡說到「是我」,宣告祂是上帝,是自由永有..
然後第13節說:「因為我耶和華,你的神必參乎你的右手,對你說不要害怕,我必幫助你.」.
我們認識的這位神不單是很有能力,還有一樣很重要的東西,就是祂必幫助我們,祂叫我們不用害怕,祂憐憫我們,祂叫我們不用害怕,祂必幫助我們..
耶穌請他們不要被祂對他們錯誤的認知擾亂了他們,祂幫助他們除去心裡的恐懼,告訴他們祂有能力拯救他們..
於是耶穌就上船,風就停了.聖經形容他們的反應是非常驚奇,是一個很強烈的字眼,十分十分驚奇..
他們很驚奇的原因是什麼呢?聖經就將分餅的事,就將這裡緊扣起來,他們原來是不明白分餅的事,所以這麼驚奇..
其實分餅的事和耶穌理海的教導,都是和摩西帶領以色列人出埃及,在曠野的經歷很相似..
分餅的事就是告訴耶穌的門徒,神在曠野用嗎哪餵飽以色列人..
而當對比一下文,繼續講下去的時候,會有餵飽四千人,之後耶穌和門徒有最後的晚餐,和他們領最後月月節的晚餐的時候,.
就讓我們很清晰知道,耶穌就是那個彌賽亞,祂擺了一個彌賽亞的筵席,是一個救恩的筵席,是一個耶穌是供應者和拯救者的意思..
門徒不明白耶穌這個身份,他們心裡如幻,如幻這個字在馬考弗音三章五節都用過,.
翻譯出來是說他們的心如石頭那麼硬,而三章五節說的那群人是法利賽人,是一群宗教領袖..
馬可現在很諷刺地用來形容這群門徒,用他們好像法利賽人一樣,心如石頭那麼硬..
其實我們看到馬可用一個很嚴厲的字眼去說他們,門徒是和耶穌一起生活的,他們不是法利賽人..
他們經常聽耶穌說天國的道理,但是他們仍然如幻,仍然是屬靈遲鈍..
如果他們能夠分辨分餅的事,背後就呈現耶穌是上帝的兒子,.
他們就不需要對於耶穌離開,上船說的話,然後叫風停止,變得那麼驚訝..
我們覺得門徒如幻,但是我們細心想想,我們其實是不是一樣?.
我們笑他們,說他們如幻,但其實我們心裡也很如幻..
在十九世紀,有一個很出名的偉大的粉卿家,他叫做卓文..
他小時候已經蒙召去服侍神,他很辛勤地服侍,但他做得很辛苦,而且處處碰釘,他很沮喪..
直到有一天,他實在太辛苦,他不想再做下去了..
當時慕迪就去到教會裡面,去分享神的話,聚會結束的時候,他就向慕迪博士分享他內心的掙扎..
他說我不想再做牧師了..
慕迪問他,如果你不想放下那麼多東西去跟從主,那你是否願意讓神去洗你原意?.
你告訴神,我現在真的不願意,但是我願意讓你作事在我身上,以至到我願意..
其實慕迪鼓勵他,你開放自己的心,讓神在你生命裡面工作..
於是卓文就想了一會,然後他說,我願意跟神祈禱,我願意讓神洗我原意..
後來神真的洗他原意,之後他更被神大大使用,成為一個著名的報導家..
奧古斯丁在這段經文的注釋裡面,他是這麼說的..
他說,神踏浪而來,在他腳下,一切的人生風險都會變成坦途,是一個平坦的路途..
然後他問信徒們,你們有什麼畏懼?.
這句話說得很真實..
在耶穌離開這件事上,我們認不出上帝,我們甚至認錯祂,我們也不明白祂..
但是當祂踏浪而來的時候,祂真的解決了門徒的問題..
祂讓門徒經歷到祂是神,祂是大能的拯救者..

$^{161}$在我們人生裡的逆境,我們的確可能會很害怕,很沮喪,很辛苦..
但是當我們越看不見神,當我們認錯對方的時候,我們就會越來越害怕..
不過,我們看不到祂,神看到,神明白,神拯救..
我們的人生風浪在神的腳下會紫色,然後我們心裡就會有那份平安,繼續有勇氣,有力量去面對很多挑戰的人生..
經文繼續下去的時候,第53節,耶穌和門徒最終不是去到白菜大,而是去到格尼撒勒這個地方..
可能是因為逆風的緣故,他們最終沒有去到加勒利東北面的白菜大,而是去了西北面的格尼撒勒這個地方..
這個地方是一個人口很多的地方,聖經說他們一上岸,那時候應該是天亮了,群眾認得出耶穌..
其實很有趣,相對於船上的門徒,他們認不出耶穌,但是這群群眾一開始就認得出耶穌..
無論耶穌去到哪裡,在村中,城裡,鄉間,所有的病人都被抬出來,放在街市裡面..
為什麼街市會放人在街市裡面呢?其實當時的街市和我們現在的街市很不同..
當時的街市是一個露天的地方,不單只是做買賣交易的地方,那裡還會是人們坐下來打牙交,聊天的地方,甚至是找工作,有人會在這裡找工作,招聘人的地方..
而這裡記載,還有病人被放在街市裡面,這是一個很普遍的現象..
當群眾見到耶穌的時候,他們其實都誤認了耶穌,他們認得那個是耶穌,不過他們不知道耶穌的身份,他們認為耶穌是什麼呢?.
我覺得在這個場景裡面,他們就當耶穌是神醫,於是那些人就抬著擔架,抬著病人,追著耶穌,去尋求他們肉身上的醫治,希望摸一摸他們身上的錘子,就能得到醫治..
同樣地,其實這群人他們不認識耶穌真正的身份,但是耶穌一樣服侍他們,憐憫他們,幫助他們..
聖經形容,當我們摸了耶穌的衣裳,的錘子,就能得到醫治,好得了..
其實同樣地,我們這群門徒,我們這群信耶穌的人,我們都縱然有我們的盲點,我們不認得上帝,我們甚至錯誤地認識上帝,但是他仍然拯救我們..
他仍然願意服侍我們,在我們做門徒的生涯中,我們願不願意這樣服侍我們身邊的人呢?.
當我們今天認清耶穌是誰的時候,其實我們的人生的焦點,我們的生活態度,其實會變得很不一樣..
我們是否真的認識他呢?我剛剛認識這半年,我想起在神學院畢業已經差不多二十年的日子,我還記得四年在神學院的生涯中,有老師說,.
我們來這間學院要讀三樣東西,不是讀書,讀三樣東西,我們是讀神,讀人,和讀自己.讀神,讀人,和讀自己.我覺得這句說話是很對的..
在神學院中,我經歷了很多,我記得在一年級的時候,我懷了孩子,不過不幸小產,與此同時,是他和我一起進入神學院的一個最好的朋友,他硬病復發..
後來在三年級的時候,我懷了孩子,我很深去體會到產後會有荷爾蒙的分泌,會影響我的情緒,會有些鬱結的情況..
在四年級的時候,因為我生小孩,所以我讀多一年,我懷著小孩,我懷著小孩的時候又作小產,我又是不能上學,與此同時,我爸爸患了重病..
我覺得讀神學院是挺辛苦的,我問上帝為什麼讀神學院,他就說讀神學院吧,我已經願意獻身給你,為什麼要捱這麼多苦,比讀書還辛苦..
很多日子我都不願意上學,因為我不是全部住在長洲,我是在九龍居住,我五點多起床上學,我五點多起床真的起不了床,我會沒有力氣,我想問上帝為什麼我會這麼辛苦,為什麼你沒有照顧我,沒有看顧我,我已經全力去讀書..
不過,回首去看的時候,我會覺得這是上帝給我的功課,我記得我們那時候經常說,我們修學分,一科三分,一科一點五分,這一科真的無限分,.
這一科是最寶貴的功課,是沒有分數的功課,但我會覺得這一科是上帝給我的最好的功課..
當時我可能學不見上帝,我會不明白上帝是怎樣的,但是我卻憑著信心去度過了這四年的神學生涯..
我也相信上帝在神學院對我的訓練,讓我有這些經歷,更加健壯我,讓我在未來這二十年的神學生涯能夠靠著祂去度過,對往後我繼續服事,繼續去做神的功,或者去經歷我人生裡不同的逆境和課題有更深的幫助..
所以弟兄姊妹,如果今天我畫一幅什麼圖畫呢?我會畫耶穌李海這一幅圖畫,聖達朗爾尼在他的腳下,一切人生的風浪都變成了坦途..
弟兄姊妹們,我們還有什麼可以怕呢?讓我們一起祈禱..
參加天父,我們謝謝你,你是大能的神,你告訴我們,當我們看不到你的時候,你看顧我們,你保守我們,你拯救我們..
許在我們弟兄姊妹的心裡,我們正在經歷人生的一些逆境,我們有時候都覺得很難捱,還有多少難處在面前等著我們..
但主要我們知道,當你坐在海上的時候,我們的人生就變成坦途,幫助我們,持守對你的信靠,等候你的出手,.
好讓我們這個門徒訓練的生命能夠更加合你的心意,更加能夠將人帶到你的根前,見證你的信實和奇妙的作為..
我們這樣禱告交託,奉主耶穌基督的聖名自祈求,阿們..
\newpage



\section{約翰福音 9:1-41}
\label{sec:qauwj9FGRZc}
\textbf{2025年8月10日崇拜直播｜梁家麟牧師｜彰顯上帝的作為｜約翰福音9\_1-41}
\newline
\newline
連結: \href{https://youtube.com/watch?v=qauwj9FGRZc}{\texttt{https://youtube.com/watch?v=qauwj9FGRZc}} ~~~~ 語音日期: 2025-08-11
\newline
\newline
\hyperref[sec:woNoWOUM52k]{\small{< < < PREV SERMON < < <}}
~
\hyperref[sec:index_chronic]{\small{[返順時目]}}
~
\hyperref[sec:iXL9M6i3TXE]{\small{> > > NEXT SERMON > > >}}
\newline
\newline
約翰福音 9:1-41
\newline
\begin{longtable}{cl}
\hline
\hline
章節 & 經文 (和合本修訂版)\\
\hline
9:1 & \begin{tabularx}{0.7\textwidth}{X} 耶穌往前走的時候，看見一個生來就失明的人。 \end{tabularx} \\ \\ \relax
9:2 & \begin{tabularx}{0.7\textwidth}{X} 門徒問耶穌：「拉比，這人生來失明，是誰犯了罪？是這人還是他的父母呢？」 \end{tabularx} \\ \\ \relax
9:3 & \begin{tabularx}{0.7\textwidth}{X} 耶穌回答：「既不是這人犯了罪，也不是他的父母，而是要在他身上顯出神的作為來。 \end{tabularx} \\ \\ \relax
9:4 & \begin{tabularx}{0.7\textwidth}{X} 趁著白日，我們必須做差我來的那位的工；黑夜來到，就沒有人能做工了。 \end{tabularx} \\ \\ \relax
9:5 & \begin{tabularx}{0.7\textwidth}{X} 我在世上的時候，是世上的光。」 \end{tabularx} \\ \\ \relax
9:6 & \begin{tabularx}{0.7\textwidth}{X} 耶穌說了這些話，就吐唾沫在地上，用唾沫和了泥抹在盲人的眼睛上， \end{tabularx} \\ \\ \relax
9:7 & \begin{tabularx}{0.7\textwidth}{X} 對他說：「你到西羅亞池子裡去洗。」（西羅亞翻出來就是「奉差遣」。）於是他去，洗了，回來就看見了。 \end{tabularx} \\ \\ \relax
9:8 & \begin{tabularx}{0.7\textwidth}{X} 他的鄰舍和素常見他討飯的人，就說：「這不是那從前坐著討飯的人嗎？」 \end{tabularx} \\ \\ \relax
9:9 & \begin{tabularx}{0.7\textwidth}{X} 有的說：「是他」；又有的說：「不是，卻是像他。」他自己說：「是我。」 \end{tabularx} \\ \\ \relax
9:10 & \begin{tabularx}{0.7\textwidth}{X} 於是他們對他說：「你的眼睛是怎麼開的呢？」 \end{tabularx} \\ \\ \relax
9:11 & \begin{tabularx}{0.7\textwidth}{X} 那人回答：「有一個名叫耶穌的，他和了泥抹我的眼睛，對我說：『你到西羅亞池子去洗。』我去一洗，就看見了。」 \end{tabularx} \\ \\ \relax
9:12 & \begin{tabularx}{0.7\textwidth}{X} 他們對他說：「那個人在哪裡？」他說：「我不知道。」 \end{tabularx} \\ \\ \relax
9:13 & \begin{tabularx}{0.7\textwidth}{X} 他們把以前失明的那個人帶到法利賽人那裡。 \end{tabularx} \\ \\ \relax
9:14 & \begin{tabularx}{0.7\textwidth}{X} 耶穌和泥開他眼睛的那一天是安息日。 \end{tabularx} \\ \\ \relax
9:15 & \begin{tabularx}{0.7\textwidth}{X} 法利賽人又問他是怎麼得看見的。他對他們說：「他把泥抹在我的眼睛上，我一洗，就看見了。」 \end{tabularx} \\ \\ \relax
9:16 & \begin{tabularx}{0.7\textwidth}{X} 於是法利賽人中有的說：「這個人不是從神來的，因為他不守安息日。」另有的說：「一個罪人怎能行這樣的神蹟呢？」他們之間就產生了分裂。 \end{tabularx} \\ \\ \relax
9:17 & \begin{tabularx}{0.7\textwidth}{X} 於是他們又對那盲人說：「他開了你的眼睛，你說他是怎樣的人呢？」他說：「他是個先知。」 \end{tabularx} \\ \\ \relax
9:18 & \begin{tabularx}{0.7\textwidth}{X} 猶太人不信他以前是失明，後來能看見的，等到叫了他的父母來， \end{tabularx} \\ \\ \relax
9:19 & \begin{tabularx}{0.7\textwidth}{X} 問他們說：「這是你們的兒子嗎？你們說他生來是失明的，現在怎麼看見了呢？」 \end{tabularx} \\ \\ \relax
9:20 & \begin{tabularx}{0.7\textwidth}{X} 他的父母就回答說：「他是我們的兒子，生來就失明，這是我們知道的。 \end{tabularx} \\ \\ \relax
9:21 & \begin{tabularx}{0.7\textwidth}{X} 至於他現在怎麼能看見，我們卻不知道；是誰開了他的眼睛，我們也不知道。他已經是成人，你們問他吧，他自己會說。」 \end{tabularx} \\ \\ \relax
9:22 & \begin{tabularx}{0.7\textwidth}{X} 他父母說這些話，是怕猶太人，因為猶太人已經商定，若有宣認耶穌是基督的，要把他趕出會堂。 \end{tabularx} \\ \\ \relax
9:23 & \begin{tabularx}{0.7\textwidth}{X} 因此他父母說「他已經是成人，你們問他吧」。 \end{tabularx} \\ \\ \relax
9:24 & \begin{tabularx}{0.7\textwidth}{X} 於是法利賽人第二次叫了那以前失明的人來，對他說：「你要將榮耀歸給神，我們知道這人是個罪人。」 \end{tabularx} \\ \\ \relax
9:25 & \begin{tabularx}{0.7\textwidth}{X} 那人就回答：「他是不是個罪人，我不知道；有一件事我知道，我本來是失明的，現在我看見了。」 \end{tabularx} \\ \\ \relax
9:26 & \begin{tabularx}{0.7\textwidth}{X} 他們就問他：「他給你做了甚麼？是怎麼開了你的眼睛？」 \end{tabularx} \\ \\ \relax
9:27 & \begin{tabularx}{0.7\textwidth}{X} 他回答他們：「我已經告訴你們了，你們不聽，為甚麼又要聽呢？難道你們也要作他的門徒嗎？」 \end{tabularx} \\ \\ \relax
9:28 & \begin{tabularx}{0.7\textwidth}{X} 他們就罵他：「你是他的門徒，而我們是摩西的門徒。 \end{tabularx} \\ \\ \relax
9:29 & \begin{tabularx}{0.7\textwidth}{X} 神對摩西說話是我們知道的，可是這個人，我們不知道他從哪裡來。」 \end{tabularx} \\ \\ \relax
9:30 & \begin{tabularx}{0.7\textwidth}{X} 那人回答，對他們說：「他開了我的眼睛，你們竟不知道他從哪裡來，這真是奇怪！ \end{tabularx} \\ \\ \relax
9:31 & \begin{tabularx}{0.7\textwidth}{X} 我們知道神不聽罪人，惟有敬奉神、遵行他旨意的，神才聽他。 \end{tabularx} \\ \\ \relax
9:32 & \begin{tabularx}{0.7\textwidth}{X} 從創世以來，未曾聽見有人開了生來就失明的人的眼睛。 \end{tabularx} \\ \\ \relax
9:33 & \begin{tabularx}{0.7\textwidth}{X} 這人若不是從神來的，甚麼也不能做。」 \end{tabularx} \\ \\ \relax
9:34 & \begin{tabularx}{0.7\textwidth}{X} 他們回答他說：「你完全是生在罪中的，還要來教訓我們嗎？」於是他們把他趕出去了。 \end{tabularx} \\ \\ \relax
9:35 & \begin{tabularx}{0.7\textwidth}{X} 耶穌聽說他們把他趕出去，就找到他，說：「你信人子嗎？」 \end{tabularx} \\ \\ \relax
9:36 & \begin{tabularx}{0.7\textwidth}{X} 那人回答說：「主啊，人子是誰？告訴我，好讓我信他。」 \end{tabularx} \\ \\ \relax
9:37 & \begin{tabularx}{0.7\textwidth}{X} 耶穌對他說：「你已經看見他，現在和你說話的就是他。」 \end{tabularx} \\ \\ \relax
9:38 & \begin{tabularx}{0.7\textwidth}{X} 他說：「主啊，我信！」他就拜耶穌。 \end{tabularx} \\ \\ \relax
9:39 & \begin{tabularx}{0.7\textwidth}{X} 耶穌說：「我為審判到這世上來，使不能看見的看見，能看見的反而失明。」 \end{tabularx} \\ \\ \relax
9:40 & \begin{tabularx}{0.7\textwidth}{X} 同他在那裡的法利賽人聽見這些話，就對他說：「難道我們也失明了嗎？」 \end{tabularx} \\ \\ \relax
9:41 & \begin{tabularx}{0.7\textwidth}{X} 耶穌對他們說：「你們若是失明的，就沒有罪了；但現在你們說『我們能看見』，你們的罪還在。」 \end{tabularx} \\ \\
[1ex]
\hline
\hline
\end{longtable}
$^{1}$約翰福音九章一至七節.
經文沒有引到秩序表,大家手上的聖經是140頁,或者看螢光幕.
約翰福音九章一至七節.
耶穌過去的時候看見一個人生來是黑眼的,門徒問耶穌說:拉比,這人生來是黑眼的,是誰犯了罪?是這人呢?是他父母呢?.
耶穌回答說:也不是這人犯了罪,也不是他父母犯了罪,是要在他身上顯出神的作為來.
趁著白日,我們必須作那差我來者的功,黑夜將到,就沒有人能作功了,我在世上的時候是世上的光.
耶穌說了者說:就套他抹在地上,用他抹和泥,抹在子的眼睛上,對他說:你往西羅亞遲子去洗,西羅亞翻出來就是鳳差險.
他去一洗,回頭就看見了,聖經讀不完,神賜福,上讀去懷疑.
我們今天說故事,故事好說,比做那些神學論辯好很多.
這是一個很長篇的故事,劇情佔了整整九章的篇幅,牽涉的人物也很多,情節也很離奇曲折,我們要用兩次才能說完.
聖經沒有說到這個故事發生的時間,或者是接著前面的八章五十九節,耶穌和人吵架,被人說要打他,他就匆匆逃跑.
逃避猶太人對他的攻擊,不過如果他和門徒匆匆離開聖殿,他怎會這麼從容在聖殿外慢慢走,還有心情討論神學問題呢?.
這真是令人有點奇怪,所以時間我們其實是不知道的,地點我們也不知道,聖經一開始九章一節是說耶穌過去,過去是什麼意思呢?.
過去,過去是什麼意思呢?過去是從哪裡經過呢?我們也是估計在聖殿,不過只是估計,我們順著這個估計猜想去說,某一次,因為不知道時間,耶穌和門徒在聖殿外碰到一個黑暗的盒子.
宗教場所門外總有一堆乞丐在等候,期待人們敬拜上帝之後善心大增,就會比較慷慨去施捨財物給他們,這個乞丐是獨自一人在這裡行乞,沒有提到他父母在身邊,估計他已經成年,可以自己自立去謀生.
他應該是慣常在聖殿門外行乞的人,因為這裡算是他的地盤,約翰特別著名,他是天生黑眼的,為什麼約翰會知道,為什麼門徒會知道呢?.
唯一的解釋就是,大概去聖殿,出入聖殿的人,常常會碰到他,對他的警方都略知一二,第八節提到有人素常見他時討飯的,即是他們都見慣他.
我們解釋這段經文仔細一點,我們再看第一節,耶穌過去的時候看見一個人生來是黑眼的,雖然第二節就講到門徒提出一個神學問題,顯然門徒都見到那個乞丐.
不過聖經作者刻意在第一節先講耶穌見到那個乞丐,所以就排除了耶穌沒有注意到那個乞丐,是門徒先發現提出問題,才讓耶穌注意他.
雖然提問題的是門徒,但耶穌早已注意那個乞丐,後面講耶穌要在那個乞丐身上顯出上帝的作為,他不是在門徒提出問題之後,然後匆匆忙忙去理會那個乞丐,順從要求做個神蹟出來.
耶穌卻早已留意他,聖經沒有明講,不過他對那個乞丐動了慈心,他在意他關懷,這是我們先講的.
這一趟門徒在見到那個乞丐的時候,突然間就爆出一個神學問題,就問他們的夫子,第二節,拉比,這人生來是乞眼的,是誰犯了罪?是這人呢?是他父母呢?.
乞丐是天生乞眼的,換言之,這排除了他本人因犯罪而導致失明的可能性,他一出生就乞盲,沒機會做某些事去整自己.
他不需要為乞丐承擔責任,除非我們用佛教輪迴的觀念,講過乞丐今生的災難是源自他前世的罪孽,隔代報應.
但基督教沒有前世今生的講法,簡單而牢固的信念,疾病是上帝的懲罰,上帝的懲罰總是因應人的罪而來,所以疾病是罪的後果.
如果這個人為自己的疾病負罪責,簡單的推論就是由他的父母來承擔,他的父母犯罪,所以造成了兒子就瞎了眼.
遇見人間的苦難,很多時候我們會感覺這個苦難是構成對我們信仰的威脅,全能的上帝怎能容許苦難發生呢?.
為了消去這個威脅,我們必須要提出一個神義論,為上帝撇清責任,苦難是人自招的,是人犯罪帶來無法避免的結果,所以人是活該的,跟慈愛的上帝無關.
此刻門徒對那個瞎眼的人,第一個反應是這個人瞎眼是自招的,不要怪上帝,然後再問要找誰負責,他們給耶穌兩個選擇,核子本人還是他的父母.
耶穌的回答令人意外,none of the above,他一樣都不選擇,既不是核子本人犯罪,又不是他的父母犯罪,第三節說,是要在他身上顯出上帝的作為來.
面對苦難,特別是人間的苦難,我們第一個反應不是尋求解釋,而是努力去解決,幫助人去解除苦難,或者最少讓他在苦難中過得好一點.
此生事外的無聊討論,idol talk是沒有意義的,身邊的人的苦難不是構成一個信仰的威脅,正如剛才說的,要我們盡快去自保,為上帝的正義辯解,而是構成一個召喚.
一個召喚,我能夠為他做些什麼,正如耶穌吩咐門徒,你們給他吃飽,給他們吃飽,耶穌說要在核子身上顯出上帝的作為來.
他的意思明顯就是要在此時此刻顯出作為,有所作為,他要做些什麼在核子身上,耶穌的作為就是上帝的作為,面對核子的苦難,耶穌準備做事了.
耶穌接著對門徒的教導也說明了這個意思,第四到第五節,趁著白日,我們必須作那差我來者的功,黑夜將到,就每人能作功了,我在世上的時候是世上的光.
這段說話的訊息有兩個,第一,趁著白日,就是把握時機的意思,只有白日能夠做事,黑夜就不能夠做事了,所以眼前的苦難就是服侍的契機.
多講無謂,行動最實際,口水改變不了這個世界,汗水和淚水才可以,這是第一個意思,做事了.
第二,耶穌說,他做的功,我們所做的功,就是那差我來者的功,我們被呼召進入人群去關愛他們,我們都是上帝的使者,耶穌的作為就是上帝的作為.
我們的作為,都是上帝的作為,我們不是做自己的功,我們是做差我來者的功,很多人在面對苦難的時候,總是問上帝在哪裡,上帝是否沉默,上帝是否缺席,甚至宣稱上帝是否死了.
耶穌提醒我們,真實的問題是,你在哪裡,基督徒在哪裡,我們在這裡,上帝就在這裡,教會在這裡,上帝就在這裡,我們用自己的作為來彰顯上帝的作為.

$^{41}$耶穌說,祂自己就是世上的光,祂也呼召我們成為世上的光,聽清二論,耶穌給門徒的答案雖然很有智慧,不過其實是迴避了問題,門徒問的是苦難的因由,而不是在面對苦難的時候應該有的態度.
耶穌是將因由轉變成態度,不過有人很奇怪,他們會將耶穌的答案勉強變成對門徒的回答,我要說,耶穌是答非所問的,如果他強行弄到耶穌不是答非所問,我們就會以為耶穌說.
這個人之所以盲眼,是因為上帝預定了他做盲眼,好讓耶穌在這個時候醫治他,即以顯出上帝的作為來,聽不聽得明?.
耶穌是預備這個盲眼的人做為一個演員,好配合耶穌顯出上帝作為的演出,教會裡面的壞鬼神學,確實常常附含這個意思,某個人的母親被汽車撞死,我們就解說上帝撞死他母親,是要讓他發現生命脆弱無常,於是他就可以放下他原有的工作,專心去傳福音.
為了讓一個人去傳道,上帝就故意撞死他母親,這個說話就算是真的也說不出口,很難說,這麼奇怪的東西怎麼說?.
是的,上帝確實可以使壞事變成好事,在苦難之中顯出他的作為,但是這個不等於上帝故意去炮製壞事來成就好事,在彰顯美善之前才製造悲慘,耶穌說要在核子身上顯出上帝的作為,意思不是說那個人的盲眼是上帝的作為,.
上帝刻意將他的眼睛弄盲,而是說上帝剛剛好在這個時候預備了他,讓他的盲眼,耶穌對他的醫治,湊合成耶穌施展神蹟大能的場景,確實沒有一個盲眼的人,耶穌怎麼醫核眼呢?.
沒有瘸腿的,耶穌怎麼使人起來行走呢?沒有傷心難過的人,耶穌怎麼可以成為人的安慰呢?不過這個不是等於耶穌為了使人有復明的神蹟,就首先上帝佈置讓這個人出生的時候,.
姑且說他如今二十歲,在二十年前就已經預先弄盲他的眼睛,不要說上帝的動機,弄盲人是不是善良呢?對二十年來都活在黑暗的人是不是公平呢?.
耶穌說上帝好像處心積慮了一點,太恐怖了,為了二十年後的神蹟,今天才弄盲了他,弄盲二十年,為了一個兩分鐘就完成的醫治事件,二十年前就開始這樣佈局,當然沒有人能否認上帝有可能這樣長期佈局的,.
預定那個人是盲的,但我不認為耶穌說顯出上帝的作為來是指著上帝早已經安排了他做盲的人,他之所以成為盲的人是上帝的作為,這個不是耶穌的意思,有沒有這樣的事實我們不確定,但耶穌的話沒有包含這樣的意思,.
我們看到耶穌的答非所問不是耶穌誤解了門徒的問題,耶穌的轉述這麼快,怎麼會誤解人呢?而是耶穌刻意要轉移話題,終止門徒就苦難的緣由,造成苦難的責任垂縮,這個討論終止它,不要讓他們糾纏下去,.
耶穌相信門徒的討論是毫無價值的,你所有的假說,又無法證明又無法否證,不可以verify,不可以falsify,所有的解釋都是沒有意義的,再說一個卸責的動機不會帶來一個負責任的人生,.
神議論的討論不會為我們建立一個正確的態度,更加不會驅動我們做正確的行動,耶穌拒絕解釋,面對問題是尋求解決,是解決不是解釋,有人可能會說沒有正確的診斷,何來正確的處方呢?.
例如有人肚痛,我們總是需要知道是哪個器官出了問題,才能對症下藥,所以為苦難尋找來源,確知問題所在,確定罪魁禍首,是不是才使我們能夠解決苦難呢?.
沒有解釋,如何解決呢?我同意在自然科學的層面,尋找問題的根源是解決問題的關鍵前提,不過我必須說,在社會科學或人民科學,包括心理學甚至神學,我非常懷疑人們所張羅的解釋,究竟是診斷還是猜想,conjecture..
某個人自我形象低落,是因為他三歲時曾經被父親罵過一句「衰仔」,你告訴我,這個推斷有證明或否證的可能性嗎?.
必須說,坊間很多心理學或神學的解釋,跟黃大仙的籤文其實沒什麼兩樣,所謂靈驗度都是靠穿鑿附會,自行對號入座而得的,靠這些神棍式的解說,要找出正確解決之道,我要說是緣木求魚..
必須特別說,神學解說的任意性其實是很顯著的,例如我患了癌症,有人告訴我這是上帝命定安排的,我不反對,我一生都在他手裡,沒有他容許,任何事情都不會發生..
不過,知道這是上帝命定的,那我還要不要治療?上帝弄到我患癌症,我是否應該信服就不要治療?你明白我的意思嗎?那要治療還是不治療?.
所以知道是上帝命定的,那又如何?治療是不是違背他的旨意呢?神學解說對我面對眼前的困境有什麼幫助呢?.
又舉多一個例子,以前的例子,真實例子,我幫一個弟兄處理他和我太太的婚姻危機,他突然問我一個問題,他是牧師,我和這個女人的婚姻會不會不是上帝的旨意呢?.
我馬上就質問他,兄弟,你問這個問題是完全沒有營養價值的,你是想為你要離婚製造藉口,都結了婚了,問是不是上帝的旨意,你想怎樣?.
我們中間在座的弟兄,你敢不敢問這個問題?你問這個問題,真是超級無聊,超級荒謬,問來幹什麼?.
我當然相信萬事萬物背後都是上帝的旨意,他是我一生的編劇和導演,但確認眼前發生的一切是出於上帝,不是要我熱衷地做各種神學的猜想和討論,而是叫我默然不語,不說話,知道一切出於他..
無論如何,很多解說真的比費特拿更廢,我們事奉的一個活著的上帝,一個living god,我們必須確認他在我們生命中的主權,進行他已經顯明的旨意,不要猜上帝在想什麼,我最討厭猜這個東西..
他生了小孩,猜什麼?你又要問他,他肯說,你就知道答案,他不肯說,你猜也沒用,上帝的旨意就是你猜不到,所以拒絕那些無謂的空談,那些god talk,是沒有營養的..
回到聖經的故事,整個故事要傳遞的訊息是,生命中有太多東西是mystical,是奧秘性的,無法在人間得到合理和充分的解釋,不一定和人的罪有關,不一定和撒旦的作為有關,不一定是上帝促成的,總之就是奧秘..
奧秘就是奧秘,企圖為生活每個事件,每個細節,尋求人間或屬靈的解釋,總是徒勞無功,徒添煩惱,所以又是無功,又是無益..
但是,上帝在這裡,上帝在這裡,他對一切有他的主權和計劃,他可以使壞事變成好事,他可以利用一件命命看來是壞事,成就他未善的旨意,顯出他的作為來..
實際上,整個醫治的神蹟,在《若人福音》這裡,確實是彰顯著上帝的作為,特別是彰顯了耶穌的作為和耶穌的身份..
正如有不少學者指出,整本舊約都沒有記載任何先知使盲眼的人服命的故事,聖經多次提到使眼睛服命是上帝自己的作為,這亦是彌賽亞的行為標記,這是彌賽亞先知反覆強調的..
所以,耶穌醫治的這個乞子,是宣示了他的彌賽亞身份,亦是暗示他的神聖,耶穌是藉著這個神蹟去說明他的彌賽亞身份..
你看,《馬太福音》11章4到5節,耶穌如何介紹自己呢?.
「你們去把所聽見所看見的事告訴約翰,就是乞子看見,缺子行走,長大麻風的潔淨,聾子聽見,死人服門,窮人有福音傳給他們.」.
乞子看見,是證明他是彌賽亞的其中一個證據.然後,耶穌開始上帝的作為,醫治就是上帝的作為..
首先,他套口水在地上,然後用口水和地上的泥,搓成一個泥團,然後用那個泥團來抹在乞子的眼睛上..
還不夠,他吩咐乞子去西羅亞的池子洗一洗,乞子聽命去洗眼,結果就看見,你不覺得奇怪嗎?.
不是耶穌一開口,乞子馬上看見,而是還要去指定地點洗眼..
你翻閱福音書,這個大概是耶穌所行的眾多神蹟裡面,手續最複雜最繁瑣的一個,無之一,最繁瑣的就是他..

$^{81}$即使耶穌叫人從死人中復活,也只是向墳墓大喊一聲,拉薩路出來,事就將成了..
耶穌有很多更偉大的神蹟,例如用五餅二魚餵飽五千人,讓彼得去捕獲滿滿兩船的魚,都不需要複雜的程序,又不需要特別的手勢,又不需要特別的咒語..
他其他的醫治個案,例如醫治好彼得的惡母,也沒有必要的手勢咒語,畢竟耶穌說有就有,命立就立..
他心想就是性.長大麻風的人怎麼跟耶穌說?你若肯,只要你肯,就能夠叫我全如.肯就行了,耶穌怎麼說?我肯,那麻風的人就全如..
那白膚長阻止耶穌不去他家,為什麼?只要你說一句話,我的僕人就會全如..
他不用去視診,不用親自去他家,不用身體接觸,什麼都不要,你說一句話,我的僕人就會全如.所以,你告訴我,死人復活容易,還是死人復明容易?.
好像死人復活難一點,是不是?如果復活也不用什麼手續,復明為什麼要用這麼多手續?我個人相信,所有什麼吐口水,和泥團,抹眼睛,洗眼睛的行動,其實都是多餘的..
不是醫治這個核子的必要步驟,耶穌之所以搞一個這麼複雜的做法,目的只有一個,就是要將這一次的醫治事件變成一個公開的醫治行動,高調的破壞安息日不能夠工作的規定..
即是對法利賽人提供一個你攻擊我的罪證,這是為後民賣下的瀑布..
是的,耶穌親手奉上讓人攻擊他的把柄,如果耶穌只是用口頭吩咐核子復明,或者耶穌好像那個女人,是有血流的,不小心,他的能力流到核子身上,那不構成工作,那耶穌就沒有安息日,沒有安息日,那後民就沒有話說了..
耶穌是故意搞這麼複雜,就是告訴你,我再犯安息日中..
西方的巫師,中國的道士在行法術的時候都需要法寶道具,需要唸咒語,需要有手勢,需要搞一些大套禮儀的,耶穌是上帝,行神蹟,根本不是能力的問題,是主權的彰顯..
所以什麼禮儀,什麼道具,全部都不用.耶穌,我以前都講過,他根本不是要證明他有醫治的能力,他要證明他是生命的主,他有權..
所以在這裡他要超視,我是安息日的主,主權是重於能力,認識耶穌的主權比認識耶穌的能力更加重要..
所以耶穌在這裡,在眾目睽睽之下,做了一個很複雜的醫治行動,我就是要你看見,我不守安息日,我不洗手,因為安息日是我的日子,我話事的日子..
故事第一部分就講完了,我們做一些信仰思考.基督教不是印度教,不是佛教,我們不鼓吹宿命論的,我們不會講凡存在就是合理的,因為這是一個被罪扭曲了的世界,這是一個世界之王在示弱的世界..
所以我們不會講這裡是有慈愛,有公義,有理性的.不過,撒旦和罪不能夠合制上帝的作為,此時此刻,應該講任何此時此刻,都是上帝彰顯他作為的奇蹟..
在一個再不合理,再被罪惡權勢,遊論踐踏的處境裡面,上帝的旨意仍然可以成就,上帝的慈愛和公義仍然可以彰顯,並且除此時此刻之外,也不要讓我們期盼,等候上帝彰顯他大能的場景..
香港何時見到上帝的作為? 今日何時見到上帝的作為? 今日沒有其他時間,沒有更多的時間.我們如何在一個不理想的現實看到上帝的作為呢?.
今日的經文告訴我們有兩個途徑,第一個是耶穌,第二個是我們自己. 耶穌,我們自己,先講耶穌,對於那個天生盲眼的人,他早已習慣活在黑暗裡,活在社會的最底層,他早已認了命,對生活沒有期待,對上帝早已絕望..
在這樣的情況下,要求他自己有所突破是不可能的,必須要有耶穌介入,闖入他的生命裡,否則他不可能就這樣讓自己的生命變成神跡,必須要有神跡施行在他身上..
如今,耶穌果然在他身上顯出上帝的作為,他的眼睛就能夠看見. 對於那些活在極度困難,覺得人力無法回天的環境的人來說,好像核子一樣,絕望的感覺是非常普遍的..
我們多了堅持,多了闖關,都無法突破困境,改變現實,我們已經心力交瘁,無力回天,你再跟我說那些哪吒的名言,我命由我不由天,都是廢話..
我這幾個月重操故業,回到國內做了好幾次培訓,是我最喜歡做的事情,我十八歲為中國獻身..
我接觸國內的家庭教會的領袖,他們很多人告訴我,他們中間有憂鬱症的是很普遍的.畢竟這麼多年,不斷面對高壓的政治迫害,折騰威嚇,並且眼前還看不到曙光..
所以心智不夠堅強的早已倒下,我分享的幾位同工都說他們自己都有憂鬱症.我聽了非常非常擔心,真的不停為他們祈禱..
我亦都希望你們為中國家庭教會的領袖祈禱,如果沒有上帝直接戒人,沒有神職的愛力幫助,他們自己要他們撐住,要他們堅強些,要他們自己解決問題,是很難的,是很難很難很難..
基督教存在在世界上,面對各種政治社會文化的敵對勢力,我們手無寸鐵,無力反抗.如果我們不相信上帝仍在,耶穌基督是教會的元首,如果我們不相信上帝戒入,神職可以期待,我們早已投降絕望了..
弟兄姊妹,在我們走投無路,自覺無力正乾坤的時候,我們呼求上帝顯出他的作為,這是我們最後的依詞,我們要在活人之地見到耶和華上帝的面,否則我們早已喪膽了..
這是第一個途徑,耶穌自己戒入,耶穌顯出作為,第二個途徑是我們自己,耶穌固然是人類拯救和出路,他可以施展神職大能改變現實,不過在我們剛才讀的聖經故事裡,耶穌不僅僅是要門徒仰賴他,期待他施展上帝的作為,.
耶穌對門徒說的是,趁著白日,我們必須作那差我來者的功,黑夜將到,就沒有人能夠作功,耶穌不是說只有他作功,他用的是我們,不是我,我們必須作上帝的功,彰顯上帝的作為..
耶穌是世界的光,他同樣告訴我們,我們都是世界的光,我們必須造炎造光,所以基督徒在人間所有的參與,包括我們和世界對著幹,和現實頑抗,本身就是上帝作為的彰顯,我們不僅僅是等待上帝的作為,我們的作為就是上帝作為的一部份,我們不是繞著手等運到,什麼都不做,.
我們的不甘心,不服輸,不隨眾,拒絕不接受現實的限定,拒絕放棄,不受敵人威嚇,堅守原則立場,通通都是上帝的作為,我們經常讀的經文,保羅在菲勒比書所說的,是一個彎曲薄茂的世代,這個彎曲薄茂的世代,我們做上帝無瑕疵的兒女,.
我們顯在這世代將好將明光照耀,將生命的道表明出來,我們活在這裡,我們的表現就是生命之道的彰顯,我們這樣做就是上帝的作為,.
這幾年,每次去英國和移居英國的弟兄姊妹說,抬起頭,露過四萬的笑臉,你們彰顯上帝的作為,我和留在香港的弟兄姊妹說,我們一起唱林子祥的歌,懷著愛與恕的心,不響夢想,不退讓,讓歌聲讓歡笑,來博之音的讚賞,.
我們同樣彰顯上帝的作為,耶穌說,我在世上的時候是世上的光,這句說話不僅適用於耶穌,也適用於我們,我在旺角的時候,我是旺角的光,我在深水Po 的時候,我是深水Po 的光,我在羅馬的時候,我是羅馬的光,你們在任何一個地步,都是那個地方的光,.
要總結了,說得很長,我們今天看了一個核子的故事,一個彰顯上帝的作為的故事,故事裡面的核子向來只關心今天有沒有人施捨,夠不夠錢吃飯,從來沒有想到他天生的殘缺得到痊癒,.
更加沒有想到他竟然成為耶穌宣示他的彌賽亞身份的一個渠道,在他身上彰顯上帝的作為,他以為他的命是一個奏坐,只是彰顯上帝的奏坐,想不到竟然是彰顯上帝的作為,.
小人物的我們,只活在當下,只能投入眼前的喜怒哀樂裡面,什麼上帝計劃,這個完全超過我們的思想範圍,在此間,一個不理想的此間,竟然是上帝施展作為的平台,這個真的讓人詫異,.
我的生命是上帝施展作為的平台,不過,還有,核子不僅僅是被動接受耶穌在他身上施展作為,我們下次會說的,他面對種種的際遇,包括來自法利賽人的威嚇,壓迫,他的堅守原則,不肯妥協,他自己都主動彰顯上帝的作為,.

$^{121}$一個無知小民,竟然在宗教權貴面前不屈不撓,這是多大的上帝作為呢?這個我們下次說,不管怎麼樣,我們是上帝施展作為的平台,我們也是主動去施展上帝作為的出口,上帝恩待我們每個..
謝謝大家.
\newpage



\section{馬可福音 7:24-37}
\label{sec:iXL9M6i3TXE}
\textbf{2025年8月17日崇拜直播｜鄭國邦傳道｜耶穌給門徒的十課－離開群眾,真認識主｜馬可福音7\_24-37}
\newline
\newline
連結: \href{https://youtube.com/watch?v=iXL9M6i3TXE}{\texttt{https://youtube.com/watch?v=iXL9M6i3TXE}} ~~~~ 語音日期: 2025-08-18
\newline
\newline
\hyperref[sec:qauwj9FGRZc]{\small{< < < PREV SERMON < < <}}
~
\hyperref[sec:index_chronic]{\small{[返順時目]}}
~
\hyperref[sec:0QX__zW2yAY]{\small{> > > NEXT SERMON > > >}}
\newline
\newline
馬可福音 7:24-37
\newline
\begin{longtable}{cl}
\hline
\hline
章節 & 經文 (和合本修訂版)\\
\hline
7:24 & \begin{tabularx}{0.7\textwidth}{X} 耶穌從那裡起身，往推羅境內去，進了一家，他不願意人知道，卻隱藏不住。 \end{tabularx} \\ \\ \relax
7:25 & \begin{tabularx}{0.7\textwidth}{X} 立刻有一個婦人，她的小女兒被污靈附著，一聽見耶穌的事，就來俯伏在他腳前。 \end{tabularx} \\ \\ \relax
7:26 & \begin{tabularx}{0.7\textwidth}{X} 這婦人是希臘人，屬敘利亞的腓尼基族。她求耶穌從她女兒身上趕出那鬼。 \end{tabularx} \\ \\ \relax
7:27 & \begin{tabularx}{0.7\textwidth}{X} 耶穌對她說：「讓孩子們先吃飽，拿孩子的餅丟給小狗吃是不妥的。」 \end{tabularx} \\ \\ \relax
7:28 & \begin{tabularx}{0.7\textwidth}{X} 婦人回答：「主啊，桌子底下的小狗也吃小孩子的碎屑呀！」 \end{tabularx} \\ \\ \relax
7:29 & \begin{tabularx}{0.7\textwidth}{X} 耶穌對她說：「憑著這句話，你回去吧，鬼已經離開你的女兒了。」 \end{tabularx} \\ \\ \relax
7:30 & \begin{tabularx}{0.7\textwidth}{X} 她就回家去，見小孩子躺在床上，鬼已經出去了。 \end{tabularx} \\ \\ \relax
7:31 & \begin{tabularx}{0.7\textwidth}{X} 耶穌又離開了推羅地區，經過西頓，就從低加坡里境內來到加利利海。 \end{tabularx} \\ \\ \relax
7:32 & \begin{tabularx}{0.7\textwidth}{X} 有人帶著一個耳聾舌結的人來見耶穌，求他為他按手。 \end{tabularx} \\ \\ \relax
7:33 & \begin{tabularx}{0.7\textwidth}{X} 耶穌領他離開眾人，到一邊去，就用指頭探他的耳朵，吐唾沫抹他的舌頭， \end{tabularx} \\ \\ \relax
7:34 & \begin{tabularx}{0.7\textwidth}{X} 望天嘆息，對他說：「以法大！」就是說「開了吧！」 \end{tabularx} \\ \\ \relax
7:35 & \begin{tabularx}{0.7\textwidth}{X} 他的耳朵立刻開了，舌結也解了，他說話也清楚了。 \end{tabularx} \\ \\ \relax
7:36 & \begin{tabularx}{0.7\textwidth}{X} 耶穌囑咐他們不要告訴人；但他越囑咐，他們越發傳揚。 \end{tabularx} \\ \\ \relax
7:37 & \begin{tabularx}{0.7\textwidth}{X} 眾人分外驚奇，說：「他所做的事樣樣都好，他甚至使聾子聽見，啞巴說話。」 \end{tabularx} \\ \\
[1ex]
\hline
\hline
\end{longtable}
$^{1}$經文是在《馬可福音》七章二十四至三十七節.
耶穌從那裡起身往推羅西頓的境內去.
進了一家不願意人知道卻隱藏不住.
當下有一個婦人她的小女兒被污鬼附著.
聽見耶穌的事就來伏伏在她腳前.
「這婦人是希尼尼人,屬敘利菲尼基族」.
她求耶穌趕出那鬼離開她的女兒.
耶穌對她說「讓兒女們先吃飽,.
不好拿兒女的餅叼給狗吃」.
婦人回答說「主啊,不錯,.
但是狗在桌子底下也吃孩子們的碎渣兒」.
耶穌對她說「因這句話,你回去吧,.
鬼已經離開你的女兒了」.
她就回家去見小孩子躺在床上.
鬼已經出去了.
耶穌又離開,又來了推羅的境界.
經過西頓,就從底加玻利境內來到加利利海.
有人帶著一個耳聾舌結的人來見耶穌.
求祂按手在他身上.
耶穌領他離開眾人到一邊去.
就用指頭探他的耳朵,套托抹他的舌頭.
望天嘆息對他說「已發大」.
就是說「開了」.
他的耳朵就開了,舌結也解了.
說話也清楚了.
耶穌祝福他們不要告訴人.
但他越發祝福,他們越發傳揚開去.
眾人憤外嬉奇,說他所作的事都好.
他連聾子也叫他們聽見,啞巴也叫他們說話.
聖經獨白.
請姐妹平安.
可以繼續和大家分享馬可科音裡面.
耶穌對門徒的一份教導.
以致讓他們真正認識耶穌.
今天在他的工作,昔日他和一班門徒.
做了很多事,究竟耶穌背後的心意是什麼呢?.
以致我們知道門徒當時對耶穌的認識.
藉著他手所作的.
這份認識慢慢變成認信.
因為認信和認識是不同的.

$^{41}$認識可能是你見到的事.
認信可能是一種從心底裡的相信.
拿今天的生活比喻.
你現在知道一號風球.
你見到風嗎?.
你見到它在哪裡嗎?.
你看天文台的apps,你見到它在這個位置.
但你真的未必親眼看到.
但你相信.
你知道天文台以往的歷史,科學等等.
你會說,是呀,正在來.
這份信也會影響你既有的行為.
或者外面的認識.
你可能知道一號風球.
你可能要帶一把傘出街.
或者在露台有植物.
你也會拿回來.
等等準備.
風來了.
或者有些轉變.
所以由認信到認識.
到一些行動.
會因著這些而有所改變.
那份認信不單止在外面.
有時候那份認信可能在裡面.
有時候不知道他心底裡會不會有些說話.
會跟你說.
我還是不行了.
我還是做得不好.
不夠厲害.
而這份認信影響了你.
可能在做事上.
或者一些抉擇.
不要試了.
我還是不行了.
而在當中也影響了你每天的行徑.
或者一些選擇.
當然除了外面的天氣.
裡面的一些感覺.
或者一些對自己的相信.

$^{81}$對人其實也很普遍.
我們也可能很經歷到這份信心.
和行為裡面的交戰.
例如什麼呢?.
可能你在職場裡面.
有人可能交一些東西給你.
特別是那些senior.
就是那些前輩.
交一些東西給你做.
不同的認信.
可以有不同的反應.
或者不同的行為.
如果你覺得那個senior.
那個前輩給你做一些事情.
你給了我一些事情做.
我學了.
其實是想我拿起來.
然後你就不用做而已.
那你通常會怎樣?.
死死地氣也要做.
因為他是你的前輩.
他是senior.
所以也要做.
但不是的.
如果你的認信也可能會想.
不是的.
可能他真的有些事情教會了我.
將來我做另一個project.
或者另一些事情.
會方便了.
好啊.
現在我去學一下.
知道了怎樣可以幫助到我.
以至這個心態.
這個認信不同了.
可能你做出來的行徑.
或者當時一直想著做的時候.
那個表現也不同.
所以這個就是剛才所說的.
你的認信或者你的相信是什麼.

$^{121}$其實很影響了.
你怎樣去認知或者認識那件事情.
以至你走出來的時候.
也會有不同的結果.
還有.
這份信的影響.
不單止在外面.
剛才說的拿傘.
做事的方式.
或者每個選擇是一些外面的事情.
其實這份信也同樣影響了你私底下.
沒有人看到的生活.
沒有人知道的選擇.
拿一個簡單的例子.
如果你身邊有些人做投資或者炒賣的時候.
你就是在升的.
買啊.
你一定會賺的.
然後你問他.
你有沒有買.
沒有啊.
你買吧.
你會看到他的信其實是說的.
他私底下也不會買的.
如果是賺的.
他就全力去做.
所以.
如果這份信只是在表面的時候.
其實也不是很深.
如果這份信深的時候.
一定會影響你私底下一些的決定.
如果將這個概念放在信仰裡面.
也是一樣的.
如果你相信耶穌.
可能你認識一些弟兄姊妹.
說我相信了.
但是在心中就行了.
你會發現.
你沒有相應的行動.
是不是相信呢?.

$^{161}$但是倒過來看.
你有這些信仰的行為.
有些人會去教會.
有些人會去小組.
或者像上星期.
聽一下培靈會等等.
這些是好的.
不是說不好的.
做了這些信仰背後.
但是在你的心底裡面.
那份認信.
有沒有產生一些改變呢?.
在你心底裡面.
甚或乎你私底下的生活.
沒有人看到的了.
不用回來.
不用交代的時候.
這個信是否在你的私底下.
也有這份影響呢?.
以前讀神學的時候.
有位老師提過.
大家回來.
可能有不同的屬靈操練.
有不同的屬靈生活.
查經,靈修等等都會做.
但是他說.
一個人是否敬虔呢?.
不是在這些場合.
而是在你自己一個人.
你去看回.
你就知道自己是否一個敬虔的人.
你自己一個人的時候.
你做什麼呢?.
沒有人知道的時候.
你在看什麼呢?.
或者在裡面.
你每一個選擇是什麼呢?.
因為好像今天.
大家在網絡上.
就看到.

$^{201}$平時為何一個人私私文文的.
為什麼上網說話這麼兇狠呢?.
或者以前覺得.
上網做了什麼都沒有人知道.
所以說很多話.
私底下他就不會說了.
當然現在有了起底組的.
大家都不敢說了.
以前跟一些弟兄.
都經歷過這樣的事情.
我不知道大家.
我中學的時候.
大家有沒有經歷過.
教會界有些聖潔運動.
大家弟兄要守潔心清.
但是那時候有一群弟兄就開玩笑.
你做這件事.
你夠膽的.
就開你電腦的瀏覽記錄給我看看.
看看你看過什麼.
當然今天現在有什麼無痕的.
所以就看不到你看過什麼.
但是那時候我覺得.
是的.
當你自己一個人.
在對著電腦.
或者一些事情的時候.
你做什麼.
沒有人知道.
但是在那個位置.
就反映你的信.
是怎麼影響你.
所以今天.
即使我們有這些信仰的行動.
或者有信仰的行徑.
這些正在做的事情.
可能我們都很想藉著這些.
看得到的事情.
跟得到的形式.
去呈現我們在裡面看不到的信心.

$^{241}$不是錯的.
真的是好的.
但是我們就好像剛才唱詩歌一樣.
我們有沒有不斷去審查這些事情呢?.
在來來回回.
在裡面我們正在做的事情.
還有我們裡面的那個人.
或多或少其實有時候你去審查.
可能是做給別人看.
也可能是做給神看.
本身就是那句.
沒有說錯的.
你做給別人看.
可能你真的是一個見證.
很想藉著你做出來的行為.
去反映你裡面看不到的信心.
本身是一件好事.
但是慢慢地.
這些外在的.
正在做的事情.
所謂的.
就是那些 doing(做).
正在做的事情.
會不會成為了你的一些誡命.
一些傳統.
或者一些框架.
一開始你很想.
是藉著守這些誡命.
做這些事情等等的形式.
來告訴自己.
是的.
我就是跟從著上帝.
但是慢慢地.
會不會為了守這些方法.
這些傳統.
反而反過來.
廢掉了誡命或者教導的原意.
B型裡面的意義.
那個價值觀.
為什麼會說這個呢.

$^{281}$如果你看回上一個段落.
今日經文的上一個段落.
其實本身.
耶穌就是在回應人.
如何認識上帝.
如何認順上帝.
當時對著一群.
法理,財人和文士的時候.
他們的認順.
就是在誡命.
傳統裡面.
所以耶穌知道.
一提醒他們.
你裡面是什麼的時候.
他們就不聽.
你不要理我.
總之我做了這些出來.
我做了這些就是敬虔.
我裡面是怎麼想你不知道的.
這些部分交給Ray.
遲些下個星期就會補充.
因為我們的次序調轉了.
但今日我們會看到.
耶穌在這群外邦人的身上.
做的是另外一些事情.
因為耶穌看到每一個人的行為.
背後內心懷著的一個什麼動機.
一個什麼價值觀.
就好像我們上正面管教.
家長教小朋友.
你不要只看冰山上面.
行為或者表現.
其實可能裡面.
勾著很多價值觀.
或者一些不安.
或者一些感受.
我們一定要挖下去.
才知道可以幫到他.
耶穌一樣.
他看到的不單是人的外面.

$^{321}$更加是你躲在.
有些太太經常說.
她丈夫躲在廁所.
不知道半小時做什麼.
想著自己的空間.
就在那個位置.
耶穌也會看到.
代表著什麼呢.
在裡面他可以知道我們.
對一些事的感覺和動機.
所以在這裡.
門徒的這一課.
是很想擴闊門徒對天國的視野.
以至在裡面對跟從主耶穌.
有一份真正的認識.
甚至要打破他們一些盲點.
真正認識耶穌.
他做這麼多看得見的行為.
醫治,趕鬼和教導.
背後他的原意是什麼呢.
如何認識他.
真正這麼愛我們每一個世上的人.
由猶太人開始.
一步一步傾流出去.
在外邦人的身上.
同樣有這一份的愛.
所以回到今天的經文.
大家如果打開聖經.
在七章裡面.
二十四節開始.
耶穌就起身.
去到推羅西頓的境內.
前期剛才說了.
就是他和一群法利賽人.
民事講完他們對神的認識.
法利賽人和民事就拿出他們的傳統.
拿出他們的律例.
我就是這樣認識上帝.
這樣做就是敬虔.
這樣來回應耶穌.

$^{361}$二節都看到了.
談不攏了.
所以耶穌就走了.
就出來.
去到一個外邦人的地方.
在猶太的歷史和信仰裡面.
我們知道推羅西頓這個地方.
不是什麼好東西.
如果你看回舊約裡面.
先知不斷罵這個地方.
如果你看回以西結書二十六章開始說的.
對這些地方的預言.
就是說我世世代代都要和推羅這個地方為敵.
因為那個地方他們很驕傲.
還有覺得自己很厲害.
所以他們欺壓猶太人等等的事情.
所以以西結先知就說.
你一定會面對這些審判.
不單止以西結.
如果你看完撒迦利亞書.
都是一樣的.
都是在罵推羅西頓那些人.
像自己厲害,豐富等等的事情.
就在那裡做很多事情.
所以舊約裡面你會看到.
都是這樣說這些人.
同樣的.
但又理解為什麼推羅西頓.
是有這種身份.
很崇高的感覺.
因為其實如果你看回以前的歷史.
推羅本身是一個島.
做很多貿易的事情.
去做染料.
一些高尚的事情.
所以當時的政府就給很多資源去發展.
剛才說了.
推羅是一個島.
它厲害在什麼呢.
如果你現在找回地圖.

$^{401}$它已經不是一個島.
是一個半島.
就是當時要發展這個地方.
所以建了一條大道.
在沿岸和推羅島.
成為了一條大通道.
所以將推羅本身是一個島.
變成了一個半島.
你想想以當時的時代.
當時要用多少錢.
用多少的物力.
去將一個島變成一個半島.
所以當時的政府.
或者當時的國王.
都很想重點投資在這個地方.
繼續發展貿易等等的事情.
所以他自然是.
當時的一些焦點.
就像今天回到香港.
我們是東方之珠.
可能他是中東之珠.
很多的焦點.
都是放在這個地方.
無論發展經濟等等的事情.
都是很蓬勃的.
所以在裡面.
他有這個身份覺得.
你們加利利的猶太人.
窮鄉僻壤.
怎麼會給得上我們這個地方.
所以是很正常.
當然正正因為這個事情.
在馬太福音裡面也有說.
他拿了一個說法.
如果耶穌的神蹟在推羅西頓行.
他們就已經不用.
像所多瑪和阿摩拉的結果.
意思就是.
他將推羅西頓和所多瑪和阿摩拉.
放在一起做一個比喻.

$^{441}$所以在這個背景下.
你就會理解猶太人.
是多麼不喜歡這一群人.
好像今天在文化上.
或者在背景上.
我們也可能有一些人.
你見到也覺得.
好像不太喜歡.
當時猶太人一樣.
見到這群人就覺得.
很多負面的感受.
所以當一個猶太人的拉比.
耶穌就是一個老師.
進入一個這樣的地方的時候.
已經是很奇怪.
為什麼耶穌跟那些法利賽人.
聊完天之後.
要去到一個這樣的地方.
你不是很不喜歡我們這群人嗎.
來到裡面想做些什麼呢.
而且來到這個地方.
進入了一個家裡.
本身耶穌在聖經裡面所說.
是想隱藏的.
但是你知道耶穌已經很出名.
那時候隱藏不了.
甚或乎是一個這樣的身份.
跟我們這麼格格不入的拉比.
來到的時候.
他想做些什麼呢.
很多人都立刻想知道.
甚或乎當時的一個婦人.
聽到這個消息.
聽到知道這個人的背景的時候.
她也很想在這個人身上.
得到一份好處.
去經歷他.
不知道大家讀這段經文的時候.
會不會把前面那些.
來幫助的人等同了.

$^{481}$想著那個婦人很慘.
女兒又被鬼附.
一定很辛苦.
很落魄地走到耶穌身邊.
衣衫褸褸地求你幫我.
但是其實如果你看回那個背景.
馬可很精準地形容這個婦人是什麼呢.
她是一個希臘人.
是屬於敘利亞肥利基族.
如果你有機會認識當時的背景.
肥利基不要以為只是一個普通的小族.
你真的上網去搜尋羅馬的肥利基戰船.
你就會知道這班人是真的很厲害.
他們是出產當時的.
有一種船真的叫肥利基船.
就是當時的科技等等的事.
讓他們可以打造到這一艘船.
在羅馬縱容.
所以這班人在重工業和經濟方面.
都是很蓬勃和有經驗的.
所以他們的背景一定不是很窮.
或者是什麼的狀況.
希臘人更加是在教育上.
希臘人都知道當時很多哲學.
很多思想很前瞻和前端的一班人.
所以當你看這些背景的時候.
就知道這個婦人不簡單.
但在這個不簡單的背後.
有一份很重要的張力.
就是要告訴我們這個婦人.
完全和猶太人沒有關係.
零關係她是一個愛邦人.
又是一個很富庶.
不是生活在猶太人的景象裡面.
或者等等的一個人.
可以說是和神的子民.
神兒女身份不單止沒有關係.
更加有著很微妙的張力.
就是不喜歡她的那種狀況.
所以當時當馬可形容.

$^{521}$這個婦人來到的時候.
其實你就會更加明白當中.
究竟耶穌為什麼要來找她.
而這個婦人來到找她的時候.
也代表著什麼呢.
還有一點要加強大家知道的.
她真的是一個有錢人.
是在說她女兒的床.
你知道為什麼那張床那麼重要嗎.
當時不是很多人睡床.
是有一張席子在地上.
但你看到醫治好了之後.
就說女兒在床上起床.
所以更加看到這個婦人的身份.
是不簡單又顯赫.
除了這些不同之外.
還有一樣東西最重要.
就是信仰的傳統和猶太人很大分別.
如果你記得前一段時間.
耶穌做很多醫治的時候.
在猶太地的時候.
人們把病人和鬼婦帶到耶穌面前.
然後耶穌就連問他們就醫治了.
你看回經文.
那個婦人有沒有帶她的女兒來.
看起來應該沒有.
應該是她來到.
然後耶穌就說你回去吧.
這個敘述就像一個羅馬的白夫長.
他的概念是如果你是一個強大的神.
你只要說一句就行.
不用帶她來那麼辛苦.
所以和當時猶太人.
或者普遍的百姓不一樣.
信仰觀眾想帶她來.
連問她.
你那麼強大.
一句就行.
所以就像剛才所說.
從種種跡象的技術.

$^{561}$告訴我們這個希臘婦人.
背景,想法,正在做的事.
完全和那群猶太人風馬牛不相及.
也是兩個世界並存的兩群人.
所以在這些不同之下.
我們就會明白.
究竟耶穌正在做的事.
背後想表達什麼呢.
當然當這個婦人來到耶穌面前的時候.
她要求或者想求耶穌.
去醫好她的女兒.
但耶穌當時也很不客氣地回應她.
做一個比喻.
如果你有餅的時候.
就先給子女吃飽.
不要拿你的兒女的餅.
去餵給狗吃.
有些色經或者大家看的時候.
會覺得為什麼這麼絕情.
這樣說一個比喻.
人家是狗.
意思就是說.
如果你有恩典或者有供應.
先給子女.
不會給狗的.
所以當時有些人當然會幫耶穌說好話.
狗可能是在說寵物狗.
或者很疼愛那些.
但其實也不用這樣.
其實你反而更加覺得.
耶穌想用一個隱晦的方式.
讓這個婦人去明白.
其實已經是滿有關懷.
如果你想想.
其實如果要表達的意思就是.
你走吧.
我不讓外邦人有救恩.
不讓外邦人有醫治.
你走就好.
直接說就好.

$^{601}$就更好.
不用用比喻來說.
但你會覺得.
為什麼耶穌用一種隱晦的方式來回答.
就是很想讓她知道.
在這個階段.
外邦人還不能夠配受.
和在時間上還不準備去領受這個救恩.
很想藉著一個隱晦的方式.
讓她明白這件事.
對於耶穌這個比喻.
當然始終用狗.
我們今天都會覺得很冒犯.
對於耶穌這麼刺激的.
帶著隱晦的回應.
你會看到婦人不會很生氣.
你想想一個有錢人.
一個有方法的人.
雖然他搞不定.
女兒被鬼附身.
但對他來說.
我來求你幫我.
你不用說這些.
不幫就不幫.
有地位找其他人幫.
但你看到婦人沒有這樣.
她不會覺得耶穌那句話.
令她腦收成.
反面的情況.
相反.
藉著這個比喻.
這個婦人還看到一線曙光.
她還看到還有機會.
在這線曙光裡.
她承認一件事.
雖然我不配.
甚或乎她卑微到一個地步.
承認我是狗.
我像狗一樣不配受.
不是兒女.

$^{641}$但狗也有機會享受.
主人家裡的豐盛.
狗也能夠藉著.
和主人的關係.
得到這個供應.
所以她想通了耶穌這個比喻.
用耶穌這個比喻.
回答她這個狀況.
以致她知道.
耶穌你說得對.
我是一個外邦人.
我是一個外族人.
甚或乎像狗一樣不配.
但無論身份.
傳統上和猶太人比不上.
但這一刻她知道.
在這個滿有恩典的上帝面前.
她仍然可以乞求這一份憐憫.
因為這一刻她知道.
自己是不配.
對她來說.
她不是做了什麼.
剛才說的doing.
她所為做了什麼.
可以配受得到耶穌的憐憫.
沒有 完全沒有.
但耶穌看到的是.
她這一份對自己的認信.
知道自己是不配.
知道在她面前的耶穌.
是有這一切的恩惠.
所以她希望可以在耶穌身上.
得到這一份恩典.
事實上這一份恩典.
不只是這一刻才說出來.
其實為什麼說還不是時候.
外邦人不是不配受.
只是耶穌還沒一步一步展現出來.
讓外邦人可以領受這件事.
其實這份恩惠和傾流.

$^{681}$早在經文已經在這裡.
在哪裡呢.
當你留意餅這個主題的時候.
其實從五餅二魚開始.
就已經在做這件事.
一直延伸到這裡去展現出來.
大家想想.
五餅二魚這件事.
其實有個彩蛋.
大家知道什麼是彩蛋嗎.
看完電影就有一段短片.
勾起你的興趣.
讓你延續一下.
知道那個續集.
或者補充一些資料.
讓你覺得很有趣味.
或者等等.
其實五餅二魚有一個彩蛋.
大家知道嗎.
五餅二魚如果你看為.
上帝很有恩典.
供應給所有人.
哇 上帝很厲害.
五個餅 兩條魚都變了這麼多.
哇 神是供應.
停在這裡就夠了.
不用留下彩蛋是什麼呢.
就是坐下來吃飽之後.
耶穌叫十月門徒出去.
收回碎餅和碎魚.
十二男子回來.
當然你聽這個的時候.
你可能就想.
是不是代表他們吃得很飽.
通常我們傳統的.
叫完一桌東西.
長者 不是長者.
是長輩.
長輩看到吃完的東西.
再叫多一個意麵.

$^{721}$有剩才代表你們飽.
沒有剩的話.
你們不夠.
你們客氣.
通常叫到一桌東西.
有剩才覺得豐盛.
可能一方面代表他們真的很飽.
吃完五個餅 兩條魚.
之後還有剩.
就分出來.
除了這個可能之外.
還有第二個可能.
就像今天所說.
這些神的恩惠豐盛.
當他們享受.
真的吃到很飽之後.
這些碎餅碎魚剩下來的.
還有十二男子.
你想想耶穌叫他們出去收.
不會說你們出去收回那十二男子.
然後扔掉.
不用扔了.
挖個洞埋了就算了.
他收得應該有下面的供應.
或者拿給一些不在場的五千人.
或者那些婦孺等等的人.
這些不在場的人.
都能夠領受五餅二魚的供應.
就是這十二男子.
同樣這個愛邦婦人.
明顯就是不在場的人.
不在場的五千個男人.
再加上婦孺.
不是猶太人.
不是根律法.
不是從摩西亞伯拉罕.
得到什麼身份.
去到神面前的人.
沒有什麼做了.
可以覺得可以得到神恩典的人.

$^{761}$在這一刻.
還不完全可以.
同樣得到耶穌的恩典.
以至得到這份幫助.
當然這個婦人.
至少也來到耶穌面前.
她的女兒更加不配.
來都沒有來.
還在家裡.
同樣這份恩典.
也臨到這個女兒身上.
所以我們在這裡看到的.
就是神的恩.
神的救贖.
絕對不是基於我們的做.
我們做了什麼.
而是神見到我們.
認信見到我們的本質.
同樣也是因為這個婦人.
認信這個上帝.
這個上帝如果真的有一位神.
是這麼有大能的時候.
是一定有慈愛和憐憫的.
她不用我宰一隻雞去還神.
才給我一些東西.
我不用燒一隻豬去做一些事.
她才給我一些東西.
所以在這個有恩典.
有慈愛的上帝面前.
同樣她對自己的認識.
本質就是我很不配.
甚或乎她真的覺得.
即使我是一隻狗.
我也在神面前有上帝的愛.
當我有困難的時候.
上帝一定不會忽視她.
而事實上也是因為這個觀念.
她見到上帝終極的計劃.
不是像法利賽人或猶太人.
或門徒那種眼光.

$^{801}$去見到這件事.
所以這個人的所視.
對上帝的人的所視.
就是耶穌很想門徒見到的事.
但再進一步.
馬可慕庭在這裡.
他再用另一個神跡醫治.
來深化對耶穌的認識.
你會看到耶穌真的很忙.
剛剛這邊想見一個婦人.
幫助到她.
就立刻離開那個地方.
去到一個加利利海的另一邊.
大家知道加利利海很大.
耶穌和猶太人就在北部.
這是一個南部.
外邦人聚集的地方.
藉著這個再一次的醫治工作.
解釋了祂更逐步去揭示.
彌賽亞的救贖工作.
以致當時的人在明和不明之間.
神仍然用祂的慈愛和忍耐.
去揭示祂在裡面.
究竟在做什麼.
兩個技術.
為什麼今天大家要放在一起去看.
其實一定很有關連.
如果看二十四節.
耶穌起身到推羅西頓.
進入了一家人裡.
不想別人知道.
但是人又知道了.
隱藏不住.
去到耶穌第二個行完的醫治之後.
祂就告訴那些人.
你不要告訴別人.
但你越叫他不要說.
越發張揚了.
那種張力很相似.
好像就在這種張力之下.

$^{841}$原來耶穌就在這種困難裡.
或者人未明的時候.
正在開展認識人和神的下一個階段.
就是如何傾流出去.
當耶穌來到這個地方.
叫做底加玻里境內的加里尼海.
要讓大家知道背景.
完全是一個外邦人的地方.
因為底加玻里是一個希臘文.
就是十個城聯合在一起的意思.
所以是一個聯合的城邦.
裡面有自己希臘化的文化自治.
等等的情況.
所以這個被醫治的人.
都可能是一個外邦人.
有人帶他到耶穌面前.
大家上星期有沒有聽梁院長說.
這個醫治說一聲就行.
或者好像剛才那個婦人.
又是說你回去就好了.
那就行.
為什麼又要做這麼多事.
又要用手指插進耳朵裡.
如果用今天的角度.
用口水再弄進去.
哇 耶穌你衛生嗎.
可能你都覺得不用做這些事.
你說一聲就醫好我就行.
為什麼耶穌還要做這些事呢?.
以至在這個做事背後一定有他的意思.
不是為了做一場秀.
讓大家知道我有多厲害.
這個刻意要帶他離開眾人.
做的這個動作.
這個觸碰遠超過只是醫治他膿.
和洩潔這個狀況.
這個觸碰不只是醫好他外面.
是醫好他裡面.
當這個觸碰.
你想想一個猶太拉比.

$^{881}$來到一個完全外邦的地方.
兩個世界完全是兩條平行線.
河水不犯井水.
現在在這一點裡相會.
如果耶穌真的說.
行了 你回去吧.
你安安回去就好了.
那種震撼不能夠出來.
但耶穌直接這樣摸他.
接納他 醫治他.
這個醫治不是醫治外面.
不是只是一個神蹟.
告訴他好了 你好了 沒事了.
是要醫治他裡面的那個Being.
他自己的價值觀.
他自己對上帝的那種隔絕.
重新再與上帝結連.
所以這個醫治不只是停留在外面.
更加在裡面.
是要讓他徹徹底底地知道.
這個神就是與他這樣結連關係.
所以後面他繼續說.
他又已開了 舌頭又解了.
還要說得清清楚楚.
代表著他是一種徹徹底底的根深.
和釋放出來.
所以在這種強烈的對比之下.
你看看猶太人怎麼說他的認識神.
怎麼說他的認信.
在七章一至二十三節.
就是用他們的Doing 律法.
正在做的事.
總之你不要想我在想什麼.
孝不孝順父母不要緊.
我做了這些事就行了.
就敬虔了.
相比之下這班外邦人.
耶穌看到他們對神的渴慕.
縱然沒有做過什麼.
去認識上帝的裡面去做.

$^{921}$但是那份渴慕那位真的神.
能夠幫助他的神.
就在耶穌的身上徹底改變過來.
本身不配的.
但是神仍然能夠在當中扭轉這些生命.
所以對當時跟著耶穌的門徒來說.
很不容易去接受.
剛才說了.
這班人他們又不喜歡.
現在神又幫助他們.
又傾流出去.
他們做了什麼.
我們做了這麼多.
他們又沒有做.
他們世世代代那種傳統等等的事.
一直深刻在他們的心裡面.
現在這一刻.
嘩!.
不同了.
所以為什麼你不難明白.
耶穌不會將這件事在猶太地做.
是要帶著這班門徒.
去到一個完全外邦的地方.
在這裡成為一種教室.
不會在會堂或者這些地方.
帶到外邦人的地方.
讓他們親身感受.
神看見的是人的內在.
不只是看見外面這一層.
可能你又會覺得很有趣.
當做完這件事之後.
為什麼剛才說36節那裡說.
耶穌做完這些事.
又叫他不要告訴別人.
誰知道他們當然不聽話.
越發地去傳揚.
可能你也會想.
耶穌你不是做這些事.
就是想讓更多人知道嗎?.
多些人知道.

$^{961}$其實可能你想像那些很客套的那些.
不用多謝我.
但其實心底裡.
你多謝我吧.
你不要說了.
然後你說出去吧.
我相信耶穌應該.
這些小小的虛偽都沒有.
就是那種狀況.
他真的不想他們說出去.
因為那個時候還沒到.
他關心的不是自己在做的事.
夠不夠多人知道.
不是關心的.
是這些事是怎樣震撼到當時的地方.
他關心的是.
當時這班人.
是否準備好去接受這份救恩.
去接受這個上帝是這樣的形式.
這樣的樣子.
不是跟他們所信開.
或者是那些形式那種的上帝.
要突破這個盲點.
去接受這個救主.
但當你看回七章第37節.
你就明白為什麼這樣.
當時的人會怎樣呢?.
眾人憤內唏氣地說.
哇!他所做的事好啊.
他連聾的都叫他們聽見.
啞的都叫他們說話.
只是一聲嘆息.
他們仍然看到耶穌的做事.
仍然看到耶穌的冰山上.
他看不到耶穌那份真正接納.
真正醫治.
真正幫助他們生命的那個位置.
所以耶穌寧願你們還不知道的.
先不要說了.
還不明白的.

$^{1001}$讓我多做一點.
明白多一點的時候.
其實因為這份張力一直延伸下去.
到第八章都保持在這裡.
所以還不是這個時候.
耶穌說不要再說了.
你們明白得醫治.
得更新就OK了.
往後去揭示彌賽亞的救贖會更加多.
所以到這裡.
對我們今天來說.
這些新約後.
知道了耶穌的工作.
明白救恩的人.
是什麼呢?.
求神給我們這份活潑的信.
這份活潑的信相對著一件事.
不是那種固態的信.
固態的信是什麼呢?.
就是一開始是很好的.
我知道這樣我是在表達我的信心.
我知道這樣是敬虔和認識上帝.
敬拜上帝.
這些做法.
無論是事奉還是其他.
我就是讓人可以在我身上見到上帝.
藉著這些做事.
我就知道自己是在敬畏神.
但好像剛才一直說.
你有沒有在審查這些做事之間.
去想想.
如果有一天.
有一個人說信耶穌.
但不是在跟隨你的做事.
不是在跟隨你所做的事情.
你會怎樣呢?.
他不屬靈.
為什麼要做這些事情?.
他不是跟隨我的敬虔方式.
他這樣做.

$^{1041}$耶穌 罰他.
如果只是這樣.
那你就是固態.
只有這種形式.
只有這種形式的信仰.
不然就不屬靈.
我這種就可以表達了.
所以當我們要不斷反思.
不單是信仰的行動.
信仰的內涵.
我們都要不斷反思.
就算我們傳福音.
可能只是要別人跟隨我.
就經歷上帝的時候.
都限制了上帝所做的事情.
以至於將福音帶給別人的時候.
我們沒有理會到人的所視.
人的內在價值.
人每一個事的動機.
當你有這種對所做的事情.
所懷著的身份.
價值和所視的時候.
這種內外.
即使你知道.
有些地方.
在裡面你自己在想什麼.
但你就會像那個婦人一樣.
你就慢慢感受到.
原來我真的這麼不配.
我這麼有限.
我這麼不配上帝的恩典.
在這種張力裡面.
你就會明白.
不只是你做了什麼.
去得著神的經歷.
而是你裡面是什麼.
但同樣的.
也因著你裡面的根深.
你所做的事情.
都帶來改變.

$^{1081}$以前張藝信牧師教我一件事.
他講了很多次.
有諸於內而仍於外.
但也不只是停留在.
我裡面有一些東西.
我再留意.
在做的時候.
你也要行在悟中.
以至悟而後行.
即是說那個doing和being.
是一路交替.
你不是想著.
我有這件事.
我做了這件事.
固態了它.
我就不用來這裡.
而是一路做.
又想我在做什麼.
當你在信仰裡面.
做的事情.
遇到一些碰板.
或者一些限制的時候.
不要這麼快就擔心.
會不會是.
我就不要了.
又或者不是自己不要.
是他不要了.
為什麼他不這樣跟著我做.
試試停一停.
想想裡面.
究竟耶穌看到我們.
裡面冰山下面的動機.
是什麼呢.
咦 你這樣做嗎.
你在想的是什麼.
我這樣做.
我想告訴你.
我是在表達什麼.
以致我們至少知道.
耶穌應該不是一個方丈.

$^{1121}$他不小器.
他不會說.
哇 只有這樣做.
你這樣做就不行了.
只要他看到我們裡面的需要.
我們裡面的所事的時候.
他一定會讓他的恩典.
臨到我們.
以致我們能夠不斷去拆身認識.
對耶穌的觀點.
最後再說一個例子.
可能大家就會明白多一點.
早前有機會去澳門.
去百革巢浸信會.
不知道大家是否認識黎斯姆.
他有一個很傳奇的.
有機會大家可以去探訪他.
黎斯姆以前在巴西宣教.
他的工作是在巴西照顧華人.
因為那時候很多人去開餐廳.
做很多苦工等等.
華人的小朋友就隨街走.
做街童.
所以當時的領袖就說.
要照顧這班孩童.
其中一個方式是做什麼.
就是開放教會.
開放教會做一些社區活動.
上來有些活動等等.
其實現在百革巢浸信會.
有機會有很多年輕人上來.
他說那時候會做什麼.
在教會一個活動中心.
就會有些聽歌.
或者你都覺得可以.
玩一下飛鏢等等.
他說那時候在教會裡.
開了一個吸煙室.
可能你馬上就說.
教會裡開吸煙室.

$^{1161}$怎麼可以這樣.
他說那時候的情況是怎樣.
是因為他知道那班人的需要.
等等的時候.
但你會看到.
在這個外在的行為.
原來慢慢會改變.
他說那些年輕人上來吸煙.
過了一段時間.
在教會門口吸.
再過一段時間.
就在教會樓下吸.
再過一段時間.
就連吃都不吃了.
這個就是由裡面.
產生到外面.
但如果你外面就已經停止了.
就產生不到在裡面.
這種的生命力.
是否在我們每一個服侍的人.
信主的人的生命裡.
都一樣求神擴闊我們的空間.
以至在我們的B型裡.
我們的鎖匙.
再尋索上帝所喜悅的生命.
我們一起祈禱吧.
多謝主.
讓我們能夠藉著經文.
再一次去看見上帝的恩典.
這樣傾流在我們.
這些不配的人的身上.
昔日的婦人.
以及耳聾舌竭的人.
經歷到上帝的恩典.
因為他們看見上帝的工作.
背後的心思.
所以你的恩典.
同樣在我們每一個生命裡觸碰.
幫助的不是我們的外表這麼簡單.
更加是我們心底裡產生變化.

$^{1201}$求主引領我們.
以至無論在人前.
在人後.
我們的生命都與你結連.
多謝你.
祈禱在奉主耶穌基督的名字而求.
阿們.
\newpage



\section{馬可福音 7:1-23}
\label{sec:0QX__zW2yAY}
\textbf{2025年8月24日崇拜直播｜陳冠成牧師｜耶穌給門徒的十課－傳道vs誡命｜馬可福音7\_1-23}
\newline
\newline
連結: \href{https://youtube.com/watch?v=0QX__zW2yAY}{\texttt{https://youtube.com/watch?v=0QX\_\_zW2yAY}} ~~~~ 語音日期: 2025-08-30
\newline
\newline
\hyperref[sec:iXL9M6i3TXE]{\small{< < < PREV SERMON < < <}}
~
\hyperref[sec:index_chronic]{\small{[返順時目]}}
~
\hyperref[sec:EMIvj5LfGAU]{\small{> > > NEXT SERMON > > >}}
\newline
\newline
馬可福音 7:1-23
\newline
\begin{longtable}{cl}
\hline
\hline
章節 & 經文 (和合本修訂版)\\
\hline
7:1 & \begin{tabularx}{0.7\textwidth}{X} 有法利賽人和幾個從耶路撒冷來的文士聚集到耶穌那裡。 \end{tabularx} \\ \\ \relax
7:2 & \begin{tabularx}{0.7\textwidth}{X} 他們曾看見他的門徒中有人用不潔淨的手，就是沒有洗的手吃飯。 \end{tabularx} \\ \\ \relax
7:3 & \begin{tabularx}{0.7\textwidth}{X} 法利賽人和所有的猶太人都拘守古人的傳統，若不按規矩洗手就不吃飯； \end{tabularx} \\ \\ \relax
7:4 & \begin{tabularx}{0.7\textwidth}{X} 從市場來，若不洗淨也不吃飯；他們還拘守好些別的規矩，如洗杯、罐、銅器、床鋪等。 \end{tabularx} \\ \\ \relax
7:5 & \begin{tabularx}{0.7\textwidth}{X} 法利賽人和文士問他說：「你的門徒為甚麼不照古人的傳統，竟然用不潔淨的手吃飯呢？」 \end{tabularx} \\ \\ \relax
7:6 & \begin{tabularx}{0.7\textwidth}{X} 耶穌對他們說：「以賽亞指著你們假冒為善的人所預言的說得好。如經上所記：『這百姓用嘴唇尊敬我，他們的心卻遠離我。 \end{tabularx} \\ \\ \relax
7:7 & \begin{tabularx}{0.7\textwidth}{X} 他們把人的規條當作教義教導人；他們拜我也是枉然。』 \end{tabularx} \\ \\ \relax
7:8 & \begin{tabularx}{0.7\textwidth}{X} 你們是離棄神的誡命，拘守人的傳統。」 \end{tabularx} \\ \\ \relax
7:9 & \begin{tabularx}{0.7\textwidth}{X} 耶穌又說：「你們誠然是廢棄神的誡命，為要守自己的傳統。 \end{tabularx} \\ \\ \relax
7:10 & \begin{tabularx}{0.7\textwidth}{X} 摩西說：『當孝敬父母』；又說：『咒罵父母的，必須處死。』 \end{tabularx} \\ \\ \relax
7:11 & \begin{tabularx}{0.7\textwidth}{X} 你們倒說：『人若對父母說：我所當供奉你的已經作了各耳板』（各耳板就是奉獻的意思）， \end{tabularx} \\ \\ \relax
7:12 & \begin{tabularx}{0.7\textwidth}{X} 你們就容許他不必再奉養父母。 \end{tabularx} \\ \\ \relax
7:13 & \begin{tabularx}{0.7\textwidth}{X} 這就是你們藉著繼承傳統，廢了神的話。你們還做許多這樣的事。」 \end{tabularx} \\ \\ \relax
7:14 & \begin{tabularx}{0.7\textwidth}{X} 耶穌又叫眾人來，對他們說：「你們都要聽我的話，也要明白。 \end{tabularx} \\ \\ \relax
7:15 & \begin{tabularx}{0.7\textwidth}{X} 從外面進去的不能玷污人，惟有從裡面出來的才玷污人。」 \end{tabularx} \\ \\ \relax
7:16 & \begin{tabularx}{0.7\textwidth}{X} 有耳可聽的，就應當聽！ \end{tabularx} \\ \\ \relax
7:17 & \begin{tabularx}{0.7\textwidth}{X} 耶穌離開眾人，進了屋子，門徒就問他這比喻的意思。 \end{tabularx} \\ \\ \relax
7:18 & \begin{tabularx}{0.7\textwidth}{X} 耶穌對他們說：「你們也是這樣不明白嗎？難道你們不了解，凡從外面進去的不能玷污人嗎？ \end{tabularx} \\ \\ \relax
7:19 & \begin{tabularx}{0.7\textwidth}{X} 因為不是進入他的心，而是進入他的肚子，又排入廁所。」（這是說，各樣的食物都是潔淨的。） \end{tabularx} \\ \\ \relax
7:20 & \begin{tabularx}{0.7\textwidth}{X} 耶穌又說：「從人裡面出來的，那才玷污人； \end{tabularx} \\ \\ \relax
7:21 & \begin{tabularx}{0.7\textwidth}{X} 因為從人心裡發出種種惡念，如淫亂、偷盜、兇殺、 \end{tabularx} \\ \\ \relax
7:22 & \begin{tabularx}{0.7\textwidth}{X} 姦淫、貪婪、邪惡、詭詐、淫蕩、嫉妒、毀謗、驕傲、狂妄。 \end{tabularx} \\ \\ \relax
7:23 & \begin{tabularx}{0.7\textwidth}{X} 這一切的惡都是從裡面出來，且能玷污人。」 \end{tabularx} \\ \\
[1ex]
\hline
\hline
\end{longtable}
$^{1}$今日經訓的內容是在新約聖經馬可福音第七章一至二十三節,亦是新約聖經第五十六頁..
馬可福音第七章第一節,有法利賽人和幾個民事從耶路撒冷來,到耶穌那裡聚集,他們曾看見他的門徒中有人用竹手,就是沒有洗的手吃飯..
原來法利賽人和猶太人都俱守古人的遺傳,若不仔細洗手,就不吃飯.從史上來,若不洗浴也不吃飯,還有好些別的規矩,他們歷代俱守,就是洗杯,灌,銅器等物..
法利賽人和民事問他:你的門徒為什麼不照古人的遺傳用竹手吃飯?耶穌說:以賽亞子著你們假冒為善之人所說的預言是不錯的..
如經上說:這百姓用嘴唇尊敬我,心卻遠離我.他們將人的吩咐當作道理教導人,所以拜我也是枉然.你們是離棄神的誡命,俱守人的遺傳..
猶說:你們誠然是廢棄神的誡命,要守自己的遺傳.摩西說:當孝敬父母.猶說:咒罵父母的必致死他..
你們賭說:人若對父母說:我說,當奉給你的已經作了國耳板.國耳板就是貢獻的儀式,以後你們就不容他再奉養父母..
這就是你們承接遺傳,廢了神的道.你們幻作許多這樣的事.耶穌又叫眾人來,對他們說:你們都要聽我的話,也要明白,從外面進去的不能污衊人,.
唯有從裡面出來的乃能污衊人.耶穌離開眾人,進了屋子,滿腹就問他這比喻的意思..
耶穌對他們說:你們也是這樣不明白嗎?豈不曉得凡從外面進入的不能污衊人,因為不是入他的心,乃是入他的屠福,又落到茅廁裡..
這時說:各樣的食物都是潔淨的.又說:從人裡面出來的乃才能污衊人,因為從裡面就是從人心裡發出惡念..
九行:偷盜,兇殺,奸淫,貪婪,邪惡,詭詐,淫蕩,嫉妒,蚌毒,驕傲,狂妄.這一切的惡都是從裡面出來,且能污衊人..
各位弟兄姊妹,平安..
如果你留意剛才各位弟兄所讀的經文,你會發覺「傳統」這個字,即是「和本亦作為傳」,其實出了五節,特別集中在上半部分,一至十三節..
法利賽人和文士指控耶穌的門徒沒有遵守猶太人所重視的洗手的傳統規矩,耶穌卻反過來斥責他們,將人定立的傳統和規矩,當成上帝誡命一樣遵守..
然後他就教訓眾人,到底真正污衊人的其實是什麼事情.事實上,這段經文不單單適用於當日的百姓,其實亦可以引導我們重新思想,教會也有不同的傳統和規矩..
我們值得透過這段經文檢視一下當下我們如何去看待這些傳統和信仰.我們有沒有過於著重,如經文所說的外在的宗教行為,卻忘記了上帝更加體重的,其實是人的內心意念..
一開始我們看到馬可記載耶穌的對手是來自耶路撒冷的精英分子,大概因為耶穌在當時已經聲名遠播,引起當時的宗教中心耶路撒冷的領袖關注,於是他們就派出他們的,我們說在信仰上的精英,.
花里塞人和文士來找耶穌,希望在耶穌當中找到一些錯處,藉此貶低耶穌的權威,甚至找機會來控告他.他們以門徒用觸手,即是說沒有洗過的手來吃飯,作為指控耶穌的理據..
馬可胸房我們不太明白當時的文化處境,你看到他刻意在第三和第四節做了一些註腳,原來花里塞人和當時的猶太人都會拘守古人留下來的一個規矩,簡單說就是吃飯前要洗手,否則就不能吃飯..
另外,他們從街上回到家中的途中,可能在過程中會沾染到一些所謂不潔之物,包括外邦人,於是他們有一個規矩,回到家中第一件事就是要自潔,要將身上的污穢洗除,然後才可以用餐..
這就是古人的傳統,他們如何在耶穌時代慢慢流傳,甚至乎盛行下來呢?我們嘗試看看舊約怎麼說,如果你有留意,特別是有參與讀經計劃,你應該已經讀完五經了,這個星期開始,恭喜大家,已經讀完五經了..
你會記得,崔一及記,上帝確實吩咐了祭司和大祭司在進入會幕或聖殿公職前,他們真的要洗手洗腳來潔淨自己,以合乎我們說的禮儀上的潔淨條例..
因為祭司和大祭司是代表以色列傳會眾來向那位聖潔的上帝獻上潔淨之物,所以他們必須在進行公職前首先自潔自己,否則無論是他們所獻上的,甚或是祭壇或是會幕都會被他們玷污,這樣就等同於褻瀆了上帝的聖潔..
所以摩西的律法明文規定寫死了,祭司和大祭司在公職前必須洗手洗腳潔淨自己,這部分一點爭議都沒有..
至於普羅百姓,其實律法沒有規定他們真的要在吃飯前洗手洗腳,但是隨著猶太人的歷史,你會知道他們亡國被擄,經歷了,之後幾百年來都不被不同的外邦政權來統治..
他們對潔淨條例的理解其實都相應出現了很多的演變,原本只適用於祭司和大祭司公職所用的洗手潔淨條例,慢慢開始擴展應用範圍,去到平民百姓身上..
要求他們在日常生活中都要像聖殿公職的祭司一樣,來去驅守洗手自潔的規矩.事實上,洗手在耶穌時代就變成了一種身份的記號,或者說是一個敬虔的標記..
讓猶太人在當時充斥著異教信仰的世界中,洗手就很清楚地讓他們分別出來,和外邦人區分出來,並且像經文所說,透過各式各樣外在潔淨的傳統規矩,不單止洗手,洗杯,洗罐,洗銅器,更加凸顯我們的身份是獨特的..
我們猶太人的身份與別不同,我們的信仰習慣也與我們身邊的外邦人不同,因為我們是上帝的子民,所以我們有別於其他不是神所揀選的族群,以凸顯他們的身份地位..
至於經文所說,法利賽人就是屈守這些傳統規矩的表表者.事實上,作為猶太教一國的派別,法利賽人出了名的就是他們對律法的態度非常嚴謹和詳盡,他們會花很多時間鑽研如何執行摩西律法的細節..
並且經文也說,他們很著重祖宗的傳統,也就是說將歷史歷代猶太人給他們口傳下來的教訓,他們會視為珍寶,甚至當成摩西的律法一樣,有同等的地位,將古人遺留下來的傳統當成上帝的誡命來遵守..
他們覺得這是一個敬虔的表現.事實上,假如當時的百姓沒有遵守這些古人傳統,法利賽人和文士也會視之為一種不敬虔,不屬靈的表達..
所以當他們看到耶穌的門徒竟然用觸手來吃飯的時候,他們就質問耶穌,耶穌你和我們一樣都是作為拉比,為什麼你可以容許你的門徒不遵守這些古人傳統,容許他們用觸手來吃飯?.
潛台詞就是為什麼你容許你的門徒對上帝不敬,不屬靈?你看到耶穌沒有向他們解釋,而是直接斥責他們..
耶穌引用以賽亞書來駁斥法利賽人和文士對他們的指控,經文特別說一個字,叫做假冒為善..
善意其實是指演員或舞台表演者的意思,耶穌斥責文士和法利賽人驅守古人傳統的行徑,其實是演戲,假裝敬拜上帝,很嚴厲的指控,.
耶穌說其實他們的行為就正如當日以賽亞先知斥責百姓一樣,嘴巴是尊敬上帝,但是內心卻是遠離上帝,所以耶穌其實諷刺法利賽人這種流於表面的敬拜行為,其實是謊言,或者用廣東話說就是浪費力氣,.
其實上帝根本不會越立,原因是他們把人的吩咐抬高到一個地位,把這些古人傳統當成真理,不單單是自己遵守,還要教訓百姓來實行,為了驅守人的吩咐這個傳統,而棄置了上帝吩咐的誡命,基本上是一種本末倒置的行為..
為了講明他們的問題所在,你會看到耶穌引用一個例證,叫做「國義版」,在9至12節這裡是國義版,國義版的意思是什麼呢?其實是在說一個獻給上帝的禮物,顧名思義你獻給上帝了,可不可以拿回?.

$^{41}$當然不行,因為已獻給上帝了,不能夠反口說要拿回,所以法利初人就教導當時的百姓,如果你把任何的物件,包括財物,拿來做了國義版,奉獻給上帝的時候,你不能夠反口,不能夠拿回來做其他用途,包括供養父母..
如果按照和合部所講,法利初人不容許做了國義版的人,他有後悔的機會,可能做國義版的人因為一時衝動,看到旁邊的人奉獻給上帝,他又想奉獻給上帝,但是沒有想清楚,他還沒有奉獻給公養父母,他拿了公養父母的錢來奉獻給上帝,但是他想反口,他想後悔,.
但是法利初人就說不行,已出之物,不能夠回頭,換句話來說,他們就變成為了遵守國義版的承諾,而沒有履行在律法當中講要孝敬父母的吩咐,而如果你參考經文這裡,都印在程序表裡面,你參考《新漢語》譯本,你就會看到另一個意思,就是法利初人容許做了國義版的人可以逃避不供養父母,.
可以有藉口,我做了國義版,所以我這個月不能拿錢來養你,換句話來說,就是說變相容許百姓可以合法或者合理地不遵守十誡裡面的第五誡,當孝敬父母,所以你看到,其實耶穌奉賜文士和法利初人,你看另一個譯本就會看到,印在程序表裡面,奉賜他們太聰明,意思是什麼呢?.
為了守住自己的傳統,竟然巧妙地廢掉了上帝的誡命,你看到其實,如果你一直看馬可福音來說,這個是到目前為止,耶穌在馬可福音裡面對宗教領袖一個最嚴厲的批判,必須留心,耶穌其實不是否定傳統規矩的作用,.
而是強調不能夠過分高舉它,不可以將傳統規矩看成為上帝的誡命,甚至乎看得比上帝的誡命更加重要,摩西的律法優先,人的傳統只是一個輔助,宗教傳統只是為了協助人更有效地去理解上帝的吩咐和執行上帝的吩咐,.
所以不能夠像我們說的,不能夠小女孩大過主人婆,將宗教傳統等同於聖經的教導,甚至為了固守一些人定立的規矩而廢棄了更加重要的上帝誡命,這個就是耶穌駁斥法利賽人固守古人傳統的一個理據..
當他駁斥完之後,接著下來,他就轉過頭來除職,提醒眾人,你看到其實經文這裡寫得很生動,耶穌又叫眾人來對他們說,你們要聽我的話也要明白我,即是說什麼?不要再聽法利賽人和文士說,聽我說,要領悟耶穌基督的教導..
耶穌間接表明,他比文士和法利賽人更有權並更有資格,或者說其實他才有資格來定義什麼使人潔淨或不潔淨..
耶穌在這個段落裡,簡單說就是重複了一個教導,從外面進來的是不能夠污衊人,唯有從裡面出來的才會污衊人..
很明顯其實耶穌這段說話,或者說這個比喻是針對文士和法利賽人過分著重洗手這個外在的宗教行為..
耶穌表示其實無論你洗不洗手,遵守與否,其實不會污衊人,其實不會影響你和上帝的關係,因為真正污衊人,破壞人和上帝的關係,其實是人內心裡出來的行動..
當耶穌講完這個比喻之後,教導完群眾之後,他們都一如以往地回到家裡,對門徒進行一課進深的課程..
門徒似乎不太明白耶穌這個比喻到底說什麼,當然也顯示在現階段,他們的屬靈領悟力似乎仍然不像耶穌預期那樣..
但你也看到耶穌仍然很有耐性地向他們講解..
耶穌說,凡從外面進入身體的,包括食物,只會進入人的肚腹,最後會去廁所出來,即是拉出來..
所以,食物乾不乾淨,吃飯前有沒有洗手,這些外在的潔淨行為,其實最多用我們今天的說法,都只是涉及衛生問題,跟屬靈無關..
你洗不洗手,不代表你有沒有敬虔一點,上帝不看這些,不抗拒這些..
上帝抗拒的是什麼呢?真正污衊人的是另一件事,你洗不洗手,因為不會顛覆你的內心..
耶穌繼續說,你見到當他刻意說完這段說話之後,有沒有留意,我本來有個括號,十九節這裡,.
「各樣的食物都是潔淨的」,這句說話大概不是耶穌基督在當時說的,如果是的話,之後彼得就不用這麼煩惱了,.
彼得在士多林傳見到天上掉下來的不潔之物,又很煩惱,吃不吃好,或者之後的士多林傳耶路撒冷會議就不需要再爭論到底哪些外邦人能吃得,或者不能吃得..
很明顯這句說話是當時作者馬可自己附加的一個注腳,所以你看到無論是和學本,或者其他的譯本,甚至英文聖經在裡面會加上一個括號,.
來表示這是馬可對當時耶穌的一個講論進一步的詮釋,很有可能是馬可很想他回應當時他一個牧養處境的需要,.
你知道馬可當時牧養的教會也有很多外邦信徒在當中,事實上各樣的食物都是潔淨,不是說真的可以不按照舊約律法的規矩什麼都能吃,.
而是強調食物本身不能污衊人,既然如此,外邦人或者馬可牧養的外邦信徒到底有沒有遵守潔淨條例,有沒有洗手,其實不影響他們的經歷,福音,不影響他們進入上帝的國度..
耶穌會繼續說到底什麼才是真正污衊人呢?20到23節,他說真正影響人的其實是人內心發出來的惡念,惡行,那些涉及道德倫理的行為,.
才會真正污衊人,才會破壞人與人之間的關係,破壞與上帝的關係..
耶穌認定內心的潔淨遠比外在的潔淨更為重要,其實也間接說明了什麼?暗諷當時的法利賽人,其實你們再驅守多少傳統都是枉然,因為在上帝的眼中,他們都是污衊,因為他們現在心存惡念,敵當耶穌..
事實上,弟兄姊妹,今天這段經文不算很複雜,很簡單,很易明白,我們做一點信仰反省吧..
你會發覺,其實在歷世歷代教會的群體,他們都會致力將聖經的教導,來應用到生活每一個層面,這方面我們都很理解..
其實在過程中,一定會發展出各式各樣的傳統與規矩,我們都不例外,我們聚會,我們崇拜都有不同的傳統,有不同的規矩..
舉一個最簡單的例子,我們的講壇,浸信會放在中間,為什麼?.
(聖經為中心的真理).
立即有人回答這個傳統的來源,所以傳統,你知道的來源是很重要,不要純粹因為它是傳統,你就不理解背後的意思而去死守..
講壇中自象徵著浸信會很強調,很看重聖經的教導,當然有它一些的歷史背景..
所以你會看到浸信會相對會少用一些禮儀的方式來表達我們的信仰,大概你記得浸信會有兩大禮儀,一個是浸禮,一個是煮餐..
但你會看到,有時候你去另一些教會的禮堂,另一些宗派,他們放在中間是什麼?.
是所謂的聖壇祭壇,講壇是兩邊的,講壇是放在另一邊的,一邊是主席,一邊是講員..
你會看到他強調的是另一件事,將聖壇祭壇放在中間,強調他們的信仰表達是向上帝獻上祭物..

$^{81}$所以相對這一類的宗派會多一些所謂禮儀的方式來表述他們的信仰..
兩種傳統各有不同,但有沒有高低之分?沒有..
大家都是在演繹他們所看重的信仰元素,目的都是怎樣?.
要幫助人更加有效地表達和實踐真理..
事實上耶穌沒有否定傳統的作用,而是提醒人怎樣?.
你不要固守了人定立的傳統和規矩,因為始終它是出於人..
既然人有軟弱有聚性,所定立的東西總會有限制,總會有瑕疵,有不完整的部分..
如果你有發覺,其實早在馬可福音第二章,你記不記得耶穌說過一個伸走舊皮袋的比喻?.
其實耶穌已經提醒當時的法利賽人,他們禁食的傳統規矩,其實是不足以承載耶穌所傳遞的福音,其實那裡已經說了..
事實上你會發現在馬可福音裡面,耶穌不斷去打破或更新他們所守的不同傳統,禁食也好,安息也好,提醒他們需要用新的皮袋來裝新的酒..
套用伸走舊皮袋的比喻,如果落入今天的經文,其實耶穌就是提醒我們,不要一成不變地去拘守傳統規矩..
最重要的是這些傳統規矩是否符合真理的教導,能否承載耶穌基督的福音..
再舉多一個例子,剛才我們說有兩大禮儀,浸禮,浸水妹很重視浸禮這個傳統,強調浸禮應該是全身滑入水中,然後從水裡面上來,.
因為昔日聖經裡面記載耶穌基督受浸的時候也是這樣,但我們也不會死守這個傳統規矩,.
假如有些姐妹因為病患的緣故,因為行動不便的緣故,她真的不適宜全身滑入水中的時候,我們不會強行來的,.
一方面我們擔心這樣反而會令她受傷,本來身體已經差了一點,你浸她下水再上來,冷親肺炎就更麻煩,.
但更重要的是我們不想因為拘守一些傳統規矩而窒礙了人經歷上帝的憐憫和關愛,這個才是更加重要,.
所以你知道我們會彈性處理,我們會為有特別需要的弟兄姊妹做灑水禮,當然我們也期望日後他們身體狀況理想的時候,.
我們可以為他們補辦一個全身滑入水中的贈禮,來盡力去持守這個傳統,事實上我很喜歡一位聖經學者,他叫David Garland,.
他說教會的信仰傳統其實好像一個蟹殼,蟹要生存要成長,其實一定要定時脫殼,.
如果過程中蟹殼太硬,就沒辦法脫殼而出,就會死在蟹殼裡面,.
但如果蟹殼太軟又會怎樣?又會不夠保護,又很容易被外來的入侵者傷害,同樣教會也是這樣,.
教會需要建立讓人感到安穩,感到受保護的信仰傳統的規範,以致我們不被世俗的歪風入侵,.
保守教會可以健康的成長,但如果這個信仰傳統變得太硬,完全沒有彈性的話,又會不知不覺扼殺了教會成長發展的空間,.
事實上忠於自己對真理的看法,捍衛他所相信的傳統,是無可厚非的,.
而且為了抵觸世俗一些不好的觀念入侵而堅守信仰傳統,也是值得尊重的,.
不過這段聖經提醒我們要小心犯上了文士和法利賽人同樣的錯誤,就是把外在的信仰規矩變成了一幅難阻人認識親近上帝的高牆,.
事實上在耶穌的時代,嚴守各式各樣的規矩,無論是安息日,無論是國禮,飲食條例,甚至是洗手,.
無疑是對法利賽人或文士一種他們覺得崇高的信仰表達,.
但是也會令那些不能夠嚴格遵守的猶太人,那些無法遵守的猶太人,將他們排外了,.
讓他們無法體驗上帝的憐憫和愛,所以弟兄姊妹當我們今天看到耶穌的教導行動的時候,.
其實也是挑戰我們,甚至乎他要打破我們的信仰傳統框框,就好像他昔日要打破猶太社會的固有傳統規範,.
特別你記得上個星期的幫傳道,我跟他反過來了一段經文講,他先講了下半節,24至27節,.
你會看到其實就是一個明顯的例子,耶穌刻意讓那些被猶太人視為不潔的外邦人,.
可以經歷到福音,經歷到上帝的恩典,你會記得耶穌怎樣跟希臘婦人互動,.
耶穌說應該先讓女兒吃飽,不應該讓本來給女兒吃的餅拿來給狗吃,.
講的是救恩的次序,先給自己人猶太人,然後再給外邦人,耶穌也有他的規矩次序,.
但是耶穌稱讚他的是什麼呢?他認同這個過程,是的,不過,你在桌下的那隻狗也可以分到女兒吃剩的餅碎,.
意思是什麼呢?他尊重這種規矩,這個救恩次序,但是上帝是有恩典可以例外處理,.
上帝不是一成不變,所以為什麼耶穌二話不說稱讚他,給他幾個讚,你女兒馬上治好,.

$^{121}$同樣的,姐妹們,如果教會今天的規矩訂得太多太死的時候,其實毫無擔心的話,.
一定會窒礙方向的擴展,讓外邊的人基本上都不能靠近我們身邊,他們無法經歷上帝的恩典,.
舉個例子,假如你看到有個陌生人,突然間上來,他來借廁所,你知道旺鎮的廁所也挺有名的,.
真的,我不是說笑的,因為我們的書店和公司清潔得很好,比起其他樓層都要優勝,所以有時候會見到一個半個,.
你明知道他不是這裡,應該是三樓那些,甚至是職員,上來借廁所,你會不會因為教會規矩不行,.
只接待會友,只接待基督徒,其他的,對不起,下去,你會不會趕他走,.
傳統規矩,不會的,又或者你會不會當其難,可以用,聽五分鐘呼音,過來過來,聽完才可以用,.
你不會這樣為了死守規矩而排外,否則你會讓人覺得教會都沒有愛心,教會什麼都講規矩,什麼都那麼硬,.
並不是說規矩不重要,但是如果這些規矩真的難阻人經歷上帝的愛和恩典接納的時候,.
我們是否需要,好像耶穌提醒我們,你要彈性,酌情處理,事實上,今天經文提醒我們,教會不是一個僵化了的組織,.
而在聖靈引導的裡面,需要不斷作出更新,變化一個好像一個有機體,.
檢視我們一直持守的信仰傳統,不是為了要推翻過往或前人的努力擺上,.
事實上是想我們謙卑,信服,誠心地審視我們過往視為一直有效的規矩,習慣,.
到今天是否仍然配合上帝對教會的心意,到今天是否仍然有效配合這個世代的福音需要,這是經文要提醒我們,.
耶穌說教會不是為了要屈守傳統規矩而存在,而是為了要回應耶穌基督的呼召,為了實踐方使命而存在,.
傳統的規矩不是純粹為了方便我們自己的需要,而是為了協助人持守真理,讓人無論信徒或未信者更有效毀身跟隨耶穌,.
所以為了這個緣故我們需要建立良好的信仰傳統習慣,一方面防止我們的信仰過於恩秦,過於內罪和排外,.
同時亦要幫助我們敢於面向這個世界,夠膽去接觸未信者,並且樂意在他們當中,.
撥出上帝的愛和基督各樣美善的旨意,所以為這個緣故,盼望今天這段經文成為我們的提醒和激勵,.
讓我們可以更新我們一直所持守的傳統規矩和習慣,我們同心祈禱,好嗎?.
你會定時適時賜下你寶貴的話語,來提醒,更新我們的信仰生命,.
確實上帝,我們好像昔日的百姓,我們在信仰生活中也有很多不同的傳統和規矩,.
這一切都是為了讓人更加去認識你,更加懂得實踐你的真理,.
然而我們也知道,因為傳統規矩往往是出於人的想法,很多時候都會有不完美,不完整,甚至乎有軟弱的地方,.
需要時常的安著性,令你提醒,來更新活化,盼望主你幫助我們,.
讓我們記顧守教會良好的傳統的同時,我們也試試去檢視,這一切到底是否仍然吻合上帝的心意,.
能否配合這個時代的福音需要,能否造就讓姐妹更加熱愛上帝的話語,進行上帝的話語,.
能否更有效地讓未信的人可以進到教會,認識福音,.
願主你藉著這段寶貴的經文,幫助我們,成為我們生活的提醒和激勵,.
我們這樣的禱告交託,奉基督耶穌您的德性的名來祈求,阿們..
(影片完結).
\newpage



\section{馬可福音 8:1-21}
\label{sec:EMIvj5LfGAU}
\textbf{2025年8月31日崇拜直播｜楊珊珊傳道｜耶穌給門徒的十課－你們還不明白嗎?｜馬可福音8\_1-21}
\newline
\newline
連結: \href{https://youtube.com/watch?v=EMIvj5LfGAU}{\texttt{https://youtube.com/watch?v=EMIvj5LfGAU}} ~~~~ 語音日期: 2025-08-31
\newline
\newline
\hyperref[sec:0QX__zW2yAY]{\small{< < < PREV SERMON < < <}}
~
\hyperref[sec:index_chronic]{\small{[返順時目]}}
~
\hyperref[sec:xAjegbjPwhU]{\small{> > > NEXT SERMON > > >}}
\newline
\newline
馬可福音 8:1-21
\newline
\begin{longtable}{cl}
\hline
\hline
章節 & 經文 (和合本修訂版)\\
\hline
8:1 & \begin{tabularx}{0.7\textwidth}{X} 那時，又有一大群人聚集，沒有甚麼吃的。耶穌叫門徒來，說： \end{tabularx} \\ \\ \relax
8:2 & \begin{tabularx}{0.7\textwidth}{X} 「我憐憫這群人，因為他們同我在這裡已經三天，沒有吃的東西了。 \end{tabularx} \\ \\ \relax
8:3 & \begin{tabularx}{0.7\textwidth}{X} 我若叫他們餓著回家，他們會在路上餓昏，因為其中有從遠處來的。」 \end{tabularx} \\ \\ \relax
8:4 & \begin{tabularx}{0.7\textwidth}{X} 門徒回答：「在這野地，從哪裡能得餅使這些人吃飽呢？」 \end{tabularx} \\ \\ \relax
8:5 & \begin{tabularx}{0.7\textwidth}{X} 耶穌問他們：「你們有多少餅？」他們說：「七個。」 \end{tabularx} \\ \\ \relax
8:6 & \begin{tabularx}{0.7\textwidth}{X} 他吩咐眾人坐在地上，就拿著這七個餅祝謝了，擘開，遞給門徒，叫他們擺開，門徒就擺在眾人面前。 \end{tabularx} \\ \\ \relax
8:7 & \begin{tabularx}{0.7\textwidth}{X} 他們還有幾條小魚；耶穌祝謝了，就吩咐也擺在眾人面前。 \end{tabularx} \\ \\ \relax
8:8 & \begin{tabularx}{0.7\textwidth}{X} 他們都吃，並且吃飽了，收拾剩下的碎屑，有七筐子。 \end{tabularx} \\ \\ \relax
8:9 & \begin{tabularx}{0.7\textwidth}{X} 人數約有四千。耶穌打發他們走了， \end{tabularx} \\ \\ \relax
8:10 & \begin{tabularx}{0.7\textwidth}{X} 隨即同門徒上船，來到大瑪努他境內。 \end{tabularx} \\ \\ \relax
8:11 & \begin{tabularx}{0.7\textwidth}{X} 法利賽人出來盤問耶穌，要求他從天上顯個神蹟給他們看，想要試探他。 \end{tabularx} \\ \\ \relax
8:12 & \begin{tabularx}{0.7\textwidth}{X} 耶穌心裡深深嘆息，說：「這世代為甚麼求神蹟呢？我實在告訴你們，沒有神蹟給這世代看。」 \end{tabularx} \\ \\ \relax
8:13 & \begin{tabularx}{0.7\textwidth}{X} 他就離開他們，又上船往海的對岸去了。 \end{tabularx} \\ \\ \relax
8:14 & \begin{tabularx}{0.7\textwidth}{X} 門徒忘了帶餅，在船上除了一個餅，沒有別的食物。 \end{tabularx} \\ \\ \relax
8:15 & \begin{tabularx}{0.7\textwidth}{X} 耶穌囑咐他們說：「你們要謹慎，要防備法利賽人的酵和希律的酵。」 \end{tabularx} \\ \\ \relax
8:16 & \begin{tabularx}{0.7\textwidth}{X} 他們彼此議論說：「這是因為我們沒有餅吧。」 \end{tabularx} \\ \\ \relax
8:17 & \begin{tabularx}{0.7\textwidth}{X} 耶穌知道了，就說：「你們為甚麼因為沒有餅就議論呢？你們還不領悟，還不明白嗎？你們的心還是愚頑嗎？ \end{tabularx} \\ \\ \relax
8:18 & \begin{tabularx}{0.7\textwidth}{X} 你們有眼睛，看不見嗎？有耳朵，聽不到嗎？也不記得嗎？ \end{tabularx} \\ \\ \relax
8:19 & \begin{tabularx}{0.7\textwidth}{X} 我擘開那五個餅分給五千人，你們收拾的碎屑裝滿了多少個籃子呢？」他們說：「十二個。」 \end{tabularx} \\ \\ \relax
8:20 & \begin{tabularx}{0.7\textwidth}{X} 「又擘開那七個餅分給四千人，你們收拾的碎屑裝滿了多少個筐子呢？」他們說：「七個。」 \end{tabularx} \\ \\ \relax
8:21 & \begin{tabularx}{0.7\textwidth}{X} 耶穌說：「你們還不明白嗎？」 \end{tabularx} \\ \\
[1ex]
\hline
\hline
\end{longtable}
$^{1}$今日的經文是《馬學福音》八章一至二十一節.
經文記載在聖經新約的部分,58至59頁.
《馬學福音》八章一至二十一節.
那時又有許多人聚集,並沒有什麼吃的..
耶穌叫門徒來說:「我憐憫者眾人,.
因為他們和我在這裡已經三天,也沒有吃的了..
我若打發他們餓著回家,就必在路上困乏,.
因為其中有從遠處來的.」.
門徒回答說:「在這野地,從哪裡能得餅,叫這些人吃飽呢?」.
耶穌問他們說:「你們有多少餅?」.
他們說:「七個.」.
祂吩咐眾人坐在地上,就拿著這七個餅,.
竹蔗了,打開第二級門徒,叫他們擺開..
門徒就擺在眾人面前,又有幾條小魚..
耶穌竹了福,就吩咐也擺在眾人面前..
眾人都吃,並且吃飽了..
收拾剩下的零碎有七康子..
人數約有四千..
耶穌打發他們走了,隨即和門徒上船,.
來到大馬魯他境內..
發理塞人出來盤問耶穌,.
求祂從天上顯個神跡給他們看,想要試探祂..
耶穌心裡深深的嘆息說:「這世代為什麼求神跡呢?.
我實在告訴你們,沒有神跡給這世代看.」.
祂就離開他們,又上船往海那邊去了..
門徒望了帶餅,在船上除了一個餅,沒有別的食物..
耶穌祝福他們說:「你們要謹慎,.
防備發理塞人的教和希律的教.」.
他們彼此議論說:「這是因為我們沒有餅吧?」.
耶穌看出來就說:「你們為什麼因為沒有餅就議論呢?.
你們還不醒悟,還不明白嗎?.
你們的心還是如幻嗎?.
你們有眼睛看不見嗎?有耳朵聽不見嗎?.
也不記得嗎?.
我打開那五個餅,分給五千人..
你們收拾的零碎,裝滿了多少男子呢?」.
他們說:「十二個.」.
又打開那七個餅,分給四千人..
你們收拾的零碎,裝滿了多少漢子呢?」.
他們說:「七個.」.

$^{41}$耶穌說:「你們還是不明白嗎?」.
聖經獨筆..
(聖經獨筆).
大英姐妹平安!.
相信大家對剛才讀過的經文的場景一點也不陌生..
如果我們將馬福音八章的一至二十一節,分成三個段落,.
第一個段落是八章的一至十節,.
即是耶穌餵飽了四千人的神蹟..
大家記得上個月有一天是十號球賽嗎?.
記得吧?.
我那天真的很記得,因為我們那天去了一個下令的營會,.
我們特意由四日三夜變成三日三晚的營會..
我記得Raymond跟大家分享過六章三十一至四十四節,.
是說五餅二魚,耶穌餵飽了五千個男人的神蹟..
之後如果我們再看八章的十一至十三節,.
就是第二個部分..
接著回想起上星期的七章一至二十三節,.
一樣不知道為什麼耶穌坐了一艘船去到另一個岸邊,.
總是有一群法利賽人等著跟祂辯論和爭執..
而第三個部分就是八章十四至二十一節,.
是有關麵包或者餅裡面的討論,.
特別是教會的討論..
幫傳到七章二十四至三十四節,.
是說在希臘有頭有臉的妻子,.
她很機智和巧妙地接住了耶穌的一句話,.
即使在桌底下的小狗都可以來小孩那裡吃餅碎..
而今天經文最後的部分,.
說的教或者考慮,.
就是耶穌提出門徒要防備法利賽人和希律的影響..
當我們拍攝六至七章和今天八章的內容時,.
你會發現馬可福音的敘事模式很相似,.
圖案好像差不多..
不知道你會不會讀到聖經發現一些重複的地方,.
你的腦就開始做一個動作,.
就是歸類,.
歸類完之後大家會怎樣?.
很安心地飛走了,讀了它就算了..
不過我們這樣處理經文,.
會讓我們錯過一些經文很重要的訊息..
就好像我們肚子痛,.

$^{81}$去看醫生,.
你會知道肚子痛可能是因為消化不良,.
又可能是因為胃癌,.
雖然兩者之間有一個相似的徵狀,.
醫生不會單單因為我們消化不良而安排我們去做化療,.
因為醫生除了研究一些相似的地方外,.
也會仔細分析不同的地方..
所以我用這麼極端的例子告訴大家,.
我們研讀神的話語也是一樣,.
我們每逢看到差不多片段的經文,.
我們就要找不同..
所以今天很想和大家帶來一個新的嚴經的招式,.
就是見相似就要找不同,見相似找不同..
首先我們看看經文的第一個部分,.
八章一至十節,.
即是耶穌用了七個餅和幾條小魚餵飽了四千人的神蹟..
第一節頭兩個字,.
「那時又有許多人聚集」,.
「那時」是指何時呢?.
弟兄姊妹,有沒有人知道?.
不知道的話我來講解一下..
首先我們可以看回上文的七章三十一節,.
我們知道耶穌在這個時候,.
在一個地方,這個地方叫做底加波里的境內..
底加波里的希臘文的意思就是十座城,.
即是指加里里海東南邊十座希臘的城市..
我再強調聖經裡面的地貌,.
Landscape是很重要的,.
因為我們發現沿海的地區,.
其實通常我們都可以想像得到,.
或者知道經濟都是比較發達,.
資訊都是比較流通,.
所以人們都會比較有錢,.
或者接受一些新的思想..
不過當我們想起五餅二魚的神蹟,.
其實耶穌都是在加里里海附近發生的..
換句話說,五餅二魚的人都是比較有錢,.
資訊都是比較流通,.
加上他們對新事物都是比較開放..
所以我們要找不同的,.

$^{121}$最不同的其實就是這十個沿海的城市,.
是希臘的城市,.
所以裡面的人大多數都是外邦人,.
很聰明,就是外邦人多,.
所以只有很少的猶太人住在裡面,.
但原來不經不覺,.
耶穌基督的工作已經開始在外邦人裡面展開了,.
由開頭七章隱藏不願意人知道的情況下,.
在一個外邦的城市家裡,.
耶穌隔空醫好這個肥膩基婦人的女兒,.
然後他親自將一個耳又聾,.
舌又打結的人帶到一個遠離人煙的地方醫治她,.
而到八章一節,.
這個那時就是耶穌望著這四千個外邦人動了慈心的時刻,.
不過這次耶穌憐憫他們的原因和上一次五餅二魚是有點不同的,.
上一次是因為他看到那一班猶太人,.
好像羊沒有牧人一般,.
還記不記得那些羊仔,.
他們會有大近視,.
會將那些六戰香口膠當作草,.
吃壞肚皮,.
都是沒有一回事,.
所以當時耶穌一坐下就教導了那一班五千個男人,.
整天,.
這次這一班的外邦人不同,.
這次這一班的外邦人跟了耶穌多少天?.
三天,.
比之前那班人足足跟了兩天,.
他們已經跟到耶穌去到一個地步,.
沒有東西吃了,.
留意上一次那個地方,.
門徒是先過耶穌察覺糧食的問題,.
所以他就主動提議,.
不如大家去附近的村莊買餅,.
不過這次卻是主耶穌親自主動告訴門徒糧食的問題,.
耶穌連問這一班的外邦人,.
很多都是很遠水路來到這個野地,.
擔心他們還沒有走到回家的半路已經暈倒在路上,.
但是這次門徒的回答,.
並不是問耶穌錢從哪裡來?.

$^{161}$要買這麼多個餅?.
你不要以為這一班門徒經歷過五餅二魚的神蹟,.
就會自動自覺記起,.
耶穌就是親自可以餵飽五千個男人的有大能的神,.
因為在聖經的六章52節,.
經文已經說了這班門徒其實不明白五餅二魚的事情,.
他們心裡還是很如幻,.
意思是他們的心比石頭還要硬,.
即使給了什麼都塞不進去,.
所以他們是沒有辦法明白耶穌可以走到這個神蹟,.
他們忘記了,.
所以在八章四節,.
看到這班門徒回答耶穌的話,.
他們仍然是沒有明白耶穌有神蹟的能力,.
或許他們忘記了,.
所以他們問了耶穌基督一個問題,.
「不是吧,耶穌,.
這班門徒是一個完完全全的野地,.
究竟我們在哪裡才可以得到這些餅,.
也不是說買,是得到,.
是令到這四千人吃飽呢?」.
其實這個問題反映出.
當時門徒一個屬靈生命的迷思或掙扎,.
明明他們以前經歷過五餅二魚,.
為什麼忘記了耶穌是有能力餵飽幾千人呢?.
明明我們以前切切實實經歷過耶穌基督的神蹟,.
為什麼今天我們只是看著困境,.
而看不到比困難更加大的那位行神跡的上帝呢?.
如果今天有弟兄姊妹在困難當中,.
我也祈求上帝憐憫和幫助你,.
讓你能夠看見,聽見並且知道.
耶穌基督是與你同在,.
祂一直在你的身邊,.
正如在那個時候,.
耶穌真的在那四千人當中,.
祂看見了這群人實際上的需要,.
在這四千人開口之先,.
耶穌已經憐憫他們,.
想著要幫他們解決他們肚餓的問題..
那麼耶穌又對這群門徒怎樣呢?.

$^{201}$其實耶穌沒有離開過這群門徒,.
即使他們的屬靈狀況不是很好,.
在八章五節,.
耶穌在那一刻只問了他們一個問題,.
就是你們有多少個餅?.
於是門徒就找了,.
有七個,.
魚呢?.
在這裡經文為了凸顯,.
這裡的魚是與之前五餅二魚的魚不同,.
大家知不知道這個魚有什麼不同?.
沒錯,是小魚,.
多了一個小字,.
這個小字很重要,.
七個餅,幾條小魚,.
餵飽了四千人,.
當中包括男女老幼,.
剩下的食物有多少個康子?.
七個康子,.
雖然剩下的數量比上次五餅二魚剩下的有多少個男子?.
十二,.
我們旺州的弟兄姊妹很聰明,.
雖然聽起來數量好像少了,.
不過這次的容器不同,.
這次的容器是康子,.
康子是比男子更大的一隻腩,.
所以這次耶穌所行的神蹟,.
裡面的數量絕對不會比上次五餅二魚行在猶太人裡面少,.
為什麼我要這麼強調這些細節,.
和這些不同?.
就是因為這些都是馬可作者深思細密的地方,.
叫我們千萬不要歸納,.
這不是一個簡單的五餅二魚的2.0,.
而這個神蹟是有它的獨特性,.
是屬於愛邦人的神蹟,.
昔日神,.
他如何在曠野裡為以色列人預備嗎哪,.
今天耶穌也為了這四千個愛邦人,.
以竹蔗和麥餅供應他們身體所需,.
這個神蹟正正告訴猶太人和愛邦人,.

$^{241}$彌賽亞的筵席從來不只是給猶太人,.
就像舊約的《以賽亞書》25章6節說,.
「萬君之耶和華,必為萬民擺設筵席」,.
而耶穌主動去憐憫這一大群愛邦人,.
正正提醒我們,.
耶穌基督的救贖從來不只是屬於我們這一群基督徒,.
同樣地,今天耶穌仍然主動記掛著一些在祂救恩以外的人,.
作為門徒的我們,.
也有繼承耶穌基督的心腸嗎?.
在過去的星期五晚,.
我聽了一對宣教士夫婦的分享,.
不知道你們有沒有聽過,.
他們的故事是這樣的,.
他們本身是一個住在普通烏克蘭國家的華人家庭,.
當2022年俄烏戰爭開始發動的時候,.
他們一家帶著一個牧師逃難到歐盟,.
而神巧妙地很快讓他們得到一個歐盟的身份,.
亦輾轉反側,這個家庭的太太在意大利讀神學,.
而當這個太太在異神畢業的時候,.
有人提議這個太太,.
「不如你家裡已經顛簸了這麼久,.
不如你留在意大利,.
意大利生活多好,安安定定,.
再加上我們這邊需要牧者.」.
這個太太想了一會兒,.
她說「沒錯,在意大利生活比較穩定,.
不過相比起意大利,.
可能一個城市沒有華人的牧者,.
但是她想去的工場有更多烏克蘭的難民,.
工場那邊可能是整個國家都沒有華人的牧者,.
再加上她會說俄語,.
因為她本身是在烏克蘭住,.
她可以去當地幫助烏克蘭難民的朋友.」.
當她想到她在逃難那段很艱難的日子,.
明明有很多人曾經為她遮風擋雨,.
所以當他們現在有能力的時候,.
為什麼她不去為那些其他有需要的人遮風擋雨呢?.
大英姐妹,讓我們記掛著,.
其實還有很多在救恩以外的人,.
耶穌基督仍然呼召他們去吃耶穌基督救恩這個很大的筵席,.

$^{281}$我們又願不願意一起去出一分力呢?.
當耶穌行完神蹟之後,.
馬可福音一貫的作風都好像跟我們香港人一樣,.
很大聲的,.
因為在八章九節,.
耶穌馬上打發了四千人離開之後,.
十節隨即和門徒上了船,.
去到另一個地方叫大瑪魯塔,.
其實是加利利海的西邊,.
也是猶太人聚集的地方,.
耶穌一來到那裡,.
就已經有一群法利賽人收到消息,.
要在那裡準備盤問耶穌,.
弟兄姊妹留意聖經在第十一節,.
用的字眼不是辯論,.
不是討論聖經,.
而是盤問,.
你不要以為這群法利賽人斯斯文文,.
內裡不是open,.
是violence,.
即是我上星期和大專生說非暴力溝通,.
我們要帶著一個沒有偏見的前設去和人溝通,.
這群法利賽人不是,.
他們是帶著前設和敵意盤問耶穌,.
聖經說這個求不是這麼簡單,.
而是要求耶穌,.
為的是他們要試探耶穌,.
當我們見到試探這個動詞,.
我們會想起在聖經裡的馬可福音,.
一章十三節,.
曾經哪一個試探耶穌?.
是撒旦,.
所以其實試探這個字出現了四次,.
一章撒旦試探耶穌,.
在這裡法利賽人試探耶穌,.
還有十章二節,.
十二章十五節,.
都是指法利賽人試探耶穌,.
後面兩次其實是想嘗試捉耶穌自殺,.
將耶穌陷入兩難的局面,.

$^{321}$所以在馬可福音裡,.
耶穌被兩類物種或人種試探,.
第一類是撒旦,.
第二類是法利賽人,.
不過我們馬上想,.
其實不用試探耶穌,.
耶穌真的行過很多神蹟,.
他叫死人復活,.
黑眼得看見,.
聾的又聽得見,.
又醫病,又講鬼,.
又有不同的神蹟,.
不過事實上,.
這一班法利賽人,.
不是我剛才所說,.
他們想要的神蹟,.
他們想要的神蹟其實,.
其實他們,.
譬如之前三章二十二節,.
他們已經承認耶穌可以講鬼,.
不過他們認為,.
耶穌是靠著撒旦的權利去講鬼,.
他們其實想要的那種神蹟,.
又或者sign是甚麼呢?.
他們想要的是天裂開了,.
聖靈彷似甲子降臨在耶穌身上,.
又有聲音從天上來說,.
你是我的愛子,我喜悅你..
很明顯,其實我這裡說的,.
大家是否很熟悉這個經文?.
其實是在說耶穌的浸泥,.
很明顯,這班法利賽人,.
是miss了耶穌的浸泥,.
而是要天上的神蹟..
他們真的不是要一個普通的神蹟,.
他們想要的證明,.
是耶穌能夠證明自己的權柄是來自上帝..
耶穌面對著他們的要求,.
我們見到十二節,.
耶穌這樣回覆,.

$^{361}$聖經說,.
耶穌深深在他的靈裡面嘆息..
這個嘆息,.
我想像未完全表達到當時耶穌的情況,.
其實這裡更加可以告訴我們,.
原來我想像耶穌在他人性的範圍裡,.
已經去到一個忍無可忍的極限..
我想像耶穌說,.
他們又來了,.
耶穌對這班法利賽人,.
已經感到很厭倦..
他發出一個看似問號,.
其實是一個感嘆號的問句,.
為什麼這個世代還要求神蹟呢?.
然後耶穌回覆他們說,.
再沒有神蹟給這個世代看..
可能你會問,.
為什麼耶穌回覆這班法利賽人的感覺,.
不像平時我們認識的耶穌呢?.
因為我們認識的耶穌,.
應該都是很有耐性才對..
沒錯,.
其實這裡說的是,.
耶穌已經對這班人很忍耐,.
忍耐到一個極限,.
是一個最大的忍耐..
弟兄姊妹,.
我們會否忽略了一個很重要的事實,.
我們有沒有忘記了神的忍耐是怎樣呢?.
沒錯,.
聖經提到神的忍耐,.
提到神的寬容,.
提到祂是長久的忍耐..
我們亦在很多經文裡,.
一次又見到神的忍耐和憐憫..
但是正正因為這樣,.
有時候我們會否開始以為,.
神的忍耐是無止境的呢?.
這個想法其實是非常危險的..
所以弟兄姊妹,.

$^{401}$我們要小心,.
當我們選擇執迷不悟的時候,.
無論神提醒過我們多少次也好,.
祂給了多少忍耐我們,.
而當我們仍然在罪裡面,.
我們便察覺不到神,.
其實祂有多忍耐我們,.
和祂有多包容我們..
我們對耶穌基督的耐性,.
是不可以奉旨的..
因為耶穌真的會離開那些.
專挑拘機,.
執迷不悟的法利賽人..
我們不可以一邊稱呼耶穌是主啊主啊,.
一邊卻選擇看不見,.
聽不到神的神蹟,.
都是以為自己那套是對的,.
看見耶穌基督的神蹟的出現,.
當作只是日常,.
根本不當耶穌的神蹟是神蹟..
最終耶穌看見這些情況,.
祂只能夠在靈裡面深深地歎息,.
不會再讓我們看見任何神蹟..
於是在十三章,.
我們可以看見耶穌選擇離開了.
這群法利賽人,.
準備乘船去下一個目的地..
而到船的上邊,.
我們就去到今天的經文第三個部分..
首先我重溫第一個部分,.
就是八章的一至十節,.
說的是一群愛邦人,.
經歷了耶穌用七個餅,.
幾條小魚餵飽四千人的神蹟..
第二個部分就是十一至十三節,.
說的是一群法利賽人,.
選擇不相信,不承認,.
他們過往一直看見和聽見,.
還有他們選擇忘記耶穌的神蹟是神蹟..
大家記得嗎?.

$^{441}$大家很聰明..
第三個部分就是十四至二十一節,.
門徒上了船,.
發現帶不夠餅,只有一個餅,.
多少人分呢?.
十三個人分,沒錯..
於是門徒就為了分餅的事而煩惱..
但這次耶穌似乎沒有像那四千人那樣動了持心,.
耶穌不在乎他們夠不夠吃,.
反而祝福他們,.
說你們要謹慎,.
要防備法利賽人的教和希律的教..
在這裡我要解釋教或是孝是一個什麼概念..
首先,教或孝是已經發了孝的麵團而組成的東西..
即是這些教我們可以在家裡做,.
我想起以前有些歐洲的朋友,.
姐妹都會做麵包,.
然後他們會回自己的娘家拿些孝帽過來,.
說這些可以用來做麵包,.
麵包就會慢慢發脹..
其實都是同一個道理,.
就是以前的麵團,.
剩下發了的教,.
你將這些教放在新的麵團裡,.
這些麵團就會在很短的時間裡慢慢長大,.
變大,然後你就可以焗,.
焗出來就有很大的麵包吃..
很多時候我們看聖經都會發現,.
在天國的比喻裡面,.
有用過這些的教作為素材,.
比喻為天國..
但是除了這個正面的象徵意義之外,.
其實教很多時候都象徵著一些不潔,.
甚至是形容罪的滲透力,.
又或者一些錯誤教導的影響..
正如在加拉太書五章九節,.
它說一些不好的教導,.
一些的面教可以令整個麵團都會完全發起來,.
不能不要的..
究竟耶穌基督在說的這些教是甚麼呢?.

$^{481}$要防備一些不好的教,.
當我們並列其他的福音書,.
我們會發現,.
耶穌說法利賽人的教是有不同種類的..
在馬太福音,.
耶穌說的法利賽人,.
連殺都哥人也包括,.
他們的教是說一些錯誤扭曲了的教導,.
而路加福音說的法利賽人的教就是假冒為善..
馬可在這裡並沒有明確說法利賽人的教是甚麼..
可能更加令我們摸不著頭腦,.
他要加上希律的教..
有些學者說,.
希律和法利賽人相似的地方很少,.
所以有些學者會歸類,.
其實兩者都可能很想得到一些他們想要的神蹟,.
其實也很符合十二節,.
我們看到關於法利賽人的特徵,.
就是他們想要耶穌基督去證明自己是神..
希律那裡其實不是很明顯,.
所以當大家讀到十二章的時候,.
馬可會再跟我們揭露多一點,.
如果大家小組查馬可福音查到十二章的時候,.
你也可以跟我分享一下,.
你找到希律的教是甚麼,.
如果你的答案跟我說是合理,.
我也會送你一份神秘的小禮物..
不過在這裡我很想說,.
法利賽人的教就是我們上一個部分,.
十一至十三節所說的,.
耶穌提醒這一班門徒,.
要防備和不要學這一班法利賽人,.
執迷不悟,.
不要只是選擇相信自己想相信的神蹟,.
而忽視了耶穌基督每一天真真正正在他們旁邊的神蹟..
法利賽人要求的是神從天而來,.
他們不是想要耶穌那一隻神蹟,.
他們想要的是權力和地位,.
但是耶穌告訴我們,.
天國的實踐從來都不是以實質的權力去霸佔或侵佔..

$^{521}$耶穌說的是我們要做一個奴僕的君王,.
他的神蹟就是神導誠肉身,.
神願意為罪人走上這條十字架受苦的道路..
弟兄姊妹,在這裡我很想說的是,.
其實耶穌最關心門徒的不是他們有沒有餅吃,.
而是因為耶穌已經是他們生命的糧,.
沒有麵包吃的困難,.
耶穌是會幫我們解決,.
因為剛剛祂先給我們看了兩個餵飽幾千人的神蹟,.
他們也身歷其境,.
耶穌最想告訴他們的是,.
我們真正要防備的困難,.
是我們要留意在我們的餅裡面,.
有沒有法利菜的酵母,.
有沒有那些酵在裡面,.
因為只要我們有一丁點那些酵母,.
這些酵母就會像老鼠屎一樣有毒,.
我們整塊麵包就可以丟掉,.
不能丟掉..
耶穌不斷提醒這班門徒,.
祂最後也問了他們兩個問題,.
就是有兩次的神蹟,.
他們真的不需要擔心餅的問題,.
反而他們要記得和承認,.
耶穌真的有走過的神蹟..
其實我不知道你們是否發現,.
這些門徒和耶穌的對答,.
其實也像那班法利菜一樣,.
因為其實那班門徒也忘記了,.
耶穌有走過的神蹟..
所以弟兄姊妹,.
我們真的要很小心,.
很小心防備,.
自己在一些什麼文化,.
意識形態,又或者朋友圈裡面,.
如果我們只和一些人,.
他們只顧著吃喝玩樂,.
我們就會很容易受他們影響,.
只顧著吃喝玩樂..
如果我們和一些充滿敵意的懷疑者相處,.

$^{561}$我們整個人也會充滿敵意,.
因為法利菜人的孝,.
他們的盲點就像毒藥,.
會慢慢一點一點地,.
沉食我們愛父的心,.
會毀滅我們..
最後,不知道大家還記不記得,.
我教大家今天一個茶經的招式,.
就是見相似找不同..
雖然門徒真的和那班法利菜人,.
他們一樣都是看不見,.
聽不見,聽不到,.
不記得,不明白,.
很相似,.
不過不同的是,.
耶穌沒有離開那班門徒,.
耶穌其實還在船上,.
叫那班門徒小心一點,.
耶穌還仍然在指出他們屬靈的盲點..
所以弟兄姊妹,.
耶穌基督是我們生命的糧,.
祂的恩典是救我們用,.
祂已經和我們經歷了,.
可能一次,兩次,很多次的神蹟,.
弟兄姊妹,我們有眼睛,.
我們願不願意看得見,.
我們有耳朵,.
願不願意聽得明白,.
而我們又願不願意用心去記得呢?.
耶穌當日的提問,.
也在問我們,.
代表著耶穌的耐性,.
我們今天的心,.
是否仍然這麼硬呢?.
盼望我們都從我們這些屬靈盲點中回轉,.
回應我們的愛主耶穌基督的耐性..
\newpage



\section{馬可福音 8:22-9:1}
\label{sec:xAjegbjPwhU}
\textbf{2025年9月7日崇拜直播｜陳冠成牧師｜耶穌給門徒的十課－基督的本質•門徒的職份｜馬可福音8\_22-9\_1}
\newline
\newline
連結: \href{https://youtube.com/watch?v=xAjegbjPwhU}{\texttt{https://youtube.com/watch?v=xAjegbjPwhU}} ~~~~ 語音日期: 2025-09-07
\newline
\newline
\hyperref[sec:EMIvj5LfGAU]{\small{< < < PREV SERMON < < <}}
~
\hyperref[sec:index_chronic]{\small{[返順時目]}}
~
\hyperref[sec:HtBaZ5KPIHs]{\small{> > > NEXT SERMON > > >}}
\newline
\newline
馬可福音 8:22-9:1
\newline
\begin{longtable}{cl}
\hline
\hline
章節 & 經文 (和合本修訂版)\\
\hline
8:22 & \begin{tabularx}{0.7\textwidth}{X} 他們來到伯賽大，有人帶一個盲人來，求耶穌摸他。 \end{tabularx} \\ \\ \relax
8:23 & \begin{tabularx}{0.7\textwidth}{X} 耶穌拉著盲人的手，領他到村外，就吐唾沫在他眼睛上，為他按手，問他：「你看見甚麼？」 \end{tabularx} \\ \\ \relax
8:24 & \begin{tabularx}{0.7\textwidth}{X} 他抬頭一看，說：「我看見人，他們好像樹木，並且行走。」 \end{tabularx} \\ \\ \relax
8:25 & \begin{tabularx}{0.7\textwidth}{X} 隨後耶穌又按手在他眼睛上，他定睛一看，就復原了，樣樣都看得清楚了。 \end{tabularx} \\ \\ \relax
8:26 & \begin{tabularx}{0.7\textwidth}{X} 耶穌打發他回家，說：「連這村子你也不要進去。」 \end{tabularx} \\ \\ \relax
8:27 & \begin{tabularx}{0.7\textwidth}{X} 耶穌和門徒出去，往凱撒利亞‧腓立比附近的村莊去。在路上，他問門徒：「人們說我是誰？」 \end{tabularx} \\ \\ \relax
8:28 & \begin{tabularx}{0.7\textwidth}{X} 他們對他說：「是施洗的約翰；有人說是以利亞；又有人說是先知中的一位。」 \end{tabularx} \\ \\ \relax
8:29 & \begin{tabularx}{0.7\textwidth}{X} 他又問他們：「你們說我是誰？」彼得回答他：「你是基督。」 \end{tabularx} \\ \\ \relax
8:30 & \begin{tabularx}{0.7\textwidth}{X} 於是耶穌切切地囑咐他們不可對任何人說起他。 \end{tabularx} \\ \\ \relax
8:31 & \begin{tabularx}{0.7\textwidth}{X} 從此，他教導他們說：「人子必須受許多的苦，被長老、祭司長和文士棄絕，並且被殺，三天後復活。」 \end{tabularx} \\ \\ \relax
8:32 & \begin{tabularx}{0.7\textwidth}{X} 耶穌明白地說了這話，彼得就拉著他，責備他。 \end{tabularx} \\ \\ \relax
8:33 & \begin{tabularx}{0.7\textwidth}{X} 耶穌轉過來看著門徒，斥責彼得說：「撒但，退到我後邊去！因為你不體會神的心意，而是體會人的意思。」 \end{tabularx} \\ \\ \relax
8:34 & \begin{tabularx}{0.7\textwidth}{X} 於是他叫眾人和門徒來，對他們說：「若有人要跟從我，就當捨己，背起自己的十字架來跟從我。 \end{tabularx} \\ \\ \relax
8:35 & \begin{tabularx}{0.7\textwidth}{X} 因為凡要救自己生命的，必喪失生命；凡為我和福音喪失生命的，必救自己的生命。 \end{tabularx} \\ \\ \relax
8:36 & \begin{tabularx}{0.7\textwidth}{X} 人就是賺得全世界，賠上自己的生命，有甚麼益處呢？ \end{tabularx} \\ \\ \relax
8:37 & \begin{tabularx}{0.7\textwidth}{X} 人還能拿甚麼換生命呢？ \end{tabularx} \\ \\ \relax
8:38 & \begin{tabularx}{0.7\textwidth}{X} 凡在這淫亂罪惡的世代，把我和我的道當作可恥的，人子在他父的榮耀裡與聖天使一同來臨的時候，也要把那人當作可恥的。」 \end{tabularx} \\ \\ \relax
9:1 & \begin{tabularx}{0.7\textwidth}{X} 耶穌又對他們說：「我實在告訴你們，站在這裡的，有人在沒經歷死亡以前，必定看見神的國帶著能力臨到。」 \end{tabularx} \\ \\
[1ex]
\hline
\hline
\end{longtable}
$^{1}$今天經文是在《馬可福音》第八章二十二節到九章一節.
在我們教會聖經的新約第五十九頁.
找到的話請聽我突出.
他們來到白菜大,有人帶一個盒子來,求耶穌摸他.
耶穌拉著盒子的手,領他到村外,就套唾沫在他眼睛上,按手在他身上,問他說:你看見什麼了?.
他就抬頭一看,說:我看見人了,他們好像樹木,並且行走.
隨後又按手在他眼睛上,他定眼一看,做復了緣,樣樣都看得清楚了.
耶穌打發他回家,說:連這村子你不要進去.
耶穌和門徒出去,坐凱撒利亞腓立比的村莊去,在路上問門徒說:人說我是誰?.
他們說:有人說是斯森的約翰,有人說是以利亞,又有人說是先知裡的一位.
又問他們說:你們說我是誰?.
彼得回答說:你是基督,耶穌就禁戒他們,不要告訴人.
從此,他教訓他們說:人子必須受許多的苦,被長老,祭司長和民事氣絕,並且被殺,過三天復活.
耶穌明明的說者說:彼得就拉著他,勸他.
耶穌轉過來,看著門徒,就責備彼得,說:撒旦,退我後面去吧.
因為你不體諒神的意思,只體貼人的意思.
於是叫眾人和門徒來,對他們說:若有人要跟從我,就當捨己,背起他的十字架來跟從我.
因為凡要救自己生命的,必喪掉生命.
凡為我和福音喪掉生命的,必救了生命.
人就算賺了全世界,賠上自己的生命,有什麼益處呢?.
人還能拿什麼換生命呢?.
凡在這淫亂惡毒的世界,把我和我的道當可恥的.
人子在他父的榮耀裡,和眾天使降臨的時候,也要把那人當可恥的.
耶穌又對他們說:我實在告誓你們,站在這裡的,有人在沒嘗死味以前,必須看見神的國大有能力降臨.
(主)請姐妹平安.
(弟弟)平安.
不知道大家有沒有留意到,特別是你常看馬福音,你會發覺其實.
無論耶穌遇到什麼奇難雜症,通常都是一剔過就能治好病人.
只要他一開口,甚至一出手,病人就會立即康復.
不過你留意剛才所讀的經文,記載著耶穌醫治的那位乞子的神蹟就很不同.
你看見耶穌把治療過程分了兩次,很特別.
並且在兩次之間,他會詢問乞子到底治療的進展如何,看見什麼.
很明顯我們知道,不是耶穌不能一次治好,耶穌絕對有能力一次治好.
而是這個乞子需要時間慢慢康復.
事實上你會看到,馬可為何把醫治乞子的神蹟記載在馬福音的上半部分的結束段落.
你知道第八章是馬福音的一半,為何在這個時候把神蹟放進去.
其實是要說明一件事,就是在目前為止,門徒仍然未能領悟到耶穌的真正身份是什麼.
他們仍然處於屬靈核眼,屬靈上的眼睛是瞎了的.
他們不太清楚耶穌是誰,特別是你記得上一回珊珊傳道跟我們說.
門徒因為忘記帶餅上船,於是被耶穌嚴厲斥責,罵到他們狗血淋頭.

$^{41}$心還是如橫,明明已經一起餵飽五千人四千人.
仍然是有眼也看不到,有耳也聽不見,你看到耶穌很生氣.
雖然門徒的屬靈表現不似預期,耶穌真的覺得恨鐵不成鋼.
但你知道耶穌不會放棄他們,所以稍後耶穌會像醫治乞子一樣.
分階段慢慢處理門徒在屬靈上的體不見,幫他們逐漸看清楚.
其實耶穌到底是誰,他的真正身份是什麼,以及門徒真正的使命到底是來做什麼.
我們看二十七節,開始看這個段落.
我們會發覺這是一個新的段落,他們準備由凱撒利亞菲勒比慢慢移去耶路撒冷.
之前你們留意,科林書一直記載耶穌在加利利一帶活躍.
耶穌和門徒來到凱撒利亞菲勒比,我們說這是一個外邦人的城市.
除了敬拜羅馬君王凱撒之外,其實也充滿了各式各樣的偶像敬拜.
耶穌特意選擇這個在異教臨立的地方,來確認自己的真正身份.
來清楚表示,在這些外邦偶像的面前,他才是人類真正的主宰.
在路上,耶穌問門徒,人們怎樣說我是誰.
門徒回答,有人說你是被希律斬首,但已經復活的施浸約翰.
他們認為,約翰的能力已經轉移到耶穌身上.
也有人認為,耶穌就像昔日的以利亞先知一樣.
是舊約,上帝預言將要差派來復興以色列的鬼位的先鋒.
當然也有人認為,耶穌就像其他先知一樣,沒有能力宣講上帝的話語.
耶穌當然知道人們怎樣去想和說他.
所以其實頭一條問題,是要來對比其實門徒對耶穌的理解.
到底是不是和外人一樣,還是有分別.
於是,耶穌就隨即問門徒第二個問題.
你們呢?你們又認為我是誰?.
你留意門徒之前,一直叫耶穌做什麼?.
就是說,夫子,即是老師,就算耶穌曾經幫他們平靜風浪.
福音書也記載,他們也很困惑,這是什麼水?.
連風浪也要聽他們說,他們仍然不明白耶穌的身份.
但現在你會看到,馬福音來到第八章,已經過了一半.
終於有人可以說出耶穌的真正身份.
你看到彼得在門徒當中脫穎而出.
他宣認耶穌是基督,是上帝所差派來的救主.
我們也知道,基督是希臘文.
等同於舊約的 Messiah,彌賽亞.
意思是一個受高者.
受高納的人,例如舊約的君王,祭司或先知.
其實就是象徵著上帝選了他,差派他去完成一些特定的使命.
事實上,在舊約先知書曾經預言.
上帝會為大衛加興起一位基督或彌賽亞.
來帶領百姓脫離罪惡的困所,復興百姓的信仰生命.

$^{81}$而隨著以色列人經歷亡國被擄.
幾百年來不斷被不同的外邦政權來統治.
其實百姓慢慢發展了一種思想.
他們會覺得上帝要興起的這位基督,這位救主.
是會帶有政治和軍事色彩.
是會帶領他們推翻那些管轄他們的政權.
幫他們,不是純粹復興信仰生命.
是要幫他們恢復那塊土地,幫他們復國.
事實上,在耶穌時代的猶太人.
包括彼得,包括所有的門徒.
他們都是這樣想的.
他們一直對基督都抱有這種的奇想.
不過,普羅大眾對基督這種一廂情願的想法.
就不等於是對的.
所以你看到耶穌很特別,耶穌沒有讚彼得.
沒有讚彼得,對的,你有慧根,答對了.
反而是禁止他們不可以向人透露耶穌真正的身份.
除了我們知道耶穌明白,因為祂的時候還未到.
祂不想引起群眾的騷動或者混亂之外.
有一個更重要的原因是.
門徒認出耶穌是基督是一回事.
不等於他們真的清楚基督的本質到底是怎樣.
他們不一定明白,其實上帝差派耶穌來是要做些什麼.
於是乎你看到,耶穌緊隨就向門徒來教導.
其實作為基督的他,將會面對一些什麼的遭遇.
第三十一節,你會看到這是整本科目書裡面.
耶穌首次預告,注意,他說了必須.
必須受許多的苦,被猶太人的宗教領袖氣絕.
並且必須被殺害,但也必須三日後復活.
我們今天當然對耶穌這番話不會有什麼疑惑.
我們是理所當然接受,因為我們是新約底下的信徒.
但你要想像對昔日的門徒來說,這番話其實是摸不著頭腦.
是匪夷所思,耶穌的教導可以說是完全顛覆了.
或者借用保羅的說法,叛倒了所有的猶太人.
叛掉了門徒心目中一直對基督的認知,特別是彼得.
你看到和合本的翻譯其實是比較溫和的.
彼得只是拉耶穌一旁勸他不要這樣.
其實表達不了他和耶穌之間的張力.
印在大家的程序表裡,新漢語譯本其實也比較到位.
當耶穌清楚宣稱他是一位必須受苦,被厭棄,被殺害的基督的時候.

$^{121}$彼得是斷然拒絕的,並且他是拉耶穌一旁.
不是勸他,而是責備他.
雖然作者沒有說責備的內容是什麼,但大概你也可以想像到.
當耶穌宣告他是一位受苦的基督的時候.
彼得第一個衝出來拉著耶穌,說耶穌不可以這樣,萬萬不可.
基督怎麼可以這樣?.
確實彼得有沒有道理呢?.
也是有的,因為彼得一直親眼看到耶穌所行的異能神蹟.
平靜風浪又見識過,開聲風浪就停止了.
開口污穢就止住了.
說一句話,摸一摸,病人就得醫治.
用幾塊餅就可以餵飽幾千人.
耶穌更新稱自己擁有赦罪的權柄.
並且他自己有權決定到底安息日,什麼應該做,什麼不應該做.
其實當你這樣想的時候,這位這麼厲害,這麼有權柄的人物.
如果他是基督的話,怎麼可能會被人棄絕,被殺和失敗收場.
不可能的,反差的.
上帝的劇本沒理由這樣寫.
所以彼得接受不了這位失敗的救主怎麼可能是上帝拯救我們的一部分.
反過來,大家試一下用一點想像力.
如果耶穌的預告內容是另一番說話.
他會跟明圖說,稍後我們就會揭竿起義.
我們會推翻羅馬帝國的管治.
你知道過程中難免會有受苦.
難免會流血,甚至有人會犧牲.
不過你放心,我們最後會搞定,我們最後會勝利.
如果耶穌是這樣預告,你猜彼得會有什麼反應.
還會不會拉耶穌一邊罵他,肯定不會.
他反而會立即響應.
主啊,我們等了這天不知道多久了.
什麼時候?在哪裡開始?.
需要我多叫一些手足來嗎?.
這個就很符合彼得心中的基督形象.
很符合普羅大眾對救主彌賽亞的看法.
所以等姐妹你會看到.
彼得之所以如此抗拒耶穌這個受苦的教導.
不是因為他對基督一點認識都沒有.
不是因為他無知.
相反是他對基督的前設太過深.
他認定了基督,你一定是這樣.

$^{161}$你一定要以勝利者的姿態出現.
你一定一出來就擺平所有惡勢力.
懲罰一切的惡人壞人.
就算遇到宗教領袖的反對.
我們知道法理賽人和文士都不太妥理.
但你既然是基督,你那麼厲害.
你是一定能搞定.
你一定有能力應付他們.
所以彼得接受不到.
為何上帝揀選的基督必須受死.
這就是他的問題癥結.
就算耶穌之後說他會復活.
大概彼得已經完全聽不進去.
彼得的信仰矛盾是.
他既承認耶穌是基督.
但同時又很想耶穌指點他.
你要怎樣做基督.
正如有不少信耶穌的朋友.
都有這個情況.
會一廂情願去想像.
耶穌應該怎樣.
耶穌應該怎樣站在我旁邊.
耶穌應該怎樣幫我滿足我的需要.
不過,你看到耶穌斬釘切鐵表明.
他要怎樣做基督.
不由門徒作主.
而是取決於上帝的旨意.
33節,耶穌嚴厲斥責彼得.
他稱呼彼得為撒旦.
撒旦,退我後面去吧.
他要和彼得在這一刻劃清界線.
因為彼得阻止耶穌受苦受死.
這個舉動其實耶穌視之為.
等同於當日他所面對撒旦的試探.
就是不看貼上帝的旨意.
看貼人的思想.
注意耶穌動作的細節.
彼得原本拉著耶穌.
和耶穌面對面罵彼得.
然後耶穌轉身.

$^{201}$用背部掌摑彼得.
然後面向所有門徒訓斥彼得.
耶穌要讓彼得和門徒分清楚.
到底誰走在前頭.
誰跟隨誰.
彼得既然認信我是基督.
你就要老老實實跟隨我.
聽我的吩咐.
所以你看到耶穌轉身的細節.
其實很有意思.
他要讓彼得弄清楚.
是讓你跟從我.
不是我聽你的.
34節耶穌繼續教導門徒和眾人.
既然我是一位受苦的基督.
那麼門徒隨之而來.
就要怎樣跟隨耶穌.
什麼才是正確的門徒職份.
34節其實包含了三件事.
首先是寫記.
如果按照上文的意思.
其實是要捨棄對耶穌一廂情願的錯誤想法和期望.
不可以像彼得那樣.
用自己的想法來盤算耶穌的行動.
懂得寫記.
其實就是懂得要求自己.
體貼上帝的意思.
而不體貼自己的意思.
這是第一件.
第二件就是背起十字架.
這裡呈現的畫面.
其實是說當時的人.
什麼時候會被人背起十字架.
就是當你被判了死囚的時候.
羅馬帝國就會要求死囚.
你背起自己的十字架.
他不是背起整個十字架兩條木.
只是背起橫的那條.
一直背著橫木.
由皆是眾直至到刑場被釘上.

$^{241}$耶穌要求門徒為福音的緣故.
為了祂的緣故.
甘願加入一直被人唾棄被人氣絕的行列.
提醒門徒要有所覺悟.
做好為主犧牲的準備.
我們也知道在第一世紀.
堅持要信耶穌.
很大機會下場就是為主殉道.
所以耶穌隨即發出第三個要求.
就是呼籲門徒要效法祂.
走祂那條受苦受死的道路.
耶穌很清楚說明.
祂不需要一群只是前呼後擁.
純粹驚嘆祂神跡奇事的跟隨者.
而是要求門徒一個一個.
背著祂的十字架.
跟隨著前面已經背著十字架的耶穌.
一路走.
一起去信服.
配合上帝的拯救計劃.
事實上想起這個畫面.
其實一點也不溫馨浪漫.
想像一下一個一個.
背著自己的橫幕.
跟著耶穌上去刑場受死.
其實是很可怕的.
會叫人卻步.
叫人不敢認耶穌為主.
不知道大家還記不記得.
幾年前有一套電影叫《沉默》.
可能你有看過或聽過.
其實是說十七世紀.
日本一段傳教歷史.
它描寫當時嚴峻的迫迫.
不斷臨到信耶穌的信徒.
敵人用盡不同的方法.
威迫你也好.
用你的家人的生命要脅你也好.
嘗試動搖他們對耶穌的忠誠.
威迫他們放棄耶穌.

$^{281}$甚至你踩耶穌的聖像.
你去攝逐耶穌.
如果不單單是你.
你的家人也會有生命危險.
要求他們不認耶穌為主.
否則就死路一條.
與此同時.
當這些信徒面對著.
最極大的迫迫時.
他們嘗試禱告呼求上帝.
但上帝卻選擇沉默而對.
無論他們如何努力祈禱.
甚至質問上帝也好.
上帝也沒有反應.
你會感受到.
當中所產生的張力.
其實就是這部電影.
要探討的一個議題.
在這個情況下.
信徒應該堅持至死忠心.
守住信仰.
認耶穌為主.
還是選擇妥協放棄一事.
放棄耶穌.
保住性命.
不受苦.
讓姐妹當身處其中.
你就會感受到那種掙扎.
你知道很難做決定.
而耶穌在第35至38節提醒我們.
不如我們一起讀.
請大家讀.
程序表內35至9章一節最後段落.
預備起.
(主席在演講中).
你看到耶穌指示那些.
將會和祂一樣.
要面對迫不及待.
甚至生命威脅的門徒.
要教他們一條.

$^{321}$其實怎樣才能保住性命.
你看到耶穌的法則很弔詭.
那些將會面對迫不及待.
仍然堅持耶穌為主.
仍然堅持傳揚耶穌福音的門徒.
雖然他們很大機會.
會失掉肉體的生命.
但他們卻必得著一個真正的生命.
就是我們所清楚明白的永生.
進入上帝的國.
相反,人若選擇迴避.
不想為主捨己背十價的吩咐.
當他面對迫不及待時.
選擇不認耶穌.
放棄這個信仰的時候.
確實他可能會避過迫不及待而來的失喪生命.
但他卻會在永恆之中.
失去了和耶穌和上帝的關係.
失去了真正的生命.
所以你看到.
耶穌隨即警告門徒.
不要因為即將要來到的迫不及待.
就臨陣退縮.
放棄信仰.
將耶穌和祂的話看為可恥.
否則,耶穌這裡說清楚.
他日他榮耀再來審判傳遞的時候.
他都會棄絕那些不認祂的人.
而為了鼓勵門徒.
要竭力與耶穌受苦受辱受死的道路有份.
耶穌最後在九章第一節.
用一個的應許來結束整個的教導.
如果我們參考上文.
神的各大有能力臨到.
其實可以解為稍後耶穌從死裡復活.
這個神蹟.
意味著上帝是會站在耶穌那一方.
證明那些對耶穌誤解.
甚至反對殺害祂的人.
其實是錯的.

$^{361}$耶穌最終會得勝.
之前祂所受的榮辱.
施禪水所受的屈辱.
通通都會轉為榮耀.
所以門徒可以放膽履行他們的職份.
效法耶穌來為福音的緣故.
受苦甚至乎犧牲.
因為最終上帝同樣會透過復活.
為這群忠信的信徒平反.
事實上,弟兄姊妹.
經文的講解大概就來到這裡.
你會發現整段的經文.
耶穌只是由一個我們說的.
資訊性的問題.
問一下你,耶穌是誰?.
一條很簡單的問題.
到最後向門徒甚至向我們.
發出一個非常靠近的挑戰.
耶穌說,要麼就不要認祂為基督.
如果認祂為基督的話.
就代表你要甘心樂意為祂受苦.
甚至乎為祂寫記.
確實,要敬拜擁抱那位.
四處傳道,醫治趕鬼.
很成功很順利的耶穌.
普遍人都會嚮往,普遍人都樂意.
但是如果要誓死跟隨一位.
必定寫記,受苦甚至死亡的耶穌.
信仰就會立刻變得很沉重.
令很多人都想馬上和耶穌拉遠距離.
甚至乎劃清界線.
彼得的表現讓我們很明白.
其實只要知道耶穌是基督.
答對了這個答案.
不足以支撐你整個信仰人生.
因為當苦難,迫不及待的時候.
我們就會散了.
確認耶穌是基督.
其實意味著我們必須徹底順服跟隨祂.
既然耶穌說了.

$^{401}$若有人要跟隨祂.
就當寫記背起十字架來跟隨祂的時候.
其實我們不能當是耶穌說說而已.
不是認真的.
我們不能只是在生活中.
少修少補,稍為調整.
來配合耶穌.
不能這樣淡化了祂對我們的嚴格要求.
耶穌說的是要我們將整個人的生命交給祂.
《石上弟兄姊妹》這段經文迫使我們重新檢視.
我們跟隨主的心態到底是怎樣.
到底是我們緊貼耶穌的吩咐.
緊貼祂的步伐.
還是不知不覺想耶穌跟我們怎樣做.
確實跟隨主的過程中.
真的涉及很多取捨和掙扎.
甚至令人感到不舒服和痛苦.
因為要為上帝的緣故.
要捨棄那些可能是你最不想放手的事.
例如亞伯拉罕要獻祂的耳朵.
又或者要為主面對一些.
你最不想淋到你身上的事.
例如耶穌要喝十字架的苦杯.
事實上,上帝要我們為祂捨棄受苦的場景.
都不一樣.
有些人可能是事業.
如果你有看《三誠心所願》.
其實很好看.
我們訪問了一個叫Neal的弟兄.
他為了主的緣故.
他捨棄自己的移民計劃.
甚至捨棄他的職業.
他竟然留在香港經營tong 房.
搞社企.
為主幫助那些有住宿需要的人.
不單給他們一個瓦遮頭.
希望給他們一個家的溫暖.
代表上帝來關懷他們.
有時間可以上YouTube重溫.
亦有人為上帝捨棄受苦的場景.

$^{441}$可能是金錢取捨.
可能是工作上的考量.
亦有人為主的緣故.
可能很簡單.
少用手機.
少玩手機.
少上網.
少看劇集.
亦有人為主放下面子.
甚至專業.
我記得有一段經歷.
有一次平日坐巴士回教會工作.
一家三口出門.
上了巴士.
你知道早上很擠迫.
我老婆和兒子先上車.
他們入車後坐.
因為我排得較後.
我遲了出門.
到我上車時.
已經去到巴士的門口位.
上到巴士便很淡定.
打算開車.
但司機突然很嚴肅地站起來.
不可以過那條黃線.
你知道巴士收銀後.
有條黃線.
原則上不可以過那條黃線.
過了那條黃線.
有權不開車.
因為涉及安危.
他說不可以過那條黃線.
不然我不開車.
我們打算說說而已.
你看到多數巴士司機都很友善.
會盡量迫.
上到便上.
他很認真.
不可以.
這個下車.

$^{481}$下車後.
我當然死都不下.
然後他說.
這個黑色頭髮長頭髮.
穿黑色衣服的那個.
即是我.
下車.
我覺得很愚蠢.
被人指名道姓地趕下車.
我覺得很不開心.
但我沒理由和他爭拗.
沒理由拖累所有人.
來不及上班.
我很不開心.
很激氣.
好,我下車.
下車後我也很激動.
我覺得有沒有搞錯.
其他司機都不是這樣.
有些司機誇張到.
開後門給你上車.
為什麼你用不著這麼嚴厲.
又不是很悲觀.
過了黃線.
但我事後會想.
其實他有權這樣做.
這是他合理的權益.
他可以叫人下車.
他又不是真的錯.
我再想.
其實我為什麼這麼生氣.
我生氣什麼.
我生氣我自己上不到車.
不是因為他錯.
他沒有錯.
如果我連為主受一點氣.
都不行的時候.
更何況我經常說.
要為主受苦.
於是我求主幫助.

$^{521}$我坐另一架車.
重新調教自己的心態.
來上班.
故事未完.
幾個星期後.
我又去坐巴士.
今次我一個人.
那天是星期六.
我一上車.
又是你.
又是那個司機.
今次我當然怕被他趕下車.
不過今次少些人.
我立即站在黃線.
沒有事.
安全.
沒事.
如是者駛到下一個站.
但有很多人上車.
又出現同一個情況.
開始被人逼.
逼過了黃線.
去到門口.
司機又來了.
不知道為什麼觸發到他.
那條黃線是不可以逾越的.
又站起來.
下車.
如果不是不開車.
最後太太說.
趕上班.
已經很餓.
不想下車.
因為我也很感同身受.
我也不想他被人趕下車.
不如逼近一點.
結果逼近了.
逼進了黃線.
沒有事.
開車.

$^{561}$後面的乘客覺得.
做司機用不用這麼兇.
用不用這麼苦.
司機說.
你覺得我很想這麼兇嗎.
為了安全.
車輛在駛.
駛到某一個燈位.
停了車.
那時車輛很靜.
我聽到司機細聲說.
其實你覺得我想這麼兇嗎.
因為我在這裡開車.
擋風玻璃都爛了.
那一刻我就明白.
原來很多時候.
所有事情都有原因.
他之所以兇.
之所以這麼嚴格執行.
不可以逾越那條黃線.
是否因為他受過傷.
可能到現在.
他還未康復.
創傷還在.
所以他為了不想.
有意外.
所以他堅持.
不可以過這條黃線.
當這樣的時候.
一方面我開始不再生氣.
我反而覺得.
很值得去同情.
去體諒.
我更奇妙的是.
其實上帝很厲害.
上帝好像向我露一手.
提醒我.
其實你凡事不要看表面.
你覺得你很生氣的事.
你覺得是別人不對的事.

$^{601}$其實對方是有負重.
上帝給我上了一課.
凡事不要急於動氣.
不要急於批評論斷.
對你不好的人.
那些臨到你身上不好的事.
有時其實.
蘊含著上帝很多福氣.
在裡面.
上帝想你明白.
不要以自己為中心.
以他為中心.
他會安慰你.
保守你.
他會替你作主.
其實上帝已經指明.
既然今天這段經文提醒我們.
慰主,受苦,捨己.
是我們每一個信徒的必修科目的時候.
其實與其你迴避.
你抗拒.
你不想它出現.
倒不如你嘗試轉一個心態.
你嘗試去細味.
去學習.
上帝為你預備.
那些捨己,受苦.
甚至受氣的人生功課.
裡面一定有些東西.
上帝是想你學的.
否則不會臨到你身上.
他期望你完成這個課程.
期望你在裡面獲益.
當然在一個今天很講求成功,舒服.
不受得氣.
我要話事的世代.
是不容易的.
但耶穌這裡清楚說明.
他要重新調教我們作為門徒應有的態度.
他提醒我們要對他委身信服.

$^{641}$所以求主幫助我們.
審查我們生命裡面.
有沒有一些想法.
有沒有一些行為.
其實我們需要去為他寫去.
又或者當我們面對一些.
為主受苦的場景的時候.
求主也幫助我們.
讓我們有更多從他而來的愛和忍耐.
來承載這些苦難.
我們經歷到上帝的祝福和保守.
求主教我們堅持走他的道路.
主席.
\newpage



\section{馬可福音 9:14-29}
\label{sec:HtBaZ5KPIHs}
\textbf{2025年9月14日崇拜直播｜陳明泉牧師｜耶穌給門徒的十課－在人不能時......｜馬可福音9\_14-29}
\newline
\newline
連結: \href{https://youtube.com/watch?v=HtBaZ5KPIHs}{\texttt{https://youtube.com/watch?v=HtBaZ5KPIHs}} ~~~~ 語音日期: 2025-09-14
\newline
\newline
\hyperref[sec:xAjegbjPwhU]{\small{< < < PREV SERMON < < <}}
~
\hyperref[sec:index_chronic]{\small{[返順時目]}}
~
\hyperref[sec:LeA_EIgcwcw]{\small{> > > NEXT SERMON > > >}}
\newline
\newline
馬可福音 9:14-29
\newline
\begin{longtable}{cl}
\hline
\hline
章節 & 經文 (和合本修訂版)\\
\hline
9:14 & \begin{tabularx}{0.7\textwidth}{X} 他們到了門徒那裡，看見有一大群人圍著他們，又有文士和他們辯論。 \end{tabularx} \\ \\ \relax
9:15 & \begin{tabularx}{0.7\textwidth}{X} 眾人一見耶穌，都很驚奇，就跑上去向他問安。 \end{tabularx} \\ \\ \relax
9:16 & \begin{tabularx}{0.7\textwidth}{X} 耶穌問他們：「你們和他們辯論甚麼？」 \end{tabularx} \\ \\ \relax
9:17 & \begin{tabularx}{0.7\textwidth}{X} 眾人中的一個回答：「老師，我帶了我的兒子到你這裡來，他被啞巴的靈附著。 \end{tabularx} \\ \\ \relax
9:18 & \begin{tabularx}{0.7\textwidth}{X} 無論在哪裡，那靈拿住他，把他摔倒，他就口吐白沫，牙關緊鎖，身體僵硬。我請過你的門徒把那靈趕出去，他們卻不能。」 \end{tabularx} \\ \\ \relax
9:19 & \begin{tabularx}{0.7\textwidth}{X} 耶穌回答：「唉！這不信的世代啊，我和你們在一起要到幾時呢？我忍耐你們要到幾時呢？把他帶到我這裡！」 \end{tabularx} \\ \\ \relax
9:20 & \begin{tabularx}{0.7\textwidth}{X} 他們就帶了他來。那靈一見耶穌，就使他重重地抽風，倒在地上，翻來覆去，口吐白沫。 \end{tabularx} \\ \\ \relax
9:21 & \begin{tabularx}{0.7\textwidth}{X} 耶穌問他父親：「他得這病有多久了呢？」父親說：「從小的時候。 \end{tabularx} \\ \\ \relax
9:22 & \begin{tabularx}{0.7\textwidth}{X} 那靈屢次把他扔在火裡、水裡，要治死他。你若能做甚麼，求你憐憫我們，幫助我們。」 \end{tabularx} \\ \\ \relax
9:23 & \begin{tabularx}{0.7\textwidth}{X} 耶穌對他說：「『你若能』，在信的人，凡事都能。」 \end{tabularx} \\ \\ \relax
9:24 & \begin{tabularx}{0.7\textwidth}{X} 孩子的父親立刻喊著說：「我信；求你幫助我的不信！」 \end{tabularx} \\ \\ \relax
9:25 & \begin{tabularx}{0.7\textwidth}{X} 耶穌看見眾人都跑上來，就斥責那污靈說：「你這聾啞的靈，我命令你從他裡頭出來，再不要進去！」 \end{tabularx} \\ \\ \relax
9:26 & \begin{tabularx}{0.7\textwidth}{X} 那靈大喊一聲，使孩子猛烈地抽了一陣風，就出來了。孩子好像死了一般，以致眾人多半說：「他死了。」 \end{tabularx} \\ \\ \relax
9:27 & \begin{tabularx}{0.7\textwidth}{X} 但耶穌拉著他的手，扶他起來，他就站起來了。 \end{tabularx} \\ \\ \relax
9:28 & \begin{tabularx}{0.7\textwidth}{X} 耶穌進了屋子，門徒就私下問他：「我們為甚麼不能趕出那靈呢？」 \end{tabularx} \\ \\ \relax
9:29 & \begin{tabularx}{0.7\textwidth}{X} 耶穌對他們說：「非用禱告，這一類的邪靈總趕不出來。」 \end{tabularx} \\ \\
[1ex]
\hline
\hline
\end{longtable}
$^{1}$今日的經訓是在《馬可福音》九章十四至二十九節新約聖經六十頁.
《馬可福音》九章十四至二十九節新約聖經六十頁.
耶穌悼念門徒那裡,看見有許多人圍著他們,又有文士和他們辯論.
眾人一見耶穌都甚稀奇,就跑上去問他的安..
耶穌問他們說:「你們和他們辯論的是什麼?」.
眾人中間有一個人回答說:「夫子,我帶了我的兒子到你這裡來,他被啞巴鬼附著,.
無論在哪裡,鬼捉弄他,把他摔倒,他就口中流沫,咬牙切齒,身體枯乾,.
我請過你的門徒把鬼趕出去,他們卻是不能.」.
耶穌說:「不信的世代,我在你們這裡要到何時呢?.
我忍耐你們要到何時呢?把他帶到我這裡來吧.」.
他們就帶了他來,他一見耶穌,鬼便叫他重重的抽搐,倒在地上,翻來覆去,口中吐沫..
耶穌問他父親說:「他得這病有多少日子呢?」.
回答說:「從小的時候,鬼屢次把他陰在火裡,水裡,要滅他..
你若能作甚麼,求你憐憫我們,幫助我們.」.
耶穌對他說:「你若能信,在信的人凡事都能.」.
孩子的父親立時喊著說:「我信,但我信不足求主幫助.」.
耶穌看見眾人都跑上來,就斥責那烏鬼說:「你這聾啞的鬼,我吩咐你從他裡頭出來,再不要進去.」.
那鬼喊叫,使孩子大大的抽了一陣風,就出來了..
孩子好像死鳥一般,以至眾人多半說他是死鳥..
但耶穌拉著他的手扶他起來,他就站起來了..
耶穌進了屋子,門徒就暗暗的問他說:「我們為甚麼不能趕出他去呢?」.
耶穌說:「非用禱告,這一類的鬼總不能出來.」.
弟兄姊妹,平安!.
我們今天繼續來到二馬賀福音來看看耶穌對門徒的十個功課裡面,其中一課,能力感..
有一個小片段,你不要看我高高大大,我從小到大都是班裡站在最後的一個,最高的..
不過從小到大我都是手腳不協調的一個人,所以從小學開始上運動堂,基本上我就不太享受..
因為打球又比同學靈活,另外力又不夠別人的力..
再加上有時候走路都會左腳踢右腳,所以就是一個很高大,但是很笨拙的人..
直到中學的時候,進了一間學校,有一班男生很喜歡打球,看到他們打球很開心,很想去..
但是第一次上場就被人笑,吃球餅,心裡面的感受是很差的..
那一刻我曾經想過,以後我都不會再做運動了,不過我看到在籃球場上的人打得很開心,我會不會有一天像他們一樣呢?.
於是放學的時候,就向學校借了一個籃球,就走到一旁,自己在那裡射籃,在那裡練,苦練..
又不會走籃,就學走籃,放學之後回家,有一個膠的紙球,自己寫完字不要的,又在那裡扮打籃球..
然後走路的時候,好像走蟹一樣,在那裡籃球的步伐,弄著弄著,身體動作就開始協調..
到最後一半年之後,經過一段日子,又再跟同學開始比賽,比賽的時候,記不記得之前我是被人笑,吃球餅..
半年之後,在同學面前就拉著一個框打籃球,那些人就說你做什麼事?.
這些就是訓練,然後就很帥,之後就很喜歡打球,然後從中學開始,天天都想著打球,上課的時候就想起Michael Jordan,識飛,I believe I can fly..
我經常都覺得自己是Michael Jordan,買了他的T恤那些東西,不過在那一刻開始,我發覺一樣道理,到今天我都覺得很有趣..
原來,以前覺得自己不行的東西,只要我心裡面認定某一些目標,能力透過訓練是可以得到的..
人的生命很有趣,能力感很多人以為是天生的一種品質,或者人的一種天賦,上帝將能力賜給人,.

$^{41}$不過有更多,其實很多能力是透過後天栽培出來,或者人有某一種信念,人就朝著即使失敗都不能放棄的一種精神,直到做到某一些目標為止..
有些心理學家說到,人的能力感來自一個叫做growth mindset,一個叫做成長的思維,成長思維就是令人願意不怕失敗繼續練習..
另外相對的就是一個叫做固定思維,就是覺得所有能力都是天賦的,人沒有能力的,天賦你有多少就有多少,所以就採取一種比較被動和消極的態度..
如果我用這個角度來看,基督教由耶穌訓練12個門徒開始,這十幾個人成為全世界改變的宗教,這一班門徒一定會有一些相當強的能力感..
耶穌如何令這一班門徒成為一班可以改變世界的一班人呢?他們的能力感是如何來的呢?今天這一段經文大概我們會看到耶穌是這樣訓練門徒,希望我們今天讀這一段馬可福音,我們生命裡面的能力感都會強一些..
我們先祈禱求上帝幫助,天父上帝願你幫助我們,讓我們直著你的話語,我們更加認定上帝在我們生命裡面的塑造,願你在我們當中此時此刻你對我們內心說話,幫助孫講的講得清楚,幫助聆聽的每位弟兄姊妹,讓他們聽得明白在生命裡面實踐你的真理,.
讓他們能夠加強我們生命裡面實踐你的真理的鬥志和決心,願你在我們當中帶領,多謝主上帝聽我們祈禱,奉靠基督耶穌得勝的聖名,Amen..
我們一起看馬可福音,在馬可福音裡面,我這次是跳了,跟上一次的經文我跳了一部分,跳了的部分就是耶穌基督登山變像的部分,在馬可福音第九章第一節裡面,耶穌跟門徒說,他們十二個門徒當中,.
有一些人未嘗死未,就會看見神的國大有能力臨到,當然這個不只是說耶穌登山變像這個奇景,不過在馬可福音裡面,他即時記載著耶穌帶著雅各,彼得和約翰,就上到山上,然後他就登山變像,然後彼得看到這樣的奇景,就說為耶穌搭座鵬,.
到最後,這段經文裡面的重點就是,天上有聲音說,這是我的愛子,你們要聽他.門徒不見一人,只見耶穌,和他們在那裡..
耶穌在那裡做一個很特別的奇景,讓他們知道上帝的榮耀,初步預想一下,上帝的那種光芒,那種榮耀臨到在人世間裡面,耶穌的其中三個門徒就見到這個奇景..
因為這個經文,我之前有說過,所以我不再重複了,我想在這個奇景之後,耶穌和門徒下山的這個趕鬼的情景,我們今天要特別去研讀的..
剛才說過,上帝的能力大有能力臨到人世間裡面,見到一個奇景而著,但是一下山的時候,就見到那些門徒的不濟,甚至這一群剩下來的九個門徒,.
在地上繼續的來應應逆逆,去實行耶穌給他們的吩咐,趕鬼,治病,傳揚天國福音這一群門徒,在這裡他們是失敗的..
在第十四節裡面,九章十四節到十九節,經文裡面這樣說,耶穌到了門徒那裡,看見有許多人圍著他們,又有民事和他們辯論,.
眾人一見耶穌就甚稀奇,跑上去問他的案,耶穌問他們說,你們和他們辯論的是什麼?.
眾人中間有一個人回答說,夫子,我帶了我的兒子到你這裡來,他被啞巴鬼附著,無論在哪裡鬼捉弄他,把他摔倒,他就口中流沫,咬牙切齒,身體枯乾..
我請過你的門徒把鬼趕出去,他們卻是不能.在這裡裡面,你看見這個記載,講到耶穌和門徒下山的時候,人們和門徒在爭辯,爭辯什麼呢?.
耶穌問他們,當時其中一個人就說,那個爸爸帶著他的兒子來找門徒,兒子被鬼附,但是沒辦法趕出來,他求耶穌的門徒來趕鬼,但是那些門徒都沒辦法把這個污鬼趕出來..
值得留意有幾個地方,那些民事和門徒辯論的內容會辯論什麼呢?大概就是之前的門徒走遍各城各鄉,被耶穌差遣到不同的地方,他們做的工作是成功的..
他們能夠做到神的國臨在的那種醫治,以及人生命得到改變的工作,到了這一刻,人們失效.如果比喻我,如果對耶穌的團隊心裡面有敵意,我一定會追著他們,.
就會說,你們這群人,又說得這麼英明神武,到了這一刻,你們什麼都做不到,那個人被鬼附,你們趕都趕不出來,你奉耶穌的名趕鬼,不是吧?如果是可以的,這個人就不用受這麼多苦了..
爭論的論點,那些民事或者法利賽人一定會就著這個失敗經歷,跟那些門徒基本上就揶揄他們,同時間,如果作為門徒的,因為他經歷過神國大能的那種權柄,.
而且,之前在馬可福音6章裡面說到,其實他們是成功的,6章7至13節,我節錄一些部分,告訴大家聽..
耶穌叫了十二個門徒來,差遣他們兩個兩個出去,也賜給他們權柄,制服污鬼.十二節,門徒就出去傳道,叫人悔改,又趕出許多的鬼,用油抹了許多的病人,治好他們..
所以在這裡,你看到那些門徒是得著耶穌的權柄,他們是有成功的經歷的,他們是做到事情的,耶穌又沒有說過要收回這個權柄,究竟有什麼出錯令他們做不到呢?.
那些門徒之前就是這樣做,到這一刻就是這樣做,但是趕都趕不到,那些門徒就要去防衛,不是耶穌的名字不行,我們當然有些人不行,你可以想像第一個不行,第二個,第三個,第四個,一直試過,每一個都不行,是沒有辦法做得到..
在經文裡面多次出現一個很重要的字眼,就是能夠與不能夠,在第十三節,九章,sorry,九章第十九節裡面,這裡特別這樣說,我請過你的門徒把鬼趕出去,他們卻是不能..
不能這個字,後來會再出現,不過在這裡,門徒似乎被耶穌差派,面對著這個孩子被鬼附的現實,他們無能為力,他們做不到任何事情來改變這個受苦的孩子..
當耶穌見到這個孩子,或者見到這個人,將他的苦況告訴耶穌的時候,然後由第十九節開始,就繼續這樣記載,我們很少見耶穌這樣說話,第十九節裡面這樣說,耶穌說,不信的世代,我載你們這類要到幾時呢?我引來你們要到幾時呢?.
把他帶到我這裡來吧.值得留意的就是,耶穌的醫病,你在《馬學福音》裡面,你一直見過的,耶穌對於那些病人,或者那些被鬼附的人,通常他都是一個很謙卑,很主動去觸摸那些很困難的人..
但是當跟門徒下山,一聽到這些被鬼附的爸爸在說這個苦況的時候,耶穌第一句某程度上叫做發悔氣,特別是在說那個不信的世代,耶穌在說誰呢?.
如果這個爸爸走過來說他不信,其實之前都沒有見過耶穌,或者沒有經歷過耶穌,你說他不信,這樣也是的,不過似乎他的語氣又重了一點,還有耶穌說我引來到你幾時呢?.
你們要在我們當中,我在你們當中要多久呢?這些是耶穌某程度上是發悔氣的話,如果對於一個病患者來找耶穌說發悔氣的話,似乎順民下理又不是很通,但是那些門徒,無端端又為什麼會被耶穌說呢?我們都不知道..
接著那些法利賽人,法利賽人通常如果耶穌要說他,他就直接說,通常都不是說他們不信的世代,通常說法利賽人或者文石是什麼?粉色的墳墓,假冒為善的人,但是在這裡,耶穌好像在發悔氣,如果身邊的人看到耶穌,真的有一點莫名其妙,.
不過我們看到這裡耶穌有一些情緒,而那種情緒帶著一種覺得有一種慨嘆,在一個這樣的狀況下,人的不信,令到好像很多的問題變得越來越複雜,.
不過當耶穌說完這句話之後,那個能與不能,和那個信心就連成了一個重要的轉折位,由人的不能,到耶穌去看是一個不信的世代..
我們繼續看第20節,20節裡面就繼續記載,他們就帶了他來,他一見耶穌,鬼便叫他重重地抽風,倒在地上,翻來覆去,口中流沫..
然後,耶穌問他父親說,他得這病有多少日子呢?回答說,從小的時候,鬼屢次把他泯在火裡,水裡,要滅他.你若能作甚麼,求你憐憫我們,幫助我們..
23節,耶穌對他說,你若能信,再信的人凡事都能.孩子的父親立時喊著說,我信,但我信不足,求主幫助..

$^{81}$耶穌看見眾人都跑上來,就斥責那污鬼說,你這聾啞的鬼,我吩咐你從他裡頭出來,再不要進去..
鬼喊叫,使孩子大大地抽了一陣風,就出來了,孩子好像死了一般,以致眾人多半說他是死了,但耶穌拉著他的手扶他起來,他就站立起來了..
在第19至27節就記載著耶穌把這個挺惡的聾啞鬼從孩子身上趕出來,經文裡面有些重要的記述.
首先說到這個孩子被聾啞的鬼依附著,這隻鬼真的挺惡的,當眾人把這個孩子和這個父親帶到耶穌面前的時候,這個孩子的狀況是差了的..
經文裡面特別說到,這個孩子一見到耶穌就重重地抽風,倒在地上,翻來覆去,口中流沫..
我希望大家不用特別想象那個情景,因為你會害怕的..
不過從字面你會看到,這個孩子去到耶穌面前,他的狀況比起之前更加差,更恐怖了..
所以在這裡,這個鬼某程度上是挺惡的,把這個孩子合制住,也是一個不容易處理的污靈..
然後耶穌問這個孩子病了什麼時候,或者被鬼附,或者有什麼狀況,有多久,這個爸爸也說不出,不過就是很小的時候..
這個鬼還要多次把他的兒子扔在火裡面,泯在水裡面,好像要危害他的性命一樣..
可想而知這個孩子真的被這個污靈糟蹋得遍體鱗傷,很痛苦的..
然後說完這一句之後,和合本的翻譯,我們值得留意的,或者在這裡記載這個爸爸所說的話,我們要讀這段經文的重點..
這個爸爸就跟耶穌說,你若能作什麼,求你憐憫我,幫助我..
面對著自己的兒子經常被這個很惡的污靈糟蹋,大概這個爸爸試過很多的方法,試過每一個方法都是沒有辦法幫助自己的兒子..
大量的失敗,大量的傷心,大量的心痛,你可以想像,從小到大,很多年了,兒子都是這樣,搞不定..
而且沒有一件事可以有任何的進展,甚至最近況的就是找到耶穌的門徒,都搞不定..
現在看到耶穌,那個情況更差了,更危急了..
所以作為這個爸爸,其實心底裡絕對是有一種完全的那種的沮喪和那種的痛心..
所以他是很痛苦地帶著一份剩餘的那份的寄望,耶穌,如果你能夠的,求你憐憫我,幫助我們..
他心裡面因為試過不同的方法,也都嘗試過所有失敗的方法,任何東西都是看到所有東西都是不行的..
所以心底裡他看耶穌都是其中一個方法,所以他跟耶穌說,如果你能做到什麼,你就幫幫我..
你若能,是在說,他是帶著一份不肯定的心,孤注一擲,甚至試試,試試不妨,反應都是這樣,找到一個人試試又何妨呢?.
不過心底裡,他的寄望,或者他心裡面,他未必是帶著一份確信可以直著耶穌,或者耶穌的權能可以醫到他自己的兒子..
所以他的說話那句就是在說,你若能,當耶穌聽到這個父親在說,你若能的時候,耶穌怎樣回應呢?.
和合本的翻譯就是,你若能信,在信的人凡事都能..
你若能信的信字在和合本,剛才我們讀的聖經,那個信字,你若能信的信字是和合本加進去的..
和合新譯本他都已經調教了,在那裡他就寫,耶穌對他說,你若能,在信的人凡事都能..
耶穌的回應很有趣,耶穌的回應,他特別回應這個沮喪的父親,在說,你若能幫到我們,你就憐憫我們..
就是說,耶穌是一個可能性,或者,如果你是可以的,你就幫幫我吧,你能夠做到什麼你就做吧..
不過耶穌是針對那個父親說,你若能這三個字,所以耶穌是重新說,你若能那句句語,然後就是回應,在信的人凡事都能..
有些英文譯本再加上的就是,你若能,問好,感嘆好,就是說,這個父親在說你若能的時候,你若能,在信的人凡事都能..
那個語氣就是這樣,在耶穌的權能裡面,不是耶穌若可以,是耶穌決定了,耶穌說可以就可以,人的回應就是信靠,信心..
你信耶穌能夠幫到你,在任何的困難當中,都能夠改變你..
二十四節,這個父親如何回應呢?和合本也是和原文有差異的..
二十四節,我們先讀和合本,孩子的父親立時喊著說,我信,但我信不足,求主幫助..
和合本的翻譯就是說,爸爸的信心不足,我信,但我信心不足,所以求主幫助..
他的意思是這樣的,和合本用這樣的角度來翻譯..
不過,和合新譯本已經更加貼近原文了.和合新譯本,和合修訂本,在二十四節裡面的翻譯就是,孩子的父親立刻喊著說,.
我信,求你幫助我的不信.原文裡面,他不是在說我信不足,原文是在說我的不信..
簡單來說,這個爸爸在說你若能信,其實他心裡面,其實他不是很信耶穌,他只是一個拼命,.

$^{121}$用不同的方法,一個沮喪的爸爸,希望幫助他的兒子,可以脫離被鬼壓制的苦況,他什麼都做過了..
耶穌只不過是在很多方法裡面的其中一個,他心裡面也有太多失敗的例子,不敢有任何寄予一些的寄望,可以成功..
他只知道,如果你做了多少,就幫多少,求你幫幫我,憐憫我..
但是耶穌的重點就是,你信,信耶穌,信靠他的時候,什麼都可以改變..
我信,但求主你幫我的不信,這個父親直認他自己心裡面其實不信上帝..
在這裡我們稍微要思想一下,在馬可福音寄述的段落裡面,或者剛才的對答裡面,.
他講了一個很簡單的邏輯,人的信,再信的人就可以凡事都能,信的人凡事都能,信和能力是掛勾的..
因為信上帝,所以就有這一份能力,在這裡裡面,整個的局面就是那些人的不能,那些人不能夠做到任何事情..
這些人的不能的原因是什麼?不信,信就凡事都能,不能就基於他們的不信..
所以整個馬可福音,特別是這一段裡面講到不信的世代,不單只是那個爸爸,不單只是那些法利賽人,.
甚至乎那些門徒,或者整個所有的人,他們都是陷在一種不信的光景裡面..
你就會說牧師,他們跟耶穌,不信?這班門徒,還有這些人如果不信,就不會來求耶穌..
在馬可裡面,對於信心的定義,其實都很mean的,信不是說這麼多事情試一下,而是毫無懷疑地全身信靠下去的..
有半點懷疑的那個都不是信,所以在馬可的分辨裡面,只有不信和信,就不會說我信一點,不信一點..
你這樣說好像很複雜,不複雜,我們平時信耶穌,當你信了主的時候,你舊有的偶像,舊有你所倚靠的東西,就要全部丟掉..
你信靠的對象只有一個,唯一一個.舊時你很多裡面信的神明,你通通都要否定它,不要它,你只是將你的信心放在一個對象當中..
有些頂子妹曾經跟我說,她未信主之前,什麼都信的,她說基督教我也信一點,甚至佛教也信,泰國的佛教也信,總之有拜錯也沒有放過,所有事情多些神明保佑自己,你說多好呢?.
所以買了很多舊時付錢買的佛牌,我說只是金牌,多少錢?幾萬元的,牧師,我說你信主,你要決定了,這些東西要丟掉,丟掉它..
那些頂子妹很掙扎的,不過到最後,如果是信靠那位舊屬你的舊主,這些東西通通都要丟掉..
那位提案說,我以前進旺針,有些頂子妹有些佛牌,之前的牧者,我剛剛進來,不懂得訂,任勞任怨,那位牧者說這些佛牌很難打,那面的牙膠很硬,丟垃圾桶又怕有人撿,.
那位牧師說Anders你處理了它吧,我說那怎麼處理?那位牧師說難不成你也會進長洲的,那我不是丟去長洲,那位牧師說不是,那我丟去哪?.
那位牧師說你在長洲坐船,你坐下層,然後把佛牌丟掉,我說牧師亂拋垃圾,佛千五元,還要丟東西進海裡,那位牧師說,安靜點,處理好了..
我心想當我做這件事的時候,好像在看Titanic,那個海洋之心丟進海裡,到最後我不是散流的燈,我說牧師不如你陪我一次進長洲,你丟吧,我不是這麼容易就範,我不是亂拋垃圾的人,.
不過那個信心只有信靠一個,沒有說信一點不信一點,耶穌不是這麼多方法裡面的其中一種,耶穌是成為了唯一的方法,.
如果你信靠他的時候,就不是你弱能,而是帶著毫不懷疑的信心,全然的去信靠,在這裡馬可福音說到一個這樣的信心,.
人的不能在於不信,而信心真正的價值就是完完全全的信靠,這個爸爸他自己知道自己生命的光景,他求耶穌幫他,他的不信在這一刻他要被更新,被改變,.
他帶著這一份信心,求耶穌讓他的兒子得到醫治,經文裡面就繼續的來到記載,其實都是繼續打硬仗的,.
然後耶穌就斥責那隻污鬼,聾啞的鬼,我吩咐你從他裡面裡頭出來,再不要進去,那隻鬼就喊叫,洗鞋子大大地抽了一陣風,.
就出來,孩子好像死了一樣,以至眾人都多半說他是死了,但耶穌拉著他的手扶他起來,他就站起來,.
在整件事裡面,這麼惡的污靈終於因著這個爸爸的決意相信,我信,信耶穌,這隻污鬼就從這個孩子趕走了出來,.
整件事,這個趕完鬼的故事就完了,然後馬可福音,因為在說門徒訓練,門徒就再問耶穌,耶穌說整件事是怎樣的,要debrief,.
耶穌也藉著這個趕鬼的事蹟,重新來幫門徒上一課,究竟這班門徒跟耶穌,除了見證著這個神蹟之外,.
對門徒自己的生命有什麼重要呢?28,29節就成為了這個段落的結語,耶穌進了屋子,門徒就暗暗地問他,.
我們為什麼不能趕出他去呢?再出現不能的字,門徒說,為什麼我們趕不出他去呢?.
耶穌說,非用禱告,這一類的鬼總不能出來,這樣就完結了,很多時候有一些前輩就會特別來引用,.
說非用禱告,有一些和合本翻譯,就是有古卷雲,加上禁食兩個字,可能大家的聖經你會看到的,.
加上禁食,所以門徒問為什麼我趕不出這個鬼,如果加上禁食兩個字,就說因為不是禁食禱告,所以趕鬼的時候就要禁食禱告,.
用這個方式來趕鬼,不禁食就沒辦法趕出這個鬼,有些前輩或者有些聖經家用某一些古抄本,.
用禁食兩個字就覺得門徒是因為沒有做禁食祈禱,而趕不到鬼,不過你這樣說,順民下理是不通的,.
再加上禁食兩個字其實不是一些可靠的古抄本,基本上最重要的抄本都是沒有禁食的字,耶穌只是說非用禱告,.

$^{161}$這一類的鬼,甚至乎鬼字都沒有寫,這一類總不能出來,不過如果你說出來都是鬼,所以和合本這裡翻譯是貼近原文的,.
但是耶穌給這個勸勉對門徒有什麼意思呢?記不記得剛才說到那個句式,再信的人凡事都能,.
而這些門徒的不能是基於什麼呢?是基於不信,而這個禱告就正正反映了這班門徒,他們是不是真的依靠神去做這個趕鬼的舉動,.
如果在六章裡面記載著門徒一直都是成功,奉耶穌的名就趕鬼,他們得著權柄,他們覺得自己擁有了,.
所以他們看到所有事情都很技術性地繼續做,之前成功的就繼續跟隨,一直做,都是沒有什麼大問題的,.
遇到這一類這麼惡的污鬼,他用技巧是處理不到的,他單單奉耶穌的名,他都沒有辦法處理到這麼惡的污靈,.
唯有藉著禱告來依靠上帝的能力,所以在這裡面的禱告,他不是說那個方法,是不是用禁食,或者某一種最虔誠的方式來祈禱,.
那個禱告的重點是有沒有信靠上帝,這個對於我們今天都很重要,很多以前我們覺得可行的方法,我們覺得都是這樣做,順著去的,.
但是如果我們生命裡面只是順著某一些技巧,順著我們的經驗,失去了對上帝的信靠,我們會經歷我們生命的不能,.
或者生命的不能更加提醒我們,其實我們是不是過著一個信靠上帝的生活,.
門徒在這裡的不能,耶穌用禱告來提醒他,告訴他你不是用你的技巧,或者用你自己一直以來的經驗,你就可以改變,.
如果我是門徒,我就會和夫子來辯論,夫子不是,我們都信你,我都是奉你的名,我不信你,怎奉你的名,是不是?.
信心,如果回到馬可福音最嚴格的定義,是唯一的道路,不是其中一個道路,雅各書說得很絕,.
他說你信神只有一位,你信的不錯,鬼魔也信,卻是賤精,簡單說什麼呢?.
你說你信神,魔鬼也信,什麼都說是信,不過真正有生命的影響力,不是像鬼魔式的信,而是一個毫不懷疑的完全信靠,.
如果你信,你就撇除所有其他的懷疑,專心一致,仰望上帝,如果這班門徒沒有進行任何禱告,他們就落在一種不信的處境,.
他們就沒有辦法實踐神國彰顯的能力,這段經文其實不是很複雜,不過當我們把焦點放回起,耶穌要強調的不信,其實是一個很嚴格的定義,.
今天我們是一個知識充滿的世代,我不是反智的,讀知識是好的,不過如果我們今天的信仰觀念純粹是由我的經驗,我的知識來經營,我跟一個未信者有什麼分別呢?.
我跟一個未信耶穌的人,我都是用那種方法來待人接物,行事為人,我跟一個不信的人有什麼分別呢?.
這段經文提醒我們,我們今天會不會說口中是信主,但過著一個不信者的生活,一個沒有神的基督徒,我們只是生活中形形異異,用世俗的智慧,用覺得有用的方法來過我們自己的生活,.
說一個負面的例子,這是我家人出現的,我一個家裡的親戚,他為了自己的子女入學,我以前是知道他的,拜到滿天神佛,都挺迷信的,.
沒多久後,他跟我說,阿全,我洗禮了,你洗禮?這麼棒?你的得救經驗是怎樣的?他說的得救經驗就是在教會裡面聽牧師講的,就是武道班的內容,.
牧師講一次的內容,我都知道,你的得救經驗是什麼?他都是講別人講的內容,當作是自己的東西,我有點不客氣,他說,那你為什麼要浸禮?你不知道嗎?有五分的,加分的,.
面試的時候,你怎樣跟校長說?我告訴他,我相信,他說,你相信?我覺得你信心有點問題,我這樣有點mean,不過今天這個世代,很多人都說,我相信的,因為五分,.
有一些前輩告訴我,進去監獄探監,他說很多囚友,他自己都是相信的,因為可以快點出來,有些真心相信的,.
但也有些可能當基督教只是其中一個方法,拿到他心目中的好處,弟兄姊妹,今天你所信的是什麼?拿到好處?拿到福利?拿到那五分?.
因為我知道,我生命只有這一位主,值得我完全信靠,只有一位耶穌能夠拯救我們的生命,我生命裡面,如果沒有他成為我們的唯一,我們的信都是枉然,.
耶穌給門徒的功課就是,門徒就要有這一份信靠,如果他沒有這一份單純決心的那種相信,他們就失去能力,他們也沒有改變世界的那種心靈,.
因為心懷熱意的人,生命是軟弱疲乏,但如果我們心靈裡面謹守我們所信,我們凡事都能做,願主繼續對我們說話,幫助我們..
(影片完結).
\newpage



\section{約翰福音 9:8-41}
\label{sec:LeA_EIgcwcw}
\textbf{2025年9月21日崇拜直播｜梁家麟牧師｜瞎子信耶穌的故事｜約翰福音9\_8-41}
\newline
\newline
連結: \href{https://youtube.com/watch?v=LeA_EIgcwcw}{\texttt{https://youtube.com/watch?v=LeA\_EIgcwcw}} ~~~~ 語音日期: 2025-09-23
\newline
\newline
\hyperref[sec:HtBaZ5KPIHs]{\small{< < < PREV SERMON < < <}}
~
\hyperref[sec:index_chronic]{\small{[返順時目]}}
~
\hyperref[sec:JJ124JmNAMw]{\small{> > > NEXT SERMON > > >}}
\newline
\newline
約翰福音 9:8-41
\newline
\begin{longtable}{cl}
\hline
\hline
章節 & 經文 (和合本修訂版)\\
\hline
9:8 & \begin{tabularx}{0.7\textwidth}{X} 他的鄰舍和素常見他討飯的人，就說：「這不是那從前坐著討飯的人嗎？」 \end{tabularx} \\ \\ \relax
9:9 & \begin{tabularx}{0.7\textwidth}{X} 有的說：「是他」；又有的說：「不是，卻是像他。」他自己說：「是我。」 \end{tabularx} \\ \\ \relax
9:10 & \begin{tabularx}{0.7\textwidth}{X} 於是他們對他說：「你的眼睛是怎麼開的呢？」 \end{tabularx} \\ \\ \relax
9:11 & \begin{tabularx}{0.7\textwidth}{X} 那人回答：「有一個名叫耶穌的，他和了泥抹我的眼睛，對我說：『你到西羅亞池子去洗。』我去一洗，就看見了。」 \end{tabularx} \\ \\ \relax
9:12 & \begin{tabularx}{0.7\textwidth}{X} 他們對他說：「那個人在哪裡？」他說：「我不知道。」 \end{tabularx} \\ \\ \relax
9:13 & \begin{tabularx}{0.7\textwidth}{X} 他們把以前失明的那個人帶到法利賽人那裡。 \end{tabularx} \\ \\ \relax
9:14 & \begin{tabularx}{0.7\textwidth}{X} 耶穌和泥開他眼睛的那一天是安息日。 \end{tabularx} \\ \\ \relax
9:15 & \begin{tabularx}{0.7\textwidth}{X} 法利賽人又問他是怎麼得看見的。他對他們說：「他把泥抹在我的眼睛上，我一洗，就看見了。」 \end{tabularx} \\ \\ \relax
9:16 & \begin{tabularx}{0.7\textwidth}{X} 於是法利賽人中有的說：「這個人不是從神來的，因為他不守安息日。」另有的說：「一個罪人怎能行這樣的神蹟呢？」他們之間就產生了分裂。 \end{tabularx} \\ \\ \relax
9:17 & \begin{tabularx}{0.7\textwidth}{X} 於是他們又對那盲人說：「他開了你的眼睛，你說他是怎樣的人呢？」他說：「他是個先知。」 \end{tabularx} \\ \\ \relax
9:18 & \begin{tabularx}{0.7\textwidth}{X} 猶太人不信他以前是失明，後來能看見的，等到叫了他的父母來， \end{tabularx} \\ \\ \relax
9:19 & \begin{tabularx}{0.7\textwidth}{X} 問他們說：「這是你們的兒子嗎？你們說他生來是失明的，現在怎麼看見了呢？」 \end{tabularx} \\ \\ \relax
9:20 & \begin{tabularx}{0.7\textwidth}{X} 他的父母就回答說：「他是我們的兒子，生來就失明，這是我們知道的。 \end{tabularx} \\ \\ \relax
9:21 & \begin{tabularx}{0.7\textwidth}{X} 至於他現在怎麼能看見，我們卻不知道；是誰開了他的眼睛，我們也不知道。他已經是成人，你們問他吧，他自己會說。」 \end{tabularx} \\ \\ \relax
9:22 & \begin{tabularx}{0.7\textwidth}{X} 他父母說這些話，是怕猶太人，因為猶太人已經商定，若有宣認耶穌是基督的，要把他趕出會堂。 \end{tabularx} \\ \\ \relax
9:23 & \begin{tabularx}{0.7\textwidth}{X} 因此他父母說「他已經是成人，你們問他吧」。 \end{tabularx} \\ \\ \relax
9:24 & \begin{tabularx}{0.7\textwidth}{X} 於是法利賽人第二次叫了那以前失明的人來，對他說：「你要將榮耀歸給神，我們知道這人是個罪人。」 \end{tabularx} \\ \\ \relax
9:25 & \begin{tabularx}{0.7\textwidth}{X} 那人就回答：「他是不是個罪人，我不知道；有一件事我知道，我本來是失明的，現在我看見了。」 \end{tabularx} \\ \\ \relax
9:26 & \begin{tabularx}{0.7\textwidth}{X} 他們就問他：「他給你做了甚麼？是怎麼開了你的眼睛？」 \end{tabularx} \\ \\ \relax
9:27 & \begin{tabularx}{0.7\textwidth}{X} 他回答他們：「我已經告訴你們了，你們不聽，為甚麼又要聽呢？難道你們也要作他的門徒嗎？」 \end{tabularx} \\ \\ \relax
9:28 & \begin{tabularx}{0.7\textwidth}{X} 他們就罵他：「你是他的門徒，而我們是摩西的門徒。 \end{tabularx} \\ \\ \relax
9:29 & \begin{tabularx}{0.7\textwidth}{X} 神對摩西說話是我們知道的，可是這個人，我們不知道他從哪裡來。」 \end{tabularx} \\ \\ \relax
9:30 & \begin{tabularx}{0.7\textwidth}{X} 那人回答，對他們說：「他開了我的眼睛，你們竟不知道他從哪裡來，這真是奇怪！ \end{tabularx} \\ \\ \relax
9:31 & \begin{tabularx}{0.7\textwidth}{X} 我們知道神不聽罪人，惟有敬奉神、遵行他旨意的，神才聽他。 \end{tabularx} \\ \\ \relax
9:32 & \begin{tabularx}{0.7\textwidth}{X} 從創世以來，未曾聽見有人開了生來就失明的人的眼睛。 \end{tabularx} \\ \\ \relax
9:33 & \begin{tabularx}{0.7\textwidth}{X} 這人若不是從神來的，甚麼也不能做。」 \end{tabularx} \\ \\ \relax
9:34 & \begin{tabularx}{0.7\textwidth}{X} 他們回答他說：「你完全是生在罪中的，還要來教訓我們嗎？」於是他們把他趕出去了。 \end{tabularx} \\ \\ \relax
9:35 & \begin{tabularx}{0.7\textwidth}{X} 耶穌聽說他們把他趕出去，就找到他，說：「你信人子嗎？」 \end{tabularx} \\ \\ \relax
9:36 & \begin{tabularx}{0.7\textwidth}{X} 那人回答說：「主啊，人子是誰？告訴我，好讓我信他。」 \end{tabularx} \\ \\ \relax
9:37 & \begin{tabularx}{0.7\textwidth}{X} 耶穌對他說：「你已經看見他，現在和你說話的就是他。」 \end{tabularx} \\ \\ \relax
9:38 & \begin{tabularx}{0.7\textwidth}{X} 他說：「主啊，我信！」他就拜耶穌。 \end{tabularx} \\ \\ \relax
9:39 & \begin{tabularx}{0.7\textwidth}{X} 耶穌說：「我為審判到這世上來，使不能看見的看見，能看見的反而失明。」 \end{tabularx} \\ \\ \relax
9:40 & \begin{tabularx}{0.7\textwidth}{X} 同他在那裡的法利賽人聽見這些話，就對他說：「難道我們也失明了嗎？」 \end{tabularx} \\ \\ \relax
9:41 & \begin{tabularx}{0.7\textwidth}{X} 耶穌對他們說：「你們若是失明的，就沒有罪了；但現在你們說『我們能看見』，你們的罪還在。」 \end{tabularx} \\ \\
[1ex]
\hline
\hline
\end{longtable}
$^{1}$以下是經訓時間.
今天我們會選讀新約約翰福音第九章十三至三十四節.
請聽我讀出.
他們把從前黑眼的人帶到法利賽人那裡.
耶穌和來開他眼睛的日子是安息日.
法利賽人也問他是怎麼得看見的.
黑子對他們說:他的泥沒在我的眼睛上.
我去一洗就看見了.
法利賽人終有的說:這個人不是從神來的.
因為他不守安息日.
又有人說:一個罪人怎能行這樣的神蹟呢.
他們就起了紛爭.
他們又對黑子說:你既然開了你的眼睛.
你說他是怎樣的人呢.
他說:事過先知.
猶太人不信他從前是黑眼.
後來能看見的.
等到叫了他的父母來問他們說.
這是你們的兒子嗎.
你們說他身內是黑眼的.
如今怎麼能看見呢.
他父母回答說.
他是我們的兒子.
身內就黑眼.
這是我們知道的.
至於他如今怎麼能看見.
我們卻不知道.
是誰開了他的眼睛.
我們也不知道.
他已經成了人.
你們問他吧.
他自己必能說.
他父母說者說是怕猶太人.
已經商議定了.
若有認耶穌基督的.
要把他趕出會堂.
因此他父母說.
他已經成了人.
你們問他吧.
所以法利賽人第二次叫了那從前黑眼的人來.

$^{41}$對他說:你該將榮耀歸給神.
我們知道這人是個罪人.
他說:他是個罪人不是我不知道.
有一件事我知道.
從前我是黑眼的.
如今能看見了.
他們問他說.
他向你作什麼.
是怎麼開了你的眼睛呢.
他回答說.
我剛才告訴你們.
你們不聽.
為什麼又要聽呢.
莫非你們也要作他的門徒嗎.
他們就罵他說.
你是他的門徒.
我們是摩西的門徒.
神對摩西說話是我們知道的.
只是這個人.
我們不知道他從哪裡來.
乃人回答說.
他開了我的眼睛.
你們竟不知道他從哪裡來.
這真是奇怪.
我們知道神不聽罪人.
唯有敬奉神.
遵行他旨意的.
神才聽他.
從創世以來.
未曾聽見有人把.
身內是黑眼的眼睛開了.
這人若不是從神來的.
什麼也不能作.
他們回答說.
你傳言身在罪孽中.
還要教訓我們嗎.
於是把他趕出去了.
領事會早晨.
今天我們會看.
差不多整張聖經.

$^{81}$讀經我們讀不完.
我們看的是八節到四十一節.
除了上一次.
一到七節之外.
全部看完.
耶穌醫治天生黑眼的人.
這個故事.
上一次我們已經說了.
不過還未說完.
應該這樣說.
叫人舉手決志.
不是很容易.
不過不是最難.
決志了之後.
他就永遠快樂地生活下去.
我們知道成為基督徒.
故事才是更長篇.
所以對這個核子來說.
他被醫治.
那一刻很快就做了.
但醫治之後.
是一個什麼故事呢?.
今天我們會看.
因為整個故事很長.
很複雜.
所以我們套用舞台劇的佈局.
將整個故事分成五幕劇.
第一幕是關連到核子.
雖然他已經醫治好.
但因為他沒有名字.
所以我們仍然叫他做核子.
雖然他曾經是核子.
不是現在還是核子.
不過我們叫他做核子.
他的鄰居和那些素常看到他行乞的人.
他們對這個核子.
突然之間能夠看見.
非常不適應.
有些人變化得這麼激烈.
是很難受的.

$^{121}$他們甚至認不出這個核子.
核眼在當時是絕症.
沒有人相信能夠醫治得好.
所以他們沒有心理準備.
去面對一個覆命的核子.
他們追問核子.
你的眼睛是怎樣覆命的呢?.
那核子是這樣說.
在十一節.
他有一個人名叫耶穌.
他和泥抹我的眼睛對我說.
你往西羅亞池子去洗.
我去洗就看見了.
核子自己都沒有想過能夠覆命.
他很被動地去接受耶穌的醫治.
這個醫生不是他主動找的.
他連耶穌是誰都不知道.
也不知道他從哪裡來.
往哪裡去.
他只是傻乎乎地.
遵從耶穌的吩咐.
去西羅亞池子一洗.
這樣做.
他唯一知道的就是.
他曾經經過一個怎樣醫治的步驟.
鄰居碰上一宗千古奇聞.
他也不知道應該怎樣去處理.
處理主要是心理上.
於是他就將核子帶到法利賽人那裡.
法利賽人.
我們常常都當他們是奸的角色.
不過他們在當時期.
是被視為最有學問的人.
希望這班法利賽人能夠解說這件事.
這是第一幕.
然後到第二幕.
核子被帶到一個法利賽人的群體.
這個有可能是猶太會堂裡.
一班宗教領袖所結集成的一個小組.
或者叫做委員會.

$^{161}$處理跟律法有關的宗教事務.
因為我們沒有聽過耶路撒冷的聖殿.
或者公會有將猶太人開除會籍.
全球封殺的做法.
我們沒有聽過.
所以我們大概不相信.
是去到耶路撒冷的公會的層次.
大概可能是一個地區性的會堂.
有一班宗教領袖.
就像我們一個堂會的執事會.
他們自覺有權將人趕出會堂.
切斷這個人跟他們的群體.
和宗教上的任何的聯繫.
法利賽人首先仍然是查問.
是值得傳譽的過程.
不過這個過程是提供不到.
破解謎團的蛛絲馬跡.
口水,泥,西羅亞廁紙.
這些元素都跟醫治毫不相干.
不信的話,你有病.
你試試把這幾樣東西加起來.
你驗滅了,是沒有關係的.
所以算來算去只能推定它是一個神蹟.
沒有其他東西可以解釋.
天生核眼的復明.
這是一個不折不扣的神蹟.
這個不可能不是一個神蹟.
但是,核子的華卻讓他們確認.
醫治事件是在安息日發生的.
施行神蹟的耶穌肯定是抵觸了.
安息日不可工作的規定.
好像我上一篇所說.
耶穌是故意的,他刻意的.
他不是不小心.
高調的違反安息日不能工作的規定.
就是要他的敵人,好像法利賽人這班人看到.
這個犯禁果然觸動了這班法利賽人的神經.
有人就下了一個這樣的結論.
這個人不可能是上帝.
因為他們不守安息日.

$^{201}$對他們來說,守安息日是最高天條.
不管基於任何理由.
哪怕是再神聖,再正義.
違反天條的人,都是不屬上帝的.
不過其他法利賽人不同意這個全清性的結論.
他們就說,一個罪人怎能夠行這樣的神蹟呢?.
結果在法利賽人中間就起了紛爭.
沒辦法達成共識,這個去到16節了.
法利賽人既然沒辦法解說.
罪人為何能夠行神蹟這個理論的困局.
於是他們就想,打下矛波.
只要證明那個盒子不是盒子,原來不是盲的.
那就不存在黑眼復明這個神蹟.
沒神蹟,就不需要面對罪人行神蹟這個困惑.
於是他們就將盒子的父母傳召來.
要他們公認,他們的兒子其實不是盲的.
這是第三幕.
盒子的父母都是猶太人社會中最底層的人.
從來沒見過達官貴人.
突然被傳召,當然是很痰痘.
面對法利賽人凶神惡煞的逼問.
他們再蠢都明白法利賽人是想引導他們.
合謀去治癒他們兒子的人的罪.
他們不敢違逆這班達官權貴.
但又不敢造假證供.
於是他們只可以扯皮,兜兜玩玩.
兒子確實生下來是盲了眼.
這個事實太容易查證.
沒辦法作偽,盲了這麼多年.
怎麼大家都知道.
所以兒子是真的盲了.
至於他怎樣復明,由於我們不在現場.
我們就不太清楚.
兒子已經成年了.
法利賽人不如直接問他.
他已經成了人.
這句說話按照當時人的標準.
過了十三歲就成年.
比我們早熟很多.
十四五歲還能結婚.

$^{241}$十三歲就成年了.
有謀生能力可以承擔責任.
不過父母說這番說話.
或不是指兒子已經成年.
而是指他已經足夠大.
可以負責自己的言行.
不用父母代勞.
聖經作者在二十二節解釋.
他父母說話是怕猶太人.
因為猶太人已經商議定了.
若要認耶穌是基督的.
要把他趕出會堂.
雖然經文記載的內容有限.
不過我們根據常識推論.
我們可以猜測到.
那位核子的家庭的境況.
第一,兒子的眼睛突然復明.
理論上是天大喜訊.
做父母的不可能完全不關心.
不可能不打爛沙盤問道督.
要查明有什麼好事發生在兒子身上.
無疑當時兒子自己也所知有限.
他是蒙查查被醫治.
所以父母所知道的.
最多是兒子所知道的.
不過要說那對父母完全不知情.
這是不可能的.
除非他們對兒子完全不關心.
怎會完全不知情?.
第二,說這對父母可能不關心自己的兒子.
這又不是完全不可能.
在古代社會.
一個家庭通常不會只有一個兒子.
這對父母自己沒有失明.
其他的兒子擁有正常視力的可能性比較高.
你說他失明是為全兒子失明.
這是不能說的.
但他沒有.
理論上他的兒子不會每個都盲.
所以多數都應該沒有盲.

$^{281}$極有可能這對父母一直以來都將注意力.
放在其他正常的兒女身上.
天生殘缺的兒子被忽略.
在家中被邊緣化.
是很自然的事.
你想想他們任由這個兒子.
甚至叫這個兒子出去行乞為生.
就知道他們沒有能力.
也沒有心力去照顧這個兒子.
你自己去吃吧.
第三,古代社會因果報應的觀念是很根深蒂固的.
疾病是由罪孽造成的.
嚴重殘疾更是上帝的刑罰.
這個想法不會只有耶穌的門徒才有.
你記得一開始耶穌的門徒討論這個問題.
誰的罪.
所以疾病是由罪而來.
這是當時人們普遍的理解.
你可以想像這個兒子的父母.
一直以來都面對身邊的人歧視的眼光.
指指點點論斷他們是罪孽之家.
父母犯罪害死兒子.
父母也很可能認為兒子是個不祥人.
是上帝用他來懲罰他們.
這樣的報應觀念和人言可畏.
使子女和父母的關係很淡薄.
他們一直視他為負累包袱.
不關心他不在乎他.
如今兒子的眼疾雖然痊癒了.
但不知為何又招來權貴的打壓.
更加印證這是一個災星.
是一個陀衰家.
兒子會為家庭帶來咒詛.
所以最好跟他撇清關係.
這可能是父母的態度.
我說這些說話你覺得很抽象.
我是家中的第一個兒子.
所以我叫梁家倫.
我爸爸以為這是上天對他的掌上.
但我三歲就患哮喘.

$^{321}$亦是一個陀衰家.
我爸爸認為上天仍在咒詛他和家族.
所以很憎恨我.
有段很長的時間.
這些觀念你覺得很奇怪.
在上一代一點也不奇怪.
很正常的想法.
由於無法從父母口中.
得到確實將耶穌定罪的口供.
所以法利賽人就第二度傳召核子應信.
這是第四幕.
見到核子之後.
法利賽人不再繞圈.
上次見了一次.
今次見第二次.
於是他們就開門見山.
將他們對耶穌的判定說出來.
二十四節.
你該將榮耀歸給上帝.
我們知道這人是個罪人.
這句說話.
兩句短短的句子其實很簡單.
不要當這個人是治好你的病.
你的服命只是上帝的作為.
你只要將榮耀歸給上帝就行了.
至於這個人.
他是個罪人.
這是我們對他的判案.
你要接受這個看法.
核子雖然沒有受過教育.
未見過世面.
卻沒有像他的父母那樣.
怕達官貴人.
或者無知者不懼.
不懂得害怕.
或者一無所有者不懼.
Nothing to lose.
難鳴一條.
所以他很直率地回答.
二十五節.

$^{361}$他是個罪人不是我不知道.
有一件事我知道.
從前我是眼下的.
如今能看見了.
前我黑眼.
今我復明.
這是鏗鏘有力的見證.
我不和你們糾纏.
這個人是不是罪人的爭議.
我只知道他治好我的眼睛.
只是講客觀的事實.
簡單明快.
事實勝於雄辯.
在生命改變.
這個壓倒性的恩情事件面前.
一切的狡辯理論.
都是軟弱無力.
甚至可笑.
法利初人死心不息.
不斷要核子回答.
耶穌究竟做了什麼在他身上.
他不是想知道他做了什麼.
是想引導核子說.
他沒有做任何事在你身上.
想證明他沒有做過事.
他們的目的是要引導核子相信.
耶穌沒有做過事.
吐口水和泥是故弄言虛.
根本沒有醫治作用.
所以耶穌不是核子的真正恩人.
他們期待核子接受這個說法.
這個說法.
那個耶穌其實和我是毫不相干.
他只是偶然路過和我搭訕了幾句.
然後是我自己心血來潮.
或者性命引導.
我自己去洗眼睛.
然後竟然就復明了.
這個完全是上帝的作為.
是上帝醫治了我.

$^{401}$是不是很簡單.
總之不關和我搭訕的耶穌事.
核子沒有被法利初人的語言藝術.
真藝的藝.
藝術套住.
沒有想勾.
雖然他也不明白.
口水和泥有什麼醫治眼疾的功效.
也不知道耶穌其實做了什麼.
背後有什麼化學作用.
有什麼化學作用.
他不知道.
但他認定.
他在遇上耶穌這個人之後.
他的殘疾才得到復原.
所以是耶穌拯救了他.
對於法利初人反覆的查問.
他感覺到很困惑.
一個簡單的問題.
答了就是答了.
為什麼要反覆再問.
於是他開玩笑的問.
你們對耶穌和他做過的什麼什麼.
這麼有興趣.
你們是不是有意成為他的門徒.
27節.
法利初人應該對我這個小人物沒什麼興趣.
他們反覆抓住我去查問耶穌的事.
他們肯定是對耶穌有興趣.
那核子當然說.
捉路了法利初人.
他們宣稱他們只是摩西的門徒.
遵守摩西所頒佈的律法.
他們不會被耶穌這個來歷不明的人吸引.
他們不會做他們的門徒.
核子在30到33節這樣回答.
他開了我的眼睛.
你們竟不知道他從哪裡來.
這真是奇怪.
我們知道上帝不聽罪人.

$^{441}$唯有經奉上帝進行他旨意的.
上帝才聽他.
從創世以來.
未曾聽見有人把生來時核子的眼睛開了.
這人若不是從上帝而來.
什麼也不能作.
能夠讓核眼的得意看見.
這個肯定是一個神通廣大的人.
這個人應該遠近馳名.
為什麼法利初人竟然不知道他的來龍去脈.
我以為你們法利初人都是上通天文下曉地理.
是無所不知的人.
為什麼你們跟我一樣孤陋寡聞.
古往今來從來沒有人試過使天生核眼的人恢復視力.
能夠做到這麼石破天驚的神跡的人.
肯定是從上帝而來.
並且一定是運用上帝賜給他的能力.
你要我論斷這個人是罪人.
不要說和我得存餘的經歷完全相衝突.
怎麼可能是一個罪人治了我.
也違反了一般的常識推理.
上帝怎麼會聽一個罪人的祈禱.
讓他有醫治的大能呢.
這個核子沒有任何社交常識和經驗.
不懂禮節不懂規矩完全率性而行.
但他憑著他講誠實話的堅持.
憑著最簡單的常識推理.
竟然就搏倒了那班通曉律法.
擅長顛倒黑白謬論連篇的法利賽人.
讓他們詞窮理屈.
在老修成奴之下.
他們決定將核子趕出去.
不僅是趕他離開他們.
而是將他趕出會堂.
即是說他們不可以再隸屬和參與猶太人的社群.
我們來到最後一幕.
耶穌聽說那個核子被趕出會堂.
孤立無援.
他有機會再遇上這個核子.
然後單刀直入呼召他.

$^{481}$35節.
你信上帝的兒子嗎?.
核子雖然未親眼見過耶穌的真容.
但從他的聲音應該認得到他.
這是洗他的眼睛復明的恩人.
有趣的是.
核子連耶穌曾經自稱是上帝的兒子都不知道.
他過去只認定耶穌醫好他的眼睛的先知.
就是從上帝而來的人.
他在36節就問.
主啊,誰是上帝的兒子? 叫我信他?.
耶穌表明身份.
說他自己就是.
核子立即跪倒在地.
堅定的回答.
主啊,我信.
沒有猶豫.
沒有保留.
這就成為了耶穌的門徒.
你看到核子的龜唇.
也看到法利賽人.
捏造謊言左之又絕的醜態.
耶穌做了一個蓋炭.
39節.
我為審判到者誓上來.
叫不能看見的可以看見.
能看見的反黑了眼.
耶穌洗一個盲眼的人看得見.
這是一個不折不扣的神蹟.
能夠行這個神蹟的必定有特殊的身份.
這麼簡單明確的事實.
致命聰明博學的法利賽人.
竟然視而不見.
他們的視力連一個盲眼的人都比不上.
耶穌是世界的光.
接受耶穌的就是光明之子.
拒絕耶穌的仍舊活在黑暗裡面.
他們在黑暗中甚麼都看不見.
核子和法利賽人.
究竟誰才是真正的盲眼呢?.

$^{521}$在場有些法利賽人.
聽到耶穌嘲諷,映射他們.
40節他們質問.
難道你是說我們盲了眼嗎?.
41節耶穌毫不客氣地回應.
你們的情況比盲了眼更糟糕.
如果真的盲了眼.
你們就不需要為自己做錯判斷負責.
不過現在你們的眼是.
看見又看不見.
有眼無珠.
拒絕上帝差來的遺旨.
拒絕祂所傳講的真理.
你們要為這個拒絕承擔後果.
地獄的永盈等待著你們.
這是耶穌在這段經文中說的最後一句話.
整個五學劇就講完了.
這個故事的主角是被耶穌覆命的核子.
他的際遇和中國一句成語.
「塞翁失馬焉知非福」非常切合.
一出世就失去視力.
靠賴行乞為生.
處於社會的最底層.
被人歧視嫌棄.
警方應該是差到不可以再差.
急至遇上耶穌.
在毫無心理準備的情況下被醫治.
這應該是天下間最大的喜訊吧.
幸福之神眷顧他.
不過不幸的是.
由於當權的宗教領袖要打壓耶穌.
而他是被耶穌醫治的.
他因此也成為了被針對的對象.
當他拒絕.
背著良心配合當權派誣告耶穌之後.
他就面對連番的傷害.
父母被欺負.
在壓力之下和他劃清界線.
他最終又被趕逐出會堂.
原本以為眼睛覆命.

$^{561}$可以使他過正常的生活.
由邊緣進入主流.
沒想到他竟然進一步被邊緣化.
更加被邊緣化.
這真是禍兮福所以.
不過在走到山窮水盡之際.
耶穌卻來呼召他成為信徒.
接受永生人間最大的禮物.
核子的命運.
你看到他像坐過山車一樣.
大喜大樂非常之疼.
核子確實是一個.
山窮水盡走投無路的境地.
才遇上耶穌.
被他呼召.
除了眼睛能夠看見之外.
他信主前的情況.
醫好到信主前的情況.
應該比覆命前還要差.
眼睛覆命了.
年輕力壯的他.
不可以維持黑食的工作.
但沒有受過教育.
沒有一技之長.
又沒有人脈關係.
他連每天去哪裡找食物都不知道.
你可以想像他連吃飯都構成問題.
加上被猶太當局針對.
被父母嫌棄.
被社群排擠.
天下之大.
他完全沒有容身之處.
你一個這樣一無所有.
一無所事的人.
要他歸順耶穌是太簡單的事.
他不用跟耶穌說.
容我先埋葬我的父親.
一無所有的人.
沒什麼要做.
容我先回家慈別父母.

$^{601}$沒有牽掛.
沒有負累.
決然一生.
改喚信仰是不用付代價的.
因為根本沒有可付的代價.
生無長物.
爛命一條.
沒什麼可以撇棄.
這令我想起很多的信徒群體.
我首先想起上世紀50年代.
中國政權易守.
大量難民逃離香港.
有一批人.
特別是前國民黨官兵.
逃到調景嶺.
他們連廣東話都不懂.
要融入香港社會是非常困難的.
他們在調景嶺.
這個山旮旯的地方.
搭著草棚為生.
然後有難民宣教士.
在中國大陸被趕出.
無地容身的宣教士.
走到這班難民的中間.
去傳難民的福音.
然後建立難民的教會.
這班難民聽福音非常容易.
當傳教士說.
「金也空 銀也空 這世界非我家」.
他們根本就立刻趴下.
哭得嘔吐.
哪用證明這世界非我家.
明明就是非我家.
「金也空 銀也空」.
也不用證明.
於是他們立刻舉手信耶穌.
並且幾天後就決志奉獻讀神學.
於是傳教士在調景嶺幫他們開了一間聖經學院.
他們就立刻讀書.
聖經都未讀完一次.

$^{641}$福音書都未讀完一次.
就讀神學.
然後就去做難民傳道.
難民始終不是只向難民傳福音.
而是向難民這樣的傳道.
一無所有的他們.
不怕吃苦.
因為都不會比現在差.
所以他們就真的咬緊牙根.
去打拼 去開荒 去建立教會.
我又想起上世紀八十年代 九十年代.
在中國河南 安徽.
這些最貧窮的地方.
福音發展得非常快.
很多信了耶穌的人.
他們就兩個兩個.
真的不用刻意去做.
就能做到耶穌.
在馬福音的說法.
不要帶兩件衣服 不要帶錢囊.
什麼都不要帶 因為沒有得帶.
他們就一路去討飯 一路去傳福音.
然後就將福音帶給很多很多地方.
一些一無所有的人.
他們就發揮了一無所有的長處.
一無掛慮 是不能掛慮的那種一無掛慮.
的人他們真的可以勇往直前.
弟兄姊妹.
我也認識很多的人.
特別是那些中年獻身做傳道.
那種糾纏掙扎.
是令到很困擾.
我有一個好朋友.
在奉獻之前.
在香港是做相當好的待遇的一個政府官員.
舊式的.
還可以派去英國讀書.
跟老婆去買那種包細食宿旅遊.
很愉快的那種.
奉獻讀神學之後.

$^{681}$生活的落差大到不得了.
所以他很多時候都忍不住會心淫.
沒有了這個 沒有了那個.
以前開的車現在開不了.
以前有的待遇現在不可以了.
我聽到他說當然我同情他.
不過我其實是共鳴不到的.
因為我從來沒有試過有過.
所以我也不知道沒有是什麼事.
所以他有掙扎我沒有.
我不知道這樣的掙扎是怎麼一回事.
所以我覺得我比他愉快很多.
因為我不用掙扎.
我想說什麼呢.
說這麼多這些.
什麼時候你覺得自己是somebody.
什麼時候你覺得自己無論生理物理心理.
覺得自己有很多東西.
你要跟隨耶穌都是難很多.
都是婆婆媽媽支支作作想來想去.
什麼時候你真正認識自己是難的一條.
是一條鋼棍.
什麼時候你就覺得無所謂.
不會輸因為不知道輸什麼.
不知道有什麼好輸.
我要說你們不容易做到的.
確實不容易做到的.
不一定說要奉獻做傳道人.
我說的是我們在生活上死心塌地跟隨耶穌.
你越覺得自己有東西.
你就越難傻乎乎地跟隨到底.
你越覺得自己聰明.
你就會像那些法利賽人那樣轉彎抹角.
做很多神學詮釋.
簡單的心靈說簡單的真理.
前我失喪,今被尋回.
耶穌踩入我的生命裡面.
翻轉了我.
這是我知道的事實.
祂是我的主,我的救主.

$^{721}$我從來不抵賴這件事.
我就跟隨祂.
我就完全壓著在祂身上.
不會左扭右扭.
扭一手,扭兩手,扭什麼扭.
是不是?.
弟兄姊妹.
我們應該不可能是天生黑眼的人.
我們不可能去到討飯的地步.
我們像財主多過像拉薩路.
成為了祂的門徒.
所以我們都不是那些進不了天國的財主.
我們自己去審查一下.
在核子的故事裡面.
在祂的際遇經歷裡面.
祂的人生態度.
祂對耶穌呼召的反應.
我們去看一看.
能不能找到自己的身影.
我不是百分百的核子.
但我這部分像祂.
我那部分的經歷跟祂差不多.
我們自己對號入座.
然後我們在這樣的基礎上.
我們去調教我們跟耶穌的態度.
耶穌,你知道我正在跟隨你.
不過我跟你跟得不算很好.
這個我自己都知道.
不用你講,不用你提醒我.
我不是全心全意的.
有時候繞一圈,有時候走回頭路.
有時候為自己的婆婆媽媽.
阿姨阿祖做些解釋.
我希望,我誠懇希望.
我可以跟你更好,更乾脆.
更撇脫.
求主幫助我們.
(影片完結).
\newpage



\section{馬可福音 10:32-45}
\label{sec:JJ124JmNAMw}
\textbf{2025年9月28日崇拜直播｜吳真真牧師｜耶穌給門徒的十課－走僕人之路｜馬可福音 10\_32-45}
\newline
\newline
連結: \href{https://youtube.com/watch?v=JJ124JmNAMw}{\texttt{https://youtube.com/watch?v=JJ124JmNAMw}} ~~~~ 語音日期: 2025-09-28
\newline
\newline
\hyperref[sec:LeA_EIgcwcw]{\small{< < < PREV SERMON < < <}}
~
\hyperref[sec:index_chronic]{\small{[返順時目]}}
~
\hyperref[sec:Jhd45blKy7w]{\small{> > > NEXT SERMON > > >}}
\newline
\newline
馬可福音 10:32-45
\newline
\begin{longtable}{cl}
\hline
\hline
章節 & 經文 (和合本修訂版)\\
\hline
10:32 & \begin{tabularx}{0.7\textwidth}{X} 他們行路上耶路撒冷去。耶穌在前頭走，他們很驚訝，跟從的人也害怕。耶穌又叫十二使徒來，把自己將要遭遇的事告訴他們， \end{tabularx} \\ \\ \relax
10:33 & \begin{tabularx}{0.7\textwidth}{X} 說：「看哪，我們上耶路撒冷去，人子將被交給祭司長和文士；他們要定他死罪，又交給外邦人。 \end{tabularx} \\ \\ \relax
10:34 & \begin{tabularx}{0.7\textwidth}{X} 他們要戲弄他，向他吐唾沫，鞭打他，殺害他；三天後，他要復活。」 \end{tabularx} \\ \\ \relax
10:35 & \begin{tabularx}{0.7\textwidth}{X} 西庇太的兒子雅各和約翰進前來，對耶穌說：「老師，我們無論求你甚麼，願你為我們做。」 \end{tabularx} \\ \\ \relax
10:36 & \begin{tabularx}{0.7\textwidth}{X} 耶穌對他們說：「要我為你們做甚麼？」 \end{tabularx} \\ \\ \relax
10:37 & \begin{tabularx}{0.7\textwidth}{X} 他們對他說：「在你的榮耀裡，請賜我們一個坐在你右邊，一個坐在你左邊。」 \end{tabularx} \\ \\ \relax
10:38 & \begin{tabularx}{0.7\textwidth}{X} 耶穌對他們說：「你們不知道所求的是甚麼。我所喝的杯，你們能喝嗎？我所受的洗，你們能受嗎？」 \end{tabularx} \\ \\ \relax
10:39 & \begin{tabularx}{0.7\textwidth}{X} 他們對他說：「我們能。」耶穌對他們說：「我所喝的杯，你們要喝；我所受的洗，你們也要受。 \end{tabularx} \\ \\ \relax
10:40 & \begin{tabularx}{0.7\textwidth}{X} 可是坐在我的左右，不是我可以賜的，而是為誰預備就賜給誰。」 \end{tabularx} \\ \\ \relax
10:41 & \begin{tabularx}{0.7\textwidth}{X} 其餘十個門徒聽見，就對雅各和約翰很生氣。 \end{tabularx} \\ \\ \relax
10:42 & \begin{tabularx}{0.7\textwidth}{X} 耶穌叫了他們來，對他們說：「你們知道，外邦人有君王作主治理他們，有大臣操權管轄他們。 \end{tabularx} \\ \\ \relax
10:43 & \begin{tabularx}{0.7\textwidth}{X} 但是在你們中間，不可這樣。你們中間誰願為大，就要作你們的用人； \end{tabularx} \\ \\ \relax
10:44 & \begin{tabularx}{0.7\textwidth}{X} 在你們中間誰願為首，就要作眾人的僕人。 \end{tabularx} \\ \\ \relax
10:45 & \begin{tabularx}{0.7\textwidth}{X} 因為人子來，並不是要受人的服事，乃是要服事人，並且要捨命作多人的贖價。」 \end{tabularx} \\ \\
[1ex]
\hline
\hline
\end{longtable}
$^{1}$(主持人:今天經訓是在馬可福音10章32至45節,新日聖經第63頁,馬可福音10章32節,請聽我讀出).
(主持人:他們行路上耶路撒冷去,耶穌在前途走,門徒就稀奇,跟隨的人也害怕.耶穌又叫過十二個門徒來,把自己將要遭遇的事告訴他們說.).
(主持人:看啊,我們上耶路撒冷去,人子將要被交給祭司長王文士,他們要定他死罪,交給外邦人,他們要戲弄他,套唾沫在他臉上,鞭打他,殺害他,過了三天,他要復活.).
(主持人:西畢泰的兒子雅各約翰進前來,對耶穌說:夫子,我們無論求你什麼,願你給我們作.耶穌說:要我給你們作什麼?他們說:賜我們在你的榮耀裡,一個在在你右邊,一個在在你左邊.).
(主持人:耶穌說:你們不知道所求的是什麼,我所學的杯,你們能學嗎?我所受的浸,你們能受嗎?他們說:我們能.).
(主持人:耶穌說:我們所學的杯,你們也要學,我所受的浸,你們也要受,只是坐在我的左右,不是我可以賜的,乃是為誰預備的,就賜給誰.).
(主持人:那十個門徒聽見,就魯魯雅各約翰,耶穌叫他們來,對他們說:你們知道,外邦人有尊為君王的,治理他們,有大臣操權管束他們,只是在你們中間,不是這樣.).
(主持人:你們中間誰願為大,就必作你們的用人;在你們中間誰願為首,就必作眾人的僕人.).
(主持人:因為人子來,並不是要受人的服侍,乃是要服侍人,並且要寫明作多人的屬也.).
聖經獨筆.
(主持人:各位弟兄姊妹平安.).
現在社會很多時候都會講生涯規劃,我們很早就看小朋友和年青人有什麼特質,有什麼興趣,於是就安排他們走一個將來要投身在什麼職業的路上..
你會問他們將來長大想做什麼,現在的答案跟我們以前傳統的做老師,醫生,警察,消防員一樣,但現在多了很多工種,你會聽到要做YouTuber,電競選手,.
或者你會聽到小朋友想做Superhero,但當我們思考的時候,我覺得有一位在美國的牧者,他的鼓勵是值得我們認真考量的..
這個牧者是Trust Rindor,他是被譽為近50年最具影響力的傳道人之一,他有一本書叫做"Improving Your Serve",他鼓勵我們換一個角度思考這條問題..
他讓我們一起想像一下,我們不如問一下耶穌,耶穌,你希望我們長大之後會成為什麼樣的人呢?.
Trust Rindor說,他相信耶穌對著我們每一個都會答同一個答案,他會說,我希望你們每一個都與別不同,我希望你們每一個都成為僕人..
我想,沒有人會想,我的生活規劃,最後我要做的是一個僕人,因為僕人看起來很卑微,僕人看起來沒有什麼被人尊重,沒有什麼尊嚴..
但是我們覺得Trust Rindor給我們的一個分享,一個鼓勵,是非常值得我們思考的,因為這正正是耶穌對我們這群門徒的吩咐..
所以今天我們就一起思考一下這個課題,究竟我們的人生如何走這個僕人的道路,我們一起祈禱..
參見天父,求你幫助我們,讓我們切實思考,當人一生可能追求的是得人尊重,有一個好的崗位,有一個被人稱許的生命的時候,你卻呼籲你的門徒去作僕人..
上帝,我祈求你幫助我們,讓我們心靈打開,讓你親自跟我們每一個人說話..
高沒僕人的口,讓我說得明白,也願意讓弟兄姊妹打開他們的心,讓他們都能夠明白你自己的真理..
我們這樣禱告,教託,奉主耶穌基督的聖名,致祈求,阿們..
大家是否記得,從七月開始,我們有一個系列,叫做耶穌給門徒的十堂課,不知不覺,現在是第十堂課,大家是否覺得呢?.
大家不覺得,我們七月開始的時候,就是馬可福音第六章,馬可福音第六章的開始,就是耶穌開始差派門徒出去傳道,然後到今天第十章,我們就說到第十堂道..
我們為什麼說第十章呢?最主要原因是因為第十一章,耶穌已經進入耶路撒冷,第十章,第六章開始,到第十章,.
都是說耶穌朝向耶路撒冷,預備在地上,來到這個世界的目標,就是為眾人死,釘身十架的這條路..
所以,他跟門徒很緊密地在這段路上,說了很多很多的言訓..
而今天我們就會說,如何走僕人的路..
其實在第八章二十七節開始,當門徒彼得說,耶穌啊,你是基督之後,耶穌就第一次預言自己要死,和他會復活..
然後,從第八章二十七節開始,到第十章,不斷重複一樣東西,就是耶穌在路上..
無論是八章,九章,十章,都不斷重複耶穌在路上,和門徒一起走..
而在這段路上的旅程當中,今次第十章,第三十二節,是第一次很清楚地說,他要去哪裡..
他的目的地是耶路撒冷..
表面上,看來耶穌和門徒一起去耶路撒冷,預備守節,沒什麼特別..
但在這段話中,耶穌的門徒和其他跟從的人,都感到稀奇和害怕..
其實,為什麼門徒會稀奇,又為什麼會害怕呢?.
當門徒已經知道,耶穌現在是用利塞亞的身份去耶路撒冷,他們聽了兩次的預言,.
耶穌說他會受苦和復活,但他們看到耶穌現在是很堅定地走向受苦的路..

$^{41}$他們可能聯想到,利塞亞先知,他預言將有利塞亞的誕生,.
而他在宣告上帝的話的時候,都不斷經歷很多苦難..
但在以賽亞書50章第7節,他說:主耶和華必幫助我,所以我不抱愧..
我眼著臉面如象堅石,我也知道我必不自蒙羞..
其實他們看到,耶穌基督作為利塞亞的身份,他要面對苦難的時候,他的意志是很堅定和勇敢..
因為這個緣故,他們覺得很震懾,覺得很稀奇..
同樣地,當中有些跟從的人,或許他們心裡不單是全敬畏的心,.
但他們也有可能因為他們而害怕..
他們害怕什麼呢?就是一直以來,在同行的過程中,.
其實開始,耶穌基督已經和一群宗教領袖有敵對的對話..
他們開始擔心,到底到了耶路撒冷會發生什麼事呢?.
耶穌是不是會受害呢?.
當門徒和跟從的人有這個稀奇和害怕,耶穌做了什麼呢?.
耶穌就叫了這12個門徒到他們面前,.
將他們要受苦和復活的預言,第三次告訴他們..
他沒有安撫他們,說不用害怕,不需要稀奇..
而是他更加鉅細無遺地將他們之前兩次受難的預言說得更清楚..
當初第一,二次,他講到他們被長老,祭司,文士氣絕..
這次他講到,如何氣絕他們呢?原來這群人是要定他們的罪..
當初他預言的時候說,我要交到別人的手裡..
這次他很清楚地說,我要交到外邦人的手裡..
而且他更加講得很清楚,在他們被殺之前,.
之前沒有說,但這次他講得很仔細,.
他們是要被戲弄,被人唾口水在他們的臉上,甚至要被鞭打..
大家聽完耶穌這一番預言,我不知道你們有沒有一個很大的感覺,.
就是耶穌很被動,他被交給猶太的領袖,.
他被定罪,他被送到外邦人的手裡..
他成為受苦的人,被唾罵,被戲弄,被鞭打..
耶穌在這一刻所形容的他自己,好像沒有說話,沒有任由別人的處置..
不過,如果我們去看第45節,我們就看到這一句經文:.
「因為人子來,並不是要受人的服侍,乃是要服侍人,並且要寫命作多人的屬駕.」.
耶穌去耶路撒冷為什麼這麼堅定?.
是因為他為了一個他來這個世界裡最終極的使命和目標,.
為世人,為罪人,附上屬駕,而堅定地去耶路撒冷..
聖經學者說,這是神聖的被動式,.
因為他之所以被人這樣對待,完全是因為他對上帝旨意主動的信服..
耶穌來是主動的服侍人,他為的是要作罪人的屬駕..
屬駕這個字是代替的意思,當人要釋放一個奴隸或罪犯的時候,.
是要交贖金的,奴隸和戰犯被扣押,他們能否得到自由,.
取決於能否得到一筆錢,或者得到一碗粥,.

$^{81}$挾持著他人的要求,以致他們能得到自由..
我們世上的人都成為罪的奴僕,我們的刑罰就是死亡,.
我們需要有人來贖回我們的自由..
耶穌基督現在就付上了這個代價,.
他不單給了屬駕,而且他更加成為屬駕,.
他為了拯救我們的生命,獻上了自己的生命..
我不知道大家能否想像得到,其實耶穌越走近耶路撒冷,.
越走近他要走到十輛壽斧的終點,他一定有很多掙扎..
但是耶穌很清楚,他走上這個苦路的目的,.
就是要成為祭物,犧牲自己的生命,去拯救世上的罪人生命,.
把罪人從罪孽綑綁釋放出來,以致我們不需要受世上的憂患,.
讓我們過著豐盛的生命..
所以耶穌沒有退縮,他勇於去面對靈慾,.
可能是藐視,但是他仍然堅定前行,.
在這條走僕人的路上堅定不移..
但是我們的信徒呢?我們的信徒的服事,.
其實是不是像耶穌那麼堅定呢?.
還是我們很容易走到一些誤區裡面,失去了焦點,.
就好像我們待會要看的一段經文,.
看雅各和約翰,他們如何走進這些事奉的誤區裡面..
其實耶穌已經不是第一次去講受苦和復活的預言,.
這次是第三次..
這次第三次,究竟那些門徒有什麼反應呢?.
第35節,西比太的兒子雅各約翰進前來,.
對耶穌說:「夫子,我們無論求什麼,願你給我們,願你給我們做.」.
大家覺不覺得很奇怪?.
耶穌說他會受苦,他會受死,.
但是現在那些門徒就跟他說:.
「耶穌啊,我們求什麼,希望你也幫我們做.」.
我覺得聽起來很嘔心,.
好像叫耶穌開張給我們任填的支票,.
其實這個心態反而是:.
「耶穌,你服從我,你聽我說.」.
其實是挺操控耶穌的..
而諷刺的是,其實當第二次耶穌講受苦的預言的時候,.
聖經告訴我們第九章的時候,.
那些門徒不明白,又不敢問,.
而且在路上又在爭論誰是大的..
耶穌教導了他們,.
「若有人願意作首先的,他必作眾人末後的,作眾人的用人.」.

$^{121}$我在想,那些門徒究竟明白這句嗎?.
他們是不是當了耶穌這樣教,.
為首的就要做人的末後,就要做眾人的僕人?.
現在耶穌是為首,是我們的首,.
也就是我們的僕人,.
於是你就要聽我們的話..
其實我們真的要小心,.
有時候我們服事的時候會掉進一個誤區,.
就是「我話事,主啊,在我服事的時候,你滿足我吧,.
你叫我服事嗎?看看我有沒有空,.
看看我有沒有興趣,或者你們能不能遷就我.」.
這是一個服事的誤區,我們一定要小心留意..
然後,耶穌當然不會說「好啊,你要什麼我給你什麼.」.
耶穌就問清楚他們,.
第36節,「你要我給你們什麼?」.
他們就說「賜我們在你的榮耀裡面,.
一個坐在你右邊,一個坐在你左邊.」.
意思就是將來耶穌統治以色列國的時候,.
當君王的時候,.
你讓我們雅各和約翰能夠分享你這份勝利和榮耀..
雅各和約翰他們兩個為什麼會有這樣的要求呢?.
因為他們兩個和彼得三個人是耶穌最親近的門徒,.
他們和其他門徒一起經歷的就是他們見證著耶穌有醫治的權柄,.
見證著耶穌餵飽五千人,餵飽四千人,.
見證耶穌平靜風浪,.
但是根據《馬可福音》的記載,.
他們三個人是比較不同的,.
他們是怎樣不同呢?.
例如當捱老女兒死的時候,.
第五章的時候說到,.
耶穌在醫治血漏病婦人之前,.
知道捱老女兒有事有病,.
但是當他醫治好血漏婦人的時候,.
有人來通知捱老的僕人,.
「你的女兒已經死了.」.
耶穌沒有帶其他人,.
祂只帶著這三個門徒,.
就是彼得,雅各和約翰,.
去到捱老的家,.
是他們三個親眼見證著耶穌叫這個女兒起身..

$^{161}$然後去到第九章,.
耶穌也是帶著這三個門徒,.
暗暗地去到高山,然後登山變像,.
他們親眼目睹耶穌登山變像..
其實在這樣的背景下,.
可能他們真的有信心,.
耶穌是一個滿有大能,.
甚至可以叫死人復活的人,.
所以就算耶穌接下來會面對什麼挑戰,.
什麼艱難,甚至耶穌真的要死,.
是的,他可以復活..
不過他們的觀念就像.
《使徒行傳》一章六節所說,.
他們聚集的時候問耶穌說,.
「主啊,你復興以色列國,就在這時候嗎?」.
現在《使徒行傳》一章六節是說.
耶穌復活,預備升天的時候,.
他們想,面前的彌賽亞得勝了,.
他真的死了,他真的復活了..
不過他們想的是,.
耶穌要建立的彌賽亞國度,.
其實在地上,是以色列國的復興..
所以他們真的想著,.
耶穌在以色列國復興的時候,.
他們可以得到這份榮耀..
而且不止這樣,.
耶穌在《馬太福音》十九章二十八節說,.
「我實在告訴你們,你們這跟從我的人,.
到復興的時候,人只坐在他榮耀的寶座上,.
你們也要坐在十二個寶座上,.
審判以色列的十二個支派.」.
耶穌真的這樣說,你們十二個門徒,.
將來跟我一起坐在寶座上,.
十二個都有份,要去審判十二個支派..
但是他們想的是,.
我比其他的門徒更加尊貴,.
我比他們厲害,我比他們更加親近你,.
所以我要坐在寶座最榮耀的位置,.
除了你是君王之外,.
我要坐在右邊,我要坐在左邊..

$^{201}$這些都是最有權柄,最被尊貴的位置..
所以我們也要跟他們一樣,.
要反思一個誤區,.
就是其實有時我們的服侍,.
有時我們都是為了滿足自己的慾望,.
自己的虛榮,.
我們的動機會不純正,.
這也是我們要小心留意的地方..
我們覺得我們比別人更有資格去得著榮耀..
然後耶穌就跟他們說,.
你們不知道自己在求什麼,.
我所學的杯,你們能學嗎?.
我所受的洗,你們能受嗎?.
耶穌沒有斥責他們,.
因為他們的確不明白耶穌受苦,.
彌賽亞不是建立在地上的國度這個道理,.
他們不知道自己在問什麼..
耶穌跟他們說,.
我所學的杯,我所受的洗,.
祂告訴你,他將要受苦..
杯可以是正面的祝福,.
但這裡是說負面的受苦,.
因為耶穌將要成為一個無辜的義人,.
替那些不義的罪人受上帝的刑罰..
洗可以是洗禮,.
但耶穌所受的洗,.
不單單是施洗約翰所施的洗禮,.
而是他被洪水淹沒殉道的含義..
所以當耶穌問他們,你們能受嗎?.
其實他是問,你們願不願意主動,.
順服地去接受從神而來受苦的服事?.
雅各和約翰說,能..
他們很快,很心急,馬上說,能..
真的可以,為何他們會這樣說?.
第一,他們未能明白耶穌得勝和得榮耀,.
所要受的苦難和挑戰是甚麼..
第二,他們也會想用自己的能力去得到榮耀..
他們相信,如果耶穌將來真的讓他們坐在榮耀寶座的左右,.
得到這份尊榮和權力的時候,.
忍受一點苦難算甚麼?.

$^{241}$但我們要留心,這是陷阱..
我們不要以為我們可以用受苦,.
用我們的服事,.
以為我們可以得到權柄和榮耀..
因為在《馬學福音》第十章第三十節,.
耶穌告訴他們,.
「跟隨他的人,在今生要受到逼迫,.
在來生就得著永生..
所以他們所受的苦,不是為了得到虛榮,.
得到他們自己的權柄,.
乃是他們所受的苦,是在天上將來得著榮耀.」.
而且,耶穌很清楚地告訴他們,.
他們是否坐在左邊或右邊,不是他們作主..
其實這是上帝的主權,是上帝的決定..
其實,祂很想告訴他們,.
這件事不是重要的,.
這件事不是你們可以得到的,.
這是上帝的恩賜..
其實,人在走事奉的路,.
或者人的心,.
在事奉的時候走入誤區,.
是很常見的事..
人的內心的私慾和虛榮感,.
常常都在吸引我們..
當我在預備這篇信息的時候,.
我在恥笑門徒,.
「為什麼你們可以這樣?」.
上帝提醒我,.
「你覺得自己很好嗎?」.
我立即想起一件事,.
我中學的時候,.
剛剛信耶穌,.
我中二的時候,缺了字之後,.
我很投入在學校裡事奉..
我們一級有七班,.
每班有一個叫做宗教部的角色,.
除了有班長,財務,康樂,.
還有宗教的角色..
我從中二開始就做班宗教部..
我們每個月都有周會,.

$^{281}$那些周會就是同一級,.
我們有七班,.
七班會給你一個主題,.
然後你們就去比賽..
每一次我帶領的班都是拿冠軍的,.
於是就很厲害了..
從中三開始,.
我就不再只是做級的宗教,.
我做社的宗教,.
好像升級一樣..
其實這種感覺挺飄搖的,.
於是我發覺,.
有一次有一個姐妹帶我們去的教會,.
是高中的時候,.
我第一次去完別人的教會,.
崇拜,.
完結之後,.
我就跟一起去的姐妹說,.
「你猜我跟你誰會快點做到崇拜的招待呢?」.
我真的覺得,.
為什麼我會這樣想呢?.
我不是回去崇拜嗎?.
為什麼我會想,.
事奉,.
做到什麼崗位就好像很厲害,.
其實就是在滿足自己的慾望,.
說到底,.
我在服侍自己..
我想起唐崇明牧師的一個分享,.
他說有一個人走到牧師面前,.
跟牧師說,.
「牧師,我被選為執事,.
我覺得自己很不配.」.
其實他在說這句話的時候,.
他心裡真正希望牧師跟他說的是什麼呢?.
「你不要這麼謙虛,.
主的恩典在你身上,.
你有你配的地方.」.
但是他猜不到,.
那個牧師跟他說,.

$^{321}$「本來你就是不配的.」.
他聽完之後,.
他的防衛機制就起來了,.
然後他就很反感,.
他就問牧師,.
「為什麼我不配?」.
然後牧師問他,.
「你剛才跟我說什麼?.
你不是說我被選為執事,.
覺得很不配嗎?.
究竟你是否誠誠實實地.
真心覺得自己不配?」.
他說,.
「你本來真的不配,.
其實我本來也不配,.
不過神之所以將這個職事交託給我們,.
不是因為我們配不配,.
而是因為祂憐憫我們,.
讓我們學習謙卑地服侍上帝.」.
後來那十個門徒聽到.
雅各和約翰這樣跟耶穌說話,.
於是他們就牢牢地對他們說,.
他們生氣不是因為.
「為什麼你們這麼自私?」.
其實他們生氣的緣故是.
他們也有這樣的要求,.
他們也很想在耶穌的左右,.
不過你說了「比我快」,.
你折竹先登,.
所以耶穌看到門徒有.
一些不正確的事奉觀念,.
於是他就像《馬可福音》十章四十二節這樣說,.
他就叫他們跟前,.
他沒有責怪他們,.
但他再一次焦聚他們,.
告訴他們門徒的事奉觀應該是怎樣,.
因為他們錯了,.
他們想錯了,.
他告訴他們,.
其實他不是罵他們,.

$^{361}$不是說「你們想錯了,.
你們不配做大人,.
你們做奴隸吧」,.
他不是這樣說,.
而是他很想告訴他們.
事奉觀跟這個世界的人的想法,.
跟天國的價值觀是不同的事,.
所以耶穌一開口就跟他們說,.
「你們知道,.
外邦人有尊為君王的,.
治理他們,.
有大臣操權管束他們.」.
耶穌用一個世界的模式跟他們作比對,.
在上的當權者用他們的權柄,.
用他們的地位去管理,.
統治屬他的子民,.
用他們在上的姿態.
要求在下的子民去臣服他們,.
去滿足他們自身的利益和榮耀,.
但是耶穌在43節這樣說,.
「只是你們中間不是這樣,.
絕對跟這個世界的價值觀不同,.
你們中間誰為大,.
就必作你們的用人,.
你們中間誰為首,.
就必作眾人的僕人.」.
他告訴我們,.
因為神國的成本觀念根本不是這樣,.
在世界上,.
當你為大衛首的時候,.
你頭上的冠冕,.
和別人稱呼你的銜頭,.
就是你的記號,.
但不像這個世界,.
在神的國度裡,.
你為大的記號,.
就是束腰的手巾和洗腳盤,.
就像耶穌一樣,.
為門徒洗腳,.
謙卑服侍..

$^{401}$上帝去衡量一個人的屬靈,.
是否達到標準,.
是否成功,.
不是衡量我們擁有什麼成就,.
不是衡量我們獎賞的數量,.
而是衡量我們有沒有將別人放在自己面前..
在神的國度裡,.
我們不是爭取,.
不是增敬,.
不是弄權,.
去拿到高位,.
在神的國度裡,.
為人的好處著想,.
成為別人的幫助,.
關心別人的需要,.
具有服侍人的心,.
才是天國裡最大的..
其實這件事一點都不容易,.
因為「獄人」這個字,.
其實是奴僕的意思,.
奴僕是主人的財產,.
他沒有自主權,.
主人叫你做什麼,.
你就做什麼,.
可以做最卑下的事情,.
正常,沒有人願意做這個階層,.
但耶穌要求我們要穿上眾人的奴僕,.
意思是我們要謙卑自己,.
去服侍人,.
這樣才是為首,.
這樣才是偉大..
有一個初來教會的弟兄,.
去到一個新的教會裡,.
剛好牧師說這篇道,.
他就去問牧師,.
「牧師,你教會有多少人?」.
牧師說,「我們大概有150個會眾左右.」.
那個人就說,.
「蛤?牧師,你不是有150個老闆?」.
這位牧師就跟他說,.

$^{441}$「不是,我只有一個老闆,.
我這個老闆就是主耶穌,.
主耶穌叫我做什麼,我就做什麼.」.
這個牧師說得很對,.
教會裡的門徒,.
我們不是分等級,.
我們不是有階級,.
誰管理誰,.
因為我們所有的門徒,.
都是聽耶穌的話,.
我們有的是與別不同,.
不同的恩慈,.
不同的崗位,.
但我們全部都是主的奴僕,.
耶穌的僕人,.
我們不是誰給誰超賢,.
我們都是為主服侍人,.
因為人子尼不是要受人的服侍,.
乃是要服侍人,.
並且要捨命作多人的贖價,.
耶穌呼籲他的門徒,.
跟隨他的腳蹤,.
除非你真的認識你跟從的耶穌,.
否則你無法做一個真正的門徒,.
耶穌解釋給他們聽,.
他的權柄和榮耀,.
不是用操控的手段,.
乃是服侍,以致犧牲,.
所以耶穌叫我們,.
要將別人的需要,.
放在自己的耳邊,.
不是我作主,.
不是求自己的榮辱,.
不是用犧牲受苦操控別人,.
達到自己的目的,.
而是像耶穌一樣,.
有清楚的使命,.
還記得我們開始的時候說,.
耶穌堅定不移地去耶路撒冷這條路上,.
他有很清楚的目的,.

$^{481}$就是為世人付上罪的贖價,.
讓人脫離罪的綑綁,.
可以過豐盛的生命,.
得到人生真正的自由,.
同樣當我們步耶穌基督的腳踏的時候,.
我們也是和耶穌一樣,.
我們要將人帶到耶穌的根前,.
讓耶穌成為他生命的主宰,.
所以我經常提醒自己,.
我的生命是否能夠幫助人,.
向耶穌走近一步,.
我是否成為了人認識耶穌的難阻,.
如果當我將人帶到我面前,.
我去服侍我自己的時候,.
這就是阻礙人去認識耶穌,.
我提醒自己,.
我們要成為流通的管子,.
將人帶到神面前,.
我們沒有分彼此,.
我們都是以耶穌為我們的主,.
我們沒有誰比誰好,.
但我們的目的,.
就是讓我們有合一的見證,.
將人帶到神的根前,.
其實走僕人的路一點也不容易,.
走路本身也不容易,.
好像我昨晚,.
我坐巴士的時候,.
下車踏在凹凸的地方,.
所以扭傷了,.
今天急著回來,.
所以我們要走這條僕人的路,.
的確會有很多失敗,.
但我們同樣看到這班門徒,.
聽完耶穌的道之後,.
他們仍然不明白偉大的意義,.
但聖經告訴我們,.
當他們一起經歷見證耶穌基督,.
釘十架,復活,.
猜拳他們實踐大使命,.

$^{521}$將人帶到耶穌的根前之後,.
雅各和約翰的生命,.
正如第39節對他們的預言一樣,.
他們真的受杯和領洗,.
使徒行傳第12章第2節告訴我們,.
雅各為了福音的緣故,.
他為主殉道,.
約翰為了傳道的緣故,.
遭受到嚴酷的逼迫,.
晚年更被放逐到拔摩島,.
但他們最終成為了耶穌的跟從者,.
成為眾人的僕人..
在十九世紀有一個年輕人,.
他十九歲的時候,.
他很清楚他一生的道路,.
是要去中國宣教,.
於是他放棄了劍橋修讀法律課程的學位,.
他只差一個學期就可以畢業,.
但他卻去了神學院裝備,.
而且也去了馬利遜在倫敦特設的華語班學中文,.
因為他領受了一個人生的目標,.
就是全為基督,全為中國的好處..
於是在1827年,.
他23歲的時候,.
他帶著他的新婚妻子去馬來西亞,.
對當地的華人進行福音宣教的工作,.
其實他很想去中國,.
但當時清朝的政府還沒有對外開放,.
所以他唯有先去馬六甲服侍,.
但1942年,南京條約簽署了,.
於是他很開心,.
他終於可以去到中國的地圖裡,.
他去了廣州,.
不過他一去就沒有回頭,.
他一去到就患了重病,.
直到他想回去新加坡醫治,.
誰知道他在澳門病逝,.
他死的時候只有39歲..
他一生中有很大的貢獻,.
雖然他沒有去中國,.

$^{561}$但他一生中最大的貢獻就是用金屬字母.
去做一些活字的印刷,.
他去翻譯聖經,.
他去推動華人的子女教育,.
他的名言是「任何不能為中國寫生的意念,.
都會令我萬分沮喪,.
我唯一的心願就是為中國而活,為中國而死.」.
他說「如果我的心可以被否開,.
那個千絲萬縷交織在一起的每一根纖維,.
都寫著中國.」.
他亦說「讓我埋葬在中國,.
如此藉著我的埋葬之地,.
贏得中國歸屬基督.」.
這個人其實就是Samuel Dyer,.
他是台約爾,.
他是誰呢?.
他其實是戴德生的岳父,.
他有一本傳記,.
由他的元孫戴紹增所寫,.
「衰至於死,台約爾傳.」.
他一生為了中國,.
他一生很想去中國傳教,.
但是他沒有機會在中國裡,.
向任何一個人傳福音,.
但是他忠心的,.
上帝要他做什麼,.
差派他做什麼,.
我看完整本傳記,.
我看到一個謙卑,信服,.
神差遣他做什麼,.
他就做什麼的一個工人..
最後,台約爾在澳門的馬來禮信墓旁被埋葬,.
他的墓碑上寫著「Humble-minded」,.
一個很謙卑的人,.
他信奉神的帶領,.
縱然不能去到他心目中的中國,.
但是他去到南洋裡面獻身寫記,.
向華人傳福音..
所以弟兄姊妹,.
我們做門徒的,.

$^{601}$一定要像耶穌基督一樣,.
帶著使命,.
有一個僕人謙卑的心,.
千萬不要走進事奉的誤區,.
要自己作主,.
要自己得榮耀,.
甚至以為用受苦可以得到榮耀和權柄,.
反而我們要用束腰的手巾和洗腳盆,.
成為我們僕人的記號,.
我們要記住別人的需要,.
放在自己耳先,.
終身將人帶到耶穌基督的根前,.
讓人得著福音..
\newpage



\section{啟示錄 2:1-7}
\label{sec:Jhd45blKy7w}
\textbf{2025年10月5日崇拜直播｜陳明泉牧師｜耶穌給教會的七封信:起初的愛｜啟示錄2\_1-7}
\newline
\newline
連結: \href{https://youtube.com/watch?v=Jhd45blKy7w}{\texttt{https://youtube.com/watch?v=Jhd45blKy7w}} ~~~~ 語音日期: 2025-10-05
\newline
\newline
\hyperref[sec:JJ124JmNAMw]{\small{< < < PREV SERMON < < <}}
~
\hyperref[sec:index_chronic]{\small{[返順時目]}}
~
\hyperref[sec:mfLhrRj9Qu0]{\small{> > > NEXT SERMON > > >}}
\newline
\newline
啟示錄 2:1-7
\newline
\begin{longtable}{cl}
\hline
\hline
章節 & 經文 (和合本修訂版)\\
\hline
2:1 & \begin{tabularx}{0.7\textwidth}{X} 「你要寫信給以弗所教會的使者，說：『那右手拿著七顆星，在七個金燈臺中間行走的這樣說： \end{tabularx} \\ \\ \relax
2:2 & \begin{tabularx}{0.7\textwidth}{X} 我知道你的行為、勞碌、忍耐，也知道你不容忍惡人。你也曾察驗那自稱為使徒卻不是使徒的，看出他們是假的。 \end{tabularx} \\ \\ \relax
2:3 & \begin{tabularx}{0.7\textwidth}{X} 你能忍耐，曾為我的名勞苦而不困倦。 \end{tabularx} \\ \\ \relax
2:4 & \begin{tabularx}{0.7\textwidth}{X} 然而，有一件事我要責備你，就是你把起初的愛心拋棄了。 \end{tabularx} \\ \\ \relax
2:5 & \begin{tabularx}{0.7\textwidth}{X} 所以你要回想你是從哪裡墜落的，並且要悔改，做起初所做的工作。你若不悔改，我要到你那裡去，把你的燈臺從原處挪去。 \end{tabularx} \\ \\ \relax
2:6 & \begin{tabularx}{0.7\textwidth}{X} 然而你還有一件可取的事，就是你恨惡尼哥拉派的行為，這種行為也是我所恨惡的。 \end{tabularx} \\ \\ \relax
2:7 & \begin{tabularx}{0.7\textwidth}{X} 凡有耳朵的都應當聽聖靈向眾教會所說的話。得勝的，我必將神樂園中生命樹的果子賜給他吃。』」 \end{tabularx} \\ \\
[1ex]
\hline
\hline
\end{longtable}
$^{1}$《新約聖經》三百五十頁 啟示錄二章一至七節.
你們要寫信給已不所教會的使者說,那右手拿著七星,在七個金燈台中間行走的,說:.
我知道你的行為牢固忍耐,也知道你不能容隱惡人,你也曾試驗那自稱為使徒卻不是使徒的..
看出他們是假的來,你也能忍耐曾為我的名奴苦,並不罰捐..
然而,有一件事我要責備你,就是你把起初的愛心離棄了,.
所以,應當回想你是從那裡墜落的,並要悔改,行起初所行的事..
你若不悔改,我就臨到你那裡,把你的燈台從原處拿去..
然而,你還是有一件可取的事,就是你恨惡尼哥拉一黨的行為,這也是我所憎惡的..
聖靈向宗教會所講說的說話,凡有異的,就應當聽..
得勝的,我必將神樂園中生命樹的果子賜給他人..
弟兄姊妹平安,我們從十月份開始,教會最後一個季度,我們用啟示錄的七個書卷貫穿成為了我們最後一季,我們一起領受上帝話語的一個講道系列..
在我講這一個書卷之前,有一個人生片段,嚴格上來說不是我的人生片段,而是我太太的人生片段..
早一段日子,街市賣菜很潮的,有個DJ在那裡叫,不知道大家有沒有去過,很型的..
其中我太太和一個朋友在那些很潮的街市裡面,在街市裡面有菜心,寫著多少錢,旁邊加了幾個字,就是有初戀感覺的菜心..
真的這樣寫的,就是多少錢一斤之外,就要有初戀的感覺..
我太太和她的朋友看完,覺得很有趣,有初戀的感覺..
賣菜的那個DJ就看到有興趣了,我太太和她的朋友,就說試試吧,太太,買回去,吃完之後,保證你會恢復那種初戀的感覺..
我太太被推銷了,然後再買一斤,買夠兩斤,平時買一斤而已..
晚上我吃的時候,一邊吃菜心,我太太就問,有什麼味道呢?菜心就是菜心味,都是那樣東西..
我問她,你吃不出有初戀的感覺嗎?我說,初戀的感覺?怎樣吃我都吃不到..
我太太就說,不是啊,一邊吃的時候,我的子女就一邊吃,一邊笑,菜心也有初戀的感覺..
其實這個銷售技巧是很高的,因為每個人的初戀味道都不同..
總之你把你的感覺吃進去,那就是初戀感覺的菜心.不過對我來說,那個都是纖維質,都是菜心..
不過今天我們說的這段經文不是回復大家初戀的感覺,不過有一個字眼是重要的..
在剛才我們所讀的經文裡面,在啟示錄裡面,那個起初的愛,它的原文所寫的叫做first love,類似我們今天所說的初戀..
這個初戀不是說人與人的那種愛情生活,令人尷尬又回味的那種感覺..
這個是直指整個教會,他們生命裡面缺乏了的一種生命素質..
今天藉著這一段經文,希望我們不單止知道在主後的時候,特別是在耳筆鎖教會,身處一個面對著很多外邦宗教衝擊的社會裡面,知道一些文化歷史..
我們更加重要的就是藉著我們讀耳筆鎖,就是耶穌基督對他說話的這個教會,同樣都給我們今天有提醒..
今天我們會不會帶著上帝的愛,繼續活著我們的信仰,也將上帝的道帶到這個世界當中..
我們在這個時間很需要上帝幫助我們,我們先求上帝幫助同心祈禱..
求主你在我們當中帶領,讓我們一起領受你話語的時候,對我們每一個內心說話..
我們祈求你,你親自臨隔在我們當中,藉著你的話語,對我們這個世代發出最重要的呼聲..
願主你在此時此刻幫助宣講的講得清楚,幫助聆聽的每位弟兄姊妹,他們不單止明白,也在生活裡面實踐你對我們的教導..
我們下來的時間,願主你賜福,大靈,獻上禱告奉基督耶穌得勝的聖名..
我們一起看看啟示錄,當很多弟兄姊妹一看啟示錄的時候,很多時候都會被那些荷里活電影,或者坊間裡面講得很神秘的一些YouTube channel等等,就會吸引了我們的眼球..
坊間裡面很多人讀啟示錄會用一個方法,就是將現代的現象讀進去啟示錄裡面..
例如有一些象徵,他用了一些飛機大炮,或者有不同的那些意象,就是在講現在這個科技很恐怖了,讀進去啟示錄,我們就覺得啟示錄所啟示的末世就越來越近了..
我們不能夠否定末世越來越近,這是一個正常現象,耶穌來的日子我們沒有辦法知道,而且他會越來越近,因為隨著時間的流動,我們基本上是等待主再回來的日子..
但是我們很難將現實生活的現象讀進去,反而比較恰當的讀法,我們是將舊約裡面有一些意象,或者舊約裡面怎樣用過的字眼,或者有一些色經的重要內容,對當代可能是很有意思的,我們就要將這一層讀進去啟示錄裡面..

$^{41}$第二件事就是啟示錄有很多象徵的言辭,待會你會讀到的,那些象徵的表達,我們最好就看上文下理,看整本聖經裡面有什麼內容是這樣說..
當我們用聖經的處境,用當代的歷史文化去看的時候,可能跟我們今天去看啟示錄很神秘,飛機大炮那些東西,可能有一個截然不同的訊息,啟示錄的訊息其實很落地的,我們一起看..
我們來看一看我們接下來的二章一節到三章二十二節,是啟示錄裡面基督書簡,就是說耶穌基督對七間亞細亞教會所發出的信件..
在七間教會的信件裡面,每一間教會的格式,耶穌對教會所說的話,是有一個固定的格式的,大家在螢光幕上會看到..
第一句通常會有一個預令,你要寫信給哪一個教會和對象,他一定是這樣寫的,每一次每一封都是這樣寫的..
然後針對教會的處境,耶穌基督就會寫信,在啟示錄裡面會描述耶穌是一個怎樣的主,是一個怎樣的角度來描述下來所說的訊息..
所以耶穌的身份,或者怎樣描述和後來所說的圈面,篤澤,或者肯定,稱讚是互相相關的..
所以這個自我,耶穌的宣告是會針對教會的處境來講他的身份的..
接下來就是那個信的主體,特別是說有些教會是被耶穌基督稱讚的,但也有更加多是被耶穌基督責備的..
又或者又有稱讚又有責備,這個就是信裡面的主體..
當講到這些教會的現象之後,耶穌基督就會作出一些提醒,要怎樣去改變,要怎樣去面對這些困難,又或者做得好的,怎樣去繼續堅持著他們做得好的地方..
最後,當你有些悔改也好,或者堅持下去也好,上帝有一個藉著聖靈帶來的應許..
最後,結語就是聖靈向眾教會所說的話,凡有異的應當聽.這個就是那個結語..
所以你用這個角度來看,其實二章第一節到三章二十二節不是很複雜..
我們每一間教會就用這個框架來看裡面的內容..
我們先看第一間教會,就是二福所教會,啟示錄二章一至到第三節..
你要寫信給二福所教會的使者說,那右手拿著七星,在七個金燈台中間行走的說,.
我知道你的行為牢固忍耐,也知道你不能容隱惡人,你也曾試驗,那自稱為使徒,卻不是使徒的,看出他們是假的來..
你也能忍耐,曾為我的命勞苦,並不罰卷..
在二章一至到三節裡面,耶穌特別說到他是一個怎樣的主呢?.
他說他右手拿著七星,在七個金燈台中間行走..
首先定義一下這兩個字眼,七個金燈台,猶太人的照明系統就是用燈台,用蠟燭台..
因為舊時的燈台有七個插蠟燭的位置,所以就用這個意象來說一個照明,也就是說這個燈台就代表著一個發光的群體..
而金燈台在早一點的蝙蝠啟示錄裡面都已經出現過了..
在啟示錄一章十二至十六節裡面,使徒約翰特別說到他見到耶穌,耶穌的呈現是怎樣呢?.
十二節那裡這樣寫:「我轉過身來,要看是誰發聲與我說話,既轉過來就看見七個金燈台,燈台中間有一個好像人子,身穿白衣,直垂到腳,空間束著金帶,他的頭與髮皆白,如白羊毛如雪,眼目如火焰,腳好像在爐中鍛鍊光明的銅,聲音如同眾誰的聲音.」.
「他右手拿著七星,從他口中出來一把兩刃利劍,面貌如同烈日放光.」.
在這個意象或者一章十二至十六節,其實說的那個燈台七星,就是耶穌基督在這些象徵物裡面出現的..
而金燈台特別說到,其實在順民夏利裡面特別說到那個燈台就是教會,所以這個金燈台的意象是說到耶穌行走在七個教會中間,教會就是成為了世上的光,上帝藉著這個教會的群體,將上帝的光明帶到這個世界裡面..
所以燈台的意象是說教會,然後第二個重要的就是右手拿著七星,七星有機會是說七點的光芒,不過有一些說法更加是說七星其實是說天象,金木水火土,太陽和月亮,這七個星體..
代表著耶穌基督掌管著整個宇宙,當然有些色經學家說的七星都是可能七個星座,無論他用哪一個解說都好,你會看到拿著七星是說耶穌基督是掌握著整個宇宙的能力..
經文再繼續說,他不是直接說耶穌基督,他在說燈台中間有一個好像人子,其實這一句整個句式是來自大義理書十章五至六節,他說到呈現出來的掌管萬物的王就是用人子來稱呼他,而所有的意象在說掌管宇宙萬物,掌管教會的那位就是耶穌基督..
所以在那個順民夏利裡面,他就用剛剛說的那些意象,耶穌在第二章裡面就說到耶穌是一個超然的主,是一個滿有能力的主,他帶著的就是權柄能力,掌握著宇宙萬物..
更加重要的就是他洞悉了整個世上裡面每一個教會裡面的實況,所以他遊走在七個燈台中間,就是因為他知道教會發生什麼事..
他是用一個全知的角度,用一個超然審判的角度,來看整個以弗所教會.他用一個這樣的角度之後,經文裡面就說耶穌知道以弗所教會是一個行為,牢立,忍耐,不能容隱惡人,能夠辨別假仕途,並且不罰券的一間教會..
簡單來說,是一間很勤力,非常優質,有智慧的正統教會,無可頂的,一間教會能夠去到一個這樣的狀況,各方面都是很有水平的..
值得留意的是以弗所這個地方,以弗所這個地方在昔日裡面其實都是一個比較文化交匯的一個海港城市,所以他們那個地方都能夠匯集到相當多的人才在裡面,而因為有人才,所以那些信了主的基督徒就不是普通的基層的百姓,即或可能有一些基層,.
但是這個城市本身都很有水平,所以那個教會的模樣,那個教會群體本身都有相當實力,而且在智慧辨別,在捍衛真道上面,他們是很特別的一個教會..
所以以弗所教會在我們世俗眼中是一個非常令人羨慕,一個有實力的教會,不過一個這麼有實力的教會,耶穌基督對他講什麼話呢?.
二章第四節到第六節就是這段經文裡面的重點,在第四節裡面就是這麼說,然而我有一件事要責備你,就是你把起初的愛心離棄了,所以應當回想你是從哪裡墜落的,並要悔改行起初所行的事..

$^{81}$你若不悔改,我就臨到你那裡,把你的燈台從原處挪去,然而你還有一件可取的事,就是你恨惡尼哥拉一黨人的行為,這也是我所恨惡的..
在二章四至六節裡面,既然有一間這麼美好的教會,一間真理這麼正統的教會,但是他面對的就是耶穌嚴厲的責備,責備他什麼呢?.
就是他把起初的愛心離棄了,我們首先要了解一下,在和合本裡面,我們寫的愛心離棄了,都明白的了,大致沒有譯錯的了..
不過,和合修定本,他把離棄的字再更加具體地說,或者更加活地呈現出來,就是把你起初的愛心拋棄了..
就是說,你當一個夫婦關係也好,或者一個很緊密的關係,他現在把那些東西丟掉了..
耶穌基督在說,這個教會很勤力,做得很多正統的,又為了上帝的名,勞碌忍耐,不能夠容隱惡人,又能夠辨別假仕途..
但是,他丟掉了最初的那份愛,起初的愛心是什麼意思呢?有兩個意思的..
第一個意思就是,有些的悉經學家在說,你剛剛信主,或者已忽所的弟兄姊妹,剛剛認識主的時候,他和上帝那種緊密的關係,.
那個著重的是那個人由罪進入到光明,剛剛接受主的時候那份熱心,類似我們今天所說的初戀的那種感覺..
在這裡就呈現了,或者這個解釋就是回復舊時的激情,或者那種對上帝的毀身,那種不計較,和那種甜蜜的感覺..
這是第一個.第二個,那個起初的愛就不是說那一刻,或者那個很短的時間,而是那個起初是說一直以來,這班人信了主之後,都是延續著那份愛的,直到有些事情發生了..
所以後來裡面所說的,從哪裡墜落,他有一個墜落點的,就是說那些基督徒一直都是這麼愛主,直到有一些事情發生,.
直到他們就將這個跟神很緊密的關係,由對神的愛,引申為對人的愛,這一種態度,丟掉了..
直到這個教會就變成了一個只有真理,但是沒有愛心的教會.第二個解釋我比較同意的,因為他強調的不是回復那種初戀的感覺,.
而是強調的那種愛心的實踐,教會有真理,但是沒有了愛心.這個似乎是在順民夏利裡面,也是讀上去的時候比較容易通順的一個理解..
如果我們繼續去看的時候,你會知道,其實這一班人將那個愛心丟棄,嚴格上來說,其實是有歷史的因素的..
因為以弗所教會,或者以弗所這個地方,有很多的文化很豐富,有很多的假使徒,有很多五花八門的危機存在在社會當中..
這班基督徒,他們要站起來捍衛他們的信仰,要為他們的信仰,能夠將真實的基督耶穌告訴世俗,或者他們的城市知道..
久而久之,在一個變度,或者在一個捍衛的姿態裡面,人就變得好像有攻擊性一樣,以致他們就將重點放在這個層面裡面..
不過我們要留意,耶穌基督特別在馬太福音裡面說,他說,在末世的時候,只因不發的時增多,許多人的愛心就會漸漸冷淡了..
耶穌說,面對著這個世界,這個末世越來越多荒謬的事,越來越多令人嘆息的事,人的內心的承載能力可能已經麻木了,習慣了,以致很多人的愛心就麻木..
耶穌就說,門徒,你千萬不要變成一個麻木的人,這個世界是這樣,但是基督徒你不可以成為一個麻木的人..
耶穌基督在約翰福音裡面再繼續說,他說,教會的印記是什麼呢?你們若彼此相愛,眾人就認出你們是基督耶穌的門徒..
在約翰福音13章35節,簡單來說,在一個荒謬的末世裡面,辨明真理是重要的,不過不要因為辨明了真理而失去了上帝整本聖經貫穿的重要又重複又重複的訊息,就是愛心..
當我們能夠贏了一場架,咬贏了別人,但是我們的心靈是冷淡,是麻木的,我們都不能夠在這個世界裡面有力的告訴我們身邊的人,我們就是一群基督徒..
耶穌基督在這裡面對二福所教會作出一個這麼嚴正的責備,就是教會除了要做辯道的工作之外,不可以失去那份起初或者一直以來實踐的那種愛心實踐..
當教會沒有愛心的時候,人認不出這個就是基督耶穌的群體.可以想像的,二福所教會是一群什麼人呢?剛才說到是一些為上帝的道勞碌忍耐,後面第六節裡面更加說到為著尼哥拉一黨的人的行為,他們很恨惡的,所以他們受到上帝的肯定,.
因為這群人都是上帝所恨惡的,不過如果你從人的角度來看,二福所教會的弟兄姊妹是一群什麼人呢?你會面對他們的時候,他們經常板起嘴臉,你做得不對,這個不適合聖經,錯的,聖經不是這樣說的..
永遠人的世俗有什麼的時候,他就會板起嘴臉,恨惡,一種勢不兩立的那種情然在那些弟兄姊妹的生命當中流露出來,這個教會是一個很正統的教會,不過這個教會卻是沒有生命的影響力..
耶穌基督傳道的時候,最討厭的那些人就是法利賽人,大概如果這個教會只有這些的變度,沒有了愛心,就好像那些法利賽人一樣,他們有真理,但是實踐不到愛..
耶穌基督給了什麼出路給這教會呢?在經文裡面,他特別的說,在二章第五節,「所以應當回想你從那裡墜落的,並要悔改,行起初所行的事.你若不悔改,我就臨到你那裡,把你的燈台從遠處挪去.」.
二章五節裡面,其實有三個很重要的動詞,就是回想,而且悔改,再加上行起初所行的事,我用三個R來幫我們記得耶穌給所教會的出路..
第一個回想,Remember,回想在哪裡墜落,正如我剛才所說,有些解經學家回想信主的時候,人家怎樣接納你,回復信主的那種感覺,這是其中一個理解的方法..
不過如果對於一個整體教會來說,或者在他們的信仰當代的處境實踐裡面,我覺得回想起在哪裡一直都是這樣實踐,但是有沒有在他們的信仰實踐裡面,有一些關鍵的時刻出現了,以致他們一群人就變成一個集體冷漠的教會,找回這一點出來..
所以從哪裡墜落,哪裡有具體,是要回想什麼時候開始變成一群沒有愛心的人..
如果用我們今天,我們弟兄姊妹的生命來讀,可能我們過往都是一群很有愛心的人,但是可能有些事件出現了,有些狀況出現了,以致我們拋棄了這種愛心的實踐..
可能有些弟兄姊妹因為在她的信仰生命實踐裡面,曾經被人傷害過,所以恨成為了她的動力,她依然很努力去實踐她的信仰,依然事奉,不過每一次事奉都是帶著一份恨意,帶著一份不甘心,帶著一份敵對的態度,來看身邊服侍的人,或者看這個世界..
可能是有這樣的狀況,又或者有些人很想證明自己,我是對的,我是重要的,不過久而久之,在一個信仰實踐裡面,我們為了要讓別人看到我們自己,我們的信仰就開始扭曲,.
以致我們不是真誠地愛身邊的人,我們的實踐大部分都是建基於為了要證明自己的存在,又或者有時候我們被人誤會,反正做什麼好事都被人誤會,我們就跌入心牆,我們就索性不做一些愛心的實踐,又或者我們避免了自己受傷,.
我們索性成為一個軍人的心牆,我們用真理武裝自己,聖經的話語成為我們的武器,我們用這件事攻擊,或者為自己製造一個防禦網,以致自己的心靈不再受傷,.

$^{121}$真的,要愛別人很多時候都會帶著傷痛,你要學習實踐愛,愛背後很多時候都會帶來一些傷害,我不愛,我就不受傷,再加上用真理來包裝,我就可以成為一個,別人看來我們都是一個很得體的基督徒,.
但心底裡面,可能我們心裡面是盲目的,是絕緣體,我們的心不再將這份愛心來實踐,耶穌基督提醒我們,我們要回想,會不會在真理一路走,漸行漸遠的時候,我們和上帝的關係都漸行漸遠,.
以致我們失去了這份起初的愛,經文第二個字眼就是悔改,repent,悔改的意思就是說,以前一直走這條道路,當我們知道了,回想知道某一些狀況出現的時候,我們就不再跟著這個方向走,我們就改換方向,我們走回上帝要我們所走的路,.
這個悔改其實需要很大的覺醒和勇氣,因為一直都是這樣走,忽然轉換跑道,轉換方向是一件很不容易的事,不過耶穌基督說,如果我們不悔改的時候,耶穌就會臨到我們當中,將燈台從遠處下去,.
簡單來說,如果我們不能夠實踐這個悔改的生命,不給一些意志來克服的時候,上帝會將我們生命的影響力,即使我們能夠做很多努力,但是我們所做的都是對生命帶來不了影響力的功夫,我們生命就好像白酒一趟一樣,.
所以這個悔改我們要立定心智,當我們知道有些東西在我們生命出現的時候,我們正視它,我們不要再走我們一直走錯的舊路,轉換我們的方向來實踐,最後,轉換了之後做些什麼呢?.
其實那個真理,那個愛心不是跟你很遠的,舊時你正在做的,你經歷到那種風盛,舊時一直在走的那個帶著愛心去做的事奉道路,你做回那些東西就可以了,不是學一些新的東西,而是做一些舊時你很經歷到上帝同在,因著愛神,愛人的那些方式,你再做回,.
所以行起初所行的事,re-do,是做一些舊時你經歷到因愛而來的動力,來做回那些事,你就行在上帝的旨意當中,可能有些頂梓梅說,牧師,我知道的,不過有些像生命裡面有些傷,出現了就出現了,那份感覺我拿不回,.
好像老夫老妻那樣經歷了幾十年之後回復初戀的感覺,去不回,其實我的意思不是叫你拿回那種感覺,不過你拿回你們一直以來怎樣互相相愛的那些,總有一些日子是一起同行的,拿回你舊時實踐的那種生命的體驗,將那些體驗放回今天當中,.
我很喜歡有一本書,那本書叫做愛的藝術,The Art of Loving,是Loving,不是Love,它說的是持續去愛的藝術,這本書有一個心理學家叫做Eric Thorm,1956年寫的,.
他在說愛是一回什麼事呢?愛不是一種感覺,愛無論你對親人,對朋友,對世界所有的愛,其實愛是一種決定,當你決定去愛的時候,你就會慢慢懂得去愛,而且這個愛是一種技術,是一種心領神會的一種體會,.
所以背後有一些主觀的成份,但更多的就是可以學習,可以回顧的,你回顧你昔日怎樣去愛身邊的人,或者你生命裡面怎樣被愛,將這份體驗不是抽象的,而是具體的,實踐在生活當中,.
你就能夠重新實踐到耶穌基督所說的起初的愛心,所以在這個裡面,其實起初的愛重點是我們有沒有決心,我們是否決定這樣去做,當你這樣去決定的時候,很多的回憶,很多的經歷就會回來,Redo你生命裡面所經歷的,.
經文裡面再繼續說二章五節,你若不悔改,我就臨到你那裡,把你的燈台從遠處挪去,這是耶穌基督對那些已忽所教會的門徒作出的警告,如果我們生命裡面失去了愛,即使做其他所有的事情,正如保羅所說的,就好像明的不響的螺一樣,只是聽到聲音,但沒有實質的,.
我們的生命失去了影響力,因為我們缺乏愛,不過耶穌繼續給一個另外的應許,他說聖靈向宗教會所說的話,凡有意的就應當聽,得性的我必將神樂園中的生命樹果子賜給他人,.
生命樹果子其實是回應在創世記裡面,人犯罪墮落的時候,上帝將人趕出伊甸園,就找基督拍封鎖著伊甸園的大門,神和人之間就開始有一個隔絕,而當人去實踐愛的時候,就好像回復去伊甸園的那種生活,.
生命樹的果子就再成為人的食物,簡單來說就是回復那種連結的生活,那種連結神和人之間的生活,那種重新連結人與人之間的生活,如果我們去愛,我們心裡面就能夠對上帝的心牆更加深刻的體會,.
愛成為我們生命裡面最大的靈修,如果我們今天要靈修,第一步要做的就是學習去愛,學習不是學新的東西,只是回復舊時你所經歷,生命裡面你怎樣去愛身邊的人,又怎樣讓身邊的人愛你,.
這個就是今天這段經文的訊息,有一個簡單的體會,近來很多機會去到醫院探望一些老人家,探望一些不同的弟兄姊妹病患,我看到一個很有趣的景象,病的人身體不舒服,一定是有些脾氣,有些不安樂,.
可能有些老人家身體離流,沒力氣,身邊的人未必很明白他的需要,我在病床的病人旁邊總會有一些家人,可能是他的配偶,或者他的兒女,或者很疼他的家人,那些家人就默默在病床旁邊陪伴著病患者,.
我去到其實很短時間,我去其實是做一下祈禱,為他們帶來一兩節的聖經,安慰一下,陪一下他們,不過我更加欣賞的就是那些家人,其實因為他很愛他,所以長時間陪伴,有時候都不知道做什麼好,所以有時候摸摸角角,又擦擦手,擦擦嘴,按摩一下,其實對那個病患者來說,那個病患者有時候不舒服,有時候會發家人的脾氣,.
不過我看到那些在病患者的家人那裡陪伴,無論那些家人怎樣發脾氣都好,他們都會忍氣吞聲,一聲就吞了他,覺得不舒服,有時候會皺眉頭,對那些家人,但家人都繼續帶著眼淚,帶著那份最大的誠意,來陪伴他,那個可能是生命最後的一刻,.
我的體會是,有時候人生去到某些狀況,那種愛心可能真的摸摸角角,有時候真的不知道怎樣做,不過他們的在這裡,無論你家人什麼處境,那些家人總在這裡,永遠不離不棄,那個就是愛心..
鼎姐妹,今天我們都依然活著行走,我們生命很有氣息,我們今天可能做了很多事情,但是我們對於這個世界,對於我們身邊的人,對於鼎姐妹,我們會不會都會有這份耐性,有這份心,與他們繼續來同行,.
我們有真理,有愛心,這個就成為了今天這個世代,基督徒重要的標誌,求主繼續來幫助我們,帶領我們走面前人生的路..
\newpage



\section{啟示錄 2:8-11}
\label{sec:mfLhrRj9Qu0}
\textbf{2025年10月12日崇拜直播｜陳冠成牧師｜耶穌給教會的七封信:至死忠心｜啟示錄2\_8-11}
\newline
\newline
連結: \href{https://youtube.com/watch?v=mfLhrRj9Qu0}{\texttt{https://youtube.com/watch?v=mfLhrRj9Qu0}} ~~~~ 語音日期: 2025-10-12
\newline
\newline
\hyperref[sec:Jhd45blKy7w]{\small{< < < PREV SERMON < < <}}
~
\hyperref[sec:index_chronic]{\small{[返順時目]}}
~
\hyperref[sec:980zAvvXkiM]{\small{> > > NEXT SERMON > > >}}
\newline
\newline
啟示錄 2:8-11
\newline
\begin{longtable}{cl}
\hline
\hline
章節 & 經文 (和合本修訂版)\\
\hline
2:8 & \begin{tabularx}{0.7\textwidth}{X} 「你要寫信給士每拿教會的使者，說：『那首先的、末後的，死過又活了的這樣說： \end{tabularx} \\ \\ \relax
2:9 & \begin{tabularx}{0.7\textwidth}{X} 我知道你的患難和貧窮—其實你卻是富足的，也知道那自稱是猶太人的所說毀謗的話，其實他們不是猶太人，而是撒但會堂的人。 \end{tabularx} \\ \\ \relax
2:10 & \begin{tabularx}{0.7\textwidth}{X} 你將要受的苦，你不用怕。看哪！魔鬼要把你們中間幾個人下在監裡，使你們受考驗，你們要遭受苦難十日。你務要至死忠心，我就賜給你那生命的冠冕。 \end{tabularx} \\ \\ \relax
2:11 & \begin{tabularx}{0.7\textwidth}{X} 凡有耳朵的都應當聽聖靈向眾教會所說的話。得勝的必不受第二次死的害。』」 \end{tabularx} \\ \\
[1ex]
\hline
\hline
\end{longtable}
$^{1}$今日經訓記載於啟示錄第二章八至十一節.
在唐用聖經新約三百五十頁.
請聽我讀出.
你要寫信給是美那教會的使者,說那首先的,末後的,死過又活的說.我知道你的患難,你的貧窮,也知道那自稱是猶太人所說的偉旁話..
其實他們不是猶太人,乃是撒旦一會的人.你將要受苦,受的苦你不再用怕.魔鬼要把你們中間幾個人下在監裡,叫你們被試煉,你們必受患難十日..
你勿要致死忠心,我就賜給你那生命的冠冕.聖靈向宗教會所說的話,凡有已的就應當停..
得勝的必不受第二次死的害.經文讀不,願神賜福常讀他話語並且遵行的人..
各位姐妹,平安.在講經文之前,不知道有沒有人去過土耳其,應該有不少吧?.
考考你們,知道土耳其三大城市是哪三個嗎?第一,伊斯坦布爾,土耳其現今人口最多,古代曾經是拜占庭帝國和鄂圖曼帝國的首都..
第二,安卡拉,土耳其現時的首都,第二大城市,可以說是政治和行政中心..
第三,伊茲密爾,土耳其工商業的重鎮,我們不熟悉這個名字,但我們熟悉它在古代的名字,就是剛才讀出經文的士美娜..
事實上,在新的時代,士美娜已經是一個繁華的城市,不過剛才玉峰弟已經讀了一段經文,士美娜教會就不太好了..
你會看到經文描述他們遭受患難,貧窮,被人毀謗,甚至有些信徒被帶到監牢裡,還要面對死亡的威嚇..
我們不禁會問,明明士美娜這個地方是欣欣向榮的,但為什麼在裡面的士美娜教會這麼多災多難?.
我們要先了解當時的背景.士美娜繁盛是因為他們擁有一個水深的港口..
在耶穌時代,這是一個很得天獨厚的地理環境,因為在古代世界沒有飛機..
要運送大量的貨物,只能靠兩個方法,一個就是陸路,所以要靠很寬闊的路,可以讓馬車通過..
所以你知道羅馬帝國有一個路,我們常說的跳條大路,是通過羅馬..
真相是,古羅馬帝國之所以強盛,是因為他們當時已經擁有很多很發達的道路網絡..
當然不靠陸路又如何?剛才大家也說到,就是水路..
靠水路,一些大的貨船,你自然會選擇哪些城市來卸貨?.
一定是有些良好的港口,那些地方來卸貨的,然後貨物從那些港口卸了,再由陸路轉運到其他城市..
所謂「近廚得食」,你自然想到不同地方的物資,自然會選擇士美娜,這個擁有良好港口的城市來卸貨..
也因為這個緣故,吸引了不同的人,來到這個城市做各種生意和貿易..
亦帶旺了士美娜這個城市,成為當時羅馬帝國亞西亞省一個很繁榮的城市..
而另一方面,這個城市之所以繁榮,亦有賴於它與當時羅馬政府的關係很友好..
士美娜城一向以效忠羅馬見稱,有些文獻亦記載,這個城市被譽為是羅馬最忠心最久遠的盟友..
事實上在耶穌時代,士美娜這個城市就與周邊其他六個城市爭奪,為當時的羅馬帝王凱撒興建一座神廟的權利..
結果當然是士美娜城市贏了,他們成功爭取到在當地興建一座敬拜凱撒的廟宇..
一方面用來感激當時羅馬帝國帶給他們的繁榮安定,而另一方面亦同時對帝國代表一個忠誠..
事實上,如果你比較熟悉當時教會的歷史,你會知道敬拜凱撒,敬拜君王,其實不是純粹一個宗教的活動..
裡面夾雜著很多與文化,生活甚至經濟息息相關的影響,可以說彼此之間形成了一種互相利益輸送的網絡..
換句話來說,敬拜凱撒其實會為生活帶來很多的方便甚至好處..
相反,不參與敬拜凱撒的行動的群體就會自然會在生活上面對很多的困難..
就像經文所說的..
第九節顯示,當時的善美娜教會面對的是貧窮..
如果有了剛才的歷史背景,你會知道,大概我們不難想像,是因為他們不願意參與敬拜凱撒的行動..
特別是當嚴苛的羅馬帝王執政的時候,他們會立法規定每個人每年都要到凱撒的神廟裡敬拜燒香,並宣誓凱撒才是他們的主人..
完成之後,政府會給予他們一張證書,像發一張良民證一樣,證明這個人已經向凱撒表忠了..
他是奉公守法的良好公民..

$^{41}$但是對善美娜教會來說,因為他們要堅守耶穌是主的信仰原則,所以他們一定會拒絕敬拜當時的羅馬君王凱撒..
所以,外邦人會視教會是一種什麼樣的群體呢?.
是一些不合群不合作的,甚至覺得他們是破壞社會秩序的一群..
所以外邦人會對基督徒進行排斥,不會跟他們做生意,不會聘請他們打工,導致當時的信徒會被邊緣化,他們受惠不了當時社會的經濟繁榮..
再加上善美娜教會不會有一張敬拜凱撒的證書,沒有那張良民證..
他們敢不敢隨便出街?.
不會的,因為很容易會被人舉報,甚至送官救治..
亦因為這樣,他們不敢貿然上街,更加不敢去找工作,甚至會因為逃避被追捕,失去了工作,失去了生計..
所以,就像經文所說,在物質上,在經濟上,漸漸變得貧窮..
不過,在耶穌基督的眼中,不是這樣看,卻是另一幅圖畫..
經文說,善美娜教會其實是富足的,是因為他們堅守信仰,效忠耶穌的緣故,所以才失去了地上的財富..
這一切,耶穌保證,最終會轉化成為天上永存的財寶..
所以,耶穌說,雖然在經濟上,在物質上,他們是缺乏貧窮,但在屬靈上,他們卻是富有的..
而另一方面,我們看到,善美娜教會的患難,部分也是來自當時猶太人對他們的毀謗..
很可能是因為善美娜教會在當地開始興起,人數慢慢增加的時候,甚至當時有很多信仰猶太教的信徒改信基督教..
因為這樣的緣故,削弱了猶太教在當地的地位和影響力,導致他們很嫉妒,甚至憎恨基督徒..
在二世紀,也有文獻曾經顯示,善美娜當地的猶太人非常敵視基督徒,甚至會聯同政府官員或群眾來挑撥排擠當時信耶穌的人..
更嚴重的是,會鼓勵他們一起迫害教會..
耶穌指責這些猶太人,實際上並非猶太人,意思是他們不是上帝的會眾,不是上帝的子民,而是撒旦的會眾..
他們為了鞏固自己在地區的勢力和利益,竟然聯同當時的羅馬政府迫害真正的上帝子民,也就是基督徒..
特別是那些脫離猶太教,改信基督教的人,被迫害得更嚴厲..
猶太人這樣做,在耶穌的眼中,就變相等同於由上帝的陣營,轉頭向魔鬼撒旦的陣營..
而另一方面,我們不難想像,為什麼猶太人在當時沒有受到羅馬政府的迫害呢?.
很大機會是因為他們在敬拜凱撒的事情上,做了一些妥協..
為了換取當地政府對他們的信任和接納,這些猶太人不惜背離上帝,宣認凱撒才是他們的主人..
所以耶穌這麼嚴厲斥責他們是屬於撒旦的一會,耶穌也暗示上帝必定會追討他們這種罪..
不過,你會留意到,耶穌似乎沒有打算在現在解決他們所面對的困難..
你看到耶穌基督是從不同的方面去勉勵甚至幫助教會,你留在苦難中吧..
我不先解決你,你先進入苦難中.耶穌在那裡等待他們..
首先,耶穌說了兩次,「我知道你的貧窮,也知道有人誹謗你.」.
表示耶穌完全理解教會當時的狀況,知道他們遭遇患難不是因為做錯了事,.
而是因為你做對了,所以面對迫不得已..
你堅守著信仰,所以你就遭遇到患難和貧窮..
耶穌的這個知道不單止在頭腦上,其實也在情感上,他很體會到善妹娜教會當時面對的艱難..
因為我們都明白耶穌在世之後,他也是這樣面對..
他也是面對身邊的人的反對和迫害,所以他能夠用我們今天的術語,.
他可以同理心到善妹娜教會的情況..
就正如我們在生活上也有經驗告訴我們,當你面對逆境困難的時候,.
不知道該怎麼辦,解決不了的時候,.
但假如身邊有一個人能夠切身處地跟你說,「我明白你的,我懂你的情況.」.

$^{81}$我能夠理解你,體恤你的時候,你會發覺縱使問題仍然未解決,.
甚至乎你見不到出路,但你內心總是會覺得安穩一點..
因為你發覺你不孤單,有人在你身邊跟你同行,安慰你,.
你會多了一份安穩和盼望,你會更加有一個韌力來應對當前的困難..
更何況經文這裡說,體恤,理解我們的是耶穌基督..
事實上,耶穌基督不單知道和體諒善妹娜教會的困境,.
其實他更能夠掌管,甚至勝過這一切的苦難..
如果你記得上星期,陳牧師跟我們說,.
耶穌基督給每一個教會的信都一定有一個自我宣告..
你看到現在耶穌基督對善妹娜教會的自我宣告,.
同樣是度身訂造,針對他們的處境..
耶穌基督對教會說,「我是首先的,末後的,死過又活的.」.
這顯示什麼呢?其實單單這幾個字,代表誰才是掌管歷史?.
是耶穌基督,他才是歷史的主宰,不是當時的政權,不是羅馬帝國..
他掌管的是過去,現在和將來,他也是那位的主人..
並且這裡說,他是從死裡復活,.
他是得勝最可怕的權勢,死亡的權勢的一位主,他是那位永活的主..
換句話來說,善妹娜教會的患難仍然在誰掌握手裡?.
在耶穌的手裡..
他最終一定會背亂反正,為教會,為聖徒伸冤,主持公道..
這是耶穌基督對當時的善妹娜教會所作出的保證,安慰和肯定..
所以,你會看到第十節,耶穌面臨當時受苦的教會,.
你不要怕將來所受的苦難..
耶穌這句話其實可圈可點,就是說苦難完結了嗎?.
還未完,還未去到最精彩的部分..
其實他向信徒預告,受苦是必然的..
並且預告了稍後受苦的情況..
經文說,猶太人會聯同當時的政府一起去迫迫基督徒,淪為了殺蛋的工具..
甚至經文說,將一些善妹娜的信徒帶到監牢裡..
你不難想像,帶到監牢,不只是等候審訊,.
在當時的情況,基督徒被帶到監牢,很大機會是等死..
而且你不知道會怎樣死,因為有不同殘忍的死法..
可能是釘你上十字架,可能是燒死你,可能是帶你到羅馬鬥獸場..
在這個時候,那種恐怖絕望的感覺,你害怕嗎?.
其實很害怕,即刻殺了你,你不會那麼害怕,你會很痛快..
但是等死,不知道什麼時候死的時候,是很令人擔心,很令人覺得恐怖..
亦因為這樣,會動搖信徒對信仰的堅持,甚至會引誘他們妥協..
你放棄,不認耶穌為主,拜一拜該殺,那不就沒事,得到釋放..
所以在這個情況下,耶穌基督就預先,我們說替善妹娜的信徒,.
先打了底,先做好心理準備,將來是會這樣..

$^{121}$背後的原因就是幫助他們,背後的目的就是幫助他們,.
你要有勇氣承受到時所面對的苦難..
而另一方面,耶穌亦提醒善妹娜教會,表面上是受迫害,是一件很負面很淒慘的事..
但耶穌卻叫他們從另一個角度看整件事,其實是被試煉..
換句話來說,上帝容許苦難逼伯臨到,是要考驗教會的忠心,.
以及信徒對耶穌的信靠..
並且,你看到耶穌這裡說了要試煉多久?十天..
耶穌表示,祂會限制教會受患難的日子..
當然我們明白,十天並不是十天,一個真實的數目,而是一段有限的日子..
十天意味著,沒錯,你所受的患難是很難受的,是不容易去面對,.
但你所受的,都仍然在耶穌的掌握裡,在上帝的保護裡..
目的是為了要煉成每一個信徒的屬靈生命..
最近我在YouTube裡面,有時在滾滾手機,看到一條YouTube影片,.
關於日本的刀匠,他怎樣鑄造日本的刀,.
怎樣將一塊像拳頭般大的鋼塊鍛煉成一塊很薄,很亮麗,很鋒利的刀..
過程中,整段影片三十多分鐘,我只看到三個動作,重重覆覆三個動作,.
你知道是什麼嗎?要鑄造刀的三個動作,就是燒它,打它,磨它,.
你很明白,燒,打,磨,那塊鐵你會很慘嗎?.
不停被燒,不停被打,不停被磨,很辛苦的..
但是正因為這種辛苦,裡面的雜質才可以拿走,沒有其他辦法..
唯有這樣,你最終才能成為刀匠心目中那把合用的刀..
同樣,善妹娜教會就好像那塊用來鑄刀的鋼石,.
每天都被上帝不斷地燒,打,磨,不好受的,.
但是從另一個角度看,是上帝為他們預備了一個成長的契機,.
而那些忠心忍耐,自死都跟隨耶穌的門徒,最終一定會通過試煉,.
不單見證到耶穌基督是真實,可靠,值得他們效忠的主,.
並且他們會得到耶穌應許要賞賜他們生命的冠冕..
意思是什麼呢?其實,確實在世人的眼中,.
基督徒這樣為主受苦,甚至殉道,是很愚蠢的..
為什麼這麼傻?妥協就行了..
是一種很失敗,甚至很羞辱的表現..
但是耶穌卻會將這種人看為失敗,羞恥的體會轉化成什麼呢?.
你是得勝的,你是榮耀的,所以你得到永恆生命的冠冕,.
背後就是你用地上的生命來換取..
並且耶穌應許當他們再來審判傳遞的時候,.
那些致死忠心的信徒,他們不會經歷第二次的死亡,.
他們不需要面對永遠與上帝隔絕的真正最大的痛苦..
換句話來說,原來在耶穌基督的掌管下,.
沒有任何的權勢,包括肉身死亡的威嚇,能夠跨勝到最後..
最終贏的是誰?是耶穌基督..

$^{161}$勝利歸於他,以及那些忠心跟隨他的人..
雖然你會看到整段經文說完了,只有四節,.
耶穌基督在啟示錄寫給善妹娜教會的短短四節經文,.
卻深深影響著當中每一個信徒..
大概在主後一百年左右,有一位叫做玻璃甲的信徒,.
他是使徒約翰的門生,亦是當時被委任為善妹娜教會的監督,.
即是今天我們稱之為主教的身份..
當時的羅馬帝王逼迫教會很厲害,.
玻璃甲寫了一封信給以弗所教會,.
勉勵他們,既然耶穌基督為了我們罪的緣故一直忍耐,.
直至到他們面對死亡..
所以現在就算我們面對掌權者的迫害,.
我們仍然不能退縮,.
反倒要為那些憎恨我們,甚至迫迫我們的人禱告..
玻璃甲致力提醒當時教會的信徒,.
你要效法耶穌,愛你的受敵..
而到了主後一百五十六年左右,當時的玻璃甲已經八十六歲,.
最後他也被捕,.
他被帶到當時的地方官員面前..
事實上當時的官員不打算將他處死,.
因為他已經八十六歲,年紀已經很大了..
於是就遊說他,.
「你講一句該殺是主,向他燒一燒香,.
那你就沒事了,有什麼問題?」.
玻璃甲的回答是,「我服侍了耶穌基督八十六年,.
祂沒有虧待過我,.
我怎能在這個時候褻瀆那位拯救我的主.」.
玻璃甲堅持堅持..
這班官員見遊說不成功,於是就恐嚇他,.
他們就恐嚇玻璃甲,.
「我們在鬥獸場有很多猛獸,.
你不改變主意,.
你不肯跪一跪,就將你扔進鬥獸場.」.
玻璃甲卻很安靜地回答,.
「你叫些猛獸來.」.
見官員見威嚇都沒有用,於是就說,.
「你不怕猛獸嗎?好啊,.
我們將你燒成灰燼.」.
玻璃甲就說,.
「你用來威嚇我的火,.

$^{201}$只不過燒個多小時,.
就會熄滅了,.
但你不知道將來要淋到一切不信者的刑罰,.
那個永遠焚燒的火,.
是多麼可怕..
來吧,不要延遲,.
照你們所做的燒死我.」.
結果最後,.
玻璃甲真的在火營裡為主殉道,.
他將耶穌基督在啟示錄對士美娜教會的提醒和吩咐,.
就是「務要忠心至死,.
必得著主所賜生命冠冕,.
有耳可聽就應當聽的教導,.
牢牢記在心裡,.
並且實踐到最後.」.
他的殉道也成為後世基督徒的激勵和幫助,.
當附近其他城市的教會都相繼式微,.
甚至乎他們的土地漸漸落入其他宗教的手上的時候,.
上帝卻保守著士美娜教會繼續在一個艱難的環境裡默默成長..
教會歷史告訴我們,.
往後的十幾個世紀,.
士美娜這個地方更是當時的東正教會,.
一個極之重要的信仰和學術中心..
所以弟兄姊妹來到這裡,.
有一些反省應用很值得我們思考..
習慣享受經濟繁榮的人,.
自然最擔心的是什麼?.
就像這段經文開始所說,.
就是窮,.
就是怕周圍的人很富庶,.
你卻是凋零樂索..
人之常情,.
每一個人總是想改善生活,.
想過得好一點,舒服一點..
不過經文提醒我們作為信徒,.
我們卻是有一些屬靈原則必須要遵守..
甚至有時為了遵守這些屬靈原則,.
你不可以不捨棄屬世的好處..
如果我們只是追求在社會裡面面俱圓,.
不得罪人,在世通達,.

$^{241}$嚴固聖經的教導,.
耶穌的吩咐,.
這樣不單止我們心中會愧對耶穌基督對我們的救恩,.
上帝也不會喜悅..
事實上,擁有豐富的物質生活,.
因為這一切都是出於祂的恩典和供應..
不過問題是,.
當你和我習慣了富裕,.
習慣了安逸的時候,.
會不會視為這一切是理所當然?.
能不能回頭?.
弟兄姊妹,我們會不會只認為上帝只會讓我們過得豐富,.
過得安逸?.
我們不再相信上帝原來也會為我們安排患難和貧窮的難度?.
原來為了我們的益處,.
耶穌基督不單止會供應安逸和豐富的生活,.
他也會安排苦難和貧窮給我們,.
很值得我們去思考..
事實上,在中國人的文化裡,.
總是充斥著富裕等於幸福,.
越多的財富,越多的物質,.
就意味著我們的生活越幸福..
所以,當面對生活中很多不滿意,.
不合自己的心意的時候,.
我們就會歸咎於窮,.
因為物質缺乏,.
於是我們就繼續向著賺多一點,.
向著這個方向走..
但這段經文提醒我們,.
原來貧窮缺乏往往是我們最親近上帝的時候,.
往往是我們最親近上帝的時候,.
當你經歷患難的時候,.
會比較親近上帝,.
還是經歷富有的時候,.
會比較親近上帝?.
很明顯是患難,艱難的日子,.
當你知道沒有上帝是不行的時候..
反過來,物質的富足,.
似乎是一種生活的幸福,.
但在經文裡說,.

$^{281}$它會成為一種試探,.
甚至和患,.
就像經文描述的猶太人那樣,.
因為要守住自己的利益,.
導致忘記了上帝,.
向世俗妥協..
事實上,.
善美娜教會和整個城市,.
他們都要面對同一個試驗,.
就是到底什麼為之貧窮,.
什麼為之富足,.
什麼為之人生的得和失..
人所賺得的,.
如果直接或間接助長了邪惡勢力的彰顯,.
將狂,.
使人受虧損甚至傷害,.
就算你再多賺,.
其實在耶穌的眼中,.
都是貧窮..
相反,人所失去的,.
如果能榮耀上帝,.
造就人的生命,.
讓人更加倚靠上帝,.
即使他再有損失,.
在耶穌基督的眼中,.
他反而是富足..
弟兄姊妹,.
經文迫使我們思考,.
到底我們今天很努力地經營,.
拼搏的,.
在耶穌的眼中,.
是得還是失,.
是貧窮還是富裕,.
又或者,.
我們有沒有將地上的財富,.
轉化為天上永恆的財寶..
而另一方面,.
你會看到,.
耶穌基督只是提醒,.
勉勵善會拿教會,.

$^{321}$前面還有更大的試煉在你面前..
耶穌沒有打算,.
要幫他們挪開,.
你不難想像,.
如果你是善會拿的信徒,.
面對當時的處境,.
甚至被迫進入絕境的時候,.
你大概都會向耶穌呼求,.
向上帝發出求救信號,.
主啊,還不夠慘嗎?.
要到什麼時候完結?.
每一個人都會問,.
或者說為什麼是我?.
他們得到的答案,.
可能也是歷世歷代信徒,.
包括我們得到的答案..
耶穌這裡說,.
你無要至死忠心,.
我就賜給你生命的冠冕..
至死忠心好像很負面,.
你試試反過來讀,.
是忠心至死,.
你對耶穌基督的忠心,.
堅持聖經的教導,.
最終走上極致的,.
就是殉道,.
在當時的處境,.
就是至死..
意思是什麼呢?.
耶穌會自有安排,.
你相信祂,.
繼續忍耐,.
堅持下去,.
不要想著幫你解決困難..
事實上,.
啟示錄一再強調,.
信徒為主受苦的真實性,.
和必然性,.
祂沒有將苦難浪漫化,.
而是告訴我們一個正確應對的方法,.

$^{361}$就是不要那麼容易妥協,.
甚至逃走,.
也不要嘗試祈求上帝,.
那麼快幫你拿走的苦難,.
主權在上帝那裡,.
可能會,.
但也可能不會..
耶穌這裡說,.
要忠心忍耐,.
和祂一起面對苦難,.
因為祂在那裡等我們,.
並且要認定,.
原來苦難裡面,.
是含有上帝美善的旨意,.
過程裡面,.
不可以缺少的,.
其實就像今天這段經文說,.
你不要只看著苦難,.
你看著耶穌,.
你要不斷用耶穌的生命,.
重新詮釋你所經歷的問題,.
特別是那些艱難的日子,.
耶穌說,.
祂就是那個首先末後,.
死過又活過的,.
其實這就是耶穌基督的生命,.
也是耶穌基督的宏大的故事,.
祂要求我們,.
要求是美那的信徒,.
以祂的掌管,.
以祂的德性和永活,.
以祂的故事和生命,.
重新詮釋你遭遇什麼,.
學習在苦難裡面,.
如何去提煉那些有意義的屬靈養分,.
事實上你會看到,.
同一個的境況裡面,.
同一個的艱難裡面,.
有人總是會將這些困難,.
敘述為一個很負面,.

$^{401}$很黑暗的故事,.
以致到很多的埋怨,.
很多的不滿,.
很多的煩躁,.
總是走不出那個困局,.
但是相反,.
屬靈人卻是能夠從這些艱難苦難裡面,.
抽取到正面的屬靈養分,.
他可以將逆境,.
轉化為敘述經歷上帝奇妙帶領的故事,.
關鍵在於我們敘述人生故事的能力,.
我們能不能夠不只是看自己,.
看環境,.
看世界,.
而是看聖經的真理,.
如何按照耶穌的教導,.
按照耶穌的生命,.
重新譜寫你的遭遇,.
所以弟兄姊妹經文提醒我們,.
不要只聚焦人的故事,.
而是聚焦自己的故事,.
別人對我有多差,.
他有多壞,.
我有多慘,.
聖經提醒我們,.
要懂得在這些人生高低起跌的裡面,.
你要看到更宏大的東西,.
看到更宏大的故事,.
就是聖經和上帝的故事,.
這種視野可以幫助我們,.
在苦難的裡面,.
在失意的裡面,.
甚至在混亂的裡面,.
你找回你的根本,.
你作為耶穌基督的信徒,.
你的本質是什麼?.
你的使命是什麼?.
以及你和上帝和耶穌的關係是什麼?.
並且認定原來,.
我經歷的問題,.

$^{441}$我經歷的困難,.
不是我人生的全部,.
我人生的全部,.
在耶穌基督的身上,.
祂的生命,.
成為我生命的道路和真理,.
唯有耶穌基督,.
才可以定義我們的日子是怎樣..
所以弟兄姊妹盼望這段經文,.
再一次提醒我們,.
特別當你面對一些困難,.
人生逆境的時候,.
你很想很快解決,.
特別如果你是一些很心急的人,.
很想你的生活井井有條的人,.
這個本質上不是錯的,.
但你會發現原來信仰也包含,.
你混亂失衡的時候,.
重點是你知道耶穌祂知道,.
祂在並且與你同行的時候,.
祂就能夠成為你穩定的牢,.
以至無論你在人生什麼逆境困境裡面,.
你都能夠有聖經裡面所說的平安和安穩,.
所以無論日子如何,.
無論得失,無論貧富,.
盼望經文提醒我們都能夠忠於跟隨耶穌,.
謹守祂的吩咐,.
我們一起禱告..
主啊,我們求你,.
堅固我們對你的忠心,.
幫助我們在世界裡面,.
我們可以因著聖經的教導,.
因著耶穌基督你的提醒,.
並且因著耶穌基督你拯救的生命,.
我們可以站立得穩..
多謝你,.
讓這段經文重新提醒我們,.
作為你的跟隨者,.
我們的本質是什麼?.
我們的本分是什麼?.

$^{481}$確實我們都希望生活得安穩一點,.
舒服一點,.
因為這一切其實都是來自上帝你的恩典供應,.
然而這一切卻不是我們人生的全部,.
因為耶穌你說過,.
我們在世上會有苦難,.
每一個忠心跟隨你的人,.
都必然面對魔鬼撒旦的試探和攻擊,.
然而這一切都能夠在耶穌你的恩典裡面,.
轉化成為屬天的平安,.
那個生命的冠冕..
所以盼望主也幫助我們,.
讓我們再重新檢視我們在生活裡面,.
我們所聚焦的,所拼搏的,.
到底在你眼中是得還是失?.
能不能夠頹利起悅,.
到底在天國裡面,.
這一切是財寶還是貧窮?.
幫助我們重新的來看著耶穌你,.
認定你是那位人生的主,歷史的主宰,.
是我們生命唯一的保障的時候,.
讓我們能夠以耶穌基督你得勝復活的生命,.
再一次來演繹我們自己的人生故事..
不是純粹聚焦在我有多困難,.
我的問題何時解決,.
而是我們可以聚焦在耶穌基督,.
你那個穩妥周詳的安排身上,.
以致我們能夠在不同的事上,.
能夠為主你多一份堅持,.
忍耐,忠心的來跟隨你..
求主你幫助我們學習這個功課,.
讓我們的生命可以像是美娜教會那樣,.
榮耀你,得到你的稱讚和掌賞..
我們的禱告交託,.
奉基督耶穌你得勝的名字來祈求..
Amen.
\newpage



\section{約翰福音 10:1-18}
\label{sec:980zAvvXkiM}
\textbf{2025年10月19日崇拜直播｜梁家麟牧師｜耶穌是好牧人｜約翰福音10\_1-18}
\newline
\newline
連結: \href{https://youtube.com/watch?v=980zAvvXkiM}{\texttt{https://youtube.com/watch?v=980zAvvXkiM}} ~~~~ 語音日期: 2025-10-19
\newline
\newline
\hyperref[sec:mfLhrRj9Qu0]{\small{< < < PREV SERMON < < <}}
~
\hyperref[sec:index_chronic]{\small{[返順時目]}}
~
\hyperref[sec:SlVA2_DGfXk]{\small{> > > NEXT SERMON > > >}}
\newline
\newline
約翰福音 10:1-18
\newline
\begin{longtable}{cl}
\hline
\hline
章節 & 經文 (和合本修訂版)\\
\hline
10:1 & \begin{tabularx}{0.7\textwidth}{X} 「我實實在在地告訴你們，那不從門進羊圈，倒從別處爬進去的，就是賊，就是強盜。 \end{tabularx} \\ \\ \relax
10:2 & \begin{tabularx}{0.7\textwidth}{X} 那從門進去的才是羊的牧人。 \end{tabularx} \\ \\ \relax
10:3 & \begin{tabularx}{0.7\textwidth}{X} 看門的給他開門，羊也聽他的聲音。他按著名叫自己的羊，把羊領出來。 \end{tabularx} \\ \\ \relax
10:4 & \begin{tabularx}{0.7\textwidth}{X} 當他把自己的羊都放出來，就走在前面，羊也跟著他，因為牠們認得他的聲音。 \end{tabularx} \\ \\ \relax
10:5 & \begin{tabularx}{0.7\textwidth}{X} 羊絕不跟陌生人，反而會逃走，因為不認得陌生人的聲音。」 \end{tabularx} \\ \\ \relax
10:6 & \begin{tabularx}{0.7\textwidth}{X} 耶穌把這比方告訴他們，但他們不明白他所說的是甚麼。 \end{tabularx} \\ \\ \relax
10:7 & \begin{tabularx}{0.7\textwidth}{X} 所以，耶穌又對他們說：「我實實在在地告訴你們，我就是羊的門。 \end{tabularx} \\ \\ \relax
10:8 & \begin{tabularx}{0.7\textwidth}{X} 凡在我以前來的都是賊，是強盜；羊沒有聽從他們。 \end{tabularx} \\ \\ \relax
10:9 & \begin{tabularx}{0.7\textwidth}{X} 我就是門，凡從我進來的，必得安全，並且可進出，找到草吃。 \end{tabularx} \\ \\ \relax
10:10 & \begin{tabularx}{0.7\textwidth}{X} 盜賊來，無非要偷竊、殺害、毀壞；我來了，是要羊得生命，並且得的更豐盛。 \end{tabularx} \\ \\ \relax
10:11 & \begin{tabularx}{0.7\textwidth}{X} 「我是好牧人，好牧人為羊捨命。 \end{tabularx} \\ \\ \relax
10:12 & \begin{tabularx}{0.7\textwidth}{X} 雇工不是牧人，羊不是他自己的，他一看見狼來，就撇下羊群逃跑；狼抓住羊，把牠們趕散。 \end{tabularx} \\ \\ \relax
10:13 & \begin{tabularx}{0.7\textwidth}{X} 雇工逃走，因為他是雇工，對羊毫不關心。 \end{tabularx} \\ \\ \relax
10:14 & \begin{tabularx}{0.7\textwidth}{X} 我是好牧人；我認識我的羊，我的羊也認識我， \end{tabularx} \\ \\ \relax
10:15 & \begin{tabularx}{0.7\textwidth}{X} 正如父認識我，我也認識父一樣；並且我為羊捨命。 \end{tabularx} \\ \\ \relax
10:16 & \begin{tabularx}{0.7\textwidth}{X} 我另外有羊，不屬這圈裡的，我必須領牠們來，牠們也要聽我的聲音，並且要合成一群，歸一個牧人。 \end{tabularx} \\ \\ \relax
10:17 & \begin{tabularx}{0.7\textwidth}{X} 為此，我父愛我，因為我把命捨去，好再取回來。 \end{tabularx} \\ \\ \relax
10:18 & \begin{tabularx}{0.7\textwidth}{X} 沒有人奪去我的命，是我自己捨的；我有權捨棄，也有權再取回。這是我從我父所受的命令。」 \end{tabularx} \\ \\
[1ex]
\hline
\hline
\end{longtable}
$^{1}$今天經課是記載於《約翰福音》第十章.
「我實實在在的告訴你們.
人進羊圈不從門進去.
倒從別處爬進去.
那人就是賊,就是強盜.
從門進去的才是羊的牧人.
開門的就給他開門.
羊也聽他的聲音.
他按著鳴叫自己的羊把羊領出來.
既放出自己的羊來.
就在前頭走.
羊也跟著他.
因為認得他的聲音.
羊不跟著生人.
因為不認得他的聲音.
必要逃跑」.
耶穌將這比喻告訴他們.
但他們不明白所說的是什麼意思.
所以耶穌又對他們說.
「我實實在在的告訴你們.
我就是羊的門.
凡在我以先來的.
那是賊,是強盜.
羊卻不聽他們.
我就是門.
凡從我進來的.
必然得救.
並且出入得草吃.
盜賊來.
無非要偷竊,殺害,毀壞.
我來了.
是要讓得生命.
並且得的更豐盛.
我是好牧人.
好牧人為羊捨命.
若是故宮不是牧人.
羊也不是他自己的.
他看見狼來了.
就撇下羊逃走.
狼抓住羊.

$^{41}$趕哨練羊群.
故宮逃走.
因為他是故宮.
並不顧念羊.
我是好牧人.
我認識我的羊.
我的羊也認識我.
正如父認識我.
我也認識父一樣.
並且我為羊捨命.
我另外有羊.
不是這圈裡的.
我必須令牠們來.
牠們也要聽我的聲音.
並且要合成一群.
歸一個牧人了.
我父愛我.
因我將命捨去.
好在取回來.
沒有人奪我的命去.
是我自己捨的.
我有權並捨了.
也有權並取回來.
這是我從我父所授的命令.
經文讀筆.
各位節目早晨.
我們繼續漫長的約翰福音之旅.
我們今天要看第十章.
這段聖經其實是很熟悉的.
不過很多人是孤立地去看第十章.
有少少無厘頭.
唯有我們將它放回原來經文的脈絡裡.
連接第九章來看.
我們才會看到耶穌說的話的份量.
緊接著第九章.
「乞子得醫治」的故事.
耶穌在第十章.
有感而發地表白自己的身份和使命.
我們記得「乞子得醫治」之後.
由於醫治他的是法利賽人所痛恨的耶穌.

$^{81}$所以他們連乞子也迫害.
將他趕出猶太的社群.
害到他幾乎走投無路.
天下之大.
竟然沒有用心的空間.
乞子到底做錯了什麼錯事.
招來這樣的慘況呢?.
沒有.
他既沒有犯罪.
也沒有違反任何宗教或道德的戒律.
他只是沒有配合法利賽人的陰謀.
拒絕說謊.
拒絕耶穌是罪人的結論而已.
但只是說事實.
只是說真話的一個小人物.
竟然被那些宗教權貴無情打壓.
要將他趕入絕境.
你可以看到.
對法利賽人來說.
守護他們心中的宗教成見.
是最高的甚至是唯一的價值.
至於人的死活.
是他們漠不關心的.
特別是像乞子這個社會最底層的賤民.
他們的生命毫無價值.
他們的感受可以不理.
這令我想起.
為了堅持某些宗教或道德的價值.
而輕視甚至殘害生命的故事.
例如有印度教徒.
因為女兒未婚懷孕.
或與低種性的人結婚.
就將女兒活活燒死.
我相信宣揚這樣的教義的人是邪惡的.
我很清楚.
過份推崇人性.
會使我們教會變成自由主義的危機.
甚麼是自由主義?不知道就無所謂.
但所有違反人性的教義或規矩.
都是邪惡的.

$^{121}$毫無例外.
完全擁抱人性有危機.
但違反人性就一定是邪惡的.
這個故事我以前跟你說過.
我有一次與某大教會的執事聊天.
教會在過去十多年間.
離開了大量的年輕人.
原因是他們無法忍受教會的崇拜禮儀.
我問執事知不知道這樣的情況.
他說當然知道.
並且他說無所謂.
他說我們的崇拜才是最符合聖禮神學的.
不跟從這個傳統的就請他們自己離開.
我當下的反應.
我之前跟你說過.
我馬上想起加拿大的三文魚.
為了產卵逆流游上石澗.
導致遍體鱗傷.
產卵後全部死亡.
動物可以為了繁衍下一代犧牲自己.
這個甚麼是聖禮神學?.
對我來說不僅僅是錯謬.
是邪惡的.
請聽清楚.
我是認真的.
我不介意有人捉住把柄去控告我.
任何不惜以趕走下一代為代價.
而承傳的神學系統都是邪惡的.
基督教有講社稷.
剛才我們師班唱的詩歌都有講到社稷這個字.
但社稷只能為了實踐.
為了完成使命而做.
而絕對不可以為了符合某一個教義或禮義的標準而做.
不可以.
教義或禮義絕對沒有最高價值.
安息日是為人而設.
當然如果我們說.
如果羅記不同意三一論.
我們就會請他離開.
但我都不會燒死他.

$^{161}$你會不會?.
不可能.
我們不會做這些事.
我們不會為了保衛某一個教義而殘害一個人的生命.
絕對不做這件事.
只有使命才有更高的價值和義.
你為了堅持信仰.
面對政權.
受到迫害.
這個我可以明白和尊重你.
但如果你在家說我為了基督的緣故.
自己釘上十字架.
我就看不到你有什麼價值.
當然你這樣做我也尊重.
但如果你氣別人這樣做.
我就視你為邪惡.
我一定要誅殺你.
對不起,你明白我的意思嗎?.
我們教會不可以容許這樣的事.
任何正確的宗教和道德.
都是為了造就和承傳生命.
絕對不可以貶抑生命,殘害生命.
這是耶穌的宣告.
安息日是為人而設.
在安息日救命是絕對沒有問題.
Elvis Weiser 史維哲說.
基督教的核心精神是尊重生命.
Reverence for life.
尊重生命是耶穌基督所堅持的基本態度.
你可以從以賽亞書.
以賽亞先知對他的預告.
這個預告是描述耶穌的.
「壓傷的爐火他不截斷,將殘的燈火他不催滅」.
得到最好的說明.
再脆弱無助的生命.
再殘缺,看起來像垃圾無價值的生命.
耶穌都不會輕忽對待.
反正你要死了,就送你一程吧.
不可以的,是不是?.
生命寶貴.

$^{201}$只要有一口氣的生命.
就值得尊重珍惜.
所以在這裡.
你看到耶穌對法利賽人顛倒是非.
踐踏人性的做法.
是多麼義憤填胸.
他作出這樣的評論.
第十節.
盜賊來,無非是要偷竊,殺害,毀壞.
我來了,是要叫讓得生命.
並且得的更豐盛.
與其說是耶穌基督對自己的吹噓.
我要使人得生命.
不如說.
他用這段說話來表達.
他對那群吻沒人性的法利賽人.
一個嚴重的自豪.
重點在於他罵那群法利賽人.
耶穌將他的使命和相應的心態.
與法利賽人作比較.
指出,核子就像一隻無助的羊.
耶穌是牧人.
而法利賽人是盜賊.
盜賊純粹是拿羊來圖利.
為了自己的利益.
就不擇手段.
可以隨意踐踏和傷害羊.
但耶穌是好牧人.
立志保護羊的生命.
並且致力使羊的生活過得更好,更豐盛.
耶穌生長在一個打魚和放牧的加利利地區.
羊群和放牧是隨處可見.
所以耶穌常常用羊和放牧來做比喻.
祂是人,是待牧羊的羊.
馬利福音9:36,大家很熟悉.
「他看見許多人,就連殞他們,因為他們困苦流離,如同羊,沒有牧人一半」.
耶穌的眼光和以賽亞的說法相對應.
「我們都如羊走迷,各人偏行幾路」.
人是迷失了的羊,迷路的羊.
這其實是舊約很多先知常用的意象.

$^{241}$在以賽亞書,耶利米書,以西結書,尼加書,撒迦利亞書.
都有大量使用牧人和放牧的技術.
耶利米書和以西結書更清楚預告.
上帝會親自做以色列人的牧人,好好地牧養他們.
好牧人和盜賊的分別.
在於他們是羊為目的或是工具.
在於他們是羊為具有終極價值或只是有工具價值,即實用性.
是羊有終極價值的,就不會為了高舉某個宗教或政治的理想.
而輕言犧牲羊的福祉甚至生命.
是羊為僅僅有工具價值的,就會宣稱為了一個更崇高的政治或宗教的目標.
就可以犧牲個別羊的利益乃至生命.
例如法利初一人認定耶穌是邪惡的.
耶穌的存在對他們的信仰系統構成威脅.
就一定要想盡一切方法去剷除它.
他們將剷除耶穌這個任務美化成為神聖的至高無上的.
為此所有代價都值得付,包括乞子的福祉,這個垃圾.
犧牲它就犧牲它,當然法利初一人未必會犧牲自己.
犧牲人的生命是很爽快的,犧牲自己就未必.
乞子願意配合法利初一人的計劃就有存活的價值.
乞子如果不肯配合就成了施展他們陰謀的障礙.
不好意思,你已經沒有存活下去的理由,死不足惜.
在除掉耶穌這個至高的宗教和政治任務的面前.
所有人愛,公義,合理這些所謂價值系統全部都可以放棄.
在法利初一人眼中,宗教系統的穩定性.
宗教規條的嚴謹性比人命更重要.
個別的人更加像螻蟻一樣.
成大事者,誰不是踐踏在很多屍體之上?.
手持心軟的人是幹不了大事的.
唯有殺伐果斷的人才可以吃大茶飯.
耶穌對人的憐憫,不過是婦人之仁.
婆婆媽媽,成事不足,真的.
耶穌很婆媽,真的太婆媽了.
不過在耶穌的眼裡,所有以色列人.
包括這個生下來就合了眼的,都是寶貴的羊.
都是值得珍惜和愛護的.
可惜,上帝原來設立用來照顧百姓的崗位.
那些君王,祭司,律法師等等.
都沒有履行上帝牧養百姓的責任.
反而好像賊一樣偷竊,殺害他們.
耶穌認為,這些人不僅沒有盡責任做應該做的事.

$^{281}$才為數餐,更加是竊奪了牧養的位置.
他們本質上不是牧人,是盜賊.
沒有做應該做的事,只是失職.
偷佔一個不屬於他們的位置.
就不僅僅是職務問題,更加是身份和資格的問題.
耶穌否定這些人有做牧人的資格.
不是職務問題,而是資格問題.
這是耶穌從一開始就說,他是羊的門的意思.
耶穌首先不是說他是好牧人.
首先第一個表白的是,他是羊的門.
1到3節,我實實在在告訴你們.
人進羊圈,不從門進去,都從別處爬進去.
那些人就是賊,就是強盜.
從門進去的才是羊的牧人.
開門的就覺得開門,羊也聽他的聲音.
他按著命,叫自己的羊,把羊領出來.
羊圈,指的是牧場主人所開設的牧場.
也可以說是牧養的基地.
每天主人所僱用的工人,就是牧人.
來到牧場,把他們自己,每個牧人要負責的羊.
把他要牧養的羊,領出來.
然後他帶這些羊去有水草的地方.
牧人是真正的工人.
是不會攀爬那些用荊棘或木條築成的矮牆.
爬進羊圈,他只會從門進入.
翻牆進羊圈的人,不可能是牧人,一定是盜賊.
所以從他們的來路,就可以知道他們的身份.
你們回來旺州,是乘電梯或走樓梯上來的.
你們不會在外牆爬上來.
如果突然我聽到槍有人敲,然後有人開槍讓我進來.
坦白說,我都會當他是來路不明的人.
好人好者,爬牆幹嘛?.
你想顯示你的飛禪走壁的能力?.
沒意思的,你不是從那裡進來的.
多數都是賊,不用說.
同樣,你無情無事,你為什麼要攀過難關走進來?.
為什麼不從門上堂堂正正進來?.
耶穌的話就是這個意思.
滿土不明白耶穌的比喻.
耶穌就解釋在第七到九節.

$^{321}$「我實實在在告訴你們,我就是羊的門.
凡在我以先來的,都是賊,都是強盜.
羊卻不聽他們,我就是門.
凡從我進來的,必然得救.
並且出入得草吃」.
耶穌是羊的門.
這段話前後指出兩個意思.
第一個關乎牧人.
從門進來的才是真牧人.
即是由耶穌核准.
通過耶穌的才是牧人.
過不到耶穌的關.
即你不是通過他,你就叫賊.
所有人都要經過耶穌的允許核准.
才有資格接觸羊群.
然後才取得羊的資格.
「我就是羊的門.
凡在我以先來的,都是賊,都是強盜.
羊卻不聽他們」.
Jesus is the gate to the sheep.
你要經過耶穌才去到羊那裡.
我要說聖經一直很強調.
那個 ownership(擁有權).
就算用耕種來做比喻.
聖經也說.
你不是看到莊稼熟了就自己去收莊稼.
你求莊稼的主給你有資格.
猜法他的工人是耶穌的工人去收莊稼.
你不去自己拿鐮刀衝去自己割.
你想死嗎.
莊稼的主權在哪裡?.
在耶穌那裡.
同樣羊的主權在哪裡?.
在耶穌那裡.
Jesus is the gate to the sheep.
羊是屬於耶穌的.
誰能夠接觸羊.
當然由耶穌說了算.
第二是關乎羊.
剛才是關乎牧人的說法.

$^{361}$第二是關乎羊的說法.
羊必須要從門出去.
才得著飽足和生命.
羊要穿過這個門才能夠出去.
進入草場和水源.
羊要藉著耶穌得救並且得草吃.
耶穌說.
我就是門.
凡我從我進來的必然得救.
並且出入得草吃.
Jesus is the gate for the sheep.
不僅是牧人要透過耶穌才能夠接觸羊.
羊也要透過耶穌才能得生命.
唯有耶穌是生命的糧.
唯有耶穌是活水.
能夠讓人喝了就不再口渴.
後面耶穌說.
只有祂才能使人得生命.
並且得得更豐盛.
耶穌是羊的門.
從門進去的才是真牧人.
從門出去的才能得生命.
這是羊的門的意思.
從舊約到新約.
聖經清楚告訴我們.
羊是屬於耶穌基督的.
所有人都是屬於耶穌基督的.
耶穌在升天前對彼得的三次託付.
都強調什麼.
你餵羊 我的羊.
你牧羊 我的羊.
羊是誰的 是耶穌的.
是耶穌的 是不是.
由於羊是屬於耶穌的.
耶穌親自挑選和差派人.
去擔當牧人的角色.
由於羊是屬於耶穌的.
只有祂才了解羊的需要.
給他們真正的幸福.
羊屬於耶穌.

$^{401}$這不僅僅是主權誰屬的問題.
也關乎以下兩個重點.
第一是責任 第二是關係.
先講責任.
耶穌接著說到故宮.
就是祂所聘用的牧人.
和羊的主人的差異.
11到13節.
我是好牧人 好牧人為羊捨命.
若是故宮不是牧人.
羊也不是他自己的.
看見狼來就撇下羊逃走.
狼抓住羊 趕散羊群.
故宮逃走 因為他是故宮.
並不顧念羊.
打工仔只有打工仔的心態.
打工仔是出賣自己的時間和勞動力.
去換取生活之資.
這是一個簡單的交換.
你出錢 我出力.
但出力可以 出命就不要了.
做工的人不會搏命.
例如風暴襲港.
會有人因為追風逐浪跌落海.
我是救護員 我接報到場營救.
能夠救的 我當然盡力救.
但我不會賠自己的命去救他.
對不起 是不是.
如果風險太大.
正如消防員的火場完全失控.
不受控制.
會吹哨叫所有消防員撤離.
不會叫他們全部在裡面燒死.
是不會的.
救護員第一任務是保護好自己.
量力而為.
如果每個消防員 每個警察都做雷鋒.
輕易殉職的話.
那浩元就沒有空位可以安葬了.
是不是 不可能的.

$^{441}$因為職務拯救墮海者要量力而為.
但如果墮海的是我的兒子.
就算風高浪急.
我都會奮不顧身.
用命跳落海.
我會不會說收工呢.
我不會的.
如果那是我的兒子就不會的.
無論能夠救或不能救.
我都會盡力去救.
這不是量力而為.
是盡力而為.
是兩碼子的事.
耶穌在比喻裡面這樣說.
遇到狼群來襲.
做故宮牧人的.
要計算能不能頂得住.
頂不住的就要保自己的命.
射羊群離去.
最多就好像舊約律法要求的.
你盡人事在狼的口裡.
拔出一隻羊耳或羊腿.
做什麼 是交數.
我有沒有跟你說過.
這個叫交數.
因為之後不見了幾隻羊.
主人說你可能吞駁.
你可能自己燒烤.
燒了來吃.
怎麼證明有狼群.
怎麼證明有獅子.
你就拿一隻羊腿出來.
你看看那個咬口.
是狼的咬口.
就好像我有沒有跟你說過.
我以前讀大學的時候.
我去托汽水.
整天會整盤汽水打爛.
打爛要做什麼.
撿一些包頭.

$^{481}$什麼叫包頭.
就是連蓋的瓶頸.
撿回去就可以交數.
打爛是不用賠的.
做事不可能不打爛的.
但是沒有包頭.
那就是你吃的.
不見一盤水就賠一盤水.
就是這樣.
如果真的被人吃了.
沒辦法.
只能夠盡力去找證據回來.
這個就是盡責.
這是《春夏及記》22章10-13節的律法規定.
不是為了拯救羊的生命.
只是為了自保.
但是羊的主人就不同.
因為羊是屬於他的.
他奮不顧身.
為了保衛羊.
不惜犧牲自己的生命.
當然現實的說法.
你要批評耶穌的說法是很容易的.
羊對主人不過是財產.
錢財是身外物.
為了拯救羊.
而連自己的生命都要賠償.
這個很蠢的.
我想起也不明白.
就算羊是我.
我也不用拼命的.
不會連自己的命都捨賣的.
所以無論是故宮或主人.
都不應該為了拯救羊而拼命.
不過耶穌在這裡是很必有所指的.
他就是你不明白為何這麼蠢的人.
他是願意為羊群捨命的人.
為了他所愛的人.
其實是受造物.
耶穌甘願降世.

$^{521}$自我犧牲.
從引領上帝來說.
我們無法理解為何上帝會.
鍾情他所做的一件事.
他不是有些戀物狂.
為何上帝會為他的受造物.
犧牲自己的生命.
在情理上完全說不過去.
我無法明白為何上帝不簡單.
消滅所有犯罪的人類.
重新再造一批完美的人出來.
唯一可以解釋的是.
上帝竟然和他所造的人類發生感情.
上帝愛世人.
愛使上帝變得愚蠢.
所以就做了蠢事.
不過即使我們不會為了保護財物.
而犧牲自己的生命.
一樣東西是屬於自己和不屬於自己.
仍然有分別.
我和太太在結婚一年後.
請了菲傭.
我們和菲傭的關係很好.
菲傭也很忠心.
我從來不罵菲傭.
多年來沒有罵過.
她很盡心竭力工作.
不過我們家的菲傭是很大手大腳的人.
無論是吸塵機或是熨斗.
很快就用爛.
用爛的只是一支膏.
不能想像.
杯,碟,鍋,打爛就更加不用說.
幸好現在這個世界物品很便宜.
所以沒有人計較.
你會發現.
公眾的物品破損通常都很高.
我在意大利神學院.
在天鵬.
我們做了十個八個浪衫架.

$^{561}$讓人可以在上面浪衫.
我發覺.
不夠一個月就爛三個四個.
兩個月就差不多要換一批.
很容易爛掉.
不知為何.
我自己在房間買了一個浪衣架.
用了這麼久都沒有爛過.
不用解釋.
公眾的和不是公眾的完全不同.
是做事的方法.
這也帶來了第二點.
第二點要說的.
關係.
羊是屬於耶穌的.
所以耶穌傾全力去拯救保護.
羊屬於耶穌的.
所以和耶穌建立了個別的深厚的感情.
14到15節說.
我是好牧人.
我認識我的羊.
我的羊也認識我.
正如父認識我.
我也認識父一樣.
並且我為羊捨命.
耶穌以他和天父的關係.
去比喻他和羊的關係.
耶穌和天父是彼此認識的.
耶穌和他所愛的人也彼此認識.
前面在3到5節.
耶穌也說.
羊也聽他的聲音.
他按著名叫自己的羊.
把羊領出來.
既放出自己的羊來.
就走到前頭.
羊也跟著他.
因為認得他的聲音.
羊不跟著生人.
因為不認得他的聲音.

$^{601}$必要逃跑.
羊是聽得出牧人的聲音.
聽到他們熟悉的聲音.
他們就會放心.
跟隨.
謹靠.
對陌生的聲音.
羊會感覺到害怕.
逃走.
中國人是重情講關係的民族.
當然明白關係的重要性.
有關係就沒有關係.
沒有關係就關係更大.
耶穌在這裡強調彼此認識.
因著認識.
對方不僅是一個空的概念.
不僅是會有學生.
人類.
而變成有血有肉的個體.
並且雙方都投注了感情在其中.
舉個最近的例子.
你知那些落後國家常常都發生反政府示威.
很多時候都會演變成暴力衝突.
人們就會乘機搶掠店舖.
到處縱火.
我們在電視新聞.
經常看到這樣的報導.
不過大部分人都是看一看就算數.
很少放在心上.
最近非洲.
印度洋的馬達加斯加島.
發生大規模的騷亂.
蔓延全國.
在那裡宣教的是.
建築學院的兩對校友夫婦.
黃聲鋒各自西宣教士.
和梁志康張佩欣宣教士.
他們都有一段時間被困在一個小鎮裡.
他們寫信息來向我求救.
我當然不是叫我救他們.

$^{641}$叫我幫他們祈禱.
我是很迫切祈禱.
很迫切祈禱.
跟某一個小島發生騷亂是兩回事.
有我所愛的人在.
我關心的兩對宣教士在.
我怎麼可能無動於衷.
馬上祈禱得很認真.
是不是.
耶穌認識我們.
耶穌在乎我們.
我們的生死和福.
我們的喜怒哀樂.
耶穌都介意.
耶穌都放在心上.
這和他和故宮的不同之處.
更加不要說和盜賊的差異.
發利賽人可以不顧核子的生死.
可以隨意踐踏他的生命.
但耶穌關心核子的命運.
關心他的情緒感受.
在結束這段說話的時候.
耶穌為了避免門徒和後人有誤會.
做了兩個保註.
兩個注釋.
第一.
耶穌不是說只有以色列人才是上帝的子民.
上帝只關心以色列人.
愛護以色列人.
其他民族和上帝無關.
上帝不管他們的生死.
在十六節.
耶穌強調.
我另外有陽.
不是只圈裡的.
我必須領他們來.
他們也要聽我的聲音.
並且要合成一群.
歸一個牧人了.
陽福音三章十六節告訴我們.

$^{681}$上帝愛世人.
不是上帝愛以色列人.
耶穌來人間不是只為以色列人犧牲.
是為所有人犧牲.
祂拯救像你和我這樣的壞人.
今日普世信徒都夢耶穌基督拯救.
歸入他的身體.
教會.
同屬群陽的大牧人照管.
所以一定要補充.
祂不只是關心以色列人.
祂是關心所有人.
第二.
耶穌是主動為拯救人類犧牲自己.
祂多番強調.
我是好牧人.
好牧人就為陽捨命.
所以祂不是因為出了什麼意外.
遭逢不測.
又或者不幸落在仇敵手中.
無辜受害.
在耶穌眼中.
猶太教當局和羅馬帝國一點都不厲害.
撒旦對耶穌也完全沒有辦法.
他們不是構成耶穌遇害的真正原因.
更加不要說是充分理由.
耶穌是自己送上門.
自己願意犧牲.
主動道成肉身.
主動迎向十字架.
一切都是為了愛.
17到18節.
我父愛我.
因我將命捨去.
好在取回來.
沒有人奪我的命去.
是我自己捨的.
我有權病捨了.
也有權病取回來.
這是我從我父所受的命令.

$^{721}$耶穌的死亡並不意味祂對生命失去控制權.
更加不代表祂失去了天父的愛.
我很喜歡《約翰福音》十章.
17到18節這兩節經文.
日後有機會再和大家在這裡思考.
我們要做一個很簡單的總結.
今天的經文傳遞了兩個很簡單.
你本來都知道的真理.
第一,耶穌是好牧人.
為我們寫命.
我們是祂的羊,屬於祂.
必須要聽祂的聲音.
唯有藉著祂,生命才得到飽足.
並且得到永遠的生命.
這是足夠重要的道理.
第二個真理.
我們要明白耶穌是多麼愛人的心腸.
祂重視每個個體的價值.
珍惜和承傳人的生命.
我們立志以祂的心為心.
尊重生命.
善待身邊每個人.
絕對不把人當作工具.
或幫我圖謀個人利益的工具.
我要說,這是這段聖經給我更重要的真理.
不是前者不重要,前者是你早就知道.
這一點你也可能知道.
不過我們再次重申.
Reference for life.
尊重每一個生命.
這是基督教信仰的一個很關鍵的精神.
\newpage



\section{啟示錄 2:12-17}
\label{sec:SlVA2_DGfXk}
\textbf{2025年10月26日崇拜直播｜羅志雄牧師｜耶穌給教會的七封信:當悔改｜啟示錄2\_12-17}
\newline
\newline
連結: \href{https://youtube.com/watch?v=SlVA2_DGfXk}{\texttt{https://youtube.com/watch?v=SlVA2\_DGfXk}} ~~~~ 語音日期: 2025-11-03
\newline
\newline
\hyperref[sec:980zAvvXkiM]{\small{< < < PREV SERMON < < <}}
~
\hyperref[sec:index_chronic]{\small{[返順時目]}}
~
\hyperref[sec:LTkMkPBreS4]{\small{> > > NEXT SERMON > > >}}
\newline
\newline
啟示錄 2:12-17
\newline
\begin{longtable}{cl}
\hline
\hline
章節 & 經文 (和合本修訂版)\\
\hline
2:12 & \begin{tabularx}{0.7\textwidth}{X} 「你要寫信給別迦摩教會的使者，說：『那有兩刃利劍的這樣說： \end{tabularx} \\ \\ \relax
2:13 & \begin{tabularx}{0.7\textwidth}{X} 我知道你的居所，就是有撒但座位之處；當我忠心的見證人安提帕在你們中間，在撒但所住的地方被殺之時，你還堅守我的名，沒有否認對我的信仰。 \end{tabularx} \\ \\ \relax
2:14 & \begin{tabularx}{0.7\textwidth}{X} 然而，有幾件事我要責備你，就是在你那裡有人服從了巴蘭的教訓；這巴蘭曾教唆巴勒將絆腳石放在以色列人面前，使他們吃祭過偶像之物，並且犯淫亂。 \end{tabularx} \\ \\ \relax
2:15 & \begin{tabularx}{0.7\textwidth}{X} 同樣，你那裡也有人服從了尼哥拉派的教訓。 \end{tabularx} \\ \\ \relax
2:16 & \begin{tabularx}{0.7\textwidth}{X} 所以，你當悔改；若不悔改，我很快就到你那裡來，用我口中的劍攻擊他們。 \end{tabularx} \\ \\ \relax
2:17 & \begin{tabularx}{0.7\textwidth}{X} 凡有耳朵的都應當聽聖靈向眾教會所說的話。得勝的，我必將那隱藏的嗎哪賜給他，並賜他一塊白石，石上寫著新的名字，除了那領受的以外，沒有人認識。』」 \end{tabularx} \\ \\
[1ex]
\hline
\hline
\end{longtable}
$^{1}$今日的經訓是啟示錄第二章十二至十七節.
聖經新約第三百五十至三百五十一頁.
這裡記載了耶穌基督透過仕途約翰.
寫給別加摩教會的信.
向他們作出鼓勵,責備,勸導和應許.
呼籲信徒悔改,更新,勸勉他們要得勝.
這裡是這樣記載的.
第十二節.
你要寫信給別加摩教會的使者.
說:那有兩人離間的說.
我知道你的居所就是有撒旦座位之處.
當我忠心的見證人安提柏在你們中間.
撒旦所住的地方被殺之時.
你還堅守我的命,沒有棄絕我的道.
然而,有幾件事我要責備你.
因為在你那裡,有人服從了巴蘭的教訓.
這巴蘭曾教導巴勒.
將半腳石放在以色列人面前.
叫他們吃掉偶像之物.
行奸淫的事.
你那裡也有人照樣服從了尼哥拉一黨人的教訓.
所以,你當悔改.
若不悔改.
我就快臨到你那裡,用我口中的劍攻擊他們.
聖靈向宗教會所說的話.
凡有耳的就應當聽.
得逞的,我必將所隱藏的嗎哪刺及他.
並刺他一塊白石.
石上寫著生命.
除了那領受的以外.
沒有人能認識.
請姐妹主內平安.
不知道大家心目中的教會是怎樣.
或者叫理想的教會是怎樣.
有一位牧師在教會的網站上發起一個投票.
如果能加入一間完美理想的教會.
這間教會應該是怎樣的呢?.
網站發出了這個問卷後.
幾天內收到很多回覆.
有人說應該有優質的講道.

$^{41}$有感動人的敬拜.
也有重視愛心和備接納.
甚至有人希望教會的設施可以現代化.
有溫暖活動多一點.
但不要批評人.
牧師很幽默地回應.
如果要找到這間完美的教會.
牧師說不要加入.
因為一加入.
這間教會就變得不完美.
不知道大家聽完之後.
會否想想我們現實.
我們很多時候都會追求理想或心目中的教會.
但我們又有多少時候去思想一下.
主耶穌心目中的教會會是怎樣呢?.
啟示錄記載著主耶穌有關七間教會評論他們的評價.
上兩星期陳明泉牧師和Ray牧都先後說過已忽所書.
已忽所的教會忠於真理.
但失去了起初的愛心.
Ray牧所說的仕美娜教會.
就在患難中至死忠心.
而在當中並沒有被責備的其中一間教會.
今天保齡姐妹和我們讀的別加摩教會又如何呢?.
在當中有什麼特質.
可以提醒我們去反思.
我們所追求的教會.
是否符合耶穌心意的教會呢?.
很想透過別加摩教會.
三個P.
主的同在.
主的保護.
主的德性.
去思想一下.
關於耶穌心目中的教會是怎樣.
我們來看第一點.
主的同在.
主耶穌是知道我們的所在.
經文第十二至十三節是這樣說.
你要寫信給別加摩教會的使者.
那有兩人離間的說.

$^{81}$我知道你的居所就是有撒旦座位之處.
當我忠心的僕人安提柏在你們中間.
撒旦所住的地方被殺之時.
你還堅守我的命,沒有棄絕我的道.
經文提到.
耶穌是以那有兩人離間.
來自稱自己.
在第一章的時候.
耶穌就是這樣.
也曾經提過.
祂是有兩人離間的.
是彰顯了祂的說話.
那種審判權柄.
和祂是去察看.
去引導和管教祂的教會.
這一句我知道短短的三個字.
在七封書信裡重複出現.
是強調主的全知.
祂不單單是觀察.
更加深刻地明白我們信徒每一個處境.
我們內心的狀況.
別加摩的信徒雖然身處在壓力的狀況之下.
但主耶穌說我知道.
我相信這就是提醒我們.
耶穌是明白.
知道我們經歷和承受的狀況.
祂會和我們一起經過.
至於耶穌對別加摩教會所說的.
我知道你的居所所指的是什麼狀況呢?.
當然按照經文裡說.
別加摩教會是身處在別加摩城.
但耶穌不單單是看這個城市.
更加看到教會裡面.
每一個群體.
每一個跟隨祂的生命居所會是怎樣.
耶穌就是這樣明白我們心靈深處裡面所有的環境.
我們的內心掙扎和壓力.
祂是親自參與在弟兄姊妹之間的生命裡面.
無論我們面對什麼狀況.
神都和我們同在.

$^{121}$這句話帶來安慰和挑戰.
安慰的是耶穌明白我們軟弱.
祂不離棄我們.
但又會帶來給我們挑戰.
因為主會察看我們的選擇和行動.
究竟我們會選擇什麼.
用什麼方式去回應神.
呼召我們在我們所處的狀況環境裡面仍然持守信仰.
所以在別加摩教會裡面.
耶穌就是傳言的.
祂知道在祂的教會裡面每一個人的狀況.
既然主知道.
祂必定了解當中的環境.
經文第十三節裡面描述到.
就是有撒旦的座位之處.
撒旦座位之處所指的是什麼呢?.
就是指別加摩城這個地方.
充滿了偶像崇拜.
和很多多元的宗教匯聚在當中.
別加摩城.
了解它的地理位置和背景.
它位於亞細亞省.
仕美那城.
以北約150公里.
一個350米的山丘之上.
這個城市是當時羅馬帝國在亞細亞省的首府.
按照很多資料.
人口有20萬.
既然一個首府有20萬的人口.
相信它在宗教文化政治上.
商業貿易交通往來.
一定是非常繁榮.
還有很多文化的參集在當中.
別加摩城被耶穌形容為撒旦的座位.
或者撒旦的地方.
如果按照當時的理解.
別加摩城有四座希臘神廟.
其中有阿斯克萊比神廟.
這個就是伊治之神.
當一些尋求伊治的信眾.

$^{161}$進入這個神廟敬拜的時候.
他們祈求有條蛇出現.
經過他們的身上.
就會得全愈.
所以那個標誌.
有條蛇纏著一個木棒.
就成為一個標記.
他們看到這個.
覺得他們就能得到伊治.
如果大家有留意.
很多醫療機構都有這個標誌.
不知道大家有沒有留意到.
另外也有雅典娜女神.
有宙斯神廟.
宙斯神廟按考古發掘出來.
原來宙斯神廟的神壇.
是一座很大的祭壇.
上了十級的樓梯.
左右兩邊就有祭壇.
看上去像是寶座.
讓誰坐呢?.
不是讓創造天地的主坐.
原來他們拜偶像撒旦.
撒旦成為宙斯神廟的寶座.
最高的主人.
另外還有一些酒神.
更加重要的是.
他們會供奉羅馬皇帝.
特別在別加摩城裡.
有奧古斯都.
羅馬帝國的第一個皇帝.
如果看路加福音.
就看到奧古斯都.
在主前二十多年.
在別加摩城興建了.
第一座敬拜奧古斯都的神廟.
所以羅馬帝國的版圖.
每一個地方都會有很多.
敬拜皇帝,敬拜該撒的神廟.
當中每一個百姓,人民.

$^{201}$都要帶著香.
到皇帝的神廟來敬拜.
並且不只是敬拜.
要敬拜.
口中也要說.
這就是他的主.
他的神.
甚至是救主.
面對這些強烈的.
宗教壓力和屬靈的黑暗的時候.
別加摩教會裡的信徒.
隨時會承受很大的壓力.
所以如果他們要妥協的時候.
他們就參與異教的崇拜.
大主仍然在別加摩城的教會裡.
別加摩的教會.
耶穌了解他們黑暗的環境.
但耶穌沒有叫這間教會.
不如就遷到別的地方.
另外找一個沒有其他信仰.
沒有其他異教的地方來敬拜.
反而耶穌仍然盼望.
這間教會在黑暗裡作光.
不要逃避這個世界.
因為耶穌和教會同在.
要令到這個教會.
能夠在艱難的環境裡.
仍然持續盼望和信心.
弟兄姊妹.
當然別加摩城這種狀態.
可能是在二千多年前的事.
但這些事在今天.
我們身處的這個世界.
這個世代.
是否很相似呢?.
可能我們所認識的環境.
沒有很大的神廟,神壇.
但當我曾經去過台灣花蓮的時候.
去探望一對宣教士夫婦.
他們都說.

$^{241}$原來當地雖然沒有很大的神壇.
但原來周圍都有很多的廟.
甚至這對宣教士夫婦.
他們本身是少數民族.
他們的少數民族都有教會.
但有些異端的基督教.
亦都在這些地方.
傳揚福音.
將他們所相信的.
但歪曲了的.
傳揚.
弟兄姊妹.
我們無法改變一些環境.
但我們可否仔細想想.
我們是靠哪一個力量.
其實是我們靠主耶穌的力量.
在困難當中.
能夠展現我們的信仰.
所以弟兄姊妹.
在你現在的職場,家庭.
甚至信仰上面.
面對挑戰的時候.
耶穌都知道你身處的環境.
不要讓你所相信的信仰.
你跟隨耶穌的行動.
被清晰.
然後接受妥協.
我們來看看.
中心見證的典範.
主繼續說.
當我中心的見證人.
安提柏在你們中間.
撒旦所住的地方.
被殺之時.
你還堅守我的名.
沒有棄絕我的道.
安提柏這位信徒.
沒有其他文獻知道他是誰.
但卻是被主耶穌所稱讚.
其實就是代表著.

$^{281}$別加摩教會裡面.
有一群仍然像安提柏.
忠於耶穌,忠於真道.
而不懼怕受逼迫的一些信徒.
換言之.
安提柏就是分享基督見證的榮耀.
亦成為信徒效法的榜樣.
見證人Marta.
後來英文翻譯為信道者.
即是說見證耶穌基督的時候.
是會有機會面對生命的危險.
正如拒絕拜當時的皇帝一樣.
隨時都會有生命的危險.
主在此稱讚別加摩教會.
是堅守神的道.
沒有棄絕主的道.
這表示了.
縱使當時的別加摩城.
這麼複雜的宗教政治上.
在一個邪惡的環境裡.
他們仍然受到社會群體的恐嚇.
因為信徒都會和一群異教徒有很多的來往.
例如商業上,人際關係.
甚至吃飯時.
有時都可能進入了他們的神廟.
在這樣的狀況下.
他們如何能夠持守住.
耶穌基督就是他們唯一的救主.
而不稱凱撒為主,神.
而不去屈膝呢?.
我相信.
弟兄姊妹.
如果我們身處其中.
我們會不會又成為.
一個安提伯一樣的見證人.
在我們生活的環境.
在我們周邊裡面.
做這個見證人呢?.
我們願不願意這樣做呢?.
弟兄姊妹在很多的伊斯蘭穆斯林的國家裡面.

$^{321}$其實都有很多這樣的見證.
我看過一份關於伊斯蘭宣教的刊物裡面.
就是有一位游索夫的一位穆者.
在他年輕的時候.
其實他一直在阿爾及利亞.
這個國家北非回教的國家.
來生長,成長.
但他偶爾在他年輕的時候.
他去到歐洲.
他第一次接觸信耶穌的人.
一班基督徒.
基督徒的生命見證很吸引他.
於是他就回家開始讀聖經.
在他兩年之後再回歐洲的時候.
他又再見到很多的基督徒.
又問了很多的問題.
結果這位年輕的游索夫.
在1980年的時候他就決志.
並且和太太結婚.
兩個都是很忠心的基督徒.
他受神的感動.
於是在1980年.
信主之後一段時間.
他就進入神學院接受裝備.
也成為了警察會的宣教士.
他就回到自己的國家.
阿爾及利亞.
在那裡宣教.
很長的時間.
1990年阿爾及利亞內戰.
游索夫和他的家人.
他沒有離開自己的國家.
他仍然來禱告.
甚至沒有停止他宣教的工作.
不過去到2006年的時候.
阿爾及利亞有一個新的宗教法的頒布.
游索夫就這樣被捕受審.
然後被判監.
經歷了一段很長的時間獲得釋放.
但是去到2000年代.

$^{361}$去到21世紀的時候.
游索夫已經在阿爾及利亞宣教.
接近36年.
在他的團體裡.
他建立了40多間的教會.
不過去到2019年的時候.
他所建立的教會.
就被執法機關強行關閉.
就說這些教會.
違反了宗教建築物的法規.
於是所有的教會都被關閉.
但是游索夫.
這位宣教士.
他沒有放棄.
仍然來跟一群弟妹去敬拜神去宣教.
不過在2023年的時候.
游索夫再次因為非法崇拜罪.
就再被判監兩年.
在他出獄之後.
游索夫這對夫婦.
仍然忠心地服侍阿爾及利亞的信徒.
在接受訪問的時候.
他就說我們遭遇了很多逼迫.
很多壓力.
很多宗教上.
信仰上的艱難.
他說神仍然與他同在.
神知道他在黑暗的狀況裡.
神也知道他作為一個忠心的見證人.
弟妹們.
游索夫的見證會不會提醒我們.
我們要有堅持和勇氣.
去面對不同的法律的壓力.
甚至教堂被關閉.
甚至有監禁的困難也好.
游索夫他仍然堅持來服侍.
體現了他忠心的服侍.
在黑暗裡發光的見證人.
弟妹們.
我們繼續看第二點.

$^{401}$如何保護他的教會呢?.
如何保護別加摩教會.
面對這麼嚴峻的異教情況呢?.
經文第十四節是這樣說.
然而有幾件事我要責備你.
因為在你那裡.
有人服從了巴蘭的教訓.
這巴蘭曾教導巴勒.
將絆腳石放在以色列人面前.
叫他們吃祭偶像之物.
行奸淫的事.
你那裡也有人照樣服從了.
尼哥拉一黨人的教訓.
所以你當悔改.
若不悔改.
我就快臨到你那裡去.
用我口中的劍攻擊他們.
別加摩教會.
就是坐落在一個撒旦之城的地方上.
主耶穌一直與這間教會同在.
他對這間教會各樣的情況瞭如指掌.
教會面對外來的逼迫.
信徒堅守主的命.
沒有棄絕神的道.
他在強硬的逼迫之下沒有妥協.
得到耶穌的稱讚.
不過在這裡教會的教導卻出現了偏差.
所以耶穌就發出了嚴厲的指責.
因為他們偏離了聖經的教導.
經文裡提到的巴蘭教訓.
和尼哥拉一黨的教訓.
他們是服從新約學者Martin Kendall.
他說教會其實有兩大敵人.
哪兩大敵人呢?.
一個就是律法主義.
一個就是自由主義.
就是覺得既然已經信了耶穌.
已經得救.
我們又會向耶穌祈禱認罪悔改的時候.
其實就是一個自由人.

$^{441}$自由的話就應該可以放縱.
所以自由主義就是以自由作為放縱的藉口.
巴蘭教訓和尼哥拉一黨的教訓.
就是屬於後者自由主義.
巴蘭的教訓如果弟子們熟悉.
可以看民數記第22章和第24章.
裡面就很詳細的說到巴蘭這位先知.
聖經裡面說他怎樣誘使以色列人犯罪.
他做很多的奏作.
各樣的狀況.
巴蘭就是因為貪圖利益.
而引導以色列人犯罪.
就好像經文裡面說的.
叫他們隨從外邦神敬拜他們.
參與他們偶像崇拜之後的聚會.
通常這些外邦的敬拜.
偶像崇拜.
一定會有一些淫亂的事.
所以經文耶穌就是責備他們.
在新約裡面彼得後書和尤大書都有說到.
指出巴蘭事件就是對信仰群體的破壞性.
去到一個嚴重的.
偏離了耶穌的教訓.
因為他們貪婪.
偏離了一些的正道.
別加摩教會的信徒.
因為要避過一些的逼迫.
他們轉移追求安逸.
享樂世俗的利益.
而妥協自己的信仰.
將信仰的標準降低.
在這裡耶穌就說.
等同是行淫的事.
這也是一直舊約.
去到耶穌時代.
一個觀念和嚴重的指控.
至於尼哥拉黨的教訓.
也影響教會.
因為他們主張信仰自由開放.
所以參與皇帝的崇拜.

$^{481}$其實沒有所謂.
只是在皇帝的神廟裡.
做一個簡單的行動.
免除很多宗教上的逼迫.
很多的麻煩.
所以在這裡.
耶穌實在是知道.
所以就告訴他們.
指責他們.
說他們偏離了教會的道.
丙姐妹.
我們今天.
身處在自由的環境裡.
我們也會看到很多不同的.
形形色色的崇拜.
偶像的崇拜.
又或者一些很特別的節日.
我相信接下來.
10月31日是什麼日子呢?.
大家都知道.
萬聖節.
這些好像對我們信仰.
沒有什麼影響.
是不是真的沒有影響呢?.
很可能這些就成為.
一些不知不覺的陷阱.
一些隱憂.
一些試探.
甚至很多時候自稱為基督徒的人.
他們選擇和罪惡妥協.
忽視了社會公義.
偏離了聖經的教導.
甚至以一些行為.
否認自己的信仰.
以致可以很方便追求安逸.
可以有個方便.
我們會不會去迎合這些.
而調整自己的信仰?.
無論是別加摩的時代.
或者現在的香港.

$^{521}$又或者其他一些城市.
信仰實踐其實是反映了我們.
生命的優先智取.
和我們的主權是屬於哪一個.
巴蘭的教訓.
尼哥拉黨的行為模式.
會不會不知不覺.
這樣成為我們的習慣.
令我們失去了警覺性.
讓姐妹這套經文提醒我們.
我們應該以神的話來警戒.
我們要自省.
回到聆聽神的聲音.
正如經文第十六節.
耶穌再一次發出警告和勸戒.
所以你當悔改.
若不悔改.
我就快臨到你那裡.
用我口中的劍攻擊他們.
耶穌很嚴正地命令別加摩教會悔改.
顯示悔改對信徒非常重要.
悔改不單止是知識上的認同.
更加是心思與行為的徹底改變.
這是法治內心憎惡的罪惡.
決意要離開錯誤歸向神.
在舊約裡.
大家最熟悉大衛.
大衛在他犯了姦淫之後.
他寫了《詩篇51篇》.
他寫了一篇凝聚悔罪詩.
是一個很好的典範.
他是真誠地懺悔.
以致神潔淨他.
他就能夠更新.
約翰·派伯曾經說過.
悔改的眼淚是有真有假.
不知道大家在神的面前悔罪的時候.
有沒有痛哭流淚.
原來約翰·派伯說.
這些眼淚有真有假.

$^{561}$有些人只是因為罪的後果而流淚.
而不是真正為罪本身的邪惡來痛悔.
有很好的例子來告訴我們.
猶大就是一個很好的例子.
他出賣了耶穌.
於是他感到後悔.
感到懊悔.
他就把三十兩的銀錢.
退還給宗教領袖.
拿回聖殿交給他們.
但猶大的悔改.
只是停留在恐懼和自責裡面.
最終他沒有歸向神.
他採取了一種行動.
來自我了斷.
相反,使徒彼得又如何呢?.
他確實是三次不認主.
因而痛悔而流淚.
但當耶穌呼召他的時候.
他就回轉去牧養神的群羊.
他就成為真正悔改的榜樣.
讓姐妹真正的悔改.
是思維和行為的改變.
不是取悅肉體的那種絲肉.
而是以神為中心.
是要結出與悔改相稱的果子.
耶穌的話就像兩刃利劍.
它可以辨明我們的心.
能夠成為教會的保護和提醒.
在希伯來書裡.
就是這樣形容.
神的道是活潑,是有功效.
比一切兩刃的利劍更快.
它能夠刺入人的魂與靈.
骨節與骨髓.
能夠否開.
連心中的思念和主義都要辨明.
讓姐妹可想而知.
神的話是多麼重要.
我們要好好重視神的話.

$^{601}$黃浸也是重視神的話的教會.
今年年初我們推出了讀經運動.
為期三年就讀完一本聖經.
差不多踏入十個月.
已經差不多完成了.
不知道大家如何看讀經運動.
如何面對每一天讀聖經.
會否成為你的重擔呢?.
不單單是等姐妹你們讀聖經.
我們同工也會安排編配來寫經文.
每一天的分享.
不知道大家有沒有在黃浸Facebook專頁看.
應該也看了十個月了.
不知道大家如何理解.
其實讀聖經就是讓大家更加深入.
理解經文的脈絡.
同時也可以應用在我們日常生活裡.
以致我們不會中了巴蘭教訓.
尼哥拉教訓.
我們能夠抵抗這些似是而非的說話.
當我們這樣做的時候.
因為耶穌已經得勝.
我們也會得勝.
在得勝的經文中.
耶穌有賞詞.
經文第十七節是這樣說.
得勝的,我必將那隱藏的嗎哪賜給他.
並賜他一塊白石.
石上寫著聖靈.
除了那領受的以外.
沒有人能知道.
如果今天耶穌說將嗎哪白石生命賜給你.
你會怎樣?.
因為你忠心.
覺得重不重視呢?.
但原來我相信.
耶穌在這裡是有祂特別的意思.
隱藏的嗎哪是主首先賜給得勝者.
表示作為耶穌的門徒.
我們要著重耶穌自己.

$^{641}$大家都知道嗎哪是出埃及的時候.
以色列人在曠野漂流四十年.
嗎哪就成為他們每天的食物.
但到了耶穌的時候.
耶穌就演繹為.
祂就是從天降下來的嗎哪.
嗎哪是代表生命的糧.
如果你有留意約翰福音的時候.
耶穌就在提醒.
但很多人不接受.
認為耶穌所說的話實在太艱難.
所以跟隨耶穌的人都離開.
耶穌說祂是嗎哪.
只要來到耶穌那裡.
耶穌就能夠讓人在靈裡得飽足.
這就是主賞賜的.
主賞賜給我們.
就是要從祂的口中而出的說話.
就成為我們生命的滿足.
當你是一個得勝者.
我們更需要神的嗎哪.
去被祂餵養.
來防備面對很多黑暗的勢力.
白石又代表什麼呢?.
原來是一種特殊的用途.
在希羅的時代.
白石在選舉的時候很有用.
他投票說.
如果你投白色的就表示贊成.
這也是用於在法庭上.
法庭上設有陪審團.
當他們去判決的時候.
究竟判這個人有罪還是無罪呢?.
白色代表無罪.
黑色代表有罪.
弟兄姊妹,原來白石是代表著.
下臣的稱讚耶穌的接納.
所以耶穌對得勝的信徒.
祂就給了這些賞賜.
你會不會去接受呢?.

$^{681}$至於生命.
祂是要寫在石頭上.
代表有朝一日.
新天新地的時候.
耶穌降臨就要將一切更新.
唯有明字被寫在高揚的生命冊上.
這個明字就是新的明字.
因為我們生命已經不再一樣.
有耶穌.
我們才能夠進入新耶路撒冷.
簡單來說.
這三重賞賜.
是代表著一個經得起考驗而得勝的信徒.
耶穌賜給他們新的身份.
並且與神有一個特別的關係上.
更加親密.
弟兄姊妹.
主耶穌對別加摩教會的說話.
今天對你和我.
仍然有提醒.
因為聖靈向宗教會所說的話.
凡有耳的就應當聽.
啟示錄的訊息.
是要提醒我們.
我們要轉化我們的視覺.
我們要尋求耶穌心目中的教會.
這間是甚麼教會呢?.
我相信在這段經文裡.
告訴我們.
就是3P這個特質的教會.
有耶穌同在.
Presence.
教會的根基.
在於主的同在.
無論在甚麼環境裡.
信徒都能夠經歷祂的引導和安慰.
教會有主的保護.
Protection.
耶穌是以祂的真理說話.
來保守教會.

$^{721}$祂會提醒我們.
在信仰上要持守純正.
不要隨波逐流.
因為主已經得勝.
所以有得勝的主在其中.
這是第三個教會的特質.
Pervail.
教會在主裡得勝.
我們能夠面對任何挑戰.
所以我們信徒能夠站立得穩.
旺朕的弟兄姊妹.
旺朕是你心目中的教會.
會不會好像.
剛才提到的那位牧師發起文卷.
說理想的教會.
你覺得旺朕是理想的教會.
旺朕是不會趕你走的.
因為你參與在其中.
很需要你在一齊攜手.
在耶穌基督裡面.
可以達成這三個P.
我們應該怎樣做呢?.
就是要調整我們的視角.
好好理解教會觀.
是否有主所喜悅的教會特質.
以及我們要實踐主的同在.
在教會的生活裡.
我們積極參與敬拜.
既然你很想教會的敬拜好的時候.
你就積極參與教會的敬拜.
有師班,有敬拜隊.
有團契,有小組.
有不同的服事.
你可以經歷主的同在.
可以彼此鼓勵同行.
甚至為教會禱告.
求主帶領教會成為三P的教會.
亦希望弟兄姊妹關心和支持教會每一個肢體.
建立彼此關懷,相愛的群體網絡.
我在旺朕13年10個月.

$^{761}$數手指計算.
原來我參與旺朕的.
三P教會已經六分一的時間.
今年應該是多少年?.
應該接近65年.
我相信上帝就是這樣.
因代旺角浸信會.
讓旺朕成為神所喜悅的教會.
透過別加摩教會.
我們好好學習和經歷三P.
總有一天會邁向成熟的三P.
我們一起來禱告.
主你的教導.
要教會行在你的心意.
因為你心中的教會.
要我們每一個在不同的經歷中.
我們仍然可以知道.
主你與我們同在.
主你保護我們.
你用你的話語來激勵我們.
縱然我們面對艱難,面對挑戰.
甚至信仰上可能去到妥協的邊緣.
主你都保護我們,挽回我們.
因為你是德性的主.
你要我們每一個跟隨你的都是德性.
你就加你的賞賜給我們.
天賦,我們有軟弱,我們有艱難.
主你全知道.
主我們祈求你與我們同在.
幫助我們進到你自己的教會.
成為你所喜悅的教會.
成為你心中的教會.
奉主耶穌基督的名.
阿們.
\newpage



\section{路加福音 17:11-19}
\label{sec:LTkMkPBreS4}
\textbf{2025年11月2日崇拜直播｜羅潔盈博士｜救恩的距離｜路加福音17\_11-19}
\newline
\newline
連結: \href{https://youtube.com/watch?v=LTkMkPBreS4}{\texttt{https://youtube.com/watch?v=LTkMkPBreS4}} ~~~~ 語音日期: 2025-11-03
\newline
\newline
\hyperref[sec:SlVA2_DGfXk]{\small{< < < PREV SERMON < < <}}
~
\hyperref[sec:index_chronic]{\small{[返順時目]}}
~
\hyperref[sec:5pSwQwIJ1AI]{\small{> > > NEXT SERMON > > >}}
\newline
\newline
路加福音 17:11-19
\newline
\begin{longtable}{cl}
\hline
\hline
章節 & 經文 (和合本修訂版)\\
\hline
17:11 & \begin{tabularx}{0.7\textwidth}{X} 耶穌往耶路撒冷去，經過撒瑪利亞和加利利中間的地區。 \end{tabularx} \\ \\ \relax
17:12 & \begin{tabularx}{0.7\textwidth}{X} 他進入一個村子，有十個痲瘋病人迎面而來，遠遠地站著， \end{tabularx} \\ \\ \relax
17:13 & \begin{tabularx}{0.7\textwidth}{X} 高聲說：「耶穌，老師啊，可憐我們吧！」 \end{tabularx} \\ \\ \relax
17:14 & \begin{tabularx}{0.7\textwidth}{X} 耶穌看見，就對他們說：「你們去，把身體給祭司檢查。」他們正去的時候就潔淨了。 \end{tabularx} \\ \\ \relax
17:15 & \begin{tabularx}{0.7\textwidth}{X} 其中有一個見自己已經好了，就回來大聲歸榮耀給神， \end{tabularx} \\ \\ \relax
17:16 & \begin{tabularx}{0.7\textwidth}{X} 又俯伏在耶穌腳前感謝他。這人是撒瑪利亞人。 \end{tabularx} \\ \\ \relax
17:17 & \begin{tabularx}{0.7\textwidth}{X} 耶穌回答說：「潔淨了的不是十個人嗎？那九個在哪裡呢？ \end{tabularx} \\ \\ \relax
17:18 & \begin{tabularx}{0.7\textwidth}{X} 除了這外族人，再沒有別人回來歸榮耀給神嗎？」 \end{tabularx} \\ \\ \relax
17:19 & \begin{tabularx}{0.7\textwidth}{X} 於是他對那人說：「起來，走吧，你的信救了你！」 \end{tabularx} \\ \\
[1ex]
\hline
\hline
\end{longtable}
$^{1}$今天宣讀的經文是記載在《新約聖經》.
《路加福音》十七章.
十一至十九節.
這裡記載.
耶穌往耶路撒冷去.
經過撒瑪利亞和加利利.
進入一個村子.
有十個長大麻風的.
迎面而來.
遠遠的站著.
高聲說:耶穌夫子可憐我們吧.
耶穌看見就對他們說.
你們去把身體給祭司擦汗.
他們去的時候就潔淨了.
內中有一個見自己已經好了.
就回來大聲歸榮耀與神.
這人是撒瑪利亞人.
耶穌說:潔淨了的不是十個人嗎.
那九個在哪裡呢.
除了這外族人.
再沒有別人回來歸榮耀與神嗎.
就對那人說:起來,走吧.
你的信救了你了.
今天我們很高興有羅傑瑩博士.
在我們當中正在講神的話語.
她今天講的是「救恩的距離」.
羅博士是浸信會神學院的助理教授.
也是我們的義務同工.
她協助我們星期六的黃昏崇拜.
我們恭敬將現在時間交給羅博士.
各位皇室弟兄姊妹主內平安.
很感恩能夠和大家一同.
在領主餐前分享上帝的恩賢.
我們也和身邊的弟兄姊妹說.
主的平安願你同在.
「Hallelujah」是什麼意思呢.
不是四字詞語.
但我們經常在教會.
第一首詩歌大家都唱得很賣力.
在舊約聖經裡有很多不同的讚美詩.

$^{41}$當我們再看今天的經文.
也是講感恩.
多數十一月在北美有感恩節.
我們都知道感恩很重要.
但面對一位什麼都有的.
可能你心目中也有一個朋友.
想送些東西給他.
他什麼都有.
那你要送什麼給他呢.
而我們今天的經文看到.
耶穌很欣賞這位愛邦的撒瑪利亞人.
因為他心被恩感之後.
回去敬拜他.
將榮耀歸給上帝.
天父也是這樣.
雖然他擁有萬有.
他也明白我們一切.
不論是需要.
或是我們未做的事情.
或是在很多困惑裡都有答案.
但天父就是欣賞每一個屬他的兒女.
將這份喜悅.
不單是透過咀唇.
也是透過我們行動去敬拜他.
所以「阿彌陀佛」不是單單四字詞語.
而是不論在什麼環境中.
我們都要讚美耶和華.
因為他配得.
聖經詩篇成書的日期超過一千年.
不是當所有事情都很好的時候.
我們才讚美.
包括詩篇二十三篇都有不同的死刃憂谷.
甚至詩篇裡有不同的哀歌.
但這些都不能難阻我們.
不去彼此對唱「阿彌陀佛」.
即使在被擄.
甚至活得不像人的境況中.
「阿彌陀佛」提醒我們.
呼籲我們的心.
要稱頌耶和華.

$^{81}$不要忘記他的一切恩惠.
但有什麼難阻我們?.
甚至在教會中.
不少弟兄姊妹可能本身有回來.
之後不見了.
有什麼能奪走他的「阿彌陀佛」呢?.
其中一個就是病痛或病患.
在存意的當下或等待的過程中.
作為一個受傷的人.
或曾經受傷的醫治者.
我們如何繼續將心理恩感的訊息.
透過生命去表達呢?.
當耶穌最後一句說.
「你的信救了你」.
我們是否真的有信心.
亦有這種傳言已靠之後的行動.
敬拜的行動去回應主的恩典呢?.
最近我去浸會醫院.
在牛頭角有一個新的專科門診.
護士長打開一個盒子.
裡面有不同類型的貼紙.
有一個由我去到那一刻開始.
我從未停過的小朋友.
眼睛發光.
根本很難選擇是哪一個.
但護士長說「你要說什麼」.
「你不說我不會給你的」.
爸爸在旁邊說「你要跟他說什麼」.
「姑娘送貼紙給你」.
他也掙扎了很久.
但實在太漂亮了.
我都想選擇.
整盒子那麼多.
護士長說「我有很多的」.
「但你不說我不給你的」.
他就很小聲說「多謝」.
之後就選了他最喜歡的.
「Hello Kitty」的貼紙.
而且是立體的.
可以按壓.

$^{121}$我們在教會裡.
在天國裡.
都像這個小朋友一樣.
需要學和上帝說「多謝」.
好像剛才年輕人.
第一句詩歌都很震撼.
「主啊我愛你」.
「I love you Lord」.
我們多久沒有跟天父說.
「我愛你」呢?.
我們多久沒有試過由衷地感恩.
去讚美他.
讚美這位天父.
因為他不單是拯救萬人.
他拯救我.
所以當我們反思我們的生活.
我們的敬拜.
敬拜這個字原文是「俯伏」的意思.
就是剛才這十個都被治癒的大麻瘋人.
其中一位.
他做的動作.
俯伏在耶穌的跟前.
今天是主餐主日.
我們都有一個片刻.
在聖靈裡.
一同安靜禱告.
求聖靈幫助我們.
打開我們的心.
願意在一切的.
可能每一天都是這樣.
或者是.
可能大家都帶著不同的病痛.
或者是需要長照.
即是長期照顧家裡不同的人.
我們都先去求聖靈幫助我們.
去領受祂的恩賢.
我們一同禱告.
天父爸爸.
多謝你.
多謝你帶領我們來到你的殿中.

$^{161}$求你藉著聖靈.
讓我們明白真理.
雖然感恩.
關於大麻瘋醫好的故事.
我們聽過很多次.
求聖靈你此刻開開我們的心.
再次的能夠.
心被恩感.
去歌頌就是你.
由願我口中的言語.
眾弟兄姊妹的心懷意念.
在你面前蒙月立.
禱告奉主耶穌基督得勝的名求.
Amen.
上回我來大家中間.
分享同一處.
耶穌要離開猶太.
因為當時斯贈約翰被下在監獄.
他想回到加利利.
到那裡傳道.
但經文很特別.
就是他必須要經過撒瑪利亞.
其實一個敬虔的猶太人.
會選擇繞過去.
不經過撒瑪利亞.
認為這個地方不潔淨.
甚至現在你搜尋撒瑪利亞.
都是一個廢墟的城市.
但我們今天再看.
比較後面.
盧伽福音十七章的經文.
就是耶穌要回到那裡.
因為他需要回到耶路撒冷.
我們看一看三個大綱.
經文很簡單.
第一就是.
耶穌經過撒瑪利亞和加利利.
中間的麻風村.
這個麻風村有十個麻風病人.
他們都不分什麼國籍.

$^{201}$因為大家都是被邊緣化.
所以裡面有猶太人和撒瑪利亞人.
他們遠遠地呼求耶穌憐憫他們.
這是一個求醫之大的距離.
遠遠地叫.
「耶穌啊!可憐我!」.
之後我們看到.
耶穌吩咐他們給祭司擦汗.
途中得到潔淨.
但祭司認可一個麻風病人.
潔淨也有一個距離.
需要再自潔和擦汗.
真的看到他們康復.
才可以回到人群當中.
相信我們經過幾年的疫情都明白.
當有人咳嗽的時候.
大家都想閃開.
同樣即使被醫治.
但都有祭司或當時文化認可的距離.
最後.
折返的撒瑪利亞人.
俯伏在耶穌腳前.
耶穌對他們說.
「你的信救了你」.
我們也會思考信心和救恩之間的距離.
當我們見到經文之前.
《路加福音》其實耶穌也說.
在整個巴勒斯坦地區.
有很多患皮膚病的人.
有很多患麻風病的人.
甚至舊約也是.
我們稍後會分享.
乃曼被治癒.
他們也同樣經歷過求醫的過程.
但耶穌在《路加福音》前面說.
那麼多人患麻風.
那麼多人我們當他們是皮膚病.
有多少人覺得醫治?.
有多少人可以痊癒到.
像雞蛋般透白.

$^{241}$像鷹蟹般的皮膚?.
只有一個人.
只有乃曼.
所以從這個角度來看.
救恩的距離.
是否千萬分之一?.
當我們對比《路加福音》十七章.
和《列王記》下的經文時.
我們發現這個距離.
有些讓我們可以更立體看這段經文.
可能在麻方村的這十個人.
都等了很多年.
一聽到耶穌要過來的時候.
懷著一個很期待的心.
因為耶穌是一個.
不單單是一個先知.
也是一個醫生.
他治好了不同的人.
黑眼的,耳聾的.
甚至啞的.
甚至瘸子能夠行走.
甚至死人都能夠復活.
同一個地方.
在撒瑪利亞這個地方.
不論猶太人或撒瑪利亞人.
都看同一本我們現在的《舊約聖經》.
都有同樣的期待.
是否真的能夠治好呢?.
我們先看一看.
乃萬去撒瑪利亞求醫的故事.
乃萬是一個將軍.
他身份尊貴.
但卻帶著一個很無助的麻風病.
以色列來的一個女子.
我們不知道她的名字.
不因為自己在苦難中.
在被擄中而失去對上帝的信心.
她很勇敢向主人介紹.
認識這位以色列的神.
可以施行醫治.

$^{281}$所以她對主母說.
「不得主人去見撒瑪利亞的先知.
就必定能治好這個大麻風」.
不過以色列王一聽到的時候很害怕.
「我又不是神.
我怎能夠行這個.
令人得醫治或死人.
可以做回生人呢?」.
他便回去尋求先知以利沙.
而求醫治的過程中.
都有一個距離.
不像這首流行曲.
叫《神愛世人》.
我父親聽說過.
神愛很多人.
很可能數我也有份.
當中去年推出的這首歌.
感動很多人.
甚至很多人說.
好像比起我們的聖詩更加感人.
裡面有三個故事.
包括第一個故事.
就是被燒傷的.
曾經在火災中.
有基督徒的老師.
救了一群小朋友.
是很久之前的事.
大家是否記得?.
那個大火燒傷的老師.
他經歷了人世間很多不同的霸凌事件.
最終他還是選擇了輕生.
第二個故事就是.
一個默默無聞的小人物.
在他的工作崗位中賣珍珠奶茶.
他要結業.
在過程中.
他又看到一個流浪漢很痛苦.
把他唯有的飯盒都給了別人.
所以整首歌感觸的地方可能就是.
我經常覺得.

$^{321}$我給了你.
我自己就沒有了.
第二就是.
可能在苦難當中.
在困難當中.
能夠拉近我們人與人之間的距離.
讓我們更加明白.
以前在旺角也有.
「信耶穌得水牛」.
「得永生」.
或者是台灣很多神愛世人.
黏在不同的燈柱.
或者是電線桿上.
但是這個神愛世人是否聽說.
好像奶萬那樣.
聽說以色列寸草不生的.
撒瑪利亞這個地方.
有一個神醫.
以利沙可以治好我的病.
那就試試吧.
而試的過程中.
他是否立刻很感恩.
很謙卑的態度呢?.
我們每一個人相信耶穌.
都有一個不同的過程.
而過程當中可能都有不同的情緒.
而奶萬一開始的情緒是發怒的.
「為何我來到.
你不是幫我做一場法事嗎?」.
或者是我們看中醫也會把握.
望聞密切.
你也見過醫生吧」.
但是以利沙沒有自己過去.
而且全程奶萬發怒的原因也是.
「為何我一定要來這裡洗澡呢?」.
因為以利沙叫他得醫治.
就是要在約旦河洗澡.
沐浴七次.
「為何我要來這個荒蕪的地方洗澡呢?」.
對於奶萬來說.

$^{361}$他很生氣.
因為期望的落空.
因為文化處境.
甚至他自己本身有的驕傲.
拒絕得醫治的方法.
但這個小女子 這個僕人.
他就勸她.
「神人以利沙也不是要你做甚麼大事」.
「他叫你去洗澡而已」.
透過僕人的勸戒.
她願意謙卑自己.
最後她也得到醫治.
以利沙全程都有一個距離.
有點像我們今天看到的經文.
耶穌沒有親自去摸這些大麻瘋的人.
因為如果摸了他.
他自己也要隔離.
透過當時猶太人結證的禮儀.
耶穌也是遠遠地跟他說.
「你走吧 你去見祭司」.
同樣求醫治的距離.
我們看到.
以利沙打發一個使者去跟乃萬說.
全程都是轉達訊息.
但也讓我們看到.
我們的上帝如果要施行醫治.
不在乎外表的儀式.
而是著重我們有沒有信服上帝.
而透過乃萬真的去約旦河.
沐浴七回之後.
他公開的見證.
除了以色列之外.
普天下沒有神.
神被恩感之後.
他的行動帶來信仰的轉變.
帶來一生人不同的敬拜讚美.
我們留意神的恩典就是這樣.
可能在日常生活當中.
在我們困難當中.
也會突然之間來到.

$^{401}$甚至是透過一些我們沒有期待的.
或是臨到一些很平凡的.
甚至在社會當中被忽略的.
這也挑戰我們.
願不願意放下自己的觀點.
虛心地去聽主的恩言.
讓祂的恩典突破我們的心房.
我們再看今天的經文.
耶穌經過撒瑪利亞.
本身有十個大麻風的人.
遠遠地向耶穌說.
「可憐我,老師,可憐我」.
「主耶穌,請你賜福我們」.
他們知道耶穌是主.
耶穌有能力.
所以他們從這幅畫中也看到.
他們遠遠地看到耶穌.
等到耶穌經過的時候.
就向耶穌去呼求.
除了我們教導小朋友要去感恩.
要去表達愛意.
小朋友有需要的時候.
我們也會鼓勵他們向大人去呼求幫助.
但我們可能年紀越來越大.
或者是因著中國文化.
我們不想麻煩其他人.
我們多久沒有向主去呼求幫助呢?.
而經文繼續.
我們看到耶穌.
面對這十個麻風的病人.
祂可憐他們.
對他們說.
「你們去,將身體給濟斯檢查」.
正去的時候就已經結症了.
而其中有一個.
見自己已經好了.
就回來大聲歸榮耀給父神.
而且還在耶穌面前.
表面的行動是.
所有人都照著耶穌的吩咐去見濟斯.

$^{441}$但十個裡面的這個撒瑪利亞人.
他和另外九個人.
一開始去的方向都不同.
因為猶太人要找猶太人的濟斯.
而撒瑪利亞人需要回到撒瑪利亞.
找自己的濟斯.
但當他發現自己已經好了的時候.
他跪在耶穌面前.
九個人反映了他們可能.
對於痊癒的恩典感到習以為常.
可能雖然一直以來.
他都期待有這個被醫治的神蹟.
但他的代辦事項太多了.
需要去找濟斯.
需要經過被潔淨的時間.
而且他們做漏了一件最重要的事.
就是耶穌今天跟我們說.
我們有沒有將榮耀歸回給父神.
有沒有折返回去.
在當中去感恩呢?.
我們每一個禱告.
不論個人或崇拜當中.
都是奉主耶穌基督的名禱告.
為甚麼呢?.
因為不奉耶穌基督的名.
我們沒有辦法去到父神那裡.
同樣不單單是身體的治癒.
在我們得救.
在我們所謂恩訓清義的過程當中.
我們亦發現.
如果不是透過大濟斯耶穌基督.
我們憑著我們自己.
是沒有辦法讓讚美,感恩.
甚至呼求去到父神面前.
而透過這九個人.
這個感恩的選擇.
我們發現透過感恩.
透過表達一種極度謙卑.
以及承認不是我有甚麼能力.
而是全是恩典.

$^{481}$很珍惜從神而來的拯救.
而且認出恩典的源頭是主耶穌.
而這大於我要去做的不同的代辦事項.
這個折返的撒瑪利亞人告訴我們.
這個回應救恩的行動.
不單回應主的恩典.
回應賜予禮物的主.
更願意與主建立更親密的關係.
感恩讓我們能夠倍增對於主的信心.
而且感恩與忘恩之間的距離.
令我們雖然身體得到醫治.
但這九個人停留在受益的層面.
心靈與主之間可能還有一段距離.
我們每個人的心都有需要被主充滿滿足的空間.
就像中世紀有神學家形容.
每一個誕生之前的靈魂或每個人.
有時心裡都有一份空虛.
為何?因為誕生之前.
上主已經向我們很深的親吻.
而這個親吻讓我們對比世上.
不論是愛情還是日光之下的追求.
心裡都覺得不夠之前的好.
從十個大麻風中只有一個回去感恩的耶穌的反應是.
我們不可以忘恩.
為何被醫治的有十個人.
只剩下你這個外族人呢?.
我們現在又是什麼心態來到父神的面前?.
讓我們再次在領主餐前再深的反思.
感恩不單單是一種信仰的實踐.
亦是讓我們心靈和上帝更親密的操練.
亦願意刻意停下來.
回顧每一天我們經歷的恩典.
讓我們更深的進入全人的敬拜和讚美.
最後以一個故事作為總結.
有一個男孩因為工作壓力很大.
他辭去工作去歐洲旅行.
途中受父母之命.
一定要去見他們的恩人一面.
在沒有辦法的情況下.
他去瑞典這個地方見恩人.

$^{521}$他希望買一份禮物給她.
但口袋裡只剩下幾元歐羅.
他經過一間水晶店.
進去和女售貨員搭訕.
店員很熱情地介紹不同的水晶.
但男孩沒有選中.
一轉身不小心.
他打破了一件很漂亮的水晶.
他心靈很沮喪.
但靈機一觸他跟店員說.
「反正都已經破爛了你又賣不出去」.
「不如你用我手上的歐羅賣給我吧」.
店員真的賣了這個破爛的水晶給他.
他很高興地帶著去見父母的恩人.
見到的時候他又故意裝作不小心.
打破了這個禮物.
恩人很溫柔地跟他說.
「不要緊,有心就行了」.
男孩打開這個包得很漂亮的禮物盒.
打開的時候竟然發現.
這個水晶是被店員很小心地.
分成三份包著.
當他想送給恩人的時候.
那一刻男孩覺得很羞愧.
因為他說謊.
本身買回來的時候.
水晶已經破爛了.
但同時他又深感受到.
這個恩人很有親和力.
亦給他「有心就行了」這一句感動.
所以男孩有點像交功課前的學生.
「拿回我的文件」.
「老師,我明天才交」.
「讓我免費幫你打工」.
「儲夠錢我再買一份給你」.
立即收起這個盒子.
去報答他的恩情.
同樣地,我們的主很希望我們.
能夠得到圓滿的救恩.
包括新心靈的治癒.

$^{561}$他很高興見到我們.
向他獻上感恩.
耶穌欣賞那個折返.
同樣也期待我們.
帶著信心喜樂.
來到他的私恩座前.
事上的與他親近.
歸榮耀於父神.
得著重新得救的生命.
阿們.
\newpage



\section{詩篇 92:12-15}
\label{sec:5pSwQwIJ1AI}
\textbf{2025年11月9日長者主日崇拜直播｜區祥江博士｜老年仍發旺-栽在聖殿中的長青樹｜詩篇92\_12-15}
\newline
\newline
連結: \href{https://youtube.com/watch?v=5pSwQwIJ1AI}{\texttt{https://youtube.com/watch?v=5pSwQwIJ1AI}} ~~~~ 語音日期: 2025-11-09
\newline
\newline
\hyperref[sec:LTkMkPBreS4]{\small{< < < PREV SERMON < < <}}
~
\hyperref[sec:index_chronic]{\small{[返順時目]}}
~
\hyperref[sec:H6vb7uu46Wk]{\small{> > > NEXT SERMON > > >}}
\newline
\newline
詩篇 92:12-15
\newline
\begin{longtable}{cl}
\hline
\hline
章節 & 經文 (和合本修訂版)\\
\hline
92:12 & \begin{tabularx}{0.7\textwidth}{X} 義人要興旺如棕樹，生長如黎巴嫩的香柏樹。 \end{tabularx} \\ \\ \relax
92:13 & \begin{tabularx}{0.7\textwidth}{X} 他們栽於耶和華的殿中，發旺在我們神的院裡。 \end{tabularx} \\ \\ \relax
92:14 & \begin{tabularx}{0.7\textwidth}{X} 他們髮白的時候仍結果子，而且鮮美多汁， \end{tabularx} \\ \\ \relax
92:15 & \begin{tabularx}{0.7\textwidth}{X} 好顯明耶和華是正直的；他是我的磐石，在他毫無不義。 \end{tabularx} \\ \\
[1ex]
\hline
\hline
\end{longtable}
$^{1}$接下來是經訓.
今天的經文是詩篇92篇.
第12到第15節.
經文已經印在你們的秩序表內.
你們可以翻開聽我重讀.
詩篇92篇 12到15節.
二人要發旺於中樹.
生長於黎巴嫩的香柏樹.
他們在於耶和華的殿中.
發旺在我們神的園裡.
他們年老的時候仍要結果子.
要滿了執帳而常發青.
很顯明耶和華是正直的.
他是我的盤石.
在他毫無不疑.
聖經獨筆.
今天我們很榮幸邀請到.
中國神學研究院的副導科教授.
歐長江博士為我們精導.
感覺港大年老仍發旺.
栽在神殿中的長青樹.
有請歐博士.
大姐妹早安.
穆姐請我來講道的時候.
也有一些想法.
平時請我來講.
夫婦主日.
這次請我講長者主日.
其中一個原因是.
我拿了旅遊卡.
做了三個女孫子的公公.
所以也是一個長者.
我想起詩篇90篇.
大家很熟悉.
我們一生的年日是70歲.
如果強壯就可以到80歲.
這段經文年代的壽命可能有些不同.
今天應該改寫.
我們一生的年日是80歲.
如果強壯就可以到90歲.

$^{41}$所以長者的人生階段可長可短.
其實是一個很大的挑戰.
不知道大家對長者的定義是什麼.
我猜65歲以上.
差不多這樣定義.
有點不好意思叫你們舉手.
這裡有多少長者呢.
很多年前我做輔導前.
我是一個職業治療師.
做人體治療師.
看一個人的身體狀況.
我們用幾個英文字母.
PIESV.
P是Physical.
I是Intellectual.
E是Emotional.
S是Social.
V是Vocational.
我看這幾個英文字母.
長者在人生階段.
比較大的挑戰可能是頭兩個.
第一個P Physical.
我們打從五六十歲之後.
多了身體的毛病.
腰酸背痛.
我想教會有時候也有挑戰.
我猜旺嬸也有挑戰.
沒有電梯是否一個挑戰.
長者要上樓梯也是一個挑戰.
可能令長者比較難回到教會.
自己的治理能力.
要不要人照顧.
需不需要坐輪椅.
我想這是其中一個挑戰.
第二個挑戰是.
Intellectual 即是智商.
其實認知障礙是近年有很多的挑戰.
我在輔導工作中.
很多時候遇到中年夫婦或大兄姐妹.
很頭痛自己年老的父母.

$^{81}$他們走失.
有沒有試過走失了父母.
或者他們很多時候忘記了一些東西.
有時候怕別人偷了他們的東西.
對費用很不信任.
到了這個年紀.
其實也容易有認知障礙的挑戰.
我找到一些台灣的數字.
認知障礙在65至75歲還算不錯.
佔認知障礙的人只有3\%.
到了75至79歲就升到7.19\%.
80至84歲就佔13\%.
90歲就佔36\%.
所以到了我們年紀大的時候.
智力上的挑戰也很大.
不過我今天很感恩.
看到長者師班.
剛才在裡面祈禱的時候.
原來有兩位都93歲.
還可以唱歌.
還挺有堅定的力量.
聲線.
所以在人生階段裡.
我們很需要在很多方面去配合.
讓長者有一個好的晚年生活.
今天很想透過一段經文.
彼此鼓勵.
我知道王震在做老人的事工上.
也下了很多心思和活動.
其實是對長者在這個階段.
能夠保持腦筋靈活.
腦筋靈活其實現在也很重要.
身體能走能走其實也很重要.
剛才數的幾樣東西.
待會也會跟大家分享.
我們看看今天的經文.
詩篇92篇.
12節.
可以把經文打出來.
二人要發旺於棕樹.

$^{121}$生長於尼巴嫩的香柏樹.
它們栽於耶和華的殿中.
發旺在我們神的院裡.
這裡說了兩種樹.
其實棕樹在香港也有.
一枝很高的.
上到頂就拆開那種.
大家知不知道那種棕樹.
我沒有拍照給大家看.
這種樹根很深.
有50尺那麼深.
在沙漠裡也能抵抗裂風.
果實也很甘甜.
我沒吃過所以不知道.
是這樣的描述.
令我印象深刻的是.
它的根很深.
老人家的信仰根很深厚.
去到很遠.
而且它們不容易被打倒.
第二種樹是香柏樹.
香柏樹在香港應該沒有.
有點像榕樹的葉.
但它的特點是巨木.
很大棵.
高到可以有120尺.
差不多有10層樓高.
而且它有1000年壽命.
令我們想起香港有什麼樹.
可以跟它們比美.
老的榕樹.
不過很可惜這次颱風.
外文村的老榕樹被打倒.
用樹來形容老人家長者也很合適.
這兩種樹一棵很高一棵很大.
一望到就知道它們已經有一年.
老人家也是.
外貌站出來其實有經歷很多風霜.
念想上也有閱歷.
一望下去就知道它們有歷史.

$^{161}$這裡是說它們發旺和生長.
即是老人家仍然有成長空間.
教會也辦很多興趣班.
做一些整腦.
很多時候年青想念工作.
想念養兒育女.
很多興趣放下了.
年長的時候可以學跳舞.
可以耍太極.
可以做餅.
其實也是一個成長和學習的機會.
它們發旺和生長.
那這裡的災在哪裡呢?.
災於耶和華的殿中.
發旺在神的院裡.
這裡特別提到.
我們說一個人的成長.
你說成長或災中是成功的土壤或健康的土壤.
事業,財富,體力和健康.
都是我們很多人的追求.
如果災中在那些地方.
也是好的.
但這裡說災中在哪裡呢?.
在耶和華的殿裡.
在神的院裡.
就等於長者們.
你們數十年.
每天來聚會嗎?.
每天來崇拜嗎?.
有吧?.
有沒有每個清晨去靈修?.
有沒有每週來教會祈禱會?.
這些都是在教會裡的紮根機會.
屬靈上的紮根.
聽神的話.
堅定了我們對信仰的把握.
我們不要小看.
每個星期,每週的參與.
而參與在神的教會裡.
這裡用災於.

$^{201}$有人種在這裡.
那是誰呢?.
是神自己.
神把我們放在教會的群體裡.
不知道大家有沒有留意.
近年世衛組織.
久不久也有關於醫療的報告.
現在最大的健康敵人是什麼呢?.
大家知道嗎?.
是寂寞.
Loneliness.
寂寞.
原來是很影響我們身心靈的健康.
六個中有一個人.
經常會經歷寂寞.
而老人家寂寞.
常見嗎?.
其實也很常見.
特別這幾年.
政治事件令很多中年夫婦.
帶著子女移民.
不知道黃浸的情況有多嚴重.
其實就留下老人家在香港.
老人家有兒孫抱抱是很開心的.
我自己可以見證.
抱著孫子什麼都不用做.
很滿足的.
雖然沒有FaceTime.
可以看螢幕.
但跟你抱著他.
給他吃東西.
那種經歷是完全不同的.
所以大家可以在這方面想想.
教會已經幫到忙.
因為回教會的過程.
每個星期都見到面.
問安.
現在有些配對計劃.
有些幼稚園會找老人家去探望他們.
抱他們.

$^{241}$其實是相得益彰.
會不會我們也有老人家.
可以參與兒童侍工的機會呢?.
可以讓他們照顧孩子.
抱抱.
其實是減少很多寂寞的感覺.
我想年青的弟兄姊妹.
其實也要有這個意識.
我們身邊的老人家.
最容易感到寂寞的是兩批人.
就是青少年和老人.
這兩批人都是我們可以刻意去想想.
在神的院裡面.
我們有一個大家庭.
我們有不同年齡階層的人.
教會很流行分齡牧養.
分開不同年齡層.
可能久不久都要打破這些界線.
不同年齡的人應該有更多接觸.
讓老人家有熟齡的兒孫在身邊.
在神的院裡面減少很多寂寞的情況.
我們看另一節經文第十四節.
第十四節.
這節經文我自己很喜歡.
不如我們一起讀吧.
第十四節開始.
他們年老的時候仍要結果子.
要滿了執像而上發青.
其實現在我們很多時候都追捧年輕.
所以年老這個問題.
aging這個問題.
就好像成為一種陰影.
我們追求進步.
所以衰退是一個陰影.
我們追求活力.
所以容易疲勞是一個陰影.
我們追求很多進步等等.
所以當老人家被人說.
你走得這麼慢.
你都沒有用.

$^{281}$你拖累人.
這些對比.
其實對老人家來說.
好像不是很好的說話.
走得慢有什麼問題.
你走得這麼快.
是不是經常都要進步.
是不是經常都要這麼精神.
老人家是多了一份懶教.
其實是會的.
是一個人生階段.
所以我們教會不應該追捧年輕.
年老有我們自己的生活節奏.
我們的特色.
這裡特別的地方是什麼呢.
年老的時候仍然要結果子.
「型」字.
年輕的時候或者成年的時候.
我們生產力很高.
做很多事情.
年老仍然要結果子.
我們不是要做很多事情.
我想結果子的做法可能有些轉變.
結果子可以是你懂得說一些安慰的話.
或者看到頂尖有需要.
你願意為他祈禱.
或者你自己生命就是一個見證.
其實老人家的結果子.
可能就是我自己.
現在這個人生階段.
我有個孫女很喜歡我.
經常叫我「公公」.
我覺得我抱著她.
給足夠我們現在流行說的attachment.
相依的關係.
我給她足夠的相依的關係.
讓她有足夠的安全感.
和我在工作裡教書或者寫作.
其實可能那個功效.
就是對我這個孫女這個人本身.

$^{321}$發揮了一個很大的作用.
和我寫書或者做什麼比較.
可能在我這個人生階段.
我更加enjoy抱著孫子.
明不明白我的意思.
其實隨著人生的閱歷不同了.
我們看的東西都不同.
這裡特別的地方就是汁醬.
原來還很juicy.
我最深刻的.
哪種水果是有這種形態的呢?.
大家有沒有喝過熱情果?.
或者叫百香果?.
百香果或者熱情果.
它的果是會皺皮的.
其實太新鮮反而不好.
皺皺的皮.
然後打開它.
裡面的百科果.
最好可以拿來開梳打水來喝.
是很好喝的.
很甜美的.
老人家有一點像這些.
雖然我們外表好像皺皺的.
有些柴柴的皮.
但裡面還是很rich的.
很豐富的.
很甜美的.
我想如果用汁醬來說.
其實我們這個人生階段.
是儲了很多人生的智慧.
只不過沒有人來問我們.
那些人不懂事.
不要來問我們.
當然我們有時候會長氣一點.
如果我們懂得很精簡地回應別人.
就不要長氣.
說來說去的那些.
其實我們儲了人生智慧.
對年輕人其實是很有幫助的.

$^{361}$是嗎?.
有些教會也是做這些配對.
找一些人生下半場已經過了的弟兄姊妹.
跟一些年輕人初出茅廬的配對.
譬如你做會計的.
找一個會計師跟你配對.
可以傳遞一下人生的閱歷.
智慧給你.
商場或者職場是怎樣.
其實一個這樣的機會.
我覺得老人家不只是拉拉筋這麼簡單.
可以做很多事.
其實可以擴闊一下.
我們能夠給年長.
貢獻他人生閱歷的地方在哪裡呢?.
所以我們要花點心思去想.
其實年老是不是一定失去了創意.
或者那種精力呢?.
我舉兩個藝術家的例子給大家聽.
一個大家很熟悉的.
叫齊白石.
知不知道齊白石.
他八十歲的時候.
轉換他的畫畫的興趣和筆鋒.
可能之前畫多點山水.
原來他八十歲之後就畫蝦.
他九十歲的時候自稱自己畫蝦是第一.
你上網打一下齊白石和蝦.
你會看到他畫的畫作是非常出色.
他的形態等等是很厲害的.
八十歲才開始畫.
厲害嗎?.
也很厲害.
另外一個就是Monette.
大家知不知道他畫瑞年外國的西洋畫家.
記載他從七十六歲開始有白內障.
這是老人家的問題.
有白內障.
誰不知白內障之後.
他的筆鋒就變了.

$^{401}$可能看得不太清楚.
所以就大筆一點.
誇張一點.
所以他畫的瑞年畫得很漂亮.
很出名的.
他說是透過他老了.
看得不太清楚.
畫出來的那些就更加接近抽象主義的傾向.
所以就成為了抽象主義的橋樑.
所以我們不要小看.
老了還有很多東西可以貢獻.
特別烹飪方面我覺得也是.
老人家煮的東西失傳了.
應該傳承下去.
其實有很多地方我們可以.
還有很多執著在這裡.
執著常青.
最後想和大家看第十五節.
第十五節這裡說.
「好顯明耶和華是正直的,他是我的盤石,在他裡面毫無不意.」.
這裡特別提到.
「好顯明耶和華」.
其實老人家站了出來.
就是一個見證.
剛才我看到長者師班站出來.
就是一個見證.
唱到這個年紀還能夠歌唱.
其實活的見證是很不容易的.
我想老人家其實是經過戰亂.
可能又經過貧窮,捱過苦,病痛.
甚至喪嘔都未定.
但仍然能夠站出來去讚美神.
他是我的盤石.
所以今天我很欣賞長者師班選的那首歌.
其實那首歌的歌詞很有意思.
「沿路是你,叫我信心更堅定,一生不變,天天賜下恭應.」.
「沿路是你,叫我信心更堅強,一生指引,幫我練強心智.」.
就是說我們信主的年日這麼長.
一直都沒有離棄神.
一直都在教會裡面生活,成長.

$^{441}$這已經是一個很大的見證.
所以長者們.
你今天應該對身邊的年輕人說.
孩子,我走過的路是證明神的信實.
你都要去信靠祂.
我們的生命就應該要發揮這個見證的作用.
有人說其實「沿路」這句話很值得我們好好珍惜.
我看過一篇新詩.
是一個巴西的詩人寫的.
他形容自己好像一個拿著一盤櫻桃的男孩.
一開始的時候,他吃這些櫻桃滿口的.
但吃到差不多的時候,他就說.
當我意識到只剩下幾顆的時候.
當然我們人的歲數總是會有限制.
可能到某個年紀,我們都很難抵擋這個衰殘的事實.
我們不是今天要永遠常青.
我看我自己的媽媽.
我媽媽90歲的時候就離世了.
她在80歲的時候,真的還是很活躍的.
經常去老人中心讀書.
老人中心又頒四方帽給她畢業.
她是很開心很活躍.
到了85歲之後.
就出現了認知障礙.
上巴士上不到,腳痛沒力氣.
臨尾腎衰竭.
我想這是人生必經階段.
我們都要面對這些身體衰殘的事實.
但在未衰殘之前.
我們還有好幾年的日子.
就好像這首詩人的說.
當我意識到我還剩下幾粒的時候.
我開始細細品味.
老人家在他離世之前的歲月.
應該是幸福的.
應該是開心快樂的.
應該是被孝順的.
應該是被服侍的.
因為我們在上半生付出了很多.
很需要被人尊榮的.

$^{481}$大家對白髮怎樣看?.
喜歡白髮嗎?.
又有人轉頭.
如果我將來白髮.
我現在也有一點,不算很多.
就算白了,我也不會贊成染髮.
一來染髮是一種味道.
二來我覺得白髮挺漂亮的.
大家有沒有覺得白髮漂亮?.
很漂亮,特別是白髮加上白鬚.
厲害.
這裡有很多長者頭髮都白了.
我很喜歡,我希望我的頭髮白得快點.
白髮是一種尊榮.
結束的時候,很想用一首.
不知道大家有沒有聽過內地的歌.
有一首叫《照照》.
寫了一首歌.
《當你老了,頭髮白了》.
有沒有聽過?.
(唱歌).
聽過吧?.
那首歌我也挺喜歡的.
那首歌是這麼說的.
當你老了,頭髮白了.
睡意紛沉.
當你老了,走不動了.
爐火旁邊打頓.
回憶青春.
多少人曾愛你.
青春歡暢的時辰.
愛慕你的美麗.
假意或真心.
其實我們老人家也有青春過的時候.
當我們老的時候.
身體狀況,有時會打瞌睡.
或者在打喊.
其實都是在回憶青春的歲月.
青春的時候好像很多人愛你.
稱讚你漂亮.

$^{521}$但那首歌說.
只有一個人還愛你虔誠的靈魂.
愛你蒼老的念想的咒文.
這首歌其實是.
有一位愛爾蘭的詩人叫葉慈.
他原作的.
《照照》只不過是改編了.
不過其實改編是改了意思.
他用一個兒子懷念媽媽的心情寫的.
但葉慈自己本身的那首歌.
是比較悲慘的.
他曾經.
可能是他老公.
曾經有一個人說愛他.
愛他虔誠的靈魂等等.
連那個人也沒有了.
葉慈那首詩.
是說孤獨的那件事是更加明顯的.
不過我自己就想借助這首歌.
鼓勵一些長者們.
只有一個人還愛你虔誠的靈魂.
愛你蒼老念想的咒文.
那個是誰?.
神自己.
神仍然很愛我們虔誠的靈魂.
多點回教會.
在神的院裡面.
我們驅散了寂寞之源.
我們的靈性根可以更加深厚.
愛你蒼老念想的咒文.
其實我們念想的咒文或白法.
是在說什麼呢?.
其實是在說神信實的痕跡.
是不是?.
神信實的痕跡.
所以不用介意念想的咒文.
那些咒文是神恩典的記號.
祝福老人家活到老.
學到老.
快樂到老.

$^{561}$多謝各位.
\newpage



\section{啟示錄 2:18-29}
\label{sec:H6vb7uu46Wk}
\textbf{2025年11月16日崇拜直播｜陳明泉牧師｜耶穌給教會的七封信:提防假先知｜啟示錄2\_18-29}
\newline
\newline
連結: \href{https://youtube.com/watch?v=H6vb7uu46Wk}{\texttt{https://youtube.com/watch?v=H6vb7uu46Wk}} ~~~~ 語音日期: 2025-11-17
\newline
\newline
\hyperref[sec:5pSwQwIJ1AI]{\small{< < < PREV SERMON < < <}}
~
\hyperref[sec:index_chronic]{\small{[返順時目]}}
~
\hyperref[sec:Ndl2CyA1nlc]{\small{> > > NEXT SERMON > > >}}
\newline
\newline
啟示錄 2:18-29
\newline
\begin{longtable}{cl}
\hline
\hline
章節 & 經文 (和合本修訂版)\\
\hline
2:18 & \begin{tabularx}{0.7\textwidth}{X} 「你要寫信給推雅推喇教會的使者，說：『神的兒子，那位眼睛如火焰、雙腳像發亮的銅的這樣說： \end{tabularx} \\ \\ \relax
2:19 & \begin{tabularx}{0.7\textwidth}{X} 我知道你的行為：愛心、信心、勤勞、忍耐；又知道你末後所行的善事比起初所行的更多。 \end{tabularx} \\ \\ \relax
2:20 & \begin{tabularx}{0.7\textwidth}{X} 然而，有一件事我要責備你，就是你容忍那自稱是先知的婦人耶洗別教唆我的僕人，引誘他們犯淫亂，吃祭過偶像之物。 \end{tabularx} \\ \\ \relax
2:21 & \begin{tabularx}{0.7\textwidth}{X} 我曾給她悔改的機會，她卻不肯悔改她的淫行。 \end{tabularx} \\ \\ \relax
2:22 & \begin{tabularx}{0.7\textwidth}{X} 看吧，我要使她病倒在床上。那些與她犯姦淫的人若不悔改他們的行為，我也要使他們同受大患難。 \end{tabularx} \\ \\ \relax
2:23 & \begin{tabularx}{0.7\textwidth}{X} 我又要殺死她的兒女，眾教會就知道，我是那察看人肺腑心腸的，我要照你們的行為報應各人。 \end{tabularx} \\ \\ \relax
2:24 & \begin{tabularx}{0.7\textwidth}{X} 至於你們其餘的推雅推喇人，就是一切不隨從這教訓，不明白他們所謂撒但深奧之理的人，我告訴你們，我不會再把別的擔子放在你們身上。 \end{tabularx} \\ \\ \relax
2:25 & \begin{tabularx}{0.7\textwidth}{X} 你們只要持守那已經有的，直到我來。 \end{tabularx} \\ \\ \relax
2:26 & \begin{tabularx}{0.7\textwidth}{X} 那得勝又遵守我命令到底的，我要賜給他權柄制伏列國； \end{tabularx} \\ \\ \relax
2:27 & \begin{tabularx}{0.7\textwidth}{X} 他必用鐵杖管轄他們，如同打碎陶器， \end{tabularx} \\ \\ \relax
2:28 & \begin{tabularx}{0.7\textwidth}{X} 像我也從我父領受了權柄一樣。我又要把晨星賜給他。 \end{tabularx} \\ \\ \relax
2:29 & \begin{tabularx}{0.7\textwidth}{X} 凡有耳朵的都應當聽聖靈向眾教會所說的話。』」 \end{tabularx} \\ \\
[1ex]
\hline
\hline
\end{longtable}
$^{1}$《聖經》最後一卷書 啟示錄.
啟示錄第二章十八至二十九節.
如果是唐坐的聖經 第351頁.
請聽我讀出.
你要寫信給推拿教會的使者說.
那眼目如火焰 腳像光明銅的神之志說.
我知道你的行為 愛心 信心 勤勞 忍耐.
又知道你末後所行的善事 比起初所行的更多.
然而 有一件事我要責備你.
就是你容讓那自稱是先知的婦人耶駛別.
教導我的僕人 引誘他們行奸淫.
吃濟偶像之物.
我曾給他悔改的機會.
他卻不肯悔改他的淫行.
看啊 我要叫他病餓在場.
那些與他行淫的人 若不悔改所行的.
我也要叫他們同受大患難.
我又要殺死他的黨內.
叫宗教會知道 我是那察看人 肺腑心腸的.
並要照你們的行為 報應你們各人.
至於你們推也推拿其餘的人.
就是一切不從那教訓.
不曉得他們素上所說撒旦深奧之道理的人.
我告訴你們 我不將別的擔子放在你們身上.
但你們已經有的 總要持守 直等到我來.
那德性 有遵守我命令到底的.
我要賜給他們權柄 制服列國.
他必用鐵匠 核管他們.
將他們如同搖滬的瓦器 打得粉碎.
將我從我父領受的權柄一樣.
我又要把神聖賜給他.
聖靈向宗教會所說的話.
凡有疑的就應當聽.
聖靈向宗教會所說的話.
凡有疑的就應當聽.
請姐妹平安.
平安.
我們今天繼續用啟示錄.
七教會的書信.
來看看耶穌基督對教會有什麼重要的訊息.

$^{41}$不單止對當代聖經的處境的教會說.
也是對我們今時今日.
香港弟兄姊妹生命中需要的訊息.
他也對我們今天宣講.
在我們看這段經文之前.
想像一個情景.
假設今天是長者日.
你的家人跟你一起去一間餐廳慶祝.
跟你去做一下節.
做一下長者節.
假設去到一間很華麗的酒店吃自助餐.
那個地方又很舒服.
坐得很豪華的坐下.
又有很熱誠的招待員來招待你.
環境又好.
還有人在拉小提琴.
是不是很棒.
然後你吃自助餐.
一直看過去.
那裡擺放食物的.
全部都是一些最好的海鮮.
不同的食物.
很開心.
你去夾的時候.
夾一塊三文魚吃.
一直吃 吃得很開心.
你吃著吃著的時候.
眼睛又再看回去.
拿東西吃的三文魚.
吃海鮮的地方.
你見到有一隻米奇.
老鼠.
在三文魚堆裡.
你吃的時候.
他又吃.
他又擔著幾塊.
然後你想吃的生蠔.
米奇在生蠔堆裡游泳.
你看到這個情景.
你會有什麼反應.

$^{81}$你會不會說繼續吃.
一定不會的.
然後你就跟當場的老闆.
或者經理來做投訴.
經理就說.
稍安毋躁.
我們這間酒店這麼大.
養幾隻可以用的.
你用得下幾隻米奇.
在我們的酒店當中.
你覺得你聽到這句話.
你會有什麼反應.
你會繼續吃.
還是你會很生氣.
離開這個地方.
這是一個虛構的情景.
不過這個情景.
你用這個角度來看.
今天我們所看的經文.
今天這段經文所說的.
一個地方的教會叫推亞推拉教會.
推亞推拉這個地方.
其實是一個小的城市.
但是他們面對著一個.
被耶穌基督責備的.
就是在說.
他容讓一個叫做.
自稱為先知的女士.
在當中帶著一份權柄.
將屬世的價值帶到教會當中.
究竟這段經文.
對於當時的教會.
或者對我們有什麼重要性.
今天希望用這段經文.
來勉勵我們去調教我們.
我們先禱告求上帝幫助我們.
同心祈禱.
求主與我們同在.
幫助我們.
讓我們一時之間去讀這段經文的時候.

$^{121}$讓我們不單止心裡面領受.
更加是時刻被你的話語.
調教我們的心思意念.
我們恭謹的去到你的根前.
幫助每一位聆聽的弟姐妹.
他們每一個能夠聽得明白.
又在生活當中實踐.
幫助宣講的講得清楚.
願你的話語帶著權能.
帶著謙卑.
在我們當中來彰顯.
求你與我們同在.
帶領我們下來領受你的話語的時間.
謙上禱告奉靠基督耶穌得勝的聖名祈求.
Amen.
我們一起看看這段經文.
這段經文裡面就說到.
推亞推拉教會使者領受耶穌基督的書卷.
首先說一下推亞推拉是甚麼地方.
其實在啟示錄裡面.
他寫的書卷其實是按著地域.
來去書寫的.
所以地方就由南去到北.
然後又再順時針地下去.
所以是一個.
好像順時針地走.
而推亞推拉教會.
就不是去到最北的那間教會.
他就比較南一點.
最北的就是別加摩教會.
南一點的就是推亞推拉這個城市.
推亞推拉這個城市其實不是一個很大的城市.
不過他有一個重要性.
就是他就是那些叫做交通的要道.
舊時的羅馬帝國.
他的首都就是.
下一次所說的宣講裡面的別加摩城.
就是小亞細亞和在羅馬帝國裡面的首都.
而這個首都以南就是推亞推拉這個城市.
所以這個地方.

$^{161}$他成為了一個好像中轉站也好.
或者一個交通的交匯點.
每一個進去的首都都要經過這個地方.
除此一般人的運輸或者做生意要經過這個地方之外.
這個地方也是守護著.
這個舊時的首都別加摩的一個重要城鎮.
如果舊時的那些外敵要攻打.
打首都的時候.
他一定要經歷過這個推亞推拉這個城市.
所以這個地方雖然不是很大.
但是他發揮了一個就好像一個大閘門一樣.
他保護著後面的首都.
同時間也是營商的一個好地方.
因為直著他就可以去到北一點的地方.
進到首都裡面做生意.
你明白了這個地理位置之後.
這個地理位置就造就了他們有一些的整個城市的發展.
其中一樣東西就是除了軍事要塞之外.
他在這裡也成為了一個商業的集散店.
在當時來說他有幾樣東西很出名的.
第一樣東西就是紫色的布匹.
紫色的布匹如果你想起.
其實在《使徒行傳》裡面有提及過的.
《使徒行傳》第十六章十四節裡面.
保羅聽到馬其頓的呼聲.
就去到歐洲裡面傳福音.
第一個見到的就是一個賣紫色布匹的婦人.
她的名字叫做呂底亞.
呂底亞是什麼人呢.
就是推亞推拉的人.
所以這個地方其實是出產很多紫色的布.
他的染料這種紫色是很上乘又很抵買的.
所以就令到這個城市專門出產這東西.
就成為了一個重要的商業的命脈.
做生意是一個很重要的幫助他們的城市發展.
第二樣東西就是這個地方會出.
那些重要的冶金的一些技術.
所以有不同的貴重金屬.
最主要是銅器.
這些金屬就幫助這個地方.

$^{201}$既是一個軍事要塞.
同時間也是藉著這些貴重金屬.
來支持他們整個社會上的命脈.
也有一些考古學家說.
這個地方也出名做一些陶瓷.
因為這個地方的地土是一些黏土.
藉著這些黏土就做這些陶瓷.
你不要看輕剛才說的這些歷史背景.
其實這些所有的歷史背景都放進去.
在剛才我們待會要讀的那段經文裡面.
他用這些背景就說到耶穌基督對這個城市的教會.
說重要的話.
我們一起知道這個這樣的狀況之後.
我們就要想這個商業城市有什麼令到上帝這麼討厭呢.
或者要令到他做商業沒有大問題.
不過在當時的社會來講.
做生意沒有今天香港這麼容易.
他們做生意不論是賣布也好.
做金屬的冶煉也好.
又或者賣陶瓷.
他們要令到那個社會大家都共同應對.
或者一個小城市大家要集合力量.
他們就會組織一些叫做商會.
進了一個商會之後.
他才可以有條件來對他做生意.
不過進了這些商會他未必一定很苛刻.
但有一樣東西很多時候當時的人做生意就避不了的.
就是很多的逢場作慶.
今日香港社會沒有那麼多了.
因為經濟不好.
八九十年代的時候香港做生意就經常逢場作慶.
在當時的推也推也有這樣的逢場作慶.
不過舊時的做法就是喝酒喝喝吃吃.
然後就去到當時的推也推也裡面的神廟.
裡面來做祭祀.
吃的肉本身就是祭祀過偶像的祭肉.
然後喝飽吃醉之後就飽暖思淫慾.
他們真的用一種叫做男女交合的方式.
來敬拜外邦的神明.
所以你看到那個生活是很荒誕的.

$^{241}$他們做生意就做生意.
不過他們要附帶就是要吃他們的祭肉.
要跟隨他們的民風習俗.
甚至乎完了這些飲飲食食之後.
還要有一些男女亂七八糟的關係.
在當時的神廟裡面有一些叫做女祭師.
這些女祭師就幫助人敬拜.
但實質上那些就是廟祭.
就是用一種這樣的方式來敬拜神.
這些林林總總成為了我們去研讀推也推也教會面對的挑戰.
一個很重要的內容.
我們了解了背景之後我們就看.
耶穌基督說那眼目如火焰.
腳像光明瞳的神之子說.
我知道你的行為.
愛心信心勤勞忍耐.
也知道你末後所行的善事.
比起初所行的更多.
值得注意的是這段經文裡面.
耶穌基督怎樣去自稱自己.
怎樣去稱呼自己.
耶穌形容自己是眼目如火焰.
腳像光明瞳的神之子.
首先講一下神之子在整個啟示錄裡面.
原則上用這個字眼只有這裡.
神之子的強調就是講到耶穌基督的超然的權柄.
耶穌基督某程度上好像坐了一隻大碟出來一樣.
他講到他是用神的兒子的角度.
來宣講以下的信息.
所以他用的規格是極之高的.
而這個極之高的規格.
他不是在講.
你想一下有時候講說話要去到一個這麼高的規格.
大概他要講的就是一些比較嚴厲狠的話.
經文裡面再繼續講.
耶穌用這個真實來自上帝的權柄.
來面對世俗的商會.
或者是當時社會裡面拜偶像.
甚至乎其他的權力.
合制住的那一群人.

$^{281}$你面對神和面對人兩種的權力.
你用什麼角度去看呢.
耶穌在這裡就講到他的權柄來自上帝.
他是神的兒子.
然後他再用兩個形容詞來形容自己.
第一就是眼目如火焰.
其實這個字眼在早一點的篇幅.
在啟示錄開頭的時候.
一章十四節已經講了.
但一章十四節裡面都是一個比較籠統.
講到末世之後.
耶穌基督所作的審判.
來到這裡再用這個字眼.
眼目如火焰就是講到一種耶穌的憤怒.
耶穌帶著一種嚴厲的角色.
來對著這群教會的使者講話.
如果你覺得眼目如火焰.
你都覺得好像很抽離.
用廣東話你會好一點.
就是眼火爆.
你記住耶穌是眼火爆.
來對著推拉推拉的弟兄姊妹.
來講這番話.
所以耶穌在這裡是惡的.
是憤怒的.
經文第二個用詞就是腳像光明銅.
同樣在開頭的篇章裡面.
一章十五節都已經有講.
都是用腳像光明銅的神之指.
或者是耶穌基督人指.
來對教會講話.
腳像光明銅是什麼意思呢.
不是腳用銅做.
而是有一些講法.
一些色經學家講到.
他的腳是用一些貴重.
硬的金屬靴.
穿上一個光明銅穿的.
很堅硬的軍靴.
來表達他的行動.

$^{321}$光明除了是講那種榮耀感之外.
更加是講到他那種堅固.
硬淨是一種能力的表現.
用這個堅固的光明銅.
散到光亮的靴子.
來踐踏仇敵.
這些用詞全部都是一些很強烈的字眼.
講到耶穌是憤怒.
帶著能力權柄.
而且是一種審判的角度.
來講接下來的內容.
講一個什麼內容呢.
19-20節.
就講到推拉教會的實況.
我知道你的行為愛心.
信心,勤勞,忍耐.
又知道你末後所行的善事.
比起初所行的更多.
然而我有一件事要責備你.
就是你容讓那自稱是先知的婦人耶洗別.
教導我的僕人.
引誘他們行奸淫.
吃偶像之物.
當你看到這個19-20節的時候.
你就掌握了.
原來推亞推拉.
耶穌要用這麼嚴厲的角度.
原來他們這個教會上上下下.
或者整個群體裡面.
他們正在面對一個挑戰.
不過在挑戰之前.
耶穌也肯定這個教會是一個什麼樣的教會呢.
他們行事是有好的行為.
他們是一群有愛心的人.
他們相信的是有信心.
又勤力,又忍耐.
這些是很好的素質.
經文再補充一句.
他們末後所行的善事比起初所行的更多.
這間教會是不斷進步.

$^{361}$不斷進步地實踐愛心,信心,勤勞,忍耐.
還有他們有美好的行為.
這間教會在眾人角度來說.
是第一.
這麼多樣好的行為.
如果你想像推亞推拉是一個什麼樣的人.
就是一個謙謙君子.
又勤力,又勞人,又愛心.
甚至乎陰森細氣.
做每一樣事默默耕耘.
你看到他永遠都是笑笑笑.
不要緊,在我身上.
OK,沒問題.
亦來信受,非常好.
不過,經文裡面.
他講到這個外在行為和他內在有一些出入.
這間教會裡面的內在信念是很缺乏的.
甚至乎他講到這間教會.
因為太好人.
好人到一個地步.
就是讓一個自稱為先知的婦人耶洗別.
來教導其他的弟兄姊妹和他們.
他在講僕人,甚至乎是教導其他的領袖.
簡單來說.
這間教會永遠都是一個好人的形象.
不過面對著是與非.
以及對於聖經的價值觀念.
他不是很持守.
大概我們今天用的直白一點的話.
就是這間教會是一班爛好人.
他們是好人,非常好的人.
但是他們對於上帝的要求和信念.
他們某程度上可能沒有勇氣直接去面對.
甚至乎他們去到一個地步.
他們不停地,教會領袖不停地退讓.
容讓一個自稱為先知的婦人.
耶洗別是耶穌基督為這個婦人給的一個名字.
待會再講解耶洗別為什麼用這個字.
但是他在講教會裡面將一個重要的教導權.
容讓一個他自稱為先知的人.

$^{401}$來成為一個重要的教導的身份.
甚至乎不單止散播畢地女來做.
而且教會要搭一個台給他.
來做一個重要的教導的角色.
耶穌基督看到這件事就真的眼火爆.
因為整個教會裡面一班好的人.
但是被一個錯誤的內容.
這個自稱為女先知的耶洗別來發揮影響力.
耶穌基督不是無的放矢地說.
不是在講他的性別問題.
而是在講這個自稱為先知的耶洗別.
他做了什麼教導呢.
引誘他們行奸淫吃製偶像之物.
簡單來說這個不是純粹吃喝的問題.
很多的色經學家他都特別去講到.
這個吃偶像之物和行奸淫.
就是回應剛才說的推亞推拉這個地方.
他們是做生意的.
他們做生意的時候就需要進入那些商會.
進入那些商會他們要有一些的關係.
有一些的網絡來連結的.
當他們連結了這些網絡的時候.
他們就飲飲食食.
而飲飲食食特別是飲的酒.
或者吃的那些製肉.
全部都是拜過當時偶像之物.
來再去吃的.
所以這件事對於早期教會來說.
或者很多時候當時的人都覺得.
拜偶像算不得什麼.
偶像是假的你拜過他就由得他.
那個物質沒有改變的.
你的心不是這樣就照吃.
但是在當時來說這種說法.
成為了一種風氣.
甚至對於上帝所說是不可以的.
在飲食在生活細節裡面.
就流露你是一個敬虔的人.
是不是從生活裡面分別出來的這些行為.
他就淡化了.

$^{441}$所以在這裡首先吃偶像之物.
是淡化了基督徒在社會上的文化角色.
第二件事更深的就是引誘他們行奸淫.
行奸淫這個字可以一字兩個意思.
第一個就是在舊約裡面說到他拜其他偶像.
所以這個其中可能那些商會他們去做了.
有機會就去拜其他的偶像.
不過更加具體的就是.
在推也推拉這個地方裡面.
有那些神廟有那些廟祭.
還有吃那些祭偶像之物.
可能帶來的不是純粹一個祭祀的宗教行為那麼簡單.
而是一個肉體行為.
所以在這個女先知說道.
這些東西算不得什麼.
將這個教導放進了教會裡面.
以致整個教會的風氣.
這一群基督徒是一群好人.
但是他們裡面的信念很薄弱.
甚至被西族裡面很多的風氣所影響.
以致他們這一群都是一群爛好人.
是一群好好對人好很有善意.
很有愛心的人.
但是對於信仰的立場他們就站立不穩.
我們在這裡稍微思考一下.
今天很多弟兄姊妹.
很多時候面對信仰都有這種掙扎.
究竟我們經常教弟妹變得有愛心.
我們要實踐愛.
我們一定不會錯的.
做對的事.
不過我們的愛心.
我們對人的那份接納.
不等於要接納所有人的行為.
特別是進入後現代.
很多的說法我們叫論述.
你接納一個人你就要接納那些人的行為.
即是說那個人做錯了事也好.
因為你要接納他.
所以你就無條件擁抱他.

$^{481}$他做什麼你都視為對的.
你要容納他.
不過接納人不等於默許某些人的行為.
你愛一個人不等於你愛那個人.
他做每一樣事情在你生命裡面都是對的.
你愛一個人.
更加是因為你知道.
你愛他.
你都想他走在對的道路裡面.
而不是他做什麼你都認同他.
那個不是愛他.
在聖經裡面說是縱容他.
那個叫做溺愛.
而不是來自上帝有真理的愛.
今天在我們的生命裡面.
我們做父母的.
我們對於子女有一些偏差行為.
因為我愛他.
所以我們會告訴他.
我們甚至乎有時候會嚴厲.
而不是純粹.
由得他.
由得他縱容他.
曾經有一些老弟兄老姐妹.
年紀大了.
子女都長大成人.
但是子女長大成人也好.
其實父母都有權柄.
不過很多時候.
有些比較年紀大的父母.
看到子女已經長大成人.
他長大成人也會走錯路.
不是小孩子.
成年人走錯路的時候.
父母覺得怎麼能說.
都不說了.
我愛他當然包他底.
成年人一直繼續走錯路.
而作為長輩的不出聲.
到最後.

$^{521}$越走越錯.
我經常勸勉.
不論是年輕的父母.
或者年長的父母.
你都是有父母的身份.
當子女無論什麼歲數.
他走錯路.
他偏離了生命的時候.
你都要提醒他.
可能你未必用嚴厲的方法也好.
但你不可以因為你愛他.
你就包容他所有錯誤的行為.
同樣在家裡是這樣.
在教會裡是這樣.
在生活不同的場景裡.
你愛一個人.
你不是容讓他.
接納他所有做錯的行為.
你接納一個人.
同時你都希望將他放回正軌.
做這件事是需要勇氣的.
我們需要有時.
從神裡面支取力量.
給我們有那份信心.
向身邊你所愛的人.
說出真心.
和真理的說話.
今日最大的問題是.
我們今日很多時候沒有這樣的勇氣.
我們覺得.
甚至乎被人愛心騎劫.
你講這些狠的話.
你就不愛我.
但我們知道.
心底裡面若是著緊一個人.
你不想他繼續這樣行差踏錯.
所以你就咬牙切齒.
很緊張.
要告訴他這些真相.
但今日這個世界.

$^{561}$好像有時顛覆了這種想法.
你接納一個人.
你就接納他所有行為.
包括錯的行為.
所以是一個彎曲薄謀的世代.
推也推拉這班的教會領袖.
我相信他們大致未必.
去到一個地步.
不知道這個.
椰洗別婦人的錯謀.
不過他們似乎.
他們只是習慣做好人.
沒有辦法.
鼓起勇氣.
去面對這個人.
甚至乎有些人可能.
不能分辨清楚.
以致為這個自稱為.
女先知的椰洗別婦人.
來作教.
經文裡面.
用椰洗別這個名號.
是有意思的.
耶穌用這個字眼描述這個女人.
首先.
究竟這個女人是甚麼人.
無從稽考.
我們沒有辦法找.
不過還要用椰洗別.
在舊約裡面的婦人.
很有意思.
椰洗別是記載在.
列王記下.
特別講到這個椰洗別.
是舊約裡面.
阿哈王的.
妻子.
王后.
一般來說在舊約裡面.
權傾朝野的.

$^{601}$女性很少.
但是椰洗別是其中一個.
權傾朝野的.
一個王后.
她的權力大到.
她的丈夫是王.
她都要點她的丈夫做事.
更加重要的是.
椰洗別將外邦的.
巴力的敬拜引入.
整個以色列民族當中.
陷整個以色列民族於不義.
這個椰洗別.
最後她要面對.
一個上帝的審判.
就是找耶護.
來到她成為了新一代的王.
接著上帝就直著.
耶護就清理了這個門戶.
這一段的歷史.
跟接下來第21節.
到第23節.
是互相相關的.
我們看完21至23節再解釋.
椰洗別.
我曾給她悔改的機會.
她卻不肯悔改.
她的淫行.
看啊!我要叫她病餓在床.
那些與她行淫的人.
若不悔改所行的.
我也要叫他們.
同受大患難.
我又要殺死她的黨類.
叫宗教會知道.
我是那察看人肺腑心腸的.
並要照你們的行為.
報應你們各人.
哇!耶穌很少這樣說話的.
殺氣騰騰.

$^{641}$為什麼用椰洗別呢?.
在聶王紀.
下第九章到第十章.
就記載著.
這個曾經這麼惡的皇后.
到最後.
上帝就直著.
耶護來清理門戶.
首先.
就是要殺死.
這個椰洗別.
在聖經裡面說到.
這個這麼惡的婦人.
朝代改變了.
上帝用椰護.
找她身邊的隨從.
即是椰洗別的隨從.
在窗那裡.
由高樓扔了她下地.
你會看到很血腥的一個場面.
甚至乎她死了之後.
她也沒有全屍.
有些野狗擔了她的屍體.
離開.
所以在這裡上帝就.
擊殺了椰洗別.
後來.
有幾十個兒女.
阿蝦和椰洗別.
上帝都直著耶護的手.
將她所有的兒女.
全部殺了.
這個典故.
其實就是用在.
二十一至二十三節.
耶穌在這裡說.
這個曾經是.
教內的人.
曾經上帝給他機會有悔改.
但是他不肯.

$^{681}$結果上帝就將.
刑罰放在他身上.
就是他病我在床.
所以他.
面對一個.
很極致的病患.
跟著跟他.
行任的人.
如果不悔改.
耶穌說都要叫他們同受大患難.
我又要.
殺死他的黨內.
黨內這個字其實原文用兒女.
用剛才說的阿蝦.
和椰洗別所生的很多的兒女.
全部都被殺.
用這個聖經典故.
就是說耶穌看到.
教會一些這樣的狀況.
他不可以視為不見.
人你容讓他.
但是耶穌是不會容讓他的.
所以.
耶穌基督在這裡.
特別說.
我是察看人.
肺腑心腸的.
並要照你們的行為報應你們各人.
可能你會說.
耶穌你不是恩典慈愛嗎.
你不是很憐憫嗎.
人有時候有罪所牽引.
你不是會赦免我們的罪嗎.
是沒錯.
不過耶穌的恩典不是.
姑息.
或者不是鼓勵人.
一路容讓罪惡.
在教會裡面來滋生.
耶穌的恩典.

$^{721}$也不是叫人在罪裡面.
繼續來維持這種現狀.
耶穌的恩典.
只是讓.
恆悔過的人.
可以重獲新的生命.
而不是縱容.
所有錯誤的謬誤.
或者這些教導.
在教會裡面來滋生.
最後.
24至28節.
最後一個段落.
至於你們推下推拉其餘的人.
就是一切不從那教訓.
不曉得他們素常所說.
撒旦森澳之理的人.
我告訴你們.
我不將別的擔子放在你們身上.
但你們已經有的.
總要持守.
直等到我來.
乃有達德性.
又遵守我命令到底的.
我要賜給他權柄.
制服列國.
他必用鐵匠.
核管他們.
將他們如同搖滬的瓦器.
打得粉碎.
像我從我父領受的權柄一樣.
我賜給他.
聖靈向宗教會所說的話.
凡有耳的應當聽.
首先.
耶穌基督.
這麼可怕的.
這麼嚴厲的面貌.
剛才對著耶西別.
一群這樣的人.

$^{761}$但接著耶穌.
他就比較放下那種怒氣.
所以他就說.
至於你們推下推拉那些人.
就是不順從那些教訓的.
絕大多數的弟兄姊妹.
這裡耶穌又說一些安慰的話.
就是你不跟那些說法.
說得故弄玄虛.
說得很深奧.
其實來自撒旦的道理的.
這一群弟兄姊妹.
我告訴你我不會將一些難擔的擔子.
放在你們身上.
你們有的繼續做.
之前做得對的.
你們的信心,行為,愛心,勤勞.
忍耐這些一切的美德.
繼續的來做.
所以在這裡.
耶穌的語氣比較緩和一點.
比較安慰性多一點.
還有一種鼓勵性會強一點.
他們做得對的事情.
繼續的來去持守.
直到耶穌基督再來的日子.
所以堅守耶穌這個命令.
或者耶穌給祝福教會的.
這些重要的價值觀念.
帶著一生的信心.
來到去持守.
不是一時的信.
而是一生的信.
來到去好好地活好.
這個信仰.
接著耶穌就給一個應許.
「那德性尤遵守我命令到底的.
我要賜給他權柄.
制服列國.
他必用鐵杖.

$^{801}$核管他們.
將他們如同搖滾的瓦器.
打得粉碎.
就像領袖的權柄一樣」.
這裡我們要小心去理解.
是不是說.
聽耶穌的話的人將來就變得很惡呢?.
其實在這段的解釋裡面.
有幾個字眼我們需要留意.
第一就是第27節裡面說到.
「他必用鐵杖.
核管他們」.
有一些和合本是括著給你的.
那個「核管」的字.
其實在原文裡面.
也在說「木樣」的字.
「木樣」的字和「核管」的字是同字的.
在舊約裡面.
所以有些地方就會譯「木樣」.
有些地方就會譯為「核管」.
和合本的翻譯.
就將「核管」這個字.
因為順民夏利是耶穌那麼堅定的.
所以他用一個比較硬朗的說法.
就是「核管」這個字.
不過我覺得如果用「木樣」這個字.
你會更加明白.
他在說將來耶穌基督要將權柄賜給他們.
賜給這班教會的弟兄姊妹.
這個權柄.
是好像耶穌將他的權力.
分享在眾教會的弟兄姊妹當中.
不過重要的就是.
他不是施行暴力.
而是好像耶穌作王一樣.
既要幫助那些國民在公義裡面.
同時間也要「木樣」這些身邊的國家.
甚至乎「木樣」是說舊時的「王統」字的時候.
也是用「木樣」這個字的.
簡單來說.

$^{841}$耶穌要將他的權力.
和這班堅守到底的弟兄姊妹.
來分享.
將耶穌的心腸.
耶穌對教會.
對他子民的帶領.
那種形式都能夠在這班推拉推拉的弟兄姊妹身上.
一齊實踐出來.
這是將來式.
不是現在馬上.
他聽話了就馬上管著人.
不是這個意思.
是將來千禧年.
或者將來的日子裡面.
耶穌將這份權力放在這班人身上.
不過重點就是.
我們今天.
如果用「木樣」這個字我們就要知道.
很多弟兄姊妹常以為.
「牧者」「木樣」是.
永遠都是笑著說.
陰森細氣.
你說什麼都好.
我為你祈禱.
求著幫助我們.
我們覺得「牧者」永遠都是這個形象.
我們覺得.
我們期望「獄者」永遠都是陰陰柔柔.
但是.
這段經文告訴我們.
「木樣」有時候是嚴厲的.
那個嚴厲不是那個「牧者」跟你.
有什麼仇口.
他要跟你對著幹.
而是.
那個硬朕是.
我們緊不緊神上帝對我們的祝福.
我們的那個「木樣」.
是要軟硬兼施的.
我們有柔的一面.

$^{881}$但同時間我們都要有.
立場清晰.
站穩陣腳的這一面.
如果我們在說今天世代.
我們要求什麼「木樣」呢.
我們經常都希望.
那些人對我很好.
要很多.
要拿很多的愛.
真的沒有說錯.
不過不可以只是純粹這一面.
這個世代是需要.
硬朕的那一面.
幫助一個人可以挺起胸膛.
幫助一個人可以頂天立地.
好好實踐你的基督教信仰.
弟兄姊妹.
這段經文說到這個應許.
說到「木樣」的理念.
今天我們怎樣看「木樣」呢.
我們就是純粹來到快樂.
舒服.
或者是這個群體.
我們需要幫助我們.
有力在這個世代裡面.
特別是世俗化的社會裡面.
我們能夠站得穩陣腳.
經文還有第二個應許的.
我又要把神聖賜給他.
聖靈向宗教會所說的話.
凡有意的就應當聽.
神聖有很多解法.
我比較接納的就是.
耶穌將他自己.
神聖就是比喻為耶穌他自己.
你要擁抱這個信仰.
不單只是說一些信條.
而是你的生命裡面活著.
就經歷基督.
和上帝有一個很深厚的關係.

$^{921}$如果你能夠.
撇棄世俗的思想.
你就能夠貼近上主的心懷.
你就明白耶穌的心腸.
你和基督耶穌就一起.
在這個世代裡面.
作美好的見證.
今天這段經文不複雜.
不過這段經文提醒我們.
我們今天的行事為人.
會不會成為一個爛好人呢?.
我們只是說愛心.
不說信仰裡面對我們的要求.
或者甚至我們對於今天.
衝擊整個信仰價值的很多價值觀念.
我們棄守了.
我們放棄了呢?.
鼓勵大家.
我們不要有好的外在行為.
卻是失去了心靈的信念.
失去了對上帝那份忠誠的信靠.
失去了這一部份的時候.
人家做到最後只是說你是一個好人.
但是經歷不到在你身上裡面.
基督教信仰是一個什麼信仰.
求主幫助我們.
今天這段經文雖然煞氣騰騰.
但也是耶穌基督對我們這個世代.
對我們每一位基督徒.
一個很嚴正.
但是很有力的勸勉和鼓勵.
我們一起學習 一起努力.
一起來努力.
\newpage



\section{約翰福音 10:19-42}
\label{sec:Ndl2CyA1nlc}
\textbf{2025年11月23日崇拜直播｜梁家麟牧師｜信耶穌到甚麼地步?｜約翰福音10\_19-42}
\newline
\newline
連結: \href{https://youtube.com/watch?v=Ndl2CyA1nlc}{\texttt{https://youtube.com/watch?v=Ndl2CyA1nlc}} ~~~~ 語音日期: 2025-11-29
\newline
\newline
\hyperref[sec:H6vb7uu46Wk]{\small{< < < PREV SERMON < < <}}
~
\hyperref[sec:index_chronic]{\small{[返順時目]}}
~
\hyperref[sec:XMVPsTUvcaY]{\small{> > > NEXT SERMON > > >}}
\newline
\newline
約翰福音 10:19-42
\newline
\begin{longtable}{cl}
\hline
\hline
章節 & 經文 (和合本修訂版)\\
\hline
10:19 & \begin{tabularx}{0.7\textwidth}{X} 猶太人為這些話又起了分裂。 \end{tabularx} \\ \\ \relax
10:20 & \begin{tabularx}{0.7\textwidth}{X} 其中有好些人說：「他是被鬼附了，而且瘋了，為甚麼聽他的呢？」 \end{tabularx} \\ \\ \relax
10:21 & \begin{tabularx}{0.7\textwidth}{X} 另有的說：「這不是被鬼附的人所說的話。鬼豈能開盲人的眼睛呢？」 \end{tabularx} \\ \\ \relax
10:22 & \begin{tabularx}{0.7\textwidth}{X} 那時正是冬天，在耶路撒冷有獻殿節。 \end{tabularx} \\ \\ \relax
10:23 & \begin{tabularx}{0.7\textwidth}{X} 耶穌在聖殿裡的所羅門廊下行走。 \end{tabularx} \\ \\ \relax
10:24 & \begin{tabularx}{0.7\textwidth}{X} 猶太人圍著他，對他說：「你讓我們猶豫不定到幾時呢？你若是基督，就明白地告訴我們。」 \end{tabularx} \\ \\ \relax
10:25 & \begin{tabularx}{0.7\textwidth}{X} 耶穌回答他們：「我已經告訴你們，你們卻不信。我奉我父的名所行的事可以為我作見證。 \end{tabularx} \\ \\ \relax
10:26 & \begin{tabularx}{0.7\textwidth}{X} 但是你們不信，因為你們不是我的羊。 \end{tabularx} \\ \\ \relax
10:27 & \begin{tabularx}{0.7\textwidth}{X} 我的羊聽我的聲音，我認識牠們，牠們也跟從我。 \end{tabularx} \\ \\ \relax
10:28 & \begin{tabularx}{0.7\textwidth}{X} 並且，我賜給他們永生；他們永不滅亡，誰也不能從我手裡把他們奪去。 \end{tabularx} \\ \\ \relax
10:29 & \begin{tabularx}{0.7\textwidth}{X} 我父所賜給我的比萬有都大，誰也不能從我父手裡把他們奪去。 \end{tabularx} \\ \\ \relax
10:30 & \begin{tabularx}{0.7\textwidth}{X} 我與父原為一。」 \end{tabularx} \\ \\ \relax
10:31 & \begin{tabularx}{0.7\textwidth}{X} 猶太人又拿起石頭來要打他。 \end{tabularx} \\ \\ \relax
10:32 & \begin{tabularx}{0.7\textwidth}{X} 耶穌回應他們：「我做了許多從父那裡來的善事給你們看，你們是為哪一件拿石頭打我呢？」 \end{tabularx} \\ \\ \relax
10:33 & \begin{tabularx}{0.7\textwidth}{X} 猶太人回答他：「我們不是為了善事拿石頭打你，而是為了你說褻瀆的話；因為你是個人，卻把自己當作神。」 \end{tabularx} \\ \\ \relax
10:34 & \begin{tabularx}{0.7\textwidth}{X} 耶穌回答他們：「你們的律法書上不是寫著『我曾說你們是諸神』嗎？ \end{tabularx} \\ \\ \relax
10:35 & \begin{tabularx}{0.7\textwidth}{X} 經上的話是不能廢的；如果那些領受神的道的人，神尚且稱他們為諸神， \end{tabularx} \\ \\ \relax
10:36 & \begin{tabularx}{0.7\textwidth}{X} 那麼父所分別為聖又差到世上來的那位說『我是神的兒子』，你們還對他說『你說褻瀆的話』嗎？ \end{tabularx} \\ \\ \relax
10:37 & \begin{tabularx}{0.7\textwidth}{X} 我若不做我父的工作，你們就不必信我； \end{tabularx} \\ \\ \relax
10:38 & \begin{tabularx}{0.7\textwidth}{X} 我若做了，你們即使不信我，也當信這些工作，好讓你們知道並且明白父在我裡面，我也在父裡面。」 \end{tabularx} \\ \\ \relax
10:39 & \begin{tabularx}{0.7\textwidth}{X} 於是，他們又要捉拿他，他卻從他們手中逃脫了。 \end{tabularx} \\ \\ \relax
10:40 & \begin{tabularx}{0.7\textwidth}{X} 耶穌又往約旦河的東邊去，到了約翰起初施洗的地方，就住在那裡。 \end{tabularx} \\ \\ \relax
10:41 & \begin{tabularx}{0.7\textwidth}{X} 有許多人來到他那裡，說：「約翰沒有行過一件神蹟，但約翰所說有關這人的一切話都是真的。」 \end{tabularx} \\ \\ \relax
10:42 & \begin{tabularx}{0.7\textwidth}{X} 在那裡，許多人信了耶穌。 \end{tabularx} \\ \\
[1ex]
\hline
\hline
\end{longtable}
$^{1}$以下我們要聽神的話.
《約翰福音》第十章二十二至三十節.
請聽弟兄讀出.
在耶路撒冷有修奠節.
是冬天的時候.
耶穌在殿裡所有門的廊下行走.
猶他人圍著祂說.
「你以為我們猶疑不定到何時呢?.
你若是基督就明明地告訴我們」.
耶穌回答說.
「我已經告訴你們,你們不信.
我奉我父之名所行的事.
可以為我作見證.
只是你們不信.
因為你們不是我的羊.
我的羊聽我的聲音.
我也認識他們.
他們也跟著我.
我又賜給他們永生.
他們永不滅亡.
誰也不能從我手裡把他們奪去.
我父把羊賜給我.
他被萬有都帶.
誰也不能從我手裡把他們奪去.
我們與父願為一」.
聖經讀筆.
弟兄姊妹早安.
今天的經文是由19節到42節.
即是第十章大部分都讀完了.
不過讀經不想讀太長.
所以我們只讀一段.
大家是基督徒.
並且我們很幸運.
我們遇到的是已經復活的耶穌.
所以我們接受耶穌基督的完全神聖.
是沒有太大的困難.
不過在耶穌在世的時代.
祂是一個充滿爭議性的人物.
唯有祂身邊的人.
對祂有非常分歧.

$^{41}$甚至是兩極化的看法.
耶穌能夠行神蹟.
能夠讓核子看見.
這證明祂有超自然的能力.
祂的教導也與眾不同.
常常有奇特的觀點.
這證明祂有非凡的智慧.
如果我們將耶穌.
對耶穌的理解停留在這樣的地步.
一個有能力又有智慧的人.
大概就是上帝差來的先知.
舊約裡面有很多先知.
以利亞就是先知.
雖然近幾百年已經不再有先知存在.
不過現在多一位先知都沒有問題.
所以對耶穌有正面看法的人.
大概都會認定耶穌是一位先知.
但進一步.
耶穌是不是舊約所要來的彌賽亞呢?.
如果祂是.
祂是不是要帶領以色列人.
從羅馬帝國的手中奪回主權.
建立一個新的國度.
是不是這樣呢?.
從祂的出身和祂的活動.
和祂所宣稱的使命.
顯然不像是一個偉大的政治家.
祂的活動範圍主要在加利利.
這個不是猶太人的權力中心.
最多是類似今天的荃灣.
離開權力中心差太遠了.
祂只是偶然去耶路撒冷.
和宗教領袖對話.
但沒有參與任何公開的政治活動.
也從來沒有挑戰過羅馬的政權.
這是什麼彌賽亞斯人呢?.
最令人困擾的是.
耶穌常常說出一些令人震驚,費解.
甚至是激怒人的說話.
祂自稱是天父上帝的兒子.

$^{81}$和上帝有不尋常的關係.
甚至是隱隱宣告祂具有神聖的地位.
這個對堅持一神論的猶太人難以接受.
因為人自稱上帝.
這是最狂妄,最褻瀆神明的說話.
耶穌可以是先知,可以不是.
耶穌可以是彌賽亞,可以不是.
無所謂的.
但耶穌自稱是上帝.
這是知事體大.
因為這搖動了猶太人信仰的根本.
他們一定要撕裂自己的衣裳.
表達份外.
問題是.
一個褻瀆神明的人.
怎可能是屬於上帝的呢?.
或者是被上帝猜測出來的人呢?.
如果我們認為耶穌褻瀆神明.
那麼耶穌就連先知都不會是.
祂只是屬於撒旦.
並且借助鬼王的能力去行神蹟.
這是反對耶穌的宗教領袖一致的看法.
當耶穌行了叫核子復明的神蹟之後.
猶太人中間就產生了極大的爭議.
有人斷言他瘋了.
並且是被鬼附的人.
所以不要聽他的說話.
但有人質疑.
被鬼附的人怎能夠行這樣的神蹟呢?.
總之眾說紛紜,意見很分歧.
你可以看到.
耶穌的言行是將人逼近牆角的.
使人沒有迴旋的餘地.
人不可以保持中立.
不可以有待考證.
你要立即表態.
如果你接受耶穌的話.
你就要奉祂為上帝.
並且必須趴在祂面前.
如果你拒絕耶穌的話.

$^{121}$你就一定要當祂是瘋子或騙徒.
褻瀆上帝,無法無天.
唯一的處理方法.
在舊約時代的理解.
就是將祂致死.
耶穌或是瘋子或是上帝.
你必須要二擇其一.
注意,在當時人的眼中.
瘋子和騙徒幾乎是同義詞.
他們相信疾病是有超自然的原因.
瘋狂是被鬼附.
不是心理問題.
是宗教問題.
騙徒是說方者.
亦與撒旦有關連.
最激烈反對耶穌的猶太教領袖.
已經做出了決定選擇.
耶穌是必須被致死的騙子.
但多數群眾仍然猶豫不定.
搖擺不定.
他們在苦惱中.
有人忍不住直接問耶穌.
二十四節.
你叫我們猶豫不定到何時呢?.
你若是基督.
就明明的告訴我們.
耶穌當時在耶路撒冷的聖殿裡.
參與了修奠節的活動.
修奠節又叫光明節.
是猶太人一個新興的宗教節日.
舊約沒有記載的.
是紀念主前二世紀.
瑪加比家族帶領以色列人.
從希臘塞留古帝國的國王安條克四世.
手中奪回耶路撒冷.
並且將第二聖殿重新獻給耶和華.
這個節日在冬天舉行.
通常為期八一.
這是一個充滿民族情懷.
致切要復興以色列的節日.

$^{161}$瑪加比王朝收回以色列的主權.
這是一個這樣的節日.
參與活動的群眾.
在聖殿的東側.
就是所羅門廊下.
遇到耶穌.
他們忍不住就撲過去問.
耶穌,你是否我們要等待的彌賽亞?.
如果是的話,請你不要閃閃縮縮.
明白清楚告訴我們吧.
你會看到提問者對耶穌.
其實有相當正面和善意的看法.
他們不是質疑耶穌.
只是想從耶穌口中得到是或不是的答案.
他們不想揣測下去.
猶豫不定,實在令他們太困苦.
或者很困擾.
抱歉,耶穌沒有像他們所願的回答.
是或不是.
我們注意到耶穌從來沒有在公開場合.
直接宣稱自己是基督.
祂只是在私下,私底下.
間接承認這個身份.
例如在厄國井旁和婦人的對話.
四章二十六節.
因為順著當時的以色列人對基督的政治性了解.
宣稱自己是基督.
等於揭竿起義.
要打倒羅馬帝國.
對耶穌來說,祂當然是基督.
不過要補充指出.
祂不是按當時的人所理解的基督.
祂是基督.
但不是你想像的基督.
是不是?.
耶穌於是這樣回應.
在一十五節.
我已經告訴你們,你們不信.
我奉我父子名所行的事,可以為我作見證.
耶穌過去一言一行.

$^{201}$包括祂的教訓,祂的神蹟.
通通都指向這個事實,祂是基督.
不過這是祂自己定義的基督.
不是革命家,不是當時的領袖.
而是一個甘願為人服侍犧牲的救世主.
所以不是耶穌故意含糊.
讓群眾感到困惑.
究竟耶穌是不是基督呢?.
而是耶穌的基督身份.
和他們的基督理解有重大差異.
或者說現實的,真實的基督.
和他們心中所期望的基督有重大差異.
以至他們無法將期望和現實對接.
他們真正要問耶穌的不是.
耶穌,你是不是基督.
而是耶穌,請你告訴我,什麼是基督.
基督是誰,是怎麼樣的.
讓耶穌自行告訴你,祂是誰.
祂承擔的是什麼樣的使命.
26,27節,耶穌接著說.
只是你們不信,因為你們不是我的羊.
我的羊聽我的聲音,我要認識他們.
他們也跟著我.
耶穌的教訓和神跡.
已經表明祂的神聖身份和使命.
不過這群人既不理解又不接受.
他們的不理解其實是不願意理解.
他們的不接受其實是不願意接受.
所以耶穌指出,你們是不信.
我們對新的事物的理解和接受.
很多時候取決於我們固有的相關認知.
如果我們已經有某一個很牢固的.
所謂前理解,pre-understanding.
先入為主的想法.
我們就沒辦法接受這個新的.
跟我們的意知相衝突的新智.
這是所謂先入為主.
在宗教信仰領域特別普遍.
你們也有這些經驗.
你們跟傳福音就有.

$^{241}$一個人如果拜觀音拜了很多年.
即使你在他面前講很多見證.
證明基督教有多好.
耶穌有多真實.
他也不一定願意改變他的信仰.
我們跟傳福音常常碰到這樣的情況.
不是我不信耶穌.
而是我沒需要信耶穌.
他們會說你信你的耶穌,我信我的觀音.
是不是這樣?.
各信各,是不是?.
人很少會無緣無故改變自己的宗教.
除非他遭遇危機.
發覺觀音其實不行的.
是不是?.
如果他很地地,覺得很靈驗.
一切順風順水.
你要他信,他不會理你.
他根本不會理你.
是不是?.
我有了我已經先入為主的東西.
我為什麼要換新的呢?.
是不是?.
所以不理解是因為我不想理解.
不接受是因為我不想接受.
信仰不是靜態的.
不僅僅是一個信念.
也是動態的.
是激發人的感情.
人的行動意志.
奧古斯丁特別指出.
信服和理解的關係.
信服是一種聆聽的方式.
出於信心的信服.
就使理解成為可能.
人在能夠去明白屬上帝的事之前.
必須要渴望去理解.
愛要在論證之前.
奧古斯丁特別從《二賽亞書》引了一段.
他最喜歡的經文.

$^{281}$這是拉丁文武家大日本的《二賽亞書》七章九節.
翻譯是:你們若不信,就斷不能理解.
人永遠是首先相信.
然後才努力尋求理解.
我做基督徒是因為我先信耶穌.
我不是先信聖經.
我都未見過五聖經.
信什麼信?.
你說他值得信就值得信.
我是先和耶穌基督相遇.
然後他拿了我的心.
他霸佔了我.
然後我在耶穌面前降服.
然後有人告訴我.
耶穌說給我作見證的就是這經.
所以我就信聖經.
但在我信聖經的時候.
我未曾看過整本聖經.
更加不要說我完全明白.
你現在也不會完全明白.
整本聖經不是全部明白.
不過我先信了.
我先接納聖經的權威.
我先信然後我努力去明白.
即使有些地方我還未明白.
我仍然咬定我先相信.
然後我接著尋求理解.
這是奧古斯丁的想法.
這也是十二世紀偉大的神學家.
安瑟倫的想法.
我信以至我能夠理解.
I believe in order to understand.
我先信然後以至我能夠明白.
我們回到耶穌的那裡.
耶穌接著再將他的話題扭轉.
講到信他的人和不信他的人的重要差異.
唯有是屬於他的人才會信他.
這是天父上帝老早選定.
那些特別是屬於耶穌的人.
28到29節.

$^{321}$我又賜給他們永生.
他們永不滅亡.
誰也不能從我手裡把他們奪去.
我父把羊賜給我.
他比萬有都大.
誰也不能從我父手裡把他們奪去.
這段話的意思很簡單.
屬於耶穌的人就屬於耶穌.
打風也打不掉.
不屬於耶穌的人就像貼索門神一樣.
怎麼也貼不上.
這是非人力所能控制.
沒辦法勉強的事.
耶穌在猶太人中間服侍了一段時間.
他們如果願意理解他.
早就理解了.
假如拖拖延延.
拖了這麼久.
一段這麼長的時間.
仍然有很多的困惑.
你就再多教導.
再多的神蹟.
再多的證據.
都不能幫助他們.
由懷疑變成相信.
所以多說無謂.
猶疑不定就猶疑不定.
耶穌在回答這個詢問的時候.
特別強調他和上帝的關係.
他如常的稱上帝為天父.
強調他所說所做的.
都是奉天父的名做的.
他是按照天父的吩咐去做.
所做的也是天父旨意的成就.
他又指出.
所有相信他的人.
都是天父給他的羊.
天父也確保.
所有他給耶穌的人.
最終都屬於耶穌.

$^{361}$得著永生.
沒有任何屬靈的力量.
能夠將他們從耶穌手中奪走.
說到這裡.
在場的猶太人.
各持或者信.
或者半信半疑的態度.
不過.
耶穌癲得得.
又不知為什麼.
加多一句.
三十節.
我與父緣為一.
唉.
這真是大件事了.
大件事了.
這句話令在場所有人大衛震驚.
與父為一.
就是與父平等.
與父緣為一.
就是耶穌從一開始.
在本質上就和上帝同等.
這幾乎等於耶穌自稱是上帝.
是不是.
好像剛才說.
聽眾或者信耶穌是上帝.
或者就要立即撕裂衣裳.
他們聽了最褻瀆性的話.
竟然有人妄稱自己是上帝.
猶太人當下的反應是後者.
他立即在地上撿起石頭.
要扔耶穌.
我要清楚說這裡.
耶穌說我與父緣為一.
這個一.
這些東西.
你飛了這一段也無所謂.
希臘文是一個中性.
就是newton的hand.
而不是陽性.

$^{401}$所以緣為一.
真正意思不是identical.
是同一.
而是同等.
equal.
同質.
神學名詞叫Oozeos.
如果我們說.
我與父緣為一是指identical.
是相同.
父就是子.
子就是父.
兩個其實是一個.
這是神學上是錯的.
這是形態論.
父子姓名是三態.
水的三態.
有時是冰.
有時是水.
有時是蒸氣.
父和子是不會同時存在的.
因為是同一個的兩個名字.
冰的時候就不能同時是蒸氣.
上帝有幾個變身.
有時變父.
有時變子.
有時變靈.
這個叫identical.
同一個.
這句經文的意思是.
我與父是同等.
同等地位.
同性質.
同神性.
同權.
同尊.
幾句的.
不過解釋一下就行了.
前兩次我都和你們說過.
耶穌多番強調他與天父密切的關係.

$^{441}$他所說的完全是天父吩咐他所說的.
他所做的完全是完成天父差派的任務.
這個說法.
不是意味著耶穌隸屬上帝.
沒有自己的想法和做法.
退一步說.
如果耶穌僅僅是人.
人本來是屬於上帝的.
人所做的.
所說的.
都是遵照上帝的吩咐去做.
這是理所當然的事.
沒有人會覺得羞恥的.
我跟著父上帝的心意去做.
不是很值得羞恥.
不過耶穌說這番說話.
真正的意思是.
他和父上帝的心意是完全一致的.
你不用問耶穌的想法.
是否符合上帝的想法.
耶穌的做法會不會和上帝對著幹.
你省一點.
你在人間看到耶穌做事.
所做的完全就是上帝旨意的兌現.
從來沒有人見過上帝.
耶穌的一言一行.
就是上帝百分百的彰顯.
耶穌說的話不是說.
第一節矮過上帝.
而是說完全一致.
就好像我和兒子說.
你不用問媽媽.
我說的是媽媽的想法.
所以不用離間我們.
不用玩兩面手法.
媽媽的想法就是我的想法.
媽媽的決定就是我的決定.
我的決定就是媽媽的決定.
他說這件事是這個意思.
他們的意見完全一致.

$^{481}$當猶太人要用石頭打耶穌的時候.
耶穌沒有立刻躲起來質問他們.
他們究竟做了什麼錯事.
以致他們要打他們.
猶太人說這個無關乎他們所做的事.
純粹是因為他們所說的話.
33節他說.
我不是為善事拿石頭打你.
是為你說遷往的話.
又為你是個人.
反將自己當作上帝.
從猶太人他們負面的反應.
不過你看這句話.
你就看到其實他們沒有錯誤理解耶穌的話.
耶穌真的說我與父原為一.
自稱為上帝.
他們生氣就生氣這裡.
你明明是人.
你自稱為上帝.
所以我就要用石頭打你.
對於猶太人的指控.
耶穌就引述舊約斯篇82篇6節.
那裡說.
我曾說你們是神.
都是至高者的兒子.
指出天父上帝曾經稱他以外的某些對象為神.
雖然那些其實是人.
不過因為他們承受了上帝的道.
成為了上帝的兒女.
就享有神聖的地位.
那我這個與眾不同.
被上帝親自差離的人間耶穌.
自稱為上帝.
又有什麼大驚小怪呢?.
詩篇82篇是一首很負爭議性.
有很多分歧性解釋的詩.
今天我沒有時間詳細解說.
簡單的理解是.
耶穌是將這首詩的這句話.
視為舊約對他自己的一個預言.

$^{521}$他是這一首詩的終極的應驗.
在舊約時代.
上帝曾經以不同的形式將他的道傳給百姓.
這個道就轉化了百姓與上帝的關係.
從而轉化了百姓的身份.
他被稱為上帝的子民.
被分別為聖.
82篇6節說的你們是神.
這個神字是英文世草的God.
就是說百姓獲得了神聖的身份和地位.
當然舊約所啟示的上帝的道是未完整的.
以色列民的神聖身份也是殘缺不全的.
所以82篇6節的應許是未兌現的要求.
耶穌自己是上帝的道.
耶穌是由上帝所分別為聖.
由差道人間來的.
所以耶穌的說法是.
我才是這一句話的真正的對言.
所以我當然就是神了.
我當然就是神了.
耶穌自稱擁有神聖的地位.
有什麼這麼潮望性呢.
耶穌詩篇82篇是支持他的上帝地位的聖經根據.
他既然是上帝的道.
能夠轉化人的生命.
使人變得神聖.
我可以轉化你使你變得神聖.
我自己是什麼呢.
當然是神聖了.
我不神聖怎麼轉化你變成神聖呢.
如果耶穌不是上帝的兒子.
怎麼帶血屬於他的人成為上帝的兒女呢.
好像希伯來書2章10節所說.
原來那為萬物所屬.
萬物所本的要領許多的兒子進榮耀裡.
使人成為上帝的兒女的肯定是上帝的兒子.
並且跟後來那些跟著他成為上帝的兒女是有懸殊差異的.
我們的上帝兒女身份是相對的.
後加的.
耶穌的上帝兒子身份是絕對的永恆.

$^{561}$耶穌最後總結說37到38節.
我若不幸我負的事,你們就不必信我.
我若幸了,你們縱然不信,也當信這些事.
又叫你們又知道又明白.
負在我裡面,我又在他裡面.
耶穌的身份和地位是由他的言行得以證明.
他所做的完全彰顯了天父的旨意.
他就是和天父上帝擁有不可分的關係.
地位同等的那一位.
負在我裡面,我也在他裡面.
這是耶穌常常說的話.
就是我與父願為一的最佳說明.
《約翰福音》裡面還有很多耶穌講論他和天父緊密不可分的關係的經文.
不過我們一後讀到的時候再說.
在場的猶太人當然無法理解和接受耶穌的話.
他們又要動手去對付耶穌.
不過39節講耶穌就躲起來了.
聖經最後說.
耶穌接著就離開耶路撒冷去到約旦河外.
從前施洗約翰事奉的地方.
我們記得施洗約翰一生最重要的工作不是替人施洗.
施洗約翰一生最重要的工作是什麼?.
作為耶穌的開路先鋒.
預備主的道,修籍主的路.
他是為耶穌作見證而生.
為耶穌作見證而活的那個.
那裡的人既然相信施洗約翰.
就相信施洗約翰所作的見證.
從而就信了耶穌.
無論如何沒有帶著各種宗教偏見的人.
理解和接受耶穌是容易很多的.
諷刺的是耶路撒冷那些學父五車的宗教領袖.
無法理解耶穌的話.
約旦河外那些無知小民就理解了.
好,經文解完了.
這麼辛苦的經文都解完了.
這段經文有很多值得深思的地方.
例如耶穌和天父的關係.
耶穌的神聖.
例如耶穌專注的事奉態度.

$^{601}$祂不強求所有人都理解和接受祂.
祂專注於上帝交給祂的羊.
這幾個主題我們日後都會有機會講的.
今天我們集中思考其中一點.
就是我們的題目.
我們要信耶穌,信到什麼地步.
耶穌是一個招來爭議的人.
很多人對祂有分歧,兩極化的看法.
即使祂有正面看法的猶太人.
也因著耶穌的言行而產生很多疑惑,很多焦慮.
耶穌是不是先知?.
耶穌是不是彌賽亞?.
耶穌究竟是誰?.
人可以喜歡耶穌.
充滿人情味的倫理教訓.
但耶穌就不太喜歡你這樣喜歡祂.
我挺囂張的.
耶穌不喜歡你這樣喜歡祂.
所以耶穌其實常常生氣.
常常故意提出一些可人聽聞的言論.
人可以喜歡耶穌的神職其事.
但耶穌接著就會講一些離經叛道的宣告.
這個叫做趕客.
我們可以喜歡耶穌的一方面.
但耶穌立即告訴我們.
祂有更多使我們受不了的地方.
好像我一直暗戀一個清麗脫俗的女生.
不過她竟然突然以蓬頭垢面.
滿口粗口的面貌向我顯現.
破壞了她在我心中的完美形象.
令人非常難受.
難受的事情怎麼去理解?.
怎麼疏解她?.
第一,我們要接受一個事實.
你喜歡一個人.
你不可以單單喜歡她某一方面.
不喜歡她其他地方.
你要喜歡就喜歡全部.
如果你只是喜歡一個女生的某一部分.
就千萬不要跟她結婚.

$^{641}$維持暗戀單戀就夠了.
對不對?.
就好像我以前都跟你講過.
你未結婚,未拍拖之前.
我們可以問一個問題.
你心目中的理想對象是怎樣的呢?.
你問我這樣的問題.
你是想我死掉嗎?.
我理想的女生是誰?.
就是我老婆.
我還有第二個選擇嗎?.
你未拍拖之前.
你可以弄張紙.
弄一大堆的check list.
她需要很大的眼睛.
她的眉毛需要兩吋長.
她需要怎樣怎樣.
瓜子面怎樣怎樣.
拿著check list就去找.
不過當你認定了一個之後.
你再拿張check list去對.
那你就要找你老婆了.
對不對?.
你想怎樣?.
你告訴我,對不對?.
如果你認定她就是你的終身伴侶.
你還拿張check list去對她.
對不對?.
你只能夠無條件擁抱她,接受她.
甚至改變你原來的想法.
改變你原來的想法,對不對?.
我不知道有沒有跟你說過.
我有一個晚輩的親戚.
他不久前跟他拍拖七,八年的女友分手.
我們問他原因.
他說他覺得她不漂亮.
我們嘴巴都凹了.
你八年後覺得她不漂亮?.
你不是吧?.
你明白嗎?.

$^{681}$你覺得很荒謬.
你現在才發現.
你以前是怎樣?.
這個不再是一個要思考討論的問題.
除非你說我是騎牛搵馬.
一直都是騎牛搵馬.
找到一個更好的就找一個舊的.
不過你省一點吧.
對不對?.
同樣的.
如果你是相信耶穌.
你就不可以只是相信他是一個倫理教師.
他是傳遞上帝信息的先知.
或者你只是喜歡他的神蹟.
不喜歡他的教導.
我們不可以將耶穌斬件.
要麼你完全信他.
死心塌地的信他.
要麼信都等於不信.
半信半疑.
和不信其實不是有分別的.
是沒什麼分別的.
好像耶穌在比喻裡面說.
那些將種子落在淺土的.
是沒辦法發芽.
紮根成長的.
不冷不熱的信仰是沒用的.
半途而廢的信仰是沒用的.
整本聖經不停告訴我們.
聽這個事實.
信耶穌如果不是完全相信.
就等於不信.
等於沒有真正相信.
是不是很霸道?.
耶穌常常作出趕客的行動.
用很激烈的要求.
例如說人子沒有拯收的地方.
變賣所有分給窮人.
愛他就要恨所有身邊的人.
撇下所有跟隨他.

$^{721}$來逼使人放棄原先跟隨他的衝動.
他還專門爆出很多極難受的話.
他宣稱自己是上帝.
宣稱他的肉可以吃.
他的血可以喝.
一聽這些你都覺得很噁心.
說這些話.
這些真是令人毛骨悚然的言論.
同樣是趕客行為.
你頂得住就繼續跟隨他.
否則離去自便.
耶穌絕對不勉強挽留你.
除了信要信全部.
信要信徹底之外.
第二點更加重要的是.
我們在信耶穌的過程裡.
必須要逐漸.
必須要慢慢.
但最後要徹底.
放下自己既有的想法.
承見偏愛,潛理解.
騰出最大的空間.
來迎接上上出我們意料之外.
驚嚇我們.
擾亂我們的知識系統.
我們的價值觀的耶穌.
正如耶穌上上說.
沒有失就沒有得.
沒有死就沒有生.
沒有捨棄就不會得著.
你放下多少.
你就得著多少.
你不要想著.
我維持我既有的所有.
信主前的所有想法做法.
而是加上.
加一點信仰在上面.
雞套不同籠的.
合不來的.
信耶穌不可以在我們.

$^{761}$原封不動的人生之上.
稍作添加.
卻是要翻轉我們既有的生命和生活.
好像我常常說的.
很無聊的抽象說話.
信仰不可以是補充性.
不是supplementary.
一定是subversive 顛覆性.
信耶穌就準備讓祂動你.
讓祂翻轉你.
你不讓祂徹底翻轉.
你的信仰.
是表面虛假的.
所有信仰困擾和衝突.
都源自我們固有的想法和期望.
和上帝引導的.
祂引導的不搭配.
我的想法和祂的引導不搭配.
我們覺得上帝辜負了我們的期望.
破壞了我們的計劃.
擾亂了我們原有的生活秩序.
上帝我明明計劃要做出一番事業去榮耀你.
為什麼你要讓我患重病.
你是不是搞錯了什麼.
如果我們願意放下固有的想法和期望.
你的信仰困擾一定減少很多.
在耶穌面前.
我們要努力放下所有既有的執著.
讓祂自由的.
隨己意的.
不按照我們自己訂單來發貨.
你告訴我你是誰.
你對我有什麼要求.
不要照我的意思.
只要照你的意思.
讓耶穌的旨意成就在我們身上.
弟兄姊妹.
你信耶穌.
信到什麼地步.
你是不是真的完全相信祂.

$^{801}$你是不是常常覺得.
你的想法和祂的想法有很多的衝突.
你放下多少你的想法.
\newpage



\section{啟示錄 3:1-6}
\label{sec:XMVPsTUvcaY}
\textbf{2025年11月30日崇拜直播｜陳冠成牧師｜耶穌給教會的七封信:務要儆醒｜啟示錄3\_1-6}
\newline
\newline
連結: \href{https://youtube.com/watch?v=XMVPsTUvcaY}{\texttt{https://youtube.com/watch?v=XMVPsTUvcaY}} ~~~~ 語音日期: 2025-12-01
\newline
\newline
\hyperref[sec:Ndl2CyA1nlc]{\small{< < < PREV SERMON < < <}}
~
\hyperref[sec:index_chronic]{\small{[返順時目]}}
~
\hyperref[sec:XJJ1wNBmbSY]{\small{> > > NEXT SERMON > > >}}
\newline
\newline
啟示錄 3:1-6
\newline
\begin{longtable}{cl}
\hline
\hline
章節 & 經文 (和合本修訂版)\\
\hline
3:1 & \begin{tabularx}{0.7\textwidth}{X} 「你要寫信給撒狄教會的使者，說：『那有神的七靈和七顆星的這樣說：我知道你的行為，就是名義上你是活的，實際上你是死的。 \end{tabularx} \\ \\ \relax
3:2 & \begin{tabularx}{0.7\textwidth}{X} 你要警醒，堅固那些剩下、快要死的，因為我發現你的行為，在我神面前沒有一樣是完全的。 \end{tabularx} \\ \\ \relax
3:3 & \begin{tabularx}{0.7\textwidth}{X} 所以，要記得你所領受和聽見的；要遵守，並要悔改。你若不警醒，我必如賊一樣來到；我幾時來到你那裡，你絕不會知道。 \end{tabularx} \\ \\ \relax
3:4 & \begin{tabularx}{0.7\textwidth}{X} 然而，在撒狄你還有幾位是未曾污穢自己衣服的，他們會穿白衣與我同行，因為他們是配穿的。 \end{tabularx} \\ \\ \relax
3:5 & \begin{tabularx}{0.7\textwidth}{X} 得勝的必這樣穿白衣，我也不從生命冊上塗去他的名；我要在我父面前，和我父的眾使者面前，宣認他的名。 \end{tabularx} \\ \\ \relax
3:6 & \begin{tabularx}{0.7\textwidth}{X} 凡有耳朵的都應當聽聖靈向眾教會所說的話。』」 \end{tabularx} \\ \\
[1ex]
\hline
\hline
\end{longtable}
$^{1}$我們今天繼續在啟示錄中聆聽聖靈對教會的說話.
今天的經分是啟示錄第三章一至六節.
在新約聖經第三百五十一頁.
找到的請聽我讀出.
「你要寫信給撒迪教會的使者說.
那有神的七靈和七星的說.
我知道你的行為.
我明你是活的.
其實是死的.
你要警醒.
堅固那餘下將要衰微的.
因我見你的行為.
在我神面前.
沒有一樣是完全的.
所以要回想你是怎樣領受.
怎樣聽見的.
又要遵守.
並要悔改.
若不警醒.
我必臨到你那裡.
如同賊一樣.
我幾時臨到.
你也決不能知道.
然而在撒迪.
你還有幾名善未曾為自己依附的.
他們要穿白衣與我同行.
因為他們是配得過的.
凡得勝的.
必這樣穿白衣.
我也必不從生命冊上.
圖污他的名.
且要在我父面前.
和我父眾使者面前.
認他的名.
聖靈向眾教會所說的話.
凡有意的就當一堂聽」.
弟兄姊妹平安.
我們先有個禱告.
父上讓我們知道.
即使洪水泛濫.

$^{41}$你仍然坐著為王.
讓我們認定.
你是那位真實與我們同在的上帝.
所以無論我們各人.
現在處於甚麼的境況.
我們帶著甚麼的心情.
讓我們都能夠定精對準你這個安穩的牢.
以致我們心得平安.
也讓我們藉著聽道.
知道上帝你要對我們做甚麼.
知道上帝你要對我們眾人說話.
盼望主求聖靈你與我們同在.
讓我們真的能夠明白.
上帝啊.
你想我們怎樣去回應你.
你想我們在這個世代裡.
怎樣依靠你.
來過每一天每一日.
為此我們同心獻上禱告.
求上帝你垂聽.
奉基督耶穌你.
得勝的名字來祈求.
阿們.
不過今天我們進入了.
啟示錄七封信的第五封.
不過在之前我想借幾條問題考考大家.
對香港的街道其實你熟不熟悉呢?.
第一條.
我想問大家知道有條叫西貢街.
西貢街在哪裡?.
在哪裡?.
油麻地佐敦.
對呀,大家都是街坊.
知道西貢街不是在西貢.
好了,深一點的第二條.
柴灣角街在哪裡?.
嘩,你們這麼厲害.
你們很多人住荃灣.
柴灣角街是一條很彎彎的街.
正常你是不會去到的.

$^{81}$柴灣角街不是在柴灣.
是在荃灣.
不要去錯地方.
好了,這條應該是最深的.
我要小心這裡.
銅鑼灣山路在哪裡?.
嘩,你們全部這麼厲害.
沒有人去過這個地方.
其實應該偶爾你經過你不會去的.
真的,銅鑼灣山道不是在銅鑼灣.
銅鑼灣山道是在沙田.
接近大衛的交界.
原來你們都喜歡四處逛逛.
你會發覺其實香港有很多街名.
我們說四個字是名不符實的.
當然街名名不符實.
我想對我們的影響不是很大.
頂多可能只是誤導你去錯地方.
或者說你浪費了一些時間或者車錢.
但是如果名不符實的是一間教會.
就不是開玩笑.
我們不能夠坐視不理.
需要立即去加以糾正.
事實上剛才經文所讀的殺敵教會.
就是一個這樣的例子.
你會留意和之前我們提過的四間教會.
以弗所,士妹娜,別加摩和推亞推拉教會相比.
殺敵教會其實很特別.
表面上它真的沒有外人的攻擊.
社群,異教徒不會對它攻擊.
你也看不到書上提及有異端的入侵.
更加反映不了他們.
好像之前提及的士妹娜教會有財政壓力.
表面上它是一間發展得不錯的教會.
但你看看耶穌基督對這教會的評價.
整體評價其實比之前那四間教會都要惡劣.
甚至你見到耶穌指控他們.
是一間有名無實的教會.
所以很盼望藉著耶穌基督對殺敵教會.
所說的這段說話成為我們的提醒和幫助.

$^{121}$特別明天就踏入十二月.
紀念耶穌基督降生的日子.
求主也幫助我們.
好像講題一樣警醒.
預備一個警醒的心.
隨時預備好迎接主的再來.
看回剛才所宣讀的經文.
你留意到耶穌基督對殺敵教會的自我宣告.
耶穌說他就是擁有上帝七靈和七星.
意思是耶穌基督對殺敵教會的一舉一動療癒指掌.
完全監察並且掌握.
並且耶穌能夠穿透教會表面的行為.
看到它裡面的屬靈狀況是怎樣.
耶穌對殺敵教會說.
我知道你的行為.
按明你是活的.
其實是死的.
什麼意思呢?.
行為可以說是教會的一些工作.
或者我們常說的事工.
按明這裡是說名義上或名譽上的意思.
耶穌基督對教會說.
殺敵教會沒錯.
你是當地擁有好的明星.
名義上你是一間很活躍的教會.
很有吸引力.
但可能是怎樣呢?.
你想像到一間教會.
為什麼會有吸引力和有明星.
可能是因為它有不少活動舉辦.
亦似乎吸引到周圍的群眾來參加.
所以當時的殺敵教會應該是頗有人氣.
外界對它一致的好評.
不過為什麼耶穌卻說他是死的呢?.
很大可能是因為.
雖然教會有很多事工發展.
但焦點不是為了彰顯耶穌基督的福音.
又或者說不是為了培育信徒.
怎樣對真理的實踐和追求.
雖然殺敵教會在當地社區有知名度.

$^{161}$但似乎發揮不到教會.
應該有的真正使命和功能.
殺敵教會的情況.
我想就好像一個人.
他又跑步又健身又注重飲食.
表面上沒有病沒有痛.
會走會走.
但當耶穌基督好像照X-ray CT.
MIR一照下去的時候.
他就照出殺敵教會其實裡面的狀況很嚴重.
可能心血管又塞了幾條.
甚至幻想癌症有抗散的跡象.
所以你看到耶穌指責他的行為.
沒錯是活的很有活力.
但從屬靈上已經是好像死了一樣.
我們知道其實在第一世紀.
基督教是一個新興的宗教.
討不討好?.
當時基督教是不討好的.
甚至乎可以說惹人反感.
特別當教會要立定心智.
他要傳揚耶穌基督的科學真理的時候.
就一定會面對到反對的聲音.
甚至乎是阻力.
再嚴重一點的就會被排擠甚至被迫害.
但留意殺敵教會反而是怎樣?.
反而是可以在當地獲得好的明星.
其實很奇怪和其他教會很不同.
隱約可能都反映了在一個異教林立的環境裡面.
殺敵教會所採取的生存之道到底是怎樣.
很大可能是盡量不要向群眾.
不要向異教徒甚至乎猶太人.
傳揚耶穌的福音.
因為他們會覺得不中聽.
他們覺得會很愚蠢.
又或者不要挑戰那些不合上帝心意的意識形態.
那些民間風俗.
因為會激起民憤.
總之就要與人為善.
這樣教會自然不會被當地人反感排斥.

$^{201}$甚至不會攻擊迫迫他們.
換句話來說.
為了和當地的社群來和睦共處.
殺敵教會似乎採取了一種與世俗妥協為由的態度.
特別在第四節你會看到.
耶穌怎樣說?.
留意到只有幾個人的衣服沒有被玷污.
很大可能整體殺敵教會都與世俗為誤.
成為了一間不熱心為耶穌基督作見證的地方.
既然如此.
魔鬼撒旦自然也沒有興趣攻擊和搞擾他.
但是殺敵教會又怎樣去證明.
正確地自己仍然有所作為呢?.
就是透過很多的活動和行為.
這是最簡單快捷的方法.
將自己變成包羅萬有.
提供了很多不同的服務.
但卻不注重福音的傳遞.
可能也有很多宗教的儀式舉行.
但卻不強調真理的實踐.
所以按世界的標準.
殺敵教會的工作行為很成功.
很有成就.
但上帝的名字卻沒有被高舉.
所以耶穌基督對殺敵教會的另一個指控是.
因我見你的行為.
在我的神面前.
沒有一樣是完全的.
完全這個字不是說完美無效.
不是說做到百分百完美.
完全這個字的意思是.
沒有完成.
沒有貫徹始終.
即是說殺敵教會有很多工作和行為.
確實贏取了很多人的口碑.
但卻贏取不了上帝的肯定和稱讚.
是因為他們沒有專心一意.
將上帝託付給他們做的每一件事.
都好好地完成.
很大可能是很多工作和行為只是留於表面.

$^{241}$很有幹勁,做了很多事情.
設計了很多事情.
但卻沒有貫徹始終.
在每件事當中按上帝的旨意.
做對它,做完它.
所以你會看到耶穌基督給他的成績表.
沒有一樣是可以通過.
沒有一樣能夠符合上帝的期望.
所以你會看到殺敵教會.
在耶穌基督的心目中.
為什麼會收到一張這樣的成績表.
所以耶穌基督也告誡殺敵教會.
他說第二節開始說.
你要警醒,堅固那些剩下將要衰微.
和學本特別括了一個括弧.
衰微的意思是死的.
即是快要死的意思.
意思是提醒殺敵教會要醒了.
不要再睡了.
不要以為自己仍然很健康,很穩妥.
其實原來在上帝的面前.
他們的屬靈狀況正在步向死亡.
要盡快鞏固那些快要沒了的生機.
你看到耶穌基督表示殺敵教會是危危乎的.
所以一定要怎樣?.
一定要盡快進行信仰鞏固的工程.
我有時也會看一套電視節目.
不知道大家有沒有看.
叫做《空間改造王》.
在某個台上,有些人有看.
挺好看的.
每次都是建築師幫那些住了幾十年的舊屋.
在日本的舊屋.
怎樣去改善空間和居住環境.
來幫他們進行翻新.
多數每一間舊屋.
當搬走所有東西.
當那些裝修拆置剩下的一間空屋.
你就會發現多數的屋都存在一個問題.
你猜猜是甚麼?.

$^{281}$一間住了幾十年的樓.
通常最大的問題是甚麼?.
是地基.
通常建築師都會發現.
當拆走所有東西.
露出那間屋的本相的時候.
就會發現原來地基用的木材.
經過這麼多年已經腐爛了.
有白蟻,可能漏水,滲壞了,破爛了.
其實很危險的.
一有地震就會倒塌.
隨時就有倒塌的危機.
所以基本上每一集.
第一件事要做的就是.
先把地基搞好.
一定是第一步.
花最多的錢.
重新去換木材.
倒水泥,來鞏固地基.
然後才在上面進行其他的事.
事實上耶穌基督在這裡呼籲.
撒迪教會同樣是.
先搞好你的信仰根基.
甚麼都不用說.
耶穌說你要回想.
即是記起教會當初是怎樣領受.
怎樣聽到福音.
裡面的內涵是甚麼.
並且當日教會是怎樣遵守.
耶穌基督的吩咐.
以至到教會當時確實被建立起來.
不過可能當運作了很多年之後.
就好像一間屋似的.
不知不覺就怎樣.
這又壞了,那又走樣,突突突.
開始走樣.
開始在上帝面前.
有很多事都不符合上帝的期望.
耶穌基督就說.
在這個時候應該怎樣.

$^{321}$當你發現的時候.
當有人提醒你的時候.
你就要立即悔改.
悔改的意思是.
停止你繼續運作的模式.
並且你要轉向正確的道路.
你要重新持守信仰.
換句話來說.
耶穌基督要教會回想這件事.
其實是在說.
撒特教會最初領受的福音內涵.
其實是純全正確的.
最初領受的東西是沒有問題的.
所以耶穌基督表示.
要解決撒特教會的現況.
它現在出現的危機.
耶穌說.
你不是要再聽多點東西.
你不要再聽多點新的道理.
不是要再學多點新的知識.
耶穌說不是這個方向.
而是將你已經有的.
最初領受了的東西.
你重新貫徹做出來.
從這個角度看.
我們就開始明白.
為什麼撒特教會的靈命.
會日漸式微.
變成有明無實.
確實當初.
他聽到福音的真理.
純全的真理.
他領受了耶穌基督的吩咐.
在地上實踐信仰.
建立教會.
只不過是怎樣.
最難的就是你怎樣去持守.
特別在一個世俗誘惑的環境裡.
你怎樣去堅持聖經的教導.
不要被世俗來同化污染.

$^{361}$他們差的大概就是這樣.
同樣的,姐妹.
我們今天絕對不缺乏聖經的真理.
你要明白聖經.
明白信仰的知識.
你在教會,在網上,在書室.
你有海量一般的信仰知識.
你隨時可以接觸到聖經真理.
我們欠缺的.
就好像耶穌基督所提醒.
一份警醒的心.
一份堅持無論什麼情況.
都實踐信仰的態度.
我們能否為上帝交託給我們的事情.
我們做到最好.
堅持到底.
以至上帝的名被高舉.
這是耶穌基督對撒迪教會的提醒.
早幾年.
我要處理睡眠窒息的問題.
我現在有差不多一噸位.
鼻鼾也很厲害.
所以要看醫生.
怎樣處理.
通常你去到公立醫院.
有一個程序.
看完之後.
建議你減肥.
減肥後.
可以改善睡眠質素.
或改善鼻鼾.
於是給了你一些建議.
減少放一些不健康的東西進肚子.
於是我每天的早餐都很好.
多士要走牛油.
不要吃沙嗲牛肉麵.
吃清一點.
吃榨菜或火腿通.
另外.
凍檸茶要走甜.

$^{401}$如果喝咖啡就要走奶和糖.
你猜行不行?.
其實真的行.
我用了大約一個月.
體重瘦了四公斤.
厲害嗎?.
我發覺很厲害.
腰部的皮帶扣進了一格以上.
真的.
我太太也看到我這個情況.
不過已經是往事.
當我發覺有明顯改善的時候.
我就會開始怎樣?.
人總是這樣.
你就會放鬆自己.
加上現在開始這個時間.
天氣冷就怎樣?.
吃火鍋.
肥牛又來了.
這樣又吃那樣又吃的時候.
覺得我有本錢去放鬆.
我可以放鬆一點的時候.
最終我又不知不覺.
我的體重又低了.
反彈了.
所以其實在這個世界上.
真的有很多事情.
會令我們對信仰鬆懈.
放軟手腳.
可能是你忙碌的工作.
可能是你很累.
累到連聖經也不想看.
甚至可能崇拜也在網上就算了.
又或者過量的娛樂.
又或者安逸的生活.
原來都會令我們對信仰疏懶.
這一切其實本來都是上帝所賜的福.
但同時也是兩刃劍.
因為人的罪性緣故.
會變成了對我們的試探.

$^{441}$引誘我們遠離了耶穌基督的吩咐.
不願意持守信仰.
不過耶穌很愛我們.
祂不會任由我們.
當我們開始墮後的時候.
甚至好像撒迪教會這裡.
開始放任的時候.
耶穌是會接手管理.
祂會提醒那些軟弱疏於防範的人.
你要醒了.
你要回轉.
正如祂提醒撒迪教會一樣.
祂這裡說.
若不警醒.
我必臨到彌蘭裡.
耶穌基督形容自己好像賊一樣.
如同賊一樣.
我何時來到.
你也決不能知道.
事實上.
警醒,疏於防範這兩個主題.
經文裡面沒有提及.
但很多聖經歷史學家都提及.
這兩個主題其實對撒迪這個城市.
一點都不陌生.
為什麼這麼說呢?.
因為原來在古時.
撒迪城其實是一個不容易攻破的城市.
因為他們擁有一個提摩留斯山.
作為一個軍事屏障.
意思是什麼呢?.
意思是撒迪城是一個平地裡興建.
但背靠著一個很陡峭的山脈.
像一個天然屏障保護著.
要去撒迪城一定要爬過這座山.
但爬過這座山很艱難.
因為一來很險要.
二來它只有一些很窄的小路.
基本上你不會察覺.
所以當時的人說.

$^{481}$基本上撒迪是一個攻不破的城市.
但很有趣的是.
在歷史上.
這個城市卻是兩次被攻陷.
而兩次被攻陷的原因.
就是因為不警醒.
就是因為疏於防範.
第一次發生在主前的550年.
當時的波斯王古烈要進攻撒迪城.
因為像剛才所說.
找不到上山的路.
於是一直圍攻著這個城市.
也沒有辦法的時候.
於是付出重傷.
希望有士兵能夠想到破敵的計謀.
最終有一個波斯士兵發覺了.
原來有一個秘密的小道.
可以上山進入撒迪城.
於是波斯大軍就漏夜抹黑.
順著山路進攻撒迪城.
而當時撒迪城的守衛.
因為自己有山勢的險要.
認為波斯的大軍絕對沒有辦法攻進來.
所以湛江的守衛是怎樣呢?.
就是先睡覺.
完全不察覺敵人已經從小路順利進城.
自以為有恃無恐的撒迪城.
結果被波斯大軍消滅了.
到了三百多年後.
到希臘王.
希臘王的安提阿普三世又進攻撒迪城.
這次的撒迪人同樣沒有吸取當日的歷史教訓.
他們仍然支持.
我們有這麼險要的山脈來保護我們.
敵人沒有辦法入侵.
於是那些守衛同樣鬆懈.
同樣睡著覺.
導致這個城市第二次被攻陷.
事實上這種沒有警醒.
疏於防備的心態.

$^{521}$現在同樣出現在撒迪教會身上.
所以你看到耶穌基督鄭重地警告他們.
不要再分碎.
別以為那是很堅固的保障.
否則耶穌就會像賊一樣.
突然在意想不到的情況下.
出現在教會面前.
要向教會問責.
要教會交賬.
提醒撒迪教會其實需要怎樣呢?.
四個字,保持警醒.
要認定主是會再回來.
是一定會再回來.
也一定要向各人.
向教會要求他們交賬.
並且在這個大前提下.
教會需要重新檢視.
到底我們的存活之道是怎樣.
如果像經文所說.
像耶穌所說.
就是你還原你的基本步.
就是記起教會當初是怎樣興旺.
當初怎樣去實踐他的信仰內涵.
你現在恢復就行了.
你現在做回當日.
你在耶穌基督身上領受的.
簡單來說有三方面.
我們常說.
向上,我們有一個忠心敬拜的生活.
對準上帝,對準耶穌.
對內,教會智力.
門徒的訓練,裝備.
以至他懂得怎樣按著聖經真理.
不單在教會.
亦在世界裡怎樣彰顯,作見證.
對外亦鼓勵弟兄姊妹出去服事.
去傳揚福音,為主作見證.
謹守著這三條我們說的信仰之柱.
你看到耶穌基督說.
這才是撒迪教會起死回生的道路.

$^{561}$當然要堅持走這條道路.
是不容易的.
因為是一條很窄的路.
當高舉耶穌基督的福音.
智力為主作見證的時候.
耶穌說必然像他一樣.
會遭受世人的挑戰.
甚至迫伯.
事實上你會看到.
撒迪教會很多人已經放棄了.
不走這條路.
不堅持和耶穌基督同行.
但是耶穌亦說.
原來有幾個人是例外的.
第四至六節.
我想請大家一起讀.
引述在程序表或螢光幕內.
預備起.
.
謝謝.
.
自己的衣服的人,才是名副其實的基督徒.
他們堅持站在上帝的那一邊,而不是站在人多那一邊.
他們按著聖經的真理來行事,而不是隨從世俗的生存法則.
他們專心尋求上帝的喜悅,多於追求人的認同和肯定.
他們確實沒有能力改變那個大環境.
他們沒有能力改革撒迪教會很多的問題.
但他們卻能夠保守自己,不被世俗污染.
我想這是對我們一個很重要的提醒.
有時要改變一些積雜了很久的問題是很難的.
是不容易的,你會覺得很無力.
但耶穌提醒我們,你不要因為這樣而放棄.
你堅持下去,堅持下去才做好你的本分.
事實上當我們每一個人都是這樣的時候.
很多我們覺得很難搞的事情,其實都可以改變.
當我們每一個人對準耶穌基督的時候.
很多我們覺得很困難的積雜,其實都可以改變過來.
耶穌表示,這幾個人是配得穿上白衣與他同行.
其實這個畫面讓當時的讀者聯想到羅馬軍隊在打聖仗的時候.
穿白衣,穿白衣在街上巡遊,象徵著海船的景象.

$^{601}$耶穌基督指這幾個撒迪教會的信徒.
也會好像和他一樣,走同一條路.
雖然會為福音的緣故遭患難受迫迫.
但最終也會和耶穌一樣在這場屬靈爭戰中得到勝利.
得到上帝對他們的肯定和認同.
所以耶穌說他們的名字會被存留在生命冊當中.
事實上在啟示錄當說到一個人的名字被存留在生命冊這個象徵表達的時候.
是代表著那個人在困難的當中仍然願意持守信仰,並且堅持到底.
致力彰顯耶穌基督賜給他的生命能力.
以致他能夠向上帝表達出那份堅貞.
而這種態度,這種緣故也會驅使當耶穌基督再來的時候.
會在父神的面前承認這個人是屬於耶穌.
所以弟兄姊妹,「佩德」這兩個字其實是在提醒我們.
確實我們會知道明白基督教是本乎於上帝的恩典.
但聖經同時告訴我們.
行事為人要與所夢的這個恩典相稱.
到底我們是不是一個名副其實「佩德」耶穌承認的基督徒.
還是只是純粹參與了教會的聚會.
其實很值得我們去思考.
有一位宣教運動的領袖,他叫做Alan Hirsch.
他曾經有一篇文章叫做「被遺忘的道路」.
英文是The Forgotten Way.
他提到教會有兩個模式.
第一個其實是在程序表的最後一欄.
一個叫做具吸引力的教會.
Attractional Church.
主要是通過聚會吸引人回到教會.
聚會的內容越是豐富.
信徒就越有信心.
讓新朋友邀請他們回來.
並且有機會接觸到福音,甚至信主.
而另一個模式叫做使命教會.
Missionary Church.
主要是通過信徒外展的服侍和見證.
強調基督徒所經歷的福音.
也應該彰顯在教會以外.
在職場的裡面,在家庭的裡面,在社群當中.
Alan Hirsch從歷史的角度來分析.
早期教會本質上應該是使命模式的教會.
當面對挑戰,面對迫迫的時候.

$^{641}$信徒於是無法在教會聚會,所以分散.
四出全道,結果福音傳開,傳得更廣.
甚至去到當時世界的中心.
但到了羅馬帝國的後期.
基督教就從一個被迫迫的宗教變成一個國家的信仰.
在那個時候,每個人基本上都應該是基督徒.
如果不是基督徒,反而是少數,反而是異類會被排擠.
而當整個帝國都變成一個基督教國家的時候.
所有人都會出席參與教會聚會的時候.
無論是教會或政府自然都會投放更多資源.
來到聚會的地方或聚會的內容身上.
教會亦因為這樣漸漸變成了一個吸引人的據點.
成為一個吸引力的關鍵點.
這種具吸引力的教會模式其實就是在這樣的歷史中誕生.
亦成為了西方教會所仿效的模式.
甚至後來華人教會所採納的一種教會模式.
它的優點是可以讓報道會的工作變得精簡.
信徒亦感覺容易達標.
因為只需要邀請人回來教會就可以了.
譬如以前我們一些大型報道會.
教會邀請了有名氣有份量的講員.
信徒就專注於邀請新朋友參與.
在這個氛圍下固然有它的好處.
但亦都會形成一種情況就是.
我們很著重如何邀請新朋友聚會.
某程度上也削弱了我們自己個人出去傳福音的動力和熱誠.
其實無論是哪一種形式都好.
其實提醒我們他們不是對立的.
首先具吸引力的教會的重點就是聚集.
因為教會我們常常強調是上帝的家.
是信徒的家.
所以很應該努力帶人回來認識福音.
這樣說無可厚非.
至於使命教會的焦點就是為分散而聚集.
我們之所以聚集不是純粹為了內聚一切.
最終的目的是為了要散開來見證上帝.
他強調聚會不是終點.
聚會只是讓信徒在屬靈上充電.
一起回來敬拜上帝.
領受上帝的恩典.

$^{681}$領受上帝的話語.
亦都一起接受裝備去建立教會.
建立上帝的國度.
亦都一起去彼此的守望.
守望教會的弟兄姊妹.
守望這個城市.
然後各人又再一次分散.
回到神給我們的每一個生活崗位裡.
參與上帝在地上的復和和更新的工作.
所以無論是哪一種模式.
他們其實是發揮一個平衡互補的作用.
讓我們的教會觀變得更加全面.
和更靈活有彈性.
所以弟兄姊妹盼望今天這段經文.
讓我們可以重新去思考.
教會存在的使命是什麼.
又或者說作為一個信徒.
我們會不會都是課真假實.
我們不單單只問.
我要為上帝做什麼事情.
我們都要去警醒思考.
其實上帝在我們當中正在進行什麼.
我作為他的僕人.
我可以怎樣裝備好自己.
配合他的科研工作.
為耶穌作見證.
讓我們同心討高.
我們知道上帝很愛教會.
很愛每一個熟悉你的人.
特別當我們有軟弱.
當我們走迷的時候.
你總是多次多方.
用不同的途徑來提醒我們回轉.
提醒我們警醒.
就正如你對撒迪教會一樣.
幫助我們一直到今天.
你對撒迪教會的提醒.
也成為我們的提醒.
讓我們回到初心.
讓我們重拾起初領受的純全信仰.

$^{721}$讓我們總是堅持.
來將耶穌基督的吩咐.
將祂託付給我們的使命.
來到地上實踐.
上帝我們知道.
我們不能靠自己.
因為我們很軟弱.
我們面對生活上不同的挑戰和誘惑.
很容易會攔阻我們走近你.
進行你的話語.
盼望聖靈幫助我們.
讓我們真的能夠.
因著你同在的緣故.
我們能夠為上帝堅持到底.
我們亦能夠彼此同心.
互相守望.
彼此督責.
以致假如我們真的有疏懶.
求你讓我們警醒回傳.
重新在信仰裡.
重拾對上帝的熱誠.
盼望主你親自使用我們.
讓我們不單在教會裡.
領受上帝你各樣的恩典和供應.
我們也走出去為上帝你作見證.
因為我們聚集原是為要被主你差遣.
進到地上不同的角落裡.
去祝福.
代表耶穌基督你去祝福身邊更多的人.
代表耶穌基督你去關顧每一個有需要的人.
盼望主讓我們看到.
你給我們這份使命.
看見你在我們當中的作供.
以致我們懂得如何配合.
懂得如何回應.
恭敬將我們每一位仰望在你的恩手.
願上帝你藉著這段經文.
繼續在生活裡繼續對我們說話.
我們這樣的禱告交託.
奉基督耶穌你得勝的名字來祈求.

$^{761}$阿們.
\newpage



\section{啟示錄 3:7-13}
\label{sec:XJJ1wNBmbSY}
\textbf{2025年12月7日主餐崇拜直播｜陳明泉牧師｜耶穌給教會的七封信:謹守與保守｜啟示錄3\_7-13}
\newline
\newline
連結: \href{https://youtube.com/watch?v=XJJ1wNBmbSY}{\texttt{https://youtube.com/watch?v=XJJ1wNBmbSY}} ~~~~ 語音日期: 2025-12-07
\newline
\newline
\hyperref[sec:XMVPsTUvcaY]{\small{< < < PREV SERMON < < <}}
~
\hyperref[sec:index_chronic]{\small{[返順時目]}}
~
\hyperref[sec:1zYtfRDiK4g]{\small{> > > NEXT SERMON > > >}}
\newline
\newline
啟示錄 3:7-13
\newline
\begin{longtable}{cl}
\hline
\hline
章節 & 經文 (和合本修訂版)\\
\hline
3:7 & \begin{tabularx}{0.7\textwidth}{X} 「你要寫信給非拉鐵非教會的使者，說：『那神聖、真實的，拿著大衛的鑰匙，開了就沒有人能關，關了就沒有人能開的這樣說： \end{tabularx} \\ \\ \relax
3:8 & \begin{tabularx}{0.7\textwidth}{X} 我知道你的行為。看哪，我在你面前給你一個敞開的門，是沒有人能關的。我知道你有一點力量，也遵守我的道，沒有否認我的名。 \end{tabularx} \\ \\ \relax
3:9 & \begin{tabularx}{0.7\textwidth}{X} 那屬撒但會堂的，自稱是猶太人，其實不是猶太人，而是說謊話的，我要使他們來到你腳前下拜，使他們知道我已經愛你了。 \end{tabularx} \\ \\ \relax
3:10 & \begin{tabularx}{0.7\textwidth}{X} 因為你遵守了我堅忍的道，我也必在普天下人受試煉的時候保守你免受試煉。 \end{tabularx} \\ \\ \relax
3:11 & \begin{tabularx}{0.7\textwidth}{X} 我必快來，你要持守你所有的，免得人奪去你的冠冕。 \end{tabularx} \\ \\ \relax
3:12 & \begin{tabularx}{0.7\textwidth}{X} 得勝的，我要使他在我神的殿中作柱子，他必不再從那裡出去。我又要把我神的名和我神城的名—從天上我神那裡降下來的新耶路撒冷，和我的新名，都寫在他上面。 \end{tabularx} \\ \\ \relax
3:13 & \begin{tabularx}{0.7\textwidth}{X} 凡有耳朵的都應當聽聖靈向眾教會所說的話。』」 \end{tabularx} \\ \\
[1ex]
\hline
\hline
\end{longtable}
$^{1}$《新約聖經》第351頁.
「也,那聖傑真實拿著大衛的藥時,開了就沒有人能關,關了就沒有人能開的.」.
「說:我知道你的行為,你略有一點力量,也曾遵守我的道,沒有棄絕我的名.」.
「看啊,我在你面前給你一個敞開的門,是無人能關的.」.
「那撒旦一會的,自稱是猶太人,其實不是猶太人,乃是說方話的.」.
「我要使他們來,在你腳前下拜,也使他們知道我是已經愛你了.」.
「你既遵守我忍耐的道,我必在普天下人受試煉的時候,保守你免去你的試煉.」.
「我必快來,你要持守你所有的,免得人奪去你的冠冕.」.
「得勝的,我要叫他在我神殿中作柱子,他也必不再從那裡出去.」.
「我要將我神的名,和我神成的名,寫成就是從天上,從我神那裡降下來的新耶路撒冷.」.
「並我的生命都寫在他上面.」.
「聖靈向宗教會所說的話,凡有耳的就應當聽.」.
弟兄姊妹平安.
平安.
我們今天繼續用啟示錄進入耶穌基督級教會的七封信系列.
藉著耶穌基督對不同教會的勸勉.
來幫助我們在這個世代.
不論是在信境還是在香港.
都陷入一種困難,大家情緒很低落的日子裡.
我們藉著上帝的話語,藉著耶穌基督的勸勉.
讓我們能夠有些力量.
來面對生命裡面的林林種種,高高低低.
在我們看這段經文之前,不知道大家有沒有這樣的經歷.
在你小時候,有沒有小時候被家人趕出門口的經歷.
我在AI裡面,我做了一幅我自己小時候的一幅相片.
可不可以音響幫我按一下powerpoint.
這就是我小時候,AI製作出來的.
所以樣子其實和我有些不相似.
但是這是一個什麼經歷呢.
小時候和我阿姨在家裡,我阿姨比我大一年.
在家裡和阿姨打架.
其實是阿姨打我多過我打阿姨.
但是因為打架,於是乎.
因為我小時候在我外婆那裡照顧.
當我和阿姨打架的時候.
阿姨打我,但是外婆是趕了我出門口.
我就站在門口裡面,很害怕.
因為我沒有試過這樣的一個經歷.
小學生的時間.
還要打架,我的臉都被打到全都損了.

$^{41}$然後外婆還說,我不要你.
我們這裡沒有這麼頑皮的小朋友.
我就在家門口裡面,其實坐了那十分鐘.
但那十分鐘實在很煎熬.
還有一種感覺,糟了,我被趕出家門.
我又和阿姨打架.
做一個這麼不理想的,這麼頑皮的行為.
我是否永遠都要坐在這裡,或者我以後回不了家.
那十分鐘就不停想這些問題.
小朋友,我想在在都會這樣.
小時候家人是用土法的方式來教.
沒有什麼正面管教.
不過這個經歷就提醒.
原來做錯了事,被驅逐.
或者被隔離的那種狀態是很慘的.
做錯事,被隔絕,或者被趕出家門.
這個是小朋友學的方法.
但如果做對了事,被驅逐.
又會有什麼樣的體會呢?.
今天所說的教會,非拉鐵非.
他們就是因為持守信仰而被驅逐出來.
究竟這段經文對於當代的非拉鐵非門徒.
有什麼重要性,也對我們今天有什麼重要性呢?.
我們今天看這段經文,希望給我們一些養分.
我們先祈禱.
上主願你在我們當中大靈.
讓我們能夠明白你的話語.
藉著你對我們的教導.
和你對我們心靈的呼聲.
讓我們在這個世代裡學懂怎樣站立得穩.
祈求你與眾弟兄姊妹同在.
幫助聆聽的每一個.
他們聽得明白,又讓孫說的說得清楚.
忠心的將你的話語帶動在眾弟兄姊妹身上.
求你與我們同在,聖靈引導我們進入真理.
獻上禱告奉求基督耶穌得勝的聖名,Amen.
我們一起看看這段經文.
啟示錄三章七至八節.
你要寫信給非拉鐵非教會的使者說.
那聖潔真實拿著大衛的約時.

$^{81}$開了就沒有人能關,關了就沒有人能開的說.
我知道你的行為,你略有一點力量也曾遵守我的道.
沒有棄絕我的名.
看啊,我在你面前給你一個敞開的門,是無人能關的.
在這裡耶穌基督用一個很重要的身份對教會說話.
耶穌基督是一個甚麼的主呢?.
他說到他是聖潔真實拿著大衛約時.
這條鎖匙開了就沒有人能關,關了就沒有人能開的身份.
對教會說以下的說話.
首先本身這樣介紹自己.
耶穌基督的自稱對於當代的非拉鐵非來說.
已經是一個很重要,很定心的一個訊息.
他說到耶穌基督是聖潔的主,分別為聖.
聖潔很多時候我們以為是道德層面不做一些污糟的事.
那些叫聖潔,聖經有說這個角度.
不過聖潔也有更加多,你看看舊約的利未記.
就說到分別出來,和這個世界分別出來的那班人分別為聖.
聖潔就是用一個這樣的角度去理解.
這個聖潔的角度是說到基督徒因為上帝是聖潔.
而基督徒被呼召出來的這班人都是和這個世界不同的.
所以這種被分別出來的心態.
其實對於非拉鐵非來說一讀上去就百般滋味在心頭.
待會會再繼續解釋,你大概知道.
第二樣是真實,真實的意思就是說他的信仰不是虛構,不是虛謬.
即使逆流而上也好,他是確確實實經歷到信仰的那種體會.
和基督耶穌同在的那種經歷.
第三個字眼就是拿著大衛的約時.
這個就特別需要解釋一下.
大衛的約時這個字是怎樣來的呢?.
其實這個字的典故就來自舊約爾賽亞書第22章20至22節.
我沒有刊登出來,不過大家記住.
爾賽亞書第22章20至22節.
在這裡面這樣說的.
爾賽亞對著當時的以色列整個猶大民族說.
到那一日我必召我僕人希勒加的兒子以利亞進來.
將你的外袍給他穿上,將你的腰帶給他繫緊.
將你的政權交在他手上.
他必作耶路撒冷居民和猶大人的父.
我必把大衛家的約時放在他肩頭上.
他開了就沒有人能關,他關了就沒有人能開.

$^{121}$你看到其實啟示錄的作者約翰就拿了爾賽亞書這個典故.
就是在說面對著亡國.
以前整個以色列人面對著國難.
但是他講到上帝定義要將以利亞進.
成為了整個猶大人希勒加的王.
所以他要將那個多麼艱難也好,多麼困難也好.
上帝定義要以他作為以色列的王.
他掌管的那條大衛的約時其實就代表著.
他就是在國家面對著困難當中.
他都依然坐著為王.
他擁有來自上帝的權柄繼續來管治以色列民族.
在這個典故裡面.
啟示錄的作者約翰就將這個大衛的約時.
就放在耶穌基督的身份上.
他講到即使那個世代.
菲拉鐵菲都面對著很多的困難.
但是他所信的主是信實,是聖潔.
而且是掌握權柄的.
所以帶著這份權柄.
他們就能夠有信心面對眼前的挑戰.
更加講到的就是那個約時.
那個守門的那個人.
他開了那道門,沒有人可以關的.
他關了那道門,就沒有人能夠開的.
即是在講的就是決定那道門開與關.
或者誰出與入,是他說了算.
就不是這個世界或者其他人可以掌握的一種權柄.
耶穌基督在這裡講一個這樣的自稱.
其實與菲拉鐵菲這個教會當時的處境是有關的.
菲拉鐵菲這個字,本身這個教會或者這個地方的名字.
英文名字就叫做Philadelphia.
如果你知道就是今天的美國有一個州.
叫做費城,就用這個字.
Philadelphia,用費城就譯不了.
Philadelphia這個字是很有意思的.
它來自兩個字根,一個就是愛,第二個就是兄弟.
所以其實Philadelphia的意思就是兄弟相愛.
這個教會在這個城市裡面.
是一個弟兄相愛.
是一個充滿了愛連結的一個小城市.

$^{161}$不過這個地方就強調了人與人之間是相親相愛.
有一份很緊密的一份關係.
雖然是這樣,但菲拉鐵菲這間教會在這個城市裡面.
其實他們正在面對一些挑戰.
在第三章第九節啟示錄裡面,他就這樣說.
撒旦一會的自稱是猶太人.
其實不是猶太人乃是說方話的.
我要使他們來在你腳前下拜.
也使他們知道我是已經愛你了.
在這裡面為何要特別將耶穌基督的身份.
和這個處境來做一個對比.
或者他要對著菲拉鐵菲教會.
他在說他們面對甚麼處境呢.
在他們面對的處境裡面特別說到.
就是撒旦一會的.
撒旦一會的意思.
他用的字眼,撒旦一會其實也不是最理想的.
有些和合修定觀經譯就是撒旦的會堂.
如果你熟悉聖經的弟兄姊妹.
你會知道會堂是甚麼人來做祈禱和學聖經的.
猶太人.
原來菲拉鐵菲裡面有很多猶太的基督徒.
他們歸順了耶穌基督.
但是當時的猶太民族或者在這個地方裡面.
有外邦人也有猶太人.
那些外邦人就覺得信耶穌的那些人很奇怪.
因為信主沒有一尊偶像在那裡拜.
也不是拜神的,覺得他們是無神論者.
而對於那些猶太人來說.
那些一直讀律法或者堅守律法的猶太人.
覺得跟耶穌的那班人是異端.
所以這班人就被驅逐出去猶太會堂.
所以他用的會堂這個字.
其實就是說那班是猶太人.
但是驅逐了那些基督徒不在會堂裡面.
不讓他們集會.
不讓他們參與在當中裡面學聖經.
或者參與崇拜的一些禮儀.
我們不要說東家不去就去西家.
這個會堂不讓我就不去.

$^{201}$不是今天,不是這樣的意思.
會堂其實是有他整個民族身份的奠定位置.
就是說如果驅逐他們離開會堂.
某程度上就是剝奪了這些人在社會上.
本身已經是邊緣化.
邊緣化裡面再被邊緣化的一個完全沒有身份的人.
因為在菲拉鐵飛這個城市是羅馬的城市.
外邦人已經排他猶太人.
猶太人當中已經有一小部分再要被排斥出來.
可想而知這班人是不是整個社會裡面.
一班極之少數的一班人.
經文裡面再繼續說.
排斥他那些人自稱是猶太人.
耶穌說那些不是猶太人.
那些不是神的選民.
那些人不是說謊的.
那些人不是真實地面對他的信仰.
不過就是按著自己的理解.
他就將這些人踢走離開會堂.
所以他第一件事說到那些人是殺蛋會堂的人.
他用那個會堂就好像一班說謊.
不承認真理的一班人.
他們的聚集卻是將真理.
或者信耶穌的一班門徒拒之於門外.
第二件事有一些文獻.
特別就說到當代第一代的基督徒.
就怎樣去面對那些猶太人.
有一個如果讀教會歷史就知道.
有一個很重要的護教學家.
亦都是當時的一個殉道者.
就叫做Justin Martyr.
這個游斯汀.
他就寫了一封信.
就是對於猶太人德里夫的對話.
在這個書卷裡面他講到.
他特別講那些當時第一代.
或者是第一世紀第二世紀的那些猶太人.
他對基督徒的那種迫害.
其中有一個在會堂裡面.
猶太人會堂裡面祈禱的.

$^{241}$學聖經的.
但是他們又加入了一個的祈禱文.
那個祈禱文就叫做.
《密達圖民眾》裡面的就坐而端.
Bakit Kamanima.
就是希伯來文.
不准都沒有人質疑的.
我讀的.
不過這個字眼.
中文的意思就是就坐而端.
在猶太人會堂裡面加入了這一句的圖文.
就坐而端.
他就坐甚麼而端呢.
他們就是就坐基督徒.
在游斯汀裡面講到.
每一日那些猶太人.
在會堂裡面除了學聖經之外.
他們就望著門外的基督徒.
就在那裡就坐他們.
求上帝降臥在他們身上.
又或者令到這些人不能夠立足.
甚至乎要讓猶太教.
或者他們這班敬拜者.
和那些基督徒要劃清界線.
在有些歷史文獻裡面.
就強調當時的基督徒.
在社會裡面.
是一班極之少數的人.
這一班人卻是沒有得到社會上的歡迎.
卻是成為社會上排斥的一班對象.
今日我想我們作為基督徒.
我們在社會上都未至於到.
好像第一代社會裡面的那種排斥感.
不過如果你對於一些比較尖銳.
或者有時候覺得坊間裡面哲學也好.
或者對於一個很開放思維的.
或者學術界裡面.
有時候對於基督徒都嗤之以鼻.
甚至乎有時候會視我們.
是一班傻的人.

$^{281}$一班不合邏輯的一班人.
我們自說自話.
又或者我們自己圍爐取暖等等.
社會上雖然沒有明顯的逼迫.
但是對於一些迷信者.
或者對於社會主流裡面.
我們很多時候都是面對一些這樣的狀況.
我記得我自己信主沒多久.
我都說過可能和大家分享過.
我在大學裡面的校園信主.
但當我信了主的時候.
對我很好的讀哲學的教授.
社會政治哲學的.
他就很抗拒基督教.
但我跟他說我很敬佩耶穌.
但聖經其他人品我全部都嗤之以鼻.
跟著他說Andas你信耶穌嗎.
不是吧 你這麼不明智的舉動你都會做.
跟著他就拿著本聖經.
他就逐句保羅寫的他就質我.
你真的信這些.
你拜一下耶穌的道理就ok.
但你真的信嗎 不是吧.
你拿他的文化內涵.
但你沒必要這樣來信.
他就繼續說.
整本聖經裡面我最不信的就是保羅.
他都是執耶穌的口水尾.
所以他所說的東西沒有動見等等.
他就一直質我的時候.
我第一個反應不是跟他互教.
剛剛信主不懂得互教.
剛剛信主的時候.
只不過是老師對自己很好.
他有一種權威性.
當他一直質我的時候.
我心裡面又不懂得回答.
到最後離開他的辦公室.
離開那個宿舍.
我就說了一句話.

$^{321}$你這樣說也是.
其實我都不是應該信保羅.
我都是拜耶穌.
就不信保羅.
聖經其實可能我們都要選擇性來閱讀.
我說完這番說話之後.
我初時是覺得自己是認同了他.
甚至我不想他對我有一些負面的詮釋.
或者對我有些成見.
所以我就迎合了他的看法.
但說完這句之後.
我一直回去自己的宿舍裡面.
我心極之不安樂.
我認同什麼呢.
為什麼有一個人在我心裡面有份量.
他一說有些質疑的話.
我就自己軟了.
我完全沒有立場.
就會完全無條件接納別人所說的話.
之後很後悔自己去祈禱.
重新去讀過保羅一次.
不過我發覺保羅不是跟人口水美.
只不過是他將耶穌基督的說話.
系統化和牧養性地去教導頂尖姐妹.
頂尖姐妹我不知道你今天會不會遇到.
好像我以前的這種情境.
有人質疑你的時候.
你為了不要令別人對你有些偏見.
又或者你要迎合這個世界對你的一些看法.
你有意無意之間你要將自己扭曲了.
符合了別人對你的期望.
不過當你這樣做的時候.
其實你會明白你的心裡面.
你會覺得不安樂.
因為這個信仰從來不是叫裡面的訊息扭曲.
來迎合這個世界.
而是正正我們活在一個扭曲的世界.
我們就要按著真理來抗衡這個世界裡面的.
那些歪曲薄茂的理論.
菲拉鐵菲弟兄姊妹姐妹其實正面對著一件事.

$^{361}$他們是社會裡面極少數的一班人.
但是他們面對著猶太人當時裡面那種的迫害.
他們依然站立得穩.
在三章第十節耶穌就這樣來稱讚教會的一班的弟兄姊妹.
第十節.
你既遵守我忍耐的道.
我必在普天下人受試煉的時候.
保守你免去你的試煉.
首先在這裡面耶穌基督是特別肯定.
這一班人其實他是遵守耶穌基督忍耐之道.
這句說話都是有一點點結結格格.
是甚麼意思呢.
其實那句更加好一點的翻譯就是.
你既遵守我的道就是忍耐.
即是用忍耐的方式來遵守逆流而上.
一個對基督教不友善.
或者猶太人對他不友善的信仰立場.
那些猶太人怎樣驅趕那些人都好.
他都依然站立得穩.
所以那個遵守這個字.
其實就是保守上帝的道.
在這個道裡面他站得穩站立得穩.
同樣的字眼.
耶穌就說既然你站立得穩.
在這個世界所有人受試煉的時候.
我都會好穩地保守你.
令你免去你的試煉.
有些識經學家或者有些其他牧者讀這裡.
都很開心 因為保守你免去你的試煉.
他就覺得免不用考試.
試煉就是考試.
好像以前有中三評核試.
免去了大家開心了.
直上中四五.
現在只有一個試沒有兩個高中的試.
又免了.
不過這裡其實不是說免去試的意思.
免去你的試煉其實那個保守.
通常如果說上帝的保守.
就不是說你沒有難關.

$^{401}$你被提了直接上天.
不是這個意思.
有些基督徒是這樣解釋的.
不過這一段的經文其實是說.
即使面對著逆流而上.
上帝都會保守你的心.
令到你在這些很兇猛的試煉.
很困難的狀況裡面.
你的心依然是穩固的.
是站立得穩的.
所以在約翰裡面用同樣的字眼.
既講到人的責任.
同時也講到上帝的保守.
他將兩個很重要的信息放在一起.
就講到飛拉鐵飛.
他既有這樣的生命的那種持守.
同時也講到上帝有這個的應許.
在這班人當中.
人的持守和上帝的保守.
在這裡有一個古舊的教父叫做愛淫瘤.
在愛淫瘤的著作裡面講到.
很多時候以為基督徒信了主.
我們信了主一次得救.
永遠得救.
我們覺得我認了信一次就行了.
然後我的人生可以很差也好.
上帝怎樣也會給我們一個天堂入場券.
但愛淫瘤的講法是.
基督徒的生命是一個持續的過程.
不是你信一刻.
你的信心不是一刻的信心.
而是一生忠心.
不斷的來持守信心.
而且是持續的結出果子.
即是用你一生人來活出一個信心的生活.
我們今天千萬不要好像很廉價.
我信了主了.
我有天堂入場券.
在整個基督教信仰裡面.
你信了主也好.

$^{441}$你的信心是要保持住的.
你要保存住的.
如果你不保存住.
你會失腳的.
如果你覺得也很抽象.
你就想起夫婦關係.
夫婦關係結了婚.
老公說我娶了老婆了.
我過了關.
這個女人是屬於我的.
然後我就任意不理她.
那段婚姻是不會幸福的.
因為你立志一生一世愛她.
那個婚禮就是一個立志.
在上帝面前立誓.
而未來人生怎樣走.
你用愛來成全對方.
你愛對方一生一世.
這個才是真正的承諾.
同樣道理.
信耶穌跟隨主也是.
一生一世.
不要讓自己被這個世界的洪流.
不要被很多話打擊.
然後你就放棄了這個信仰.
有時候我們做牧者.
有時候會看到很多頂尖姐妹出出入入.
有時候會有一份唏噓感.
因為可能很多年輕信主.
堅心信弟兄姊妹姐妹.
到一路成長的時候.
世界或者這個職場.
甚至這個世代.
有很多東西衝擊著自己.
有時候覺得那個信仰不夠用.
以致放棄了.
又或者更加吸引他的.
或者星期六星期日這些時間.
有更加多興趣.
他可以實踐.

$^{481}$他想開心.
於是就放棄了他的信仰生活.
我們以為信了耶穌就行了.
但我們沒有想過.
這個信仰不是信了就行.
其實那個信心是要不斷不斷.
這樣來保守下去.
菲拉鐵飛的這一班門徒.
再艱難也好.
他都繼續來到去遵守這個信心.
所以上帝也保守他們的心.
在這個試煉有多麼難的苦難.
災難當中.
他們的信心都不會有所動搖.
這個就是上帝對他們的承諾.
我們再繼續看.
第十一節到第十三節.
這是繼續下來的經文的應許部分.
「我必快來,你要持守你所有的,免得人奪去你的冠冕..
得勝的我要叫他在我神殿中作處子,他也必不在從那裡出去..
我又要將我神的名和我神聖的名.
這成就是從天上從我神那裡降下來的新耶路撒冷,並我的生命都寫在他上面..
聖靈向宗教會所說的話,凡有異的就應當聽.」.
在第十一節到十三節.
也是一個應許的延伸.
既是保守當下那些即使面對著大災難的菲拉鐵飛門徒.
他們的心被堅固.
在將來的日子上帝會賜給他們有冠冕.
但第一件事要告訴他們.
要記住要繼續持守你所有的.
所以你以前信對了,做對了的事就要繼續去做.
不要說上帝有應許然後又放棄了.
不是,保持著,免得有人奪去你的冠冕.
即是你的冠冕要延續到最後的信心你才能得到.
但上帝已經為你預備了冠冕.
接著你用一生來活出這個信仰的時候.
在將來的日子上帝要叫這班人.
即是一個意象來的.
就是在說作神殿中的儲子.
用中國人所說的話就是中流砥儲.

$^{521}$這班人昔日用忠心來演繹他的信仰.
在將來的日子裡面這班人.
就成為上帝新耶路撒冷裡面的中流砥儲.
這班人真的用他的生命來活出信仰.
接著上帝再繼續說.
這個儲子的價值在於什麼呢.
除了他穩固之外.
更加重要的是上帝有三重的名字.
放在這條儲子上.
有上帝的名字.
有上帝神耶路撒冷城的名字.
和上帝新的名字.
用這個叫做三重銘文.
即是說重要的東西說三次.
他說到上帝用一個蓋章蓋章再蓋章.
來說這條儲子是一個經得起風浪.
硬朗有力量的儲子.
坊間的認證算不得什麼.
這條儲子是上帝親自認證.
是這班人用生命努力逆流而上.
來活出中流砥儲的生命.
耶穌基督說你們也不再從那裡出去.
用剛才那些猶太人.
那些人驅逐這些人出去會堂.
但耶穌說這些人將來就是成為了.
新城耶路撒冷神殿裡面的中流砥儲.
永遠都是與上帝同在.
這班人不會被驅逐出去.
也不會被社會排斥.
反而是被上帝稱許.
是一個永恆榮耀的身份.
這班忠心的門徒.
他們的生命是永遠屬於上帝.
今日段經文其實.
如果你去解釋背景其實不是很複雜.
不過對我們今天有意義.
而且是很深刻的意義.
今日我們的身份認同.
即是我們怎樣去界定自己的身份.
你向一個陌生人.

$^{561}$你向身邊的人介紹自己.
你會用什麼身份去介紹呢.
有些頂展會.
以前年輕的時候會說基督徒是我的身份.
但為了要謀生.
或者在這個世界裡面.
令到別人不會有一些負面標籤.
基督徒的身份就放到很低.
不過今日這段經文提醒我們.
如果我們真誠相信.
跟隨基督的這班門徒就是我們的身份.
無論這個世代有多難.
無論我們面對的處境有多困難.
我們都是用這個身份.
用上帝給我們永恆的.
這個榮耀的身份認同.
來好好過生活.
當這個世界有不同的價值觀念的時候.
我們就活出基督徒的價值觀.
當這個世界有不同的學術.
說不同的東西的時候.
我們就用基督徒的眼光來看這個學術世界.
當這個商業社會有多麼狡詐的時候.
我們就用基督徒的身份.
來好好營商.
好好過生活.
在我預備這篇講章的時候.
剛剛澳門發生一件事令我很感觸.
澳門有一個年輕人在我過去講到禮拜五的時候.
晚上凌晨就交通意外突然間去世.
我自己都很悲傷.
我覺得二十多歲的年輕人.
很年輕.
每一次見到他很殷勤服侍.
我是很喜悅.
但剛剛的禮拜五就是最後一次服侍.
很悲傷自己.
昨天在祈禱的時候.
他的女朋友.
有些教會的同工去探訪他.

$^{601}$他女朋友說這個年輕的弟兄有一個心願.
她說她很想用一個基督徒的身份.
建立一間影音的公司.
因為她讀影音,讀媒體.
她的志願是用基督徒的角度.
用基督徒事奉的心.
來建立一間這樣的公司.
她鼓勵她的兄弟.
他們籌組一間這樣的公司.
雖然可能他已經回到天家.
但這個心志用上帝的角度.
作為基督徒的這種價值觀.
來成為他未來人生的志業.
既是心裡面有一種很大的悲傷.
但亦都是提醒我們眾人.
帶著基督徒這個身份好好地過生活.
世界可能不理解你.
但上帝知道.
上帝會將永恆這個身份賦予給你.
弟兄姊妹,這個世代需要硬朗的基督徒.
持守你的信心.
上帝將來會保守你的心.
亦都應許在亂世裡你站立得穩.
將來你就被稱為中流砥柱.
上帝的名字印在你的心上.
願上帝繼續激勵我們,幫助我們.
\newpage



\section{啟示錄 3:14-22}
\label{sec:1zYtfRDiK4g}
\textbf{2025年12月14日崇拜直播｜陳冠成牧師｜耶穌給教會的七封信:重拾愛主的心｜啟示錄3\_14-22}
\newline
\newline
連結: \href{https://youtube.com/watch?v=1zYtfRDiK4g}{\texttt{https://youtube.com/watch?v=1zYtfRDiK4g}} ~~~~ 語音日期: 2025-12-14
\newline
\newline
\hyperref[sec:XJJ1wNBmbSY]{\small{< < < PREV SERMON < < <}}
~
\hyperref[sec:index_chronic]{\small{[返順時目]}}
~
\hyperref[sec:wYKTi6uXSJI]{\small{> > > NEXT SERMON > > >}}
\newline
\newline
啟示錄 3:14-22
\newline
\begin{longtable}{cl}
\hline
\hline
章節 & 經文 (和合本修訂版)\\
\hline
3:14 & \begin{tabularx}{0.7\textwidth}{X} 「你要寫信給老底嘉教會的使者，說：『那位阿們、誠信真實的見證者、神創造的根源這樣說： \end{tabularx} \\ \\ \relax
3:15 & \begin{tabularx}{0.7\textwidth}{X} 我知道你的行為，你也不冷也不熱；我巴不得你或冷或熱。 \end{tabularx} \\ \\ \relax
3:16 & \begin{tabularx}{0.7\textwidth}{X} 既然你如溫水，也不冷也不熱，我要從我口中把你吐出去。 \end{tabularx} \\ \\ \relax
3:17 & \begin{tabularx}{0.7\textwidth}{X} 你說：我是富足的，已經發了財，一樣都不缺，卻不知道你是困苦、可憐、貧窮、瞎眼、赤身的。 \end{tabularx} \\ \\ \relax
3:18 & \begin{tabularx}{0.7\textwidth}{X} 我勸你向我買從火中鍛鍊出來的金子，使你富足；又買白衣穿上，使你赤身的羞恥不露出來；又買眼藥抹你的眼睛，使你能看見。 \end{tabularx} \\ \\ \relax
3:19 & \begin{tabularx}{0.7\textwidth}{X} 凡我所疼愛的，我就責備管教。所以，你要發熱心，也要悔改。 \end{tabularx} \\ \\ \relax
3:20 & \begin{tabularx}{0.7\textwidth}{X} 看哪，我站在門外叩門，若有聽見我聲音而開門的，我要進到他那裡去，我與他，他與我一起吃飯。 \end{tabularx} \\ \\ \relax
3:21 & \begin{tabularx}{0.7\textwidth}{X} 得勝的，我要賜他在我寶座上與我同坐，就如我得了勝，在我父的寶座上與他同坐一般。 \end{tabularx} \\ \\ \relax
3:22 & \begin{tabularx}{0.7\textwidth}{X} 凡有耳朵的都應當聽聖靈向眾教會所說的話。』」 \end{tabularx} \\ \\
[1ex]
\hline
\hline
\end{longtable}
$^{1}$《新約聖經》第3章14-22 新約聖經352頁 《老弟家教會的信》.
你要寫信給老弟家教會的使者,說:那為阿們的,為誠信真實見證的,在神創造萬物之上為元首的,說:.
我知道你的行為,你也不冷也不熱,我巴不得你或冷或熱,你既如溫水,也不冷也不熱,所以我必從我口中把你吐出去..
你說:我是富足,已經發了財,一樣都不缺,卻不知道你是那困苦,可憐,貧窮,乞眼,赤身的..
我勸你向我買火鍊的金紙,叫你富足,又買白衣穿上,叫你赤身的羞恥不露出來,又買眼藥擦你的眼睛,使你能看見..
凡我所疼愛的,我就責備管教他,所以你要發熱心,也要悔改..
看啊,我站在門外叩門,若有聽見我聲音就開門的,我要進到他那裡去,我與他,他與我,一同坐直..
得成的,我要賜他在我寶座上與我同坐,就如我得了性,在我父的寶座上與他同坐一半..
聖靈向眾教會所說的話,凡有疑的,就應當聽.凡有疑的,就應當聽..
請姐妹平安..
終於來到這個系列最後一次宣講,我們一起看看耶穌基督給老弟家教會的信.你會留意,從剛才的經文可以看到,其實在啟示錄七間教會裡面,老弟家教會似乎是最富裕,最富足的一間..
但是你會看到在耶穌基督的眼中,他卻是最佛善足晨的一個.你看到耶穌對這間教會只有責備,完全看不到任何的稱讚.到底老弟家教會出了什麼問題,以至到耶穌給了他一個這樣的負評呢?.
我們先從耶穌基督對這間教會的自我宣稱入手了解一下.你會留意到,在耶穌一開始就會說,他是那位阿們的,忠信真實的見證者,是神的創造裡面的原手..
耶穌說他是那位阿們的,其實意思是什麼呢?阿們的意思是真實可信的.耶穌強調他才是真正可靠,某程度上其實是在暗示老弟家教會現在所依靠的對象,其實不是真實,不是可靠的..
另一方面,耶穌說他是那位忠信真實的見證者,見證者其實有另一個意思,是在說殉道者.我們知道耶穌基督為了要忠心誠實,為真理,為上帝作見證的時候,他不惜犧牲自己..
但是作為他的跟隨者,老弟家教會卻是沉醉於他現時的富足安逸的生活裡面,不再願意為信仰來付上代價.耶穌最後說,他是那位在上帝創造中作元首的,意思是提醒老弟家教會,他才是教會的主,他才是萬有的元首,教會的主宰..
換句話來說,誰可以在教會主導呢?不是老弟家的信徒,而是耶穌本人,教會必須信服於耶穌基督的主權和大令下..
透過耶穌基督以上對老弟家教會的自我宣稱,你會發現似乎有一個景象出現在教會當中,大概是因為現在生活已經富裕,自己覺得已經一無所缺,所以老弟家教會也覺得前途已經掌握在自己手裡,於是開始有些東西走歪了..
所謂的偏行幾路,逐漸將信靠的對象從耶穌身上轉移到物質的財富身上,亦因為這樣,教會的行為也陸陸續續走樣,發揮不了應有的功用..
所以你看到經文繼續說,第十五節,耶穌指責老弟家教會的行為,既不冷也不熱.留意這裡所說的冷和熱,不是在說屬靈的溫度,不是在說冷對耶穌冷淡,熱就是對耶穌熱心..
耶穌很明顯不是這個意思,因為下一句你看到耶穌說什麼?他說:我巴不得你一是冷,一是熱,冷熱也好,但是你既如溫水,又不冷又不熱,所以我必從我口中吐出來吧你..
假如冷熱是屬靈溫度的話,耶穌沒理由說我巴不得你一是冷一是熱,沒可能,沒理由耶穌說你屬靈火熱,一是屬靈冰冷,耶穌不會這樣說,耶穌也不可能說屬靈的冷漠比屬靈的溫暖好..
所以你看到冷和熱不是溫度,到底是什麼意思呢?我們必須留意老弟家教會附近一帶的地理狀況到底是什麼.事實上根據考古的資料,你會知道老弟家教會附近有兩個城市,一個叫希拉波納,另一個叫我們比較熟悉的哥羅西..
他們的位置在哪裡呢?希拉波納在這裡,哥羅西在那邊,老弟家在那邊,好像一個三角形..
在當時古羅馬的時代,希拉波納這個地方可以說是一個水療的聖地,就像我們今天一樣,一個泡溫泉的地方,很出名的.因為當地擁有一些很高溫,含有碳酸鈣的熱泉水,被視為極具醫療的功效..
至於哥羅西,他沒有熱泉水,但反過來,他有一些很冰冷的泉水,很清澈的,很新鮮的,可以直接被人飲用來解渴..
至於老弟家,他沒有自己的水源,那怎麼辦呢?他也要喝水,他就要透過一些水管將外面的水引進來.所以在老弟家,他會用一些石造的水管將幾公里以外的熱泉水引進來當地..
那麼你知道那些熱泉水經過幾公里引進來,就會變成什麼呢?還沒變成冷水,變成暖水,暖暖的.並且你知道當中的礦物沉澱的時候,會有些化學反應..
所以會令到那些溫水有股味道,很難入口.如果你直接飲用的話,就會像經文說的,會反胃,立刻想吐出來..
所以當地人一定會將那些從外面引進來的泉水,先放在一個儲水缸,讓它一方面放涼一點,真的放涼了,另一方面讓那些氣味散走了,才可以飲用..
換句話來說,經文這裡說冷熱,不是說淑寧的溫度,而是說熱水或冷水背後所代表的功效..
耶穌指責老弟家教會的工作就像暖水一樣,沒有功效.它不像希拉玻那那樣的熱泉水可以發揮治療的作用,也不像哥羅西那樣的冷泉水可以給人飲用解渴..
老弟家教會的工作就像當地那些帶有難聞氣味的溫水,一喝進口就反胃,想吐出來,毫無功效可言..
就正如大家有去茶餐廳,揭開餐牌,你會看到有熱飲,冷飲,熱奶茶,冷奶茶,熱咖啡,冷咖啡,熱檸樂,冷檸樂,你會不會看到有暖飲賣給你?有沒有暖咖啡?暖奶茶?暖檸樂?沒有的,為什麼?.
因為東西不冷不熱,喝下去總是覺得感覺一般般,特別今天天氣冷了,你當然想喝一杯熱飲,為身體補溫,或者夏天熱熱的,你就很想喝一杯冷飲,特別如果你有做運動打完球,.
你就很想消暑解渴,透過一杯冷飲,我還記得我以前中學時代,夏天打完籃球,跟同學,我們一定會做一件事,就是去七十一,知不知道買什麼?知不知道?你們這麼厲害,現在還有沒有?還有的?已經三十多四十年了,還有的?.
真的買斯諾冰飲的,因為你一喝上去,第一件事就是那些冷的感覺就會上腦,冷到頭殼都冷了,所謂的冷上腦,但你又會覺得真的很透心涼,不過當然你長大了就知道這樣是很傷身的,千萬不要做..
你會發覺其實耶穌基督責備老弟家教會,要麼你就做熱水,要麼你就做冷水,兩次打一樣你都會提供到功效,發揮到教會的功用,但你偏偏兩頭不到岸,你做暖水,你發揮不到教會應有的功能..
事實上老弟家這個地方比起鄰近的希拉波納和哥羅西,物質上都要富裕,教會建立在當地理論上都應該比其他地方的教會更有作為,但耶穌卻給了他一個負評或者甚至劣評,說他沒有功用,一無是處..
原因是當我們再看下文的時候,你會發覺很明顯他們沒有正確看待上帝賜給他們的那種付足.我們繼續看,17至18節,我想大家再讀一次,好嗎?.

$^{41}$第17至18節,預備起..
(教會禮儀).
你看到耶穌表示物質上的那種富足,物質上的一無所缺,似乎就是老弟家教會的問題所在..
特別經文提到三樣東西,金紙,衣服或者是眼藥,其實這三樣東西就是令老弟家覺得自豪,覺得富足的一些事物.事實上根據聖經學者巴萊的資料,老弟家教會是當時其中一個最富有的城市之一..
它的金融業經濟相當之發達,如何發達呢?當地設有所謂的票據交易所,意思是假如有些人要去老弟家旅行或者進行一些商貿交易的話,其實他不用帶太多現金,他只需要出示一張所謂的信用證明書..
就好像今天那些,你拿了旅行支票或者本票,你拿著這張東西,你就可以在老弟家當地的票據交易所提取現金,在當時其實是一個很先進的設施..
而另一方面,在當時主後的61年,老弟家曾經發生了一場大地震,整個城市遭到嚴重破壞,但是由於當地人真的非常富裕,他們完全有能力,有資源,獨力承擔災後的重建,所以他們不需要依賴羅馬政府對他們的經濟援助..
他們堅持我們自己用錢去重建這個城市,成為了當時一時的計劃.這件事其實也反映了老弟家人內心的自信或者自豪感,而老弟家教會大概也沾染了這種自負的態度..
所以你看到他們自誇的是什麼?我們是富足的,我們是發了財,我們已經富起來,我們也一無所缺,但是你看到在耶穌的眼中,老弟家教會,他形容老弟家教會是貧苦,可憐和窮困..
沒錯,物質上他們的確非常富裕,但是原來在靈性上他們真的窮得可憐,或者用以前很流行的說法就是老弟家教會窮得只剩下錢..
而另一方面,為什麼在經文中提到要買白衣呢?因為老弟家教會其實也是當時製衣業頗為繁榮的中心,當地以盛產一些優質的黑色羊毛來見稱,並且他們不單止盛產羊毛,他們還懂得用這些羊毛來製作一些嫁嫣物美的外衣,很深受當地人的愛戴和穿上..
所以老弟家人也會以自己擁有一個很好的製衣業來自居,但是耶穌卻說老弟家人你引以為傲的這件外衣,製造衣服的行業,其實沒有遮醜的功能,你表面上可以穿得很剛鮮,但是實際上耶穌形容你是赤身老體..
另一方面,老弟家人對他們的眼科保健很有自信,因為當地原來也設有一些頗有規模的醫藥中心,以盛產眼油來出名,知道為什麼眼油在當時這麼重要嗎?因為古代的道路不會鋪石屎,不會鋪蠟青,因為沒有東西鋪,走下去就會沙塵滾滾..
沙塵氧起的時候,風吹過的時候,或者馬車經過的時候,沙塵氧起,你的眼睛就會容易入沙,容易生病或者乾澀.所以眼油在當時基本上是一個經常需要的物品..
耶穌卻指責,沒錯,老弟家有眼藥,你們對眼保健得很好,但是你們在靈性上是盲的,看不到東西,看不清楚自己的屬靈真實狀況是怎樣..
等於我們來到這裡,我們會看到,你想像一下一個城市,如果無論金融,製衣業,製藥業都發展蓬勃,你怎樣看這個城市?是不是一個富裕城市?.
不要說昔日的標準,就算放在今天的標準,都必然是一個令人感到自豪,令人感到自信的一個富裕城市..
而老弟家教會似乎也感染了這個城市給他的那種優越感..
你想像一下,特別當我們之前看了其他教會的狀況,當其他教會都面臨困乏,面臨逼迫,奉獻又缺少的時候,總是遇到不同的壓力,特別是財政的壓力的時候,很多的士工都沒有辦法推展,但是老弟家教會怎樣?.
我們沒有這些問題,我們富足到一個地步,我們想做什麼,就可以馬上實踐,沒有什麼是我們搞不定的,這種自信,滿有把握的心態,也變成了一種我們說的屬靈上的驕傲..
而到了教會,他們開始覺得我們富裕,我們富足,就等於我們很棒,等於我們已經很安穩,他們覺得富足就是他們最大的保障,完全察覺不到其實富足反而是他們最大的問題,最大的危機..
當然,來到這裡我們必須說明,聖經不是說人不能生活富足,特別是在舊約裡面,在猶太的傳統裡面,他們都會認為富足的生活確實是上帝賜福的一個記號,遵守上帝的吩咐,上帝就賜下豐足的生活,來顯明那個人行在上帝的旨意裡面,.
上帝也喜悅你,我們認識的亞伯拉罕,很多的以色列人,譬如大衛,他們的生活都是富足的,不過聖經告訴我們,這個不是必然的法則,聖經也記載了很多上帝喜悅的人,他們的物質生活是困乏的,最佳的例子就是耶穌和他的門徒..
特別是在福音書裡面,耶穌對於地上的富足,多數都是採取一個有保留,甚至乎是負面的態度,特別我們記得耶穌經常說什麼?越有錢的人,下一句呢,就越難進天國..
原因不是富足本身,而是富足會不知不覺成為了人,依靠甚至服侍的對象和焦點,偷偷地佔據了我們的心,取締了耶穌基督在我們身上的地位..
所以你看到,當老弟家教會自誇自己富足一無所缺的時候,耶穌就直接指出他們的屬靈真實狀況是怎樣,就像這裡說的,是困苦,可憐,貧窮,乞眼,赤身..
於是耶穌就勸他們,老弟家教會要向耶穌買三樣東西,經文說第一個就是火煉的金紙,意思就是那些經得起任何試煉考驗的天上財富..
如果對應啟示錄當時寫作的背景,火煉的金紙的意思就是即使面對逼迫,面對苦難,仍然持守真理的這種純存的信仰生命..
這種富足對老弟家的信徒來說是陌生的,但耶穌說這才是真正的富足..
第二樣耶穌要求老弟教會向他們買的是白衣,要買一件白衣穿上,意思是當審判來到的時候,老弟家的信徒他們在地上,他們所依靠所誇耀的那件羊毛衣服,其實沒有辦法替他們遮醜了..
因為當主來臨的時候,當他審判全地的時候,每一個人都要在主的面前闖開,打開這件衣服看看裡面是什麼..
而那些穿上耶穌所說的白衣的人,即是那些和耶穌在屬靈爭戰裡面一同獲得勝利的信徒,他們是沒有事的,他們不會赤身露體,他們不會露出羞恥,因為有耶穌這件德性的白衣遮蓋了他們..
他們不單止不會羞愧,反而會因為得到耶穌認同的緣故,他們可以分享到耶穌那份榮耀..
最後一樣要買的就是眼藥,這個意思比較清楚簡單,就是要他們看清楚,老弟家他們信徒他們需要看清楚現在他們所誇耀所倚重的一切,.
其實反而是帶著他們步向滅亡,他們要看清楚自己在上帝面前的光景是怎樣,唯有這樣才意識到自己其實是需要回轉悔改..
正所謂愛之深自然責之切,雖然耶穌嚴厲責備老弟家教會,但不代表耶穌不愛這間教會,你會看到耶穌沒有放棄過他們,特別在最後一個段落,19-22節,你看到耶穌不單止指正他們的問題,耶穌其實是內心在等他們回轉,等他們悔改..
特別在這個段落你會看到經文呈現了我們說的三個圖畫,三個景象,這三個景象其實反映著耶穌到底有多愛老弟家教會..
第一個影像就是看到耶穌站在門外叩門,我們經常會覺得這段經文是對那些未信者說的,.
譬如有時報道會對未信者說,現在耶穌站在你的心門外叩門,你肯不肯開放自己等耶穌進來,我們會這樣理解這一節的經文,但你看原本的意思,原來不是這樣,這句話不是給未信者,而是給信徒,因為把耶穌拒諸門外的是老弟家教會,.
明明錯的是他們,但反而是誰低聲下氣,好像在求他們,明明應該是老弟家教會向耶穌赴京請罪,但反而是耶穌放下身段,主動勸他們悔改,.

$^{81}$你看到耶穌的語氣不是很嚴肅很威嚇,很開門,再不開門,我會把門爆開,耶穌不是這樣,耶穌充滿了關懷和愛意,我在你的門外叩你門,你聽到嗎?.
如果你聽到,你就應,你就開門,你看到耶穌把他的姿態放到多低,明明他是對的一方,耶穌對老弟家教會的愛,不單止驅使他,我們說主動尋找迷洋,主動來敲老弟家教會的門,並且願意耐心等候教會回應他,等他自己情願開門,.
第二幅圖畫就是主要和他一同坐直,坐直的意思是吃晚飯,原來在當時的文化,每個人都像我們現在這樣,應該吃三餐,我們也有三餐,早餐,午餐,晚餐,.
早餐你會吃得豐富一點還是不豐富?豐富一點,早餐應該吃最飽,午餐就吃少一點,晚餐因為要睡覺,所以吃最少,但他們反過來,早餐可能他們最簡單,一塊乾的麵包,有時候可能沾點橄欖油,沾點酒,吃完就走了,因為要工作,至於午餐呢?.
他們不會像我們這樣,如果你近家就回家吃,他們多數不會回家吃午餐,而是預先一早包好午餐的份量,然後在戶外吃,至於晚餐呢?晚餐反而是他們全日最重視的一餐,經過一日勞碌,是時候和家人享受一個歡聚,歡愉的時刻,所以晚餐的時間是用得最久的,.
耶穌說他要進來,不是進來省老弟家教會一餐,而是和他們一起吃晚餐,並且不是吃完就立刻走,而是要留下來和他們團契相交,意思其實是什麼?修復和修補那個破爛了的關係,東方人會比較明白這個意境,.
很多時候當一段人際關係出了問題的時候,如果願意一起坐下來吃一頓飯的時候,我們說好像擺和頭酒,一起願意吃一頓飯其實就是一起想修復關係,想和好,最後一個景象就是不單止吃一餐,而是怎麼樣?.
耶穌說他要賜老弟家教會同坐在寶座上,我不知道當你想到這句的時候你有什麼反應,剛才才罵到老弟家教會飛起,你覺得按照老弟家教會現時的屬靈情況,你能不能想像到他將來可以坐在耶穌旁邊?.
其實你想像不到,落差太大了,你很難想像他們將來可以和耶穌同坐在寶座上,分享上帝給他們的榮耀和權柄,但是你看到耶穌的愛驅使他,即使一間破爛的教會,耶穌還看好他們,他覺得老弟家教會將來是有前景的,.
所以他才苦口婆心地勸他們回頭,耶穌盼望有一天他們醒過來,真的會回轉,可以和他們一起走上德性的道路,一起同坐德性榮耀的寶座..
弟姐們,經文大概講解到這裡,有兩個分享,我想很值得我們互相提醒,你留意這裡,耶穌是鄭重提醒教會,提醒信徒,要向他買東西,向他買金紙,買外衣,買眼藥,首先就是說,耶穌所購買的才是正貨,.
等於你去藥房買東西,你會找一間有個R字的,你就知道是正貨,有個正字的,也不是正貨,有個正字的,只是不能配藥,有個R字的就更高一層,因為有藥劑師駐場..
耶穌說他賣給我們的才是正貨,換句話來說,老弟家教會手上所持有的那些富足是什麼?假的,仿製品,A貨,無論他們在地上怎樣誇耀他們的金融,他們的制衣,他們的制約,他們的教會誇耀他們的戶口有多多錢,.
耶穌說什麼?將來在天上全部不計算,全部不計算,因為是假的,就正如我們,譬如當你升上小學的時候,讀到高小的時候,你還會不會計算你幼稚園的成績?你不會吧,你不會再誇耀你幼稚園拿過操行獎,拿過最佳模範生獎,不會吧,你升到高中,你都不會誇耀你小學考第一的了,.
你相信出來工作,你不會再誇耀,你看看我大學多威風,拿過多少一級榮譽,你不會再誇耀,同樣你去到天家,你不會誇耀,你看看我地上多麼富有,你不會,因為你知道耶穌都不計算這些東西,耶穌不計算這些東西,耶穌關心的不是我們地上的財富,不是地上的成績表,而是你天上的成績單是怎樣的,.
耶穌不關心我們地上有多少財富,有多富足,而是關心我們能否有智慧,努力地將地上的財富轉到天上的戶口,這才是耶穌關心的,我們所擁有的財富,弟兄姊妹,我們需要檢視,是不是耶穌認可,是不是耶穌眼中的證貨,是不是真的能夠轉賬到天上,.
另一方面,你見到耶穌提醒我們,我們要向祂買東西,買東西即是要付錢,即是要付代價,耶穌表示跟隨祂的人都要向祂付代價,.
當然我們很明白,救恩是白白無條件的賜給我們,但是每一個堅持跟隨耶穌的人都要隨時為主犧牲,隨時為主寫記,付這個代價,不知道大家有沒有看過潘博華有一本著作叫做追隨基督,或者有聽過,.
他的英文很有意思,叫做The Cause of Discipleship,即是說做門徒有代價,是有代價,裡面特別提到,假如信耶穌這個字很容易讓我們誤解是不需要付代價的話,我想我們需要互相提醒,.
他講到很極致,其實基督徒就是耶穌的跟隨者,跟著耶穌去哪裡?就是去受死,當耶穌呼召一個人的時候,他是呼召他們為耶穌死,為耶穌犧牲寫記,正如耶穌也說,若有人要跟從我,就當寫記天天背起他的十字架來跟從我,.
他在說的就是,我們不只是缺了自信耶穌,然後就回到世界裡面,找我們的財富,耶穌說你每一天都要為你的信仰,為上帝付上成長的代價,耶穌要求我們讓他居首位,而不是只將耶穌,將真理當是你生活的其中一部分,.
他提醒我們要向他買火煉的金紙,而不是向這個世界買能扭壞的財富,最後當我們看完啟示佬這七間教會的時候,我們會發覺教會是不是有很多問題?.
教會其實不是我們在現在家司看到的示範單位,大家有沒有以前拿現在家司的那本簡介,方方地現在還有沒有派?現在可能網上,上世紀好像很遠,十多年前還是很流行的,一本一本的,你去到現在就會拿一本然後揭一揭,你會發覺裡面每一頁都怎樣?.
飯廳也好,廚房也好,洗手間也好,每一頁都乾乾淨淨的,每一頁都是整整齊齊的,你會發覺只要我的房子是這樣,我就可以隨時叫人來坐坐,隨時邀請人來,不過抱歉教會不是這樣,教會不是現在家司的示範單位,教會也不是百貨公司的展示的櫥窗,.
教會其實是你跟我日常起居生活的客飯廳,甚至乎可能是洗手間,因為裡面住的是罪人,所以就會怎樣?就會污穢,就會髒,甚至乎會很雜亂,不做統計,有沒有人說你的衣服是到處丟的,在客廳,當沒有人來的時候,桌面又有吃完飯送的即未抹,.
地板又鋪滿頭髮,鋪滿塵,廁所不說了,你會發覺總會有這些混亂的東西,因為沒有人來看,你自己忍受得到,你住得越久,發覺這些積雜的問題就會越來越多,除非你定期的打掃,徹底的清潔乾淨,不過人總是有惰性和罪性,.
你明明看到問題,明明看到髒東西,你都會怎樣?你不搞它,不想搞它,因為很煩,很忙的,遲些再搞,等我搞定我手頭上的東西,有空再搞好它,又或者說因為人的罪性,會將問題淡化了,或者說正常化了,我都不是很差,你看看隔壁那個比我還差,我們向下看齊的,總是喜歡,.
事實上,這些不單單是啟示錄那七間教會所呈現的問題,其實也代表了普世教會,包括我們,會出現的狀況,但是好消息是,耶穌是教會的主,祂比我們任何一個人更加關心神家的事情,教會出現了混亂,出現了問題,走在一些不合上帝的旨意上,.
耶穌不會不理,耶穌一定會出手,來執整教會,執整我們每一個人的生命,所以有人說得很好,對一個人最大的懲罰,其實不是嚴厲責備他,而是什麼呢?.
對,不理他,放縱他,任由他繼續走錯,讓他自食其果,甚至乎回不了身,這個才是最大的懲罰,但是耶穌沒有這樣打算對待老弟家教會,也沒有這樣打算對待我們,因為我們都是祂所愛的,祂用重價將我們買回來,祂也是我們挽回的對手,而祂也是會挽回我們,.
正如耶穌所說,祂所愛的,祂一定會管教,一定會責備,所以弟兄姊妹盼望這個系列的宣講,讓我們重新認真看待耶穌對我們的提醒,甚至祂對我們的責備和管教,因為是出於愛,祂不想我們走錯路,.
希望我們珍惜上帝對我們這份不離不棄的愛和恩典,也求主監察我們,就像祂監察這七間教會一樣,拿走上帝不喜悅的心思意念,重新眺望我們對主的愛,對主的熱心,以致我們真的願意順服,配合耶穌基督的大靈,我們同心圖告..
並且將來當我們真的走迷,行差踏錯的時候,你也要拯救我們,主啊,因為你愛我們,希望我們真的得著你所認定的真正的豐盛和富足,主啊,就正如你對待老弟家教會一樣,你會因為我們行差踏錯而責備我們,管教我們,但這一切都是為了我們的好處,為了我們的益處而做的..
盼望主啊,當我們面對你的時候,我們真的能夠謙卑,柔和起來,我們的心真的能夠被你軟化,願意真的去,如經文所說,有耳可聽的,應當聽..
盼望主啊,你親自藉著這一系列的宣講,對教會,對我們每一個說話,讓我們真的願意為了信仰,為了上帝,為了真理的緣故,我們負上代價,願意我們所擁有的那些地上的財富,都能夠因著聖靈的帶領,轉化為屬天的寶藏,.
以致我們真的可以,到主在內的時候,我們真的可以身穿那件德性的帕袍,和耶穌一同的坐直,分享耶穌那份德性的榮耀,為此我們恭敬將每一位交託在你的恩手,願上帝你繼續親自藉著你的話語對我們說話,讓我們行事為人,能夠與所蒙的恩相稱..
願上我們的禱告,奉基督耶穌的德性的名字來祈求..
阿們..
\newpage



\section{箴言 8:12-31}
\label{sec:wYKTi6uXSJI}
\textbf{2025年12月21日聖誕節主日崇拜直播｜高銘謙牧師｜愛智慧｜箴言8\_12-31}
\newline
\newline
連結: \href{https://youtube.com/watch?v=wYKTi6uXSJI}{\texttt{https://youtube.com/watch?v=wYKTi6uXSJI}} ~~~~ 語音日期: 2025-12-21
\newline
\newline
\hyperref[sec:1zYtfRDiK4g]{\small{< < < PREV SERMON < < <}}
~
\hyperref[sec:index_chronic]{\small{[返順時目]}}
~
\hyperref[sec:yqcP_iDkMnA]{\small{> > > NEXT SERMON > > >}}
\newline
\newline
箴言 8:12-31
\newline
\begin{longtable}{cl}
\hline
\hline
章節 & 經文 (和合本修訂版)\\
\hline
8:12 & \begin{tabularx}{0.7\textwidth}{X} 我－智慧以靈巧為居所，又尋得知識和智謀。 \end{tabularx} \\ \\ \relax
8:13 & \begin{tabularx}{0.7\textwidth}{X} 敬畏耶和華就是恨惡邪惡；我恨惡驕傲、狂妄、惡道，和乖謬的口。 \end{tabularx} \\ \\ \relax
8:14 & \begin{tabularx}{0.7\textwidth}{X} 我有策略和健全的知識，我聰明，又有能力。 \end{tabularx} \\ \\ \relax
8:15 & \begin{tabularx}{0.7\textwidth}{X} 君王藉我治國，王子藉我定公平， \end{tabularx} \\ \\ \relax
8:16 & \begin{tabularx}{0.7\textwidth}{X} 王公貴族，所有公義的審判官，都藉我掌權。 \end{tabularx} \\ \\ \relax
8:17 & \begin{tabularx}{0.7\textwidth}{X} 愛我的，我也愛他，懇切尋求我的，必尋見。 \end{tabularx} \\ \\ \relax
8:18 & \begin{tabularx}{0.7\textwidth}{X} 財富和尊榮在我，恆久的財寶和繁榮也在我。 \end{tabularx} \\ \\ \relax
8:19 & \begin{tabularx}{0.7\textwidth}{X} 我的果實勝過金子，強如純金，我的出產超乎上選的銀子。 \end{tabularx} \\ \\ \relax
8:20 & \begin{tabularx}{0.7\textwidth}{X} 我在公義的道上走，在公平的路中行， \end{tabularx} \\ \\ \relax
8:21 & \begin{tabularx}{0.7\textwidth}{X} 使愛我的承受財產，充滿他們的庫房。 \end{tabularx} \\ \\ \relax
8:22 & \begin{tabularx}{0.7\textwidth}{X} 「耶和華在造化的起頭，在太初創造萬物之先，就有了我。 \end{tabularx} \\ \\ \relax
8:23 & \begin{tabularx}{0.7\textwidth}{X} 從亙古，從太初，未有大地以前，我已被立。 \end{tabularx} \\ \\ \relax
8:24 & \begin{tabularx}{0.7\textwidth}{X} 沒有深淵，沒有大水的泉源，我已出生。 \end{tabularx} \\ \\ \relax
8:25 & \begin{tabularx}{0.7\textwidth}{X} 大山未曾奠定，小山未有之先，我已出生。 \end{tabularx} \\ \\ \relax
8:26 & \begin{tabularx}{0.7\textwidth}{X} 那時，他還沒有創造大地和田野，並世上頭一撮塵土。 \end{tabularx} \\ \\ \relax
8:27 & \begin{tabularx}{0.7\textwidth}{X} 他立高天，我在那裡，他在淵面的周圍劃出圓圈， \end{tabularx} \\ \\ \relax
8:28 & \begin{tabularx}{0.7\textwidth}{X} 上使穹蒼堅硬，下使淵源穩固， \end{tabularx} \\ \\ \relax
8:29 & \begin{tabularx}{0.7\textwidth}{X} 為滄海定出範圍，使水不越過界限，奠定大地的根基。 \end{tabularx} \\ \\ \relax
8:30 & \begin{tabularx}{0.7\textwidth}{X} 那時，我在他旁邊為工程師，天天充滿喜樂，時時在他面前歡笑， \end{tabularx} \\ \\ \relax
8:31 & \begin{tabularx}{0.7\textwidth}{X} 在他的全地歡笑，喜愛住在人世間。 \end{tabularx} \\ \\
[1ex]
\hline
\hline
\end{longtable}
$^{1}$今日我們選讀的經文是《舊約聖經》的針言.
《舊約聖經》針言第八章十二至三十一節.
請聽我讀出.
我自謂以靈命為居所.
又尋得知識和謀略.
敬畏耶和華.
在乎恨惡邪惡.
那驕傲狂妄並惡毒.
以及乖貓的口都為我所恨惡.
我有謀略和真知識.
我乃聰明我有能力.
帝王藉我坐國位.
君王藉我定公平.
王子和首領.
世上一切的審判官.
都是藉我掌權.
愛我的我也愛他.
懇切尋求我的必尋得見.
豐富尊榮在我.
恨狗的才並公義也在我.
我的果實勝過黃金.
強如精金.
我的出產超乎高銀.
我在公義的道上走.
在公平的路中行.
使愛我的承受貨才.
並充滿他們的付付.
在耶和華造化的起頭.
我在太初創造萬物之先就有了我.
從亙古從太初.
未有世界以前我已被立.
沒有深淵沒有大水的泉源.
我已伸出.
大山未曾奠定.
小山未有之先.
我已伸出.
耶和華還沒有創造大地和田野.
並世上的土質.
我已伸出.
他立高天.

$^{41}$我在那裡.
他在冤滅的周圍.
劃出圓圈.
上使窮滄堅硬.
下使深淵穩固.
為滄海訂出界限.
使水不越過他的命令.
立定大地的根基.
那時我在他那裡為公司.
日日為他所喜愛.
上上在他面前踴躍.
踴躍在他為人預備可住的地.
也喜悅住在世人之間.
各位頂子們早安.
和大家一起度過聖誕節的崇拜.
聖誕節的崇拜應該要和大家分享新約.
不過我心想可能舊約這一段也可能是適合.
因為近來也有機會教聖經神學的時候.
原來愛智慧這東西.
也是我女兒的英文名字.
Sophia這個英文名字.
就是來自希臘文.
希臘文的意思就是智慧的意思.
我想問這裡在座有沒有姐妹.
或者可能是有弟兄.
是叫做Sophia的.
有沒有.
值得一想的.
你可以改這個名字.
因為有智慧.
愛智慧這個主題也等於愛耶穌.
耶穌就是神的智慧.
這是哥林多前書第一章二十多節.
清晰地提到.
智慧就是耶穌.
耶穌就是智慧的化身.
我們祈求主慶祝耶穌降世為人.
並且道成肉身.
降在馬槽的節日裡.
我們認識什麼叫做愛智慧.

$^{81}$我們低頭祈禱.
慈悲仁愛的天父.
多謝你讓我們今早一同來到你面前.
很想認識這個智慧.
主你就是智慧的化身.
讓我們不單止能夠認識你.
相信我們每一個都能夠成為一個智者.
活在我們的人生.
活出你那份愛.
活出真正的智慧.
幫助孩子講解你的說話清楚明白.
我們一同去到謙卑領受.
奉耶穌基督的名字祈禱.
阿們.
今天這段經文就是來自箴言剛才讀到那段經文.
剛才讀到這段經文.
先看一下.
剛才讀到這段經文.
可以分為三個段落.
第一個段落就是由十二到第十六節的經文.
這段段落就是告訴我們.
愛智慧的第一個好處到底是什麼.
並且清晰地定義什麼叫做有智慧.
智慧我想大家都知道.
在箴言裡面敬威耶和華是智慧的開端.
智慧是和敬威耶和華有關的.
這是第一點我們會留意.
第二點.
他用愛字告訴我們.
愛智慧帶給你第二個好處.
就是給你有尊榮這件事情.
第三點.
他的話題搖身一變.
告訴我們原來智慧不是一個被造物.
智慧是和創造主一起創造這個世界.
後來我們也知道在歌羅西書裡面.
告訴我們這個智慧其實就是耶穌.
是創造天地萬物的那一位.
所以第三點.
跟我們說智慧和創造秩序之間.

$^{121}$他們的關係到底是怎麼樣.
透過這樣我們用另外一個方式.
去默想在這個聖誕節裡面.
怎樣可以愛智慧.
所以首先我們去到第一點.
第一點你會發現.
智慧是用第一生的我去說的.
也就是說智慧不是一個死物.
不是一個概念.
不是一門學說.
也不是一個學位.
智慧是一個有位格的個體.
他會跟你說話.
他用我去進行稱呼.
告訴你我是怎樣.
並且你如果愛他的話.
他會帶給你會是怎樣.
所以第一個.
他就說我智慧是以靈命為居所.
有尋得知識和謀略.
然後十三節是清晰地告訴我們.
敬畏耶和華究竟是怎樣.
因為我們剛才讀到九章第十節.
也是我自己那間中學的校訓.
敬畏耶和華是智慧的開端.
但是我們不是很清楚.
什麼叫做敬畏耶和華.
這裡就非常具體告訴你.
敬畏耶和華就一定有這個表現.
這個表現就是恨惡.
恨惡這個字就在十三節那裡.
總共出現兩次這麼多.
第一次就叫做恨惡邪惡.
第二次是有四樣東西是你要恨惡的.
第一驕傲.
第二狂妄.
接著惡道.
然後就是乖妙的口.
看上去有數得計.
一數就是有五樣東西這麼多.

$^{161}$你是要恨惡的.
但其實應該是四樣東西.
因為第一個恨惡邪惡.
那個邪惡是一個籠統的類別.
那個類別包含了四樣東西.
驕傲狂妄.
惡道和乖妙的口.
讓我逐一去解釋.
特別驕傲和狂妄一定是打孖的解釋.
因為驕傲那個字的字根.
是解作高這樣的意思.
這個字在以賽亞書的點擊率也挺高的.
特別你會看到.
以賽亞書的第六章.
也就是說以賽亞先知蒙召的那一章.
他蒙召的時候就見到一個異象.
就說我看見主坐在高高的寶座上.
你聽到和學本的做高高.
原文真的有兩個高.
那裡第一個高是主動的.
第二個高是避抬高.
也就是被動的.
無論主動還是被動.
高這個位置是專屬給耶和華的.
以賽亞書的第二章也告訴我們.
很多人很想成為高.
包括那些狂妄人.
包括那些惡人.
很想成為高.
但如果你想成為高的話.
上帝就將你降低了.
凡至高上帝必降低.
自卑就會被升高.
這是耶穌基督在不同的福音書裡.
都有記載的道理.
所以做人千萬不要高.
聖高那些靈界的客.
做人千萬不要高.
但是告訴我們.
原來敬畏耶和華的人.

$^{201}$你是恨惡什麼.
恨惡高.
就是說我不喜歡自己高.
我不會朝朝夕夕.
高也可以有很多表現.
第一 財富高.
希望大家未必有.
不過可能是地位高.
可能是名氣高.
但通通都是他恨惡.
恨惡的意思.
不是代表你有錢就有問題.
而是說那種態度.
那種態度就好像這個世界.
沒有你不行.
所以恨惡這件事情.
就是敬畏耶和華的一個表現.
第二 第三和第四.
就是邪惡的惡道.
第四就是乖妙的口.
特別道路真的不是指屯門公路.
而是指人生道德的取向.
一個叫善道.
一個叫惡道的意思.
而你真的看看嘴巴說什麼.
就知道你是不是敬畏耶和華的人.
不敬畏耶和華的人.
是有乖妙的口.
乖妙的口是什麼呢.
其實就是那些.
就是七成真 三成假.
說些不說些.
說六成真的東西.
搭棚衝就是四成假.
類似這樣.
說些不說些.
去到誤導別人.
又或者保持自己良好的形象.
目的的底下.
所以他說了一些.

$^{241}$似是而非的說話.
類似這些圓滑的舌頭.
我們就叫這一種.
叫做乖妙的口.
如果你離開這些事情的時候.
你就會發現在第十四節那裡.
你就會發現同樣地.
都是有四樣東西.
是智慧帶給你的.
那四樣東西.
包括第一是謀略.
第二是知識.
第三是聰明.
第四是能力.
頭三樣全部鎖上門都有的.
最後那樣是能力.
你可以解作財力.
可以解作軍事勢力.
可以解作智力.
都可以.
能力就是這樣的表現.
所以你會發現.
智慧.
Wisdom.
就和我們IQ高還是低呢.
是沒有關係的.
IQ高不代表有Wisdom.
IQ低是可以有智慧的.
智慧是來自敬畏耶和華.
所以你會發現有些IQ高.
他一不敬畏耶和華就等於什麼.
等於沒有智慧.
這個正正應驗了.
哥林多前書保羅所提到的那段經文.
那段經文就是說.
猶太人是要神蹟.
希臘人是求智慧.
我們是卻全十字架的基督.
在猶太人是半腳式.
接著下一句就是.

$^{281}$在愛邦人是愚蠢.
就是說世人覺得這個智慧是愚蠢.
但是這個智慧.
但在那蒙召的無論是猶太人希臘人.
基督總是神的能力.
神的智慧.
在世人覺得是蠢.
但是在真正有智慧人的眼界.
是一個如珠如寶的構成.
所以第一點我們會發現.
如果是這樣的時候.
就會給你一個真正的智慧.
話說幾個月前.
我在八月份.
八月份我就答應了一間教會.
逢禮拜二.
有三個連續的禮拜二晚上.
就舉行一個Zoom的課程.
就是有一個線上的課程.
所以在家裡進行.
有五十多個弟兄姊妹報名.
晚上八時就講到九時半.
已經進行了第一個禮拜.
接著就到第二個禮拜.
第二個禮拜的禮拜二.
那一天真的不知道為什麼.
我是完全忘記了那一晚要教書.
那一晚還約了一家人一起吃一頓飯.
為什麼呢.
因為我女兒正要升Poly一年級.
所以想著進入宿舍之前.
吃一頓最後晚餐.
就叫老婆煮.
我老婆煮是很好吃的.
總之我老婆煮就是很好吃的.
一家人一家三口.
和那兩個女人一起吃.
七點兩個字開飯.
吃得很開心.
很開心.

$^{321}$接著就吃到七點九.
我的手機就擺到充電.
並且調了震機.
七點九就進去切水果.
切了幾顆橙出來.
就說一起吃.
吃到八點一個字.
然後就說我洗碗.
接著就進去洗碗.
一洗就洗到八點九個字.
倒垃圾.
洗完手.
想著檢查一下手機有什麼訊息.
就發現有幾十個追魂call.
每個人都在尋找你.
然後傳道人就說.
其實我們都等了半個小時.
到八點半zoom link都close了.
有五十多個在等.
接著堂主任都出動了.
就在找你.
都找不到.
我覺得很面目.
比爵利你有什麼感受.
其實我也覺得很面目.
接著我也是很抱歉.
因為我恐慌.
真的覺得原來那間教會.
覺得我不重視這件事之類.
很不開心.
那天就很想找弄倦.
不過你始終是成人.
都要和人交代.
看看怎樣可以和人補課之類.
的確是補課了.
那一晚就很面目.
就要傳WhatsApp給堂主任.
傳的時候.
突然之間就有聲音和你說話.
高明軒 你這樣做實在太面目了.

$^{361}$不如你怎樣都要說一些話.
為的目的就是要保持良好的形象.
那聲音就教你怎樣說話.
不如這樣.
你忘記了也是事實.
不過你會不會捏造一些藉口.
告訴你都是情有可原.
這樣會好一點.
接著我又不理那聲音.
那聲音繼續說話.
就說不如這樣吧.
不是叫你說謊的.
聲音很狡猾.
就跟你說你現在是不是有鼻水.
我說是 我有鼻水.
不如這樣吧.
你跟別人說你有鼻水.
然後就說你忘記了.
那個因果關他自己串.
你不是說謊的.
一卵嘴這樣說.
很多這些事情.
剛才我們提到那把聲音.
那個邏輯.
那個說話的方式.
我們叫這一種.
叫乖妙的喉.
這不是上帝的自慰.
後來我就膽粗粗.
不理那把聲音.
就如實報道.
如實報道是不是很蠢.
在世俗的眼光就是蠢.
你為什麼要告訴別人.
你跟女兒吃最後晚餐.
為什麼要告訴別人.
你真的忘記了.
你是蠢.
是蠢 如實報道.
但是如果是這樣.

$^{401}$我是站在真相 真理裡面.
這個才是有自慰.
世俗的眼光就覺得什麼.
蠢.
世俗覺得什麼叫聰明.
就是剛才那個.
乖妙的喉才叫聰明.
是顛倒了那個邏輯.
和不同的價值觀和看法.
正正就是這件事.
是保羅用了.
保羅在哥林多前書用了.
希臘那些的自慰.
看上去是聰明.
但是耶穌基督十字架的自慰是什麼.
簡直蠢到+01.
但是這個就是神的自慰.
那天我就如實地.
和那個堂主任報道.
不好意思 我真是大頭蝦.
我完全忘記了.
因為我掛念切剷.
之類那些.
完全是這樣.
然後報道了之後.
我都真的很想運動卷.
但是.
我覺得生命裡面有個喜樂.
我知道.
敬畏耶和華的這個態度.
是一定要高過.
怕別人怎麼看我們.
是不是.
有很多時候.
我怕別人怎麼看我.
是大過我們怕上帝.
如果是這樣.
我們就採用了.
別人的觀點和看法.
而那件事可能就是.

$^{441}$活出一個不怕上帝的概念.
所以第一點.
我們會留意.
最終整件事.
最多補課.
其實沒什麼.
沒什麼.
但是別人都體諒.
不知是客氣還是不客氣.
起碼他不在我面前說.
不過anyway.
客氣的就說.
是不是高牧師.
你說了講究影音培靈會.
辛苦得很.
忘記了.
他都幫我解釋.
不過我也覺得有點客氣.
我就說不是.
不是.
其實真的忘記了.
所以第一點.
就是要敬畏耶和華.
在乎.
恨惡.
那個驕傲.
和狂妄的事情.
難怪林肯總統.
有句話.
這句話很出名.
就是不要問.
上帝是否站在我們這一邊.
要問.
我們是否站在.
上帝的那一邊.
如果要站在上帝的那一邊.
就一定要配合上帝的智慧.
上帝的做法.
上帝的看法.
原來我們不是.

$^{481}$我們倒轉.
我們在想.
上帝是否站在我這一邊.
配合我的agenda.
我的題目.
我的做法.
這是第一點.
接著第二點.
17到21節.
你會看到17節.
是用愛我的做開始.
21節是用愛我的做結束.
看到21節都有嗎.
所以愛是首尾呼應.
成為一個主題.
告訴你.
愛他.
是你會有什麼回饋給你的.
第一.
就是愛.
愛我的.
就是愛的字位.
得來的回應.
就是我也愛他.
是不是很簡單的邏輯報應.
你做好事就有好報.
你愛他.
他就愛你.
你想你老婆愛你.
很簡單.
你愛你老婆.
你就一定有報應.
那個報應就是.
你老婆愛你.
很簡單.
你尋求這個智慧.
你就必能夠尋得見.
這是不是耶穌說過的.
你叩門.
就必為你開門.

$^{521}$尋找就要尋見.
其實你要知道.
這是新約.
引用舊約.
即是耶穌很有意識.
知道舊約所說的智慧.
就是主.
去到這裡.
我們發現一個簡單的報應觀.
難怪有一位學者.
在1993年出版了一份國際期刊.
去描述真言的智慧的時候.
那個報應觀.
他用了一個比喻去描述.
那個比喻就叫做回力鏢效應.
大家有沒有玩過回力鏢呢.
你一飛出去的時候.
那個回力鏢一定會回來.
有時回力鏢就中你幾刃.
這個叫回力鏢.
回力鏢的效應就是.
在創造秩序下.
智慧是真言的角色.
善就有善報.
惡就有惡報.
你背讀下一句.
若然未報.
事實未到.
即是代表遲早都報.
類似是這樣.
這是一個簡單的.
我們叫做智慧的回力鏢效應.
無論怎樣也好.
你如果真的做了一件事.
在你的人生裡.
而你又活在日光之下.
你又活在創造秩序當中的時候.
而這個創造秩序又是滿有神的智慧的時候.
智慧其中一個法則就是回力鏢的效應.
你就變聖所有的事情.

$^{561}$很多時候都會有一個報應.
會臨到你的身上.
你溫柔對你老婆.
愛心多一點.
你一定有報應.
那個報應就是.
你老婆會愛你的.
你細心養育你的子女.
做一個好的榜樣.
她不差得去哪裡.
反觀.
若然你高度借貸.
洗腳不抹腳.
是一定有一個報應臨到.
一定要摸就是破產.
要摸就是家破人亡.
這個不能談的.
無論你是敬畏主還是不敬畏主.
無論你今早有沒有靈修.
這件事都一定會回頭.
我認識一個人.
他有份愛情.
他被人看到.
跟他聊天.
聊天的時候.
他說牧師.
其實我都看到旁邊那位弟兄也有份愛情.
不過那位弟兄蠢.
為何蠢呢.
因為他跟那個小三.
出街的時候沒有戴戴戴戴帽和口罩.
所以就被人看到.
他看到這個案件.
他不產生敬畏.
他不產生警惕.
他就說當初我有份愛情.
我看到這個弟兄有份愛情.
這樣的份愛情.
我就聰明一點.
他每逢出街的時候都怎樣.

$^{601}$戴戴戴帽和口罩.
變成了.
就是希望萬無一失.
誰不知這個迴力鏢都怎樣.
一定會回來.
一定會被人撞破.
就是吃糖水的時候.
你都要脫下口罩.
就跟他一起吃的時候.
就被人看到了.
就拍了一張照片.
就變成了整件事.
就是避無可避.
類似是這樣.
做那件事.
每一樣事.
是有一個後果.
那個後果會回頭.
有時候是上帝的恩典.
希望你能夠悔改.
所以那個後果沒那麼快回頭.
希望你及早悔改.
以致那個後果不會淋在你的身上.
但有很多時候.
我們每一次又一次.
一次又一次做那件事.
然後我們忘記了.
原來那個智慧法則.
在這個創造秩序當中.
那個迴力鏢一定會回來.
我前陣子在Netflix.
看了一個沒有什麼人看的韓片.
一聽名字就不想看.
叫做《陽光先生》.
這個韓片.
有沒有人看的?.
我們不看的.
這個是說120年前的韓國.
當時有奴隸制度.
有一個主角是9歲的小朋友.

$^{641}$他是一個奴隸生出來的小朋友.
他父母都是奴隸.
因為做錯了一件事.
就被他主人打死了他父母.
最後這個9歲的小朋友逃走.
因為很多人追殺他.
然後就去到一個特定的地方.
就遇到一個美國的傳教士.
就上了一艘船.
帶了他去三藩市.
隨後9歲的小朋友就在那裡長大.
長大到20年之後.
搖身一變.
他就成為一個美軍.
美國的軍隊的將領.
然後被服役被指派.
派去哪裡呢?.
就是派去韓國.
去到進行服役.
派去韓國.
你知不知道他想做什麼?.
就很想找回那隊的主人.
就幫他去到報仇.
20年後搖身一變.
成為一個大約30歲的壯年男丁.
還是美軍的將領.
就踩進去那個主人的家裡.
那個主人都認不出他.
後來有些特徵認出他的時候.
那個主人嚇得臉都青了.
類似是這樣.
20年那個報應都回來了.
所以第二點就是想和大家說這件事情.
好了,接著最後.
最後就是智慧和創造秩序之間的關係.
這段經文非常詳細.
就是說了上帝怎樣創造天地萬物.
首先我讀一讀給大家聽.
歌羅西書一章十六節.
他說因為萬有都是靠他做的.

$^{681}$那個他就是愛子.
耶穌基督.
靠他做的無論天上的地上的.
能看見不能看見.
有為主子執政掌權.
一切都是藉著他做的.
又是為他做的.
他在萬有之先萬有也靠他而立.
保羅在歌羅西書這一段.
是採用了那個概念.
是採用了真言的第八章.
22到31節.
為什麼這樣呢.
因為真言很清晰告訴你.
未創造世界而先就有了我.
我就是什麼.
智慧Sophia.
從亙古從太宗未落納世界.
以前就已經有這個智慧.
所以22到26節是說了智慧不是被造物.
它是一個創造主同級的一件.
一樣的位格.
去創造我們天父一同創造這個世界.
27到31節就說明他怎樣去創造這個世界.
並且在30節你會發現.
我用紅色highlight有個字.
去形容這個智慧.
和合本就譯作公司.
其實不是company.
是說engineer.
也就是應該譯作工程師.
智慧在創造的時候.
稱呼自己是創造秩序的工程師.
也就是說耶穌基督這個智慧.
在創造天地萬物的時候.
正正應驗保羅在歌羅西書一章的16節所說.
萬物都是藉著他造的.
又是為他造的.
他在萬有而先.
萬有是靠著他而立的.

$^{721}$工程師看回希伯來文.
原來他那個字根.
你懂得讀的.
就跟我一起讀好不好.
Amen.
是不是你每天都讀的.
奇元圖就說了.
Amen這個意思是解作真實確立真相真理的意思.
也就是說這個世界.
什麼叫真相什麼叫真理.
並不是我們人類用我們嘴巴說了就算.
而是有一個客觀既定的真理.
植入了這個創造秩序當中.
好像一個晶片一樣放了進去.
在整個創造的DNA.
所以你不會說我想這樣對就等於對.
因為這個世界真是一個客觀的存在.
的真相和真理.
譬如舉個例子.
物理原則其中一樣叫地心吸力.
地心吸力如果你很看不開.
由30樓跳下去的話.
你就一定怎樣.
就頭破血流.
你就一定死的.
不是的高牧師.
如果我有錢一點是不是不一樣.
有錢一點是沒有分別的.
因為是無差別的.
所有人跳下去都是一定死的.
不是的我早上用了已到自見靈修會不會不一樣.
答案是不會的.
不會的.
一定是會跳下去一定是死的.
我們會發現在創造秩序面前.
人是何等渺小.
我們沒有辦法逾越和突破創造秩序的法則.
無論你有錢有地位有身份.
甚至有資源甚至我更熟稔更健全.
都是沒有辦法有任何的豁免權.

$^{761}$我們人一定是生老病死.
我們人一定是沒有辦法超越.
這個創造秩序給我們的限定和法則.
如果是這樣我們就帶著一個敬畏的心.
與其和這個創造秩序對著幹.
倒不如迎合這個創造秩序的法則.
以至你能夠活得有智慧有價值.
拒絕愚昧的事情.
所以當那些惡人也問也沒有.
就用一些說話告訴你.
這樣才聰明兄弟兄弟.
這樣才聰明姐姐.
這樣的時候.
你會用真言的智慧進行審視.
就會知道別人認為聰明的事情.
未必真的聰明.
可能真的愚昧也未必.
成為我們的榜樣和提醒.
在今天的聖誕節.
聖誕節我們都知道.
紀念耶穌基督降生為人.
道成肉身.
降載馬槽的一個偉大的福音.
而這個福音原來冥冥中已經預備好.
是何時預備好.
是創造世界的時候已經預備好.
而耶穌基督是神.
是藉著耶穌創造世界.
在這個法則裡面.
祂不是選擇用一個高高在上的姿態.
描繪祂的出生.
而是選擇用世上漢為愚絕的方式.
道成肉身.
馬槽一塊布一個嬰孩.
成為一個記號.
道成肉身.
讓我們明白什麼叫做真正的智慧.
讓我們在這個聖誕節裡面.
能夠開始過敬畏上帝的生活.
求主幫助我們一起祈禱.

$^{801}$慈悲人愛的天父多謝你.
透過《箴言》第八章.
我們看到耶穌你就是智慧的化身.
你創造天地萬物.
你創造了我們每一個人的生命.
叫我們不再沉淪過活.
叫我們能夠成為一個智者.
幫助我們拒絕愚昧的方式.
和那種生活的價值觀.
改造我們的生命.
迎合你那份智慧的向導.
以致我們成為一個追隨基督的人.
活出基督的樣式.
成為人們認識主的渠道.
幫助我們每一個.
奉耶穌基督的名祈禱.
阿們.
謝謝大家.
(咳嗽).
\newpage



\section{約翰福音 11:1-27}
\label{sec:yqcP_iDkMnA}
\textbf{2025年12月28日崇拜直播｜梁家麟牧師｜復活和生命在我｜約翰福音11\_1-27}
\newline
\newline
連結: \href{https://youtube.com/watch?v=yqcP_iDkMnA}{\texttt{https://youtube.com/watch?v=yqcP\_iDkMnA}} ~~~~ 語音日期: 2025-12-28
\newline
\newline
\hyperref[sec:wYKTi6uXSJI]{\small{< < < PREV SERMON < < <}}
~
\hyperref[sec:index_chronic]{\small{[返順時目]}}
~
\hyperref[sec:DESvXfy__1I]{\small{> > > NEXT SERMON > > >}}
\newline
\newline
約翰福音 11:1-27
\newline
\begin{longtable}{cl}
\hline
\hline
章節 & 經文 (和合本修訂版)\\
\hline
11:1 & \begin{tabularx}{0.7\textwidth}{X} 有一個患病的人，名叫拉撒路，住在伯大尼，就是馬利亞和她姐姐馬大的村莊。 \end{tabularx} \\ \\ \relax
11:2 & \begin{tabularx}{0.7\textwidth}{X} 這馬利亞就是那用香膏抹主，又用頭髮擦他腳的；患病的拉撒路是她的弟弟。 \end{tabularx} \\ \\ \relax
11:3 & \begin{tabularx}{0.7\textwidth}{X} 姊妹兩個就打發人去見耶穌，說：「主啊，你所愛的人病了。」 \end{tabularx} \\ \\ \relax
11:4 & \begin{tabularx}{0.7\textwidth}{X} 耶穌聽見後卻說：「這病不至於死，而是為了神的榮耀，為要使神的兒子藉此得榮耀。」 \end{tabularx} \\ \\ \relax
11:5 & \begin{tabularx}{0.7\textwidth}{X} 耶穌素來愛馬大和她妹妹，以及拉撒路。 \end{tabularx} \\ \\ \relax
11:6 & \begin{tabularx}{0.7\textwidth}{X} 他聽見拉撒路病了，仍在原地住了兩天， \end{tabularx} \\ \\ \relax
11:7 & \begin{tabularx}{0.7\textwidth}{X} 然後對門徒說：「我們再到猶太去吧！」 \end{tabularx} \\ \\ \relax
11:8 & \begin{tabularx}{0.7\textwidth}{X} 門徒對他說：「拉比，猶太人近來要拿石頭打你，你還再到那裡去嗎？」 \end{tabularx} \\ \\ \relax
11:9 & \begin{tabularx}{0.7\textwidth}{X} 耶穌回答：「白天不是有十二小時嗎？人若在白天行走，就不致跌倒，因為他看見這世上的光。 \end{tabularx} \\ \\ \relax
11:10 & \begin{tabularx}{0.7\textwidth}{X} 人若在黑夜行走，就會跌倒，因為他沒有光。」 \end{tabularx} \\ \\ \relax
11:11 & \begin{tabularx}{0.7\textwidth}{X} 耶穌說了這些話，隨後對他們說：「我們的朋友拉撒路睡了，我去叫醒他。」 \end{tabularx} \\ \\ \relax
11:12 & \begin{tabularx}{0.7\textwidth}{X} 門徒就說：「主啊，他若睡了，就會好的。」 \end{tabularx} \\ \\ \relax
11:13 & \begin{tabularx}{0.7\textwidth}{X} 耶穌說這話是指拉撒路死了，他們卻以為他是指通常的睡眠。 \end{tabularx} \\ \\ \relax
11:14 & \begin{tabularx}{0.7\textwidth}{X} 於是耶穌就明白地告訴他們：「拉撒路死了。 \end{tabularx} \\ \\ \relax
11:15 & \begin{tabularx}{0.7\textwidth}{X} 為了你們的緣故，我不在那裡反而歡喜，為要使你們信。現在我們到他那裡去吧。」 \end{tabularx} \\ \\ \relax
11:16 & \begin{tabularx}{0.7\textwidth}{X} 於是那稱為低土馬的多馬對其他的門徒說：「我們也去和他同死吧！」 \end{tabularx} \\ \\ \relax
11:17 & \begin{tabularx}{0.7\textwidth}{X} 耶穌到了，知道拉撒路在墳墓裡已經四天了。 \end{tabularx} \\ \\ \relax
11:18 & \begin{tabularx}{0.7\textwidth}{X} 伯大尼離耶路撒冷不遠，約有六里路。 \end{tabularx} \\ \\ \relax
11:19 & \begin{tabularx}{0.7\textwidth}{X} 有好些猶太人來看馬大和馬利亞，要為她們弟弟的緣故安慰她們。 \end{tabularx} \\ \\ \relax
11:20 & \begin{tabularx}{0.7\textwidth}{X} 馬大聽見耶穌來了，就出去迎接他；馬利亞卻仍然坐在家裡。 \end{tabularx} \\ \\ \relax
11:21 & \begin{tabularx}{0.7\textwidth}{X} 馬大對耶穌說：「主啊，你若早在這裡，我弟弟就不會死了。 \end{tabularx} \\ \\ \relax
11:22 & \begin{tabularx}{0.7\textwidth}{X} 我也知道，即使現在，你無論向神求甚麼，神也必賜給你。」 \end{tabularx} \\ \\ \relax
11:23 & \begin{tabularx}{0.7\textwidth}{X} 耶穌對她說：「你弟弟會復活的。」 \end{tabularx} \\ \\ \relax
11:24 & \begin{tabularx}{0.7\textwidth}{X} 馬大對他說：「我知道在末日復活的時候，他會復活。」 \end{tabularx} \\ \\ \relax
11:25 & \begin{tabularx}{0.7\textwidth}{X} 耶穌對她說：「復活在我，生命也在我。信我的人雖然死了，也必復活。 \end{tabularx} \\ \\ \relax
11:26 & \begin{tabularx}{0.7\textwidth}{X} 凡活著信我的人必永遠不死。你信這話嗎？」 \end{tabularx} \\ \\ \relax
11:27 & \begin{tabularx}{0.7\textwidth}{X} 馬大對他說：「主啊，是的。我信你是基督，是神的兒子，就是那要臨到世界的。」 \end{tabularx} \\ \\
[1ex]
\hline
\hline
\end{longtable}
$^{1}$今日經訓在《約翰福音》11章1至27節.
在《新約聖經》143至144頁.
《約翰福音》11章1至27節.
請聽我讀出第11章第一節.
有一個患病的人名叫拉薩盧.
住在伯特尼.
就是瑪利亞和他姐姐瑪大的村莊.
這瑪利亞就是那榮香膏末主.
又用頭髮擦他腳的.
患病的拉薩盧是他的兄弟.
他姐妹兩個就打發人去見耶穌.
說:主啊,你所愛的人病了.
耶穌聽見就說:.
這病不至於死.
乃是為神的榮耀.
叫神的兒子因此得榮耀.
耶穌數來阿瑪大和他妹子並拉薩盧.
聽見拉薩盧病了.
就在所居之地營住了兩天.
然後對門徒說:我們再往猶太去吧.
門徒說:拉比.
猶太人近來要拿石頭打你.
你還往哪裡去嗎?.
耶穌回答說:白日不是有十二小時嗎?.
人在白日走路.
就不至跌倒.
因為看見者世上的光.
若在黑夜走路.
就必跌倒.
因為他沒有光.
耶穌說了者說:隨後對他們說.
我們的朋友拉薩盧睡了.
我去叫醒他.
門徒說:主啊.
他若睡了就必好了.
耶穌這話是指著他死說的.
他們卻以為是說照常睡了.
耶穌就明明地告訴他們說.
拉薩盧死了.
我沒有在那裡就歡喜.

$^{41}$這是為你們的緣故.
好叫你們相信.
如今我們可以往他那裡去吧.
多瑪又稱為底土瑪.
就對那同作門徒的說.
我們也去和他同死吧.
耶穌到了.
就知道拉薩盧在墳墓裡.
已經四天了.
不帶離耶路撒冷不遠.
約有六里路.
有好些猶太人來看馬大和瑪利亞.
要為他們的兄弟安.
為他們.
馬大聽見耶穌來了.
就出去迎接他.
瑪利亞卻依然在座家裡.
馬大對耶穌說.
主啊,你若早在這裡.
我兄弟必不死.
就是現在我也知道.
你無論向神求什麼.
神也必賜給你.
耶穌說.
你兄弟必然復活.
馬大說.
我知道在末日復活的時候.
他必復活.
耶穌對他說.
復活在我.
生命也在我.
信我的人雖然死了.
也必復活.
凡活著信我的人.
必永遠不死.
你信這話嗎?.
馬大說.
主啊,是的.
我信你是基督.
是神的兒子.

$^{81}$就是那要臨到世界的.
祝聖誕快樂.
又預祝新年快樂.
我們繼續《若人福音》.
很長的故事.
很長的旅程.
我們來到第十一章.
我很少跟大家說.
《若人福音》結構的問題.
因為我覺得我們不是去.
研讀一本書.
我們是藉著這本書.
問上帝今天跟我們說什麼.
所以信息是比書卷本身的結構重要.
不過提一提.
很多人將一章十九節.
一到十八節是前言.
太初有道的重要前言.
然後一章十九節到十章四十二節.
視為《若人福音》第一個大段落.
並且這個大段落的開始和結尾.
都是用施西約翰做主題的.
以施西約翰開始.
然後接著到十章四十二節.
就說到耶穌和他的門徒.
為了逃避在耶路撒冷的猶太人的追殺.
於是就退到了施西約翰事奉的地方.
這是第一個大段落.
所以十一章就開始另一個大段落.
耶穌基督在人間醫病趕鬼.
行了很多偉大的神蹟.
不過沒有任何神蹟彼得上讓死人復活更加震撼人心.
因為復活就是跨越了生死的界線.
將明明是一條單程的不歸路扭轉過來.
我們安慰死者家人通常說的是.
人死不能復生.
你們要節哀順變.
所以死人復活就是打破了人死不能復生這一條定律.
《若人福音》十一章記載了耶穌洗拉撒路復活的故事.
故事很長.

$^{121}$所以我們要分兩次來講.
這個故事首先讓我們感到困擾的是.
只有《若人福音》記載了這個故事.
《三卷符類福音》隻字未提.
馬太和路加都有提到耶穌洗死人復活.
並且通常是放在耶穌的口中.
在耶穌的自我宣稱.
證明祂自己的神聖身份.
如何證明祂是基督呢.
你告訴約翰聽我做過什麼.
其中包括洗死人復活.
不過卻沒有相關的復活事件記載.
按道理死人復活這麼震撼性的事蹟是不可能忽略不記的.
除非是聖經作者不知道這件事.
不過如果我們只看約翰的記述.
耶穌行這個神蹟是在一個很公開的場合做的.
祂的門徒在那裡.
拉撒路的鄰居在那裡.
還有很多其他人都在那裡.
在故事末段36到43節更提到.
有人在現場看到這件事之後.
就將這件事報告給耶路撒冷的祭司團隊.
結果大祭司蓋亞法下決心要殺死耶穌.
連猶太教最高當局都知道這件事.
你告訴我門徒為什麼不知道.
故事發生的地點是在伯大利.
離耶路撒冷只有幾公里遠.
所以連猶太教領袖都知道不是很奇怪的事.
問題是這麼佳智學問的一件神蹟.
並且在時間上是很靠近耶穌上十字架.
已經是最高峰的時候.
為什麼馬太馬可路加都不記呢.
你說篇幅有限.
福音書作者只能夠選最重要的去記.
但是聖父都不夠重要.
沒理由不記 沒理由不提.
就《芙蕾福音》缺了這個故事.
學者做了很多的討論爭論.
自由派的學者自然就說.
因為這件事不是真實發生的.

$^{161}$是作的 是約翰作出來的.
所以是他獨家報導.
我們當然不會接納.
否定聖經的歷史性和事實性的解說.
一個比較可取.
比較多福音派的學者認同的解說是.
馬可和馬太這兩卷福音書.
比較多其實是受彼得的影響.
馬可寄信的多數是來自彼得的回憶.
馬太是依據馬可的版本.
做一個擴充版的技術.
彼得很可能不是這次死人復活的現場.
沒有目睹拉薩路復活的神蹟.
我們如果看約翰福音.
我們就發現從六章六十八節開始.
一直到十三章六節.
這六章的聖經都沒有提到彼得的名字.
或者彼得是因某些緣故.
沒有在耶穌身邊請了長假回大陸.
排隊堵塞很久.
所以不能回來.
一個小小的旁證.
在十一章這一段聖經的技術裡面.
耶穌帶同門徒前往伯大利.
途中有些討論.
其中沒有提到彼得的名字.
但提到唯一具名字的發言人是多瑪.
在其他聖經的地方.
通常最勇於發言.
會記名的門徒是誰?.
是彼得.
為什麼會是多瑪呢?.
我語無倫次亂說話.
達哥在的時候.
只是他在說話.
他不在的時候.
舒金就說話.
如果達哥在的話.
怎麼輪到他說話呢?.
就是這樣.

$^{201}$為什麼多瑪可以說話呢?.
或者記載他說話呢?.
可能是大哥不在.
彼得缺席了.
這是一個推論.
無疑這個解釋.
也不能完全消去所有的困惑.
彼得不在場.
事後其他門徒.
也應該會把這件事告訴他.
這麼轟動的事.
他不可能永遠一無所知.
另外.
路加在馬可之外.
所依據的史事來源.
也應該有提到這個死人復活的故事吧.
無論如何.
三卷伏特哲福音.
沒有提到這個故事.
總是令人有些困惑.
但是.
我們一定要很簡單.
做個結論.
這個無損我們對約翰這卷福音的.
真實性的信任.
我們一起去看.
約翰的獨家報導.
故事發生在伯大利.
這是鑒嵐山東南一個小的村落.
離耶路撒冷只有三公里.
聖經說六里不是公里.
是很短的.
所以三公里.
走一個小時就應該到了.
走幾萬也要到.
三公里就到了.
那裡住了一個家庭.
三姊弟.
瑪大,瑪利亞和拉薩路.
三個人似乎都沒有結婚.

$^{241}$否則他們不會住在一起的.
耶穌和這個家庭是非常熟悉的.
由於伯大利是耶路撒冷去耶利哥的必經之路.
耶穌和門徒每一次往來耶路撒冷.
都有可能記住在這個家庭裡面.
瑪大和瑪利亞兩姐妹.
在路加福音十章三十八節都有提過.
你記得那個故事.
一個煮飯一個聽道.
這個故事你很熟悉的.
但是在那個故事裡面.
沒有提到他們有弟弟.
所以提到他們有弟弟的.
只有在這裡.
在《約翰福音》十一章和十二章.
此地三人和耶穌的關係都很好.
聖經講第五節.
耶穌掃來愛瑪大和他妹子並拉薩路.
這一段的主角是拉薩路.
在第三節也有提到.
他是耶穌所愛的人.
耶穌在第六節又稱他為「我們的朋友」.
「我們」是指耶穌和門徒.
拉薩路是他們整個團隊的朋友.
故事一開始就說拉薩路病了.
沒有提到什麼病.
病得很重.
無藥可治是肯定的.
瑪大和瑪利亞於是想找耶穌幫忙.
他們的警察去找耶穌.
故事說當耶穌和門徒到達的時候.
拉薩路已經死了並且埋葬了四天.
第十七節.
聖經沒有提到耶穌當時在哪裡.
如果我們接著《約翰福音》十章.
第四十節的講論.
耶穌和門徒為了躲避猶太教當局.
即耶穌教領袖的追殺.
於是就退到約旦河外.
即斯西約翰服侍的地方.

$^{281}$如果我們就約旦河外來看.
這個地方是呼應一章二十八節.
那裡寫到約旦河外的伯大利.
聖經很麻煩的有兩個伯大利.
一個是現在這裡的地方.
即拉薩路的地方.
另一個伯大利在一百五十公里以外.
就是斯西約翰服侍的地方.
那是今天加利地東北部一個叫做巴塔尼亞的地方.
距離耶路撒冷約一百五十公里遠.
距離南部的猶太.
需要大概行三天的路程.
一天行一百五十公里差不多.
一個小時行四五公里.
你也要行十幾個小時的了.
也很辛苦的了.
信差由伯大利出發找到耶穌的時候.
是在三天之後.
當耶穌聽到拉薩路的病訊.
從耶穌的反應顯示拉薩路還沒死的.
那時候還沒死的.
在第四節.
但耶穌就刻意擔言兩天.
當他決定動身去伯大利的時候.
就宣告還沒出發已經宣告拉薩路已經死了.
顯然拉薩路是在耶穌擔國的這兩天死的.
接著耶穌和門徒去伯大利需要走三天的路程.
到達伯大利的時候.
聖經說拉薩路已經埋葬了四天.
時間剛剛好.
三天走路.
四天.
就是剛剛好在耶穌擔國的兩天之間.
大概是三加一.
就是四天.
無論如何.
我講這些你們聽了也覺得很煩的事情.
我想講一件事.
從耶穌的住處到達伯大利.
由於需要三天的時間.

$^{321}$即使耶穌聽到訊息.
馬上第一時間趕路.
也趕不上的.
因為死了四天.
要走三天.
三天除非他跑.
否則的話.
走三天是趕不上的.
四天的路程.
所以拉薩路的死亡和耶穌的擔國是沒有關係的.
第六第七節說.
耶穌聽見拉薩路病了.
就在所居之地營住了兩天.
然後對門徒說.
我們再往猶太去吧.
當然你可以問.
賓喪也可以早一點去.
為什麼要拖延兩天呢?.
耶穌的擔言是故意的.
要等到拉薩路病死了.
並且在他還沒走之前.
已經明確告訴門徒.
拉薩路已經死了.
然後他才動身.
原因是.
耶穌不打算在門徒中間.
再多走一個醫病的神蹟.
他準備要走的一個死人復活的神蹟.
這個人是死了.
連屍體也埋葬好.
並且是足足四天才發生的復活事件.
所以期間.
應該不會存在假死,誤判的可能性.
你埋葬了四天就餓死了你.
對不對?.
就算還沒死.
四天不喝水,不吃東西.
都一定死了.
是不是?.
在十五節.

$^{361}$耶穌對門徒說.
我沒有再那裡就歡喜.
只是為你們的緣故.
好叫你們相信.
下次我會再解釋這段經文.
注意在第七節.
耶穌的說法是.
我們再往猶太去吧.
他不是說.
他去伯大利這條村莊.
而是強調他去南部的猶太.
他知道猶太是代表什麼?.
是代表耶路撒冷.
就是猶太教的權威所在.
也知道這裡意味著政治的風險.
前往猶太.
耶路撒冷的猶太教領袖.
已經準備好要殺害耶穌.
你去到一個距離耶路撒冷只有三公里遠的伯大利.
是很危險的事.
耶穌說了這番話之後.
門徒立刻就有反對的意見.
門徒說:拉比.
猶太人近來要拿石頭打你.
你橫往那裡去嗎?.
第八節.
耶穌這樣回答.
白日不是有十二小時嗎?.
人在白日走路就不至跌倒.
因為看見這世上的光.
若在黑夜走路就必跌倒.
因為他沒有光.
這段說話有兩個可能的意思.
第一.
簡單的.
我們要趁著白日去做事.
把握做事的時機.
不要瞻前顧後.
這與前面九章四節.
我們也看過.

$^{401}$耶穌的說話意思是一樣的.
趁著白日.
我們必須作那差我來者的功.
黑夜將到.
就沒有人能作功了.
趁著還有白日.
還有光.
就做事吧.
這是第一個可能的解釋.
第二個可能的解釋.
是耶穌是世界的光.
如果他不去需要光的人那裡.
他們就沒有光了.
所以儘管前路凶險.
為了有需要的人.
耶穌仍然需要冒險前往.
我們知道.
如果耶穌懼怕上十字架.
拒絕道成肉身.
人間根本就再沒有世上的光.
這是約翰福音的前言強調.
如果耶穌不道成肉身.
人間怎會見到光明呢?.
趨吉避凶是人之常情.
但完成使命比保存性命更加重要.
耶穌向門徒宣告拉撒路已經死了.
祂首先用睡了這個說法.
在門徒誤會睡了就是睡覺之後.
祂再澄清拉撒路已經死了.
距離現場有幾天的步程.
百多公里.
但耶穌竟然已經預知.
祂感知了這個死訊.
當然這是一個超自然的能力.
門徒多瑪聽到拉撒路已經死了之後.
耶穌卻仍然要街頭門徒前往猶太.
感到非常不解.
祂對身邊的人嘰嘰咕咕地說.
我們也去和祂同死吧.
這句說話不應該是滿懷信心和勇氣的宣告.

$^{441}$是 我們和耶穌一起去慷慨就義.
而拉撒路已經死了.
他們為何要去呢?.
如果祂還未死就去救祂.
救了性命也要拼一拼.
死了純粹是奔喪.
即是說我陪祂一起死.
這是完全多餘無聊的行動.
不過耶穌說要去就跟著去死.
在懷疑和困惑中.
門徒跟隨耶穌前往伯大利.
在這個時候拉撒路的喪禮已經做完很久了.
因為在近東的地方天氣很熱.
屍體很容易發臭很容易腐爛.
所以通常是隨死隨埋葬.
不會等的.
不會等親戚來才舉行國安式禮.
通常都不是的.
就近的人就有參加.
喪禮就參加.
遠的人來民信趕到都已經葬了.
不會等的 是不能等的.
所以已經死了.
並且馬上葬了.
已經有四天的時間.
在聖經著名埋葬了四天.
為何是四天呢?.
有很多說法.
搞一些很複雜的東西.
對不起.
有幾個說法.
第一個我講兩個比較有趣的說法.
第一.
猶太人有一個信念.
人死後.
靈魂就會在肉身旁邊徘徊三天.
看看有沒有機會還陽.
其實我們中國人都是這個說法.
所以可以招魂的.
三天之後才去陰間.

$^{481}$所以埋葬了四天.
就杜絕了一切僥倖的可能性.
靈魂都走了.
根本就不在身邊.
這是第一個說法.
第二個說法是.
在史海古卷的其中一份文獻.
告訴你也沒什麼意思.
4Q521的文獻裡面.
曾經提到當時的人.
相信真正的彌賽亞必須要行四個神蹟.
提到預告彌賽亞有四大神蹟.
包括伊茲大麻風.
趕走龍亞人的鬼.
醫好生來的黑眼的.
和死了最少三天的人復活.
如果這份文獻的記載確實.
使徒約翰在記述耶穌生平的時候.
刻意地要將這四個事件納入耶穌的生平裡面.
所以刻意要突出.
耶穌在拉撒路復活的事件.
是在祂死了超過三天之後發生的.
這些複雜的東西.
輕輕提一下就算了.
回來這個故事.
陸續有遠方的親友文信前來.
安慰死者家人.
所以現場就很熱鬧.
第十八到十九節.
伯大利離耶路撒冷不遠.
約有六里路.
我再說不是六公里.
是三公里左右.
有好些猶太人來看馬大和瑪利亞.
要為他們的兄弟安慰他們.
這段經文引來很多《識經》上的爭論.
《識經》說有很多來自耶路撒冷的猶太人.
來探望安慰喪家.
住在耶路撒冷裡面的人.
通常都是社會地位比較高.

$^{521}$比較富有的人.
特別是四十六節說.
這些來探望的人.
後來就去了見法利賽人.
最終驚動了祭史長.
該阿發.
所以我們就這段經文去估計.
馬大和瑪利亞的親友圈裡面.
應該有一些權貴人物.
所以才回到耶路撒冷.
能夠接觸到高層.
並且他們是來自耶路撒冷的.
是這樣我們都可以推斷.
各從其類.
馬大和瑪利亞應該是有錢和身份的.
馬大和瑪利亞應該是有錢人.
有兩個證據支持這個推論.
第一.
馬大常常接待耶穌和門徒.
和今天田先生經常請客是兩回事.
古代社會食物很昂貴.
能夠經常請客.
通常最少要付一個風雨人的錢才能做到.
一般人很懶的.
省著做 省著花.
第二個理由.
是跟著.
在十一章二節提到.
瑪利亞是那個用香膏羔耶穌的女人.
跟著在十二章就記述這個故事.
瑪利亞用香膏羔耶穌.
能夠用三十兩的銀錢.
一個工人一年的工資.
一年的工資.
用來羔耶穌.
這麼豪奢的行為.
關鍵不是在於.
你對耶穌是否願意傳言奉獻.
我想傳言奉獻.
但我沒有家底傳言奉獻.

$^{561}$你明不明.
我只能用蒸蒸豉油來羔耶穌.
真拿大香膏我沒有錢買.
我根本沒有.
沒有就沒有.
能夠豪到出手.
用這樣的香膏來羔耶穌的腳.
這個一定是有家底的人.
沒有家底怎麼做呢.
當然也有學者不同意.
有人認為.
瑪利亞不是原生家庭有錢.
不是咬著銀錢出世.
而是多年賣淫.
有人猜測他做不正當的工作.
所以積累了一筆錢.
這個說法將瑪利亞.
馬大瑪利亞.
撒拉薩路爾家庭的瑪利亞.
和莫大拉瑪利亞等同.
你知道聖經有很多個瑪利亞.
這是兩個學派的爭論.
到底這裡的瑪利亞.
是不是莫大拉瑪利亞.
因為莫大拉瑪利亞.
又有高耶穌的一次.
是不是兩個瑪利亞分別高耶穌呢.
如果是同一個故事的話.
兩個人就是同一個人.
這些扭曲的事.
我們留下十二章再說.
反對馬大和瑪利亞很有錢的人.
所精忍的一個理由是.
耶穌曾經說過一個財主和拉薩路的故事.
你記不記得.
不太記得吧.
路加福音十六章十九節二十節.
拉薩路在比喻裡.
是一個窮困到行乞渡日的人.
行乞渡日家裡不可能有錢的.

$^{601}$這裡又牽涉到另一個複雜問題.
剛才的問題是.
這個瑪利亞是不是那個瑪利亞.
這裡的問題是.
這個拉薩路是不是那個拉薩路.
耶穌說的拉薩路這個比喻.
他隨便找個名字.
小強.
志明.
還是真的就是這個拉薩路呢.
這個就很複雜了.
瑪利亞.
瑪大瑪利亞是不是很有錢.
這個瑪利亞是不是那個瑪利亞.
這個拉薩路是不是那個拉薩路.
你覺不覺得解經很麻煩.
真的很麻煩.
回歸正傳.
很多人來調驗.
馬大爺是行動派.
收拾悲痛的心情.
做接待親友的工作.
有人報信.
耶穌和門徒前來探望.
他立刻飛奔跑到村子的外頭.
去迎接耶穌.
他看到耶穌的第一句話是.
二十一節.
主啊,你若早在這裡.
我兄弟必不死.
意思就引人說的就是.
耶穌,你來得太遲了.
接著一句話.
不容易理解.
二十二節.
就算現在我也知道.
你無論向上帝求什麼.
上帝也必賜給你.
一個比較合理的解釋是.
馬大爺說.

$^{641}$耶穌,可惜你來遲了.
如果你早點來.
我弟弟就不會死了.
因為我相信.
只要你幫我們向上帝祈求.
無論求什麼.
上帝都會答應.
你要我弟弟得醫治.
他肯定就得痊癒.
就不會造成今天死亡的情況.
不過你會看到.
如果真是這個意思的話.
瑪利亞是相信.
耶穌的祈求是什麼都可以.
但是有條紅線.
死了就不行了.
死了就一了百了.
死了就不能再求了.
其實這是人之常情的想法.
你記得大衛為他的兒子祈求.
他死了.
洗澡,吃飯.
因為都沒用了.
再求都失氣了.
死了就是死了.
所以死是條紅線.
過了就什麼都不能再求了.
你親人病了.
你可以中醫,西醫,無醫,泰醫.
什麼都給他.
把多少錢都放在他身上.
他斷氣了.
還有什麼好做.
大家不做了.
就完了.
所以耶穌的祈求有效.
不過只於死亡線.
耶穌回答馬大.
跟他說.
你兄弟必然復活.

$^{681}$馬大回應.
我知道在末日復活的時候.
他必復活.
猶太人中的法利賽派是相信.
在末日死人要復活的.
撒刀該派的人是不相信復活的事件的.
馬大理解耶穌的話是指著末日的復活.
這是一般人對亡者家屬慣常的安慰詞.
就像我們常常安慰死者的家人.
不要太難過.
我們將來在天家還是可以重逢的.
但是這個安慰的話.
仍然無法掩蓋眼前冰冷的事實.
我們和死者在人間是永缺.
人間永別.
天家重逢.
耶穌知道馬大不明白他的說話.
於是鄭重跟他說.
25,26節.
復活在我,生命也在我.
信我的人雖然死了,也必復活.
凡活著信我的人必永遠不死.
你信這話嗎?.
耶穌的說話是個非常偉大的宣告.
復活在我,生命也在我.
這句話中文翻譯是有少許遷就的.
耶穌的原話其實是.
我是復活,我是生命.
就像他說我是生命的糧.
我是羊的門一樣.
我是,我是.
我是復活,我是生命.
是這樣的意思.
他強調.
他不僅僅能夠賜人復活和生命.
他說,我就是復活.
我就是生命.
我以前也和你說過這篇道.
主權善於大能.
耶穌的這句話是非常霸氣的.

$^{721}$他的意思是什麼呢?.
我是復活和生命的主.
我掌管人的生命.
我要人生就生,要人死就死.
我說了算.
他強調的不是能力,是主權.
或者說主權在先,能力在後.
耶穌說了算.
除了超視主權之外.
耶穌的說話也揭示了人類唯一的出路.
聖經說,罪的公家乃是死.
人人都可以死.
自從罪進入人間之後.
死亡就是所有人的命運.
不僅僅是人間的死,並且是永恆的死.
只有耶穌能夠逆轉這個命運.
信我的人,雖然死了,也必復活.
凡活著信我的人,必永遠不死.
按著定命,每個人都會面對人間的死亡.
基督徒都無法豁免的.
不過,信耶穌的人.
人間的死亡只是一道門.
通過了,就由此進入永生.
只要我們在人間相信耶穌.
就得著永生.
對耶穌這番說話.
馬大的回應在二十七節.
主啊,是的,我信你是基督,是上帝的兒子.
就算那要臨到世界的.
嘿嘿嘿,明顯的,你看這句話就知道.
馬大其實聽不明白耶穌說什麼.
他聽不明白耶穌說什麼.
所以他的回應是牛頭不對馬嘴的.
不過,我也要說.
馬大的說話,也是對耶穌所有教導,所有行動.
最恰當的回應.
而馬大聽不明白.
復活在我,生命也在我,是什麼意思.
但是他從一開始就認定耶穌是基督,是上帝的兒子.
所以耶穌說什麼都是正確的.

$^{761}$都肯定會應驗的.
我信你,這裡用的是過去式.
表示他一早就信了,不是現在才信.
就算他弟弟如今失救死了.
也沒有損他對耶穌的信心.
我信你是基督,是上帝的兒子.
我們記得,當耶穌考問門徒.
你們說我是誰的時候.
彼得就給了一個正確答案.
就是你是基督,是永生上帝的兒子.
馬大福音十六章十六節,你記得吧.
這段聖經是天主教的壓陣經文.
教皇的權威就建在這篇經文裡.
耶穌在讚賞彼得的答案正確之餘.
宣告教會要建在這個認信的基礎上.
所以,這個給我們大家一個備給,袋落袋了.
我信你是基督,是上帝的兒子.
這句話是一切信仰問題的標準答案.
你信不信耶穌能夠解決當前的問題呢?.
你回答,耶穌是基督,是永生上帝的兒子.
有人問你,你信不信旺朕有美好的前景呢?.
你就回答,耶穌是基督,是永生上帝的兒子.
有人問你,你對香港的未來還有沒有信心呢?.
你就回答,耶穌是基督,是永生上帝的兒子.
不管問題是什麼,你這樣回答都不能說是錯的.
你明不明白?我不知道人間這個問題,那個問題是怎麼解決的.
我只知道,耶穌是基督,是永生上帝的兒子.
這句話萬事萬靈,你考試的時候.
所有問題都不會回答,你就回答這句,填這句進去.
可能你的老師會打大毛給你分數的.
有人問你,姐妹你的病能不能夠醫好?.
弟兄你的問題能不能夠解決?.
姐妹,面對連番打擊,你仍然對耶穌有信心嗎?.
我們的答案都是,我信耶穌是基督,是上帝的兒子.
我們聖經就看到這裡,就看完了.
我們來到應用部分.
我們專注思考耶穌那個偉大的宣告.
福活在我,生命也在我.
我跟你們說過不止一次.
主權先於能力,主權重於能力.

$^{801}$上帝有很多種不同的屬性.
祂慈愛,祂公義,一大堆,全知全能.
但是上帝最基本的屬性,最重要的屬性是什麼?.
Sovereignty,主權,上帝是上帝.
上帝有權,想刺的就是祂,收取的就是祂.
祂有權,祂可以使人生,可以使人死.
祂有權,上帝是上帝.
認識耶穌的主權比相信耶穌的能力更加重要.
相信耶穌有能力,期待祂有能力解決我們的問題.
滿足我們的需要.
這個對耶穌的理解,其實不會比我們信大慈大悲,觀世音菩薩.
或者有求必應的黃大仙高明很多.
你可能會說,不是的,我們信的耶穌能力比觀音,比黃大仙要大.
對我這個小人物,這句話是沒有意思的.
畢竟作為小老百姓的我,我要求不高.
我心胸都不說,我不需要神明能不能移動天體,星壽.
能不能止息,世界和平,和戰爭,省一點吧.
我只需要神明能移動那六個波波的號碼,對嗎?.
三個,不用六個,我沒那麼貪心,是嗎?.
中過三獎也好,對嗎?.
我的腰很小,所以少許能力也足夠.
不用拜很大的神,對嗎?.
我不要求什麼和平,我要求我家宅平安就夠了,對嗎?.
耶穌不僅是滿天神佛的其中一個.
耶穌不僅是位列先班中的守衛,最有能力的那個.
他是創天造地,掌管宇宙人生,掌管人類歷史發展.
掌管你和我生命的那一位,他是上帝,他有權.
上帝掌握一切主權,萬有都出於他,屬於他,聽任他的擺佈.
耶穌不是有解決問題的答案,他是答案.
耶穌不僅能夠幫助我們找到出路,他就是出路.
耶穌不僅有能力滿足我們的需要,他就是我們的滿足.
耶穌不僅有方法藉著他的血來洗淨我們的罪.
我們經常說這件事,耶穌是很直接的.
在《瑪利福音》九章二字說,你的罪赦了.
要搞很多手續才能赦你的罪,你的罪赦了.
他有權,他有權,他有權.
所以,他不僅有能力洗罪,他有權柄赦罪.
耶穌不僅帶領我們上天堂享永福.
他是呼召人進入他的國度,在他的主權之下成為他的子民.
實踐他所吩咐的使命.

$^{841}$耶穌所宣講的是國度的福音,國度是關乎上帝在耶穌的主權.
在這裡不跟大家說那些複雜的神學,那些保羅新冠的說法.
總之,不僅是你得救,是強調我們進入他的權威,在他的治權之下.
福音無疑包括了在世得享長壽,在世得享永生的英雄.
但耶穌的福音呼召我們成為他的子民,踐行他的旨意.
向他的旨意在我們生命和生活裡面兌現.
就好像今天我們這些父親全部都很陰春.
父親只是給靈用的那個,我們根本不會覺得父親有權威.
他沒有牙老虎力,不過你用同樣的想法去看耶穌就死定了.
耶穌不僅是有求必應的那個.
耶穌是擺明了,我經常說這些粗俗的話.
要動你,要動你,要動你,要動你,要動你.
他是我們的主,是我們的主.
在個人層面是真的,認識耶穌的主權.
確認他在自己生命裡面的主權.
在口頭宣認基督是主的同時.
在行事為人每一方面去兌現基督是主是很關鍵的.
我要說,這是真基督徒和假基督徒的區別所在.
真基督徒是當基督是主.
當他是我的真正的老闆,他話事,他話事.
我們所信仰的是主耶穌基督,我信我主耶穌基督.
讓我們再一次確認基督的主權.
福活在我,生命也在我.
他掌控我們的生死禍福,一切都是他說了算.
讓我們真誠的祈禱,然而不要照我們的意思.
只要照你的意思.
願耶穌的旨意成就在我們身上.
在旺鎮身上,在香港身上.
在中國之上,在世界之上.
今年,明年,都是耶穌基督的年頭.
\newpage



\section{歌羅西書 1:1-14}
\label{sec:DESvXfy__1I}
\textbf{2026年1月4日主餐崇拜直播｜陳明泉牧師｜門徒DNA:福音介定的身份｜歌羅西書1\_1-14}
\newline
\newline
連結: \href{https://youtube.com/watch?v=DESvXfy__1I}{\texttt{https://youtube.com/watch?v=DESvXfy\_\_1I}} ~~~~ 語音日期: 2026-01-04
\newline
\newline
\hyperref[sec:yqcP_iDkMnA]{\small{< < < PREV SERMON < < <}}
~
\hyperref[sec:index_chronic]{\small{[返順時目]}}
~
\hyperref[sec:eZQl9Arirgg]{\small{> > > NEXT SERMON > > >}}
\newline
\newline
歌羅西書 1:1-14
\newline
\begin{longtable}{cl}
\hline
\hline
章節 & 經文 (和合本修訂版)\\
\hline
1:1 & \begin{tabularx}{0.7\textwidth}{X} 奉神旨意，作基督耶穌使徒的保羅，和我們的弟兄提摩太， \end{tabularx} \\ \\ \relax
1:2 & \begin{tabularx}{0.7\textwidth}{X} 寫信給歌羅西的聖徒，在基督裡忠心的弟兄。願恩惠、平安從我們的父神歸給你們！ \end{tabularx} \\ \\ \relax
1:3 & \begin{tabularx}{0.7\textwidth}{X} 我們為你們禱告的時候，常常感謝我們主耶穌基督的父神， \end{tabularx} \\ \\ \relax
1:4 & \begin{tabularx}{0.7\textwidth}{X} 因為聽見你們對基督耶穌的信心，並對眾聖徒有的愛心。 \end{tabularx} \\ \\ \relax
1:5 & \begin{tabularx}{0.7\textwidth}{X} 這都是因著那給你們存在天上的盼望，它就是你們從前所聽見真理的道，就是福音； \end{tabularx} \\ \\ \relax
1:6 & \begin{tabularx}{0.7\textwidth}{X} 這福音傳到你們那裡，也傳到普天下，並且繼續增長，不斷結果，正如自從你們聽見福音，真正知道神恩惠的日子起，在你們中間也是這樣。 \end{tabularx} \\ \\ \relax
1:7 & \begin{tabularx}{0.7\textwidth}{X} 這福音是你們從我們所親愛、一同作僕人的以巴弗學到的。他為我們作了基督的忠心僕役， \end{tabularx} \\ \\ \relax
1:8 & \begin{tabularx}{0.7\textwidth}{X} 也把聖靈賜給你們的愛告訴我們。 \end{tabularx} \\ \\ \relax
1:9 & \begin{tabularx}{0.7\textwidth}{X} 因此，我們自從聽見的日子就不住地為你們禱告和祈求，願你們滿有一切屬靈的智慧和悟性，真正知道神的旨意， \end{tabularx} \\ \\ \relax
1:10 & \begin{tabularx}{0.7\textwidth}{X} 好使你們行事為人對得起主，凡事蒙他喜悅，在一切善事上結果子，對神的認識更有長進。 \end{tabularx} \\ \\ \relax
1:11 & \begin{tabularx}{0.7\textwidth}{X} 願你們從他榮耀的權能中，得以在一切事上力上加力，好使你們凡事歡歡喜喜地忍耐寬容， \end{tabularx} \\ \\ \relax
1:12 & \begin{tabularx}{0.7\textwidth}{X} 又感謝父，使你們配與眾聖徒在光明中分享基業。 \end{tabularx} \\ \\ \relax
1:13 & \begin{tabularx}{0.7\textwidth}{X} 他救了我們脫離黑暗的權勢，遷移到他愛子的國度裡。 \end{tabularx} \\ \\ \relax
1:14 & \begin{tabularx}{0.7\textwidth}{X} 藉著他的愛子，我們得蒙救贖，罪得赦免。 \end{tabularx} \\ \\
[1ex]
\hline
\hline
\end{longtable}
$^{1}$今日的經訓是記載在.
歌羅西書第一章一至十四節.
歌羅西書一章一至十四節.
在聖經的二百七十九頁.
奉神旨意作基督耶穌使徒的保羅和兄弟提摩泰.
寫信給哥羅西的聖徒.
在基督裡有忠心的弟兄.
願因為平安從神我們的父歸於你們.
上上為你們禱告.
恩慶見你們在基督耶穌裡的信心.
並向眾聖徒的愛心.
是為那給你們存在天上的盼望.
這盼望就是你們從前在福音真理的道上所聽見的.
這福音傳到你們那裡也傳到普天之下.
並且結果增長.
如同在你們中間.
自從你們聽見福音真知道神恩惠的日子一樣.
正如你們從我們所親愛.
一同作僕人的爾巴夫所學的.
他為我們作了基督忠心的執事.
也把你們因聖靈所傳的愛心告訴了我們.
因此我們自從聽見這日子.
也就為你們不住的禱告祈求.
願你們在一切屬靈的智慧上,悟性上.
滿心知道神的旨意.
好叫你們行事為人對得起主.
凡事蒙他喜悅.
在一切善事上結果子.
漸漸多知道神.
照他榮耀的權能得以在各樣的力上加力.
好叫你們凡事歡歡喜喜的忍耐寬容.
叫我們能與眾聖徒在光明中同得基業.
他救了我們脫離黑暗的權勢.
把我們遷到他愛子的國內.
我們在愛子裡得蒙救贖.
罪過得以赦免.
弟兄姊妹新年快樂.
祝願新的一年裡每位弟兄姊妹.
我們經歷上帝更深.
經歷一個蒙福的生活.

$^{41}$在我們今年第一個.
2026年我們開始的主題.
我們希望用歌羅西書.
貫穿我們這幾次的講道.
今年整個教會我們希望重新思考和實踐門徒的生活.
其實早幾年前我們都以門徒不斷作為我們的主題.
但是我們希望今年開始我們會進深一個叫2.0.
這個在教會內部的架構.
我們希望不同的事奉崗位.
可以幫助弟兄姊妹成為一個門徒.
更加重要的就是門徒這個身份.
不只是在教會的事奉崗位裡面.
在我們的生命所有的全部.
我們都要活出一個門徒的身份.
藉著我們的生命被上帝大大來使用.
這個是我們的願景.
也希望我們一起來彼此努力.
說起剛才的門徒的身份.
有些身份很容易會察覺到.
特別是一些可能要穿制服的紀律部隊.
或者有些平時很容易在某個位置裡面看到的.
紀律部隊可能看到他穿著的制服.
你認出他這個身份.
又或者他們可能在學校裡面.
有年紀的差距.
譬如在小學裡面一個大人.
很多時候那個就是老師.
不過有些身份是不容易察覺的.
有一個這樣的例子.
這個是讀神學的例子.
我問准了他.
一位很外向的同學.
名字叫吳珍珍.
很外向.
這位很外向的同學.
他吃飯的時候.
因為我們神學院吃飯.
在飯堂裡面所有人都很隨意.
剛剛入學.
於是這位很外向的同學.

$^{81}$到處跟別人搭訕.
在飯桌當中就看到有一位年長一點的男士.
珍珍就跟這位男士談得很投契.
談得很興高采烈.
大家知道他的性格.
他的人設是這樣的.
然後一直談的時候.
談到吃完飯.
他就問我怎樣稱呼你.
這位弟兄就說.
這裡每個人都叫我大叔.
他說是呀.
他有點尷尬.
不過他說謝謝你大叔.
談得很開心.
然後吃完飯.
隔了幾天就上課了.
珍珍又遇見這位男士.
他就說大叔.
你今天又來上課了.
大叔就說.
是呀 我很多時候都會上課的.
當進到課室的時候.
他就坐下.
這位大叔就站在講台那裡.
其實這位大叔是我們舊約的教授.
不過珍珍就.
看到他站在台上.
就叫他坐下來.
我們當然就笑到拍桌子.
有些身份是不察覺的.
如果他平時不修邊幅.
我真的會當他是一個大叔.
但是當他站在講台上.
他將他自己所學的.
或者累積的經驗.
和大家分享.
他就是我們神學院的老師.
有些身份是容易察覺.
有些身份是不容易察覺的.

$^{121}$直到在我們生命裡面.
有一些表彰.
或者在某些處境當中.
這個身份就顯露出來了.
門徒這個身份我們容易察覺嗎?.
今天你身處在你的生活場景裡面.
別人察覺到你是一個基督的門徒嗎?.
今天我們就用.
歌羅西書第一章一至十四節.
我們思考這個題目.
我們先求主幫助我們.
我們同心祈禱.
求主與我們同在.
幫助我們.
我們求你在我們當中.
顯明你的話語.
讓我們能夠明白真理.
又讓眾靈姐妹在接受的過程裡面.
他們每一個不單聽得明白.
更加被你的話語所感動激勵.
讓我們在這個時代裡面.
我們挺起胸膛成為你所使用的門徒.
幫助孫講的講得清楚.
忠心的和清晰的.
將你的話語帶到弟兄姊妹當中.
願你的聖靈與我們同在.
因為我們深信一切的明白.
或者進入真理.
全部都是你聖靈的工作.
願你此刻帶領我們.
祝福我們下來領受你的話語的時間.
獻上禱告奉求基督耶穌得勝的聖名.
Amen.
我們一起看看歌羅西書第一章.
我們先看第一個段落.
第一節至第八節.
我們首先看第一二節.
奉神旨意作基督耶穌使徒的.
保羅和兄弟提摩泰.
寫信給哥羅西的聖徒.

$^{161}$在基督裡有忠心的弟兄.
願因為平安.
從神我們的父歸於你們.
在這裡很清晰.
整個書卷裡面.
他所說的受信人.
就是一群在哥羅西的聖徒.
或者是哥羅西教會的弟兄姊妹.
哥羅西教會是一個原則上.
大部分都是外邦人的教會.
相對於猶太人.
這些可能是在羅馬社會裡面.
總之不是以色列裔的基督徒.
這個教會是由以巴弗來建立的.
不是保羅一手一腳來建立的教會.
甚至保羅可能是差遣.
或者有不同的門徒.
去到傳福音.
去到外邦.
去到哥羅西這個地方.
就將人帶到上帝面前.
他們接受福音就組成了這個教會.
值得留意的是.
經文裡面保羅稱呼這一群弟兄姊妹.
做什麼呢?.
就是聖徒和忠心的弟兄.
值得留意的是.
聖徒這個字不只是偶然出現的.
在第二節出現過.
在第四節裡面再出現.
到了第十二節還繼續出現.
連續三次.
保羅用這個聖徒的身份.
來稱呼這一群接受書卷的一眾基督徒.
經文裡面用聖徒這個字眼.
他特別強調的就是.
這群人不只是信了耶穌.
他更加講到這一群人.
是被上帝分別出來.
分別為聖.

$^{201}$所以那個聖徒的字本身.
其實帶著一個回應書信裡面.
這一群人面對的處境.
待會兒會再仔細講一下他們面對什麼處境.
不過在這裡.
他特別講到這一群受信的基督徒.
是一群被上帝恩典所選照出來.
他們分別出來的一群門徒.
另外他用另外一個字眼.
就是在基督裡有忠心的弟兄.
有忠心的弟兄.
我們以為有些不忠心.
有些忠心的意思.
其實他在講的就是在基督裡.
篤信的弟兄.
這一群是信了主.
接受耶穌成為他生命救主的人.
他用忠心這個字.
同樣忠心這個字又是出現幾次.
一章第七節.
四章九節.
四章七節.
忠心這個字全部都是放在形容.
那些執事.
傳福音使命裡面.
上帝扶托了給他.
他就忠心將他所領受的.
帶去給他身邊的人.
忠心這個字和學本譯得好的.
不單只是信了.
更加是上帝交付給他的使命.
他都繼續來實踐.
保羅用這兩個字眼.
聖徒和忠心的弟兄.
來描述這一間教會.
接受信件的弟兄姊妹.
第三節繼續講.
我們主耶穌基督的父.
上上為你們禱告.
因聽見你們在基督耶穌裡的信心.

$^{241}$並向眾聖徒的愛心.
是為那給你們存在天上的盼望.
這盼望就是你們從前在福音真理的道上所聽見的.
這福音傳到你們那裡.
也傳到普天之下.
並且結果增長.
如同你們中間.
自從你們聽見福音.
真知道神恩惠的日子一樣.
正如你們從我們所親愛一同作僕人的以巴弗所學的.
他為我們作了基督忠心的執事.
也把你們因聖靈所傳的愛心告訴了我們.
在第三節至第八節.
某程度上是保羅的土民.
保羅為著教會.
或整個團隊.
為教會的發展獻上禱告.
以及不住的守望.
他特別提到這班弟兄姊妹.
用一個信愛忘的格式.
來描述他們的處境.
首先他們提到他們聽見在基督耶穌裡的信心.
即是他們聽見福音.
他們相信.
並向眾聖徒的愛心.
這個愛心是一種奉獻的表達.
即是他們有信心.
以及有愛心來供養.
福音使命.
即是拆全工作等等.
他們帶著信心和愛心來實踐.
他用信愛忘.
其實他用盼望的篇幅來講解多一點.
接下來他就講到第五節.
他為了給你們存在天上的盼望.
這個盼望就是你們從前在福音真理的道上所聽見.
這福音傳到你們那裡也傳到普天之下.
並且結果增長如同在你們中間.
自從你們聽見福音.
真知道神恩惠的日子一樣.

$^{281}$這句話很長.
是在講福音盼望.
他不是純粹講一種主觀的明天會更好.
這個盼望其實是講到.
因為有上帝與他們同在.
他們領受了福音.
他們的生命因著福音帶來生命一個新的開始.
而這個新的開始或者這種積極的態度.
就一直在他們當中不斷地繁衍.
這班弟兄姊妹不單止聽到福音.
因為他們有來自上帝的盼望.
他們的心裡面很踏實安穩.
以致他們就將這個踏實安穩的訊息帶到身邊的人當中.
還有一個字眼很有趣.
就是第五節裡面講到存在天上的盼望.
什麼意思呢?.
是不是很離地?.
那個盼望.
其實那個具體的意思就是說.
他們的生命在上帝當中是被紀念的.
所以存在天上的盼望.
更加比較口語化的解釋就是很穩陣的盼望.
他們的生命是上帝所保存的.
地上的生命,地上的經歷是不會動搖的.
那個存在天上的盼望.
一個來自上帝的應許.
也都來自上帝的保證.
以致他們的心有力量來實踐他們的信仰.
這個就是信望外.
特別是盼望.
他用多一點的篇幅來講的.
如果你看保羅書信你就知道.
我們有一個叫鏡觀的角度.
通常保羅講多一點的內容.
你就要留意了.
因為在順民夏利裡面.
他講多一點,描述多一點的.
往往都是他們面對群體的挑戰或困難.
而這種盼望的體會.
正正就是歌羅西書.

$^{321}$或者哥羅西門徒裡面.
他們所面對的一種困難.
保羅特別將盼望和福音兩者緊扣在一起.
因為如果離棄了福音訊息.
或者離棄了福音.
他們的生命就沒有盼望了.
又或者有一些教導.
在當時社會裡面.
可能有一些將福音溝淡了.
或者將它異化了.
或者扭曲了.
他們生命裡面的那份很穩陣的盼望感.
就會消失了.
保羅在這裡裡面.
他再強調這個福音是為他們帶來盼望.
而這份盼望令到他們可以積極地用信心和愛心.
來回應上帝的道.
經文裡面再繼續說.
這個盼望或者這個福音是誰帶給他的呢?.
剛才說過.
不是保羅一手一腳將福音帶到這班門徒那裡.
是透過僕人爾巴弗所學的.
這個爾巴弗在聖經裡面描述他的篇幅不是很清楚.
不過可以肯定的就是.
他不像一般的旅行報道.
或者聖經裡面說到的阿波羅那些那麼厲害.
那麼有口才的領袖.
他是一個忠心的執事.
福音的執事.
他未必是很有魅力的人.
不過保羅怎樣教他.
他怎樣從保羅裡面領受這個福音.
他就如實地栽培和幫助弟兄姊妹.
所以在這裡講到爾巴弗是代表了保羅將福音帶到他們當中.
爾巴弗代表著一份正統.
這個是純正的福音.
將這個純正的福音帶到他們那裡.
而且這班人聽到了.
他們還可以結果增長.
這個福音其實本身在他們的群體裡面.

$^{361}$滿有一份的影響力.
不過去到第九節.
你就更加看到保羅他要講的內容是甚麼.
因此我們自從聽見的日子.
也就為你們不住禱告祈求.
願你們在一切屬靈的智慧誤性上滿心知道神的旨意.
好叫你們行事為人對得起主.
凡事蒙他喜悅.
在一切善事上結果子.
漸漸地多知道神.
召他榮耀的權能.
得以在各樣的力上加力.
好叫你們凡事歡歡喜喜地忍耐寬容.
去到這裡開始有些見真章.
保羅要講些甚麼呢.
他再繼續講為這班哥羅西門徒不住的來禱告.
他講我求上帝給你們有聰明誤性有智慧誤性.
這個智慧誤性不是令到自己要墊高自己.
或者令到自己被人覺得你很叻很聰明.
而是這種的智慧是訴諸於認識上帝的旨意.
所以這個智慧誤性是要幫助那些門徒.
更加知道上帝的心意.
第十節裡面他特別的來到他講.
好叫你們行事為人對得起主.
凡事蒙他喜悅.
在一切善事上結果子漸漸地多知道神.
為甚麼保羅要寫這個歌羅西書呢.
哥羅西的教會正面對一個甚麼的挑戰呢.
在後來的篇幅.
在歌羅西書二章八節他就有講.
原來在當時社會裡面.
有一些人用他們的理學和虛空的妄言.
二章八節裡面這樣寫.
基督乃是照著人間的遺傳和世上的小學.
把你們擄去.
簡單來說是一個外邦教會.
這個外邦教會除了接受福音之外.
他們本土裡面有一些文化思想.
有一些思潮.
而這些思潮.

$^{401}$加上一些猶太教的傳統.
大雜亂炒雜碎.
就影響了本來接受很好的基督徒.
他們接受了基督信仰.
但同時一路實踐的時候.
文化有些東西衝擊他們.
而當時最重要衝擊他們的.
就是一個早期的洛斯底主義.
這個洛斯底主義很複雜.
簡單來說就是人要追求一種神秘的經驗.
或者神秘的知識.
這個知識和剛才說的屬靈知識是相對的.
這個神秘的知識令到你可以明白.
蘊含在這個宇宙萬物當中.
是隱藏不可見的.
但要透過一些神秘的操練.
以致到你經歷了.
這樣你人生就可以提升一個層次.
令到你天地共融.
甚至你可以神人合一.
這個說法和我們之前所說的福音完全背道而馳.
因為在聖經裡面說.
這班人已經領受了福音.
已經知道了.
不需要再額外拿一些神秘的經驗.
來補足福音.
保羅其實很認真地說.
他很擔心.
在當代裡面有些早期的律師底.
或者世間裡面.
有些文化裡面對於信仰.
或者對於上帝.
有一些猜測性的.
神秘性的.
甚至引人入性.
但是又摸不著頭腦的.
有一些所謂的神秘學問.
這種神秘學問.
是在教會的群體.
或者在周遭的文化當中.

$^{441}$以至到保羅要求上帝幫助他們.
他們要拿的是真實的智慧.
而這個是明白上帝的智慧.
以至到他們凡事.
行事為人.
要對得起主.
和合本用這個詞.
很傳神.
就是你做的事.
你對不對得起我.
他有這個.
好像領受了別人的恩惠.
但是你怎樣回應.
和合本就是這樣的翻譯.
不過其實原文用的字眼.
不是對得起.
而是配得起.
即是說你所領受的福音.
的恩典是這麼多.
或者你所經歷的是這麼豐富.
你的實踐.
是否配合.
是否對等.
當然人沒辦法用他自己的行為.
來換取上帝的恩典.
不過當我們領受了上帝的恩典的時候.
我們會不會.
因為這個恩典實在太大了.
所以我們的做事方式.
我們願意自己生命.
來作一些更新.
配合我們所領受的恩惠.
所以恩典有多大.
我們回應的方式.
我們是否相應地.
來回應上帝呢.
舉個不得體的例子.
我今天這件西裝.
不是新的.
是新的.

$^{481}$但是這件西裝.
我好好保護.
每年一月才穿的.
因為是童工.
以前送給我的.
穿了這件西裝.
我覺得肩膀挺一點.
你覺不覺得.
肩膀挺好像厲害一點.
我穿了這個.
我沒理由穿褲子.
或者穿一雙白飯裙.
為了配得起.
這個童工送給我的西裝.
我就要穿一條正裝一點的.
得體一點的西褲和皮鞋.
我覺得才合襯.
大家看上去才舒服.
如果我穿這件西裝.
裡面穿一件利工紋.
穿一條短褲.
穿一雙白飯裙.
你覺得你搞什麼啊牧師.
如果你覺得這樣.
就叫配得起.
同樣我們生命裡領受了的恩典.
領受了上帝給我們的恩惠.
我們的行事為人.
我們是不是和我所領受的.
這個恩典的外衣是一致的.
如果一致的時候.
我們生活不會興奮的.
我們不會亂相信其他東西的.
不過寶萊這裡說.
就是正正有些門徒.
有些弟兄姊妹開始.
已經受著那些異端.
或者世間上的小學.
或者引人入勝的.
一些神秘的知識所吸引.

$^{521}$以致他們就沒辦法好好過一個.
配得起上帝恩典的生活.
經文裡再繼續說.
如果你生命裡配得起上帝這個恩典.
你自動而言.
你的生命就會能夠結出果子的.
而且你就能夠過一個.
土神喜悅的那種生活.
而且一直實踐的時候.
你會慢慢明白上帝.
這條路是一直保羅教他們的.
好好的東西你繼續實踐.
你就不要忽然間覺得搞搞不懂意思.
在福音的根基裡再去找一些新的東西.
侵淡了那個福音.
以致他們就反而處於一種卡住的狀態.
十一節.
這個福音不單止是幫助他們成長.
或者幫他們知道上帝.
更加是賦予他們生命的力量.
召他榮耀的權能得以在各樣的力上加力.
好叫你們凡事歡歡喜喜地忍耐寬容.
這裡特別說到.
他們的生命領受了這個福音的時候.
他們生命是有力量的.
相反來說.
他接受了一個淡化了的福音.
他們的生命是軟弱無力的.
他們面對著那個世界.
是不懂得回應的.
你打算跟這個世界看齊多一點.
接近多一點.
融入多一點不同的想法.
但是到最後你的信仰.
賦予你的能力就會越來越弱.
以致到這一班的門徒.
就處於一種左右不是人.
而且有一種迷失的狀態.
這個像不像我們今天面對的那種處境.
說一個尖銳一點的例子.

$^{561}$我知道有些影子們今天覺得.
我們自己的信仰零收.
有時候想追求一些新意.
更加深入.
本來這個動機是很好的.
非常好的.
不過為了要令自己的信仰.
變得更加提升或更加深入.
就會借助坊間有不同的方法.
其中一種就是修禪.
不知道大家有沒有聽過.
就是用佛學的.
佛教的修煉方式.
融入基督教.
即是基督教的禪修.
用這些修禪的方法放進去.
以致令自己的信仰.
好像有些激盪,有些新意.
又或者令自己可以更加安靜.
這些做法.
如果我們不明白修禪的意思.
我們就隨意去覺得.
這件事挺好的.
又或者用一些叫做新紀元.
New Age.
即是做瑜伽.
我做那個動作.
我是想著基督耶穌.
這樣抽取放下來的時候.
在這個做法裡面.
我們出現兩大短板位.
第一.
我們覺得那些東西.
跟我們的信仰一定是相合的.
但我們自己沒有仔細去檢視過.
會不會有些本質上有衝突的東西.
第二就是.
我們經常覺得我們自己所信的信仰.
不夠用.
我覺得不夠用.

$^{601}$那個才是最大的重點.
在歌羅西書裡面.
他講到一個最重要.
這個福音不夠你用.
這個福音就是為你生命帶來盼望的福音.
你經常覺得不夠用.
然後就拿其他東西來補足.
其實不是這個信仰.
或者這個聖經呈現的福音不夠用.
只不過是我們沒有好好地進入.
去實踐.
又或者我們深夜.
我們覺得人家都挺好的.
我們就拿這些東西放入了.
在我們的信仰觀念裡面.
弟兄姊妹.
我真是很誠懇的.
我是很鄭重的.
鼓勵大家.
如果我們成為一個門徒.
我們不要心懷異議.
我們的操練.
在聖經裡面.
或者在基督教的操練裡面.
我們二千多年.
或者算上舊約這麼長的導師.
是足夠我們.
而且是完備的.
來建立我們的生命.
不要再假借其他.
我們看完馬可福音.
又再和佛經來互相交流.
看看其他的佛經.
看看會不會有新的激盪.
那個帶來的不是我們生命更加踏實.
只是會帶來我們生命變得更加軟弱無力.
我們更加迷失.
我們究竟所信的是什麼.
今天是一個炒雜碎的世代.
喜歡的東西就手尖拿來.

$^{641}$但是我們真是自己的本業.
或者我們自己本身的信仰.
我們卻是反而錯失了.
保羅在這裡面.
其實他是在說.
這個聖經.
這個信仰是足夠你使用.
上帝給你福音的知識.
是足夠有餘.
如果你有知識的焦慮.
其實你要走向的不是走向外面.
而是走回到上帝.
走回到聖經裡面來尋索.
走回到基督教的傳統裡面.
來重新思考你們的生命.
我們繼續看一章十二至十四節.
他救了我們脫離黑暗的權勢,.
把我們遷到他愛子的國裡,.
我們在愛子裡得蒙救贖,.
罪過得以赦免.」.
在這裡面其實是一個小結.
他講到這個祈禱.
他再繼續延續.
他說多謝上帝.
多謝上帝甚麼呢.
全部只有一個主詞.
就是耶穌基督.
耶穌基督將這班信徒或者聖徒.
賜予他們有光明的基業.
耶穌基督救了他們脫離黑暗的權勢.
耶穌基督幫助他們.
讓他們遷到愛子的國度當中.
耶穌基督令到這班人蒙救贖.
罪過得以赦免.
全部的主詞只有一個.
就是耶穌基督.
這個以基督為中心.
或者基督所做的.
他羅列所有的工作.
他就要告訴那些門徒.

$^{681}$那些哥羅西的弟兄姊妹.
你的生命是上帝藉著他的愛子耶穌.
做了這些工作.
而不是你有些神秘知識幫你提升.
或者幫你去做到一些.
能人所不能的那些能力.
那些不是你需要重要考慮.
最重要的就是.
讓你從黑暗當中走進光明.
他特別做一個比喻.
就是在光明當中.
同德基業.
對比於在黑暗當中的權勢.
昔日這班弟兄姊妹未接受福音.
他們活在黑暗權勢當中.
被黑暗勢力籠罩.
現在他們已經進入光明.
進入光明你又再進入黑暗做什麼呢.
為什麼又要走回舊時的回頭路呢.
所以他強調.
你生命是光明之子.
一個光明裡面同德基業.
所以你們就可以得著.
上帝所給你的這份這麼榮耀的職事.
更加重要的就是.
在第十一節那裡.
他這裡特別說到.
「召他榮耀的權能得以在各樣的事上力上加力.
好叫你們凡事歡歡喜喜地忍耐.
又叫我們能與眾聖徒在光明中同德基業」.
有一些譯本.
特別是和修本.
他譯了一些和合本沒有譯的字眼.
和修本就譯了.
在光明中分享基業」.
剛才的十二節十一節那裡說.
基業或上帝給你的那份.
為什麼你得到呢?.
是上帝給的.
而你有資格.

$^{721}$配得取這個基業.
是說上帝給你qualify(獲得).
你有資格.
本來以前是沒有資格的.
現在是有資格.
因著基督耶穌.
令到每一個人已經拿到這個資格.
你已經足夠有餘.
去經歷這個生命的豐盛.
既然你又有資格.
你又有能力.
你所接受的是正統.
你還拿這些來做什麼?.
所以在這裡其實.
整個的段落.
其實要強化這班信徒.
沒有信錯主.
不過外面的文化.
或世間的思潮實在太大了.
以至他們開始有些迷失.
在這段經文裡.
有一個很重要的觀點.
之後也會再說.
就是耶穌基督帶來的福音.
不是他們信仰起點這麼簡單.
我們經常覺得信了主.
作為一個起點.
信了主之後.
我們喜歡怎樣就怎樣.
在歌羅西書裡面.
他在說一個重要的神學是.
基督不單是我們的起點.
而是我們的全部.
是我們生命所有豐盛.
我們所信主的東西.
我們覺得好的.
豐盛的東西.
全部都包攬在.
在基督耶穌的生命.
他的恩典當中.

$^{761}$所以我們有耶穌.
我們就不假外求.
足夠有餘過我們自己的生活.
今天這段經文其實是說.
今天你過的信仰生活.
弟兄姊妹.
你覺得你自己的信仰.
足不足夠.
正如剛才所說.
很多弟兄姊妹找外面的東西.
想著充實一下自己的信仰生命.
我們會借助其他的宗教.
其他的修煉方法.
修禪.
讀佛經.
我們不需要踩低別人.
不過在我們基督教裡面.
其實上帝已經很清楚說到.
我們所信的主.
或者這個福音已經足夠完備.
我們信不信得過.
這個所啟示出來的信仰.
基督耶穌已經包攬了.
世界上所有的豐富.
我們信不信得過我們所信的主.
是一個豐富超乎我們想像的主.
還是我們生命裡面.
信了主之後我們還覺得不夠.
我們還要借助.
如來佛祖有些東西幫幫我.
令到我生命更加豐盛呢.
歌羅西書.
他要將我們的焦點.
放回基督耶穌那裡.
最後一個這樣的例子.
有一次我去吃日本菜.
去吃日本菜的時候.
我們吃那些叫做.
Omakase.
吃的時候我覺得那些東西不夠味.

$^{801}$想加一點調味料.
問那個廚師.
你可不可以給我一點調味料.
那個廚師聽到這句話.
馬上黑臉.
你明白的吧.
他覺得你不懂得吃.
你吃Omakase.
你問他拿調味料.
他煮出來.
放上什麼你吃什麼.
你信得過他.
你就全盤接受.
這個味道不是淡口味.
不是他的東西淡.
而是你的舌頭鈍.
是不是.
你不懂得吃.
所以你問他拿調味料.
其實對他是一個冒犯.
所以我們不懂得吃.
我們就想加一些調味料.
令我們覺得合適自己的口味.
我們今天信主.
就像我之前吃我媽媽姐姐.
我們覺得有些東西.
不太合我口味.
我加一些不是基督教的東西.
加一點鹽和醋.
最好是加多一點七味粉.
夠辣.
我們覺得這個合我口味.
我們想拿一個合我口味的信仰.
不過我們沒有想過.
這個是冒犯了我們的主.
沒有想過我們這樣做.
原來吃不到食物的真味.
求主幫助我們.
今天我們所領受的福音.
足夠有餘.

$^{841}$我們生命要經歷上帝的恩典.
去吸真下這個信仰的味道.
好好地去聖經那裡.
經歷上帝所啟示豐富的耶穌.
門徒DNA第一件事.
什麼界定你的身份.
找回信仰的真味道.
\newpage



\section{約翰福音 11:28-57}
\label{sec:eZQl9Arirgg}
\textbf{2026年1月11日崇拜直播｜梁家麟牧師｜相信復活的主｜約翰福音11\_28-57}
\newline
\newline
連結: \href{https://youtube.com/watch?v=eZQl9Arirgg}{\texttt{https://youtube.com/watch?v=eZQl9Arirgg}} ~~~~ 語音日期: 2026-01-14
\newline
\newline
\hyperref[sec:DESvXfy__1I]{\small{< < < PREV SERMON < < <}}
~
\hyperref[sec:index_chronic]{\small{[返順時目]}}
~
\hyperref[sec:RDxRU6iPPzM]{\small{> > > NEXT SERMON > > >}}
\newline
\newline
約翰福音 11:28-57
\newline
\begin{longtable}{cl}
\hline
\hline
章節 & 經文 (和合本修訂版)\\
\hline
11:28 & \begin{tabularx}{0.7\textwidth}{X} 馬大說了這話就回去，叫她妹妹馬利亞，私下說：「老師來了，他在叫你。」 \end{tabularx} \\ \\ \relax
11:29 & \begin{tabularx}{0.7\textwidth}{X} 馬利亞聽見了，急忙起來，到耶穌那裡去。 \end{tabularx} \\ \\ \relax
11:30 & \begin{tabularx}{0.7\textwidth}{X} 那時，耶穌還沒有進村子，仍在馬大迎接他的地方。 \end{tabularx} \\ \\ \relax
11:31 & \begin{tabularx}{0.7\textwidth}{X} 那些同馬利亞在家裡安慰她的猶太人，見她急忙起來，出去，就跟著她，以為她要往墳墓那裡去哭。 \end{tabularx} \\ \\ \relax
11:32 & \begin{tabularx}{0.7\textwidth}{X} 馬利亞到了耶穌那裡，看見他，就俯伏在他腳前，對他說：「主啊，你若早在這裡，我弟弟就不會死了。」 \end{tabularx} \\ \\ \relax
11:33 & \begin{tabularx}{0.7\textwidth}{X} 耶穌看見她哭，並看見與她同來的猶太人也哭，就心裡悲嘆，又甚憂愁， \end{tabularx} \\ \\ \relax
11:34 & \begin{tabularx}{0.7\textwidth}{X} 就說：「你們把他安放在哪裡？」他們對他說：「主啊，請你來看。」 \end{tabularx} \\ \\ \relax
11:35 & \begin{tabularx}{0.7\textwidth}{X} 耶穌哭了。 \end{tabularx} \\ \\ \relax
11:36 & \begin{tabularx}{0.7\textwidth}{X} 猶太人就說：「你看，他多麼愛他！」 \end{tabularx} \\ \\ \relax
11:37 & \begin{tabularx}{0.7\textwidth}{X} 其中有人說：「他既然開了盲人的眼睛，難道不能叫這人不死嗎？」 \end{tabularx} \\ \\ \relax
11:38 & \begin{tabularx}{0.7\textwidth}{X} 耶穌又心裡悲嘆，來到墳墓前。那墳墓是個穴，有一塊石頭擋著。 \end{tabularx} \\ \\ \relax
11:39 & \begin{tabularx}{0.7\textwidth}{X} 耶穌說：「把石頭挪開！」那死者的姐姐馬大對他說：「主啊，他現在必定臭了，因為他已經死了四天了。」 \end{tabularx} \\ \\ \relax
11:40 & \begin{tabularx}{0.7\textwidth}{X} 耶穌對她說：「我不是對你說過，你若信就必看見神的榮耀嗎？」 \end{tabularx} \\ \\ \relax
11:41 & \begin{tabularx}{0.7\textwidth}{X} 於是他們把石頭挪開。耶穌舉目望天，說：「父啊，我感謝你，因為你已經聽了我。 \end{tabularx} \\ \\ \relax
11:42 & \begin{tabularx}{0.7\textwidth}{X} 我知道你常常聽我，但我說這話是為了周圍站著的眾人，要使他們信是你差了我來的。」 \end{tabularx} \\ \\ \relax
11:43 & \begin{tabularx}{0.7\textwidth}{X} 說了這些話，他大聲呼叫說：「拉撒路，出來！」 \end{tabularx} \\ \\ \relax
11:44 & \begin{tabularx}{0.7\textwidth}{X} 那死了的人就出來了，手腳都裹著布，臉上包著頭巾。耶穌對他們說：「解開他，讓他走！」 \end{tabularx} \\ \\ \relax
11:45 & \begin{tabularx}{0.7\textwidth}{X} 於是來看馬利亞的猶太人中，有很多人見了耶穌所做的事，就信了他。 \end{tabularx} \\ \\ \relax
11:46 & \begin{tabularx}{0.7\textwidth}{X} 但其中也有人去見法利賽人，把耶穌所做的事告訴他們。 \end{tabularx} \\ \\ \relax
11:47 & \begin{tabularx}{0.7\textwidth}{X} 祭司長和法利賽人召開議會，說：「這人行好些神蹟，我們怎麼辦呢？ \end{tabularx} \\ \\ \relax
11:48 & \begin{tabularx}{0.7\textwidth}{X} 若讓他這樣做，人人都要信他；羅馬人也要來毀滅我們的聖殿和我們的民族。」 \end{tabularx} \\ \\ \relax
11:49 & \begin{tabularx}{0.7\textwidth}{X} 其中有一個人，名叫該亞法，那年當大祭司，對他們說：「你們甚麼都不知道， \end{tabularx} \\ \\ \relax
11:50 & \begin{tabularx}{0.7\textwidth}{X} 也不想想，一個人替百姓死，免得整個民族滅亡，這對你們是有利的。」 \end{tabularx} \\ \\ \relax
11:51 & \begin{tabularx}{0.7\textwidth}{X} 他這話不是出於自己的意思，而是因他那年當大祭司，所以預言耶穌將為這民族而死。 \end{tabularx} \\ \\ \relax
11:52 & \begin{tabularx}{0.7\textwidth}{X} 他不但替這民族死，還要把神四散的兒女都聚集起來，合成一群。 \end{tabularx} \\ \\ \relax
11:53 & \begin{tabularx}{0.7\textwidth}{X} 從那日起，他們就商議要殺耶穌。 \end{tabularx} \\ \\ \relax
11:54 & \begin{tabularx}{0.7\textwidth}{X} 所以，耶穌不再公開在猶太人中走動，卻離開那裡，往靠近曠野的鄉間去，到了一座城，名叫以法蓮，就在那裡和門徒住下來。 \end{tabularx} \\ \\ \relax
11:55 & \begin{tabularx}{0.7\textwidth}{X} 猶太人的逾越節近了，有許多人從鄉下上耶路撒冷去，要在過節前潔淨自己。 \end{tabularx} \\ \\ \relax
11:56 & \begin{tabularx}{0.7\textwidth}{X} 於是他們尋找耶穌，站在聖殿裡彼此說：「你們認為怎樣，他不會來過節吧？」 \end{tabularx} \\ \\ \relax
11:57 & \begin{tabularx}{0.7\textwidth}{X} 那時，祭司長和法利賽人早已下令，若有人知道耶穌的下落，就要報告，他們好去捉拿他。 \end{tabularx} \\ \\
[1ex]
\hline
\hline
\end{longtable}
$^{1}$《約翰福音》第十一章二十八至四十四節.
《聖經新約》第一百四十四頁.
《約翰福音》第十一章二十八節.
馬大說了這話,就回去暗暗的叫他妹子瑪利亞說:夫子來了,叫你..
瑪利亞聽見了,就急忙起來到耶穌那裡去..
那時,耶穌還沒有進村子,仍在馬大迎接他的地方..
那時,和瑪利亞在家裡安慰他的猶太人,.
見他急忙起來出去,就跟著他,以為他要往墳墓那裡去哭..
瑪利亞到了耶穌那裡,看見他,就俯伏在他腳前,說:.
主啊,你若早在這裡,我兄弟必不死..
耶穌看見他哭,並看見與他同來的猶太人也哭,就心裡悲嘆,又甚憂愁,便說:.
你們把他安放在那裡?.
他們回答說:請主來看..
耶穌哭了,猶太人就說:你看他愛這人是何等懇切?.
其中有人說:他既然開了核子的眼睛,豈不能叫這人不死?.
耶穌又心裡悲嘆,來到墳墓前,那墳墓是個洞,有一塊石頭擋著..
耶穌說:你們把石頭挪開..
那死人的子子馬大對他說:主啊,他現在必是臭了,因為他死了已經四天了..
耶穌說:我不是對你說過,你若信,就必看見神的榮耀嗎?.
但我說者說:是為周圍站著的眾人,叫他們信,是你差了我來..
說了者說,到大聲呼叫說:拉撒路出來!那死人就出來了,手腳裹著布,臉上包著手巾..
耶穌對他們說:解開,叫他走..
主啊,你現在必是臭了,因為你已經聽我,我也知道你常聽我..
主啊,你現在必是臭了,因為你已經聽我..
我們將拉撒路的故事講完..
耶穌和門徒來到伯大利的時候,拉撒路已經下葬了四天的時間..
很多親友,民信從各地趕來吊唁..
馬大聽說耶穌來了,趕出村口去接待耶穌..
瑪利亞沒有跟著來,馬大和耶穌談了一會兒,就趕回家裡,叫瑪利亞出來見耶穌..
耶穌似乎不是很急於進村,進入馬大的屋..
當馬大轉身回家的時候,耶穌和門徒沒有跟著,仍然留在村外..
或者原因很簡單,屋裡擠滿了人,耶穌和門徒不想跟他們一起擠,所以就不進去..
也有另一個更可能的原因,或者馬大知道他妹妹的情緒特別受弟弟的死亡困擾,.
他應該是比較脆弱的人,希望耶穌找機會跟瑪利亞單獨聊天,開解和安慰她..
所以馬大就進屋叫瑪利亞出來,聽到耶穌的召喚,瑪利亞就急忙出來..
不過原來在屋裡陪伴他的親友也跟著他一起出來..
他們以為瑪利亞要去墳墓再去哭一場,就陪他們一起去..
結果整個故事就是瑪利亞仍然沒有機會跟耶穌單獨聊天..
瑪利亞看到耶穌就伏伏在他腳前哭著說:.
「主啊,你若早在這裡,我兄弟必不死.」.

$^{41}$這句話跟之前21節馬大所說的一樣,都是為耶穌的遲到感到遺憾..
他們相信耶穌能夠治病,再嚴重的病都能夠治,.
他們沒有想過耶穌能夠使死人復活..
瑪利亞在哭,伴隨他的人也一起哭..
有些悉經者估計跟瑪利亞一起哭的人,可能是職業的送葬人,送殯人..
當時確實有一個習俗,有一點錢的人辦喪事,.
會請幾個擅長豪哭的婦女在會場豪圖大哭,.
希望製造一個悲情的氣氛..
不過由於喪事已經過了四天,.
喪家不一定會請送殯人這麼長的時間..
並且聖經明說,這裡陪哭的是猶太人,.
就是前面所說從耶路撒冷來的猶太人,.
他們應該是比較有身份地位的,.
應該不會從事送殯人的行業,.
這個行業比較卑賤..
所以我們簡單說,.
應該將他們理解為跟瑪利亞很有感情的親戚..
一哭還哭,耶穌看到現場一片哭聲,.
心裡突然感觸,也哭了一份,.
在35節說耶穌哭了..
先不要用字面解釋,直接說,.
如果真是說耶穌哭,那耶穌為什麼哭呢?.
這是這段經文一個懸疑性的問題,.
不少人信理成章地推論,.
反映了耶穌是一個會哭的正常人,.
不過耶穌具有人性能夠哭,.
不等於在這個場合,在某個場合就要哭,.
我有人性,但我現在不哭,我沒有哭的感覺..
瑪利亞哭,其他人哭,.
是因為他們失去了親人,失去了好友,.
傷心難過,.
耶穌早就知道這個噩耗,.
所以不存在驚嚇的問題,.
不是經文噩耗,.
他早就知道,不會驚慌失措,.
並且耶穌是故意遲到,.
所以也不存在遲到,錯失了救治的時機,.
而有的懊悔,遺憾的困擾,.
沒有這些問題,耶穌沒有這些困擾的問題..
還有,你記得前面的經文說,.

$^{81}$耶穌聽到拉撒路死的訊息,.
最初的反應是,.
知道拉撒路死亡,而他不在現場,.
他為此感到高興,.
甚至他說,我沒有在那裡就歡喜,.
這是為你們的緣故,.
好叫你們相信,.
他歡喜,因為拉撒路死亡是一個展現他生命主權的機會,.
如果前面的經文說,他很高興,.
為什麼又要哭呢?.
耶穌早就定義,此行是要讓拉撒路復活的,.
並且宣告,生與死是由他輸了算,.
你記得我上次說的,.
所以,成竹在胸的他,.
根本沒有哭的動機,沒有哭的必要,.
耶穌的哭,肯定不是為了傷心難過,.
不是懊悔,不是無助,不是絕望,.
有人說,耶穌在喪禮的場合,.
哭是合適的行為,.
這個配合環境,.
又和身邊的人的情緒反應看齊,.
不過,如果這樣解釋,.
就是禮節的問題,.
在某個場合,有合宜的表現,.
我要說,這些合宜的表現,.
和感情沒有什麼關係,.
職業宋座的人哭得夠大聲,.
不見得他們很有情感,.
這是職業需要,.
我在參加喪禮的時候,.
勉強擠幾滴眼淚出來,.
湊合現場的氣氛,.
也不一定是真情流露,.
對不對?.
聖經確實有教導,.
與喜樂的人同樂,與哀哭的人同哭,.
這是羅馬書十二章五節的教導,.
注意,這是保羅的說話,.
是保羅對門徒的要求,.
不是對耶穌的要求,.

$^{121}$不是說耶穌在相同的場景,.
就要和我們有相同的情緒反應,.
沒有這樣的事,.
正如過去的事件,.
我們沒有看到耶穌有流露和身邊的人,.
相同的情緒反應的做法,.
門徒軟弱,他就軟弱,.
門徒驚慌,他就驚慌,.
正義表達,他認同他們,.
和他們同行,.
在東宇的加利利湖上,.
耶穌有沒有和船的門徒一起害怕?.
沒有,.
在門徒為缺乏食物感到擔憂的時候,.
耶穌有沒有表達相同的擔憂?.
沒有,.
耶穌總是有一些超然,.
告訴他們,.
不用怕,不用困擾,.
是不是?.
我們看看,.
耶穌哭了這個哭字,.
這個哭字,.
和前面馬尼亞和眾人的哭字,.
其實不是同一個希臘文,.
這個字不一定是指哭泣,.
更多是指情緒受困擾,.
有激動,.
甚至是憤怒的反應,.
這是一個複雜,.
複合的情緒反應,.
有悲傷,也有憤怒,.
英文的翻譯就是,.
Jesus was troubled,.
Troubled,.
就是耶穌被干擾,困擾,.
很講究困擾,.
面對人間的苦難,.
面對人生的悲苦,.
我們都會有類似的複雜的反應,.

$^{161}$就好像你不久前見到那場大火,.
你除了感覺到傷心難過之外,.
我不知道你有沒有那個反應,.
很自然的反應是憤怒,.
看到死了那麼多人,.
有激動,很憤怒的感覺,.
有沒有搞錯,.
怎能讓這樣的災難發生呢?.
這不應該只是天災,.
一定有人禍的成份,.
一定有些人做錯了,.
要負責的,.
是政府部門,是相關的單位,.
總之不應該是這樣,.
我不知道你會不會很自然,.
就有這個情緒的反應,.
不僅僅是難過,.
也有激動,有憤怒,.
網上甚至有人批評消防員救火不力,.
或者用了一些不恰當的救火方法,.
這些批評指責不一定合理,.
不過是人之常情,.
somebody must be blamed,.
有些人要負責的,.
關鍵的是,.
who is the easy scapegoat,.
誰是容易被人套上的代罪羔羊,.
總之這是很自然的反應,.
耶穌為什麼會有情緒的波動,.
接下來的經文,.
我們看到約翰記載這件事的意圖,.
周圍的猶太人看到耶穌的情緒反應,.
就說三十六節,.
你看他愛這人是何等懇切,.
耶穌的激動表達了他對死者和死者家屬的同情憐憫,.
他關心,他在乎,他介意,他愛他們,.
當然,猶太人在這裡誤會了耶穌,.
他們以為耶穌是為死了的拉撒路激動,.
可能是愛他英年早逝,.
不過耶穌其實不是,.

$^{201}$對於很快將他復活過來的人,.
有什麼好哭的,.
死了一會兒,只是碎了,有什麼值得哭,.
耶穌是為當時在場的生人去慨嘆,.
特別是為在他面前的瑪利亞,.
耶穌憐憫他們當下那種困苦無助,.
人生無常,生命脆弱,.
特別是當死神召喚的時候,.
即使你再財雄勢大,叱tak風雲的人,.
亞歷山大大帝,拿破崙,.
都只能乖乖跟著走,.
何況那些無權無勢的小老百姓,.
除了哭之外,根本沒有任何頑抗的能力,.
在電視新聞裡面,.
看到加薩地帶被以色列轟炸之後的死亡的寢席,.
生還者抱著死去的家人,.
嚎啕大哭的時候,.
你都有那種強烈的,測然不忍的心,.
很無能的,很無助的感覺,.
覺得一切都沒有什麼好做,.
耶穌當然有逆轉生死契闊的能力,.
他們人間只是為了要打敗罪和死的權勢,.
耶穌成就的救恩,.
讓保羅可以豪情的高喊,.
死呀,你得勝的權勢在哪裡,.
死呀,你的毒鉤在哪裡,.
講那麼多前書十五章,五十五節,.
不過,耶穌不僅僅是要解決死的權勢的問題,.
他要解救那些在死的權勢之下,困苦無助的人,.
他的焦點從來不在問題,.
是在人那裡,.
他要拯救的人,.
總是投注他的感情,.
在醫治幫助他們之前,.
聖經總是說,.
情了持心,先有感情,然後才有行動,.
耶穌看到那些缺乏人生方向,.
不知道為了什麼而活的人,.
聖經說,就憐憫他們,.
先不說拯救他們,幫助他們,.

$^{241}$先說他憐憫他們,.
感情先於行動,.
耶穌在解決人死亡問題之前,.
憐憫在死亡面前困苦無助的人,.
包括他眼前的瑪利亞,.
無情人間有情上帝,.
上帝有能力拯救世人,.
但是他更關愛在最終的世人,.
我們大家最熟悉,.
很喜歡,經常背誦,經常說的,.
《約翰福音》三章十六節,.
整個耶穌基督度身的救贖行動,.
大前提,原因是上帝愛世人,.
上帝愛世人,上帝在乎,上帝關心,.
上帝關心比上帝拯救更讓人詫異,.
關他什麼事?為什麼他會關心?.
我們這些無人生命,.
死十個都不厭多,.
為什麼他在乎,他關心?.
這段聖經兩次提到耶穌在喪街現場.
心裡悲嘆,33節38節,.
33節更加說,心裡悲嘆有甚憂愁,.
充分表達耶穌憐憫關愛嘆息,.
在瑪利亞人的帶領下,.
耶穌終於來到拉薩路墳墓前,.
這是用一個山洞作為墳墓,.
有塊石頭在墓門上擋住,封住,.
耶穌開口吩咐,39節,你們把石頭挪開,.
但提醒主,他現在必是臭了,.
因為他死了已經四天了,.
在耶穌堅持下,你們就把石頭挪開吧,.
耶穌隨即做了一個祈禱,41到42節,.
我知道你常聽我,.
但我說者說,是為周圍站著的眾人,.
叫他們信你是差了我來..
與其說是一個祈禱,不如說是一個宣告,.
從一開始,耶穌不是祈求天父聽他祈禱,.
使拉薩路復活,而是指出天父早已聽他祈禱,.
正如過去天父從來都聽他祈禱一樣,.
過去在講約婚姻時,我都兩次說過,.

$^{281}$耶穌強調他總是遵從天父的旨意,.
他做的決定從來沒有離開過天父的決定,.
我們不要因此錯誤理解,.
以為耶穌只是跟班,只是弟弟,只是一個傀儡,.
完全沒有自己的想法,只能夠執行天父的旨意,.
他的地位肯定在天父之下,.
耶穌的意思其實是,他跟天父是不可能有對立的想法,.
天父的想法就是他的想法,他的想法就是天父的想法,.
由於我們在人間是看不到天父的想法的,.
你只能看到耶穌基督的想法,看到耶穌基督的做法,.
你就專注在耶穌身上,從耶穌身上,.
充分完全百分百去了解天父的心意,.
從來沒有人見過上帝,.
只有他傅懷禮的獨生子將他表明出來,.
耶穌基督是可見的上帝,那個reveal God,.
我們看耶穌就完全知道什麼叫天父的旨意,.
你不能想像耶穌的想法會跟天父有任何的差距,.
在這裡耶穌的宣告同樣表明,天父總是聽他的話,.
但不是因為他夠兇,所以天父聽他的話,.
而是因為根本上他的想法跟天父的想法是沒有任何差異的,.
所以他的想法就是天父的想法,.
對人來說,祈禱主要有兩個功能,.
第一,我們向上帝講出我們的心意,.
我想我兒子入藍白,入藍白,入藍白,.
即是祈求上帝兌現,.
第二,我們在禱告中學習尋求明白上帝的心意,.
所以是雙向的,我們讓上帝明白我們的想法,.
我們也想明白上帝的想法,.
但耶穌這番說話既不是祈求天父隨聽他的心意,.
他的心意就是天父的心意,根本不用求,.
也不是想藉著祈禱去明白天父的心意,.
也不需要的,他怎會不知道天父的心意呢?.
所以你要問這是不是祈禱呢?.
不如說耶穌跟天父說兩句話,聊聊天,.
這不算是一個嚴格的祈禱,.
不過耶穌指出這個聊天有特別的功能,.
就是講給在場的人聽,.
但我說者說,是為周圍站著的眾人,.
叫他們信是你差了我來,.
耶穌要證明他跟天父的緊密關係,.

$^{321}$再藉以證明他的特殊身份和地位,.
就是他的神聖來源,.
他是被父差來證明他的神聖,.
證明他的神聖來源,也證明他的神聖,.
在對話裡面,耶穌沒有求天父賜他洗人復活的權柄和能力,.
因為復活和生命都在他那裡,.
他本來就有這個權柄和能力,不需要另外賜予的,.
接著耶穌就展開復活的行動,.
他大聲向墳墓呼叫說,拉薩路出來,.
沒有儀式,沒有咒語,不用做什麼法事,.
就是一個吩咐,一個命令,.
不需要用什麼方法步驟促成,.
純粹是我說了話,我的旨意就兌現,.
正如耶穌從來都不是用法力去止住風浪,.
而是命令風和雨止歇,.
如果舊約,我們會找到一些洗人復活的故事,.
以利亞和以利沙都有洗人復活,.
不過兩個舊約才知道,.
要一次幾次伏在屍體身上,.
然後才看到死者還陽,.
以利亞就伏了三次,.
以利沙更加要眼對眼,口對口貼近死者,.
你看到聖經的技術是儀式感十足,.
耶穌不需要做什麼洗拉薩路復活,.
他只是命令已經死了的拉薩路復活過來,.
正如我以前經常說,這是橫斑辦的,.
耶穌要你活,你死了也不肯跟我爬起來,.
我要你活,你就不可以死,.
耶穌只是要表現他的主權,.
詩就這樣成了,.
全身纏著果施布的拉薩路從墳墓裡面自己走出來,.
樣子一定是很嚇人的,.
耶穌再吩咐人幫忙解開纏在他身上的布,.
讓拉薩路自由活動,拉薩路就這樣活過來了,.
讓死人復活的技術就停在這裡,.
沒有再提到馬大和馬利亞的反應,.
沒有提到耶穌就此復活事件,.
對門徒,對周圍圍觀的人所做的一些總結性的教導,.
彷彿耶穌做完這件神蹟就立即收工,.
唯一記下的後果是,.

$^{361}$圍觀的人因著木島這個驚心動魄的神蹟就相信了耶穌,.
45節,那些來看馬利亞的猶太人看見耶穌,.
看到耶穌所作的事就多有信他的,.
約翰跟著記述的都是那些從耶路撒冷來探望馬大馬利亞的猶太人,.
他們將這個木島的復活事件回到耶路撒冷之後,.
告訴那些有資格參加工會的會議的法利賽人,.
結果這個故事很快就傳到大祭司該阿發的耳中,.
猶太的領袖在工會開會討論這件事,.
他們對耶穌這個神蹟非常焦慮,.
47,48節,只能行好些神蹟,我們怎麼辦呢?.
若這樣猶他,人人都要信他,羅馬人也要來奪我們的地土和我們的百姓,.
復活的神蹟帶來兩個風險,.
第一,很多人會因神蹟信耶穌,.
我見到我也信,能夠使死人復活,為什麼可以不信呢?.
第二,耶穌或者會掀起一個社會的風潮,造成騷亂,.
惹來羅馬帝國出兵鎮壓,.
如果羅馬人真的出兵的話,對猶太人大禍臨頭,.
當然第二個風險只是推測,.
可能是這班猶太教的領袖危言聳聽,.
不過工會的猶太人顯然認同這個推論分析,.
在診斷之後處方,.
該亞發大祭司一錘定音,.
說:四十九五十節,你們不知道什麼,.
獨不想一個人替百姓死,免得通國滅亡,.
就是你們的益處..
既然耶穌是一個危險人物,會讓整個民族陷在死地,.
為了消除讓所有人死的風險,就先整死這個人,對不對?.
言下之意,防範於未然,.
不管耶穌是不是真的該死,.
不管耶穌是不是真的有值得死的罪狀,.
為了大眾的好處,.
這是必須被犧牲的棋子,.
一定要將耶穌殺死,快狠準,.
該亞發果斷的處方,統一了工會的思想,目標既定,.
聖經就說他們開展謀害耶穌的行動,.
讓死人復活,成了讓耶穌上十字架的關鍵事件,.
這就是後話..
整個故事就說完了,我們做些反省和應用..
其實張聖經主要是突顯兩個重點主題,兩個概念,.
第一是復活,第二是相信..

$^{401}$先說復活,無論是記載這個故事的約翰,.
還是故事裡面的耶穌,.
都是不停地將焦點集中在.
這是一件非比尋常的死人復活事件..
約翰記載的手法,無論是門徒和耶穌的對話,.
因為馬大和瑪利亞和耶穌的對話都是高度選擇性的,.
除了耶穌死人復活的行動,.
和這件事直接相關之外,.
其餘的對話一概沒有記載..
在整張聖經裡面,多瑪講了一句,.
其他門徒講了兩句,瑪利亞講了一句,.
馬大講話最多,講了五句,.
所有這些說話幾乎都是一樣的通,.
消極性,表達對拉撒路死亡的難過,.
相信就算耶穌去,都再無可為之處..
所以,所有他們的說話都只反映一個事實,.
沒有人在事先相信死人復活,.
也沒有人期待死人復活..
門徒和多瑪覺得拉撒路都死了,.
再去伯特利都是多餘,.
還冒生命危險,他們已經餓了,不想去..
然後,馬大和瑪利亞都只是相信耶穌能夠治病,.
沒有相信耶穌能夠使死人復活,.
所以都慨嘆耶穌來遲了..
耶穌在沒有人相信和期待的前提下,.
施行這個復活的神蹟..
在故事裡面,.
耶穌刻意延後前往伯特利的行程,.
直拖到他宣佈拉撒路的死訊,.
說明他從一開始就不打算行醫治的神蹟,.
只是打算行一個復活的神蹟..
他對馬大做的偉大的宣告,.
我們上次重點說的,.
復活在我,生命也在我,.
宣示了他掌控人的生與死,.
他不僅有能力使人復活,.
他更是生命的主..
所以復活是整個故事的主題..
在神蹟之後,.
約翰專注記載復活事件所帶來的影響,.

$^{441}$其他不相關的就省略,.
沒有提到馬大馬利亞在看到拉撒路復活後有什麼反應,.
全部不用說了..
所以復活是一個很清楚的主題..
然後我們說相信,.
相信什麼?.
首先當然是相信耶穌是復活的主,.
生命的主,他有能力,有主權..
接著就是相信耶穌的神聖來源和神聖地位,.
他與天父有緊密不可分的關係,.
他是由天父差來的,.
我們剛才說了,.
他在行神之前對天父的祈禱宣告,.
因為你已經聽我的,.
我知道你常聽我的,.
但我說者說,.
是為周圍站著的眾人,.
叫他們信是你差了我來,.
他的神聖來源,.
神聖地位..
最後,.
因著確認耶穌的神聖身份,.
就將榮耀歸給耶穌..
耶穌在聽見那些報信者帶來拉撒路病危的消息之後,.
在第四節,.
我們上次說的,.
回應說,.
這病不止死,乃是為上帝的榮耀,.
叫上帝的兒子得榮耀..
在拉撒路的墓前,.
耶穌對馬大這樣說,.
我不是對你說過,.
你若信,就必看見上帝的榮耀嗎?.
第四十節,.
注意,這兩處的聖經,.
耶穌已經沒有隱藏,.
沒有客氣,.
在說到讓上帝得榮耀的同時,.
他竟然宣稱,.
讓上帝的兒子,即是自己,得榮耀..

$^{481}$如果他純粹是人,.
這句話是極其褻瀆性的,.
我們誰也不說,.
我希望你們榮耀我,你們想死嗎?.
只是榮耀上帝,沒理由,.
榮耀上帝的同時,榮耀梁家倫,.
即是你立即趕我出會,對不對?.
不可能的,.
他竟然大肆說,.
父上帝得榮耀,.
叫他的兒子,即是我,.
都得榮耀,.
耶穌得榮耀,.
就是父上帝得榮耀,.
父上帝得榮耀,.
就是耶穌得榮耀,.
耶穌和父上帝完全平等,.
同樣是世人榮耀的對象..
在宣告拉薩路的死訊時,.
剛才也引述了,.
耶穌說了一段奇怪的話,.
祂說祂自己不在現場,.
感覺到歡喜,.
祂對門徒說了十五節,.
我沒有在哪裡就歡喜,.
只是為你們的緣故,.
好叫你們相信,.
耶穌藉著死人復活的神蹟,.
讓門徒相信祂,.
當然門徒早就相信了祂,.
否則他們也不會成為門徒,.
但是起初的相信,.
和見到復活後的相信,.
應該是有不同的,.
過去門徒,.
或者是接受耶穌的教訓,.
隱隱相信耶穌是舊約所應許的彌賽亞,.
但是耶穌堅持,.
祂最重要的身份,.
是天父上帝的兒子,.

$^{521}$他是聖子,.
是和父上帝有最特殊和親密的關係,.
並且是完全平等,.
不可以分開的,.
就這一點,.
門徒其實聽來是最多半信半疑,.
或者不是不信,.
而是不完全明白,.
新的經歷,.
新的知識,.
帶來深刻的相信,.
即使復活的行動,.
耶穌確立了他的神聖來源和神聖地位,.
他是復活和生命的主,.
耶穌既然能夠克勝死亡,.
連死亡也征服不了他,.
連死亡他也不用怕,.
一個真是不死的小強,.
他真是天下無敵,.
連死亡也不怕,.
天下無敵,.
整張聖經也會有,.
整張聖經也圍繞著相信這個主題,.
耶穌期望馬大和瑪利亞相信他,.
猶太人在目睹神蹟之後,.
有很多人相信了他,.
耶穌耶路撒冷公會的領袖擔心,.
復活的神蹟會使很多人相信耶穌,.
至於相信耶穌會有什麼後果,.
這個就回到第一個主題,.
耶穌的偉大宣告所帶來的應許,.
復活在我,生命也在我,.
信我的人,雖然死了,也必復活,.
兩個主題互相緊扣,.
我們看到復活,我們就相信他,.
我們相信他,我們也同時能夠復活,.
相信復活的主就得著復活,.
相信生命的主就得著生命,.
兩個很簡單的主題,.
大家都很熟悉,.

$^{561}$今天也沒有太多時間將它詳細的引申,.
復活的盼望是我們每個基督徒都知道的,.
不過我希望它真的成為了我們信仰生命的重要部分,.
重要的元素,.
正如我以前說的,.
我接受了信徒只聽了三個禮拜的道理,.
他們最深刻記住的是,.
這是一個復活的主,.
會為他們帶來復活的生命,.
我希望我們每個基督徒都確認這一點,.
怎樣證明我們確認?.
今天只能說簡單一點的,.
灑脫一點,.
從容一點,.
我們的格局大一點,.
不要婆婆媽媽,.
不要孜孜不孜,.
不要太執著,.
不要為人生的小成小敗太糾纏,.
我們相信好戲在後頭,.
keep the fork,.
還有一條筆,.
拿著叉子,.
不要這麼快不愛它,.
好事在後頭,.
復活是我們的盼望,.
生命不固定在人間,.
這幾十寒暑,.
我們有復活的盼望,.
第二,相信,.
相信其實是最難說的,.
任何對相信,對信,.
說任何事都會偏離相信,.
我們要信得牢固一點,.
信得硬淨一點,.
我們的念力要強一點,.
其實是將信變成一個行為,.
信其實要和行為相對,.
即是無依無靠,.
你經常想做的事有渣拿,.

$^{601}$信就是完全沒有渣拿,.
這就叫信,.
那怎麼辦呢?.
我怎知道我相信呢?.
我說我相信,.
但說著說著怎知道都是呢?.
真的,.
信本身沒有什麼好證明的,.
我們永遠是後來,.
事後孔明的,.
藉著我們的關係,.
藉著我們的行動表現,.
你看到我們相信,.
是不是?.
你真的很愛耶穌,.
真的很愛,很愛祂,.
我和祂可以有那種黏黏的關係,.
我生命裡面真的離不開祂,.
這個沒有多少外在的表現,.
是你自己知道的,.
你自己知道,.
它是不是真的matter,.
是不是真的有關係,.
你會知道,.
然後看看我們生活的行為,.
是不是活出一個相信的人生,.
在各樣的事上表明我們是上帝的用人,.
我們的生活,.
是不是真的燃出一個信仰的生命,.
最後我自己也有些不滿意,.
覺得很抽象,.
不過,復活和相信是整張聖經的主題,.
也希望是我們每一個基督徒,.
我們生命的兩大主題,.
我們旺朕,.
我不管我們今天是什麼題目,.
我希望我們更加認定有復活的盼望,.
我們更加相信我們的主,.
多謝大家收看,我們下期再見.
\newpage



\section{歌羅西書 1:15-23}
\label{sec:RDxRU6iPPzM}
\textbf{2026年1月18日崇拜直播｜陳冠成牧師｜門徒DNA:連於元首基督｜歌羅西書1\_15-23}
\newline
\newline
連結: \href{https://youtube.com/watch?v=RDxRU6iPPzM}{\texttt{https://youtube.com/watch?v=RDxRU6iPPzM}} ~~~~ 語音日期: 2026-01-18
\newline
\newline
\hyperref[sec:eZQl9Arirgg]{\small{< < < PREV SERMON < < <}}
~
\hyperref[sec:index_chronic]{\small{[返順時目]}}
~
\hyperref[sec:sjw5HKR2iIk]{\small{> > > NEXT SERMON > > >}}
\newline
\newline
歌羅西書 1:15-23
\newline
\begin{longtable}{cl}
\hline
\hline
章節 & 經文 (和合本修訂版)\\
\hline
1:15 & \begin{tabularx}{0.7\textwidth}{X} 愛子是那看不見的神之像，是首生的，在一切被造的以先。 \end{tabularx} \\ \\ \relax
1:16 & \begin{tabularx}{0.7\textwidth}{X} 因為萬有都是在他裡面造的，無論是天上的、地上的，能看見的、不能看見的，或是有權位的、統治的，或是執政的、掌權的，一概都是藉著他為著他造的。 \end{tabularx} \\ \\ \relax
1:17 & \begin{tabularx}{0.7\textwidth}{X} 他在萬有之先；萬有也靠他而存在。 \end{tabularx} \\ \\ \relax
1:18 & \begin{tabularx}{0.7\textwidth}{X} 他是身體（教會）的頭；他是元始，是從死人中復活的首生者，好讓他在萬有中居首位。 \end{tabularx} \\ \\ \relax
1:19 & \begin{tabularx}{0.7\textwidth}{X} 因為神喜歡使一切的豐盛在他裡面居住， \end{tabularx} \\ \\ \relax
1:20 & \begin{tabularx}{0.7\textwidth}{X} 藉著他，神使萬有與自己和好，無論是地上的、天上的，都藉著他在十字架上所流的血促成了和平。 \end{tabularx} \\ \\ \relax
1:21 & \begin{tabularx}{0.7\textwidth}{X} 從前你們與神隔絕，心思上與他為敵，行為邪惡； \end{tabularx} \\ \\ \relax
1:22 & \begin{tabularx}{0.7\textwidth}{X} 但如今，他藉著他兒子肉身的死，已經使你們與他自己和好了，把你們獻在他的面前，成為聖潔，沒有瑕疵，無可指責。 \end{tabularx} \\ \\ \relax
1:23 & \begin{tabularx}{0.7\textwidth}{X} 只要你們持守信仰，根基穩固，堅定不移，不致動搖，離開了你們從前所聽見的福音的盼望；這福音也是傳給天下一切被造之物的，我—保羅作了這福音的僕役。 \end{tabularx} \\ \\
[1ex]
\hline
\hline
\end{longtable}
$^{1}$主席.
《聖經》第十五節第279至280頁.
「愛子是那不能看見之神的像,是守生的,在一切被做的以先,因為萬有都是靠他做的..
無論是天上的,地上的,能看見的,不能看見的,或是有為的,主治的,執政的,掌權的,一概都是藉著他做的..
又是為他做的,他在萬有之先,萬有也靠他而立,他也是教會全體之首,他是原始,.
是從死裡首先復生的,使他可以在凡事上居首位,因為父喜歡叫一切的豐盛在他裡面居住..
既然藉著他在十字架上所流的血,成就了和平,便藉著他叫萬有,無論是地上的,天上的,都與自己和好了..
你們從前與神隔絕,因為惡行,心裡與他為敵,但如今他藉著基督的肉身受死,叫你們與自己和好,都成了聖潔,沒有瑕疵..
無可責備,把你們引到自己面前,只要你們在所信的道上恆心根基穩固,堅定不移,不致被引動失去..
福音的盼望,這福音就是你們所聽過的,也是傳語普天下萬人聽的.我保羅也作了這福音的執事..
願神賜福上天,讀他話語並且遵行的人..
弟兄姊妹平安..
我們說,如果要說歌羅西書對基督教信仰的其中一個最重要的神學貢獻,毫無疑問就是剛才所讀的這段基督論,也就是來自剛才15至20節,很多時候被稱為基督頌詞的經文..
他幫我們更深入了解耶穌基督的真正身份,和他為我們進行了什麼事情..
事實上,如果耶穌只是在第一世紀的其中一個最偉大的教師或先知,又或者祂只是一個最偉大的受造之物,更甚如果祂的犧牲只是一個最偉大的人寫幾個人的精神,那怎麼可能祂在十字架上的死會產生一個這麼不可思議的大能?.
不單止使得全人類的罪都得到赦免,甚至整個宇宙的所有受造物都因為耶穌犧牲的緣故能與上帝和好..
換句話說,我們身為凡人,你跟我的善行就算再多再偉大,我想都不可能有這個效果,頂多你能夠為自己贖罪,救你自己,你救不了所有人..
所以唯一合理的解釋就是,耶穌基督擁有一個超越了萬有的身份和本質,所以祂才有資格和能力完滿解決一切受造物和上帝之間的破壞關係..
事實上,保羅就是透過剛才玉峰弟兄所讀的這段經文,顯示了耶穌基督這種超越性,.
為了幫助大家盡快掌握經文脈絡,其實在廣道大綱已經印了剛才所讀的15至23節的分段,包括如何從創造的角度看到基督的超越,.
然後保羅幫我們從救世的角度看到耶穌的超越,到最後因為耶穌這兩方面的超越的時候,我們作為他的跟隨者,作為信徒,應該如何合而回應..
我會先按著這個脈絡來將經文講解,不過過程中因為涉及很多保羅整個神學思路的了解,所以大家要給一點耐性..
我想這一篇也是我來往朕這麼久,這一篇最神學的道,我們一起禱告求上帝幫助我們去理解,同心祈禱..
聖靈求你,內署在我們的生命當中,讓我們有從你而來的悟性和智慧,可以認識耶穌基督,祂在這個世界,這個宇宙,在我們生命當中的作為,.
以致我們體會到耶穌那種超越,並且我們真的能夠合而的,以相稱的言行來配合跟隨這位救主,這位創造的主..
求聖靈與我們同在,誰聽我們的禱告,奉基督耶穌你的名字來祈求,阿們..
你看到第十五至一開始,保羅首先從創造的角度來解釋為什麼耶穌基督有這種超越性..
你看到這段經文出現了很多跟創造和受造有關的字,重複出現..
一方面揭示了耶穌基督在創造裡面的重要身份和角色是什麼..
保羅首先開宗明義,愛子就是那看不見之神的像..
他不單單表明了耶穌基督與上帝的關係就是子與父的關係,是如此親密..
他更加強調,耶穌基督不是純粹複製了父神的形象,然後在世間顯現,耶穌不是上帝的複製品,.
耶穌也不是父上帝形象的一個投影或投射,保羅說,愛子就是神的像,神的形象..
如果你看NYV很清楚,The Son Is,就是說耶穌基督的本質就是上帝的本質,.
所以他能夠將父上帝的形象甚至本質正確無誤地顯明出來,.
人透過這位看得見的耶穌,就可以認識那一位看不見的上帝,.
也就像約翰福音所說,保羅單單用這一句就揭示了耶穌基督擁有神聖的身份..
他繼續說,愛子與受造物的關係到底是怎樣的呢?.
這裡和學本說,愛子是守生的,是一切被造的以先..
和學本這句亦有時會讓人誤會,耶穌基督是一切受造物裡第一個出生的..

$^{41}$其實新譯本也印在秩序表裡,他的翻譯比較準確,.
其實他是說愛子是首先的,是一切受造的之上,.
其實是要強調耶穌基督的地位超越了萬有的超越性..
十六字保羅繼續解釋,為什麼耶穌擁有這種超越的身份,原因是什麼?.
原文其實是用了三個前置詞或介詞,英文就是preposition,.
來表達到底耶穌基督在創造過程中扮演了什麼重要角色..
和學本翻譯不出來的,特別是第一個他翻譯不到,.
首先保羅這裡說,萬有其實是在他裡面而造,英文是in him,.
和學本只是翻譯成靠他,其實不太準確,應該是在他裡面而造,.
強調萬物的受造必須在耶穌裡面發生,不可以離開耶穌基督..
你看到保羅很霸氣地宣告,簡單說就是沒有耶穌基督的參與,.
創造根本不能夠成事,就是這個意思..
耶穌基督就是整個創造的關鍵,key person..
第二,保羅說萬有都是藉著他造的,和學本翻譯到,.
through him,就是說耶穌基督在創造過程中是一個中介的角色,.
上帝透過他來完成整個創造工作,.
特別從下文我們可以看到,整個創造工作什麼時候會完成呢?.
就是當耶穌基督在十字架上成就了救贖的工作之後,.
整個創造工作才能夠順利完成..
第三,耶穌的超越性在於萬有是為他而造,for him..
耶穌基督不單單是創造的中介人,.
他更加是創造本身的終極目的和目標,.
就是說一切受造之物,無論天上的地上的能看見不能看見,.
有為自主的執政的掌權的,.
所有無論在物質界的受造物,.
又或者靈界的事物,或者權勢,.
通通都有他本身的目的和意義而被造,.
都是要為了耶穌基督而造,.
就是說他們都是必須為耶穌基督效力,.
他們必須朝向耶穌基督的榮耀的最終結局前進..
簡單來說,為什麼耶穌基督可以超越萬有?.
因為他本來就說,他不單是受造物,.
他根本是創造主,.
他是整個宇宙創造的核心,.
萬物不是無緣無故被奪進世間,.
而是在他裡面被造,.
是藉著他被造,也是為了他而造..
第十七節,保羅繼續說耶穌基督的超越性,.
他說「他在萬有之先,萬有也靠他而立」,.
這次不是說地位上,是先,.

$^{81}$是說時間上,耶穌基督是早於受造物先存在世上,.
沒有任何受造物,耶穌基督已經存在,.
和上帝一樣..
並且在萬物創造之前,.
耶穌基督早就已經存在,.
他不單單參與整個創造工作,.
萬物創造之後,他如何維繫,.
如何繼續存活,.
經文說,都是有賴耶穌基督繼續之後的護佑工作,.
使萬物得以有秩序,.
按上帝的旨意,存活,維繫下去..
來到這裡,你會看到,.
保羅整個創造的神學思路,.
萬有不單單是剛才所說的,.
看下文,你會看到,.
保羅心目中萬有,包括教會這個群體,.
所以第十八節,.
保羅不是突然間轉了話題,.
由宇宙萬物突然間跳到教會,.
其實他是從創造的角度,.
來理解教會的出現,教會的存在,.
其實根本就是上帝創造計劃的一部分..
保羅用了一個我們很熟悉的比喻,.
頭和身體,.
用這個關係來比喻耶穌基督,.
祂就是教會生命的源頭,.
教會的設立,教會的存活,.
甚至往後的發展,.
完全是出於基督的主宰,.
祂的供應和帶領..
而頭和身體這個比喻,.
也強調耶穌基督的地位超越在教會上面,.
祂之所以被稱為教會的元首,.
是因為經文說,.
祂是第一個從死人中復活的,.
意思是因為耶穌基督從死裡復活,.
祂開創了一條新的道路,.
就是一條可以復活得性死亡的道路,.
耶穌基督創建了,.
或者說引進了這條道路,.

$^{121}$也讓之後每一個跟隨者,.
包括你和我,.
可以跟著這條通往永生的道路,.
我們說成為一個出死入生的新人類..
換句話來說,.
在保羅眼中,.
教會群體就是,.
父上帝藉著耶穌基督的死而復活,.
所創造出來的新群體..
神學一點說,.
教會就是神在基督裡面的新創造,.
沒有了耶穌基督的死而復活,.
沒有了耶穌基督戰勝死亡的權勢,.
就不會有今天教會群體的存在,.
甚至存活..
所以你會看到,.
保羅為什麼說耶穌基督能夠居首位,.
一方面是因為他在救贖的範疇裡面,.
他是超越的,.
因為他是首先第一個戰勝了死亡的核制,.
開創了教會這個合資格承受永生的群體的第一個人..
而另一方面,.
他也在創造的範疇裡面,.
他是超越的,.
他是居首位的,.
因為他是萬物創造的經手人,.
那個維繫者,.
他更加是創造的目標和目的..
你會看到,.
我們說保羅就是用了創造和救贖,.
這兩個基督教最重要的神學範疇,.
來建構了屬於他自己的基督論..
他強調耶穌基督,.
凡事,.
即是每一個範疇裡面,.
他都是居首..
我們沒有一件事,.
不是在耶穌基督的掌管掌握的裡面,.
他在你任何的生活領域裡面,.
他都要作主,.

$^{161}$他都要掌管..
他不只是救了你,.
然後讓我們無人駕駛,.
耶穌基督救了你,.
是要掌管我們,.
是要讓我們回復到上帝創造的旨意裡面..
所以,.
無論是原先已經發生了的創造,.
以及之後在屬靈上教會的新創造,.
保羅說,.
耶穌基督都是居首位..
再看,.
去到19-20節,.
保羅進一步解釋,.
基督為什麼在凡事上居首位,.
有兩個原因,.
第一個是,.
因為上帝喜悅,.
經文說,.
因為父神喜悅,.
叫一切的風聖,.
居住在基督裡面,.
所以他凡事居首位..
第二個原因是,.
因為父藉著耶穌基督,.
使萬有與他和好,.
所以耶穌基督凡事居首位..
先解釋第一個原因,.
居住,.
我想讓我們想起舊約,.
上帝住在哪裡?.
讓我看看,.
在受約裡面,.
上帝通常住在聖殿,.
或者說,.
之前就是會幕,.
那就是上帝的居所,.
人如果要敬拜親近上帝,.
就要來到,.
以前的會幕,.

$^{201}$或者大衛時期的聖殿,.
他才可以親近上帝,.
經歷上帝的福氣,.
風聖,.
但現在說什麼?.
上帝喜歡居住的地方,.
不是聖殿,.
他喜歡住在哪裡?.
住在他的愛子那裡,.
他喜歡住在他的愛子裡面,.
所以今天我們每一個,.
新約底下的信徒,.
我們要親近上帝,.
經歷他的福氣,.
你不需要再去找耶路撒冷的聖殿,.
因為都不存在了,.
你只需要住在耶穌基督裡面,.
也就是說,.
你一生去認識,.
跟隨耶穌基督和他的吩咐,.
保羅說,.
你就能夠經歷上帝一切的風聖,.
這裡說一切的風聖,.
不是純粹說人間的福氣,.
不是人找到你生活一切的滿足感,.
一切的風聖,.
其實有一個更高的層次是,.
你找回你做人的意義是為了什麼,.
靠我們自己找不到,.
因為我們是受造物,.
你唯有連結創造你的上帝,.
你才知道你受造的意義是什麼,.
你為什麼為了什麼來到這個世界,.
這是一切風聖的意思,.
保羅說,.
你只要有耶穌住在他裡面,.
就有了所有答案,.
你找到一切的風聖,.
所以,.
哥羅西的信徒不需要擔心,.

$^{241}$自己已經接受了正統福音信仰,.
仍然有不足之處,.
簡單說,.
當你有病,.
或者需要吃一些補充物的時候,.
耶穌就是那一顆,.
一顆完成的藥,.
你吃一顆就完成,.
你不需要在這一顆以外,.
再吃其他補充劑或者藥,.
沒用的了,.
那一顆已經完成了,.
信徒不需要在基督以外,.
另外找所謂完美風聖,.
當然,.
保羅說,.
耶穌基督擁有上帝一切的風聖,.
也是為了他之後的一些勸勉,.
賣下伏筆,.
他提醒哥羅西的信徒,.
其實你的信仰根基已經很好,.
你只要專心跟隨耶穌,.
順服他的吩咐,.
你就不需要聽其他假教師的異端邪說,.
你也不需要在基督以外,.
尋求屬靈的風聖,.
你不會找到,.
因為上帝已經命定了,.
一切住在基督的裡面,.
而基督居首位的另一個原因,.
二十字這裡說,.
是因為父神,.
他藉著他在十字架上流出的寶血,.
使萬有與上帝和好,.
他強調耶穌基督的救贖,.
所帶來的那種關係的復和,.
他不只是說,.
耶穌基督救贖你,.
讓我們可以回到天家,.
他是說一段關係的和好,.

$^{281}$簡單說就是因為罪進入世界,.
令人和上帝原本和諧的關係破裂了,.
也間接令整個大地,.
整個創造受到虧損,.
所有秩序都混亂了,.
而上帝最終透過耶穌基督的寶血,.
讓一切的受造物,.
都能夠因為這個救恩,.
與上帝和好,.
之前被罪所破壞的創造秩序,.
也會因為耶穌救贖的緣故,.
一一恢復過來,.
最終上帝整個的創造計劃,.
會因為耶穌的救贖,.
而達至完全,.
甚至於圓滿,.
換句話來說,.
在上帝的角度來看,.
對上帝而言,.
萬物與他和好,.
其實已經因為耶穌在十字架上,.
成就的救恩,.
已經成為了客觀既定的事實,.
在上帝眼中已經完成了,.
再沒有任何的勢力,.
任何的事物,.
可以難阻上帝與人恢復和好關係,.
因為已經完成了,.
藉著耶穌基督,.
但有趣的是,.
這個和好的效果對人來說,.
對我們來說,.
它又不會自動發生,.
它不會自動發生,.
唯一能夠發生的是什麼?.
就是我們經常說,.
要用信心來回應,.
上帝的救恩已經預備了,.
和好的救恩已經預備了,.
但你要憑信心來知取,.

$^{321}$你要接受耶穌基督這個和好信息,.
或者用保羅的話所說,.
就是接受和好的福音,.
與上帝和好的福音,.
唯有這樣,.
你才能夠與上帝復和,.
經歷到當中所帶來,.
什麼是豐盛的生命,.
所以你看到,.
保羅就是這樣,.
單單建構了整個基督論述,.
也將這個論述,.
應用實踐在,.
哥羅西信徒的身上,.
最終這一切的神學論述,.
是要落地的,.
是要實踐的,.
所以他勉勵,.
他的聽眾,.
他的讀者,.
甚至包括今天的我們,.
你要對你所聽見,.
聽到的福音真理,.
你要有相應合宜的回應,.
保羅在最後的段落,.
他提醒,.
既然耶穌那麼厲害,.
既然有耶穌處理好,.
他那麼超越的時候,.
你跟隨了他,.
你認順了他,.
你聽了他的福音,.
領受了他的吩咐,.
那你要怎樣回應?.
就是剩下21至23節,.
保羅要交代,.
首先保羅表示,.
哥羅西信徒,.
循前的光景是同神隔絕的,.
就像剛才所分享,.

$^{361}$他們的心思,.
他們的言行,.
這裡說得很嚴重,.
是與神為敵,.
保羅強調,.
這種跟上帝敵黨的態度,.
不是無心之失,.
不是不經意得罪了上帝,.
而是人有意識地背離,.
背叛上帝,.
所以哥羅西的信徒,.
或者包括我們,.
我們沒有能力靠自己,.
撥亂反正,.
搞好我們跟上帝的關係,.
人沒有能力靠自己跟上帝和好,.
唯有上帝親自透過他的愛子,.
促成這件事,.
所以保羅繼續說,.
從前不能逆轉的情況,.
現在有轉機了,.
因為耶穌基督,.
寫命的緣故,.
哥羅西的信徒,.
可以跟上帝和好,.
而這個和好是有目的的,.
是要使他們成為聖傑,.
就像舊約那些,.
無瑕疵,.
毫無指責的濟生一樣,.
可以成為一個活祭,.
獻給上帝,.
獻到上帝面前,.
當然必須強調,.
信徒不會一信主,.
就會立即變成,.
沒有瑕疵,.
無可指責,.
因為這是一生的過程,.
也就是強調教會,.

$^{401}$強調信徒在信主之後,.
其實你要不斷促進這段關係,.
你要不斷在生命裡面,.
改善,.
改進你的生命,.
最終在信服基督的事情上,.
你達至無可挑剔的地步,.
換句話來說,.
保羅在這裡期望,.
哥羅西的信徒,.
雖然他們已經擁有一個很好的根基,.
特別像上一次陳牧師所說,.
他們已經結果只有增長,.
他們在屬靈上,.
已經有一個很好的基礎,.
但是保羅強調,.
你們要繼續成長,.
不要滿足於現狀,.
不要安於現狀,.
因為不進則退,.
所以,.
哥羅西信徒的當前責任,.
保羅在這裡很簡單地說,.
你繼續持守你的信仰,.
如果按照上文,.
你繼續做好你本身已經有的屬靈基礎,.
在那裡繼續建上去,.
你繼續活出和福音相稱的.
美善品德和行為,.
讓教會無論在質量,.
或者人數上都能夠增長,.
如果按照下文,.
堅守信仰就是提醒信徒,.
在過程裡面,.
撒旦不會輕易放過我們,.
他會透過其他人,.
其他的異端邪術,.
來去搞擾教會,.
他會動搖教會已經奠定好的信仰根基,.
他會引誘教會偏離原本已經有的福音盼望,.

$^{441}$保羅就是這樣勉勵哥羅西的信徒,.
你要堅定不移,.
這樣去相信和實踐,.
你已經聽了的福音,.
他強調,.
這就是他在耶穌基督裡面,.
找到的人生意義,.
保羅就是因為這個福音的緣故,.
因為這個和好的福音的緣故,.
他知道人生受造的目的,.
意義是甚麼,.
首先就是,.
你要和上帝和好,.
不是你做了多少善行,.
是你和上帝和好,.
藉著耶穌基督和上帝和好,.
這是每一個人或每一個罪人,.
他在地上第一個的意義和目的,.
沒有這東西,.
你所有東西都是彎的,.
所有東西都是向著死亡前進,.
所以保羅很明白,.
原來與神和好,.
不單單只是他人生的目的和意義,.
不單單只是他受造的目的和意義,.
並且他知道他要承擔這個和好的積分,.
他要繼續去延展這個和好的信息,.
讓其他人知道,.
所以你看到為何保羅打後,.
他會刁氣萬事,.
只追求一個目標,.
就是奮身成為福音的執事或是僕役,.
因為他明白到人生唯一的目的,.
不單單只是與神和好,.
並且是勸人與神和好,.
好了,終於說完了,.
還有大概十分鐘左右,.
我們可以做一些應用反省,.
如果根據剛才保羅的精彩基督論,.
假如有人問我們,.

$^{481}$到底耶穌基督為何要來到這個世界上?.
到底耶穌基督為何要替我們受死?.
我們的答案會是甚麼?.
為了拯救我,為了幫我贖罪,.
這是正確的答案,.
這個標準答案是對的,.
不過如果按照保羅在歌羅西書所說的基督論,.
他只答對一半,.
耶穌基督拯救我們,替我們贖罪,.
最終的目的是為了使萬物與上帝和好,.
他提醒作為耶穌門徒的我們,.
原來我們在上帝的復和計劃裡,.
是有一個重要的角色,.
我們需要像保羅那樣,.
去完成這個勸人和好的職份,.
因為在創造的萬物當中,.
其實你想想,.
誰首先與上帝和好?.
在宇宙創造萬物當中,.
誰首先與上帝和好?.
誰呀?.
就是基督徒,.
就是信了耶穌的人,.
首先與上帝和好,.
然後他們都領受了這個和好的職份,.
繼續在上帝的創造旨意裡,.
去傳揚這個福音,.
最終使萬物在基督的裡面,.
都與上帝和好,.
簡單來說就是,.
原來我們人生的意義,.
不是單單與上帝和好,.
這是第一因,.
與上帝和好了之後,.
你要繼續將這個和好信息,.
傳給你身邊的人,.
最終達至,.
好像保羅所說,.
萬物藉著基督與上帝和好,.
所以教會群體的召命,.

$^{521}$就是回應這位創造主,.
和救贖主,.
他已經成就了的和好工作,.
所以我們都要肩負起這個職份,.
我們要為到福音,.
見證福音,.
來做好本份,.
向世人展示出基督和好的福音,.
將上帝的豐盛實踐出來,.
事實上,.
保羅這個和好的信息,.
提醒了我們,.
我們與上帝和好,.
不是之後各自各精彩,.
我們很多時候搞錯了,.
我信了耶穌就行了,.
耶穌是我的恩主,.
中國人的觀念,.
有人幫了你,.
你稱他為什麼?.
向你施恩惠的時候,.
你叫他恩公,.
你怎麼回應?.
報了恩就行了,.
你會不會之後再以身相許?.
大概不會,.
很多時候我們會用恩公恩主的觀念,.
來信耶穌,.
我信了你,.
我做一些事情來報答你,.
就拉平了,.
之後就怎樣?.
我有我的生活,.
你有你的忙碌,.
如果有事的時候,.
我有問題,.
我就會找恩公來,.
又或者,.
我們以為我們與神和好,.
上帝拯救了我們之後,.

$^{561}$我只需要在我原有的生活方式和態度裡面,.
我再額外做一些宗教行為,.
或者說信仰聚會,.
這樣就行了,.
搞錯了,.
與神和好不是在我們已有的生活上,.
已有的做人態度裡面,.
我們外加一些宗教行為,.
與神和好其實,.
如果用我們這一系列的講題,.
就是上帝要改變你的DNA,.
上帝要改變我們底層的顏色,.
不想我們只留於表面的敬虔,.
因為耶穌基督來到世上,.
不是純粹解決我們生活的煩惱,.
不是滿足我們個人的喜好需要,.
耶穌基督救我們,.
是要將我們帶回到,.
歸回到上帝創造的旨意裡面,.
所有蒙基督救贖的人,.
都要忠心履行,.
耶穌基督在創造的時候,.
已經賦予了給我們的任務和責任,.
就是成為一個大地的管家,.
特別是在新約裡面,.
就是成為福音的培育執事管家,.
換句話來說,.
與神和好是提醒我們,.
你回到原廠設定,.
你要回歸耶穌基督這位創造主,.
對你的設定,.
意思是什麼呢?.
可能大家也有試過,.
你手機電腦用著用著,.
又會怎樣?.
用得久,就會開始慢,.
背不起,.
或者不停地load,.
所謂的死機,.
很多的bug,.

$^{601}$這個時候,.
除了你買了一塊布料,.
還要怎樣?.
丟掉去買布料,.
就要拿回去維修,.
或者說拿回去reset,.
有個reset按鈕,.
就回復到你出廠的時候,.
那個正常正確的狀況,.
其實上帝也是藉著耶穌基督的拯救,.
恢復我們原先被賦予的目的和功能,.
我們起初在哪裡來,.
我們最終會去哪裡,.
在中間這個過程,.
我們在人生,.
在世有什麼目的意義,.
保羅說答案就在耶穌裡面,.
你不和耶穌結連,.
你不和上帝和好,.
你自己去找的話,.
你永遠都會迷失,.
永遠都會焦慮,.
永遠都不會滿足,.
因為唯一能夠讓我們豐盛的是耶穌基督,.
和祂的旨意,.
而和好也提醒了我們,.
不要和神和好就算了,.
意思是什麼呢?.
舉例,一對情侶,一對夫婦吵架,.
和好之後會怎樣?.
會不會固態服盟?.
我和我老婆吵架,我惹她生氣,.
她原諒我,我又原諒她,.
我又繼續做我的惡行,.
我又繼續做我的漏襲,.
下一次又怎樣?.
又惹她生氣,又罵她,.
然後又和好?.
循環不息?.
不會的,.

$^{641}$和好其實是有後續的,.
和好了,你不會當兩個人沒事發生,.
你不會固態服盟,.
和好意味著我們一段關係裡,.
在態度,在行為上,.
你要成長,你要改變,.
或者說夫婦關係之所以能夠維繫更親密,.
是因為雙方都願意因為愛對方的緣故而改變,.
我們和上帝的關係其實都是這樣,.
和好意味我們要繼續和上帝的關係有增值,.
有增長,.
你要認識上帝的喜好是什麼,.
你要改變你自己的漏襲,.
不要再做一些上帝憎惡的事,.
另一方面你要致力做一些上帝喜悅稱讚的事情,.
或者說用耶穌所說,.
你要將你的財富轉到天上,.
你要持續在你和上帝的關係裡,.
你要存錢,.
你要增值,.
所以頂枝末與神和好,.
其實只是一個起始點,.
或者說在你這一生裡,.
你都要繼續向前走的方向,.
提醒我們你要持續依靠耶穌,.
致力活出一個和上帝連結的生活,.
而最後你會發現,.
在整段經文裡,.
保羅他由萬物的創造,.
最終去到教會這個屬靈的新創造裡,.
他顯示了上帝和基督,.
他在歷史上,.
在人類歷史上,.
他是有絕對的主權,.
對有些人來說,.
歷史只不過是一個周瑜復始,.
一個沒有意義的循環,.
人在裡面除了吃喝玩樂,.
除了滿足一些人的需要之外,.
其實不知道為了什麼,.

$^{681}$特別大多數人會在歷史裡,.
覺得自己只是一個小角色,.
我是一個微不足道的人,.
我顧慮自己已經很好,.
但是保羅其實這裡向我們宣告,.
耶穌基督來到這個世上,.
一方面他肯定,.
這個世界的歷史有開始,.
也有終結,.
掌管在他的手裡,.
在當中每個人都被賦予了,.
上帝要他履行的角色和身份,.
也給了他們能力,.
給了他們恩賜,.
並且耶穌當他再來的時候,.
他會向所有的人問責,.
無論是信徒或非信徒,.
所以作為信徒,.
我們不能夠再說自己,.
我們只是小角色,.
只是小薯條,.
我們是微不足道,.
沒有影響力,.
因為上帝不是這樣看我們,.
耶穌說每一個凡相信,.
他福音真理的人,.
已經被賦予了一個重要的,.
角色和崗位,.
好像保羅一樣,.
好像哥羅西的信徒一樣,.
就是要出去,.
圍到基督這位創造主,.
這位救贖主,.
來作見證,.
將這個與人和好的積分,.
這個福音,.
一路去見證,.
延展開去,.
藉著這樣,.
我們的生命,.

$^{721}$就能夠不斷,.
連於我們的元首基督,.
所以,弟子們,.
盼望這段經文,.
能夠再一次提醒我們,.
其實裡面有很多東西我們是知道,.
不過我們未必有時間去整合,.
和想想怎樣去實踐去形容,.
但是今天,.
無論你理解到多少都好,.
你明白一件事,.
就是你受造是有目的的,.
目的的意義,.
來自於耶穌基督賦予給你,.
如果用回這段經文,.
就是你要致力一生都與神和好,.
因為這個不是說那一刻,.
而是說你一生都要靠主,.
與神和好,.
意思就是,.
你要建立和上帝的關係,.
聽祂的吩咐,.
不斷去將這個和好的福音,.
傳遞開去,.
直到祂再來的日子,.
讓我們同心同意..
耶穌說我們多謝你,.
你不單單只是那位,.
拯救我們出死入生的主,.
為我們打開了永生的道路,.
其實你更加是那位,.
創造我們生命的主,.
我們的人生其實,.
原本就是掌握在你的裡面,.
你提醒我們,.
不是無緣無故來到這個世界,.
我們是有目的被做的,.
我們是有意義被做,.
我們也是因為你愛我們的緣故而被做,.
上帝是盼望藉著今天這段經文,.

$^{761}$再一次喚醒我們每一個,.
我們要歸回你,.
我們明白我們已經因為耶穌,.
你的緣故與上帝和好,.
並且我們知道和好,.
意味著我們要繼續去成長,.
在與上帝的關係裡面,.
不斷去增值,.
不斷的討你喜悅,.
亦都可以承擔這個勸人和好的職份,.
讓更多的人藉著耶穌基督,.
你可以與神和好,.
回到上帝你創造他們的旨意裡面,.
下王主你幫助我們,.
讓我們將已有的信仰知識,.
我們能夠轉化成為我們生活的一個堅實的信念,.
並且我們致力在我們的言行,.
在我們生活不同的崗位裡面,.
實踐出來,.
讓人看見我們,.
就能夠看見上帝的榮美,.
就能夠認識到耶穌基督你裡頭的豐盛,.
為此我們同心獻上祈禱,.
奉基督耶穌你得勝的名字,.
來祈求,.
阿們..
以我們的名字祈求.
阿們.
\newpage



\section{歌羅西書 1:24-2:5}
\label{sec:sjw5HKR2iIk}
\textbf{2026年1月25日崇拜直播｜吳真真牧師｜門徒DNA:在基督裡完全｜歌羅西書1\_24-2\_5}
\newline
\newline
連結: \href{https://youtube.com/watch?v=sjw5HKR2iIk}{\texttt{https://youtube.com/watch?v=sjw5HKR2iIk}} ~~~~ 語音日期: 2026-01-25
\newline
\newline
\hyperref[sec:RDxRU6iPPzM]{\small{< < < PREV SERMON < < <}}
~
\hyperref[sec:index_chronic]{\small{[返順時目]}}
~
\hyperref[sec:_GdsDmJHJDI]{\small{> > > NEXT SERMON > > >}}
\newline
\newline
歌羅西書 1:24-2:5
\newline
\begin{longtable}{cl}
\hline
\hline
章節 & 經文 (和合本修訂版)\\
\hline
1:24 & \begin{tabularx}{0.7\textwidth}{X} 現在我為你們受苦，倒很快樂；並且為基督的身體，就是為教會，我要在自己的肉身上補滿基督未盡的苦難。 \end{tabularx} \\ \\ \relax
1:25 & \begin{tabularx}{0.7\textwidth}{X} 我照神為你們所賜我的職分作了教會的僕役，要把神的道傳得完滿； \end{tabularx} \\ \\ \relax
1:26 & \begin{tabularx}{0.7\textwidth}{X} 這道就是歷世歷代所隱藏的奧祕，但如今向他的聖徒顯明了。 \end{tabularx} \\ \\ \relax
1:27 & \begin{tabularx}{0.7\textwidth}{X} 神要讓他們知道，這奧祕在外邦人中有何等豐盛的榮耀；就是基督在你們心裡成了得榮耀的盼望。 \end{tabularx} \\ \\ \relax
1:28 & \begin{tabularx}{0.7\textwidth}{X} 我們傳揚他，是用諸般的智慧，勸戒各人，教導各人，要把各人在基督裡完完全全地獻上。 \end{tabularx} \\ \\ \relax
1:29 & \begin{tabularx}{0.7\textwidth}{X} 我也為此勞苦，照著他在我裡面運用的大能盡心竭力。 \end{tabularx} \\ \\ \relax
2:1 & \begin{tabularx}{0.7\textwidth}{X} 我要你們知道，我為你們和老底嘉人，和所有沒有與我見過面的人，是何等地勤奮； \end{tabularx} \\ \\ \relax
2:2 & \begin{tabularx}{0.7\textwidth}{X} 為要使他們的心得安慰，因愛心互相聯絡，以致有從確實了解所產生的豐盛，好深知神的奧祕，就是基督； \end{tabularx} \\ \\ \relax
2:3 & \begin{tabularx}{0.7\textwidth}{X} 在他裡面蘊藏著一切智慧和知識。 \end{tabularx} \\ \\ \relax
2:4 & \begin{tabularx}{0.7\textwidth}{X} 我說這話，免得有人用花言巧語迷惑你們。 \end{tabularx} \\ \\ \relax
2:5 & \begin{tabularx}{0.7\textwidth}{X} 雖然我身體不在你們那裡，心卻與你們同在，很高興見你們循規蹈矩，對基督的信心也堅固。 \end{tabularx} \\ \\
[1ex]
\hline
\hline
\end{longtable}
$^{1}$今日經份是記載在新約聖經.
歌羅西書一章二十四節至二章五節.
我們教會的聖經新約的二百八十頁.
歌羅西書一章二十四節.
現在我為你們受苦 賭博歡樂.
並且為基督的身體.
就是為教會要在我肉身上.
保門基督患難的缺欠.
我召神為你們所賜我的職份.
作了教會的執事.
要把神的道理傳得傳遍.
這道理就是歷世歷代所隱藏的奧秘.
但如今向他的聖號顯明了.
神願意讓他們知道.
這奧秘在愛邦人中有何等豐盛的榮耀.
就是基督在你們心裡成了有榮耀的盼望.
我們傳揚它.
是用諸般的智慧.
勸誡各人 教導各人.
要把各人在基督裡.
完完全全地引到神面前.
我也為此勞苦.
召著他在我裡面運用的大能.
盡心竭力.
我願你們曉得.
我為你們和老弟家人.
並一切沒有與我親自見面的人.
是何等的盡心竭力.
要他們的心得著安慰.
因愛心互相聯絡.
以致豐豐足足在悟性中有充足的信心.
使他們真知道.
神的奧秘就是基督.
所積蓄的一切智慧知識.
都在他裡面藏著.
我說者說.
免得有人用花言巧語迷惑你們.
我申指雖然與你們相離.
心卻與你們同在.
見你們循規蹈矩.

$^{41}$信基督之心也堅固.
我就歡喜了.
各位弟兄姊妹.
今天來到門徒DNA的第三講.
在基督裡源泉.
承接上次Raymond和大家的宣講.
我們知道基督使我們這群信徒.
蒙救贖.
以致我們能夠同神和好.
這個過程是一個進行式.
信徒不應該安於現狀.
要不斷地長進.
上上思念一下.
我們的行為是否與領受的愛相稱.
否則會很容易被引渡離開福音的盼望.
保羅在第一章23節說明了.
他就是福音的執事.
今天我們說到一章24至二章5節.
保羅說他如何承擔教會的執事.
全然投身在傳揚福音的事宮上.
不辭勞苦.
使各人在基督裡源泉.
同時我們一起讀這段經文的時候.
我們也想一想.
旺朕如何因著這段經文被激勵.
在基督裡我們一起邁向成熟源泉.
讓我們一同祈禱.
親愛的主.
是你愛我們.
你賜下這個奇妙的救恩.
在我們的身上.
上帝啊.
我們真的不能離開你.
因為離開你.
我們人生就會毀了.
我們不能像保羅所說.
我們完完全全獻情給你.
我求你去教導我們.
願聖靈引導我們.
去留心你的話語.

$^{81}$留心你寶貴的真理.
以致我們心靈能夠去確切去回應.
讓這個門徒的DNA.
在我們身上不斷發生過效.
我們這樣禱告交託.
奉主耶穌基督得勝名字祈求.
阿們.
一章24-25節.
讓我們一起讀出.
現在我為你們受苦,倒閣,歡樂.
並且為基督的身體.
就是為教會.
要在我肉身上.
補滿基督患難的缺陷.
我照神為你們所賜我的職份.
作了教會的執事.
要把神的道理傳得傳遍.
上次Raymond說.
我們就要看看.
保羅如何奮身完成這個使命.
他這個奮身是怎樣呢?.
第24節說到.
原來他做基督身體.
這個教會的執事.
是要受苦的.
不過他說.
雖然受苦.
但我是歡樂的.
原因是因為.
這樣在我肉身上.
能夠補滿基督患難的缺陷.
這個片語.
我覺得是很難明白的.
我們要理解一下.
究竟基督患難的缺陷.
又或者在和守版裡.
譯作基督未盡的苦難.
是為甚麼呢?.
這是一個猶太人的觀點.
在新天新地之前.

$^{121}$將會有很多很多的苦難.
在新約.
馬太福音,馬可福音,啟示錄.
都曾經說過.
在彌賽亞的時期裡.
在末後的日子裡.
是會有限額的苦難.
所以保羅說.
他要補滿的意思.
是既然神預定彌賽亞的時期.
會有這麼多苦難.
就讓他們一起去承受.
一起承受多一點.
好讓這個苦難能夠早日滿額.
耶穌基督的聖功.
就可以早日成全.
簡單來說.
保羅要以自己的生命.
去延續耶穌基督的一生的苦難.
不過我們一定要.
小心留意一點.
保羅受的苦.
並不是基督代我們死.
為我們贖罪的苦.
他自己不可以像基督那樣.
將人從罪和死亡裡拯救出來.
而是他作為一個基督的執事.
因為他有這個使命在身的緣故.
他受監禁.
受皮肉之苦.
受在教會裡被攻擊.
抵毀誣衊的困苦.
都是為了教會的信仰的使命.
可以延續下去受苦.
而保羅盡心竭力地實踐這個使命的時候.
又能夠幫助到人.
從這些死亡和罪的困苦裡拯救出來.
當我們明白了.
這個寶滿基督未盡的苦難的時候.
我們也要知道一件事.

$^{161}$就是保羅並不是說.
我去寶滿這件事.
他就是在抬高自己的身家地位.
或者他要誇耀自己.
我在這個計劃裡有多獨特的角色.
有多超然的地位.
因為我們看到25節.
他說:我照神為你們所賜的職份.
作了教會的執事.
神所賜的職份.
很清楚地說明.
是神叫他做這個職事.
是神賜給他的.
就是說他這個崗位.
他這個職份.
他這個身份並不是出於仁義.
不是被人欽點.
不是自己要爭著做的.
他是按著神的旨意.
是神自己親自呼召他.
「執事」這個字.
「愚民直譯」是「瀑譯」.
他是受神差遣.
是聽命於神去服事.
他說他為你們去做這個職份.
你們的意思就是基督的教會.
基督的身體教會.
亦在經文裡提到.
就是上帝的聖徒.
基督的聖徒.
一起跟從基督.
從這個世界分別出來的聖徒.
我們可以肯定.
他雖然受託.
受任命去做這個職事.
但是他甘心承擔這個苦難.
因為聖經裡面說.
「我們受患難,原是命令的」.
這是在帖撒羅尼加前書第三章第三節.
登山寶訓亦說.

$^{201}$「我們為而受逼迫,是有福的」.
全神的道是一場屬靈的真戰.
戰爭從來都是遇了受苦.
打仗永遠都不會舒服.
我們是要將人從撒旦那邊搶過來.
所以我們是遇了受苦.
當我們在教會裡面.
我們作為傳道.
我們作為領袖.
我們在教會裡面.
團長,導師,組長,會務委員.
部長,部員等等.
我們的事奉都是神賜給我們的職份.
我們是否都有著保羅這份在教會裡的使命.
我們是否清晰地在作這個身份裡面.
我們有甚麼清晰的目標和方向.
還是我們會容易,很殷勤地去做.
或者我們很健功地去做這個職事呢?.
近幾年我有機會和一些教會的領袖去聊天.
在當中我都很受他們的激勵和感動.
因為和他們聊天.
他們去擔當一些職事的時候.
或許他們都會很忙.
他們工作,家庭,照顧家人等等都很忙.
他們都會擔心自己未能夠承擔.
又或者他們又覺得自己的能力很不足,很有限.
但是當我們發出邀請的時候.
他們都會謙卑地向神祈禱.
尋求神清晰的帶領.
求主的加力引導.
這幾年我們和他們一起事奉的時候.
我們看到神和他們同在.
賜他們力量.
而且他們常常都思考.
如何在服事裡面更加合乎神的心意.
我很希望這也成為旺朕的福氣.
既然保羅為教會願意甘心勞苦去做服役的時候.
其實他最終的使命是什麼呢?.
第25節的下半節.
他講到要把神的道理傳得傳遍.

$^{241}$他要把福音講清講楚.
要把福音完整地分享.
而且要把福音的果效發揮出來.
第29節也提到.
原來傳福音的目的就是.
把各人完完全全地引到神面前.
引到這個字是解為「呈獻」.
是呈獻給神.
什麼叫把人完完全全地呈獻到神面前呢?.
其實這是一個舊約獻祭的觀念.
是把沒有殘缺的守身的祭物獻上.
完完全全的意思就像22節那裡提到.
就是把你們獻在基督的面前.
成為聖潔.
沒有瑕疵,無可指責.
這個獻上有兩個意義.
第一個意義是完完全全地提到.
一個屬靈生命的成長.
是關乎到信徒生命品格的問題.
第二個意義是關乎.
當他作為瀑翼的時候.
這個使命能不能夠完成任務呢?.
我們知道我們沒有人.
能夠憑著自己的力量.
在這方面做到完美.
但唯有在基督裡面.
保羅很明確地表達.
他的清楚目標.
就是要帶領每一個人.
經歷一個屬靈的轉變過程.
以至到在審判的日子.
他們能夠被上帝認出.
保羅要瀑翼在教會裡.
幫助眾聖徒在基督裡面.
他們的生命能夠完全被塑造.
被基督認出.
所以在基督裡面.
完全要使基督徒成為一個.
在主裡面成為一個成熟的人.
所以教會不單單只是.

$^{281}$我把福音說了.
弟兄姊妹信了主.
受了洗這麼簡單.
而是關乎到福音.
落實在一個人的生命裡.
他到底在屬靈的生命裡.
經歷了什麼改變.
什麼轉化.
以至他一生都邁向成熟呢?.
我們一起看看.
究竟怎樣可以做到.
在第26至27節.
讓我們一起讀.
「這道理就是歷世歷代所隱藏的奧秘.
但如今向他的聖徒顯明了.
神願意叫他們知道.
這奧秘在愛邦人中.
有何等豐盛的榮耀.
就是基督在你們心裡.
成了有榮耀的盼望」.
作為教會的瀑易.
要清楚傳揚神的道.
完全是建基於基督的救贖上.
對於猶太人來說.
他們很難想像.
這個救贖可以淋到愛邦人當中.
所以保羅指出.
神的福音和救贖這件事.
是歷世歷代的奧秘.
這個奧秘不是在說什麼神秘主義.
只是在舊約中.
這個奧秘被隱藏了.
當耶穌基督被釘十字架.
受死復活後.
這個新約的救贖計劃顯明出來.
從前只有猶太人是神的子民.
但今天愛邦人.
包括你和我.
都可以因為耶穌基督在我們心裡.
在我們一切當中的緣故.

$^{321}$成為神的兒女.
這是極大的榮耀.
而他們的盼望.
就是在基督將來榮耀再來的時候.
信徒都同樣可以分享這份榮耀.
第27節的豐盛榮耀.
根據希臘的原文翻譯.
有其他的譯本.
例如新翰語譯本.
環球聖經譯本.
都將它翻譯為「榮耀的豐盛」.
它將核心重點放在豐盛上.
世界上的每一個人.
不只是猶太人.
不是有些特權的人.
而是每一個人.
都可以因為基督而得著救贖.
是可以得著這份榮耀.
你想像一下.
本來我們一群滿身醉污的人.
不配有神在我們的生命裡.
但就是因為這個恩典的救贖.
我們的心裡有耶穌基督和我們同行.
我睡醒有耶穌和我們同行.
我面對困難有耶穌和我們同行.
我歡樂有耶穌和我們同行.
這是多麼大的豐盛.
所以除了在末世榮耀的盼望裡.
還有今生豐盛的生命在等著我們.
這個豐盛的生命.
一切都在基督裡.
讓我們一起讀第一至五節.
二章一至五節.
「我願意你們曉得.
我為你們和老弟家人.
並一切沒有與我親自見面的人.
是何等地盡心竭力.
要叫他們的心得安慰.
恩愛心互相聯絡.
以致豐豐足足在五姓中有充足的信心.

$^{361}$使他們真知神的奧秘.
就是基督所積蓄的一切知識.
都在他裡面藏著.
我說者說.
免得有人用花言巧語迷惑你們.
我心指需與你們相離.
心卻與你們同在.
見你們循規蹈矩.
信基督的心也堅固.
我就歡喜了」.
一章到第二章有一個特別的地方.
就是二章一至五節.
他將我們得蒙這份舊屬的豐盛恩典.
的普世價值.
說完之後.
他就將焦點放在.
哥羅西和老弟家教會當中.
這兩間教會和第四章第十三節的希拉玻利教會.
都是位於利古斯河谷的教會.
或許因為他們地理上比較相近.
而且亦有多往來.
甚至他們面對的問題也有類似.
他將關於得熟.
使他們在基督裡面完全的教導.
很落實很貼地針對教會的問題去說出來.
而這個教會正面對什麼問題呢?.
第二章第四節說.
「免得有人用花言巧語迷惑你們」.
哥羅西和老弟家教會可能正面對一個挑戰.
就是他們在基督教義的根基上不夠穩固.
以至有些人會用一些好聽浮誇.
花言巧語,巧言善辯的技巧.
亦帶著不正確的動機去迷惑他們離開基督.
所以寶萊要寫這封信.
就是為了要他們真知道奧秘.
知道他們這個得熟的身份.
而且是在基督裡面.
而在基督裡面所積蓄的一切智慧知識.
都在他裡面藏著.
積蓄的一切就是一切的寶庫.

$^{401}$他很想告訴你.
在基督裡面是隱藏著一切的智慧和知識的寶庫.
在面對異端的迷惑.
可能那些假教師自以為自己有知識有學問.
但是寶萊告訴他們.
我們並不需要在基督以外再去尋找這些知識或智慧.
這個奧秘的答案.
就是那位在被釘十字架.
受死復活的基督耶穌裡面可以找到.
而事實上這個寶庫是需要這班讀者去發掘.
去經驗才能夠知道.
所以保羅和教會的群體.
要藉著基督裡面的奧秘知識和寶庫.
繼續在他們生命的成熟裡面邁向成熟和完全.
到底要他們成熟和邁向完全會有什麼特質呢?.
讓我們一起看第二節.
要叫他們的心得著安慰.
因為愛心互相聯絡.
以至豐豐足足在悟性上有充足的信心.
使他們得安慰這個字.
原文是說到受激勵和得著鼓舞.
一些希臘的歷史學家會用一個例子去解釋這個字.
他說如果有一營希臘的兵.
他們處於一個極度失望和沮喪的惡劣環境當中.
他們的將軍知道這件事.
就會派一位領袖向他們作一番訓話.
然後可以令到這班士兵振作起來.
以致恢復信心和鬥志.
所以一個能夠被激勵的教會.
或者一個被激勵的信徒.
他們會有一個能力.
在面對這個世界現實的挑戰困難的時候.
他們會有勇氣有信心去面對.
以致能夠去應付任何的環境.
第二,保羅和江南西教會其實是素未謀面的.
這裡提到他們是素未謀面.
而且他們的肉身不在他們當中.
這個教會也不是他們親手建立的教會.
但是他們沒有堆起領袖的款項.
用大道理去壓下去.

$^{441}$而是用愛.
是因著愛.
他們對他們有一個有溫度的關顧.
而且他們不單單去關顧那些信徒的個人屬靈光景.
他們也很重視教會裡面的群體性.
和他們成員之間的關係.
因為信徒之間是因為什麼而聯繫起來呢?.
就是因為愛心.
如果沒有愛就不是教會.
我們教會的行政.
教會的傳統.
教會的做法.
都會隨著世代的根系時間的變遷而有變動.
但是愛是教會的標誌.
無論是對上帝的愛.
或者是弟兄姊妹之間的愛.
我最近看了一個節目叫《經慧線》.
是LOW TV 裡面的一個節目.
去探討一個問題.
可不可以用AI去做輔導治療.
去探討的最終答案是.
其實我們知道在文獻裡面說到.
輔導能夠真是達到目標.
一個很重要的關鍵.
就是輔導員和受助者之間的關係.
才是重要的因素.
所以AI不可以取代.
當然在某些危急的情況下.
它可以做到一些協助.
但是原來人與人之間的相處溫度.
那份愛的關係.
那份信任的關係.
才是人的成長.
人的治療.
人的屬靈成長.
心靈的成長.
治療裡面的一個很重要的關係.
我們留意到.
「得激勵」和「愛中聯繫」.
這個字都是一個被動的詞.

$^{481}$所以我們很清楚知道.
這兩樣東西能夠做到.
完全是神的工作.
而「得激勵」這個延伸的結果.
就是豐豐足足在悟性上有充足的信心.
又或者《環球聖經》本有云曰.
「從而得到一切的豐盛」.
我們說的是豐盛的生命.
我們從而得到一切的豐盛.
是由領悟而來的確信.
這個「傳揚真理得以完備」.
不單止是理性的認知.
而是經歷神的傳避和豐富.
以致我們在當中得到一些領悟.
有確切的了解.
我常常看到一些熟齡成熟的肢體.
他們固之然很努力去查考聖經.
去明白真理.
但是他們是有一個表現的.
他們縱然面對一些很艱難的挑戰.
困難.
但他們仍然能夠看出上帝在背後.
對他們美好的心意.
對他們熟齡成熟的操練.
可能有些人會覺得.
好像很阿Q.
但其實這些從來都是神而來的信心.
有肢體告訴我們.
他在詩篇84篇第11節說.
他未嘗留下一件好處.
不給那些行動正直的人.
他相信上帝對那些行為純正.
心裡全正直的人.
會賜下豐盛的恩典和福分.
而且是毫無保留.
所以當他們遇到生命裡.
不同的經歷和難處挑戰的時候.
他們都會對這些真理有更深的領悟.
有更深的了解和確信.
所以在2章5節說.

$^{521}$保羅看見一群信徒循規蹈矩.
信基督的心堅固.
他就歡喜.
循規蹈矩這個字.
漁民有幾個意思.
譬如他會形容在聚會裡按次序.
或者祭司按更次.
或者大祭司的等次.
這些都是循規蹈矩這個字.
其實如果我們能夠明白這個字.
我們就會知道保羅看見什麼呢.
保羅看見一群信徒.
他們做他們應該做的事.
按著上帝的吩咐行事.
不是用自己的方法.
不是用自己的目標去做.
是一個信服基督的人.
以至到他的生命堅定不移的表現.
所以當我們說一個人邁向成熟.
邁向完全的時候.
無論他人生裡遇到多少挑戰.
他仍然學習用勇氣,愛心去面對.
有悟性去理解基督待贖.
使他成為神兒女這個寶貴的身份.
以至他有充足的信心去面對一切.
當我們說我們在基督裡面.
完全是我們生命中的屬靈成熟.
我們的品格能夠成熟之外.
第二就是我們如何去回應.
作為教會的培育身份.
去完成這個使命任務.
我們能不能夠去完成呢?.
我們一起看第一章28至29節.
讓我們一起讀.
「我們傳揚他,是用諸般的智慧.
勸戒各人,教導各人.
要把各人在基督裡.
完完全全地引到神面前.
我又為此勞苦.
召著他在我們裡的大能.

$^{561}$盡心竭力」.
我們看到我們要在這個教會的.
培育身份裡我們要怎樣做呢?.
第28節說「勸戒各人,教導各人.
要把各人在基督裡.
完完全全地引到神面前」.
「各人」這個字出現了三次.
它強調了一件事.
其實教會裡有所有的人.
「各人」就是教會裡的所有人.
而他們都會很獨特.
他們有各式各樣的人.
我記得梁家倫牧師曾經說過.
教會裡藏著很多「雜不撚」的人.
什麼人都有.
我們會有不同的性格.
有不同的處境.
也有不同的身份.
但我們都是各式各樣的人.
然後第28節一開始說.
「我們傳揚他」.
本來用「我們」.
「我們」.
正正因為我們有各式各樣的人.
所以福音傳揚的工作.
不是一個人,一個超級的使徒.
一個牧師或一個偉大的報道家的事.
而是一個對公.
我記得當中有很多雜不撚的人.
但我發覺原來我們都要互相配搭去服侍.
而且用「諸般的智慧」去勸戒和教導.
「諸般」這個字.
我們常說「諸一般」.
我們的智慧怎能「諸一般」呢?.
其實「諸般」這個字.
跟「各人」這個字是一樣的.
即是說是各種的智慧.
所以我們去傳揚,去勸戒,去教導的時候.
我們絕對不可以千篇一律.
一本通書看到老.

$^{601}$我們要勸戒,去教導.
這兩個字本來的意義是很相近的.
但他們也有些不同.
因為勸戒這個字是有些撥亂反正的意味.
所以勸戒就是說那些人走歪了.
所以我們要警告他,責備他.
而教導這個字就是一個教育的過程.
是將一個人他不知道這件事.
是對一些必須要明白的觀念去作出教導.
避免他成為一個無知者.
既然我們要用各式各樣的智慧去勸戒和教導人.
我就想起在我們當中.
我們會用不同的方法.
我們要用一個對攻.
我就想起最近我們開始有三福的訓練.
三福的隊員在這幾個月裡.
我們跟差不多20個朋友去傳揚福音.
在當中我們接觸到一些在兒童主教學.
或者兒童崇拜裡面.
有小羊天地的家長.
我們又接觸過一些很有智慧的長輩.
有一些很年輕的年輕人.
或者是個性和身份上很獨特的朋友.
又或者是我們一班新移民婦女.
其實我們去跟他們傳福音的時候.
都真的要有不同的方法.
昨天我們在哈拿小組.
不知道大家知不知道我們有個哈拿小組.
哈拿小組的獨特之處就是.
我們其實是來自一班新移民的婦女.
他們最近正在上一個傳福音的工具叫Alpha Course.
我們叫啟發課程.
昨天我們分享到.
你會比較接納別人.
用什麼方法向他傳福音的時候.
他們說了一件事就是要有關係.
因為現在這個世代實在有太多太多騙局.
你對一個人不熟悉.
你見不到他生命的見證.
你聽不到他的故事.

$^{641}$你見不到他如何以身作則的時候.
你很難去信服他所說的話.
所以他說一定要有關係.
而我們在傳福音的過程中.
也會面對不同的人.
我很記得為什麼我說要是一個對工.
因為我很記得早幾年我們接觸過一位姐妹.
她以前是曾經頑皮過.
曾經誤入歧途.
當我和這位姐妹接觸的時候.
其實我飽受了她不少的粗言穢語的對待.
這也是要用很多的愛心去忍耐和引導她.
但很有趣的.
大家是否認識我們教會有一位福音幹事.
她是葛錦婷.
錦婷是我們的福音幹事.
她也是一個曾經遇到艱難的基層婦女.
在家庭各樣的生活艱難中.
她也經歷了上帝的很多.
但這個曾經誤入歧途的姐妹.
從來對錦婷不會說一句粗言穢語.
因為她知道她之前生命中經歷的艱難.
她對這個婦女有極大的尊重.
所以我們傳福音真的要一起去做.
因為教會裡真的有很多人.
我們要用不同的智慧.
用不同的方法去勸戒他們.
去教導他們.
另外我們要留意的是.
我們要竭盡全力.
第29節和第2章第1節.
都提到保羅傳福音時是盡心竭力.
這個字是運動員或比武的人.
盡全力去爭取勝利的意思.
又或者可以指竭力去奮鬥.
去完成一個事件或目標的意思.
其實保羅在傳福音這件事上.
真的盡力為了真性.
正如我說的,是一個屬靈的真戰.
要將人從魔鬼撒旦救出來的時候.

$^{681}$他必須勝利.
所以為了達到這個目的.
甚至他在獄中什麼都不能做.
他面對的是將來.
他寫這封信的時候.
在羅馬他面對將要被裁決甚至處死.
他都仍然竭力寫信給哥羅西教會的信徒.
而且他竭力為他們祈禱.
在什麼都不能做的時候.
他仍然竭力地為他們擺上.
但不要緊.
因為竭力裡面有靠主的大能.
照著他在我裡面運用的大能.
他不是靠自己的能力.
歸根究底是基督的大能在他裡面運行.
這個運用的大能.
運行的能力是描繪著神自己親自的工作.
甚至發生超自然的事情.
所以他能夠憑著神給他的力量.
做到這個傳揚的功夫.
幾年前我們曾經在誠心所願訪問過蘇永睿牧師.
我記得他常常提起一個故事去激勵我們.
他說有一次有一個牧師坐飛機.
但很感恩很幸運地他的機位升級.
可以坐商務客位.
他跟牧師坐的時候.
旁邊是一個企業的總裁.
他們兩個就聊天.
企業總裁問牧師.
你做什麼的?你從事什麼工作?.
牧師就說我從事人心靈的跨國企業.
總裁就說.
我們的公司在世界各地30個國家有辦事處.
不知道你的跨國企業有多少辦事處呢?.
他就說我們差不多在世界每一個國家都有辦事處.
總裁就問他.
我們的員工有2萬人.
牧師你從事的工作有多少人?.
在這個事業裡從事的真的有數以億計.
總裁又好奇你們公司的資產有多少?.

$^{721}$總裁就說我們有100億美金.
牧師就說我沒有正式統計過.
不過以我們多年的經驗.
只要我們跟我們的老闆一話要lor .
有需要的時候.
他一定給我們足夠.
而且足夠有餘.
所以弟兄姊妹.
在我們的事奉裡.
我們真的不是靠自己的能力.
是靠著這位滿有大能的上帝.
總結.
可能長了一點.
不過我也要說.
其實我們傳揚他.
是用諸般的智慧.
勸誡各人.
教導各人.
要把各人在基督裡完完全全的引導神面前.
在我剛出來做廚道人.
牧養的其中一個團契.
一群年輕人的團契.
他們的團份.
那時候每個星期.
我都要跟他們一起去背這句金句.
我跟他們做了七年多的導師.
他們是一群真的去衷心去完成.
上帝託付的年輕人.
當然現在都長大了.
不過他們常常提醒自己.
他們常常都有一個故事提醒自己.
我們不要成為俱樂部.
我們是帶著使命的群體.
當我想到他們的時候.
我就想到旺鎮.
我來到旺鎮已經八年多了.
羅牧師在去年退休之後.
我就成為最年長的全職的傳道人.
我不怕說的.
因為大家都說我很年輕.

$^{761}$我自己一直在想旺鎮的時候.
當保羅作為教會的執事.
為了福音的緣故.
他能夠延續不辭勞苦.
盡心竭力的時候.
我都要問我或者旺鎮.
我們可以怎樣去回應上帝託付我們的使命.
在基督裡面完全呢?.
其實是百感交集的.
不過我也很想勸勉大家.
我們在個人層面.
或者在靈命成熟的層面裡.
我們在弟兄姊妹裡的靈命塑造又是怎樣呢?.
其實剛才提到我們要在勸戒和教導上.
是怎樣呢?.
我們的愛心.
我們怎樣彼此去守望.
我們怎樣在這條路上加油打氣.
我們是否能夠鼓勵激勵弟兄姊妹.
信服我們的上帝.
是否有一個受教的心.
這都是我們常常要反思和反省.
另外我們在作傳福音的事供裡.
我們能不能夠清楚我們的定位呢?.
我們能否幫助人從死亡和罪惡的綑綁裡救贖出來呢?.
我們又有多麼的甘心受苦的心智呢?.
在我們傳福音的事上.
其實今年我們會有一個很清楚的主題.
我們要門徒不斷.
不知道大家有沒有看星期三的誠心所願.
剛剛的誠心所願就是門徒要這樣訓練出來.
嘉賓就是陳啟興牧師.
我們分享到在門徒的生命裡.
怎樣用神的真理去建立這些生命.
還有在傳福音的事上我們怎樣盡心竭力呢?.
特別在我們今年傳福音的群體.
其實有很多對象是我們很想傳福音的.
包括在培導小學裡.
我們有幾班小六的學生.
我們有進去教他們聖經課.

$^{801}$去跟他們分享福音,建立關係.
又或者我們有在一班家長的群體裡.
去教導他們怎樣做家長,有正面管教班.
又或者我們有哈拿小組裡的新移民婦女.
星期二,星期五我們有50家退休人士和長者的服務.
我們有些弟兄姊妹去到福音的機構.
無家者協會等等的機構裡.
一起在宿舍裡做福音的事工.
弟兄姊妹,我們這些工作其實很少人做.
我在這裡也很想挑戰大家.
很想幫助大家一起回歸到.
上帝想我們在裡面完全的使命.
很願意邀請大家.
無論在弟兄姊妹的守望裡.
無論在我們福音事工的傳揚裡.
我們願意旺朕都成為合神心意的一間教會.
我們一起祈禱.
親愛的天父,在你裡面有無限的豐富.
在資源上,在我們生命裡的經歷.
你也有寶藏在裡面讓我們去發掘.
但上帝求你幫助我們,不要坐在這裡.
而是我們能夠真的聽到你的激勵.
我們帶著士氣,帶著勇敢走出去.
我們回應你的吩咐,按你的心意去踏出去.
我們要繼續成為你中心的福音使者.
成為你的門徒.
在這個世界上繼續為你作炎作光.
這樣禱告,交託,是奉主耶穌基督得勝的明智祈求.
阿們.
歌羅西書這樣說.
基督在我們心裡成為有榮耀的盼望.
我們是否願意堅固我們的信心.
將全人獻上盡心竭力對神說.
憑你而行.
阿們.
陳寗泉牧師.
陳明泉牧師在下星期的活動中,請大家留意!.
親愛的恩主,求你顧念我們,祝福我們,讓旺朕不斷反思,被你激勵,.
好使我們在地上能夠成為一間合理心意,並且幫助人認識你,.
在生命上能夠邁向成熟完全的一間教會,幫助我們看得到我們可以作的,.

$^{841}$幫助我們在困難和缺陷中,你親自在裡面加能賜力,.
保守我們有這份忠心,我又懇求我們的弟兄姊妹,一起在這條天路上努力同奔,.
又願我主耶穌基督的恩惠,天父上帝的慈愛,聖靈的感動和引導,.
讓我們在地上與眾弟兄姊妹同在,從今時直到永遠.阿們..
讓我們一同唱差遣詩,回不祝福..
《差遣詩》.
今將散去,叩諸之房,喜愛歡樂滿心理..
叩知我眾信贈哀愁,一口救贖人得死..
今天生理,今天生理,令我眼眼光的黑..
願主解痛,願主解痛,滿於我眾生命中..
讓我們請坐,讓我們一同聆聽電樂《蜜荼》..
《蜜荼》.
願主祝福大家.
練劍.
別動.
\newpage



\section{歌羅西書 2:6-23}
\label{sec:_GdsDmJHJDI}
\textbf{2026年2月1日主餐崇拜直播｜陳明泉牧師｜門徒DNA:在基督裡的智慧｜歌羅西書2\_6-23}
\newline
\newline
連結: \href{https://youtube.com/watch?v=_GdsDmJHJDI}{\texttt{https://youtube.com/watch?v=\_GdsDmJHJDI}} ~~~~ 語音日期: 2026-02-02
\newline
\newline
\hyperref[sec:sjw5HKR2iIk]{\small{< < < PREV SERMON < < <}}
~
\hyperref[sec:index_chronic]{\small{[返順時目]}}
~
\hyperref[sec:397tpwFaRoA]{\small{> > > NEXT SERMON > > >}}
\newline
\newline
歌羅西書 2:6-23
\newline
\begin{longtable}{cl}
\hline
\hline
章節 & 經文 (和合本修訂版)\\
\hline
2:6 & \begin{tabularx}{0.7\textwidth}{X} 既然你們接受了主基督耶穌，就要靠著他而生活， \end{tabularx} \\ \\ \relax
2:7 & \begin{tabularx}{0.7\textwidth}{X} 照著你們所領受的教導，在他裡面生根建造，信心堅固，充滿著感謝的心。 \end{tabularx} \\ \\ \relax
2:8 & \begin{tabularx}{0.7\textwidth}{X} 你們要謹慎，免得有人用他的哲學和虛空的廢話，不照著基督，而是照人間的傳統和世上粗淺的學說，把你們擄去。 \end{tabularx} \\ \\ \relax
2:9 & \begin{tabularx}{0.7\textwidth}{X} 因為神本性一切的豐盛都有形有體地居住在基督裡面； \end{tabularx} \\ \\ \relax
2:10 & \begin{tabularx}{0.7\textwidth}{X} 你們在他裡面也已經成為豐盛。他是所有執政掌權者的元首。 \end{tabularx} \\ \\ \relax
2:11 & \begin{tabularx}{0.7\textwidth}{X} 你們也在他裡面受了不是人手所行的割禮，而是使你們脫去肉體情慾的基督的割禮。 \end{tabularx} \\ \\ \relax
2:12 & \begin{tabularx}{0.7\textwidth}{X} 你們既受洗與他一同埋葬，也就在此禮上，因信那使他從死人中復活的神的作為跟他一同復活。 \end{tabularx} \\ \\ \relax
2:13 & \begin{tabularx}{0.7\textwidth}{X} 你們從前在過犯和未受割禮的肉體中死了，神卻赦免了你們一切的過犯，使你們與基督一同活過來， \end{tabularx} \\ \\ \relax
2:14 & \begin{tabularx}{0.7\textwidth}{X} 塗去了在律例上所寫、敵對我們、束縛我們的字據，把它撤去，釘在十字架上。 \end{tabularx} \\ \\ \relax
2:15 & \begin{tabularx}{0.7\textwidth}{X} 基督既將一切執政者、掌權者的權勢解除了，就在凱旋的行列中，將他們公開示眾，仗著十字架誇勝。 \end{tabularx} \\ \\ \relax
2:16 & \begin{tabularx}{0.7\textwidth}{X} 所以，不要讓任何人在飲食上，或節期、初一、安息日等事上評斷你們。 \end{tabularx} \\ \\ \relax
2:17 & \begin{tabularx}{0.7\textwidth}{X} 這些原是未來的事的影子，真體卻是屬基督的。 \end{tabularx} \\ \\ \relax
2:18 & \begin{tabularx}{0.7\textwidth}{X} 不要讓人藉著故作謙虛和敬拜天使奪去你們的獎賞。這等人拘泥在所見過的幻象，隨著自己的慾望無故地自高自大， \end{tabularx} \\ \\ \relax
2:19 & \begin{tabularx}{0.7\textwidth}{X} 不緊隨元首；其實，由於他全身藉著關節筋絡才得到滋養，互相聯絡，靠神所賜的成長而成長。 \end{tabularx} \\ \\ \relax
2:20 & \begin{tabularx}{0.7\textwidth}{X} 既然你們與基督同死而脫離了世上粗淺的學說，為甚麼仍像生活在世俗中一樣，去服從那「不可拿、不可嘗、不可摸」等類的規條呢？ \end{tabularx} \\ \\ \relax
2:21 & \begin{tabularx}{0.7\textwidth}{X} 見上節 \end{tabularx} \\ \\ \relax
2:22 & \begin{tabularx}{0.7\textwidth}{X} 這些都是根據人的命令和教導，論到這一切都是一經使用就都敗壞了。 \end{tabularx} \\ \\ \relax
2:23 & \begin{tabularx}{0.7\textwidth}{X} 這些規條使人徒有智慧之名，用私意崇拜，自表謙卑，苦待己身，其實在克制肉體的情慾上毫無功效。 \end{tabularx} \\ \\
[1ex]
\hline
\hline
\end{longtable}
$^{1}$以下我們要聽神對我們的說話.
神的說話,今天要我們聽的.
是記載在歌羅西書第二章.
六至二十三節.
聽聽小弟讀出.
你們雖然接受了主耶穌基督.
就當遵祂而行.
在祂裡面生根建造,信心堅固.
正如你們所領的教訓.
你們要謹慎.
恐怕有人用他的理學和虛空的妄言.
不照著基督.
乃照人間的遺傳和世上的小學.
把你們擄去.
因為神本性一切的豐盛.
都有形有體的居住在基督裡面.
你們在祂裡面也得了豐盛.
祂是各樣執政掌權者的元首.
你們在祂裡面也受了不是人手所行的割禮.
乃是基督使你們脫去肉身情慾的割禮.
你們既受浸與祂一同埋葬.
也就在此與祂一同復活.
都信乃叫祂從死裡復活.
神的公用.
你們從前在割飯或未受割禮的肉體中死了.
神赦免了你們一切割飯.
便叫你們與基督一同活過來.
有圖沒了.
在律例上所寫.
攻擊我們.
有礙於我們的自據.
把他撤去.
釘在十字架上.
既將一切執政的掌權的腦內.
明顯吸中人汗.
就將十字架跨性.
所以不拘在飲食上.
或節日,月朔,安息日.
都不可讓人論斷你們.
這些源自後世的影兒.

$^{41}$那形體卻是基督.
不可讓人因為故意謙虛和敬拜天使.
就奪去你們的掌上.
這等人俱例在所見過的.
隨著自己的肉心.
無辜的至高至大.
不辭元首.
全身既然靠著他.
跟節得以相助聯絡.
就因神大德長進.
你們若是與基督同死.
脫離了世上的小學.
為什麼形象在世俗中活著.
服從那一不可靠,不可上,不可我.
等類的規條呢?.
這都是照人所吩咐,所教導的.
說到這一切正用的時候.
都敗壞了.
這些規條使人徒有自衛之名.
用施以崇拜.
自表謙卑.
苦待己身.
其實在黑者肉體上的承諾上.
是毫無功效.
聖經讀筆.
能夠在教會每個星期見面.
這是上帝給我們的恩典.
不知道大家有沒有看過一本書.
這本書有一個文學家.
大文豪.
叫做C.S. Lewis 魯易斯.
其中一本書叫做《地獄來鴻》.
有些譯法叫做《魔鬼書信》.
我見大家都有些茫然.
這本書很有趣.
這本書其實是編輯了舊時的一些散文.
記載著大約三十一篇的信件.
C.S. Lewis 是一個文學家.
他用則寫的角度.
來講述上帝的工作.

$^{81}$他講述的書信是.
魔鬼叔叔.
叫做 Screwtapes.
即是大魔鬼.
寫給他的侄兒.
他的侄兒.
名字叫做 Wormwood.
即是兩叔侄來寫.
在地上怎樣去迷惑.
每一個地上的人類.
魔鬼寫信給他的侄兒.
其中一封信是這樣講的.
他叫他的侄兒不要喚起人類的理性.
要幫他們盡量將他們的腦的思想.
放在人的感官體驗和意識當中.
還要將它美化成為最真實的生活.
簡單來說就是叫人不要用腦.
用多一些他的感覺.
這個就是他最真實的生活.
然後叔叔這個大魔鬼再繼續講.
上上之策就是.
任何科學都不要讓人類來接觸.
給他一個籠統的抽象的念頭.
讓他自己覺得什麼都懂.
偶爾閒聊和閱讀的時候.
拿到一些零零碎碎的東西.
就將它成為一個叫當代研究成果.
最重要的是什麼呢?.
這個叔叔提醒那個魔鬼仔.
千萬要記住.
你的目標就是令到人類.
讓他們覺得自己醉生夢死.
在這個書信其實是一些諷刺的.
一個魔鬼要迷惑人類.
他要令到人不懂得用腦.
第二件事就是要覺得.
什麼為之真實呢?.
就是他的感官.
看到的,經歷到的東西為之真實.
看不到的東西就傳寄予一種否定.

$^{121}$更加重要的是.
令到人類覺得自己很厲害.
這個就是魔鬼的手段.
讓人逐步逐步離開上帝.
今天歌羅西書其實講的內容.
直指就是要回應.
今天可能我們生活在地上.
有些東西會拉走我們.
讓我們離開這個信仰.
有時候可能這些學問或掌握的知識.
會令到我們覺得自己很厲害.
不過可能我們為自己的信仰.
埋下一個很重要的炸彈.
今天我們就藉著.
歌羅西書某程度上幾重點的經文.
我們重新思考.
今天我們在信的是什麼.
我們今天實踐的信仰是怎樣的信仰.
我們一起祈禱吧.
求上帝幫助.
願主與我們同在.
幫助我們讓我們能夠明白你的話語.
藉著你的訊息.
藉著你對我們所說的話.
調教我們.
讓我們能夠行事為人都謹慎.
成為你所喜悅的子民.
致福給每一位聽道的弟兄姊妹.
讓他們不單止聽得明白.
也都讓你的話語進到他們的心靈當中.
調教,更新.
讓他們每一個都走在上帝的旨意裡.
幫助孫講的講得清楚.
願你的話語忠心向弟兄姊妹孫講.
祈求你與我們同在.
獻上禱告奉耶穌名頭.
Amen.
我們先看第一個段落.
二章六至到第七節.
很多學者說這兩節經文.

$^{161}$就是歌羅西書裡面的一個很重點.
那個叫做鑰節.
你理解整個歌羅西書.
你拿著這兩節.
大概你會明白整個歌羅西書說什麼.
二章六至七節裡面怎樣說呢.
你們既然接受了主基督耶穌.
就當遵祂而行.
在祂裡面生根建造.
信心堅固.
正如你們所領的教訓.
你說牧師.
很平常的說話.
有什麼重點呢.
在和合本裡面.
有些地方譯的時候.
有些地方失去了原文的本意.
後來那個和合新譯本.
他已經有些調教了.
和合新譯本裡面這樣寫.
所以你們既然接受基督耶穌為主.
就應當在祂裡面行事.
也就是在祂裡面生根.
被建造.
在信仰.
就是你們所受的教導上.
這個應該是新漢語譯本.
新漢語譯本裡面.
在這裡面這樣寫.
他告訴我們有幾樣東西.
是之前的和合本你看不到的.
第一樣東西就是.
你們既然接受了主耶穌基督這個和合本.
但是在新漢語譯本裡面就說到.
你們接受了基督耶穌為主.
簡單來說.
不是拿個信仰.
拿個宗教來給自己.
令自己感覺有種心靈慰藉這麼簡單.
這個是有的.

$^{201}$信耶穌一定有心靈慰藉.
有份平安感.
這個是一個很重要的信仰經歷.
不過他強調的就是.
不單止拿個宗教來信.
更加重要的就是.
以基督耶穌成為你生命的主宰.
所以換言之.
我不是純粹拿個救恩.
拿個宗教感覺.
而是我生命裡面.
要有一個主宰.
帶領我們走人生的道路.
靈修學家A.W. Tozer圖書.
他就說過一個很重要的一句話.
他說今個世代裡面.
最大的異端是什麼呢.
就是將救贖的主.
和生命的主兩者分割了.
簡單來說就是.
你信耶穌不單止是說.
我要耶穌的救恩.
而更加重要的就是.
當你拿到救恩的時候.
你更加必須要以耶穌成為你生命的主宰.
你心靈裡面握著的.
以前是用你自己的感覺.
或者其他的信仰,神明.
成為你內心裡面握著的寶座.
他坐在那裡的.
現在你要讓出這個寶座給他.
耶穌基督成為你生命的主宰.
按著耶穌的心意.
來走人生的道路.
所以在這裡面.
其實重點就是.
這班信徒既然以耶穌基督.
你口稱裡面是主的時候.
你生命的一切活著.
或者你人生價值觀念.

$^{241}$都要以在基督耶穌裡.
以耶穌為主的生命方式.
來度過生活.
在第七節裡面講幾樣東西.
第一在他裡面生根或者紮根.
第二就是建造.
第三就是在信仰.
即是所接受的教導上得以堅固.
第四就是帶著感恩的心.
這些都不是一些很複雜的東西.
不過我想多一點的重點提出.
就是那個紮根的意思.
那個紮根的意思就是.
很多弟兄姊妹有時候信仰.
我們有時候都會的.
我們會覺得那個信仰.
是我們生命的其中一部份.
我禮拜六日上教會.
我滿足一個信仰生活.
我們就覺得這個是一個紮根的生活.
不過在聖經裡面所講的紮根的生活.
就是講到你的信仰.
其實要落到你生命的深處.
用我們今天廣東話的俗語.
或者香港的俗語來講.
要入血.
你行事為人.
你做每一樣東西.
你都要以上帝的心.
上帝是怎樣想的呢?.
耶穌基督怎樣提醒我們呢?.
用這個角度來思考人生的問題.
你做每一個的抉擇.
你都要想這是不是上帝所喜悅的呢?.
所以那個紮根的意思就是強調.
我們人生裡面方方面面.
不只是在禮拜六日或者周末的時間回來教會.
而是人生所有的決定.
每一件事都是從上帝的心腸.
或者用耶穌基督成為我們一個最重要的.

$^{281}$考慮的因素.
你覺得這樣講都很抽象.
其實不抽象的.
我們很多弟兄姊妹都經歷過.
你做父母你就知道.
如果特別在一些新手父母.
或者你回憶你舊時.
剛剛你做新手父母的時候.
你做每一樣東西.
你都是為著剛剛出生的baby來去想的.
那個廁所板太大了.
它會掉下去.
你要買隻鴨仔給它.
它要吃竹仔.
你的心裡面.
你上班都是想著兒子的.
你每時每刻.
你快點想回去見到新生的baby.
覺得它很可愛.
覺得它很滿足.
其實那個體會就是類似一個這樣的體會.
如果將我們對一個新生baby.
我們都能夠有一種本能性的做法.
我們對上帝在經文裡面更加講.
我們領受了這個恩惠.
我們無時無刻都要將上帝放在我們的心靈當中.
然後不單止放在心靈裡面.
更加要建立我們的信仰生活.
那個信仰生活都是那一樣東西.
除了聚會之外.
你要將信仰帶到你生命的方方面面裡面.
你要讓上帝在你生命.
在你的人生當中要儲東西.
儲一些經歷.
有些弟兄姊妹就曾經和我聊天.
我說有沒有見證你們的生命.
你說見證嗎?.
睡晚感恩那些.
我說除了睡晚可以感恩.
其實你每天都可以感恩.

$^{321}$有些弟兄姊妹就好像輕了一點.
好像生命裡面都是這樣過.
我說你要讓上帝介入你生命當中.
你生命每一天經歷過不同的事.
其實你只要用信仰.
用上帝的角度來看你每一天的經歷.
一定會有很多見證可以分享.
可能那些未必是石破天驚.
很宏大.
很令人蕩氣回腸的一些經歷.
不過有些小經歷.
其實是上帝為你生命譜寫的.
那些都是你生命裡面一些重要的一頁.
你不要容讓這些生命就這樣流失了.
每一天你用時間來想一想.
你讓上帝介入在你的生命當中.
亦都讓你的生活呈現有上帝的生活.
你生命就會有很多見證.
有很豐富的經歷.
相反來說.
如果我們的信仰只是一個聚會式的信仰.
你所講的或者你所經歷的東西.
就會變得很暗淡.
或者你生命裡面經歷上帝的事物.
你覺得很遠.
上帝和你好像一個很抽象的信念.
不過其實上帝和你同行.
不是一個這樣的信仰.
堅固就是保守你的信仰.
不要失腳.
你要持守你所信的.
所以這個堅固和那個教導是重要的.
因為上帝怎樣教.
聖經裡面怎樣啟示.
你就按照聖經的教導來實踐.
最後你的心裡面常常都有基督的愛.
和上帝的恩惠.
所以在這裡裡面他講到.
你以基督耶穌為主.
在基督裡你就會有剛才那四方面.

$^{361}$你的生命就會有這樣的實踐.
做了這個立論之後.
接著保羅就要用一個反義.
就是講到理想的基督徒生命是一個這樣的生命.
不過他們那些信徒.
哥羅西教會的信徒面對的挑戰是甚麼呢.
第八節開始就告訴你.
直到第十五節我們一起聽聽那段經文.
你們要謹慎恐怕有人用他的理學和虛空的妄言不照著基督.
乃照著人間遺傳和世上的小學打你們老去.
因為神本性一切的豐盛.
都有形有體地居住在基督裡面.
你們在他裡面得以豐盛.
他是各樣執政掌權者的元首.
你們在他裡面也受了不是人手所行的割禮.
乃是基督使你們脫去肉體情慾的割禮.
你們既受洗與他一同埋葬.
就在此與他一同復活.
都因信那叫他從死裡復活.
神的公用.
你們從前在過犯和未受割禮的肉體中死了.
神赦免了你們一切過犯.
要你們與基督一同活過來.
又圖沒了在律法在律例上所寫.
攻擊我們有礙於我們的自居.
把他撤去釘在十字架上.
既將一切執政的掌權的老來.
明顯地給世人看.
給眾人看就像十字架跨性.
這裡好像很長的篇幅.
其實有幾樣東西你值得留意.
我沒有辦法逐節來跟大家解釋.
重要的觀念跟大家講.
第一.
剛才講到二章六至七節.
就是理想基督徒的生命.
是一個以基督為主的生命.
不過在他們的實踐信仰裡.
有東西拉走他們.
什麼東西拉走他們呢?.

$^{401}$就是人的理學和虛空的妄言.
人的理學.
用原文的譯本用的字就是philosophy.
就是哲學.
人間的哲學.
不是在講大學哲學系.
當然可能都有這些成份.
不過在這裡所講的就是.
是人虛構想像出來的世界.
這個就叫做philosophy.
和虛空的妄言.
在這裡虛空的妄言究竟講什麼呢?.
讀上文下理你會知道.
其實就是講到人間的遺傳.
又或者那些在哥羅西.
這個外邦人身處的城市裡面的.
那些生鏽的現象.
所以那個拉走那些基督徒的.
就是這兩樣東西.
特別講一下那個世上的小學.
剛才在第八節那裡我們讀過了.
其實那個虛空的妄言.
和那個世上的小學是有關係的.
那個世上的小學.
就是那個原文用那個字叫做Stoichion.
Stoichion在聖經裡面只有這裡.
在歌羅西書才出現.
是外邦人的地方才出現.
大概就是講到那些隱蔽的奧秘的靈性.
夜觀星象.
人的運程是怎樣.
在那些希臘神話裡面去看那些星座.
用那些星座來決定人的命運.
又或者在星座裡面決定人的道德.
這件事和中國人所講的就快過年了.
拿著一本通勝和大家講.
接下來馬年運程是甚麼.
一脈相承都是那樣東西.
那些世上的小學那些有些神秘性的.
你為了令到你自己比較順一點.

$^{441}$在門口放一隻琴鎧在那裡.
琴鎧是甲鬼.
總之你放一隻青蛙在那裡.
你放不到也不要緊.
放一隻琴鎧在那裡也可以.
他們講過的.
總之放一隻青蛙就可以了.
可以擋煞等等.
你是不是真的很正.
他們就會有這種說法.
在哥羅西面對外邦的很多的那些習俗.
人世間裡面的這些習俗.
再加上這一班人剛剛接受基督.
因為當時看的聖經就是看舊約聖經.
所以他們理解的.
就會跟隨猶太人的傳統.
所以哥羅西教會裡面那些人.
受著兩個很重要的文化來影響.
一方面就是那些猶太人舊約的文化.
第二方面就是他們自己本土生來的.
斯羅文化那些迷信 星獸那些東西.
兩樣東西夾雜著.
以致到他們就經常都受著很多.
這些說話來影響他們.
保羅說你要謹慎.
如果你的人生.
你的信仰.
你信耶穌基督為主.
但同時間你又夾雜著很多猶太人的傳統.
又有你中國人.
或者你外邦人的那些星獸習俗的時候.
你會很迷失的.
你不會知道跟哪一個的.
在這裡面他特別說到.
用那個字眼就是老去.
俘虜了你的意思.
你不要想著我相信那個.
拿回那個宗族就可以了.
那些我看一看.
那些看一看.

$^{481}$他會彈進來的.
我不知道是不是我自己的問題.
我近來我發覺年尾的時候.
我看YouTube那些.
有大量那些運程那些東西.
另外如果有些國家裡面打仗.
就會有些印度神童說話.
你們有沒有的.
有些印度神童.
又說阿根廷有個女巫.
又或者等等那些東西.
就會出現在你平時接收的資訊裡面.
接著你又讀聖經.
接著你又看看這些.
這些又調和聖經所說的話.
你越調和到最後.
其實你自己就很混淆.
你不知道自己何去何從.
保羅說如果你不清晰的時候.
就會被這些言論.
這些東西來老去.
不是純粹令到你覺得.
好像娛樂那樣看.
而是會敗壞了.
會迷惑了你的心.
第十節裡面這樣說.
你們在他裡面.
這個他就是在基督裡面.
得了豐盛.
他是各樣執政掌權者的元首.
你們在他裡面.
也受了不是人手所行的國禮.
乃是基督使你們脫去肉體情慾的國禮.
你們既受洗與他一同埋葬.
也就與他一同復活.
都因信那叫他從死裡復活神的功用.
簡單來說.
你們信了主的人.
你在耶穌基督裡面.
就得到一切的豐盛.

$^{521}$你得到這些豐盛的東西.
你就不需要再看外面的那些東西.
來補足你的信仰.
你的信仰生命.
他這裡面說到.
你所信的上帝是一個很厲害的上帝.
你說今天你想聽坊間裡面說的那些東西.
在這裡面他用一個層級的說法.
世界上掌管的那些執政掌權者.
都要拜服在基督耶穌之下.
現在你的直線和耶穌來聯絡.
你還要透過他們.
還要回頭問他們.
所以在這裡面.
對於保羅來說.
如果那些人去求外面的東西.
來經歷信仰.
他覺得他們是傻的.
他覺得他們不能夠活出豐盛的生命.
這個豐盛的生命這個字都很重要.
你信了主之後.
你就擁有了一切感覺上很足夠.
那種感覺.
這一種是相對於地上.
我們人生裡面那種的貴乏.
你有了上帝.
你心靈就會踏實滿足.
你以前沒有上帝的時候.
你什麼都覺得沒有.
但信了主之後.
你覺得即使你什麼都沒有都好.
你的心靈都會很踏實.
很安穩.
這個就是那種豐盛的感覺.
用一個不算貼切的例子.
如果你在生命當中.
或者你人生裡面.
你在家裡煮菜.
你的爸爸媽媽.
或者你家裡面的太太.

$^{561}$是一個下廚的高手.
每一道菜都煮得很好吃.
比外面街頭吃的好吃.
或者吃快餐.
然後你跟家裡煮菜的那個人說.
我今天就想改變口味.
我覺得你吃的東西.
吃得太好了.
我覺得很厭.
每一餐都這麼有營養.
我想吃些不健康的東西.
我想吃快餐.
偶爾一兩次都可以.
如果你每天都去吃快餐.
你對於家裡煮好菜的那個人.
是一種冒犯.
你寧願吃快餐.
你都不吃家裡為你精心預備.
最好的最有營養又最好吃的食物.
在這裡其實是一個這樣的.
在信仰裡面.
將最豐盛的東西全部給你.
然後你說你不要.
我要外面吃.
吃快餐.
吃些低階的東西.
到最後其實就是將你拿走.
當然在這班哥羅西的信徒裡面.
可能有些人基於敬虔的動機.
他們想追溯猶太傳統來實踐的信仰.
甚至乎有些外邦人會行割禮.
割禮在舊約裡面就是分別為聖.
英俠出世八日之後.
就在他的生殖器裡面割去那個包皮.
這個是一個分別為聖的以色列人的記號.
但是外邦人經歷了上帝的恩典.
就不用做這些肉體規條.
而是要藉著基督耶穌.
來實踐生命分別為聖的那個經歷.
所以在經文裡面講到.

$^{601}$你們這些外邦哥羅西門徒.
你們都行了割禮.
不過割禮不是猶太人那些BB仔出世.
那八日當然要做的那些.
你是心靈做了割禮.
是上帝幫你做的.
這個割禮就是幫你脫去了情慾的割禮.
令到恩信你就是有這個割禮的功用.
再加上你們這班人是恩信而受贈.
已經表達與基督同死同埋葬同復活.
你們每一個已經經歷了.
要做的禮儀已經做了.
還要走回去做舊時的禮儀.
所以不用怕那些人來質問你.
說你做得不完全.
在基督裡面你已經完美成就了所有的一切.
這班人藉著基督耶穌生命已經重新活過來.
他們成為一個新做的人.
再加上在第十四節裡面說到.
律例上說到舊時的以色列人沒有做這些事就要被剪除.
但這班外邦人不會的了.
也沒有人可以用這些律例規條來攻擊你.
這些有礙舊時這些人信仰的字句.
耶穌基督已經廢除了.
摔去了.
釘在十字架上.
他將這些舊有的已經成為了過去式.
現在就是藉著基督耶穌生命得著激勵.
所以在這一段經文裡面.
其實他要說到拉走人的那些信仰.
或者那些傳統習俗.
會令到人反而會失去了信仰的重心.
在這裡裡面保羅就要把他們拉回來.
我們繼續看第十六節.
所以不拘在飲食或節期.
月數即是初一.
安息日都不可讓人論斷你們.
這些完事後事的影兒.
那形體卻是基督.
不可讓人因著故意謙虛和敬拜天使.

$^{641}$就奪去你們的獎賞.
這等人驅離在所見過的.
隨著自己的慾心.
無故地自高自大.
不持定元首.
全身寄言靠著他.
根折得以相助聯絡.
就因神大得獎進.
在第十六節其實他要撞一撞那些.
哥羅西門徒的膽子.
有些人可能質疑你.
你沒守這些傳統.
又或者你不跟這個習俗.
保羅說你不用怕他們.
你不用在這些驅離.
在這些節期.
或者那些不同的猶太人的文化實踐裡面.
來被別人論斷你.
你不用怕.
因為那些其實舊約聖經指著的是一個影.
真正的實體是耶穌基督.
你現在已經信了基督耶穌.
你是拿著這個信仰的核心.
而不是拿著那些影.
影照耶穌基督工作神心意的影兒.
再追溯到上帝那裡.
你現在一步已經到位了.
因信基督耶穌.
你已經拿了生命一切的豐盛.
你就不用走回頭.
也不怕別人去質疑你.
論斷你.
因為你知道你心裡面.
你有上帝.
一切都能夠滿足.
十八節裡面說到有些人.
那個信仰論斷那些人.
通常他們會怎樣做呢.
他們會故意謙虛.
會令到自己好像很謙卑一樣.

$^{681}$去敬拜神.
或者敬拜天使.
敬拜天使的傳統就是.
神你拜的嗎.
我們極其拜的就是拜天使.
透過天使去到彌賽亞.
再由彌賽亞去到父神那裡.
是有層級的.
我們這些人低級一點.
我們就拜一些低級一點的神.
大概有一個這樣的層級的觀念.
聖經說從來都沒有這樣說過.
耶穌叫人以祂為主.
我們直接.
Direct Report.
我們叫這個.
我們生命直接向上帝交仗.
耶穌就成為我們生命的中介.
唯一的中介.
所以你說說到自己那麼謙卑.
到最後其實只不過是要.
令到自己自高自大.
這種這樣的宗教實踐.
令到有意無意之間.
用謙卑的方式來自高.
簡單來說.
其實謙卑成為了不是生命的那種表徵.
而是成為了一個手段.
讓人覺得你很謙卑.
你很自大.
所以在經文裡面.
其實保羅他要守護住這一班的信徒.
不要迷失.
讓他們不要陷在.
這個信仰的那種迷失.
和別人的控訴裡面.
我們再看20至22節.
然後我們有些總結.
你們若是與基督同死.
脫離世上的小學.

$^{721}$為甚麼形象在世俗中活著.
服從那不可拿,不可嘗,不可摸等類的規條呢?.
這都是人所吩咐,所教導的.
說到這一切.
正用的時候都敗壞了.
這些規條使人徒有自衛之名.
用施義敬拜自表謙虛.
輔待己身.
其實在克制肉體的情慾上.
卻無果後.
簡單來說在這裡裡面.
其實有一個這段落的結尾.
他就特別再呼籲.
剛才所說的那些星獸.
猶太人傳統.
那些東西.
放下它吧,弟兄姊妹.
不要再驅離這些東西.
要活著.
就是要以基督為心.
以基督耶穌在主裡面活著.
不要再走向.
拿那些舊約的規條成為了信仰實踐.
我在這裡裡面大膽的跟大家說.
今日坊間裡面有些宗派.
其實他很著重.
要回到舊約那種實踐裡面.
基督教的教會.
他們覺得那裡才.
那些節旗.
守那些節旗.
才明白信仰.
你說偶爾做一下逾越節.
晚餐做一下鑄盤節.
我們當活動這樣OK的.
但如果你將它成為了教會裡面.
那個章則裡面.
必須要守的那些內容.
這個就有機會違反了.
哥羅西保羅對門徒的那些教導.

$^{761}$事實上我們拿到真正信仰的養分.
不是從那些律例.
而是從基督耶穌裡面.
從神學裡面來說.
我們是由新約理解舊約.
我們是由耶穌基督的生命.
在主裡面來理解舊約.
所以原則上來說.
舊約的節旗對我們來說.
可有可無的.
即使是有.
有的節旗其實都是一些教導性質.
而不是要透過這些信仰.
藉著那個規條來分別為勝.
那些規條是分別不了你為勝的.
唯有基督耶穌才可以將你分別為勝.
弟兄姊妹,做那些方法是很吸引人的.
甚至乎有時候會很有趣,很好玩的.
我們帶著經歷一下.
用一個學習態度模仿.
但是如果將它成為了我們信仰裡面的規則.
或者必然需要遵守的那些內容.
可能我們的生命.
就會陷在哥羅西那些門徒.
面對的那種艱難的處境裡面.
經文裡面提到.
有很多人做這些實踐.
好像很有智慧.
其實這些人徒有智慧之名.
其實到最後只是用他自己的角度.
來理解崇拜.
表面上好像很謙卑.
即是做很多攻克己身的一些活動.
但是到最後可能他自己.
沒有辦法實踐教導之外.
甚至乎可能自己都敗壞了.
整個經文我們大概理解了.
整段的內容.
它要告訴你今天我們的基督徒信仰.
我們信些什麼呢?.

$^{801}$如果我們的信仰.
我們不是以基督為焦點和中心.
甚至乎在基督以外.
我們加插了很多這個世界的文化.
又或者猶太人的一些傳統.
然後炒在一起.
到最後我們自己就不知道.
我們生命的重點是什麼.
帶來的就是很多信仰的壓力.
又或者面對著有很多.
心靈裡面的重擔.
我很喜歡一位神學家.
應該一兩年前過世.
Tim Keller.
他在一個福音講座裡面講到.
真正的福音信仰.
信仰是叫人得自由.
而不是賦予人有更加多的壓力.
如果一個人在信仰裡面.
他越來越覺得自己信仰很有壓力.
他只有兩件事.
第一可能他對信仰不是很真實了解.
有很多的繁文縟節.
加了在心靈當中.
以至到他覺得很有壓力.
第二件事.
可能在他的牧養環境裡面.
加諸了大量.
這個不是信仰本身要求他的實踐.
在他的生命當中.
這個人就充滿了壓力.
如果你真是一個領受福音的人.
你的心靈是很輕散的.
你因為心靈裡面有富足.
所以你願意和別人分享.
你願意奉獻.
因為你知道耶穌很愛你.
你心靈裡面因為被愛而激勵.
有一份滿足感.
所以你就很樂意來愛身邊的人.

$^{841}$即或那些人不是很可愛.
即或那些人都是三尖八角.
但是因為你知道耶穌都愛你.
這個三尖八角的人.
所以你就愛這些三尖八角的人.
你知道耶穌去實踐.
去落手落腳去服侍.
不計任何代價來付出.
圍住我們.
我們因為有這份豐盛的救恩.
所以我們都不計較.
學習來施與.
今日這段經文好像神學性比較強.
不過我很想我們弟兄姊妹.
我們知道我們所信的基督是誰.
你信了他.
你生命裡面你以他為主.
你生命經歷這份豐盛.
用這個準繩來看你平時的生活.
看這個文化.
你的心靈就會輕散.
以及你就會清晰.
你自己人生道路和方向.
求主繼續來幫助我們.
激勵我們.
\newpage



\section{歌羅西書 3:1-15}
\label{sec:397tpwFaRoA}
\textbf{2026年2月8日崇拜直播｜陳冠成牧師｜門徒DNA:棄舊更新｜歌羅西書3\_1-15}
\newline
\newline
連結: \href{https://youtube.com/watch?v=397tpwFaRoA}{\texttt{https://youtube.com/watch?v=397tpwFaRoA}} ~~~~ 語音日期: 2026-02-09
\newline
\newline
\hyperref[sec:_GdsDmJHJDI]{\small{< < < PREV SERMON < < <}}
~
\hyperref[sec:index_chronic]{\small{[返順時目]}}
~
\hyperref[sec:jemYOQSwajw]{\small{> > > NEXT SERMON > > >}}
\newline
\newline
歌羅西書 3:1-15
\newline
\begin{longtable}{cl}
\hline
\hline
章節 & 經文 (和合本修訂版)\\
\hline
3:1 & \begin{tabularx}{0.7\textwidth}{X} 所以，既然你們已經與基督一同復活，就當求上面的事；那裡有基督，坐在神的右邊。 \end{tabularx} \\ \\ \relax
3:2 & \begin{tabularx}{0.7\textwidth}{X} 你們要思考上面的事，不要思考地上的事。 \end{tabularx} \\ \\ \relax
3:3 & \begin{tabularx}{0.7\textwidth}{X} 因為你們已經死了，你們的生命與基督一同藏在神裡面。 \end{tabularx} \\ \\ \relax
3:4 & \begin{tabularx}{0.7\textwidth}{X} 基督是你們的生命，他顯現的時候，你們也要與他一同在榮耀裡顯現。 \end{tabularx} \\ \\ \relax
3:5 & \begin{tabularx}{0.7\textwidth}{X} 所以，要治死你們在地上的肢體；就如淫亂、污穢、邪情、惡慾和貪婪—貪婪就是拜偶像。 \end{tabularx} \\ \\ \relax
3:6 & \begin{tabularx}{0.7\textwidth}{X} 因這些事，神的憤怒必臨到那些悖逆的人。 \end{tabularx} \\ \\ \relax
3:7 & \begin{tabularx}{0.7\textwidth}{X} 當你們在這些事中活著的時候，你們的行為也曾是這樣的。 \end{tabularx} \\ \\ \relax
3:8 & \begin{tabularx}{0.7\textwidth}{X} 但現在你們要棄絕這一切的事，就是惱恨、憤怒、惡毒、毀謗和口中污穢的言語。 \end{tabularx} \\ \\ \relax
3:9 & \begin{tabularx}{0.7\textwidth}{X} 不要彼此說謊，因為你們已經脫去舊人和舊人的行為， \end{tabularx} \\ \\ \relax
3:10 & \begin{tabularx}{0.7\textwidth}{X} 穿上了新人，這新人照著造他的主的形像在知識上不斷地更新。 \end{tabularx} \\ \\ \relax
3:11 & \begin{tabularx}{0.7\textwidth}{X} 在這事上並不分希臘人和猶太人，受割禮的和未受割禮的，未開化的人、西古提人、為奴的、自主的；惟獨基督是一切，又在一切之內。 \end{tabularx} \\ \\ \relax
3:12 & \begin{tabularx}{0.7\textwidth}{X} 所以，你們既是神的選民，聖潔、蒙愛的人，要穿上憐憫、恩慈、謙虛、溫柔和忍耐。 \end{tabularx} \\ \\ \relax
3:13 & \begin{tabularx}{0.7\textwidth}{X} 倘若這人與那人有嫌隙，總要彼此容忍，彼此饒恕；主怎樣饒恕了你們，你們也要怎樣饒恕人。 \end{tabularx} \\ \\ \relax
3:14 & \begin{tabularx}{0.7\textwidth}{X} 除此以外，還要穿上愛心，因為愛是貫通全德的。 \end{tabularx} \\ \\ \relax
3:15 & \begin{tabularx}{0.7\textwidth}{X} 你們要讓基督所賜的和平在你們心裡作主，也為此蒙召，歸為一體。你們還要存感謝的心。 \end{tabularx} \\ \\
[1ex]
\hline
\hline
\end{longtable}
$^{1}$今日的經訓是記載在.
《歌羅西書》三章一至十五節.
《新約聖經》二百八十一至二百八十二頁.
《歌羅西書》三章一至十五節.
所以你們若真與基督一同復活.
就當求在上面的事.
那裡有基督坐在神的右邊.
你們要思念上面的事.
不要思念地上的事.
因為你們已經死了.
你們的生命與基督一同藏在神裡面.
基督是我們的生命.
祂顯現的時候.
你們也要與祂一同顯現在榮耀裡.
所以要致死你們在地上的肢體.
就如淫亂,污穢,邪情,惡慾和貪婪.
貪婪就如拜偶像一樣.
因這些事.
神的憤怒必臨到那活躍之子.
當你們在這些事中活著的時候.
也曾這樣行過.
但現在你們要棄絕這一切的事.
以及怒恨,憤怒,惡毒,誹謗並口中污穢的言語.
不要彼此說謊.
因你們已經脫去舊人和舊人的行為.
穿上了新人.
這新人在知識上漸漸更新.
正如做他主的形象.
在此並不分希利尼人,猶太人.
受割禮的,不受割禮的.
化外人,西古提人為奴的,自主的.
唯有基督是包括一切.
又在各人之內.
所以你們既是神的選民.
聖潔無礙的人.
就要存憐憫,恩慈,謙虛,溫柔,忍耐的心.
倘若這人與那人有嫌隙.
總要彼此包容,彼此饒恕.
主怎樣饒恕了你們.
你們也要怎樣饒恕人.

$^{41}$在這一切之外.
要存著愛心.
愛心就是聯絡傳德的.
又要叫基督的平安在你們心裡作主.
你們也為此蒙召,歸為一體.
.
氣救更新,好嗎?.
特別我們去回應.
好像今天這段經文所說.
脫去舊人和舊人的行為.
並且穿上新人.
如果你發覺其實.
打從上一章開始.
即是第二章.
保羅不斷去重複宣講一個主題.
就是說既然信徒已經和基督聯合了.
我們很熟悉的經常說.
我們和基督同死同埋葬同復活了.
既然如此.
我們應該有什麼相稱的行為呢?.
保羅不斷環繞著這個主題來宣講.
我們會發覺.
上一次陳牧師已經和我們分享了.
不要再受那些外在的宗教行為來限制你.
不需要額外你透過守節旗又或者割禮.
用這種外在的方式來找多一點屬靈的豐富.
保羅說是行不通的.
因為基督已經有一切的豐富在裡面.
你追求基督就已經搞定.
所以保羅提醒你要擺脫這些虛妄的學說.
現在來到第三章.
你會發現保羅很正面地說.
擺脫了舊又一些東西.
那你要怎樣去開始追求屬靈的成長.
向著一個什麼的方向.
你見到剛才一開始讀經文.
其實保羅用了一個命令的語氣.
來吩咐哥羅西的信徒.
你應該去追求天上的事.
思念天上的事.

$^{81}$理由很簡單.
他首先說.
耶穌基督現在在哪裡.
看一看.
就升到天上.
坐在了父神的右邊.
就是說所有.
一切屬天的權柄都聚焦在基督的裡面.
用我們一個簡單的概念.
就是我們和基督聯合了.
那耶穌在哪裡我們就怎樣.
你就在哪裡.
你的心思意念就對著在哪裡.
不難理解的 從一個地域空間的角度去想.
事實上和諧本也沒有那麼到位.
和諧本也只是做求.
只是祈求.
或者說思念.
想一想.
其實你看看.
印給大家看到大綱裡面.
《新漢語》,《日本》或者《NIV》.
也更加貼近原文.
說的是追求.
或者說思想.
甚至是思考.
英文是set your heart.
或者說set your mind.
你看到意思其實一方面.
這個不是一個三分鐘熱度的追求.
或者說想法.
你覺得set是什麼意思.
你刻意去做的.
你是有目標去做.
你是有意志.
你是決心去操練.
你要一生.
你所有的事情都是想著.
怎樣去追求上面的事.
怎樣去思考.

$^{121}$怎樣去實踐上面的事.
並且在這裡保羅很強調的是.
是你們.
即是說是整個教會群體.
不只是傳道人.
不只是教會領袖.
是所有信徒.
在一個身體裡面.
一起去思考,尋求.
怎樣追隨基督.
落實祂的吩咐.
保羅你看到他這裡繼續說.
他說因為信徒已經是.
既然和基督同死.
他們的生命已經和耶穌基督.
一起藏在上帝的裡面.
所以他們不再需要思念地上的事.
當然保羅這裡不是說.
我們信了耶穌之後.
你不需要再想日常生活是怎樣過.
他不是這個意思.
其實你可以看根據常文.
地上的事.
就是針對那些假教師.
那些虛妄的教導.
那些就是屬於地上的事.
屬於肉體的事.
保羅說扔下它吧.
你不需要再思念這些言論.
因為基督就已經是你.
整個人的生命.
是一個整全豐盛的生命.
將來他再來的時候.
我們都會和基督同享.
分享他的榮耀.
而按照下文你會發現.
接下來我們看的就是.
地上的事就是指那些.
會破壞人和上帝的關係.
或者說破壞基督身體.

$^{161}$這些事情.
你看到保羅用了兩個.
都挺嚴厲的命令的語句.
他要我們嚴肅去處理.
包括你治死他.
你棄絕他.
包括一些什麼呢?.
第五節這裡說了幾個事情.
淫亂污穢邪情惡慾和貪婪.
頭四個你一眼就看到.
是和一些情慾.
或者說和性有關的罪行.
保羅很強調.
你們從前.
哥羅西的信徒.
從前未信主的時候.
就是活在這種狀況.
甚至乎做了這些情況.
因為在當時來說.
在一個希臘文化裡面.
其實沒有所謂的.
不犯法的.
很正常的.
每個人都是這樣的.
但是保羅說.
你已經在基督裡面了.
你要脫去這一切的污穢.
脫去這一切上帝不喜悅的事情.
至於貪欲.
或者說貪婪.
為什麼保羅擺到最後.
還刻意說貪婪是太偶像呢?.
他highlight了這個事情.
事實上貪婪我們可以總括為.
人將自己的慾望.
包括剛才說的那四項.
和性和情慾有關的慾望.
你只要將它無限放大.
無限高舉.
高舉到取締了上帝的地位.

$^{201}$你優先追求這些慾望.
多於優先追求上帝.
無論如何你都要得到手的話.
保羅就說.
貪婪或者說貪戀這些情慾.
就變成了你拜偶像.
你所貪戀的事物.
或者你的貪念.
就成為了人去敬拜.
事奉的一個對象.
其實已故的神學家.
Tim Keller.
提摩太海納牧師.
他說過.
貪婪其實.
我們不是那麼容易發現.
這個罪行的.
為什麼呢?.
貪婪它自己很懂得包裝.
隱藏自己.
意思是什麼呢?.
無論人擁有的事物.
到底是否足夠都好.
你想想.
你向上一去比較的時候.
你總會發現身邊有人.
會比你擁有得更多.
擁有得更豐富.
所以一比之下.
你就會開始有個想法.
我現在的生活質素.
我現在擁有的.
其實都不算是什麼.
比起某某人.
其實小小的.
小無見大無.
然後我們會有種想法.
我不是貪.
我只是.
我想要多一點.

$^{241}$我想拿多一點.
我想擁有多一點.
我想享受多一點.
不算過份了.
我又沒有犯法.
我又沒有錯.
你會發現.
人有種特性.
會將自己的貪念合理化.
或者甚至正常化.
Tim Keller 繼續說.
人有另一個特性就是.
他會不斷為自己製造.
和追求心目中的那種偶像.
來滿足自己無窮無盡的貪念.
他表示.
假如我們覺得.
任何比上帝更吸引.
更重要的事物.
我們很想從他們身上.
得到比上帝更多更大的滿足.
一旦失去了這些事物.
你又會忐忑不安.
甚至你覺得.
人生不知道做什麼好.
失去了方向和意義的時候.
Tim Keller 說.
其實他們就是我們從貪念.
所衍生製造出來的偶像.
聖經亦一再強調.
上帝憎惡他的子民.
這種敬拜偶像的行為.
所以你看到保羅這套命令信徒.
你要自死一切這一類.
好像貪婪一樣的拜偶像行為.
徹底和他隔離.
不要給任何機會.
在我們身上滋長.
事實上.
信徒既然已經藉著.

$^{281}$耶穌基督的拯救.
和上帝和好.
保羅再一次提醒.
康復了.
那就不會再做一些.
破壞和上帝關係.
有損基督身體的事情.
所以你看到.
保羅開始轉變他的思路.
不單單說人和上帝的關係.
第八到九節.
其實你看到.
保羅也命令信徒.
要棄絕那些.
會損害基督身體的事情.
損害信徒在基督裡合一的事情.
他同樣.
好像一個對比.
他又說了五件事.
首先是老恨.
最簡單的理解就是.
是一個不適當的情況.
來發脾氣.
至於憤怒.
其實意思是.
你容易暴躁.
特別是暴躁.
一碰你就著.
控制不了.
惡毒.
傷害人的念頭.
毀謗.
對人的一些惡意詆毀.
甚至於咒罵別人.
可能在心裡也不奇怪.
口中的污穢.
就是那些比較低俗.
猥褻甚至帶侮辱性的言語.
難不難想像.
如果這五件事.

$^{321}$都發生在教會的群體.
甚至蔓延的話.
一定會破壞.
信徒在基督裡的合一.
令到基督的身體受到損害.
至於最後一個.
吩咐保羅勸誡.
他們要去扔走的.
你看到.
保羅很鄭重地說.
其實用了另一個命令語句.
他在說.
「不要彼此說謊」.
其實意思是.
你一定要避開他.
你不可以讓他發生.
我們會奇怪.
為什麼雪芳反而.
保羅這麼特別重視呢?.
是不是因為哥羅西教會的人.
經常互相說謊.
經常互相欺騙.
經常互相騙騙呢?.
大概未必是這個向度.
其實我們可以想像.
在教會的場景.
有時說謊的人的動機.
不一定是要全心傷害別人.
不一定是要全心說「我要欺控你」.
反而是怎樣呢?.
說謊或者不說真話.
是不想面對衝突.
是不想破壞大家的關係.
令到對方和自己都不開心.
說謊的另一個面向.
就是隱瞞.
就是不說自己心底裡的真實需要.
就是不為真理去堅持.
所以其實.
謊言.

$^{361}$或者這裡所說的謊話.
雖然它沒有像之前的四種行為.
殺傷力這麼明顯直接.
怒恨你,誹謗你.
這些是直接的攻擊.
你一定是感受到的.
但是這類的謊言.
其實不是這麼容易察覺的.
它沒有這麼直接.
這類謊言的破壞力是間接.
很隱藏的.
但是又影響深遠.
意思是什麼呢?.
可能用一個場景類比.
就是當我們犯錯.
或者你小時候做錯事.
怕父母不開心.
怕他們不高興.
或者甚至乎你擔心被人罵.
被人責罰.
然後你會怎樣呢?.
就不要說.
不要說真相.
你用一個謊言去掩飾你的問題.
但是你背後的動機是.
你想維繫,換取彼此的和諧.
不要衝突.
又或者可能配偶.
假如老公做錯事.
怕被老婆罵.
你還是不要說.
老公明明是想外出吃的.
老婆怎麼都說要在家裡住.
然後你不想面對衝突又怎樣呢?.
隨便吧,但你的臉是黑黑的.
這類其實都會破壞人與人之間的誠信.
有時候我們為了不想說自己真實的需要.
不想破壞一些和諧的關係.
甚至怕說出來.
怕我坦白的時候對方會反對.

$^{401}$甚至乎批評.
於是我們採取的方法就是.
你迎合對方.
強行迎合對方.
你壓制著你內心真正的想法和需要.
我們說這類的謊言其實表面上沒有殺傷力.
即時可能都化解了一些矛盾衝突.
不過其實它是會不知不覺地破壞彼此的互信.
它會漸漸令到那段關係開始冷漠疏離.
因為其實你想想.
每一次我們跟別人說謊言.
其實背後是告訴自己什麼?.
我不是很信得過對方會體諒,接受,理解我的軟弱.
所以我選擇不說真話,說謊話.
我不是很相信對方會理解,接納我一些真實的需要.
我也很擔心我說真話.
對方會否定我,批評我.
甚至乎攻擊我.
所以怎樣?.
我寧願說謊話.
為了不想破壞大家的關係.
我選擇不說話.
但是姐妹保羅她很明白.
我們說好像這類軟性謊言那種殺傷力.
所以你看到她放在最後.
她提醒我們不要做,不要教,不要碰這件事.
因為會破壞我們在基督裡面的合一.
所以你看到保羅他命令我們要停止彼此說謊.
棄絕一切會破壞基督身體合一的事情.
因為保羅接著說.
我們已經脫去舊人和舊人的行為.
並且已經穿上了新人.
保羅用一個我們每天都做的類比.
脫下這件衣服.
穿上這件衣服.
不過脫下那件是一件叫做舊我的衣服.
英文譯做一個old self.
將你舊的我.
即從前未蒙基督拯救,未與神和好的你.
脫下它.

$^{441}$為什麼要脫下呢?.
因為要脫下舊的我.
那些不好的行為.
即上帝不喜悅我們做的事情.
並且也要穿上新的我.
new self.
即是說在基督裡面新創造的你和我.
這個新我並且會不斷地更新成長.
它會漸漸恢復裡面創造主賦予它的形象.
使它在這個更新的過程裡面.
它會越來越認識耶穌.
它也對耶穌的理解更加豐富.
更有力去執行耶穌的吩咐.
所以第十一節.
保羅這裡說這個新我不單止是我們個別的信徒.
也涵蓋整個耶穌基督的新創造.
整個教會群體.
而在裡面保羅說.
不再區分到底是外邦人.
到底是猶太人.
到底是受割禮還是未受割禮.
開化了還是未開化.
到底是圍爐還是自主.
因為在基督這個新創造的教會群體裡面.
無論是種族無論是文化程度.
無論是社會地位.
保羅說已經失去了它原先有的重要性.
因為唯有基督在這個群體裡面.
祂就是一切.
祂最大.
所有人都聽祂的話事.
以祂為一個元首.
所有人都住在祂的裡面.
要聽從祂的帶領和掌管.
換句話來說.
保羅這裡說教會的蒙召.
教會群體的蒙召.
就是要在一個充滿複雜.
多樣性的成員組織底下.
我們要操練在基督裡合一.

$^{481}$保羅這裡用了一個很有趣的字.
他說「蒙召」.
就是說我們一班蒙恩基督教的罪人.
我們走在一起.
互相建立關係成為一體.
其實是上帝簡算了.
上帝叫我們這樣做.
本身有沒有難度?.
已經有了.
一班罪人要走在一起.
在基督裡合一.
不容易的.
你要靠基督.
我們靠自己一定是不容易的.
再加上每一個人都有他不同的背景.
不同的複雜性.
就更加難.
常加難.
所以第十二節.
保羅他很鄭重提醒我們.
他說我們之所以走在一起成為一個身體.
是因為上帝愛我們.
所以他選擇我們.
並且將這個分別為性的合一重任.
這個這麼有難度的重任.
交給你和我.
所以也為這個緣故.
保羅他提醒我們.
當我們要實踐合一的時候.
你一定要穿上這五樣的屬靈品格.
同樣是五樣.
好像一個對比.
平衡這樣.
第一個就是一個有憐憫的心腸.
就是說我們願意去體恤人的軟弱.
或者說人一些的困境.
第二個是恩慈.
對人.
人慈.
但是要留意.

$^{521}$恩慈不是無限的人慈.
在基督裡的人慈.
是有屬靈原則的.
不會過分的感情用事.
第三是謙卑.
謙卑是承認.
我們每一個人在上帝面前.
本來就是有罪.
也都完全要依靠耶穌基督的恩典.
所以本質上.
我們不給別人優勝.
不給別人優越.
溫柔.
溫柔其實很有意思.
他不是純粹說你好好.
他是說在適當的時候.
你對著適當的人.
你因為適當的原因.
表達你適當的情緒.
這個叫溫柔.
就是說溫柔不是一味做一個順民.
不是一味順人.
溫柔是我們會遊走在.
你不要過分的憤怒.
也都不要過分不懂得憤怒.
因為憤怒是上帝給我們的一個前水.
你是在中間.
因著基督的緣故.
因著聖經的教導.
你是在裡面拿捏的.
耶穌都會發怒的.
你要懂得在兩邊.
去採取一個平衡.
中庸之道.
這個為之溫柔.
最後忍耐.
我們很明白這個是來自上帝.
他自己本身的屬性.
我們經常說.
上帝是一個不輕易發怒的上帝.

$^{561}$你不容易激怒祂的.
祂會生氣.
不過不是那麼容易被你激怒.
忍耐是說.
我們明白神對我們的忍耐.
也都成為我們對人.
恆久忍耐的一個基礎.
和原因.
當然忍耐不是說.
我們死忍,啞忍.
隨它吧.
其實忍耐背後有個含義.
就是忍耐是我不急於.
馬上解決所有問題.
包括人際的問題.
我們說忍耐其實是by time.
你爭取時間.
你先預備好自己.
你都求上帝預備好對方.
然後配合上帝的時機.
逐漸改善那些不理想的情況.
那些不合上帝心意的關係狀況.
其實這個才是忍耐的極致.
不是純粹死忍.
保羅其實他用了一個例子.
我們都會明白.
如果信徒之間有嫌隙的時候.
那怎麼辦呢?.
如果用從前那個舊的我.
用你Old self的時候.
你就好像八字到九字那樣.
那些老奴,惡毒都出來了.
保羅說現在既然我們已經穿上了.
那個新的我.
就要跟新的方式來應對.
保羅這裡說.
總要互相的寬容.
彼此的饒恕主.
怎樣饒恕你們.
要怎樣饒恕人.

$^{601}$這裡提醒我們.
彼此饒恕一個的基礎是.
我們自己每一個本身都是怎樣.
我們都是夢耶穌基督.
施恩寬恕的罪人.
所以當我們彼此有嫌隙的時候.
我們自然也要效法基督.
以恩慈來相待.
去互相的寬容.
最後保羅他說.
穿上了.
除了說穿上了這五個的屬靈品格.
保羅強調最重要的是什麼呢?.
新漢語本來也是他比較到位的.
是穿上愛.
要穿上這件衣服.
就不是純粹要愛.
是穿上.
你明白我的意思了.
就是說我們自己做不到的.
你穿上基督給你的愛.
穿上上帝給你的愛.
因為愛就是教會連結在一起的原動力.
我們很明白有一本書說.
我們愛因為什麼?.
神才愛你.
你沒有離開上帝的愛.
你基本上沒有力量去愛其他人.
又或者說.
神的愛就是信徒合一和人際關係那種的編合劑.
一切的屬靈品格都是源於愛.
我們也記得勝靈九果子第一個是什麼?.
就是人愛.
或者說保羅說.
要有信,有望,有愛最大就是什麼?.
其實保羅他自己也是這麼說的.
你會發覺事實上上帝的愛.
好像一種膠水一樣.
當你快要散散地了.
有裂的時候怎麼辦?.

$^{641}$上帝的愛就包著.
就包著這個身體.
讓它會有不同.
甚至乎會有不衝突.
但是不要粉碎.
不要爛了.
讓每一個信徒能夠彼此尊重.
也能夠接納到對方和我是不同的.
剛才保羅也說了.
我們是來自不同的層面.
有不同的性格.
不同的看法.
其實就在上帝的愛底下.
在聖經的真理底下.
我們怎樣夢照去努力求同存異.
在基督裡合一.
甚至保羅繼續說.
他說除了愛之外還有一樣東西很重要.
你要讓基督的平安掌管你的心.
意思就是說.
你要用那個心.
他不是純粹在說這裡.
沒有這麼虛的.
他是在說.
掌管我們的思想.
我們的情感.
和我們的意志.
你要有基督的和平在裡面.
來引導你做正確的決定.
特別當教會的群體.
出現可能意見不同.
甚至乎所謂不和衝突的時候.
每一個人首先需要有來自耶穌的平安.
來維持群體的平安.
也需要有基督的平安.
來處理人與人之間的矛盾.
甚至乎嫌隙.
也憑著基督所賜的平安.
在裡面做合適的判斷.
甚至乎決定.

$^{681}$我記得以前在神學院.
有一位我很尊敬的老師教我們.
提醒我們.
譬如像今天這樣.
你有沒有看天氣報告做人?.
有沒有穿多了衣服?.
有,但是你發覺不是這麼冷.
為什麼?.
經常都是這樣.
我們會看天色做人.
看天氣做人.
但是那位老師提醒我們.
我們有沒有看上帝面色做人?.
他提醒我們.
有時我們特別可能做信徒.
我們不知不覺會看人面色做人.
比看上帝面色做人多.
因為上帝你看不見.
上帝不輕易發怒.
你有時覺得算了.
上帝得失了.
上帝暫時沒事的.
但得失人就馬上有後果.
所以有時我們會讓人的面色.
主導你做正確的決定.
那位老師提醒我們每一個.
假如真的在討上帝喜悅和討人喜悅.
兩者之間不能夠兼得的時候.
那怎麼做?.
耶穌說了.
耶穌自己在哈西瑪利園實踐這件事.
我不能夠在我的意願.
耶穌說和上帝的意願之間.
最後他掙扎完之後.
照你的意思.
其實就是這樣.
我們固然很想兩者兼得.
這是最理想的狀況.
又可以討好人,更可以取悅上帝.
但如果真的有些情況.

$^{721}$不可以兩者兼得.
上帝就期望我們以祂優先.
在你做判斷的時候.
你堅持先過到上帝.
再過到人.
當然是不容易的.
所以為什麼普羅打了底.
上帝呼召我們做這件事.
上帝愛我們才做這件事.
所以有時上帝讓你遇到一些近的關係.
你試試反過來想.
上帝呼召你.
上帝愛你.
其實你會有些轉化.
你會有些韌力.
好好去相處.
當它是一個藝術.
有來有往.
普羅說其實在這些時候.
我們真的很需要基督所給我們的平安.
因為可能在和對方堅持.
甚至在討論爭拗過程.
其實大家都不舒服.
大家都會不開心.
你又會覺得難受.
對方又會覺得難受.
在這個時候.
基督的平安額外顯得重要.
一方面.
要符合雙方的內心.
是靠基督.
另一方面.
也因為這份內在平安.
讓我們可以確信.
當我們真的沒有辦法.
有時不能夠認同的時候.
當我們要本著真理做決定的時候.
你確信上帝是會喜悅你.
上帝是會對你和顏悅色.
祂是會給你讚好.

$^{761}$支持你繼續堅持.
那時候上帝的平安就很重要.
讓你做正確上帝喜悅的決定.
所以普羅這裡說.
基督的愛和平安就成為了.
信徒之間和而不同.
成為一體的關鍵.
普羅再一次提醒.
我們呼召.
我們夢召.
就是為了這樣.
最後.
普羅再說一樣.
我們還要心存感恩.
既有愛.
最重要是愛.
還要加上平安.
還要有心存感恩.
感恩我們當然明白.
對象就是上帝.
沒有焦點的.
那個不一樣的你和我.
放在一起.
就是學習在基督裡.
這個成長合一的功課.
然後我們在裡面.
上帝的形象.
就會呈現出來.
事實上普羅這個提醒很重要.
或者說在實際上我們都會明白.
一定是怎樣?.
沒有那麼暴躁的.
比起一個常常抱怨.
甚至乎發怒批評的人.
容易相處一點.
如果那個人那麼暴躁.
你碰一碰就著.
你一定會跟他保持十尺遠.
你不會去碰他.
其實一個常常感恩的人.

$^{801}$他會很自然地流露出.
基督的愛和基督的寬容.
或者說面對一些相處上的問題.
其實他可以用感恩來化解.
轉化了他.
可能裡面會很生氣.
甚至乎裡面也會生氣.
不過他知道是上帝促成.
容許這種狀況.
要在裡面學習功課的時候.
他會用感恩來轉化這種負面的怒氣.
事實上感恩最核心的關鍵是.
我們每一個已經.
藉著基督的拯救和神和好.
得到永生的盼望.
徹底解決了.
你和我終極要面對最大的問題.
但是自己又沒有能力解決那個問題.
就是死亡和上帝隔絕.
我們感恩上帝幫我們解決了.
你人生最大的難題.
所以相比之下.
其他一切人際的問題.
教會弟兄姊妹相交裡面.
可能會碰一碰的問題.
那些就怎麼樣?.
就可以輕看.
所以做基督徒.
一個感恩的基督徒.
理論上應該慢慢學習.
不用那麼多計較.
有時候不需要太過跟著計較.
太過執著.
在一些事情上.
弟兄姊妹其實.
本來這裡說感恩.
他給了我們一個很正面的能量.
一個很正確的做人態度.
促進我們和上帝.
甚至和弟兄姊妹的關係.

$^{841}$時間關係.
其實本來還有兩節我要說.
不過說不了.
因為會過期.
我留給大家在查經的時候.
再和弟兄姊妹互動.
今天你會發覺這段經文.
涉及很多和關係有關的東西.
我們做一點應用反省.
首先當然一定是.
檢視我們和上帝的關係.
在日常生活裡面.
很真誠地問自己.
留心一下.
有沒有一些你正在做的行為.
甚至想法.
其實不知不覺.
正在破壞我們和上帝的關係呢?.
這是很直接的應用.
有沒有一些事情.
你明知道不可以做.
不過你都在做呢?.
中國人有一句話.
物因小善而不為.
另一句就是物因小惡而為之.
就是說如果應用在我們的信仰裡面.
他提醒我們.
在上帝的眼中.
即使善行再微小.
你實踐出來又如何?.
你就一定建立和上帝的親密關係.
相反.
在上帝眼中的惡行.
無論多小,多不察覺.
你做出來又如何?.
一定破壞你和上帝已經好的關係.
所以弟兄姊妹.
其實上帝既然這樣愛我們.
我們很自然也要以愛來回應他.
從今天起.

$^{881}$特別趁著快過年了.
我們會不會努力為上帝的緣故里.
棄掉一個不好的習慣?.
棄掉一些不太適合聖經真理的想法?.
並且為上帝的緣故里建立一個新的行為.
一個新的習慣.
讓你和上帝的關係更加美好.
當然我們要明白.
我們的靈命不會突然有一天突然改變.
不會的.
而是某一天.
你下定決心.
你終於願意為主的緣故.
靠主來做改變.
當你願意這樣做的時候.
你就是在說.
我不再將我有限的人生和精力.
耗費在那些沒什麼益處.
沒什麼養分.
甚至上帝不喜悅的事情裡.
如果你現在開始打掃你的屬靈心靈的時候.
將來你就會有更多空間.
讓耶穌住在裡面.
你裡面上帝的形象.
就會一路一路呈現出來.
而第二方面.
其實經文也提醒我們.
在人際關係裡.
我們要做一些稀有更新的事情.
檢視我們和人的關係距離是否恰當.
我們有沒有適當地為每一段的人際關係.
來設立一個適當的界線.
既可以彼此靠近合作.
甚至分享.
同時也為有需要的時候.
雙方有一個安全的空間距離.
特別為耶穌基督的緣故.
為真理的緣故.
可能為一些合理需要的緣故.
我們要懂得溫柔.

$^{921}$但是堅定地拒絕別人一些期望和要求.
甚至是一些不合理的要求.
適當地,幽默地.
好像幾年前有一本書叫做《被討厭的勇氣》.
有沒有聽過或者看過?.
其實如果用回屬靈裡.
有時候我們拒絕別人.
你是要有少許幽默化解它.
適當地為主堅持.
將來如果你有信主.
他有信主的時候將來都會怎樣?.
在上面見.
不會散了.
這個才是最重要的.
但是你適當地為主堅持.
幽默地化解了它.
這樣就可以彼此仍然在群體裡.
不會各生東西.
事實上弟兄姊妹學習在基督裡合一.
是需要情理兼備的.
更加需要在真理的光照裡做一些溝通.
讓彼此的對話首先更符合上帝的心意.
當然一定會遇到意見不同.
甚至乎所謂不理想的情況.
你需要適當地保持距離.
讓對方都可以有一個冷靜.
也要維護雙方的界線和尊重.
事實上所謂獨屬靈智慧.
其實就是我們不會立即將彼此的差異.
敵對化了它.
我們不會純粹因為對方和我不同.
你就是錯.
不會因為對方和我不同.
我就和他死拗.
差異很多時候只是.
兩種不同角度的良好意願.
剛好走在一起.
屬靈智慧也是懂得.
在衝突中你會調教自己.
你會修正你自己.

$^{961}$你會將人與人之間的摩擦.
用屬靈的角度去看.
你會變成一個磨練.
你會視彼此的差異.
是一個靈命成長的階梯.
最後這段經文也提醒我們.
你為了你和對方的不同.
確實教會的成員.
是會因為彼此的差異和不同.
走在一起特別要埋身配搭的時候.
是會有些碰撞.
甚至所謂有爭拗,不和.
是會很激氣.
但同樣也是因為這份差異.
我們才可以彼此互相造就.
上帝是否很奇妙?.
正正這就是他智慧的創造.
從宏觀來說.
《創世記》說哪人獨居不好.
其實上帝已經表明了.
我們每一個都需要透過關係.
透過群體.
透過彼此的差異.
來完成祂的吩咐.
在這個前提下.
我們只有一條路.
有時會不舒服.
但從上帝的角度.
知取恩典,知取愛,感恩.
我們在裡面學習.
互相去理解.
嘗試欣賞對方.
嘗試去感恩.
慶幸祂將不同的你和我放在一起.
變成你和我的人生.
不是只有一種顏色.
不是只有一把聲音.
在信仰的路上.
我們欣賞這種多樣性.
原來可以彼此豐富對方的生命.

$^{1001}$會激發更多的潛能.
並且互相保衛.
弟兄姊妹盼望.
藉著今天這段經文.
我們將這個願意互相理解.
欣賞,感恩的態度.
不只是用在教會的場景.
其實會明白.
特別下一章保羅就說.
用在你其他生活場景.
在你的夫妻關係裡面.
在你的親子互動裡面.
在你和人.
和家人的關係.
和朋友的關係.
多一點的欣賞.
多一點的稱讚.
其實讓你所身處的群體.
成長得.
按著上帝的旨意.
成長得更好.
發展得更美.
滿有耶穌基督長大成人.
那個的新娘.
我們同心祈禱.
主耶穌我們多謝你.
將我們從罪類徹底拯救.
我們都感恩.
因為我們每一個都是蒙簡選.
在上帝的旨意.
在教會這個群體裡面.
去學習彼此的合一.
亦都由此延伸到我們每一個身處.
你託付給我們的崗位和角色.
包括可能我們的家庭.
我們工作的地方.
讓我們在裡頭.
藉著耶穌基督的愛.
藉著你所賜的平安.
所以我們可以在當中.

$^{1041}$學習怎樣去和人相處.
我們在當中怎樣維繫彼此.
基督裡同一個身體.
上帝幫助我們.
讓我們藉著今天的經文.
我們能夠去審查.
審查和你的關係.
和人的關係.
在裡頭有沒有一些.
不單止是不能夠理解.
甚至乎是破壞.
近體合一的一些情況.
願上帝你光照我們.
也願上帝你慢慢引導我們.
讓我們真的順著你的帶領.
學習彼此的合一.
求主你與我們同在.
因為我們知道.
這是我們一生的功課.
讓我們都敏銳.
上帝你預備的時機和安排.
在過程的裡面.
讓我們知道.
我們夢照就是要維持.
為的是要將裡頭.
耶穌基督.
你所賦予我們上帝的形象.
展現出來.
求主幫助我們.
讓我們真的.
藉著新一年的來到.
我們的靈命都能夠氣候更新.
多謝主.
誰聽我們的禱告.
奉耶穌你的名去祈求.
Amen.
\newpage



\section{歌羅西書 3:18-4:5}
\label{sec:jemYOQSwajw}
\textbf{2026年2月15青少主日崇拜直播｜楊珊珊傳道｜門徒DNA:我們這一家｜歌羅西書3\_18-4\_5}
\newline
\newline
連結: \href{https://youtube.com/watch?v=jemYOQSwajw}{\texttt{https://youtube.com/watch?v=jemYOQSwajw}} ~~~~ 語音日期: 2026-02-16
\newline
\newline
\hyperref[sec:397tpwFaRoA]{\small{< < < PREV SERMON < < <}}
~
\hyperref[sec:index_chronic]{\small{[返順時目]}}
~
\hyperref[sec:NhZRAevmHb4]{\small{> > > NEXT SERMON > > >}}
\newline
\newline
歌羅西書 3:18-4:5
\newline
\begin{longtable}{cl}
\hline
\hline
章節 & 經文 (和合本修訂版)\\
\hline
3:18 & \begin{tabularx}{0.7\textwidth}{X} 你們作妻子的，要順服自己的丈夫，這在主裡面是合宜的。 \end{tabularx} \\ \\ \relax
3:19 & \begin{tabularx}{0.7\textwidth}{X} 你們作丈夫的，要愛你們的妻子，不可虐待她們。 \end{tabularx} \\ \\ \relax
3:20 & \begin{tabularx}{0.7\textwidth}{X} 你們作兒女的，要凡事聽從父母，因為這是主所喜悅的。 \end{tabularx} \\ \\ \relax
3:21 & \begin{tabularx}{0.7\textwidth}{X} 你們作父親的，不要惹兒女生氣，恐怕他們會灰心。 \end{tabularx} \\ \\ \relax
3:22 & \begin{tabularx}{0.7\textwidth}{X} 你們作僕人的，要凡事聽從你們肉身的主人，不要只在眼前服事，像是討人喜歡的，總要心存誠實，因為你們敬畏主。 \end{tabularx} \\ \\ \relax
3:23 & \begin{tabularx}{0.7\textwidth}{X} 你們無論做甚麼，都要從心裡做，像是為主做的，不是為人做的； \end{tabularx} \\ \\ \relax
3:24 & \begin{tabularx}{0.7\textwidth}{X} 因為你們知道，從主那裡必得著基業作為賞賜。你們要服侍的是主基督。 \end{tabularx} \\ \\ \relax
3:25 & \begin{tabularx}{0.7\textwidth}{X} 行不義的人必受不義的報應；主並不偏待人。 \end{tabularx} \\ \\ \relax
4:1 & \begin{tabularx}{0.7\textwidth}{X} 你們作主人的，待僕人要公正，因為知道，你們也有一位主在天上。 \end{tabularx} \\ \\ \relax
4:2 & \begin{tabularx}{0.7\textwidth}{X} 你們要恆切禱告，在禱告中警醒感恩。 \end{tabularx} \\ \\ \relax
4:3 & \begin{tabularx}{0.7\textwidth}{X} 同時，也要為我們禱告，求神給我們開傳道的門，能宣講基督的奧祕， \end{tabularx} \\ \\ \relax
4:4 & \begin{tabularx}{0.7\textwidth}{X} 使我能按著所該說的話將這奧祕顯明出來，我為此而被捆鎖。 \end{tabularx} \\ \\ \relax
4:5 & \begin{tabularx}{0.7\textwidth}{X} 你們要把握時機，用智慧與外人來往。 \end{tabularx} \\ \\
[1ex]
\hline
\hline
\end{longtable}
$^{1}$我們以聆聽聖道敬拜.
今日的經訓是在.
《歌羅西書》三章十八至四章五節.
你們作妻子的.
當信服自己的丈夫.
這在主裡面是相宜的.
你們作丈夫的.
要愛你們的妻子.
不可輔待他們.
你們作兒女的.
要凡事聽從父母.
因為這是主所喜悅的.
你們作父親的.
不要惹兒女的氣.
恐怕他們失了志氣.
你們作婦人的.
要凡事聽從你們育身的主人.
不要只在眼前事奉.
像是土人喜歡的.
總要全心誠實敬畏主.
無論作甚麼都要從心裡作.
像是給主作的.
不是給人作的.
因你們知道從主那裡.
必得著基業為賞賜.
你們所事奉的乃是主基督.
那行不義的必受不義的報應.
主並不偏待人.
你們作主人的.
要公公平平的待僕人.
因為知道你們也有一位主在天上.
你們要行切禱告.
在此警醒感恩.
也要為我們禱告.
求神給我們開傳道的門.
能以講基督的奧秘.
叫我按著所該說的話.
將這奧秘發明出來.
你們要愛惜光陰.
用智慧與外人相交.

$^{41}$願神的話語常在我們心裡.
弟兄姊妹平安.
活出基督的生命.
不是靠知道多少的神學理論.
而是在日常生活中.
在最親近的人的面前.
很切實地活出主的見證.
上個主日Ray牧師提醒我們.
既然與基督一同復活.
就要脫去舊人.
穿上新人.
讓基督的平安作決定.
今天我們看《歌羅西書》三章十八節.
至四章一節.
上半部分保羅說.
這個新人的樣式.
就是要從家庭生活開始實踐出來.
為什麼呢?.
因為我們在家裡.
沒有辦法裝模作樣.
因為只有當我們的信仰.
能夠在最接近我們身上.
即身邊最親近的家人.
活出來.
這個信仰才是真實.
你是否在基督裡面.
是一個新造的人呢?.
看你家裡的表現.
就應該知道了.
保羅說完最貼身的家人.
就推到最遠的外人.
非信徒.
今天下半部分.
會說四章的二至六節.
多一節.
提醒我們.
要專心圍到宣教的門.
禱告.
不過今天的主題是.
「我們這一家」.

$^{81}$說家庭.
所以讓第二部分.
我們會集中看第五至第六節.
我相信每一位弟兄姊妹回來.
應該身邊也會有一些未信主的親友.
讓我們一起跟隨保羅學習.
如何跟外人,非信徒交往.
讓我們在聽到之前.
一起有個祈禱.
親愛的天父.
做你兒女的福分.
求聖靈幫助我們脫虛救人.
穿上新人的時候.
我們能夠在家庭裡.
以耶穌基督為主.
活出屬你群體的樣式.
求你預備我們眾人的心.
幫助孩子能夠講解得清楚.
也幫助聽到的能夠明白.
讓我們都能夠被你的話語根深.
奉主耶穌基督得勝的明智祈求.
阿們.
保羅在《歌羅西書》.
將焦點放在一個根本性的問題.
就是在家裡究竟誰才是真正的主呢?.
《歌羅西書》三章十八節至四章一節.
也有人說這是一個家庭的守則.
保羅要我們看見.
當基督成為一家之主的時候.
家庭成員之間的關係.
是會怎樣去改變的.
其實在希羅文化的社會裡.
一個完整的家庭.
他們會有奴隸和自由人去組成.
家庭裡最基本有三組的關係.
就是丈夫和妻子.
父親和兒女.
主人和奴僕.
當時家庭的秩序.
其實和我們中國人的三綱五倫很相似.

$^{121}$都是先提及地位較高的一方.
就是主人,父親和丈夫.
然後先提及經常被忽略.
或者地位較低的那一方.
就是僕人,兒女或妻子.
但保羅在這裡卻是完全相反.
他先提到妻子.
才提到丈夫.
先提到兒女.
接著才提及父親.
才提及僕人.
接著再提及主人.
這個順序其實本身已經透露了.
保羅的關注點.
其實是那些在家裡看似次要的.
軟弱的.
又或者是無權的人.
我們一起看看第一對的關係.
就是在說妻子和丈夫.
三章的十八字.
新譯本是這樣說.
作妻子的要馴服你們的丈夫.
在主裡是合二的.
馴服是有受管轄的意思.
在當時羅馬的家庭.
其實管轄妻子本身就是每一個丈夫的權力.
如果妻子不馴服的話.
其實當時的丈夫是可以體罰.
經濟制裁就是不給家用.
又或者是休妻.
所以那個是古希羅文化的世界.
不是香港.
所以大家千萬不要這樣做.
保羅如果只是要求妻子去馴服丈夫的話.
根本是多此一舉.
因為社會就是這樣要求.
所以告訴我們.
其實這段經文的重點.
不是在前面半句的命令.
而是在後面半句的理由.

$^{161}$保羅沒有說.
妻子馴服丈夫是社會合二的.
而保羅就是說.
這個在主裡是合二的.
這種馴服的動機是很不同的.
因為保羅對太太的姐妹說.
其實標準不是由社會而來.
而是基督那裡來的.
焦點不是丈夫.
而是主基督.
妻子不是為了討好.
又或者是很害怕丈夫才去馴服.
而是為了主的緣故而馴服.
在主裡也代表了馴服是會有條線的.
有個界限的.
當丈夫的要求違背了主的心意的時候.
妻子首先是要馴服主.
弟兄姊妹,我們要留意保羅在這裡用詞的區別.
有些不同的.
保羅對妻子說要馴服.
但對兒女和僕人說的是什麼?.
要凡事聽從.
這兩個字是不一樣的.
聽從是一個單方面.
是一個全面的馴服.
全面的服從.
但馴服是不同的.
保羅在《爾忽所書》五章二十一節.
他說丈夫和妻子要全敬畏基督的心.
彼此馴服.
所以馴服是可以彼此的.
但聽從是不可以彼此的.
這種馴服其實不是一種單方向的服從.
不是屈服.
而是一種自我的限制.
妻子主動將丈夫視為自己的界限.
在重大的決定上.
願意考慮丈夫的大齡.
我認識很多姐妹很有才幹.
美貌與智慧並重.

$^{201}$甚至是專業人士.
但很多時候她們都願意選擇.
不用這些來壓過她們的丈夫.
貶低丈夫,操控丈夫.
保羅勸勉妻子.
千萬不要出於勉強去馴服丈夫.
而是因著基督的緣故.
馴服這個一家之主.
那作為丈夫的又應該怎樣做呢?.
第十九節.
作為丈夫的要愛你的妻子.
不可以苛待她們.
聖經的教導非常清楚.
不可以苛待你的妻子.
如果我們將基督與教會的關係.
類比成為丈夫與妻子的關係.
你就會發現.
基督是不會虐待教會的.
相反,基督是愛教會.
甚至是為了救贖她而捨己.
走上十字架的道路.
這種愛的呼喚.
是叫我們放下自己的需要.
自己的慾望.
不再以自己的滿足擺在先.
這種愛不是佔有.
而是捨己.
不是講求從對方的身上取得甚麼.
而是講求自己能夠付出甚麼.
這種持久衛生的愛.
是不會自然地發生.
而是靠每一天的意志去實踐.
這種就是令妻子樂於馴服的愛.
我聽過一個傳道人的見證.
有一次,這個傳道人和他太太吵架得很厲害.
不過他仍然要事奉.
他仍然去舉辦報道會.
宣講神與人和好的訊息.
講完之後,還有一大堆人信主.
有些人立即站起來去陪談.

$^{241}$當時他坐下來.
心裡很不舒服.
因為他和師母吵架了.
還未和好.
他想起他剛剛講的訊息.
是一個神與人和好的訊息.
他發現自己和自己最親近的人.
還在吵架.
他吐告之後,打電話給他太太.
她說「對不起,我沒有好好地去愛你」.
「我在外面做傳道人」.
「其實我回到家裡,也應該做一個傳道人」.
「我在外面教人怎樣去愛」.
「但是我自己卻不知道怎樣去愛人」.
「對不起,我做錯了」.
「我需要改變」.
「你可不可以幫助我學習怎樣去愛你?」.
師母聽到這裡,她流著眼淚.
她說「老公,其實我不需要你的錢」.
「也不需要一直去旅行,買東西,吃東西」.
那一刻,那個傳道人才發現.
原來他太太需要的是很簡單.
就是她能夠順靠主.
按著福音的教導,好好帶領這個家.
不過這件事,她卻未能做好.
弟兄姊妹,怎樣去愛一個人?.
就是當我們活在一個愛的狀態裡.
我們才懂得怎樣去愛.
即使我有情緒上的需要.
有不同的慾望.
想控制我的伴侶來滿足我自己.
但是因為我心裡有神的愛.
我也看到另外一個人.
對方是按著神的形象去做的.
對方也是神所愛的孩子.
於是這位傳道人跟他的妻子說.
「親愛的,主耶穌很愛我們」.
「我們裡面有神的愛」.
「我們依靠祂,我們會過得到的」.
其實不是因為他們兩個人的生命很好.

$^{281}$是完全相反.
他們兩個人的生命簡直是一塌糊塗.
但如果沒有神的愛.
他們甚麼都做不到.
當丈夫和妻子一起跟著主耶穌.
一起依靠祂.
他們就能夠走在同一條路上.
結伴同行.
夫婦就成為了很好的夥伴.
很好的朋友.
一起跟著神的呼召.
一起變得更像主耶穌基督.
弟兄姊妹,我們愛.
因為我們選擇信服基督.
所以愛和信服同樣是一個命令.
這個命令本身都是愛.
神愛我們.
所以祂這樣命令我們.
我們信服這個命令.
因為這是愛的真諦.
第二對的關係是兒女和父親.
又或者可能是兒女和父母.
又或者有不同的家庭組合.
可能是兒女和母親.
或者是兒女和家長.
作為兒女我們要事事聽從父母.
因為這是主所喜愛的.
保羅在這裡勸勉作兒女.
要凡事聽從父母.
沒錯,這次用的字眼不是信服.
而是聽從.
即是聽話,亦是服從.
是單向地遵從.
在上位的父母權威的命令.
而且不只是在他們很有道理的時候.
聽他們說的話,而是凡事都聽從.
當然這裡有一個假設.
假設就是父母要按著.
《申命記》四章八節和六章七節的吩咐.
藉著新教和賢教.

$^{321}$將神的律法傳遞給子女.
就好像大家記得上年.
其實青少主日我也在這裡和在座的家長和大人呼籲過.
我們就是年輕人的第一本聖經.
雖然聽從父母符合社會標準的規範.
亦是我們中國人常說的很孝順的表現.
不過更加重要的是我們背後的動機.
因為這是主所喜悅的.
信徒做事和對待其他人的原則.
就是看主究竟喜不喜悅.
主所喜悅的我們就做.
主所不喜悅的我們就不做.
所以除非我們的父母.
他們要求我們不要主的喜悅.
例如叫我們做壞事.
又或者叫我們不要信主.
否則只要你屬於這個家庭.
作為兒女的重要是凡事聽從父母.
不是視乎父母的性格或他們的要求.
因為父母和兒女有好的關係.
這也是我們有的神好喜悅.
所以保羅再次將信徒的品格和生活.
與基督連繫起來.
帶出基督是這個家庭的主權.
作為父親,父母,母親,家長又應該怎樣呢?.
聖經說不要激怒兒女.
免得他們灰心喪志.
在保羅的時代希羅的世界.
父親在家中負責教導兒女.
所以保羅沒有提及母親的責任.
因為主要教導的工作都放在父親身上.
但來到21世紀的香港.
我們很少有這種角色上的定型.
所以父母雙方都有責任共同教育子女.
作為父親不要激怒兒女.
其實大家是否覺得被家人激怒.
是一件非常容易發生的事情.
我相信每一個家庭都會發生.
兒子或女兒激怒父母.
或父母激怒兒子或女兒.

$^{361}$所以保羅的重點不是叫我們不要激怒兒女.
而是背後有一個原因.
原因是我們的兒女會灰心喪志.
不知道大家被人批評一次.
通常都沒有問題.
我們不會灰心.
但如果我們被人批評很多次.
特別是那些拆毀的聲音.
是來自身邊最親的人.
我們就會開始質疑自己的價值.
從而對事,對人或對生命失去動力.
年輕人其實很重視父母怎樣看待他們.
如果兒女已經盡力.
但只能緩和父母的怒氣.
永遠得不到父母的微笑.
如果兒女已經盡力.
但得不到父母的喜悅.
父母仍然覺得.
A或B的兒女比自己的兒女好的時候.
兒女就會灰心喪志.
大受打擊.
幾年前香港有一套電影《年少日記》.
講述一位哥哥.
無論他多努力讀書彈琴.
換來的只是母親的哭泣聲.
還有父親的藤條蚊豬肉.
他不能像表現優秀的弟弟般.
最終這個孩子對生命失去了盼望.
造成了一個悲劇.
弟兄姊妹.
這不只是電影劇情.
而是一些家庭的真實寫照.
雖然古代的父親擁有絕對的權柄.
去管治家庭成員.
但保羅並非吩咐父親運用這些權柄.
去管治兒女.
令他們聽父母的吩咐.
保羅對父親的訓練.
就是他們應該從兒女的利益出發.
為父親在兒女身上運用權力擺一條線.

$^{401}$父母一定要謹慎.
讓子女配得自己的聽從.
我們成為父母或做父母的弟兄姊妹.
都可以想想辦法.
讓子女聽自己說是一件容易的事情.
要避免對兒女有一些不合理的要求.
或者太過苛刻的管教.
令他們很灰心和喪志.
而不願意繼續遵從父母的教導.
最慘的是逐漸失去了.
在主內應該開心的一種做事的動力.
第三對的關係是講到僕人和主人.
這裡描述不單是上司和下屬的關係.
而是基督教家庭的一種關係.
一家之主不是建基於盟約或血緣關係上.
其實勞僕和主人的關係是羅馬制度獨有的關係.
保羅是無意推翻社會的制度.
但他比其他關係多了四字經文.
去鼓勵或勉勵在家庭權力階級最低層的僕人.
為什麼呢?.
因為在哥羅西教會的弟兄姊妹當中.
都是一些勞僕的信徒.
例如阿尼西姆是肥理門的奴隸.
與妻子和兒女不同.
奴隸當時被視為有生命的財產.
其實是走不走得的物件和財產.
主人可以任意對待他們.
他們亦經常遇到不公平和不公義的對待.
無論付出了多少都不會有獎賞.
不過保羅在二十四節勉勵這些僕人.
「你們是為主作公,所以一定會從基督裡得著,積業著,獎賞」.
反之,一些行不義的人.
神一定會按著他們的不義去報應在他們身上.
所以做主人要公平公正地對待僕人.
後面的原因讓我們一起讀出四章一節.
準備,3,2,1,起.
「做主人的,對待你們的僕人要公正公正.
要知道你們也有一位主在天上」.
謝謝弟兄姊妹.
這個原因讓我們看見保羅最終想說的訊息.

$^{441}$家庭的規章最重要的重點就是七個字.
「基督是家庭的主」.
基督徒其實不應該有兩個標準.
我們不可以用一個標準去看待其他人.
但我們又期望我們的主用另外一個標準去看待我們.
因為基督是一家之主.
主內的夫妻,親子和主僕關係才會出現變化.
保羅給了一家之主很多個不可以,不好.
其實是要防止他們濫權.
來到21世紀,丈夫和父親其實很少都是這種至高無上的存在.
至少我在家中,我爸爸其實是一家之主.
不過是煮飯的主.
不過我們也會去想.
其實會不會有其他家庭成員成為了萬人之上沒有王館的一家之主呢?.
我們其實已經沒有羅馬時代的主僕關係.
但我們在家庭中也一樣有同一個主人.
無論我們是什麼角色也好.
那個主人是主耶穌基督.
我不知道大家在家中有沒有一塊牌寫著七個字.
「基督是我家之主」.
有沒有弟兄姊妹有這個牌?.
有吧.
我有一位信義代的好姊妹.
她也是一個傳道人.
她家門口其實也有這塊「基督是我家之主」的裝飾.
我記得她跟我分享過她大學四年級時的一個故事.
她說小時候她每次做錯事.
她都會跟媽媽說對不起.
之後媽媽又會帶著她跟天父說對不起.
因為看到他們犯錯,做錯事.
魔鬼最開心.
天父就最不開心.
長大後不知變了什麼時候.
當我這個朋友跟別人說對不起的時候.
都覺得說了就算了.
真的大了,當作沒有這回事.
有一天她爸爸去了美國公幹.
她,姐姐和媽媽竟然鬧了一場史無前例這麼大的架.
於是媽媽就要遙距Zoom邀請爸爸去主持一個遙距的會議.
所以我這個朋友在晚上很不服氣.

$^{481}$就在客廳開了爸爸在美國舉辦的家庭視像會議.
於是他們就在說著說著.
搬了很多他們不和的東西說出來.
他們逐一陳詞,逐一道歉.
其實每個人都是各執己見.
很無奈和很不爽地看著對方.
於是爸爸一開始就提醒他們.
其實我們家門口掛著「基德是我家之主」.
我這個全良朋友心裡想又來了.
之後她爸爸又搬了另外一句.
我們要小心我們眼裡的糧和刺.
她心裡想我也真的明白.
不過別人的那根刺在我眼裡真的很刺.
五塊錢不行.
結果經過一輪很激烈的討論和道歉之後.
他們又拿著對方一些前言不對後語的供詞.
死都要別人服輸.
於是又再鬧起來.
這個時候爸爸突然之間拍手.
聲音顫抖地說「我代魔鬼拍」.
我朋友當然不理會.
心想可不可以少搬一天天賦和魔鬼的議事論事呢.
忽然之間他聽到電腦另一邊傳來一些抽搐聲.
其實他爸爸開始哭.
越哭越厲害.
我朋友告訴我那次是他爸爸第一次在他面前哭.
於是他也開始哭.
他姐姐也開始哭.
連媽媽也哭了一份.
因為他們都聽到魔鬼笑得很開心.
家裡有很多個一家之主.
所以我朋友願意跟天父說.
天父我做錯了.
對不起.
基督是我家之主.
不是用來掛的.
是要每一個人願意在家裡讓步.
讓天父進來家裡做主.
這七個字才有意思.
保羅他整個的教導.

$^{521}$是要指出家中只有一個主人.
就是基督.
因此無論是家庭裡邊一對的關係.
關係裡邊一個方.
如果這個成員只想著怎樣堅持自己的權利.
這個家很快會凹凹亂亂散.
如果每個人在家裡都可以宣認.
唯有基督是我家之主的時候.
這個家的主權就是在主耶穌基督手上.
大家只要照顧自己或對其他成員的責任.
其實每一個成員的權利都會自動被保障.
我們繼續看第二個部分.
當我們領受了保羅對家庭中怎樣活出新人樣式的提醒之後.
他就勸信徒要專心警醒禱告.
並且要打開宣教的門去祈禱.
我剛才也說了今天不會特別去講二字四字.
因為我很想繼續講家庭的主線.
我們直接去看五至六節.
程序表我都印了在後面.
讓我用一個message的譯本去讀出五至六節給大家聽.
其實在愛人當中生活和工作的時候.
這個譯本很棒.
要用你們的腦袋.
不要錯過任何機會.
要善用每一個機會.
說話要有恩典.
目標就是在對話中將別人最好的一面帶出來.
不是貶低他們.
也不是排斥他們.
第五節,在這裡的愛人其實是指非信徒.
在座各位有沒有一些弟兄姊妹家裡是有非信徒的?.
有沒有的?.
是不是有的?.
幸好我不是唯一一個.
保羅提醒我們.
其實我們要用智慧去跟這些親友交往.
所謂的交往.
其實原文叫做walk tour.
是指我們要主動走向一些微信者.
保羅在這裡給了兩個錦囊.

$^{561}$第一,我們要好好把握時機.
第二,我們要用智慧.
說話的時候要帶著恩典.
要將別人最好的一面帶出來.
不是貶低或排斥別人.
非常直接,清楚的教導.
讓我分享最後一個姐妹的故事.
她也是傳道人.
目前是家裡唯一一個基督徒.
她的父母其實是比較寡言的.
所以從小到大就不懂得關心她或跟她聊天.
所以儘管她的父母不擅長用言語跟她溝通.
但也會用煮飯幫她收拾東西去表示她的愛意.
但我的朋友不太喜歡她的家人在未得到她同意下.
以幫她收拾為理由去翻閱她的東西.
她覺得很不被尊重.
讀大學期間,她的家人問她成績如何.
她如實告知全A.
但其實她的家人未必知道大學的A有多難拿.
以為這是一個很正常的發揮.
所以她家人的回應就是跟她說.
大學生只懂讀書,連簡單的生活技能也不懂.
其實我的朋友面對這些評語.
她也覺得很委屈和很難受.
這些一點一滴都讓她很難跟父母溝通.
在家裡也沒有什麼空間.
於是她就跟我分享她想搬出來的想法.
沒多久,她的教會裡有位阿姨很有心很想支持傳道人.
以一個低於市價的租金租了現在的單位給我和我的同事.
所以我們也勒紙好希望這間恩典屋可以接待不同的弟兄姊妹.
對她來說最感恩的事情就是.
我們現在住的地方和她祖母所住的老人院.
只是步行距離有幾分鐘.
所以她可以經常去探望祖母.
而且當我們有了自己的地方.
多了一些私人的空間.
有一個舒適合適的人際距離的時候.
我朋友和她父母的關係慢慢變好.
她也很樂意邀請父母來我們家吃飯.
最近她告訴我.

$^{601}$年廿九那晚,即是什麼時候?.
明晚,沒錯.
她已經幫祖母向老人院申請請假.
她父母會帶祖母上來.
她親手煮飯,還叫其他親戚和祖母一起吃團年飯.
原來自從祖母進了老人院之後.
她一家人已經好幾年沒有這樣坐在一起吃飯.
她在家族裡,背份最低,又是唯一一個基督徒.
不過她善用上帝給我們的恩典.
她在找時機,她藉著款待.
將家人重新連結在一起.
我相信她未信主的家人.
不僅見到她的主動和孝順.
更加將會看見福音的大能.
因為其實先是主耶穌.
祂融化了人和人之間的冷漠.
兒女和父母之間的誤會.
因著上帝的恩典,我的朋友有勇氣和她父母復和.
不僅如此,她更有能力為身邊的親友多走一步.
帶著恩典前行.
她的見證也很激勵我.
因為其實我也不是做得特別好.
我剛才和大家分享的第一個傳道人的見證也是一樣.
這篇道理也是我自己說給自己聽.
我也受著這三個傳道人的見證去激勵.
有時候其實我的生命也是這樣一塌糊塗.
我和我家人的相處仍然會撞板.
但是我知道,主耶穌基督.
即使在我家人還未信主的時候.
我也可以邀請祂成為我的一家之主.
我很希望旺朕的大家.
只要心裡尊主耶穌基督為聖.
常常做好準備.
無論誰要求我們交代我們心裡有盼望的原因的時候.
我們就抓緊時機.
帶著恩典,溫柔而堅定.
去和他們一起分享主耶穌基督的福音.
\newpage



\section{約翰福音 12:1-8}
\label{sec:NhZRAevmHb4}
\textbf{2026年2月22日崇拜直播｜梁家麟牧師｜腳的事奉｜約翰福音12\_1-8}
\newline
\newline
連結: \href{https://youtube.com/watch?v=NhZRAevmHb4}{\texttt{https://youtube.com/watch?v=NhZRAevmHb4}} ~~~~ 語音日期: 2026-02-22
\newline
\newline
\hyperref[sec:jemYOQSwajw]{\small{< < < PREV SERMON < < <}}
~
\hyperref[sec:index_chronic]{\small{[返順時目]}}
~
\hyperref[sec:8qvQyuuLcyA]{\small{> > > NEXT SERMON > > >}}
\newline
\newline
約翰福音 12:1-8
\newline
\begin{longtable}{cl}
\hline
\hline
章節 & 經文 (和合本修訂版)\\
\hline
12:1 & \begin{tabularx}{0.7\textwidth}{X} 逾越節前六天，耶穌來到伯大尼，就是他使拉撒路從死人中復活的地方。 \end{tabularx} \\ \\ \relax
12:2 & \begin{tabularx}{0.7\textwidth}{X} 有人在那裡為耶穌預備宴席；馬大伺候，拉撒路也在同耶穌坐席的人中間。 \end{tabularx} \\ \\ \relax
12:3 & \begin{tabularx}{0.7\textwidth}{X} 馬利亞拿著一斤極貴的純哪噠香膏，抹耶穌的腳，又用自己頭髮去擦，屋裡充滿了膏的香氣。 \end{tabularx} \\ \\ \relax
12:4 & \begin{tabularx}{0.7\textwidth}{X} 有一個門徒，就是那將要出賣耶穌的加略人猶大，說： \end{tabularx} \\ \\ \relax
12:5 & \begin{tabularx}{0.7\textwidth}{X} 「為甚麼不把這香膏賣三百個銀幣去賙濟窮人呢？」 \end{tabularx} \\ \\ \relax
12:6 & \begin{tabularx}{0.7\textwidth}{X} 他說這話，並不是關心窮人，而是因為他是個賊，又管錢囊，常偷取錢囊中所存的。 \end{tabularx} \\ \\ \relax
12:7 & \begin{tabularx}{0.7\textwidth}{X} 耶穌說：「由她吧！她這香膏本是為我的安葬之日留著的。 \end{tabularx} \\ \\ \relax
12:8 & \begin{tabularx}{0.7\textwidth}{X} 因為常有窮人和你們在一起，但是你們不常有我。」 \end{tabularx} \\ \\
[1ex]
\hline
\hline
\end{longtable}
$^{1}$今日聆聽的聖道是記載在約翰福音十二章第一至第八節.
在新約聖經一百四十五頁.
逾越節前六日,耶穌來到伯大利.
就是他叫拉撒路從死裡復活之處.
有人在那裡給耶穌預備筵席.
馬大侍候,拉撒路也在那和耶穌坐直的人中.
瑪利亞就拿著一斤極貴的珍拿達香膏.
撥耶穌的腳,又用自己頭髮去擦.
屋裡就滿了高的香氣.
有一個門徒,就是那將要罵耶穌的加略人猶大.
說:這香膏為什麼不賣三十兩銀子周濟窮人呢?.
他說者說:並不是掛念窮人,乃因他是個賊.
又帶著錢囊,尚取其中所存的.
耶穌說:猶他吧,他是為我安葬之日存留的.
因為常有窮人和你們同在.
只是你們不常有我.
各位聽眾朋友,新年快樂.
我們來到耶穌基督在人間的生涯最後.
和最重要的一個星期.
耶穌再次來到伯大利.
這是祂前往耶路撒冷的必經之地.
此行祂即將面對的是十字架的犧牲.
這是《約翰福音》第十二章.
你會發現,使徒約翰用了整本福音書.
差不多一半的篇幅.
來記載耶穌人間生涯最後一個星期.
一個星期用一半的篇幅,十幾章來記載.
比例是很大的.
和其他福音書比較,也是最高的比例.
顯示對約翰來說.
耶穌預備上十字架的事件.
比任何其他事蹟都更重要.
耶穌在這個星期的言行.
對門徒的教導.
在他心中也留下了最深刻的印象.
耶穌走上十字架道路的第一個事件.
是瑪利亞告末祂.
接著才是終之星期日,就是Palm Sunday.
他榮耀進入耶路撒冷.
聖經一開始說,預節前六日.

$^{41}$也就是耶穌死前一個星期.
在伯大利,耶穌和門徒一起參與一場的筵席.
約翰沒有說這場筵席是誰主持.
不過記載相類似故事的馬可和馬太.
說明這個筵席是設在長大麻風西門的家.
奇怪的是,馬太和馬可記載了長大麻風西門.
即是說那個筵席的主人的名字.
但是沒有提到馬大瑪利亞拉撒路一家的名字.
連要告末耶穌的瑪利亞的名字都沒有提到.
只是說有一個女人做了這樣的事.
所以各自留記了一些我覺得都很奇怪的事.
為什麼會不記馬利亞的名字呢?.
這場筵席或許是為了慶祝拉撒路從死裡復生而設的.
所以約翰註明拉撒路也在席間.
當時的社會其實並不平等,男女不平等.
女性通常是不能參與筵席的.
所以馬大就以幫忙接待客人的身份出現在現場.
瑪利亞就完全沒有角色.
在這個時候耶路撒冷的工會已經下達要逮捕耶穌的命令.
所以耶穌的朋友,就是長大麻風的西門.
願意公開接待他,不可以說是很有勇氣的.
即使參加筵席的其他賓客,也很可能會受株連.
在故事的結尾,約翰記述第十節.
但祭司長商議連拉撒路也要殺了.
所以凡是跟耶穌有聯繫的人.
都要面對猶太領袖的迫害.
風聲鶴唳,山雨欲來風滿流.
在筵席間,馬利亞就出場了.
第三節說,馬利亞拿著一斤極貴的珍拿達香膏.
抹耶穌的腳,又用自己的頭髮去擦.
屋內就滿了膏的香氣.
這是一個極其不尋常的行動.
馬利亞為什麼要這樣做呢?.
你說她愛耶穌,你只說她愛耶穌,其實不足夠解釋.
你們都愛我是嗎?.
你們有沒有將所有的香水倒在我的腳上?.
不要說一斤或一瓶,一滴都沒有.
你們苦心自問,你們是不是很不愛心?.
不要說香水,倒一瓶生抽也好.
十年八輩都捨不得.

$^{81}$所以愛跟她做這件事沒有必然的關係.
多數識經者都主張.
馬利亞是相信耶穌即將被殺死亡.
然後她決定這樣做.
預先膏耶穌尚未死亡的身體.
這也是耶穌說的第七節.
七世為我安葬之日存留的.
一般來說,香膏是膏在屍體的身上.
但是馬利亞想,這個小女子應該沒有機會參加耶穌的葬禮.
不如就預先做這件事.
圓一個心願,表達對耶穌的敬意.
身後事變成身前事.
做法是極其不道地的.
不過可以勉強理解.
耶穌在這件事之前的好幾個場合.
都預告了祂即將面對被釘十字架的命運.
但是在約翰福音就有三章十四節,八章二十八節,十章十一節,十五節.
其他福音書記載的就更多了.
例如在出發前往耶路撒冷之前.
就是沒多久前發生的事.
耶穌就對門徒做了一個預告.
在馬福音十章三十三到三十四節說.
看吧,我不想耶路撒冷去.
人子將要被交給祭司長和文士.
他們要定他的死罪.
交給外邦人他們要戲弄他.
吐唾沫在他臉上鞭打他殺害他.
過了三天他要復活.
這個預告不可以說不詳盡.
連他要受怎樣的苦都有清楚說明.
不過聖經沒有記載門徒對耶穌的預告有什麼反應.
連一個回應都沒有.
反倒馬上說門徒爭取在天上的位置和榮耀.
有這樣的爭論.
耶穌上十字架的預告是主要跟門徒說的.
但他們聽了之後都沒什麼反應.
大概他們不當是一回事.
耶穌常常說一些門徒不太明白的事.
常常做一些很驚嚇性的言論.
門徒聽多了,嚇大了就不以為意.

$^{121}$反倒我們的夫子喜歡語出驚人,語無倫次.
聽一下不要當真就沒事了.
這個大概都是我們很多人.
在讀聖經或者聽講道的時候的態度.
聖經總是喜歡唱高調的.
說一些不吃人間煙火的要求.
高舉一個理想主義的價值觀.
例如什麼愛你的仇敵.
為耶穌捨棄一切.
不要愛世界等等.
這些不合心意的說話.
聽習慣了就沒什麼特別的反應.
反正聖經就是喜歡講大的.
講到很大的,什麼都講大碎.
牧者喜歡慣常說一些屬靈口號.
只要你不當真就沒事了.
瑪利亞卻跟多數的門徒不一樣.
他或者在沒多久之前.
聽到耶路撒冷的猶太公會.
要謀害耶穌.
他也都聽到耶穌.
曾經自己預告.
他要上十字架這樣的說話.
單純的他聽了耶穌講就相信了.
他沒能力力挽狂瀾去拯救耶穌.
保護耶穌.
只能夠盡他能力所及的.
表達對耶穌的一些心意.
至於問為什麼要高耶穌的腳呢.
猶太人如果高一個生人.
是高他的頭的.
流到鬍鬚又流到衣襟.
是倒在頭上的.
如果高屍體死了的.
就應該是高全身的.
而不會只是高腳的.
對於這個問題.
不同的學者有不同的講法.
不過最多數的講法都是說.
瑪利亞謙卑.

$^{161}$所以不高頭就高腳.
不過我想來想去.
我都不覺得為什麼會謙卑.
我覺得最簡單的答案.
就是這個是瑪利亞竭盡所能.
所能夠企及的高度.
瑪利亞是一個很矮小的女人.
矮小不是指她的身材.
而是指她的社會地位和知識水平.
她怎樣都不能夠高攀.
和耶穌平頭並肩.
怎樣抬高她的腳尖.
都無法摸到耶穌的頭.
她只能夠接受.
她能夠靠近的就是耶穌腳的位置.
事實上是真的.
他們吃筵席.
大家都是這樣靠著靠著坐著.
腳伸出來.
瑪利亞又沒資格坐過去.
只能夠摸到耶穌伸出來的腳.
再逼近一點.
在聖經我們見到.
瑪利亞總是在耶穌的腳前.
耶穌和門徒在他家作客.
耶穌坐著和門徒談道.
瑪利亞都沒資格.
也不能夠坐在門徒那裡.
她就只能夠坐在耶穌的腳前.
路加福音是這樣記載的.
路加福音十章三十九節.
馬大有一個妹子名叫瑪利亞.
在耶穌腳前坐著聽她的道.
瑪利亞所能夠企及的.
是耶穌的腳的高度.
所能夠觸摸的是耶穌的腳.
她就接受這個現實的限制.
在她觸手可及的範圍裡去事奉耶穌.
高不到頭就隨便高一隻腳.
弟兄姊妹.

$^{201}$我是很想和你們分享聖經.
最深層的真理.
不過我資質有限.
我只能夠摸到耶穌的腳眼.
我認命.
我就盡力在腳眼的高度去理解耶穌.
然後和你們分享.
在腳眼高度的耶穌的故事.
我已經盡力了.
我高度不夠.
請你們原諒.
你們可能比我好一點.
有人摸到耶穌的膝頭.
有人高攀到耶穌的腰.
不過我不相信你們有人可以高度摸到耶穌的頭.
無論如何.
你們在你們力所能及的部位.
就將香膏倒在耶穌的身上.
這是我對這段經文.
腳眼高度的默想.
並不是一個嚴謹的聖經.
純粹參考一下.
瑪利亞這個豪奢的奉獻.
惹來眾人的憎惡.
投尼的眼神多數不是欣賞.
而是不理解,不認同甚至反感.
連耶穌的門徒加略人猶大都忍不住批評說.
這香膏為什麼不賣三十兩銀子周濟窮人呢?.
在出賣耶穌以前.
加略人猶大在福音書是第一次有對白.
之前他沒有對白.
從來都沒有對白.
必須要說.
加略人猶大這個反應是正常,合理的.
當你看到一些超乎尋常的豪奢行為.
你心裡總會有一些厭惡的反感.
有錢人用銀子去點煙.
你會不會覺得很側目?.
有錢人不是這樣用的.
這樣做是想顯示什麼?.

$^{241}$對不對?.
吃自助餐.
你看到有人一整盤壽司拿回來.
只吃了上面那塊壽司.
全部飯都丟掉.
你都會有些覺得.
搞什麼?這樣都行?.
這麼折墮?.
你都會覺得不應該這樣浪費東西.
所以一瓶貴重的香膏用來塗腳.
就好像今天有人用眼霜去塗腳一樣.
你都會覺得很荒誕.
這麼貴,要不要這樣塗?.
不過聖經作者卻是說.
加略人猶大是一個有私心的人.
他的批評不是出於正義.
是出於私心.
第六節.
他說者說,並不是掛念窮人.
乃因他是國賊,又帶著錢囊.
尚取其中所存.
約翰這個寶注.
顯示他對猶大,加略人猶大.
是非常反感的.
他直指猶大不僅僅是.
謀害耶穌的時候才做壞事.
卻一直是壞蛋.
一直做壞事.
他貪財.
出賣耶穌正正是由於他貪財.
約翰這個解釋.
或許是為日後猶大背叛耶穌.
去做一個解說的鋪路.
因為其他福音書.
在記載猶大賣主的時候.
只是提到撒旦.
就突然進入他的心中.
沒有提到猶大本身的性格.
和他跟耶穌的關係.
一個跟隨耶穌多年的門徒.

$^{281}$竟然決定賣主.
總是需要一個原因的.
說他一直都很貪財.
就比較合理.
必須要指出.
約翰對猶大.
對加略人猶大的評論.
其實不是很恰當的.
第六節的寶注也不算很相干.
加略人猶大.
或者真是一個貪財的人.
但這跟他批評瑪利亞浪費金錢.
沒有什麼關係.
他只是建議.
將這個香膏.
用來賣了周濟窮人.
他沒有叫瑪利亞變賣香膏.
然後把錢送給他.
他不會因這個批評得到什麼好處.
我們不可以因為一個人施得不好.
就因人廢言否定他所說的一切.
他說的話有道理還是沒有道理.
不是由他個人的品德去決定.
不過.
加略人猶大對瑪利亞評論.
可能有道理.
但仍然掩蓋不了.
這番說話出於他的口.
所產生的荒誕感.
聖經如果說他一個貪財的人.
我們可以推算.
他不會是一個關心窮人的人.
一個自己不關心窮人的人.
卻批評其他人不關心窮人.
這就很荒謬.
我們在香港.
我們遇上過很多.
說這樣荒謬說話的人.
他們喜歡站在道德高地.
堅持最高的道德標準.

$^{321}$然後批評其他人不道德.
不過他們只是扮演批評者反對黨的角色.
沒有打算身體力行.
將對別人的要求.
兌現在自己身上.
在經濟不景氣的時候.
有人要求教會.
賣掉教產去分給窮人.
不過做這樣要求的人.
卻不會掏空自己的錢包.
或者賣掉自己的手機去周濟窮人.
我媽媽對這些人的判定是.
只會罵人不會罵自己.
文巢巢一些的說法就是.
嚴於律人就寬以律己.
耶穌的判語就是.
能說不能行.
當然.
加里蘭猶大說這番說話是荒謬的.
不等於他對馬尼亞的批評是荒謬的.
前面我們已經說過不可以因人廢言.
耶穌也沒有因法利賽人的假冒為善.
門徒不要聽法律上的說話.
所以我們怎麼看.
加里蘭猶大對馬尼亞的指責呢.
為什麼不做上帝金錢的好管家.
將錢財用在能夠產生最大效用.
幫助最多人的事上.
反而高耶穌的腳這麼浪費呢.
不用我們去將來解釋.
耶穌已經給了答案在第七第八節.
他親自為馬尼亞做一個辯解.
耶穌的辯解有兩點.
第一.
馬尼亞這樣做是他預知耶穌即將死亡.
所以提前為他做喪事.
用香膏膏他的身體.
即使你說安葬禮都不用用這麼昂貴的珍拿大香膏.
你可以這樣說.
但是馬尼亞行動就不再是完全沒有意義不合理.

$^{361}$畢竟喪禮是人生大事.
有人要求封安大葬.
用幾十萬的名貴棺木.
其他人都不可以說太多.
我是喜歡用很貴很貴的東西去籌備喪禮.
是可以的.
這是第一.
第二.
耶穌的提法是.
不要將服侍耶穌和周濟窮人做一個簡單的二婚對立.
服侍了耶穌就幫助不到窮人.
於是我們只能夠在這兩者二擇其一.
周濟窮人是什麼時候都可以做的事.
窮人很多到處都是.
今天不做明天做.
我們沒辦法想像每一天全世界的窮人都消失了.
你想周濟窮人都不知道在什麼地方去周濟.
沒有的.
不可能的.
所以在適當的時候做適當的事.
在迫切的時候做迫切的事.
不用做那些無謂的對比.
錢花在這裡還可以賺回來.
再用在別處.
是不是.
宏福縣火災.
你馬上想到可以捐一些錢去幫助他們.
身邊有人就跟你說.
非洲的災民更慘.
烏克蘭的難民更慘.
所以不應該幫助香港的災民.
這樣的言論其實沒有什麼營養價值.
你覺得烏克蘭難民更值得幫助.
那你就全力幫他.
何必批評那些想幫香港災民的人.
各人有不同的負擔和領袖.
各人只需要順著自己的領袖.
各做各.
是不是很好.
我們很難完全比較不同事物的價值.

$^{401}$聖誕節和家人去酒店吃自助餐.
一邊吃一邊想到世界上仍有十億吃不飽的人.
然後質問自己.
為什麼不吃麥當勞.
省點錢去救濟基民.
於是一邊吃一邊流淚.
一邊吃一邊內疚.
你告訴我你還怎麼活下去.
對不對.
情人節買一束玫瑰花送給女朋友.
人家就告訴你.
一束玫瑰花的價錢可以買十斤臘腸.
是不是更實惠一些.
所以乾脆送臘味給女朋友算了.
不過如果你這樣做.
我就祝你好運.
買十斤臘腸.
尚有窮人和你們同在.
只是你不常有我.
耶穌的說話兼顧了常規責任和特殊考慮.
兼顧了人情和道理.
耶穌認同照顧窮人的重要性.
又肯定了瑪利亞為他所做的不尋常的舉動.
我們回去說瑪利亞.
雖然她是被加勒人猶大批評的人.
不過她自己其實是沒什麼能力去參與.
耶穌和加勒人猶大的辯論.
不要說她自己從來沒想過.
那個私藏的真拿大香膏要怎樣用.
才達到最大的效益.
最討上帝的喜悅.
她從來沒有將香膏抹在耶穌的腳.
和將它變賣周濟窮人做一個比較.
判斷哪樣東西才是異所當為.
她甚至沒想過一定要做一些.
最有用,最有效益,最有影響力的事.
一個小女子從來沒想過要拯救或改變世界.
她的心掛著耶穌.
為祂宣告即將死亡.
而覺得自己很不安樂.

$^{441}$很想做一些事在祂身上.
想起家裡藏有一瓶極珍貴的香膏.
就決心用它在耶穌的身上.
她用香膏高耶穌的腳.
或者跟她記住耶穌即將死亡有關.
但卻不一定連繫到耶穌的安葬禮儀.
「祂是為我安葬之日存留的」.
這是耶穌的說法,耶穌的解讀.
但不一定是瑪利亞的想法.
畢竟將喪葬禮儀.
做在一個還沒死的生人身上.
總是不可思議的事.
你家裡有個病人還沒斷氣.
你會不會預先跟他做死人化妝.
穿好壽衣?不會的.
大吉利是嗎?.
我個人相信.
小女子瑪利亞只是隨感覺行事.
沒有太多複雜周詳的想法.
所做的事不一定合理.
不一定有道理.
要她自己去解釋.
她也不見得能說出什麼道理.
但就是這樣跟隨感覺走.
她很愛耶穌.
很想做些什麼在祂身上.
很想表達對耶穌完全的付出和捨棄.
就這樣豪奢一次.
傻一次.
不行嗎?.
為什麼一定要有道理?.
一定要有神學理據?.
一定要能夠說服其他人?.
不一定需要的.
我自己也試過不止一次.
在香港或者有時候去外地出差的時候.
想買一件禮物給太太.
就突然間豪情.
就找了能夠拿出來的錢.
想都不想就買了一件.

$^{481}$糟糕的是.
我買了回來送給太太之後.
才發覺我買的不是她喜歡的東西.
例如有一次我買了一串黑色的珍珠鏈給她.
我都沒見過她怎麼戴過.
她說不太好襯衣服.
不過人就是這樣.
不一定考慮周全.
不一定是精明的選擇.
不一定合情合理.
就是真心,真意,真情流露.
就是這樣.
你也知道有個弟兄跟我說.
他很想在上帝面前.
跳一支舞去敬拜祂.
雖然我想不到.
跳舞跟報到直堂.
宣教大使命有什麼直接的關係.
不過他想跳就跳吧.
對不對.
有個姐妹每個禮拜.
認真插盤花.
呈獻給上帝.
雖然我不確定.
上帝是不是真的會欣賞.
那些什麼小園樓,草原樓.
那些什麼花洞.
不過她想插就插吧.
不用說那麼多什麼跳舞神學.
花道神學.
對不對.
不用那麼吵.
就是喜歡做.
覺得我很想做些美好的事給上帝.
想做些美好的事給祂.
那就做吧.
瑪利亞是一個做事不夠理性.
考慮不夠周全的人.
要高耶穌.
高生人是高過頭的.

$^{521}$高死人是高全身的.
怎麼能夠只高一雙腳呢.
總不能夠好像我前面說.
因為他唯一能夠摸到的是耶穌的腳.
那就事事淡淡.
張張就求其高腳.
你要瑪利亞做神學解釋.
為什麼高一雙腳.
我猜她也啞口無言.
她沒辦法高其他部位.
事事淡淡.
但耶穌基督就是護著瑪利亞.
和她辯解.
擋住其他人對她的批評.
而這個顯示.
耶穌閱立瑪利亞所做的一切.
就這一點.
瑪利亞就贏了所有人.
包括批評者.
加略人猶大 為她擺設筵席的長大麻風的西門.
甚至奔走服侍的馬大等等.
聖經沒有記載耶穌對他們的讚賞.
獨獨就讚賞了瑪利亞.
馬可福音的記載就更加誇張.
耶穌說了什麼話.
我實在告訴你們.
普天之下.
無論在什麼地方傳這福音.
有述說者女人所作的以為紀念.
讚得他很盡.
我們很愛耶穌.
一心想做些事在耶穌身上.
一心想討他喜悅.
我們做的事是對是錯.
唯一判斷的標準.
就是耶穌受不受.
耶穌閱立不閱立.
耶穌受 耶穌閱立.
那就贏了.
加略人猶大 受不住.

$^{561}$身邊的人看不過眼.
事事私意又有什麼所謂呢.
他們是閒雜人等.
是無關痛癢的路人甲.
我和太太說情話.
你在身邊無意聽到.
覺得很噁心 很肉麻.
我就會說關你屁事.
我太太受不受是最重要的.
其他人有什麼所謂.
弟兄姊妹.
很多時候我們的事奉.
是做在人前.
多過做在耶穌面前.
我們在崇拜的公土.
是對人說而不是對上帝說.
我們做的鋪張華麗的服事.
是想在眾人面前.
顯示我們的專業.
我們的專業.
我們的優秀.
我們做很多事情.
總是瞻前顧後.
考慮最多的是其他人的看法和評價.
我們介意人的想法.
過於專注地去討上帝的喜悅.
保羅在這份所書六章六到七節說.
必要旨在眼前.
事奉像是討人喜歡的.
要像基督的僕人.
從心裡遵行上帝的旨意.
甘心事奉好像服事主不像服事人.
這是一個非常非常重要的教導.
你一生跟隨耶穌.
一生服事祂.
不停問自己.
耶穌喜不喜歡我做的事.
我所做的事.
祂悅立不悅立.
我是否專注用心地去哄祂開心.

$^{601}$討祂喜悅.
這永遠是我不停調教自己.
問的問題是什麼.
當然我不會完全不理黃震弟兄姊妹妹的感受.
但我真的仍然.
我第一個要問的問題.
不停問的問題是.
上帝喜不喜歡我今天的服事.
你悅不悅立.
你收不收貨.
這永遠是對我最重要的提醒.
專注基督.
專心哄祂開心.
討祂喜悅.
無疑我們事奉主也要運用理性.
要用的.
盡心盡性盡義盡力.
不過有時候傻傻地.
做一些出格不尋常的事也沒有所謂.
畢竟你們記得我曾經說過不止一次.
做一個愛者是不可以太精明的.
我經常說.
做丈夫不要太精明.
做太太不要太精明.
做父母不要太精明.
做耶穌的門徒也不需要太精明.
力所能及.
摸不到腳眼.
摸腳板底也無所謂.
我們努力.
做我們所能夠攀及的.
腳眼的事奉.
腳的服侍.
\newpage



\section{詩篇 22:1-31}
\label{sec:8qvQyuuLcyA}
\textbf{2026年3月1日主餐崇拜直播｜陳明泉牧師｜基督之詩篇:我的神!為什麼離棄我?｜詩篇22\_1-31}
\newline
\newline
連結: \href{https://youtube.com/watch?v=8qvQyuuLcyA}{\texttt{https://youtube.com/watch?v=8qvQyuuLcyA}} ~~~~ 語音日期: 2026-03-02
\newline
\newline
\hyperref[sec:NhZRAevmHb4]{\small{< < < PREV SERMON < < <}}
~
\hyperref[sec:index_chronic]{\small{[返順時目]}}
~
\hyperref[sec:yBE6W2fCqSM]{\small{> > > NEXT SERMON > > >}}
\newline
\newline
詩篇 22:1-31
\newline
\begin{longtable}{cl}
\hline
\hline
章節 & 經文 (和合本修訂版)\\
\hline
22:1 & \begin{tabularx}{0.7\textwidth}{X} 我的神，我的神，為甚麼離棄我？為甚麼遠離不救我，不聽我的呻吟？ \end{tabularx} \\ \\ \relax
22:2 & \begin{tabularx}{0.7\textwidth}{X} 我的神啊，我白日呼求，你不應允；夜間呼求，也不得安寧。 \end{tabularx} \\ \\ \relax
22:3 & \begin{tabularx}{0.7\textwidth}{X} 但你是神聖的，用以色列的讚美為寶座。 \end{tabularx} \\ \\ \relax
22:4 & \begin{tabularx}{0.7\textwidth}{X} 我們的祖宗倚靠你；他們倚靠你，你解救他們。 \end{tabularx} \\ \\ \relax
22:5 & \begin{tabularx}{0.7\textwidth}{X} 他們哀求你，就蒙解救；他們倚靠你，就不羞愧。 \end{tabularx} \\ \\ \relax
22:6 & \begin{tabularx}{0.7\textwidth}{X} 但我是蟲，不是人，被眾人羞辱，被百姓藐視。 \end{tabularx} \\ \\ \relax
22:7 & \begin{tabularx}{0.7\textwidth}{X} 凡看見我的都嗤笑我；他們撇嘴搖頭： \end{tabularx} \\ \\ \relax
22:8 & \begin{tabularx}{0.7\textwidth}{X} 「他把自己交託給耶和華，讓耶和華救他吧！耶和華既喜愛他，可以搭救他吧！」 \end{tabularx} \\ \\ \relax
22:9 & \begin{tabularx}{0.7\textwidth}{X} 但你是叫我出母腹的，我在母懷裡，你就使我有倚靠的心。 \end{tabularx} \\ \\ \relax
22:10 & \begin{tabularx}{0.7\textwidth}{X} 我自出母胎就交在你手裡，自我出母腹，你就是我的神。 \end{tabularx} \\ \\ \relax
22:11 & \begin{tabularx}{0.7\textwidth}{X} 求你不要遠離我！因為災難臨頭，無人幫助。 \end{tabularx} \\ \\ \relax
22:12 & \begin{tabularx}{0.7\textwidth}{X} 許多公牛環繞我，巴珊大力的公牛四面圍困我。 \end{tabularx} \\ \\ \relax
22:13 & \begin{tabularx}{0.7\textwidth}{X} 牠們向我張口，好像獵食吼叫的獅子。 \end{tabularx} \\ \\ \relax
22:14 & \begin{tabularx}{0.7\textwidth}{X} 我如水被倒出，我的骨頭都脫了節，我的心如蠟，在我裡面熔化。 \end{tabularx} \\ \\ \relax
22:15 & \begin{tabularx}{0.7\textwidth}{X} 我的精力枯乾，如同瓦片，我的舌頭緊貼上顎。你將我安置在死灰中。 \end{tabularx} \\ \\ \relax
22:16 & \begin{tabularx}{0.7\textwidth}{X} 犬類圍著我，惡黨環繞我；他們扎了我的手、我的腳。 \end{tabularx} \\ \\ \relax
22:17 & \begin{tabularx}{0.7\textwidth}{X} 我數遍我的骨頭；他們瞪著眼看我。 \end{tabularx} \\ \\ \relax
22:18 & \begin{tabularx}{0.7\textwidth}{X} 他們分我的外衣，為我的內衣抽籤。 \end{tabularx} \\ \\ \relax
22:19 & \begin{tabularx}{0.7\textwidth}{X} 耶和華啊，求你不要遠離我！我的救主啊，求你快來幫助我！ \end{tabularx} \\ \\ \relax
22:20 & \begin{tabularx}{0.7\textwidth}{X} 求你救我的性命脫離刀劍，使我僅有的脫離犬類， \end{tabularx} \\ \\ \relax
22:21 & \begin{tabularx}{0.7\textwidth}{X} 求你救我脫離獅子的口；你已經應允我，使我脫離野牛的角。 \end{tabularx} \\ \\ \relax
22:22 & \begin{tabularx}{0.7\textwidth}{X} 我要將你的名傳給我的弟兄，在會眾中我要讚美你。 \end{tabularx} \\ \\ \relax
22:23 & \begin{tabularx}{0.7\textwidth}{X} 敬畏耶和華的人哪，要讚美他！雅各的後裔啊，要榮耀他！以色列的後裔啊，要懼怕他！ \end{tabularx} \\ \\ \relax
22:24 & \begin{tabularx}{0.7\textwidth}{X} 因為他沒有藐視、憎惡受苦的人，也沒有轉臉不顧他們；那受苦之人呼求的時候，他就垂聽。 \end{tabularx} \\ \\ \relax
22:25 & \begin{tabularx}{0.7\textwidth}{X} 我在大會中讚美你的話是從你而來，我要在敬畏耶和華的人面前還我的願。 \end{tabularx} \\ \\ \relax
22:26 & \begin{tabularx}{0.7\textwidth}{X} 願困苦的人吃得飽足，願尋求耶和華的人讚美他。願你們的心永遠活著！ \end{tabularx} \\ \\ \relax
22:27 & \begin{tabularx}{0.7\textwidth}{X} 地的四極都要想念耶和華，並且歸順他，列國的萬族都要在你面前敬拜。 \end{tabularx} \\ \\ \relax
22:28 & \begin{tabularx}{0.7\textwidth}{X} 因為國度屬於耶和華，他是管理列國的。 \end{tabularx} \\ \\ \relax
22:29 & \begin{tabularx}{0.7\textwidth}{X} 地上富足的人都必吃喝而敬拜，凡下到塵土中不能存活自己性命的人，都要在他面前下拜； \end{tabularx} \\ \\ \relax
22:30 & \begin{tabularx}{0.7\textwidth}{X} 必有後裔事奉他，主所做的事必傳給後代。 \end{tabularx} \\ \\ \relax
22:31 & \begin{tabularx}{0.7\textwidth}{X} 他們必來傳他的公義給尚未出生的子民，這是他的作為。 \end{tabularx} \\ \\
[1ex]
\hline
\hline
\end{longtable}
$^{1}$以下是經訓的時間.
今天我們選讀的經文是舊約聖經的詩篇第22篇.
一首大衛的詩.
詩篇第22篇 請聽我讀出.
我的神 我的神 為什麼離棄我.
為什麼遠離不救我 不聽我哀憨的言語.
我的神 我白日呼求 理不應運.
夜間呼求 並不駐聲.
但你是聖潔的 是用以色列的讚美為寶座的.
我們的祖宗倚靠你.
他們倚靠你 你便解救他們.
他們哀求你 便蒙解救.
他們倚靠你 就不羞愧.
但我是蟲 不是人 被眾人羞辱 被百姓藐視.
凡看見我的 都竊笑我 他們撇嘴搖頭說.
他把自己交託耶和華 耶和華可以救他吧.
耶和華既喜悅他 可以答救他吧.
但你是叫我出無福的.
我在無福的裡 你就使我有倚靠的心.
我自出母胎 就被交在你手裡.
從我母親生我 你就是我的神.
求你不要遠離我 因為急難臨近了.
沒有人幫助我 有許多公牛圍繞我.
巴山大力的牛工 四面困住我.
他們向我張口 好像找屍吼叫的獅子.
我如水被倒出來 我的骨頭都脫了節.
我心在我裡面如蠟融化.
我的精力枯乾如同瓦片 我的舌頭貼在我牙床上.
你將我安置在死地的塵土中.
犬類圍著我 惡黨環繞我.
他們紮了我的手 我的腳 我的骨頭 我都能數過.
他們瞪著眼看我 他們分我的外衣 為我的女衣沾溝.
耶和華 求你不要遠離我.
我的救主 求你快來幫助我.
求你救我的靈魂脫離刀劍.
救我的生命脫離犬類 救我脫離獅子的口.
你已經應運了我 使我脫離野牛的國.
我要將你的名傳於我的弟兄.
在會中我要讚美你.
你們敬畏耶和華的人要讚美他.

$^{41}$雅各的後裔都要榮耀他.
以色列的後裔都要懼怕他.
因為他沒有藐視珍惡受苦的人.
也沒有向他掩面.
那受苦之人呼籲的時候 他就垂聽.
我在大會中讚美你的話 是從你而來的.
我要在敬畏耶和華的人面前還我的怨.
謙卑的人必吃得飽足.
尋求耶和華的人必讚美他.
願你們的心永遠活著.
地的四極都要想念耶和華 並且歸順他.
列國的萬族都要在你面前敬拜.
因為國權是耶和華的 他是管理萬國的.
地上一切豐肥的人必吃渴而敬拜.
凡下到塵土中不能存活自己性命的人.
都要在他面前下拜.
他必有後裔事奉他.
主所行的事必傳於後代.
他們必來把他的公義傳給將要新的民.
賢明者是 是他所行的.
弟兄姊妹平安.
我們進入三月份.
我們開始藉著這幾個星期.
我們一齊默想基督耶穌.
接下來幾次的講道.
我們一個小的主題.
我們以基督的詩篇成為了接下來幾次講道的主要系列.
在這個系列裡我們會集中.
用舊約的詩篇.
從這個角度在詩篇裡面其實已經預表了.
甚至乎耶穌基督在十字架上.
在他的受輔當中引用的詩篇.
接下來這幾次我們會將那幾篇詩篇.
跟大家一起重讀一次.
在這個重讀過程裡面.
我們不單止明白了聖經的意思.
更加是我們更加體會到基督耶穌的心腸.
在他受輔當中.
他面對著十字架的酷刑的呼喊.
也在他的生命當中.

$^{81}$幫助我們重新去檢視自己的人生.
今天我們第一篇所看的就是詩篇第22篇.
傳統以來或者我們今天的世代.
我們讀很多的就是詩篇第23篇.
耶和華是我的牧者.
我們經常讀的.
不過我們很少留意.
原來這篇詩篇在千多年的教會文化裡面.
是一個主要重讀的詩篇.
在詩篇第22篇裡面.
甚至乎有些早期的教父說過.
他說這篇詩篇應該是每個基督徒必要讀的詩篇.
因為在這裡面.
他特別描述耶穌基督的受苦.
他預表了耶穌基督在十字架將要所承受的苦難.
亦都在這篇詩篇裡面體會到一個受傷的基督.
一個受苦的基督.
他對於他自己那種痛苦的負陷.
所以今天我們嘗試用這篇詩篇.
去了解一下大衛所寫的心境.
然後再用耶穌基督為何要引用這篇詩篇.
用這個角度來重新演繹.
不過在我們看這篇詩篇的時候.
剛才Pelton為我們讀的時候.
我們會發現一個很重要.
亦都幾難處理的問題.
就是這篇詩篇裡面大約是三十一節的經文.
你會發覺上半部.
用了大量的篇幅.
講到私人面對著的那種苦境.
是一個個人的哀歌.
他那裡面是講他自己很慘情.
他面對著一個很困難的境況.
下半部特別是由二十二節開始.
是一個頌讚的詩歌.
由哀歌到頌讚的詩歌.
你會發現中間有沒有一些轉捩點呢?.
在過程裡面.
為何私人會轉念呢?.
那個轉念的形式是怎樣呢?.

$^{121}$這是這篇詩篇裡難讀的地方.
今天我們向難度挑戰一下.
或者我們更加要從一個私人的心態裡面出發.
如果你有一個新的角度來詮釋.
可能這篇詩篇會為我們帶來.
我們今天很複雜的人生帶來一些出路.
我們先有個祈禱.
求上帝幫助我們同心祈禱.
求主你與我們同在.
讓我們能夠明白你對我們所講的話.
亦都讓你的話語成為我們心中裡面的聲音.
在人生每個階段的道路當中.
你的話語成為我們人生的出路和幫助.
恭敬將來的時間交託給上帝你.
願你親自保守眾弟兄姊妹.
讓他們不單止聽得明白.
願你的話語進到他們的心靈當中.
幫助孫講的講得清楚.
願你的話語忠心的.
亦都帶著上帝的亮光.
讓我們一齊來領受.
我們恭敬將來的時間交託給你.
願主你親自帶領.
獻上禱告奉基督耶穌得勝的聖名.
Amen.
我們一齊來看看這篇詩篇.
第一第二節.
我的神我的神為什麼離棄我.
為什麼遠離不救我.
不聽我哀慳的言語.
我的神我白日呼求你不應允.
夜間呼求並不助聲.
這個第一第二節.
在標題裡面我們知道.
是大衛裡面他自己經歷了一些苦難.
他寫下來的一個詩篇.
當然叫做大衛之詩.
是大衛自己寫的.
又或者用大衛的名義.
一篇關於大衛的詩.

$^{161}$又或者有人借助大衛的名氣.
大衛一個詩人一個君王.
他的身份來作這個詩篇.
無論怎樣都好.
在這個詩篇裡面.
我們就用大衛的角度來看.
大衛他都經歷過一些的苦難.
而且那個都是一個很不容易的.
人生一段很艱難的日子.
在他艱難的日子裡面.
他面對生命裡面有這麼多的痛苦.
他的禱告是.
上帝為什麼離棄我.
為什麼遠離我.
為什麼不聽我哀慳的言語.
在這裡裡面.
大衛問上帝為什麼.
當我們今天.
我們面對著一些弟兄姊妹.
人生裡面有一些的痛苦.
我們有很多時候就會說.
你不要問為什麼.
你解釋不了.
所以你應該問如何.
不要問為什麼.
所以今天這個世界的解讀就是.
我的神我的神.
如果被弟兄姊妹勸勉了之後.
如何令到我不被你離棄.
如何令到我跟你緊密結連.
如何令我的聲音.
可以被你聆聽到.
但是你讀上去.
你會發覺如果所有東西都變了.
如何呢.
你會完全體會不到.
在這裡大衛很直率地.
將他心靈的呼喊.
去問上帝.
他面對著生命裡面很多的艱難.

$^{201}$他問上帝為什麼會出現這些狀況.
上帝為什麼你不聽我.
甚至上帝為什麼你不救我.
一問為什麼的時候.
不是一個不神性的一個發問.
為什麼很多時候.
除了問這個事件為什麼有這樣的成因.
有一個源頭之外.
更加重要的就是.
上帝我在這一刻裡面面對的苦難.
我的意義在於什麼呢.
苦難和一般人生裡面的.
那些叫做逆境際遇.
有些不同的.
苦難很多時候動搖的.
不是在說你理性明白了那個事件.
你就可以繼續走下去.
苦難很多時候動搖的.
就是你對上帝的信心.
動搖了你人生的存在價值和意義.
所以有這些事件出現的時候.
是會令到你心裡面有一個重大的打擊.
甚至這件事的出現.
會令到你整個人倒塌下來.
大衛在這裡某程度上也是這樣說.
上帝的離棄.
上帝的不拯救和上帝的不聽.
好像孤身一人面對著這樣的狀況.
在第二節裡面他再繼續說.
他日間又呼求上帝不應允.
夜間又呼求不停止的來祈禱.
但是他所禱告的對象沉默不言.
上帝不出聲.
他不是否定上帝的存在.
他只是說他的禱告上帝沒有回應.
不出聲.
他就好像對著空氣來講他自己的苦況.
他覺得自己孤身一人.
來面對著生命的這種狀態.
為什麼是成為禱告的一部份.

$^{241}$為什麼不是質問上帝.
也不是帶著一份不敬.
為什麼是私人在當下的處境.
和他能夠有力量面對的處境.
成為了一個很大的差距.
這種落差令他不能自拔.
心體裡面有一種很強烈的痛苦.
然後由第三節開始.
詩篇裡面再繼續用一個.
我們今天會經常用的音樂.
那個叫做福音.
不是我們傳的福音.
而是音樂裡面有一種音樂不同的聲部.
好像在唱兩首歌放在一首歌裡面的那種.
互相的不同的聲音.
在同一首的樂曲裡面出現.
在第三節裡面這樣說.
但你是聖傑的.
是用以色列的讚美為寶座的.
我們的祖宗倚靠你.
他們倚靠你.
你便解救他們.
他們哀求你便解救.
他們倚靠你.
就不羞愧.
第六節.
我是蟲不是人.
被眾人羞辱被百姓藐視.
看見我的刀刺少我.
他們撇嘴搖頭說.
他把自己交託耶和華.
耶和華可以救他吧.
耶和華既喜悅他可以答救他吧.
第九節.
但你是叫我出母腹的.
我在母腹裡你就使我有一有倚靠的心.
我引用這一段.
他見到一個很重要的句式.
我間中讓大家看到.
就是但你時.

$^{281}$第六節是但我時.
第九節是但你時.
在這段經文裡面.
其實是但你但我.
但來但去.
但來但去為何是這樣的原因呢.
原來在一個私人的痛苦的歷程當中.
他正在看上帝.
在他歷世歷代裡面.
就是那個保守人的拯救人的上帝.
如何保守他的祖宗.
跨越了經歷了上帝而來的解救.
大衛是讀舊約.
或者在上帝的工作裡面.
他知道上帝一直這樣去看顧.
但他又看回自己.
我的處境是被人羞辱.
應該想著上帝會答救.
但上帝沒有答救.
想著上帝會出手幫助.
但上帝沒有幫助.
他帶來的反而是羞辱,藐視,被人恥笑.
甚至被人搖頭去責他.
他把自己交託耶和華.
耶和華可以救他.
如果耶和華喜悅他.
可以答救他.
那些身邊的人,敵人.
用上帝不應允許他,不答救他.
成為了恥笑他,藐視他的一些內容.
他看上帝過去歷史歷代.
上帝這樣保守別人.
看回自己.
他不是這樣經歷的.
然後他又彈回上帝那裡.
但你叫我出母腹的時候.
我在母腹裡.
使我有一個倚靠的心.
你會見到在這裡.
有種心靈的交集.

$^{321}$看著上帝.
看到上帝很偉大的工作.
看著自己.
覺得自己不是經歷這樣的處境.
然後彈來彈去.
代表著心靈裡的一種複雜性.
近代很多的識經學家.
特別是用詩篇作為模樣的學者.
他就用這種彈來彈去.
那種心靈裡面.
看著上帝又看著自己.
那種複雜的.
就叫這種叫做福音性.
正如我剛才解釋的所說.
這種福音性.
就是講到在生命當中裡面.
他面對著的就是一種很大的困難.
而這種困難和上帝的那種保守.
是同時間在他的心靈當中.
同時間出現的.
兩把聲音都是這麼強烈.
兩個的處境都是真實在他的生命當中.
值得留意的就是.
在經文裡面.
他用這樣的福音性.
同樣都應用了在馬太福音.
在耶穌被釘十架的時候.
馬太的作者都用這個言辭.
來描述耶穌所經歷的處境.
在十六節那裡.
耶穌的手和腳被扎上十字架裡面.
他們扎了我的手和腳.
其實在詩篇第十六節.
就回應了在馬太福音第27章35節.
同樣的他們分我的外衣.
為我的女衣縫製.
在第十八節.
詩篇第十八節裡面.
講到一個這樣的情景.
其實大衛真的沒有經歷過這樣的狀況.

$^{361}$他用私人很誇張的手法.
來描述他被剝奪的痛苦.
但是在馬太福音.
耶穌就成就了這種被剝奪.
在十字架上的痛苦.
馬太福音第27章35節.
這種被嘲笑.
耶和華說可以救他的.
你就可以答救他吧.
後來都繼續的.
來嘲笑耶穌在十字架上.
這種福音.
那種雙重的聲音.
既是有上帝昔日的帶領.
上帝的工作.
但同時當下心靈的那種痛苦.
兩樣東西在一個受苦的人生命裡面.
同時間出現.
我們叫這個福音性.
傳統以來.
我們通常安慰人.
我們就會壓低了個人.
人生命裡面那種痛苦的聲音.
我們有時候會覺得不耐煩.
對於有些人面對著苦難.
我們覺得那些聲音不是來自上帝的.
我們甚至覺得那些人生命裡面的那種痛處.
就是那個人的信心.
或者他的屬靈深度不夠.
以致他沒有辦法去克服眼前這些困難.
不過詩篇22篇裡面.
你同意那個叫做福音性.
他當下的痛苦.
上帝基本上是沒有否定到他的.
甚至乎人生命裡面當下面對的那種苦難.
和上帝昔日裡面的帶領.
在一個人生命當中是同時存在的.
我們不需要壓低了我們自己個人的那種苦況.
又或者覺得那樣東西是不神聖.
反而是當下.

$^{401}$在這個處境當中.
我們的心靈那種痛苦.
上帝是聆聽.
詩人是這樣來祈禱.
耶穌也在十字架裡面作出這樣的呼喊.
而我們作為一個跟隨耶穌基督的人.
我們的生命都是面對著同樣的艱難.
人生裡面充滿了這種福音.
我們一方面相信上帝.
在我們人生裡面.
或者見到其他弟兄姊妹.
如何安然地渡過一個很難的處境.
但回望自己有時候.
可能在這些關口.
在一些苦難當中.
我依然是很痛苦,很掙扎.
我們面對著的就是一個信仰.
和我們的生命裡面一個很大的落差.
兩把聲音,兩個處境都是真實.
在我們的生命當中.
我們每一個人都在經歷這樣的狀況.
早前我和太太去探訪教會的弟兄.
我看著他的時候.
我很明白一個這樣的狀況.
這個老人家在療養院.
因為他癌症已經到達末期.
我去探望他的時候.
我第一件事問他.
細伯你覺得腳怎樣?.
你的身體有沒有地方很痛?.
很少人在癌症裡問他痛不痛.
只是問他的信心堅固嗎?.
第一件事我問他腳痛不痛.
或者身體有哪個地方痛.
他就好像很飾然地說.
他的癌細胞已經去到骨.
所以他的膝蓋很痛.
所以他覺得很痛苦.
在這麼痛苦的情況下.
因為我是牧者,我去探望他.

$^{441}$他硬著頭皮來招呼我.
其實我覺得不用招呼.
我說你累了就閉上眼睛就行.
但是他強行睜大雙眼來跟我聊天.
我知道他其實很疲累.
但對於一個老人家來說.
一個病痛纏繞生命的人.
生命最軟弱的人.
他其實沒有必要跟你說謊.
來哄你.
他在那一刻裡面.
我說細伯我跟你讀一段聖經好不好?.
我就跟他讀一段聖經.
讀聖經的時候.
他會說一句.
他就很大力地說「是的,沒錯,是的」.
我看到他的腳是捲曲的,是痛的.
但每一次說聖經的時候.
他說「是的」.
他的痛楚和他的信心是連在一起.
又或者正正是因為這種痛楚.
令他的信心變得更加真實.
他的痛楚不會因為我認信.
那樣東西就消失了.
他的痛楚依然在他的認信.
說阿們或者聆聽聖經.
他認同上帝所說的每句話的時候.
那個痛楚依然在他入心入肺地痛.
但是在痛楚裡面.
他依然是認信.
大衛在這裡面就是在說這樣的狀況.
他經歷了上帝神聖的拯救.
或者他的祖先.
他的過去是這樣經歷拯救.
但當下他面對的痛苦.
在人生裡面我們就一起這樣經歷.
所以這裡面說的那個福音.
就是人生裡面的兩把聲音.
在你的心靈當中作出呼喊.
在這裡面大衛就用這個詩篇.

$^{481}$來描述一個這樣的狀況.
直至到第十九節.
二十二篇第十九節.
經文裡面就作出一個和合本.
你看不到那個翻譯.
用新譯本讀給大家聽那個內容.
新譯本裡面這樣寫.
至於你 耶和華.
求你不要遠離我.
我的力量 求你快來幫助我.
求你搭救我的性命脫離刀劍.
搭救我的生命脫離惡狗的爪.
求你拯救我脫離獅子的口.
拯救我脫離野牛的角.
你已經英雲了我.
整個詩篇裡面的轉捩點.
就是在第十九節到第二十一節.
如果你看和合本.
你看不到的和合本沒有翻譯到.
那三個字就是至於你.
其實那個至於你.
就是剛才裡面所說的.
彈來彈去那個字.
原文用的是彈你.
不過新譯本為了要.
好像有些轉捩點要告訴你.
有些不同的.
他就寫至於你.
其實那個意思是彈你.
他又彈回到上帝那裡.
他就求上帝不要遠離他.
求上帝來幫助他.
搭救他脫離刀劍.
脫離惡狗的爪.
脫離獅子的口.
脫離野牛的角.
這個是一個求上帝祈求來的.
不過後面加了一句.
我們完全覺得是很詫異的一句.
你已經英雲了我.

$^{521}$接著下面所說的就是一個讚歌.
而在這個讚歌裡面.
或者這個頌讚的詩篇.
後面都沒有說到.
這個詩人脫離了什麼苦境.
所以在這個祈求.
到上帝的英雲.
其實那個英雲不是說.
上帝立即將這個人.
從一個苦難裡面.
拉走.
上帝就平息了所有的困難.
或者上帝將這些刀劍遠離了他.
這個詩人依然是面對著同一個處境.
在這篇詩篇裡面.
那個轉捩點就是你已經英雲我.
這個你已經英雲我是什麼意思呢.
有些的識經學家用這一句話.
就是在說一個信心完成式.
什麼意思呢.
就是說他用一個英文裡面說的完成式.
就是說那件事已經完成了.
那就是完成了那件事.
不過在詩人面對的當下處境.
他依然面對著苦難.
面對著那種困難.
但是他的信心.
令到他看見上帝已經答救了他.
上帝已經幫助他.
所以很多的舊約學者.
就用這個這樣的句式.
其實在先知書在舊約裡面.
經常都出現一個這樣的句式.
就是說現實的處境是這樣.
但是因為他的信心.
超越了他的環境.
成就了一個叫做信心完成式.
所以整個詩篇裡面.
他講到這個詩人.
他彈來彈去.

$^{561}$但是到最後.
他可以走出這個困境的出路.
是在於他的觀念改變了.
眼前依然充滿了挑戰很多的困難.
但是因為他的意念.
他對上帝的心.
他帶著信心來仰望上帝.
這一份信心幫他超越了眼前的環境.
苦難能夠勝過.
就是因為你的信心.
善於環境的改變.
環境還沒有改.
但是你相信上帝已經幫助了你.
你的生命就可以渡過一個最低谷.
最低下的最困難的處境.
這種狀況正正是很多時候.
面對著人生與死的時候.
人心靈裡面能夠經歷到上帝.
人認順上帝.
甚至知道人生的意義.
人生命的價值.
不是肉體那幾十年的日子.
而是可以看到肉體之後.
上帝救贖的日子.
淋到他心靈當中.
他就有一份踏實安穩.
來渡過眼前即使都是病患.
又或者滿佈很多困難的痛苦經歷.
我們教會裡面有很多弟兄姊妹.
曾經都經歷過一些這樣的狀況.
當我和他們去探訪.
當我們一起去祈求上帝的時候.
我看到很多的家庭.
甚至有很多未必一定是生與死的.
很多是經歷經濟困難的弟兄姊妹.
他們其實是足襟見爪的.
不過他說祈禱之後.
交託了他就可以安心.
因為他知道上帝已經為他預備.
祈禱之後他不是立刻銀行戶口多了很多錢.

$^{601}$又或者眼前所有的困難都沒有了.
不過就是單單一份他對上帝的認訊和信心.
那份倚靠就可以幫他勝過眼前這些苦難的處境和能力.
所以在經文裡面.
詩人就是一個這樣的認訊.
令他由一個愛歌變成一個讚歌.
關鍵就是他有一個這樣心裡的信心.
接著去到第22節打後.
就是整個詩篇裡面讚美的那部分.
我讀22-26節給大家聽.
我要將你的名傳於我的弟兄.
在會中我要讚美你.
你們敬畏耶和華的人要讚美他.
雅各的後裔都要榮耀他.
以色列的後裔都要懼怕他.
因為他沒有藐視憎惡仇父的人.
也沒有向他掩面.
那受苦之人呼籲的時候他就垂聽.
我在大會中讚美你的話是從你而來的.
我要在敬畏耶和華的人面前還我的怨.
謙卑的人必吃得飽足.
尋求耶和華的人必讚美他.
願你們的心永遠活著.
在22-26節裡面.
詩人到上帝的會堂大會裡面讚美他.
然後講到雅各的後裔.
以色列的後裔都要敬畏耶和華.
特別在第24節裡面這樣說.
他沒有藐視憎惡仇父的人.
也沒有向他掩面.
在這裡面特別強調.
上帝沉默的狀況.
不是代表著上帝對世人面對的苦難.
掩面不看或藐視這些人.
甚至這些受苦的人覺得他不耐煩.
上帝體恤人的痛苦.
上帝明人內深深處的那種負憾.
他在呼籲人生命面對著苦難的時候.
轉念以信心來回應.
但上帝從來不會用.

$^{641}$你沒有信心所以你面對著這些事.
你沒有力量來面對.
用一個優勝劣敗的角度.
來看一個面對著苦難的人.
相反上帝在呼籲每一個人.
在你苦難當中.
你仰望這個他不稀罕你.
面對著生命所有困惑.
動搖你生命很多意義的處境.
將你的生命焦點放在上帝那裡.
上帝從來沒有藐視你.
也沒有憎惡你.
沒有向你掩面.
反而當你呼籲的時候.
上帝必然來到你垂聽.
當然垂聽的結果.
未必一定是將你眼前那種苦況.
馬上挪走的.
可能有些弟兄姊妹是這樣.
但有更多的可能是.
你的實況依然繼續.
不過你可以在患難當中.
你可以仰望天上的主.
你能夠有力量.
可以面對著生命每一種的艱難.
很多人面對著艱難的時候.
他會失去那種自尊.
他會覺得自己在世界上.
是一個最低下的人.
甚至覺得自己是一個沒用.
甚至可能被人藐視的一個人.
我想起我家裡有一個這樣的經歷.
我最小的妹妹.
我和大家說過.
我爸爸在我十六歲時去世.
我妹妹三,四歲.
很小的時候.
在她成長的時候.
大部分時間爸爸是缺席的.
我記得有一次她哭著回家.

$^{681}$她和我在那裡說.
有一篇作文是關於爸爸的.
她自己亂作.
作了她爸爸的一些東西.
但她在生命裡沒有經歷過爸爸.
所以她亂作.
作完之後.
不知道是否得到高分.
但回家她很不開心.
因為她知道那件事不是真實.
然後我說為什麼你不跟老師說.
寫不到爸爸就寫媽媽.
或者寫其他家人.
為什麼要說謊呢.
她覺得很羞恥.
她覺得她沒有爸爸.
她失去了父愛.
在她小時候沒有爸爸照顧的生命.
她覺得被人好像有種缺憾一樣.
她的生命好像有一種虧欠.
總之很難說出來的那種很難過.
又覺得很尷尬.
不懂得怎麼說.
她說原來她大部分時間都是經歷這樣的階段.
我記得那時候我都不懂得回應她.
我媽媽回應她.
我媽媽說傻女.
你爸爸又不是做壞事.
他是病離開這個世界.
他又沒有做一些對不起家庭的事.
為什麼你會用羞愧來描述他的離開呢.
大膽一點跟別人說.
跟別人說分享你的生命就是這樣來經歷.
每個人都有眼淚的.
不要因為這些困難而眼淚.
就否定自己.
我們很多弟兄姊妹.
包括我自己都有這樣的經歷.
困難,苦難.
或者看到一些事情幫不了忙.

$^{721}$幫不了忙.
就會覺得自己低人一下.
羞愧,沒力的那種感覺.
覺得上帝離開自己.
不過這篇詩篇裡面說到.
上帝沒有輕看人的那種痛苦的吶喊和那種呼求.
上帝沒有輕看我們.
反而我們學會尋求他.
你找他的時候.
你的生命就得到那種踏實,安穩.
耶穌為什麼在十字架上引用這篇詩篇呢.
耶穌引用這篇詩篇.
不是純粹為了滿足我們識經的慾望.
耶穌大概在當下.
在十字架上他都在經歷一個.
身心靈受到極大痛苦的那種掙扎.
那種壓迫感.
他向上帝呼喊.
我的神,我的神為什麼離棄我.
不是純粹是一個悲壯式的演繹.
而是在一個痛苦裡面.
他和我們血肉之軀一樣.
他都會向上帝呼喊.
當然耶穌基督所引用的第一句.
其實在很多的識經.
或者我們都知道.
引用第一句就是引用整首的詩篇.
耶穌在痛苦當中.
他依然信靠上帝.
依然來到學習,來到去信從.
希伯來書第二章九至十節裡面這樣說.
唯獨見那成為比天使少一點的耶穌.
因他受死的苦.
就得了尊貴榮耀為冠冕.
叫他恩著神的恩.
為人人嘗了死味.
原來那為萬物所屬.
為萬物所本的.
要領許多的兒子進榮耀裡去.
使救他的元帥.

$^{761}$因受苦難得以完全.
本是合而的.
在這裡希伯來書說到耶穌基督受苦難.
因受苦難得以完全.
不是說耶穌是神.
本來不是神受苦難之後成為神.
不是這個意思.
而是說使命完全.
救贖將神的兒女帶到榮耀當中.
這個救贖計劃.
因著耶穌這種受苦難.
這種痛苦.
救贖工作就在耶穌的受難當中.
來成就出來.
苦難對於我們這個世代.
對於很多人來說是一個荒謬的事件.
不過基督教和其他宗教不同.
我們所信的上帝是進入苦難當中.
我們所信的上帝是承受苦難的耶穌基督.
我們所信的主.
他自己親身嘗過痛苦的滋味.
苦難成為了他救贖恩典的一部分.
苦難也成為了他神聖的一部分.
苦難也成為了.
藉著痛苦明白耶穌的心腸.
多一些的途徑.
所以弟兄姊妹.
當我們看詩篇第22篇.
我們明白.
原來詩人在一個這麼痛苦的經歷當中.
他好像親身經歷說得大了.
而他用詩的角度來描述的時候.
在耶穌基督就傳言來實現.
不過重點是.
詩人面對的苦難.
同樣是基督和耶穌面對的苦難.
正正因為耶穌基督面對的苦難.
我們世人每一個人.
我們人生都會經歷的.
或大或小我們不需要比較.

$^{801}$每一個人的苦難是不同的.
但是有一樣東西是共同的.
當我們去尋求上帝.
當我們人生面對大大小小的苦難.
當我們向祂呼求.
用信心仰望祂的時候.
我們的人生就有力量來跨勝.
我們的生命可以超越眼前所有的境界.
最後說一個例子.
在教會裡服事.
以前有些神學生.
我記得有一個神學生問我.
牧師我做喪禮的時候失敗了.
我說失敗了什麼.
他說在喪禮裡看到主家和家庭.
我心裡很感動.
我在喪禮裡哭了.
然後我流出眼淚.
我有一點失態.
我覺得自己失敗了.
他就想找我談談.
我說流眼淚有什麼問題.
他說我不專業.
如果我自己流眼淚我怎樣安慰別人.
我跟他說你找錯人了.
我說我做了十幾二十年.
基本上我每次喪禮都會哭.
不過我想說.
我們牧會和一般的服務人員不同.
我們會動情的.
那些是我們所愛的人.
我們有感受.
人就有感受.
因為愛就有感受.
上帝不是一個冷冰冰坐在天上.
看你人間疾苦的上帝.
我們所信的主是有情的上帝.
耶穌基督在天上.
都是用祂的慈愛和情.
來和我們一同感受人生在地上所經歷.

$^{841}$所以弟兄姊妹.
我不是叫你哭啼啼走過人生的每一步.
不過我們不用覺得有情感.
就覺得自己不能夠安慰別人.
反而與哀哭的人哀哭.
與歡笑的人同笑.
這是我們世代需要的一種生命.
耶穌在十誡經歷患難.
詩篇二十二篇描述患難.
我們今天可能會在不同的苦難當中.
但讓我們帶著信心.
帶著上帝的應許.
我們有力量走人生的道路.
願主幫助我們.
(影片完結).
\newpage



\section{馬可福音 4:1-20}
\label{sec:yBE6W2fCqSM}
\textbf{2026年3月8日崇拜直播｜黃偉立牧師｜門徒必備的那種態度｜馬可福音4\_1-20}
\newline
\newline
連結: \href{https://youtube.com/watch?v=yBE6W2fCqSM}{\texttt{https://youtube.com/watch?v=yBE6W2fCqSM}} ~~~~ 語音日期: 2026-03-09
\newline
\newline
\hyperref[sec:8qvQyuuLcyA]{\small{< < < PREV SERMON < < <}}
~
\hyperref[sec:index_chronic]{\small{[返順時目]}}
~
\hyperref[sec:PGMHgAUtHoA]{\small{> > > NEXT SERMON > > >}}
\newline
\newline
馬可福音 4:1-20
\newline
\begin{longtable}{cl}
\hline
\hline
章節 & 經文 (和合本修訂版)\\
\hline
4:1 & \begin{tabularx}{0.7\textwidth}{X} 耶穌又在海邊教導人。有一大群人到他那裡聚集，他只好上船坐下。船在海裡，眾人都靠近海，站在岸上。 \end{tabularx} \\ \\ \relax
4:2 & \begin{tabularx}{0.7\textwidth}{X} 耶穌就用許多比喻教導他們。在教導的時候，他對他們說： \end{tabularx} \\ \\ \relax
4:3 & \begin{tabularx}{0.7\textwidth}{X} 「你們聽啊，有一個撒種的出去撒種。 \end{tabularx} \\ \\ \relax
4:4 & \begin{tabularx}{0.7\textwidth}{X} 他撒的時候，有的落在路旁，飛鳥來把它吃掉了。 \end{tabularx} \\ \\ \relax
4:5 & \begin{tabularx}{0.7\textwidth}{X} 有的落在土淺的石頭地上，因為土不深，很快就長出苗來， \end{tabularx} \\ \\ \relax
4:6 & \begin{tabularx}{0.7\textwidth}{X} 太陽出來一曬，因為沒有根就枯乾了。 \end{tabularx} \\ \\ \relax
4:7 & \begin{tabularx}{0.7\textwidth}{X} 有的落在荊棘裡，荊棘長起來，把它擠住了，就結不出果實。 \end{tabularx} \\ \\ \relax
4:8 & \begin{tabularx}{0.7\textwidth}{X} 又有的落在好土裡，就發芽長大，結出果實，有三十倍的，有六十倍的，有一百倍的。」 \end{tabularx} \\ \\ \relax
4:9 & \begin{tabularx}{0.7\textwidth}{X} 耶穌又說：「有耳可聽的，就應當聽！」 \end{tabularx} \\ \\ \relax
4:10 & \begin{tabularx}{0.7\textwidth}{X} 耶穌獨自一人的時候，跟隨他的人和十二使徒問他這些比喻的意思。 \end{tabularx} \\ \\ \relax
4:11 & \begin{tabularx}{0.7\textwidth}{X} 耶穌對他們說：「神國的奧祕只讓你們知道，若是對外人講，凡事就用比喻， \end{tabularx} \\ \\ \relax
4:12 & \begin{tabularx}{0.7\textwidth}{X} 要他們看了又看，卻看不清，聽了又聽，卻不明白，免得他們回轉過來，獲得赦免。」 \end{tabularx} \\ \\ \relax
4:13 & \begin{tabularx}{0.7\textwidth}{X} 耶穌又對他們說：「你們不明白這比喻嗎？這樣怎能明白一切的比喻呢？ \end{tabularx} \\ \\ \relax
4:14 & \begin{tabularx}{0.7\textwidth}{X} 撒種的人所撒的就是道。 \end{tabularx} \\ \\ \relax
4:15 & \begin{tabularx}{0.7\textwidth}{X} 那撒在路旁的種子，就是人聽了道，撒但立刻來，把撒在他們心裡的道奪了去。 \end{tabularx} \\ \\ \relax
4:16 & \begin{tabularx}{0.7\textwidth}{X} 那撒在石頭地上的，就是人聽了道，立刻歡喜領受， \end{tabularx} \\ \\ \relax
4:17 & \begin{tabularx}{0.7\textwidth}{X} 因心裡沒有根，不過是暫時的，一旦為道遭受患難或迫害，立刻就跌倒。 \end{tabularx} \\ \\ \relax
4:18 & \begin{tabularx}{0.7\textwidth}{X} 還有那撒在荊棘裡的，就是人聽了道， \end{tabularx} \\ \\ \relax
4:19 & \begin{tabularx}{0.7\textwidth}{X} 後來有世上的憂慮、錢財的迷惑，和別樣的私慾進來，把道擠住了，結不出果實。 \end{tabularx} \\ \\ \relax
4:20 & \begin{tabularx}{0.7\textwidth}{X} 那撒在好土裡的，就是人聽了道，領受了，並且結了果實，有三十倍的，有六十倍的，有一百倍的。」 \end{tabularx} \\ \\
[1ex]
\hline
\hline
\end{longtable}
$^{1}$請聽我讀出.
耶穌又在海邊教訓人.
有許多人.
到他那裡聚集.
他只得上船坐下.
船在海裡.
眾人都靠近岸.
站在岸上.
耶穌就用比喻.
教訓他們許多道理.
在教訓之間.
他對他們說.
你們聽哦,有一個殺種的.
出去殺種.
殺的時候有落在路旁的.
飛了來吃盡了.
有落在土淺石頭地上的.
土技不深.
發苗最快.
日頭出來一曬.
因為沒有根就枯乾了.
有落在荊棘裡的.
荊棘長起來.
把它擠住了.
就不結實.
又有落在好土裡的.
結實有三十倍的.
有六十倍的.
有一百倍的.
又說有耳可聽的.
就應當聽.
無人的時候.
跟隨耶穌的人和十二個門徒.
問他這比喻的意思.
耶穌對他們說.
神哥的奧秘.
只叫你們知道.
若是對外人說.
凡事就用比喻.
叫他們看似看見.

$^{41}$卻不曉得.
聽似聽見卻不明白.
恐怕他們回轉過來.
就得赦免.
又對他們說.
你們不明白.
這比喻嗎.
這樣怎能明白.
一切的比喻呢.
殺中之人所殺的就是道.
那殺在路旁的.
就是人聽了道.
殺但立即來.
那殺在他心裡的.
就是道掉了去.
那殺在石頭地上的.
就是人聽了道.
立即歡喜領受.
但他心裡沒有根.
不過是暫時的.
以及為道遭了患難.
或是受了迫迫.
立即就跌倒了.
還有那殺在荊棘裡的.
就是人聽了道.
後來有世上的.
思慮.
錢財的遺惑.
一樣的私慾盡來.
把道祭住了.
就不能結實.
那殺在好土地上的.
就是人聽了道.
又領受並且結實.
有三十倍的.
有六十倍的.
有一百倍的.
今天為我們證道的.
大家都很熟悉他.
他之前是翻案證的.

$^{81}$現在已經成為牧師.
是中華宣導會友愛台牧師.
黃偉立牧師.
今天的證道題目是.
門徒必備的那種態度.
我把時間交給他.
各位靈姐妹早安.
其實我本身預計.
我一講就會哭.
一回來.
我現在就哭不出來.
因為我可以說是欲哭無淚.
因為剛才我本來在屯門出來.
我打算開車過來.
一開車出去.
一出停車場去到路面.
突然聽到啪一聲.
我車的tie 盤就斷了.
然後車就扭不到彎.
我就整輛車塞在門口.
很危急.
因為我預租了時間.
我也感恩.
但我又要立即找人幫忙.
不知車房怎麼處理.
總之弄了很久.
又發訊息給Anders.
為我祈禱.
終於來得到.
但驚魂未定.
到現在心還是在跳.
所以一會兒.
如果我說話漏漏氣.
不知道自己在說什麼.
可能不太順暢.
大家就多多包涵.
今天真的很開心.
可以回到旺朕的家.
見到大家.
給我機會跟大家分享神的話語.

$^{121}$我覺得回到這裡.
一磚一瓦.
都是我的回憶.
我以前坐在台下.
聽Anders 林艷芳牧師.
梁念安牧師講道.
我最有印象.
就是以前的林凱勝牧師.
可能大家都有印象.
他每次上來.
第一句就是.
碧兒!碧玲!碧玲!.
我杯水呢?.
他經常都很口乾.
經常呼喚碧玲.
碧玲姨.
這些印象好像只是昨天的事.
非常開心.
今天在這裡.
餵養我很多的地方.
我今天在這裡可以講道.
這是我自己的榮幸.
數一數手指.
我已經十二年沒有在台上.
對上一次最有印象.
的那一天.
我就在這裡.
做崇拜主席.
一晃眼就回到這裡.
身份是不同了.
不過我依然都是.
當初那個傻乎乎.
很單純很想去服侍主的.
那個阿娜.
我沒有忘記過.
我是一個旺朕的一份子.
只不過是回頭看.
我已經是老旺朕了.
剛才想一想.
其實我回到這裡.

$^{161}$我沒有覺得自己離開過.
只不過是沒有空出來.
屯門要出來旺角是很遠的.
是很多里路的.
我都希望我自己再回來的時候.
就好像網上的一句名言.
願我出走半生.
歸來仍是少年.
見到大家.
我覺得很多很熟悉的面孔.
我覺得很親切很感恩.
但我更感恩的是很多我不熟悉的面孔.
旺朕依然是一間很有活力.
很有吸引力的教會.
吸引很多人在這裡一起去敬拜他.
認識他.
可以為這件事感恩.
願神繼續去賜福教會.
繼續使用教會.
今天想和大家看的經文.
可能大家都很熟悉.
就是殺種比喻.
不過熟悉歸熟悉.
很多人對這段經文其實都有誤解.
首先有人看到比喻裡面形容殺種的人.
他殺種的方法.
馬上批評他一定是一個外行人.
因為稍為有知識的都知道.
種子一定要殺在泥土裡面.
才會成長.
這件事我兩個女兒.
在幼稚園的時候已經學會了.
我想大家做過父母的話.
都會有這個經驗.
家裡突然多了很多盆栽.
就是學校給他們的.
要拿回家照顧盆栽.
就是要他們學會這件事.
雖然我們的經驗都告訴我們.
通常最後照顧他的肯定都是父母.

$^{201}$最後淋水的都是做父母.
因為他們三分鐘熱度之後.
最後就是我們跟他一起學習這件事.
但是這件事依然是很重要的.
因為家裡可以多一點綠色.
小朋友多一份知識.
我們都要多謝學校.
但是比喻裡面殺種的人.
他就好像真的亂來.
因為他殺去的地方.
有些根本是沒有泥的.
是不會成長的.
是浪費了種子.
究竟這個人是不是真的不懂呢.
我們不知道.
但是有聖經的考古學者跟我們證明.
古代是真的有人這樣殺種的.
至少不可以用我們現代人.
很有系統去耕種的方法.
去理解當時怎樣殺種.
不過最重要最重要都是.
要看耶穌他怎樣解釋.
這個殺種的人.
這個殺種的人其實是代表著.
傳播神的道.
他可以是神或者耶穌.
就好像耶穌現在講這個比喻的時候.
他就是在傳播神的道.
他也可以是你和我.
我們都會傳福音.
我們會將神的道傳出去.
無論是哪一個都好.
無論是神是耶穌還是我們都好.
傳神的道的人是會亂傳的嗎.
不會的.
所以比喻就還比喻.
我們不可以想著每一個細節.
都要對應一些意思出來.
只能夠從整體的比喻裡面.
去理解一下重點是什麼.

$^{241}$這裡的重點很明顯是在講.
無論聽道的人他的態度是怎樣.
神的道.
是為所有人而傳開的.
任何人都有機會聽到.
就好像種子.
是可以殺在不同的地面.
只不過是會產生出不同的效果.
第二個誤解.
有人認為這個比喻.
其實不是教導.
只不過是一個事實的陳述.
就好像有人跟你講.
一年有365日.
一年有四季春夏秋冬.
一日有24小時.
一個星期有七天.
所以這裡很直白地告訴你.
這個世界有四種土地.
有些是路邊.
有些叫土千石頭地.
有些叫荊棘.
有些叫好土.
比喻不同的人聽到神的道.
就有不同的反應.
沒有其他意思的.
就是這樣.
當然不是.
我們看回整個的解釋.
很明顯耶穌的殺種比喻.
是一個徹頭徹尾的教導.
是對我們有要求.
是對每一個聽到這個比喻的人.
是有一個要求.
是想我們要有改變.
是有不同的.
事實上這個比喻故事的背景.
一開始已經告訴我.
這個是耶穌給當時的人的一個教訓.
有一天.

$^{281}$耶穌去到海邊.
當時他已經是一個很出名的人.
所以一如以往有很多人聚到他那裡.
但人實在太多了.
耶穌唯有走上一艘船.
上面的所有人.
用比喻去教導他們.
不過我們可能都會問.
當時真的有這麼多人想聽耶穌說話嗎?.
我們都有印象.
曾經有一班人見過五餅二魚之後.
跟著耶穌.
耶穌說你們想來找我吃東西.
所以究竟是不是真的有這麼多人.
想聽耶穌說話呢?.
或者耶穌說的殺種比喻.
正是想說出一些真相.
如果種子落在沒有泥的路上.
因為發不了芽.
所以一有鳥飛過.
就會吃掉.
耶穌後來解釋.
這樣的情況就像一個沒有心聽道的人.
聽了道就像被魔鬼偷走了.
不會相信.
也不會成長.
如果種子落在有一點泥的土淺石頭地上.
因為有泥.
所以會發芽.
但因為根不夠深.
所以不會成長得好.
只要太陽曬得猛一點.
就會曬死他們.
覺得是很好的道理.
又感動又導人向善.
就可以很開心地去信耶穌.
但如果他其實是沒有心的.
他沒有心是想認識多一點這個信仰.
後來發現.
原來信耶穌還要受這麼多苦.

$^{321}$原來不是這麼懵的.
還要承受因為信耶穌而帶來的.
那些很苦澀的經歷.
就是耶穌說的.
為道遇到患難.
或是受了逼迫.
例如可能在朋友.
或是工作的地方提出一些我們信仰的原則.
不是的.
我們可以對同事多一點愛和包容.
類似這些說法.
可能會被人取笑.
我以前也試過.
我還在新城電台工作的時候.
有一天.
有一班同事在說色情笑話.
色情笑話當然我會聽.
說著說著.
大聲地笑.
突然有一個女士.
她看著我.
說:阿立.
你們教徒不可以分前行為.
這樣就很辛苦了.
你聽到這樣的問題.
你會怎樣回答她?.
我真的不懂回答她,我呆了.
難道你會回答她.
我們很辛苦的,多謝你諒解.
多謝關心.
你不會這樣回答的.
所以那一刻我真的不知道該說什麼.
幾秒之後,我不知道從哪裡來的智慧.
我就大聲地對她說.
你不覺得.
精粹是很重要的嗎?.
她就.
哈哈哈哈.
就這樣走了.
很明顯對於未信的人來說.

$^{361}$我們這些想法是很好笑的.
又或者.
平日可能回教會而錯失了很多.
跟朋友相聚.
或者搏殺上位做事的機會.
他們沒有想過原來信耶穌要付代價.
所以一受到小小的苦.
就放棄了.
當初信耶穌也是為了開心.
你現在要跟我說要受苦.
那就別搞了.
我記得我以前還是年輕人的時候.
在旺鎮這裡聚會.
每年都會參加一些暑期的夏令營.
到現在還是有的.
我記得有一年參加完之後.
通常都會去一堆山旮旯的地方.
通常星期日.
下午回來就會.
無論如何都要上來教會這裡.
有一點最後的分享.
大家在五樓那裡找個房間分享.
其中一個.
是一位新來的姊妹.
參加完這個營會之後.
就覺得跟我們一樣.
在營會裡面無論是訊息.
說什麼,唱什麼.
都覺得很感動,很有意思,很有得著.
突然間姊妹就說.
我覺得你們這群人都很好.
我也很想繼續信耶穌.
聽到就很開心.
突然間她說了一句.
所以我想回去.
我們覺得奇怪.
我們剛剛才放下行李在四樓.
為什麼突然間要回去呢?.
是不是漏了東西呢?.
原來她的意思是.

$^{401}$很想繼續回教會.
我們聽到就覺得.
我們一起努力.
接下這個基因.
過了幾個星期之後.
我就從來沒有見過她.
可能是在開學之後.
很多不同的東西去追求.
她就沒有再出現在這裡.
她可能是有一點泥的石地.
福音的種子.
在那些泥土裡面成長.
當然她也可能是因為下一種情況.
而放棄信仰.
就是種子落在有荊棘的泥土裡面.
如果是這樣的話.
種子一發芽就會被那些密密麻麻的荊棘堵住.
意思就是.
壓到它窒息而成長不了.
情形就好像.
耶穌所解釋的那樣.
聽到的人明顯不會有一個結界.
或者一個保護網.
保護他不受到誘惑.
反而在成長的時候.
有各種的思慮.
錢財的迷惑和別樣的私慾.
例如怕不夠錢花.
或者想讓自己生活過得好一點.
很想賺到第一桶金.
快點.
和滿足一些不同的慾望.
等等這些的誘惑.
經常都會出現在人的身邊.
大家打開不同的社交平台.
例如IG.
Facebook.
無論什麼都好.
裡面都充斥著很多的訊息.
就是告訴你去旅行很開心.

$^{441}$這件事很好.
不斷去誘惑我們.
要去消費.
要去做很多不同的事情.
很多的慾望在裡面.
我相信.
我們身邊還有一些人是這樣的.
就是覺得這些東西很棒.
就全心去追求.
就沒有心去追求這個信仰.
可能你身邊也有一些人.
以前信得很好.
回教會也回得很好.
但是現在呢.
可能會跟你說.
先做好事吧.
我正在搏殺.
這件事很好玩.
先做好事吧.
到我死了的時候.
或者退了休的時候.
我就信耶穌.
不過就好像《雅各書》那樣說.
其實明天如何.
你們還不知道.
又有誰可以肯定.
我們臨死之前.
還有機會.
還有時間.
可以信回耶穌呢?.
唯有落在好肚子裡的種子.
才能夠茁壯成長.
就好像一個人聽了道之後.
願意接受從神而來的教導.
讓神掌管和轉化他自己的生命.
他們這樣的屬靈生命.
才有機會成長.
可以有三十倍.
六十倍.
甚至一百倍.

$^{481}$這麼誇張的成長.
不過如果大家心水清的話.
說到這裡.
你會發現.
耶穌對這四種土地的解釋.
當中應該最重要的是什麼土.
是好土.
但是他的解釋.
簡直是寥寥可數幾個字.
非常簡單.
對比起之前三種土地的解釋.
簡直好像沒有解釋過.
如果從字面.
從字數去解釋的話.
為什麼會這樣呢?.
因為好土的真正意思.
不在於要用言語.
或者道理去理解.
而是在於.
聽到這個比喻的人.
他的真實反應.
可以去顯明出來.
這也是耶穌說這個比喻的真正意思.
什麼意思呢?.
因為耶穌經常用比喻.
我們有些書.
整本都是講耶穌的比喻.
但是大家想清楚.
其實耶穌很少自己解釋這個比喻.
而這個可以說是.
少數耶穌會自己解釋的比喻.
原因是因為有人主動問他為什麼.
第四章第十節說.
沒有人的時候.
即是當耶穌說完要說的話之後.
所有人都散了.
只有耶穌自己一個.
其他福音書記載.
門徒就趁著這個時候去找耶穌.
問他這個比喻其實是什麼意思呢?.

$^{521}$惟有《馬可福音》這裡特別提到.
還有多一班人.
叫做跟隨耶穌的人.
或者就是想表達.
在聽耶穌教導的人裡面.
有些人和門徒一樣.
是有心追求耶穌的.
同時門徒不是特別蠢而不明白.
而是當時很多聽到的人.
其實他們都不明白.
只不過有些人.
選擇就這樣聽完不明白就走了.
而這兩班人.
他們選擇留下.
問耶穌其實這個比喻是什麼意思.
這個就是他們能夠明白比喻的關鍵.
更加重要的是.
他們這個想知道更多的心.
將他們和那些無心的人劃分了出來.
在耶穌眼中.
他們被分為.
你們即是跟隨耶穌的人.
和他們即是所謂的愛人.
我不是說這裡.
你知道這是一個象徵.
他們 愛人.
原來要和耶穌做自己人.
不在乎血緣,身份和地位.
也不在乎耶穌特別喜歡誰.
而只是在乎聽到祂的道的人.
就算不明白.
但最重要是很想追求.
很想知道更多.
這種態度成為一個人.
是不是好土.
生命能不能夠成長的關鍵.
他們和那些無心追求的人.
聽不明白就算的愛人是不同的.
因為那些愛人就是.
撒種比喻頭三種土地的真實寫照.

$^{561}$有趣的是.
耶穌說祂用比喻.
竟然是為了要這群愛人聽明白.
四章第十一節說.
神國的奧秘只叫你們知道.
若是對愛人說.
凡事就用比喻.
這句原文的意思.
直譯的話就是.
神國的奧秘只會賜給你們.
而對外面的人.
凡事就會變成比喻.
那是怎樣呢.
神國的奧秘.
從來都是由神去啟示人.
人才有機會明白和知道.
而在這裡.
這個奧秘是特別給那些跟隨耶穌的人.
是因為他們有心接受.
和想知道更多.
明白了的比喻.
就成為了他們成長的養分.
就像撒種比喻裡的好土的說法.
至於那些無心的愛人.
他們不問.
聽不明白就算.
那些比喻.
就只是一堆聽不明白的道理.
耶穌還引用了.
以賽亞書一段聽起來很囂張的經文.
去強調他的說法.
「叫他們看事看見卻不曉得.
聽事聽見卻不明白.
恐怕他們回轉過來就得赦免」.
即是耶穌是故意令他們不明白.
免得有人回轉相信他.
當然不是.
他來到這個世界.
就是要人明白天國的事.
因為這樣而能得到豐盛的生命.

$^{601}$甚至第三和第九節.
其實他已經向現場所有人說.
凡有耳朵的都要聽.
所以耶穌這樣說.
其實是帶著一種諷刺.
甚至好像說反話的意味.
去說出一些真相.
其實無論耶穌說什麼.
無心聽的人都不會想知道為什麼.
所以神的道在他們生命裡.
是不會發揮作用.
既然如此.
耶穌會繼續向所有人說.
不過祂希望那些懂得欣賞的人.
會接受和追求更多.
祂就會用更多心機令他們明白.
在這個層面上.
耶穌用比喻.
不單只是用日常生活中.
為人熟悉的東西.
去類比一些道理這麼簡單.
比喻更加是一種對人的考驗.
希望有心的人.
會懷著想知道更多的心去追求他.
生命才能得益.
至於無心的人.
他們就會很快.
因為聽不明白.
又不理會.
而錯過了生命中最重要的東西.
我有一個朋友.
他是自己開公司的.
專門提供一些.
需要很多心機.
很多創意的服務.
他做得很成功.
很多大公司都會找他合作.
除了因為他工作很高質.
很好.
而他行銷的方法.

$^{641}$經營的方法.
都挺不同的.
他跟我說.
如果知道客人找他.
不是真的欣賞他.
可能他做的東西.
也是一些很麻煩.
很多要求.
不會讓他有什麼發揮的話.
他就會設一個.
高到完全不合理的價錢給他.
如果他知道客人.
真的喜歡他的東西.
讓他可以任意發揮的話.
他一元不收.
都可以.
一元都可以不收.
他完全免費跟他做都可以.
他還有一點點囂張.
他說這些我拿來拿獎的.
我說真的.
他做完之後真的拿獎的.
為什麼他會這樣做呢?.
你就這樣聽.
他是趕客的.
不過他想說的是.
如果有人這麼高的價錢都肯做.
那些東西就算不想做.
我死都死定給你.
不過他的創意和心機.
是想留給那些.
懂得真心欣賞他的人.
可能你都會覺得耶穌是趕客的.
但其實他是希望.
有心去追求的人會繼續追求.
生命才有機會成長.
成為更加好的門徒.
在解釋撒種比喻之前.
耶穌問了門徒一個問題.
你們不明白這比喻嗎?.

$^{681}$這樣怎能明白一切的比喻呢?.
什麼意思呢?.
我的理解是.
只要繼續用.
你很想去明白撒種比喻這個心.
去理解其他的比喻.
神就會有機會.
亦都會用心機.
令我們能夠明白一切關於神的道理.
丁姐妹.
這一刻.
你是有心追求耶穌的門徒.
還是聽而就算的愛人呢?.
說起門徒.
我想起我以前在旺鎮.
都有參與過門徒訓練.
早前收拾房子.
那張證書還在.
珍而重之.
而我現在服侍的教會.
友愛堂.
都很重視門徒訓練.
只不過我們有少少不同的是.
很多教會都會將門訓.
門徒訓練.
看成是一個短期的課程.
不過我們教會.
就很喜歡.
亦都特別將它定性為一個.
一年的屬靈訓練.
在這一年的裡面.
門訓每個禮拜.
都會在師父帶著弟兄姊妹的情況下.
反思信仰.
查考聖經.
行山.
去露營.
或者一起去完成一些屬靈目標.
例如可能一年就要看完新約聖經.
又或者.

$^{721}$一起去參與服侍.
一起去參與一些社區的服務等等.
還會寫一個叫做良心審查.
不知道大家有沒有聽過.
類似一些屬靈摘記的東西.
你就這樣聽.
你會覺得.
這些和我們在教會.
平時參加團契小組.
又有什麼分別呢?.
的確是的.
如果你計算每一個項目的話.
其實就沒什麼特別的.
因為門徒訓練的這些項目.
一向都散落在教會不同的場景裡面.
就好像我們今天這樣.
其實我們在聽道.
都是門徒訓練.
你們的靈命.
我們一起都一起成長.
一起聽道就一起成長.
這些都是門徒訓練.
所以我們的門訓講的是.
特別用一年的時間.
擺明車馬.
排除萬難.
加大力度.
開行渦輪.
這樣一起操練我們的靈命.
務求要有所成長.
成為名副其實的門徒.
我們教會過去已經開過兩代門訓.
至少用了兩年的時間.
最近正在開第三代.
在我所帶過的徒弟裡面.
有些人一開始就問我.
會不會太辛苦呢?.
我跟他說.
對比起青少年門訓.
我們即稱門訓真的舒服很多.

$^{761}$青少年門訓會怎樣呢?.
以前就給我們的童主任牧師.
當時他還是一個傳道人.
就會帶一些年青人去海邊.
去碼頭丟他們下水.
不會游泳都會丟.
不過救他們上來的意思.
我的理解可能不是.
可能其實他們會游泳.
但總之丟他們下水.
讓他們很有經歷.
其實他們是這樣的.
還有行山.
我聽過一個故事.
他們通常都是裙邊行山.
走著走著.
傳道那時傳道.
傳道前面沒有路了.
等一下.
他就會自己躺下去.
中間空了的位置.
變了一條橋.
從我背後走過去.
是這樣做訓練的.
所以這些年青人特別辛苦.
但印象很深刻.
不過我們即稱門訓.
不要搞這些.
我的腰都受不了.
要我做一條橋的話.
不過即稱又不可以這麼舒服.
我們做得門徒訓練.
你又不要想著.
上完水後.
你不要分別.
晚上回來又玩又親.
又開心又不行.
這樣都是訓練.
我們就要一點辛苦.
就好像做健身一樣.

$^{801}$去做健美團契.
以前Anders帶我們去做健美團契.
就是一起做健身.
做完健身之後.
大家有經驗的話.
怎樣做得對呢?.
就是肌肉會痛.
肌肉痛為何是對的?.
因為肌肉要撕裂.
才會成長.
那個肌肉才會操得到.
所以門徒訓練都應該是這樣.
即稱不會太辛苦.
不過亦都不可以太舒服.
有少許辛苦.
有少許痛.
這樣就對了.
但後來我見到有很多人.
他們都很忙的.
都願意擠點時間出來.
即是每個星期相聚.
一起去查經.
一起去投入不同的屬靈操練.
還有人每一次都會預備了一本簿.
遇到任何.
聽到任何有啟發的東西.
就會寫下來.
他說免得忘記了.
有些人一開始.
即是他去背金句.
好像很辛苦的樣子.
做苦力都沒有那麼辛苦.
即是逐粒逐粒字掙出來.
還要翻手很多次.
才完整背到一句.
去到最後.
他已經是背金句當食生菜.
很容易的.
見到他們那麼堅持.
有成長.

$^{841}$是令我覺得很安慰.
亦都覺得很值得.
而門訓後期最明顯的不同.
就是每一次的門訓時間.
都越來越長.
因為他們想知的東西越來越多.
很多東西想問.
他們可以互相鼓勵的東西.
都越來越深入.
看得出他們真的很想知道.
關於多些神的東西.
以致他們可以切切實實地成長.
當然不是每一間教會.
都要一樣的做法.
不一定要好像友愛堂那樣.
亦都好像剛才那樣說.
其實門徒訓練.
根本就已經散落在.
我們教會不同的環節裡面.
只要大家很用心地去參加.
就可以了.
最重要不是甚麼形式.
而是我們的心.
是不是真的很想知道.
多些關於神的東西.
這個心.
事實上.
往往很多屬靈的傳統.
都是和這個想法.
一出一轍.
一脈相承.
例如我以前經常聽到.
你不是主要學的老師.
就是主要學的學生.
正正就是這個說法.
你要麼就好好教導人.
要麼就自己好好學習.
沒有人可以卸造一個基督徒.
要不斷地成長.
不斷地追求.

$^{881}$我經常覺得.
屬靈危機.
很多時候都不是在患難.
患病.
甚至生命不如意的時候.
這些都會.
但是通常時間都很短.
很快會過的.
但是我覺得屬靈危機.
真正會出現的.
是當我們的心覺得.
我們已經聽夠的時候.
我們覺得聽來聽去.
都是那些東西.
這些東西我都懂.
甚至可能剛才你聽到.
我一開始就說.
又是講撒種比喻.
這些時候我們就要小心.
因為可能我們的屬靈危機就出現.
我們不想再追求.
我們不想再知多一點.
我們慢慢開始生命萎縮.
最後甚至死亡.
神的道是偉大.
浩瀚.
有深度的.
不同的時候.
你聽應該是有不同的想法.
不同的領受.
不存在.
我們有些時候原來已經聽夠.
而不需要再認識多一點.
我們要保持渴慕.
不斷更新.
屬靈生命.
是何時都可以成長的.
蘋果的創辦人.
喬布斯Steve Jobs.
他都有句名言.

$^{921}$Stay hungry.
Stay foolish.
他的意思就是.
保持饑渴.
保持謙卑.
這個是他對發明的態度.
但是我們每一個門徒.
面對著神的道.
我們都要有這樣的態度.
各位弟兄姊妹.
想知更多.
是我們每一個門徒.
能夠成長的關鍵.
我們要不斷追求成長.
蛻變.
我們的生命.
是可以不斷成長.
凡有耳的.
就應當聽.
我們一起祈禱.
親愛的天父上帝.
我們多謝你.
你的道是真實.
你是那位真正將那個.
有很多很深入的生命之道.
擺在我們生命裡面的那位神.
主是你的道.
是活潑.
是永恆.
無論在什麼時候.
我們都可以有領受.
主我們祈求我們每一個的耐心.
被你親自改變.
我們要真正繼續追求你.
無論我們回了多少年教會.
無論我們信了多少年耶穌.
我們最重要的是.
我們不滿足.
我們要繼續追求你.
我們要繼續認識你多些.

$^{961}$這個心正正就是你想我們想要的.
藉著撒忠的比喻.
你要提醒我們每一個.
凡有耳的我們都要聽.
就是我們每一個人都要行動.
要用心地追求你.
求上帝你幫助我們.
當我們用心追求你的時候.
求你顯明你自己的心意.
承認你在我們心裡面.
不斷地工作.
提醒.
以致我們能夠更加明白你的道.
你的道在我們生命裡面產生果效.
以致我們的生命.
就算這個世界多麼黑暗.
就算這個世界變化多麼大.
但是我們都可以.
好像一個鬧.
就像在海裡站得那麼穩.
我們在你的裡面能夠站立得穩.
能夠不斷成長.
求主幫助我們.
與我們同在.
我們多謝你聽我們祈禱.
奉主名求.
Amen.
\newpage



\section{詩篇 69:1-36}
\label{sec:PGMHgAUtHoA}
\textbf{2026年3月15日崇拜直播｜陳冠成牧師｜基督之詩篇:急難中的禱告｜詩篇69\_1-36}
\newline
\newline
連結: \href{https://youtube.com/watch?v=PGMHgAUtHoA}{\texttt{https://youtube.com/watch?v=PGMHgAUtHoA}} ~~~~ 語音日期: 2026-03-16
\newline
\newline
\hyperref[sec:yBE6W2fCqSM]{\small{< < < PREV SERMON < < <}}
~
\hyperref[sec:index_chronic]{\small{[返順時目]}}
~
\hyperref[sec:p8dj6FAIn5Y]{\small{> > > NEXT SERMON > > >}}
\newline
\newline
詩篇 69:1-36
\newline
\begin{longtable}{cl}
\hline
\hline
章節 & 經文 (和合本修訂版)\\
\hline
69:1 & \begin{tabularx}{0.7\textwidth}{X} 神啊，求你救我！因為眾水就要淹沒我。 \end{tabularx} \\ \\ \relax
69:2 & \begin{tabularx}{0.7\textwidth}{X} 我深陷在淤泥中，沒有立腳之地；我到了深水之中，波濤漫過我身。 \end{tabularx} \\ \\ \relax
69:3 & \begin{tabularx}{0.7\textwidth}{X} 我因呼求困乏，喉嚨發乾；我因等候神，眼睛失明。 \end{tabularx} \\ \\ \relax
69:4 & \begin{tabularx}{0.7\textwidth}{X} 無故恨我的，比我的頭髮還多；無理與我為仇、要把我剪除的，甚為強盛。我沒有搶奪，他們竟然要我償還！ \end{tabularx} \\ \\ \relax
69:5 & \begin{tabularx}{0.7\textwidth}{X} 神啊，我的愚昧，你原知道，我的罪愆不能向你隱瞞。 \end{tabularx} \\ \\ \relax
69:6 & \begin{tabularx}{0.7\textwidth}{X} 萬軍之主耶和華啊，求你不要讓那等候你的因我蒙羞！以色列的神啊，求你不要讓那尋求你的因我受辱！ \end{tabularx} \\ \\ \relax
69:7 & \begin{tabularx}{0.7\textwidth}{X} 因我為你的緣故受了辱罵，滿面羞愧。 \end{tabularx} \\ \\ \relax
69:8 & \begin{tabularx}{0.7\textwidth}{X} 我的兄弟把我當陌生人，我母親的兒子把我當外邦人。 \end{tabularx} \\ \\ \relax
69:9 & \begin{tabularx}{0.7\textwidth}{X} 因我為你的殿心裡焦急，如同火燒，並且辱罵你的人的辱罵都落在我身上。 \end{tabularx} \\ \\ \relax
69:10 & \begin{tabularx}{0.7\textwidth}{X} 我哭泣，以禁食刻苦我心；這倒成了我的羞辱。 \end{tabularx} \\ \\ \relax
69:11 & \begin{tabularx}{0.7\textwidth}{X} 我拿麻布當衣裳，卻成了他們的笑柄。 \end{tabularx} \\ \\ \relax
69:12 & \begin{tabularx}{0.7\textwidth}{X} 坐在城門口的談論我，酒徒也以我為歌曲。 \end{tabularx} \\ \\ \relax
69:13 & \begin{tabularx}{0.7\textwidth}{X} 至於我，耶和華啊，在悅納的時候我向你祈禱。神啊，求你按你豐盛的慈愛，憑你拯救的信實應允我！ \end{tabularx} \\ \\ \relax
69:14 & \begin{tabularx}{0.7\textwidth}{X} 求你搭救我脫離淤泥，不叫我陷在其中；求你使我脫離那些恨我的人，使我脫離深水。 \end{tabularx} \\ \\ \relax
69:15 & \begin{tabularx}{0.7\textwidth}{X} 求你不容波濤漫過我，不容深淵吞滅我，不容深坑在我以上合口。 \end{tabularx} \\ \\ \relax
69:16 & \begin{tabularx}{0.7\textwidth}{X} 耶和華啊，求你應允我！因為你的慈愛本為美好；求你按你豐盛的憐憫轉回眷顧我！ \end{tabularx} \\ \\ \relax
69:17 & \begin{tabularx}{0.7\textwidth}{X} 不要轉臉不顧你的僕人；我在急難之中，求你速速應允我！ \end{tabularx} \\ \\ \relax
69:18 & \begin{tabularx}{0.7\textwidth}{X} 求你親近我，救贖我！求你因我仇敵的緣故將我贖回！ \end{tabularx} \\ \\ \relax
69:19 & \begin{tabularx}{0.7\textwidth}{X} 你知道我所受的辱罵、欺凌、羞辱；我的敵人都在你面前。 \end{tabularx} \\ \\ \relax
69:20 & \begin{tabularx}{0.7\textwidth}{X} 辱罵刺傷我的心，使我憂愁。我指望有人體恤，卻沒有一個；指望有人安慰，卻找不著一個。 \end{tabularx} \\ \\ \relax
69:21 & \begin{tabularx}{0.7\textwidth}{X} 他們拿苦膽給我當食物；我渴了，他們拿醋給我喝。 \end{tabularx} \\ \\ \relax
69:22 & \begin{tabularx}{0.7\textwidth}{X} 願他們的筵席在他們面前變為羅網，在他們平安的時候變為圈套。 \end{tabularx} \\ \\ \relax
69:23 & \begin{tabularx}{0.7\textwidth}{X} 願他們的眼睛昏花，看不見；求你使他們的腰常常戰抖。 \end{tabularx} \\ \\ \relax
69:24 & \begin{tabularx}{0.7\textwidth}{X} 求你將你的惱恨倒在他們身上，使你的烈怒追上他們。 \end{tabularx} \\ \\ \relax
69:25 & \begin{tabularx}{0.7\textwidth}{X} 願他們的住處變為廢墟，他們的帳棚無人居住。 \end{tabularx} \\ \\ \relax
69:26 & \begin{tabularx}{0.7\textwidth}{X} 因為你所擊打的，他們就迫害；你所擊傷的，他們述說他的愁苦。 \end{tabularx} \\ \\ \relax
69:27 & \begin{tabularx}{0.7\textwidth}{X} 求你使他們罪上加罪，不容他們在你面前稱義。 \end{tabularx} \\ \\ \relax
69:28 & \begin{tabularx}{0.7\textwidth}{X} 願他們從生命冊上被塗去，不得名列在義人之中。 \end{tabularx} \\ \\ \relax
69:29 & \begin{tabularx}{0.7\textwidth}{X} 但我困苦憂傷；神啊，願你的救恩將我安置在高處。 \end{tabularx} \\ \\ \relax
69:30 & \begin{tabularx}{0.7\textwidth}{X} 我要以詩歌讚美神的名，以感謝尊他為大！ \end{tabularx} \\ \\ \relax
69:31 & \begin{tabularx}{0.7\textwidth}{X} 這就讓耶和華喜悅，勝似獻牛，獻有角有蹄的公牛。 \end{tabularx} \\ \\ \relax
69:32 & \begin{tabularx}{0.7\textwidth}{X} 謙卑的人看見了就喜樂；尋求神的人，願你們的心甦醒。 \end{tabularx} \\ \\ \relax
69:33 & \begin{tabularx}{0.7\textwidth}{X} 因為耶和華聽了窮乏的人，不藐視被囚的人。 \end{tabularx} \\ \\ \relax
69:34 & \begin{tabularx}{0.7\textwidth}{X} 願天和地、海洋和其中一切的動物都讚美他！ \end{tabularx} \\ \\ \relax
69:35 & \begin{tabularx}{0.7\textwidth}{X} 因為神要拯救錫安，建造猶大的城鎮；他的子民要在那裡居住，得地為業。 \end{tabularx} \\ \\ \relax
69:36 & \begin{tabularx}{0.7\textwidth}{X} 他僕人的後裔要承受這地，愛他名的人要住在其中。 \end{tabularx} \\ \\
[1ex]
\hline
\hline
\end{longtable}
$^{1}$以下以聆聽聖道敬拜上主.
今日經費在詩篇69篇1至21節.
在教會聖經舊約742至744頁.
請聽弟兄讀出.
神啊,求你救我,因為縱水要淹沒我.
我陷在深淤泥中,沒有立腳之地.
我到了深水中,大水漫過我身.
我因呼求困乏,喉嚨發乾.
我因等候神,眼睛失明.
無辜恨我的比我頭髮還多.
無理與我為仇,要把我剪除的甚為強盛.
我沒有搶奪的,要叫我償還.
神啊,我的愚昧你原知道.
我的罪險不能隱瞞.
萬君的主耶和華啊,求你叫那等候你的.
不要因我蒙羞.
以色列的神啊,求你叫那尋求你的.
不要因我受辱.
因為我為你的緣故受了辱罵,滿面羞愧.
我的弟兄看我為外路人.
我的同胞看我為外邦人.
因我為你的殿心裡焦急,如同火燒.
並且辱罵你人的辱罵,都落在我身上.
我哭泣,以禁食刻苦我心.
這倒算為我的羞辱.
我罵罵暴動衣裳,就成了他們的笑談.
坐在城門口的談論我.
酒徒也以我為歌曲.
但我在月立的時候,向你耶和華祈禱.
神啊,求你按你豐盛的慈愛.
憑你拯救的誠實,應運我.
求你搭救我出來於泥.
不叫我陷在其中.
求你使我脫離那些恨我的人.
使我出來深水.
求你不用帶水漫過我.
不用深淵吞滅我.
不用坑坑在我耳上合口.
耶和華,求你應運我.
因為你的慈愛本為美好.

$^{41}$求你按你豐盛的慈悲.
迴轉眷顧我.
不要掩面不顧你的僕人.
我是在急難之中.
求你速速的應運我.
求你親近我,救贖我.
求你因我的受敵,把我贖回.
你知道我受的辱罵,欺凌,羞辱.
我的敵人都在你面前.
我又辱罵傷破了我的心.
我又滿了憂愁.
我只望有人體恤,卻沒有一個.
我只望有人安慰,卻找不著一個.
他們拿苦膽給我當食物.
我喝了,他們拿醋給我喝.
《聖經》讀得.
請姐妹平安.
剛才所讀的詩篇69篇.
可以說是在《新玉聖經》裡面.
引用得最多的其中一個詩篇.
它和前一星期陳明全牧師.
所宣講的詩篇22篇.
同樣是因為預表了耶穌基督的苦難.
所以學者會給它一個名稱.
叫做彌賽亞的詩篇或者基督的詩篇.
譬如主耶穌,我們都記得.
當祂在釘十字架的時候.
曾經有人拿苦膽調和的酒給祂喝.
拿醋給祂喝,我們都記得.
其實正好應驗了剛才所讀的最後一節經文.
第21節那裡說.
他們拿苦膽給我當食物.
我喝了他們拿醋給我喝.
所以當我們去閱讀這篇詩篇的時候.
其實我們可以嘗試將耶穌基督.
祂的事蹟代入在其中.
了解到底上帝的靈.
如何透過詩篇來幫助我們.
來進到耶穌基督在苦難裡面的心路歷程.
我們開始逐個段落來看整個經文.

$^{81}$首先你會發覺第一個段落是1至12節.
你會看到,在廣道大綱亦都印了.
或者在螢光幕都看到.
詩人反反覆覆提到誰?.
提到自己,我.
用很多個我字,將焦點先在自己那裡.
向上帝訴苦.
他首先很概括形容他所遭遇的困境.
就好像整個人踩進泥沼.
找不到立腳的地方,所以怎樣?.
一路沉下去,上不來.
詩人又形容他現在的遭遇.
好像被洪水淹沒一樣.
隨時會有生命危險.
然後詩人再繼續仔細地描述.
到底他的痛苦是來自甚麼.
首先第三節說,詩人在困苦裡面.
聲嘶力竭地向神去呼求.
他形容自己喊到喉嚨都乾了.
他等候上帝的出手幫助.
但一路望,望到怎樣?.
眼睛好像發盲了,看不到.
當然這是一個誇張手法.
但要強調的是甚麼?.
他一路去求上帝幫助,但上帝不出聲.
上帝沒有反應.
而另一方面,四至六節,詩人聲稱.
自己是無緣無故被人整,被人害.
他強調自己沒有犯錯.
他說如果我真的有犯罪.
他一定瞞不過上帝.
上帝只得一清二楚.
換句話來說,詩人在這裡申訴的理據是.
上帝,你知道我是清白的.
你但是不理會我的呼求.
上帝,如果這樣會導致其他.
好像我一樣,來尋求等候你的人.
他們會誤以為.
當他們有需要上帝幫助的時候.
上帝都會怎樣?.

$^{121}$你都會好像對我一樣.
對他們不聞不問.
任由他們被人羞辱,被人欺壓.
這是第二方面的痛苦.
第三方面是七至十二節.
詩人表示他的冤屈跟誰有關?.
跟上帝有關.
他是為了上帝的緣故.
為了上帝的事情著急.
為了上帝的事情熱心.
他在做正確的事.
但反而怎樣?.
他的同胞跟他割席.
辱罵他,甚至將他成為取笑的對象.
但你見到詩人透過這短短十二節的經文.
他在說他在受三方面的痛苦.
第一是甚麼?.
不斷的呼求上帝,但上帝沒有回應.
第二方面,他是無緣無故被人憎惡.
但上帝沒有即時替他平反.
讓他繼續含冤受屈.
第三方面,他是為上帝發熱心.
但反而被自己人排擠,甚至攻擊.
再加上甚麼?.
有一樣東西叫孤單感.
他只有一個.
攻害他的人,陷害他的人有無數.
寡不敵眾.
他根本無能力擺脫眼前極大的痛苦.
事實上你會發現在詩篇裡.
很多篇章都是在極大的痛苦,屈辱.
甚至創傷的情況下寫成.
但很奇妙,經過一輪的掙扎,一輪的禱告後.
詩人通常最後又會怎樣?.
我們未讀的,稍後會一起讀.
他最後又會用讚美,稱頌上帝的說話來總結他的遭遇.
同樣今天我們看詩篇69篇的結構其實都是這樣.
在痛苦,苦難裡面禱告.
讓詩人對他所信的上帝有更深入的認識和信靠.
亦為他帶來一趟成長的旅程.

$^{161}$正如大家是否認識C.S. Lewis.
都聽過或看過他的文學作品.
甚至看過由他的文學作品改編成的一些電影.
C.S. Lewis曾經說過.
他按他的自身經歷體會到.
苦難本身其實沒有甚麼益處.
如果真的要有,其實就是迫使我們怎樣?.
你看到自己的限制,看見自己的脆弱.
從而你懂得找回上帝.
你懂得去呼求他.
老奕斯自己說,在他晚年62歲的時候.
他失去了他最心愛的妻子.
他感到很痛苦.
喪妻的痛楚一直揮之不去.
於是他不斷質問上帝.
「上帝,你是在哪裡?」.
「如果你是全能全善的」.
「為甚麼你要讓我生活得這麼痛苦?」.
老奕斯他之所以不斷追問上帝.
不是質問他,不是對上帝失去信心.
相反是他確信上帝會應門.
就好像他不斷敲上帝的門.
他相信最終有一天上帝會開門來回應他.
而最終老奕斯他說.
上帝給他的答案非常獨特.
他形容,其實他最終明白.
上帝的門根本沒有關.
上帝一直開著門.
他看到上帝仍然沉默.
但絕不是那種冷眼旁觀.
上帝好像搖頭.
但又不是拒絕他.
而是將他很在意的苦難問題放在一旁.
他形容上帝的沉默就好像對他說.
「孩子,其實有些事你現在是沒有辦法明白的」.
「但你可以相信我」.
老奕斯他終於意識到.
原來自己的痛苦受苦的過程.
為何這麼漫長,越來越艱難.
是因為他將焦點放在哪裡.

$^{201}$放在自己,不是專注於上帝.
原來當我們過度專注於自身的苦難的時候.
其實會削弱了我們承受苦難的能力.
所以老奕斯表示.
他說上帝不會製造苦難.
但上帝的美善卻會轉化苦難.
來成就他的心意.
所以我們看到苦難同樣迫使私人.
正視自己的處境和限制.
也迫使他要轉向上帝.
呼求上帝尋求他的幫助和帶領.
所以我們看到私人的禱告.
他不只是停留在訴苦的階段.
當我們再看第十三節的時候.
私人說到「但我」這個字.
有沒有留意?其實是一種轉折.
即是下一回合.
發出他第二輪的禱告.
焦點,這次我降低了給大家.
焦點是求上帝應允他.
在急難裡面的禱告.
出現了三次,在十三節到十八節裡.
大家留意十三節到十八節.
每一節這次不是說我這麼多.
是說誰?求你,求上帝你.
求上帝你,求上帝你,求上帝你.
很迫切地呼求上帝拯救他脫離苦難.
不要讓仇敵吞滅他.
之所以私人不斷呼求上帝幫助.
好像在掠上帝一樣.
不是因為祈求多點,禱告多點.
用力點,效果就會多點.
而是通過不斷呼求.
將焦點由自己轉移到誰身上.
回到上帝,回到上帝的屬性裡.
所以你看到私人的禱告.
他的呼求是有根有機.
神的信實,神的豐盛慈愛.
神的慈悲憐憫.
就是私人他在急難中持續呼求的一個確據.

$^{241}$他確信上帝始終會眷顧他.
最後無論如何都會應記.
他不會對他掩面不顧.
你再看第十九至二十一節.
私人他表示.
神啊,你知道,你知道.
你知道我有多辛苦.
我在極大的憂愁當中.
我只望身邊有一個人.
搭我肩膀,擁抱我.
說一兩句安慰體諒的話都好.
不過,現實很殘酷.
一個都沒有.
不單止這樣.
我的敵人還要落井下石.
神啊,我有多麼的孤單.
我有多麼的無助.
我有多麼的冤屈.
神都看見了.
在你眼裡發生.
但是上帝你沒有即時解除我的痛苦.
你反而容許我繼續停留在痛苦掙扎的裡面.
馬丁路德我們都認識.
一個很偉大的宗教改革家.
我們亦都知道.
他在宗教改革的事情上.
受到當時的教廷嚴厲的迫害.
大概他都好像私人一樣.
被人逼迫眾叛親離.
遭受極大的困苦和掙扎.
據說有一天.
馬丁路德的太太.
看見馬丁路德這麼愁.
於是很想鼓勵他重新振作.
他特意穿了一些當時黑色的衣服.
從頭到尾都穿黑色的衣服.
即是有一個黑沙.
即是代表甚麼.
去喪禮.
來吸引馬丁路德的注意.

$^{281}$馬丁路德一看見他.
第一句就問.
誰去世了.
誰回天家.
你要去喪禮.
他太太就回答.
就是上帝.
上帝死了.
不在了.
走了.
你知道.
督正馬丁路德的死穴.
他就馬上彈起.
很憤怒地亂說.
上帝怎會死.
上帝是永活的.
他太太就笑了.
暗笑.
上吊了.
你都知道.
上帝是永活.
沒有死.
他仍然掌管萬有.
掌管你的生命.
他只是暫時不出手.
他最終會幫你帶領到困境.
所以你現在.
用不著這麼愁.
一個很有智慧的提醒.
上帝還在.
可能你片天空倒塌了.
但是上帝掌管的大宇宙.
還在.
很穩固.
同樣斯人他明白.
上帝你知道我清白.
上帝你理解我的痛苦.
上帝你又是慈愛.
滿有豐盛.
你願意施恩.

$^{321}$我又知道你是有能力.
馬上拯救人脫離困境.
但問題是什麼.
你不出手.
我一直求你.
你也沒有反應.
你也不出手幫我.
反而是.
你任由我繼續受苦.
弟兄姊妹.
不知道你是否身處這種境況.
斯人這種情況.
不知道能給你什麼提醒.
或者上帝藉著斯人這種狀況.
想讓我們知道什麼.
想怎樣認識他.
我想有兩件事很重要.
在這篇詩篇裡.
第一.
我們要明白.
上帝要求我們承認祂的主權.
祂真的有權這樣做.
祂有權不按人的意願行事.
因為祂是上帝.
甚至祂不會問過你和我.
可能你和我都經歷過.
無怨無辜.
你有沒有做錯.
你有沒有犯罪.
突然間有病患苦難.
闖進你的安穩生活裡.
搞亂了你.
但是聖經同樣強調.
上帝不是不愛我們.
豐盛慈愛私恩憐憫.
仍然是上帝的本質.
因為祂說了是這樣.
祂是這樣的.
祂也願意按祂的時間.
出手幫助我們.

$^{361}$所以也帶出第二個提醒.
就是上帝不單只期望我們.
尊重祂的主權.
上帝也要求我們相信.
祂的出手時間才是怎樣.
祂的介入才是最完美的.
才是最好的.
確實你和我可能在人生裡.
有很多困難.
有很多甚至所謂的痛苦.
神真的沒有馬上幫你拿走.
不是好像你看醫生.
很痛的時候.
給你一粒止痛藥.
給你一粒消炎藥.
你馬上兩天就沒事.
不會這樣的.
現實很多事.
但是聖經強調這樣.
不代表我們人生玩完了.
因為為什麼.
上帝還在.
上帝沒有玩完.
我們的生命好像倒了.
但不代表上帝.
祂所掌管的東西都倒了.
我們的生命依然安穩在上帝的手裡.
祂最終會有妥善的安排.
這是這篇詩篇的上半部分.
給我們每一個很重要的提醒.
再繼續看下去.
有下半節.
我們剛才還沒有讀.
你會發現敵人的落井下石的行為.
令到詩人的傷害和創傷加劇.
很大機會因為忍無可忍.
所以他有第三輪的呼求.
不過我們看到.
詩人這次不是為自己求助.
而是怎樣.

$^{401}$而是他很希望那些苦害他的敵人.
他們有報應.
我要請鼎姐妹.
讀之前試試代入詩人的角色.
代入他的狀況.
你帶著情感來讀22至28節.
慢慢讀不要急.
預備起.
上上箭頭.
便為方牆.
願他們的帳幕無人居住.
因為你所擊打的.
他們就迫迫.
你所擊傷的.
他們棄說他的仇.
願你在他們的罪上加罪.
不容他們在你門前輕易.
願他們從生命冊上被塗抹.
不得記錄在二人之中.
大家都遷就著讀.
很溫柔.
可能大家都不習慣.
這樣對你的敵人說話.
所以大家都會很輕.
其實剛才我們所讀的.
可以說是詩篇.
一種很特殊的哀歌.
我們通常叫他做什麼.
叫做就坐詩.
很多時候是因為上帝的指紋.
他遭遇到極大的不公義.
被人欺負到無法還擊.
有冤無路訴.
受到極盡的委屈的時候.
向上帝發出一些很激烈的言語.
你看到詩人在這一節的內容.
祈求上帝懲罰惡人.
讓他們有應得的報應.
你看到詩人很多次都是.
願他們有什麼下場.

$^{441}$願上帝怎樣對付他們.
其實如果真實.
可能情緒激烈很多.
沒有我們剛才讀得那麼從容.
事實上我想讀到就坐詩的時候.
可能很多弟兄姊妹都會有一個困惑.
羅馬書在新約裡面.
特別十二章.
保羅不是明確說到.
逼迫我們那些要怎樣.
要為他祝福.
還要多說一句.
怕死我們不知道.
只要祝福不可什麼.
是的,你記得了.
只要祝福不可就坐.
是一個硬命令.
你不用查經都聽得明白.
每個字都做不做到是另一回事.
問題來了.
就坐詩是不是和保羅的教導.
和新約聖經的教導有衝突有矛盾.
確實我們很明白.
新約都強調.
耶穌都強調.
舌頭是很麻煩的東西.
我們要好好控制管理.
不能夠因為舊約聖經.
就坐詩用它來做藉口.
聖經都這樣.
就坐人我都這樣去就坐人.
不能夠因為這樣的緣故.
用惡毒的言語去直接.
或者你在網上留下一些.
惡毒的言語來辱罵別人.
不可以這樣.
這個我們很清楚.
所以你留意到.
就坐詩人在就坐的內容.
是對誰說的.

$^{481}$對那個仇人說.
還是對上帝說.
是對上帝說.
不是面對面.
去到那個仇人面前就罵他.
是自己和上帝的對話.
或者說得實際一點.
就是他在他的禱告裡面.
在上帝面前.
將傷害他的人.
透過就坐詩.
拉在一起.
擺在上帝的面前.
因為詩人的內心實在太痛苦.
沒有辦法去釋懷.
你想想那個情況.
不斷被人欺負.
欺負到心口.
他渴望的公義又怎樣.
一直都沒有出現.
在沒有出路的情況下.
他唯有能夠做的就是.
我回到上帝的面前.
就坐我的敵人.
我將審判懲罰的權利.
交回給誰.
是交回給上帝.
因為詩人知道他沒有權.
越過上帝自己去報復.
他不可以這樣做.
因為他確信唯有上帝.
才有權去懲治審判惡人.
所以詩人這裡通過就坐他的仇敵.
他是期望上帝最終會為他主持公道.
因為詩人確信上帝是公義.
公義的意思就是.
上帝不會誤判.
不會錯判.
換句話來說.
其實詩篇這裡的就坐部分.

$^{521}$其實是幫詩人排毒.
排去他內心那些屈住的苦毒.
讓自己的生命可以變為輕散.
可以稍稍得到釋放和安息.
不知道大家.
先問問大家.
有沒有試過去刮痧.
或者拔罐.
有吧.
痛不痛.
那個過程.
很痛的.
拔罐或者酒罐那些.
很痛的.
刮痧就更痛.
我也很怕.
但是刮痧我都冒冷汗.
特別是我最近的手.
應該嚴重勞損發炎.
手肘這裡.
又去試試推拿.
又去拔罐.
都好像搞不定.
還要吃消炎藥.
頂住.
因為我知道裡面應該有些東西塞住了.
瘀塞.
所以希望透過不同方法.
將裡面的瘀血清走.
因為中醫都說.
痛則如何.
痛就不痛.
痛的時候就如何.
就不痛了.
同樣其實.
咒坐師的功效就是這樣.
就好像你去刮痧.
你去拔罐.
就將私人體內那些.
深藏.

$^{561}$深入在皮膚.
深入地方那些憤恨.
那些怨恨.
刮下刮下.
拔下就如何.
由深層.
拔出來表面.
然後慢慢排走.
排走的時候就會沒那麼痛.
沒那麼辛苦.
又或者.
私人其實透過咒坐師.
就好像我們經常說.
情緒倒垃圾一樣.
但是他是來到上帝的面前.
將這些苦毒.
一一的傾倒.
不容許這些毒素繼續.
存留在心裡面.
蠶食自己的生命.
因為這樣.
不單單破壞他和人的關係.
最後這些苦毒都會影響.
他和上帝的親密.
所以其實.
師人他通過這段咒坐的禱告.
其實是向我們示範.
如何去處理.
我們在內心那些冤屈.
那些苦澀.
就是你坦誠來到上帝面前.
你承認.
有時你真的會被人氣死.
害得很慘.
你是有一些負面情緒.
甚至想法.
師人叫我們.
不是要否定他.
不是要壓制他.
不是當沒事發生.

$^{601}$我OK.
不是要扮強.
因為其實無論我們.
有沒有將心裡面那些怨恨.
說出口.
上帝知不知道?.
其實上帝知道一清二楚.
但是當我們願意.
向上帝敞開.
說出來的時候.
其實就是在幫我們自己.
提醒我們自己.
上帝是信得過的.
無論我們變成怎樣.
上帝仍然可以親近.
上帝是會接受那個軟弱的我.
那個不堪的我.
上帝是會憐憫.
我真的有一些幽暗面.
不能夠對別人說.
但是可以對上帝說.
上帝是願意對我施恩.
他最終會給我有出路.
這就是師人在上帝面前.
發出咒詛言語.
最終希望得到的效果.
我曾經記得.
有一個姐妹.
因為婚姻很觸礁.
很痛苦很難過.
於是我們去傾談的時候.
她一路去訴苦.
說到哭哭啼啼的時候.
她說我不想撐下去.
我想上帝拿走我.
拿走我是什麼意思?.
想死是嗎?.
我們一聽到就想死.
當然我確認她不是真的想死.
我跟她都熟悉.

$^{641}$而是怎樣?.
而是裡面實在太痛.
沒辦法解脫.
所以才希望這樣有出路.
所以我沒有立刻阻止她.
你不要這樣想.
不合上帝心意.
這個其實我們每一個都明白.
反而像這套詩篇一樣.
我鼓勵你.
這些都是你很真實的狀況.
你跟上帝說.
因為上帝是做你的.
創造你的.
上帝是愛你的.
祂會拯救你.
你就算有多負面的想法情緒.
上帝都願意聽.
上帝都承載到.
並且祂會幫助我們面對前路.
所以我引導她在困境裡面.
重新發現上帝.
上帝還在.
雖然可能婚姻不如意.
但是那位設立婚姻的上帝.
祂還在.
我先幫助這位弟兄姊妹.
將她的負面情緒.
安穩在上帝的裡面.
來提升她的靈性.
提升她的忍耐力.
那種的韌力.
然後才慢慢去用時間.
去實際處理.
或者說修補她的夫婦關係.
同樣弟兄姊妹.
私人你會發現.
她咒詛而仇敵之後.
正如她自己所說.
問題有沒有馬上解決.

$^{681}$沒有吧.
你看到二十九節那裡.
但我是困乎憂傷.
沒有什麼實際的功效.
似乎.
但是你會見到.
她的靈性是提升了.
當她正確在上帝面前.
宣洩她內心的冤屈之後.
她最後反而可以真誠的敬拜上帝.
排毒之後.
跟上帝的關係反而近了.
正面了.
於是她發出最後一輪禱告.
請大家又再讀.
第二十九至三十六節.
讀完最後這幾節經文.
預備起.
我要以詩歌讚美神的名.
這便叫耶和華喜悅.
聖恥獻牛.
或是獻有國有題的公牛.
見彼的人看見了就喜樂.
尋求神的人.
願你們的心甦醒.
因為耶和華聽了.
同佛人不藐視被囚的人.
願天和地.
海洋和其中一切的動物都讚美他.
因為神要拯救錫安.
建造猶大的城邑.
他的民要在那裡居住.
得以為宴.
他的人的後裔要承受為宴.
愛他命的人也要住在其中.
你發現很奇妙.
詩人在禱告的結尾部分.
其實他轉念.
雖然他仍然憂傷困苦.
但同時他又可以向上帝.

$^{721}$發出稱讚和歌頌.
表面上好像很分裂.
很為和.
但其實不難理解.
你和我都做得到.
就是當我們在苦難裡.
我們不再選擇高舉自己.
我有多慘,多辛苦,多可憐.
不是純粹將焦點放在自己無限放大.
而是你開始回到真理.
你選擇高舉上帝的美善.
上帝的信實.
上帝的拯救和同在的時候.
你和我都會像詩人一樣.
靈性得到提升.
既是一面受苦.
但同時又可以一面稱讚美上帝.
確實我們很明白詩人.
他身處的那片小天地.
真的好像倒塌了.
讓他的生存空間不斷壓縮.
更難受的是上帝彷似對他掩面不顧.
但詩人的決心.
我仍然是投靠上帝.
這是我唯一的出路.
唯一的指望.
因為他對上帝的救恩.
對上帝的信實仍然很有信心.
他對未來仍然有期盼.
我們看到他在最後那段的禱文裡.
他明白人的安危和福祉.
其實不是建基在人間的力量裡.
而是建基在永生上帝的權能和應許當中.
所以等於姐妹.
你見到詩人在掙扎的過程裡.
他透過禱告.
他的生命.
他的思想.
其實是不斷轉化.
由起初我們記得他是怎樣.

$^{761}$他來到神面前申訴他的苦情.
然後祈求上帝出手幫助.
在最後真的頂不住的時候.
他走助他的敵人.
求上帝主持公義.
當這一切完了之後.
最後的結局是.
他越來越接近上帝.
他越來越向上帝闖開.
最後發出讚美和稱謝.
他是用上帝賜給他的永恆盼望.
用一個英文字cover.
即是遮蓋了他當下短暫的困難.
用永恆遮蓋短暫.
他是用上帝的完備拯救計劃.
來cover自己有限的生命.
就是這樣一步一步經歷到.
上帝很真實的同在和心靈的醫治.
其實當中給了我們幾個很重要的屬靈教訓.
首先我們明白.
斯人不是特別喜歡面對苦難.
他和我們一樣都很怕遇到苦難.
但是他卻告訴我們.
如何在苦難中分別出來.
和世人有分別.
如何超越苦難的本身.
維持在苦難背後那位至聖至善的上帝.
至聖至善的上帝.
將苦難轉化為一個親近上帝.
一種屬靈成長的契機.
所以你看到其實斯人.
他不是為了我們現在所經歷的苦難.
提供一個標準的答案.
即時拿開你所有的問題.
而是邀請正在經歷苦難的人.
你回到上帝那裡.
你和上帝同行.
在苦惱的裡面.
你和上帝一起去探索.
當中上帝給你甚麼屬靈的意義.

$^{801}$你見到斯人之所以沒有被苦難打垮.
是因為他懂得運用這一切的屬靈資源.
他將苦難轉化為一趟屬靈的旅程.
他透過在上帝面前很坦誠地禱告.
去保持和上帝連繫.
他去傾訴他自己想法感受的同時.
他亦都願意順靠上帝對他的應許.
順從上帝在當下.
在他生命裡面一些.
他覺得很不舒服的安排.
就是這樣很奇妙.
這個盼望就是這樣來的.
所以這段經文提醒我們.
面對苦難的時候.
除了我們要關注自己的情緒感受.
我們的情緒心靈健康的時候.
他更加提醒我們.
你一定要抓緊那個刺人心靈的柱.
不要本末倒置.
我們一定需要上帝的介入.
需要真理重新調教我們.
應對苦難的思維和態度.
就正如有一位很出名的清教徒的領袖.
叫做巴赫斯特 Richard Webster.
在苦難裡面他亦都誤出一些道理.
他表示聖經強調的順服忍耐.
並不是要人去抹殺裡面內心的感受.
因為這些都是上帝賜給人一些美好的情感.
是對人有益處.
特別當面對逆境的時候.
身體的心靈反應其實就是代表一種訊號.
而這個訊號是提醒我們.
你要回歸上帝.
你要去呼求他.
就好像詩篇一樣.
所以我們作為上帝而來的好管家.
我們都有責任去好好管理自己的心靈.
包括管理我們可能那些憤怒,哀傷,恐懼或者怨毒的情緒.
以免造成失控,失去秩序.
影響我們和上帝的關係.

$^{841}$產生一些不合上帝心意的行為.
並且在苦難裡面.
我們需要認定上帝的主權,上帝的公義.
上帝的智慧和慈愛.
仍然是有效.
有一些東西我們是不明白.
上帝亦都不給答案.
甚至乎不即時解決.
但是我們仍然相信.
我們仍然在祂的拯救,微善旨意裡面.
所以你見到.
私人的禱告其實是一種應對苦難人生的智慧.
他勉勵我們.
與其你追求一個沒有痛苦的人生.
倒不如怎樣?.
你回歸上帝裡面.
你去培養足夠承載苦難的能力.
畢竟你和我都會很明白.
在現實裡面.
苦難是會以不同的時間.
在不同的時間.
用不同的方式.
未經你批准.
就是怎樣?.
進入你的生命裡面.
或者進入你身邊的人的生命裡面.
打亂你原本很穩定的生活秩序.
所以那些懂得依靠耶穌跨越面前苦難的人.
他們反而明白.
真正的智慧是在於.
不是過一個沒有痛苦的生活.
而是明白痛苦.
本來就是我們整全生命不缺少的一部份.
但是上帝會在當中給我們出路.
事實上.
如果我們還記得.
耶穌在踏上苦難之前.
祂和門徒說.
很真實地說.
在世上我們會怎樣?.

$^{881}$有苦難的.
耶穌沒有違反商品說明條例.
真的會有.
沒有苦難才是假象.
耶穌說了是會有.
但是下一句很重要.
我們可以放心.
因為苦難不是我們的結果.
耶穌勝過世界.
才是我們終極的盼望.
提醒我們.
特別在禦苦期的這段時間.
我們不是單單紀念耶穌.
怎樣為人犧牲寫命.
我們也要同時深思.
作為耶穌的門徒.
我們怎樣可以和受苦的主同行.
無論是個人或群體的苦難.
我們都要相信.
是通過上帝才能臨到.
出現在我們身上.
在裡面.
我們怎樣可以效法耶穌.
展現出對上帝那種信靠,信服.
怎樣讓我們的生命.
更加貼近耶穌.
邁向那種豐盛圓滿.
事實上.
弟兄姊妹耶穌已經透過十字架的苦難.
告訴我們.
甚麼才是整全的人生.
不是一個沒有掙扎.
沒有苦楚的人生.
而是一個更有從神而來.
韌力可以應對不同生活光景.
包括苦難的人生.
所以假如我們真的在面對苦難.
求主席今天這段經文幫助我們.
轉變思維和態度.
不要總是想.

$^{921}$沒有辦法了.
我甚麼都做不到.
而是反而專注.
我仍然在這個時況.
能夠為主做些甚麼.
亦謹記苦難確實.
有它一定程度的破壞力,影響力.
但不會大於上帝的權能.
苦難無法攔阻我們走近上帝.
服侍祂,討祂的喜悅.
而另一方面.
在苦難裡.
你和我亦需要邀請弟兄姊妹去同行.
為你代禱守望.
這樣我們才更有韌力和耐性.
應對苦難所帶來的衝擊和挑戰.
又或者假如上帝是使用你.
你去陪伴身邊只經歷苦困的人的時候.
你亦不要忘記.
你要時刻保持和上帝禱告.
將你和對方一同帶到神的面前.
尋求醫治和安穩.
因為很可能不單單是對方.
作為同行者的你.
亦會被對方的困難,痛苦所牽動.
你都會受傷,你都會難過.
你都會和對方陷入苦困裡.
甚至出不來.
在這個時候.
我們更需要像私人一樣.
透過禱告.
讓你和對方都進入上帝避難所裡.
一起去經歷.
祂給予我們靈性的安穩和成長.
而最後.
這篇詩篇提醒我們.
永遠都要認定.
上帝是永遠同在和掌管.
你都要確信凡事.
上帝一定會有最好的安排.

$^{961}$只不過不一定按我們時間發生.
並且耶穌亦都應許.
有一天祂會來接我們.
回到永恆福樂的天間裡.
所以我們可以帶著這一切的應許和期盼.
好好為主而活.
好好過我們的每一天.
我們同心禱告.
因主願你藉著這篇詩篇.
讓我們走近耶穌你.
幫助我們特別在遇苦期裡.
我們能夠與你這位受苦的主同行.
作為你的門徒.
在人世間我們同樣.
因為種種的緣故.
我們經歷不同程度的傷痛.
不同程度的困難.
我們亦都會像詩人一樣.
在裡面有很多疑惑.
很多掙扎.
甚至紛紛不平.
我們亦都可能像詩人一樣.
一路呼求.
一路希望上帝來幫助.
但仍然在等候當中.
求主你藉著.
你自己的話語.
堅固我們對你的信靠.
讓我們透過禱告.
在苦難當中與你結連.
讓我們總是仰望你.
等候你.
亦都幫助我們.
在當中找到同行者.
來為我們守望.
讓我們有更多.
從上帝你而來那份的忍耐.
來在苦難當中榮耀你.
討你的喜悅.
亦都盼望主你幫助我們.

$^{1001}$可能我們身邊.
我們所照顧的親友家人.
同樣在艱難的當中.
我們亦都會覺得很乏力.
很無助.
甚至我們亦都受傷.
盼望主你同樣的.
一一去醫治去幫助我們.
讓我們能夠有.
因著你的緣故.
更多承載別人生命的能力.
讓我們亦都在過程當中.
我們有恩典.
我們能夠讓自己和對方.
慢慢靠著主你.
將苦楚一路一路轉化為.
對你的讚美和稱頌.
為此我求主你.
與我們每一個同在.
與我們所關愛的人同在.
奉基督耶穌你得勝的名字.
你亦都可以祈求.
Amen.
\newpage



\section{約書亞記 1:1-9}
\label{sec:p8dj6FAIn5Y}
\textbf{2026年3月22日崇拜直播｜蔡康怡牧師｜超越功能價值 立足聖約人生｜約書亞記1\_1-9}
\newline
\newline
連結: \href{https://youtube.com/watch?v=p8dj6FAIn5Y}{\texttt{https://youtube.com/watch?v=p8dj6FAIn5Y}} ~~~~ 語音日期: 2026-03-24
\newline
\newline
\hyperref[sec:PGMHgAUtHoA]{\small{< < < PREV SERMON < < <}}
~
\hyperref[sec:index_chronic]{\small{[返順時目]}}
~
\hyperref[sec:3zRaVYVNQY0]{\small{> > > NEXT SERMON > > >}}
\newline
\newline
約書亞記 1:1-9
\newline
\begin{longtable}{cl}
\hline
\hline
章節 & 經文 (和合本修訂版)\\
\hline
1:1 & \begin{tabularx}{0.7\textwidth}{X} 耶和華的僕人摩西死了以後，耶和華對摩西的助手嫩的兒子約書亞說： \end{tabularx} \\ \\ \relax
1:2 & \begin{tabularx}{0.7\textwidth}{X} 「我的僕人摩西死了。現在你要起來，和眾百姓過這約旦河，往我所要賜給以色列人的地去。 \end{tabularx} \\ \\ \relax
1:3 & \begin{tabularx}{0.7\textwidth}{X} 凡你們腳掌所踏之地，我都照我所應許摩西的話賜給你們了。 \end{tabularx} \\ \\ \relax
1:4 & \begin{tabularx}{0.7\textwidth}{X} 從曠野和這黎巴嫩，直到大河，就是幼發拉底河，赫人的全地，又到大海日落的方向，都要作你們的疆土。 \end{tabularx} \\ \\ \relax
1:5 & \begin{tabularx}{0.7\textwidth}{X} 你一生的日子，必無人能在你面前站立得住。我怎樣與摩西同在，也必照樣與你同在；我必不撇下你，也不丟棄你。 \end{tabularx} \\ \\ \relax
1:6 & \begin{tabularx}{0.7\textwidth}{X} 你當剛強壯膽，因為你必使這百姓承受那地為業，就是我向他們列祖起誓要給他們的地。 \end{tabularx} \\ \\ \relax
1:7 & \begin{tabularx}{0.7\textwidth}{X} 只要剛強，大大壯膽，謹守遵行我僕人摩西所吩咐你的一切律法，不可偏離左右，使你無論往哪裡去，都可以順利。 \end{tabularx} \\ \\ \relax
1:8 & \begin{tabularx}{0.7\textwidth}{X} 這律法書不可離開你的口，總要晝夜思想，好使你謹守遵行這書上所寫的一切話。如此，你的道路就可以亨通，凡事順利。 \end{tabularx} \\ \\ \relax
1:9 & \begin{tabularx}{0.7\textwidth}{X} 我豈沒有吩咐你嗎？你當剛強壯膽，不要懼怕，也不要驚惶，因為你無論往哪裡去，耶和華你的神必與你同在。」 \end{tabularx} \\ \\
[1ex]
\hline
\hline
\end{longtable}
$^{1}$《約書亞記》一章一至九節.
我踏之地 我都照著我所應許摩西的話賜給你們了.
從曠野和者黎巴嫩 直到帕拉大河 客人的傳遞.
又到大海日落之處 都要作你們的警戒.
你平生的日子 必無一人能在你面前站立得住.
我怎樣與摩西同在 也必照樣與你同在.
我必不撇下你 也不丟棄你.
你當剛強壯膽 因為你必使這百姓承受那地為業.
就是我向他們列祖起誓 應許賜給他們的地.
只要剛強大大壯膽 謹守遵行我僕人摩西所吩咐你的一切律法.
不可偏離左右 使你無論往哪裡去 都可以順利.
這律法書不可離開你的口 總要晝夜思想.
好使你謹守遵行這書上所寫的一切話.
如此你的道路就可以亨通 凡事順利.
我豈沒有吩咐你嗎 你當剛強壯膽.
不要懼怕 也不要驚慌 因為你無論往哪裡去.
耶和華你的神必與你同在 聖經六品.
今天很高興邀請到蔡康儀牧師與我們分享訊息.
蔡牧師是香港亞洲歸屬協會的總幹事.
今天與我們分享的題目是超越功能價值 立足聖若人生.
我們請蔡牧師.
請姐妹平安.
很感恩這次是第一次在大家當中分享神的話語.
雖然是第一次分享上帝的說話.
但是旺角浸信會給我的印象都已經很深刻.
因為無論我認識的牧者.
或者是我去認識大家的事工.
我發現這間教會都是一間心胸很廣闊.
很願意去聆聽這個世代的需要.
並且是一間很有活力的教會.
所以每一次當我想起旺角浸信會的時候.
我感受到為大家很感恩.
首先介紹一下我們機構.
我在香港亞洲歸屬協會服務.
我們機構是一間大使命的教會.
我們知道現在我們面對種種的困難和挑戰.
其實我們更加知道需要實踐大使命.
但是我們又知道這麼艱難的時代.
實踐大使命一點都不容易.
我們機構一直都很著力去提倡什麼呢?.

$^{41}$大使命不是你生活上的其中一部分.
我們知道生活有家庭,職業,教會各方面.
當順境的時候我們可能每一項都兼顧得很好.
但是一到很忙碌,一到艱難的時間.
很多時候我們都會犧牲我們娛樂消遣的時間.
下一步我們都知道大使命很重要.
但是卻是有時有心無力.
我們機構一直都很想去實踐一個藍圖.
就是從左邊這個圖變成右邊.
就是將大使命融入在你的生活每一個部分.
所以你不是多做一件事.
而是你做每一件事的時候.
其實你都是指向大使命.
但是的確今時今日面對的挑戰越來越大.
最近我們機構有一個市集.
所以我們因為很想將生活和信仰融入.
所以在大家不是很遠的地方.
在何文田我們有一個好感生活體驗館.
我們就是這樣實踐這個理念.
最近就有一個市集.
市集完有些電影節目幫我們拍完照片回來之後.
發一些照片給我們.
就是這些照片來的.
就不是實體照片.
當我看著這些照片的時候.
我都不得不慨嘆今時今日做人真的非常艱難.
特別是下一代我們現在面對的挑戰是什麼呢.
人不單止要和人去競爭.
現在還要和AI去競爭.
所以究竟在這麼艱難這麼多挑戰的世代裡.
我們怎樣可以站立得穩.
去持守著大使命一直向前行.
今天我就很想和大家分享.
我們看看這段經文我相信大家都不陌生.
一開始的時候就說到.
耶和華的僕人摩西死了以後.
一代偉大的領袖他離開這個世界之後.
我們知道神就差派了亂的兒子約書亞.
來接續摩西的工作.
一代過去代表一代興起.

$^{81}$代表著一個時代的轉變.
同樣地當中我們都不難去思想到.
其實當中時代的更替的轉變有多大.
對於約書亞來說她面對的挑戰是什麼呢.
我相信大家都很熟悉聖經.
你看看整本聖經描述約書亞的地方多不多.
其實都不是很多.
所以我們都不太知道她是一個怎樣的人.
但描述摩西的地方多不多呢.
就非常之多了.
我們知道摩西的出身.
她有給上帝很好的塗灶.
所以面對著一個這麼偉大的領袖的接班.
我相信她的壓力一點都不小.
我自己就很明白.
因為我在亞洲歌唱會事奉了六年的時間.
但我對上一任的總歌手事奉了多久呢.
事奉了47年.
我就很知道約書亞所面對的壓力有多大.
接著不單止面對著自己接班的壓力.
約書亞還面對著民族方面的差異.
我們知道摩西帶領著一班以色列人.
都是經歷曠野.
雖然都艱難.
但都是單日都是以色列的民族.
但當約書亞她要帶領以色列人進入迦南的時候.
就變了多元的民族.
並且第一代和第二代的差異亦都很明顯.
我們都可想而知約書亞實在面對的挑戰不小.
不單止民族的挑戰.
環境,宗教的都變得多元.
都變得侵襲的時候.
約書亞要帶領著一班以色列人.
在這樣的環境下來履行上帝的使命.
一點都不容易.
正如今時今日剛才說過.
我們現在面對AI很大的挑戰.
在這樣的情況下.
於是我就怎樣呢.
既然AI好像很厲害.

$^{121}$有一天我就問AI.
AI你這麼厲害.
究竟有什麼你是做不到的.
為什麼我會問這個問題呢.
當我們問到AI它做不到的東西.
我們做得到.
我們就不會被AI取代.
一問之下.
完完全全AI就告訴我們.
它有四點.
主要有四點它是做不到的.
它說什麼呢.
第一就是創造力和藝術的表達.
它說AI生成的藝術,音樂或者文學作品.
依賴於已有的數據和模式.
缺乏真正原創的情感驅動和內在靈感.
認不認同.
大家都認同.
接著我們又再看.
情感的支援.
AI又不是被人做得很到位.
它說AI雖然可以模擬情感和提供支持.
但人類在面對複雜的情感交流的時候.
擁有的同理心和直覺.
這是AI沒有辦法能夠達到的心動.
大家又認不認同呢.
都認同.
接著第三.
道德倫理價值的判斷.
AI的抉擇基於算法和預設的規則.
缺乏主觀道德的意識和價值觀.
人類可以根據文化,哲學和道德框架.
對複雜的倫理問題進行深刻的思考.
而這不是單純的邏輯運算所能夠實現的.
所以你問AI.
它沒有什麼判斷.
什麼都認同你.
接著到第四.
創新和問題解決.
AI擅長解決已有數據中的模式問題.

$^{161}$但對於全新的無法依據過去經驗解決的問題.
仍然需要人類創新的思維.
人類能夠從零開始設計和構想解決方案.
而AI通常是依靠已有的數據和框架.
不知道這一點大家又同不同意.
如果以上四點大家都同意的話.
換言之在今時今日.
AI都逐漸取代了人類.
最近我有個弟兄告訴我.
他在銀行工作.
他說銀行現在都越來越大挑戰.
要裁員.
裁員的時候.
他怎樣將員工分成等級.
分成三級.
怎樣分呢.
第一級就叫做Better than AI.
第二級就叫做Equals to AI.
即是和AI同等.
第三級就叫做Poorer than AI.
比AI更加差.
銀行不夠錢.
第一批裁員是甚麼人呢.
當然是裁員Poorer than AI.
第二批是甚麼呢.
就是Equals to AI.
所以換句話來說.
可以留下的就是Better than AI.
但要Better than AI容易嗎.
一點都不容易.
所以剛才那四點就給我們一個很具體的方向.
究竟我們怎樣成為一個Better than AI的人呢.
但是究竟人還做得到嗎.
最近我又看了一篇文章.
他說老師還未被AI取代.
但你們已經將他變得可以被取代.
我看完這篇文章之後.
網上的文章.
我覺得很發人深省.
不單做老師.

$^{201}$做人的工作.
其實我們都要聽聽.
我讀一讀給大家聽.
讀一部分.
老師就問AI.
你覺得自己可不可以取代老師呢.
AI就回答.
看你怎樣定義老師.
於是老師就皺眉.
你覺得現在的社會是怎樣定義老師的呢.
AI就說.
一個負責知識傳遞.
內容講解.
評量出題.
成績輸出的人.
停了一停之後.
AI就繼續說.
但這些功能我都能夠承認.
並且全年無休.
不抱怨成本低.
你不用給幾萬元請老師.
你便宜點請AI都可以.
但AI就沉默了一下.
就說.
但老師不單是這樣.
AI就停頓了一會.
語氣冷靜但精準.
他說我知道老師也曾經是情感的支援者.
價值的引導者.
以及信念的建構者.
但你們已經很少讓他們有時間做這些事.
老師就問.
你這樣說是甚麼意思.
AI就回答.
當老師花更多時間在填寫資料.
處理行政.
應對考題.
配合教學的績效.
而不是和學生建立真實的連結.
他已經慢慢接近我AI.

$^{241}$老師繼續回應.
老師和你還是不同的.
他們會傾聽.
會等待孩子成長.
會容許他們卡住.
AI就回答.
是呀但你們的制度.
不獎勵等待.
不容許卡住.
不承認那些看不見的價值.
久而久之.
他們只學會快速輸出.
標準流程.
講重點.
看分數.
試問你們還需不需要一個會等待的人呢.
還是只需要我這些快速的AI就搞定呢.
老師聽到這裡好像沒有辦法反駁.
於是老師就說.
你認為未來你會不會完全取代老師呢.
AI就回答.
如果你繼續讓老師只剩下我能夠做的事.
答案是肯定的.
他說我不是要來打敗老師.
是你們先將老師變成可以讓我打敗.
這樣的樣子.
弟兄姊妹我不知道你們聽完之後.
有什麼感受.
我看完這段文章.
我的感受很深.
的確我不是做老師.
但我是做牧師.
的確有很多行政.
有很多各方面.
事務上的處理.
都叫我們忙得不可開交.
慢慢我們人就變成一個機械人.
但是我們又不讓AI精準.
於是我們就被淘汰.
所以今時今日.

$^{281}$我們似乎沒有地方站了.
究竟我們站在一個什麼的立足點上.
很感恩.
透過上帝的說話.
我們看一看.
這段經文我們如果問大家.
這段經文的主題是什麼呢.
我想大家一定會.
首先想起四個字.
就是剛強壯膽.
因為這段經文不斷出現了很多次.
但是你要看到.
每一次神叫若書雅剛強壯膽.
祂不是好像啦啦隊式那樣.
叫大家你行的.
加油啊.
你做得到的.
不是這樣的.
而是每一次.
神說你剛強壯膽.
背後都有一個很清晰的原因.
譬如六節.
祂說你當剛強壯膽.
因為....
第七節.
只要剛強大大壯膽.
....
使你無論往那裡去.
都可以順利.
第九節.
祂說我喜沒有吩咐你嗎.
你當剛強壯膽.
不要懼怕.
也不要驚慌.
因為....
所以換句話來說.
神不是空口說白話.
不是只是為若書雅打打氣那麼簡單.
而是祂很清楚告訴你.
因為....

$^{321}$所以你就可以剛強壯膽.
因為什麼呢.
神很清楚.
和若書雅重申.
因為我和你立的約.
所以我們抓住上帝的約.
就可以立足在一個不敗之地.
就可以剛強壯膽.
我們一起看看.
這個約背後.
告訴我們什麼.
以致我們可以抓得緊.
我們可以站得穩.
我們一起看看.
第一.
我們看到.
上帝是叫若書雅立足應許.
以致建立信念.
成為一個信念的建構者.
第三節.
祂說凡你們腳掌所踏之地.
我都照著我所應許摩西的話.
賜給你們了.
第六節說.
你當剛強壯膽.
因為你必使這百姓承受那地為業.
就是我向他們列祖起誓.
應許賜給他們的地.
當若書雅很害怕.
背負著很大的挑戰的時候.
神跟她說什麼.
我不是告訴過你嗎.
我應許過摩西的話.
我現在賜給你.
我應許過以色列人.
你的列祖要賜給你.
我就會賜給你.
所以祂叫若書雅立足什麼.
是立足著上帝的應許.
而這個應許是祂的列祖.

$^{361}$早已經知道得清清楚楚.
有一本書叫All the Promises in the Bible.
作者告訴我們.
究竟聖經裡面有多少個應許呢.
原來算著算著.
有超過八千個應許.
好了.
可能你會說.
八千個應許很久都記不清.
這八千個應許不是說.
八千個獨立不同的應許.
否則你這樣記就很長時間記.
而是重重複複.
不斷重複上帝跟人立約.
所帶給人的應許.
為什麼聖經的作者.
不斷重複這麼冤悶.
重重複複這樣呢.
如果我們這樣看.
我們不明白上帝的苦心.
我們回想在艱難的日子裡.
其實很多時候我們做決定.
是憑著什麼去做決定呢.
你想想你做決定.
越來越艱難.
你會憑什麼去做決定呢.
你會不會憑著你過去的經驗.
來做決定呢.
為什麼呢.
因為這些經驗很穩定.
你抓緊過去做過的.
你現在再做.
現在這麼艱難.
不要這麼冒險.
不如穩穩陣陣就算了.
是不是這樣呢.
所以當艱難的日子.
我們的決定就越來越穩定.
但是我想說.
我們的經驗有多少呢.

$^{401}$我們的經驗大不大.
AI那些大數據.
一定大不過.
所以當我們只是抓緊.
我們這些狹小的經驗.
其實我們就穩穩陣陣.
一切都不出所料.
但是上帝跟我們說.
我們可以不用抓緊.
我們的經驗這麼狹窄.
我們只要抓緊上帝的應許.
應許是什麼呢.
是那些還沒發生的事.
但是上帝答應了的事.
你抓得緊的時候.
你就方言去看看.
神.
祂怎樣令你.
出人意外的結果.
所以當我們抓緊上帝的應許.
我們就明白.
神知道我們都要有些渣拿.
我們人的經驗很有限.
我們要等待上帝的應許.
都還沒見到.
心裡都很怕.
所以祂就在聖經裡面.
前人呢.
神怎樣跟他立約.
怎樣答應過他什麼.
在聖經裡面記載得清清楚楚.
重複又重複.
當我們在聖經裡面.
一直看著神這些不變的應許.
在不同人身上.
都來到淺行的時候.
這個就成為比起AI的大數據.
更加厲害的大數據.
我們抓緊這些應許的時候.
這些應許就成為我們堅穩的信念.

$^{441}$無論是多少的未知.
我們抓緊這些信念.
就成為了一個.
真真實實信念的建構者.
今時今日.
我們不是被逆境打敗.
是被我們薄弱的信念打敗.
我們左搖右擺.
所以搞得我們很迷失.
搞得我們不知道應該怎樣做.
所以是時候我們回去.
抓緊神的應許.
建立穩固的信念.
以致我們可以剛強壯膽.
面對前面的道路.
第二點.
接著我們再看.
在這個約裡面.
不單單是應許這麼簡單.
第五節告訴我們.
神說:你平生的日子.
必無一人能在你面前站立得住.
我怎樣與摩西同在.
也必照樣與你同在.
我必不撇下你.
也不丟棄你.
神和人的立約不是很冷冰冰.
簽張合同.
你做你的事.
我履行我的責任.
這麼冷冰冰.
其實約裡最重要最重要的精神是什麼.
是神要和我們建立那個關係.
所以神就提醒若舒雅.
當上帝怎樣與摩西同在.
聖經告訴我們.
上帝怎樣與摩西同在.
摩西可以面對面和神聊天.
好像朋友一樣.
所以神怎樣與摩西同在.

$^{481}$神也告訴若舒雅.
我都這樣與你同在.
我必不撇下你.
也不丟棄你.
當我們真真實實的.
與上帝建立一個這樣的關係.
不離不棄的關係.
我們還有什麼需要害怕呢.
今時今日.
我們的心臟.
弱一點也不行.
每天都會面對很多.
突如其來叫我們很擔心的消息.
打仗.
前一陣子.
上年年尾的時候.
發生了雲福縣的大火.
當我們看到這些消息的時候.
其實我們心裡面都很不舒服.
我們都很想安慰.
那些死難者的家屬.
但有時我們都不知道怎樣安慰.
我看到很多新聞.
那些家屬是與死難者.
談著電話.
聽著.
他的家人.
就在火場裡面.
慢慢沒有了消息.
如果我們簡單一點問AI.
我們很想安慰這些弟兄姊妹.
這些死難者的家屬.
不如你給我們一些安慰說話.
我想AI會傳遞很多.
這些說話出來.
但這些人不是盜位.
這些人不是盜獄.
這些人不是入心.
這些人不是真真正正能夠安慰到人.
有時我作為牧者.

$^{521}$在那段時間我都詞窮.
我相信他們正在面對這麼大的家破人亡.
其實我們告訴他們.
神與你同在.
不單止沒有作用.
可能他們覺得比粗口更難受.
有時我們想想.
真真正正他們怎樣才得到安慰.
我想起自己在艱難的時候.
怎樣才真正得到安慰.
是因為我想起我平時.
神怎樣不離不棄的.
與自己同在.
怎樣沒有放棄過自己.
想起這些經歷的時候.
就在這些艱難的日子裡.
就成為了我很大的安慰.
所以同神建立緊密的關係.
時刻感受祂的不離不棄.
是非常重要的.
今時今日怎樣幫到我們.
去面對這麼艱難的逆境.
最近有個弟兄跟我說.
他說他的女兒是中學生.
有一天問他女兒.
你將來長大後想做什麼.
他女兒就告訴他.
我想做獸醫.
然後爸爸就問他.
你做獸醫.
你知不知道有什麼資格.
於是女兒就去做資料搜集.
看看究竟怎樣才能讀獸醫.
他就告訴他爸爸.
原來讀獸醫最好是在澳洲讀.
怎樣才能讀獸醫呢.
原來現在60\%是看你學術的成績.
即是本科的成績.
然後40\%就要考一個試.
是叫什麼試呢.

$^{561}$就是CASPER的試.
這個試是考什麼呢.
就是考你那些溝通能力.
同理心 道德 適應能力.
判斷能力 團隊合作等等.
於是他爸爸知道他女兒要考CASPER的試.
於是就拿一些試卷來問他.
那些試是什麼呢.
舉個例子.
譬如就說.
當你看到一個人在便利店那裡.
好像沒有給錢.
就拿東西放在袋子裡.
那你應該怎樣做呢.
就類似這些問題.
你就要回答.
然後他就問他女兒.
他女兒就回答.
當然是告發他.
就去跟職員說.
告發他.
如果這樣回答.
那你那40分就沒有了.
那你要怎樣回答呢.
你又要很有同理心.
才不會有判斷.
然後你不要以為他不給錢就放在袋子裡.
一定是偷東西.
如果你前次也是偷東西.
那你的做法就很不同了.
你要持平一點.
然後用同理心去代入他的心態.
然後又要好好地.
用他的角度去跟他溝通.
於是又要用辦法引導他.
去做正確的事等等等等.
用完這些方法.
那你要回答這條問題了.
那我想說什麼呢.
其實很多時候我們教會.

$^{601}$經常都很想吸引年輕人回教會.
通常因為我們機構和很多教會合作.
做年輕人的工作.
問教會的時候.
我就問他.
通常你怎樣去吸引年輕人回教會.
通常很多教務同工就跟我說.
我們都搞一些興趣班.
搞一些活動.
例如音樂.
夾樂隊.
打球.
那我通常就問他.
那通常果效大不大呢.
他說通常搞這些班的時候都很多人來的.
但一進到信仰的時候.
大家就好像拉牛上樹那樣.
他們覺得果效都不大.
我想說的是什麼呢.
其實很多時候我們吸引來.
去認識我們的信仰.
我們就要用一些.
和我們信仰沒有什麼關係的東西去吸引他.
例如打球.
不是完全沒有關係.
但關係不大.
那好了.
我們用這些東西吸引他之後.
然後我們就要他引入信仰的時候.
其實就隔了叢沙.
我們就做得很辛苦.
其實我們有沒有想過.
其實我們信仰就是自大吸引力.
其實我們不用這樣拉牛上樹.
不停去推銷別人.
你去回教會.
你去信耶穌.
為什麼這樣說呢.
其實今時今日最需要.
就是剛才說的軟實力.

$^{641}$你要有同理心.
你要懂得和別人溝通.
你要懂得團隊合作.
這些軟實力的高手.
你在外面沒有課程可以學的.
很難有課程可以教你.
那你在哪裡學是最適合的呢.
你跟耶穌學.
就最適合.
耶穌就是這些軟實力的高手.
你看一看福音書.
他怎樣處理.
怎樣有同理心.
怎樣和你溝通.
怎樣引導別人去思考.
所以如果我們教會擺明.
是用耶穌這些軟實力.
去讓大家知道.
當你能夠學上耶穌的時候.
你就不單止能夠成為教會裡面.
像耶穌的一個門徒.
你更加能夠成為在外面.
不會被AI取代.
那些優秀人才.
今時今日讀獸醫這些理科.
都要懂得和別人溝通.
都要懂得和別人相處.
所以如果我們真真正正明白.
和耶穌建立一個緊密的關係.
我們領受他怎樣和別人相交之道.
那我們就可以站穩陣腳.
第三我們再看.
第七第八節告訴我們.
神說只要剛強大大壯膽.
謹守進行.
我僕人摩西所吩咐你的一切律法.
不可偏離左右.
使你無論往哪裡去都可以順利.
這律法書不可離開你的口.
總要晝夜思想.

$^{681}$好使你謹守進行.
這書上所寫的一切話.
如此你的道路就可以亨通.
凡事順利.
這裡說到在這個鑰句裡面.
不單止神給我們應許.
不單止神要和我們建立關係.
神還賜下寶貴的律法.
叫我們謹守進行.
這裡不單止.
highlight就不單止謹守進行.
還要晝夜思想.
當我們又晝夜思想.
又謹守進行上帝的律法的時候.
神就應許什麼.
我們道路就可以亨通.
凡事順利.
厲害吧.
道路亨通.
凡事順利.
我們真是夢寐以求.
但怎樣做到.
我們就要晝夜思想.
謹守進行上帝一切的律法.
好了 弟兄姊妹.
你相不相信當我們真的這樣晝夜思想.
謹守進行上帝的律法.
我們的道路就可以亨通呢.
我看到大家有些迷惘.
是不是呢.
有時即使做基督徒.
有時都會質疑一下究竟是不是呢.
大家質疑的時候.
我就給大家看.
前一陣子.
我就發覺這段時間.
如果你有留意.
你會發覺.
有很多這些的報道.
就說到無論在英國,美國.

$^{721}$甚至歐洲的地方.
過去一直以來.
信主的年青人越來越低.
但是呢.
現在越來越多的統計.
說一班Gen Z世代的年青人.
突然間喜歡看聖經.
好了 因為我機構是做生涯規劃.
我很多時問年青人.
你將來大概想做什麼.
十個大概有八個.
你知不知道他們想做什麼.
知不知道.
就是想做YouTuber.
是啦.
那你要做YouTuber.
你要拍些片.
點擊率要高.
過去你拍些什麼片.
點擊率會高呢.
就是那些飲飲食食.
去旅行.
開開心心的片.
點擊率高.
但是現在你留意一下YouTube.
留意一下TikTok.
有很多年青人拍的.
歐美的年青人拍的片.
他在拍.
他就說.
其實我買了一個新的版本的聖經.
然後我看完之後.
很啟發我.
很有領受.
這樣就點擊率很高.
好了.
這些的報道.
我有留意.
然後記者就問下這些.
就是Gen Z.

$^{761}$英美的Gen Z.
為什麼你們突然間這麼喜歡看聖經呢.
我先買一買個關子.
因為.
這裡是英美.
這裡不是英美.
這裡是香港.
是不是.
香港未必是這樣的.
好了.
大家如果這樣想.
我想給大家看一幅圖.
大家知不知道這裡是哪裡呢.
應該都不知道.
因為大家沒來過.
就是我家.
這裡.
有一班Gen Z.
他們是一班大學生.
其實他們在做什麼.
就是在查聖經.
那我一個月就會.
這個是其中一個群組.
我就會在家裡查聖經.
好了.
那.
嘩.
就是我現在年紀.
我二十多年了.
我事奉了.
不得不強調.
我年紀都越來越大了.
好了.
他們查聖經.
六點多在我家裡.
一路查聖經.
你知不知道查到幾點才走呢.
查到十二點多都不走.
好了.
那接著.

$^{801}$剛才我說了年紀越來越大.
做了機構之後.
就跟木櫃都是一天放假.
做了機構之後.
一至禮拜日都要.
上班就回教會.
接著又這麼晚.
十二點多.
那時候又給他們一些說話.
我就說.
我澳洲信主.
我就說.
外國的教會生活.
跟香港真的很不同.
就是你外國又夠大.
那你在家裡查經查得晚.
聚會聚得晚.
不要緊的.
你可以借宿一宵.
家裡有很多地方.
但是你不借宿一宵都不要緊.
你在家裡.
你自己開車來的.
外國大家自己開車.
你可以自己回家.
但是香港就不同了.
你看家裡這麼小.
哪有得給你借宿一宵.
又沒有被子.
接著我住調景嶺.
大家有些人住粉嶺上水.
現在十二點多.
快沒有地鐵了.
不如我們不要停在這裡.
好的東西不用磨爛.
下次再查.
好了.
那你知不知道年輕人說什麼呢.
有什麼回應呢.
他說其實我們都可以在這裡睡.

$^{841}$接著我就說.
你們很喜歡在別人家裡睡嗎.
我說不是啊.
但不知道為什麼.
當我常被我們渴望了聖經之後.
我們就很嚮往那些早期教會的生活.
凡物公用.
一起祈禱.
一起讀神的話語.
我看完這班年輕人.
我都嘆為觀止.
於是我不得不.
我覺得我是幫全香港的教會問的一個問題.
問什麼問題我就問.
你們這些年輕人對我們這些牧者來說.
其實感覺是什麼呢.
你們很忙的.
又艱難的.
讀書又忙.
又沒有錢.
你們要做兼職.
做兼職又忙.
接著要上莊又忙.
教會事奉又忙.
很忙.
我和你不是認識了很久.
你來家裡查經又查到三更半夜.
我想問他們什麼問題呢.
我就想問.
究竟教會有什麼聚會.
是可以令到你們.
磨爛脊賴死不走的.
你覺得我是不是幫全香港的教會問呢.
如果大家都知道你的答案.
照著去做.
那全香港的教會.
那些年輕人都興旺了.
那些年輕人就很認真地想.
想完之後.
他就回答了我兩個字.

$^{881}$大家猜一下是哪兩個字.
應該都是猜不到的.
因為我也猜不到.
但是覺得很合理.
他們就回答了我純粹這兩個字.
那我就叫他們說.
不如你們多說一點什麼叫純粹.
他說其實我們現在這班年輕人.
正在面對GEN-C.
現在GEN-α.
面對一個什麼世代呢.
是一個非常迷失.
極度迷失的世代.
因為他們每一天接收的資訊都非常多.
有統計一個年輕人一天接收的資訊.
等於十六七歲的人.
一個人一世接收的資訊.
所以現在很辛苦.
每一天就等於人家一世接收這麼多資訊.
當你接收這麼多資訊的時候.
你就會怎樣呢.
你就會選擇困難.
A好像對 B好像對.
於是大家就選擇困難.
又很怕做錯選擇.
又不知道怎樣做選擇.
又看到很多假的東西.
所以他說過去他們不想給信仰框住自己.
但原來不給信仰框住自己是更迷失的.
他還不知道怎樣選擇.
所以為什麼他們這麼認真.
他說如果有一間教會是認認真真.
讓他們深度認識聖經.
認識真理.
認識耶穌.
這間教會就吸引他們回來.
因為他們深度認識了耶穌.
認識了真理之後.
他們就有依歸可以做選擇.
這班年青人只是少數.

$^{921}$接著我聽完之後.
我周圍看到年青人就問.
其實答案是大同小異.
所以今時今日.
其實我們有沒有掌握到我們信仰的優勢.
神早已經給我們立這個約.
在N這麼多年前已經立了這個約.
祂叫我們立足上帝給我們的應許.
以至我們抓緊這個應許.
我們就可以建構很穩固的信念.
在這個左搖右擺的世代裡面.
我們可以站立得穩.
上帝跟我們建立緊密的關係.
去猜探耶穌道成肉身來這個世界.
祂怎樣活一次.
我們照著來跟祂去學習祂的榜樣.
我們跟祂建立緊密愛的關係.
我們就能夠成為別人真正情感的引導者.
我們抓緊著上帝的話語.
我們就可以做到價值判斷.
成為別人價值的引導者.
當我們真正成為這三項的強者的時候.
我們怎會被AI淘汰.
我們真的做回一個人.
活出人應有的價值.
我們就能夠真真實實.
並且很輕省去見證主在我們生命裡面.
那份的人性.
那份的光輝.
那份的吸引力.
但願上帝賜福黃浸每一位頂子妹.
我們好好抓緊上帝的約.
成為一個上帝所應許豐盛的人生.
願主賜福大家.
\newpage



\section{詩篇 16:1-11}
\label{sec:3zRaVYVNQY0}
\textbf{2026年3月29日崇拜直播｜鄭國邦傳道｜基督之詩篇:不讓你的聖者遇見朽壞｜詩篇16\_1-11}
\newline
\newline
連結: \href{https://youtube.com/watch?v=3zRaVYVNQY0}{\texttt{https://youtube.com/watch?v=3zRaVYVNQY0}} ~~~~ 語音日期: 2026-03-30
\newline
\newline
\hyperref[sec:p8dj6FAIn5Y]{\small{< < < PREV SERMON < < <}}
~
\hyperref[sec:index_chronic]{\small{[返順時目]}}
~
\hyperref[sec:ltS5lDDLFjU]{\small{> > > NEXT SERMON > > >}}
\newline
\newline
詩篇 16:1-11
\newline
\begin{longtable}{cl}
\hline
\hline
章節 & 經文 (和合本修訂版)\\
\hline
16:1 & \begin{tabularx}{0.7\textwidth}{X} 神啊，求你保佑我，因為我投靠你。 \end{tabularx} \\ \\ \relax
16:2 & \begin{tabularx}{0.7\textwidth}{X} 我曾對耶和華說：「你是我的主，我的福氣惟獨從你而來。」 \end{tabularx} \\ \\ \relax
16:3 & \begin{tabularx}{0.7\textwidth}{X} 論到世上的聖民，他們是尊貴的人，是我最喜悅的。 \end{tabularx} \\ \\ \relax
16:4 & \begin{tabularx}{0.7\textwidth}{X} 追逐別神的，他們的愁苦必增加；他們所澆奠的血我不獻上，我嘴唇也不提別神的名號。 \end{tabularx} \\ \\ \relax
16:5 & \begin{tabularx}{0.7\textwidth}{X} 耶和華是我的產業，是我杯中的福分；我所得的，你為我持守。 \end{tabularx} \\ \\ \relax
16:6 & \begin{tabularx}{0.7\textwidth}{X} 用繩量給我的地界，坐落在佳美之處；我的產業實在美好。 \end{tabularx} \\ \\ \relax
16:7 & \begin{tabularx}{0.7\textwidth}{X} 我要稱頌那指引我的耶和華，在夜間我的心腸也指教我。 \end{tabularx} \\ \\ \relax
16:8 & \begin{tabularx}{0.7\textwidth}{X} 我讓耶和華常在我面前，因他在我右邊，我就不致動搖。 \end{tabularx} \\ \\ \relax
16:9 & \begin{tabularx}{0.7\textwidth}{X} 因此，我的心歡喜，我的靈快樂；我的肉身也要安然居住。 \end{tabularx} \\ \\ \relax
16:10 & \begin{tabularx}{0.7\textwidth}{X} 因為你必不將我的靈魂撇在陰間，也不讓你的聖者見地府。 \end{tabularx} \\ \\ \relax
16:11 & \begin{tabularx}{0.7\textwidth}{X} 你必將生命的道路指示我。在你面前有滿足的喜樂，在你右手中有永遠的福樂。 \end{tabularx} \\ \\
[1ex]
\hline
\hline
\end{longtable}
$^{1}$詩篇第十六篇.
神啊求你保佑我.
因為我投靠你.
我的心啊.
你曾對耶和華說.
你是我的主.
我的好處不在你以外.
令到世上的聖民.
他們又美又善.
是我最喜悅的.
以別臣代替耶和華的.
他們的受苦必加增.
他們所澆點的血.
我不欠賞.
我的嘴唇也不提別臣的名號.
耶和華是我的產業.
是我杯中的糞.
我所得的.
你為我賜受.
用乘涼給我的地界.
坐落在皆美之處.
我的產業實在美好.
我必稱頌那只有我的耶和華.
我的心牆在夜間也警戒我.
我將耶和華常擺在我面前.
因祂在我右邊.
我便不自搖動.
因此我的心歡喜.
我的靈快樂.
我的肉身也要安然居住.
因為你必不將我的靈魂.
撇在陰間.
也不叫你的聖者見紐壞.
你必將生命的道路指示我.
在你面前有滿足的喜樂.
在你右手中有永遠的福樂.
這是上帝的話語.
很偉大等等.
最重要的是.
我們要將自己的生命也讀進去.

$^{41}$以至讓上帝的話.
和我們的生命.
都能有這份潔潔.
才能真正體驗到.
上帝如何今天和我們同行.
而上帝所做的事.
不只是停留在二千多年前.
為我們做的.
而是今天也同樣.
在我們生命裡進行.
所以當我們要讀入.
今天的基督之詩.
在裡面去明白上帝的話的時候.
我們也嘗試將自己的生命.
投入在當中去想一想.
先給大家一個問題去想一想.
去回想你從小到大.
有沒有經歷過被人扔下.
這個經歷可能是很傷痛.
但先試想一想.
當一說扔下這件事.
我一定會回想起.
我一個中學同學的一個經歷.
是什麼呢.
他長大了.
說回一次經歷.
他有一次跟中學的 field trip.
即是出去考察.
去到國內一個地方.
他說是一個茶園.
即是跟大隊十幾人.
一起去到一個茶園去考察.
很大的.
茶園整個山頭.
他就說去完茶園.
品完茶.
即是一班人和老師.
大家要走了.
喝完茶.
特別可能急一點.

$^{81}$他就去廁所.
去完廁所.
人去樓空.
整個地方沒有人.
原來旅遊巴已經開走了.
只剩下他一個人.
扔下他一個人.
接著可能你都會想.
很慌.
在國內山頭野嶺.
被人扔下很害怕.
他說他一點都不害怕.
他還坐在一塊石頭上.
看著茶園在等.
都幾好風景.
你想想那時候一點都不害怕.
他在想什麼呢.
因為他在想.
你們數人數的時候.
都會知道少了一個人.
少了一個人的時候.
一定會回來找我.
所以我都沒什麼好怕.
就坐著等.
與其在慌.
走來走去.
不如坐著等.
正如他所想.
半個小時之後.
都等了半個小時.
之後旅遊巴就隨隨的.
回來接他.
當時我們問他.
你害不害怕.
他說不害怕.
我知道他不會扔下我.
他們不夠人.
一定都要回來.
下一個景點都要數齊人.
只不過可能遲了一點.

$^{121}$上車的時候.
大家就想著要走了.
原來這個經歷.
好看的就是.
可能扔下是一件事.
但是背後經歷這件事的信念.
和你的想法.
可能會影響你當時的處境.
如何回應.
當然剛才想.
你從小到大.
有沒有被人扔下的經歷.
未必好像這個同學.
真的被人肉身上扔下在茶園.
可能是在心靈裡.
可能在關係上.
可能在小時候.
拍拖被人扔了.
又或者好像被人離棄了.
又或者可能那件事很簡單.
可能對方只是沒有回覆你訊息.
你已經感到.
他扔下我了.
他不理我了.
或者等等.
所以這個經驗.
其實是很見仁見智的.
但是那個感受.
可能都是很真實.
出現在你的生命裡.
而當中人生裡.
我們或多或少.
一定會經歷過這些時刻.
可能被人背叛.
或者被人出賣.
又或者好端端的一個關係.
突然之間對方斷亂了.
很傷那種感覺.
當然對方可能有種種原因去解釋.
我們不理對方是怎樣.

$^{161}$但是那個被扔下的人.
那種受傷害的感覺是真實的.
可能在受這種傷害的時候.
你可以埋怨很多事情.
阿神為什麼你讓這件事發生.
或者埋怨那個人.
你為什麼這樣做.
又或者埋怨自己.
為什麼我這麼蠢.
這樣也被他傷害到.
很多這些的可能去出現.
但是當我們.
好像剛才所說.
經歷被撇下.
被遺忘.
這種經歷.
有實在的那種感受的時候.
你背後的信念是什麼呢.
經歷著那件事的時候.
你怎樣理解這個狀況呢.
也會影響到今天你怎樣去看.
上帝在你生命裡面的工作.
當我們看回詩篇第十六篇.
是大衛的金詩.
如果你看回細節的時候.
金詩其實不是很多手.
這是其中之一.
裡面就是說大衛在經歷過危難的時候.
得到上帝的幫助.
以至他能夠經歷過這一切的時候.
告訴我們.
他怎樣憑著信心.
背後是什麼信念.
以至他能夠捉緊上帝去經歷.
同樣的.
為什麼說這首是基督之詩呢.
我們稍後的後半部分就會說.
就是大衛同樣受到這份感動.
他能夠說出一些.
不是當時的人能夠理解到的一些說話.

$^{201}$以至他帶出.
原來上帝真正的那份拯救和幫助.
是在耶穌基督的身上去成就.
所以我們先去看回.
十六篇的內文.
理解了內容的時候.
我們就可以在後半部分.
讀入在新約裡面耶穌基督所做的事.
連繫到今天我們怎樣應用和經歷上帝.
先看回詩篇十六篇一至二節.
神啊 求你保佑我.
因為我投靠你.
我的心啊 你曾對耶和華說.
你是我的主.
我的好處不在你以外.
這個是大衛一開始的時候的命題.
如果你教同學作文.
通常第一句就是最重要的.
那句話你一看.
後面那些不看也不要緊.
最重要是看頭和尾.
第一句就是劈頭.
第一句的命題的時候.
很重要的.
也是大衛的宣告.
向人說究竟他所信的神.
是一個怎樣的神.
他說什麼呢.
我的主啊 神啊 求你保佑我.
因為我投靠你.
這個心態很重要.
投靠那兩個字.
投靠本身原文裡面是說.
好像如果大家記得舊約裡面.
有那些屠城.
設立一些制度的時候.
就是你走進去一個地方.
去得到幫助.
這個就是投靠的意思.
今天不知道大家那種.

$^{241}$投靠的心是怎樣.
有兩種可能.
一種就是.
神啊 你那麼厲害.
我真的都很信你.
我在這裡.
你幫我吧.
我有困難.
你幫我吧.
我現在需要你幫助.
這種就是其中一種呼求.
但是私人大衛說那種是什麼呢.
神啊 我搞不定.
我要走去你那裡得到幫助.
大家明白兩個分別嗎.
一個是站在那裡.
神啊 你幫我吧.
你那麼厲害.
你在這裡幫我吧.
另一種是.
我真的搞不定 神啊.
我要投靠你.
大衛是後者.
今天幫我們去呼求神.
去認定這是幫助我們神的時候.
我們那種心態是哪一種呢.
大衛就說.
他用這種尋求庇護的行動.
去告訴我們.
就是這種態度.
而他的心呢.
就在第二節說.
我的心啊.
你曾對耶和華說.
你是我的主.
我的好處不在你以外.
他那種行動.
不只是心裡面相信神有大能.
而且在他的生命.
在他的生活裡面.

$^{281}$都是這種投靠的心.
所以你會看到.
大衛一開始就說.
這種裡外一致的信仰宣告.
所以今天當我們.
去到上帝面前的時候.
我們是哪一種呢.
我是那個主.
我有困難.
你來幫我.
好像那些大佬那樣.
叫齊那些人.
來幫我.
我們去吧.
還是不是.
我們是做跟班那個.
我有事.
我要走去我大佬那裡.
我的主那裡.
去找他幫助我.
以致當我們面對著.
無論那種急難是怎樣.
在這些困難裡面去找他.
當順境的時候.
你也不會只是坐著說.
我現在一切都安好了.
我不當你存在了.
或者不用找你了.
順境的時候.
你又怎樣裡外一致.
去表達出.
這個是你的主呢.
因為在順境的時候.
他也說.
我的好處都不在你以外.
就算今天我在一切的安穩.
順境裡面的時候.
私人都說.
我的一切好處.
都不在你之外.

$^{321}$所以你看到.
大衛不是那些.
有事才穩實.
沒事就算了.
他是知道一切都是出於實.
但最慘的是什麼呢.
有事都不穩實.
沒事不穩實.
都正常.
有事都不穩實.
為什麼呢.
因為都是靠自己.
有事我自己搞得定.
我有手有腳.
或者我有這些東西.
有這些資源.
我解決到.
真的解決不到才找你.
真的解決不到的那個.
就已經不知道是什麼時候了.
而且可能你心底裡.
不是真的只是想.
我真的解決不到才找神.
可能再埋深一點的想法是.
找神真的解決到嗎.
我有手有腳都解決不到.
神你真的可以解決到這件事嗎.
我自己現在都沒有把握了.
我找到神的時候.
這個主權交給你.
我就更加沒有把握.
我怎樣可以信靠你呢.
甚或乎可能覺得.
找神是很羞愧的.
要神幫好像不行.
或者很不濟的時候才幫.
我初信主的時候.
我爸爸跟我說什麼呢.
你為什麼信耶穌.
你又不是吸毒.

$^{361}$你又不是犯罪.
你又不是做了這些.
又不是沒有手沒有腳.
為什麼要信耶穌.
因為在我爸爸的概念裡面.
他的概念是.
那些人在生命走到絕處.
很差的時候才信耶穌.
可能你都會想.
我現在還沒有到那個情況.
不用找上帝了.
我靠自己就解決了.
那究竟什麼時候.
我們才有那個心態.
我要走去神那裡.
是解決不到才找.
還是我覺得我自己都還沒解決到.
所以不用這樣住.
而使之我們還沒經歷到徹底的投靠神的心.
詩人在這裡告訴我們.
這份徹底投靠神是怎麼樣.
就是一切的好處.
一切好的事都不在上帝以外.
上帝不是我們其中一個錦囊.
上帝不是我們其中一個選擇.
是一切的好處都不在上帝以外.
是我們的唯一.
所以當讀到這裡的時候.
想想我的心裡面.
怎麼看這個上帝呢?.
在心裡面你現在那個投靠神的心有多高呢?.
以致你是不是真的在相信著上帝的身上.
你得到你需要的東西呢?.
接著繼續看下去.
當我們能夠認信和定神的那個.
在我們生命裡面的位置的時候.
他說論到世上的聖文.
他們又美又善是我最喜悅的.
接著就說.
追逐著別神的.

$^{401}$他們的仇負一定會加增.
他們所囂電的血我不獻上.
我的嘴唇也不提別神的名號.
先看第三節.
聖文不是說那些一出來就有光那些.
聖文在聖經裡面不斷說.
不是那種很聖潔無瑕.
是說那些願意分別出來.
歸順屬於上帝的一群人.
以致在裡面不只是說那些外表.
好像模特兒一樣.
好像很漂亮.
或者你看那些照片.
很漂亮的那種.
而是那些在生命裡面.
願意去屬於上帝.
完全屬於上帝的人.
而那個又美又善.
不只是說那個外表.
而是在說那個Excellence.
那種很原理很嚮往.
平時在聖經裡面又美又善.
有時候會形容一些水源.
又或者會形容一些貴族的一種狀態.
所以那一種生命流露.
不一定只是在外表上那種.
而是那種生命力的流露.
是一種生命的狀態.
是使人嚮往著那種.
嘩!這種狀況不是只是外表.
是裡面也有力量.
而且這種又美又善.
是相對著後面第四節所說的那種對比.
不屬神的人會是怎樣呢?.
裡面就是那種愁煩.
愁眉慘目或者等等.
而那種又美又善.
就是那種反而.
嘩!一出來是有那種力量.
是有那種盼望在裡面當中.

$^{441}$這一種就是私人所嚮往的那種生命狀況.
所以在裡面一開始我們就想.
嘩!原來上帝是這樣信靠你.
接著再去到下一步.
原來信靠主的人是這樣的.
今天有沒有嚮往過這種生命呢?.
還是可能再看下去第四節.
在裡面說以別神代替耶和華.
其實在原文裡面也沒有說是別的神.
只是說別的而已.
找其他的東西去代替神.
所以不一定是在說神明.
可能是你其他的東西.
將你剛才說的那種生命力.
全神貫注在一些東西上面去代替了神.
已經是在說好像第四節那種.
所以如果你的生命有別的事.
去取代了上帝.
可能你的仇苦就會這樣加震.
你的別的未必是一定壞事.
可能你的別的是在說.
你很著緊著你的工作或者學業.
或者是一些緊張的事.
甚至你的投資.
嘩!最近全都掉下去了.
嘩!很緊張啊!我的別的就是.
嘩!怎麼辦呢?.
或者那個別的可能是你的家人.
或者等等的事.
本身不是壞事.
那個別的未必一定是殺人放火.
壞事做盡講的那些東西.
而那個別的的事.
可能將你的所有專注和心力.
都投放在裡面.
比上帝更加重要.
這個狀況已經令你.
好像第四節所說那種.
仇苦就不斷加震.
簡單來說.

$^{481}$大衛很想說.
這個世界上有兩個選擇.
一個是屬於上帝的.
一個就是將你的心力去消電.
獻上給其他的東西當中.
而你又是在哪一種狀況裡面呢?.
可能你會覺得.
大衛這樣寫好像很極端.
嘩!將這個世界分做.
一是屬神一是不屬神.
當然大衛就是想用這種形態.
就是告訴我們.
今天究竟你是在什麼狀況.
而我們會不會在另一個極端裡面.
完全不分這兩樣東西.
都沒有所謂.
不要說分不分屬神.
還是屬這個世界.
沒有所謂.
所有東西都沒有所謂.
我都是這樣生活.
走向了另一個極端.
而大衛為什麼用那邊來說.
只有兩種.
一是屬神一是不屬神.
就是想提醒我們.
今天我們的焦點.
是放在什麼裡面呢?.
是放在別的裡面.
還是放在唯一的上帝裡面呢?.
以致當你放在唯一的上帝裡面的時候.
你有沒有產生著那種嚮往.
嘩!原來屬於神的人真的這麼好.
又美又善.
還是會變成將這份最喜悅的.
最嚮往的.
放在其他的焦點裡面.
可能不知道你有沒有想過.
其實不上教會的人都挺開心的.
有沒有見過?.

$^{521}$有沒有這樣想過?.
嘩!其實他們都挺好的.
星期日又多了一點時間.
其實他們又不用想那麼多限制上的東西.
其實他們都挺好的.
嘩!喜歡打牌又可以.
喜歡說什麼話又可以.
自己自由.
或者你在看別人的社交媒體.
嘩!其實都挺好的.
這些生活都挺好的.
花天酒地又好.
有空又換一下跑車.
有空又去一下這裡旅行.
都挺好的.
嘩!好像很想他們這樣的生活.
可能你總是覺得.
其他的會比你所擁有的好.
甚或乎不是比較外在的生活.
可能比較家裡.
嘩!你看人家老公這麼好.
為什麼我老公這樣.
嘩!人家老婆精打細算.
為什麼我老婆這樣.
又或者.
嘩!人家的家人是這樣的.
嘩!我覺得真的不好.
當這些其他的東西不斷去佔據你.
還有去告訴你他好像比你好.
但是詩人就說.
有什麼才是你真正嚮往的呢?.
他真正嚮往的就是那班屬於上帝的聖民.
他們的生命力是又美又善.
不是那種生活形態.
不是那種生活方式或者生活的質素.
是那種生命力.
讓他嚮往著這一種的生活.
所以當回顧一至四節.
當他定睛在上帝是一個怎樣的上帝.
當他再定睛他今天所認信.

$^{561}$他想覺得最好的選擇是什麼的時候.
他就再進一步去第五第六節.
我有這些認信有這些肯定的時候.
怎麼看今天我擁有的一切的事呢?.
第五節說耶和華是我的產業.
是我杯中的糞.
我所得的你為我持守.
醒糧給我的地界坐落在街尾之處.
我的產業實在美好.
簡單去剖析一下這段經文就是說.
他杯中的糞.
他所得的用醒糧給我的.
很簡單的概念就是.
所有都是被動.
杯中的糞的意思是什麼呢?.
就是當你去喝.
或者別人請你作客的時候.
你拿個杯出來.
他就倒多少給你.
那你會不會跟主人說多點多點多點.
不夠不夠多點.
這樣不會的.
客氣都會說夠了夠了.
杯中的糞就是他給你的.
有多少是他決定給你的.
他所糧給你的地界.
是很美好.
就是多大就多大.
好像當時十二支派那樣.
摩西領袖說糧給他們多少就多少.
會不會反抗會開個庭.
沒有的.
他們都是馴服著說多少就多少.
以致當他重看的時候.
他是很被動的知道.
今天我擁有的都是上帝所糧給我的.
我不會覺得上帝你好像差一點.
這個可以給我.
又或者這個其實有這個會更好的.
上帝不如你給我.

$^{601}$斯人說沒有.
糧多少就是多少.
住多少就是多少.
當然如果你是一個生活比較基層.
或者窮的人.
可能你都會覺得.
命裡有時終須有.
命裡無時莫強求.
即是一個窮的人.
你這樣唱都正常.
你沒有.
你沒有這些的時候.
當然有就有.
沒有就沒有.
但當你有少少能力.
有少少知識.
和有少少基本的財富的時候.
你就會想我要多一點.
或者好像可以再增值多一點.
甚麼乎其實大衛不是一個窮人.
大衛就算這篇詩未做君王之前.
他都是素羅的女婿.
都不是一個窮人.
所以其實大衛絕對不是一個窮人.
他就說有就有.
沒有就沒有.
其實我都是很被動的.
他是一個有能力有錢的人.
甚或乎都經歷過年少氣盛.
闖蕩江湖.
去做很多事的時候.
到這一刻他都是這樣說.
你糧給我的.
你針給我的.
我所擁有的.
都不是靠我自己.
是上帝你給我.
在裡面第五和第六節.
其實你會見到一頭一尾是一個呼應.
他說耶和華是我的產業.

$^{641}$到第六節尾.
你就知道我的產業實在美好.
很簡單教同學.
A等於B.
B等於C.
A是不是等於C.
就是.
所以最美好是甚麼呢.
就是耶和華就是最美好.
這個就是第五第六節裡面.
他告訴我們的一件事.
可能你說.
今日我擁有的東西.
是否真的上帝完全給我的呢.
可能你未完全調教到.
今日究竟我們怎樣看呢.
有機會你又可以找一下.
洪福堂的老闆.
即是CEO.
使徒永富的見證.
又看一下.
他不是窮人.
他不會說有就有.
沒有就沒有.
他亦都是一個經歷過高高低低的商人.
在裡面去怎樣定精在上帝的裡面.
以致他能夠過往叱吒風雲.
擁有很多在金融界的資財或者財富.
到遼島然後就創立了洪福堂的生意.
用上帝的心意去建立這個生意的時候.
你就會見到今日.
其實這件事不只是大衛時代的事.
是今日我們都可以這樣去看我們所擁有的一切的事.
不是我們自己爭取回來的.
是上帝糧給我們每一個.
當上帝糧給我們每一個都是這樣多的時候.
你有沒有爭攤真是很好.
上帝你給我的很好.
哪怕你可能只是住在百幾二百呎的tong 房.
都幾好.

$^{681}$都起碼可以住.
可能你在萬幾二萬呎的家裡都很好.
而不是只是不斷覺得這樣會再好一點.
其他的可能你嚮往的可能要更大更好.
現在已經是最好了.
當你有一個這樣供應你的上帝.
或者當你定精在上帝是這樣的上帝.
我擁有的一切都是上帝的時候.
你會怎樣去讚美這個上帝呢.
第七節繼續說.
我要稱頌那指教我的耶和華.
在夜間我的心牆也都指教我.
第七節裡面說什麼呢.
我要稱頌那個指教我的神.
突然之間好像畫風一轉.
當你如果擁有這一切都美好.
當你這樣定精上帝的時候.
你如果平時祈禱的時候.
上帝你是一個聽祈禱的上帝.
上帝你是一個滿有榮耀的上帝.
上帝你是一個很有能力的上帝.
但大衛在這裡特別去形容那個上帝是什麼呢.
那個指教我的上帝.
那個指教我的上帝就是在說.
今日當我有這一切.
我有這個認順.
你怎樣在每日的生活裡面去做決定呢.
當你相信這一切都是出於上帝的時候.
那上帝在你下一步.
在你未來的每一個決定裡面.
有沒有成為你的提醒和幫助呢.
接著第八節就說.
我要讓耶和華上在我面前.
因祂在我右邊我就不自動搖.
什麼是常常放在上帝面前呢.
和指教我怎樣走在一起呢.
早一陣子和一個小組玩了一些遊戲.
我們其中一個遊戲呢.
就要拿一個電話出來.
展示一個桌布給大家看.

$^{721}$究竟你的桌布是放了什麼.
成為你的那個桌面.
就是你的電話是可以調較桌面的.
可能有些就會放你的子女.
你的孫子或者是一些重要的.
可能有些人放跑車.
有些人可能放他的旅行的照片.
有些人可能是放你男朋友女朋友的照片.
這個就是你將什麼常常放在你面前.
為什麼你放這些東西在桌布上呢.
就是因為你喜歡嘛.
你見到的時候一打開一滾開.
就見到它就很開心了.
有沒有人放經文的.
應該都有的可能有的.
有沒有人放耶穌的樣子在這裡呢.
或者有沒有人放一些.
一些屬靈的東西在這裡呢.
不是對錯.
純粹這樣說而已.
究竟你放什麼在你面前.
就是代表著你著什麼.
你如果放你伴侶在裡面.
就是你很著緊你的伴侶.
上次一滾開就見到他.
然後你放這些東西.
就是常常放在你面前.
就是你緊張的東西.
當然在裡面怎樣成為一個提醒呢.
他就說當我有這個認信.
當我願意按著神的心意.
去這樣選擇和擁有這一切的時候.
這一個才是最大最大的安全感.
即是說你正在經歷.
好像被人丟下.
即是說你正在經歷.
好像你真的什麼都沒有.
但是有著剛才這一切的時候.
這份安全感才穩如泰山.
以致好像詩人所說.

$^{761}$不自動搖.
當我們去經歷很多事情的時候.
你有沒有想過你生命的安全感.
建立在什麼上面呢.
剛才說了你嚮往的東西是什麼.
你擺在面前的東西是什麼.
這些就是你安全感勾緊的東西.
而當我們正在上正面管教的時候.
有一個slide那個show出來.
我們教會在學校.
或者平時我們家長裡面.
都有上正面管教班.
這份安全感你不要忽略.
這份安全感原來在小朋友身上.
已經很重要.
當你看slide的時候.
左手邊就說一個小朋友缺乏安全感.
會哭鬧.
會有很多情緒等等的事.
可能今天你看到一個小朋友的狀況.
其實都是源於他沒有那份安全感.
或他和父母分離.
或他可能家庭裡面都沒有那份安全感.
或許很多事情令他覺得.
很不感覺到安穩.
所以他就會去做很多表演.
去想要人的注意.
最大的一句就是.
其實這份父母給他的安全感.
或環境給他的安全感.
是一切發展的基礎.
以至能夠幫助他自信和獨立.
其實回想自己也是.
當你在一個群體或家裡.
你夠安全感的時候.
你就會去探索.
你就會願意離開父母多一步.
因為你知道父母不會離開.
你就多走一步去看看這個世界.
他還在這裡.

$^{801}$你就多走一步去看看這個世界.
因為你有安全感去做這件事.
但如果沒有的時候.
他就會很害怕做錯.
一做錯原來父母就離開了.
我當然不要做錯.
形成可能不知道是不是你的工作環境.
就是有一句.
多做多錯.
少做少錯.
不做就當然不會錯.
所以大家要躺平.
什麼都不要做.
但如果你夠安全感.
你不會只是想這件事.
我不做錯就夠.
如果你夠安全感.
你就會去想.
我怎樣可以做好一點.
上帝如果你這麼好.
我又這樣放你在我生命裡的時候.
我怎樣可以做好一點.
去活出這份生命力.
而不是只是想.
打牌又不對.
衝紅燈又不對.
這樣又不對.
不會只是想怎樣才做不錯的生命.
而是做好我們的生命.
去流露出來.
所以當我們繼續看最後.
9至11節.
當你有這個認信和經歷的時候.
你就會知道.
第九節說.
我的心就會歡喜.
我的靈就會快樂.
我的肉身也安然居住.
因為你一定不會將我的靈魂.
撇在陰間.

$^{841}$也不會將你的聖者.
見扭壞.
你一定會將生命的道路指示我.
在你面前有滿足的喜樂.
在你右手裡有永遠的福樂.
因著一開始的時候.
你認信上帝是一個怎樣的上帝.
當這樣認信的時候.
我怎樣選擇.
我嚮往的生命.
以至當我擺上帝在我面前.
每日的生活.
我知道怎樣去做的時候.
你就能夠經歷第九和十一節.
九至十一節裡面.
那種生命力.
那個生命力是怎樣呢.
他說我的心歡喜.
我的靈.
你見到那個原文就是說.
那種榮耀.
那種榮耀很快樂.
很輕飄飄的一種感覺.
以至他的肉身都安然居住.
那種狀態不是只是在說身體上.
是在說身心靈.
都是一種很開心.
很喜悅的狀況.
那種榮耀是在說什麼呢.
就是在說你不會覺得羞恥.
你不會覺得說.
你整天擺神在後面.
或者你醜不醜.
你靠你自己.
你不要整天說神.
你不會有那種羞辱感.
你的靈一發到.
我就是很嚮往.
我就是覺得這種生命力是吸引我.
我有力這樣去生活.

$^{881}$以至當我們認定上帝這份美善的時候.
他一定不會丟下我們.
這個認信是很重要的.
如果那個只是一個有能力的神.
但是他五時花六時變.
你都捉不到他.
或者都不知道他怎樣想的時候.
這個神多厲害都對你來說沒有幫助.
因為你沒有那份安全感在你的生命裡面.
但是聖經裡面就是不斷的允許.
上帝不會丟下你.
如果你記得上個星期的賽穆斯說約書亞那段經文.
其實都有說他不會丟下我們.
其實聖經裡面都不斷不斷的告訴你.
他從來都沒有丟下過他所愛的人.
只是反而反過來.
我們有時候看不見離開了.
他這份丟下呢.
不單止不會丟下.
不單止不會放下你.
亦都不會見你口壞.
不會讓你有你就算.
因為上帝說過.
他從來都沒有忘記你.
丟下最害怕的一個狀況就是.
被人忘記了.
那種忘記有時候看電影.
可能你都會覺得.
如果有一天我失智.
這個世界忘記一切的時候.
都幾心酸.
但是上帝不會失智.
上帝不會忘記我們每一個人.
就算我們經歷在什麼狀況之下.
他都不會忘記我們.
不會讓我們的身體這一切見扭壞.
而這份真正的安全感.
可以改寫我們的生命.
從今天開始.
我們不只是怕做錯了什麼.

$^{921}$而是我們生命既然有價值.
有能力的時候.
有意義的時候.
我們怎樣去做一些事情.
換取上帝的喜悅呢?.
正如剛才所說.
當我們有這份力量的時候.
我們就能夠經歷到大衛這份實在.
當然詩篇已經解釋完.
再跳下去.
和耶穌基督有什麼關係呢?.
這個看起來像是大衛人生經歷.
他的剖白就是這樣信靠神.
但是其實你一邊聽.
大衛你是被人追殺.
是經歷很多的事情.
你都未至於死到陰間.
死了就應該寫不到這篇詩.
捨得就應該未死透.
還在經歷中.
所以為什麼會說得這麼誇張.
又落到陰間呢?.
當然另一邊廂.
他說經歷到那份平安.
或者那份喜樂.
但是你看回大衛的經歷.
他那個被掃羅追殺完.
追殺完了.
沒事了.
然後他作王了.
然後又被他的兒子追殺.
又被非利士人.
其實他那份平安不是永遠的.
也不是永遠的.
為什麼第九第十一節.
說得這麼誇張呢?.
你那份喜樂那份平安.
真的直到永遠嗎?.
又不像啊.
其實你會看到大衛.

$^{961}$不單止將他自己的生命寫下去.
他更加是按著聖靈的感動.
明白到真正所嚮往的那種生命.
是怎麼樣.
甚至將耶穌基督.
在之後所作的事.
都寫入了九至十一節裡面.
為什麼這麼說呢?.
其實當你打開新約市圖行傳的時候.
這一篇詩篇都被引述過兩次.
你想想.
那時候經常說.
聖經不是印出來的.
是手抄的.
如果那樣東西不重要.
其實不會寫那麼多次.
也不會重複一樣的寫下去.
因為要節省筆墨和篇幅.
為什麼兩次要引述這篇詩篇十六篇呢?.
一開始的時候就在市圖行傳二章裡面說.
當彼得在五場節按著聖靈的感動.
領受聖靈的時候就去講述.
在二章二十五至三十六節.
都說到關於耶穌基督的事.
三十一節說什麼呢?.
就預先看明這件事.
論到耶穌基督復活的時候.
祂的靈魂不撇在陰間.
祂的肉身也不見了壞.
第三十六節的結論是什麼呢?.
故此以色列全家當確實知道.
你們釘在十字架上的耶穌.
神已經納祂為王為基督了.
在頭一次引述十六篇.
它的重點放多一點.
就是說上帝沒有將耶穌基督撇在陰間.
當這些猶太人.
想著迎接祂來作君王.
在祂來到耶路撒冷的時候.
揮了宗旨迎接祂.

$^{1001}$誰知失望.
也不是我們所要的.
然後就釘祂在十字架上.
甚或乎在聖經裡形容.
就是那個匠人所棄的那塊石頭.
就是這些猶太人撇棄的耶穌基督.
神納祂為救贖主.
成為救贖的事供.
所以當彼得引述的時候.
他說當你們以為被神撇棄.
被你們撇棄.
將一切世界都撇棄的耶穌基督.
神從來都沒有撇棄祂.
祂要成就救恩.
藉著祂去成就.
在裡面.
甚或乎當時的人都會笑耶穌.
你這個拆毀聖殿的.
三日又說重建.
你可以救救自己.
在裡面.
神就說祂沒有撇棄.
是你們這些猶太人撇棄的.
第二次.
當保羅在十三章裡引述.
他說十三章二十六節四十三節.
我選了幾節經文去說.
三十節就說.
神卻叫他從死裡復活.
三十七至三十八就說.
唯獨神所復活的.
他並未見紐懷.
所以弟兄們.
你們當曉得赦罪之道.
是由這個人傳給你們.
當保羅引述這個詩篇十六篇的時候.
他就強調復活這個向度.
就是說.
如果上帝藉著耶穌基督去成就這份救恩.
他一定要經歷這個死亡.

$^{1041}$以至復活帶來的那種生命力.
能夠將赦罪之道帶來給我們.
以至這個生命之道.
就經歷到真正永遠的福樂.
剛才說了.
大衛的時代他不知道什麼是永遠.
因為他的經歷都是在被追殺.
或者是得救的那一個.
但是藉著耶穌基督.
我們知道什麼是永遠的福樂.
所以當整合了所有的東西的時候.
回想起一開始的經歷.
有沒有被人丟下過.
甚至你現在會不會覺得.
都好像有一些感受.
就是好像被人.
世界都離我而去了.
上帝沒有丟下過我們.
看似耶穌基督被人撇棄.
從來都沒有撇棄.
反而藉著這件事成就.
得著永遠的福樂.
今天你的福樂.
你所嚮往的是什麼呢?.
保羅在哥林多前書都有說過.
若基督沒有復活.
我們所傳的便是謊言.
你們所信的也都是謊言.
今天在復活節之前的一個星期.
我們有沒有去想這個復活的大能.
如何進入我們的生命裡面.
去更新我們.
將這些說話進入我們的生命裡面.
去改變 去提醒我們.
以至今天我們的認訊是什麼呢?.
我們的生命的焦點是什麼呢?.
最終我們可以同樣靠著.
好像大衛的經歷一樣.
去經歷這個生命.
捉緊著上帝的安全感.

$^{1081}$成為我們真正的避難所.
讓上帝的平安取代我們的擔心.
讓上帝而來的盼望取代我們的焦慮.
讓恩著信靠上帝取代我們的牽掛.
我們一起祈禱.
親愛的天父.
我們謝謝你.
我們知道此時此刻.
這裡你就與我們同在.
就在我們身邊.
我們要將我們的生命.
完全交託在你主的手中.
讓我們每一刻都更加意識到.
我們今天擁有的.
我們今天所作的.
我們接下來的計劃.
都有你參與當中.
所以我們仰望上帝.
你成為我們的保障.
成為我們的盾牌.
成為我們的避難所.
成為我們最大的保護者.
以至當我們在這種狀態之下.
恢復我們生命的那份喜樂.
幫助我們活在你的恩典和憐憫當中.
幫助我們在困難裡面.
仍然定精順靠你.
我們祈禱奉主耶穌基督的名字而求.
Amen.
\newpage



\section{馬太福音 28:16-20}
\label{sec:ltS5lDDLFjU}
\textbf{2026年4月5日復活節暨主餐崇拜直播｜陳明泉牧師｜告別浮淺,門徒進深:疑惑裡的差遣｜馬太福音28\_16-20}
\newline
\newline
連結: \href{https://youtube.com/watch?v=ltS5lDDLFjU}{\texttt{https://youtube.com/watch?v=ltS5lDDLFjU}} ~~~~ 語音日期: 2026-04-05
\newline
\newline
\hyperref[sec:3zRaVYVNQY0]{\small{< < < PREV SERMON < < <}}
~
\hyperref[sec:index_chronic]{\small{[返順時目]}}
~
\hyperref[sec:rsvExPNsP2I]{\small{> > > NEXT SERMON > > >}}
\newline
\newline
馬太福音 28:16-20
\newline
\begin{longtable}{cl}
\hline
\hline
章節 & 經文 (和合本修訂版)\\
\hline
28:16 & \begin{tabularx}{0.7\textwidth}{X} 十一個門徒往加利利去，到了耶穌指定他們去的山上。 \end{tabularx} \\ \\ \relax
28:17 & \begin{tabularx}{0.7\textwidth}{X} 他們見了耶穌就拜他，然而還有人疑惑。 \end{tabularx} \\ \\ \relax
28:18 & \begin{tabularx}{0.7\textwidth}{X} 耶穌進前來，對他們說：「天上地下所有的權柄都賜給我了。 \end{tabularx} \\ \\ \relax
28:19 & \begin{tabularx}{0.7\textwidth}{X} 所以，你們要去，使萬民作我的門徒，奉父、子、聖靈的名給他們施洗， \end{tabularx} \\ \\ \relax
28:20 & \begin{tabularx}{0.7\textwidth}{X} 凡我所吩咐你們的，都教導他們遵守。看哪，我天天與你們同在，直到世代的終結。」 \end{tabularx} \\ \\
[1ex]
\hline
\hline
\end{longtable}
$^{1}$今天的經訓是記載在新約聖經.
《馬太福音》28章16-20節.
十一個門徒往加里里去.
到了耶穌約定的山上.
他們見了耶穌就拜祂.
然而還有人疑惑.
耶穌進前來對他們說.
天上地下所有的權柄都賜給我了.
所以你們要去.
使萬民作我的門徒.
奉父子聖靈的名.
給他們施浸.
凡我所吩咐你們的.
都教訓他們遵守.
我就常與你們同在.
直到世界的末了.
弟兄姊妹平安.
平安.
今天是復活節.
我們帶著感恩.
帶著歡喜雀躍的心一起來慶祝耶穌基督.
從死裡復活.
這一份快樂.
不僅是耶穌從十字架受苦脫離.
更加是對我們每一個相信耶穌基督的基督徒.
每一個門徒.
是我們生命裡最大的信息.
因為我們知道.
我們所信的主連死亡都可以戰勝.
我們生命裡沒有一件事可以難了我們.
我們借著這個救拔我們的主.
我們人生就可以開展新的一頁.
今天我們和大家去看的這段經文.
也是開展我們接下來這個新的講道系列.
今天我們所看的大使命這段經文.
其實對大家來說都一點都不陌生.
不過過往裡面我們去看這段經文.
我們只是集中看後面那兩節.
所以我們有時候對於整個聖經裡面.
呈現出來的耶穌復活的圖畫.

$^{41}$我們有一點偏頗.
希望今天我們可以比較仔細一點.
來重溫耶穌基督如何復活之後.
第一樣對門徒所吩咐的東西.
成為了我們今天的激勵.
也是在我們讀這段經文的時候.
更加提升我們.
更加進入門徒的心靈當中.
以致我們能夠更加明白上帝對這個世代.
對我們每一個弟兄姊妹的生命的計劃.
我很記得有一個很喜愛讀的神學家.
叫做Dallas Willard.
中文的譯名叫做魏諾德.
他在說今天作為門徒.
或者基督教裡面最大的敵人.
不是世俗的敗壞.
這個世界的敗壞不是教會裡面.
要面對真正的敵人.
他說到今天真正我們要注意的.
就是我們基督徒的生命.
或者我們作為一個門徒.
我們的生命是否面對著一種浮淺的狀態.
所以他有一句名言.
教會最大的危機不是世界敗壞.
而是屬靈的膚淺.
我們在今天過的是一個有深度的門徒信仰.
又或者是一個很浮淺的.
他的深度 說的是浮淺和有深度.
其實是界定於生活的行為.
和我們跟著基督.
如果我們繼續用浮淺的方式.
來經歷我們的信仰生活.
我們的生命是沒有影響力的.
希望接下來這幾次.
我們這個季度.
我們不同同工一起來盡心.
挑戰一下我們的內心.
也重新去思考基督耶穌的生命.
讓我們重新立志要跟隨他.
在這個時間我們一起有祈禱.

$^{81}$求上帝幫助我們.
求主你與我們同在.
讓我們能夠明白你的話語.
藉著你對我們所宣講的.
讓我們每一個弟兄姊妹.
我們的生命更加踏實地跟隨你.
幫助宣講的.
讓你的話語.
衷心地向弟兄姊妹解明.
讓眾人得著好處.
在牧羊當中生命得著成長.
又為眾弟兄姊妹交託給你.
願你讓他們不單止聽得明白.
讓他們明白過後.
讓他們每一個在生活當中.
緊緊地跟隨耶穌基督的腳蹤.
讓我們整個群體.
成為一個耶穌所喜悅的群體.
恭敬將來的時間交託給你.
願你賜福.
獻上我們的禱告.
奉靠基督耶穌得勝的聖名.
Amen.
我們一起看看聖經.
我們看回我們最耳熟能詳的大使命.
在第十九節至到第二十節.
我們不會讀了.
所以你們要去.
使萬人作我的門徒.
奉父子聖靈的名.
給他們施浸.
我所吩咐你們的.
都教訓他們遵守.
在這兩節經文裡面.
我們過往很多的前輩.
或者傳統裡面.
我們很強調整個經文裡面.
那個最重要的動詞就是.
所以你們要去.
我們覺得那個去字.

$^{121}$出去.
就是一個最重要的動詞.
不過在這段聖經裡面.
其實那個語法.
那個句式.
那個命令語氣.
最重要的那個字眼.
是使萬民作我的門徒.
這個動詞.
是整句說話裡面.
一個叫做最重要的命令語氣.
而其他幾個動詞.
包括去 另外是施浸.
教訓 遵守.
這幾個其他的動詞.
它是在那個使萬民作我門徒.
這個大命令之下.
所作的那些實踐.
換句話來講.
耶穌基督頒佈的這個大使命.
他要門徒最重要最重要的核心.
不是就是蒙頭烏形的出去.
或者去做一些禮儀.
又或者不是很說教.
好像法利賽人那樣講講講教訓.
最重要的就是.
要讓到這個世界有更加多人.
成為基督耶穌的門徒.
這十一個門徒.
他就肩擔住這個最重要的命令.
耶穌就差遣他們.
去到萬民當中.
去到這個世界裡面.
就是要建立更加多的門徒.
而去施浸 教訓 遵守.
就是朝著叫人作門徒的這個核心.
一齊的來去匯聚.
以致到一個弟兄姊妹的生命.
就可以得著牧羊.
得著來去教導.

$^{161}$讓他成為一個上帝所使用的門徒.
今天我們有一個很重要的誤點.
我們覺得整個福音.
我們講的福音信息.
我們只是單從是由人從罪惡裡面.
來去救拔出來.
這個是福音的重點.
是真的.
不過福音不單止是叫人從罪裡面.
脫離罪惡得生命.
福音其實除了處理人的罪之外.
更加重要的就是經歷上帝的封城.
經歷天國.
藉著人的生命.
將上帝的封城帶到在人群當中.
所以你留意一下.
《馬太福音》裡面特別講到耶穌基督.
祂所講的那個是天國的福音.
祂不單止講人的罪得到赦免.
人的罪被赦免之後.
人怎樣去活出一個封城的生命呢?.
這個是耶穌基督那個重要的關懷.
福音不單止是講拯救人的罪.
福音更加是叫人得著封城的生命.
如果沒有了第一步.
就走不到第二步.
但是如果只是走第一步叫人悔罪.
你都未有真正去經歷到這個榮耀的福音.
所以坊間裡面有時候會將信徒分成兩個等級.
第一個等級就叫基督徒.
基督徒的意思就是說.
我認信 我認罪了.
耶穌救贖了我.
所以我就成為一個基督徒.
然後聖女一些的.
我要盡心一些 我要追求多一些的.
參與門徒訓練.
所以我們就成為門徒了.
不過在聖經裡面不是這樣分的.
在聖經裡面從來都沒有講過.

$^{201}$就是說你信主就可以了.
就保持著這個狀態.
你有心你就再追求就成為門徒.
今天的課程可能是這樣.
不過在聖經裡面從來講.
上帝所呼召的或者使徒他們所承擔的.
就是驕慢民作 耶穌基督的門徒.
我們每一個不單止被耶穌基督拯救.
更加是被耶穌訓練成為跟隨祂的人.
在生活各方面裡面.
我們都能夠活出一個上帝所喜悅的樣式.
更加大的挑戰是.
耶穌在這裡面去吩咐那些門徒.
凡我所教導你們的都教訓他們遵守.
就是說沒有選擇性的.
耶穌教什麼.
那些門徒如實地將所有東西都要教給其他人.
他不是純粹教的.
還要叫他們生命裡面落實來實踐.
今天有時候我們會覺得.
這些教導是好的.
但叫到你實踐的時候.
我們聽的時候我們的心裡面是雀躍的.
但到了生活裡面.
我們就沒有辦法貫徹成為我們生命的一部分.
門徒的生命就是講到不是選擇性地看聖經.
或者覺得這件事適合我聽.
我們就實踐多一點.
有些東西很難我們就放在一邊.
而是我們勉勵用我們的人生.
將我們整個生命交付在上帝的手裡面.
覺得困難的地方我們更加要鍛鍊自己.
以致我們行在上帝的旨意當中.
所以第一個重點.
我們看見對於大使命.
其實最重點的就是叫人成為門徒.
亦都我們每一個弟兄姊妹.
我們領受了的不是一個分了階層的福音.
當我們既稱耶穌基督為主的時候.
我們每一個都要成為一個跟從主的門徒.

$^{241}$經文裡面我們又再回頭去看.
我們明白了這一句之後.
我們再看回十六節到第十七節.
這兩節經文我們平時比較少讀.
十一個門徒往加利利去.
到了耶穌指定他們去的山上.
他們見了耶穌就拜他.
然而還有人疑惑.
在這兩節經文裡面描述了.
這個門徒他帶著耶穌的約定.
去到加利利和耶穌再遇.
值得留意的就是加利利這個地方.
這個加利利這個地方在《馬太福音》裡面.
或者在福音書裡面講的是一個很重要的地方.
為什麼這麼重要呢?.
在第四章十八節《馬太福音》.
它講到這是耶穌呼召門徒的起點.
耶穌就在那裡焦聚他所選召的門徒.
加利利的海岸邊有很多的暴雨.
他只是挑選耶穌心儀的.
和耶穌所揀選的那幾個門徒.
這個是成為了他們的起點.
然後耶穌全道的生活除了圍繞著加利利.
然後再走遍各城各鄉之後.
到了耶穌他預言自己將要受害.
《馬太福音》第二十六章三十六節裡面.
講到耶穌和門徒守完這個雨節晚餐之後.
耶穌就這樣講.
那時耶穌對他們說.
今夜你們為我的緣故都要跌倒.
因為經上記著說.
我要擊打牧人 羊就分散了.
但我復活以後要在你們以先往加利利去.
耶穌強調他將要在雨節晚餐之後.
講到他自己將要面對著被擊打.
那些門徒會四散的.
不過他講到在受害之後.
耶穌說他復活之後.
他要讓那些門徒先去加利利.
他在加利利那裡等那些門徒.

$^{281}$這個是耶穌對門徒的約定.
在那裡來呼召他.
亦都講到復活之後就在加利利那裡.
我們會合了.
當耶穌復活的時候.
28章第7節裡面.
空墳墓的婦女去到空墳墓裡面找耶穌.
那些婦女找不到耶穌的屍體.
因為耶穌已經復活了.
28章第7節裡面.
那個天使向那班婦女顯現.
他特別的來告訴那班婦女.
耶穌基督已經復活了.
並且你們先往加利利去.
在那裡你們要見他.
我已經告訴你們了.
天使將耶穌之前和門徒的約定.
再一次告訴那班婦女.
藉著那班婦女通知那班門徒.
你要記得這是耶穌復活之後與門徒的約定.
所以加利利這個地方是一個約定.
然後耶穌都害怕.
那班婦女就擔心驚惶失措.
耶穌再一次顯現在那些婦女當中.
上一次是天使.
現在就是耶穌顯現28章第10節馬太福音.
經文裡面耶穌對他們說.
不要害怕.
你們去告訴我的弟兄.
叫他們往加利利去.
在那裡必見我.
你看到加利利這個地方.
是一個很重要領受上帝使命.
是耶穌和那班門徒的約定.
即使經歷耶穌十字架上的門徒有多麼軟弱.
但是耶穌約定了他們.
一定在那裡見面.
他復活過後就在那裡約定他.
如果你作為門徒你會不會去呢?.
兩天之前可能還是很軟弱.

$^{321}$還要出賣耶穌.
甚至乎在耶穌被捉拿被釘的時候.
這班門徒四散.
其實他們面對著信心最軟弱的情景.
如果這班婦女提醒他.
耶穌在那裡等你.
又或者耶穌再三的來講這個重要的約定.
這班人他會不會應約呢?.
如果對於我們今天來說.
我們覺得很尷尬.
又或者是不是真的.
又或者覺得自己很卑微.
不配.
甚至乎對於耶穌復活的訊息.
自己都好像有點害怕.
難以接受的一個訊息.
耶穌在講這個約定是不是真的成就呢?.
對於很多人來說.
可能都是一個很大的疑問.
對於那些門徒來說.
當這些婦女通知他的時候.
在經文裡面講到.
帶著懷疑的態度.
帶著忐忑不安的心.
這班門徒到了最後是往加里尼去.
到耶穌指定會合的那個山上面.
當這班門徒見到耶穌的時候.
就俯伏來去拜祂.
然而經文裡面講.
「還有人疑惑」.
「還有人」是誰呢?.
不知道.
又或者「還有人」有很多的意思.
「還有人疑惑」的意思就是說.
那些十一個門徒每個人都帶著疑惑.
去到上帝面前.
就是說那個疑惑的程度.
每個人都有.
或者是有些人充滿信心.
帶著無比的決心來見耶穌.

$^{361}$但有一兩個可能就有點害怕不去.
我比較接納的解釋是.
十一個門徒每個人都帶著一份疑惑.
去到加里尼來會合耶穌.
不過眼前見到的這個耶穌.
已經是一個復活了的耶穌.
祂的身體完全可見.
所以那個疑惑.
祂是不是復活了.
那個疑惑已經解除了.
不過聖經裡面講.
這班人還繼續帶著一份疑惑的心.
但同時間這班門徒也見到耶穌.
就去俯伏敬拜祂.
敬拜祂所眼見的主.
但仍然帶著疑惑的態度.
這個就是馬太福音.
特意製造一個這樣的章歷.
告訴讀聖經的人知道.
其實這個也是我們生命的寫照.
疑惑和敬拜一起出現.
不只是這裡.
其實在馬太福音裡面.
一早已經埋下伏線.
這兩個字或者這樣的句式.
一早已經出現.
在馬太福音第十四章.
三十一節至三十三節.
我沒有印給大家.
我告訴你那個是什麼情景.
大家很熟悉的.
這個就是講到耶穌基督.
和那些門徒一起.
來到平靜風和海的那段.
那段經文講到.
耶穌和那些門徒一起出海.
夜裡四更天的時候.
耶穌在海面上行走.
去門徒那裡.
門徒見到耶穌在水面上走的時候.

$^{401}$就很驚慌.
說是一個鬼怪來的.
就害怕就喊叫起來.
耶穌連忙對他們說.
你們放心.
是我不要怕.
彼得說.
主如果是你的.
請你叫我從水上走到你那裡去.
耶穌說.
你來吧.
彼得就從船上下去.
在水面上走.
要到耶穌那裡去.
你記得這一段經文.
彼得問耶穌可不可以在水上走.
耶穌說可以.
你過來吧.
彼得就走了.
最重要的.
我想講的那段經文是.
只因見風甚大.
就害怕.
將要沉下去.
便喊著說.
主呀救我.
耶穌趕緊伸手拉住他說.
你這小信的人為什麼疑惑呢.
是同一個字.
他們上了船.
風就住了.
在船上的人都敬拜他.
說.
你真是神的兒子.
既疑惑.
又認定耶穌基督是神的兒子.
在《馬太福音》裡面.
他講到.
原來門徒的生命就是這樣的一種生命.
如果你從這個角度來講.

$^{441}$他們不會疑惑耶穌是否復活了.
他不是在講耶穌復活的事蹟的那種疑惑.
疑惑的是指向自己的生命.
彼得走這個水面的時候.
是不是可以的.
看到那些風浪來到.
自己害怕.
然後就沉下去.
所以耶穌說你為什麼要疑惑呢.
那種不信得過上帝幫助他.
或者讓他經歷這個神蹟的那種疑惑.
去到在這個大使命裡面.
復活的處顯現在這班門徒當中.
這班人心裡面帶著一份疑惑.
這份疑惑不是指向耶穌基督復活的事實.
而是指向他們自己的心靈深處.
他還有沒有面子去接受耶穌基督的差遣.
眼前的這個夫子.
是他們跟足三年的一個夫子.
但是看回他們心靈的那種狀況.
他們每一個都好像虧負了這個夫子一樣.
在他的生命當中.
好像留下一個很重大的污點.
兩日前他見到耶穌這樣的狀況.
三日之後耶穌復活了.
好像沒有面子見耶穌一樣.
這個疑惑是指向我們心靈裡面的那種虧欠感.
那種不足感.
甚至乎覺得自己虧負上帝.
馬太福音裡面特別講到這個描述.
他講到那些門徒.
雖然在馬太福音流傳的時代.
那些門徒已經是一班領袖.
是教會的領袖.
不過馬太特別加入這個註腳.
講到他們疑惑.
其實講到這一班的門徒.
不是想像中之前那麼威猛.
他們心靈深處裡面都是帶著那份愧疚.
有一種質疑.

$^{481}$以致到他們覺得是不是真的繼續這樣走下去.
但是耶穌沒有因為這一班人的不足.
這一班門徒的軟弱.
耶穌說這班婦女比你還厲害.
將大使命交給那些婦女.
就不交給那些門徒.
撕掉他們那幾個門徒的柴.
沒有.
耶穌知道他們的軟弱.
但是都依然將這個職事.
將這個大使命交付給他.
經文第十八節後半段繼續這樣寫.
耶穌進前來對他們說.
天上地下所有的權柄都賜給我了.
所以你們要去.
使萬民作我的門徒.
奉父子聖靈的名.
給他們施浸.
凡我所吩咐你們的.
都教訓他們遵守.
我就賞與你們同在.
直到世界的末了.
在這裡既然門徒有一種疑惑.
耶穌怎樣去回應他的疑惑.
他怎樣去面對那些門徒覺得自己不行.
耶穌沒有特別去說.
不是你可以的.
不是不停地說這些安慰的話.
某程度上耶穌繼續說.
天上地下所有的權柄都賜給我了.
耶穌用他將他自己的權柄提升.
他沒有否定那些門徒.
自己心靈軟弱的那些過去.
不過這個大使命.
從來不是說因為你很厲害.
所以我交給你.
就是正正因為你不夠厲害.
上帝有天上地下的權柄.
所以上帝可以.
你不行不要緊.

$^{521}$但是我不會因為這樣而放棄了你.
反而將一個上帝同在.
以及滿有能力的差遣.
賦予給這一班的門徒.
所以這一班門徒.
他既是看到自己生命的那種不濟.
他的疑惑是對自己可能有一些的質疑.
但同時間耶穌沒有質疑他們.
反而是說用他的權柄.
補足了人生命裡面的不足.
補足了生命裡面.
我們覺得不理想的那個地方.
這一班的門徒不是一班所謂的信心強人.
在經歷耶穌受難之後的復活.
這班人和我們一樣都是血肉之軀.
我們不明白.
我們也沒有那種膽識.
甚至乎我們都是軟弱不堪.
在人生裡面.
可能我們都曾經好像這些門徒一樣.
有一些時間我們好像背棄了耶穌一樣.
甚至乎有時候我們口裡面說到自己是很厲害的.
但真實的自己其實是軟弱不堪的一個人.
曾經有些弟兄姊妹和我說.
牧師你不要交付這些責任給我.
有時候我們很喜歡.
不是很喜歡.
其實我們很想動員多些弟兄姊妹參與服事.
但當我們動員弟兄姊妹的時候.
有些弟兄姊妹很多時候會推我們.
為什麼呢.
牧師這件事我不做了.
因為我覺得我自己不配.
我覺得我自己爛身爛勢.
當想起做一些這麼重要的事工.
或者肩負起一些這麼重要的責任.
我看到我自己是一個軟弱不堪的人.
又或者在信仰裡面我們會挑戰這些弟兄姊妹.
你走前多一步.
你承擔多一些.

$^{561}$我們既肯定她的能力.
同時間我也知道她的心氣願意跟隨主.
但那些弟兄姊妹想一想就會放棄.
因為她就好像那些門徒一樣.
懷疑自己的能力.
亦都看回自己的chat record.
可能很多地方都是很不濟.
我很記得我自己都經歷過這樣的狀況.
作為牧者.
我不是說我自己都是很厲害的人.
我自己都有很多很軟弱的地方.
生命裡面都是千瘡百孔.
我記得有一個這樣的片段.
是很不樂意的.
話說在我讀神學的時候.
在我讀神學的時候.
神學院的經濟有一段時間不太好.
一般來說.
你可以理解.
如果神學院的經濟不太好.
通常籌款的是哪些人呢.
就是院長 老師.
或者有些拓展部.
那次神學院就交了一個榮耀的職事給我.
他說你是學生會的.
不如你組織一下同學一起籌款.
我說讀書要做籌款的.
我那時候都覺得有一種詫異.
他說你動員同學一起來做.
他不是說服我的.
而是困難.
當我聽到這個說法的時候.
心裡面就覺得很怪.
慢慢地想答應 要想一下.
不過盡力而為.
心裡面有一點不心甘情願.
心裡面有一個卡住.
我很記得當我心裡面有這個卡住的時候.
那天正在上一個去年已故的牧者.
曾立華牧師上他的課.

$^{601}$他說牧者今天或者我們作門徒的.
基本上不是做每一樣東西都是你喜歡的.
你說你自己是蒙召讀神學也好.
出來牧會你都有很多地方.
不是做你自己喜歡的東西.
總會有一些東西你會打一個突.
或者有些東西不是心甘情願.
不是很樂意去做的.
不過面對這樣的情景 他說了一句很重要的話.
他說眾人要為美的事要留心去做.
然後他再說 曾牧師他說.
他說神學院事奉了幾十年.
他說有些崗位他是很喜歡做的.
但是有一些又要籌款.
又要走去實習.
很微細的那些東西.
他不喜歡的性格不是這樣的.
但是眾人都要為美的事.
他就逼自己 勉勵自己來做.
他說自己有時候都不是很有恩賜的.
他就教我們 你沒有恩賜不要緊.
若望賜恩的主 上帝就是這樣幫你.
今天我們很多時候.
我們經常覺得我疑惑 我不好.
我就推了 我不做了.
我安安分分做回我的基督徒就OK了.
不過耶穌基督從來都不是想.
我們的信仰生命是一個這樣的狀態.
我們拿了得救的入場券.
我們上天堂穩入了.
我們就繼續過我們自己的信仰生活.
耶穌基督正在挑戰我們地上的每一個.
不單止信祂幫我們從罪惡裡面拯救出來.
祂更加呼召我們成為這個世代的門徒.
成為門徒不代表做每一樣事都是那麼純心.
那麼喜歡 那麼好玩 那麼有刺激性.
總會有一些是令到你要學的.
你要祈禱的.
你要在上帝面前不斷又不斷立志的.
這一班門徒面對著耶穌基督頒佈下來的大使命.

$^{641}$大概都可以說我們生命不濟.
我們不配領受這麼榮耀的即時.
你叫那個莫大拉瑪利亞做吧.
不是 聖經說到這個大使命既然是命令如戲.
我們沒得避的.
我們知道自己的內心.
我們就祈禱吧.
在上帝面前有足夠的勇氣.
承擔上帝要你 要我一起成就的那個即時.
未必每個弟兄姊妹一定要在教會裡面有個崗位.
那個才叫門徒生命.
可能在你的職場裡面.
在你的家庭當中.
在你的生活各方面.
上帝呼召你.
你進入在那個真實當中.
你體會到那個需要.
你心裡面有感動.
你知道你不是最完美的一個人.
但上帝將這個處境放在你的眼前.
他在呼籲你.
一起來努力.
門徒的疑惑在聖經裡面的記載.
他不是要合理化這些門徒.
人是這麼軟弱的.
或者人就是這樣.
他是想說即使人這麼軟弱也好.
上帝這個即時都是交付在每一個領受恩典的基督徒當中.
有些人說基督教不可信的.
那個教會亂七八糟 搞成這樣.
那些人又這樣.
有些人是這樣說.
不過你可以得意一點說.
即使亂七八糟.
即使我們生命裡面破破爛爛.
上帝都將這個即時交給教會.
將這個榮耀的福音.
藉著教會來傳到這個世界當中.
兄弟姐妹.
不要輕看上帝給你門徒這個身份.

$^{681}$不要覺得生命裡面有些不理想的地方.
你就停留在某一個地步.
這個榮耀即時其實是鼓勵我們繼續前進.
這個大使命成為了我們門徒的命令式.
我們跟隨它.
當我們的生命真的不夠能力的時候.
我們就仰望這個次因的主.
耶穌基督已經復活了.
他復活第一個交付給門徒的.
就是這個差遣的大使命.
今日當我們慶祝復活節.
我們不要忘記.
上帝同樣將這個即時交給我們眾人.
求主繼續帶領我們.
\newpage



\section{馬可福音 1:35-39}
\label{sec:rsvExPNsP2I}
\textbf{2026年4月12日崇拜直播｜何梓峰傳道｜告別浮淺,門徒進深:退到曠野為要看到你｜馬可福音1\_35-39}
\newline
\newline
連結: \href{https://youtube.com/watch?v=rsvExPNsP2I}{\texttt{https://youtube.com/watch?v=rsvExPNsP2I}} ~~~~ 語音日期: 2026-04-12
\newline
\newline
\hyperref[sec:ltS5lDDLFjU]{\small{< < < PREV SERMON < < <}}
~
\hyperref[sec:index_chronic]{\small{[返順時目]}}
~
\hyperref[sec:Xc4fVvDZWRE]{\small{> > > NEXT SERMON > > >}}
\newline
\newline
馬可福音 1:35-39
\newline
\begin{longtable}{cl}
\hline
\hline
章節 & 經文 (和合本修訂版)\\
\hline
1:35 & \begin{tabularx}{0.7\textwidth}{X} 次日早晨，天未亮的時候，耶穌起來，到曠野地方去，在那裡禱告。 \end{tabularx} \\ \\ \relax
1:36 & \begin{tabularx}{0.7\textwidth}{X} 西門和同伴出去找他， \end{tabularx} \\ \\ \relax
1:37 & \begin{tabularx}{0.7\textwidth}{X} 找到了就對他說：「眾人都在找你！」 \end{tabularx} \\ \\ \relax
1:38 & \begin{tabularx}{0.7\textwidth}{X} 耶穌對他們說：「讓我們往別處去，到鄰近的鄉村，我也好在那裡傳道，因為我是為這事出來的。」 \end{tabularx} \\ \\ \relax
1:39 & \begin{tabularx}{0.7\textwidth}{X} 於是他走遍全加利利，在他們的會堂傳道，並且趕鬼。 \end{tabularx} \\ \\
[1ex]
\hline
\hline
\end{longtable}
$^{1}$今日的經訓是《馬可福音》第一章三十五至三十九節.
聖經新約四十八頁.
「次日早晨,天未亮的時候,耶穌起來到曠野地方去,在那裡禱告..
西門和同伴追了他去,遇見了就對他說:「眾人都找你.」.
耶穌對他們說:「我們可以往別處去,到鄰近的鄉村,我也好在那裡傳道,.
因為我是為這事出來的,於是在加利利全地盡了會堂,傳道講鬼.」.
弟兄姊妹早晨.
早晨.
上星期陳牧師引用了耶穌給門徒的大使命.
提到我們在事奉中如果遇到困惑的時候.
看不看到神的同在,神的恩典.
其中他用到《馬太福音》第28章那裡說到的大使命.
大家如果有印象的話,第20節耶穌是這樣說的.
「凡我所吩咐你們的,都教他們遵守,我就常與你們同在,直到世界的末了.」.
耶穌親自去到應許,祂會常與我們同在.
但是問心那句,如果我今天去問大家.
你和神的關係,十分的話,你會給自己多少分呢?.
有些人可能是給三分,有些五分,有些七分.
有不同的分數.
有時候我們在現實裡面,我們會不會看到一個狀況.
就是事奉得越多,心裡面和神的關係好像越來越遠.
反而沒有一個更加親近的感覺.
不論你是做師士也好,師班,團體的導師.
甚至是帶敬拜的,又或者今天回來崇拜大家一起.
在當中敬拜的事奉者,即是大家.
我們回教會,會不會變成了一個義務,或者一種責任呢?.
今天我們會一起去看《馬學福音》的一章35至39節.
我們拆解一下耶穌事奉的源頭.
也讓我們一起思考一下,如何在一個忙碌的事奉生活裡面.
我們真真正正地去到經歷和耶穌的同行.
我們一起去祈禱吧.
求主你開我們的心,開我們的耳朵.
讓你的福音流到我們心中,讓你的教訓親自臨到我們當中.
讓我們今天都可以看清楚事奉的源頭.
也可以在當中更加火熱地去事奉你.
我們這樣去不完全禱告,以上禱告奉主耶穌的聖名祈求.
好,我們一起去看看今天這一段的經文.
第35節他這樣去提到就是次日早晨,天未亮的時候.
這個時間大概就是凌晨三至六點的時候.
有沒有人凌晨三至六點已經起身的?正常,有人舉手.

$^{41}$我今天問何先生的時候,他都說三四點就起身了.
原來都有人是三四點起身的.
經文描述到耶穌在這個時候去曠野去祈禱.
按照道理來說,耶穌是應該覺得很累的.
因為在前一天裡面,耶穌在加伯隆這個地方做了很多不同的事奉.
由早上在會堂開始去教訓人,然後就醫治好西門岳母的業病.
然後晚上還要醫治好群眾的病,甚至做一些講鬼的工作.
經歷了一整天這麼高強度,這麼激烈的一個事奉的話.
如果換成是你,你都應該會很想休息多一點,或者睡多一點.
即是給我起碼睡到明天十二點.
但是耶穌的選擇,祂不是選擇一個充分的休息.
而是選擇去曠野.
這裡我們要去問一個問題,就是為什麼是曠野呢?.
為什麼不在會堂裡面,為什麼不在家裡,為什麼不在家裡外面.
為什麼不在廁所呢?.
曠野,你知道的,什麼都沒有的,沒有人,沒有水,甚至沒有床.
但正正是因為什麼都沒有的時候,才沒有人能夠去打擾祂.
連門徒都要追著祂去找祂.
這裡的追是一個很主動的,甚至是可以翻譯做追捕的意思.
即是門徒在追捕耶穌,到處找耶穌.
是一個很主動的追的動詞.
我們可以見得到,原來耶穌真的將祂和天父的時間.
看得比任何的東西都更加重要.
曠野在聖經裡面有很多不同的意思.
它不單只是代表了一個試煉或者試探的地方.
我們經常聽到耶穌在曠野受到40天的試探.
它也是一個耶穌去迴避群眾或者休息的地方.
曠野也是象徵著和上帝的同行.
就好像在舊約裡面.
我們會見到摩西在曠野那裡燃燒的荊棘.
我們會見到以利亞在曠野那裡遇見神.
我們會見到施洗約翰.
在曠野裡面是怎樣去到預備主的道路.
怎樣去叫人去悔改.
曠野就是一個和神彼此結連.
彼此去到建立關係的一個場景.
而原來透過祈禱,透過安靜.
耶穌和神有一個連結.
並且知道天父的心意在哪裡.
所以當門徒找到耶穌的時候.

$^{81}$就跟祂說「大家都在找你,你去了哪裡?」.
但是耶穌的回答是出乎每一個人的意料之外的.
祂說「我們可以往別處去,到鄰近的昌村.
我也好在那裡傳到,因為我是為這事出來的」.
耶穌知道,清楚地明白到.
祂的使命是來自於天父.
就是將福音傳遍.
但是弟兄姊妹,有沒有想過.
如果今天換成是你的話.
你會不會跟耶穌做出一個同樣的選擇呢?.
我嘗試將這段經文改寫.
不知道你看完之後會有什麼感覺.
讀一次給大家聽.
「次日早晨,天未亮的時候.
耶穌起身,就為醫治和講鬼作預備.
耶穌和同伴見到祂的時候就說.
「眾人都找你」.
那耶穌就跟他們說.
「好了,新的一天的事奉又開始了」.
如果是這樣的情況.
你會不會更加欣賞這個勤力的耶穌呢?.
有人點頭,有人搖頭.
耶穌沒有浪費任何的時間.
祂幫助到更多的人.
祂將群眾的需要放在自己的心裡面.
這個耶穌值得大家效法嗎?.
點頭,搖頭.
這個是我們今時今日很多事奉者的寫照.
事奉基本上不會有停的一天.
就好像你這個星期大小早.
下個星期又是到你大.
然後下個星期不是我大.
但下個星期又是我大.
或者練完敬拜之後.
你又要犯罪.
然後又要練敬拜.
然後就可以上台表演.
上台敬拜.
然後又要再練個敬拜.
甚至有時候小組也會聽到一些聲音.

$^{121}$「為什麼這麼快又輪到我們小組做私事?」.
「剛剛才做完沒多久又到我們小組做?」.
我曾經也有一段這樣去到服事的時間.
不是現在服事,不是現在教會服事.
讀神學之前的服事就是.
星期三晚上就會大祈禱會.
然後星期四就會上一些聖經的課.
星期五又是另一個祈禱會.
然後星期六就是自己的小組.
星期日就是一個中學的團契.
我的生活就是預備,事奉,預備,事奉.
預備,事奉,預備,事奉,預備,事奉.
聽上去是不是很正,很火熱?.
是不是很正?.
可以為神工作了這麼多.
可以為神這麼努力的時候.
你覺得是不是很正,很火熱?.
這麼忙碌的事奉生活.
是不是代表著他的屬靈生命很好?.
他參加了很多不同的服事.
他的出席率很高.
他的屬靈生命是不是很好?.
其實未必的.
在這段經文裡面.
有一樣東西是沒有描述到的.
就是事奉的源頭.
事奉的源頭.
就是說有什麼是推動到我們去事奉呢?.
耶穌在加拔隆裡面大受歡迎.
其實他大可以留在這個地方.
繼續在這裡服事.
每個人都需要耶穌,是不是?.
但是耶穌沒有這樣去做.
耶穌知道他事奉的源頭是上帝.
他不是要滿足當地人的需要.
就是說耶穌知道他離開.
是會令一班的群眾.
或者門徒失望或者埋怨.
但是耶穌仍然選擇離開.
就算耶穌他不在這裡服事.

$^{161}$他仍然在其他地方服事同樣的東西.
他仍然選擇離開.
因為他知道他要完成天父的託付.
事奉的源頭.
對一個事奉者來說是很重要的.
如果我們的事奉不是出自於神的話.
我們很容易會走錯了一個方向.
而當我們走錯了這個方向的時候.
事奉本身是一件很美好的事情.
都會變到成為和神的關係越來越遠.
是不是?.
本身我們可以造就到其他人.
本身甚至可以造就到自己.
但是就因為走錯方向.
和神的關係好像很平淡.
甚至是越走越遠.
今天我很想和大家分享六個事奉.
一個迷失.
當我們去到迷失的時候會怎樣呢?.
當我們去講這幾個迷失的時候.
你不需要去評價其他人.
就是說這個人一定是這樣.
事奉一定是這樣.
不需要.
我們去到反思一下自己的內心.
究竟是不是真的有這些感覺呢?.
第一就是責任感和義務的感覺.
有時候我會問一些導師或者組長.
為什麼你會做導師或者組長呢?.
他就會說因為沒有人做.
因為沒有人做.
因為我是這裡裡面最久的一個.
所以我就自薦而言.
就是成為了組長.
就是因為這份單純沒有人做.
或者因為自己在教會裡面留了很長的時間.
不好意思去到拒絕.
都是時候要拿回一些東西.
我都不可以一輩子都是做組員.
於是你就承擔了這些事工.

$^{201}$這些事奉對你來說.
可能更加像一份工作.
是基督徒的責任.
但是和神沒有什麼關係.
當生活已經很繁忙.
很忙碌的時候.
你還要加上事奉.
事奉對這個事奉者而言.
它是一個負擔.
一個包袱.
他不能夠很全然地投身在這個事奉當中.
就好像有時候我回教會.
你都說是基督徒的義務.
有時候還是一個對神的回應呢?.
第二就是被認同和被看見.
有時候我們都會聽到一些欣賞和讚賞.
就是有你在真的不一樣了.
這個事奉沒有了你不行.
有時候這些說話都會令人飄飄然.
我不是說事奉者不能獲得讚賞或者欣賞.
我通常都會鼓勵大家多點讚賞和欣賞.
因為這是對人一個肯定和鼓勵.
但是如果我們的事奉.
只是想獲得其他人的肯定和欣賞的時候.
這一種純粹是一種被需要的感覺.
或者是被肯定的感覺.
在事奉當中.
你留意的不是你事奉的對象.
而是你留意著自己有沒有被人讚賞.
有沒有被人看到.
就是有人讚賞你的時候.
你就會事奉得很開心.
如果有一天沒有人讚賞你.
你的心就會很失落.
很不開心.
事奉變成了餵養這種感覺的一種工具.
而不是用來榮耀神.
第三,一個的成就感.
有些人事奉是因為能夠在那裡獲得一個成功的快樂.
當你是一個很有能力的人的時候.

$^{241}$可能在社會也好,公司也好,學校也好.
你已經習慣了一個很高效率的運作模式.
你自然會將這個標準帶回教會當中.
你所追求的是一個很完美的敬拜.
一個很順暢的流程.
一個很華麗的演出.
甚至一些很好的行政操作.
如果事情辦得很順利的時候.
你就會很有滿足感.
但是一旦這個事奉出錯了.
或者是它有些差池的時候.
你就會覺得很焦急,很煩躁.
特別是很多有能力的人的時候.
他會將教會的事奉當作KPI.
他會做得好就開心,做不好就焦慮.
但是忘記了神是要與他同行.
而不是要將事情做得很漂亮.
第四,就是填補生活的空虛.
有時候一個人真的很悶,很寂寞.
有時候事奉就成為了一個這樣的靈性的逃脫.
就是當現實生活,家庭的關係.
或者工作裡面遇到一些困難的時候.
可能你會覺得很孤單,很窒息.
教會剛好有事奉.
於是你就在當中去服侍.
事奉成為了一個避風港.
在事奉裡面你仍然可以認識到很多朋友.
可以認識到很多不同的事務.
可以令你覺得生活好像很充實,很有意思.
但是這種忙碌其實是一種轉移的焦點.
是掩蓋了你內心的空虛.
你的事奉不是為了讓你覺得不孤單.
而是因為被神的愛而去激勵.
結果當你越服侍的時候.
你安靜的時間反而是越來越少.
第五就是一個交易的心態.
有些人事奉是出於一個不敢不做的心態.
就是說我作為基督徒.
如果我不事奉的時候.
我怕神會懲罰我.

$^{281}$或者是我越做多些事奉.
就是一個愛神的表現.
如果我不事奉就是一個不合格的基督徒.
就會失去了神的恩典.
這種事奉的感覺更加類似還債.
就是為了安撫那些不做會內疚的恐懼感.
在這些情況下.
事奉已經不單是和上帝沒有了交流.
並且只是一個指令,一個標準.
而和神那份關係完全脫節.
而去到最後就是位置和影響力.
在我們事奉的裡面.
我們有時也會沉溺在那種.
很棒,可以命令很多人.
可以發號施令的感覺.
但是在我們那種被尊重的感覺裡面.
我們被擁有話事權的感覺裡面.
我們就是那種控制慾就出來了.
就是當有人不同意見的時候.
你就會覺得好像被人冒犯了.
很不安的感覺.
而你不是在尋求上帝的旨意.
你只是在保護你自己的一個地盤.
在這種事奉的源頭下.
所操控你的就是那種權力慾.
你服侍的就是你的影響力.
而不是神.
神只是一個架空了的名義上的領導.
事實上,我剛剛所說的這些迷失.
其實是可以幫助到我們的.
一個欣賞,一個讚賞.
可以幫事奉者得到一些鼓勵.
一個有能力去影響人的人.
是能夠帶領一個群體.
是可以有更加好的一個事奉.
是不是?.
一個很專注在一個程序裡面的時候.
可以令整個事奉更加有系統.
更加有效率.
其實這些迷失.

$^{321}$都能夠成為我們事奉一個的動力.
一個的幫忙.
但是這裡裡面最重要的就是.
你的源頭究竟在哪裡.
如果我們的源頭都是在這裡的時候.
我們和神的關係又在哪裡呢?.
你看不看到神在當中呢?.
如果你今天事奉的裡面.
你覺得和神的關係好像越來越遠.
這個是一個警號.
這個亦是一個提醒.
去告訴你.
是時候要處理你和神的關係.
不要等到真的惡劣的時候.
情況很嚴重的時候.
才去處理.
如果到時才處理.
就真的很遲.
曾經我事奉的時候.
我在公司.
我和老闆很親密.
他經常都會和我說一個叫木桶理論的東西.
不知道大家有沒有聽過木桶理論.
木桶就是由很多的木條組成.
綁在一起又形成一個木桶.
他經常都會這樣說.
Toby,我們是一個團隊.
我們每個人就好像一塊木條.
怎樣才能災得多些水.
就是大家都要差不多的程度.
如果其中一條短了會怎樣.
那你的水只可以去到那個位置.
你不會再高.
即是其他人再高也好.
因為你的關係.
你都只會在那個水位.
他經常這樣形容我.
你不要做一條這麼短的木.
即是你不要拉低整個團隊的效率.
善意提醒.

$^{361}$我將這個東西放在我的事奉裡.
我會很逼迫.
那時候我是很逼迫我的組長.
為甚麼他會這樣做.
我們有沒有方法可以做好一點.
我們是一條團隊.
我會用很多不同的方法.
叫大家去事奉.
直到我發現有一刻是很不妥的.
就是大家是很枯乾.
他們不願意去事奉.
或者是需要休息幾個星期.
直到有一天在我靈修的時候.
神這樣跟我說.
這一間不是你的教會.
這一間是我的教會.
原來在教會的事奉裡.
原來在教會的當中.
我們都要明白到.
我們的源頭是來自於哪裡.
這裡是上帝的教會.
我們所事奉的.
都是因為我們和上帝的關係.
以致有一個回應.
如果事奉是為了叫人去看到上帝.
但是事奉者那個人.
看不到上帝的存在.
這是一件很諷刺的事情.
前幾天當我預備這篇道.
嘉威就問我.
這篇道裡最深感動你的是甚麼.
或者你最想說的是甚麼.
我就很想說很可惜.
事奉本身是一個愛教會的表現.
不愛教會.
你不會花時間在教會的事奉裡面.
你不會回來崇拜.
你不會回來事奉.
你就是因為喜歡這個地方.
喜歡這裡的人.

$^{401}$所以你們每一刻.
每一秒的投入.
都是代表你們有多愛教會.
你們是教會的每一個弟兄姊妹.
亦都是每一個成員.
但是這些迷失.
亦都是令事奉者透不過氣.
有人會著重這件事做得好不好.
有人會著重他自己的價值和能力.
有人著重他的話事權.
本應該是一個很喜樂的事奉.
反而是給身邊的事奉者帶來一種壓力.
這種的無奈.
這種的精神很累.
很疲倦.
甚至是大家互相埋怨.
大家互相對罵.
屬靈生命一塌糊塗.
有人因為這樣離開了教會.
這種的事奉源頭.
不單只是影響你自己.
亦都是影響你身邊的人.
聖經裡面有一對姐妹.
姐姐同樣是忙於事奉的裡面.
她看不到更加重要的事情.
這一對姐妹就是瑪大和瑪利亞.
在路加福音去到描述.
就是當姐姐瑪大因為她操持家務.
去招待客人而感到心靈很煩擾.
甚至和耶穌去到抱怨.
她妹妹瑪利亞不幫忙.
因為她妹妹當時安靜地.
坐在耶穌的腳前聽到.
這樣形成了一個很大的對比.
姐姐很忙碌.
而妹妹就很安靜地聽到.
耶穌這樣提醒瑪大.
這裡用了新譯本.
他說瑪大瑪大.
你為許多事操心忙碌.

$^{441}$但最需要的是只有一件事.
瑪利亞已經選擇了.
那個最上好的福分.
是不能從他奪去的.
今天的你是瑪大還是瑪利亞呢.
有時候我們都好像瑪大一樣.
忙碌在事奉的裡面.
但我們忘記了那個源頭.
以致我們在事奉的當中.
很忙亂很困惑.
但耶穌這裡說了一個.
最需要的一件事情.
就是安靜在神的裡面.
去聽他的說話.
和他建立關係.
我們能夠事奉是一種福氣.
是一種恩典.
但有一樣更加好的東西.
你是可以得到的.
就是和神的關係.
耶穌這樣去形容.
瑪利亞已經得到他最上好的福分.
英文是The Best Portion.
最好的那一部分.
就是和神的關係.
作為耶穌的門徒.
我們有時候都需要.
暫停一下手上的工作.
去看到神那份的同在.
路雲在建立生命即是的這本書裡.
他提到事奉和屬靈生命.
是不可以分割的.
即是兩件事是不可以分開.
事奉不是一種朝九晚五的工作.
而是本質上一個生活的方式.
在你生活裡面已經呈現了事奉.
是叫別人去到明白.
去到看得見.
是叫別人得釋放自由.
路雲形容.

$^{481}$我們今日人渴望一種很暫身的屬靈生命.
就是在生命裡面去到經歷上帝.
事奉者可能都會問.
Toby你叫我們休息.
有可能嗎.
事奉是沒有可能停的.
好像你剛剛所說.
就是因為教會沒有人做.
我才要去做.
我要怎樣才能夠找到這個源頭呢.
所以我們就是要在日常的生活裡面.
我們就是要花時間去找到上帝.
花時間去找到上帝.
耶穌已經為我們作出了一個示範.
就是找到你的曠野.
找到你和神的曠野.
按照馬可的描述.
耶穌不單止是在這一次裡面祈禱.
祂甚至是在五餅二魚.
餵飽了五千人之後.
祂也是辭別了門徒.
然後一個人在山上祈禱.
在赫西瑪利園那裡.
耶穌也是吩咐門徒去等祂.
然後祂一個人走前一點.
一個人去祈禱.
找到你和神的曠野.
是今天給大家的一個挑戰.
不單止是事奉者.
而是我們每一個門徒都很需要這樣去做.
但是我們要怎樣去找到一個曠野的時刻呢.
這個曠野時刻好聽.
第一就是我們要定一個曠野時刻.
就好像在耶穌那樣.
在一天裡面有沒有一個最安靜.
最不容易被人打擾的時間呢.
可能不用很長.
一開始五分鐘.
接著十分鐘.
接著十五二十分鐘.

$^{521}$你可以嘗試將你的電話調成飛行模式.
關掉所有的通知.
然後告訴自己.
這段時間只是屬於你和上帝.
只是屬於你和神.
既然你找到時間.
那你就要找地方.
所以第二步就是找一個屬於你的曠野.
不一定要去郊外公園這些地方.
你可以在家裡的廁所也行.
房間也行.
教會的祈禱室也行.
車裡面也行.
可能你上班的時候.
你午餐時間會有很多急房.
你可以在那些急房裡面成為你的曠野.
這個地方只需要沒有太多人打擾.
沒有太多人去騷擾你.
沒有太多人去煩你的地方.
所以你找到時間找到地方之後.
你就可以做一個屬靈的呼吸.
其實我們和神的關係很簡單.
就是父子.
就是父親和兒女的關係.
我們從來都不是一個主僕的關係.
當我們事奉到很累.
很大壓力.
甚至是不知道有什麼意義的時候.
不妨將這些東西和神去說.
就好像我們平時呼吸那樣.
呼出你的疲倦.
呼出你的壓力.
呼出你說我要對這件事一定要好的一些焦慮.
呼出別人的評價.
甚至是呼出你對事奉成果的執著.
將這些呼出的時候.
就在安靜的裡面去放手.
向上帝的旨意降伏.
在當中放下自己和上帝契合.
讓上帝在當中去轉化你.

$^{561}$我們同時間也可以吸入一些東西.
比如吸入上帝的愛.
祂的同在.
祂的平安.
祂的肯定.
要我們今天怎樣走這條道路呢.
怎樣去做呢.
這段時間你可以看看聖經.
你可以祈禱.
聽詩歌.
甚至是純粹一個很安靜的時間.
就是單單只是你和神一個相遇的時間.
我也有個建議給事奉者.
就是當我們每一次事奉完之後.
我們會不會有一些時間.
都可以問一下自己這三條問題呢.
就是在這個事奉的裡面.
我們有沒有感受到神的同在.
我們有沒有看到欣賞的事情呢.
我們會不會看到神在當中的工作.
而我們下一次的事奉裡面.
我們有沒有更多的空間去交給神呢.
這個是給大家一個挑戰.
你說為什麼是一個挑戰.
因為我們香港人真的很忙.
有時候Me time都沒有時間.
就是個人時間都沒有了.
你有五至十分鐘的時間.
你寧願做什麼呢.
看看YouTube.
錄一下IG.
大家有沒有去看《讀經計劃》.
好了,不問了.
五至十分鐘你寧願去看電視劇.
讀一下書.
不如想像一下.
我們給自己一個挑戰.
在這個星期的裡面.
有沒有一個時間可以和神相交.
有一個屬於你和神的曠野呢.

$^{601}$你不要回答星期六.
高博士那個《讀經日》.
那個不算的.
我們一至星期五的裡面.
在當中有沒有一個時間呢.
我想邀請大家閉上雙眼.
我們現在就想一下.
我們花一點時間去想一下.
你在當中.
這個星期.
接下來這個星期裡面.
有沒有一個時間.
地方是馬上想到的.
是可以成為了你和神的一個曠野.
留給祂呢.
好,我相信你們心裡面應該有一個時間.
有一個地點.
希望大家可以在這個星期裡面.
堅持下去.
一起和神有一個更多的結連.
最後我想分享一件事.
在我讀神學院的時候.
也有一種感覺就是.
我已經全身事奉了.
我已經投身在這個事奉裡面的時候.
但是為什麼我和神的關係.
越來越遠呢.
於是有一天.
走著走著.
我就去了神學院下面一個.
叫做藍tum 的地方.
你看到這兩個都是同一個地方.
這邊就是右邊.
這邊就是左邊.
你看到這些地方.
基本上我都坐過的.
無論是每一塊石頭.
甚至去到最遠那些.
我都坐過這些石頭.
或者坐過這些間條.

$^{641}$在當中裡面.
在這個曠野的時間裡面.
我是會去祈禱.
會聽詩歌.
甚至是會唱詩歌.
真的和神有一個親近.
當時神透過一首詩歌.
給了一個很大的鼓勵.
就是一首英文的詩歌.
叫做Here Again.
即是再次在這裡.
其中一句歌詞翻譯成中文.
就是這樣的.
它就說.
當我感到被離棄的時候.
主就在這裡.
祂能夠使孵肝的骨去復活.
主就在這裡.
很奇妙的.
透過詩歌.
透過聖經.
透過經文.
祂是讓我和神的關係.
去到越拉越近.
祂也是讓我見到.
神在我身上.
想看到的工作是怎樣.
以及堅持在這個事奉的裡面.
這個地方成為了.
我和神的一個曠野.
在我們的生活裡面.
我們很容易會因為家裡.
工作.
家庭.
去忙到很忙.
忙到不可開交.
但是就忘記了我們.
創造出來是為了和神.
去建立一份的關係.
亦都忘記了原來我們離開了神.

$^{681}$我們永遠沒有辦法去到獲得.
一份真正的滿足和成功.
如果耶穌.
一個神的兒子.
都需要時間和天父獨處.
才能夠完成他的人生使命的時候.
那我們作為門徒.
不是應該更需要多一點的時間.
去和神去到獨處嗎.
祝福大家在這個禮拜裡面.
在曠野.
亦都是見到神.
並且事奉得力.
讓我們一起和心禱告.
親愛的天父.
招聚我們在教會的當中.
亦都是因著我們的能力.
讓我們有事奉的機會.
是,主.
我們有時候在事奉的當中.
我們都會覺得很剛潔.
我們有時候都會被一些不同的源頭.
以致讓我們迷失了在事奉的當中.
求主你將我們去認清.
讓我們都知道.
我們事奉的源頭是來自於上帝.
是源自於你的託付.
求主你都是這樣去開我們的眼睛.
讓我們在接下來和你獨處的裡面的時候.
你讓我們可以看見你在我們身上的工作.
求主你這樣帶領,保守,祝福.
以上禱告奉主耶穌得勝的明智祈求.
阿們.
\newpage



\section{馬太福音 8:27-35}
\label{sec:Xc4fVvDZWRE}
\textbf{2026年4月19日崇拜直播｜陳冠成牧師｜告別浮淺,門徒進深:拒絕成功主義｜馬太福音8\_27-35}
\newline
\newline
連結: \href{https://youtube.com/watch?v=Xc4fVvDZWRE}{\texttt{https://youtube.com/watch?v=Xc4fVvDZWRE}} ~~~~ 語音日期: 2026-04-19
\newline
\newline
\hyperref[sec:rsvExPNsP2I]{\small{< < < PREV SERMON < < <}}
~
\hyperref[sec:index_chronic]{\small{[返順時目]}}
~
\hyperref[sec:3bVTrDHjdBs]{\small{> > > NEXT SERMON > > >}}
\newline
\newline
馬太福音 8:27-35
\newline
\begin{longtable}{cl}
\hline
\hline
章節 & 經文 (和合本修訂版)\\
\hline
8:27 & \begin{tabularx}{0.7\textwidth}{X} 眾人驚訝地說：「這是怎樣的一個人？連風和海都聽從他。」 \end{tabularx} \\ \\
[1ex]
\hline
\hline
\end{longtable}
$^{1}$《一路人情》 詞:陳家麗 曲:陳家麗.
一路人情 你是人 附著我走情路.
天天教我 怎約求吻 永不離婚夫的.
做我心知己知 做我甘愛你曾甘.
做我孤獨並犧牲 不變甘愛負責.
吻住情愛你.
一路人情 你是人 附著我走情路.
天天教我 怎約求吻 永不離婚夫的.
做我心知己知 做我甘愛你曾甘.
做我孤獨並犧牲 不變甘愛負責.
吻住情愛你.
我冥冥地追魂鎖 忘詞真自由 從誰淚光釋放.
做我心知己知 做我甘愛你曾甘.
做我孤獨並犧牲 不變甘愛負責.
吻住情愛你.
我冥冥地追魂鎖 忘詞真自由 從誰淚光釋放.
今天經訓是《馬可福音》8章27-35節 新約聖經59頁.
耶穌和門徒出去往愛薩利亞 菲勒比的村莊去.
在路上問門徒說 人說我是誰 他們說.
有人說是施浸的約翰 有人說是以利亞.
又有人說是先知女的一位 又問他們說 你們說我是誰.
彼得回答說 你是基督 耶穌就禁戒他們 不要告訴人.
從此他教訓他們說 人只必須受許多的苦.
被長老 祭司長和民事氣絕 並且被殺 過三天復活.
耶穌明明的說者說 彼得就拉著他 勸他.
耶穌轉過來向著門徒 就責備彼得說 撒旦 退我後邊去吧.
因為你不體貼神的意思 只體貼人的意思.
於是叫眾人和門徒來 對他們說 若有人要跟從我 就當社己.
背起他的十字架來跟從我 因為凡要救自己生命的 必喪掉生命.
凡為我和福音喪掉生命的 必救了生命.
弟兄姊妹平安.
就著今天這個講題 想先問大家一個問題.
你覺得你是否一個成功的人士?.
通常沒有人夠膽回答.
沒有人夠膽這麼承認自己是成功人士.
轉到另一個問題.
在你而言 你覺得甚麼為之成功?.
心裡可能開始有點盤算.
可能我要有甚麼物質或身份 就代表成功.
就著這個題目 我問AI.

$^{41}$大家幾個字都同意 2006年成功定義排行榜.
你猜第一位是甚麼?.
你猜猜.
猜到嗎?.
成功的定義 第一位是身心健康和快樂.
你覺不覺得?.
你覺不覺得現在多了很多講.
裡面這幾年講的不是身體健康這麼簡單.
心靈健康 或者是情緒也要健康.
現在很流行的一個字叫情緒.
給你一些情緒價值.
又或者幸福感這幾個字在這幾年很流行.
這是第一位.
第二位知不知道是甚麼?.
第二位就是工作和生活平衡.
我們用英文講就是甚麼?.
work life balance.
原來不只是工作 不是工作最重要.
原來我個人的需要或者家裡的需要都很重要.
所以你發現一有假期香港人就會怎樣?.
大家都答到了.
有人說可能你一年多過四次旅行.
有沒有人? 你已經是一個成功人士.
如果按照這件事講.
你去六次 你超級成功.
第三 第三是甚麼?.
第三就是財務穩定和自由.
財富依然是重要的.
不過目標不是為了大富大貴.
而是為了甚麼?.
夠就夠了 是的 我財富自由.
下一句是甚麼?.
我可以實現一下自己的人生夢想.
所以你又發現現在有時電視.
我昨天無端端看電視講怎樣退休.
講去哪裡退休比較划算.
講的人你知不知道多少歲?.
五十歲都未夠.
三十多歲就想退休.
即是想我不一定要很多錢.

$^{81}$如果在某個地方退休.
我有某個數目就可以了.
不需要很富貴 但是財務自由.
不需要為了房子做奴隸.
不需要為了工作做奴隸.
第四 第四是甚麼?.
第四即是傳統一點 甚麼為之成功?.
就是對社會有影響力.
得到人認同.
比較近我們傳統所講.
你事業有成了.
可以獲得人的尊重.
和發揮到你在那個界別的影響力.
幫到人.
你會發覺其實疫情過後.
其實傳統我們說成功的定義是甚麼?.
財富 你的地位 你的事業.
通常這些都是以前我們說.
成功追求的對象.
但是疫情過後其實開始轉了.
你有沒有發覺?.
成功的定義變得比較個人化了.
不只是物質或者財富的累積.
講求個人的健康或者家庭需要.
甚至是個人願望的實現.
不過不知道以上這些成功的定義.
是不是都是你所追求呢?.
又或者說以上這些對成功的追求.
會不會都影響我們對信仰的理解?.
會不會對我們來說都有一些偏頗呢?.
事實上今天這個題目.
成功主義或者說成功提神學.
其實就是說將世俗定義一些的成功.
譬如剛才所說的事業 財富.
甚至家庭幸福健康.
將它視為上帝唯一給你的祝福.
並不是這件事不好.
但是如果你只是單單.
上帝就是要這樣祝福我的時候.
問題就在這裡.

$^{121}$因為這樣會導致我們將信仰慢慢淪為.
追求我個人幸福.
追求我個人利益需要的一個工具或者手段.
純粹我只是希望在信仰裡面得到好處.
我只是希望上帝幫我實現我的夢想.
所以其實成功主義某程度上.
會令一個信徒他很難去接受.
甚至乎理解耶穌所強調一些信仰的核心價值.
就好像我們剛才所讀的那一段經文.
耶穌提到「捨己,背起十價,跟隨他」.
我們看回剛才所讀的經文.
你會看到彼得雖然是率先去宣認耶穌是基督.
他承認耶穌是上帝差派來的救主.
不過其實你會發覺.
他承認耶穌是基督.
原來不代表他願意擁抱耶穌基督的十字架.
是有一個落差的.
所以你會看到當耶穌宣告自己.
作為基督必須受苦.
被宗教領袖氣絕 甚至乎殺害的時候.
彼得的反應是超大的.
按了他的按鈕.
沒錯 他承認耶穌基督是上帝差派來的救主.
但是他對於耶穌基督怎樣執行他的拯救工作.
他很有意見.
他有自己一套的想法.
這一套想法也受到當時由他人的一些思想所影響.
就是他們會認為基督就是那位帶領他們.
成功推翻當時的羅馬政權.
幫助他們復國的那一位.
事實上當時無論是彼得或者其他門徒.
對耶穌基督都有這種成功的期望.
畢竟他們真的有些證據.
他們看到耶穌所做的成不成功.
真的行得通.
耶穌在他們眼前行了很多神蹟其事.
一句說話制伏了烏鬼.
一句說話讓大自然可以風浪平靜.
說一句話摸一摸那個人就醫好了.
又或者他可以用幾個餅 兩條魚.

$^{161}$就餵飽了五千人.
耶穌更新稱他比宗教領袖更厲害.
他擁有卸罪的權柄.
他有權去定義安息日做什麼不做什麼.
耶穌厲害嗎.
厲害 凡凡都可以.
在門徒的眼中.
試問一個這麼成功的人士.
怎麼可能他最後會被人氣絕.
甚至乎被殺害 失敗收場.
所以你見到耶穌十字架受苦的預告.
其實是完全顛覆了當時的門徒.
他們對成功的基督的認知.
全部顛覆了.
反過來 如果我們嘗試想像.
耶穌如果不是好像經文裡面所說.
而是對著門徒和彼得說.
稍後我們會揭竿起義.
我們會推翻羅馬的政權.
過程你知難免會有死傷.
會流血 甚至乎會有人犧牲性命.
不過不要緊 最後我們贏.
我們會成功.
你猜彼得最後會不會投訴.
應該不會了.
彼得大概不會再責備耶穌.
而是立即響應主啊.
我等了這一天很久了.
何時開始 哪裡行動.
要不要我多派一些人來幫忙.
因為這個完全符合彼得心目中.
那個成功救主的形象.
事實上你會看到到目前為止.
耶穌確實令彼得和所有門徒.
他們是享有更多的能力.
接著耶穌出入就怎樣.
更厲害 是不是.
人氣更好 影響力更大.
但是這一切的成功.
隨著耶穌這個受苦受死的預告.

$^{201}$意味著怎樣.
沒有了 完全沒有了.
所以你就明白為什麼彼得責備耶穌.
直接要指點耶穌.
你不可以這樣做基督.
基督應該不是這樣.
你看到彼得無疑是撇下所有.
撇下了他如夫的專業.
跟隨耶穌.
這一點毋庸置疑.
但是他另一方面又很想.
用自己的成功的看法.
來定義耶穌.
你應該怎樣做基督.
事實上我想我們都有這個情況.
我們可能當中都有人.
會立志跟隨耶穌.
但是有時遇到這些情況.
我們又會不知不覺.
套用了世上那個成功的標準.
甚至乎成功的法則.
來定義耶穌你應該幫我這樣.
耶穌我期望你這樣.
來達成我的心願.
不過你會見到耶穌卻很清楚.
斬釘截鐵地跟彼得說.
他怎樣去做基督.
不是人去定義.
甚至乎他將會說.
門徒要怎樣做門徒.
也不是人去定義.
不是按照世上的法則.
而是取決於天父的旨意.
正如耶穌都經常說.
先求的是什麼.
祂的國和義.
而不是求自身的好處.
所以你見到三十四節.
耶穌又說了一句.
我們很熟悉的聖經今句.

$^{241}$若有人要跟隨祂.
就要操練三件事.
第一件是寫己.
寫己顧名思義.
如果對應上文彼得說的.
就是你要捨棄那些.
耶穌錯誤的期望和看法.
不要像彼得那樣.
用你自己的需要和想法.
去限制耶穌怎樣做.
簡單說寫己就是.
你懂得要求自己.
體貼上帝的意思.
不體貼自己的意思.
這就是寫己.
你就是這樣操練.
學習為上帝的緣故.
放下一些地上暫時的好處.
選擇追求那些.
上帝看為更美好的事情.
這就是寫己.
第二個.
我們都很熟悉背起十字架.
簡單說就是效法耶穌.
稍後為了成就上帝的拯救計劃.
而輕看自己眼前的屈辱.
甚至失敗.
控制十字架的橫木.
從街上一直走走走.
走到營場.
當中一路示狀給人羞辱.
耶穌就是這樣要求門徒.
為福音的緣故.
加入這個好像失敗.
不成功的行列.
來操練順服上帝的心智.
特別如果你看路加.
他說天天背起你的十字架.
不一定說你沒有了一條命.
天天背起你的十字架.

$^{281}$就是你每天做好準備.
隨時有機會.
上帝要用你的時候.
你要擺上犧牲.
第三.
跟隨主.
這三個字我們都明白.
不過我們可能總是有選擇.
好像門徒這樣.
我們會很樂意跟隨那一位.
四維幫人.
很正面.
事事順利.
贏得很多人氣.
很成功的那一位主.
但是當耶穌也要求我們.
要誓死跟隨那一位.
要捨己.
要受苦.
甚至乎要受辱.
很失敗的主的時候.
人就未必願意.
我們未必願意.
特別我們可能.
如果你會受到坊間一些.
民間信仰的影響.
每個人拜他所相信的神.
或者偶像.
都是期望那個神明會怎樣.
你保佑我的.
你幫我的.
我拜你.
我信你.
你不幫我不只.
還要令我失敗.
那就怎樣.
轉台.
不要的.
是不是.
所以你會明白.

$^{321}$如果我們要誓死跟隨那一位.
捨己受苦.
好像失敗的主.
其實我們的信仰.
會馬上變得很沉重.
我們會很想.
跟耶穌拉遠一點距離.
甚至乎劃清界線.
我們拒絕.
將我們的生命主權.
交給那一位.
要求我們吃苦的主.
事實上.
耶穌今天這篇教導.
就是針對.
特別是我們今天.
一個追求很成功.
很舒適.
很自主.
不想捱苦的一個世代.
耶穌要求我們.
重新聚焦在十字架上.
十字架不只是掛在口邊.
也不是掛在頸上.
十字架是要來實踐.
要求我們活出門徒.
應該有那種屬靈深度.
事實上耶穌這裡說.
若有人要跟隨他.
就要捨己.
天天背起他的十字架.
跟隨他.
耶穌不只是在預告.
他自己將來會怎樣.
其實也在預告.
跟隨耶穌的人.
將來走一條什麼路.
但是這一條是成功的路.
意思就是說.
在耶穌裡面所說.

$^{361}$什麼叫做成功.
就是你回應.
上帝給你的呼召.
你並且要按照上帝的方式.
和按照上帝的時間表.
完成上帝所吩咐的.
耶穌就是這樣.
你會看到.
他在十字架上就是這樣得勝.
不是按自己的意思.
是按照上帝的方式和時間表.
來得勝.
其實不只是十字架上.
你會看到.
耶穌在他在加利利開展.
福音侍功之前.
他已經是這樣做的.
如果你留意我們回一回帶.
還記不記得.
耶穌正式出道之前.
他要通過一個考驗.
就是在曠野做什麼.
接受試探.
你留意.
那個是聖靈催耶穌去曠野.
接受試探的.
這個是一個考驗.
你特別打開程序表.
你會看到.
其實在馬太福音裡面記載了.
魔鬼撒旦在曠野.
對耶穌作出第三個試探的場景.
還記不記得.
這裡說.
耶穌又帶他.
就是耶穌.
不是 魔鬼又帶耶穌.
上了一座最高的山.
將地上的萬國和萬國的榮耀.
都指給耶穌看.

$^{401}$然後對耶穌說.
你如果去苦伏敬拜我.
我就將這一切.
萬國和萬國的榮耀.
都賜給你.
耶穌怎麼說.
撒旦退去吧.
因為經上記著說.
當拜主你的神.
丹要事奉他.
你留意其實.
耶穌面對這個試探.
同時也在教導我們.
什麼是真正的成功.
撒旦引誘耶穌給了他一條.
我們說是一條捷徑.
一條救贖的捷徑.
你要很快速地成功得著萬民.
拯救他們.
有辦法.
魔鬼撒旦對耶穌說.
你只要拜一拜就行了.
拜一拜的話.
你擁有全世界.
得著所有人的心.
確實是假如.
耶穌真的太軟弱.
屈從了撒旦這個試探.
確實可能他真的很快速.
贏得全世界.
這裡說了.
不過實際上.
耶穌是不是成功.
耶穌是徹底失敗.
因為他違背了.
上帝給他的時間表.
上帝說不是這個方式.
也不是在這個時候.
所以這段經文.
間接給我們看到.

$^{441}$原來每一個看上去.
好像是拓展.
上帝聖功的一個機會.
甚至方法.
不一定是出於上帝的心意.
我們要懂得去分辨.
正如耶穌斬釘切鐵地拒絕了.
撒旦這個成功的方案.
耶穌表示真正的成功.
就好像他最後.
勝過這個試探所說.
真正的成功是你專一.
敬拜天父.
單單事奉祂.
而不是事奉自己的需要.
弟兄姊妹.
其實好像我們反省.
我們會不會有時.
因為急於我們想獲取.
自己心目中的那些成就.
或者那些成功.
我們很想急於去贏.
而迴避了甚至拒絕.
上帝安排給我們一些.
艱巨一點的任務.
我們很想走short cut捷徑.
而迴避了吃苦.
我們不想面對上帝.
為我們擺在眼前.
那些挫敗.
那些煩惱.
但上帝說這些是對我們的生命有益處.
所以你見到耶穌.
無論祂在.
瑪厄福音.
或者馬太福音.
其實祂言行身教.
教我們你放心.
讓上帝的心意落實.
在你自己身上.

$^{481}$讓上帝在我們的生命當中得勝.
而不是急於因為要求勝.
不想吃苦而忽視了.
上帝的旨意.
上帝給我們的時間表.
而是要設法.
讓耶穌祂這個社紀.
背十字架的得勝的方式.
成為我們.
在生活裡面的成功法則.
簡單說.
耶穌的成功法則.
也是門徒的成功法則.
對我們每一個來說什麼是成功.
就是我們成就.
上帝的心意.
你回應上帝的呼召.
不過記住不是按你和我的方式和時間表.
是好像耶穌那樣.
按上帝給他的時機.
和方式.
來完成上帝.
託付給你的事物.
可能是一段關係.
可能是一個事奉.
可能是一些事情.
而另一方面.
就好像我們今天這個講題所講.
我們怎樣去拒絕.
成功主義.
入侵我們的信仰.
影響我們的生命.
我們不妨透過.
三個的操練.
三個在主裡的操練.
來開始.
經文也印在.
港澳大綱裡面.
第一個的在主裡是.
你要在主裡放鬆.

$^{521}$有時這個動作容易嗎.
放鬆.
不是說你放鬆身體.
做運動那些放鬆那麼簡單.
而是真的你.
面對一些很緊張的場景.
你都可以處之泰然.
那個叫放鬆.
如果你有留意.
我們看另一段福音書.
一樣是福音.
讓福音記載當耶穌他.
餵飽了五千人.
經歷這個神蹟之後.
你記得福音書記載有很多人即刻怎樣.
耶穌那麼成功那麼厲害.
有著落.
有大批人很想跟隨耶穌.
以祂為主.
以祂為首.
但是當耶穌開始宣告.
祂就是生命的糧.
人們可以吃祂的肉.
喝祂的血的時候.
群眾就怎樣.
耶穌你說什麼.
說一些離地都不知道你說什麼.
雖然這些是天國真理.
但是當時大部分的群眾都接受不了.
於是乎.
人們就紛紛議論耶穌.
甚至離開祂.
一下子走了幾千人.
不過你留意過程的裡面.
耶穌講了三句很重要的說話.
每一句.
其實都讓我們明白.
怎樣可以在天父裡面.
放鬆.
我想大家一起讀.

$^{561}$其實印在圓光幕都有.
這三句就是當時.
耶穌在養福音裡面所講的.
說話 我們逐句開始讀.
預備起.
(唸).
是 你有沒有留意這三句.
其實都是重複講同一件事.
人們為什麼來耶穌那裡.
是不是因為耶穌厲害.
不是 耶穌自己都不是這樣講.
耶穌說 負給我的.
負交多少給我.
是祂的意思.
留不留得住.
都是負的旨意.
你看到其實面對.
那裡講的可能幾千人突然間離場.
耶穌怎樣.
依然是慫容的.
因為.
祂知道祂的福音工作.
不是.
直根於人.
人為的事情.
而是深深直根於天父的主權.
和計劃.
耶穌明白 始終是天父.
招聚人前來.
也是天父親自留住他們.
所以無論當下的景況.
是怎樣 突然間不見了.
幾千人耶穌.
祂都仍然深信怎樣.
天父必定會將.
對的人在對的時候.
帶到你的身邊.
來完成福音使命.
而耶穌.
第一要做的就是怎樣.

$^{601}$當下按照上帝給予的時間表.
來宣講相關的訊息.
哪怕祂講完之後怎樣.
嚇倒很多人.
甚至乎我們逐一講.
所謂趕客走.
耶穌不介意 因為祂在執行.
上帝的旨意.
所以有這個結果.
就是說做對了事情.
但原來不一定是所謂成功.
成功的關鍵是.
做對了天父的旨意.
按照天父的心意來做.
不容易啊.
只會是你放我們回到這個場景.
如果我們面對類似的場景.
我們能不能像耶穌.
那樣處之泰然.
很難的.
不要講王札.
任何一間香港的教會.
你的聚會人數.
突然間不見了幾百.
或者說你在網頁的瀏覽人數.
明明每個禮拜都過千.
突然間那一個月只剩幾十.
你會怎樣?.
你會立即很緊張.
因為你覺得.
數字是一個很鎖定的東西.
數字跌了就代表什麼?.
代表我失敗.
代表我一定有些問題出了.
所以我們的邏輯是這樣.
於是教會就會很緊張.
立即開會去檢討.
發生什麼事 出了什麼問題.
然後我們就要.
做點什麼.

$^{641}$做點東西.
做點東西來.
谷起人氣.
谷起羞恥.
這個就是人的成功法則.
讓姐妹事實.
真的會這樣.
如果你留意.
在剛才約翰福音記載第六章後.
第七章耶穌的親兄弟又怎樣?.
叫他快點上耶路撒冷.
為什麼?.
說你人氣跌了這麼多.
不見了幾千人.
說道也好 行神蹟也好.
治病趕鬼也好.
把你的人氣拿回來.
召集所有人回來.
但是你看到耶穌拒絕這樣做.
耶穌沒有按世人.
那個成功的常理.
去出牌.
耶穌是拒絕追逐人氣.
追求成功主義.
耶穌反而是從容不迫.
按照上帝的時間表.
用上帝的方式.
去遵行上帝的旨意.
至於在日常生活中.
我們怎樣像耶穌.
這樣從容不迫.
我們怎樣可以.
在耶穌裡面放鬆.
其實有幾樣東西.
你可以不妨提醒自己.
假如你現在面對一些.
人生的失意 卡住卡住.
或者逆境的時候.
會不會難阻你享受親近耶穌?.
照理應該是不會的.

$^{681}$如果你知道.
耶穌比萬有都大的時候.
照理他應該比你身處的逆境大.
你不會因為純粹見到逆境.
而見不到耶穌.
你其實可以.
將你的憂慮交托.
憂慮一定是會有的.
但是關鍵是怎樣交托.
不會纏累你和耶穌的親密.
又或者說.
假如你是處於.
在壓力當中.
你不能躲避.
你能不能夠保持.
同樣的步伐.
你知道其實祂和你一起工作.
我們說信徒不是.
work life balance.
而是你的work.
和你的worship要balance.
你的工作就是你的事奉.
你的工作就是你的敬拜.
兩樣東西應該沒有抵觸.
應該拿到平衡.
不會因為你身處壓力.
我就捨棄了耶穌.
捨棄了你的信仰生活.
又或者說.
我們過份承擔了.
別人的責任.
那些責任其實是上帝良給.
對方的份.
但是很多時候因為我們.
抵不住脖子.
他這樣不行.
我要怎樣.
我要幫他.
我要出手抵不住脖子.
抱著他來做.

$^{721}$但誰知對方已經幾十歲了.
你要幫到他什麼時候.
你交回他給上帝.
讓上帝.
搞定他.
又或者說.
你願不願意接受上帝給你.
當下的一些局限.
你不會因為這樣而死命.
和上帝拼過.
就是不肯.
你堅持要你自己心目中的事情.
來去成就.
甚至不惜和上帝對著幹.
又或者說.
你能不能夠.
當下順著上帝給你一些供應.
給你一些安排.
你放膽去做人生.
每一個大小決定.
說起做人生決定.
不知道大家怎麼看.
我最近有些弟兄姊妹去說.
關於類似找工作.
我現在在做一份兼職的工作.
但是又有一份兼職.
兼職那份可能人工少一點.
但是穩定一點.
但是兼職那份.
雖然人工都不夠.
不過我又可以抽時間照顧家人.
可以選擇嗎.
抖不抖結.
我們明明在生活裡面.
經常有句口號.
有得選擇就是什麼.
有得選擇就是老闆.
有得選擇是好的.
但是原來生活真實裡面.
有得選擇都很煩.

$^{761}$於是乎那位弟兄姊妹就問.
死了會不會選錯呢.
哪一個才是上帝的旨意.
結果就很抖結.
弟兄姊妹你會怎樣去提醒與勉勵他.
其實沒人知道.
沒人知道你選擇完之後.
上班之後.
上了新工之後.
會不會遇到順境或者逆境.
因為那個不是重點.
但是很多時候我們就會.
想我選擇之後有什麼不好.
有什麼好純粹這樣想.
而去想這個是不是上帝的旨意呢.
這個是不是呢.
其實上帝不是用這個方式.
上帝讓你選擇.
這個方式是可以的.
沒有違背上帝的旨意.
都是幫助到人的.
上帝說你就放鬆去選擇.
選擇完記住.
不代表一帆風順.
選擇完之後都一樣會遇到困難.
一樣會遇到不安和憂慮.
但是你知道.
上帝還在.
那你就有平安繼續去走.
這個才是重點.
在耶穌裡面放鬆.
因為一切無論你怎樣選擇.
都在上帝的手裡.
除非你選擇不相信他.
那就沒辦法了.
至於第二個的操練.
是在說.
我們怎樣在上帝裡面放手.
意思就是.
我們為了滿足上帝.

$^{801}$為了福音的緣故.
你放下一些眼前個人的得失.
或者需要.
就好像《路加福音》裡面說.
最後一段經文.
有一個人就說.
主啊 我要跟隨你.
好不好啊 好啊.
不過他有個條件.
容許我先回去和家人道別.
你看看耶穌怎樣回答他.
耶穌說.
手拿著一個泥巴.
但是有不斷向後望的人.
新譯本譯得對一點.
不是說他不能救.
是說你和神的國不相配.
不合用.
這裡在說什麼呢.
我和家人道別也不行.
其實這裡的家人.
原本沒有人字.
你看英文清楚一點.
是house.
就是什麼 屋.
屋而已 那個物業.
你說都可以.
我搞定那層樓.
那就去讀神學 有什麼問題呢.
說不定他真的搞定之後跟隨耶穌呢.
但是.
搞多久.
會不會在搞的過程裡面.
捨不得.
其實這就是耶穌.
他所說的.
你會不會捨不得放手.
捨不得為了上帝的緣故.
放下你眼前的得失.
這段經文令我想起.

$^{841}$很久以前.
教會的一個弟兄姊妹.
他和我分享.
他很想讀神學 真的想讀神學.
不過.
他有一樣東西卡住.
這是不是上帝的心意呢.
這是不是上帝的心意.
其實怎樣去驗證.
其實很難.
因為客觀條件來講.
只要你有平安.
你的牧者又覺得.
你成熟.
加上你入表神學院修理.
那就是了.
之後的東西一路去驗證.
但是這位弟兄姊妹他有個難處.
我很怕如果不是.
我選錯的時候.
我沒有那份豐厚的.
公積金.
這才是重點.
講這件事是二十年前.
到最後有一天我遇到他.
他的臉.
都是悠悠愁愁.
當然我也知道他沒有選擇.
全職事奉這條路.
弟兄姊妹我想有很多事情.
我們很想為上帝去做.
但是我們又不捨得.
在地上一些的財寶.
我們不捨得那些的保障.
我們會擔心.
會不會我沒有這些的時候.
我就選錯.
我就失敗.
正如耶穌這裡所說.
你明明拿著泥巴.

$^{881}$你應該向前看一路去耕田.
你一路向後看的時候.
其實你都不知道在爬什麼.
所以耶穌說.
不好意思.
你太眷戀你的家業.
你不捨得放手既然是這樣.
就不能夠為主所用.
所以弟兄姊妹其實在主裡面.
放手的重點是說.
我們有時候真的要明白.
如魚紅掌.
是不可以兼得的.
所謂的屬靈.
分辨的能力其實.
不單單只是說.
我們懂得分辨善惡.
好壞.
更是我們懂得在兩件好的.
事情裡面怎樣.
選更好.
這個才是屬靈.
兩件都是好.
就像上帝給你的旨意.
選更好的福分.
放下次好.
事實上放手我想.
對不同的人在不同的年紀.
都是一個學習 都是一個功課.
很多時候不會一切都搞定.
而是你用你的一生.
在不同的生命階段.
學會放手.
意思就是你通過不斷的取捨.
慢慢讓上帝和祂的旨意.
我們經常說的.
放在你人生的首位.
又或者好像耶穌.
在赫西瑪利園禱告那樣說.
放手就是照上帝的意思.

$^{921}$不照我們的意思.
所以你看到耶穌.
祂不單單只是放鬆.
祂也都示範怎樣.
放手.
而最後一個要操練的.
就是怎樣在主裡面聽從.
回到我們最初那一段.
經文某一副音 當耶穌.
預告了自己受難之後.
下一個段落就是.
祂帶著三個門徒 彼得,雅各,約翰.
就去經歷登山.
變相這個神蹟.
記不記得 有點印象.
彼得面對 又見到耶穌.
和以利亞和摩西去.
去傾談的時候 祂很害怕.
其實經文形容祂 不知道發生什麼事.
不知道怎麼辦 於是符咒.
Do something.
不如我為耶穌理 為摩西.
和以利亞建三座壇.
證證上帝這個奇妙的作為.
你有沒有留意.
是上帝主動.
馬上打停彼得.
上帝命令他.
當下你只要做一件事.
彼得.
聽耶穌講.
聽耶穌講什麼.
如果你跟上下文.
就是耶穌之前說完.
祂會受死受苦這件事.
彼得有沒有聽.
聽了.
但是沒有聽從.
現在上帝又說.
你要聽從耶穌講.

$^{961}$耶穌之後還會有第二次.
第三次預言的受苦.
你要聽從耶穌講.
不要當耳邊風.
不要按自己一些成功的想法.
去忽視了耶穌.
這個重要的教導.
事實上給了我們.
很多的反省.
無疑彼得.
也認定他所做的.
一切都是出於好意.
問題是.
他有沒有對準耶穌.
對準上帝的吩咐.
會不會純粹因為出於好意.
讓自己甚至別人.
帶進一個祈禱裡面.
這個是很值得我們反省.
上帝固然願意.
我們為祂做一些事情.
但是祂更期望我們聽清楚.
仔細聽清楚.
聽從祂的吩咐和提醒.
正如我們都很明白.
上帝喜悅的是什麼.
是聽命勝於憲制.
不是做一些事情.
先聽了.
其實不難理解.
我想我們中間有很多.
很能幹的弟兄姊妹.
做事是很厲害的.
一呼百應.
但是能聽的弟兄姊妹.
多不多.
能聽從的弟兄姊妹.
更加難能可貴.
所以你看到.
上帝對彼得所說的.

$^{1001}$其實也是對我們所說.
先聽耶穌說.
我們可能有很多.
美好的計劃.
很想為耶穌實現.
但是原來跟隨主手要的任務.
不是首先為主做些什麼.
而是先聽了耶穌說些什麼.
這樣我們才能夠.
真正和耶穌同工.
正如昨天有沒有參加.
到今天.
昨天我們上半天就是操練.
什麼叫做聽.
真是用心去聽.
聽聽你裡面的心聲.
聽聽耶穌怎麼跟你說.
不純粹只是做些什麼.
事實上我們很明白.
跟耶穌同父一丫.
不是各有各做.
我們是配合耶穌.
不單止祂的旨意.
還要是祂的步伐.
雖然跟父一丫.
今天我們很講快速.
講求即食的世代.
我們有時會不會.
忘記了.
我們要配合耶穌的步伐.
我們有沒有好好操練.
怎樣在主裡面去聽從祂.
我們能不能夠.
像這裡所說的操練聆聽的功課.
放慢你生活的腳步.
以至到你內心.
心靈有更多的空間.
敏銳耶穌祂想你怎樣.
因應耶穌的需要.
有時可能你真的要.

$^{1041}$加快你的步伐.
但有時也可能放慢.
甚至停下來.
以至到我們能活出.
上帝的心意.
就好像施班今天獻重的.
副歌所說.
求耶穌帶領我們.
幫助我們深知祂的旨意.
幫助我們更愛耶穌的親近.
幫助我們更多.
好像耶穌那樣犧牲.
不變更愛父神.
求主帶領我們.
我們一起禱告.
主啊藉著這段經文.
幫助我們重新去定義.
什麼才是在主裡面的成功.
就是我們願意去回應.
上帝你給我們的照明.
並且按照上帝你的方式.
上帝你的時間表.
來去成就上帝你所教託的.
這一切可能是.
我們的生活.
我們的工作.
我們在教會裡面的服侍.
又或者我們和人的關係.
上帝我們知道這個.
不容易因為在今天這個.
講求快速的世代.
講求自主的世代.
我們很多時候都想.
走得比耶穌快.
想超前你.
甚至乎有時意見會和你相反.
我們很想自己贏.
但是我們沒有想過.
真正要贏的是耶穌你.
是耶穌你贏取我們的生命.

$^{1081}$是耶穌你在我們的身上.
獲得勝利.
意味著我們需要學習.
降伏的功課.
在主裡面學習放鬆.
因為知道萬事.
有上帝你托住.
有上帝你包底和帶領.
在耶穌你裡面學習放手.
不要去執著於.
我們自己.
很想成就的事情.
但卻不是上帝.
你的心意.
也在耶穌你裡面學習.
聆聽聽從.
聽從耶穌你的吩咐.
聽從你的幫助.
挽回和提醒.
以致我們無論人生.
是順是逆是高是低.
原來我們都是在主裡面.
經歷你給我們的成長路.
拜望主你幫助我們.
或許我們都有.
不同的狀況.
我們都需要去教託你.
需要去教託你.
願主你按照.
我們每一個獨特的生命.
來拖帶我們.
讓我們真的能夠走完.
耶穌你為我們所命定.
帶領的人生路.
我們這樣的禱告教託.
奉主耶穌基督你得勝的名.
來祈求 阿們.
\newpage



\section{歷代志上 17:15-27}
\label{sec:3bVTrDHjdBs}
\textbf{2026年4月26日崇拜直播｜鍾家威傳道｜告別浮淺,門徒進深:擁抱限制中的祝福｜歷代志上17\_15-27}
\newline
\newline
連結: \href{https://youtube.com/watch?v=3bVTrDHjdBs}{\texttt{https://youtube.com/watch?v=3bVTrDHjdBs}} ~~~~ 語音日期: 2026-04-28
\newline
\newline
\hyperref[sec:Xc4fVvDZWRE]{\small{< < < PREV SERMON < < <}}
~
\hyperref[sec:index_chronic]{\small{[返順時目]}}
~
\hyperref[sec:U29QDtiVI5U]{\small{> > > NEXT SERMON > > >}}
\newline
\newline
歷代志上 17:15-27
\newline
\begin{longtable}{cl}
\hline
\hline
章節 & 經文 (和合本修訂版)\\
\hline
17:15 & \begin{tabularx}{0.7\textwidth}{X} 拿單就按這一切話，照這一切異象告訴大衛。 \end{tabularx} \\ \\ \relax
17:16 & \begin{tabularx}{0.7\textwidth}{X} 於是大衛王進去，坐在耶和華面前，說：「耶和華神啊，我是誰，我的家算甚麼，你竟帶領我到這地步呢？ \end{tabularx} \\ \\ \relax
17:17 & \begin{tabularx}{0.7\textwidth}{X} 神啊，這在你眼中還看為小，你又說到你僕人的家將來的情況。耶和華神啊，你看顧我好像看顧尊貴的人。 \end{tabularx} \\ \\ \relax
17:18 & \begin{tabularx}{0.7\textwidth}{X} 你加於僕人的尊榮，大衛還有甚麼可以對你說呢？你是知道你僕人的。 \end{tabularx} \\ \\ \relax
17:19 & \begin{tabularx}{0.7\textwidth}{X} 耶和華啊，因你僕人的緣故，也照著你的心意，你行這一切大事，為了顯明這一切偉大的事。 \end{tabularx} \\ \\ \relax
17:20 & \begin{tabularx}{0.7\textwidth}{X} 耶和華啊，照我們耳中一切所聽見的，沒有可比你的，除你以外再沒有神。 \end{tabularx} \\ \\ \relax
17:21 & \begin{tabularx}{0.7\textwidth}{X} 世上有何國能比你的百姓以色列呢？神親自去救贖世上的一國，作自己的子民，又行大而可畏的事，顯出你的大名，在你從埃及贖出來的子民面前驅逐了列國。 \end{tabularx} \\ \\ \relax
17:22 & \begin{tabularx}{0.7\textwidth}{X} 你使你的百姓以色列作你的子民，直到永遠；你－耶和華也作他們的神。 \end{tabularx} \\ \\ \relax
17:23 & \begin{tabularx}{0.7\textwidth}{X} 現在，耶和華啊，你所應許僕人和僕人家的話，求你堅定，直到永遠；求你照你所說的而行。 \end{tabularx} \\ \\ \relax
17:24 & \begin{tabularx}{0.7\textwidth}{X} 願你的名永遠堅立，被尊為大，人要說：『萬軍之耶和華－以色列的神，是以色列的神。』這樣，你僕人大衛的家必在你面前堅立。 \end{tabularx} \\ \\ \relax
17:25 & \begin{tabularx}{0.7\textwidth}{X} 我的神啊，因你啟示你的僕人，要為他建立家室，所以僕人大膽在你面前祈禱。 \end{tabularx} \\ \\ \relax
17:26 & \begin{tabularx}{0.7\textwidth}{X} 現在，耶和華啊，惟有你是神！你應許將這福氣賜給僕人。 \end{tabularx} \\ \\ \relax
17:27 & \begin{tabularx}{0.7\textwidth}{X} 現在，你喜悅賜福給僕人的家，可以永存在你面前。耶和華啊，因你已經賜福，還要賜福到永遠。」 \end{tabularx} \\ \\
[1ex]
\hline
\hline
\end{longtable}
$^{1}$以下是以聆聽聖道敬拜.
經文是選至歷代至上十七章十五至二十七節.
是在教會聖經舊約的五百一十八頁.
請聽弟兄讀出.
歷代至上第十七章第十五節.
拿丹就案者一切話.
召者默事告訴大衛.
於是大衛王進去.
坐在耶和華面前說.
「耶和華神啊!我是誰?.
我的家算甚麼?.
你竟使我到這地步呢?.
神啊!這在你眼中還看為少.
又應許你僕人的家至於久遠.
耶和華神啊!.
你看顧我好像看顧高貴的人.
你駕於僕人的尊榮.
我還有何言可說呢?.
因為你知道你的僕人」.
耶和華啊!.
你行了這大事.
並且顯明出來.
是因你僕人的緣故.
也是照你的心意.
耶和華啊!.
照我們耳中聽見.
沒有可比你的.
除你以外再無神.
世上有何民能比你的民以色列呢?.
你神從埃及救贖他們作自己的子民.
又在你贖出來的民面前.
行大而可畏的事.
驅逐列邦人.
顯出你的大名.
你使以色列人作你的子民.
直到永遠.
你耶和華也作他們的神.
耶和華啊!.
你所應許僕人和僕人家的話.
求你堅定.

$^{41}$直到永遠.
照你所說的而行.
願你的名永遠建立.
被尊為大.
說:萬君之耶和華是以色列的神.
是致你以色列的神.
這樣你僕人大位的家.
必在你面前建立.
我的神啊!.
因你啟示僕人說.
我必為你建立家室.
所以僕人大膽在你面前祈禱.
耶和華啊!.
唯有你是神.
你也應許將這福氣賜給僕人.
現在你喜悅賜福於僕人的家.
可以永存在你面前.
耶和華啊!.
你已經賜福.
還要賜福到永遠.
各位弟兄姊妹平安.
平安.
擁抱限制終的祝福.
我已經很經歷了.
第一次在這裡和大家分享訊息.
主日的時候.
有點緊張和興奮和感動.
但是剛剛一開始祈禱的時候.
又發現原來也有不少的事奉人員.
今天也是第一次.
參與各項的事奉.
給了我一份同行的祝福.
希望以下的訊息.
都能夠我們在上帝裡面一同領受.
我們先一起去祈禱.
今天我們來到告別浮踐門徒進心的第四課.
題目是「擁抱限制的祝福」.
之前兩星期我們分別學習到.
在上帝裡面要有獨處的時間.
和上星期知道.

$^{81}$如何在基督裡面有個成功.
今天我想和大家分享有關限制的課題.
說到限制.
代表我們人身為受造物.
我們都是很有限的.
無論我們的能力.
性格.
家庭背景.
時間.
對神的認識.
都很有限.
但是很多人對於限制.
很多一些反感厭惡.
覺得是一種軟弱.
失敗.
因為我們社會的價值觀.
不斷在灌輸我們.
例如我們看到運動品牌.
Like it, just do it.
無論你天份是怎樣.
多累都好.
你就盡力去做.
一定會成功的.
Nothing is impossible.
Red Bull送你一對翼.
人類不應該被限制.
現實去束縛.
我們應該飛起來.
去超越極限.
這些的廣告.
這些的意識文化.
將身體很累.
受家庭的限制.
平庸的才華描繪成.
我們需要克服的敵人.
而不是我們生命裡面的一部份.
我們不需要去接納.
而是用來打破.
彷彿教導我們.
覺得一定要超越限制.

$^{121}$那些人才能夠去被看見.
被尊重.
導致我們去面對.
我們生命裡面很多限制.
一些軟弱的時候.
我們會有一種挫敗感.
羞恥感.
有時我們在信仰裡面.
都會有一種浮淺狀態.
有些人會去比較事奉.
比較什麼呢?.
他們比較什麼呢?.
比較自己付出得多.
犧牲得多.
覺得這樣就是越來越愛神.
這樣其實是什麼一個正確的心明呢?.
不是說超越限制.
追求卓越.
為主犧牲有問題.
但今天我們要反思的是.
作為基督徒.
我們有沒有想清楚.
我們當中是按自己的心意而行.
還是按神的心意而行.
是滿足自己.
還是滿足神呢?.
甚至有些人不承認自己有限.
拒絕讓神進入自己生命裡面.
用很多人和事.
別人的需要去滿足自己.
令到最後操勞過度.
burn out(消失).
在家庭.
身體.
情緒.
人際關係.
甚至靈性都出現了很大的問題.
今天我們如何在限制裡面去自處呢?.
今天選了歷代至上十七章的經文.
大衛他在人生裡面都遇上了一個限制.

$^{161}$他很熱心.
很渴望去為神見電.
但是他被神拒絕.
他的心意被神拒絕.
他從神而來領受的限制就是.
讓我們去知道.
人究竟願不願意.
馴服在上帝的主權底下呢?.
大衛的反應是怎樣呢?.
神如何幫助他去面對呢?.
就是我們今天想去學習.
一起去看.
剛才大家讀的經文.
其實只是整個大衛獻電士的一個下半集.
我們其實今天會看整個十七章.
但是因為篇幅太長.
我會在內文裡面和大家補充.
大家先幫我去讀十七章一召而接的經文.
預備起.
這裡其實本身是一個很感人.
很鼓舞的一個畫面.
當時大衛處於人生的高峰.
他住在一個名貴很舒服的香柏木宮裡面.
生活安穩.
他看著自己住得這麼舒服.
再想想象徵和神同在的藥櫃.
竟然在簡陋的萬紙帳篷裡面.
他內心很不安樂.
他很想去為神建殿.
認為這樣對神是好.
可能他看到當時古近東文化裡面.
其他的神明都有屬於他們自己很大的殿.
唯獨耶和華沒有.
他很想去為神做這件事.
他內心很好.
他不是想高舉自己.
告訴別人他有多厲害多厲害.
而是純粹想表達同神的敬畏和感恩.
他的熱心連拿丹仙子都認同.
所以都支持他的做法.

$^{201}$另外我們都看到有一些客觀的條件.
都支持他這樣做.
十一章指出.
他當時已經統一了以色列國.
國力強盛.
十四章又指出.
太耳王希蘭都將那些香柏木運到大衛殿裡.
運到大衛那裡.
又差遣使者石匠木匠去給大衛建造宮殿.
一些人,一些資源都已經預備得很好.
大衛又是神揀選的君王.
整件事情很合情合法合理.
我想我們都會很支持大衛的做法.
試想一下如果在座當中有一個弟兄姊妹.
有一天睡在你那張又大又舒服的海馬床褥.
突然之間很感動.
想為教會的籌款去盡一分力.
就跟Anders說.
我最近有一點點感動去包起了籌款的餘數.
我想大家都會很贊同.
覺得是一件好事.
不過上帝卻在這個時候對大衛說NO.
神的計劃不像大衛預期那樣.
他為大衛設了一個限制.
他借那丹仙招和大衛說.
你不可去建造殿宇.
在第十七章五至六節解釋了一些原因.
他說神自從帶領以色列民出埃及以來.
他都是過著流動的生活.
和那些以色列民一樣.
住在移動的帳幕裡面.
象徵著其實神的同在是不需要被固定地限制.
神不需要一個固定的殿.
而且神更加在第六節指出.
他未曾吩咐任何一個先知去為他建殿.
當然在後面的經文我們都知道.
神都有解釋大衛因為流很多人的血.
因為不潔而被限制.
我們看回整件事裡面.
大衛其實有很好的出發點.

$^{241}$很熱心去為神做事.
神也不是不喜悅大衛的心意.
不過我想大家都知道有時.
一份禮物是不是屬於好呢?.
其實不在於那個禮物的價值.
有時更加不在於那個送禮人的動機.
是看收禮物那個人怎麼看我.
好像我以前經常買些小禮物去送給我老婆.
一開始都在不斷摸索.
經常送完給她她就回我一句.
「你不知道我喜歡什麼嗎?」.
然後下刷一萬字就說.
「我說了很多次了」.
「這麼多年你還不清楚我是什麼人嗎?」.
大衛遇到的狀況.
同樣反映了他對神的無知.
他以為神很需要.
以為神和其他神明一樣.
需要大的宮殿.
但是神的心意不是這樣.
大衛遇到的限制不是他沒有錢.
沒有人才,沒有能力.
而是他未清楚他的主人是誰.
他未清楚主權在哪裡.
神的心意主權並不是這樣.
他有更高更遠的計劃.
大衛無意中將自己.
成為服侍神的那個人.
變為去供應神的那個人.
覺得上帝需要自己去供應.
他用自己的方式.
藍圖去取代了神的計劃.
他的熱誠.
遮蓋了他對上帝主權的敬畏.
但我們看到大衛他沒有一意而行.
而是他學習怎樣去面對.
我們回到經文看到.
在這個限制裡面我們看到.
大衛教我們怎樣透過神的話.
去重整我們的生命.

$^{281}$當上帝去透過.
先知拿丹對大衛的計劃.
say no的時候.
他沒有埋怨.
他沒有自憐.
他沒有陷入一種自我否定裡面.
他沒有繼續按自己的心意去起電.
他怎樣去面對呢?.
他做了一個很關鍵的動作.
十六節 大家剛才都有讀.
他說「於是大衛王進去坐在瑤華面前」.
進去是在回望前面坐在.
代表很謙卑屈膝在上帝的躍跪當中.
是一種完全降伏放鬆的姿態.
大衛停下來.
安靜去到神面前.
他重新去校正自己對上帝主權的看法.
去領受祝福.
他的生命被重整.
他重新去看自己.
重新去看神.
重新去看未來.
他怎樣去看自己呢?.
用了三個英文字給大家去.
今天做一個檢視.
他looks more(看得更清楚).
他接受自己的限制.
謙卑馴服在上帝的主權裡面.
去學習放下自己.
其實神在第七節.
一早已經跟大衛說.
提醒他的出身.
神是在陽圈裡面選擇他出來.
所以大衛坐下來.
第一句說話就跟神說.
「耶和華神啊,我是誰?」.
「我的家算什麼?」.
「你竟使我到這地步?」.
「我是什麼人?」.
「我做了什麼?」.

$^{321}$「我的身份又是什麼?」.
「我的家庭有什麼背景?」.
「神的恩典竟然臨到我當中?」.
他覺得自己很不配.
然後大衛就在整個禱告裡面.
他繼續去承接不配這個主題.
去不斷強調自己僕人的身份.
大家幫我讀17至19節.
預備起.
23至27節.
繼續不斷重複僕人,僕人,僕人.
大衛就是上帝的僕人.
他看清楚自己的身份.
在限制裡他看清楚自己的身份.
看到自己是多麼渺小.
多麼有限.
多麼不配.
大衛心明他不是自己有多厲害.
才有今天的地位.
完全是神主動無條件地揀選.
他很看到自己有多小.
今天當我們願意看到自己這個身份的時候.
我們就會學習到去接受自己的不足.
接受自己的一切.
我們不需要去按著人的需要.
去追逐在一些忙碌裡面.
我們接納自己的好.
不用透過一些方式去逞強.
我們接納自己的軟弱.
明白人都是需要限制的.
不能夠去認清我們僕人的身份.
其實會帶來很多的影響.
不能夠認清自己的身份.
其實不能夠認清自己的限制.
其實都會對不太好.
我想在我心靈裡面.
最近有一件小事.
就是早前我工人姐姐放假.
放一個多兩個月.
我剛好放大假.

$^{361}$在來這間教會之前.
我就一心做一個絕世好巴.
打算她每天上學的飯盒.
就幫我女兒去處理.
當然最後我是保留不住的.
但是大家猜一下.
我這個絕世好巴.
做一個做飯盒的人.
可以保留多久?.
保留了多久?.
一天?.
一兩個星期?.
其實我只是做了兩天.
我女兒就說.
爸爸你不如都是買外賣給我算了.
不承認自己的限制.
原來有時會有一些不好的情況出現.
在我們的人生命裡面.
我們有沒有承認我們有限制.
承認我們是主僕人的身份呢?.
如果我們忘記了這個身份.
我們很可能就常常以自己的角度.
自己的標準去看一些事情.
就好像在大衛不能夠成功見電一樣.
他可以去選擇.
很失望去問自己.
為什麼不行?.
為什麼不是我?.
可以自責去問.
為什麼是不是我做得不夠好?.
我做錯了什麼?.
或者他很憤怒.
將一些責任推給別人去埋怨神.
甚至乎有些人很想去追求一些解釋.
問個明白.
究竟我們當中一起所行的.
是按照我們自己的心意.
還是上帝的心意呢?.
還是我們有時其實很想去支配控制.
操控神呢?.

$^{401}$今天我們在面對限制裡頭.
我們是不是在拒絕去承認一些限制.
按自己的心意而行呢?.
聖經裡面第一件去拒絕承認自己有限制的事.
大家知不知道是什麼?.
就是我們的始祖.
亞當夏娃.
他違反了上帝為他們設立的限制.
去偷吃分別善惡樹的果子.
形成罪.
在我們的人生裡面.
我們有沒有沒有正視接納自己的限制.
不斷在超越.
甚至乎形成了一些很壞的惡果呢?.
今天我們會不會都.
我們可能好像大衛一樣.
很想為神.
為我們的家庭.
為我們的事業.
去做很多的事情.
不知道大家身邊有沒有一些人.
是以沒有時間放假為榮的.
鬥多那些AL.
明明很忙.
滿口很多的埋怨.
但是仍然不斷去接觸.
去填滿自己的時間表.
不知道這些人的內心.
其實反映著什麼呢?.
有時我們不眠不休.
我們拒絕承認身體的疲勞.
情緒的警號.
忘記了身邊和家人的重要性.
更加遺忘了神.
其實有時我們做傳道人.
有很多事情想做.
很多事情想認識.
特別在新的環境裡面.
都會因為適應這個身份.
而有些爪狂.

$^{441}$特別在這次預備信息的時候.
我也在裡面很經歷到上帝.
我這次呢.
我太太說你特別緊張.
平時你不會是這樣的.
為什麼你會改來改去.
又會問問別人.
又會睡不著覺呢?.
我自己也在反省.
究竟我是否看到自己很大呢?.
我很介意一篇信息.
或者別人怎麼看自己.
而沒有願意放下.
去正視自己的身份.
看到自己很小.
將主權.
將一切都交給上帝呢?.
今天大家又在一個什麼的境況裡面呢?.
有一個情緒健康的門徒測試.
有幾條問題有些評估.
反映了我們有沒有一些.
不肯承認自己.
生命裡面有限的狀況.
例如你會不會在別人.
叫你去請求的時候.
你會懂得say no.
不會讓自己操勞過度呢?.
你很明白上帝給你.
什麼是有所謂.
有所不為的呢?.
當別人有重擔的時候.
你可不可以很清晰地去分辨.
什麼時候出手.
什麼時候放手呢?.
你又能不能夠去察覺自己的情緒.
人際關係.
身體.
靈性的狀況呢?.
身邊的人.
會不會都很肯定你在家庭休息.

$^{481}$工作各方面都可以取得平衡.
還是另一方面.
我們不懂得去跟別人say no.
我們平衡不了我們很多的事情.
我們在很多的事情裡面找狂.
不能夠放下呢?.
今天大家是否在一個健康的狀況裡面呢?.
等於我們要認知我們是受造物.
我們要承認我們有限.
有容量.
就好像大家的手提電話那樣.
當容量超標的時候.
不知道大家有沒有試過.
會怎樣?.
可能很多功能都function不到.
拍不到照.
load得很慢.
甚至乎自己關了手機.
運作不了.
我們有沒有去打開自己心靈裡面的容量.
去檢查一下.
究竟會不會我們已經超越了我們的所限呢?.
以致到我們在影響我們心靈裡面各個層面.
我們的個人生命.
家庭.
事奉.
甚至乎跟上帝的關係呢?.
上帝很想我們在限制裡面謙卑下來.
更加重要的是.
我們在限制裡面.
我們可以看到上帝有多大.
「look big your God」.
我們怎樣可以在限制裡面去學習敬畏我們的上帝呢?.
我們再一起去看經文裡面所說的話.
上帝在第九和第十節提及過「我必為我們以色列選一個地方栽種他們.
使他們住自己的地方不再遷移.
凶惡之子也不像從前要害他們.
並不像我命是施至你我們以色列的時候一樣.
我必賜福你的一切仇敵.
並且我耶和華應許你.

$^{521}$必為你建立家室」.
上帝展示了他對以色列的拯救.
他會有自己的地方.
會為大衛去建立家室.
大衛是想為上帝去建什麼?.
建殿.
但是神卻為他去建立家室.
神展現他應許的超越性.
大衛自己以為可以去為神建殿嗎?.
但是神的計劃去超過人的所想所求.
大衛他怎樣回應呢?.
大家替我讀20-22節.
大衛回想過去.
他看到上帝對以色列人奇妙偉大的拯救.
他宣告「沒有可比你的.
除你以外再無別神」.
他宣告神的偉大.
他肯定神的同在.
今天我們有沒有很意識到.
上帝和我們同在呢?.
還是我們覺得神很遠.
覺得神幫助我們.
不能解決我們生命裡面的問題.
以前我讀神學.
一年級已經有一個屬靈的操練.
就是和耶穌去行長洲.
去哪裡呢?.
可能去吃東西.
去百佳買東西.
去墳場也試過.
去思想一下.
如果上帝一直和你一起的時候.
你的生命會有什麼轉變.
的確你真的會有很多意識.
會不同了.
可能對人會有禮貌一點.
我去到墳場以前很害怕.
沒有那份恐懼.
今天我們有沒有這個意識.
上帝其實一直和我們同在呢?.

$^{561}$還是我們將上帝放在一旁呢?.
人生的限制出現.
可能是一些我們能力的限制.
又或者是我們在生命裡面.
遇上一些不同的經歷.
一些遺憾.
這些種種的出現.
可能我們未必意識到.
但是在裡面我們教人去看到.
神的同在.
或者都有機會讓我們去.
重新去認識神.
看到上帝的超越性.
我在考A級.
大家知不知道.
應該大家都知道什麼是A級.
在A級的時候.
我有一個很沉重的經歷.
就是我要去重讀.
就是要重讀.
看到身邊的人都去上大學.
自己的內心很不好受.
更加我選擇自己用自修身的形式去重讀.
那一年要自己去安排自己讀書的計劃.
每天去自修室.
沒有人陪我.
沒有老師去幫我.
我覺得很挫敗.
很孤單.
還有我很怕在街上.
看到之前的老師和同學.
試過在馬路對面.
我看到有相關的人.
我會走開去避開他們.
因為我覺得這個經歷很失敗.
很醜.
但最後是什麼幫助我經歷這個艱難呢?.
有一位老師很好地提醒我.
在我全年自修.
不需要去上學的過程裡面.

$^{601}$他鼓勵我每天先去靈修.
再去溫習.
在那段時間裡面.
我真的能夠和上帝建立一段很好的關係.
對他認識更加多.
也幫助到我去面對我心裡面的挫敗.
我每天都和上帝很近.
這一年幫助到我跨過失敗的經歷.
更加這個經歷成為我日後蒙召的打底.
因為他幫助到我.
經歷到原來有神同在那份美好.
弟兄姊妹在一切的限制裡頭.
或者在我們日常生命裡面.
我們有沒有遺忘了上帝.
其實一直都和我們一起呢?.
祂的計劃.
祂的安慰.
祂的同在.
甚至超越我們的所想所求.
最後是大衛怎樣去重整.
他怎樣看未來.
他Look forward.
The God Plan.
他能夠學習.
用上帝的眼光去看當下的事情.
在十一到第十四節.
沒有擺出來.
神去強調他永恆裡面的應許.
而至到大衛在他自己禱告裡面的回應.
他都很大膽去和上帝求這個福分.
他求神堅定他所說的話.
要上帝一再去保證.
他不斷去說.
他說求你堅定直到永遠.
最後在二十七節.
他說你已經賜福.
還要賜福到永遠.
大衛沒有停留在.
起不了電的不開心失落裡面.
他告別了他屬靈的短視.

$^{641}$他調整他的視野.
去仰望永恆.
他擁抱神給他這個.
俗稱為的大衛之約.
他相信這個只是一個短暫的限制.
而是為了成就永遠的榮耀和祝福.
他可以去學習去靠著上帝的恩典.
他的應許去前行.
我們都知道大衛最後是沒有見電的.
但是他參與在整個見電的計劃裡面.
他沒有停下.
他很賣力在當中去預備材料.
將聖殿的藍圖傳給所羅門.
又幫忙召集不同的人去建立聖殿.
人生裡面很多的限制不似預期.
有時我們覺得可能是一些失去挫敗.
但是其實我們想深一層.
我們已經掌握到.
和上帝有永恆的上帝一起.
是不是已經是一份最好的福氣呢?.
有永恆的視野幫助我們.
不會停留在眼前.
不會停留在眼前的人和事.
逆境裡面.
我們不需要無限放大一些不信心的事情.
甚至乎不斷在一些負面的情緒裡面.
去陷入去.
我們有沒有這個視野呢?.
耶穌基督祂自己都自限了祂的身份.
祂為祂自己設限.
耶穌祂本身是無限的神.
祂是因為愛而選擇自限.
祂主動的道成肉身.
釘死了在十字架上面.
這個自限最終成就了福音.
為上帝為人的罪去犧牲.
祝福了每一個人.
在人的處境裡面.
我們都要學習去設限.
否則我們不能夠去承擔.

$^{681}$我們生命裡面不同的福份.
我們不能夠去承載很多不同的生命.
包括我們的家庭.
我們的休息.
我們的工作.
我們的情緒.
我們的靈命等等.
能夠去為自己的生命設限.
能夠承載更多的祝福.
大家想想如果一張白紙.
它可以裝到的水有多少?.
可能裝不到很多.
那些水就會離開了.
滲走了.
但如果將那張紙有些限.
接起了它.
成為了一個杯.
你就能夠去承載一些東西.
今天我們的生命有沒有設限的深度呢?.
我們不願意去設限.
我們就不能夠去承載.
我們生命其他的東西.
弟兄姊妹今天和大家分享.
我們告別浮淺門徒盡心.
在很多人生不似預期.
很多限制的事情裡面.
其實可能那個也是.
上帝和我們相遇的地方.
我記得第一次上台介紹自己的時候.
和大家分享一宗四件.
一家四口.
不過大家知不知道在我小兒子出生之前.
其實我們還經歷過一段很黑暗的日子.
話說在我大女兒出生之後.
我們一家都很想再有第二個小朋友.
因為很想再生一個小朋友去陪我的大女兒.
介紹一下.
將來我可能也會說.
我的大女兒叫貝貝.
很想去陪伴這個小寶貝.

$^{721}$讓她在生命成長裡面.
以及將來長大後都有伴.
所以我們很想有第二個小朋友.
計劃了很久.
最後我太太.
在一次驗孕的過程中.
真的發現她懷孕了.
我還很記得她發現自己懷孕的時候.
很開心很興奮地打給我.
和我說一個消息.
聲音是很雀躍的.
由心出發.
我們立刻去計劃怎樣去迎接她.
家裡有什麼安排.
幫她改名.
教傭人姐姐怎樣去照顧她.
以為一切都在自己的計劃裡面.
不過幾天之後.
太太發現自己有流血的跡象.
我們立刻去醫院檢查.
那一刻我們才發現.
原來那個胚胎是一個.
不正常的地方受精.
俗稱「宮外人」.
我們太太要馬上做手術.
去拿她出來終止懷孕.
那種由很雀躍到很悲哀的失落.
成為了一個很大的限制.
給我太太和我們一家很大的哀痛.
我們發現原來生命不在我們掌握裡面.
在那個限制裡面.
我和太太學習安靜.
我們坐在神面前.
我們發現自己其實真的很渺小.
我們掌握不到生命的去留.
但我們看到.
順帶讓我們經歷到祂的大是多大.
神雖然拒絕了我們那次的計劃.
但祂沒有離開過我們.
祂在那次的經歷裡面.

$^{761}$祂插派了很多.
原來身邊有很多人有共同的經歷.
祂這些共同經歷.
能夠安慰了我太太.
幫助到她度過一段不容易的日子.
而後來我們更加發現.
原來身邊的人.
都有這些共同的經歷出現了.
但他們不願意和別人分享.
但他們看到我太太的分享之後.
他們都能夠去學習和別人分享.
更加我太太都成為他們的陪伴者.
幫助到他們去經歷不同的難關.
很感恩我之後有了我現在的小兒子.
不過我想告訴大家.
真正的祝福不是在我最後得了兒子.
而是在那段艱難的時刻裡面.
上帝介入了我的生命.
祂擴闊了我生命的容量.
讓我們的家庭能夠成為其他人的祝福.
原來在我們很多的限制裡面.
有時我們看似接受不到.
不能夠去承載.
或者叫別人找狂.
或許是我們看到我們自己各行的有限.
我們的能力.
我們的背景.
我們一些的遺憾.
但原來在這些限制裡面.
神會擁抱我們的軟弱.
神會在裡面和我們相遇.
我們不需要用人的方式去作出回應.
我們不需要為這些限制去找狂.
神會在那裡陪伴我們去走.
今天希望我們一起去看.
我們如何在限制裡面去看回上帝.
看回我們自己.
以及靠著上帝的恩典去走過將來.
願意讓我們今天的訊息能夠幫助到大家.
我們一起去祈禱.

$^{801}$透過你的說話.
讓我們今天去正視自己的生命.
我們有沒有以自己的意念去蓋過你.
看不到上帝你在我們人生裡面的主權.
願意你幫助我們.
我們放下我們自己.
看我們看得越來越小.
去高舉你.
謙卑自己.
亦都靠著你去走過前路.
人生裡面我們有很多的限制.
但這些不是一些絕對的失敗,負面.
而是可能很多時候.
上帝你就在限制裡面與我們每一個相遇.
願你幫助我們每一個.
禱告奉主耶穌基督的名字而求.
阿們.
我們一起來唱一首回應詩歌.
誰曾應許.
人生裡面的困難.
我想大家都會經歷很多的事.
每個人都要面對信徒.
亦都不例外.
但在我們的生命裡面.
有神與我們同行.
是一個最大的安慰和力量.
誰曾應許.
\newpage



\section{哥林多後書 4:7-18}
\label{sec:U29QDtiVI5U}
\textbf{2026年5月3日主餐崇拜直播｜陳明泉牧師｜告別浮淺,門徒進深:發掘在失意裡的珍寶｜哥林多後書4\_7-18}
\newline
\newline
連結: \href{https://youtube.com/watch?v=U29QDtiVI5U}{\texttt{https://youtube.com/watch?v=U29QDtiVI5U}} ~~~~ 語音日期: 2026-05-07
\newline
\newline
\hyperref[sec:3bVTrDHjdBs]{\small{< < < PREV SERMON < < <}}
~
\hyperref[sec:index_chronic]{\small{[返順時目]}}
~
\hyperref[sec:wdtTzXf4txM]{\small{> > > NEXT SERMON > > >}}
\newline
\newline
哥林多後書 4:7-18
\newline
\begin{longtable}{cl}
\hline
\hline
章節 & 經文 (和合本修訂版)\\
\hline
4:7 & \begin{tabularx}{0.7\textwidth}{X} 我們有這寶貝放在瓦器裡，為要顯明這莫大的能力是出於神，不是出於我們。 \end{tabularx} \\ \\ \relax
4:8 & \begin{tabularx}{0.7\textwidth}{X} 我們處處受困，卻不被捆住；內心困擾，卻沒有絕望； \end{tabularx} \\ \\ \relax
4:9 & \begin{tabularx}{0.7\textwidth}{X} 遭受迫害，卻不被撇棄；擊倒在地，卻不致滅亡。 \end{tabularx} \\ \\ \relax
4:10 & \begin{tabularx}{0.7\textwidth}{X} 我們身上常帶著耶穌的死，使耶穌的生也在我們身上顯明。 \end{tabularx} \\ \\ \relax
4:11 & \begin{tabularx}{0.7\textwidth}{X} 因為我們這活著的人常為耶穌被置於死地，使耶穌的生命在我們這必死的人身上顯明出來。 \end{tabularx} \\ \\ \relax
4:12 & \begin{tabularx}{0.7\textwidth}{X} 這樣看來，死是在我們身上運作，生卻在你們身上運作。 \end{tabularx} \\ \\ \relax
4:13 & \begin{tabularx}{0.7\textwidth}{X} 但我們既然有從同一位靈而來的信心，正如經上記著：「我信，故我說話」，我們也信，所以也說話； \end{tabularx} \\ \\ \relax
4:14 & \begin{tabularx}{0.7\textwidth}{X} 因為知道，那使主耶穌復活的也必使我們與耶穌一同復活，並且使我們與你們一起站在他面前。 \end{tabularx} \\ \\ \relax
4:15 & \begin{tabularx}{0.7\textwidth}{X} 凡事都是為了你們，好使恩惠既藉著更多的人而加增，感恩也格外顯多，好歸榮耀給神。 \end{tabularx} \\ \\ \relax
4:16 & \begin{tabularx}{0.7\textwidth}{X} 所以，我們不喪膽。雖然我們外在的人日漸朽壞，內在的人卻日日更新。 \end{tabularx} \\ \\ \relax
4:17 & \begin{tabularx}{0.7\textwidth}{X} 我們這短暫而輕微的苦楚要為我們成就極重、無比、永遠的榮耀。 \end{tabularx} \\ \\ \relax
4:18 & \begin{tabularx}{0.7\textwidth}{X} 因為我們不是顧念看得見的，而是顧念看不見的；原來看得見的是暫時的，看不見的才是永遠的。 \end{tabularx} \\ \\
[1ex]
\hline
\hline
\end{longtable}
$^{1}$主持人:現在是經訓時間,今日的經訓是記載在哥林多後書第四章七至十八節,在唐容的聖經是二百四十七至二百四十八頁.請聽我讀出:.
主持人:我們有這寶貝放在瓦器裡,要顯明這莫大的能力是出於神,不是出於我們.我們四面受敵,卻不被困住..
心裡作難,卻不致失望;遭迫迫,卻不致被凋棄;打倒了,卻不致死亡.身上常帶著耶穌的死,使耶穌的生也顯明在我們身上..
因為我們這活著的人時常為耶穌被交於死地,使耶穌的生在我們這必死的身上顯明出來..
這讓看來,死是在我們身上發動,生卻在你們身上發動.但我們既有信心,正如經上記著說:.
「我因信,所以而處如此說話;我們也信,所以也說話.自己知道那叫主耶穌復活的,也必叫我們與主耶穌一同復活,並且叫我們與你們一同站在他面前,凡事都是為你們..
我們這至暫至輕的苦楚,要為我們成就極重無比永遠的榮耀.原來我們不是顧念所見的,乃是顧念所不見的.因為所見的是暫時的,所不見的是永遠的.」.
願神賜福常讀他話語並且進行的人.
弟兄姊妹平安!.
我們今天繼續用這個系列,我們希望可以幫助我們浮淺的生活,進入一個有深度的門徒生活..
今天我們要探討的主題就是人生中總會有一些失意,或流眼淚,甚至我們不想回想的,有一些痛苦的經歷或經驗..
這些經驗對我們的生命,或我們可以怎樣面對它,以致讓我們的生命更有力量走未來人生的路..
現在有一套很有趣的外國真人秀,節目名叫作《壯亂變現今》,這個節目很有趣的,.
它講到有一些被視為爛屋,有些人可能住在那間屋,有很多囤積了很多雜物,節目主持人會帶著一些專業人士去幫他清理,.
一般來說我們想這些清理節目就是扔掉所有東西,蛇蟲鼠蟻,但是這個節目就有一點不同,.
這個節目除了幫人扔東西,其實就在那間屋裡堆滿的雜物當中,就會找到一些充滿了價值,.
或者有一些人可能因為埋藏在一些垃圾,亂七八糟的雜物裡面,可能隱藏著很多很有價值的東西,.
所以到最後這些舊屋,有些人整間屋很大,到最後只有一張床來睡,其他都是雜物,煙頭,到處扔,.
一間這樣人家覺得快要不要的屋,可以幫他清理之後還可以賺到很多錢,.
有一個屋主在這些變賣和在這個執拾過程裡面找到38萬美金,不是在那裡有些錢,而是將那些東西變賣,.
到最後原來舊時是污穢的雜物的,亂七八糟的地方,原來有很多價值連城的東西在裡面,.
在這個節目裡面最吸引我的或者令到我很深體會的就是,我們經常在人生裡面覺得有些位置我不想碰的,.
在我心靈的角落裡面可能有一些記憶我們放在那裡,我們不想碰它,不想理會它,甚至我們不想想它,.
就放在我們心靈某一個角落當中,如果我們不碰它的時候,我們健康的人生會好好的,我們這樣來生活的,.
但當我們人生裡面經歷著一些低潮,或者在我們人生裡面面對著一些狀況的時候,.
在這些角落裡面的記憶就會出來,會令到我們很受困擾或者情緒很低落,.
很多弟兄姊妹在我們的成長裡面,或者在屬靈的人生裡面,正正因為我們沒有處理在心靈角落這些雜物,.
我們經常覺得那些東西是不好的,我們不再提及,我們覺得不理它就沒事,不過那些東西依然在,.
我們今天怎樣去處理它呢?今天我們藉著保羅他自己的心靈,或者用一個叫做創傷的角度,.
來去重讀保羅在哥林多後書,特別是第四章他所講的這一番話,讓我們重新去打理我們自己的心靈,.
可能這些事情積累在你的心靈很久,又或者近來發生,無論怎樣都好,希望今天我們去讀這段經文,.
我用幾個F的英文字母來作為幫我們去記,希望我們生活裡面拿到生活來應用,.
在我們一起看聖經之前我們一起祈禱,求上帝幫助我們,上主願你幫助我們,.
願我們每一個同在,讓我們能夠不單止明白你的華語,也藉著你的華語,藉著你的僕人所寫下的信息,.
能夠可以讓我們每一個弟兄姊妹,我們的生命可以再發出光芒,讓我們可以擺脫過去生命裡面很多的困難,.
求你與我們同在,幫助聆聽的每一個,不單止聽得明白,也進入我們的內心,扶持我們走未來人生的路,.
幫助孫講的講得清楚,將你的華語中心帶到弟兄姊妹當中,讓我們一同領受,一起謙卑受教,.
一起按著你的旨意來對行,我們祈求你與我們同在,獻上禱告奉求基督耶穌得勝的聖名,Amen.
我們一起看看哥林多後書第四章,在哥林多後書裡面,其實這一段,或者整個書卷,.
其實是保羅面對著一個生命裡面一個很大的低潮,他寫下給教會的一個書卷,.

$^{41}$其中一個很重要的面對的處境就是,因為保羅面對著哥林多教會,對他的非議,甚至拒絕保羅的教導,.
對於一個一手一腳建立教會的牧者來說,那些羊仔拒絕他,甚至覺得敵黨他,覺得他的教導不中聽,.
甚至乎他自己的經歷,和他所說的榮譽福音,是一個很不對稱的那種生活,.
所以保羅帶著的就是一個很強烈的情緒,來將這一段書卷,或者帶著一個很抽心的那種感覺來寫下,.
希望整個哥林多教會都能夠可以回轉過來,其中在第四章第七節那裡,.
保羅他在說他自己,或者在說每個人生是一個怎樣的實況呢?.
我們有著寶貝放在瓦器裡,要顯明這莫大的能力是出於神,不是出於我們..
在這裡裡面我們經常都會背的,這一句的金句,我們覺得有個寶貝藏在瓦器裡面,.
我們會想像有個罌裡面放一些很貴重的東西,不過其實這段經文他除了說有一個很貴重的心靈裡面,.
有一個很貴重的上帝的靈,或者我們所傳的福音之外,.
或者在現代中文譯本裡面,他比我們有更加多具體裡面所說的意思,.
經文裡面是這樣譯的,可視擁有者屬靈寶物的我們,不過是像普通的瓦器,.
為要證明這無上的能力是出於上帝,不是屬於我們..
簡單來說,現代中文譯本他更加貼近原文的意思,.
他強調的就是我們今天這一班人,我們信主,我們擁有這個福音,.
哇,這件事這麼正,不過擁有這個福音的我們,我們這一班的基督徒,.
其實只不過是一個普通的瓦器,這個普通的瓦器在當代社會裡面就好像一個杯,.
或者一個很易碎的器皿,舊時是沒有塑膠的,沒有那些膠杯紙杯的,.
他們用這些瓦器就某程度上是用完即棄的一種器皿,.
所以說一個很平常,甚至很易碎,沒有什麼價值的器皿,.
竟然承載著上帝寶貴的福音,.
保羅這裡其實要強調的是人生,我們每一個就像這個瓦器一樣,.
很普通的,甚至我們沒有什麼很了不起,.
我們得到這個福音並不是因為我們自己很厲害,很有智慧而得到這個信仰,.
相反來說,我們每一個都是血肉之軀,我們都是一個很平凡,.
甚至再認真去看的時候,我們的生命其實是一個很脆弱的一班人,.
正正是因為我們這麼脆弱,這個福音內在我們的生命當中,.
所以我們人生裡面能夠有力量面對所有的挑戰,.
並不是因為我們自己的能力有多強,而是內在我們生命裡面的福音,.
或者上帝的靈住在我們的生命當中,.
以致我們有能力面對著生命的林林種種,.
所有的高低起伏都是來自上帝,.
所以在這裡裡面,保羅他要強調,我們的生命很脆弱,.
但是擁有或者我們生命裡面盛載著的卻是這個寶貝的福音,.
這個寶貝的福音卻是為我們生命帶來一個很大的改變,.
承認人生是脆弱是很重要的一回事,.
我們今天這個社會或者這個文化裡面,.
他要教曉我們,或者我們的受訓練,.
我們經常都要令自己變得很強大,我們要裝強,.
不過在聖經裡面,特別是保羅自己,.

$^{81}$他自己面對著一個很痛苦的傳福音的經歷,.
被人拒絕的經歷,他很直接去講,他都是一個很脆弱的人,.
他面對著生命裡面這些林林種種,.
如果按他自己血肉之軀,他頂不住,.
唯有透過上帝,他才有力量去面對,.
正如我剛才所講,我們今天這個世界的文化,.
他會教曉你,或者你要怎樣做都好,.
你要令到你自己看上去是強的,.
在我讀大學的時候,我們經常都要present(展示)的,.
我的老師,我的lecturer(教師)說,.
如果你去做簡報,講到你不會講,你不會回答別人的問題,.
總之你就要淡定,總之要裝模作樣,輸人就不要輸陣,.
經常都是這樣教的,所以不懂得回答都好,.
你都要微笑,然後挺起胸膛,然後你還要質疑別人的眼睛,.
別人就會覺得別人問了一些蠢問題,別人就會知難而退,.
其實你沒有實力的,不過你裝得很有實力,.
你不尷尬,別人就很尷尬的了,就是這樣教,.
以前我們都是這樣過我們自己的生活,.
不過在聖經裡面,它提醒我們,我們要承認自己,.
我們不是想像中那麼強,甚至我們在整個信仰裡面,.
我們不是要裝得自己很堅強,.
而是我們才可以順應著上帝的計劃,.
我們牧者探病很多時候很有趣的,特別是男士,.
男士很少認自己弱的,探病呢,姐妹呢,通常你去探他,.
姐妹不舒服,他會告訴你,我很辛苦,很不舒服,.
男士通常不說的,去到一些地步最極致的,.
我去探過一些男的老人家,他可能身體都離流的了,.
你覺得怎樣?沒事,很快出院的了,.
我覺得他是安慰我,但同時間也安慰自己,.
甚至為自己打打氣,在最軟弱的日子裡面,.
他怎樣都不會流露自己半點的軟弱,.
不過在裡面保羅說到,他自己或者我們整個人生都是啞氣,.
然後去到第八節,然後北風就轉了,.
本來很脆弱的自己的生命,去到第八第九節他怎樣說呢?.
我們四面受敵,卻不被困住,心裡作難,卻不自失望,.
遭逼迫,卻不被凋棄,打倒了,卻不致死亡,.
你留意在這裡兩節的經文,他講了四次卻不,.
四面受敵,不會困住,心裡作難,不會失望,.
遭逼迫,卻不被凋棄,打倒了,卻不致死亡,.
和剛才所講的軟弱是很不同的,剛才講人生很軟弱,.

$^{121}$人有上帝的寶貝,內在我們的生命當中,.
因為有上帝這一份的能力,所以我們人生裡面面對著四面受敵也好,.
心裡作難也好,遭逼迫,甚至被打倒,這些一切都不會走向絕路,.
有幾個字和大家仔細研讀一下,你就更加明白他的意思,.
他講的四面受敵,保羅講的四面受敵,不是周圍都是敵人,.
他不是要周圍樹立敵人的人,所以其實他講的四面受敵,.
是講四方八面來到的苦難,那些不是他計劃之內,.
甚至乎如果按著人的力量是招架不住,.
四方八面湧過來的那些苦難,這個叫做四面受敵,.
卻不被困住,困住的字在和合本譯的字眼沒有那麼有力量,.
原文用困住的字是壓扁的意思,四方八面來到的那些苦難,.
不會將我壓扁,不會將我打碎,這個就是他第一個的豪情壯語,.
然後心裡作難,心裡作難是什麼意思呢?呂鎮忠的本譯就是計謀絕了,.
比較間接,計謀絕了,廣東話會譯得比較傳神,沒辦法,.
即是我面對著生命裡面所有的難處,我投降了,沒辦法,.
沒有計劃,沒有辦法來招架,但是在沒辦法,.
到絕地的時候,我內心裡面不會失去盼望,.
遭逼迫,這個字也很直接,也不需要再有補充的空間,.
不過卻不被丟棄,那個被丟棄的意思就是說,.
因為我們面對著逼迫,這些外在的困難那麼大,.
也不會質疑了我們自己本身既有的身份,.
即是說人面對多少困難,絕路也好,.
我不會質疑自己的,我不會否定了自己我是一個什麼樣的人,.
在這個世界裡面有時候人經歷很多苦況,.
連自己是誰可能都會有一份迷失感,.
但是保羅說不會迷失,因為有上帝在我心靈當中,.
打倒了就好像打拳一樣,砰一聲打下去,.
但是被打的時候卻不致死亡,意思就是說打倒了之後,.
他可以爬起來,他還沒死,他還一息尚存,.
所以在整個處於逆境,一種艱難階段裡面,.
保羅說到因為面對著那麼大的困難,.
但正正因為有上帝在心靈當中,.
上帝令到一個血肉之軀,一個軟弱的人,.
他可以在這個逆境裡面有一種堅毅力,.
而他可以在人生裡面不怕,.
經文裡面保羅在第十節開始解釋,.
面對著那麼多苦難,很多人基本上已經倒下了,.
很多的色經學家也好,很多心理專家都去研究保羅的心態,.
第十節至第十三節是一個最重點,.
來講為什麼他有這一份能力,第十節,.

$^{161}$身上常帶著耶穌的死,使耶穌的生也顯明在我們身上,.
因為我們這活著的人,是常為耶穌被交於死地,.
使耶穌的生在我們這必死的身上顯明出來,.
這樣看來,死是在我們身上發動,.
生卻在你們身上發動,.
但我們既有信心,正如經上記著說,.
我因信,所以如此說話,我們也信,所以也說話,.
在這裡有幾個重點,保羅為什麼他面對著那麼多危難,.
他也有這份堅毅力,首先他將他自己所傳的福音,.
和他的際遇,和耶穌基督的生命,他畫成了一個等號,.
簡單來說就是,因為他傳福音,他才面對這樣的逼迫,.
原本他是社會裡面一個上等人,他是羅馬公民,.
他還要擁有加瑪列門下的律法知識,原則上來說,.
他當時來說,他是高人一等的,比普通平民百姓,.
不過,他將這些東西丟下了,.
他現在就將他畢生的努力,就傳耶穌基督的福音,.
所以耶穌基督所經歷的際遇,和他所說的內容,.
甚至乎他個人自己的經歷,就畫成了一個等號,.
沒理由,他純粹傳耶穌很受苦,他自己享受世界,.
他正正在傳一個受苦的基督,所以傳的人,.
他自己也好像耶穌一樣,甚至乎他身上就帶著耶穌基督的死,.
耶穌的死,代表著他由受苦到死亡,死在十字架上的經歷,.
他就成為了他生命的一部份,他也用這個方式,.
來看自己整個的使命,所以既然耶穌受苦,被逼迫,.
死在十字架上,之後復活,所以保羅看自己,既然耶穌會復活,.
所以他自己生命裡面即使經歷著這些逼迫,.
他也會看在逼迫過後,在苦難過後,耶穌的生命就會呈現在他的人生當中,.
有一些的《釋經書》甚至說,保羅有時候他寧願受苦受到一個地步,.
他寧願死了他也好過,不過他為了福音的緣故,.
他想起耶穌基督,他知道活著能夠繼續成為福音的使者,.
所以他就咬緊牙關,看著耶穌的生命是怎樣,.
他就繼續走這一條道路,整個歷程不是被迫的,.
而是他仰望著主,而至到他就將耶穌的生命呈現在他自己的一生當中,.
所以經文裡面,首先第一樣東西,他的動力,.
是將他自己的際遇和他所傳的信息,畫一個等號,.
第二樣東西,耶穌的死或者這個苦難並不是終點,.
今日這個世界看人生,人生的死亡就是人世間的終點,.
不過在保羅裡面看,死亡只不過是一站,.
永生,生命才是每一個基督徒的終點,.
亦都是耶穌基督的終點,所以他看死又何干,.

$^{201}$死亡只不過是經過人生的一站,不是終點,.
只要他知道他的終點是邁向永生,所以他就有力量去面對,.
第三方面,保羅引用詩篇第116篇第8至10節,.
他引用舊約詩篇,他不是解釋那段詩篇,.
他只是引用那句說話,詩篇116篇第8至10節是這樣寫的,.
主啊你救我的命免了死亡,救我的眼免了眼淚,.
救我的腳免了跌倒,我要在耶和華面前行活人之路,.
我因順所以如此說話,我受了極大的困苦..
這個詩篇其實都是一個哀歌的形式,.
他在說我很痛苦,不過上帝救我脫離了死亡,.
他幫我抹乾眼淚,我的腳跌倒沒力,他扶起我,.
為什麼我有這樣的能力呢?因為我信,.
我信才能夠說得出這番說話,不過後面再補充,.
我是受了極大的困苦,保羅就用這一句說話,.
我信所以如此說話,他不單只是說一個外在的知識,.
而是說這是他生命裡面所信靠的主,.
他因著這個信心,所以他就有力量去面對這些艱難..
整個經文裡面其實說到保羅那種能夠在苦難裡面有一種堅毅,.
其實他將自己和耶穌的際遇畫了等號,.
第二件事他看死亡只不過是一站,最後他深信耶穌的復活,.
也為他自己的生命帶來真正的出路,.
所以他憑著信心來做一個這樣的宣講..
我們繼續看第十五節,.
凡事都是為你們,好叫因為因人,多月發加增,.
所以我們不喪膽,外體雖然委為,內心卻一天,身馳一天,.
我們這站至輕的苦楚,要為我們成就極重無比,永遠的榮耀..
在這裡裡面保羅他繼續說,.
做這麼多事,受這麼多苦,他的目的不是為了自己,.
原來是為了別人,為了什麼人呢?.
首先為了身邊的人聽到福音的好處,.
所以他講到很多人因為上帝的恩惠,.
可以越來越多人得著上帝的恩典,.
多人得著福音,所以上帝得榮耀,.
為了別人得到福音的好處,.
然後間接地反映上帝的榮耀,.
所以這是為了上帝得到榮耀,也為了生命其他人得到好處,.
第二件事,他再看自己,.
今天他看自己其實很大的苦難,四方八面湧過來,.
不過他怎樣描述這些苦難呢?.
至暫至輕,短暫的,.

$^{241}$如果有個時限和永生相比,.
可能是一刻,可能是一段很短的時間,.
本來用一個時間無限的角度來對比他當下的痛苦,.
所以他描述眼前所面對的只不過是一個至暫至輕的苦楚,.
當他將來所承受的就是一個極重無比永遠的榮耀,.
這個榮耀不是為了自己,.
而是一個上帝榮耀得著彰顯的日子,.
所以整件事,他的努力,.
或者他心裡面正面的動力就是看到上帝的榮耀,.
和人得著福音的好處,.
在保守自己的心懷裡面去看自己的死,.
所有事情都不是一個終點站,.
他因著信心,他就有力量來走面前人生的路,.
苦難對於我們來說很多時候是荒謬的,.
我們有時候要解釋苦難為我們帶來什麼好處,.
有時候真的說不出,.
甚至我們逼別人說,我們覺得是一個很冒犯的舉動,.
我們不需要為身邊所有的苦難賦予一個合理的解釋,.
有一些經歷來到真的無法解釋,.
不過我們可以善用在這些日子裡面,.
上帝容許這些事出現在我們的心靈當中,.
在我們的人生裡面出現了這些事,.
這些苦難卻可以塑造我們生命,.
成為我們生命更加成熟的一個機遇,.
有一個早期教父,他的名字叫做愛淫柳,.
愛淫柳有一個說法很有意思的,.
他講到人出世,或者每一個人,.
我們每一個人都是按著上帝的形象而做的,.
所以他就用《創世記》裡面講,.
人是按著上帝形象來做的,.
不過這個形象不等於是最完美的一個形式,.
人最完美的,或者人最終的成長,.
就要成為了按著上帝的樣式,.
形象就是出生的時候,.
樣式就是你生命裡面最成熟的那種表現,.
當然有不同的《釋經》解法,.
有些人將形象和樣式等同一個字,.
不過愛淫柳的解法是兩個意思,.
形象就是像上帝,.
但樣式就好像成熟的人,.

$^{281}$接近上帝的那種表現,.
而形象到樣式中間,.
是要透過什麼來塑造一個人呢?.
苦難,人生裡面很多的美德,.
人的成熟是透過生命這些苦難,.
這些逆流來去鍛造一個生命,.
在愛淫柳裡面繼續講,.
很多的美德其實是由苦難鍛鍊出來的,.
例如忍耐就是由苦難裡面鍛鍊出來的,.
信心也是由逆流而上的處境裡面鍛鍊出來的,.
甚至愛心也是在一些不可愛的人裡面,.
因著你要操練你的愛,.
在一些不可愛,不體面的人裡面,.
一路一路實踐出來,.
當我們人生裡面經歷著一些逆流,.
這些苦難,這些艱難,.
我們生命就可以慢慢地成熟,.
朝著上帝的樣式來進發,.
如果沒有這些東西,.
我們永遠都是在地上裡面浮淺的一個人,.
當然愛淫柳的講法不是要你合理化你生命裡面所有的苦難,.
不過我們也不能否定在我們生命裡面出現了的這些際遇,.
就好像一間破屋裡面一樣,.
我們在當中撥開了很多的垃圾,.
找回一些生命裡面最珍貴的東西,.
很多人視為垃圾,.
甚至我們覺得埋藏了它,我們不理它,.
但當我們梳理,撥開了這些不起眼的,.
或者這些垃圾堆裡面,.
我們找回生命裡面一些很重要的寶物,.
這個就是我們自己成長裡面最重要的珍寶,.
或者正正是在這些苦難當中鍛鍊我們自己的生命,.
幫助我們,讓我們成為一個更加成熟的人..
有一本書叫做《門徒的EQ覺醒》,.
這本書我們整個同工團隊一起來看,.
在這本書裡面,它圍著我們今天,.
我們如果真的很想在人生裡面找到這些這樣的際遇,.
你的負面經歷,你可以怎樣做呢?.
它給出三部曲,我又用了F字母來幫大家記得,.
第一件事你要做的,.

$^{321}$在你人生裡面面對著這些艱難,.
你面對它,你不要逃避它,.
如果你真的痛苦的,你不要轉過頭來,沒事,.
否定它,平外傷而已,沒事的,.
面對它,正視它,.
真的受傷的,你說出來,.
分享你自己心靈之間的那種痛處,.
不要掩飾它,不要裝作若無其事,.
正正是因為面對這些傷痛,我們才能夠仰望上帝,.
如果我們生命裡面經常覺得好好的,.
我們就沒有一個動力,.
我們會覺得透過自己的能力去解決眼前這些艱難,.
這本書的作者提醒我們,面對它,.
不需要擔心的,.
在上帝的懷抱,在上帝的愛當中,.
面對這些聖靈裡面的苦難,.
面對了它之後,我們接下來所做的就是等候,.
在過渡期中來到靜候,Foot Pause,.
我們不要急於為自己解決這些感受,.
立即為自己製造一大堆的行動清單,.
我做這些事情,我可以解決了,.
或者找一本坊間裡面的心理書看看,.
然後自學怎樣幫自己,.
人生有些歷程是要熬的,.
或者正正是在這些傷痛裡面,.
將你的經歷帶到在上帝面前,.
仰望這個慈愛的主,.
可能在你的傷痛裡面,可能上帝顯得沉默,.
但正正就是這一種停下來,.
在這個生命裡面的那種仰望,.
我們聖靈才可以得到一個不是廉價的出路,.
而是一個真是有踏實,經歷上帝同在的一個經歷,.
相反來說,如果我們用一些坊間裡面的方法,.
我們沒有停下來,我們就好像在用一個止痛藥一樣,.
當你發燒的時候,你立即按下那個燒,.
以為就沒事了,那些燒退了,以為那個人就沒事,.
其實可能你的人,那個體質已經虛到不得了,.
因為你用那些退燒藥來按下自己的發燒,.
你覺得沒事了,沒有發燒就OK了,.
不過可能你自己真正令到你發病的,.

$^{361}$或者那種痛苦,你都沒有正視它,.
不要按低病症,而是要正視心靈深處裡面的這些傷,.
以及這些生命裡面的體會..
最後就是發現新的面向,Find,.
在這些心靈痛處裡面,嘗試找一找上帝的恩典,.
嘗試找到一下自己,原來在這些處境當中,.
原來自己會有這一面的,認識自己,.
看到自己心靈的價值,如果能夠在痛苦裡面,.
真正地看到真實的自己,.
你的人生,你過你自己未來人生的生活,.
你就變得踏實,有力量,相反來說,.
如果你所採用的都是一些廉價的方法,.
或者你沒有在這些歷程裡面幫助自己,.
重新認識自己,認識上帝,我們所有東西都是很浮淺,.
我們在生命裡面第二個際遇來到的時候,.
壓在我們內心裡面的這些負面記憶,又再爆發出來..
在我預備這篇講章的時候,我想起一個經歷,.
在我爸爸媽媽過世的時候,.
我爸爸媽媽所住的家,我一直不敢進去,.
因為很多記憶,有很多很抽心的歷程,.
一直都不夠膽面對,終於有一天,.
我自己靜靜的上去了,我記得我特別想去看的,.
我特別進去的,不是進去看那間屋,.
而是我想拉開我沒有鬥三十多年的我爸爸的衣櫃,.
其實對很多人來說,沒有什麼是衣櫃,.
但那是在說我和爸爸之間的關係,.
我鼓起勇氣拉開,那些衣服可能我媽媽生前還會拿來看,.
所以摺得很整齊,我爸爸生前只穿幾件衣服,.
如果在在你就知道,穿蒙特嬌的,大家知不知道?.
我爸爸穿蒙特嬌的,穿蒙特嬌絲質的襯衫,.
有幾件蒙特嬌就摺在那裡,.
我打開的時候想,我爸爸去世的時候,.
大概我現在四十多五十歲的年紀,.
我爸爸在這段時間在想什麼呢?.
就是在那裡很多勇氣,很多體會,.
拿一件灰色的蒙特嬌想起他,其實沒有什麼結論,.
不過我覺得這個經歷提醒我,面對,.
可能人開始大,開始要整合,回顧以前年輕的時候想衝,.
或者做一些事務可以抵消,或者分散了一些注意力,.
但人生去到某個階段的時候,.

$^{401}$我們很想拿回那個缺失了的那個洞洞,.
很想找回這些失去了的記憶,.
那件蒙特嬌拆出來,.
我心裡面回味爸爸昔日的日子,.
也提醒我,今天我作為一個中年人,.
我可以怎樣祝福我這個家庭,.
我怎樣有力量走未來人生的道路..
弟兄姊妹,在我們人生裡面,可能我們都有這些角落,.
有這些我們不敢打開的抽屜,.
甚至可能生命裡面有很多東西壓住了,.
我們遮蓋了,看不到生命的寶物,.
但當我們去面對它,重新靜下來,.
那個歷程未必是快的,每個人的時間不同,.
但願意靜下來,重新再去發掘的時候,.
你會發現你的生命不再一樣,.
求主幫助我們,如果我們門徒的生命要盡心,.
我們這些記憶,我們需要在上帝的恩典當中,.
有信心,有力量,來重新面對,.
求主繼續幫助我們,帶領我們走未來人生的道路..
\newpage



\section{路得記 1:19-22}
\label{sec:wdtTzXf4txM}
\textbf{2026年5月10日崇拜直播｜吳真真牧師｜告別浮淺,門徒進深:在痛苦中找到珍寶母親｜路得記1\_19-22}
\newline
\newline
連結: \href{https://youtube.com/watch?v=wdtTzXf4txM}{\texttt{https://youtube.com/watch?v=wdtTzXf4txM}} ~~~~ 語音日期: 2026-05-10
\newline
\newline
\hyperref[sec:U29QDtiVI5U]{\small{< < < PREV SERMON < < <}}
~
\hyperref[sec:index_chronic]{\small{[返順時目]}}
~
\hyperref[sec:9iKJWzoiFEk]{\small{> > > NEXT SERMON > > >}}
\newline
\newline
路得記 1:19-22
\newline
\begin{longtable}{cl}
\hline
\hline
章節 & 經文 (和合本修訂版)\\
\hline
1:19 & \begin{tabularx}{0.7\textwidth}{X} 於是二人同行，來到伯利恆。她們到了伯利恆，全城因她們騷動起來。婦女們說：「這是拿娥米嗎？」 \end{tabularx} \\ \\ \relax
1:20 & \begin{tabularx}{0.7\textwidth}{X} 拿娥米對她們說：「不要叫我拿娥米，要叫我瑪拉，因為全能者使我受盡了苦。 \end{tabularx} \\ \\ \relax
1:21 & \begin{tabularx}{0.7\textwidth}{X} 我滿滿地出去，耶和華使我空空地回來。耶和華使我受苦，全能者降禍於我。你們為何還叫我拿娥米呢？」 \end{tabularx} \\ \\ \relax
1:22 & \begin{tabularx}{0.7\textwidth}{X} 拿娥米從摩押地回來了，她的媳婦摩押女子路得跟她在一起。她們到了伯利恆，正是開始收割大麥的時候。 \end{tabularx} \\ \\
[1ex]
\hline
\hline
\end{longtable}
$^{1}$今日選讀的經訓是記在勞德記第一章19-22節.
於是二人同行來到伯利行.
他們到了伯利行.
合乘的人就都驚訝.
婦女們說:這是拿俄米嗎?.
拿俄米對他們說:.
不要叫我拿俄米,要叫我馬拉..
因為全能者使我受了大苦,.
我滿滿的出去,.
也或說使我空空的回來,.
也或說降我與我,.
全能者使我受苦..
既是這樣,.
你們為何還叫我拿俄米呢?.
拿俄米和他兒婦摩訶女子勞德.
從摩訶地回來到伯利行.
正是動手割大麥的時候..
.
各位弟兄姊妹平安.
也特別恭祝各位母親.
願上帝的恩典厚厚傾注於大家身上.
今天是母親節.
我們今天講的主題.
也是繼續接觸教會的主題.
就是告別浮淺門徒盡心系列.
回應上星期陳牧師跟我們分享.
發掘在失意中的珍寶.
他提到我們人生充滿苦難.
在苦難中我們必須找到正面的意義.
其中的方法是要我們正視這些痛苦.
不要逃避.
讓我們能夠停留在這些時間中.
去經歷神給我們的恩典.
找到當中的意義.
於是就啟發了我今天講的母親節訊息.
今天講的母親節訊息.
跟我們平時慶祝母親節有些不同.
因為我們很多時候說.
母親節快樂.
我們一起去慶祝.

$^{41}$是一個讓我們的子女向母親表達敬意的時候.
不過在我反思這個題目的時候.
我想到一件事.
很多時候在母親節.
我們都會忽略了.
或許有些母親.
只是處於一些失意,痛苦,難受的時候.
或許他們剛有親人離世.
或許在我們當中有些母親已經不在了.
或許他們的童年成長.
他們經歷不到母親的愛.
或許種種原因都是一些傷心失意的時候.
所以我今天想用一個特別的例子.
去講如何紀念母親節.
或在母親節中找一些新的意義.
其實今天的媽媽一點都不簡單.
是承受著很多壓力.
在很多調查中告訴我們.
現在的媽媽在情緒上.
所承受的困擾和壓力.
是遠遠比以前大.
特別是他們要面對社會或對他們的期待.
那種無力感,無助感.
的確讓他們處於艱難的時候.
我在想.
如果我們能夠一起去明白.
痛苦失意的人的處境.
當中會經歷一些什麼心路歷程.
從中得到智慧.
知道如何與他們相處.
甚至與自己好好相處.
好好照顧自己.
我相信這是一份很有意義的母親節禮物.
讓我們一同去禱告.
將我們的心交託給我們的主.
你愛我們.
你的厚恩圍繞著我們.
不過有些時候.
的確當我們面對失意,痛苦的時候.
我們真的感受不到這份恩情.

$^{81}$我們會迷失.
我們會痛苦.
我們想逃避.
但主,我們祈求你.
願意幫助我們.
打開我們的心去研讀你的話語.
去明白你在我們痛苦的過程中.
你的珍寶,你的美好心意.
求你幫助我們弟兄姊妹.
一同去領受這份額外隱藏的祝福和恩情.
我們這樣禱告交託.
奉主耶穌基督的聖名自祈求.
阿們.
我們今天會從一位舊約聖經人物.
拿訛米.
這位我們說她是一個傷失的母親.
她究竟如何在她極度痛苦和失意的過程中.
找到上帝的恩典.
找到出路.
這個故事大家很耳熟能詳.
但我相信有些朋友未必很熟悉這個故事.
所以我也要有背景的簡介.
我不讀這裡的經文.
大家可以透過螢光幕上顯示的經文.
我們一起去了解背景.
這個故事的主角是拿訛米.
她是在士思時代的一個傷失的母親.
這裡說她失去了丈夫和兒子.
路得記一開始就介紹她的背景.
原來拿訛米的故鄉是在伯利恆.
亦等於以法他.
原來因為當時饑荒.
所以她的丈夫帶著她和兩個兒子.
去到一個外邦人的地方摩崖.
希望能夠有足夠的糧食.
但後來她的丈夫死了.
於是她失去了丈夫.
成為了寡婦.
不過她有兩個兒子.
而這兩個兒子又各自娶了兩個摩崖女子為妻.

$^{121}$之後他們就在那裡定居了大約十年.
不過好景不常.
兩個兒子都死了.
她剛剛失去了丈夫.
現在又失去了兒子.
我們看第三到第五節會發現.
在這節經文的開始.
就將之前的主角以利米.
轉為她的妻子拿訛米.
我們就將目光注意到拿訛米的身上.
在經文裡我們看到三節和五節.
是一個漸進的描述.
第三節說以利米勒的丈夫死了.
然後第五節就說她的兩個兒子馬倫和基連也死了.
第三節說她只剩下自己和兩個兒子.
但第五節就說只剩下拿訛米自己.
而且強調她失去了丈夫和兒子.
她從兩個階段的變化.
由人妻到人母.
再變到一無所有.
她失去了一切賦予她身份的東西.
她一無所有.
其實拿訛米的處境真的一點也不容易.
失去丈夫令她變成寡婦.
再失去兩個兒子.
再加上她兩個兒子沒有後繼.
她沒有再生.
所以她是失下由虛.
以至她的家庭絕後.
拿訛米真的處於一個婦女當中最惡劣的處境.
她寄居外邦成為寡婦.
無兒無女.
亦沒有後代.
其實她的指望是甚麼呢?.
在她的人生裡的確找不到盼望.
拿訛米的處境令我想起一套電影.
不知道大家有沒有看過.
叫做《永遠的愛麗絲》.
有沒有人看過?.
有的舉手.

$^{161}$有少許朋友看過.
這套電影由朱利安摩亞飾演一位年輕的語言學教授.
她是一個媽媽.
不過我怎樣說她年輕呢?.
年輕不是因為她五十歲.
而是因為她五十歲就罹患了一個病.
是早發性的亞茲海默症.
即是說她五十歲開始已經腦退化.
已經患了失智症.
而她就開始過著一個失去的人生.
這個失去是慢慢逐漸的失去.
由她的認知,記憶,語言能力.
甚至是自我意識慢慢地失去.
但因為這個全方位的崩壞.
在這個過程裡她重新定義了自己.
再次與家人連結.
影片的最後告訴我們一件事.
縱然她的知識,腦的知性部分消逝.
但她的情感連結依然是人類存在的最後證明.
她一直以來的關係是很疏離,僵化的一個女兒.
是一個很反叛的女兒.
因為這個緣故.
大家同樣因為失去而產生共鳴.
繼而他們之間的關係得到修補和連結.
其實,弟兄姊妹.
我不知道你們是否覺得失去這件事.
是我們人生的常態.
可能我們會失去至親.
我們失去親人.
我們可以失去一份工作.
我們會失去我們的健康.
又或者我們內化一點.
會失去我們的尊嚴.
可能失去一些別人的信任.
失去我們的滿足感,榮譽感.
這些其實都會帶來痛苦.
會帶來給我們有哀傷,恐懼.
讓我們覺得前面沒有盼望.
這些情緒會帶來生活日常的變化.
讓我們有一種很強烈不能掌控的感覺.

$^{201}$不過,其實這一切都在告訴我們.
我們人生真正的需要.
我們需要依伴.
我們需要與人連結.
就好像撈外賣一樣.
最終他決定歸回.
我們看看第六至第七節.
讓我們一起中讀.
(梁美芬讀).
當初撈米一家因為面對大地饑荒離開自己的家鄉.
去一個有糧食提供的地方.
能夠給他們一個溫飽的地方去生活.
現在這一刻,摩訶不是饑荒.
不是沒有東西吃.
而且那兩個媳婦又是摩訶人.
其實他們大可以留在那裡.
繼續生活.
讓兩個媳婦繼續供養他們.
不過這一刻他們決定要歸回本鄉.
並不是因為糧食問題.
不是沒得吃.
而是他們要歸回眷顧他們百姓的神那裡.
歸回這個字其實我沒有列出來.
但它單單在第一章出現了十二次.
是路得記一個非常重要的主題.
除了剛才我們說第六第七節.
講到歸回或回猶大.
這裡出現了兩次之外.
很主要是回去歸回這個字.
就是要他們的媳婦回到自己的家鄉.
「你們回去吧」.
主要是用在這裡的對話.
其實歸回的意思是在某些過範當中回轉.
在一些過範當中.
在一些做錯了的事情上可以回轉.
其實當娜牛,米血兩個媳婦不要跟著自己.
要歸回自己摩訶的本鄉的時候.
除了一來他們說你們可以找回你們的夫家裡.
可以找回一個新夫家.
有一個好歸宿.

$^{241}$還有一句就是他們會歸回他們的神那裡.
其實是有一個宗教的意思.
最後他的媳婦娥爾巴決定了歸回她的娘家.
代表了她會歸回她摩訶的愛邦神那裡.
而她是離開了耶和華.
但這次勞德的歸回就是要回到那位眷顧他的耶和華那裡.
他要跟從耶和華.
「眷顧」這個動詞是指臨在,到訪的意思.
是形容耶和華和他子民的關係.
這個字是有立約的含義.
針對不同的情況其實會帶出不同的意思.
如果當他的子民手約的話.
這個字就會譯作「眷顧」.
但當他的子民背約的時候.
這個「臨在」就是帶來審判.
所以原來上帝的臨在.
他跟我們的同在.
是跟我們有互動的.
他的反應是關乎我們的態度.
從南岸禮在第六節所說.
我們知道他對上帝有清楚的認識.
上帝是臨在的.
上帝是跟他的子民同在的.
上帝是掌權的.
上帝是監察人心的.
當與他立約的子民順服的時候.
上帝就賜福供應.
事實上在事司的時代.
這群神的子民.
當他們遇到外敵有問題的時候.
他們就會向上帝呼求.
神就會興起事司去拯救他們.
所以納俄米就好像當時其他事司記裡的.
以色列文一樣.
很清楚上帝是一個掌權的上帝.
是一個賞善罰惡的上帝.
是一個在他的子民困難迴轉的時候.
會施行拯救的上帝.
於是當他知道自己一無所有.
一無所依的時候.

$^{281}$他還可以投靠誰呢?.
於是他就會回歸到.
眷顧他的子民的耶和華那裡.
他會歸回神的懷抱.
其實很多時候失去.
或者我們說我們遇到一些艱難的處境.
本來是上帝對我們信仰的考驗.
就好像伊利米勒一家一樣.
他們面對著生命一些挑戰和艱難.
我們很多時候因為怕在過程裡面那種痛苦.
我們會找一些比較容易的方法去解決.
比較容易的方法和出路.
既然糧食有問題.
我們所處的地方有饑荒.
我們就不要捱餓.
最有效的方法就是去外面找東西吃.
但是這個情況下.
他們就離開了眷顧他們的神.
而且還長居於外邦的地方.
其實我們這個告別浮淺門徒進心的系列.
是啟發自一本書叫做《門徒的EQ覺醒》.
書裡面提到一件事.
我們會因為失去帶來的痛苦.
令我們覺得失去掌控.
人生會進入悲傷.
而這些悲傷會擾亂了我們的人生.
所以我們很多時候都會趨吉避凶.
捨難取義.
接著我們會否定傷心的重要.
輕視了哀傷對我們生命的重要性.
我們很快就會做決定.
很快就會行動.
我們要過新的生活.
但是其實在我們心靈的深處.
還有很長的路要走.
所以在這個世界裡面.
有很多很快速,很容易.
我們以為會幫我們逃避這些傷痛的方式.
所以沉迷網絡世界,煲劇.
或者一些成癮的問題.

$^{321}$性的問題,賭的問題,毒的問題.
就在這個世界裡面一直充斥著.
我相信那河米知道其實他要歸回.
他要將以前一個自己做得不理想的決定.
去作出一個回轉的行動.
但是作出這個行動.
他的掙扎是沒有停止的.
讓我們一起從第十九至二十一節.
去明白一個哀傷失意的人.
心裡面究竟會有甚麼狀態呢?.
我們就會讀第十九節.
「於是二人同行來到伯利恆.
他們到了伯利恆.
合成的人都驚訝.
古裡門說:這是那河米嗎?」.
那河米一回到伯利恆.
他就發現這一群曾經.
和他一樣在伯利恆的鄉下.
他們都經歷過饑荒.
但在神的眷有下.
這群婦女仍然生存著.
這群婦女看到那河米的時候非常驚訝.
他們就很自然地說:.
「這是那河米嗎?」.
當然一個離開了家鄉這麼多年的人.
意想不到地出現.
經歷了這麼多.
這麼久沒有見面.
再加上經歷了這麼多災難.
他們的樣子可能都會有點不同.
他們那個時代沒有SK-II.
沒有醫美.
所以可能那河米在他們面前出現.
他們就稀稀地看著.
「咦?它好像是那河米」.
於是他們就說:「這是那河米嗎?」.
為城裡帶來一些騷動和回響.
但殊不知他們一說.
那河米的反應很大.
心裡一直抑壓著很大的傷感怨氣.

$^{361}$盡數出來.
而且這個怨氣聽起來都很可怕.
二十節.
「你們不要叫我那河米.
要叫我馬拉.
因為全能者使我受了大苦.
我滿滿地出去.
耶和華使我空空地回來.
耶和華降災與我.
全能者使我受苦.
既是這樣.
你們為什麼還要叫我那河米?」.
是的.
那河米的意思是美好.
有快樂和甘甜的意思.
其實當他們在說一件.
現在不再在他們身上發生的事的時候.
就觸動了他們的痛處.
將他們心情裡最真實的一面展露出來.
他用一個字「苦」來做總結.
所以他叫他們做馬拉.
因為這個字跟「苦」希臘文的字根是一樣的.
他說「全能者使他吃苦」.
為什麼他會吃苦呢?.
他的歸因是跟他的行惡有關.
因為第二十節的受苦就是跟行惡有關的苦.
是神降禍於他.
意思是神作證地控告我.
懲罰我.
所以他現在所經歷的一切.
都是耶和華對他的控訴.
他成為了被告.
他被耶和華定罪.
神要懲罰他.
就正如第十三節他跟他的媳婦說.
「耶和華伸手攻擊我」.
那河米認為這是神作在他身上.
恩爵控告他要去受苦.
是神的心意.
然後我們又會看為什麼他會這樣想呢?.

$^{401}$第二十至二十一節.
大家留意到藍色字的「全能者」.
他為什麼多次在《創世記》出現.
而這個字是跟列祖的事蹟很有密切的關係.
特別是用在這位上帝對列祖的祝福和咒詛.
表達了對上帝的威榮.
同時亦表達了上帝的審判.
按照著那河米現在的處境.
這個字「全能者」去形容神是最貼切的.
因為他現在的困難就是.
這個「全能者」會祝福人.
讓人後裔眾多.
但是他卻沒有後裔.
這個「全能者」會幫人在處境裡.
賜給他們一個新的名字.
現在那河米也得到一個新的名字.
但不是一個祝福的名字.
而是一個咒詛的名字.
他相信上帝掌管一切.
管理萬有.
他是幕後的導演.
如果沒有他的許可.
事情是不會發生的.
他知道上帝可以祝福人.
但現在他的祝福我沒有份.
對他來說是一個咒詛.
有時我們失去有苦難.
我們真的會問為甚麼.
有時的確是我們做錯事.
我們要負上我們的責任.
好像那河米他決定歸回耶和華.
這是一個非常正確的決定.
不過我們更加要小心.
當我們陷入這些困難.
我們要避免被一些罪咎感纏繞.
令我們在情緒裡果足不前.
我們的確會常常問.
為甚麼這個災難要加在我身上.
為甚麼我要受這些苦.
在四月份的《浸神》的院訓裡.

$^{441}$主題是「災禍苦難,神學反思」.
當中舊約教授黃福光博士分享.
其實聖經對於災禍苦難有很多解釋.
包括人的罪,神的審判.
人的愚昧決定,邪靈的攻擊.
他講到墮落的世界裡.
災難是常態.
神可能要透過這些災難教導我們.
或者神要我們為蒙召.
祂呼籲我們為神的國受苦.
但其實很多時候我們都未必有能力辯明.
到底在這個處境裡.
誰才是正確的答案.
而事實上有些時候我們會搞錯.
所以不要被那些解釋不到的罪咎感一路纏繞.
但我們如果盡心去找的時候.
我們會發覺不斷地問為甚麼.
其實是一個心靈痛苦的表達.
未必是要一個確實的答案.
而是需要理解和陪伴.
因為在這些失落裡.
我們跌落了心靈的深淵.
就好像撈俄米一樣.
他在二十一節裡說.
「我滿滿地出去,耶和華時空空地回來」.
撈俄米強調自己當初離開白令衡的時候是滿滿的.
其實他心裡想的是.
我們一家齊齊整整.
有丈夫,有兒子.
但現在他回來的時候.
他是空空地回來.
因為他已經沒有丈夫,沒有兒子.
生活無依.
他一直叫苦,一直抱怨.
他向他的鄉里那班婦女說出他最心痛的事.
就是他原來的幸福家庭已經化為烏有.
他已經失去一切.
其實當我們在痛苦的深淵的時候.
很多時候我們只覺得只有自己一個.
我們很孤單,很無助.

$^{481}$找不到出路.
我們很容易將焦點放在自己身上.
就好像在十一至十三節短短三節的經文裡.
撈俄米說了八個「我」字.
這段說話是撈俄米對著兩個媳婦.
要求他們不要和他一起回來猶大地.
他拒絕他們.
表面上他在解釋為甚麼.
是為你們好.
但其實他下面的潛台詞是說一件事.
他就是說「是的,你們是失去的.
你們都失去了丈夫.
你們是慘的.
不過你們仍然有盼望.
但我呢?.
我為你們的緣故甚是受苦」.
和修版的譯法是「我比你們更苦」.
當跌落情感深淵的時候.
孤單,無助,沒有出路.
就是內心的掙扎.
其實我們面對苦難,哀傷.
很多時候有以下的反應.
首先第一件事.
我們會否認.
我們會說「沒事的.
神是信實的神.
神是私人的神.
祂會保守我的.
沒事的,沒事的」.
但漸漸地我們可能會有憤怒.
我們會對傷害我們的人.
甚至我們對施害於我們的上帝感到憤怒.
我們會很生氣.
我們會有怨氣.
但下一步我們會有懇求.
我們希望神給我們轉機.
有新的突破.
設下新的出路.
希望有改變.
但當沒有轉變的時候.

$^{521}$我們又會陷入沮喪.
心情會跌落低谷.
感到無助.
慢慢地我們可能會接受.
這是事實.
這個歷程.
幫助我們一樣很重要的事.
就是尋找到意義.
讓我們一起去看第二十二節.
在過程中上帝如何施行他的恩惠.
第二十二節我們一起讀.
「那我米和他兒父摩訶類子路德.
從摩訶地回來的伯尼行.
正是動手割大物的時候」.
那我米的自白說完了.
我覺得當時的婦女.
這裡聖經沒有記載他們有任何回應.
因為被他們的苦水嚇倒了.
不過很奇妙.
聖經記載第二十二節.
就是一個最好的回應.
他告訴她.
「你不是空空的你」.
空空的回來不是真實.
他補充一句.
「他有路德和他一起」.
這個摩訶的女子路德.
是和他一起回來的.
路德的另一個主題.
除了說歸回之外就是恩惠.
上帝沒有叫那我米一直陷在深淵.
一直在他的痛苦裡.
他眷顧手藥.
跟從他的百姓.
他決定讓這一位跟隨他回去的摩訶女子路德.
要和那我米一起領受這個恩惠.
其實那我米知道上帝是恩代的神.
大家看看一章第八第九節.
那我米叫他的姨父回娘家的時候.
他是怨耶和華恩代他們.

$^{561}$他知道神是恩代的神.
但當他陷在深淵裡的時候.
他見不到這個恩代的神.
同樣可以施恩惠在他身上.
他因著他的內疚.
因著他的掙扎.
因著他的自怨自艾.
無視了身邊的人和事.
看不到上帝的恩惠.
事實上他並不孤單.
當他歸回上帝的時候.
神就用豐厚的恩惠恩代他.
不過這個恩代的方法並不是他自己想像的方法.
是用一個更奇妙的方法.
就是這節經文的下半部.
經文告訴我們.
他們回去的時候就是動手割大麥的時候.
這個就是恩惠的開始.
就是神施恩的開始.
這個時候路德就可以去找糧食.
不單單解決他們的生計.
而且在當中就找出上帝要給那我米生命的珍寶.
就是隱藏在以色列文當中的恩惠.
其實在以色列人的律法裡面.
大家會覺得律例裡面好像很嚴苛.
很嚴厲.
但其實裡面是隱藏著上帝很寶貴的恩惠.
他對孤兒寡婦的恩典.
以至他們可以取地上的莊稼得飽足.
寡婦的丈夫至親的親屬.
可以娶他為夫.
然後為他立候.
這些就是上帝在他的律例訓詞裡面.
所隱藏的寶藏.
直到第四章我們就會看到.
上帝的恩惠不單解決了維持生計.
更重要的是讓那我米經歷神母撇下的恩典.
他真的不是空空的來.
他的息父路德比起七個兒子更好.
而且路德和幫阿斯所生的俄彼得這個養子.

$^{601}$就是大衛的祖父.
是耶穌的祖先.
以至到救恩由他們出來.
弟兄姊妹.
信仰裡的確我們會一路去查經.
回來教會.
我們對上帝有一些認識.
不過很可惜有時我們會停留在知性的層面.
但當我們遇到失意苦難.
正正就是神給我們信仰考驗的時候.
當我們正視這些痛苦.
明辨和關注自己內心最深埋的情緒的時候.
我們就會看到一樣很重要的東西.
我們會知道自己有多軟弱.
有多需要神的眷顧.
然後去經歷他的真實.
到時我們信仰.
我們對神的認識.
就不是再是一句口號.
而是我們的經歷感悟當中.
對神真正的認識.
其實《路得記》這本書.
表面上神沒有直接參與人的活動之中.
但正正就是講一樣很重要的東西.
背後有神的掌管.
拿訛米深信神的眷顧.
他是眷顧我們的神.
亦會給苦楚思念我們的上帝.
縱然不明白.
但當拿訛米甚至波亞斯路德.
在不明白的處境中仍然守約.
盡上本分.
就會看到上帝口口恩惠的心意.
以及他美好的計劃.
當他經歷上帝觸摸他內心的深處.
經歷安慰,醫治,脫離苦境.
他沉著珍寶.
到這時候他沒有再問「為什麼」.
他不計之前的苦.
因為上帝本身就是他的滿足.

$^{641}$所以我在這篇信息中.
有幾個應用想跟大家分享.
第一,我們要認定神掌管一切.
苦難是不能避免的.
逃避或會令我們錯過了認識神和成長的機會.
第二,當我們靜候神的作為.
我們要靜靜地看祂的作為.
認定神的律例對我們的恩惠.
只要我們跟從神做好本分.
祂就會眷顧我們.
第三,痛苦容易令我們孤立自己.
千萬不要忽略.
其實神差遣很多小天使圍繞在你的身邊.
第四,我們陪伴一些在經歷失意,痛苦的人的身邊.
有時我們覺得很不容易.
很想快點把他拔出來.
但是我們要明白正在經歷痛苦的人.
心情的確有起伏.
終而決意歸回.
行動上回到伯利恆.
但是心情仍然有起伏.
我們願意學習用愛去同行.
我想分享一個小小的經驗.
是我的心聲,是我的歷程.
其實我們守本分等候神出手並不是容易的事.
我記得以前經歷了十多年家庭中的失序,失去秩序.
以致帶來很多痛苦,很多哀傷.
我記得有一個朋友送了一張這樣的卡給我.
很小張,你們看不到.
是一張英文的卡.
這張卡當時陪伴了我度過很多日子.
在苦難中陪伴著我,提醒著我.
我把它翻譯成中文.
這封信是一位朋友給我的一封信.
這封信是這樣寫的.
「今天早上你醒來的時候.
我透過你的窗戶為你帶來一輪圈例的日出.
希望能夠引起你的注意.
但是你卻匆匆離去,甚至沒有察覺.
後來我發現你和朋友們一起散步.

$^{681}$於是我讓你沐浴在溫暖的陽光下.
讓空氣中瀰漫著大自然的芬芳.
但你依然沒有注意到我.
當你經過的時候,我在雷雨中向你呼喊.
並在天空中畫出一道美麗的彩虹.
而你甚至沒有看一眼.
傍晚我將月光灑在你的臉上.
送來一陣清涼的微風,讓你安然入睡.
當你睡著的時候,我守護著你.
與你分享彼此的想法.
但你絲毫沒有察覺,我就在你身邊.
我揀選了你,希望你能盡快跟我說話.
在那之前,我會一直陪伴在你身邊.
我是你的朋友,我很愛你,耶穌」.
這封是耶穌給我寫的信.
這張卡我已經很久沒有找到.
我用了今年農曆年收拾東西時撿到.
我發覺很多時候我忘記了神.
其實在我的苦難中,祂叫我要仰望祂,留意祂.
在這些苦難中,祂就在我身邊.
祂很想我知道.
希望透過這一點點的經歷.
鼓勵自己在做媽媽的艱難中.
神一直看顧著我.
亦很想鼓勵我們當中的弟兄姊妹.
讓我們一起在苦難中找到上帝給我們的真報.
我們一起祈禱.
你的恩典沒有離開過我們.
只是我們有時陷在自己的苦難中.
陷在自己的情緒中.
我們不察覺你的同在.
不察覺你的保守.
不察覺你的恩典從沒有離開.
上帝,我們願意繼續守在自己的本份中.
按你的心意.
行你的真理.
等候你的出手.
讓我們能夠經歷在恩典中.
從你而來得到的滿足,平安和快樂.
我恭敬將眾弟兄姊妹交在你的手中.

$^{721}$求你與他們同行.
禱告交托.
奉主耶穌基督的聖名自祈求.
阿們.
謝謝大家.
\newpage



\section{約翰福音 13:1-35}
\label{sec:9iKJWzoiFEk}
\textbf{2026年5月17日崇拜直播｜陳冠成牧師｜告別浮淺,門徒進深:靈命成熟,以愛為尺｜約翰福音13\_1-35}
\newline
\newline
連結: \href{https://youtube.com/watch?v=9iKJWzoiFEk}{\texttt{https://youtube.com/watch?v=9iKJWzoiFEk}} ~~~~ 語音日期: 2026-05-18
\newline
\newline
\hyperref[sec:wdtTzXf4txM]{\small{< < < PREV SERMON < < <}}
~
\hyperref[sec:index_chronic]{\small{[返順時目]}}
~
\hyperref[sec:0dlSUYhlo5s]{\small{> > > NEXT SERMON > > >}}
\newline
\newline
約翰福音 13:1-35
\newline
\begin{longtable}{cl}
\hline
\hline
章節 & 經文 (和合本修訂版)\\
\hline
13:1 & \begin{tabularx}{0.7\textwidth}{X} 逾越節以前，耶穌知道自己離世歸父的時候到了。他一向愛世間屬自己的人，就愛他們到底。 \end{tabularx} \\ \\ \relax
13:2 & \begin{tabularx}{0.7\textwidth}{X} 晚餐的時候，魔鬼已把出賣耶穌的意思放在加略人西門的兒子猶大心裡。 \end{tabularx} \\ \\ \relax
13:3 & \begin{tabularx}{0.7\textwidth}{X} 耶穌知道父已把萬有交在他手裡，且知道自己是從神出來的，又要回到神那裡去， \end{tabularx} \\ \\ \relax
13:4 & \begin{tabularx}{0.7\textwidth}{X} 就離席站起來，脫了衣服，拿一條手巾束腰， \end{tabularx} \\ \\ \relax
13:5 & \begin{tabularx}{0.7\textwidth}{X} 隨後把水倒在盆裡，開始洗門徒的腳，並用束腰的手巾擦乾。 \end{tabularx} \\ \\ \relax
13:6 & \begin{tabularx}{0.7\textwidth}{X} 到了西門‧彼得跟前，彼得對他說：「主啊，你洗我的腳嗎？」 \end{tabularx} \\ \\ \relax
13:7 & \begin{tabularx}{0.7\textwidth}{X} 耶穌回答他說：「我所做的，你現在不知道，但以後會明白。」 \end{tabularx} \\ \\ \relax
13:8 & \begin{tabularx}{0.7\textwidth}{X} 彼得對他說：「你絕對不可以洗我的腳！」耶穌回答他：「我若不洗你，你就與我無份了。」 \end{tabularx} \\ \\ \relax
13:9 & \begin{tabularx}{0.7\textwidth}{X} 西門‧彼得對他說：「主啊，不僅是我的腳，連手和頭也要洗！」 \end{tabularx} \\ \\ \relax
13:10 & \begin{tabularx}{0.7\textwidth}{X} 耶穌對他說：「凡洗過澡的人不需要再洗，只要把腳一洗，全身就乾淨了。你們是乾淨的，然而不都是乾淨的。」 \end{tabularx} \\ \\ \relax
13:11 & \begin{tabularx}{0.7\textwidth}{X} 耶穌已知道要出賣他的是誰，因此說「你們不都是乾淨的」。 \end{tabularx} \\ \\ \relax
13:12 & \begin{tabularx}{0.7\textwidth}{X} 耶穌洗完了他們的腳，就穿上衣服，又坐下，對他們說：「我為你們所做的，你們明白嗎？ \end{tabularx} \\ \\ \relax
13:13 & \begin{tabularx}{0.7\textwidth}{X} 你們稱呼我老師，稱呼我主，你們說的不錯，我本來就是。 \end{tabularx} \\ \\ \relax
13:14 & \begin{tabularx}{0.7\textwidth}{X} 我是你們的主，你們的老師，尚且洗你們的腳，你們也應當彼此洗腳。 \end{tabularx} \\ \\ \relax
13:15 & \begin{tabularx}{0.7\textwidth}{X} 我給你們作了榜樣，為要你們照著我為你們所做的去做。 \end{tabularx} \\ \\ \relax
13:16 & \begin{tabularx}{0.7\textwidth}{X} 我實實在在地告訴你們，僕人不大於主人；奉差的人也不大於差他的人。 \end{tabularx} \\ \\ \relax
13:17 & \begin{tabularx}{0.7\textwidth}{X} 你們既知道這些事，若是去實行就有福了。 \end{tabularx} \\ \\ \relax
13:18 & \begin{tabularx}{0.7\textwidth}{X} 我不是指著你們眾人說的，我知道我所揀選的是誰；但是要應驗經上的話：『吃我飯的人用腳踢我。』 \end{tabularx} \\ \\ \relax
13:19 & \begin{tabularx}{0.7\textwidth}{X} 事情還沒有發生，我現在先告訴你們，讓你們到事情發生的時候好信我就是那位。 \end{tabularx} \\ \\ \relax
13:20 & \begin{tabularx}{0.7\textwidth}{X} 我實實在在地告訴你們，接納我所差遣的就是接納我；接納我的就是接納差遣我的那位。」 \end{tabularx} \\ \\ \relax
13:21 & \begin{tabularx}{0.7\textwidth}{X} 耶穌說了這些話，心裡憂愁，於是明確地說：「我實實在在地告訴你們，你們中間有一個人要出賣我。」 \end{tabularx} \\ \\ \relax
13:22 & \begin{tabularx}{0.7\textwidth}{X} 門徒彼此相看，猜不出他說的是誰。 \end{tabularx} \\ \\ \relax
13:23 & \begin{tabularx}{0.7\textwidth}{X} 門徒中有一個人，是耶穌所愛的，側身挨近耶穌的胸懷。 \end{tabularx} \\ \\ \relax
13:24 & \begin{tabularx}{0.7\textwidth}{X} 西門‧彼得就對這個人示意，要問耶穌是指著誰說的。 \end{tabularx} \\ \\ \relax
13:25 & \begin{tabularx}{0.7\textwidth}{X} 於是那人緊靠著耶穌的胸膛，問他：「主啊，是誰呢？」 \end{tabularx} \\ \\ \relax
13:26 & \begin{tabularx}{0.7\textwidth}{X} 耶穌回答：「我蘸一點餅給誰，就是誰。」耶穌就蘸了一點餅，遞給加略人西門的兒子猶大。 \end{tabularx} \\ \\ \relax
13:27 & \begin{tabularx}{0.7\textwidth}{X} 他接了那餅以後，撒但就進入他的心。於是耶穌對他說：「你要做的，快做吧！」 \end{tabularx} \\ \\ \relax
13:28 & \begin{tabularx}{0.7\textwidth}{X} 同席的人沒有一個知道耶穌為甚麼對他說這話。 \end{tabularx} \\ \\ \relax
13:29 & \begin{tabularx}{0.7\textwidth}{X} 有人因猶大管錢囊，以為耶穌是對他說「你去買我們過節所需要的東西」，或是叫他拿些甚麼給窮人。 \end{tabularx} \\ \\ \relax
13:30 & \begin{tabularx}{0.7\textwidth}{X} 猶大受了那點餅以後立刻出去。那時候是夜間了。 \end{tabularx} \\ \\ \relax
13:31 & \begin{tabularx}{0.7\textwidth}{X} 猶大出去後，耶穌說：「如今人子得了榮耀，神在人子身上也得了榮耀。 \end{tabularx} \\ \\ \relax
13:32 & \begin{tabularx}{0.7\textwidth}{X} 如果神因人子得了榮耀，神也要因自己榮耀人子，並且要立刻榮耀他。 \end{tabularx} \\ \\ \relax
13:33 & \begin{tabularx}{0.7\textwidth}{X} 孩子們！我與你們同在的時候不多了；你們會找我，但我所去的地方，你們不能去。這話我曾對猶太人說過，現在也照樣對你們說。 \end{tabularx} \\ \\ \relax
13:34 & \begin{tabularx}{0.7\textwidth}{X} 我賜給你們一條新命令，乃是叫你們彼此相愛；我怎樣愛你們，你們也要怎樣彼此相愛。 \end{tabularx} \\ \\ \relax
13:35 & \begin{tabularx}{0.7\textwidth}{X} 你們若彼此相愛，眾人因此就認出你們是我的門徒了。」 \end{tabularx} \\ \\
[1ex]
\hline
\hline
\end{longtable}
$^{1}$《約翰福音》十三章一至五節和十二至十七節.
在聖經的一百四十七至一百四十八頁.
逾越節以前,耶穌知道自己離世歸父的時候到了.
他既然愛世間屬自己的人,就愛他們到底.
在犯的時候,魔鬼已將賣耶穌的意思放在西門的兒子加略人猶大心裡.
耶穌知道父已將萬有交在他手裡.
且知道自己是從神出來的,又要歸到神那裡去.
就離席站起來,脫了衣服,拿了一條手巾束腰.
隨後把水倒在盆裡,就洗門徒的腳.
並用自己所束的手巾擦乾.
十二節.
耶穌洗完了他們的腳,就穿上衣服,又坐下對他們說.
「我向你們所作的,你們明白嗎?.
你們稱呼我夫子,稱呼我主.
你們說的不錯,我本來是.
我是你們的主,你們的夫子.
尚且洗你們的腳,你們也當彼此洗腳.
我給你們作了榜樣,叫你們照著我向你們所作的去作」.
我實實在在地告訴你們.
僕人不能大於主人,差人也不能大於差他的人.
你們既知道這事,若是去行,就有福了.
聖經讀不完,神賜福,上讀他話,如有遵行的人.
請姐妹平安.
我們先有個禱告,一起預備好我們的新心靈.
聖靈願你與我們每一個同在,開啟我們的心靈.
讓我們藉著聖經,藉著耶穌基督昔日所作的.
讓我們不單止經歷到上帝你的愛.
並且我們通過愛的實踐,讓我們的生命更加成熟.
可以成為美好的果子獻給上帝你.
為此我們恭敬獻上禱告.
奉基督耶穌你得勝的名字來祈求,Amen.
好,問問大家,大家都有去街市買過水果.
你會怎樣選擇水果到底熟不熟.
你會用什麼方法,過兩招給我看.
你會怎樣,先看價錢還是怎樣.
你會有些準則,首先你看一下它的顏色.
有沒有變色,譬如你買香蕉.
除非你不等於吃,不然你會買什麼顏色.
黃黃的,你不會買青色的.
或者你買鳳梨,你想馬上吃,你應該買什麼顏色.

$^{41}$都是黃色,你不會買青色的.
你看外殼有沒有變色,黃色,紅色,橙色.
你就知道它熟不熟.
又或者,還有第二招是什麼.
拿上去聞一聞.
譬如有些水果一熟就會有香味.
譬如是什麼,芒果.
或者有沒有人喜歡吃番石榴.
現在吹山的那些沒有那麼香.
但如果你發覺樹上那些小小的那些.
香到放在屋上,你會發覺到處都聞到.
你聞一下它的香味,你就知道它熟不熟.
又或者還有第三招是什麼.
這個不是每個人都可以.
有些街市一盤一盤做好就不要摸,不要選.
是這盤就這盤,指一指它給你.
但如果可以的話,你會發覺摸一摸它的果實.
怎樣辨別它熟不熟.
如果買奇異果就是.
軟一點點.
軟一點點,又不要一拿上去就嗶.
又不可以太軟.
軟硬適中的,你買水蜜桃或者買奇異果你就知道.
那我們的靈命怎樣分熟不熟.
怎樣分熟不熟.
其實和水果有點像,和你選水果.
你看上去那個人有沒有喜樂.
有沒有帶著神的愛,那個眼神會不同一點.
或者聞一下它,當然它很難聞到的.
我們好像一個形象化.
聞一下它有沒有基督的香氣.
你不要靠近去聞它,別人會罵你.
又或者其實一個成熟的靈命應該都是.
又不會太林志志.
又不會太硬邦邦.
應該是光柔兼備.
當然最實際聖經裡面說.
來衡量我們的屬靈生命是否成熟.
很多時候都和水果有關.
你有沒有結出聖靈的果子.

$^{81}$那九個的特質,記不記得.
通常我們記得頭幾個.
人愛喜樂和平,忍耐恩慈良善,信實溫柔節制.
其實他們不是說九種的果子.
不是說有九粒.
其實他們是說一個整體.
不過有九個的面向.
並且以人愛或者以愛為首,為核心.
其他那八種的果子.
其實可以說是愛的延伸.
所以其實就著今天的講題.
其實我們看到成熟的靈命.
其實和實踐愛是有關係的.
很直接的.
而在今天我們看的經文約翰福音裡面.
你會看到耶穌亦都示範.
如何維持一個很真實的愛的實踐.
經文裡面說.
耶穌知道自己離世.
即是說其實準備去受死受難的時候到了.
他愛世上屬於他的人.
留意這裡,耶穌說愛他們到底.
你見到約翰以耶穌對門徒的愛.
作為他受難受死的序幕.
一個愛到底的故事.
當吃如意戰晚餐的時候.
耶穌就突然離開他的坐直.
然後站起來脫去他的外衣.
然後用毛巾束腰.
然後倒水,倒一盆水來替門徒洗腳.
並且用束腰的毛巾幫門徒擦乾淨雙腳.
你見到整個行動.
為什麼說好像一個愛到底的象徵呢?.
當然我們首先要了解當時的一些背景.
洗腳在當時的時期.
其實是一種很慣常接待客人的風俗習慣.
因為你知道古代的路面.
沒有鋪了石屎.
通常都是一些泥路充滿沙塵.
而以前的人通常穿什麼鞋?.

$^{121}$涼鞋或者草鞋.
即是腳趾露出來.
所以你不難想像.
走過一些路腳一定會沾滿了很多塵土.
甚至乎污垢.
所以當時每家每戶基本上有一個習慣.
就是在門口放一個洗腳盆.
無論是客人還是自己人.
你入屋之前一定要洗乾淨你的腳.
洗乾淨那些塵土污糟的東西.
你才可以入屋.
所以其實一定不會吃飯後就去洗腳.
我們要先明白這個典故.
換句話來說.
其實耶穌在吃飯的時候.
突然起來替門徒洗腳.
其實不是處理衛生問題.
很明顯.
而是要在祂的離世之前.
藉著這個我們說不尋常的舉動.
就好像師父和他的徒弟.
要好好道別.
透過這個動作.
我們可以好好地說再見或者說暫別.
讓門徒知道原來他們的夫子耶穌.
是怎樣愛他們.
他們就是耶穌愛到底的對象.
事實上我們都知道猶太人.
他們覺得幫人洗腳.
是一個極之卑賤的事情.
基本上他們不會自己來.
一定是交給奴僕來處理.
但你見到耶穌卻不顧他們的身世.
他於尊降貴.
逐一來替他們的門徒洗腳.
你見到身為老師的他竟然自願.
現在變聖所有門徒的僕人.
或者說奴僕.
全因為就像經文所說.
一個字.

$^{161}$愛並且愛到底.
耶穌視他們為自己人.
所以他們才可以享受到.
耶穌這個尊貴特別的服侍.
其實不難理解.
除非你是做陪月.
或者從事幼兒照顧.
否則你會否喜歡幫其他人的寶寶來換尿片.
喜歡嗎?.
有沒有人很喜歡幫.
不是你自己的寶寶.
是幫你隔壁左右的弟兄姊妹.
有尿片要換就幫你換.
不會吧?你會覺得是骯髒的.
你不會見到有人嘔你主動去清理.
除非那個是你生的小朋友.
其實我們之所以願意替自己所生的.
自己所愛的寶寶來清潔.
是因為什麼?.
是因為愛.
一個字就是這麼簡單.
你所生的.
那個是你自己人.
你因為愛.
所以你甘願成為一個僕人.
承擔別人覺得很骯髒.
很厭惡的服侍工作.
其實耶穌基督就是這樣.
所以你會見到.
為什麼耶穌透過替門徒洗腳的動作.
經文中描述.
是一個愛到底的行動.
一個意義.
我們又看看當時的門徒有什麼反應.
面對著耶穌洗腳的服侍.
經文只是集中記載彼得.
彼得第一個反應是.
主啊,你要洗我的腳嗎?.
第六節這裡說.
然後再重複.

$^{201}$你不可以洗我的腳.
你一定不可以.
很明顯是因為尷尬.
特別是男生和男生.
如果突然間.
珍珍牧師不在.
如果珍珍牧師突然說道下來.
替你洗腳.
你也會退縮.
為什麼呢?.
更何況如果是一個男人和另一個男人.
突然洗Matthew的腳,你也會馬上退縮.
所以其實出於尷尬是很自然的.
而且始終那個是自己的老師.
尊敬的老師.
怎麼可能讓耶穌這麼委屈.
所以彼得拒絕了讓耶穌洗腳.
但耶穌就回答了.
如果你不讓我洗腳.
你就與我無份.
按照上文是無份.
這裡的「份」其實很自然就是說.
耶穌之後的拯救.
因為我們知道經文稍後發展.
就是耶穌會在十字架上.
透過祂的犧牲,寫命.
將人從罪惡裡徹底拯救出來.
當然也包括了他的門徒.
所以其實洗腳.
這裡說的那個「份」.
是說與耶穌的救恩有份.
換句話來說.
替門徒洗腳.
也是一個象徵性的行動.
就是凡被耶穌洗過腳的.
都象徵著他們的罪污.
也會被耶穌的寶血洗潔淨.
所以耶穌其實是與彼得說.
彼得如果你不讓我洗你的腳.
你就與我的愛.

$^{241}$你就與我的救恩無關係.
你就不是屬於我的人.
所以彼得當然不會完全明白.
耶穌當時這個動作.
到底是什麼意思.
不過他明白一件事.
就是洗腳對他來說是有好處的.
所以你看到彼得接著怎樣.
不要只是洗腳.
還要洗什麼.
洗頭.
洗全身.
總之露出來的那些.
耶穌你也幫我洗一洗.
大概他認為.
耶穌洗的範圍越多.
自己的福氣就越大.
當然這是誤解了.
耶穌洗腳的核心意義.
所以耶穌也馬上糾正他.
不用這樣.
只是洗腳,你全身就潔淨了.
接著你又看到.
耶穌對著所有的門徒說.
祂說你們都是潔淨的.
因為耶穌你知道稍後.
祂會逐一替所有的門徒洗腳.
也象徵著他們的罪惡.
會被耶穌一一洗淨.
成為天國裡的自己人.
但是你留意耶穌下一句說什麼.
耶穌說.
門徒不都是潔淨的.
因為耶穌原本知道.
當中有一個人會出賣耶穌.
這個我們當然知道是猶大.
但是你會發現.
此時此刻.
耶穌仍然當猶大是自己人.
仍然甘願為猶大洗腳.

$^{281}$很奇怪吧.
換作是你.
如果你是耶穌.
你明知猶大稍後會出賣你.
很大機會.
你會不會替祂洗腳.
你會不會愛祂到底.
大概不會吧.
很真實.
我們通常應該怎麼做.
如果換作是你.
你一定是先洗乾淨猶大.
叫祂買豉油也好.
買水也好.
洗乾淨之後.
接著怎樣.
接著就替其他11個.
來來來.
便宜你便宜你.
我替你洗腳.
正常應該是這樣吧.
沒有理由.
明知那個人想謀害自己.
我為什麼還要對他這麼好.
不正常.
但是.
經文就是要打破這種常理.
經文刻意告訴我們.
原來打算出賣耶穌.
傷害耶穌的猶大.
竟然都是耶穌.
他甘心服侍.
甘心去愛的對象.
你看到耶穌基督的愛.
是不是很不可思議.
事實上.
我們回到當時的場景.
當耶穌對所有門徒說.
你們不都是潔淨的.
其實我相信.

$^{321}$基本上在場所有人.
都聽不明白耶穌這句說什麼.
除了一個人.
是誰.
誰大概聽得明白耶穌說什麼.
就是想出賣耶穌那個.
其實猶大會心知肚明.
耶穌是在說祂.
祂知道耶穌已經在洞悉祂.
圖謀不軌.
但是耶穌卻沒有直接揭發祂的惡行.
耶穌的心意其實很明白.
就是怎樣.
我不直接揭穿你.
我是期望你.
收手.
我是期望猶大你收手.
這一刻你仍然是我愛的對象.
我希望你可以收手回轉.
當然我們最後知道.
猶大始終沒有領耶穌的情.
他仍然最後決定.
要出賣耶穌.
我們說糟蹋了耶穌這份對他的愛.
所以弟兄姊妹.
停在這裡我們看到.
耶穌為門徒洗腳.
其實是向我們顯示.
什麼叫做好像經文所說.
愛到底.
愛到底不是純粹說在時間上.
直到那個人死那一刻.
仍然愛對方到底.
愛到底其實這裡有個更深遠的意思.
就是愛到一個極致.
極致到怎樣.
連害你的那個都愛.
連猶大.
我們說這個這麼壞.
出賣師父的人.

$^{361}$這個敗類.
耶穌說那一刻.
都是他愛到底的對象.
耶穌願意給他機會.
希望他悔改.
盡他的努力去挽回猶大.
事實上你會看到.
再看闊一點.
耶穌不只是洞悉猶大的問題.
其實耶穌也知道.
稍後還有什麼人會離開他.
耶穌稍後會預告彼得.
你不要裝聰明.
你之後三次不認我.
其他的門徒呢?.
當在克西瑪利園.
耶穌被捕的時候.
他們全部都四散離棄耶穌.
但是.
耶穌仍然沒有放棄他們每一個.
仍然甘心樂意逐一為他們洗腳.
向我們展示出.
原來耶穌體一個愛的關係.
到底是怎樣.
即使對方可能會有軟弱.
會有跌倒.
甚至不是很可靠.
還說有機會背離自己.
傷害自己.
但是耶穌仍然堅持.
接納寬恕他們.
愛他們到底.
你會見到愛驅使耶穌.
他甘願承擔被拒絕.
受傷害的風險.
這個為之愛到底.
我們也有類似的經歷.
特別在座已經結婚的弟兄姊妹.
回想起.
你當日怎樣和你的配偶示愛.

$^{401}$還記得嗎?.
可能很久了.
你怎樣當時確認.
彼此是男女朋友的關係.
其實你有沒有發覺.
你需要很大的勇氣.
特別可能是弟兄.
通常是弟兄示愛.
很少姊妹示愛.
姊妹比較含蓄.
弟兄你有沒有發覺.
很大的勇氣.
當你向對方示愛的時候.
因為在你示愛之後.
其實你不知道對方有什麼反應.
我們當然會期待.
對方一定會喜歡我.
答應我的.
這是我們一廂情願的想法.
但對方亦可能很嚴肅地.
拒絕你.
或者很冷淡地對待你.
走了.
你會看到其實.
向別人展示愛.
其實就等於你毫無保留.
展露你脆弱的一面.
你攤開了出來.
沒有防範.
等對方回應.
你也期待.
可能預期會在情感上.
受到傷害.
除非你不表達愛意.
否則你一定會冒險.
當然你不去表達愛意.
不去示愛.
你永遠都不會被別人傷害.
不會被別人拒絕.
但也可能因為這樣.

$^{441}$你就找不到你旁邊的摯愛.
所以你看到愛會驅使你.
多了很多勇氣.
來做一些很傻瓜的事.
我還記得我.
差不多二十多年前.
我追求女生的時候.
我向太太表達愛的時候.
其實我都很忐忑.
我以為一定成功.
但又很怕被拒絕.
你知道拒絕了就很尷尬.
日後大家怎樣相見.
當時還是同事.
會不會嚇走對方.
對方辭職之後.
結果.
你猜我成不成功.
我當然不是one take pass.
你都記得我說過.
我一示愛的時候.
對方就送了一段經文給我.
信與不信不能跟父親騙.
即是說yes or no.
即是no.
即是你沒有行,節省一點.
很慘受傷害.
當然最後沒有放棄.
可能大家都聽過.
最後我又再次鼓起勇氣.
隔了一段時間.
又再去示愛的時候.
那次就成功打動對方.
但弟兄姊妹你會看到.
原來愛人示愛是有風險的.
你見到耶穌明知道.
不是所有人都會接受祂的愛.
有些人一開始就對祂嗤之以鼻.
有些人一開始就想除滅祂.
又或者你見到有些人.

$^{481}$起初很擁戴耶穌.
跟隨祂.
但之後開始覺得.
耶穌都不是很幫到忙.
又不是很照我的意思做.
於是就離棄耶穌.
當然有些人發誓.
我至死不渝.
我跟硬你,耶穌.
但到了禍患臨到的時候.
卻選擇離棄祂.
但你發覺即使如此.
耶穌都很堅持.
祂會冒險去愛人.
不在於人對祂怎樣.
而在於愛就是祂的本質.
耶穌仍然選擇去愛.
世上的人到底.
特別在這段經文你會看到.
在臨別之前.
祂將一個最重要的教導.
傳授給祂所愛的門徒.
其實就像我們中國人看的功夫片.
師傅要跟徒弟入室弟子道別的時候.
通常會怎樣?.
「我過了最後一招給你」.
「我將這個大招傳給你」.
「我只是給自己人」.
「例如上興趣班的我不會給他們」.
「入室弟子我才會給他們」.
其實情況就是這樣.
我想大家一起讀.
最後的段落.
我截錄了14至15和34至35.
或者你看PowerPoint都可以.
預備起.
「我賜給你們一條新命令」.
「乃是叫你們彼此相愛」.
「我怎樣愛你們」.
「你們也要怎樣相愛」.

$^{521}$「你們若有彼此相愛的心」.
「眾人因此就認出彼此愛的門徒了」.
你看到耶穌說得很清楚.
怎樣讓人認出.
我們是祂的門徒.
看到那個成熟的靈命.
就是愛人.
你通過愛你身邊的人彼此相愛.
你就展現出你的靈命.
不過回想起.
較早之前.
我們還記得耶穌預告自己.
即將受死受難的時候.
還記得門徒當時的反應是怎樣?.
他們不是為耶穌分憂.
爭著認誰做最大.
所以其實耶穌藉著彼此洗腳的吩咐.
來要求門徒.
要愛就首先要放下你真正的心.
並且要效法他.
你寫記放下身段.
將耶穌對你的愛.
你回饋給身邊的人.
透過彼此洗腳的實踐.
其實就是學習彼此相愛.
並且怎樣愛到底.
愛到極致.
事實上耶穌本人就是這樣.
愛門徒到底.
甚至稍後不惜愛到十字架上.
犧牲自己的生命.
來成全門徒的生命.
並且不單在十字架上.
你會看到當耶穌復活的時候.
他亦三次問彼得.
你愛我嗎?.
大家有沒有留意當耶穌問你愛我嗎?.
其實這三個字的時候.
其實耶穌是首先向彼得示愛.
如果耶穌不是首先仍然愛彼得的時候.

$^{561}$他不會說這三個字.
就好像一對情侶吵架的時候.
如果其中一方問另一個.
你還愛不愛我?.
其實就是代表他仍然愛對方.
如果不是他不會這樣說.
所以當耶穌三次問彼得.
你愛我嗎?.
其實背後有一個很深的含義.
就是說彼得雖然你三次不認我.
你不愛我.
但是我從來沒有不認你.
到現在我仍然愛你.
並且我不止一次饒恕你,接納你.
我還會一次,兩次,三次.
無數次愛你到底.
我還要把我的羊托付給你.
那麼彼得,你願意愛我嗎?.
弟兄姊妹,是否容易說出來.
其實耶穌這個愛的呼籲和邀請.
其實是很難的.
正所謂中國人有時都說.
一次不忠又如何?.
百次不容.
你被人害過,被人離棄過.
你不是那麼輕易去相信一個人.
你要有很大的愛意.
所以你看到耶穌是因為愛的緣故.
似乎做了一件很多人都覺得是很傻的事.
竟然縱容一個曾經失信,背棄他的人.
但是弟兄姊妹,耶穌就是叫我們這樣效法他.
耶穌何嘗不是這樣對我們?.
耶穌何嘗不是這樣愛那個經常失信於他的我們?.
所以這段經文提醒我們.
為了愛的緣故.
我們又願不願意效法耶穌.
做一些別人認為很傻瓜的事情.
譬如像耶穌那樣.
為了愛的緣故.
給多次機會那些犯錯的人.

$^{601}$讓軟弱跌倒的人.
他們仍然有改善成長的空間.
又或者為了愛的緣故.
我們會不會少一點記性?.
不要經常翻舊帳.
記住別人怎樣傷害過你.
別人怎樣對過你不好.
又或者為了愛的緣故.
會不會少一點計較?.
包容身邊的人不容易相處.
有時也挺難受的.
但是我們會不會為了神的緣故.
嘗試用神的愛去接納他?.
又或者為了愛的緣故.
我們會不會甘願犧牲,吃虧一點?.
在別人的需要上.
我們不只是看到.
我有什麼吃虧.
而是看到我有什麼責任.
我怎樣按著神的供應.
來扶他一把?.
又或者為了愛的緣故.
像耶穌那樣.
放膽去冒險.
也就是說可能你付出真誠,付出愛.
未必會馬上收到正面的回報.
但是你堅持為了上帝的緣故.
你這樣去做.
正如C.S. Lewis 魯易斯也說過.
他怎樣看待愛呢?.
愛其實就是你遇上了傷害.
很實際吧?.
有時有些流行曲.
有愛就會有痛.
好像掛鉤.
C.S. Lewis 表示.
愛的付出其實不是一定賺的投資.
因為實踐愛一定存在著風險.
事實上當我們選擇去愛人的時候.
你就可能會令自己很心煩,很糾結.

$^{641}$甚至你難過,會受傷.
有時會很想退縮.
不想繼續去愛.
但這個時候你就會明白.
為什麼耶穌仍然要我們堅持.
去實踐愛,彼此相愛.
因為唯有這樣.
我們才真正體會到.
愛的本質是什麼.
開始發現.
其實上帝就是這樣一直愛我們.
當我們愛到受不了,很想放棄的時候.
你回想起上帝.
其實他就是這樣堅持.
唯有你實踐愛的時候.
你才明白上帝怎樣恆久忍耐我們.
唯有通過愛的實踐.
我們才明白.
原來耶穌就是這樣愛我們到底.
最後我們會否因為愛的緣故.
好好管理自己的言行.
事實上成熟的靈命.
應該好像成熟的水果.
吃下去會怎樣.
會讓人覺得甜.
會讓人覺得滋潤.
會讓人覺得愉快.
如果不成熟的水果吃下去會怎樣.
會澀,會苦.
會不舒服想吐出來.
我們的靈命.
有時會否都是這樣.
我曾經聽過一個.
很快人心醒的一段說話.
有時和人相處最大的毛病是什麼.
就是我們將一些最差的脾氣.
最壞的情緒.
留給你最親密的人.
你覺不覺得都很好提醒.
越近的人有時會怎樣.

$^{681}$你會越被他刺激得厲害.
反而你對陌生的人好多.
我們有時將忍耐,寬容和溫柔.
留給那些陌生的外人.
這個情況在現實裡面很常見.
可能有時我們會在外面.
人人都叫我們你是好先生.
你是好女士.
脾氣又好又有耐性又平易近人.
但一回到家可能就不是這樣.
可能你會變成另一個人.
因為熟悉的緣故.
因為好多原因.
你會對你最熟悉的人.
最親近的人.
你會好多懵憎.
你會好多不滿.
你會甚至乎躁底.
忘記了原來上帝設立家庭.
是一個講愛的地方.
是一個孕育愛,實踐愛的場景.
所以弟兄姊妹其實既然.
靈命的成熟和我們實踐愛息息相關.
我們的言語亦需要帶著好多甜味.
好多愛意.
讓人感受到被造就,被滋潤.
而不是被冒犯,被攻擊.
所以很值得透過這段經文重新檢視.
無論在家庭,職場或教會.
我相信這三個場景.
都佔據了我們絕大部份的生活時間.
在裡頭我們如何像耶穌所說.
去實踐愛,彼此相愛.
我們愛的言語到底足不足夠.
又或者我們一些相處的模式.
會不會難阻我們去實踐愛.
好像經文最後這裡說.
有三樣東西我們需要去提防.
這三句說話很相似.
第一句是什麼.

$^{721}$無話可說.
不是沒話可說.
是在說什麼.
是在說一段關係裡.
再沒有什麼深入的交流.
對話的內容只是很務實,很功能.
譬如你見到小朋友上學.
第一句就問什麼.
做好功課了嗎.
可能回到家就問吃飯了嗎.
很多事務性的東西.
或者說大部份的時間.
只是享受回到家自己滾手機.
那個Me time.
躲在房間裡打機的時間.
開始覺得不和你身邊很親密的人交談.
都沒什麼所謂.
這是第一樣要留意的.
第二個就是.
沒有好話說.
就是說說話沒有好話.
一出口可能就有很多的諷刺.
批評,潑冷水.
讓人總是覺得.
為什麼我做什麼.
都好像被你否定,被你彈.
感受不到我們的愛意.
感受不到我們對對方那種正能量.
最後一句就是.
好的說話忘了說.
好話沒說.
可能習慣了沉默.
或者說很吝嗇.
來對別人表達稱讚.
對別人表達肯定.
對別人鼓勵.
甚至對別人關心.
這三樣東西我們要留意.
因為都很破壞我們去實踐愛.
如果真的有的話.

$^{761}$求主幫助我們去打破.
這些相處的局面.
為了愛神,愛人的緣故.
我們需要改變自己.
無論在你工作之前.
回到辦公室之前.
回到家之前.
或者回到教會之前.
將你的情緒,將你的言行.
將你的舌頭交給耶穌管理.
讓耶穌的愛充滿你.
也讓聖靈去管束.
調教我們的言行.
去解除那些可能會叛倒人.
傷害人的說話.
也致力活出那些會榮耀神.
造就人的言行.
求主幫助我們.
讓我們依靠祂繼續的成長.
滿有耶穌基督.
祂長大成人的新娘.
我們一起禱告.
因主我願你幫助我們.
讓我們藉著今天.
耶穌基督你這段愛到底的經文.
我們重新去回想.
耶穌其實你是怎樣愛我們每一個.
你不單單只愛到一個地步.
為我們犧牲,捨命.
你更加愛我們到底.
就是到如今.
我們可能已經認順你.
但我們仍有虧欠,仍有失信的時候.
有軟弱的時候.
耶穌說你總是一次又一次.
用不同的方式來挽回我們.
再一次將我們帶到你的根前.
帶到你的愛裡面.
重新的學習.
怎樣去愛上帝.

$^{801}$愛我們身邊的人.
主呀幫助我們.
我們確實很想.
立定心智去跟隨你.
然而我們亦都會.
好像門徒一樣.
我們會有軟弱.
有愛到疲倦.
愛到很煩躁.
愛到不想繼續的時候.
所以求主幫助我們.
讓我們重新回到你的懷裡.
讓你的愛包圍我們.
擁抱我們.
以致我們重新有力量.
去向別人實踐愛.
幫助我們無論在家裡.
無論在工作環境.
又或者在教會當中.
讓我們的言語.
能夠被你潔淨.
讓我們能夠謹慎.
不說那些會拆毀人.
不說那些沒有造就性的話.
讓我們的言語.
就好像一個成熟的果子.
不只是榮耀你.
亦都能夠成為別人的滋潤.
成為別人的甘甜和養分.
願主幫助我們.
讓我們藉著彼此相愛.
來彰顯我們成熟的靈明.
亦都顯明我們對上帝的愛.
恭敬獻上禱告.
奉基督耶穌你.
得勝的名字來祈求.
Amen.
\newpage



\section{約翰福音 19:1-22}
\label{sec:0dlSUYhlo5s}
\textbf{2026年5月24日崇拜直播｜盧國偉先生｜在黑暗中靜待黎明｜約翰福音19\_1-22}
\newline
\newline
連結: \href{https://youtube.com/watch?v=0dlSUYhlo5s}{\texttt{https://youtube.com/watch?v=0dlSUYhlo5s}} ~~~~ 語音日期: 2026-05-24
\newline
\newline
\hyperref[sec:9iKJWzoiFEk]{\small{< < < PREV SERMON < < <}}
~
\hyperref[sec:index_chronic]{\small{[返順時目]}}
~
\hyperref[sec:dKG666yV4KA]{\small{> > > NEXT SERMON > > >}}
\newline
\newline
約翰福音 19:1-22
\newline
\begin{longtable}{cl}
\hline
\hline
章節 & 經文 (和合本修訂版)\\
\hline
19:1 & \begin{tabularx}{0.7\textwidth}{X} 於是，彼拉多命令把耶穌帶去鞭打了。 \end{tabularx} \\ \\ \relax
19:2 & \begin{tabularx}{0.7\textwidth}{X} 士兵用荊棘編了冠冕，戴在他頭上，給他穿上紫袍， \end{tabularx} \\ \\ \relax
19:3 & \begin{tabularx}{0.7\textwidth}{X} 又走到他面前，說：「萬歲，猶太人的王！」他們就打他耳光。 \end{tabularx} \\ \\ \relax
19:4 & \begin{tabularx}{0.7\textwidth}{X} 彼拉多又出來對眾人說：「看，我帶他出來見你們，讓你們知道我查不出他有甚麼罪狀。」 \end{tabularx} \\ \\ \relax
19:5 & \begin{tabularx}{0.7\textwidth}{X} 耶穌出來，戴著荊棘冠冕，穿著紫袍。彼拉多對他們說：「看哪，這個人！」 \end{tabularx} \\ \\ \relax
19:6 & \begin{tabularx}{0.7\textwidth}{X} 祭司長和聖殿警衛看見他，就喊著說：「釘十字架！釘十字架！」彼拉多對他們說：「你們自己把他帶去釘十字架吧！我查不出他有甚麼罪狀。」 \end{tabularx} \\ \\ \relax
19:7 & \begin{tabularx}{0.7\textwidth}{X} 猶太人回答他：「我們有律法，按照律法，他是該死的，因為他自以為是神的兒子。」 \end{tabularx} \\ \\ \relax
19:8 & \begin{tabularx}{0.7\textwidth}{X} 彼拉多聽見這話，越發害怕， \end{tabularx} \\ \\ \relax
19:9 & \begin{tabularx}{0.7\textwidth}{X} 又進了總督府，對耶穌說：「你是哪裡來的？」耶穌卻不回答。 \end{tabularx} \\ \\ \relax
19:10 & \begin{tabularx}{0.7\textwidth}{X} 於是彼拉多對他說：「你不對我說話嗎？難道你不知道我有權柄釋放你，也有權柄把你釘十字架嗎？」 \end{tabularx} \\ \\ \relax
19:11 & \begin{tabularx}{0.7\textwidth}{X} 耶穌回答他：「若不是從上頭賜給你的，你就毫無權柄辦我，所以，把我交給你的那人罪更重了。」 \end{tabularx} \\ \\ \relax
19:12 & \begin{tabularx}{0.7\textwidth}{X} 從此，彼拉多想要釋放耶穌，無奈猶太人喊著說：「你若釋放這個人，你就不是凱撒的忠臣。凡自立為王的就是背叛凱撒。」 \end{tabularx} \\ \\ \relax
19:13 & \begin{tabularx}{0.7\textwidth}{X} 彼拉多聽見這些話，就帶耶穌出來，到了一個地方，叫「鋪華石處」，希伯來話叫厄巴大，就在那裡坐堂。 \end{tabularx} \\ \\ \relax
19:14 & \begin{tabularx}{0.7\textwidth}{X} 那日是逾越節的預備日，約在正午。彼拉多對猶太人說：「看哪，你們的王！」 \end{tabularx} \\ \\ \relax
19:15 & \begin{tabularx}{0.7\textwidth}{X} 他們就喊著：「除掉他！除掉他！把他釘十字架！」彼拉多對他們說：「要我把你們的王釘十字架嗎？」祭司長回答：「除了凱撒，我們沒有王。」 \end{tabularx} \\ \\ \relax
19:16 & \begin{tabularx}{0.7\textwidth}{X} 於是彼拉多把耶穌交給他們去釘十字架。 \end{tabularx} \\ \\ \relax
19:17 & \begin{tabularx}{0.7\textwidth}{X} 耶穌背著自己的十字架出來，到了一個地方，名叫「髑髏地」，希伯來話叫各各他。 \end{tabularx} \\ \\ \relax
19:18 & \begin{tabularx}{0.7\textwidth}{X} 他們就在那裡把他釘在十字架上，還有兩個人和他一同被釘，一邊一個，耶穌在中間。 \end{tabularx} \\ \\ \relax
19:19 & \begin{tabularx}{0.7\textwidth}{X} 彼拉多又寫了一個牌子，釘在十字架上，寫的是：「猶太人的王，拿撒勒人耶穌。」 \end{tabularx} \\ \\ \relax
19:20 & \begin{tabularx}{0.7\textwidth}{X} 有許多猶太人念這牌子，因為耶穌被釘十字架的地方靠近城，而且牌子是用希伯來、羅馬、希臘三種文字寫的。 \end{tabularx} \\ \\ \relax
19:21 & \begin{tabularx}{0.7\textwidth}{X} 猶太人的祭司長就對彼拉多說：「不要寫『猶太人的王』，要寫『那人說：我是猶太人的王』。」 \end{tabularx} \\ \\ \relax
19:22 & \begin{tabularx}{0.7\textwidth}{X} 彼拉多回答：「我寫了就寫了。」 \end{tabularx} \\ \\
[1ex]
\hline
\hline
\end{longtable}
$^{1}$讓我們以聆聽聖道繼續我們的敬拜.
今日我會讀出的經文是在新約聖經的約翰福音十九章.
約翰福音第十九章十二到十六節.
從此彼拉多想要釋放耶穌.
無奈猶太人喊著說:你若釋放這個人就不是該殺的中神.
凡以自己為王的就是背叛該殺了.
彼拉多聽見這話就帶耶穌出來.
到了一個地方名叫鋪華錫柱.
希伯來話叫阿巴達就在那裡坐堂.
那日是預備雨節的日子若有政誤.
彼拉多對猶太人說:看吶,這是你們的王.
他們喊著說:除掉他,除掉他,釘他在十字架上.
彼拉多說:我可以把你們的王釘十字架嗎?.
祭司長回答說:除了該殺,我們沒有王.
於是彼拉多將耶穌交給他們去釘十字架.
好,頂子妹,早安.
很特別,我去很多不同的教會講道.
很少見到有首領隊.
所以為教會感恩.
因為能夠建立一隊首領隊並不容易.
因為一來要有一定的音樂設備.
另外也要有人懂得玩首領.
我第一次認識首領就是在二十多年前.
在長洲當時在神學院認識到首領這個樂器.
其實你見到首領很特別.
因為每一個人只負責一至兩個音符.
你想想剛才不計指揮,有九個人.
沒有其中一個人,整個樂章就不完美.
其實教會也一樣,我們也是不同的性情.
我們中間有不同年紀,有不同的背景.
我們除了宗教信仰同樣之外.
我們有很多生活的經驗.
對很多事情的看法都不一樣.
但我們就像一隊首領隊一樣.
我們不可以沒有哪一個.
你試想想剛才那九個弟兄.
沒有其中一個,那個樂章就不完美了.
同樣教會也一樣.
教會的挑戰就是我們一班不同背景.
有不同的社會經驗.

$^{41}$對很多事情的看法都很不一樣.
在不同行業,甚至連信主的經歷也有不同.
但我們走在一起學習.
成為一個優美的樂章去榮耀神.
來到旺角浸信會宣講神的話.
其實每一次踏上台都很緊張.
所以鼓勵你們為教會的牧者去禱告.
支持他們.
可能你說,不是的,你們一班受過神學訓練的人.
應該都沒有問題.
但我想說,在今時今日這個資訊氾濫的日子裡面.
站上台講話是很大壓力的.
所以當教會問,盧先生你那段講道可以放在網上嗎?.
我就想,為什麼會問呢?.
我就明白,今時今日在網上的言論.
都要很小心.
而且現在講求私隱.
在港台講話也要很小心.
跟二,三十年前很不同.
以前可能一講完,就沒有人知道你說了什麼.
但現在你有錄影,有錄播.
你在座有百多二百部手提電話.
現在每個人都有一個手機.
錄音,錄影,錄像都很容易.
訊息就可以瞬間傳千里.
所以在港台的每一位同工都很需要大家去紀念和支持.
讓我們都能夠專心,按著聖經去宣講.
亦都盼望每一次講的內容.
都能夠幫助弟兄姊妹去對應今時今日.
此時此地香港的處境.
講到今天是個很特別的日子.
當然不是說我來到港島很特別.
不知道大家知道今晚有什麼事發生.
第一個土生土長的香港人.
上太空,我想這個都是歷史.
第一個在香港長大.
在香港受教育的人.
離開地球.
我都很肯定應該是第一次.
所以今天是個很特別的日子.

$^{81}$剛才來到教會時我看了大家的周刊.
看到今年的少年部夏令營.
知道這是什麼主題嗎?.
先代大家做宣傳.
AI訓練營.
如果有家長都可以留意.
我想今時今日.
我們經常聽到兩個英文字.
就是A和I.
AI,人工智能.
其實為什麼今天很多人講人工智能呢?.
其實是因為我們的科技很進步.
電腦的發明不是今時今日.
已經是好幾十年的事.
個人電腦應該是幾十年前的發明.
但今天那種突破就是那種算力.
不知道大家聽過一個字叫算力嗎?.
叫computational power.
即是說今天運算的速度.
是嚇死你的.
是說每秒幾兆.
不是幾億,是幾兆.
甚至是更大的運算力.
當它運算力越大的時候.
它就可以瞬間製造很多的影像.
我們很多時候上網除了看文字之外.
是看到很多影像.
以前我們說有圖有真相.
今天就算你兒子或女兒打給你.
你見到他的樣子你都不要完全相信.
「媽媽,我被人抓了」.
「你快點匯多少錢來」.
你真的千萬不要相信.
你說「不是,她的聲音和樣子是她」.
「是阿強仔,是阿小玲」.
但我告訴你,今天很容易做得到.
不是什麼很高科技.
只要它拿了你的相片.
拿了你一段的片段.
拿了你的語音.

$^{121}$很容易就可以複製.
全靠算力的突破.
以致很多以前做不到的事.
今時今日都做得到.
當然AI帶來方便.
但也帶來焦慮.
聽過很多家長說.
很擔心.
以前說做會計師很有前途.
讀翻譯的,專業.
但今天.
我不知道在在沒頂尖做翻譯的專業.
現在用人工智能做翻譯很快.
會計師很有保障.
但我自己也認識一些.
在會計師行業做了很多年的人.
都做到拍檔.
在會計師行做拍檔.
他說現在新入行的很難.
因為你知道很多入門的工種.
已經不用人手去做.
當然你簽名.
所謂負責的都是以老闆做簽名.
但很多中低層的工種.
其實現在很簡單.
很多時候可以用人工智能取代.
很多行政支援的工夫.
很多不同的行業都用人工智能.
現在甚至政府部門.
都用了一些無人機.
用了很多人工智能的東西去幫忙.
所以會擔心將來的子女怎麼辦.
就用電腦去讀.
就像現在的太空人.
以前在香港大學都是讀計算機科學.
電腦科學.
但都不一定.
可能將來不一定是這個行業.
所以我們都不知道將來會怎樣.
我相信社會一定會變得很快.

$^{161}$但唯獨是一樣東西不變.
我很記得.
我以前的老師.
其實都是貴教會的顧問牧師.
梁家良牧師.
他以前都教我們神學.
他說的話我覺得都很有意思.
他說人生很多事都變得很急.
你是不能預知將來.
但唯一一個不變的.
就是我們的主.
就好像我們中學.
應該很多人都學過二次方程式.
X是多少 Y是多少.
叫你計算一下 X是什麼 Y是什麼.
如果整條數都沒有任何的資料給你.
根本計算不到X和Y.
那些我們不知道的東西.
唯有怎樣計算得到.
他一定會給你一個提示.
就是有一個常數.
叫Constant.
Given一個Constant.
就叫你慢慢透過這些已知的資料.
慢慢一步步去計算.
X是什麼 Y是什麼.
其實人生都是.
X和Y變得很快.
現在都不是X和Y了.
現在GEN-C.
GEN-α.
甚至是GEN-β.
快到你同學都不知道.
X和Y是怎樣變的.
但當你整條數學.
有一個叫常數.
叫Constant.
Given了給我們.
就是我們的主.
不變的主的時候.

$^{201}$你就可以去面對一個幻變的世代.
我都盼望當我們.
香港人今天很開心.
我們有人上到太空了.
科技今天很發達.
AI很先進.
可能令到我們都覺得很焦慮.
很怕工作不知道能不能保住.
不知道將來下一代.
怎樣賺錢.
怎樣算的時候.
社會變得很急.
但我們只有一個Given.
給了一個不變的常數.
就是我們的主.
跟它應許的時候.
我們就能夠.
憑信去面對將來.
幻變的社會.
今天我想跟大家看一些經文.
麻煩可否幫我出一下PowerPoint.
不好意思.
要看回自己的樣子.
好.
OK 好 謝謝.
我相信大家在早前的壽苦節.
都應該曾經看過這段經文.
在樣樣form裡面.
記載耶穌基督怎樣一步一步.
踏上十字架的路.
今天我想跟大家重新看一段經文.
在樣樣form十一章.
我不知道大家記不記得發生了一件事.
耶穌叫馬大和瑪利亞的弟弟拉撒路復活.
我相信你們頂尖很多都很熟悉這件事.
聖經裡面說.
馬大和瑪利亞的弟弟拉撒路死了.
但我們看回經文裡面.
耶穌祂叫死人復活.
當然我們如果看到一個生命得到保留.

$^{241}$祂能夠死而復生.
大家都會很興奮.
你看回經文裡面說.
樣樣form十一章四十五節的時候.
祂說.
於是離開瑪利亞的猶太人中.
有很多人看到耶穌所做的事.
就相信了主 相信了耶穌.
在很多人的眼中.
這件事當然是值得興奮.
就算你不是當事人的親屬.
你看到一個生命得到保留.
人最寶貴的是什麼?.
都是一條命.
很老實.
除了謀生.
除了很多其他的事情.
最寶貴都是一條命.
當看到耶穌醫治拉撒路的時候.
很多人都很興奮.
也因為這件事看到耶穌的能力.
而相信了耶穌.
但看到第四十六節裡面這樣說.
但其中也有人去見法利塞人.
把耶穌做這件事告訴他們.
在四十七節你看到經文裡面說.
這些祭司長和法利塞人去開會.
他們說.
這個人.
這個耶穌.
做了這些事.
我們怎麼辦呢?.
可能你會覺得奇怪.
不是一件好事嗎?.
救了一個人的生命.
但這些有權力.
我們這樣說是一班既得利益的人.
他們竟然第一個反應是.
這個耶穌做了這些事.
我們怎麼辦呢?.

$^{281}$在四十八節聖經很清楚地說.
如果這班祭司長和法利塞人繼續說.
如果讓他這樣繼續做.
如果人人都要信他.
可能到時連羅馬人都要來毀滅.
我們的聖殿和我們的民族.
大家都知道在這個時候.
去到耶穌的時代.
猶太人已經沒有自己的國家.
已經在幾百年前.
被亞述,巴比倫分別滅了他們的國家.
接著他們就流散.
跟著很多不同的外族,外邦去管治他們.
亞述,巴比倫,馬代,波斯,希臘.
到如今的羅馬.
猶太人當時不是一個獨立的政治實體.
他們只是在羅馬帝國裡.
有限度自治的一個小區域.
當時管治他們的是一個.
羅馬授權的地方的長官.
叫彼拉多.
這班宗教領袖.
他們知道這件事.
他們不單止會因為一個生命得到保留.
一個年輕人得到復活而興奮.
他們反而就會想.
如果這個人,這個耶穌.
繼續在做這些事.
而越來越多人信了這個耶穌的時候.
那我們怎麼辦呢?.
甚至他們會想.
羅馬人會不會都要來毀滅我們的聖殿.
剷平我們的地方和我們的民族.
可能你會問.
為什麼耶穌治了這個人.
為什麼又關到羅馬的事呢?.
因為他們就會想.
耶穌說自己是王.
很多人看到他跟從了他.
會不會令到中央.

$^{321}$令到羅馬帝國覺得我們這班猶太人.
又想去搞事,又想去反叛.
不知道大家知不知道在.
主前二世紀.
大概主前一百六十幾年的時候.
有一個叫馬加比的革命.
當時猶太人他們很想脫離.
希臘人的統治.
他們起義叛變.
帶來了短暫的一個獨立.
而在主後大概七十年.
大家都知道.
第二聖殿也都被羅馬軍隊剷平.
因為當時猶太人也都想去起義.
脫離羅馬的統治.
所以你會看到整個的情況就是.
這些猶太民族.
他們在羅馬的整個的地域.
整個領土的裡邊.
而羅馬人也都知道這些猶太人.
他們可能心裡一直都不想被羅馬統治.
所以這班的世師長.
這些犯下錯誤的人就會想.
甚至他們想多了.
他們說.
與在上這個耶穌.
自稱是猶太人的王.
他在那裡醫治人.
又講很多道理.
而越來越多人信了他的時候.
那時不只是宗教問題.
會不會有時我們這些宗教領袖.
都失去了我們的權位.
你知道在當時這些的祭師長.
這些的法利賽人.
宗教領袖.
他們所擁有的資源.
不單止是宗教層面.
講個例子你就明白了.
就等於我們很多鄉郊.

$^{361}$因為我以前都是住新界北區.
很多鄉郊 很多的圍村.
很多的鄉事委員會.
他們不單止去處理很多宗族的問題.
或者地方的問題.
宗教問題.
甚至某程度上是管理地區的資源的分配.
同樣當時的這些祭師長和法利賽人.
他們不單止好像今天的教會的.
執事或者是牧者.
他們當時都有一定的權力.
可以去調撥某些的資源.
甚至在猶太民族裡面.
他們都可能有一定程度的管治權.
所以當時就想到.
當耶穌越來越多人認識.
甚至有些人講到.
去信了耶穌的時候.
變相就是越來越少人去聽這些祭師長.
或者法利賽人的話.
他們甚至會想到.
會不會到時羅馬人又會以為我們去到反叛叛變.
而很直接.
就連我們這些的宗教領袖.
或者一些在建制裡面的人.
的既得的利益都會受影響的.
49節其中有一個人.
是當時的大祭師.
叫蓋亞法.
他竟然講出這樣的說話.
在49節裡面.
我們見到聖經說.
「你們甚麼都不知道也不想像.
一個人替百姓死.
免得整個民族滅亡.
這對你們是有利的」.
然後第51節.
聖經自己去作了一個詮釋.
它說.
蓋亞法他這樣講.

$^{401}$原來不是出於他自己的意思.
而是因為他當時是當年的大祭師.
所以預言了耶穌將為這個民族而犧牲.
有時神很奇妙的神的工作.
他透過一個人的口.
似乎預言了耶穌將要為整個民族犧牲.
而52節他說.
耶穌不單止替當時的猶太民族去死.
而是要將上帝四散的兒女聚集起來.
合成一群.
但在53節.
聖經當時就說.
從那天起.
他們.
他們就指那些宗教領袖.
大祭師.
或者是法利賽人.
他們就要相宜.
要殺耶穌.
原來殺機就從那時開始.
可能大家都記得.
到後來他們是接近瘋狂.
當那些官員問.
我放這個好還是放那個好.
那個明明是十惡不赦的人.
他都說.
放了他.
我們要釘死他.
是近乎失了理性.
瘋狂地要釘死耶穌.
但原來整個的殺機.
早在這裡已經開始.
在51節他說.
當時大祭師說的那句話.
這句話不是出於他自己的意思.
當然不是說他突然好像.
突然上身那樣說了一些東西.
而是原來他背後這句話.
似乎就是神透過他的口.
將耶穌將要為人類犧牲這件事.

$^{441}$去預告出來.
當然在當時他們的視野裡.
他們是不知道的.
在當時他們有限的視野裡.
當時他們未有新約.
他們都未知道完全整個救恩的內容.
但當時他們只知道.
因為他們的利益將會受影響.
他們不想耶穌繼續在他們中間出現.
可能原先都是一個宗教的問題.
因為你知道猶太人一直等候彌賽亞.
他們作為宗教領袖.
他們都要負責守護猶太人的宗教傳統.
當後來他們發覺.
耶穌好像不太對路.
好像慢慢和我們心目中期待著的彌賽亞.
有些出入.
但到了後來.
就牽涉到資源分配的問題.
他們會想.
耶穌在我們中間往來.
會不會慢慢人們一個個不聽我們的指摘.
會不會他們一個個就不聽我們.
甚至走去跟耶穌.
另起了一堂.
即是所謂另起了一個山頭.
好像另起了一個不知道什麼的組織.
那我們怎麼辦.
所以他們就起了一個殺機.
你有沒有想過一個宗教的領袖.
原本是要守護著民族的宗教.
甚至在地區裡面作一定的資源的分配.
在整個羅馬帝國統治之下.
他們在猶太這個地區.
去幫助去管理.
幫助猶太的民族.
但竟然這班的領袖.
他們因為個人的利益.
受到挑戰.
他們竟然去到最後.

$^{481}$起了殺機.
來到第十二章.
第九節的時候.
聖經繼續記載.
他說有大群猶太人知道耶穌在那裡.
就來了.
當然當時也說到.
很多人開始跟從耶穌.
很多猶太人開始發覺.
耶穌真是一個很好的人.
他不單只是好人.
他更加有著一個不一樣的能力.
他們當時未完全清楚知道救恩.
未知道耶穌是上帝的兒子.
成就十家救恩等等.
他們未完全知道這個啟示.
但開始發覺耶穌是很好.
很值得跟從.
所以發覺來到十二章第九節.
這班猶太人知道耶穌在那裡.
就去到了.
他們不單止為了耶穌.
更加是去看看.
以之前在第十一章.
說曾經有一個年輕人.
死而復活.
他們很想去看看.
很想去看看拉薩路這件事.
好像等於我們.
有什麼歌手來到.
我們很想去看看他的真面目一樣.
但來到十二章十節.
他說:於是眾祭司長商議.
連拉薩路都要殺了.
你見到嗎?.
剛才第十一章仍然說.
這個耶穌不可以再留在這裡.
不是他死就是我們亡.
但竟然來到第十二章十節.
他說:連拉薩路都要殺了.

$^{521}$你認真想想有多恐怖.
拉薩路做了什麼錯事?.
你說:耶穌都是阻礙別人開飯.
你這麼張揚.
影響了別人.
別人當然對你不客氣.
你這麼張揚.
鼓動群眾.
影響了大祭司.
發揮錯人的權位.
他當然會對你不客氣.
但拉薩路是無辜的.
他只是因為被耶穌施行神蹟.
他重新撿回一條命.
但竟然這班宗教領袖.
說連拉薩路都要殺.
有時你認真讀入聖經時.
你就會看到人性的黑暗.
為什麼連拉薩路都要殺?.
就算把耶穌關在牢裡.
把他踢走殺了他也算了.
因為聖經十二章十一節說.
因為有許多猶太人.
為了拉薩路的緣故.
開始背離他們.
背離這班宗教領袖去信耶穌.
所以這班宗教領袖說.
連拉薩路都要死.
因為這個年輕人.
而走去信了耶穌.
當你投入去看一段經文.
你就會看到人性是多麼恐怖.
這班人不是作奸犯科的大賊.
是當時的一些很尊貴.
有地位的猶太人中間.
往來的宗教領袖.
是被選立出來的一班大祭司.
都是一些有學識,受尊重的一班人.
但想不到他們有這樣的想法.
當然經文來到第十九章.

$^{561}$就是我們很多時候在壽苦節都會看的經文.
來到第十九章第六節的時候.
祭司長和差役見到耶穌.
當時他們已經進入瘋狂狀態的群眾.
他們就說「釘個十字架」.
其實你見到彼拉多是很被動.
如果你能夠細心看那段經文你就知道.
其實當時彼拉多多次的去表達.
「我真的查不到他有什麼罪」.
「你們猶太人的事你們自己管理」.
用現代的術語來說.
「這件事真的不是我跟開的」.
「你們自己民族的事你們民族處理了它」.
但你發覺第十九章第七節.
這些猶太人是怎樣回答的.
當彼拉多說「你們自己處理吧」.
「我真的看不到他有什麼罪」.
雖然說你砌他生豬肉.
也要有羅馬的法律才能夠進去.
「我真的打完他一身之後」.
「我真的沒有什麼可以做」.
當時猶太人在第七節這樣說.
「我們有律法 按那律法」.
「他(耶穌)是要死的」.
「因為他以自己為神的兒子」.
「今天你和我很重視很珍貴耶穌基督的身份」.
「祂是神的兒子」.
竟然當時成為一個入祂罪的原因.
就是因為祂說祂是神的兒子.
所以祂要死.
因為猶太人的觀念說.
「你是一個人來的」.
「你說你是神的兒子」.
那不就是很殲滅了耶和華上主的榮耀.
你怎可以說你是神的兒子.
隨便你說就算了.
所以他們說這個人一定要死.
但為什麼他們又要借助彼拉多的手呢.
因為經文很清楚地說.
因為按照當時羅馬帝國的法律.

$^{601}$猶太人是沒有權執行死刑.
就是說用現在法律的框架.
他們要去將他提升到更高一級的法院.
香港也是一樣 裁判所 地區法院.
高等法院 終審法院等等.
當時也有他們的法律框架.
他們很清楚地說.
因為按照羅馬的法律.
我們猶太人沒有權執行死刑.
言下之意就是要借助彼拉多的手.
但你看到彼拉多一直都很被動.
細心留意經文 你記得.
彼拉多曾經幾次跟猶太人說.
「你們自己釘他十字架」.
「我真的查不到他有什麼罪」.
「算了 你都是憎恨他」.
「我心想 你不就打他一身」.
「踢他走 叫他以後都不要回來搞事就算了」.
「真是現實地說 你好 我好」.
「我又做過地方官 我又安安樂樂」.
「你們就繼續在聖殿那裡」.
但當時猶太人說 他要死.
言下之意就是說.
我們要借助你的手 要他死.
可能經文今天我們看的時候.
你覺得耶穌就是一步步踏上救恩 犧牲.
但回顧背後 原來古往今來.
人性是很黑暗.
人可以因為權力.
因為既得利益而被影響.
而做出很多我們想不到的事.
第十九章十二節 他說.
從此彼拉多要釋放耶穌.
你看到經文有幾次他都說彼拉多.
「我查不到他有什麼罪」.
「可不可以算了 不如你們做回吧」.
「我真的查不到 我打了他一頓」.
「你都出了一條氣」.
但在約翰福音十九章十二節.
這些宗教領袖竟然說.

$^{641}$「大人 彼拉多 如果你釋放這個人」.
「這個耶穌 你就不是凱撒的中神」.
「凡以自己為王」.
這個耶穌說自己是猶太人的王.
他就是背叛凱撒 叛國.
你明白背後的手腕有多恐怖嗎?.
他們想利用彼拉多的權力.
去將耶穌置諸死地.
他們就對彼拉多說.
「你釋放他 你有沒有想過」.
「你釋放了他 你就不是你的老闆」.
「凱撒大帝的忠實地方官」.
「你有沒有想過你烏紗不保」.
「凡以自己為王」.
你知道這裡是誰管的嗎?.
羅馬嘛 權力是來自凱撒.
這個耶穌說自己是王.
你如果釋放了他.
那就是縱容這個叛黨.
到了這個時候彼拉多也沒有辦法.
接著他們就進入了瘋狂的狀態.
我們都記得.
「除掉他 釘他十字架」.
彼拉多某程度上我都覺得他很可憐.
不要說可憐 也很無奈.
「我可不可以將你們的王釘在十字架」.
「我真的可以嗎」.
「我雖然是一個有權力的地方官」.
「但真的沒有一個法律可以說他被判死刑」.
世師長再加一句.
「在這裡 我除了凱撒 我們沒有王的」.
就是說這個耶穌走出來.
說是猶太人的王.
我們是猶太人.
不過我們沒有王的.
除了凱撒之外.
用回今天的說話.
他真的是政治很正確.
到了這個時候.
彼拉多知道他已經沒有辦法再抽離這個局.

$^{681}$他知道自己可能會殺了一個無辜的人.
但他也沒有辦法.
到了最後大家都記得.
彼拉多用了一個牌.
我相信是一個木牌.
寫了一個名號.
大家都記得在十字架上有一個牌.
寫的是「猶太人的王 拿撒勒耶穌」.
我不知道當時是不是罪證.
你也要找一條罪才能安他.
這個人說是王.
那就要死了.
可能他也覺得有點牽強.
但你發覺聖經繼續說.
二十節有很多猶太人唸這個名號.
因為耶穌被釘十字架的地方.
是和城很近.
就像旺角那樣.
很多人往來都會很容易看到.
當時這個牌是用了三個語文來寫的.
希伯來文.
羅馬文.
希臘文.
或者希尼尼語.
翻譯就是.
等於今天在港鐵.
聽廣播.
或者在酒吧聽廣播是怎樣的.
有三種語文.
廣東話.
普通話.
英語.
為什麼要說三次.
因為他要讓當時.
這三種語言的人都知道這個訊息.
當時.
希伯來文很重要.
因為這個也叫做猶太人的地方.
可能很多人的母語是希伯來文.
為什麼用羅馬文.

$^{721}$很簡單.
這裡是羅馬的地方.
當然要用官方語言.
希臘語.
或者叫希尼尼語.
因為當時如果大家記得.
在羅馬統治這個地方之前.
是希臘人管治的.
所以很多人可能都會懂希臘語.
所以他就用了三種語言去寫出這句話.
猶太人的王.
拿撒勒賢耶穌.
但有時候你留意經文.
很有趣.
二十一節.
這些猶太人的祭司長看到就很緊張.
他說.
彼拉多.
你不要寫猶太人的王.
你說.
是他自己說.
我是猶太人的王.
你有沒有留意到分別.
就是說我們沒有說他是王.
你弄清楚.
我們只是告他.
他自己說他是猶太人的王.
我們沒有王的.
除了凱撒之外.
但見到這個時候.
彼拉多也有一個很有趣的回應.
二十二節.
彼拉多說.
我寫就寫了.
沒有補充.
我可以推斷當時彼拉多可能都覺得.
現在誰做官了.
你叫我A就A.
叫我B就B.
我已經寫了.

$^{761}$或者這樣說.
我不會改了.
你們自己再自己斟酌.
你可以看到整個事件.
除了我們看到耶穌一步一步登上十字架.
犧牲之外.
看到中間很多不同的權力.
很多的人性的黑暗.
在人類歷史裡面.
我們都會不斷為了資源而爭戰.
今天你看到國與國之間很多的戰爭.
其實說到尾都是為了爭奪資源.
當然資源牽涉到錢.
牽涉到權力.
其實古往今來.
人類之間都會因為一些利益.
權力而彼此去到.
可以這樣說.
人與人之間的打仗.
我們看到原來在耶穌時代.
同樣都是面對著一個這樣的情況.
在這幅圖.
當然可能這些字比較小.
我嘗試去整合了整個《約翰福音》.
這套記載故事.
表面上看好像是一場宗教的審判.
因為耶穌說了一個不合當時猶太教.
舊約以色列人的一個宗教上的事.
他竟然說自己是神的兒子.
僭越了耶和華上主的榮耀.
但發覺背後其實是一個很人性的一件事.
他們為了保存自己的權力和地位.
他們勾結了羅馬一些地方的權力.
最後動了殺機.
要將耶穌置諸四地.
當然當時彼拉多也是沒有辦法.
他知道現在已經騎虎難下了.
當你們猶太民族的領袖說到叛亂.
又說到凱撒.
你知道他都已經沒有選擇.

$^{801}$最後耶穌基督就是犧牲在他們手上.
當然我這樣說不是要大家今天很憤怒猶太人.
我知道在歐洲曾經在二次大戰的時候.
猶太人是受到很大的災禍.
我不知道你們知不知道在二次大戰的時候.
猶太民族在整個歐洲經歷很大的禍患.
集中營很多上百萬的猶太人被送入毒氣室.
其中一個為什麼能夠說服到當時教會去所謂默認這件事呢.
如果你認識二次大戰的歷史.
他們說猶太人就是殺死你們宗教的耶穌那個人.
猶太人是不是該死呢.
希望今天不斷說一些宗教領袖.
猶太教的領袖.
舊約的祭司長不好.
今天我們對現在的猶太民族或者對整個猶太人有種敵意.
其實這個只是反映了人性.
不是說只是宗教領袖,祭司長是多麼大奸大惡.
其實如果換成我和你是當時.
有可能我們都會有這樣的反應.
當人擁有越多的時候.
他就不捨得去失去.
當人擁有的時候他就很想去保護他擁有的東西.
當人有一定的權位的時候.
他就會做出一些決定.
去到為了自己或者為了身邊的人的利益.
今天頂尖可能你在工作的裡面.
你在職場的裡面.
其實在任何的圈子,朋友裡面.
都會有這樣的情況.
今天你和我能否守住自己的初心.
能否去時時保守自己的心.
這個都很重要.
所以我強調今天的訊息不是說我們要仇恨這些猶太人.
那些猶太人是活該的.
我們都是因為猶太人才犧牲的.
當然也不要另一個極端.
他們都是舊恩的參與者.
好像別人踢球一樣.
你也是為了球給他才入球.
沒有猶太人,沒有宗教領袖,耶穌也不能上十字架.

$^{841}$這也是另一個極端的錯誤的想法.
我們知道宗賢上主讓他的兒子犧牲在十字架上.
但在人間其實看到當時充滿了很多人性和罪惡.
在這裡多說一點.
你有沒有幻想一下當時一群跟從耶穌的人.
當耶穌一步一步地犧牲.
在十字架上寫信的時候.
他們有什麼感受?.
當然今天.
例如我們在指教會上會有受苦節崇拜.
也有復活節崇拜.
通常兩個崇拜是很接近的.
受苦節崇拜我們習慣在星期五.
復活節崇拜則在星期日.
其實中間只是隔了一天.
我們今天沒有什麼特別的感覺.
還好,有幾天假期.
幾天假期,就是受苦節,復活節.
中間那天我們不是有什麼太大的體會.
但我不知道你有沒有聽過一個叫做.
Holy Saturday.
聖周六.
你幻想一下當時.
這群跟從耶穌的人.
看到他們一直跟從的耶穌.
當然他們不會知道之後如何有基督教.
慢慢基督教就發展.
接著就有耶穌復活那些.
但你嘗試.
好像坐時光機.
就好像我們那時看《尋秦記》那套電影.
你幻想一下你和我坐著時光機.
去到當時《約翰福音》十九章這個場景的時候.
你看到耶穌.
最後釘在十字架上犧牲的時候.
一切都沒有了.
沒有了.
你嘗試暫時忘記了.
你認知那個復活事件之後.
你當時看到是一片黑暗.

$^{881}$霧門關上.
舊大石.
把墓室都封住了.
接著那些人說.
快點回家吧.
安息日就到了.
之前耶穌在地上很多的教導.
對耶穌很多的希望.
拉撒路復活那位救主.
好像帶來一種希望.
完全都沒有了.
正所謂人死不能復生.
那祂真的死了.
你可以幻想那種.
很不明白.
很沒有希望.
很無力的那種感覺.
所以有人曾經說過.
我不知道大家會不會經常背《仕途信經》.
在我自己的宗派教會.
我們每個崇拜都會背《仕途信經》.
我相信頂尖的你都應該認識《仕途信經》.
就算教會裡未必每個星期去背.
它說.
受苦埋葬.
降在陰間.
第三天從死裡復活.
升天.
坐在全能負上帝的右邊.
將來要從那裡降臨.
審判活人死人.
我每個星期都會背.
所以我很記得.
你有沒有留意.
當《仕途信經》對耶穌基督的記載時.
它說.
再被拉多手下遇難.
被釘在十字架上.
受死埋葬.
降在陰間.

$^{921}$就是這個位置.
聖周六就是降在陰間的位置.
可能只是時間短短的.
可能安息日.
中間.
可能都是二十個小時.
三十個小時的時間.
他們是完全沒有希望的.
但正正就是這段.
好像在人看來沒有希望.
很黑暗.
很無力的日子.
耶穌就去到陰間.
當然我們知道.
聖經對陰間不是太多的記載.
我們今天也不可能對陰間有太多的詮釋.
我們不可能有太多的.
過分的詮釋去幻想.
但《仕途信經》很清楚地說.
受死埋葬.
降在陰間.
我們就知道.
哪怕去到陰間那裡.
耶穌基督也在那裡.
在人類最黑暗的時間.
在大地無光的日子.
耶穌基督仍然在那裡.
所以今天可能你在生活的裡邊.
在你工作.
在你個人.
或者在各方面.
你經歷那種.
好像看似沒有希望.
好像你覺得很困難.
很無力的時候.
都盼望你仍記得.
在那裡.
耶穌也在.
有人說過一個比喻.
如果你覺得人生時很困難.

$^{961}$面對的困難.
好像一個看不盡頭的隧道.
那耶穌在哪裡?.
你猜耶穌在哪裡?.
那你就在隧道的盡頭.
離遠一點光那裡.
就是祂了.
忍著黑暗裡面前行.
不是啊.
耶穌就在你旁邊.
祂不是在隧道口的洞裡等你.
加油啊.
多走一段日子就到了.
耶穌在你旁邊.
我不知道你記不記得.
在我小時候有一個叫Footprint.
竹印的那個比喻.
沙灘裡面為何只有一對竹印.
我記得我八十年代上教會.
很感動的例子就是.
為何只有一對竹印.
因為祂背著你.
另外一個版本就是說.
因為祂走一步.
你跟著祂一步.
你就踩在祂的腳上.
所以回望祂只有一對竹印.
因為你是跟著祂的腳步.
祂踏在那裡.
安全.
你踩下去.
安全.
所以結果只有一對竹印.
同樣我告訴你.
在隧道裡面.
你說.
哎 我很苦啊 耶穌.
我現在身體又不好.
經濟又困難.
社會又亂.

$^{1001}$很多很多的事情覺得很艱難.
原來在那條.
好像看不到出口的隧道.
離遠那點光.
耶穌不是在那點光裡等待.
耶穌在我們的旁邊.
在我們的旁邊.
甚至你可以想.
是祂背著你走.
那個場景不是沙灘.
是一個你覺得很黑暗.
無盡的隧道.
原來在受難和復活之間.
存在著一個容易.
我們被遺忘的空間.
在《仕途信經》裡面就這樣說.
它說.
受死埋葬降在陰間.
第三天從死裡復活.
升天.
坐在全能附上帝的右邊.
將來要從那裡降臨.
審判活人死人.
《仕途信經》是歷代教會.
一個很珍貴的禮物.
它將很多重要的真理.
將它壓縮成一個.
我們叫信經.
讓我們能夠.
將整個基督教核心去記得.
所以我們今天同樣都要記住.
縱然有時我們個別的困難.
好像很艱難.
可能你說.
不是的.
以前那些門徒兩天就解決了.
三天而已.
理智就是短.
但當時在他們來說.
他們真的不知道.

$^{1041}$所以他們的心一樣很焦慮.
很無力.
很無助.
同樣今天我們不知道.
是三天還是三年.
還是像梁朝偉那樣說.
三年又三年.
我們不知道經濟何時會好轉.
我們不知道很多問題何時會解決.
但我們知道.
在這一刻.
耶穌仍然都一齊同在.
在我們中間.
縱然有時客觀環境.
令我們感到困難,疲倦.
但是我們知道.
正如昔日的門徒一樣.
縱然他們見到耶穌.
好像已經離開了.
好像已經沒有辦法了.
那個大石都已經封了.
那還有什麼好說呢?.
原來黎明將近.
仍然我們彭訊仍然有希望.
我剛才都說到.
XYZ我真的計算不來.
錢又計算不來.
健康又計算不來.
很多東西都計算不來.
什麼都計算不來.
AI又來了.
我真的搞不定了.
很焦慮.
你只要記得.
耶穌在這裡.
我們仍然有希望.
盼望我們中間頂智媒.
可能對很多事都有不同的看法.
就像我說的守靈隊.
可能都有不同的性情.

$^{1081}$不同的高矮.
但是走在一起的時候.
我們就成為一個很美麗的樂章.
我們在中間唯一的最大的公因素.
就是我們都相信主耶穌基督.
我們是同一個信仰,同一個主.
一信一洗.
既然是這樣的時候.
我們就持著這個盼望.
去走每一天.
並且當我們身邊有.
我們認識的人.
可能他都覺得很無力.
很累.
覺得很多無奈的時候.
你都嘗試提上一杯涼水.
最後我想分享一下.
我本身教會在大埔.
我想你大埔也都知道.
在幾個月前發生了一件很恐怖的事.
其實這件事繼續在發生.
而我們教會中間有幾個家庭.
都是黃福苑的災民.
當然我不知道應該說感恩還是不感恩.
他們沒有事.
他們是.
我不要說幸運.
如果這樣說.
好像其他人就不幸運.
所以我都很小心不要這樣說.
他們是.
剛剛那段時間是出去了.
那就是沒有事.
但不可以說幸運.
當然不可以這樣說.
所以我們教會對這件事.
是很貼身的感受.
因為這件事不是香港的事.
是我的社區.
等於旺角浸信會.

$^{1121}$如果你在旺角做街坊.
博鞋街,西洋菜街有事.
幾乎就是你的事一樣.
所以這件事對於我.
或者對我的教會是一個很大的震撼.
很大的震撼.
所以我們會特別關注.
當然中間有很多東西可以去判斷.
很多東西可以去說.
但我覺得這段時間是很多眼淚.
很多哀傷.
在香港人裡面有很多的眼淚.
很多的傷害.
盼望在這些這麼困難的日子.
好像一條隧道.
我都說不到是何時才完.
單說黃福苑這件民生的事.
都不知道要搞多久.
都不知道要搞多久才完.
那些災民的傷痛.
其實你真的沒有辦法去幻想.
你很難去想得到那種傷痛.
但是我們教會有牧者去陪著他們.
我們教會是有兩位牧者.
他們是私底下.
不是教會是他們私底下.
因為他們都是大埔街坊.
他們都是去陪著這些人.
可能只是遞包紙巾.
其實沒什麼做得到.
因為作為教會.
我又沒得陪給你.
我們也沒有這樣的權力去做什麼的仲裁.
但就是那一刻.
我們遞上一杯涼水.
希望我們一起走過這段不容易的路.
我們一起去祈禱.
各位頂智媒體.
如果你同樣面對很多的困難.
或者你是經濟的困難.

$^{1161}$或者你是其他的困難.
可能是個人的身體健康.
或者是很多不同的原因.
令到你覺得很無力.
很不容易的時候.
你都可以去祈禱.
求神幫助我們.
是一個看似困難.
甚至覺得沒辦法了.
沒話說了.
覺得很不公平.
或者覺得痛到你不懂得去形容的時候.
你都可以去祈禱.
求主幫助我們.
求神幫助我們每一個.
不管是弘福縣的災民.
是我們頂智媒.
是旺角針訊會的每一個頂智媒.
你都去祈禱.
求耶穌幫助我們.
讓我們能夠繼續有力.
走過.
走得完這條隧道.
因為相信隧道雖然黑暗.
好像很長.
但是一步一步走.
你總是會走到出口的一天.
我們一起去祈禱.
神求你幫助.
當我們面對著世界的紛亂.
面對AI時代的來臨.
我們都有很多的焦慮.
或者我們都有很多的不知道.
我們每天打開新聞.
看到的新聞.
十單有八單都是令人不開心的事.
不是天災就人禍.
不是各與各之間的爭鬥.
就是很多無奈的事情.
神求你幫助我們.

$^{1201}$讓我們能夠仍然心存盼望.
縱然客觀的環境好像給我們一個很無希望.
就好像昔日的門徒.
當見到連主你自己都釘死了.
好像黑暗掌權.
門徒要四散.
怕被牽連的時候.
仍然心存盼望.
我們不知道黎明何時來到.
亦都不知道地上的紛亂何日終止.
但是讓我們仍然心存盼望.
以致我們仍然在地上.
努力活出一點光.
為身邊的人.
如有力的話.
遞上一杯的涼水.
或者只有半杯就一人一口.
跟他們分享.
成為他們彼此的鼓勵.
多謝你聽我們祈禱.
奉耶穌的名.
阿們.
\newpage



\section{創世記 }
\label{sec:dKG666yV4KA}
\textbf{2026年5月31日崇拜直播｜楊珊珊傳道｜告別浮淺,門徒進深:衝破過去的藩籬｜創世記 四十一 46-52,四十五 1-8}
\newline
\newline
連結: \href{https://youtube.com/watch?v=dKG666yV4KA}{\texttt{https://youtube.com/watch?v=dKG666yV4KA}} ~~~~ 語音日期: 2026-05-31
\newline
\newline
\hyperref[sec:0dlSUYhlo5s]{\small{< < < PREV SERMON < < <}}
~
\hyperref[sec:index_chronic]{\small{[返順時目]}}
~
\hyperref[sec:NSYx8y4kLPg]{\small{> > > NEXT SERMON > > >}}
\newline
\newline
$^{1}$提交文件:創世記第41章46-52節,第45章1-8節.
提交文件:約瑟見埃及王法老的時候,年三十歲,他從法老面前出去,遍行埃及全地七個豐年之內,地的出產極豐極盛..
約瑟重唸埃及地七個豐年一切的糧食,把糧食積存在各城裡,各城周圍田地的糧食都積存在本城裡..
約瑟積蓄五谷甚多,如同海邊的沙,無法計算,因為谷不可聲數..
方年未到以前,安城的祭司波提菲拉的女兒阿西納給約瑟生了兩個兒子,約瑟給長子喜名叫嗎哪西..
因為他說:「神使我忘了一切的困苦,和我父的全家.」他給次子喜名叫爾法蓮,因為他說:「神使我在受苦的地方昌盛.」.
第二段經文在四十五章一至八節,約瑟在左右站著的人面前情不自禁,吩咐一聲說:「人都要離我出去.」.
約瑟和弟兄們相認的時候,並沒有一人站在他面前,他就放聲大哭,埃及人和法老家中的人都聽見了,約瑟對他弟兄們說:「我是約瑟,我的父親還在嗎?」.
他弟兄不能回答,因為在他面前刀驚惶,約瑟又對他弟兄們說:「請你們近前來.」他們就近前來,他說:「我是你們的兄弟約瑟,就是你們所邁到埃及的,現在不要因為把我邁到這裡,自憂自恨.」.
約瑟說:「這是神差我在你們以先來,為要保存生命.現在這地的饑荒已經二年了,還有五年,不能耕種,不能收成.神差我在你們以先來,為要給你們存留餘種在世上,又要大師拯救,保存你們的生命.」.
約瑟說:「這樣看來,差我到這裡來的,不是你們,乃是神.祂又使我如法老的父,作祂全家的主,並埃及全地的宰相.」.
弟兄姊妹平安!在跟隨耶穌的路上,我們經常聽到一句:「舊事已過,都變成新的了.」.
但在很多信仰的路上,弟兄姊妹,特別是那些在原生家庭受過創傷,又或者成全了家族幾代的拉扯的弟兄姊妹,可能心裡經常都會有一個痛..
痛就是:如果我已經是一個新造的人,為何過去的陰霾和痛苦仍然拉扯著我們呢?.
究竟怎樣才算是真正的德性呢?約瑟的生命正是給我們每一位門徒一個赤裸而又真實的答案..
約瑟的家庭組合可能跟我們這些普通人有點不同,但相同的都是家家有本蘭殿的經.他的父親雅各昔日是澳門賭王,有四房的妻子..
雅各最愛的二房,拉傑,就是約瑟的親生母親.當年他母親生了他第二個親生弟弟,變亞敏的時候,拉傑因為難產而去世..
約瑟在年幼的時候,喪失了母親.一個小朋友,瞬間失去了最重要的母愛和保護網..
更慘的是,上一代四房的爭寵,明爭暗鬥的政治糾紛,直接轉移到下一代12個兒子身上..
約瑟在他的家排行第十一,因為他是他父親和他最愛的女人的潔青品,所以雅各極度偏愛他,甚至做了一件獨一無二彩色的大外套送給他..
其他哥哥是沒有份的.你可以想像到,約瑟從小到大都在家裡承受著很多其他來自三房的哥哥的冷嘲熱諷和排擠..
在17歲那一年,他聽了父親的吩咐,去探望他的哥哥.約瑟這個年輕人懷著善意去找他的哥哥,誰知等待約瑟的竟然是一群親生哥哥,因為妒忌和怨恨,於是出賣和拋棄他..
十個哥哥當著他的面前,強行奪去了他父親給他那件很漂亮的七彩大衣,他們不留情面將約瑟推到一個曠野的深坑裡..
約瑟在坑裡歇斯底里哭,哀求,但上面最親的骨肉卻裝作聽不到,看不見,還要很冷血地在坑邊吃飯,甚至在商量怎樣去殺死這個弟弟..
最後這群哥哥以20個銀錢將他當成商品賣給外族,在那一刻,17歲的約瑟徹底被自己原生家庭遺棄..
他瞬間由一個德寵的少爺仔,變成販賣到埃及異鄉的奴隸,他人生的安全感,主導權,身份,甚至連言語,在一夜之間完全清零..
約瑟輾轉被賣到法老護衛長波提弗的家,他很努力去生存,他是一個很謙虛,很威神的人,他很忠心去工作,他的老闆甚至將他家裡所有東西交給他打點..
誰知他的惡夢未完,因為女主人波太太看中了他的美色,日夜對他進行權力不對等的色誘和騷擾..
當波太太強行捉住約瑟的外衣時,約瑟內心一定極其驚恐,那一刻可能觸發了他在17歲時的創傷記憶和本能反應..
可能他腦海裡也在想昔日他的哥哥們如何瘋狂地去搶他外衣的恐怖畫面,你們要那件衣服嗎?我給你,求你放過我.於是他在那一刻就把外衣丟下給波太太,很赤裸裸,很驚惶地走了出去..
有些學者解釋說這件大外衣代表著約瑟的身份和權力,他敬畏神,不可以背叛他的主人..
所以他丟下這件外衣代表約瑟主動放下自己的身份和權力,他再不是那個不知天高地厚的少爺仔..
約瑟堅守了對神對人的道德,不過現實換來的卻是惡意的插贓嫁禍..
他在毫無自辯機會的情況下含冤入獄,由一個受害者變成一個階下囚,被關在一個冷冰冰皇宮的監牢裡..
他在監牢裡等了一年又一年,他很深為法老的酒精解夢,奢望高官可以幫他出獄後伸冤..
但酒精的官職恢復後,就把約瑟忘記一乾二淨,於是習慣了等待的約瑟,在監牢裡又再苦等了兩年的時間..
直到有一天,法老發了一個夢,在埃及裡沒有任何術士可以解釋這個夢,酒精才突然想起有一個少年人叫約瑟可以解夢..
這次約瑟自己梳洗乾淨,換好衣服,站在法老面前,神透過他幫法老解夢,更賜給約瑟有治國智慧,成功幫助法老化解國家糧食饑荒危機..
轉眼間,就來到今天的經脈,創世記四十一章..
如今三十歲的約瑟,經歷了整整十三年血淚交織,老弟和囚犯黑暗的日子..

$^{41}$他搖身一變,變成了埃及的宰相,一人之下,萬人之上..
這個時候,約瑟掌管著大權,但他仍然很小心翼翼,決定不再走父親雅各的舊路..
即是他父親娶了很多老婆,他娶了一個老婆,他父親偏愛他的兒子,他兩個兒子都很疼愛..
他為兩個兒子改了兩個名字,這兩個兒子的名字正正就是記錄了約瑟如何將這十三年來,.
他的原生家庭和他的際遇帶給他的創傷,交回給神,並且容讓神去翻轉,去重寫他的生命的劇本..
這兩個兒子的名字也是我們作為耶穌基督的門徒去轉化生命,打破過去的一些詛咒的兩大靈性的里程碑..
他第一個兒子的名字叫嗎哪西,這是他第一個里程碑..
剛才我們讀到《創世記》的41章46至51節,大兒子出世的時候,約瑟三十多歲,.
剛剛在埃及執掌大權,正正進入他人生最輝煌,薄煞的年頭,他幫他大兒子改名為嗎哪西..
這個名字很特別,它不是一個埃及的名字,它是希伯來文的名字,意思就是使之忘記,.
即是神讓他去忘記一切的困苦,以及他父的全家..
在門徒訓練的視覺裡,這裡所說的忘記不是指失憶,更加不是那些阿Q精神,要自我麻醉..
因為我相信約瑟很難去忘記自己被親生哥哥扔到坑裡,.
在監倉裡含冤受屈十多年的經歷和歷史.約瑟的忘記到底是指甚麼呢?.
其實他的忘記是指他的人生,他願意忘記由傷害他的人去定義他自己..
這是指生命主權的交托,將生命的主權完完全全surrender給神去掌管..
約瑟的決定是他不再容許受害者的心態去綁架他的現在和未來..
他每當過去的委屈,自親的出賣,那些因為他人失職而對他不公平的對待,.
有時候真的會像海浪一樣,一而再,再而三地湧到他的心頭上..
不過約瑟是靠著神,他可以見到自己可以做出一個截然不同的選擇..
他可以選擇不再孤零零一個無能為力的受害者,他可以選擇去依靠神..
他可以將這個痛交給上帝,他對神說:.
「我的人生是由上帝你來定義,而不是由傷害我的人去定義.」.
這個就是因為約瑟看見,其實他自己是可以有選擇的..
這一份忘記就是門徒對生命主權的放手..
想像一下如果有一個露宿者,他在街頭流浪很多年,.
他身邊總是有一袋二袋在街頭撿回來的舊衣和雜物,.
這些就是他的全部家當..
後來有社工幫他,終於排到可以上公屋,可以上樓..
而教會的弟兄姊妹都很熱心,幫他找到一些傢俬,家電,.
打理得妥妥當當,準備好一個很乾淨,有瓦遮頭的家給他入住..
但這個流浪漢入住了這間新屋之後,他的生活竟然和他以前在街頭流浪是一模一樣..
他依舊不洗澡,每天都繼續在街頭撿那些一袋二袋的垃圾回家..
轉眼間,原來一個很乾淨,很舒服的新屋,.
被他堆到很骯髒,很臭,變成了另外一個垃圾缸..
弟兄姊妹,我想問這個流浪漢怎樣才可以真正享受到他新屋的生活?.
唯一的方法就是他要作出一個選擇,.
就是他要去發現,去察覺,並且承認他自己已經不是睡街了,.
他有一個新的地方,他的環境不同了..
他要選擇脫離過往流浪的生活,.

$^{81}$選擇去洗澡,選擇親手將以前在街頭當作寶物,.
但其實今天已經沒有用舊家當逐樣去斷捨離..
只有他願意作出這個選擇,他的新屋才不會繼續發臭,.
他才可以舒舒服服地住在這個新的家,有一個新的開始..
弟兄姊妹,藥室的生命故事,或者是我們的生命,.
有時候都好像這個流浪漢一樣,.
可能過去有一些傷害你的人,其實他已經不再在我們身邊,.
又或者他已經不可以做什麼去傷害我們,.
可能有一些欺壓,又或者一些不公義的事情,.
已經發生了很多年,上帝其實已經帶領我們去到一個新的階段,.
但問題是,你的肉身走了到新屋裡,.
你的心靈還是否選擇流浪在街上呢?.
如果我們一邊很想跟隨主耶穌過一些新的生活,.
一邊我們的頸有介紹就轉向後面,越轉越後,.
永遠執著在過往誰對我們不住,.
誰出賣過我們,我們只會原地踏步,.
不斷地被拉扯,甚至會讓自己去跌倒,.
一個眼睛望向後面的人,是沒有辦法好好走前面的路的..
弟兄姊妹,今天我們不需要再被過往的創傷困死,.
我們可以選擇不再做一個活在過去的受害者,.
我們可以選擇轉身去倚靠上帝,向前邁進..
在我們每個人的心裡,可能都會有一本寫滿了怨恨的小氣堡,.
但作為耶穌基督的門徒,我們可以靠著主的恩典,.
將心裡積怨很多年的傷害,不公義,冤屈,.
甚至是每一道未康復的傷痕,逐樣逐樣拿出來,.
選擇放下給上帝,讓祂為我們作主,.
醫治甚至更新..
今天,你願不願意靠主,換一個比回憶更廣闊的視角呢?.
第二個里程碑關乎他的第二個兒子,就是以法連..
不久後,約瑟又生了第二個兒子,.
他替他的小兒子改名為以法連,.
意思是使之昌盛,讓他繁衍..
約瑟為這個名字做了一個很深刻的註腳,.
在經文裡也印了,「神使我在受苦的地方昌盛.」.
如果我們代入約瑟的處境,用平常人的眼光去看,.
埃及這個地方到底代表什麼呢?.
埃及是約瑟人生所有的黑歷史,所有惡夢的源頭,.
他在那裡被當成貨物販賣,.
那裡也面對了一些無恥的性騷擾,.
在那裡甚至被含冤去誣告,更加坐了十多年的冤獄..

$^{121}$如果按照我們人的本能和報復的心理,.
一個人一旦得勢力,他就掌大權,.
最想做的事情通常會有兩樣,.
第一,馬上離開傷心地,永遠都再不回來;.
第二,他用盡所有手頭上的權力去向當日整蠱他,.
踐踏過他的人,瘋狂地報復,.
讓對方也試試自己受苦的滋味..
不過,我們作為跟隨主的門徒,.
約瑟讓我們看見一個完全相反的選擇,.
當日叫他最痛,流得最多眼淚的那個受苦之地,.
今日竟然成為了他昌盛的起點..
上帝是他沒有抹走,也沒有拿走約瑟做過奴隸,.
坐過監的那段黑歷史..
如果沒有當初在波提弗家做管家的磨練,.
約瑟根本未必學會如何管理一個龐大機構,.
這裡說的是法老的政府,一個政權,一個國家..
如果他沒有在皇宮的監倉和高官囚犯經年累月的相處,.
他根本沒有可能培養出他的政治智慧..
人以為被糟蹋,被浪費的年日,.
在上帝的手中,竟然全部都翻轉過來,.
變成日後管理埃及,化解全國的饑荒,.
甚至拯救無數生命最紮實的智慧和底氣..
約瑟站在埃及這個受苦之地,.
他看見自己原來有第三個選擇,.
他不但可以選擇忘記放下過去的怨恨,.
馬來西亞這個名字,.
他還可以選擇留在那個曾經傷害過他的地方,.
繼續忠心,繼續去愛,繼續去作神的器皿..
而整個故事最啟發我的是約瑟那幾次拒絕扮演受害者的心態和行動..
約瑟第一,他面對引誘,他堅決說不,.
他面對波太妻子權力的壓迫和引誘,.
他打破沉默,說他不要,這是需要極大的勇氣..
第二,約瑟拒絕自憐,他以一個很頹廢的狀態去示人,.
當法老在監倉裡叫他,召喚他的時候,.
他沒有披頭散髮去博同情,.
而是他剃頭,剃鬚,更換乾淨的衣服,.
他拒絕扮演一個很自憐的受害者..
第三,約瑟無論在環境悲戰或是富足,.
他在護衛長的家,他在監倉,他在皇宮,.
他都選擇去服侍,去成長..

$^{161}$約瑟他在監倉十幾年,其實是他生命中最重要的蛻變期,.
創傷奪去了他的天真,.
但奪不走他的善良和對神的忠心,.
苦難抹去了他的傲氣,.
但讓他穿上謙卑..
在監倉裡,監獄長他很相信約瑟,.
他將囚犯交給約瑟去管理,.
約瑟他在那裡聽了不同人苦難的故事,.
他發現可能世上真的有很多冤獄,.
他的傷痕在服侍其他人裡,.
他漸漸都被轉化和療癒..
聖經記載,當約瑟見到酒精和善長面帶受用的時候,.
約瑟主動關心他們兩個,.
「為什麼你們今天戴著面帶受用?」.
他不再沉溺在自身的痛苦,.
而是成為一個能夠分擔別人痛苦的聆聽者..
第四,約瑟用愛切斷負面連鎖反應..
當饑荒來到的時候,.
約瑟和他當年傷害過的哥哥重逢,.
在這個時候,約瑟已經手握大權,.
只要他揮一揮手,.
就可以將他以前受過的痛苦,.
雙倍奉還給這班哥哥..
如果約瑟選擇報復,.
他只是在重演他爺爺,爸爸那種世世代代傳下來的真經和報復..
但約瑟選擇了一個完全相反的路,.
這是門徒生命最關鍵的一步,.
就是他打破了幾代相傳的詛咒,.
那些傷痛不代存..
弟兄姊妹,我們每一個人都有我們自己成長的歷史,.
最深影響我們的性格,.
往往都是小時候我們在我們家人,父母身上看到,學到的東西..
如果父母以前經常用打打鬧鬧,.
用一些很難聽,有骨的說話去管教我們,.
其實我們長大後很自然,.
我們都會用這一套,.
用一些有刺,有骨的說話去對待我們的下一代..
聖經裡面,大衛王就是一個歷史的教訓,.
大衛自己犯了錯,.
沒有好好處理,.

$^{201}$結果他的兒子就學他的壞事,.
家裡出現了兄弟相殘,甚至背叛,.
痛苦一代傳一代..
家原本是神所設立,.
是讓人得到溫暖,安全,成長的地方,.
但在這個破碎的世界裡面,.
家庭往往就成為了傷害人最深的地方..
更加危險的是,.
如果我們不小心,.
這條傷痛代傳的枷鎖,這條鎖鏈,.
不單止捆綁了我們自己的家庭,.
還會帶入這個屬靈的家庭,即是教會..
如果我們帶著以往受過的傷害,.
又不去處理,.
結果我們就無意間將我們在原生家庭.
學到那一套的操控慾,冷漠,.
一些惡性的競爭,比較,.
甚至是喜歡批評人的態度,.
傳染給教會裡面的弟兄姊妹,年輕人,.
或者是我們帶領著的門徒..
我們試想像一下,.
如果我們手上有一個黑漆漆的水,.
滴了一些墨進去,.
可以下一頁..
我們倒進去另外一杯的水裡面,.
第二杯的清水都會變黑,.
再倒進去第三杯,第四杯都會變黑..
可能到了我們這一代,.
我們很努力加很多清水,.
令到杯水變成淺灰色,.
可能比起以前我們在父母身上所承受的,.
已經淡了很多,.
溫和了很多..
不過問題是,.
杯水仍然是骯髒,.
裡面的墨水仍然是有毒,.
是不能喝的..
只要我們一天繼續倒以前那杯水,.
我們的下一代,.
教會的弟兄姊妹那杯水,.

$^{241}$都是永遠無可能清澈的..
只要我們做什麼呢?.
就是我們阻止污染,.
唯一的辦法就是我們停手不再倒,.
將整杯本身的那杯污水倒掉,.
徹底換一個新的水源,.
我們才可以有一杯乾淨的水,.
再給其他人去喝..
弟兄姊妹,.
藥室就是那個決定停手不倒,.
更換水源的人..
藥室曾經受過很多傷害,.
但是她沒有將她的怨恨,.
這些污水倒給她的下一代,.
也沒有倒給她的哥哥,.
她將她內心那杯污水倒掉,.
她知道她生命的泉源,.
不是那杯污水,.
而是耶和華真神..
聖經給了我們很大的影響,.
在基督裡,我們都是新造的人,.
在主裡面,真理才會叫人得自由,.
我們不需要再去無能為力地.
複製過往的痛苦給其他人,.
藥室切斷了仇恨的鎖鏈,.
她成為了一個帶著傷痕的醫治者,.
一個wounded healer,.
她被上帝治好的傷痕,.
沒有變成污水,污染人的墨水,.
反而化作了一個清泉,.
她醫治了她整個家族的未來,.
所以當她們兄弟相認的時候,.
藥室沒有倒出那些怨恨的污水,.
她反而是流下眼淚,.
藥室對那群戰戰兢兢的哥哥,.
說了一句很溫柔的話,.
在《創世記》45章5節,.
這個直譯是這樣說的,.
現在哥哥你不要將賣我到這裡,.
你自怨自艾,自暴自棄,.

$^{281}$她改寫了她自己人生的劇本,.
如果我們在藥室煉腦的時候,.
請她寫下自己的回憶錄,.
我猜她的生命故事大概是這樣的,.
她這樣說,.
我藥室17歲來到埃及圍爐,.
無辜地坐了十幾年牢,.
欣著神的同在,.
我為酒精和法老解夢,.
30歲那年走出去監倉,.
成為埃及的宰相,.
法老賜給我一個新的名字,.
叫做薩法那特巴萊亞,.
意思是時代的拯救者,.
我娶了祭司的女兒阿西勒,.
我在三十幾歲的時候生了兩個兒子,.
大兒子叫嗎哪西,.
因為神使我忘記一切的創傷和我父全家,.
二兒子叫爾法蓮,.
因為神使我在受苦的地方昌盛,.
每一天兩個兒子的名字都提醒著我,.
是神的話語和祂的同在托住了我的生命,.
我年輕的時候我發過夢,.
人稱為我做做夢的人,.
直到我四十歲和家人重逢的時候,.
那天那些夢才美夢成真,.
原來神一直記得我二十幾年前失落的夢,.
當我聽到哥哥他為了當年的惡行懊悔的時候,.
我的眼淚都已經不能止住,.
如果沒有神,原諒和盼望根本沒有可能,.
後來我安頓了我的爸爸和我的哥哥,.
我的爸爸也為了法老去祝福,.
我的一生就在埃及和家人相伴中度過..
當我們談論到門徒訓練的時候,.
其實我們不是在說如何去做一個完美,.
一定不會受傷的門徒,.
約瑟他前後哭過好幾次,.
如果你回去再細看約瑟的故事,.
這個代表其實他心仍然會痛,.
不過門徒和其他不信者不同之處,.

$^{321}$就是我們不需要獨自去承受一切,.
我們可以有一個交托的對象,.
最後我想分享我在神學院一位好朋友的故事,.
我這個朋友也是一個傳道人,.
我畢業就來了旺鎮,.
他畢業就選擇去了神希戒毒會裡事奉,.
因為在那裡正是他一家認識神的起點,.
我這個朋友他說他回首過往,.
如果沒有神的恩典,.
他和他的家人或者已經早已陷入了一個難以想像的破碎之中,.
我這個朋友他憶起童年的時候,.
他和爺爺一起住,.
父親的身影很少會出現,.
記憶之中他會記得媽媽經常一手抱著妹妹,.
一隻手會牽著他,.
搭著一些很顛簸的大飛,.
就會去到一個小島和他爸爸相會,.
每當他抵達岸邊的時候,.
他都會忍不住奔向父親的懷抱,.
然而這些短暫的相聚之後,.
回程就只有他們三個人,.
看著在岸上父親的身影慢慢地延續,.
很細小的他其實心裡也充滿不捨和疑惑,.
而他發現他的母親雙眼總是會泛紅,.
很想哭,.
直到我這個朋友,.
其實她是一位姊妹,.
去到小學六年級的時候,.
她偶爾發現了她父親在島上寫過一本日記,.
原來她才知道原來她父親那時候正在戒毒,.
所以她無法經常陪著他們兩個,.
直到她小六升中的階段,.
就是她成長的一個轉捩點,.
那一年除了她知道她父親的過去,.
最疼愛她的爺爺就診斷出末期的癌症,.
而這些腫瘤讓爺爺痛得有幾次都很想輕生,.
我朋友感到極度無助,.
於是有一天醫生就告訴他們一家,.
叫他們準備,.
那一晚我這個朋友就跪在床邊祈禱,.

$^{361}$去哀求天父可否感動爺爺去信主,.
因為她很怕永遠都失去了爺爺,.
如果天父你答應她的話,.
她就願意一生去事奉天父,.
而奇蹟真的發生了,.
在那個週末的夜晚,.
她的爺爺竟然主動跟她說,.
孫女,這個星期我應該可以跟你去教會,.
醫生批准我出院,.
於是我的朋友立刻很感動,.
她立刻哭了,.
因為她知道是上帝讓這件事發生,.
後來我的朋友的爺爺不但決志信主,.
他更加領受了洗禮,.
這讓我朋友很深切地體會到,.
神真的很真實,.
事奉天父這個承諾,.
我朋友一直沒有忘記,.
但在真正被神使用之前,.
神希望她先經歷被神醫治,.
在她參加洗禮班的過程中,.
我朋友發現她多次表達她不願意,.
因為原來在很多事情中,.
她仍然想靠自己,.
她不願意很多的不願意,.
是因為她害怕的根源,.
很擔心讓其他人知道她家的過去,.
她很想守護她的家人,.
很想把家庭的秘密深深地收藏,.
亦不願意分享,.
不過神醫治的恩仇,.
卻不斷帶來安慰和勇氣給我朋友,.
我朋友浸禮半年後的一個主日,.
她有一天乘坐小巴,.
突然發了一個夢,.
夢中她看見一個很大發光的十字架,.
並且在夢中聽到一句,.
「你爸爸是在這裡迴轉的.」.
於是她醒來後,.
她決定去神希島一次,.

$^{401}$她很記得她出發的早上,.
當船再次靠岸邊時,.
她看見一塊很大的石頭,.
寫著「耶穌愛你」,.
簡單四個字,.
就很深深地觸動了她,.
讓她明白到,.
其實上帝招聚了很多被社會視為沒有希望的人,.
但是神卻是用福音去救他們,.
神很愛這個群體,.
而我朋友的父親就是其中一員,.
我朋友踏上神希島那一次,.
他是第一次覺得可以很安然地放下自己的羞恥感,.
而在整個參觀神希島的過程中,.
當他再登上山峰時,.
他看見一個很高聳入魂很大的十字架,.
就是這個神希島的十字架,.
這個場景正正就是他在小巴發夢看見的場景,.
這個正午陽光的十字架,.
陽光曬在十字架上很耀眼,.
而十字架寫著的正正就是羅馬書一章十六節,.
「我不以福音為恥」,.
這句說話立即擊中他的內心,.
我朋友向神禱告,.
他願意選擇將多年埋藏在心裡的羞恥感,.
完完全全放在神的恩典裡面,.
在那一刻他立志,.
我朋友願意為這個群體付出,.
因為他知道神要用他和他的家人在生命裡面的經歷,.
這份榮耀去見證上帝,.
回程之後,.
我這個姐妹鼓起勇氣和她爸爸媽媽分享,.
她沒有想到她爸爸竟然很喜樂地問起小島的變化,.
媽媽亦益述當日怎樣去見證父親受洗的片段,.
這是她一家人第一次很坦然分享他們的往事和秘密,.
這份恩典也成為了我這位姐妹踏上讀神學的確據,.
她在小島長洲的見度神學院接受裝備,.
她期盼像她爸爸一樣,.
在小島有一個安靜的時候,.
可以認識神,認識自己,.

$^{441}$今天神帶領我這位朋友回到神希島,.
即神希會的機構成為同工,.
是一班和她有相同背景,.
被神所愛的人當中學習牧羊,.
心中她感到無比的充實,.
弟兄姊妹,不要將過去去定義你,.
你過去可能經歷了一些家庭的破碎,.
父母的缺席,.
甚至是一些不為人知的侵害和委屈,.
這一事的確發生過,.
也深刻地改變了你,.
但今天當你和我來到神的面前,.
我們說很想做耶穌基督的門徒,.
我們願不願意像藥室一樣,.
將那件滋滿血淚的彩衣交到上帝的手中,.
你願不願意跟神說,.
「主啊,我知道這一切都是很痛,.
我都不明白為什麼,.
不過我選擇不再自己捉緊這一份的恨,.
我選擇不再回頭望,.
求你的同在進入我的埃及,.
在我的痛處裡面建立我的馬來西亞和以法連.」.
上帝邀請我們衝破過去的乏力,.
迎向美好的未來,.
這是一個漫長的過程,.
然而當我們踏出信心的步伐,.
我們轉身願意跟隨神的時候,.
我們就會發現,.
神會運用我們過去不堪回首的傷痛,.
開創出我們可以奉獻給這個世界最美好的禮物..
願神賜福給大家!.
\newpage



\section{腓立比書 4:1-14}
\label{sec:NSYx8y4kLPg}
\textbf{2026年6月7日主餐崇拜直播｜陳明泉牧師｜告別浮淺,門徒進深:喜樂的負傷者｜腓立比書4\_1-14}
\newline
\newline
連結: \href{https://youtube.com/watch?v=NSYx8y4kLPg}{\texttt{https://youtube.com/watch?v=NSYx8y4kLPg}} ~~~~ 語音日期: 2026-06-07
\newline
\newline
\hyperref[sec:dKG666yV4KA]{\small{< < < PREV SERMON < < <}}
~
\hyperref[sec:index_chronic]{\small{[返順時目]}}
~
\hyperref[sec:wKhk31WV7Os]{\small{> > > NEXT SERMON > > >}}
\newline
\newline
腓立比書 4:1-14
\newline
\begin{longtable}{cl}
\hline
\hline
章節 & 經文 (和合本修訂版)\\
\hline
4:1 & \begin{tabularx}{0.7\textwidth}{X} 我所親愛、所想念的弟兄們，你們就是我的喜樂，我的冠冕。我親愛的，你們應當靠主站立得穩。 \end{tabularx} \\ \\ \relax
4:2 & \begin{tabularx}{0.7\textwidth}{X} 我勸友阿蝶和循都基要在主裡同心。 \end{tabularx} \\ \\ \relax
4:3 & \begin{tabularx}{0.7\textwidth}{X} 我也求你這真實同負一軛的，要幫助這兩個女人，因為她們在福音上曾與我、革利免和我其餘的同工一同勞苦，他們的名字都在生命冊上。 \end{tabularx} \\ \\ \relax
4:4 & \begin{tabularx}{0.7\textwidth}{X} 你們要靠主常常喜樂。我再說，你們要喜樂。 \end{tabularx} \\ \\ \relax
4:5 & \begin{tabularx}{0.7\textwidth}{X} 要讓眾人知道你們謙讓的心。主已經近了。 \end{tabularx} \\ \\ \relax
4:6 & \begin{tabularx}{0.7\textwidth}{X} 應當一無掛慮，只要凡事藉著禱告、祈求和感謝，將你們所要的告訴神。 \end{tabularx} \\ \\ \relax
4:7 & \begin{tabularx}{0.7\textwidth}{X} 神所賜那超越人所能了解的平安，必在基督耶穌裡，保守你們的心懷意念。 \end{tabularx} \\ \\ \relax
4:8 & \begin{tabularx}{0.7\textwidth}{X} 末了，弟兄們，凡是真實的、凡是可敬的、凡是公義的、凡是清潔的、凡是可愛的、凡是有美名的，若有甚麼德行，若有甚麼稱讚，你們都要留意。 \end{tabularx} \\ \\ \relax
4:9 & \begin{tabularx}{0.7\textwidth}{X} 你們從我所學習的，所領受的，所聽見的，所看見的事，你們都要繼續去做，賜平安的神就必與你們同在。 \end{tabularx} \\ \\ \relax
4:10 & \begin{tabularx}{0.7\textwidth}{X} 我靠主大大喜樂，因為你們關懷我的心如今又表現了出來；其實你們一直都關懷我，只是沒有機會罷了。 \end{tabularx} \\ \\ \relax
4:11 & \begin{tabularx}{0.7\textwidth}{X} 我並不是因缺乏而說這話，因為我已經學會無論在甚麼景況都可以知足。 \end{tabularx} \\ \\ \relax
4:12 & \begin{tabularx}{0.7\textwidth}{X} 我知道怎樣處卑賤，也知道怎樣處豐富；或飽足或飢餓，或有餘或缺乏，任何事情，任何景況，我都得了祕訣。 \end{tabularx} \\ \\ \relax
4:13 & \begin{tabularx}{0.7\textwidth}{X} 我靠著那加給我力量的，凡事都能做。 \end{tabularx} \\ \\ \relax
4:14 & \begin{tabularx}{0.7\textwidth}{X} 然而，你們能和我分擔憂患是一件好事。 \end{tabularx} \\ \\
[1ex]
\hline
\hline
\end{longtable}
$^{1}$今天要讀出的經訓是保羅書信中的《斐納比書》第四章一至十四節,在新約聖經的二百七十六頁..
我所親愛,所想念的弟兄們,你們就是我的喜樂,我的冠冕.我親愛的弟兄,你們應當靠主站立得穩..
我勸友愛的,和秦都基,要在主裡同心.我有求你這真實和負一額的幫助者兩個女人,因為她們在福音上曾與我一同勞苦..
民有生氣面,並其如和我一同作供的,她們的名字都在生命冊上.你們要靠主常常喜樂.我再說,你們要喜樂..
弟兄們,我還有未盡的話.凡是真實的,可敬的,公義的,清潔的,可愛的,有美名的,若有什麼德行,若有什麼稱讚,這些事你們都要思念..
你們在我身上所學習的,所領受的,所聽見的,所看見的,這些事你們都要去行,此平安的神就必與你們同在..
我靠主大大的喜樂,因為你們思念我的心,如今又發生.你們向來就思念我,只是沒得機會.我並不是因缺乏說這話.我無論在什麼情況都可以知足,這事我已經學會了..
我知道怎樣處悲戰,也知道怎樣處豐富,或飽足或飢餓,或有餘或缺乏,隨是隨在,我都得了被缺..
我靠著那加給我力量的,凡事都能做.然而,你們和我同受患難,完是美事..
我們常覺得一些靈性很高的人,他們的情感波動會比較小,甚至是一條水平,甚至是一條圓環..
有一些屬靈傳統真的強調我們太過情緒高高低低是不好,.
不過我們未必一定要按著某一種屬靈傳統,覺得屬靈的人通常沒所謂,沒什麼開心或不開心..
另外有一些來自哲學家的說法,他說其實人生命裡對於快樂或人的深度和痛苦兩者之間沒有完全的關係..
簡單來說,我們今天經常在坊間說我們要快樂是怎樣呢?少一點痛苦,少一點不開心的東西..
總之我們人生活在一個很平順的日子,我們的生命就快樂了..
但是有些哲學家就說他們再去深思去發覺原來人的快樂和人所經歷的痛苦兩者可以並存在人生命當中..
所以在你的心靈裡既有快樂的記憶,但同時間也有一些痛苦..
不過這幾件事有什麼關係呢?我說幾個哲學家很有意思的..
他說人的不開心不是因為苦難多,而是人的痛苦來自於人的無知..
這就是蘇格拉底所說的,人的理性發揮不了果效,以致人活在一種想不通的地步..
所以蘇格拉底鼓勵人真的不是避免地上的受苦,而是要追求人的智慧,人的理性不斷地提升..
藉著美德和藉著人的選擇,可以將我們人生裡經歷一些逆流或者不開心的事情想得通,那我就可以快樂了..
所以在蘇格拉底他覺得快樂就要成為一個義人,而且是懂得思考的一個人..
第二個人叫帕拉圖,這個帕拉圖大家都聽過了,不過他講了一個比喻很有意思的,很好笑的..
他說人生像什麼呢?像嘔痕,什麼嘔痕呢?.
他說原來人處於一種不平衡的狀態就會有不開心,或者人開心就要處於平衡的狀態..
他是用嘔痕來作為一個比喻,當你的皮膚很癢的時候,有一個動作會做的就是幫它嘔..
嘔這個動作就令到你身體取得一種平衡,其實嘔本身刮在皮膚裡面是產生痛楚的,癢癢和痛楚中和了,平衡了,那你就得到舒暢了..
他用這個角度再去延伸,其實人生裡面最重要的不是沒有痛苦,而是能夠在人生裡面產生一種和諧..
你有其他的東西來抵消,或者藉著你的理性,你的理智來將一些東西調和,這個就是帕拉圖看的人生的快樂了..
第三個,這個是近代的哲學家叫做尼采,尼采就完全不同意蘇格拉底和帕拉圖所講的東西..
他覺得痛苦是人生命的催化劑,他反對人世間活在一種紙醉金迷,或者滿足肉體的快樂..
他覺得快樂和痛苦是一銀兩面,是孿生子女,所以你要快樂的時候,你必須要承受同樣份量的痛苦..
他特別講到,痛苦的解決帶來的那種成就感,或者你面對挑戰的時候,你努力跨越過了,那種滿足,那種快樂是世間其他東西不能夠媲美的..
不過在尼采的這個講法裡面,人要成為一個怎樣的人呢?人要成為一個超人,你要快樂,你就要成為一個超人,你要克服所有困難,你就快樂了..
剛才林林總總的幾個哲學家講的,這些叫做人本的角度來看快樂,人本的角度到最後是講人自己製造快樂給自己..
不論是理智的識想,或者用人的方法來調和,甚至你要令到你自己有一個強力的意志,克服所有的困難,令到你人生命裡面有一個快樂..
可能會有一些提醒,或者有一些啟發的,不過聖經似乎和剛才那幾個哲學家,甚至乎坊間裡面所講的快樂,是截然不同的一個不同的向度..
今天我們看的菲勒比書,菲勒比書裡面有重複的字眼,就是快樂,喜樂這個字..
今天我們就用快樂神學這個角度,重新看一下,一個門徒盡心的人,一個寫書卷的寶囊,面對著牢獄之苦,面對福音的時候帶來大量的逼迫,.

$^{41}$外面肉體不斷的受到侵害,但是心底裡面有快樂,究竟是怎樣做得到?.
他的神學怎樣去令到他成為一個快樂的人?.
希望在他身上,在聖經裡面我們拿到養分,我們先祈禱..
求上主幫助我們,讓我們直著你的話語,我們生命被建立,求你與我們同在,讓我們每一個領受你的話語的時候,.
我們不單止聽得明白,更加讓我們內心好像有一團火提升,提升我們,讓我們的心靈可以更加接近你..
幫助孫講的,讓弟兄姊妹聽得明白,讓我們在這個時間一起領受你的話語,.
祈求你與我們同在,帶領我們,獻上禱告奉基督耶穌得勝的聖名..
今天我所講的肥勒比書其實很豐富的內容,其實我們講道裡面講過很多次了,.
不過今天我用一個比較神學的角度,我又怕很心餓,所以我用了四個詩來講四個嚮導,讓大家知道..
第一個詩就叫做Communal Commitment,四章第一節..
我所親愛,所賞念的弟兄們,你們就是我的喜樂,我的冠冕,我親愛的弟兄..
你們應當靠主站立得穩..
我們再看第十節..
我靠主大大地喜樂,因為你們思念我的心如今又發生,你們向來就思念我,只是沒得機會..
我們再看第三節..
我也求你這真實同乎一額的幫助這兩個女人,因為她們在福音上曾與我一同勞苦,.
患有隔離免,並其餘和我一同作供,她們的名字都在生命冊上..
這幾節經文其實都是在同一個篇幅裡面的..
它講道其實保羅的快樂來自什麼呢?.
不是靠自己會想,原來上帝所命定的,或者上帝的心意在快樂是透過群體裡面來產生的..
所以快樂不是在一個孤單人的生命當中實踐,.
也不是一個我純粹自己抽離了一個信仰群體,自己明白聖經就可以得到快樂..
保羅在這裡特別斬釘截鐵地講道,.
他特別講道他的喜樂就是講道這一班屬靈群體肥納比教會的弟兄姊妹就是他的喜樂,.
不用他們做些什麼,保羅想起他們就是喜樂..
當然有些歷史的,因為肥納比教會是保羅在歐洲裡面第一個建立的信仰群體..
某程度上對於保羅來說,這個群體很有標誌性的,很有回憶的..
但是保羅他寫出的這句話,他講道這班群體是保羅除了一手一腳建立,.
他不是很強調自己做些什麼,.
他只是講道保羅想起他們,他的心就快樂了..
當然這一班的群體,他們也有一些努力為保羅來做的,.
保羅不是純粹想起他的,也有一些很具體的,.
例如在經文裡面講道肥納比教會為保羅宣教的需要,他供應保羅所需..
你要記住保羅其實他不是一個駐堂牧者,他是一個宣教士..
他建立了一個群體之後,他就不會純粹停留在一個群體裡面,.
他就離開繼續去承擔福音照明,但是他離開不等於他們的關係就斷裂了..
保羅去到不同的地方,去到傳福音的時候,他想起這一班的群體,.
想起這一班弟兄姊妹,繼續的有胸襟的將他們僅有的資源,.
那時候的教會不是很有錢,將僅有的資源聚合,.
然後供讓保羅所需,可能是一些金錢,可能是一些實質的需要,.

$^{81}$可能是一些衣物等等,這些都成為了保羅點滴在心頭的關懷..
所以對保羅來說,這個群體對他帶來的就是莫大的喜樂..
更加重要的就是,這班群體他雖然是一個可以跟他有一個供應,.
但是他更加講到這班群體即使在地域上面跟保羅,.
未必跟著保羅到處隨走隨存,但是這一班人對於福音的照明,.
他們是同一個心的,同一個心智的..
所以在經文裡面他特別講到,他說是同父一額的,.
他講到這個群體跟他自己是背擔著同一個福音照明,.
一起來努力,而且為著福音的緣故一同勞苦..
所以不單止這個群體供應他,保羅想起他,.
想起這班群體,跟他是同波段的,是一起努力的..
我在旺朕有時候想起弟兄姊妹有些情景,.
我在這篇講章預備的時候想起一些圖畫,.
想起以前一起在街上被人傳福音,街頭騙傳,一起吃檸檬,.
又或者我們一起去到學校裡面,一班弟兄姊妹,.
一起來出隊,缺支人數不是很多,但是我們一班人一起努力,.
一起來打點,又想起我們在旺朕裡面,.
我們在這裡一起來做一些福音聚會..
我記得有時候,剛剛我事奉母女,我們少年人組織一個報道會,.
有一個年輕人說,不如這次玩大他,我說怎樣玩大他?.
他說名字叫玩大他,於是我們在四樓做一個叫玩大他的報道會,.
玩些什麼呢?我們用舊時運貨的小車,那塊木板車,.
前面放一些膠的東西,將人一樣掃下去,人肉保齡球玩這些,.
將我們平時玩的那些層層疊放大,那些叫玩大他,.
一直玩得很開心,一起去做這些工作..
你說福音果後是不是很多呢?有一些已經繼續在我們當中,.
不過重要的就是一起努力,一起來做..
今日我們的社會,今日的教會,我們很需要這些動力,.
福音工作不是一小撮人或者幾個人給他們自己回憶,.
不要這麼不公道,快樂是應該一起來做,在一個群體裡面,.
一起來彼此配搭,一起來努力,捏起衣袖一起做,.
這個就是我們教會,我們有的DNA..
保羅想起這一班人都是一起為福音老虎供應他所需,.
而且這個關係令到他帶來心裡面有喜樂,.
不過這個群體也不是完美的群體,.
保羅再特別提到兩個名字,「友阿諜」和「秦都姬」,.
說這兩個姐妹在主裡面要同心,.
幾個字你大概體會到肥勒比書裡面的寫作處境,.
弟兄姊妹之間可能各自熱心,但是有一些爭拗,.
有一些不同心,或者大家有各執一詞,有些堅持..

$^{121}$保羅他不看這些堅持就覺得很煩,不要煩,不要叫,.
他不是,他用一個很正面,很喜樂的角度,.
大家一起事奉,大家一起在主裡面彼此同心,.
即使大家有不同意見,但是都可以同心,.
關鍵是我們有沒有願意放我們自己的心力,.
來彼此建立,所以這個群體不是完美的群體,.
這個群體雖然有供養保羅,.
但是上帝就命定了人的快樂在這個群體裡面發生,.
我們就必須要尊重,.
我們甚至乎要毀身,投資時間心力在一個群體當中,.
這個是第一個詩,喜樂神學不是靠自己,.
是靠群體,在群體裡面建立一種喜樂的關係..
第二個段落四至七字,你們要靠主常常喜樂,.
我再說,你們要喜樂,.
當叫眾人知道你們謙讓的心,主已經近了,.
將你們所有的告訴神,神所賜出人意外的平安,.
必在基督耶穌女保守你們的心懷意念..
第二個段落最重要的字眼,.
或者其實整個篇章裡面最重要的字眼就是,.
那個喜樂不是自己,而是靠主喜樂,.
保羅在這裡面多次強調,你的喜樂是要靠上帝,.
和人間的智慧或者人為操作不同,.
喜樂是來自上帝的,.
人是要透過上帝,要依靠上帝,.
生命才能產生喜樂這種態度和心內的體會..
經文裡面他再繼續說,.
他怎樣去演繹那個靠主喜樂的那種內涵呢?.
在第五節裡面他特別強調的就是,.
當叫眾人知道你們謙讓的心,主已經近了..
在這裡面要解釋一下的,.
首先,他講到你的謙讓,.
謙讓的意思其實回應菲勒比書講到,.
耶穌基督就是一個最謙卑榜樣的那種生命,.
所以你時刻以基督耶穌的心為心,.
體會耶穌基督由高降到卑,.
然後上帝又將他由最低升為最高,.
這個這樣的V型反彈,.
他特別講到耶穌自己反彈的那個是上帝給他的,.
但是下到去最謙卑的是他自己的努力,.
同樣他也鼓勵我們每一個弟兄姊妹都帶著這一份心態,.

$^{161}$成為一個謙讓的人,一個謙卑的人,.
或者你覺得謙讓,謙卑都好像太過文雅,.
其實他原文用的字眼其實是成為一個溫柔的人,.
成為一個有胸襟,大度的人,.
簡單來說你的胸襟越廣闊,.
你不再那麼多的介懷,介意,.
你越放得自己越低,.
你所得的喜樂就越高,.
而且這一種不是一種表現,.
當叫眾人知道你謙讓的心,.
是不是在說做一場戲給人家,.
不是,在群體當中,.
在人與人之間的關係,.
你就實踐那種的溫柔,.
那種的大度,那種心裡面的那種,.
有一種的胸襟去容納不同的人,.
這個就是一個很重要的,.
靠主喜樂裡面的那種生命表現,.
藉著一個你靠主的人,.
你就看耶穌的榜樣,.
然後我們就成為耶穌那個心為心的一個信徒,.
經文裡面後面加個字眼很有意思,.
他說主已經近了,無端端有時候讀,.
以為是不是主已經近了,是說後面的,.
不是,原來是同一句,.
我們要問主已經近了,.
關我們謙讓是什麼事呢?.
有幾個層面來講,.
這些色經學家給一個很重要的,.
首先他講到主已經近了,.
他講到人不能夠謙讓,.
原因就是我們寸土不讓,.
我們整天覺得自己最有道理,.
我們會瓦起我們的爭頭,.
不過如果在我們的角度裡面覺得主已經近了,.
就是說上帝審判的時間到了,.
你不用為自己做什麼的申冤,.
你不用為自己來抗辯,.
上帝已經在這裡了,.
上帝會為你作主,.

$^{201}$所以你就不用耿耿於懷,.
覺得別人欺負你,.
在時間上面上帝回來的時間越來越近了,.
上帝會幫你申冤的了,.
所以我們就可以放下那些爭頭,.
這個在時間上面主已經近了,.
就是說耶穌越來越近,.
回來這個時間,.
所以我們只管放下自己,.
不用整天瓦起爭頭,.
第二個是空間上所講上主已經近了,.
就是說其實上帝一直在我們身邊,.
耶穌基督與我們同在的,.
我們整天都很想為自己抗辯,.
覺得含冤受屈,.
我們很想告訴別人,.
我們據理力爭,.
不過上帝與我們身邊,.
上帝對我們自己一清二楚,.
有時候我們覺得大聲講道理,.
就等於道理在我們那邊,.
不過在菲勒比書裡面說,.
不是道理在我那邊,.
而是上帝在我那邊,.
我們就可以放下自己,.
我們不用耿耿於懷,.
變得大道,.
就是在心裡面知道,.
耶穌就近在我趾尺,.
我就可以帶著這一份有信心,.
有底氣,.
我們可以不需要斤斤計較,.
第三,更加講道的就是,.
讓別人看到我們的生命主已經近了,.
既然上帝與我們同在,.
又一起回來了,.
接下來的日子越來越近了,.
我們更加要讓別人看到,.
在我們的生命裡面,.
活出被主體見,.

$^{241}$或者活出主已經到的一種生命表現,.
所以保羅在這裡說的,.
一人小句,.
退一步海闊天空,.
他不再這麼浮淺說這些了,.
他再用一個神學,.
用耶穌基督的生命,.
用上帝的計劃,.
來調教他的心思意念,.
以至到他的心靈裡面,.
活出一種坦蕩蕩和寬容,.
人與人之間就不再變得這麼繃緊,.
經文裡面再繼續說,.
除了呈現一種謙讓的心之外,.
他更加講道,.
人需要藉著禱告來交托,.
應當一無掛慮,.
將你所有的告訴神,.
神所賜出人以外的平安,.
必在基督耶穌女保守你們的心懷意念,.
你不要讀錯這句經文,.
將你所有的告訴神,.
你有些頂梓梅就想著後半句,.
神就將你所求的一定必給你,.
經文不是這麼說,.
他不是說你求什麼他就給你什麼,.
他只是說,.
你想要的東西告訴上帝,.
給不給是上帝的主權,.
最重要的就是,.
當你祈求交托的時候,.
上帝會將一份出人意外,.
地上沒有東西能夠媲美的平安,.
賜給你和我,.
我們內心裡面,.
當我們有這份平安,.
你就會對照一下,.
你所求的東西是不是這麼重要的東西,.
當有上帝而來的平安,.
你就對照一下,.

$^{281}$今天我們內心裡面耿耿於懷的事,.
令到我們不開心,.
相對於那份平安,.
可能那件事變得微不足道,.
所以經文裡面說到,.
我們的心懷意念,.
藉著這些禱告,.
其實練成了我們對善物,.
對我們所渴求的那些事,.
那種更新和調教,.
如果沒有這些調教,.
我們經常都是那樣東西,.
到最後我們心靈不開心,.
因為太過執著,.
太過覺得我們自己要的東西,.
上帝一定要給我們,.
但上帝給的可能超乎我們所想像,.
甚至上帝將我們現在這一刻.
未覺得好的東西,.
但將來會覺得好的,.
祂會給我們,.
求主幫助我們,.
所以第二個,.
第二個詩,.
你要記住,.
就是Cling to the Lord,.
勾住上帝,.
我們用這個詞來說,.
我們要靠主而起來,.
第三個,.
第八節,.
弟兄們,.
我還有未盡的話,.
凡是真實的,.
可敬的,.
公義的,.
聖潔的,.
可愛的,.
有美名的,.
若有什麼德行,.

$^{321}$若有什麼稱讚,.
這些事你們都要思念,.
你們在我身上所學的,.
所領受,.
所聽見的,.
所看見的,.
這些事你們都要去行,.
最平安的神,就必與你們同在,.
在這裡,.
你會見到,.
保羅再強調,.
我還未說完,.
我還有話要說,.
他就是這個意思,.
他叫弟兄姊妹,.
腦裡要念記一些事,.
要思想,.
或者要重新想起這些事,.
想起什麼呢?.
真實的,.
即是有真理的,.
可敬的,.
即是令人望而生畏的,.
公義的,.
清潔的,.
可愛的,.
即是令你賞心悅目,.
有美名的,.
即是說他活出一個美德的生命,.
有德行,.
稱讚,.
這些你們都要思念,.
簡單來說,.
就是我們的內心裡面,.
要儲多些正面的經驗,.
要珍惜這些體會,.
珍惜,.
在我們人生當中,.
雖然這個世界很黑暗,.
很多時候,.

$^{361}$弟兄姊妹的歌有時候都這樣唱,.
黑暗世界,.
正正這個世界黑暗,.
我們的視野,.
不是只看到這個黑暗的世界,.
我們和微信者有什麼分別呢?.
如果我們看到這個世界真是黑暗,.
就是這個世界有黑暗,.
但我們看到依然上帝掌管,.
我們依然看到上帝.
還有作功在這個世界當中,.
我們的視野看見的是上帝,.
將恩典,.
將美好的事物,.
依然在一個我們覺得敗壞,.
不理想的世界裡面,.
而且這些東西不是看過就過眼雲煙,.
我們看過這些事,.
我們要銘記內心,.
我很記得我剛剛出道,.
我們舊時的主任牧師,.
梁應安牧師,.
他教我這個身仔,.
他說,.
明泉,.
你有沒有想過,.
你將來牧會久了的時候,.
弟兄姊妹期望你給他們一些東西,.
期望你貪病,.
給他們一些金句,.
或者在人生某些道理裡面,.
給他們一些智慧,.
這些東西是怎樣來的?.
我說讀聖經,.
讀完聖經,.
怎樣令他們明白你所說的話呢?.
還有,.
經常每個人都問你,.
你真的有這麼多東西給人嗎?.
那一刻啞口無言,.

$^{401}$真的不知道怎樣回答牧師,.
這條題目這麼深,.
他說其實不是很深,.
你要讓你的群體牧養你,.
他們未必說到一些至理名言,.
他叫我觀察一下,.
在弟兄姊妹生命當中,.
看到上帝的作為,.
你捕捉弟兄姊妹生命裡面,.
去經歷到神的那些工作,.
那些弟兄姊妹說不到給你聽,.
但是你的眼可以看到,.
是可以整合的,.
他說這個叫做inward journey,.
就是說在你的生命裡面,.
你要做到長做長有,.
關鍵不是你讀多少書,.
關鍵是你能不能夠從弟兄姊妹生命裡面,.
發現到上帝的足跡,.
在這些生命裡面,.
找回一些牧養你自己的養分,.
這些拿到的養分,.
就是你用來牧養弟兄姊妹的資源,.
你這樣做就能夠長做長有,.
簡單來說,你覺得你自己很有需要,.
或者很有責任感,.
將你自己學到的東西全部分享了,.
你有一天會乾枯的,.
一個乾枯的牧者,.
和一個做極都有課的牧者的分別就是,.
在弟兄姊妹在群體裡面,.
看到上帝的足跡,.
他給了我這句,.
二十多年了,.
我今天依然受用,.
在大家的生命當中,.
我看到很多上帝的作為,.
很多的小故事,.
你們生命裡面的經歷,.
其實在牧養我,.

$^{441}$所以那種珍惜這種的關係,.
珍惜我們生命裡面所有的,.
這些就是成為了我們生命裡面一些很重要的養分,.
還有一樣東西很好笑的,.
我的同道經常都聽到我說,.
Anders你真是很hard sell,經常都說黃浸很厲害,.
說到所有的弟兄姊妹,.
又可以上刀山下來游,.
什麼都懂的,什麼都那麼厲害,.
我說是的,我的弟兄姊妹很厲害的,.
幫我們手打理我們大堂,.
我們場地的那個姐姐,.
很厲害的,.
有什麼厲害的?.
不是很厲害的,.
是國家隊運動員來的,.
之前年輕,.
大家知不知道的?.
不知道的?.
我們稍後訪問一下她,.
國家隊的運動員代表過國家出去參加奧運的,.
嘩,.
這些你說我說出來是不是很厲害?.
不是我自己很厲害,.
而是上帝將一些不同故事的人放在我們的群體當中,.
我每一次一說起,.
我就眉飛色舞,.
然後我說是的,黃震真的很厲害,我們這個群體很厲害的,.
然後他另外跟我說一個童話,.
他說我教會不是這樣的,.
我說不是,你教會不是這樣的,.
只不過你沒有發現他們生命裡面上帝奇妙的作為,.
我們不是要比較這些人厲害些,那些人什麼,.
而是我們沒有珍惜生命裡面那些故事,.
看到那些弟兄姊妹生命最良藝的地方,.
我們將這些故事成為了我自己故事的一部分,.
我們就對我們自己的群體,.
對我們所身處的這種掛會關係,.
我們就會有很多很正面,.
也都很激勵自己的這種體會,.

$^{481}$最後,十至十四節,.
又是繼續重複,我靠主大大地起來,.
因為你們思念我的心如今又發生,.
你們向來就思念我,只是沒得機會,.
我並不是因為缺乏說者說,.
我無論在什麼境況都可以知足,.
這是我已經學會了,.
我知道怎樣處悲戰,.
也知道怎樣處豐富或飽足或飢餓,.
或有餘或缺乏,.
隨時隨在,我都得了被缺,.
我靠著那加給我力量的,.
凡事都能作,然而你們和我同受患難,.
原是美是..
在這裡某程度上作為段落的一個總結,.
也是保羅的豪情壯語,.
他再強調他的喜樂是靠主而來的,.
他來自上帝的喜樂,.
令到他的心裡有一份知足,.
有一份滿足感,contentment,.
他覺得他內心充滿了內容,.
在他的信仰生命裡不是蒼白一片,.
而是他經歷到心靈裡因為上帝的公認,.
藉著群體的紀念他,.
他想起教會時刻都想起他,.
就是一個這麼簡單的想念,.
對保羅來說,.
當他自己面對著逼迫的時候,.
他想起教會弟兄姊妹會擔心他,.
會想起他,會繼續支持他,.
他就有力量,.
即使多缺乏也好,.
他心靈的富足就可以抵禦眼前的困苦和困難,.
保羅再說,.
我不是因為沒有錢跟你說這番話,.
而是說到他什麼狀況,.
他的心裡都可以滿足,.
因為上帝讓他看到他的計劃,.
上帝把他成為福音的執事,.
他明白耶穌基督的心腸,.

$^{521}$他更加體會到上帝設立的教會.
和他之間這麼緊密的情誼,.
所以他心裡所看的自己的人生,.
雖然面對著大量的困苦,大量的苦難,.
但他看自己人生不是貴乏的人,.
相反來說,今天我們可以很有錢,.
我們可以吃過世的,.
但我心靈也可以很貴乏,.
重點是因為我體會到.
我生命的豐富從來不是由物質而來,.
心靈的豐富是來自上帝,.
看到上帝在我生命裡所建立的關係,.
建立了生命裡很多的恩典,.
我生命就是由這些積惑有苦難,.
都有大量的恩典交織而成的美好的生命,.
回顧去看,一切都滿足,.
福福何求?.
帶著這份心態,.
缺乏又如何?.
鞭打又如何?.
坐牢又如何?.
總之心裡踏實,安穩,.
一切都是足夠,.
生命就可以豁然開朗,.
時刻就快樂..
弟兄姊妹,我不知道大家每一個的故事,.
不過今日社會瀰漫著一個.
整天都很低沉,.
貧窮的時候又說怕沒錢,.
現在好一點又說怕保不保得住,.
樓價升又擔心,樓價跌又擔心,.
我們基本上就被這個社會的風氣.
影響了我們內心裡的快樂,.
不過在聖經裡說你快不快樂,.
根本就和經濟環境沒有絕對的關係,.
甚至你也不應該看經濟環境,.
看社會來定義自己是否快樂的人,.
即使去到保羅,面對著全福音的患難,.
面對著福音帶來的很多傷害,.
心靈依然可以豁達,.

$^{561}$因為最重要的是人的快樂不是靠自己,.
不是靠環境,.
我們只是靠主純常喜樂,.
弟兄姊妹,鼓勵大家,.
今天聽完這篇訊息,.
我們立志,我們要成為一個快樂的人,.
我們更加認定快樂不是要做到超人,.
或者要做到很理性很強的人,.
我們只是做一個倚靠上帝的人,.
我們按著今天所教導的,.
我們就可以快樂,.
其實就可以在我們生命當中,.
可以很近,求主幫助我們,.
願我們同在..
願我們同唱「主活著」.
在《世紀眾讚》的第205首,.
願我們知道當主活在我們心的時候,.
我們的生命可以得著滿足..
\newpage



\section{約翰福音 12:9-19}
\label{sec:wKhk31WV7Os}
\textbf{2026年6月14日崇拜直播｜梁家麟牧師｜尊祂為王｜約翰福音12\_9-19}
\newline
\newline
連結: \href{https://youtube.com/watch?v=wKhk31WV7Os}{\texttt{https://youtube.com/watch?v=wKhk31WV7Os}} ~~~~ 語音日期: 2026-06-15
\newline
\newline
\hyperref[sec:NSYx8y4kLPg]{\small{< < < PREV SERMON < < <}}
~
\hyperref[sec:index_chronic]{\small{[返順時目]}}
~
\hyperref[sec:c5rLyKd8AcQ]{\small{> > > NEXT SERMON > > >}}
\newline
\newline
約翰福音 12:9-19
\newline
\begin{longtable}{cl}
\hline
\hline
章節 & 經文 (和合本修訂版)\\
\hline
12:9 & \begin{tabularx}{0.7\textwidth}{X} 有一大群猶太人知道耶穌在那裡，就來了，不但是為耶穌的緣故，也是要看耶穌使他從死人中復活的拉撒路。 \end{tabularx} \\ \\ \relax
12:10 & \begin{tabularx}{0.7\textwidth}{X} 於是眾祭司長商議連拉撒路也要殺了， \end{tabularx} \\ \\ \relax
12:11 & \begin{tabularx}{0.7\textwidth}{X} 因為有許多猶太人為了拉撒路的緣故，開始背離他們，信了耶穌。 \end{tabularx} \\ \\ \relax
12:12 & \begin{tabularx}{0.7\textwidth}{X} 第二天，有一大群上來過節的人聽見耶穌要來耶路撒冷， \end{tabularx} \\ \\ \relax
12:13 & \begin{tabularx}{0.7\textwidth}{X} 就拿著棕樹枝出去迎接他，喊著：「和散那，以色列的王！奉主名來的是應當稱頌的！」 \end{tabularx} \\ \\ \relax
12:14 & \begin{tabularx}{0.7\textwidth}{X} 耶穌找到了一匹驢駒，就騎上，如經上所記： \end{tabularx} \\ \\ \relax
12:15 & \begin{tabularx}{0.7\textwidth}{X} 「錫安的兒女啊，不要懼怕！看哪，你的王來了；他騎在驢駒上。」 \end{tabularx} \\ \\ \relax
12:16 & \begin{tabularx}{0.7\textwidth}{X} 門徒當初不明白這些事，等到耶穌得了榮耀後才想起這些話是指他寫的，並且人們果然對他做了這些事。 \end{tabularx} \\ \\ \relax
12:17 & \begin{tabularx}{0.7\textwidth}{X} 當耶穌呼喚拉撒路，使他從死人中復活出墳墓的時候，同耶穌在那裡的眾人就作見證。 \end{tabularx} \\ \\ \relax
12:18 & \begin{tabularx}{0.7\textwidth}{X} 眾人因聽見耶穌行了這神蹟，就去迎接他。 \end{tabularx} \\ \\ \relax
12:19 & \begin{tabularx}{0.7\textwidth}{X} 法利賽人彼此說：「你們看，你們一事無成，世人都隨著他去了。」 \end{tabularx} \\ \\
[1ex]
\hline
\hline
\end{longtable}
$^{1}$主持人:以下是聆聽聖導的時間,今天所選讀的經文是在約翰福音12章9至19節,在新約聖經的145至146頁,請聽我讀出..
主持人:約翰他從死裡所復活的拉薩路,但祭司長雙爾連拉薩路也要殺了,因為好些猶太人為死裡所復活的拉薩路的緣故回去信了耶穌..
主持人:逐中樹枝出去迎接他,喊著說:「和撒拿,奉主名來的以色列王,是應當稱頌的.」.
主持人:耶穌得了一個盧驅,就騎上如經上所記的說:「錫安的民,不要懼怕,你的王騎著盧驅來了.」.
主持人:這些事門徒起先不明白,等到耶穌得了榮耀以後,才想起這話是指著他捨的,並且眾人果願向他這樣行了..
主持人:當耶穌呼喚拉薩路,叫他從死裡復活出墳墓的時候,和耶穌在那裡的眾人就作見證.眾人因聽見耶穌行了這神蹟,就去迎接他..
主持人:法利賽人彼此說:「漢啊,你們是徒勞無益,世人都隨從他去了.」聖經獨筆,願神賜福上讀他話語並且遵行的人..
主持人:各位主持人,早安!很久不見,繼續漫長的約翰福音之旅..
主持人:其實我們只去了一半,如果看耶穌的故事的話,我們已經來到耶穌在人間事奉生涯的最後一段..
主持人:我們這一段聖經的開始是說到耶穌在人間最後一個禮拜的故事..
主持人:最後一個禮拜又叫聖周,Holy Week..
主持人:這個禮拜是從耶穌基督榮耀進入耶路撒冷開始的..
主持人:這個入城的故事是非常重要的,四卷福音書都有記載,不過約翰的記載和三卷符類福音有很大的不同..
主持人:首先約翰是很重視拉撒路復活的事件的,十二章一到八節,瑪利亞用香膏勾耶穌的腳,.
主持人:以至這一段耶穌榮耀進入耶路撒冷,約翰將這兩件事都拉在一起,.
主持人:看是和拉撒路的復活有關的,可以說都是拉撒路的復活事件的延續性的影響..
主持人:耶穌在伯大利使拉撒路復活,這個死人復活的事件太過轟動了..
主持人:不少從耶路撒冷來參加拉撒路的喪禮的人,都因為目睹這個神蹟而信了耶穌,十一節..
主持人:當然他們的相信是不是真是死心塌地徹底的相信,這個我們永遠可以再講..
主持人:不過我真的要講,對我們這些小老百姓來說,我們的信心永遠都是高高低低,.
主持人:只是一時一樣的,什麼時候真的說很清楚從一而終呢?.
主持人:通常都不是的,所以也不用太過批評他們,他們說他們信耶穌,是真的信的..
主持人:他們最少一定確定耶穌沒有辦法拒絕否定的法力,.
主持人:他一定是一個擁有某個神聖身份的人,大概就是舊約應許要來的那個彌賽亞..
主持人:並且整件事還在不斷發酵,十七到十八節那裡說,當耶穌呼喚拉撒路叫他從死裡復活出墳墓的時候,.
主持人:和耶穌在那裡的眾人就作見證,眾人因聽見耶穌行了這神蹟就去迎接他..
主持人:伯大利的目睹者將他們看到的神蹟帶回耶路撒冷到處告訴別人,他們不一定是很正面的作見證,.
主持人:只不過是這麼過癮的事,為什麼不說呢?.
主持人:你知道古代社會,一個非常封閉的社會,沒有什麼新聞,沒有什麼話題,.
主持人:人見到面就是聊天,聊天不是說是非有什麼好說的,基本上難得有件過癮的事出來說,.
主持人:你可以想像他們一定是說來說去,並且一定是一傳十,十傳百,以訛傳訛..
主持人:就算完全沒有目睹的,我間接聽三首先聽回來的故事,我也會說給第四手聽的..
主持人:所以整件事就是這樣,不停以說是非的方法,越說越軟,越說越軟..
主持人:引起很多人的關注,聽到的人都想一睹施行神蹟的耶穌的風采..
主持人:所以當一天之後,耶穌從伯大利轉到耶路撒冷的時候,.
主持人:竟然就有大批的群眾聚集在城門口,揮著鐘樹枝向耶穌歡呼..
主持人:永耀進入耶路撒冷是拉撒路事件的一個後續的發展,.
主持人:並且都連接著耶穌的受苦和死亡,因為約翰刻意記載,.
主持人:是因為耶穌引起了一個哄動,就加強了祭司長要殺害他的動機..
主持人:所以復活事件是因,永耀進城是果,永耀進城是因,.

$^{41}$主持人:祭司長要殺害耶穌是果..
主持人:我們再看看耶穌進城的經歷,.
主持人:我列了幾段婦女福音的記載,大家看看比較一下..
主持人:馬可的記載是很簡單的,他們把爐區牽到耶穌那裡,把自己的衣服搭在上面,.
主持人:耶穌就騎上,有許多人把衣服鋪在路上,有人把田間的樹枝撼下來,鋪在路上,.
主持人:前行後鎚的人都喊著說:和珊,奉主命來的是應當稱頌的,.
主持人:那將要來的我祖大位極國,是應當稱頌的,高高在上,和珊..
主持人:馬太的記述和馬可是相類似的,牽了爐和爐區來,把自己的衣服搭在上面,.
主持人:耶穌就騎上,眾人多半把衣服鋪在路上,或者是撼下樹枝來,鋪在路上,.
主持人:前行後鎚的眾人喊著說:和珊是歸於大位的子孫,奉主命來的是應當稱頌的,高高在上,和珊..
主持人:路加就有這樣的記述:他們牽到耶穌那裡,把自己的衣服搭在上面,扶著耶穌騎上,走的時候,.
主持人:眾人把衣服鋪在路上,將近耶路撒冷,正下橄欖山的時候,眾門徒因所見過一切的異能,.
主持人:就歡呼起來,大聲讚美上帝說:奉主命來的王是應當稱頌的,在天上有和平,在至高之處有榮光..
主持人:我們看完這三段符類福音,我們再比較一下約翰的記述,第十二十三節:.
主持人:第二天有許多上來過節的人聽見耶穌將到耶路撒冷之後拿著鐘樹枝出去迎接他,.
主持人:喊著說:和珊吶,奉主命來的以色列王是應當稱頌的..
主持人:簡單說,和珊吶就是拯救我們的意思,不過在這段聖經,他的意境應該不是一個祈求,而是一個讚美,.
主持人:就是讚美這個拯救者..
主持人:好,對照三卷的符類福音,你看見約翰的記述是最簡約的,並且漏了許多我們看來都算是頗重要的細節,.
主持人:除了他沒有提到耶穌在進城之前先差遣兩個門徒去為他納一隻驢驢回來,.
主持人:你也要問他那隻驢驢是怎麼來的,三卷符類福音是有記載的,.
主持人:就是怎麼差遣兩個門徒去拿一隻驢驢回來,約翰也沒有記載群眾是將衣服鋪在驢驢身上,.
主持人:鋪在路上作為紅地毯去歡迎耶穌,都沒有記載這件事,所以記載得很簡單,.
主持人:就算是群眾對耶穌的讚頌,吶喊也是最簡單的記述,.
主持人:約翰卻是將耶穌進入聖城的事件連上拉薩路的復活事件,.
主持人:這個佈局是為耶穌受耶路撒冷群眾夾道歡迎的待遇提供了一個很合理的原因,.
主持人:因為如果不是因為拉薩路復活事件所引起的哄動,群眾是沒有理由這麼熱情去款待他的,.
主持人:耶穌又不是第一次來耶路撒冷,耶穌一年來幾次的,.
主持人:就是首節,以前沒有聽過這麼多人去歡迎他,耶穌沒有這麼多粉絲,.
主持人:為什麼突然間這次這麼多粉絲呢?.
主持人:很明顯,一定是和拉薩路復活事件引起的關注和哄動有關係,.
主持人:所以三權俘虜福音沒有將這件事拉上拉薩路的復活事件來說,.
主持人:所以就解釋不了為什麼突然間這次耶穌會這麼受歡迎,.
主持人:但是約翰就將那個原因提了出來,.
主持人:四卷福音書都有提到耶穌騎驢進入耶路撒冷,.
主持人:不同三卷俘虜福音,約翰只是簡單說耶穌得了一個驢驢就騎上,.
主持人:沒有交代這隻驢驢是怎麼獲得的,.
主持人:俘虜福音詳細記述,.
主持人:耶穌是刻意打發兩個門徒前往一條村子,.
主持人:將一條山在柱子上的驢驢子牽來,.

$^{81}$主持人:有人還來查問門徒為什麼要拿走那隻驢驢,.
主持人:耶穌就預先吩咐門徒作答,主要用它,.
主持人:記不記得這幾個字,很多人將它說了很多的解說,.
主持人:馬太宗特別記述,門徒拉過來的不只是一頭驢驢,.
主持人:而是驢驢連同他媽媽,一隻大驢加一隻小驢驢,.
主持人:所以是拉來兩隻的生畜,.
主持人:路加就強調這隻驢驢是過去從來沒有人騎過的,.
主持人:這些都是《三卷弗雷福音》有說,.
主持人:但是在約翰是完全不提的,.
主持人:約翰沒有提到驢驢從何而得,.
主持人:但是卻沒有漏記到騎驢驢進城的事件,.
主持人:因為這是一個很關鍵的證據,.
主持人:要證明耶穌的王者身份,.
主持人:十四十五節,耶穌得了一個驢驢就騎上,.
主持人:如經上所記的說,石安的民那不要懼怕,.
主持人:你的王騎著驢驢來了,.
主持人:騎驢進城是耶穌的王者身份的證明,.
主持人:這是舊約聖經一早預言了的,.
主持人:群眾稱耶穌為奉主名利的以色列王十三節,.
主持人:這一句原來是記載在詩篇一百一十八篇二十六節,.
主持人:奉耶華命來的應當稱頌的,.
主持人:不過詩篇一百一十八篇沒有提到以色列王這個稱呼,.
主持人:所以約翰是將這一節詩篇連上兩段經文,.
主持人:一段是撒迦利亞書九章九節,.
主持人:石安的民那應當大大喜樂,.
主持人:以色列耶路撒冷的民那應當歡呼,.
主持人:看那你的王來到你這裡,.
主持人:他是公義的,並且施行拯救,.
主持人:謙謙和和的騎著驢,就是騎著驢的驢子..
主持人:第二段是西班牙書三章十五節,.
主持人:耶華已經除去你的刑罰,.
主持人:趕出你的仇敵,以色列的王耶華在你中間,.
主持人:你必不再懼怕災禍..
主持人:兩段經文都有提到王,或者是以色列的王,.
主持人:約翰將騎著驢的驢子和以色列的王將兩個驢子複合起來,.
主持人:變成耶穌進入聖城的景象是舊約的預言,.
主持人:而騎著驢驅就是他王者身份的證據..
主持人:我們要留意,約翰福音除了很強調耶穌是上帝的道,.
主持人:上帝的兒子之外,.
主持人:約翰福音其實是非常強調耶穌的王者身份的,.

$^{121}$主持人:所以如果我們要看約翰福音對耶穌的敘述的話,.
主持人:耶穌是王,這是約翰福音常常放重點去記述的..
主持人:例如耶穌在呼召拿丹葉的時候,.
主持人:拿丹葉就已經認出耶穌的身份,.
主持人:拉比你是上帝的兒子,你是以色列的王,一章四十九節..
主持人:一口群眾認定耶穌是王,又強迫耶穌作王..
主持人:耶穌在彼拉多面前被指控的也是他自認為猶太人的王,.
主持人:這是講了很多很多次的..
主持人:耶穌對這個身份,包括這個身份帶來的罪名,沒有直接肯定,也從來沒有否定過..
主持人:最終他都是含蓄的間接承認了,.
主持人:你是我是王,我為此而生,也為此來到世間,特為給真理作見證,凡屬真理的人就聽我的話..
主持人:另外,耶穌又稱撒旦做這世界的王,並且宣告這個世界的王已經被驅逐,.
主持人:他沒有主宰耶穌的能力,所以耶穌是用王者的身份去驅逐撒旦這個世界的王..
主持人:所以,約翰是很強調耶穌是王,是以色列的王..
主持人:他刻意記載耶穌騎驢驅進城的唯一目的,就是證明他是以色列的王..
主持人:約翰還補充,門徒最初不明白這件事,不知道為什麼耶穌要騎上這條驢..
主持人:不過,在耶穌復活之後,他們就想起舊約的聖經,然後就知道,原來這不是一件偶然的事件,.
主持人:耶穌是刻意用這條驢來騎,要去兌現舊約的預言..
主持人:不過,我們這樣說,除了這是舊約的預言之外,騎驢驅進城和耶穌的王者身份,兩者的關聯,其實是有點勉強的..
主持人:古代的君王都是騎馬的,不會騎驢的,因為高頭進馬和矮小的驢子,其實差距很大..
主持人:驢的形狀很怪異,頭很大,身圓,腳很短,走路也很慢..
主持人:要他們去輔載寵物都還可以,載人就有些有礙觀瞻,根本無法顯出君王的威武..
主持人:撒迦利亞書九章九節預告這個王是謙謙和和的騎驢,就是騎驢的驅子..
撒迦利亞的預告其實是要強調,這個王是與眾不同的,是很謙和的,一個謙謙和和的君王..
和人間所有君王,包括舊約的大衛,所羅門這些王都很不一樣..
主持人:耶穌所騎的不是一隻比較大的驢子,而是尚未成年的驢驅,通常叫驢驅就是指著他兩歲以下..
主持人:驢的特性就是他負重能力很強的,所以雖然一隻驢子都未長成,他仍然能夠背得起耶穌這個人..
主持人:當然耶穌很瘦,我也猜到,所以應該不是很重..
主持人:不過你仍然可以想像一個成年男人騎在一隻小驢驅之上,他兩隻腳幾乎都碰到地,不縮起基本上就會擦地而行..
主持人:所以賣相實在很滑稽,絕對不會給人一個英明神武的形象..
主持人:雖然騎驢驅的賣相有些寒酸,不過對耶穌來說,這已經是少有的VIP待遇..
主持人:跟同時期的平民百姓一樣,耶穌絕大多數的時間,除了這次之外,可能就是他媽媽帶他去埃及..
主持人:我們看圖畫,可能是騎著一隻驢.除了這之外,我也沒有聽過耶穌有溝通工具..
主持人:他每次坐都是11號,基本上都是走路,也很少找生畜代步..
主持人:當然不算坐船,橫過加利利是船來的,是溝通工具..
主持人:即使從加利利去耶路撒冷,單程也要180公里,走路要走三天..
主持人:不過耶穌幾乎都是走路去,沒有試過騎馬,騎驢..
主持人:耶穌在他在世最後一個星期,用一個謙和的方法展示他的王者身份..
主持人:這個既符合舊約的預言,也顯露了他在人間最尊貴的面貌..
主持人:在他三年半的傳道生涯裡,曾經不止一次有人認出他應該是王,不過他也拒絕..

$^{161}$主持人:有人想擁納他做王,他也拒絕,並且躲開..
主持人:每一次他被擁納為王,都是因為他行了一件很厲害的神蹟,通常都是這樣的..
主持人:包括了第一件偉大的神蹟,就是五餅二魚,餵飽五千人以上..
主持人:這是一件極其轟動的神蹟..
主持人:革命就是為了請客吃飯,一個人不用籌糧都可以讓人吃得飽..
主持人:你告訴我,不找他做王還找誰?.
主持人:所以五餅二魚神蹟之後,群眾就擁納他做王,他就躲開..
主持人:第二次群眾要擁納他做王,就是他行神蹟,治好拉薩路之後,就是這裡..
主持人:耶穌沒有刻意表明他自己的王者身份,但是他有時候不小心露出一手,就顯出他非凡的身份和能力..
主持人:群眾認出或者猜到,這是一件不是耶穌預計的事..
主持人:不過這次,我相信耶穌不是被冤枉,而是他選擇自己坐上路區..
主持人:大概他知道他的時間已經到,他不再遮掩他的王者身份,他和敵人對決的時間到..
主持人:他在人間也只剩下五天都不夠的時間,是不是?.
主持人:耶穌降世為人,謙卑虛己,從來沒有以王者身份顯現..
主持人:即使有人猜到他是王,或者期待他成為王,他也會拒絕他們的擁戴..
主持人:但是他生命最後一個星期,就是日後我們叫做終之主日的時候,他以王者的身份登場..
主持人:騎著驢驅,我們覺得榮耀的進入耶路撒冷,接受群眾對他歡呼喝彩..
主持人:喝彩的群眾,有少部分在伯大利目睹拉薩路復活事件的神蹟..
主持人:但大部分是在耶路撒冷聽信目睹者的傳言,抱著臭熱鬧追捧明星的心態..
主持人:他們對耶穌的認識非常有限,談不上真的知道他是王,或者王者的身份是怎樣的一件事..
主持人:他們只是知道耶穌很厲害,很有能力,或者可能是從天上下來的..
主持人:這樣的非凡人物,應該可以帶領他們去推翻羅馬政權,復興以色列國..
主持人:這明顯不是耶穌的降世計劃,所以這群群眾其實是不了解耶穌的真正身份..
主持人:他們就算看到耶穌是王,也不知道耶穌是王是什麼一件事..
主持人:因著誤會而結合,因著了解而分開..
主持人:我們知道這些歡迎耶穌的群眾,很大部分在幾天之後,就在彼拉多的虎牙面前高喊要釘死耶穌..
主持人:所以這一刻他們對耶穌的熱情不是太持久,他們對耶穌的擁戴是過眼雲煙..
主持人:很多解說這段經文的人,對夾道歡迎耶穌的群眾都有比較負面的評價..
主持人:但是在這篇簡單講述的應用部分,我也想用一個比較正面的角度來認同他們此刻的想法和做法..
主持人:不管那是一個偶然的誤會,還是一個心思熟慮的反應也好..
主持人:並且我希望我們代入他們,設想成為他們之間的一份子..
主持人:想想當日的場景,耶穌基督進入聖城耶路撒冷,騎著一頭驢驅,兩旁圍上一大堆的群眾..
主持人:他們對耶穌大聲歡呼,高喊奉主名來的王是應當稱頌的..
主持人:然後他們將自己的衣服鋪在地上,讓耶穌的坐騎踩過,以表示對他的尊敬和歡迎..
主持人:你可以看到一個相當詭異的圖畫,一個沒有階型美容,一個在窮鄉僻壤出生成長,.
主持人:在他人生絕大部分的日子都是在相當偏僻,遠離政治中心的加利利湖邊,.
主持人:和一大堆深水Po ,石硤尾的小百姓相處的人,吃一下腸粉,吃一下豬肝麵..
主持人:這樣的人,竟然我們要認出他是王,他是特首,他是普世的王..
主持人:你覺不覺得這是一件很古怪的事?.
主持人:如果不是他釋放出來的神蹟,基本上左看右看,前看後看,你都看不出他和王有任何的關聯..

$^{201}$主持人:並且他這次進入耶路撒冷,其實也沒有什麼很了不起的,沒有華美衣冠,沒有旌旗飄揚,.
主持人:也沒有護從,沒有侍衛,不是騎著一頭高頭進馬,昂首前進,.
主持人:而是不知何時偷了一隻小驢驢來,即是一下一下地走..
主持人:並且這隻驢驢還是臨時張羅回來的..
主持人:你告訴我,這個是王,如果我在群眾中間,我靠近他,看什麼看什麼,看王呀看王呀,.
主持人:一指指著這樣的這個,我都覺得有點滑稽,望之不似人棍..
主持人:但是,如果我們不斤斤計較接待的規格,.
主持人:我們都可以說,耶穌在這裡所接受的,就是君王式的歡呼儀式..
主持人:群眾夾道歡迎,大聲讚美說,奉主名來的王是應當稱頌的,.
主持人:並且主動將他們身上的衣服脫下來鋪上,在地上被耶穌的座騎踩過..
主持人:一切是自願的,是自然的,沒有人動員他們,勉強他們..
主持人:他們這樣做也不會有任何報酬,他們是自發的,.
主持人:心甘情願的去到城門口,去奉獻自己的衣裳出來..
主持人:你不要小看,對小老百姓來說,衣服是很貴的,.
主持人:所以你記得那些死囚脫下來的衣服,要抽籤拿走的,不會浪費的,衣服是貴的東西..
主持人:耶穌在現實生活裡面沒有王者的身份,他們尊奉祂做王..
主持人:祂沒有紅地毯,有群眾用衣服鋪上成為的地毯..
主持人:祂沒有任何御用的儀仗隊,但是有群眾烏合組成的儀仗隊..
主持人:耶穌基督不是君王,群眾卻認出祂是王,.
主持人:並且用一個他們能力所及,接待王的方式去接待祂..
主持人:我不知道你看這個景象,或者你幻想這樣的景象,.
主持人:你心中有什麼樣的感覺?我自己一邊看,一邊很受觸動..
主持人:我們的情況其實都是一樣的..
主持人:說說我們教會,一個舊式商業大廈的單位,.
主持人:我們認出這裡是永生上帝的居所..
主持人:我們每個星期一,我們懷著恭敬的心來敬拜祂..
主持人:彷彿這個樓底九呎的房間,我們更加在這裡遭遇上帝的同在..
主持人:我們身邊的弟兄姊妹只是一群庸碌平凡的人,.
主持人:類似的人在香港最少有七百萬,.
主持人:但是我們認出他們是上帝所愛的兒女..
主持人:他們也是上帝對我們最大的厚切,值得我們恭敬地去接待他們..
主持人:我們所做的事情其實都是極其細碎,卑微的工作..
主持人:在主日學,為一群頑劣的小朋友安排聚會,.
主持人:順受一個上奏,認出我們視之為上帝交託給我們最偉大的使命..
主持人:我們必須要嚴肅從事..
主持人:我剛剛從意大利回來,數了數第五年在意大利的服侍..
主持人:有人問我這五年有什麼感觸..
主持人:我其實也說不出,我一方面認出這是上帝奇妙的作為,.
主持人:因為如果沒有祂的作為,我根本不可能熬過這五年的時間..
主持人:但是你問這五年做了什麼呢?其實一點都不偉大..

$^{241}$主持人:我們剛剛有第二屆的畢業生,六個,就是六個而已..
主持人:第一屆就七個,第二屆六個..
主持人:這兩年,其實應該說做了五年,製造了十三個畢業生..
主持人:這十三個畢業生絕大部分都沒有你們這麼好樣的..
主持人:大概有一半是有學位的,但是其他就沒有的..
主持人:其實學歷是我完全無法預計的..
主持人:一方面我不可以假設他們懂英文,因為意大利英文不是第二語言..
主持人:所以我除了用中文教他們之外,我沒有什麼懂得教..
主持人:而再糟糕的就是,其實他們的中文程度也不是很好..
主持人:就算有大學位的中文程度也不是很好..
主持人:我剛剛在五月安排幫他們考試,我教系統神學要考試..
主持人:本來是Close Book,不能用書的,當然不能用電腦,不能用手機..
主持人:我就給他們一張紙牌,但是他們竟然臨場有人跟我說他們不會寫字..
主持人:不奇怪,因為他們只會用拼音..
主持人:其實你用電腦用多了也不是很會寫字..
主持人:叫你拿一支筆來寫字,他們真的寫不到..
主持人:有很多個說寫不到字,就寫不到字..
主持人:於是沒辦法,只能給他們用電腦,就叫他們自己在上帝面前誓願不要在電腦遊走..
主持人:每一個同學在學院三年的時間,你問我做了什麼?.
主持人:和他們一起,由吃早餐到晚上九點十點,在我的房間九前,我處理過一個弟兄和他拋棄了他很多年的爸爸,說要再見他..
主持人:他一方面說不在乎,另一方面又很緊張,問我可以怎樣處理,去應對他見見爸爸的事..
主持人:我和一對夫婦處理他們,整天都路過大風霜,不停地喊著交,好像大家除了張開喉嚨罵對方之外都不懂得溝通的感覺..
主持人:我處理一個小姐妹,被家人遺棄了很多年,這個不奇怪的..
主持人:在歐洲,父母生了小孩之後,要看護,要做生意,沒有力氣照顧,於是把她扔回大陸給奶奶照顧,照顧了十幾年,和父母基本上沒有任何感情..
主持人:我想說什麼,都是這些瑣瑣碎碎,全部都是沒有任何能說出神聖意義的東西,這是我日常的工作..
主持人:當然我們知道,事奉無大事,你領受上帝的呼召的時候,你要怎樣說,我要為主作見證,我要將呼音傳遍地極..
主持人:你每一天日常的工作,都是和婆婆聊天,送走那個老伯,都是做一些瑣瑣碎碎的事情..
主持人:然後,我們都照鏡自知,我們不做這樣的事,那些偉大的事又怎樣輪到我們做呢?.
主持人:我們說將我們自己獻上,我們照鏡自問,不就是爛命一條,就是一條這樣的破破碎碎的生命..
主持人:然後,我們用信仰的眼睛,我們認定,我眼前這個小弟兄,這個小姐妹,就是耶穌的化身,他是那個化妝的基督..
主持人:King in disguise,一個化妝的君王,我們的主常常微服出巡,常常化妝..
主持人:你記得馬太二十五那個比喻,無論是你給他麵包他吃,沒有給他麵包他吃,給他水他喝,沒有給他水他喝,去探他病,沒有探他病,去探他監,沒有探他監..
主持人:耶穌事後說,你有做,沒有做在這個小子身上,就是有做,沒有做在我身上..
主持人:我們身邊,就是我們左望右望,都只不過是極其普通庸碌的人,竟然可以是耶穌的化身..
主持人:因為你和我,憑信心認出耶穌來,然後我們就專心致意地,我們就將自己埋在對方身上,專注地服侍他,服侍他..
主持人:我們能夠付出的也都很悠閒,身上一件爛衣服,我們只能夠做的,就是張開我們的喉嚨,大聲地吶喊,奉主名來的王..
主持人:我們能夠叫的聲音悠閒,能夠作出的影響悠閒,但是我們就是這樣去實踐我們的信仰,去唱剛才我們唱的詩歌,擁戴我們的主為王..
主持人:我們的主,一個萬王之王,萬主之主,是我們這一班這樣的人,我們在旺角,我們在九龍,我們去尊奉他為王..
主持人:我謝謝剛才服侍的弟兄姊妹,謝謝剛才的師班,他們是用他們最大的可能去向這個世界宣告耶穌基督,值得我們趴在他們面前,專心服侍的各位..
主持人:你和我,是可能出於誤會,最少不是一個周詳的思考,但是我們今天就在這個位置,我們在這個位置,我們去尊奉耶穌基督,和珊吶,高高在上,和珊吶,從天上而來的王,是應當尊崇的..

$^{281}$主持人:耶穌基督是我們應信的主,耶穌基督是我們應信的王,你和我,看看你有什麼,就把衣服脫下來,鋪在地上,宣告,那個騎著驢驅,兩隻腳搖搖晃晃的,是萬王之王,萬主之主..
\newpage



\section{詩篇 50:10-17}
\label{sec:c5rLyKd8AcQ}
\textbf{2026年6月21日父親節主日崇拜直播｜鄭國邦傳道｜告別浮淺,門徒進深:如釋重「父」｜詩篇50\_10-17}
\newline
\newline
連結: \href{https://youtube.com/watch?v=c5rLyKd8AcQ}{\texttt{https://youtube.com/watch?v=c5rLyKd8AcQ}} ~~~~ 語音日期: 2026-06-22
\newline
\newline
\hyperref[sec:wKhk31WV7Os]{\small{< < < PREV SERMON < < <}}
~
\hyperref[sec:index_chronic]{\small{[返順時目]}}
~
\hyperref[sec:i6Mr5bzJFRQ]{\small{> > > NEXT SERMON > > >}}
\newline
\newline
詩篇 50:10-17
\newline
\begin{longtable}{cl}
\hline
\hline
章節 & 經文 (和合本修訂版)\\
\hline
50:10 & \begin{tabularx}{0.7\textwidth}{X} 因為，林中的百獸是我的，千山的牲畜也是我的。 \end{tabularx} \\ \\ \relax
50:11 & \begin{tabularx}{0.7\textwidth}{X} 山中的飛鳥，我都知道，田野的走獸也都屬我。 \end{tabularx} \\ \\ \relax
50:12 & \begin{tabularx}{0.7\textwidth}{X} 「我若是飢餓，不用告訴你，因為世界和其中所充滿的都是我的。 \end{tabularx} \\ \\ \relax
50:13 & \begin{tabularx}{0.7\textwidth}{X} 我豈吃公牛的肉呢？我豈喝公山羊的血呢？ \end{tabularx} \\ \\ \relax
50:14 & \begin{tabularx}{0.7\textwidth}{X} 你們要以感謝為祭獻給神，又要向至高者還你的願， \end{tabularx} \\ \\ \relax
50:15 & \begin{tabularx}{0.7\textwidth}{X} 並要在患難之日求告我，我必搭救你，你也要榮耀我。」 \end{tabularx} \\ \\ \relax
50:16 & \begin{tabularx}{0.7\textwidth}{X} 但神對惡人說：「你怎敢傳講我的律例，口中提到我的約呢？ \end{tabularx} \\ \\ \relax
50:17 & \begin{tabularx}{0.7\textwidth}{X} 其實你恨惡管教，將我的言語拋在腦後。 \end{tabularx} \\ \\
[1ex]
\hline
\hline
\end{longtable}
$^{1}$以下經訓是在詩篇51篇10-17節.
在聖經舊約726-727節.
詩篇51篇10-17節.
請會眾聽我讀出.
神啊!求你為我做清潔的心.
使我裡面重新有淨植的靈.
不要凋棄我,使我離開你的面.
不要從我收回你的聖靈.
求你使我贏得救恩之樂.
賜我樂意的靈扶持我.
我就把你的道指教有過犯的人.
罪人必歸順你.
神啊!你是拯救我的神.
求你救我脫離流人血的罪.
我的舌頭就高聲歌唱你的公義.
主啊!求你使我嘴唇張開.
我的口便全揚讚美你的話.
你本不起來祭物.
若起來我就獻上凡祭.
你也不起願.
神所要的祭就是憂傷的靈.
神啊!憂傷痛悔的心.
你必不慶賀聖經獨筆.
願神賜上讀去懷念這些人.
弟兄姊妹平安.
平安.
今天是父親節.
所以在這裡先祝各位父親.
父親的父親.
媽媽的父親.
各位父親都有這份喜樂的心.
當然了.
今天也很想藉著這個訊息.
在這裡都可以成為大家的祝福.
就好像今年上次珍珍牧師說.
母親節教會沒什麼特別的禮物送給大家.
今天父親節很公平的.
也是沒什麼特別的禮物送給大家.
但是希望這個訊息.
都可以成為大家的祝福之餘.

$^{41}$我今天也有一份非物質的禮物送給大家.
就是答應大家早點講完道.
早點去喝茶或者有活動.
所以希望大家都能夠在裡面.
真的去想一下怎樣可以成為.
身邊一些父親的祝福.
當然偷換一下剛才那段經文的概念.
平時你可能對父母沒什麼心思.
父親節,母親節才送一份小禮物.
其實你也知道.
怎麼會只是喜悅這個小小的禮物.
最重要是那個心.
你對父母,對爸爸的那份心意.
所以當然我知道在座當中.
又不是每一個都是父親.
有姊妹等等.
但有一樣東西可以肯定的.
就是我們每一個人都一定有父親.
當然親和疏還沒有討論.
或者是不是健在.
都可能是另一個故事.
但是說多說少.
我都知道大家都有個爸爸.
當想起這個爸爸的時候.
通常可能大家都是存著一份尊敬的心.
欣賞爸爸可能付出.
特別可能我這一輩.
爸爸那輩都是很忙著去工作.
我爸爸就會經常回內地工作.
去開廠.
所以兩邊走.
所以小時候不太看到爸爸的形象.
但說多說少都會覺得.
爸爸是默默耕耘的.
為工作去很努力.
當然想起可能我們對社會.
對爸爸的印象.
或者有一定的形象.
又或者可能爸爸都是一種.
令你很自豪.

$^{81}$或者是一個很有成就的代表.
不知是不是.
去衡量你都會有時候想.
如果我好像爸爸這樣就好了.
無論弟兄姊妹都可能會有這個想法.
因為爸爸很容易成為一個印象.
就是如果好像他這樣就好了.
我以前聽過一個前輩分享一件事.
就是他說我出糧第一件事.
你知不知道我做什麼.
你想想你自己出糧第一件事做什麼.
他就是托冷氣機回家.
知不知道為什麼托冷氣機回家.
就是因為以前家裡沒有冷氣機.
爸爸又不肯給錢買.
所以出糧第一件事就是買個冷氣機回去.
好像光宗耀祖那樣.
爸爸你做不到的.
我做到了.
我買個冷氣機回來了.
所以可能爸爸的形象.
都是這樣去呈現出來.
一種成就 一種肯定.
但是反過來.
如果你的爸爸是一個很軟弱.
很無能.
很無作為的形象的時候.
你會怎樣呢.
可能如果你爸爸是那種.
這種又軟弱又無能.
可能你覺得他諸千里.
或者都不覺得這個爸爸.
是很令你很有face很有面子.
甚或乎如果.
如果你爸爸是一個人君子.
你爸爸是一個囚犯.
那會怎樣呢.
你怎樣去接納這個爸爸呢.
當然這些軟弱的狀況.
不容易接納的.

$^{121}$但是當然也是一個極端的例子.
但是你想一下.
其實每個爸爸.
都一定有他長處和短處的地方.
他一定有些做得很好的地方.
有一些做得不好的地方.
那究竟今天你怎樣看這個爸爸呢.
因為當我們慢慢理解到的時候.
其實今天你怎樣看.
你和爸爸的關係呢.
其實可能正在影響你人生很多的方向.
通常何時會發現這個狀況呢.
就是當你可能很努力了.
很努力建立你的事業.
你的工作.
很努力地建立你的家庭.
一個可能你很想期望的美滿家庭.
又或者過著你想著.
一些關係可以很順暢很開心的時候.
不知道有沒有去到一個位置.
你總是覺得我怎樣努力呢.
都好像卡住了.
都不滿足的.
都覺得好像差一點點.
究竟怎樣才是最好的呢.
為甚麼會有關係呢.
其實今天當我們去認識.
自己的原生家庭.
原生家庭的意思就是你爸爸媽媽.
一起去成長.
當你建立了一個新的家庭.
那個就是新生家庭.
原生家庭就是說你爸爸媽媽.
和你一起成長的那個.
原來都會直接影響.
今天我們的人際關係.
我們價值觀.
和內在的那種幸不幸福的感覺.
以致在裡面我們知道.
這個和爸爸的關係.

$^{161}$不只是一種情感上的依附.
原來是一個深層.
我們今天的性格塑造的一個歷程.
我不知道有沒有聽過一些小朋友會說.
人家的爸爸就好了 為甚麼我爸爸會這樣.
這些說話.
又或者慢慢一路去成長的時候.
其實爸爸的印象.
印在我們裡面.
去告訴我們甚麼是對甚麼是錯.
再繼續如果我們成為爸爸之後.
其實很大程度.
我們都是在學這個爸爸.
給我們的影響.
這些在之前我去上那些好爸爸的課程.
或者等等 遲些教會有機會.
又可以和大家一起去研究.
走在一起去想.
原來這個爸爸的形象.
都一路影響著我們去建立.
自己成為爸爸的時候.
怎樣去作每一個決定.
為甚麼只提爸爸.
因為媽媽的形象通常是另一方面.
就是講關於愛 自論 或者那種親密.
爸爸好像剛才所說.
通常當我們在工作裡面.
要做一些決定.
一種決策力 或者一種方向.
或者面對挑戰的時候.
當如果你和爸爸的那個關係.
越緊密的時候.
原來回到這些場所裡面的時候.
你那種力量去面對這些挑戰.
是大一點的.
那些研究都會說.
原來和爸爸的這種關係.
都會影響著那個人的自信心.
和那種毅力.
當然 今天當你想起爸爸的時候.

$^{201}$可以有很多不同的片段.
不是強迫大家.
今天你就要支持你爸爸.
而是想你回想.
今天當你看到.
想起爸爸這個形象的時候.
你怎樣看和你怎樣接納.
接納不是那種.
要馬上擺和頭酒.
一會兒的那頓飯就是和頭酒.
不是這個意思.
而是在裡面.
你怎樣去找到.
這個形象對你的影響是什麼呢?.
當你很拒絕這個形象.
在你生活裡面出來的時候.
原來某程度上.
你都在否定自己的那個形象.
那個圓滿在裡面成為了一個缺口.
而致令到你今天.
在生活裡面.
有時候可能有些狀況就是.
為什麼我堅持不了.
這麼容易放棄.
或者等等.
所以今天當我們要去面對這個關係的時候.
真正的那個圓滿.
或者那份滿足.
不是只怪爸爸媽媽.
為什麼會這樣.
或者是要去另一個極端.
就是馬上什麼都不管.
而是帶著一種根深的眼光.
究竟上帝為什麼讓我在這個家庭裡面去成長呢?.
我爸爸有好有不好的地方.
今天在我的生命裡面.
我怎樣去轉化呢?.
剛才說了.
很多家庭輔導研究都告訴我們.
這個原生家庭對我們來說是很大的影響.

$^{241}$但是亦都不是代表著.
總之有什麼就是我爸媽錯.
有什麼事今天做得不好.
就是他們不好.
其實我們仍然可以有轉化.
心理學上可能都會說.
童年決定論.
不知道大家有沒有聽過.
你的童年.
如果用一個貼合一點.
大家有沒有聽過三歲定八十.
你小時候是怎樣呢?.
已經定了你將來是怎樣的.
這個童年決定論.
是心理學上的想法.
但是在我們每一個跟從上帝.
認識耶穌的人來說.
上帝才是我們最大的牧者.
祂過往許可我們在這個家庭.
這個盆栽裡面去成長.
有祂的心意.
當我們再展望前面.
怎樣去建立上帝給我們的形象的時候.
相信上帝會給我們一個.
根深的眼光去轉化.
所以今天為何一開始就問.
你怎樣看你爸爸呢?.
因為同時間反映著一樣東西.
就是你怎樣去接納一個.
又好又不好.
又堅強又軟弱的一個父親.
在我們裡面.
我們當然不能夠改變我們爸爸的形象.
或者改變他.
沒有理由一個軟弱的爸爸.
可以由我們的手將他變成剛強.
因為這個是爸爸他自己的功課.
我們自己的功課呢?.
我們要學的課題.
就是我們怎樣去面對著.

$^{281}$在這個家庭裡面的成長.
所以首先第一件事.
當要接受爸爸都有軟弱的時候.
你知不知道原來我們自己.
都是軟弱的.
我們有期望是正常的.
而期望和現實不符合的時候.
有失望都是正常的.
但是究竟怎樣可以求主.
去更新我們的眼光.
去調節到這件事呢?.
就是希望我們今天可以學到的功課.
事實上我們一開始說的時候是明白的.
每個人都不會是完全的.
這個大家都明白.
但是當放到親人或者身邊的人身上的時候.
總是不容易去接受的.
甚或乎一個再好的人.
我們都知道他一定有不好的地方.
甚至我們很多時候知道.
一些可能表面好像很好先生.
或者很好的一個男人.
哇 誰知道背後原來是這樣的.
通常那些揭秘.
那些黑歷史的醜聞就是這樣.
哇 在人前有一個這樣的狀況.
原來背後是這樣的.
但是其實即使今天.
我們沒有被人揭穿這一面.
我們誠實面對自己的時候.
其實知道自己都有這些的地方.
最大的問題就是.
為什麼在今天我們的文化裡面.
我們不容易去揭露.
或者去展現脆弱的一面給人看.
通常我們都很容易就是說.
男人來的嘛 要面子的嘛.
你不可以這樣去說他的嘛.
為什麼呢?.
為什麼當我們在人前的時候.

$^{321}$不願意去表露出.
我們都有脆弱的一面呢?.
也許可能是一直成長.
我們的文化裡面都覺得.
我們要讓家人自豪的.
要做一些東西.
好像剛才說的.
買個冷氣機回去.
或者做一些光宗耀祖.
我們才覺得是最好的.
又或者做社會裡面.
希望你做到的人.
或者做一些應該做的事.
可能某程度上你都會想.
不要令人擔心.
我自己的問題就自己處理.
我自己的困難或者困惑.
不要成為了別人的負擔.
本身這樣想是好的.
但慢慢就會扭曲了.
你都不知道.
又或者將那些脆弱或者軟弱的事.
一些錯的想法.
都藏在一些隱密的地方.
沒有人去幫助你.
試想一下一個經歷.
我不知道你小時候.
有沒有試過張開喉嚨.
很慘地去叫爸爸媽媽.
你幫我吧.
我很想要這件事.
或者很想做到這件事.
你想起這個片段的時候.
那時候你是否很脆弱.
很無助.
好像世界只有這個選擇.
我要這件事.
為什麼說這件事呢.
上年年尾的時候.
我經歷了一段時間很低落.

$^{361}$失眠了一個月.
在晚上的時候.
我經常出現這個圖像.
這個圖像是什麼呢.
就是張開喉嚨在爸爸媽媽的房間門口敲門.
讓我進去和你睡吧.
那個時候我的印象應該是三,四歲左右.
讓我進去和你睡吧.
我很怕黑.
很怕.
那個時候我很軟弱.
很無助.
很脆弱.
彷彿整個世界拋棄了自己在門外.
但爸爸媽媽又不讓我進去.
當然我跳出來想.
就知道爸爸媽媽當然想我好.
學習獨立.
可能你們都經歷過.
不可以這樣依附著的.
自己要成長.
自己要學習.
所以在那個黑漆漆的晚上的空間裡.
我就在外面哀哭切齒.
經歷那種悲傷.
當你經歷過這個狀況.
你知道什麼是真正脆弱.
當然你的經歷未必是這樣.
但你回想起來.
最無助,哀哭的時候.
那個就是最脆弱的你.
通常在這種狀況下.
哭著哭著.
那個脆弱你就會想收起來.
如果按照我剛才的經歷.
開始慢慢你就會變陣.
得不到你想要的東西.
我展現我的脆弱出來.
不可以改變的時候.
開始就會發脾氣.

$^{401}$你都不愛我.
你都不疼我.
接著就可能又會轉變.
就會變成憤怒.
不單止發脾氣.
很生氣.
很生氣爸爸媽媽.
為什麼會這樣對我.
又或者.
可能你不用這種呈現.
你就會寄存了這份脆弱的痛苦.
往後.
你不會再想讓自己.
落入這種脆弱不堪的狀況.
那會怎樣呢.
開始你就會麻木.
開始你就會收起你的情緒.
否認你有這些脆弱.
其實本身.
在你的腦裡面是在幫你.
因為在這些脆弱裡面的時候.
通常很哀哭.
接著腦裡就會幫你去填補.
讓你不要跌入這些悲哀的裡面.
最後慢慢長大了.
會變成就是沒什麼.
我不怕.
我沒事會怕的.
我不會有脆弱的一面.
展現這個假裝沒事的樣式.
去遮蓋這些脆弱.
或者在裡面你不會再相信.
有人會接受到自己脆弱的這一面.
本身如果大家信耶穌的時候.
是一個轉機.
當信耶穌的時候.
可以將這些脆弱.
這些過犯.
去到神面前去展現.
但可惜的是什麼呢.

$^{441}$我們很容易忽略了.
要面對自己這些生命.
那份洞察力.
究竟這些軟弱.
我是不是真的完全呈現在上帝面前呢.
甚或乎可能在教會裡面.
我們很容易將.
你還有這些軟弱.
還有這些脆弱.
還有這些脾氣.
是不是你不夠信心.
你不夠在裡面.
不夠屬靈.
或者等等這些狀況.
令你又再一次.
藏起了自己的軟弱.
或者這些困難.
甚或乎可能在普遍的一些文化裡面.
就是推崇這些權力.
影響力成功的感覺.
擺在信仰裡面.
以致令到即使信了耶穌.
回到我們這個群體裡面.
都一樣會迴避著軟弱.
迴避著受傷害.
這些狀況.
甚或乎當你成為了一些.
很重要的位份的時候.
更加難去接受到自己.
或者其他人有軟弱.
受傷的一面.
情況回到我們剛才所說.
當你怎樣接受.
爸爸這個形象.
有軟弱,有堅強的地方的時候.
原來回到教會.
你對領袖的期望.
都會投射在裡面.
當你面對著原來領袖.
都會有軟弱的時候.

$^{481}$嘩!我接受不到.
嘩!原來是這樣的.
原來將這件事都會投射在當中.
當大家回到教會裡面的時候.
可能你都會在裡面.
分享很多信仰上的經歷.
很多感恩的事.
但你再想想.
你回到群體裡面.
在團契裡面.
對上一次.
去分享著那些傷心失意.
很困惑的事是何時?.
你是會收起那些傷心的事.
困惑的事.
還是都會躺開呢?.
你介不介意去請求別人去幫助你.
還是你會很少想人幫你.
總是覺得我在這裡要支持人.
要去愛人,去幫助人.
當遇上一些很難解釋在你生活上的事的時候.
你會不會馬上告訴別人.
我真的接受不了.
我需要屈膝在神面前去求.
神明顯顯明祂的旨意在我生命裡.
還是你會繼續假裝沒事.
信靠神.
沒事的.
繼續用這種堅強或者挺胸的方向.
在危機裡面仍然展露著那種.
沒事的.
總之有這份堅定就可以了.
好讓身邊的人見到你的信心.
或者不用為你擔心.
前者好像剛才聽到很軟弱.
又分享這些狀況等等.
後者好像很堅強.
但其實這兩個形象.
當你放在耶穌的身上.
其實耶穌在赫西瑪利園.

$^{521}$都是呈現著這份軟弱.
在很多的困惑裡面.
祂和祂的門徒都是分享著.
我現在很困惑.
我們一起去禱告.
甚或乎當不明白神的心意的時候.
去屈膝在石頭裡去禱告.
求神去顯明.
相比今天.
當我們的社會好像要很強.
甚至我們那個年代要聽強人事理.
能飛天遁地.
那種形象.
哪一種才是讓我們真正.
可以在神面前去求主更生呢?.
甚或乎當看到今天的經文裡面.
在舊約裡面.
這個一家之主.
這個重要的群體領袖.
一國之君大衛.
他都去學習去說出自己的失敗和掙扎.
當然你看回聖經.
大衛經常說是最合神心意的人.
但不是說他是一個完美的人.
他都會有做錯或做得不好的地方.
他不是一個很好的先生.
在他那次嚴重的失德裡面.
他就是和那個婦人.
拔事八通姦.
還設局害死了他老公烏尼亞.
當先知將這件事在他前面呈現的時候.
他可以有很多的解說.
又或者像那些男人那樣.
要面子 要死撐.
有很多的原因去解釋這件做錯的事.
但你看到大衛.
他沒有否認這件事.
當先知拿丹來到他面前的時候.
他懊悔 他馬上懺悔.
而這個詩篇51篇就是他的懺悔詩.

$^{561}$就是他告訴當時的人知道.
或者之後的人去看的時候.
他因著這個錯去面對自己這個軟弱.
事實上好像剛才說.
大衛那時候的身份已經是一國之君.
已經掌握了所有大權.
出去打仗都可以有兵兵出去打仗.
他其實在那一刻.
他做什麼都沒有人可以管到他.
因為他已經是最大的那個.
已經是君王.
他有能力去遮蓋一切的指控.
或者否定所有的事.
又或者可能只是簡單說一句.
感情缺失.
就遮蓋了他做錯的事.
但是大衛他沒有這樣.
他同樣去說出這些失敗和掙扎.
很重要.
因為他明白他都不是一個完美的人.
他要訴說出這些狀況.
讓人知道他都不是完美.
更加令到整個群體都知道.
原來上帝不是只是要一個完美的人.
去成就他的心意.
當你看回詩篇51篇第一至九節的時候.
其實你會看到大衛很明白一個狀況.
你如果看回51篇一至九節裡面.
不斷重複的就是在說罪和過犯這件事.
這個就是大衛去知道人的本相.
亦都是他自己的狀況.
在這些罪和過犯去營繞他裡面的時候.
他知道自己其實就是這麼軟弱.
當上帝在《創世記》裡面.
讓亞當夏娃他們在犯罪過後.
啟動了這個墮落的後果之後.
原來就是為了讓我們知道.
我們是軟弱的.
我們要尋求他.
並且意識到他才是我們真正的幫助.

$^{601}$真正的救主.
所以一至九節剛才說了.
除了說他罪和過犯最多之外.
另一個出現得最多的字就是「求你」.
求你去幫助我.
求你在裡面按你的慈愛連率我.
將我的罪孽去洗淨.
而這個求你不單止在一至九節.
他一直延續到今天的那段經文.
十節到十七節裡面.
亦都有很多個求你去說明.
究竟在這種狀況之下.
人的出路是怎麼算呢?.
第十節裡面說:神啊.
求你為我做清潔的心.
使我裡面重新有靜寂的靈.
原來當人的軟弱裡面的時候.
當你真正懊悔.
真正知道自己軟弱.
沒有出路的時候.
你不單止只是會求神.
或者求對方去赦免你的罪這麼簡單.
如果你真是能夠面對自己的時候.
你在裡面會求神.
讓你重新可以建立一個清潔的心.
我兒子這陣子經常說一番說話.
是同學抄回來的.
他未知內容的.
他說有一個同學跟他說.
對不起是沒有用的.
他這樣教導我兒子.
我又覺得對.
如果你做錯了一件事.
只會說對不起.
但你不改.
其實那個對不起是沒有用的.
而事實上在裡面都在說這件事.
當我們要求神去幫助我們這個軟弱的時候.
不只是處理了那件事.
而是你知不知道.

$^{641}$你自己都會繼續有這個軟弱的心.
繼續在裡面會犯錯.
直到你願意被神去更新改變.
而這個清潔的心.
正直的靈的意思就是那種從內而外的更新.
得到的正直.
對神有著那份單純.
毫無隱藏的心的那種狀況.
第二個繼續下去.
他說什麼呢.
不要丟棄我.
使我離開你的面.
不要從我收回你的聖靈.
就好像剛才所說.
你最脆弱的時候是什麼時候.
是不是那種.
你不要丟棄我.
你丟棄我就沒有了.
就好像我剛才所說的那個狀況.
你不要不讓我進去.
我很害怕.
我沒有其他方法了.
大衛就是落到這種最脆弱的狀況.
在神面前就是這樣.
哀求不要丟棄我.
你不要丟棄我.
你丟棄我就沒有了.
這個就是當時大衛展現出來的狀況.
在裡面他說不要收回你的聖靈.
說的就是不要結束了.
那種我們人和神之間的關係.
這種的盟約的關係.
如果神收回這個靈的時候.
就會結束.
也都代表著神和他一刀兩斷.
所以在這種的狀況之下.
大衛苦苦哀求你不要丟棄我.
不要要我離開你的面.
也都不要收回你的聖靈.
當大衛在這種狀況之下.

$^{681}$那麼脆弱.
有什麼出路呢?.
第十二節就說.
求你使我贏得救恩之囊.
賜我樂意的靈扶持我.
原來當你落到這種脆弱和軟弱裡面的時候.
不一定只是優秀.
不一定只是很醜.
將我自己軟弱或者最脆弱的地方展現出來.
可能我們就會有這種害怕.
但是大衛說.
原來當你這樣展露出來的時候.
可以得到一份的歡樂.
就是真正救恩的歡樂.
一種將你唯一的盼望放在上帝那裡.
和將你那些固有以前的面子工程.
形象建立的東西.
可以全部放下.
不需要再靠自己的力去撐著這些面子.
以至他能夠在上帝面前.
得到這份真正的滿足和真正的出路.
而他說什麼呢.
當他得到這份真正的救恩之樂之後.
第十三節就說.
我就會將你的道去指教那些同樣有過犯的人.
罪人一定會歸順於你.
原來這種釋放.
這種將這些負擔放到神面前的時候.
這種生命力的自由.
才能夠成為最大的吸引力.
真正的喜樂就由這裡發生.
當我們回想今天.
當整個的社會世界價值觀.
都會覺得軟弱失敗是一種負擔.
你有這一面很醜.
你收起來不要讓人知道的時候.
原來上帝不是要我們這樣.
當我們在這個世界被唾棄被拒絕的時候.
上帝的眼中卻看為一個時機.
一個時機讓我們真正展現這份脆弱的時候.

$^{721}$明白到神的恩典.
而上帝一直在等我們.
要讓我們知道的訊息.
就是祂接納我們.
將我們能夠走回上帝應該要走的路上.
十三節就說.
祂就會將你的道指教給過犯的人.
其他的罪人一定會歸順你.
二至十四節就說.
神啊,你是拯救我的神.
求你救我脫離流人血的罪.
我的舌頭就高聲歌唱你的公義.
所以當我們在這些脆弱軟弱的時候.
其實也是上帝重建我們.
在這些位置當我們面對上帝的時候.
這個面對上帝不是說.
我們承認軟弱就不需要為我們的過犯去承擔.
但是當我們呈現了這個軟弱的時候.
你要去承擔這一切的事.
就變了不是一個懲罰.
是一個上帝修正我們.
更新我們的歷程.
所以今天無論你是爸爸或者不是爸爸也好.
神要的真的不是我們要成為一個完美的人.
不是要成為一個完美的領袖.
或者在家裡的一家之主.
也不是需要一個完美的獻祭.
神才願立.
在裡面最後說什麼呢?.
第十七節說.
神要的祭就是憂傷的力.
神啊,憂傷痛悔的心.
你必不輕寒.
其實聖經裡面不是這一次說這件事.
神不是要那些凡祭或者獻祭的美物.
正如剛才一開始所說.
難道神我們今天做爸爸媽媽的.
真的很想要子女送給我們.
那些金鏈或者那些領帶.
或者一餐飯嗎?.

$^{761}$其實這些都不是最重要.
而是那個心.
上帝也一樣.
祂不是要我們那些獻祭.
或者所獻上的那些禮物有多完美.
但是在這裡就是說.
我們對著上帝有沒有這一份.
最脆弱憂傷的靈.
呈現在祂面前呢?.
這個就是聖經裡面.
想展現的圖畫告訴你.
在我身上是一樣的.
當我信了主之後.
我都以為上帝要醫治我幫助.
修補我的缺點.
我的脾氣或者軟弱.
我都要靠著神去勝過的.
但是原來我沒想過這些軟弱.
或者這些憂傷或者這些缺點.
其實可能是上帝一條線.
讓我去到祂的根前.
讓我知道神的計劃.
不是只是要我剋星.
或者遮蓋著這些.
而是我要知道.
我自己都是一個軟弱的人.
當我以為信靠上帝.
我就會成為一個很好的先生.
很好的老公.
很好的爸爸.
原來都不是的.
原來我都會有軟弱.
做得不好的地方.
當我每一次.
在面對教養的時候.
我都會有失控.
都會有發脾氣.
可能都會拍桌子等等的時候.
我不能夠不承認.
原來我都真的會有搞不定的時候.

$^{801}$直到我以為那些.
靠著上帝.
家裡會很和諧.
很喜樂.
這些幻象過去之後.
開始很多的自我的指責會出現.
我都是做得不好.
我原來都是很沒用的.
在這些黑暗面.
不斷控訴自己的時候.
我才知道原來我是很沒用的.
原來我真的不是靠著.
我自己一番的意志.
或者努力.
可以改變到一些東西.
去到很多時候出問題的時候.
去到一晚.
我真的受不了.
哭著和我太太說.
我真的受不了.
我為什麼會這樣.
我為什麼會這麼差.
我頭腦上知道我要這樣做好.
但為什麼我真的做不了.
我真的很沒用.
那一刻原來才釋放.
我才明白到.
原來我要承認我自己的軟弱.
多過我要繼續撐著.
我行的.
我可以的.
我可以做一個美好的形象.
原來我們每一個人.
都有這一面.
我們都要放下.
苦苦撐著的這些外表.
承認著我們有這個軟弱.
以致我們可以將這一面.
呈現在上帝面前.
這一份憂傷痛悔的心.

$^{841}$上帝必不輕看.
今天我們想想自己的狀況.
當你面對著自己很多.
有缺失或者不完美的地方.
當你在想的時候.
你是否已經充滿了解釋.
防衛或者戒備呢.
當聽到有人提醒你的時候.
你第一時間就一聽到.
你這樣說有問題.
找他的破綻.
或者找一些想駁回他的地方.
有沒有一些時候.
你總是想找別人的缺點.
或者錯失不完美的地方呢.
當遇到一些事.
你是否會很急著.
要去解決它找改善的方案呢.
在別人面前.
你是否不敢說我不知道.
我不懂得做.
我不知道怎麼辦.
當你願意承認自己的軟弱.
面對著這些狀況.
一定勝過你扮出來的那種美好.
以致我們今天.
同樣面對著.
可能爸爸有軟弱的一面.
領袖有軟弱的一面.
自己有軟弱的一面的時候.
原來這些都是上帝留給我們的機會.
回到將我們這份憂傷痛悔的心.
去呈現在上帝的面前.
所以今天各位爸爸.
OK的.
你覺得不OK是OK的.
各位子女.
你爸爸有不OK的地方.
是OK的.
是可以的.

$^{881}$怎樣去能夠成為這個祝福呢.
就將這些困難和脆弱的一面.
呈現給上帝.
上帝一定不會輕看.
我們一起祈禱吧.
親愛的主.
求你幫助我們.
是我們有時候知道.
這個社會告訴我們.
我們要將最美好.
最剛強的一面去呈現.
但我們忽略了.
上帝不是要這些東西.
也不是要我們將這份完美去呈現.
是你要我們像大衛一樣.
在你面前.
懊悔,憂傷.
以致承認.
我們都只不過是人.
有軟弱的地方.
所以願主你祝福我們每一個.
以致我們能夠在你的面前.
在人的面前都同樣.
有這份信心去承認.
自己的限制.
我們這樣祈禱.
是奉主耶穌基督的名字而求.
阿們.
\newpage



\section{列王記上 19:1-18}
\label{sec:i6Mr5bzJFRQ}
\textbf{2026年6月28日崇拜直播｜鍾家威傳道｜告別浮淺,門徒進深:在低谷中重新站起來｜列王記上19\_1-18}
\newline
\newline
連結: \href{https://youtube.com/watch?v=i6Mr5bzJFRQ}{\texttt{https://youtube.com/watch?v=i6Mr5bzJFRQ}} ~~~~ 語音日期: 2026-06-28
\newline
\newline
\hyperref[sec:c5rLyKd8AcQ]{\small{< < < PREV SERMON < < <}}
~
\hyperref[sec:index_chronic]{\small{[返順時目]}}
~
\hyperref[sec:code]{\small{> > > NEXT SERMON > > >}}
\newline
\newline
列王記上 19:1-18
\newline
\begin{longtable}{cl}
\hline
\hline
章節 & 經文 (和合本修訂版)\\
\hline
19:1 & \begin{tabularx}{0.7\textwidth}{X} 亞哈把以利亞一切所做的和他用刀殺眾先知的事都告訴耶洗別。 \end{tabularx} \\ \\ \relax
19:2 & \begin{tabularx}{0.7\textwidth}{X} 耶洗別就派使者到以利亞那裡，說：「明日約這時候，我若不使你的性命像那些人的性命一樣，願神明重重懲罰我。」 \end{tabularx} \\ \\ \relax
19:3 & \begin{tabularx}{0.7\textwidth}{X} 以利亞害怕，就起來逃命，到了猶大的別是巴，把僕人留在那裡。 \end{tabularx} \\ \\ \relax
19:4 & \begin{tabularx}{0.7\textwidth}{X} 他自己在曠野走了一日的路程，來到一棵羅騰樹下，就坐在那裡求死，說：「耶和華啊，現在夠了！求你取我的性命吧，因為我不比我的祖先好。」 \end{tabularx} \\ \\ \relax
19:5 & \begin{tabularx}{0.7\textwidth}{X} 他躺在羅騰樹下睡著了。看哪，有一個天使拍他，對他說：「起來吃吧！」 \end{tabularx} \\ \\ \relax
19:6 & \begin{tabularx}{0.7\textwidth}{X} 他觀看，看哪，頭旁有燒熱的石頭烤的餅和一壺水，他就吃了喝了，又再躺下。 \end{tabularx} \\ \\ \relax
19:7 & \begin{tabularx}{0.7\textwidth}{X} 耶和華的使者回來，第二次拍他，說：「起來吃吧！因為你要走的路很遠。」 \end{tabularx} \\ \\ \relax
19:8 & \begin{tabularx}{0.7\textwidth}{X} 他就起來吃了喝了，仗著這飲食的力走了四十晝夜，到了神的山，就是何烈山。 \end{tabularx} \\ \\ \relax
19:9 & \begin{tabularx}{0.7\textwidth}{X} 他在那裡進了一個洞，在洞中過夜。看哪，耶和華的話臨到他，說：「以利亞，你在這裡做甚麼？」 \end{tabularx} \\ \\ \relax
19:10 & \begin{tabularx}{0.7\textwidth}{X} 他說：「我為耶和華－萬軍之神大發熱心，因為以色列人背棄了你的約，毀壞了你的壇，用刀殺了你的先知，只剩下我一人，他們還要追殺我。」 \end{tabularx} \\ \\ \relax
19:11 & \begin{tabularx}{0.7\textwidth}{X} 耶和華說：「你出來站在山上，在耶和華面前。」看哪，耶和華從那裡經過。在耶和華面前有烈風大作，山崩石裂，耶和華卻不在風中；風後有地震，耶和華也不在其中； \end{tabularx} \\ \\ \relax
19:12 & \begin{tabularx}{0.7\textwidth}{X} 地震後有火，耶和華也不在火中；火以後，有輕微細小的聲音。 \end{tabularx} \\ \\ \relax
19:13 & \begin{tabularx}{0.7\textwidth}{X} 以利亞聽見，就用外衣蒙臉，出來站在洞口。聽啊，有聲音向他說：「以利亞，你在這裡做甚麼？」 \end{tabularx} \\ \\ \relax
19:14 & \begin{tabularx}{0.7\textwidth}{X} 他說：「我為耶和華－萬軍之神大發熱心，因為以色列人背棄了你的約，毀壞了你的壇，用刀殺了你的先知，只剩下我一人，他們還要追殺我。」 \end{tabularx} \\ \\ \relax
19:15 & \begin{tabularx}{0.7\textwidth}{X} 耶和華對他說：「去吧，從原路回去，往大馬士革的曠野去。到了那裡，你要膏哈薛作亞蘭王， \end{tabularx} \\ \\ \relax
19:16 & \begin{tabularx}{0.7\textwidth}{X} 又膏寧示的孫子耶戶作以色列王，並膏亞伯‧米何拉人沙法的兒子以利沙作先知接續你。 \end{tabularx} \\ \\ \relax
19:17 & \begin{tabularx}{0.7\textwidth}{X} 將來逃過哈薛之刀的，必被耶戶所殺；逃過耶戶之刀的，必被以利沙所殺。 \end{tabularx} \\ \\ \relax
19:18 & \begin{tabularx}{0.7\textwidth}{X} 但我在以色列中留下七千人，是未曾向巴力屈膝，未曾親吻巴力的。」 \end{tabularx} \\ \\
[1ex]
\hline
\hline
\end{longtable}
$^{1}$各位弟兄姊妹平安.
平安.
很感恩有機會跟大家分享神的話.
很特別的.
今天的講題是.
低谷中重新去站立起來.
一開始想用一個.
大家最近應該都有留意的.
世界盃.
大家有沒有看世界盃?.
大家剛剛有沒有留意呢?.
我們台上剛好有十一位弟兄.
事而復我們.
你連我連獨京.
剛剛十一位弟兄.
十一位弟兄.
事而復我們數一下你連我連獨京.
剛剛十一位.
不知道大家世界盃支持哪一隊呢?.
我支持阿根廷.
因為有美斯.
美斯很厲害.
他四年之前就拿了世界盃冠軍.
今年他現在三十九歲.
兩場球進了五球.
在神射手榜的第一位.
不過不知道大家知不知道.
其實在十年之前.
美斯曾經退出過國家隊.
因為在十年之前他試過連續三年輸了決賽.
他又進入了人生的低谷.
他又不願意去面對國家隊的賽事.
他宣佈永久退出國家隊.
經過年年的失敗.
他在國家隊的賽事裡面.
他進入了低谷.
很多的情緒他不想再面對.
幸好他最後走出來.
以致我們今天仍然可以欣賞他的球技.
關於人生的低谷.

$^{41}$我不知道大家會想起什麼時間呢?.
可能是一次失敗經歷.
或許是一次在職場上的挫折.
可能是一些突然的意外.
身體的危機.
又或者是一段令你很難堪破裂的關係.
人生的低谷.
不知道大家會有什麼反應.
和用什麼情緒去面對呢?.
今天很想和大家分享一個聖經人物.
他在人生裡面低谷的經歷.
他叫做以利亞.
以利亞本身的名字就是.
我的神是耶和華的意思.
代表他和上帝的關係很好.
而且他很有能力.
和很有自信的一個人.
我先簡單交代一下以利亞的背景.
以利亞身處在.
所羅門王死後.
南北國分裂的時代.
他身處在第七任以色列王阿哈的年代.
《列王記上》十六章.
是這樣形容阿哈王的.
阿哈行耶和華眼中看為惡的事.
比他以前的列王更甚.
又娶了一個愛邦的女子耶洗別為妻.
去事奉敬拜巴力.
更加指他所行的耶和華的烈怒.
比以前的諸王更甚.
可以說當時北國以色列的.
屬靈的狀況是歷史的最低點.
正正在這個處境下.
以利亞出現.
他一出場就表現到他和上帝有美好的關係.
透過很多神蹟騎士表現他有多厲害.
在十七章那裡.
他一出場就說.
我指著所事奉的永生耶和華.
以色列的上帝起誓.

$^{81}$這幾年我若不禱告.
必不降露水也不下雨.
他表明他和上帝的關係.
神才是那位的真神.
降雨下露的是耶和華神.
而不是被視為暴風雷電神的巴力.
其後他更加經歷.
好像五餅如魚那種的公認.
至致到譚面不短缺.
平柔不剛確.
更加實現了聖經裡第一個死人復活的神蹟.
讓婦人的孩子死而復活.
而且他更加在加密山上.
用一個人擊殺了四百五十個巴力先知.
可以看到以利亞是一個超級英雄人物.
可能說是一個美斯級數的先知.
但原來這麼偉大的先知.
經歷了這麼多的神蹟.
信心更堅定的屬靈偉人.
他也會有崩潰耗盡人生有低谷的時間.
弟兄姊妹.
今天我們每一個人都會面對一些挫敗.
甚至走入人生的低谷.
很多不同的情緒反應.
而且這些低谷是每個人都會經歷的.
每個人生階段都會有.
無論你可能剛剛出來工作.
還是出來工作好幾年.
在事業上開始有些困境.
又或者你退休.
有很多不同的挫折.
每個人都會有人生的低谷.
情緒崩潰的時間.
今天我們就從以利亞去面對人生的低谷.
他有三種不是太健康的反應.
我們去作出反思.
並且我們看看.
神如何在低谷裡面去尋找我們每一個.
讓我們重新去站立起來.
今天我的鋪排會先講.

$^{121}$人在低谷裡面的一些反應.
然後再講神的回應.
所以會很穿雜去講一些經文.
所以大家更加要去留心.
如果你有程序表的話.
可以看著程序表的大綱.
會比你更加容易去聽訊息.
在開始聽訊息之前.
我們先一起去祈禱.
七位天父多謝你.
讓我們有聽你話語的時間.
將我們每一個人的心交到你的根前.
讓我們在一息間的時間.
專心去聆聽你的話.
求你幫助我們每一個.
讓你的話去觸摸我們的生命.
並且帶出從你而來的安慰提醒和管教.
將孩子交上.
讓孩子能夠講得清楚說得明白.
願一切榮耀都歸給你.
禱告奉主耶穌基督的名字而求.
阿們.
首先我們來看看.
以利亞為什麼會是一個失敗的景況.
剛剛明明讀到十八章.
他一個人擊殺了四百五十個巴黎先知.
為什麼他又會進入人生的低谷呢?.
剛才大家讀的經文.
我想大家有留意到主要有三個人物.
以利亞就是先知.
阿哈就是昏君.
耶洗別就是昏君的外邦皇后.
是敬拜賈神的.
在一至三節我們看到.
以利亞本身擊敗了巴黎先知.
但是身為以色列王的阿哈.
他不單沒有恩賞.
更加回到耶洗別那裡.
他告發他.
他將他擊殺巴黎先知的事.

$^{161}$告發給耶洗別.
然後耶洗別就說一定要去報復殺了以利亞.
以利亞很害怕.
所以他就走了.
第三節裡面.
以利亞見者光景就起來逃命.
其實裡面有兩個意思.
第一個是說以利亞害怕.
指的是他面對死亡的恐懼.
一種人性的軟弱.
內心出現那份抖震.
所以他離開.
第二個就好像經文這裡的直譯.
他說以利亞見者光景.
代表以利亞對現實的光景很失望.
很灰心.
他對阿哈很失望.
明明你已經見證了耶和華的神蹟.
但是你仍然不願意去相信祂.
他對自己也很失望.
他認為自己可以帶領以色列國.
去重返耶和華身邊.
但是他沒有做到.
可能他對神也很失望.
他不明白為什麼神會容許這件事去發生.
他的努力被徹底地否定.
無視.
特別是有一種預期上的落差.
給他很大的挫敗.
無力感.
很沮喪,很焦慮.
記得我小時候也有一次類似的經歷.
由高跌到低.
那一年我很幸運地在自己中學的陸運會那裡.
拿了全場總冠軍.
很開心地到處找人拍照.
好像明星一樣.
但是回到家裡.
父母有一句無心的回應.
讓我跌到低谷.

$^{201}$可能那個年代他們是流行用激將法.
他們說金牌不是真的.
無法保護錢.
這些無心的話.
讓我跌到低谷.
明明很開心地在坐車回來之前.
回到家裡就很沮喪.
不被明白.
被忽視.
我想以利亞有一點就是這裡.
由高跌到低.
在這種狀態下他有什麼行動.
有什麼反應呢?.
今天我看到他們.
在低谷裡面有三種.
值得我們反思的反應.
從而幫助我們如何去面對.
第一.
他盲目地比較.
他陷入了自我的否定.
經文第三至第四節.
大家幫我再讀一次好嗎?.
預備起.
面對著森林的低谷.
以利亞選擇逃跑.
他去到猶大的別士巴.
猶大的別士巴是什麼意思呢?.
就是南國最遠的地方.
其實他本身只是離開了.
耶駒別管轄的地方.
仍然留在北國已經足夠了.
但是他不斷地跑.
超過了自己的國界.
去到南國.
更加去到一個邊境的地方.
表示他很想遠遠地避開那件事.
他不想再面對當下的場景.
更加想放棄自己原有的先知的積分.
之後他更加走到更加遠的曠野裡面.
離開他的僕人.

$^{241}$自己一個在羅藤樹下求死.
以利亞說罷了.
就是夠了.
Enough 的意思.
他所表達的是.
無論他在肉體上的那種疲乏,飢餓.
還是他對當時的現況都說夠了.
他受夠了.
求上帝取他的性命.
他指的是.
特別我們留意他的原因是什麼?.
他說耶駒說.
罷了,求你取我的性命.
因為我不勝於我的列祖.
他覺得他不夠他的列祖好.
那些列祖好可能是他自己心目中.
裡面一些的屬靈位人.
摩西,約書亞.
曾經帶領以色列國.
走過艱難日子的人.
他覺得自己的事奉遠遠比不上他們.
所以他很想放棄自己.
弟兄姊妹那一刻他信錯了教.
不知道大家有沒有聽過.
我們生命裡面.
生命的邪教.
其中一個邪教是什麼?.
就是比教.
這個教是怎麼樣呢?.
這個教的價值取決於我們和人的對比.
而不是取決於我們對自身的認定.
更加不是取決於上帝怎麼看待我們.
以利亞剛剛明明在加密山.
戰勝了四百五十個巴力先知.
剛剛拿完世界盃.
很有能力.
很能經歷上帝的恩典.
但是當他比較的時候.
他就否定了自己的成就.
自己的價值.

$^{281}$更加忘記了上帝的愛.
不知道有時候我們會不會進入這個困境裡面呢?.
我們否定了自己的價值.
忘記了我們是上帝所喜愛的兒女.
彷彿他進入到一個環境就是.
他不斷背後有一把聲音跟他說.
我做得不夠好.
的確適當的比較可以幫我們去調整一些目標.
幫我們有些動力去追求一個更加好的情況.
但是如果盲目地去比較.
很可能會讓我們更加不斷的迷失.
定義不了自己.
甚至乎越來越灰沉.
越來越自卑.
不斷去否定自己.
甚至乎當我們如果整天去想自己跟其他人的差異的時候.
會影響我們自己和其他人的關係.
我們會妒忌.
會猜疑.
會憎恨.
會覺得自己比不上別人.
而有些羞愧.
很討厭自己.
不懂得欣賞自己.
落入不斷去否定.
不斷讓自己挫敗的陰霾裡面.
弟兄姊妹,今天我們要學習去「甩膠」.
我們要放棄那些盲目的比較.
很欣賞李靖.
有沒有聽過他之前拿獎的一句金句.
就是最重要多謝自己.
我們要學習去欣賞自己的好.
去接納自己的價值.
欣賞自己的付出.
我們不需要透過跟別人的比較去肯定自己.
就算如果真的在一些情況裡面你覺得自己不行.
也好.
正如邦傳上星期和我們分享.
不OK也是OK的.
承認我們的軟弱總比裝出來的好更加好.

$^{321}$我們要認定基督徒的價值是由神去定義.
我們不需要盲目去比較.
我們再繼續看下去.
第二個讓我們在失敗低谷裡面更加挫敗的狀況是什麼呢?.
就是放大自我.
陷入孤立自憐的困境裡面.
以利亞在羅騰樹下有一些經歷之後.
他就去到河列山.
在河列山裡面分別有兩次對話.
第一次對話跟上帝有對話之後.
神透過微小的聲音去引導他.
然後再有第二次的對話.
我們在這裡試試對比兩次對話.
看看有什麼發言好不好?.
大家幫我讀9-10節.
我讀13-14節.
預備起!.
13-14節.
以利亞聽見就用愛依蒙上念.
出來站在洞口.
有聲音向他說.
以利亞你在這裡做什麼?.
他說:我為耶和華萬軍之神大發熱心.
因為以色列人背棄了你的約.
毀壞了你的壇.
用刀殺了你的先知.
只剩下我一人.
他們還要尋索我的命.
我們會發現在這兩次對話裡面.
以利亞的回應是怎樣的?.
一模一樣的.
為什麼會這樣呢?.
首先神為什麼會問以利亞你在這裡做什麼呢?.
神在兩次的提問裡都有提及他的名字.
記得我一開始說過.
以利亞這個名字是什麼意思呢?.
就是我的神是耶和華.
彷彿再透過後面.
以利亞對神的回應.
我們可以了解到.

$^{361}$神這樣的問其實是回應他核心的身份.
他先知的角色.
他說:以利亞你是帶領人宣認.
耶和華就是神的先知.
你現在在做什麼呢?.
所以以利亞就針對他先知的身份去作出自辯.
他說:我很熱心.
但是那些以色列人很差.
然後他只剩下自己一個.
還要被人追殺.
他這樣去回應.
然後神就透過微小的聲音去作出引導.
再次神提問.
以利亞又再次用同樣的方式去回應.
為什麼會這樣呢?.
以利亞為什麼無動於衷呢?.
我們發現.
在以利亞的提問裡面.
有一個字重複了很多次.
不知道大家有沒有發現.
是什麼字呢?.
是一個「我」字.
他說:我為你大大付出的是我.
仍然堅持的.
只剩下我.
受害被追殺的也是我.
以利亞他將自己放大.
他將自己陷入了一個.
感情的孤島裡面.
他只是強調他自己.
他看不到其他人.
聽不到上帝的聲音.
他讓自己孤立了.
明明上帝在中間有微小的聲音去引導他.
他聽到火後面有聲音.
因為他自己也懂得用外衣去遮掩.
代表他知道上帝與他同在.
但他仍然用同一番說話去回應.
原來人生在一些低谷的位置裡面.
會有一種狀況.

$^{401}$就是看到自己.
覺得自己做很多事情都沒有被人看到.
人生有很多埋怨.
對人有很多批評.
設了很多框框去定義.
有很多應該的事情要做.
有很多不應該的事情是不可以做的.
很多的自憐.
覺得自己就是受害者.
對於一些可能與自己想法不同的人.
即時去斷絕一些關係.
認為是別人不好.
這種情況我想無論對自己.
還是對與其他人的關係.
都會帶來一定的傷害.
我給大家猜一下這個字.
猜到這是甚麼.
甲骨文來的.
這是甚麼字.
猜到嗎.
神字.
神字我就不懂了.
這個呢.
都是一個「我」字來的.
原來這個「我」字其實像甚麼.
其實像一把類似斧頭的武器.
叫「戈」的一個武器.
其實「我」字原始字形.
其實是一個武器來的.
那時候人們舉戈很振奮士氣在兵隊裡面.
去區別敵我雙方.
慢慢才演變成今天「我」的代名詞.
不過這個甲骨文的「我」.
給我很大的反省.
原來我們不好好處理我們的「我」.
我們很可能有機會去傷害自己.
或者傷害其他人.
今天我們在一些情緒狀況的時候.
會不會都是這樣呢.
第三個不良的反應呢.

$^{441}$就是他過度的思慮.
陷入了一個精神的內耗裡面.
剛剛提及了以利亞他兩次的回應.
其中一樣東西很重要.
他兩次都提及「只剩下我一人」.
不過「只剩下我一人」.
其實不是一個事實來的.
而是他被當時的恐懼壓倒性了.
當時的事實其實有很多先知.
藏起來.
而且他是知道的.
但是他放大了恐懼.
在第十八章.
有一個人叫俄巴底.
他把先知藏起來了.
俄巴底將一百個先知藏了.
每五十人藏在一個洞裡.
拿餅和水供養他們.
其實當時這番說話是俄巴底.
他跟以利亞說的.
以利亞是知道他不是只有自己一個.
不過在一個困境狀態裡面.
他極度的恐懼和焦慮.
甚至乎他可能燒毀了.
他很容易進入了一個過度思慮的黑洞裡面.
不知道大家有沒有看過這部電影《Inside Out》.
中文是什麼?《玩轉老朋友2》.
裡面有一個角色叫焦裕.
他有一部機器製造很多負面的幻覺出來.
他不斷在腦裡幻想很多的一些場景.
一些最壞的結果.
可能去到派對.
不受歡迎他就會永遠被孤立.
可能在比賽裡面他射不進那球.
他就會逐出球隊.
有時候我們都可能會這樣的狀況.
覺得做錯一些事就永不翻身.
糟了我做錯了某些事.
以後沒得做回好的了.
糟了我剛剛跟別人打招呼.

$^{481}$沒有人回應我.
是不是那些人生氣了?.
糟了可能我帶詩歌的時候.
有些失誤.
會不會那些人覺得我很沒用,很廢?.
糟了我最近很累.
是不是患了一些絕症?.
昨天我太太才有一個這樣的狀況.
她叫我以後不要跑步.
為什麼呢?.
因為她最近看到一些新聞.
一些不幸的新聞.
就是中暑之後跑完步過身.
她說你以後不要跑步了.
跑步會死的.
我就慢慢地跟她解釋.
其實這些不是絕對的.
可能多喝點水.
或者我會經常留意自己的身體狀況.
不要太熱去跑.
就解決了一些問題.
的確有時候我們為了將來有不同的打算.
有很多不同的預計是好的.
可以幫我們預測風險.
做好準備去做一些事.
但是如果形成一種很消極,很負面的狀態能量.
就會給我們很多精神內耗.
我們要留意.
我不知道你是不是一個很多小劇場的人.
做一些事情就會想很多的畫面.
很多一些狀況出來的人.
或許我們人生過去有不同的經歷.
都會讓我們作出很多的計謀.
計劃,打算.
很多的小劇場出現.
可能我們原生家庭.
可能我們過去一次的挫敗.
都會引致很多不同的小劇場出現.
但是記得珊珊姑娘要跟我們分享.
不要讓我們過去去否定我們自己.

$^{521}$反而我們要從實.
神帶領我們向前走那份恩典.
等一會我們看到.
以利亞在低谷裡面那份佛力.
不單單是因為他對耶駛別那種恐懼.
更加陷入了這三個內心的黑洞裡面.
他越比較就越自卑.
他越看自己就越孤單.
他越想得多就越絕望.
可能今天我們自己.
或者身邊的朋友.
都會有機會在這種狀態裡面.
但是弟兄姊妹我想跟大家說.
神一定會和我們一起去面對人生的低谷.
他會用不同的方式去帶領.
去參與在我們的生命裡面.
我們可以看到.
神用不同的方式去參與在我們的生命裡面.
第一.
神是一個體恤人軟弱.
願意被人及時供應的神.
大家幫我讀五到第七節.
在羅藤樹下的經歷.
預備起.
第八節我再讀一次.
「他就起來吃了喝了.
將著者飲食的力.
走了四十晝夜.
到了上帝的山就是荷列山」.
記得以利亞去求死的時候.
在曠野.
他身心都疲乏.
神就派天使去給他有實際和及時的供應.
給他餅和水.
其實兩次的提問,兩次的供應.
以利亞的反應都有不同.
第一次不知道大家有沒有留意.
第一次他沒有特別提及他起來.
在第八節這裡沒有顯示.
第八節是說他起來吃了喝了.

$^{561}$但是第六節他就說他吃了喝了.
仍然躺下.
第一次他不知道有沒有起來.
有些學者認為他繼續躺下去吃和喝.
所反映的就是.
他沒有完全去遵行天使給他的命令.
他依然顧我.
他仍然繼續想放棄自己.
但是在這個狀態裡面.
神沒有放棄他.
他再一次去供應,去鼓勵他.
再次叫他起來.
並且加多了一句.
第七節他說.
「因為你當走的路甚遠」.
這句話是什麼意思呢?.
「甚遠」這字其實和.
以利亞剛剛說的「罷了」.
即是「夠了」那一字是相對應的.
是神正在回應他的遭遇.
早前以利亞說「夠了」.
在肉體上.
無論是指肉體上的疲乏.
還是對現況的失望.
他都表示夠了.
所以他求神去取他的性命.
神就在這裡回應他的遭遇.
NIV的譯本.
用「too much for you」去解釋這個字.
幫助我們更加容易理解.
神想表達的意思.
神在這裡說.
以利亞你所走過的路.
你所過去的經歷.
實在都令你受救了.
你所遭遇的.
我都看見.
我都明白.
我都理解.
神在這裡表達.

$^{601}$他體恤他的狀況.
他同情他的遭遇.
並且作出及時的回應.
所以以利亞帶來一定的安慰.
第二次的時候.
他就起落.
去吃,去喝.
並且走向神的山.
弟子姊妹.
今天神都呼喚我們每一個的名字.
跟我們說.
too much for you.
今天可能我們正在經歷不同的.
困擾,低谷.
可能在家庭裡.
可能在職場裡.
可能在人際關係裡.
一些不同的狀況都令我們很辛苦.
但神跟我們說.
too much for you.
他很了解.
很體恤我們的需要.
他會作出及時的回應.
羅騰樹.
羅騰樹下是以利亞在人生裡.
那個低谷裡的一個狀況.
羅騰樹.
剛剛大家讀經的時候.
有一個字眼說是一個小樹.
是一個很矮的樹.
樹葉很稀疏的.
在曠野這個沙漠的狀況裡.
是唯一能夠供應很小的遮陰的地方.
所象徵的就是那種絕望和荒涼.
但羅騰樹下除了是他求死的記號.
也是上帝恩典臨到的記號.
上帝就在當中體恤他的軟弱.
給他餅和水.
讓他重新得力走人生的路.
分享一個我在羅騰樹下的記號.

$^{641}$不知道大家有沒有留意.
認識我的都知道.
我很喜歡跑馬拉松.
我每年都會跑馬拉松.
馬拉松除了是我在運動上的一個興趣.
其實在我屬靈生命裡.
也是一個記號.
我人生跑第一隻馬拉松.
正正是我人生一個很低谷的狀態.
當時我媽媽有一個重病.
都不知道會不會有生命的危險.
我在職場和在事奉的路上.
都在困擾我夠不夠緊.
應該去全時間事奉上帝呢?.
而且當年我本身不是那些長跑運動員.
我沒有練習去跑馬拉松.
所以帶著很多的焦慮恐懼.
去跑馬拉松.
早上起床還在想去不去.
但最後我都去.
去到的時候.
上帝給我看到很多的恩典.
就是當我放下選手包的時候.
突然之間後面有一個人去拍我.
是我十幾年沒見的大學師弟.
不是很熟悉的.
他說你又跑馬拉松.
不如我們一起跑.
這個師弟其實是長跑高手.
他可以跑很多的時間.
就是很好的成績.
就完成這個賽事.
但是在那場比賽裡面.
他就陪我一路走著跑著.
結果大概五六個小時.
去完成我人生第一隻馬拉松.
在那個過程很痛苦.
但是我在裡面很深經歷.
上帝告訴我.
我人生很多的困境.

$^{681}$但上帝會為我預備.
人在我身邊及時作出陪伴.
這個馬拉松的記號.
也是我其中一個.
我全時間願意回應上帝呼召的記號.
我看到在人生裡面.
神會有公認陪伴的上帝.
第二我們看到神.
還會透過什麼方法去陪伴我們呢?.
神會調整我們人生的焦點.
用微小聲音作引導的上帝.
我們將鏡頭由駱騰樹下.
轉回在荷列山上.
那兩次對話中間.
剛剛大家沒有讀到的.
就是神微小聲音.
十一至十二節我讀給大家聽.
「而我話說你出來站在山上.
在我面前.
那時而我話從那裡經過.
在面前有裂風大作.
崩山碎石.
而我卻不在風中.
風後地震.
而我卻不在其中.
地震後有火.
而我話也不在火中.
火後有微小的聲音」.
這裡有很多的自然現象.
對於當時的人來說.
有裂風,有地震,有火.
其實是象徵神同在的一個表徵.
大經文卻三次表示.
「而我話不在其中」.
微小的聲音.
其實有溫柔,靜止的意思.
剛剛那些超自然現象.
是一些很轟烈,很震撼的現象.
對比的是溫柔,靜止的聲音.
想帶出什麼呢?.

$^{721}$神想告訴以利亞.
神的作為不一定是轟轟烈烈.
那些震撼的事裡面去成就.
有時神反而透過一些微小的聲音.
微小的事情去做祂的事.
去跟人說話.
以利亞用自己的想法.
過去的一些經歷去理解神.
覺得要做很多的事.
要摩西過紅海.
要做很多很震撼的事情.
才可以拯救國家.
他認為這樣才彰顯神的作為.
這樣才能告訴別人耶和華是真神.
所以以色列人仍然不信.
對他帶來很大的失望.
但是神提醒他.
神有時不是用人期待的方式去呈現出來.
反而這裡所強調的是.
神的話,神的聲音才是最重要.
神會常常提醒我們.
調整我們生命的焦點.
神會透過微小的聲音去引導我們.
我們常常期待神在我們生命裡面有什麼作為.
是不是那些即時的轉變.
一些很震撼的見證才能夠去見證到神呢?.
其實不一定是.
神會用微小的聲音在我們生命裡面作出引導.
神的話今天在我們的生命裡面.
我們有沒有敏感去留意呢?.
Toby傳道曾經提及過我們.
要推到曠野去遇見神.
我們有多久.
我們有沒有留意.
我們要靜下來去聽神在我們生命裡面的帶領呢?.
可能當我們去安靜的時候.
神就有帶領幫助我們走過人生裡面的低谷.
最後我們更加看到神很回應.
以利亞那種孤單感.
當以利亞上上提及.

$^{761}$他自己一個人去面對不同的事情的時候.
神作出親自的拯救和預備.
十五至十八節.
神為他預備了不同的人去走這段路.
裡面提及哈舍,耶護,以利沙.
其實這三個人都是後來幫他對抗巴力的很重要的人物.
而他們分別有不同的角色.
他們分別是愛邦人的王.
以色列的王.
和接替他使命的先知.
十七節特別說那些敬拜巴力的人.
將來無一可以脫離神的審判.
他們沒有人可以走過這三個人的手裡面.
以利亞上上強調自己只有一個人.
神用一個數字去回應他.
第十八節可能很小.
大家幫我一起讀.
一起讀,預備起.
「但我在以色列人中留下七千人」.
「是未曾向巴力屈膝的」.
「未曾與巴力親水的」.
七千人的準備.
這七千這個數字未必是一個實際的數字.
但更加可能是象徵完全.
象徵神會主動.
留下是主動.
很主動去為人預備.
足夠完整的群體.
去與以利亞經歷對抗巴力.
弟兄姊妹神是會包底的神.
神會為我們在生命裡面預備不同的人.
不同的崗位,不同的方式.
去成就神的美事.
有時當我們的焦點.
放在很多不同的事.
和人的罪性裡面.
我們會有很多的失望,很多的焦慮.
甚至會走入人生的低谷.
忘記了上帝是願意與我們同在.
包底的神.

$^{801}$而在低谷裡面.
神很好.
很真實地為我們預備了教會群體.
身邊的弟兄姊妹.
成為我們很寶貴的禮物.
去幫助我們同行.
分享一個我在上帝裡面.
很經歷到教會群體和我同行的經歷.
在我讀神學之前.
我是做中學老師.
我第一年教書.
就經歷了人生一個很困擾的時間.
我是教體育的.
出來工作第一年.
我帶著很熱血.
每一次都很振奮.
很多的預備.
每一次上課都很大聲.
街坊和家長都說.
你學校來了.
老師很吵.
但每一次同學都很享受上我的課.
其中只有一班同學.
我是很不願意去面對.
有時人就是這樣.
可能一百個人喜歡你.
但一個人對你有一些不好的說話.
你都會很受影響.
那時候有一班同學.
我每次都帶著很多焦慮和不平安去到他們當中.
特別是因為他們是.
我男女一起教的那班.
很多的女同學.
跟女同學上PE課對於她來說是很大壓力的.
男同學我就可以鐵拳教育.
軍訓一些都可以.
但對著女同學.
我真的毫無方法.
首先他們本身未必很喜歡做運動.
另外他們有很多的手段.

$^{841}$都不是手段.
很多方式我都操作不了.
譬如他們經常說自己不舒服.
又或者我很想他們準時換衣服.
他們就在女更衣室.
我沒有辦法進去再提醒他們.
很多這些焦慮.
我會覺得自己管理課堂不好.
覺得自己很失敗.
自己做不到一個很好的老師.
甚至乎之後影響到我每一次上這一課的時候.
我都很大的壓力.
甚至乎會用一些小動作.
可能特意遲一點上課.
希望可以減低上課的時間逃避去面對這件事.
又或者我不知道是不是.
我可能也有一些生理反應.
在上這一課之前那一兩天.
很容易作病.
真的試過上他們那一課的時候請了病假.
很多的困擾.
情況呢.
整整半年都是這樣.
直到有一次我在教會的導師會裡面.
去和弟兄姊妹分享.
本身那個會其實是在討論很重要的事情.
但是當我分享這件事的時候.
傳道人和一些導師都即時停下那些商討.
他們很多個開聲為我祈禱.
他們呢.
我很經歷到他們那份支持和同行.
當中我們很多的眼淚.
在之後我每一次再走進課堂的時候.
我都會想起有他們的祈禱.
成為我的同伴去托住我.
所以我就多了一份力量平安去面對.
慢慢呢.
又真的.
當我自己願意去面對.
去願意克服一些狀態的時候.

$^{881}$那班同學又慢慢衝上我的體育課.
最後我們都建立了一個很好的關係.
甚至乎現在有時候見到.
我們都會去.
還記得大家去打招呼.
我想我看到原來在屬靈群體裡面.
上帝給我們很重要的一份禮物.
是幫助我們同行.
去經歷很多人生艱難的地方.
過了這個星期.
其實我剛剛來到黃浸五個月了.
時間過得很快.
我覺得教會給我很深印象.
我經歷到那份溫情.
那份很多的體諒.
很多的支持和鼓勵.
所以我都很容易去投入.
參與其中.
我相信我們這間教會.
都可以形成很大的力量.
去幫助很多的人去重新站立起來.
弟兄姊妹.
今天藉著這個題目.
上帝告訴我們.
可能在一些人生的低谷困境裡面.
我們可能很容易會走入一個惡性循環.
我們會去比較.
我們會想很多的自己.
我們會很多overthink.
為我們帶來很多的負能量.
情緒.
但是這些情緒.
這些負能量.
並不是壓倒性在我們生命裡面.
神才是我們人生的關鍵.
其實在十九章裡面.
大家可以回去看看.
「起來」這字不斷重複.
神不斷鼓勵我們重新去站立起來.
祂體恤我們的軟弱.

$^{921}$給我們及時的供應.
更加用微小的聲音去引導我們.
給我們身邊.
你們身邊.
很多的同行者.
去跨過這些難關.
願意神幫助我們每一個.
今天是我們門徒進心系列的講道.
最後一個講.
在整個系列裡面.
我相信我們教母同工.
很想我們一起去學習.
做一個情緒健康的門徒.
我們要學習去面對.
去承認自己在人生裡面.
有不同的情緒狀態.
我們可以有疑惑.
恐懼,焦慮,傷心,難過,失意.
但我們千萬不要讓這些情緒去定義我們自己.
不要讓它們壓過了他們的生命.
希望我們一起去努力.
之前有一個調查報告.
很多人信主.
很多名義的基督徒.
但真正很具體.
很真實.
作主門徒的人有多少呢?.
願意我們一起去努力.
告別浮淺.
做一個進心.
真實跟隨主的門徒.
我們一起祈禱.
藉著今天的訊息.
你提醒我們.
你愛我們每一個.
就算我們在一些不同的困境.
低谷裡面.
你一次,兩次.
不斷地去鼓勵,提醒我們.
靠著你重新站立起來.

$^{961}$願意幫助我們每一個.
可能我們生命裡有不同的狀況.
願你的愛去填滿我們每一個.
更加讓我們透過你的說話.
在我們生命站立起來.
做一個真實,進心的門徒.
讓更加多人去認識你.
為我們教會每一個弟兄姊妹的生命祈禱.
願意我們都成為這個世代的光.
讓我們身邊的人.
可能有不同的陰霾也好.
我們幫助他們重新站立起來.
我們這樣祈禱.
奉主耶穌基督名字而求.
阿們.
\newpage

\allsectionsfont{\centering}

\setlength\parindent{0pt}
\setlength{\columnsep}{1.25em}
\setlength{\parfillskip}{0pt}
\setlength{\tabcolsep}{1em}
\raggedbottom

\pagenumbering{gobble}


\newfontfamily\leftfont[Path=../fonts/fell_french_canon/, Ligatures=TeX, ItalicFont=IMFeFCit29C.otf, BoldFont=AveriaLibre-Bold.ttf]{IMFeFCrm29C.otf}
\newfontfamily\leftcitationfont[Path=../fonts/frankruehl/]{FrankRuehlCLM-Medium.ttf}
\newfontfamily\centerfont[Path=../fonts/garamond/, Ligatures=TeX, ItalicFont=EBGaramond-SemiBoldItalic.ttf]{EBGaramond-SemiBold.ttf}
\newfontfamily\rightfont[Path=../fonts/averia/, Ligatures=TeX, ItalicFont=AveriaLibre-RegularItalic.ttf, BoldFont=AveriaLibre-Bold.ttf, BoldItalicFont=AveriaLibre-BoldItalic.ttf]{AveriaLibre-Light.ttf}
\newfontfamily\rightcitationfont[Path=../fonts/rashi/]{Mekorot-Rashi.ttf}
\definecolor{hcolor}{HTML}{D3230C}
\definecolor{rcolor}{HTML}{D36F0C}
\newcommand{\chfont}[1]{\centerfont{\huge\textcolor{hcolor}{#1}}}
\newcommand{\leftcitation}[1]{\leftcitationfont{\Large\textcolor{hcolor}{#1}}}
\newcommand{\rightcitation}[1]{\rightcitationfont{\normalsize\textcolor{rcolor}{#1}}}
\newfontfamily\flowerfont[Path=../fonts/fell_flowers/]{IMFeFlow2.otf}

\begin{sloppypar}

\chapter*{\chfont{編按結語}}

\columnratio{0.5,0.5}\begin{paracol}{2}

\fontsize{11}{13}\leftfont \Large \leftcitation{א} \leftfont 余少好文．宏志博覽群書而不忘．善存藏經籍文獻備後時之用。\leftcitation{ב} \leftfont 歸主年時．受友所薦．聞道網海．\switchcolumn\fontsize{11}{13}\rightfont \Large \leftcitation{ח} \rightfont 有見粵道之危．國之封講道千言亦將就至．急之何則為？\leftcitation{ט} \rightfont 嘗聞猶太者之傳承．在其力守口述之

\end{paracol}


\columnratio{0.32,0.32,0.32}\begin{paracol}{3}

\fontsize{11}{13}\leftfont \Large 尤以吳約翰遜者 \switchcolumn[2]\fontsize{11}{13}\rightfont \Large 統．以煉千載不

\end{paracol}

\columnratio{0.32,0.32,0.32}
\begin{paracol}{3}\fontsize{11}{13}\leftfont \Large 為重．其載上之粵語講道緩緩入耳．收之藏其音頻．善妥整存．反復而嚼．受益無窮。\leftcitation{ג} \leftfont 我城我國既限．歷一四一九之不測．肺疫延年．信徒靈長屢受圍創．神州燈臺數盡指日可待．粵道之求與日俱增。\leftcitation{ד} \leftfont 觀乎社、經、法、媒、言、信、網之地．愈趨受鋤．自翔不果．授受壓力．粵道聖言亦愈漸艱難。\leftcitation{ה} \leftfont 況崇基例乎．學苑講道屢逆權勢者．其言末強受壓．舊章盡刪以存其身。講道釋數失傳．徒嘆奈何。

\switchcolumn

\fontsize{11}{13}\centerfont 
\begin{tikzpicture}
    \node (0,0) [xshift=-0.10cm, yshift=-1.0cm, opacity=0.10]{\includegraphics[width=0.30\textwidth]{../ot_frontcover.png}} ;
    \node (0,0) [xshift=+0.20cm, yshift=+2.0cm, opacity=0.10]{\includegraphics[width=0.20\textwidth]{../christ_on_cross.png}} ;
\end{tikzpicture}
\Large 

\leftcitation{ס} \centerfont 詩百又廿七載：
\leftcitation{ע} \centerfont 非耶和華建屋宇.則匠人之經營徒.
\leftcitation{פ} \centerfont 非耶和華衛城邑.則守者之儆醒徒.
\leftcitation{צ} \centerfont 余獻是卷予華人社區．願為福音流通之器．願獻斯微材為祭榮耀上帝．
\leftcitation{ק} \centerfont 阿門

\switchcolumn

\fontsize{11}{13}\rightfont \Large 滅．時越次聖殿期及當今。\leftcitation{י} \rightfont 猶太者力廣納之．筆錄以卷軸．便以傳、閱、頌、攜、守、鎖、抄、譯、釋、編，得書塔木德、密示拿等經傳．家喻戶曉．傳流若芳。\leftcitation{כ} \rightfont 猶太者文以載道．傳其口述．今我輩粵道之傳應當作如是．遂力行粵音識辨之法．載言載道．以盡忠傳粵道以待興。\leftcitation{ל} \rightfont 蒙下賜恩惠．無畏海量字音文書．既馭上帝之道．今廣及粵語講道．重駛編程之技．匯導粵音遂字稿．重塑講道現場．以傚猶太卷軸之舉便以傳流。\leftcitation{מ} \rightfont 是卷乃粵音口述傳之屬．莫通華文白話之語．

\end{paracol}

\columnratio{0.5,0.5}
\begin{paracol}{2}\fontsize{11}{13}\leftfont \Large \leftcitation{ו} \leftfont 斯殺一違儆百逆．既禁壓之．我輩聞風無奈．在所難免。\leftcitation{ז} \leftfont 另有異人例乎．以版權之名．脅網絡頻道之舉．同授礙予粵道之存流。

\switchcolumn

\fontsize{11}{13}\rightfont \Large 惟待後繼來者之傚．以譯釋傳之於神州華文地。\leftcitation{נ} \rightfont 今能排程驅馭圖靈以編彙文檔，其碼長共數千千亦無逢大礙．全蒙上帝保守。

\end{paracol}



\columnratio{1}\begin{paracol}{1}

\fontsize{11}{13}\rightfont \Large
~~~~~~~~~~~~~~~~~~~~~~~~~~~~~~~~~~~~~~~~~~~~~~~~~~~~~~~~~~~~~~~~~~~~~~~~~~~~~~~\leftcitation{ר} \rightfont 二零二三年二月一日

~~~~~~~~~~~~~~~~~~~~~~~~~~~~~~~~~~~~~~~~~~~~~~~~~~~~~~~~~~~~~~~~~~~~~~~~~~~~~~~\leftcitation{ש} \rightfont 米迦勒

~~~~~~~~~~~~~~~~~~~~~~~~~~~~~~~~~~~~~~~~~~~~~~~~~~~~~~~~~~~~~~~~~~~~~~~~~~~~~~~\leftcitation{ת} \rightfont 書於香港

\end{paracol}

\end{sloppypar}
\end{document}
